% Editable manuscript. Figure code is external and loaded with short input calls.
% No empirical-study names or confidential source observations are included.
\documentclass[11pt,letterpaper]{article}
\usepackage[T1]{fontenc}
\usepackage[utf8]{inputenc}
\usepackage{lmodern}
\usepackage[margin=1in]{geometry}
\usepackage{amsmath,amssymb,booktabs,array,longtable,graphicx,pgf}
\usepackage{pdflscape,caption,microtype}
\usepackage[hidelinks]{hyperref}
\usepackage{enumitem}
\setlist{nosep,leftmargin=*}
\setlength{\parindent}{0pt}
\setlength{\parskip}{0.55em}
\setlength{\emergencystretch}{2em}
\widowpenalty=10000
\clubpenalty=10000
\captionsetup{font=small,labelfont=bf}
\hypersetup{pdftitle={Two-Way Fixed Effects and Modern Difference-in-Differences Estimates},pdfauthor={},pdfsubject={Anonymized research draft, 8 October 2026}}
\newcommand{\TWFE}{\mathrm{TWFE}}
\newcommand{\CSA}{\mathrm{CSA}}
\newcommand{\BJS}{\mathrm{BJS}}
\newcommand{\SE}{\operatorname{SE}}
\newcommand{\Cov}{\operatorname{Cov}}
\newsavebox{\writeupfigure}
\begin{document}
% Keep the main text together without changing the author's prose or font size.
\newgeometry{left=1in,right=1in,top=0.75in,bottom=0.75in}
\setlength{\parskip}{0.2em}
\enlargethispage{0.5\baselineskip}
\begin{center}
{\Large\bfseries Two-Way Fixed Effects and Modern
Difference-in-Differences: \\[3pt] An Empirical Comparison}\par
\vspace{0.4em}
{\small Peter Hull, Brown University\footnote{I thank Jiafeng Chen and Jonathan Roth for helpful comments. OpenAI's ChatGPT did all of the heavy lifting. All errors in prompting or review are my own. Replication data and code is available at \url{https://github.com/peterdhull/twfe-replication}.}}\par
\vspace{0.4em}
{\small 8 October 2026}
\end{center}

This note compares two-way fixed effects (TWFE) estimates from top general-interest journals in economics with estimates from the group-time difference-in-differences estimator of Callaway and Sant'Anna (2021; CSA) and the imputation estimator of Borusyak, Jaravel, and Spiess (2024; BJS). We use recovered observation-level data and focus on binary treatments with staggered adoption. 

\paragraph{Sample and implementation.}
We (the author and ChatGPT Astra) screened 5,014 bibliographic records from 2016--2025 in five economics journals: AER, Econometrica, JPE, QJE, and ReStud. Screening was for published TWFE or CSA/BJS estimates with a single binary treatment and staggered adoption; this advanced 53 papers to package and data review; four further papers remained held for author verification. We collected qualifying main-text results from tables and figures, allowing panels, repeated cross sections, and fixed baseline characteristics interacted with time effects; other time-varying controls were excluded. The final sample includes 335 specifications from 22 papers with available replication data.

Each scalar regression contributes one comparison; each event study contributes a post-treatment average over commonly estimable published horizons or bins and an analogous pre-treatment average when available. Baseline TWFE retains the eligible published specification and sample, including always-treated observations, subject to documented exceptions. Modern estimators use their own identifiable treated populations. We also considered an ``aligned'' comparison which refits TWFE on the exact union of observations used by the corresponding modern estimator. We generally rely on published standard errors; covariances between estimates was estimated by a weighted/Bayesian bootstrap. 
% AUTHOR CHECK (left unchanged at your request): The current figures use joint-bootstrap SEs for both axis scaling and paired inference; the sentence about published standard errors differs from the implementation and the appendix below.

Some specifications were excluded due to small effective sample sizes (i.e., high leverage). This leaves 275 distinct specifications from 20 studies: 247 scalar comparisons and 28 post-treatment event-study averages. The event studies also contribute 22 separate pre-treatment averages, yielding 297 summaries. These include 263 CSA comparisons and 288 BJS comparisons. The appendix details the search and screening.

\paragraph{Results.}
Figures~\ref{fig:baseline} and~\ref{fig:aligned} plot each modern estimate against TWFE. Both axes scale by the TWFE standard error; coordinates outside $[-5,5]$ are capped. Filled markers identify a pointwise significant paired difference at the 5\% level; crosses indicate high-leverage estimates. Circles denote static specifications; squares denote event-study averages, with pre-treatment and post-treatment averages entering separately.

In the baseline sample, opposite signs occur in 12.2\% of CSA pairs and 6.9\% of BJS pairs; significant differences occur in 8.4\% and 14.6\%, respectively. Both conditions hold in 1.1\% and 1.0\%. With aligned samples, the corresponding shares are 15.2\% and 5.6\% for opposite signs, 8.0\% and 10.4\% for significant differences, and 0.0\% and 1.0\% for both. These shares weight retained pairs equally, including pre-treatment averages and warning-marked pairs in the denominator. Tests are not adjusted for multiple comparisons. Specifications within a study are dependent and studies contribute unequal numbers of pairs. Both figure families retain the same pairs; one event-study average also changes its horizon set when samples are aligned (see the appendix).


\clearpage
\setlength{\parskip}{0.55em}
\newgeometry{margin=0.5in}
\begin{landscape}
\begin{figure}[p]
\centering
\caption{Baseline comparison: eligible original samples}
\label{fig:baseline}
\begin{lrbox}{\writeupfigure}
%% Creator: Matplotlib, PGF backend
%%
%% To include the figure in your LaTeX document, write
%%   \input{<filename>.pgf}
%%
%% Make sure the required packages are loaded in your preamble
%%   \usepackage{pgf}
%%
%% Also ensure that all the required font packages are loaded; for instance,
%% the lmodern package is sometimes necessary when using math font.
%%   \usepackage{lmodern}
%%
%% Figures using additional raster images can only be included by \input if
%% they are in the same directory as the main LaTeX file. For loading figures
%% from other directories you can use the `import` package
%%   \usepackage{import}
%%
%% and then include the figures with
%%   \import{<path to file>}{<filename>.pgf}
%%
%% Matplotlib used the following preamble
%%   \def\mathdefault#1{#1}
%%   \everymath=\expandafter{\the\everymath\displaystyle}
%%   \IfFileExists{scrextend.sty}{
%%     \usepackage[fontsize=9.000000pt]{scrextend}
%%   }{
%%     \renewcommand{\normalsize}{\fontsize{9.000000}{10.800000}\selectfont}
%%     \normalsize
%%   }
%%   
%%   \makeatletter\@ifpackageloaded{underscore}{}{\usepackage[strings]{underscore}}\makeatother
%%
\begingroup%
\makeatletter%
\begin{pgfpicture}%
\pgfpathrectangle{\pgfpointorigin}{\pgfqpoint{14.000000in}{10.800000in}}%
\pgfusepath{use as bounding box, clip}%
\begin{pgfscope}%
\pgfsetbuttcap%
\pgfsetmiterjoin%
\definecolor{currentfill}{rgb}{1.000000,1.000000,1.000000}%
\pgfsetfillcolor{currentfill}%
\pgfsetlinewidth{0.000000pt}%
\definecolor{currentstroke}{rgb}{1.000000,1.000000,1.000000}%
\pgfsetstrokecolor{currentstroke}%
\pgfsetdash{}{0pt}%
\pgfpathmoveto{\pgfqpoint{0.000000in}{0.000000in}}%
\pgfpathlineto{\pgfqpoint{14.000000in}{0.000000in}}%
\pgfpathlineto{\pgfqpoint{14.000000in}{10.800000in}}%
\pgfpathlineto{\pgfqpoint{0.000000in}{10.800000in}}%
\pgfpathlineto{\pgfqpoint{0.000000in}{0.000000in}}%
\pgfpathclose%
\pgfusepath{fill}%
\end{pgfscope}%
\begin{pgfscope}%
\pgfsetbuttcap%
\pgfsetmiterjoin%
\definecolor{currentfill}{rgb}{1.000000,1.000000,1.000000}%
\pgfsetfillcolor{currentfill}%
\pgfsetlinewidth{0.000000pt}%
\definecolor{currentstroke}{rgb}{0.000000,0.000000,0.000000}%
\pgfsetstrokecolor{currentstroke}%
\pgfsetstrokeopacity{0.000000}%
\pgfsetdash{}{0pt}%
\pgfpathmoveto{\pgfqpoint{1.109909in}{4.050000in}}%
\pgfpathlineto{\pgfqpoint{6.509909in}{4.050000in}}%
\pgfpathlineto{\pgfqpoint{6.509909in}{9.450000in}}%
\pgfpathlineto{\pgfqpoint{1.109909in}{9.450000in}}%
\pgfpathlineto{\pgfqpoint{1.109909in}{4.050000in}}%
\pgfpathclose%
\pgfusepath{fill}%
\end{pgfscope}%
\begin{pgfscope}%
\pgfpathrectangle{\pgfqpoint{1.109909in}{4.050000in}}{\pgfqpoint{5.400000in}{5.400000in}}%
\pgfusepath{clip}%
\pgfsetrectcap%
\pgfsetroundjoin%
\pgfsetlinewidth{0.602250pt}%
\definecolor{currentstroke}{rgb}{0.898039,0.913725,0.929412}%
\pgfsetstrokecolor{currentstroke}%
\pgfsetdash{}{0pt}%
\pgfpathmoveto{\pgfqpoint{1.109909in}{4.050000in}}%
\pgfpathlineto{\pgfqpoint{1.109909in}{9.450000in}}%
\pgfusepath{stroke}%
\end{pgfscope}%
\begin{pgfscope}%
\pgfsetbuttcap%
\pgfsetroundjoin%
\definecolor{currentfill}{rgb}{0.090196,0.168627,0.227451}%
\pgfsetfillcolor{currentfill}%
\pgfsetlinewidth{0.803000pt}%
\definecolor{currentstroke}{rgb}{0.090196,0.168627,0.227451}%
\pgfsetstrokecolor{currentstroke}%
\pgfsetdash{}{0pt}%
\pgfsys@defobject{currentmarker}{\pgfqpoint{0.000000in}{-0.048611in}}{\pgfqpoint{0.000000in}{0.000000in}}{%
\pgfpathmoveto{\pgfqpoint{0.000000in}{0.000000in}}%
\pgfpathlineto{\pgfqpoint{0.000000in}{-0.048611in}}%
\pgfusepath{stroke,fill}%
}%
\begin{pgfscope}%
\pgfsys@transformshift{1.109909in}{4.050000in}%
\pgfsys@useobject{currentmarker}{}%
\end{pgfscope}%
\end{pgfscope}%
\begin{pgfscope}%
\pgfsetbuttcap%
\pgfsetroundjoin%
\definecolor{currentfill}{rgb}{0.090196,0.168627,0.227451}%
\pgfsetfillcolor{currentfill}%
\pgfsetlinewidth{0.000000pt}%
\definecolor{currentstroke}{rgb}{0.090196,0.168627,0.227451}%
\pgfsetstrokecolor{currentstroke}%
\pgfsetdash{}{0pt}%
\pgfpathmoveto{\pgfqpoint{0.968151in}{3.902182in}}%
\pgfpathlineto{\pgfqpoint{1.046393in}{3.902182in}}%
\pgfpathlineto{\pgfqpoint{1.046393in}{3.891811in}}%
\pgfpathlineto{\pgfqpoint{0.968151in}{3.891811in}}%
\pgfpathlineto{\pgfqpoint{0.968151in}{3.902182in}}%
\pgfpathlineto{\pgfqpoint{0.968151in}{3.902182in}}%
\pgfpathclose%
\pgfpathmoveto{\pgfqpoint{1.073151in}{3.948940in}}%
\pgfpathlineto{\pgfqpoint{1.121550in}{3.948940in}}%
\pgfpathlineto{\pgfqpoint{1.121550in}{3.938549in}}%
\pgfpathlineto{\pgfqpoint{1.084440in}{3.938549in}}%
\pgfpathlineto{\pgfqpoint{1.084440in}{3.916225in}}%
\pgfpathquadraticcurveto{\pgfqpoint{1.087116in}{3.917143in}}{\pgfqpoint{1.089792in}{3.917592in}}%
\pgfpathquadraticcurveto{\pgfqpoint{1.092487in}{3.918041in}}{\pgfqpoint{1.095183in}{3.918041in}}%
\pgfpathquadraticcurveto{\pgfqpoint{1.110436in}{3.918041in}}{\pgfqpoint{1.119343in}{3.909682in}}%
\pgfpathquadraticcurveto{\pgfqpoint{1.128268in}{3.901323in}}{\pgfqpoint{1.128268in}{3.887045in}}%
\pgfpathquadraticcurveto{\pgfqpoint{1.128268in}{3.872338in}}{\pgfqpoint{1.119108in}{3.864174in}}%
\pgfpathquadraticcurveto{\pgfqpoint{1.109948in}{3.856030in}}{\pgfqpoint{1.093288in}{3.856030in}}%
\pgfpathquadraticcurveto{\pgfqpoint{1.087546in}{3.856030in}}{\pgfqpoint{1.081589in}{3.857006in}}%
\pgfpathquadraticcurveto{\pgfqpoint{1.075651in}{3.857983in}}{\pgfqpoint{1.069304in}{3.859936in}}%
\pgfpathlineto{\pgfqpoint{1.069304in}{3.872338in}}%
\pgfpathquadraticcurveto{\pgfqpoint{1.074792in}{3.869350in}}{\pgfqpoint{1.080651in}{3.867885in}}%
\pgfpathquadraticcurveto{\pgfqpoint{1.086511in}{3.866420in}}{\pgfqpoint{1.093034in}{3.866420in}}%
\pgfpathquadraticcurveto{\pgfqpoint{1.103600in}{3.866420in}}{\pgfqpoint{1.109753in}{3.871967in}}%
\pgfpathquadraticcurveto{\pgfqpoint{1.115925in}{3.877514in}}{\pgfqpoint{1.115925in}{3.887045in}}%
\pgfpathquadraticcurveto{\pgfqpoint{1.115925in}{3.896557in}}{\pgfqpoint{1.109753in}{3.902104in}}%
\pgfpathquadraticcurveto{\pgfqpoint{1.103600in}{3.907670in}}{\pgfqpoint{1.093034in}{3.907670in}}%
\pgfpathquadraticcurveto{\pgfqpoint{1.088093in}{3.907670in}}{\pgfqpoint{1.083171in}{3.906577in}}%
\pgfpathquadraticcurveto{\pgfqpoint{1.078268in}{3.905483in}}{\pgfqpoint{1.073151in}{3.903159in}}%
\pgfpathlineto{\pgfqpoint{1.073151in}{3.948940in}}%
\pgfpathlineto{\pgfqpoint{1.073151in}{3.948940in}}%
\pgfpathclose%
\pgfpathmoveto{\pgfqpoint{1.152546in}{3.873315in}}%
\pgfpathlineto{\pgfqpoint{1.165436in}{3.873315in}}%
\pgfpathlineto{\pgfqpoint{1.165436in}{3.857807in}}%
\pgfpathlineto{\pgfqpoint{1.152546in}{3.857807in}}%
\pgfpathlineto{\pgfqpoint{1.152546in}{3.873315in}}%
\pgfpathlineto{\pgfqpoint{1.152546in}{3.873315in}}%
\pgfpathclose%
\pgfpathmoveto{\pgfqpoint{1.218640in}{3.940815in}}%
\pgfpathquadraticcurveto{\pgfqpoint{1.209128in}{3.940815in}}{\pgfqpoint{1.204323in}{3.931440in}}%
\pgfpathquadraticcurveto{\pgfqpoint{1.199538in}{3.922084in}}{\pgfqpoint{1.199538in}{3.903276in}}%
\pgfpathquadraticcurveto{\pgfqpoint{1.199538in}{3.884545in}}{\pgfqpoint{1.204323in}{3.875170in}}%
\pgfpathquadraticcurveto{\pgfqpoint{1.209128in}{3.865795in}}{\pgfqpoint{1.218640in}{3.865795in}}%
\pgfpathquadraticcurveto{\pgfqpoint{1.228229in}{3.865795in}}{\pgfqpoint{1.233015in}{3.875170in}}%
\pgfpathquadraticcurveto{\pgfqpoint{1.237819in}{3.884545in}}{\pgfqpoint{1.237819in}{3.903276in}}%
\pgfpathquadraticcurveto{\pgfqpoint{1.237819in}{3.922084in}}{\pgfqpoint{1.233015in}{3.931440in}}%
\pgfpathquadraticcurveto{\pgfqpoint{1.228229in}{3.940815in}}{\pgfqpoint{1.218640in}{3.940815in}}%
\pgfpathlineto{\pgfqpoint{1.218640in}{3.940815in}}%
\pgfpathclose%
\pgfpathmoveto{\pgfqpoint{1.218640in}{3.950581in}}%
\pgfpathquadraticcurveto{\pgfqpoint{1.233972in}{3.950581in}}{\pgfqpoint{1.242058in}{3.938452in}}%
\pgfpathquadraticcurveto{\pgfqpoint{1.250143in}{3.926342in}}{\pgfqpoint{1.250143in}{3.903276in}}%
\pgfpathquadraticcurveto{\pgfqpoint{1.250143in}{3.880268in}}{\pgfqpoint{1.242058in}{3.868139in}}%
\pgfpathquadraticcurveto{\pgfqpoint{1.233972in}{3.856030in}}{\pgfqpoint{1.218640in}{3.856030in}}%
\pgfpathquadraticcurveto{\pgfqpoint{1.203327in}{3.856030in}}{\pgfqpoint{1.195241in}{3.868139in}}%
\pgfpathquadraticcurveto{\pgfqpoint{1.187155in}{3.880268in}}{\pgfqpoint{1.187155in}{3.903276in}}%
\pgfpathquadraticcurveto{\pgfqpoint{1.187155in}{3.926342in}}{\pgfqpoint{1.195241in}{3.938452in}}%
\pgfpathquadraticcurveto{\pgfqpoint{1.203327in}{3.950581in}}{\pgfqpoint{1.218640in}{3.950581in}}%
\pgfpathlineto{\pgfqpoint{1.218640in}{3.950581in}}%
\pgfpathclose%
\pgfusepath{fill}%
\end{pgfscope}%
\begin{pgfscope}%
\pgfpathrectangle{\pgfqpoint{1.109909in}{4.050000in}}{\pgfqpoint{5.400000in}{5.400000in}}%
\pgfusepath{clip}%
\pgfsetrectcap%
\pgfsetroundjoin%
\pgfsetlinewidth{0.602250pt}%
\definecolor{currentstroke}{rgb}{0.898039,0.913725,0.929412}%
\pgfsetstrokecolor{currentstroke}%
\pgfsetdash{}{0pt}%
\pgfpathmoveto{\pgfqpoint{2.459909in}{4.050000in}}%
\pgfpathlineto{\pgfqpoint{2.459909in}{9.450000in}}%
\pgfusepath{stroke}%
\end{pgfscope}%
\begin{pgfscope}%
\pgfsetbuttcap%
\pgfsetroundjoin%
\definecolor{currentfill}{rgb}{0.090196,0.168627,0.227451}%
\pgfsetfillcolor{currentfill}%
\pgfsetlinewidth{0.803000pt}%
\definecolor{currentstroke}{rgb}{0.090196,0.168627,0.227451}%
\pgfsetstrokecolor{currentstroke}%
\pgfsetdash{}{0pt}%
\pgfsys@defobject{currentmarker}{\pgfqpoint{0.000000in}{-0.048611in}}{\pgfqpoint{0.000000in}{0.000000in}}{%
\pgfpathmoveto{\pgfqpoint{0.000000in}{0.000000in}}%
\pgfpathlineto{\pgfqpoint{0.000000in}{-0.048611in}}%
\pgfusepath{stroke,fill}%
}%
\begin{pgfscope}%
\pgfsys@transformshift{2.459909in}{4.050000in}%
\pgfsys@useobject{currentmarker}{}%
\end{pgfscope}%
\end{pgfscope}%
\begin{pgfscope}%
\pgfsetbuttcap%
\pgfsetroundjoin%
\definecolor{currentfill}{rgb}{0.090196,0.168627,0.227451}%
\pgfsetfillcolor{currentfill}%
\pgfsetlinewidth{0.000000pt}%
\definecolor{currentstroke}{rgb}{0.090196,0.168627,0.227451}%
\pgfsetstrokecolor{currentstroke}%
\pgfsetdash{}{0pt}%
\pgfpathmoveto{\pgfqpoint{2.318151in}{3.902182in}}%
\pgfpathlineto{\pgfqpoint{2.396393in}{3.902182in}}%
\pgfpathlineto{\pgfqpoint{2.396393in}{3.891811in}}%
\pgfpathlineto{\pgfqpoint{2.318151in}{3.891811in}}%
\pgfpathlineto{\pgfqpoint{2.318151in}{3.902182in}}%
\pgfpathlineto{\pgfqpoint{2.318151in}{3.902182in}}%
\pgfpathclose%
\pgfpathmoveto{\pgfqpoint{2.433640in}{3.868178in}}%
\pgfpathlineto{\pgfqpoint{2.476667in}{3.868178in}}%
\pgfpathlineto{\pgfqpoint{2.476667in}{3.857807in}}%
\pgfpathlineto{\pgfqpoint{2.418815in}{3.857807in}}%
\pgfpathlineto{\pgfqpoint{2.418815in}{3.868178in}}%
\pgfpathquadraticcurveto{\pgfqpoint{2.425827in}{3.875444in}}{\pgfqpoint{2.437936in}{3.887670in}}%
\pgfpathquadraticcurveto{\pgfqpoint{2.450065in}{3.899916in}}{\pgfqpoint{2.453171in}{3.903471in}}%
\pgfpathquadraticcurveto{\pgfqpoint{2.459089in}{3.910112in}}{\pgfqpoint{2.461433in}{3.914721in}}%
\pgfpathquadraticcurveto{\pgfqpoint{2.463796in}{3.919331in}}{\pgfqpoint{2.463796in}{3.923784in}}%
\pgfpathquadraticcurveto{\pgfqpoint{2.463796in}{3.931049in}}{\pgfqpoint{2.458698in}{3.935620in}}%
\pgfpathquadraticcurveto{\pgfqpoint{2.453600in}{3.940209in}}{\pgfqpoint{2.445417in}{3.940209in}}%
\pgfpathquadraticcurveto{\pgfqpoint{2.439616in}{3.940209in}}{\pgfqpoint{2.433171in}{3.938198in}}%
\pgfpathquadraticcurveto{\pgfqpoint{2.426745in}{3.936186in}}{\pgfqpoint{2.419421in}{3.932084in}}%
\pgfpathlineto{\pgfqpoint{2.419421in}{3.944545in}}%
\pgfpathquadraticcurveto{\pgfqpoint{2.426862in}{3.947534in}}{\pgfqpoint{2.433327in}{3.949057in}}%
\pgfpathquadraticcurveto{\pgfqpoint{2.439811in}{3.950581in}}{\pgfqpoint{2.445183in}{3.950581in}}%
\pgfpathquadraticcurveto{\pgfqpoint{2.459343in}{3.950581in}}{\pgfqpoint{2.467761in}{3.943491in}}%
\pgfpathquadraticcurveto{\pgfqpoint{2.476179in}{3.936420in}}{\pgfqpoint{2.476179in}{3.924584in}}%
\pgfpathquadraticcurveto{\pgfqpoint{2.476179in}{3.918959in}}{\pgfqpoint{2.474069in}{3.913920in}}%
\pgfpathquadraticcurveto{\pgfqpoint{2.471979in}{3.908901in}}{\pgfqpoint{2.466413in}{3.902065in}}%
\pgfpathquadraticcurveto{\pgfqpoint{2.464890in}{3.900288in}}{\pgfqpoint{2.456706in}{3.891831in}}%
\pgfpathquadraticcurveto{\pgfqpoint{2.448542in}{3.883373in}}{\pgfqpoint{2.433640in}{3.868178in}}%
\pgfpathlineto{\pgfqpoint{2.433640in}{3.868178in}}%
\pgfpathclose%
\pgfpathmoveto{\pgfqpoint{2.502546in}{3.873315in}}%
\pgfpathlineto{\pgfqpoint{2.515436in}{3.873315in}}%
\pgfpathlineto{\pgfqpoint{2.515436in}{3.857807in}}%
\pgfpathlineto{\pgfqpoint{2.502546in}{3.857807in}}%
\pgfpathlineto{\pgfqpoint{2.502546in}{3.873315in}}%
\pgfpathlineto{\pgfqpoint{2.502546in}{3.873315in}}%
\pgfpathclose%
\pgfpathmoveto{\pgfqpoint{2.542409in}{3.948940in}}%
\pgfpathlineto{\pgfqpoint{2.590808in}{3.948940in}}%
\pgfpathlineto{\pgfqpoint{2.590808in}{3.938549in}}%
\pgfpathlineto{\pgfqpoint{2.553698in}{3.938549in}}%
\pgfpathlineto{\pgfqpoint{2.553698in}{3.916225in}}%
\pgfpathquadraticcurveto{\pgfqpoint{2.556374in}{3.917143in}}{\pgfqpoint{2.559050in}{3.917592in}}%
\pgfpathquadraticcurveto{\pgfqpoint{2.561745in}{3.918041in}}{\pgfqpoint{2.564440in}{3.918041in}}%
\pgfpathquadraticcurveto{\pgfqpoint{2.579694in}{3.918041in}}{\pgfqpoint{2.588600in}{3.909682in}}%
\pgfpathquadraticcurveto{\pgfqpoint{2.597526in}{3.901323in}}{\pgfqpoint{2.597526in}{3.887045in}}%
\pgfpathquadraticcurveto{\pgfqpoint{2.597526in}{3.872338in}}{\pgfqpoint{2.588366in}{3.864174in}}%
\pgfpathquadraticcurveto{\pgfqpoint{2.579206in}{3.856030in}}{\pgfqpoint{2.562546in}{3.856030in}}%
\pgfpathquadraticcurveto{\pgfqpoint{2.556804in}{3.856030in}}{\pgfqpoint{2.550847in}{3.857006in}}%
\pgfpathquadraticcurveto{\pgfqpoint{2.544909in}{3.857983in}}{\pgfqpoint{2.538561in}{3.859936in}}%
\pgfpathlineto{\pgfqpoint{2.538561in}{3.872338in}}%
\pgfpathquadraticcurveto{\pgfqpoint{2.544050in}{3.869350in}}{\pgfqpoint{2.549909in}{3.867885in}}%
\pgfpathquadraticcurveto{\pgfqpoint{2.555768in}{3.866420in}}{\pgfqpoint{2.562292in}{3.866420in}}%
\pgfpathquadraticcurveto{\pgfqpoint{2.572858in}{3.866420in}}{\pgfqpoint{2.579011in}{3.871967in}}%
\pgfpathquadraticcurveto{\pgfqpoint{2.585183in}{3.877514in}}{\pgfqpoint{2.585183in}{3.887045in}}%
\pgfpathquadraticcurveto{\pgfqpoint{2.585183in}{3.896557in}}{\pgfqpoint{2.579011in}{3.902104in}}%
\pgfpathquadraticcurveto{\pgfqpoint{2.572858in}{3.907670in}}{\pgfqpoint{2.562292in}{3.907670in}}%
\pgfpathquadraticcurveto{\pgfqpoint{2.557350in}{3.907670in}}{\pgfqpoint{2.552429in}{3.906577in}}%
\pgfpathquadraticcurveto{\pgfqpoint{2.547526in}{3.905483in}}{\pgfqpoint{2.542409in}{3.903159in}}%
\pgfpathlineto{\pgfqpoint{2.542409in}{3.948940in}}%
\pgfpathlineto{\pgfqpoint{2.542409in}{3.948940in}}%
\pgfpathclose%
\pgfusepath{fill}%
\end{pgfscope}%
\begin{pgfscope}%
\pgfpathrectangle{\pgfqpoint{1.109909in}{4.050000in}}{\pgfqpoint{5.400000in}{5.400000in}}%
\pgfusepath{clip}%
\pgfsetrectcap%
\pgfsetroundjoin%
\pgfsetlinewidth{0.602250pt}%
\definecolor{currentstroke}{rgb}{0.898039,0.913725,0.929412}%
\pgfsetstrokecolor{currentstroke}%
\pgfsetdash{}{0pt}%
\pgfpathmoveto{\pgfqpoint{3.809909in}{4.050000in}}%
\pgfpathlineto{\pgfqpoint{3.809909in}{9.450000in}}%
\pgfusepath{stroke}%
\end{pgfscope}%
\begin{pgfscope}%
\pgfsetbuttcap%
\pgfsetroundjoin%
\definecolor{currentfill}{rgb}{0.090196,0.168627,0.227451}%
\pgfsetfillcolor{currentfill}%
\pgfsetlinewidth{0.803000pt}%
\definecolor{currentstroke}{rgb}{0.090196,0.168627,0.227451}%
\pgfsetstrokecolor{currentstroke}%
\pgfsetdash{}{0pt}%
\pgfsys@defobject{currentmarker}{\pgfqpoint{0.000000in}{-0.048611in}}{\pgfqpoint{0.000000in}{0.000000in}}{%
\pgfpathmoveto{\pgfqpoint{0.000000in}{0.000000in}}%
\pgfpathlineto{\pgfqpoint{0.000000in}{-0.048611in}}%
\pgfusepath{stroke,fill}%
}%
\begin{pgfscope}%
\pgfsys@transformshift{3.809909in}{4.050000in}%
\pgfsys@useobject{currentmarker}{}%
\end{pgfscope}%
\end{pgfscope}%
\begin{pgfscope}%
\pgfsetbuttcap%
\pgfsetroundjoin%
\definecolor{currentfill}{rgb}{0.090196,0.168627,0.227451}%
\pgfsetfillcolor{currentfill}%
\pgfsetlinewidth{0.000000pt}%
\definecolor{currentstroke}{rgb}{0.090196,0.168627,0.227451}%
\pgfsetstrokecolor{currentstroke}%
\pgfsetdash{}{0pt}%
\pgfpathmoveto{\pgfqpoint{3.749636in}{3.940815in}}%
\pgfpathquadraticcurveto{\pgfqpoint{3.740124in}{3.940815in}}{\pgfqpoint{3.735319in}{3.931440in}}%
\pgfpathquadraticcurveto{\pgfqpoint{3.730534in}{3.922084in}}{\pgfqpoint{3.730534in}{3.903276in}}%
\pgfpathquadraticcurveto{\pgfqpoint{3.730534in}{3.884545in}}{\pgfqpoint{3.735319in}{3.875170in}}%
\pgfpathquadraticcurveto{\pgfqpoint{3.740124in}{3.865795in}}{\pgfqpoint{3.749636in}{3.865795in}}%
\pgfpathquadraticcurveto{\pgfqpoint{3.759225in}{3.865795in}}{\pgfqpoint{3.764011in}{3.875170in}}%
\pgfpathquadraticcurveto{\pgfqpoint{3.768815in}{3.884545in}}{\pgfqpoint{3.768815in}{3.903276in}}%
\pgfpathquadraticcurveto{\pgfqpoint{3.768815in}{3.922084in}}{\pgfqpoint{3.764011in}{3.931440in}}%
\pgfpathquadraticcurveto{\pgfqpoint{3.759225in}{3.940815in}}{\pgfqpoint{3.749636in}{3.940815in}}%
\pgfpathlineto{\pgfqpoint{3.749636in}{3.940815in}}%
\pgfpathclose%
\pgfpathmoveto{\pgfqpoint{3.749636in}{3.950581in}}%
\pgfpathquadraticcurveto{\pgfqpoint{3.764968in}{3.950581in}}{\pgfqpoint{3.773054in}{3.938452in}}%
\pgfpathquadraticcurveto{\pgfqpoint{3.781140in}{3.926342in}}{\pgfqpoint{3.781140in}{3.903276in}}%
\pgfpathquadraticcurveto{\pgfqpoint{3.781140in}{3.880268in}}{\pgfqpoint{3.773054in}{3.868139in}}%
\pgfpathquadraticcurveto{\pgfqpoint{3.764968in}{3.856030in}}{\pgfqpoint{3.749636in}{3.856030in}}%
\pgfpathquadraticcurveto{\pgfqpoint{3.734323in}{3.856030in}}{\pgfqpoint{3.726237in}{3.868139in}}%
\pgfpathquadraticcurveto{\pgfqpoint{3.718151in}{3.880268in}}{\pgfqpoint{3.718151in}{3.903276in}}%
\pgfpathquadraticcurveto{\pgfqpoint{3.718151in}{3.926342in}}{\pgfqpoint{3.726237in}{3.938452in}}%
\pgfpathquadraticcurveto{\pgfqpoint{3.734323in}{3.950581in}}{\pgfqpoint{3.749636in}{3.950581in}}%
\pgfpathlineto{\pgfqpoint{3.749636in}{3.950581in}}%
\pgfpathclose%
\pgfpathmoveto{\pgfqpoint{3.802800in}{3.873315in}}%
\pgfpathlineto{\pgfqpoint{3.815690in}{3.873315in}}%
\pgfpathlineto{\pgfqpoint{3.815690in}{3.857807in}}%
\pgfpathlineto{\pgfqpoint{3.802800in}{3.857807in}}%
\pgfpathlineto{\pgfqpoint{3.802800in}{3.873315in}}%
\pgfpathlineto{\pgfqpoint{3.802800in}{3.873315in}}%
\pgfpathclose%
\pgfpathmoveto{\pgfqpoint{3.868893in}{3.940815in}}%
\pgfpathquadraticcurveto{\pgfqpoint{3.859382in}{3.940815in}}{\pgfqpoint{3.854577in}{3.931440in}}%
\pgfpathquadraticcurveto{\pgfqpoint{3.849792in}{3.922084in}}{\pgfqpoint{3.849792in}{3.903276in}}%
\pgfpathquadraticcurveto{\pgfqpoint{3.849792in}{3.884545in}}{\pgfqpoint{3.854577in}{3.875170in}}%
\pgfpathquadraticcurveto{\pgfqpoint{3.859382in}{3.865795in}}{\pgfqpoint{3.868893in}{3.865795in}}%
\pgfpathquadraticcurveto{\pgfqpoint{3.878483in}{3.865795in}}{\pgfqpoint{3.883268in}{3.875170in}}%
\pgfpathquadraticcurveto{\pgfqpoint{3.888073in}{3.884545in}}{\pgfqpoint{3.888073in}{3.903276in}}%
\pgfpathquadraticcurveto{\pgfqpoint{3.888073in}{3.922084in}}{\pgfqpoint{3.883268in}{3.931440in}}%
\pgfpathquadraticcurveto{\pgfqpoint{3.878483in}{3.940815in}}{\pgfqpoint{3.868893in}{3.940815in}}%
\pgfpathlineto{\pgfqpoint{3.868893in}{3.940815in}}%
\pgfpathclose%
\pgfpathmoveto{\pgfqpoint{3.868893in}{3.950581in}}%
\pgfpathquadraticcurveto{\pgfqpoint{3.884225in}{3.950581in}}{\pgfqpoint{3.892311in}{3.938452in}}%
\pgfpathquadraticcurveto{\pgfqpoint{3.900397in}{3.926342in}}{\pgfqpoint{3.900397in}{3.903276in}}%
\pgfpathquadraticcurveto{\pgfqpoint{3.900397in}{3.880268in}}{\pgfqpoint{3.892311in}{3.868139in}}%
\pgfpathquadraticcurveto{\pgfqpoint{3.884225in}{3.856030in}}{\pgfqpoint{3.868893in}{3.856030in}}%
\pgfpathquadraticcurveto{\pgfqpoint{3.853581in}{3.856030in}}{\pgfqpoint{3.845495in}{3.868139in}}%
\pgfpathquadraticcurveto{\pgfqpoint{3.837409in}{3.880268in}}{\pgfqpoint{3.837409in}{3.903276in}}%
\pgfpathquadraticcurveto{\pgfqpoint{3.837409in}{3.926342in}}{\pgfqpoint{3.845495in}{3.938452in}}%
\pgfpathquadraticcurveto{\pgfqpoint{3.853581in}{3.950581in}}{\pgfqpoint{3.868893in}{3.950581in}}%
\pgfpathlineto{\pgfqpoint{3.868893in}{3.950581in}}%
\pgfpathclose%
\pgfusepath{fill}%
\end{pgfscope}%
\begin{pgfscope}%
\pgfpathrectangle{\pgfqpoint{1.109909in}{4.050000in}}{\pgfqpoint{5.400000in}{5.400000in}}%
\pgfusepath{clip}%
\pgfsetrectcap%
\pgfsetroundjoin%
\pgfsetlinewidth{0.602250pt}%
\definecolor{currentstroke}{rgb}{0.898039,0.913725,0.929412}%
\pgfsetstrokecolor{currentstroke}%
\pgfsetdash{}{0pt}%
\pgfpathmoveto{\pgfqpoint{5.159909in}{4.050000in}}%
\pgfpathlineto{\pgfqpoint{5.159909in}{9.450000in}}%
\pgfusepath{stroke}%
\end{pgfscope}%
\begin{pgfscope}%
\pgfsetbuttcap%
\pgfsetroundjoin%
\definecolor{currentfill}{rgb}{0.090196,0.168627,0.227451}%
\pgfsetfillcolor{currentfill}%
\pgfsetlinewidth{0.803000pt}%
\definecolor{currentstroke}{rgb}{0.090196,0.168627,0.227451}%
\pgfsetstrokecolor{currentstroke}%
\pgfsetdash{}{0pt}%
\pgfsys@defobject{currentmarker}{\pgfqpoint{0.000000in}{-0.048611in}}{\pgfqpoint{0.000000in}{0.000000in}}{%
\pgfpathmoveto{\pgfqpoint{0.000000in}{0.000000in}}%
\pgfpathlineto{\pgfqpoint{0.000000in}{-0.048611in}}%
\pgfusepath{stroke,fill}%
}%
\begin{pgfscope}%
\pgfsys@transformshift{5.159909in}{4.050000in}%
\pgfsys@useobject{currentmarker}{}%
\end{pgfscope}%
\end{pgfscope}%
\begin{pgfscope}%
\pgfsetbuttcap%
\pgfsetroundjoin%
\definecolor{currentfill}{rgb}{0.090196,0.168627,0.227451}%
\pgfsetfillcolor{currentfill}%
\pgfsetlinewidth{0.000000pt}%
\definecolor{currentstroke}{rgb}{0.090196,0.168627,0.227451}%
\pgfsetstrokecolor{currentstroke}%
\pgfsetdash{}{0pt}%
\pgfpathmoveto{\pgfqpoint{5.083893in}{3.868178in}}%
\pgfpathlineto{\pgfqpoint{5.126921in}{3.868178in}}%
\pgfpathlineto{\pgfqpoint{5.126921in}{3.857807in}}%
\pgfpathlineto{\pgfqpoint{5.069069in}{3.857807in}}%
\pgfpathlineto{\pgfqpoint{5.069069in}{3.868178in}}%
\pgfpathquadraticcurveto{\pgfqpoint{5.076081in}{3.875444in}}{\pgfqpoint{5.088190in}{3.887670in}}%
\pgfpathquadraticcurveto{\pgfqpoint{5.100319in}{3.899916in}}{\pgfqpoint{5.103425in}{3.903471in}}%
\pgfpathquadraticcurveto{\pgfqpoint{5.109343in}{3.910112in}}{\pgfqpoint{5.111686in}{3.914721in}}%
\pgfpathquadraticcurveto{\pgfqpoint{5.114050in}{3.919331in}}{\pgfqpoint{5.114050in}{3.923784in}}%
\pgfpathquadraticcurveto{\pgfqpoint{5.114050in}{3.931049in}}{\pgfqpoint{5.108952in}{3.935620in}}%
\pgfpathquadraticcurveto{\pgfqpoint{5.103854in}{3.940209in}}{\pgfqpoint{5.095671in}{3.940209in}}%
\pgfpathquadraticcurveto{\pgfqpoint{5.089870in}{3.940209in}}{\pgfqpoint{5.083425in}{3.938198in}}%
\pgfpathquadraticcurveto{\pgfqpoint{5.076999in}{3.936186in}}{\pgfqpoint{5.069675in}{3.932084in}}%
\pgfpathlineto{\pgfqpoint{5.069675in}{3.944545in}}%
\pgfpathquadraticcurveto{\pgfqpoint{5.077116in}{3.947534in}}{\pgfqpoint{5.083581in}{3.949057in}}%
\pgfpathquadraticcurveto{\pgfqpoint{5.090065in}{3.950581in}}{\pgfqpoint{5.095436in}{3.950581in}}%
\pgfpathquadraticcurveto{\pgfqpoint{5.109597in}{3.950581in}}{\pgfqpoint{5.118015in}{3.943491in}}%
\pgfpathquadraticcurveto{\pgfqpoint{5.126433in}{3.936420in}}{\pgfqpoint{5.126433in}{3.924584in}}%
\pgfpathquadraticcurveto{\pgfqpoint{5.126433in}{3.918959in}}{\pgfqpoint{5.124323in}{3.913920in}}%
\pgfpathquadraticcurveto{\pgfqpoint{5.122233in}{3.908901in}}{\pgfqpoint{5.116667in}{3.902065in}}%
\pgfpathquadraticcurveto{\pgfqpoint{5.115143in}{3.900288in}}{\pgfqpoint{5.106960in}{3.891831in}}%
\pgfpathquadraticcurveto{\pgfqpoint{5.098796in}{3.883373in}}{\pgfqpoint{5.083893in}{3.868178in}}%
\pgfpathlineto{\pgfqpoint{5.083893in}{3.868178in}}%
\pgfpathclose%
\pgfpathmoveto{\pgfqpoint{5.152800in}{3.873315in}}%
\pgfpathlineto{\pgfqpoint{5.165690in}{3.873315in}}%
\pgfpathlineto{\pgfqpoint{5.165690in}{3.857807in}}%
\pgfpathlineto{\pgfqpoint{5.152800in}{3.857807in}}%
\pgfpathlineto{\pgfqpoint{5.152800in}{3.873315in}}%
\pgfpathlineto{\pgfqpoint{5.152800in}{3.873315in}}%
\pgfpathclose%
\pgfpathmoveto{\pgfqpoint{5.192663in}{3.948940in}}%
\pgfpathlineto{\pgfqpoint{5.241061in}{3.948940in}}%
\pgfpathlineto{\pgfqpoint{5.241061in}{3.938549in}}%
\pgfpathlineto{\pgfqpoint{5.203952in}{3.938549in}}%
\pgfpathlineto{\pgfqpoint{5.203952in}{3.916225in}}%
\pgfpathquadraticcurveto{\pgfqpoint{5.206628in}{3.917143in}}{\pgfqpoint{5.209304in}{3.917592in}}%
\pgfpathquadraticcurveto{\pgfqpoint{5.211999in}{3.918041in}}{\pgfqpoint{5.214694in}{3.918041in}}%
\pgfpathquadraticcurveto{\pgfqpoint{5.229948in}{3.918041in}}{\pgfqpoint{5.238854in}{3.909682in}}%
\pgfpathquadraticcurveto{\pgfqpoint{5.247780in}{3.901323in}}{\pgfqpoint{5.247780in}{3.887045in}}%
\pgfpathquadraticcurveto{\pgfqpoint{5.247780in}{3.872338in}}{\pgfqpoint{5.238620in}{3.864174in}}%
\pgfpathquadraticcurveto{\pgfqpoint{5.229460in}{3.856030in}}{\pgfqpoint{5.212800in}{3.856030in}}%
\pgfpathquadraticcurveto{\pgfqpoint{5.207058in}{3.856030in}}{\pgfqpoint{5.201100in}{3.857006in}}%
\pgfpathquadraticcurveto{\pgfqpoint{5.195163in}{3.857983in}}{\pgfqpoint{5.188815in}{3.859936in}}%
\pgfpathlineto{\pgfqpoint{5.188815in}{3.872338in}}%
\pgfpathquadraticcurveto{\pgfqpoint{5.194304in}{3.869350in}}{\pgfqpoint{5.200163in}{3.867885in}}%
\pgfpathquadraticcurveto{\pgfqpoint{5.206022in}{3.866420in}}{\pgfqpoint{5.212546in}{3.866420in}}%
\pgfpathquadraticcurveto{\pgfqpoint{5.223112in}{3.866420in}}{\pgfqpoint{5.229265in}{3.871967in}}%
\pgfpathquadraticcurveto{\pgfqpoint{5.235436in}{3.877514in}}{\pgfqpoint{5.235436in}{3.887045in}}%
\pgfpathquadraticcurveto{\pgfqpoint{5.235436in}{3.896557in}}{\pgfqpoint{5.229265in}{3.902104in}}%
\pgfpathquadraticcurveto{\pgfqpoint{5.223112in}{3.907670in}}{\pgfqpoint{5.212546in}{3.907670in}}%
\pgfpathquadraticcurveto{\pgfqpoint{5.207604in}{3.907670in}}{\pgfqpoint{5.202683in}{3.906577in}}%
\pgfpathquadraticcurveto{\pgfqpoint{5.197780in}{3.905483in}}{\pgfqpoint{5.192663in}{3.903159in}}%
\pgfpathlineto{\pgfqpoint{5.192663in}{3.948940in}}%
\pgfpathlineto{\pgfqpoint{5.192663in}{3.948940in}}%
\pgfpathclose%
\pgfusepath{fill}%
\end{pgfscope}%
\begin{pgfscope}%
\pgfpathrectangle{\pgfqpoint{1.109909in}{4.050000in}}{\pgfqpoint{5.400000in}{5.400000in}}%
\pgfusepath{clip}%
\pgfsetrectcap%
\pgfsetroundjoin%
\pgfsetlinewidth{0.602250pt}%
\definecolor{currentstroke}{rgb}{0.898039,0.913725,0.929412}%
\pgfsetstrokecolor{currentstroke}%
\pgfsetdash{}{0pt}%
\pgfpathmoveto{\pgfqpoint{6.509909in}{4.050000in}}%
\pgfpathlineto{\pgfqpoint{6.509909in}{9.450000in}}%
\pgfusepath{stroke}%
\end{pgfscope}%
\begin{pgfscope}%
\pgfsetbuttcap%
\pgfsetroundjoin%
\definecolor{currentfill}{rgb}{0.090196,0.168627,0.227451}%
\pgfsetfillcolor{currentfill}%
\pgfsetlinewidth{0.803000pt}%
\definecolor{currentstroke}{rgb}{0.090196,0.168627,0.227451}%
\pgfsetstrokecolor{currentstroke}%
\pgfsetdash{}{0pt}%
\pgfsys@defobject{currentmarker}{\pgfqpoint{0.000000in}{-0.048611in}}{\pgfqpoint{0.000000in}{0.000000in}}{%
\pgfpathmoveto{\pgfqpoint{0.000000in}{0.000000in}}%
\pgfpathlineto{\pgfqpoint{0.000000in}{-0.048611in}}%
\pgfusepath{stroke,fill}%
}%
\begin{pgfscope}%
\pgfsys@transformshift{6.509909in}{4.050000in}%
\pgfsys@useobject{currentmarker}{}%
\end{pgfscope}%
\end{pgfscope}%
\begin{pgfscope}%
\pgfsetbuttcap%
\pgfsetroundjoin%
\definecolor{currentfill}{rgb}{0.090196,0.168627,0.227451}%
\pgfsetfillcolor{currentfill}%
\pgfsetlinewidth{0.000000pt}%
\definecolor{currentstroke}{rgb}{0.090196,0.168627,0.227451}%
\pgfsetstrokecolor{currentstroke}%
\pgfsetdash{}{0pt}%
\pgfpathmoveto{\pgfqpoint{6.423405in}{3.948940in}}%
\pgfpathlineto{\pgfqpoint{6.471804in}{3.948940in}}%
\pgfpathlineto{\pgfqpoint{6.471804in}{3.938549in}}%
\pgfpathlineto{\pgfqpoint{6.434694in}{3.938549in}}%
\pgfpathlineto{\pgfqpoint{6.434694in}{3.916225in}}%
\pgfpathquadraticcurveto{\pgfqpoint{6.437370in}{3.917143in}}{\pgfqpoint{6.440046in}{3.917592in}}%
\pgfpathquadraticcurveto{\pgfqpoint{6.442741in}{3.918041in}}{\pgfqpoint{6.445436in}{3.918041in}}%
\pgfpathquadraticcurveto{\pgfqpoint{6.460690in}{3.918041in}}{\pgfqpoint{6.469597in}{3.909682in}}%
\pgfpathquadraticcurveto{\pgfqpoint{6.478522in}{3.901323in}}{\pgfqpoint{6.478522in}{3.887045in}}%
\pgfpathquadraticcurveto{\pgfqpoint{6.478522in}{3.872338in}}{\pgfqpoint{6.469362in}{3.864174in}}%
\pgfpathquadraticcurveto{\pgfqpoint{6.460202in}{3.856030in}}{\pgfqpoint{6.443542in}{3.856030in}}%
\pgfpathquadraticcurveto{\pgfqpoint{6.437800in}{3.856030in}}{\pgfqpoint{6.431843in}{3.857006in}}%
\pgfpathquadraticcurveto{\pgfqpoint{6.425905in}{3.857983in}}{\pgfqpoint{6.419558in}{3.859936in}}%
\pgfpathlineto{\pgfqpoint{6.419558in}{3.872338in}}%
\pgfpathquadraticcurveto{\pgfqpoint{6.425046in}{3.869350in}}{\pgfqpoint{6.430905in}{3.867885in}}%
\pgfpathquadraticcurveto{\pgfqpoint{6.436765in}{3.866420in}}{\pgfqpoint{6.443288in}{3.866420in}}%
\pgfpathquadraticcurveto{\pgfqpoint{6.453854in}{3.866420in}}{\pgfqpoint{6.460007in}{3.871967in}}%
\pgfpathquadraticcurveto{\pgfqpoint{6.466179in}{3.877514in}}{\pgfqpoint{6.466179in}{3.887045in}}%
\pgfpathquadraticcurveto{\pgfqpoint{6.466179in}{3.896557in}}{\pgfqpoint{6.460007in}{3.902104in}}%
\pgfpathquadraticcurveto{\pgfqpoint{6.453854in}{3.907670in}}{\pgfqpoint{6.443288in}{3.907670in}}%
\pgfpathquadraticcurveto{\pgfqpoint{6.438347in}{3.907670in}}{\pgfqpoint{6.433425in}{3.906577in}}%
\pgfpathquadraticcurveto{\pgfqpoint{6.428522in}{3.905483in}}{\pgfqpoint{6.423405in}{3.903159in}}%
\pgfpathlineto{\pgfqpoint{6.423405in}{3.948940in}}%
\pgfpathlineto{\pgfqpoint{6.423405in}{3.948940in}}%
\pgfpathclose%
\pgfpathmoveto{\pgfqpoint{6.502800in}{3.873315in}}%
\pgfpathlineto{\pgfqpoint{6.515690in}{3.873315in}}%
\pgfpathlineto{\pgfqpoint{6.515690in}{3.857807in}}%
\pgfpathlineto{\pgfqpoint{6.502800in}{3.857807in}}%
\pgfpathlineto{\pgfqpoint{6.502800in}{3.873315in}}%
\pgfpathlineto{\pgfqpoint{6.502800in}{3.873315in}}%
\pgfpathclose%
\pgfpathmoveto{\pgfqpoint{6.568893in}{3.940815in}}%
\pgfpathquadraticcurveto{\pgfqpoint{6.559382in}{3.940815in}}{\pgfqpoint{6.554577in}{3.931440in}}%
\pgfpathquadraticcurveto{\pgfqpoint{6.549792in}{3.922084in}}{\pgfqpoint{6.549792in}{3.903276in}}%
\pgfpathquadraticcurveto{\pgfqpoint{6.549792in}{3.884545in}}{\pgfqpoint{6.554577in}{3.875170in}}%
\pgfpathquadraticcurveto{\pgfqpoint{6.559382in}{3.865795in}}{\pgfqpoint{6.568893in}{3.865795in}}%
\pgfpathquadraticcurveto{\pgfqpoint{6.578483in}{3.865795in}}{\pgfqpoint{6.583268in}{3.875170in}}%
\pgfpathquadraticcurveto{\pgfqpoint{6.588073in}{3.884545in}}{\pgfqpoint{6.588073in}{3.903276in}}%
\pgfpathquadraticcurveto{\pgfqpoint{6.588073in}{3.922084in}}{\pgfqpoint{6.583268in}{3.931440in}}%
\pgfpathquadraticcurveto{\pgfqpoint{6.578483in}{3.940815in}}{\pgfqpoint{6.568893in}{3.940815in}}%
\pgfpathlineto{\pgfqpoint{6.568893in}{3.940815in}}%
\pgfpathclose%
\pgfpathmoveto{\pgfqpoint{6.568893in}{3.950581in}}%
\pgfpathquadraticcurveto{\pgfqpoint{6.584225in}{3.950581in}}{\pgfqpoint{6.592311in}{3.938452in}}%
\pgfpathquadraticcurveto{\pgfqpoint{6.600397in}{3.926342in}}{\pgfqpoint{6.600397in}{3.903276in}}%
\pgfpathquadraticcurveto{\pgfqpoint{6.600397in}{3.880268in}}{\pgfqpoint{6.592311in}{3.868139in}}%
\pgfpathquadraticcurveto{\pgfqpoint{6.584225in}{3.856030in}}{\pgfqpoint{6.568893in}{3.856030in}}%
\pgfpathquadraticcurveto{\pgfqpoint{6.553581in}{3.856030in}}{\pgfqpoint{6.545495in}{3.868139in}}%
\pgfpathquadraticcurveto{\pgfqpoint{6.537409in}{3.880268in}}{\pgfqpoint{6.537409in}{3.903276in}}%
\pgfpathquadraticcurveto{\pgfqpoint{6.537409in}{3.926342in}}{\pgfqpoint{6.545495in}{3.938452in}}%
\pgfpathquadraticcurveto{\pgfqpoint{6.553581in}{3.950581in}}{\pgfqpoint{6.568893in}{3.950581in}}%
\pgfpathlineto{\pgfqpoint{6.568893in}{3.950581in}}%
\pgfpathclose%
\pgfusepath{fill}%
\end{pgfscope}%
\begin{pgfscope}%
\pgfsetbuttcap%
\pgfsetroundjoin%
\definecolor{currentfill}{rgb}{0.090196,0.168627,0.227451}%
\pgfsetfillcolor{currentfill}%
\pgfsetlinewidth{0.000000pt}%
\definecolor{currentstroke}{rgb}{0.090196,0.168627,0.227451}%
\pgfsetstrokecolor{currentstroke}%
\pgfsetdash{}{0pt}%
\pgfpathmoveto{\pgfqpoint{2.789456in}{3.698087in}}%
\pgfpathlineto{\pgfqpoint{2.883676in}{3.698087in}}%
\pgfpathlineto{\pgfqpoint{2.883676in}{3.685387in}}%
\pgfpathlineto{\pgfqpoint{2.844145in}{3.685387in}}%
\pgfpathlineto{\pgfqpoint{2.844145in}{3.586702in}}%
\pgfpathlineto{\pgfqpoint{2.829011in}{3.586702in}}%
\pgfpathlineto{\pgfqpoint{2.829011in}{3.685387in}}%
\pgfpathlineto{\pgfqpoint{2.789456in}{3.685387in}}%
\pgfpathlineto{\pgfqpoint{2.789456in}{3.698087in}}%
\pgfpathlineto{\pgfqpoint{2.789456in}{3.698087in}}%
\pgfpathclose%
\pgfpathmoveto{\pgfqpoint{2.888308in}{3.698087in}}%
\pgfpathlineto{\pgfqpoint{2.903514in}{3.698087in}}%
\pgfpathlineto{\pgfqpoint{2.926932in}{3.603938in}}%
\pgfpathlineto{\pgfqpoint{2.950278in}{3.698087in}}%
\pgfpathlineto{\pgfqpoint{2.967227in}{3.698087in}}%
\pgfpathlineto{\pgfqpoint{2.990645in}{3.603938in}}%
\pgfpathlineto{\pgfqpoint{3.013991in}{3.698087in}}%
\pgfpathlineto{\pgfqpoint{3.029293in}{3.698087in}}%
\pgfpathlineto{\pgfqpoint{3.001315in}{3.586702in}}%
\pgfpathlineto{\pgfqpoint{2.982361in}{3.586702in}}%
\pgfpathlineto{\pgfqpoint{2.958872in}{3.683382in}}%
\pgfpathlineto{\pgfqpoint{2.935143in}{3.586702in}}%
\pgfpathlineto{\pgfqpoint{2.916189in}{3.586702in}}%
\pgfpathlineto{\pgfqpoint{2.888308in}{3.698087in}}%
\pgfpathlineto{\pgfqpoint{2.888308in}{3.698087in}}%
\pgfpathclose%
\pgfpathmoveto{\pgfqpoint{3.049273in}{3.698087in}}%
\pgfpathlineto{\pgfqpoint{3.113273in}{3.698087in}}%
\pgfpathlineto{\pgfqpoint{3.113273in}{3.685387in}}%
\pgfpathlineto{\pgfqpoint{3.064336in}{3.685387in}}%
\pgfpathlineto{\pgfqpoint{3.064336in}{3.652564in}}%
\pgfpathlineto{\pgfqpoint{3.108499in}{3.652564in}}%
\pgfpathlineto{\pgfqpoint{3.108499in}{3.639888in}}%
\pgfpathlineto{\pgfqpoint{3.064336in}{3.639888in}}%
\pgfpathlineto{\pgfqpoint{3.064336in}{3.586702in}}%
\pgfpathlineto{\pgfqpoint{3.049273in}{3.586702in}}%
\pgfpathlineto{\pgfqpoint{3.049273in}{3.698087in}}%
\pgfpathlineto{\pgfqpoint{3.049273in}{3.698087in}}%
\pgfpathclose%
\pgfpathmoveto{\pgfqpoint{3.137144in}{3.698087in}}%
\pgfpathlineto{\pgfqpoint{3.207565in}{3.698087in}}%
\pgfpathlineto{\pgfqpoint{3.207565in}{3.685387in}}%
\pgfpathlineto{\pgfqpoint{3.152207in}{3.685387in}}%
\pgfpathlineto{\pgfqpoint{3.152207in}{3.652421in}}%
\pgfpathlineto{\pgfqpoint{3.205250in}{3.652421in}}%
\pgfpathlineto{\pgfqpoint{3.205250in}{3.639745in}}%
\pgfpathlineto{\pgfqpoint{3.152207in}{3.639745in}}%
\pgfpathlineto{\pgfqpoint{3.152207in}{3.599378in}}%
\pgfpathlineto{\pgfqpoint{3.208902in}{3.599378in}}%
\pgfpathlineto{\pgfqpoint{3.208902in}{3.586702in}}%
\pgfpathlineto{\pgfqpoint{3.137144in}{3.586702in}}%
\pgfpathlineto{\pgfqpoint{3.137144in}{3.698087in}}%
\pgfpathlineto{\pgfqpoint{3.137144in}{3.698087in}}%
\pgfpathclose%
\pgfpathmoveto{\pgfqpoint{3.353110in}{3.631915in}}%
\pgfpathlineto{\pgfqpoint{3.353110in}{3.625207in}}%
\pgfpathlineto{\pgfqpoint{3.289994in}{3.625207in}}%
\pgfpathquadraticcurveto{\pgfqpoint{3.290901in}{3.611028in}}{\pgfqpoint{3.298540in}{3.603604in}}%
\pgfpathquadraticcurveto{\pgfqpoint{3.306179in}{3.596179in}}{\pgfqpoint{3.319833in}{3.596179in}}%
\pgfpathquadraticcurveto{\pgfqpoint{3.327735in}{3.596179in}}{\pgfqpoint{3.335159in}{3.598113in}}%
\pgfpathquadraticcurveto{\pgfqpoint{3.342583in}{3.600047in}}{\pgfqpoint{3.349911in}{3.603938in}}%
\pgfpathlineto{\pgfqpoint{3.349911in}{3.590952in}}%
\pgfpathquadraticcurveto{\pgfqpoint{3.342511in}{3.587824in}}{\pgfqpoint{3.334753in}{3.586177in}}%
\pgfpathquadraticcurveto{\pgfqpoint{3.326995in}{3.584530in}}{\pgfqpoint{3.319022in}{3.584530in}}%
\pgfpathquadraticcurveto{\pgfqpoint{3.299017in}{3.584530in}}{\pgfqpoint{3.287344in}{3.596156in}}%
\pgfpathquadraticcurveto{\pgfqpoint{3.275671in}{3.607805in}}{\pgfqpoint{3.275671in}{3.627666in}}%
\pgfpathquadraticcurveto{\pgfqpoint{3.275671in}{3.648172in}}{\pgfqpoint{3.286747in}{3.660203in}}%
\pgfpathquadraticcurveto{\pgfqpoint{3.297824in}{3.672258in}}{\pgfqpoint{3.316634in}{3.672258in}}%
\pgfpathquadraticcurveto{\pgfqpoint{3.333488in}{3.672258in}}{\pgfqpoint{3.343299in}{3.661396in}}%
\pgfpathquadraticcurveto{\pgfqpoint{3.353110in}{3.650559in}}{\pgfqpoint{3.353110in}{3.631915in}}%
\pgfpathlineto{\pgfqpoint{3.353110in}{3.631915in}}%
\pgfpathclose%
\pgfpathmoveto{\pgfqpoint{3.339384in}{3.635949in}}%
\pgfpathquadraticcurveto{\pgfqpoint{3.339241in}{3.647193in}}{\pgfqpoint{3.333082in}{3.653901in}}%
\pgfpathquadraticcurveto{\pgfqpoint{3.326923in}{3.660633in}}{\pgfqpoint{3.316778in}{3.660633in}}%
\pgfpathquadraticcurveto{\pgfqpoint{3.305295in}{3.660633in}}{\pgfqpoint{3.298397in}{3.654140in}}%
\pgfpathquadraticcurveto{\pgfqpoint{3.291498in}{3.647646in}}{\pgfqpoint{3.290447in}{3.635854in}}%
\pgfpathlineto{\pgfqpoint{3.339384in}{3.635949in}}%
\pgfpathlineto{\pgfqpoint{3.339384in}{3.635949in}}%
\pgfpathclose%
\pgfpathmoveto{\pgfqpoint{3.428902in}{3.667794in}}%
\pgfpathlineto{\pgfqpoint{3.428902in}{3.654808in}}%
\pgfpathquadraticcurveto{\pgfqpoint{3.423101in}{3.657792in}}{\pgfqpoint{3.416823in}{3.659272in}}%
\pgfpathquadraticcurveto{\pgfqpoint{3.410569in}{3.660776in}}{\pgfqpoint{3.403837in}{3.660776in}}%
\pgfpathquadraticcurveto{\pgfqpoint{3.393620in}{3.660776in}}{\pgfqpoint{3.388512in}{3.657649in}}%
\pgfpathquadraticcurveto{\pgfqpoint{3.383403in}{3.654521in}}{\pgfqpoint{3.383403in}{3.648243in}}%
\pgfpathquadraticcurveto{\pgfqpoint{3.383403in}{3.643469in}}{\pgfqpoint{3.387055in}{3.640748in}}%
\pgfpathquadraticcurveto{\pgfqpoint{3.390708in}{3.638026in}}{\pgfqpoint{3.401760in}{3.635567in}}%
\pgfpathlineto{\pgfqpoint{3.406463in}{3.634517in}}%
\pgfpathquadraticcurveto{\pgfqpoint{3.421072in}{3.631390in}}{\pgfqpoint{3.427231in}{3.625685in}}%
\pgfpathquadraticcurveto{\pgfqpoint{3.433390in}{3.619979in}}{\pgfqpoint{3.433390in}{3.609762in}}%
\pgfpathquadraticcurveto{\pgfqpoint{3.433390in}{3.598113in}}{\pgfqpoint{3.424176in}{3.591310in}}%
\pgfpathquadraticcurveto{\pgfqpoint{3.414961in}{3.584530in}}{\pgfqpoint{3.398848in}{3.584530in}}%
\pgfpathquadraticcurveto{\pgfqpoint{3.392140in}{3.584530in}}{\pgfqpoint{3.384859in}{3.585843in}}%
\pgfpathquadraticcurveto{\pgfqpoint{3.377578in}{3.587156in}}{\pgfqpoint{3.369534in}{3.589758in}}%
\pgfpathlineto{\pgfqpoint{3.369534in}{3.603938in}}%
\pgfpathquadraticcurveto{\pgfqpoint{3.377149in}{3.599975in}}{\pgfqpoint{3.384525in}{3.597994in}}%
\pgfpathquadraticcurveto{\pgfqpoint{3.391901in}{3.596036in}}{\pgfqpoint{3.399158in}{3.596036in}}%
\pgfpathquadraticcurveto{\pgfqpoint{3.408850in}{3.596036in}}{\pgfqpoint{3.414054in}{3.599354in}}%
\pgfpathquadraticcurveto{\pgfqpoint{3.419282in}{3.602673in}}{\pgfqpoint{3.419282in}{3.608712in}}%
\pgfpathquadraticcurveto{\pgfqpoint{3.419282in}{3.614298in}}{\pgfqpoint{3.415510in}{3.617282in}}%
\pgfpathquadraticcurveto{\pgfqpoint{3.411762in}{3.620266in}}{\pgfqpoint{3.398991in}{3.623035in}}%
\pgfpathlineto{\pgfqpoint{3.394217in}{3.624157in}}%
\pgfpathquadraticcurveto{\pgfqpoint{3.381469in}{3.626831in}}{\pgfqpoint{3.375788in}{3.632393in}}%
\pgfpathquadraticcurveto{\pgfqpoint{3.370130in}{3.637955in}}{\pgfqpoint{3.370130in}{3.647646in}}%
\pgfpathquadraticcurveto{\pgfqpoint{3.370130in}{3.659439in}}{\pgfqpoint{3.378485in}{3.665837in}}%
\pgfpathquadraticcurveto{\pgfqpoint{3.386841in}{3.672258in}}{\pgfqpoint{3.402214in}{3.672258in}}%
\pgfpathquadraticcurveto{\pgfqpoint{3.409805in}{3.672258in}}{\pgfqpoint{3.416513in}{3.671136in}}%
\pgfpathquadraticcurveto{\pgfqpoint{3.423245in}{3.670038in}}{\pgfqpoint{3.428902in}{3.667794in}}%
\pgfpathlineto{\pgfqpoint{3.428902in}{3.667794in}}%
\pgfpathclose%
\pgfpathmoveto{\pgfqpoint{3.468815in}{3.693981in}}%
\pgfpathlineto{\pgfqpoint{3.468815in}{3.670253in}}%
\pgfpathlineto{\pgfqpoint{3.497079in}{3.670253in}}%
\pgfpathlineto{\pgfqpoint{3.497079in}{3.659582in}}%
\pgfpathlineto{\pgfqpoint{3.468815in}{3.659582in}}%
\pgfpathlineto{\pgfqpoint{3.468815in}{3.614226in}}%
\pgfpathquadraticcurveto{\pgfqpoint{3.468815in}{3.604009in}}{\pgfqpoint{3.471608in}{3.601097in}}%
\pgfpathquadraticcurveto{\pgfqpoint{3.474401in}{3.598185in}}{\pgfqpoint{3.482995in}{3.598185in}}%
\pgfpathlineto{\pgfqpoint{3.497079in}{3.598185in}}%
\pgfpathlineto{\pgfqpoint{3.497079in}{3.586702in}}%
\pgfpathlineto{\pgfqpoint{3.482995in}{3.586702in}}%
\pgfpathquadraticcurveto{\pgfqpoint{3.467097in}{3.586702in}}{\pgfqpoint{3.461057in}{3.592623in}}%
\pgfpathquadraticcurveto{\pgfqpoint{3.455018in}{3.598567in}}{\pgfqpoint{3.455018in}{3.614226in}}%
\pgfpathlineto{\pgfqpoint{3.455018in}{3.659582in}}%
\pgfpathlineto{\pgfqpoint{3.444944in}{3.659582in}}%
\pgfpathlineto{\pgfqpoint{3.444944in}{3.670253in}}%
\pgfpathlineto{\pgfqpoint{3.455018in}{3.670253in}}%
\pgfpathlineto{\pgfqpoint{3.455018in}{3.693981in}}%
\pgfpathlineto{\pgfqpoint{3.468815in}{3.693981in}}%
\pgfpathlineto{\pgfqpoint{3.468815in}{3.693981in}}%
\pgfpathclose%
\pgfpathmoveto{\pgfqpoint{3.515126in}{3.670253in}}%
\pgfpathlineto{\pgfqpoint{3.528852in}{3.670253in}}%
\pgfpathlineto{\pgfqpoint{3.528852in}{3.586702in}}%
\pgfpathlineto{\pgfqpoint{3.515126in}{3.586702in}}%
\pgfpathlineto{\pgfqpoint{3.515126in}{3.670253in}}%
\pgfpathlineto{\pgfqpoint{3.515126in}{3.670253in}}%
\pgfpathclose%
\pgfpathmoveto{\pgfqpoint{3.515126in}{3.702790in}}%
\pgfpathlineto{\pgfqpoint{3.528852in}{3.702790in}}%
\pgfpathlineto{\pgfqpoint{3.528852in}{3.685387in}}%
\pgfpathlineto{\pgfqpoint{3.515126in}{3.685387in}}%
\pgfpathlineto{\pgfqpoint{3.515126in}{3.702790in}}%
\pgfpathlineto{\pgfqpoint{3.515126in}{3.702790in}}%
\pgfpathclose%
\pgfpathmoveto{\pgfqpoint{3.622620in}{3.654211in}}%
\pgfpathquadraticcurveto{\pgfqpoint{3.627776in}{3.663473in}}{\pgfqpoint{3.634937in}{3.667866in}}%
\pgfpathquadraticcurveto{\pgfqpoint{3.642099in}{3.672258in}}{\pgfqpoint{3.651791in}{3.672258in}}%
\pgfpathquadraticcurveto{\pgfqpoint{3.664848in}{3.672258in}}{\pgfqpoint{3.671938in}{3.663115in}}%
\pgfpathquadraticcurveto{\pgfqpoint{3.679028in}{3.653996in}}{\pgfqpoint{3.679028in}{3.637143in}}%
\pgfpathlineto{\pgfqpoint{3.679028in}{3.586702in}}%
\pgfpathlineto{\pgfqpoint{3.665230in}{3.586702in}}%
\pgfpathlineto{\pgfqpoint{3.665230in}{3.636689in}}%
\pgfpathquadraticcurveto{\pgfqpoint{3.665230in}{3.648697in}}{\pgfqpoint{3.660957in}{3.654498in}}%
\pgfpathquadraticcurveto{\pgfqpoint{3.656708in}{3.660322in}}{\pgfqpoint{3.647995in}{3.660322in}}%
\pgfpathquadraticcurveto{\pgfqpoint{3.637324in}{3.660322in}}{\pgfqpoint{3.631118in}{3.653232in}}%
\pgfpathquadraticcurveto{\pgfqpoint{3.624935in}{3.646166in}}{\pgfqpoint{3.624935in}{3.633920in}}%
\pgfpathlineto{\pgfqpoint{3.624935in}{3.586702in}}%
\pgfpathlineto{\pgfqpoint{3.611137in}{3.586702in}}%
\pgfpathlineto{\pgfqpoint{3.611137in}{3.636689in}}%
\pgfpathquadraticcurveto{\pgfqpoint{3.611137in}{3.648768in}}{\pgfqpoint{3.606888in}{3.654545in}}%
\pgfpathquadraticcurveto{\pgfqpoint{3.602639in}{3.660322in}}{\pgfqpoint{3.593759in}{3.660322in}}%
\pgfpathquadraticcurveto{\pgfqpoint{3.583232in}{3.660322in}}{\pgfqpoint{3.577025in}{3.653209in}}%
\pgfpathquadraticcurveto{\pgfqpoint{3.570842in}{3.646095in}}{\pgfqpoint{3.570842in}{3.633920in}}%
\pgfpathlineto{\pgfqpoint{3.570842in}{3.586702in}}%
\pgfpathlineto{\pgfqpoint{3.557045in}{3.586702in}}%
\pgfpathlineto{\pgfqpoint{3.557045in}{3.670253in}}%
\pgfpathlineto{\pgfqpoint{3.570842in}{3.670253in}}%
\pgfpathlineto{\pgfqpoint{3.570842in}{3.657267in}}%
\pgfpathquadraticcurveto{\pgfqpoint{3.575545in}{3.664953in}}{\pgfqpoint{3.582110in}{3.668606in}}%
\pgfpathquadraticcurveto{\pgfqpoint{3.588674in}{3.672258in}}{\pgfqpoint{3.597698in}{3.672258in}}%
\pgfpathquadraticcurveto{\pgfqpoint{3.606817in}{3.672258in}}{\pgfqpoint{3.613190in}{3.667627in}}%
\pgfpathquadraticcurveto{\pgfqpoint{3.619564in}{3.663020in}}{\pgfqpoint{3.622620in}{3.654211in}}%
\pgfpathlineto{\pgfqpoint{3.622620in}{3.654211in}}%
\pgfpathclose%
\pgfpathmoveto{\pgfqpoint{3.744364in}{3.628692in}}%
\pgfpathquadraticcurveto{\pgfqpoint{3.727726in}{3.628692in}}{\pgfqpoint{3.721304in}{3.624897in}}%
\pgfpathquadraticcurveto{\pgfqpoint{3.714883in}{3.621101in}}{\pgfqpoint{3.714883in}{3.611911in}}%
\pgfpathquadraticcurveto{\pgfqpoint{3.714883in}{3.604606in}}{\pgfqpoint{3.719705in}{3.600309in}}%
\pgfpathquadraticcurveto{\pgfqpoint{3.724527in}{3.596036in}}{\pgfqpoint{3.732787in}{3.596036in}}%
\pgfpathquadraticcurveto{\pgfqpoint{3.744221in}{3.596036in}}{\pgfqpoint{3.751120in}{3.604129in}}%
\pgfpathquadraticcurveto{\pgfqpoint{3.758019in}{3.612221in}}{\pgfqpoint{3.758019in}{3.625637in}}%
\pgfpathlineto{\pgfqpoint{3.758019in}{3.628692in}}%
\pgfpathlineto{\pgfqpoint{3.744364in}{3.628692in}}%
\pgfpathlineto{\pgfqpoint{3.744364in}{3.628692in}}%
\pgfpathclose%
\pgfpathmoveto{\pgfqpoint{3.771745in}{3.634374in}}%
\pgfpathlineto{\pgfqpoint{3.771745in}{3.586702in}}%
\pgfpathlineto{\pgfqpoint{3.758019in}{3.586702in}}%
\pgfpathlineto{\pgfqpoint{3.758019in}{3.599378in}}%
\pgfpathquadraticcurveto{\pgfqpoint{3.753316in}{3.591787in}}{\pgfqpoint{3.746298in}{3.588159in}}%
\pgfpathquadraticcurveto{\pgfqpoint{3.739280in}{3.584530in}}{\pgfqpoint{3.729134in}{3.584530in}}%
\pgfpathquadraticcurveto{\pgfqpoint{3.716315in}{3.584530in}}{\pgfqpoint{3.708724in}{3.591739in}}%
\pgfpathquadraticcurveto{\pgfqpoint{3.701157in}{3.598949in}}{\pgfqpoint{3.701157in}{3.611028in}}%
\pgfpathquadraticcurveto{\pgfqpoint{3.701157in}{3.625112in}}{\pgfqpoint{3.710586in}{3.632273in}}%
\pgfpathquadraticcurveto{\pgfqpoint{3.720039in}{3.639435in}}{\pgfqpoint{3.738755in}{3.639435in}}%
\pgfpathlineto{\pgfqpoint{3.758019in}{3.639435in}}%
\pgfpathlineto{\pgfqpoint{3.758019in}{3.640795in}}%
\pgfpathquadraticcurveto{\pgfqpoint{3.758019in}{3.650272in}}{\pgfqpoint{3.751788in}{3.655452in}}%
\pgfpathquadraticcurveto{\pgfqpoint{3.745558in}{3.660633in}}{\pgfqpoint{3.734291in}{3.660633in}}%
\pgfpathquadraticcurveto{\pgfqpoint{3.727129in}{3.660633in}}{\pgfqpoint{3.720326in}{3.658914in}}%
\pgfpathquadraticcurveto{\pgfqpoint{3.713546in}{3.657195in}}{\pgfqpoint{3.707292in}{3.653758in}}%
\pgfpathlineto{\pgfqpoint{3.707292in}{3.666457in}}%
\pgfpathquadraticcurveto{\pgfqpoint{3.714811in}{3.669370in}}{\pgfqpoint{3.721901in}{3.670802in}}%
\pgfpathquadraticcurveto{\pgfqpoint{3.728991in}{3.672258in}}{\pgfqpoint{3.735699in}{3.672258in}}%
\pgfpathquadraticcurveto{\pgfqpoint{3.753841in}{3.672258in}}{\pgfqpoint{3.762793in}{3.662853in}}%
\pgfpathquadraticcurveto{\pgfqpoint{3.771745in}{3.653471in}}{\pgfqpoint{3.771745in}{3.634374in}}%
\pgfpathlineto{\pgfqpoint{3.771745in}{3.634374in}}%
\pgfpathclose%
\pgfpathmoveto{\pgfqpoint{3.813592in}{3.693981in}}%
\pgfpathlineto{\pgfqpoint{3.813592in}{3.670253in}}%
\pgfpathlineto{\pgfqpoint{3.841856in}{3.670253in}}%
\pgfpathlineto{\pgfqpoint{3.841856in}{3.659582in}}%
\pgfpathlineto{\pgfqpoint{3.813592in}{3.659582in}}%
\pgfpathlineto{\pgfqpoint{3.813592in}{3.614226in}}%
\pgfpathquadraticcurveto{\pgfqpoint{3.813592in}{3.604009in}}{\pgfqpoint{3.816385in}{3.601097in}}%
\pgfpathquadraticcurveto{\pgfqpoint{3.819178in}{3.598185in}}{\pgfqpoint{3.827772in}{3.598185in}}%
\pgfpathlineto{\pgfqpoint{3.841856in}{3.598185in}}%
\pgfpathlineto{\pgfqpoint{3.841856in}{3.586702in}}%
\pgfpathlineto{\pgfqpoint{3.827772in}{3.586702in}}%
\pgfpathquadraticcurveto{\pgfqpoint{3.811873in}{3.586702in}}{\pgfqpoint{3.805834in}{3.592623in}}%
\pgfpathquadraticcurveto{\pgfqpoint{3.799794in}{3.598567in}}{\pgfqpoint{3.799794in}{3.614226in}}%
\pgfpathlineto{\pgfqpoint{3.799794in}{3.659582in}}%
\pgfpathlineto{\pgfqpoint{3.789720in}{3.659582in}}%
\pgfpathlineto{\pgfqpoint{3.789720in}{3.670253in}}%
\pgfpathlineto{\pgfqpoint{3.799794in}{3.670253in}}%
\pgfpathlineto{\pgfqpoint{3.799794in}{3.693981in}}%
\pgfpathlineto{\pgfqpoint{3.813592in}{3.693981in}}%
\pgfpathlineto{\pgfqpoint{3.813592in}{3.693981in}}%
\pgfpathclose%
\pgfpathmoveto{\pgfqpoint{3.931374in}{3.631915in}}%
\pgfpathlineto{\pgfqpoint{3.931374in}{3.625207in}}%
\pgfpathlineto{\pgfqpoint{3.868258in}{3.625207in}}%
\pgfpathquadraticcurveto{\pgfqpoint{3.869165in}{3.611028in}}{\pgfqpoint{3.876804in}{3.603604in}}%
\pgfpathquadraticcurveto{\pgfqpoint{3.884443in}{3.596179in}}{\pgfqpoint{3.898097in}{3.596179in}}%
\pgfpathquadraticcurveto{\pgfqpoint{3.905999in}{3.596179in}}{\pgfqpoint{3.913423in}{3.598113in}}%
\pgfpathquadraticcurveto{\pgfqpoint{3.920847in}{3.600047in}}{\pgfqpoint{3.928175in}{3.603938in}}%
\pgfpathlineto{\pgfqpoint{3.928175in}{3.590952in}}%
\pgfpathquadraticcurveto{\pgfqpoint{3.920775in}{3.587824in}}{\pgfqpoint{3.913017in}{3.586177in}}%
\pgfpathquadraticcurveto{\pgfqpoint{3.905258in}{3.584530in}}{\pgfqpoint{3.897285in}{3.584530in}}%
\pgfpathquadraticcurveto{\pgfqpoint{3.877281in}{3.584530in}}{\pgfqpoint{3.865608in}{3.596156in}}%
\pgfpathquadraticcurveto{\pgfqpoint{3.853935in}{3.607805in}}{\pgfqpoint{3.853935in}{3.627666in}}%
\pgfpathquadraticcurveto{\pgfqpoint{3.853935in}{3.648172in}}{\pgfqpoint{3.865011in}{3.660203in}}%
\pgfpathquadraticcurveto{\pgfqpoint{3.876087in}{3.672258in}}{\pgfqpoint{3.894898in}{3.672258in}}%
\pgfpathquadraticcurveto{\pgfqpoint{3.911752in}{3.672258in}}{\pgfqpoint{3.921563in}{3.661396in}}%
\pgfpathquadraticcurveto{\pgfqpoint{3.931374in}{3.650559in}}{\pgfqpoint{3.931374in}{3.631915in}}%
\pgfpathlineto{\pgfqpoint{3.931374in}{3.631915in}}%
\pgfpathclose%
\pgfpathmoveto{\pgfqpoint{3.917648in}{3.635949in}}%
\pgfpathquadraticcurveto{\pgfqpoint{3.917505in}{3.647193in}}{\pgfqpoint{3.911346in}{3.653901in}}%
\pgfpathquadraticcurveto{\pgfqpoint{3.905187in}{3.660633in}}{\pgfqpoint{3.895041in}{3.660633in}}%
\pgfpathquadraticcurveto{\pgfqpoint{3.883559in}{3.660633in}}{\pgfqpoint{3.876660in}{3.654140in}}%
\pgfpathquadraticcurveto{\pgfqpoint{3.869762in}{3.647646in}}{\pgfqpoint{3.868711in}{3.635854in}}%
\pgfpathlineto{\pgfqpoint{3.917648in}{3.635949in}}%
\pgfpathlineto{\pgfqpoint{3.917648in}{3.635949in}}%
\pgfpathclose%
\pgfpathmoveto{\pgfqpoint{4.035430in}{3.702623in}}%
\pgfpathquadraticcurveto{\pgfqpoint{4.025452in}{3.685483in}}{\pgfqpoint{4.020582in}{3.668677in}}%
\pgfpathquadraticcurveto{\pgfqpoint{4.015736in}{3.651896in}}{\pgfqpoint{4.015736in}{3.634660in}}%
\pgfpathquadraticcurveto{\pgfqpoint{4.015736in}{3.617449in}}{\pgfqpoint{4.020630in}{3.600548in}}%
\pgfpathquadraticcurveto{\pgfqpoint{4.025523in}{3.583647in}}{\pgfqpoint{4.035430in}{3.566555in}}%
\pgfpathlineto{\pgfqpoint{4.023494in}{3.566555in}}%
\pgfpathquadraticcurveto{\pgfqpoint{4.012322in}{3.584100in}}{\pgfqpoint{4.006760in}{3.601025in}}%
\pgfpathquadraticcurveto{\pgfqpoint{4.001198in}{3.617950in}}{\pgfqpoint{4.001198in}{3.634660in}}%
\pgfpathquadraticcurveto{\pgfqpoint{4.001198in}{3.651299in}}{\pgfqpoint{4.006712in}{3.668152in}}%
\pgfpathquadraticcurveto{\pgfqpoint{4.012251in}{3.685029in}}{\pgfqpoint{4.023494in}{3.702623in}}%
\pgfpathlineto{\pgfqpoint{4.035430in}{3.702623in}}%
\pgfpathlineto{\pgfqpoint{4.035430in}{3.702623in}}%
\pgfpathclose%
\pgfpathmoveto{\pgfqpoint{4.047222in}{3.698087in}}%
\pgfpathlineto{\pgfqpoint{4.141443in}{3.698087in}}%
\pgfpathlineto{\pgfqpoint{4.141443in}{3.685387in}}%
\pgfpathlineto{\pgfqpoint{4.101912in}{3.685387in}}%
\pgfpathlineto{\pgfqpoint{4.101912in}{3.586702in}}%
\pgfpathlineto{\pgfqpoint{4.086778in}{3.586702in}}%
\pgfpathlineto{\pgfqpoint{4.086778in}{3.685387in}}%
\pgfpathlineto{\pgfqpoint{4.047222in}{3.685387in}}%
\pgfpathlineto{\pgfqpoint{4.047222in}{3.698087in}}%
\pgfpathlineto{\pgfqpoint{4.047222in}{3.698087in}}%
\pgfpathclose%
\pgfpathmoveto{\pgfqpoint{4.146074in}{3.698087in}}%
\pgfpathlineto{\pgfqpoint{4.161281in}{3.698087in}}%
\pgfpathlineto{\pgfqpoint{4.184699in}{3.603938in}}%
\pgfpathlineto{\pgfqpoint{4.208045in}{3.698087in}}%
\pgfpathlineto{\pgfqpoint{4.224994in}{3.698087in}}%
\pgfpathlineto{\pgfqpoint{4.248412in}{3.603938in}}%
\pgfpathlineto{\pgfqpoint{4.271758in}{3.698087in}}%
\pgfpathlineto{\pgfqpoint{4.287060in}{3.698087in}}%
\pgfpathlineto{\pgfqpoint{4.259082in}{3.586702in}}%
\pgfpathlineto{\pgfqpoint{4.240128in}{3.586702in}}%
\pgfpathlineto{\pgfqpoint{4.216639in}{3.683382in}}%
\pgfpathlineto{\pgfqpoint{4.192910in}{3.586702in}}%
\pgfpathlineto{\pgfqpoint{4.173956in}{3.586702in}}%
\pgfpathlineto{\pgfqpoint{4.146074in}{3.698087in}}%
\pgfpathlineto{\pgfqpoint{4.146074in}{3.698087in}}%
\pgfpathclose%
\pgfpathmoveto{\pgfqpoint{4.307040in}{3.698087in}}%
\pgfpathlineto{\pgfqpoint{4.371040in}{3.698087in}}%
\pgfpathlineto{\pgfqpoint{4.371040in}{3.685387in}}%
\pgfpathlineto{\pgfqpoint{4.322103in}{3.685387in}}%
\pgfpathlineto{\pgfqpoint{4.322103in}{3.652564in}}%
\pgfpathlineto{\pgfqpoint{4.366265in}{3.652564in}}%
\pgfpathlineto{\pgfqpoint{4.366265in}{3.639888in}}%
\pgfpathlineto{\pgfqpoint{4.322103in}{3.639888in}}%
\pgfpathlineto{\pgfqpoint{4.322103in}{3.586702in}}%
\pgfpathlineto{\pgfqpoint{4.307040in}{3.586702in}}%
\pgfpathlineto{\pgfqpoint{4.307040in}{3.698087in}}%
\pgfpathlineto{\pgfqpoint{4.307040in}{3.698087in}}%
\pgfpathclose%
\pgfpathmoveto{\pgfqpoint{4.394911in}{3.698087in}}%
\pgfpathlineto{\pgfqpoint{4.465332in}{3.698087in}}%
\pgfpathlineto{\pgfqpoint{4.465332in}{3.685387in}}%
\pgfpathlineto{\pgfqpoint{4.409974in}{3.685387in}}%
\pgfpathlineto{\pgfqpoint{4.409974in}{3.652421in}}%
\pgfpathlineto{\pgfqpoint{4.463017in}{3.652421in}}%
\pgfpathlineto{\pgfqpoint{4.463017in}{3.639745in}}%
\pgfpathlineto{\pgfqpoint{4.409974in}{3.639745in}}%
\pgfpathlineto{\pgfqpoint{4.409974in}{3.599378in}}%
\pgfpathlineto{\pgfqpoint{4.466669in}{3.599378in}}%
\pgfpathlineto{\pgfqpoint{4.466669in}{3.586702in}}%
\pgfpathlineto{\pgfqpoint{4.394911in}{3.586702in}}%
\pgfpathlineto{\pgfqpoint{4.394911in}{3.698087in}}%
\pgfpathlineto{\pgfqpoint{4.394911in}{3.698087in}}%
\pgfpathclose%
\pgfpathmoveto{\pgfqpoint{4.606771in}{3.694435in}}%
\pgfpathlineto{\pgfqpoint{4.606771in}{3.679730in}}%
\pgfpathquadraticcurveto{\pgfqpoint{4.598201in}{3.683836in}}{\pgfqpoint{4.590586in}{3.685841in}}%
\pgfpathquadraticcurveto{\pgfqpoint{4.582971in}{3.687870in}}{\pgfqpoint{4.575881in}{3.687870in}}%
\pgfpathquadraticcurveto{\pgfqpoint{4.563587in}{3.687870in}}{\pgfqpoint{4.556903in}{3.683096in}}%
\pgfpathquadraticcurveto{\pgfqpoint{4.550219in}{3.678321in}}{\pgfqpoint{4.550219in}{3.669513in}}%
\pgfpathquadraticcurveto{\pgfqpoint{4.550219in}{3.662113in}}{\pgfqpoint{4.554660in}{3.658341in}}%
\pgfpathquadraticcurveto{\pgfqpoint{4.559100in}{3.654593in}}{\pgfqpoint{4.571489in}{3.652278in}}%
\pgfpathlineto{\pgfqpoint{4.580584in}{3.650416in}}%
\pgfpathquadraticcurveto{\pgfqpoint{4.597437in}{3.647193in}}{\pgfqpoint{4.605458in}{3.639100in}}%
\pgfpathquadraticcurveto{\pgfqpoint{4.613479in}{3.631008in}}{\pgfqpoint{4.613479in}{3.617449in}}%
\pgfpathquadraticcurveto{\pgfqpoint{4.613479in}{3.601240in}}{\pgfqpoint{4.602617in}{3.592885in}}%
\pgfpathquadraticcurveto{\pgfqpoint{4.591780in}{3.584530in}}{\pgfqpoint{4.570821in}{3.584530in}}%
\pgfpathquadraticcurveto{\pgfqpoint{4.562919in}{3.584530in}}{\pgfqpoint{4.553991in}{3.586321in}}%
\pgfpathquadraticcurveto{\pgfqpoint{4.545087in}{3.588111in}}{\pgfqpoint{4.535538in}{3.591620in}}%
\pgfpathlineto{\pgfqpoint{4.535538in}{3.607137in}}%
\pgfpathquadraticcurveto{\pgfqpoint{4.544705in}{3.602004in}}{\pgfqpoint{4.553514in}{3.599378in}}%
\pgfpathquadraticcurveto{\pgfqpoint{4.562322in}{3.596776in}}{\pgfqpoint{4.570821in}{3.596776in}}%
\pgfpathquadraticcurveto{\pgfqpoint{4.583711in}{3.596776in}}{\pgfqpoint{4.590729in}{3.601837in}}%
\pgfpathquadraticcurveto{\pgfqpoint{4.597748in}{3.606922in}}{\pgfqpoint{4.597748in}{3.616327in}}%
\pgfpathquadraticcurveto{\pgfqpoint{4.597748in}{3.624515in}}{\pgfqpoint{4.592711in}{3.629146in}}%
\pgfpathquadraticcurveto{\pgfqpoint{4.587674in}{3.633777in}}{\pgfqpoint{4.576192in}{3.636093in}}%
\pgfpathlineto{\pgfqpoint{4.567001in}{3.637883in}}%
\pgfpathquadraticcurveto{\pgfqpoint{4.550148in}{3.641225in}}{\pgfqpoint{4.542604in}{3.648387in}}%
\pgfpathquadraticcurveto{\pgfqpoint{4.535085in}{3.655548in}}{\pgfqpoint{4.535085in}{3.668319in}}%
\pgfpathquadraticcurveto{\pgfqpoint{4.535085in}{3.683096in}}{\pgfqpoint{4.545493in}{3.691594in}}%
\pgfpathquadraticcurveto{\pgfqpoint{4.555901in}{3.700092in}}{\pgfqpoint{4.574163in}{3.700092in}}%
\pgfpathquadraticcurveto{\pgfqpoint{4.582016in}{3.700092in}}{\pgfqpoint{4.590133in}{3.698660in}}%
\pgfpathquadraticcurveto{\pgfqpoint{4.598273in}{3.697252in}}{\pgfqpoint{4.606771in}{3.694435in}}%
\pgfpathlineto{\pgfqpoint{4.606771in}{3.694435in}}%
\pgfpathclose%
\pgfpathmoveto{\pgfqpoint{4.636992in}{3.698087in}}%
\pgfpathlineto{\pgfqpoint{4.707413in}{3.698087in}}%
\pgfpathlineto{\pgfqpoint{4.707413in}{3.685387in}}%
\pgfpathlineto{\pgfqpoint{4.652055in}{3.685387in}}%
\pgfpathlineto{\pgfqpoint{4.652055in}{3.652421in}}%
\pgfpathlineto{\pgfqpoint{4.705098in}{3.652421in}}%
\pgfpathlineto{\pgfqpoint{4.705098in}{3.639745in}}%
\pgfpathlineto{\pgfqpoint{4.652055in}{3.639745in}}%
\pgfpathlineto{\pgfqpoint{4.652055in}{3.599378in}}%
\pgfpathlineto{\pgfqpoint{4.708750in}{3.599378in}}%
\pgfpathlineto{\pgfqpoint{4.708750in}{3.586702in}}%
\pgfpathlineto{\pgfqpoint{4.636992in}{3.586702in}}%
\pgfpathlineto{\pgfqpoint{4.636992in}{3.698087in}}%
\pgfpathlineto{\pgfqpoint{4.636992in}{3.698087in}}%
\pgfpathclose%
\pgfpathmoveto{\pgfqpoint{4.786189in}{3.667794in}}%
\pgfpathlineto{\pgfqpoint{4.786189in}{3.654808in}}%
\pgfpathquadraticcurveto{\pgfqpoint{4.780389in}{3.657792in}}{\pgfqpoint{4.774110in}{3.659272in}}%
\pgfpathquadraticcurveto{\pgfqpoint{4.767856in}{3.660776in}}{\pgfqpoint{4.761124in}{3.660776in}}%
\pgfpathquadraticcurveto{\pgfqpoint{4.750907in}{3.660776in}}{\pgfqpoint{4.745799in}{3.657649in}}%
\pgfpathquadraticcurveto{\pgfqpoint{4.740690in}{3.654521in}}{\pgfqpoint{4.740690in}{3.648243in}}%
\pgfpathquadraticcurveto{\pgfqpoint{4.740690in}{3.643469in}}{\pgfqpoint{4.744343in}{3.640748in}}%
\pgfpathquadraticcurveto{\pgfqpoint{4.747995in}{3.638026in}}{\pgfqpoint{4.759048in}{3.635567in}}%
\pgfpathlineto{\pgfqpoint{4.763750in}{3.634517in}}%
\pgfpathquadraticcurveto{\pgfqpoint{4.778360in}{3.631390in}}{\pgfqpoint{4.784518in}{3.625685in}}%
\pgfpathquadraticcurveto{\pgfqpoint{4.790677in}{3.619979in}}{\pgfqpoint{4.790677in}{3.609762in}}%
\pgfpathquadraticcurveto{\pgfqpoint{4.790677in}{3.598113in}}{\pgfqpoint{4.781463in}{3.591310in}}%
\pgfpathquadraticcurveto{\pgfqpoint{4.772249in}{3.584530in}}{\pgfqpoint{4.756135in}{3.584530in}}%
\pgfpathquadraticcurveto{\pgfqpoint{4.749427in}{3.584530in}}{\pgfqpoint{4.742147in}{3.585843in}}%
\pgfpathquadraticcurveto{\pgfqpoint{4.734866in}{3.587156in}}{\pgfqpoint{4.726821in}{3.589758in}}%
\pgfpathlineto{\pgfqpoint{4.726821in}{3.603938in}}%
\pgfpathquadraticcurveto{\pgfqpoint{4.734436in}{3.599975in}}{\pgfqpoint{4.741812in}{3.597994in}}%
\pgfpathquadraticcurveto{\pgfqpoint{4.749189in}{3.596036in}}{\pgfqpoint{4.756446in}{3.596036in}}%
\pgfpathquadraticcurveto{\pgfqpoint{4.766137in}{3.596036in}}{\pgfqpoint{4.771341in}{3.599354in}}%
\pgfpathquadraticcurveto{\pgfqpoint{4.776569in}{3.602673in}}{\pgfqpoint{4.776569in}{3.608712in}}%
\pgfpathquadraticcurveto{\pgfqpoint{4.776569in}{3.614298in}}{\pgfqpoint{4.772798in}{3.617282in}}%
\pgfpathquadraticcurveto{\pgfqpoint{4.769050in}{3.620266in}}{\pgfqpoint{4.756278in}{3.623035in}}%
\pgfpathlineto{\pgfqpoint{4.751504in}{3.624157in}}%
\pgfpathquadraticcurveto{\pgfqpoint{4.738757in}{3.626831in}}{\pgfqpoint{4.733075in}{3.632393in}}%
\pgfpathquadraticcurveto{\pgfqpoint{4.727418in}{3.637955in}}{\pgfqpoint{4.727418in}{3.647646in}}%
\pgfpathquadraticcurveto{\pgfqpoint{4.727418in}{3.659439in}}{\pgfqpoint{4.735773in}{3.665837in}}%
\pgfpathquadraticcurveto{\pgfqpoint{4.744128in}{3.672258in}}{\pgfqpoint{4.759501in}{3.672258in}}%
\pgfpathquadraticcurveto{\pgfqpoint{4.767092in}{3.672258in}}{\pgfqpoint{4.773800in}{3.671136in}}%
\pgfpathquadraticcurveto{\pgfqpoint{4.780532in}{3.670038in}}{\pgfqpoint{4.786189in}{3.667794in}}%
\pgfpathlineto{\pgfqpoint{4.786189in}{3.667794in}}%
\pgfpathclose%
\pgfpathmoveto{\pgfqpoint{4.810371in}{3.702623in}}%
\pgfpathlineto{\pgfqpoint{4.822307in}{3.702623in}}%
\pgfpathquadraticcurveto{\pgfqpoint{4.833479in}{3.685029in}}{\pgfqpoint{4.839041in}{3.668152in}}%
\pgfpathquadraticcurveto{\pgfqpoint{4.844603in}{3.651299in}}{\pgfqpoint{4.844603in}{3.634660in}}%
\pgfpathquadraticcurveto{\pgfqpoint{4.844603in}{3.617950in}}{\pgfqpoint{4.839041in}{3.601025in}}%
\pgfpathquadraticcurveto{\pgfqpoint{4.833479in}{3.584100in}}{\pgfqpoint{4.822307in}{3.566555in}}%
\pgfpathlineto{\pgfqpoint{4.810371in}{3.566555in}}%
\pgfpathquadraticcurveto{\pgfqpoint{4.820278in}{3.583647in}}{\pgfqpoint{4.825172in}{3.600548in}}%
\pgfpathquadraticcurveto{\pgfqpoint{4.830065in}{3.617449in}}{\pgfqpoint{4.830065in}{3.634660in}}%
\pgfpathquadraticcurveto{\pgfqpoint{4.830065in}{3.651896in}}{\pgfqpoint{4.825172in}{3.668677in}}%
\pgfpathquadraticcurveto{\pgfqpoint{4.820278in}{3.685483in}}{\pgfqpoint{4.810371in}{3.702623in}}%
\pgfpathlineto{\pgfqpoint{4.810371in}{3.702623in}}%
\pgfpathclose%
\pgfusepath{fill}%
\end{pgfscope}%
\begin{pgfscope}%
\pgfpathrectangle{\pgfqpoint{1.109909in}{4.050000in}}{\pgfqpoint{5.400000in}{5.400000in}}%
\pgfusepath{clip}%
\pgfsetrectcap%
\pgfsetroundjoin%
\pgfsetlinewidth{0.602250pt}%
\definecolor{currentstroke}{rgb}{0.898039,0.913725,0.929412}%
\pgfsetstrokecolor{currentstroke}%
\pgfsetdash{}{0pt}%
\pgfpathmoveto{\pgfqpoint{1.109909in}{4.050000in}}%
\pgfpathlineto{\pgfqpoint{6.509909in}{4.050000in}}%
\pgfusepath{stroke}%
\end{pgfscope}%
\begin{pgfscope}%
\pgfsetbuttcap%
\pgfsetroundjoin%
\definecolor{currentfill}{rgb}{0.090196,0.168627,0.227451}%
\pgfsetfillcolor{currentfill}%
\pgfsetlinewidth{0.803000pt}%
\definecolor{currentstroke}{rgb}{0.090196,0.168627,0.227451}%
\pgfsetstrokecolor{currentstroke}%
\pgfsetdash{}{0pt}%
\pgfsys@defobject{currentmarker}{\pgfqpoint{-0.048611in}{0.000000in}}{\pgfqpoint{-0.000000in}{0.000000in}}{%
\pgfpathmoveto{\pgfqpoint{-0.000000in}{0.000000in}}%
\pgfpathlineto{\pgfqpoint{-0.048611in}{0.000000in}}%
\pgfusepath{stroke,fill}%
}%
\begin{pgfscope}%
\pgfsys@transformshift{1.109909in}{4.050000in}%
\pgfsys@useobject{currentmarker}{}%
\end{pgfscope}%
\end{pgfscope}%
\begin{pgfscope}%
\pgfsetbuttcap%
\pgfsetroundjoin%
\definecolor{currentfill}{rgb}{0.090196,0.168627,0.227451}%
\pgfsetfillcolor{currentfill}%
\pgfsetlinewidth{0.000000pt}%
\definecolor{currentstroke}{rgb}{0.090196,0.168627,0.227451}%
\pgfsetstrokecolor{currentstroke}%
\pgfsetdash{}{0pt}%
\pgfpathmoveto{\pgfqpoint{0.715929in}{4.046890in}}%
\pgfpathlineto{\pgfqpoint{0.794171in}{4.046890in}}%
\pgfpathlineto{\pgfqpoint{0.794171in}{4.036519in}}%
\pgfpathlineto{\pgfqpoint{0.715929in}{4.036519in}}%
\pgfpathlineto{\pgfqpoint{0.715929in}{4.046890in}}%
\pgfpathlineto{\pgfqpoint{0.715929in}{4.046890in}}%
\pgfpathclose%
\pgfpathmoveto{\pgfqpoint{0.820929in}{4.093647in}}%
\pgfpathlineto{\pgfqpoint{0.869327in}{4.093647in}}%
\pgfpathlineto{\pgfqpoint{0.869327in}{4.083257in}}%
\pgfpathlineto{\pgfqpoint{0.832218in}{4.083257in}}%
\pgfpathlineto{\pgfqpoint{0.832218in}{4.060933in}}%
\pgfpathquadraticcurveto{\pgfqpoint{0.834894in}{4.061851in}}{\pgfqpoint{0.837570in}{4.062300in}}%
\pgfpathquadraticcurveto{\pgfqpoint{0.840265in}{4.062749in}}{\pgfqpoint{0.842960in}{4.062749in}}%
\pgfpathquadraticcurveto{\pgfqpoint{0.858214in}{4.062749in}}{\pgfqpoint{0.867120in}{4.054390in}}%
\pgfpathquadraticcurveto{\pgfqpoint{0.876046in}{4.046030in}}{\pgfqpoint{0.876046in}{4.031753in}}%
\pgfpathquadraticcurveto{\pgfqpoint{0.876046in}{4.017046in}}{\pgfqpoint{0.866886in}{4.008882in}}%
\pgfpathquadraticcurveto{\pgfqpoint{0.857726in}{4.000737in}}{\pgfqpoint{0.841066in}{4.000737in}}%
\pgfpathquadraticcurveto{\pgfqpoint{0.835324in}{4.000737in}}{\pgfqpoint{0.829367in}{4.001714in}}%
\pgfpathquadraticcurveto{\pgfqpoint{0.823429in}{4.002690in}}{\pgfqpoint{0.817081in}{4.004644in}}%
\pgfpathlineto{\pgfqpoint{0.817081in}{4.017046in}}%
\pgfpathquadraticcurveto{\pgfqpoint{0.822570in}{4.014058in}}{\pgfqpoint{0.828429in}{4.012593in}}%
\pgfpathquadraticcurveto{\pgfqpoint{0.834288in}{4.011128in}}{\pgfqpoint{0.840812in}{4.011128in}}%
\pgfpathquadraticcurveto{\pgfqpoint{0.851378in}{4.011128in}}{\pgfqpoint{0.857531in}{4.016675in}}%
\pgfpathquadraticcurveto{\pgfqpoint{0.863702in}{4.022222in}}{\pgfqpoint{0.863702in}{4.031753in}}%
\pgfpathquadraticcurveto{\pgfqpoint{0.863702in}{4.041265in}}{\pgfqpoint{0.857531in}{4.046812in}}%
\pgfpathquadraticcurveto{\pgfqpoint{0.851378in}{4.052378in}}{\pgfqpoint{0.840812in}{4.052378in}}%
\pgfpathquadraticcurveto{\pgfqpoint{0.835870in}{4.052378in}}{\pgfqpoint{0.830949in}{4.051284in}}%
\pgfpathquadraticcurveto{\pgfqpoint{0.826046in}{4.050190in}}{\pgfqpoint{0.820929in}{4.047866in}}%
\pgfpathlineto{\pgfqpoint{0.820929in}{4.093647in}}%
\pgfpathlineto{\pgfqpoint{0.820929in}{4.093647in}}%
\pgfpathclose%
\pgfpathmoveto{\pgfqpoint{0.900324in}{4.018022in}}%
\pgfpathlineto{\pgfqpoint{0.913214in}{4.018022in}}%
\pgfpathlineto{\pgfqpoint{0.913214in}{4.002515in}}%
\pgfpathlineto{\pgfqpoint{0.900324in}{4.002515in}}%
\pgfpathlineto{\pgfqpoint{0.900324in}{4.018022in}}%
\pgfpathlineto{\pgfqpoint{0.900324in}{4.018022in}}%
\pgfpathclose%
\pgfpathmoveto{\pgfqpoint{0.966417in}{4.085522in}}%
\pgfpathquadraticcurveto{\pgfqpoint{0.956906in}{4.085522in}}{\pgfqpoint{0.952101in}{4.076147in}}%
\pgfpathquadraticcurveto{\pgfqpoint{0.947316in}{4.066792in}}{\pgfqpoint{0.947316in}{4.047983in}}%
\pgfpathquadraticcurveto{\pgfqpoint{0.947316in}{4.029253in}}{\pgfqpoint{0.952101in}{4.019878in}}%
\pgfpathquadraticcurveto{\pgfqpoint{0.956906in}{4.010503in}}{\pgfqpoint{0.966417in}{4.010503in}}%
\pgfpathquadraticcurveto{\pgfqpoint{0.976007in}{4.010503in}}{\pgfqpoint{0.980792in}{4.019878in}}%
\pgfpathquadraticcurveto{\pgfqpoint{0.985597in}{4.029253in}}{\pgfqpoint{0.985597in}{4.047983in}}%
\pgfpathquadraticcurveto{\pgfqpoint{0.985597in}{4.066792in}}{\pgfqpoint{0.980792in}{4.076147in}}%
\pgfpathquadraticcurveto{\pgfqpoint{0.976007in}{4.085522in}}{\pgfqpoint{0.966417in}{4.085522in}}%
\pgfpathlineto{\pgfqpoint{0.966417in}{4.085522in}}%
\pgfpathclose%
\pgfpathmoveto{\pgfqpoint{0.966417in}{4.095288in}}%
\pgfpathquadraticcurveto{\pgfqpoint{0.981749in}{4.095288in}}{\pgfqpoint{0.989835in}{4.083159in}}%
\pgfpathquadraticcurveto{\pgfqpoint{0.997921in}{4.071050in}}{\pgfqpoint{0.997921in}{4.047983in}}%
\pgfpathquadraticcurveto{\pgfqpoint{0.997921in}{4.024976in}}{\pgfqpoint{0.989835in}{4.012847in}}%
\pgfpathquadraticcurveto{\pgfqpoint{0.981749in}{4.000737in}}{\pgfqpoint{0.966417in}{4.000737in}}%
\pgfpathquadraticcurveto{\pgfqpoint{0.951105in}{4.000737in}}{\pgfqpoint{0.943019in}{4.012847in}}%
\pgfpathquadraticcurveto{\pgfqpoint{0.934933in}{4.024976in}}{\pgfqpoint{0.934933in}{4.047983in}}%
\pgfpathquadraticcurveto{\pgfqpoint{0.934933in}{4.071050in}}{\pgfqpoint{0.943019in}{4.083159in}}%
\pgfpathquadraticcurveto{\pgfqpoint{0.951105in}{4.095288in}}{\pgfqpoint{0.966417in}{4.095288in}}%
\pgfpathlineto{\pgfqpoint{0.966417in}{4.095288in}}%
\pgfpathclose%
\pgfusepath{fill}%
\end{pgfscope}%
\begin{pgfscope}%
\pgfpathrectangle{\pgfqpoint{1.109909in}{4.050000in}}{\pgfqpoint{5.400000in}{5.400000in}}%
\pgfusepath{clip}%
\pgfsetrectcap%
\pgfsetroundjoin%
\pgfsetlinewidth{0.602250pt}%
\definecolor{currentstroke}{rgb}{0.898039,0.913725,0.929412}%
\pgfsetstrokecolor{currentstroke}%
\pgfsetdash{}{0pt}%
\pgfpathmoveto{\pgfqpoint{1.109909in}{5.400000in}}%
\pgfpathlineto{\pgfqpoint{6.509909in}{5.400000in}}%
\pgfusepath{stroke}%
\end{pgfscope}%
\begin{pgfscope}%
\pgfsetbuttcap%
\pgfsetroundjoin%
\definecolor{currentfill}{rgb}{0.090196,0.168627,0.227451}%
\pgfsetfillcolor{currentfill}%
\pgfsetlinewidth{0.803000pt}%
\definecolor{currentstroke}{rgb}{0.090196,0.168627,0.227451}%
\pgfsetstrokecolor{currentstroke}%
\pgfsetdash{}{0pt}%
\pgfsys@defobject{currentmarker}{\pgfqpoint{-0.048611in}{0.000000in}}{\pgfqpoint{-0.000000in}{0.000000in}}{%
\pgfpathmoveto{\pgfqpoint{-0.000000in}{0.000000in}}%
\pgfpathlineto{\pgfqpoint{-0.048611in}{0.000000in}}%
\pgfusepath{stroke,fill}%
}%
\begin{pgfscope}%
\pgfsys@transformshift{1.109909in}{5.400000in}%
\pgfsys@useobject{currentmarker}{}%
\end{pgfscope}%
\end{pgfscope}%
\begin{pgfscope}%
\pgfsetbuttcap%
\pgfsetroundjoin%
\definecolor{currentfill}{rgb}{0.090196,0.168627,0.227451}%
\pgfsetfillcolor{currentfill}%
\pgfsetlinewidth{0.000000pt}%
\definecolor{currentstroke}{rgb}{0.090196,0.168627,0.227451}%
\pgfsetstrokecolor{currentstroke}%
\pgfsetdash{}{0pt}%
\pgfpathmoveto{\pgfqpoint{0.715929in}{5.396890in}}%
\pgfpathlineto{\pgfqpoint{0.794171in}{5.396890in}}%
\pgfpathlineto{\pgfqpoint{0.794171in}{5.386519in}}%
\pgfpathlineto{\pgfqpoint{0.715929in}{5.386519in}}%
\pgfpathlineto{\pgfqpoint{0.715929in}{5.396890in}}%
\pgfpathlineto{\pgfqpoint{0.715929in}{5.396890in}}%
\pgfpathclose%
\pgfpathmoveto{\pgfqpoint{0.831417in}{5.362886in}}%
\pgfpathlineto{\pgfqpoint{0.874445in}{5.362886in}}%
\pgfpathlineto{\pgfqpoint{0.874445in}{5.352515in}}%
\pgfpathlineto{\pgfqpoint{0.816593in}{5.352515in}}%
\pgfpathlineto{\pgfqpoint{0.816593in}{5.362886in}}%
\pgfpathquadraticcurveto{\pgfqpoint{0.823605in}{5.370151in}}{\pgfqpoint{0.835714in}{5.382378in}}%
\pgfpathquadraticcurveto{\pgfqpoint{0.847843in}{5.394624in}}{\pgfqpoint{0.850949in}{5.398179in}}%
\pgfpathquadraticcurveto{\pgfqpoint{0.856867in}{5.404819in}}{\pgfqpoint{0.859210in}{5.409429in}}%
\pgfpathquadraticcurveto{\pgfqpoint{0.861574in}{5.414038in}}{\pgfqpoint{0.861574in}{5.418491in}}%
\pgfpathquadraticcurveto{\pgfqpoint{0.861574in}{5.425757in}}{\pgfqpoint{0.856476in}{5.430327in}}%
\pgfpathquadraticcurveto{\pgfqpoint{0.851378in}{5.434917in}}{\pgfqpoint{0.843195in}{5.434917in}}%
\pgfpathquadraticcurveto{\pgfqpoint{0.837394in}{5.434917in}}{\pgfqpoint{0.830949in}{5.432905in}}%
\pgfpathquadraticcurveto{\pgfqpoint{0.824523in}{5.430894in}}{\pgfqpoint{0.817199in}{5.426792in}}%
\pgfpathlineto{\pgfqpoint{0.817199in}{5.439253in}}%
\pgfpathquadraticcurveto{\pgfqpoint{0.824640in}{5.442241in}}{\pgfqpoint{0.831105in}{5.443765in}}%
\pgfpathquadraticcurveto{\pgfqpoint{0.837589in}{5.445288in}}{\pgfqpoint{0.842960in}{5.445288in}}%
\pgfpathquadraticcurveto{\pgfqpoint{0.857120in}{5.445288in}}{\pgfqpoint{0.865538in}{5.438198in}}%
\pgfpathquadraticcurveto{\pgfqpoint{0.873956in}{5.431128in}}{\pgfqpoint{0.873956in}{5.419292in}}%
\pgfpathquadraticcurveto{\pgfqpoint{0.873956in}{5.413667in}}{\pgfqpoint{0.871847in}{5.408628in}}%
\pgfpathquadraticcurveto{\pgfqpoint{0.869757in}{5.403608in}}{\pgfqpoint{0.864191in}{5.396772in}}%
\pgfpathquadraticcurveto{\pgfqpoint{0.862667in}{5.394995in}}{\pgfqpoint{0.854484in}{5.386538in}}%
\pgfpathquadraticcurveto{\pgfqpoint{0.846320in}{5.378081in}}{\pgfqpoint{0.831417in}{5.362886in}}%
\pgfpathlineto{\pgfqpoint{0.831417in}{5.362886in}}%
\pgfpathclose%
\pgfpathmoveto{\pgfqpoint{0.900324in}{5.368022in}}%
\pgfpathlineto{\pgfqpoint{0.913214in}{5.368022in}}%
\pgfpathlineto{\pgfqpoint{0.913214in}{5.352515in}}%
\pgfpathlineto{\pgfqpoint{0.900324in}{5.352515in}}%
\pgfpathlineto{\pgfqpoint{0.900324in}{5.368022in}}%
\pgfpathlineto{\pgfqpoint{0.900324in}{5.368022in}}%
\pgfpathclose%
\pgfpathmoveto{\pgfqpoint{0.940187in}{5.443647in}}%
\pgfpathlineto{\pgfqpoint{0.988585in}{5.443647in}}%
\pgfpathlineto{\pgfqpoint{0.988585in}{5.433257in}}%
\pgfpathlineto{\pgfqpoint{0.951476in}{5.433257in}}%
\pgfpathlineto{\pgfqpoint{0.951476in}{5.410933in}}%
\pgfpathquadraticcurveto{\pgfqpoint{0.954152in}{5.411851in}}{\pgfqpoint{0.956827in}{5.412300in}}%
\pgfpathquadraticcurveto{\pgfqpoint{0.959523in}{5.412749in}}{\pgfqpoint{0.962218in}{5.412749in}}%
\pgfpathquadraticcurveto{\pgfqpoint{0.977472in}{5.412749in}}{\pgfqpoint{0.986378in}{5.404390in}}%
\pgfpathquadraticcurveto{\pgfqpoint{0.995304in}{5.396030in}}{\pgfqpoint{0.995304in}{5.381753in}}%
\pgfpathquadraticcurveto{\pgfqpoint{0.995304in}{5.367046in}}{\pgfqpoint{0.986144in}{5.358882in}}%
\pgfpathquadraticcurveto{\pgfqpoint{0.976984in}{5.350737in}}{\pgfqpoint{0.960324in}{5.350737in}}%
\pgfpathquadraticcurveto{\pgfqpoint{0.954581in}{5.350737in}}{\pgfqpoint{0.948624in}{5.351714in}}%
\pgfpathquadraticcurveto{\pgfqpoint{0.942687in}{5.352690in}}{\pgfqpoint{0.936339in}{5.354644in}}%
\pgfpathlineto{\pgfqpoint{0.936339in}{5.367046in}}%
\pgfpathquadraticcurveto{\pgfqpoint{0.941827in}{5.364058in}}{\pgfqpoint{0.947687in}{5.362593in}}%
\pgfpathquadraticcurveto{\pgfqpoint{0.953546in}{5.361128in}}{\pgfqpoint{0.960070in}{5.361128in}}%
\pgfpathquadraticcurveto{\pgfqpoint{0.970636in}{5.361128in}}{\pgfqpoint{0.976788in}{5.366675in}}%
\pgfpathquadraticcurveto{\pgfqpoint{0.982960in}{5.372222in}}{\pgfqpoint{0.982960in}{5.381753in}}%
\pgfpathquadraticcurveto{\pgfqpoint{0.982960in}{5.391265in}}{\pgfqpoint{0.976788in}{5.396812in}}%
\pgfpathquadraticcurveto{\pgfqpoint{0.970636in}{5.402378in}}{\pgfqpoint{0.960070in}{5.402378in}}%
\pgfpathquadraticcurveto{\pgfqpoint{0.955128in}{5.402378in}}{\pgfqpoint{0.950206in}{5.401284in}}%
\pgfpathquadraticcurveto{\pgfqpoint{0.945304in}{5.400190in}}{\pgfqpoint{0.940187in}{5.397866in}}%
\pgfpathlineto{\pgfqpoint{0.940187in}{5.443647in}}%
\pgfpathlineto{\pgfqpoint{0.940187in}{5.443647in}}%
\pgfpathclose%
\pgfusepath{fill}%
\end{pgfscope}%
\begin{pgfscope}%
\pgfpathrectangle{\pgfqpoint{1.109909in}{4.050000in}}{\pgfqpoint{5.400000in}{5.400000in}}%
\pgfusepath{clip}%
\pgfsetrectcap%
\pgfsetroundjoin%
\pgfsetlinewidth{0.602250pt}%
\definecolor{currentstroke}{rgb}{0.898039,0.913725,0.929412}%
\pgfsetstrokecolor{currentstroke}%
\pgfsetdash{}{0pt}%
\pgfpathmoveto{\pgfqpoint{1.109909in}{6.750000in}}%
\pgfpathlineto{\pgfqpoint{6.509909in}{6.750000in}}%
\pgfusepath{stroke}%
\end{pgfscope}%
\begin{pgfscope}%
\pgfsetbuttcap%
\pgfsetroundjoin%
\definecolor{currentfill}{rgb}{0.090196,0.168627,0.227451}%
\pgfsetfillcolor{currentfill}%
\pgfsetlinewidth{0.803000pt}%
\definecolor{currentstroke}{rgb}{0.090196,0.168627,0.227451}%
\pgfsetstrokecolor{currentstroke}%
\pgfsetdash{}{0pt}%
\pgfsys@defobject{currentmarker}{\pgfqpoint{-0.048611in}{0.000000in}}{\pgfqpoint{-0.000000in}{0.000000in}}{%
\pgfpathmoveto{\pgfqpoint{-0.000000in}{0.000000in}}%
\pgfpathlineto{\pgfqpoint{-0.048611in}{0.000000in}}%
\pgfusepath{stroke,fill}%
}%
\begin{pgfscope}%
\pgfsys@transformshift{1.109909in}{6.750000in}%
\pgfsys@useobject{currentmarker}{}%
\end{pgfscope}%
\end{pgfscope}%
\begin{pgfscope}%
\pgfsetbuttcap%
\pgfsetroundjoin%
\definecolor{currentfill}{rgb}{0.090196,0.168627,0.227451}%
\pgfsetfillcolor{currentfill}%
\pgfsetlinewidth{0.000000pt}%
\definecolor{currentstroke}{rgb}{0.090196,0.168627,0.227451}%
\pgfsetstrokecolor{currentstroke}%
\pgfsetdash{}{0pt}%
\pgfpathmoveto{\pgfqpoint{0.852413in}{6.785522in}}%
\pgfpathquadraticcurveto{\pgfqpoint{0.842902in}{6.785522in}}{\pgfqpoint{0.838097in}{6.776147in}}%
\pgfpathquadraticcurveto{\pgfqpoint{0.833312in}{6.766792in}}{\pgfqpoint{0.833312in}{6.747983in}}%
\pgfpathquadraticcurveto{\pgfqpoint{0.833312in}{6.729253in}}{\pgfqpoint{0.838097in}{6.719878in}}%
\pgfpathquadraticcurveto{\pgfqpoint{0.842902in}{6.710503in}}{\pgfqpoint{0.852413in}{6.710503in}}%
\pgfpathquadraticcurveto{\pgfqpoint{0.862003in}{6.710503in}}{\pgfqpoint{0.866788in}{6.719878in}}%
\pgfpathquadraticcurveto{\pgfqpoint{0.871593in}{6.729253in}}{\pgfqpoint{0.871593in}{6.747983in}}%
\pgfpathquadraticcurveto{\pgfqpoint{0.871593in}{6.766792in}}{\pgfqpoint{0.866788in}{6.776147in}}%
\pgfpathquadraticcurveto{\pgfqpoint{0.862003in}{6.785522in}}{\pgfqpoint{0.852413in}{6.785522in}}%
\pgfpathlineto{\pgfqpoint{0.852413in}{6.785522in}}%
\pgfpathclose%
\pgfpathmoveto{\pgfqpoint{0.852413in}{6.795288in}}%
\pgfpathquadraticcurveto{\pgfqpoint{0.867745in}{6.795288in}}{\pgfqpoint{0.875831in}{6.783159in}}%
\pgfpathquadraticcurveto{\pgfqpoint{0.883917in}{6.771050in}}{\pgfqpoint{0.883917in}{6.747983in}}%
\pgfpathquadraticcurveto{\pgfqpoint{0.883917in}{6.724976in}}{\pgfqpoint{0.875831in}{6.712847in}}%
\pgfpathquadraticcurveto{\pgfqpoint{0.867745in}{6.700737in}}{\pgfqpoint{0.852413in}{6.700737in}}%
\pgfpathquadraticcurveto{\pgfqpoint{0.837101in}{6.700737in}}{\pgfqpoint{0.829015in}{6.712847in}}%
\pgfpathquadraticcurveto{\pgfqpoint{0.820929in}{6.724976in}}{\pgfqpoint{0.820929in}{6.747983in}}%
\pgfpathquadraticcurveto{\pgfqpoint{0.820929in}{6.771050in}}{\pgfqpoint{0.829015in}{6.783159in}}%
\pgfpathquadraticcurveto{\pgfqpoint{0.837101in}{6.795288in}}{\pgfqpoint{0.852413in}{6.795288in}}%
\pgfpathlineto{\pgfqpoint{0.852413in}{6.795288in}}%
\pgfpathclose%
\pgfpathmoveto{\pgfqpoint{0.905577in}{6.718022in}}%
\pgfpathlineto{\pgfqpoint{0.918468in}{6.718022in}}%
\pgfpathlineto{\pgfqpoint{0.918468in}{6.702515in}}%
\pgfpathlineto{\pgfqpoint{0.905577in}{6.702515in}}%
\pgfpathlineto{\pgfqpoint{0.905577in}{6.718022in}}%
\pgfpathlineto{\pgfqpoint{0.905577in}{6.718022in}}%
\pgfpathclose%
\pgfpathmoveto{\pgfqpoint{0.971671in}{6.785522in}}%
\pgfpathquadraticcurveto{\pgfqpoint{0.962160in}{6.785522in}}{\pgfqpoint{0.957355in}{6.776147in}}%
\pgfpathquadraticcurveto{\pgfqpoint{0.952570in}{6.766792in}}{\pgfqpoint{0.952570in}{6.747983in}}%
\pgfpathquadraticcurveto{\pgfqpoint{0.952570in}{6.729253in}}{\pgfqpoint{0.957355in}{6.719878in}}%
\pgfpathquadraticcurveto{\pgfqpoint{0.962160in}{6.710503in}}{\pgfqpoint{0.971671in}{6.710503in}}%
\pgfpathquadraticcurveto{\pgfqpoint{0.981261in}{6.710503in}}{\pgfqpoint{0.986046in}{6.719878in}}%
\pgfpathquadraticcurveto{\pgfqpoint{0.990851in}{6.729253in}}{\pgfqpoint{0.990851in}{6.747983in}}%
\pgfpathquadraticcurveto{\pgfqpoint{0.990851in}{6.766792in}}{\pgfqpoint{0.986046in}{6.776147in}}%
\pgfpathquadraticcurveto{\pgfqpoint{0.981261in}{6.785522in}}{\pgfqpoint{0.971671in}{6.785522in}}%
\pgfpathlineto{\pgfqpoint{0.971671in}{6.785522in}}%
\pgfpathclose%
\pgfpathmoveto{\pgfqpoint{0.971671in}{6.795288in}}%
\pgfpathquadraticcurveto{\pgfqpoint{0.987003in}{6.795288in}}{\pgfqpoint{0.995089in}{6.783159in}}%
\pgfpathquadraticcurveto{\pgfqpoint{1.003175in}{6.771050in}}{\pgfqpoint{1.003175in}{6.747983in}}%
\pgfpathquadraticcurveto{\pgfqpoint{1.003175in}{6.724976in}}{\pgfqpoint{0.995089in}{6.712847in}}%
\pgfpathquadraticcurveto{\pgfqpoint{0.987003in}{6.700737in}}{\pgfqpoint{0.971671in}{6.700737in}}%
\pgfpathquadraticcurveto{\pgfqpoint{0.956359in}{6.700737in}}{\pgfqpoint{0.948273in}{6.712847in}}%
\pgfpathquadraticcurveto{\pgfqpoint{0.940187in}{6.724976in}}{\pgfqpoint{0.940187in}{6.747983in}}%
\pgfpathquadraticcurveto{\pgfqpoint{0.940187in}{6.771050in}}{\pgfqpoint{0.948273in}{6.783159in}}%
\pgfpathquadraticcurveto{\pgfqpoint{0.956359in}{6.795288in}}{\pgfqpoint{0.971671in}{6.795288in}}%
\pgfpathlineto{\pgfqpoint{0.971671in}{6.795288in}}%
\pgfpathclose%
\pgfusepath{fill}%
\end{pgfscope}%
\begin{pgfscope}%
\pgfpathrectangle{\pgfqpoint{1.109909in}{4.050000in}}{\pgfqpoint{5.400000in}{5.400000in}}%
\pgfusepath{clip}%
\pgfsetrectcap%
\pgfsetroundjoin%
\pgfsetlinewidth{0.602250pt}%
\definecolor{currentstroke}{rgb}{0.898039,0.913725,0.929412}%
\pgfsetstrokecolor{currentstroke}%
\pgfsetdash{}{0pt}%
\pgfpathmoveto{\pgfqpoint{1.109909in}{8.100000in}}%
\pgfpathlineto{\pgfqpoint{6.509909in}{8.100000in}}%
\pgfusepath{stroke}%
\end{pgfscope}%
\begin{pgfscope}%
\pgfsetbuttcap%
\pgfsetroundjoin%
\definecolor{currentfill}{rgb}{0.090196,0.168627,0.227451}%
\pgfsetfillcolor{currentfill}%
\pgfsetlinewidth{0.803000pt}%
\definecolor{currentstroke}{rgb}{0.090196,0.168627,0.227451}%
\pgfsetstrokecolor{currentstroke}%
\pgfsetdash{}{0pt}%
\pgfsys@defobject{currentmarker}{\pgfqpoint{-0.048611in}{0.000000in}}{\pgfqpoint{-0.000000in}{0.000000in}}{%
\pgfpathmoveto{\pgfqpoint{-0.000000in}{0.000000in}}%
\pgfpathlineto{\pgfqpoint{-0.048611in}{0.000000in}}%
\pgfusepath{stroke,fill}%
}%
\begin{pgfscope}%
\pgfsys@transformshift{1.109909in}{8.100000in}%
\pgfsys@useobject{currentmarker}{}%
\end{pgfscope}%
\end{pgfscope}%
\begin{pgfscope}%
\pgfsetbuttcap%
\pgfsetroundjoin%
\definecolor{currentfill}{rgb}{0.090196,0.168627,0.227451}%
\pgfsetfillcolor{currentfill}%
\pgfsetlinewidth{0.000000pt}%
\definecolor{currentstroke}{rgb}{0.090196,0.168627,0.227451}%
\pgfsetstrokecolor{currentstroke}%
\pgfsetdash{}{0pt}%
\pgfpathmoveto{\pgfqpoint{0.836671in}{8.062886in}}%
\pgfpathlineto{\pgfqpoint{0.879699in}{8.062886in}}%
\pgfpathlineto{\pgfqpoint{0.879699in}{8.052515in}}%
\pgfpathlineto{\pgfqpoint{0.821847in}{8.052515in}}%
\pgfpathlineto{\pgfqpoint{0.821847in}{8.062886in}}%
\pgfpathquadraticcurveto{\pgfqpoint{0.828859in}{8.070151in}}{\pgfqpoint{0.840968in}{8.082378in}}%
\pgfpathquadraticcurveto{\pgfqpoint{0.853097in}{8.094624in}}{\pgfqpoint{0.856202in}{8.098179in}}%
\pgfpathquadraticcurveto{\pgfqpoint{0.862120in}{8.104819in}}{\pgfqpoint{0.864464in}{8.109429in}}%
\pgfpathquadraticcurveto{\pgfqpoint{0.866827in}{8.114038in}}{\pgfqpoint{0.866827in}{8.118491in}}%
\pgfpathquadraticcurveto{\pgfqpoint{0.866827in}{8.125757in}}{\pgfqpoint{0.861730in}{8.130327in}}%
\pgfpathquadraticcurveto{\pgfqpoint{0.856632in}{8.134917in}}{\pgfqpoint{0.848449in}{8.134917in}}%
\pgfpathquadraticcurveto{\pgfqpoint{0.842648in}{8.134917in}}{\pgfqpoint{0.836202in}{8.132905in}}%
\pgfpathquadraticcurveto{\pgfqpoint{0.829777in}{8.130894in}}{\pgfqpoint{0.822452in}{8.126792in}}%
\pgfpathlineto{\pgfqpoint{0.822452in}{8.139253in}}%
\pgfpathquadraticcurveto{\pgfqpoint{0.829894in}{8.142241in}}{\pgfqpoint{0.836359in}{8.143765in}}%
\pgfpathquadraticcurveto{\pgfqpoint{0.842843in}{8.145288in}}{\pgfqpoint{0.848214in}{8.145288in}}%
\pgfpathquadraticcurveto{\pgfqpoint{0.862374in}{8.145288in}}{\pgfqpoint{0.870792in}{8.138198in}}%
\pgfpathquadraticcurveto{\pgfqpoint{0.879210in}{8.131128in}}{\pgfqpoint{0.879210in}{8.119292in}}%
\pgfpathquadraticcurveto{\pgfqpoint{0.879210in}{8.113667in}}{\pgfqpoint{0.877101in}{8.108628in}}%
\pgfpathquadraticcurveto{\pgfqpoint{0.875011in}{8.103608in}}{\pgfqpoint{0.869445in}{8.096772in}}%
\pgfpathquadraticcurveto{\pgfqpoint{0.867921in}{8.094995in}}{\pgfqpoint{0.859738in}{8.086538in}}%
\pgfpathquadraticcurveto{\pgfqpoint{0.851574in}{8.078081in}}{\pgfqpoint{0.836671in}{8.062886in}}%
\pgfpathlineto{\pgfqpoint{0.836671in}{8.062886in}}%
\pgfpathclose%
\pgfpathmoveto{\pgfqpoint{0.905577in}{8.068022in}}%
\pgfpathlineto{\pgfqpoint{0.918468in}{8.068022in}}%
\pgfpathlineto{\pgfqpoint{0.918468in}{8.052515in}}%
\pgfpathlineto{\pgfqpoint{0.905577in}{8.052515in}}%
\pgfpathlineto{\pgfqpoint{0.905577in}{8.068022in}}%
\pgfpathlineto{\pgfqpoint{0.905577in}{8.068022in}}%
\pgfpathclose%
\pgfpathmoveto{\pgfqpoint{0.945441in}{8.143647in}}%
\pgfpathlineto{\pgfqpoint{0.993839in}{8.143647in}}%
\pgfpathlineto{\pgfqpoint{0.993839in}{8.133257in}}%
\pgfpathlineto{\pgfqpoint{0.956730in}{8.133257in}}%
\pgfpathlineto{\pgfqpoint{0.956730in}{8.110933in}}%
\pgfpathquadraticcurveto{\pgfqpoint{0.959406in}{8.111851in}}{\pgfqpoint{0.962081in}{8.112300in}}%
\pgfpathquadraticcurveto{\pgfqpoint{0.964777in}{8.112749in}}{\pgfqpoint{0.967472in}{8.112749in}}%
\pgfpathquadraticcurveto{\pgfqpoint{0.982726in}{8.112749in}}{\pgfqpoint{0.991632in}{8.104390in}}%
\pgfpathquadraticcurveto{\pgfqpoint{1.000558in}{8.096030in}}{\pgfqpoint{1.000558in}{8.081753in}}%
\pgfpathquadraticcurveto{\pgfqpoint{1.000558in}{8.067046in}}{\pgfqpoint{0.991398in}{8.058882in}}%
\pgfpathquadraticcurveto{\pgfqpoint{0.982238in}{8.050737in}}{\pgfqpoint{0.965577in}{8.050737in}}%
\pgfpathquadraticcurveto{\pgfqpoint{0.959835in}{8.050737in}}{\pgfqpoint{0.953878in}{8.051714in}}%
\pgfpathquadraticcurveto{\pgfqpoint{0.947941in}{8.052690in}}{\pgfqpoint{0.941593in}{8.054644in}}%
\pgfpathlineto{\pgfqpoint{0.941593in}{8.067046in}}%
\pgfpathquadraticcurveto{\pgfqpoint{0.947081in}{8.064058in}}{\pgfqpoint{0.952941in}{8.062593in}}%
\pgfpathquadraticcurveto{\pgfqpoint{0.958800in}{8.061128in}}{\pgfqpoint{0.965324in}{8.061128in}}%
\pgfpathquadraticcurveto{\pgfqpoint{0.975890in}{8.061128in}}{\pgfqpoint{0.982042in}{8.066675in}}%
\pgfpathquadraticcurveto{\pgfqpoint{0.988214in}{8.072222in}}{\pgfqpoint{0.988214in}{8.081753in}}%
\pgfpathquadraticcurveto{\pgfqpoint{0.988214in}{8.091265in}}{\pgfqpoint{0.982042in}{8.096812in}}%
\pgfpathquadraticcurveto{\pgfqpoint{0.975890in}{8.102378in}}{\pgfqpoint{0.965324in}{8.102378in}}%
\pgfpathquadraticcurveto{\pgfqpoint{0.960382in}{8.102378in}}{\pgfqpoint{0.955460in}{8.101284in}}%
\pgfpathquadraticcurveto{\pgfqpoint{0.950558in}{8.100190in}}{\pgfqpoint{0.945441in}{8.097866in}}%
\pgfpathlineto{\pgfqpoint{0.945441in}{8.143647in}}%
\pgfpathlineto{\pgfqpoint{0.945441in}{8.143647in}}%
\pgfpathclose%
\pgfusepath{fill}%
\end{pgfscope}%
\begin{pgfscope}%
\pgfpathrectangle{\pgfqpoint{1.109909in}{4.050000in}}{\pgfqpoint{5.400000in}{5.400000in}}%
\pgfusepath{clip}%
\pgfsetrectcap%
\pgfsetroundjoin%
\pgfsetlinewidth{0.602250pt}%
\definecolor{currentstroke}{rgb}{0.898039,0.913725,0.929412}%
\pgfsetstrokecolor{currentstroke}%
\pgfsetdash{}{0pt}%
\pgfpathmoveto{\pgfqpoint{1.109909in}{9.450000in}}%
\pgfpathlineto{\pgfqpoint{6.509909in}{9.450000in}}%
\pgfusepath{stroke}%
\end{pgfscope}%
\begin{pgfscope}%
\pgfsetbuttcap%
\pgfsetroundjoin%
\definecolor{currentfill}{rgb}{0.090196,0.168627,0.227451}%
\pgfsetfillcolor{currentfill}%
\pgfsetlinewidth{0.803000pt}%
\definecolor{currentstroke}{rgb}{0.090196,0.168627,0.227451}%
\pgfsetstrokecolor{currentstroke}%
\pgfsetdash{}{0pt}%
\pgfsys@defobject{currentmarker}{\pgfqpoint{-0.048611in}{0.000000in}}{\pgfqpoint{-0.000000in}{0.000000in}}{%
\pgfpathmoveto{\pgfqpoint{-0.000000in}{0.000000in}}%
\pgfpathlineto{\pgfqpoint{-0.048611in}{0.000000in}}%
\pgfusepath{stroke,fill}%
}%
\begin{pgfscope}%
\pgfsys@transformshift{1.109909in}{9.450000in}%
\pgfsys@useobject{currentmarker}{}%
\end{pgfscope}%
\end{pgfscope}%
\begin{pgfscope}%
\pgfsetbuttcap%
\pgfsetroundjoin%
\definecolor{currentfill}{rgb}{0.090196,0.168627,0.227451}%
\pgfsetfillcolor{currentfill}%
\pgfsetlinewidth{0.000000pt}%
\definecolor{currentstroke}{rgb}{0.090196,0.168627,0.227451}%
\pgfsetstrokecolor{currentstroke}%
\pgfsetdash{}{0pt}%
\pgfpathmoveto{\pgfqpoint{0.826183in}{9.493647in}}%
\pgfpathlineto{\pgfqpoint{0.874581in}{9.493647in}}%
\pgfpathlineto{\pgfqpoint{0.874581in}{9.483257in}}%
\pgfpathlineto{\pgfqpoint{0.837472in}{9.483257in}}%
\pgfpathlineto{\pgfqpoint{0.837472in}{9.460933in}}%
\pgfpathquadraticcurveto{\pgfqpoint{0.840148in}{9.461851in}}{\pgfqpoint{0.842824in}{9.462300in}}%
\pgfpathquadraticcurveto{\pgfqpoint{0.845519in}{9.462749in}}{\pgfqpoint{0.848214in}{9.462749in}}%
\pgfpathquadraticcurveto{\pgfqpoint{0.863468in}{9.462749in}}{\pgfqpoint{0.872374in}{9.454390in}}%
\pgfpathquadraticcurveto{\pgfqpoint{0.881300in}{9.446030in}}{\pgfqpoint{0.881300in}{9.431753in}}%
\pgfpathquadraticcurveto{\pgfqpoint{0.881300in}{9.417046in}}{\pgfqpoint{0.872140in}{9.408882in}}%
\pgfpathquadraticcurveto{\pgfqpoint{0.862980in}{9.400737in}}{\pgfqpoint{0.846320in}{9.400737in}}%
\pgfpathquadraticcurveto{\pgfqpoint{0.840577in}{9.400737in}}{\pgfqpoint{0.834620in}{9.401714in}}%
\pgfpathquadraticcurveto{\pgfqpoint{0.828683in}{9.402690in}}{\pgfqpoint{0.822335in}{9.404644in}}%
\pgfpathlineto{\pgfqpoint{0.822335in}{9.417046in}}%
\pgfpathquadraticcurveto{\pgfqpoint{0.827824in}{9.414058in}}{\pgfqpoint{0.833683in}{9.412593in}}%
\pgfpathquadraticcurveto{\pgfqpoint{0.839542in}{9.411128in}}{\pgfqpoint{0.846066in}{9.411128in}}%
\pgfpathquadraticcurveto{\pgfqpoint{0.856632in}{9.411128in}}{\pgfqpoint{0.862785in}{9.416675in}}%
\pgfpathquadraticcurveto{\pgfqpoint{0.868956in}{9.422222in}}{\pgfqpoint{0.868956in}{9.431753in}}%
\pgfpathquadraticcurveto{\pgfqpoint{0.868956in}{9.441265in}}{\pgfqpoint{0.862785in}{9.446812in}}%
\pgfpathquadraticcurveto{\pgfqpoint{0.856632in}{9.452378in}}{\pgfqpoint{0.846066in}{9.452378in}}%
\pgfpathquadraticcurveto{\pgfqpoint{0.841124in}{9.452378in}}{\pgfqpoint{0.836202in}{9.451284in}}%
\pgfpathquadraticcurveto{\pgfqpoint{0.831300in}{9.450190in}}{\pgfqpoint{0.826183in}{9.447866in}}%
\pgfpathlineto{\pgfqpoint{0.826183in}{9.493647in}}%
\pgfpathlineto{\pgfqpoint{0.826183in}{9.493647in}}%
\pgfpathclose%
\pgfpathmoveto{\pgfqpoint{0.905577in}{9.418022in}}%
\pgfpathlineto{\pgfqpoint{0.918468in}{9.418022in}}%
\pgfpathlineto{\pgfqpoint{0.918468in}{9.402515in}}%
\pgfpathlineto{\pgfqpoint{0.905577in}{9.402515in}}%
\pgfpathlineto{\pgfqpoint{0.905577in}{9.418022in}}%
\pgfpathlineto{\pgfqpoint{0.905577in}{9.418022in}}%
\pgfpathclose%
\pgfpathmoveto{\pgfqpoint{0.971671in}{9.485522in}}%
\pgfpathquadraticcurveto{\pgfqpoint{0.962160in}{9.485522in}}{\pgfqpoint{0.957355in}{9.476147in}}%
\pgfpathquadraticcurveto{\pgfqpoint{0.952570in}{9.466792in}}{\pgfqpoint{0.952570in}{9.447983in}}%
\pgfpathquadraticcurveto{\pgfqpoint{0.952570in}{9.429253in}}{\pgfqpoint{0.957355in}{9.419878in}}%
\pgfpathquadraticcurveto{\pgfqpoint{0.962160in}{9.410503in}}{\pgfqpoint{0.971671in}{9.410503in}}%
\pgfpathquadraticcurveto{\pgfqpoint{0.981261in}{9.410503in}}{\pgfqpoint{0.986046in}{9.419878in}}%
\pgfpathquadraticcurveto{\pgfqpoint{0.990851in}{9.429253in}}{\pgfqpoint{0.990851in}{9.447983in}}%
\pgfpathquadraticcurveto{\pgfqpoint{0.990851in}{9.466792in}}{\pgfqpoint{0.986046in}{9.476147in}}%
\pgfpathquadraticcurveto{\pgfqpoint{0.981261in}{9.485522in}}{\pgfqpoint{0.971671in}{9.485522in}}%
\pgfpathlineto{\pgfqpoint{0.971671in}{9.485522in}}%
\pgfpathclose%
\pgfpathmoveto{\pgfqpoint{0.971671in}{9.495288in}}%
\pgfpathquadraticcurveto{\pgfqpoint{0.987003in}{9.495288in}}{\pgfqpoint{0.995089in}{9.483159in}}%
\pgfpathquadraticcurveto{\pgfqpoint{1.003175in}{9.471050in}}{\pgfqpoint{1.003175in}{9.447983in}}%
\pgfpathquadraticcurveto{\pgfqpoint{1.003175in}{9.424976in}}{\pgfqpoint{0.995089in}{9.412847in}}%
\pgfpathquadraticcurveto{\pgfqpoint{0.987003in}{9.400737in}}{\pgfqpoint{0.971671in}{9.400737in}}%
\pgfpathquadraticcurveto{\pgfqpoint{0.956359in}{9.400737in}}{\pgfqpoint{0.948273in}{9.412847in}}%
\pgfpathquadraticcurveto{\pgfqpoint{0.940187in}{9.424976in}}{\pgfqpoint{0.940187in}{9.447983in}}%
\pgfpathquadraticcurveto{\pgfqpoint{0.940187in}{9.471050in}}{\pgfqpoint{0.948273in}{9.483159in}}%
\pgfpathquadraticcurveto{\pgfqpoint{0.956359in}{9.495288in}}{\pgfqpoint{0.971671in}{9.495288in}}%
\pgfpathlineto{\pgfqpoint{0.971671in}{9.495288in}}%
\pgfpathclose%
\pgfusepath{fill}%
\end{pgfscope}%
\begin{pgfscope}%
\pgfsetbuttcap%
\pgfsetroundjoin%
\definecolor{currentfill}{rgb}{0.090196,0.168627,0.227451}%
\pgfsetfillcolor{currentfill}%
\pgfsetlinewidth{0.000000pt}%
\definecolor{currentstroke}{rgb}{0.090196,0.168627,0.227451}%
\pgfsetstrokecolor{currentstroke}%
\pgfsetdash{}{0pt}%
\pgfpathmoveto{\pgfqpoint{0.452082in}{5.886992in}}%
\pgfpathlineto{\pgfqpoint{0.467957in}{5.886992in}}%
\pgfpathquadraticcurveto{\pgfqpoint{0.460867in}{5.879377in}}{\pgfqpoint{0.457382in}{5.870760in}}%
\pgfpathquadraticcurveto{\pgfqpoint{0.453873in}{5.862142in}}{\pgfqpoint{0.453873in}{5.852450in}}%
\pgfpathquadraticcurveto{\pgfqpoint{0.453873in}{5.833353in}}{\pgfqpoint{0.465546in}{5.823207in}}%
\pgfpathquadraticcurveto{\pgfqpoint{0.477219in}{5.813062in}}{\pgfqpoint{0.499300in}{5.813062in}}%
\pgfpathquadraticcurveto{\pgfqpoint{0.521310in}{5.813062in}}{\pgfqpoint{0.532983in}{5.823207in}}%
\pgfpathquadraticcurveto{\pgfqpoint{0.544656in}{5.833353in}}{\pgfqpoint{0.544656in}{5.852450in}}%
\pgfpathquadraticcurveto{\pgfqpoint{0.544656in}{5.862142in}}{\pgfqpoint{0.541147in}{5.870760in}}%
\pgfpathquadraticcurveto{\pgfqpoint{0.537638in}{5.879377in}}{\pgfqpoint{0.530548in}{5.886992in}}%
\pgfpathlineto{\pgfqpoint{0.546303in}{5.886992in}}%
\pgfpathquadraticcurveto{\pgfqpoint{0.551674in}{5.879091in}}{\pgfqpoint{0.554372in}{5.870234in}}%
\pgfpathquadraticcurveto{\pgfqpoint{0.557046in}{5.861402in}}{\pgfqpoint{0.557046in}{5.851567in}}%
\pgfpathquadraticcurveto{\pgfqpoint{0.557046in}{5.826263in}}{\pgfqpoint{0.541577in}{5.811701in}}%
\pgfpathquadraticcurveto{\pgfqpoint{0.526084in}{5.797164in}}{\pgfqpoint{0.499300in}{5.797164in}}%
\pgfpathquadraticcurveto{\pgfqpoint{0.472445in}{5.797164in}}{\pgfqpoint{0.456976in}{5.811701in}}%
\pgfpathquadraticcurveto{\pgfqpoint{0.441484in}{5.826263in}}{\pgfqpoint{0.441484in}{5.851567in}}%
\pgfpathquadraticcurveto{\pgfqpoint{0.441484in}{5.861545in}}{\pgfqpoint{0.444133in}{5.870378in}}%
\pgfpathquadraticcurveto{\pgfqpoint{0.446783in}{5.879234in}}{\pgfqpoint{0.452082in}{5.886992in}}%
\pgfpathlineto{\pgfqpoint{0.452082in}{5.886992in}}%
\pgfpathclose%
\pgfpathmoveto{\pgfqpoint{0.447141in}{5.977036in}}%
\pgfpathlineto{\pgfqpoint{0.461846in}{5.977036in}}%
\pgfpathquadraticcurveto{\pgfqpoint{0.457740in}{5.968466in}}{\pgfqpoint{0.455735in}{5.960851in}}%
\pgfpathquadraticcurveto{\pgfqpoint{0.453706in}{5.953236in}}{\pgfqpoint{0.453706in}{5.946146in}}%
\pgfpathquadraticcurveto{\pgfqpoint{0.453706in}{5.933852in}}{\pgfqpoint{0.458480in}{5.927168in}}%
\pgfpathquadraticcurveto{\pgfqpoint{0.463254in}{5.920484in}}{\pgfqpoint{0.472063in}{5.920484in}}%
\pgfpathquadraticcurveto{\pgfqpoint{0.479463in}{5.920484in}}{\pgfqpoint{0.483235in}{5.924924in}}%
\pgfpathquadraticcurveto{\pgfqpoint{0.486983in}{5.929364in}}{\pgfqpoint{0.489298in}{5.941753in}}%
\pgfpathlineto{\pgfqpoint{0.491160in}{5.950849in}}%
\pgfpathquadraticcurveto{\pgfqpoint{0.494383in}{5.967702in}}{\pgfqpoint{0.502475in}{5.975723in}}%
\pgfpathquadraticcurveto{\pgfqpoint{0.510568in}{5.983743in}}{\pgfqpoint{0.524127in}{5.983743in}}%
\pgfpathquadraticcurveto{\pgfqpoint{0.540336in}{5.983743in}}{\pgfqpoint{0.548691in}{5.972882in}}%
\pgfpathquadraticcurveto{\pgfqpoint{0.557046in}{5.962044in}}{\pgfqpoint{0.557046in}{5.941085in}}%
\pgfpathquadraticcurveto{\pgfqpoint{0.557046in}{5.933184in}}{\pgfqpoint{0.555255in}{5.924256in}}%
\pgfpathquadraticcurveto{\pgfqpoint{0.553465in}{5.915352in}}{\pgfqpoint{0.549956in}{5.905803in}}%
\pgfpathlineto{\pgfqpoint{0.534439in}{5.905803in}}%
\pgfpathquadraticcurveto{\pgfqpoint{0.539572in}{5.914970in}}{\pgfqpoint{0.542198in}{5.923778in}}%
\pgfpathquadraticcurveto{\pgfqpoint{0.544799in}{5.932587in}}{\pgfqpoint{0.544799in}{5.941085in}}%
\pgfpathquadraticcurveto{\pgfqpoint{0.544799in}{5.953976in}}{\pgfqpoint{0.539739in}{5.960994in}}%
\pgfpathquadraticcurveto{\pgfqpoint{0.534654in}{5.968012in}}{\pgfqpoint{0.525249in}{5.968012in}}%
\pgfpathquadraticcurveto{\pgfqpoint{0.517061in}{5.968012in}}{\pgfqpoint{0.512430in}{5.962975in}}%
\pgfpathquadraticcurveto{\pgfqpoint{0.507799in}{5.957938in}}{\pgfqpoint{0.505483in}{5.946456in}}%
\pgfpathlineto{\pgfqpoint{0.503693in}{5.937266in}}%
\pgfpathquadraticcurveto{\pgfqpoint{0.500351in}{5.920412in}}{\pgfqpoint{0.493189in}{5.912869in}}%
\pgfpathquadraticcurveto{\pgfqpoint{0.486028in}{5.905349in}}{\pgfqpoint{0.473257in}{5.905349in}}%
\pgfpathquadraticcurveto{\pgfqpoint{0.458480in}{5.905349in}}{\pgfqpoint{0.449982in}{5.915757in}}%
\pgfpathquadraticcurveto{\pgfqpoint{0.441484in}{5.926165in}}{\pgfqpoint{0.441484in}{5.944427in}}%
\pgfpathquadraticcurveto{\pgfqpoint{0.441484in}{5.952281in}}{\pgfqpoint{0.442916in}{5.960397in}}%
\pgfpathquadraticcurveto{\pgfqpoint{0.444324in}{5.968537in}}{\pgfqpoint{0.447141in}{5.977036in}}%
\pgfpathlineto{\pgfqpoint{0.447141in}{5.977036in}}%
\pgfpathclose%
\pgfpathmoveto{\pgfqpoint{0.458337in}{6.047337in}}%
\pgfpathlineto{\pgfqpoint{0.513767in}{6.026879in}}%
\pgfpathlineto{\pgfqpoint{0.513767in}{6.067843in}}%
\pgfpathlineto{\pgfqpoint{0.458337in}{6.047337in}}%
\pgfpathlineto{\pgfqpoint{0.458337in}{6.047337in}}%
\pgfpathclose%
\pgfpathmoveto{\pgfqpoint{0.443489in}{6.038815in}}%
\pgfpathlineto{\pgfqpoint{0.443489in}{6.055907in}}%
\pgfpathlineto{\pgfqpoint{0.554873in}{6.098351in}}%
\pgfpathlineto{\pgfqpoint{0.554873in}{6.082691in}}%
\pgfpathlineto{\pgfqpoint{0.526299in}{6.072546in}}%
\pgfpathlineto{\pgfqpoint{0.526299in}{6.022344in}}%
\pgfpathlineto{\pgfqpoint{0.554873in}{6.012198in}}%
\pgfpathlineto{\pgfqpoint{0.554873in}{5.996300in}}%
\pgfpathlineto{\pgfqpoint{0.443489in}{6.038815in}}%
\pgfpathlineto{\pgfqpoint{0.443489in}{6.038815in}}%
\pgfpathclose%
\pgfpathmoveto{\pgfqpoint{0.509661in}{6.234036in}}%
\pgfpathlineto{\pgfqpoint{0.516369in}{6.234036in}}%
\pgfpathlineto{\pgfqpoint{0.516369in}{6.170920in}}%
\pgfpathquadraticcurveto{\pgfqpoint{0.530548in}{6.171827in}}{\pgfqpoint{0.537972in}{6.179466in}}%
\pgfpathquadraticcurveto{\pgfqpoint{0.545396in}{6.187105in}}{\pgfqpoint{0.545396in}{6.200760in}}%
\pgfpathquadraticcurveto{\pgfqpoint{0.545396in}{6.208661in}}{\pgfqpoint{0.543463in}{6.216085in}}%
\pgfpathquadraticcurveto{\pgfqpoint{0.541529in}{6.223509in}}{\pgfqpoint{0.537638in}{6.230838in}}%
\pgfpathlineto{\pgfqpoint{0.550624in}{6.230838in}}%
\pgfpathquadraticcurveto{\pgfqpoint{0.553751in}{6.223438in}}{\pgfqpoint{0.555398in}{6.215679in}}%
\pgfpathquadraticcurveto{\pgfqpoint{0.557046in}{6.207921in}}{\pgfqpoint{0.557046in}{6.199948in}}%
\pgfpathquadraticcurveto{\pgfqpoint{0.557046in}{6.179944in}}{\pgfqpoint{0.545420in}{6.168270in}}%
\pgfpathquadraticcurveto{\pgfqpoint{0.533771in}{6.156597in}}{\pgfqpoint{0.513910in}{6.156597in}}%
\pgfpathquadraticcurveto{\pgfqpoint{0.493404in}{6.156597in}}{\pgfqpoint{0.481373in}{6.167674in}}%
\pgfpathquadraticcurveto{\pgfqpoint{0.469318in}{6.178750in}}{\pgfqpoint{0.469318in}{6.197561in}}%
\pgfpathquadraticcurveto{\pgfqpoint{0.469318in}{6.214414in}}{\pgfqpoint{0.480179in}{6.224225in}}%
\pgfpathquadraticcurveto{\pgfqpoint{0.491017in}{6.234036in}}{\pgfqpoint{0.509661in}{6.234036in}}%
\pgfpathlineto{\pgfqpoint{0.509661in}{6.234036in}}%
\pgfpathclose%
\pgfpathmoveto{\pgfqpoint{0.505626in}{6.220310in}}%
\pgfpathquadraticcurveto{\pgfqpoint{0.494383in}{6.220167in}}{\pgfqpoint{0.487675in}{6.214008in}}%
\pgfpathquadraticcurveto{\pgfqpoint{0.480943in}{6.207849in}}{\pgfqpoint{0.480943in}{6.197704in}}%
\pgfpathquadraticcurveto{\pgfqpoint{0.480943in}{6.186222in}}{\pgfqpoint{0.487436in}{6.179323in}}%
\pgfpathquadraticcurveto{\pgfqpoint{0.493929in}{6.172424in}}{\pgfqpoint{0.505722in}{6.171374in}}%
\pgfpathlineto{\pgfqpoint{0.505626in}{6.220310in}}%
\pgfpathlineto{\pgfqpoint{0.505626in}{6.220310in}}%
\pgfpathclose%
\pgfpathmoveto{\pgfqpoint{0.473782in}{6.309829in}}%
\pgfpathlineto{\pgfqpoint{0.486768in}{6.309829in}}%
\pgfpathquadraticcurveto{\pgfqpoint{0.483784in}{6.304028in}}{\pgfqpoint{0.482304in}{6.297750in}}%
\pgfpathquadraticcurveto{\pgfqpoint{0.480800in}{6.291495in}}{\pgfqpoint{0.480800in}{6.284763in}}%
\pgfpathquadraticcurveto{\pgfqpoint{0.480800in}{6.274546in}}{\pgfqpoint{0.483927in}{6.269438in}}%
\pgfpathquadraticcurveto{\pgfqpoint{0.487054in}{6.264329in}}{\pgfqpoint{0.493332in}{6.264329in}}%
\pgfpathquadraticcurveto{\pgfqpoint{0.498107in}{6.264329in}}{\pgfqpoint{0.500828in}{6.267982in}}%
\pgfpathquadraticcurveto{\pgfqpoint{0.503549in}{6.271634in}}{\pgfqpoint{0.506008in}{6.282687in}}%
\pgfpathlineto{\pgfqpoint{0.507059in}{6.287389in}}%
\pgfpathquadraticcurveto{\pgfqpoint{0.510186in}{6.301999in}}{\pgfqpoint{0.515891in}{6.308158in}}%
\pgfpathquadraticcurveto{\pgfqpoint{0.521596in}{6.314316in}}{\pgfqpoint{0.531813in}{6.314316in}}%
\pgfpathquadraticcurveto{\pgfqpoint{0.543463in}{6.314316in}}{\pgfqpoint{0.550266in}{6.305102in}}%
\pgfpathquadraticcurveto{\pgfqpoint{0.557046in}{6.295888in}}{\pgfqpoint{0.557046in}{6.279774in}}%
\pgfpathquadraticcurveto{\pgfqpoint{0.557046in}{6.273066in}}{\pgfqpoint{0.555733in}{6.265786in}}%
\pgfpathquadraticcurveto{\pgfqpoint{0.554420in}{6.258505in}}{\pgfqpoint{0.551818in}{6.250460in}}%
\pgfpathlineto{\pgfqpoint{0.537638in}{6.250460in}}%
\pgfpathquadraticcurveto{\pgfqpoint{0.541601in}{6.258075in}}{\pgfqpoint{0.543582in}{6.265451in}}%
\pgfpathquadraticcurveto{\pgfqpoint{0.545540in}{6.272828in}}{\pgfqpoint{0.545540in}{6.280085in}}%
\pgfpathquadraticcurveto{\pgfqpoint{0.545540in}{6.289776in}}{\pgfqpoint{0.542221in}{6.294980in}}%
\pgfpathquadraticcurveto{\pgfqpoint{0.538903in}{6.300208in}}{\pgfqpoint{0.532864in}{6.300208in}}%
\pgfpathquadraticcurveto{\pgfqpoint{0.527278in}{6.300208in}}{\pgfqpoint{0.524294in}{6.296437in}}%
\pgfpathquadraticcurveto{\pgfqpoint{0.521310in}{6.292689in}}{\pgfqpoint{0.518541in}{6.279918in}}%
\pgfpathlineto{\pgfqpoint{0.517419in}{6.275143in}}%
\pgfpathquadraticcurveto{\pgfqpoint{0.514745in}{6.262396in}}{\pgfqpoint{0.509183in}{6.256714in}}%
\pgfpathquadraticcurveto{\pgfqpoint{0.503621in}{6.251057in}}{\pgfqpoint{0.493929in}{6.251057in}}%
\pgfpathquadraticcurveto{\pgfqpoint{0.482137in}{6.251057in}}{\pgfqpoint{0.475739in}{6.259412in}}%
\pgfpathquadraticcurveto{\pgfqpoint{0.469318in}{6.267767in}}{\pgfqpoint{0.469318in}{6.283140in}}%
\pgfpathquadraticcurveto{\pgfqpoint{0.469318in}{6.290731in}}{\pgfqpoint{0.470440in}{6.297439in}}%
\pgfpathquadraticcurveto{\pgfqpoint{0.471538in}{6.304171in}}{\pgfqpoint{0.473782in}{6.309829in}}%
\pgfpathlineto{\pgfqpoint{0.473782in}{6.309829in}}%
\pgfpathclose%
\pgfpathmoveto{\pgfqpoint{0.447595in}{6.349742in}}%
\pgfpathlineto{\pgfqpoint{0.471323in}{6.349742in}}%
\pgfpathlineto{\pgfqpoint{0.471323in}{6.378006in}}%
\pgfpathlineto{\pgfqpoint{0.481994in}{6.378006in}}%
\pgfpathlineto{\pgfqpoint{0.481994in}{6.349742in}}%
\pgfpathlineto{\pgfqpoint{0.527349in}{6.349742in}}%
\pgfpathquadraticcurveto{\pgfqpoint{0.537566in}{6.349742in}}{\pgfqpoint{0.540479in}{6.352535in}}%
\pgfpathquadraticcurveto{\pgfqpoint{0.543391in}{6.355328in}}{\pgfqpoint{0.543391in}{6.363921in}}%
\pgfpathlineto{\pgfqpoint{0.543391in}{6.378006in}}%
\pgfpathlineto{\pgfqpoint{0.554873in}{6.378006in}}%
\pgfpathlineto{\pgfqpoint{0.554873in}{6.363921in}}%
\pgfpathquadraticcurveto{\pgfqpoint{0.554873in}{6.348023in}}{\pgfqpoint{0.548953in}{6.341984in}}%
\pgfpathquadraticcurveto{\pgfqpoint{0.543009in}{6.335944in}}{\pgfqpoint{0.527349in}{6.335944in}}%
\pgfpathlineto{\pgfqpoint{0.481994in}{6.335944in}}%
\pgfpathlineto{\pgfqpoint{0.481994in}{6.325870in}}%
\pgfpathlineto{\pgfqpoint{0.471323in}{6.325870in}}%
\pgfpathlineto{\pgfqpoint{0.471323in}{6.335944in}}%
\pgfpathlineto{\pgfqpoint{0.447595in}{6.335944in}}%
\pgfpathlineto{\pgfqpoint{0.447595in}{6.349742in}}%
\pgfpathlineto{\pgfqpoint{0.447595in}{6.349742in}}%
\pgfpathclose%
\pgfpathmoveto{\pgfqpoint{0.471323in}{6.396053in}}%
\pgfpathlineto{\pgfqpoint{0.471323in}{6.409779in}}%
\pgfpathlineto{\pgfqpoint{0.554873in}{6.409779in}}%
\pgfpathlineto{\pgfqpoint{0.554873in}{6.396053in}}%
\pgfpathlineto{\pgfqpoint{0.471323in}{6.396053in}}%
\pgfpathlineto{\pgfqpoint{0.471323in}{6.396053in}}%
\pgfpathclose%
\pgfpathmoveto{\pgfqpoint{0.438786in}{6.396053in}}%
\pgfpathlineto{\pgfqpoint{0.438786in}{6.409779in}}%
\pgfpathlineto{\pgfqpoint{0.456188in}{6.409779in}}%
\pgfpathlineto{\pgfqpoint{0.456188in}{6.396053in}}%
\pgfpathlineto{\pgfqpoint{0.438786in}{6.396053in}}%
\pgfpathlineto{\pgfqpoint{0.438786in}{6.396053in}}%
\pgfpathclose%
\pgfpathmoveto{\pgfqpoint{0.487365in}{6.503546in}}%
\pgfpathquadraticcurveto{\pgfqpoint{0.478102in}{6.508702in}}{\pgfqpoint{0.473710in}{6.515864in}}%
\pgfpathquadraticcurveto{\pgfqpoint{0.469318in}{6.523025in}}{\pgfqpoint{0.469318in}{6.532717in}}%
\pgfpathquadraticcurveto{\pgfqpoint{0.469318in}{6.545775in}}{\pgfqpoint{0.478461in}{6.552865in}}%
\pgfpathquadraticcurveto{\pgfqpoint{0.487579in}{6.559954in}}{\pgfqpoint{0.504433in}{6.559954in}}%
\pgfpathlineto{\pgfqpoint{0.554873in}{6.559954in}}%
\pgfpathlineto{\pgfqpoint{0.554873in}{6.546157in}}%
\pgfpathlineto{\pgfqpoint{0.504886in}{6.546157in}}%
\pgfpathquadraticcurveto{\pgfqpoint{0.492879in}{6.546157in}}{\pgfqpoint{0.487078in}{6.541884in}}%
\pgfpathquadraticcurveto{\pgfqpoint{0.481253in}{6.537635in}}{\pgfqpoint{0.481253in}{6.528921in}}%
\pgfpathquadraticcurveto{\pgfqpoint{0.481253in}{6.518251in}}{\pgfqpoint{0.488343in}{6.512044in}}%
\pgfpathquadraticcurveto{\pgfqpoint{0.495409in}{6.505862in}}{\pgfqpoint{0.507655in}{6.505862in}}%
\pgfpathlineto{\pgfqpoint{0.554873in}{6.505862in}}%
\pgfpathlineto{\pgfqpoint{0.554873in}{6.492064in}}%
\pgfpathlineto{\pgfqpoint{0.504886in}{6.492064in}}%
\pgfpathquadraticcurveto{\pgfqpoint{0.492807in}{6.492064in}}{\pgfqpoint{0.487030in}{6.487815in}}%
\pgfpathquadraticcurveto{\pgfqpoint{0.481253in}{6.483566in}}{\pgfqpoint{0.481253in}{6.474685in}}%
\pgfpathquadraticcurveto{\pgfqpoint{0.481253in}{6.464158in}}{\pgfqpoint{0.488367in}{6.457951in}}%
\pgfpathquadraticcurveto{\pgfqpoint{0.495481in}{6.451769in}}{\pgfqpoint{0.507655in}{6.451769in}}%
\pgfpathlineto{\pgfqpoint{0.554873in}{6.451769in}}%
\pgfpathlineto{\pgfqpoint{0.554873in}{6.437971in}}%
\pgfpathlineto{\pgfqpoint{0.471323in}{6.437971in}}%
\pgfpathlineto{\pgfqpoint{0.471323in}{6.451769in}}%
\pgfpathlineto{\pgfqpoint{0.484309in}{6.451769in}}%
\pgfpathquadraticcurveto{\pgfqpoint{0.476622in}{6.456471in}}{\pgfqpoint{0.472970in}{6.463036in}}%
\pgfpathquadraticcurveto{\pgfqpoint{0.469318in}{6.469601in}}{\pgfqpoint{0.469318in}{6.478624in}}%
\pgfpathquadraticcurveto{\pgfqpoint{0.469318in}{6.487743in}}{\pgfqpoint{0.473949in}{6.494117in}}%
\pgfpathquadraticcurveto{\pgfqpoint{0.478556in}{6.500490in}}{\pgfqpoint{0.487365in}{6.503546in}}%
\pgfpathlineto{\pgfqpoint{0.487365in}{6.503546in}}%
\pgfpathclose%
\pgfpathmoveto{\pgfqpoint{0.512883in}{6.625291in}}%
\pgfpathquadraticcurveto{\pgfqpoint{0.512883in}{6.608652in}}{\pgfqpoint{0.516679in}{6.602231in}}%
\pgfpathquadraticcurveto{\pgfqpoint{0.520474in}{6.595809in}}{\pgfqpoint{0.529665in}{6.595809in}}%
\pgfpathquadraticcurveto{\pgfqpoint{0.536970in}{6.595809in}}{\pgfqpoint{0.541267in}{6.600632in}}%
\pgfpathquadraticcurveto{\pgfqpoint{0.545540in}{6.605454in}}{\pgfqpoint{0.545540in}{6.613713in}}%
\pgfpathquadraticcurveto{\pgfqpoint{0.545540in}{6.625148in}}{\pgfqpoint{0.537447in}{6.632046in}}%
\pgfpathquadraticcurveto{\pgfqpoint{0.529355in}{6.638945in}}{\pgfqpoint{0.515939in}{6.638945in}}%
\pgfpathlineto{\pgfqpoint{0.512883in}{6.638945in}}%
\pgfpathlineto{\pgfqpoint{0.512883in}{6.625291in}}%
\pgfpathlineto{\pgfqpoint{0.512883in}{6.625291in}}%
\pgfpathclose%
\pgfpathmoveto{\pgfqpoint{0.507202in}{6.652671in}}%
\pgfpathlineto{\pgfqpoint{0.554873in}{6.652671in}}%
\pgfpathlineto{\pgfqpoint{0.554873in}{6.638945in}}%
\pgfpathlineto{\pgfqpoint{0.542198in}{6.638945in}}%
\pgfpathquadraticcurveto{\pgfqpoint{0.549789in}{6.634243in}}{\pgfqpoint{0.553417in}{6.627224in}}%
\pgfpathquadraticcurveto{\pgfqpoint{0.557046in}{6.620206in}}{\pgfqpoint{0.557046in}{6.610061in}}%
\pgfpathquadraticcurveto{\pgfqpoint{0.557046in}{6.597242in}}{\pgfqpoint{0.549836in}{6.589651in}}%
\pgfpathquadraticcurveto{\pgfqpoint{0.542627in}{6.582083in}}{\pgfqpoint{0.530548in}{6.582083in}}%
\pgfpathquadraticcurveto{\pgfqpoint{0.516464in}{6.582083in}}{\pgfqpoint{0.509303in}{6.591513in}}%
\pgfpathquadraticcurveto{\pgfqpoint{0.502141in}{6.600966in}}{\pgfqpoint{0.502141in}{6.619681in}}%
\pgfpathlineto{\pgfqpoint{0.502141in}{6.638945in}}%
\pgfpathlineto{\pgfqpoint{0.500780in}{6.638945in}}%
\pgfpathquadraticcurveto{\pgfqpoint{0.491303in}{6.638945in}}{\pgfqpoint{0.486123in}{6.632715in}}%
\pgfpathquadraticcurveto{\pgfqpoint{0.480943in}{6.626484in}}{\pgfqpoint{0.480943in}{6.615217in}}%
\pgfpathquadraticcurveto{\pgfqpoint{0.480943in}{6.608056in}}{\pgfqpoint{0.482662in}{6.601252in}}%
\pgfpathquadraticcurveto{\pgfqpoint{0.484381in}{6.594473in}}{\pgfqpoint{0.487818in}{6.588218in}}%
\pgfpathlineto{\pgfqpoint{0.475119in}{6.588218in}}%
\pgfpathquadraticcurveto{\pgfqpoint{0.472206in}{6.595738in}}{\pgfqpoint{0.470774in}{6.602828in}}%
\pgfpathquadraticcurveto{\pgfqpoint{0.469318in}{6.609918in}}{\pgfqpoint{0.469318in}{6.616625in}}%
\pgfpathquadraticcurveto{\pgfqpoint{0.469318in}{6.634768in}}{\pgfqpoint{0.478723in}{6.643720in}}%
\pgfpathquadraticcurveto{\pgfqpoint{0.488105in}{6.652671in}}{\pgfqpoint{0.507202in}{6.652671in}}%
\pgfpathlineto{\pgfqpoint{0.507202in}{6.652671in}}%
\pgfpathclose%
\pgfpathmoveto{\pgfqpoint{0.447595in}{6.694518in}}%
\pgfpathlineto{\pgfqpoint{0.471323in}{6.694518in}}%
\pgfpathlineto{\pgfqpoint{0.471323in}{6.722782in}}%
\pgfpathlineto{\pgfqpoint{0.481994in}{6.722782in}}%
\pgfpathlineto{\pgfqpoint{0.481994in}{6.694518in}}%
\pgfpathlineto{\pgfqpoint{0.527349in}{6.694518in}}%
\pgfpathquadraticcurveto{\pgfqpoint{0.537566in}{6.694518in}}{\pgfqpoint{0.540479in}{6.697311in}}%
\pgfpathquadraticcurveto{\pgfqpoint{0.543391in}{6.700104in}}{\pgfqpoint{0.543391in}{6.708698in}}%
\pgfpathlineto{\pgfqpoint{0.543391in}{6.722782in}}%
\pgfpathlineto{\pgfqpoint{0.554873in}{6.722782in}}%
\pgfpathlineto{\pgfqpoint{0.554873in}{6.708698in}}%
\pgfpathquadraticcurveto{\pgfqpoint{0.554873in}{6.692799in}}{\pgfqpoint{0.548953in}{6.686760in}}%
\pgfpathquadraticcurveto{\pgfqpoint{0.543009in}{6.680720in}}{\pgfqpoint{0.527349in}{6.680720in}}%
\pgfpathlineto{\pgfqpoint{0.481994in}{6.680720in}}%
\pgfpathlineto{\pgfqpoint{0.481994in}{6.670647in}}%
\pgfpathlineto{\pgfqpoint{0.471323in}{6.670647in}}%
\pgfpathlineto{\pgfqpoint{0.471323in}{6.680720in}}%
\pgfpathlineto{\pgfqpoint{0.447595in}{6.680720in}}%
\pgfpathlineto{\pgfqpoint{0.447595in}{6.694518in}}%
\pgfpathlineto{\pgfqpoint{0.447595in}{6.694518in}}%
\pgfpathclose%
\pgfpathmoveto{\pgfqpoint{0.509661in}{6.812300in}}%
\pgfpathlineto{\pgfqpoint{0.516369in}{6.812300in}}%
\pgfpathlineto{\pgfqpoint{0.516369in}{6.749184in}}%
\pgfpathquadraticcurveto{\pgfqpoint{0.530548in}{6.750091in}}{\pgfqpoint{0.537972in}{6.757730in}}%
\pgfpathquadraticcurveto{\pgfqpoint{0.545396in}{6.765369in}}{\pgfqpoint{0.545396in}{6.779023in}}%
\pgfpathquadraticcurveto{\pgfqpoint{0.545396in}{6.786925in}}{\pgfqpoint{0.543463in}{6.794349in}}%
\pgfpathquadraticcurveto{\pgfqpoint{0.541529in}{6.801773in}}{\pgfqpoint{0.537638in}{6.809102in}}%
\pgfpathlineto{\pgfqpoint{0.550624in}{6.809102in}}%
\pgfpathquadraticcurveto{\pgfqpoint{0.553751in}{6.801701in}}{\pgfqpoint{0.555398in}{6.793943in}}%
\pgfpathquadraticcurveto{\pgfqpoint{0.557046in}{6.786185in}}{\pgfqpoint{0.557046in}{6.778212in}}%
\pgfpathquadraticcurveto{\pgfqpoint{0.557046in}{6.758207in}}{\pgfqpoint{0.545420in}{6.746534in}}%
\pgfpathquadraticcurveto{\pgfqpoint{0.533771in}{6.734861in}}{\pgfqpoint{0.513910in}{6.734861in}}%
\pgfpathquadraticcurveto{\pgfqpoint{0.493404in}{6.734861in}}{\pgfqpoint{0.481373in}{6.745938in}}%
\pgfpathquadraticcurveto{\pgfqpoint{0.469318in}{6.757014in}}{\pgfqpoint{0.469318in}{6.775825in}}%
\pgfpathquadraticcurveto{\pgfqpoint{0.469318in}{6.792678in}}{\pgfqpoint{0.480179in}{6.802489in}}%
\pgfpathquadraticcurveto{\pgfqpoint{0.491017in}{6.812300in}}{\pgfqpoint{0.509661in}{6.812300in}}%
\pgfpathlineto{\pgfqpoint{0.509661in}{6.812300in}}%
\pgfpathclose%
\pgfpathmoveto{\pgfqpoint{0.505626in}{6.798574in}}%
\pgfpathquadraticcurveto{\pgfqpoint{0.494383in}{6.798431in}}{\pgfqpoint{0.487675in}{6.792272in}}%
\pgfpathquadraticcurveto{\pgfqpoint{0.480943in}{6.786113in}}{\pgfqpoint{0.480943in}{6.775968in}}%
\pgfpathquadraticcurveto{\pgfqpoint{0.480943in}{6.764486in}}{\pgfqpoint{0.487436in}{6.757587in}}%
\pgfpathquadraticcurveto{\pgfqpoint{0.493929in}{6.750688in}}{\pgfqpoint{0.505722in}{6.749638in}}%
\pgfpathlineto{\pgfqpoint{0.505626in}{6.798574in}}%
\pgfpathlineto{\pgfqpoint{0.505626in}{6.798574in}}%
\pgfpathclose%
\pgfpathmoveto{\pgfqpoint{0.438953in}{6.916356in}}%
\pgfpathquadraticcurveto{\pgfqpoint{0.456093in}{6.906378in}}{\pgfqpoint{0.472898in}{6.901508in}}%
\pgfpathquadraticcurveto{\pgfqpoint{0.489680in}{6.896662in}}{\pgfqpoint{0.506915in}{6.896662in}}%
\pgfpathquadraticcurveto{\pgfqpoint{0.524127in}{6.896662in}}{\pgfqpoint{0.541028in}{6.901556in}}%
\pgfpathquadraticcurveto{\pgfqpoint{0.557929in}{6.906450in}}{\pgfqpoint{0.575021in}{6.916356in}}%
\pgfpathlineto{\pgfqpoint{0.575021in}{6.904421in}}%
\pgfpathquadraticcurveto{\pgfqpoint{0.557475in}{6.893249in}}{\pgfqpoint{0.540550in}{6.887687in}}%
\pgfpathquadraticcurveto{\pgfqpoint{0.523625in}{6.882125in}}{\pgfqpoint{0.506915in}{6.882125in}}%
\pgfpathquadraticcurveto{\pgfqpoint{0.490277in}{6.882125in}}{\pgfqpoint{0.473424in}{6.887639in}}%
\pgfpathquadraticcurveto{\pgfqpoint{0.456546in}{6.893177in}}{\pgfqpoint{0.438953in}{6.904421in}}%
\pgfpathlineto{\pgfqpoint{0.438953in}{6.916356in}}%
\pgfpathlineto{\pgfqpoint{0.438953in}{6.916356in}}%
\pgfpathclose%
\pgfpathmoveto{\pgfqpoint{0.443489in}{6.928149in}}%
\pgfpathlineto{\pgfqpoint{0.443489in}{7.022370in}}%
\pgfpathlineto{\pgfqpoint{0.456188in}{7.022370in}}%
\pgfpathlineto{\pgfqpoint{0.456188in}{6.982839in}}%
\pgfpathlineto{\pgfqpoint{0.554873in}{6.982839in}}%
\pgfpathlineto{\pgfqpoint{0.554873in}{6.967704in}}%
\pgfpathlineto{\pgfqpoint{0.456188in}{6.967704in}}%
\pgfpathlineto{\pgfqpoint{0.456188in}{6.928149in}}%
\pgfpathlineto{\pgfqpoint{0.443489in}{6.928149in}}%
\pgfpathlineto{\pgfqpoint{0.443489in}{6.928149in}}%
\pgfpathclose%
\pgfpathmoveto{\pgfqpoint{0.443489in}{7.027001in}}%
\pgfpathlineto{\pgfqpoint{0.443489in}{7.042207in}}%
\pgfpathlineto{\pgfqpoint{0.537638in}{7.065625in}}%
\pgfpathlineto{\pgfqpoint{0.443489in}{7.088971in}}%
\pgfpathlineto{\pgfqpoint{0.443489in}{7.105920in}}%
\pgfpathlineto{\pgfqpoint{0.537638in}{7.129338in}}%
\pgfpathlineto{\pgfqpoint{0.443489in}{7.152684in}}%
\pgfpathlineto{\pgfqpoint{0.443489in}{7.167986in}}%
\pgfpathlineto{\pgfqpoint{0.554873in}{7.140009in}}%
\pgfpathlineto{\pgfqpoint{0.554873in}{7.121055in}}%
\pgfpathlineto{\pgfqpoint{0.458194in}{7.097565in}}%
\pgfpathlineto{\pgfqpoint{0.554873in}{7.073837in}}%
\pgfpathlineto{\pgfqpoint{0.554873in}{7.054883in}}%
\pgfpathlineto{\pgfqpoint{0.443489in}{7.027001in}}%
\pgfpathlineto{\pgfqpoint{0.443489in}{7.027001in}}%
\pgfpathclose%
\pgfpathmoveto{\pgfqpoint{0.443489in}{7.187967in}}%
\pgfpathlineto{\pgfqpoint{0.443489in}{7.251966in}}%
\pgfpathlineto{\pgfqpoint{0.456188in}{7.251966in}}%
\pgfpathlineto{\pgfqpoint{0.456188in}{7.203030in}}%
\pgfpathlineto{\pgfqpoint{0.489012in}{7.203030in}}%
\pgfpathlineto{\pgfqpoint{0.489012in}{7.247192in}}%
\pgfpathlineto{\pgfqpoint{0.501688in}{7.247192in}}%
\pgfpathlineto{\pgfqpoint{0.501688in}{7.203030in}}%
\pgfpathlineto{\pgfqpoint{0.554873in}{7.203030in}}%
\pgfpathlineto{\pgfqpoint{0.554873in}{7.187967in}}%
\pgfpathlineto{\pgfqpoint{0.443489in}{7.187967in}}%
\pgfpathlineto{\pgfqpoint{0.443489in}{7.187967in}}%
\pgfpathclose%
\pgfpathmoveto{\pgfqpoint{0.443489in}{7.275838in}}%
\pgfpathlineto{\pgfqpoint{0.443489in}{7.346259in}}%
\pgfpathlineto{\pgfqpoint{0.456188in}{7.346259in}}%
\pgfpathlineto{\pgfqpoint{0.456188in}{7.290901in}}%
\pgfpathlineto{\pgfqpoint{0.489155in}{7.290901in}}%
\pgfpathlineto{\pgfqpoint{0.489155in}{7.343943in}}%
\pgfpathlineto{\pgfqpoint{0.501831in}{7.343943in}}%
\pgfpathlineto{\pgfqpoint{0.501831in}{7.290901in}}%
\pgfpathlineto{\pgfqpoint{0.542198in}{7.290901in}}%
\pgfpathlineto{\pgfqpoint{0.542198in}{7.347595in}}%
\pgfpathlineto{\pgfqpoint{0.554873in}{7.347595in}}%
\pgfpathlineto{\pgfqpoint{0.554873in}{7.275838in}}%
\pgfpathlineto{\pgfqpoint{0.443489in}{7.275838in}}%
\pgfpathlineto{\pgfqpoint{0.443489in}{7.275838in}}%
\pgfpathclose%
\pgfpathmoveto{\pgfqpoint{0.447141in}{7.487697in}}%
\pgfpathlineto{\pgfqpoint{0.461846in}{7.487697in}}%
\pgfpathquadraticcurveto{\pgfqpoint{0.457740in}{7.479128in}}{\pgfqpoint{0.455735in}{7.471513in}}%
\pgfpathquadraticcurveto{\pgfqpoint{0.453706in}{7.463898in}}{\pgfqpoint{0.453706in}{7.456808in}}%
\pgfpathquadraticcurveto{\pgfqpoint{0.453706in}{7.444514in}}{\pgfqpoint{0.458480in}{7.437830in}}%
\pgfpathquadraticcurveto{\pgfqpoint{0.463254in}{7.431146in}}{\pgfqpoint{0.472063in}{7.431146in}}%
\pgfpathquadraticcurveto{\pgfqpoint{0.479463in}{7.431146in}}{\pgfqpoint{0.483235in}{7.435586in}}%
\pgfpathquadraticcurveto{\pgfqpoint{0.486983in}{7.440026in}}{\pgfqpoint{0.489298in}{7.452415in}}%
\pgfpathlineto{\pgfqpoint{0.491160in}{7.461510in}}%
\pgfpathquadraticcurveto{\pgfqpoint{0.494383in}{7.478364in}}{\pgfqpoint{0.502475in}{7.486385in}}%
\pgfpathquadraticcurveto{\pgfqpoint{0.510568in}{7.494405in}}{\pgfqpoint{0.524127in}{7.494405in}}%
\pgfpathquadraticcurveto{\pgfqpoint{0.540336in}{7.494405in}}{\pgfqpoint{0.548691in}{7.483544in}}%
\pgfpathquadraticcurveto{\pgfqpoint{0.557046in}{7.472706in}}{\pgfqpoint{0.557046in}{7.451747in}}%
\pgfpathquadraticcurveto{\pgfqpoint{0.557046in}{7.443845in}}{\pgfqpoint{0.555255in}{7.434918in}}%
\pgfpathquadraticcurveto{\pgfqpoint{0.553465in}{7.426013in}}{\pgfqpoint{0.549956in}{7.416465in}}%
\pgfpathlineto{\pgfqpoint{0.534439in}{7.416465in}}%
\pgfpathquadraticcurveto{\pgfqpoint{0.539572in}{7.425632in}}{\pgfqpoint{0.542198in}{7.434440in}}%
\pgfpathquadraticcurveto{\pgfqpoint{0.544799in}{7.443249in}}{\pgfqpoint{0.544799in}{7.451747in}}%
\pgfpathquadraticcurveto{\pgfqpoint{0.544799in}{7.464638in}}{\pgfqpoint{0.539739in}{7.471656in}}%
\pgfpathquadraticcurveto{\pgfqpoint{0.534654in}{7.478674in}}{\pgfqpoint{0.525249in}{7.478674in}}%
\pgfpathquadraticcurveto{\pgfqpoint{0.517061in}{7.478674in}}{\pgfqpoint{0.512430in}{7.473637in}}%
\pgfpathquadraticcurveto{\pgfqpoint{0.507799in}{7.468600in}}{\pgfqpoint{0.505483in}{7.457118in}}%
\pgfpathlineto{\pgfqpoint{0.503693in}{7.447928in}}%
\pgfpathquadraticcurveto{\pgfqpoint{0.500351in}{7.431074in}}{\pgfqpoint{0.493189in}{7.423531in}}%
\pgfpathquadraticcurveto{\pgfqpoint{0.486028in}{7.416011in}}{\pgfqpoint{0.473257in}{7.416011in}}%
\pgfpathquadraticcurveto{\pgfqpoint{0.458480in}{7.416011in}}{\pgfqpoint{0.449982in}{7.426419in}}%
\pgfpathquadraticcurveto{\pgfqpoint{0.441484in}{7.436827in}}{\pgfqpoint{0.441484in}{7.455089in}}%
\pgfpathquadraticcurveto{\pgfqpoint{0.441484in}{7.462943in}}{\pgfqpoint{0.442916in}{7.471059in}}%
\pgfpathquadraticcurveto{\pgfqpoint{0.444324in}{7.479199in}}{\pgfqpoint{0.447141in}{7.487697in}}%
\pgfpathlineto{\pgfqpoint{0.447141in}{7.487697in}}%
\pgfpathclose%
\pgfpathmoveto{\pgfqpoint{0.443489in}{7.517919in}}%
\pgfpathlineto{\pgfqpoint{0.443489in}{7.588340in}}%
\pgfpathlineto{\pgfqpoint{0.456188in}{7.588340in}}%
\pgfpathlineto{\pgfqpoint{0.456188in}{7.532982in}}%
\pgfpathlineto{\pgfqpoint{0.489155in}{7.532982in}}%
\pgfpathlineto{\pgfqpoint{0.489155in}{7.586024in}}%
\pgfpathlineto{\pgfqpoint{0.501831in}{7.586024in}}%
\pgfpathlineto{\pgfqpoint{0.501831in}{7.532982in}}%
\pgfpathlineto{\pgfqpoint{0.542198in}{7.532982in}}%
\pgfpathlineto{\pgfqpoint{0.542198in}{7.589677in}}%
\pgfpathlineto{\pgfqpoint{0.554873in}{7.589677in}}%
\pgfpathlineto{\pgfqpoint{0.554873in}{7.517919in}}%
\pgfpathlineto{\pgfqpoint{0.443489in}{7.517919in}}%
\pgfpathlineto{\pgfqpoint{0.443489in}{7.517919in}}%
\pgfpathclose%
\pgfpathmoveto{\pgfqpoint{0.473782in}{7.667116in}}%
\pgfpathlineto{\pgfqpoint{0.486768in}{7.667116in}}%
\pgfpathquadraticcurveto{\pgfqpoint{0.483784in}{7.661315in}}{\pgfqpoint{0.482304in}{7.655037in}}%
\pgfpathquadraticcurveto{\pgfqpoint{0.480800in}{7.648783in}}{\pgfqpoint{0.480800in}{7.642051in}}%
\pgfpathquadraticcurveto{\pgfqpoint{0.480800in}{7.631834in}}{\pgfqpoint{0.483927in}{7.626725in}}%
\pgfpathquadraticcurveto{\pgfqpoint{0.487054in}{7.621617in}}{\pgfqpoint{0.493332in}{7.621617in}}%
\pgfpathquadraticcurveto{\pgfqpoint{0.498107in}{7.621617in}}{\pgfqpoint{0.500828in}{7.625269in}}%
\pgfpathquadraticcurveto{\pgfqpoint{0.503549in}{7.628921in}}{\pgfqpoint{0.506008in}{7.639974in}}%
\pgfpathlineto{\pgfqpoint{0.507059in}{7.644677in}}%
\pgfpathquadraticcurveto{\pgfqpoint{0.510186in}{7.659286in}}{\pgfqpoint{0.515891in}{7.665445in}}%
\pgfpathquadraticcurveto{\pgfqpoint{0.521596in}{7.671604in}}{\pgfqpoint{0.531813in}{7.671604in}}%
\pgfpathquadraticcurveto{\pgfqpoint{0.543463in}{7.671604in}}{\pgfqpoint{0.550266in}{7.662389in}}%
\pgfpathquadraticcurveto{\pgfqpoint{0.557046in}{7.653175in}}{\pgfqpoint{0.557046in}{7.637062in}}%
\pgfpathquadraticcurveto{\pgfqpoint{0.557046in}{7.630354in}}{\pgfqpoint{0.555733in}{7.623073in}}%
\pgfpathquadraticcurveto{\pgfqpoint{0.554420in}{7.615792in}}{\pgfqpoint{0.551818in}{7.607747in}}%
\pgfpathlineto{\pgfqpoint{0.537638in}{7.607747in}}%
\pgfpathquadraticcurveto{\pgfqpoint{0.541601in}{7.615362in}}{\pgfqpoint{0.543582in}{7.622739in}}%
\pgfpathquadraticcurveto{\pgfqpoint{0.545540in}{7.630115in}}{\pgfqpoint{0.545540in}{7.637372in}}%
\pgfpathquadraticcurveto{\pgfqpoint{0.545540in}{7.647064in}}{\pgfqpoint{0.542221in}{7.652268in}}%
\pgfpathquadraticcurveto{\pgfqpoint{0.538903in}{7.657496in}}{\pgfqpoint{0.532864in}{7.657496in}}%
\pgfpathquadraticcurveto{\pgfqpoint{0.527278in}{7.657496in}}{\pgfqpoint{0.524294in}{7.653724in}}%
\pgfpathquadraticcurveto{\pgfqpoint{0.521310in}{7.649976in}}{\pgfqpoint{0.518541in}{7.637205in}}%
\pgfpathlineto{\pgfqpoint{0.517419in}{7.632431in}}%
\pgfpathquadraticcurveto{\pgfqpoint{0.514745in}{7.619683in}}{\pgfqpoint{0.509183in}{7.614002in}}%
\pgfpathquadraticcurveto{\pgfqpoint{0.503621in}{7.608344in}}{\pgfqpoint{0.493929in}{7.608344in}}%
\pgfpathquadraticcurveto{\pgfqpoint{0.482137in}{7.608344in}}{\pgfqpoint{0.475739in}{7.616699in}}%
\pgfpathquadraticcurveto{\pgfqpoint{0.469318in}{7.625054in}}{\pgfqpoint{0.469318in}{7.640428in}}%
\pgfpathquadraticcurveto{\pgfqpoint{0.469318in}{7.648019in}}{\pgfqpoint{0.470440in}{7.654727in}}%
\pgfpathquadraticcurveto{\pgfqpoint{0.471538in}{7.661458in}}{\pgfqpoint{0.473782in}{7.667116in}}%
\pgfpathlineto{\pgfqpoint{0.473782in}{7.667116in}}%
\pgfpathclose%
\pgfpathmoveto{\pgfqpoint{0.438953in}{7.691298in}}%
\pgfpathlineto{\pgfqpoint{0.438953in}{7.703234in}}%
\pgfpathquadraticcurveto{\pgfqpoint{0.456546in}{7.714405in}}{\pgfqpoint{0.473424in}{7.719967in}}%
\pgfpathquadraticcurveto{\pgfqpoint{0.490277in}{7.725530in}}{\pgfqpoint{0.506915in}{7.725530in}}%
\pgfpathquadraticcurveto{\pgfqpoint{0.523625in}{7.725530in}}{\pgfqpoint{0.540550in}{7.719967in}}%
\pgfpathquadraticcurveto{\pgfqpoint{0.557475in}{7.714405in}}{\pgfqpoint{0.575021in}{7.703234in}}%
\pgfpathlineto{\pgfqpoint{0.575021in}{7.691298in}}%
\pgfpathquadraticcurveto{\pgfqpoint{0.557929in}{7.701204in}}{\pgfqpoint{0.541028in}{7.706098in}}%
\pgfpathquadraticcurveto{\pgfqpoint{0.524127in}{7.710992in}}{\pgfqpoint{0.506915in}{7.710992in}}%
\pgfpathquadraticcurveto{\pgfqpoint{0.489680in}{7.710992in}}{\pgfqpoint{0.472898in}{7.706098in}}%
\pgfpathquadraticcurveto{\pgfqpoint{0.456093in}{7.701204in}}{\pgfqpoint{0.438953in}{7.691298in}}%
\pgfpathlineto{\pgfqpoint{0.438953in}{7.691298in}}%
\pgfpathclose%
\pgfusepath{fill}%
\end{pgfscope}%
\begin{pgfscope}%
\pgfpathrectangle{\pgfqpoint{1.109909in}{4.050000in}}{\pgfqpoint{5.400000in}{5.400000in}}%
\pgfusepath{clip}%
\pgfsetrectcap%
\pgfsetroundjoin%
\pgfsetlinewidth{0.853187pt}%
\definecolor{currentstroke}{rgb}{0.745098,0.776471,0.803922}%
\pgfsetstrokecolor{currentstroke}%
\pgfsetdash{}{0pt}%
\pgfpathmoveto{\pgfqpoint{1.109909in}{6.750000in}}%
\pgfpathlineto{\pgfqpoint{6.509909in}{6.750000in}}%
\pgfusepath{stroke}%
\end{pgfscope}%
\begin{pgfscope}%
\pgfpathrectangle{\pgfqpoint{1.109909in}{4.050000in}}{\pgfqpoint{5.400000in}{5.400000in}}%
\pgfusepath{clip}%
\pgfsetrectcap%
\pgfsetroundjoin%
\pgfsetlinewidth{0.853187pt}%
\definecolor{currentstroke}{rgb}{0.745098,0.776471,0.803922}%
\pgfsetstrokecolor{currentstroke}%
\pgfsetdash{}{0pt}%
\pgfpathmoveto{\pgfqpoint{3.809909in}{4.050000in}}%
\pgfpathlineto{\pgfqpoint{3.809909in}{9.450000in}}%
\pgfusepath{stroke}%
\end{pgfscope}%
\begin{pgfscope}%
\pgfpathrectangle{\pgfqpoint{1.109909in}{4.050000in}}{\pgfqpoint{5.400000in}{5.400000in}}%
\pgfusepath{clip}%
\pgfsetbuttcap%
\pgfsetroundjoin%
\pgfsetlinewidth{1.104125pt}%
\definecolor{currentstroke}{rgb}{0.203922,0.243137,0.278431}%
\pgfsetstrokecolor{currentstroke}%
\pgfsetdash{{5.500000pt}{4.400000pt}}{0.000000pt}%
\pgfpathmoveto{\pgfqpoint{1.109909in}{4.050000in}}%
\pgfpathlineto{\pgfqpoint{6.509909in}{9.450000in}}%
\pgfusepath{stroke}%
\end{pgfscope}%
\begin{pgfscope}%
\pgfsetrectcap%
\pgfsetmiterjoin%
\pgfsetlinewidth{0.803000pt}%
\definecolor{currentstroke}{rgb}{0.627451,0.670588,0.705882}%
\pgfsetstrokecolor{currentstroke}%
\pgfsetdash{}{0pt}%
\pgfpathmoveto{\pgfqpoint{1.109909in}{4.050000in}}%
\pgfpathlineto{\pgfqpoint{1.109909in}{9.450000in}}%
\pgfusepath{stroke}%
\end{pgfscope}%
\begin{pgfscope}%
\pgfsetrectcap%
\pgfsetmiterjoin%
\pgfsetlinewidth{0.803000pt}%
\definecolor{currentstroke}{rgb}{0.627451,0.670588,0.705882}%
\pgfsetstrokecolor{currentstroke}%
\pgfsetdash{}{0pt}%
\pgfpathmoveto{\pgfqpoint{1.109909in}{4.050000in}}%
\pgfpathlineto{\pgfqpoint{6.509909in}{4.050000in}}%
\pgfusepath{stroke}%
\end{pgfscope}%
\begin{pgfscope}%
\pgfpathrectangle{\pgfqpoint{1.109909in}{4.050000in}}{\pgfqpoint{5.400000in}{5.400000in}}%
\pgfusepath{clip}%
\pgfsetbuttcap%
\pgfsetroundjoin%
\pgfsetlinewidth{1.003750pt}%
\definecolor{currentstroke}{rgb}{0.541176,0.380392,0.658824}%
\pgfsetstrokecolor{currentstroke}%
\pgfsetstrokeopacity{0.900000}%
\pgfsetdash{}{0pt}%
\pgfpathmoveto{\pgfqpoint{5.043551in}{8.576759in}}%
\pgfpathcurveto{\pgfqpoint{5.053805in}{8.576759in}}{\pgfqpoint{5.063640in}{8.580833in}}{\pgfqpoint{5.070891in}{8.588084in}}%
\pgfpathcurveto{\pgfqpoint{5.078142in}{8.595335in}}{\pgfqpoint{5.082216in}{8.605170in}}{\pgfqpoint{5.082216in}{8.615424in}}%
\pgfpathcurveto{\pgfqpoint{5.082216in}{8.625679in}}{\pgfqpoint{5.078142in}{8.635514in}}{\pgfqpoint{5.070891in}{8.642765in}}%
\pgfpathcurveto{\pgfqpoint{5.063640in}{8.650015in}}{\pgfqpoint{5.053805in}{8.654089in}}{\pgfqpoint{5.043551in}{8.654089in}}%
\pgfpathcurveto{\pgfqpoint{5.033297in}{8.654089in}}{\pgfqpoint{5.023461in}{8.650015in}}{\pgfqpoint{5.016210in}{8.642765in}}%
\pgfpathcurveto{\pgfqpoint{5.008960in}{8.635514in}}{\pgfqpoint{5.004886in}{8.625679in}}{\pgfqpoint{5.004886in}{8.615424in}}%
\pgfpathcurveto{\pgfqpoint{5.004886in}{8.605170in}}{\pgfqpoint{5.008960in}{8.595335in}}{\pgfqpoint{5.016210in}{8.588084in}}%
\pgfpathcurveto{\pgfqpoint{5.023461in}{8.580833in}}{\pgfqpoint{5.033297in}{8.576759in}}{\pgfqpoint{5.043551in}{8.576759in}}%
\pgfpathlineto{\pgfqpoint{5.043551in}{8.576759in}}%
\pgfpathclose%
\pgfusepath{stroke}%
\end{pgfscope}%
\begin{pgfscope}%
\pgfpathrectangle{\pgfqpoint{1.109909in}{4.050000in}}{\pgfqpoint{5.400000in}{5.400000in}}%
\pgfusepath{clip}%
\pgfsetbuttcap%
\pgfsetroundjoin%
\pgfsetlinewidth{1.003750pt}%
\definecolor{currentstroke}{rgb}{0.541176,0.380392,0.658824}%
\pgfsetstrokecolor{currentstroke}%
\pgfsetstrokeopacity{0.900000}%
\pgfsetdash{}{0pt}%
\pgfpathmoveto{\pgfqpoint{4.702823in}{8.163420in}}%
\pgfpathcurveto{\pgfqpoint{4.713077in}{8.163420in}}{\pgfqpoint{4.722912in}{8.167494in}}{\pgfqpoint{4.730163in}{8.174745in}}%
\pgfpathcurveto{\pgfqpoint{4.737414in}{8.181995in}}{\pgfqpoint{4.741488in}{8.191831in}}{\pgfqpoint{4.741488in}{8.202085in}}%
\pgfpathcurveto{\pgfqpoint{4.741488in}{8.212339in}}{\pgfqpoint{4.737414in}{8.222174in}}{\pgfqpoint{4.730163in}{8.229425in}}%
\pgfpathcurveto{\pgfqpoint{4.722912in}{8.236676in}}{\pgfqpoint{4.713077in}{8.240750in}}{\pgfqpoint{4.702823in}{8.240750in}}%
\pgfpathcurveto{\pgfqpoint{4.692569in}{8.240750in}}{\pgfqpoint{4.682733in}{8.236676in}}{\pgfqpoint{4.675482in}{8.229425in}}%
\pgfpathcurveto{\pgfqpoint{4.668232in}{8.222174in}}{\pgfqpoint{4.664158in}{8.212339in}}{\pgfqpoint{4.664158in}{8.202085in}}%
\pgfpathcurveto{\pgfqpoint{4.664158in}{8.191831in}}{\pgfqpoint{4.668232in}{8.181995in}}{\pgfqpoint{4.675482in}{8.174745in}}%
\pgfpathcurveto{\pgfqpoint{4.682733in}{8.167494in}}{\pgfqpoint{4.692569in}{8.163420in}}{\pgfqpoint{4.702823in}{8.163420in}}%
\pgfpathlineto{\pgfqpoint{4.702823in}{8.163420in}}%
\pgfpathclose%
\pgfusepath{stroke}%
\end{pgfscope}%
\begin{pgfscope}%
\pgfpathrectangle{\pgfqpoint{1.109909in}{4.050000in}}{\pgfqpoint{5.400000in}{5.400000in}}%
\pgfusepath{clip}%
\pgfsetbuttcap%
\pgfsetroundjoin%
\pgfsetlinewidth{1.003750pt}%
\definecolor{currentstroke}{rgb}{0.541176,0.380392,0.658824}%
\pgfsetstrokecolor{currentstroke}%
\pgfsetstrokeopacity{0.900000}%
\pgfsetdash{}{0pt}%
\pgfpathmoveto{\pgfqpoint{5.589912in}{8.598223in}}%
\pgfpathcurveto{\pgfqpoint{5.600166in}{8.598223in}}{\pgfqpoint{5.610001in}{8.602297in}}{\pgfqpoint{5.617252in}{8.609548in}}%
\pgfpathcurveto{\pgfqpoint{5.624503in}{8.616799in}}{\pgfqpoint{5.628577in}{8.626634in}}{\pgfqpoint{5.628577in}{8.636888in}}%
\pgfpathcurveto{\pgfqpoint{5.628577in}{8.647142in}}{\pgfqpoint{5.624503in}{8.656978in}}{\pgfqpoint{5.617252in}{8.664229in}}%
\pgfpathcurveto{\pgfqpoint{5.610001in}{8.671479in}}{\pgfqpoint{5.600166in}{8.675553in}}{\pgfqpoint{5.589912in}{8.675553in}}%
\pgfpathcurveto{\pgfqpoint{5.579658in}{8.675553in}}{\pgfqpoint{5.569822in}{8.671479in}}{\pgfqpoint{5.562571in}{8.664229in}}%
\pgfpathcurveto{\pgfqpoint{5.555321in}{8.656978in}}{\pgfqpoint{5.551247in}{8.647142in}}{\pgfqpoint{5.551247in}{8.636888in}}%
\pgfpathcurveto{\pgfqpoint{5.551247in}{8.626634in}}{\pgfqpoint{5.555321in}{8.616799in}}{\pgfqpoint{5.562571in}{8.609548in}}%
\pgfpathcurveto{\pgfqpoint{5.569822in}{8.602297in}}{\pgfqpoint{5.579658in}{8.598223in}}{\pgfqpoint{5.589912in}{8.598223in}}%
\pgfpathlineto{\pgfqpoint{5.589912in}{8.598223in}}%
\pgfpathclose%
\pgfusepath{stroke}%
\end{pgfscope}%
\begin{pgfscope}%
\pgfpathrectangle{\pgfqpoint{1.109909in}{4.050000in}}{\pgfqpoint{5.400000in}{5.400000in}}%
\pgfusepath{clip}%
\pgfsetbuttcap%
\pgfsetroundjoin%
\pgfsetlinewidth{1.003750pt}%
\definecolor{currentstroke}{rgb}{0.541176,0.380392,0.658824}%
\pgfsetstrokecolor{currentstroke}%
\pgfsetstrokeopacity{0.900000}%
\pgfsetdash{}{0pt}%
\pgfpathmoveto{\pgfqpoint{4.527112in}{7.433096in}}%
\pgfpathcurveto{\pgfqpoint{4.537366in}{7.433096in}}{\pgfqpoint{4.547201in}{7.437170in}}{\pgfqpoint{4.554452in}{7.444421in}}%
\pgfpathcurveto{\pgfqpoint{4.561703in}{7.451671in}}{\pgfqpoint{4.565777in}{7.461507in}}{\pgfqpoint{4.565777in}{7.471761in}}%
\pgfpathcurveto{\pgfqpoint{4.565777in}{7.482015in}}{\pgfqpoint{4.561703in}{7.491850in}}{\pgfqpoint{4.554452in}{7.499101in}}%
\pgfpathcurveto{\pgfqpoint{4.547201in}{7.506352in}}{\pgfqpoint{4.537366in}{7.510426in}}{\pgfqpoint{4.527112in}{7.510426in}}%
\pgfpathcurveto{\pgfqpoint{4.516858in}{7.510426in}}{\pgfqpoint{4.507022in}{7.506352in}}{\pgfqpoint{4.499772in}{7.499101in}}%
\pgfpathcurveto{\pgfqpoint{4.492521in}{7.491850in}}{\pgfqpoint{4.488447in}{7.482015in}}{\pgfqpoint{4.488447in}{7.471761in}}%
\pgfpathcurveto{\pgfqpoint{4.488447in}{7.461507in}}{\pgfqpoint{4.492521in}{7.451671in}}{\pgfqpoint{4.499772in}{7.444421in}}%
\pgfpathcurveto{\pgfqpoint{4.507022in}{7.437170in}}{\pgfqpoint{4.516858in}{7.433096in}}{\pgfqpoint{4.527112in}{7.433096in}}%
\pgfpathlineto{\pgfqpoint{4.527112in}{7.433096in}}%
\pgfpathclose%
\pgfusepath{stroke}%
\end{pgfscope}%
\begin{pgfscope}%
\pgfpathrectangle{\pgfqpoint{1.109909in}{4.050000in}}{\pgfqpoint{5.400000in}{5.400000in}}%
\pgfusepath{clip}%
\pgfsetbuttcap%
\pgfsetroundjoin%
\pgfsetlinewidth{1.003750pt}%
\definecolor{currentstroke}{rgb}{0.541176,0.380392,0.658824}%
\pgfsetstrokecolor{currentstroke}%
\pgfsetstrokeopacity{0.900000}%
\pgfsetdash{}{0pt}%
\pgfpathmoveto{\pgfqpoint{5.188432in}{8.693200in}}%
\pgfpathcurveto{\pgfqpoint{5.198686in}{8.693200in}}{\pgfqpoint{5.208522in}{8.697274in}}{\pgfqpoint{5.215773in}{8.704525in}}%
\pgfpathcurveto{\pgfqpoint{5.223023in}{8.711775in}}{\pgfqpoint{5.227097in}{8.721611in}}{\pgfqpoint{5.227097in}{8.731865in}}%
\pgfpathcurveto{\pgfqpoint{5.227097in}{8.742119in}}{\pgfqpoint{5.223023in}{8.751954in}}{\pgfqpoint{5.215773in}{8.759205in}}%
\pgfpathcurveto{\pgfqpoint{5.208522in}{8.766456in}}{\pgfqpoint{5.198686in}{8.770530in}}{\pgfqpoint{5.188432in}{8.770530in}}%
\pgfpathcurveto{\pgfqpoint{5.178178in}{8.770530in}}{\pgfqpoint{5.168343in}{8.766456in}}{\pgfqpoint{5.161092in}{8.759205in}}%
\pgfpathcurveto{\pgfqpoint{5.153841in}{8.751954in}}{\pgfqpoint{5.149767in}{8.742119in}}{\pgfqpoint{5.149767in}{8.731865in}}%
\pgfpathcurveto{\pgfqpoint{5.149767in}{8.721611in}}{\pgfqpoint{5.153841in}{8.711775in}}{\pgfqpoint{5.161092in}{8.704525in}}%
\pgfpathcurveto{\pgfqpoint{5.168343in}{8.697274in}}{\pgfqpoint{5.178178in}{8.693200in}}{\pgfqpoint{5.188432in}{8.693200in}}%
\pgfpathlineto{\pgfqpoint{5.188432in}{8.693200in}}%
\pgfpathclose%
\pgfusepath{stroke}%
\end{pgfscope}%
\begin{pgfscope}%
\pgfpathrectangle{\pgfqpoint{1.109909in}{4.050000in}}{\pgfqpoint{5.400000in}{5.400000in}}%
\pgfusepath{clip}%
\pgfsetbuttcap%
\pgfsetroundjoin%
\pgfsetlinewidth{1.003750pt}%
\definecolor{currentstroke}{rgb}{0.541176,0.380392,0.658824}%
\pgfsetstrokecolor{currentstroke}%
\pgfsetstrokeopacity{0.900000}%
\pgfsetdash{}{0pt}%
\pgfpathmoveto{\pgfqpoint{4.731819in}{7.368039in}}%
\pgfpathcurveto{\pgfqpoint{4.742073in}{7.368039in}}{\pgfqpoint{4.751909in}{7.372113in}}{\pgfqpoint{4.759160in}{7.379363in}}%
\pgfpathcurveto{\pgfqpoint{4.766410in}{7.386614in}}{\pgfqpoint{4.770484in}{7.396450in}}{\pgfqpoint{4.770484in}{7.406704in}}%
\pgfpathcurveto{\pgfqpoint{4.770484in}{7.416958in}}{\pgfqpoint{4.766410in}{7.426793in}}{\pgfqpoint{4.759160in}{7.434044in}}%
\pgfpathcurveto{\pgfqpoint{4.751909in}{7.441295in}}{\pgfqpoint{4.742073in}{7.445369in}}{\pgfqpoint{4.731819in}{7.445369in}}%
\pgfpathcurveto{\pgfqpoint{4.721565in}{7.445369in}}{\pgfqpoint{4.711730in}{7.441295in}}{\pgfqpoint{4.704479in}{7.434044in}}%
\pgfpathcurveto{\pgfqpoint{4.697228in}{7.426793in}}{\pgfqpoint{4.693154in}{7.416958in}}{\pgfqpoint{4.693154in}{7.406704in}}%
\pgfpathcurveto{\pgfqpoint{4.693154in}{7.396450in}}{\pgfqpoint{4.697228in}{7.386614in}}{\pgfqpoint{4.704479in}{7.379363in}}%
\pgfpathcurveto{\pgfqpoint{4.711730in}{7.372113in}}{\pgfqpoint{4.721565in}{7.368039in}}{\pgfqpoint{4.731819in}{7.368039in}}%
\pgfpathlineto{\pgfqpoint{4.731819in}{7.368039in}}%
\pgfpathclose%
\pgfusepath{stroke}%
\end{pgfscope}%
\begin{pgfscope}%
\pgfpathrectangle{\pgfqpoint{1.109909in}{4.050000in}}{\pgfqpoint{5.400000in}{5.400000in}}%
\pgfusepath{clip}%
\pgfsetbuttcap%
\pgfsetroundjoin%
\pgfsetlinewidth{1.003750pt}%
\definecolor{currentstroke}{rgb}{0.541176,0.380392,0.658824}%
\pgfsetstrokecolor{currentstroke}%
\pgfsetstrokeopacity{0.900000}%
\pgfsetdash{}{0pt}%
\pgfpathmoveto{\pgfqpoint{5.275237in}{8.405781in}}%
\pgfpathcurveto{\pgfqpoint{5.285491in}{8.405781in}}{\pgfqpoint{5.295327in}{8.409855in}}{\pgfqpoint{5.302578in}{8.417106in}}%
\pgfpathcurveto{\pgfqpoint{5.309828in}{8.424357in}}{\pgfqpoint{5.313902in}{8.434192in}}{\pgfqpoint{5.313902in}{8.444446in}}%
\pgfpathcurveto{\pgfqpoint{5.313902in}{8.454700in}}{\pgfqpoint{5.309828in}{8.464536in}}{\pgfqpoint{5.302578in}{8.471786in}}%
\pgfpathcurveto{\pgfqpoint{5.295327in}{8.479037in}}{\pgfqpoint{5.285491in}{8.483111in}}{\pgfqpoint{5.275237in}{8.483111in}}%
\pgfpathcurveto{\pgfqpoint{5.264983in}{8.483111in}}{\pgfqpoint{5.255148in}{8.479037in}}{\pgfqpoint{5.247897in}{8.471786in}}%
\pgfpathcurveto{\pgfqpoint{5.240646in}{8.464536in}}{\pgfqpoint{5.236572in}{8.454700in}}{\pgfqpoint{5.236572in}{8.444446in}}%
\pgfpathcurveto{\pgfqpoint{5.236572in}{8.434192in}}{\pgfqpoint{5.240646in}{8.424357in}}{\pgfqpoint{5.247897in}{8.417106in}}%
\pgfpathcurveto{\pgfqpoint{5.255148in}{8.409855in}}{\pgfqpoint{5.264983in}{8.405781in}}{\pgfqpoint{5.275237in}{8.405781in}}%
\pgfpathlineto{\pgfqpoint{5.275237in}{8.405781in}}%
\pgfpathclose%
\pgfusepath{stroke}%
\end{pgfscope}%
\begin{pgfscope}%
\pgfpathrectangle{\pgfqpoint{1.109909in}{4.050000in}}{\pgfqpoint{5.400000in}{5.400000in}}%
\pgfusepath{clip}%
\pgfsetbuttcap%
\pgfsetroundjoin%
\pgfsetlinewidth{1.003750pt}%
\definecolor{currentstroke}{rgb}{0.541176,0.380392,0.658824}%
\pgfsetstrokecolor{currentstroke}%
\pgfsetstrokeopacity{0.900000}%
\pgfsetdash{}{0pt}%
\pgfpathmoveto{\pgfqpoint{3.804486in}{7.094566in}}%
\pgfpathcurveto{\pgfqpoint{3.814740in}{7.094566in}}{\pgfqpoint{3.824575in}{7.098640in}}{\pgfqpoint{3.831826in}{7.105891in}}%
\pgfpathcurveto{\pgfqpoint{3.839077in}{7.113142in}}{\pgfqpoint{3.843151in}{7.122977in}}{\pgfqpoint{3.843151in}{7.133231in}}%
\pgfpathcurveto{\pgfqpoint{3.843151in}{7.143486in}}{\pgfqpoint{3.839077in}{7.153321in}}{\pgfqpoint{3.831826in}{7.160572in}}%
\pgfpathcurveto{\pgfqpoint{3.824575in}{7.167823in}}{\pgfqpoint{3.814740in}{7.171897in}}{\pgfqpoint{3.804486in}{7.171897in}}%
\pgfpathcurveto{\pgfqpoint{3.794232in}{7.171897in}}{\pgfqpoint{3.784396in}{7.167823in}}{\pgfqpoint{3.777146in}{7.160572in}}%
\pgfpathcurveto{\pgfqpoint{3.769895in}{7.153321in}}{\pgfqpoint{3.765821in}{7.143486in}}{\pgfqpoint{3.765821in}{7.133231in}}%
\pgfpathcurveto{\pgfqpoint{3.765821in}{7.122977in}}{\pgfqpoint{3.769895in}{7.113142in}}{\pgfqpoint{3.777146in}{7.105891in}}%
\pgfpathcurveto{\pgfqpoint{3.784396in}{7.098640in}}{\pgfqpoint{3.794232in}{7.094566in}}{\pgfqpoint{3.804486in}{7.094566in}}%
\pgfpathlineto{\pgfqpoint{3.804486in}{7.094566in}}%
\pgfpathclose%
\pgfusepath{stroke}%
\end{pgfscope}%
\begin{pgfscope}%
\pgfpathrectangle{\pgfqpoint{1.109909in}{4.050000in}}{\pgfqpoint{5.400000in}{5.400000in}}%
\pgfusepath{clip}%
\pgfsetbuttcap%
\pgfsetroundjoin%
\pgfsetlinewidth{1.003750pt}%
\definecolor{currentstroke}{rgb}{0.541176,0.380392,0.658824}%
\pgfsetstrokecolor{currentstroke}%
\pgfsetstrokeopacity{0.900000}%
\pgfsetdash{}{0pt}%
\pgfpathmoveto{\pgfqpoint{4.614842in}{6.366205in}}%
\pgfpathcurveto{\pgfqpoint{4.625096in}{6.366205in}}{\pgfqpoint{4.634931in}{6.370279in}}{\pgfqpoint{4.642182in}{6.377530in}}%
\pgfpathcurveto{\pgfqpoint{4.649433in}{6.384780in}}{\pgfqpoint{4.653507in}{6.394616in}}{\pgfqpoint{4.653507in}{6.404870in}}%
\pgfpathcurveto{\pgfqpoint{4.653507in}{6.415124in}}{\pgfqpoint{4.649433in}{6.424960in}}{\pgfqpoint{4.642182in}{6.432210in}}%
\pgfpathcurveto{\pgfqpoint{4.634931in}{6.439461in}}{\pgfqpoint{4.625096in}{6.443535in}}{\pgfqpoint{4.614842in}{6.443535in}}%
\pgfpathcurveto{\pgfqpoint{4.604588in}{6.443535in}}{\pgfqpoint{4.594752in}{6.439461in}}{\pgfqpoint{4.587502in}{6.432210in}}%
\pgfpathcurveto{\pgfqpoint{4.580251in}{6.424960in}}{\pgfqpoint{4.576177in}{6.415124in}}{\pgfqpoint{4.576177in}{6.404870in}}%
\pgfpathcurveto{\pgfqpoint{4.576177in}{6.394616in}}{\pgfqpoint{4.580251in}{6.384780in}}{\pgfqpoint{4.587502in}{6.377530in}}%
\pgfpathcurveto{\pgfqpoint{4.594752in}{6.370279in}}{\pgfqpoint{4.604588in}{6.366205in}}{\pgfqpoint{4.614842in}{6.366205in}}%
\pgfpathlineto{\pgfqpoint{4.614842in}{6.366205in}}%
\pgfpathclose%
\pgfusepath{stroke}%
\end{pgfscope}%
\begin{pgfscope}%
\pgfpathrectangle{\pgfqpoint{1.109909in}{4.050000in}}{\pgfqpoint{5.400000in}{5.400000in}}%
\pgfusepath{clip}%
\pgfsetbuttcap%
\pgfsetroundjoin%
\pgfsetlinewidth{1.003750pt}%
\definecolor{currentstroke}{rgb}{0.541176,0.380392,0.658824}%
\pgfsetstrokecolor{currentstroke}%
\pgfsetstrokeopacity{0.900000}%
\pgfsetdash{}{0pt}%
\pgfpathmoveto{\pgfqpoint{5.200068in}{7.982702in}}%
\pgfpathcurveto{\pgfqpoint{5.210322in}{7.982702in}}{\pgfqpoint{5.220157in}{7.986776in}}{\pgfqpoint{5.227408in}{7.994027in}}%
\pgfpathcurveto{\pgfqpoint{5.234659in}{8.001278in}}{\pgfqpoint{5.238733in}{8.011113in}}{\pgfqpoint{5.238733in}{8.021367in}}%
\pgfpathcurveto{\pgfqpoint{5.238733in}{8.031621in}}{\pgfqpoint{5.234659in}{8.041457in}}{\pgfqpoint{5.227408in}{8.048708in}}%
\pgfpathcurveto{\pgfqpoint{5.220157in}{8.055958in}}{\pgfqpoint{5.210322in}{8.060032in}}{\pgfqpoint{5.200068in}{8.060032in}}%
\pgfpathcurveto{\pgfqpoint{5.189814in}{8.060032in}}{\pgfqpoint{5.179978in}{8.055958in}}{\pgfqpoint{5.172727in}{8.048708in}}%
\pgfpathcurveto{\pgfqpoint{5.165477in}{8.041457in}}{\pgfqpoint{5.161403in}{8.031621in}}{\pgfqpoint{5.161403in}{8.021367in}}%
\pgfpathcurveto{\pgfqpoint{5.161403in}{8.011113in}}{\pgfqpoint{5.165477in}{8.001278in}}{\pgfqpoint{5.172727in}{7.994027in}}%
\pgfpathcurveto{\pgfqpoint{5.179978in}{7.986776in}}{\pgfqpoint{5.189814in}{7.982702in}}{\pgfqpoint{5.200068in}{7.982702in}}%
\pgfpathlineto{\pgfqpoint{5.200068in}{7.982702in}}%
\pgfpathclose%
\pgfusepath{stroke}%
\end{pgfscope}%
\begin{pgfscope}%
\pgfpathrectangle{\pgfqpoint{1.109909in}{4.050000in}}{\pgfqpoint{5.400000in}{5.400000in}}%
\pgfusepath{clip}%
\pgfsetbuttcap%
\pgfsetroundjoin%
\pgfsetlinewidth{1.003750pt}%
\definecolor{currentstroke}{rgb}{0.541176,0.380392,0.658824}%
\pgfsetstrokecolor{currentstroke}%
\pgfsetstrokeopacity{0.900000}%
\pgfsetdash{}{0pt}%
\pgfpathmoveto{\pgfqpoint{4.619955in}{6.614121in}}%
\pgfpathcurveto{\pgfqpoint{4.630209in}{6.614121in}}{\pgfqpoint{4.640044in}{6.618195in}}{\pgfqpoint{4.647295in}{6.625445in}}%
\pgfpathcurveto{\pgfqpoint{4.654546in}{6.632696in}}{\pgfqpoint{4.658620in}{6.642532in}}{\pgfqpoint{4.658620in}{6.652786in}}%
\pgfpathcurveto{\pgfqpoint{4.658620in}{6.663040in}}{\pgfqpoint{4.654546in}{6.672875in}}{\pgfqpoint{4.647295in}{6.680126in}}%
\pgfpathcurveto{\pgfqpoint{4.640044in}{6.687377in}}{\pgfqpoint{4.630209in}{6.691451in}}{\pgfqpoint{4.619955in}{6.691451in}}%
\pgfpathcurveto{\pgfqpoint{4.609701in}{6.691451in}}{\pgfqpoint{4.599865in}{6.687377in}}{\pgfqpoint{4.592614in}{6.680126in}}%
\pgfpathcurveto{\pgfqpoint{4.585364in}{6.672875in}}{\pgfqpoint{4.581290in}{6.663040in}}{\pgfqpoint{4.581290in}{6.652786in}}%
\pgfpathcurveto{\pgfqpoint{4.581290in}{6.642532in}}{\pgfqpoint{4.585364in}{6.632696in}}{\pgfqpoint{4.592614in}{6.625445in}}%
\pgfpathcurveto{\pgfqpoint{4.599865in}{6.618195in}}{\pgfqpoint{4.609701in}{6.614121in}}{\pgfqpoint{4.619955in}{6.614121in}}%
\pgfpathlineto{\pgfqpoint{4.619955in}{6.614121in}}%
\pgfpathclose%
\pgfusepath{stroke}%
\end{pgfscope}%
\begin{pgfscope}%
\pgfpathrectangle{\pgfqpoint{1.109909in}{4.050000in}}{\pgfqpoint{5.400000in}{5.400000in}}%
\pgfusepath{clip}%
\pgfsetbuttcap%
\pgfsetroundjoin%
\pgfsetlinewidth{1.003750pt}%
\definecolor{currentstroke}{rgb}{0.541176,0.380392,0.658824}%
\pgfsetstrokecolor{currentstroke}%
\pgfsetstrokeopacity{0.900000}%
\pgfsetdash{}{0pt}%
\pgfpathmoveto{\pgfqpoint{4.887800in}{8.030041in}}%
\pgfpathcurveto{\pgfqpoint{4.898055in}{8.030041in}}{\pgfqpoint{4.907890in}{8.034115in}}{\pgfqpoint{4.915141in}{8.041366in}}%
\pgfpathcurveto{\pgfqpoint{4.922392in}{8.048616in}}{\pgfqpoint{4.926465in}{8.058452in}}{\pgfqpoint{4.926465in}{8.068706in}}%
\pgfpathcurveto{\pgfqpoint{4.926465in}{8.078960in}}{\pgfqpoint{4.922392in}{8.088796in}}{\pgfqpoint{4.915141in}{8.096046in}}%
\pgfpathcurveto{\pgfqpoint{4.907890in}{8.103297in}}{\pgfqpoint{4.898055in}{8.107371in}}{\pgfqpoint{4.887800in}{8.107371in}}%
\pgfpathcurveto{\pgfqpoint{4.877546in}{8.107371in}}{\pgfqpoint{4.867711in}{8.103297in}}{\pgfqpoint{4.860460in}{8.096046in}}%
\pgfpathcurveto{\pgfqpoint{4.853209in}{8.088796in}}{\pgfqpoint{4.849135in}{8.078960in}}{\pgfqpoint{4.849135in}{8.068706in}}%
\pgfpathcurveto{\pgfqpoint{4.849135in}{8.058452in}}{\pgfqpoint{4.853209in}{8.048616in}}{\pgfqpoint{4.860460in}{8.041366in}}%
\pgfpathcurveto{\pgfqpoint{4.867711in}{8.034115in}}{\pgfqpoint{4.877546in}{8.030041in}}{\pgfqpoint{4.887800in}{8.030041in}}%
\pgfpathlineto{\pgfqpoint{4.887800in}{8.030041in}}%
\pgfpathclose%
\pgfusepath{stroke}%
\end{pgfscope}%
\begin{pgfscope}%
\pgfpathrectangle{\pgfqpoint{1.109909in}{4.050000in}}{\pgfqpoint{5.400000in}{5.400000in}}%
\pgfusepath{clip}%
\pgfsetbuttcap%
\pgfsetroundjoin%
\pgfsetlinewidth{1.003750pt}%
\definecolor{currentstroke}{rgb}{0.541176,0.380392,0.658824}%
\pgfsetstrokecolor{currentstroke}%
\pgfsetstrokeopacity{0.900000}%
\pgfsetdash{}{0pt}%
\pgfpathmoveto{\pgfqpoint{4.674953in}{7.579523in}}%
\pgfpathcurveto{\pgfqpoint{4.685207in}{7.579523in}}{\pgfqpoint{4.695042in}{7.583597in}}{\pgfqpoint{4.702293in}{7.590848in}}%
\pgfpathcurveto{\pgfqpoint{4.709544in}{7.598099in}}{\pgfqpoint{4.713618in}{7.607934in}}{\pgfqpoint{4.713618in}{7.618188in}}%
\pgfpathcurveto{\pgfqpoint{4.713618in}{7.628442in}}{\pgfqpoint{4.709544in}{7.638278in}}{\pgfqpoint{4.702293in}{7.645529in}}%
\pgfpathcurveto{\pgfqpoint{4.695042in}{7.652779in}}{\pgfqpoint{4.685207in}{7.656853in}}{\pgfqpoint{4.674953in}{7.656853in}}%
\pgfpathcurveto{\pgfqpoint{4.664699in}{7.656853in}}{\pgfqpoint{4.654863in}{7.652779in}}{\pgfqpoint{4.647612in}{7.645529in}}%
\pgfpathcurveto{\pgfqpoint{4.640362in}{7.638278in}}{\pgfqpoint{4.636288in}{7.628442in}}{\pgfqpoint{4.636288in}{7.618188in}}%
\pgfpathcurveto{\pgfqpoint{4.636288in}{7.607934in}}{\pgfqpoint{4.640362in}{7.598099in}}{\pgfqpoint{4.647612in}{7.590848in}}%
\pgfpathcurveto{\pgfqpoint{4.654863in}{7.583597in}}{\pgfqpoint{4.664699in}{7.579523in}}{\pgfqpoint{4.674953in}{7.579523in}}%
\pgfpathlineto{\pgfqpoint{4.674953in}{7.579523in}}%
\pgfpathclose%
\pgfusepath{stroke}%
\end{pgfscope}%
\begin{pgfscope}%
\pgfpathrectangle{\pgfqpoint{1.109909in}{4.050000in}}{\pgfqpoint{5.400000in}{5.400000in}}%
\pgfusepath{clip}%
\pgfsetbuttcap%
\pgfsetroundjoin%
\pgfsetlinewidth{1.003750pt}%
\definecolor{currentstroke}{rgb}{0.541176,0.380392,0.658824}%
\pgfsetstrokecolor{currentstroke}%
\pgfsetstrokeopacity{0.900000}%
\pgfsetdash{}{0pt}%
\pgfpathmoveto{\pgfqpoint{4.704749in}{7.339891in}}%
\pgfpathcurveto{\pgfqpoint{4.715003in}{7.339891in}}{\pgfqpoint{4.724838in}{7.343965in}}{\pgfqpoint{4.732089in}{7.351216in}}%
\pgfpathcurveto{\pgfqpoint{4.739340in}{7.358467in}}{\pgfqpoint{4.743414in}{7.368302in}}{\pgfqpoint{4.743414in}{7.378556in}}%
\pgfpathcurveto{\pgfqpoint{4.743414in}{7.388810in}}{\pgfqpoint{4.739340in}{7.398646in}}{\pgfqpoint{4.732089in}{7.405897in}}%
\pgfpathcurveto{\pgfqpoint{4.724838in}{7.413147in}}{\pgfqpoint{4.715003in}{7.417221in}}{\pgfqpoint{4.704749in}{7.417221in}}%
\pgfpathcurveto{\pgfqpoint{4.694494in}{7.417221in}}{\pgfqpoint{4.684659in}{7.413147in}}{\pgfqpoint{4.677408in}{7.405897in}}%
\pgfpathcurveto{\pgfqpoint{4.670158in}{7.398646in}}{\pgfqpoint{4.666084in}{7.388810in}}{\pgfqpoint{4.666084in}{7.378556in}}%
\pgfpathcurveto{\pgfqpoint{4.666084in}{7.368302in}}{\pgfqpoint{4.670158in}{7.358467in}}{\pgfqpoint{4.677408in}{7.351216in}}%
\pgfpathcurveto{\pgfqpoint{4.684659in}{7.343965in}}{\pgfqpoint{4.694494in}{7.339891in}}{\pgfqpoint{4.704749in}{7.339891in}}%
\pgfpathlineto{\pgfqpoint{4.704749in}{7.339891in}}%
\pgfpathclose%
\pgfusepath{stroke}%
\end{pgfscope}%
\begin{pgfscope}%
\pgfpathrectangle{\pgfqpoint{1.109909in}{4.050000in}}{\pgfqpoint{5.400000in}{5.400000in}}%
\pgfusepath{clip}%
\pgfsetbuttcap%
\pgfsetroundjoin%
\pgfsetlinewidth{1.003750pt}%
\definecolor{currentstroke}{rgb}{0.541176,0.380392,0.658824}%
\pgfsetstrokecolor{currentstroke}%
\pgfsetstrokeopacity{0.900000}%
\pgfsetdash{}{0pt}%
\pgfpathmoveto{\pgfqpoint{3.789252in}{6.106416in}}%
\pgfpathcurveto{\pgfqpoint{3.799506in}{6.106416in}}{\pgfqpoint{3.809341in}{6.110490in}}{\pgfqpoint{3.816592in}{6.117741in}}%
\pgfpathcurveto{\pgfqpoint{3.823843in}{6.124992in}}{\pgfqpoint{3.827917in}{6.134827in}}{\pgfqpoint{3.827917in}{6.145081in}}%
\pgfpathcurveto{\pgfqpoint{3.827917in}{6.155336in}}{\pgfqpoint{3.823843in}{6.165171in}}{\pgfqpoint{3.816592in}{6.172422in}}%
\pgfpathcurveto{\pgfqpoint{3.809341in}{6.179673in}}{\pgfqpoint{3.799506in}{6.183747in}}{\pgfqpoint{3.789252in}{6.183747in}}%
\pgfpathcurveto{\pgfqpoint{3.778998in}{6.183747in}}{\pgfqpoint{3.769162in}{6.179673in}}{\pgfqpoint{3.761911in}{6.172422in}}%
\pgfpathcurveto{\pgfqpoint{3.754661in}{6.165171in}}{\pgfqpoint{3.750587in}{6.155336in}}{\pgfqpoint{3.750587in}{6.145081in}}%
\pgfpathcurveto{\pgfqpoint{3.750587in}{6.134827in}}{\pgfqpoint{3.754661in}{6.124992in}}{\pgfqpoint{3.761911in}{6.117741in}}%
\pgfpathcurveto{\pgfqpoint{3.769162in}{6.110490in}}{\pgfqpoint{3.778998in}{6.106416in}}{\pgfqpoint{3.789252in}{6.106416in}}%
\pgfpathlineto{\pgfqpoint{3.789252in}{6.106416in}}%
\pgfpathclose%
\pgfusepath{stroke}%
\end{pgfscope}%
\begin{pgfscope}%
\pgfpathrectangle{\pgfqpoint{1.109909in}{4.050000in}}{\pgfqpoint{5.400000in}{5.400000in}}%
\pgfusepath{clip}%
\pgfsetbuttcap%
\pgfsetroundjoin%
\pgfsetlinewidth{1.003750pt}%
\definecolor{currentstroke}{rgb}{0.541176,0.380392,0.658824}%
\pgfsetstrokecolor{currentstroke}%
\pgfsetstrokeopacity{0.900000}%
\pgfsetdash{}{0pt}%
\pgfpathmoveto{\pgfqpoint{4.055423in}{7.516099in}}%
\pgfpathcurveto{\pgfqpoint{4.065678in}{7.516099in}}{\pgfqpoint{4.075513in}{7.520173in}}{\pgfqpoint{4.082764in}{7.527424in}}%
\pgfpathcurveto{\pgfqpoint{4.090014in}{7.534675in}}{\pgfqpoint{4.094088in}{7.544510in}}{\pgfqpoint{4.094088in}{7.554764in}}%
\pgfpathcurveto{\pgfqpoint{4.094088in}{7.565018in}}{\pgfqpoint{4.090014in}{7.574854in}}{\pgfqpoint{4.082764in}{7.582104in}}%
\pgfpathcurveto{\pgfqpoint{4.075513in}{7.589355in}}{\pgfqpoint{4.065678in}{7.593429in}}{\pgfqpoint{4.055423in}{7.593429in}}%
\pgfpathcurveto{\pgfqpoint{4.045169in}{7.593429in}}{\pgfqpoint{4.035334in}{7.589355in}}{\pgfqpoint{4.028083in}{7.582104in}}%
\pgfpathcurveto{\pgfqpoint{4.020832in}{7.574854in}}{\pgfqpoint{4.016758in}{7.565018in}}{\pgfqpoint{4.016758in}{7.554764in}}%
\pgfpathcurveto{\pgfqpoint{4.016758in}{7.544510in}}{\pgfqpoint{4.020832in}{7.534675in}}{\pgfqpoint{4.028083in}{7.527424in}}%
\pgfpathcurveto{\pgfqpoint{4.035334in}{7.520173in}}{\pgfqpoint{4.045169in}{7.516099in}}{\pgfqpoint{4.055423in}{7.516099in}}%
\pgfpathlineto{\pgfqpoint{4.055423in}{7.516099in}}%
\pgfpathclose%
\pgfusepath{stroke}%
\end{pgfscope}%
\begin{pgfscope}%
\pgfpathrectangle{\pgfqpoint{1.109909in}{4.050000in}}{\pgfqpoint{5.400000in}{5.400000in}}%
\pgfusepath{clip}%
\pgfsetbuttcap%
\pgfsetroundjoin%
\pgfsetlinewidth{1.003750pt}%
\definecolor{currentstroke}{rgb}{0.541176,0.380392,0.658824}%
\pgfsetstrokecolor{currentstroke}%
\pgfsetstrokeopacity{0.900000}%
\pgfsetdash{}{0pt}%
\pgfpathmoveto{\pgfqpoint{4.580418in}{7.419579in}}%
\pgfpathcurveto{\pgfqpoint{4.590672in}{7.419579in}}{\pgfqpoint{4.600507in}{7.423653in}}{\pgfqpoint{4.607758in}{7.430904in}}%
\pgfpathcurveto{\pgfqpoint{4.615009in}{7.438155in}}{\pgfqpoint{4.619083in}{7.447990in}}{\pgfqpoint{4.619083in}{7.458244in}}%
\pgfpathcurveto{\pgfqpoint{4.619083in}{7.468498in}}{\pgfqpoint{4.615009in}{7.478334in}}{\pgfqpoint{4.607758in}{7.485585in}}%
\pgfpathcurveto{\pgfqpoint{4.600507in}{7.492835in}}{\pgfqpoint{4.590672in}{7.496909in}}{\pgfqpoint{4.580418in}{7.496909in}}%
\pgfpathcurveto{\pgfqpoint{4.570164in}{7.496909in}}{\pgfqpoint{4.560328in}{7.492835in}}{\pgfqpoint{4.553078in}{7.485585in}}%
\pgfpathcurveto{\pgfqpoint{4.545827in}{7.478334in}}{\pgfqpoint{4.541753in}{7.468498in}}{\pgfqpoint{4.541753in}{7.458244in}}%
\pgfpathcurveto{\pgfqpoint{4.541753in}{7.447990in}}{\pgfqpoint{4.545827in}{7.438155in}}{\pgfqpoint{4.553078in}{7.430904in}}%
\pgfpathcurveto{\pgfqpoint{4.560328in}{7.423653in}}{\pgfqpoint{4.570164in}{7.419579in}}{\pgfqpoint{4.580418in}{7.419579in}}%
\pgfpathlineto{\pgfqpoint{4.580418in}{7.419579in}}%
\pgfpathclose%
\pgfusepath{stroke}%
\end{pgfscope}%
\begin{pgfscope}%
\pgfpathrectangle{\pgfqpoint{1.109909in}{4.050000in}}{\pgfqpoint{5.400000in}{5.400000in}}%
\pgfusepath{clip}%
\pgfsetbuttcap%
\pgfsetroundjoin%
\pgfsetlinewidth{1.003750pt}%
\definecolor{currentstroke}{rgb}{0.541176,0.380392,0.658824}%
\pgfsetstrokecolor{currentstroke}%
\pgfsetstrokeopacity{0.900000}%
\pgfsetdash{}{0pt}%
\pgfpathmoveto{\pgfqpoint{3.750245in}{6.537397in}}%
\pgfpathcurveto{\pgfqpoint{3.760499in}{6.537397in}}{\pgfqpoint{3.770335in}{6.541471in}}{\pgfqpoint{3.777585in}{6.548722in}}%
\pgfpathcurveto{\pgfqpoint{3.784836in}{6.555973in}}{\pgfqpoint{3.788910in}{6.565808in}}{\pgfqpoint{3.788910in}{6.576062in}}%
\pgfpathcurveto{\pgfqpoint{3.788910in}{6.586316in}}{\pgfqpoint{3.784836in}{6.596152in}}{\pgfqpoint{3.777585in}{6.603403in}}%
\pgfpathcurveto{\pgfqpoint{3.770335in}{6.610653in}}{\pgfqpoint{3.760499in}{6.614727in}}{\pgfqpoint{3.750245in}{6.614727in}}%
\pgfpathcurveto{\pgfqpoint{3.739991in}{6.614727in}}{\pgfqpoint{3.730156in}{6.610653in}}{\pgfqpoint{3.722905in}{6.603403in}}%
\pgfpathcurveto{\pgfqpoint{3.715654in}{6.596152in}}{\pgfqpoint{3.711580in}{6.586316in}}{\pgfqpoint{3.711580in}{6.576062in}}%
\pgfpathcurveto{\pgfqpoint{3.711580in}{6.565808in}}{\pgfqpoint{3.715654in}{6.555973in}}{\pgfqpoint{3.722905in}{6.548722in}}%
\pgfpathcurveto{\pgfqpoint{3.730156in}{6.541471in}}{\pgfqpoint{3.739991in}{6.537397in}}{\pgfqpoint{3.750245in}{6.537397in}}%
\pgfpathlineto{\pgfqpoint{3.750245in}{6.537397in}}%
\pgfpathclose%
\pgfusepath{stroke}%
\end{pgfscope}%
\begin{pgfscope}%
\pgfpathrectangle{\pgfqpoint{1.109909in}{4.050000in}}{\pgfqpoint{5.400000in}{5.400000in}}%
\pgfusepath{clip}%
\pgfsetbuttcap%
\pgfsetroundjoin%
\pgfsetlinewidth{1.003750pt}%
\definecolor{currentstroke}{rgb}{0.541176,0.380392,0.658824}%
\pgfsetstrokecolor{currentstroke}%
\pgfsetstrokeopacity{0.900000}%
\pgfsetdash{}{0pt}%
\pgfpathmoveto{\pgfqpoint{5.096846in}{7.666997in}}%
\pgfpathcurveto{\pgfqpoint{5.107100in}{7.666997in}}{\pgfqpoint{5.116935in}{7.671071in}}{\pgfqpoint{5.124186in}{7.678322in}}%
\pgfpathcurveto{\pgfqpoint{5.131437in}{7.685573in}}{\pgfqpoint{5.135511in}{7.695408in}}{\pgfqpoint{5.135511in}{7.705662in}}%
\pgfpathcurveto{\pgfqpoint{5.135511in}{7.715916in}}{\pgfqpoint{5.131437in}{7.725752in}}{\pgfqpoint{5.124186in}{7.733003in}}%
\pgfpathcurveto{\pgfqpoint{5.116935in}{7.740253in}}{\pgfqpoint{5.107100in}{7.744327in}}{\pgfqpoint{5.096846in}{7.744327in}}%
\pgfpathcurveto{\pgfqpoint{5.086592in}{7.744327in}}{\pgfqpoint{5.076756in}{7.740253in}}{\pgfqpoint{5.069505in}{7.733003in}}%
\pgfpathcurveto{\pgfqpoint{5.062255in}{7.725752in}}{\pgfqpoint{5.058181in}{7.715916in}}{\pgfqpoint{5.058181in}{7.705662in}}%
\pgfpathcurveto{\pgfqpoint{5.058181in}{7.695408in}}{\pgfqpoint{5.062255in}{7.685573in}}{\pgfqpoint{5.069505in}{7.678322in}}%
\pgfpathcurveto{\pgfqpoint{5.076756in}{7.671071in}}{\pgfqpoint{5.086592in}{7.666997in}}{\pgfqpoint{5.096846in}{7.666997in}}%
\pgfpathlineto{\pgfqpoint{5.096846in}{7.666997in}}%
\pgfpathclose%
\pgfusepath{stroke}%
\end{pgfscope}%
\begin{pgfscope}%
\pgfpathrectangle{\pgfqpoint{1.109909in}{4.050000in}}{\pgfqpoint{5.400000in}{5.400000in}}%
\pgfusepath{clip}%
\pgfsetbuttcap%
\pgfsetroundjoin%
\pgfsetlinewidth{1.003750pt}%
\definecolor{currentstroke}{rgb}{0.541176,0.380392,0.658824}%
\pgfsetstrokecolor{currentstroke}%
\pgfsetstrokeopacity{0.900000}%
\pgfsetdash{}{0pt}%
\pgfpathmoveto{\pgfqpoint{4.828288in}{7.575316in}}%
\pgfpathcurveto{\pgfqpoint{4.838542in}{7.575316in}}{\pgfqpoint{4.848378in}{7.579390in}}{\pgfqpoint{4.855628in}{7.586641in}}%
\pgfpathcurveto{\pgfqpoint{4.862879in}{7.593892in}}{\pgfqpoint{4.866953in}{7.603727in}}{\pgfqpoint{4.866953in}{7.613981in}}%
\pgfpathcurveto{\pgfqpoint{4.866953in}{7.624235in}}{\pgfqpoint{4.862879in}{7.634071in}}{\pgfqpoint{4.855628in}{7.641321in}}%
\pgfpathcurveto{\pgfqpoint{4.848378in}{7.648572in}}{\pgfqpoint{4.838542in}{7.652646in}}{\pgfqpoint{4.828288in}{7.652646in}}%
\pgfpathcurveto{\pgfqpoint{4.818034in}{7.652646in}}{\pgfqpoint{4.808198in}{7.648572in}}{\pgfqpoint{4.800948in}{7.641321in}}%
\pgfpathcurveto{\pgfqpoint{4.793697in}{7.634071in}}{\pgfqpoint{4.789623in}{7.624235in}}{\pgfqpoint{4.789623in}{7.613981in}}%
\pgfpathcurveto{\pgfqpoint{4.789623in}{7.603727in}}{\pgfqpoint{4.793697in}{7.593892in}}{\pgfqpoint{4.800948in}{7.586641in}}%
\pgfpathcurveto{\pgfqpoint{4.808198in}{7.579390in}}{\pgfqpoint{4.818034in}{7.575316in}}{\pgfqpoint{4.828288in}{7.575316in}}%
\pgfpathlineto{\pgfqpoint{4.828288in}{7.575316in}}%
\pgfpathclose%
\pgfusepath{stroke}%
\end{pgfscope}%
\begin{pgfscope}%
\pgfpathrectangle{\pgfqpoint{1.109909in}{4.050000in}}{\pgfqpoint{5.400000in}{5.400000in}}%
\pgfusepath{clip}%
\pgfsetbuttcap%
\pgfsetroundjoin%
\pgfsetlinewidth{1.003750pt}%
\definecolor{currentstroke}{rgb}{0.541176,0.380392,0.658824}%
\pgfsetstrokecolor{currentstroke}%
\pgfsetstrokeopacity{0.900000}%
\pgfsetdash{}{0pt}%
\pgfpathmoveto{\pgfqpoint{5.211822in}{6.816682in}}%
\pgfpathcurveto{\pgfqpoint{5.222076in}{6.816682in}}{\pgfqpoint{5.231912in}{6.820756in}}{\pgfqpoint{5.239163in}{6.828007in}}%
\pgfpathcurveto{\pgfqpoint{5.246413in}{6.835258in}}{\pgfqpoint{5.250487in}{6.845093in}}{\pgfqpoint{5.250487in}{6.855347in}}%
\pgfpathcurveto{\pgfqpoint{5.250487in}{6.865601in}}{\pgfqpoint{5.246413in}{6.875437in}}{\pgfqpoint{5.239163in}{6.882687in}}%
\pgfpathcurveto{\pgfqpoint{5.231912in}{6.889938in}}{\pgfqpoint{5.222076in}{6.894012in}}{\pgfqpoint{5.211822in}{6.894012in}}%
\pgfpathcurveto{\pgfqpoint{5.201568in}{6.894012in}}{\pgfqpoint{5.191733in}{6.889938in}}{\pgfqpoint{5.184482in}{6.882687in}}%
\pgfpathcurveto{\pgfqpoint{5.177231in}{6.875437in}}{\pgfqpoint{5.173157in}{6.865601in}}{\pgfqpoint{5.173157in}{6.855347in}}%
\pgfpathcurveto{\pgfqpoint{5.173157in}{6.845093in}}{\pgfqpoint{5.177231in}{6.835258in}}{\pgfqpoint{5.184482in}{6.828007in}}%
\pgfpathcurveto{\pgfqpoint{5.191733in}{6.820756in}}{\pgfqpoint{5.201568in}{6.816682in}}{\pgfqpoint{5.211822in}{6.816682in}}%
\pgfpathlineto{\pgfqpoint{5.211822in}{6.816682in}}%
\pgfpathclose%
\pgfusepath{stroke}%
\end{pgfscope}%
\begin{pgfscope}%
\pgfpathrectangle{\pgfqpoint{1.109909in}{4.050000in}}{\pgfqpoint{5.400000in}{5.400000in}}%
\pgfusepath{clip}%
\pgfsetbuttcap%
\pgfsetroundjoin%
\pgfsetlinewidth{1.003750pt}%
\definecolor{currentstroke}{rgb}{0.541176,0.380392,0.658824}%
\pgfsetstrokecolor{currentstroke}%
\pgfsetstrokeopacity{0.900000}%
\pgfsetdash{}{0pt}%
\pgfpathmoveto{\pgfqpoint{5.289251in}{7.167715in}}%
\pgfpathcurveto{\pgfqpoint{5.299505in}{7.167715in}}{\pgfqpoint{5.309340in}{7.171789in}}{\pgfqpoint{5.316591in}{7.179040in}}%
\pgfpathcurveto{\pgfqpoint{5.323842in}{7.186291in}}{\pgfqpoint{5.327916in}{7.196126in}}{\pgfqpoint{5.327916in}{7.206380in}}%
\pgfpathcurveto{\pgfqpoint{5.327916in}{7.216634in}}{\pgfqpoint{5.323842in}{7.226470in}}{\pgfqpoint{5.316591in}{7.233721in}}%
\pgfpathcurveto{\pgfqpoint{5.309340in}{7.240971in}}{\pgfqpoint{5.299505in}{7.245045in}}{\pgfqpoint{5.289251in}{7.245045in}}%
\pgfpathcurveto{\pgfqpoint{5.278996in}{7.245045in}}{\pgfqpoint{5.269161in}{7.240971in}}{\pgfqpoint{5.261910in}{7.233721in}}%
\pgfpathcurveto{\pgfqpoint{5.254659in}{7.226470in}}{\pgfqpoint{5.250585in}{7.216634in}}{\pgfqpoint{5.250585in}{7.206380in}}%
\pgfpathcurveto{\pgfqpoint{5.250585in}{7.196126in}}{\pgfqpoint{5.254659in}{7.186291in}}{\pgfqpoint{5.261910in}{7.179040in}}%
\pgfpathcurveto{\pgfqpoint{5.269161in}{7.171789in}}{\pgfqpoint{5.278996in}{7.167715in}}{\pgfqpoint{5.289251in}{7.167715in}}%
\pgfpathlineto{\pgfqpoint{5.289251in}{7.167715in}}%
\pgfpathclose%
\pgfusepath{stroke}%
\end{pgfscope}%
\begin{pgfscope}%
\pgfpathrectangle{\pgfqpoint{1.109909in}{4.050000in}}{\pgfqpoint{5.400000in}{5.400000in}}%
\pgfusepath{clip}%
\pgfsetbuttcap%
\pgfsetroundjoin%
\pgfsetlinewidth{1.003750pt}%
\definecolor{currentstroke}{rgb}{0.541176,0.380392,0.658824}%
\pgfsetstrokecolor{currentstroke}%
\pgfsetstrokeopacity{0.900000}%
\pgfsetdash{}{0pt}%
\pgfpathmoveto{\pgfqpoint{5.446328in}{9.401911in}}%
\pgfpathcurveto{\pgfqpoint{5.456582in}{9.401911in}}{\pgfqpoint{5.466417in}{9.405985in}}{\pgfqpoint{5.473668in}{9.413236in}}%
\pgfpathcurveto{\pgfqpoint{5.480919in}{9.420487in}}{\pgfqpoint{5.484993in}{9.430322in}}{\pgfqpoint{5.484993in}{9.440576in}}%
\pgfpathcurveto{\pgfqpoint{5.484993in}{9.450831in}}{\pgfqpoint{5.480919in}{9.460666in}}{\pgfqpoint{5.473668in}{9.467917in}}%
\pgfpathcurveto{\pgfqpoint{5.466417in}{9.475168in}}{\pgfqpoint{5.456582in}{9.479242in}}{\pgfqpoint{5.446328in}{9.479242in}}%
\pgfpathcurveto{\pgfqpoint{5.436073in}{9.479242in}}{\pgfqpoint{5.426238in}{9.475168in}}{\pgfqpoint{5.418987in}{9.467917in}}%
\pgfpathcurveto{\pgfqpoint{5.411737in}{9.460666in}}{\pgfqpoint{5.407663in}{9.450831in}}{\pgfqpoint{5.407663in}{9.440576in}}%
\pgfpathcurveto{\pgfqpoint{5.407663in}{9.430322in}}{\pgfqpoint{5.411737in}{9.420487in}}{\pgfqpoint{5.418987in}{9.413236in}}%
\pgfpathcurveto{\pgfqpoint{5.426238in}{9.405985in}}{\pgfqpoint{5.436073in}{9.401911in}}{\pgfqpoint{5.446328in}{9.401911in}}%
\pgfpathlineto{\pgfqpoint{5.446328in}{9.401911in}}%
\pgfpathclose%
\pgfusepath{stroke}%
\end{pgfscope}%
\begin{pgfscope}%
\pgfpathrectangle{\pgfqpoint{1.109909in}{4.050000in}}{\pgfqpoint{5.400000in}{5.400000in}}%
\pgfusepath{clip}%
\pgfsetbuttcap%
\pgfsetroundjoin%
\pgfsetlinewidth{1.003750pt}%
\definecolor{currentstroke}{rgb}{0.541176,0.380392,0.658824}%
\pgfsetstrokecolor{currentstroke}%
\pgfsetstrokeopacity{0.900000}%
\pgfsetdash{}{0pt}%
\pgfpathmoveto{\pgfqpoint{3.993260in}{6.895074in}}%
\pgfpathcurveto{\pgfqpoint{4.003514in}{6.895074in}}{\pgfqpoint{4.013349in}{6.899147in}}{\pgfqpoint{4.020600in}{6.906398in}}%
\pgfpathcurveto{\pgfqpoint{4.027851in}{6.913649in}}{\pgfqpoint{4.031925in}{6.923484in}}{\pgfqpoint{4.031925in}{6.933739in}}%
\pgfpathcurveto{\pgfqpoint{4.031925in}{6.943993in}}{\pgfqpoint{4.027851in}{6.953828in}}{\pgfqpoint{4.020600in}{6.961079in}}%
\pgfpathcurveto{\pgfqpoint{4.013349in}{6.968330in}}{\pgfqpoint{4.003514in}{6.972404in}}{\pgfqpoint{3.993260in}{6.972404in}}%
\pgfpathcurveto{\pgfqpoint{3.983006in}{6.972404in}}{\pgfqpoint{3.973170in}{6.968330in}}{\pgfqpoint{3.965919in}{6.961079in}}%
\pgfpathcurveto{\pgfqpoint{3.958669in}{6.953828in}}{\pgfqpoint{3.954595in}{6.943993in}}{\pgfqpoint{3.954595in}{6.933739in}}%
\pgfpathcurveto{\pgfqpoint{3.954595in}{6.923484in}}{\pgfqpoint{3.958669in}{6.913649in}}{\pgfqpoint{3.965919in}{6.906398in}}%
\pgfpathcurveto{\pgfqpoint{3.973170in}{6.899147in}}{\pgfqpoint{3.983006in}{6.895074in}}{\pgfqpoint{3.993260in}{6.895074in}}%
\pgfpathlineto{\pgfqpoint{3.993260in}{6.895074in}}%
\pgfpathclose%
\pgfusepath{stroke}%
\end{pgfscope}%
\begin{pgfscope}%
\pgfpathrectangle{\pgfqpoint{1.109909in}{4.050000in}}{\pgfqpoint{5.400000in}{5.400000in}}%
\pgfusepath{clip}%
\pgfsetbuttcap%
\pgfsetroundjoin%
\pgfsetlinewidth{1.003750pt}%
\definecolor{currentstroke}{rgb}{0.541176,0.380392,0.658824}%
\pgfsetstrokecolor{currentstroke}%
\pgfsetstrokeopacity{0.900000}%
\pgfsetdash{}{0pt}%
\pgfpathmoveto{\pgfqpoint{3.950792in}{8.046169in}}%
\pgfpathcurveto{\pgfqpoint{3.961047in}{8.046169in}}{\pgfqpoint{3.970882in}{8.050243in}}{\pgfqpoint{3.978133in}{8.057494in}}%
\pgfpathcurveto{\pgfqpoint{3.985384in}{8.064744in}}{\pgfqpoint{3.989458in}{8.074580in}}{\pgfqpoint{3.989458in}{8.084834in}}%
\pgfpathcurveto{\pgfqpoint{3.989458in}{8.095088in}}{\pgfqpoint{3.985384in}{8.104924in}}{\pgfqpoint{3.978133in}{8.112174in}}%
\pgfpathcurveto{\pgfqpoint{3.970882in}{8.119425in}}{\pgfqpoint{3.961047in}{8.123499in}}{\pgfqpoint{3.950792in}{8.123499in}}%
\pgfpathcurveto{\pgfqpoint{3.940538in}{8.123499in}}{\pgfqpoint{3.930703in}{8.119425in}}{\pgfqpoint{3.923452in}{8.112174in}}%
\pgfpathcurveto{\pgfqpoint{3.916201in}{8.104924in}}{\pgfqpoint{3.912127in}{8.095088in}}{\pgfqpoint{3.912127in}{8.084834in}}%
\pgfpathcurveto{\pgfqpoint{3.912127in}{8.074580in}}{\pgfqpoint{3.916201in}{8.064744in}}{\pgfqpoint{3.923452in}{8.057494in}}%
\pgfpathcurveto{\pgfqpoint{3.930703in}{8.050243in}}{\pgfqpoint{3.940538in}{8.046169in}}{\pgfqpoint{3.950792in}{8.046169in}}%
\pgfpathlineto{\pgfqpoint{3.950792in}{8.046169in}}%
\pgfpathclose%
\pgfusepath{stroke}%
\end{pgfscope}%
\begin{pgfscope}%
\pgfpathrectangle{\pgfqpoint{1.109909in}{4.050000in}}{\pgfqpoint{5.400000in}{5.400000in}}%
\pgfusepath{clip}%
\pgfsetbuttcap%
\pgfsetroundjoin%
\pgfsetlinewidth{1.003750pt}%
\definecolor{currentstroke}{rgb}{0.541176,0.380392,0.658824}%
\pgfsetstrokecolor{currentstroke}%
\pgfsetstrokeopacity{0.900000}%
\pgfsetdash{}{0pt}%
\pgfpathmoveto{\pgfqpoint{4.717950in}{7.464530in}}%
\pgfpathcurveto{\pgfqpoint{4.728204in}{7.464530in}}{\pgfqpoint{4.738039in}{7.468604in}}{\pgfqpoint{4.745290in}{7.475855in}}%
\pgfpathcurveto{\pgfqpoint{4.752541in}{7.483105in}}{\pgfqpoint{4.756615in}{7.492941in}}{\pgfqpoint{4.756615in}{7.503195in}}%
\pgfpathcurveto{\pgfqpoint{4.756615in}{7.513449in}}{\pgfqpoint{4.752541in}{7.523285in}}{\pgfqpoint{4.745290in}{7.530535in}}%
\pgfpathcurveto{\pgfqpoint{4.738039in}{7.537786in}}{\pgfqpoint{4.728204in}{7.541860in}}{\pgfqpoint{4.717950in}{7.541860in}}%
\pgfpathcurveto{\pgfqpoint{4.707695in}{7.541860in}}{\pgfqpoint{4.697860in}{7.537786in}}{\pgfqpoint{4.690609in}{7.530535in}}%
\pgfpathcurveto{\pgfqpoint{4.683358in}{7.523285in}}{\pgfqpoint{4.679284in}{7.513449in}}{\pgfqpoint{4.679284in}{7.503195in}}%
\pgfpathcurveto{\pgfqpoint{4.679284in}{7.492941in}}{\pgfqpoint{4.683358in}{7.483105in}}{\pgfqpoint{4.690609in}{7.475855in}}%
\pgfpathcurveto{\pgfqpoint{4.697860in}{7.468604in}}{\pgfqpoint{4.707695in}{7.464530in}}{\pgfqpoint{4.717950in}{7.464530in}}%
\pgfpathlineto{\pgfqpoint{4.717950in}{7.464530in}}%
\pgfpathclose%
\pgfusepath{stroke}%
\end{pgfscope}%
\begin{pgfscope}%
\pgfpathrectangle{\pgfqpoint{1.109909in}{4.050000in}}{\pgfqpoint{5.400000in}{5.400000in}}%
\pgfusepath{clip}%
\pgfsetbuttcap%
\pgfsetroundjoin%
\pgfsetlinewidth{1.003750pt}%
\definecolor{currentstroke}{rgb}{0.541176,0.380392,0.658824}%
\pgfsetstrokecolor{currentstroke}%
\pgfsetstrokeopacity{0.900000}%
\pgfsetdash{}{0pt}%
\pgfpathmoveto{\pgfqpoint{4.545796in}{8.160866in}}%
\pgfpathcurveto{\pgfqpoint{4.556050in}{8.160866in}}{\pgfqpoint{4.565885in}{8.164940in}}{\pgfqpoint{4.573136in}{8.172191in}}%
\pgfpathcurveto{\pgfqpoint{4.580387in}{8.179441in}}{\pgfqpoint{4.584461in}{8.189277in}}{\pgfqpoint{4.584461in}{8.199531in}}%
\pgfpathcurveto{\pgfqpoint{4.584461in}{8.209785in}}{\pgfqpoint{4.580387in}{8.219620in}}{\pgfqpoint{4.573136in}{8.226871in}}%
\pgfpathcurveto{\pgfqpoint{4.565885in}{8.234122in}}{\pgfqpoint{4.556050in}{8.238196in}}{\pgfqpoint{4.545796in}{8.238196in}}%
\pgfpathcurveto{\pgfqpoint{4.535542in}{8.238196in}}{\pgfqpoint{4.525706in}{8.234122in}}{\pgfqpoint{4.518455in}{8.226871in}}%
\pgfpathcurveto{\pgfqpoint{4.511205in}{8.219620in}}{\pgfqpoint{4.507131in}{8.209785in}}{\pgfqpoint{4.507131in}{8.199531in}}%
\pgfpathcurveto{\pgfqpoint{4.507131in}{8.189277in}}{\pgfqpoint{4.511205in}{8.179441in}}{\pgfqpoint{4.518455in}{8.172191in}}%
\pgfpathcurveto{\pgfqpoint{4.525706in}{8.164940in}}{\pgfqpoint{4.535542in}{8.160866in}}{\pgfqpoint{4.545796in}{8.160866in}}%
\pgfpathlineto{\pgfqpoint{4.545796in}{8.160866in}}%
\pgfpathclose%
\pgfusepath{stroke}%
\end{pgfscope}%
\begin{pgfscope}%
\pgfpathrectangle{\pgfqpoint{1.109909in}{4.050000in}}{\pgfqpoint{5.400000in}{5.400000in}}%
\pgfusepath{clip}%
\pgfsetbuttcap%
\pgfsetroundjoin%
\pgfsetlinewidth{1.003750pt}%
\definecolor{currentstroke}{rgb}{0.541176,0.380392,0.658824}%
\pgfsetstrokecolor{currentstroke}%
\pgfsetstrokeopacity{0.900000}%
\pgfsetdash{}{0pt}%
\pgfpathmoveto{\pgfqpoint{4.801160in}{8.583431in}}%
\pgfpathcurveto{\pgfqpoint{4.811414in}{8.583431in}}{\pgfqpoint{4.821250in}{8.587505in}}{\pgfqpoint{4.828500in}{8.594756in}}%
\pgfpathcurveto{\pgfqpoint{4.835751in}{8.602007in}}{\pgfqpoint{4.839825in}{8.611842in}}{\pgfqpoint{4.839825in}{8.622097in}}%
\pgfpathcurveto{\pgfqpoint{4.839825in}{8.632351in}}{\pgfqpoint{4.835751in}{8.642186in}}{\pgfqpoint{4.828500in}{8.649437in}}%
\pgfpathcurveto{\pgfqpoint{4.821250in}{8.656688in}}{\pgfqpoint{4.811414in}{8.660762in}}{\pgfqpoint{4.801160in}{8.660762in}}%
\pgfpathcurveto{\pgfqpoint{4.790906in}{8.660762in}}{\pgfqpoint{4.781071in}{8.656688in}}{\pgfqpoint{4.773820in}{8.649437in}}%
\pgfpathcurveto{\pgfqpoint{4.766569in}{8.642186in}}{\pgfqpoint{4.762495in}{8.632351in}}{\pgfqpoint{4.762495in}{8.622097in}}%
\pgfpathcurveto{\pgfqpoint{4.762495in}{8.611842in}}{\pgfqpoint{4.766569in}{8.602007in}}{\pgfqpoint{4.773820in}{8.594756in}}%
\pgfpathcurveto{\pgfqpoint{4.781071in}{8.587505in}}{\pgfqpoint{4.790906in}{8.583431in}}{\pgfqpoint{4.801160in}{8.583431in}}%
\pgfpathlineto{\pgfqpoint{4.801160in}{8.583431in}}%
\pgfpathclose%
\pgfusepath{stroke}%
\end{pgfscope}%
\begin{pgfscope}%
\pgfpathrectangle{\pgfqpoint{1.109909in}{4.050000in}}{\pgfqpoint{5.400000in}{5.400000in}}%
\pgfusepath{clip}%
\pgfsetbuttcap%
\pgfsetroundjoin%
\pgfsetlinewidth{1.003750pt}%
\definecolor{currentstroke}{rgb}{0.541176,0.380392,0.658824}%
\pgfsetstrokecolor{currentstroke}%
\pgfsetstrokeopacity{0.900000}%
\pgfsetdash{}{0pt}%
\pgfpathmoveto{\pgfqpoint{4.557042in}{8.439796in}}%
\pgfpathcurveto{\pgfqpoint{4.567296in}{8.439796in}}{\pgfqpoint{4.577132in}{8.443870in}}{\pgfqpoint{4.584383in}{8.451121in}}%
\pgfpathcurveto{\pgfqpoint{4.591633in}{8.458372in}}{\pgfqpoint{4.595707in}{8.468207in}}{\pgfqpoint{4.595707in}{8.478461in}}%
\pgfpathcurveto{\pgfqpoint{4.595707in}{8.488715in}}{\pgfqpoint{4.591633in}{8.498551in}}{\pgfqpoint{4.584383in}{8.505801in}}%
\pgfpathcurveto{\pgfqpoint{4.577132in}{8.513052in}}{\pgfqpoint{4.567296in}{8.517126in}}{\pgfqpoint{4.557042in}{8.517126in}}%
\pgfpathcurveto{\pgfqpoint{4.546788in}{8.517126in}}{\pgfqpoint{4.536953in}{8.513052in}}{\pgfqpoint{4.529702in}{8.505801in}}%
\pgfpathcurveto{\pgfqpoint{4.522451in}{8.498551in}}{\pgfqpoint{4.518377in}{8.488715in}}{\pgfqpoint{4.518377in}{8.478461in}}%
\pgfpathcurveto{\pgfqpoint{4.518377in}{8.468207in}}{\pgfqpoint{4.522451in}{8.458372in}}{\pgfqpoint{4.529702in}{8.451121in}}%
\pgfpathcurveto{\pgfqpoint{4.536953in}{8.443870in}}{\pgfqpoint{4.546788in}{8.439796in}}{\pgfqpoint{4.557042in}{8.439796in}}%
\pgfpathlineto{\pgfqpoint{4.557042in}{8.439796in}}%
\pgfpathclose%
\pgfusepath{stroke}%
\end{pgfscope}%
\begin{pgfscope}%
\pgfpathrectangle{\pgfqpoint{1.109909in}{4.050000in}}{\pgfqpoint{5.400000in}{5.400000in}}%
\pgfusepath{clip}%
\pgfsetbuttcap%
\pgfsetroundjoin%
\pgfsetlinewidth{1.003750pt}%
\definecolor{currentstroke}{rgb}{0.541176,0.380392,0.658824}%
\pgfsetstrokecolor{currentstroke}%
\pgfsetstrokeopacity{0.900000}%
\pgfsetdash{}{0pt}%
\pgfpathmoveto{\pgfqpoint{4.980286in}{8.347834in}}%
\pgfpathcurveto{\pgfqpoint{4.990540in}{8.347834in}}{\pgfqpoint{5.000375in}{8.351908in}}{\pgfqpoint{5.007626in}{8.359159in}}%
\pgfpathcurveto{\pgfqpoint{5.014877in}{8.366409in}}{\pgfqpoint{5.018951in}{8.376245in}}{\pgfqpoint{5.018951in}{8.386499in}}%
\pgfpathcurveto{\pgfqpoint{5.018951in}{8.396753in}}{\pgfqpoint{5.014877in}{8.406589in}}{\pgfqpoint{5.007626in}{8.413839in}}%
\pgfpathcurveto{\pgfqpoint{5.000375in}{8.421090in}}{\pgfqpoint{4.990540in}{8.425164in}}{\pgfqpoint{4.980286in}{8.425164in}}%
\pgfpathcurveto{\pgfqpoint{4.970032in}{8.425164in}}{\pgfqpoint{4.960196in}{8.421090in}}{\pgfqpoint{4.952945in}{8.413839in}}%
\pgfpathcurveto{\pgfqpoint{4.945695in}{8.406589in}}{\pgfqpoint{4.941621in}{8.396753in}}{\pgfqpoint{4.941621in}{8.386499in}}%
\pgfpathcurveto{\pgfqpoint{4.941621in}{8.376245in}}{\pgfqpoint{4.945695in}{8.366409in}}{\pgfqpoint{4.952945in}{8.359159in}}%
\pgfpathcurveto{\pgfqpoint{4.960196in}{8.351908in}}{\pgfqpoint{4.970032in}{8.347834in}}{\pgfqpoint{4.980286in}{8.347834in}}%
\pgfpathlineto{\pgfqpoint{4.980286in}{8.347834in}}%
\pgfpathclose%
\pgfusepath{stroke}%
\end{pgfscope}%
\begin{pgfscope}%
\pgfpathrectangle{\pgfqpoint{1.109909in}{4.050000in}}{\pgfqpoint{5.400000in}{5.400000in}}%
\pgfusepath{clip}%
\pgfsetbuttcap%
\pgfsetroundjoin%
\pgfsetlinewidth{1.003750pt}%
\definecolor{currentstroke}{rgb}{0.541176,0.380392,0.658824}%
\pgfsetstrokecolor{currentstroke}%
\pgfsetstrokeopacity{0.900000}%
\pgfsetdash{}{0pt}%
\pgfpathmoveto{\pgfqpoint{2.790911in}{6.788590in}}%
\pgfpathcurveto{\pgfqpoint{2.801165in}{6.788590in}}{\pgfqpoint{2.811001in}{6.792664in}}{\pgfqpoint{2.818251in}{6.799915in}}%
\pgfpathcurveto{\pgfqpoint{2.825502in}{6.807166in}}{\pgfqpoint{2.829576in}{6.817001in}}{\pgfqpoint{2.829576in}{6.827256in}}%
\pgfpathcurveto{\pgfqpoint{2.829576in}{6.837510in}}{\pgfqpoint{2.825502in}{6.847345in}}{\pgfqpoint{2.818251in}{6.854596in}}%
\pgfpathcurveto{\pgfqpoint{2.811001in}{6.861847in}}{\pgfqpoint{2.801165in}{6.865921in}}{\pgfqpoint{2.790911in}{6.865921in}}%
\pgfpathcurveto{\pgfqpoint{2.780657in}{6.865921in}}{\pgfqpoint{2.770822in}{6.861847in}}{\pgfqpoint{2.763571in}{6.854596in}}%
\pgfpathcurveto{\pgfqpoint{2.756320in}{6.847345in}}{\pgfqpoint{2.752246in}{6.837510in}}{\pgfqpoint{2.752246in}{6.827256in}}%
\pgfpathcurveto{\pgfqpoint{2.752246in}{6.817001in}}{\pgfqpoint{2.756320in}{6.807166in}}{\pgfqpoint{2.763571in}{6.799915in}}%
\pgfpathcurveto{\pgfqpoint{2.770822in}{6.792664in}}{\pgfqpoint{2.780657in}{6.788590in}}{\pgfqpoint{2.790911in}{6.788590in}}%
\pgfpathlineto{\pgfqpoint{2.790911in}{6.788590in}}%
\pgfpathclose%
\pgfusepath{stroke}%
\end{pgfscope}%
\begin{pgfscope}%
\pgfpathrectangle{\pgfqpoint{1.109909in}{4.050000in}}{\pgfqpoint{5.400000in}{5.400000in}}%
\pgfusepath{clip}%
\pgfsetbuttcap%
\pgfsetroundjoin%
\pgfsetlinewidth{1.003750pt}%
\definecolor{currentstroke}{rgb}{0.541176,0.380392,0.658824}%
\pgfsetstrokecolor{currentstroke}%
\pgfsetstrokeopacity{0.900000}%
\pgfsetdash{}{0pt}%
\pgfpathmoveto{\pgfqpoint{4.624943in}{7.364497in}}%
\pgfpathcurveto{\pgfqpoint{4.635197in}{7.364497in}}{\pgfqpoint{4.645032in}{7.368571in}}{\pgfqpoint{4.652283in}{7.375822in}}%
\pgfpathcurveto{\pgfqpoint{4.659534in}{7.383073in}}{\pgfqpoint{4.663608in}{7.392908in}}{\pgfqpoint{4.663608in}{7.403163in}}%
\pgfpathcurveto{\pgfqpoint{4.663608in}{7.413417in}}{\pgfqpoint{4.659534in}{7.423252in}}{\pgfqpoint{4.652283in}{7.430503in}}%
\pgfpathcurveto{\pgfqpoint{4.645032in}{7.437754in}}{\pgfqpoint{4.635197in}{7.441828in}}{\pgfqpoint{4.624943in}{7.441828in}}%
\pgfpathcurveto{\pgfqpoint{4.614689in}{7.441828in}}{\pgfqpoint{4.604853in}{7.437754in}}{\pgfqpoint{4.597602in}{7.430503in}}%
\pgfpathcurveto{\pgfqpoint{4.590352in}{7.423252in}}{\pgfqpoint{4.586278in}{7.413417in}}{\pgfqpoint{4.586278in}{7.403163in}}%
\pgfpathcurveto{\pgfqpoint{4.586278in}{7.392908in}}{\pgfqpoint{4.590352in}{7.383073in}}{\pgfqpoint{4.597602in}{7.375822in}}%
\pgfpathcurveto{\pgfqpoint{4.604853in}{7.368571in}}{\pgfqpoint{4.614689in}{7.364497in}}{\pgfqpoint{4.624943in}{7.364497in}}%
\pgfpathlineto{\pgfqpoint{4.624943in}{7.364497in}}%
\pgfpathclose%
\pgfusepath{stroke}%
\end{pgfscope}%
\begin{pgfscope}%
\pgfpathrectangle{\pgfqpoint{1.109909in}{4.050000in}}{\pgfqpoint{5.400000in}{5.400000in}}%
\pgfusepath{clip}%
\pgfsetbuttcap%
\pgfsetroundjoin%
\pgfsetlinewidth{1.003750pt}%
\definecolor{currentstroke}{rgb}{0.541176,0.380392,0.658824}%
\pgfsetstrokecolor{currentstroke}%
\pgfsetstrokeopacity{0.900000}%
\pgfsetdash{}{0pt}%
\pgfpathmoveto{\pgfqpoint{3.716348in}{7.081861in}}%
\pgfpathcurveto{\pgfqpoint{3.726602in}{7.081861in}}{\pgfqpoint{3.736437in}{7.085935in}}{\pgfqpoint{3.743688in}{7.093186in}}%
\pgfpathcurveto{\pgfqpoint{3.750939in}{7.100437in}}{\pgfqpoint{3.755013in}{7.110272in}}{\pgfqpoint{3.755013in}{7.120526in}}%
\pgfpathcurveto{\pgfqpoint{3.755013in}{7.130780in}}{\pgfqpoint{3.750939in}{7.140616in}}{\pgfqpoint{3.743688in}{7.147867in}}%
\pgfpathcurveto{\pgfqpoint{3.736437in}{7.155117in}}{\pgfqpoint{3.726602in}{7.159191in}}{\pgfqpoint{3.716348in}{7.159191in}}%
\pgfpathcurveto{\pgfqpoint{3.706094in}{7.159191in}}{\pgfqpoint{3.696258in}{7.155117in}}{\pgfqpoint{3.689007in}{7.147867in}}%
\pgfpathcurveto{\pgfqpoint{3.681757in}{7.140616in}}{\pgfqpoint{3.677683in}{7.130780in}}{\pgfqpoint{3.677683in}{7.120526in}}%
\pgfpathcurveto{\pgfqpoint{3.677683in}{7.110272in}}{\pgfqpoint{3.681757in}{7.100437in}}{\pgfqpoint{3.689007in}{7.093186in}}%
\pgfpathcurveto{\pgfqpoint{3.696258in}{7.085935in}}{\pgfqpoint{3.706094in}{7.081861in}}{\pgfqpoint{3.716348in}{7.081861in}}%
\pgfpathlineto{\pgfqpoint{3.716348in}{7.081861in}}%
\pgfpathclose%
\pgfusepath{stroke}%
\end{pgfscope}%
\begin{pgfscope}%
\pgfpathrectangle{\pgfqpoint{1.109909in}{4.050000in}}{\pgfqpoint{5.400000in}{5.400000in}}%
\pgfusepath{clip}%
\pgfsetbuttcap%
\pgfsetroundjoin%
\pgfsetlinewidth{1.003750pt}%
\definecolor{currentstroke}{rgb}{0.541176,0.380392,0.658824}%
\pgfsetstrokecolor{currentstroke}%
\pgfsetstrokeopacity{0.900000}%
\pgfsetdash{}{0pt}%
\pgfpathmoveto{\pgfqpoint{4.912794in}{8.046450in}}%
\pgfpathcurveto{\pgfqpoint{4.923048in}{8.046450in}}{\pgfqpoint{4.932883in}{8.050524in}}{\pgfqpoint{4.940134in}{8.057775in}}%
\pgfpathcurveto{\pgfqpoint{4.947385in}{8.065026in}}{\pgfqpoint{4.951459in}{8.074861in}}{\pgfqpoint{4.951459in}{8.085115in}}%
\pgfpathcurveto{\pgfqpoint{4.951459in}{8.095370in}}{\pgfqpoint{4.947385in}{8.105205in}}{\pgfqpoint{4.940134in}{8.112456in}}%
\pgfpathcurveto{\pgfqpoint{4.932883in}{8.119706in}}{\pgfqpoint{4.923048in}{8.123780in}}{\pgfqpoint{4.912794in}{8.123780in}}%
\pgfpathcurveto{\pgfqpoint{4.902539in}{8.123780in}}{\pgfqpoint{4.892704in}{8.119706in}}{\pgfqpoint{4.885453in}{8.112456in}}%
\pgfpathcurveto{\pgfqpoint{4.878202in}{8.105205in}}{\pgfqpoint{4.874128in}{8.095370in}}{\pgfqpoint{4.874128in}{8.085115in}}%
\pgfpathcurveto{\pgfqpoint{4.874128in}{8.074861in}}{\pgfqpoint{4.878202in}{8.065026in}}{\pgfqpoint{4.885453in}{8.057775in}}%
\pgfpathcurveto{\pgfqpoint{4.892704in}{8.050524in}}{\pgfqpoint{4.902539in}{8.046450in}}{\pgfqpoint{4.912794in}{8.046450in}}%
\pgfpathlineto{\pgfqpoint{4.912794in}{8.046450in}}%
\pgfpathclose%
\pgfusepath{stroke}%
\end{pgfscope}%
\begin{pgfscope}%
\pgfpathrectangle{\pgfqpoint{1.109909in}{4.050000in}}{\pgfqpoint{5.400000in}{5.400000in}}%
\pgfusepath{clip}%
\pgfsetbuttcap%
\pgfsetroundjoin%
\pgfsetlinewidth{1.003750pt}%
\definecolor{currentstroke}{rgb}{0.541176,0.380392,0.658824}%
\pgfsetstrokecolor{currentstroke}%
\pgfsetstrokeopacity{0.900000}%
\pgfsetdash{}{0pt}%
\pgfpathmoveto{\pgfqpoint{4.249382in}{7.760512in}}%
\pgfpathcurveto{\pgfqpoint{4.259636in}{7.760512in}}{\pgfqpoint{4.269472in}{7.764586in}}{\pgfqpoint{4.276723in}{7.771837in}}%
\pgfpathcurveto{\pgfqpoint{4.283973in}{7.779088in}}{\pgfqpoint{4.288047in}{7.788923in}}{\pgfqpoint{4.288047in}{7.799177in}}%
\pgfpathcurveto{\pgfqpoint{4.288047in}{7.809431in}}{\pgfqpoint{4.283973in}{7.819267in}}{\pgfqpoint{4.276723in}{7.826517in}}%
\pgfpathcurveto{\pgfqpoint{4.269472in}{7.833768in}}{\pgfqpoint{4.259636in}{7.837842in}}{\pgfqpoint{4.249382in}{7.837842in}}%
\pgfpathcurveto{\pgfqpoint{4.239128in}{7.837842in}}{\pgfqpoint{4.229293in}{7.833768in}}{\pgfqpoint{4.222042in}{7.826517in}}%
\pgfpathcurveto{\pgfqpoint{4.214791in}{7.819267in}}{\pgfqpoint{4.210717in}{7.809431in}}{\pgfqpoint{4.210717in}{7.799177in}}%
\pgfpathcurveto{\pgfqpoint{4.210717in}{7.788923in}}{\pgfqpoint{4.214791in}{7.779088in}}{\pgfqpoint{4.222042in}{7.771837in}}%
\pgfpathcurveto{\pgfqpoint{4.229293in}{7.764586in}}{\pgfqpoint{4.239128in}{7.760512in}}{\pgfqpoint{4.249382in}{7.760512in}}%
\pgfpathlineto{\pgfqpoint{4.249382in}{7.760512in}}%
\pgfpathclose%
\pgfusepath{stroke}%
\end{pgfscope}%
\begin{pgfscope}%
\pgfpathrectangle{\pgfqpoint{1.109909in}{4.050000in}}{\pgfqpoint{5.400000in}{5.400000in}}%
\pgfusepath{clip}%
\pgfsetbuttcap%
\pgfsetroundjoin%
\pgfsetlinewidth{1.003750pt}%
\definecolor{currentstroke}{rgb}{0.541176,0.380392,0.658824}%
\pgfsetstrokecolor{currentstroke}%
\pgfsetstrokeopacity{0.900000}%
\pgfsetdash{}{0pt}%
\pgfpathmoveto{\pgfqpoint{4.162220in}{6.535484in}}%
\pgfpathcurveto{\pgfqpoint{4.172474in}{6.535484in}}{\pgfqpoint{4.182309in}{6.539558in}}{\pgfqpoint{4.189560in}{6.546808in}}%
\pgfpathcurveto{\pgfqpoint{4.196811in}{6.554059in}}{\pgfqpoint{4.200885in}{6.563895in}}{\pgfqpoint{4.200885in}{6.574149in}}%
\pgfpathcurveto{\pgfqpoint{4.200885in}{6.584403in}}{\pgfqpoint{4.196811in}{6.594238in}}{\pgfqpoint{4.189560in}{6.601489in}}%
\pgfpathcurveto{\pgfqpoint{4.182309in}{6.608740in}}{\pgfqpoint{4.172474in}{6.612814in}}{\pgfqpoint{4.162220in}{6.612814in}}%
\pgfpathcurveto{\pgfqpoint{4.151966in}{6.612814in}}{\pgfqpoint{4.142130in}{6.608740in}}{\pgfqpoint{4.134879in}{6.601489in}}%
\pgfpathcurveto{\pgfqpoint{4.127629in}{6.594238in}}{\pgfqpoint{4.123555in}{6.584403in}}{\pgfqpoint{4.123555in}{6.574149in}}%
\pgfpathcurveto{\pgfqpoint{4.123555in}{6.563895in}}{\pgfqpoint{4.127629in}{6.554059in}}{\pgfqpoint{4.134879in}{6.546808in}}%
\pgfpathcurveto{\pgfqpoint{4.142130in}{6.539558in}}{\pgfqpoint{4.151966in}{6.535484in}}{\pgfqpoint{4.162220in}{6.535484in}}%
\pgfpathlineto{\pgfqpoint{4.162220in}{6.535484in}}%
\pgfpathclose%
\pgfusepath{stroke}%
\end{pgfscope}%
\begin{pgfscope}%
\pgfpathrectangle{\pgfqpoint{1.109909in}{4.050000in}}{\pgfqpoint{5.400000in}{5.400000in}}%
\pgfusepath{clip}%
\pgfsetbuttcap%
\pgfsetroundjoin%
\pgfsetlinewidth{1.003750pt}%
\definecolor{currentstroke}{rgb}{0.541176,0.380392,0.658824}%
\pgfsetstrokecolor{currentstroke}%
\pgfsetstrokeopacity{0.900000}%
\pgfsetdash{}{0pt}%
\pgfpathmoveto{\pgfqpoint{2.975455in}{5.240758in}}%
\pgfpathcurveto{\pgfqpoint{2.985709in}{5.240758in}}{\pgfqpoint{2.995545in}{5.244832in}}{\pgfqpoint{3.002795in}{5.252083in}}%
\pgfpathcurveto{\pgfqpoint{3.010046in}{5.259334in}}{\pgfqpoint{3.014120in}{5.269169in}}{\pgfqpoint{3.014120in}{5.279423in}}%
\pgfpathcurveto{\pgfqpoint{3.014120in}{5.289677in}}{\pgfqpoint{3.010046in}{5.299513in}}{\pgfqpoint{3.002795in}{5.306764in}}%
\pgfpathcurveto{\pgfqpoint{2.995545in}{5.314014in}}{\pgfqpoint{2.985709in}{5.318088in}}{\pgfqpoint{2.975455in}{5.318088in}}%
\pgfpathcurveto{\pgfqpoint{2.965201in}{5.318088in}}{\pgfqpoint{2.955366in}{5.314014in}}{\pgfqpoint{2.948115in}{5.306764in}}%
\pgfpathcurveto{\pgfqpoint{2.940864in}{5.299513in}}{\pgfqpoint{2.936790in}{5.289677in}}{\pgfqpoint{2.936790in}{5.279423in}}%
\pgfpathcurveto{\pgfqpoint{2.936790in}{5.269169in}}{\pgfqpoint{2.940864in}{5.259334in}}{\pgfqpoint{2.948115in}{5.252083in}}%
\pgfpathcurveto{\pgfqpoint{2.955366in}{5.244832in}}{\pgfqpoint{2.965201in}{5.240758in}}{\pgfqpoint{2.975455in}{5.240758in}}%
\pgfpathlineto{\pgfqpoint{2.975455in}{5.240758in}}%
\pgfpathclose%
\pgfusepath{stroke}%
\end{pgfscope}%
\begin{pgfscope}%
\pgfpathrectangle{\pgfqpoint{1.109909in}{4.050000in}}{\pgfqpoint{5.400000in}{5.400000in}}%
\pgfusepath{clip}%
\pgfsetbuttcap%
\pgfsetroundjoin%
\pgfsetlinewidth{1.003750pt}%
\definecolor{currentstroke}{rgb}{0.541176,0.380392,0.658824}%
\pgfsetstrokecolor{currentstroke}%
\pgfsetstrokeopacity{0.900000}%
\pgfsetdash{}{0pt}%
\pgfpathmoveto{\pgfqpoint{3.161596in}{6.432515in}}%
\pgfpathcurveto{\pgfqpoint{3.171850in}{6.432515in}}{\pgfqpoint{3.181686in}{6.436589in}}{\pgfqpoint{3.188936in}{6.443840in}}%
\pgfpathcurveto{\pgfqpoint{3.196187in}{6.451091in}}{\pgfqpoint{3.200261in}{6.460926in}}{\pgfqpoint{3.200261in}{6.471180in}}%
\pgfpathcurveto{\pgfqpoint{3.200261in}{6.481434in}}{\pgfqpoint{3.196187in}{6.491270in}}{\pgfqpoint{3.188936in}{6.498521in}}%
\pgfpathcurveto{\pgfqpoint{3.181686in}{6.505771in}}{\pgfqpoint{3.171850in}{6.509845in}}{\pgfqpoint{3.161596in}{6.509845in}}%
\pgfpathcurveto{\pgfqpoint{3.151342in}{6.509845in}}{\pgfqpoint{3.141506in}{6.505771in}}{\pgfqpoint{3.134256in}{6.498521in}}%
\pgfpathcurveto{\pgfqpoint{3.127005in}{6.491270in}}{\pgfqpoint{3.122931in}{6.481434in}}{\pgfqpoint{3.122931in}{6.471180in}}%
\pgfpathcurveto{\pgfqpoint{3.122931in}{6.460926in}}{\pgfqpoint{3.127005in}{6.451091in}}{\pgfqpoint{3.134256in}{6.443840in}}%
\pgfpathcurveto{\pgfqpoint{3.141506in}{6.436589in}}{\pgfqpoint{3.151342in}{6.432515in}}{\pgfqpoint{3.161596in}{6.432515in}}%
\pgfpathlineto{\pgfqpoint{3.161596in}{6.432515in}}%
\pgfpathclose%
\pgfusepath{stroke}%
\end{pgfscope}%
\begin{pgfscope}%
\pgfpathrectangle{\pgfqpoint{1.109909in}{4.050000in}}{\pgfqpoint{5.400000in}{5.400000in}}%
\pgfusepath{clip}%
\pgfsetbuttcap%
\pgfsetroundjoin%
\pgfsetlinewidth{1.003750pt}%
\definecolor{currentstroke}{rgb}{0.541176,0.380392,0.658824}%
\pgfsetstrokecolor{currentstroke}%
\pgfsetstrokeopacity{0.900000}%
\pgfsetdash{}{0pt}%
\pgfpathmoveto{\pgfqpoint{3.395395in}{5.918499in}}%
\pgfpathcurveto{\pgfqpoint{3.405650in}{5.918499in}}{\pgfqpoint{3.415485in}{5.922573in}}{\pgfqpoint{3.422736in}{5.929823in}}%
\pgfpathcurveto{\pgfqpoint{3.429987in}{5.937074in}}{\pgfqpoint{3.434061in}{5.946910in}}{\pgfqpoint{3.434061in}{5.957164in}}%
\pgfpathcurveto{\pgfqpoint{3.434061in}{5.967418in}}{\pgfqpoint{3.429987in}{5.977253in}}{\pgfqpoint{3.422736in}{5.984504in}}%
\pgfpathcurveto{\pgfqpoint{3.415485in}{5.991755in}}{\pgfqpoint{3.405650in}{5.995829in}}{\pgfqpoint{3.395395in}{5.995829in}}%
\pgfpathcurveto{\pgfqpoint{3.385141in}{5.995829in}}{\pgfqpoint{3.375306in}{5.991755in}}{\pgfqpoint{3.368055in}{5.984504in}}%
\pgfpathcurveto{\pgfqpoint{3.360804in}{5.977253in}}{\pgfqpoint{3.356730in}{5.967418in}}{\pgfqpoint{3.356730in}{5.957164in}}%
\pgfpathcurveto{\pgfqpoint{3.356730in}{5.946910in}}{\pgfqpoint{3.360804in}{5.937074in}}{\pgfqpoint{3.368055in}{5.929823in}}%
\pgfpathcurveto{\pgfqpoint{3.375306in}{5.922573in}}{\pgfqpoint{3.385141in}{5.918499in}}{\pgfqpoint{3.395395in}{5.918499in}}%
\pgfpathlineto{\pgfqpoint{3.395395in}{5.918499in}}%
\pgfpathclose%
\pgfusepath{stroke}%
\end{pgfscope}%
\begin{pgfscope}%
\pgfpathrectangle{\pgfqpoint{1.109909in}{4.050000in}}{\pgfqpoint{5.400000in}{5.400000in}}%
\pgfusepath{clip}%
\pgfsetbuttcap%
\pgfsetroundjoin%
\pgfsetlinewidth{1.003750pt}%
\definecolor{currentstroke}{rgb}{0.541176,0.380392,0.658824}%
\pgfsetstrokecolor{currentstroke}%
\pgfsetstrokeopacity{0.900000}%
\pgfsetdash{}{0pt}%
\pgfpathmoveto{\pgfqpoint{2.954013in}{5.522240in}}%
\pgfpathcurveto{\pgfqpoint{2.964267in}{5.522240in}}{\pgfqpoint{2.974102in}{5.526314in}}{\pgfqpoint{2.981353in}{5.533565in}}%
\pgfpathcurveto{\pgfqpoint{2.988604in}{5.540816in}}{\pgfqpoint{2.992678in}{5.550651in}}{\pgfqpoint{2.992678in}{5.560905in}}%
\pgfpathcurveto{\pgfqpoint{2.992678in}{5.571160in}}{\pgfqpoint{2.988604in}{5.580995in}}{\pgfqpoint{2.981353in}{5.588246in}}%
\pgfpathcurveto{\pgfqpoint{2.974102in}{5.595497in}}{\pgfqpoint{2.964267in}{5.599571in}}{\pgfqpoint{2.954013in}{5.599571in}}%
\pgfpathcurveto{\pgfqpoint{2.943759in}{5.599571in}}{\pgfqpoint{2.933923in}{5.595497in}}{\pgfqpoint{2.926672in}{5.588246in}}%
\pgfpathcurveto{\pgfqpoint{2.919422in}{5.580995in}}{\pgfqpoint{2.915348in}{5.571160in}}{\pgfqpoint{2.915348in}{5.560905in}}%
\pgfpathcurveto{\pgfqpoint{2.915348in}{5.550651in}}{\pgfqpoint{2.919422in}{5.540816in}}{\pgfqpoint{2.926672in}{5.533565in}}%
\pgfpathcurveto{\pgfqpoint{2.933923in}{5.526314in}}{\pgfqpoint{2.943759in}{5.522240in}}{\pgfqpoint{2.954013in}{5.522240in}}%
\pgfpathlineto{\pgfqpoint{2.954013in}{5.522240in}}%
\pgfpathclose%
\pgfusepath{stroke}%
\end{pgfscope}%
\begin{pgfscope}%
\pgfpathrectangle{\pgfqpoint{1.109909in}{4.050000in}}{\pgfqpoint{5.400000in}{5.400000in}}%
\pgfusepath{clip}%
\pgfsetbuttcap%
\pgfsetroundjoin%
\pgfsetlinewidth{1.003750pt}%
\definecolor{currentstroke}{rgb}{0.541176,0.380392,0.658824}%
\pgfsetstrokecolor{currentstroke}%
\pgfsetstrokeopacity{0.900000}%
\pgfsetdash{}{0pt}%
\pgfpathmoveto{\pgfqpoint{3.381345in}{5.767861in}}%
\pgfpathcurveto{\pgfqpoint{3.391599in}{5.767861in}}{\pgfqpoint{3.401435in}{5.771935in}}{\pgfqpoint{3.408686in}{5.779186in}}%
\pgfpathcurveto{\pgfqpoint{3.415936in}{5.786437in}}{\pgfqpoint{3.420010in}{5.796272in}}{\pgfqpoint{3.420010in}{5.806526in}}%
\pgfpathcurveto{\pgfqpoint{3.420010in}{5.816780in}}{\pgfqpoint{3.415936in}{5.826616in}}{\pgfqpoint{3.408686in}{5.833867in}}%
\pgfpathcurveto{\pgfqpoint{3.401435in}{5.841117in}}{\pgfqpoint{3.391599in}{5.845191in}}{\pgfqpoint{3.381345in}{5.845191in}}%
\pgfpathcurveto{\pgfqpoint{3.371091in}{5.845191in}}{\pgfqpoint{3.361256in}{5.841117in}}{\pgfqpoint{3.354005in}{5.833867in}}%
\pgfpathcurveto{\pgfqpoint{3.346754in}{5.826616in}}{\pgfqpoint{3.342680in}{5.816780in}}{\pgfqpoint{3.342680in}{5.806526in}}%
\pgfpathcurveto{\pgfqpoint{3.342680in}{5.796272in}}{\pgfqpoint{3.346754in}{5.786437in}}{\pgfqpoint{3.354005in}{5.779186in}}%
\pgfpathcurveto{\pgfqpoint{3.361256in}{5.771935in}}{\pgfqpoint{3.371091in}{5.767861in}}{\pgfqpoint{3.381345in}{5.767861in}}%
\pgfpathlineto{\pgfqpoint{3.381345in}{5.767861in}}%
\pgfpathclose%
\pgfusepath{stroke}%
\end{pgfscope}%
\begin{pgfscope}%
\pgfpathrectangle{\pgfqpoint{1.109909in}{4.050000in}}{\pgfqpoint{5.400000in}{5.400000in}}%
\pgfusepath{clip}%
\pgfsetbuttcap%
\pgfsetroundjoin%
\pgfsetlinewidth{1.003750pt}%
\definecolor{currentstroke}{rgb}{0.541176,0.380392,0.658824}%
\pgfsetstrokecolor{currentstroke}%
\pgfsetstrokeopacity{0.900000}%
\pgfsetdash{}{0pt}%
\pgfpathmoveto{\pgfqpoint{3.876176in}{6.634396in}}%
\pgfpathcurveto{\pgfqpoint{3.886430in}{6.634396in}}{\pgfqpoint{3.896265in}{6.638470in}}{\pgfqpoint{3.903516in}{6.645721in}}%
\pgfpathcurveto{\pgfqpoint{3.910767in}{6.652971in}}{\pgfqpoint{3.914841in}{6.662807in}}{\pgfqpoint{3.914841in}{6.673061in}}%
\pgfpathcurveto{\pgfqpoint{3.914841in}{6.683315in}}{\pgfqpoint{3.910767in}{6.693150in}}{\pgfqpoint{3.903516in}{6.700401in}}%
\pgfpathcurveto{\pgfqpoint{3.896265in}{6.707652in}}{\pgfqpoint{3.886430in}{6.711726in}}{\pgfqpoint{3.876176in}{6.711726in}}%
\pgfpathcurveto{\pgfqpoint{3.865922in}{6.711726in}}{\pgfqpoint{3.856086in}{6.707652in}}{\pgfqpoint{3.848835in}{6.700401in}}%
\pgfpathcurveto{\pgfqpoint{3.841585in}{6.693150in}}{\pgfqpoint{3.837511in}{6.683315in}}{\pgfqpoint{3.837511in}{6.673061in}}%
\pgfpathcurveto{\pgfqpoint{3.837511in}{6.662807in}}{\pgfqpoint{3.841585in}{6.652971in}}{\pgfqpoint{3.848835in}{6.645721in}}%
\pgfpathcurveto{\pgfqpoint{3.856086in}{6.638470in}}{\pgfqpoint{3.865922in}{6.634396in}}{\pgfqpoint{3.876176in}{6.634396in}}%
\pgfpathlineto{\pgfqpoint{3.876176in}{6.634396in}}%
\pgfpathclose%
\pgfusepath{stroke}%
\end{pgfscope}%
\begin{pgfscope}%
\pgfpathrectangle{\pgfqpoint{1.109909in}{4.050000in}}{\pgfqpoint{5.400000in}{5.400000in}}%
\pgfusepath{clip}%
\pgfsetbuttcap%
\pgfsetroundjoin%
\pgfsetlinewidth{1.003750pt}%
\definecolor{currentstroke}{rgb}{0.541176,0.380392,0.658824}%
\pgfsetstrokecolor{currentstroke}%
\pgfsetstrokeopacity{0.900000}%
\pgfsetdash{}{0pt}%
\pgfpathmoveto{\pgfqpoint{4.233335in}{6.644954in}}%
\pgfpathcurveto{\pgfqpoint{4.243589in}{6.644954in}}{\pgfqpoint{4.253424in}{6.649028in}}{\pgfqpoint{4.260675in}{6.656279in}}%
\pgfpathcurveto{\pgfqpoint{4.267926in}{6.663529in}}{\pgfqpoint{4.272000in}{6.673365in}}{\pgfqpoint{4.272000in}{6.683619in}}%
\pgfpathcurveto{\pgfqpoint{4.272000in}{6.693873in}}{\pgfqpoint{4.267926in}{6.703709in}}{\pgfqpoint{4.260675in}{6.710959in}}%
\pgfpathcurveto{\pgfqpoint{4.253424in}{6.718210in}}{\pgfqpoint{4.243589in}{6.722284in}}{\pgfqpoint{4.233335in}{6.722284in}}%
\pgfpathcurveto{\pgfqpoint{4.223081in}{6.722284in}}{\pgfqpoint{4.213245in}{6.718210in}}{\pgfqpoint{4.205995in}{6.710959in}}%
\pgfpathcurveto{\pgfqpoint{4.198744in}{6.703709in}}{\pgfqpoint{4.194670in}{6.693873in}}{\pgfqpoint{4.194670in}{6.683619in}}%
\pgfpathcurveto{\pgfqpoint{4.194670in}{6.673365in}}{\pgfqpoint{4.198744in}{6.663529in}}{\pgfqpoint{4.205995in}{6.656279in}}%
\pgfpathcurveto{\pgfqpoint{4.213245in}{6.649028in}}{\pgfqpoint{4.223081in}{6.644954in}}{\pgfqpoint{4.233335in}{6.644954in}}%
\pgfpathlineto{\pgfqpoint{4.233335in}{6.644954in}}%
\pgfpathclose%
\pgfusepath{stroke}%
\end{pgfscope}%
\begin{pgfscope}%
\pgfpathrectangle{\pgfqpoint{1.109909in}{4.050000in}}{\pgfqpoint{5.400000in}{5.400000in}}%
\pgfusepath{clip}%
\pgfsetbuttcap%
\pgfsetroundjoin%
\pgfsetlinewidth{1.003750pt}%
\definecolor{currentstroke}{rgb}{0.541176,0.380392,0.658824}%
\pgfsetstrokecolor{currentstroke}%
\pgfsetstrokeopacity{0.900000}%
\pgfsetdash{}{0pt}%
\pgfpathmoveto{\pgfqpoint{4.686104in}{7.846921in}}%
\pgfpathcurveto{\pgfqpoint{4.696358in}{7.846921in}}{\pgfqpoint{4.706193in}{7.850995in}}{\pgfqpoint{4.713444in}{7.858246in}}%
\pgfpathcurveto{\pgfqpoint{4.720695in}{7.865497in}}{\pgfqpoint{4.724769in}{7.875332in}}{\pgfqpoint{4.724769in}{7.885586in}}%
\pgfpathcurveto{\pgfqpoint{4.724769in}{7.895840in}}{\pgfqpoint{4.720695in}{7.905676in}}{\pgfqpoint{4.713444in}{7.912927in}}%
\pgfpathcurveto{\pgfqpoint{4.706193in}{7.920177in}}{\pgfqpoint{4.696358in}{7.924251in}}{\pgfqpoint{4.686104in}{7.924251in}}%
\pgfpathcurveto{\pgfqpoint{4.675849in}{7.924251in}}{\pgfqpoint{4.666014in}{7.920177in}}{\pgfqpoint{4.658763in}{7.912927in}}%
\pgfpathcurveto{\pgfqpoint{4.651512in}{7.905676in}}{\pgfqpoint{4.647438in}{7.895840in}}{\pgfqpoint{4.647438in}{7.885586in}}%
\pgfpathcurveto{\pgfqpoint{4.647438in}{7.875332in}}{\pgfqpoint{4.651512in}{7.865497in}}{\pgfqpoint{4.658763in}{7.858246in}}%
\pgfpathcurveto{\pgfqpoint{4.666014in}{7.850995in}}{\pgfqpoint{4.675849in}{7.846921in}}{\pgfqpoint{4.686104in}{7.846921in}}%
\pgfpathlineto{\pgfqpoint{4.686104in}{7.846921in}}%
\pgfpathclose%
\pgfusepath{stroke}%
\end{pgfscope}%
\begin{pgfscope}%
\pgfpathrectangle{\pgfqpoint{1.109909in}{4.050000in}}{\pgfqpoint{5.400000in}{5.400000in}}%
\pgfusepath{clip}%
\pgfsetbuttcap%
\pgfsetroundjoin%
\pgfsetlinewidth{1.003750pt}%
\definecolor{currentstroke}{rgb}{0.541176,0.380392,0.658824}%
\pgfsetstrokecolor{currentstroke}%
\pgfsetstrokeopacity{0.900000}%
\pgfsetdash{}{0pt}%
\pgfpathmoveto{\pgfqpoint{4.297443in}{6.753311in}}%
\pgfpathcurveto{\pgfqpoint{4.307697in}{6.753311in}}{\pgfqpoint{4.317532in}{6.757385in}}{\pgfqpoint{4.324783in}{6.764636in}}%
\pgfpathcurveto{\pgfqpoint{4.332034in}{6.771886in}}{\pgfqpoint{4.336108in}{6.781722in}}{\pgfqpoint{4.336108in}{6.791976in}}%
\pgfpathcurveto{\pgfqpoint{4.336108in}{6.802230in}}{\pgfqpoint{4.332034in}{6.812066in}}{\pgfqpoint{4.324783in}{6.819316in}}%
\pgfpathcurveto{\pgfqpoint{4.317532in}{6.826567in}}{\pgfqpoint{4.307697in}{6.830641in}}{\pgfqpoint{4.297443in}{6.830641in}}%
\pgfpathcurveto{\pgfqpoint{4.287189in}{6.830641in}}{\pgfqpoint{4.277353in}{6.826567in}}{\pgfqpoint{4.270102in}{6.819316in}}%
\pgfpathcurveto{\pgfqpoint{4.262852in}{6.812066in}}{\pgfqpoint{4.258778in}{6.802230in}}{\pgfqpoint{4.258778in}{6.791976in}}%
\pgfpathcurveto{\pgfqpoint{4.258778in}{6.781722in}}{\pgfqpoint{4.262852in}{6.771886in}}{\pgfqpoint{4.270102in}{6.764636in}}%
\pgfpathcurveto{\pgfqpoint{4.277353in}{6.757385in}}{\pgfqpoint{4.287189in}{6.753311in}}{\pgfqpoint{4.297443in}{6.753311in}}%
\pgfpathlineto{\pgfqpoint{4.297443in}{6.753311in}}%
\pgfpathclose%
\pgfusepath{stroke}%
\end{pgfscope}%
\begin{pgfscope}%
\pgfpathrectangle{\pgfqpoint{1.109909in}{4.050000in}}{\pgfqpoint{5.400000in}{5.400000in}}%
\pgfusepath{clip}%
\pgfsetbuttcap%
\pgfsetroundjoin%
\pgfsetlinewidth{1.003750pt}%
\definecolor{currentstroke}{rgb}{0.541176,0.380392,0.658824}%
\pgfsetstrokecolor{currentstroke}%
\pgfsetstrokeopacity{0.900000}%
\pgfsetdash{}{0pt}%
\pgfpathmoveto{\pgfqpoint{4.952343in}{7.536182in}}%
\pgfpathcurveto{\pgfqpoint{4.962597in}{7.536182in}}{\pgfqpoint{4.972432in}{7.540256in}}{\pgfqpoint{4.979683in}{7.547507in}}%
\pgfpathcurveto{\pgfqpoint{4.986934in}{7.554758in}}{\pgfqpoint{4.991008in}{7.564593in}}{\pgfqpoint{4.991008in}{7.574847in}}%
\pgfpathcurveto{\pgfqpoint{4.991008in}{7.585101in}}{\pgfqpoint{4.986934in}{7.594937in}}{\pgfqpoint{4.979683in}{7.602188in}}%
\pgfpathcurveto{\pgfqpoint{4.972432in}{7.609438in}}{\pgfqpoint{4.962597in}{7.613512in}}{\pgfqpoint{4.952343in}{7.613512in}}%
\pgfpathcurveto{\pgfqpoint{4.942088in}{7.613512in}}{\pgfqpoint{4.932253in}{7.609438in}}{\pgfqpoint{4.925002in}{7.602188in}}%
\pgfpathcurveto{\pgfqpoint{4.917752in}{7.594937in}}{\pgfqpoint{4.913678in}{7.585101in}}{\pgfqpoint{4.913678in}{7.574847in}}%
\pgfpathcurveto{\pgfqpoint{4.913678in}{7.564593in}}{\pgfqpoint{4.917752in}{7.554758in}}{\pgfqpoint{4.925002in}{7.547507in}}%
\pgfpathcurveto{\pgfqpoint{4.932253in}{7.540256in}}{\pgfqpoint{4.942088in}{7.536182in}}{\pgfqpoint{4.952343in}{7.536182in}}%
\pgfpathlineto{\pgfqpoint{4.952343in}{7.536182in}}%
\pgfpathclose%
\pgfusepath{stroke}%
\end{pgfscope}%
\begin{pgfscope}%
\pgfpathrectangle{\pgfqpoint{1.109909in}{4.050000in}}{\pgfqpoint{5.400000in}{5.400000in}}%
\pgfusepath{clip}%
\pgfsetbuttcap%
\pgfsetroundjoin%
\pgfsetlinewidth{1.003750pt}%
\definecolor{currentstroke}{rgb}{0.541176,0.380392,0.658824}%
\pgfsetstrokecolor{currentstroke}%
\pgfsetstrokeopacity{0.900000}%
\pgfsetdash{}{0pt}%
\pgfpathmoveto{\pgfqpoint{4.741417in}{6.839179in}}%
\pgfpathcurveto{\pgfqpoint{4.751671in}{6.839179in}}{\pgfqpoint{4.761506in}{6.843253in}}{\pgfqpoint{4.768757in}{6.850503in}}%
\pgfpathcurveto{\pgfqpoint{4.776008in}{6.857754in}}{\pgfqpoint{4.780082in}{6.867590in}}{\pgfqpoint{4.780082in}{6.877844in}}%
\pgfpathcurveto{\pgfqpoint{4.780082in}{6.888098in}}{\pgfqpoint{4.776008in}{6.897933in}}{\pgfqpoint{4.768757in}{6.905184in}}%
\pgfpathcurveto{\pgfqpoint{4.761506in}{6.912435in}}{\pgfqpoint{4.751671in}{6.916509in}}{\pgfqpoint{4.741417in}{6.916509in}}%
\pgfpathcurveto{\pgfqpoint{4.731163in}{6.916509in}}{\pgfqpoint{4.721327in}{6.912435in}}{\pgfqpoint{4.714076in}{6.905184in}}%
\pgfpathcurveto{\pgfqpoint{4.706826in}{6.897933in}}{\pgfqpoint{4.702752in}{6.888098in}}{\pgfqpoint{4.702752in}{6.877844in}}%
\pgfpathcurveto{\pgfqpoint{4.702752in}{6.867590in}}{\pgfqpoint{4.706826in}{6.857754in}}{\pgfqpoint{4.714076in}{6.850503in}}%
\pgfpathcurveto{\pgfqpoint{4.721327in}{6.843253in}}{\pgfqpoint{4.731163in}{6.839179in}}{\pgfqpoint{4.741417in}{6.839179in}}%
\pgfpathlineto{\pgfqpoint{4.741417in}{6.839179in}}%
\pgfpathclose%
\pgfusepath{stroke}%
\end{pgfscope}%
\begin{pgfscope}%
\pgfpathrectangle{\pgfqpoint{1.109909in}{4.050000in}}{\pgfqpoint{5.400000in}{5.400000in}}%
\pgfusepath{clip}%
\pgfsetbuttcap%
\pgfsetroundjoin%
\pgfsetlinewidth{1.003750pt}%
\definecolor{currentstroke}{rgb}{0.541176,0.380392,0.658824}%
\pgfsetstrokecolor{currentstroke}%
\pgfsetstrokeopacity{0.900000}%
\pgfsetdash{}{0pt}%
\pgfpathmoveto{\pgfqpoint{5.265361in}{7.383555in}}%
\pgfpathcurveto{\pgfqpoint{5.275615in}{7.383555in}}{\pgfqpoint{5.285450in}{7.387629in}}{\pgfqpoint{5.292701in}{7.394880in}}%
\pgfpathcurveto{\pgfqpoint{5.299952in}{7.402130in}}{\pgfqpoint{5.304026in}{7.411966in}}{\pgfqpoint{5.304026in}{7.422220in}}%
\pgfpathcurveto{\pgfqpoint{5.304026in}{7.432474in}}{\pgfqpoint{5.299952in}{7.442309in}}{\pgfqpoint{5.292701in}{7.449560in}}%
\pgfpathcurveto{\pgfqpoint{5.285450in}{7.456811in}}{\pgfqpoint{5.275615in}{7.460885in}}{\pgfqpoint{5.265361in}{7.460885in}}%
\pgfpathcurveto{\pgfqpoint{5.255107in}{7.460885in}}{\pgfqpoint{5.245271in}{7.456811in}}{\pgfqpoint{5.238021in}{7.449560in}}%
\pgfpathcurveto{\pgfqpoint{5.230770in}{7.442309in}}{\pgfqpoint{5.226696in}{7.432474in}}{\pgfqpoint{5.226696in}{7.422220in}}%
\pgfpathcurveto{\pgfqpoint{5.226696in}{7.411966in}}{\pgfqpoint{5.230770in}{7.402130in}}{\pgfqpoint{5.238021in}{7.394880in}}%
\pgfpathcurveto{\pgfqpoint{5.245271in}{7.387629in}}{\pgfqpoint{5.255107in}{7.383555in}}{\pgfqpoint{5.265361in}{7.383555in}}%
\pgfpathlineto{\pgfqpoint{5.265361in}{7.383555in}}%
\pgfpathclose%
\pgfusepath{stroke}%
\end{pgfscope}%
\begin{pgfscope}%
\pgfpathrectangle{\pgfqpoint{1.109909in}{4.050000in}}{\pgfqpoint{5.400000in}{5.400000in}}%
\pgfusepath{clip}%
\pgfsetbuttcap%
\pgfsetroundjoin%
\pgfsetlinewidth{1.003750pt}%
\definecolor{currentstroke}{rgb}{0.541176,0.380392,0.658824}%
\pgfsetstrokecolor{currentstroke}%
\pgfsetstrokeopacity{0.900000}%
\pgfsetdash{}{0pt}%
\pgfpathmoveto{\pgfqpoint{3.838843in}{5.929654in}}%
\pgfpathcurveto{\pgfqpoint{3.849097in}{5.929654in}}{\pgfqpoint{3.858932in}{5.933728in}}{\pgfqpoint{3.866183in}{5.940979in}}%
\pgfpathcurveto{\pgfqpoint{3.873434in}{5.948229in}}{\pgfqpoint{3.877508in}{5.958065in}}{\pgfqpoint{3.877508in}{5.968319in}}%
\pgfpathcurveto{\pgfqpoint{3.877508in}{5.978573in}}{\pgfqpoint{3.873434in}{5.988409in}}{\pgfqpoint{3.866183in}{5.995659in}}%
\pgfpathcurveto{\pgfqpoint{3.858932in}{6.002910in}}{\pgfqpoint{3.849097in}{6.006984in}}{\pgfqpoint{3.838843in}{6.006984in}}%
\pgfpathcurveto{\pgfqpoint{3.828589in}{6.006984in}}{\pgfqpoint{3.818753in}{6.002910in}}{\pgfqpoint{3.811502in}{5.995659in}}%
\pgfpathcurveto{\pgfqpoint{3.804252in}{5.988409in}}{\pgfqpoint{3.800178in}{5.978573in}}{\pgfqpoint{3.800178in}{5.968319in}}%
\pgfpathcurveto{\pgfqpoint{3.800178in}{5.958065in}}{\pgfqpoint{3.804252in}{5.948229in}}{\pgfqpoint{3.811502in}{5.940979in}}%
\pgfpathcurveto{\pgfqpoint{3.818753in}{5.933728in}}{\pgfqpoint{3.828589in}{5.929654in}}{\pgfqpoint{3.838843in}{5.929654in}}%
\pgfpathlineto{\pgfqpoint{3.838843in}{5.929654in}}%
\pgfpathclose%
\pgfusepath{stroke}%
\end{pgfscope}%
\begin{pgfscope}%
\pgfpathrectangle{\pgfqpoint{1.109909in}{4.050000in}}{\pgfqpoint{5.400000in}{5.400000in}}%
\pgfusepath{clip}%
\pgfsetbuttcap%
\pgfsetroundjoin%
\pgfsetlinewidth{1.003750pt}%
\definecolor{currentstroke}{rgb}{0.541176,0.380392,0.658824}%
\pgfsetstrokecolor{currentstroke}%
\pgfsetstrokeopacity{0.900000}%
\pgfsetdash{}{0pt}%
\pgfpathmoveto{\pgfqpoint{5.421135in}{8.559833in}}%
\pgfpathcurveto{\pgfqpoint{5.431389in}{8.559833in}}{\pgfqpoint{5.441224in}{8.563907in}}{\pgfqpoint{5.448475in}{8.571158in}}%
\pgfpathcurveto{\pgfqpoint{5.455726in}{8.578408in}}{\pgfqpoint{5.459800in}{8.588244in}}{\pgfqpoint{5.459800in}{8.598498in}}%
\pgfpathcurveto{\pgfqpoint{5.459800in}{8.608752in}}{\pgfqpoint{5.455726in}{8.618587in}}{\pgfqpoint{5.448475in}{8.625838in}}%
\pgfpathcurveto{\pgfqpoint{5.441224in}{8.633089in}}{\pgfqpoint{5.431389in}{8.637163in}}{\pgfqpoint{5.421135in}{8.637163in}}%
\pgfpathcurveto{\pgfqpoint{5.410881in}{8.637163in}}{\pgfqpoint{5.401045in}{8.633089in}}{\pgfqpoint{5.393794in}{8.625838in}}%
\pgfpathcurveto{\pgfqpoint{5.386544in}{8.618587in}}{\pgfqpoint{5.382470in}{8.608752in}}{\pgfqpoint{5.382470in}{8.598498in}}%
\pgfpathcurveto{\pgfqpoint{5.382470in}{8.588244in}}{\pgfqpoint{5.386544in}{8.578408in}}{\pgfqpoint{5.393794in}{8.571158in}}%
\pgfpathcurveto{\pgfqpoint{5.401045in}{8.563907in}}{\pgfqpoint{5.410881in}{8.559833in}}{\pgfqpoint{5.421135in}{8.559833in}}%
\pgfpathlineto{\pgfqpoint{5.421135in}{8.559833in}}%
\pgfpathclose%
\pgfusepath{stroke}%
\end{pgfscope}%
\begin{pgfscope}%
\pgfpathrectangle{\pgfqpoint{1.109909in}{4.050000in}}{\pgfqpoint{5.400000in}{5.400000in}}%
\pgfusepath{clip}%
\pgfsetbuttcap%
\pgfsetroundjoin%
\pgfsetlinewidth{1.003750pt}%
\definecolor{currentstroke}{rgb}{0.541176,0.380392,0.658824}%
\pgfsetstrokecolor{currentstroke}%
\pgfsetstrokeopacity{0.900000}%
\pgfsetdash{}{0pt}%
\pgfpathmoveto{\pgfqpoint{5.259903in}{8.576765in}}%
\pgfpathcurveto{\pgfqpoint{5.270157in}{8.576765in}}{\pgfqpoint{5.279992in}{8.580839in}}{\pgfqpoint{5.287243in}{8.588090in}}%
\pgfpathcurveto{\pgfqpoint{5.294494in}{8.595340in}}{\pgfqpoint{5.298568in}{8.605176in}}{\pgfqpoint{5.298568in}{8.615430in}}%
\pgfpathcurveto{\pgfqpoint{5.298568in}{8.625684in}}{\pgfqpoint{5.294494in}{8.635520in}}{\pgfqpoint{5.287243in}{8.642770in}}%
\pgfpathcurveto{\pgfqpoint{5.279992in}{8.650021in}}{\pgfqpoint{5.270157in}{8.654095in}}{\pgfqpoint{5.259903in}{8.654095in}}%
\pgfpathcurveto{\pgfqpoint{5.249649in}{8.654095in}}{\pgfqpoint{5.239813in}{8.650021in}}{\pgfqpoint{5.232562in}{8.642770in}}%
\pgfpathcurveto{\pgfqpoint{5.225312in}{8.635520in}}{\pgfqpoint{5.221238in}{8.625684in}}{\pgfqpoint{5.221238in}{8.615430in}}%
\pgfpathcurveto{\pgfqpoint{5.221238in}{8.605176in}}{\pgfqpoint{5.225312in}{8.595340in}}{\pgfqpoint{5.232562in}{8.588090in}}%
\pgfpathcurveto{\pgfqpoint{5.239813in}{8.580839in}}{\pgfqpoint{5.249649in}{8.576765in}}{\pgfqpoint{5.259903in}{8.576765in}}%
\pgfpathlineto{\pgfqpoint{5.259903in}{8.576765in}}%
\pgfpathclose%
\pgfusepath{stroke}%
\end{pgfscope}%
\begin{pgfscope}%
\pgfpathrectangle{\pgfqpoint{1.109909in}{4.050000in}}{\pgfqpoint{5.400000in}{5.400000in}}%
\pgfusepath{clip}%
\pgfsetbuttcap%
\pgfsetroundjoin%
\pgfsetlinewidth{1.003750pt}%
\definecolor{currentstroke}{rgb}{0.541176,0.380392,0.658824}%
\pgfsetstrokecolor{currentstroke}%
\pgfsetstrokeopacity{0.900000}%
\pgfsetdash{}{0pt}%
\pgfpathmoveto{\pgfqpoint{4.917170in}{7.544694in}}%
\pgfpathcurveto{\pgfqpoint{4.927424in}{7.544694in}}{\pgfqpoint{4.937260in}{7.548768in}}{\pgfqpoint{4.944511in}{7.556019in}}%
\pgfpathcurveto{\pgfqpoint{4.951761in}{7.563269in}}{\pgfqpoint{4.955835in}{7.573105in}}{\pgfqpoint{4.955835in}{7.583359in}}%
\pgfpathcurveto{\pgfqpoint{4.955835in}{7.593613in}}{\pgfqpoint{4.951761in}{7.603449in}}{\pgfqpoint{4.944511in}{7.610699in}}%
\pgfpathcurveto{\pgfqpoint{4.937260in}{7.617950in}}{\pgfqpoint{4.927424in}{7.622024in}}{\pgfqpoint{4.917170in}{7.622024in}}%
\pgfpathcurveto{\pgfqpoint{4.906916in}{7.622024in}}{\pgfqpoint{4.897081in}{7.617950in}}{\pgfqpoint{4.889830in}{7.610699in}}%
\pgfpathcurveto{\pgfqpoint{4.882579in}{7.603449in}}{\pgfqpoint{4.878505in}{7.593613in}}{\pgfqpoint{4.878505in}{7.583359in}}%
\pgfpathcurveto{\pgfqpoint{4.878505in}{7.573105in}}{\pgfqpoint{4.882579in}{7.563269in}}{\pgfqpoint{4.889830in}{7.556019in}}%
\pgfpathcurveto{\pgfqpoint{4.897081in}{7.548768in}}{\pgfqpoint{4.906916in}{7.544694in}}{\pgfqpoint{4.917170in}{7.544694in}}%
\pgfpathlineto{\pgfqpoint{4.917170in}{7.544694in}}%
\pgfpathclose%
\pgfusepath{stroke}%
\end{pgfscope}%
\begin{pgfscope}%
\pgfpathrectangle{\pgfqpoint{1.109909in}{4.050000in}}{\pgfqpoint{5.400000in}{5.400000in}}%
\pgfusepath{clip}%
\pgfsetbuttcap%
\pgfsetroundjoin%
\pgfsetlinewidth{1.003750pt}%
\definecolor{currentstroke}{rgb}{0.541176,0.380392,0.658824}%
\pgfsetstrokecolor{currentstroke}%
\pgfsetstrokeopacity{0.900000}%
\pgfsetdash{}{0pt}%
\pgfpathmoveto{\pgfqpoint{4.579500in}{7.952222in}}%
\pgfpathcurveto{\pgfqpoint{4.589754in}{7.952222in}}{\pgfqpoint{4.599590in}{7.956296in}}{\pgfqpoint{4.606841in}{7.963546in}}%
\pgfpathcurveto{\pgfqpoint{4.614091in}{7.970797in}}{\pgfqpoint{4.618165in}{7.980632in}}{\pgfqpoint{4.618165in}{7.990887in}}%
\pgfpathcurveto{\pgfqpoint{4.618165in}{8.001141in}}{\pgfqpoint{4.614091in}{8.010976in}}{\pgfqpoint{4.606841in}{8.018227in}}%
\pgfpathcurveto{\pgfqpoint{4.599590in}{8.025478in}}{\pgfqpoint{4.589754in}{8.029552in}}{\pgfqpoint{4.579500in}{8.029552in}}%
\pgfpathcurveto{\pgfqpoint{4.569246in}{8.029552in}}{\pgfqpoint{4.559411in}{8.025478in}}{\pgfqpoint{4.552160in}{8.018227in}}%
\pgfpathcurveto{\pgfqpoint{4.544909in}{8.010976in}}{\pgfqpoint{4.540835in}{8.001141in}}{\pgfqpoint{4.540835in}{7.990887in}}%
\pgfpathcurveto{\pgfqpoint{4.540835in}{7.980632in}}{\pgfqpoint{4.544909in}{7.970797in}}{\pgfqpoint{4.552160in}{7.963546in}}%
\pgfpathcurveto{\pgfqpoint{4.559411in}{7.956296in}}{\pgfqpoint{4.569246in}{7.952222in}}{\pgfqpoint{4.579500in}{7.952222in}}%
\pgfpathlineto{\pgfqpoint{4.579500in}{7.952222in}}%
\pgfpathclose%
\pgfusepath{stroke}%
\end{pgfscope}%
\begin{pgfscope}%
\pgfpathrectangle{\pgfqpoint{1.109909in}{4.050000in}}{\pgfqpoint{5.400000in}{5.400000in}}%
\pgfusepath{clip}%
\pgfsetbuttcap%
\pgfsetroundjoin%
\pgfsetlinewidth{1.003750pt}%
\definecolor{currentstroke}{rgb}{0.541176,0.380392,0.658824}%
\pgfsetstrokecolor{currentstroke}%
\pgfsetstrokeopacity{0.900000}%
\pgfsetdash{}{0pt}%
\pgfpathmoveto{\pgfqpoint{5.290372in}{8.441682in}}%
\pgfpathcurveto{\pgfqpoint{5.300626in}{8.441682in}}{\pgfqpoint{5.310462in}{8.445756in}}{\pgfqpoint{5.317713in}{8.453007in}}%
\pgfpathcurveto{\pgfqpoint{5.324963in}{8.460257in}}{\pgfqpoint{5.329037in}{8.470093in}}{\pgfqpoint{5.329037in}{8.480347in}}%
\pgfpathcurveto{\pgfqpoint{5.329037in}{8.490601in}}{\pgfqpoint{5.324963in}{8.500436in}}{\pgfqpoint{5.317713in}{8.507687in}}%
\pgfpathcurveto{\pgfqpoint{5.310462in}{8.514938in}}{\pgfqpoint{5.300626in}{8.519012in}}{\pgfqpoint{5.290372in}{8.519012in}}%
\pgfpathcurveto{\pgfqpoint{5.280118in}{8.519012in}}{\pgfqpoint{5.270283in}{8.514938in}}{\pgfqpoint{5.263032in}{8.507687in}}%
\pgfpathcurveto{\pgfqpoint{5.255781in}{8.500436in}}{\pgfqpoint{5.251707in}{8.490601in}}{\pgfqpoint{5.251707in}{8.480347in}}%
\pgfpathcurveto{\pgfqpoint{5.251707in}{8.470093in}}{\pgfqpoint{5.255781in}{8.460257in}}{\pgfqpoint{5.263032in}{8.453007in}}%
\pgfpathcurveto{\pgfqpoint{5.270283in}{8.445756in}}{\pgfqpoint{5.280118in}{8.441682in}}{\pgfqpoint{5.290372in}{8.441682in}}%
\pgfpathlineto{\pgfqpoint{5.290372in}{8.441682in}}%
\pgfpathclose%
\pgfusepath{stroke}%
\end{pgfscope}%
\begin{pgfscope}%
\pgfpathrectangle{\pgfqpoint{1.109909in}{4.050000in}}{\pgfqpoint{5.400000in}{5.400000in}}%
\pgfusepath{clip}%
\pgfsetbuttcap%
\pgfsetroundjoin%
\pgfsetlinewidth{1.003750pt}%
\definecolor{currentstroke}{rgb}{0.541176,0.380392,0.658824}%
\pgfsetstrokecolor{currentstroke}%
\pgfsetstrokeopacity{0.900000}%
\pgfsetdash{}{0pt}%
\pgfpathmoveto{\pgfqpoint{4.169811in}{7.364099in}}%
\pgfpathcurveto{\pgfqpoint{4.180065in}{7.364099in}}{\pgfqpoint{4.189900in}{7.368173in}}{\pgfqpoint{4.197151in}{7.375424in}}%
\pgfpathcurveto{\pgfqpoint{4.204402in}{7.382675in}}{\pgfqpoint{4.208476in}{7.392510in}}{\pgfqpoint{4.208476in}{7.402764in}}%
\pgfpathcurveto{\pgfqpoint{4.208476in}{7.413018in}}{\pgfqpoint{4.204402in}{7.422854in}}{\pgfqpoint{4.197151in}{7.430105in}}%
\pgfpathcurveto{\pgfqpoint{4.189900in}{7.437355in}}{\pgfqpoint{4.180065in}{7.441429in}}{\pgfqpoint{4.169811in}{7.441429in}}%
\pgfpathcurveto{\pgfqpoint{4.159556in}{7.441429in}}{\pgfqpoint{4.149721in}{7.437355in}}{\pgfqpoint{4.142470in}{7.430105in}}%
\pgfpathcurveto{\pgfqpoint{4.135220in}{7.422854in}}{\pgfqpoint{4.131146in}{7.413018in}}{\pgfqpoint{4.131146in}{7.402764in}}%
\pgfpathcurveto{\pgfqpoint{4.131146in}{7.392510in}}{\pgfqpoint{4.135220in}{7.382675in}}{\pgfqpoint{4.142470in}{7.375424in}}%
\pgfpathcurveto{\pgfqpoint{4.149721in}{7.368173in}}{\pgfqpoint{4.159556in}{7.364099in}}{\pgfqpoint{4.169811in}{7.364099in}}%
\pgfpathlineto{\pgfqpoint{4.169811in}{7.364099in}}%
\pgfpathclose%
\pgfusepath{stroke}%
\end{pgfscope}%
\begin{pgfscope}%
\pgfpathrectangle{\pgfqpoint{1.109909in}{4.050000in}}{\pgfqpoint{5.400000in}{5.400000in}}%
\pgfusepath{clip}%
\pgfsetbuttcap%
\pgfsetroundjoin%
\pgfsetlinewidth{1.003750pt}%
\definecolor{currentstroke}{rgb}{0.541176,0.380392,0.658824}%
\pgfsetstrokecolor{currentstroke}%
\pgfsetstrokeopacity{0.900000}%
\pgfsetdash{}{0pt}%
\pgfpathmoveto{\pgfqpoint{4.477958in}{8.048929in}}%
\pgfpathcurveto{\pgfqpoint{4.488212in}{8.048929in}}{\pgfqpoint{4.498048in}{8.053003in}}{\pgfqpoint{4.505298in}{8.060254in}}%
\pgfpathcurveto{\pgfqpoint{4.512549in}{8.067504in}}{\pgfqpoint{4.516623in}{8.077340in}}{\pgfqpoint{4.516623in}{8.087594in}}%
\pgfpathcurveto{\pgfqpoint{4.516623in}{8.097848in}}{\pgfqpoint{4.512549in}{8.107683in}}{\pgfqpoint{4.505298in}{8.114934in}}%
\pgfpathcurveto{\pgfqpoint{4.498048in}{8.122185in}}{\pgfqpoint{4.488212in}{8.126259in}}{\pgfqpoint{4.477958in}{8.126259in}}%
\pgfpathcurveto{\pgfqpoint{4.467704in}{8.126259in}}{\pgfqpoint{4.457868in}{8.122185in}}{\pgfqpoint{4.450618in}{8.114934in}}%
\pgfpathcurveto{\pgfqpoint{4.443367in}{8.107683in}}{\pgfqpoint{4.439293in}{8.097848in}}{\pgfqpoint{4.439293in}{8.087594in}}%
\pgfpathcurveto{\pgfqpoint{4.439293in}{8.077340in}}{\pgfqpoint{4.443367in}{8.067504in}}{\pgfqpoint{4.450618in}{8.060254in}}%
\pgfpathcurveto{\pgfqpoint{4.457868in}{8.053003in}}{\pgfqpoint{4.467704in}{8.048929in}}{\pgfqpoint{4.477958in}{8.048929in}}%
\pgfpathlineto{\pgfqpoint{4.477958in}{8.048929in}}%
\pgfpathclose%
\pgfusepath{stroke}%
\end{pgfscope}%
\begin{pgfscope}%
\pgfpathrectangle{\pgfqpoint{1.109909in}{4.050000in}}{\pgfqpoint{5.400000in}{5.400000in}}%
\pgfusepath{clip}%
\pgfsetbuttcap%
\pgfsetroundjoin%
\pgfsetlinewidth{1.003750pt}%
\definecolor{currentstroke}{rgb}{0.541176,0.380392,0.658824}%
\pgfsetstrokecolor{currentstroke}%
\pgfsetstrokeopacity{0.900000}%
\pgfsetdash{}{0pt}%
\pgfpathmoveto{\pgfqpoint{4.384748in}{6.502656in}}%
\pgfpathcurveto{\pgfqpoint{4.395003in}{6.502656in}}{\pgfqpoint{4.404838in}{6.506730in}}{\pgfqpoint{4.412089in}{6.513981in}}%
\pgfpathcurveto{\pgfqpoint{4.419339in}{6.521231in}}{\pgfqpoint{4.423413in}{6.531067in}}{\pgfqpoint{4.423413in}{6.541321in}}%
\pgfpathcurveto{\pgfqpoint{4.423413in}{6.551575in}}{\pgfqpoint{4.419339in}{6.561411in}}{\pgfqpoint{4.412089in}{6.568661in}}%
\pgfpathcurveto{\pgfqpoint{4.404838in}{6.575912in}}{\pgfqpoint{4.395003in}{6.579986in}}{\pgfqpoint{4.384748in}{6.579986in}}%
\pgfpathcurveto{\pgfqpoint{4.374494in}{6.579986in}}{\pgfqpoint{4.364659in}{6.575912in}}{\pgfqpoint{4.357408in}{6.568661in}}%
\pgfpathcurveto{\pgfqpoint{4.350157in}{6.561411in}}{\pgfqpoint{4.346083in}{6.551575in}}{\pgfqpoint{4.346083in}{6.541321in}}%
\pgfpathcurveto{\pgfqpoint{4.346083in}{6.531067in}}{\pgfqpoint{4.350157in}{6.521231in}}{\pgfqpoint{4.357408in}{6.513981in}}%
\pgfpathcurveto{\pgfqpoint{4.364659in}{6.506730in}}{\pgfqpoint{4.374494in}{6.502656in}}{\pgfqpoint{4.384748in}{6.502656in}}%
\pgfpathlineto{\pgfqpoint{4.384748in}{6.502656in}}%
\pgfpathclose%
\pgfusepath{stroke}%
\end{pgfscope}%
\begin{pgfscope}%
\pgfpathrectangle{\pgfqpoint{1.109909in}{4.050000in}}{\pgfqpoint{5.400000in}{5.400000in}}%
\pgfusepath{clip}%
\pgfsetbuttcap%
\pgfsetroundjoin%
\pgfsetlinewidth{1.003750pt}%
\definecolor{currentstroke}{rgb}{0.541176,0.380392,0.658824}%
\pgfsetstrokecolor{currentstroke}%
\pgfsetstrokeopacity{0.900000}%
\pgfsetdash{}{0pt}%
\pgfpathmoveto{\pgfqpoint{3.242720in}{5.972222in}}%
\pgfpathcurveto{\pgfqpoint{3.252974in}{5.972222in}}{\pgfqpoint{3.262809in}{5.976296in}}{\pgfqpoint{3.270060in}{5.983547in}}%
\pgfpathcurveto{\pgfqpoint{3.277311in}{5.990798in}}{\pgfqpoint{3.281385in}{6.000633in}}{\pgfqpoint{3.281385in}{6.010887in}}%
\pgfpathcurveto{\pgfqpoint{3.281385in}{6.021141in}}{\pgfqpoint{3.277311in}{6.030977in}}{\pgfqpoint{3.270060in}{6.038227in}}%
\pgfpathcurveto{\pgfqpoint{3.262809in}{6.045478in}}{\pgfqpoint{3.252974in}{6.049552in}}{\pgfqpoint{3.242720in}{6.049552in}}%
\pgfpathcurveto{\pgfqpoint{3.232466in}{6.049552in}}{\pgfqpoint{3.222630in}{6.045478in}}{\pgfqpoint{3.215379in}{6.038227in}}%
\pgfpathcurveto{\pgfqpoint{3.208129in}{6.030977in}}{\pgfqpoint{3.204055in}{6.021141in}}{\pgfqpoint{3.204055in}{6.010887in}}%
\pgfpathcurveto{\pgfqpoint{3.204055in}{6.000633in}}{\pgfqpoint{3.208129in}{5.990798in}}{\pgfqpoint{3.215379in}{5.983547in}}%
\pgfpathcurveto{\pgfqpoint{3.222630in}{5.976296in}}{\pgfqpoint{3.232466in}{5.972222in}}{\pgfqpoint{3.242720in}{5.972222in}}%
\pgfpathlineto{\pgfqpoint{3.242720in}{5.972222in}}%
\pgfpathclose%
\pgfusepath{stroke}%
\end{pgfscope}%
\begin{pgfscope}%
\pgfpathrectangle{\pgfqpoint{1.109909in}{4.050000in}}{\pgfqpoint{5.400000in}{5.400000in}}%
\pgfusepath{clip}%
\pgfsetbuttcap%
\pgfsetroundjoin%
\pgfsetlinewidth{1.003750pt}%
\definecolor{currentstroke}{rgb}{0.541176,0.380392,0.658824}%
\pgfsetstrokecolor{currentstroke}%
\pgfsetstrokeopacity{0.900000}%
\pgfsetdash{}{0pt}%
\pgfpathmoveto{\pgfqpoint{4.005481in}{6.471551in}}%
\pgfpathcurveto{\pgfqpoint{4.015735in}{6.471551in}}{\pgfqpoint{4.025571in}{6.475625in}}{\pgfqpoint{4.032821in}{6.482876in}}%
\pgfpathcurveto{\pgfqpoint{4.040072in}{6.490127in}}{\pgfqpoint{4.044146in}{6.499962in}}{\pgfqpoint{4.044146in}{6.510216in}}%
\pgfpathcurveto{\pgfqpoint{4.044146in}{6.520470in}}{\pgfqpoint{4.040072in}{6.530306in}}{\pgfqpoint{4.032821in}{6.537557in}}%
\pgfpathcurveto{\pgfqpoint{4.025571in}{6.544807in}}{\pgfqpoint{4.015735in}{6.548881in}}{\pgfqpoint{4.005481in}{6.548881in}}%
\pgfpathcurveto{\pgfqpoint{3.995227in}{6.548881in}}{\pgfqpoint{3.985392in}{6.544807in}}{\pgfqpoint{3.978141in}{6.537557in}}%
\pgfpathcurveto{\pgfqpoint{3.970890in}{6.530306in}}{\pgfqpoint{3.966816in}{6.520470in}}{\pgfqpoint{3.966816in}{6.510216in}}%
\pgfpathcurveto{\pgfqpoint{3.966816in}{6.499962in}}{\pgfqpoint{3.970890in}{6.490127in}}{\pgfqpoint{3.978141in}{6.482876in}}%
\pgfpathcurveto{\pgfqpoint{3.985392in}{6.475625in}}{\pgfqpoint{3.995227in}{6.471551in}}{\pgfqpoint{4.005481in}{6.471551in}}%
\pgfpathlineto{\pgfqpoint{4.005481in}{6.471551in}}%
\pgfpathclose%
\pgfusepath{stroke}%
\end{pgfscope}%
\begin{pgfscope}%
\pgfpathrectangle{\pgfqpoint{1.109909in}{4.050000in}}{\pgfqpoint{5.400000in}{5.400000in}}%
\pgfusepath{clip}%
\pgfsetbuttcap%
\pgfsetroundjoin%
\pgfsetlinewidth{1.003750pt}%
\definecolor{currentstroke}{rgb}{0.541176,0.380392,0.658824}%
\pgfsetstrokecolor{currentstroke}%
\pgfsetstrokeopacity{0.900000}%
\pgfsetdash{}{0pt}%
\pgfpathmoveto{\pgfqpoint{4.644293in}{7.182666in}}%
\pgfpathcurveto{\pgfqpoint{4.654547in}{7.182666in}}{\pgfqpoint{4.664382in}{7.186740in}}{\pgfqpoint{4.671633in}{7.193991in}}%
\pgfpathcurveto{\pgfqpoint{4.678884in}{7.201242in}}{\pgfqpoint{4.682958in}{7.211077in}}{\pgfqpoint{4.682958in}{7.221331in}}%
\pgfpathcurveto{\pgfqpoint{4.682958in}{7.231585in}}{\pgfqpoint{4.678884in}{7.241421in}}{\pgfqpoint{4.671633in}{7.248672in}}%
\pgfpathcurveto{\pgfqpoint{4.664382in}{7.255922in}}{\pgfqpoint{4.654547in}{7.259996in}}{\pgfqpoint{4.644293in}{7.259996in}}%
\pgfpathcurveto{\pgfqpoint{4.634039in}{7.259996in}}{\pgfqpoint{4.624203in}{7.255922in}}{\pgfqpoint{4.616952in}{7.248672in}}%
\pgfpathcurveto{\pgfqpoint{4.609702in}{7.241421in}}{\pgfqpoint{4.605628in}{7.231585in}}{\pgfqpoint{4.605628in}{7.221331in}}%
\pgfpathcurveto{\pgfqpoint{4.605628in}{7.211077in}}{\pgfqpoint{4.609702in}{7.201242in}}{\pgfqpoint{4.616952in}{7.193991in}}%
\pgfpathcurveto{\pgfqpoint{4.624203in}{7.186740in}}{\pgfqpoint{4.634039in}{7.182666in}}{\pgfqpoint{4.644293in}{7.182666in}}%
\pgfpathlineto{\pgfqpoint{4.644293in}{7.182666in}}%
\pgfpathclose%
\pgfusepath{stroke}%
\end{pgfscope}%
\begin{pgfscope}%
\pgfpathrectangle{\pgfqpoint{1.109909in}{4.050000in}}{\pgfqpoint{5.400000in}{5.400000in}}%
\pgfusepath{clip}%
\pgfsetbuttcap%
\pgfsetroundjoin%
\pgfsetlinewidth{1.003750pt}%
\definecolor{currentstroke}{rgb}{0.541176,0.380392,0.658824}%
\pgfsetstrokecolor{currentstroke}%
\pgfsetstrokeopacity{0.900000}%
\pgfsetdash{}{0pt}%
\pgfpathmoveto{\pgfqpoint{3.628864in}{6.233386in}}%
\pgfpathcurveto{\pgfqpoint{3.639118in}{6.233386in}}{\pgfqpoint{3.648954in}{6.237460in}}{\pgfqpoint{3.656205in}{6.244710in}}%
\pgfpathcurveto{\pgfqpoint{3.663455in}{6.251961in}}{\pgfqpoint{3.667529in}{6.261796in}}{\pgfqpoint{3.667529in}{6.272051in}}%
\pgfpathcurveto{\pgfqpoint{3.667529in}{6.282305in}}{\pgfqpoint{3.663455in}{6.292140in}}{\pgfqpoint{3.656205in}{6.299391in}}%
\pgfpathcurveto{\pgfqpoint{3.648954in}{6.306642in}}{\pgfqpoint{3.639118in}{6.310716in}}{\pgfqpoint{3.628864in}{6.310716in}}%
\pgfpathcurveto{\pgfqpoint{3.618610in}{6.310716in}}{\pgfqpoint{3.608775in}{6.306642in}}{\pgfqpoint{3.601524in}{6.299391in}}%
\pgfpathcurveto{\pgfqpoint{3.594273in}{6.292140in}}{\pgfqpoint{3.590199in}{6.282305in}}{\pgfqpoint{3.590199in}{6.272051in}}%
\pgfpathcurveto{\pgfqpoint{3.590199in}{6.261796in}}{\pgfqpoint{3.594273in}{6.251961in}}{\pgfqpoint{3.601524in}{6.244710in}}%
\pgfpathcurveto{\pgfqpoint{3.608775in}{6.237460in}}{\pgfqpoint{3.618610in}{6.233386in}}{\pgfqpoint{3.628864in}{6.233386in}}%
\pgfpathlineto{\pgfqpoint{3.628864in}{6.233386in}}%
\pgfpathclose%
\pgfusepath{stroke}%
\end{pgfscope}%
\begin{pgfscope}%
\pgfpathrectangle{\pgfqpoint{1.109909in}{4.050000in}}{\pgfqpoint{5.400000in}{5.400000in}}%
\pgfusepath{clip}%
\pgfsetbuttcap%
\pgfsetroundjoin%
\pgfsetlinewidth{1.003750pt}%
\definecolor{currentstroke}{rgb}{0.541176,0.380392,0.658824}%
\pgfsetstrokecolor{currentstroke}%
\pgfsetstrokeopacity{0.900000}%
\pgfsetdash{}{0pt}%
\pgfpathmoveto{\pgfqpoint{5.491189in}{7.987202in}}%
\pgfpathcurveto{\pgfqpoint{5.501443in}{7.987202in}}{\pgfqpoint{5.511278in}{7.991276in}}{\pgfqpoint{5.518529in}{7.998526in}}%
\pgfpathcurveto{\pgfqpoint{5.525780in}{8.005777in}}{\pgfqpoint{5.529854in}{8.015613in}}{\pgfqpoint{5.529854in}{8.025867in}}%
\pgfpathcurveto{\pgfqpoint{5.529854in}{8.036121in}}{\pgfqpoint{5.525780in}{8.045956in}}{\pgfqpoint{5.518529in}{8.053207in}}%
\pgfpathcurveto{\pgfqpoint{5.511278in}{8.060458in}}{\pgfqpoint{5.501443in}{8.064532in}}{\pgfqpoint{5.491189in}{8.064532in}}%
\pgfpathcurveto{\pgfqpoint{5.480935in}{8.064532in}}{\pgfqpoint{5.471099in}{8.060458in}}{\pgfqpoint{5.463848in}{8.053207in}}%
\pgfpathcurveto{\pgfqpoint{5.456598in}{8.045956in}}{\pgfqpoint{5.452524in}{8.036121in}}{\pgfqpoint{5.452524in}{8.025867in}}%
\pgfpathcurveto{\pgfqpoint{5.452524in}{8.015613in}}{\pgfqpoint{5.456598in}{8.005777in}}{\pgfqpoint{5.463848in}{7.998526in}}%
\pgfpathcurveto{\pgfqpoint{5.471099in}{7.991276in}}{\pgfqpoint{5.480935in}{7.987202in}}{\pgfqpoint{5.491189in}{7.987202in}}%
\pgfpathlineto{\pgfqpoint{5.491189in}{7.987202in}}%
\pgfpathclose%
\pgfusepath{stroke}%
\end{pgfscope}%
\begin{pgfscope}%
\pgfpathrectangle{\pgfqpoint{1.109909in}{4.050000in}}{\pgfqpoint{5.400000in}{5.400000in}}%
\pgfusepath{clip}%
\pgfsetbuttcap%
\pgfsetroundjoin%
\pgfsetlinewidth{1.003750pt}%
\definecolor{currentstroke}{rgb}{0.541176,0.380392,0.658824}%
\pgfsetstrokecolor{currentstroke}%
\pgfsetstrokeopacity{0.900000}%
\pgfsetdash{}{0pt}%
\pgfpathmoveto{\pgfqpoint{4.826526in}{6.764565in}}%
\pgfpathcurveto{\pgfqpoint{4.836780in}{6.764565in}}{\pgfqpoint{4.846616in}{6.768639in}}{\pgfqpoint{4.853867in}{6.775890in}}%
\pgfpathcurveto{\pgfqpoint{4.861117in}{6.783141in}}{\pgfqpoint{4.865191in}{6.792976in}}{\pgfqpoint{4.865191in}{6.803230in}}%
\pgfpathcurveto{\pgfqpoint{4.865191in}{6.813484in}}{\pgfqpoint{4.861117in}{6.823320in}}{\pgfqpoint{4.853867in}{6.830571in}}%
\pgfpathcurveto{\pgfqpoint{4.846616in}{6.837821in}}{\pgfqpoint{4.836780in}{6.841895in}}{\pgfqpoint{4.826526in}{6.841895in}}%
\pgfpathcurveto{\pgfqpoint{4.816272in}{6.841895in}}{\pgfqpoint{4.806437in}{6.837821in}}{\pgfqpoint{4.799186in}{6.830571in}}%
\pgfpathcurveto{\pgfqpoint{4.791935in}{6.823320in}}{\pgfqpoint{4.787861in}{6.813484in}}{\pgfqpoint{4.787861in}{6.803230in}}%
\pgfpathcurveto{\pgfqpoint{4.787861in}{6.792976in}}{\pgfqpoint{4.791935in}{6.783141in}}{\pgfqpoint{4.799186in}{6.775890in}}%
\pgfpathcurveto{\pgfqpoint{4.806437in}{6.768639in}}{\pgfqpoint{4.816272in}{6.764565in}}{\pgfqpoint{4.826526in}{6.764565in}}%
\pgfpathlineto{\pgfqpoint{4.826526in}{6.764565in}}%
\pgfpathclose%
\pgfusepath{stroke}%
\end{pgfscope}%
\begin{pgfscope}%
\pgfpathrectangle{\pgfqpoint{1.109909in}{4.050000in}}{\pgfqpoint{5.400000in}{5.400000in}}%
\pgfusepath{clip}%
\pgfsetbuttcap%
\pgfsetroundjoin%
\pgfsetlinewidth{1.003750pt}%
\definecolor{currentstroke}{rgb}{0.541176,0.380392,0.658824}%
\pgfsetstrokecolor{currentstroke}%
\pgfsetstrokeopacity{0.900000}%
\pgfsetdash{}{0pt}%
\pgfpathmoveto{\pgfqpoint{4.469512in}{7.864920in}}%
\pgfpathcurveto{\pgfqpoint{4.479766in}{7.864920in}}{\pgfqpoint{4.489602in}{7.868994in}}{\pgfqpoint{4.496853in}{7.876244in}}%
\pgfpathcurveto{\pgfqpoint{4.504103in}{7.883495in}}{\pgfqpoint{4.508177in}{7.893331in}}{\pgfqpoint{4.508177in}{7.903585in}}%
\pgfpathcurveto{\pgfqpoint{4.508177in}{7.913839in}}{\pgfqpoint{4.504103in}{7.923674in}}{\pgfqpoint{4.496853in}{7.930925in}}%
\pgfpathcurveto{\pgfqpoint{4.489602in}{7.938176in}}{\pgfqpoint{4.479766in}{7.942250in}}{\pgfqpoint{4.469512in}{7.942250in}}%
\pgfpathcurveto{\pgfqpoint{4.459258in}{7.942250in}}{\pgfqpoint{4.449423in}{7.938176in}}{\pgfqpoint{4.442172in}{7.930925in}}%
\pgfpathcurveto{\pgfqpoint{4.434921in}{7.923674in}}{\pgfqpoint{4.430847in}{7.913839in}}{\pgfqpoint{4.430847in}{7.903585in}}%
\pgfpathcurveto{\pgfqpoint{4.430847in}{7.893331in}}{\pgfqpoint{4.434921in}{7.883495in}}{\pgfqpoint{4.442172in}{7.876244in}}%
\pgfpathcurveto{\pgfqpoint{4.449423in}{7.868994in}}{\pgfqpoint{4.459258in}{7.864920in}}{\pgfqpoint{4.469512in}{7.864920in}}%
\pgfpathlineto{\pgfqpoint{4.469512in}{7.864920in}}%
\pgfpathclose%
\pgfusepath{stroke}%
\end{pgfscope}%
\begin{pgfscope}%
\pgfpathrectangle{\pgfqpoint{1.109909in}{4.050000in}}{\pgfqpoint{5.400000in}{5.400000in}}%
\pgfusepath{clip}%
\pgfsetbuttcap%
\pgfsetroundjoin%
\pgfsetlinewidth{1.003750pt}%
\definecolor{currentstroke}{rgb}{0.541176,0.380392,0.658824}%
\pgfsetstrokecolor{currentstroke}%
\pgfsetstrokeopacity{0.900000}%
\pgfsetdash{}{0pt}%
\pgfpathmoveto{\pgfqpoint{3.446414in}{6.953505in}}%
\pgfpathcurveto{\pgfqpoint{3.456669in}{6.953505in}}{\pgfqpoint{3.466504in}{6.957579in}}{\pgfqpoint{3.473755in}{6.964830in}}%
\pgfpathcurveto{\pgfqpoint{3.481006in}{6.972081in}}{\pgfqpoint{3.485079in}{6.981916in}}{\pgfqpoint{3.485079in}{6.992170in}}%
\pgfpathcurveto{\pgfqpoint{3.485079in}{7.002424in}}{\pgfqpoint{3.481006in}{7.012260in}}{\pgfqpoint{3.473755in}{7.019511in}}%
\pgfpathcurveto{\pgfqpoint{3.466504in}{7.026761in}}{\pgfqpoint{3.456669in}{7.030835in}}{\pgfqpoint{3.446414in}{7.030835in}}%
\pgfpathcurveto{\pgfqpoint{3.436160in}{7.030835in}}{\pgfqpoint{3.426325in}{7.026761in}}{\pgfqpoint{3.419074in}{7.019511in}}%
\pgfpathcurveto{\pgfqpoint{3.411823in}{7.012260in}}{\pgfqpoint{3.407749in}{7.002424in}}{\pgfqpoint{3.407749in}{6.992170in}}%
\pgfpathcurveto{\pgfqpoint{3.407749in}{6.981916in}}{\pgfqpoint{3.411823in}{6.972081in}}{\pgfqpoint{3.419074in}{6.964830in}}%
\pgfpathcurveto{\pgfqpoint{3.426325in}{6.957579in}}{\pgfqpoint{3.436160in}{6.953505in}}{\pgfqpoint{3.446414in}{6.953505in}}%
\pgfpathlineto{\pgfqpoint{3.446414in}{6.953505in}}%
\pgfpathclose%
\pgfusepath{stroke}%
\end{pgfscope}%
\begin{pgfscope}%
\pgfpathrectangle{\pgfqpoint{1.109909in}{4.050000in}}{\pgfqpoint{5.400000in}{5.400000in}}%
\pgfusepath{clip}%
\pgfsetbuttcap%
\pgfsetroundjoin%
\pgfsetlinewidth{1.003750pt}%
\definecolor{currentstroke}{rgb}{0.541176,0.380392,0.658824}%
\pgfsetstrokecolor{currentstroke}%
\pgfsetstrokeopacity{0.900000}%
\pgfsetdash{}{0pt}%
\pgfpathmoveto{\pgfqpoint{5.352256in}{7.453373in}}%
\pgfpathcurveto{\pgfqpoint{5.362510in}{7.453373in}}{\pgfqpoint{5.372345in}{7.457447in}}{\pgfqpoint{5.379596in}{7.464698in}}%
\pgfpathcurveto{\pgfqpoint{5.386847in}{7.471949in}}{\pgfqpoint{5.390921in}{7.481784in}}{\pgfqpoint{5.390921in}{7.492038in}}%
\pgfpathcurveto{\pgfqpoint{5.390921in}{7.502292in}}{\pgfqpoint{5.386847in}{7.512128in}}{\pgfqpoint{5.379596in}{7.519379in}}%
\pgfpathcurveto{\pgfqpoint{5.372345in}{7.526629in}}{\pgfqpoint{5.362510in}{7.530703in}}{\pgfqpoint{5.352256in}{7.530703in}}%
\pgfpathcurveto{\pgfqpoint{5.342002in}{7.530703in}}{\pgfqpoint{5.332166in}{7.526629in}}{\pgfqpoint{5.324915in}{7.519379in}}%
\pgfpathcurveto{\pgfqpoint{5.317665in}{7.512128in}}{\pgfqpoint{5.313591in}{7.502292in}}{\pgfqpoint{5.313591in}{7.492038in}}%
\pgfpathcurveto{\pgfqpoint{5.313591in}{7.481784in}}{\pgfqpoint{5.317665in}{7.471949in}}{\pgfqpoint{5.324915in}{7.464698in}}%
\pgfpathcurveto{\pgfqpoint{5.332166in}{7.457447in}}{\pgfqpoint{5.342002in}{7.453373in}}{\pgfqpoint{5.352256in}{7.453373in}}%
\pgfpathlineto{\pgfqpoint{5.352256in}{7.453373in}}%
\pgfpathclose%
\pgfusepath{stroke}%
\end{pgfscope}%
\begin{pgfscope}%
\pgfpathrectangle{\pgfqpoint{1.109909in}{4.050000in}}{\pgfqpoint{5.400000in}{5.400000in}}%
\pgfusepath{clip}%
\pgfsetbuttcap%
\pgfsetroundjoin%
\pgfsetlinewidth{1.003750pt}%
\definecolor{currentstroke}{rgb}{0.541176,0.380392,0.658824}%
\pgfsetstrokecolor{currentstroke}%
\pgfsetstrokeopacity{0.900000}%
\pgfsetdash{}{0pt}%
\pgfpathmoveto{\pgfqpoint{3.756043in}{7.147343in}}%
\pgfpathcurveto{\pgfqpoint{3.766297in}{7.147343in}}{\pgfqpoint{3.776132in}{7.151417in}}{\pgfqpoint{3.783383in}{7.158668in}}%
\pgfpathcurveto{\pgfqpoint{3.790634in}{7.165919in}}{\pgfqpoint{3.794708in}{7.175754in}}{\pgfqpoint{3.794708in}{7.186008in}}%
\pgfpathcurveto{\pgfqpoint{3.794708in}{7.196262in}}{\pgfqpoint{3.790634in}{7.206098in}}{\pgfqpoint{3.783383in}{7.213349in}}%
\pgfpathcurveto{\pgfqpoint{3.776132in}{7.220599in}}{\pgfqpoint{3.766297in}{7.224673in}}{\pgfqpoint{3.756043in}{7.224673in}}%
\pgfpathcurveto{\pgfqpoint{3.745789in}{7.224673in}}{\pgfqpoint{3.735953in}{7.220599in}}{\pgfqpoint{3.728702in}{7.213349in}}%
\pgfpathcurveto{\pgfqpoint{3.721452in}{7.206098in}}{\pgfqpoint{3.717378in}{7.196262in}}{\pgfqpoint{3.717378in}{7.186008in}}%
\pgfpathcurveto{\pgfqpoint{3.717378in}{7.175754in}}{\pgfqpoint{3.721452in}{7.165919in}}{\pgfqpoint{3.728702in}{7.158668in}}%
\pgfpathcurveto{\pgfqpoint{3.735953in}{7.151417in}}{\pgfqpoint{3.745789in}{7.147343in}}{\pgfqpoint{3.756043in}{7.147343in}}%
\pgfpathlineto{\pgfqpoint{3.756043in}{7.147343in}}%
\pgfpathclose%
\pgfusepath{stroke}%
\end{pgfscope}%
\begin{pgfscope}%
\pgfpathrectangle{\pgfqpoint{1.109909in}{4.050000in}}{\pgfqpoint{5.400000in}{5.400000in}}%
\pgfusepath{clip}%
\pgfsetbuttcap%
\pgfsetroundjoin%
\pgfsetlinewidth{1.003750pt}%
\definecolor{currentstroke}{rgb}{0.541176,0.380392,0.658824}%
\pgfsetstrokecolor{currentstroke}%
\pgfsetstrokeopacity{0.900000}%
\pgfsetdash{}{0pt}%
\pgfpathmoveto{\pgfqpoint{3.248845in}{6.127868in}}%
\pgfpathcurveto{\pgfqpoint{3.259099in}{6.127868in}}{\pgfqpoint{3.268935in}{6.131942in}}{\pgfqpoint{3.276185in}{6.139192in}}%
\pgfpathcurveto{\pgfqpoint{3.283436in}{6.146443in}}{\pgfqpoint{3.287510in}{6.156278in}}{\pgfqpoint{3.287510in}{6.166533in}}%
\pgfpathcurveto{\pgfqpoint{3.287510in}{6.176787in}}{\pgfqpoint{3.283436in}{6.186622in}}{\pgfqpoint{3.276185in}{6.193873in}}%
\pgfpathcurveto{\pgfqpoint{3.268935in}{6.201124in}}{\pgfqpoint{3.259099in}{6.205198in}}{\pgfqpoint{3.248845in}{6.205198in}}%
\pgfpathcurveto{\pgfqpoint{3.238591in}{6.205198in}}{\pgfqpoint{3.228755in}{6.201124in}}{\pgfqpoint{3.221505in}{6.193873in}}%
\pgfpathcurveto{\pgfqpoint{3.214254in}{6.186622in}}{\pgfqpoint{3.210180in}{6.176787in}}{\pgfqpoint{3.210180in}{6.166533in}}%
\pgfpathcurveto{\pgfqpoint{3.210180in}{6.156278in}}{\pgfqpoint{3.214254in}{6.146443in}}{\pgfqpoint{3.221505in}{6.139192in}}%
\pgfpathcurveto{\pgfqpoint{3.228755in}{6.131942in}}{\pgfqpoint{3.238591in}{6.127868in}}{\pgfqpoint{3.248845in}{6.127868in}}%
\pgfpathlineto{\pgfqpoint{3.248845in}{6.127868in}}%
\pgfpathclose%
\pgfusepath{stroke}%
\end{pgfscope}%
\begin{pgfscope}%
\pgfpathrectangle{\pgfqpoint{1.109909in}{4.050000in}}{\pgfqpoint{5.400000in}{5.400000in}}%
\pgfusepath{clip}%
\pgfsetbuttcap%
\pgfsetroundjoin%
\pgfsetlinewidth{1.003750pt}%
\definecolor{currentstroke}{rgb}{0.541176,0.380392,0.658824}%
\pgfsetstrokecolor{currentstroke}%
\pgfsetstrokeopacity{0.900000}%
\pgfsetdash{}{0pt}%
\pgfpathmoveto{\pgfqpoint{3.890069in}{6.567410in}}%
\pgfpathcurveto{\pgfqpoint{3.900323in}{6.567410in}}{\pgfqpoint{3.910158in}{6.571484in}}{\pgfqpoint{3.917409in}{6.578735in}}%
\pgfpathcurveto{\pgfqpoint{3.924660in}{6.585985in}}{\pgfqpoint{3.928734in}{6.595821in}}{\pgfqpoint{3.928734in}{6.606075in}}%
\pgfpathcurveto{\pgfqpoint{3.928734in}{6.616329in}}{\pgfqpoint{3.924660in}{6.626165in}}{\pgfqpoint{3.917409in}{6.633415in}}%
\pgfpathcurveto{\pgfqpoint{3.910158in}{6.640666in}}{\pgfqpoint{3.900323in}{6.644740in}}{\pgfqpoint{3.890069in}{6.644740in}}%
\pgfpathcurveto{\pgfqpoint{3.879815in}{6.644740in}}{\pgfqpoint{3.869979in}{6.640666in}}{\pgfqpoint{3.862728in}{6.633415in}}%
\pgfpathcurveto{\pgfqpoint{3.855478in}{6.626165in}}{\pgfqpoint{3.851404in}{6.616329in}}{\pgfqpoint{3.851404in}{6.606075in}}%
\pgfpathcurveto{\pgfqpoint{3.851404in}{6.595821in}}{\pgfqpoint{3.855478in}{6.585985in}}{\pgfqpoint{3.862728in}{6.578735in}}%
\pgfpathcurveto{\pgfqpoint{3.869979in}{6.571484in}}{\pgfqpoint{3.879815in}{6.567410in}}{\pgfqpoint{3.890069in}{6.567410in}}%
\pgfpathlineto{\pgfqpoint{3.890069in}{6.567410in}}%
\pgfpathclose%
\pgfusepath{stroke}%
\end{pgfscope}%
\begin{pgfscope}%
\pgfpathrectangle{\pgfqpoint{1.109909in}{4.050000in}}{\pgfqpoint{5.400000in}{5.400000in}}%
\pgfusepath{clip}%
\pgfsetbuttcap%
\pgfsetroundjoin%
\pgfsetlinewidth{1.003750pt}%
\definecolor{currentstroke}{rgb}{0.541176,0.380392,0.658824}%
\pgfsetstrokecolor{currentstroke}%
\pgfsetstrokeopacity{0.900000}%
\pgfsetdash{}{0pt}%
\pgfpathmoveto{\pgfqpoint{4.512269in}{7.475585in}}%
\pgfpathcurveto{\pgfqpoint{4.522523in}{7.475585in}}{\pgfqpoint{4.532358in}{7.479659in}}{\pgfqpoint{4.539609in}{7.486910in}}%
\pgfpathcurveto{\pgfqpoint{4.546860in}{7.494161in}}{\pgfqpoint{4.550934in}{7.503996in}}{\pgfqpoint{4.550934in}{7.514250in}}%
\pgfpathcurveto{\pgfqpoint{4.550934in}{7.524504in}}{\pgfqpoint{4.546860in}{7.534340in}}{\pgfqpoint{4.539609in}{7.541591in}}%
\pgfpathcurveto{\pgfqpoint{4.532358in}{7.548841in}}{\pgfqpoint{4.522523in}{7.552915in}}{\pgfqpoint{4.512269in}{7.552915in}}%
\pgfpathcurveto{\pgfqpoint{4.502015in}{7.552915in}}{\pgfqpoint{4.492179in}{7.548841in}}{\pgfqpoint{4.484928in}{7.541591in}}%
\pgfpathcurveto{\pgfqpoint{4.477678in}{7.534340in}}{\pgfqpoint{4.473604in}{7.524504in}}{\pgfqpoint{4.473604in}{7.514250in}}%
\pgfpathcurveto{\pgfqpoint{4.473604in}{7.503996in}}{\pgfqpoint{4.477678in}{7.494161in}}{\pgfqpoint{4.484928in}{7.486910in}}%
\pgfpathcurveto{\pgfqpoint{4.492179in}{7.479659in}}{\pgfqpoint{4.502015in}{7.475585in}}{\pgfqpoint{4.512269in}{7.475585in}}%
\pgfpathlineto{\pgfqpoint{4.512269in}{7.475585in}}%
\pgfpathclose%
\pgfusepath{stroke}%
\end{pgfscope}%
\begin{pgfscope}%
\pgfpathrectangle{\pgfqpoint{1.109909in}{4.050000in}}{\pgfqpoint{5.400000in}{5.400000in}}%
\pgfusepath{clip}%
\pgfsetbuttcap%
\pgfsetroundjoin%
\pgfsetlinewidth{1.003750pt}%
\definecolor{currentstroke}{rgb}{0.541176,0.380392,0.658824}%
\pgfsetstrokecolor{currentstroke}%
\pgfsetstrokeopacity{0.900000}%
\pgfsetdash{}{0pt}%
\pgfpathmoveto{\pgfqpoint{4.979839in}{6.591059in}}%
\pgfpathcurveto{\pgfqpoint{4.990093in}{6.591059in}}{\pgfqpoint{4.999929in}{6.595133in}}{\pgfqpoint{5.007179in}{6.602384in}}%
\pgfpathcurveto{\pgfqpoint{5.014430in}{6.609635in}}{\pgfqpoint{5.018504in}{6.619470in}}{\pgfqpoint{5.018504in}{6.629725in}}%
\pgfpathcurveto{\pgfqpoint{5.018504in}{6.639979in}}{\pgfqpoint{5.014430in}{6.649814in}}{\pgfqpoint{5.007179in}{6.657065in}}%
\pgfpathcurveto{\pgfqpoint{4.999929in}{6.664316in}}{\pgfqpoint{4.990093in}{6.668390in}}{\pgfqpoint{4.979839in}{6.668390in}}%
\pgfpathcurveto{\pgfqpoint{4.969585in}{6.668390in}}{\pgfqpoint{4.959749in}{6.664316in}}{\pgfqpoint{4.952499in}{6.657065in}}%
\pgfpathcurveto{\pgfqpoint{4.945248in}{6.649814in}}{\pgfqpoint{4.941174in}{6.639979in}}{\pgfqpoint{4.941174in}{6.629725in}}%
\pgfpathcurveto{\pgfqpoint{4.941174in}{6.619470in}}{\pgfqpoint{4.945248in}{6.609635in}}{\pgfqpoint{4.952499in}{6.602384in}}%
\pgfpathcurveto{\pgfqpoint{4.959749in}{6.595133in}}{\pgfqpoint{4.969585in}{6.591059in}}{\pgfqpoint{4.979839in}{6.591059in}}%
\pgfpathlineto{\pgfqpoint{4.979839in}{6.591059in}}%
\pgfpathclose%
\pgfusepath{stroke}%
\end{pgfscope}%
\begin{pgfscope}%
\pgfpathrectangle{\pgfqpoint{1.109909in}{4.050000in}}{\pgfqpoint{5.400000in}{5.400000in}}%
\pgfusepath{clip}%
\pgfsetbuttcap%
\pgfsetroundjoin%
\pgfsetlinewidth{1.003750pt}%
\definecolor{currentstroke}{rgb}{0.541176,0.380392,0.658824}%
\pgfsetstrokecolor{currentstroke}%
\pgfsetstrokeopacity{0.900000}%
\pgfsetdash{}{0pt}%
\pgfpathmoveto{\pgfqpoint{4.767453in}{6.676501in}}%
\pgfpathcurveto{\pgfqpoint{4.777707in}{6.676501in}}{\pgfqpoint{4.787542in}{6.680575in}}{\pgfqpoint{4.794793in}{6.687826in}}%
\pgfpathcurveto{\pgfqpoint{4.802044in}{6.695077in}}{\pgfqpoint{4.806118in}{6.704912in}}{\pgfqpoint{4.806118in}{6.715166in}}%
\pgfpathcurveto{\pgfqpoint{4.806118in}{6.725420in}}{\pgfqpoint{4.802044in}{6.735256in}}{\pgfqpoint{4.794793in}{6.742507in}}%
\pgfpathcurveto{\pgfqpoint{4.787542in}{6.749757in}}{\pgfqpoint{4.777707in}{6.753831in}}{\pgfqpoint{4.767453in}{6.753831in}}%
\pgfpathcurveto{\pgfqpoint{4.757199in}{6.753831in}}{\pgfqpoint{4.747363in}{6.749757in}}{\pgfqpoint{4.740113in}{6.742507in}}%
\pgfpathcurveto{\pgfqpoint{4.732862in}{6.735256in}}{\pgfqpoint{4.728788in}{6.725420in}}{\pgfqpoint{4.728788in}{6.715166in}}%
\pgfpathcurveto{\pgfqpoint{4.728788in}{6.704912in}}{\pgfqpoint{4.732862in}{6.695077in}}{\pgfqpoint{4.740113in}{6.687826in}}%
\pgfpathcurveto{\pgfqpoint{4.747363in}{6.680575in}}{\pgfqpoint{4.757199in}{6.676501in}}{\pgfqpoint{4.767453in}{6.676501in}}%
\pgfpathlineto{\pgfqpoint{4.767453in}{6.676501in}}%
\pgfpathclose%
\pgfusepath{stroke}%
\end{pgfscope}%
\begin{pgfscope}%
\pgfpathrectangle{\pgfqpoint{1.109909in}{4.050000in}}{\pgfqpoint{5.400000in}{5.400000in}}%
\pgfusepath{clip}%
\pgfsetbuttcap%
\pgfsetroundjoin%
\pgfsetlinewidth{1.003750pt}%
\definecolor{currentstroke}{rgb}{0.541176,0.380392,0.658824}%
\pgfsetstrokecolor{currentstroke}%
\pgfsetstrokeopacity{0.900000}%
\pgfsetdash{}{0pt}%
\pgfpathmoveto{\pgfqpoint{5.060910in}{8.007266in}}%
\pgfpathcurveto{\pgfqpoint{5.071164in}{8.007266in}}{\pgfqpoint{5.081000in}{8.011340in}}{\pgfqpoint{5.088250in}{8.018591in}}%
\pgfpathcurveto{\pgfqpoint{5.095501in}{8.025841in}}{\pgfqpoint{5.099575in}{8.035677in}}{\pgfqpoint{5.099575in}{8.045931in}}%
\pgfpathcurveto{\pgfqpoint{5.099575in}{8.056185in}}{\pgfqpoint{5.095501in}{8.066020in}}{\pgfqpoint{5.088250in}{8.073271in}}%
\pgfpathcurveto{\pgfqpoint{5.081000in}{8.080522in}}{\pgfqpoint{5.071164in}{8.084596in}}{\pgfqpoint{5.060910in}{8.084596in}}%
\pgfpathcurveto{\pgfqpoint{5.050656in}{8.084596in}}{\pgfqpoint{5.040820in}{8.080522in}}{\pgfqpoint{5.033570in}{8.073271in}}%
\pgfpathcurveto{\pgfqpoint{5.026319in}{8.066020in}}{\pgfqpoint{5.022245in}{8.056185in}}{\pgfqpoint{5.022245in}{8.045931in}}%
\pgfpathcurveto{\pgfqpoint{5.022245in}{8.035677in}}{\pgfqpoint{5.026319in}{8.025841in}}{\pgfqpoint{5.033570in}{8.018591in}}%
\pgfpathcurveto{\pgfqpoint{5.040820in}{8.011340in}}{\pgfqpoint{5.050656in}{8.007266in}}{\pgfqpoint{5.060910in}{8.007266in}}%
\pgfpathlineto{\pgfqpoint{5.060910in}{8.007266in}}%
\pgfpathclose%
\pgfusepath{stroke}%
\end{pgfscope}%
\begin{pgfscope}%
\pgfpathrectangle{\pgfqpoint{1.109909in}{4.050000in}}{\pgfqpoint{5.400000in}{5.400000in}}%
\pgfusepath{clip}%
\pgfsetbuttcap%
\pgfsetroundjoin%
\pgfsetlinewidth{1.003750pt}%
\definecolor{currentstroke}{rgb}{0.541176,0.380392,0.658824}%
\pgfsetstrokecolor{currentstroke}%
\pgfsetstrokeopacity{0.900000}%
\pgfsetdash{}{0pt}%
\pgfpathmoveto{\pgfqpoint{4.495274in}{7.442201in}}%
\pgfpathcurveto{\pgfqpoint{4.505528in}{7.442201in}}{\pgfqpoint{4.515363in}{7.446275in}}{\pgfqpoint{4.522614in}{7.453526in}}%
\pgfpathcurveto{\pgfqpoint{4.529865in}{7.460777in}}{\pgfqpoint{4.533939in}{7.470612in}}{\pgfqpoint{4.533939in}{7.480866in}}%
\pgfpathcurveto{\pgfqpoint{4.533939in}{7.491120in}}{\pgfqpoint{4.529865in}{7.500956in}}{\pgfqpoint{4.522614in}{7.508206in}}%
\pgfpathcurveto{\pgfqpoint{4.515363in}{7.515457in}}{\pgfqpoint{4.505528in}{7.519531in}}{\pgfqpoint{4.495274in}{7.519531in}}%
\pgfpathcurveto{\pgfqpoint{4.485020in}{7.519531in}}{\pgfqpoint{4.475184in}{7.515457in}}{\pgfqpoint{4.467933in}{7.508206in}}%
\pgfpathcurveto{\pgfqpoint{4.460683in}{7.500956in}}{\pgfqpoint{4.456609in}{7.491120in}}{\pgfqpoint{4.456609in}{7.480866in}}%
\pgfpathcurveto{\pgfqpoint{4.456609in}{7.470612in}}{\pgfqpoint{4.460683in}{7.460777in}}{\pgfqpoint{4.467933in}{7.453526in}}%
\pgfpathcurveto{\pgfqpoint{4.475184in}{7.446275in}}{\pgfqpoint{4.485020in}{7.442201in}}{\pgfqpoint{4.495274in}{7.442201in}}%
\pgfpathlineto{\pgfqpoint{4.495274in}{7.442201in}}%
\pgfpathclose%
\pgfusepath{stroke}%
\end{pgfscope}%
\begin{pgfscope}%
\pgfpathrectangle{\pgfqpoint{1.109909in}{4.050000in}}{\pgfqpoint{5.400000in}{5.400000in}}%
\pgfusepath{clip}%
\pgfsetbuttcap%
\pgfsetroundjoin%
\pgfsetlinewidth{1.003750pt}%
\definecolor{currentstroke}{rgb}{0.541176,0.380392,0.658824}%
\pgfsetstrokecolor{currentstroke}%
\pgfsetstrokeopacity{0.900000}%
\pgfsetdash{}{0pt}%
\pgfpathmoveto{\pgfqpoint{5.189977in}{8.151443in}}%
\pgfpathcurveto{\pgfqpoint{5.200231in}{8.151443in}}{\pgfqpoint{5.210067in}{8.155517in}}{\pgfqpoint{5.217317in}{8.162768in}}%
\pgfpathcurveto{\pgfqpoint{5.224568in}{8.170018in}}{\pgfqpoint{5.228642in}{8.179854in}}{\pgfqpoint{5.228642in}{8.190108in}}%
\pgfpathcurveto{\pgfqpoint{5.228642in}{8.200362in}}{\pgfqpoint{5.224568in}{8.210198in}}{\pgfqpoint{5.217317in}{8.217448in}}%
\pgfpathcurveto{\pgfqpoint{5.210067in}{8.224699in}}{\pgfqpoint{5.200231in}{8.228773in}}{\pgfqpoint{5.189977in}{8.228773in}}%
\pgfpathcurveto{\pgfqpoint{5.179723in}{8.228773in}}{\pgfqpoint{5.169887in}{8.224699in}}{\pgfqpoint{5.162637in}{8.217448in}}%
\pgfpathcurveto{\pgfqpoint{5.155386in}{8.210198in}}{\pgfqpoint{5.151312in}{8.200362in}}{\pgfqpoint{5.151312in}{8.190108in}}%
\pgfpathcurveto{\pgfqpoint{5.151312in}{8.179854in}}{\pgfqpoint{5.155386in}{8.170018in}}{\pgfqpoint{5.162637in}{8.162768in}}%
\pgfpathcurveto{\pgfqpoint{5.169887in}{8.155517in}}{\pgfqpoint{5.179723in}{8.151443in}}{\pgfqpoint{5.189977in}{8.151443in}}%
\pgfpathlineto{\pgfqpoint{5.189977in}{8.151443in}}%
\pgfpathclose%
\pgfusepath{stroke}%
\end{pgfscope}%
\begin{pgfscope}%
\pgfpathrectangle{\pgfqpoint{1.109909in}{4.050000in}}{\pgfqpoint{5.400000in}{5.400000in}}%
\pgfusepath{clip}%
\pgfsetbuttcap%
\pgfsetroundjoin%
\pgfsetlinewidth{1.003750pt}%
\definecolor{currentstroke}{rgb}{0.541176,0.380392,0.658824}%
\pgfsetstrokecolor{currentstroke}%
\pgfsetstrokeopacity{0.900000}%
\pgfsetdash{}{0pt}%
\pgfpathmoveto{\pgfqpoint{3.763222in}{5.608904in}}%
\pgfpathcurveto{\pgfqpoint{3.773476in}{5.608904in}}{\pgfqpoint{3.783311in}{5.612978in}}{\pgfqpoint{3.790562in}{5.620229in}}%
\pgfpathcurveto{\pgfqpoint{3.797813in}{5.627480in}}{\pgfqpoint{3.801887in}{5.637315in}}{\pgfqpoint{3.801887in}{5.647569in}}%
\pgfpathcurveto{\pgfqpoint{3.801887in}{5.657824in}}{\pgfqpoint{3.797813in}{5.667659in}}{\pgfqpoint{3.790562in}{5.674910in}}%
\pgfpathcurveto{\pgfqpoint{3.783311in}{5.682161in}}{\pgfqpoint{3.773476in}{5.686235in}}{\pgfqpoint{3.763222in}{5.686235in}}%
\pgfpathcurveto{\pgfqpoint{3.752968in}{5.686235in}}{\pgfqpoint{3.743132in}{5.682161in}}{\pgfqpoint{3.735881in}{5.674910in}}%
\pgfpathcurveto{\pgfqpoint{3.728631in}{5.667659in}}{\pgfqpoint{3.724557in}{5.657824in}}{\pgfqpoint{3.724557in}{5.647569in}}%
\pgfpathcurveto{\pgfqpoint{3.724557in}{5.637315in}}{\pgfqpoint{3.728631in}{5.627480in}}{\pgfqpoint{3.735881in}{5.620229in}}%
\pgfpathcurveto{\pgfqpoint{3.743132in}{5.612978in}}{\pgfqpoint{3.752968in}{5.608904in}}{\pgfqpoint{3.763222in}{5.608904in}}%
\pgfpathlineto{\pgfqpoint{3.763222in}{5.608904in}}%
\pgfpathclose%
\pgfusepath{stroke}%
\end{pgfscope}%
\begin{pgfscope}%
\pgfpathrectangle{\pgfqpoint{1.109909in}{4.050000in}}{\pgfqpoint{5.400000in}{5.400000in}}%
\pgfusepath{clip}%
\pgfsetbuttcap%
\pgfsetroundjoin%
\pgfsetlinewidth{1.003750pt}%
\definecolor{currentstroke}{rgb}{0.541176,0.380392,0.658824}%
\pgfsetstrokecolor{currentstroke}%
\pgfsetstrokeopacity{0.900000}%
\pgfsetdash{}{0pt}%
\pgfpathmoveto{\pgfqpoint{4.667432in}{7.392531in}}%
\pgfpathcurveto{\pgfqpoint{4.677686in}{7.392531in}}{\pgfqpoint{4.687522in}{7.396605in}}{\pgfqpoint{4.694772in}{7.403855in}}%
\pgfpathcurveto{\pgfqpoint{4.702023in}{7.411106in}}{\pgfqpoint{4.706097in}{7.420942in}}{\pgfqpoint{4.706097in}{7.431196in}}%
\pgfpathcurveto{\pgfqpoint{4.706097in}{7.441450in}}{\pgfqpoint{4.702023in}{7.451285in}}{\pgfqpoint{4.694772in}{7.458536in}}%
\pgfpathcurveto{\pgfqpoint{4.687522in}{7.465787in}}{\pgfqpoint{4.677686in}{7.469861in}}{\pgfqpoint{4.667432in}{7.469861in}}%
\pgfpathcurveto{\pgfqpoint{4.657178in}{7.469861in}}{\pgfqpoint{4.647342in}{7.465787in}}{\pgfqpoint{4.640092in}{7.458536in}}%
\pgfpathcurveto{\pgfqpoint{4.632841in}{7.451285in}}{\pgfqpoint{4.628767in}{7.441450in}}{\pgfqpoint{4.628767in}{7.431196in}}%
\pgfpathcurveto{\pgfqpoint{4.628767in}{7.420942in}}{\pgfqpoint{4.632841in}{7.411106in}}{\pgfqpoint{4.640092in}{7.403855in}}%
\pgfpathcurveto{\pgfqpoint{4.647342in}{7.396605in}}{\pgfqpoint{4.657178in}{7.392531in}}{\pgfqpoint{4.667432in}{7.392531in}}%
\pgfpathlineto{\pgfqpoint{4.667432in}{7.392531in}}%
\pgfpathclose%
\pgfusepath{stroke}%
\end{pgfscope}%
\begin{pgfscope}%
\pgfsetbuttcap%
\pgfsetroundjoin%
\pgfsetlinewidth{1.003750pt}%
\definecolor{currentstroke}{rgb}{0.541176,0.380392,0.658824}%
\pgfsetstrokecolor{currentstroke}%
\pgfsetstrokeopacity{0.900000}%
\pgfsetdash{}{0pt}%
\pgfpathmoveto{\pgfqpoint{6.509909in}{9.390262in}}%
\pgfpathlineto{\pgfqpoint{6.569647in}{9.450000in}}%
\pgfpathlineto{\pgfqpoint{6.509909in}{9.509738in}}%
\pgfpathlineto{\pgfqpoint{6.450171in}{9.450000in}}%
\pgfpathlineto{\pgfqpoint{6.509909in}{9.390262in}}%
\pgfpathclose%
\pgfusepath{stroke}%
\end{pgfscope}%
\begin{pgfscope}%
\pgfpathrectangle{\pgfqpoint{1.109909in}{4.050000in}}{\pgfqpoint{5.400000in}{5.400000in}}%
\pgfusepath{clip}%
\pgfsetbuttcap%
\pgfsetroundjoin%
\pgfsetlinewidth{1.003750pt}%
\definecolor{currentstroke}{rgb}{0.321569,0.321569,0.321569}%
\pgfsetstrokecolor{currentstroke}%
\pgfsetstrokeopacity{0.900000}%
\pgfsetdash{}{0pt}%
\pgfpathmoveto{\pgfqpoint{2.896767in}{5.732081in}}%
\pgfpathcurveto{\pgfqpoint{2.907021in}{5.732081in}}{\pgfqpoint{2.916856in}{5.736155in}}{\pgfqpoint{2.924107in}{5.743406in}}%
\pgfpathcurveto{\pgfqpoint{2.931358in}{5.750657in}}{\pgfqpoint{2.935432in}{5.760492in}}{\pgfqpoint{2.935432in}{5.770746in}}%
\pgfpathcurveto{\pgfqpoint{2.935432in}{5.781000in}}{\pgfqpoint{2.931358in}{5.790836in}}{\pgfqpoint{2.924107in}{5.798086in}}%
\pgfpathcurveto{\pgfqpoint{2.916856in}{5.805337in}}{\pgfqpoint{2.907021in}{5.809411in}}{\pgfqpoint{2.896767in}{5.809411in}}%
\pgfpathcurveto{\pgfqpoint{2.886513in}{5.809411in}}{\pgfqpoint{2.876677in}{5.805337in}}{\pgfqpoint{2.869426in}{5.798086in}}%
\pgfpathcurveto{\pgfqpoint{2.862176in}{5.790836in}}{\pgfqpoint{2.858102in}{5.781000in}}{\pgfqpoint{2.858102in}{5.770746in}}%
\pgfpathcurveto{\pgfqpoint{2.858102in}{5.760492in}}{\pgfqpoint{2.862176in}{5.750657in}}{\pgfqpoint{2.869426in}{5.743406in}}%
\pgfpathcurveto{\pgfqpoint{2.876677in}{5.736155in}}{\pgfqpoint{2.886513in}{5.732081in}}{\pgfqpoint{2.896767in}{5.732081in}}%
\pgfpathlineto{\pgfqpoint{2.896767in}{5.732081in}}%
\pgfpathclose%
\pgfusepath{stroke}%
\end{pgfscope}%
\begin{pgfscope}%
\pgfpathrectangle{\pgfqpoint{1.109909in}{4.050000in}}{\pgfqpoint{5.400000in}{5.400000in}}%
\pgfusepath{clip}%
\pgfsetbuttcap%
\pgfsetroundjoin%
\pgfsetlinewidth{1.003750pt}%
\definecolor{currentstroke}{rgb}{0.321569,0.321569,0.321569}%
\pgfsetstrokecolor{currentstroke}%
\pgfsetstrokeopacity{0.900000}%
\pgfsetdash{}{0pt}%
\pgfpathmoveto{\pgfqpoint{2.509379in}{5.237038in}}%
\pgfpathcurveto{\pgfqpoint{2.519633in}{5.237038in}}{\pgfqpoint{2.529469in}{5.241112in}}{\pgfqpoint{2.536719in}{5.248363in}}%
\pgfpathcurveto{\pgfqpoint{2.543970in}{5.255614in}}{\pgfqpoint{2.548044in}{5.265449in}}{\pgfqpoint{2.548044in}{5.275703in}}%
\pgfpathcurveto{\pgfqpoint{2.548044in}{5.285958in}}{\pgfqpoint{2.543970in}{5.295793in}}{\pgfqpoint{2.536719in}{5.303044in}}%
\pgfpathcurveto{\pgfqpoint{2.529469in}{5.310295in}}{\pgfqpoint{2.519633in}{5.314369in}}{\pgfqpoint{2.509379in}{5.314369in}}%
\pgfpathcurveto{\pgfqpoint{2.499125in}{5.314369in}}{\pgfqpoint{2.489290in}{5.310295in}}{\pgfqpoint{2.482039in}{5.303044in}}%
\pgfpathcurveto{\pgfqpoint{2.474788in}{5.295793in}}{\pgfqpoint{2.470714in}{5.285958in}}{\pgfqpoint{2.470714in}{5.275703in}}%
\pgfpathcurveto{\pgfqpoint{2.470714in}{5.265449in}}{\pgfqpoint{2.474788in}{5.255614in}}{\pgfqpoint{2.482039in}{5.248363in}}%
\pgfpathcurveto{\pgfqpoint{2.489290in}{5.241112in}}{\pgfqpoint{2.499125in}{5.237038in}}{\pgfqpoint{2.509379in}{5.237038in}}%
\pgfpathlineto{\pgfqpoint{2.509379in}{5.237038in}}%
\pgfpathclose%
\pgfusepath{stroke}%
\end{pgfscope}%
\begin{pgfscope}%
\pgfpathrectangle{\pgfqpoint{1.109909in}{4.050000in}}{\pgfqpoint{5.400000in}{5.400000in}}%
\pgfusepath{clip}%
\pgfsetbuttcap%
\pgfsetroundjoin%
\pgfsetlinewidth{1.003750pt}%
\definecolor{currentstroke}{rgb}{0.321569,0.321569,0.321569}%
\pgfsetstrokecolor{currentstroke}%
\pgfsetstrokeopacity{0.900000}%
\pgfsetdash{}{0pt}%
\pgfpathmoveto{\pgfqpoint{2.572917in}{5.851928in}}%
\pgfpathcurveto{\pgfqpoint{2.583171in}{5.851928in}}{\pgfqpoint{2.593007in}{5.856002in}}{\pgfqpoint{2.600257in}{5.863253in}}%
\pgfpathcurveto{\pgfqpoint{2.607508in}{5.870504in}}{\pgfqpoint{2.611582in}{5.880339in}}{\pgfqpoint{2.611582in}{5.890593in}}%
\pgfpathcurveto{\pgfqpoint{2.611582in}{5.900847in}}{\pgfqpoint{2.607508in}{5.910683in}}{\pgfqpoint{2.600257in}{5.917934in}}%
\pgfpathcurveto{\pgfqpoint{2.593007in}{5.925184in}}{\pgfqpoint{2.583171in}{5.929258in}}{\pgfqpoint{2.572917in}{5.929258in}}%
\pgfpathcurveto{\pgfqpoint{2.562663in}{5.929258in}}{\pgfqpoint{2.552827in}{5.925184in}}{\pgfqpoint{2.545577in}{5.917934in}}%
\pgfpathcurveto{\pgfqpoint{2.538326in}{5.910683in}}{\pgfqpoint{2.534252in}{5.900847in}}{\pgfqpoint{2.534252in}{5.890593in}}%
\pgfpathcurveto{\pgfqpoint{2.534252in}{5.880339in}}{\pgfqpoint{2.538326in}{5.870504in}}{\pgfqpoint{2.545577in}{5.863253in}}%
\pgfpathcurveto{\pgfqpoint{2.552827in}{5.856002in}}{\pgfqpoint{2.562663in}{5.851928in}}{\pgfqpoint{2.572917in}{5.851928in}}%
\pgfpathlineto{\pgfqpoint{2.572917in}{5.851928in}}%
\pgfpathclose%
\pgfusepath{stroke}%
\end{pgfscope}%
\begin{pgfscope}%
\pgfpathrectangle{\pgfqpoint{1.109909in}{4.050000in}}{\pgfqpoint{5.400000in}{5.400000in}}%
\pgfusepath{clip}%
\pgfsetbuttcap%
\pgfsetroundjoin%
\pgfsetlinewidth{1.003750pt}%
\definecolor{currentstroke}{rgb}{0.321569,0.321569,0.321569}%
\pgfsetstrokecolor{currentstroke}%
\pgfsetstrokeopacity{0.900000}%
\pgfsetdash{}{0pt}%
\pgfpathmoveto{\pgfqpoint{2.124150in}{4.941137in}}%
\pgfpathcurveto{\pgfqpoint{2.134405in}{4.941137in}}{\pgfqpoint{2.144240in}{4.945211in}}{\pgfqpoint{2.151491in}{4.952462in}}%
\pgfpathcurveto{\pgfqpoint{2.158741in}{4.959713in}}{\pgfqpoint{2.162815in}{4.969548in}}{\pgfqpoint{2.162815in}{4.979802in}}%
\pgfpathcurveto{\pgfqpoint{2.162815in}{4.990057in}}{\pgfqpoint{2.158741in}{4.999892in}}{\pgfqpoint{2.151491in}{5.007143in}}%
\pgfpathcurveto{\pgfqpoint{2.144240in}{5.014393in}}{\pgfqpoint{2.134405in}{5.018467in}}{\pgfqpoint{2.124150in}{5.018467in}}%
\pgfpathcurveto{\pgfqpoint{2.113896in}{5.018467in}}{\pgfqpoint{2.104061in}{5.014393in}}{\pgfqpoint{2.096810in}{5.007143in}}%
\pgfpathcurveto{\pgfqpoint{2.089559in}{4.999892in}}{\pgfqpoint{2.085485in}{4.990057in}}{\pgfqpoint{2.085485in}{4.979802in}}%
\pgfpathcurveto{\pgfqpoint{2.085485in}{4.969548in}}{\pgfqpoint{2.089559in}{4.959713in}}{\pgfqpoint{2.096810in}{4.952462in}}%
\pgfpathcurveto{\pgfqpoint{2.104061in}{4.945211in}}{\pgfqpoint{2.113896in}{4.941137in}}{\pgfqpoint{2.124150in}{4.941137in}}%
\pgfpathlineto{\pgfqpoint{2.124150in}{4.941137in}}%
\pgfpathclose%
\pgfusepath{stroke}%
\end{pgfscope}%
\begin{pgfscope}%
\pgfpathrectangle{\pgfqpoint{1.109909in}{4.050000in}}{\pgfqpoint{5.400000in}{5.400000in}}%
\pgfusepath{clip}%
\pgfsetbuttcap%
\pgfsetroundjoin%
\pgfsetlinewidth{1.003750pt}%
\definecolor{currentstroke}{rgb}{0.321569,0.321569,0.321569}%
\pgfsetstrokecolor{currentstroke}%
\pgfsetstrokeopacity{0.900000}%
\pgfsetdash{}{0pt}%
\pgfpathmoveto{\pgfqpoint{2.298862in}{5.564593in}}%
\pgfpathcurveto{\pgfqpoint{2.309116in}{5.564593in}}{\pgfqpoint{2.318952in}{5.568667in}}{\pgfqpoint{2.326203in}{5.575917in}}%
\pgfpathcurveto{\pgfqpoint{2.333453in}{5.583168in}}{\pgfqpoint{2.337527in}{5.593004in}}{\pgfqpoint{2.337527in}{5.603258in}}%
\pgfpathcurveto{\pgfqpoint{2.337527in}{5.613512in}}{\pgfqpoint{2.333453in}{5.623347in}}{\pgfqpoint{2.326203in}{5.630598in}}%
\pgfpathcurveto{\pgfqpoint{2.318952in}{5.637849in}}{\pgfqpoint{2.309116in}{5.641923in}}{\pgfqpoint{2.298862in}{5.641923in}}%
\pgfpathcurveto{\pgfqpoint{2.288608in}{5.641923in}}{\pgfqpoint{2.278773in}{5.637849in}}{\pgfqpoint{2.271522in}{5.630598in}}%
\pgfpathcurveto{\pgfqpoint{2.264271in}{5.623347in}}{\pgfqpoint{2.260197in}{5.613512in}}{\pgfqpoint{2.260197in}{5.603258in}}%
\pgfpathcurveto{\pgfqpoint{2.260197in}{5.593004in}}{\pgfqpoint{2.264271in}{5.583168in}}{\pgfqpoint{2.271522in}{5.575917in}}%
\pgfpathcurveto{\pgfqpoint{2.278773in}{5.568667in}}{\pgfqpoint{2.288608in}{5.564593in}}{\pgfqpoint{2.298862in}{5.564593in}}%
\pgfpathlineto{\pgfqpoint{2.298862in}{5.564593in}}%
\pgfpathclose%
\pgfusepath{stroke}%
\end{pgfscope}%
\begin{pgfscope}%
\pgfpathrectangle{\pgfqpoint{1.109909in}{4.050000in}}{\pgfqpoint{5.400000in}{5.400000in}}%
\pgfusepath{clip}%
\pgfsetbuttcap%
\pgfsetroundjoin%
\pgfsetlinewidth{1.003750pt}%
\definecolor{currentstroke}{rgb}{0.321569,0.321569,0.321569}%
\pgfsetstrokecolor{currentstroke}%
\pgfsetstrokeopacity{0.900000}%
\pgfsetdash{}{0pt}%
\pgfpathmoveto{\pgfqpoint{4.637424in}{7.856149in}}%
\pgfpathcurveto{\pgfqpoint{4.647678in}{7.856149in}}{\pgfqpoint{4.657514in}{7.860223in}}{\pgfqpoint{4.664764in}{7.867473in}}%
\pgfpathcurveto{\pgfqpoint{4.672015in}{7.874724in}}{\pgfqpoint{4.676089in}{7.884560in}}{\pgfqpoint{4.676089in}{7.894814in}}%
\pgfpathcurveto{\pgfqpoint{4.676089in}{7.905068in}}{\pgfqpoint{4.672015in}{7.914903in}}{\pgfqpoint{4.664764in}{7.922154in}}%
\pgfpathcurveto{\pgfqpoint{4.657514in}{7.929405in}}{\pgfqpoint{4.647678in}{7.933479in}}{\pgfqpoint{4.637424in}{7.933479in}}%
\pgfpathcurveto{\pgfqpoint{4.627170in}{7.933479in}}{\pgfqpoint{4.617334in}{7.929405in}}{\pgfqpoint{4.610084in}{7.922154in}}%
\pgfpathcurveto{\pgfqpoint{4.602833in}{7.914903in}}{\pgfqpoint{4.598759in}{7.905068in}}{\pgfqpoint{4.598759in}{7.894814in}}%
\pgfpathcurveto{\pgfqpoint{4.598759in}{7.884560in}}{\pgfqpoint{4.602833in}{7.874724in}}{\pgfqpoint{4.610084in}{7.867473in}}%
\pgfpathcurveto{\pgfqpoint{4.617334in}{7.860223in}}{\pgfqpoint{4.627170in}{7.856149in}}{\pgfqpoint{4.637424in}{7.856149in}}%
\pgfpathlineto{\pgfqpoint{4.637424in}{7.856149in}}%
\pgfpathclose%
\pgfusepath{stroke}%
\end{pgfscope}%
\begin{pgfscope}%
\pgfpathrectangle{\pgfqpoint{1.109909in}{4.050000in}}{\pgfqpoint{5.400000in}{5.400000in}}%
\pgfusepath{clip}%
\pgfsetbuttcap%
\pgfsetroundjoin%
\pgfsetlinewidth{1.003750pt}%
\definecolor{currentstroke}{rgb}{0.321569,0.321569,0.321569}%
\pgfsetstrokecolor{currentstroke}%
\pgfsetstrokeopacity{0.900000}%
\pgfsetdash{}{0pt}%
\pgfpathmoveto{\pgfqpoint{3.125470in}{6.183529in}}%
\pgfpathcurveto{\pgfqpoint{3.135724in}{6.183529in}}{\pgfqpoint{3.145559in}{6.187603in}}{\pgfqpoint{3.152810in}{6.194854in}}%
\pgfpathcurveto{\pgfqpoint{3.160061in}{6.202105in}}{\pgfqpoint{3.164135in}{6.211940in}}{\pgfqpoint{3.164135in}{6.222194in}}%
\pgfpathcurveto{\pgfqpoint{3.164135in}{6.232448in}}{\pgfqpoint{3.160061in}{6.242284in}}{\pgfqpoint{3.152810in}{6.249534in}}%
\pgfpathcurveto{\pgfqpoint{3.145559in}{6.256785in}}{\pgfqpoint{3.135724in}{6.260859in}}{\pgfqpoint{3.125470in}{6.260859in}}%
\pgfpathcurveto{\pgfqpoint{3.115216in}{6.260859in}}{\pgfqpoint{3.105380in}{6.256785in}}{\pgfqpoint{3.098129in}{6.249534in}}%
\pgfpathcurveto{\pgfqpoint{3.090879in}{6.242284in}}{\pgfqpoint{3.086805in}{6.232448in}}{\pgfqpoint{3.086805in}{6.222194in}}%
\pgfpathcurveto{\pgfqpoint{3.086805in}{6.211940in}}{\pgfqpoint{3.090879in}{6.202105in}}{\pgfqpoint{3.098129in}{6.194854in}}%
\pgfpathcurveto{\pgfqpoint{3.105380in}{6.187603in}}{\pgfqpoint{3.115216in}{6.183529in}}{\pgfqpoint{3.125470in}{6.183529in}}%
\pgfpathlineto{\pgfqpoint{3.125470in}{6.183529in}}%
\pgfpathclose%
\pgfusepath{stroke}%
\end{pgfscope}%
\begin{pgfscope}%
\pgfpathrectangle{\pgfqpoint{1.109909in}{4.050000in}}{\pgfqpoint{5.400000in}{5.400000in}}%
\pgfusepath{clip}%
\pgfsetbuttcap%
\pgfsetroundjoin%
\pgfsetlinewidth{1.003750pt}%
\definecolor{currentstroke}{rgb}{0.321569,0.321569,0.321569}%
\pgfsetstrokecolor{currentstroke}%
\pgfsetstrokeopacity{0.900000}%
\pgfsetdash{}{0pt}%
\pgfpathmoveto{\pgfqpoint{3.136434in}{6.127601in}}%
\pgfpathcurveto{\pgfqpoint{3.146689in}{6.127601in}}{\pgfqpoint{3.156524in}{6.131675in}}{\pgfqpoint{3.163775in}{6.138926in}}%
\pgfpathcurveto{\pgfqpoint{3.171025in}{6.146176in}}{\pgfqpoint{3.175099in}{6.156012in}}{\pgfqpoint{3.175099in}{6.166266in}}%
\pgfpathcurveto{\pgfqpoint{3.175099in}{6.176520in}}{\pgfqpoint{3.171025in}{6.186355in}}{\pgfqpoint{3.163775in}{6.193606in}}%
\pgfpathcurveto{\pgfqpoint{3.156524in}{6.200857in}}{\pgfqpoint{3.146689in}{6.204931in}}{\pgfqpoint{3.136434in}{6.204931in}}%
\pgfpathcurveto{\pgfqpoint{3.126180in}{6.204931in}}{\pgfqpoint{3.116345in}{6.200857in}}{\pgfqpoint{3.109094in}{6.193606in}}%
\pgfpathcurveto{\pgfqpoint{3.101843in}{6.186355in}}{\pgfqpoint{3.097769in}{6.176520in}}{\pgfqpoint{3.097769in}{6.166266in}}%
\pgfpathcurveto{\pgfqpoint{3.097769in}{6.156012in}}{\pgfqpoint{3.101843in}{6.146176in}}{\pgfqpoint{3.109094in}{6.138926in}}%
\pgfpathcurveto{\pgfqpoint{3.116345in}{6.131675in}}{\pgfqpoint{3.126180in}{6.127601in}}{\pgfqpoint{3.136434in}{6.127601in}}%
\pgfpathlineto{\pgfqpoint{3.136434in}{6.127601in}}%
\pgfpathclose%
\pgfusepath{stroke}%
\end{pgfscope}%
\begin{pgfscope}%
\pgfpathrectangle{\pgfqpoint{1.109909in}{4.050000in}}{\pgfqpoint{5.400000in}{5.400000in}}%
\pgfusepath{clip}%
\pgfsetbuttcap%
\pgfsetroundjoin%
\pgfsetlinewidth{1.003750pt}%
\definecolor{currentstroke}{rgb}{0.321569,0.321569,0.321569}%
\pgfsetstrokecolor{currentstroke}%
\pgfsetstrokeopacity{0.900000}%
\pgfsetdash{}{0pt}%
\pgfpathmoveto{\pgfqpoint{2.405662in}{5.946888in}}%
\pgfpathcurveto{\pgfqpoint{2.415916in}{5.946888in}}{\pgfqpoint{2.425751in}{5.950962in}}{\pgfqpoint{2.433002in}{5.958213in}}%
\pgfpathcurveto{\pgfqpoint{2.440253in}{5.965463in}}{\pgfqpoint{2.444327in}{5.975299in}}{\pgfqpoint{2.444327in}{5.985553in}}%
\pgfpathcurveto{\pgfqpoint{2.444327in}{5.995807in}}{\pgfqpoint{2.440253in}{6.005643in}}{\pgfqpoint{2.433002in}{6.012893in}}%
\pgfpathcurveto{\pgfqpoint{2.425751in}{6.020144in}}{\pgfqpoint{2.415916in}{6.024218in}}{\pgfqpoint{2.405662in}{6.024218in}}%
\pgfpathcurveto{\pgfqpoint{2.395407in}{6.024218in}}{\pgfqpoint{2.385572in}{6.020144in}}{\pgfqpoint{2.378321in}{6.012893in}}%
\pgfpathcurveto{\pgfqpoint{2.371071in}{6.005643in}}{\pgfqpoint{2.366997in}{5.995807in}}{\pgfqpoint{2.366997in}{5.985553in}}%
\pgfpathcurveto{\pgfqpoint{2.366997in}{5.975299in}}{\pgfqpoint{2.371071in}{5.965463in}}{\pgfqpoint{2.378321in}{5.958213in}}%
\pgfpathcurveto{\pgfqpoint{2.385572in}{5.950962in}}{\pgfqpoint{2.395407in}{5.946888in}}{\pgfqpoint{2.405662in}{5.946888in}}%
\pgfpathlineto{\pgfqpoint{2.405662in}{5.946888in}}%
\pgfpathclose%
\pgfusepath{stroke}%
\end{pgfscope}%
\begin{pgfscope}%
\pgfpathrectangle{\pgfqpoint{1.109909in}{4.050000in}}{\pgfqpoint{5.400000in}{5.400000in}}%
\pgfusepath{clip}%
\pgfsetbuttcap%
\pgfsetroundjoin%
\pgfsetlinewidth{1.003750pt}%
\definecolor{currentstroke}{rgb}{0.321569,0.321569,0.321569}%
\pgfsetstrokecolor{currentstroke}%
\pgfsetstrokeopacity{0.900000}%
\pgfsetdash{}{0pt}%
\pgfpathmoveto{\pgfqpoint{2.677698in}{6.230355in}}%
\pgfpathcurveto{\pgfqpoint{2.687952in}{6.230355in}}{\pgfqpoint{2.697787in}{6.234429in}}{\pgfqpoint{2.705038in}{6.241679in}}%
\pgfpathcurveto{\pgfqpoint{2.712289in}{6.248930in}}{\pgfqpoint{2.716363in}{6.258766in}}{\pgfqpoint{2.716363in}{6.269020in}}%
\pgfpathcurveto{\pgfqpoint{2.716363in}{6.279274in}}{\pgfqpoint{2.712289in}{6.289109in}}{\pgfqpoint{2.705038in}{6.296360in}}%
\pgfpathcurveto{\pgfqpoint{2.697787in}{6.303611in}}{\pgfqpoint{2.687952in}{6.307685in}}{\pgfqpoint{2.677698in}{6.307685in}}%
\pgfpathcurveto{\pgfqpoint{2.667444in}{6.307685in}}{\pgfqpoint{2.657608in}{6.303611in}}{\pgfqpoint{2.650357in}{6.296360in}}%
\pgfpathcurveto{\pgfqpoint{2.643107in}{6.289109in}}{\pgfqpoint{2.639033in}{6.279274in}}{\pgfqpoint{2.639033in}{6.269020in}}%
\pgfpathcurveto{\pgfqpoint{2.639033in}{6.258766in}}{\pgfqpoint{2.643107in}{6.248930in}}{\pgfqpoint{2.650357in}{6.241679in}}%
\pgfpathcurveto{\pgfqpoint{2.657608in}{6.234429in}}{\pgfqpoint{2.667444in}{6.230355in}}{\pgfqpoint{2.677698in}{6.230355in}}%
\pgfpathlineto{\pgfqpoint{2.677698in}{6.230355in}}%
\pgfpathclose%
\pgfusepath{stroke}%
\end{pgfscope}%
\begin{pgfscope}%
\pgfpathrectangle{\pgfqpoint{1.109909in}{4.050000in}}{\pgfqpoint{5.400000in}{5.400000in}}%
\pgfusepath{clip}%
\pgfsetbuttcap%
\pgfsetroundjoin%
\pgfsetlinewidth{1.003750pt}%
\definecolor{currentstroke}{rgb}{0.321569,0.321569,0.321569}%
\pgfsetstrokecolor{currentstroke}%
\pgfsetstrokeopacity{0.900000}%
\pgfsetdash{}{0pt}%
\pgfpathmoveto{\pgfqpoint{2.081715in}{5.758707in}}%
\pgfpathcurveto{\pgfqpoint{2.091969in}{5.758707in}}{\pgfqpoint{2.101805in}{5.762781in}}{\pgfqpoint{2.109056in}{5.770032in}}%
\pgfpathcurveto{\pgfqpoint{2.116306in}{5.777282in}}{\pgfqpoint{2.120380in}{5.787118in}}{\pgfqpoint{2.120380in}{5.797372in}}%
\pgfpathcurveto{\pgfqpoint{2.120380in}{5.807626in}}{\pgfqpoint{2.116306in}{5.817461in}}{\pgfqpoint{2.109056in}{5.824712in}}%
\pgfpathcurveto{\pgfqpoint{2.101805in}{5.831963in}}{\pgfqpoint{2.091969in}{5.836037in}}{\pgfqpoint{2.081715in}{5.836037in}}%
\pgfpathcurveto{\pgfqpoint{2.071461in}{5.836037in}}{\pgfqpoint{2.061626in}{5.831963in}}{\pgfqpoint{2.054375in}{5.824712in}}%
\pgfpathcurveto{\pgfqpoint{2.047124in}{5.817461in}}{\pgfqpoint{2.043050in}{5.807626in}}{\pgfqpoint{2.043050in}{5.797372in}}%
\pgfpathcurveto{\pgfqpoint{2.043050in}{5.787118in}}{\pgfqpoint{2.047124in}{5.777282in}}{\pgfqpoint{2.054375in}{5.770032in}}%
\pgfpathcurveto{\pgfqpoint{2.061626in}{5.762781in}}{\pgfqpoint{2.071461in}{5.758707in}}{\pgfqpoint{2.081715in}{5.758707in}}%
\pgfpathlineto{\pgfqpoint{2.081715in}{5.758707in}}%
\pgfpathclose%
\pgfusepath{stroke}%
\end{pgfscope}%
\begin{pgfscope}%
\pgfpathrectangle{\pgfqpoint{1.109909in}{4.050000in}}{\pgfqpoint{5.400000in}{5.400000in}}%
\pgfusepath{clip}%
\pgfsetbuttcap%
\pgfsetroundjoin%
\pgfsetlinewidth{1.003750pt}%
\definecolor{currentstroke}{rgb}{0.321569,0.321569,0.321569}%
\pgfsetstrokecolor{currentstroke}%
\pgfsetstrokeopacity{0.900000}%
\pgfsetdash{}{0pt}%
\pgfpathmoveto{\pgfqpoint{2.375252in}{5.937764in}}%
\pgfpathcurveto{\pgfqpoint{2.385507in}{5.937764in}}{\pgfqpoint{2.395342in}{5.941838in}}{\pgfqpoint{2.402593in}{5.949089in}}%
\pgfpathcurveto{\pgfqpoint{2.409843in}{5.956339in}}{\pgfqpoint{2.413917in}{5.966175in}}{\pgfqpoint{2.413917in}{5.976429in}}%
\pgfpathcurveto{\pgfqpoint{2.413917in}{5.986683in}}{\pgfqpoint{2.409843in}{5.996519in}}{\pgfqpoint{2.402593in}{6.003769in}}%
\pgfpathcurveto{\pgfqpoint{2.395342in}{6.011020in}}{\pgfqpoint{2.385507in}{6.015094in}}{\pgfqpoint{2.375252in}{6.015094in}}%
\pgfpathcurveto{\pgfqpoint{2.364998in}{6.015094in}}{\pgfqpoint{2.355163in}{6.011020in}}{\pgfqpoint{2.347912in}{6.003769in}}%
\pgfpathcurveto{\pgfqpoint{2.340661in}{5.996519in}}{\pgfqpoint{2.336587in}{5.986683in}}{\pgfqpoint{2.336587in}{5.976429in}}%
\pgfpathcurveto{\pgfqpoint{2.336587in}{5.966175in}}{\pgfqpoint{2.340661in}{5.956339in}}{\pgfqpoint{2.347912in}{5.949089in}}%
\pgfpathcurveto{\pgfqpoint{2.355163in}{5.941838in}}{\pgfqpoint{2.364998in}{5.937764in}}{\pgfqpoint{2.375252in}{5.937764in}}%
\pgfpathlineto{\pgfqpoint{2.375252in}{5.937764in}}%
\pgfpathclose%
\pgfusepath{stroke}%
\end{pgfscope}%
\begin{pgfscope}%
\pgfpathrectangle{\pgfqpoint{1.109909in}{4.050000in}}{\pgfqpoint{5.400000in}{5.400000in}}%
\pgfusepath{clip}%
\pgfsetbuttcap%
\pgfsetroundjoin%
\pgfsetlinewidth{1.003750pt}%
\definecolor{currentstroke}{rgb}{0.321569,0.321569,0.321569}%
\pgfsetstrokecolor{currentstroke}%
\pgfsetstrokeopacity{0.900000}%
\pgfsetdash{}{0pt}%
\pgfpathmoveto{\pgfqpoint{3.017169in}{6.466603in}}%
\pgfpathcurveto{\pgfqpoint{3.027423in}{6.466603in}}{\pgfqpoint{3.037259in}{6.470677in}}{\pgfqpoint{3.044510in}{6.477927in}}%
\pgfpathcurveto{\pgfqpoint{3.051760in}{6.485178in}}{\pgfqpoint{3.055834in}{6.495014in}}{\pgfqpoint{3.055834in}{6.505268in}}%
\pgfpathcurveto{\pgfqpoint{3.055834in}{6.515522in}}{\pgfqpoint{3.051760in}{6.525357in}}{\pgfqpoint{3.044510in}{6.532608in}}%
\pgfpathcurveto{\pgfqpoint{3.037259in}{6.539859in}}{\pgfqpoint{3.027423in}{6.543933in}}{\pgfqpoint{3.017169in}{6.543933in}}%
\pgfpathcurveto{\pgfqpoint{3.006915in}{6.543933in}}{\pgfqpoint{2.997080in}{6.539859in}}{\pgfqpoint{2.989829in}{6.532608in}}%
\pgfpathcurveto{\pgfqpoint{2.982578in}{6.525357in}}{\pgfqpoint{2.978504in}{6.515522in}}{\pgfqpoint{2.978504in}{6.505268in}}%
\pgfpathcurveto{\pgfqpoint{2.978504in}{6.495014in}}{\pgfqpoint{2.982578in}{6.485178in}}{\pgfqpoint{2.989829in}{6.477927in}}%
\pgfpathcurveto{\pgfqpoint{2.997080in}{6.470677in}}{\pgfqpoint{3.006915in}{6.466603in}}{\pgfqpoint{3.017169in}{6.466603in}}%
\pgfpathlineto{\pgfqpoint{3.017169in}{6.466603in}}%
\pgfpathclose%
\pgfusepath{stroke}%
\end{pgfscope}%
\begin{pgfscope}%
\pgfpathrectangle{\pgfqpoint{1.109909in}{4.050000in}}{\pgfqpoint{5.400000in}{5.400000in}}%
\pgfusepath{clip}%
\pgfsetbuttcap%
\pgfsetroundjoin%
\pgfsetlinewidth{1.003750pt}%
\definecolor{currentstroke}{rgb}{0.321569,0.321569,0.321569}%
\pgfsetstrokecolor{currentstroke}%
\pgfsetstrokeopacity{0.900000}%
\pgfsetdash{}{0pt}%
\pgfpathmoveto{\pgfqpoint{2.912955in}{5.715049in}}%
\pgfpathcurveto{\pgfqpoint{2.923210in}{5.715049in}}{\pgfqpoint{2.933045in}{5.719123in}}{\pgfqpoint{2.940296in}{5.726373in}}%
\pgfpathcurveto{\pgfqpoint{2.947547in}{5.733624in}}{\pgfqpoint{2.951621in}{5.743460in}}{\pgfqpoint{2.951621in}{5.753714in}}%
\pgfpathcurveto{\pgfqpoint{2.951621in}{5.763968in}}{\pgfqpoint{2.947547in}{5.773803in}}{\pgfqpoint{2.940296in}{5.781054in}}%
\pgfpathcurveto{\pgfqpoint{2.933045in}{5.788305in}}{\pgfqpoint{2.923210in}{5.792379in}}{\pgfqpoint{2.912955in}{5.792379in}}%
\pgfpathcurveto{\pgfqpoint{2.902701in}{5.792379in}}{\pgfqpoint{2.892866in}{5.788305in}}{\pgfqpoint{2.885615in}{5.781054in}}%
\pgfpathcurveto{\pgfqpoint{2.878364in}{5.773803in}}{\pgfqpoint{2.874290in}{5.763968in}}{\pgfqpoint{2.874290in}{5.753714in}}%
\pgfpathcurveto{\pgfqpoint{2.874290in}{5.743460in}}{\pgfqpoint{2.878364in}{5.733624in}}{\pgfqpoint{2.885615in}{5.726373in}}%
\pgfpathcurveto{\pgfqpoint{2.892866in}{5.719123in}}{\pgfqpoint{2.902701in}{5.715049in}}{\pgfqpoint{2.912955in}{5.715049in}}%
\pgfpathlineto{\pgfqpoint{2.912955in}{5.715049in}}%
\pgfpathclose%
\pgfusepath{stroke}%
\end{pgfscope}%
\begin{pgfscope}%
\pgfpathrectangle{\pgfqpoint{1.109909in}{4.050000in}}{\pgfqpoint{5.400000in}{5.400000in}}%
\pgfusepath{clip}%
\pgfsetbuttcap%
\pgfsetroundjoin%
\pgfsetlinewidth{1.003750pt}%
\definecolor{currentstroke}{rgb}{0.321569,0.321569,0.321569}%
\pgfsetstrokecolor{currentstroke}%
\pgfsetstrokeopacity{0.900000}%
\pgfsetdash{}{0pt}%
\pgfpathmoveto{\pgfqpoint{3.153354in}{5.910398in}}%
\pgfpathcurveto{\pgfqpoint{3.163609in}{5.910398in}}{\pgfqpoint{3.173444in}{5.914472in}}{\pgfqpoint{3.180695in}{5.921722in}}%
\pgfpathcurveto{\pgfqpoint{3.187946in}{5.928973in}}{\pgfqpoint{3.192020in}{5.938809in}}{\pgfqpoint{3.192020in}{5.949063in}}%
\pgfpathcurveto{\pgfqpoint{3.192020in}{5.959317in}}{\pgfqpoint{3.187946in}{5.969152in}}{\pgfqpoint{3.180695in}{5.976403in}}%
\pgfpathcurveto{\pgfqpoint{3.173444in}{5.983654in}}{\pgfqpoint{3.163609in}{5.987728in}}{\pgfqpoint{3.153354in}{5.987728in}}%
\pgfpathcurveto{\pgfqpoint{3.143100in}{5.987728in}}{\pgfqpoint{3.133265in}{5.983654in}}{\pgfqpoint{3.126014in}{5.976403in}}%
\pgfpathcurveto{\pgfqpoint{3.118763in}{5.969152in}}{\pgfqpoint{3.114689in}{5.959317in}}{\pgfqpoint{3.114689in}{5.949063in}}%
\pgfpathcurveto{\pgfqpoint{3.114689in}{5.938809in}}{\pgfqpoint{3.118763in}{5.928973in}}{\pgfqpoint{3.126014in}{5.921722in}}%
\pgfpathcurveto{\pgfqpoint{3.133265in}{5.914472in}}{\pgfqpoint{3.143100in}{5.910398in}}{\pgfqpoint{3.153354in}{5.910398in}}%
\pgfpathlineto{\pgfqpoint{3.153354in}{5.910398in}}%
\pgfpathclose%
\pgfusepath{stroke}%
\end{pgfscope}%
\begin{pgfscope}%
\pgfpathrectangle{\pgfqpoint{1.109909in}{4.050000in}}{\pgfqpoint{5.400000in}{5.400000in}}%
\pgfusepath{clip}%
\pgfsetbuttcap%
\pgfsetroundjoin%
\pgfsetlinewidth{1.003750pt}%
\definecolor{currentstroke}{rgb}{0.321569,0.321569,0.321569}%
\pgfsetstrokecolor{currentstroke}%
\pgfsetstrokeopacity{0.900000}%
\pgfsetdash{}{0pt}%
\pgfpathmoveto{\pgfqpoint{2.883340in}{6.420814in}}%
\pgfpathcurveto{\pgfqpoint{2.893594in}{6.420814in}}{\pgfqpoint{2.903429in}{6.424888in}}{\pgfqpoint{2.910680in}{6.432139in}}%
\pgfpathcurveto{\pgfqpoint{2.917931in}{6.439390in}}{\pgfqpoint{2.922005in}{6.449225in}}{\pgfqpoint{2.922005in}{6.459479in}}%
\pgfpathcurveto{\pgfqpoint{2.922005in}{6.469733in}}{\pgfqpoint{2.917931in}{6.479569in}}{\pgfqpoint{2.910680in}{6.486820in}}%
\pgfpathcurveto{\pgfqpoint{2.903429in}{6.494070in}}{\pgfqpoint{2.893594in}{6.498144in}}{\pgfqpoint{2.883340in}{6.498144in}}%
\pgfpathcurveto{\pgfqpoint{2.873086in}{6.498144in}}{\pgfqpoint{2.863250in}{6.494070in}}{\pgfqpoint{2.855999in}{6.486820in}}%
\pgfpathcurveto{\pgfqpoint{2.848749in}{6.479569in}}{\pgfqpoint{2.844675in}{6.469733in}}{\pgfqpoint{2.844675in}{6.459479in}}%
\pgfpathcurveto{\pgfqpoint{2.844675in}{6.449225in}}{\pgfqpoint{2.848749in}{6.439390in}}{\pgfqpoint{2.855999in}{6.432139in}}%
\pgfpathcurveto{\pgfqpoint{2.863250in}{6.424888in}}{\pgfqpoint{2.873086in}{6.420814in}}{\pgfqpoint{2.883340in}{6.420814in}}%
\pgfpathlineto{\pgfqpoint{2.883340in}{6.420814in}}%
\pgfpathclose%
\pgfusepath{stroke}%
\end{pgfscope}%
\begin{pgfscope}%
\pgfpathrectangle{\pgfqpoint{1.109909in}{4.050000in}}{\pgfqpoint{5.400000in}{5.400000in}}%
\pgfusepath{clip}%
\pgfsetbuttcap%
\pgfsetroundjoin%
\pgfsetlinewidth{1.003750pt}%
\definecolor{currentstroke}{rgb}{0.321569,0.321569,0.321569}%
\pgfsetstrokecolor{currentstroke}%
\pgfsetstrokeopacity{0.900000}%
\pgfsetdash{}{0pt}%
\pgfpathmoveto{\pgfqpoint{2.406228in}{5.363449in}}%
\pgfpathcurveto{\pgfqpoint{2.416482in}{5.363449in}}{\pgfqpoint{2.426317in}{5.367523in}}{\pgfqpoint{2.433568in}{5.374773in}}%
\pgfpathcurveto{\pgfqpoint{2.440819in}{5.382024in}}{\pgfqpoint{2.444893in}{5.391859in}}{\pgfqpoint{2.444893in}{5.402114in}}%
\pgfpathcurveto{\pgfqpoint{2.444893in}{5.412368in}}{\pgfqpoint{2.440819in}{5.422203in}}{\pgfqpoint{2.433568in}{5.429454in}}%
\pgfpathcurveto{\pgfqpoint{2.426317in}{5.436705in}}{\pgfqpoint{2.416482in}{5.440779in}}{\pgfqpoint{2.406228in}{5.440779in}}%
\pgfpathcurveto{\pgfqpoint{2.395974in}{5.440779in}}{\pgfqpoint{2.386138in}{5.436705in}}{\pgfqpoint{2.378887in}{5.429454in}}%
\pgfpathcurveto{\pgfqpoint{2.371637in}{5.422203in}}{\pgfqpoint{2.367563in}{5.412368in}}{\pgfqpoint{2.367563in}{5.402114in}}%
\pgfpathcurveto{\pgfqpoint{2.367563in}{5.391859in}}{\pgfqpoint{2.371637in}{5.382024in}}{\pgfqpoint{2.378887in}{5.374773in}}%
\pgfpathcurveto{\pgfqpoint{2.386138in}{5.367523in}}{\pgfqpoint{2.395974in}{5.363449in}}{\pgfqpoint{2.406228in}{5.363449in}}%
\pgfpathlineto{\pgfqpoint{2.406228in}{5.363449in}}%
\pgfpathclose%
\pgfusepath{stroke}%
\end{pgfscope}%
\begin{pgfscope}%
\pgfpathrectangle{\pgfqpoint{1.109909in}{4.050000in}}{\pgfqpoint{5.400000in}{5.400000in}}%
\pgfusepath{clip}%
\pgfsetbuttcap%
\pgfsetroundjoin%
\pgfsetlinewidth{1.003750pt}%
\definecolor{currentstroke}{rgb}{0.321569,0.321569,0.321569}%
\pgfsetstrokecolor{currentstroke}%
\pgfsetstrokeopacity{0.900000}%
\pgfsetdash{}{0pt}%
\pgfpathmoveto{\pgfqpoint{2.529534in}{5.827003in}}%
\pgfpathcurveto{\pgfqpoint{2.539788in}{5.827003in}}{\pgfqpoint{2.549624in}{5.831077in}}{\pgfqpoint{2.556875in}{5.838328in}}%
\pgfpathcurveto{\pgfqpoint{2.564125in}{5.845579in}}{\pgfqpoint{2.568199in}{5.855414in}}{\pgfqpoint{2.568199in}{5.865668in}}%
\pgfpathcurveto{\pgfqpoint{2.568199in}{5.875922in}}{\pgfqpoint{2.564125in}{5.885758in}}{\pgfqpoint{2.556875in}{5.893008in}}%
\pgfpathcurveto{\pgfqpoint{2.549624in}{5.900259in}}{\pgfqpoint{2.539788in}{5.904333in}}{\pgfqpoint{2.529534in}{5.904333in}}%
\pgfpathcurveto{\pgfqpoint{2.519280in}{5.904333in}}{\pgfqpoint{2.509445in}{5.900259in}}{\pgfqpoint{2.502194in}{5.893008in}}%
\pgfpathcurveto{\pgfqpoint{2.494943in}{5.885758in}}{\pgfqpoint{2.490869in}{5.875922in}}{\pgfqpoint{2.490869in}{5.865668in}}%
\pgfpathcurveto{\pgfqpoint{2.490869in}{5.855414in}}{\pgfqpoint{2.494943in}{5.845579in}}{\pgfqpoint{2.502194in}{5.838328in}}%
\pgfpathcurveto{\pgfqpoint{2.509445in}{5.831077in}}{\pgfqpoint{2.519280in}{5.827003in}}{\pgfqpoint{2.529534in}{5.827003in}}%
\pgfpathlineto{\pgfqpoint{2.529534in}{5.827003in}}%
\pgfpathclose%
\pgfusepath{stroke}%
\end{pgfscope}%
\begin{pgfscope}%
\pgfpathrectangle{\pgfqpoint{1.109909in}{4.050000in}}{\pgfqpoint{5.400000in}{5.400000in}}%
\pgfusepath{clip}%
\pgfsetbuttcap%
\pgfsetroundjoin%
\pgfsetlinewidth{1.003750pt}%
\definecolor{currentstroke}{rgb}{0.321569,0.321569,0.321569}%
\pgfsetstrokecolor{currentstroke}%
\pgfsetstrokeopacity{0.900000}%
\pgfsetdash{}{0pt}%
\pgfpathmoveto{\pgfqpoint{2.227219in}{5.447358in}}%
\pgfpathcurveto{\pgfqpoint{2.237473in}{5.447358in}}{\pgfqpoint{2.247309in}{5.451432in}}{\pgfqpoint{2.254559in}{5.458683in}}%
\pgfpathcurveto{\pgfqpoint{2.261810in}{5.465934in}}{\pgfqpoint{2.265884in}{5.475769in}}{\pgfqpoint{2.265884in}{5.486023in}}%
\pgfpathcurveto{\pgfqpoint{2.265884in}{5.496277in}}{\pgfqpoint{2.261810in}{5.506113in}}{\pgfqpoint{2.254559in}{5.513364in}}%
\pgfpathcurveto{\pgfqpoint{2.247309in}{5.520614in}}{\pgfqpoint{2.237473in}{5.524688in}}{\pgfqpoint{2.227219in}{5.524688in}}%
\pgfpathcurveto{\pgfqpoint{2.216965in}{5.524688in}}{\pgfqpoint{2.207129in}{5.520614in}}{\pgfqpoint{2.199879in}{5.513364in}}%
\pgfpathcurveto{\pgfqpoint{2.192628in}{5.506113in}}{\pgfqpoint{2.188554in}{5.496277in}}{\pgfqpoint{2.188554in}{5.486023in}}%
\pgfpathcurveto{\pgfqpoint{2.188554in}{5.475769in}}{\pgfqpoint{2.192628in}{5.465934in}}{\pgfqpoint{2.199879in}{5.458683in}}%
\pgfpathcurveto{\pgfqpoint{2.207129in}{5.451432in}}{\pgfqpoint{2.216965in}{5.447358in}}{\pgfqpoint{2.227219in}{5.447358in}}%
\pgfpathlineto{\pgfqpoint{2.227219in}{5.447358in}}%
\pgfpathclose%
\pgfusepath{stroke}%
\end{pgfscope}%
\begin{pgfscope}%
\pgfpathrectangle{\pgfqpoint{1.109909in}{4.050000in}}{\pgfqpoint{5.400000in}{5.400000in}}%
\pgfusepath{clip}%
\pgfsetbuttcap%
\pgfsetroundjoin%
\pgfsetlinewidth{1.003750pt}%
\definecolor{currentstroke}{rgb}{0.321569,0.321569,0.321569}%
\pgfsetstrokecolor{currentstroke}%
\pgfsetstrokeopacity{0.900000}%
\pgfsetdash{}{0pt}%
\pgfpathmoveto{\pgfqpoint{2.962151in}{6.133640in}}%
\pgfpathcurveto{\pgfqpoint{2.972405in}{6.133640in}}{\pgfqpoint{2.982241in}{6.137714in}}{\pgfqpoint{2.989492in}{6.144964in}}%
\pgfpathcurveto{\pgfqpoint{2.996742in}{6.152215in}}{\pgfqpoint{3.000816in}{6.162051in}}{\pgfqpoint{3.000816in}{6.172305in}}%
\pgfpathcurveto{\pgfqpoint{3.000816in}{6.182559in}}{\pgfqpoint{2.996742in}{6.192394in}}{\pgfqpoint{2.989492in}{6.199645in}}%
\pgfpathcurveto{\pgfqpoint{2.982241in}{6.206896in}}{\pgfqpoint{2.972405in}{6.210970in}}{\pgfqpoint{2.962151in}{6.210970in}}%
\pgfpathcurveto{\pgfqpoint{2.951897in}{6.210970in}}{\pgfqpoint{2.942062in}{6.206896in}}{\pgfqpoint{2.934811in}{6.199645in}}%
\pgfpathcurveto{\pgfqpoint{2.927560in}{6.192394in}}{\pgfqpoint{2.923486in}{6.182559in}}{\pgfqpoint{2.923486in}{6.172305in}}%
\pgfpathcurveto{\pgfqpoint{2.923486in}{6.162051in}}{\pgfqpoint{2.927560in}{6.152215in}}{\pgfqpoint{2.934811in}{6.144964in}}%
\pgfpathcurveto{\pgfqpoint{2.942062in}{6.137714in}}{\pgfqpoint{2.951897in}{6.133640in}}{\pgfqpoint{2.962151in}{6.133640in}}%
\pgfpathlineto{\pgfqpoint{2.962151in}{6.133640in}}%
\pgfpathclose%
\pgfusepath{stroke}%
\end{pgfscope}%
\begin{pgfscope}%
\pgfpathrectangle{\pgfqpoint{1.109909in}{4.050000in}}{\pgfqpoint{5.400000in}{5.400000in}}%
\pgfusepath{clip}%
\pgfsetbuttcap%
\pgfsetroundjoin%
\pgfsetlinewidth{1.003750pt}%
\definecolor{currentstroke}{rgb}{0.321569,0.321569,0.321569}%
\pgfsetstrokecolor{currentstroke}%
\pgfsetstrokeopacity{0.900000}%
\pgfsetdash{}{0pt}%
\pgfpathmoveto{\pgfqpoint{4.649867in}{7.578043in}}%
\pgfpathcurveto{\pgfqpoint{4.660121in}{7.578043in}}{\pgfqpoint{4.669957in}{7.582117in}}{\pgfqpoint{4.677208in}{7.589367in}}%
\pgfpathcurveto{\pgfqpoint{4.684458in}{7.596618in}}{\pgfqpoint{4.688532in}{7.606454in}}{\pgfqpoint{4.688532in}{7.616708in}}%
\pgfpathcurveto{\pgfqpoint{4.688532in}{7.626962in}}{\pgfqpoint{4.684458in}{7.636797in}}{\pgfqpoint{4.677208in}{7.644048in}}%
\pgfpathcurveto{\pgfqpoint{4.669957in}{7.651299in}}{\pgfqpoint{4.660121in}{7.655373in}}{\pgfqpoint{4.649867in}{7.655373in}}%
\pgfpathcurveto{\pgfqpoint{4.639613in}{7.655373in}}{\pgfqpoint{4.629778in}{7.651299in}}{\pgfqpoint{4.622527in}{7.644048in}}%
\pgfpathcurveto{\pgfqpoint{4.615276in}{7.636797in}}{\pgfqpoint{4.611202in}{7.626962in}}{\pgfqpoint{4.611202in}{7.616708in}}%
\pgfpathcurveto{\pgfqpoint{4.611202in}{7.606454in}}{\pgfqpoint{4.615276in}{7.596618in}}{\pgfqpoint{4.622527in}{7.589367in}}%
\pgfpathcurveto{\pgfqpoint{4.629778in}{7.582117in}}{\pgfqpoint{4.639613in}{7.578043in}}{\pgfqpoint{4.649867in}{7.578043in}}%
\pgfpathlineto{\pgfqpoint{4.649867in}{7.578043in}}%
\pgfpathclose%
\pgfusepath{stroke}%
\end{pgfscope}%
\begin{pgfscope}%
\pgfpathrectangle{\pgfqpoint{1.109909in}{4.050000in}}{\pgfqpoint{5.400000in}{5.400000in}}%
\pgfusepath{clip}%
\pgfsetbuttcap%
\pgfsetroundjoin%
\pgfsetlinewidth{1.003750pt}%
\definecolor{currentstroke}{rgb}{0.321569,0.321569,0.321569}%
\pgfsetstrokecolor{currentstroke}%
\pgfsetstrokeopacity{0.900000}%
\pgfsetdash{}{0pt}%
\pgfpathmoveto{\pgfqpoint{2.078843in}{4.893654in}}%
\pgfpathcurveto{\pgfqpoint{2.089097in}{4.893654in}}{\pgfqpoint{2.098932in}{4.897728in}}{\pgfqpoint{2.106183in}{4.904979in}}%
\pgfpathcurveto{\pgfqpoint{2.113434in}{4.912230in}}{\pgfqpoint{2.117508in}{4.922065in}}{\pgfqpoint{2.117508in}{4.932319in}}%
\pgfpathcurveto{\pgfqpoint{2.117508in}{4.942573in}}{\pgfqpoint{2.113434in}{4.952409in}}{\pgfqpoint{2.106183in}{4.959660in}}%
\pgfpathcurveto{\pgfqpoint{2.098932in}{4.966910in}}{\pgfqpoint{2.089097in}{4.970984in}}{\pgfqpoint{2.078843in}{4.970984in}}%
\pgfpathcurveto{\pgfqpoint{2.068588in}{4.970984in}}{\pgfqpoint{2.058753in}{4.966910in}}{\pgfqpoint{2.051502in}{4.959660in}}%
\pgfpathcurveto{\pgfqpoint{2.044252in}{4.952409in}}{\pgfqpoint{2.040178in}{4.942573in}}{\pgfqpoint{2.040178in}{4.932319in}}%
\pgfpathcurveto{\pgfqpoint{2.040178in}{4.922065in}}{\pgfqpoint{2.044252in}{4.912230in}}{\pgfqpoint{2.051502in}{4.904979in}}%
\pgfpathcurveto{\pgfqpoint{2.058753in}{4.897728in}}{\pgfqpoint{2.068588in}{4.893654in}}{\pgfqpoint{2.078843in}{4.893654in}}%
\pgfpathlineto{\pgfqpoint{2.078843in}{4.893654in}}%
\pgfpathclose%
\pgfusepath{stroke}%
\end{pgfscope}%
\begin{pgfscope}%
\pgfpathrectangle{\pgfqpoint{1.109909in}{4.050000in}}{\pgfqpoint{5.400000in}{5.400000in}}%
\pgfusepath{clip}%
\pgfsetbuttcap%
\pgfsetroundjoin%
\pgfsetlinewidth{1.003750pt}%
\definecolor{currentstroke}{rgb}{0.321569,0.321569,0.321569}%
\pgfsetstrokecolor{currentstroke}%
\pgfsetstrokeopacity{0.900000}%
\pgfsetdash{}{0pt}%
\pgfpathmoveto{\pgfqpoint{4.262444in}{7.650613in}}%
\pgfpathcurveto{\pgfqpoint{4.272698in}{7.650613in}}{\pgfqpoint{4.282533in}{7.654687in}}{\pgfqpoint{4.289784in}{7.661937in}}%
\pgfpathcurveto{\pgfqpoint{4.297035in}{7.669188in}}{\pgfqpoint{4.301109in}{7.679023in}}{\pgfqpoint{4.301109in}{7.689278in}}%
\pgfpathcurveto{\pgfqpoint{4.301109in}{7.699532in}}{\pgfqpoint{4.297035in}{7.709367in}}{\pgfqpoint{4.289784in}{7.716618in}}%
\pgfpathcurveto{\pgfqpoint{4.282533in}{7.723869in}}{\pgfqpoint{4.272698in}{7.727943in}}{\pgfqpoint{4.262444in}{7.727943in}}%
\pgfpathcurveto{\pgfqpoint{4.252189in}{7.727943in}}{\pgfqpoint{4.242354in}{7.723869in}}{\pgfqpoint{4.235103in}{7.716618in}}%
\pgfpathcurveto{\pgfqpoint{4.227852in}{7.709367in}}{\pgfqpoint{4.223778in}{7.699532in}}{\pgfqpoint{4.223778in}{7.689278in}}%
\pgfpathcurveto{\pgfqpoint{4.223778in}{7.679023in}}{\pgfqpoint{4.227852in}{7.669188in}}{\pgfqpoint{4.235103in}{7.661937in}}%
\pgfpathcurveto{\pgfqpoint{4.242354in}{7.654687in}}{\pgfqpoint{4.252189in}{7.650613in}}{\pgfqpoint{4.262444in}{7.650613in}}%
\pgfpathlineto{\pgfqpoint{4.262444in}{7.650613in}}%
\pgfpathclose%
\pgfusepath{stroke}%
\end{pgfscope}%
\begin{pgfscope}%
\pgfpathrectangle{\pgfqpoint{1.109909in}{4.050000in}}{\pgfqpoint{5.400000in}{5.400000in}}%
\pgfusepath{clip}%
\pgfsetbuttcap%
\pgfsetroundjoin%
\pgfsetlinewidth{1.003750pt}%
\definecolor{currentstroke}{rgb}{0.321569,0.321569,0.321569}%
\pgfsetstrokecolor{currentstroke}%
\pgfsetstrokeopacity{0.900000}%
\pgfsetdash{}{0pt}%
\pgfpathmoveto{\pgfqpoint{3.391060in}{5.893897in}}%
\pgfpathcurveto{\pgfqpoint{3.401314in}{5.893897in}}{\pgfqpoint{3.411150in}{5.897971in}}{\pgfqpoint{3.418401in}{5.905222in}}%
\pgfpathcurveto{\pgfqpoint{3.425651in}{5.912472in}}{\pgfqpoint{3.429725in}{5.922308in}}{\pgfqpoint{3.429725in}{5.932562in}}%
\pgfpathcurveto{\pgfqpoint{3.429725in}{5.942816in}}{\pgfqpoint{3.425651in}{5.952652in}}{\pgfqpoint{3.418401in}{5.959902in}}%
\pgfpathcurveto{\pgfqpoint{3.411150in}{5.967153in}}{\pgfqpoint{3.401314in}{5.971227in}}{\pgfqpoint{3.391060in}{5.971227in}}%
\pgfpathcurveto{\pgfqpoint{3.380806in}{5.971227in}}{\pgfqpoint{3.370971in}{5.967153in}}{\pgfqpoint{3.363720in}{5.959902in}}%
\pgfpathcurveto{\pgfqpoint{3.356469in}{5.952652in}}{\pgfqpoint{3.352395in}{5.942816in}}{\pgfqpoint{3.352395in}{5.932562in}}%
\pgfpathcurveto{\pgfqpoint{3.352395in}{5.922308in}}{\pgfqpoint{3.356469in}{5.912472in}}{\pgfqpoint{3.363720in}{5.905222in}}%
\pgfpathcurveto{\pgfqpoint{3.370971in}{5.897971in}}{\pgfqpoint{3.380806in}{5.893897in}}{\pgfqpoint{3.391060in}{5.893897in}}%
\pgfpathlineto{\pgfqpoint{3.391060in}{5.893897in}}%
\pgfpathclose%
\pgfusepath{stroke}%
\end{pgfscope}%
\begin{pgfscope}%
\pgfsetbuttcap%
\pgfsetroundjoin%
\pgfsetlinewidth{1.003750pt}%
\definecolor{currentstroke}{rgb}{0.321569,0.321569,0.321569}%
\pgfsetstrokecolor{currentstroke}%
\pgfsetstrokeopacity{0.900000}%
\pgfsetdash{}{0pt}%
\pgfpathmoveto{\pgfqpoint{1.109909in}{3.990262in}}%
\pgfpathlineto{\pgfqpoint{1.169647in}{4.050000in}}%
\pgfpathlineto{\pgfqpoint{1.109909in}{4.109738in}}%
\pgfpathlineto{\pgfqpoint{1.050171in}{4.050000in}}%
\pgfpathlineto{\pgfqpoint{1.109909in}{3.990262in}}%
\pgfpathclose%
\pgfusepath{stroke}%
\end{pgfscope}%
\begin{pgfscope}%
\pgfsetbuttcap%
\pgfsetroundjoin%
\pgfsetlinewidth{1.003750pt}%
\definecolor{currentstroke}{rgb}{0.321569,0.321569,0.321569}%
\pgfsetstrokecolor{currentstroke}%
\pgfsetstrokeopacity{0.900000}%
\pgfsetdash{}{0pt}%
\pgfpathmoveto{\pgfqpoint{1.109909in}{3.990262in}}%
\pgfpathlineto{\pgfqpoint{1.169647in}{4.050000in}}%
\pgfpathlineto{\pgfqpoint{1.109909in}{4.109738in}}%
\pgfpathlineto{\pgfqpoint{1.050171in}{4.050000in}}%
\pgfpathlineto{\pgfqpoint{1.109909in}{3.990262in}}%
\pgfpathclose%
\pgfusepath{stroke}%
\end{pgfscope}%
\begin{pgfscope}%
\pgfsetbuttcap%
\pgfsetroundjoin%
\pgfsetlinewidth{1.003750pt}%
\definecolor{currentstroke}{rgb}{0.321569,0.321569,0.321569}%
\pgfsetstrokecolor{currentstroke}%
\pgfsetstrokeopacity{0.900000}%
\pgfsetdash{}{0pt}%
\pgfpathmoveto{\pgfqpoint{1.109909in}{3.990262in}}%
\pgfpathlineto{\pgfqpoint{1.169647in}{4.050000in}}%
\pgfpathlineto{\pgfqpoint{1.109909in}{4.109738in}}%
\pgfpathlineto{\pgfqpoint{1.050171in}{4.050000in}}%
\pgfpathlineto{\pgfqpoint{1.109909in}{3.990262in}}%
\pgfpathclose%
\pgfusepath{stroke}%
\end{pgfscope}%
\begin{pgfscope}%
\pgfsetbuttcap%
\pgfsetroundjoin%
\pgfsetlinewidth{1.003750pt}%
\definecolor{currentstroke}{rgb}{0.321569,0.321569,0.321569}%
\pgfsetstrokecolor{currentstroke}%
\pgfsetstrokeopacity{0.900000}%
\pgfsetdash{}{0pt}%
\pgfpathmoveto{\pgfqpoint{1.109909in}{4.243287in}}%
\pgfpathlineto{\pgfqpoint{1.169647in}{4.303025in}}%
\pgfpathlineto{\pgfqpoint{1.109909in}{4.362763in}}%
\pgfpathlineto{\pgfqpoint{1.050171in}{4.303025in}}%
\pgfpathlineto{\pgfqpoint{1.109909in}{4.243287in}}%
\pgfpathclose%
\pgfusepath{stroke}%
\end{pgfscope}%
\begin{pgfscope}%
\pgfsetbuttcap%
\pgfsetroundjoin%
\pgfsetlinewidth{1.003750pt}%
\definecolor{currentstroke}{rgb}{0.321569,0.321569,0.321569}%
\pgfsetstrokecolor{currentstroke}%
\pgfsetstrokeopacity{0.900000}%
\pgfsetdash{}{0pt}%
\pgfpathmoveto{\pgfqpoint{1.109909in}{4.398923in}}%
\pgfpathlineto{\pgfqpoint{1.169647in}{4.458661in}}%
\pgfpathlineto{\pgfqpoint{1.109909in}{4.518399in}}%
\pgfpathlineto{\pgfqpoint{1.050171in}{4.458661in}}%
\pgfpathlineto{\pgfqpoint{1.109909in}{4.398923in}}%
\pgfpathclose%
\pgfusepath{stroke}%
\end{pgfscope}%
\begin{pgfscope}%
\pgfsetbuttcap%
\pgfsetroundjoin%
\pgfsetlinewidth{1.003750pt}%
\definecolor{currentstroke}{rgb}{0.321569,0.321569,0.321569}%
\pgfsetstrokecolor{currentstroke}%
\pgfsetstrokeopacity{0.900000}%
\pgfsetdash{}{0pt}%
\pgfpathmoveto{\pgfqpoint{1.109909in}{3.990262in}}%
\pgfpathlineto{\pgfqpoint{1.169647in}{4.050000in}}%
\pgfpathlineto{\pgfqpoint{1.109909in}{4.109738in}}%
\pgfpathlineto{\pgfqpoint{1.050171in}{4.050000in}}%
\pgfpathlineto{\pgfqpoint{1.109909in}{3.990262in}}%
\pgfpathclose%
\pgfusepath{stroke}%
\end{pgfscope}%
\begin{pgfscope}%
\pgfpathrectangle{\pgfqpoint{1.109909in}{4.050000in}}{\pgfqpoint{5.400000in}{5.400000in}}%
\pgfusepath{clip}%
\pgfsetbuttcap%
\pgfsetroundjoin%
\pgfsetlinewidth{1.003750pt}%
\definecolor{currentstroke}{rgb}{0.768627,0.603922,0.000000}%
\pgfsetstrokecolor{currentstroke}%
\pgfsetstrokeopacity{0.900000}%
\pgfsetdash{}{0pt}%
\pgfpathmoveto{\pgfqpoint{1.626682in}{8.832877in}}%
\pgfpathcurveto{\pgfqpoint{1.636936in}{8.832877in}}{\pgfqpoint{1.646772in}{8.836951in}}{\pgfqpoint{1.654022in}{8.844201in}}%
\pgfpathcurveto{\pgfqpoint{1.661273in}{8.851452in}}{\pgfqpoint{1.665347in}{8.861287in}}{\pgfqpoint{1.665347in}{8.871542in}}%
\pgfpathcurveto{\pgfqpoint{1.665347in}{8.881796in}}{\pgfqpoint{1.661273in}{8.891631in}}{\pgfqpoint{1.654022in}{8.898882in}}%
\pgfpathcurveto{\pgfqpoint{1.646772in}{8.906133in}}{\pgfqpoint{1.636936in}{8.910207in}}{\pgfqpoint{1.626682in}{8.910207in}}%
\pgfpathcurveto{\pgfqpoint{1.616428in}{8.910207in}}{\pgfqpoint{1.606593in}{8.906133in}}{\pgfqpoint{1.599342in}{8.898882in}}%
\pgfpathcurveto{\pgfqpoint{1.592091in}{8.891631in}}{\pgfqpoint{1.588017in}{8.881796in}}{\pgfqpoint{1.588017in}{8.871542in}}%
\pgfpathcurveto{\pgfqpoint{1.588017in}{8.861287in}}{\pgfqpoint{1.592091in}{8.851452in}}{\pgfqpoint{1.599342in}{8.844201in}}%
\pgfpathcurveto{\pgfqpoint{1.606593in}{8.836951in}}{\pgfqpoint{1.616428in}{8.832877in}}{\pgfqpoint{1.626682in}{8.832877in}}%
\pgfpathlineto{\pgfqpoint{1.626682in}{8.832877in}}%
\pgfpathclose%
\pgfusepath{stroke}%
\end{pgfscope}%
\begin{pgfscope}%
\pgfpathrectangle{\pgfqpoint{1.109909in}{4.050000in}}{\pgfqpoint{5.400000in}{5.400000in}}%
\pgfusepath{clip}%
\pgfsetbuttcap%
\pgfsetroundjoin%
\pgfsetlinewidth{1.003750pt}%
\definecolor{currentstroke}{rgb}{0.768627,0.603922,0.000000}%
\pgfsetstrokecolor{currentstroke}%
\pgfsetstrokeopacity{0.900000}%
\pgfsetdash{}{0pt}%
\pgfpathmoveto{\pgfqpoint{4.125473in}{7.558326in}}%
\pgfpathcurveto{\pgfqpoint{4.135727in}{7.558326in}}{\pgfqpoint{4.145563in}{7.562400in}}{\pgfqpoint{4.152814in}{7.569651in}}%
\pgfpathcurveto{\pgfqpoint{4.160064in}{7.576901in}}{\pgfqpoint{4.164138in}{7.586737in}}{\pgfqpoint{4.164138in}{7.596991in}}%
\pgfpathcurveto{\pgfqpoint{4.164138in}{7.607245in}}{\pgfqpoint{4.160064in}{7.617080in}}{\pgfqpoint{4.152814in}{7.624331in}}%
\pgfpathcurveto{\pgfqpoint{4.145563in}{7.631582in}}{\pgfqpoint{4.135727in}{7.635656in}}{\pgfqpoint{4.125473in}{7.635656in}}%
\pgfpathcurveto{\pgfqpoint{4.115219in}{7.635656in}}{\pgfqpoint{4.105384in}{7.631582in}}{\pgfqpoint{4.098133in}{7.624331in}}%
\pgfpathcurveto{\pgfqpoint{4.090882in}{7.617080in}}{\pgfqpoint{4.086808in}{7.607245in}}{\pgfqpoint{4.086808in}{7.596991in}}%
\pgfpathcurveto{\pgfqpoint{4.086808in}{7.586737in}}{\pgfqpoint{4.090882in}{7.576901in}}{\pgfqpoint{4.098133in}{7.569651in}}%
\pgfpathcurveto{\pgfqpoint{4.105384in}{7.562400in}}{\pgfqpoint{4.115219in}{7.558326in}}{\pgfqpoint{4.125473in}{7.558326in}}%
\pgfpathlineto{\pgfqpoint{4.125473in}{7.558326in}}%
\pgfpathclose%
\pgfusepath{stroke}%
\end{pgfscope}%
\begin{pgfscope}%
\pgfpathrectangle{\pgfqpoint{1.109909in}{4.050000in}}{\pgfqpoint{5.400000in}{5.400000in}}%
\pgfusepath{clip}%
\pgfsetbuttcap%
\pgfsetroundjoin%
\pgfsetlinewidth{1.003750pt}%
\definecolor{currentstroke}{rgb}{0.768627,0.603922,0.000000}%
\pgfsetstrokecolor{currentstroke}%
\pgfsetstrokeopacity{0.900000}%
\pgfsetdash{}{0pt}%
\pgfpathmoveto{\pgfqpoint{5.900179in}{7.791661in}}%
\pgfpathcurveto{\pgfqpoint{5.910433in}{7.791661in}}{\pgfqpoint{5.920268in}{7.795735in}}{\pgfqpoint{5.927519in}{7.802986in}}%
\pgfpathcurveto{\pgfqpoint{5.934770in}{7.810237in}}{\pgfqpoint{5.938844in}{7.820072in}}{\pgfqpoint{5.938844in}{7.830327in}}%
\pgfpathcurveto{\pgfqpoint{5.938844in}{7.840581in}}{\pgfqpoint{5.934770in}{7.850416in}}{\pgfqpoint{5.927519in}{7.857667in}}%
\pgfpathcurveto{\pgfqpoint{5.920268in}{7.864918in}}{\pgfqpoint{5.910433in}{7.868992in}}{\pgfqpoint{5.900179in}{7.868992in}}%
\pgfpathcurveto{\pgfqpoint{5.889925in}{7.868992in}}{\pgfqpoint{5.880089in}{7.864918in}}{\pgfqpoint{5.872838in}{7.857667in}}%
\pgfpathcurveto{\pgfqpoint{5.865588in}{7.850416in}}{\pgfqpoint{5.861514in}{7.840581in}}{\pgfqpoint{5.861514in}{7.830327in}}%
\pgfpathcurveto{\pgfqpoint{5.861514in}{7.820072in}}{\pgfqpoint{5.865588in}{7.810237in}}{\pgfqpoint{5.872838in}{7.802986in}}%
\pgfpathcurveto{\pgfqpoint{5.880089in}{7.795735in}}{\pgfqpoint{5.889925in}{7.791661in}}{\pgfqpoint{5.900179in}{7.791661in}}%
\pgfpathlineto{\pgfqpoint{5.900179in}{7.791661in}}%
\pgfpathclose%
\pgfusepath{stroke}%
\end{pgfscope}%
\begin{pgfscope}%
\pgfpathrectangle{\pgfqpoint{1.109909in}{4.050000in}}{\pgfqpoint{5.400000in}{5.400000in}}%
\pgfusepath{clip}%
\pgfsetbuttcap%
\pgfsetroundjoin%
\pgfsetlinewidth{1.003750pt}%
\definecolor{currentstroke}{rgb}{0.768627,0.603922,0.000000}%
\pgfsetstrokecolor{currentstroke}%
\pgfsetstrokeopacity{0.900000}%
\pgfsetdash{}{0pt}%
\pgfpathmoveto{\pgfqpoint{5.853300in}{7.967190in}}%
\pgfpathcurveto{\pgfqpoint{5.863554in}{7.967190in}}{\pgfqpoint{5.873390in}{7.971264in}}{\pgfqpoint{5.880640in}{7.978515in}}%
\pgfpathcurveto{\pgfqpoint{5.887891in}{7.985766in}}{\pgfqpoint{5.891965in}{7.995601in}}{\pgfqpoint{5.891965in}{8.005855in}}%
\pgfpathcurveto{\pgfqpoint{5.891965in}{8.016109in}}{\pgfqpoint{5.887891in}{8.025945in}}{\pgfqpoint{5.880640in}{8.033195in}}%
\pgfpathcurveto{\pgfqpoint{5.873390in}{8.040446in}}{\pgfqpoint{5.863554in}{8.044520in}}{\pgfqpoint{5.853300in}{8.044520in}}%
\pgfpathcurveto{\pgfqpoint{5.843046in}{8.044520in}}{\pgfqpoint{5.833211in}{8.040446in}}{\pgfqpoint{5.825960in}{8.033195in}}%
\pgfpathcurveto{\pgfqpoint{5.818709in}{8.025945in}}{\pgfqpoint{5.814635in}{8.016109in}}{\pgfqpoint{5.814635in}{8.005855in}}%
\pgfpathcurveto{\pgfqpoint{5.814635in}{7.995601in}}{\pgfqpoint{5.818709in}{7.985766in}}{\pgfqpoint{5.825960in}{7.978515in}}%
\pgfpathcurveto{\pgfqpoint{5.833211in}{7.971264in}}{\pgfqpoint{5.843046in}{7.967190in}}{\pgfqpoint{5.853300in}{7.967190in}}%
\pgfpathlineto{\pgfqpoint{5.853300in}{7.967190in}}%
\pgfpathclose%
\pgfusepath{stroke}%
\end{pgfscope}%
\begin{pgfscope}%
\pgfpathrectangle{\pgfqpoint{1.109909in}{4.050000in}}{\pgfqpoint{5.400000in}{5.400000in}}%
\pgfusepath{clip}%
\pgfsetbuttcap%
\pgfsetroundjoin%
\pgfsetlinewidth{1.003750pt}%
\definecolor{currentstroke}{rgb}{0.768627,0.603922,0.000000}%
\pgfsetstrokecolor{currentstroke}%
\pgfsetstrokeopacity{0.900000}%
\pgfsetdash{}{0pt}%
\pgfpathmoveto{\pgfqpoint{5.697154in}{7.893921in}}%
\pgfpathcurveto{\pgfqpoint{5.707408in}{7.893921in}}{\pgfqpoint{5.717243in}{7.897995in}}{\pgfqpoint{5.724494in}{7.905246in}}%
\pgfpathcurveto{\pgfqpoint{5.731745in}{7.912496in}}{\pgfqpoint{5.735819in}{7.922332in}}{\pgfqpoint{5.735819in}{7.932586in}}%
\pgfpathcurveto{\pgfqpoint{5.735819in}{7.942840in}}{\pgfqpoint{5.731745in}{7.952675in}}{\pgfqpoint{5.724494in}{7.959926in}}%
\pgfpathcurveto{\pgfqpoint{5.717243in}{7.967177in}}{\pgfqpoint{5.707408in}{7.971251in}}{\pgfqpoint{5.697154in}{7.971251in}}%
\pgfpathcurveto{\pgfqpoint{5.686900in}{7.971251in}}{\pgfqpoint{5.677064in}{7.967177in}}{\pgfqpoint{5.669813in}{7.959926in}}%
\pgfpathcurveto{\pgfqpoint{5.662563in}{7.952675in}}{\pgfqpoint{5.658489in}{7.942840in}}{\pgfqpoint{5.658489in}{7.932586in}}%
\pgfpathcurveto{\pgfqpoint{5.658489in}{7.922332in}}{\pgfqpoint{5.662563in}{7.912496in}}{\pgfqpoint{5.669813in}{7.905246in}}%
\pgfpathcurveto{\pgfqpoint{5.677064in}{7.897995in}}{\pgfqpoint{5.686900in}{7.893921in}}{\pgfqpoint{5.697154in}{7.893921in}}%
\pgfpathlineto{\pgfqpoint{5.697154in}{7.893921in}}%
\pgfpathclose%
\pgfusepath{stroke}%
\end{pgfscope}%
\begin{pgfscope}%
\pgfpathrectangle{\pgfqpoint{1.109909in}{4.050000in}}{\pgfqpoint{5.400000in}{5.400000in}}%
\pgfusepath{clip}%
\pgfsetbuttcap%
\pgfsetroundjoin%
\pgfsetlinewidth{1.003750pt}%
\definecolor{currentstroke}{rgb}{0.768627,0.603922,0.000000}%
\pgfsetstrokecolor{currentstroke}%
\pgfsetstrokeopacity{0.900000}%
\pgfsetdash{}{0pt}%
\pgfpathmoveto{\pgfqpoint{5.421457in}{7.470901in}}%
\pgfpathcurveto{\pgfqpoint{5.431711in}{7.470901in}}{\pgfqpoint{5.441546in}{7.474975in}}{\pgfqpoint{5.448797in}{7.482226in}}%
\pgfpathcurveto{\pgfqpoint{5.456048in}{7.489476in}}{\pgfqpoint{5.460122in}{7.499312in}}{\pgfqpoint{5.460122in}{7.509566in}}%
\pgfpathcurveto{\pgfqpoint{5.460122in}{7.519820in}}{\pgfqpoint{5.456048in}{7.529655in}}{\pgfqpoint{5.448797in}{7.536906in}}%
\pgfpathcurveto{\pgfqpoint{5.441546in}{7.544157in}}{\pgfqpoint{5.431711in}{7.548231in}}{\pgfqpoint{5.421457in}{7.548231in}}%
\pgfpathcurveto{\pgfqpoint{5.411203in}{7.548231in}}{\pgfqpoint{5.401367in}{7.544157in}}{\pgfqpoint{5.394116in}{7.536906in}}%
\pgfpathcurveto{\pgfqpoint{5.386866in}{7.529655in}}{\pgfqpoint{5.382792in}{7.519820in}}{\pgfqpoint{5.382792in}{7.509566in}}%
\pgfpathcurveto{\pgfqpoint{5.382792in}{7.499312in}}{\pgfqpoint{5.386866in}{7.489476in}}{\pgfqpoint{5.394116in}{7.482226in}}%
\pgfpathcurveto{\pgfqpoint{5.401367in}{7.474975in}}{\pgfqpoint{5.411203in}{7.470901in}}{\pgfqpoint{5.421457in}{7.470901in}}%
\pgfpathlineto{\pgfqpoint{5.421457in}{7.470901in}}%
\pgfpathclose%
\pgfusepath{stroke}%
\end{pgfscope}%
\begin{pgfscope}%
\pgfsetbuttcap%
\pgfsetroundjoin%
\pgfsetlinewidth{1.003750pt}%
\definecolor{currentstroke}{rgb}{0.768627,0.603922,0.000000}%
\pgfsetstrokecolor{currentstroke}%
\pgfsetstrokeopacity{0.900000}%
\pgfsetdash{}{0pt}%
\pgfpathmoveto{\pgfqpoint{1.109909in}{4.900495in}}%
\pgfpathlineto{\pgfqpoint{1.169647in}{4.960234in}}%
\pgfpathlineto{\pgfqpoint{1.109909in}{5.019972in}}%
\pgfpathlineto{\pgfqpoint{1.050171in}{4.960234in}}%
\pgfpathlineto{\pgfqpoint{1.109909in}{4.900495in}}%
\pgfpathclose%
\pgfusepath{stroke}%
\end{pgfscope}%
\begin{pgfscope}%
\pgfpathrectangle{\pgfqpoint{1.109909in}{4.050000in}}{\pgfqpoint{5.400000in}{5.400000in}}%
\pgfusepath{clip}%
\pgfsetbuttcap%
\pgfsetroundjoin%
\pgfsetlinewidth{1.003750pt}%
\definecolor{currentstroke}{rgb}{0.200000,0.133333,0.533333}%
\pgfsetstrokecolor{currentstroke}%
\pgfsetstrokeopacity{0.900000}%
\pgfsetdash{}{0pt}%
\pgfpathmoveto{\pgfqpoint{2.056079in}{4.717590in}}%
\pgfpathcurveto{\pgfqpoint{2.066333in}{4.717590in}}{\pgfqpoint{2.076168in}{4.721664in}}{\pgfqpoint{2.083419in}{4.728915in}}%
\pgfpathcurveto{\pgfqpoint{2.090670in}{4.736166in}}{\pgfqpoint{2.094744in}{4.746001in}}{\pgfqpoint{2.094744in}{4.756255in}}%
\pgfpathcurveto{\pgfqpoint{2.094744in}{4.766510in}}{\pgfqpoint{2.090670in}{4.776345in}}{\pgfqpoint{2.083419in}{4.783596in}}%
\pgfpathcurveto{\pgfqpoint{2.076168in}{4.790846in}}{\pgfqpoint{2.066333in}{4.794920in}}{\pgfqpoint{2.056079in}{4.794920in}}%
\pgfpathcurveto{\pgfqpoint{2.045825in}{4.794920in}}{\pgfqpoint{2.035989in}{4.790846in}}{\pgfqpoint{2.028738in}{4.783596in}}%
\pgfpathcurveto{\pgfqpoint{2.021488in}{4.776345in}}{\pgfqpoint{2.017414in}{4.766510in}}{\pgfqpoint{2.017414in}{4.756255in}}%
\pgfpathcurveto{\pgfqpoint{2.017414in}{4.746001in}}{\pgfqpoint{2.021488in}{4.736166in}}{\pgfqpoint{2.028738in}{4.728915in}}%
\pgfpathcurveto{\pgfqpoint{2.035989in}{4.721664in}}{\pgfqpoint{2.045825in}{4.717590in}}{\pgfqpoint{2.056079in}{4.717590in}}%
\pgfpathlineto{\pgfqpoint{2.056079in}{4.717590in}}%
\pgfpathclose%
\pgfusepath{stroke}%
\end{pgfscope}%
\begin{pgfscope}%
\pgfpathrectangle{\pgfqpoint{1.109909in}{4.050000in}}{\pgfqpoint{5.400000in}{5.400000in}}%
\pgfusepath{clip}%
\pgfsetbuttcap%
\pgfsetroundjoin%
\pgfsetlinewidth{1.003750pt}%
\definecolor{currentstroke}{rgb}{0.200000,0.133333,0.533333}%
\pgfsetstrokecolor{currentstroke}%
\pgfsetstrokeopacity{0.900000}%
\pgfsetdash{}{0pt}%
\pgfpathmoveto{\pgfqpoint{2.264572in}{4.728543in}}%
\pgfpathcurveto{\pgfqpoint{2.274826in}{4.728543in}}{\pgfqpoint{2.284662in}{4.732616in}}{\pgfqpoint{2.291912in}{4.739867in}}%
\pgfpathcurveto{\pgfqpoint{2.299163in}{4.747118in}}{\pgfqpoint{2.303237in}{4.756953in}}{\pgfqpoint{2.303237in}{4.767208in}}%
\pgfpathcurveto{\pgfqpoint{2.303237in}{4.777462in}}{\pgfqpoint{2.299163in}{4.787297in}}{\pgfqpoint{2.291912in}{4.794548in}}%
\pgfpathcurveto{\pgfqpoint{2.284662in}{4.801799in}}{\pgfqpoint{2.274826in}{4.805873in}}{\pgfqpoint{2.264572in}{4.805873in}}%
\pgfpathcurveto{\pgfqpoint{2.254318in}{4.805873in}}{\pgfqpoint{2.244483in}{4.801799in}}{\pgfqpoint{2.237232in}{4.794548in}}%
\pgfpathcurveto{\pgfqpoint{2.229981in}{4.787297in}}{\pgfqpoint{2.225907in}{4.777462in}}{\pgfqpoint{2.225907in}{4.767208in}}%
\pgfpathcurveto{\pgfqpoint{2.225907in}{4.756953in}}{\pgfqpoint{2.229981in}{4.747118in}}{\pgfqpoint{2.237232in}{4.739867in}}%
\pgfpathcurveto{\pgfqpoint{2.244483in}{4.732616in}}{\pgfqpoint{2.254318in}{4.728543in}}{\pgfqpoint{2.264572in}{4.728543in}}%
\pgfpathlineto{\pgfqpoint{2.264572in}{4.728543in}}%
\pgfpathclose%
\pgfusepath{stroke}%
\end{pgfscope}%
\begin{pgfscope}%
\pgfpathrectangle{\pgfqpoint{1.109909in}{4.050000in}}{\pgfqpoint{5.400000in}{5.400000in}}%
\pgfusepath{clip}%
\pgfsetbuttcap%
\pgfsetroundjoin%
\pgfsetlinewidth{1.003750pt}%
\definecolor{currentstroke}{rgb}{0.200000,0.133333,0.533333}%
\pgfsetstrokecolor{currentstroke}%
\pgfsetstrokeopacity{0.900000}%
\pgfsetdash{}{0pt}%
\pgfpathmoveto{\pgfqpoint{2.286435in}{4.719296in}}%
\pgfpathcurveto{\pgfqpoint{2.296689in}{4.719296in}}{\pgfqpoint{2.306524in}{4.723370in}}{\pgfqpoint{2.313775in}{4.730621in}}%
\pgfpathcurveto{\pgfqpoint{2.321026in}{4.737871in}}{\pgfqpoint{2.325100in}{4.747707in}}{\pgfqpoint{2.325100in}{4.757961in}}%
\pgfpathcurveto{\pgfqpoint{2.325100in}{4.768215in}}{\pgfqpoint{2.321026in}{4.778051in}}{\pgfqpoint{2.313775in}{4.785301in}}%
\pgfpathcurveto{\pgfqpoint{2.306524in}{4.792552in}}{\pgfqpoint{2.296689in}{4.796626in}}{\pgfqpoint{2.286435in}{4.796626in}}%
\pgfpathcurveto{\pgfqpoint{2.276180in}{4.796626in}}{\pgfqpoint{2.266345in}{4.792552in}}{\pgfqpoint{2.259094in}{4.785301in}}%
\pgfpathcurveto{\pgfqpoint{2.251843in}{4.778051in}}{\pgfqpoint{2.247769in}{4.768215in}}{\pgfqpoint{2.247769in}{4.757961in}}%
\pgfpathcurveto{\pgfqpoint{2.247769in}{4.747707in}}{\pgfqpoint{2.251843in}{4.737871in}}{\pgfqpoint{2.259094in}{4.730621in}}%
\pgfpathcurveto{\pgfqpoint{2.266345in}{4.723370in}}{\pgfqpoint{2.276180in}{4.719296in}}{\pgfqpoint{2.286435in}{4.719296in}}%
\pgfpathlineto{\pgfqpoint{2.286435in}{4.719296in}}%
\pgfpathclose%
\pgfusepath{stroke}%
\end{pgfscope}%
\begin{pgfscope}%
\pgfpathrectangle{\pgfqpoint{1.109909in}{4.050000in}}{\pgfqpoint{5.400000in}{5.400000in}}%
\pgfusepath{clip}%
\pgfsetbuttcap%
\pgfsetroundjoin%
\pgfsetlinewidth{1.003750pt}%
\definecolor{currentstroke}{rgb}{0.200000,0.133333,0.533333}%
\pgfsetstrokecolor{currentstroke}%
\pgfsetstrokeopacity{0.900000}%
\pgfsetdash{}{0pt}%
\pgfpathmoveto{\pgfqpoint{2.866694in}{5.223825in}}%
\pgfpathcurveto{\pgfqpoint{2.876948in}{5.223825in}}{\pgfqpoint{2.886784in}{5.227899in}}{\pgfqpoint{2.894035in}{5.235150in}}%
\pgfpathcurveto{\pgfqpoint{2.901285in}{5.242401in}}{\pgfqpoint{2.905359in}{5.252236in}}{\pgfqpoint{2.905359in}{5.262490in}}%
\pgfpathcurveto{\pgfqpoint{2.905359in}{5.272745in}}{\pgfqpoint{2.901285in}{5.282580in}}{\pgfqpoint{2.894035in}{5.289831in}}%
\pgfpathcurveto{\pgfqpoint{2.886784in}{5.297082in}}{\pgfqpoint{2.876948in}{5.301155in}}{\pgfqpoint{2.866694in}{5.301155in}}%
\pgfpathcurveto{\pgfqpoint{2.856440in}{5.301155in}}{\pgfqpoint{2.846605in}{5.297082in}}{\pgfqpoint{2.839354in}{5.289831in}}%
\pgfpathcurveto{\pgfqpoint{2.832103in}{5.282580in}}{\pgfqpoint{2.828029in}{5.272745in}}{\pgfqpoint{2.828029in}{5.262490in}}%
\pgfpathcurveto{\pgfqpoint{2.828029in}{5.252236in}}{\pgfqpoint{2.832103in}{5.242401in}}{\pgfqpoint{2.839354in}{5.235150in}}%
\pgfpathcurveto{\pgfqpoint{2.846605in}{5.227899in}}{\pgfqpoint{2.856440in}{5.223825in}}{\pgfqpoint{2.866694in}{5.223825in}}%
\pgfpathlineto{\pgfqpoint{2.866694in}{5.223825in}}%
\pgfpathclose%
\pgfusepath{stroke}%
\end{pgfscope}%
\begin{pgfscope}%
\pgfpathrectangle{\pgfqpoint{1.109909in}{4.050000in}}{\pgfqpoint{5.400000in}{5.400000in}}%
\pgfusepath{clip}%
\pgfsetbuttcap%
\pgfsetroundjoin%
\pgfsetlinewidth{1.003750pt}%
\definecolor{currentstroke}{rgb}{0.200000,0.133333,0.533333}%
\pgfsetstrokecolor{currentstroke}%
\pgfsetstrokeopacity{0.900000}%
\pgfsetdash{}{0pt}%
\pgfpathmoveto{\pgfqpoint{2.106825in}{4.550626in}}%
\pgfpathcurveto{\pgfqpoint{2.117079in}{4.550626in}}{\pgfqpoint{2.126914in}{4.554700in}}{\pgfqpoint{2.134165in}{4.561951in}}%
\pgfpathcurveto{\pgfqpoint{2.141416in}{4.569201in}}{\pgfqpoint{2.145490in}{4.579037in}}{\pgfqpoint{2.145490in}{4.589291in}}%
\pgfpathcurveto{\pgfqpoint{2.145490in}{4.599545in}}{\pgfqpoint{2.141416in}{4.609380in}}{\pgfqpoint{2.134165in}{4.616631in}}%
\pgfpathcurveto{\pgfqpoint{2.126914in}{4.623882in}}{\pgfqpoint{2.117079in}{4.627956in}}{\pgfqpoint{2.106825in}{4.627956in}}%
\pgfpathcurveto{\pgfqpoint{2.096571in}{4.627956in}}{\pgfqpoint{2.086735in}{4.623882in}}{\pgfqpoint{2.079484in}{4.616631in}}%
\pgfpathcurveto{\pgfqpoint{2.072234in}{4.609380in}}{\pgfqpoint{2.068160in}{4.599545in}}{\pgfqpoint{2.068160in}{4.589291in}}%
\pgfpathcurveto{\pgfqpoint{2.068160in}{4.579037in}}{\pgfqpoint{2.072234in}{4.569201in}}{\pgfqpoint{2.079484in}{4.561951in}}%
\pgfpathcurveto{\pgfqpoint{2.086735in}{4.554700in}}{\pgfqpoint{2.096571in}{4.550626in}}{\pgfqpoint{2.106825in}{4.550626in}}%
\pgfpathlineto{\pgfqpoint{2.106825in}{4.550626in}}%
\pgfpathclose%
\pgfusepath{stroke}%
\end{pgfscope}%
\begin{pgfscope}%
\pgfpathrectangle{\pgfqpoint{1.109909in}{4.050000in}}{\pgfqpoint{5.400000in}{5.400000in}}%
\pgfusepath{clip}%
\pgfsetbuttcap%
\pgfsetroundjoin%
\pgfsetlinewidth{1.003750pt}%
\definecolor{currentstroke}{rgb}{0.200000,0.133333,0.533333}%
\pgfsetstrokecolor{currentstroke}%
\pgfsetstrokeopacity{0.900000}%
\pgfsetdash{}{0pt}%
\pgfpathmoveto{\pgfqpoint{2.758336in}{5.169575in}}%
\pgfpathcurveto{\pgfqpoint{2.768590in}{5.169575in}}{\pgfqpoint{2.778426in}{5.173649in}}{\pgfqpoint{2.785676in}{5.180900in}}%
\pgfpathcurveto{\pgfqpoint{2.792927in}{5.188150in}}{\pgfqpoint{2.797001in}{5.197986in}}{\pgfqpoint{2.797001in}{5.208240in}}%
\pgfpathcurveto{\pgfqpoint{2.797001in}{5.218494in}}{\pgfqpoint{2.792927in}{5.228330in}}{\pgfqpoint{2.785676in}{5.235580in}}%
\pgfpathcurveto{\pgfqpoint{2.778426in}{5.242831in}}{\pgfqpoint{2.768590in}{5.246905in}}{\pgfqpoint{2.758336in}{5.246905in}}%
\pgfpathcurveto{\pgfqpoint{2.748082in}{5.246905in}}{\pgfqpoint{2.738247in}{5.242831in}}{\pgfqpoint{2.730996in}{5.235580in}}%
\pgfpathcurveto{\pgfqpoint{2.723745in}{5.228330in}}{\pgfqpoint{2.719671in}{5.218494in}}{\pgfqpoint{2.719671in}{5.208240in}}%
\pgfpathcurveto{\pgfqpoint{2.719671in}{5.197986in}}{\pgfqpoint{2.723745in}{5.188150in}}{\pgfqpoint{2.730996in}{5.180900in}}%
\pgfpathcurveto{\pgfqpoint{2.738247in}{5.173649in}}{\pgfqpoint{2.748082in}{5.169575in}}{\pgfqpoint{2.758336in}{5.169575in}}%
\pgfpathlineto{\pgfqpoint{2.758336in}{5.169575in}}%
\pgfpathclose%
\pgfusepath{stroke}%
\end{pgfscope}%
\begin{pgfscope}%
\pgfpathrectangle{\pgfqpoint{1.109909in}{4.050000in}}{\pgfqpoint{5.400000in}{5.400000in}}%
\pgfusepath{clip}%
\pgfsetbuttcap%
\pgfsetroundjoin%
\pgfsetlinewidth{1.003750pt}%
\definecolor{currentstroke}{rgb}{0.176471,0.713725,0.643137}%
\pgfsetstrokecolor{currentstroke}%
\pgfsetstrokeopacity{0.900000}%
\pgfsetdash{}{0pt}%
\pgfpathmoveto{\pgfqpoint{5.906232in}{8.448250in}}%
\pgfpathcurveto{\pgfqpoint{5.916486in}{8.448250in}}{\pgfqpoint{5.926322in}{8.452324in}}{\pgfqpoint{5.933572in}{8.459575in}}%
\pgfpathcurveto{\pgfqpoint{5.940823in}{8.466826in}}{\pgfqpoint{5.944897in}{8.476661in}}{\pgfqpoint{5.944897in}{8.486915in}}%
\pgfpathcurveto{\pgfqpoint{5.944897in}{8.497169in}}{\pgfqpoint{5.940823in}{8.507005in}}{\pgfqpoint{5.933572in}{8.514256in}}%
\pgfpathcurveto{\pgfqpoint{5.926322in}{8.521506in}}{\pgfqpoint{5.916486in}{8.525580in}}{\pgfqpoint{5.906232in}{8.525580in}}%
\pgfpathcurveto{\pgfqpoint{5.895978in}{8.525580in}}{\pgfqpoint{5.886143in}{8.521506in}}{\pgfqpoint{5.878892in}{8.514256in}}%
\pgfpathcurveto{\pgfqpoint{5.871641in}{8.507005in}}{\pgfqpoint{5.867567in}{8.497169in}}{\pgfqpoint{5.867567in}{8.486915in}}%
\pgfpathcurveto{\pgfqpoint{5.867567in}{8.476661in}}{\pgfqpoint{5.871641in}{8.466826in}}{\pgfqpoint{5.878892in}{8.459575in}}%
\pgfpathcurveto{\pgfqpoint{5.886143in}{8.452324in}}{\pgfqpoint{5.895978in}{8.448250in}}{\pgfqpoint{5.906232in}{8.448250in}}%
\pgfpathlineto{\pgfqpoint{5.906232in}{8.448250in}}%
\pgfpathclose%
\pgfusepath{stroke}%
\end{pgfscope}%
\begin{pgfscope}%
\pgfpathrectangle{\pgfqpoint{1.109909in}{4.050000in}}{\pgfqpoint{5.400000in}{5.400000in}}%
\pgfusepath{clip}%
\pgfsetbuttcap%
\pgfsetroundjoin%
\pgfsetlinewidth{1.003750pt}%
\definecolor{currentstroke}{rgb}{0.176471,0.713725,0.643137}%
\pgfsetstrokecolor{currentstroke}%
\pgfsetstrokeopacity{0.900000}%
\pgfsetdash{}{0pt}%
\pgfpathmoveto{\pgfqpoint{5.882356in}{8.677178in}}%
\pgfpathcurveto{\pgfqpoint{5.892610in}{8.677178in}}{\pgfqpoint{5.902445in}{8.681252in}}{\pgfqpoint{5.909696in}{8.688503in}}%
\pgfpathcurveto{\pgfqpoint{5.916947in}{8.695753in}}{\pgfqpoint{5.921021in}{8.705589in}}{\pgfqpoint{5.921021in}{8.715843in}}%
\pgfpathcurveto{\pgfqpoint{5.921021in}{8.726097in}}{\pgfqpoint{5.916947in}{8.735933in}}{\pgfqpoint{5.909696in}{8.743183in}}%
\pgfpathcurveto{\pgfqpoint{5.902445in}{8.750434in}}{\pgfqpoint{5.892610in}{8.754508in}}{\pgfqpoint{5.882356in}{8.754508in}}%
\pgfpathcurveto{\pgfqpoint{5.872102in}{8.754508in}}{\pgfqpoint{5.862266in}{8.750434in}}{\pgfqpoint{5.855015in}{8.743183in}}%
\pgfpathcurveto{\pgfqpoint{5.847765in}{8.735933in}}{\pgfqpoint{5.843691in}{8.726097in}}{\pgfqpoint{5.843691in}{8.715843in}}%
\pgfpathcurveto{\pgfqpoint{5.843691in}{8.705589in}}{\pgfqpoint{5.847765in}{8.695753in}}{\pgfqpoint{5.855015in}{8.688503in}}%
\pgfpathcurveto{\pgfqpoint{5.862266in}{8.681252in}}{\pgfqpoint{5.872102in}{8.677178in}}{\pgfqpoint{5.882356in}{8.677178in}}%
\pgfpathlineto{\pgfqpoint{5.882356in}{8.677178in}}%
\pgfpathclose%
\pgfusepath{stroke}%
\end{pgfscope}%
\begin{pgfscope}%
\pgfpathrectangle{\pgfqpoint{1.109909in}{4.050000in}}{\pgfqpoint{5.400000in}{5.400000in}}%
\pgfusepath{clip}%
\pgfsetbuttcap%
\pgfsetroundjoin%
\pgfsetlinewidth{1.003750pt}%
\definecolor{currentstroke}{rgb}{0.176471,0.713725,0.643137}%
\pgfsetstrokecolor{currentstroke}%
\pgfsetstrokeopacity{0.900000}%
\pgfsetdash{}{0pt}%
\pgfpathmoveto{\pgfqpoint{5.940281in}{8.645761in}}%
\pgfpathcurveto{\pgfqpoint{5.950536in}{8.645761in}}{\pgfqpoint{5.960371in}{8.649835in}}{\pgfqpoint{5.967622in}{8.657086in}}%
\pgfpathcurveto{\pgfqpoint{5.974872in}{8.664337in}}{\pgfqpoint{5.978946in}{8.674172in}}{\pgfqpoint{5.978946in}{8.684426in}}%
\pgfpathcurveto{\pgfqpoint{5.978946in}{8.694680in}}{\pgfqpoint{5.974872in}{8.704516in}}{\pgfqpoint{5.967622in}{8.711767in}}%
\pgfpathcurveto{\pgfqpoint{5.960371in}{8.719017in}}{\pgfqpoint{5.950536in}{8.723091in}}{\pgfqpoint{5.940281in}{8.723091in}}%
\pgfpathcurveto{\pgfqpoint{5.930027in}{8.723091in}}{\pgfqpoint{5.920192in}{8.719017in}}{\pgfqpoint{5.912941in}{8.711767in}}%
\pgfpathcurveto{\pgfqpoint{5.905690in}{8.704516in}}{\pgfqpoint{5.901616in}{8.694680in}}{\pgfqpoint{5.901616in}{8.684426in}}%
\pgfpathcurveto{\pgfqpoint{5.901616in}{8.674172in}}{\pgfqpoint{5.905690in}{8.664337in}}{\pgfqpoint{5.912941in}{8.657086in}}%
\pgfpathcurveto{\pgfqpoint{5.920192in}{8.649835in}}{\pgfqpoint{5.930027in}{8.645761in}}{\pgfqpoint{5.940281in}{8.645761in}}%
\pgfpathlineto{\pgfqpoint{5.940281in}{8.645761in}}%
\pgfpathclose%
\pgfusepath{stroke}%
\end{pgfscope}%
\begin{pgfscope}%
\pgfsetbuttcap%
\pgfsetroundjoin%
\pgfsetlinewidth{1.003750pt}%
\definecolor{currentstroke}{rgb}{0.176471,0.713725,0.643137}%
\pgfsetstrokecolor{currentstroke}%
\pgfsetstrokeopacity{0.900000}%
\pgfsetdash{}{0pt}%
\pgfpathmoveto{\pgfqpoint{6.179266in}{9.390262in}}%
\pgfpathlineto{\pgfqpoint{6.239004in}{9.450000in}}%
\pgfpathlineto{\pgfqpoint{6.179266in}{9.509738in}}%
\pgfpathlineto{\pgfqpoint{6.119527in}{9.450000in}}%
\pgfpathlineto{\pgfqpoint{6.179266in}{9.390262in}}%
\pgfpathclose%
\pgfusepath{stroke}%
\end{pgfscope}%
\begin{pgfscope}%
\pgfpathrectangle{\pgfqpoint{1.109909in}{4.050000in}}{\pgfqpoint{5.400000in}{5.400000in}}%
\pgfusepath{clip}%
\pgfsetbuttcap%
\pgfsetroundjoin%
\pgfsetlinewidth{1.003750pt}%
\definecolor{currentstroke}{rgb}{0.000000,0.447059,0.698039}%
\pgfsetstrokecolor{currentstroke}%
\pgfsetstrokeopacity{0.900000}%
\pgfsetdash{}{0pt}%
\pgfpathmoveto{\pgfqpoint{3.722823in}{7.038932in}}%
\pgfpathcurveto{\pgfqpoint{3.733078in}{7.038932in}}{\pgfqpoint{3.742913in}{7.043006in}}{\pgfqpoint{3.750164in}{7.050257in}}%
\pgfpathcurveto{\pgfqpoint{3.757414in}{7.057507in}}{\pgfqpoint{3.761488in}{7.067343in}}{\pgfqpoint{3.761488in}{7.077597in}}%
\pgfpathcurveto{\pgfqpoint{3.761488in}{7.087851in}}{\pgfqpoint{3.757414in}{7.097686in}}{\pgfqpoint{3.750164in}{7.104937in}}%
\pgfpathcurveto{\pgfqpoint{3.742913in}{7.112188in}}{\pgfqpoint{3.733078in}{7.116262in}}{\pgfqpoint{3.722823in}{7.116262in}}%
\pgfpathcurveto{\pgfqpoint{3.712569in}{7.116262in}}{\pgfqpoint{3.702734in}{7.112188in}}{\pgfqpoint{3.695483in}{7.104937in}}%
\pgfpathcurveto{\pgfqpoint{3.688232in}{7.097686in}}{\pgfqpoint{3.684158in}{7.087851in}}{\pgfqpoint{3.684158in}{7.077597in}}%
\pgfpathcurveto{\pgfqpoint{3.684158in}{7.067343in}}{\pgfqpoint{3.688232in}{7.057507in}}{\pgfqpoint{3.695483in}{7.050257in}}%
\pgfpathcurveto{\pgfqpoint{3.702734in}{7.043006in}}{\pgfqpoint{3.712569in}{7.038932in}}{\pgfqpoint{3.722823in}{7.038932in}}%
\pgfpathlineto{\pgfqpoint{3.722823in}{7.038932in}}%
\pgfpathclose%
\pgfusepath{stroke}%
\end{pgfscope}%
\begin{pgfscope}%
\pgfpathrectangle{\pgfqpoint{1.109909in}{4.050000in}}{\pgfqpoint{5.400000in}{5.400000in}}%
\pgfusepath{clip}%
\pgfsetbuttcap%
\pgfsetroundjoin%
\pgfsetlinewidth{1.003750pt}%
\definecolor{currentstroke}{rgb}{0.901961,0.403922,0.619608}%
\pgfsetstrokecolor{currentstroke}%
\pgfsetstrokeopacity{0.900000}%
\pgfsetdash{}{0pt}%
\pgfpathmoveto{\pgfqpoint{3.280281in}{4.364849in}}%
\pgfpathcurveto{\pgfqpoint{3.290535in}{4.364849in}}{\pgfqpoint{3.300370in}{4.368923in}}{\pgfqpoint{3.307621in}{4.376174in}}%
\pgfpathcurveto{\pgfqpoint{3.314872in}{4.383424in}}{\pgfqpoint{3.318946in}{4.393260in}}{\pgfqpoint{3.318946in}{4.403514in}}%
\pgfpathcurveto{\pgfqpoint{3.318946in}{4.413768in}}{\pgfqpoint{3.314872in}{4.423603in}}{\pgfqpoint{3.307621in}{4.430854in}}%
\pgfpathcurveto{\pgfqpoint{3.300370in}{4.438105in}}{\pgfqpoint{3.290535in}{4.442179in}}{\pgfqpoint{3.280281in}{4.442179in}}%
\pgfpathcurveto{\pgfqpoint{3.270027in}{4.442179in}}{\pgfqpoint{3.260191in}{4.438105in}}{\pgfqpoint{3.252941in}{4.430854in}}%
\pgfpathcurveto{\pgfqpoint{3.245690in}{4.423603in}}{\pgfqpoint{3.241616in}{4.413768in}}{\pgfqpoint{3.241616in}{4.403514in}}%
\pgfpathcurveto{\pgfqpoint{3.241616in}{4.393260in}}{\pgfqpoint{3.245690in}{4.383424in}}{\pgfqpoint{3.252941in}{4.376174in}}%
\pgfpathcurveto{\pgfqpoint{3.260191in}{4.368923in}}{\pgfqpoint{3.270027in}{4.364849in}}{\pgfqpoint{3.280281in}{4.364849in}}%
\pgfpathlineto{\pgfqpoint{3.280281in}{4.364849in}}%
\pgfpathclose%
\pgfusepath{stroke}%
\end{pgfscope}%
\begin{pgfscope}%
\pgfpathrectangle{\pgfqpoint{1.109909in}{4.050000in}}{\pgfqpoint{5.400000in}{5.400000in}}%
\pgfusepath{clip}%
\pgfsetbuttcap%
\pgfsetroundjoin%
\pgfsetlinewidth{1.003750pt}%
\definecolor{currentstroke}{rgb}{0.901961,0.403922,0.619608}%
\pgfsetstrokecolor{currentstroke}%
\pgfsetstrokeopacity{0.900000}%
\pgfsetdash{}{0pt}%
\pgfpathmoveto{\pgfqpoint{4.932680in}{8.278461in}}%
\pgfpathcurveto{\pgfqpoint{4.942934in}{8.278461in}}{\pgfqpoint{4.952769in}{8.282535in}}{\pgfqpoint{4.960020in}{8.289786in}}%
\pgfpathcurveto{\pgfqpoint{4.967271in}{8.297037in}}{\pgfqpoint{4.971345in}{8.306872in}}{\pgfqpoint{4.971345in}{8.317126in}}%
\pgfpathcurveto{\pgfqpoint{4.971345in}{8.327380in}}{\pgfqpoint{4.967271in}{8.337216in}}{\pgfqpoint{4.960020in}{8.344466in}}%
\pgfpathcurveto{\pgfqpoint{4.952769in}{8.351717in}}{\pgfqpoint{4.942934in}{8.355791in}}{\pgfqpoint{4.932680in}{8.355791in}}%
\pgfpathcurveto{\pgfqpoint{4.922425in}{8.355791in}}{\pgfqpoint{4.912590in}{8.351717in}}{\pgfqpoint{4.905339in}{8.344466in}}%
\pgfpathcurveto{\pgfqpoint{4.898088in}{8.337216in}}{\pgfqpoint{4.894014in}{8.327380in}}{\pgfqpoint{4.894014in}{8.317126in}}%
\pgfpathcurveto{\pgfqpoint{4.894014in}{8.306872in}}{\pgfqpoint{4.898088in}{8.297037in}}{\pgfqpoint{4.905339in}{8.289786in}}%
\pgfpathcurveto{\pgfqpoint{4.912590in}{8.282535in}}{\pgfqpoint{4.922425in}{8.278461in}}{\pgfqpoint{4.932680in}{8.278461in}}%
\pgfpathlineto{\pgfqpoint{4.932680in}{8.278461in}}%
\pgfpathclose%
\pgfusepath{stroke}%
\end{pgfscope}%
\begin{pgfscope}%
\pgfpathrectangle{\pgfqpoint{1.109909in}{4.050000in}}{\pgfqpoint{5.400000in}{5.400000in}}%
\pgfusepath{clip}%
\pgfsetbuttcap%
\pgfsetroundjoin%
\pgfsetlinewidth{1.003750pt}%
\definecolor{currentstroke}{rgb}{0.901961,0.403922,0.619608}%
\pgfsetstrokecolor{currentstroke}%
\pgfsetstrokeopacity{0.900000}%
\pgfsetdash{}{0pt}%
\pgfpathmoveto{\pgfqpoint{4.891079in}{6.929789in}}%
\pgfpathcurveto{\pgfqpoint{4.901333in}{6.929789in}}{\pgfqpoint{4.911168in}{6.933863in}}{\pgfqpoint{4.918419in}{6.941114in}}%
\pgfpathcurveto{\pgfqpoint{4.925670in}{6.948364in}}{\pgfqpoint{4.929744in}{6.958200in}}{\pgfqpoint{4.929744in}{6.968454in}}%
\pgfpathcurveto{\pgfqpoint{4.929744in}{6.978708in}}{\pgfqpoint{4.925670in}{6.988543in}}{\pgfqpoint{4.918419in}{6.995794in}}%
\pgfpathcurveto{\pgfqpoint{4.911168in}{7.003045in}}{\pgfqpoint{4.901333in}{7.007119in}}{\pgfqpoint{4.891079in}{7.007119in}}%
\pgfpathcurveto{\pgfqpoint{4.880825in}{7.007119in}}{\pgfqpoint{4.870989in}{7.003045in}}{\pgfqpoint{4.863738in}{6.995794in}}%
\pgfpathcurveto{\pgfqpoint{4.856488in}{6.988543in}}{\pgfqpoint{4.852414in}{6.978708in}}{\pgfqpoint{4.852414in}{6.968454in}}%
\pgfpathcurveto{\pgfqpoint{4.852414in}{6.958200in}}{\pgfqpoint{4.856488in}{6.948364in}}{\pgfqpoint{4.863738in}{6.941114in}}%
\pgfpathcurveto{\pgfqpoint{4.870989in}{6.933863in}}{\pgfqpoint{4.880825in}{6.929789in}}{\pgfqpoint{4.891079in}{6.929789in}}%
\pgfpathlineto{\pgfqpoint{4.891079in}{6.929789in}}%
\pgfpathclose%
\pgfusepath{stroke}%
\end{pgfscope}%
\begin{pgfscope}%
\pgfpathrectangle{\pgfqpoint{1.109909in}{4.050000in}}{\pgfqpoint{5.400000in}{5.400000in}}%
\pgfusepath{clip}%
\pgfsetbuttcap%
\pgfsetroundjoin%
\pgfsetlinewidth{1.003750pt}%
\definecolor{currentstroke}{rgb}{0.941176,0.584314,0.337255}%
\pgfsetstrokecolor{currentstroke}%
\pgfsetstrokeopacity{0.900000}%
\pgfsetdash{}{0pt}%
\pgfpathmoveto{\pgfqpoint{6.217054in}{8.515620in}}%
\pgfpathcurveto{\pgfqpoint{6.227308in}{8.515620in}}{\pgfqpoint{6.237144in}{8.519694in}}{\pgfqpoint{6.244395in}{8.526944in}}%
\pgfpathcurveto{\pgfqpoint{6.251645in}{8.534195in}}{\pgfqpoint{6.255719in}{8.544030in}}{\pgfqpoint{6.255719in}{8.554285in}}%
\pgfpathcurveto{\pgfqpoint{6.255719in}{8.564539in}}{\pgfqpoint{6.251645in}{8.574374in}}{\pgfqpoint{6.244395in}{8.581625in}}%
\pgfpathcurveto{\pgfqpoint{6.237144in}{8.588876in}}{\pgfqpoint{6.227308in}{8.592950in}}{\pgfqpoint{6.217054in}{8.592950in}}%
\pgfpathcurveto{\pgfqpoint{6.206800in}{8.592950in}}{\pgfqpoint{6.196965in}{8.588876in}}{\pgfqpoint{6.189714in}{8.581625in}}%
\pgfpathcurveto{\pgfqpoint{6.182463in}{8.574374in}}{\pgfqpoint{6.178389in}{8.564539in}}{\pgfqpoint{6.178389in}{8.554285in}}%
\pgfpathcurveto{\pgfqpoint{6.178389in}{8.544030in}}{\pgfqpoint{6.182463in}{8.534195in}}{\pgfqpoint{6.189714in}{8.526944in}}%
\pgfpathcurveto{\pgfqpoint{6.196965in}{8.519694in}}{\pgfqpoint{6.206800in}{8.515620in}}{\pgfqpoint{6.217054in}{8.515620in}}%
\pgfpathlineto{\pgfqpoint{6.217054in}{8.515620in}}%
\pgfpathclose%
\pgfusepath{stroke}%
\end{pgfscope}%
\begin{pgfscope}%
\pgfpathrectangle{\pgfqpoint{1.109909in}{4.050000in}}{\pgfqpoint{5.400000in}{5.400000in}}%
\pgfusepath{clip}%
\pgfsetbuttcap%
\pgfsetroundjoin%
\pgfsetlinewidth{1.003750pt}%
\definecolor{currentstroke}{rgb}{0.941176,0.584314,0.337255}%
\pgfsetstrokecolor{currentstroke}%
\pgfsetstrokeopacity{0.900000}%
\pgfsetdash{}{0pt}%
\pgfpathmoveto{\pgfqpoint{6.191285in}{8.060323in}}%
\pgfpathcurveto{\pgfqpoint{6.201539in}{8.060323in}}{\pgfqpoint{6.211375in}{8.064397in}}{\pgfqpoint{6.218626in}{8.071648in}}%
\pgfpathcurveto{\pgfqpoint{6.225876in}{8.078898in}}{\pgfqpoint{6.229950in}{8.088734in}}{\pgfqpoint{6.229950in}{8.098988in}}%
\pgfpathcurveto{\pgfqpoint{6.229950in}{8.109242in}}{\pgfqpoint{6.225876in}{8.119077in}}{\pgfqpoint{6.218626in}{8.126328in}}%
\pgfpathcurveto{\pgfqpoint{6.211375in}{8.133579in}}{\pgfqpoint{6.201539in}{8.137653in}}{\pgfqpoint{6.191285in}{8.137653in}}%
\pgfpathcurveto{\pgfqpoint{6.181031in}{8.137653in}}{\pgfqpoint{6.171196in}{8.133579in}}{\pgfqpoint{6.163945in}{8.126328in}}%
\pgfpathcurveto{\pgfqpoint{6.156694in}{8.119077in}}{\pgfqpoint{6.152620in}{8.109242in}}{\pgfqpoint{6.152620in}{8.098988in}}%
\pgfpathcurveto{\pgfqpoint{6.152620in}{8.088734in}}{\pgfqpoint{6.156694in}{8.078898in}}{\pgfqpoint{6.163945in}{8.071648in}}%
\pgfpathcurveto{\pgfqpoint{6.171196in}{8.064397in}}{\pgfqpoint{6.181031in}{8.060323in}}{\pgfqpoint{6.191285in}{8.060323in}}%
\pgfpathlineto{\pgfqpoint{6.191285in}{8.060323in}}%
\pgfpathclose%
\pgfusepath{stroke}%
\end{pgfscope}%
\begin{pgfscope}%
\pgfsetbuttcap%
\pgfsetroundjoin%
\pgfsetlinewidth{1.003750pt}%
\definecolor{currentstroke}{rgb}{0.941176,0.584314,0.337255}%
\pgfsetstrokecolor{currentstroke}%
\pgfsetstrokeopacity{0.900000}%
\pgfsetdash{}{0pt}%
\pgfpathmoveto{\pgfqpoint{6.509909in}{9.390262in}}%
\pgfpathlineto{\pgfqpoint{6.569647in}{9.450000in}}%
\pgfpathlineto{\pgfqpoint{6.509909in}{9.509738in}}%
\pgfpathlineto{\pgfqpoint{6.450171in}{9.450000in}}%
\pgfpathlineto{\pgfqpoint{6.509909in}{9.390262in}}%
\pgfpathclose%
\pgfusepath{stroke}%
\end{pgfscope}%
\begin{pgfscope}%
\pgfsetbuttcap%
\pgfsetroundjoin%
\pgfsetlinewidth{1.003750pt}%
\definecolor{currentstroke}{rgb}{0.941176,0.584314,0.337255}%
\pgfsetstrokecolor{currentstroke}%
\pgfsetstrokeopacity{0.900000}%
\pgfsetdash{}{0pt}%
\pgfpathmoveto{\pgfqpoint{6.509909in}{9.390262in}}%
\pgfpathlineto{\pgfqpoint{6.569647in}{9.450000in}}%
\pgfpathlineto{\pgfqpoint{6.509909in}{9.509738in}}%
\pgfpathlineto{\pgfqpoint{6.450171in}{9.450000in}}%
\pgfpathlineto{\pgfqpoint{6.509909in}{9.390262in}}%
\pgfpathclose%
\pgfusepath{stroke}%
\end{pgfscope}%
\begin{pgfscope}%
\pgfpathrectangle{\pgfqpoint{1.109909in}{4.050000in}}{\pgfqpoint{5.400000in}{5.400000in}}%
\pgfusepath{clip}%
\pgfsetbuttcap%
\pgfsetroundjoin%
\pgfsetlinewidth{1.003750pt}%
\definecolor{currentstroke}{rgb}{0.835294,0.368627,0.000000}%
\pgfsetstrokecolor{currentstroke}%
\pgfsetstrokeopacity{0.900000}%
\pgfsetdash{}{0pt}%
\pgfpathmoveto{\pgfqpoint{5.286424in}{7.671795in}}%
\pgfpathcurveto{\pgfqpoint{5.296678in}{7.671795in}}{\pgfqpoint{5.306513in}{7.675869in}}{\pgfqpoint{5.313764in}{7.683119in}}%
\pgfpathcurveto{\pgfqpoint{5.321015in}{7.690370in}}{\pgfqpoint{5.325089in}{7.700206in}}{\pgfqpoint{5.325089in}{7.710460in}}%
\pgfpathcurveto{\pgfqpoint{5.325089in}{7.720714in}}{\pgfqpoint{5.321015in}{7.730549in}}{\pgfqpoint{5.313764in}{7.737800in}}%
\pgfpathcurveto{\pgfqpoint{5.306513in}{7.745051in}}{\pgfqpoint{5.296678in}{7.749125in}}{\pgfqpoint{5.286424in}{7.749125in}}%
\pgfpathcurveto{\pgfqpoint{5.276170in}{7.749125in}}{\pgfqpoint{5.266334in}{7.745051in}}{\pgfqpoint{5.259083in}{7.737800in}}%
\pgfpathcurveto{\pgfqpoint{5.251833in}{7.730549in}}{\pgfqpoint{5.247759in}{7.720714in}}{\pgfqpoint{5.247759in}{7.710460in}}%
\pgfpathcurveto{\pgfqpoint{5.247759in}{7.700206in}}{\pgfqpoint{5.251833in}{7.690370in}}{\pgfqpoint{5.259083in}{7.683119in}}%
\pgfpathcurveto{\pgfqpoint{5.266334in}{7.675869in}}{\pgfqpoint{5.276170in}{7.671795in}}{\pgfqpoint{5.286424in}{7.671795in}}%
\pgfpathlineto{\pgfqpoint{5.286424in}{7.671795in}}%
\pgfpathclose%
\pgfusepath{stroke}%
\end{pgfscope}%
\begin{pgfscope}%
\pgfpathrectangle{\pgfqpoint{1.109909in}{4.050000in}}{\pgfqpoint{5.400000in}{5.400000in}}%
\pgfusepath{clip}%
\pgfsetbuttcap%
\pgfsetroundjoin%
\pgfsetlinewidth{1.003750pt}%
\definecolor{currentstroke}{rgb}{0.835294,0.368627,0.000000}%
\pgfsetstrokecolor{currentstroke}%
\pgfsetstrokeopacity{0.900000}%
\pgfsetdash{}{0pt}%
\pgfpathmoveto{\pgfqpoint{4.979358in}{7.374667in}}%
\pgfpathcurveto{\pgfqpoint{4.989612in}{7.374667in}}{\pgfqpoint{4.999448in}{7.378741in}}{\pgfqpoint{5.006698in}{7.385992in}}%
\pgfpathcurveto{\pgfqpoint{5.013949in}{7.393243in}}{\pgfqpoint{5.018023in}{7.403078in}}{\pgfqpoint{5.018023in}{7.413332in}}%
\pgfpathcurveto{\pgfqpoint{5.018023in}{7.423586in}}{\pgfqpoint{5.013949in}{7.433422in}}{\pgfqpoint{5.006698in}{7.440673in}}%
\pgfpathcurveto{\pgfqpoint{4.999448in}{7.447923in}}{\pgfqpoint{4.989612in}{7.451997in}}{\pgfqpoint{4.979358in}{7.451997in}}%
\pgfpathcurveto{\pgfqpoint{4.969104in}{7.451997in}}{\pgfqpoint{4.959269in}{7.447923in}}{\pgfqpoint{4.952018in}{7.440673in}}%
\pgfpathcurveto{\pgfqpoint{4.944767in}{7.433422in}}{\pgfqpoint{4.940693in}{7.423586in}}{\pgfqpoint{4.940693in}{7.413332in}}%
\pgfpathcurveto{\pgfqpoint{4.940693in}{7.403078in}}{\pgfqpoint{4.944767in}{7.393243in}}{\pgfqpoint{4.952018in}{7.385992in}}%
\pgfpathcurveto{\pgfqpoint{4.959269in}{7.378741in}}{\pgfqpoint{4.969104in}{7.374667in}}{\pgfqpoint{4.979358in}{7.374667in}}%
\pgfpathlineto{\pgfqpoint{4.979358in}{7.374667in}}%
\pgfpathclose%
\pgfusepath{stroke}%
\end{pgfscope}%
\begin{pgfscope}%
\pgfpathrectangle{\pgfqpoint{1.109909in}{4.050000in}}{\pgfqpoint{5.400000in}{5.400000in}}%
\pgfusepath{clip}%
\pgfsetbuttcap%
\pgfsetroundjoin%
\pgfsetlinewidth{1.003750pt}%
\definecolor{currentstroke}{rgb}{0.549020,0.160784,0.505882}%
\pgfsetstrokecolor{currentstroke}%
\pgfsetstrokeopacity{0.900000}%
\pgfsetdash{}{0pt}%
\pgfpathmoveto{\pgfqpoint{5.420768in}{8.208836in}}%
\pgfpathcurveto{\pgfqpoint{5.431022in}{8.208836in}}{\pgfqpoint{5.440858in}{8.212910in}}{\pgfqpoint{5.448109in}{8.220160in}}%
\pgfpathcurveto{\pgfqpoint{5.455359in}{8.227411in}}{\pgfqpoint{5.459433in}{8.237247in}}{\pgfqpoint{5.459433in}{8.247501in}}%
\pgfpathcurveto{\pgfqpoint{5.459433in}{8.257755in}}{\pgfqpoint{5.455359in}{8.267590in}}{\pgfqpoint{5.448109in}{8.274841in}}%
\pgfpathcurveto{\pgfqpoint{5.440858in}{8.282092in}}{\pgfqpoint{5.431022in}{8.286166in}}{\pgfqpoint{5.420768in}{8.286166in}}%
\pgfpathcurveto{\pgfqpoint{5.410514in}{8.286166in}}{\pgfqpoint{5.400679in}{8.282092in}}{\pgfqpoint{5.393428in}{8.274841in}}%
\pgfpathcurveto{\pgfqpoint{5.386177in}{8.267590in}}{\pgfqpoint{5.382103in}{8.257755in}}{\pgfqpoint{5.382103in}{8.247501in}}%
\pgfpathcurveto{\pgfqpoint{5.382103in}{8.237247in}}{\pgfqpoint{5.386177in}{8.227411in}}{\pgfqpoint{5.393428in}{8.220160in}}%
\pgfpathcurveto{\pgfqpoint{5.400679in}{8.212910in}}{\pgfqpoint{5.410514in}{8.208836in}}{\pgfqpoint{5.420768in}{8.208836in}}%
\pgfpathlineto{\pgfqpoint{5.420768in}{8.208836in}}%
\pgfpathclose%
\pgfusepath{stroke}%
\end{pgfscope}%
\begin{pgfscope}%
\pgfpathrectangle{\pgfqpoint{1.109909in}{4.050000in}}{\pgfqpoint{5.400000in}{5.400000in}}%
\pgfusepath{clip}%
\pgfsetbuttcap%
\pgfsetroundjoin%
\pgfsetlinewidth{1.003750pt}%
\definecolor{currentstroke}{rgb}{0.149020,0.274510,0.325490}%
\pgfsetstrokecolor{currentstroke}%
\pgfsetstrokeopacity{0.900000}%
\pgfsetdash{}{0pt}%
\pgfpathmoveto{\pgfqpoint{4.567144in}{7.450560in}}%
\pgfpathcurveto{\pgfqpoint{4.577398in}{7.450560in}}{\pgfqpoint{4.587233in}{7.454634in}}{\pgfqpoint{4.594484in}{7.461885in}}%
\pgfpathcurveto{\pgfqpoint{4.601735in}{7.469135in}}{\pgfqpoint{4.605809in}{7.478971in}}{\pgfqpoint{4.605809in}{7.489225in}}%
\pgfpathcurveto{\pgfqpoint{4.605809in}{7.499479in}}{\pgfqpoint{4.601735in}{7.509314in}}{\pgfqpoint{4.594484in}{7.516565in}}%
\pgfpathcurveto{\pgfqpoint{4.587233in}{7.523816in}}{\pgfqpoint{4.577398in}{7.527890in}}{\pgfqpoint{4.567144in}{7.527890in}}%
\pgfpathcurveto{\pgfqpoint{4.556890in}{7.527890in}}{\pgfqpoint{4.547054in}{7.523816in}}{\pgfqpoint{4.539804in}{7.516565in}}%
\pgfpathcurveto{\pgfqpoint{4.532553in}{7.509314in}}{\pgfqpoint{4.528479in}{7.499479in}}{\pgfqpoint{4.528479in}{7.489225in}}%
\pgfpathcurveto{\pgfqpoint{4.528479in}{7.478971in}}{\pgfqpoint{4.532553in}{7.469135in}}{\pgfqpoint{4.539804in}{7.461885in}}%
\pgfpathcurveto{\pgfqpoint{4.547054in}{7.454634in}}{\pgfqpoint{4.556890in}{7.450560in}}{\pgfqpoint{4.567144in}{7.450560in}}%
\pgfpathlineto{\pgfqpoint{4.567144in}{7.450560in}}%
\pgfpathclose%
\pgfusepath{stroke}%
\end{pgfscope}%
\begin{pgfscope}%
\pgfpathrectangle{\pgfqpoint{1.109909in}{4.050000in}}{\pgfqpoint{5.400000in}{5.400000in}}%
\pgfusepath{clip}%
\pgfsetbuttcap%
\pgfsetroundjoin%
\pgfsetlinewidth{1.003750pt}%
\definecolor{currentstroke}{rgb}{0.149020,0.274510,0.325490}%
\pgfsetstrokecolor{currentstroke}%
\pgfsetstrokeopacity{0.900000}%
\pgfsetdash{}{0pt}%
\pgfpathmoveto{\pgfqpoint{4.808235in}{7.648475in}}%
\pgfpathcurveto{\pgfqpoint{4.818489in}{7.648475in}}{\pgfqpoint{4.828325in}{7.652549in}}{\pgfqpoint{4.835575in}{7.659800in}}%
\pgfpathcurveto{\pgfqpoint{4.842826in}{7.667050in}}{\pgfqpoint{4.846900in}{7.676886in}}{\pgfqpoint{4.846900in}{7.687140in}}%
\pgfpathcurveto{\pgfqpoint{4.846900in}{7.697394in}}{\pgfqpoint{4.842826in}{7.707230in}}{\pgfqpoint{4.835575in}{7.714480in}}%
\pgfpathcurveto{\pgfqpoint{4.828325in}{7.721731in}}{\pgfqpoint{4.818489in}{7.725805in}}{\pgfqpoint{4.808235in}{7.725805in}}%
\pgfpathcurveto{\pgfqpoint{4.797981in}{7.725805in}}{\pgfqpoint{4.788146in}{7.721731in}}{\pgfqpoint{4.780895in}{7.714480in}}%
\pgfpathcurveto{\pgfqpoint{4.773644in}{7.707230in}}{\pgfqpoint{4.769570in}{7.697394in}}{\pgfqpoint{4.769570in}{7.687140in}}%
\pgfpathcurveto{\pgfqpoint{4.769570in}{7.676886in}}{\pgfqpoint{4.773644in}{7.667050in}}{\pgfqpoint{4.780895in}{7.659800in}}%
\pgfpathcurveto{\pgfqpoint{4.788146in}{7.652549in}}{\pgfqpoint{4.797981in}{7.648475in}}{\pgfqpoint{4.808235in}{7.648475in}}%
\pgfpathlineto{\pgfqpoint{4.808235in}{7.648475in}}%
\pgfpathclose%
\pgfusepath{stroke}%
\end{pgfscope}%
\begin{pgfscope}%
\pgfpathrectangle{\pgfqpoint{1.109909in}{4.050000in}}{\pgfqpoint{5.400000in}{5.400000in}}%
\pgfusepath{clip}%
\pgfsetbuttcap%
\pgfsetroundjoin%
\pgfsetlinewidth{1.003750pt}%
\definecolor{currentstroke}{rgb}{0.149020,0.274510,0.325490}%
\pgfsetstrokecolor{currentstroke}%
\pgfsetstrokeopacity{0.900000}%
\pgfsetdash{}{0pt}%
\pgfpathmoveto{\pgfqpoint{5.005902in}{7.866990in}}%
\pgfpathcurveto{\pgfqpoint{5.016156in}{7.866990in}}{\pgfqpoint{5.025992in}{7.871064in}}{\pgfqpoint{5.033243in}{7.878315in}}%
\pgfpathcurveto{\pgfqpoint{5.040493in}{7.885566in}}{\pgfqpoint{5.044567in}{7.895401in}}{\pgfqpoint{5.044567in}{7.905655in}}%
\pgfpathcurveto{\pgfqpoint{5.044567in}{7.915909in}}{\pgfqpoint{5.040493in}{7.925745in}}{\pgfqpoint{5.033243in}{7.932996in}}%
\pgfpathcurveto{\pgfqpoint{5.025992in}{7.940246in}}{\pgfqpoint{5.016156in}{7.944320in}}{\pgfqpoint{5.005902in}{7.944320in}}%
\pgfpathcurveto{\pgfqpoint{4.995648in}{7.944320in}}{\pgfqpoint{4.985813in}{7.940246in}}{\pgfqpoint{4.978562in}{7.932996in}}%
\pgfpathcurveto{\pgfqpoint{4.971311in}{7.925745in}}{\pgfqpoint{4.967237in}{7.915909in}}{\pgfqpoint{4.967237in}{7.905655in}}%
\pgfpathcurveto{\pgfqpoint{4.967237in}{7.895401in}}{\pgfqpoint{4.971311in}{7.885566in}}{\pgfqpoint{4.978562in}{7.878315in}}%
\pgfpathcurveto{\pgfqpoint{4.985813in}{7.871064in}}{\pgfqpoint{4.995648in}{7.866990in}}{\pgfqpoint{5.005902in}{7.866990in}}%
\pgfpathlineto{\pgfqpoint{5.005902in}{7.866990in}}%
\pgfpathclose%
\pgfusepath{stroke}%
\end{pgfscope}%
\begin{pgfscope}%
\pgfpathrectangle{\pgfqpoint{1.109909in}{4.050000in}}{\pgfqpoint{5.400000in}{5.400000in}}%
\pgfusepath{clip}%
\pgfsetbuttcap%
\pgfsetroundjoin%
\pgfsetlinewidth{1.003750pt}%
\definecolor{currentstroke}{rgb}{0.000000,0.619608,0.450980}%
\pgfsetstrokecolor{currentstroke}%
\pgfsetstrokeopacity{0.900000}%
\pgfsetdash{}{0pt}%
\pgfpathmoveto{\pgfqpoint{5.353423in}{7.701937in}}%
\pgfpathcurveto{\pgfqpoint{5.363677in}{7.701937in}}{\pgfqpoint{5.373513in}{7.706011in}}{\pgfqpoint{5.380763in}{7.713262in}}%
\pgfpathcurveto{\pgfqpoint{5.388014in}{7.720512in}}{\pgfqpoint{5.392088in}{7.730348in}}{\pgfqpoint{5.392088in}{7.740602in}}%
\pgfpathcurveto{\pgfqpoint{5.392088in}{7.750856in}}{\pgfqpoint{5.388014in}{7.760691in}}{\pgfqpoint{5.380763in}{7.767942in}}%
\pgfpathcurveto{\pgfqpoint{5.373513in}{7.775193in}}{\pgfqpoint{5.363677in}{7.779267in}}{\pgfqpoint{5.353423in}{7.779267in}}%
\pgfpathcurveto{\pgfqpoint{5.343169in}{7.779267in}}{\pgfqpoint{5.333334in}{7.775193in}}{\pgfqpoint{5.326083in}{7.767942in}}%
\pgfpathcurveto{\pgfqpoint{5.318832in}{7.760691in}}{\pgfqpoint{5.314758in}{7.750856in}}{\pgfqpoint{5.314758in}{7.740602in}}%
\pgfpathcurveto{\pgfqpoint{5.314758in}{7.730348in}}{\pgfqpoint{5.318832in}{7.720512in}}{\pgfqpoint{5.326083in}{7.713262in}}%
\pgfpathcurveto{\pgfqpoint{5.333334in}{7.706011in}}{\pgfqpoint{5.343169in}{7.701937in}}{\pgfqpoint{5.353423in}{7.701937in}}%
\pgfpathlineto{\pgfqpoint{5.353423in}{7.701937in}}%
\pgfpathclose%
\pgfusepath{stroke}%
\end{pgfscope}%
\begin{pgfscope}%
\pgfpathrectangle{\pgfqpoint{1.109909in}{4.050000in}}{\pgfqpoint{5.400000in}{5.400000in}}%
\pgfusepath{clip}%
\pgfsetbuttcap%
\pgfsetroundjoin%
\pgfsetlinewidth{1.003750pt}%
\definecolor{currentstroke}{rgb}{0.000000,0.619608,0.450980}%
\pgfsetstrokecolor{currentstroke}%
\pgfsetstrokeopacity{0.900000}%
\pgfsetdash{}{0pt}%
\pgfpathmoveto{\pgfqpoint{5.621920in}{8.354679in}}%
\pgfpathcurveto{\pgfqpoint{5.632174in}{8.354679in}}{\pgfqpoint{5.642010in}{8.358753in}}{\pgfqpoint{5.649260in}{8.366004in}}%
\pgfpathcurveto{\pgfqpoint{5.656511in}{8.373255in}}{\pgfqpoint{5.660585in}{8.383090in}}{\pgfqpoint{5.660585in}{8.393344in}}%
\pgfpathcurveto{\pgfqpoint{5.660585in}{8.403598in}}{\pgfqpoint{5.656511in}{8.413434in}}{\pgfqpoint{5.649260in}{8.420685in}}%
\pgfpathcurveto{\pgfqpoint{5.642010in}{8.427935in}}{\pgfqpoint{5.632174in}{8.432009in}}{\pgfqpoint{5.621920in}{8.432009in}}%
\pgfpathcurveto{\pgfqpoint{5.611666in}{8.432009in}}{\pgfqpoint{5.601830in}{8.427935in}}{\pgfqpoint{5.594580in}{8.420685in}}%
\pgfpathcurveto{\pgfqpoint{5.587329in}{8.413434in}}{\pgfqpoint{5.583255in}{8.403598in}}{\pgfqpoint{5.583255in}{8.393344in}}%
\pgfpathcurveto{\pgfqpoint{5.583255in}{8.383090in}}{\pgfqpoint{5.587329in}{8.373255in}}{\pgfqpoint{5.594580in}{8.366004in}}%
\pgfpathcurveto{\pgfqpoint{5.601830in}{8.358753in}}{\pgfqpoint{5.611666in}{8.354679in}}{\pgfqpoint{5.621920in}{8.354679in}}%
\pgfpathlineto{\pgfqpoint{5.621920in}{8.354679in}}%
\pgfpathclose%
\pgfusepath{stroke}%
\end{pgfscope}%
\begin{pgfscope}%
\pgfpathrectangle{\pgfqpoint{1.109909in}{4.050000in}}{\pgfqpoint{5.400000in}{5.400000in}}%
\pgfusepath{clip}%
\pgfsetbuttcap%
\pgfsetroundjoin%
\pgfsetlinewidth{1.003750pt}%
\definecolor{currentstroke}{rgb}{0.549020,0.337255,0.294118}%
\pgfsetstrokecolor{currentstroke}%
\pgfsetstrokeopacity{0.900000}%
\pgfsetdash{}{0pt}%
\pgfpathmoveto{\pgfqpoint{5.605880in}{5.268369in}}%
\pgfpathcurveto{\pgfqpoint{5.616134in}{5.268369in}}{\pgfqpoint{5.625969in}{5.272443in}}{\pgfqpoint{5.633220in}{5.279693in}}%
\pgfpathcurveto{\pgfqpoint{5.640471in}{5.286944in}}{\pgfqpoint{5.644545in}{5.296779in}}{\pgfqpoint{5.644545in}{5.307034in}}%
\pgfpathcurveto{\pgfqpoint{5.644545in}{5.317288in}}{\pgfqpoint{5.640471in}{5.327123in}}{\pgfqpoint{5.633220in}{5.334374in}}%
\pgfpathcurveto{\pgfqpoint{5.625969in}{5.341625in}}{\pgfqpoint{5.616134in}{5.345699in}}{\pgfqpoint{5.605880in}{5.345699in}}%
\pgfpathcurveto{\pgfqpoint{5.595626in}{5.345699in}}{\pgfqpoint{5.585790in}{5.341625in}}{\pgfqpoint{5.578539in}{5.334374in}}%
\pgfpathcurveto{\pgfqpoint{5.571289in}{5.327123in}}{\pgfqpoint{5.567215in}{5.317288in}}{\pgfqpoint{5.567215in}{5.307034in}}%
\pgfpathcurveto{\pgfqpoint{5.567215in}{5.296779in}}{\pgfqpoint{5.571289in}{5.286944in}}{\pgfqpoint{5.578539in}{5.279693in}}%
\pgfpathcurveto{\pgfqpoint{5.585790in}{5.272443in}}{\pgfqpoint{5.595626in}{5.268369in}}{\pgfqpoint{5.605880in}{5.268369in}}%
\pgfpathlineto{\pgfqpoint{5.605880in}{5.268369in}}%
\pgfpathclose%
\pgfusepath{stroke}%
\end{pgfscope}%
\begin{pgfscope}%
\pgfsetbuttcap%
\pgfsetroundjoin%
\definecolor{currentfill}{rgb}{0.352941,0.407843,0.450980}%
\pgfsetfillcolor{currentfill}%
\pgfsetlinewidth{0.000000pt}%
\definecolor{currentstroke}{rgb}{0.352941,0.407843,0.450980}%
\pgfsetstrokecolor{currentstroke}%
\pgfsetdash{}{0pt}%
\pgfpathmoveto{\pgfqpoint{1.135226in}{9.595947in}}%
\pgfpathlineto{\pgfqpoint{1.180644in}{9.595947in}}%
\pgfpathlineto{\pgfqpoint{1.180644in}{9.585000in}}%
\pgfpathlineto{\pgfqpoint{1.119578in}{9.585000in}}%
\pgfpathlineto{\pgfqpoint{1.119578in}{9.595947in}}%
\pgfpathquadraticcurveto{\pgfqpoint{1.126979in}{9.603617in}}{\pgfqpoint{1.139762in}{9.616522in}}%
\pgfpathquadraticcurveto{\pgfqpoint{1.152564in}{9.629449in}}{\pgfqpoint{1.155842in}{9.633201in}}%
\pgfpathquadraticcurveto{\pgfqpoint{1.162089in}{9.640211in}}{\pgfqpoint{1.164563in}{9.645076in}}%
\pgfpathquadraticcurveto{\pgfqpoint{1.167058in}{9.649941in}}{\pgfqpoint{1.167058in}{9.654642in}}%
\pgfpathquadraticcurveto{\pgfqpoint{1.167058in}{9.662311in}}{\pgfqpoint{1.161677in}{9.667135in}}%
\pgfpathquadraticcurveto{\pgfqpoint{1.156296in}{9.671980in}}{\pgfqpoint{1.147658in}{9.671980in}}%
\pgfpathquadraticcurveto{\pgfqpoint{1.141535in}{9.671980in}}{\pgfqpoint{1.134731in}{9.669857in}}%
\pgfpathquadraticcurveto{\pgfqpoint{1.127948in}{9.667733in}}{\pgfqpoint{1.120217in}{9.663404in}}%
\pgfpathlineto{\pgfqpoint{1.120217in}{9.676557in}}%
\pgfpathquadraticcurveto{\pgfqpoint{1.128072in}{9.679711in}}{\pgfqpoint{1.134896in}{9.681319in}}%
\pgfpathquadraticcurveto{\pgfqpoint{1.141741in}{9.682928in}}{\pgfqpoint{1.147410in}{9.682928in}}%
\pgfpathquadraticcurveto{\pgfqpoint{1.162357in}{9.682928in}}{\pgfqpoint{1.171243in}{9.675444in}}%
\pgfpathquadraticcurveto{\pgfqpoint{1.180128in}{9.667981in}}{\pgfqpoint{1.180128in}{9.655487in}}%
\pgfpathquadraticcurveto{\pgfqpoint{1.180128in}{9.649550in}}{\pgfqpoint{1.177902in}{9.644231in}}%
\pgfpathquadraticcurveto{\pgfqpoint{1.175696in}{9.638932in}}{\pgfqpoint{1.169820in}{9.631717in}}%
\pgfpathquadraticcurveto{\pgfqpoint{1.168212in}{9.629840in}}{\pgfqpoint{1.159574in}{9.620914in}}%
\pgfpathquadraticcurveto{\pgfqpoint{1.150956in}{9.611987in}}{\pgfqpoint{1.135226in}{9.595947in}}%
\pgfpathlineto{\pgfqpoint{1.135226in}{9.595947in}}%
\pgfpathclose%
\pgfpathmoveto{\pgfqpoint{1.237421in}{9.638273in}}%
\pgfpathquadraticcurveto{\pgfqpoint{1.228659in}{9.638273in}}{\pgfqpoint{1.223526in}{9.632273in}}%
\pgfpathquadraticcurveto{\pgfqpoint{1.218413in}{9.626294in}}{\pgfqpoint{1.218413in}{9.615863in}}%
\pgfpathquadraticcurveto{\pgfqpoint{1.218413in}{9.605493in}}{\pgfqpoint{1.223526in}{9.599452in}}%
\pgfpathquadraticcurveto{\pgfqpoint{1.228659in}{9.593432in}}{\pgfqpoint{1.237421in}{9.593432in}}%
\pgfpathquadraticcurveto{\pgfqpoint{1.246183in}{9.593432in}}{\pgfqpoint{1.251296in}{9.599452in}}%
\pgfpathquadraticcurveto{\pgfqpoint{1.256409in}{9.605493in}}{\pgfqpoint{1.256409in}{9.615863in}}%
\pgfpathquadraticcurveto{\pgfqpoint{1.256409in}{9.626294in}}{\pgfqpoint{1.251296in}{9.632273in}}%
\pgfpathquadraticcurveto{\pgfqpoint{1.246183in}{9.638273in}}{\pgfqpoint{1.237421in}{9.638273in}}%
\pgfpathlineto{\pgfqpoint{1.237421in}{9.638273in}}%
\pgfpathclose%
\pgfpathmoveto{\pgfqpoint{1.263253in}{9.679072in}}%
\pgfpathlineto{\pgfqpoint{1.263253in}{9.667218in}}%
\pgfpathquadraticcurveto{\pgfqpoint{1.258347in}{9.669527in}}{\pgfqpoint{1.253357in}{9.670743in}}%
\pgfpathquadraticcurveto{\pgfqpoint{1.248368in}{9.671980in}}{\pgfqpoint{1.243462in}{9.671980in}}%
\pgfpathquadraticcurveto{\pgfqpoint{1.230576in}{9.671980in}}{\pgfqpoint{1.223773in}{9.663280in}}%
\pgfpathquadraticcurveto{\pgfqpoint{1.216990in}{9.654580in}}{\pgfqpoint{1.216021in}{9.636994in}}%
\pgfpathquadraticcurveto{\pgfqpoint{1.219815in}{9.642602in}}{\pgfqpoint{1.225546in}{9.645591in}}%
\pgfpathquadraticcurveto{\pgfqpoint{1.231298in}{9.648581in}}{\pgfqpoint{1.238184in}{9.648581in}}%
\pgfpathquadraticcurveto{\pgfqpoint{1.252677in}{9.648581in}}{\pgfqpoint{1.261089in}{9.639778in}}%
\pgfpathquadraticcurveto{\pgfqpoint{1.269500in}{9.630995in}}{\pgfqpoint{1.269500in}{9.615863in}}%
\pgfpathquadraticcurveto{\pgfqpoint{1.269500in}{9.601039in}}{\pgfqpoint{1.260738in}{9.592071in}}%
\pgfpathquadraticcurveto{\pgfqpoint{1.251976in}{9.583124in}}{\pgfqpoint{1.237421in}{9.583124in}}%
\pgfpathquadraticcurveto{\pgfqpoint{1.220722in}{9.583124in}}{\pgfqpoint{1.211898in}{9.595906in}}%
\pgfpathquadraticcurveto{\pgfqpoint{1.203074in}{9.608709in}}{\pgfqpoint{1.203074in}{9.632995in}}%
\pgfpathquadraticcurveto{\pgfqpoint{1.203074in}{9.655796in}}{\pgfqpoint{1.213898in}{9.669362in}}%
\pgfpathquadraticcurveto{\pgfqpoint{1.224721in}{9.682928in}}{\pgfqpoint{1.242946in}{9.682928in}}%
\pgfpathquadraticcurveto{\pgfqpoint{1.247853in}{9.682928in}}{\pgfqpoint{1.252842in}{9.681959in}}%
\pgfpathquadraticcurveto{\pgfqpoint{1.257831in}{9.680990in}}{\pgfqpoint{1.263253in}{9.679072in}}%
\pgfpathlineto{\pgfqpoint{1.263253in}{9.679072in}}%
\pgfpathclose%
\pgfpathmoveto{\pgfqpoint{1.331349in}{9.636871in}}%
\pgfpathquadraticcurveto{\pgfqpoint{1.340688in}{9.634871in}}{\pgfqpoint{1.345925in}{9.628542in}}%
\pgfpathquadraticcurveto{\pgfqpoint{1.351182in}{9.622233in}}{\pgfqpoint{1.351182in}{9.612956in}}%
\pgfpathquadraticcurveto{\pgfqpoint{1.351182in}{9.598730in}}{\pgfqpoint{1.341389in}{9.590917in}}%
\pgfpathquadraticcurveto{\pgfqpoint{1.331596in}{9.583124in}}{\pgfqpoint{1.313557in}{9.583124in}}%
\pgfpathquadraticcurveto{\pgfqpoint{1.307517in}{9.583124in}}{\pgfqpoint{1.301105in}{9.584320in}}%
\pgfpathquadraticcurveto{\pgfqpoint{1.294693in}{9.585515in}}{\pgfqpoint{1.287869in}{9.587907in}}%
\pgfpathlineto{\pgfqpoint{1.287869in}{9.600462in}}%
\pgfpathquadraticcurveto{\pgfqpoint{1.293271in}{9.597308in}}{\pgfqpoint{1.299703in}{9.595700in}}%
\pgfpathquadraticcurveto{\pgfqpoint{1.306156in}{9.594092in}}{\pgfqpoint{1.313186in}{9.594092in}}%
\pgfpathquadraticcurveto{\pgfqpoint{1.325411in}{9.594092in}}{\pgfqpoint{1.331823in}{9.598916in}}%
\pgfpathquadraticcurveto{\pgfqpoint{1.338235in}{9.603740in}}{\pgfqpoint{1.338235in}{9.612956in}}%
\pgfpathquadraticcurveto{\pgfqpoint{1.338235in}{9.621470in}}{\pgfqpoint{1.332277in}{9.626253in}}%
\pgfpathquadraticcurveto{\pgfqpoint{1.326319in}{9.631057in}}{\pgfqpoint{1.315701in}{9.631057in}}%
\pgfpathlineto{\pgfqpoint{1.304486in}{9.631057in}}%
\pgfpathlineto{\pgfqpoint{1.304486in}{9.641757in}}%
\pgfpathlineto{\pgfqpoint{1.316217in}{9.641757in}}%
\pgfpathquadraticcurveto{\pgfqpoint{1.325803in}{9.641757in}}{\pgfqpoint{1.330895in}{9.645591in}}%
\pgfpathquadraticcurveto{\pgfqpoint{1.335988in}{9.649426in}}{\pgfqpoint{1.335988in}{9.656642in}}%
\pgfpathquadraticcurveto{\pgfqpoint{1.335988in}{9.664043in}}{\pgfqpoint{1.330730in}{9.668001in}}%
\pgfpathquadraticcurveto{\pgfqpoint{1.325494in}{9.671980in}}{\pgfqpoint{1.315701in}{9.671980in}}%
\pgfpathquadraticcurveto{\pgfqpoint{1.310341in}{9.671980in}}{\pgfqpoint{1.304218in}{9.670805in}}%
\pgfpathquadraticcurveto{\pgfqpoint{1.298095in}{9.669651in}}{\pgfqpoint{1.290755in}{9.667218in}}%
\pgfpathlineto{\pgfqpoint{1.290755in}{9.678804in}}%
\pgfpathquadraticcurveto{\pgfqpoint{1.298177in}{9.680866in}}{\pgfqpoint{1.304651in}{9.681897in}}%
\pgfpathquadraticcurveto{\pgfqpoint{1.311124in}{9.682928in}}{\pgfqpoint{1.316856in}{9.682928in}}%
\pgfpathquadraticcurveto{\pgfqpoint{1.331679in}{9.682928in}}{\pgfqpoint{1.340296in}{9.676186in}}%
\pgfpathquadraticcurveto{\pgfqpoint{1.348935in}{9.669465in}}{\pgfqpoint{1.348935in}{9.658002in}}%
\pgfpathquadraticcurveto{\pgfqpoint{1.348935in}{9.650003in}}{\pgfqpoint{1.344358in}{9.644499in}}%
\pgfpathquadraticcurveto{\pgfqpoint{1.339781in}{9.638994in}}{\pgfqpoint{1.331349in}{9.636871in}}%
\pgfpathlineto{\pgfqpoint{1.331349in}{9.636871in}}%
\pgfpathclose%
\pgfpathmoveto{\pgfqpoint{1.457933in}{9.646086in}}%
\pgfpathquadraticcurveto{\pgfqpoint{1.455933in}{9.647241in}}{\pgfqpoint{1.453583in}{9.647777in}}%
\pgfpathquadraticcurveto{\pgfqpoint{1.451233in}{9.648333in}}{\pgfqpoint{1.448408in}{9.648333in}}%
\pgfpathquadraticcurveto{\pgfqpoint{1.438348in}{9.648333in}}{\pgfqpoint{1.432967in}{9.641798in}}%
\pgfpathquadraticcurveto{\pgfqpoint{1.427586in}{9.635263in}}{\pgfqpoint{1.427586in}{9.623016in}}%
\pgfpathlineto{\pgfqpoint{1.427586in}{9.585000in}}%
\pgfpathlineto{\pgfqpoint{1.415670in}{9.585000in}}%
\pgfpathlineto{\pgfqpoint{1.415670in}{9.657157in}}%
\pgfpathlineto{\pgfqpoint{1.427586in}{9.657157in}}%
\pgfpathlineto{\pgfqpoint{1.427586in}{9.645942in}}%
\pgfpathquadraticcurveto{\pgfqpoint{1.431338in}{9.652518in}}{\pgfqpoint{1.437317in}{9.655693in}}%
\pgfpathquadraticcurveto{\pgfqpoint{1.443316in}{9.658889in}}{\pgfqpoint{1.451893in}{9.658889in}}%
\pgfpathquadraticcurveto{\pgfqpoint{1.453109in}{9.658889in}}{\pgfqpoint{1.454593in}{9.658724in}}%
\pgfpathquadraticcurveto{\pgfqpoint{1.456078in}{9.658580in}}{\pgfqpoint{1.457871in}{9.658250in}}%
\pgfpathlineto{\pgfqpoint{1.457933in}{9.646086in}}%
\pgfpathlineto{\pgfqpoint{1.457933in}{9.646086in}}%
\pgfpathclose%
\pgfpathmoveto{\pgfqpoint{1.529183in}{9.624047in}}%
\pgfpathlineto{\pgfqpoint{1.529183in}{9.618254in}}%
\pgfpathlineto{\pgfqpoint{1.474674in}{9.618254in}}%
\pgfpathquadraticcurveto{\pgfqpoint{1.475457in}{9.606008in}}{\pgfqpoint{1.482054in}{9.599596in}}%
\pgfpathquadraticcurveto{\pgfqpoint{1.488651in}{9.593185in}}{\pgfqpoint{1.500444in}{9.593185in}}%
\pgfpathquadraticcurveto{\pgfqpoint{1.507268in}{9.593185in}}{\pgfqpoint{1.513680in}{9.594855in}}%
\pgfpathquadraticcurveto{\pgfqpoint{1.520091in}{9.596525in}}{\pgfqpoint{1.526421in}{9.599885in}}%
\pgfpathlineto{\pgfqpoint{1.526421in}{9.588670in}}%
\pgfpathquadraticcurveto{\pgfqpoint{1.520030in}{9.585969in}}{\pgfqpoint{1.513329in}{9.584546in}}%
\pgfpathquadraticcurveto{\pgfqpoint{1.506629in}{9.583124in}}{\pgfqpoint{1.499743in}{9.583124in}}%
\pgfpathquadraticcurveto{\pgfqpoint{1.482467in}{9.583124in}}{\pgfqpoint{1.472385in}{9.593164in}}%
\pgfpathquadraticcurveto{\pgfqpoint{1.462304in}{9.603225in}}{\pgfqpoint{1.462304in}{9.620378in}}%
\pgfpathquadraticcurveto{\pgfqpoint{1.462304in}{9.638087in}}{\pgfqpoint{1.471870in}{9.648478in}}%
\pgfpathquadraticcurveto{\pgfqpoint{1.481436in}{9.658889in}}{\pgfqpoint{1.497681in}{9.658889in}}%
\pgfpathquadraticcurveto{\pgfqpoint{1.512237in}{9.658889in}}{\pgfqpoint{1.520710in}{9.649508in}}%
\pgfpathquadraticcurveto{\pgfqpoint{1.529183in}{9.640149in}}{\pgfqpoint{1.529183in}{9.624047in}}%
\pgfpathlineto{\pgfqpoint{1.529183in}{9.624047in}}%
\pgfpathclose%
\pgfpathmoveto{\pgfqpoint{1.517329in}{9.627531in}}%
\pgfpathquadraticcurveto{\pgfqpoint{1.517205in}{9.637242in}}{\pgfqpoint{1.511886in}{9.643035in}}%
\pgfpathquadraticcurveto{\pgfqpoint{1.506567in}{9.648849in}}{\pgfqpoint{1.497805in}{9.648849in}}%
\pgfpathquadraticcurveto{\pgfqpoint{1.487889in}{9.648849in}}{\pgfqpoint{1.481931in}{9.643241in}}%
\pgfpathquadraticcurveto{\pgfqpoint{1.475972in}{9.637633in}}{\pgfqpoint{1.475065in}{9.627449in}}%
\pgfpathlineto{\pgfqpoint{1.517329in}{9.627531in}}%
\pgfpathlineto{\pgfqpoint{1.517329in}{9.627531in}}%
\pgfpathclose%
\pgfpathmoveto{\pgfqpoint{1.560376in}{9.677650in}}%
\pgfpathlineto{\pgfqpoint{1.560376in}{9.657157in}}%
\pgfpathlineto{\pgfqpoint{1.584785in}{9.657157in}}%
\pgfpathlineto{\pgfqpoint{1.584785in}{9.647942in}}%
\pgfpathlineto{\pgfqpoint{1.560376in}{9.647942in}}%
\pgfpathlineto{\pgfqpoint{1.560376in}{9.608771in}}%
\pgfpathquadraticcurveto{\pgfqpoint{1.560376in}{9.599947in}}{\pgfqpoint{1.562788in}{9.597432in}}%
\pgfpathquadraticcurveto{\pgfqpoint{1.565200in}{9.594916in}}{\pgfqpoint{1.572622in}{9.594916in}}%
\pgfpathlineto{\pgfqpoint{1.584785in}{9.594916in}}%
\pgfpathlineto{\pgfqpoint{1.584785in}{9.585000in}}%
\pgfpathlineto{\pgfqpoint{1.572622in}{9.585000in}}%
\pgfpathquadraticcurveto{\pgfqpoint{1.558891in}{9.585000in}}{\pgfqpoint{1.553675in}{9.590113in}}%
\pgfpathquadraticcurveto{\pgfqpoint{1.548459in}{9.595246in}}{\pgfqpoint{1.548459in}{9.608771in}}%
\pgfpathlineto{\pgfqpoint{1.548459in}{9.647942in}}%
\pgfpathlineto{\pgfqpoint{1.539759in}{9.647942in}}%
\pgfpathlineto{\pgfqpoint{1.539759in}{9.657157in}}%
\pgfpathlineto{\pgfqpoint{1.548459in}{9.657157in}}%
\pgfpathlineto{\pgfqpoint{1.548459in}{9.677650in}}%
\pgfpathlineto{\pgfqpoint{1.560376in}{9.677650in}}%
\pgfpathlineto{\pgfqpoint{1.560376in}{9.677650in}}%
\pgfpathclose%
\pgfpathmoveto{\pgfqpoint{1.633172in}{9.621264in}}%
\pgfpathquadraticcurveto{\pgfqpoint{1.618802in}{9.621264in}}{\pgfqpoint{1.613257in}{9.617986in}}%
\pgfpathquadraticcurveto{\pgfqpoint{1.607711in}{9.614708in}}{\pgfqpoint{1.607711in}{9.606771in}}%
\pgfpathquadraticcurveto{\pgfqpoint{1.607711in}{9.600462in}}{\pgfqpoint{1.611875in}{9.596751in}}%
\pgfpathquadraticcurveto{\pgfqpoint{1.616040in}{9.593061in}}{\pgfqpoint{1.623173in}{9.593061in}}%
\pgfpathquadraticcurveto{\pgfqpoint{1.633048in}{9.593061in}}{\pgfqpoint{1.639006in}{9.600050in}}%
\pgfpathquadraticcurveto{\pgfqpoint{1.644964in}{9.607039in}}{\pgfqpoint{1.644964in}{9.618625in}}%
\pgfpathlineto{\pgfqpoint{1.644964in}{9.621264in}}%
\pgfpathlineto{\pgfqpoint{1.633172in}{9.621264in}}%
\pgfpathlineto{\pgfqpoint{1.633172in}{9.621264in}}%
\pgfpathclose%
\pgfpathmoveto{\pgfqpoint{1.656819in}{9.626171in}}%
\pgfpathlineto{\pgfqpoint{1.656819in}{9.585000in}}%
\pgfpathlineto{\pgfqpoint{1.644964in}{9.585000in}}%
\pgfpathlineto{\pgfqpoint{1.644964in}{9.595947in}}%
\pgfpathquadraticcurveto{\pgfqpoint{1.640903in}{9.589391in}}{\pgfqpoint{1.634842in}{9.586258in}}%
\pgfpathquadraticcurveto{\pgfqpoint{1.628781in}{9.583124in}}{\pgfqpoint{1.620019in}{9.583124in}}%
\pgfpathquadraticcurveto{\pgfqpoint{1.608948in}{9.583124in}}{\pgfqpoint{1.602392in}{9.589350in}}%
\pgfpathquadraticcurveto{\pgfqpoint{1.595856in}{9.595576in}}{\pgfqpoint{1.595856in}{9.606008in}}%
\pgfpathquadraticcurveto{\pgfqpoint{1.595856in}{9.618172in}}{\pgfqpoint{1.604000in}{9.624357in}}%
\pgfpathquadraticcurveto{\pgfqpoint{1.612164in}{9.630541in}}{\pgfqpoint{1.628327in}{9.630541in}}%
\pgfpathlineto{\pgfqpoint{1.644964in}{9.630541in}}%
\pgfpathlineto{\pgfqpoint{1.644964in}{9.631717in}}%
\pgfpathquadraticcurveto{\pgfqpoint{1.644964in}{9.639901in}}{\pgfqpoint{1.639584in}{9.644375in}}%
\pgfpathquadraticcurveto{\pgfqpoint{1.634203in}{9.648849in}}{\pgfqpoint{1.624472in}{9.648849in}}%
\pgfpathquadraticcurveto{\pgfqpoint{1.618287in}{9.648849in}}{\pgfqpoint{1.612411in}{9.647364in}}%
\pgfpathquadraticcurveto{\pgfqpoint{1.606556in}{9.645880in}}{\pgfqpoint{1.601155in}{9.642911in}}%
\pgfpathlineto{\pgfqpoint{1.601155in}{9.653879in}}%
\pgfpathquadraticcurveto{\pgfqpoint{1.607649in}{9.656394in}}{\pgfqpoint{1.613772in}{9.657631in}}%
\pgfpathquadraticcurveto{\pgfqpoint{1.619895in}{9.658889in}}{\pgfqpoint{1.625688in}{9.658889in}}%
\pgfpathquadraticcurveto{\pgfqpoint{1.641357in}{9.658889in}}{\pgfqpoint{1.649088in}{9.650766in}}%
\pgfpathquadraticcurveto{\pgfqpoint{1.656819in}{9.642664in}}{\pgfqpoint{1.656819in}{9.626171in}}%
\pgfpathlineto{\pgfqpoint{1.656819in}{9.626171in}}%
\pgfpathclose%
\pgfpathmoveto{\pgfqpoint{1.681229in}{9.657157in}}%
\pgfpathlineto{\pgfqpoint{1.693083in}{9.657157in}}%
\pgfpathlineto{\pgfqpoint{1.693083in}{9.585000in}}%
\pgfpathlineto{\pgfqpoint{1.681229in}{9.585000in}}%
\pgfpathlineto{\pgfqpoint{1.681229in}{9.657157in}}%
\pgfpathlineto{\pgfqpoint{1.681229in}{9.657157in}}%
\pgfpathclose%
\pgfpathmoveto{\pgfqpoint{1.681229in}{9.685257in}}%
\pgfpathlineto{\pgfqpoint{1.693083in}{9.685257in}}%
\pgfpathlineto{\pgfqpoint{1.693083in}{9.670228in}}%
\pgfpathlineto{\pgfqpoint{1.681229in}{9.670228in}}%
\pgfpathlineto{\pgfqpoint{1.681229in}{9.685257in}}%
\pgfpathlineto{\pgfqpoint{1.681229in}{9.685257in}}%
\pgfpathclose%
\pgfpathmoveto{\pgfqpoint{1.777878in}{9.628562in}}%
\pgfpathlineto{\pgfqpoint{1.777878in}{9.585000in}}%
\pgfpathlineto{\pgfqpoint{1.766023in}{9.585000in}}%
\pgfpathlineto{\pgfqpoint{1.766023in}{9.628171in}}%
\pgfpathquadraticcurveto{\pgfqpoint{1.766023in}{9.638417in}}{\pgfqpoint{1.762024in}{9.643488in}}%
\pgfpathquadraticcurveto{\pgfqpoint{1.758024in}{9.648581in}}{\pgfqpoint{1.750046in}{9.648581in}}%
\pgfpathquadraticcurveto{\pgfqpoint{1.740439in}{9.648581in}}{\pgfqpoint{1.734893in}{9.642458in}}%
\pgfpathquadraticcurveto{\pgfqpoint{1.729347in}{9.636355in}}{\pgfqpoint{1.729347in}{9.625779in}}%
\pgfpathlineto{\pgfqpoint{1.729347in}{9.585000in}}%
\pgfpathlineto{\pgfqpoint{1.717431in}{9.585000in}}%
\pgfpathlineto{\pgfqpoint{1.717431in}{9.657157in}}%
\pgfpathlineto{\pgfqpoint{1.729347in}{9.657157in}}%
\pgfpathlineto{\pgfqpoint{1.729347in}{9.645942in}}%
\pgfpathquadraticcurveto{\pgfqpoint{1.733615in}{9.652457in}}{\pgfqpoint{1.739367in}{9.655673in}}%
\pgfpathquadraticcurveto{\pgfqpoint{1.745139in}{9.658889in}}{\pgfqpoint{1.752685in}{9.658889in}}%
\pgfpathquadraticcurveto{\pgfqpoint{1.765116in}{9.658889in}}{\pgfqpoint{1.771487in}{9.651199in}}%
\pgfpathquadraticcurveto{\pgfqpoint{1.777878in}{9.643509in}}{\pgfqpoint{1.777878in}{9.628562in}}%
\pgfpathlineto{\pgfqpoint{1.777878in}{9.628562in}}%
\pgfpathclose%
\pgfpathmoveto{\pgfqpoint{1.863229in}{9.624047in}}%
\pgfpathlineto{\pgfqpoint{1.863229in}{9.618254in}}%
\pgfpathlineto{\pgfqpoint{1.808720in}{9.618254in}}%
\pgfpathquadraticcurveto{\pgfqpoint{1.809503in}{9.606008in}}{\pgfqpoint{1.816100in}{9.599596in}}%
\pgfpathquadraticcurveto{\pgfqpoint{1.822698in}{9.593185in}}{\pgfqpoint{1.834490in}{9.593185in}}%
\pgfpathquadraticcurveto{\pgfqpoint{1.841314in}{9.593185in}}{\pgfqpoint{1.847726in}{9.594855in}}%
\pgfpathquadraticcurveto{\pgfqpoint{1.854138in}{9.596525in}}{\pgfqpoint{1.860467in}{9.599885in}}%
\pgfpathlineto{\pgfqpoint{1.860467in}{9.588670in}}%
\pgfpathquadraticcurveto{\pgfqpoint{1.854076in}{9.585969in}}{\pgfqpoint{1.847375in}{9.584546in}}%
\pgfpathquadraticcurveto{\pgfqpoint{1.840675in}{9.583124in}}{\pgfqpoint{1.833789in}{9.583124in}}%
\pgfpathquadraticcurveto{\pgfqpoint{1.816513in}{9.583124in}}{\pgfqpoint{1.806431in}{9.593164in}}%
\pgfpathquadraticcurveto{\pgfqpoint{1.796350in}{9.603225in}}{\pgfqpoint{1.796350in}{9.620378in}}%
\pgfpathquadraticcurveto{\pgfqpoint{1.796350in}{9.638087in}}{\pgfqpoint{1.805916in}{9.648478in}}%
\pgfpathquadraticcurveto{\pgfqpoint{1.815482in}{9.658889in}}{\pgfqpoint{1.831728in}{9.658889in}}%
\pgfpathquadraticcurveto{\pgfqpoint{1.846283in}{9.658889in}}{\pgfqpoint{1.854756in}{9.649508in}}%
\pgfpathquadraticcurveto{\pgfqpoint{1.863229in}{9.640149in}}{\pgfqpoint{1.863229in}{9.624047in}}%
\pgfpathlineto{\pgfqpoint{1.863229in}{9.624047in}}%
\pgfpathclose%
\pgfpathmoveto{\pgfqpoint{1.851375in}{9.627531in}}%
\pgfpathquadraticcurveto{\pgfqpoint{1.851251in}{9.637242in}}{\pgfqpoint{1.845932in}{9.643035in}}%
\pgfpathquadraticcurveto{\pgfqpoint{1.840613in}{9.648849in}}{\pgfqpoint{1.831851in}{9.648849in}}%
\pgfpathquadraticcurveto{\pgfqpoint{1.821935in}{9.648849in}}{\pgfqpoint{1.815977in}{9.643241in}}%
\pgfpathquadraticcurveto{\pgfqpoint{1.810019in}{9.637633in}}{\pgfqpoint{1.809112in}{9.627449in}}%
\pgfpathlineto{\pgfqpoint{1.851375in}{9.627531in}}%
\pgfpathlineto{\pgfqpoint{1.851375in}{9.627531in}}%
\pgfpathclose%
\pgfpathmoveto{\pgfqpoint{1.930171in}{9.646210in}}%
\pgfpathlineto{\pgfqpoint{1.930171in}{9.685257in}}%
\pgfpathlineto{\pgfqpoint{1.942025in}{9.685257in}}%
\pgfpathlineto{\pgfqpoint{1.942025in}{9.585000in}}%
\pgfpathlineto{\pgfqpoint{1.930171in}{9.585000in}}%
\pgfpathlineto{\pgfqpoint{1.930171in}{9.595824in}}%
\pgfpathquadraticcurveto{\pgfqpoint{1.926439in}{9.589391in}}{\pgfqpoint{1.920728in}{9.586258in}}%
\pgfpathquadraticcurveto{\pgfqpoint{1.915038in}{9.583124in}}{\pgfqpoint{1.907039in}{9.583124in}}%
\pgfpathquadraticcurveto{\pgfqpoint{1.893968in}{9.583124in}}{\pgfqpoint{1.885742in}{9.593556in}}%
\pgfpathquadraticcurveto{\pgfqpoint{1.877537in}{9.604008in}}{\pgfqpoint{1.877537in}{9.621017in}}%
\pgfpathquadraticcurveto{\pgfqpoint{1.877537in}{9.638025in}}{\pgfqpoint{1.885742in}{9.648457in}}%
\pgfpathquadraticcurveto{\pgfqpoint{1.893968in}{9.658889in}}{\pgfqpoint{1.907039in}{9.658889in}}%
\pgfpathquadraticcurveto{\pgfqpoint{1.915038in}{9.658889in}}{\pgfqpoint{1.920728in}{9.655755in}}%
\pgfpathquadraticcurveto{\pgfqpoint{1.926439in}{9.652642in}}{\pgfqpoint{1.930171in}{9.646210in}}%
\pgfpathlineto{\pgfqpoint{1.930171in}{9.646210in}}%
\pgfpathclose%
\pgfpathmoveto{\pgfqpoint{1.889783in}{9.621017in}}%
\pgfpathquadraticcurveto{\pgfqpoint{1.889783in}{9.607946in}}{\pgfqpoint{1.895164in}{9.600503in}}%
\pgfpathquadraticcurveto{\pgfqpoint{1.900545in}{9.593061in}}{\pgfqpoint{1.909946in}{9.593061in}}%
\pgfpathquadraticcurveto{\pgfqpoint{1.919347in}{9.593061in}}{\pgfqpoint{1.924749in}{9.600503in}}%
\pgfpathquadraticcurveto{\pgfqpoint{1.930171in}{9.607946in}}{\pgfqpoint{1.930171in}{9.621017in}}%
\pgfpathquadraticcurveto{\pgfqpoint{1.930171in}{9.634087in}}{\pgfqpoint{1.924749in}{9.641530in}}%
\pgfpathquadraticcurveto{\pgfqpoint{1.919347in}{9.648972in}}{\pgfqpoint{1.909946in}{9.648972in}}%
\pgfpathquadraticcurveto{\pgfqpoint{1.900545in}{9.648972in}}{\pgfqpoint{1.895164in}{9.641530in}}%
\pgfpathquadraticcurveto{\pgfqpoint{1.889783in}{9.634087in}}{\pgfqpoint{1.889783in}{9.621017in}}%
\pgfpathlineto{\pgfqpoint{1.889783in}{9.621017in}}%
\pgfpathclose%
\pgfpathmoveto{\pgfqpoint{2.029459in}{9.681196in}}%
\pgfpathlineto{\pgfqpoint{2.040406in}{9.681196in}}%
\pgfpathlineto{\pgfqpoint{2.006905in}{9.572754in}}%
\pgfpathlineto{\pgfqpoint{1.995957in}{9.572754in}}%
\pgfpathlineto{\pgfqpoint{2.029459in}{9.681196in}}%
\pgfpathlineto{\pgfqpoint{2.029459in}{9.681196in}}%
\pgfpathclose%
\pgfpathmoveto{\pgfqpoint{2.135880in}{9.636871in}}%
\pgfpathquadraticcurveto{\pgfqpoint{2.145219in}{9.634871in}}{\pgfqpoint{2.150456in}{9.628542in}}%
\pgfpathquadraticcurveto{\pgfqpoint{2.155713in}{9.622233in}}{\pgfqpoint{2.155713in}{9.612956in}}%
\pgfpathquadraticcurveto{\pgfqpoint{2.155713in}{9.598730in}}{\pgfqpoint{2.145920in}{9.590917in}}%
\pgfpathquadraticcurveto{\pgfqpoint{2.136128in}{9.583124in}}{\pgfqpoint{2.118088in}{9.583124in}}%
\pgfpathquadraticcurveto{\pgfqpoint{2.112048in}{9.583124in}}{\pgfqpoint{2.105636in}{9.584320in}}%
\pgfpathquadraticcurveto{\pgfqpoint{2.099224in}{9.585515in}}{\pgfqpoint{2.092400in}{9.587907in}}%
\pgfpathlineto{\pgfqpoint{2.092400in}{9.600462in}}%
\pgfpathquadraticcurveto{\pgfqpoint{2.097802in}{9.597308in}}{\pgfqpoint{2.104234in}{9.595700in}}%
\pgfpathquadraticcurveto{\pgfqpoint{2.110687in}{9.594092in}}{\pgfqpoint{2.117717in}{9.594092in}}%
\pgfpathquadraticcurveto{\pgfqpoint{2.129943in}{9.594092in}}{\pgfqpoint{2.136354in}{9.598916in}}%
\pgfpathquadraticcurveto{\pgfqpoint{2.142766in}{9.603740in}}{\pgfqpoint{2.142766in}{9.612956in}}%
\pgfpathquadraticcurveto{\pgfqpoint{2.142766in}{9.621470in}}{\pgfqpoint{2.136808in}{9.626253in}}%
\pgfpathquadraticcurveto{\pgfqpoint{2.130850in}{9.631057in}}{\pgfqpoint{2.120232in}{9.631057in}}%
\pgfpathlineto{\pgfqpoint{2.109017in}{9.631057in}}%
\pgfpathlineto{\pgfqpoint{2.109017in}{9.641757in}}%
\pgfpathlineto{\pgfqpoint{2.120748in}{9.641757in}}%
\pgfpathquadraticcurveto{\pgfqpoint{2.130334in}{9.641757in}}{\pgfqpoint{2.135427in}{9.645591in}}%
\pgfpathquadraticcurveto{\pgfqpoint{2.140519in}{9.649426in}}{\pgfqpoint{2.140519in}{9.656642in}}%
\pgfpathquadraticcurveto{\pgfqpoint{2.140519in}{9.664043in}}{\pgfqpoint{2.135262in}{9.668001in}}%
\pgfpathquadraticcurveto{\pgfqpoint{2.130025in}{9.671980in}}{\pgfqpoint{2.120232in}{9.671980in}}%
\pgfpathquadraticcurveto{\pgfqpoint{2.114872in}{9.671980in}}{\pgfqpoint{2.108749in}{9.670805in}}%
\pgfpathquadraticcurveto{\pgfqpoint{2.102626in}{9.669651in}}{\pgfqpoint{2.095287in}{9.667218in}}%
\pgfpathlineto{\pgfqpoint{2.095287in}{9.678804in}}%
\pgfpathquadraticcurveto{\pgfqpoint{2.102709in}{9.680866in}}{\pgfqpoint{2.109182in}{9.681897in}}%
\pgfpathquadraticcurveto{\pgfqpoint{2.115656in}{9.682928in}}{\pgfqpoint{2.121387in}{9.682928in}}%
\pgfpathquadraticcurveto{\pgfqpoint{2.136210in}{9.682928in}}{\pgfqpoint{2.144828in}{9.676186in}}%
\pgfpathquadraticcurveto{\pgfqpoint{2.153466in}{9.669465in}}{\pgfqpoint{2.153466in}{9.658002in}}%
\pgfpathquadraticcurveto{\pgfqpoint{2.153466in}{9.650003in}}{\pgfqpoint{2.148889in}{9.644499in}}%
\pgfpathquadraticcurveto{\pgfqpoint{2.144312in}{9.638994in}}{\pgfqpoint{2.135880in}{9.636871in}}%
\pgfpathlineto{\pgfqpoint{2.135880in}{9.636871in}}%
\pgfpathclose%
\pgfpathmoveto{\pgfqpoint{2.180535in}{9.681196in}}%
\pgfpathlineto{\pgfqpoint{2.231622in}{9.681196in}}%
\pgfpathlineto{\pgfqpoint{2.231622in}{9.670228in}}%
\pgfpathlineto{\pgfqpoint{2.192451in}{9.670228in}}%
\pgfpathlineto{\pgfqpoint{2.192451in}{9.646663in}}%
\pgfpathquadraticcurveto{\pgfqpoint{2.195276in}{9.647632in}}{\pgfqpoint{2.198100in}{9.648107in}}%
\pgfpathquadraticcurveto{\pgfqpoint{2.200945in}{9.648581in}}{\pgfqpoint{2.203790in}{9.648581in}}%
\pgfpathquadraticcurveto{\pgfqpoint{2.219892in}{9.648581in}}{\pgfqpoint{2.229293in}{9.639757in}}%
\pgfpathquadraticcurveto{\pgfqpoint{2.238714in}{9.630933in}}{\pgfqpoint{2.238714in}{9.615863in}}%
\pgfpathquadraticcurveto{\pgfqpoint{2.238714in}{9.600339in}}{\pgfqpoint{2.229045in}{9.591721in}}%
\pgfpathquadraticcurveto{\pgfqpoint{2.219376in}{9.583124in}}{\pgfqpoint{2.201791in}{9.583124in}}%
\pgfpathquadraticcurveto{\pgfqpoint{2.195729in}{9.583124in}}{\pgfqpoint{2.189441in}{9.584155in}}%
\pgfpathquadraticcurveto{\pgfqpoint{2.183174in}{9.585186in}}{\pgfqpoint{2.176474in}{9.587247in}}%
\pgfpathlineto{\pgfqpoint{2.176474in}{9.600339in}}%
\pgfpathquadraticcurveto{\pgfqpoint{2.182267in}{9.597184in}}{\pgfqpoint{2.188452in}{9.595638in}}%
\pgfpathquadraticcurveto{\pgfqpoint{2.194637in}{9.594092in}}{\pgfqpoint{2.201523in}{9.594092in}}%
\pgfpathquadraticcurveto{\pgfqpoint{2.212676in}{9.594092in}}{\pgfqpoint{2.219170in}{9.599947in}}%
\pgfpathquadraticcurveto{\pgfqpoint{2.225685in}{9.605802in}}{\pgfqpoint{2.225685in}{9.615863in}}%
\pgfpathquadraticcurveto{\pgfqpoint{2.225685in}{9.625903in}}{\pgfqpoint{2.219170in}{9.631758in}}%
\pgfpathquadraticcurveto{\pgfqpoint{2.212676in}{9.637633in}}{\pgfqpoint{2.201523in}{9.637633in}}%
\pgfpathquadraticcurveto{\pgfqpoint{2.196307in}{9.637633in}}{\pgfqpoint{2.191111in}{9.636479in}}%
\pgfpathquadraticcurveto{\pgfqpoint{2.185937in}{9.635324in}}{\pgfqpoint{2.180535in}{9.632871in}}%
\pgfpathlineto{\pgfqpoint{2.180535in}{9.681196in}}%
\pgfpathlineto{\pgfqpoint{2.180535in}{9.681196in}}%
\pgfpathclose%
\pgfpathmoveto{\pgfqpoint{2.303780in}{9.636871in}}%
\pgfpathquadraticcurveto{\pgfqpoint{2.313119in}{9.634871in}}{\pgfqpoint{2.318355in}{9.628542in}}%
\pgfpathquadraticcurveto{\pgfqpoint{2.323612in}{9.622233in}}{\pgfqpoint{2.323612in}{9.612956in}}%
\pgfpathquadraticcurveto{\pgfqpoint{2.323612in}{9.598730in}}{\pgfqpoint{2.313820in}{9.590917in}}%
\pgfpathquadraticcurveto{\pgfqpoint{2.304027in}{9.583124in}}{\pgfqpoint{2.285988in}{9.583124in}}%
\pgfpathquadraticcurveto{\pgfqpoint{2.279947in}{9.583124in}}{\pgfqpoint{2.273535in}{9.584320in}}%
\pgfpathquadraticcurveto{\pgfqpoint{2.267124in}{9.585515in}}{\pgfqpoint{2.260300in}{9.587907in}}%
\pgfpathlineto{\pgfqpoint{2.260300in}{9.600462in}}%
\pgfpathquadraticcurveto{\pgfqpoint{2.265701in}{9.597308in}}{\pgfqpoint{2.272133in}{9.595700in}}%
\pgfpathquadraticcurveto{\pgfqpoint{2.278586in}{9.594092in}}{\pgfqpoint{2.285617in}{9.594092in}}%
\pgfpathquadraticcurveto{\pgfqpoint{2.297842in}{9.594092in}}{\pgfqpoint{2.304254in}{9.598916in}}%
\pgfpathquadraticcurveto{\pgfqpoint{2.310665in}{9.603740in}}{\pgfqpoint{2.310665in}{9.612956in}}%
\pgfpathquadraticcurveto{\pgfqpoint{2.310665in}{9.621470in}}{\pgfqpoint{2.304707in}{9.626253in}}%
\pgfpathquadraticcurveto{\pgfqpoint{2.298749in}{9.631057in}}{\pgfqpoint{2.288132in}{9.631057in}}%
\pgfpathlineto{\pgfqpoint{2.276916in}{9.631057in}}%
\pgfpathlineto{\pgfqpoint{2.276916in}{9.641757in}}%
\pgfpathlineto{\pgfqpoint{2.288647in}{9.641757in}}%
\pgfpathquadraticcurveto{\pgfqpoint{2.298234in}{9.641757in}}{\pgfqpoint{2.303326in}{9.645591in}}%
\pgfpathquadraticcurveto{\pgfqpoint{2.308418in}{9.649426in}}{\pgfqpoint{2.308418in}{9.656642in}}%
\pgfpathquadraticcurveto{\pgfqpoint{2.308418in}{9.664043in}}{\pgfqpoint{2.303161in}{9.668001in}}%
\pgfpathquadraticcurveto{\pgfqpoint{2.297924in}{9.671980in}}{\pgfqpoint{2.288132in}{9.671980in}}%
\pgfpathquadraticcurveto{\pgfqpoint{2.282772in}{9.671980in}}{\pgfqpoint{2.276648in}{9.670805in}}%
\pgfpathquadraticcurveto{\pgfqpoint{2.270525in}{9.669651in}}{\pgfqpoint{2.263186in}{9.667218in}}%
\pgfpathlineto{\pgfqpoint{2.263186in}{9.678804in}}%
\pgfpathquadraticcurveto{\pgfqpoint{2.270608in}{9.680866in}}{\pgfqpoint{2.277081in}{9.681897in}}%
\pgfpathquadraticcurveto{\pgfqpoint{2.283555in}{9.682928in}}{\pgfqpoint{2.289286in}{9.682928in}}%
\pgfpathquadraticcurveto{\pgfqpoint{2.304109in}{9.682928in}}{\pgfqpoint{2.312727in}{9.676186in}}%
\pgfpathquadraticcurveto{\pgfqpoint{2.321365in}{9.669465in}}{\pgfqpoint{2.321365in}{9.658002in}}%
\pgfpathquadraticcurveto{\pgfqpoint{2.321365in}{9.650003in}}{\pgfqpoint{2.316788in}{9.644499in}}%
\pgfpathquadraticcurveto{\pgfqpoint{2.312212in}{9.638994in}}{\pgfqpoint{2.303780in}{9.636871in}}%
\pgfpathlineto{\pgfqpoint{2.303780in}{9.636871in}}%
\pgfpathclose%
\pgfpathmoveto{\pgfqpoint{2.440486in}{9.654395in}}%
\pgfpathlineto{\pgfqpoint{2.440486in}{9.643303in}}%
\pgfpathquadraticcurveto{\pgfqpoint{2.435456in}{9.646086in}}{\pgfqpoint{2.430405in}{9.647467in}}%
\pgfpathquadraticcurveto{\pgfqpoint{2.425354in}{9.648849in}}{\pgfqpoint{2.420200in}{9.648849in}}%
\pgfpathquadraticcurveto{\pgfqpoint{2.408655in}{9.648849in}}{\pgfqpoint{2.402264in}{9.641530in}}%
\pgfpathquadraticcurveto{\pgfqpoint{2.395893in}{9.634232in}}{\pgfqpoint{2.395893in}{9.621017in}}%
\pgfpathquadraticcurveto{\pgfqpoint{2.395893in}{9.607802in}}{\pgfqpoint{2.402264in}{9.600483in}}%
\pgfpathquadraticcurveto{\pgfqpoint{2.408655in}{9.593185in}}{\pgfqpoint{2.420200in}{9.593185in}}%
\pgfpathquadraticcurveto{\pgfqpoint{2.425354in}{9.593185in}}{\pgfqpoint{2.430405in}{9.594566in}}%
\pgfpathquadraticcurveto{\pgfqpoint{2.435456in}{9.595947in}}{\pgfqpoint{2.440486in}{9.598730in}}%
\pgfpathlineto{\pgfqpoint{2.440486in}{9.587763in}}%
\pgfpathquadraticcurveto{\pgfqpoint{2.435518in}{9.585454in}}{\pgfqpoint{2.430199in}{9.584299in}}%
\pgfpathquadraticcurveto{\pgfqpoint{2.424900in}{9.583124in}}{\pgfqpoint{2.418901in}{9.583124in}}%
\pgfpathquadraticcurveto{\pgfqpoint{2.402594in}{9.583124in}}{\pgfqpoint{2.392986in}{9.593370in}}%
\pgfpathquadraticcurveto{\pgfqpoint{2.383400in}{9.603617in}}{\pgfqpoint{2.383400in}{9.621017in}}%
\pgfpathquadraticcurveto{\pgfqpoint{2.383400in}{9.638664in}}{\pgfqpoint{2.393089in}{9.648766in}}%
\pgfpathquadraticcurveto{\pgfqpoint{2.402800in}{9.658889in}}{\pgfqpoint{2.419684in}{9.658889in}}%
\pgfpathquadraticcurveto{\pgfqpoint{2.425148in}{9.658889in}}{\pgfqpoint{2.430364in}{9.657755in}}%
\pgfpathquadraticcurveto{\pgfqpoint{2.435580in}{9.656642in}}{\pgfqpoint{2.440486in}{9.654395in}}%
\pgfpathlineto{\pgfqpoint{2.440486in}{9.654395in}}%
\pgfpathclose%
\pgfpathmoveto{\pgfqpoint{2.489058in}{9.648849in}}%
\pgfpathquadraticcurveto{\pgfqpoint{2.479534in}{9.648849in}}{\pgfqpoint{2.473988in}{9.641406in}}%
\pgfpathquadraticcurveto{\pgfqpoint{2.468442in}{9.633964in}}{\pgfqpoint{2.468442in}{9.621017in}}%
\pgfpathquadraticcurveto{\pgfqpoint{2.468442in}{9.608070in}}{\pgfqpoint{2.473947in}{9.600627in}}%
\pgfpathquadraticcurveto{\pgfqpoint{2.479472in}{9.593185in}}{\pgfqpoint{2.489058in}{9.593185in}}%
\pgfpathquadraticcurveto{\pgfqpoint{2.498542in}{9.593185in}}{\pgfqpoint{2.504067in}{9.600648in}}%
\pgfpathquadraticcurveto{\pgfqpoint{2.509613in}{9.608132in}}{\pgfqpoint{2.509613in}{9.621017in}}%
\pgfpathquadraticcurveto{\pgfqpoint{2.509613in}{9.633840in}}{\pgfqpoint{2.504067in}{9.641344in}}%
\pgfpathquadraticcurveto{\pgfqpoint{2.498542in}{9.648849in}}{\pgfqpoint{2.489058in}{9.648849in}}%
\pgfpathlineto{\pgfqpoint{2.489058in}{9.648849in}}%
\pgfpathclose%
\pgfpathmoveto{\pgfqpoint{2.489058in}{9.658889in}}%
\pgfpathquadraticcurveto{\pgfqpoint{2.504521in}{9.658889in}}{\pgfqpoint{2.513344in}{9.648828in}}%
\pgfpathquadraticcurveto{\pgfqpoint{2.522189in}{9.638788in}}{\pgfqpoint{2.522189in}{9.621017in}}%
\pgfpathquadraticcurveto{\pgfqpoint{2.522189in}{9.603307in}}{\pgfqpoint{2.513344in}{9.593205in}}%
\pgfpathquadraticcurveto{\pgfqpoint{2.504521in}{9.583124in}}{\pgfqpoint{2.489058in}{9.583124in}}%
\pgfpathquadraticcurveto{\pgfqpoint{2.473534in}{9.583124in}}{\pgfqpoint{2.464731in}{9.593205in}}%
\pgfpathquadraticcurveto{\pgfqpoint{2.455949in}{9.603307in}}{\pgfqpoint{2.455949in}{9.621017in}}%
\pgfpathquadraticcurveto{\pgfqpoint{2.455949in}{9.638788in}}{\pgfqpoint{2.464731in}{9.648828in}}%
\pgfpathquadraticcurveto{\pgfqpoint{2.473534in}{9.658889in}}{\pgfqpoint{2.489058in}{9.658889in}}%
\pgfpathlineto{\pgfqpoint{2.489058in}{9.658889in}}%
\pgfpathclose%
\pgfpathmoveto{\pgfqpoint{2.598016in}{9.643303in}}%
\pgfpathquadraticcurveto{\pgfqpoint{2.602469in}{9.651302in}}{\pgfqpoint{2.608654in}{9.655095in}}%
\pgfpathquadraticcurveto{\pgfqpoint{2.614839in}{9.658889in}}{\pgfqpoint{2.623209in}{9.658889in}}%
\pgfpathquadraticcurveto{\pgfqpoint{2.634486in}{9.658889in}}{\pgfqpoint{2.640609in}{9.650993in}}%
\pgfpathquadraticcurveto{\pgfqpoint{2.646732in}{9.643117in}}{\pgfqpoint{2.646732in}{9.628562in}}%
\pgfpathlineto{\pgfqpoint{2.646732in}{9.585000in}}%
\pgfpathlineto{\pgfqpoint{2.634816in}{9.585000in}}%
\pgfpathlineto{\pgfqpoint{2.634816in}{9.628171in}}%
\pgfpathquadraticcurveto{\pgfqpoint{2.634816in}{9.638541in}}{\pgfqpoint{2.631125in}{9.643550in}}%
\pgfpathquadraticcurveto{\pgfqpoint{2.627456in}{9.648581in}}{\pgfqpoint{2.619931in}{9.648581in}}%
\pgfpathquadraticcurveto{\pgfqpoint{2.610715in}{9.648581in}}{\pgfqpoint{2.605355in}{9.642458in}}%
\pgfpathquadraticcurveto{\pgfqpoint{2.600015in}{9.636355in}}{\pgfqpoint{2.600015in}{9.625779in}}%
\pgfpathlineto{\pgfqpoint{2.600015in}{9.585000in}}%
\pgfpathlineto{\pgfqpoint{2.588099in}{9.585000in}}%
\pgfpathlineto{\pgfqpoint{2.588099in}{9.628171in}}%
\pgfpathquadraticcurveto{\pgfqpoint{2.588099in}{9.638602in}}{\pgfqpoint{2.584429in}{9.643592in}}%
\pgfpathquadraticcurveto{\pgfqpoint{2.580760in}{9.648581in}}{\pgfqpoint{2.573091in}{9.648581in}}%
\pgfpathquadraticcurveto{\pgfqpoint{2.563999in}{9.648581in}}{\pgfqpoint{2.558638in}{9.642437in}}%
\pgfpathquadraticcurveto{\pgfqpoint{2.553299in}{9.636293in}}{\pgfqpoint{2.553299in}{9.625779in}}%
\pgfpathlineto{\pgfqpoint{2.553299in}{9.585000in}}%
\pgfpathlineto{\pgfqpoint{2.541383in}{9.585000in}}%
\pgfpathlineto{\pgfqpoint{2.541383in}{9.657157in}}%
\pgfpathlineto{\pgfqpoint{2.553299in}{9.657157in}}%
\pgfpathlineto{\pgfqpoint{2.553299in}{9.645942in}}%
\pgfpathquadraticcurveto{\pgfqpoint{2.557360in}{9.652580in}}{\pgfqpoint{2.563030in}{9.655735in}}%
\pgfpathquadraticcurveto{\pgfqpoint{2.568699in}{9.658889in}}{\pgfqpoint{2.576492in}{9.658889in}}%
\pgfpathquadraticcurveto{\pgfqpoint{2.584368in}{9.658889in}}{\pgfqpoint{2.589872in}{9.654889in}}%
\pgfpathquadraticcurveto{\pgfqpoint{2.595377in}{9.650910in}}{\pgfqpoint{2.598016in}{9.643303in}}%
\pgfpathlineto{\pgfqpoint{2.598016in}{9.643303in}}%
\pgfpathclose%
\pgfpathmoveto{\pgfqpoint{2.681821in}{9.595824in}}%
\pgfpathlineto{\pgfqpoint{2.681821in}{9.557560in}}%
\pgfpathlineto{\pgfqpoint{2.669905in}{9.557560in}}%
\pgfpathlineto{\pgfqpoint{2.669905in}{9.657157in}}%
\pgfpathlineto{\pgfqpoint{2.681821in}{9.657157in}}%
\pgfpathlineto{\pgfqpoint{2.681821in}{9.646210in}}%
\pgfpathquadraticcurveto{\pgfqpoint{2.685573in}{9.652642in}}{\pgfqpoint{2.691263in}{9.655755in}}%
\pgfpathquadraticcurveto{\pgfqpoint{2.696974in}{9.658889in}}{\pgfqpoint{2.704891in}{9.658889in}}%
\pgfpathquadraticcurveto{\pgfqpoint{2.718044in}{9.658889in}}{\pgfqpoint{2.726249in}{9.648457in}}%
\pgfpathquadraticcurveto{\pgfqpoint{2.734475in}{9.638025in}}{\pgfqpoint{2.734475in}{9.621017in}}%
\pgfpathquadraticcurveto{\pgfqpoint{2.734475in}{9.604008in}}{\pgfqpoint{2.726249in}{9.593556in}}%
\pgfpathquadraticcurveto{\pgfqpoint{2.718044in}{9.583124in}}{\pgfqpoint{2.704891in}{9.583124in}}%
\pgfpathquadraticcurveto{\pgfqpoint{2.696974in}{9.583124in}}{\pgfqpoint{2.691263in}{9.586258in}}%
\pgfpathquadraticcurveto{\pgfqpoint{2.685573in}{9.589391in}}{\pgfqpoint{2.681821in}{9.595824in}}%
\pgfpathlineto{\pgfqpoint{2.681821in}{9.595824in}}%
\pgfpathclose%
\pgfpathmoveto{\pgfqpoint{2.722167in}{9.621017in}}%
\pgfpathquadraticcurveto{\pgfqpoint{2.722167in}{9.634087in}}{\pgfqpoint{2.716786in}{9.641530in}}%
\pgfpathquadraticcurveto{\pgfqpoint{2.711405in}{9.648972in}}{\pgfqpoint{2.702004in}{9.648972in}}%
\pgfpathquadraticcurveto{\pgfqpoint{2.692583in}{9.648972in}}{\pgfqpoint{2.687202in}{9.641530in}}%
\pgfpathquadraticcurveto{\pgfqpoint{2.681821in}{9.634087in}}{\pgfqpoint{2.681821in}{9.621017in}}%
\pgfpathquadraticcurveto{\pgfqpoint{2.681821in}{9.607946in}}{\pgfqpoint{2.687202in}{9.600503in}}%
\pgfpathquadraticcurveto{\pgfqpoint{2.692583in}{9.593061in}}{\pgfqpoint{2.702004in}{9.593061in}}%
\pgfpathquadraticcurveto{\pgfqpoint{2.711405in}{9.593061in}}{\pgfqpoint{2.716786in}{9.600503in}}%
\pgfpathquadraticcurveto{\pgfqpoint{2.722167in}{9.607946in}}{\pgfqpoint{2.722167in}{9.621017in}}%
\pgfpathlineto{\pgfqpoint{2.722167in}{9.621017in}}%
\pgfpathclose%
\pgfpathmoveto{\pgfqpoint{2.786923in}{9.621264in}}%
\pgfpathquadraticcurveto{\pgfqpoint{2.772553in}{9.621264in}}{\pgfqpoint{2.767008in}{9.617986in}}%
\pgfpathquadraticcurveto{\pgfqpoint{2.761462in}{9.614708in}}{\pgfqpoint{2.761462in}{9.606771in}}%
\pgfpathquadraticcurveto{\pgfqpoint{2.761462in}{9.600462in}}{\pgfqpoint{2.765626in}{9.596751in}}%
\pgfpathquadraticcurveto{\pgfqpoint{2.769791in}{9.593061in}}{\pgfqpoint{2.776924in}{9.593061in}}%
\pgfpathquadraticcurveto{\pgfqpoint{2.786799in}{9.593061in}}{\pgfqpoint{2.792757in}{9.600050in}}%
\pgfpathquadraticcurveto{\pgfqpoint{2.798716in}{9.607039in}}{\pgfqpoint{2.798716in}{9.618625in}}%
\pgfpathlineto{\pgfqpoint{2.798716in}{9.621264in}}%
\pgfpathlineto{\pgfqpoint{2.786923in}{9.621264in}}%
\pgfpathlineto{\pgfqpoint{2.786923in}{9.621264in}}%
\pgfpathclose%
\pgfpathmoveto{\pgfqpoint{2.810570in}{9.626171in}}%
\pgfpathlineto{\pgfqpoint{2.810570in}{9.585000in}}%
\pgfpathlineto{\pgfqpoint{2.798716in}{9.585000in}}%
\pgfpathlineto{\pgfqpoint{2.798716in}{9.595947in}}%
\pgfpathquadraticcurveto{\pgfqpoint{2.794654in}{9.589391in}}{\pgfqpoint{2.788593in}{9.586258in}}%
\pgfpathquadraticcurveto{\pgfqpoint{2.782532in}{9.583124in}}{\pgfqpoint{2.773770in}{9.583124in}}%
\pgfpathquadraticcurveto{\pgfqpoint{2.762699in}{9.583124in}}{\pgfqpoint{2.756143in}{9.589350in}}%
\pgfpathquadraticcurveto{\pgfqpoint{2.749607in}{9.595576in}}{\pgfqpoint{2.749607in}{9.606008in}}%
\pgfpathquadraticcurveto{\pgfqpoint{2.749607in}{9.618172in}}{\pgfqpoint{2.757751in}{9.624357in}}%
\pgfpathquadraticcurveto{\pgfqpoint{2.765915in}{9.630541in}}{\pgfqpoint{2.782078in}{9.630541in}}%
\pgfpathlineto{\pgfqpoint{2.798716in}{9.630541in}}%
\pgfpathlineto{\pgfqpoint{2.798716in}{9.631717in}}%
\pgfpathquadraticcurveto{\pgfqpoint{2.798716in}{9.639901in}}{\pgfqpoint{2.793335in}{9.644375in}}%
\pgfpathquadraticcurveto{\pgfqpoint{2.787954in}{9.648849in}}{\pgfqpoint{2.778223in}{9.648849in}}%
\pgfpathquadraticcurveto{\pgfqpoint{2.772038in}{9.648849in}}{\pgfqpoint{2.766162in}{9.647364in}}%
\pgfpathquadraticcurveto{\pgfqpoint{2.760307in}{9.645880in}}{\pgfqpoint{2.754906in}{9.642911in}}%
\pgfpathlineto{\pgfqpoint{2.754906in}{9.653879in}}%
\pgfpathquadraticcurveto{\pgfqpoint{2.761400in}{9.656394in}}{\pgfqpoint{2.767523in}{9.657631in}}%
\pgfpathquadraticcurveto{\pgfqpoint{2.773646in}{9.658889in}}{\pgfqpoint{2.779439in}{9.658889in}}%
\pgfpathquadraticcurveto{\pgfqpoint{2.795108in}{9.658889in}}{\pgfqpoint{2.802839in}{9.650766in}}%
\pgfpathquadraticcurveto{\pgfqpoint{2.810570in}{9.642664in}}{\pgfqpoint{2.810570in}{9.626171in}}%
\pgfpathlineto{\pgfqpoint{2.810570in}{9.626171in}}%
\pgfpathclose%
\pgfpathmoveto{\pgfqpoint{2.876790in}{9.646086in}}%
\pgfpathquadraticcurveto{\pgfqpoint{2.874790in}{9.647241in}}{\pgfqpoint{2.872439in}{9.647777in}}%
\pgfpathquadraticcurveto{\pgfqpoint{2.870089in}{9.648333in}}{\pgfqpoint{2.867265in}{9.648333in}}%
\pgfpathquadraticcurveto{\pgfqpoint{2.857204in}{9.648333in}}{\pgfqpoint{2.851823in}{9.641798in}}%
\pgfpathquadraticcurveto{\pgfqpoint{2.846442in}{9.635263in}}{\pgfqpoint{2.846442in}{9.623016in}}%
\pgfpathlineto{\pgfqpoint{2.846442in}{9.585000in}}%
\pgfpathlineto{\pgfqpoint{2.834526in}{9.585000in}}%
\pgfpathlineto{\pgfqpoint{2.834526in}{9.657157in}}%
\pgfpathlineto{\pgfqpoint{2.846442in}{9.657157in}}%
\pgfpathlineto{\pgfqpoint{2.846442in}{9.645942in}}%
\pgfpathquadraticcurveto{\pgfqpoint{2.850194in}{9.652518in}}{\pgfqpoint{2.856173in}{9.655693in}}%
\pgfpathquadraticcurveto{\pgfqpoint{2.862173in}{9.658889in}}{\pgfqpoint{2.870749in}{9.658889in}}%
\pgfpathquadraticcurveto{\pgfqpoint{2.871965in}{9.658889in}}{\pgfqpoint{2.873450in}{9.658724in}}%
\pgfpathquadraticcurveto{\pgfqpoint{2.874934in}{9.658580in}}{\pgfqpoint{2.876728in}{9.658250in}}%
\pgfpathlineto{\pgfqpoint{2.876790in}{9.646086in}}%
\pgfpathlineto{\pgfqpoint{2.876790in}{9.646086in}}%
\pgfpathclose%
\pgfpathmoveto{\pgfqpoint{2.889221in}{9.657157in}}%
\pgfpathlineto{\pgfqpoint{2.901076in}{9.657157in}}%
\pgfpathlineto{\pgfqpoint{2.901076in}{9.585000in}}%
\pgfpathlineto{\pgfqpoint{2.889221in}{9.585000in}}%
\pgfpathlineto{\pgfqpoint{2.889221in}{9.657157in}}%
\pgfpathlineto{\pgfqpoint{2.889221in}{9.657157in}}%
\pgfpathclose%
\pgfpathmoveto{\pgfqpoint{2.889221in}{9.685257in}}%
\pgfpathlineto{\pgfqpoint{2.901076in}{9.685257in}}%
\pgfpathlineto{\pgfqpoint{2.901076in}{9.670228in}}%
\pgfpathlineto{\pgfqpoint{2.889221in}{9.670228in}}%
\pgfpathlineto{\pgfqpoint{2.889221in}{9.685257in}}%
\pgfpathlineto{\pgfqpoint{2.889221in}{9.685257in}}%
\pgfpathclose%
\pgfpathmoveto{\pgfqpoint{2.971872in}{9.655034in}}%
\pgfpathlineto{\pgfqpoint{2.971872in}{9.643818in}}%
\pgfpathquadraticcurveto{\pgfqpoint{2.966862in}{9.646395in}}{\pgfqpoint{2.961440in}{9.647674in}}%
\pgfpathquadraticcurveto{\pgfqpoint{2.956039in}{9.648972in}}{\pgfqpoint{2.950225in}{9.648972in}}%
\pgfpathquadraticcurveto{\pgfqpoint{2.941401in}{9.648972in}}{\pgfqpoint{2.936989in}{9.646272in}}%
\pgfpathquadraticcurveto{\pgfqpoint{2.932577in}{9.643571in}}{\pgfqpoint{2.932577in}{9.638149in}}%
\pgfpathquadraticcurveto{\pgfqpoint{2.932577in}{9.634026in}}{\pgfqpoint{2.935732in}{9.631675in}}%
\pgfpathquadraticcurveto{\pgfqpoint{2.938886in}{9.629325in}}{\pgfqpoint{2.948431in}{9.627202in}}%
\pgfpathlineto{\pgfqpoint{2.952493in}{9.626294in}}%
\pgfpathquadraticcurveto{\pgfqpoint{2.965110in}{9.623594in}}{\pgfqpoint{2.970429in}{9.618666in}}%
\pgfpathquadraticcurveto{\pgfqpoint{2.975748in}{9.613739in}}{\pgfqpoint{2.975748in}{9.604915in}}%
\pgfpathquadraticcurveto{\pgfqpoint{2.975748in}{9.594855in}}{\pgfqpoint{2.967790in}{9.588979in}}%
\pgfpathquadraticcurveto{\pgfqpoint{2.959832in}{9.583124in}}{\pgfqpoint{2.945916in}{9.583124in}}%
\pgfpathquadraticcurveto{\pgfqpoint{2.940123in}{9.583124in}}{\pgfqpoint{2.933835in}{9.584258in}}%
\pgfpathquadraticcurveto{\pgfqpoint{2.927547in}{9.585392in}}{\pgfqpoint{2.920599in}{9.587639in}}%
\pgfpathlineto{\pgfqpoint{2.920599in}{9.599885in}}%
\pgfpathquadraticcurveto{\pgfqpoint{2.927176in}{9.596463in}}{\pgfqpoint{2.933546in}{9.594752in}}%
\pgfpathquadraticcurveto{\pgfqpoint{2.939917in}{9.593061in}}{\pgfqpoint{2.946184in}{9.593061in}}%
\pgfpathquadraticcurveto{\pgfqpoint{2.954554in}{9.593061in}}{\pgfqpoint{2.959049in}{9.595927in}}%
\pgfpathquadraticcurveto{\pgfqpoint{2.963564in}{9.598792in}}{\pgfqpoint{2.963564in}{9.604008in}}%
\pgfpathquadraticcurveto{\pgfqpoint{2.963564in}{9.608832in}}{\pgfqpoint{2.960306in}{9.611410in}}%
\pgfpathquadraticcurveto{\pgfqpoint{2.957069in}{9.613987in}}{\pgfqpoint{2.946040in}{9.616378in}}%
\pgfpathlineto{\pgfqpoint{2.941916in}{9.617347in}}%
\pgfpathquadraticcurveto{\pgfqpoint{2.930907in}{9.619656in}}{\pgfqpoint{2.926001in}{9.624460in}}%
\pgfpathquadraticcurveto{\pgfqpoint{2.921115in}{9.629263in}}{\pgfqpoint{2.921115in}{9.637633in}}%
\pgfpathquadraticcurveto{\pgfqpoint{2.921115in}{9.647818in}}{\pgfqpoint{2.928330in}{9.653343in}}%
\pgfpathquadraticcurveto{\pgfqpoint{2.935546in}{9.658889in}}{\pgfqpoint{2.948823in}{9.658889in}}%
\pgfpathquadraticcurveto{\pgfqpoint{2.955379in}{9.658889in}}{\pgfqpoint{2.961172in}{9.657920in}}%
\pgfpathquadraticcurveto{\pgfqpoint{2.966986in}{9.656972in}}{\pgfqpoint{2.971872in}{9.655034in}}%
\pgfpathlineto{\pgfqpoint{2.971872in}{9.655034in}}%
\pgfpathclose%
\pgfpathmoveto{\pgfqpoint{3.022568in}{9.648849in}}%
\pgfpathquadraticcurveto{\pgfqpoint{3.013043in}{9.648849in}}{\pgfqpoint{3.007497in}{9.641406in}}%
\pgfpathquadraticcurveto{\pgfqpoint{3.001951in}{9.633964in}}{\pgfqpoint{3.001951in}{9.621017in}}%
\pgfpathquadraticcurveto{\pgfqpoint{3.001951in}{9.608070in}}{\pgfqpoint{3.007456in}{9.600627in}}%
\pgfpathquadraticcurveto{\pgfqpoint{3.012981in}{9.593185in}}{\pgfqpoint{3.022568in}{9.593185in}}%
\pgfpathquadraticcurveto{\pgfqpoint{3.032051in}{9.593185in}}{\pgfqpoint{3.037576in}{9.600648in}}%
\pgfpathquadraticcurveto{\pgfqpoint{3.043122in}{9.608132in}}{\pgfqpoint{3.043122in}{9.621017in}}%
\pgfpathquadraticcurveto{\pgfqpoint{3.043122in}{9.633840in}}{\pgfqpoint{3.037576in}{9.641344in}}%
\pgfpathquadraticcurveto{\pgfqpoint{3.032051in}{9.648849in}}{\pgfqpoint{3.022568in}{9.648849in}}%
\pgfpathlineto{\pgfqpoint{3.022568in}{9.648849in}}%
\pgfpathclose%
\pgfpathmoveto{\pgfqpoint{3.022568in}{9.658889in}}%
\pgfpathquadraticcurveto{\pgfqpoint{3.038030in}{9.658889in}}{\pgfqpoint{3.046854in}{9.648828in}}%
\pgfpathquadraticcurveto{\pgfqpoint{3.055698in}{9.638788in}}{\pgfqpoint{3.055698in}{9.621017in}}%
\pgfpathquadraticcurveto{\pgfqpoint{3.055698in}{9.603307in}}{\pgfqpoint{3.046854in}{9.593205in}}%
\pgfpathquadraticcurveto{\pgfqpoint{3.038030in}{9.583124in}}{\pgfqpoint{3.022568in}{9.583124in}}%
\pgfpathquadraticcurveto{\pgfqpoint{3.007043in}{9.583124in}}{\pgfqpoint{2.998240in}{9.593205in}}%
\pgfpathquadraticcurveto{\pgfqpoint{2.989458in}{9.603307in}}{\pgfqpoint{2.989458in}{9.621017in}}%
\pgfpathquadraticcurveto{\pgfqpoint{2.989458in}{9.638788in}}{\pgfqpoint{2.998240in}{9.648828in}}%
\pgfpathquadraticcurveto{\pgfqpoint{3.007043in}{9.658889in}}{\pgfqpoint{3.022568in}{9.658889in}}%
\pgfpathlineto{\pgfqpoint{3.022568in}{9.658889in}}%
\pgfpathclose%
\pgfpathmoveto{\pgfqpoint{3.135339in}{9.628562in}}%
\pgfpathlineto{\pgfqpoint{3.135339in}{9.585000in}}%
\pgfpathlineto{\pgfqpoint{3.123484in}{9.585000in}}%
\pgfpathlineto{\pgfqpoint{3.123484in}{9.628171in}}%
\pgfpathquadraticcurveto{\pgfqpoint{3.123484in}{9.638417in}}{\pgfqpoint{3.119485in}{9.643488in}}%
\pgfpathquadraticcurveto{\pgfqpoint{3.115485in}{9.648581in}}{\pgfqpoint{3.107507in}{9.648581in}}%
\pgfpathquadraticcurveto{\pgfqpoint{3.097900in}{9.648581in}}{\pgfqpoint{3.092354in}{9.642458in}}%
\pgfpathquadraticcurveto{\pgfqpoint{3.086808in}{9.636355in}}{\pgfqpoint{3.086808in}{9.625779in}}%
\pgfpathlineto{\pgfqpoint{3.086808in}{9.585000in}}%
\pgfpathlineto{\pgfqpoint{3.074892in}{9.585000in}}%
\pgfpathlineto{\pgfqpoint{3.074892in}{9.657157in}}%
\pgfpathlineto{\pgfqpoint{3.086808in}{9.657157in}}%
\pgfpathlineto{\pgfqpoint{3.086808in}{9.645942in}}%
\pgfpathquadraticcurveto{\pgfqpoint{3.091076in}{9.652457in}}{\pgfqpoint{3.096827in}{9.655673in}}%
\pgfpathquadraticcurveto{\pgfqpoint{3.102600in}{9.658889in}}{\pgfqpoint{3.110146in}{9.658889in}}%
\pgfpathquadraticcurveto{\pgfqpoint{3.122577in}{9.658889in}}{\pgfqpoint{3.128948in}{9.651199in}}%
\pgfpathquadraticcurveto{\pgfqpoint{3.135339in}{9.643509in}}{\pgfqpoint{3.135339in}{9.628562in}}%
\pgfpathlineto{\pgfqpoint{3.135339in}{9.628562in}}%
\pgfpathclose%
\pgfpathmoveto{\pgfqpoint{3.204960in}{9.655034in}}%
\pgfpathlineto{\pgfqpoint{3.204960in}{9.643818in}}%
\pgfpathquadraticcurveto{\pgfqpoint{3.199950in}{9.646395in}}{\pgfqpoint{3.194528in}{9.647674in}}%
\pgfpathquadraticcurveto{\pgfqpoint{3.189127in}{9.648972in}}{\pgfqpoint{3.183313in}{9.648972in}}%
\pgfpathquadraticcurveto{\pgfqpoint{3.174489in}{9.648972in}}{\pgfqpoint{3.170077in}{9.646272in}}%
\pgfpathquadraticcurveto{\pgfqpoint{3.165665in}{9.643571in}}{\pgfqpoint{3.165665in}{9.638149in}}%
\pgfpathquadraticcurveto{\pgfqpoint{3.165665in}{9.634026in}}{\pgfqpoint{3.168820in}{9.631675in}}%
\pgfpathquadraticcurveto{\pgfqpoint{3.171974in}{9.629325in}}{\pgfqpoint{3.181519in}{9.627202in}}%
\pgfpathlineto{\pgfqpoint{3.185581in}{9.626294in}}%
\pgfpathquadraticcurveto{\pgfqpoint{3.198198in}{9.623594in}}{\pgfqpoint{3.203517in}{9.618666in}}%
\pgfpathquadraticcurveto{\pgfqpoint{3.208836in}{9.613739in}}{\pgfqpoint{3.208836in}{9.604915in}}%
\pgfpathquadraticcurveto{\pgfqpoint{3.208836in}{9.594855in}}{\pgfqpoint{3.200878in}{9.588979in}}%
\pgfpathquadraticcurveto{\pgfqpoint{3.192920in}{9.583124in}}{\pgfqpoint{3.179004in}{9.583124in}}%
\pgfpathquadraticcurveto{\pgfqpoint{3.173211in}{9.583124in}}{\pgfqpoint{3.166923in}{9.584258in}}%
\pgfpathquadraticcurveto{\pgfqpoint{3.160635in}{9.585392in}}{\pgfqpoint{3.153687in}{9.587639in}}%
\pgfpathlineto{\pgfqpoint{3.153687in}{9.599885in}}%
\pgfpathquadraticcurveto{\pgfqpoint{3.160264in}{9.596463in}}{\pgfqpoint{3.166634in}{9.594752in}}%
\pgfpathquadraticcurveto{\pgfqpoint{3.173005in}{9.593061in}}{\pgfqpoint{3.179272in}{9.593061in}}%
\pgfpathquadraticcurveto{\pgfqpoint{3.187642in}{9.593061in}}{\pgfqpoint{3.192137in}{9.595927in}}%
\pgfpathquadraticcurveto{\pgfqpoint{3.196652in}{9.598792in}}{\pgfqpoint{3.196652in}{9.604008in}}%
\pgfpathquadraticcurveto{\pgfqpoint{3.196652in}{9.608832in}}{\pgfqpoint{3.193394in}{9.611410in}}%
\pgfpathquadraticcurveto{\pgfqpoint{3.190158in}{9.613987in}}{\pgfqpoint{3.179128in}{9.616378in}}%
\pgfpathlineto{\pgfqpoint{3.175005in}{9.617347in}}%
\pgfpathquadraticcurveto{\pgfqpoint{3.163995in}{9.619656in}}{\pgfqpoint{3.159089in}{9.624460in}}%
\pgfpathquadraticcurveto{\pgfqpoint{3.154203in}{9.629263in}}{\pgfqpoint{3.154203in}{9.637633in}}%
\pgfpathquadraticcurveto{\pgfqpoint{3.154203in}{9.647818in}}{\pgfqpoint{3.161418in}{9.653343in}}%
\pgfpathquadraticcurveto{\pgfqpoint{3.168634in}{9.658889in}}{\pgfqpoint{3.181911in}{9.658889in}}%
\pgfpathquadraticcurveto{\pgfqpoint{3.188467in}{9.658889in}}{\pgfqpoint{3.194260in}{9.657920in}}%
\pgfpathquadraticcurveto{\pgfqpoint{3.200074in}{9.656972in}}{\pgfqpoint{3.204960in}{9.655034in}}%
\pgfpathlineto{\pgfqpoint{3.204960in}{9.655034in}}%
\pgfpathclose%
\pgfpathmoveto{\pgfqpoint{3.230730in}{9.653219in}}%
\pgfpathlineto{\pgfqpoint{3.244317in}{9.653219in}}%
\pgfpathlineto{\pgfqpoint{3.244317in}{9.636871in}}%
\pgfpathlineto{\pgfqpoint{3.230730in}{9.636871in}}%
\pgfpathlineto{\pgfqpoint{3.230730in}{9.653219in}}%
\pgfpathlineto{\pgfqpoint{3.230730in}{9.653219in}}%
\pgfpathclose%
\pgfpathmoveto{\pgfqpoint{3.230730in}{9.601369in}}%
\pgfpathlineto{\pgfqpoint{3.244317in}{9.601369in}}%
\pgfpathlineto{\pgfqpoint{3.244317in}{9.590278in}}%
\pgfpathlineto{\pgfqpoint{3.233761in}{9.569661in}}%
\pgfpathlineto{\pgfqpoint{3.225453in}{9.569661in}}%
\pgfpathlineto{\pgfqpoint{3.230730in}{9.590278in}}%
\pgfpathlineto{\pgfqpoint{3.230730in}{9.601369in}}%
\pgfpathlineto{\pgfqpoint{3.230730in}{9.601369in}}%
\pgfpathclose%
\pgfpathmoveto{\pgfqpoint{3.355191in}{9.636871in}}%
\pgfpathquadraticcurveto{\pgfqpoint{3.364530in}{9.634871in}}{\pgfqpoint{3.369767in}{9.628542in}}%
\pgfpathquadraticcurveto{\pgfqpoint{3.375024in}{9.622233in}}{\pgfqpoint{3.375024in}{9.612956in}}%
\pgfpathquadraticcurveto{\pgfqpoint{3.375024in}{9.598730in}}{\pgfqpoint{3.365231in}{9.590917in}}%
\pgfpathquadraticcurveto{\pgfqpoint{3.355439in}{9.583124in}}{\pgfqpoint{3.337399in}{9.583124in}}%
\pgfpathquadraticcurveto{\pgfqpoint{3.331359in}{9.583124in}}{\pgfqpoint{3.324947in}{9.584320in}}%
\pgfpathquadraticcurveto{\pgfqpoint{3.318535in}{9.585515in}}{\pgfqpoint{3.311711in}{9.587907in}}%
\pgfpathlineto{\pgfqpoint{3.311711in}{9.600462in}}%
\pgfpathquadraticcurveto{\pgfqpoint{3.317113in}{9.597308in}}{\pgfqpoint{3.323545in}{9.595700in}}%
\pgfpathquadraticcurveto{\pgfqpoint{3.329998in}{9.594092in}}{\pgfqpoint{3.337028in}{9.594092in}}%
\pgfpathquadraticcurveto{\pgfqpoint{3.349254in}{9.594092in}}{\pgfqpoint{3.355665in}{9.598916in}}%
\pgfpathquadraticcurveto{\pgfqpoint{3.362077in}{9.603740in}}{\pgfqpoint{3.362077in}{9.612956in}}%
\pgfpathquadraticcurveto{\pgfqpoint{3.362077in}{9.621470in}}{\pgfqpoint{3.356119in}{9.626253in}}%
\pgfpathquadraticcurveto{\pgfqpoint{3.350161in}{9.631057in}}{\pgfqpoint{3.339543in}{9.631057in}}%
\pgfpathlineto{\pgfqpoint{3.328328in}{9.631057in}}%
\pgfpathlineto{\pgfqpoint{3.328328in}{9.641757in}}%
\pgfpathlineto{\pgfqpoint{3.340059in}{9.641757in}}%
\pgfpathquadraticcurveto{\pgfqpoint{3.349645in}{9.641757in}}{\pgfqpoint{3.354738in}{9.645591in}}%
\pgfpathquadraticcurveto{\pgfqpoint{3.359830in}{9.649426in}}{\pgfqpoint{3.359830in}{9.656642in}}%
\pgfpathquadraticcurveto{\pgfqpoint{3.359830in}{9.664043in}}{\pgfqpoint{3.354573in}{9.668001in}}%
\pgfpathquadraticcurveto{\pgfqpoint{3.349336in}{9.671980in}}{\pgfqpoint{3.339543in}{9.671980in}}%
\pgfpathquadraticcurveto{\pgfqpoint{3.334183in}{9.671980in}}{\pgfqpoint{3.328060in}{9.670805in}}%
\pgfpathquadraticcurveto{\pgfqpoint{3.321937in}{9.669651in}}{\pgfqpoint{3.314598in}{9.667218in}}%
\pgfpathlineto{\pgfqpoint{3.314598in}{9.678804in}}%
\pgfpathquadraticcurveto{\pgfqpoint{3.322020in}{9.680866in}}{\pgfqpoint{3.328493in}{9.681897in}}%
\pgfpathquadraticcurveto{\pgfqpoint{3.334967in}{9.682928in}}{\pgfqpoint{3.340698in}{9.682928in}}%
\pgfpathquadraticcurveto{\pgfqpoint{3.355521in}{9.682928in}}{\pgfqpoint{3.364139in}{9.676186in}}%
\pgfpathquadraticcurveto{\pgfqpoint{3.372777in}{9.669465in}}{\pgfqpoint{3.372777in}{9.658002in}}%
\pgfpathquadraticcurveto{\pgfqpoint{3.372777in}{9.650003in}}{\pgfqpoint{3.368200in}{9.644499in}}%
\pgfpathquadraticcurveto{\pgfqpoint{3.363623in}{9.638994in}}{\pgfqpoint{3.355191in}{9.636871in}}%
\pgfpathlineto{\pgfqpoint{3.355191in}{9.636871in}}%
\pgfpathclose%
\pgfpathmoveto{\pgfqpoint{3.410917in}{9.595947in}}%
\pgfpathlineto{\pgfqpoint{3.456335in}{9.595947in}}%
\pgfpathlineto{\pgfqpoint{3.456335in}{9.585000in}}%
\pgfpathlineto{\pgfqpoint{3.395269in}{9.585000in}}%
\pgfpathlineto{\pgfqpoint{3.395269in}{9.595947in}}%
\pgfpathquadraticcurveto{\pgfqpoint{3.402671in}{9.603617in}}{\pgfqpoint{3.415453in}{9.616522in}}%
\pgfpathquadraticcurveto{\pgfqpoint{3.428255in}{9.629449in}}{\pgfqpoint{3.431533in}{9.633201in}}%
\pgfpathquadraticcurveto{\pgfqpoint{3.437780in}{9.640211in}}{\pgfqpoint{3.440254in}{9.645076in}}%
\pgfpathquadraticcurveto{\pgfqpoint{3.442749in}{9.649941in}}{\pgfqpoint{3.442749in}{9.654642in}}%
\pgfpathquadraticcurveto{\pgfqpoint{3.442749in}{9.662311in}}{\pgfqpoint{3.437368in}{9.667135in}}%
\pgfpathquadraticcurveto{\pgfqpoint{3.431987in}{9.671980in}}{\pgfqpoint{3.423349in}{9.671980in}}%
\pgfpathquadraticcurveto{\pgfqpoint{3.417226in}{9.671980in}}{\pgfqpoint{3.410422in}{9.669857in}}%
\pgfpathquadraticcurveto{\pgfqpoint{3.403640in}{9.667733in}}{\pgfqpoint{3.395908in}{9.663404in}}%
\pgfpathlineto{\pgfqpoint{3.395908in}{9.676557in}}%
\pgfpathquadraticcurveto{\pgfqpoint{3.403763in}{9.679711in}}{\pgfqpoint{3.410587in}{9.681319in}}%
\pgfpathquadraticcurveto{\pgfqpoint{3.417432in}{9.682928in}}{\pgfqpoint{3.423101in}{9.682928in}}%
\pgfpathquadraticcurveto{\pgfqpoint{3.438048in}{9.682928in}}{\pgfqpoint{3.446934in}{9.675444in}}%
\pgfpathquadraticcurveto{\pgfqpoint{3.455819in}{9.667981in}}{\pgfqpoint{3.455819in}{9.655487in}}%
\pgfpathquadraticcurveto{\pgfqpoint{3.455819in}{9.649550in}}{\pgfqpoint{3.453593in}{9.644231in}}%
\pgfpathquadraticcurveto{\pgfqpoint{3.451387in}{9.638932in}}{\pgfqpoint{3.445511in}{9.631717in}}%
\pgfpathquadraticcurveto{\pgfqpoint{3.443903in}{9.629840in}}{\pgfqpoint{3.435265in}{9.620914in}}%
\pgfpathquadraticcurveto{\pgfqpoint{3.426647in}{9.611987in}}{\pgfqpoint{3.410917in}{9.595947in}}%
\pgfpathlineto{\pgfqpoint{3.410917in}{9.595947in}}%
\pgfpathclose%
\pgfpathmoveto{\pgfqpoint{3.575848in}{9.654395in}}%
\pgfpathlineto{\pgfqpoint{3.575848in}{9.643303in}}%
\pgfpathquadraticcurveto{\pgfqpoint{3.570817in}{9.646086in}}{\pgfqpoint{3.565766in}{9.647467in}}%
\pgfpathquadraticcurveto{\pgfqpoint{3.560715in}{9.648849in}}{\pgfqpoint{3.555561in}{9.648849in}}%
\pgfpathquadraticcurveto{\pgfqpoint{3.544016in}{9.648849in}}{\pgfqpoint{3.537625in}{9.641530in}}%
\pgfpathquadraticcurveto{\pgfqpoint{3.531255in}{9.634232in}}{\pgfqpoint{3.531255in}{9.621017in}}%
\pgfpathquadraticcurveto{\pgfqpoint{3.531255in}{9.607802in}}{\pgfqpoint{3.537625in}{9.600483in}}%
\pgfpathquadraticcurveto{\pgfqpoint{3.544016in}{9.593185in}}{\pgfqpoint{3.555561in}{9.593185in}}%
\pgfpathquadraticcurveto{\pgfqpoint{3.560715in}{9.593185in}}{\pgfqpoint{3.565766in}{9.594566in}}%
\pgfpathquadraticcurveto{\pgfqpoint{3.570817in}{9.595947in}}{\pgfqpoint{3.575848in}{9.598730in}}%
\pgfpathlineto{\pgfqpoint{3.575848in}{9.587763in}}%
\pgfpathquadraticcurveto{\pgfqpoint{3.570879in}{9.585454in}}{\pgfqpoint{3.565560in}{9.584299in}}%
\pgfpathquadraticcurveto{\pgfqpoint{3.560262in}{9.583124in}}{\pgfqpoint{3.554262in}{9.583124in}}%
\pgfpathquadraticcurveto{\pgfqpoint{3.537955in}{9.583124in}}{\pgfqpoint{3.528348in}{9.593370in}}%
\pgfpathquadraticcurveto{\pgfqpoint{3.518761in}{9.603617in}}{\pgfqpoint{3.518761in}{9.621017in}}%
\pgfpathquadraticcurveto{\pgfqpoint{3.518761in}{9.638664in}}{\pgfqpoint{3.528451in}{9.648766in}}%
\pgfpathquadraticcurveto{\pgfqpoint{3.538161in}{9.658889in}}{\pgfqpoint{3.555046in}{9.658889in}}%
\pgfpathquadraticcurveto{\pgfqpoint{3.560509in}{9.658889in}}{\pgfqpoint{3.565725in}{9.657755in}}%
\pgfpathquadraticcurveto{\pgfqpoint{3.570941in}{9.656642in}}{\pgfqpoint{3.575848in}{9.654395in}}%
\pgfpathlineto{\pgfqpoint{3.575848in}{9.654395in}}%
\pgfpathclose%
\pgfpathmoveto{\pgfqpoint{3.629265in}{9.621264in}}%
\pgfpathquadraticcurveto{\pgfqpoint{3.614895in}{9.621264in}}{\pgfqpoint{3.609349in}{9.617986in}}%
\pgfpathquadraticcurveto{\pgfqpoint{3.603803in}{9.614708in}}{\pgfqpoint{3.603803in}{9.606771in}}%
\pgfpathquadraticcurveto{\pgfqpoint{3.603803in}{9.600462in}}{\pgfqpoint{3.607968in}{9.596751in}}%
\pgfpathquadraticcurveto{\pgfqpoint{3.612132in}{9.593061in}}{\pgfqpoint{3.619266in}{9.593061in}}%
\pgfpathquadraticcurveto{\pgfqpoint{3.629141in}{9.593061in}}{\pgfqpoint{3.635099in}{9.600050in}}%
\pgfpathquadraticcurveto{\pgfqpoint{3.641057in}{9.607039in}}{\pgfqpoint{3.641057in}{9.618625in}}%
\pgfpathlineto{\pgfqpoint{3.641057in}{9.621264in}}%
\pgfpathlineto{\pgfqpoint{3.629265in}{9.621264in}}%
\pgfpathlineto{\pgfqpoint{3.629265in}{9.621264in}}%
\pgfpathclose%
\pgfpathmoveto{\pgfqpoint{3.652911in}{9.626171in}}%
\pgfpathlineto{\pgfqpoint{3.652911in}{9.585000in}}%
\pgfpathlineto{\pgfqpoint{3.641057in}{9.585000in}}%
\pgfpathlineto{\pgfqpoint{3.641057in}{9.595947in}}%
\pgfpathquadraticcurveto{\pgfqpoint{3.636996in}{9.589391in}}{\pgfqpoint{3.630934in}{9.586258in}}%
\pgfpathquadraticcurveto{\pgfqpoint{3.624873in}{9.583124in}}{\pgfqpoint{3.616111in}{9.583124in}}%
\pgfpathquadraticcurveto{\pgfqpoint{3.605040in}{9.583124in}}{\pgfqpoint{3.598484in}{9.589350in}}%
\pgfpathquadraticcurveto{\pgfqpoint{3.591949in}{9.595576in}}{\pgfqpoint{3.591949in}{9.606008in}}%
\pgfpathquadraticcurveto{\pgfqpoint{3.591949in}{9.618172in}}{\pgfqpoint{3.600092in}{9.624357in}}%
\pgfpathquadraticcurveto{\pgfqpoint{3.608257in}{9.630541in}}{\pgfqpoint{3.624420in}{9.630541in}}%
\pgfpathlineto{\pgfqpoint{3.641057in}{9.630541in}}%
\pgfpathlineto{\pgfqpoint{3.641057in}{9.631717in}}%
\pgfpathquadraticcurveto{\pgfqpoint{3.641057in}{9.639901in}}{\pgfqpoint{3.635676in}{9.644375in}}%
\pgfpathquadraticcurveto{\pgfqpoint{3.630295in}{9.648849in}}{\pgfqpoint{3.620564in}{9.648849in}}%
\pgfpathquadraticcurveto{\pgfqpoint{3.614380in}{9.648849in}}{\pgfqpoint{3.608504in}{9.647364in}}%
\pgfpathquadraticcurveto{\pgfqpoint{3.602649in}{9.645880in}}{\pgfqpoint{3.597247in}{9.642911in}}%
\pgfpathlineto{\pgfqpoint{3.597247in}{9.653879in}}%
\pgfpathquadraticcurveto{\pgfqpoint{3.603742in}{9.656394in}}{\pgfqpoint{3.609865in}{9.657631in}}%
\pgfpathquadraticcurveto{\pgfqpoint{3.615988in}{9.658889in}}{\pgfqpoint{3.621781in}{9.658889in}}%
\pgfpathquadraticcurveto{\pgfqpoint{3.637449in}{9.658889in}}{\pgfqpoint{3.645180in}{9.650766in}}%
\pgfpathquadraticcurveto{\pgfqpoint{3.652911in}{9.642664in}}{\pgfqpoint{3.652911in}{9.626171in}}%
\pgfpathlineto{\pgfqpoint{3.652911in}{9.626171in}}%
\pgfpathclose%
\pgfpathmoveto{\pgfqpoint{3.688784in}{9.595824in}}%
\pgfpathlineto{\pgfqpoint{3.688784in}{9.557560in}}%
\pgfpathlineto{\pgfqpoint{3.676868in}{9.557560in}}%
\pgfpathlineto{\pgfqpoint{3.676868in}{9.657157in}}%
\pgfpathlineto{\pgfqpoint{3.688784in}{9.657157in}}%
\pgfpathlineto{\pgfqpoint{3.688784in}{9.646210in}}%
\pgfpathquadraticcurveto{\pgfqpoint{3.692536in}{9.652642in}}{\pgfqpoint{3.698226in}{9.655755in}}%
\pgfpathquadraticcurveto{\pgfqpoint{3.703937in}{9.658889in}}{\pgfqpoint{3.711854in}{9.658889in}}%
\pgfpathquadraticcurveto{\pgfqpoint{3.725007in}{9.658889in}}{\pgfqpoint{3.733212in}{9.648457in}}%
\pgfpathquadraticcurveto{\pgfqpoint{3.741438in}{9.638025in}}{\pgfqpoint{3.741438in}{9.621017in}}%
\pgfpathquadraticcurveto{\pgfqpoint{3.741438in}{9.604008in}}{\pgfqpoint{3.733212in}{9.593556in}}%
\pgfpathquadraticcurveto{\pgfqpoint{3.725007in}{9.583124in}}{\pgfqpoint{3.711854in}{9.583124in}}%
\pgfpathquadraticcurveto{\pgfqpoint{3.703937in}{9.583124in}}{\pgfqpoint{3.698226in}{9.586258in}}%
\pgfpathquadraticcurveto{\pgfqpoint{3.692536in}{9.589391in}}{\pgfqpoint{3.688784in}{9.595824in}}%
\pgfpathlineto{\pgfqpoint{3.688784in}{9.595824in}}%
\pgfpathclose%
\pgfpathmoveto{\pgfqpoint{3.729130in}{9.621017in}}%
\pgfpathquadraticcurveto{\pgfqpoint{3.729130in}{9.634087in}}{\pgfqpoint{3.723749in}{9.641530in}}%
\pgfpathquadraticcurveto{\pgfqpoint{3.718368in}{9.648972in}}{\pgfqpoint{3.708967in}{9.648972in}}%
\pgfpathquadraticcurveto{\pgfqpoint{3.699546in}{9.648972in}}{\pgfqpoint{3.694165in}{9.641530in}}%
\pgfpathquadraticcurveto{\pgfqpoint{3.688784in}{9.634087in}}{\pgfqpoint{3.688784in}{9.621017in}}%
\pgfpathquadraticcurveto{\pgfqpoint{3.688784in}{9.607946in}}{\pgfqpoint{3.694165in}{9.600503in}}%
\pgfpathquadraticcurveto{\pgfqpoint{3.699546in}{9.593061in}}{\pgfqpoint{3.708967in}{9.593061in}}%
\pgfpathquadraticcurveto{\pgfqpoint{3.718368in}{9.593061in}}{\pgfqpoint{3.723749in}{9.600503in}}%
\pgfpathquadraticcurveto{\pgfqpoint{3.729130in}{9.607946in}}{\pgfqpoint{3.729130in}{9.621017in}}%
\pgfpathlineto{\pgfqpoint{3.729130in}{9.621017in}}%
\pgfpathclose%
\pgfpathmoveto{\pgfqpoint{3.772548in}{9.595824in}}%
\pgfpathlineto{\pgfqpoint{3.772548in}{9.557560in}}%
\pgfpathlineto{\pgfqpoint{3.760632in}{9.557560in}}%
\pgfpathlineto{\pgfqpoint{3.760632in}{9.657157in}}%
\pgfpathlineto{\pgfqpoint{3.772548in}{9.657157in}}%
\pgfpathlineto{\pgfqpoint{3.772548in}{9.646210in}}%
\pgfpathquadraticcurveto{\pgfqpoint{3.776300in}{9.652642in}}{\pgfqpoint{3.781990in}{9.655755in}}%
\pgfpathquadraticcurveto{\pgfqpoint{3.787701in}{9.658889in}}{\pgfqpoint{3.795618in}{9.658889in}}%
\pgfpathquadraticcurveto{\pgfqpoint{3.808771in}{9.658889in}}{\pgfqpoint{3.816976in}{9.648457in}}%
\pgfpathquadraticcurveto{\pgfqpoint{3.825202in}{9.638025in}}{\pgfqpoint{3.825202in}{9.621017in}}%
\pgfpathquadraticcurveto{\pgfqpoint{3.825202in}{9.604008in}}{\pgfqpoint{3.816976in}{9.593556in}}%
\pgfpathquadraticcurveto{\pgfqpoint{3.808771in}{9.583124in}}{\pgfqpoint{3.795618in}{9.583124in}}%
\pgfpathquadraticcurveto{\pgfqpoint{3.787701in}{9.583124in}}{\pgfqpoint{3.781990in}{9.586258in}}%
\pgfpathquadraticcurveto{\pgfqpoint{3.776300in}{9.589391in}}{\pgfqpoint{3.772548in}{9.595824in}}%
\pgfpathlineto{\pgfqpoint{3.772548in}{9.595824in}}%
\pgfpathclose%
\pgfpathmoveto{\pgfqpoint{3.812894in}{9.621017in}}%
\pgfpathquadraticcurveto{\pgfqpoint{3.812894in}{9.634087in}}{\pgfqpoint{3.807513in}{9.641530in}}%
\pgfpathquadraticcurveto{\pgfqpoint{3.802132in}{9.648972in}}{\pgfqpoint{3.792731in}{9.648972in}}%
\pgfpathquadraticcurveto{\pgfqpoint{3.783310in}{9.648972in}}{\pgfqpoint{3.777929in}{9.641530in}}%
\pgfpathquadraticcurveto{\pgfqpoint{3.772548in}{9.634087in}}{\pgfqpoint{3.772548in}{9.621017in}}%
\pgfpathquadraticcurveto{\pgfqpoint{3.772548in}{9.607946in}}{\pgfqpoint{3.777929in}{9.600503in}}%
\pgfpathquadraticcurveto{\pgfqpoint{3.783310in}{9.593061in}}{\pgfqpoint{3.792731in}{9.593061in}}%
\pgfpathquadraticcurveto{\pgfqpoint{3.802132in}{9.593061in}}{\pgfqpoint{3.807513in}{9.600503in}}%
\pgfpathquadraticcurveto{\pgfqpoint{3.812894in}{9.607946in}}{\pgfqpoint{3.812894in}{9.621017in}}%
\pgfpathlineto{\pgfqpoint{3.812894in}{9.621017in}}%
\pgfpathclose%
\pgfpathmoveto{\pgfqpoint{3.906575in}{9.624047in}}%
\pgfpathlineto{\pgfqpoint{3.906575in}{9.618254in}}%
\pgfpathlineto{\pgfqpoint{3.852065in}{9.618254in}}%
\pgfpathquadraticcurveto{\pgfqpoint{3.852849in}{9.606008in}}{\pgfqpoint{3.859446in}{9.599596in}}%
\pgfpathquadraticcurveto{\pgfqpoint{3.866043in}{9.593185in}}{\pgfqpoint{3.877836in}{9.593185in}}%
\pgfpathquadraticcurveto{\pgfqpoint{3.884660in}{9.593185in}}{\pgfqpoint{3.891071in}{9.594855in}}%
\pgfpathquadraticcurveto{\pgfqpoint{3.897483in}{9.596525in}}{\pgfqpoint{3.903812in}{9.599885in}}%
\pgfpathlineto{\pgfqpoint{3.903812in}{9.588670in}}%
\pgfpathquadraticcurveto{\pgfqpoint{3.897421in}{9.585969in}}{\pgfqpoint{3.890721in}{9.584546in}}%
\pgfpathquadraticcurveto{\pgfqpoint{3.884020in}{9.583124in}}{\pgfqpoint{3.877135in}{9.583124in}}%
\pgfpathquadraticcurveto{\pgfqpoint{3.859858in}{9.583124in}}{\pgfqpoint{3.849777in}{9.593164in}}%
\pgfpathquadraticcurveto{\pgfqpoint{3.839695in}{9.603225in}}{\pgfqpoint{3.839695in}{9.620378in}}%
\pgfpathquadraticcurveto{\pgfqpoint{3.839695in}{9.638087in}}{\pgfqpoint{3.849261in}{9.648478in}}%
\pgfpathquadraticcurveto{\pgfqpoint{3.858827in}{9.658889in}}{\pgfqpoint{3.875073in}{9.658889in}}%
\pgfpathquadraticcurveto{\pgfqpoint{3.889628in}{9.658889in}}{\pgfqpoint{3.898101in}{9.649508in}}%
\pgfpathquadraticcurveto{\pgfqpoint{3.906575in}{9.640149in}}{\pgfqpoint{3.906575in}{9.624047in}}%
\pgfpathlineto{\pgfqpoint{3.906575in}{9.624047in}}%
\pgfpathclose%
\pgfpathmoveto{\pgfqpoint{3.894720in}{9.627531in}}%
\pgfpathquadraticcurveto{\pgfqpoint{3.894597in}{9.637242in}}{\pgfqpoint{3.889278in}{9.643035in}}%
\pgfpathquadraticcurveto{\pgfqpoint{3.883959in}{9.648849in}}{\pgfqpoint{3.875197in}{9.648849in}}%
\pgfpathquadraticcurveto{\pgfqpoint{3.865280in}{9.648849in}}{\pgfqpoint{3.859322in}{9.643241in}}%
\pgfpathquadraticcurveto{\pgfqpoint{3.853364in}{9.637633in}}{\pgfqpoint{3.852457in}{9.627449in}}%
\pgfpathlineto{\pgfqpoint{3.894720in}{9.627531in}}%
\pgfpathlineto{\pgfqpoint{3.894720in}{9.627531in}}%
\pgfpathclose%
\pgfpathmoveto{\pgfqpoint{3.973516in}{9.646210in}}%
\pgfpathlineto{\pgfqpoint{3.973516in}{9.685257in}}%
\pgfpathlineto{\pgfqpoint{3.985370in}{9.685257in}}%
\pgfpathlineto{\pgfqpoint{3.985370in}{9.585000in}}%
\pgfpathlineto{\pgfqpoint{3.973516in}{9.585000in}}%
\pgfpathlineto{\pgfqpoint{3.973516in}{9.595824in}}%
\pgfpathquadraticcurveto{\pgfqpoint{3.969784in}{9.589391in}}{\pgfqpoint{3.964074in}{9.586258in}}%
\pgfpathquadraticcurveto{\pgfqpoint{3.958383in}{9.583124in}}{\pgfqpoint{3.950384in}{9.583124in}}%
\pgfpathquadraticcurveto{\pgfqpoint{3.937314in}{9.583124in}}{\pgfqpoint{3.929088in}{9.593556in}}%
\pgfpathquadraticcurveto{\pgfqpoint{3.920882in}{9.604008in}}{\pgfqpoint{3.920882in}{9.621017in}}%
\pgfpathquadraticcurveto{\pgfqpoint{3.920882in}{9.638025in}}{\pgfqpoint{3.929088in}{9.648457in}}%
\pgfpathquadraticcurveto{\pgfqpoint{3.937314in}{9.658889in}}{\pgfqpoint{3.950384in}{9.658889in}}%
\pgfpathquadraticcurveto{\pgfqpoint{3.958383in}{9.658889in}}{\pgfqpoint{3.964074in}{9.655755in}}%
\pgfpathquadraticcurveto{\pgfqpoint{3.969784in}{9.652642in}}{\pgfqpoint{3.973516in}{9.646210in}}%
\pgfpathlineto{\pgfqpoint{3.973516in}{9.646210in}}%
\pgfpathclose%
\pgfpathmoveto{\pgfqpoint{3.933128in}{9.621017in}}%
\pgfpathquadraticcurveto{\pgfqpoint{3.933128in}{9.607946in}}{\pgfqpoint{3.938509in}{9.600503in}}%
\pgfpathquadraticcurveto{\pgfqpoint{3.943890in}{9.593061in}}{\pgfqpoint{3.953291in}{9.593061in}}%
\pgfpathquadraticcurveto{\pgfqpoint{3.962692in}{9.593061in}}{\pgfqpoint{3.968094in}{9.600503in}}%
\pgfpathquadraticcurveto{\pgfqpoint{3.973516in}{9.607946in}}{\pgfqpoint{3.973516in}{9.621017in}}%
\pgfpathquadraticcurveto{\pgfqpoint{3.973516in}{9.634087in}}{\pgfqpoint{3.968094in}{9.641530in}}%
\pgfpathquadraticcurveto{\pgfqpoint{3.962692in}{9.648972in}}{\pgfqpoint{3.953291in}{9.648972in}}%
\pgfpathquadraticcurveto{\pgfqpoint{3.943890in}{9.648972in}}{\pgfqpoint{3.938509in}{9.641530in}}%
\pgfpathquadraticcurveto{\pgfqpoint{3.933128in}{9.634087in}}{\pgfqpoint{3.933128in}{9.621017in}}%
\pgfpathlineto{\pgfqpoint{3.933128in}{9.621017in}}%
\pgfpathclose%
\pgfpathmoveto{\pgfqpoint{4.080205in}{9.685113in}}%
\pgfpathquadraticcurveto{\pgfqpoint{4.071588in}{9.670310in}}{\pgfqpoint{4.067382in}{9.655796in}}%
\pgfpathquadraticcurveto{\pgfqpoint{4.063197in}{9.641303in}}{\pgfqpoint{4.063197in}{9.626418in}}%
\pgfpathquadraticcurveto{\pgfqpoint{4.063197in}{9.611554in}}{\pgfqpoint{4.067423in}{9.596957in}}%
\pgfpathquadraticcurveto{\pgfqpoint{4.071650in}{9.582361in}}{\pgfqpoint{4.080205in}{9.567600in}}%
\pgfpathlineto{\pgfqpoint{4.069897in}{9.567600in}}%
\pgfpathquadraticcurveto{\pgfqpoint{4.060249in}{9.582753in}}{\pgfqpoint{4.055445in}{9.597370in}}%
\pgfpathquadraticcurveto{\pgfqpoint{4.050642in}{9.611987in}}{\pgfqpoint{4.050642in}{9.626418in}}%
\pgfpathquadraticcurveto{\pgfqpoint{4.050642in}{9.640788in}}{\pgfqpoint{4.055404in}{9.655343in}}%
\pgfpathquadraticcurveto{\pgfqpoint{4.060187in}{9.669919in}}{\pgfqpoint{4.069897in}{9.685113in}}%
\pgfpathlineto{\pgfqpoint{4.080205in}{9.685113in}}%
\pgfpathlineto{\pgfqpoint{4.080205in}{9.685113in}}%
\pgfpathclose%
\pgfpathmoveto{\pgfqpoint{4.150693in}{9.646210in}}%
\pgfpathlineto{\pgfqpoint{4.150693in}{9.685257in}}%
\pgfpathlineto{\pgfqpoint{4.162547in}{9.685257in}}%
\pgfpathlineto{\pgfqpoint{4.162547in}{9.585000in}}%
\pgfpathlineto{\pgfqpoint{4.150693in}{9.585000in}}%
\pgfpathlineto{\pgfqpoint{4.150693in}{9.595824in}}%
\pgfpathquadraticcurveto{\pgfqpoint{4.146961in}{9.589391in}}{\pgfqpoint{4.141250in}{9.586258in}}%
\pgfpathquadraticcurveto{\pgfqpoint{4.135560in}{9.583124in}}{\pgfqpoint{4.127561in}{9.583124in}}%
\pgfpathquadraticcurveto{\pgfqpoint{4.114490in}{9.583124in}}{\pgfqpoint{4.106264in}{9.593556in}}%
\pgfpathquadraticcurveto{\pgfqpoint{4.098059in}{9.604008in}}{\pgfqpoint{4.098059in}{9.621017in}}%
\pgfpathquadraticcurveto{\pgfqpoint{4.098059in}{9.638025in}}{\pgfqpoint{4.106264in}{9.648457in}}%
\pgfpathquadraticcurveto{\pgfqpoint{4.114490in}{9.658889in}}{\pgfqpoint{4.127561in}{9.658889in}}%
\pgfpathquadraticcurveto{\pgfqpoint{4.135560in}{9.658889in}}{\pgfqpoint{4.141250in}{9.655755in}}%
\pgfpathquadraticcurveto{\pgfqpoint{4.146961in}{9.652642in}}{\pgfqpoint{4.150693in}{9.646210in}}%
\pgfpathlineto{\pgfqpoint{4.150693in}{9.646210in}}%
\pgfpathclose%
\pgfpathmoveto{\pgfqpoint{4.110305in}{9.621017in}}%
\pgfpathquadraticcurveto{\pgfqpoint{4.110305in}{9.607946in}}{\pgfqpoint{4.115686in}{9.600503in}}%
\pgfpathquadraticcurveto{\pgfqpoint{4.121067in}{9.593061in}}{\pgfqpoint{4.130468in}{9.593061in}}%
\pgfpathquadraticcurveto{\pgfqpoint{4.139869in}{9.593061in}}{\pgfqpoint{4.145270in}{9.600503in}}%
\pgfpathquadraticcurveto{\pgfqpoint{4.150693in}{9.607946in}}{\pgfqpoint{4.150693in}{9.621017in}}%
\pgfpathquadraticcurveto{\pgfqpoint{4.150693in}{9.634087in}}{\pgfqpoint{4.145270in}{9.641530in}}%
\pgfpathquadraticcurveto{\pgfqpoint{4.139869in}{9.648972in}}{\pgfqpoint{4.130468in}{9.648972in}}%
\pgfpathquadraticcurveto{\pgfqpoint{4.121067in}{9.648972in}}{\pgfqpoint{4.115686in}{9.641530in}}%
\pgfpathquadraticcurveto{\pgfqpoint{4.110305in}{9.634087in}}{\pgfqpoint{4.110305in}{9.621017in}}%
\pgfpathlineto{\pgfqpoint{4.110305in}{9.621017in}}%
\pgfpathclose%
\pgfpathmoveto{\pgfqpoint{4.186977in}{9.657157in}}%
\pgfpathlineto{\pgfqpoint{4.198832in}{9.657157in}}%
\pgfpathlineto{\pgfqpoint{4.198832in}{9.585000in}}%
\pgfpathlineto{\pgfqpoint{4.186977in}{9.585000in}}%
\pgfpathlineto{\pgfqpoint{4.186977in}{9.657157in}}%
\pgfpathlineto{\pgfqpoint{4.186977in}{9.657157in}}%
\pgfpathclose%
\pgfpathmoveto{\pgfqpoint{4.186977in}{9.685257in}}%
\pgfpathlineto{\pgfqpoint{4.198832in}{9.685257in}}%
\pgfpathlineto{\pgfqpoint{4.198832in}{9.670228in}}%
\pgfpathlineto{\pgfqpoint{4.186977in}{9.670228in}}%
\pgfpathlineto{\pgfqpoint{4.186977in}{9.685257in}}%
\pgfpathlineto{\pgfqpoint{4.186977in}{9.685257in}}%
\pgfpathclose%
\pgfpathmoveto{\pgfqpoint{4.256434in}{9.621264in}}%
\pgfpathquadraticcurveto{\pgfqpoint{4.242064in}{9.621264in}}{\pgfqpoint{4.236518in}{9.617986in}}%
\pgfpathquadraticcurveto{\pgfqpoint{4.230972in}{9.614708in}}{\pgfqpoint{4.230972in}{9.606771in}}%
\pgfpathquadraticcurveto{\pgfqpoint{4.230972in}{9.600462in}}{\pgfqpoint{4.235137in}{9.596751in}}%
\pgfpathquadraticcurveto{\pgfqpoint{4.239301in}{9.593061in}}{\pgfqpoint{4.246435in}{9.593061in}}%
\pgfpathquadraticcurveto{\pgfqpoint{4.256310in}{9.593061in}}{\pgfqpoint{4.262268in}{9.600050in}}%
\pgfpathquadraticcurveto{\pgfqpoint{4.268226in}{9.607039in}}{\pgfqpoint{4.268226in}{9.618625in}}%
\pgfpathlineto{\pgfqpoint{4.268226in}{9.621264in}}%
\pgfpathlineto{\pgfqpoint{4.256434in}{9.621264in}}%
\pgfpathlineto{\pgfqpoint{4.256434in}{9.621264in}}%
\pgfpathclose%
\pgfpathmoveto{\pgfqpoint{4.280081in}{9.626171in}}%
\pgfpathlineto{\pgfqpoint{4.280081in}{9.585000in}}%
\pgfpathlineto{\pgfqpoint{4.268226in}{9.585000in}}%
\pgfpathlineto{\pgfqpoint{4.268226in}{9.595947in}}%
\pgfpathquadraticcurveto{\pgfqpoint{4.264165in}{9.589391in}}{\pgfqpoint{4.258104in}{9.586258in}}%
\pgfpathquadraticcurveto{\pgfqpoint{4.252042in}{9.583124in}}{\pgfqpoint{4.243280in}{9.583124in}}%
\pgfpathquadraticcurveto{\pgfqpoint{4.232209in}{9.583124in}}{\pgfqpoint{4.225653in}{9.589350in}}%
\pgfpathquadraticcurveto{\pgfqpoint{4.219118in}{9.595576in}}{\pgfqpoint{4.219118in}{9.606008in}}%
\pgfpathquadraticcurveto{\pgfqpoint{4.219118in}{9.618172in}}{\pgfqpoint{4.227262in}{9.624357in}}%
\pgfpathquadraticcurveto{\pgfqpoint{4.235426in}{9.630541in}}{\pgfqpoint{4.251589in}{9.630541in}}%
\pgfpathlineto{\pgfqpoint{4.268226in}{9.630541in}}%
\pgfpathlineto{\pgfqpoint{4.268226in}{9.631717in}}%
\pgfpathquadraticcurveto{\pgfqpoint{4.268226in}{9.639901in}}{\pgfqpoint{4.262845in}{9.644375in}}%
\pgfpathquadraticcurveto{\pgfqpoint{4.257464in}{9.648849in}}{\pgfqpoint{4.247734in}{9.648849in}}%
\pgfpathquadraticcurveto{\pgfqpoint{4.241549in}{9.648849in}}{\pgfqpoint{4.235673in}{9.647364in}}%
\pgfpathquadraticcurveto{\pgfqpoint{4.229818in}{9.645880in}}{\pgfqpoint{4.224416in}{9.642911in}}%
\pgfpathlineto{\pgfqpoint{4.224416in}{9.653879in}}%
\pgfpathquadraticcurveto{\pgfqpoint{4.230911in}{9.656394in}}{\pgfqpoint{4.237034in}{9.657631in}}%
\pgfpathquadraticcurveto{\pgfqpoint{4.243157in}{9.658889in}}{\pgfqpoint{4.248950in}{9.658889in}}%
\pgfpathquadraticcurveto{\pgfqpoint{4.264618in}{9.658889in}}{\pgfqpoint{4.272349in}{9.650766in}}%
\pgfpathquadraticcurveto{\pgfqpoint{4.280081in}{9.642664in}}{\pgfqpoint{4.280081in}{9.626171in}}%
\pgfpathlineto{\pgfqpoint{4.280081in}{9.626171in}}%
\pgfpathclose%
\pgfpathmoveto{\pgfqpoint{4.360670in}{9.643303in}}%
\pgfpathquadraticcurveto{\pgfqpoint{4.365123in}{9.651302in}}{\pgfqpoint{4.371308in}{9.655095in}}%
\pgfpathquadraticcurveto{\pgfqpoint{4.377493in}{9.658889in}}{\pgfqpoint{4.385863in}{9.658889in}}%
\pgfpathquadraticcurveto{\pgfqpoint{4.397140in}{9.658889in}}{\pgfqpoint{4.403263in}{9.650993in}}%
\pgfpathquadraticcurveto{\pgfqpoint{4.409386in}{9.643117in}}{\pgfqpoint{4.409386in}{9.628562in}}%
\pgfpathlineto{\pgfqpoint{4.409386in}{9.585000in}}%
\pgfpathlineto{\pgfqpoint{4.397470in}{9.585000in}}%
\pgfpathlineto{\pgfqpoint{4.397470in}{9.628171in}}%
\pgfpathquadraticcurveto{\pgfqpoint{4.397470in}{9.638541in}}{\pgfqpoint{4.393780in}{9.643550in}}%
\pgfpathquadraticcurveto{\pgfqpoint{4.390110in}{9.648581in}}{\pgfqpoint{4.382585in}{9.648581in}}%
\pgfpathquadraticcurveto{\pgfqpoint{4.373369in}{9.648581in}}{\pgfqpoint{4.368009in}{9.642458in}}%
\pgfpathquadraticcurveto{\pgfqpoint{4.362670in}{9.636355in}}{\pgfqpoint{4.362670in}{9.625779in}}%
\pgfpathlineto{\pgfqpoint{4.362670in}{9.585000in}}%
\pgfpathlineto{\pgfqpoint{4.350753in}{9.585000in}}%
\pgfpathlineto{\pgfqpoint{4.350753in}{9.628171in}}%
\pgfpathquadraticcurveto{\pgfqpoint{4.350753in}{9.638602in}}{\pgfqpoint{4.347084in}{9.643592in}}%
\pgfpathquadraticcurveto{\pgfqpoint{4.343414in}{9.648581in}}{\pgfqpoint{4.335745in}{9.648581in}}%
\pgfpathquadraticcurveto{\pgfqpoint{4.326653in}{9.648581in}}{\pgfqpoint{4.321293in}{9.642437in}}%
\pgfpathquadraticcurveto{\pgfqpoint{4.315953in}{9.636293in}}{\pgfqpoint{4.315953in}{9.625779in}}%
\pgfpathlineto{\pgfqpoint{4.315953in}{9.585000in}}%
\pgfpathlineto{\pgfqpoint{4.304037in}{9.585000in}}%
\pgfpathlineto{\pgfqpoint{4.304037in}{9.657157in}}%
\pgfpathlineto{\pgfqpoint{4.315953in}{9.657157in}}%
\pgfpathlineto{\pgfqpoint{4.315953in}{9.645942in}}%
\pgfpathquadraticcurveto{\pgfqpoint{4.320014in}{9.652580in}}{\pgfqpoint{4.325684in}{9.655735in}}%
\pgfpathquadraticcurveto{\pgfqpoint{4.331353in}{9.658889in}}{\pgfqpoint{4.339146in}{9.658889in}}%
\pgfpathquadraticcurveto{\pgfqpoint{4.347022in}{9.658889in}}{\pgfqpoint{4.352526in}{9.654889in}}%
\pgfpathquadraticcurveto{\pgfqpoint{4.358031in}{9.650910in}}{\pgfqpoint{4.360670in}{9.643303in}}%
\pgfpathlineto{\pgfqpoint{4.360670in}{9.643303in}}%
\pgfpathclose%
\pgfpathmoveto{\pgfqpoint{4.460968in}{9.648849in}}%
\pgfpathquadraticcurveto{\pgfqpoint{4.451443in}{9.648849in}}{\pgfqpoint{4.445898in}{9.641406in}}%
\pgfpathquadraticcurveto{\pgfqpoint{4.440352in}{9.633964in}}{\pgfqpoint{4.440352in}{9.621017in}}%
\pgfpathquadraticcurveto{\pgfqpoint{4.440352in}{9.608070in}}{\pgfqpoint{4.445856in}{9.600627in}}%
\pgfpathquadraticcurveto{\pgfqpoint{4.451382in}{9.593185in}}{\pgfqpoint{4.460968in}{9.593185in}}%
\pgfpathquadraticcurveto{\pgfqpoint{4.470452in}{9.593185in}}{\pgfqpoint{4.475977in}{9.600648in}}%
\pgfpathquadraticcurveto{\pgfqpoint{4.481523in}{9.608132in}}{\pgfqpoint{4.481523in}{9.621017in}}%
\pgfpathquadraticcurveto{\pgfqpoint{4.481523in}{9.633840in}}{\pgfqpoint{4.475977in}{9.641344in}}%
\pgfpathquadraticcurveto{\pgfqpoint{4.470452in}{9.648849in}}{\pgfqpoint{4.460968in}{9.648849in}}%
\pgfpathlineto{\pgfqpoint{4.460968in}{9.648849in}}%
\pgfpathclose%
\pgfpathmoveto{\pgfqpoint{4.460968in}{9.658889in}}%
\pgfpathquadraticcurveto{\pgfqpoint{4.476430in}{9.658889in}}{\pgfqpoint{4.485254in}{9.648828in}}%
\pgfpathquadraticcurveto{\pgfqpoint{4.494099in}{9.638788in}}{\pgfqpoint{4.494099in}{9.621017in}}%
\pgfpathquadraticcurveto{\pgfqpoint{4.494099in}{9.603307in}}{\pgfqpoint{4.485254in}{9.593205in}}%
\pgfpathquadraticcurveto{\pgfqpoint{4.476430in}{9.583124in}}{\pgfqpoint{4.460968in}{9.583124in}}%
\pgfpathquadraticcurveto{\pgfqpoint{4.445444in}{9.583124in}}{\pgfqpoint{4.436641in}{9.593205in}}%
\pgfpathquadraticcurveto{\pgfqpoint{4.427858in}{9.603307in}}{\pgfqpoint{4.427858in}{9.621017in}}%
\pgfpathquadraticcurveto{\pgfqpoint{4.427858in}{9.638788in}}{\pgfqpoint{4.436641in}{9.648828in}}%
\pgfpathquadraticcurveto{\pgfqpoint{4.445444in}{9.658889in}}{\pgfqpoint{4.460968in}{9.658889in}}%
\pgfpathlineto{\pgfqpoint{4.460968in}{9.658889in}}%
\pgfpathclose%
\pgfpathmoveto{\pgfqpoint{4.573739in}{9.628562in}}%
\pgfpathlineto{\pgfqpoint{4.573739in}{9.585000in}}%
\pgfpathlineto{\pgfqpoint{4.561885in}{9.585000in}}%
\pgfpathlineto{\pgfqpoint{4.561885in}{9.628171in}}%
\pgfpathquadraticcurveto{\pgfqpoint{4.561885in}{9.638417in}}{\pgfqpoint{4.557885in}{9.643488in}}%
\pgfpathquadraticcurveto{\pgfqpoint{4.553886in}{9.648581in}}{\pgfqpoint{4.545907in}{9.648581in}}%
\pgfpathquadraticcurveto{\pgfqpoint{4.536300in}{9.648581in}}{\pgfqpoint{4.530754in}{9.642458in}}%
\pgfpathquadraticcurveto{\pgfqpoint{4.525209in}{9.636355in}}{\pgfqpoint{4.525209in}{9.625779in}}%
\pgfpathlineto{\pgfqpoint{4.525209in}{9.585000in}}%
\pgfpathlineto{\pgfqpoint{4.513292in}{9.585000in}}%
\pgfpathlineto{\pgfqpoint{4.513292in}{9.657157in}}%
\pgfpathlineto{\pgfqpoint{4.525209in}{9.657157in}}%
\pgfpathlineto{\pgfqpoint{4.525209in}{9.645942in}}%
\pgfpathquadraticcurveto{\pgfqpoint{4.529476in}{9.652457in}}{\pgfqpoint{4.535228in}{9.655673in}}%
\pgfpathquadraticcurveto{\pgfqpoint{4.541001in}{9.658889in}}{\pgfqpoint{4.548546in}{9.658889in}}%
\pgfpathquadraticcurveto{\pgfqpoint{4.560978in}{9.658889in}}{\pgfqpoint{4.567348in}{9.651199in}}%
\pgfpathquadraticcurveto{\pgfqpoint{4.573739in}{9.643509in}}{\pgfqpoint{4.573739in}{9.628562in}}%
\pgfpathlineto{\pgfqpoint{4.573739in}{9.628562in}}%
\pgfpathclose%
\pgfpathmoveto{\pgfqpoint{4.644845in}{9.646210in}}%
\pgfpathlineto{\pgfqpoint{4.644845in}{9.685257in}}%
\pgfpathlineto{\pgfqpoint{4.656699in}{9.685257in}}%
\pgfpathlineto{\pgfqpoint{4.656699in}{9.585000in}}%
\pgfpathlineto{\pgfqpoint{4.644845in}{9.585000in}}%
\pgfpathlineto{\pgfqpoint{4.644845in}{9.595824in}}%
\pgfpathquadraticcurveto{\pgfqpoint{4.641114in}{9.589391in}}{\pgfqpoint{4.635403in}{9.586258in}}%
\pgfpathquadraticcurveto{\pgfqpoint{4.629713in}{9.583124in}}{\pgfqpoint{4.621714in}{9.583124in}}%
\pgfpathquadraticcurveto{\pgfqpoint{4.608643in}{9.583124in}}{\pgfqpoint{4.600417in}{9.593556in}}%
\pgfpathquadraticcurveto{\pgfqpoint{4.592212in}{9.604008in}}{\pgfqpoint{4.592212in}{9.621017in}}%
\pgfpathquadraticcurveto{\pgfqpoint{4.592212in}{9.638025in}}{\pgfqpoint{4.600417in}{9.648457in}}%
\pgfpathquadraticcurveto{\pgfqpoint{4.608643in}{9.658889in}}{\pgfqpoint{4.621714in}{9.658889in}}%
\pgfpathquadraticcurveto{\pgfqpoint{4.629713in}{9.658889in}}{\pgfqpoint{4.635403in}{9.655755in}}%
\pgfpathquadraticcurveto{\pgfqpoint{4.641114in}{9.652642in}}{\pgfqpoint{4.644845in}{9.646210in}}%
\pgfpathlineto{\pgfqpoint{4.644845in}{9.646210in}}%
\pgfpathclose%
\pgfpathmoveto{\pgfqpoint{4.604458in}{9.621017in}}%
\pgfpathquadraticcurveto{\pgfqpoint{4.604458in}{9.607946in}}{\pgfqpoint{4.609839in}{9.600503in}}%
\pgfpathquadraticcurveto{\pgfqpoint{4.615219in}{9.593061in}}{\pgfqpoint{4.624620in}{9.593061in}}%
\pgfpathquadraticcurveto{\pgfqpoint{4.634022in}{9.593061in}}{\pgfqpoint{4.639423in}{9.600503in}}%
\pgfpathquadraticcurveto{\pgfqpoint{4.644845in}{9.607946in}}{\pgfqpoint{4.644845in}{9.621017in}}%
\pgfpathquadraticcurveto{\pgfqpoint{4.644845in}{9.634087in}}{\pgfqpoint{4.639423in}{9.641530in}}%
\pgfpathquadraticcurveto{\pgfqpoint{4.634022in}{9.648972in}}{\pgfqpoint{4.624620in}{9.648972in}}%
\pgfpathquadraticcurveto{\pgfqpoint{4.615219in}{9.648972in}}{\pgfqpoint{4.609839in}{9.641530in}}%
\pgfpathquadraticcurveto{\pgfqpoint{4.604458in}{9.634087in}}{\pgfqpoint{4.604458in}{9.621017in}}%
\pgfpathlineto{\pgfqpoint{4.604458in}{9.621017in}}%
\pgfpathclose%
\pgfpathmoveto{\pgfqpoint{4.727125in}{9.655034in}}%
\pgfpathlineto{\pgfqpoint{4.727125in}{9.643818in}}%
\pgfpathquadraticcurveto{\pgfqpoint{4.722115in}{9.646395in}}{\pgfqpoint{4.716693in}{9.647674in}}%
\pgfpathquadraticcurveto{\pgfqpoint{4.711291in}{9.648972in}}{\pgfqpoint{4.705478in}{9.648972in}}%
\pgfpathquadraticcurveto{\pgfqpoint{4.696654in}{9.648972in}}{\pgfqpoint{4.692242in}{9.646272in}}%
\pgfpathquadraticcurveto{\pgfqpoint{4.687830in}{9.643571in}}{\pgfqpoint{4.687830in}{9.638149in}}%
\pgfpathquadraticcurveto{\pgfqpoint{4.687830in}{9.634026in}}{\pgfqpoint{4.690984in}{9.631675in}}%
\pgfpathquadraticcurveto{\pgfqpoint{4.694139in}{9.629325in}}{\pgfqpoint{4.703684in}{9.627202in}}%
\pgfpathlineto{\pgfqpoint{4.707745in}{9.626294in}}%
\pgfpathquadraticcurveto{\pgfqpoint{4.720363in}{9.623594in}}{\pgfqpoint{4.725682in}{9.618666in}}%
\pgfpathquadraticcurveto{\pgfqpoint{4.731001in}{9.613739in}}{\pgfqpoint{4.731001in}{9.604915in}}%
\pgfpathquadraticcurveto{\pgfqpoint{4.731001in}{9.594855in}}{\pgfqpoint{4.723043in}{9.588979in}}%
\pgfpathquadraticcurveto{\pgfqpoint{4.715085in}{9.583124in}}{\pgfqpoint{4.701169in}{9.583124in}}%
\pgfpathquadraticcurveto{\pgfqpoint{4.695376in}{9.583124in}}{\pgfqpoint{4.689088in}{9.584258in}}%
\pgfpathquadraticcurveto{\pgfqpoint{4.682800in}{9.585392in}}{\pgfqpoint{4.675852in}{9.587639in}}%
\pgfpathlineto{\pgfqpoint{4.675852in}{9.599885in}}%
\pgfpathquadraticcurveto{\pgfqpoint{4.682429in}{9.596463in}}{\pgfqpoint{4.688799in}{9.594752in}}%
\pgfpathquadraticcurveto{\pgfqpoint{4.695170in}{9.593061in}}{\pgfqpoint{4.701437in}{9.593061in}}%
\pgfpathquadraticcurveto{\pgfqpoint{4.709807in}{9.593061in}}{\pgfqpoint{4.714301in}{9.595927in}}%
\pgfpathquadraticcurveto{\pgfqpoint{4.718816in}{9.598792in}}{\pgfqpoint{4.718816in}{9.604008in}}%
\pgfpathquadraticcurveto{\pgfqpoint{4.718816in}{9.608832in}}{\pgfqpoint{4.715559in}{9.611410in}}%
\pgfpathquadraticcurveto{\pgfqpoint{4.712322in}{9.613987in}}{\pgfqpoint{4.701293in}{9.616378in}}%
\pgfpathlineto{\pgfqpoint{4.697169in}{9.617347in}}%
\pgfpathquadraticcurveto{\pgfqpoint{4.686160in}{9.619656in}}{\pgfqpoint{4.681253in}{9.624460in}}%
\pgfpathquadraticcurveto{\pgfqpoint{4.676367in}{9.629263in}}{\pgfqpoint{4.676367in}{9.637633in}}%
\pgfpathquadraticcurveto{\pgfqpoint{4.676367in}{9.647818in}}{\pgfqpoint{4.683583in}{9.653343in}}%
\pgfpathquadraticcurveto{\pgfqpoint{4.690799in}{9.658889in}}{\pgfqpoint{4.704076in}{9.658889in}}%
\pgfpathquadraticcurveto{\pgfqpoint{4.710632in}{9.658889in}}{\pgfqpoint{4.716425in}{9.657920in}}%
\pgfpathquadraticcurveto{\pgfqpoint{4.722239in}{9.656972in}}{\pgfqpoint{4.727125in}{9.655034in}}%
\pgfpathlineto{\pgfqpoint{4.727125in}{9.655034in}}%
\pgfpathclose%
\pgfpathmoveto{\pgfqpoint{4.748009in}{9.685113in}}%
\pgfpathlineto{\pgfqpoint{4.758317in}{9.685113in}}%
\pgfpathquadraticcurveto{\pgfqpoint{4.767966in}{9.669919in}}{\pgfqpoint{4.772769in}{9.655343in}}%
\pgfpathquadraticcurveto{\pgfqpoint{4.777573in}{9.640788in}}{\pgfqpoint{4.777573in}{9.626418in}}%
\pgfpathquadraticcurveto{\pgfqpoint{4.777573in}{9.611987in}}{\pgfqpoint{4.772769in}{9.597370in}}%
\pgfpathquadraticcurveto{\pgfqpoint{4.767966in}{9.582753in}}{\pgfqpoint{4.758317in}{9.567600in}}%
\pgfpathlineto{\pgfqpoint{4.748009in}{9.567600in}}%
\pgfpathquadraticcurveto{\pgfqpoint{4.756565in}{9.582361in}}{\pgfqpoint{4.760791in}{9.596957in}}%
\pgfpathquadraticcurveto{\pgfqpoint{4.765018in}{9.611554in}}{\pgfqpoint{4.765018in}{9.626418in}}%
\pgfpathquadraticcurveto{\pgfqpoint{4.765018in}{9.641303in}}{\pgfqpoint{4.760791in}{9.655796in}}%
\pgfpathquadraticcurveto{\pgfqpoint{4.756565in}{9.670310in}}{\pgfqpoint{4.748009in}{9.685113in}}%
\pgfpathlineto{\pgfqpoint{4.748009in}{9.685113in}}%
\pgfpathclose%
\pgfusepath{fill}%
\end{pgfscope}%
\begin{pgfscope}%
\pgfsetbuttcap%
\pgfsetroundjoin%
\definecolor{currentfill}{rgb}{0.090196,0.168627,0.227451}%
\pgfsetfillcolor{currentfill}%
\pgfsetlinewidth{0.000000pt}%
\definecolor{currentstroke}{rgb}{0.090196,0.168627,0.227451}%
\pgfsetstrokecolor{currentstroke}%
\pgfsetdash{}{0pt}%
\pgfpathmoveto{\pgfqpoint{1.240187in}{9.846667in}}%
\pgfpathquadraticcurveto{\pgfqpoint{1.230100in}{9.841441in}}{\pgfqpoint{1.219193in}{9.838798in}}%
\pgfpathquadraticcurveto{\pgfqpoint{1.208286in}{9.836124in}}{\pgfqpoint{1.196406in}{9.836124in}}%
\pgfpathquadraticcurveto{\pgfqpoint{1.160981in}{9.836124in}}{\pgfqpoint{1.140291in}{9.855933in}}%
\pgfpathquadraticcurveto{\pgfqpoint{1.119601in}{9.875742in}}{\pgfqpoint{1.119601in}{9.909618in}}%
\pgfpathquadraticcurveto{\pgfqpoint{1.119601in}{9.943615in}}{\pgfqpoint{1.140291in}{9.963394in}}%
\pgfpathquadraticcurveto{\pgfqpoint{1.160981in}{9.983203in}}{\pgfqpoint{1.196406in}{9.983203in}}%
\pgfpathquadraticcurveto{\pgfqpoint{1.208286in}{9.983203in}}{\pgfqpoint{1.219193in}{9.980530in}}%
\pgfpathquadraticcurveto{\pgfqpoint{1.230100in}{9.977886in}}{\pgfqpoint{1.240187in}{9.972661in}}%
\pgfpathlineto{\pgfqpoint{1.240187in}{9.943342in}}%
\pgfpathquadraticcurveto{\pgfqpoint{1.230009in}{9.950269in}}{\pgfqpoint{1.220135in}{9.953490in}}%
\pgfpathquadraticcurveto{\pgfqpoint{1.210261in}{9.956710in}}{\pgfqpoint{1.199354in}{9.956710in}}%
\pgfpathquadraticcurveto{\pgfqpoint{1.179788in}{9.956710in}}{\pgfqpoint{1.168577in}{9.944162in}}%
\pgfpathquadraticcurveto{\pgfqpoint{1.157396in}{9.931645in}}{\pgfqpoint{1.157396in}{9.909618in}}%
\pgfpathquadraticcurveto{\pgfqpoint{1.157396in}{9.887682in}}{\pgfqpoint{1.168577in}{9.875135in}}%
\pgfpathquadraticcurveto{\pgfqpoint{1.179788in}{9.862617in}}{\pgfqpoint{1.199354in}{9.862617in}}%
\pgfpathquadraticcurveto{\pgfqpoint{1.210261in}{9.862617in}}{\pgfqpoint{1.220135in}{9.865838in}}%
\pgfpathquadraticcurveto{\pgfqpoint{1.230009in}{9.869089in}}{\pgfqpoint{1.240187in}{9.876016in}}%
\pgfpathlineto{\pgfqpoint{1.240187in}{9.846667in}}%
\pgfpathlineto{\pgfqpoint{1.240187in}{9.846667in}}%
\pgfpathclose%
\pgfpathmoveto{\pgfqpoint{1.372713in}{9.976185in}}%
\pgfpathlineto{\pgfqpoint{1.372713in}{9.946168in}}%
\pgfpathquadraticcurveto{\pgfqpoint{1.361046in}{9.951393in}}{\pgfqpoint{1.349926in}{9.954036in}}%
\pgfpathquadraticcurveto{\pgfqpoint{1.338837in}{9.956710in}}{\pgfqpoint{1.328963in}{9.956710in}}%
\pgfpathquadraticcurveto{\pgfqpoint{1.315868in}{9.956710in}}{\pgfqpoint{1.309579in}{9.953095in}}%
\pgfpathquadraticcurveto{\pgfqpoint{1.303321in}{9.949510in}}{\pgfqpoint{1.303321in}{9.941914in}}%
\pgfpathquadraticcurveto{\pgfqpoint{1.303321in}{9.936202in}}{\pgfqpoint{1.307544in}{9.933012in}}%
\pgfpathquadraticcurveto{\pgfqpoint{1.311767in}{9.929852in}}{\pgfqpoint{1.322887in}{9.927574in}}%
\pgfpathlineto{\pgfqpoint{1.338442in}{9.924444in}}%
\pgfpathquadraticcurveto{\pgfqpoint{1.362079in}{9.919674in}}{\pgfqpoint{1.372045in}{9.909983in}}%
\pgfpathquadraticcurveto{\pgfqpoint{1.382040in}{9.900321in}}{\pgfqpoint{1.382040in}{9.882457in}}%
\pgfpathquadraticcurveto{\pgfqpoint{1.382040in}{9.859032in}}{\pgfqpoint{1.368125in}{9.847578in}}%
\pgfpathquadraticcurveto{\pgfqpoint{1.354210in}{9.836124in}}{\pgfqpoint{1.325621in}{9.836124in}}%
\pgfpathquadraticcurveto{\pgfqpoint{1.312162in}{9.836124in}}{\pgfqpoint{1.298581in}{9.838707in}}%
\pgfpathquadraticcurveto{\pgfqpoint{1.285000in}{9.841259in}}{\pgfqpoint{1.271420in}{9.846302in}}%
\pgfpathlineto{\pgfqpoint{1.271420in}{9.877140in}}%
\pgfpathquadraticcurveto{\pgfqpoint{1.285000in}{9.869939in}}{\pgfqpoint{1.297670in}{9.866263in}}%
\pgfpathquadraticcurveto{\pgfqpoint{1.310339in}{9.862617in}}{\pgfqpoint{1.322127in}{9.862617in}}%
\pgfpathquadraticcurveto{\pgfqpoint{1.334097in}{9.862617in}}{\pgfqpoint{1.340447in}{9.866597in}}%
\pgfpathquadraticcurveto{\pgfqpoint{1.346797in}{9.870608in}}{\pgfqpoint{1.346797in}{9.878021in}}%
\pgfpathquadraticcurveto{\pgfqpoint{1.346797in}{9.884644in}}{\pgfqpoint{1.342483in}{9.888260in}}%
\pgfpathquadraticcurveto{\pgfqpoint{1.338169in}{9.891875in}}{\pgfqpoint{1.325256in}{9.894731in}}%
\pgfpathlineto{\pgfqpoint{1.311098in}{9.897860in}}%
\pgfpathquadraticcurveto{\pgfqpoint{1.289831in}{9.902418in}}{\pgfqpoint{1.279987in}{9.912383in}}%
\pgfpathquadraticcurveto{\pgfqpoint{1.270174in}{9.922348in}}{\pgfqpoint{1.270174in}{9.939240in}}%
\pgfpathquadraticcurveto{\pgfqpoint{1.270174in}{9.960417in}}{\pgfqpoint{1.283846in}{9.971810in}}%
\pgfpathquadraticcurveto{\pgfqpoint{1.297518in}{9.983203in}}{\pgfqpoint{1.323160in}{9.983203in}}%
\pgfpathquadraticcurveto{\pgfqpoint{1.334857in}{9.983203in}}{\pgfqpoint{1.347192in}{9.981441in}}%
\pgfpathquadraticcurveto{\pgfqpoint{1.359527in}{9.979679in}}{\pgfqpoint{1.372713in}{9.976185in}}%
\pgfpathlineto{\pgfqpoint{1.372713in}{9.976185in}}%
\pgfpathclose%
\pgfpathmoveto{\pgfqpoint{1.500135in}{9.864714in}}%
\pgfpathlineto{\pgfqpoint{1.442986in}{9.864714in}}%
\pgfpathlineto{\pgfqpoint{1.433963in}{9.838889in}}%
\pgfpathlineto{\pgfqpoint{1.397201in}{9.838889in}}%
\pgfpathlineto{\pgfqpoint{1.449701in}{9.980651in}}%
\pgfpathlineto{\pgfqpoint{1.493299in}{9.980651in}}%
\pgfpathlineto{\pgfqpoint{1.545799in}{9.838889in}}%
\pgfpathlineto{\pgfqpoint{1.509067in}{9.838889in}}%
\pgfpathlineto{\pgfqpoint{1.500135in}{9.864714in}}%
\pgfpathlineto{\pgfqpoint{1.500135in}{9.864714in}}%
\pgfpathclose%
\pgfpathmoveto{\pgfqpoint{1.452101in}{9.891024in}}%
\pgfpathlineto{\pgfqpoint{1.490929in}{9.891024in}}%
\pgfpathlineto{\pgfqpoint{1.471545in}{9.947413in}}%
\pgfpathlineto{\pgfqpoint{1.452101in}{9.891024in}}%
\pgfpathlineto{\pgfqpoint{1.452101in}{9.891024in}}%
\pgfpathclose%
\pgfpathmoveto{\pgfqpoint{1.617379in}{9.945226in}}%
\pgfpathlineto{\pgfqpoint{1.651376in}{9.945226in}}%
\pgfpathlineto{\pgfqpoint{1.677869in}{9.871732in}}%
\pgfpathlineto{\pgfqpoint{1.704241in}{9.945226in}}%
\pgfpathlineto{\pgfqpoint{1.738329in}{9.945226in}}%
\pgfpathlineto{\pgfqpoint{1.696463in}{9.838889in}}%
\pgfpathlineto{\pgfqpoint{1.659154in}{9.838889in}}%
\pgfpathlineto{\pgfqpoint{1.617379in}{9.945226in}}%
\pgfpathlineto{\pgfqpoint{1.617379in}{9.945226in}}%
\pgfpathclose%
\pgfpathmoveto{\pgfqpoint{1.840595in}{9.941914in}}%
\pgfpathlineto{\pgfqpoint{1.840595in}{9.916089in}}%
\pgfpathquadraticcurveto{\pgfqpoint{1.829688in}{9.920647in}}{\pgfqpoint{1.819510in}{9.922925in}}%
\pgfpathquadraticcurveto{\pgfqpoint{1.809362in}{9.925204in}}{\pgfqpoint{1.800339in}{9.925204in}}%
\pgfpathquadraticcurveto{\pgfqpoint{1.790647in}{9.925204in}}{\pgfqpoint{1.785938in}{9.922773in}}%
\pgfpathquadraticcurveto{\pgfqpoint{1.781259in}{9.920343in}}{\pgfqpoint{1.781259in}{9.915330in}}%
\pgfpathquadraticcurveto{\pgfqpoint{1.781259in}{9.911228in}}{\pgfqpoint{1.784814in}{9.909041in}}%
\pgfpathquadraticcurveto{\pgfqpoint{1.788368in}{9.906884in}}{\pgfqpoint{1.797574in}{9.905820in}}%
\pgfpathlineto{\pgfqpoint{1.803559in}{9.904970in}}%
\pgfpathquadraticcurveto{\pgfqpoint{1.829688in}{9.901658in}}{\pgfqpoint{1.838681in}{9.894063in}}%
\pgfpathquadraticcurveto{\pgfqpoint{1.847704in}{9.886467in}}{\pgfqpoint{1.847704in}{9.870213in}}%
\pgfpathquadraticcurveto{\pgfqpoint{1.847704in}{9.853229in}}{\pgfqpoint{1.835156in}{9.844661in}}%
\pgfpathquadraticcurveto{\pgfqpoint{1.822639in}{9.836124in}}{\pgfqpoint{1.797787in}{9.836124in}}%
\pgfpathquadraticcurveto{\pgfqpoint{1.787244in}{9.836124in}}{\pgfqpoint{1.775972in}{9.837795in}}%
\pgfpathquadraticcurveto{\pgfqpoint{1.764731in}{9.839466in}}{\pgfqpoint{1.752852in}{9.842778in}}%
\pgfpathlineto{\pgfqpoint{1.752852in}{9.868602in}}%
\pgfpathquadraticcurveto{\pgfqpoint{1.763030in}{9.863681in}}{\pgfqpoint{1.773694in}{9.861189in}}%
\pgfpathquadraticcurveto{\pgfqpoint{1.784388in}{9.858728in}}{\pgfqpoint{1.795387in}{9.858728in}}%
\pgfpathquadraticcurveto{\pgfqpoint{1.805382in}{9.858728in}}{\pgfqpoint{1.810395in}{9.861463in}}%
\pgfpathquadraticcurveto{\pgfqpoint{1.815439in}{9.864227in}}{\pgfqpoint{1.815439in}{9.869666in}}%
\pgfpathquadraticcurveto{\pgfqpoint{1.815439in}{9.874223in}}{\pgfqpoint{1.811975in}{9.876441in}}%
\pgfpathquadraticcurveto{\pgfqpoint{1.808512in}{9.878659in}}{\pgfqpoint{1.798151in}{9.879905in}}%
\pgfpathlineto{\pgfqpoint{1.792166in}{9.880664in}}%
\pgfpathquadraticcurveto{\pgfqpoint{1.769471in}{9.883520in}}{\pgfqpoint{1.760356in}{9.891207in}}%
\pgfpathquadraticcurveto{\pgfqpoint{1.751242in}{9.898893in}}{\pgfqpoint{1.751242in}{9.914570in}}%
\pgfpathquadraticcurveto{\pgfqpoint{1.751242in}{9.931463in}}{\pgfqpoint{1.762817in}{9.939605in}}%
\pgfpathquadraticcurveto{\pgfqpoint{1.774423in}{9.947778in}}{\pgfqpoint{1.798334in}{9.947778in}}%
\pgfpathquadraticcurveto{\pgfqpoint{1.807752in}{9.947778in}}{\pgfqpoint{1.818082in}{9.946350in}}%
\pgfpathquadraticcurveto{\pgfqpoint{1.828442in}{9.944952in}}{\pgfqpoint{1.840595in}{9.941914in}}%
\pgfpathlineto{\pgfqpoint{1.840595in}{9.941914in}}%
\pgfpathclose%
\pgfpathmoveto{\pgfqpoint{1.925543in}{9.980651in}}%
\pgfpathlineto{\pgfqpoint{2.056185in}{9.980651in}}%
\pgfpathlineto{\pgfqpoint{2.056185in}{9.953003in}}%
\pgfpathlineto{\pgfqpoint{2.009184in}{9.953003in}}%
\pgfpathlineto{\pgfqpoint{2.009184in}{9.838889in}}%
\pgfpathlineto{\pgfqpoint{1.972635in}{9.838889in}}%
\pgfpathlineto{\pgfqpoint{1.972635in}{9.953003in}}%
\pgfpathlineto{\pgfqpoint{1.925543in}{9.953003in}}%
\pgfpathlineto{\pgfqpoint{1.925543in}{9.980651in}}%
\pgfpathlineto{\pgfqpoint{1.925543in}{9.980651in}}%
\pgfpathclose%
\pgfpathmoveto{\pgfqpoint{2.063051in}{9.980651in}}%
\pgfpathlineto{\pgfqpoint{2.098082in}{9.980651in}}%
\pgfpathlineto{\pgfqpoint{2.122570in}{9.877626in}}%
\pgfpathlineto{\pgfqpoint{2.146875in}{9.980651in}}%
\pgfpathlineto{\pgfqpoint{2.182088in}{9.980651in}}%
\pgfpathlineto{\pgfqpoint{2.206393in}{9.877626in}}%
\pgfpathlineto{\pgfqpoint{2.230912in}{9.980651in}}%
\pgfpathlineto{\pgfqpoint{2.265638in}{9.980651in}}%
\pgfpathlineto{\pgfqpoint{2.232218in}{9.838889in}}%
\pgfpathlineto{\pgfqpoint{2.190078in}{9.838889in}}%
\pgfpathlineto{\pgfqpoint{2.164345in}{9.946654in}}%
\pgfpathlineto{\pgfqpoint{2.138915in}{9.838889in}}%
\pgfpathlineto{\pgfqpoint{2.096745in}{9.838889in}}%
\pgfpathlineto{\pgfqpoint{2.063051in}{9.980651in}}%
\pgfpathlineto{\pgfqpoint{2.063051in}{9.980651in}}%
\pgfpathclose%
\pgfpathmoveto{\pgfqpoint{2.289579in}{9.980651in}}%
\pgfpathlineto{\pgfqpoint{2.388199in}{9.980651in}}%
\pgfpathlineto{\pgfqpoint{2.388199in}{9.953003in}}%
\pgfpathlineto{\pgfqpoint{2.326129in}{9.953003in}}%
\pgfpathlineto{\pgfqpoint{2.326129in}{9.926632in}}%
\pgfpathlineto{\pgfqpoint{2.384523in}{9.926632in}}%
\pgfpathlineto{\pgfqpoint{2.384523in}{9.898984in}}%
\pgfpathlineto{\pgfqpoint{2.326129in}{9.898984in}}%
\pgfpathlineto{\pgfqpoint{2.326129in}{9.838889in}}%
\pgfpathlineto{\pgfqpoint{2.289579in}{9.838889in}}%
\pgfpathlineto{\pgfqpoint{2.289579in}{9.980651in}}%
\pgfpathlineto{\pgfqpoint{2.289579in}{9.980651in}}%
\pgfpathclose%
\pgfpathmoveto{\pgfqpoint{2.422409in}{9.980651in}}%
\pgfpathlineto{\pgfqpoint{2.521029in}{9.980651in}}%
\pgfpathlineto{\pgfqpoint{2.521029in}{9.953003in}}%
\pgfpathlineto{\pgfqpoint{2.458959in}{9.953003in}}%
\pgfpathlineto{\pgfqpoint{2.458959in}{9.926632in}}%
\pgfpathlineto{\pgfqpoint{2.517353in}{9.926632in}}%
\pgfpathlineto{\pgfqpoint{2.517353in}{9.898984in}}%
\pgfpathlineto{\pgfqpoint{2.458959in}{9.898984in}}%
\pgfpathlineto{\pgfqpoint{2.458959in}{9.866506in}}%
\pgfpathlineto{\pgfqpoint{2.523125in}{9.866506in}}%
\pgfpathlineto{\pgfqpoint{2.523125in}{9.838889in}}%
\pgfpathlineto{\pgfqpoint{2.422409in}{9.838889in}}%
\pgfpathlineto{\pgfqpoint{2.422409in}{9.980651in}}%
\pgfpathlineto{\pgfqpoint{2.422409in}{9.980651in}}%
\pgfpathclose%
\pgfusepath{fill}%
\end{pgfscope}%
\begin{pgfscope}%
\pgfpathrectangle{\pgfqpoint{1.109909in}{4.050000in}}{\pgfqpoint{5.400000in}{5.400000in}}%
\pgfusepath{clip}%
\pgfsetbuttcap%
\pgfsetroundjoin%
\definecolor{currentfill}{rgb}{0.541176,0.380392,0.658824}%
\pgfsetfillcolor{currentfill}%
\pgfsetfillopacity{0.900000}%
\pgfsetlinewidth{0.803000pt}%
\definecolor{currentstroke}{rgb}{0.541176,0.380392,0.658824}%
\pgfsetstrokecolor{currentstroke}%
\pgfsetstrokeopacity{0.900000}%
\pgfsetdash{}{0pt}%
\pgfsys@defobject{currentmarker}{\pgfqpoint{-0.038665in}{-0.038665in}}{\pgfqpoint{0.038665in}{0.038665in}}{%
\pgfpathmoveto{\pgfqpoint{0.000000in}{-0.038665in}}%
\pgfpathcurveto{\pgfqpoint{0.010254in}{-0.038665in}}{\pgfqpoint{0.020090in}{-0.034591in}}{\pgfqpoint{0.027340in}{-0.027340in}}%
\pgfpathcurveto{\pgfqpoint{0.034591in}{-0.020090in}}{\pgfqpoint{0.038665in}{-0.010254in}}{\pgfqpoint{0.038665in}{0.000000in}}%
\pgfpathcurveto{\pgfqpoint{0.038665in}{0.010254in}}{\pgfqpoint{0.034591in}{0.020090in}}{\pgfqpoint{0.027340in}{0.027340in}}%
\pgfpathcurveto{\pgfqpoint{0.020090in}{0.034591in}}{\pgfqpoint{0.010254in}{0.038665in}}{\pgfqpoint{0.000000in}{0.038665in}}%
\pgfpathcurveto{\pgfqpoint{-0.010254in}{0.038665in}}{\pgfqpoint{-0.020090in}{0.034591in}}{\pgfqpoint{-0.027340in}{0.027340in}}%
\pgfpathcurveto{\pgfqpoint{-0.034591in}{0.020090in}}{\pgfqpoint{-0.038665in}{0.010254in}}{\pgfqpoint{-0.038665in}{0.000000in}}%
\pgfpathcurveto{\pgfqpoint{-0.038665in}{-0.010254in}}{\pgfqpoint{-0.034591in}{-0.020090in}}{\pgfqpoint{-0.027340in}{-0.027340in}}%
\pgfpathcurveto{\pgfqpoint{-0.020090in}{-0.034591in}}{\pgfqpoint{-0.010254in}{-0.038665in}}{\pgfqpoint{0.000000in}{-0.038665in}}%
\pgfpathlineto{\pgfqpoint{0.000000in}{-0.038665in}}%
\pgfpathclose%
\pgfusepath{stroke,fill}%
}%
\begin{pgfscope}%
\pgfsys@transformshift{5.485502in}{6.858186in}%
\pgfsys@useobject{currentmarker}{}%
\end{pgfscope}%
\end{pgfscope}%
\begin{pgfscope}%
\pgfpathrectangle{\pgfqpoint{1.109909in}{4.050000in}}{\pgfqpoint{5.400000in}{5.400000in}}%
\pgfusepath{clip}%
\pgfsetbuttcap%
\pgfsetroundjoin%
\definecolor{currentfill}{rgb}{0.321569,0.321569,0.321569}%
\pgfsetfillcolor{currentfill}%
\pgfsetfillopacity{0.900000}%
\pgfsetlinewidth{0.803000pt}%
\definecolor{currentstroke}{rgb}{0.321569,0.321569,0.321569}%
\pgfsetstrokecolor{currentstroke}%
\pgfsetstrokeopacity{0.900000}%
\pgfsetdash{}{0pt}%
\pgfsys@defobject{currentmarker}{\pgfqpoint{-0.038665in}{-0.038665in}}{\pgfqpoint{0.038665in}{0.038665in}}{%
\pgfpathmoveto{\pgfqpoint{0.000000in}{-0.038665in}}%
\pgfpathcurveto{\pgfqpoint{0.010254in}{-0.038665in}}{\pgfqpoint{0.020090in}{-0.034591in}}{\pgfqpoint{0.027340in}{-0.027340in}}%
\pgfpathcurveto{\pgfqpoint{0.034591in}{-0.020090in}}{\pgfqpoint{0.038665in}{-0.010254in}}{\pgfqpoint{0.038665in}{0.000000in}}%
\pgfpathcurveto{\pgfqpoint{0.038665in}{0.010254in}}{\pgfqpoint{0.034591in}{0.020090in}}{\pgfqpoint{0.027340in}{0.027340in}}%
\pgfpathcurveto{\pgfqpoint{0.020090in}{0.034591in}}{\pgfqpoint{0.010254in}{0.038665in}}{\pgfqpoint{0.000000in}{0.038665in}}%
\pgfpathcurveto{\pgfqpoint{-0.010254in}{0.038665in}}{\pgfqpoint{-0.020090in}{0.034591in}}{\pgfqpoint{-0.027340in}{0.027340in}}%
\pgfpathcurveto{\pgfqpoint{-0.034591in}{0.020090in}}{\pgfqpoint{-0.038665in}{0.010254in}}{\pgfqpoint{-0.038665in}{0.000000in}}%
\pgfpathcurveto{\pgfqpoint{-0.038665in}{-0.010254in}}{\pgfqpoint{-0.034591in}{-0.020090in}}{\pgfqpoint{-0.027340in}{-0.027340in}}%
\pgfpathcurveto{\pgfqpoint{-0.020090in}{-0.034591in}}{\pgfqpoint{-0.010254in}{-0.038665in}}{\pgfqpoint{0.000000in}{-0.038665in}}%
\pgfpathlineto{\pgfqpoint{0.000000in}{-0.038665in}}%
\pgfpathclose%
\pgfusepath{stroke,fill}%
}%
\begin{pgfscope}%
\pgfsys@transformshift{2.715439in}{6.567949in}%
\pgfsys@useobject{currentmarker}{}%
\end{pgfscope}%
\end{pgfscope}%
\begin{pgfscope}%
\pgfsetbuttcap%
\pgfsetroundjoin%
\definecolor{currentfill}{rgb}{0.321569,0.321569,0.321569}%
\pgfsetfillcolor{currentfill}%
\pgfsetfillopacity{0.900000}%
\pgfsetlinewidth{0.803000pt}%
\definecolor{currentstroke}{rgb}{0.321569,0.321569,0.321569}%
\pgfsetstrokecolor{currentstroke}%
\pgfsetstrokeopacity{0.900000}%
\pgfsetdash{}{0pt}%
\pgfsys@defobject{currentmarker}{\pgfqpoint{-0.059738in}{-0.059738in}}{\pgfqpoint{0.059738in}{0.059738in}}{%
\pgfpathmoveto{\pgfqpoint{0.000000in}{-0.059738in}}%
\pgfpathlineto{\pgfqpoint{0.059738in}{0.000000in}}%
\pgfpathlineto{\pgfqpoint{0.000000in}{0.059738in}}%
\pgfpathlineto{\pgfqpoint{-0.059738in}{0.000000in}}%
\pgfpathlineto{\pgfqpoint{0.000000in}{-0.059738in}}%
\pgfpathclose%
\pgfusepath{stroke,fill}%
}%
\begin{pgfscope}%
\pgfsys@transformshift{1.109909in}{4.513452in}%
\pgfsys@useobject{currentmarker}{}%
\end{pgfscope}%
\begin{pgfscope}%
\pgfsys@transformshift{1.109909in}{4.560820in}%
\pgfsys@useobject{currentmarker}{}%
\end{pgfscope}%
\begin{pgfscope}%
\pgfsys@transformshift{1.109909in}{4.056751in}%
\pgfsys@useobject{currentmarker}{}%
\end{pgfscope}%
\end{pgfscope}%
\begin{pgfscope}%
\pgfsetbuttcap%
\pgfsetroundjoin%
\definecolor{currentfill}{rgb}{0.768627,0.603922,0.000000}%
\pgfsetfillcolor{currentfill}%
\pgfsetfillopacity{0.900000}%
\pgfsetlinewidth{0.803000pt}%
\definecolor{currentstroke}{rgb}{0.768627,0.603922,0.000000}%
\pgfsetstrokecolor{currentstroke}%
\pgfsetstrokeopacity{0.900000}%
\pgfsetdash{}{0pt}%
\pgfsys@defobject{currentmarker}{\pgfqpoint{-0.059738in}{-0.059738in}}{\pgfqpoint{0.059738in}{0.059738in}}{%
\pgfpathmoveto{\pgfqpoint{0.000000in}{-0.059738in}}%
\pgfpathlineto{\pgfqpoint{0.059738in}{0.000000in}}%
\pgfpathlineto{\pgfqpoint{0.000000in}{0.059738in}}%
\pgfpathlineto{\pgfqpoint{-0.059738in}{0.000000in}}%
\pgfpathlineto{\pgfqpoint{0.000000in}{-0.059738in}}%
\pgfpathclose%
\pgfusepath{stroke,fill}%
}%
\begin{pgfscope}%
\pgfsys@transformshift{1.109909in}{6.282727in}%
\pgfsys@useobject{currentmarker}{}%
\end{pgfscope}%
\end{pgfscope}%
\begin{pgfscope}%
\pgfpathrectangle{\pgfqpoint{1.109909in}{4.050000in}}{\pgfqpoint{5.400000in}{5.400000in}}%
\pgfusepath{clip}%
\pgfsetbuttcap%
\pgfsetroundjoin%
\definecolor{currentfill}{rgb}{0.176471,0.713725,0.643137}%
\pgfsetfillcolor{currentfill}%
\pgfsetfillopacity{0.900000}%
\pgfsetlinewidth{0.803000pt}%
\definecolor{currentstroke}{rgb}{0.176471,0.713725,0.643137}%
\pgfsetstrokecolor{currentstroke}%
\pgfsetstrokeopacity{0.900000}%
\pgfsetdash{}{0pt}%
\pgfsys@defobject{currentmarker}{\pgfqpoint{-0.038665in}{-0.038665in}}{\pgfqpoint{0.038665in}{0.038665in}}{%
\pgfpathmoveto{\pgfqpoint{0.000000in}{-0.038665in}}%
\pgfpathcurveto{\pgfqpoint{0.010254in}{-0.038665in}}{\pgfqpoint{0.020090in}{-0.034591in}}{\pgfqpoint{0.027340in}{-0.027340in}}%
\pgfpathcurveto{\pgfqpoint{0.034591in}{-0.020090in}}{\pgfqpoint{0.038665in}{-0.010254in}}{\pgfqpoint{0.038665in}{0.000000in}}%
\pgfpathcurveto{\pgfqpoint{0.038665in}{0.010254in}}{\pgfqpoint{0.034591in}{0.020090in}}{\pgfqpoint{0.027340in}{0.027340in}}%
\pgfpathcurveto{\pgfqpoint{0.020090in}{0.034591in}}{\pgfqpoint{0.010254in}{0.038665in}}{\pgfqpoint{0.000000in}{0.038665in}}%
\pgfpathcurveto{\pgfqpoint{-0.010254in}{0.038665in}}{\pgfqpoint{-0.020090in}{0.034591in}}{\pgfqpoint{-0.027340in}{0.027340in}}%
\pgfpathcurveto{\pgfqpoint{-0.034591in}{0.020090in}}{\pgfqpoint{-0.038665in}{0.010254in}}{\pgfqpoint{-0.038665in}{0.000000in}}%
\pgfpathcurveto{\pgfqpoint{-0.038665in}{-0.010254in}}{\pgfqpoint{-0.034591in}{-0.020090in}}{\pgfqpoint{-0.027340in}{-0.027340in}}%
\pgfpathcurveto{\pgfqpoint{-0.020090in}{-0.034591in}}{\pgfqpoint{-0.010254in}{-0.038665in}}{\pgfqpoint{0.000000in}{-0.038665in}}%
\pgfpathlineto{\pgfqpoint{0.000000in}{-0.038665in}}%
\pgfpathclose%
\pgfusepath{stroke,fill}%
}%
\begin{pgfscope}%
\pgfsys@transformshift{4.553259in}{9.244124in}%
\pgfsys@useobject{currentmarker}{}%
\end{pgfscope}%
\begin{pgfscope}%
\pgfsys@transformshift{5.668580in}{6.179358in}%
\pgfsys@useobject{currentmarker}{}%
\end{pgfscope}%
\end{pgfscope}%
\begin{pgfscope}%
\pgfsetbuttcap%
\pgfsetroundjoin%
\definecolor{currentfill}{rgb}{0.000000,0.447059,0.698039}%
\pgfsetfillcolor{currentfill}%
\pgfsetfillopacity{0.900000}%
\pgfsetlinewidth{0.803000pt}%
\definecolor{currentstroke}{rgb}{0.000000,0.447059,0.698039}%
\pgfsetstrokecolor{currentstroke}%
\pgfsetstrokeopacity{0.900000}%
\pgfsetdash{}{0pt}%
\pgfsys@defobject{currentmarker}{\pgfqpoint{-0.059738in}{-0.059738in}}{\pgfqpoint{0.059738in}{0.059738in}}{%
\pgfpathmoveto{\pgfqpoint{0.000000in}{-0.059738in}}%
\pgfpathlineto{\pgfqpoint{0.059738in}{0.000000in}}%
\pgfpathlineto{\pgfqpoint{0.000000in}{0.059738in}}%
\pgfpathlineto{\pgfqpoint{-0.059738in}{0.000000in}}%
\pgfpathlineto{\pgfqpoint{0.000000in}{-0.059738in}}%
\pgfpathclose%
\pgfusepath{stroke,fill}%
}%
\begin{pgfscope}%
\pgfsys@transformshift{3.828074in}{9.450000in}%
\pgfsys@useobject{currentmarker}{}%
\end{pgfscope}%
\begin{pgfscope}%
\pgfsys@transformshift{2.335642in}{4.050000in}%
\pgfsys@useobject{currentmarker}{}%
\end{pgfscope}%
\begin{pgfscope}%
\pgfsys@transformshift{2.356385in}{4.050000in}%
\pgfsys@useobject{currentmarker}{}%
\end{pgfscope}%
\begin{pgfscope}%
\pgfsys@transformshift{2.793156in}{4.050000in}%
\pgfsys@useobject{currentmarker}{}%
\end{pgfscope}%
\begin{pgfscope}%
\pgfsys@transformshift{3.746914in}{4.050000in}%
\pgfsys@useobject{currentmarker}{}%
\end{pgfscope}%
\begin{pgfscope}%
\pgfsys@transformshift{2.853497in}{4.050000in}%
\pgfsys@useobject{currentmarker}{}%
\end{pgfscope}%
\end{pgfscope}%
\begin{pgfscope}%
\pgfsetbuttcap%
\pgfsetroundjoin%
\definecolor{currentfill}{rgb}{0.901961,0.196078,0.156863}%
\pgfsetfillcolor{currentfill}%
\pgfsetfillopacity{0.900000}%
\pgfsetlinewidth{0.803000pt}%
\definecolor{currentstroke}{rgb}{0.901961,0.196078,0.156863}%
\pgfsetstrokecolor{currentstroke}%
\pgfsetstrokeopacity{0.900000}%
\pgfsetdash{}{0pt}%
\pgfsys@defobject{currentmarker}{\pgfqpoint{-0.059738in}{-0.059738in}}{\pgfqpoint{0.059738in}{0.059738in}}{%
\pgfpathmoveto{\pgfqpoint{0.000000in}{-0.059738in}}%
\pgfpathlineto{\pgfqpoint{0.059738in}{0.000000in}}%
\pgfpathlineto{\pgfqpoint{0.000000in}{0.059738in}}%
\pgfpathlineto{\pgfqpoint{-0.059738in}{0.000000in}}%
\pgfpathlineto{\pgfqpoint{0.000000in}{-0.059738in}}%
\pgfpathclose%
\pgfusepath{stroke,fill}%
}%
\begin{pgfscope}%
\pgfsys@transformshift{6.509909in}{9.450000in}%
\pgfsys@useobject{currentmarker}{}%
\end{pgfscope}%
\end{pgfscope}%
\begin{pgfscope}%
\pgfpathrectangle{\pgfqpoint{1.109909in}{4.050000in}}{\pgfqpoint{5.400000in}{5.400000in}}%
\pgfusepath{clip}%
\pgfsetbuttcap%
\pgfsetroundjoin%
\definecolor{currentfill}{rgb}{0.901961,0.403922,0.619608}%
\pgfsetfillcolor{currentfill}%
\pgfsetfillopacity{0.900000}%
\pgfsetlinewidth{0.803000pt}%
\definecolor{currentstroke}{rgb}{0.901961,0.403922,0.619608}%
\pgfsetstrokecolor{currentstroke}%
\pgfsetstrokeopacity{0.900000}%
\pgfsetdash{}{0pt}%
\pgfsys@defobject{currentmarker}{\pgfqpoint{-0.038665in}{-0.038665in}}{\pgfqpoint{0.038665in}{0.038665in}}{%
\pgfpathmoveto{\pgfqpoint{0.000000in}{-0.038665in}}%
\pgfpathcurveto{\pgfqpoint{0.010254in}{-0.038665in}}{\pgfqpoint{0.020090in}{-0.034591in}}{\pgfqpoint{0.027340in}{-0.027340in}}%
\pgfpathcurveto{\pgfqpoint{0.034591in}{-0.020090in}}{\pgfqpoint{0.038665in}{-0.010254in}}{\pgfqpoint{0.038665in}{0.000000in}}%
\pgfpathcurveto{\pgfqpoint{0.038665in}{0.010254in}}{\pgfqpoint{0.034591in}{0.020090in}}{\pgfqpoint{0.027340in}{0.027340in}}%
\pgfpathcurveto{\pgfqpoint{0.020090in}{0.034591in}}{\pgfqpoint{0.010254in}{0.038665in}}{\pgfqpoint{0.000000in}{0.038665in}}%
\pgfpathcurveto{\pgfqpoint{-0.010254in}{0.038665in}}{\pgfqpoint{-0.020090in}{0.034591in}}{\pgfqpoint{-0.027340in}{0.027340in}}%
\pgfpathcurveto{\pgfqpoint{-0.034591in}{0.020090in}}{\pgfqpoint{-0.038665in}{0.010254in}}{\pgfqpoint{-0.038665in}{0.000000in}}%
\pgfpathcurveto{\pgfqpoint{-0.038665in}{-0.010254in}}{\pgfqpoint{-0.034591in}{-0.020090in}}{\pgfqpoint{-0.027340in}{-0.027340in}}%
\pgfpathcurveto{\pgfqpoint{-0.020090in}{-0.034591in}}{\pgfqpoint{-0.010254in}{-0.038665in}}{\pgfqpoint{0.000000in}{-0.038665in}}%
\pgfpathlineto{\pgfqpoint{0.000000in}{-0.038665in}}%
\pgfpathclose%
\pgfusepath{stroke,fill}%
}%
\begin{pgfscope}%
\pgfsys@transformshift{5.006097in}{6.484513in}%
\pgfsys@useobject{currentmarker}{}%
\end{pgfscope}%
\end{pgfscope}%
\begin{pgfscope}%
\pgfsetbuttcap%
\pgfsetroundjoin%
\definecolor{currentfill}{rgb}{0.549020,0.160784,0.505882}%
\pgfsetfillcolor{currentfill}%
\pgfsetfillopacity{0.900000}%
\pgfsetlinewidth{0.803000pt}%
\definecolor{currentstroke}{rgb}{0.549020,0.160784,0.505882}%
\pgfsetstrokecolor{currentstroke}%
\pgfsetstrokeopacity{0.900000}%
\pgfsetdash{}{0pt}%
\pgfsys@defobject{currentmarker}{\pgfqpoint{-0.059738in}{-0.059738in}}{\pgfqpoint{0.059738in}{0.059738in}}{%
\pgfpathmoveto{\pgfqpoint{0.000000in}{-0.059738in}}%
\pgfpathlineto{\pgfqpoint{0.059738in}{0.000000in}}%
\pgfpathlineto{\pgfqpoint{0.000000in}{0.059738in}}%
\pgfpathlineto{\pgfqpoint{-0.059738in}{0.000000in}}%
\pgfpathlineto{\pgfqpoint{0.000000in}{-0.059738in}}%
\pgfpathclose%
\pgfusepath{stroke,fill}%
}%
\begin{pgfscope}%
\pgfsys@transformshift{6.509909in}{9.450000in}%
\pgfsys@useobject{currentmarker}{}%
\end{pgfscope}%
\begin{pgfscope}%
\pgfsys@transformshift{6.509909in}{9.450000in}%
\pgfsys@useobject{currentmarker}{}%
\end{pgfscope}%
\end{pgfscope}%
\begin{pgfscope}%
\pgfpathrectangle{\pgfqpoint{1.109909in}{4.050000in}}{\pgfqpoint{5.400000in}{5.400000in}}%
\pgfusepath{clip}%
\pgfsetbuttcap%
\pgfsetroundjoin%
\definecolor{currentfill}{rgb}{0.000000,0.560784,0.611765}%
\pgfsetfillcolor{currentfill}%
\pgfsetfillopacity{0.900000}%
\pgfsetlinewidth{0.803000pt}%
\definecolor{currentstroke}{rgb}{0.000000,0.560784,0.611765}%
\pgfsetstrokecolor{currentstroke}%
\pgfsetstrokeopacity{0.900000}%
\pgfsetdash{}{0pt}%
\pgfsys@defobject{currentmarker}{\pgfqpoint{-0.038665in}{-0.038665in}}{\pgfqpoint{0.038665in}{0.038665in}}{%
\pgfpathmoveto{\pgfqpoint{0.000000in}{-0.038665in}}%
\pgfpathcurveto{\pgfqpoint{0.010254in}{-0.038665in}}{\pgfqpoint{0.020090in}{-0.034591in}}{\pgfqpoint{0.027340in}{-0.027340in}}%
\pgfpathcurveto{\pgfqpoint{0.034591in}{-0.020090in}}{\pgfqpoint{0.038665in}{-0.010254in}}{\pgfqpoint{0.038665in}{0.000000in}}%
\pgfpathcurveto{\pgfqpoint{0.038665in}{0.010254in}}{\pgfqpoint{0.034591in}{0.020090in}}{\pgfqpoint{0.027340in}{0.027340in}}%
\pgfpathcurveto{\pgfqpoint{0.020090in}{0.034591in}}{\pgfqpoint{0.010254in}{0.038665in}}{\pgfqpoint{0.000000in}{0.038665in}}%
\pgfpathcurveto{\pgfqpoint{-0.010254in}{0.038665in}}{\pgfqpoint{-0.020090in}{0.034591in}}{\pgfqpoint{-0.027340in}{0.027340in}}%
\pgfpathcurveto{\pgfqpoint{-0.034591in}{0.020090in}}{\pgfqpoint{-0.038665in}{0.010254in}}{\pgfqpoint{-0.038665in}{0.000000in}}%
\pgfpathcurveto{\pgfqpoint{-0.038665in}{-0.010254in}}{\pgfqpoint{-0.034591in}{-0.020090in}}{\pgfqpoint{-0.027340in}{-0.027340in}}%
\pgfpathcurveto{\pgfqpoint{-0.020090in}{-0.034591in}}{\pgfqpoint{-0.010254in}{-0.038665in}}{\pgfqpoint{0.000000in}{-0.038665in}}%
\pgfpathlineto{\pgfqpoint{0.000000in}{-0.038665in}}%
\pgfpathclose%
\pgfusepath{stroke,fill}%
}%
\begin{pgfscope}%
\pgfsys@transformshift{1.789047in}{6.175558in}%
\pgfsys@useobject{currentmarker}{}%
\end{pgfscope}%
\end{pgfscope}%
\begin{pgfscope}%
\pgfsetbuttcap%
\pgfsetroundjoin%
\definecolor{currentfill}{rgb}{0.000000,0.560784,0.611765}%
\pgfsetfillcolor{currentfill}%
\pgfsetfillopacity{0.900000}%
\pgfsetlinewidth{0.803000pt}%
\definecolor{currentstroke}{rgb}{0.000000,0.560784,0.611765}%
\pgfsetstrokecolor{currentstroke}%
\pgfsetstrokeopacity{0.900000}%
\pgfsetdash{}{0pt}%
\pgfsys@defobject{currentmarker}{\pgfqpoint{-0.059738in}{-0.059738in}}{\pgfqpoint{0.059738in}{0.059738in}}{%
\pgfpathmoveto{\pgfqpoint{0.000000in}{-0.059738in}}%
\pgfpathlineto{\pgfqpoint{0.059738in}{0.000000in}}%
\pgfpathlineto{\pgfqpoint{0.000000in}{0.059738in}}%
\pgfpathlineto{\pgfqpoint{-0.059738in}{0.000000in}}%
\pgfpathlineto{\pgfqpoint{0.000000in}{-0.059738in}}%
\pgfpathclose%
\pgfusepath{stroke,fill}%
}%
\begin{pgfscope}%
\pgfsys@transformshift{1.109909in}{4.270378in}%
\pgfsys@useobject{currentmarker}{}%
\end{pgfscope}%
\end{pgfscope}%
\begin{pgfscope}%
\pgfsetbuttcap%
\pgfsetroundjoin%
\pgfsetlinewidth{0.903375pt}%
\definecolor{currentstroke}{rgb}{0.321569,0.321569,0.321569}%
\pgfsetstrokecolor{currentstroke}%
\pgfsetdash{}{0pt}%
\pgfpathmoveto{\pgfqpoint{1.109909in}{3.977166in}}%
\pgfpathlineto{\pgfqpoint{1.182743in}{4.050000in}}%
\pgfpathlineto{\pgfqpoint{1.109909in}{4.122834in}}%
\pgfpathlineto{\pgfqpoint{1.037075in}{4.050000in}}%
\pgfpathlineto{\pgfqpoint{1.109909in}{3.977166in}}%
\pgfpathclose%
\pgfusepath{stroke}%
\end{pgfscope}%
\begin{pgfscope}%
\pgfsetbuttcap%
\pgfsetroundjoin%
\pgfsetlinewidth{0.903375pt}%
\definecolor{currentstroke}{rgb}{0.321569,0.321569,0.321569}%
\pgfsetstrokecolor{currentstroke}%
\pgfsetdash{}{0pt}%
\pgfpathmoveto{\pgfqpoint{1.109909in}{4.541182in}}%
\pgfpathlineto{\pgfqpoint{1.182743in}{4.614015in}}%
\pgfpathlineto{\pgfqpoint{1.109909in}{4.686849in}}%
\pgfpathlineto{\pgfqpoint{1.037075in}{4.614015in}}%
\pgfpathlineto{\pgfqpoint{1.109909in}{4.541182in}}%
\pgfpathclose%
\pgfusepath{stroke}%
\end{pgfscope}%
\begin{pgfscope}%
\pgfsetbuttcap%
\pgfsetroundjoin%
\pgfsetlinewidth{0.903375pt}%
\definecolor{currentstroke}{rgb}{0.321569,0.321569,0.321569}%
\pgfsetstrokecolor{currentstroke}%
\pgfsetdash{}{0pt}%
\pgfpathmoveto{\pgfqpoint{1.109909in}{3.977166in}}%
\pgfpathlineto{\pgfqpoint{1.182743in}{4.050000in}}%
\pgfpathlineto{\pgfqpoint{1.109909in}{4.122834in}}%
\pgfpathlineto{\pgfqpoint{1.037075in}{4.050000in}}%
\pgfpathlineto{\pgfqpoint{1.109909in}{3.977166in}}%
\pgfpathclose%
\pgfusepath{stroke}%
\end{pgfscope}%
\begin{pgfscope}%
\pgfpathrectangle{\pgfqpoint{1.109909in}{4.050000in}}{\pgfqpoint{5.400000in}{5.400000in}}%
\pgfusepath{clip}%
\pgfsetbuttcap%
\pgfsetroundjoin%
\pgfsetlinewidth{0.803000pt}%
\definecolor{currentstroke}{rgb}{0.768627,0.603922,0.000000}%
\pgfsetstrokecolor{currentstroke}%
\pgfsetdash{}{0pt}%
\pgfpathmoveto{\pgfqpoint{5.277397in}{6.952231in}}%
\pgfpathlineto{\pgfqpoint{5.370567in}{6.952231in}}%
\pgfpathlineto{\pgfqpoint{5.370567in}{7.045401in}}%
\pgfpathlineto{\pgfqpoint{5.277397in}{7.045401in}}%
\pgfpathlineto{\pgfqpoint{5.277397in}{6.952231in}}%
\pgfpathclose%
\pgfusepath{stroke}%
\end{pgfscope}%
\begin{pgfscope}%
\pgfpathrectangle{\pgfqpoint{1.109909in}{4.050000in}}{\pgfqpoint{5.400000in}{5.400000in}}%
\pgfusepath{clip}%
\pgfsetbuttcap%
\pgfsetroundjoin%
\pgfsetlinewidth{0.803000pt}%
\definecolor{currentstroke}{rgb}{0.768627,0.603922,0.000000}%
\pgfsetstrokecolor{currentstroke}%
\pgfsetdash{}{0pt}%
\pgfpathmoveto{\pgfqpoint{5.306915in}{6.863795in}}%
\pgfpathlineto{\pgfqpoint{5.400085in}{6.863795in}}%
\pgfpathlineto{\pgfqpoint{5.400085in}{6.956964in}}%
\pgfpathlineto{\pgfqpoint{5.306915in}{6.956964in}}%
\pgfpathlineto{\pgfqpoint{5.306915in}{6.863795in}}%
\pgfpathclose%
\pgfusepath{stroke}%
\end{pgfscope}%
\begin{pgfscope}%
\pgfpathrectangle{\pgfqpoint{1.109909in}{4.050000in}}{\pgfqpoint{5.400000in}{5.400000in}}%
\pgfusepath{clip}%
\pgfsetbuttcap%
\pgfsetroundjoin%
\pgfsetlinewidth{0.803000pt}%
\definecolor{currentstroke}{rgb}{0.768627,0.603922,0.000000}%
\pgfsetstrokecolor{currentstroke}%
\pgfsetdash{}{0pt}%
\pgfpathmoveto{\pgfqpoint{5.239558in}{7.443911in}}%
\pgfpathlineto{\pgfqpoint{5.332728in}{7.443911in}}%
\pgfpathlineto{\pgfqpoint{5.332728in}{7.537080in}}%
\pgfpathlineto{\pgfqpoint{5.239558in}{7.537080in}}%
\pgfpathlineto{\pgfqpoint{5.239558in}{7.443911in}}%
\pgfpathclose%
\pgfusepath{stroke}%
\end{pgfscope}%
\begin{pgfscope}%
\pgfpathrectangle{\pgfqpoint{1.109909in}{4.050000in}}{\pgfqpoint{5.400000in}{5.400000in}}%
\pgfusepath{clip}%
\pgfsetbuttcap%
\pgfsetroundjoin%
\pgfsetlinewidth{0.803000pt}%
\definecolor{currentstroke}{rgb}{0.768627,0.603922,0.000000}%
\pgfsetstrokecolor{currentstroke}%
\pgfsetdash{}{0pt}%
\pgfpathmoveto{\pgfqpoint{4.993053in}{6.596691in}}%
\pgfpathlineto{\pgfqpoint{5.086223in}{6.596691in}}%
\pgfpathlineto{\pgfqpoint{5.086223in}{6.689861in}}%
\pgfpathlineto{\pgfqpoint{4.993053in}{6.689861in}}%
\pgfpathlineto{\pgfqpoint{4.993053in}{6.596691in}}%
\pgfpathclose%
\pgfusepath{stroke}%
\end{pgfscope}%
\begin{pgfscope}%
\pgfpathrectangle{\pgfqpoint{1.109909in}{4.050000in}}{\pgfqpoint{5.400000in}{5.400000in}}%
\pgfusepath{clip}%
\pgfsetbuttcap%
\pgfsetroundjoin%
\pgfsetlinewidth{0.803000pt}%
\definecolor{currentstroke}{rgb}{0.768627,0.603922,0.000000}%
\pgfsetstrokecolor{currentstroke}%
\pgfsetdash{}{0pt}%
\pgfpathmoveto{\pgfqpoint{3.523162in}{6.322828in}}%
\pgfpathlineto{\pgfqpoint{3.616332in}{6.322828in}}%
\pgfpathlineto{\pgfqpoint{3.616332in}{6.415998in}}%
\pgfpathlineto{\pgfqpoint{3.523162in}{6.415998in}}%
\pgfpathlineto{\pgfqpoint{3.523162in}{6.322828in}}%
\pgfpathclose%
\pgfusepath{stroke}%
\end{pgfscope}%
\begin{pgfscope}%
\pgfpathrectangle{\pgfqpoint{1.109909in}{4.050000in}}{\pgfqpoint{5.400000in}{5.400000in}}%
\pgfusepath{clip}%
\pgfsetbuttcap%
\pgfsetroundjoin%
\pgfsetlinewidth{0.803000pt}%
\definecolor{currentstroke}{rgb}{0.768627,0.603922,0.000000}%
\pgfsetstrokecolor{currentstroke}%
\pgfsetdash{}{0pt}%
\pgfpathmoveto{\pgfqpoint{3.397441in}{6.147879in}}%
\pgfpathlineto{\pgfqpoint{3.490610in}{6.147879in}}%
\pgfpathlineto{\pgfqpoint{3.490610in}{6.241049in}}%
\pgfpathlineto{\pgfqpoint{3.397441in}{6.241049in}}%
\pgfpathlineto{\pgfqpoint{3.397441in}{6.147879in}}%
\pgfpathclose%
\pgfusepath{stroke}%
\end{pgfscope}%
\begin{pgfscope}%
\pgfsetbuttcap%
\pgfsetroundjoin%
\pgfsetlinewidth{0.903375pt}%
\definecolor{currentstroke}{rgb}{0.337255,0.705882,0.913725}%
\pgfsetstrokecolor{currentstroke}%
\pgfsetdash{}{0pt}%
\pgfpathmoveto{\pgfqpoint{6.506560in}{9.377166in}}%
\pgfpathlineto{\pgfqpoint{6.579394in}{9.450000in}}%
\pgfpathlineto{\pgfqpoint{6.506560in}{9.522834in}}%
\pgfpathlineto{\pgfqpoint{6.433726in}{9.450000in}}%
\pgfpathlineto{\pgfqpoint{6.506560in}{9.377166in}}%
\pgfpathclose%
\pgfusepath{stroke}%
\end{pgfscope}%
\begin{pgfscope}%
\pgfpathrectangle{\pgfqpoint{1.109909in}{4.050000in}}{\pgfqpoint{5.400000in}{5.400000in}}%
\pgfusepath{clip}%
\pgfsetbuttcap%
\pgfsetroundjoin%
\pgfsetlinewidth{0.803000pt}%
\definecolor{currentstroke}{rgb}{0.901961,0.196078,0.156863}%
\pgfsetstrokecolor{currentstroke}%
\pgfsetdash{}{0pt}%
\pgfpathmoveto{\pgfqpoint{5.186421in}{8.651910in}}%
\pgfpathlineto{\pgfqpoint{5.279590in}{8.651910in}}%
\pgfpathlineto{\pgfqpoint{5.279590in}{8.745080in}}%
\pgfpathlineto{\pgfqpoint{5.186421in}{8.745080in}}%
\pgfpathlineto{\pgfqpoint{5.186421in}{8.651910in}}%
\pgfpathclose%
\pgfusepath{stroke}%
\end{pgfscope}%
\begin{pgfscope}%
\pgfpathrectangle{\pgfqpoint{1.109909in}{4.050000in}}{\pgfqpoint{5.400000in}{5.400000in}}%
\pgfusepath{clip}%
\pgfsetbuttcap%
\pgfsetroundjoin%
\pgfsetlinewidth{0.803000pt}%
\definecolor{currentstroke}{rgb}{0.901961,0.196078,0.156863}%
\pgfsetstrokecolor{currentstroke}%
\pgfsetdash{}{0pt}%
\pgfpathmoveto{\pgfqpoint{3.969540in}{7.012826in}}%
\pgfpathlineto{\pgfqpoint{4.062709in}{7.012826in}}%
\pgfpathlineto{\pgfqpoint{4.062709in}{7.105995in}}%
\pgfpathlineto{\pgfqpoint{3.969540in}{7.105995in}}%
\pgfpathlineto{\pgfqpoint{3.969540in}{7.012826in}}%
\pgfpathclose%
\pgfusepath{stroke}%
\end{pgfscope}%
\begin{pgfscope}%
\pgfsetbuttcap%
\pgfsetroundjoin%
\pgfsetlinewidth{0.803000pt}%
\definecolor{currentstroke}{rgb}{0.901961,0.196078,0.156863}%
\pgfsetstrokecolor{currentstroke}%
\pgfsetdash{}{0pt}%
\pgfpathmoveto{\pgfqpoint{6.509909in}{9.167565in}}%
\pgfpathlineto{\pgfqpoint{6.603595in}{9.261250in}}%
\pgfpathlineto{\pgfqpoint{6.509909in}{9.354936in}}%
\pgfpathlineto{\pgfqpoint{6.416223in}{9.261250in}}%
\pgfpathlineto{\pgfqpoint{6.509909in}{9.167565in}}%
\pgfpathclose%
\pgfusepath{stroke}%
\end{pgfscope}%
\begin{pgfscope}%
\pgfsetbuttcap%
\pgfsetroundjoin%
\pgfsetlinewidth{0.803000pt}%
\definecolor{currentstroke}{rgb}{0.901961,0.196078,0.156863}%
\pgfsetstrokecolor{currentstroke}%
\pgfsetdash{}{0pt}%
\pgfpathmoveto{\pgfqpoint{6.463324in}{9.214665in}}%
\pgfpathlineto{\pgfqpoint{6.556494in}{9.214665in}}%
\pgfpathlineto{\pgfqpoint{6.556494in}{9.307835in}}%
\pgfpathlineto{\pgfqpoint{6.463324in}{9.307835in}}%
\pgfpathlineto{\pgfqpoint{6.463324in}{9.214665in}}%
\pgfpathclose%
\pgfusepath{stroke}%
\end{pgfscope}%
\begin{pgfscope}%
\pgfsetbuttcap%
\pgfsetroundjoin%
\pgfsetlinewidth{0.803000pt}%
\definecolor{currentstroke}{rgb}{0.901961,0.196078,0.156863}%
\pgfsetstrokecolor{currentstroke}%
\pgfsetdash{}{0pt}%
\pgfpathmoveto{\pgfqpoint{6.147143in}{9.356314in}}%
\pgfpathlineto{\pgfqpoint{6.240829in}{9.450000in}}%
\pgfpathlineto{\pgfqpoint{6.147143in}{9.543686in}}%
\pgfpathlineto{\pgfqpoint{6.053457in}{9.450000in}}%
\pgfpathlineto{\pgfqpoint{6.147143in}{9.356314in}}%
\pgfpathclose%
\pgfusepath{stroke}%
\end{pgfscope}%
\begin{pgfscope}%
\pgfsetbuttcap%
\pgfsetroundjoin%
\pgfsetlinewidth{0.803000pt}%
\definecolor{currentstroke}{rgb}{0.498039,0.560784,0.678431}%
\pgfsetstrokecolor{currentstroke}%
\pgfsetdash{}{0pt}%
\pgfpathmoveto{\pgfqpoint{1.109909in}{4.002032in}}%
\pgfpathlineto{\pgfqpoint{1.203595in}{4.095717in}}%
\pgfpathlineto{\pgfqpoint{1.109909in}{4.189403in}}%
\pgfpathlineto{\pgfqpoint{1.016223in}{4.095717in}}%
\pgfpathlineto{\pgfqpoint{1.109909in}{4.002032in}}%
\pgfpathclose%
\pgfusepath{stroke}%
\end{pgfscope}%
\begin{pgfscope}%
\pgfpathrectangle{\pgfqpoint{1.109909in}{4.050000in}}{\pgfqpoint{5.400000in}{5.400000in}}%
\pgfusepath{clip}%
\pgfsetbuttcap%
\pgfsetroundjoin%
\definecolor{currentfill}{rgb}{0.541176,0.380392,0.658824}%
\pgfsetfillcolor{currentfill}%
\pgfsetlinewidth{1.254687pt}%
\definecolor{currentstroke}{rgb}{0.541176,0.380392,0.658824}%
\pgfsetstrokecolor{currentstroke}%
\pgfsetdash{}{0pt}%
\pgfsys@defobject{currentmarker}{\pgfqpoint{-0.043368in}{-0.043368in}}{\pgfqpoint{0.043368in}{0.043368in}}{%
\pgfpathmoveto{\pgfqpoint{-0.043368in}{-0.043368in}}%
\pgfpathlineto{\pgfqpoint{0.043368in}{0.043368in}}%
\pgfpathmoveto{\pgfqpoint{-0.043368in}{0.043368in}}%
\pgfpathlineto{\pgfqpoint{0.043368in}{-0.043368in}}%
\pgfusepath{stroke,fill}%
}%
\begin{pgfscope}%
\pgfsys@transformshift{4.256096in}{7.893276in}%
\pgfsys@useobject{currentmarker}{}%
\end{pgfscope}%
\begin{pgfscope}%
\pgfsys@transformshift{3.156299in}{6.159509in}%
\pgfsys@useobject{currentmarker}{}%
\end{pgfscope}%
\begin{pgfscope}%
\pgfsys@transformshift{3.799135in}{6.087055in}%
\pgfsys@useobject{currentmarker}{}%
\end{pgfscope}%
\end{pgfscope}%
\begin{pgfscope}%
\pgfpathrectangle{\pgfqpoint{1.109909in}{4.050000in}}{\pgfqpoint{5.400000in}{5.400000in}}%
\pgfusepath{clip}%
\pgfsetbuttcap%
\pgfsetroundjoin%
\definecolor{currentfill}{rgb}{0.321569,0.321569,0.321569}%
\pgfsetfillcolor{currentfill}%
\pgfsetlinewidth{1.254687pt}%
\definecolor{currentstroke}{rgb}{0.321569,0.321569,0.321569}%
\pgfsetstrokecolor{currentstroke}%
\pgfsetdash{}{0pt}%
\pgfsys@defobject{currentmarker}{\pgfqpoint{-0.043368in}{-0.043368in}}{\pgfqpoint{0.043368in}{0.043368in}}{%
\pgfpathmoveto{\pgfqpoint{-0.043368in}{-0.043368in}}%
\pgfpathlineto{\pgfqpoint{0.043368in}{0.043368in}}%
\pgfpathmoveto{\pgfqpoint{-0.043368in}{0.043368in}}%
\pgfpathlineto{\pgfqpoint{0.043368in}{-0.043368in}}%
\pgfusepath{stroke,fill}%
}%
\begin{pgfscope}%
\pgfsys@transformshift{2.516781in}{4.169533in}%
\pgfsys@useobject{currentmarker}{}%
\end{pgfscope}%
\begin{pgfscope}%
\pgfsys@transformshift{2.377783in}{4.867435in}%
\pgfsys@useobject{currentmarker}{}%
\end{pgfscope}%
\begin{pgfscope}%
\pgfsys@transformshift{3.085305in}{5.994774in}%
\pgfsys@useobject{currentmarker}{}%
\end{pgfscope}%
\begin{pgfscope}%
\pgfsys@transformshift{2.372997in}{7.241779in}%
\pgfsys@useobject{currentmarker}{}%
\end{pgfscope}%
\begin{pgfscope}%
\pgfsys@transformshift{2.558051in}{6.653504in}%
\pgfsys@useobject{currentmarker}{}%
\end{pgfscope}%
\begin{pgfscope}%
\pgfsys@transformshift{2.701521in}{6.441911in}%
\pgfsys@useobject{currentmarker}{}%
\end{pgfscope}%
\begin{pgfscope}%
\pgfsys@transformshift{3.135744in}{6.548892in}%
\pgfsys@useobject{currentmarker}{}%
\end{pgfscope}%
\begin{pgfscope}%
\pgfsys@transformshift{4.824820in}{7.947941in}%
\pgfsys@useobject{currentmarker}{}%
\end{pgfscope}%
\begin{pgfscope}%
\pgfsys@transformshift{4.986314in}{7.730241in}%
\pgfsys@useobject{currentmarker}{}%
\end{pgfscope}%
\begin{pgfscope}%
\pgfsys@transformshift{5.126840in}{7.817369in}%
\pgfsys@useobject{currentmarker}{}%
\end{pgfscope}%
\begin{pgfscope}%
\pgfsys@transformshift{5.326870in}{8.098603in}%
\pgfsys@useobject{currentmarker}{}%
\end{pgfscope}%
\begin{pgfscope}%
\pgfsys@transformshift{5.000613in}{7.745674in}%
\pgfsys@useobject{currentmarker}{}%
\end{pgfscope}%
\begin{pgfscope}%
\pgfsys@transformshift{5.403789in}{7.722213in}%
\pgfsys@useobject{currentmarker}{}%
\end{pgfscope}%
\begin{pgfscope}%
\pgfsys@transformshift{5.084116in}{7.662473in}%
\pgfsys@useobject{currentmarker}{}%
\end{pgfscope}%
\begin{pgfscope}%
\pgfsys@transformshift{5.330586in}{7.719748in}%
\pgfsys@useobject{currentmarker}{}%
\end{pgfscope}%
\begin{pgfscope}%
\pgfsys@transformshift{5.979646in}{9.132322in}%
\pgfsys@useobject{currentmarker}{}%
\end{pgfscope}%
\begin{pgfscope}%
\pgfsys@transformshift{5.461111in}{8.562867in}%
\pgfsys@useobject{currentmarker}{}%
\end{pgfscope}%
\begin{pgfscope}%
\pgfsys@transformshift{2.563110in}{5.970834in}%
\pgfsys@useobject{currentmarker}{}%
\end{pgfscope}%
\begin{pgfscope}%
\pgfsys@transformshift{2.982456in}{6.196579in}%
\pgfsys@useobject{currentmarker}{}%
\end{pgfscope}%
\begin{pgfscope}%
\pgfsys@transformshift{2.497801in}{5.859820in}%
\pgfsys@useobject{currentmarker}{}%
\end{pgfscope}%
\begin{pgfscope}%
\pgfsys@transformshift{2.664702in}{6.110016in}%
\pgfsys@useobject{currentmarker}{}%
\end{pgfscope}%
\begin{pgfscope}%
\pgfsys@transformshift{2.628748in}{7.095888in}%
\pgfsys@useobject{currentmarker}{}%
\end{pgfscope}%
\begin{pgfscope}%
\pgfsys@transformshift{2.870857in}{6.927409in}%
\pgfsys@useobject{currentmarker}{}%
\end{pgfscope}%
\begin{pgfscope}%
\pgfsys@transformshift{2.524839in}{4.746731in}%
\pgfsys@useobject{currentmarker}{}%
\end{pgfscope}%
\begin{pgfscope}%
\pgfsys@transformshift{3.452575in}{6.324784in}%
\pgfsys@useobject{currentmarker}{}%
\end{pgfscope}%
\begin{pgfscope}%
\pgfsys@transformshift{2.991326in}{6.307716in}%
\pgfsys@useobject{currentmarker}{}%
\end{pgfscope}%
\begin{pgfscope}%
\pgfsys@transformshift{2.314533in}{6.670153in}%
\pgfsys@useobject{currentmarker}{}%
\end{pgfscope}%
\begin{pgfscope}%
\pgfsys@transformshift{2.631112in}{6.772558in}%
\pgfsys@useobject{currentmarker}{}%
\end{pgfscope}%
\begin{pgfscope}%
\pgfsys@transformshift{4.341546in}{8.032677in}%
\pgfsys@useobject{currentmarker}{}%
\end{pgfscope}%
\begin{pgfscope}%
\pgfsys@transformshift{4.698332in}{8.044531in}%
\pgfsys@useobject{currentmarker}{}%
\end{pgfscope}%
\begin{pgfscope}%
\pgfsys@transformshift{2.387607in}{5.091740in}%
\pgfsys@useobject{currentmarker}{}%
\end{pgfscope}%
\begin{pgfscope}%
\pgfsys@transformshift{4.719868in}{7.963388in}%
\pgfsys@useobject{currentmarker}{}%
\end{pgfscope}%
\begin{pgfscope}%
\pgfsys@transformshift{4.840741in}{8.006626in}%
\pgfsys@useobject{currentmarker}{}%
\end{pgfscope}%
\begin{pgfscope}%
\pgfsys@transformshift{4.942178in}{8.081004in}%
\pgfsys@useobject{currentmarker}{}%
\end{pgfscope}%
\begin{pgfscope}%
\pgfsys@transformshift{5.187717in}{8.313906in}%
\pgfsys@useobject{currentmarker}{}%
\end{pgfscope}%
\begin{pgfscope}%
\pgfsys@transformshift{2.283149in}{5.532419in}%
\pgfsys@useobject{currentmarker}{}%
\end{pgfscope}%
\begin{pgfscope}%
\pgfsys@transformshift{2.173829in}{4.983989in}%
\pgfsys@useobject{currentmarker}{}%
\end{pgfscope}%
\begin{pgfscope}%
\pgfsys@transformshift{4.474776in}{7.124548in}%
\pgfsys@useobject{currentmarker}{}%
\end{pgfscope}%
\begin{pgfscope}%
\pgfsys@transformshift{4.780807in}{6.861523in}%
\pgfsys@useobject{currentmarker}{}%
\end{pgfscope}%
\begin{pgfscope}%
\pgfsys@transformshift{6.508073in}{7.475589in}%
\pgfsys@useobject{currentmarker}{}%
\end{pgfscope}%
\begin{pgfscope}%
\pgfsys@transformshift{6.358741in}{8.126442in}%
\pgfsys@useobject{currentmarker}{}%
\end{pgfscope}%
\end{pgfscope}%
\begin{pgfscope}%
\pgfsetbuttcap%
\pgfsetroundjoin%
\definecolor{currentfill}{rgb}{0.321569,0.321569,0.321569}%
\pgfsetfillcolor{currentfill}%
\pgfsetlinewidth{1.254687pt}%
\definecolor{currentstroke}{rgb}{0.321569,0.321569,0.321569}%
\pgfsetstrokecolor{currentstroke}%
\pgfsetdash{}{0pt}%
\pgfsys@defobject{currentmarker}{\pgfqpoint{-0.043368in}{-0.043368in}}{\pgfqpoint{0.043368in}{0.043368in}}{%
\pgfpathmoveto{\pgfqpoint{-0.043368in}{-0.043368in}}%
\pgfpathlineto{\pgfqpoint{0.043368in}{0.043368in}}%
\pgfpathmoveto{\pgfqpoint{-0.043368in}{0.043368in}}%
\pgfpathlineto{\pgfqpoint{0.043368in}{-0.043368in}}%
\pgfusepath{stroke,fill}%
}%
\begin{pgfscope}%
\pgfsys@transformshift{1.109909in}{4.050000in}%
\pgfsys@useobject{currentmarker}{}%
\end{pgfscope}%
\begin{pgfscope}%
\pgfsys@transformshift{1.109909in}{4.614015in}%
\pgfsys@useobject{currentmarker}{}%
\end{pgfscope}%
\begin{pgfscope}%
\pgfsys@transformshift{1.109909in}{4.050000in}%
\pgfsys@useobject{currentmarker}{}%
\end{pgfscope}%
\end{pgfscope}%
\begin{pgfscope}%
\pgfpathrectangle{\pgfqpoint{1.109909in}{4.050000in}}{\pgfqpoint{5.400000in}{5.400000in}}%
\pgfusepath{clip}%
\pgfsetbuttcap%
\pgfsetroundjoin%
\definecolor{currentfill}{rgb}{0.768627,0.603922,0.000000}%
\pgfsetfillcolor{currentfill}%
\pgfsetlinewidth{1.254687pt}%
\definecolor{currentstroke}{rgb}{0.768627,0.603922,0.000000}%
\pgfsetstrokecolor{currentstroke}%
\pgfsetdash{}{0pt}%
\pgfsys@defobject{currentmarker}{\pgfqpoint{-0.043368in}{-0.043368in}}{\pgfqpoint{0.043368in}{0.043368in}}{%
\pgfpathmoveto{\pgfqpoint{-0.043368in}{-0.043368in}}%
\pgfpathlineto{\pgfqpoint{0.043368in}{0.043368in}}%
\pgfpathmoveto{\pgfqpoint{-0.043368in}{0.043368in}}%
\pgfpathlineto{\pgfqpoint{0.043368in}{-0.043368in}}%
\pgfusepath{stroke,fill}%
}%
\begin{pgfscope}%
\pgfsys@transformshift{5.323982in}{6.998816in}%
\pgfsys@useobject{currentmarker}{}%
\end{pgfscope}%
\begin{pgfscope}%
\pgfsys@transformshift{5.353500in}{6.910379in}%
\pgfsys@useobject{currentmarker}{}%
\end{pgfscope}%
\begin{pgfscope}%
\pgfsys@transformshift{5.286143in}{7.490496in}%
\pgfsys@useobject{currentmarker}{}%
\end{pgfscope}%
\begin{pgfscope}%
\pgfsys@transformshift{5.039638in}{6.643276in}%
\pgfsys@useobject{currentmarker}{}%
\end{pgfscope}%
\end{pgfscope}%
\begin{pgfscope}%
\pgfpathrectangle{\pgfqpoint{1.109909in}{4.050000in}}{\pgfqpoint{5.400000in}{5.400000in}}%
\pgfusepath{clip}%
\pgfsetbuttcap%
\pgfsetroundjoin%
\definecolor{currentfill}{rgb}{0.768627,0.603922,0.000000}%
\pgfsetfillcolor{currentfill}%
\pgfsetlinewidth{1.254687pt}%
\definecolor{currentstroke}{rgb}{0.768627,0.603922,0.000000}%
\pgfsetstrokecolor{currentstroke}%
\pgfsetdash{}{0pt}%
\pgfsys@defobject{currentmarker}{\pgfqpoint{-0.043368in}{-0.043368in}}{\pgfqpoint{0.043368in}{0.043368in}}{%
\pgfpathmoveto{\pgfqpoint{-0.043368in}{-0.043368in}}%
\pgfpathlineto{\pgfqpoint{0.043368in}{0.043368in}}%
\pgfpathmoveto{\pgfqpoint{-0.043368in}{0.043368in}}%
\pgfpathlineto{\pgfqpoint{0.043368in}{-0.043368in}}%
\pgfusepath{stroke,fill}%
}%
\begin{pgfscope}%
\pgfsys@transformshift{3.569747in}{6.369413in}%
\pgfsys@useobject{currentmarker}{}%
\end{pgfscope}%
\begin{pgfscope}%
\pgfsys@transformshift{3.444025in}{6.194464in}%
\pgfsys@useobject{currentmarker}{}%
\end{pgfscope}%
\end{pgfscope}%
\begin{pgfscope}%
\pgfpathrectangle{\pgfqpoint{1.109909in}{4.050000in}}{\pgfqpoint{5.400000in}{5.400000in}}%
\pgfusepath{clip}%
\pgfsetbuttcap%
\pgfsetroundjoin%
\definecolor{currentfill}{rgb}{0.768627,0.603922,0.000000}%
\pgfsetfillcolor{currentfill}%
\pgfsetlinewidth{1.254687pt}%
\definecolor{currentstroke}{rgb}{0.768627,0.603922,0.000000}%
\pgfsetstrokecolor{currentstroke}%
\pgfsetdash{}{0pt}%
\pgfsys@defobject{currentmarker}{\pgfqpoint{-0.043368in}{-0.043368in}}{\pgfqpoint{0.043368in}{0.043368in}}{%
\pgfpathmoveto{\pgfqpoint{-0.043368in}{-0.043368in}}%
\pgfpathlineto{\pgfqpoint{0.043368in}{0.043368in}}%
\pgfpathmoveto{\pgfqpoint{-0.043368in}{0.043368in}}%
\pgfpathlineto{\pgfqpoint{0.043368in}{-0.043368in}}%
\pgfusepath{stroke,fill}%
}%
\begin{pgfscope}%
\pgfsys@transformshift{5.049865in}{7.731584in}%
\pgfsys@useobject{currentmarker}{}%
\end{pgfscope}%
\begin{pgfscope}%
\pgfsys@transformshift{4.588717in}{7.202925in}%
\pgfsys@useobject{currentmarker}{}%
\end{pgfscope}%
\end{pgfscope}%
\begin{pgfscope}%
\pgfsetbuttcap%
\pgfsetroundjoin%
\definecolor{currentfill}{rgb}{0.337255,0.705882,0.913725}%
\pgfsetfillcolor{currentfill}%
\pgfsetlinewidth{1.254687pt}%
\definecolor{currentstroke}{rgb}{0.337255,0.705882,0.913725}%
\pgfsetstrokecolor{currentstroke}%
\pgfsetdash{}{0pt}%
\pgfsys@defobject{currentmarker}{\pgfqpoint{-0.043368in}{-0.043368in}}{\pgfqpoint{0.043368in}{0.043368in}}{%
\pgfpathmoveto{\pgfqpoint{-0.043368in}{-0.043368in}}%
\pgfpathlineto{\pgfqpoint{0.043368in}{0.043368in}}%
\pgfpathmoveto{\pgfqpoint{-0.043368in}{0.043368in}}%
\pgfpathlineto{\pgfqpoint{0.043368in}{-0.043368in}}%
\pgfusepath{stroke,fill}%
}%
\begin{pgfscope}%
\pgfsys@transformshift{6.506560in}{9.450000in}%
\pgfsys@useobject{currentmarker}{}%
\end{pgfscope}%
\end{pgfscope}%
\begin{pgfscope}%
\pgfpathrectangle{\pgfqpoint{1.109909in}{4.050000in}}{\pgfqpoint{5.400000in}{5.400000in}}%
\pgfusepath{clip}%
\pgfsetbuttcap%
\pgfsetroundjoin%
\definecolor{currentfill}{rgb}{0.901961,0.196078,0.156863}%
\pgfsetfillcolor{currentfill}%
\pgfsetlinewidth{1.254687pt}%
\definecolor{currentstroke}{rgb}{0.901961,0.196078,0.156863}%
\pgfsetstrokecolor{currentstroke}%
\pgfsetdash{}{0pt}%
\pgfsys@defobject{currentmarker}{\pgfqpoint{-0.043368in}{-0.043368in}}{\pgfqpoint{0.043368in}{0.043368in}}{%
\pgfpathmoveto{\pgfqpoint{-0.043368in}{-0.043368in}}%
\pgfpathlineto{\pgfqpoint{0.043368in}{0.043368in}}%
\pgfpathmoveto{\pgfqpoint{-0.043368in}{0.043368in}}%
\pgfpathlineto{\pgfqpoint{0.043368in}{-0.043368in}}%
\pgfusepath{stroke,fill}%
}%
\begin{pgfscope}%
\pgfsys@transformshift{5.233005in}{8.698495in}%
\pgfsys@useobject{currentmarker}{}%
\end{pgfscope}%
\begin{pgfscope}%
\pgfsys@transformshift{4.016124in}{7.059411in}%
\pgfsys@useobject{currentmarker}{}%
\end{pgfscope}%
\end{pgfscope}%
\begin{pgfscope}%
\pgfsetbuttcap%
\pgfsetroundjoin%
\definecolor{currentfill}{rgb}{0.901961,0.196078,0.156863}%
\pgfsetfillcolor{currentfill}%
\pgfsetlinewidth{1.254687pt}%
\definecolor{currentstroke}{rgb}{0.901961,0.196078,0.156863}%
\pgfsetstrokecolor{currentstroke}%
\pgfsetdash{}{0pt}%
\pgfsys@defobject{currentmarker}{\pgfqpoint{-0.043368in}{-0.043368in}}{\pgfqpoint{0.043368in}{0.043368in}}{%
\pgfpathmoveto{\pgfqpoint{-0.043368in}{-0.043368in}}%
\pgfpathlineto{\pgfqpoint{0.043368in}{0.043368in}}%
\pgfpathmoveto{\pgfqpoint{-0.043368in}{0.043368in}}%
\pgfpathlineto{\pgfqpoint{0.043368in}{-0.043368in}}%
\pgfusepath{stroke,fill}%
}%
\begin{pgfscope}%
\pgfsys@transformshift{6.509909in}{9.261250in}%
\pgfsys@useobject{currentmarker}{}%
\end{pgfscope}%
\end{pgfscope}%
\begin{pgfscope}%
\pgfpathrectangle{\pgfqpoint{1.109909in}{4.050000in}}{\pgfqpoint{5.400000in}{5.400000in}}%
\pgfusepath{clip}%
\pgfsetbuttcap%
\pgfsetroundjoin%
\definecolor{currentfill}{rgb}{0.835294,0.368627,0.000000}%
\pgfsetfillcolor{currentfill}%
\pgfsetlinewidth{1.254687pt}%
\definecolor{currentstroke}{rgb}{0.835294,0.368627,0.000000}%
\pgfsetstrokecolor{currentstroke}%
\pgfsetdash{}{0pt}%
\pgfsys@defobject{currentmarker}{\pgfqpoint{-0.043368in}{-0.043368in}}{\pgfqpoint{0.043368in}{0.043368in}}{%
\pgfpathmoveto{\pgfqpoint{-0.043368in}{-0.043368in}}%
\pgfpathlineto{\pgfqpoint{0.043368in}{0.043368in}}%
\pgfpathmoveto{\pgfqpoint{-0.043368in}{0.043368in}}%
\pgfpathlineto{\pgfqpoint{0.043368in}{-0.043368in}}%
\pgfusepath{stroke,fill}%
}%
\begin{pgfscope}%
\pgfsys@transformshift{4.924005in}{6.573326in}%
\pgfsys@useobject{currentmarker}{}%
\end{pgfscope}%
\end{pgfscope}%
\begin{pgfscope}%
\pgfpathrectangle{\pgfqpoint{1.109909in}{4.050000in}}{\pgfqpoint{5.400000in}{5.400000in}}%
\pgfusepath{clip}%
\pgfsetbuttcap%
\pgfsetroundjoin%
\pgfsetlinewidth{1.003750pt}%
\definecolor{currentstroke}{rgb}{0.541176,0.380392,0.658824}%
\pgfsetstrokecolor{currentstroke}%
\pgfsetstrokeopacity{0.900000}%
\pgfsetdash{}{0pt}%
\pgfpathmoveto{\pgfqpoint{4.989018in}{8.297326in}}%
\pgfpathlineto{\pgfqpoint{5.073501in}{8.297326in}}%
\pgfpathlineto{\pgfqpoint{5.073501in}{8.381809in}}%
\pgfpathlineto{\pgfqpoint{4.989018in}{8.381809in}}%
\pgfpathlineto{\pgfqpoint{4.989018in}{8.297326in}}%
\pgfpathclose%
\pgfusepath{stroke}%
\end{pgfscope}%
\begin{pgfscope}%
\pgfpathrectangle{\pgfqpoint{1.109909in}{4.050000in}}{\pgfqpoint{5.400000in}{5.400000in}}%
\pgfusepath{clip}%
\pgfsetbuttcap%
\pgfsetroundjoin%
\pgfsetlinewidth{1.003750pt}%
\definecolor{currentstroke}{rgb}{0.541176,0.380392,0.658824}%
\pgfsetstrokecolor{currentstroke}%
\pgfsetstrokeopacity{0.900000}%
\pgfsetdash{}{0pt}%
\pgfpathmoveto{\pgfqpoint{4.849701in}{7.586784in}}%
\pgfpathlineto{\pgfqpoint{4.934184in}{7.586784in}}%
\pgfpathlineto{\pgfqpoint{4.934184in}{7.671267in}}%
\pgfpathlineto{\pgfqpoint{4.849701in}{7.671267in}}%
\pgfpathlineto{\pgfqpoint{4.849701in}{7.586784in}}%
\pgfpathclose%
\pgfusepath{stroke}%
\end{pgfscope}%
\begin{pgfscope}%
\pgfpathrectangle{\pgfqpoint{1.109909in}{4.050000in}}{\pgfqpoint{5.400000in}{5.400000in}}%
\pgfusepath{clip}%
\pgfsetbuttcap%
\pgfsetroundjoin%
\pgfsetlinewidth{1.003750pt}%
\definecolor{currentstroke}{rgb}{0.541176,0.380392,0.658824}%
\pgfsetstrokecolor{currentstroke}%
\pgfsetstrokeopacity{0.900000}%
\pgfsetdash{}{0pt}%
\pgfpathmoveto{\pgfqpoint{5.158233in}{7.210262in}}%
\pgfpathlineto{\pgfqpoint{5.242716in}{7.210262in}}%
\pgfpathlineto{\pgfqpoint{5.242716in}{7.294744in}}%
\pgfpathlineto{\pgfqpoint{5.158233in}{7.294744in}}%
\pgfpathlineto{\pgfqpoint{5.158233in}{7.210262in}}%
\pgfpathclose%
\pgfusepath{stroke}%
\end{pgfscope}%
\begin{pgfscope}%
\pgfpathrectangle{\pgfqpoint{1.109909in}{4.050000in}}{\pgfqpoint{5.400000in}{5.400000in}}%
\pgfusepath{clip}%
\pgfsetbuttcap%
\pgfsetroundjoin%
\pgfsetlinewidth{1.003750pt}%
\definecolor{currentstroke}{rgb}{0.541176,0.380392,0.658824}%
\pgfsetstrokecolor{currentstroke}%
\pgfsetstrokeopacity{0.900000}%
\pgfsetdash{}{0pt}%
\pgfpathmoveto{\pgfqpoint{4.936433in}{8.558592in}}%
\pgfpathlineto{\pgfqpoint{5.020916in}{8.558592in}}%
\pgfpathlineto{\pgfqpoint{5.020916in}{8.643075in}}%
\pgfpathlineto{\pgfqpoint{4.936433in}{8.643075in}}%
\pgfpathlineto{\pgfqpoint{4.936433in}{8.558592in}}%
\pgfpathclose%
\pgfusepath{stroke}%
\end{pgfscope}%
\begin{pgfscope}%
\pgfpathrectangle{\pgfqpoint{1.109909in}{4.050000in}}{\pgfqpoint{5.400000in}{5.400000in}}%
\pgfusepath{clip}%
\pgfsetbuttcap%
\pgfsetroundjoin%
\pgfsetlinewidth{1.003750pt}%
\definecolor{currentstroke}{rgb}{0.541176,0.380392,0.658824}%
\pgfsetstrokecolor{currentstroke}%
\pgfsetstrokeopacity{0.900000}%
\pgfsetdash{}{0pt}%
\pgfpathmoveto{\pgfqpoint{4.480377in}{8.293545in}}%
\pgfpathlineto{\pgfqpoint{4.564860in}{8.293545in}}%
\pgfpathlineto{\pgfqpoint{4.564860in}{8.378028in}}%
\pgfpathlineto{\pgfqpoint{4.480377in}{8.378028in}}%
\pgfpathlineto{\pgfqpoint{4.480377in}{8.293545in}}%
\pgfpathclose%
\pgfusepath{stroke}%
\end{pgfscope}%
\begin{pgfscope}%
\pgfpathrectangle{\pgfqpoint{1.109909in}{4.050000in}}{\pgfqpoint{5.400000in}{5.400000in}}%
\pgfusepath{clip}%
\pgfsetbuttcap%
\pgfsetroundjoin%
\pgfsetlinewidth{1.003750pt}%
\definecolor{currentstroke}{rgb}{0.541176,0.380392,0.658824}%
\pgfsetstrokecolor{currentstroke}%
\pgfsetstrokeopacity{0.900000}%
\pgfsetdash{}{0pt}%
\pgfpathmoveto{\pgfqpoint{5.430957in}{8.382940in}}%
\pgfpathlineto{\pgfqpoint{5.515440in}{8.382940in}}%
\pgfpathlineto{\pgfqpoint{5.515440in}{8.467423in}}%
\pgfpathlineto{\pgfqpoint{5.430957in}{8.467423in}}%
\pgfpathlineto{\pgfqpoint{5.430957in}{8.382940in}}%
\pgfpathclose%
\pgfusepath{stroke}%
\end{pgfscope}%
\begin{pgfscope}%
\pgfpathrectangle{\pgfqpoint{1.109909in}{4.050000in}}{\pgfqpoint{5.400000in}{5.400000in}}%
\pgfusepath{clip}%
\pgfsetbuttcap%
\pgfsetroundjoin%
\pgfsetlinewidth{1.003750pt}%
\definecolor{currentstroke}{rgb}{0.541176,0.380392,0.658824}%
\pgfsetstrokecolor{currentstroke}%
\pgfsetstrokeopacity{0.900000}%
\pgfsetdash{}{0pt}%
\pgfpathmoveto{\pgfqpoint{4.099474in}{7.712172in}}%
\pgfpathlineto{\pgfqpoint{4.183957in}{7.712172in}}%
\pgfpathlineto{\pgfqpoint{4.183957in}{7.796655in}}%
\pgfpathlineto{\pgfqpoint{4.099474in}{7.796655in}}%
\pgfpathlineto{\pgfqpoint{4.099474in}{7.712172in}}%
\pgfpathclose%
\pgfusepath{stroke}%
\end{pgfscope}%
\begin{pgfscope}%
\pgfpathrectangle{\pgfqpoint{1.109909in}{4.050000in}}{\pgfqpoint{5.400000in}{5.400000in}}%
\pgfusepath{clip}%
\pgfsetbuttcap%
\pgfsetroundjoin%
\pgfsetlinewidth{1.003750pt}%
\definecolor{currentstroke}{rgb}{0.541176,0.380392,0.658824}%
\pgfsetstrokecolor{currentstroke}%
\pgfsetstrokeopacity{0.900000}%
\pgfsetdash{}{0pt}%
\pgfpathmoveto{\pgfqpoint{3.606780in}{6.555414in}}%
\pgfpathlineto{\pgfqpoint{3.691263in}{6.555414in}}%
\pgfpathlineto{\pgfqpoint{3.691263in}{6.639896in}}%
\pgfpathlineto{\pgfqpoint{3.606780in}{6.639896in}}%
\pgfpathlineto{\pgfqpoint{3.606780in}{6.555414in}}%
\pgfpathclose%
\pgfusepath{stroke}%
\end{pgfscope}%
\begin{pgfscope}%
\pgfpathrectangle{\pgfqpoint{1.109909in}{4.050000in}}{\pgfqpoint{5.400000in}{5.400000in}}%
\pgfusepath{clip}%
\pgfsetbuttcap%
\pgfsetroundjoin%
\pgfsetlinewidth{1.003750pt}%
\definecolor{currentstroke}{rgb}{0.541176,0.380392,0.658824}%
\pgfsetstrokecolor{currentstroke}%
\pgfsetstrokeopacity{0.900000}%
\pgfsetdash{}{0pt}%
\pgfpathmoveto{\pgfqpoint{2.481071in}{5.588147in}}%
\pgfpathlineto{\pgfqpoint{2.565554in}{5.588147in}}%
\pgfpathlineto{\pgfqpoint{2.565554in}{5.672629in}}%
\pgfpathlineto{\pgfqpoint{2.481071in}{5.672629in}}%
\pgfpathlineto{\pgfqpoint{2.481071in}{5.588147in}}%
\pgfpathclose%
\pgfusepath{stroke}%
\end{pgfscope}%
\begin{pgfscope}%
\pgfpathrectangle{\pgfqpoint{1.109909in}{4.050000in}}{\pgfqpoint{5.400000in}{5.400000in}}%
\pgfusepath{clip}%
\pgfsetbuttcap%
\pgfsetroundjoin%
\pgfsetlinewidth{1.003750pt}%
\definecolor{currentstroke}{rgb}{0.541176,0.380392,0.658824}%
\pgfsetstrokecolor{currentstroke}%
\pgfsetstrokeopacity{0.900000}%
\pgfsetdash{}{0pt}%
\pgfpathmoveto{\pgfqpoint{4.055392in}{6.739341in}}%
\pgfpathlineto{\pgfqpoint{4.139875in}{6.739341in}}%
\pgfpathlineto{\pgfqpoint{4.139875in}{6.823824in}}%
\pgfpathlineto{\pgfqpoint{4.055392in}{6.823824in}}%
\pgfpathlineto{\pgfqpoint{4.055392in}{6.739341in}}%
\pgfpathclose%
\pgfusepath{stroke}%
\end{pgfscope}%
\begin{pgfscope}%
\pgfpathrectangle{\pgfqpoint{1.109909in}{4.050000in}}{\pgfqpoint{5.400000in}{5.400000in}}%
\pgfusepath{clip}%
\pgfsetbuttcap%
\pgfsetroundjoin%
\pgfsetlinewidth{1.003750pt}%
\definecolor{currentstroke}{rgb}{0.541176,0.380392,0.658824}%
\pgfsetstrokecolor{currentstroke}%
\pgfsetstrokeopacity{0.900000}%
\pgfsetdash{}{0pt}%
\pgfpathmoveto{\pgfqpoint{5.038672in}{7.330026in}}%
\pgfpathlineto{\pgfqpoint{5.123155in}{7.330026in}}%
\pgfpathlineto{\pgfqpoint{5.123155in}{7.414509in}}%
\pgfpathlineto{\pgfqpoint{5.038672in}{7.414509in}}%
\pgfpathlineto{\pgfqpoint{5.038672in}{7.330026in}}%
\pgfpathclose%
\pgfusepath{stroke}%
\end{pgfscope}%
\begin{pgfscope}%
\pgfpathrectangle{\pgfqpoint{1.109909in}{4.050000in}}{\pgfqpoint{5.400000in}{5.400000in}}%
\pgfusepath{clip}%
\pgfsetbuttcap%
\pgfsetroundjoin%
\pgfsetlinewidth{1.003750pt}%
\definecolor{currentstroke}{rgb}{0.541176,0.380392,0.658824}%
\pgfsetstrokecolor{currentstroke}%
\pgfsetstrokeopacity{0.900000}%
\pgfsetdash{}{0pt}%
\pgfpathmoveto{\pgfqpoint{3.535271in}{5.936701in}}%
\pgfpathlineto{\pgfqpoint{3.619754in}{5.936701in}}%
\pgfpathlineto{\pgfqpoint{3.619754in}{6.021184in}}%
\pgfpathlineto{\pgfqpoint{3.535271in}{6.021184in}}%
\pgfpathlineto{\pgfqpoint{3.535271in}{5.936701in}}%
\pgfpathclose%
\pgfusepath{stroke}%
\end{pgfscope}%
\begin{pgfscope}%
\pgfpathrectangle{\pgfqpoint{1.109909in}{4.050000in}}{\pgfqpoint{5.400000in}{5.400000in}}%
\pgfusepath{clip}%
\pgfsetbuttcap%
\pgfsetroundjoin%
\pgfsetlinewidth{1.003750pt}%
\definecolor{currentstroke}{rgb}{0.541176,0.380392,0.658824}%
\pgfsetstrokecolor{currentstroke}%
\pgfsetstrokeopacity{0.900000}%
\pgfsetdash{}{0pt}%
\pgfpathmoveto{\pgfqpoint{3.502361in}{5.525869in}}%
\pgfpathlineto{\pgfqpoint{3.586844in}{5.525869in}}%
\pgfpathlineto{\pgfqpoint{3.586844in}{5.610352in}}%
\pgfpathlineto{\pgfqpoint{3.502361in}{5.610352in}}%
\pgfpathlineto{\pgfqpoint{3.502361in}{5.525869in}}%
\pgfpathclose%
\pgfusepath{stroke}%
\end{pgfscope}%
\begin{pgfscope}%
\pgfpathrectangle{\pgfqpoint{1.109909in}{4.050000in}}{\pgfqpoint{5.400000in}{5.400000in}}%
\pgfusepath{clip}%
\pgfsetbuttcap%
\pgfsetroundjoin%
\pgfsetlinewidth{1.003750pt}%
\definecolor{currentstroke}{rgb}{0.541176,0.380392,0.658824}%
\pgfsetstrokecolor{currentstroke}%
\pgfsetstrokeopacity{0.900000}%
\pgfsetdash{}{0pt}%
\pgfpathmoveto{\pgfqpoint{3.208139in}{5.335598in}}%
\pgfpathlineto{\pgfqpoint{3.292622in}{5.335598in}}%
\pgfpathlineto{\pgfqpoint{3.292622in}{5.420081in}}%
\pgfpathlineto{\pgfqpoint{3.208139in}{5.420081in}}%
\pgfpathlineto{\pgfqpoint{3.208139in}{5.335598in}}%
\pgfpathclose%
\pgfusepath{stroke}%
\end{pgfscope}%
\begin{pgfscope}%
\pgfpathrectangle{\pgfqpoint{1.109909in}{4.050000in}}{\pgfqpoint{5.400000in}{5.400000in}}%
\pgfusepath{clip}%
\pgfsetbuttcap%
\pgfsetroundjoin%
\pgfsetlinewidth{1.003750pt}%
\definecolor{currentstroke}{rgb}{0.541176,0.380392,0.658824}%
\pgfsetstrokecolor{currentstroke}%
\pgfsetstrokeopacity{0.900000}%
\pgfsetdash{}{0pt}%
\pgfpathmoveto{\pgfqpoint{2.610664in}{5.179465in}}%
\pgfpathlineto{\pgfqpoint{2.695147in}{5.179465in}}%
\pgfpathlineto{\pgfqpoint{2.695147in}{5.263948in}}%
\pgfpathlineto{\pgfqpoint{2.610664in}{5.263948in}}%
\pgfpathlineto{\pgfqpoint{2.610664in}{5.179465in}}%
\pgfpathclose%
\pgfusepath{stroke}%
\end{pgfscope}%
\begin{pgfscope}%
\pgfpathrectangle{\pgfqpoint{1.109909in}{4.050000in}}{\pgfqpoint{5.400000in}{5.400000in}}%
\pgfusepath{clip}%
\pgfsetbuttcap%
\pgfsetroundjoin%
\pgfsetlinewidth{1.003750pt}%
\definecolor{currentstroke}{rgb}{0.541176,0.380392,0.658824}%
\pgfsetstrokecolor{currentstroke}%
\pgfsetstrokeopacity{0.900000}%
\pgfsetdash{}{0pt}%
\pgfpathmoveto{\pgfqpoint{4.763176in}{7.726518in}}%
\pgfpathlineto{\pgfqpoint{4.847659in}{7.726518in}}%
\pgfpathlineto{\pgfqpoint{4.847659in}{7.811001in}}%
\pgfpathlineto{\pgfqpoint{4.763176in}{7.811001in}}%
\pgfpathlineto{\pgfqpoint{4.763176in}{7.726518in}}%
\pgfpathclose%
\pgfusepath{stroke}%
\end{pgfscope}%
\begin{pgfscope}%
\pgfpathrectangle{\pgfqpoint{1.109909in}{4.050000in}}{\pgfqpoint{5.400000in}{5.400000in}}%
\pgfusepath{clip}%
\pgfsetbuttcap%
\pgfsetroundjoin%
\pgfsetlinewidth{1.003750pt}%
\definecolor{currentstroke}{rgb}{0.541176,0.380392,0.658824}%
\pgfsetstrokecolor{currentstroke}%
\pgfsetstrokeopacity{0.900000}%
\pgfsetdash{}{0pt}%
\pgfpathmoveto{\pgfqpoint{4.352926in}{6.957558in}}%
\pgfpathlineto{\pgfqpoint{4.437409in}{6.957558in}}%
\pgfpathlineto{\pgfqpoint{4.437409in}{7.042041in}}%
\pgfpathlineto{\pgfqpoint{4.352926in}{7.042041in}}%
\pgfpathlineto{\pgfqpoint{4.352926in}{6.957558in}}%
\pgfpathclose%
\pgfusepath{stroke}%
\end{pgfscope}%
\begin{pgfscope}%
\pgfpathrectangle{\pgfqpoint{1.109909in}{4.050000in}}{\pgfqpoint{5.400000in}{5.400000in}}%
\pgfusepath{clip}%
\pgfsetbuttcap%
\pgfsetroundjoin%
\pgfsetlinewidth{1.003750pt}%
\definecolor{currentstroke}{rgb}{0.541176,0.380392,0.658824}%
\pgfsetstrokecolor{currentstroke}%
\pgfsetstrokeopacity{0.900000}%
\pgfsetdash{}{0pt}%
\pgfpathmoveto{\pgfqpoint{4.109384in}{8.326995in}}%
\pgfpathlineto{\pgfqpoint{4.193867in}{8.326995in}}%
\pgfpathlineto{\pgfqpoint{4.193867in}{8.411477in}}%
\pgfpathlineto{\pgfqpoint{4.109384in}{8.411477in}}%
\pgfpathlineto{\pgfqpoint{4.109384in}{8.326995in}}%
\pgfpathclose%
\pgfusepath{stroke}%
\end{pgfscope}%
\begin{pgfscope}%
\pgfpathrectangle{\pgfqpoint{1.109909in}{4.050000in}}{\pgfqpoint{5.400000in}{5.400000in}}%
\pgfusepath{clip}%
\pgfsetbuttcap%
\pgfsetroundjoin%
\pgfsetlinewidth{1.003750pt}%
\definecolor{currentstroke}{rgb}{0.541176,0.380392,0.658824}%
\pgfsetstrokecolor{currentstroke}%
\pgfsetstrokeopacity{0.900000}%
\pgfsetdash{}{0pt}%
\pgfpathmoveto{\pgfqpoint{3.797609in}{5.557686in}}%
\pgfpathlineto{\pgfqpoint{3.882092in}{5.557686in}}%
\pgfpathlineto{\pgfqpoint{3.882092in}{5.642168in}}%
\pgfpathlineto{\pgfqpoint{3.797609in}{5.642168in}}%
\pgfpathlineto{\pgfqpoint{3.797609in}{5.557686in}}%
\pgfpathclose%
\pgfusepath{stroke}%
\end{pgfscope}%
\begin{pgfscope}%
\pgfpathrectangle{\pgfqpoint{1.109909in}{4.050000in}}{\pgfqpoint{5.400000in}{5.400000in}}%
\pgfusepath{clip}%
\pgfsetbuttcap%
\pgfsetroundjoin%
\pgfsetlinewidth{1.003750pt}%
\definecolor{currentstroke}{rgb}{0.541176,0.380392,0.658824}%
\pgfsetstrokecolor{currentstroke}%
\pgfsetstrokeopacity{0.900000}%
\pgfsetdash{}{0pt}%
\pgfpathmoveto{\pgfqpoint{3.194372in}{5.141574in}}%
\pgfpathlineto{\pgfqpoint{3.278855in}{5.141574in}}%
\pgfpathlineto{\pgfqpoint{3.278855in}{5.226056in}}%
\pgfpathlineto{\pgfqpoint{3.194372in}{5.226056in}}%
\pgfpathlineto{\pgfqpoint{3.194372in}{5.141574in}}%
\pgfpathclose%
\pgfusepath{stroke}%
\end{pgfscope}%
\begin{pgfscope}%
\pgfpathrectangle{\pgfqpoint{1.109909in}{4.050000in}}{\pgfqpoint{5.400000in}{5.400000in}}%
\pgfusepath{clip}%
\pgfsetbuttcap%
\pgfsetroundjoin%
\pgfsetlinewidth{1.003750pt}%
\definecolor{currentstroke}{rgb}{0.541176,0.380392,0.658824}%
\pgfsetstrokecolor{currentstroke}%
\pgfsetstrokeopacity{0.900000}%
\pgfsetdash{}{0pt}%
\pgfpathmoveto{\pgfqpoint{3.559677in}{6.135995in}}%
\pgfpathlineto{\pgfqpoint{3.644160in}{6.135995in}}%
\pgfpathlineto{\pgfqpoint{3.644160in}{6.220477in}}%
\pgfpathlineto{\pgfqpoint{3.559677in}{6.220477in}}%
\pgfpathlineto{\pgfqpoint{3.559677in}{6.135995in}}%
\pgfpathclose%
\pgfusepath{stroke}%
\end{pgfscope}%
\begin{pgfscope}%
\pgfpathrectangle{\pgfqpoint{1.109909in}{4.050000in}}{\pgfqpoint{5.400000in}{5.400000in}}%
\pgfusepath{clip}%
\pgfsetbuttcap%
\pgfsetroundjoin%
\pgfsetlinewidth{1.003750pt}%
\definecolor{currentstroke}{rgb}{0.541176,0.380392,0.658824}%
\pgfsetstrokecolor{currentstroke}%
\pgfsetstrokeopacity{0.900000}%
\pgfsetdash{}{0pt}%
\pgfpathmoveto{\pgfqpoint{2.840063in}{5.084429in}}%
\pgfpathlineto{\pgfqpoint{2.924546in}{5.084429in}}%
\pgfpathlineto{\pgfqpoint{2.924546in}{5.168912in}}%
\pgfpathlineto{\pgfqpoint{2.840063in}{5.168912in}}%
\pgfpathlineto{\pgfqpoint{2.840063in}{5.084429in}}%
\pgfpathclose%
\pgfusepath{stroke}%
\end{pgfscope}%
\begin{pgfscope}%
\pgfpathrectangle{\pgfqpoint{1.109909in}{4.050000in}}{\pgfqpoint{5.400000in}{5.400000in}}%
\pgfusepath{clip}%
\pgfsetbuttcap%
\pgfsetroundjoin%
\pgfsetlinewidth{1.003750pt}%
\definecolor{currentstroke}{rgb}{0.541176,0.380392,0.658824}%
\pgfsetstrokecolor{currentstroke}%
\pgfsetstrokeopacity{0.900000}%
\pgfsetdash{}{0pt}%
\pgfpathmoveto{\pgfqpoint{3.384281in}{6.983714in}}%
\pgfpathlineto{\pgfqpoint{3.468763in}{6.983714in}}%
\pgfpathlineto{\pgfqpoint{3.468763in}{7.068197in}}%
\pgfpathlineto{\pgfqpoint{3.384281in}{7.068197in}}%
\pgfpathlineto{\pgfqpoint{3.384281in}{6.983714in}}%
\pgfpathclose%
\pgfusepath{stroke}%
\end{pgfscope}%
\begin{pgfscope}%
\pgfpathrectangle{\pgfqpoint{1.109909in}{4.050000in}}{\pgfqpoint{5.400000in}{5.400000in}}%
\pgfusepath{clip}%
\pgfsetbuttcap%
\pgfsetroundjoin%
\pgfsetlinewidth{1.003750pt}%
\definecolor{currentstroke}{rgb}{0.541176,0.380392,0.658824}%
\pgfsetstrokecolor{currentstroke}%
\pgfsetstrokeopacity{0.900000}%
\pgfsetdash{}{0pt}%
\pgfpathmoveto{\pgfqpoint{3.369680in}{6.258189in}}%
\pgfpathlineto{\pgfqpoint{3.454163in}{6.258189in}}%
\pgfpathlineto{\pgfqpoint{3.454163in}{6.342672in}}%
\pgfpathlineto{\pgfqpoint{3.369680in}{6.342672in}}%
\pgfpathlineto{\pgfqpoint{3.369680in}{6.258189in}}%
\pgfpathclose%
\pgfusepath{stroke}%
\end{pgfscope}%
\begin{pgfscope}%
\pgfpathrectangle{\pgfqpoint{1.109909in}{4.050000in}}{\pgfqpoint{5.400000in}{5.400000in}}%
\pgfusepath{clip}%
\pgfsetbuttcap%
\pgfsetroundjoin%
\pgfsetlinewidth{1.003750pt}%
\definecolor{currentstroke}{rgb}{0.768627,0.603922,0.000000}%
\pgfsetstrokecolor{currentstroke}%
\pgfsetstrokeopacity{0.900000}%
\pgfsetdash{}{0pt}%
\pgfpathmoveto{\pgfqpoint{3.447419in}{6.401781in}}%
\pgfpathlineto{\pgfqpoint{3.531902in}{6.401781in}}%
\pgfpathlineto{\pgfqpoint{3.531902in}{6.486264in}}%
\pgfpathlineto{\pgfqpoint{3.447419in}{6.486264in}}%
\pgfpathlineto{\pgfqpoint{3.447419in}{6.401781in}}%
\pgfpathclose%
\pgfusepath{stroke}%
\end{pgfscope}%
\begin{pgfscope}%
\pgfpathrectangle{\pgfqpoint{1.109909in}{4.050000in}}{\pgfqpoint{5.400000in}{5.400000in}}%
\pgfusepath{clip}%
\pgfsetbuttcap%
\pgfsetroundjoin%
\pgfsetlinewidth{1.003750pt}%
\definecolor{currentstroke}{rgb}{0.768627,0.603922,0.000000}%
\pgfsetstrokecolor{currentstroke}%
\pgfsetstrokeopacity{0.900000}%
\pgfsetdash{}{0pt}%
\pgfpathmoveto{\pgfqpoint{3.746269in}{6.820534in}}%
\pgfpathlineto{\pgfqpoint{3.830752in}{6.820534in}}%
\pgfpathlineto{\pgfqpoint{3.830752in}{6.905017in}}%
\pgfpathlineto{\pgfqpoint{3.746269in}{6.905017in}}%
\pgfpathlineto{\pgfqpoint{3.746269in}{6.820534in}}%
\pgfpathclose%
\pgfusepath{stroke}%
\end{pgfscope}%
\begin{pgfscope}%
\pgfpathrectangle{\pgfqpoint{1.109909in}{4.050000in}}{\pgfqpoint{5.400000in}{5.400000in}}%
\pgfusepath{clip}%
\pgfsetbuttcap%
\pgfsetroundjoin%
\pgfsetlinewidth{1.003750pt}%
\definecolor{currentstroke}{rgb}{0.200000,0.133333,0.533333}%
\pgfsetstrokecolor{currentstroke}%
\pgfsetstrokeopacity{0.900000}%
\pgfsetdash{}{0pt}%
\pgfpathmoveto{\pgfqpoint{2.642018in}{5.292716in}}%
\pgfpathlineto{\pgfqpoint{2.726500in}{5.292716in}}%
\pgfpathlineto{\pgfqpoint{2.726500in}{5.377199in}}%
\pgfpathlineto{\pgfqpoint{2.642018in}{5.377199in}}%
\pgfpathlineto{\pgfqpoint{2.642018in}{5.292716in}}%
\pgfpathclose%
\pgfusepath{stroke}%
\end{pgfscope}%
\begin{pgfscope}%
\pgfpathrectangle{\pgfqpoint{1.109909in}{4.050000in}}{\pgfqpoint{5.400000in}{5.400000in}}%
\pgfusepath{clip}%
\pgfsetbuttcap%
\pgfsetroundjoin%
\pgfsetlinewidth{1.003750pt}%
\definecolor{currentstroke}{rgb}{0.200000,0.133333,0.533333}%
\pgfsetstrokecolor{currentstroke}%
\pgfsetstrokeopacity{0.900000}%
\pgfsetdash{}{0pt}%
\pgfpathmoveto{\pgfqpoint{2.652444in}{5.503698in}}%
\pgfpathlineto{\pgfqpoint{2.736927in}{5.503698in}}%
\pgfpathlineto{\pgfqpoint{2.736927in}{5.588180in}}%
\pgfpathlineto{\pgfqpoint{2.652444in}{5.588180in}}%
\pgfpathlineto{\pgfqpoint{2.652444in}{5.503698in}}%
\pgfpathclose%
\pgfusepath{stroke}%
\end{pgfscope}%
\begin{pgfscope}%
\pgfpathrectangle{\pgfqpoint{1.109909in}{4.050000in}}{\pgfqpoint{5.400000in}{5.400000in}}%
\pgfusepath{clip}%
\pgfsetbuttcap%
\pgfsetroundjoin%
\pgfsetlinewidth{1.003750pt}%
\definecolor{currentstroke}{rgb}{0.200000,0.133333,0.533333}%
\pgfsetstrokecolor{currentstroke}%
\pgfsetstrokeopacity{0.900000}%
\pgfsetdash{}{0pt}%
\pgfpathmoveto{\pgfqpoint{3.077285in}{6.040599in}}%
\pgfpathlineto{\pgfqpoint{3.161768in}{6.040599in}}%
\pgfpathlineto{\pgfqpoint{3.161768in}{6.125082in}}%
\pgfpathlineto{\pgfqpoint{3.077285in}{6.125082in}}%
\pgfpathlineto{\pgfqpoint{3.077285in}{6.040599in}}%
\pgfpathclose%
\pgfusepath{stroke}%
\end{pgfscope}%
\begin{pgfscope}%
\pgfpathrectangle{\pgfqpoint{1.109909in}{4.050000in}}{\pgfqpoint{5.400000in}{5.400000in}}%
\pgfusepath{clip}%
\pgfsetbuttcap%
\pgfsetroundjoin%
\pgfsetlinewidth{1.003750pt}%
\definecolor{currentstroke}{rgb}{0.200000,0.133333,0.533333}%
\pgfsetstrokecolor{currentstroke}%
\pgfsetstrokeopacity{0.900000}%
\pgfsetdash{}{0pt}%
\pgfpathmoveto{\pgfqpoint{4.205223in}{7.201825in}}%
\pgfpathlineto{\pgfqpoint{4.289706in}{7.201825in}}%
\pgfpathlineto{\pgfqpoint{4.289706in}{7.286308in}}%
\pgfpathlineto{\pgfqpoint{4.205223in}{7.286308in}}%
\pgfpathlineto{\pgfqpoint{4.205223in}{7.201825in}}%
\pgfpathclose%
\pgfusepath{stroke}%
\end{pgfscope}%
\begin{pgfscope}%
\pgfpathrectangle{\pgfqpoint{1.109909in}{4.050000in}}{\pgfqpoint{5.400000in}{5.400000in}}%
\pgfusepath{clip}%
\pgfsetbuttcap%
\pgfsetroundjoin%
\pgfsetlinewidth{1.003750pt}%
\definecolor{currentstroke}{rgb}{0.200000,0.133333,0.533333}%
\pgfsetstrokecolor{currentstroke}%
\pgfsetstrokeopacity{0.900000}%
\pgfsetdash{}{0pt}%
\pgfpathmoveto{\pgfqpoint{4.001422in}{7.045685in}}%
\pgfpathlineto{\pgfqpoint{4.085905in}{7.045685in}}%
\pgfpathlineto{\pgfqpoint{4.085905in}{7.130167in}}%
\pgfpathlineto{\pgfqpoint{4.001422in}{7.130167in}}%
\pgfpathlineto{\pgfqpoint{4.001422in}{7.045685in}}%
\pgfpathclose%
\pgfusepath{stroke}%
\end{pgfscope}%
\begin{pgfscope}%
\pgfpathrectangle{\pgfqpoint{1.109909in}{4.050000in}}{\pgfqpoint{5.400000in}{5.400000in}}%
\pgfusepath{clip}%
\pgfsetbuttcap%
\pgfsetroundjoin%
\pgfsetlinewidth{1.003750pt}%
\definecolor{currentstroke}{rgb}{0.200000,0.133333,0.533333}%
\pgfsetstrokecolor{currentstroke}%
\pgfsetstrokeopacity{0.900000}%
\pgfsetdash{}{0pt}%
\pgfpathmoveto{\pgfqpoint{3.723240in}{6.716811in}}%
\pgfpathlineto{\pgfqpoint{3.807723in}{6.716811in}}%
\pgfpathlineto{\pgfqpoint{3.807723in}{6.801294in}}%
\pgfpathlineto{\pgfqpoint{3.723240in}{6.801294in}}%
\pgfpathlineto{\pgfqpoint{3.723240in}{6.716811in}}%
\pgfpathclose%
\pgfusepath{stroke}%
\end{pgfscope}%
\begin{pgfscope}%
\pgfpathrectangle{\pgfqpoint{1.109909in}{4.050000in}}{\pgfqpoint{5.400000in}{5.400000in}}%
\pgfusepath{clip}%
\pgfsetbuttcap%
\pgfsetroundjoin%
\definecolor{currentfill}{rgb}{0.176471,0.713725,0.643137}%
\pgfsetfillcolor{currentfill}%
\pgfsetfillopacity{0.900000}%
\pgfsetlinewidth{0.803000pt}%
\definecolor{currentstroke}{rgb}{0.176471,0.713725,0.643137}%
\pgfsetstrokecolor{currentstroke}%
\pgfsetstrokeopacity{0.900000}%
\pgfsetdash{}{0pt}%
\pgfsys@defobject{currentmarker}{\pgfqpoint{-0.042241in}{-0.042241in}}{\pgfqpoint{0.042241in}{0.042241in}}{%
\pgfpathmoveto{\pgfqpoint{-0.042241in}{-0.042241in}}%
\pgfpathlineto{\pgfqpoint{0.042241in}{-0.042241in}}%
\pgfpathlineto{\pgfqpoint{0.042241in}{0.042241in}}%
\pgfpathlineto{\pgfqpoint{-0.042241in}{0.042241in}}%
\pgfpathlineto{\pgfqpoint{-0.042241in}{-0.042241in}}%
\pgfpathclose%
\pgfusepath{stroke,fill}%
}%
\begin{pgfscope}%
\pgfsys@transformshift{5.535325in}{6.090324in}%
\pgfsys@useobject{currentmarker}{}%
\end{pgfscope}%
\end{pgfscope}%
\begin{pgfscope}%
\pgfpathrectangle{\pgfqpoint{1.109909in}{4.050000in}}{\pgfqpoint{5.400000in}{5.400000in}}%
\pgfusepath{clip}%
\pgfsetbuttcap%
\pgfsetroundjoin%
\pgfsetlinewidth{1.003750pt}%
\definecolor{currentstroke}{rgb}{0.176471,0.713725,0.643137}%
\pgfsetstrokecolor{currentstroke}%
\pgfsetstrokeopacity{0.900000}%
\pgfsetdash{}{0pt}%
\pgfpathmoveto{\pgfqpoint{5.141474in}{7.763158in}}%
\pgfpathlineto{\pgfqpoint{5.225957in}{7.763158in}}%
\pgfpathlineto{\pgfqpoint{5.225957in}{7.847641in}}%
\pgfpathlineto{\pgfqpoint{5.141474in}{7.847641in}}%
\pgfpathlineto{\pgfqpoint{5.141474in}{7.763158in}}%
\pgfpathclose%
\pgfusepath{stroke}%
\end{pgfscope}%
\begin{pgfscope}%
\pgfpathrectangle{\pgfqpoint{1.109909in}{4.050000in}}{\pgfqpoint{5.400000in}{5.400000in}}%
\pgfusepath{clip}%
\pgfsetbuttcap%
\pgfsetroundjoin%
\pgfsetlinewidth{1.003750pt}%
\definecolor{currentstroke}{rgb}{0.176471,0.713725,0.643137}%
\pgfsetstrokecolor{currentstroke}%
\pgfsetstrokeopacity{0.900000}%
\pgfsetdash{}{0pt}%
\pgfpathmoveto{\pgfqpoint{5.396797in}{8.156839in}}%
\pgfpathlineto{\pgfqpoint{5.481280in}{8.156839in}}%
\pgfpathlineto{\pgfqpoint{5.481280in}{8.241322in}}%
\pgfpathlineto{\pgfqpoint{5.396797in}{8.241322in}}%
\pgfpathlineto{\pgfqpoint{5.396797in}{8.156839in}}%
\pgfpathclose%
\pgfusepath{stroke}%
\end{pgfscope}%
\begin{pgfscope}%
\pgfpathrectangle{\pgfqpoint{1.109909in}{4.050000in}}{\pgfqpoint{5.400000in}{5.400000in}}%
\pgfusepath{clip}%
\pgfsetbuttcap%
\pgfsetroundjoin%
\pgfsetlinewidth{1.003750pt}%
\definecolor{currentstroke}{rgb}{0.176471,0.713725,0.643137}%
\pgfsetstrokecolor{currentstroke}%
\pgfsetstrokeopacity{0.900000}%
\pgfsetdash{}{0pt}%
\pgfpathmoveto{\pgfqpoint{4.946151in}{7.674753in}}%
\pgfpathlineto{\pgfqpoint{5.030634in}{7.674753in}}%
\pgfpathlineto{\pgfqpoint{5.030634in}{7.759235in}}%
\pgfpathlineto{\pgfqpoint{4.946151in}{7.759235in}}%
\pgfpathlineto{\pgfqpoint{4.946151in}{7.674753in}}%
\pgfpathclose%
\pgfusepath{stroke}%
\end{pgfscope}%
\begin{pgfscope}%
\pgfpathrectangle{\pgfqpoint{1.109909in}{4.050000in}}{\pgfqpoint{5.400000in}{5.400000in}}%
\pgfusepath{clip}%
\pgfsetbuttcap%
\pgfsetroundjoin%
\pgfsetlinewidth{1.003750pt}%
\definecolor{currentstroke}{rgb}{0.176471,0.713725,0.643137}%
\pgfsetstrokecolor{currentstroke}%
\pgfsetstrokeopacity{0.900000}%
\pgfsetdash{}{0pt}%
\pgfpathmoveto{\pgfqpoint{3.187133in}{6.574501in}}%
\pgfpathlineto{\pgfqpoint{3.271616in}{6.574501in}}%
\pgfpathlineto{\pgfqpoint{3.271616in}{6.658984in}}%
\pgfpathlineto{\pgfqpoint{3.187133in}{6.658984in}}%
\pgfpathlineto{\pgfqpoint{3.187133in}{6.574501in}}%
\pgfpathclose%
\pgfusepath{stroke}%
\end{pgfscope}%
\begin{pgfscope}%
\pgfsetbuttcap%
\pgfsetroundjoin%
\pgfsetlinewidth{1.003750pt}%
\definecolor{currentstroke}{rgb}{0.901961,0.196078,0.156863}%
\pgfsetstrokecolor{currentstroke}%
\pgfsetstrokeopacity{0.900000}%
\pgfsetdash{}{0pt}%
\pgfpathmoveto{\pgfqpoint{6.104902in}{9.407759in}}%
\pgfpathlineto{\pgfqpoint{6.189385in}{9.407759in}}%
\pgfpathlineto{\pgfqpoint{6.189385in}{9.492241in}}%
\pgfpathlineto{\pgfqpoint{6.104902in}{9.492241in}}%
\pgfpathlineto{\pgfqpoint{6.104902in}{9.407759in}}%
\pgfpathclose%
\pgfusepath{stroke}%
\end{pgfscope}%
\begin{pgfscope}%
\pgfpathrectangle{\pgfqpoint{1.109909in}{4.050000in}}{\pgfqpoint{5.400000in}{5.400000in}}%
\pgfusepath{clip}%
\pgfsetbuttcap%
\pgfsetroundjoin%
\definecolor{currentfill}{rgb}{0.498039,0.560784,0.678431}%
\pgfsetfillcolor{currentfill}%
\pgfsetfillopacity{0.900000}%
\pgfsetlinewidth{0.803000pt}%
\definecolor{currentstroke}{rgb}{0.498039,0.560784,0.678431}%
\pgfsetstrokecolor{currentstroke}%
\pgfsetstrokeopacity{0.900000}%
\pgfsetdash{}{0pt}%
\pgfsys@defobject{currentmarker}{\pgfqpoint{-0.042241in}{-0.042241in}}{\pgfqpoint{0.042241in}{0.042241in}}{%
\pgfpathmoveto{\pgfqpoint{-0.042241in}{-0.042241in}}%
\pgfpathlineto{\pgfqpoint{0.042241in}{-0.042241in}}%
\pgfpathlineto{\pgfqpoint{0.042241in}{0.042241in}}%
\pgfpathlineto{\pgfqpoint{-0.042241in}{0.042241in}}%
\pgfpathlineto{\pgfqpoint{-0.042241in}{-0.042241in}}%
\pgfpathclose%
\pgfusepath{stroke,fill}%
}%
\begin{pgfscope}%
\pgfsys@transformshift{5.060214in}{7.133796in}%
\pgfsys@useobject{currentmarker}{}%
\end{pgfscope}%
\end{pgfscope}%
\begin{pgfscope}%
\pgfsetbuttcap%
\pgfsetroundjoin%
\pgfsetlinewidth{1.003750pt}%
\definecolor{currentstroke}{rgb}{0.498039,0.560784,0.678431}%
\pgfsetstrokecolor{currentstroke}%
\pgfsetstrokeopacity{0.900000}%
\pgfsetdash{}{0pt}%
\pgfpathmoveto{\pgfqpoint{1.067668in}{4.053476in}}%
\pgfpathlineto{\pgfqpoint{1.152150in}{4.053476in}}%
\pgfpathlineto{\pgfqpoint{1.152150in}{4.137959in}}%
\pgfpathlineto{\pgfqpoint{1.067668in}{4.137959in}}%
\pgfpathlineto{\pgfqpoint{1.067668in}{4.053476in}}%
\pgfpathclose%
\pgfusepath{stroke}%
\end{pgfscope}%
\begin{pgfscope}%
\pgfpathrectangle{\pgfqpoint{1.109909in}{4.050000in}}{\pgfqpoint{5.400000in}{5.400000in}}%
\pgfusepath{clip}%
\pgfsetbuttcap%
\pgfsetroundjoin%
\pgfsetlinewidth{1.003750pt}%
\definecolor{currentstroke}{rgb}{0.709804,0.121569,0.270588}%
\pgfsetstrokecolor{currentstroke}%
\pgfsetstrokeopacity{0.900000}%
\pgfsetdash{}{0pt}%
\pgfpathmoveto{\pgfqpoint{6.220372in}{9.309057in}}%
\pgfpathlineto{\pgfqpoint{6.304855in}{9.309057in}}%
\pgfpathlineto{\pgfqpoint{6.304855in}{9.393540in}}%
\pgfpathlineto{\pgfqpoint{6.220372in}{9.393540in}}%
\pgfpathlineto{\pgfqpoint{6.220372in}{9.309057in}}%
\pgfpathclose%
\pgfusepath{stroke}%
\end{pgfscope}%
\begin{pgfscope}%
\pgfsetbuttcap%
\pgfsetmiterjoin%
\definecolor{currentfill}{rgb}{1.000000,1.000000,1.000000}%
\pgfsetfillcolor{currentfill}%
\pgfsetfillopacity{0.970000}%
\pgfsetlinewidth{0.702625pt}%
\definecolor{currentstroke}{rgb}{0.796078,0.827451,0.850980}%
\pgfsetstrokecolor{currentstroke}%
\pgfsetstrokeopacity{0.970000}%
\pgfsetdash{}{0pt}%
\pgfpathmoveto{\pgfqpoint{1.298909in}{7.579333in}}%
\pgfpathlineto{\pgfqpoint{3.325784in}{7.579333in}}%
\pgfpathquadraticcurveto{\pgfqpoint{3.403701in}{7.579333in}}{\pgfqpoint{3.403701in}{7.657250in}}%
\pgfpathlineto{\pgfqpoint{3.403701in}{8.316000in}}%
\pgfpathquadraticcurveto{\pgfqpoint{3.403701in}{8.393917in}}{\pgfqpoint{3.325784in}{8.393917in}}%
\pgfpathlineto{\pgfqpoint{1.298909in}{8.393917in}}%
\pgfpathquadraticcurveto{\pgfqpoint{1.220992in}{8.393917in}}{\pgfqpoint{1.220992in}{8.316000in}}%
\pgfpathlineto{\pgfqpoint{1.220992in}{7.657250in}}%
\pgfpathquadraticcurveto{\pgfqpoint{1.220992in}{7.579333in}}{\pgfqpoint{1.298909in}{7.579333in}}%
\pgfpathlineto{\pgfqpoint{1.298909in}{7.579333in}}%
\pgfpathclose%
\pgfusepath{stroke,fill}%
\end{pgfscope}%
\begin{pgfscope}%
\pgfsetbuttcap%
\pgfsetroundjoin%
\definecolor{currentfill}{rgb}{0.090196,0.168627,0.227451}%
\pgfsetfillcolor{currentfill}%
\pgfsetlinewidth{0.000000pt}%
\definecolor{currentstroke}{rgb}{0.090196,0.168627,0.227451}%
\pgfsetstrokecolor{currentstroke}%
\pgfsetdash{}{0pt}%
\pgfpathmoveto{\pgfqpoint{1.374723in}{8.266105in}}%
\pgfpathlineto{\pgfqpoint{1.374723in}{8.252470in}}%
\pgfpathquadraticcurveto{\pgfqpoint{1.366776in}{8.256277in}}{\pgfqpoint{1.359715in}{8.258137in}}%
\pgfpathquadraticcurveto{\pgfqpoint{1.352654in}{8.260018in}}{\pgfqpoint{1.346080in}{8.260018in}}%
\pgfpathquadraticcurveto{\pgfqpoint{1.334680in}{8.260018in}}{\pgfqpoint{1.328482in}{8.255591in}}%
\pgfpathquadraticcurveto{\pgfqpoint{1.322284in}{8.251164in}}{\pgfqpoint{1.322284in}{8.242996in}}%
\pgfpathquadraticcurveto{\pgfqpoint{1.322284in}{8.236134in}}{\pgfqpoint{1.326401in}{8.232637in}}%
\pgfpathquadraticcurveto{\pgfqpoint{1.330518in}{8.229161in}}{\pgfqpoint{1.342007in}{8.227014in}}%
\pgfpathlineto{\pgfqpoint{1.350440in}{8.225288in}}%
\pgfpathquadraticcurveto{\pgfqpoint{1.366068in}{8.222299in}}{\pgfqpoint{1.373505in}{8.214796in}}%
\pgfpathquadraticcurveto{\pgfqpoint{1.380943in}{8.207292in}}{\pgfqpoint{1.380943in}{8.194719in}}%
\pgfpathquadraticcurveto{\pgfqpoint{1.380943in}{8.179689in}}{\pgfqpoint{1.370871in}{8.171941in}}%
\pgfpathquadraticcurveto{\pgfqpoint{1.360822in}{8.164194in}}{\pgfqpoint{1.341387in}{8.164194in}}%
\pgfpathquadraticcurveto{\pgfqpoint{1.334060in}{8.164194in}}{\pgfqpoint{1.325781in}{8.165854in}}%
\pgfpathquadraticcurveto{\pgfqpoint{1.317525in}{8.167514in}}{\pgfqpoint{1.308671in}{8.170768in}}%
\pgfpathlineto{\pgfqpoint{1.308671in}{8.185156in}}%
\pgfpathquadraticcurveto{\pgfqpoint{1.317171in}{8.180397in}}{\pgfqpoint{1.325339in}{8.177962in}}%
\pgfpathquadraticcurveto{\pgfqpoint{1.333507in}{8.175549in}}{\pgfqpoint{1.341387in}{8.175549in}}%
\pgfpathquadraticcurveto{\pgfqpoint{1.353340in}{8.175549in}}{\pgfqpoint{1.359848in}{8.180242in}}%
\pgfpathquadraticcurveto{\pgfqpoint{1.366356in}{8.184957in}}{\pgfqpoint{1.366356in}{8.193678in}}%
\pgfpathquadraticcurveto{\pgfqpoint{1.366356in}{8.201271in}}{\pgfqpoint{1.361685in}{8.205565in}}%
\pgfpathquadraticcurveto{\pgfqpoint{1.357015in}{8.209859in}}{\pgfqpoint{1.346367in}{8.212007in}}%
\pgfpathlineto{\pgfqpoint{1.337845in}{8.213667in}}%
\pgfpathquadraticcurveto{\pgfqpoint{1.322218in}{8.216766in}}{\pgfqpoint{1.315223in}{8.223406in}}%
\pgfpathquadraticcurveto{\pgfqpoint{1.308250in}{8.230047in}}{\pgfqpoint{1.308250in}{8.241889in}}%
\pgfpathquadraticcurveto{\pgfqpoint{1.308250in}{8.255591in}}{\pgfqpoint{1.317901in}{8.263471in}}%
\pgfpathquadraticcurveto{\pgfqpoint{1.327552in}{8.271352in}}{\pgfqpoint{1.344486in}{8.271352in}}%
\pgfpathquadraticcurveto{\pgfqpoint{1.351768in}{8.271352in}}{\pgfqpoint{1.359295in}{8.270023in}}%
\pgfpathquadraticcurveto{\pgfqpoint{1.366843in}{8.268717in}}{\pgfqpoint{1.374723in}{8.266105in}}%
\pgfpathlineto{\pgfqpoint{1.374723in}{8.266105in}}%
\pgfpathclose%
\pgfpathmoveto{\pgfqpoint{1.402193in}{8.243682in}}%
\pgfpathlineto{\pgfqpoint{1.414921in}{8.243682in}}%
\pgfpathlineto{\pgfqpoint{1.414921in}{8.166208in}}%
\pgfpathlineto{\pgfqpoint{1.402193in}{8.166208in}}%
\pgfpathlineto{\pgfqpoint{1.402193in}{8.243682in}}%
\pgfpathlineto{\pgfqpoint{1.402193in}{8.243682in}}%
\pgfpathclose%
\pgfpathmoveto{\pgfqpoint{1.402193in}{8.273853in}}%
\pgfpathlineto{\pgfqpoint{1.414921in}{8.273853in}}%
\pgfpathlineto{\pgfqpoint{1.414921in}{8.257716in}}%
\pgfpathlineto{\pgfqpoint{1.402193in}{8.257716in}}%
\pgfpathlineto{\pgfqpoint{1.402193in}{8.273853in}}%
\pgfpathlineto{\pgfqpoint{1.402193in}{8.273853in}}%
\pgfpathclose%
\pgfpathmoveto{\pgfqpoint{1.492528in}{8.205853in}}%
\pgfpathquadraticcurveto{\pgfqpoint{1.492528in}{8.219688in}}{\pgfqpoint{1.486817in}{8.227280in}}%
\pgfpathquadraticcurveto{\pgfqpoint{1.481128in}{8.234895in}}{\pgfqpoint{1.470813in}{8.234895in}}%
\pgfpathquadraticcurveto{\pgfqpoint{1.460586in}{8.234895in}}{\pgfqpoint{1.454875in}{8.227280in}}%
\pgfpathquadraticcurveto{\pgfqpoint{1.449164in}{8.219688in}}{\pgfqpoint{1.449164in}{8.205853in}}%
\pgfpathquadraticcurveto{\pgfqpoint{1.449164in}{8.192085in}}{\pgfqpoint{1.454875in}{8.184470in}}%
\pgfpathquadraticcurveto{\pgfqpoint{1.460586in}{8.176855in}}{\pgfqpoint{1.470813in}{8.176855in}}%
\pgfpathquadraticcurveto{\pgfqpoint{1.481128in}{8.176855in}}{\pgfqpoint{1.486817in}{8.184470in}}%
\pgfpathquadraticcurveto{\pgfqpoint{1.492528in}{8.192085in}}{\pgfqpoint{1.492528in}{8.205853in}}%
\pgfpathlineto{\pgfqpoint{1.492528in}{8.205853in}}%
\pgfpathclose%
\pgfpathmoveto{\pgfqpoint{1.505255in}{8.175815in}}%
\pgfpathquadraticcurveto{\pgfqpoint{1.505255in}{8.156048in}}{\pgfqpoint{1.496468in}{8.146397in}}%
\pgfpathquadraticcurveto{\pgfqpoint{1.487702in}{8.136746in}}{\pgfqpoint{1.469573in}{8.136746in}}%
\pgfpathquadraticcurveto{\pgfqpoint{1.462866in}{8.136746in}}{\pgfqpoint{1.456912in}{8.137742in}}%
\pgfpathquadraticcurveto{\pgfqpoint{1.450957in}{8.138738in}}{\pgfqpoint{1.445357in}{8.140819in}}%
\pgfpathlineto{\pgfqpoint{1.445357in}{8.153193in}}%
\pgfpathquadraticcurveto{\pgfqpoint{1.450957in}{8.150160in}}{\pgfqpoint{1.456425in}{8.148721in}}%
\pgfpathquadraticcurveto{\pgfqpoint{1.461892in}{8.147260in}}{\pgfqpoint{1.467559in}{8.147260in}}%
\pgfpathquadraticcurveto{\pgfqpoint{1.480087in}{8.147260in}}{\pgfqpoint{1.486308in}{8.153790in}}%
\pgfpathquadraticcurveto{\pgfqpoint{1.492528in}{8.160320in}}{\pgfqpoint{1.492528in}{8.173535in}}%
\pgfpathlineto{\pgfqpoint{1.492528in}{8.179844in}}%
\pgfpathquadraticcurveto{\pgfqpoint{1.488587in}{8.172982in}}{\pgfqpoint{1.482434in}{8.169595in}}%
\pgfpathquadraticcurveto{\pgfqpoint{1.476280in}{8.166208in}}{\pgfqpoint{1.467692in}{8.166208in}}%
\pgfpathquadraticcurveto{\pgfqpoint{1.453459in}{8.166208in}}{\pgfqpoint{1.444737in}{8.177055in}}%
\pgfpathquadraticcurveto{\pgfqpoint{1.436016in}{8.187923in}}{\pgfqpoint{1.436016in}{8.205853in}}%
\pgfpathquadraticcurveto{\pgfqpoint{1.436016in}{8.223827in}}{\pgfqpoint{1.444737in}{8.234673in}}%
\pgfpathquadraticcurveto{\pgfqpoint{1.453459in}{8.245542in}}{\pgfqpoint{1.467692in}{8.245542in}}%
\pgfpathquadraticcurveto{\pgfqpoint{1.476280in}{8.245542in}}{\pgfqpoint{1.482434in}{8.242155in}}%
\pgfpathquadraticcurveto{\pgfqpoint{1.488587in}{8.238768in}}{\pgfqpoint{1.492528in}{8.231928in}}%
\pgfpathlineto{\pgfqpoint{1.492528in}{8.243682in}}%
\pgfpathlineto{\pgfqpoint{1.505255in}{8.243682in}}%
\pgfpathlineto{\pgfqpoint{1.505255in}{8.175815in}}%
\pgfpathlineto{\pgfqpoint{1.505255in}{8.175815in}}%
\pgfpathclose%
\pgfpathmoveto{\pgfqpoint{1.595900in}{8.212980in}}%
\pgfpathlineto{\pgfqpoint{1.595900in}{8.166208in}}%
\pgfpathlineto{\pgfqpoint{1.583172in}{8.166208in}}%
\pgfpathlineto{\pgfqpoint{1.583172in}{8.212560in}}%
\pgfpathquadraticcurveto{\pgfqpoint{1.583172in}{8.223561in}}{\pgfqpoint{1.578878in}{8.229007in}}%
\pgfpathquadraticcurveto{\pgfqpoint{1.574584in}{8.234474in}}{\pgfqpoint{1.566017in}{8.234474in}}%
\pgfpathquadraticcurveto{\pgfqpoint{1.555702in}{8.234474in}}{\pgfqpoint{1.549748in}{8.227900in}}%
\pgfpathquadraticcurveto{\pgfqpoint{1.543793in}{8.221348in}}{\pgfqpoint{1.543793in}{8.209992in}}%
\pgfpathlineto{\pgfqpoint{1.543793in}{8.166208in}}%
\pgfpathlineto{\pgfqpoint{1.530999in}{8.166208in}}%
\pgfpathlineto{\pgfqpoint{1.530999in}{8.243682in}}%
\pgfpathlineto{\pgfqpoint{1.543793in}{8.243682in}}%
\pgfpathlineto{\pgfqpoint{1.543793in}{8.231641in}}%
\pgfpathquadraticcurveto{\pgfqpoint{1.548375in}{8.238635in}}{\pgfqpoint{1.554551in}{8.242089in}}%
\pgfpathquadraticcurveto{\pgfqpoint{1.560749in}{8.245542in}}{\pgfqpoint{1.568850in}{8.245542in}}%
\pgfpathquadraticcurveto{\pgfqpoint{1.582198in}{8.245542in}}{\pgfqpoint{1.589038in}{8.237285in}}%
\pgfpathquadraticcurveto{\pgfqpoint{1.595900in}{8.229029in}}{\pgfqpoint{1.595900in}{8.212980in}}%
\pgfpathlineto{\pgfqpoint{1.595900in}{8.212980in}}%
\pgfpathclose%
\pgfpathmoveto{\pgfqpoint{1.614848in}{8.210678in}}%
\pgfpathlineto{\pgfqpoint{1.652124in}{8.210678in}}%
\pgfpathlineto{\pgfqpoint{1.652124in}{8.199345in}}%
\pgfpathlineto{\pgfqpoint{1.614848in}{8.199345in}}%
\pgfpathlineto{\pgfqpoint{1.614848in}{8.210678in}}%
\pgfpathlineto{\pgfqpoint{1.614848in}{8.210678in}}%
\pgfpathclose%
\pgfpathmoveto{\pgfqpoint{1.717268in}{8.231796in}}%
\pgfpathquadraticcurveto{\pgfqpoint{1.715121in}{8.233035in}}{\pgfqpoint{1.712598in}{8.233611in}}%
\pgfpathquadraticcurveto{\pgfqpoint{1.710074in}{8.234208in}}{\pgfqpoint{1.707042in}{8.234208in}}%
\pgfpathquadraticcurveto{\pgfqpoint{1.696240in}{8.234208in}}{\pgfqpoint{1.690462in}{8.227191in}}%
\pgfpathquadraticcurveto{\pgfqpoint{1.684685in}{8.220174in}}{\pgfqpoint{1.684685in}{8.207026in}}%
\pgfpathlineto{\pgfqpoint{1.684685in}{8.166208in}}%
\pgfpathlineto{\pgfqpoint{1.671891in}{8.166208in}}%
\pgfpathlineto{\pgfqpoint{1.671891in}{8.243682in}}%
\pgfpathlineto{\pgfqpoint{1.684685in}{8.243682in}}%
\pgfpathlineto{\pgfqpoint{1.684685in}{8.231641in}}%
\pgfpathquadraticcurveto{\pgfqpoint{1.688714in}{8.238702in}}{\pgfqpoint{1.695133in}{8.242111in}}%
\pgfpathquadraticcurveto{\pgfqpoint{1.701574in}{8.245542in}}{\pgfqpoint{1.710783in}{8.245542in}}%
\pgfpathquadraticcurveto{\pgfqpoint{1.712089in}{8.245542in}}{\pgfqpoint{1.713683in}{8.245365in}}%
\pgfpathquadraticcurveto{\pgfqpoint{1.715276in}{8.245210in}}{\pgfqpoint{1.717202in}{8.244855in}}%
\pgfpathlineto{\pgfqpoint{1.717268in}{8.231796in}}%
\pgfpathlineto{\pgfqpoint{1.717268in}{8.231796in}}%
\pgfpathclose%
\pgfpathmoveto{\pgfqpoint{1.793768in}{8.208133in}}%
\pgfpathlineto{\pgfqpoint{1.793768in}{8.201913in}}%
\pgfpathlineto{\pgfqpoint{1.735242in}{8.201913in}}%
\pgfpathquadraticcurveto{\pgfqpoint{1.736084in}{8.188764in}}{\pgfqpoint{1.743167in}{8.181880in}}%
\pgfpathquadraticcurveto{\pgfqpoint{1.750250in}{8.174996in}}{\pgfqpoint{1.762912in}{8.174996in}}%
\pgfpathquadraticcurveto{\pgfqpoint{1.770239in}{8.174996in}}{\pgfqpoint{1.777123in}{8.176789in}}%
\pgfpathquadraticcurveto{\pgfqpoint{1.784007in}{8.178582in}}{\pgfqpoint{1.790802in}{8.182190in}}%
\pgfpathlineto{\pgfqpoint{1.790802in}{8.170148in}}%
\pgfpathquadraticcurveto{\pgfqpoint{1.783940in}{8.167249in}}{\pgfqpoint{1.776746in}{8.165721in}}%
\pgfpathquadraticcurveto{\pgfqpoint{1.769552in}{8.164194in}}{\pgfqpoint{1.762159in}{8.164194in}}%
\pgfpathquadraticcurveto{\pgfqpoint{1.743610in}{8.164194in}}{\pgfqpoint{1.732785in}{8.174974in}}%
\pgfpathquadraticcurveto{\pgfqpoint{1.721961in}{8.185776in}}{\pgfqpoint{1.721961in}{8.204193in}}%
\pgfpathquadraticcurveto{\pgfqpoint{1.721961in}{8.223207in}}{\pgfqpoint{1.732232in}{8.234363in}}%
\pgfpathquadraticcurveto{\pgfqpoint{1.742503in}{8.245542in}}{\pgfqpoint{1.759946in}{8.245542in}}%
\pgfpathquadraticcurveto{\pgfqpoint{1.775573in}{8.245542in}}{\pgfqpoint{1.784671in}{8.235470in}}%
\pgfpathquadraticcurveto{\pgfqpoint{1.793768in}{8.225421in}}{\pgfqpoint{1.793768in}{8.208133in}}%
\pgfpathlineto{\pgfqpoint{1.793768in}{8.208133in}}%
\pgfpathclose%
\pgfpathmoveto{\pgfqpoint{1.781041in}{8.211874in}}%
\pgfpathquadraticcurveto{\pgfqpoint{1.780908in}{8.222299in}}{\pgfqpoint{1.775197in}{8.228520in}}%
\pgfpathquadraticcurveto{\pgfqpoint{1.769486in}{8.234762in}}{\pgfqpoint{1.760078in}{8.234762in}}%
\pgfpathquadraticcurveto{\pgfqpoint{1.749431in}{8.234762in}}{\pgfqpoint{1.743034in}{8.228741in}}%
\pgfpathquadraticcurveto{\pgfqpoint{1.736637in}{8.222720in}}{\pgfqpoint{1.735663in}{8.211785in}}%
\pgfpathlineto{\pgfqpoint{1.781041in}{8.211874in}}%
\pgfpathlineto{\pgfqpoint{1.781041in}{8.211874in}}%
\pgfpathclose%
\pgfpathmoveto{\pgfqpoint{1.805545in}{8.243682in}}%
\pgfpathlineto{\pgfqpoint{1.819025in}{8.243682in}}%
\pgfpathlineto{\pgfqpoint{1.843241in}{8.178671in}}%
\pgfpathlineto{\pgfqpoint{1.867457in}{8.243682in}}%
\pgfpathlineto{\pgfqpoint{1.880938in}{8.243682in}}%
\pgfpathlineto{\pgfqpoint{1.851874in}{8.166208in}}%
\pgfpathlineto{\pgfqpoint{1.834586in}{8.166208in}}%
\pgfpathlineto{\pgfqpoint{1.805545in}{8.243682in}}%
\pgfpathlineto{\pgfqpoint{1.805545in}{8.243682in}}%
\pgfpathclose%
\pgfpathmoveto{\pgfqpoint{1.964787in}{8.208133in}}%
\pgfpathlineto{\pgfqpoint{1.964787in}{8.201913in}}%
\pgfpathlineto{\pgfqpoint{1.906261in}{8.201913in}}%
\pgfpathquadraticcurveto{\pgfqpoint{1.907102in}{8.188764in}}{\pgfqpoint{1.914185in}{8.181880in}}%
\pgfpathquadraticcurveto{\pgfqpoint{1.921268in}{8.174996in}}{\pgfqpoint{1.933930in}{8.174996in}}%
\pgfpathquadraticcurveto{\pgfqpoint{1.941257in}{8.174996in}}{\pgfqpoint{1.948141in}{8.176789in}}%
\pgfpathquadraticcurveto{\pgfqpoint{1.955025in}{8.178582in}}{\pgfqpoint{1.961821in}{8.182190in}}%
\pgfpathlineto{\pgfqpoint{1.961821in}{8.170148in}}%
\pgfpathquadraticcurveto{\pgfqpoint{1.954959in}{8.167249in}}{\pgfqpoint{1.947765in}{8.165721in}}%
\pgfpathquadraticcurveto{\pgfqpoint{1.940571in}{8.164194in}}{\pgfqpoint{1.933177in}{8.164194in}}%
\pgfpathquadraticcurveto{\pgfqpoint{1.914628in}{8.164194in}}{\pgfqpoint{1.903804in}{8.174974in}}%
\pgfpathquadraticcurveto{\pgfqpoint{1.892979in}{8.185776in}}{\pgfqpoint{1.892979in}{8.204193in}}%
\pgfpathquadraticcurveto{\pgfqpoint{1.892979in}{8.223207in}}{\pgfqpoint{1.903250in}{8.234363in}}%
\pgfpathquadraticcurveto{\pgfqpoint{1.913521in}{8.245542in}}{\pgfqpoint{1.930964in}{8.245542in}}%
\pgfpathquadraticcurveto{\pgfqpoint{1.946591in}{8.245542in}}{\pgfqpoint{1.955689in}{8.235470in}}%
\pgfpathquadraticcurveto{\pgfqpoint{1.964787in}{8.225421in}}{\pgfqpoint{1.964787in}{8.208133in}}%
\pgfpathlineto{\pgfqpoint{1.964787in}{8.208133in}}%
\pgfpathclose%
\pgfpathmoveto{\pgfqpoint{1.952059in}{8.211874in}}%
\pgfpathquadraticcurveto{\pgfqpoint{1.951926in}{8.222299in}}{\pgfqpoint{1.946215in}{8.228520in}}%
\pgfpathquadraticcurveto{\pgfqpoint{1.940504in}{8.234762in}}{\pgfqpoint{1.931097in}{8.234762in}}%
\pgfpathquadraticcurveto{\pgfqpoint{1.920449in}{8.234762in}}{\pgfqpoint{1.914052in}{8.228741in}}%
\pgfpathquadraticcurveto{\pgfqpoint{1.907655in}{8.222720in}}{\pgfqpoint{1.906681in}{8.211785in}}%
\pgfpathlineto{\pgfqpoint{1.952059in}{8.211874in}}%
\pgfpathlineto{\pgfqpoint{1.952059in}{8.211874in}}%
\pgfpathclose%
\pgfpathmoveto{\pgfqpoint{2.030573in}{8.231796in}}%
\pgfpathquadraticcurveto{\pgfqpoint{2.028426in}{8.233035in}}{\pgfqpoint{2.025903in}{8.233611in}}%
\pgfpathquadraticcurveto{\pgfqpoint{2.023379in}{8.234208in}}{\pgfqpoint{2.020347in}{8.234208in}}%
\pgfpathquadraticcurveto{\pgfqpoint{2.009545in}{8.234208in}}{\pgfqpoint{2.003767in}{8.227191in}}%
\pgfpathquadraticcurveto{\pgfqpoint{1.997990in}{8.220174in}}{\pgfqpoint{1.997990in}{8.207026in}}%
\pgfpathlineto{\pgfqpoint{1.997990in}{8.166208in}}%
\pgfpathlineto{\pgfqpoint{1.985196in}{8.166208in}}%
\pgfpathlineto{\pgfqpoint{1.985196in}{8.243682in}}%
\pgfpathlineto{\pgfqpoint{1.997990in}{8.243682in}}%
\pgfpathlineto{\pgfqpoint{1.997990in}{8.231641in}}%
\pgfpathquadraticcurveto{\pgfqpoint{2.002018in}{8.238702in}}{\pgfqpoint{2.008438in}{8.242111in}}%
\pgfpathquadraticcurveto{\pgfqpoint{2.014879in}{8.245542in}}{\pgfqpoint{2.024087in}{8.245542in}}%
\pgfpathquadraticcurveto{\pgfqpoint{2.025393in}{8.245542in}}{\pgfqpoint{2.026987in}{8.245365in}}%
\pgfpathquadraticcurveto{\pgfqpoint{2.028581in}{8.245210in}}{\pgfqpoint{2.030507in}{8.244855in}}%
\pgfpathlineto{\pgfqpoint{2.030573in}{8.231796in}}%
\pgfpathlineto{\pgfqpoint{2.030573in}{8.231796in}}%
\pgfpathclose%
\pgfpathmoveto{\pgfqpoint{2.093305in}{8.241402in}}%
\pgfpathlineto{\pgfqpoint{2.093305in}{8.229361in}}%
\pgfpathquadraticcurveto{\pgfqpoint{2.087926in}{8.232128in}}{\pgfqpoint{2.082104in}{8.233500in}}%
\pgfpathquadraticcurveto{\pgfqpoint{2.076305in}{8.234895in}}{\pgfqpoint{2.070063in}{8.234895in}}%
\pgfpathquadraticcurveto{\pgfqpoint{2.060589in}{8.234895in}}{\pgfqpoint{2.055852in}{8.231995in}}%
\pgfpathquadraticcurveto{\pgfqpoint{2.051115in}{8.229095in}}{\pgfqpoint{2.051115in}{8.223273in}}%
\pgfpathquadraticcurveto{\pgfqpoint{2.051115in}{8.218846in}}{\pgfqpoint{2.054502in}{8.216323in}}%
\pgfpathquadraticcurveto{\pgfqpoint{2.057888in}{8.213799in}}{\pgfqpoint{2.068137in}{8.211520in}}%
\pgfpathlineto{\pgfqpoint{2.072498in}{8.210546in}}%
\pgfpathquadraticcurveto{\pgfqpoint{2.086045in}{8.207646in}}{\pgfqpoint{2.091755in}{8.202355in}}%
\pgfpathquadraticcurveto{\pgfqpoint{2.097466in}{8.197065in}}{\pgfqpoint{2.097466in}{8.187591in}}%
\pgfpathquadraticcurveto{\pgfqpoint{2.097466in}{8.176789in}}{\pgfqpoint{2.088922in}{8.170480in}}%
\pgfpathquadraticcurveto{\pgfqpoint{2.080378in}{8.164194in}}{\pgfqpoint{2.065436in}{8.164194in}}%
\pgfpathquadraticcurveto{\pgfqpoint{2.059216in}{8.164194in}}{\pgfqpoint{2.052465in}{8.165411in}}%
\pgfpathquadraticcurveto{\pgfqpoint{2.045714in}{8.166629in}}{\pgfqpoint{2.038254in}{8.169042in}}%
\pgfpathlineto{\pgfqpoint{2.038254in}{8.182190in}}%
\pgfpathquadraticcurveto{\pgfqpoint{2.045315in}{8.178516in}}{\pgfqpoint{2.052155in}{8.176678in}}%
\pgfpathquadraticcurveto{\pgfqpoint{2.058995in}{8.174863in}}{\pgfqpoint{2.065724in}{8.174863in}}%
\pgfpathquadraticcurveto{\pgfqpoint{2.074711in}{8.174863in}}{\pgfqpoint{2.079537in}{8.177940in}}%
\pgfpathquadraticcurveto{\pgfqpoint{2.084384in}{8.181017in}}{\pgfqpoint{2.084384in}{8.186617in}}%
\pgfpathquadraticcurveto{\pgfqpoint{2.084384in}{8.191797in}}{\pgfqpoint{2.080887in}{8.194564in}}%
\pgfpathquadraticcurveto{\pgfqpoint{2.077412in}{8.197331in}}{\pgfqpoint{2.065569in}{8.199898in}}%
\pgfpathlineto{\pgfqpoint{2.061142in}{8.200939in}}%
\pgfpathquadraticcurveto{\pgfqpoint{2.049322in}{8.203418in}}{\pgfqpoint{2.044054in}{8.208576in}}%
\pgfpathquadraticcurveto{\pgfqpoint{2.038808in}{8.213733in}}{\pgfqpoint{2.038808in}{8.222720in}}%
\pgfpathquadraticcurveto{\pgfqpoint{2.038808in}{8.233655in}}{\pgfqpoint{2.046555in}{8.239587in}}%
\pgfpathquadraticcurveto{\pgfqpoint{2.054302in}{8.245542in}}{\pgfqpoint{2.068558in}{8.245542in}}%
\pgfpathquadraticcurveto{\pgfqpoint{2.075597in}{8.245542in}}{\pgfqpoint{2.081817in}{8.244501in}}%
\pgfpathquadraticcurveto{\pgfqpoint{2.088059in}{8.243483in}}{\pgfqpoint{2.093305in}{8.241402in}}%
\pgfpathlineto{\pgfqpoint{2.093305in}{8.241402in}}%
\pgfpathclose%
\pgfpathmoveto{\pgfqpoint{2.152938in}{8.205145in}}%
\pgfpathquadraticcurveto{\pgfqpoint{2.137509in}{8.205145in}}{\pgfqpoint{2.131555in}{8.201625in}}%
\pgfpathquadraticcurveto{\pgfqpoint{2.125600in}{8.198105in}}{\pgfqpoint{2.125600in}{8.189583in}}%
\pgfpathquadraticcurveto{\pgfqpoint{2.125600in}{8.182810in}}{\pgfqpoint{2.130072in}{8.178826in}}%
\pgfpathquadraticcurveto{\pgfqpoint{2.134543in}{8.174863in}}{\pgfqpoint{2.142202in}{8.174863in}}%
\pgfpathquadraticcurveto{\pgfqpoint{2.152805in}{8.174863in}}{\pgfqpoint{2.159202in}{8.182367in}}%
\pgfpathquadraticcurveto{\pgfqpoint{2.165599in}{8.189871in}}{\pgfqpoint{2.165599in}{8.202311in}}%
\pgfpathlineto{\pgfqpoint{2.165599in}{8.205145in}}%
\pgfpathlineto{\pgfqpoint{2.152938in}{8.205145in}}%
\pgfpathlineto{\pgfqpoint{2.152938in}{8.205145in}}%
\pgfpathclose%
\pgfpathmoveto{\pgfqpoint{2.178327in}{8.210413in}}%
\pgfpathlineto{\pgfqpoint{2.178327in}{8.166208in}}%
\pgfpathlineto{\pgfqpoint{2.165599in}{8.166208in}}%
\pgfpathlineto{\pgfqpoint{2.165599in}{8.177962in}}%
\pgfpathquadraticcurveto{\pgfqpoint{2.161239in}{8.170923in}}{\pgfqpoint{2.154731in}{8.167559in}}%
\pgfpathquadraticcurveto{\pgfqpoint{2.148223in}{8.164194in}}{\pgfqpoint{2.138815in}{8.164194in}}%
\pgfpathquadraticcurveto{\pgfqpoint{2.126929in}{8.164194in}}{\pgfqpoint{2.119890in}{8.170879in}}%
\pgfpathquadraticcurveto{\pgfqpoint{2.112873in}{8.177564in}}{\pgfqpoint{2.112873in}{8.188764in}}%
\pgfpathquadraticcurveto{\pgfqpoint{2.112873in}{8.201824in}}{\pgfqpoint{2.121616in}{8.208465in}}%
\pgfpathquadraticcurveto{\pgfqpoint{2.130382in}{8.215105in}}{\pgfqpoint{2.147736in}{8.215105in}}%
\pgfpathlineto{\pgfqpoint{2.165599in}{8.215105in}}%
\pgfpathlineto{\pgfqpoint{2.165599in}{8.216367in}}%
\pgfpathquadraticcurveto{\pgfqpoint{2.165599in}{8.225155in}}{\pgfqpoint{2.159822in}{8.229958in}}%
\pgfpathquadraticcurveto{\pgfqpoint{2.154045in}{8.234762in}}{\pgfqpoint{2.143597in}{8.234762in}}%
\pgfpathquadraticcurveto{\pgfqpoint{2.136956in}{8.234762in}}{\pgfqpoint{2.130647in}{8.233168in}}%
\pgfpathquadraticcurveto{\pgfqpoint{2.124361in}{8.231574in}}{\pgfqpoint{2.118561in}{8.228387in}}%
\pgfpathlineto{\pgfqpoint{2.118561in}{8.240163in}}%
\pgfpathquadraticcurveto{\pgfqpoint{2.125534in}{8.242863in}}{\pgfqpoint{2.132108in}{8.244191in}}%
\pgfpathquadraticcurveto{\pgfqpoint{2.138683in}{8.245542in}}{\pgfqpoint{2.144903in}{8.245542in}}%
\pgfpathquadraticcurveto{\pgfqpoint{2.161725in}{8.245542in}}{\pgfqpoint{2.170026in}{8.236820in}}%
\pgfpathquadraticcurveto{\pgfqpoint{2.178327in}{8.228121in}}{\pgfqpoint{2.178327in}{8.210413in}}%
\pgfpathlineto{\pgfqpoint{2.178327in}{8.210413in}}%
\pgfpathclose%
\pgfpathmoveto{\pgfqpoint{2.204535in}{8.273853in}}%
\pgfpathlineto{\pgfqpoint{2.217263in}{8.273853in}}%
\pgfpathlineto{\pgfqpoint{2.217263in}{8.166208in}}%
\pgfpathlineto{\pgfqpoint{2.204535in}{8.166208in}}%
\pgfpathlineto{\pgfqpoint{2.204535in}{8.273853in}}%
\pgfpathlineto{\pgfqpoint{2.204535in}{8.273853in}}%
\pgfpathclose%
\pgfpathmoveto{\pgfqpoint{2.293276in}{8.241402in}}%
\pgfpathlineto{\pgfqpoint{2.293276in}{8.229361in}}%
\pgfpathquadraticcurveto{\pgfqpoint{2.287897in}{8.232128in}}{\pgfqpoint{2.282076in}{8.233500in}}%
\pgfpathquadraticcurveto{\pgfqpoint{2.276276in}{8.234895in}}{\pgfqpoint{2.270034in}{8.234895in}}%
\pgfpathquadraticcurveto{\pgfqpoint{2.260560in}{8.234895in}}{\pgfqpoint{2.255823in}{8.231995in}}%
\pgfpathquadraticcurveto{\pgfqpoint{2.251086in}{8.229095in}}{\pgfqpoint{2.251086in}{8.223273in}}%
\pgfpathquadraticcurveto{\pgfqpoint{2.251086in}{8.218846in}}{\pgfqpoint{2.254473in}{8.216323in}}%
\pgfpathquadraticcurveto{\pgfqpoint{2.257860in}{8.213799in}}{\pgfqpoint{2.268108in}{8.211520in}}%
\pgfpathlineto{\pgfqpoint{2.272469in}{8.210546in}}%
\pgfpathquadraticcurveto{\pgfqpoint{2.286016in}{8.207646in}}{\pgfqpoint{2.291727in}{8.202355in}}%
\pgfpathquadraticcurveto{\pgfqpoint{2.297438in}{8.197065in}}{\pgfqpoint{2.297438in}{8.187591in}}%
\pgfpathquadraticcurveto{\pgfqpoint{2.297438in}{8.176789in}}{\pgfqpoint{2.288893in}{8.170480in}}%
\pgfpathquadraticcurveto{\pgfqpoint{2.280349in}{8.164194in}}{\pgfqpoint{2.265408in}{8.164194in}}%
\pgfpathquadraticcurveto{\pgfqpoint{2.259188in}{8.164194in}}{\pgfqpoint{2.252436in}{8.165411in}}%
\pgfpathquadraticcurveto{\pgfqpoint{2.245685in}{8.166629in}}{\pgfqpoint{2.238225in}{8.169042in}}%
\pgfpathlineto{\pgfqpoint{2.238225in}{8.182190in}}%
\pgfpathquadraticcurveto{\pgfqpoint{2.245287in}{8.178516in}}{\pgfqpoint{2.252127in}{8.176678in}}%
\pgfpathquadraticcurveto{\pgfqpoint{2.258966in}{8.174863in}}{\pgfqpoint{2.265696in}{8.174863in}}%
\pgfpathquadraticcurveto{\pgfqpoint{2.274683in}{8.174863in}}{\pgfqpoint{2.279508in}{8.177940in}}%
\pgfpathquadraticcurveto{\pgfqpoint{2.284356in}{8.181017in}}{\pgfqpoint{2.284356in}{8.186617in}}%
\pgfpathquadraticcurveto{\pgfqpoint{2.284356in}{8.191797in}}{\pgfqpoint{2.280858in}{8.194564in}}%
\pgfpathquadraticcurveto{\pgfqpoint{2.277383in}{8.197331in}}{\pgfqpoint{2.265541in}{8.199898in}}%
\pgfpathlineto{\pgfqpoint{2.261114in}{8.200939in}}%
\pgfpathquadraticcurveto{\pgfqpoint{2.249293in}{8.203418in}}{\pgfqpoint{2.244025in}{8.208576in}}%
\pgfpathquadraticcurveto{\pgfqpoint{2.238779in}{8.213733in}}{\pgfqpoint{2.238779in}{8.222720in}}%
\pgfpathquadraticcurveto{\pgfqpoint{2.238779in}{8.233655in}}{\pgfqpoint{2.246526in}{8.239587in}}%
\pgfpathquadraticcurveto{\pgfqpoint{2.254274in}{8.245542in}}{\pgfqpoint{2.268529in}{8.245542in}}%
\pgfpathquadraticcurveto{\pgfqpoint{2.275568in}{8.245542in}}{\pgfqpoint{2.281788in}{8.244501in}}%
\pgfpathquadraticcurveto{\pgfqpoint{2.288030in}{8.243483in}}{\pgfqpoint{2.293276in}{8.241402in}}%
\pgfpathlineto{\pgfqpoint{2.293276in}{8.241402in}}%
\pgfpathclose%
\pgfpathmoveto{\pgfqpoint{2.320946in}{8.183784in}}%
\pgfpathlineto{\pgfqpoint{2.335533in}{8.183784in}}%
\pgfpathlineto{\pgfqpoint{2.335533in}{8.166208in}}%
\pgfpathlineto{\pgfqpoint{2.320946in}{8.166208in}}%
\pgfpathlineto{\pgfqpoint{2.320946in}{8.183784in}}%
\pgfpathlineto{\pgfqpoint{2.320946in}{8.183784in}}%
\pgfpathclose%
\pgfpathmoveto{\pgfqpoint{2.320946in}{8.239454in}}%
\pgfpathlineto{\pgfqpoint{2.335533in}{8.239454in}}%
\pgfpathlineto{\pgfqpoint{2.335533in}{8.221901in}}%
\pgfpathlineto{\pgfqpoint{2.320946in}{8.221901in}}%
\pgfpathlineto{\pgfqpoint{2.320946in}{8.239454in}}%
\pgfpathlineto{\pgfqpoint{2.320946in}{8.239454in}}%
\pgfpathclose%
\pgfpathmoveto{\pgfqpoint{2.414667in}{8.177962in}}%
\pgfpathlineto{\pgfqpoint{2.437489in}{8.177962in}}%
\pgfpathlineto{\pgfqpoint{2.437489in}{8.256764in}}%
\pgfpathlineto{\pgfqpoint{2.412653in}{8.251784in}}%
\pgfpathlineto{\pgfqpoint{2.412653in}{8.264512in}}%
\pgfpathlineto{\pgfqpoint{2.437356in}{8.269492in}}%
\pgfpathlineto{\pgfqpoint{2.451323in}{8.269492in}}%
\pgfpathlineto{\pgfqpoint{2.451323in}{8.177962in}}%
\pgfpathlineto{\pgfqpoint{2.474145in}{8.177962in}}%
\pgfpathlineto{\pgfqpoint{2.474145in}{8.166208in}}%
\pgfpathlineto{\pgfqpoint{2.414667in}{8.166208in}}%
\pgfpathlineto{\pgfqpoint{2.414667in}{8.177962in}}%
\pgfpathlineto{\pgfqpoint{2.414667in}{8.177962in}}%
\pgfpathclose%
\pgfpathmoveto{\pgfqpoint{2.514409in}{8.177962in}}%
\pgfpathlineto{\pgfqpoint{2.563173in}{8.177962in}}%
\pgfpathlineto{\pgfqpoint{2.563173in}{8.166208in}}%
\pgfpathlineto{\pgfqpoint{2.497608in}{8.166208in}}%
\pgfpathlineto{\pgfqpoint{2.497608in}{8.177962in}}%
\pgfpathquadraticcurveto{\pgfqpoint{2.505555in}{8.186197in}}{\pgfqpoint{2.519279in}{8.200053in}}%
\pgfpathquadraticcurveto{\pgfqpoint{2.533025in}{8.213932in}}{\pgfqpoint{2.536545in}{8.217961in}}%
\pgfpathquadraticcurveto{\pgfqpoint{2.543252in}{8.225487in}}{\pgfqpoint{2.545908in}{8.230711in}}%
\pgfpathquadraticcurveto{\pgfqpoint{2.548586in}{8.235935in}}{\pgfqpoint{2.548586in}{8.240982in}}%
\pgfpathquadraticcurveto{\pgfqpoint{2.548586in}{8.249216in}}{\pgfqpoint{2.542809in}{8.254396in}}%
\pgfpathquadraticcurveto{\pgfqpoint{2.537031in}{8.259598in}}{\pgfqpoint{2.527757in}{8.259598in}}%
\pgfpathquadraticcurveto{\pgfqpoint{2.521183in}{8.259598in}}{\pgfqpoint{2.513878in}{8.257318in}}%
\pgfpathquadraticcurveto{\pgfqpoint{2.506595in}{8.255038in}}{\pgfqpoint{2.498295in}{8.250389in}}%
\pgfpathlineto{\pgfqpoint{2.498295in}{8.264512in}}%
\pgfpathquadraticcurveto{\pgfqpoint{2.506728in}{8.267898in}}{\pgfqpoint{2.514055in}{8.269625in}}%
\pgfpathquadraticcurveto{\pgfqpoint{2.521404in}{8.271352in}}{\pgfqpoint{2.527491in}{8.271352in}}%
\pgfpathquadraticcurveto{\pgfqpoint{2.543539in}{8.271352in}}{\pgfqpoint{2.553080in}{8.263316in}}%
\pgfpathquadraticcurveto{\pgfqpoint{2.562620in}{8.255303in}}{\pgfqpoint{2.562620in}{8.241889in}}%
\pgfpathquadraticcurveto{\pgfqpoint{2.562620in}{8.235514in}}{\pgfqpoint{2.560229in}{8.229803in}}%
\pgfpathquadraticcurveto{\pgfqpoint{2.557861in}{8.224115in}}{\pgfqpoint{2.551552in}{8.216367in}}%
\pgfpathquadraticcurveto{\pgfqpoint{2.549826in}{8.214353in}}{\pgfqpoint{2.540551in}{8.204768in}}%
\pgfpathquadraticcurveto{\pgfqpoint{2.531298in}{8.195184in}}{\pgfqpoint{2.514409in}{8.177962in}}%
\pgfpathlineto{\pgfqpoint{2.514409in}{8.177962in}}%
\pgfpathclose%
\pgfpathmoveto{\pgfqpoint{2.592503in}{8.183784in}}%
\pgfpathlineto{\pgfqpoint{2.607112in}{8.183784in}}%
\pgfpathlineto{\pgfqpoint{2.607112in}{8.166208in}}%
\pgfpathlineto{\pgfqpoint{2.592503in}{8.166208in}}%
\pgfpathlineto{\pgfqpoint{2.592503in}{8.183784in}}%
\pgfpathlineto{\pgfqpoint{2.592503in}{8.183784in}}%
\pgfpathclose%
\pgfpathmoveto{\pgfqpoint{2.649568in}{8.177962in}}%
\pgfpathlineto{\pgfqpoint{2.698332in}{8.177962in}}%
\pgfpathlineto{\pgfqpoint{2.698332in}{8.166208in}}%
\pgfpathlineto{\pgfqpoint{2.632767in}{8.166208in}}%
\pgfpathlineto{\pgfqpoint{2.632767in}{8.177962in}}%
\pgfpathquadraticcurveto{\pgfqpoint{2.640714in}{8.186197in}}{\pgfqpoint{2.654438in}{8.200053in}}%
\pgfpathquadraticcurveto{\pgfqpoint{2.668184in}{8.213932in}}{\pgfqpoint{2.671703in}{8.217961in}}%
\pgfpathquadraticcurveto{\pgfqpoint{2.678410in}{8.225487in}}{\pgfqpoint{2.681067in}{8.230711in}}%
\pgfpathquadraticcurveto{\pgfqpoint{2.683745in}{8.235935in}}{\pgfqpoint{2.683745in}{8.240982in}}%
\pgfpathquadraticcurveto{\pgfqpoint{2.683745in}{8.249216in}}{\pgfqpoint{2.677968in}{8.254396in}}%
\pgfpathquadraticcurveto{\pgfqpoint{2.672190in}{8.259598in}}{\pgfqpoint{2.662916in}{8.259598in}}%
\pgfpathquadraticcurveto{\pgfqpoint{2.656341in}{8.259598in}}{\pgfqpoint{2.649037in}{8.257318in}}%
\pgfpathquadraticcurveto{\pgfqpoint{2.641754in}{8.255038in}}{\pgfqpoint{2.633453in}{8.250389in}}%
\pgfpathlineto{\pgfqpoint{2.633453in}{8.264512in}}%
\pgfpathquadraticcurveto{\pgfqpoint{2.641887in}{8.267898in}}{\pgfqpoint{2.649214in}{8.269625in}}%
\pgfpathquadraticcurveto{\pgfqpoint{2.656563in}{8.271352in}}{\pgfqpoint{2.662650in}{8.271352in}}%
\pgfpathquadraticcurveto{\pgfqpoint{2.678698in}{8.271352in}}{\pgfqpoint{2.688239in}{8.263316in}}%
\pgfpathquadraticcurveto{\pgfqpoint{2.697779in}{8.255303in}}{\pgfqpoint{2.697779in}{8.241889in}}%
\pgfpathquadraticcurveto{\pgfqpoint{2.697779in}{8.235514in}}{\pgfqpoint{2.695388in}{8.229803in}}%
\pgfpathquadraticcurveto{\pgfqpoint{2.693020in}{8.224115in}}{\pgfqpoint{2.686711in}{8.216367in}}%
\pgfpathquadraticcurveto{\pgfqpoint{2.684985in}{8.214353in}}{\pgfqpoint{2.675710in}{8.204768in}}%
\pgfpathquadraticcurveto{\pgfqpoint{2.666457in}{8.195184in}}{\pgfqpoint{2.649568in}{8.177962in}}%
\pgfpathlineto{\pgfqpoint{2.649568in}{8.177962in}}%
\pgfpathclose%
\pgfpathmoveto{\pgfqpoint{2.815517in}{8.211652in}}%
\pgfpathquadraticcurveto{\pgfqpoint{2.809496in}{8.211652in}}{\pgfqpoint{2.806065in}{8.206539in}}%
\pgfpathquadraticcurveto{\pgfqpoint{2.802656in}{8.201426in}}{\pgfqpoint{2.802656in}{8.192284in}}%
\pgfpathquadraticcurveto{\pgfqpoint{2.802656in}{8.183297in}}{\pgfqpoint{2.806065in}{8.178139in}}%
\pgfpathquadraticcurveto{\pgfqpoint{2.809496in}{8.172982in}}{\pgfqpoint{2.815517in}{8.172982in}}%
\pgfpathquadraticcurveto{\pgfqpoint{2.821405in}{8.172982in}}{\pgfqpoint{2.824814in}{8.178139in}}%
\pgfpathquadraticcurveto{\pgfqpoint{2.828245in}{8.183297in}}{\pgfqpoint{2.828245in}{8.192284in}}%
\pgfpathquadraticcurveto{\pgfqpoint{2.828245in}{8.201359in}}{\pgfqpoint{2.824814in}{8.206495in}}%
\pgfpathquadraticcurveto{\pgfqpoint{2.821405in}{8.211652in}}{\pgfqpoint{2.815517in}{8.211652in}}%
\pgfpathlineto{\pgfqpoint{2.815517in}{8.211652in}}%
\pgfpathclose%
\pgfpathmoveto{\pgfqpoint{2.815517in}{8.220440in}}%
\pgfpathquadraticcurveto{\pgfqpoint{2.826452in}{8.220440in}}{\pgfqpoint{2.832871in}{8.212826in}}%
\pgfpathquadraticcurveto{\pgfqpoint{2.839313in}{8.205233in}}{\pgfqpoint{2.839313in}{8.192284in}}%
\pgfpathquadraticcurveto{\pgfqpoint{2.839313in}{8.179357in}}{\pgfqpoint{2.832849in}{8.171764in}}%
\pgfpathquadraticcurveto{\pgfqpoint{2.826386in}{8.164194in}}{\pgfqpoint{2.815517in}{8.164194in}}%
\pgfpathquadraticcurveto{\pgfqpoint{2.804449in}{8.164194in}}{\pgfqpoint{2.798008in}{8.171764in}}%
\pgfpathquadraticcurveto{\pgfqpoint{2.791589in}{8.179357in}}{\pgfqpoint{2.791589in}{8.192284in}}%
\pgfpathquadraticcurveto{\pgfqpoint{2.791589in}{8.205299in}}{\pgfqpoint{2.798052in}{8.212870in}}%
\pgfpathquadraticcurveto{\pgfqpoint{2.804516in}{8.220440in}}{\pgfqpoint{2.815517in}{8.220440in}}%
\pgfpathlineto{\pgfqpoint{2.815517in}{8.220440in}}%
\pgfpathclose%
\pgfpathmoveto{\pgfqpoint{2.744130in}{8.262564in}}%
\pgfpathquadraticcurveto{\pgfqpoint{2.738176in}{8.262564in}}{\pgfqpoint{2.734745in}{8.257406in}}%
\pgfpathquadraticcurveto{\pgfqpoint{2.731336in}{8.252271in}}{\pgfqpoint{2.731336in}{8.243262in}}%
\pgfpathquadraticcurveto{\pgfqpoint{2.731336in}{8.234142in}}{\pgfqpoint{2.734723in}{8.229007in}}%
\pgfpathquadraticcurveto{\pgfqpoint{2.738110in}{8.223893in}}{\pgfqpoint{2.744130in}{8.223893in}}%
\pgfpathquadraticcurveto{\pgfqpoint{2.750151in}{8.223893in}}{\pgfqpoint{2.753560in}{8.229007in}}%
\pgfpathquadraticcurveto{\pgfqpoint{2.756991in}{8.234142in}}{\pgfqpoint{2.756991in}{8.243262in}}%
\pgfpathquadraticcurveto{\pgfqpoint{2.756991in}{8.252182in}}{\pgfqpoint{2.753538in}{8.257362in}}%
\pgfpathquadraticcurveto{\pgfqpoint{2.750085in}{8.262564in}}{\pgfqpoint{2.744130in}{8.262564in}}%
\pgfpathlineto{\pgfqpoint{2.744130in}{8.262564in}}%
\pgfpathclose%
\pgfpathmoveto{\pgfqpoint{2.806597in}{8.271352in}}%
\pgfpathlineto{\pgfqpoint{2.817664in}{8.271352in}}%
\pgfpathlineto{\pgfqpoint{2.753051in}{8.164194in}}%
\pgfpathlineto{\pgfqpoint{2.741983in}{8.164194in}}%
\pgfpathlineto{\pgfqpoint{2.806597in}{8.271352in}}%
\pgfpathlineto{\pgfqpoint{2.806597in}{8.271352in}}%
\pgfpathclose%
\pgfpathmoveto{\pgfqpoint{2.744130in}{8.271352in}}%
\pgfpathquadraticcurveto{\pgfqpoint{2.755065in}{8.271352in}}{\pgfqpoint{2.761551in}{8.263781in}}%
\pgfpathquadraticcurveto{\pgfqpoint{2.768059in}{8.256211in}}{\pgfqpoint{2.768059in}{8.243262in}}%
\pgfpathquadraticcurveto{\pgfqpoint{2.768059in}{8.230202in}}{\pgfqpoint{2.761595in}{8.222654in}}%
\pgfpathquadraticcurveto{\pgfqpoint{2.755132in}{8.215105in}}{\pgfqpoint{2.744130in}{8.215105in}}%
\pgfpathquadraticcurveto{\pgfqpoint{2.733129in}{8.215105in}}{\pgfqpoint{2.726732in}{8.222676in}}%
\pgfpathquadraticcurveto{\pgfqpoint{2.720335in}{8.230268in}}{\pgfqpoint{2.720335in}{8.243262in}}%
\pgfpathquadraticcurveto{\pgfqpoint{2.720335in}{8.256145in}}{\pgfqpoint{2.726754in}{8.263737in}}%
\pgfpathquadraticcurveto{\pgfqpoint{2.733196in}{8.271352in}}{\pgfqpoint{2.744130in}{8.271352in}}%
\pgfpathlineto{\pgfqpoint{2.744130in}{8.271352in}}%
\pgfpathclose%
\pgfusepath{fill}%
\end{pgfscope}%
\begin{pgfscope}%
\pgfsetbuttcap%
\pgfsetroundjoin%
\definecolor{currentfill}{rgb}{0.090196,0.168627,0.227451}%
\pgfsetfillcolor{currentfill}%
\pgfsetlinewidth{0.000000pt}%
\definecolor{currentstroke}{rgb}{0.090196,0.168627,0.227451}%
\pgfsetstrokecolor{currentstroke}%
\pgfsetdash{}{0pt}%
\pgfpathmoveto{\pgfqpoint{1.374723in}{8.046522in}}%
\pgfpathlineto{\pgfqpoint{1.374723in}{8.032887in}}%
\pgfpathquadraticcurveto{\pgfqpoint{1.366776in}{8.036694in}}{\pgfqpoint{1.359715in}{8.038553in}}%
\pgfpathquadraticcurveto{\pgfqpoint{1.352654in}{8.040435in}}{\pgfqpoint{1.346080in}{8.040435in}}%
\pgfpathquadraticcurveto{\pgfqpoint{1.334680in}{8.040435in}}{\pgfqpoint{1.328482in}{8.036008in}}%
\pgfpathquadraticcurveto{\pgfqpoint{1.322284in}{8.031581in}}{\pgfqpoint{1.322284in}{8.023413in}}%
\pgfpathquadraticcurveto{\pgfqpoint{1.322284in}{8.016551in}}{\pgfqpoint{1.326401in}{8.013053in}}%
\pgfpathquadraticcurveto{\pgfqpoint{1.330518in}{8.009578in}}{\pgfqpoint{1.342007in}{8.007431in}}%
\pgfpathlineto{\pgfqpoint{1.350440in}{8.005704in}}%
\pgfpathquadraticcurveto{\pgfqpoint{1.366068in}{8.002716in}}{\pgfqpoint{1.373505in}{7.995212in}}%
\pgfpathquadraticcurveto{\pgfqpoint{1.380943in}{7.987708in}}{\pgfqpoint{1.380943in}{7.975135in}}%
\pgfpathquadraticcurveto{\pgfqpoint{1.380943in}{7.960105in}}{\pgfqpoint{1.370871in}{7.952358in}}%
\pgfpathquadraticcurveto{\pgfqpoint{1.360822in}{7.944611in}}{\pgfqpoint{1.341387in}{7.944611in}}%
\pgfpathquadraticcurveto{\pgfqpoint{1.334060in}{7.944611in}}{\pgfqpoint{1.325781in}{7.946271in}}%
\pgfpathquadraticcurveto{\pgfqpoint{1.317525in}{7.947931in}}{\pgfqpoint{1.308671in}{7.951185in}}%
\pgfpathlineto{\pgfqpoint{1.308671in}{7.965573in}}%
\pgfpathquadraticcurveto{\pgfqpoint{1.317171in}{7.960814in}}{\pgfqpoint{1.325339in}{7.958379in}}%
\pgfpathquadraticcurveto{\pgfqpoint{1.333507in}{7.955966in}}{\pgfqpoint{1.341387in}{7.955966in}}%
\pgfpathquadraticcurveto{\pgfqpoint{1.353340in}{7.955966in}}{\pgfqpoint{1.359848in}{7.960659in}}%
\pgfpathquadraticcurveto{\pgfqpoint{1.366356in}{7.965374in}}{\pgfqpoint{1.366356in}{7.974095in}}%
\pgfpathquadraticcurveto{\pgfqpoint{1.366356in}{7.981688in}}{\pgfqpoint{1.361685in}{7.985982in}}%
\pgfpathquadraticcurveto{\pgfqpoint{1.357015in}{7.990276in}}{\pgfqpoint{1.346367in}{7.992423in}}%
\pgfpathlineto{\pgfqpoint{1.337845in}{7.994083in}}%
\pgfpathquadraticcurveto{\pgfqpoint{1.322218in}{7.997182in}}{\pgfqpoint{1.315223in}{8.003823in}}%
\pgfpathquadraticcurveto{\pgfqpoint{1.308250in}{8.010464in}}{\pgfqpoint{1.308250in}{8.022306in}}%
\pgfpathquadraticcurveto{\pgfqpoint{1.308250in}{8.036008in}}{\pgfqpoint{1.317901in}{8.043888in}}%
\pgfpathquadraticcurveto{\pgfqpoint{1.327552in}{8.051768in}}{\pgfqpoint{1.344486in}{8.051768in}}%
\pgfpathquadraticcurveto{\pgfqpoint{1.351768in}{8.051768in}}{\pgfqpoint{1.359295in}{8.050440in}}%
\pgfpathquadraticcurveto{\pgfqpoint{1.366843in}{8.049134in}}{\pgfqpoint{1.374723in}{8.046522in}}%
\pgfpathlineto{\pgfqpoint{1.374723in}{8.046522in}}%
\pgfpathclose%
\pgfpathmoveto{\pgfqpoint{1.402193in}{8.024099in}}%
\pgfpathlineto{\pgfqpoint{1.414921in}{8.024099in}}%
\pgfpathlineto{\pgfqpoint{1.414921in}{7.946625in}}%
\pgfpathlineto{\pgfqpoint{1.402193in}{7.946625in}}%
\pgfpathlineto{\pgfqpoint{1.402193in}{8.024099in}}%
\pgfpathlineto{\pgfqpoint{1.402193in}{8.024099in}}%
\pgfpathclose%
\pgfpathmoveto{\pgfqpoint{1.402193in}{8.054270in}}%
\pgfpathlineto{\pgfqpoint{1.414921in}{8.054270in}}%
\pgfpathlineto{\pgfqpoint{1.414921in}{8.038133in}}%
\pgfpathlineto{\pgfqpoint{1.402193in}{8.038133in}}%
\pgfpathlineto{\pgfqpoint{1.402193in}{8.054270in}}%
\pgfpathlineto{\pgfqpoint{1.402193in}{8.054270in}}%
\pgfpathclose%
\pgfpathmoveto{\pgfqpoint{1.492528in}{7.986270in}}%
\pgfpathquadraticcurveto{\pgfqpoint{1.492528in}{8.000104in}}{\pgfqpoint{1.486817in}{8.007697in}}%
\pgfpathquadraticcurveto{\pgfqpoint{1.481128in}{8.015311in}}{\pgfqpoint{1.470813in}{8.015311in}}%
\pgfpathquadraticcurveto{\pgfqpoint{1.460586in}{8.015311in}}{\pgfqpoint{1.454875in}{8.007697in}}%
\pgfpathquadraticcurveto{\pgfqpoint{1.449164in}{8.000104in}}{\pgfqpoint{1.449164in}{7.986270in}}%
\pgfpathquadraticcurveto{\pgfqpoint{1.449164in}{7.972501in}}{\pgfqpoint{1.454875in}{7.964887in}}%
\pgfpathquadraticcurveto{\pgfqpoint{1.460586in}{7.957272in}}{\pgfqpoint{1.470813in}{7.957272in}}%
\pgfpathquadraticcurveto{\pgfqpoint{1.481128in}{7.957272in}}{\pgfqpoint{1.486817in}{7.964887in}}%
\pgfpathquadraticcurveto{\pgfqpoint{1.492528in}{7.972501in}}{\pgfqpoint{1.492528in}{7.986270in}}%
\pgfpathlineto{\pgfqpoint{1.492528in}{7.986270in}}%
\pgfpathclose%
\pgfpathmoveto{\pgfqpoint{1.505255in}{7.956232in}}%
\pgfpathquadraticcurveto{\pgfqpoint{1.505255in}{7.936465in}}{\pgfqpoint{1.496468in}{7.926814in}}%
\pgfpathquadraticcurveto{\pgfqpoint{1.487702in}{7.917163in}}{\pgfqpoint{1.469573in}{7.917163in}}%
\pgfpathquadraticcurveto{\pgfqpoint{1.462866in}{7.917163in}}{\pgfqpoint{1.456912in}{7.918159in}}%
\pgfpathquadraticcurveto{\pgfqpoint{1.450957in}{7.919155in}}{\pgfqpoint{1.445357in}{7.921236in}}%
\pgfpathlineto{\pgfqpoint{1.445357in}{7.933609in}}%
\pgfpathquadraticcurveto{\pgfqpoint{1.450957in}{7.930577in}}{\pgfqpoint{1.456425in}{7.929138in}}%
\pgfpathquadraticcurveto{\pgfqpoint{1.461892in}{7.927677in}}{\pgfqpoint{1.467559in}{7.927677in}}%
\pgfpathquadraticcurveto{\pgfqpoint{1.480087in}{7.927677in}}{\pgfqpoint{1.486308in}{7.934207in}}%
\pgfpathquadraticcurveto{\pgfqpoint{1.492528in}{7.940737in}}{\pgfqpoint{1.492528in}{7.953952in}}%
\pgfpathlineto{\pgfqpoint{1.492528in}{7.960260in}}%
\pgfpathquadraticcurveto{\pgfqpoint{1.488587in}{7.953398in}}{\pgfqpoint{1.482434in}{7.950012in}}%
\pgfpathquadraticcurveto{\pgfqpoint{1.476280in}{7.946625in}}{\pgfqpoint{1.467692in}{7.946625in}}%
\pgfpathquadraticcurveto{\pgfqpoint{1.453459in}{7.946625in}}{\pgfqpoint{1.444737in}{7.957471in}}%
\pgfpathquadraticcurveto{\pgfqpoint{1.436016in}{7.968340in}}{\pgfqpoint{1.436016in}{7.986270in}}%
\pgfpathquadraticcurveto{\pgfqpoint{1.436016in}{8.004243in}}{\pgfqpoint{1.444737in}{8.015090in}}%
\pgfpathquadraticcurveto{\pgfqpoint{1.453459in}{8.025958in}}{\pgfqpoint{1.467692in}{8.025958in}}%
\pgfpathquadraticcurveto{\pgfqpoint{1.476280in}{8.025958in}}{\pgfqpoint{1.482434in}{8.022572in}}%
\pgfpathquadraticcurveto{\pgfqpoint{1.488587in}{8.019185in}}{\pgfqpoint{1.492528in}{8.012345in}}%
\pgfpathlineto{\pgfqpoint{1.492528in}{8.024099in}}%
\pgfpathlineto{\pgfqpoint{1.505255in}{8.024099in}}%
\pgfpathlineto{\pgfqpoint{1.505255in}{7.956232in}}%
\pgfpathlineto{\pgfqpoint{1.505255in}{7.956232in}}%
\pgfpathclose%
\pgfpathmoveto{\pgfqpoint{1.595900in}{7.993397in}}%
\pgfpathlineto{\pgfqpoint{1.595900in}{7.946625in}}%
\pgfpathlineto{\pgfqpoint{1.583172in}{7.946625in}}%
\pgfpathlineto{\pgfqpoint{1.583172in}{7.992977in}}%
\pgfpathquadraticcurveto{\pgfqpoint{1.583172in}{8.003978in}}{\pgfqpoint{1.578878in}{8.009423in}}%
\pgfpathquadraticcurveto{\pgfqpoint{1.574584in}{8.014891in}}{\pgfqpoint{1.566017in}{8.014891in}}%
\pgfpathquadraticcurveto{\pgfqpoint{1.555702in}{8.014891in}}{\pgfqpoint{1.549748in}{8.008316in}}%
\pgfpathquadraticcurveto{\pgfqpoint{1.543793in}{8.001764in}}{\pgfqpoint{1.543793in}{7.990409in}}%
\pgfpathlineto{\pgfqpoint{1.543793in}{7.946625in}}%
\pgfpathlineto{\pgfqpoint{1.530999in}{7.946625in}}%
\pgfpathlineto{\pgfqpoint{1.530999in}{8.024099in}}%
\pgfpathlineto{\pgfqpoint{1.543793in}{8.024099in}}%
\pgfpathlineto{\pgfqpoint{1.543793in}{8.012057in}}%
\pgfpathquadraticcurveto{\pgfqpoint{1.548375in}{8.019052in}}{\pgfqpoint{1.554551in}{8.022505in}}%
\pgfpathquadraticcurveto{\pgfqpoint{1.560749in}{8.025958in}}{\pgfqpoint{1.568850in}{8.025958in}}%
\pgfpathquadraticcurveto{\pgfqpoint{1.582198in}{8.025958in}}{\pgfqpoint{1.589038in}{8.017702in}}%
\pgfpathquadraticcurveto{\pgfqpoint{1.595900in}{8.009445in}}{\pgfqpoint{1.595900in}{7.993397in}}%
\pgfpathlineto{\pgfqpoint{1.595900in}{7.993397in}}%
\pgfpathclose%
\pgfpathmoveto{\pgfqpoint{1.621267in}{8.024099in}}%
\pgfpathlineto{\pgfqpoint{1.633995in}{8.024099in}}%
\pgfpathlineto{\pgfqpoint{1.633995in}{7.946625in}}%
\pgfpathlineto{\pgfqpoint{1.621267in}{7.946625in}}%
\pgfpathlineto{\pgfqpoint{1.621267in}{8.024099in}}%
\pgfpathlineto{\pgfqpoint{1.621267in}{8.024099in}}%
\pgfpathclose%
\pgfpathmoveto{\pgfqpoint{1.621267in}{8.054270in}}%
\pgfpathlineto{\pgfqpoint{1.633995in}{8.054270in}}%
\pgfpathlineto{\pgfqpoint{1.633995in}{8.038133in}}%
\pgfpathlineto{\pgfqpoint{1.621267in}{8.038133in}}%
\pgfpathlineto{\pgfqpoint{1.621267in}{8.054270in}}%
\pgfpathlineto{\pgfqpoint{1.621267in}{8.054270in}}%
\pgfpathclose%
\pgfpathmoveto{\pgfqpoint{1.723223in}{8.024099in}}%
\pgfpathlineto{\pgfqpoint{1.723223in}{7.946625in}}%
\pgfpathlineto{\pgfqpoint{1.710429in}{7.946625in}}%
\pgfpathlineto{\pgfqpoint{1.710429in}{8.014204in}}%
\pgfpathlineto{\pgfqpoint{1.675499in}{8.014204in}}%
\pgfpathlineto{\pgfqpoint{1.675499in}{7.946625in}}%
\pgfpathlineto{\pgfqpoint{1.662705in}{7.946625in}}%
\pgfpathlineto{\pgfqpoint{1.662705in}{8.014204in}}%
\pgfpathlineto{\pgfqpoint{1.650530in}{8.014204in}}%
\pgfpathlineto{\pgfqpoint{1.650530in}{8.024099in}}%
\pgfpathlineto{\pgfqpoint{1.662705in}{8.024099in}}%
\pgfpathlineto{\pgfqpoint{1.662705in}{8.029500in}}%
\pgfpathquadraticcurveto{\pgfqpoint{1.662705in}{8.042161in}}{\pgfqpoint{1.668681in}{8.048204in}}%
\pgfpathquadraticcurveto{\pgfqpoint{1.674680in}{8.054270in}}{\pgfqpoint{1.687054in}{8.054270in}}%
\pgfpathlineto{\pgfqpoint{1.699848in}{8.054270in}}%
\pgfpathlineto{\pgfqpoint{1.699848in}{8.043667in}}%
\pgfpathlineto{\pgfqpoint{1.687673in}{8.043667in}}%
\pgfpathquadraticcurveto{\pgfqpoint{1.680834in}{8.043667in}}{\pgfqpoint{1.678155in}{8.040900in}}%
\pgfpathquadraticcurveto{\pgfqpoint{1.675499in}{8.038133in}}{\pgfqpoint{1.675499in}{8.030939in}}%
\pgfpathlineto{\pgfqpoint{1.675499in}{8.024099in}}%
\pgfpathlineto{\pgfqpoint{1.723223in}{8.024099in}}%
\pgfpathlineto{\pgfqpoint{1.723223in}{8.024099in}}%
\pgfpathclose%
\pgfpathmoveto{\pgfqpoint{1.710429in}{8.054115in}}%
\pgfpathlineto{\pgfqpoint{1.723223in}{8.054115in}}%
\pgfpathlineto{\pgfqpoint{1.723223in}{8.038000in}}%
\pgfpathlineto{\pgfqpoint{1.710429in}{8.038000in}}%
\pgfpathlineto{\pgfqpoint{1.710429in}{8.054115in}}%
\pgfpathlineto{\pgfqpoint{1.710429in}{8.054115in}}%
\pgfpathclose%
\pgfpathmoveto{\pgfqpoint{1.805611in}{8.021133in}}%
\pgfpathlineto{\pgfqpoint{1.805611in}{8.009224in}}%
\pgfpathquadraticcurveto{\pgfqpoint{1.800210in}{8.012212in}}{\pgfqpoint{1.794787in}{8.013695in}}%
\pgfpathquadraticcurveto{\pgfqpoint{1.789364in}{8.015178in}}{\pgfqpoint{1.783830in}{8.015178in}}%
\pgfpathquadraticcurveto{\pgfqpoint{1.771434in}{8.015178in}}{\pgfqpoint{1.764572in}{8.007320in}}%
\pgfpathquadraticcurveto{\pgfqpoint{1.757732in}{7.999484in}}{\pgfqpoint{1.757732in}{7.985296in}}%
\pgfpathquadraticcurveto{\pgfqpoint{1.757732in}{7.971107in}}{\pgfqpoint{1.764572in}{7.963249in}}%
\pgfpathquadraticcurveto{\pgfqpoint{1.771434in}{7.955413in}}{\pgfqpoint{1.783830in}{7.955413in}}%
\pgfpathquadraticcurveto{\pgfqpoint{1.789364in}{7.955413in}}{\pgfqpoint{1.794787in}{7.956896in}}%
\pgfpathquadraticcurveto{\pgfqpoint{1.800210in}{7.958379in}}{\pgfqpoint{1.805611in}{7.961367in}}%
\pgfpathlineto{\pgfqpoint{1.805611in}{7.949591in}}%
\pgfpathquadraticcurveto{\pgfqpoint{1.800276in}{7.947112in}}{\pgfqpoint{1.794565in}{7.945872in}}%
\pgfpathquadraticcurveto{\pgfqpoint{1.788877in}{7.944611in}}{\pgfqpoint{1.782435in}{7.944611in}}%
\pgfpathquadraticcurveto{\pgfqpoint{1.764926in}{7.944611in}}{\pgfqpoint{1.754611in}{7.955612in}}%
\pgfpathquadraticcurveto{\pgfqpoint{1.744318in}{7.966613in}}{\pgfqpoint{1.744318in}{7.985296in}}%
\pgfpathquadraticcurveto{\pgfqpoint{1.744318in}{8.004243in}}{\pgfqpoint{1.754722in}{8.015090in}}%
\pgfpathquadraticcurveto{\pgfqpoint{1.765147in}{8.025958in}}{\pgfqpoint{1.783276in}{8.025958in}}%
\pgfpathquadraticcurveto{\pgfqpoint{1.789142in}{8.025958in}}{\pgfqpoint{1.794742in}{8.024741in}}%
\pgfpathquadraticcurveto{\pgfqpoint{1.800343in}{8.023546in}}{\pgfqpoint{1.805611in}{8.021133in}}%
\pgfpathlineto{\pgfqpoint{1.805611in}{8.021133in}}%
\pgfpathclose%
\pgfpathmoveto{\pgfqpoint{1.862964in}{7.985561in}}%
\pgfpathquadraticcurveto{\pgfqpoint{1.847535in}{7.985561in}}{\pgfqpoint{1.841581in}{7.982042in}}%
\pgfpathquadraticcurveto{\pgfqpoint{1.835627in}{7.978522in}}{\pgfqpoint{1.835627in}{7.970000in}}%
\pgfpathquadraticcurveto{\pgfqpoint{1.835627in}{7.963227in}}{\pgfqpoint{1.840098in}{7.959242in}}%
\pgfpathquadraticcurveto{\pgfqpoint{1.844569in}{7.955280in}}{\pgfqpoint{1.852228in}{7.955280in}}%
\pgfpathquadraticcurveto{\pgfqpoint{1.862831in}{7.955280in}}{\pgfqpoint{1.869228in}{7.962784in}}%
\pgfpathquadraticcurveto{\pgfqpoint{1.875625in}{7.970288in}}{\pgfqpoint{1.875625in}{7.982728in}}%
\pgfpathlineto{\pgfqpoint{1.875625in}{7.985561in}}%
\pgfpathlineto{\pgfqpoint{1.862964in}{7.985561in}}%
\pgfpathlineto{\pgfqpoint{1.862964in}{7.985561in}}%
\pgfpathclose%
\pgfpathmoveto{\pgfqpoint{1.888353in}{7.990829in}}%
\pgfpathlineto{\pgfqpoint{1.888353in}{7.946625in}}%
\pgfpathlineto{\pgfqpoint{1.875625in}{7.946625in}}%
\pgfpathlineto{\pgfqpoint{1.875625in}{7.958379in}}%
\pgfpathquadraticcurveto{\pgfqpoint{1.871265in}{7.951340in}}{\pgfqpoint{1.864757in}{7.947975in}}%
\pgfpathquadraticcurveto{\pgfqpoint{1.858249in}{7.944611in}}{\pgfqpoint{1.848841in}{7.944611in}}%
\pgfpathquadraticcurveto{\pgfqpoint{1.836955in}{7.944611in}}{\pgfqpoint{1.829916in}{7.951296in}}%
\pgfpathquadraticcurveto{\pgfqpoint{1.822899in}{7.957980in}}{\pgfqpoint{1.822899in}{7.969181in}}%
\pgfpathquadraticcurveto{\pgfqpoint{1.822899in}{7.982241in}}{\pgfqpoint{1.831642in}{7.988882in}}%
\pgfpathquadraticcurveto{\pgfqpoint{1.840408in}{7.995522in}}{\pgfqpoint{1.857762in}{7.995522in}}%
\pgfpathlineto{\pgfqpoint{1.875625in}{7.995522in}}%
\pgfpathlineto{\pgfqpoint{1.875625in}{7.996784in}}%
\pgfpathquadraticcurveto{\pgfqpoint{1.875625in}{8.005572in}}{\pgfqpoint{1.869848in}{8.010375in}}%
\pgfpathquadraticcurveto{\pgfqpoint{1.864071in}{8.015178in}}{\pgfqpoint{1.853623in}{8.015178in}}%
\pgfpathquadraticcurveto{\pgfqpoint{1.846982in}{8.015178in}}{\pgfqpoint{1.840673in}{8.013585in}}%
\pgfpathquadraticcurveto{\pgfqpoint{1.834387in}{8.011991in}}{\pgfqpoint{1.828587in}{8.008803in}}%
\pgfpathlineto{\pgfqpoint{1.828587in}{8.020579in}}%
\pgfpathquadraticcurveto{\pgfqpoint{1.835560in}{8.023280in}}{\pgfqpoint{1.842134in}{8.024608in}}%
\pgfpathquadraticcurveto{\pgfqpoint{1.848709in}{8.025958in}}{\pgfqpoint{1.854929in}{8.025958in}}%
\pgfpathquadraticcurveto{\pgfqpoint{1.871752in}{8.025958in}}{\pgfqpoint{1.880052in}{8.017237in}}%
\pgfpathquadraticcurveto{\pgfqpoint{1.888353in}{8.008538in}}{\pgfqpoint{1.888353in}{7.990829in}}%
\pgfpathlineto{\pgfqpoint{1.888353in}{7.990829in}}%
\pgfpathclose%
\pgfpathmoveto{\pgfqpoint{1.978975in}{7.993397in}}%
\pgfpathlineto{\pgfqpoint{1.978975in}{7.946625in}}%
\pgfpathlineto{\pgfqpoint{1.966248in}{7.946625in}}%
\pgfpathlineto{\pgfqpoint{1.966248in}{7.992977in}}%
\pgfpathquadraticcurveto{\pgfqpoint{1.966248in}{8.003978in}}{\pgfqpoint{1.961953in}{8.009423in}}%
\pgfpathquadraticcurveto{\pgfqpoint{1.957659in}{8.014891in}}{\pgfqpoint{1.949093in}{8.014891in}}%
\pgfpathquadraticcurveto{\pgfqpoint{1.938778in}{8.014891in}}{\pgfqpoint{1.932823in}{8.008316in}}%
\pgfpathquadraticcurveto{\pgfqpoint{1.926869in}{8.001764in}}{\pgfqpoint{1.926869in}{7.990409in}}%
\pgfpathlineto{\pgfqpoint{1.926869in}{7.946625in}}%
\pgfpathlineto{\pgfqpoint{1.914074in}{7.946625in}}%
\pgfpathlineto{\pgfqpoint{1.914074in}{8.024099in}}%
\pgfpathlineto{\pgfqpoint{1.926869in}{8.024099in}}%
\pgfpathlineto{\pgfqpoint{1.926869in}{8.012057in}}%
\pgfpathquadraticcurveto{\pgfqpoint{1.931451in}{8.019052in}}{\pgfqpoint{1.937627in}{8.022505in}}%
\pgfpathquadraticcurveto{\pgfqpoint{1.943824in}{8.025958in}}{\pgfqpoint{1.951926in}{8.025958in}}%
\pgfpathquadraticcurveto{\pgfqpoint{1.965274in}{8.025958in}}{\pgfqpoint{1.972114in}{8.017702in}}%
\pgfpathquadraticcurveto{\pgfqpoint{1.978975in}{8.009445in}}{\pgfqpoint{1.978975in}{7.993397in}}%
\pgfpathlineto{\pgfqpoint{1.978975in}{7.993397in}}%
\pgfpathclose%
\pgfpathmoveto{\pgfqpoint{2.016938in}{8.046102in}}%
\pgfpathlineto{\pgfqpoint{2.016938in}{8.024099in}}%
\pgfpathlineto{\pgfqpoint{2.043146in}{8.024099in}}%
\pgfpathlineto{\pgfqpoint{2.043146in}{8.014204in}}%
\pgfpathlineto{\pgfqpoint{2.016938in}{8.014204in}}%
\pgfpathlineto{\pgfqpoint{2.016938in}{7.972147in}}%
\pgfpathquadraticcurveto{\pgfqpoint{2.016938in}{7.962673in}}{\pgfqpoint{2.019528in}{7.959973in}}%
\pgfpathquadraticcurveto{\pgfqpoint{2.022117in}{7.957272in}}{\pgfqpoint{2.030086in}{7.957272in}}%
\pgfpathlineto{\pgfqpoint{2.043146in}{7.957272in}}%
\pgfpathlineto{\pgfqpoint{2.043146in}{7.946625in}}%
\pgfpathlineto{\pgfqpoint{2.030086in}{7.946625in}}%
\pgfpathquadraticcurveto{\pgfqpoint{2.015344in}{7.946625in}}{\pgfqpoint{2.009744in}{7.952115in}}%
\pgfpathquadraticcurveto{\pgfqpoint{2.004143in}{7.957626in}}{\pgfqpoint{2.004143in}{7.972147in}}%
\pgfpathlineto{\pgfqpoint{2.004143in}{8.014204in}}%
\pgfpathlineto{\pgfqpoint{1.994802in}{8.014204in}}%
\pgfpathlineto{\pgfqpoint{1.994802in}{8.024099in}}%
\pgfpathlineto{\pgfqpoint{2.004143in}{8.024099in}}%
\pgfpathlineto{\pgfqpoint{2.004143in}{8.046102in}}%
\pgfpathlineto{\pgfqpoint{2.016938in}{8.046102in}}%
\pgfpathlineto{\pgfqpoint{2.016938in}{8.046102in}}%
\pgfpathclose%
\pgfpathmoveto{\pgfqpoint{2.155882in}{8.012345in}}%
\pgfpathlineto{\pgfqpoint{2.155882in}{8.054270in}}%
\pgfpathlineto{\pgfqpoint{2.168610in}{8.054270in}}%
\pgfpathlineto{\pgfqpoint{2.168610in}{7.946625in}}%
\pgfpathlineto{\pgfqpoint{2.155882in}{7.946625in}}%
\pgfpathlineto{\pgfqpoint{2.155882in}{7.958246in}}%
\pgfpathquadraticcurveto{\pgfqpoint{2.151875in}{7.951340in}}{\pgfqpoint{2.145744in}{7.947975in}}%
\pgfpathquadraticcurveto{\pgfqpoint{2.139634in}{7.944611in}}{\pgfqpoint{2.131046in}{7.944611in}}%
\pgfpathquadraticcurveto{\pgfqpoint{2.117012in}{7.944611in}}{\pgfqpoint{2.108180in}{7.955811in}}%
\pgfpathquadraticcurveto{\pgfqpoint{2.099370in}{7.967034in}}{\pgfqpoint{2.099370in}{7.985296in}}%
\pgfpathquadraticcurveto{\pgfqpoint{2.099370in}{8.003557in}}{\pgfqpoint{2.108180in}{8.014758in}}%
\pgfpathquadraticcurveto{\pgfqpoint{2.117012in}{8.025958in}}{\pgfqpoint{2.131046in}{8.025958in}}%
\pgfpathquadraticcurveto{\pgfqpoint{2.139634in}{8.025958in}}{\pgfqpoint{2.145744in}{8.022594in}}%
\pgfpathquadraticcurveto{\pgfqpoint{2.151875in}{8.019251in}}{\pgfqpoint{2.155882in}{8.012345in}}%
\pgfpathlineto{\pgfqpoint{2.155882in}{8.012345in}}%
\pgfpathclose%
\pgfpathmoveto{\pgfqpoint{2.112518in}{7.985296in}}%
\pgfpathquadraticcurveto{\pgfqpoint{2.112518in}{7.971262in}}{\pgfqpoint{2.118296in}{7.963271in}}%
\pgfpathquadraticcurveto{\pgfqpoint{2.124073in}{7.955280in}}{\pgfqpoint{2.134167in}{7.955280in}}%
\pgfpathquadraticcurveto{\pgfqpoint{2.144261in}{7.955280in}}{\pgfqpoint{2.150060in}{7.963271in}}%
\pgfpathquadraticcurveto{\pgfqpoint{2.155882in}{7.971262in}}{\pgfqpoint{2.155882in}{7.985296in}}%
\pgfpathquadraticcurveto{\pgfqpoint{2.155882in}{7.999329in}}{\pgfqpoint{2.150060in}{8.007320in}}%
\pgfpathquadraticcurveto{\pgfqpoint{2.144261in}{8.015311in}}{\pgfqpoint{2.134167in}{8.015311in}}%
\pgfpathquadraticcurveto{\pgfqpoint{2.124073in}{8.015311in}}{\pgfqpoint{2.118296in}{8.007320in}}%
\pgfpathquadraticcurveto{\pgfqpoint{2.112518in}{7.999329in}}{\pgfqpoint{2.112518in}{7.985296in}}%
\pgfpathlineto{\pgfqpoint{2.112518in}{7.985296in}}%
\pgfpathclose%
\pgfpathmoveto{\pgfqpoint{2.194840in}{8.024099in}}%
\pgfpathlineto{\pgfqpoint{2.207568in}{8.024099in}}%
\pgfpathlineto{\pgfqpoint{2.207568in}{7.946625in}}%
\pgfpathlineto{\pgfqpoint{2.194840in}{7.946625in}}%
\pgfpathlineto{\pgfqpoint{2.194840in}{8.024099in}}%
\pgfpathlineto{\pgfqpoint{2.194840in}{8.024099in}}%
\pgfpathclose%
\pgfpathmoveto{\pgfqpoint{2.194840in}{8.054270in}}%
\pgfpathlineto{\pgfqpoint{2.207568in}{8.054270in}}%
\pgfpathlineto{\pgfqpoint{2.207568in}{8.038133in}}%
\pgfpathlineto{\pgfqpoint{2.194840in}{8.038133in}}%
\pgfpathlineto{\pgfqpoint{2.194840in}{8.054270in}}%
\pgfpathlineto{\pgfqpoint{2.194840in}{8.054270in}}%
\pgfpathclose%
\pgfpathmoveto{\pgfqpoint{2.321145in}{8.054270in}}%
\pgfpathlineto{\pgfqpoint{2.321145in}{8.043667in}}%
\pgfpathlineto{\pgfqpoint{2.308970in}{8.043667in}}%
\pgfpathquadraticcurveto{\pgfqpoint{2.302130in}{8.043667in}}{\pgfqpoint{2.299452in}{8.040900in}}%
\pgfpathquadraticcurveto{\pgfqpoint{2.296796in}{8.038133in}}{\pgfqpoint{2.296796in}{8.030939in}}%
\pgfpathlineto{\pgfqpoint{2.296796in}{8.024099in}}%
\pgfpathlineto{\pgfqpoint{2.317758in}{8.024099in}}%
\pgfpathlineto{\pgfqpoint{2.317758in}{8.014204in}}%
\pgfpathlineto{\pgfqpoint{2.296796in}{8.014204in}}%
\pgfpathlineto{\pgfqpoint{2.296796in}{7.946625in}}%
\pgfpathlineto{\pgfqpoint{2.284002in}{7.946625in}}%
\pgfpathlineto{\pgfqpoint{2.284002in}{8.014204in}}%
\pgfpathlineto{\pgfqpoint{2.249072in}{8.014204in}}%
\pgfpathlineto{\pgfqpoint{2.249072in}{7.946625in}}%
\pgfpathlineto{\pgfqpoint{2.236278in}{7.946625in}}%
\pgfpathlineto{\pgfqpoint{2.236278in}{8.014204in}}%
\pgfpathlineto{\pgfqpoint{2.224103in}{8.014204in}}%
\pgfpathlineto{\pgfqpoint{2.224103in}{8.024099in}}%
\pgfpathlineto{\pgfqpoint{2.236278in}{8.024099in}}%
\pgfpathlineto{\pgfqpoint{2.236278in}{8.029500in}}%
\pgfpathquadraticcurveto{\pgfqpoint{2.236278in}{8.042427in}}{\pgfqpoint{2.242298in}{8.048337in}}%
\pgfpathquadraticcurveto{\pgfqpoint{2.248319in}{8.054270in}}{\pgfqpoint{2.261379in}{8.054270in}}%
\pgfpathlineto{\pgfqpoint{2.273421in}{8.054270in}}%
\pgfpathlineto{\pgfqpoint{2.273421in}{8.043667in}}%
\pgfpathlineto{\pgfqpoint{2.261246in}{8.043667in}}%
\pgfpathquadraticcurveto{\pgfqpoint{2.254406in}{8.043667in}}{\pgfqpoint{2.251706in}{8.040900in}}%
\pgfpathquadraticcurveto{\pgfqpoint{2.249072in}{8.038133in}}{\pgfqpoint{2.249072in}{8.030939in}}%
\pgfpathlineto{\pgfqpoint{2.249072in}{8.024099in}}%
\pgfpathlineto{\pgfqpoint{2.284002in}{8.024099in}}%
\pgfpathlineto{\pgfqpoint{2.284002in}{8.029500in}}%
\pgfpathquadraticcurveto{\pgfqpoint{2.284002in}{8.042427in}}{\pgfqpoint{2.290022in}{8.048337in}}%
\pgfpathquadraticcurveto{\pgfqpoint{2.296043in}{8.054270in}}{\pgfqpoint{2.309125in}{8.054270in}}%
\pgfpathlineto{\pgfqpoint{2.321145in}{8.054270in}}%
\pgfpathlineto{\pgfqpoint{2.321145in}{8.054270in}}%
\pgfpathclose%
\pgfpathmoveto{\pgfqpoint{2.398065in}{7.988549in}}%
\pgfpathlineto{\pgfqpoint{2.398065in}{7.982329in}}%
\pgfpathlineto{\pgfqpoint{2.339539in}{7.982329in}}%
\pgfpathquadraticcurveto{\pgfqpoint{2.340380in}{7.969181in}}{\pgfqpoint{2.347464in}{7.962297in}}%
\pgfpathquadraticcurveto{\pgfqpoint{2.354547in}{7.955413in}}{\pgfqpoint{2.367209in}{7.955413in}}%
\pgfpathquadraticcurveto{\pgfqpoint{2.374535in}{7.955413in}}{\pgfqpoint{2.381420in}{7.957206in}}%
\pgfpathquadraticcurveto{\pgfqpoint{2.388304in}{7.958999in}}{\pgfqpoint{2.395099in}{7.962607in}}%
\pgfpathlineto{\pgfqpoint{2.395099in}{7.950565in}}%
\pgfpathquadraticcurveto{\pgfqpoint{2.388237in}{7.947665in}}{\pgfqpoint{2.381043in}{7.946138in}}%
\pgfpathquadraticcurveto{\pgfqpoint{2.373849in}{7.944611in}}{\pgfqpoint{2.366456in}{7.944611in}}%
\pgfpathquadraticcurveto{\pgfqpoint{2.347906in}{7.944611in}}{\pgfqpoint{2.337082in}{7.955391in}}%
\pgfpathquadraticcurveto{\pgfqpoint{2.326258in}{7.966193in}}{\pgfqpoint{2.326258in}{7.984609in}}%
\pgfpathquadraticcurveto{\pgfqpoint{2.326258in}{8.003624in}}{\pgfqpoint{2.336529in}{8.014780in}}%
\pgfpathquadraticcurveto{\pgfqpoint{2.346800in}{8.025958in}}{\pgfqpoint{2.364242in}{8.025958in}}%
\pgfpathquadraticcurveto{\pgfqpoint{2.379870in}{8.025958in}}{\pgfqpoint{2.388968in}{8.015887in}}%
\pgfpathquadraticcurveto{\pgfqpoint{2.398065in}{8.005837in}}{\pgfqpoint{2.398065in}{7.988549in}}%
\pgfpathlineto{\pgfqpoint{2.398065in}{7.988549in}}%
\pgfpathclose%
\pgfpathmoveto{\pgfqpoint{2.385337in}{7.992290in}}%
\pgfpathquadraticcurveto{\pgfqpoint{2.385205in}{8.002716in}}{\pgfqpoint{2.379494in}{8.008936in}}%
\pgfpathquadraticcurveto{\pgfqpoint{2.373783in}{8.015178in}}{\pgfqpoint{2.364375in}{8.015178in}}%
\pgfpathquadraticcurveto{\pgfqpoint{2.353728in}{8.015178in}}{\pgfqpoint{2.347331in}{8.009158in}}%
\pgfpathquadraticcurveto{\pgfqpoint{2.340934in}{8.003137in}}{\pgfqpoint{2.339960in}{7.992202in}}%
\pgfpathlineto{\pgfqpoint{2.385337in}{7.992290in}}%
\pgfpathlineto{\pgfqpoint{2.385337in}{7.992290in}}%
\pgfpathclose%
\pgfpathmoveto{\pgfqpoint{2.463852in}{8.012212in}}%
\pgfpathquadraticcurveto{\pgfqpoint{2.461705in}{8.013452in}}{\pgfqpoint{2.459181in}{8.014027in}}%
\pgfpathquadraticcurveto{\pgfqpoint{2.456658in}{8.014625in}}{\pgfqpoint{2.453625in}{8.014625in}}%
\pgfpathquadraticcurveto{\pgfqpoint{2.442823in}{8.014625in}}{\pgfqpoint{2.437046in}{8.007608in}}%
\pgfpathquadraticcurveto{\pgfqpoint{2.431268in}{8.000591in}}{\pgfqpoint{2.431268in}{7.987443in}}%
\pgfpathlineto{\pgfqpoint{2.431268in}{7.946625in}}%
\pgfpathlineto{\pgfqpoint{2.418474in}{7.946625in}}%
\pgfpathlineto{\pgfqpoint{2.418474in}{8.024099in}}%
\pgfpathlineto{\pgfqpoint{2.431268in}{8.024099in}}%
\pgfpathlineto{\pgfqpoint{2.431268in}{8.012057in}}%
\pgfpathquadraticcurveto{\pgfqpoint{2.435297in}{8.019118in}}{\pgfqpoint{2.441716in}{8.022527in}}%
\pgfpathquadraticcurveto{\pgfqpoint{2.448158in}{8.025958in}}{\pgfqpoint{2.457366in}{8.025958in}}%
\pgfpathquadraticcurveto{\pgfqpoint{2.458672in}{8.025958in}}{\pgfqpoint{2.460266in}{8.025781in}}%
\pgfpathquadraticcurveto{\pgfqpoint{2.461860in}{8.025626in}}{\pgfqpoint{2.463785in}{8.025272in}}%
\pgfpathlineto{\pgfqpoint{2.463852in}{8.012212in}}%
\pgfpathlineto{\pgfqpoint{2.463852in}{8.012212in}}%
\pgfpathclose%
\pgfpathmoveto{\pgfqpoint{2.540352in}{7.988549in}}%
\pgfpathlineto{\pgfqpoint{2.540352in}{7.982329in}}%
\pgfpathlineto{\pgfqpoint{2.481826in}{7.982329in}}%
\pgfpathquadraticcurveto{\pgfqpoint{2.482667in}{7.969181in}}{\pgfqpoint{2.489750in}{7.962297in}}%
\pgfpathquadraticcurveto{\pgfqpoint{2.496834in}{7.955413in}}{\pgfqpoint{2.509495in}{7.955413in}}%
\pgfpathquadraticcurveto{\pgfqpoint{2.516822in}{7.955413in}}{\pgfqpoint{2.523706in}{7.957206in}}%
\pgfpathquadraticcurveto{\pgfqpoint{2.530590in}{7.958999in}}{\pgfqpoint{2.537386in}{7.962607in}}%
\pgfpathlineto{\pgfqpoint{2.537386in}{7.950565in}}%
\pgfpathquadraticcurveto{\pgfqpoint{2.530524in}{7.947665in}}{\pgfqpoint{2.523330in}{7.946138in}}%
\pgfpathquadraticcurveto{\pgfqpoint{2.516136in}{7.944611in}}{\pgfqpoint{2.508742in}{7.944611in}}%
\pgfpathquadraticcurveto{\pgfqpoint{2.490193in}{7.944611in}}{\pgfqpoint{2.479369in}{7.955391in}}%
\pgfpathquadraticcurveto{\pgfqpoint{2.468545in}{7.966193in}}{\pgfqpoint{2.468545in}{7.984609in}}%
\pgfpathquadraticcurveto{\pgfqpoint{2.468545in}{8.003624in}}{\pgfqpoint{2.478815in}{8.014780in}}%
\pgfpathquadraticcurveto{\pgfqpoint{2.489086in}{8.025958in}}{\pgfqpoint{2.506529in}{8.025958in}}%
\pgfpathquadraticcurveto{\pgfqpoint{2.522156in}{8.025958in}}{\pgfqpoint{2.531254in}{8.015887in}}%
\pgfpathquadraticcurveto{\pgfqpoint{2.540352in}{8.005837in}}{\pgfqpoint{2.540352in}{7.988549in}}%
\pgfpathlineto{\pgfqpoint{2.540352in}{7.988549in}}%
\pgfpathclose%
\pgfpathmoveto{\pgfqpoint{2.527624in}{7.992290in}}%
\pgfpathquadraticcurveto{\pgfqpoint{2.527491in}{8.002716in}}{\pgfqpoint{2.521780in}{8.008936in}}%
\pgfpathquadraticcurveto{\pgfqpoint{2.516069in}{8.015178in}}{\pgfqpoint{2.506662in}{8.015178in}}%
\pgfpathquadraticcurveto{\pgfqpoint{2.496015in}{8.015178in}}{\pgfqpoint{2.489617in}{8.009158in}}%
\pgfpathquadraticcurveto{\pgfqpoint{2.483220in}{8.003137in}}{\pgfqpoint{2.482246in}{7.992202in}}%
\pgfpathlineto{\pgfqpoint{2.527624in}{7.992290in}}%
\pgfpathlineto{\pgfqpoint{2.527624in}{7.992290in}}%
\pgfpathclose%
\pgfpathmoveto{\pgfqpoint{2.625662in}{7.993397in}}%
\pgfpathlineto{\pgfqpoint{2.625662in}{7.946625in}}%
\pgfpathlineto{\pgfqpoint{2.612934in}{7.946625in}}%
\pgfpathlineto{\pgfqpoint{2.612934in}{7.992977in}}%
\pgfpathquadraticcurveto{\pgfqpoint{2.612934in}{8.003978in}}{\pgfqpoint{2.608640in}{8.009423in}}%
\pgfpathquadraticcurveto{\pgfqpoint{2.604345in}{8.014891in}}{\pgfqpoint{2.595779in}{8.014891in}}%
\pgfpathquadraticcurveto{\pgfqpoint{2.585464in}{8.014891in}}{\pgfqpoint{2.579509in}{8.008316in}}%
\pgfpathquadraticcurveto{\pgfqpoint{2.573555in}{8.001764in}}{\pgfqpoint{2.573555in}{7.990409in}}%
\pgfpathlineto{\pgfqpoint{2.573555in}{7.946625in}}%
\pgfpathlineto{\pgfqpoint{2.560761in}{7.946625in}}%
\pgfpathlineto{\pgfqpoint{2.560761in}{8.024099in}}%
\pgfpathlineto{\pgfqpoint{2.573555in}{8.024099in}}%
\pgfpathlineto{\pgfqpoint{2.573555in}{8.012057in}}%
\pgfpathquadraticcurveto{\pgfqpoint{2.578137in}{8.019052in}}{\pgfqpoint{2.584313in}{8.022505in}}%
\pgfpathquadraticcurveto{\pgfqpoint{2.590511in}{8.025958in}}{\pgfqpoint{2.598612in}{8.025958in}}%
\pgfpathquadraticcurveto{\pgfqpoint{2.611960in}{8.025958in}}{\pgfqpoint{2.618800in}{8.017702in}}%
\pgfpathquadraticcurveto{\pgfqpoint{2.625662in}{8.009445in}}{\pgfqpoint{2.625662in}{7.993397in}}%
\pgfpathlineto{\pgfqpoint{2.625662in}{7.993397in}}%
\pgfpathclose%
\pgfpathmoveto{\pgfqpoint{2.706788in}{8.021133in}}%
\pgfpathlineto{\pgfqpoint{2.706788in}{8.009224in}}%
\pgfpathquadraticcurveto{\pgfqpoint{2.701387in}{8.012212in}}{\pgfqpoint{2.695964in}{8.013695in}}%
\pgfpathquadraticcurveto{\pgfqpoint{2.690541in}{8.015178in}}{\pgfqpoint{2.685007in}{8.015178in}}%
\pgfpathquadraticcurveto{\pgfqpoint{2.672611in}{8.015178in}}{\pgfqpoint{2.665749in}{8.007320in}}%
\pgfpathquadraticcurveto{\pgfqpoint{2.658909in}{7.999484in}}{\pgfqpoint{2.658909in}{7.985296in}}%
\pgfpathquadraticcurveto{\pgfqpoint{2.658909in}{7.971107in}}{\pgfqpoint{2.665749in}{7.963249in}}%
\pgfpathquadraticcurveto{\pgfqpoint{2.672611in}{7.955413in}}{\pgfqpoint{2.685007in}{7.955413in}}%
\pgfpathquadraticcurveto{\pgfqpoint{2.690541in}{7.955413in}}{\pgfqpoint{2.695964in}{7.956896in}}%
\pgfpathquadraticcurveto{\pgfqpoint{2.701387in}{7.958379in}}{\pgfqpoint{2.706788in}{7.961367in}}%
\pgfpathlineto{\pgfqpoint{2.706788in}{7.949591in}}%
\pgfpathquadraticcurveto{\pgfqpoint{2.701453in}{7.947112in}}{\pgfqpoint{2.695742in}{7.945872in}}%
\pgfpathquadraticcurveto{\pgfqpoint{2.690054in}{7.944611in}}{\pgfqpoint{2.683612in}{7.944611in}}%
\pgfpathquadraticcurveto{\pgfqpoint{2.666103in}{7.944611in}}{\pgfqpoint{2.655788in}{7.955612in}}%
\pgfpathquadraticcurveto{\pgfqpoint{2.645495in}{7.966613in}}{\pgfqpoint{2.645495in}{7.985296in}}%
\pgfpathquadraticcurveto{\pgfqpoint{2.645495in}{8.004243in}}{\pgfqpoint{2.655899in}{8.015090in}}%
\pgfpathquadraticcurveto{\pgfqpoint{2.666324in}{8.025958in}}{\pgfqpoint{2.684453in}{8.025958in}}%
\pgfpathquadraticcurveto{\pgfqpoint{2.690319in}{8.025958in}}{\pgfqpoint{2.695920in}{8.024741in}}%
\pgfpathquadraticcurveto{\pgfqpoint{2.701520in}{8.023546in}}{\pgfqpoint{2.706788in}{8.021133in}}%
\pgfpathlineto{\pgfqpoint{2.706788in}{8.021133in}}%
\pgfpathclose%
\pgfpathmoveto{\pgfqpoint{2.795197in}{7.988549in}}%
\pgfpathlineto{\pgfqpoint{2.795197in}{7.982329in}}%
\pgfpathlineto{\pgfqpoint{2.736671in}{7.982329in}}%
\pgfpathquadraticcurveto{\pgfqpoint{2.737512in}{7.969181in}}{\pgfqpoint{2.744595in}{7.962297in}}%
\pgfpathquadraticcurveto{\pgfqpoint{2.751679in}{7.955413in}}{\pgfqpoint{2.764340in}{7.955413in}}%
\pgfpathquadraticcurveto{\pgfqpoint{2.771667in}{7.955413in}}{\pgfqpoint{2.778551in}{7.957206in}}%
\pgfpathquadraticcurveto{\pgfqpoint{2.785435in}{7.958999in}}{\pgfqpoint{2.792231in}{7.962607in}}%
\pgfpathlineto{\pgfqpoint{2.792231in}{7.950565in}}%
\pgfpathquadraticcurveto{\pgfqpoint{2.785369in}{7.947665in}}{\pgfqpoint{2.778175in}{7.946138in}}%
\pgfpathquadraticcurveto{\pgfqpoint{2.770981in}{7.944611in}}{\pgfqpoint{2.763587in}{7.944611in}}%
\pgfpathquadraticcurveto{\pgfqpoint{2.745038in}{7.944611in}}{\pgfqpoint{2.734214in}{7.955391in}}%
\pgfpathquadraticcurveto{\pgfqpoint{2.723390in}{7.966193in}}{\pgfqpoint{2.723390in}{7.984609in}}%
\pgfpathquadraticcurveto{\pgfqpoint{2.723390in}{8.003624in}}{\pgfqpoint{2.733660in}{8.014780in}}%
\pgfpathquadraticcurveto{\pgfqpoint{2.743931in}{8.025958in}}{\pgfqpoint{2.761374in}{8.025958in}}%
\pgfpathquadraticcurveto{\pgfqpoint{2.777002in}{8.025958in}}{\pgfqpoint{2.786099in}{8.015887in}}%
\pgfpathquadraticcurveto{\pgfqpoint{2.795197in}{8.005837in}}{\pgfqpoint{2.795197in}{7.988549in}}%
\pgfpathlineto{\pgfqpoint{2.795197in}{7.988549in}}%
\pgfpathclose%
\pgfpathmoveto{\pgfqpoint{2.782469in}{7.992290in}}%
\pgfpathquadraticcurveto{\pgfqpoint{2.782336in}{8.002716in}}{\pgfqpoint{2.776625in}{8.008936in}}%
\pgfpathquadraticcurveto{\pgfqpoint{2.770914in}{8.015178in}}{\pgfqpoint{2.761507in}{8.015178in}}%
\pgfpathquadraticcurveto{\pgfqpoint{2.750860in}{8.015178in}}{\pgfqpoint{2.744462in}{8.009158in}}%
\pgfpathquadraticcurveto{\pgfqpoint{2.738065in}{8.003137in}}{\pgfqpoint{2.737091in}{7.992202in}}%
\pgfpathlineto{\pgfqpoint{2.782469in}{7.992290in}}%
\pgfpathlineto{\pgfqpoint{2.782469in}{7.992290in}}%
\pgfpathclose%
\pgfpathmoveto{\pgfqpoint{2.865477in}{8.021819in}}%
\pgfpathlineto{\pgfqpoint{2.865477in}{8.009777in}}%
\pgfpathquadraticcurveto{\pgfqpoint{2.860098in}{8.012544in}}{\pgfqpoint{2.854276in}{8.013917in}}%
\pgfpathquadraticcurveto{\pgfqpoint{2.848477in}{8.015311in}}{\pgfqpoint{2.842235in}{8.015311in}}%
\pgfpathquadraticcurveto{\pgfqpoint{2.832761in}{8.015311in}}{\pgfqpoint{2.828024in}{8.012411in}}%
\pgfpathquadraticcurveto{\pgfqpoint{2.823287in}{8.009512in}}{\pgfqpoint{2.823287in}{8.003690in}}%
\pgfpathquadraticcurveto{\pgfqpoint{2.823287in}{7.999263in}}{\pgfqpoint{2.826673in}{7.996740in}}%
\pgfpathquadraticcurveto{\pgfqpoint{2.830060in}{7.994216in}}{\pgfqpoint{2.840309in}{7.991936in}}%
\pgfpathlineto{\pgfqpoint{2.844670in}{7.990962in}}%
\pgfpathquadraticcurveto{\pgfqpoint{2.858216in}{7.988062in}}{\pgfqpoint{2.863927in}{7.982772in}}%
\pgfpathquadraticcurveto{\pgfqpoint{2.869638in}{7.977482in}}{\pgfqpoint{2.869638in}{7.968008in}}%
\pgfpathquadraticcurveto{\pgfqpoint{2.869638in}{7.957206in}}{\pgfqpoint{2.861094in}{7.950897in}}%
\pgfpathquadraticcurveto{\pgfqpoint{2.852550in}{7.944611in}}{\pgfqpoint{2.837608in}{7.944611in}}%
\pgfpathquadraticcurveto{\pgfqpoint{2.831388in}{7.944611in}}{\pgfqpoint{2.824637in}{7.945828in}}%
\pgfpathquadraticcurveto{\pgfqpoint{2.817886in}{7.947046in}}{\pgfqpoint{2.810426in}{7.949458in}}%
\pgfpathlineto{\pgfqpoint{2.810426in}{7.962607in}}%
\pgfpathquadraticcurveto{\pgfqpoint{2.817487in}{7.958932in}}{\pgfqpoint{2.824327in}{7.957095in}}%
\pgfpathquadraticcurveto{\pgfqpoint{2.831167in}{7.955280in}}{\pgfqpoint{2.837896in}{7.955280in}}%
\pgfpathquadraticcurveto{\pgfqpoint{2.846883in}{7.955280in}}{\pgfqpoint{2.851709in}{7.958357in}}%
\pgfpathquadraticcurveto{\pgfqpoint{2.856556in}{7.961434in}}{\pgfqpoint{2.856556in}{7.967034in}}%
\pgfpathquadraticcurveto{\pgfqpoint{2.856556in}{7.972214in}}{\pgfqpoint{2.853059in}{7.974980in}}%
\pgfpathquadraticcurveto{\pgfqpoint{2.849584in}{7.977747in}}{\pgfqpoint{2.837741in}{7.980315in}}%
\pgfpathlineto{\pgfqpoint{2.833314in}{7.981355in}}%
\pgfpathquadraticcurveto{\pgfqpoint{2.821494in}{7.983835in}}{\pgfqpoint{2.816225in}{7.988992in}}%
\pgfpathquadraticcurveto{\pgfqpoint{2.810979in}{7.994150in}}{\pgfqpoint{2.810979in}{8.003137in}}%
\pgfpathquadraticcurveto{\pgfqpoint{2.810979in}{8.014072in}}{\pgfqpoint{2.818727in}{8.020004in}}%
\pgfpathquadraticcurveto{\pgfqpoint{2.826474in}{8.025958in}}{\pgfqpoint{2.840729in}{8.025958in}}%
\pgfpathquadraticcurveto{\pgfqpoint{2.847768in}{8.025958in}}{\pgfqpoint{2.853989in}{8.024918in}}%
\pgfpathquadraticcurveto{\pgfqpoint{2.860231in}{8.023900in}}{\pgfqpoint{2.865477in}{8.021819in}}%
\pgfpathlineto{\pgfqpoint{2.865477in}{8.021819in}}%
\pgfpathclose%
\pgfpathmoveto{\pgfqpoint{2.893146in}{7.964201in}}%
\pgfpathlineto{\pgfqpoint{2.907733in}{7.964201in}}%
\pgfpathlineto{\pgfqpoint{2.907733in}{7.946625in}}%
\pgfpathlineto{\pgfqpoint{2.893146in}{7.946625in}}%
\pgfpathlineto{\pgfqpoint{2.893146in}{7.964201in}}%
\pgfpathlineto{\pgfqpoint{2.893146in}{7.964201in}}%
\pgfpathclose%
\pgfpathmoveto{\pgfqpoint{2.893146in}{8.019871in}}%
\pgfpathlineto{\pgfqpoint{2.907733in}{8.019871in}}%
\pgfpathlineto{\pgfqpoint{2.907733in}{8.002318in}}%
\pgfpathlineto{\pgfqpoint{2.893146in}{8.002318in}}%
\pgfpathlineto{\pgfqpoint{2.893146in}{8.019871in}}%
\pgfpathlineto{\pgfqpoint{2.893146in}{8.019871in}}%
\pgfpathclose%
\pgfpathmoveto{\pgfqpoint{3.014315in}{7.995677in}}%
\pgfpathquadraticcurveto{\pgfqpoint{3.004354in}{7.995677in}}{\pgfqpoint{2.998643in}{7.990342in}}%
\pgfpathquadraticcurveto{\pgfqpoint{2.992955in}{7.985008in}}{\pgfqpoint{2.992955in}{7.975689in}}%
\pgfpathquadraticcurveto{\pgfqpoint{2.992955in}{7.966348in}}{\pgfqpoint{2.998643in}{7.961013in}}%
\pgfpathquadraticcurveto{\pgfqpoint{3.004354in}{7.955678in}}{\pgfqpoint{3.014315in}{7.955678in}}%
\pgfpathquadraticcurveto{\pgfqpoint{3.024276in}{7.955678in}}{\pgfqpoint{3.030009in}{7.961035in}}%
\pgfpathquadraticcurveto{\pgfqpoint{3.035765in}{7.966414in}}{\pgfqpoint{3.035765in}{7.975689in}}%
\pgfpathquadraticcurveto{\pgfqpoint{3.035765in}{7.985008in}}{\pgfqpoint{3.030054in}{7.990342in}}%
\pgfpathquadraticcurveto{\pgfqpoint{3.024365in}{7.995677in}}{\pgfqpoint{3.014315in}{7.995677in}}%
\pgfpathlineto{\pgfqpoint{3.014315in}{7.995677in}}%
\pgfpathclose%
\pgfpathmoveto{\pgfqpoint{3.000348in}{8.001609in}}%
\pgfpathquadraticcurveto{\pgfqpoint{2.991361in}{8.003823in}}{\pgfqpoint{2.986336in}{8.009977in}}%
\pgfpathquadraticcurveto{\pgfqpoint{2.981334in}{8.016152in}}{\pgfqpoint{2.981334in}{8.025007in}}%
\pgfpathquadraticcurveto{\pgfqpoint{2.981334in}{8.037380in}}{\pgfqpoint{2.990143in}{8.044574in}}%
\pgfpathquadraticcurveto{\pgfqpoint{2.998975in}{8.051768in}}{\pgfqpoint{3.014315in}{8.051768in}}%
\pgfpathquadraticcurveto{\pgfqpoint{3.029744in}{8.051768in}}{\pgfqpoint{3.038531in}{8.044574in}}%
\pgfpathquadraticcurveto{\pgfqpoint{3.047319in}{8.037380in}}{\pgfqpoint{3.047319in}{8.025007in}}%
\pgfpathquadraticcurveto{\pgfqpoint{3.047319in}{8.016152in}}{\pgfqpoint{3.042295in}{8.009977in}}%
\pgfpathquadraticcurveto{\pgfqpoint{3.037292in}{8.003823in}}{\pgfqpoint{3.028371in}{8.001609in}}%
\pgfpathquadraticcurveto{\pgfqpoint{3.038465in}{7.999263in}}{\pgfqpoint{3.044087in}{7.992401in}}%
\pgfpathquadraticcurveto{\pgfqpoint{3.049732in}{7.985561in}}{\pgfqpoint{3.049732in}{7.975689in}}%
\pgfpathquadraticcurveto{\pgfqpoint{3.049732in}{7.960659in}}{\pgfqpoint{3.040568in}{7.952624in}}%
\pgfpathquadraticcurveto{\pgfqpoint{3.031404in}{7.944611in}}{\pgfqpoint{3.014315in}{7.944611in}}%
\pgfpathquadraticcurveto{\pgfqpoint{2.997249in}{7.944611in}}{\pgfqpoint{2.988063in}{7.952624in}}%
\pgfpathquadraticcurveto{\pgfqpoint{2.978899in}{7.960659in}}{\pgfqpoint{2.978899in}{7.975689in}}%
\pgfpathquadraticcurveto{\pgfqpoint{2.978899in}{7.985561in}}{\pgfqpoint{2.984565in}{7.992401in}}%
\pgfpathquadraticcurveto{\pgfqpoint{2.990254in}{7.999263in}}{\pgfqpoint{3.000348in}{8.001609in}}%
\pgfpathlineto{\pgfqpoint{3.000348in}{8.001609in}}%
\pgfpathclose%
\pgfpathmoveto{\pgfqpoint{2.995235in}{8.023678in}}%
\pgfpathquadraticcurveto{\pgfqpoint{2.995235in}{8.015665in}}{\pgfqpoint{3.000237in}{8.011172in}}%
\pgfpathquadraticcurveto{\pgfqpoint{3.005262in}{8.006678in}}{\pgfqpoint{3.014315in}{8.006678in}}%
\pgfpathquadraticcurveto{\pgfqpoint{3.023324in}{8.006678in}}{\pgfqpoint{3.028393in}{8.011172in}}%
\pgfpathquadraticcurveto{\pgfqpoint{3.033485in}{8.015665in}}{\pgfqpoint{3.033485in}{8.023678in}}%
\pgfpathquadraticcurveto{\pgfqpoint{3.033485in}{8.031714in}}{\pgfqpoint{3.028393in}{8.036207in}}%
\pgfpathquadraticcurveto{\pgfqpoint{3.023324in}{8.040701in}}{\pgfqpoint{3.014315in}{8.040701in}}%
\pgfpathquadraticcurveto{\pgfqpoint{3.005262in}{8.040701in}}{\pgfqpoint{3.000237in}{8.036207in}}%
\pgfpathquadraticcurveto{\pgfqpoint{2.995235in}{8.031714in}}{\pgfqpoint{2.995235in}{8.023678in}}%
\pgfpathlineto{\pgfqpoint{2.995235in}{8.023678in}}%
\pgfpathclose%
\pgfpathmoveto{\pgfqpoint{3.074568in}{7.964201in}}%
\pgfpathlineto{\pgfqpoint{3.089177in}{7.964201in}}%
\pgfpathlineto{\pgfqpoint{3.089177in}{7.946625in}}%
\pgfpathlineto{\pgfqpoint{3.074568in}{7.946625in}}%
\pgfpathlineto{\pgfqpoint{3.074568in}{7.964201in}}%
\pgfpathlineto{\pgfqpoint{3.074568in}{7.964201in}}%
\pgfpathclose%
\pgfpathmoveto{\pgfqpoint{3.157996in}{8.037734in}}%
\pgfpathlineto{\pgfqpoint{3.122712in}{7.982595in}}%
\pgfpathlineto{\pgfqpoint{3.157996in}{7.982595in}}%
\pgfpathlineto{\pgfqpoint{3.157996in}{8.037734in}}%
\pgfpathlineto{\pgfqpoint{3.157996in}{8.037734in}}%
\pgfpathclose%
\pgfpathmoveto{\pgfqpoint{3.154322in}{8.049909in}}%
\pgfpathlineto{\pgfqpoint{3.171897in}{8.049909in}}%
\pgfpathlineto{\pgfqpoint{3.171897in}{7.982595in}}%
\pgfpathlineto{\pgfqpoint{3.186640in}{7.982595in}}%
\pgfpathlineto{\pgfqpoint{3.186640in}{7.970974in}}%
\pgfpathlineto{\pgfqpoint{3.171897in}{7.970974in}}%
\pgfpathlineto{\pgfqpoint{3.171897in}{7.946625in}}%
\pgfpathlineto{\pgfqpoint{3.157996in}{7.946625in}}%
\pgfpathlineto{\pgfqpoint{3.157996in}{7.970974in}}%
\pgfpathlineto{\pgfqpoint{3.111379in}{7.970974in}}%
\pgfpathlineto{\pgfqpoint{3.111379in}{7.984454in}}%
\pgfpathlineto{\pgfqpoint{3.154322in}{8.049909in}}%
\pgfpathlineto{\pgfqpoint{3.154322in}{8.049909in}}%
\pgfpathclose%
\pgfpathmoveto{\pgfqpoint{3.297582in}{7.992069in}}%
\pgfpathquadraticcurveto{\pgfqpoint{3.291561in}{7.992069in}}{\pgfqpoint{3.288130in}{7.986956in}}%
\pgfpathquadraticcurveto{\pgfqpoint{3.284722in}{7.981842in}}{\pgfqpoint{3.284722in}{7.972701in}}%
\pgfpathquadraticcurveto{\pgfqpoint{3.284722in}{7.963714in}}{\pgfqpoint{3.288130in}{7.958556in}}%
\pgfpathquadraticcurveto{\pgfqpoint{3.291561in}{7.953398in}}{\pgfqpoint{3.297582in}{7.953398in}}%
\pgfpathquadraticcurveto{\pgfqpoint{3.303470in}{7.953398in}}{\pgfqpoint{3.306879in}{7.958556in}}%
\pgfpathquadraticcurveto{\pgfqpoint{3.310310in}{7.963714in}}{\pgfqpoint{3.310310in}{7.972701in}}%
\pgfpathquadraticcurveto{\pgfqpoint{3.310310in}{7.981776in}}{\pgfqpoint{3.306879in}{7.986911in}}%
\pgfpathquadraticcurveto{\pgfqpoint{3.303470in}{7.992069in}}{\pgfqpoint{3.297582in}{7.992069in}}%
\pgfpathlineto{\pgfqpoint{3.297582in}{7.992069in}}%
\pgfpathclose%
\pgfpathmoveto{\pgfqpoint{3.297582in}{8.000857in}}%
\pgfpathquadraticcurveto{\pgfqpoint{3.308517in}{8.000857in}}{\pgfqpoint{3.314936in}{7.993242in}}%
\pgfpathquadraticcurveto{\pgfqpoint{3.321378in}{7.985650in}}{\pgfqpoint{3.321378in}{7.972701in}}%
\pgfpathquadraticcurveto{\pgfqpoint{3.321378in}{7.959773in}}{\pgfqpoint{3.314914in}{7.952181in}}%
\pgfpathquadraticcurveto{\pgfqpoint{3.308451in}{7.944611in}}{\pgfqpoint{3.297582in}{7.944611in}}%
\pgfpathquadraticcurveto{\pgfqpoint{3.286515in}{7.944611in}}{\pgfqpoint{3.280073in}{7.952181in}}%
\pgfpathquadraticcurveto{\pgfqpoint{3.273654in}{7.959773in}}{\pgfqpoint{3.273654in}{7.972701in}}%
\pgfpathquadraticcurveto{\pgfqpoint{3.273654in}{7.985716in}}{\pgfqpoint{3.280117in}{7.993286in}}%
\pgfpathquadraticcurveto{\pgfqpoint{3.286581in}{8.000857in}}{\pgfqpoint{3.297582in}{8.000857in}}%
\pgfpathlineto{\pgfqpoint{3.297582in}{8.000857in}}%
\pgfpathclose%
\pgfpathmoveto{\pgfqpoint{3.226196in}{8.042980in}}%
\pgfpathquadraticcurveto{\pgfqpoint{3.220241in}{8.042980in}}{\pgfqpoint{3.216810in}{8.037823in}}%
\pgfpathquadraticcurveto{\pgfqpoint{3.213401in}{8.032688in}}{\pgfqpoint{3.213401in}{8.023678in}}%
\pgfpathquadraticcurveto{\pgfqpoint{3.213401in}{8.014559in}}{\pgfqpoint{3.216788in}{8.009423in}}%
\pgfpathquadraticcurveto{\pgfqpoint{3.220175in}{8.004310in}}{\pgfqpoint{3.226196in}{8.004310in}}%
\pgfpathquadraticcurveto{\pgfqpoint{3.232216in}{8.004310in}}{\pgfqpoint{3.235625in}{8.009423in}}%
\pgfpathquadraticcurveto{\pgfqpoint{3.239056in}{8.014559in}}{\pgfqpoint{3.239056in}{8.023678in}}%
\pgfpathquadraticcurveto{\pgfqpoint{3.239056in}{8.032599in}}{\pgfqpoint{3.235603in}{8.037779in}}%
\pgfpathquadraticcurveto{\pgfqpoint{3.232150in}{8.042980in}}{\pgfqpoint{3.226196in}{8.042980in}}%
\pgfpathlineto{\pgfqpoint{3.226196in}{8.042980in}}%
\pgfpathclose%
\pgfpathmoveto{\pgfqpoint{3.288662in}{8.051768in}}%
\pgfpathlineto{\pgfqpoint{3.299729in}{8.051768in}}%
\pgfpathlineto{\pgfqpoint{3.235116in}{7.944611in}}%
\pgfpathlineto{\pgfqpoint{3.224048in}{7.944611in}}%
\pgfpathlineto{\pgfqpoint{3.288662in}{8.051768in}}%
\pgfpathlineto{\pgfqpoint{3.288662in}{8.051768in}}%
\pgfpathclose%
\pgfpathmoveto{\pgfqpoint{3.226196in}{8.051768in}}%
\pgfpathquadraticcurveto{\pgfqpoint{3.237130in}{8.051768in}}{\pgfqpoint{3.243616in}{8.044198in}}%
\pgfpathquadraticcurveto{\pgfqpoint{3.250124in}{8.036628in}}{\pgfqpoint{3.250124in}{8.023678in}}%
\pgfpathquadraticcurveto{\pgfqpoint{3.250124in}{8.010618in}}{\pgfqpoint{3.243660in}{8.003070in}}%
\pgfpathquadraticcurveto{\pgfqpoint{3.237197in}{7.995522in}}{\pgfqpoint{3.226196in}{7.995522in}}%
\pgfpathquadraticcurveto{\pgfqpoint{3.215194in}{7.995522in}}{\pgfqpoint{3.208797in}{8.003092in}}%
\pgfpathquadraticcurveto{\pgfqpoint{3.202400in}{8.010685in}}{\pgfqpoint{3.202400in}{8.023678in}}%
\pgfpathquadraticcurveto{\pgfqpoint{3.202400in}{8.036561in}}{\pgfqpoint{3.208819in}{8.044154in}}%
\pgfpathquadraticcurveto{\pgfqpoint{3.215261in}{8.051768in}}{\pgfqpoint{3.226196in}{8.051768in}}%
\pgfpathlineto{\pgfqpoint{3.226196in}{8.051768in}}%
\pgfpathclose%
\pgfusepath{fill}%
\end{pgfscope}%
\begin{pgfscope}%
\pgfsetbuttcap%
\pgfsetroundjoin%
\definecolor{currentfill}{rgb}{0.090196,0.168627,0.227451}%
\pgfsetfillcolor{currentfill}%
\pgfsetlinewidth{0.000000pt}%
\definecolor{currentstroke}{rgb}{0.090196,0.168627,0.227451}%
\pgfsetstrokecolor{currentstroke}%
\pgfsetdash{}{0pt}%
\pgfpathmoveto{\pgfqpoint{1.374723in}{7.826939in}}%
\pgfpathlineto{\pgfqpoint{1.374723in}{7.813303in}}%
\pgfpathquadraticcurveto{\pgfqpoint{1.366776in}{7.817111in}}{\pgfqpoint{1.359715in}{7.818970in}}%
\pgfpathquadraticcurveto{\pgfqpoint{1.352654in}{7.820852in}}{\pgfqpoint{1.346080in}{7.820852in}}%
\pgfpathquadraticcurveto{\pgfqpoint{1.334680in}{7.820852in}}{\pgfqpoint{1.328482in}{7.816424in}}%
\pgfpathquadraticcurveto{\pgfqpoint{1.322284in}{7.811997in}}{\pgfqpoint{1.322284in}{7.803829in}}%
\pgfpathquadraticcurveto{\pgfqpoint{1.322284in}{7.796967in}}{\pgfqpoint{1.326401in}{7.793470in}}%
\pgfpathquadraticcurveto{\pgfqpoint{1.330518in}{7.789995in}}{\pgfqpoint{1.342007in}{7.787848in}}%
\pgfpathlineto{\pgfqpoint{1.350440in}{7.786121in}}%
\pgfpathquadraticcurveto{\pgfqpoint{1.366068in}{7.783133in}}{\pgfqpoint{1.373505in}{7.775629in}}%
\pgfpathquadraticcurveto{\pgfqpoint{1.380943in}{7.768125in}}{\pgfqpoint{1.380943in}{7.755552in}}%
\pgfpathquadraticcurveto{\pgfqpoint{1.380943in}{7.740522in}}{\pgfqpoint{1.370871in}{7.732775in}}%
\pgfpathquadraticcurveto{\pgfqpoint{1.360822in}{7.725027in}}{\pgfqpoint{1.341387in}{7.725027in}}%
\pgfpathquadraticcurveto{\pgfqpoint{1.334060in}{7.725027in}}{\pgfqpoint{1.325781in}{7.726687in}}%
\pgfpathquadraticcurveto{\pgfqpoint{1.317525in}{7.728348in}}{\pgfqpoint{1.308671in}{7.731602in}}%
\pgfpathlineto{\pgfqpoint{1.308671in}{7.745990in}}%
\pgfpathquadraticcurveto{\pgfqpoint{1.317171in}{7.741230in}}{\pgfqpoint{1.325339in}{7.738796in}}%
\pgfpathquadraticcurveto{\pgfqpoint{1.333507in}{7.736383in}}{\pgfqpoint{1.341387in}{7.736383in}}%
\pgfpathquadraticcurveto{\pgfqpoint{1.353340in}{7.736383in}}{\pgfqpoint{1.359848in}{7.741076in}}%
\pgfpathquadraticcurveto{\pgfqpoint{1.366356in}{7.745790in}}{\pgfqpoint{1.366356in}{7.754512in}}%
\pgfpathquadraticcurveto{\pgfqpoint{1.366356in}{7.762104in}}{\pgfqpoint{1.361685in}{7.766398in}}%
\pgfpathquadraticcurveto{\pgfqpoint{1.357015in}{7.770693in}}{\pgfqpoint{1.346367in}{7.772840in}}%
\pgfpathlineto{\pgfqpoint{1.337845in}{7.774500in}}%
\pgfpathquadraticcurveto{\pgfqpoint{1.322218in}{7.777599in}}{\pgfqpoint{1.315223in}{7.784240in}}%
\pgfpathquadraticcurveto{\pgfqpoint{1.308250in}{7.790880in}}{\pgfqpoint{1.308250in}{7.802723in}}%
\pgfpathquadraticcurveto{\pgfqpoint{1.308250in}{7.816424in}}{\pgfqpoint{1.317901in}{7.824305in}}%
\pgfpathquadraticcurveto{\pgfqpoint{1.327552in}{7.832185in}}{\pgfqpoint{1.344486in}{7.832185in}}%
\pgfpathquadraticcurveto{\pgfqpoint{1.351768in}{7.832185in}}{\pgfqpoint{1.359295in}{7.830857in}}%
\pgfpathquadraticcurveto{\pgfqpoint{1.366843in}{7.829551in}}{\pgfqpoint{1.374723in}{7.826939in}}%
\pgfpathlineto{\pgfqpoint{1.374723in}{7.826939in}}%
\pgfpathclose%
\pgfpathmoveto{\pgfqpoint{1.402193in}{7.804516in}}%
\pgfpathlineto{\pgfqpoint{1.414921in}{7.804516in}}%
\pgfpathlineto{\pgfqpoint{1.414921in}{7.727042in}}%
\pgfpathlineto{\pgfqpoint{1.402193in}{7.727042in}}%
\pgfpathlineto{\pgfqpoint{1.402193in}{7.804516in}}%
\pgfpathlineto{\pgfqpoint{1.402193in}{7.804516in}}%
\pgfpathclose%
\pgfpathmoveto{\pgfqpoint{1.402193in}{7.834686in}}%
\pgfpathlineto{\pgfqpoint{1.414921in}{7.834686in}}%
\pgfpathlineto{\pgfqpoint{1.414921in}{7.818549in}}%
\pgfpathlineto{\pgfqpoint{1.402193in}{7.818549in}}%
\pgfpathlineto{\pgfqpoint{1.402193in}{7.834686in}}%
\pgfpathlineto{\pgfqpoint{1.402193in}{7.834686in}}%
\pgfpathclose%
\pgfpathmoveto{\pgfqpoint{1.492528in}{7.766686in}}%
\pgfpathquadraticcurveto{\pgfqpoint{1.492528in}{7.780521in}}{\pgfqpoint{1.486817in}{7.788113in}}%
\pgfpathquadraticcurveto{\pgfqpoint{1.481128in}{7.795728in}}{\pgfqpoint{1.470813in}{7.795728in}}%
\pgfpathquadraticcurveto{\pgfqpoint{1.460586in}{7.795728in}}{\pgfqpoint{1.454875in}{7.788113in}}%
\pgfpathquadraticcurveto{\pgfqpoint{1.449164in}{7.780521in}}{\pgfqpoint{1.449164in}{7.766686in}}%
\pgfpathquadraticcurveto{\pgfqpoint{1.449164in}{7.752918in}}{\pgfqpoint{1.454875in}{7.745303in}}%
\pgfpathquadraticcurveto{\pgfqpoint{1.460586in}{7.737689in}}{\pgfqpoint{1.470813in}{7.737689in}}%
\pgfpathquadraticcurveto{\pgfqpoint{1.481128in}{7.737689in}}{\pgfqpoint{1.486817in}{7.745303in}}%
\pgfpathquadraticcurveto{\pgfqpoint{1.492528in}{7.752918in}}{\pgfqpoint{1.492528in}{7.766686in}}%
\pgfpathlineto{\pgfqpoint{1.492528in}{7.766686in}}%
\pgfpathclose%
\pgfpathmoveto{\pgfqpoint{1.505255in}{7.736648in}}%
\pgfpathquadraticcurveto{\pgfqpoint{1.505255in}{7.716882in}}{\pgfqpoint{1.496468in}{7.707230in}}%
\pgfpathquadraticcurveto{\pgfqpoint{1.487702in}{7.697579in}}{\pgfqpoint{1.469573in}{7.697579in}}%
\pgfpathquadraticcurveto{\pgfqpoint{1.462866in}{7.697579in}}{\pgfqpoint{1.456912in}{7.698576in}}%
\pgfpathquadraticcurveto{\pgfqpoint{1.450957in}{7.699572in}}{\pgfqpoint{1.445357in}{7.701652in}}%
\pgfpathlineto{\pgfqpoint{1.445357in}{7.714026in}}%
\pgfpathquadraticcurveto{\pgfqpoint{1.450957in}{7.710993in}}{\pgfqpoint{1.456425in}{7.709555in}}%
\pgfpathquadraticcurveto{\pgfqpoint{1.461892in}{7.708094in}}{\pgfqpoint{1.467559in}{7.708094in}}%
\pgfpathquadraticcurveto{\pgfqpoint{1.480087in}{7.708094in}}{\pgfqpoint{1.486308in}{7.714624in}}%
\pgfpathquadraticcurveto{\pgfqpoint{1.492528in}{7.721154in}}{\pgfqpoint{1.492528in}{7.734368in}}%
\pgfpathlineto{\pgfqpoint{1.492528in}{7.740677in}}%
\pgfpathquadraticcurveto{\pgfqpoint{1.488587in}{7.733815in}}{\pgfqpoint{1.482434in}{7.730428in}}%
\pgfpathquadraticcurveto{\pgfqpoint{1.476280in}{7.727042in}}{\pgfqpoint{1.467692in}{7.727042in}}%
\pgfpathquadraticcurveto{\pgfqpoint{1.453459in}{7.727042in}}{\pgfqpoint{1.444737in}{7.737888in}}%
\pgfpathquadraticcurveto{\pgfqpoint{1.436016in}{7.748757in}}{\pgfqpoint{1.436016in}{7.766686in}}%
\pgfpathquadraticcurveto{\pgfqpoint{1.436016in}{7.784660in}}{\pgfqpoint{1.444737in}{7.795507in}}%
\pgfpathquadraticcurveto{\pgfqpoint{1.453459in}{7.806375in}}{\pgfqpoint{1.467692in}{7.806375in}}%
\pgfpathquadraticcurveto{\pgfqpoint{1.476280in}{7.806375in}}{\pgfqpoint{1.482434in}{7.802988in}}%
\pgfpathquadraticcurveto{\pgfqpoint{1.488587in}{7.799602in}}{\pgfqpoint{1.492528in}{7.792762in}}%
\pgfpathlineto{\pgfqpoint{1.492528in}{7.804516in}}%
\pgfpathlineto{\pgfqpoint{1.505255in}{7.804516in}}%
\pgfpathlineto{\pgfqpoint{1.505255in}{7.736648in}}%
\pgfpathlineto{\pgfqpoint{1.505255in}{7.736648in}}%
\pgfpathclose%
\pgfpathmoveto{\pgfqpoint{1.595900in}{7.773814in}}%
\pgfpathlineto{\pgfqpoint{1.595900in}{7.727042in}}%
\pgfpathlineto{\pgfqpoint{1.583172in}{7.727042in}}%
\pgfpathlineto{\pgfqpoint{1.583172in}{7.773393in}}%
\pgfpathquadraticcurveto{\pgfqpoint{1.583172in}{7.784395in}}{\pgfqpoint{1.578878in}{7.789840in}}%
\pgfpathquadraticcurveto{\pgfqpoint{1.574584in}{7.795307in}}{\pgfqpoint{1.566017in}{7.795307in}}%
\pgfpathquadraticcurveto{\pgfqpoint{1.555702in}{7.795307in}}{\pgfqpoint{1.549748in}{7.788733in}}%
\pgfpathquadraticcurveto{\pgfqpoint{1.543793in}{7.782181in}}{\pgfqpoint{1.543793in}{7.770826in}}%
\pgfpathlineto{\pgfqpoint{1.543793in}{7.727042in}}%
\pgfpathlineto{\pgfqpoint{1.530999in}{7.727042in}}%
\pgfpathlineto{\pgfqpoint{1.530999in}{7.804516in}}%
\pgfpathlineto{\pgfqpoint{1.543793in}{7.804516in}}%
\pgfpathlineto{\pgfqpoint{1.543793in}{7.792474in}}%
\pgfpathquadraticcurveto{\pgfqpoint{1.548375in}{7.799469in}}{\pgfqpoint{1.554551in}{7.802922in}}%
\pgfpathquadraticcurveto{\pgfqpoint{1.560749in}{7.806375in}}{\pgfqpoint{1.568850in}{7.806375in}}%
\pgfpathquadraticcurveto{\pgfqpoint{1.582198in}{7.806375in}}{\pgfqpoint{1.589038in}{7.798118in}}%
\pgfpathquadraticcurveto{\pgfqpoint{1.595900in}{7.789862in}}{\pgfqpoint{1.595900in}{7.773814in}}%
\pgfpathlineto{\pgfqpoint{1.595900in}{7.773814in}}%
\pgfpathclose%
\pgfpathmoveto{\pgfqpoint{1.621267in}{7.804516in}}%
\pgfpathlineto{\pgfqpoint{1.633995in}{7.804516in}}%
\pgfpathlineto{\pgfqpoint{1.633995in}{7.727042in}}%
\pgfpathlineto{\pgfqpoint{1.621267in}{7.727042in}}%
\pgfpathlineto{\pgfqpoint{1.621267in}{7.804516in}}%
\pgfpathlineto{\pgfqpoint{1.621267in}{7.804516in}}%
\pgfpathclose%
\pgfpathmoveto{\pgfqpoint{1.621267in}{7.834686in}}%
\pgfpathlineto{\pgfqpoint{1.633995in}{7.834686in}}%
\pgfpathlineto{\pgfqpoint{1.633995in}{7.818549in}}%
\pgfpathlineto{\pgfqpoint{1.621267in}{7.818549in}}%
\pgfpathlineto{\pgfqpoint{1.621267in}{7.834686in}}%
\pgfpathlineto{\pgfqpoint{1.621267in}{7.834686in}}%
\pgfpathclose%
\pgfpathmoveto{\pgfqpoint{1.723223in}{7.804516in}}%
\pgfpathlineto{\pgfqpoint{1.723223in}{7.727042in}}%
\pgfpathlineto{\pgfqpoint{1.710429in}{7.727042in}}%
\pgfpathlineto{\pgfqpoint{1.710429in}{7.794621in}}%
\pgfpathlineto{\pgfqpoint{1.675499in}{7.794621in}}%
\pgfpathlineto{\pgfqpoint{1.675499in}{7.727042in}}%
\pgfpathlineto{\pgfqpoint{1.662705in}{7.727042in}}%
\pgfpathlineto{\pgfqpoint{1.662705in}{7.794621in}}%
\pgfpathlineto{\pgfqpoint{1.650530in}{7.794621in}}%
\pgfpathlineto{\pgfqpoint{1.650530in}{7.804516in}}%
\pgfpathlineto{\pgfqpoint{1.662705in}{7.804516in}}%
\pgfpathlineto{\pgfqpoint{1.662705in}{7.809917in}}%
\pgfpathquadraticcurveto{\pgfqpoint{1.662705in}{7.822578in}}{\pgfqpoint{1.668681in}{7.828621in}}%
\pgfpathquadraticcurveto{\pgfqpoint{1.674680in}{7.834686in}}{\pgfqpoint{1.687054in}{7.834686in}}%
\pgfpathlineto{\pgfqpoint{1.699848in}{7.834686in}}%
\pgfpathlineto{\pgfqpoint{1.699848in}{7.824083in}}%
\pgfpathlineto{\pgfqpoint{1.687673in}{7.824083in}}%
\pgfpathquadraticcurveto{\pgfqpoint{1.680834in}{7.824083in}}{\pgfqpoint{1.678155in}{7.821316in}}%
\pgfpathquadraticcurveto{\pgfqpoint{1.675499in}{7.818549in}}{\pgfqpoint{1.675499in}{7.811355in}}%
\pgfpathlineto{\pgfqpoint{1.675499in}{7.804516in}}%
\pgfpathlineto{\pgfqpoint{1.723223in}{7.804516in}}%
\pgfpathlineto{\pgfqpoint{1.723223in}{7.804516in}}%
\pgfpathclose%
\pgfpathmoveto{\pgfqpoint{1.710429in}{7.834531in}}%
\pgfpathlineto{\pgfqpoint{1.723223in}{7.834531in}}%
\pgfpathlineto{\pgfqpoint{1.723223in}{7.818417in}}%
\pgfpathlineto{\pgfqpoint{1.710429in}{7.818417in}}%
\pgfpathlineto{\pgfqpoint{1.710429in}{7.834531in}}%
\pgfpathlineto{\pgfqpoint{1.710429in}{7.834531in}}%
\pgfpathclose%
\pgfpathmoveto{\pgfqpoint{1.805611in}{7.801549in}}%
\pgfpathlineto{\pgfqpoint{1.805611in}{7.789641in}}%
\pgfpathquadraticcurveto{\pgfqpoint{1.800210in}{7.792629in}}{\pgfqpoint{1.794787in}{7.794112in}}%
\pgfpathquadraticcurveto{\pgfqpoint{1.789364in}{7.795595in}}{\pgfqpoint{1.783830in}{7.795595in}}%
\pgfpathquadraticcurveto{\pgfqpoint{1.771434in}{7.795595in}}{\pgfqpoint{1.764572in}{7.787737in}}%
\pgfpathquadraticcurveto{\pgfqpoint{1.757732in}{7.779901in}}{\pgfqpoint{1.757732in}{7.765712in}}%
\pgfpathquadraticcurveto{\pgfqpoint{1.757732in}{7.751523in}}{\pgfqpoint{1.764572in}{7.743665in}}%
\pgfpathquadraticcurveto{\pgfqpoint{1.771434in}{7.735829in}}{\pgfqpoint{1.783830in}{7.735829in}}%
\pgfpathquadraticcurveto{\pgfqpoint{1.789364in}{7.735829in}}{\pgfqpoint{1.794787in}{7.737312in}}%
\pgfpathquadraticcurveto{\pgfqpoint{1.800210in}{7.738796in}}{\pgfqpoint{1.805611in}{7.741784in}}%
\pgfpathlineto{\pgfqpoint{1.805611in}{7.730008in}}%
\pgfpathquadraticcurveto{\pgfqpoint{1.800276in}{7.727529in}}{\pgfqpoint{1.794565in}{7.726289in}}%
\pgfpathquadraticcurveto{\pgfqpoint{1.788877in}{7.725027in}}{\pgfqpoint{1.782435in}{7.725027in}}%
\pgfpathquadraticcurveto{\pgfqpoint{1.764926in}{7.725027in}}{\pgfqpoint{1.754611in}{7.736029in}}%
\pgfpathquadraticcurveto{\pgfqpoint{1.744318in}{7.747030in}}{\pgfqpoint{1.744318in}{7.765712in}}%
\pgfpathquadraticcurveto{\pgfqpoint{1.744318in}{7.784660in}}{\pgfqpoint{1.754722in}{7.795507in}}%
\pgfpathquadraticcurveto{\pgfqpoint{1.765147in}{7.806375in}}{\pgfqpoint{1.783276in}{7.806375in}}%
\pgfpathquadraticcurveto{\pgfqpoint{1.789142in}{7.806375in}}{\pgfqpoint{1.794742in}{7.805158in}}%
\pgfpathquadraticcurveto{\pgfqpoint{1.800343in}{7.803962in}}{\pgfqpoint{1.805611in}{7.801549in}}%
\pgfpathlineto{\pgfqpoint{1.805611in}{7.801549in}}%
\pgfpathclose%
\pgfpathmoveto{\pgfqpoint{1.862964in}{7.765978in}}%
\pgfpathquadraticcurveto{\pgfqpoint{1.847535in}{7.765978in}}{\pgfqpoint{1.841581in}{7.762458in}}%
\pgfpathquadraticcurveto{\pgfqpoint{1.835627in}{7.758939in}}{\pgfqpoint{1.835627in}{7.750417in}}%
\pgfpathquadraticcurveto{\pgfqpoint{1.835627in}{7.743643in}}{\pgfqpoint{1.840098in}{7.739659in}}%
\pgfpathquadraticcurveto{\pgfqpoint{1.844569in}{7.735697in}}{\pgfqpoint{1.852228in}{7.735697in}}%
\pgfpathquadraticcurveto{\pgfqpoint{1.862831in}{7.735697in}}{\pgfqpoint{1.869228in}{7.743201in}}%
\pgfpathquadraticcurveto{\pgfqpoint{1.875625in}{7.750704in}}{\pgfqpoint{1.875625in}{7.763145in}}%
\pgfpathlineto{\pgfqpoint{1.875625in}{7.765978in}}%
\pgfpathlineto{\pgfqpoint{1.862964in}{7.765978in}}%
\pgfpathlineto{\pgfqpoint{1.862964in}{7.765978in}}%
\pgfpathclose%
\pgfpathmoveto{\pgfqpoint{1.888353in}{7.771246in}}%
\pgfpathlineto{\pgfqpoint{1.888353in}{7.727042in}}%
\pgfpathlineto{\pgfqpoint{1.875625in}{7.727042in}}%
\pgfpathlineto{\pgfqpoint{1.875625in}{7.738796in}}%
\pgfpathquadraticcurveto{\pgfqpoint{1.871265in}{7.731757in}}{\pgfqpoint{1.864757in}{7.728392in}}%
\pgfpathquadraticcurveto{\pgfqpoint{1.858249in}{7.725027in}}{\pgfqpoint{1.848841in}{7.725027in}}%
\pgfpathquadraticcurveto{\pgfqpoint{1.836955in}{7.725027in}}{\pgfqpoint{1.829916in}{7.731712in}}%
\pgfpathquadraticcurveto{\pgfqpoint{1.822899in}{7.738397in}}{\pgfqpoint{1.822899in}{7.749598in}}%
\pgfpathquadraticcurveto{\pgfqpoint{1.822899in}{7.762658in}}{\pgfqpoint{1.831642in}{7.769298in}}%
\pgfpathquadraticcurveto{\pgfqpoint{1.840408in}{7.775939in}}{\pgfqpoint{1.857762in}{7.775939in}}%
\pgfpathlineto{\pgfqpoint{1.875625in}{7.775939in}}%
\pgfpathlineto{\pgfqpoint{1.875625in}{7.777201in}}%
\pgfpathquadraticcurveto{\pgfqpoint{1.875625in}{7.785988in}}{\pgfqpoint{1.869848in}{7.790792in}}%
\pgfpathquadraticcurveto{\pgfqpoint{1.864071in}{7.795595in}}{\pgfqpoint{1.853623in}{7.795595in}}%
\pgfpathquadraticcurveto{\pgfqpoint{1.846982in}{7.795595in}}{\pgfqpoint{1.840673in}{7.794001in}}%
\pgfpathquadraticcurveto{\pgfqpoint{1.834387in}{7.792408in}}{\pgfqpoint{1.828587in}{7.789220in}}%
\pgfpathlineto{\pgfqpoint{1.828587in}{7.800996in}}%
\pgfpathquadraticcurveto{\pgfqpoint{1.835560in}{7.803697in}}{\pgfqpoint{1.842134in}{7.805025in}}%
\pgfpathquadraticcurveto{\pgfqpoint{1.848709in}{7.806375in}}{\pgfqpoint{1.854929in}{7.806375in}}%
\pgfpathquadraticcurveto{\pgfqpoint{1.871752in}{7.806375in}}{\pgfqpoint{1.880052in}{7.797654in}}%
\pgfpathquadraticcurveto{\pgfqpoint{1.888353in}{7.788954in}}{\pgfqpoint{1.888353in}{7.771246in}}%
\pgfpathlineto{\pgfqpoint{1.888353in}{7.771246in}}%
\pgfpathclose%
\pgfpathmoveto{\pgfqpoint{1.978975in}{7.773814in}}%
\pgfpathlineto{\pgfqpoint{1.978975in}{7.727042in}}%
\pgfpathlineto{\pgfqpoint{1.966248in}{7.727042in}}%
\pgfpathlineto{\pgfqpoint{1.966248in}{7.773393in}}%
\pgfpathquadraticcurveto{\pgfqpoint{1.966248in}{7.784395in}}{\pgfqpoint{1.961953in}{7.789840in}}%
\pgfpathquadraticcurveto{\pgfqpoint{1.957659in}{7.795307in}}{\pgfqpoint{1.949093in}{7.795307in}}%
\pgfpathquadraticcurveto{\pgfqpoint{1.938778in}{7.795307in}}{\pgfqpoint{1.932823in}{7.788733in}}%
\pgfpathquadraticcurveto{\pgfqpoint{1.926869in}{7.782181in}}{\pgfqpoint{1.926869in}{7.770826in}}%
\pgfpathlineto{\pgfqpoint{1.926869in}{7.727042in}}%
\pgfpathlineto{\pgfqpoint{1.914074in}{7.727042in}}%
\pgfpathlineto{\pgfqpoint{1.914074in}{7.804516in}}%
\pgfpathlineto{\pgfqpoint{1.926869in}{7.804516in}}%
\pgfpathlineto{\pgfqpoint{1.926869in}{7.792474in}}%
\pgfpathquadraticcurveto{\pgfqpoint{1.931451in}{7.799469in}}{\pgfqpoint{1.937627in}{7.802922in}}%
\pgfpathquadraticcurveto{\pgfqpoint{1.943824in}{7.806375in}}{\pgfqpoint{1.951926in}{7.806375in}}%
\pgfpathquadraticcurveto{\pgfqpoint{1.965274in}{7.806375in}}{\pgfqpoint{1.972114in}{7.798118in}}%
\pgfpathquadraticcurveto{\pgfqpoint{1.978975in}{7.789862in}}{\pgfqpoint{1.978975in}{7.773814in}}%
\pgfpathlineto{\pgfqpoint{1.978975in}{7.773814in}}%
\pgfpathclose%
\pgfpathmoveto{\pgfqpoint{2.016938in}{7.826518in}}%
\pgfpathlineto{\pgfqpoint{2.016938in}{7.804516in}}%
\pgfpathlineto{\pgfqpoint{2.043146in}{7.804516in}}%
\pgfpathlineto{\pgfqpoint{2.043146in}{7.794621in}}%
\pgfpathlineto{\pgfqpoint{2.016938in}{7.794621in}}%
\pgfpathlineto{\pgfqpoint{2.016938in}{7.752564in}}%
\pgfpathquadraticcurveto{\pgfqpoint{2.016938in}{7.743090in}}{\pgfqpoint{2.019528in}{7.740389in}}%
\pgfpathquadraticcurveto{\pgfqpoint{2.022117in}{7.737689in}}{\pgfqpoint{2.030086in}{7.737689in}}%
\pgfpathlineto{\pgfqpoint{2.043146in}{7.737689in}}%
\pgfpathlineto{\pgfqpoint{2.043146in}{7.727042in}}%
\pgfpathlineto{\pgfqpoint{2.030086in}{7.727042in}}%
\pgfpathquadraticcurveto{\pgfqpoint{2.015344in}{7.727042in}}{\pgfqpoint{2.009744in}{7.732531in}}%
\pgfpathquadraticcurveto{\pgfqpoint{2.004143in}{7.738043in}}{\pgfqpoint{2.004143in}{7.752564in}}%
\pgfpathlineto{\pgfqpoint{2.004143in}{7.794621in}}%
\pgfpathlineto{\pgfqpoint{1.994802in}{7.794621in}}%
\pgfpathlineto{\pgfqpoint{1.994802in}{7.804516in}}%
\pgfpathlineto{\pgfqpoint{2.004143in}{7.804516in}}%
\pgfpathlineto{\pgfqpoint{2.004143in}{7.826518in}}%
\pgfpathlineto{\pgfqpoint{2.016938in}{7.826518in}}%
\pgfpathlineto{\pgfqpoint{2.016938in}{7.826518in}}%
\pgfpathclose%
\pgfpathmoveto{\pgfqpoint{2.156723in}{7.815871in}}%
\pgfpathlineto{\pgfqpoint{2.156723in}{7.777333in}}%
\pgfpathlineto{\pgfqpoint{2.195239in}{7.777333in}}%
\pgfpathlineto{\pgfqpoint{2.195239in}{7.765579in}}%
\pgfpathlineto{\pgfqpoint{2.156723in}{7.765579in}}%
\pgfpathlineto{\pgfqpoint{2.156723in}{7.727042in}}%
\pgfpathlineto{\pgfqpoint{2.145102in}{7.727042in}}%
\pgfpathlineto{\pgfqpoint{2.145102in}{7.765579in}}%
\pgfpathlineto{\pgfqpoint{2.106564in}{7.765579in}}%
\pgfpathlineto{\pgfqpoint{2.106564in}{7.777333in}}%
\pgfpathlineto{\pgfqpoint{2.145102in}{7.777333in}}%
\pgfpathlineto{\pgfqpoint{2.145102in}{7.815871in}}%
\pgfpathlineto{\pgfqpoint{2.156723in}{7.815871in}}%
\pgfpathlineto{\pgfqpoint{2.156723in}{7.815871in}}%
\pgfpathclose%
\pgfpathmoveto{\pgfqpoint{2.318179in}{7.775474in}}%
\pgfpathquadraticcurveto{\pgfqpoint{2.322672in}{7.773947in}}{\pgfqpoint{2.326922in}{7.768966in}}%
\pgfpathquadraticcurveto{\pgfqpoint{2.331172in}{7.763986in}}{\pgfqpoint{2.335466in}{7.755264in}}%
\pgfpathlineto{\pgfqpoint{2.349655in}{7.727042in}}%
\pgfpathlineto{\pgfqpoint{2.334625in}{7.727042in}}%
\pgfpathlineto{\pgfqpoint{2.321433in}{7.753538in}}%
\pgfpathquadraticcurveto{\pgfqpoint{2.316297in}{7.763919in}}{\pgfqpoint{2.311494in}{7.767306in}}%
\pgfpathquadraticcurveto{\pgfqpoint{2.306690in}{7.770693in}}{\pgfqpoint{2.298390in}{7.770693in}}%
\pgfpathlineto{\pgfqpoint{2.283160in}{7.770693in}}%
\pgfpathlineto{\pgfqpoint{2.283160in}{7.727042in}}%
\pgfpathlineto{\pgfqpoint{2.269193in}{7.727042in}}%
\pgfpathlineto{\pgfqpoint{2.269193in}{7.830326in}}%
\pgfpathlineto{\pgfqpoint{2.300736in}{7.830326in}}%
\pgfpathquadraticcurveto{\pgfqpoint{2.318444in}{7.830326in}}{\pgfqpoint{2.327166in}{7.822910in}}%
\pgfpathquadraticcurveto{\pgfqpoint{2.335887in}{7.815517in}}{\pgfqpoint{2.335887in}{7.800576in}}%
\pgfpathquadraticcurveto{\pgfqpoint{2.335887in}{7.790814in}}{\pgfqpoint{2.331349in}{7.784372in}}%
\pgfpathquadraticcurveto{\pgfqpoint{2.326811in}{7.777953in}}{\pgfqpoint{2.318179in}{7.775474in}}%
\pgfpathlineto{\pgfqpoint{2.318179in}{7.775474in}}%
\pgfpathclose%
\pgfpathmoveto{\pgfqpoint{2.283160in}{7.818837in}}%
\pgfpathlineto{\pgfqpoint{2.283160in}{7.782181in}}%
\pgfpathlineto{\pgfqpoint{2.300736in}{7.782181in}}%
\pgfpathquadraticcurveto{\pgfqpoint{2.310830in}{7.782181in}}{\pgfqpoint{2.315987in}{7.786852in}}%
\pgfpathquadraticcurveto{\pgfqpoint{2.321145in}{7.791522in}}{\pgfqpoint{2.321145in}{7.800576in}}%
\pgfpathquadraticcurveto{\pgfqpoint{2.321145in}{7.809629in}}{\pgfqpoint{2.315987in}{7.814233in}}%
\pgfpathquadraticcurveto{\pgfqpoint{2.310830in}{7.818837in}}{\pgfqpoint{2.300736in}{7.818837in}}%
\pgfpathlineto{\pgfqpoint{2.283160in}{7.818837in}}%
\pgfpathlineto{\pgfqpoint{2.283160in}{7.818837in}}%
\pgfpathclose%
\pgfpathmoveto{\pgfqpoint{2.426996in}{7.768966in}}%
\pgfpathlineto{\pgfqpoint{2.426996in}{7.762746in}}%
\pgfpathlineto{\pgfqpoint{2.368470in}{7.762746in}}%
\pgfpathquadraticcurveto{\pgfqpoint{2.369311in}{7.749598in}}{\pgfqpoint{2.376395in}{7.742714in}}%
\pgfpathquadraticcurveto{\pgfqpoint{2.383478in}{7.735829in}}{\pgfqpoint{2.396140in}{7.735829in}}%
\pgfpathquadraticcurveto{\pgfqpoint{2.403466in}{7.735829in}}{\pgfqpoint{2.410350in}{7.737622in}}%
\pgfpathquadraticcurveto{\pgfqpoint{2.417235in}{7.739415in}}{\pgfqpoint{2.424030in}{7.743023in}}%
\pgfpathlineto{\pgfqpoint{2.424030in}{7.730982in}}%
\pgfpathquadraticcurveto{\pgfqpoint{2.417168in}{7.728082in}}{\pgfqpoint{2.409974in}{7.726555in}}%
\pgfpathquadraticcurveto{\pgfqpoint{2.402780in}{7.725027in}}{\pgfqpoint{2.395387in}{7.725027in}}%
\pgfpathquadraticcurveto{\pgfqpoint{2.376837in}{7.725027in}}{\pgfqpoint{2.366013in}{7.735807in}}%
\pgfpathquadraticcurveto{\pgfqpoint{2.355189in}{7.746609in}}{\pgfqpoint{2.355189in}{7.765026in}}%
\pgfpathquadraticcurveto{\pgfqpoint{2.355189in}{7.784040in}}{\pgfqpoint{2.365460in}{7.795197in}}%
\pgfpathquadraticcurveto{\pgfqpoint{2.375731in}{7.806375in}}{\pgfqpoint{2.393173in}{7.806375in}}%
\pgfpathquadraticcurveto{\pgfqpoint{2.408801in}{7.806375in}}{\pgfqpoint{2.417899in}{7.796303in}}%
\pgfpathquadraticcurveto{\pgfqpoint{2.426996in}{7.786254in}}{\pgfqpoint{2.426996in}{7.768966in}}%
\pgfpathlineto{\pgfqpoint{2.426996in}{7.768966in}}%
\pgfpathclose%
\pgfpathmoveto{\pgfqpoint{2.414268in}{7.772707in}}%
\pgfpathquadraticcurveto{\pgfqpoint{2.414136in}{7.783133in}}{\pgfqpoint{2.408425in}{7.789353in}}%
\pgfpathquadraticcurveto{\pgfqpoint{2.402714in}{7.795595in}}{\pgfqpoint{2.393306in}{7.795595in}}%
\pgfpathquadraticcurveto{\pgfqpoint{2.382659in}{7.795595in}}{\pgfqpoint{2.376262in}{7.789574in}}%
\pgfpathquadraticcurveto{\pgfqpoint{2.369865in}{7.783553in}}{\pgfqpoint{2.368891in}{7.772618in}}%
\pgfpathlineto{\pgfqpoint{2.414268in}{7.772707in}}%
\pgfpathlineto{\pgfqpoint{2.414268in}{7.772707in}}%
\pgfpathclose%
\pgfpathmoveto{\pgfqpoint{2.438772in}{7.804516in}}%
\pgfpathlineto{\pgfqpoint{2.452253in}{7.804516in}}%
\pgfpathlineto{\pgfqpoint{2.476469in}{7.739504in}}%
\pgfpathlineto{\pgfqpoint{2.500685in}{7.804516in}}%
\pgfpathlineto{\pgfqpoint{2.514166in}{7.804516in}}%
\pgfpathlineto{\pgfqpoint{2.485102in}{7.727042in}}%
\pgfpathlineto{\pgfqpoint{2.467814in}{7.727042in}}%
\pgfpathlineto{\pgfqpoint{2.438772in}{7.804516in}}%
\pgfpathlineto{\pgfqpoint{2.438772in}{7.804516in}}%
\pgfpathclose%
\pgfpathmoveto{\pgfqpoint{2.598015in}{7.768966in}}%
\pgfpathlineto{\pgfqpoint{2.598015in}{7.762746in}}%
\pgfpathlineto{\pgfqpoint{2.539489in}{7.762746in}}%
\pgfpathquadraticcurveto{\pgfqpoint{2.540330in}{7.749598in}}{\pgfqpoint{2.547413in}{7.742714in}}%
\pgfpathquadraticcurveto{\pgfqpoint{2.554496in}{7.735829in}}{\pgfqpoint{2.567158in}{7.735829in}}%
\pgfpathquadraticcurveto{\pgfqpoint{2.574485in}{7.735829in}}{\pgfqpoint{2.581369in}{7.737622in}}%
\pgfpathquadraticcurveto{\pgfqpoint{2.588253in}{7.739415in}}{\pgfqpoint{2.595048in}{7.743023in}}%
\pgfpathlineto{\pgfqpoint{2.595048in}{7.730982in}}%
\pgfpathquadraticcurveto{\pgfqpoint{2.588186in}{7.728082in}}{\pgfqpoint{2.580992in}{7.726555in}}%
\pgfpathquadraticcurveto{\pgfqpoint{2.573798in}{7.725027in}}{\pgfqpoint{2.566405in}{7.725027in}}%
\pgfpathquadraticcurveto{\pgfqpoint{2.547856in}{7.725027in}}{\pgfqpoint{2.537031in}{7.735807in}}%
\pgfpathquadraticcurveto{\pgfqpoint{2.526207in}{7.746609in}}{\pgfqpoint{2.526207in}{7.765026in}}%
\pgfpathquadraticcurveto{\pgfqpoint{2.526207in}{7.784040in}}{\pgfqpoint{2.536478in}{7.795197in}}%
\pgfpathquadraticcurveto{\pgfqpoint{2.546749in}{7.806375in}}{\pgfqpoint{2.564192in}{7.806375in}}%
\pgfpathquadraticcurveto{\pgfqpoint{2.579819in}{7.806375in}}{\pgfqpoint{2.588917in}{7.796303in}}%
\pgfpathquadraticcurveto{\pgfqpoint{2.598015in}{7.786254in}}{\pgfqpoint{2.598015in}{7.768966in}}%
\pgfpathlineto{\pgfqpoint{2.598015in}{7.768966in}}%
\pgfpathclose%
\pgfpathmoveto{\pgfqpoint{2.585287in}{7.772707in}}%
\pgfpathquadraticcurveto{\pgfqpoint{2.585154in}{7.783133in}}{\pgfqpoint{2.579443in}{7.789353in}}%
\pgfpathquadraticcurveto{\pgfqpoint{2.573732in}{7.795595in}}{\pgfqpoint{2.564324in}{7.795595in}}%
\pgfpathquadraticcurveto{\pgfqpoint{2.553677in}{7.795595in}}{\pgfqpoint{2.547280in}{7.789574in}}%
\pgfpathquadraticcurveto{\pgfqpoint{2.540883in}{7.783553in}}{\pgfqpoint{2.539909in}{7.772618in}}%
\pgfpathlineto{\pgfqpoint{2.585287in}{7.772707in}}%
\pgfpathlineto{\pgfqpoint{2.585287in}{7.772707in}}%
\pgfpathclose%
\pgfpathmoveto{\pgfqpoint{2.663801in}{7.792629in}}%
\pgfpathquadraticcurveto{\pgfqpoint{2.661654in}{7.793868in}}{\pgfqpoint{2.659130in}{7.794444in}}%
\pgfpathquadraticcurveto{\pgfqpoint{2.656607in}{7.795042in}}{\pgfqpoint{2.653574in}{7.795042in}}%
\pgfpathquadraticcurveto{\pgfqpoint{2.642772in}{7.795042in}}{\pgfqpoint{2.636995in}{7.788025in}}%
\pgfpathquadraticcurveto{\pgfqpoint{2.631218in}{7.781008in}}{\pgfqpoint{2.631218in}{7.767859in}}%
\pgfpathlineto{\pgfqpoint{2.631218in}{7.727042in}}%
\pgfpathlineto{\pgfqpoint{2.618423in}{7.727042in}}%
\pgfpathlineto{\pgfqpoint{2.618423in}{7.804516in}}%
\pgfpathlineto{\pgfqpoint{2.631218in}{7.804516in}}%
\pgfpathlineto{\pgfqpoint{2.631218in}{7.792474in}}%
\pgfpathquadraticcurveto{\pgfqpoint{2.635246in}{7.799535in}}{\pgfqpoint{2.641666in}{7.802944in}}%
\pgfpathquadraticcurveto{\pgfqpoint{2.648107in}{7.806375in}}{\pgfqpoint{2.657315in}{7.806375in}}%
\pgfpathquadraticcurveto{\pgfqpoint{2.658621in}{7.806375in}}{\pgfqpoint{2.660215in}{7.806198in}}%
\pgfpathquadraticcurveto{\pgfqpoint{2.661809in}{7.806043in}}{\pgfqpoint{2.663735in}{7.805689in}}%
\pgfpathlineto{\pgfqpoint{2.663801in}{7.792629in}}%
\pgfpathlineto{\pgfqpoint{2.663801in}{7.792629in}}%
\pgfpathclose%
\pgfpathmoveto{\pgfqpoint{2.726533in}{7.802236in}}%
\pgfpathlineto{\pgfqpoint{2.726533in}{7.790194in}}%
\pgfpathquadraticcurveto{\pgfqpoint{2.721154in}{7.792961in}}{\pgfqpoint{2.715332in}{7.794333in}}%
\pgfpathquadraticcurveto{\pgfqpoint{2.709533in}{7.795728in}}{\pgfqpoint{2.703291in}{7.795728in}}%
\pgfpathquadraticcurveto{\pgfqpoint{2.693817in}{7.795728in}}{\pgfqpoint{2.689080in}{7.792828in}}%
\pgfpathquadraticcurveto{\pgfqpoint{2.684343in}{7.789928in}}{\pgfqpoint{2.684343in}{7.784107in}}%
\pgfpathquadraticcurveto{\pgfqpoint{2.684343in}{7.779680in}}{\pgfqpoint{2.687729in}{7.777156in}}%
\pgfpathquadraticcurveto{\pgfqpoint{2.691116in}{7.774633in}}{\pgfqpoint{2.701365in}{7.772353in}}%
\pgfpathlineto{\pgfqpoint{2.705725in}{7.771379in}}%
\pgfpathquadraticcurveto{\pgfqpoint{2.719272in}{7.768479in}}{\pgfqpoint{2.724983in}{7.763189in}}%
\pgfpathquadraticcurveto{\pgfqpoint{2.730694in}{7.757898in}}{\pgfqpoint{2.730694in}{7.748424in}}%
\pgfpathquadraticcurveto{\pgfqpoint{2.730694in}{7.737622in}}{\pgfqpoint{2.722150in}{7.731314in}}%
\pgfpathquadraticcurveto{\pgfqpoint{2.713606in}{7.725027in}}{\pgfqpoint{2.698664in}{7.725027in}}%
\pgfpathquadraticcurveto{\pgfqpoint{2.692444in}{7.725027in}}{\pgfqpoint{2.685693in}{7.726245in}}%
\pgfpathquadraticcurveto{\pgfqpoint{2.678942in}{7.727462in}}{\pgfqpoint{2.671482in}{7.729875in}}%
\pgfpathlineto{\pgfqpoint{2.671482in}{7.743023in}}%
\pgfpathquadraticcurveto{\pgfqpoint{2.678543in}{7.739349in}}{\pgfqpoint{2.685383in}{7.737512in}}%
\pgfpathquadraticcurveto{\pgfqpoint{2.692223in}{7.735697in}}{\pgfqpoint{2.698952in}{7.735697in}}%
\pgfpathquadraticcurveto{\pgfqpoint{2.707939in}{7.735697in}}{\pgfqpoint{2.712765in}{7.738773in}}%
\pgfpathquadraticcurveto{\pgfqpoint{2.717612in}{7.741850in}}{\pgfqpoint{2.717612in}{7.747451in}}%
\pgfpathquadraticcurveto{\pgfqpoint{2.717612in}{7.752630in}}{\pgfqpoint{2.714115in}{7.755397in}}%
\pgfpathquadraticcurveto{\pgfqpoint{2.710640in}{7.758164in}}{\pgfqpoint{2.698797in}{7.760732in}}%
\pgfpathlineto{\pgfqpoint{2.694370in}{7.761772in}}%
\pgfpathquadraticcurveto{\pgfqpoint{2.682550in}{7.764251in}}{\pgfqpoint{2.677281in}{7.769409in}}%
\pgfpathquadraticcurveto{\pgfqpoint{2.672035in}{7.774566in}}{\pgfqpoint{2.672035in}{7.783553in}}%
\pgfpathquadraticcurveto{\pgfqpoint{2.672035in}{7.794488in}}{\pgfqpoint{2.679783in}{7.800421in}}%
\pgfpathquadraticcurveto{\pgfqpoint{2.687530in}{7.806375in}}{\pgfqpoint{2.701785in}{7.806375in}}%
\pgfpathquadraticcurveto{\pgfqpoint{2.708824in}{7.806375in}}{\pgfqpoint{2.715045in}{7.805335in}}%
\pgfpathquadraticcurveto{\pgfqpoint{2.721287in}{7.804316in}}{\pgfqpoint{2.726533in}{7.802236in}}%
\pgfpathlineto{\pgfqpoint{2.726533in}{7.802236in}}%
\pgfpathclose%
\pgfpathmoveto{\pgfqpoint{2.786166in}{7.765978in}}%
\pgfpathquadraticcurveto{\pgfqpoint{2.770737in}{7.765978in}}{\pgfqpoint{2.764783in}{7.762458in}}%
\pgfpathquadraticcurveto{\pgfqpoint{2.758828in}{7.758939in}}{\pgfqpoint{2.758828in}{7.750417in}}%
\pgfpathquadraticcurveto{\pgfqpoint{2.758828in}{7.743643in}}{\pgfqpoint{2.763300in}{7.739659in}}%
\pgfpathquadraticcurveto{\pgfqpoint{2.767771in}{7.735697in}}{\pgfqpoint{2.775430in}{7.735697in}}%
\pgfpathquadraticcurveto{\pgfqpoint{2.786033in}{7.735697in}}{\pgfqpoint{2.792430in}{7.743201in}}%
\pgfpathquadraticcurveto{\pgfqpoint{2.798827in}{7.750704in}}{\pgfqpoint{2.798827in}{7.763145in}}%
\pgfpathlineto{\pgfqpoint{2.798827in}{7.765978in}}%
\pgfpathlineto{\pgfqpoint{2.786166in}{7.765978in}}%
\pgfpathlineto{\pgfqpoint{2.786166in}{7.765978in}}%
\pgfpathclose%
\pgfpathmoveto{\pgfqpoint{2.811555in}{7.771246in}}%
\pgfpathlineto{\pgfqpoint{2.811555in}{7.727042in}}%
\pgfpathlineto{\pgfqpoint{2.798827in}{7.727042in}}%
\pgfpathlineto{\pgfqpoint{2.798827in}{7.738796in}}%
\pgfpathquadraticcurveto{\pgfqpoint{2.794466in}{7.731757in}}{\pgfqpoint{2.787959in}{7.728392in}}%
\pgfpathquadraticcurveto{\pgfqpoint{2.781451in}{7.725027in}}{\pgfqpoint{2.772043in}{7.725027in}}%
\pgfpathquadraticcurveto{\pgfqpoint{2.760156in}{7.725027in}}{\pgfqpoint{2.753117in}{7.731712in}}%
\pgfpathquadraticcurveto{\pgfqpoint{2.746100in}{7.738397in}}{\pgfqpoint{2.746100in}{7.749598in}}%
\pgfpathquadraticcurveto{\pgfqpoint{2.746100in}{7.762658in}}{\pgfqpoint{2.754844in}{7.769298in}}%
\pgfpathquadraticcurveto{\pgfqpoint{2.763610in}{7.775939in}}{\pgfqpoint{2.780964in}{7.775939in}}%
\pgfpathlineto{\pgfqpoint{2.798827in}{7.775939in}}%
\pgfpathlineto{\pgfqpoint{2.798827in}{7.777201in}}%
\pgfpathquadraticcurveto{\pgfqpoint{2.798827in}{7.785988in}}{\pgfqpoint{2.793050in}{7.790792in}}%
\pgfpathquadraticcurveto{\pgfqpoint{2.787272in}{7.795595in}}{\pgfqpoint{2.776824in}{7.795595in}}%
\pgfpathquadraticcurveto{\pgfqpoint{2.770184in}{7.795595in}}{\pgfqpoint{2.763875in}{7.794001in}}%
\pgfpathquadraticcurveto{\pgfqpoint{2.757589in}{7.792408in}}{\pgfqpoint{2.751789in}{7.789220in}}%
\pgfpathlineto{\pgfqpoint{2.751789in}{7.800996in}}%
\pgfpathquadraticcurveto{\pgfqpoint{2.758762in}{7.803697in}}{\pgfqpoint{2.765336in}{7.805025in}}%
\pgfpathquadraticcurveto{\pgfqpoint{2.771910in}{7.806375in}}{\pgfqpoint{2.778130in}{7.806375in}}%
\pgfpathquadraticcurveto{\pgfqpoint{2.794953in}{7.806375in}}{\pgfqpoint{2.803254in}{7.797654in}}%
\pgfpathquadraticcurveto{\pgfqpoint{2.811555in}{7.788954in}}{\pgfqpoint{2.811555in}{7.771246in}}%
\pgfpathlineto{\pgfqpoint{2.811555in}{7.771246in}}%
\pgfpathclose%
\pgfpathmoveto{\pgfqpoint{2.837763in}{7.834686in}}%
\pgfpathlineto{\pgfqpoint{2.850491in}{7.834686in}}%
\pgfpathlineto{\pgfqpoint{2.850491in}{7.727042in}}%
\pgfpathlineto{\pgfqpoint{2.837763in}{7.727042in}}%
\pgfpathlineto{\pgfqpoint{2.837763in}{7.834686in}}%
\pgfpathlineto{\pgfqpoint{2.837763in}{7.834686in}}%
\pgfpathclose%
\pgfpathmoveto{\pgfqpoint{2.880374in}{7.744617in}}%
\pgfpathlineto{\pgfqpoint{2.894961in}{7.744617in}}%
\pgfpathlineto{\pgfqpoint{2.894961in}{7.727042in}}%
\pgfpathlineto{\pgfqpoint{2.880374in}{7.727042in}}%
\pgfpathlineto{\pgfqpoint{2.880374in}{7.744617in}}%
\pgfpathlineto{\pgfqpoint{2.880374in}{7.744617in}}%
\pgfpathclose%
\pgfpathmoveto{\pgfqpoint{2.880374in}{7.800288in}}%
\pgfpathlineto{\pgfqpoint{2.894961in}{7.800288in}}%
\pgfpathlineto{\pgfqpoint{2.894961in}{7.782734in}}%
\pgfpathlineto{\pgfqpoint{2.880374in}{7.782734in}}%
\pgfpathlineto{\pgfqpoint{2.880374in}{7.800288in}}%
\pgfpathlineto{\pgfqpoint{2.880374in}{7.800288in}}%
\pgfpathclose%
\pgfpathmoveto{\pgfqpoint{2.974095in}{7.738796in}}%
\pgfpathlineto{\pgfqpoint{2.996917in}{7.738796in}}%
\pgfpathlineto{\pgfqpoint{2.996917in}{7.817598in}}%
\pgfpathlineto{\pgfqpoint{2.972081in}{7.812617in}}%
\pgfpathlineto{\pgfqpoint{2.972081in}{7.825345in}}%
\pgfpathlineto{\pgfqpoint{2.996784in}{7.830326in}}%
\pgfpathlineto{\pgfqpoint{3.010752in}{7.830326in}}%
\pgfpathlineto{\pgfqpoint{3.010752in}{7.738796in}}%
\pgfpathlineto{\pgfqpoint{3.033573in}{7.738796in}}%
\pgfpathlineto{\pgfqpoint{3.033573in}{7.727042in}}%
\pgfpathlineto{\pgfqpoint{2.974095in}{7.727042in}}%
\pgfpathlineto{\pgfqpoint{2.974095in}{7.738796in}}%
\pgfpathlineto{\pgfqpoint{2.974095in}{7.738796in}}%
\pgfpathclose%
\pgfpathmoveto{\pgfqpoint{3.061796in}{7.744617in}}%
\pgfpathlineto{\pgfqpoint{3.076405in}{7.744617in}}%
\pgfpathlineto{\pgfqpoint{3.076405in}{7.727042in}}%
\pgfpathlineto{\pgfqpoint{3.061796in}{7.727042in}}%
\pgfpathlineto{\pgfqpoint{3.061796in}{7.744617in}}%
\pgfpathlineto{\pgfqpoint{3.061796in}{7.744617in}}%
\pgfpathclose%
\pgfpathmoveto{\pgfqpoint{3.109254in}{7.738796in}}%
\pgfpathlineto{\pgfqpoint{3.132076in}{7.738796in}}%
\pgfpathlineto{\pgfqpoint{3.132076in}{7.817598in}}%
\pgfpathlineto{\pgfqpoint{3.107240in}{7.812617in}}%
\pgfpathlineto{\pgfqpoint{3.107240in}{7.825345in}}%
\pgfpathlineto{\pgfqpoint{3.131943in}{7.830326in}}%
\pgfpathlineto{\pgfqpoint{3.145910in}{7.830326in}}%
\pgfpathlineto{\pgfqpoint{3.145910in}{7.738796in}}%
\pgfpathlineto{\pgfqpoint{3.168732in}{7.738796in}}%
\pgfpathlineto{\pgfqpoint{3.168732in}{7.727042in}}%
\pgfpathlineto{\pgfqpoint{3.109254in}{7.727042in}}%
\pgfpathlineto{\pgfqpoint{3.109254in}{7.738796in}}%
\pgfpathlineto{\pgfqpoint{3.109254in}{7.738796in}}%
\pgfpathclose%
\pgfpathmoveto{\pgfqpoint{3.284810in}{7.772486in}}%
\pgfpathquadraticcurveto{\pgfqpoint{3.278789in}{7.772486in}}{\pgfqpoint{3.275358in}{7.767372in}}%
\pgfpathquadraticcurveto{\pgfqpoint{3.271949in}{7.762259in}}{\pgfqpoint{3.271949in}{7.753117in}}%
\pgfpathquadraticcurveto{\pgfqpoint{3.271949in}{7.744130in}}{\pgfqpoint{3.275358in}{7.738973in}}%
\pgfpathquadraticcurveto{\pgfqpoint{3.278789in}{7.733815in}}{\pgfqpoint{3.284810in}{7.733815in}}%
\pgfpathquadraticcurveto{\pgfqpoint{3.290698in}{7.733815in}}{\pgfqpoint{3.294107in}{7.738973in}}%
\pgfpathquadraticcurveto{\pgfqpoint{3.297538in}{7.744130in}}{\pgfqpoint{3.297538in}{7.753117in}}%
\pgfpathquadraticcurveto{\pgfqpoint{3.297538in}{7.762193in}}{\pgfqpoint{3.294107in}{7.767328in}}%
\pgfpathquadraticcurveto{\pgfqpoint{3.290698in}{7.772486in}}{\pgfqpoint{3.284810in}{7.772486in}}%
\pgfpathlineto{\pgfqpoint{3.284810in}{7.772486in}}%
\pgfpathclose%
\pgfpathmoveto{\pgfqpoint{3.284810in}{7.781273in}}%
\pgfpathquadraticcurveto{\pgfqpoint{3.295745in}{7.781273in}}{\pgfqpoint{3.302164in}{7.773659in}}%
\pgfpathquadraticcurveto{\pgfqpoint{3.308606in}{7.766066in}}{\pgfqpoint{3.308606in}{7.753117in}}%
\pgfpathquadraticcurveto{\pgfqpoint{3.308606in}{7.740190in}}{\pgfqpoint{3.302142in}{7.732598in}}%
\pgfpathquadraticcurveto{\pgfqpoint{3.295679in}{7.725027in}}{\pgfqpoint{3.284810in}{7.725027in}}%
\pgfpathquadraticcurveto{\pgfqpoint{3.273742in}{7.725027in}}{\pgfqpoint{3.267301in}{7.732598in}}%
\pgfpathquadraticcurveto{\pgfqpoint{3.260882in}{7.740190in}}{\pgfqpoint{3.260882in}{7.753117in}}%
\pgfpathquadraticcurveto{\pgfqpoint{3.260882in}{7.766133in}}{\pgfqpoint{3.267345in}{7.773703in}}%
\pgfpathquadraticcurveto{\pgfqpoint{3.273809in}{7.781273in}}{\pgfqpoint{3.284810in}{7.781273in}}%
\pgfpathlineto{\pgfqpoint{3.284810in}{7.781273in}}%
\pgfpathclose%
\pgfpathmoveto{\pgfqpoint{3.213423in}{7.823397in}}%
\pgfpathquadraticcurveto{\pgfqpoint{3.207469in}{7.823397in}}{\pgfqpoint{3.204038in}{7.818240in}}%
\pgfpathquadraticcurveto{\pgfqpoint{3.200629in}{7.813104in}}{\pgfqpoint{3.200629in}{7.804095in}}%
\pgfpathquadraticcurveto{\pgfqpoint{3.200629in}{7.794975in}}{\pgfqpoint{3.204016in}{7.789840in}}%
\pgfpathquadraticcurveto{\pgfqpoint{3.207403in}{7.784727in}}{\pgfqpoint{3.213423in}{7.784727in}}%
\pgfpathquadraticcurveto{\pgfqpoint{3.219444in}{7.784727in}}{\pgfqpoint{3.222853in}{7.789840in}}%
\pgfpathquadraticcurveto{\pgfqpoint{3.226284in}{7.794975in}}{\pgfqpoint{3.226284in}{7.804095in}}%
\pgfpathquadraticcurveto{\pgfqpoint{3.226284in}{7.813016in}}{\pgfqpoint{3.222831in}{7.818195in}}%
\pgfpathquadraticcurveto{\pgfqpoint{3.219378in}{7.823397in}}{\pgfqpoint{3.213423in}{7.823397in}}%
\pgfpathlineto{\pgfqpoint{3.213423in}{7.823397in}}%
\pgfpathclose%
\pgfpathmoveto{\pgfqpoint{3.275890in}{7.832185in}}%
\pgfpathlineto{\pgfqpoint{3.286957in}{7.832185in}}%
\pgfpathlineto{\pgfqpoint{3.222344in}{7.725027in}}%
\pgfpathlineto{\pgfqpoint{3.211276in}{7.725027in}}%
\pgfpathlineto{\pgfqpoint{3.275890in}{7.832185in}}%
\pgfpathlineto{\pgfqpoint{3.275890in}{7.832185in}}%
\pgfpathclose%
\pgfpathmoveto{\pgfqpoint{3.213423in}{7.832185in}}%
\pgfpathquadraticcurveto{\pgfqpoint{3.224358in}{7.832185in}}{\pgfqpoint{3.230844in}{7.824615in}}%
\pgfpathquadraticcurveto{\pgfqpoint{3.237352in}{7.817044in}}{\pgfqpoint{3.237352in}{7.804095in}}%
\pgfpathquadraticcurveto{\pgfqpoint{3.237352in}{7.791035in}}{\pgfqpoint{3.230888in}{7.783487in}}%
\pgfpathquadraticcurveto{\pgfqpoint{3.224425in}{7.775939in}}{\pgfqpoint{3.213423in}{7.775939in}}%
\pgfpathquadraticcurveto{\pgfqpoint{3.202422in}{7.775939in}}{\pgfqpoint{3.196025in}{7.783509in}}%
\pgfpathquadraticcurveto{\pgfqpoint{3.189628in}{7.791102in}}{\pgfqpoint{3.189628in}{7.804095in}}%
\pgfpathquadraticcurveto{\pgfqpoint{3.189628in}{7.816978in}}{\pgfqpoint{3.196047in}{7.824570in}}%
\pgfpathquadraticcurveto{\pgfqpoint{3.202489in}{7.832185in}}{\pgfqpoint{3.213423in}{7.832185in}}%
\pgfpathlineto{\pgfqpoint{3.213423in}{7.832185in}}%
\pgfpathclose%
\pgfusepath{fill}%
\end{pgfscope}%
\begin{pgfscope}%
\pgfsetbuttcap%
\pgfsetmiterjoin%
\definecolor{currentfill}{rgb}{1.000000,1.000000,1.000000}%
\pgfsetfillcolor{currentfill}%
\pgfsetlinewidth{0.000000pt}%
\definecolor{currentstroke}{rgb}{0.000000,0.000000,0.000000}%
\pgfsetstrokecolor{currentstroke}%
\pgfsetstrokeopacity{0.000000}%
\pgfsetdash{}{0pt}%
\pgfpathmoveto{\pgfqpoint{8.036091in}{4.050000in}}%
\pgfpathlineto{\pgfqpoint{13.436091in}{4.050000in}}%
\pgfpathlineto{\pgfqpoint{13.436091in}{9.450000in}}%
\pgfpathlineto{\pgfqpoint{8.036091in}{9.450000in}}%
\pgfpathlineto{\pgfqpoint{8.036091in}{4.050000in}}%
\pgfpathclose%
\pgfusepath{fill}%
\end{pgfscope}%
\begin{pgfscope}%
\pgfpathrectangle{\pgfqpoint{8.036091in}{4.050000in}}{\pgfqpoint{5.400000in}{5.400000in}}%
\pgfusepath{clip}%
\pgfsetrectcap%
\pgfsetroundjoin%
\pgfsetlinewidth{0.602250pt}%
\definecolor{currentstroke}{rgb}{0.898039,0.913725,0.929412}%
\pgfsetstrokecolor{currentstroke}%
\pgfsetdash{}{0pt}%
\pgfpathmoveto{\pgfqpoint{8.036091in}{4.050000in}}%
\pgfpathlineto{\pgfqpoint{8.036091in}{9.450000in}}%
\pgfusepath{stroke}%
\end{pgfscope}%
\begin{pgfscope}%
\pgfsetbuttcap%
\pgfsetroundjoin%
\definecolor{currentfill}{rgb}{0.090196,0.168627,0.227451}%
\pgfsetfillcolor{currentfill}%
\pgfsetlinewidth{0.803000pt}%
\definecolor{currentstroke}{rgb}{0.090196,0.168627,0.227451}%
\pgfsetstrokecolor{currentstroke}%
\pgfsetdash{}{0pt}%
\pgfsys@defobject{currentmarker}{\pgfqpoint{0.000000in}{-0.048611in}}{\pgfqpoint{0.000000in}{0.000000in}}{%
\pgfpathmoveto{\pgfqpoint{0.000000in}{0.000000in}}%
\pgfpathlineto{\pgfqpoint{0.000000in}{-0.048611in}}%
\pgfusepath{stroke,fill}%
}%
\begin{pgfscope}%
\pgfsys@transformshift{8.036091in}{4.050000in}%
\pgfsys@useobject{currentmarker}{}%
\end{pgfscope}%
\end{pgfscope}%
\begin{pgfscope}%
\pgfsetbuttcap%
\pgfsetroundjoin%
\definecolor{currentfill}{rgb}{0.090196,0.168627,0.227451}%
\pgfsetfillcolor{currentfill}%
\pgfsetlinewidth{0.000000pt}%
\definecolor{currentstroke}{rgb}{0.090196,0.168627,0.227451}%
\pgfsetstrokecolor{currentstroke}%
\pgfsetdash{}{0pt}%
\pgfpathmoveto{\pgfqpoint{7.894333in}{3.902182in}}%
\pgfpathlineto{\pgfqpoint{7.972575in}{3.902182in}}%
\pgfpathlineto{\pgfqpoint{7.972575in}{3.891811in}}%
\pgfpathlineto{\pgfqpoint{7.894333in}{3.891811in}}%
\pgfpathlineto{\pgfqpoint{7.894333in}{3.902182in}}%
\pgfpathlineto{\pgfqpoint{7.894333in}{3.902182in}}%
\pgfpathclose%
\pgfpathmoveto{\pgfqpoint{7.999333in}{3.948940in}}%
\pgfpathlineto{\pgfqpoint{8.047732in}{3.948940in}}%
\pgfpathlineto{\pgfqpoint{8.047732in}{3.938549in}}%
\pgfpathlineto{\pgfqpoint{8.010622in}{3.938549in}}%
\pgfpathlineto{\pgfqpoint{8.010622in}{3.916225in}}%
\pgfpathquadraticcurveto{\pgfqpoint{8.013298in}{3.917143in}}{\pgfqpoint{8.015974in}{3.917592in}}%
\pgfpathquadraticcurveto{\pgfqpoint{8.018669in}{3.918041in}}{\pgfqpoint{8.021364in}{3.918041in}}%
\pgfpathquadraticcurveto{\pgfqpoint{8.036618in}{3.918041in}}{\pgfqpoint{8.045525in}{3.909682in}}%
\pgfpathquadraticcurveto{\pgfqpoint{8.054450in}{3.901323in}}{\pgfqpoint{8.054450in}{3.887045in}}%
\pgfpathquadraticcurveto{\pgfqpoint{8.054450in}{3.872338in}}{\pgfqpoint{8.045290in}{3.864174in}}%
\pgfpathquadraticcurveto{\pgfqpoint{8.036130in}{3.856030in}}{\pgfqpoint{8.019470in}{3.856030in}}%
\pgfpathquadraticcurveto{\pgfqpoint{8.013728in}{3.856030in}}{\pgfqpoint{8.007771in}{3.857006in}}%
\pgfpathquadraticcurveto{\pgfqpoint{8.001833in}{3.857983in}}{\pgfqpoint{7.995485in}{3.859936in}}%
\pgfpathlineto{\pgfqpoint{7.995485in}{3.872338in}}%
\pgfpathquadraticcurveto{\pgfqpoint{8.000974in}{3.869350in}}{\pgfqpoint{8.006833in}{3.867885in}}%
\pgfpathquadraticcurveto{\pgfqpoint{8.012692in}{3.866420in}}{\pgfqpoint{8.019216in}{3.866420in}}%
\pgfpathquadraticcurveto{\pgfqpoint{8.029782in}{3.866420in}}{\pgfqpoint{8.035935in}{3.871967in}}%
\pgfpathquadraticcurveto{\pgfqpoint{8.042107in}{3.877514in}}{\pgfqpoint{8.042107in}{3.887045in}}%
\pgfpathquadraticcurveto{\pgfqpoint{8.042107in}{3.896557in}}{\pgfqpoint{8.035935in}{3.902104in}}%
\pgfpathquadraticcurveto{\pgfqpoint{8.029782in}{3.907670in}}{\pgfqpoint{8.019216in}{3.907670in}}%
\pgfpathquadraticcurveto{\pgfqpoint{8.014275in}{3.907670in}}{\pgfqpoint{8.009353in}{3.906577in}}%
\pgfpathquadraticcurveto{\pgfqpoint{8.004450in}{3.905483in}}{\pgfqpoint{7.999333in}{3.903159in}}%
\pgfpathlineto{\pgfqpoint{7.999333in}{3.948940in}}%
\pgfpathlineto{\pgfqpoint{7.999333in}{3.948940in}}%
\pgfpathclose%
\pgfpathmoveto{\pgfqpoint{8.078728in}{3.873315in}}%
\pgfpathlineto{\pgfqpoint{8.091618in}{3.873315in}}%
\pgfpathlineto{\pgfqpoint{8.091618in}{3.857807in}}%
\pgfpathlineto{\pgfqpoint{8.078728in}{3.857807in}}%
\pgfpathlineto{\pgfqpoint{8.078728in}{3.873315in}}%
\pgfpathlineto{\pgfqpoint{8.078728in}{3.873315in}}%
\pgfpathclose%
\pgfpathmoveto{\pgfqpoint{8.144821in}{3.940815in}}%
\pgfpathquadraticcurveto{\pgfqpoint{8.135310in}{3.940815in}}{\pgfqpoint{8.130505in}{3.931440in}}%
\pgfpathquadraticcurveto{\pgfqpoint{8.125720in}{3.922084in}}{\pgfqpoint{8.125720in}{3.903276in}}%
\pgfpathquadraticcurveto{\pgfqpoint{8.125720in}{3.884545in}}{\pgfqpoint{8.130505in}{3.875170in}}%
\pgfpathquadraticcurveto{\pgfqpoint{8.135310in}{3.865795in}}{\pgfqpoint{8.144821in}{3.865795in}}%
\pgfpathquadraticcurveto{\pgfqpoint{8.154411in}{3.865795in}}{\pgfqpoint{8.159196in}{3.875170in}}%
\pgfpathquadraticcurveto{\pgfqpoint{8.164001in}{3.884545in}}{\pgfqpoint{8.164001in}{3.903276in}}%
\pgfpathquadraticcurveto{\pgfqpoint{8.164001in}{3.922084in}}{\pgfqpoint{8.159196in}{3.931440in}}%
\pgfpathquadraticcurveto{\pgfqpoint{8.154411in}{3.940815in}}{\pgfqpoint{8.144821in}{3.940815in}}%
\pgfpathlineto{\pgfqpoint{8.144821in}{3.940815in}}%
\pgfpathclose%
\pgfpathmoveto{\pgfqpoint{8.144821in}{3.950581in}}%
\pgfpathquadraticcurveto{\pgfqpoint{8.160153in}{3.950581in}}{\pgfqpoint{8.168239in}{3.938452in}}%
\pgfpathquadraticcurveto{\pgfqpoint{8.176325in}{3.926342in}}{\pgfqpoint{8.176325in}{3.903276in}}%
\pgfpathquadraticcurveto{\pgfqpoint{8.176325in}{3.880268in}}{\pgfqpoint{8.168239in}{3.868139in}}%
\pgfpathquadraticcurveto{\pgfqpoint{8.160153in}{3.856030in}}{\pgfqpoint{8.144821in}{3.856030in}}%
\pgfpathquadraticcurveto{\pgfqpoint{8.129509in}{3.856030in}}{\pgfqpoint{8.121423in}{3.868139in}}%
\pgfpathquadraticcurveto{\pgfqpoint{8.113337in}{3.880268in}}{\pgfqpoint{8.113337in}{3.903276in}}%
\pgfpathquadraticcurveto{\pgfqpoint{8.113337in}{3.926342in}}{\pgfqpoint{8.121423in}{3.938452in}}%
\pgfpathquadraticcurveto{\pgfqpoint{8.129509in}{3.950581in}}{\pgfqpoint{8.144821in}{3.950581in}}%
\pgfpathlineto{\pgfqpoint{8.144821in}{3.950581in}}%
\pgfpathclose%
\pgfusepath{fill}%
\end{pgfscope}%
\begin{pgfscope}%
\pgfpathrectangle{\pgfqpoint{8.036091in}{4.050000in}}{\pgfqpoint{5.400000in}{5.400000in}}%
\pgfusepath{clip}%
\pgfsetrectcap%
\pgfsetroundjoin%
\pgfsetlinewidth{0.602250pt}%
\definecolor{currentstroke}{rgb}{0.898039,0.913725,0.929412}%
\pgfsetstrokecolor{currentstroke}%
\pgfsetdash{}{0pt}%
\pgfpathmoveto{\pgfqpoint{9.386091in}{4.050000in}}%
\pgfpathlineto{\pgfqpoint{9.386091in}{9.450000in}}%
\pgfusepath{stroke}%
\end{pgfscope}%
\begin{pgfscope}%
\pgfsetbuttcap%
\pgfsetroundjoin%
\definecolor{currentfill}{rgb}{0.090196,0.168627,0.227451}%
\pgfsetfillcolor{currentfill}%
\pgfsetlinewidth{0.803000pt}%
\definecolor{currentstroke}{rgb}{0.090196,0.168627,0.227451}%
\pgfsetstrokecolor{currentstroke}%
\pgfsetdash{}{0pt}%
\pgfsys@defobject{currentmarker}{\pgfqpoint{0.000000in}{-0.048611in}}{\pgfqpoint{0.000000in}{0.000000in}}{%
\pgfpathmoveto{\pgfqpoint{0.000000in}{0.000000in}}%
\pgfpathlineto{\pgfqpoint{0.000000in}{-0.048611in}}%
\pgfusepath{stroke,fill}%
}%
\begin{pgfscope}%
\pgfsys@transformshift{9.386091in}{4.050000in}%
\pgfsys@useobject{currentmarker}{}%
\end{pgfscope}%
\end{pgfscope}%
\begin{pgfscope}%
\pgfsetbuttcap%
\pgfsetroundjoin%
\definecolor{currentfill}{rgb}{0.090196,0.168627,0.227451}%
\pgfsetfillcolor{currentfill}%
\pgfsetlinewidth{0.000000pt}%
\definecolor{currentstroke}{rgb}{0.090196,0.168627,0.227451}%
\pgfsetstrokecolor{currentstroke}%
\pgfsetdash{}{0pt}%
\pgfpathmoveto{\pgfqpoint{9.244333in}{3.902182in}}%
\pgfpathlineto{\pgfqpoint{9.322575in}{3.902182in}}%
\pgfpathlineto{\pgfqpoint{9.322575in}{3.891811in}}%
\pgfpathlineto{\pgfqpoint{9.244333in}{3.891811in}}%
\pgfpathlineto{\pgfqpoint{9.244333in}{3.902182in}}%
\pgfpathlineto{\pgfqpoint{9.244333in}{3.902182in}}%
\pgfpathclose%
\pgfpathmoveto{\pgfqpoint{9.359821in}{3.868178in}}%
\pgfpathlineto{\pgfqpoint{9.402849in}{3.868178in}}%
\pgfpathlineto{\pgfqpoint{9.402849in}{3.857807in}}%
\pgfpathlineto{\pgfqpoint{9.344997in}{3.857807in}}%
\pgfpathlineto{\pgfqpoint{9.344997in}{3.868178in}}%
\pgfpathquadraticcurveto{\pgfqpoint{9.352009in}{3.875444in}}{\pgfqpoint{9.364118in}{3.887670in}}%
\pgfpathquadraticcurveto{\pgfqpoint{9.376247in}{3.899916in}}{\pgfqpoint{9.379353in}{3.903471in}}%
\pgfpathquadraticcurveto{\pgfqpoint{9.385271in}{3.910112in}}{\pgfqpoint{9.387614in}{3.914721in}}%
\pgfpathquadraticcurveto{\pgfqpoint{9.389978in}{3.919331in}}{\pgfqpoint{9.389978in}{3.923784in}}%
\pgfpathquadraticcurveto{\pgfqpoint{9.389978in}{3.931049in}}{\pgfqpoint{9.384880in}{3.935620in}}%
\pgfpathquadraticcurveto{\pgfqpoint{9.379782in}{3.940209in}}{\pgfqpoint{9.371599in}{3.940209in}}%
\pgfpathquadraticcurveto{\pgfqpoint{9.365798in}{3.940209in}}{\pgfqpoint{9.359353in}{3.938198in}}%
\pgfpathquadraticcurveto{\pgfqpoint{9.352927in}{3.936186in}}{\pgfqpoint{9.345603in}{3.932084in}}%
\pgfpathlineto{\pgfqpoint{9.345603in}{3.944545in}}%
\pgfpathquadraticcurveto{\pgfqpoint{9.353044in}{3.947534in}}{\pgfqpoint{9.359509in}{3.949057in}}%
\pgfpathquadraticcurveto{\pgfqpoint{9.365993in}{3.950581in}}{\pgfqpoint{9.371364in}{3.950581in}}%
\pgfpathquadraticcurveto{\pgfqpoint{9.385525in}{3.950581in}}{\pgfqpoint{9.393942in}{3.943491in}}%
\pgfpathquadraticcurveto{\pgfqpoint{9.402360in}{3.936420in}}{\pgfqpoint{9.402360in}{3.924584in}}%
\pgfpathquadraticcurveto{\pgfqpoint{9.402360in}{3.918959in}}{\pgfqpoint{9.400251in}{3.913920in}}%
\pgfpathquadraticcurveto{\pgfqpoint{9.398161in}{3.908901in}}{\pgfqpoint{9.392595in}{3.902065in}}%
\pgfpathquadraticcurveto{\pgfqpoint{9.391071in}{3.900288in}}{\pgfqpoint{9.382888in}{3.891831in}}%
\pgfpathquadraticcurveto{\pgfqpoint{9.374724in}{3.883373in}}{\pgfqpoint{9.359821in}{3.868178in}}%
\pgfpathlineto{\pgfqpoint{9.359821in}{3.868178in}}%
\pgfpathclose%
\pgfpathmoveto{\pgfqpoint{9.428728in}{3.873315in}}%
\pgfpathlineto{\pgfqpoint{9.441618in}{3.873315in}}%
\pgfpathlineto{\pgfqpoint{9.441618in}{3.857807in}}%
\pgfpathlineto{\pgfqpoint{9.428728in}{3.857807in}}%
\pgfpathlineto{\pgfqpoint{9.428728in}{3.873315in}}%
\pgfpathlineto{\pgfqpoint{9.428728in}{3.873315in}}%
\pgfpathclose%
\pgfpathmoveto{\pgfqpoint{9.468591in}{3.948940in}}%
\pgfpathlineto{\pgfqpoint{9.516989in}{3.948940in}}%
\pgfpathlineto{\pgfqpoint{9.516989in}{3.938549in}}%
\pgfpathlineto{\pgfqpoint{9.479880in}{3.938549in}}%
\pgfpathlineto{\pgfqpoint{9.479880in}{3.916225in}}%
\pgfpathquadraticcurveto{\pgfqpoint{9.482556in}{3.917143in}}{\pgfqpoint{9.485232in}{3.917592in}}%
\pgfpathquadraticcurveto{\pgfqpoint{9.487927in}{3.918041in}}{\pgfqpoint{9.490622in}{3.918041in}}%
\pgfpathquadraticcurveto{\pgfqpoint{9.505876in}{3.918041in}}{\pgfqpoint{9.514782in}{3.909682in}}%
\pgfpathquadraticcurveto{\pgfqpoint{9.523708in}{3.901323in}}{\pgfqpoint{9.523708in}{3.887045in}}%
\pgfpathquadraticcurveto{\pgfqpoint{9.523708in}{3.872338in}}{\pgfqpoint{9.514548in}{3.864174in}}%
\pgfpathquadraticcurveto{\pgfqpoint{9.505388in}{3.856030in}}{\pgfqpoint{9.488728in}{3.856030in}}%
\pgfpathquadraticcurveto{\pgfqpoint{9.482985in}{3.856030in}}{\pgfqpoint{9.477028in}{3.857006in}}%
\pgfpathquadraticcurveto{\pgfqpoint{9.471091in}{3.857983in}}{\pgfqpoint{9.464743in}{3.859936in}}%
\pgfpathlineto{\pgfqpoint{9.464743in}{3.872338in}}%
\pgfpathquadraticcurveto{\pgfqpoint{9.470232in}{3.869350in}}{\pgfqpoint{9.476091in}{3.867885in}}%
\pgfpathquadraticcurveto{\pgfqpoint{9.481950in}{3.866420in}}{\pgfqpoint{9.488474in}{3.866420in}}%
\pgfpathquadraticcurveto{\pgfqpoint{9.499040in}{3.866420in}}{\pgfqpoint{9.505192in}{3.871967in}}%
\pgfpathquadraticcurveto{\pgfqpoint{9.511364in}{3.877514in}}{\pgfqpoint{9.511364in}{3.887045in}}%
\pgfpathquadraticcurveto{\pgfqpoint{9.511364in}{3.896557in}}{\pgfqpoint{9.505192in}{3.902104in}}%
\pgfpathquadraticcurveto{\pgfqpoint{9.499040in}{3.907670in}}{\pgfqpoint{9.488474in}{3.907670in}}%
\pgfpathquadraticcurveto{\pgfqpoint{9.483532in}{3.907670in}}{\pgfqpoint{9.478610in}{3.906577in}}%
\pgfpathquadraticcurveto{\pgfqpoint{9.473708in}{3.905483in}}{\pgfqpoint{9.468591in}{3.903159in}}%
\pgfpathlineto{\pgfqpoint{9.468591in}{3.948940in}}%
\pgfpathlineto{\pgfqpoint{9.468591in}{3.948940in}}%
\pgfpathclose%
\pgfusepath{fill}%
\end{pgfscope}%
\begin{pgfscope}%
\pgfpathrectangle{\pgfqpoint{8.036091in}{4.050000in}}{\pgfqpoint{5.400000in}{5.400000in}}%
\pgfusepath{clip}%
\pgfsetrectcap%
\pgfsetroundjoin%
\pgfsetlinewidth{0.602250pt}%
\definecolor{currentstroke}{rgb}{0.898039,0.913725,0.929412}%
\pgfsetstrokecolor{currentstroke}%
\pgfsetdash{}{0pt}%
\pgfpathmoveto{\pgfqpoint{10.736091in}{4.050000in}}%
\pgfpathlineto{\pgfqpoint{10.736091in}{9.450000in}}%
\pgfusepath{stroke}%
\end{pgfscope}%
\begin{pgfscope}%
\pgfsetbuttcap%
\pgfsetroundjoin%
\definecolor{currentfill}{rgb}{0.090196,0.168627,0.227451}%
\pgfsetfillcolor{currentfill}%
\pgfsetlinewidth{0.803000pt}%
\definecolor{currentstroke}{rgb}{0.090196,0.168627,0.227451}%
\pgfsetstrokecolor{currentstroke}%
\pgfsetdash{}{0pt}%
\pgfsys@defobject{currentmarker}{\pgfqpoint{0.000000in}{-0.048611in}}{\pgfqpoint{0.000000in}{0.000000in}}{%
\pgfpathmoveto{\pgfqpoint{0.000000in}{0.000000in}}%
\pgfpathlineto{\pgfqpoint{0.000000in}{-0.048611in}}%
\pgfusepath{stroke,fill}%
}%
\begin{pgfscope}%
\pgfsys@transformshift{10.736091in}{4.050000in}%
\pgfsys@useobject{currentmarker}{}%
\end{pgfscope}%
\end{pgfscope}%
\begin{pgfscope}%
\pgfsetbuttcap%
\pgfsetroundjoin%
\definecolor{currentfill}{rgb}{0.090196,0.168627,0.227451}%
\pgfsetfillcolor{currentfill}%
\pgfsetlinewidth{0.000000pt}%
\definecolor{currentstroke}{rgb}{0.090196,0.168627,0.227451}%
\pgfsetstrokecolor{currentstroke}%
\pgfsetdash{}{0pt}%
\pgfpathmoveto{\pgfqpoint{10.675817in}{3.940815in}}%
\pgfpathquadraticcurveto{\pgfqpoint{10.666306in}{3.940815in}}{\pgfqpoint{10.661501in}{3.931440in}}%
\pgfpathquadraticcurveto{\pgfqpoint{10.656716in}{3.922084in}}{\pgfqpoint{10.656716in}{3.903276in}}%
\pgfpathquadraticcurveto{\pgfqpoint{10.656716in}{3.884545in}}{\pgfqpoint{10.661501in}{3.875170in}}%
\pgfpathquadraticcurveto{\pgfqpoint{10.666306in}{3.865795in}}{\pgfqpoint{10.675817in}{3.865795in}}%
\pgfpathquadraticcurveto{\pgfqpoint{10.685407in}{3.865795in}}{\pgfqpoint{10.690192in}{3.875170in}}%
\pgfpathquadraticcurveto{\pgfqpoint{10.694997in}{3.884545in}}{\pgfqpoint{10.694997in}{3.903276in}}%
\pgfpathquadraticcurveto{\pgfqpoint{10.694997in}{3.922084in}}{\pgfqpoint{10.690192in}{3.931440in}}%
\pgfpathquadraticcurveto{\pgfqpoint{10.685407in}{3.940815in}}{\pgfqpoint{10.675817in}{3.940815in}}%
\pgfpathlineto{\pgfqpoint{10.675817in}{3.940815in}}%
\pgfpathclose%
\pgfpathmoveto{\pgfqpoint{10.675817in}{3.950581in}}%
\pgfpathquadraticcurveto{\pgfqpoint{10.691150in}{3.950581in}}{\pgfqpoint{10.699235in}{3.938452in}}%
\pgfpathquadraticcurveto{\pgfqpoint{10.707321in}{3.926342in}}{\pgfqpoint{10.707321in}{3.903276in}}%
\pgfpathquadraticcurveto{\pgfqpoint{10.707321in}{3.880268in}}{\pgfqpoint{10.699235in}{3.868139in}}%
\pgfpathquadraticcurveto{\pgfqpoint{10.691150in}{3.856030in}}{\pgfqpoint{10.675817in}{3.856030in}}%
\pgfpathquadraticcurveto{\pgfqpoint{10.660505in}{3.856030in}}{\pgfqpoint{10.652419in}{3.868139in}}%
\pgfpathquadraticcurveto{\pgfqpoint{10.644333in}{3.880268in}}{\pgfqpoint{10.644333in}{3.903276in}}%
\pgfpathquadraticcurveto{\pgfqpoint{10.644333in}{3.926342in}}{\pgfqpoint{10.652419in}{3.938452in}}%
\pgfpathquadraticcurveto{\pgfqpoint{10.660505in}{3.950581in}}{\pgfqpoint{10.675817in}{3.950581in}}%
\pgfpathlineto{\pgfqpoint{10.675817in}{3.950581in}}%
\pgfpathclose%
\pgfpathmoveto{\pgfqpoint{10.728982in}{3.873315in}}%
\pgfpathlineto{\pgfqpoint{10.741872in}{3.873315in}}%
\pgfpathlineto{\pgfqpoint{10.741872in}{3.857807in}}%
\pgfpathlineto{\pgfqpoint{10.728982in}{3.857807in}}%
\pgfpathlineto{\pgfqpoint{10.728982in}{3.873315in}}%
\pgfpathlineto{\pgfqpoint{10.728982in}{3.873315in}}%
\pgfpathclose%
\pgfpathmoveto{\pgfqpoint{10.795075in}{3.940815in}}%
\pgfpathquadraticcurveto{\pgfqpoint{10.785564in}{3.940815in}}{\pgfqpoint{10.780759in}{3.931440in}}%
\pgfpathquadraticcurveto{\pgfqpoint{10.775974in}{3.922084in}}{\pgfqpoint{10.775974in}{3.903276in}}%
\pgfpathquadraticcurveto{\pgfqpoint{10.775974in}{3.884545in}}{\pgfqpoint{10.780759in}{3.875170in}}%
\pgfpathquadraticcurveto{\pgfqpoint{10.785564in}{3.865795in}}{\pgfqpoint{10.795075in}{3.865795in}}%
\pgfpathquadraticcurveto{\pgfqpoint{10.804665in}{3.865795in}}{\pgfqpoint{10.809450in}{3.875170in}}%
\pgfpathquadraticcurveto{\pgfqpoint{10.814255in}{3.884545in}}{\pgfqpoint{10.814255in}{3.903276in}}%
\pgfpathquadraticcurveto{\pgfqpoint{10.814255in}{3.922084in}}{\pgfqpoint{10.809450in}{3.931440in}}%
\pgfpathquadraticcurveto{\pgfqpoint{10.804665in}{3.940815in}}{\pgfqpoint{10.795075in}{3.940815in}}%
\pgfpathlineto{\pgfqpoint{10.795075in}{3.940815in}}%
\pgfpathclose%
\pgfpathmoveto{\pgfqpoint{10.795075in}{3.950581in}}%
\pgfpathquadraticcurveto{\pgfqpoint{10.810407in}{3.950581in}}{\pgfqpoint{10.818493in}{3.938452in}}%
\pgfpathquadraticcurveto{\pgfqpoint{10.826579in}{3.926342in}}{\pgfqpoint{10.826579in}{3.903276in}}%
\pgfpathquadraticcurveto{\pgfqpoint{10.826579in}{3.880268in}}{\pgfqpoint{10.818493in}{3.868139in}}%
\pgfpathquadraticcurveto{\pgfqpoint{10.810407in}{3.856030in}}{\pgfqpoint{10.795075in}{3.856030in}}%
\pgfpathquadraticcurveto{\pgfqpoint{10.779763in}{3.856030in}}{\pgfqpoint{10.771677in}{3.868139in}}%
\pgfpathquadraticcurveto{\pgfqpoint{10.763591in}{3.880268in}}{\pgfqpoint{10.763591in}{3.903276in}}%
\pgfpathquadraticcurveto{\pgfqpoint{10.763591in}{3.926342in}}{\pgfqpoint{10.771677in}{3.938452in}}%
\pgfpathquadraticcurveto{\pgfqpoint{10.779763in}{3.950581in}}{\pgfqpoint{10.795075in}{3.950581in}}%
\pgfpathlineto{\pgfqpoint{10.795075in}{3.950581in}}%
\pgfpathclose%
\pgfusepath{fill}%
\end{pgfscope}%
\begin{pgfscope}%
\pgfpathrectangle{\pgfqpoint{8.036091in}{4.050000in}}{\pgfqpoint{5.400000in}{5.400000in}}%
\pgfusepath{clip}%
\pgfsetrectcap%
\pgfsetroundjoin%
\pgfsetlinewidth{0.602250pt}%
\definecolor{currentstroke}{rgb}{0.898039,0.913725,0.929412}%
\pgfsetstrokecolor{currentstroke}%
\pgfsetdash{}{0pt}%
\pgfpathmoveto{\pgfqpoint{12.086091in}{4.050000in}}%
\pgfpathlineto{\pgfqpoint{12.086091in}{9.450000in}}%
\pgfusepath{stroke}%
\end{pgfscope}%
\begin{pgfscope}%
\pgfsetbuttcap%
\pgfsetroundjoin%
\definecolor{currentfill}{rgb}{0.090196,0.168627,0.227451}%
\pgfsetfillcolor{currentfill}%
\pgfsetlinewidth{0.803000pt}%
\definecolor{currentstroke}{rgb}{0.090196,0.168627,0.227451}%
\pgfsetstrokecolor{currentstroke}%
\pgfsetdash{}{0pt}%
\pgfsys@defobject{currentmarker}{\pgfqpoint{0.000000in}{-0.048611in}}{\pgfqpoint{0.000000in}{0.000000in}}{%
\pgfpathmoveto{\pgfqpoint{0.000000in}{0.000000in}}%
\pgfpathlineto{\pgfqpoint{0.000000in}{-0.048611in}}%
\pgfusepath{stroke,fill}%
}%
\begin{pgfscope}%
\pgfsys@transformshift{12.086091in}{4.050000in}%
\pgfsys@useobject{currentmarker}{}%
\end{pgfscope}%
\end{pgfscope}%
\begin{pgfscope}%
\pgfsetbuttcap%
\pgfsetroundjoin%
\definecolor{currentfill}{rgb}{0.090196,0.168627,0.227451}%
\pgfsetfillcolor{currentfill}%
\pgfsetlinewidth{0.000000pt}%
\definecolor{currentstroke}{rgb}{0.090196,0.168627,0.227451}%
\pgfsetstrokecolor{currentstroke}%
\pgfsetdash{}{0pt}%
\pgfpathmoveto{\pgfqpoint{12.010075in}{3.868178in}}%
\pgfpathlineto{\pgfqpoint{12.053103in}{3.868178in}}%
\pgfpathlineto{\pgfqpoint{12.053103in}{3.857807in}}%
\pgfpathlineto{\pgfqpoint{11.995251in}{3.857807in}}%
\pgfpathlineto{\pgfqpoint{11.995251in}{3.868178in}}%
\pgfpathquadraticcurveto{\pgfqpoint{12.002263in}{3.875444in}}{\pgfqpoint{12.014372in}{3.887670in}}%
\pgfpathquadraticcurveto{\pgfqpoint{12.026501in}{3.899916in}}{\pgfqpoint{12.029607in}{3.903471in}}%
\pgfpathquadraticcurveto{\pgfqpoint{12.035525in}{3.910112in}}{\pgfqpoint{12.037868in}{3.914721in}}%
\pgfpathquadraticcurveto{\pgfqpoint{12.040232in}{3.919331in}}{\pgfqpoint{12.040232in}{3.923784in}}%
\pgfpathquadraticcurveto{\pgfqpoint{12.040232in}{3.931049in}}{\pgfqpoint{12.035134in}{3.935620in}}%
\pgfpathquadraticcurveto{\pgfqpoint{12.030036in}{3.940209in}}{\pgfqpoint{12.021853in}{3.940209in}}%
\pgfpathquadraticcurveto{\pgfqpoint{12.016052in}{3.940209in}}{\pgfqpoint{12.009607in}{3.938198in}}%
\pgfpathquadraticcurveto{\pgfqpoint{12.003181in}{3.936186in}}{\pgfqpoint{11.995857in}{3.932084in}}%
\pgfpathlineto{\pgfqpoint{11.995857in}{3.944545in}}%
\pgfpathquadraticcurveto{\pgfqpoint{12.003298in}{3.947534in}}{\pgfqpoint{12.009763in}{3.949057in}}%
\pgfpathquadraticcurveto{\pgfqpoint{12.016247in}{3.950581in}}{\pgfqpoint{12.021618in}{3.950581in}}%
\pgfpathquadraticcurveto{\pgfqpoint{12.035778in}{3.950581in}}{\pgfqpoint{12.044196in}{3.943491in}}%
\pgfpathquadraticcurveto{\pgfqpoint{12.052614in}{3.936420in}}{\pgfqpoint{12.052614in}{3.924584in}}%
\pgfpathquadraticcurveto{\pgfqpoint{12.052614in}{3.918959in}}{\pgfqpoint{12.050505in}{3.913920in}}%
\pgfpathquadraticcurveto{\pgfqpoint{12.048415in}{3.908901in}}{\pgfqpoint{12.042849in}{3.902065in}}%
\pgfpathquadraticcurveto{\pgfqpoint{12.041325in}{3.900288in}}{\pgfqpoint{12.033142in}{3.891831in}}%
\pgfpathquadraticcurveto{\pgfqpoint{12.024978in}{3.883373in}}{\pgfqpoint{12.010075in}{3.868178in}}%
\pgfpathlineto{\pgfqpoint{12.010075in}{3.868178in}}%
\pgfpathclose%
\pgfpathmoveto{\pgfqpoint{12.078982in}{3.873315in}}%
\pgfpathlineto{\pgfqpoint{12.091872in}{3.873315in}}%
\pgfpathlineto{\pgfqpoint{12.091872in}{3.857807in}}%
\pgfpathlineto{\pgfqpoint{12.078982in}{3.857807in}}%
\pgfpathlineto{\pgfqpoint{12.078982in}{3.873315in}}%
\pgfpathlineto{\pgfqpoint{12.078982in}{3.873315in}}%
\pgfpathclose%
\pgfpathmoveto{\pgfqpoint{12.118845in}{3.948940in}}%
\pgfpathlineto{\pgfqpoint{12.167243in}{3.948940in}}%
\pgfpathlineto{\pgfqpoint{12.167243in}{3.938549in}}%
\pgfpathlineto{\pgfqpoint{12.130134in}{3.938549in}}%
\pgfpathlineto{\pgfqpoint{12.130134in}{3.916225in}}%
\pgfpathquadraticcurveto{\pgfqpoint{12.132810in}{3.917143in}}{\pgfqpoint{12.135485in}{3.917592in}}%
\pgfpathquadraticcurveto{\pgfqpoint{12.138181in}{3.918041in}}{\pgfqpoint{12.140876in}{3.918041in}}%
\pgfpathquadraticcurveto{\pgfqpoint{12.156130in}{3.918041in}}{\pgfqpoint{12.165036in}{3.909682in}}%
\pgfpathquadraticcurveto{\pgfqpoint{12.173962in}{3.901323in}}{\pgfqpoint{12.173962in}{3.887045in}}%
\pgfpathquadraticcurveto{\pgfqpoint{12.173962in}{3.872338in}}{\pgfqpoint{12.164802in}{3.864174in}}%
\pgfpathquadraticcurveto{\pgfqpoint{12.155642in}{3.856030in}}{\pgfqpoint{12.138982in}{3.856030in}}%
\pgfpathquadraticcurveto{\pgfqpoint{12.133239in}{3.856030in}}{\pgfqpoint{12.127282in}{3.857006in}}%
\pgfpathquadraticcurveto{\pgfqpoint{12.121345in}{3.857983in}}{\pgfqpoint{12.114997in}{3.859936in}}%
\pgfpathlineto{\pgfqpoint{12.114997in}{3.872338in}}%
\pgfpathquadraticcurveto{\pgfqpoint{12.120485in}{3.869350in}}{\pgfqpoint{12.126345in}{3.867885in}}%
\pgfpathquadraticcurveto{\pgfqpoint{12.132204in}{3.866420in}}{\pgfqpoint{12.138728in}{3.866420in}}%
\pgfpathquadraticcurveto{\pgfqpoint{12.149294in}{3.866420in}}{\pgfqpoint{12.155446in}{3.871967in}}%
\pgfpathquadraticcurveto{\pgfqpoint{12.161618in}{3.877514in}}{\pgfqpoint{12.161618in}{3.887045in}}%
\pgfpathquadraticcurveto{\pgfqpoint{12.161618in}{3.896557in}}{\pgfqpoint{12.155446in}{3.902104in}}%
\pgfpathquadraticcurveto{\pgfqpoint{12.149294in}{3.907670in}}{\pgfqpoint{12.138728in}{3.907670in}}%
\pgfpathquadraticcurveto{\pgfqpoint{12.133786in}{3.907670in}}{\pgfqpoint{12.128864in}{3.906577in}}%
\pgfpathquadraticcurveto{\pgfqpoint{12.123962in}{3.905483in}}{\pgfqpoint{12.118845in}{3.903159in}}%
\pgfpathlineto{\pgfqpoint{12.118845in}{3.948940in}}%
\pgfpathlineto{\pgfqpoint{12.118845in}{3.948940in}}%
\pgfpathclose%
\pgfusepath{fill}%
\end{pgfscope}%
\begin{pgfscope}%
\pgfpathrectangle{\pgfqpoint{8.036091in}{4.050000in}}{\pgfqpoint{5.400000in}{5.400000in}}%
\pgfusepath{clip}%
\pgfsetrectcap%
\pgfsetroundjoin%
\pgfsetlinewidth{0.602250pt}%
\definecolor{currentstroke}{rgb}{0.898039,0.913725,0.929412}%
\pgfsetstrokecolor{currentstroke}%
\pgfsetdash{}{0pt}%
\pgfpathmoveto{\pgfqpoint{13.436091in}{4.050000in}}%
\pgfpathlineto{\pgfqpoint{13.436091in}{9.450000in}}%
\pgfusepath{stroke}%
\end{pgfscope}%
\begin{pgfscope}%
\pgfsetbuttcap%
\pgfsetroundjoin%
\definecolor{currentfill}{rgb}{0.090196,0.168627,0.227451}%
\pgfsetfillcolor{currentfill}%
\pgfsetlinewidth{0.803000pt}%
\definecolor{currentstroke}{rgb}{0.090196,0.168627,0.227451}%
\pgfsetstrokecolor{currentstroke}%
\pgfsetdash{}{0pt}%
\pgfsys@defobject{currentmarker}{\pgfqpoint{0.000000in}{-0.048611in}}{\pgfqpoint{0.000000in}{0.000000in}}{%
\pgfpathmoveto{\pgfqpoint{0.000000in}{0.000000in}}%
\pgfpathlineto{\pgfqpoint{0.000000in}{-0.048611in}}%
\pgfusepath{stroke,fill}%
}%
\begin{pgfscope}%
\pgfsys@transformshift{13.436091in}{4.050000in}%
\pgfsys@useobject{currentmarker}{}%
\end{pgfscope}%
\end{pgfscope}%
\begin{pgfscope}%
\pgfsetbuttcap%
\pgfsetroundjoin%
\definecolor{currentfill}{rgb}{0.090196,0.168627,0.227451}%
\pgfsetfillcolor{currentfill}%
\pgfsetlinewidth{0.000000pt}%
\definecolor{currentstroke}{rgb}{0.090196,0.168627,0.227451}%
\pgfsetstrokecolor{currentstroke}%
\pgfsetdash{}{0pt}%
\pgfpathmoveto{\pgfqpoint{13.349587in}{3.948940in}}%
\pgfpathlineto{\pgfqpoint{13.397985in}{3.948940in}}%
\pgfpathlineto{\pgfqpoint{13.397985in}{3.938549in}}%
\pgfpathlineto{\pgfqpoint{13.360876in}{3.938549in}}%
\pgfpathlineto{\pgfqpoint{13.360876in}{3.916225in}}%
\pgfpathquadraticcurveto{\pgfqpoint{13.363552in}{3.917143in}}{\pgfqpoint{13.366228in}{3.917592in}}%
\pgfpathquadraticcurveto{\pgfqpoint{13.368923in}{3.918041in}}{\pgfqpoint{13.371618in}{3.918041in}}%
\pgfpathquadraticcurveto{\pgfqpoint{13.386872in}{3.918041in}}{\pgfqpoint{13.395778in}{3.909682in}}%
\pgfpathquadraticcurveto{\pgfqpoint{13.404704in}{3.901323in}}{\pgfqpoint{13.404704in}{3.887045in}}%
\pgfpathquadraticcurveto{\pgfqpoint{13.404704in}{3.872338in}}{\pgfqpoint{13.395544in}{3.864174in}}%
\pgfpathquadraticcurveto{\pgfqpoint{13.386384in}{3.856030in}}{\pgfqpoint{13.369724in}{3.856030in}}%
\pgfpathquadraticcurveto{\pgfqpoint{13.363982in}{3.856030in}}{\pgfqpoint{13.358025in}{3.857006in}}%
\pgfpathquadraticcurveto{\pgfqpoint{13.352087in}{3.857983in}}{\pgfqpoint{13.345739in}{3.859936in}}%
\pgfpathlineto{\pgfqpoint{13.345739in}{3.872338in}}%
\pgfpathquadraticcurveto{\pgfqpoint{13.351228in}{3.869350in}}{\pgfqpoint{13.357087in}{3.867885in}}%
\pgfpathquadraticcurveto{\pgfqpoint{13.362946in}{3.866420in}}{\pgfqpoint{13.369470in}{3.866420in}}%
\pgfpathquadraticcurveto{\pgfqpoint{13.380036in}{3.866420in}}{\pgfqpoint{13.386189in}{3.871967in}}%
\pgfpathquadraticcurveto{\pgfqpoint{13.392360in}{3.877514in}}{\pgfqpoint{13.392360in}{3.887045in}}%
\pgfpathquadraticcurveto{\pgfqpoint{13.392360in}{3.896557in}}{\pgfqpoint{13.386189in}{3.902104in}}%
\pgfpathquadraticcurveto{\pgfqpoint{13.380036in}{3.907670in}}{\pgfqpoint{13.369470in}{3.907670in}}%
\pgfpathquadraticcurveto{\pgfqpoint{13.364528in}{3.907670in}}{\pgfqpoint{13.359607in}{3.906577in}}%
\pgfpathquadraticcurveto{\pgfqpoint{13.354704in}{3.905483in}}{\pgfqpoint{13.349587in}{3.903159in}}%
\pgfpathlineto{\pgfqpoint{13.349587in}{3.948940in}}%
\pgfpathlineto{\pgfqpoint{13.349587in}{3.948940in}}%
\pgfpathclose%
\pgfpathmoveto{\pgfqpoint{13.428982in}{3.873315in}}%
\pgfpathlineto{\pgfqpoint{13.441872in}{3.873315in}}%
\pgfpathlineto{\pgfqpoint{13.441872in}{3.857807in}}%
\pgfpathlineto{\pgfqpoint{13.428982in}{3.857807in}}%
\pgfpathlineto{\pgfqpoint{13.428982in}{3.873315in}}%
\pgfpathlineto{\pgfqpoint{13.428982in}{3.873315in}}%
\pgfpathclose%
\pgfpathmoveto{\pgfqpoint{13.495075in}{3.940815in}}%
\pgfpathquadraticcurveto{\pgfqpoint{13.485564in}{3.940815in}}{\pgfqpoint{13.480759in}{3.931440in}}%
\pgfpathquadraticcurveto{\pgfqpoint{13.475974in}{3.922084in}}{\pgfqpoint{13.475974in}{3.903276in}}%
\pgfpathquadraticcurveto{\pgfqpoint{13.475974in}{3.884545in}}{\pgfqpoint{13.480759in}{3.875170in}}%
\pgfpathquadraticcurveto{\pgfqpoint{13.485564in}{3.865795in}}{\pgfqpoint{13.495075in}{3.865795in}}%
\pgfpathquadraticcurveto{\pgfqpoint{13.504665in}{3.865795in}}{\pgfqpoint{13.509450in}{3.875170in}}%
\pgfpathquadraticcurveto{\pgfqpoint{13.514255in}{3.884545in}}{\pgfqpoint{13.514255in}{3.903276in}}%
\pgfpathquadraticcurveto{\pgfqpoint{13.514255in}{3.922084in}}{\pgfqpoint{13.509450in}{3.931440in}}%
\pgfpathquadraticcurveto{\pgfqpoint{13.504665in}{3.940815in}}{\pgfqpoint{13.495075in}{3.940815in}}%
\pgfpathlineto{\pgfqpoint{13.495075in}{3.940815in}}%
\pgfpathclose%
\pgfpathmoveto{\pgfqpoint{13.495075in}{3.950581in}}%
\pgfpathquadraticcurveto{\pgfqpoint{13.510407in}{3.950581in}}{\pgfqpoint{13.518493in}{3.938452in}}%
\pgfpathquadraticcurveto{\pgfqpoint{13.526579in}{3.926342in}}{\pgfqpoint{13.526579in}{3.903276in}}%
\pgfpathquadraticcurveto{\pgfqpoint{13.526579in}{3.880268in}}{\pgfqpoint{13.518493in}{3.868139in}}%
\pgfpathquadraticcurveto{\pgfqpoint{13.510407in}{3.856030in}}{\pgfqpoint{13.495075in}{3.856030in}}%
\pgfpathquadraticcurveto{\pgfqpoint{13.479763in}{3.856030in}}{\pgfqpoint{13.471677in}{3.868139in}}%
\pgfpathquadraticcurveto{\pgfqpoint{13.463591in}{3.880268in}}{\pgfqpoint{13.463591in}{3.903276in}}%
\pgfpathquadraticcurveto{\pgfqpoint{13.463591in}{3.926342in}}{\pgfqpoint{13.471677in}{3.938452in}}%
\pgfpathquadraticcurveto{\pgfqpoint{13.479763in}{3.950581in}}{\pgfqpoint{13.495075in}{3.950581in}}%
\pgfpathlineto{\pgfqpoint{13.495075in}{3.950581in}}%
\pgfpathclose%
\pgfusepath{fill}%
\end{pgfscope}%
\begin{pgfscope}%
\pgfsetbuttcap%
\pgfsetroundjoin%
\definecolor{currentfill}{rgb}{0.090196,0.168627,0.227451}%
\pgfsetfillcolor{currentfill}%
\pgfsetlinewidth{0.000000pt}%
\definecolor{currentstroke}{rgb}{0.090196,0.168627,0.227451}%
\pgfsetstrokecolor{currentstroke}%
\pgfsetdash{}{0pt}%
\pgfpathmoveto{\pgfqpoint{9.715637in}{3.698087in}}%
\pgfpathlineto{\pgfqpoint{9.809858in}{3.698087in}}%
\pgfpathlineto{\pgfqpoint{9.809858in}{3.685387in}}%
\pgfpathlineto{\pgfqpoint{9.770327in}{3.685387in}}%
\pgfpathlineto{\pgfqpoint{9.770327in}{3.586702in}}%
\pgfpathlineto{\pgfqpoint{9.755192in}{3.586702in}}%
\pgfpathlineto{\pgfqpoint{9.755192in}{3.685387in}}%
\pgfpathlineto{\pgfqpoint{9.715637in}{3.685387in}}%
\pgfpathlineto{\pgfqpoint{9.715637in}{3.698087in}}%
\pgfpathlineto{\pgfqpoint{9.715637in}{3.698087in}}%
\pgfpathclose%
\pgfpathmoveto{\pgfqpoint{9.814489in}{3.698087in}}%
\pgfpathlineto{\pgfqpoint{9.829696in}{3.698087in}}%
\pgfpathlineto{\pgfqpoint{9.853113in}{3.603938in}}%
\pgfpathlineto{\pgfqpoint{9.876460in}{3.698087in}}%
\pgfpathlineto{\pgfqpoint{9.893409in}{3.698087in}}%
\pgfpathlineto{\pgfqpoint{9.916827in}{3.603938in}}%
\pgfpathlineto{\pgfqpoint{9.940173in}{3.698087in}}%
\pgfpathlineto{\pgfqpoint{9.955475in}{3.698087in}}%
\pgfpathlineto{\pgfqpoint{9.927497in}{3.586702in}}%
\pgfpathlineto{\pgfqpoint{9.908543in}{3.586702in}}%
\pgfpathlineto{\pgfqpoint{9.885054in}{3.683382in}}%
\pgfpathlineto{\pgfqpoint{9.861325in}{3.586702in}}%
\pgfpathlineto{\pgfqpoint{9.842371in}{3.586702in}}%
\pgfpathlineto{\pgfqpoint{9.814489in}{3.698087in}}%
\pgfpathlineto{\pgfqpoint{9.814489in}{3.698087in}}%
\pgfpathclose%
\pgfpathmoveto{\pgfqpoint{9.975455in}{3.698087in}}%
\pgfpathlineto{\pgfqpoint{10.039455in}{3.698087in}}%
\pgfpathlineto{\pgfqpoint{10.039455in}{3.685387in}}%
\pgfpathlineto{\pgfqpoint{9.990518in}{3.685387in}}%
\pgfpathlineto{\pgfqpoint{9.990518in}{3.652564in}}%
\pgfpathlineto{\pgfqpoint{10.034680in}{3.652564in}}%
\pgfpathlineto{\pgfqpoint{10.034680in}{3.639888in}}%
\pgfpathlineto{\pgfqpoint{9.990518in}{3.639888in}}%
\pgfpathlineto{\pgfqpoint{9.990518in}{3.586702in}}%
\pgfpathlineto{\pgfqpoint{9.975455in}{3.586702in}}%
\pgfpathlineto{\pgfqpoint{9.975455in}{3.698087in}}%
\pgfpathlineto{\pgfqpoint{9.975455in}{3.698087in}}%
\pgfpathclose%
\pgfpathmoveto{\pgfqpoint{10.063326in}{3.698087in}}%
\pgfpathlineto{\pgfqpoint{10.133747in}{3.698087in}}%
\pgfpathlineto{\pgfqpoint{10.133747in}{3.685387in}}%
\pgfpathlineto{\pgfqpoint{10.078389in}{3.685387in}}%
\pgfpathlineto{\pgfqpoint{10.078389in}{3.652421in}}%
\pgfpathlineto{\pgfqpoint{10.131432in}{3.652421in}}%
\pgfpathlineto{\pgfqpoint{10.131432in}{3.639745in}}%
\pgfpathlineto{\pgfqpoint{10.078389in}{3.639745in}}%
\pgfpathlineto{\pgfqpoint{10.078389in}{3.599378in}}%
\pgfpathlineto{\pgfqpoint{10.135084in}{3.599378in}}%
\pgfpathlineto{\pgfqpoint{10.135084in}{3.586702in}}%
\pgfpathlineto{\pgfqpoint{10.063326in}{3.586702in}}%
\pgfpathlineto{\pgfqpoint{10.063326in}{3.698087in}}%
\pgfpathlineto{\pgfqpoint{10.063326in}{3.698087in}}%
\pgfpathclose%
\pgfpathmoveto{\pgfqpoint{10.279292in}{3.631915in}}%
\pgfpathlineto{\pgfqpoint{10.279292in}{3.625207in}}%
\pgfpathlineto{\pgfqpoint{10.216176in}{3.625207in}}%
\pgfpathquadraticcurveto{\pgfqpoint{10.217083in}{3.611028in}}{\pgfqpoint{10.224722in}{3.603604in}}%
\pgfpathquadraticcurveto{\pgfqpoint{10.232360in}{3.596179in}}{\pgfqpoint{10.246015in}{3.596179in}}%
\pgfpathquadraticcurveto{\pgfqpoint{10.253916in}{3.596179in}}{\pgfqpoint{10.261340in}{3.598113in}}%
\pgfpathquadraticcurveto{\pgfqpoint{10.268765in}{3.600047in}}{\pgfqpoint{10.276093in}{3.603938in}}%
\pgfpathlineto{\pgfqpoint{10.276093in}{3.590952in}}%
\pgfpathquadraticcurveto{\pgfqpoint{10.268693in}{3.587824in}}{\pgfqpoint{10.260935in}{3.586177in}}%
\pgfpathquadraticcurveto{\pgfqpoint{10.253176in}{3.584530in}}{\pgfqpoint{10.245203in}{3.584530in}}%
\pgfpathquadraticcurveto{\pgfqpoint{10.225199in}{3.584530in}}{\pgfqpoint{10.213526in}{3.596156in}}%
\pgfpathquadraticcurveto{\pgfqpoint{10.201853in}{3.607805in}}{\pgfqpoint{10.201853in}{3.627666in}}%
\pgfpathquadraticcurveto{\pgfqpoint{10.201853in}{3.648172in}}{\pgfqpoint{10.212929in}{3.660203in}}%
\pgfpathquadraticcurveto{\pgfqpoint{10.224005in}{3.672258in}}{\pgfqpoint{10.242816in}{3.672258in}}%
\pgfpathquadraticcurveto{\pgfqpoint{10.259669in}{3.672258in}}{\pgfqpoint{10.269481in}{3.661396in}}%
\pgfpathquadraticcurveto{\pgfqpoint{10.279292in}{3.650559in}}{\pgfqpoint{10.279292in}{3.631915in}}%
\pgfpathlineto{\pgfqpoint{10.279292in}{3.631915in}}%
\pgfpathclose%
\pgfpathmoveto{\pgfqpoint{10.265566in}{3.635949in}}%
\pgfpathquadraticcurveto{\pgfqpoint{10.265423in}{3.647193in}}{\pgfqpoint{10.259264in}{3.653901in}}%
\pgfpathquadraticcurveto{\pgfqpoint{10.253105in}{3.660633in}}{\pgfqpoint{10.242959in}{3.660633in}}%
\pgfpathquadraticcurveto{\pgfqpoint{10.231477in}{3.660633in}}{\pgfqpoint{10.224578in}{3.654140in}}%
\pgfpathquadraticcurveto{\pgfqpoint{10.217679in}{3.647646in}}{\pgfqpoint{10.216629in}{3.635854in}}%
\pgfpathlineto{\pgfqpoint{10.265566in}{3.635949in}}%
\pgfpathlineto{\pgfqpoint{10.265566in}{3.635949in}}%
\pgfpathclose%
\pgfpathmoveto{\pgfqpoint{10.355084in}{3.667794in}}%
\pgfpathlineto{\pgfqpoint{10.355084in}{3.654808in}}%
\pgfpathquadraticcurveto{\pgfqpoint{10.349283in}{3.657792in}}{\pgfqpoint{10.343005in}{3.659272in}}%
\pgfpathquadraticcurveto{\pgfqpoint{10.336751in}{3.660776in}}{\pgfqpoint{10.330019in}{3.660776in}}%
\pgfpathquadraticcurveto{\pgfqpoint{10.319802in}{3.660776in}}{\pgfqpoint{10.314693in}{3.657649in}}%
\pgfpathquadraticcurveto{\pgfqpoint{10.309585in}{3.654521in}}{\pgfqpoint{10.309585in}{3.648243in}}%
\pgfpathquadraticcurveto{\pgfqpoint{10.309585in}{3.643469in}}{\pgfqpoint{10.313237in}{3.640748in}}%
\pgfpathquadraticcurveto{\pgfqpoint{10.316890in}{3.638026in}}{\pgfqpoint{10.327942in}{3.635567in}}%
\pgfpathlineto{\pgfqpoint{10.332645in}{3.634517in}}%
\pgfpathquadraticcurveto{\pgfqpoint{10.347254in}{3.631390in}}{\pgfqpoint{10.353413in}{3.625685in}}%
\pgfpathquadraticcurveto{\pgfqpoint{10.359572in}{3.619979in}}{\pgfqpoint{10.359572in}{3.609762in}}%
\pgfpathquadraticcurveto{\pgfqpoint{10.359572in}{3.598113in}}{\pgfqpoint{10.350357in}{3.591310in}}%
\pgfpathquadraticcurveto{\pgfqpoint{10.341143in}{3.584530in}}{\pgfqpoint{10.325030in}{3.584530in}}%
\pgfpathquadraticcurveto{\pgfqpoint{10.318322in}{3.584530in}}{\pgfqpoint{10.311041in}{3.585843in}}%
\pgfpathquadraticcurveto{\pgfqpoint{10.303760in}{3.587156in}}{\pgfqpoint{10.295715in}{3.589758in}}%
\pgfpathlineto{\pgfqpoint{10.295715in}{3.603938in}}%
\pgfpathquadraticcurveto{\pgfqpoint{10.303330in}{3.599975in}}{\pgfqpoint{10.310707in}{3.597994in}}%
\pgfpathquadraticcurveto{\pgfqpoint{10.318083in}{3.596036in}}{\pgfqpoint{10.325340in}{3.596036in}}%
\pgfpathquadraticcurveto{\pgfqpoint{10.335032in}{3.596036in}}{\pgfqpoint{10.340236in}{3.599354in}}%
\pgfpathquadraticcurveto{\pgfqpoint{10.345464in}{3.602673in}}{\pgfqpoint{10.345464in}{3.608712in}}%
\pgfpathquadraticcurveto{\pgfqpoint{10.345464in}{3.614298in}}{\pgfqpoint{10.341692in}{3.617282in}}%
\pgfpathquadraticcurveto{\pgfqpoint{10.337944in}{3.620266in}}{\pgfqpoint{10.325173in}{3.623035in}}%
\pgfpathlineto{\pgfqpoint{10.320399in}{3.624157in}}%
\pgfpathquadraticcurveto{\pgfqpoint{10.307651in}{3.626831in}}{\pgfqpoint{10.301970in}{3.632393in}}%
\pgfpathquadraticcurveto{\pgfqpoint{10.296312in}{3.637955in}}{\pgfqpoint{10.296312in}{3.647646in}}%
\pgfpathquadraticcurveto{\pgfqpoint{10.296312in}{3.659439in}}{\pgfqpoint{10.304667in}{3.665837in}}%
\pgfpathquadraticcurveto{\pgfqpoint{10.313022in}{3.672258in}}{\pgfqpoint{10.328396in}{3.672258in}}%
\pgfpathquadraticcurveto{\pgfqpoint{10.335987in}{3.672258in}}{\pgfqpoint{10.342695in}{3.671136in}}%
\pgfpathquadraticcurveto{\pgfqpoint{10.349426in}{3.670038in}}{\pgfqpoint{10.355084in}{3.667794in}}%
\pgfpathlineto{\pgfqpoint{10.355084in}{3.667794in}}%
\pgfpathclose%
\pgfpathmoveto{\pgfqpoint{10.394997in}{3.693981in}}%
\pgfpathlineto{\pgfqpoint{10.394997in}{3.670253in}}%
\pgfpathlineto{\pgfqpoint{10.423261in}{3.670253in}}%
\pgfpathlineto{\pgfqpoint{10.423261in}{3.659582in}}%
\pgfpathlineto{\pgfqpoint{10.394997in}{3.659582in}}%
\pgfpathlineto{\pgfqpoint{10.394997in}{3.614226in}}%
\pgfpathquadraticcurveto{\pgfqpoint{10.394997in}{3.604009in}}{\pgfqpoint{10.397790in}{3.601097in}}%
\pgfpathquadraticcurveto{\pgfqpoint{10.400583in}{3.598185in}}{\pgfqpoint{10.409177in}{3.598185in}}%
\pgfpathlineto{\pgfqpoint{10.423261in}{3.598185in}}%
\pgfpathlineto{\pgfqpoint{10.423261in}{3.586702in}}%
\pgfpathlineto{\pgfqpoint{10.409177in}{3.586702in}}%
\pgfpathquadraticcurveto{\pgfqpoint{10.393278in}{3.586702in}}{\pgfqpoint{10.387239in}{3.592623in}}%
\pgfpathquadraticcurveto{\pgfqpoint{10.381199in}{3.598567in}}{\pgfqpoint{10.381199in}{3.614226in}}%
\pgfpathlineto{\pgfqpoint{10.381199in}{3.659582in}}%
\pgfpathlineto{\pgfqpoint{10.371126in}{3.659582in}}%
\pgfpathlineto{\pgfqpoint{10.371126in}{3.670253in}}%
\pgfpathlineto{\pgfqpoint{10.381199in}{3.670253in}}%
\pgfpathlineto{\pgfqpoint{10.381199in}{3.693981in}}%
\pgfpathlineto{\pgfqpoint{10.394997in}{3.693981in}}%
\pgfpathlineto{\pgfqpoint{10.394997in}{3.693981in}}%
\pgfpathclose%
\pgfpathmoveto{\pgfqpoint{10.441308in}{3.670253in}}%
\pgfpathlineto{\pgfqpoint{10.455034in}{3.670253in}}%
\pgfpathlineto{\pgfqpoint{10.455034in}{3.586702in}}%
\pgfpathlineto{\pgfqpoint{10.441308in}{3.586702in}}%
\pgfpathlineto{\pgfqpoint{10.441308in}{3.670253in}}%
\pgfpathlineto{\pgfqpoint{10.441308in}{3.670253in}}%
\pgfpathclose%
\pgfpathmoveto{\pgfqpoint{10.441308in}{3.702790in}}%
\pgfpathlineto{\pgfqpoint{10.455034in}{3.702790in}}%
\pgfpathlineto{\pgfqpoint{10.455034in}{3.685387in}}%
\pgfpathlineto{\pgfqpoint{10.441308in}{3.685387in}}%
\pgfpathlineto{\pgfqpoint{10.441308in}{3.702790in}}%
\pgfpathlineto{\pgfqpoint{10.441308in}{3.702790in}}%
\pgfpathclose%
\pgfpathmoveto{\pgfqpoint{10.548801in}{3.654211in}}%
\pgfpathquadraticcurveto{\pgfqpoint{10.553958in}{3.663473in}}{\pgfqpoint{10.561119in}{3.667866in}}%
\pgfpathquadraticcurveto{\pgfqpoint{10.568281in}{3.672258in}}{\pgfqpoint{10.577972in}{3.672258in}}%
\pgfpathquadraticcurveto{\pgfqpoint{10.591030in}{3.672258in}}{\pgfqpoint{10.598120in}{3.663115in}}%
\pgfpathquadraticcurveto{\pgfqpoint{10.605210in}{3.653996in}}{\pgfqpoint{10.605210in}{3.637143in}}%
\pgfpathlineto{\pgfqpoint{10.605210in}{3.586702in}}%
\pgfpathlineto{\pgfqpoint{10.591412in}{3.586702in}}%
\pgfpathlineto{\pgfqpoint{10.591412in}{3.636689in}}%
\pgfpathquadraticcurveto{\pgfqpoint{10.591412in}{3.648697in}}{\pgfqpoint{10.587139in}{3.654498in}}%
\pgfpathquadraticcurveto{\pgfqpoint{10.582890in}{3.660322in}}{\pgfqpoint{10.574177in}{3.660322in}}%
\pgfpathquadraticcurveto{\pgfqpoint{10.563506in}{3.660322in}}{\pgfqpoint{10.557300in}{3.653232in}}%
\pgfpathquadraticcurveto{\pgfqpoint{10.551117in}{3.646166in}}{\pgfqpoint{10.551117in}{3.633920in}}%
\pgfpathlineto{\pgfqpoint{10.551117in}{3.586702in}}%
\pgfpathlineto{\pgfqpoint{10.537319in}{3.586702in}}%
\pgfpathlineto{\pgfqpoint{10.537319in}{3.636689in}}%
\pgfpathquadraticcurveto{\pgfqpoint{10.537319in}{3.648768in}}{\pgfqpoint{10.533070in}{3.654545in}}%
\pgfpathquadraticcurveto{\pgfqpoint{10.528821in}{3.660322in}}{\pgfqpoint{10.519941in}{3.660322in}}%
\pgfpathquadraticcurveto{\pgfqpoint{10.509413in}{3.660322in}}{\pgfqpoint{10.503207in}{3.653209in}}%
\pgfpathquadraticcurveto{\pgfqpoint{10.497024in}{3.646095in}}{\pgfqpoint{10.497024in}{3.633920in}}%
\pgfpathlineto{\pgfqpoint{10.497024in}{3.586702in}}%
\pgfpathlineto{\pgfqpoint{10.483226in}{3.586702in}}%
\pgfpathlineto{\pgfqpoint{10.483226in}{3.670253in}}%
\pgfpathlineto{\pgfqpoint{10.497024in}{3.670253in}}%
\pgfpathlineto{\pgfqpoint{10.497024in}{3.657267in}}%
\pgfpathquadraticcurveto{\pgfqpoint{10.501727in}{3.664953in}}{\pgfqpoint{10.508291in}{3.668606in}}%
\pgfpathquadraticcurveto{\pgfqpoint{10.514856in}{3.672258in}}{\pgfqpoint{10.523880in}{3.672258in}}%
\pgfpathquadraticcurveto{\pgfqpoint{10.532998in}{3.672258in}}{\pgfqpoint{10.539372in}{3.667627in}}%
\pgfpathquadraticcurveto{\pgfqpoint{10.545746in}{3.663020in}}{\pgfqpoint{10.548801in}{3.654211in}}%
\pgfpathlineto{\pgfqpoint{10.548801in}{3.654211in}}%
\pgfpathclose%
\pgfpathmoveto{\pgfqpoint{10.670546in}{3.628692in}}%
\pgfpathquadraticcurveto{\pgfqpoint{10.653908in}{3.628692in}}{\pgfqpoint{10.647486in}{3.624897in}}%
\pgfpathquadraticcurveto{\pgfqpoint{10.641065in}{3.621101in}}{\pgfqpoint{10.641065in}{3.611911in}}%
\pgfpathquadraticcurveto{\pgfqpoint{10.641065in}{3.604606in}}{\pgfqpoint{10.645887in}{3.600309in}}%
\pgfpathquadraticcurveto{\pgfqpoint{10.650709in}{3.596036in}}{\pgfqpoint{10.658969in}{3.596036in}}%
\pgfpathquadraticcurveto{\pgfqpoint{10.670403in}{3.596036in}}{\pgfqpoint{10.677302in}{3.604129in}}%
\pgfpathquadraticcurveto{\pgfqpoint{10.684201in}{3.612221in}}{\pgfqpoint{10.684201in}{3.625637in}}%
\pgfpathlineto{\pgfqpoint{10.684201in}{3.628692in}}%
\pgfpathlineto{\pgfqpoint{10.670546in}{3.628692in}}%
\pgfpathlineto{\pgfqpoint{10.670546in}{3.628692in}}%
\pgfpathclose%
\pgfpathmoveto{\pgfqpoint{10.697927in}{3.634374in}}%
\pgfpathlineto{\pgfqpoint{10.697927in}{3.586702in}}%
\pgfpathlineto{\pgfqpoint{10.684201in}{3.586702in}}%
\pgfpathlineto{\pgfqpoint{10.684201in}{3.599378in}}%
\pgfpathquadraticcurveto{\pgfqpoint{10.679498in}{3.591787in}}{\pgfqpoint{10.672480in}{3.588159in}}%
\pgfpathquadraticcurveto{\pgfqpoint{10.665462in}{3.584530in}}{\pgfqpoint{10.655316in}{3.584530in}}%
\pgfpathquadraticcurveto{\pgfqpoint{10.642497in}{3.584530in}}{\pgfqpoint{10.634906in}{3.591739in}}%
\pgfpathquadraticcurveto{\pgfqpoint{10.627339in}{3.598949in}}{\pgfqpoint{10.627339in}{3.611028in}}%
\pgfpathquadraticcurveto{\pgfqpoint{10.627339in}{3.625112in}}{\pgfqpoint{10.636768in}{3.632273in}}%
\pgfpathquadraticcurveto{\pgfqpoint{10.646221in}{3.639435in}}{\pgfqpoint{10.664936in}{3.639435in}}%
\pgfpathlineto{\pgfqpoint{10.684201in}{3.639435in}}%
\pgfpathlineto{\pgfqpoint{10.684201in}{3.640795in}}%
\pgfpathquadraticcurveto{\pgfqpoint{10.684201in}{3.650272in}}{\pgfqpoint{10.677970in}{3.655452in}}%
\pgfpathquadraticcurveto{\pgfqpoint{10.671740in}{3.660633in}}{\pgfqpoint{10.660472in}{3.660633in}}%
\pgfpathquadraticcurveto{\pgfqpoint{10.653311in}{3.660633in}}{\pgfqpoint{10.646508in}{3.658914in}}%
\pgfpathquadraticcurveto{\pgfqpoint{10.639728in}{3.657195in}}{\pgfqpoint{10.633474in}{3.653758in}}%
\pgfpathlineto{\pgfqpoint{10.633474in}{3.666457in}}%
\pgfpathquadraticcurveto{\pgfqpoint{10.640993in}{3.669370in}}{\pgfqpoint{10.648083in}{3.670802in}}%
\pgfpathquadraticcurveto{\pgfqpoint{10.655173in}{3.672258in}}{\pgfqpoint{10.661881in}{3.672258in}}%
\pgfpathquadraticcurveto{\pgfqpoint{10.680023in}{3.672258in}}{\pgfqpoint{10.688975in}{3.662853in}}%
\pgfpathquadraticcurveto{\pgfqpoint{10.697927in}{3.653471in}}{\pgfqpoint{10.697927in}{3.634374in}}%
\pgfpathlineto{\pgfqpoint{10.697927in}{3.634374in}}%
\pgfpathclose%
\pgfpathmoveto{\pgfqpoint{10.739774in}{3.693981in}}%
\pgfpathlineto{\pgfqpoint{10.739774in}{3.670253in}}%
\pgfpathlineto{\pgfqpoint{10.768038in}{3.670253in}}%
\pgfpathlineto{\pgfqpoint{10.768038in}{3.659582in}}%
\pgfpathlineto{\pgfqpoint{10.739774in}{3.659582in}}%
\pgfpathlineto{\pgfqpoint{10.739774in}{3.614226in}}%
\pgfpathquadraticcurveto{\pgfqpoint{10.739774in}{3.604009in}}{\pgfqpoint{10.742567in}{3.601097in}}%
\pgfpathquadraticcurveto{\pgfqpoint{10.745360in}{3.598185in}}{\pgfqpoint{10.753953in}{3.598185in}}%
\pgfpathlineto{\pgfqpoint{10.768038in}{3.598185in}}%
\pgfpathlineto{\pgfqpoint{10.768038in}{3.586702in}}%
\pgfpathlineto{\pgfqpoint{10.753953in}{3.586702in}}%
\pgfpathquadraticcurveto{\pgfqpoint{10.738055in}{3.586702in}}{\pgfqpoint{10.732015in}{3.592623in}}%
\pgfpathquadraticcurveto{\pgfqpoint{10.725976in}{3.598567in}}{\pgfqpoint{10.725976in}{3.614226in}}%
\pgfpathlineto{\pgfqpoint{10.725976in}{3.659582in}}%
\pgfpathlineto{\pgfqpoint{10.715902in}{3.659582in}}%
\pgfpathlineto{\pgfqpoint{10.715902in}{3.670253in}}%
\pgfpathlineto{\pgfqpoint{10.725976in}{3.670253in}}%
\pgfpathlineto{\pgfqpoint{10.725976in}{3.693981in}}%
\pgfpathlineto{\pgfqpoint{10.739774in}{3.693981in}}%
\pgfpathlineto{\pgfqpoint{10.739774in}{3.693981in}}%
\pgfpathclose%
\pgfpathmoveto{\pgfqpoint{10.857556in}{3.631915in}}%
\pgfpathlineto{\pgfqpoint{10.857556in}{3.625207in}}%
\pgfpathlineto{\pgfqpoint{10.794439in}{3.625207in}}%
\pgfpathquadraticcurveto{\pgfqpoint{10.795347in}{3.611028in}}{\pgfqpoint{10.802985in}{3.603604in}}%
\pgfpathquadraticcurveto{\pgfqpoint{10.810624in}{3.596179in}}{\pgfqpoint{10.824279in}{3.596179in}}%
\pgfpathquadraticcurveto{\pgfqpoint{10.832180in}{3.596179in}}{\pgfqpoint{10.839604in}{3.598113in}}%
\pgfpathquadraticcurveto{\pgfqpoint{10.847028in}{3.600047in}}{\pgfqpoint{10.854357in}{3.603938in}}%
\pgfpathlineto{\pgfqpoint{10.854357in}{3.590952in}}%
\pgfpathquadraticcurveto{\pgfqpoint{10.846957in}{3.587824in}}{\pgfqpoint{10.839199in}{3.586177in}}%
\pgfpathquadraticcurveto{\pgfqpoint{10.831440in}{3.584530in}}{\pgfqpoint{10.823467in}{3.584530in}}%
\pgfpathquadraticcurveto{\pgfqpoint{10.803463in}{3.584530in}}{\pgfqpoint{10.791790in}{3.596156in}}%
\pgfpathquadraticcurveto{\pgfqpoint{10.780117in}{3.607805in}}{\pgfqpoint{10.780117in}{3.627666in}}%
\pgfpathquadraticcurveto{\pgfqpoint{10.780117in}{3.648172in}}{\pgfqpoint{10.791193in}{3.660203in}}%
\pgfpathquadraticcurveto{\pgfqpoint{10.802269in}{3.672258in}}{\pgfqpoint{10.821080in}{3.672258in}}%
\pgfpathquadraticcurveto{\pgfqpoint{10.837933in}{3.672258in}}{\pgfqpoint{10.847745in}{3.661396in}}%
\pgfpathquadraticcurveto{\pgfqpoint{10.857556in}{3.650559in}}{\pgfqpoint{10.857556in}{3.631915in}}%
\pgfpathlineto{\pgfqpoint{10.857556in}{3.631915in}}%
\pgfpathclose%
\pgfpathmoveto{\pgfqpoint{10.843830in}{3.635949in}}%
\pgfpathquadraticcurveto{\pgfqpoint{10.843686in}{3.647193in}}{\pgfqpoint{10.837528in}{3.653901in}}%
\pgfpathquadraticcurveto{\pgfqpoint{10.831369in}{3.660633in}}{\pgfqpoint{10.821223in}{3.660633in}}%
\pgfpathquadraticcurveto{\pgfqpoint{10.809741in}{3.660633in}}{\pgfqpoint{10.802842in}{3.654140in}}%
\pgfpathquadraticcurveto{\pgfqpoint{10.795943in}{3.647646in}}{\pgfqpoint{10.794893in}{3.635854in}}%
\pgfpathlineto{\pgfqpoint{10.843830in}{3.635949in}}%
\pgfpathlineto{\pgfqpoint{10.843830in}{3.635949in}}%
\pgfpathclose%
\pgfpathmoveto{\pgfqpoint{10.961612in}{3.702623in}}%
\pgfpathquadraticcurveto{\pgfqpoint{10.951633in}{3.685483in}}{\pgfqpoint{10.946764in}{3.668677in}}%
\pgfpathquadraticcurveto{\pgfqpoint{10.941918in}{3.651896in}}{\pgfqpoint{10.941918in}{3.634660in}}%
\pgfpathquadraticcurveto{\pgfqpoint{10.941918in}{3.617449in}}{\pgfqpoint{10.946811in}{3.600548in}}%
\pgfpathquadraticcurveto{\pgfqpoint{10.951705in}{3.583647in}}{\pgfqpoint{10.961612in}{3.566555in}}%
\pgfpathlineto{\pgfqpoint{10.949676in}{3.566555in}}%
\pgfpathquadraticcurveto{\pgfqpoint{10.938504in}{3.584100in}}{\pgfqpoint{10.932942in}{3.601025in}}%
\pgfpathquadraticcurveto{\pgfqpoint{10.927380in}{3.617950in}}{\pgfqpoint{10.927380in}{3.634660in}}%
\pgfpathquadraticcurveto{\pgfqpoint{10.927380in}{3.651299in}}{\pgfqpoint{10.932894in}{3.668152in}}%
\pgfpathquadraticcurveto{\pgfqpoint{10.938432in}{3.685029in}}{\pgfqpoint{10.949676in}{3.702623in}}%
\pgfpathlineto{\pgfqpoint{10.961612in}{3.702623in}}%
\pgfpathlineto{\pgfqpoint{10.961612in}{3.702623in}}%
\pgfpathclose%
\pgfpathmoveto{\pgfqpoint{10.973404in}{3.698087in}}%
\pgfpathlineto{\pgfqpoint{11.067625in}{3.698087in}}%
\pgfpathlineto{\pgfqpoint{11.067625in}{3.685387in}}%
\pgfpathlineto{\pgfqpoint{11.028094in}{3.685387in}}%
\pgfpathlineto{\pgfqpoint{11.028094in}{3.586702in}}%
\pgfpathlineto{\pgfqpoint{11.012959in}{3.586702in}}%
\pgfpathlineto{\pgfqpoint{11.012959in}{3.685387in}}%
\pgfpathlineto{\pgfqpoint{10.973404in}{3.685387in}}%
\pgfpathlineto{\pgfqpoint{10.973404in}{3.698087in}}%
\pgfpathlineto{\pgfqpoint{10.973404in}{3.698087in}}%
\pgfpathclose%
\pgfpathmoveto{\pgfqpoint{11.072256in}{3.698087in}}%
\pgfpathlineto{\pgfqpoint{11.087462in}{3.698087in}}%
\pgfpathlineto{\pgfqpoint{11.110880in}{3.603938in}}%
\pgfpathlineto{\pgfqpoint{11.134227in}{3.698087in}}%
\pgfpathlineto{\pgfqpoint{11.151176in}{3.698087in}}%
\pgfpathlineto{\pgfqpoint{11.174594in}{3.603938in}}%
\pgfpathlineto{\pgfqpoint{11.197940in}{3.698087in}}%
\pgfpathlineto{\pgfqpoint{11.213242in}{3.698087in}}%
\pgfpathlineto{\pgfqpoint{11.185264in}{3.586702in}}%
\pgfpathlineto{\pgfqpoint{11.166310in}{3.586702in}}%
\pgfpathlineto{\pgfqpoint{11.142821in}{3.683382in}}%
\pgfpathlineto{\pgfqpoint{11.119092in}{3.586702in}}%
\pgfpathlineto{\pgfqpoint{11.100138in}{3.586702in}}%
\pgfpathlineto{\pgfqpoint{11.072256in}{3.698087in}}%
\pgfpathlineto{\pgfqpoint{11.072256in}{3.698087in}}%
\pgfpathclose%
\pgfpathmoveto{\pgfqpoint{11.233222in}{3.698087in}}%
\pgfpathlineto{\pgfqpoint{11.297222in}{3.698087in}}%
\pgfpathlineto{\pgfqpoint{11.297222in}{3.685387in}}%
\pgfpathlineto{\pgfqpoint{11.248285in}{3.685387in}}%
\pgfpathlineto{\pgfqpoint{11.248285in}{3.652564in}}%
\pgfpathlineto{\pgfqpoint{11.292447in}{3.652564in}}%
\pgfpathlineto{\pgfqpoint{11.292447in}{3.639888in}}%
\pgfpathlineto{\pgfqpoint{11.248285in}{3.639888in}}%
\pgfpathlineto{\pgfqpoint{11.248285in}{3.586702in}}%
\pgfpathlineto{\pgfqpoint{11.233222in}{3.586702in}}%
\pgfpathlineto{\pgfqpoint{11.233222in}{3.698087in}}%
\pgfpathlineto{\pgfqpoint{11.233222in}{3.698087in}}%
\pgfpathclose%
\pgfpathmoveto{\pgfqpoint{11.321093in}{3.698087in}}%
\pgfpathlineto{\pgfqpoint{11.391514in}{3.698087in}}%
\pgfpathlineto{\pgfqpoint{11.391514in}{3.685387in}}%
\pgfpathlineto{\pgfqpoint{11.336156in}{3.685387in}}%
\pgfpathlineto{\pgfqpoint{11.336156in}{3.652421in}}%
\pgfpathlineto{\pgfqpoint{11.389199in}{3.652421in}}%
\pgfpathlineto{\pgfqpoint{11.389199in}{3.639745in}}%
\pgfpathlineto{\pgfqpoint{11.336156in}{3.639745in}}%
\pgfpathlineto{\pgfqpoint{11.336156in}{3.599378in}}%
\pgfpathlineto{\pgfqpoint{11.392851in}{3.599378in}}%
\pgfpathlineto{\pgfqpoint{11.392851in}{3.586702in}}%
\pgfpathlineto{\pgfqpoint{11.321093in}{3.586702in}}%
\pgfpathlineto{\pgfqpoint{11.321093in}{3.698087in}}%
\pgfpathlineto{\pgfqpoint{11.321093in}{3.698087in}}%
\pgfpathclose%
\pgfpathmoveto{\pgfqpoint{11.532953in}{3.694435in}}%
\pgfpathlineto{\pgfqpoint{11.532953in}{3.679730in}}%
\pgfpathquadraticcurveto{\pgfqpoint{11.524383in}{3.683836in}}{\pgfqpoint{11.516768in}{3.685841in}}%
\pgfpathquadraticcurveto{\pgfqpoint{11.509153in}{3.687870in}}{\pgfqpoint{11.502063in}{3.687870in}}%
\pgfpathquadraticcurveto{\pgfqpoint{11.489769in}{3.687870in}}{\pgfqpoint{11.483085in}{3.683096in}}%
\pgfpathquadraticcurveto{\pgfqpoint{11.476401in}{3.678321in}}{\pgfqpoint{11.476401in}{3.669513in}}%
\pgfpathquadraticcurveto{\pgfqpoint{11.476401in}{3.662113in}}{\pgfqpoint{11.480841in}{3.658341in}}%
\pgfpathquadraticcurveto{\pgfqpoint{11.485281in}{3.654593in}}{\pgfqpoint{11.497671in}{3.652278in}}%
\pgfpathlineto{\pgfqpoint{11.506766in}{3.650416in}}%
\pgfpathquadraticcurveto{\pgfqpoint{11.523619in}{3.647193in}}{\pgfqpoint{11.531640in}{3.639100in}}%
\pgfpathquadraticcurveto{\pgfqpoint{11.539661in}{3.631008in}}{\pgfqpoint{11.539661in}{3.617449in}}%
\pgfpathquadraticcurveto{\pgfqpoint{11.539661in}{3.601240in}}{\pgfqpoint{11.528799in}{3.592885in}}%
\pgfpathquadraticcurveto{\pgfqpoint{11.517962in}{3.584530in}}{\pgfqpoint{11.497002in}{3.584530in}}%
\pgfpathquadraticcurveto{\pgfqpoint{11.489101in}{3.584530in}}{\pgfqpoint{11.480173in}{3.586321in}}%
\pgfpathquadraticcurveto{\pgfqpoint{11.471269in}{3.588111in}}{\pgfqpoint{11.461720in}{3.591620in}}%
\pgfpathlineto{\pgfqpoint{11.461720in}{3.607137in}}%
\pgfpathquadraticcurveto{\pgfqpoint{11.470887in}{3.602004in}}{\pgfqpoint{11.479696in}{3.599378in}}%
\pgfpathquadraticcurveto{\pgfqpoint{11.488504in}{3.596776in}}{\pgfqpoint{11.497002in}{3.596776in}}%
\pgfpathquadraticcurveto{\pgfqpoint{11.509893in}{3.596776in}}{\pgfqpoint{11.516911in}{3.601837in}}%
\pgfpathquadraticcurveto{\pgfqpoint{11.523929in}{3.606922in}}{\pgfqpoint{11.523929in}{3.616327in}}%
\pgfpathquadraticcurveto{\pgfqpoint{11.523929in}{3.624515in}}{\pgfqpoint{11.518893in}{3.629146in}}%
\pgfpathquadraticcurveto{\pgfqpoint{11.513856in}{3.633777in}}{\pgfqpoint{11.502373in}{3.636093in}}%
\pgfpathlineto{\pgfqpoint{11.493183in}{3.637883in}}%
\pgfpathquadraticcurveto{\pgfqpoint{11.476330in}{3.641225in}}{\pgfqpoint{11.468786in}{3.648387in}}%
\pgfpathquadraticcurveto{\pgfqpoint{11.461267in}{3.655548in}}{\pgfqpoint{11.461267in}{3.668319in}}%
\pgfpathquadraticcurveto{\pgfqpoint{11.461267in}{3.683096in}}{\pgfqpoint{11.471675in}{3.691594in}}%
\pgfpathquadraticcurveto{\pgfqpoint{11.482083in}{3.700092in}}{\pgfqpoint{11.500344in}{3.700092in}}%
\pgfpathquadraticcurveto{\pgfqpoint{11.508198in}{3.700092in}}{\pgfqpoint{11.516314in}{3.698660in}}%
\pgfpathquadraticcurveto{\pgfqpoint{11.524455in}{3.697252in}}{\pgfqpoint{11.532953in}{3.694435in}}%
\pgfpathlineto{\pgfqpoint{11.532953in}{3.694435in}}%
\pgfpathclose%
\pgfpathmoveto{\pgfqpoint{11.563174in}{3.698087in}}%
\pgfpathlineto{\pgfqpoint{11.633595in}{3.698087in}}%
\pgfpathlineto{\pgfqpoint{11.633595in}{3.685387in}}%
\pgfpathlineto{\pgfqpoint{11.578237in}{3.685387in}}%
\pgfpathlineto{\pgfqpoint{11.578237in}{3.652421in}}%
\pgfpathlineto{\pgfqpoint{11.631280in}{3.652421in}}%
\pgfpathlineto{\pgfqpoint{11.631280in}{3.639745in}}%
\pgfpathlineto{\pgfqpoint{11.578237in}{3.639745in}}%
\pgfpathlineto{\pgfqpoint{11.578237in}{3.599378in}}%
\pgfpathlineto{\pgfqpoint{11.634932in}{3.599378in}}%
\pgfpathlineto{\pgfqpoint{11.634932in}{3.586702in}}%
\pgfpathlineto{\pgfqpoint{11.563174in}{3.586702in}}%
\pgfpathlineto{\pgfqpoint{11.563174in}{3.698087in}}%
\pgfpathlineto{\pgfqpoint{11.563174in}{3.698087in}}%
\pgfpathclose%
\pgfpathmoveto{\pgfqpoint{11.712371in}{3.667794in}}%
\pgfpathlineto{\pgfqpoint{11.712371in}{3.654808in}}%
\pgfpathquadraticcurveto{\pgfqpoint{11.706571in}{3.657792in}}{\pgfqpoint{11.700292in}{3.659272in}}%
\pgfpathquadraticcurveto{\pgfqpoint{11.694038in}{3.660776in}}{\pgfqpoint{11.687306in}{3.660776in}}%
\pgfpathquadraticcurveto{\pgfqpoint{11.677089in}{3.660776in}}{\pgfqpoint{11.671981in}{3.657649in}}%
\pgfpathquadraticcurveto{\pgfqpoint{11.666872in}{3.654521in}}{\pgfqpoint{11.666872in}{3.648243in}}%
\pgfpathquadraticcurveto{\pgfqpoint{11.666872in}{3.643469in}}{\pgfqpoint{11.670525in}{3.640748in}}%
\pgfpathquadraticcurveto{\pgfqpoint{11.674177in}{3.638026in}}{\pgfqpoint{11.685229in}{3.635567in}}%
\pgfpathlineto{\pgfqpoint{11.689932in}{3.634517in}}%
\pgfpathquadraticcurveto{\pgfqpoint{11.704541in}{3.631390in}}{\pgfqpoint{11.710700in}{3.625685in}}%
\pgfpathquadraticcurveto{\pgfqpoint{11.716859in}{3.619979in}}{\pgfqpoint{11.716859in}{3.609762in}}%
\pgfpathquadraticcurveto{\pgfqpoint{11.716859in}{3.598113in}}{\pgfqpoint{11.707645in}{3.591310in}}%
\pgfpathquadraticcurveto{\pgfqpoint{11.698430in}{3.584530in}}{\pgfqpoint{11.682317in}{3.584530in}}%
\pgfpathquadraticcurveto{\pgfqpoint{11.675609in}{3.584530in}}{\pgfqpoint{11.668328in}{3.585843in}}%
\pgfpathquadraticcurveto{\pgfqpoint{11.661048in}{3.587156in}}{\pgfqpoint{11.653003in}{3.589758in}}%
\pgfpathlineto{\pgfqpoint{11.653003in}{3.603938in}}%
\pgfpathquadraticcurveto{\pgfqpoint{11.660618in}{3.599975in}}{\pgfqpoint{11.667994in}{3.597994in}}%
\pgfpathquadraticcurveto{\pgfqpoint{11.675370in}{3.596036in}}{\pgfqpoint{11.682627in}{3.596036in}}%
\pgfpathquadraticcurveto{\pgfqpoint{11.692319in}{3.596036in}}{\pgfqpoint{11.697523in}{3.599354in}}%
\pgfpathquadraticcurveto{\pgfqpoint{11.702751in}{3.602673in}}{\pgfqpoint{11.702751in}{3.608712in}}%
\pgfpathquadraticcurveto{\pgfqpoint{11.702751in}{3.614298in}}{\pgfqpoint{11.698979in}{3.617282in}}%
\pgfpathquadraticcurveto{\pgfqpoint{11.695232in}{3.620266in}}{\pgfqpoint{11.682460in}{3.623035in}}%
\pgfpathlineto{\pgfqpoint{11.677686in}{3.624157in}}%
\pgfpathquadraticcurveto{\pgfqpoint{11.664939in}{3.626831in}}{\pgfqpoint{11.659257in}{3.632393in}}%
\pgfpathquadraticcurveto{\pgfqpoint{11.653600in}{3.637955in}}{\pgfqpoint{11.653600in}{3.647646in}}%
\pgfpathquadraticcurveto{\pgfqpoint{11.653600in}{3.659439in}}{\pgfqpoint{11.661955in}{3.665837in}}%
\pgfpathquadraticcurveto{\pgfqpoint{11.670310in}{3.672258in}}{\pgfqpoint{11.685683in}{3.672258in}}%
\pgfpathquadraticcurveto{\pgfqpoint{11.693274in}{3.672258in}}{\pgfqpoint{11.699982in}{3.671136in}}%
\pgfpathquadraticcurveto{\pgfqpoint{11.706714in}{3.670038in}}{\pgfqpoint{11.712371in}{3.667794in}}%
\pgfpathlineto{\pgfqpoint{11.712371in}{3.667794in}}%
\pgfpathclose%
\pgfpathmoveto{\pgfqpoint{11.736553in}{3.702623in}}%
\pgfpathlineto{\pgfqpoint{11.748489in}{3.702623in}}%
\pgfpathquadraticcurveto{\pgfqpoint{11.759661in}{3.685029in}}{\pgfqpoint{11.765223in}{3.668152in}}%
\pgfpathquadraticcurveto{\pgfqpoint{11.770785in}{3.651299in}}{\pgfqpoint{11.770785in}{3.634660in}}%
\pgfpathquadraticcurveto{\pgfqpoint{11.770785in}{3.617950in}}{\pgfqpoint{11.765223in}{3.601025in}}%
\pgfpathquadraticcurveto{\pgfqpoint{11.759661in}{3.584100in}}{\pgfqpoint{11.748489in}{3.566555in}}%
\pgfpathlineto{\pgfqpoint{11.736553in}{3.566555in}}%
\pgfpathquadraticcurveto{\pgfqpoint{11.746460in}{3.583647in}}{\pgfqpoint{11.751353in}{3.600548in}}%
\pgfpathquadraticcurveto{\pgfqpoint{11.756247in}{3.617449in}}{\pgfqpoint{11.756247in}{3.634660in}}%
\pgfpathquadraticcurveto{\pgfqpoint{11.756247in}{3.651896in}}{\pgfqpoint{11.751353in}{3.668677in}}%
\pgfpathquadraticcurveto{\pgfqpoint{11.746460in}{3.685483in}}{\pgfqpoint{11.736553in}{3.702623in}}%
\pgfpathlineto{\pgfqpoint{11.736553in}{3.702623in}}%
\pgfpathclose%
\pgfusepath{fill}%
\end{pgfscope}%
\begin{pgfscope}%
\pgfpathrectangle{\pgfqpoint{8.036091in}{4.050000in}}{\pgfqpoint{5.400000in}{5.400000in}}%
\pgfusepath{clip}%
\pgfsetrectcap%
\pgfsetroundjoin%
\pgfsetlinewidth{0.602250pt}%
\definecolor{currentstroke}{rgb}{0.898039,0.913725,0.929412}%
\pgfsetstrokecolor{currentstroke}%
\pgfsetdash{}{0pt}%
\pgfpathmoveto{\pgfqpoint{8.036091in}{4.050000in}}%
\pgfpathlineto{\pgfqpoint{13.436091in}{4.050000in}}%
\pgfusepath{stroke}%
\end{pgfscope}%
\begin{pgfscope}%
\pgfsetbuttcap%
\pgfsetroundjoin%
\definecolor{currentfill}{rgb}{0.090196,0.168627,0.227451}%
\pgfsetfillcolor{currentfill}%
\pgfsetlinewidth{0.803000pt}%
\definecolor{currentstroke}{rgb}{0.090196,0.168627,0.227451}%
\pgfsetstrokecolor{currentstroke}%
\pgfsetdash{}{0pt}%
\pgfsys@defobject{currentmarker}{\pgfqpoint{-0.048611in}{0.000000in}}{\pgfqpoint{-0.000000in}{0.000000in}}{%
\pgfpathmoveto{\pgfqpoint{-0.000000in}{0.000000in}}%
\pgfpathlineto{\pgfqpoint{-0.048611in}{0.000000in}}%
\pgfusepath{stroke,fill}%
}%
\begin{pgfscope}%
\pgfsys@transformshift{8.036091in}{4.050000in}%
\pgfsys@useobject{currentmarker}{}%
\end{pgfscope}%
\end{pgfscope}%
\begin{pgfscope}%
\pgfsetbuttcap%
\pgfsetroundjoin%
\definecolor{currentfill}{rgb}{0.090196,0.168627,0.227451}%
\pgfsetfillcolor{currentfill}%
\pgfsetlinewidth{0.000000pt}%
\definecolor{currentstroke}{rgb}{0.090196,0.168627,0.227451}%
\pgfsetstrokecolor{currentstroke}%
\pgfsetdash{}{0pt}%
\pgfpathmoveto{\pgfqpoint{7.642111in}{4.046890in}}%
\pgfpathlineto{\pgfqpoint{7.720353in}{4.046890in}}%
\pgfpathlineto{\pgfqpoint{7.720353in}{4.036519in}}%
\pgfpathlineto{\pgfqpoint{7.642111in}{4.036519in}}%
\pgfpathlineto{\pgfqpoint{7.642111in}{4.046890in}}%
\pgfpathlineto{\pgfqpoint{7.642111in}{4.046890in}}%
\pgfpathclose%
\pgfpathmoveto{\pgfqpoint{7.747111in}{4.093647in}}%
\pgfpathlineto{\pgfqpoint{7.795509in}{4.093647in}}%
\pgfpathlineto{\pgfqpoint{7.795509in}{4.083257in}}%
\pgfpathlineto{\pgfqpoint{7.758400in}{4.083257in}}%
\pgfpathlineto{\pgfqpoint{7.758400in}{4.060933in}}%
\pgfpathquadraticcurveto{\pgfqpoint{7.761076in}{4.061851in}}{\pgfqpoint{7.763751in}{4.062300in}}%
\pgfpathquadraticcurveto{\pgfqpoint{7.766447in}{4.062749in}}{\pgfqpoint{7.769142in}{4.062749in}}%
\pgfpathquadraticcurveto{\pgfqpoint{7.784396in}{4.062749in}}{\pgfqpoint{7.793302in}{4.054390in}}%
\pgfpathquadraticcurveto{\pgfqpoint{7.802228in}{4.046030in}}{\pgfqpoint{7.802228in}{4.031753in}}%
\pgfpathquadraticcurveto{\pgfqpoint{7.802228in}{4.017046in}}{\pgfqpoint{7.793068in}{4.008882in}}%
\pgfpathquadraticcurveto{\pgfqpoint{7.783908in}{4.000737in}}{\pgfqpoint{7.767248in}{4.000737in}}%
\pgfpathquadraticcurveto{\pgfqpoint{7.761505in}{4.000737in}}{\pgfqpoint{7.755548in}{4.001714in}}%
\pgfpathquadraticcurveto{\pgfqpoint{7.749611in}{4.002690in}}{\pgfqpoint{7.743263in}{4.004644in}}%
\pgfpathlineto{\pgfqpoint{7.743263in}{4.017046in}}%
\pgfpathquadraticcurveto{\pgfqpoint{7.748751in}{4.014058in}}{\pgfqpoint{7.754611in}{4.012593in}}%
\pgfpathquadraticcurveto{\pgfqpoint{7.760470in}{4.011128in}}{\pgfqpoint{7.766994in}{4.011128in}}%
\pgfpathquadraticcurveto{\pgfqpoint{7.777560in}{4.011128in}}{\pgfqpoint{7.783712in}{4.016675in}}%
\pgfpathquadraticcurveto{\pgfqpoint{7.789884in}{4.022222in}}{\pgfqpoint{7.789884in}{4.031753in}}%
\pgfpathquadraticcurveto{\pgfqpoint{7.789884in}{4.041265in}}{\pgfqpoint{7.783712in}{4.046812in}}%
\pgfpathquadraticcurveto{\pgfqpoint{7.777560in}{4.052378in}}{\pgfqpoint{7.766994in}{4.052378in}}%
\pgfpathquadraticcurveto{\pgfqpoint{7.762052in}{4.052378in}}{\pgfqpoint{7.757130in}{4.051284in}}%
\pgfpathquadraticcurveto{\pgfqpoint{7.752228in}{4.050190in}}{\pgfqpoint{7.747111in}{4.047866in}}%
\pgfpathlineto{\pgfqpoint{7.747111in}{4.093647in}}%
\pgfpathlineto{\pgfqpoint{7.747111in}{4.093647in}}%
\pgfpathclose%
\pgfpathmoveto{\pgfqpoint{7.826505in}{4.018022in}}%
\pgfpathlineto{\pgfqpoint{7.839396in}{4.018022in}}%
\pgfpathlineto{\pgfqpoint{7.839396in}{4.002515in}}%
\pgfpathlineto{\pgfqpoint{7.826505in}{4.002515in}}%
\pgfpathlineto{\pgfqpoint{7.826505in}{4.018022in}}%
\pgfpathlineto{\pgfqpoint{7.826505in}{4.018022in}}%
\pgfpathclose%
\pgfpathmoveto{\pgfqpoint{7.892599in}{4.085522in}}%
\pgfpathquadraticcurveto{\pgfqpoint{7.883087in}{4.085522in}}{\pgfqpoint{7.878283in}{4.076147in}}%
\pgfpathquadraticcurveto{\pgfqpoint{7.873498in}{4.066792in}}{\pgfqpoint{7.873498in}{4.047983in}}%
\pgfpathquadraticcurveto{\pgfqpoint{7.873498in}{4.029253in}}{\pgfqpoint{7.878283in}{4.019878in}}%
\pgfpathquadraticcurveto{\pgfqpoint{7.883087in}{4.010503in}}{\pgfqpoint{7.892599in}{4.010503in}}%
\pgfpathquadraticcurveto{\pgfqpoint{7.902189in}{4.010503in}}{\pgfqpoint{7.906974in}{4.019878in}}%
\pgfpathquadraticcurveto{\pgfqpoint{7.911779in}{4.029253in}}{\pgfqpoint{7.911779in}{4.047983in}}%
\pgfpathquadraticcurveto{\pgfqpoint{7.911779in}{4.066792in}}{\pgfqpoint{7.906974in}{4.076147in}}%
\pgfpathquadraticcurveto{\pgfqpoint{7.902189in}{4.085522in}}{\pgfqpoint{7.892599in}{4.085522in}}%
\pgfpathlineto{\pgfqpoint{7.892599in}{4.085522in}}%
\pgfpathclose%
\pgfpathmoveto{\pgfqpoint{7.892599in}{4.095288in}}%
\pgfpathquadraticcurveto{\pgfqpoint{7.907931in}{4.095288in}}{\pgfqpoint{7.916017in}{4.083159in}}%
\pgfpathquadraticcurveto{\pgfqpoint{7.924103in}{4.071050in}}{\pgfqpoint{7.924103in}{4.047983in}}%
\pgfpathquadraticcurveto{\pgfqpoint{7.924103in}{4.024976in}}{\pgfqpoint{7.916017in}{4.012847in}}%
\pgfpathquadraticcurveto{\pgfqpoint{7.907931in}{4.000737in}}{\pgfqpoint{7.892599in}{4.000737in}}%
\pgfpathquadraticcurveto{\pgfqpoint{7.877287in}{4.000737in}}{\pgfqpoint{7.869201in}{4.012847in}}%
\pgfpathquadraticcurveto{\pgfqpoint{7.861115in}{4.024976in}}{\pgfqpoint{7.861115in}{4.047983in}}%
\pgfpathquadraticcurveto{\pgfqpoint{7.861115in}{4.071050in}}{\pgfqpoint{7.869201in}{4.083159in}}%
\pgfpathquadraticcurveto{\pgfqpoint{7.877287in}{4.095288in}}{\pgfqpoint{7.892599in}{4.095288in}}%
\pgfpathlineto{\pgfqpoint{7.892599in}{4.095288in}}%
\pgfpathclose%
\pgfusepath{fill}%
\end{pgfscope}%
\begin{pgfscope}%
\pgfpathrectangle{\pgfqpoint{8.036091in}{4.050000in}}{\pgfqpoint{5.400000in}{5.400000in}}%
\pgfusepath{clip}%
\pgfsetrectcap%
\pgfsetroundjoin%
\pgfsetlinewidth{0.602250pt}%
\definecolor{currentstroke}{rgb}{0.898039,0.913725,0.929412}%
\pgfsetstrokecolor{currentstroke}%
\pgfsetdash{}{0pt}%
\pgfpathmoveto{\pgfqpoint{8.036091in}{5.400000in}}%
\pgfpathlineto{\pgfqpoint{13.436091in}{5.400000in}}%
\pgfusepath{stroke}%
\end{pgfscope}%
\begin{pgfscope}%
\pgfsetbuttcap%
\pgfsetroundjoin%
\definecolor{currentfill}{rgb}{0.090196,0.168627,0.227451}%
\pgfsetfillcolor{currentfill}%
\pgfsetlinewidth{0.803000pt}%
\definecolor{currentstroke}{rgb}{0.090196,0.168627,0.227451}%
\pgfsetstrokecolor{currentstroke}%
\pgfsetdash{}{0pt}%
\pgfsys@defobject{currentmarker}{\pgfqpoint{-0.048611in}{0.000000in}}{\pgfqpoint{-0.000000in}{0.000000in}}{%
\pgfpathmoveto{\pgfqpoint{-0.000000in}{0.000000in}}%
\pgfpathlineto{\pgfqpoint{-0.048611in}{0.000000in}}%
\pgfusepath{stroke,fill}%
}%
\begin{pgfscope}%
\pgfsys@transformshift{8.036091in}{5.400000in}%
\pgfsys@useobject{currentmarker}{}%
\end{pgfscope}%
\end{pgfscope}%
\begin{pgfscope}%
\pgfsetbuttcap%
\pgfsetroundjoin%
\definecolor{currentfill}{rgb}{0.090196,0.168627,0.227451}%
\pgfsetfillcolor{currentfill}%
\pgfsetlinewidth{0.000000pt}%
\definecolor{currentstroke}{rgb}{0.090196,0.168627,0.227451}%
\pgfsetstrokecolor{currentstroke}%
\pgfsetdash{}{0pt}%
\pgfpathmoveto{\pgfqpoint{7.642111in}{5.396890in}}%
\pgfpathlineto{\pgfqpoint{7.720353in}{5.396890in}}%
\pgfpathlineto{\pgfqpoint{7.720353in}{5.386519in}}%
\pgfpathlineto{\pgfqpoint{7.642111in}{5.386519in}}%
\pgfpathlineto{\pgfqpoint{7.642111in}{5.396890in}}%
\pgfpathlineto{\pgfqpoint{7.642111in}{5.396890in}}%
\pgfpathclose%
\pgfpathmoveto{\pgfqpoint{7.757599in}{5.362886in}}%
\pgfpathlineto{\pgfqpoint{7.800626in}{5.362886in}}%
\pgfpathlineto{\pgfqpoint{7.800626in}{5.352515in}}%
\pgfpathlineto{\pgfqpoint{7.742775in}{5.352515in}}%
\pgfpathlineto{\pgfqpoint{7.742775in}{5.362886in}}%
\pgfpathquadraticcurveto{\pgfqpoint{7.749787in}{5.370151in}}{\pgfqpoint{7.761896in}{5.382378in}}%
\pgfpathquadraticcurveto{\pgfqpoint{7.774025in}{5.394624in}}{\pgfqpoint{7.777130in}{5.398179in}}%
\pgfpathquadraticcurveto{\pgfqpoint{7.783048in}{5.404819in}}{\pgfqpoint{7.785392in}{5.409429in}}%
\pgfpathquadraticcurveto{\pgfqpoint{7.787755in}{5.414038in}}{\pgfqpoint{7.787755in}{5.418491in}}%
\pgfpathquadraticcurveto{\pgfqpoint{7.787755in}{5.425757in}}{\pgfqpoint{7.782658in}{5.430327in}}%
\pgfpathquadraticcurveto{\pgfqpoint{7.777560in}{5.434917in}}{\pgfqpoint{7.769376in}{5.434917in}}%
\pgfpathquadraticcurveto{\pgfqpoint{7.763576in}{5.434917in}}{\pgfqpoint{7.757130in}{5.432905in}}%
\pgfpathquadraticcurveto{\pgfqpoint{7.750705in}{5.430894in}}{\pgfqpoint{7.743380in}{5.426792in}}%
\pgfpathlineto{\pgfqpoint{7.743380in}{5.439253in}}%
\pgfpathquadraticcurveto{\pgfqpoint{7.750822in}{5.442241in}}{\pgfqpoint{7.757287in}{5.443765in}}%
\pgfpathquadraticcurveto{\pgfqpoint{7.763771in}{5.445288in}}{\pgfqpoint{7.769142in}{5.445288in}}%
\pgfpathquadraticcurveto{\pgfqpoint{7.783302in}{5.445288in}}{\pgfqpoint{7.791720in}{5.438198in}}%
\pgfpathquadraticcurveto{\pgfqpoint{7.800138in}{5.431128in}}{\pgfqpoint{7.800138in}{5.419292in}}%
\pgfpathquadraticcurveto{\pgfqpoint{7.800138in}{5.413667in}}{\pgfqpoint{7.798029in}{5.408628in}}%
\pgfpathquadraticcurveto{\pgfqpoint{7.795939in}{5.403608in}}{\pgfqpoint{7.790373in}{5.396772in}}%
\pgfpathquadraticcurveto{\pgfqpoint{7.788849in}{5.394995in}}{\pgfqpoint{7.780666in}{5.386538in}}%
\pgfpathquadraticcurveto{\pgfqpoint{7.772501in}{5.378081in}}{\pgfqpoint{7.757599in}{5.362886in}}%
\pgfpathlineto{\pgfqpoint{7.757599in}{5.362886in}}%
\pgfpathclose%
\pgfpathmoveto{\pgfqpoint{7.826505in}{5.368022in}}%
\pgfpathlineto{\pgfqpoint{7.839396in}{5.368022in}}%
\pgfpathlineto{\pgfqpoint{7.839396in}{5.352515in}}%
\pgfpathlineto{\pgfqpoint{7.826505in}{5.352515in}}%
\pgfpathlineto{\pgfqpoint{7.826505in}{5.368022in}}%
\pgfpathlineto{\pgfqpoint{7.826505in}{5.368022in}}%
\pgfpathclose%
\pgfpathmoveto{\pgfqpoint{7.866369in}{5.443647in}}%
\pgfpathlineto{\pgfqpoint{7.914767in}{5.443647in}}%
\pgfpathlineto{\pgfqpoint{7.914767in}{5.433257in}}%
\pgfpathlineto{\pgfqpoint{7.877658in}{5.433257in}}%
\pgfpathlineto{\pgfqpoint{7.877658in}{5.410933in}}%
\pgfpathquadraticcurveto{\pgfqpoint{7.880334in}{5.411851in}}{\pgfqpoint{7.883009in}{5.412300in}}%
\pgfpathquadraticcurveto{\pgfqpoint{7.885705in}{5.412749in}}{\pgfqpoint{7.888400in}{5.412749in}}%
\pgfpathquadraticcurveto{\pgfqpoint{7.903654in}{5.412749in}}{\pgfqpoint{7.912560in}{5.404390in}}%
\pgfpathquadraticcurveto{\pgfqpoint{7.921486in}{5.396030in}}{\pgfqpoint{7.921486in}{5.381753in}}%
\pgfpathquadraticcurveto{\pgfqpoint{7.921486in}{5.367046in}}{\pgfqpoint{7.912326in}{5.358882in}}%
\pgfpathquadraticcurveto{\pgfqpoint{7.903166in}{5.350737in}}{\pgfqpoint{7.886505in}{5.350737in}}%
\pgfpathquadraticcurveto{\pgfqpoint{7.880763in}{5.350737in}}{\pgfqpoint{7.874806in}{5.351714in}}%
\pgfpathquadraticcurveto{\pgfqpoint{7.868869in}{5.352690in}}{\pgfqpoint{7.862521in}{5.354644in}}%
\pgfpathlineto{\pgfqpoint{7.862521in}{5.367046in}}%
\pgfpathquadraticcurveto{\pgfqpoint{7.868009in}{5.364058in}}{\pgfqpoint{7.873869in}{5.362593in}}%
\pgfpathquadraticcurveto{\pgfqpoint{7.879728in}{5.361128in}}{\pgfqpoint{7.886251in}{5.361128in}}%
\pgfpathquadraticcurveto{\pgfqpoint{7.896818in}{5.361128in}}{\pgfqpoint{7.902970in}{5.366675in}}%
\pgfpathquadraticcurveto{\pgfqpoint{7.909142in}{5.372222in}}{\pgfqpoint{7.909142in}{5.381753in}}%
\pgfpathquadraticcurveto{\pgfqpoint{7.909142in}{5.391265in}}{\pgfqpoint{7.902970in}{5.396812in}}%
\pgfpathquadraticcurveto{\pgfqpoint{7.896818in}{5.402378in}}{\pgfqpoint{7.886251in}{5.402378in}}%
\pgfpathquadraticcurveto{\pgfqpoint{7.881310in}{5.402378in}}{\pgfqpoint{7.876388in}{5.401284in}}%
\pgfpathquadraticcurveto{\pgfqpoint{7.871486in}{5.400190in}}{\pgfqpoint{7.866369in}{5.397866in}}%
\pgfpathlineto{\pgfqpoint{7.866369in}{5.443647in}}%
\pgfpathlineto{\pgfqpoint{7.866369in}{5.443647in}}%
\pgfpathclose%
\pgfusepath{fill}%
\end{pgfscope}%
\begin{pgfscope}%
\pgfpathrectangle{\pgfqpoint{8.036091in}{4.050000in}}{\pgfqpoint{5.400000in}{5.400000in}}%
\pgfusepath{clip}%
\pgfsetrectcap%
\pgfsetroundjoin%
\pgfsetlinewidth{0.602250pt}%
\definecolor{currentstroke}{rgb}{0.898039,0.913725,0.929412}%
\pgfsetstrokecolor{currentstroke}%
\pgfsetdash{}{0pt}%
\pgfpathmoveto{\pgfqpoint{8.036091in}{6.750000in}}%
\pgfpathlineto{\pgfqpoint{13.436091in}{6.750000in}}%
\pgfusepath{stroke}%
\end{pgfscope}%
\begin{pgfscope}%
\pgfsetbuttcap%
\pgfsetroundjoin%
\definecolor{currentfill}{rgb}{0.090196,0.168627,0.227451}%
\pgfsetfillcolor{currentfill}%
\pgfsetlinewidth{0.803000pt}%
\definecolor{currentstroke}{rgb}{0.090196,0.168627,0.227451}%
\pgfsetstrokecolor{currentstroke}%
\pgfsetdash{}{0pt}%
\pgfsys@defobject{currentmarker}{\pgfqpoint{-0.048611in}{0.000000in}}{\pgfqpoint{-0.000000in}{0.000000in}}{%
\pgfpathmoveto{\pgfqpoint{-0.000000in}{0.000000in}}%
\pgfpathlineto{\pgfqpoint{-0.048611in}{0.000000in}}%
\pgfusepath{stroke,fill}%
}%
\begin{pgfscope}%
\pgfsys@transformshift{8.036091in}{6.750000in}%
\pgfsys@useobject{currentmarker}{}%
\end{pgfscope}%
\end{pgfscope}%
\begin{pgfscope}%
\pgfsetbuttcap%
\pgfsetroundjoin%
\definecolor{currentfill}{rgb}{0.090196,0.168627,0.227451}%
\pgfsetfillcolor{currentfill}%
\pgfsetlinewidth{0.000000pt}%
\definecolor{currentstroke}{rgb}{0.090196,0.168627,0.227451}%
\pgfsetstrokecolor{currentstroke}%
\pgfsetdash{}{0pt}%
\pgfpathmoveto{\pgfqpoint{7.778595in}{6.785522in}}%
\pgfpathquadraticcurveto{\pgfqpoint{7.769084in}{6.785522in}}{\pgfqpoint{7.764279in}{6.776147in}}%
\pgfpathquadraticcurveto{\pgfqpoint{7.759494in}{6.766792in}}{\pgfqpoint{7.759494in}{6.747983in}}%
\pgfpathquadraticcurveto{\pgfqpoint{7.759494in}{6.729253in}}{\pgfqpoint{7.764279in}{6.719878in}}%
\pgfpathquadraticcurveto{\pgfqpoint{7.769084in}{6.710503in}}{\pgfqpoint{7.778595in}{6.710503in}}%
\pgfpathquadraticcurveto{\pgfqpoint{7.788185in}{6.710503in}}{\pgfqpoint{7.792970in}{6.719878in}}%
\pgfpathquadraticcurveto{\pgfqpoint{7.797775in}{6.729253in}}{\pgfqpoint{7.797775in}{6.747983in}}%
\pgfpathquadraticcurveto{\pgfqpoint{7.797775in}{6.766792in}}{\pgfqpoint{7.792970in}{6.776147in}}%
\pgfpathquadraticcurveto{\pgfqpoint{7.788185in}{6.785522in}}{\pgfqpoint{7.778595in}{6.785522in}}%
\pgfpathlineto{\pgfqpoint{7.778595in}{6.785522in}}%
\pgfpathclose%
\pgfpathmoveto{\pgfqpoint{7.778595in}{6.795288in}}%
\pgfpathquadraticcurveto{\pgfqpoint{7.793927in}{6.795288in}}{\pgfqpoint{7.802013in}{6.783159in}}%
\pgfpathquadraticcurveto{\pgfqpoint{7.810099in}{6.771050in}}{\pgfqpoint{7.810099in}{6.747983in}}%
\pgfpathquadraticcurveto{\pgfqpoint{7.810099in}{6.724976in}}{\pgfqpoint{7.802013in}{6.712847in}}%
\pgfpathquadraticcurveto{\pgfqpoint{7.793927in}{6.700737in}}{\pgfqpoint{7.778595in}{6.700737in}}%
\pgfpathquadraticcurveto{\pgfqpoint{7.763283in}{6.700737in}}{\pgfqpoint{7.755197in}{6.712847in}}%
\pgfpathquadraticcurveto{\pgfqpoint{7.747111in}{6.724976in}}{\pgfqpoint{7.747111in}{6.747983in}}%
\pgfpathquadraticcurveto{\pgfqpoint{7.747111in}{6.771050in}}{\pgfqpoint{7.755197in}{6.783159in}}%
\pgfpathquadraticcurveto{\pgfqpoint{7.763283in}{6.795288in}}{\pgfqpoint{7.778595in}{6.795288in}}%
\pgfpathlineto{\pgfqpoint{7.778595in}{6.795288in}}%
\pgfpathclose%
\pgfpathmoveto{\pgfqpoint{7.831759in}{6.718022in}}%
\pgfpathlineto{\pgfqpoint{7.844650in}{6.718022in}}%
\pgfpathlineto{\pgfqpoint{7.844650in}{6.702515in}}%
\pgfpathlineto{\pgfqpoint{7.831759in}{6.702515in}}%
\pgfpathlineto{\pgfqpoint{7.831759in}{6.718022in}}%
\pgfpathlineto{\pgfqpoint{7.831759in}{6.718022in}}%
\pgfpathclose%
\pgfpathmoveto{\pgfqpoint{7.897853in}{6.785522in}}%
\pgfpathquadraticcurveto{\pgfqpoint{7.888341in}{6.785522in}}{\pgfqpoint{7.883537in}{6.776147in}}%
\pgfpathquadraticcurveto{\pgfqpoint{7.878751in}{6.766792in}}{\pgfqpoint{7.878751in}{6.747983in}}%
\pgfpathquadraticcurveto{\pgfqpoint{7.878751in}{6.729253in}}{\pgfqpoint{7.883537in}{6.719878in}}%
\pgfpathquadraticcurveto{\pgfqpoint{7.888341in}{6.710503in}}{\pgfqpoint{7.897853in}{6.710503in}}%
\pgfpathquadraticcurveto{\pgfqpoint{7.907443in}{6.710503in}}{\pgfqpoint{7.912228in}{6.719878in}}%
\pgfpathquadraticcurveto{\pgfqpoint{7.917033in}{6.729253in}}{\pgfqpoint{7.917033in}{6.747983in}}%
\pgfpathquadraticcurveto{\pgfqpoint{7.917033in}{6.766792in}}{\pgfqpoint{7.912228in}{6.776147in}}%
\pgfpathquadraticcurveto{\pgfqpoint{7.907443in}{6.785522in}}{\pgfqpoint{7.897853in}{6.785522in}}%
\pgfpathlineto{\pgfqpoint{7.897853in}{6.785522in}}%
\pgfpathclose%
\pgfpathmoveto{\pgfqpoint{7.897853in}{6.795288in}}%
\pgfpathquadraticcurveto{\pgfqpoint{7.913185in}{6.795288in}}{\pgfqpoint{7.921271in}{6.783159in}}%
\pgfpathquadraticcurveto{\pgfqpoint{7.929357in}{6.771050in}}{\pgfqpoint{7.929357in}{6.747983in}}%
\pgfpathquadraticcurveto{\pgfqpoint{7.929357in}{6.724976in}}{\pgfqpoint{7.921271in}{6.712847in}}%
\pgfpathquadraticcurveto{\pgfqpoint{7.913185in}{6.700737in}}{\pgfqpoint{7.897853in}{6.700737in}}%
\pgfpathquadraticcurveto{\pgfqpoint{7.882541in}{6.700737in}}{\pgfqpoint{7.874455in}{6.712847in}}%
\pgfpathquadraticcurveto{\pgfqpoint{7.866369in}{6.724976in}}{\pgfqpoint{7.866369in}{6.747983in}}%
\pgfpathquadraticcurveto{\pgfqpoint{7.866369in}{6.771050in}}{\pgfqpoint{7.874455in}{6.783159in}}%
\pgfpathquadraticcurveto{\pgfqpoint{7.882541in}{6.795288in}}{\pgfqpoint{7.897853in}{6.795288in}}%
\pgfpathlineto{\pgfqpoint{7.897853in}{6.795288in}}%
\pgfpathclose%
\pgfusepath{fill}%
\end{pgfscope}%
\begin{pgfscope}%
\pgfpathrectangle{\pgfqpoint{8.036091in}{4.050000in}}{\pgfqpoint{5.400000in}{5.400000in}}%
\pgfusepath{clip}%
\pgfsetrectcap%
\pgfsetroundjoin%
\pgfsetlinewidth{0.602250pt}%
\definecolor{currentstroke}{rgb}{0.898039,0.913725,0.929412}%
\pgfsetstrokecolor{currentstroke}%
\pgfsetdash{}{0pt}%
\pgfpathmoveto{\pgfqpoint{8.036091in}{8.100000in}}%
\pgfpathlineto{\pgfqpoint{13.436091in}{8.100000in}}%
\pgfusepath{stroke}%
\end{pgfscope}%
\begin{pgfscope}%
\pgfsetbuttcap%
\pgfsetroundjoin%
\definecolor{currentfill}{rgb}{0.090196,0.168627,0.227451}%
\pgfsetfillcolor{currentfill}%
\pgfsetlinewidth{0.803000pt}%
\definecolor{currentstroke}{rgb}{0.090196,0.168627,0.227451}%
\pgfsetstrokecolor{currentstroke}%
\pgfsetdash{}{0pt}%
\pgfsys@defobject{currentmarker}{\pgfqpoint{-0.048611in}{0.000000in}}{\pgfqpoint{-0.000000in}{0.000000in}}{%
\pgfpathmoveto{\pgfqpoint{-0.000000in}{0.000000in}}%
\pgfpathlineto{\pgfqpoint{-0.048611in}{0.000000in}}%
\pgfusepath{stroke,fill}%
}%
\begin{pgfscope}%
\pgfsys@transformshift{8.036091in}{8.100000in}%
\pgfsys@useobject{currentmarker}{}%
\end{pgfscope}%
\end{pgfscope}%
\begin{pgfscope}%
\pgfsetbuttcap%
\pgfsetroundjoin%
\definecolor{currentfill}{rgb}{0.090196,0.168627,0.227451}%
\pgfsetfillcolor{currentfill}%
\pgfsetlinewidth{0.000000pt}%
\definecolor{currentstroke}{rgb}{0.090196,0.168627,0.227451}%
\pgfsetstrokecolor{currentstroke}%
\pgfsetdash{}{0pt}%
\pgfpathmoveto{\pgfqpoint{7.762853in}{8.062886in}}%
\pgfpathlineto{\pgfqpoint{7.805880in}{8.062886in}}%
\pgfpathlineto{\pgfqpoint{7.805880in}{8.052515in}}%
\pgfpathlineto{\pgfqpoint{7.748029in}{8.052515in}}%
\pgfpathlineto{\pgfqpoint{7.748029in}{8.062886in}}%
\pgfpathquadraticcurveto{\pgfqpoint{7.755041in}{8.070151in}}{\pgfqpoint{7.767150in}{8.082378in}}%
\pgfpathquadraticcurveto{\pgfqpoint{7.779279in}{8.094624in}}{\pgfqpoint{7.782384in}{8.098179in}}%
\pgfpathquadraticcurveto{\pgfqpoint{7.788302in}{8.104819in}}{\pgfqpoint{7.790646in}{8.109429in}}%
\pgfpathquadraticcurveto{\pgfqpoint{7.793009in}{8.114038in}}{\pgfqpoint{7.793009in}{8.118491in}}%
\pgfpathquadraticcurveto{\pgfqpoint{7.793009in}{8.125757in}}{\pgfqpoint{7.787912in}{8.130327in}}%
\pgfpathquadraticcurveto{\pgfqpoint{7.782814in}{8.134917in}}{\pgfqpoint{7.774630in}{8.134917in}}%
\pgfpathquadraticcurveto{\pgfqpoint{7.768830in}{8.134917in}}{\pgfqpoint{7.762384in}{8.132905in}}%
\pgfpathquadraticcurveto{\pgfqpoint{7.755959in}{8.130894in}}{\pgfqpoint{7.748634in}{8.126792in}}%
\pgfpathlineto{\pgfqpoint{7.748634in}{8.139253in}}%
\pgfpathquadraticcurveto{\pgfqpoint{7.756076in}{8.142241in}}{\pgfqpoint{7.762541in}{8.143765in}}%
\pgfpathquadraticcurveto{\pgfqpoint{7.769025in}{8.145288in}}{\pgfqpoint{7.774396in}{8.145288in}}%
\pgfpathquadraticcurveto{\pgfqpoint{7.788556in}{8.145288in}}{\pgfqpoint{7.796974in}{8.138198in}}%
\pgfpathquadraticcurveto{\pgfqpoint{7.805392in}{8.131128in}}{\pgfqpoint{7.805392in}{8.119292in}}%
\pgfpathquadraticcurveto{\pgfqpoint{7.805392in}{8.113667in}}{\pgfqpoint{7.803283in}{8.108628in}}%
\pgfpathquadraticcurveto{\pgfqpoint{7.801193in}{8.103608in}}{\pgfqpoint{7.795626in}{8.096772in}}%
\pgfpathquadraticcurveto{\pgfqpoint{7.794103in}{8.094995in}}{\pgfqpoint{7.785919in}{8.086538in}}%
\pgfpathquadraticcurveto{\pgfqpoint{7.777755in}{8.078081in}}{\pgfqpoint{7.762853in}{8.062886in}}%
\pgfpathlineto{\pgfqpoint{7.762853in}{8.062886in}}%
\pgfpathclose%
\pgfpathmoveto{\pgfqpoint{7.831759in}{8.068022in}}%
\pgfpathlineto{\pgfqpoint{7.844650in}{8.068022in}}%
\pgfpathlineto{\pgfqpoint{7.844650in}{8.052515in}}%
\pgfpathlineto{\pgfqpoint{7.831759in}{8.052515in}}%
\pgfpathlineto{\pgfqpoint{7.831759in}{8.068022in}}%
\pgfpathlineto{\pgfqpoint{7.831759in}{8.068022in}}%
\pgfpathclose%
\pgfpathmoveto{\pgfqpoint{7.871623in}{8.143647in}}%
\pgfpathlineto{\pgfqpoint{7.920021in}{8.143647in}}%
\pgfpathlineto{\pgfqpoint{7.920021in}{8.133257in}}%
\pgfpathlineto{\pgfqpoint{7.882912in}{8.133257in}}%
\pgfpathlineto{\pgfqpoint{7.882912in}{8.110933in}}%
\pgfpathquadraticcurveto{\pgfqpoint{7.885587in}{8.111851in}}{\pgfqpoint{7.888263in}{8.112300in}}%
\pgfpathquadraticcurveto{\pgfqpoint{7.890959in}{8.112749in}}{\pgfqpoint{7.893654in}{8.112749in}}%
\pgfpathquadraticcurveto{\pgfqpoint{7.908908in}{8.112749in}}{\pgfqpoint{7.917814in}{8.104390in}}%
\pgfpathquadraticcurveto{\pgfqpoint{7.926740in}{8.096030in}}{\pgfqpoint{7.926740in}{8.081753in}}%
\pgfpathquadraticcurveto{\pgfqpoint{7.926740in}{8.067046in}}{\pgfqpoint{7.917580in}{8.058882in}}%
\pgfpathquadraticcurveto{\pgfqpoint{7.908419in}{8.050737in}}{\pgfqpoint{7.891759in}{8.050737in}}%
\pgfpathquadraticcurveto{\pgfqpoint{7.886017in}{8.050737in}}{\pgfqpoint{7.880060in}{8.051714in}}%
\pgfpathquadraticcurveto{\pgfqpoint{7.874123in}{8.052690in}}{\pgfqpoint{7.867775in}{8.054644in}}%
\pgfpathlineto{\pgfqpoint{7.867775in}{8.067046in}}%
\pgfpathquadraticcurveto{\pgfqpoint{7.873263in}{8.064058in}}{\pgfqpoint{7.879123in}{8.062593in}}%
\pgfpathquadraticcurveto{\pgfqpoint{7.884982in}{8.061128in}}{\pgfqpoint{7.891505in}{8.061128in}}%
\pgfpathquadraticcurveto{\pgfqpoint{7.902072in}{8.061128in}}{\pgfqpoint{7.908224in}{8.066675in}}%
\pgfpathquadraticcurveto{\pgfqpoint{7.914396in}{8.072222in}}{\pgfqpoint{7.914396in}{8.081753in}}%
\pgfpathquadraticcurveto{\pgfqpoint{7.914396in}{8.091265in}}{\pgfqpoint{7.908224in}{8.096812in}}%
\pgfpathquadraticcurveto{\pgfqpoint{7.902072in}{8.102378in}}{\pgfqpoint{7.891505in}{8.102378in}}%
\pgfpathquadraticcurveto{\pgfqpoint{7.886564in}{8.102378in}}{\pgfqpoint{7.881642in}{8.101284in}}%
\pgfpathquadraticcurveto{\pgfqpoint{7.876740in}{8.100190in}}{\pgfqpoint{7.871623in}{8.097866in}}%
\pgfpathlineto{\pgfqpoint{7.871623in}{8.143647in}}%
\pgfpathlineto{\pgfqpoint{7.871623in}{8.143647in}}%
\pgfpathclose%
\pgfusepath{fill}%
\end{pgfscope}%
\begin{pgfscope}%
\pgfpathrectangle{\pgfqpoint{8.036091in}{4.050000in}}{\pgfqpoint{5.400000in}{5.400000in}}%
\pgfusepath{clip}%
\pgfsetrectcap%
\pgfsetroundjoin%
\pgfsetlinewidth{0.602250pt}%
\definecolor{currentstroke}{rgb}{0.898039,0.913725,0.929412}%
\pgfsetstrokecolor{currentstroke}%
\pgfsetdash{}{0pt}%
\pgfpathmoveto{\pgfqpoint{8.036091in}{9.450000in}}%
\pgfpathlineto{\pgfqpoint{13.436091in}{9.450000in}}%
\pgfusepath{stroke}%
\end{pgfscope}%
\begin{pgfscope}%
\pgfsetbuttcap%
\pgfsetroundjoin%
\definecolor{currentfill}{rgb}{0.090196,0.168627,0.227451}%
\pgfsetfillcolor{currentfill}%
\pgfsetlinewidth{0.803000pt}%
\definecolor{currentstroke}{rgb}{0.090196,0.168627,0.227451}%
\pgfsetstrokecolor{currentstroke}%
\pgfsetdash{}{0pt}%
\pgfsys@defobject{currentmarker}{\pgfqpoint{-0.048611in}{0.000000in}}{\pgfqpoint{-0.000000in}{0.000000in}}{%
\pgfpathmoveto{\pgfqpoint{-0.000000in}{0.000000in}}%
\pgfpathlineto{\pgfqpoint{-0.048611in}{0.000000in}}%
\pgfusepath{stroke,fill}%
}%
\begin{pgfscope}%
\pgfsys@transformshift{8.036091in}{9.450000in}%
\pgfsys@useobject{currentmarker}{}%
\end{pgfscope}%
\end{pgfscope}%
\begin{pgfscope}%
\pgfsetbuttcap%
\pgfsetroundjoin%
\definecolor{currentfill}{rgb}{0.090196,0.168627,0.227451}%
\pgfsetfillcolor{currentfill}%
\pgfsetlinewidth{0.000000pt}%
\definecolor{currentstroke}{rgb}{0.090196,0.168627,0.227451}%
\pgfsetstrokecolor{currentstroke}%
\pgfsetdash{}{0pt}%
\pgfpathmoveto{\pgfqpoint{7.752365in}{9.493647in}}%
\pgfpathlineto{\pgfqpoint{7.800763in}{9.493647in}}%
\pgfpathlineto{\pgfqpoint{7.800763in}{9.483257in}}%
\pgfpathlineto{\pgfqpoint{7.763654in}{9.483257in}}%
\pgfpathlineto{\pgfqpoint{7.763654in}{9.460933in}}%
\pgfpathquadraticcurveto{\pgfqpoint{7.766330in}{9.461851in}}{\pgfqpoint{7.769005in}{9.462300in}}%
\pgfpathquadraticcurveto{\pgfqpoint{7.771701in}{9.462749in}}{\pgfqpoint{7.774396in}{9.462749in}}%
\pgfpathquadraticcurveto{\pgfqpoint{7.789650in}{9.462749in}}{\pgfqpoint{7.798556in}{9.454390in}}%
\pgfpathquadraticcurveto{\pgfqpoint{7.807482in}{9.446030in}}{\pgfqpoint{7.807482in}{9.431753in}}%
\pgfpathquadraticcurveto{\pgfqpoint{7.807482in}{9.417046in}}{\pgfqpoint{7.798322in}{9.408882in}}%
\pgfpathquadraticcurveto{\pgfqpoint{7.789162in}{9.400737in}}{\pgfqpoint{7.772501in}{9.400737in}}%
\pgfpathquadraticcurveto{\pgfqpoint{7.766759in}{9.400737in}}{\pgfqpoint{7.760802in}{9.401714in}}%
\pgfpathquadraticcurveto{\pgfqpoint{7.754865in}{9.402690in}}{\pgfqpoint{7.748517in}{9.404644in}}%
\pgfpathlineto{\pgfqpoint{7.748517in}{9.417046in}}%
\pgfpathquadraticcurveto{\pgfqpoint{7.754005in}{9.414058in}}{\pgfqpoint{7.759865in}{9.412593in}}%
\pgfpathquadraticcurveto{\pgfqpoint{7.765724in}{9.411128in}}{\pgfqpoint{7.772248in}{9.411128in}}%
\pgfpathquadraticcurveto{\pgfqpoint{7.782814in}{9.411128in}}{\pgfqpoint{7.788966in}{9.416675in}}%
\pgfpathquadraticcurveto{\pgfqpoint{7.795138in}{9.422222in}}{\pgfqpoint{7.795138in}{9.431753in}}%
\pgfpathquadraticcurveto{\pgfqpoint{7.795138in}{9.441265in}}{\pgfqpoint{7.788966in}{9.446812in}}%
\pgfpathquadraticcurveto{\pgfqpoint{7.782814in}{9.452378in}}{\pgfqpoint{7.772248in}{9.452378in}}%
\pgfpathquadraticcurveto{\pgfqpoint{7.767306in}{9.452378in}}{\pgfqpoint{7.762384in}{9.451284in}}%
\pgfpathquadraticcurveto{\pgfqpoint{7.757482in}{9.450190in}}{\pgfqpoint{7.752365in}{9.447866in}}%
\pgfpathlineto{\pgfqpoint{7.752365in}{9.493647in}}%
\pgfpathlineto{\pgfqpoint{7.752365in}{9.493647in}}%
\pgfpathclose%
\pgfpathmoveto{\pgfqpoint{7.831759in}{9.418022in}}%
\pgfpathlineto{\pgfqpoint{7.844650in}{9.418022in}}%
\pgfpathlineto{\pgfqpoint{7.844650in}{9.402515in}}%
\pgfpathlineto{\pgfqpoint{7.831759in}{9.402515in}}%
\pgfpathlineto{\pgfqpoint{7.831759in}{9.418022in}}%
\pgfpathlineto{\pgfqpoint{7.831759in}{9.418022in}}%
\pgfpathclose%
\pgfpathmoveto{\pgfqpoint{7.897853in}{9.485522in}}%
\pgfpathquadraticcurveto{\pgfqpoint{7.888341in}{9.485522in}}{\pgfqpoint{7.883537in}{9.476147in}}%
\pgfpathquadraticcurveto{\pgfqpoint{7.878751in}{9.466792in}}{\pgfqpoint{7.878751in}{9.447983in}}%
\pgfpathquadraticcurveto{\pgfqpoint{7.878751in}{9.429253in}}{\pgfqpoint{7.883537in}{9.419878in}}%
\pgfpathquadraticcurveto{\pgfqpoint{7.888341in}{9.410503in}}{\pgfqpoint{7.897853in}{9.410503in}}%
\pgfpathquadraticcurveto{\pgfqpoint{7.907443in}{9.410503in}}{\pgfqpoint{7.912228in}{9.419878in}}%
\pgfpathquadraticcurveto{\pgfqpoint{7.917033in}{9.429253in}}{\pgfqpoint{7.917033in}{9.447983in}}%
\pgfpathquadraticcurveto{\pgfqpoint{7.917033in}{9.466792in}}{\pgfqpoint{7.912228in}{9.476147in}}%
\pgfpathquadraticcurveto{\pgfqpoint{7.907443in}{9.485522in}}{\pgfqpoint{7.897853in}{9.485522in}}%
\pgfpathlineto{\pgfqpoint{7.897853in}{9.485522in}}%
\pgfpathclose%
\pgfpathmoveto{\pgfqpoint{7.897853in}{9.495288in}}%
\pgfpathquadraticcurveto{\pgfqpoint{7.913185in}{9.495288in}}{\pgfqpoint{7.921271in}{9.483159in}}%
\pgfpathquadraticcurveto{\pgfqpoint{7.929357in}{9.471050in}}{\pgfqpoint{7.929357in}{9.447983in}}%
\pgfpathquadraticcurveto{\pgfqpoint{7.929357in}{9.424976in}}{\pgfqpoint{7.921271in}{9.412847in}}%
\pgfpathquadraticcurveto{\pgfqpoint{7.913185in}{9.400737in}}{\pgfqpoint{7.897853in}{9.400737in}}%
\pgfpathquadraticcurveto{\pgfqpoint{7.882541in}{9.400737in}}{\pgfqpoint{7.874455in}{9.412847in}}%
\pgfpathquadraticcurveto{\pgfqpoint{7.866369in}{9.424976in}}{\pgfqpoint{7.866369in}{9.447983in}}%
\pgfpathquadraticcurveto{\pgfqpoint{7.866369in}{9.471050in}}{\pgfqpoint{7.874455in}{9.483159in}}%
\pgfpathquadraticcurveto{\pgfqpoint{7.882541in}{9.495288in}}{\pgfqpoint{7.897853in}{9.495288in}}%
\pgfpathlineto{\pgfqpoint{7.897853in}{9.495288in}}%
\pgfpathclose%
\pgfusepath{fill}%
\end{pgfscope}%
\begin{pgfscope}%
\pgfsetbuttcap%
\pgfsetroundjoin%
\definecolor{currentfill}{rgb}{0.090196,0.168627,0.227451}%
\pgfsetfillcolor{currentfill}%
\pgfsetlinewidth{0.000000pt}%
\definecolor{currentstroke}{rgb}{0.090196,0.168627,0.227451}%
\pgfsetstrokecolor{currentstroke}%
\pgfsetdash{}{0pt}%
\pgfpathmoveto{\pgfqpoint{7.427869in}{5.850054in}}%
\pgfpathlineto{\pgfqpoint{7.468666in}{5.850054in}}%
\pgfpathlineto{\pgfqpoint{7.468666in}{5.874236in}}%
\pgfpathquadraticcurveto{\pgfqpoint{7.468666in}{5.886387in}}{\pgfqpoint{7.463629in}{5.892235in}}%
\pgfpathquadraticcurveto{\pgfqpoint{7.458592in}{5.898108in}}{\pgfqpoint{7.448232in}{5.898108in}}%
\pgfpathquadraticcurveto{\pgfqpoint{7.437776in}{5.898108in}}{\pgfqpoint{7.432835in}{5.892235in}}%
\pgfpathquadraticcurveto{\pgfqpoint{7.427869in}{5.886387in}}{\pgfqpoint{7.427869in}{5.874236in}}%
\pgfpathlineto{\pgfqpoint{7.427869in}{5.850054in}}%
\pgfpathlineto{\pgfqpoint{7.427869in}{5.850054in}}%
\pgfpathclose%
\pgfpathmoveto{\pgfqpoint{7.382060in}{5.850054in}}%
\pgfpathlineto{\pgfqpoint{7.415623in}{5.850054in}}%
\pgfpathlineto{\pgfqpoint{7.415623in}{5.872374in}}%
\pgfpathquadraticcurveto{\pgfqpoint{7.415623in}{5.883403in}}{\pgfqpoint{7.411493in}{5.888798in}}%
\pgfpathquadraticcurveto{\pgfqpoint{7.407340in}{5.894217in}}{\pgfqpoint{7.398842in}{5.894217in}}%
\pgfpathquadraticcurveto{\pgfqpoint{7.390415in}{5.894217in}}{\pgfqpoint{7.386237in}{5.888798in}}%
\pgfpathquadraticcurveto{\pgfqpoint{7.382060in}{5.883403in}}{\pgfqpoint{7.382060in}{5.872374in}}%
\pgfpathlineto{\pgfqpoint{7.382060in}{5.850054in}}%
\pgfpathlineto{\pgfqpoint{7.382060in}{5.850054in}}%
\pgfpathclose%
\pgfpathmoveto{\pgfqpoint{7.369671in}{5.834991in}}%
\pgfpathlineto{\pgfqpoint{7.369671in}{5.873496in}}%
\pgfpathquadraticcurveto{\pgfqpoint{7.369671in}{5.890731in}}{\pgfqpoint{7.376832in}{5.900041in}}%
\pgfpathquadraticcurveto{\pgfqpoint{7.383993in}{5.909375in}}{\pgfqpoint{7.397194in}{5.909375in}}%
\pgfpathquadraticcurveto{\pgfqpoint{7.407435in}{5.909375in}}{\pgfqpoint{7.413475in}{5.904601in}}%
\pgfpathquadraticcurveto{\pgfqpoint{7.419514in}{5.899826in}}{\pgfqpoint{7.420994in}{5.890564in}}%
\pgfpathquadraticcurveto{\pgfqpoint{7.423381in}{5.901688in}}{\pgfqpoint{7.430973in}{5.907847in}}%
\pgfpathquadraticcurveto{\pgfqpoint{7.438540in}{5.914006in}}{\pgfqpoint{7.449879in}{5.914006in}}%
\pgfpathquadraticcurveto{\pgfqpoint{7.464799in}{5.914006in}}{\pgfqpoint{7.472939in}{5.903861in}}%
\pgfpathquadraticcurveto{\pgfqpoint{7.481055in}{5.893715in}}{\pgfqpoint{7.481055in}{5.874976in}}%
\pgfpathlineto{\pgfqpoint{7.481055in}{5.834991in}}%
\pgfpathlineto{\pgfqpoint{7.369671in}{5.834991in}}%
\pgfpathlineto{\pgfqpoint{7.369671in}{5.834991in}}%
\pgfpathclose%
\pgfpathmoveto{\pgfqpoint{7.369671in}{5.939811in}}%
\pgfpathlineto{\pgfqpoint{7.369671in}{5.954874in}}%
\pgfpathlineto{\pgfqpoint{7.473297in}{5.954874in}}%
\pgfpathquadraticcurveto{\pgfqpoint{7.493444in}{5.954874in}}{\pgfqpoint{7.502539in}{5.947235in}}%
\pgfpathquadraticcurveto{\pgfqpoint{7.511635in}{5.939596in}}{\pgfqpoint{7.511635in}{5.922648in}}%
\pgfpathlineto{\pgfqpoint{7.511635in}{5.916918in}}%
\pgfpathlineto{\pgfqpoint{7.498959in}{5.916918in}}%
\pgfpathlineto{\pgfqpoint{7.498959in}{5.921621in}}%
\pgfpathquadraticcurveto{\pgfqpoint{7.498959in}{5.931599in}}{\pgfqpoint{7.493349in}{5.935705in}}%
\pgfpathquadraticcurveto{\pgfqpoint{7.487763in}{5.939811in}}{\pgfqpoint{7.473297in}{5.939811in}}%
\pgfpathlineto{\pgfqpoint{7.369671in}{5.939811in}}%
\pgfpathlineto{\pgfqpoint{7.369671in}{5.939811in}}%
\pgfpathclose%
\pgfpathmoveto{\pgfqpoint{7.373323in}{6.051649in}}%
\pgfpathlineto{\pgfqpoint{7.388028in}{6.051649in}}%
\pgfpathquadraticcurveto{\pgfqpoint{7.383922in}{6.043079in}}{\pgfqpoint{7.381917in}{6.035464in}}%
\pgfpathquadraticcurveto{\pgfqpoint{7.379888in}{6.027849in}}{\pgfqpoint{7.379888in}{6.020760in}}%
\pgfpathquadraticcurveto{\pgfqpoint{7.379888in}{6.008466in}}{\pgfqpoint{7.384662in}{6.001782in}}%
\pgfpathquadraticcurveto{\pgfqpoint{7.389436in}{5.995098in}}{\pgfqpoint{7.398245in}{5.995098in}}%
\pgfpathquadraticcurveto{\pgfqpoint{7.405645in}{5.995098in}}{\pgfqpoint{7.409417in}{5.999538in}}%
\pgfpathquadraticcurveto{\pgfqpoint{7.413164in}{6.003978in}}{\pgfqpoint{7.415480in}{6.016367in}}%
\pgfpathlineto{\pgfqpoint{7.417342in}{6.025462in}}%
\pgfpathquadraticcurveto{\pgfqpoint{7.420565in}{6.042316in}}{\pgfqpoint{7.428657in}{6.050336in}}%
\pgfpathquadraticcurveto{\pgfqpoint{7.436750in}{6.058357in}}{\pgfqpoint{7.450309in}{6.058357in}}%
\pgfpathquadraticcurveto{\pgfqpoint{7.466517in}{6.058357in}}{\pgfqpoint{7.474872in}{6.047496in}}%
\pgfpathquadraticcurveto{\pgfqpoint{7.483227in}{6.036658in}}{\pgfqpoint{7.483227in}{6.015699in}}%
\pgfpathquadraticcurveto{\pgfqpoint{7.483227in}{6.007797in}}{\pgfqpoint{7.481437in}{5.998869in}}%
\pgfpathquadraticcurveto{\pgfqpoint{7.479647in}{5.989965in}}{\pgfqpoint{7.476138in}{5.980417in}}%
\pgfpathlineto{\pgfqpoint{7.460621in}{5.980417in}}%
\pgfpathquadraticcurveto{\pgfqpoint{7.465753in}{5.989583in}}{\pgfqpoint{7.468379in}{5.998392in}}%
\pgfpathquadraticcurveto{\pgfqpoint{7.470981in}{6.007201in}}{\pgfqpoint{7.470981in}{6.015699in}}%
\pgfpathquadraticcurveto{\pgfqpoint{7.470981in}{6.028589in}}{\pgfqpoint{7.465921in}{6.035608in}}%
\pgfpathquadraticcurveto{\pgfqpoint{7.460836in}{6.042626in}}{\pgfqpoint{7.451431in}{6.042626in}}%
\pgfpathquadraticcurveto{\pgfqpoint{7.443243in}{6.042626in}}{\pgfqpoint{7.438612in}{6.037589in}}%
\pgfpathquadraticcurveto{\pgfqpoint{7.433980in}{6.032552in}}{\pgfqpoint{7.431665in}{6.021070in}}%
\pgfpathlineto{\pgfqpoint{7.429875in}{6.011879in}}%
\pgfpathquadraticcurveto{\pgfqpoint{7.426533in}{5.995026in}}{\pgfqpoint{7.419371in}{5.987483in}}%
\pgfpathquadraticcurveto{\pgfqpoint{7.412210in}{5.979963in}}{\pgfqpoint{7.399438in}{5.979963in}}%
\pgfpathquadraticcurveto{\pgfqpoint{7.384662in}{5.979963in}}{\pgfqpoint{7.376164in}{5.990371in}}%
\pgfpathquadraticcurveto{\pgfqpoint{7.367665in}{6.000779in}}{\pgfqpoint{7.367665in}{6.019041in}}%
\pgfpathquadraticcurveto{\pgfqpoint{7.367665in}{6.026895in}}{\pgfqpoint{7.369098in}{6.035011in}}%
\pgfpathquadraticcurveto{\pgfqpoint{7.370506in}{6.043151in}}{\pgfqpoint{7.373323in}{6.051649in}}%
\pgfpathlineto{\pgfqpoint{7.373323in}{6.051649in}}%
\pgfpathclose%
\pgfpathmoveto{\pgfqpoint{7.435842in}{6.201300in}}%
\pgfpathlineto{\pgfqpoint{7.442550in}{6.201300in}}%
\pgfpathlineto{\pgfqpoint{7.442550in}{6.138184in}}%
\pgfpathquadraticcurveto{\pgfqpoint{7.456730in}{6.139091in}}{\pgfqpoint{7.464154in}{6.146730in}}%
\pgfpathquadraticcurveto{\pgfqpoint{7.471578in}{6.154368in}}{\pgfqpoint{7.471578in}{6.168023in}}%
\pgfpathquadraticcurveto{\pgfqpoint{7.471578in}{6.175924in}}{\pgfqpoint{7.469645in}{6.183349in}}%
\pgfpathquadraticcurveto{\pgfqpoint{7.467711in}{6.190773in}}{\pgfqpoint{7.463820in}{6.198101in}}%
\pgfpathlineto{\pgfqpoint{7.476806in}{6.198101in}}%
\pgfpathquadraticcurveto{\pgfqpoint{7.479933in}{6.190701in}}{\pgfqpoint{7.481580in}{6.182943in}}%
\pgfpathquadraticcurveto{\pgfqpoint{7.483227in}{6.175184in}}{\pgfqpoint{7.483227in}{6.167211in}}%
\pgfpathquadraticcurveto{\pgfqpoint{7.483227in}{6.147207in}}{\pgfqpoint{7.471602in}{6.135534in}}%
\pgfpathquadraticcurveto{\pgfqpoint{7.459953in}{6.123861in}}{\pgfqpoint{7.440092in}{6.123861in}}%
\pgfpathquadraticcurveto{\pgfqpoint{7.419586in}{6.123861in}}{\pgfqpoint{7.407555in}{6.134937in}}%
\pgfpathquadraticcurveto{\pgfqpoint{7.395500in}{6.146013in}}{\pgfqpoint{7.395500in}{6.164824in}}%
\pgfpathquadraticcurveto{\pgfqpoint{7.395500in}{6.181678in}}{\pgfqpoint{7.406361in}{6.191489in}}%
\pgfpathquadraticcurveto{\pgfqpoint{7.417199in}{6.201300in}}{\pgfqpoint{7.435842in}{6.201300in}}%
\pgfpathlineto{\pgfqpoint{7.435842in}{6.201300in}}%
\pgfpathclose%
\pgfpathmoveto{\pgfqpoint{7.431808in}{6.187574in}}%
\pgfpathquadraticcurveto{\pgfqpoint{7.420565in}{6.187431in}}{\pgfqpoint{7.413857in}{6.181272in}}%
\pgfpathquadraticcurveto{\pgfqpoint{7.407125in}{6.175113in}}{\pgfqpoint{7.407125in}{6.164967in}}%
\pgfpathquadraticcurveto{\pgfqpoint{7.407125in}{6.153485in}}{\pgfqpoint{7.413618in}{6.146586in}}%
\pgfpathquadraticcurveto{\pgfqpoint{7.420111in}{6.139687in}}{\pgfqpoint{7.431904in}{6.138637in}}%
\pgfpathlineto{\pgfqpoint{7.431808in}{6.187574in}}%
\pgfpathlineto{\pgfqpoint{7.431808in}{6.187574in}}%
\pgfpathclose%
\pgfpathmoveto{\pgfqpoint{7.399964in}{6.277092in}}%
\pgfpathlineto{\pgfqpoint{7.412950in}{6.277092in}}%
\pgfpathquadraticcurveto{\pgfqpoint{7.409966in}{6.271291in}}{\pgfqpoint{7.408486in}{6.265013in}}%
\pgfpathquadraticcurveto{\pgfqpoint{7.406982in}{6.258759in}}{\pgfqpoint{7.406982in}{6.252027in}}%
\pgfpathquadraticcurveto{\pgfqpoint{7.406982in}{6.241810in}}{\pgfqpoint{7.410109in}{6.236701in}}%
\pgfpathquadraticcurveto{\pgfqpoint{7.413236in}{6.231593in}}{\pgfqpoint{7.419514in}{6.231593in}}%
\pgfpathquadraticcurveto{\pgfqpoint{7.424289in}{6.231593in}}{\pgfqpoint{7.427010in}{6.235245in}}%
\pgfpathquadraticcurveto{\pgfqpoint{7.429731in}{6.238898in}}{\pgfqpoint{7.432190in}{6.249950in}}%
\pgfpathlineto{\pgfqpoint{7.433240in}{6.254653in}}%
\pgfpathquadraticcurveto{\pgfqpoint{7.436368in}{6.269262in}}{\pgfqpoint{7.442073in}{6.275421in}}%
\pgfpathquadraticcurveto{\pgfqpoint{7.447778in}{6.281580in}}{\pgfqpoint{7.457995in}{6.281580in}}%
\pgfpathquadraticcurveto{\pgfqpoint{7.469645in}{6.281580in}}{\pgfqpoint{7.476448in}{6.272365in}}%
\pgfpathquadraticcurveto{\pgfqpoint{7.483227in}{6.263151in}}{\pgfqpoint{7.483227in}{6.247038in}}%
\pgfpathquadraticcurveto{\pgfqpoint{7.483227in}{6.240330in}}{\pgfqpoint{7.481914in}{6.233049in}}%
\pgfpathquadraticcurveto{\pgfqpoint{7.480602in}{6.225768in}}{\pgfqpoint{7.478000in}{6.217724in}}%
\pgfpathlineto{\pgfqpoint{7.463820in}{6.217724in}}%
\pgfpathquadraticcurveto{\pgfqpoint{7.467783in}{6.225339in}}{\pgfqpoint{7.469764in}{6.232715in}}%
\pgfpathquadraticcurveto{\pgfqpoint{7.471721in}{6.240091in}}{\pgfqpoint{7.471721in}{6.247348in}}%
\pgfpathquadraticcurveto{\pgfqpoint{7.471721in}{6.257040in}}{\pgfqpoint{7.468403in}{6.262244in}}%
\pgfpathquadraticcurveto{\pgfqpoint{7.465085in}{6.267472in}}{\pgfqpoint{7.459046in}{6.267472in}}%
\pgfpathquadraticcurveto{\pgfqpoint{7.453460in}{6.267472in}}{\pgfqpoint{7.450476in}{6.263700in}}%
\pgfpathquadraticcurveto{\pgfqpoint{7.447492in}{6.259952in}}{\pgfqpoint{7.444723in}{6.247181in}}%
\pgfpathlineto{\pgfqpoint{7.443601in}{6.242407in}}%
\pgfpathquadraticcurveto{\pgfqpoint{7.440927in}{6.229659in}}{\pgfqpoint{7.435365in}{6.223978in}}%
\pgfpathquadraticcurveto{\pgfqpoint{7.429803in}{6.218320in}}{\pgfqpoint{7.420111in}{6.218320in}}%
\pgfpathquadraticcurveto{\pgfqpoint{7.408319in}{6.218320in}}{\pgfqpoint{7.401921in}{6.226675in}}%
\pgfpathquadraticcurveto{\pgfqpoint{7.395500in}{6.235030in}}{\pgfqpoint{7.395500in}{6.250404in}}%
\pgfpathquadraticcurveto{\pgfqpoint{7.395500in}{6.257995in}}{\pgfqpoint{7.396622in}{6.264703in}}%
\pgfpathquadraticcurveto{\pgfqpoint{7.397720in}{6.271434in}}{\pgfqpoint{7.399964in}{6.277092in}}%
\pgfpathlineto{\pgfqpoint{7.399964in}{6.277092in}}%
\pgfpathclose%
\pgfpathmoveto{\pgfqpoint{7.373776in}{6.317005in}}%
\pgfpathlineto{\pgfqpoint{7.397505in}{6.317005in}}%
\pgfpathlineto{\pgfqpoint{7.397505in}{6.345269in}}%
\pgfpathlineto{\pgfqpoint{7.408175in}{6.345269in}}%
\pgfpathlineto{\pgfqpoint{7.408175in}{6.317005in}}%
\pgfpathlineto{\pgfqpoint{7.453531in}{6.317005in}}%
\pgfpathquadraticcurveto{\pgfqpoint{7.463748in}{6.317005in}}{\pgfqpoint{7.466661in}{6.319798in}}%
\pgfpathquadraticcurveto{\pgfqpoint{7.469573in}{6.322591in}}{\pgfqpoint{7.469573in}{6.331185in}}%
\pgfpathlineto{\pgfqpoint{7.469573in}{6.345269in}}%
\pgfpathlineto{\pgfqpoint{7.481055in}{6.345269in}}%
\pgfpathlineto{\pgfqpoint{7.481055in}{6.331185in}}%
\pgfpathquadraticcurveto{\pgfqpoint{7.481055in}{6.315286in}}{\pgfqpoint{7.475135in}{6.309247in}}%
\pgfpathquadraticcurveto{\pgfqpoint{7.469191in}{6.303207in}}{\pgfqpoint{7.453531in}{6.303207in}}%
\pgfpathlineto{\pgfqpoint{7.408175in}{6.303207in}}%
\pgfpathlineto{\pgfqpoint{7.408175in}{6.293134in}}%
\pgfpathlineto{\pgfqpoint{7.397505in}{6.293134in}}%
\pgfpathlineto{\pgfqpoint{7.397505in}{6.303207in}}%
\pgfpathlineto{\pgfqpoint{7.373776in}{6.303207in}}%
\pgfpathlineto{\pgfqpoint{7.373776in}{6.317005in}}%
\pgfpathlineto{\pgfqpoint{7.373776in}{6.317005in}}%
\pgfpathclose%
\pgfpathmoveto{\pgfqpoint{7.397505in}{6.363316in}}%
\pgfpathlineto{\pgfqpoint{7.397505in}{6.377042in}}%
\pgfpathlineto{\pgfqpoint{7.481055in}{6.377042in}}%
\pgfpathlineto{\pgfqpoint{7.481055in}{6.363316in}}%
\pgfpathlineto{\pgfqpoint{7.397505in}{6.363316in}}%
\pgfpathlineto{\pgfqpoint{7.397505in}{6.363316in}}%
\pgfpathclose%
\pgfpathmoveto{\pgfqpoint{7.364968in}{6.363316in}}%
\pgfpathlineto{\pgfqpoint{7.364968in}{6.377042in}}%
\pgfpathlineto{\pgfqpoint{7.382370in}{6.377042in}}%
\pgfpathlineto{\pgfqpoint{7.382370in}{6.363316in}}%
\pgfpathlineto{\pgfqpoint{7.364968in}{6.363316in}}%
\pgfpathlineto{\pgfqpoint{7.364968in}{6.363316in}}%
\pgfpathclose%
\pgfpathmoveto{\pgfqpoint{7.413546in}{6.470809in}}%
\pgfpathquadraticcurveto{\pgfqpoint{7.404284in}{6.475966in}}{\pgfqpoint{7.399892in}{6.483127in}}%
\pgfpathquadraticcurveto{\pgfqpoint{7.395500in}{6.490289in}}{\pgfqpoint{7.395500in}{6.499980in}}%
\pgfpathquadraticcurveto{\pgfqpoint{7.395500in}{6.513038in}}{\pgfqpoint{7.404642in}{6.520128in}}%
\pgfpathquadraticcurveto{\pgfqpoint{7.413761in}{6.527218in}}{\pgfqpoint{7.430615in}{6.527218in}}%
\pgfpathlineto{\pgfqpoint{7.481055in}{6.527218in}}%
\pgfpathlineto{\pgfqpoint{7.481055in}{6.513420in}}%
\pgfpathlineto{\pgfqpoint{7.431068in}{6.513420in}}%
\pgfpathquadraticcurveto{\pgfqpoint{7.419061in}{6.513420in}}{\pgfqpoint{7.413260in}{6.509147in}}%
\pgfpathquadraticcurveto{\pgfqpoint{7.407435in}{6.504898in}}{\pgfqpoint{7.407435in}{6.496185in}}%
\pgfpathquadraticcurveto{\pgfqpoint{7.407435in}{6.485514in}}{\pgfqpoint{7.414525in}{6.479308in}}%
\pgfpathquadraticcurveto{\pgfqpoint{7.421591in}{6.473125in}}{\pgfqpoint{7.433837in}{6.473125in}}%
\pgfpathlineto{\pgfqpoint{7.481055in}{6.473125in}}%
\pgfpathlineto{\pgfqpoint{7.481055in}{6.459327in}}%
\pgfpathlineto{\pgfqpoint{7.431068in}{6.459327in}}%
\pgfpathquadraticcurveto{\pgfqpoint{7.418989in}{6.459327in}}{\pgfqpoint{7.413212in}{6.455078in}}%
\pgfpathquadraticcurveto{\pgfqpoint{7.407435in}{6.450829in}}{\pgfqpoint{7.407435in}{6.441949in}}%
\pgfpathquadraticcurveto{\pgfqpoint{7.407435in}{6.431421in}}{\pgfqpoint{7.414549in}{6.425215in}}%
\pgfpathquadraticcurveto{\pgfqpoint{7.421663in}{6.419032in}}{\pgfqpoint{7.433837in}{6.419032in}}%
\pgfpathlineto{\pgfqpoint{7.481055in}{6.419032in}}%
\pgfpathlineto{\pgfqpoint{7.481055in}{6.405234in}}%
\pgfpathlineto{\pgfqpoint{7.397505in}{6.405234in}}%
\pgfpathlineto{\pgfqpoint{7.397505in}{6.419032in}}%
\pgfpathlineto{\pgfqpoint{7.410491in}{6.419032in}}%
\pgfpathquadraticcurveto{\pgfqpoint{7.402804in}{6.423735in}}{\pgfqpoint{7.399152in}{6.430299in}}%
\pgfpathquadraticcurveto{\pgfqpoint{7.395500in}{6.436864in}}{\pgfqpoint{7.395500in}{6.445888in}}%
\pgfpathquadraticcurveto{\pgfqpoint{7.395500in}{6.455007in}}{\pgfqpoint{7.400131in}{6.461380in}}%
\pgfpathquadraticcurveto{\pgfqpoint{7.404738in}{6.467754in}}{\pgfqpoint{7.413546in}{6.470809in}}%
\pgfpathlineto{\pgfqpoint{7.413546in}{6.470809in}}%
\pgfpathclose%
\pgfpathmoveto{\pgfqpoint{7.439065in}{6.592554in}}%
\pgfpathquadraticcurveto{\pgfqpoint{7.439065in}{6.575916in}}{\pgfqpoint{7.442861in}{6.569494in}}%
\pgfpathquadraticcurveto{\pgfqpoint{7.446656in}{6.563073in}}{\pgfqpoint{7.455847in}{6.563073in}}%
\pgfpathquadraticcurveto{\pgfqpoint{7.463151in}{6.563073in}}{\pgfqpoint{7.467448in}{6.567895in}}%
\pgfpathquadraticcurveto{\pgfqpoint{7.471721in}{6.572717in}}{\pgfqpoint{7.471721in}{6.580977in}}%
\pgfpathquadraticcurveto{\pgfqpoint{7.471721in}{6.592411in}}{\pgfqpoint{7.463629in}{6.599310in}}%
\pgfpathquadraticcurveto{\pgfqpoint{7.455536in}{6.606209in}}{\pgfqpoint{7.442121in}{6.606209in}}%
\pgfpathlineto{\pgfqpoint{7.439065in}{6.606209in}}%
\pgfpathlineto{\pgfqpoint{7.439065in}{6.592554in}}%
\pgfpathlineto{\pgfqpoint{7.439065in}{6.592554in}}%
\pgfpathclose%
\pgfpathmoveto{\pgfqpoint{7.433384in}{6.619935in}}%
\pgfpathlineto{\pgfqpoint{7.481055in}{6.619935in}}%
\pgfpathlineto{\pgfqpoint{7.481055in}{6.606209in}}%
\pgfpathlineto{\pgfqpoint{7.468379in}{6.606209in}}%
\pgfpathquadraticcurveto{\pgfqpoint{7.475970in}{6.601506in}}{\pgfqpoint{7.479599in}{6.594488in}}%
\pgfpathquadraticcurveto{\pgfqpoint{7.483227in}{6.587470in}}{\pgfqpoint{7.483227in}{6.577324in}}%
\pgfpathquadraticcurveto{\pgfqpoint{7.483227in}{6.564505in}}{\pgfqpoint{7.476018in}{6.556914in}}%
\pgfpathquadraticcurveto{\pgfqpoint{7.468809in}{6.549347in}}{\pgfqpoint{7.456730in}{6.549347in}}%
\pgfpathquadraticcurveto{\pgfqpoint{7.442646in}{6.549347in}}{\pgfqpoint{7.435484in}{6.558776in}}%
\pgfpathquadraticcurveto{\pgfqpoint{7.428323in}{6.568229in}}{\pgfqpoint{7.428323in}{6.586944in}}%
\pgfpathlineto{\pgfqpoint{7.428323in}{6.606209in}}%
\pgfpathlineto{\pgfqpoint{7.426962in}{6.606209in}}%
\pgfpathquadraticcurveto{\pgfqpoint{7.417485in}{6.606209in}}{\pgfqpoint{7.412305in}{6.599978in}}%
\pgfpathquadraticcurveto{\pgfqpoint{7.407125in}{6.593748in}}{\pgfqpoint{7.407125in}{6.582480in}}%
\pgfpathquadraticcurveto{\pgfqpoint{7.407125in}{6.575319in}}{\pgfqpoint{7.408844in}{6.568516in}}%
\pgfpathquadraticcurveto{\pgfqpoint{7.410562in}{6.561736in}}{\pgfqpoint{7.414000in}{6.555482in}}%
\pgfpathlineto{\pgfqpoint{7.401300in}{6.555482in}}%
\pgfpathquadraticcurveto{\pgfqpoint{7.398388in}{6.563001in}}{\pgfqpoint{7.396956in}{6.570091in}}%
\pgfpathquadraticcurveto{\pgfqpoint{7.395500in}{6.577181in}}{\pgfqpoint{7.395500in}{6.583889in}}%
\pgfpathquadraticcurveto{\pgfqpoint{7.395500in}{6.602031in}}{\pgfqpoint{7.404905in}{6.610983in}}%
\pgfpathquadraticcurveto{\pgfqpoint{7.414286in}{6.619935in}}{\pgfqpoint{7.433384in}{6.619935in}}%
\pgfpathlineto{\pgfqpoint{7.433384in}{6.619935in}}%
\pgfpathclose%
\pgfpathmoveto{\pgfqpoint{7.373776in}{6.661782in}}%
\pgfpathlineto{\pgfqpoint{7.397505in}{6.661782in}}%
\pgfpathlineto{\pgfqpoint{7.397505in}{6.690046in}}%
\pgfpathlineto{\pgfqpoint{7.408175in}{6.690046in}}%
\pgfpathlineto{\pgfqpoint{7.408175in}{6.661782in}}%
\pgfpathlineto{\pgfqpoint{7.453531in}{6.661782in}}%
\pgfpathquadraticcurveto{\pgfqpoint{7.463748in}{6.661782in}}{\pgfqpoint{7.466661in}{6.664575in}}%
\pgfpathquadraticcurveto{\pgfqpoint{7.469573in}{6.667368in}}{\pgfqpoint{7.469573in}{6.675961in}}%
\pgfpathlineto{\pgfqpoint{7.469573in}{6.690046in}}%
\pgfpathlineto{\pgfqpoint{7.481055in}{6.690046in}}%
\pgfpathlineto{\pgfqpoint{7.481055in}{6.675961in}}%
\pgfpathquadraticcurveto{\pgfqpoint{7.481055in}{6.660063in}}{\pgfqpoint{7.475135in}{6.654023in}}%
\pgfpathquadraticcurveto{\pgfqpoint{7.469191in}{6.647984in}}{\pgfqpoint{7.453531in}{6.647984in}}%
\pgfpathlineto{\pgfqpoint{7.408175in}{6.647984in}}%
\pgfpathlineto{\pgfqpoint{7.408175in}{6.637910in}}%
\pgfpathlineto{\pgfqpoint{7.397505in}{6.637910in}}%
\pgfpathlineto{\pgfqpoint{7.397505in}{6.647984in}}%
\pgfpathlineto{\pgfqpoint{7.373776in}{6.647984in}}%
\pgfpathlineto{\pgfqpoint{7.373776in}{6.661782in}}%
\pgfpathlineto{\pgfqpoint{7.373776in}{6.661782in}}%
\pgfpathclose%
\pgfpathmoveto{\pgfqpoint{7.435842in}{6.779564in}}%
\pgfpathlineto{\pgfqpoint{7.442550in}{6.779564in}}%
\pgfpathlineto{\pgfqpoint{7.442550in}{6.716447in}}%
\pgfpathquadraticcurveto{\pgfqpoint{7.456730in}{6.717355in}}{\pgfqpoint{7.464154in}{6.724993in}}%
\pgfpathquadraticcurveto{\pgfqpoint{7.471578in}{6.732632in}}{\pgfqpoint{7.471578in}{6.746287in}}%
\pgfpathquadraticcurveto{\pgfqpoint{7.471578in}{6.754188in}}{\pgfqpoint{7.469645in}{6.761612in}}%
\pgfpathquadraticcurveto{\pgfqpoint{7.467711in}{6.769036in}}{\pgfqpoint{7.463820in}{6.776365in}}%
\pgfpathlineto{\pgfqpoint{7.476806in}{6.776365in}}%
\pgfpathquadraticcurveto{\pgfqpoint{7.479933in}{6.768965in}}{\pgfqpoint{7.481580in}{6.761207in}}%
\pgfpathquadraticcurveto{\pgfqpoint{7.483227in}{6.753448in}}{\pgfqpoint{7.483227in}{6.745475in}}%
\pgfpathquadraticcurveto{\pgfqpoint{7.483227in}{6.725471in}}{\pgfqpoint{7.471602in}{6.713798in}}%
\pgfpathquadraticcurveto{\pgfqpoint{7.459953in}{6.702125in}}{\pgfqpoint{7.440092in}{6.702125in}}%
\pgfpathquadraticcurveto{\pgfqpoint{7.419586in}{6.702125in}}{\pgfqpoint{7.407555in}{6.713201in}}%
\pgfpathquadraticcurveto{\pgfqpoint{7.395500in}{6.724277in}}{\pgfqpoint{7.395500in}{6.743088in}}%
\pgfpathquadraticcurveto{\pgfqpoint{7.395500in}{6.759941in}}{\pgfqpoint{7.406361in}{6.769753in}}%
\pgfpathquadraticcurveto{\pgfqpoint{7.417199in}{6.779564in}}{\pgfqpoint{7.435842in}{6.779564in}}%
\pgfpathlineto{\pgfqpoint{7.435842in}{6.779564in}}%
\pgfpathclose%
\pgfpathmoveto{\pgfqpoint{7.431808in}{6.765838in}}%
\pgfpathquadraticcurveto{\pgfqpoint{7.420565in}{6.765694in}}{\pgfqpoint{7.413857in}{6.759536in}}%
\pgfpathquadraticcurveto{\pgfqpoint{7.407125in}{6.753377in}}{\pgfqpoint{7.407125in}{6.743231in}}%
\pgfpathquadraticcurveto{\pgfqpoint{7.407125in}{6.731749in}}{\pgfqpoint{7.413618in}{6.724850in}}%
\pgfpathquadraticcurveto{\pgfqpoint{7.420111in}{6.717951in}}{\pgfqpoint{7.431904in}{6.716901in}}%
\pgfpathlineto{\pgfqpoint{7.431808in}{6.765838in}}%
\pgfpathlineto{\pgfqpoint{7.431808in}{6.765838in}}%
\pgfpathclose%
\pgfpathmoveto{\pgfqpoint{7.365135in}{6.883620in}}%
\pgfpathquadraticcurveto{\pgfqpoint{7.382275in}{6.873641in}}{\pgfqpoint{7.399080in}{6.868772in}}%
\pgfpathquadraticcurveto{\pgfqpoint{7.415862in}{6.863926in}}{\pgfqpoint{7.433097in}{6.863926in}}%
\pgfpathquadraticcurveto{\pgfqpoint{7.450309in}{6.863926in}}{\pgfqpoint{7.467210in}{6.868819in}}%
\pgfpathquadraticcurveto{\pgfqpoint{7.484111in}{6.873713in}}{\pgfqpoint{7.501203in}{6.883620in}}%
\pgfpathlineto{\pgfqpoint{7.501203in}{6.871684in}}%
\pgfpathquadraticcurveto{\pgfqpoint{7.483657in}{6.860512in}}{\pgfqpoint{7.466732in}{6.854950in}}%
\pgfpathquadraticcurveto{\pgfqpoint{7.449807in}{6.849388in}}{\pgfqpoint{7.433097in}{6.849388in}}%
\pgfpathquadraticcurveto{\pgfqpoint{7.416459in}{6.849388in}}{\pgfqpoint{7.399605in}{6.854902in}}%
\pgfpathquadraticcurveto{\pgfqpoint{7.382728in}{6.860441in}}{\pgfqpoint{7.365135in}{6.871684in}}%
\pgfpathlineto{\pgfqpoint{7.365135in}{6.883620in}}%
\pgfpathlineto{\pgfqpoint{7.365135in}{6.883620in}}%
\pgfpathclose%
\pgfpathmoveto{\pgfqpoint{7.369671in}{6.895412in}}%
\pgfpathlineto{\pgfqpoint{7.369671in}{6.989633in}}%
\pgfpathlineto{\pgfqpoint{7.382370in}{6.989633in}}%
\pgfpathlineto{\pgfqpoint{7.382370in}{6.950102in}}%
\pgfpathlineto{\pgfqpoint{7.481055in}{6.950102in}}%
\pgfpathlineto{\pgfqpoint{7.481055in}{6.934967in}}%
\pgfpathlineto{\pgfqpoint{7.382370in}{6.934967in}}%
\pgfpathlineto{\pgfqpoint{7.382370in}{6.895412in}}%
\pgfpathlineto{\pgfqpoint{7.369671in}{6.895412in}}%
\pgfpathlineto{\pgfqpoint{7.369671in}{6.895412in}}%
\pgfpathclose%
\pgfpathmoveto{\pgfqpoint{7.369671in}{6.994264in}}%
\pgfpathlineto{\pgfqpoint{7.369671in}{7.009470in}}%
\pgfpathlineto{\pgfqpoint{7.463820in}{7.032888in}}%
\pgfpathlineto{\pgfqpoint{7.369671in}{7.056235in}}%
\pgfpathlineto{\pgfqpoint{7.369671in}{7.073184in}}%
\pgfpathlineto{\pgfqpoint{7.463820in}{7.096602in}}%
\pgfpathlineto{\pgfqpoint{7.369671in}{7.119948in}}%
\pgfpathlineto{\pgfqpoint{7.369671in}{7.135250in}}%
\pgfpathlineto{\pgfqpoint{7.481055in}{7.107272in}}%
\pgfpathlineto{\pgfqpoint{7.481055in}{7.088318in}}%
\pgfpathlineto{\pgfqpoint{7.384375in}{7.064829in}}%
\pgfpathlineto{\pgfqpoint{7.481055in}{7.041100in}}%
\pgfpathlineto{\pgfqpoint{7.481055in}{7.022146in}}%
\pgfpathlineto{\pgfqpoint{7.369671in}{6.994264in}}%
\pgfpathlineto{\pgfqpoint{7.369671in}{6.994264in}}%
\pgfpathclose%
\pgfpathmoveto{\pgfqpoint{7.369671in}{7.155230in}}%
\pgfpathlineto{\pgfqpoint{7.369671in}{7.219230in}}%
\pgfpathlineto{\pgfqpoint{7.382370in}{7.219230in}}%
\pgfpathlineto{\pgfqpoint{7.382370in}{7.170293in}}%
\pgfpathlineto{\pgfqpoint{7.415194in}{7.170293in}}%
\pgfpathlineto{\pgfqpoint{7.415194in}{7.214455in}}%
\pgfpathlineto{\pgfqpoint{7.427869in}{7.214455in}}%
\pgfpathlineto{\pgfqpoint{7.427869in}{7.170293in}}%
\pgfpathlineto{\pgfqpoint{7.481055in}{7.170293in}}%
\pgfpathlineto{\pgfqpoint{7.481055in}{7.155230in}}%
\pgfpathlineto{\pgfqpoint{7.369671in}{7.155230in}}%
\pgfpathlineto{\pgfqpoint{7.369671in}{7.155230in}}%
\pgfpathclose%
\pgfpathmoveto{\pgfqpoint{7.369671in}{7.243101in}}%
\pgfpathlineto{\pgfqpoint{7.369671in}{7.313522in}}%
\pgfpathlineto{\pgfqpoint{7.382370in}{7.313522in}}%
\pgfpathlineto{\pgfqpoint{7.382370in}{7.258164in}}%
\pgfpathlineto{\pgfqpoint{7.415337in}{7.258164in}}%
\pgfpathlineto{\pgfqpoint{7.415337in}{7.311207in}}%
\pgfpathlineto{\pgfqpoint{7.428013in}{7.311207in}}%
\pgfpathlineto{\pgfqpoint{7.428013in}{7.258164in}}%
\pgfpathlineto{\pgfqpoint{7.468379in}{7.258164in}}%
\pgfpathlineto{\pgfqpoint{7.468379in}{7.314859in}}%
\pgfpathlineto{\pgfqpoint{7.481055in}{7.314859in}}%
\pgfpathlineto{\pgfqpoint{7.481055in}{7.243101in}}%
\pgfpathlineto{\pgfqpoint{7.369671in}{7.243101in}}%
\pgfpathlineto{\pgfqpoint{7.369671in}{7.243101in}}%
\pgfpathclose%
\pgfpathmoveto{\pgfqpoint{7.373323in}{7.454961in}}%
\pgfpathlineto{\pgfqpoint{7.388028in}{7.454961in}}%
\pgfpathquadraticcurveto{\pgfqpoint{7.383922in}{7.446391in}}{\pgfqpoint{7.381917in}{7.438776in}}%
\pgfpathquadraticcurveto{\pgfqpoint{7.379888in}{7.431161in}}{\pgfqpoint{7.379888in}{7.424071in}}%
\pgfpathquadraticcurveto{\pgfqpoint{7.379888in}{7.411777in}}{\pgfqpoint{7.384662in}{7.405093in}}%
\pgfpathquadraticcurveto{\pgfqpoint{7.389436in}{7.398409in}}{\pgfqpoint{7.398245in}{7.398409in}}%
\pgfpathquadraticcurveto{\pgfqpoint{7.405645in}{7.398409in}}{\pgfqpoint{7.409417in}{7.402849in}}%
\pgfpathquadraticcurveto{\pgfqpoint{7.413164in}{7.407289in}}{\pgfqpoint{7.415480in}{7.419679in}}%
\pgfpathlineto{\pgfqpoint{7.417342in}{7.428774in}}%
\pgfpathquadraticcurveto{\pgfqpoint{7.420565in}{7.445627in}}{\pgfqpoint{7.428657in}{7.453648in}}%
\pgfpathquadraticcurveto{\pgfqpoint{7.436750in}{7.461669in}}{\pgfqpoint{7.450309in}{7.461669in}}%
\pgfpathquadraticcurveto{\pgfqpoint{7.466517in}{7.461669in}}{\pgfqpoint{7.474872in}{7.450807in}}%
\pgfpathquadraticcurveto{\pgfqpoint{7.483227in}{7.439970in}}{\pgfqpoint{7.483227in}{7.419010in}}%
\pgfpathquadraticcurveto{\pgfqpoint{7.483227in}{7.411109in}}{\pgfqpoint{7.481437in}{7.402181in}}%
\pgfpathquadraticcurveto{\pgfqpoint{7.479647in}{7.393277in}}{\pgfqpoint{7.476138in}{7.383728in}}%
\pgfpathlineto{\pgfqpoint{7.460621in}{7.383728in}}%
\pgfpathquadraticcurveto{\pgfqpoint{7.465753in}{7.392895in}}{\pgfqpoint{7.468379in}{7.401704in}}%
\pgfpathquadraticcurveto{\pgfqpoint{7.470981in}{7.410512in}}{\pgfqpoint{7.470981in}{7.419010in}}%
\pgfpathquadraticcurveto{\pgfqpoint{7.470981in}{7.431901in}}{\pgfqpoint{7.465921in}{7.438919in}}%
\pgfpathquadraticcurveto{\pgfqpoint{7.460836in}{7.445938in}}{\pgfqpoint{7.451431in}{7.445938in}}%
\pgfpathquadraticcurveto{\pgfqpoint{7.443243in}{7.445938in}}{\pgfqpoint{7.438612in}{7.440901in}}%
\pgfpathquadraticcurveto{\pgfqpoint{7.433980in}{7.435864in}}{\pgfqpoint{7.431665in}{7.424382in}}%
\pgfpathlineto{\pgfqpoint{7.429875in}{7.415191in}}%
\pgfpathquadraticcurveto{\pgfqpoint{7.426533in}{7.398338in}}{\pgfqpoint{7.419371in}{7.390794in}}%
\pgfpathquadraticcurveto{\pgfqpoint{7.412210in}{7.383275in}}{\pgfqpoint{7.399438in}{7.383275in}}%
\pgfpathquadraticcurveto{\pgfqpoint{7.384662in}{7.383275in}}{\pgfqpoint{7.376164in}{7.393683in}}%
\pgfpathquadraticcurveto{\pgfqpoint{7.367665in}{7.404091in}}{\pgfqpoint{7.367665in}{7.422352in}}%
\pgfpathquadraticcurveto{\pgfqpoint{7.367665in}{7.430206in}}{\pgfqpoint{7.369098in}{7.438322in}}%
\pgfpathquadraticcurveto{\pgfqpoint{7.370506in}{7.446463in}}{\pgfqpoint{7.373323in}{7.454961in}}%
\pgfpathlineto{\pgfqpoint{7.373323in}{7.454961in}}%
\pgfpathclose%
\pgfpathmoveto{\pgfqpoint{7.369671in}{7.485182in}}%
\pgfpathlineto{\pgfqpoint{7.369671in}{7.555603in}}%
\pgfpathlineto{\pgfqpoint{7.382370in}{7.555603in}}%
\pgfpathlineto{\pgfqpoint{7.382370in}{7.500245in}}%
\pgfpathlineto{\pgfqpoint{7.415337in}{7.500245in}}%
\pgfpathlineto{\pgfqpoint{7.415337in}{7.553288in}}%
\pgfpathlineto{\pgfqpoint{7.428013in}{7.553288in}}%
\pgfpathlineto{\pgfqpoint{7.428013in}{7.500245in}}%
\pgfpathlineto{\pgfqpoint{7.468379in}{7.500245in}}%
\pgfpathlineto{\pgfqpoint{7.468379in}{7.556940in}}%
\pgfpathlineto{\pgfqpoint{7.481055in}{7.556940in}}%
\pgfpathlineto{\pgfqpoint{7.481055in}{7.485182in}}%
\pgfpathlineto{\pgfqpoint{7.369671in}{7.485182in}}%
\pgfpathlineto{\pgfqpoint{7.369671in}{7.485182in}}%
\pgfpathclose%
\pgfpathmoveto{\pgfqpoint{7.399964in}{7.634379in}}%
\pgfpathlineto{\pgfqpoint{7.412950in}{7.634379in}}%
\pgfpathquadraticcurveto{\pgfqpoint{7.409966in}{7.628579in}}{\pgfqpoint{7.408486in}{7.622300in}}%
\pgfpathquadraticcurveto{\pgfqpoint{7.406982in}{7.616046in}}{\pgfqpoint{7.406982in}{7.609314in}}%
\pgfpathquadraticcurveto{\pgfqpoint{7.406982in}{7.599097in}}{\pgfqpoint{7.410109in}{7.593989in}}%
\pgfpathquadraticcurveto{\pgfqpoint{7.413236in}{7.588880in}}{\pgfqpoint{7.419514in}{7.588880in}}%
\pgfpathquadraticcurveto{\pgfqpoint{7.424289in}{7.588880in}}{\pgfqpoint{7.427010in}{7.592533in}}%
\pgfpathquadraticcurveto{\pgfqpoint{7.429731in}{7.596185in}}{\pgfqpoint{7.432190in}{7.607237in}}%
\pgfpathlineto{\pgfqpoint{7.433240in}{7.611940in}}%
\pgfpathquadraticcurveto{\pgfqpoint{7.436368in}{7.626549in}}{\pgfqpoint{7.442073in}{7.632708in}}%
\pgfpathquadraticcurveto{\pgfqpoint{7.447778in}{7.638867in}}{\pgfqpoint{7.457995in}{7.638867in}}%
\pgfpathquadraticcurveto{\pgfqpoint{7.469645in}{7.638867in}}{\pgfqpoint{7.476448in}{7.629653in}}%
\pgfpathquadraticcurveto{\pgfqpoint{7.483227in}{7.620438in}}{\pgfqpoint{7.483227in}{7.604325in}}%
\pgfpathquadraticcurveto{\pgfqpoint{7.483227in}{7.597617in}}{\pgfqpoint{7.481914in}{7.590336in}}%
\pgfpathquadraticcurveto{\pgfqpoint{7.480602in}{7.583056in}}{\pgfqpoint{7.478000in}{7.575011in}}%
\pgfpathlineto{\pgfqpoint{7.463820in}{7.575011in}}%
\pgfpathquadraticcurveto{\pgfqpoint{7.467783in}{7.582626in}}{\pgfqpoint{7.469764in}{7.590002in}}%
\pgfpathquadraticcurveto{\pgfqpoint{7.471721in}{7.597378in}}{\pgfqpoint{7.471721in}{7.604635in}}%
\pgfpathquadraticcurveto{\pgfqpoint{7.471721in}{7.614327in}}{\pgfqpoint{7.468403in}{7.619531in}}%
\pgfpathquadraticcurveto{\pgfqpoint{7.465085in}{7.624759in}}{\pgfqpoint{7.459046in}{7.624759in}}%
\pgfpathquadraticcurveto{\pgfqpoint{7.453460in}{7.624759in}}{\pgfqpoint{7.450476in}{7.620987in}}%
\pgfpathquadraticcurveto{\pgfqpoint{7.447492in}{7.617240in}}{\pgfqpoint{7.444723in}{7.604468in}}%
\pgfpathlineto{\pgfqpoint{7.443601in}{7.599694in}}%
\pgfpathquadraticcurveto{\pgfqpoint{7.440927in}{7.586947in}}{\pgfqpoint{7.435365in}{7.581265in}}%
\pgfpathquadraticcurveto{\pgfqpoint{7.429803in}{7.575608in}}{\pgfqpoint{7.420111in}{7.575608in}}%
\pgfpathquadraticcurveto{\pgfqpoint{7.408319in}{7.575608in}}{\pgfqpoint{7.401921in}{7.583963in}}%
\pgfpathquadraticcurveto{\pgfqpoint{7.395500in}{7.592318in}}{\pgfqpoint{7.395500in}{7.607691in}}%
\pgfpathquadraticcurveto{\pgfqpoint{7.395500in}{7.615282in}}{\pgfqpoint{7.396622in}{7.621990in}}%
\pgfpathquadraticcurveto{\pgfqpoint{7.397720in}{7.628722in}}{\pgfqpoint{7.399964in}{7.634379in}}%
\pgfpathlineto{\pgfqpoint{7.399964in}{7.634379in}}%
\pgfpathclose%
\pgfpathmoveto{\pgfqpoint{7.365135in}{7.658561in}}%
\pgfpathlineto{\pgfqpoint{7.365135in}{7.670497in}}%
\pgfpathquadraticcurveto{\pgfqpoint{7.382728in}{7.681669in}}{\pgfqpoint{7.399605in}{7.687231in}}%
\pgfpathquadraticcurveto{\pgfqpoint{7.416459in}{7.692793in}}{\pgfqpoint{7.433097in}{7.692793in}}%
\pgfpathquadraticcurveto{\pgfqpoint{7.449807in}{7.692793in}}{\pgfqpoint{7.466732in}{7.687231in}}%
\pgfpathquadraticcurveto{\pgfqpoint{7.483657in}{7.681669in}}{\pgfqpoint{7.501203in}{7.670497in}}%
\pgfpathlineto{\pgfqpoint{7.501203in}{7.658561in}}%
\pgfpathquadraticcurveto{\pgfqpoint{7.484111in}{7.668468in}}{\pgfqpoint{7.467210in}{7.673362in}}%
\pgfpathquadraticcurveto{\pgfqpoint{7.450309in}{7.678255in}}{\pgfqpoint{7.433097in}{7.678255in}}%
\pgfpathquadraticcurveto{\pgfqpoint{7.415862in}{7.678255in}}{\pgfqpoint{7.399080in}{7.673362in}}%
\pgfpathquadraticcurveto{\pgfqpoint{7.382275in}{7.668468in}}{\pgfqpoint{7.365135in}{7.658561in}}%
\pgfpathlineto{\pgfqpoint{7.365135in}{7.658561in}}%
\pgfpathclose%
\pgfusepath{fill}%
\end{pgfscope}%
\begin{pgfscope}%
\pgfpathrectangle{\pgfqpoint{8.036091in}{4.050000in}}{\pgfqpoint{5.400000in}{5.400000in}}%
\pgfusepath{clip}%
\pgfsetrectcap%
\pgfsetroundjoin%
\pgfsetlinewidth{0.853187pt}%
\definecolor{currentstroke}{rgb}{0.745098,0.776471,0.803922}%
\pgfsetstrokecolor{currentstroke}%
\pgfsetdash{}{0pt}%
\pgfpathmoveto{\pgfqpoint{8.036091in}{6.750000in}}%
\pgfpathlineto{\pgfqpoint{13.436091in}{6.750000in}}%
\pgfusepath{stroke}%
\end{pgfscope}%
\begin{pgfscope}%
\pgfpathrectangle{\pgfqpoint{8.036091in}{4.050000in}}{\pgfqpoint{5.400000in}{5.400000in}}%
\pgfusepath{clip}%
\pgfsetrectcap%
\pgfsetroundjoin%
\pgfsetlinewidth{0.853187pt}%
\definecolor{currentstroke}{rgb}{0.745098,0.776471,0.803922}%
\pgfsetstrokecolor{currentstroke}%
\pgfsetdash{}{0pt}%
\pgfpathmoveto{\pgfqpoint{10.736091in}{4.050000in}}%
\pgfpathlineto{\pgfqpoint{10.736091in}{9.450000in}}%
\pgfusepath{stroke}%
\end{pgfscope}%
\begin{pgfscope}%
\pgfpathrectangle{\pgfqpoint{8.036091in}{4.050000in}}{\pgfqpoint{5.400000in}{5.400000in}}%
\pgfusepath{clip}%
\pgfsetbuttcap%
\pgfsetroundjoin%
\pgfsetlinewidth{1.104125pt}%
\definecolor{currentstroke}{rgb}{0.203922,0.243137,0.278431}%
\pgfsetstrokecolor{currentstroke}%
\pgfsetdash{{5.500000pt}{4.400000pt}}{0.000000pt}%
\pgfpathmoveto{\pgfqpoint{8.036091in}{4.050000in}}%
\pgfpathlineto{\pgfqpoint{13.436091in}{9.450000in}}%
\pgfusepath{stroke}%
\end{pgfscope}%
\begin{pgfscope}%
\pgfsetrectcap%
\pgfsetmiterjoin%
\pgfsetlinewidth{0.803000pt}%
\definecolor{currentstroke}{rgb}{0.627451,0.670588,0.705882}%
\pgfsetstrokecolor{currentstroke}%
\pgfsetdash{}{0pt}%
\pgfpathmoveto{\pgfqpoint{8.036091in}{4.050000in}}%
\pgfpathlineto{\pgfqpoint{8.036091in}{9.450000in}}%
\pgfusepath{stroke}%
\end{pgfscope}%
\begin{pgfscope}%
\pgfsetrectcap%
\pgfsetmiterjoin%
\pgfsetlinewidth{0.803000pt}%
\definecolor{currentstroke}{rgb}{0.627451,0.670588,0.705882}%
\pgfsetstrokecolor{currentstroke}%
\pgfsetdash{}{0pt}%
\pgfpathmoveto{\pgfqpoint{8.036091in}{4.050000in}}%
\pgfpathlineto{\pgfqpoint{13.436091in}{4.050000in}}%
\pgfusepath{stroke}%
\end{pgfscope}%
\begin{pgfscope}%
\pgfpathrectangle{\pgfqpoint{8.036091in}{4.050000in}}{\pgfqpoint{5.400000in}{5.400000in}}%
\pgfusepath{clip}%
\pgfsetbuttcap%
\pgfsetroundjoin%
\pgfsetlinewidth{1.003750pt}%
\definecolor{currentstroke}{rgb}{0.541176,0.380392,0.658824}%
\pgfsetstrokecolor{currentstroke}%
\pgfsetstrokeopacity{0.900000}%
\pgfsetdash{}{0pt}%
\pgfpathmoveto{\pgfqpoint{11.969733in}{8.797932in}}%
\pgfpathcurveto{\pgfqpoint{11.979987in}{8.797932in}}{\pgfqpoint{11.989822in}{8.802006in}}{\pgfqpoint{11.997073in}{8.809256in}}%
\pgfpathcurveto{\pgfqpoint{12.004324in}{8.816507in}}{\pgfqpoint{12.008398in}{8.826343in}}{\pgfqpoint{12.008398in}{8.836597in}}%
\pgfpathcurveto{\pgfqpoint{12.008398in}{8.846851in}}{\pgfqpoint{12.004324in}{8.856686in}}{\pgfqpoint{11.997073in}{8.863937in}}%
\pgfpathcurveto{\pgfqpoint{11.989822in}{8.871188in}}{\pgfqpoint{11.979987in}{8.875262in}}{\pgfqpoint{11.969733in}{8.875262in}}%
\pgfpathcurveto{\pgfqpoint{11.959478in}{8.875262in}}{\pgfqpoint{11.949643in}{8.871188in}}{\pgfqpoint{11.942392in}{8.863937in}}%
\pgfpathcurveto{\pgfqpoint{11.935142in}{8.856686in}}{\pgfqpoint{11.931068in}{8.846851in}}{\pgfqpoint{11.931068in}{8.836597in}}%
\pgfpathcurveto{\pgfqpoint{11.931068in}{8.826343in}}{\pgfqpoint{11.935142in}{8.816507in}}{\pgfqpoint{11.942392in}{8.809256in}}%
\pgfpathcurveto{\pgfqpoint{11.949643in}{8.802006in}}{\pgfqpoint{11.959478in}{8.797932in}}{\pgfqpoint{11.969733in}{8.797932in}}%
\pgfpathlineto{\pgfqpoint{11.969733in}{8.797932in}}%
\pgfpathclose%
\pgfusepath{stroke}%
\end{pgfscope}%
\begin{pgfscope}%
\pgfpathrectangle{\pgfqpoint{8.036091in}{4.050000in}}{\pgfqpoint{5.400000in}{5.400000in}}%
\pgfusepath{clip}%
\pgfsetbuttcap%
\pgfsetroundjoin%
\pgfsetlinewidth{1.003750pt}%
\definecolor{currentstroke}{rgb}{0.541176,0.380392,0.658824}%
\pgfsetstrokecolor{currentstroke}%
\pgfsetstrokeopacity{0.900000}%
\pgfsetdash{}{0pt}%
\pgfpathmoveto{\pgfqpoint{11.629004in}{8.344831in}}%
\pgfpathcurveto{\pgfqpoint{11.639259in}{8.344831in}}{\pgfqpoint{11.649094in}{8.348905in}}{\pgfqpoint{11.656345in}{8.356156in}}%
\pgfpathcurveto{\pgfqpoint{11.663596in}{8.363407in}}{\pgfqpoint{11.667670in}{8.373242in}}{\pgfqpoint{11.667670in}{8.383496in}}%
\pgfpathcurveto{\pgfqpoint{11.667670in}{8.393750in}}{\pgfqpoint{11.663596in}{8.403586in}}{\pgfqpoint{11.656345in}{8.410837in}}%
\pgfpathcurveto{\pgfqpoint{11.649094in}{8.418087in}}{\pgfqpoint{11.639259in}{8.422161in}}{\pgfqpoint{11.629004in}{8.422161in}}%
\pgfpathcurveto{\pgfqpoint{11.618750in}{8.422161in}}{\pgfqpoint{11.608915in}{8.418087in}}{\pgfqpoint{11.601664in}{8.410837in}}%
\pgfpathcurveto{\pgfqpoint{11.594413in}{8.403586in}}{\pgfqpoint{11.590339in}{8.393750in}}{\pgfqpoint{11.590339in}{8.383496in}}%
\pgfpathcurveto{\pgfqpoint{11.590339in}{8.373242in}}{\pgfqpoint{11.594413in}{8.363407in}}{\pgfqpoint{11.601664in}{8.356156in}}%
\pgfpathcurveto{\pgfqpoint{11.608915in}{8.348905in}}{\pgfqpoint{11.618750in}{8.344831in}}{\pgfqpoint{11.629004in}{8.344831in}}%
\pgfpathlineto{\pgfqpoint{11.629004in}{8.344831in}}%
\pgfpathclose%
\pgfusepath{stroke}%
\end{pgfscope}%
\begin{pgfscope}%
\pgfpathrectangle{\pgfqpoint{8.036091in}{4.050000in}}{\pgfqpoint{5.400000in}{5.400000in}}%
\pgfusepath{clip}%
\pgfsetbuttcap%
\pgfsetroundjoin%
\pgfsetlinewidth{1.003750pt}%
\definecolor{currentstroke}{rgb}{0.541176,0.380392,0.658824}%
\pgfsetstrokecolor{currentstroke}%
\pgfsetstrokeopacity{0.900000}%
\pgfsetdash{}{0pt}%
\pgfpathmoveto{\pgfqpoint{12.516094in}{8.815966in}}%
\pgfpathcurveto{\pgfqpoint{12.526348in}{8.815966in}}{\pgfqpoint{12.536183in}{8.820040in}}{\pgfqpoint{12.543434in}{8.827291in}}%
\pgfpathcurveto{\pgfqpoint{12.550685in}{8.834542in}}{\pgfqpoint{12.554759in}{8.844377in}}{\pgfqpoint{12.554759in}{8.854631in}}%
\pgfpathcurveto{\pgfqpoint{12.554759in}{8.864885in}}{\pgfqpoint{12.550685in}{8.874721in}}{\pgfqpoint{12.543434in}{8.881972in}}%
\pgfpathcurveto{\pgfqpoint{12.536183in}{8.889222in}}{\pgfqpoint{12.526348in}{8.893296in}}{\pgfqpoint{12.516094in}{8.893296in}}%
\pgfpathcurveto{\pgfqpoint{12.505839in}{8.893296in}}{\pgfqpoint{12.496004in}{8.889222in}}{\pgfqpoint{12.488753in}{8.881972in}}%
\pgfpathcurveto{\pgfqpoint{12.481503in}{8.874721in}}{\pgfqpoint{12.477429in}{8.864885in}}{\pgfqpoint{12.477429in}{8.854631in}}%
\pgfpathcurveto{\pgfqpoint{12.477429in}{8.844377in}}{\pgfqpoint{12.481503in}{8.834542in}}{\pgfqpoint{12.488753in}{8.827291in}}%
\pgfpathcurveto{\pgfqpoint{12.496004in}{8.820040in}}{\pgfqpoint{12.505839in}{8.815966in}}{\pgfqpoint{12.516094in}{8.815966in}}%
\pgfpathlineto{\pgfqpoint{12.516094in}{8.815966in}}%
\pgfpathclose%
\pgfusepath{stroke}%
\end{pgfscope}%
\begin{pgfscope}%
\pgfpathrectangle{\pgfqpoint{8.036091in}{4.050000in}}{\pgfqpoint{5.400000in}{5.400000in}}%
\pgfusepath{clip}%
\pgfsetbuttcap%
\pgfsetroundjoin%
\pgfsetlinewidth{1.003750pt}%
\definecolor{currentstroke}{rgb}{0.541176,0.380392,0.658824}%
\pgfsetstrokecolor{currentstroke}%
\pgfsetstrokeopacity{0.900000}%
\pgfsetdash{}{0pt}%
\pgfpathmoveto{\pgfqpoint{11.453294in}{7.604987in}}%
\pgfpathcurveto{\pgfqpoint{11.463548in}{7.604987in}}{\pgfqpoint{11.473383in}{7.609061in}}{\pgfqpoint{11.480634in}{7.616311in}}%
\pgfpathcurveto{\pgfqpoint{11.487885in}{7.623562in}}{\pgfqpoint{11.491959in}{7.633397in}}{\pgfqpoint{11.491959in}{7.643652in}}%
\pgfpathcurveto{\pgfqpoint{11.491959in}{7.653906in}}{\pgfqpoint{11.487885in}{7.663741in}}{\pgfqpoint{11.480634in}{7.670992in}}%
\pgfpathcurveto{\pgfqpoint{11.473383in}{7.678243in}}{\pgfqpoint{11.463548in}{7.682317in}}{\pgfqpoint{11.453294in}{7.682317in}}%
\pgfpathcurveto{\pgfqpoint{11.443040in}{7.682317in}}{\pgfqpoint{11.433204in}{7.678243in}}{\pgfqpoint{11.425953in}{7.670992in}}%
\pgfpathcurveto{\pgfqpoint{11.418703in}{7.663741in}}{\pgfqpoint{11.414629in}{7.653906in}}{\pgfqpoint{11.414629in}{7.643652in}}%
\pgfpathcurveto{\pgfqpoint{11.414629in}{7.633397in}}{\pgfqpoint{11.418703in}{7.623562in}}{\pgfqpoint{11.425953in}{7.616311in}}%
\pgfpathcurveto{\pgfqpoint{11.433204in}{7.609061in}}{\pgfqpoint{11.443040in}{7.604987in}}{\pgfqpoint{11.453294in}{7.604987in}}%
\pgfpathlineto{\pgfqpoint{11.453294in}{7.604987in}}%
\pgfpathclose%
\pgfusepath{stroke}%
\end{pgfscope}%
\begin{pgfscope}%
\pgfpathrectangle{\pgfqpoint{8.036091in}{4.050000in}}{\pgfqpoint{5.400000in}{5.400000in}}%
\pgfusepath{clip}%
\pgfsetbuttcap%
\pgfsetroundjoin%
\pgfsetlinewidth{1.003750pt}%
\definecolor{currentstroke}{rgb}{0.541176,0.380392,0.658824}%
\pgfsetstrokecolor{currentstroke}%
\pgfsetstrokeopacity{0.900000}%
\pgfsetdash{}{0pt}%
\pgfpathmoveto{\pgfqpoint{12.114614in}{8.953262in}}%
\pgfpathcurveto{\pgfqpoint{12.124868in}{8.953262in}}{\pgfqpoint{12.134704in}{8.957336in}}{\pgfqpoint{12.141954in}{8.964587in}}%
\pgfpathcurveto{\pgfqpoint{12.149205in}{8.971837in}}{\pgfqpoint{12.153279in}{8.981673in}}{\pgfqpoint{12.153279in}{8.991927in}}%
\pgfpathcurveto{\pgfqpoint{12.153279in}{9.002181in}}{\pgfqpoint{12.149205in}{9.012017in}}{\pgfqpoint{12.141954in}{9.019267in}}%
\pgfpathcurveto{\pgfqpoint{12.134704in}{9.026518in}}{\pgfqpoint{12.124868in}{9.030592in}}{\pgfqpoint{12.114614in}{9.030592in}}%
\pgfpathcurveto{\pgfqpoint{12.104360in}{9.030592in}}{\pgfqpoint{12.094525in}{9.026518in}}{\pgfqpoint{12.087274in}{9.019267in}}%
\pgfpathcurveto{\pgfqpoint{12.080023in}{9.012017in}}{\pgfqpoint{12.075949in}{9.002181in}}{\pgfqpoint{12.075949in}{8.991927in}}%
\pgfpathcurveto{\pgfqpoint{12.075949in}{8.981673in}}{\pgfqpoint{12.080023in}{8.971837in}}{\pgfqpoint{12.087274in}{8.964587in}}%
\pgfpathcurveto{\pgfqpoint{12.094525in}{8.957336in}}{\pgfqpoint{12.104360in}{8.953262in}}{\pgfqpoint{12.114614in}{8.953262in}}%
\pgfpathlineto{\pgfqpoint{12.114614in}{8.953262in}}%
\pgfpathclose%
\pgfusepath{stroke}%
\end{pgfscope}%
\begin{pgfscope}%
\pgfpathrectangle{\pgfqpoint{8.036091in}{4.050000in}}{\pgfqpoint{5.400000in}{5.400000in}}%
\pgfusepath{clip}%
\pgfsetbuttcap%
\pgfsetroundjoin%
\pgfsetlinewidth{1.003750pt}%
\definecolor{currentstroke}{rgb}{0.541176,0.380392,0.658824}%
\pgfsetstrokecolor{currentstroke}%
\pgfsetstrokeopacity{0.900000}%
\pgfsetdash{}{0pt}%
\pgfpathmoveto{\pgfqpoint{11.658001in}{7.468230in}}%
\pgfpathcurveto{\pgfqpoint{11.668255in}{7.468230in}}{\pgfqpoint{11.678091in}{7.472304in}}{\pgfqpoint{11.685342in}{7.479554in}}%
\pgfpathcurveto{\pgfqpoint{11.692592in}{7.486805in}}{\pgfqpoint{11.696666in}{7.496641in}}{\pgfqpoint{11.696666in}{7.506895in}}%
\pgfpathcurveto{\pgfqpoint{11.696666in}{7.517149in}}{\pgfqpoint{11.692592in}{7.526984in}}{\pgfqpoint{11.685342in}{7.534235in}}%
\pgfpathcurveto{\pgfqpoint{11.678091in}{7.541486in}}{\pgfqpoint{11.668255in}{7.545560in}}{\pgfqpoint{11.658001in}{7.545560in}}%
\pgfpathcurveto{\pgfqpoint{11.647747in}{7.545560in}}{\pgfqpoint{11.637912in}{7.541486in}}{\pgfqpoint{11.630661in}{7.534235in}}%
\pgfpathcurveto{\pgfqpoint{11.623410in}{7.526984in}}{\pgfqpoint{11.619336in}{7.517149in}}{\pgfqpoint{11.619336in}{7.506895in}}%
\pgfpathcurveto{\pgfqpoint{11.619336in}{7.496641in}}{\pgfqpoint{11.623410in}{7.486805in}}{\pgfqpoint{11.630661in}{7.479554in}}%
\pgfpathcurveto{\pgfqpoint{11.637912in}{7.472304in}}{\pgfqpoint{11.647747in}{7.468230in}}{\pgfqpoint{11.658001in}{7.468230in}}%
\pgfpathlineto{\pgfqpoint{11.658001in}{7.468230in}}%
\pgfpathclose%
\pgfusepath{stroke}%
\end{pgfscope}%
\begin{pgfscope}%
\pgfpathrectangle{\pgfqpoint{8.036091in}{4.050000in}}{\pgfqpoint{5.400000in}{5.400000in}}%
\pgfusepath{clip}%
\pgfsetbuttcap%
\pgfsetroundjoin%
\pgfsetlinewidth{1.003750pt}%
\definecolor{currentstroke}{rgb}{0.541176,0.380392,0.658824}%
\pgfsetstrokecolor{currentstroke}%
\pgfsetstrokeopacity{0.900000}%
\pgfsetdash{}{0pt}%
\pgfpathmoveto{\pgfqpoint{12.201419in}{8.641661in}}%
\pgfpathcurveto{\pgfqpoint{12.211673in}{8.641661in}}{\pgfqpoint{12.221509in}{8.645735in}}{\pgfqpoint{12.228759in}{8.652986in}}%
\pgfpathcurveto{\pgfqpoint{12.236010in}{8.660236in}}{\pgfqpoint{12.240084in}{8.670072in}}{\pgfqpoint{12.240084in}{8.680326in}}%
\pgfpathcurveto{\pgfqpoint{12.240084in}{8.690580in}}{\pgfqpoint{12.236010in}{8.700416in}}{\pgfqpoint{12.228759in}{8.707666in}}%
\pgfpathcurveto{\pgfqpoint{12.221509in}{8.714917in}}{\pgfqpoint{12.211673in}{8.718991in}}{\pgfqpoint{12.201419in}{8.718991in}}%
\pgfpathcurveto{\pgfqpoint{12.191165in}{8.718991in}}{\pgfqpoint{12.181330in}{8.714917in}}{\pgfqpoint{12.174079in}{8.707666in}}%
\pgfpathcurveto{\pgfqpoint{12.166828in}{8.700416in}}{\pgfqpoint{12.162754in}{8.690580in}}{\pgfqpoint{12.162754in}{8.680326in}}%
\pgfpathcurveto{\pgfqpoint{12.162754in}{8.670072in}}{\pgfqpoint{12.166828in}{8.660236in}}{\pgfqpoint{12.174079in}{8.652986in}}%
\pgfpathcurveto{\pgfqpoint{12.181330in}{8.645735in}}{\pgfqpoint{12.191165in}{8.641661in}}{\pgfqpoint{12.201419in}{8.641661in}}%
\pgfpathlineto{\pgfqpoint{12.201419in}{8.641661in}}%
\pgfpathclose%
\pgfusepath{stroke}%
\end{pgfscope}%
\begin{pgfscope}%
\pgfpathrectangle{\pgfqpoint{8.036091in}{4.050000in}}{\pgfqpoint{5.400000in}{5.400000in}}%
\pgfusepath{clip}%
\pgfsetbuttcap%
\pgfsetroundjoin%
\pgfsetlinewidth{1.003750pt}%
\definecolor{currentstroke}{rgb}{0.541176,0.380392,0.658824}%
\pgfsetstrokecolor{currentstroke}%
\pgfsetstrokeopacity{0.900000}%
\pgfsetdash{}{0pt}%
\pgfpathmoveto{\pgfqpoint{10.730668in}{7.135660in}}%
\pgfpathcurveto{\pgfqpoint{10.740922in}{7.135660in}}{\pgfqpoint{10.750757in}{7.139734in}}{\pgfqpoint{10.758008in}{7.146985in}}%
\pgfpathcurveto{\pgfqpoint{10.765259in}{7.154236in}}{\pgfqpoint{10.769333in}{7.164071in}}{\pgfqpoint{10.769333in}{7.174325in}}%
\pgfpathcurveto{\pgfqpoint{10.769333in}{7.184579in}}{\pgfqpoint{10.765259in}{7.194415in}}{\pgfqpoint{10.758008in}{7.201665in}}%
\pgfpathcurveto{\pgfqpoint{10.750757in}{7.208916in}}{\pgfqpoint{10.740922in}{7.212990in}}{\pgfqpoint{10.730668in}{7.212990in}}%
\pgfpathcurveto{\pgfqpoint{10.720414in}{7.212990in}}{\pgfqpoint{10.710578in}{7.208916in}}{\pgfqpoint{10.703327in}{7.201665in}}%
\pgfpathcurveto{\pgfqpoint{10.696077in}{7.194415in}}{\pgfqpoint{10.692003in}{7.184579in}}{\pgfqpoint{10.692003in}{7.174325in}}%
\pgfpathcurveto{\pgfqpoint{10.692003in}{7.164071in}}{\pgfqpoint{10.696077in}{7.154236in}}{\pgfqpoint{10.703327in}{7.146985in}}%
\pgfpathcurveto{\pgfqpoint{10.710578in}{7.139734in}}{\pgfqpoint{10.720414in}{7.135660in}}{\pgfqpoint{10.730668in}{7.135660in}}%
\pgfpathlineto{\pgfqpoint{10.730668in}{7.135660in}}%
\pgfpathclose%
\pgfusepath{stroke}%
\end{pgfscope}%
\begin{pgfscope}%
\pgfpathrectangle{\pgfqpoint{8.036091in}{4.050000in}}{\pgfqpoint{5.400000in}{5.400000in}}%
\pgfusepath{clip}%
\pgfsetbuttcap%
\pgfsetroundjoin%
\pgfsetlinewidth{1.003750pt}%
\definecolor{currentstroke}{rgb}{0.541176,0.380392,0.658824}%
\pgfsetstrokecolor{currentstroke}%
\pgfsetstrokeopacity{0.900000}%
\pgfsetdash{}{0pt}%
\pgfpathmoveto{\pgfqpoint{11.541024in}{6.368898in}}%
\pgfpathcurveto{\pgfqpoint{11.551278in}{6.368898in}}{\pgfqpoint{11.561113in}{6.372972in}}{\pgfqpoint{11.568364in}{6.380223in}}%
\pgfpathcurveto{\pgfqpoint{11.575615in}{6.387474in}}{\pgfqpoint{11.579689in}{6.397309in}}{\pgfqpoint{11.579689in}{6.407563in}}%
\pgfpathcurveto{\pgfqpoint{11.579689in}{6.417817in}}{\pgfqpoint{11.575615in}{6.427653in}}{\pgfqpoint{11.568364in}{6.434904in}}%
\pgfpathcurveto{\pgfqpoint{11.561113in}{6.442154in}}{\pgfqpoint{11.551278in}{6.446228in}}{\pgfqpoint{11.541024in}{6.446228in}}%
\pgfpathcurveto{\pgfqpoint{11.530770in}{6.446228in}}{\pgfqpoint{11.520934in}{6.442154in}}{\pgfqpoint{11.513683in}{6.434904in}}%
\pgfpathcurveto{\pgfqpoint{11.506433in}{6.427653in}}{\pgfqpoint{11.502359in}{6.417817in}}{\pgfqpoint{11.502359in}{6.407563in}}%
\pgfpathcurveto{\pgfqpoint{11.502359in}{6.397309in}}{\pgfqpoint{11.506433in}{6.387474in}}{\pgfqpoint{11.513683in}{6.380223in}}%
\pgfpathcurveto{\pgfqpoint{11.520934in}{6.372972in}}{\pgfqpoint{11.530770in}{6.368898in}}{\pgfqpoint{11.541024in}{6.368898in}}%
\pgfpathlineto{\pgfqpoint{11.541024in}{6.368898in}}%
\pgfpathclose%
\pgfusepath{stroke}%
\end{pgfscope}%
\begin{pgfscope}%
\pgfpathrectangle{\pgfqpoint{8.036091in}{4.050000in}}{\pgfqpoint{5.400000in}{5.400000in}}%
\pgfusepath{clip}%
\pgfsetbuttcap%
\pgfsetroundjoin%
\pgfsetlinewidth{1.003750pt}%
\definecolor{currentstroke}{rgb}{0.541176,0.380392,0.658824}%
\pgfsetstrokecolor{currentstroke}%
\pgfsetstrokeopacity{0.900000}%
\pgfsetdash{}{0pt}%
\pgfpathmoveto{\pgfqpoint{12.126249in}{7.843425in}}%
\pgfpathcurveto{\pgfqpoint{12.136504in}{7.843425in}}{\pgfqpoint{12.146339in}{7.847499in}}{\pgfqpoint{12.153590in}{7.854749in}}%
\pgfpathcurveto{\pgfqpoint{12.160841in}{7.862000in}}{\pgfqpoint{12.164915in}{7.871836in}}{\pgfqpoint{12.164915in}{7.882090in}}%
\pgfpathcurveto{\pgfqpoint{12.164915in}{7.892344in}}{\pgfqpoint{12.160841in}{7.902179in}}{\pgfqpoint{12.153590in}{7.909430in}}%
\pgfpathcurveto{\pgfqpoint{12.146339in}{7.916681in}}{\pgfqpoint{12.136504in}{7.920755in}}{\pgfqpoint{12.126249in}{7.920755in}}%
\pgfpathcurveto{\pgfqpoint{12.115995in}{7.920755in}}{\pgfqpoint{12.106160in}{7.916681in}}{\pgfqpoint{12.098909in}{7.909430in}}%
\pgfpathcurveto{\pgfqpoint{12.091658in}{7.902179in}}{\pgfqpoint{12.087584in}{7.892344in}}{\pgfqpoint{12.087584in}{7.882090in}}%
\pgfpathcurveto{\pgfqpoint{12.087584in}{7.871836in}}{\pgfqpoint{12.091658in}{7.862000in}}{\pgfqpoint{12.098909in}{7.854749in}}%
\pgfpathcurveto{\pgfqpoint{12.106160in}{7.847499in}}{\pgfqpoint{12.115995in}{7.843425in}}{\pgfqpoint{12.126249in}{7.843425in}}%
\pgfpathlineto{\pgfqpoint{12.126249in}{7.843425in}}%
\pgfpathclose%
\pgfusepath{stroke}%
\end{pgfscope}%
\begin{pgfscope}%
\pgfpathrectangle{\pgfqpoint{8.036091in}{4.050000in}}{\pgfqpoint{5.400000in}{5.400000in}}%
\pgfusepath{clip}%
\pgfsetbuttcap%
\pgfsetroundjoin%
\pgfsetlinewidth{1.003750pt}%
\definecolor{currentstroke}{rgb}{0.541176,0.380392,0.658824}%
\pgfsetstrokecolor{currentstroke}%
\pgfsetstrokeopacity{0.900000}%
\pgfsetdash{}{0pt}%
\pgfpathmoveto{\pgfqpoint{11.546137in}{6.894818in}}%
\pgfpathcurveto{\pgfqpoint{11.556391in}{6.894818in}}{\pgfqpoint{11.566226in}{6.898892in}}{\pgfqpoint{11.573477in}{6.906143in}}%
\pgfpathcurveto{\pgfqpoint{11.580728in}{6.913394in}}{\pgfqpoint{11.584802in}{6.923229in}}{\pgfqpoint{11.584802in}{6.933483in}}%
\pgfpathcurveto{\pgfqpoint{11.584802in}{6.943738in}}{\pgfqpoint{11.580728in}{6.953573in}}{\pgfqpoint{11.573477in}{6.960824in}}%
\pgfpathcurveto{\pgfqpoint{11.566226in}{6.968075in}}{\pgfqpoint{11.556391in}{6.972149in}}{\pgfqpoint{11.546137in}{6.972149in}}%
\pgfpathcurveto{\pgfqpoint{11.535882in}{6.972149in}}{\pgfqpoint{11.526047in}{6.968075in}}{\pgfqpoint{11.518796in}{6.960824in}}%
\pgfpathcurveto{\pgfqpoint{11.511546in}{6.953573in}}{\pgfqpoint{11.507472in}{6.943738in}}{\pgfqpoint{11.507472in}{6.933483in}}%
\pgfpathcurveto{\pgfqpoint{11.507472in}{6.923229in}}{\pgfqpoint{11.511546in}{6.913394in}}{\pgfqpoint{11.518796in}{6.906143in}}%
\pgfpathcurveto{\pgfqpoint{11.526047in}{6.898892in}}{\pgfqpoint{11.535882in}{6.894818in}}{\pgfqpoint{11.546137in}{6.894818in}}%
\pgfpathlineto{\pgfqpoint{11.546137in}{6.894818in}}%
\pgfpathclose%
\pgfusepath{stroke}%
\end{pgfscope}%
\begin{pgfscope}%
\pgfpathrectangle{\pgfqpoint{8.036091in}{4.050000in}}{\pgfqpoint{5.400000in}{5.400000in}}%
\pgfusepath{clip}%
\pgfsetbuttcap%
\pgfsetroundjoin%
\pgfsetlinewidth{1.003750pt}%
\definecolor{currentstroke}{rgb}{0.541176,0.380392,0.658824}%
\pgfsetstrokecolor{currentstroke}%
\pgfsetstrokeopacity{0.900000}%
\pgfsetdash{}{0pt}%
\pgfpathmoveto{\pgfqpoint{11.813982in}{8.243192in}}%
\pgfpathcurveto{\pgfqpoint{11.824236in}{8.243192in}}{\pgfqpoint{11.834072in}{8.247266in}}{\pgfqpoint{11.841323in}{8.254516in}}%
\pgfpathcurveto{\pgfqpoint{11.848573in}{8.261767in}}{\pgfqpoint{11.852647in}{8.271603in}}{\pgfqpoint{11.852647in}{8.281857in}}%
\pgfpathcurveto{\pgfqpoint{11.852647in}{8.292111in}}{\pgfqpoint{11.848573in}{8.301946in}}{\pgfqpoint{11.841323in}{8.309197in}}%
\pgfpathcurveto{\pgfqpoint{11.834072in}{8.316448in}}{\pgfqpoint{11.824236in}{8.320522in}}{\pgfqpoint{11.813982in}{8.320522in}}%
\pgfpathcurveto{\pgfqpoint{11.803728in}{8.320522in}}{\pgfqpoint{11.793893in}{8.316448in}}{\pgfqpoint{11.786642in}{8.309197in}}%
\pgfpathcurveto{\pgfqpoint{11.779391in}{8.301946in}}{\pgfqpoint{11.775317in}{8.292111in}}{\pgfqpoint{11.775317in}{8.281857in}}%
\pgfpathcurveto{\pgfqpoint{11.775317in}{8.271603in}}{\pgfqpoint{11.779391in}{8.261767in}}{\pgfqpoint{11.786642in}{8.254516in}}%
\pgfpathcurveto{\pgfqpoint{11.793893in}{8.247266in}}{\pgfqpoint{11.803728in}{8.243192in}}{\pgfqpoint{11.813982in}{8.243192in}}%
\pgfpathlineto{\pgfqpoint{11.813982in}{8.243192in}}%
\pgfpathclose%
\pgfusepath{stroke}%
\end{pgfscope}%
\begin{pgfscope}%
\pgfpathrectangle{\pgfqpoint{8.036091in}{4.050000in}}{\pgfqpoint{5.400000in}{5.400000in}}%
\pgfusepath{clip}%
\pgfsetbuttcap%
\pgfsetroundjoin%
\pgfsetlinewidth{1.003750pt}%
\definecolor{currentstroke}{rgb}{0.541176,0.380392,0.658824}%
\pgfsetstrokecolor{currentstroke}%
\pgfsetstrokeopacity{0.900000}%
\pgfsetdash{}{0pt}%
\pgfpathmoveto{\pgfqpoint{11.601135in}{7.786150in}}%
\pgfpathcurveto{\pgfqpoint{11.611389in}{7.786150in}}{\pgfqpoint{11.621224in}{7.790224in}}{\pgfqpoint{11.628475in}{7.797475in}}%
\pgfpathcurveto{\pgfqpoint{11.635726in}{7.804725in}}{\pgfqpoint{11.639800in}{7.814561in}}{\pgfqpoint{11.639800in}{7.824815in}}%
\pgfpathcurveto{\pgfqpoint{11.639800in}{7.835069in}}{\pgfqpoint{11.635726in}{7.844905in}}{\pgfqpoint{11.628475in}{7.852155in}}%
\pgfpathcurveto{\pgfqpoint{11.621224in}{7.859406in}}{\pgfqpoint{11.611389in}{7.863480in}}{\pgfqpoint{11.601135in}{7.863480in}}%
\pgfpathcurveto{\pgfqpoint{11.590881in}{7.863480in}}{\pgfqpoint{11.581045in}{7.859406in}}{\pgfqpoint{11.573794in}{7.852155in}}%
\pgfpathcurveto{\pgfqpoint{11.566544in}{7.844905in}}{\pgfqpoint{11.562470in}{7.835069in}}{\pgfqpoint{11.562470in}{7.824815in}}%
\pgfpathcurveto{\pgfqpoint{11.562470in}{7.814561in}}{\pgfqpoint{11.566544in}{7.804725in}}{\pgfqpoint{11.573794in}{7.797475in}}%
\pgfpathcurveto{\pgfqpoint{11.581045in}{7.790224in}}{\pgfqpoint{11.590881in}{7.786150in}}{\pgfqpoint{11.601135in}{7.786150in}}%
\pgfpathlineto{\pgfqpoint{11.601135in}{7.786150in}}%
\pgfpathclose%
\pgfusepath{stroke}%
\end{pgfscope}%
\begin{pgfscope}%
\pgfpathrectangle{\pgfqpoint{8.036091in}{4.050000in}}{\pgfqpoint{5.400000in}{5.400000in}}%
\pgfusepath{clip}%
\pgfsetbuttcap%
\pgfsetroundjoin%
\pgfsetlinewidth{1.003750pt}%
\definecolor{currentstroke}{rgb}{0.541176,0.380392,0.658824}%
\pgfsetstrokecolor{currentstroke}%
\pgfsetstrokeopacity{0.900000}%
\pgfsetdash{}{0pt}%
\pgfpathmoveto{\pgfqpoint{11.630930in}{7.440174in}}%
\pgfpathcurveto{\pgfqpoint{11.641184in}{7.440174in}}{\pgfqpoint{11.651020in}{7.444248in}}{\pgfqpoint{11.658271in}{7.451499in}}%
\pgfpathcurveto{\pgfqpoint{11.665521in}{7.458750in}}{\pgfqpoint{11.669595in}{7.468585in}}{\pgfqpoint{11.669595in}{7.478839in}}%
\pgfpathcurveto{\pgfqpoint{11.669595in}{7.489093in}}{\pgfqpoint{11.665521in}{7.498929in}}{\pgfqpoint{11.658271in}{7.506180in}}%
\pgfpathcurveto{\pgfqpoint{11.651020in}{7.513430in}}{\pgfqpoint{11.641184in}{7.517504in}}{\pgfqpoint{11.630930in}{7.517504in}}%
\pgfpathcurveto{\pgfqpoint{11.620676in}{7.517504in}}{\pgfqpoint{11.610841in}{7.513430in}}{\pgfqpoint{11.603590in}{7.506180in}}%
\pgfpathcurveto{\pgfqpoint{11.596339in}{7.498929in}}{\pgfqpoint{11.592265in}{7.489093in}}{\pgfqpoint{11.592265in}{7.478839in}}%
\pgfpathcurveto{\pgfqpoint{11.592265in}{7.468585in}}{\pgfqpoint{11.596339in}{7.458750in}}{\pgfqpoint{11.603590in}{7.451499in}}%
\pgfpathcurveto{\pgfqpoint{11.610841in}{7.444248in}}{\pgfqpoint{11.620676in}{7.440174in}}{\pgfqpoint{11.630930in}{7.440174in}}%
\pgfpathlineto{\pgfqpoint{11.630930in}{7.440174in}}%
\pgfpathclose%
\pgfusepath{stroke}%
\end{pgfscope}%
\begin{pgfscope}%
\pgfpathrectangle{\pgfqpoint{8.036091in}{4.050000in}}{\pgfqpoint{5.400000in}{5.400000in}}%
\pgfusepath{clip}%
\pgfsetbuttcap%
\pgfsetroundjoin%
\pgfsetlinewidth{1.003750pt}%
\definecolor{currentstroke}{rgb}{0.541176,0.380392,0.658824}%
\pgfsetstrokecolor{currentstroke}%
\pgfsetstrokeopacity{0.900000}%
\pgfsetdash{}{0pt}%
\pgfpathmoveto{\pgfqpoint{10.715434in}{6.348694in}}%
\pgfpathcurveto{\pgfqpoint{10.725688in}{6.348694in}}{\pgfqpoint{10.735523in}{6.352768in}}{\pgfqpoint{10.742774in}{6.360019in}}%
\pgfpathcurveto{\pgfqpoint{10.750025in}{6.367270in}}{\pgfqpoint{10.754099in}{6.377105in}}{\pgfqpoint{10.754099in}{6.387359in}}%
\pgfpathcurveto{\pgfqpoint{10.754099in}{6.397613in}}{\pgfqpoint{10.750025in}{6.407449in}}{\pgfqpoint{10.742774in}{6.414699in}}%
\pgfpathcurveto{\pgfqpoint{10.735523in}{6.421950in}}{\pgfqpoint{10.725688in}{6.426024in}}{\pgfqpoint{10.715434in}{6.426024in}}%
\pgfpathcurveto{\pgfqpoint{10.705179in}{6.426024in}}{\pgfqpoint{10.695344in}{6.421950in}}{\pgfqpoint{10.688093in}{6.414699in}}%
\pgfpathcurveto{\pgfqpoint{10.680842in}{6.407449in}}{\pgfqpoint{10.676769in}{6.397613in}}{\pgfqpoint{10.676769in}{6.387359in}}%
\pgfpathcurveto{\pgfqpoint{10.676769in}{6.377105in}}{\pgfqpoint{10.680842in}{6.367270in}}{\pgfqpoint{10.688093in}{6.360019in}}%
\pgfpathcurveto{\pgfqpoint{10.695344in}{6.352768in}}{\pgfqpoint{10.705179in}{6.348694in}}{\pgfqpoint{10.715434in}{6.348694in}}%
\pgfpathlineto{\pgfqpoint{10.715434in}{6.348694in}}%
\pgfpathclose%
\pgfusepath{stroke}%
\end{pgfscope}%
\begin{pgfscope}%
\pgfpathrectangle{\pgfqpoint{8.036091in}{4.050000in}}{\pgfqpoint{5.400000in}{5.400000in}}%
\pgfusepath{clip}%
\pgfsetbuttcap%
\pgfsetroundjoin%
\pgfsetlinewidth{1.003750pt}%
\definecolor{currentstroke}{rgb}{0.541176,0.380392,0.658824}%
\pgfsetstrokecolor{currentstroke}%
\pgfsetstrokeopacity{0.900000}%
\pgfsetdash{}{0pt}%
\pgfpathmoveto{\pgfqpoint{10.981605in}{7.529936in}}%
\pgfpathcurveto{\pgfqpoint{10.991859in}{7.529936in}}{\pgfqpoint{11.001695in}{7.534010in}}{\pgfqpoint{11.008946in}{7.541261in}}%
\pgfpathcurveto{\pgfqpoint{11.016196in}{7.548512in}}{\pgfqpoint{11.020270in}{7.558347in}}{\pgfqpoint{11.020270in}{7.568601in}}%
\pgfpathcurveto{\pgfqpoint{11.020270in}{7.578855in}}{\pgfqpoint{11.016196in}{7.588691in}}{\pgfqpoint{11.008946in}{7.595942in}}%
\pgfpathcurveto{\pgfqpoint{11.001695in}{7.603192in}}{\pgfqpoint{10.991859in}{7.607266in}}{\pgfqpoint{10.981605in}{7.607266in}}%
\pgfpathcurveto{\pgfqpoint{10.971351in}{7.607266in}}{\pgfqpoint{10.961516in}{7.603192in}}{\pgfqpoint{10.954265in}{7.595942in}}%
\pgfpathcurveto{\pgfqpoint{10.947014in}{7.588691in}}{\pgfqpoint{10.942940in}{7.578855in}}{\pgfqpoint{10.942940in}{7.568601in}}%
\pgfpathcurveto{\pgfqpoint{10.942940in}{7.558347in}}{\pgfqpoint{10.947014in}{7.548512in}}{\pgfqpoint{10.954265in}{7.541261in}}%
\pgfpathcurveto{\pgfqpoint{10.961516in}{7.534010in}}{\pgfqpoint{10.971351in}{7.529936in}}{\pgfqpoint{10.981605in}{7.529936in}}%
\pgfpathlineto{\pgfqpoint{10.981605in}{7.529936in}}%
\pgfpathclose%
\pgfusepath{stroke}%
\end{pgfscope}%
\begin{pgfscope}%
\pgfpathrectangle{\pgfqpoint{8.036091in}{4.050000in}}{\pgfqpoint{5.400000in}{5.400000in}}%
\pgfusepath{clip}%
\pgfsetbuttcap%
\pgfsetroundjoin%
\pgfsetlinewidth{1.003750pt}%
\definecolor{currentstroke}{rgb}{0.541176,0.380392,0.658824}%
\pgfsetstrokecolor{currentstroke}%
\pgfsetstrokeopacity{0.900000}%
\pgfsetdash{}{0pt}%
\pgfpathmoveto{\pgfqpoint{11.506600in}{7.340659in}}%
\pgfpathcurveto{\pgfqpoint{11.516854in}{7.340659in}}{\pgfqpoint{11.526689in}{7.344733in}}{\pgfqpoint{11.533940in}{7.351984in}}%
\pgfpathcurveto{\pgfqpoint{11.541191in}{7.359235in}}{\pgfqpoint{11.545265in}{7.369070in}}{\pgfqpoint{11.545265in}{7.379324in}}%
\pgfpathcurveto{\pgfqpoint{11.545265in}{7.389578in}}{\pgfqpoint{11.541191in}{7.399414in}}{\pgfqpoint{11.533940in}{7.406665in}}%
\pgfpathcurveto{\pgfqpoint{11.526689in}{7.413915in}}{\pgfqpoint{11.516854in}{7.417989in}}{\pgfqpoint{11.506600in}{7.417989in}}%
\pgfpathcurveto{\pgfqpoint{11.496346in}{7.417989in}}{\pgfqpoint{11.486510in}{7.413915in}}{\pgfqpoint{11.479259in}{7.406665in}}%
\pgfpathcurveto{\pgfqpoint{11.472009in}{7.399414in}}{\pgfqpoint{11.467935in}{7.389578in}}{\pgfqpoint{11.467935in}{7.379324in}}%
\pgfpathcurveto{\pgfqpoint{11.467935in}{7.369070in}}{\pgfqpoint{11.472009in}{7.359235in}}{\pgfqpoint{11.479259in}{7.351984in}}%
\pgfpathcurveto{\pgfqpoint{11.486510in}{7.344733in}}{\pgfqpoint{11.496346in}{7.340659in}}{\pgfqpoint{11.506600in}{7.340659in}}%
\pgfpathlineto{\pgfqpoint{11.506600in}{7.340659in}}%
\pgfpathclose%
\pgfusepath{stroke}%
\end{pgfscope}%
\begin{pgfscope}%
\pgfpathrectangle{\pgfqpoint{8.036091in}{4.050000in}}{\pgfqpoint{5.400000in}{5.400000in}}%
\pgfusepath{clip}%
\pgfsetbuttcap%
\pgfsetroundjoin%
\pgfsetlinewidth{1.003750pt}%
\definecolor{currentstroke}{rgb}{0.541176,0.380392,0.658824}%
\pgfsetstrokecolor{currentstroke}%
\pgfsetstrokeopacity{0.900000}%
\pgfsetdash{}{0pt}%
\pgfpathmoveto{\pgfqpoint{10.676427in}{6.173686in}}%
\pgfpathcurveto{\pgfqpoint{10.686681in}{6.173686in}}{\pgfqpoint{10.696516in}{6.177760in}}{\pgfqpoint{10.703767in}{6.185010in}}%
\pgfpathcurveto{\pgfqpoint{10.711018in}{6.192261in}}{\pgfqpoint{10.715092in}{6.202097in}}{\pgfqpoint{10.715092in}{6.212351in}}%
\pgfpathcurveto{\pgfqpoint{10.715092in}{6.222605in}}{\pgfqpoint{10.711018in}{6.232440in}}{\pgfqpoint{10.703767in}{6.239691in}}%
\pgfpathcurveto{\pgfqpoint{10.696516in}{6.246942in}}{\pgfqpoint{10.686681in}{6.251016in}}{\pgfqpoint{10.676427in}{6.251016in}}%
\pgfpathcurveto{\pgfqpoint{10.666173in}{6.251016in}}{\pgfqpoint{10.656337in}{6.246942in}}{\pgfqpoint{10.649087in}{6.239691in}}%
\pgfpathcurveto{\pgfqpoint{10.641836in}{6.232440in}}{\pgfqpoint{10.637762in}{6.222605in}}{\pgfqpoint{10.637762in}{6.212351in}}%
\pgfpathcurveto{\pgfqpoint{10.637762in}{6.202097in}}{\pgfqpoint{10.641836in}{6.192261in}}{\pgfqpoint{10.649087in}{6.185010in}}%
\pgfpathcurveto{\pgfqpoint{10.656337in}{6.177760in}}{\pgfqpoint{10.666173in}{6.173686in}}{\pgfqpoint{10.676427in}{6.173686in}}%
\pgfpathlineto{\pgfqpoint{10.676427in}{6.173686in}}%
\pgfpathclose%
\pgfusepath{stroke}%
\end{pgfscope}%
\begin{pgfscope}%
\pgfpathrectangle{\pgfqpoint{8.036091in}{4.050000in}}{\pgfqpoint{5.400000in}{5.400000in}}%
\pgfusepath{clip}%
\pgfsetbuttcap%
\pgfsetroundjoin%
\pgfsetlinewidth{1.003750pt}%
\definecolor{currentstroke}{rgb}{0.541176,0.380392,0.658824}%
\pgfsetstrokecolor{currentstroke}%
\pgfsetstrokeopacity{0.900000}%
\pgfsetdash{}{0pt}%
\pgfpathmoveto{\pgfqpoint{12.023028in}{7.618075in}}%
\pgfpathcurveto{\pgfqpoint{12.033282in}{7.618075in}}{\pgfqpoint{12.043117in}{7.622149in}}{\pgfqpoint{12.050368in}{7.629400in}}%
\pgfpathcurveto{\pgfqpoint{12.057619in}{7.636651in}}{\pgfqpoint{12.061693in}{7.646486in}}{\pgfqpoint{12.061693in}{7.656740in}}%
\pgfpathcurveto{\pgfqpoint{12.061693in}{7.666994in}}{\pgfqpoint{12.057619in}{7.676830in}}{\pgfqpoint{12.050368in}{7.684081in}}%
\pgfpathcurveto{\pgfqpoint{12.043117in}{7.691331in}}{\pgfqpoint{12.033282in}{7.695405in}}{\pgfqpoint{12.023028in}{7.695405in}}%
\pgfpathcurveto{\pgfqpoint{12.012773in}{7.695405in}}{\pgfqpoint{12.002938in}{7.691331in}}{\pgfqpoint{11.995687in}{7.684081in}}%
\pgfpathcurveto{\pgfqpoint{11.988437in}{7.676830in}}{\pgfqpoint{11.984363in}{7.666994in}}{\pgfqpoint{11.984363in}{7.656740in}}%
\pgfpathcurveto{\pgfqpoint{11.984363in}{7.646486in}}{\pgfqpoint{11.988437in}{7.636651in}}{\pgfqpoint{11.995687in}{7.629400in}}%
\pgfpathcurveto{\pgfqpoint{12.002938in}{7.622149in}}{\pgfqpoint{12.012773in}{7.618075in}}{\pgfqpoint{12.023028in}{7.618075in}}%
\pgfpathlineto{\pgfqpoint{12.023028in}{7.618075in}}%
\pgfpathclose%
\pgfusepath{stroke}%
\end{pgfscope}%
\begin{pgfscope}%
\pgfpathrectangle{\pgfqpoint{8.036091in}{4.050000in}}{\pgfqpoint{5.400000in}{5.400000in}}%
\pgfusepath{clip}%
\pgfsetbuttcap%
\pgfsetroundjoin%
\pgfsetlinewidth{1.003750pt}%
\definecolor{currentstroke}{rgb}{0.541176,0.380392,0.658824}%
\pgfsetstrokecolor{currentstroke}%
\pgfsetstrokeopacity{0.900000}%
\pgfsetdash{}{0pt}%
\pgfpathmoveto{\pgfqpoint{11.754470in}{7.635806in}}%
\pgfpathcurveto{\pgfqpoint{11.764724in}{7.635806in}}{\pgfqpoint{11.774559in}{7.639880in}}{\pgfqpoint{11.781810in}{7.647130in}}%
\pgfpathcurveto{\pgfqpoint{11.789061in}{7.654381in}}{\pgfqpoint{11.793135in}{7.664217in}}{\pgfqpoint{11.793135in}{7.674471in}}%
\pgfpathcurveto{\pgfqpoint{11.793135in}{7.684725in}}{\pgfqpoint{11.789061in}{7.694560in}}{\pgfqpoint{11.781810in}{7.701811in}}%
\pgfpathcurveto{\pgfqpoint{11.774559in}{7.709062in}}{\pgfqpoint{11.764724in}{7.713136in}}{\pgfqpoint{11.754470in}{7.713136in}}%
\pgfpathcurveto{\pgfqpoint{11.744216in}{7.713136in}}{\pgfqpoint{11.734380in}{7.709062in}}{\pgfqpoint{11.727129in}{7.701811in}}%
\pgfpathcurveto{\pgfqpoint{11.719879in}{7.694560in}}{\pgfqpoint{11.715805in}{7.684725in}}{\pgfqpoint{11.715805in}{7.674471in}}%
\pgfpathcurveto{\pgfqpoint{11.715805in}{7.664217in}}{\pgfqpoint{11.719879in}{7.654381in}}{\pgfqpoint{11.727129in}{7.647130in}}%
\pgfpathcurveto{\pgfqpoint{11.734380in}{7.639880in}}{\pgfqpoint{11.744216in}{7.635806in}}{\pgfqpoint{11.754470in}{7.635806in}}%
\pgfpathlineto{\pgfqpoint{11.754470in}{7.635806in}}%
\pgfpathclose%
\pgfusepath{stroke}%
\end{pgfscope}%
\begin{pgfscope}%
\pgfpathrectangle{\pgfqpoint{8.036091in}{4.050000in}}{\pgfqpoint{5.400000in}{5.400000in}}%
\pgfusepath{clip}%
\pgfsetbuttcap%
\pgfsetroundjoin%
\pgfsetlinewidth{1.003750pt}%
\definecolor{currentstroke}{rgb}{0.541176,0.380392,0.658824}%
\pgfsetstrokecolor{currentstroke}%
\pgfsetstrokeopacity{0.900000}%
\pgfsetdash{}{0pt}%
\pgfpathmoveto{\pgfqpoint{12.138004in}{7.035549in}}%
\pgfpathcurveto{\pgfqpoint{12.148258in}{7.035549in}}{\pgfqpoint{12.158094in}{7.039623in}}{\pgfqpoint{12.165344in}{7.046874in}}%
\pgfpathcurveto{\pgfqpoint{12.172595in}{7.054125in}}{\pgfqpoint{12.176669in}{7.063960in}}{\pgfqpoint{12.176669in}{7.074214in}}%
\pgfpathcurveto{\pgfqpoint{12.176669in}{7.084468in}}{\pgfqpoint{12.172595in}{7.094304in}}{\pgfqpoint{12.165344in}{7.101555in}}%
\pgfpathcurveto{\pgfqpoint{12.158094in}{7.108805in}}{\pgfqpoint{12.148258in}{7.112879in}}{\pgfqpoint{12.138004in}{7.112879in}}%
\pgfpathcurveto{\pgfqpoint{12.127750in}{7.112879in}}{\pgfqpoint{12.117915in}{7.108805in}}{\pgfqpoint{12.110664in}{7.101555in}}%
\pgfpathcurveto{\pgfqpoint{12.103413in}{7.094304in}}{\pgfqpoint{12.099339in}{7.084468in}}{\pgfqpoint{12.099339in}{7.074214in}}%
\pgfpathcurveto{\pgfqpoint{12.099339in}{7.063960in}}{\pgfqpoint{12.103413in}{7.054125in}}{\pgfqpoint{12.110664in}{7.046874in}}%
\pgfpathcurveto{\pgfqpoint{12.117915in}{7.039623in}}{\pgfqpoint{12.127750in}{7.035549in}}{\pgfqpoint{12.138004in}{7.035549in}}%
\pgfpathlineto{\pgfqpoint{12.138004in}{7.035549in}}%
\pgfpathclose%
\pgfusepath{stroke}%
\end{pgfscope}%
\begin{pgfscope}%
\pgfpathrectangle{\pgfqpoint{8.036091in}{4.050000in}}{\pgfqpoint{5.400000in}{5.400000in}}%
\pgfusepath{clip}%
\pgfsetbuttcap%
\pgfsetroundjoin%
\pgfsetlinewidth{1.003750pt}%
\definecolor{currentstroke}{rgb}{0.541176,0.380392,0.658824}%
\pgfsetstrokecolor{currentstroke}%
\pgfsetstrokeopacity{0.900000}%
\pgfsetdash{}{0pt}%
\pgfpathmoveto{\pgfqpoint{12.215432in}{7.333379in}}%
\pgfpathcurveto{\pgfqpoint{12.225686in}{7.333379in}}{\pgfqpoint{12.235522in}{7.337453in}}{\pgfqpoint{12.242773in}{7.344704in}}%
\pgfpathcurveto{\pgfqpoint{12.250023in}{7.351954in}}{\pgfqpoint{12.254097in}{7.361790in}}{\pgfqpoint{12.254097in}{7.372044in}}%
\pgfpathcurveto{\pgfqpoint{12.254097in}{7.382298in}}{\pgfqpoint{12.250023in}{7.392134in}}{\pgfqpoint{12.242773in}{7.399384in}}%
\pgfpathcurveto{\pgfqpoint{12.235522in}{7.406635in}}{\pgfqpoint{12.225686in}{7.410709in}}{\pgfqpoint{12.215432in}{7.410709in}}%
\pgfpathcurveto{\pgfqpoint{12.205178in}{7.410709in}}{\pgfqpoint{12.195343in}{7.406635in}}{\pgfqpoint{12.188092in}{7.399384in}}%
\pgfpathcurveto{\pgfqpoint{12.180841in}{7.392134in}}{\pgfqpoint{12.176767in}{7.382298in}}{\pgfqpoint{12.176767in}{7.372044in}}%
\pgfpathcurveto{\pgfqpoint{12.176767in}{7.361790in}}{\pgfqpoint{12.180841in}{7.351954in}}{\pgfqpoint{12.188092in}{7.344704in}}%
\pgfpathcurveto{\pgfqpoint{12.195343in}{7.337453in}}{\pgfqpoint{12.205178in}{7.333379in}}{\pgfqpoint{12.215432in}{7.333379in}}%
\pgfpathlineto{\pgfqpoint{12.215432in}{7.333379in}}%
\pgfpathclose%
\pgfusepath{stroke}%
\end{pgfscope}%
\begin{pgfscope}%
\pgfpathrectangle{\pgfqpoint{8.036091in}{4.050000in}}{\pgfqpoint{5.400000in}{5.400000in}}%
\pgfusepath{clip}%
\pgfsetbuttcap%
\pgfsetroundjoin%
\pgfsetlinewidth{1.003750pt}%
\definecolor{currentstroke}{rgb}{0.541176,0.380392,0.658824}%
\pgfsetstrokecolor{currentstroke}%
\pgfsetstrokeopacity{0.900000}%
\pgfsetdash{}{0pt}%
\pgfpathmoveto{\pgfqpoint{10.919442in}{6.881247in}}%
\pgfpathcurveto{\pgfqpoint{10.929696in}{6.881247in}}{\pgfqpoint{10.939531in}{6.885321in}}{\pgfqpoint{10.946782in}{6.892572in}}%
\pgfpathcurveto{\pgfqpoint{10.954033in}{6.899823in}}{\pgfqpoint{10.958107in}{6.909658in}}{\pgfqpoint{10.958107in}{6.919913in}}%
\pgfpathcurveto{\pgfqpoint{10.958107in}{6.930167in}}{\pgfqpoint{10.954033in}{6.940002in}}{\pgfqpoint{10.946782in}{6.947253in}}%
\pgfpathcurveto{\pgfqpoint{10.939531in}{6.954504in}}{\pgfqpoint{10.929696in}{6.958578in}}{\pgfqpoint{10.919442in}{6.958578in}}%
\pgfpathcurveto{\pgfqpoint{10.909187in}{6.958578in}}{\pgfqpoint{10.899352in}{6.954504in}}{\pgfqpoint{10.892101in}{6.947253in}}%
\pgfpathcurveto{\pgfqpoint{10.884851in}{6.940002in}}{\pgfqpoint{10.880777in}{6.930167in}}{\pgfqpoint{10.880777in}{6.919913in}}%
\pgfpathcurveto{\pgfqpoint{10.880777in}{6.909658in}}{\pgfqpoint{10.884851in}{6.899823in}}{\pgfqpoint{10.892101in}{6.892572in}}%
\pgfpathcurveto{\pgfqpoint{10.899352in}{6.885321in}}{\pgfqpoint{10.909187in}{6.881247in}}{\pgfqpoint{10.919442in}{6.881247in}}%
\pgfpathlineto{\pgfqpoint{10.919442in}{6.881247in}}%
\pgfpathclose%
\pgfusepath{stroke}%
\end{pgfscope}%
\begin{pgfscope}%
\pgfpathrectangle{\pgfqpoint{8.036091in}{4.050000in}}{\pgfqpoint{5.400000in}{5.400000in}}%
\pgfusepath{clip}%
\pgfsetbuttcap%
\pgfsetroundjoin%
\pgfsetlinewidth{1.003750pt}%
\definecolor{currentstroke}{rgb}{0.541176,0.380392,0.658824}%
\pgfsetstrokecolor{currentstroke}%
\pgfsetstrokeopacity{0.900000}%
\pgfsetdash{}{0pt}%
\pgfpathmoveto{\pgfqpoint{11.644131in}{7.541775in}}%
\pgfpathcurveto{\pgfqpoint{11.654385in}{7.541775in}}{\pgfqpoint{11.664221in}{7.545849in}}{\pgfqpoint{11.671472in}{7.553099in}}%
\pgfpathcurveto{\pgfqpoint{11.678722in}{7.560350in}}{\pgfqpoint{11.682796in}{7.570186in}}{\pgfqpoint{11.682796in}{7.580440in}}%
\pgfpathcurveto{\pgfqpoint{11.682796in}{7.590694in}}{\pgfqpoint{11.678722in}{7.600529in}}{\pgfqpoint{11.671472in}{7.607780in}}%
\pgfpathcurveto{\pgfqpoint{11.664221in}{7.615031in}}{\pgfqpoint{11.654385in}{7.619105in}}{\pgfqpoint{11.644131in}{7.619105in}}%
\pgfpathcurveto{\pgfqpoint{11.633877in}{7.619105in}}{\pgfqpoint{11.624042in}{7.615031in}}{\pgfqpoint{11.616791in}{7.607780in}}%
\pgfpathcurveto{\pgfqpoint{11.609540in}{7.600529in}}{\pgfqpoint{11.605466in}{7.590694in}}{\pgfqpoint{11.605466in}{7.580440in}}%
\pgfpathcurveto{\pgfqpoint{11.605466in}{7.570186in}}{\pgfqpoint{11.609540in}{7.560350in}}{\pgfqpoint{11.616791in}{7.553099in}}%
\pgfpathcurveto{\pgfqpoint{11.624042in}{7.545849in}}{\pgfqpoint{11.633877in}{7.541775in}}{\pgfqpoint{11.644131in}{7.541775in}}%
\pgfpathlineto{\pgfqpoint{11.644131in}{7.541775in}}%
\pgfpathclose%
\pgfusepath{stroke}%
\end{pgfscope}%
\begin{pgfscope}%
\pgfpathrectangle{\pgfqpoint{8.036091in}{4.050000in}}{\pgfqpoint{5.400000in}{5.400000in}}%
\pgfusepath{clip}%
\pgfsetbuttcap%
\pgfsetroundjoin%
\pgfsetlinewidth{1.003750pt}%
\definecolor{currentstroke}{rgb}{0.541176,0.380392,0.658824}%
\pgfsetstrokecolor{currentstroke}%
\pgfsetstrokeopacity{0.900000}%
\pgfsetdash{}{0pt}%
\pgfpathmoveto{\pgfqpoint{11.471978in}{8.194617in}}%
\pgfpathcurveto{\pgfqpoint{11.482232in}{8.194617in}}{\pgfqpoint{11.492067in}{8.198691in}}{\pgfqpoint{11.499318in}{8.205941in}}%
\pgfpathcurveto{\pgfqpoint{11.506569in}{8.213192in}}{\pgfqpoint{11.510643in}{8.223028in}}{\pgfqpoint{11.510643in}{8.233282in}}%
\pgfpathcurveto{\pgfqpoint{11.510643in}{8.243536in}}{\pgfqpoint{11.506569in}{8.253371in}}{\pgfqpoint{11.499318in}{8.260622in}}%
\pgfpathcurveto{\pgfqpoint{11.492067in}{8.267873in}}{\pgfqpoint{11.482232in}{8.271947in}}{\pgfqpoint{11.471978in}{8.271947in}}%
\pgfpathcurveto{\pgfqpoint{11.461724in}{8.271947in}}{\pgfqpoint{11.451888in}{8.267873in}}{\pgfqpoint{11.444637in}{8.260622in}}%
\pgfpathcurveto{\pgfqpoint{11.437387in}{8.253371in}}{\pgfqpoint{11.433313in}{8.243536in}}{\pgfqpoint{11.433313in}{8.233282in}}%
\pgfpathcurveto{\pgfqpoint{11.433313in}{8.223028in}}{\pgfqpoint{11.437387in}{8.213192in}}{\pgfqpoint{11.444637in}{8.205941in}}%
\pgfpathcurveto{\pgfqpoint{11.451888in}{8.198691in}}{\pgfqpoint{11.461724in}{8.194617in}}{\pgfqpoint{11.471978in}{8.194617in}}%
\pgfpathlineto{\pgfqpoint{11.471978in}{8.194617in}}%
\pgfpathclose%
\pgfusepath{stroke}%
\end{pgfscope}%
\begin{pgfscope}%
\pgfpathrectangle{\pgfqpoint{8.036091in}{4.050000in}}{\pgfqpoint{5.400000in}{5.400000in}}%
\pgfusepath{clip}%
\pgfsetbuttcap%
\pgfsetroundjoin%
\pgfsetlinewidth{1.003750pt}%
\definecolor{currentstroke}{rgb}{0.541176,0.380392,0.658824}%
\pgfsetstrokecolor{currentstroke}%
\pgfsetstrokeopacity{0.900000}%
\pgfsetdash{}{0pt}%
\pgfpathmoveto{\pgfqpoint{11.727342in}{8.736117in}}%
\pgfpathcurveto{\pgfqpoint{11.737596in}{8.736117in}}{\pgfqpoint{11.747432in}{8.740191in}}{\pgfqpoint{11.754682in}{8.747442in}}%
\pgfpathcurveto{\pgfqpoint{11.761933in}{8.754693in}}{\pgfqpoint{11.766007in}{8.764528in}}{\pgfqpoint{11.766007in}{8.774783in}}%
\pgfpathcurveto{\pgfqpoint{11.766007in}{8.785037in}}{\pgfqpoint{11.761933in}{8.794872in}}{\pgfqpoint{11.754682in}{8.802123in}}%
\pgfpathcurveto{\pgfqpoint{11.747432in}{8.809374in}}{\pgfqpoint{11.737596in}{8.813448in}}{\pgfqpoint{11.727342in}{8.813448in}}%
\pgfpathcurveto{\pgfqpoint{11.717088in}{8.813448in}}{\pgfqpoint{11.707252in}{8.809374in}}{\pgfqpoint{11.700002in}{8.802123in}}%
\pgfpathcurveto{\pgfqpoint{11.692751in}{8.794872in}}{\pgfqpoint{11.688677in}{8.785037in}}{\pgfqpoint{11.688677in}{8.774783in}}%
\pgfpathcurveto{\pgfqpoint{11.688677in}{8.764528in}}{\pgfqpoint{11.692751in}{8.754693in}}{\pgfqpoint{11.700002in}{8.747442in}}%
\pgfpathcurveto{\pgfqpoint{11.707252in}{8.740191in}}{\pgfqpoint{11.717088in}{8.736117in}}{\pgfqpoint{11.727342in}{8.736117in}}%
\pgfpathlineto{\pgfqpoint{11.727342in}{8.736117in}}%
\pgfpathclose%
\pgfusepath{stroke}%
\end{pgfscope}%
\begin{pgfscope}%
\pgfpathrectangle{\pgfqpoint{8.036091in}{4.050000in}}{\pgfqpoint{5.400000in}{5.400000in}}%
\pgfusepath{clip}%
\pgfsetbuttcap%
\pgfsetroundjoin%
\pgfsetlinewidth{1.003750pt}%
\definecolor{currentstroke}{rgb}{0.541176,0.380392,0.658824}%
\pgfsetstrokecolor{currentstroke}%
\pgfsetstrokeopacity{0.900000}%
\pgfsetdash{}{0pt}%
\pgfpathmoveto{\pgfqpoint{11.483224in}{8.034880in}}%
\pgfpathcurveto{\pgfqpoint{11.493478in}{8.034880in}}{\pgfqpoint{11.503314in}{8.038954in}}{\pgfqpoint{11.510565in}{8.046205in}}%
\pgfpathcurveto{\pgfqpoint{11.517815in}{8.053456in}}{\pgfqpoint{11.521889in}{8.063291in}}{\pgfqpoint{11.521889in}{8.073545in}}%
\pgfpathcurveto{\pgfqpoint{11.521889in}{8.083799in}}{\pgfqpoint{11.517815in}{8.093635in}}{\pgfqpoint{11.510565in}{8.100886in}}%
\pgfpathcurveto{\pgfqpoint{11.503314in}{8.108136in}}{\pgfqpoint{11.493478in}{8.112210in}}{\pgfqpoint{11.483224in}{8.112210in}}%
\pgfpathcurveto{\pgfqpoint{11.472970in}{8.112210in}}{\pgfqpoint{11.463135in}{8.108136in}}{\pgfqpoint{11.455884in}{8.100886in}}%
\pgfpathcurveto{\pgfqpoint{11.448633in}{8.093635in}}{\pgfqpoint{11.444559in}{8.083799in}}{\pgfqpoint{11.444559in}{8.073545in}}%
\pgfpathcurveto{\pgfqpoint{11.444559in}{8.063291in}}{\pgfqpoint{11.448633in}{8.053456in}}{\pgfqpoint{11.455884in}{8.046205in}}%
\pgfpathcurveto{\pgfqpoint{11.463135in}{8.038954in}}{\pgfqpoint{11.472970in}{8.034880in}}{\pgfqpoint{11.483224in}{8.034880in}}%
\pgfpathlineto{\pgfqpoint{11.483224in}{8.034880in}}%
\pgfpathclose%
\pgfusepath{stroke}%
\end{pgfscope}%
\begin{pgfscope}%
\pgfpathrectangle{\pgfqpoint{8.036091in}{4.050000in}}{\pgfqpoint{5.400000in}{5.400000in}}%
\pgfusepath{clip}%
\pgfsetbuttcap%
\pgfsetroundjoin%
\pgfsetlinewidth{1.003750pt}%
\definecolor{currentstroke}{rgb}{0.541176,0.380392,0.658824}%
\pgfsetstrokecolor{currentstroke}%
\pgfsetstrokeopacity{0.900000}%
\pgfsetdash{}{0pt}%
\pgfpathmoveto{\pgfqpoint{11.906468in}{8.168182in}}%
\pgfpathcurveto{\pgfqpoint{11.916722in}{8.168182in}}{\pgfqpoint{11.926557in}{8.172256in}}{\pgfqpoint{11.933808in}{8.179506in}}%
\pgfpathcurveto{\pgfqpoint{11.941059in}{8.186757in}}{\pgfqpoint{11.945133in}{8.196593in}}{\pgfqpoint{11.945133in}{8.206847in}}%
\pgfpathcurveto{\pgfqpoint{11.945133in}{8.217101in}}{\pgfqpoint{11.941059in}{8.226936in}}{\pgfqpoint{11.933808in}{8.234187in}}%
\pgfpathcurveto{\pgfqpoint{11.926557in}{8.241438in}}{\pgfqpoint{11.916722in}{8.245512in}}{\pgfqpoint{11.906468in}{8.245512in}}%
\pgfpathcurveto{\pgfqpoint{11.896214in}{8.245512in}}{\pgfqpoint{11.886378in}{8.241438in}}{\pgfqpoint{11.879127in}{8.234187in}}%
\pgfpathcurveto{\pgfqpoint{11.871877in}{8.226936in}}{\pgfqpoint{11.867803in}{8.217101in}}{\pgfqpoint{11.867803in}{8.206847in}}%
\pgfpathcurveto{\pgfqpoint{11.867803in}{8.196593in}}{\pgfqpoint{11.871877in}{8.186757in}}{\pgfqpoint{11.879127in}{8.179506in}}%
\pgfpathcurveto{\pgfqpoint{11.886378in}{8.172256in}}{\pgfqpoint{11.896214in}{8.168182in}}{\pgfqpoint{11.906468in}{8.168182in}}%
\pgfpathlineto{\pgfqpoint{11.906468in}{8.168182in}}%
\pgfpathclose%
\pgfusepath{stroke}%
\end{pgfscope}%
\begin{pgfscope}%
\pgfpathrectangle{\pgfqpoint{8.036091in}{4.050000in}}{\pgfqpoint{5.400000in}{5.400000in}}%
\pgfusepath{clip}%
\pgfsetbuttcap%
\pgfsetroundjoin%
\pgfsetlinewidth{1.003750pt}%
\definecolor{currentstroke}{rgb}{0.541176,0.380392,0.658824}%
\pgfsetstrokecolor{currentstroke}%
\pgfsetstrokeopacity{0.900000}%
\pgfsetdash{}{0pt}%
\pgfpathmoveto{\pgfqpoint{9.717093in}{6.377954in}}%
\pgfpathcurveto{\pgfqpoint{9.727347in}{6.377954in}}{\pgfqpoint{9.737182in}{6.382028in}}{\pgfqpoint{9.744433in}{6.389278in}}%
\pgfpathcurveto{\pgfqpoint{9.751684in}{6.396529in}}{\pgfqpoint{9.755758in}{6.406364in}}{\pgfqpoint{9.755758in}{6.416619in}}%
\pgfpathcurveto{\pgfqpoint{9.755758in}{6.426873in}}{\pgfqpoint{9.751684in}{6.436708in}}{\pgfqpoint{9.744433in}{6.443959in}}%
\pgfpathcurveto{\pgfqpoint{9.737182in}{6.451210in}}{\pgfqpoint{9.727347in}{6.455284in}}{\pgfqpoint{9.717093in}{6.455284in}}%
\pgfpathcurveto{\pgfqpoint{9.706839in}{6.455284in}}{\pgfqpoint{9.697003in}{6.451210in}}{\pgfqpoint{9.689753in}{6.443959in}}%
\pgfpathcurveto{\pgfqpoint{9.682502in}{6.436708in}}{\pgfqpoint{9.678428in}{6.426873in}}{\pgfqpoint{9.678428in}{6.416619in}}%
\pgfpathcurveto{\pgfqpoint{9.678428in}{6.406364in}}{\pgfqpoint{9.682502in}{6.396529in}}{\pgfqpoint{9.689753in}{6.389278in}}%
\pgfpathcurveto{\pgfqpoint{9.697003in}{6.382028in}}{\pgfqpoint{9.706839in}{6.377954in}}{\pgfqpoint{9.717093in}{6.377954in}}%
\pgfpathlineto{\pgfqpoint{9.717093in}{6.377954in}}%
\pgfpathclose%
\pgfusepath{stroke}%
\end{pgfscope}%
\begin{pgfscope}%
\pgfpathrectangle{\pgfqpoint{8.036091in}{4.050000in}}{\pgfqpoint{5.400000in}{5.400000in}}%
\pgfusepath{clip}%
\pgfsetbuttcap%
\pgfsetroundjoin%
\pgfsetlinewidth{1.003750pt}%
\definecolor{currentstroke}{rgb}{0.541176,0.380392,0.658824}%
\pgfsetstrokecolor{currentstroke}%
\pgfsetstrokeopacity{0.900000}%
\pgfsetdash{}{0pt}%
\pgfpathmoveto{\pgfqpoint{11.551124in}{7.503075in}}%
\pgfpathcurveto{\pgfqpoint{11.561379in}{7.503075in}}{\pgfqpoint{11.571214in}{7.507149in}}{\pgfqpoint{11.578465in}{7.514400in}}%
\pgfpathcurveto{\pgfqpoint{11.585715in}{7.521651in}}{\pgfqpoint{11.589789in}{7.531486in}}{\pgfqpoint{11.589789in}{7.541740in}}%
\pgfpathcurveto{\pgfqpoint{11.589789in}{7.551994in}}{\pgfqpoint{11.585715in}{7.561830in}}{\pgfqpoint{11.578465in}{7.569081in}}%
\pgfpathcurveto{\pgfqpoint{11.571214in}{7.576331in}}{\pgfqpoint{11.561379in}{7.580405in}}{\pgfqpoint{11.551124in}{7.580405in}}%
\pgfpathcurveto{\pgfqpoint{11.540870in}{7.580405in}}{\pgfqpoint{11.531035in}{7.576331in}}{\pgfqpoint{11.523784in}{7.569081in}}%
\pgfpathcurveto{\pgfqpoint{11.516533in}{7.561830in}}{\pgfqpoint{11.512459in}{7.551994in}}{\pgfqpoint{11.512459in}{7.541740in}}%
\pgfpathcurveto{\pgfqpoint{11.512459in}{7.531486in}}{\pgfqpoint{11.516533in}{7.521651in}}{\pgfqpoint{11.523784in}{7.514400in}}%
\pgfpathcurveto{\pgfqpoint{11.531035in}{7.507149in}}{\pgfqpoint{11.540870in}{7.503075in}}{\pgfqpoint{11.551124in}{7.503075in}}%
\pgfpathlineto{\pgfqpoint{11.551124in}{7.503075in}}%
\pgfpathclose%
\pgfusepath{stroke}%
\end{pgfscope}%
\begin{pgfscope}%
\pgfpathrectangle{\pgfqpoint{8.036091in}{4.050000in}}{\pgfqpoint{5.400000in}{5.400000in}}%
\pgfusepath{clip}%
\pgfsetbuttcap%
\pgfsetroundjoin%
\pgfsetlinewidth{1.003750pt}%
\definecolor{currentstroke}{rgb}{0.541176,0.380392,0.658824}%
\pgfsetstrokecolor{currentstroke}%
\pgfsetstrokeopacity{0.900000}%
\pgfsetdash{}{0pt}%
\pgfpathmoveto{\pgfqpoint{10.642530in}{7.312239in}}%
\pgfpathcurveto{\pgfqpoint{10.652784in}{7.312239in}}{\pgfqpoint{10.662619in}{7.316313in}}{\pgfqpoint{10.669870in}{7.323564in}}%
\pgfpathcurveto{\pgfqpoint{10.677121in}{7.330815in}}{\pgfqpoint{10.681195in}{7.340650in}}{\pgfqpoint{10.681195in}{7.350904in}}%
\pgfpathcurveto{\pgfqpoint{10.681195in}{7.361158in}}{\pgfqpoint{10.677121in}{7.370994in}}{\pgfqpoint{10.669870in}{7.378245in}}%
\pgfpathcurveto{\pgfqpoint{10.662619in}{7.385495in}}{\pgfqpoint{10.652784in}{7.389569in}}{\pgfqpoint{10.642530in}{7.389569in}}%
\pgfpathcurveto{\pgfqpoint{10.632275in}{7.389569in}}{\pgfqpoint{10.622440in}{7.385495in}}{\pgfqpoint{10.615189in}{7.378245in}}%
\pgfpathcurveto{\pgfqpoint{10.607938in}{7.370994in}}{\pgfqpoint{10.603864in}{7.361158in}}{\pgfqpoint{10.603864in}{7.350904in}}%
\pgfpathcurveto{\pgfqpoint{10.603864in}{7.340650in}}{\pgfqpoint{10.607938in}{7.330815in}}{\pgfqpoint{10.615189in}{7.323564in}}%
\pgfpathcurveto{\pgfqpoint{10.622440in}{7.316313in}}{\pgfqpoint{10.632275in}{7.312239in}}{\pgfqpoint{10.642530in}{7.312239in}}%
\pgfpathlineto{\pgfqpoint{10.642530in}{7.312239in}}%
\pgfpathclose%
\pgfusepath{stroke}%
\end{pgfscope}%
\begin{pgfscope}%
\pgfpathrectangle{\pgfqpoint{8.036091in}{4.050000in}}{\pgfqpoint{5.400000in}{5.400000in}}%
\pgfusepath{clip}%
\pgfsetbuttcap%
\pgfsetroundjoin%
\pgfsetlinewidth{1.003750pt}%
\definecolor{currentstroke}{rgb}{0.541176,0.380392,0.658824}%
\pgfsetstrokecolor{currentstroke}%
\pgfsetstrokeopacity{0.900000}%
\pgfsetdash{}{0pt}%
\pgfpathmoveto{\pgfqpoint{11.838975in}{8.097504in}}%
\pgfpathcurveto{\pgfqpoint{11.849229in}{8.097504in}}{\pgfqpoint{11.859065in}{8.101578in}}{\pgfqpoint{11.866316in}{8.108829in}}%
\pgfpathcurveto{\pgfqpoint{11.873566in}{8.116079in}}{\pgfqpoint{11.877640in}{8.125915in}}{\pgfqpoint{11.877640in}{8.136169in}}%
\pgfpathcurveto{\pgfqpoint{11.877640in}{8.146423in}}{\pgfqpoint{11.873566in}{8.156259in}}{\pgfqpoint{11.866316in}{8.163509in}}%
\pgfpathcurveto{\pgfqpoint{11.859065in}{8.170760in}}{\pgfqpoint{11.849229in}{8.174834in}}{\pgfqpoint{11.838975in}{8.174834in}}%
\pgfpathcurveto{\pgfqpoint{11.828721in}{8.174834in}}{\pgfqpoint{11.818886in}{8.170760in}}{\pgfqpoint{11.811635in}{8.163509in}}%
\pgfpathcurveto{\pgfqpoint{11.804384in}{8.156259in}}{\pgfqpoint{11.800310in}{8.146423in}}{\pgfqpoint{11.800310in}{8.136169in}}%
\pgfpathcurveto{\pgfqpoint{11.800310in}{8.125915in}}{\pgfqpoint{11.804384in}{8.116079in}}{\pgfqpoint{11.811635in}{8.108829in}}%
\pgfpathcurveto{\pgfqpoint{11.818886in}{8.101578in}}{\pgfqpoint{11.828721in}{8.097504in}}{\pgfqpoint{11.838975in}{8.097504in}}%
\pgfpathlineto{\pgfqpoint{11.838975in}{8.097504in}}%
\pgfpathclose%
\pgfusepath{stroke}%
\end{pgfscope}%
\begin{pgfscope}%
\pgfpathrectangle{\pgfqpoint{8.036091in}{4.050000in}}{\pgfqpoint{5.400000in}{5.400000in}}%
\pgfusepath{clip}%
\pgfsetbuttcap%
\pgfsetroundjoin%
\pgfsetlinewidth{1.003750pt}%
\definecolor{currentstroke}{rgb}{0.541176,0.380392,0.658824}%
\pgfsetstrokecolor{currentstroke}%
\pgfsetstrokeopacity{0.900000}%
\pgfsetdash{}{0pt}%
\pgfpathmoveto{\pgfqpoint{11.175564in}{7.814747in}}%
\pgfpathcurveto{\pgfqpoint{11.185818in}{7.814747in}}{\pgfqpoint{11.195654in}{7.818821in}}{\pgfqpoint{11.202905in}{7.826072in}}%
\pgfpathcurveto{\pgfqpoint{11.210155in}{7.833323in}}{\pgfqpoint{11.214229in}{7.843158in}}{\pgfqpoint{11.214229in}{7.853412in}}%
\pgfpathcurveto{\pgfqpoint{11.214229in}{7.863666in}}{\pgfqpoint{11.210155in}{7.873502in}}{\pgfqpoint{11.202905in}{7.880753in}}%
\pgfpathcurveto{\pgfqpoint{11.195654in}{7.888003in}}{\pgfqpoint{11.185818in}{7.892077in}}{\pgfqpoint{11.175564in}{7.892077in}}%
\pgfpathcurveto{\pgfqpoint{11.165310in}{7.892077in}}{\pgfqpoint{11.155475in}{7.888003in}}{\pgfqpoint{11.148224in}{7.880753in}}%
\pgfpathcurveto{\pgfqpoint{11.140973in}{7.873502in}}{\pgfqpoint{11.136899in}{7.863666in}}{\pgfqpoint{11.136899in}{7.853412in}}%
\pgfpathcurveto{\pgfqpoint{11.136899in}{7.843158in}}{\pgfqpoint{11.140973in}{7.833323in}}{\pgfqpoint{11.148224in}{7.826072in}}%
\pgfpathcurveto{\pgfqpoint{11.155475in}{7.818821in}}{\pgfqpoint{11.165310in}{7.814747in}}{\pgfqpoint{11.175564in}{7.814747in}}%
\pgfpathlineto{\pgfqpoint{11.175564in}{7.814747in}}%
\pgfpathclose%
\pgfusepath{stroke}%
\end{pgfscope}%
\begin{pgfscope}%
\pgfpathrectangle{\pgfqpoint{8.036091in}{4.050000in}}{\pgfqpoint{5.400000in}{5.400000in}}%
\pgfusepath{clip}%
\pgfsetbuttcap%
\pgfsetroundjoin%
\pgfsetlinewidth{1.003750pt}%
\definecolor{currentstroke}{rgb}{0.541176,0.380392,0.658824}%
\pgfsetstrokecolor{currentstroke}%
\pgfsetstrokeopacity{0.900000}%
\pgfsetdash{}{0pt}%
\pgfpathmoveto{\pgfqpoint{11.088401in}{6.905893in}}%
\pgfpathcurveto{\pgfqpoint{11.098656in}{6.905893in}}{\pgfqpoint{11.108491in}{6.909967in}}{\pgfqpoint{11.115742in}{6.917218in}}%
\pgfpathcurveto{\pgfqpoint{11.122993in}{6.924468in}}{\pgfqpoint{11.127067in}{6.934304in}}{\pgfqpoint{11.127067in}{6.944558in}}%
\pgfpathcurveto{\pgfqpoint{11.127067in}{6.954812in}}{\pgfqpoint{11.122993in}{6.964647in}}{\pgfqpoint{11.115742in}{6.971898in}}%
\pgfpathcurveto{\pgfqpoint{11.108491in}{6.979149in}}{\pgfqpoint{11.098656in}{6.983223in}}{\pgfqpoint{11.088401in}{6.983223in}}%
\pgfpathcurveto{\pgfqpoint{11.078147in}{6.983223in}}{\pgfqpoint{11.068312in}{6.979149in}}{\pgfqpoint{11.061061in}{6.971898in}}%
\pgfpathcurveto{\pgfqpoint{11.053810in}{6.964647in}}{\pgfqpoint{11.049736in}{6.954812in}}{\pgfqpoint{11.049736in}{6.944558in}}%
\pgfpathcurveto{\pgfqpoint{11.049736in}{6.934304in}}{\pgfqpoint{11.053810in}{6.924468in}}{\pgfqpoint{11.061061in}{6.917218in}}%
\pgfpathcurveto{\pgfqpoint{11.068312in}{6.909967in}}{\pgfqpoint{11.078147in}{6.905893in}}{\pgfqpoint{11.088401in}{6.905893in}}%
\pgfpathlineto{\pgfqpoint{11.088401in}{6.905893in}}%
\pgfpathclose%
\pgfusepath{stroke}%
\end{pgfscope}%
\begin{pgfscope}%
\pgfpathrectangle{\pgfqpoint{8.036091in}{4.050000in}}{\pgfqpoint{5.400000in}{5.400000in}}%
\pgfusepath{clip}%
\pgfsetbuttcap%
\pgfsetroundjoin%
\pgfsetlinewidth{1.003750pt}%
\definecolor{currentstroke}{rgb}{0.541176,0.380392,0.658824}%
\pgfsetstrokecolor{currentstroke}%
\pgfsetstrokeopacity{0.900000}%
\pgfsetdash{}{0pt}%
\pgfpathmoveto{\pgfqpoint{9.901637in}{4.917246in}}%
\pgfpathcurveto{\pgfqpoint{9.911891in}{4.917246in}}{\pgfqpoint{9.921727in}{4.921320in}}{\pgfqpoint{9.928977in}{4.928570in}}%
\pgfpathcurveto{\pgfqpoint{9.936228in}{4.935821in}}{\pgfqpoint{9.940302in}{4.945657in}}{\pgfqpoint{9.940302in}{4.955911in}}%
\pgfpathcurveto{\pgfqpoint{9.940302in}{4.966165in}}{\pgfqpoint{9.936228in}{4.976000in}}{\pgfqpoint{9.928977in}{4.983251in}}%
\pgfpathcurveto{\pgfqpoint{9.921727in}{4.990502in}}{\pgfqpoint{9.911891in}{4.994576in}}{\pgfqpoint{9.901637in}{4.994576in}}%
\pgfpathcurveto{\pgfqpoint{9.891383in}{4.994576in}}{\pgfqpoint{9.881547in}{4.990502in}}{\pgfqpoint{9.874297in}{4.983251in}}%
\pgfpathcurveto{\pgfqpoint{9.867046in}{4.976000in}}{\pgfqpoint{9.862972in}{4.966165in}}{\pgfqpoint{9.862972in}{4.955911in}}%
\pgfpathcurveto{\pgfqpoint{9.862972in}{4.945657in}}{\pgfqpoint{9.867046in}{4.935821in}}{\pgfqpoint{9.874297in}{4.928570in}}%
\pgfpathcurveto{\pgfqpoint{9.881547in}{4.921320in}}{\pgfqpoint{9.891383in}{4.917246in}}{\pgfqpoint{9.901637in}{4.917246in}}%
\pgfpathlineto{\pgfqpoint{9.901637in}{4.917246in}}%
\pgfpathclose%
\pgfusepath{stroke}%
\end{pgfscope}%
\begin{pgfscope}%
\pgfpathrectangle{\pgfqpoint{8.036091in}{4.050000in}}{\pgfqpoint{5.400000in}{5.400000in}}%
\pgfusepath{clip}%
\pgfsetbuttcap%
\pgfsetroundjoin%
\pgfsetlinewidth{1.003750pt}%
\definecolor{currentstroke}{rgb}{0.541176,0.380392,0.658824}%
\pgfsetstrokecolor{currentstroke}%
\pgfsetstrokeopacity{0.900000}%
\pgfsetdash{}{0pt}%
\pgfpathmoveto{\pgfqpoint{10.087778in}{6.542632in}}%
\pgfpathcurveto{\pgfqpoint{10.098032in}{6.542632in}}{\pgfqpoint{10.107867in}{6.546706in}}{\pgfqpoint{10.115118in}{6.553957in}}%
\pgfpathcurveto{\pgfqpoint{10.122369in}{6.561207in}}{\pgfqpoint{10.126443in}{6.571043in}}{\pgfqpoint{10.126443in}{6.581297in}}%
\pgfpathcurveto{\pgfqpoint{10.126443in}{6.591551in}}{\pgfqpoint{10.122369in}{6.601387in}}{\pgfqpoint{10.115118in}{6.608637in}}%
\pgfpathcurveto{\pgfqpoint{10.107867in}{6.615888in}}{\pgfqpoint{10.098032in}{6.619962in}}{\pgfqpoint{10.087778in}{6.619962in}}%
\pgfpathcurveto{\pgfqpoint{10.077524in}{6.619962in}}{\pgfqpoint{10.067688in}{6.615888in}}{\pgfqpoint{10.060438in}{6.608637in}}%
\pgfpathcurveto{\pgfqpoint{10.053187in}{6.601387in}}{\pgfqpoint{10.049113in}{6.591551in}}{\pgfqpoint{10.049113in}{6.581297in}}%
\pgfpathcurveto{\pgfqpoint{10.049113in}{6.571043in}}{\pgfqpoint{10.053187in}{6.561207in}}{\pgfqpoint{10.060438in}{6.553957in}}%
\pgfpathcurveto{\pgfqpoint{10.067688in}{6.546706in}}{\pgfqpoint{10.077524in}{6.542632in}}{\pgfqpoint{10.087778in}{6.542632in}}%
\pgfpathlineto{\pgfqpoint{10.087778in}{6.542632in}}%
\pgfpathclose%
\pgfusepath{stroke}%
\end{pgfscope}%
\begin{pgfscope}%
\pgfpathrectangle{\pgfqpoint{8.036091in}{4.050000in}}{\pgfqpoint{5.400000in}{5.400000in}}%
\pgfusepath{clip}%
\pgfsetbuttcap%
\pgfsetroundjoin%
\pgfsetlinewidth{1.003750pt}%
\definecolor{currentstroke}{rgb}{0.541176,0.380392,0.658824}%
\pgfsetstrokecolor{currentstroke}%
\pgfsetstrokeopacity{0.900000}%
\pgfsetdash{}{0pt}%
\pgfpathmoveto{\pgfqpoint{10.321577in}{5.887536in}}%
\pgfpathcurveto{\pgfqpoint{10.331831in}{5.887536in}}{\pgfqpoint{10.341667in}{5.891610in}}{\pgfqpoint{10.348918in}{5.898861in}}%
\pgfpathcurveto{\pgfqpoint{10.356168in}{5.906112in}}{\pgfqpoint{10.360242in}{5.915947in}}{\pgfqpoint{10.360242in}{5.926201in}}%
\pgfpathcurveto{\pgfqpoint{10.360242in}{5.936456in}}{\pgfqpoint{10.356168in}{5.946291in}}{\pgfqpoint{10.348918in}{5.953542in}}%
\pgfpathcurveto{\pgfqpoint{10.341667in}{5.960793in}}{\pgfqpoint{10.331831in}{5.964867in}}{\pgfqpoint{10.321577in}{5.964867in}}%
\pgfpathcurveto{\pgfqpoint{10.311323in}{5.964867in}}{\pgfqpoint{10.301488in}{5.960793in}}{\pgfqpoint{10.294237in}{5.953542in}}%
\pgfpathcurveto{\pgfqpoint{10.286986in}{5.946291in}}{\pgfqpoint{10.282912in}{5.936456in}}{\pgfqpoint{10.282912in}{5.926201in}}%
\pgfpathcurveto{\pgfqpoint{10.282912in}{5.915947in}}{\pgfqpoint{10.286986in}{5.906112in}}{\pgfqpoint{10.294237in}{5.898861in}}%
\pgfpathcurveto{\pgfqpoint{10.301488in}{5.891610in}}{\pgfqpoint{10.311323in}{5.887536in}}{\pgfqpoint{10.321577in}{5.887536in}}%
\pgfpathlineto{\pgfqpoint{10.321577in}{5.887536in}}%
\pgfpathclose%
\pgfusepath{stroke}%
\end{pgfscope}%
\begin{pgfscope}%
\pgfpathrectangle{\pgfqpoint{8.036091in}{4.050000in}}{\pgfqpoint{5.400000in}{5.400000in}}%
\pgfusepath{clip}%
\pgfsetbuttcap%
\pgfsetroundjoin%
\pgfsetlinewidth{1.003750pt}%
\definecolor{currentstroke}{rgb}{0.541176,0.380392,0.658824}%
\pgfsetstrokecolor{currentstroke}%
\pgfsetstrokeopacity{0.900000}%
\pgfsetdash{}{0pt}%
\pgfpathmoveto{\pgfqpoint{9.880194in}{5.460470in}}%
\pgfpathcurveto{\pgfqpoint{9.890449in}{5.460470in}}{\pgfqpoint{9.900284in}{5.464544in}}{\pgfqpoint{9.907535in}{5.471795in}}%
\pgfpathcurveto{\pgfqpoint{9.914786in}{5.479046in}}{\pgfqpoint{9.918859in}{5.488881in}}{\pgfqpoint{9.918859in}{5.499135in}}%
\pgfpathcurveto{\pgfqpoint{9.918859in}{5.509389in}}{\pgfqpoint{9.914786in}{5.519225in}}{\pgfqpoint{9.907535in}{5.526476in}}%
\pgfpathcurveto{\pgfqpoint{9.900284in}{5.533726in}}{\pgfqpoint{9.890449in}{5.537800in}}{\pgfqpoint{9.880194in}{5.537800in}}%
\pgfpathcurveto{\pgfqpoint{9.869940in}{5.537800in}}{\pgfqpoint{9.860105in}{5.533726in}}{\pgfqpoint{9.852854in}{5.526476in}}%
\pgfpathcurveto{\pgfqpoint{9.845603in}{5.519225in}}{\pgfqpoint{9.841529in}{5.509389in}}{\pgfqpoint{9.841529in}{5.499135in}}%
\pgfpathcurveto{\pgfqpoint{9.841529in}{5.488881in}}{\pgfqpoint{9.845603in}{5.479046in}}{\pgfqpoint{9.852854in}{5.471795in}}%
\pgfpathcurveto{\pgfqpoint{9.860105in}{5.464544in}}{\pgfqpoint{9.869940in}{5.460470in}}{\pgfqpoint{9.880194in}{5.460470in}}%
\pgfpathlineto{\pgfqpoint{9.880194in}{5.460470in}}%
\pgfpathclose%
\pgfusepath{stroke}%
\end{pgfscope}%
\begin{pgfscope}%
\pgfpathrectangle{\pgfqpoint{8.036091in}{4.050000in}}{\pgfqpoint{5.400000in}{5.400000in}}%
\pgfusepath{clip}%
\pgfsetbuttcap%
\pgfsetroundjoin%
\pgfsetlinewidth{1.003750pt}%
\definecolor{currentstroke}{rgb}{0.541176,0.380392,0.658824}%
\pgfsetstrokecolor{currentstroke}%
\pgfsetstrokeopacity{0.900000}%
\pgfsetdash{}{0pt}%
\pgfpathmoveto{\pgfqpoint{10.307527in}{5.981720in}}%
\pgfpathcurveto{\pgfqpoint{10.317781in}{5.981720in}}{\pgfqpoint{10.327617in}{5.985794in}}{\pgfqpoint{10.334867in}{5.993045in}}%
\pgfpathcurveto{\pgfqpoint{10.342118in}{6.000296in}}{\pgfqpoint{10.346192in}{6.010131in}}{\pgfqpoint{10.346192in}{6.020385in}}%
\pgfpathcurveto{\pgfqpoint{10.346192in}{6.030640in}}{\pgfqpoint{10.342118in}{6.040475in}}{\pgfqpoint{10.334867in}{6.047726in}}%
\pgfpathcurveto{\pgfqpoint{10.327617in}{6.054976in}}{\pgfqpoint{10.317781in}{6.059050in}}{\pgfqpoint{10.307527in}{6.059050in}}%
\pgfpathcurveto{\pgfqpoint{10.297273in}{6.059050in}}{\pgfqpoint{10.287438in}{6.054976in}}{\pgfqpoint{10.280187in}{6.047726in}}%
\pgfpathcurveto{\pgfqpoint{10.272936in}{6.040475in}}{\pgfqpoint{10.268862in}{6.030640in}}{\pgfqpoint{10.268862in}{6.020385in}}%
\pgfpathcurveto{\pgfqpoint{10.268862in}{6.010131in}}{\pgfqpoint{10.272936in}{6.000296in}}{\pgfqpoint{10.280187in}{5.993045in}}%
\pgfpathcurveto{\pgfqpoint{10.287438in}{5.985794in}}{\pgfqpoint{10.297273in}{5.981720in}}{\pgfqpoint{10.307527in}{5.981720in}}%
\pgfpathlineto{\pgfqpoint{10.307527in}{5.981720in}}%
\pgfpathclose%
\pgfusepath{stroke}%
\end{pgfscope}%
\begin{pgfscope}%
\pgfpathrectangle{\pgfqpoint{8.036091in}{4.050000in}}{\pgfqpoint{5.400000in}{5.400000in}}%
\pgfusepath{clip}%
\pgfsetbuttcap%
\pgfsetroundjoin%
\pgfsetlinewidth{1.003750pt}%
\definecolor{currentstroke}{rgb}{0.541176,0.380392,0.658824}%
\pgfsetstrokecolor{currentstroke}%
\pgfsetstrokeopacity{0.900000}%
\pgfsetdash{}{0pt}%
\pgfpathmoveto{\pgfqpoint{10.802358in}{6.630592in}}%
\pgfpathcurveto{\pgfqpoint{10.812612in}{6.630592in}}{\pgfqpoint{10.822447in}{6.634666in}}{\pgfqpoint{10.829698in}{6.641917in}}%
\pgfpathcurveto{\pgfqpoint{10.836949in}{6.649167in}}{\pgfqpoint{10.841023in}{6.659003in}}{\pgfqpoint{10.841023in}{6.669257in}}%
\pgfpathcurveto{\pgfqpoint{10.841023in}{6.679511in}}{\pgfqpoint{10.836949in}{6.689346in}}{\pgfqpoint{10.829698in}{6.696597in}}%
\pgfpathcurveto{\pgfqpoint{10.822447in}{6.703848in}}{\pgfqpoint{10.812612in}{6.707922in}}{\pgfqpoint{10.802358in}{6.707922in}}%
\pgfpathcurveto{\pgfqpoint{10.792103in}{6.707922in}}{\pgfqpoint{10.782268in}{6.703848in}}{\pgfqpoint{10.775017in}{6.696597in}}%
\pgfpathcurveto{\pgfqpoint{10.767766in}{6.689346in}}{\pgfqpoint{10.763692in}{6.679511in}}{\pgfqpoint{10.763692in}{6.669257in}}%
\pgfpathcurveto{\pgfqpoint{10.763692in}{6.659003in}}{\pgfqpoint{10.767766in}{6.649167in}}{\pgfqpoint{10.775017in}{6.641917in}}%
\pgfpathcurveto{\pgfqpoint{10.782268in}{6.634666in}}{\pgfqpoint{10.792103in}{6.630592in}}{\pgfqpoint{10.802358in}{6.630592in}}%
\pgfpathlineto{\pgfqpoint{10.802358in}{6.630592in}}%
\pgfpathclose%
\pgfusepath{stroke}%
\end{pgfscope}%
\begin{pgfscope}%
\pgfpathrectangle{\pgfqpoint{8.036091in}{4.050000in}}{\pgfqpoint{5.400000in}{5.400000in}}%
\pgfusepath{clip}%
\pgfsetbuttcap%
\pgfsetroundjoin%
\pgfsetlinewidth{1.003750pt}%
\definecolor{currentstroke}{rgb}{0.541176,0.380392,0.658824}%
\pgfsetstrokecolor{currentstroke}%
\pgfsetstrokeopacity{0.900000}%
\pgfsetdash{}{0pt}%
\pgfpathmoveto{\pgfqpoint{11.159517in}{7.095414in}}%
\pgfpathcurveto{\pgfqpoint{11.169771in}{7.095414in}}{\pgfqpoint{11.179606in}{7.099488in}}{\pgfqpoint{11.186857in}{7.106739in}}%
\pgfpathcurveto{\pgfqpoint{11.194108in}{7.113990in}}{\pgfqpoint{11.198182in}{7.123825in}}{\pgfqpoint{11.198182in}{7.134079in}}%
\pgfpathcurveto{\pgfqpoint{11.198182in}{7.144333in}}{\pgfqpoint{11.194108in}{7.154169in}}{\pgfqpoint{11.186857in}{7.161420in}}%
\pgfpathcurveto{\pgfqpoint{11.179606in}{7.168670in}}{\pgfqpoint{11.169771in}{7.172744in}}{\pgfqpoint{11.159517in}{7.172744in}}%
\pgfpathcurveto{\pgfqpoint{11.149263in}{7.172744in}}{\pgfqpoint{11.139427in}{7.168670in}}{\pgfqpoint{11.132176in}{7.161420in}}%
\pgfpathcurveto{\pgfqpoint{11.124926in}{7.154169in}}{\pgfqpoint{11.120852in}{7.144333in}}{\pgfqpoint{11.120852in}{7.134079in}}%
\pgfpathcurveto{\pgfqpoint{11.120852in}{7.123825in}}{\pgfqpoint{11.124926in}{7.113990in}}{\pgfqpoint{11.132176in}{7.106739in}}%
\pgfpathcurveto{\pgfqpoint{11.139427in}{7.099488in}}{\pgfqpoint{11.149263in}{7.095414in}}{\pgfqpoint{11.159517in}{7.095414in}}%
\pgfpathlineto{\pgfqpoint{11.159517in}{7.095414in}}%
\pgfpathclose%
\pgfusepath{stroke}%
\end{pgfscope}%
\begin{pgfscope}%
\pgfpathrectangle{\pgfqpoint{8.036091in}{4.050000in}}{\pgfqpoint{5.400000in}{5.400000in}}%
\pgfusepath{clip}%
\pgfsetbuttcap%
\pgfsetroundjoin%
\pgfsetlinewidth{1.003750pt}%
\definecolor{currentstroke}{rgb}{0.541176,0.380392,0.658824}%
\pgfsetstrokecolor{currentstroke}%
\pgfsetstrokeopacity{0.900000}%
\pgfsetdash{}{0pt}%
\pgfpathmoveto{\pgfqpoint{11.612285in}{7.815149in}}%
\pgfpathcurveto{\pgfqpoint{11.622539in}{7.815149in}}{\pgfqpoint{11.632375in}{7.819223in}}{\pgfqpoint{11.639626in}{7.826474in}}%
\pgfpathcurveto{\pgfqpoint{11.646876in}{7.833725in}}{\pgfqpoint{11.650950in}{7.843560in}}{\pgfqpoint{11.650950in}{7.853814in}}%
\pgfpathcurveto{\pgfqpoint{11.650950in}{7.864068in}}{\pgfqpoint{11.646876in}{7.873904in}}{\pgfqpoint{11.639626in}{7.881154in}}%
\pgfpathcurveto{\pgfqpoint{11.632375in}{7.888405in}}{\pgfqpoint{11.622539in}{7.892479in}}{\pgfqpoint{11.612285in}{7.892479in}}%
\pgfpathcurveto{\pgfqpoint{11.602031in}{7.892479in}}{\pgfqpoint{11.592196in}{7.888405in}}{\pgfqpoint{11.584945in}{7.881154in}}%
\pgfpathcurveto{\pgfqpoint{11.577694in}{7.873904in}}{\pgfqpoint{11.573620in}{7.864068in}}{\pgfqpoint{11.573620in}{7.853814in}}%
\pgfpathcurveto{\pgfqpoint{11.573620in}{7.843560in}}{\pgfqpoint{11.577694in}{7.833725in}}{\pgfqpoint{11.584945in}{7.826474in}}%
\pgfpathcurveto{\pgfqpoint{11.592196in}{7.819223in}}{\pgfqpoint{11.602031in}{7.815149in}}{\pgfqpoint{11.612285in}{7.815149in}}%
\pgfpathlineto{\pgfqpoint{11.612285in}{7.815149in}}%
\pgfpathclose%
\pgfusepath{stroke}%
\end{pgfscope}%
\begin{pgfscope}%
\pgfpathrectangle{\pgfqpoint{8.036091in}{4.050000in}}{\pgfqpoint{5.400000in}{5.400000in}}%
\pgfusepath{clip}%
\pgfsetbuttcap%
\pgfsetroundjoin%
\pgfsetlinewidth{1.003750pt}%
\definecolor{currentstroke}{rgb}{0.541176,0.380392,0.658824}%
\pgfsetstrokecolor{currentstroke}%
\pgfsetstrokeopacity{0.900000}%
\pgfsetdash{}{0pt}%
\pgfpathmoveto{\pgfqpoint{11.223624in}{6.962694in}}%
\pgfpathcurveto{\pgfqpoint{11.233879in}{6.962694in}}{\pgfqpoint{11.243714in}{6.966768in}}{\pgfqpoint{11.250965in}{6.974019in}}%
\pgfpathcurveto{\pgfqpoint{11.258216in}{6.981270in}}{\pgfqpoint{11.262290in}{6.991105in}}{\pgfqpoint{11.262290in}{7.001360in}}%
\pgfpathcurveto{\pgfqpoint{11.262290in}{7.011614in}}{\pgfqpoint{11.258216in}{7.021449in}}{\pgfqpoint{11.250965in}{7.028700in}}%
\pgfpathcurveto{\pgfqpoint{11.243714in}{7.035951in}}{\pgfqpoint{11.233879in}{7.040025in}}{\pgfqpoint{11.223624in}{7.040025in}}%
\pgfpathcurveto{\pgfqpoint{11.213370in}{7.040025in}}{\pgfqpoint{11.203535in}{7.035951in}}{\pgfqpoint{11.196284in}{7.028700in}}%
\pgfpathcurveto{\pgfqpoint{11.189033in}{7.021449in}}{\pgfqpoint{11.184959in}{7.011614in}}{\pgfqpoint{11.184959in}{7.001360in}}%
\pgfpathcurveto{\pgfqpoint{11.184959in}{6.991105in}}{\pgfqpoint{11.189033in}{6.981270in}}{\pgfqpoint{11.196284in}{6.974019in}}%
\pgfpathcurveto{\pgfqpoint{11.203535in}{6.966768in}}{\pgfqpoint{11.213370in}{6.962694in}}{\pgfqpoint{11.223624in}{6.962694in}}%
\pgfpathlineto{\pgfqpoint{11.223624in}{6.962694in}}%
\pgfpathclose%
\pgfusepath{stroke}%
\end{pgfscope}%
\begin{pgfscope}%
\pgfpathrectangle{\pgfqpoint{8.036091in}{4.050000in}}{\pgfqpoint{5.400000in}{5.400000in}}%
\pgfusepath{clip}%
\pgfsetbuttcap%
\pgfsetroundjoin%
\pgfsetlinewidth{1.003750pt}%
\definecolor{currentstroke}{rgb}{0.541176,0.380392,0.658824}%
\pgfsetstrokecolor{currentstroke}%
\pgfsetstrokeopacity{0.900000}%
\pgfsetdash{}{0pt}%
\pgfpathmoveto{\pgfqpoint{11.878524in}{7.510572in}}%
\pgfpathcurveto{\pgfqpoint{11.888778in}{7.510572in}}{\pgfqpoint{11.898614in}{7.514646in}}{\pgfqpoint{11.905865in}{7.521896in}}%
\pgfpathcurveto{\pgfqpoint{11.913115in}{7.529147in}}{\pgfqpoint{11.917189in}{7.538983in}}{\pgfqpoint{11.917189in}{7.549237in}}%
\pgfpathcurveto{\pgfqpoint{11.917189in}{7.559491in}}{\pgfqpoint{11.913115in}{7.569326in}}{\pgfqpoint{11.905865in}{7.576577in}}%
\pgfpathcurveto{\pgfqpoint{11.898614in}{7.583828in}}{\pgfqpoint{11.888778in}{7.587902in}}{\pgfqpoint{11.878524in}{7.587902in}}%
\pgfpathcurveto{\pgfqpoint{11.868270in}{7.587902in}}{\pgfqpoint{11.858435in}{7.583828in}}{\pgfqpoint{11.851184in}{7.576577in}}%
\pgfpathcurveto{\pgfqpoint{11.843933in}{7.569326in}}{\pgfqpoint{11.839859in}{7.559491in}}{\pgfqpoint{11.839859in}{7.549237in}}%
\pgfpathcurveto{\pgfqpoint{11.839859in}{7.538983in}}{\pgfqpoint{11.843933in}{7.529147in}}{\pgfqpoint{11.851184in}{7.521896in}}%
\pgfpathcurveto{\pgfqpoint{11.858435in}{7.514646in}}{\pgfqpoint{11.868270in}{7.510572in}}{\pgfqpoint{11.878524in}{7.510572in}}%
\pgfpathlineto{\pgfqpoint{11.878524in}{7.510572in}}%
\pgfpathclose%
\pgfusepath{stroke}%
\end{pgfscope}%
\begin{pgfscope}%
\pgfpathrectangle{\pgfqpoint{8.036091in}{4.050000in}}{\pgfqpoint{5.400000in}{5.400000in}}%
\pgfusepath{clip}%
\pgfsetbuttcap%
\pgfsetroundjoin%
\pgfsetlinewidth{1.003750pt}%
\definecolor{currentstroke}{rgb}{0.541176,0.380392,0.658824}%
\pgfsetstrokecolor{currentstroke}%
\pgfsetstrokeopacity{0.900000}%
\pgfsetdash{}{0pt}%
\pgfpathmoveto{\pgfqpoint{11.667599in}{7.253077in}}%
\pgfpathcurveto{\pgfqpoint{11.677853in}{7.253077in}}{\pgfqpoint{11.687688in}{7.257151in}}{\pgfqpoint{11.694939in}{7.264402in}}%
\pgfpathcurveto{\pgfqpoint{11.702190in}{7.271653in}}{\pgfqpoint{11.706264in}{7.281488in}}{\pgfqpoint{11.706264in}{7.291742in}}%
\pgfpathcurveto{\pgfqpoint{11.706264in}{7.301996in}}{\pgfqpoint{11.702190in}{7.311832in}}{\pgfqpoint{11.694939in}{7.319082in}}%
\pgfpathcurveto{\pgfqpoint{11.687688in}{7.326333in}}{\pgfqpoint{11.677853in}{7.330407in}}{\pgfqpoint{11.667599in}{7.330407in}}%
\pgfpathcurveto{\pgfqpoint{11.657344in}{7.330407in}}{\pgfqpoint{11.647509in}{7.326333in}}{\pgfqpoint{11.640258in}{7.319082in}}%
\pgfpathcurveto{\pgfqpoint{11.633007in}{7.311832in}}{\pgfqpoint{11.628933in}{7.301996in}}{\pgfqpoint{11.628933in}{7.291742in}}%
\pgfpathcurveto{\pgfqpoint{11.628933in}{7.281488in}}{\pgfqpoint{11.633007in}{7.271653in}}{\pgfqpoint{11.640258in}{7.264402in}}%
\pgfpathcurveto{\pgfqpoint{11.647509in}{7.257151in}}{\pgfqpoint{11.657344in}{7.253077in}}{\pgfqpoint{11.667599in}{7.253077in}}%
\pgfpathlineto{\pgfqpoint{11.667599in}{7.253077in}}%
\pgfpathclose%
\pgfusepath{stroke}%
\end{pgfscope}%
\begin{pgfscope}%
\pgfpathrectangle{\pgfqpoint{8.036091in}{4.050000in}}{\pgfqpoint{5.400000in}{5.400000in}}%
\pgfusepath{clip}%
\pgfsetbuttcap%
\pgfsetroundjoin%
\pgfsetlinewidth{1.003750pt}%
\definecolor{currentstroke}{rgb}{0.541176,0.380392,0.658824}%
\pgfsetstrokecolor{currentstroke}%
\pgfsetstrokeopacity{0.900000}%
\pgfsetdash{}{0pt}%
\pgfpathmoveto{\pgfqpoint{12.191543in}{7.653143in}}%
\pgfpathcurveto{\pgfqpoint{12.201797in}{7.653143in}}{\pgfqpoint{12.211632in}{7.657217in}}{\pgfqpoint{12.218883in}{7.664467in}}%
\pgfpathcurveto{\pgfqpoint{12.226134in}{7.671718in}}{\pgfqpoint{12.230208in}{7.681553in}}{\pgfqpoint{12.230208in}{7.691808in}}%
\pgfpathcurveto{\pgfqpoint{12.230208in}{7.702062in}}{\pgfqpoint{12.226134in}{7.711897in}}{\pgfqpoint{12.218883in}{7.719148in}}%
\pgfpathcurveto{\pgfqpoint{12.211632in}{7.726399in}}{\pgfqpoint{12.201797in}{7.730473in}}{\pgfqpoint{12.191543in}{7.730473in}}%
\pgfpathcurveto{\pgfqpoint{12.181289in}{7.730473in}}{\pgfqpoint{12.171453in}{7.726399in}}{\pgfqpoint{12.164202in}{7.719148in}}%
\pgfpathcurveto{\pgfqpoint{12.156952in}{7.711897in}}{\pgfqpoint{12.152878in}{7.702062in}}{\pgfqpoint{12.152878in}{7.691808in}}%
\pgfpathcurveto{\pgfqpoint{12.152878in}{7.681553in}}{\pgfqpoint{12.156952in}{7.671718in}}{\pgfqpoint{12.164202in}{7.664467in}}%
\pgfpathcurveto{\pgfqpoint{12.171453in}{7.657217in}}{\pgfqpoint{12.181289in}{7.653143in}}{\pgfqpoint{12.191543in}{7.653143in}}%
\pgfpathlineto{\pgfqpoint{12.191543in}{7.653143in}}%
\pgfpathclose%
\pgfusepath{stroke}%
\end{pgfscope}%
\begin{pgfscope}%
\pgfpathrectangle{\pgfqpoint{8.036091in}{4.050000in}}{\pgfqpoint{5.400000in}{5.400000in}}%
\pgfusepath{clip}%
\pgfsetbuttcap%
\pgfsetroundjoin%
\pgfsetlinewidth{1.003750pt}%
\definecolor{currentstroke}{rgb}{0.541176,0.380392,0.658824}%
\pgfsetstrokecolor{currentstroke}%
\pgfsetstrokeopacity{0.900000}%
\pgfsetdash{}{0pt}%
\pgfpathmoveto{\pgfqpoint{10.765025in}{5.816990in}}%
\pgfpathcurveto{\pgfqpoint{10.775279in}{5.816990in}}{\pgfqpoint{10.785114in}{5.821064in}}{\pgfqpoint{10.792365in}{5.828315in}}%
\pgfpathcurveto{\pgfqpoint{10.799616in}{5.835566in}}{\pgfqpoint{10.803690in}{5.845401in}}{\pgfqpoint{10.803690in}{5.855655in}}%
\pgfpathcurveto{\pgfqpoint{10.803690in}{5.865909in}}{\pgfqpoint{10.799616in}{5.875745in}}{\pgfqpoint{10.792365in}{5.882996in}}%
\pgfpathcurveto{\pgfqpoint{10.785114in}{5.890246in}}{\pgfqpoint{10.775279in}{5.894320in}}{\pgfqpoint{10.765025in}{5.894320in}}%
\pgfpathcurveto{\pgfqpoint{10.754770in}{5.894320in}}{\pgfqpoint{10.744935in}{5.890246in}}{\pgfqpoint{10.737684in}{5.882996in}}%
\pgfpathcurveto{\pgfqpoint{10.730434in}{5.875745in}}{\pgfqpoint{10.726360in}{5.865909in}}{\pgfqpoint{10.726360in}{5.855655in}}%
\pgfpathcurveto{\pgfqpoint{10.726360in}{5.845401in}}{\pgfqpoint{10.730434in}{5.835566in}}{\pgfqpoint{10.737684in}{5.828315in}}%
\pgfpathcurveto{\pgfqpoint{10.744935in}{5.821064in}}{\pgfqpoint{10.754770in}{5.816990in}}{\pgfqpoint{10.765025in}{5.816990in}}%
\pgfpathlineto{\pgfqpoint{10.765025in}{5.816990in}}%
\pgfpathclose%
\pgfusepath{stroke}%
\end{pgfscope}%
\begin{pgfscope}%
\pgfpathrectangle{\pgfqpoint{8.036091in}{4.050000in}}{\pgfqpoint{5.400000in}{5.400000in}}%
\pgfusepath{clip}%
\pgfsetbuttcap%
\pgfsetroundjoin%
\pgfsetlinewidth{1.003750pt}%
\definecolor{currentstroke}{rgb}{0.541176,0.380392,0.658824}%
\pgfsetstrokecolor{currentstroke}%
\pgfsetstrokeopacity{0.900000}%
\pgfsetdash{}{0pt}%
\pgfpathmoveto{\pgfqpoint{12.347317in}{8.682643in}}%
\pgfpathcurveto{\pgfqpoint{12.357571in}{8.682643in}}{\pgfqpoint{12.367406in}{8.686717in}}{\pgfqpoint{12.374657in}{8.693968in}}%
\pgfpathcurveto{\pgfqpoint{12.381908in}{8.701219in}}{\pgfqpoint{12.385982in}{8.711054in}}{\pgfqpoint{12.385982in}{8.721308in}}%
\pgfpathcurveto{\pgfqpoint{12.385982in}{8.731562in}}{\pgfqpoint{12.381908in}{8.741398in}}{\pgfqpoint{12.374657in}{8.748649in}}%
\pgfpathcurveto{\pgfqpoint{12.367406in}{8.755899in}}{\pgfqpoint{12.357571in}{8.759973in}}{\pgfqpoint{12.347317in}{8.759973in}}%
\pgfpathcurveto{\pgfqpoint{12.337062in}{8.759973in}}{\pgfqpoint{12.327227in}{8.755899in}}{\pgfqpoint{12.319976in}{8.748649in}}%
\pgfpathcurveto{\pgfqpoint{12.312726in}{8.741398in}}{\pgfqpoint{12.308652in}{8.731562in}}{\pgfqpoint{12.308652in}{8.721308in}}%
\pgfpathcurveto{\pgfqpoint{12.308652in}{8.711054in}}{\pgfqpoint{12.312726in}{8.701219in}}{\pgfqpoint{12.319976in}{8.693968in}}%
\pgfpathcurveto{\pgfqpoint{12.327227in}{8.686717in}}{\pgfqpoint{12.337062in}{8.682643in}}{\pgfqpoint{12.347317in}{8.682643in}}%
\pgfpathlineto{\pgfqpoint{12.347317in}{8.682643in}}%
\pgfpathclose%
\pgfusepath{stroke}%
\end{pgfscope}%
\begin{pgfscope}%
\pgfpathrectangle{\pgfqpoint{8.036091in}{4.050000in}}{\pgfqpoint{5.400000in}{5.400000in}}%
\pgfusepath{clip}%
\pgfsetbuttcap%
\pgfsetroundjoin%
\pgfsetlinewidth{1.003750pt}%
\definecolor{currentstroke}{rgb}{0.541176,0.380392,0.658824}%
\pgfsetstrokecolor{currentstroke}%
\pgfsetstrokeopacity{0.900000}%
\pgfsetdash{}{0pt}%
\pgfpathmoveto{\pgfqpoint{12.186085in}{8.664545in}}%
\pgfpathcurveto{\pgfqpoint{12.196339in}{8.664545in}}{\pgfqpoint{12.206174in}{8.668619in}}{\pgfqpoint{12.213425in}{8.675870in}}%
\pgfpathcurveto{\pgfqpoint{12.220676in}{8.683120in}}{\pgfqpoint{12.224750in}{8.692956in}}{\pgfqpoint{12.224750in}{8.703210in}}%
\pgfpathcurveto{\pgfqpoint{12.224750in}{8.713464in}}{\pgfqpoint{12.220676in}{8.723299in}}{\pgfqpoint{12.213425in}{8.730550in}}%
\pgfpathcurveto{\pgfqpoint{12.206174in}{8.737801in}}{\pgfqpoint{12.196339in}{8.741875in}}{\pgfqpoint{12.186085in}{8.741875in}}%
\pgfpathcurveto{\pgfqpoint{12.175830in}{8.741875in}}{\pgfqpoint{12.165995in}{8.737801in}}{\pgfqpoint{12.158744in}{8.730550in}}%
\pgfpathcurveto{\pgfqpoint{12.151494in}{8.723299in}}{\pgfqpoint{12.147420in}{8.713464in}}{\pgfqpoint{12.147420in}{8.703210in}}%
\pgfpathcurveto{\pgfqpoint{12.147420in}{8.692956in}}{\pgfqpoint{12.151494in}{8.683120in}}{\pgfqpoint{12.158744in}{8.675870in}}%
\pgfpathcurveto{\pgfqpoint{12.165995in}{8.668619in}}{\pgfqpoint{12.175830in}{8.664545in}}{\pgfqpoint{12.186085in}{8.664545in}}%
\pgfpathlineto{\pgfqpoint{12.186085in}{8.664545in}}%
\pgfpathclose%
\pgfusepath{stroke}%
\end{pgfscope}%
\begin{pgfscope}%
\pgfpathrectangle{\pgfqpoint{8.036091in}{4.050000in}}{\pgfqpoint{5.400000in}{5.400000in}}%
\pgfusepath{clip}%
\pgfsetbuttcap%
\pgfsetroundjoin%
\pgfsetlinewidth{1.003750pt}%
\definecolor{currentstroke}{rgb}{0.541176,0.380392,0.658824}%
\pgfsetstrokecolor{currentstroke}%
\pgfsetstrokeopacity{0.900000}%
\pgfsetdash{}{0pt}%
\pgfpathmoveto{\pgfqpoint{11.843352in}{7.654500in}}%
\pgfpathcurveto{\pgfqpoint{11.853606in}{7.654500in}}{\pgfqpoint{11.863442in}{7.658574in}}{\pgfqpoint{11.870692in}{7.665825in}}%
\pgfpathcurveto{\pgfqpoint{11.877943in}{7.673076in}}{\pgfqpoint{11.882017in}{7.682911in}}{\pgfqpoint{11.882017in}{7.693165in}}%
\pgfpathcurveto{\pgfqpoint{11.882017in}{7.703419in}}{\pgfqpoint{11.877943in}{7.713255in}}{\pgfqpoint{11.870692in}{7.720505in}}%
\pgfpathcurveto{\pgfqpoint{11.863442in}{7.727756in}}{\pgfqpoint{11.853606in}{7.731830in}}{\pgfqpoint{11.843352in}{7.731830in}}%
\pgfpathcurveto{\pgfqpoint{11.833098in}{7.731830in}}{\pgfqpoint{11.823263in}{7.727756in}}{\pgfqpoint{11.816012in}{7.720505in}}%
\pgfpathcurveto{\pgfqpoint{11.808761in}{7.713255in}}{\pgfqpoint{11.804687in}{7.703419in}}{\pgfqpoint{11.804687in}{7.693165in}}%
\pgfpathcurveto{\pgfqpoint{11.804687in}{7.682911in}}{\pgfqpoint{11.808761in}{7.673076in}}{\pgfqpoint{11.816012in}{7.665825in}}%
\pgfpathcurveto{\pgfqpoint{11.823263in}{7.658574in}}{\pgfqpoint{11.833098in}{7.654500in}}{\pgfqpoint{11.843352in}{7.654500in}}%
\pgfpathlineto{\pgfqpoint{11.843352in}{7.654500in}}%
\pgfpathclose%
\pgfusepath{stroke}%
\end{pgfscope}%
\begin{pgfscope}%
\pgfpathrectangle{\pgfqpoint{8.036091in}{4.050000in}}{\pgfqpoint{5.400000in}{5.400000in}}%
\pgfusepath{clip}%
\pgfsetbuttcap%
\pgfsetroundjoin%
\pgfsetlinewidth{1.003750pt}%
\definecolor{currentstroke}{rgb}{0.541176,0.380392,0.658824}%
\pgfsetstrokecolor{currentstroke}%
\pgfsetstrokeopacity{0.900000}%
\pgfsetdash{}{0pt}%
\pgfpathmoveto{\pgfqpoint{11.505682in}{7.767102in}}%
\pgfpathcurveto{\pgfqpoint{11.515936in}{7.767102in}}{\pgfqpoint{11.525772in}{7.771176in}}{\pgfqpoint{11.533022in}{7.778426in}}%
\pgfpathcurveto{\pgfqpoint{11.540273in}{7.785677in}}{\pgfqpoint{11.544347in}{7.795513in}}{\pgfqpoint{11.544347in}{7.805767in}}%
\pgfpathcurveto{\pgfqpoint{11.544347in}{7.816021in}}{\pgfqpoint{11.540273in}{7.825856in}}{\pgfqpoint{11.533022in}{7.833107in}}%
\pgfpathcurveto{\pgfqpoint{11.525772in}{7.840358in}}{\pgfqpoint{11.515936in}{7.844432in}}{\pgfqpoint{11.505682in}{7.844432in}}%
\pgfpathcurveto{\pgfqpoint{11.495428in}{7.844432in}}{\pgfqpoint{11.485592in}{7.840358in}}{\pgfqpoint{11.478342in}{7.833107in}}%
\pgfpathcurveto{\pgfqpoint{11.471091in}{7.825856in}}{\pgfqpoint{11.467017in}{7.816021in}}{\pgfqpoint{11.467017in}{7.805767in}}%
\pgfpathcurveto{\pgfqpoint{11.467017in}{7.795513in}}{\pgfqpoint{11.471091in}{7.785677in}}{\pgfqpoint{11.478342in}{7.778426in}}%
\pgfpathcurveto{\pgfqpoint{11.485592in}{7.771176in}}{\pgfqpoint{11.495428in}{7.767102in}}{\pgfqpoint{11.505682in}{7.767102in}}%
\pgfpathlineto{\pgfqpoint{11.505682in}{7.767102in}}%
\pgfpathclose%
\pgfusepath{stroke}%
\end{pgfscope}%
\begin{pgfscope}%
\pgfpathrectangle{\pgfqpoint{8.036091in}{4.050000in}}{\pgfqpoint{5.400000in}{5.400000in}}%
\pgfusepath{clip}%
\pgfsetbuttcap%
\pgfsetroundjoin%
\pgfsetlinewidth{1.003750pt}%
\definecolor{currentstroke}{rgb}{0.541176,0.380392,0.658824}%
\pgfsetstrokecolor{currentstroke}%
\pgfsetstrokeopacity{0.900000}%
\pgfsetdash{}{0pt}%
\pgfpathmoveto{\pgfqpoint{12.216554in}{8.364018in}}%
\pgfpathcurveto{\pgfqpoint{12.226808in}{8.364018in}}{\pgfqpoint{12.236644in}{8.368092in}}{\pgfqpoint{12.243894in}{8.375343in}}%
\pgfpathcurveto{\pgfqpoint{12.251145in}{8.382594in}}{\pgfqpoint{12.255219in}{8.392429in}}{\pgfqpoint{12.255219in}{8.402683in}}%
\pgfpathcurveto{\pgfqpoint{12.255219in}{8.412937in}}{\pgfqpoint{12.251145in}{8.422773in}}{\pgfqpoint{12.243894in}{8.430023in}}%
\pgfpathcurveto{\pgfqpoint{12.236644in}{8.437274in}}{\pgfqpoint{12.226808in}{8.441348in}}{\pgfqpoint{12.216554in}{8.441348in}}%
\pgfpathcurveto{\pgfqpoint{12.206300in}{8.441348in}}{\pgfqpoint{12.196464in}{8.437274in}}{\pgfqpoint{12.189214in}{8.430023in}}%
\pgfpathcurveto{\pgfqpoint{12.181963in}{8.422773in}}{\pgfqpoint{12.177889in}{8.412937in}}{\pgfqpoint{12.177889in}{8.402683in}}%
\pgfpathcurveto{\pgfqpoint{12.177889in}{8.392429in}}{\pgfqpoint{12.181963in}{8.382594in}}{\pgfqpoint{12.189214in}{8.375343in}}%
\pgfpathcurveto{\pgfqpoint{12.196464in}{8.368092in}}{\pgfqpoint{12.206300in}{8.364018in}}{\pgfqpoint{12.216554in}{8.364018in}}%
\pgfpathlineto{\pgfqpoint{12.216554in}{8.364018in}}%
\pgfpathclose%
\pgfusepath{stroke}%
\end{pgfscope}%
\begin{pgfscope}%
\pgfpathrectangle{\pgfqpoint{8.036091in}{4.050000in}}{\pgfqpoint{5.400000in}{5.400000in}}%
\pgfusepath{clip}%
\pgfsetbuttcap%
\pgfsetroundjoin%
\pgfsetlinewidth{1.003750pt}%
\definecolor{currentstroke}{rgb}{0.541176,0.380392,0.658824}%
\pgfsetstrokecolor{currentstroke}%
\pgfsetstrokeopacity{0.900000}%
\pgfsetdash{}{0pt}%
\pgfpathmoveto{\pgfqpoint{11.095992in}{6.967041in}}%
\pgfpathcurveto{\pgfqpoint{11.106246in}{6.967041in}}{\pgfqpoint{11.116082in}{6.971115in}}{\pgfqpoint{11.123333in}{6.978365in}}%
\pgfpathcurveto{\pgfqpoint{11.130583in}{6.985616in}}{\pgfqpoint{11.134657in}{6.995452in}}{\pgfqpoint{11.134657in}{7.005706in}}%
\pgfpathcurveto{\pgfqpoint{11.134657in}{7.015960in}}{\pgfqpoint{11.130583in}{7.025795in}}{\pgfqpoint{11.123333in}{7.033046in}}%
\pgfpathcurveto{\pgfqpoint{11.116082in}{7.040297in}}{\pgfqpoint{11.106246in}{7.044371in}}{\pgfqpoint{11.095992in}{7.044371in}}%
\pgfpathcurveto{\pgfqpoint{11.085738in}{7.044371in}}{\pgfqpoint{11.075903in}{7.040297in}}{\pgfqpoint{11.068652in}{7.033046in}}%
\pgfpathcurveto{\pgfqpoint{11.061401in}{7.025795in}}{\pgfqpoint{11.057327in}{7.015960in}}{\pgfqpoint{11.057327in}{7.005706in}}%
\pgfpathcurveto{\pgfqpoint{11.057327in}{6.995452in}}{\pgfqpoint{11.061401in}{6.985616in}}{\pgfqpoint{11.068652in}{6.978365in}}%
\pgfpathcurveto{\pgfqpoint{11.075903in}{6.971115in}}{\pgfqpoint{11.085738in}{6.967041in}}{\pgfqpoint{11.095992in}{6.967041in}}%
\pgfpathlineto{\pgfqpoint{11.095992in}{6.967041in}}%
\pgfpathclose%
\pgfusepath{stroke}%
\end{pgfscope}%
\begin{pgfscope}%
\pgfpathrectangle{\pgfqpoint{8.036091in}{4.050000in}}{\pgfqpoint{5.400000in}{5.400000in}}%
\pgfusepath{clip}%
\pgfsetbuttcap%
\pgfsetroundjoin%
\pgfsetlinewidth{1.003750pt}%
\definecolor{currentstroke}{rgb}{0.541176,0.380392,0.658824}%
\pgfsetstrokecolor{currentstroke}%
\pgfsetstrokeopacity{0.900000}%
\pgfsetdash{}{0pt}%
\pgfpathmoveto{\pgfqpoint{11.404140in}{7.886615in}}%
\pgfpathcurveto{\pgfqpoint{11.414394in}{7.886615in}}{\pgfqpoint{11.424229in}{7.890689in}}{\pgfqpoint{11.431480in}{7.897940in}}%
\pgfpathcurveto{\pgfqpoint{11.438731in}{7.905191in}}{\pgfqpoint{11.442805in}{7.915026in}}{\pgfqpoint{11.442805in}{7.925280in}}%
\pgfpathcurveto{\pgfqpoint{11.442805in}{7.935534in}}{\pgfqpoint{11.438731in}{7.945370in}}{\pgfqpoint{11.431480in}{7.952621in}}%
\pgfpathcurveto{\pgfqpoint{11.424229in}{7.959871in}}{\pgfqpoint{11.414394in}{7.963945in}}{\pgfqpoint{11.404140in}{7.963945in}}%
\pgfpathcurveto{\pgfqpoint{11.393886in}{7.963945in}}{\pgfqpoint{11.384050in}{7.959871in}}{\pgfqpoint{11.376800in}{7.952621in}}%
\pgfpathcurveto{\pgfqpoint{11.369549in}{7.945370in}}{\pgfqpoint{11.365475in}{7.935534in}}{\pgfqpoint{11.365475in}{7.925280in}}%
\pgfpathcurveto{\pgfqpoint{11.365475in}{7.915026in}}{\pgfqpoint{11.369549in}{7.905191in}}{\pgfqpoint{11.376800in}{7.897940in}}%
\pgfpathcurveto{\pgfqpoint{11.384050in}{7.890689in}}{\pgfqpoint{11.393886in}{7.886615in}}{\pgfqpoint{11.404140in}{7.886615in}}%
\pgfpathlineto{\pgfqpoint{11.404140in}{7.886615in}}%
\pgfpathclose%
\pgfusepath{stroke}%
\end{pgfscope}%
\begin{pgfscope}%
\pgfpathrectangle{\pgfqpoint{8.036091in}{4.050000in}}{\pgfqpoint{5.400000in}{5.400000in}}%
\pgfusepath{clip}%
\pgfsetbuttcap%
\pgfsetroundjoin%
\pgfsetlinewidth{1.003750pt}%
\definecolor{currentstroke}{rgb}{0.541176,0.380392,0.658824}%
\pgfsetstrokecolor{currentstroke}%
\pgfsetstrokeopacity{0.900000}%
\pgfsetdash{}{0pt}%
\pgfpathmoveto{\pgfqpoint{11.310930in}{6.831036in}}%
\pgfpathcurveto{\pgfqpoint{11.321184in}{6.831036in}}{\pgfqpoint{11.331020in}{6.835110in}}{\pgfqpoint{11.338271in}{6.842361in}}%
\pgfpathcurveto{\pgfqpoint{11.345521in}{6.849611in}}{\pgfqpoint{11.349595in}{6.859447in}}{\pgfqpoint{11.349595in}{6.869701in}}%
\pgfpathcurveto{\pgfqpoint{11.349595in}{6.879955in}}{\pgfqpoint{11.345521in}{6.889791in}}{\pgfqpoint{11.338271in}{6.897041in}}%
\pgfpathcurveto{\pgfqpoint{11.331020in}{6.904292in}}{\pgfqpoint{11.321184in}{6.908366in}}{\pgfqpoint{11.310930in}{6.908366in}}%
\pgfpathcurveto{\pgfqpoint{11.300676in}{6.908366in}}{\pgfqpoint{11.290841in}{6.904292in}}{\pgfqpoint{11.283590in}{6.897041in}}%
\pgfpathcurveto{\pgfqpoint{11.276339in}{6.889791in}}{\pgfqpoint{11.272265in}{6.879955in}}{\pgfqpoint{11.272265in}{6.869701in}}%
\pgfpathcurveto{\pgfqpoint{11.272265in}{6.859447in}}{\pgfqpoint{11.276339in}{6.849611in}}{\pgfqpoint{11.283590in}{6.842361in}}%
\pgfpathcurveto{\pgfqpoint{11.290841in}{6.835110in}}{\pgfqpoint{11.300676in}{6.831036in}}{\pgfqpoint{11.310930in}{6.831036in}}%
\pgfpathlineto{\pgfqpoint{11.310930in}{6.831036in}}%
\pgfpathclose%
\pgfusepath{stroke}%
\end{pgfscope}%
\begin{pgfscope}%
\pgfpathrectangle{\pgfqpoint{8.036091in}{4.050000in}}{\pgfqpoint{5.400000in}{5.400000in}}%
\pgfusepath{clip}%
\pgfsetbuttcap%
\pgfsetroundjoin%
\pgfsetlinewidth{1.003750pt}%
\definecolor{currentstroke}{rgb}{0.541176,0.380392,0.658824}%
\pgfsetstrokecolor{currentstroke}%
\pgfsetstrokeopacity{0.900000}%
\pgfsetdash{}{0pt}%
\pgfpathmoveto{\pgfqpoint{10.168902in}{5.945136in}}%
\pgfpathcurveto{\pgfqpoint{10.179156in}{5.945136in}}{\pgfqpoint{10.188991in}{5.949210in}}{\pgfqpoint{10.196242in}{5.956461in}}%
\pgfpathcurveto{\pgfqpoint{10.203493in}{5.963712in}}{\pgfqpoint{10.207567in}{5.973547in}}{\pgfqpoint{10.207567in}{5.983801in}}%
\pgfpathcurveto{\pgfqpoint{10.207567in}{5.994055in}}{\pgfqpoint{10.203493in}{6.003891in}}{\pgfqpoint{10.196242in}{6.011141in}}%
\pgfpathcurveto{\pgfqpoint{10.188991in}{6.018392in}}{\pgfqpoint{10.179156in}{6.022466in}}{\pgfqpoint{10.168902in}{6.022466in}}%
\pgfpathcurveto{\pgfqpoint{10.158647in}{6.022466in}}{\pgfqpoint{10.148812in}{6.018392in}}{\pgfqpoint{10.141561in}{6.011141in}}%
\pgfpathcurveto{\pgfqpoint{10.134311in}{6.003891in}}{\pgfqpoint{10.130237in}{5.994055in}}{\pgfqpoint{10.130237in}{5.983801in}}%
\pgfpathcurveto{\pgfqpoint{10.130237in}{5.973547in}}{\pgfqpoint{10.134311in}{5.963712in}}{\pgfqpoint{10.141561in}{5.956461in}}%
\pgfpathcurveto{\pgfqpoint{10.148812in}{5.949210in}}{\pgfqpoint{10.158647in}{5.945136in}}{\pgfqpoint{10.168902in}{5.945136in}}%
\pgfpathlineto{\pgfqpoint{10.168902in}{5.945136in}}%
\pgfpathclose%
\pgfusepath{stroke}%
\end{pgfscope}%
\begin{pgfscope}%
\pgfpathrectangle{\pgfqpoint{8.036091in}{4.050000in}}{\pgfqpoint{5.400000in}{5.400000in}}%
\pgfusepath{clip}%
\pgfsetbuttcap%
\pgfsetroundjoin%
\pgfsetlinewidth{1.003750pt}%
\definecolor{currentstroke}{rgb}{0.541176,0.380392,0.658824}%
\pgfsetstrokecolor{currentstroke}%
\pgfsetstrokeopacity{0.900000}%
\pgfsetdash{}{0pt}%
\pgfpathmoveto{\pgfqpoint{10.931663in}{6.382159in}}%
\pgfpathcurveto{\pgfqpoint{10.941917in}{6.382159in}}{\pgfqpoint{10.951752in}{6.386233in}}{\pgfqpoint{10.959003in}{6.393484in}}%
\pgfpathcurveto{\pgfqpoint{10.966254in}{6.400734in}}{\pgfqpoint{10.970328in}{6.410570in}}{\pgfqpoint{10.970328in}{6.420824in}}%
\pgfpathcurveto{\pgfqpoint{10.970328in}{6.431078in}}{\pgfqpoint{10.966254in}{6.440913in}}{\pgfqpoint{10.959003in}{6.448164in}}%
\pgfpathcurveto{\pgfqpoint{10.951752in}{6.455415in}}{\pgfqpoint{10.941917in}{6.459489in}}{\pgfqpoint{10.931663in}{6.459489in}}%
\pgfpathcurveto{\pgfqpoint{10.921409in}{6.459489in}}{\pgfqpoint{10.911573in}{6.455415in}}{\pgfqpoint{10.904323in}{6.448164in}}%
\pgfpathcurveto{\pgfqpoint{10.897072in}{6.440913in}}{\pgfqpoint{10.892998in}{6.431078in}}{\pgfqpoint{10.892998in}{6.420824in}}%
\pgfpathcurveto{\pgfqpoint{10.892998in}{6.410570in}}{\pgfqpoint{10.897072in}{6.400734in}}{\pgfqpoint{10.904323in}{6.393484in}}%
\pgfpathcurveto{\pgfqpoint{10.911573in}{6.386233in}}{\pgfqpoint{10.921409in}{6.382159in}}{\pgfqpoint{10.931663in}{6.382159in}}%
\pgfpathlineto{\pgfqpoint{10.931663in}{6.382159in}}%
\pgfpathclose%
\pgfusepath{stroke}%
\end{pgfscope}%
\begin{pgfscope}%
\pgfpathrectangle{\pgfqpoint{8.036091in}{4.050000in}}{\pgfqpoint{5.400000in}{5.400000in}}%
\pgfusepath{clip}%
\pgfsetbuttcap%
\pgfsetroundjoin%
\pgfsetlinewidth{1.003750pt}%
\definecolor{currentstroke}{rgb}{0.541176,0.380392,0.658824}%
\pgfsetstrokecolor{currentstroke}%
\pgfsetstrokeopacity{0.900000}%
\pgfsetdash{}{0pt}%
\pgfpathmoveto{\pgfqpoint{11.570475in}{7.509021in}}%
\pgfpathcurveto{\pgfqpoint{11.580729in}{7.509021in}}{\pgfqpoint{11.590564in}{7.513095in}}{\pgfqpoint{11.597815in}{7.520346in}}%
\pgfpathcurveto{\pgfqpoint{11.605066in}{7.527596in}}{\pgfqpoint{11.609140in}{7.537432in}}{\pgfqpoint{11.609140in}{7.547686in}}%
\pgfpathcurveto{\pgfqpoint{11.609140in}{7.557940in}}{\pgfqpoint{11.605066in}{7.567776in}}{\pgfqpoint{11.597815in}{7.575026in}}%
\pgfpathcurveto{\pgfqpoint{11.590564in}{7.582277in}}{\pgfqpoint{11.580729in}{7.586351in}}{\pgfqpoint{11.570475in}{7.586351in}}%
\pgfpathcurveto{\pgfqpoint{11.560221in}{7.586351in}}{\pgfqpoint{11.550385in}{7.582277in}}{\pgfqpoint{11.543134in}{7.575026in}}%
\pgfpathcurveto{\pgfqpoint{11.535884in}{7.567776in}}{\pgfqpoint{11.531810in}{7.557940in}}{\pgfqpoint{11.531810in}{7.547686in}}%
\pgfpathcurveto{\pgfqpoint{11.531810in}{7.537432in}}{\pgfqpoint{11.535884in}{7.527596in}}{\pgfqpoint{11.543134in}{7.520346in}}%
\pgfpathcurveto{\pgfqpoint{11.550385in}{7.513095in}}{\pgfqpoint{11.560221in}{7.509021in}}{\pgfqpoint{11.570475in}{7.509021in}}%
\pgfpathlineto{\pgfqpoint{11.570475in}{7.509021in}}%
\pgfpathclose%
\pgfusepath{stroke}%
\end{pgfscope}%
\begin{pgfscope}%
\pgfpathrectangle{\pgfqpoint{8.036091in}{4.050000in}}{\pgfqpoint{5.400000in}{5.400000in}}%
\pgfusepath{clip}%
\pgfsetbuttcap%
\pgfsetroundjoin%
\pgfsetlinewidth{1.003750pt}%
\definecolor{currentstroke}{rgb}{0.541176,0.380392,0.658824}%
\pgfsetstrokecolor{currentstroke}%
\pgfsetstrokeopacity{0.900000}%
\pgfsetdash{}{0pt}%
\pgfpathmoveto{\pgfqpoint{10.555046in}{5.944036in}}%
\pgfpathcurveto{\pgfqpoint{10.565300in}{5.944036in}}{\pgfqpoint{10.575136in}{5.948110in}}{\pgfqpoint{10.582386in}{5.955361in}}%
\pgfpathcurveto{\pgfqpoint{10.589637in}{5.962612in}}{\pgfqpoint{10.593711in}{5.972447in}}{\pgfqpoint{10.593711in}{5.982701in}}%
\pgfpathcurveto{\pgfqpoint{10.593711in}{5.992955in}}{\pgfqpoint{10.589637in}{6.002791in}}{\pgfqpoint{10.582386in}{6.010041in}}%
\pgfpathcurveto{\pgfqpoint{10.575136in}{6.017292in}}{\pgfqpoint{10.565300in}{6.021366in}}{\pgfqpoint{10.555046in}{6.021366in}}%
\pgfpathcurveto{\pgfqpoint{10.544792in}{6.021366in}}{\pgfqpoint{10.534957in}{6.017292in}}{\pgfqpoint{10.527706in}{6.010041in}}%
\pgfpathcurveto{\pgfqpoint{10.520455in}{6.002791in}}{\pgfqpoint{10.516381in}{5.992955in}}{\pgfqpoint{10.516381in}{5.982701in}}%
\pgfpathcurveto{\pgfqpoint{10.516381in}{5.972447in}}{\pgfqpoint{10.520455in}{5.962612in}}{\pgfqpoint{10.527706in}{5.955361in}}%
\pgfpathcurveto{\pgfqpoint{10.534957in}{5.948110in}}{\pgfqpoint{10.544792in}{5.944036in}}{\pgfqpoint{10.555046in}{5.944036in}}%
\pgfpathlineto{\pgfqpoint{10.555046in}{5.944036in}}%
\pgfpathclose%
\pgfusepath{stroke}%
\end{pgfscope}%
\begin{pgfscope}%
\pgfpathrectangle{\pgfqpoint{8.036091in}{4.050000in}}{\pgfqpoint{5.400000in}{5.400000in}}%
\pgfusepath{clip}%
\pgfsetbuttcap%
\pgfsetroundjoin%
\pgfsetlinewidth{1.003750pt}%
\definecolor{currentstroke}{rgb}{0.541176,0.380392,0.658824}%
\pgfsetstrokecolor{currentstroke}%
\pgfsetstrokeopacity{0.900000}%
\pgfsetdash{}{0pt}%
\pgfpathmoveto{\pgfqpoint{12.417370in}{8.763497in}}%
\pgfpathcurveto{\pgfqpoint{12.427625in}{8.763497in}}{\pgfqpoint{12.437460in}{8.767571in}}{\pgfqpoint{12.444711in}{8.774822in}}%
\pgfpathcurveto{\pgfqpoint{12.451962in}{8.782073in}}{\pgfqpoint{12.456035in}{8.791908in}}{\pgfqpoint{12.456035in}{8.802162in}}%
\pgfpathcurveto{\pgfqpoint{12.456035in}{8.812417in}}{\pgfqpoint{12.451962in}{8.822252in}}{\pgfqpoint{12.444711in}{8.829503in}}%
\pgfpathcurveto{\pgfqpoint{12.437460in}{8.836753in}}{\pgfqpoint{12.427625in}{8.840827in}}{\pgfqpoint{12.417370in}{8.840827in}}%
\pgfpathcurveto{\pgfqpoint{12.407116in}{8.840827in}}{\pgfqpoint{12.397281in}{8.836753in}}{\pgfqpoint{12.390030in}{8.829503in}}%
\pgfpathcurveto{\pgfqpoint{12.382779in}{8.822252in}}{\pgfqpoint{12.378705in}{8.812417in}}{\pgfqpoint{12.378705in}{8.802162in}}%
\pgfpathcurveto{\pgfqpoint{12.378705in}{8.791908in}}{\pgfqpoint{12.382779in}{8.782073in}}{\pgfqpoint{12.390030in}{8.774822in}}%
\pgfpathcurveto{\pgfqpoint{12.397281in}{8.767571in}}{\pgfqpoint{12.407116in}{8.763497in}}{\pgfqpoint{12.417370in}{8.763497in}}%
\pgfpathlineto{\pgfqpoint{12.417370in}{8.763497in}}%
\pgfpathclose%
\pgfusepath{stroke}%
\end{pgfscope}%
\begin{pgfscope}%
\pgfpathrectangle{\pgfqpoint{8.036091in}{4.050000in}}{\pgfqpoint{5.400000in}{5.400000in}}%
\pgfusepath{clip}%
\pgfsetbuttcap%
\pgfsetroundjoin%
\pgfsetlinewidth{1.003750pt}%
\definecolor{currentstroke}{rgb}{0.541176,0.380392,0.658824}%
\pgfsetstrokecolor{currentstroke}%
\pgfsetstrokeopacity{0.900000}%
\pgfsetdash{}{0pt}%
\pgfpathmoveto{\pgfqpoint{11.752708in}{7.017656in}}%
\pgfpathcurveto{\pgfqpoint{11.762962in}{7.017656in}}{\pgfqpoint{11.772798in}{7.021730in}}{\pgfqpoint{11.780048in}{7.028980in}}%
\pgfpathcurveto{\pgfqpoint{11.787299in}{7.036231in}}{\pgfqpoint{11.791373in}{7.046067in}}{\pgfqpoint{11.791373in}{7.056321in}}%
\pgfpathcurveto{\pgfqpoint{11.791373in}{7.066575in}}{\pgfqpoint{11.787299in}{7.076410in}}{\pgfqpoint{11.780048in}{7.083661in}}%
\pgfpathcurveto{\pgfqpoint{11.772798in}{7.090912in}}{\pgfqpoint{11.762962in}{7.094986in}}{\pgfqpoint{11.752708in}{7.094986in}}%
\pgfpathcurveto{\pgfqpoint{11.742454in}{7.094986in}}{\pgfqpoint{11.732618in}{7.090912in}}{\pgfqpoint{11.725368in}{7.083661in}}%
\pgfpathcurveto{\pgfqpoint{11.718117in}{7.076410in}}{\pgfqpoint{11.714043in}{7.066575in}}{\pgfqpoint{11.714043in}{7.056321in}}%
\pgfpathcurveto{\pgfqpoint{11.714043in}{7.046067in}}{\pgfqpoint{11.718117in}{7.036231in}}{\pgfqpoint{11.725368in}{7.028980in}}%
\pgfpathcurveto{\pgfqpoint{11.732618in}{7.021730in}}{\pgfqpoint{11.742454in}{7.017656in}}{\pgfqpoint{11.752708in}{7.017656in}}%
\pgfpathlineto{\pgfqpoint{11.752708in}{7.017656in}}%
\pgfpathclose%
\pgfusepath{stroke}%
\end{pgfscope}%
\begin{pgfscope}%
\pgfpathrectangle{\pgfqpoint{8.036091in}{4.050000in}}{\pgfqpoint{5.400000in}{5.400000in}}%
\pgfusepath{clip}%
\pgfsetbuttcap%
\pgfsetroundjoin%
\pgfsetlinewidth{1.003750pt}%
\definecolor{currentstroke}{rgb}{0.541176,0.380392,0.658824}%
\pgfsetstrokecolor{currentstroke}%
\pgfsetstrokeopacity{0.900000}%
\pgfsetdash{}{0pt}%
\pgfpathmoveto{\pgfqpoint{12.411683in}{7.462755in}}%
\pgfpathcurveto{\pgfqpoint{12.421937in}{7.462755in}}{\pgfqpoint{12.431773in}{7.466829in}}{\pgfqpoint{12.439024in}{7.474079in}}%
\pgfpathcurveto{\pgfqpoint{12.446274in}{7.481330in}}{\pgfqpoint{12.450348in}{7.491166in}}{\pgfqpoint{12.450348in}{7.501420in}}%
\pgfpathcurveto{\pgfqpoint{12.450348in}{7.511674in}}{\pgfqpoint{12.446274in}{7.521509in}}{\pgfqpoint{12.439024in}{7.528760in}}%
\pgfpathcurveto{\pgfqpoint{12.431773in}{7.536011in}}{\pgfqpoint{12.421937in}{7.540085in}}{\pgfqpoint{12.411683in}{7.540085in}}%
\pgfpathcurveto{\pgfqpoint{12.401429in}{7.540085in}}{\pgfqpoint{12.391594in}{7.536011in}}{\pgfqpoint{12.384343in}{7.528760in}}%
\pgfpathcurveto{\pgfqpoint{12.377092in}{7.521509in}}{\pgfqpoint{12.373018in}{7.511674in}}{\pgfqpoint{12.373018in}{7.501420in}}%
\pgfpathcurveto{\pgfqpoint{12.373018in}{7.491166in}}{\pgfqpoint{12.377092in}{7.481330in}}{\pgfqpoint{12.384343in}{7.474079in}}%
\pgfpathcurveto{\pgfqpoint{12.391594in}{7.466829in}}{\pgfqpoint{12.401429in}{7.462755in}}{\pgfqpoint{12.411683in}{7.462755in}}%
\pgfpathlineto{\pgfqpoint{12.411683in}{7.462755in}}%
\pgfpathclose%
\pgfusepath{stroke}%
\end{pgfscope}%
\begin{pgfscope}%
\pgfpathrectangle{\pgfqpoint{8.036091in}{4.050000in}}{\pgfqpoint{5.400000in}{5.400000in}}%
\pgfusepath{clip}%
\pgfsetbuttcap%
\pgfsetroundjoin%
\pgfsetlinewidth{1.003750pt}%
\definecolor{currentstroke}{rgb}{0.541176,0.380392,0.658824}%
\pgfsetstrokecolor{currentstroke}%
\pgfsetstrokeopacity{0.900000}%
\pgfsetdash{}{0pt}%
\pgfpathmoveto{\pgfqpoint{11.395694in}{8.213282in}}%
\pgfpathcurveto{\pgfqpoint{11.405948in}{8.213282in}}{\pgfqpoint{11.415784in}{8.217356in}}{\pgfqpoint{11.423035in}{8.224606in}}%
\pgfpathcurveto{\pgfqpoint{11.430285in}{8.231857in}}{\pgfqpoint{11.434359in}{8.241693in}}{\pgfqpoint{11.434359in}{8.251947in}}%
\pgfpathcurveto{\pgfqpoint{11.434359in}{8.262201in}}{\pgfqpoint{11.430285in}{8.272036in}}{\pgfqpoint{11.423035in}{8.279287in}}%
\pgfpathcurveto{\pgfqpoint{11.415784in}{8.286538in}}{\pgfqpoint{11.405948in}{8.290612in}}{\pgfqpoint{11.395694in}{8.290612in}}%
\pgfpathcurveto{\pgfqpoint{11.385440in}{8.290612in}}{\pgfqpoint{11.375605in}{8.286538in}}{\pgfqpoint{11.368354in}{8.279287in}}%
\pgfpathcurveto{\pgfqpoint{11.361103in}{8.272036in}}{\pgfqpoint{11.357029in}{8.262201in}}{\pgfqpoint{11.357029in}{8.251947in}}%
\pgfpathcurveto{\pgfqpoint{11.357029in}{8.241693in}}{\pgfqpoint{11.361103in}{8.231857in}}{\pgfqpoint{11.368354in}{8.224606in}}%
\pgfpathcurveto{\pgfqpoint{11.375605in}{8.217356in}}{\pgfqpoint{11.385440in}{8.213282in}}{\pgfqpoint{11.395694in}{8.213282in}}%
\pgfpathlineto{\pgfqpoint{11.395694in}{8.213282in}}%
\pgfpathclose%
\pgfusepath{stroke}%
\end{pgfscope}%
\begin{pgfscope}%
\pgfpathrectangle{\pgfqpoint{8.036091in}{4.050000in}}{\pgfqpoint{5.400000in}{5.400000in}}%
\pgfusepath{clip}%
\pgfsetbuttcap%
\pgfsetroundjoin%
\pgfsetlinewidth{1.003750pt}%
\definecolor{currentstroke}{rgb}{0.541176,0.380392,0.658824}%
\pgfsetstrokecolor{currentstroke}%
\pgfsetstrokeopacity{0.900000}%
\pgfsetdash{}{0pt}%
\pgfpathmoveto{\pgfqpoint{10.372596in}{6.603283in}}%
\pgfpathcurveto{\pgfqpoint{10.382850in}{6.603283in}}{\pgfqpoint{10.392686in}{6.607357in}}{\pgfqpoint{10.399937in}{6.614608in}}%
\pgfpathcurveto{\pgfqpoint{10.407187in}{6.621858in}}{\pgfqpoint{10.411261in}{6.631694in}}{\pgfqpoint{10.411261in}{6.641948in}}%
\pgfpathcurveto{\pgfqpoint{10.411261in}{6.652202in}}{\pgfqpoint{10.407187in}{6.662037in}}{\pgfqpoint{10.399937in}{6.669288in}}%
\pgfpathcurveto{\pgfqpoint{10.392686in}{6.676539in}}{\pgfqpoint{10.382850in}{6.680613in}}{\pgfqpoint{10.372596in}{6.680613in}}%
\pgfpathcurveto{\pgfqpoint{10.362342in}{6.680613in}}{\pgfqpoint{10.352507in}{6.676539in}}{\pgfqpoint{10.345256in}{6.669288in}}%
\pgfpathcurveto{\pgfqpoint{10.338005in}{6.662037in}}{\pgfqpoint{10.333931in}{6.652202in}}{\pgfqpoint{10.333931in}{6.641948in}}%
\pgfpathcurveto{\pgfqpoint{10.333931in}{6.631694in}}{\pgfqpoint{10.338005in}{6.621858in}}{\pgfqpoint{10.345256in}{6.614608in}}%
\pgfpathcurveto{\pgfqpoint{10.352507in}{6.607357in}}{\pgfqpoint{10.362342in}{6.603283in}}{\pgfqpoint{10.372596in}{6.603283in}}%
\pgfpathlineto{\pgfqpoint{10.372596in}{6.603283in}}%
\pgfpathclose%
\pgfusepath{stroke}%
\end{pgfscope}%
\begin{pgfscope}%
\pgfpathrectangle{\pgfqpoint{8.036091in}{4.050000in}}{\pgfqpoint{5.400000in}{5.400000in}}%
\pgfusepath{clip}%
\pgfsetbuttcap%
\pgfsetroundjoin%
\pgfsetlinewidth{1.003750pt}%
\definecolor{currentstroke}{rgb}{0.541176,0.380392,0.658824}%
\pgfsetstrokecolor{currentstroke}%
\pgfsetstrokeopacity{0.900000}%
\pgfsetdash{}{0pt}%
\pgfpathmoveto{\pgfqpoint{12.278437in}{7.973763in}}%
\pgfpathcurveto{\pgfqpoint{12.288692in}{7.973763in}}{\pgfqpoint{12.298527in}{7.977837in}}{\pgfqpoint{12.305778in}{7.985087in}}%
\pgfpathcurveto{\pgfqpoint{12.313028in}{7.992338in}}{\pgfqpoint{12.317102in}{8.002173in}}{\pgfqpoint{12.317102in}{8.012428in}}%
\pgfpathcurveto{\pgfqpoint{12.317102in}{8.022682in}}{\pgfqpoint{12.313028in}{8.032517in}}{\pgfqpoint{12.305778in}{8.039768in}}%
\pgfpathcurveto{\pgfqpoint{12.298527in}{8.047019in}}{\pgfqpoint{12.288692in}{8.051093in}}{\pgfqpoint{12.278437in}{8.051093in}}%
\pgfpathcurveto{\pgfqpoint{12.268183in}{8.051093in}}{\pgfqpoint{12.258348in}{8.047019in}}{\pgfqpoint{12.251097in}{8.039768in}}%
\pgfpathcurveto{\pgfqpoint{12.243846in}{8.032517in}}{\pgfqpoint{12.239772in}{8.022682in}}{\pgfqpoint{12.239772in}{8.012428in}}%
\pgfpathcurveto{\pgfqpoint{12.239772in}{8.002173in}}{\pgfqpoint{12.243846in}{7.992338in}}{\pgfqpoint{12.251097in}{7.985087in}}%
\pgfpathcurveto{\pgfqpoint{12.258348in}{7.977837in}}{\pgfqpoint{12.268183in}{7.973763in}}{\pgfqpoint{12.278437in}{7.973763in}}%
\pgfpathlineto{\pgfqpoint{12.278437in}{7.973763in}}%
\pgfpathclose%
\pgfusepath{stroke}%
\end{pgfscope}%
\begin{pgfscope}%
\pgfpathrectangle{\pgfqpoint{8.036091in}{4.050000in}}{\pgfqpoint{5.400000in}{5.400000in}}%
\pgfusepath{clip}%
\pgfsetbuttcap%
\pgfsetroundjoin%
\pgfsetlinewidth{1.003750pt}%
\definecolor{currentstroke}{rgb}{0.541176,0.380392,0.658824}%
\pgfsetstrokecolor{currentstroke}%
\pgfsetstrokeopacity{0.900000}%
\pgfsetdash{}{0pt}%
\pgfpathmoveto{\pgfqpoint{10.682224in}{7.557099in}}%
\pgfpathcurveto{\pgfqpoint{10.692479in}{7.557099in}}{\pgfqpoint{10.702314in}{7.561173in}}{\pgfqpoint{10.709565in}{7.568423in}}%
\pgfpathcurveto{\pgfqpoint{10.716816in}{7.575674in}}{\pgfqpoint{10.720889in}{7.585510in}}{\pgfqpoint{10.720889in}{7.595764in}}%
\pgfpathcurveto{\pgfqpoint{10.720889in}{7.606018in}}{\pgfqpoint{10.716816in}{7.615853in}}{\pgfqpoint{10.709565in}{7.623104in}}%
\pgfpathcurveto{\pgfqpoint{10.702314in}{7.630355in}}{\pgfqpoint{10.692479in}{7.634429in}}{\pgfqpoint{10.682224in}{7.634429in}}%
\pgfpathcurveto{\pgfqpoint{10.671970in}{7.634429in}}{\pgfqpoint{10.662135in}{7.630355in}}{\pgfqpoint{10.654884in}{7.623104in}}%
\pgfpathcurveto{\pgfqpoint{10.647633in}{7.615853in}}{\pgfqpoint{10.643559in}{7.606018in}}{\pgfqpoint{10.643559in}{7.595764in}}%
\pgfpathcurveto{\pgfqpoint{10.643559in}{7.585510in}}{\pgfqpoint{10.647633in}{7.575674in}}{\pgfqpoint{10.654884in}{7.568423in}}%
\pgfpathcurveto{\pgfqpoint{10.662135in}{7.561173in}}{\pgfqpoint{10.671970in}{7.557099in}}{\pgfqpoint{10.682224in}{7.557099in}}%
\pgfpathlineto{\pgfqpoint{10.682224in}{7.557099in}}%
\pgfpathclose%
\pgfusepath{stroke}%
\end{pgfscope}%
\begin{pgfscope}%
\pgfpathrectangle{\pgfqpoint{8.036091in}{4.050000in}}{\pgfqpoint{5.400000in}{5.400000in}}%
\pgfusepath{clip}%
\pgfsetbuttcap%
\pgfsetroundjoin%
\pgfsetlinewidth{1.003750pt}%
\definecolor{currentstroke}{rgb}{0.541176,0.380392,0.658824}%
\pgfsetstrokecolor{currentstroke}%
\pgfsetstrokeopacity{0.900000}%
\pgfsetdash{}{0pt}%
\pgfpathmoveto{\pgfqpoint{10.175027in}{6.229411in}}%
\pgfpathcurveto{\pgfqpoint{10.185281in}{6.229411in}}{\pgfqpoint{10.195116in}{6.233485in}}{\pgfqpoint{10.202367in}{6.240736in}}%
\pgfpathcurveto{\pgfqpoint{10.209618in}{6.247986in}}{\pgfqpoint{10.213692in}{6.257822in}}{\pgfqpoint{10.213692in}{6.268076in}}%
\pgfpathcurveto{\pgfqpoint{10.213692in}{6.278330in}}{\pgfqpoint{10.209618in}{6.288165in}}{\pgfqpoint{10.202367in}{6.295416in}}%
\pgfpathcurveto{\pgfqpoint{10.195116in}{6.302667in}}{\pgfqpoint{10.185281in}{6.306741in}}{\pgfqpoint{10.175027in}{6.306741in}}%
\pgfpathcurveto{\pgfqpoint{10.164773in}{6.306741in}}{\pgfqpoint{10.154937in}{6.302667in}}{\pgfqpoint{10.147687in}{6.295416in}}%
\pgfpathcurveto{\pgfqpoint{10.140436in}{6.288165in}}{\pgfqpoint{10.136362in}{6.278330in}}{\pgfqpoint{10.136362in}{6.268076in}}%
\pgfpathcurveto{\pgfqpoint{10.136362in}{6.257822in}}{\pgfqpoint{10.140436in}{6.247986in}}{\pgfqpoint{10.147687in}{6.240736in}}%
\pgfpathcurveto{\pgfqpoint{10.154937in}{6.233485in}}{\pgfqpoint{10.164773in}{6.229411in}}{\pgfqpoint{10.175027in}{6.229411in}}%
\pgfpathlineto{\pgfqpoint{10.175027in}{6.229411in}}%
\pgfpathclose%
\pgfusepath{stroke}%
\end{pgfscope}%
\begin{pgfscope}%
\pgfpathrectangle{\pgfqpoint{8.036091in}{4.050000in}}{\pgfqpoint{5.400000in}{5.400000in}}%
\pgfusepath{clip}%
\pgfsetbuttcap%
\pgfsetroundjoin%
\pgfsetlinewidth{1.003750pt}%
\definecolor{currentstroke}{rgb}{0.541176,0.380392,0.658824}%
\pgfsetstrokecolor{currentstroke}%
\pgfsetstrokeopacity{0.900000}%
\pgfsetdash{}{0pt}%
\pgfpathmoveto{\pgfqpoint{10.816251in}{7.034515in}}%
\pgfpathcurveto{\pgfqpoint{10.826505in}{7.034515in}}{\pgfqpoint{10.836340in}{7.038589in}}{\pgfqpoint{10.843591in}{7.045840in}}%
\pgfpathcurveto{\pgfqpoint{10.850842in}{7.053091in}}{\pgfqpoint{10.854916in}{7.062926in}}{\pgfqpoint{10.854916in}{7.073181in}}%
\pgfpathcurveto{\pgfqpoint{10.854916in}{7.083435in}}{\pgfqpoint{10.850842in}{7.093270in}}{\pgfqpoint{10.843591in}{7.100521in}}%
\pgfpathcurveto{\pgfqpoint{10.836340in}{7.107772in}}{\pgfqpoint{10.826505in}{7.111846in}}{\pgfqpoint{10.816251in}{7.111846in}}%
\pgfpathcurveto{\pgfqpoint{10.805997in}{7.111846in}}{\pgfqpoint{10.796161in}{7.107772in}}{\pgfqpoint{10.788910in}{7.100521in}}%
\pgfpathcurveto{\pgfqpoint{10.781660in}{7.093270in}}{\pgfqpoint{10.777586in}{7.083435in}}{\pgfqpoint{10.777586in}{7.073181in}}%
\pgfpathcurveto{\pgfqpoint{10.777586in}{7.062926in}}{\pgfqpoint{10.781660in}{7.053091in}}{\pgfqpoint{10.788910in}{7.045840in}}%
\pgfpathcurveto{\pgfqpoint{10.796161in}{7.038589in}}{\pgfqpoint{10.805997in}{7.034515in}}{\pgfqpoint{10.816251in}{7.034515in}}%
\pgfpathlineto{\pgfqpoint{10.816251in}{7.034515in}}%
\pgfpathclose%
\pgfusepath{stroke}%
\end{pgfscope}%
\begin{pgfscope}%
\pgfpathrectangle{\pgfqpoint{8.036091in}{4.050000in}}{\pgfqpoint{5.400000in}{5.400000in}}%
\pgfusepath{clip}%
\pgfsetbuttcap%
\pgfsetroundjoin%
\pgfsetlinewidth{1.003750pt}%
\definecolor{currentstroke}{rgb}{0.541176,0.380392,0.658824}%
\pgfsetstrokecolor{currentstroke}%
\pgfsetstrokeopacity{0.900000}%
\pgfsetdash{}{0pt}%
\pgfpathmoveto{\pgfqpoint{11.438451in}{7.824533in}}%
\pgfpathcurveto{\pgfqpoint{11.448705in}{7.824533in}}{\pgfqpoint{11.458540in}{7.828607in}}{\pgfqpoint{11.465791in}{7.835858in}}%
\pgfpathcurveto{\pgfqpoint{11.473042in}{7.843109in}}{\pgfqpoint{11.477116in}{7.852944in}}{\pgfqpoint{11.477116in}{7.863198in}}%
\pgfpathcurveto{\pgfqpoint{11.477116in}{7.873453in}}{\pgfqpoint{11.473042in}{7.883288in}}{\pgfqpoint{11.465791in}{7.890539in}}%
\pgfpathcurveto{\pgfqpoint{11.458540in}{7.897789in}}{\pgfqpoint{11.448705in}{7.901863in}}{\pgfqpoint{11.438451in}{7.901863in}}%
\pgfpathcurveto{\pgfqpoint{11.428196in}{7.901863in}}{\pgfqpoint{11.418361in}{7.897789in}}{\pgfqpoint{11.411110in}{7.890539in}}%
\pgfpathcurveto{\pgfqpoint{11.403859in}{7.883288in}}{\pgfqpoint{11.399785in}{7.873453in}}{\pgfqpoint{11.399785in}{7.863198in}}%
\pgfpathcurveto{\pgfqpoint{11.399785in}{7.852944in}}{\pgfqpoint{11.403859in}{7.843109in}}{\pgfqpoint{11.411110in}{7.835858in}}%
\pgfpathcurveto{\pgfqpoint{11.418361in}{7.828607in}}{\pgfqpoint{11.428196in}{7.824533in}}{\pgfqpoint{11.438451in}{7.824533in}}%
\pgfpathlineto{\pgfqpoint{11.438451in}{7.824533in}}%
\pgfpathclose%
\pgfusepath{stroke}%
\end{pgfscope}%
\begin{pgfscope}%
\pgfpathrectangle{\pgfqpoint{8.036091in}{4.050000in}}{\pgfqpoint{5.400000in}{5.400000in}}%
\pgfusepath{clip}%
\pgfsetbuttcap%
\pgfsetroundjoin%
\pgfsetlinewidth{1.003750pt}%
\definecolor{currentstroke}{rgb}{0.541176,0.380392,0.658824}%
\pgfsetstrokecolor{currentstroke}%
\pgfsetstrokeopacity{0.900000}%
\pgfsetdash{}{0pt}%
\pgfpathmoveto{\pgfqpoint{11.906021in}{6.861880in}}%
\pgfpathcurveto{\pgfqpoint{11.916275in}{6.861880in}}{\pgfqpoint{11.926110in}{6.865954in}}{\pgfqpoint{11.933361in}{6.873204in}}%
\pgfpathcurveto{\pgfqpoint{11.940612in}{6.880455in}}{\pgfqpoint{11.944686in}{6.890291in}}{\pgfqpoint{11.944686in}{6.900545in}}%
\pgfpathcurveto{\pgfqpoint{11.944686in}{6.910799in}}{\pgfqpoint{11.940612in}{6.920634in}}{\pgfqpoint{11.933361in}{6.927885in}}%
\pgfpathcurveto{\pgfqpoint{11.926110in}{6.935136in}}{\pgfqpoint{11.916275in}{6.939210in}}{\pgfqpoint{11.906021in}{6.939210in}}%
\pgfpathcurveto{\pgfqpoint{11.895767in}{6.939210in}}{\pgfqpoint{11.885931in}{6.935136in}}{\pgfqpoint{11.878681in}{6.927885in}}%
\pgfpathcurveto{\pgfqpoint{11.871430in}{6.920634in}}{\pgfqpoint{11.867356in}{6.910799in}}{\pgfqpoint{11.867356in}{6.900545in}}%
\pgfpathcurveto{\pgfqpoint{11.867356in}{6.890291in}}{\pgfqpoint{11.871430in}{6.880455in}}{\pgfqpoint{11.878681in}{6.873204in}}%
\pgfpathcurveto{\pgfqpoint{11.885931in}{6.865954in}}{\pgfqpoint{11.895767in}{6.861880in}}{\pgfqpoint{11.906021in}{6.861880in}}%
\pgfpathlineto{\pgfqpoint{11.906021in}{6.861880in}}%
\pgfpathclose%
\pgfusepath{stroke}%
\end{pgfscope}%
\begin{pgfscope}%
\pgfpathrectangle{\pgfqpoint{8.036091in}{4.050000in}}{\pgfqpoint{5.400000in}{5.400000in}}%
\pgfusepath{clip}%
\pgfsetbuttcap%
\pgfsetroundjoin%
\pgfsetlinewidth{1.003750pt}%
\definecolor{currentstroke}{rgb}{0.541176,0.380392,0.658824}%
\pgfsetstrokecolor{currentstroke}%
\pgfsetstrokeopacity{0.900000}%
\pgfsetdash{}{0pt}%
\pgfpathmoveto{\pgfqpoint{11.693635in}{6.977808in}}%
\pgfpathcurveto{\pgfqpoint{11.703889in}{6.977808in}}{\pgfqpoint{11.713724in}{6.981882in}}{\pgfqpoint{11.720975in}{6.989132in}}%
\pgfpathcurveto{\pgfqpoint{11.728226in}{6.996383in}}{\pgfqpoint{11.732300in}{7.006219in}}{\pgfqpoint{11.732300in}{7.016473in}}%
\pgfpathcurveto{\pgfqpoint{11.732300in}{7.026727in}}{\pgfqpoint{11.728226in}{7.036562in}}{\pgfqpoint{11.720975in}{7.043813in}}%
\pgfpathcurveto{\pgfqpoint{11.713724in}{7.051064in}}{\pgfqpoint{11.703889in}{7.055138in}}{\pgfqpoint{11.693635in}{7.055138in}}%
\pgfpathcurveto{\pgfqpoint{11.683381in}{7.055138in}}{\pgfqpoint{11.673545in}{7.051064in}}{\pgfqpoint{11.666294in}{7.043813in}}%
\pgfpathcurveto{\pgfqpoint{11.659044in}{7.036562in}}{\pgfqpoint{11.654970in}{7.026727in}}{\pgfqpoint{11.654970in}{7.016473in}}%
\pgfpathcurveto{\pgfqpoint{11.654970in}{7.006219in}}{\pgfqpoint{11.659044in}{6.996383in}}{\pgfqpoint{11.666294in}{6.989132in}}%
\pgfpathcurveto{\pgfqpoint{11.673545in}{6.981882in}}{\pgfqpoint{11.683381in}{6.977808in}}{\pgfqpoint{11.693635in}{6.977808in}}%
\pgfpathlineto{\pgfqpoint{11.693635in}{6.977808in}}%
\pgfpathclose%
\pgfusepath{stroke}%
\end{pgfscope}%
\begin{pgfscope}%
\pgfpathrectangle{\pgfqpoint{8.036091in}{4.050000in}}{\pgfqpoint{5.400000in}{5.400000in}}%
\pgfusepath{clip}%
\pgfsetbuttcap%
\pgfsetroundjoin%
\pgfsetlinewidth{1.003750pt}%
\definecolor{currentstroke}{rgb}{0.541176,0.380392,0.658824}%
\pgfsetstrokecolor{currentstroke}%
\pgfsetstrokeopacity{0.900000}%
\pgfsetdash{}{0pt}%
\pgfpathmoveto{\pgfqpoint{11.987092in}{8.133581in}}%
\pgfpathcurveto{\pgfqpoint{11.997346in}{8.133581in}}{\pgfqpoint{12.007181in}{8.137655in}}{\pgfqpoint{12.014432in}{8.144906in}}%
\pgfpathcurveto{\pgfqpoint{12.021683in}{8.152157in}}{\pgfqpoint{12.025757in}{8.161992in}}{\pgfqpoint{12.025757in}{8.172246in}}%
\pgfpathcurveto{\pgfqpoint{12.025757in}{8.182501in}}{\pgfqpoint{12.021683in}{8.192336in}}{\pgfqpoint{12.014432in}{8.199587in}}%
\pgfpathcurveto{\pgfqpoint{12.007181in}{8.206837in}}{\pgfqpoint{11.997346in}{8.210911in}}{\pgfqpoint{11.987092in}{8.210911in}}%
\pgfpathcurveto{\pgfqpoint{11.976838in}{8.210911in}}{\pgfqpoint{11.967002in}{8.206837in}}{\pgfqpoint{11.959752in}{8.199587in}}%
\pgfpathcurveto{\pgfqpoint{11.952501in}{8.192336in}}{\pgfqpoint{11.948427in}{8.182501in}}{\pgfqpoint{11.948427in}{8.172246in}}%
\pgfpathcurveto{\pgfqpoint{11.948427in}{8.161992in}}{\pgfqpoint{11.952501in}{8.152157in}}{\pgfqpoint{11.959752in}{8.144906in}}%
\pgfpathcurveto{\pgfqpoint{11.967002in}{8.137655in}}{\pgfqpoint{11.976838in}{8.133581in}}{\pgfqpoint{11.987092in}{8.133581in}}%
\pgfpathlineto{\pgfqpoint{11.987092in}{8.133581in}}%
\pgfpathclose%
\pgfusepath{stroke}%
\end{pgfscope}%
\begin{pgfscope}%
\pgfpathrectangle{\pgfqpoint{8.036091in}{4.050000in}}{\pgfqpoint{5.400000in}{5.400000in}}%
\pgfusepath{clip}%
\pgfsetbuttcap%
\pgfsetroundjoin%
\pgfsetlinewidth{1.003750pt}%
\definecolor{currentstroke}{rgb}{0.541176,0.380392,0.658824}%
\pgfsetstrokecolor{currentstroke}%
\pgfsetstrokeopacity{0.900000}%
\pgfsetdash{}{0pt}%
\pgfpathmoveto{\pgfqpoint{11.421456in}{7.896910in}}%
\pgfpathcurveto{\pgfqpoint{11.431710in}{7.896910in}}{\pgfqpoint{11.441545in}{7.900984in}}{\pgfqpoint{11.448796in}{7.908235in}}%
\pgfpathcurveto{\pgfqpoint{11.456047in}{7.915485in}}{\pgfqpoint{11.460121in}{7.925321in}}{\pgfqpoint{11.460121in}{7.935575in}}%
\pgfpathcurveto{\pgfqpoint{11.460121in}{7.945829in}}{\pgfqpoint{11.456047in}{7.955664in}}{\pgfqpoint{11.448796in}{7.962915in}}%
\pgfpathcurveto{\pgfqpoint{11.441545in}{7.970166in}}{\pgfqpoint{11.431710in}{7.974240in}}{\pgfqpoint{11.421456in}{7.974240in}}%
\pgfpathcurveto{\pgfqpoint{11.411202in}{7.974240in}}{\pgfqpoint{11.401366in}{7.970166in}}{\pgfqpoint{11.394115in}{7.962915in}}%
\pgfpathcurveto{\pgfqpoint{11.386865in}{7.955664in}}{\pgfqpoint{11.382791in}{7.945829in}}{\pgfqpoint{11.382791in}{7.935575in}}%
\pgfpathcurveto{\pgfqpoint{11.382791in}{7.925321in}}{\pgfqpoint{11.386865in}{7.915485in}}{\pgfqpoint{11.394115in}{7.908235in}}%
\pgfpathcurveto{\pgfqpoint{11.401366in}{7.900984in}}{\pgfqpoint{11.411202in}{7.896910in}}{\pgfqpoint{11.421456in}{7.896910in}}%
\pgfpathlineto{\pgfqpoint{11.421456in}{7.896910in}}%
\pgfpathclose%
\pgfusepath{stroke}%
\end{pgfscope}%
\begin{pgfscope}%
\pgfpathrectangle{\pgfqpoint{8.036091in}{4.050000in}}{\pgfqpoint{5.400000in}{5.400000in}}%
\pgfusepath{clip}%
\pgfsetbuttcap%
\pgfsetroundjoin%
\pgfsetlinewidth{1.003750pt}%
\definecolor{currentstroke}{rgb}{0.541176,0.380392,0.658824}%
\pgfsetstrokecolor{currentstroke}%
\pgfsetstrokeopacity{0.900000}%
\pgfsetdash{}{0pt}%
\pgfpathmoveto{\pgfqpoint{12.116159in}{8.581303in}}%
\pgfpathcurveto{\pgfqpoint{12.126413in}{8.581303in}}{\pgfqpoint{12.136248in}{8.585377in}}{\pgfqpoint{12.143499in}{8.592627in}}%
\pgfpathcurveto{\pgfqpoint{12.150750in}{8.599878in}}{\pgfqpoint{12.154824in}{8.609713in}}{\pgfqpoint{12.154824in}{8.619968in}}%
\pgfpathcurveto{\pgfqpoint{12.154824in}{8.630222in}}{\pgfqpoint{12.150750in}{8.640057in}}{\pgfqpoint{12.143499in}{8.647308in}}%
\pgfpathcurveto{\pgfqpoint{12.136248in}{8.654559in}}{\pgfqpoint{12.126413in}{8.658633in}}{\pgfqpoint{12.116159in}{8.658633in}}%
\pgfpathcurveto{\pgfqpoint{12.105905in}{8.658633in}}{\pgfqpoint{12.096069in}{8.654559in}}{\pgfqpoint{12.088819in}{8.647308in}}%
\pgfpathcurveto{\pgfqpoint{12.081568in}{8.640057in}}{\pgfqpoint{12.077494in}{8.630222in}}{\pgfqpoint{12.077494in}{8.619968in}}%
\pgfpathcurveto{\pgfqpoint{12.077494in}{8.609713in}}{\pgfqpoint{12.081568in}{8.599878in}}{\pgfqpoint{12.088819in}{8.592627in}}%
\pgfpathcurveto{\pgfqpoint{12.096069in}{8.585377in}}{\pgfqpoint{12.105905in}{8.581303in}}{\pgfqpoint{12.116159in}{8.581303in}}%
\pgfpathlineto{\pgfqpoint{12.116159in}{8.581303in}}%
\pgfpathclose%
\pgfusepath{stroke}%
\end{pgfscope}%
\begin{pgfscope}%
\pgfpathrectangle{\pgfqpoint{8.036091in}{4.050000in}}{\pgfqpoint{5.400000in}{5.400000in}}%
\pgfusepath{clip}%
\pgfsetbuttcap%
\pgfsetroundjoin%
\pgfsetlinewidth{1.003750pt}%
\definecolor{currentstroke}{rgb}{0.541176,0.380392,0.658824}%
\pgfsetstrokecolor{currentstroke}%
\pgfsetstrokeopacity{0.900000}%
\pgfsetdash{}{0pt}%
\pgfpathmoveto{\pgfqpoint{10.689403in}{5.888266in}}%
\pgfpathcurveto{\pgfqpoint{10.699658in}{5.888266in}}{\pgfqpoint{10.709493in}{5.892340in}}{\pgfqpoint{10.716744in}{5.899591in}}%
\pgfpathcurveto{\pgfqpoint{10.723994in}{5.906842in}}{\pgfqpoint{10.728068in}{5.916677in}}{\pgfqpoint{10.728068in}{5.926931in}}%
\pgfpathcurveto{\pgfqpoint{10.728068in}{5.937185in}}{\pgfqpoint{10.723994in}{5.947021in}}{\pgfqpoint{10.716744in}{5.954272in}}%
\pgfpathcurveto{\pgfqpoint{10.709493in}{5.961522in}}{\pgfqpoint{10.699658in}{5.965596in}}{\pgfqpoint{10.689403in}{5.965596in}}%
\pgfpathcurveto{\pgfqpoint{10.679149in}{5.965596in}}{\pgfqpoint{10.669314in}{5.961522in}}{\pgfqpoint{10.662063in}{5.954272in}}%
\pgfpathcurveto{\pgfqpoint{10.654812in}{5.947021in}}{\pgfqpoint{10.650738in}{5.937185in}}{\pgfqpoint{10.650738in}{5.926931in}}%
\pgfpathcurveto{\pgfqpoint{10.650738in}{5.916677in}}{\pgfqpoint{10.654812in}{5.906842in}}{\pgfqpoint{10.662063in}{5.899591in}}%
\pgfpathcurveto{\pgfqpoint{10.669314in}{5.892340in}}{\pgfqpoint{10.679149in}{5.888266in}}{\pgfqpoint{10.689403in}{5.888266in}}%
\pgfpathlineto{\pgfqpoint{10.689403in}{5.888266in}}%
\pgfpathclose%
\pgfusepath{stroke}%
\end{pgfscope}%
\begin{pgfscope}%
\pgfpathrectangle{\pgfqpoint{8.036091in}{4.050000in}}{\pgfqpoint{5.400000in}{5.400000in}}%
\pgfusepath{clip}%
\pgfsetbuttcap%
\pgfsetroundjoin%
\pgfsetlinewidth{1.003750pt}%
\definecolor{currentstroke}{rgb}{0.541176,0.380392,0.658824}%
\pgfsetstrokecolor{currentstroke}%
\pgfsetstrokeopacity{0.900000}%
\pgfsetdash{}{0pt}%
\pgfpathmoveto{\pgfqpoint{11.593614in}{7.657104in}}%
\pgfpathcurveto{\pgfqpoint{11.603868in}{7.657104in}}{\pgfqpoint{11.613703in}{7.661178in}}{\pgfqpoint{11.620954in}{7.668428in}}%
\pgfpathcurveto{\pgfqpoint{11.628205in}{7.675679in}}{\pgfqpoint{11.632279in}{7.685515in}}{\pgfqpoint{11.632279in}{7.695769in}}%
\pgfpathcurveto{\pgfqpoint{11.632279in}{7.706023in}}{\pgfqpoint{11.628205in}{7.715858in}}{\pgfqpoint{11.620954in}{7.723109in}}%
\pgfpathcurveto{\pgfqpoint{11.613703in}{7.730360in}}{\pgfqpoint{11.603868in}{7.734434in}}{\pgfqpoint{11.593614in}{7.734434in}}%
\pgfpathcurveto{\pgfqpoint{11.583360in}{7.734434in}}{\pgfqpoint{11.573524in}{7.730360in}}{\pgfqpoint{11.566274in}{7.723109in}}%
\pgfpathcurveto{\pgfqpoint{11.559023in}{7.715858in}}{\pgfqpoint{11.554949in}{7.706023in}}{\pgfqpoint{11.554949in}{7.695769in}}%
\pgfpathcurveto{\pgfqpoint{11.554949in}{7.685515in}}{\pgfqpoint{11.559023in}{7.675679in}}{\pgfqpoint{11.566274in}{7.668428in}}%
\pgfpathcurveto{\pgfqpoint{11.573524in}{7.661178in}}{\pgfqpoint{11.583360in}{7.657104in}}{\pgfqpoint{11.593614in}{7.657104in}}%
\pgfpathlineto{\pgfqpoint{11.593614in}{7.657104in}}%
\pgfpathclose%
\pgfusepath{stroke}%
\end{pgfscope}%
\begin{pgfscope}%
\pgfsetbuttcap%
\pgfsetroundjoin%
\pgfsetlinewidth{1.003750pt}%
\definecolor{currentstroke}{rgb}{0.541176,0.380392,0.658824}%
\pgfsetstrokecolor{currentstroke}%
\pgfsetstrokeopacity{0.900000}%
\pgfsetdash{}{0pt}%
\pgfpathmoveto{\pgfqpoint{13.436091in}{9.390262in}}%
\pgfpathlineto{\pgfqpoint{13.495829in}{9.450000in}}%
\pgfpathlineto{\pgfqpoint{13.436091in}{9.509738in}}%
\pgfpathlineto{\pgfqpoint{13.376353in}{9.450000in}}%
\pgfpathlineto{\pgfqpoint{13.436091in}{9.390262in}}%
\pgfpathclose%
\pgfusepath{stroke}%
\end{pgfscope}%
\begin{pgfscope}%
\pgfpathrectangle{\pgfqpoint{8.036091in}{4.050000in}}{\pgfqpoint{5.400000in}{5.400000in}}%
\pgfusepath{clip}%
\pgfsetbuttcap%
\pgfsetroundjoin%
\pgfsetlinewidth{1.003750pt}%
\definecolor{currentstroke}{rgb}{0.321569,0.321569,0.321569}%
\pgfsetstrokecolor{currentstroke}%
\pgfsetstrokeopacity{0.900000}%
\pgfsetdash{}{0pt}%
\pgfpathmoveto{\pgfqpoint{9.435561in}{5.451830in}}%
\pgfpathcurveto{\pgfqpoint{9.445815in}{5.451830in}}{\pgfqpoint{9.455651in}{5.455904in}}{\pgfqpoint{9.462901in}{5.463154in}}%
\pgfpathcurveto{\pgfqpoint{9.470152in}{5.470405in}}{\pgfqpoint{9.474226in}{5.480241in}}{\pgfqpoint{9.474226in}{5.490495in}}%
\pgfpathcurveto{\pgfqpoint{9.474226in}{5.500749in}}{\pgfqpoint{9.470152in}{5.510584in}}{\pgfqpoint{9.462901in}{5.517835in}}%
\pgfpathcurveto{\pgfqpoint{9.455651in}{5.525086in}}{\pgfqpoint{9.445815in}{5.529160in}}{\pgfqpoint{9.435561in}{5.529160in}}%
\pgfpathcurveto{\pgfqpoint{9.425307in}{5.529160in}}{\pgfqpoint{9.415471in}{5.525086in}}{\pgfqpoint{9.408221in}{5.517835in}}%
\pgfpathcurveto{\pgfqpoint{9.400970in}{5.510584in}}{\pgfqpoint{9.396896in}{5.500749in}}{\pgfqpoint{9.396896in}{5.490495in}}%
\pgfpathcurveto{\pgfqpoint{9.396896in}{5.480241in}}{\pgfqpoint{9.400970in}{5.470405in}}{\pgfqpoint{9.408221in}{5.463154in}}%
\pgfpathcurveto{\pgfqpoint{9.415471in}{5.455904in}}{\pgfqpoint{9.425307in}{5.451830in}}{\pgfqpoint{9.435561in}{5.451830in}}%
\pgfpathlineto{\pgfqpoint{9.435561in}{5.451830in}}%
\pgfpathclose%
\pgfusepath{stroke}%
\end{pgfscope}%
\begin{pgfscope}%
\pgfpathrectangle{\pgfqpoint{8.036091in}{4.050000in}}{\pgfqpoint{5.400000in}{5.400000in}}%
\pgfusepath{clip}%
\pgfsetbuttcap%
\pgfsetroundjoin%
\pgfsetlinewidth{1.003750pt}%
\definecolor{currentstroke}{rgb}{0.321569,0.321569,0.321569}%
\pgfsetstrokecolor{currentstroke}%
\pgfsetstrokeopacity{0.900000}%
\pgfsetdash{}{0pt}%
\pgfpathmoveto{\pgfqpoint{9.499099in}{5.491570in}}%
\pgfpathcurveto{\pgfqpoint{9.509353in}{5.491570in}}{\pgfqpoint{9.519188in}{5.495644in}}{\pgfqpoint{9.526439in}{5.502895in}}%
\pgfpathcurveto{\pgfqpoint{9.533690in}{5.510145in}}{\pgfqpoint{9.537764in}{5.519981in}}{\pgfqpoint{9.537764in}{5.530235in}}%
\pgfpathcurveto{\pgfqpoint{9.537764in}{5.540489in}}{\pgfqpoint{9.533690in}{5.550324in}}{\pgfqpoint{9.526439in}{5.557575in}}%
\pgfpathcurveto{\pgfqpoint{9.519188in}{5.564826in}}{\pgfqpoint{9.509353in}{5.568900in}}{\pgfqpoint{9.499099in}{5.568900in}}%
\pgfpathcurveto{\pgfqpoint{9.488845in}{5.568900in}}{\pgfqpoint{9.479009in}{5.564826in}}{\pgfqpoint{9.471759in}{5.557575in}}%
\pgfpathcurveto{\pgfqpoint{9.464508in}{5.550324in}}{\pgfqpoint{9.460434in}{5.540489in}}{\pgfqpoint{9.460434in}{5.530235in}}%
\pgfpathcurveto{\pgfqpoint{9.460434in}{5.519981in}}{\pgfqpoint{9.464508in}{5.510145in}}{\pgfqpoint{9.471759in}{5.502895in}}%
\pgfpathcurveto{\pgfqpoint{9.479009in}{5.495644in}}{\pgfqpoint{9.488845in}{5.491570in}}{\pgfqpoint{9.499099in}{5.491570in}}%
\pgfpathlineto{\pgfqpoint{9.499099in}{5.491570in}}%
\pgfpathclose%
\pgfusepath{stroke}%
\end{pgfscope}%
\begin{pgfscope}%
\pgfpathrectangle{\pgfqpoint{8.036091in}{4.050000in}}{\pgfqpoint{5.400000in}{5.400000in}}%
\pgfusepath{clip}%
\pgfsetbuttcap%
\pgfsetroundjoin%
\pgfsetlinewidth{1.003750pt}%
\definecolor{currentstroke}{rgb}{0.321569,0.321569,0.321569}%
\pgfsetstrokecolor{currentstroke}%
\pgfsetstrokeopacity{0.900000}%
\pgfsetdash{}{0pt}%
\pgfpathmoveto{\pgfqpoint{9.050332in}{4.997999in}}%
\pgfpathcurveto{\pgfqpoint{9.060586in}{4.997999in}}{\pgfqpoint{9.070422in}{5.002073in}}{\pgfqpoint{9.077673in}{5.009324in}}%
\pgfpathcurveto{\pgfqpoint{9.084923in}{5.016575in}}{\pgfqpoint{9.088997in}{5.026410in}}{\pgfqpoint{9.088997in}{5.036664in}}%
\pgfpathcurveto{\pgfqpoint{9.088997in}{5.046918in}}{\pgfqpoint{9.084923in}{5.056754in}}{\pgfqpoint{9.077673in}{5.064005in}}%
\pgfpathcurveto{\pgfqpoint{9.070422in}{5.071255in}}{\pgfqpoint{9.060586in}{5.075329in}}{\pgfqpoint{9.050332in}{5.075329in}}%
\pgfpathcurveto{\pgfqpoint{9.040078in}{5.075329in}}{\pgfqpoint{9.030243in}{5.071255in}}{\pgfqpoint{9.022992in}{5.064005in}}%
\pgfpathcurveto{\pgfqpoint{9.015741in}{5.056754in}}{\pgfqpoint{9.011667in}{5.046918in}}{\pgfqpoint{9.011667in}{5.036664in}}%
\pgfpathcurveto{\pgfqpoint{9.011667in}{5.026410in}}{\pgfqpoint{9.015741in}{5.016575in}}{\pgfqpoint{9.022992in}{5.009324in}}%
\pgfpathcurveto{\pgfqpoint{9.030243in}{5.002073in}}{\pgfqpoint{9.040078in}{4.997999in}}{\pgfqpoint{9.050332in}{4.997999in}}%
\pgfpathlineto{\pgfqpoint{9.050332in}{4.997999in}}%
\pgfpathclose%
\pgfusepath{stroke}%
\end{pgfscope}%
\begin{pgfscope}%
\pgfpathrectangle{\pgfqpoint{8.036091in}{4.050000in}}{\pgfqpoint{5.400000in}{5.400000in}}%
\pgfusepath{clip}%
\pgfsetbuttcap%
\pgfsetroundjoin%
\pgfsetlinewidth{1.003750pt}%
\definecolor{currentstroke}{rgb}{0.321569,0.321569,0.321569}%
\pgfsetstrokecolor{currentstroke}%
\pgfsetstrokeopacity{0.900000}%
\pgfsetdash{}{0pt}%
\pgfpathmoveto{\pgfqpoint{9.225044in}{5.143534in}}%
\pgfpathcurveto{\pgfqpoint{9.235298in}{5.143534in}}{\pgfqpoint{9.245134in}{5.147608in}}{\pgfqpoint{9.252384in}{5.154858in}}%
\pgfpathcurveto{\pgfqpoint{9.259635in}{5.162109in}}{\pgfqpoint{9.263709in}{5.171944in}}{\pgfqpoint{9.263709in}{5.182199in}}%
\pgfpathcurveto{\pgfqpoint{9.263709in}{5.192453in}}{\pgfqpoint{9.259635in}{5.202288in}}{\pgfqpoint{9.252384in}{5.209539in}}%
\pgfpathcurveto{\pgfqpoint{9.245134in}{5.216790in}}{\pgfqpoint{9.235298in}{5.220864in}}{\pgfqpoint{9.225044in}{5.220864in}}%
\pgfpathcurveto{\pgfqpoint{9.214790in}{5.220864in}}{\pgfqpoint{9.204955in}{5.216790in}}{\pgfqpoint{9.197704in}{5.209539in}}%
\pgfpathcurveto{\pgfqpoint{9.190453in}{5.202288in}}{\pgfqpoint{9.186379in}{5.192453in}}{\pgfqpoint{9.186379in}{5.182199in}}%
\pgfpathcurveto{\pgfqpoint{9.186379in}{5.171944in}}{\pgfqpoint{9.190453in}{5.162109in}}{\pgfqpoint{9.197704in}{5.154858in}}%
\pgfpathcurveto{\pgfqpoint{9.204955in}{5.147608in}}{\pgfqpoint{9.214790in}{5.143534in}}{\pgfqpoint{9.225044in}{5.143534in}}%
\pgfpathlineto{\pgfqpoint{9.225044in}{5.143534in}}%
\pgfpathclose%
\pgfusepath{stroke}%
\end{pgfscope}%
\begin{pgfscope}%
\pgfpathrectangle{\pgfqpoint{8.036091in}{4.050000in}}{\pgfqpoint{5.400000in}{5.400000in}}%
\pgfusepath{clip}%
\pgfsetbuttcap%
\pgfsetroundjoin%
\pgfsetlinewidth{1.003750pt}%
\definecolor{currentstroke}{rgb}{0.321569,0.321569,0.321569}%
\pgfsetstrokecolor{currentstroke}%
\pgfsetstrokeopacity{0.900000}%
\pgfsetdash{}{0pt}%
\pgfpathmoveto{\pgfqpoint{12.905828in}{9.160611in}}%
\pgfpathcurveto{\pgfqpoint{12.916082in}{9.160611in}}{\pgfqpoint{12.925917in}{9.164685in}}{\pgfqpoint{12.933168in}{9.171936in}}%
\pgfpathcurveto{\pgfqpoint{12.940419in}{9.179187in}}{\pgfqpoint{12.944493in}{9.189022in}}{\pgfqpoint{12.944493in}{9.199276in}}%
\pgfpathcurveto{\pgfqpoint{12.944493in}{9.209530in}}{\pgfqpoint{12.940419in}{9.219366in}}{\pgfqpoint{12.933168in}{9.226617in}}%
\pgfpathcurveto{\pgfqpoint{12.925917in}{9.233867in}}{\pgfqpoint{12.916082in}{9.237941in}}{\pgfqpoint{12.905828in}{9.237941in}}%
\pgfpathcurveto{\pgfqpoint{12.895574in}{9.237941in}}{\pgfqpoint{12.885738in}{9.233867in}}{\pgfqpoint{12.878487in}{9.226617in}}%
\pgfpathcurveto{\pgfqpoint{12.871237in}{9.219366in}}{\pgfqpoint{12.867163in}{9.209530in}}{\pgfqpoint{12.867163in}{9.199276in}}%
\pgfpathcurveto{\pgfqpoint{12.867163in}{9.189022in}}{\pgfqpoint{12.871237in}{9.179187in}}{\pgfqpoint{12.878487in}{9.171936in}}%
\pgfpathcurveto{\pgfqpoint{12.885738in}{9.164685in}}{\pgfqpoint{12.895574in}{9.160611in}}{\pgfqpoint{12.905828in}{9.160611in}}%
\pgfpathlineto{\pgfqpoint{12.905828in}{9.160611in}}%
\pgfpathclose%
\pgfusepath{stroke}%
\end{pgfscope}%
\begin{pgfscope}%
\pgfpathrectangle{\pgfqpoint{8.036091in}{4.050000in}}{\pgfqpoint{5.400000in}{5.400000in}}%
\pgfusepath{clip}%
\pgfsetbuttcap%
\pgfsetroundjoin%
\pgfsetlinewidth{1.003750pt}%
\definecolor{currentstroke}{rgb}{0.321569,0.321569,0.321569}%
\pgfsetstrokecolor{currentstroke}%
\pgfsetstrokeopacity{0.900000}%
\pgfsetdash{}{0pt}%
\pgfpathmoveto{\pgfqpoint{12.387293in}{8.805295in}}%
\pgfpathcurveto{\pgfqpoint{12.397547in}{8.805295in}}{\pgfqpoint{12.407382in}{8.809369in}}{\pgfqpoint{12.414633in}{8.816620in}}%
\pgfpathcurveto{\pgfqpoint{12.421884in}{8.823871in}}{\pgfqpoint{12.425958in}{8.833706in}}{\pgfqpoint{12.425958in}{8.843960in}}%
\pgfpathcurveto{\pgfqpoint{12.425958in}{8.854214in}}{\pgfqpoint{12.421884in}{8.864050in}}{\pgfqpoint{12.414633in}{8.871300in}}%
\pgfpathcurveto{\pgfqpoint{12.407382in}{8.878551in}}{\pgfqpoint{12.397547in}{8.882625in}}{\pgfqpoint{12.387293in}{8.882625in}}%
\pgfpathcurveto{\pgfqpoint{12.377039in}{8.882625in}}{\pgfqpoint{12.367203in}{8.878551in}}{\pgfqpoint{12.359952in}{8.871300in}}%
\pgfpathcurveto{\pgfqpoint{12.352702in}{8.864050in}}{\pgfqpoint{12.348628in}{8.854214in}}{\pgfqpoint{12.348628in}{8.843960in}}%
\pgfpathcurveto{\pgfqpoint{12.348628in}{8.833706in}}{\pgfqpoint{12.352702in}{8.823871in}}{\pgfqpoint{12.359952in}{8.816620in}}%
\pgfpathcurveto{\pgfqpoint{12.367203in}{8.809369in}}{\pgfqpoint{12.377039in}{8.805295in}}{\pgfqpoint{12.387293in}{8.805295in}}%
\pgfpathlineto{\pgfqpoint{12.387293in}{8.805295in}}%
\pgfpathclose%
\pgfusepath{stroke}%
\end{pgfscope}%
\begin{pgfscope}%
\pgfpathrectangle{\pgfqpoint{8.036091in}{4.050000in}}{\pgfqpoint{5.400000in}{5.400000in}}%
\pgfusepath{clip}%
\pgfsetbuttcap%
\pgfsetroundjoin%
\pgfsetlinewidth{1.003750pt}%
\definecolor{currentstroke}{rgb}{0.321569,0.321569,0.321569}%
\pgfsetstrokecolor{currentstroke}%
\pgfsetstrokeopacity{0.900000}%
\pgfsetdash{}{0pt}%
\pgfpathmoveto{\pgfqpoint{10.051652in}{5.910447in}}%
\pgfpathcurveto{\pgfqpoint{10.061906in}{5.910447in}}{\pgfqpoint{10.071741in}{5.914521in}}{\pgfqpoint{10.078992in}{5.921772in}}%
\pgfpathcurveto{\pgfqpoint{10.086243in}{5.929022in}}{\pgfqpoint{10.090317in}{5.938858in}}{\pgfqpoint{10.090317in}{5.949112in}}%
\pgfpathcurveto{\pgfqpoint{10.090317in}{5.959366in}}{\pgfqpoint{10.086243in}{5.969202in}}{\pgfqpoint{10.078992in}{5.976452in}}%
\pgfpathcurveto{\pgfqpoint{10.071741in}{5.983703in}}{\pgfqpoint{10.061906in}{5.987777in}}{\pgfqpoint{10.051652in}{5.987777in}}%
\pgfpathcurveto{\pgfqpoint{10.041398in}{5.987777in}}{\pgfqpoint{10.031562in}{5.983703in}}{\pgfqpoint{10.024311in}{5.976452in}}%
\pgfpathcurveto{\pgfqpoint{10.017061in}{5.969202in}}{\pgfqpoint{10.012987in}{5.959366in}}{\pgfqpoint{10.012987in}{5.949112in}}%
\pgfpathcurveto{\pgfqpoint{10.012987in}{5.938858in}}{\pgfqpoint{10.017061in}{5.929022in}}{\pgfqpoint{10.024311in}{5.921772in}}%
\pgfpathcurveto{\pgfqpoint{10.031562in}{5.914521in}}{\pgfqpoint{10.041398in}{5.910447in}}{\pgfqpoint{10.051652in}{5.910447in}}%
\pgfpathlineto{\pgfqpoint{10.051652in}{5.910447in}}%
\pgfpathclose%
\pgfusepath{stroke}%
\end{pgfscope}%
\begin{pgfscope}%
\pgfpathrectangle{\pgfqpoint{8.036091in}{4.050000in}}{\pgfqpoint{5.400000in}{5.400000in}}%
\pgfusepath{clip}%
\pgfsetbuttcap%
\pgfsetroundjoin%
\pgfsetlinewidth{1.003750pt}%
\definecolor{currentstroke}{rgb}{0.321569,0.321569,0.321569}%
\pgfsetstrokecolor{currentstroke}%
\pgfsetstrokeopacity{0.900000}%
\pgfsetdash{}{0pt}%
\pgfpathmoveto{\pgfqpoint{10.062616in}{6.024285in}}%
\pgfpathcurveto{\pgfqpoint{10.072870in}{6.024285in}}{\pgfqpoint{10.082706in}{6.028359in}}{\pgfqpoint{10.089957in}{6.035610in}}%
\pgfpathcurveto{\pgfqpoint{10.097207in}{6.042861in}}{\pgfqpoint{10.101281in}{6.052696in}}{\pgfqpoint{10.101281in}{6.062950in}}%
\pgfpathcurveto{\pgfqpoint{10.101281in}{6.073204in}}{\pgfqpoint{10.097207in}{6.083040in}}{\pgfqpoint{10.089957in}{6.090291in}}%
\pgfpathcurveto{\pgfqpoint{10.082706in}{6.097541in}}{\pgfqpoint{10.072870in}{6.101615in}}{\pgfqpoint{10.062616in}{6.101615in}}%
\pgfpathcurveto{\pgfqpoint{10.052362in}{6.101615in}}{\pgfqpoint{10.042527in}{6.097541in}}{\pgfqpoint{10.035276in}{6.090291in}}%
\pgfpathcurveto{\pgfqpoint{10.028025in}{6.083040in}}{\pgfqpoint{10.023951in}{6.073204in}}{\pgfqpoint{10.023951in}{6.062950in}}%
\pgfpathcurveto{\pgfqpoint{10.023951in}{6.052696in}}{\pgfqpoint{10.028025in}{6.042861in}}{\pgfqpoint{10.035276in}{6.035610in}}%
\pgfpathcurveto{\pgfqpoint{10.042527in}{6.028359in}}{\pgfqpoint{10.052362in}{6.024285in}}{\pgfqpoint{10.062616in}{6.024285in}}%
\pgfpathlineto{\pgfqpoint{10.062616in}{6.024285in}}%
\pgfpathclose%
\pgfusepath{stroke}%
\end{pgfscope}%
\begin{pgfscope}%
\pgfpathrectangle{\pgfqpoint{8.036091in}{4.050000in}}{\pgfqpoint{5.400000in}{5.400000in}}%
\pgfusepath{clip}%
\pgfsetbuttcap%
\pgfsetroundjoin%
\pgfsetlinewidth{1.003750pt}%
\definecolor{currentstroke}{rgb}{0.321569,0.321569,0.321569}%
\pgfsetstrokecolor{currentstroke}%
\pgfsetstrokeopacity{0.900000}%
\pgfsetdash{}{0pt}%
\pgfpathmoveto{\pgfqpoint{9.331843in}{5.352726in}}%
\pgfpathcurveto{\pgfqpoint{9.342097in}{5.352726in}}{\pgfqpoint{9.351933in}{5.356800in}}{\pgfqpoint{9.359184in}{5.364051in}}%
\pgfpathcurveto{\pgfqpoint{9.366434in}{5.371301in}}{\pgfqpoint{9.370508in}{5.381137in}}{\pgfqpoint{9.370508in}{5.391391in}}%
\pgfpathcurveto{\pgfqpoint{9.370508in}{5.401645in}}{\pgfqpoint{9.366434in}{5.411481in}}{\pgfqpoint{9.359184in}{5.418731in}}%
\pgfpathcurveto{\pgfqpoint{9.351933in}{5.425982in}}{\pgfqpoint{9.342097in}{5.430056in}}{\pgfqpoint{9.331843in}{5.430056in}}%
\pgfpathcurveto{\pgfqpoint{9.321589in}{5.430056in}}{\pgfqpoint{9.311754in}{5.425982in}}{\pgfqpoint{9.304503in}{5.418731in}}%
\pgfpathcurveto{\pgfqpoint{9.297252in}{5.411481in}}{\pgfqpoint{9.293178in}{5.401645in}}{\pgfqpoint{9.293178in}{5.391391in}}%
\pgfpathcurveto{\pgfqpoint{9.293178in}{5.381137in}}{\pgfqpoint{9.297252in}{5.371301in}}{\pgfqpoint{9.304503in}{5.364051in}}%
\pgfpathcurveto{\pgfqpoint{9.311754in}{5.356800in}}{\pgfqpoint{9.321589in}{5.352726in}}{\pgfqpoint{9.331843in}{5.352726in}}%
\pgfpathlineto{\pgfqpoint{9.331843in}{5.352726in}}%
\pgfpathclose%
\pgfusepath{stroke}%
\end{pgfscope}%
\begin{pgfscope}%
\pgfpathrectangle{\pgfqpoint{8.036091in}{4.050000in}}{\pgfqpoint{5.400000in}{5.400000in}}%
\pgfusepath{clip}%
\pgfsetbuttcap%
\pgfsetroundjoin%
\pgfsetlinewidth{1.003750pt}%
\definecolor{currentstroke}{rgb}{0.321569,0.321569,0.321569}%
\pgfsetstrokecolor{currentstroke}%
\pgfsetstrokeopacity{0.900000}%
\pgfsetdash{}{0pt}%
\pgfpathmoveto{\pgfqpoint{9.603880in}{5.599630in}}%
\pgfpathcurveto{\pgfqpoint{9.614134in}{5.599630in}}{\pgfqpoint{9.623969in}{5.603704in}}{\pgfqpoint{9.631220in}{5.610954in}}%
\pgfpathcurveto{\pgfqpoint{9.638471in}{5.618205in}}{\pgfqpoint{9.642545in}{5.628041in}}{\pgfqpoint{9.642545in}{5.638295in}}%
\pgfpathcurveto{\pgfqpoint{9.642545in}{5.648549in}}{\pgfqpoint{9.638471in}{5.658384in}}{\pgfqpoint{9.631220in}{5.665635in}}%
\pgfpathcurveto{\pgfqpoint{9.623969in}{5.672886in}}{\pgfqpoint{9.614134in}{5.676960in}}{\pgfqpoint{9.603880in}{5.676960in}}%
\pgfpathcurveto{\pgfqpoint{9.593625in}{5.676960in}}{\pgfqpoint{9.583790in}{5.672886in}}{\pgfqpoint{9.576539in}{5.665635in}}%
\pgfpathcurveto{\pgfqpoint{9.569289in}{5.658384in}}{\pgfqpoint{9.565215in}{5.648549in}}{\pgfqpoint{9.565215in}{5.638295in}}%
\pgfpathcurveto{\pgfqpoint{9.565215in}{5.628041in}}{\pgfqpoint{9.569289in}{5.618205in}}{\pgfqpoint{9.576539in}{5.610954in}}%
\pgfpathcurveto{\pgfqpoint{9.583790in}{5.603704in}}{\pgfqpoint{9.593625in}{5.599630in}}{\pgfqpoint{9.603880in}{5.599630in}}%
\pgfpathlineto{\pgfqpoint{9.603880in}{5.599630in}}%
\pgfpathclose%
\pgfusepath{stroke}%
\end{pgfscope}%
\begin{pgfscope}%
\pgfpathrectangle{\pgfqpoint{8.036091in}{4.050000in}}{\pgfqpoint{5.400000in}{5.400000in}}%
\pgfusepath{clip}%
\pgfsetbuttcap%
\pgfsetroundjoin%
\pgfsetlinewidth{1.003750pt}%
\definecolor{currentstroke}{rgb}{0.321569,0.321569,0.321569}%
\pgfsetstrokecolor{currentstroke}%
\pgfsetstrokeopacity{0.900000}%
\pgfsetdash{}{0pt}%
\pgfpathmoveto{\pgfqpoint{9.007897in}{5.075690in}}%
\pgfpathcurveto{\pgfqpoint{9.018151in}{5.075690in}}{\pgfqpoint{9.027987in}{5.079764in}}{\pgfqpoint{9.035237in}{5.087014in}}%
\pgfpathcurveto{\pgfqpoint{9.042488in}{5.094265in}}{\pgfqpoint{9.046562in}{5.104101in}}{\pgfqpoint{9.046562in}{5.114355in}}%
\pgfpathcurveto{\pgfqpoint{9.046562in}{5.124609in}}{\pgfqpoint{9.042488in}{5.134444in}}{\pgfqpoint{9.035237in}{5.141695in}}%
\pgfpathcurveto{\pgfqpoint{9.027987in}{5.148946in}}{\pgfqpoint{9.018151in}{5.153020in}}{\pgfqpoint{9.007897in}{5.153020in}}%
\pgfpathcurveto{\pgfqpoint{8.997643in}{5.153020in}}{\pgfqpoint{8.987807in}{5.148946in}}{\pgfqpoint{8.980557in}{5.141695in}}%
\pgfpathcurveto{\pgfqpoint{8.973306in}{5.134444in}}{\pgfqpoint{8.969232in}{5.124609in}}{\pgfqpoint{8.969232in}{5.114355in}}%
\pgfpathcurveto{\pgfqpoint{8.969232in}{5.104101in}}{\pgfqpoint{8.973306in}{5.094265in}}{\pgfqpoint{8.980557in}{5.087014in}}%
\pgfpathcurveto{\pgfqpoint{8.987807in}{5.079764in}}{\pgfqpoint{8.997643in}{5.075690in}}{\pgfqpoint{9.007897in}{5.075690in}}%
\pgfpathlineto{\pgfqpoint{9.007897in}{5.075690in}}%
\pgfpathclose%
\pgfusepath{stroke}%
\end{pgfscope}%
\begin{pgfscope}%
\pgfpathrectangle{\pgfqpoint{8.036091in}{4.050000in}}{\pgfqpoint{5.400000in}{5.400000in}}%
\pgfusepath{clip}%
\pgfsetbuttcap%
\pgfsetroundjoin%
\pgfsetlinewidth{1.003750pt}%
\definecolor{currentstroke}{rgb}{0.321569,0.321569,0.321569}%
\pgfsetstrokecolor{currentstroke}%
\pgfsetstrokeopacity{0.900000}%
\pgfsetdash{}{0pt}%
\pgfpathmoveto{\pgfqpoint{9.301434in}{5.307745in}}%
\pgfpathcurveto{\pgfqpoint{9.311688in}{5.307745in}}{\pgfqpoint{9.321524in}{5.311819in}}{\pgfqpoint{9.328775in}{5.319069in}}%
\pgfpathcurveto{\pgfqpoint{9.336025in}{5.326320in}}{\pgfqpoint{9.340099in}{5.336156in}}{\pgfqpoint{9.340099in}{5.346410in}}%
\pgfpathcurveto{\pgfqpoint{9.340099in}{5.356664in}}{\pgfqpoint{9.336025in}{5.366499in}}{\pgfqpoint{9.328775in}{5.373750in}}%
\pgfpathcurveto{\pgfqpoint{9.321524in}{5.381001in}}{\pgfqpoint{9.311688in}{5.385075in}}{\pgfqpoint{9.301434in}{5.385075in}}%
\pgfpathcurveto{\pgfqpoint{9.291180in}{5.385075in}}{\pgfqpoint{9.281345in}{5.381001in}}{\pgfqpoint{9.274094in}{5.373750in}}%
\pgfpathcurveto{\pgfqpoint{9.266843in}{5.366499in}}{\pgfqpoint{9.262769in}{5.356664in}}{\pgfqpoint{9.262769in}{5.346410in}}%
\pgfpathcurveto{\pgfqpoint{9.262769in}{5.336156in}}{\pgfqpoint{9.266843in}{5.326320in}}{\pgfqpoint{9.274094in}{5.319069in}}%
\pgfpathcurveto{\pgfqpoint{9.281345in}{5.311819in}}{\pgfqpoint{9.291180in}{5.307745in}}{\pgfqpoint{9.301434in}{5.307745in}}%
\pgfpathlineto{\pgfqpoint{9.301434in}{5.307745in}}%
\pgfpathclose%
\pgfusepath{stroke}%
\end{pgfscope}%
\begin{pgfscope}%
\pgfpathrectangle{\pgfqpoint{8.036091in}{4.050000in}}{\pgfqpoint{5.400000in}{5.400000in}}%
\pgfusepath{clip}%
\pgfsetbuttcap%
\pgfsetroundjoin%
\pgfsetlinewidth{1.003750pt}%
\definecolor{currentstroke}{rgb}{0.321569,0.321569,0.321569}%
\pgfsetstrokecolor{currentstroke}%
\pgfsetstrokeopacity{0.900000}%
\pgfsetdash{}{0pt}%
\pgfpathmoveto{\pgfqpoint{9.809521in}{5.582491in}}%
\pgfpathcurveto{\pgfqpoint{9.819776in}{5.582491in}}{\pgfqpoint{9.829611in}{5.586565in}}{\pgfqpoint{9.836862in}{5.593816in}}%
\pgfpathcurveto{\pgfqpoint{9.844112in}{5.601067in}}{\pgfqpoint{9.848186in}{5.610902in}}{\pgfqpoint{9.848186in}{5.621156in}}%
\pgfpathcurveto{\pgfqpoint{9.848186in}{5.631411in}}{\pgfqpoint{9.844112in}{5.641246in}}{\pgfqpoint{9.836862in}{5.648497in}}%
\pgfpathcurveto{\pgfqpoint{9.829611in}{5.655747in}}{\pgfqpoint{9.819776in}{5.659821in}}{\pgfqpoint{9.809521in}{5.659821in}}%
\pgfpathcurveto{\pgfqpoint{9.799267in}{5.659821in}}{\pgfqpoint{9.789432in}{5.655747in}}{\pgfqpoint{9.782181in}{5.648497in}}%
\pgfpathcurveto{\pgfqpoint{9.774930in}{5.641246in}}{\pgfqpoint{9.770856in}{5.631411in}}{\pgfqpoint{9.770856in}{5.621156in}}%
\pgfpathcurveto{\pgfqpoint{9.770856in}{5.610902in}}{\pgfqpoint{9.774930in}{5.601067in}}{\pgfqpoint{9.782181in}{5.593816in}}%
\pgfpathcurveto{\pgfqpoint{9.789432in}{5.586565in}}{\pgfqpoint{9.799267in}{5.582491in}}{\pgfqpoint{9.809521in}{5.582491in}}%
\pgfpathlineto{\pgfqpoint{9.809521in}{5.582491in}}%
\pgfpathclose%
\pgfusepath{stroke}%
\end{pgfscope}%
\begin{pgfscope}%
\pgfpathrectangle{\pgfqpoint{8.036091in}{4.050000in}}{\pgfqpoint{5.400000in}{5.400000in}}%
\pgfusepath{clip}%
\pgfsetbuttcap%
\pgfsetroundjoin%
\pgfsetlinewidth{1.003750pt}%
\definecolor{currentstroke}{rgb}{0.321569,0.321569,0.321569}%
\pgfsetstrokecolor{currentstroke}%
\pgfsetstrokeopacity{0.900000}%
\pgfsetdash{}{0pt}%
\pgfpathmoveto{\pgfqpoint{9.332409in}{5.621097in}}%
\pgfpathcurveto{\pgfqpoint{9.342664in}{5.621097in}}{\pgfqpoint{9.352499in}{5.625171in}}{\pgfqpoint{9.359750in}{5.632421in}}%
\pgfpathcurveto{\pgfqpoint{9.367001in}{5.639672in}}{\pgfqpoint{9.371075in}{5.649508in}}{\pgfqpoint{9.371075in}{5.659762in}}%
\pgfpathcurveto{\pgfqpoint{9.371075in}{5.670016in}}{\pgfqpoint{9.367001in}{5.679851in}}{\pgfqpoint{9.359750in}{5.687102in}}%
\pgfpathcurveto{\pgfqpoint{9.352499in}{5.694353in}}{\pgfqpoint{9.342664in}{5.698427in}}{\pgfqpoint{9.332409in}{5.698427in}}%
\pgfpathcurveto{\pgfqpoint{9.322155in}{5.698427in}}{\pgfqpoint{9.312320in}{5.694353in}}{\pgfqpoint{9.305069in}{5.687102in}}%
\pgfpathcurveto{\pgfqpoint{9.297818in}{5.679851in}}{\pgfqpoint{9.293744in}{5.670016in}}{\pgfqpoint{9.293744in}{5.659762in}}%
\pgfpathcurveto{\pgfqpoint{9.293744in}{5.649508in}}{\pgfqpoint{9.297818in}{5.639672in}}{\pgfqpoint{9.305069in}{5.632421in}}%
\pgfpathcurveto{\pgfqpoint{9.312320in}{5.625171in}}{\pgfqpoint{9.322155in}{5.621097in}}{\pgfqpoint{9.332409in}{5.621097in}}%
\pgfpathlineto{\pgfqpoint{9.332409in}{5.621097in}}%
\pgfpathclose%
\pgfusepath{stroke}%
\end{pgfscope}%
\begin{pgfscope}%
\pgfpathrectangle{\pgfqpoint{8.036091in}{4.050000in}}{\pgfqpoint{5.400000in}{5.400000in}}%
\pgfusepath{clip}%
\pgfsetbuttcap%
\pgfsetroundjoin%
\pgfsetlinewidth{1.003750pt}%
\definecolor{currentstroke}{rgb}{0.321569,0.321569,0.321569}%
\pgfsetstrokecolor{currentstroke}%
\pgfsetstrokeopacity{0.900000}%
\pgfsetdash{}{0pt}%
\pgfpathmoveto{\pgfqpoint{9.455716in}{5.650831in}}%
\pgfpathcurveto{\pgfqpoint{9.465970in}{5.650831in}}{\pgfqpoint{9.475806in}{5.654905in}}{\pgfqpoint{9.483056in}{5.662156in}}%
\pgfpathcurveto{\pgfqpoint{9.490307in}{5.669406in}}{\pgfqpoint{9.494381in}{5.679242in}}{\pgfqpoint{9.494381in}{5.689496in}}%
\pgfpathcurveto{\pgfqpoint{9.494381in}{5.699750in}}{\pgfqpoint{9.490307in}{5.709586in}}{\pgfqpoint{9.483056in}{5.716836in}}%
\pgfpathcurveto{\pgfqpoint{9.475806in}{5.724087in}}{\pgfqpoint{9.465970in}{5.728161in}}{\pgfqpoint{9.455716in}{5.728161in}}%
\pgfpathcurveto{\pgfqpoint{9.445462in}{5.728161in}}{\pgfqpoint{9.435626in}{5.724087in}}{\pgfqpoint{9.428376in}{5.716836in}}%
\pgfpathcurveto{\pgfqpoint{9.421125in}{5.709586in}}{\pgfqpoint{9.417051in}{5.699750in}}{\pgfqpoint{9.417051in}{5.689496in}}%
\pgfpathcurveto{\pgfqpoint{9.417051in}{5.679242in}}{\pgfqpoint{9.421125in}{5.669406in}}{\pgfqpoint{9.428376in}{5.662156in}}%
\pgfpathcurveto{\pgfqpoint{9.435626in}{5.654905in}}{\pgfqpoint{9.445462in}{5.650831in}}{\pgfqpoint{9.455716in}{5.650831in}}%
\pgfpathlineto{\pgfqpoint{9.455716in}{5.650831in}}%
\pgfpathclose%
\pgfusepath{stroke}%
\end{pgfscope}%
\begin{pgfscope}%
\pgfpathrectangle{\pgfqpoint{8.036091in}{4.050000in}}{\pgfqpoint{5.400000in}{5.400000in}}%
\pgfusepath{clip}%
\pgfsetbuttcap%
\pgfsetroundjoin%
\pgfsetlinewidth{1.003750pt}%
\definecolor{currentstroke}{rgb}{0.321569,0.321569,0.321569}%
\pgfsetstrokecolor{currentstroke}%
\pgfsetstrokeopacity{0.900000}%
\pgfsetdash{}{0pt}%
\pgfpathmoveto{\pgfqpoint{9.153401in}{5.222809in}}%
\pgfpathcurveto{\pgfqpoint{9.163655in}{5.222809in}}{\pgfqpoint{9.173490in}{5.226883in}}{\pgfqpoint{9.180741in}{5.234134in}}%
\pgfpathcurveto{\pgfqpoint{9.187992in}{5.241385in}}{\pgfqpoint{9.192066in}{5.251220in}}{\pgfqpoint{9.192066in}{5.261474in}}%
\pgfpathcurveto{\pgfqpoint{9.192066in}{5.271728in}}{\pgfqpoint{9.187992in}{5.281564in}}{\pgfqpoint{9.180741in}{5.288814in}}%
\pgfpathcurveto{\pgfqpoint{9.173490in}{5.296065in}}{\pgfqpoint{9.163655in}{5.300139in}}{\pgfqpoint{9.153401in}{5.300139in}}%
\pgfpathcurveto{\pgfqpoint{9.143147in}{5.300139in}}{\pgfqpoint{9.133311in}{5.296065in}}{\pgfqpoint{9.126060in}{5.288814in}}%
\pgfpathcurveto{\pgfqpoint{9.118810in}{5.281564in}}{\pgfqpoint{9.114736in}{5.271728in}}{\pgfqpoint{9.114736in}{5.261474in}}%
\pgfpathcurveto{\pgfqpoint{9.114736in}{5.251220in}}{\pgfqpoint{9.118810in}{5.241385in}}{\pgfqpoint{9.126060in}{5.234134in}}%
\pgfpathcurveto{\pgfqpoint{9.133311in}{5.226883in}}{\pgfqpoint{9.143147in}{5.222809in}}{\pgfqpoint{9.153401in}{5.222809in}}%
\pgfpathlineto{\pgfqpoint{9.153401in}{5.222809in}}%
\pgfpathclose%
\pgfusepath{stroke}%
\end{pgfscope}%
\begin{pgfscope}%
\pgfpathrectangle{\pgfqpoint{8.036091in}{4.050000in}}{\pgfqpoint{5.400000in}{5.400000in}}%
\pgfusepath{clip}%
\pgfsetbuttcap%
\pgfsetroundjoin%
\pgfsetlinewidth{1.003750pt}%
\definecolor{currentstroke}{rgb}{0.321569,0.321569,0.321569}%
\pgfsetstrokecolor{currentstroke}%
\pgfsetstrokeopacity{0.900000}%
\pgfsetdash{}{0pt}%
\pgfpathmoveto{\pgfqpoint{9.005024in}{5.153847in}}%
\pgfpathcurveto{\pgfqpoint{9.015278in}{5.153847in}}{\pgfqpoint{9.025114in}{5.157921in}}{\pgfqpoint{9.032365in}{5.165171in}}%
\pgfpathcurveto{\pgfqpoint{9.039615in}{5.172422in}}{\pgfqpoint{9.043689in}{5.182258in}}{\pgfqpoint{9.043689in}{5.192512in}}%
\pgfpathcurveto{\pgfqpoint{9.043689in}{5.202766in}}{\pgfqpoint{9.039615in}{5.212601in}}{\pgfqpoint{9.032365in}{5.219852in}}%
\pgfpathcurveto{\pgfqpoint{9.025114in}{5.227103in}}{\pgfqpoint{9.015278in}{5.231177in}}{\pgfqpoint{9.005024in}{5.231177in}}%
\pgfpathcurveto{\pgfqpoint{8.994770in}{5.231177in}}{\pgfqpoint{8.984935in}{5.227103in}}{\pgfqpoint{8.977684in}{5.219852in}}%
\pgfpathcurveto{\pgfqpoint{8.970433in}{5.212601in}}{\pgfqpoint{8.966359in}{5.202766in}}{\pgfqpoint{8.966359in}{5.192512in}}%
\pgfpathcurveto{\pgfqpoint{8.966359in}{5.182258in}}{\pgfqpoint{8.970433in}{5.172422in}}{\pgfqpoint{8.977684in}{5.165171in}}%
\pgfpathcurveto{\pgfqpoint{8.984935in}{5.157921in}}{\pgfqpoint{8.994770in}{5.153847in}}{\pgfqpoint{9.005024in}{5.153847in}}%
\pgfpathlineto{\pgfqpoint{9.005024in}{5.153847in}}%
\pgfpathclose%
\pgfusepath{stroke}%
\end{pgfscope}%
\begin{pgfscope}%
\pgfpathrectangle{\pgfqpoint{8.036091in}{4.050000in}}{\pgfqpoint{5.400000in}{5.400000in}}%
\pgfusepath{clip}%
\pgfsetbuttcap%
\pgfsetroundjoin%
\pgfsetlinewidth{1.003750pt}%
\definecolor{currentstroke}{rgb}{0.321569,0.321569,0.321569}%
\pgfsetstrokecolor{currentstroke}%
\pgfsetstrokeopacity{0.900000}%
\pgfsetdash{}{0pt}%
\pgfpathmoveto{\pgfqpoint{11.646050in}{7.324405in}}%
\pgfpathcurveto{\pgfqpoint{11.656304in}{7.324405in}}{\pgfqpoint{11.666140in}{7.328479in}}{\pgfqpoint{11.673390in}{7.335730in}}%
\pgfpathcurveto{\pgfqpoint{11.680641in}{7.342980in}}{\pgfqpoint{11.684715in}{7.352816in}}{\pgfqpoint{11.684715in}{7.363070in}}%
\pgfpathcurveto{\pgfqpoint{11.684715in}{7.373324in}}{\pgfqpoint{11.680641in}{7.383160in}}{\pgfqpoint{11.673390in}{7.390410in}}%
\pgfpathcurveto{\pgfqpoint{11.666140in}{7.397661in}}{\pgfqpoint{11.656304in}{7.401735in}}{\pgfqpoint{11.646050in}{7.401735in}}%
\pgfpathcurveto{\pgfqpoint{11.635796in}{7.401735in}}{\pgfqpoint{11.625961in}{7.397661in}}{\pgfqpoint{11.618710in}{7.390410in}}%
\pgfpathcurveto{\pgfqpoint{11.611459in}{7.383160in}}{\pgfqpoint{11.607385in}{7.373324in}}{\pgfqpoint{11.607385in}{7.363070in}}%
\pgfpathcurveto{\pgfqpoint{11.607385in}{7.352816in}}{\pgfqpoint{11.611459in}{7.342980in}}{\pgfqpoint{11.618710in}{7.335730in}}%
\pgfpathcurveto{\pgfqpoint{11.625961in}{7.328479in}}{\pgfqpoint{11.635796in}{7.324405in}}{\pgfqpoint{11.646050in}{7.324405in}}%
\pgfpathlineto{\pgfqpoint{11.646050in}{7.324405in}}%
\pgfpathclose%
\pgfusepath{stroke}%
\end{pgfscope}%
\begin{pgfscope}%
\pgfpathrectangle{\pgfqpoint{8.036091in}{4.050000in}}{\pgfqpoint{5.400000in}{5.400000in}}%
\pgfusepath{clip}%
\pgfsetbuttcap%
\pgfsetroundjoin%
\pgfsetlinewidth{1.003750pt}%
\definecolor{currentstroke}{rgb}{0.321569,0.321569,0.321569}%
\pgfsetstrokecolor{currentstroke}%
\pgfsetstrokeopacity{0.900000}%
\pgfsetdash{}{0pt}%
\pgfpathmoveto{\pgfqpoint{11.766923in}{7.428621in}}%
\pgfpathcurveto{\pgfqpoint{11.777177in}{7.428621in}}{\pgfqpoint{11.787012in}{7.432695in}}{\pgfqpoint{11.794263in}{7.439946in}}%
\pgfpathcurveto{\pgfqpoint{11.801514in}{7.447197in}}{\pgfqpoint{11.805588in}{7.457032in}}{\pgfqpoint{11.805588in}{7.467286in}}%
\pgfpathcurveto{\pgfqpoint{11.805588in}{7.477540in}}{\pgfqpoint{11.801514in}{7.487376in}}{\pgfqpoint{11.794263in}{7.494627in}}%
\pgfpathcurveto{\pgfqpoint{11.787012in}{7.501877in}}{\pgfqpoint{11.777177in}{7.505951in}}{\pgfqpoint{11.766923in}{7.505951in}}%
\pgfpathcurveto{\pgfqpoint{11.756669in}{7.505951in}}{\pgfqpoint{11.746833in}{7.501877in}}{\pgfqpoint{11.739583in}{7.494627in}}%
\pgfpathcurveto{\pgfqpoint{11.732332in}{7.487376in}}{\pgfqpoint{11.728258in}{7.477540in}}{\pgfqpoint{11.728258in}{7.467286in}}%
\pgfpathcurveto{\pgfqpoint{11.728258in}{7.457032in}}{\pgfqpoint{11.732332in}{7.447197in}}{\pgfqpoint{11.739583in}{7.439946in}}%
\pgfpathcurveto{\pgfqpoint{11.746833in}{7.432695in}}{\pgfqpoint{11.756669in}{7.428621in}}{\pgfqpoint{11.766923in}{7.428621in}}%
\pgfpathlineto{\pgfqpoint{11.766923in}{7.428621in}}%
\pgfpathclose%
\pgfusepath{stroke}%
\end{pgfscope}%
\begin{pgfscope}%
\pgfpathrectangle{\pgfqpoint{8.036091in}{4.050000in}}{\pgfqpoint{5.400000in}{5.400000in}}%
\pgfusepath{clip}%
\pgfsetbuttcap%
\pgfsetroundjoin%
\pgfsetlinewidth{1.003750pt}%
\definecolor{currentstroke}{rgb}{0.321569,0.321569,0.321569}%
\pgfsetstrokecolor{currentstroke}%
\pgfsetstrokeopacity{0.900000}%
\pgfsetdash{}{0pt}%
\pgfpathmoveto{\pgfqpoint{11.868360in}{7.527313in}}%
\pgfpathcurveto{\pgfqpoint{11.878614in}{7.527313in}}{\pgfqpoint{11.888450in}{7.531387in}}{\pgfqpoint{11.895700in}{7.538638in}}%
\pgfpathcurveto{\pgfqpoint{11.902951in}{7.545889in}}{\pgfqpoint{11.907025in}{7.555724in}}{\pgfqpoint{11.907025in}{7.565978in}}%
\pgfpathcurveto{\pgfqpoint{11.907025in}{7.576232in}}{\pgfqpoint{11.902951in}{7.586068in}}{\pgfqpoint{11.895700in}{7.593319in}}%
\pgfpathcurveto{\pgfqpoint{11.888450in}{7.600569in}}{\pgfqpoint{11.878614in}{7.604643in}}{\pgfqpoint{11.868360in}{7.604643in}}%
\pgfpathcurveto{\pgfqpoint{11.858106in}{7.604643in}}{\pgfqpoint{11.848271in}{7.600569in}}{\pgfqpoint{11.841020in}{7.593319in}}%
\pgfpathcurveto{\pgfqpoint{11.833769in}{7.586068in}}{\pgfqpoint{11.829695in}{7.576232in}}{\pgfqpoint{11.829695in}{7.565978in}}%
\pgfpathcurveto{\pgfqpoint{11.829695in}{7.555724in}}{\pgfqpoint{11.833769in}{7.545889in}}{\pgfqpoint{11.841020in}{7.538638in}}%
\pgfpathcurveto{\pgfqpoint{11.848271in}{7.531387in}}{\pgfqpoint{11.858106in}{7.527313in}}{\pgfqpoint{11.868360in}{7.527313in}}%
\pgfpathlineto{\pgfqpoint{11.868360in}{7.527313in}}%
\pgfpathclose%
\pgfusepath{stroke}%
\end{pgfscope}%
\begin{pgfscope}%
\pgfpathrectangle{\pgfqpoint{8.036091in}{4.050000in}}{\pgfqpoint{5.400000in}{5.400000in}}%
\pgfusepath{clip}%
\pgfsetbuttcap%
\pgfsetroundjoin%
\pgfsetlinewidth{1.003750pt}%
\definecolor{currentstroke}{rgb}{0.321569,0.321569,0.321569}%
\pgfsetstrokecolor{currentstroke}%
\pgfsetstrokeopacity{0.900000}%
\pgfsetdash{}{0pt}%
\pgfpathmoveto{\pgfqpoint{12.113899in}{7.797045in}}%
\pgfpathcurveto{\pgfqpoint{12.124153in}{7.797045in}}{\pgfqpoint{12.133988in}{7.801119in}}{\pgfqpoint{12.141239in}{7.808370in}}%
\pgfpathcurveto{\pgfqpoint{12.148490in}{7.815621in}}{\pgfqpoint{12.152564in}{7.825456in}}{\pgfqpoint{12.152564in}{7.835710in}}%
\pgfpathcurveto{\pgfqpoint{12.152564in}{7.845964in}}{\pgfqpoint{12.148490in}{7.855800in}}{\pgfqpoint{12.141239in}{7.863051in}}%
\pgfpathcurveto{\pgfqpoint{12.133988in}{7.870301in}}{\pgfqpoint{12.124153in}{7.874375in}}{\pgfqpoint{12.113899in}{7.874375in}}%
\pgfpathcurveto{\pgfqpoint{12.103645in}{7.874375in}}{\pgfqpoint{12.093809in}{7.870301in}}{\pgfqpoint{12.086558in}{7.863051in}}%
\pgfpathcurveto{\pgfqpoint{12.079308in}{7.855800in}}{\pgfqpoint{12.075234in}{7.845964in}}{\pgfqpoint{12.075234in}{7.835710in}}%
\pgfpathcurveto{\pgfqpoint{12.075234in}{7.825456in}}{\pgfqpoint{12.079308in}{7.815621in}}{\pgfqpoint{12.086558in}{7.808370in}}%
\pgfpathcurveto{\pgfqpoint{12.093809in}{7.801119in}}{\pgfqpoint{12.103645in}{7.797045in}}{\pgfqpoint{12.113899in}{7.797045in}}%
\pgfpathlineto{\pgfqpoint{12.113899in}{7.797045in}}%
\pgfpathclose%
\pgfusepath{stroke}%
\end{pgfscope}%
\begin{pgfscope}%
\pgfsetbuttcap%
\pgfsetroundjoin%
\pgfsetlinewidth{1.003750pt}%
\definecolor{currentstroke}{rgb}{0.321569,0.321569,0.321569}%
\pgfsetstrokecolor{currentstroke}%
\pgfsetstrokeopacity{0.900000}%
\pgfsetdash{}{0pt}%
\pgfpathmoveto{\pgfqpoint{8.036091in}{3.990262in}}%
\pgfpathlineto{\pgfqpoint{8.095829in}{4.050000in}}%
\pgfpathlineto{\pgfqpoint{8.036091in}{4.109738in}}%
\pgfpathlineto{\pgfqpoint{7.976353in}{4.050000in}}%
\pgfpathlineto{\pgfqpoint{8.036091in}{3.990262in}}%
\pgfpathclose%
\pgfusepath{stroke}%
\end{pgfscope}%
\begin{pgfscope}%
\pgfsetbuttcap%
\pgfsetroundjoin%
\pgfsetlinewidth{1.003750pt}%
\definecolor{currentstroke}{rgb}{0.321569,0.321569,0.321569}%
\pgfsetstrokecolor{currentstroke}%
\pgfsetstrokeopacity{0.900000}%
\pgfsetdash{}{0pt}%
\pgfpathmoveto{\pgfqpoint{8.036091in}{3.990262in}}%
\pgfpathlineto{\pgfqpoint{8.095829in}{4.050000in}}%
\pgfpathlineto{\pgfqpoint{8.036091in}{4.109738in}}%
\pgfpathlineto{\pgfqpoint{7.976353in}{4.050000in}}%
\pgfpathlineto{\pgfqpoint{8.036091in}{3.990262in}}%
\pgfpathclose%
\pgfusepath{stroke}%
\end{pgfscope}%
\begin{pgfscope}%
\pgfsetbuttcap%
\pgfsetroundjoin%
\pgfsetlinewidth{1.003750pt}%
\definecolor{currentstroke}{rgb}{0.321569,0.321569,0.321569}%
\pgfsetstrokecolor{currentstroke}%
\pgfsetstrokeopacity{0.900000}%
\pgfsetdash{}{0pt}%
\pgfpathmoveto{\pgfqpoint{8.036091in}{3.990262in}}%
\pgfpathlineto{\pgfqpoint{8.095829in}{4.050000in}}%
\pgfpathlineto{\pgfqpoint{8.036091in}{4.109738in}}%
\pgfpathlineto{\pgfqpoint{7.976353in}{4.050000in}}%
\pgfpathlineto{\pgfqpoint{8.036091in}{3.990262in}}%
\pgfpathclose%
\pgfusepath{stroke}%
\end{pgfscope}%
\begin{pgfscope}%
\pgfsetbuttcap%
\pgfsetroundjoin%
\pgfsetlinewidth{1.003750pt}%
\definecolor{currentstroke}{rgb}{0.321569,0.321569,0.321569}%
\pgfsetstrokecolor{currentstroke}%
\pgfsetstrokeopacity{0.900000}%
\pgfsetdash{}{0pt}%
\pgfpathmoveto{\pgfqpoint{8.036091in}{3.990262in}}%
\pgfpathlineto{\pgfqpoint{8.095829in}{4.050000in}}%
\pgfpathlineto{\pgfqpoint{8.036091in}{4.109738in}}%
\pgfpathlineto{\pgfqpoint{7.976353in}{4.050000in}}%
\pgfpathlineto{\pgfqpoint{8.036091in}{3.990262in}}%
\pgfpathclose%
\pgfusepath{stroke}%
\end{pgfscope}%
\begin{pgfscope}%
\pgfsetbuttcap%
\pgfsetroundjoin%
\pgfsetlinewidth{1.003750pt}%
\definecolor{currentstroke}{rgb}{0.321569,0.321569,0.321569}%
\pgfsetstrokecolor{currentstroke}%
\pgfsetstrokeopacity{0.900000}%
\pgfsetdash{}{0pt}%
\pgfpathmoveto{\pgfqpoint{8.036091in}{3.990262in}}%
\pgfpathlineto{\pgfqpoint{8.095829in}{4.050000in}}%
\pgfpathlineto{\pgfqpoint{8.036091in}{4.109738in}}%
\pgfpathlineto{\pgfqpoint{7.976353in}{4.050000in}}%
\pgfpathlineto{\pgfqpoint{8.036091in}{3.990262in}}%
\pgfpathclose%
\pgfusepath{stroke}%
\end{pgfscope}%
\begin{pgfscope}%
\pgfsetbuttcap%
\pgfsetroundjoin%
\pgfsetlinewidth{1.003750pt}%
\definecolor{currentstroke}{rgb}{0.321569,0.321569,0.321569}%
\pgfsetstrokecolor{currentstroke}%
\pgfsetstrokeopacity{0.900000}%
\pgfsetdash{}{0pt}%
\pgfpathmoveto{\pgfqpoint{8.036091in}{3.990262in}}%
\pgfpathlineto{\pgfqpoint{8.095829in}{4.050000in}}%
\pgfpathlineto{\pgfqpoint{8.036091in}{4.109738in}}%
\pgfpathlineto{\pgfqpoint{7.976353in}{4.050000in}}%
\pgfpathlineto{\pgfqpoint{8.036091in}{3.990262in}}%
\pgfpathclose%
\pgfusepath{stroke}%
\end{pgfscope}%
\begin{pgfscope}%
\pgfsetbuttcap%
\pgfsetroundjoin%
\pgfsetlinewidth{1.003750pt}%
\definecolor{currentstroke}{rgb}{0.321569,0.321569,0.321569}%
\pgfsetstrokecolor{currentstroke}%
\pgfsetstrokeopacity{0.900000}%
\pgfsetdash{}{0pt}%
\pgfpathmoveto{\pgfqpoint{8.036091in}{3.990262in}}%
\pgfpathlineto{\pgfqpoint{8.095829in}{4.050000in}}%
\pgfpathlineto{\pgfqpoint{8.036091in}{4.109738in}}%
\pgfpathlineto{\pgfqpoint{7.976353in}{4.050000in}}%
\pgfpathlineto{\pgfqpoint{8.036091in}{3.990262in}}%
\pgfpathclose%
\pgfusepath{stroke}%
\end{pgfscope}%
\begin{pgfscope}%
\pgfpathrectangle{\pgfqpoint{8.036091in}{4.050000in}}{\pgfqpoint{5.400000in}{5.400000in}}%
\pgfusepath{clip}%
\pgfsetbuttcap%
\pgfsetroundjoin%
\pgfsetlinewidth{1.003750pt}%
\definecolor{currentstroke}{rgb}{0.176471,0.713725,0.643137}%
\pgfsetstrokecolor{currentstroke}%
\pgfsetstrokeopacity{0.900000}%
\pgfsetdash{}{0pt}%
\pgfpathmoveto{\pgfqpoint{12.832414in}{9.155940in}}%
\pgfpathcurveto{\pgfqpoint{12.842668in}{9.155940in}}{\pgfqpoint{12.852504in}{9.160014in}}{\pgfqpoint{12.859754in}{9.167265in}}%
\pgfpathcurveto{\pgfqpoint{12.867005in}{9.174515in}}{\pgfqpoint{12.871079in}{9.184351in}}{\pgfqpoint{12.871079in}{9.194605in}}%
\pgfpathcurveto{\pgfqpoint{12.871079in}{9.204859in}}{\pgfqpoint{12.867005in}{9.214695in}}{\pgfqpoint{12.859754in}{9.221945in}}%
\pgfpathcurveto{\pgfqpoint{12.852504in}{9.229196in}}{\pgfqpoint{12.842668in}{9.233270in}}{\pgfqpoint{12.832414in}{9.233270in}}%
\pgfpathcurveto{\pgfqpoint{12.822160in}{9.233270in}}{\pgfqpoint{12.812324in}{9.229196in}}{\pgfqpoint{12.805074in}{9.221945in}}%
\pgfpathcurveto{\pgfqpoint{12.797823in}{9.214695in}}{\pgfqpoint{12.793749in}{9.204859in}}{\pgfqpoint{12.793749in}{9.194605in}}%
\pgfpathcurveto{\pgfqpoint{12.793749in}{9.184351in}}{\pgfqpoint{12.797823in}{9.174515in}}{\pgfqpoint{12.805074in}{9.167265in}}%
\pgfpathcurveto{\pgfqpoint{12.812324in}{9.160014in}}{\pgfqpoint{12.822160in}{9.155940in}}{\pgfqpoint{12.832414in}{9.155940in}}%
\pgfpathlineto{\pgfqpoint{12.832414in}{9.155940in}}%
\pgfpathclose%
\pgfusepath{stroke}%
\end{pgfscope}%
\begin{pgfscope}%
\pgfpathrectangle{\pgfqpoint{8.036091in}{4.050000in}}{\pgfqpoint{5.400000in}{5.400000in}}%
\pgfusepath{clip}%
\pgfsetbuttcap%
\pgfsetroundjoin%
\pgfsetlinewidth{1.003750pt}%
\definecolor{currentstroke}{rgb}{0.176471,0.713725,0.643137}%
\pgfsetstrokecolor{currentstroke}%
\pgfsetstrokeopacity{0.900000}%
\pgfsetdash{}{0pt}%
\pgfpathmoveto{\pgfqpoint{12.808538in}{9.116846in}}%
\pgfpathcurveto{\pgfqpoint{12.818792in}{9.116846in}}{\pgfqpoint{12.828627in}{9.120920in}}{\pgfqpoint{12.835878in}{9.128171in}}%
\pgfpathcurveto{\pgfqpoint{12.843129in}{9.135422in}}{\pgfqpoint{12.847203in}{9.145257in}}{\pgfqpoint{12.847203in}{9.155512in}}%
\pgfpathcurveto{\pgfqpoint{12.847203in}{9.165766in}}{\pgfqpoint{12.843129in}{9.175601in}}{\pgfqpoint{12.835878in}{9.182852in}}%
\pgfpathcurveto{\pgfqpoint{12.828627in}{9.190103in}}{\pgfqpoint{12.818792in}{9.194177in}}{\pgfqpoint{12.808538in}{9.194177in}}%
\pgfpathcurveto{\pgfqpoint{12.798283in}{9.194177in}}{\pgfqpoint{12.788448in}{9.190103in}}{\pgfqpoint{12.781197in}{9.182852in}}%
\pgfpathcurveto{\pgfqpoint{12.773947in}{9.175601in}}{\pgfqpoint{12.769873in}{9.165766in}}{\pgfqpoint{12.769873in}{9.155512in}}%
\pgfpathcurveto{\pgfqpoint{12.769873in}{9.145257in}}{\pgfqpoint{12.773947in}{9.135422in}}{\pgfqpoint{12.781197in}{9.128171in}}%
\pgfpathcurveto{\pgfqpoint{12.788448in}{9.120920in}}{\pgfqpoint{12.798283in}{9.116846in}}{\pgfqpoint{12.808538in}{9.116846in}}%
\pgfpathlineto{\pgfqpoint{12.808538in}{9.116846in}}%
\pgfpathclose%
\pgfusepath{stroke}%
\end{pgfscope}%
\begin{pgfscope}%
\pgfpathrectangle{\pgfqpoint{8.036091in}{4.050000in}}{\pgfqpoint{5.400000in}{5.400000in}}%
\pgfusepath{clip}%
\pgfsetbuttcap%
\pgfsetroundjoin%
\pgfsetlinewidth{1.003750pt}%
\definecolor{currentstroke}{rgb}{0.176471,0.713725,0.643137}%
\pgfsetstrokecolor{currentstroke}%
\pgfsetstrokeopacity{0.900000}%
\pgfsetdash{}{0pt}%
\pgfpathmoveto{\pgfqpoint{12.866463in}{9.111727in}}%
\pgfpathcurveto{\pgfqpoint{12.876717in}{9.111727in}}{\pgfqpoint{12.886553in}{9.115801in}}{\pgfqpoint{12.893804in}{9.123052in}}%
\pgfpathcurveto{\pgfqpoint{12.901054in}{9.130302in}}{\pgfqpoint{12.905128in}{9.140138in}}{\pgfqpoint{12.905128in}{9.150392in}}%
\pgfpathcurveto{\pgfqpoint{12.905128in}{9.160646in}}{\pgfqpoint{12.901054in}{9.170481in}}{\pgfqpoint{12.893804in}{9.177732in}}%
\pgfpathcurveto{\pgfqpoint{12.886553in}{9.184983in}}{\pgfqpoint{12.876717in}{9.189057in}}{\pgfqpoint{12.866463in}{9.189057in}}%
\pgfpathcurveto{\pgfqpoint{12.856209in}{9.189057in}}{\pgfqpoint{12.846374in}{9.184983in}}{\pgfqpoint{12.839123in}{9.177732in}}%
\pgfpathcurveto{\pgfqpoint{12.831872in}{9.170481in}}{\pgfqpoint{12.827798in}{9.160646in}}{\pgfqpoint{12.827798in}{9.150392in}}%
\pgfpathcurveto{\pgfqpoint{12.827798in}{9.140138in}}{\pgfqpoint{12.831872in}{9.130302in}}{\pgfqpoint{12.839123in}{9.123052in}}%
\pgfpathcurveto{\pgfqpoint{12.846374in}{9.115801in}}{\pgfqpoint{12.856209in}{9.111727in}}{\pgfqpoint{12.866463in}{9.111727in}}%
\pgfpathlineto{\pgfqpoint{12.866463in}{9.111727in}}%
\pgfpathclose%
\pgfusepath{stroke}%
\end{pgfscope}%
\begin{pgfscope}%
\pgfpathrectangle{\pgfqpoint{8.036091in}{4.050000in}}{\pgfqpoint{5.400000in}{5.400000in}}%
\pgfusepath{clip}%
\pgfsetbuttcap%
\pgfsetroundjoin%
\pgfsetlinewidth{1.003750pt}%
\definecolor{currentstroke}{rgb}{0.176471,0.713725,0.643137}%
\pgfsetstrokecolor{currentstroke}%
\pgfsetstrokeopacity{0.900000}%
\pgfsetdash{}{0pt}%
\pgfpathmoveto{\pgfqpoint{13.105447in}{9.058890in}}%
\pgfpathcurveto{\pgfqpoint{13.115702in}{9.058890in}}{\pgfqpoint{13.125537in}{9.062964in}}{\pgfqpoint{13.132788in}{9.070215in}}%
\pgfpathcurveto{\pgfqpoint{13.140038in}{9.077466in}}{\pgfqpoint{13.144112in}{9.087301in}}{\pgfqpoint{13.144112in}{9.097555in}}%
\pgfpathcurveto{\pgfqpoint{13.144112in}{9.107809in}}{\pgfqpoint{13.140038in}{9.117645in}}{\pgfqpoint{13.132788in}{9.124896in}}%
\pgfpathcurveto{\pgfqpoint{13.125537in}{9.132146in}}{\pgfqpoint{13.115702in}{9.136220in}}{\pgfqpoint{13.105447in}{9.136220in}}%
\pgfpathcurveto{\pgfqpoint{13.095193in}{9.136220in}}{\pgfqpoint{13.085358in}{9.132146in}}{\pgfqpoint{13.078107in}{9.124896in}}%
\pgfpathcurveto{\pgfqpoint{13.070856in}{9.117645in}}{\pgfqpoint{13.066782in}{9.107809in}}{\pgfqpoint{13.066782in}{9.097555in}}%
\pgfpathcurveto{\pgfqpoint{13.066782in}{9.087301in}}{\pgfqpoint{13.070856in}{9.077466in}}{\pgfqpoint{13.078107in}{9.070215in}}%
\pgfpathcurveto{\pgfqpoint{13.085358in}{9.062964in}}{\pgfqpoint{13.095193in}{9.058890in}}{\pgfqpoint{13.105447in}{9.058890in}}%
\pgfpathlineto{\pgfqpoint{13.105447in}{9.058890in}}%
\pgfpathclose%
\pgfusepath{stroke}%
\end{pgfscope}%
\begin{pgfscope}%
\pgfpathrectangle{\pgfqpoint{8.036091in}{4.050000in}}{\pgfqpoint{5.400000in}{5.400000in}}%
\pgfusepath{clip}%
\pgfsetbuttcap%
\pgfsetroundjoin%
\pgfsetlinewidth{1.003750pt}%
\definecolor{currentstroke}{rgb}{0.176471,0.713725,0.643137}%
\pgfsetstrokecolor{currentstroke}%
\pgfsetstrokeopacity{0.900000}%
\pgfsetdash{}{0pt}%
\pgfpathmoveto{\pgfqpoint{12.594762in}{8.886052in}}%
\pgfpathcurveto{\pgfqpoint{12.605016in}{8.886052in}}{\pgfqpoint{12.614852in}{8.890126in}}{\pgfqpoint{12.622102in}{8.897377in}}%
\pgfpathcurveto{\pgfqpoint{12.629353in}{8.904628in}}{\pgfqpoint{12.633427in}{8.914463in}}{\pgfqpoint{12.633427in}{8.924717in}}%
\pgfpathcurveto{\pgfqpoint{12.633427in}{8.934971in}}{\pgfqpoint{12.629353in}{8.944807in}}{\pgfqpoint{12.622102in}{8.952058in}}%
\pgfpathcurveto{\pgfqpoint{12.614852in}{8.959308in}}{\pgfqpoint{12.605016in}{8.963382in}}{\pgfqpoint{12.594762in}{8.963382in}}%
\pgfpathcurveto{\pgfqpoint{12.584508in}{8.963382in}}{\pgfqpoint{12.574673in}{8.959308in}}{\pgfqpoint{12.567422in}{8.952058in}}%
\pgfpathcurveto{\pgfqpoint{12.560171in}{8.944807in}}{\pgfqpoint{12.556097in}{8.934971in}}{\pgfqpoint{12.556097in}{8.924717in}}%
\pgfpathcurveto{\pgfqpoint{12.556097in}{8.914463in}}{\pgfqpoint{12.560171in}{8.904628in}}{\pgfqpoint{12.567422in}{8.897377in}}%
\pgfpathcurveto{\pgfqpoint{12.574673in}{8.890126in}}{\pgfqpoint{12.584508in}{8.886052in}}{\pgfqpoint{12.594762in}{8.886052in}}%
\pgfpathlineto{\pgfqpoint{12.594762in}{8.886052in}}%
\pgfpathclose%
\pgfusepath{stroke}%
\end{pgfscope}%
\begin{pgfscope}%
\pgfpathrectangle{\pgfqpoint{8.036091in}{4.050000in}}{\pgfqpoint{5.400000in}{5.400000in}}%
\pgfusepath{clip}%
\pgfsetbuttcap%
\pgfsetroundjoin%
\pgfsetlinewidth{1.003750pt}%
\definecolor{currentstroke}{rgb}{0.000000,0.447059,0.698039}%
\pgfsetstrokecolor{currentstroke}%
\pgfsetstrokeopacity{0.900000}%
\pgfsetdash{}{0pt}%
\pgfpathmoveto{\pgfqpoint{10.649005in}{6.696987in}}%
\pgfpathcurveto{\pgfqpoint{10.659259in}{6.696987in}}{\pgfqpoint{10.669095in}{6.701061in}}{\pgfqpoint{10.676346in}{6.708312in}}%
\pgfpathcurveto{\pgfqpoint{10.683596in}{6.715562in}}{\pgfqpoint{10.687670in}{6.725398in}}{\pgfqpoint{10.687670in}{6.735652in}}%
\pgfpathcurveto{\pgfqpoint{10.687670in}{6.745906in}}{\pgfqpoint{10.683596in}{6.755741in}}{\pgfqpoint{10.676346in}{6.762992in}}%
\pgfpathcurveto{\pgfqpoint{10.669095in}{6.770243in}}{\pgfqpoint{10.659259in}{6.774317in}}{\pgfqpoint{10.649005in}{6.774317in}}%
\pgfpathcurveto{\pgfqpoint{10.638751in}{6.774317in}}{\pgfqpoint{10.628916in}{6.770243in}}{\pgfqpoint{10.621665in}{6.762992in}}%
\pgfpathcurveto{\pgfqpoint{10.614414in}{6.755741in}}{\pgfqpoint{10.610340in}{6.745906in}}{\pgfqpoint{10.610340in}{6.735652in}}%
\pgfpathcurveto{\pgfqpoint{10.610340in}{6.725398in}}{\pgfqpoint{10.614414in}{6.715562in}}{\pgfqpoint{10.621665in}{6.708312in}}%
\pgfpathcurveto{\pgfqpoint{10.628916in}{6.701061in}}{\pgfqpoint{10.638751in}{6.696987in}}{\pgfqpoint{10.649005in}{6.696987in}}%
\pgfpathlineto{\pgfqpoint{10.649005in}{6.696987in}}%
\pgfpathclose%
\pgfusepath{stroke}%
\end{pgfscope}%
\begin{pgfscope}%
\pgfpathrectangle{\pgfqpoint{8.036091in}{4.050000in}}{\pgfqpoint{5.400000in}{5.400000in}}%
\pgfusepath{clip}%
\pgfsetbuttcap%
\pgfsetroundjoin%
\pgfsetlinewidth{1.003750pt}%
\definecolor{currentstroke}{rgb}{0.000000,0.447059,0.698039}%
\pgfsetstrokecolor{currentstroke}%
\pgfsetstrokeopacity{0.900000}%
\pgfsetdash{}{0pt}%
\pgfpathmoveto{\pgfqpoint{10.673096in}{5.993991in}}%
\pgfpathcurveto{\pgfqpoint{10.683350in}{5.993991in}}{\pgfqpoint{10.693186in}{5.998065in}}{\pgfqpoint{10.700437in}{6.005315in}}%
\pgfpathcurveto{\pgfqpoint{10.707687in}{6.012566in}}{\pgfqpoint{10.711761in}{6.022401in}}{\pgfqpoint{10.711761in}{6.032656in}}%
\pgfpathcurveto{\pgfqpoint{10.711761in}{6.042910in}}{\pgfqpoint{10.707687in}{6.052745in}}{\pgfqpoint{10.700437in}{6.059996in}}%
\pgfpathcurveto{\pgfqpoint{10.693186in}{6.067247in}}{\pgfqpoint{10.683350in}{6.071321in}}{\pgfqpoint{10.673096in}{6.071321in}}%
\pgfpathcurveto{\pgfqpoint{10.662842in}{6.071321in}}{\pgfqpoint{10.653007in}{6.067247in}}{\pgfqpoint{10.645756in}{6.059996in}}%
\pgfpathcurveto{\pgfqpoint{10.638505in}{6.052745in}}{\pgfqpoint{10.634431in}{6.042910in}}{\pgfqpoint{10.634431in}{6.032656in}}%
\pgfpathcurveto{\pgfqpoint{10.634431in}{6.022401in}}{\pgfqpoint{10.638505in}{6.012566in}}{\pgfqpoint{10.645756in}{6.005315in}}%
\pgfpathcurveto{\pgfqpoint{10.653007in}{5.998065in}}{\pgfqpoint{10.662842in}{5.993991in}}{\pgfqpoint{10.673096in}{5.993991in}}%
\pgfpathlineto{\pgfqpoint{10.673096in}{5.993991in}}%
\pgfpathclose%
\pgfusepath{stroke}%
\end{pgfscope}%
\begin{pgfscope}%
\pgfpathrectangle{\pgfqpoint{8.036091in}{4.050000in}}{\pgfqpoint{5.400000in}{5.400000in}}%
\pgfusepath{clip}%
\pgfsetbuttcap%
\pgfsetroundjoin%
\pgfsetlinewidth{1.003750pt}%
\definecolor{currentstroke}{rgb}{0.337255,0.705882,0.913725}%
\pgfsetstrokecolor{currentstroke}%
\pgfsetstrokeopacity{0.900000}%
\pgfsetdash{}{0pt}%
\pgfpathmoveto{\pgfqpoint{9.011668in}{4.408018in}}%
\pgfpathcurveto{\pgfqpoint{9.021922in}{4.408018in}}{\pgfqpoint{9.031757in}{4.412092in}}{\pgfqpoint{9.039008in}{4.419343in}}%
\pgfpathcurveto{\pgfqpoint{9.046259in}{4.426593in}}{\pgfqpoint{9.050333in}{4.436429in}}{\pgfqpoint{9.050333in}{4.446683in}}%
\pgfpathcurveto{\pgfqpoint{9.050333in}{4.456937in}}{\pgfqpoint{9.046259in}{4.466772in}}{\pgfqpoint{9.039008in}{4.474023in}}%
\pgfpathcurveto{\pgfqpoint{9.031757in}{4.481274in}}{\pgfqpoint{9.021922in}{4.485348in}}{\pgfqpoint{9.011668in}{4.485348in}}%
\pgfpathcurveto{\pgfqpoint{9.001414in}{4.485348in}}{\pgfqpoint{8.991578in}{4.481274in}}{\pgfqpoint{8.984327in}{4.474023in}}%
\pgfpathcurveto{\pgfqpoint{8.977077in}{4.466772in}}{\pgfqpoint{8.973003in}{4.456937in}}{\pgfqpoint{8.973003in}{4.446683in}}%
\pgfpathcurveto{\pgfqpoint{8.973003in}{4.436429in}}{\pgfqpoint{8.977077in}{4.426593in}}{\pgfqpoint{8.984327in}{4.419343in}}%
\pgfpathcurveto{\pgfqpoint{8.991578in}{4.412092in}}{\pgfqpoint{9.001414in}{4.408018in}}{\pgfqpoint{9.011668in}{4.408018in}}%
\pgfpathlineto{\pgfqpoint{9.011668in}{4.408018in}}%
\pgfpathclose%
\pgfusepath{stroke}%
\end{pgfscope}%
\begin{pgfscope}%
\pgfpathrectangle{\pgfqpoint{8.036091in}{4.050000in}}{\pgfqpoint{5.400000in}{5.400000in}}%
\pgfusepath{clip}%
\pgfsetbuttcap%
\pgfsetroundjoin%
\pgfsetlinewidth{1.003750pt}%
\definecolor{currentstroke}{rgb}{0.337255,0.705882,0.913725}%
\pgfsetstrokecolor{currentstroke}%
\pgfsetstrokeopacity{0.900000}%
\pgfsetdash{}{0pt}%
\pgfpathmoveto{\pgfqpoint{9.726806in}{5.096390in}}%
\pgfpathcurveto{\pgfqpoint{9.737060in}{5.096390in}}{\pgfqpoint{9.746895in}{5.100464in}}{\pgfqpoint{9.754146in}{5.107714in}}%
\pgfpathcurveto{\pgfqpoint{9.761397in}{5.114965in}}{\pgfqpoint{9.765471in}{5.124800in}}{\pgfqpoint{9.765471in}{5.135055in}}%
\pgfpathcurveto{\pgfqpoint{9.765471in}{5.145309in}}{\pgfqpoint{9.761397in}{5.155144in}}{\pgfqpoint{9.754146in}{5.162395in}}%
\pgfpathcurveto{\pgfqpoint{9.746895in}{5.169646in}}{\pgfqpoint{9.737060in}{5.173720in}}{\pgfqpoint{9.726806in}{5.173720in}}%
\pgfpathcurveto{\pgfqpoint{9.716552in}{5.173720in}}{\pgfqpoint{9.706716in}{5.169646in}}{\pgfqpoint{9.699466in}{5.162395in}}%
\pgfpathcurveto{\pgfqpoint{9.692215in}{5.155144in}}{\pgfqpoint{9.688141in}{5.145309in}}{\pgfqpoint{9.688141in}{5.135055in}}%
\pgfpathcurveto{\pgfqpoint{9.688141in}{5.124800in}}{\pgfqpoint{9.692215in}{5.114965in}}{\pgfqpoint{9.699466in}{5.107714in}}%
\pgfpathcurveto{\pgfqpoint{9.706716in}{5.100464in}}{\pgfqpoint{9.716552in}{5.096390in}}{\pgfqpoint{9.726806in}{5.096390in}}%
\pgfpathlineto{\pgfqpoint{9.726806in}{5.096390in}}%
\pgfpathclose%
\pgfusepath{stroke}%
\end{pgfscope}%
\begin{pgfscope}%
\pgfsetbuttcap%
\pgfsetroundjoin%
\pgfsetlinewidth{1.003750pt}%
\definecolor{currentstroke}{rgb}{0.337255,0.705882,0.913725}%
\pgfsetstrokecolor{currentstroke}%
\pgfsetstrokeopacity{0.900000}%
\pgfsetdash{}{0pt}%
\pgfpathmoveto{\pgfqpoint{13.432742in}{9.390262in}}%
\pgfpathlineto{\pgfqpoint{13.492480in}{9.450000in}}%
\pgfpathlineto{\pgfqpoint{13.432742in}{9.509738in}}%
\pgfpathlineto{\pgfqpoint{13.373004in}{9.450000in}}%
\pgfpathlineto{\pgfqpoint{13.432742in}{9.390262in}}%
\pgfpathclose%
\pgfusepath{stroke}%
\end{pgfscope}%
\begin{pgfscope}%
\pgfsetbuttcap%
\pgfsetroundjoin%
\pgfsetlinewidth{1.003750pt}%
\definecolor{currentstroke}{rgb}{0.337255,0.705882,0.913725}%
\pgfsetstrokecolor{currentstroke}%
\pgfsetstrokeopacity{0.900000}%
\pgfsetdash{}{0pt}%
\pgfpathmoveto{\pgfqpoint{8.036091in}{3.990262in}}%
\pgfpathlineto{\pgfqpoint{8.095829in}{4.050000in}}%
\pgfpathlineto{\pgfqpoint{8.036091in}{4.109738in}}%
\pgfpathlineto{\pgfqpoint{7.976353in}{4.050000in}}%
\pgfpathlineto{\pgfqpoint{8.036091in}{3.990262in}}%
\pgfpathclose%
\pgfusepath{stroke}%
\end{pgfscope}%
\begin{pgfscope}%
\pgfsetbuttcap%
\pgfsetroundjoin%
\pgfsetlinewidth{1.003750pt}%
\definecolor{currentstroke}{rgb}{0.901961,0.196078,0.156863}%
\pgfsetstrokecolor{currentstroke}%
\pgfsetstrokeopacity{0.900000}%
\pgfsetdash{}{0pt}%
\pgfpathmoveto{\pgfqpoint{13.436091in}{9.390262in}}%
\pgfpathlineto{\pgfqpoint{13.495829in}{9.450000in}}%
\pgfpathlineto{\pgfqpoint{13.436091in}{9.509738in}}%
\pgfpathlineto{\pgfqpoint{13.376353in}{9.450000in}}%
\pgfpathlineto{\pgfqpoint{13.436091in}{9.390262in}}%
\pgfpathclose%
\pgfusepath{stroke}%
\end{pgfscope}%
\begin{pgfscope}%
\pgfpathrectangle{\pgfqpoint{8.036091in}{4.050000in}}{\pgfqpoint{5.400000in}{5.400000in}}%
\pgfusepath{clip}%
\pgfsetbuttcap%
\pgfsetroundjoin%
\pgfsetlinewidth{1.003750pt}%
\definecolor{currentstroke}{rgb}{0.662745,0.368627,0.000000}%
\pgfsetstrokecolor{currentstroke}%
\pgfsetstrokeopacity{0.900000}%
\pgfsetdash{}{0pt}%
\pgfpathmoveto{\pgfqpoint{12.885083in}{9.038197in}}%
\pgfpathcurveto{\pgfqpoint{12.895337in}{9.038197in}}{\pgfqpoint{12.905173in}{9.042271in}}{\pgfqpoint{12.912424in}{9.049521in}}%
\pgfpathcurveto{\pgfqpoint{12.919674in}{9.056772in}}{\pgfqpoint{12.923748in}{9.066608in}}{\pgfqpoint{12.923748in}{9.076862in}}%
\pgfpathcurveto{\pgfqpoint{12.923748in}{9.087116in}}{\pgfqpoint{12.919674in}{9.096951in}}{\pgfqpoint{12.912424in}{9.104202in}}%
\pgfpathcurveto{\pgfqpoint{12.905173in}{9.111453in}}{\pgfqpoint{12.895337in}{9.115527in}}{\pgfqpoint{12.885083in}{9.115527in}}%
\pgfpathcurveto{\pgfqpoint{12.874829in}{9.115527in}}{\pgfqpoint{12.864994in}{9.111453in}}{\pgfqpoint{12.857743in}{9.104202in}}%
\pgfpathcurveto{\pgfqpoint{12.850492in}{9.096951in}}{\pgfqpoint{12.846418in}{9.087116in}}{\pgfqpoint{12.846418in}{9.076862in}}%
\pgfpathcurveto{\pgfqpoint{12.846418in}{9.066608in}}{\pgfqpoint{12.850492in}{9.056772in}}{\pgfqpoint{12.857743in}{9.049521in}}%
\pgfpathcurveto{\pgfqpoint{12.864994in}{9.042271in}}{\pgfqpoint{12.874829in}{9.038197in}}{\pgfqpoint{12.885083in}{9.038197in}}%
\pgfpathlineto{\pgfqpoint{12.885083in}{9.038197in}}%
\pgfpathclose%
\pgfusepath{stroke}%
\end{pgfscope}%
\begin{pgfscope}%
\pgfpathrectangle{\pgfqpoint{8.036091in}{4.050000in}}{\pgfqpoint{5.400000in}{5.400000in}}%
\pgfusepath{clip}%
\pgfsetbuttcap%
\pgfsetroundjoin%
\pgfsetlinewidth{1.003750pt}%
\definecolor{currentstroke}{rgb}{0.662745,0.368627,0.000000}%
\pgfsetstrokecolor{currentstroke}%
\pgfsetstrokeopacity{0.900000}%
\pgfsetdash{}{0pt}%
\pgfpathmoveto{\pgfqpoint{12.505713in}{8.658033in}}%
\pgfpathcurveto{\pgfqpoint{12.515967in}{8.658033in}}{\pgfqpoint{12.525803in}{8.662107in}}{\pgfqpoint{12.533054in}{8.669358in}}%
\pgfpathcurveto{\pgfqpoint{12.540304in}{8.676609in}}{\pgfqpoint{12.544378in}{8.686444in}}{\pgfqpoint{12.544378in}{8.696698in}}%
\pgfpathcurveto{\pgfqpoint{12.544378in}{8.706952in}}{\pgfqpoint{12.540304in}{8.716788in}}{\pgfqpoint{12.533054in}{8.724039in}}%
\pgfpathcurveto{\pgfqpoint{12.525803in}{8.731289in}}{\pgfqpoint{12.515967in}{8.735363in}}{\pgfqpoint{12.505713in}{8.735363in}}%
\pgfpathcurveto{\pgfqpoint{12.495459in}{8.735363in}}{\pgfqpoint{12.485624in}{8.731289in}}{\pgfqpoint{12.478373in}{8.724039in}}%
\pgfpathcurveto{\pgfqpoint{12.471122in}{8.716788in}}{\pgfqpoint{12.467048in}{8.706952in}}{\pgfqpoint{12.467048in}{8.696698in}}%
\pgfpathcurveto{\pgfqpoint{12.467048in}{8.686444in}}{\pgfqpoint{12.471122in}{8.676609in}}{\pgfqpoint{12.478373in}{8.669358in}}%
\pgfpathcurveto{\pgfqpoint{12.485624in}{8.662107in}}{\pgfqpoint{12.495459in}{8.658033in}}{\pgfqpoint{12.505713in}{8.658033in}}%
\pgfpathlineto{\pgfqpoint{12.505713in}{8.658033in}}%
\pgfpathclose%
\pgfusepath{stroke}%
\end{pgfscope}%
\begin{pgfscope}%
\pgfsetbuttcap%
\pgfsetroundjoin%
\pgfsetlinewidth{1.003750pt}%
\definecolor{currentstroke}{rgb}{0.662745,0.368627,0.000000}%
\pgfsetstrokecolor{currentstroke}%
\pgfsetstrokeopacity{0.900000}%
\pgfsetdash{}{0pt}%
\pgfpathmoveto{\pgfqpoint{13.436091in}{9.390262in}}%
\pgfpathlineto{\pgfqpoint{13.495829in}{9.450000in}}%
\pgfpathlineto{\pgfqpoint{13.436091in}{9.509738in}}%
\pgfpathlineto{\pgfqpoint{13.376353in}{9.450000in}}%
\pgfpathlineto{\pgfqpoint{13.436091in}{9.390262in}}%
\pgfpathclose%
\pgfusepath{stroke}%
\end{pgfscope}%
\begin{pgfscope}%
\pgfpathrectangle{\pgfqpoint{8.036091in}{4.050000in}}{\pgfqpoint{5.400000in}{5.400000in}}%
\pgfusepath{clip}%
\pgfsetbuttcap%
\pgfsetroundjoin%
\pgfsetlinewidth{1.003750pt}%
\definecolor{currentstroke}{rgb}{0.901961,0.403922,0.619608}%
\pgfsetstrokecolor{currentstroke}%
\pgfsetstrokeopacity{0.900000}%
\pgfsetdash{}{0pt}%
\pgfpathmoveto{\pgfqpoint{10.206463in}{5.803782in}}%
\pgfpathcurveto{\pgfqpoint{10.216717in}{5.803782in}}{\pgfqpoint{10.226552in}{5.807856in}}{\pgfqpoint{10.233803in}{5.815106in}}%
\pgfpathcurveto{\pgfqpoint{10.241054in}{5.822357in}}{\pgfqpoint{10.245128in}{5.832192in}}{\pgfqpoint{10.245128in}{5.842447in}}%
\pgfpathcurveto{\pgfqpoint{10.245128in}{5.852701in}}{\pgfqpoint{10.241054in}{5.862536in}}{\pgfqpoint{10.233803in}{5.869787in}}%
\pgfpathcurveto{\pgfqpoint{10.226552in}{5.877038in}}{\pgfqpoint{10.216717in}{5.881112in}}{\pgfqpoint{10.206463in}{5.881112in}}%
\pgfpathcurveto{\pgfqpoint{10.196209in}{5.881112in}}{\pgfqpoint{10.186373in}{5.877038in}}{\pgfqpoint{10.179122in}{5.869787in}}%
\pgfpathcurveto{\pgfqpoint{10.171872in}{5.862536in}}{\pgfqpoint{10.167798in}{5.852701in}}{\pgfqpoint{10.167798in}{5.842447in}}%
\pgfpathcurveto{\pgfqpoint{10.167798in}{5.832192in}}{\pgfqpoint{10.171872in}{5.822357in}}{\pgfqpoint{10.179122in}{5.815106in}}%
\pgfpathcurveto{\pgfqpoint{10.186373in}{5.807856in}}{\pgfqpoint{10.196209in}{5.803782in}}{\pgfqpoint{10.206463in}{5.803782in}}%
\pgfpathlineto{\pgfqpoint{10.206463in}{5.803782in}}%
\pgfpathclose%
\pgfusepath{stroke}%
\end{pgfscope}%
\begin{pgfscope}%
\pgfpathrectangle{\pgfqpoint{8.036091in}{4.050000in}}{\pgfqpoint{5.400000in}{5.400000in}}%
\pgfusepath{clip}%
\pgfsetbuttcap%
\pgfsetroundjoin%
\pgfsetlinewidth{1.003750pt}%
\definecolor{currentstroke}{rgb}{0.901961,0.403922,0.619608}%
\pgfsetstrokecolor{currentstroke}%
\pgfsetstrokeopacity{0.900000}%
\pgfsetdash{}{0pt}%
\pgfpathmoveto{\pgfqpoint{11.858861in}{7.942463in}}%
\pgfpathcurveto{\pgfqpoint{11.869115in}{7.942463in}}{\pgfqpoint{11.878951in}{7.946537in}}{\pgfqpoint{11.886202in}{7.953788in}}%
\pgfpathcurveto{\pgfqpoint{11.893452in}{7.961039in}}{\pgfqpoint{11.897526in}{7.970874in}}{\pgfqpoint{11.897526in}{7.981128in}}%
\pgfpathcurveto{\pgfqpoint{11.897526in}{7.991382in}}{\pgfqpoint{11.893452in}{8.001218in}}{\pgfqpoint{11.886202in}{8.008469in}}%
\pgfpathcurveto{\pgfqpoint{11.878951in}{8.015719in}}{\pgfqpoint{11.869115in}{8.019793in}}{\pgfqpoint{11.858861in}{8.019793in}}%
\pgfpathcurveto{\pgfqpoint{11.848607in}{8.019793in}}{\pgfqpoint{11.838772in}{8.015719in}}{\pgfqpoint{11.831521in}{8.008469in}}%
\pgfpathcurveto{\pgfqpoint{11.824270in}{8.001218in}}{\pgfqpoint{11.820196in}{7.991382in}}{\pgfqpoint{11.820196in}{7.981128in}}%
\pgfpathcurveto{\pgfqpoint{11.820196in}{7.970874in}}{\pgfqpoint{11.824270in}{7.961039in}}{\pgfqpoint{11.831521in}{7.953788in}}%
\pgfpathcurveto{\pgfqpoint{11.838772in}{7.946537in}}{\pgfqpoint{11.848607in}{7.942463in}}{\pgfqpoint{11.858861in}{7.942463in}}%
\pgfpathlineto{\pgfqpoint{11.858861in}{7.942463in}}%
\pgfpathclose%
\pgfusepath{stroke}%
\end{pgfscope}%
\begin{pgfscope}%
\pgfpathrectangle{\pgfqpoint{8.036091in}{4.050000in}}{\pgfqpoint{5.400000in}{5.400000in}}%
\pgfusepath{clip}%
\pgfsetbuttcap%
\pgfsetroundjoin%
\pgfsetlinewidth{1.003750pt}%
\definecolor{currentstroke}{rgb}{0.901961,0.403922,0.619608}%
\pgfsetstrokecolor{currentstroke}%
\pgfsetstrokeopacity{0.900000}%
\pgfsetdash{}{0pt}%
\pgfpathmoveto{\pgfqpoint{11.817260in}{6.904203in}}%
\pgfpathcurveto{\pgfqpoint{11.827514in}{6.904203in}}{\pgfqpoint{11.837350in}{6.908277in}}{\pgfqpoint{11.844601in}{6.915527in}}%
\pgfpathcurveto{\pgfqpoint{11.851851in}{6.922778in}}{\pgfqpoint{11.855925in}{6.932614in}}{\pgfqpoint{11.855925in}{6.942868in}}%
\pgfpathcurveto{\pgfqpoint{11.855925in}{6.953122in}}{\pgfqpoint{11.851851in}{6.962957in}}{\pgfqpoint{11.844601in}{6.970208in}}%
\pgfpathcurveto{\pgfqpoint{11.837350in}{6.977459in}}{\pgfqpoint{11.827514in}{6.981533in}}{\pgfqpoint{11.817260in}{6.981533in}}%
\pgfpathcurveto{\pgfqpoint{11.807006in}{6.981533in}}{\pgfqpoint{11.797171in}{6.977459in}}{\pgfqpoint{11.789920in}{6.970208in}}%
\pgfpathcurveto{\pgfqpoint{11.782669in}{6.962957in}}{\pgfqpoint{11.778595in}{6.953122in}}{\pgfqpoint{11.778595in}{6.942868in}}%
\pgfpathcurveto{\pgfqpoint{11.778595in}{6.932614in}}{\pgfqpoint{11.782669in}{6.922778in}}{\pgfqpoint{11.789920in}{6.915527in}}%
\pgfpathcurveto{\pgfqpoint{11.797171in}{6.908277in}}{\pgfqpoint{11.807006in}{6.904203in}}{\pgfqpoint{11.817260in}{6.904203in}}%
\pgfpathlineto{\pgfqpoint{11.817260in}{6.904203in}}%
\pgfpathclose%
\pgfusepath{stroke}%
\end{pgfscope}%
\begin{pgfscope}%
\pgfpathrectangle{\pgfqpoint{8.036091in}{4.050000in}}{\pgfqpoint{5.400000in}{5.400000in}}%
\pgfusepath{clip}%
\pgfsetbuttcap%
\pgfsetroundjoin%
\pgfsetlinewidth{1.003750pt}%
\definecolor{currentstroke}{rgb}{0.941176,0.584314,0.337255}%
\pgfsetstrokecolor{currentstroke}%
\pgfsetstrokeopacity{0.900000}%
\pgfsetdash{}{0pt}%
\pgfpathmoveto{\pgfqpoint{13.117467in}{9.192593in}}%
\pgfpathcurveto{\pgfqpoint{13.127721in}{9.192593in}}{\pgfqpoint{13.137557in}{9.196667in}}{\pgfqpoint{13.144808in}{9.203918in}}%
\pgfpathcurveto{\pgfqpoint{13.152058in}{9.211169in}}{\pgfqpoint{13.156132in}{9.221004in}}{\pgfqpoint{13.156132in}{9.231259in}}%
\pgfpathcurveto{\pgfqpoint{13.156132in}{9.241513in}}{\pgfqpoint{13.152058in}{9.251348in}}{\pgfqpoint{13.144808in}{9.258599in}}%
\pgfpathcurveto{\pgfqpoint{13.137557in}{9.265850in}}{\pgfqpoint{13.127721in}{9.269924in}}{\pgfqpoint{13.117467in}{9.269924in}}%
\pgfpathcurveto{\pgfqpoint{13.107213in}{9.269924in}}{\pgfqpoint{13.097378in}{9.265850in}}{\pgfqpoint{13.090127in}{9.258599in}}%
\pgfpathcurveto{\pgfqpoint{13.082876in}{9.251348in}}{\pgfqpoint{13.078802in}{9.241513in}}{\pgfqpoint{13.078802in}{9.231259in}}%
\pgfpathcurveto{\pgfqpoint{13.078802in}{9.221004in}}{\pgfqpoint{13.082876in}{9.211169in}}{\pgfqpoint{13.090127in}{9.203918in}}%
\pgfpathcurveto{\pgfqpoint{13.097378in}{9.196667in}}{\pgfqpoint{13.107213in}{9.192593in}}{\pgfqpoint{13.117467in}{9.192593in}}%
\pgfpathlineto{\pgfqpoint{13.117467in}{9.192593in}}%
\pgfpathclose%
\pgfusepath{stroke}%
\end{pgfscope}%
\begin{pgfscope}%
\pgfpathrectangle{\pgfqpoint{8.036091in}{4.050000in}}{\pgfqpoint{5.400000in}{5.400000in}}%
\pgfusepath{clip}%
\pgfsetbuttcap%
\pgfsetroundjoin%
\pgfsetlinewidth{1.003750pt}%
\definecolor{currentstroke}{rgb}{0.835294,0.368627,0.000000}%
\pgfsetstrokecolor{currentstroke}%
\pgfsetstrokeopacity{0.900000}%
\pgfsetdash{}{0pt}%
\pgfpathmoveto{\pgfqpoint{12.212606in}{8.097912in}}%
\pgfpathcurveto{\pgfqpoint{12.222860in}{8.097912in}}{\pgfqpoint{12.232695in}{8.101986in}}{\pgfqpoint{12.239946in}{8.109237in}}%
\pgfpathcurveto{\pgfqpoint{12.247197in}{8.116488in}}{\pgfqpoint{12.251271in}{8.126323in}}{\pgfqpoint{12.251271in}{8.136577in}}%
\pgfpathcurveto{\pgfqpoint{12.251271in}{8.146831in}}{\pgfqpoint{12.247197in}{8.156667in}}{\pgfqpoint{12.239946in}{8.163917in}}%
\pgfpathcurveto{\pgfqpoint{12.232695in}{8.171168in}}{\pgfqpoint{12.222860in}{8.175242in}}{\pgfqpoint{12.212606in}{8.175242in}}%
\pgfpathcurveto{\pgfqpoint{12.202352in}{8.175242in}}{\pgfqpoint{12.192516in}{8.171168in}}{\pgfqpoint{12.185265in}{8.163917in}}%
\pgfpathcurveto{\pgfqpoint{12.178015in}{8.156667in}}{\pgfqpoint{12.173941in}{8.146831in}}{\pgfqpoint{12.173941in}{8.136577in}}%
\pgfpathcurveto{\pgfqpoint{12.173941in}{8.126323in}}{\pgfqpoint{12.178015in}{8.116488in}}{\pgfqpoint{12.185265in}{8.109237in}}%
\pgfpathcurveto{\pgfqpoint{12.192516in}{8.101986in}}{\pgfqpoint{12.202352in}{8.097912in}}{\pgfqpoint{12.212606in}{8.097912in}}%
\pgfpathlineto{\pgfqpoint{12.212606in}{8.097912in}}%
\pgfpathclose%
\pgfusepath{stroke}%
\end{pgfscope}%
\begin{pgfscope}%
\pgfpathrectangle{\pgfqpoint{8.036091in}{4.050000in}}{\pgfqpoint{5.400000in}{5.400000in}}%
\pgfusepath{clip}%
\pgfsetbuttcap%
\pgfsetroundjoin%
\pgfsetlinewidth{1.003750pt}%
\definecolor{currentstroke}{rgb}{0.835294,0.368627,0.000000}%
\pgfsetstrokecolor{currentstroke}%
\pgfsetstrokeopacity{0.900000}%
\pgfsetdash{}{0pt}%
\pgfpathmoveto{\pgfqpoint{11.905540in}{7.967727in}}%
\pgfpathcurveto{\pgfqpoint{11.915794in}{7.967727in}}{\pgfqpoint{11.925630in}{7.971801in}}{\pgfqpoint{11.932880in}{7.979052in}}%
\pgfpathcurveto{\pgfqpoint{11.940131in}{7.986303in}}{\pgfqpoint{11.944205in}{7.996138in}}{\pgfqpoint{11.944205in}{8.006392in}}%
\pgfpathcurveto{\pgfqpoint{11.944205in}{8.016646in}}{\pgfqpoint{11.940131in}{8.026482in}}{\pgfqpoint{11.932880in}{8.033733in}}%
\pgfpathcurveto{\pgfqpoint{11.925630in}{8.040983in}}{\pgfqpoint{11.915794in}{8.045057in}}{\pgfqpoint{11.905540in}{8.045057in}}%
\pgfpathcurveto{\pgfqpoint{11.895286in}{8.045057in}}{\pgfqpoint{11.885450in}{8.040983in}}{\pgfqpoint{11.878200in}{8.033733in}}%
\pgfpathcurveto{\pgfqpoint{11.870949in}{8.026482in}}{\pgfqpoint{11.866875in}{8.016646in}}{\pgfqpoint{11.866875in}{8.006392in}}%
\pgfpathcurveto{\pgfqpoint{11.866875in}{7.996138in}}{\pgfqpoint{11.870949in}{7.986303in}}{\pgfqpoint{11.878200in}{7.979052in}}%
\pgfpathcurveto{\pgfqpoint{11.885450in}{7.971801in}}{\pgfqpoint{11.895286in}{7.967727in}}{\pgfqpoint{11.905540in}{7.967727in}}%
\pgfpathlineto{\pgfqpoint{11.905540in}{7.967727in}}%
\pgfpathclose%
\pgfusepath{stroke}%
\end{pgfscope}%
\begin{pgfscope}%
\pgfpathrectangle{\pgfqpoint{8.036091in}{4.050000in}}{\pgfqpoint{5.400000in}{5.400000in}}%
\pgfusepath{clip}%
\pgfsetbuttcap%
\pgfsetroundjoin%
\pgfsetlinewidth{1.003750pt}%
\definecolor{currentstroke}{rgb}{0.149020,0.274510,0.325490}%
\pgfsetstrokecolor{currentstroke}%
\pgfsetstrokeopacity{0.900000}%
\pgfsetdash{}{0pt}%
\pgfpathmoveto{\pgfqpoint{11.493326in}{7.525468in}}%
\pgfpathcurveto{\pgfqpoint{11.503580in}{7.525468in}}{\pgfqpoint{11.513415in}{7.529542in}}{\pgfqpoint{11.520666in}{7.536793in}}%
\pgfpathcurveto{\pgfqpoint{11.527917in}{7.544044in}}{\pgfqpoint{11.531991in}{7.553879in}}{\pgfqpoint{11.531991in}{7.564133in}}%
\pgfpathcurveto{\pgfqpoint{11.531991in}{7.574387in}}{\pgfqpoint{11.527917in}{7.584223in}}{\pgfqpoint{11.520666in}{7.591474in}}%
\pgfpathcurveto{\pgfqpoint{11.513415in}{7.598724in}}{\pgfqpoint{11.503580in}{7.602798in}}{\pgfqpoint{11.493326in}{7.602798in}}%
\pgfpathcurveto{\pgfqpoint{11.483072in}{7.602798in}}{\pgfqpoint{11.473236in}{7.598724in}}{\pgfqpoint{11.465985in}{7.591474in}}%
\pgfpathcurveto{\pgfqpoint{11.458735in}{7.584223in}}{\pgfqpoint{11.454661in}{7.574387in}}{\pgfqpoint{11.454661in}{7.564133in}}%
\pgfpathcurveto{\pgfqpoint{11.454661in}{7.553879in}}{\pgfqpoint{11.458735in}{7.544044in}}{\pgfqpoint{11.465985in}{7.536793in}}%
\pgfpathcurveto{\pgfqpoint{11.473236in}{7.529542in}}{\pgfqpoint{11.483072in}{7.525468in}}{\pgfqpoint{11.493326in}{7.525468in}}%
\pgfpathlineto{\pgfqpoint{11.493326in}{7.525468in}}%
\pgfpathclose%
\pgfusepath{stroke}%
\end{pgfscope}%
\begin{pgfscope}%
\pgfpathrectangle{\pgfqpoint{8.036091in}{4.050000in}}{\pgfqpoint{5.400000in}{5.400000in}}%
\pgfusepath{clip}%
\pgfsetbuttcap%
\pgfsetroundjoin%
\pgfsetlinewidth{1.003750pt}%
\definecolor{currentstroke}{rgb}{0.149020,0.274510,0.325490}%
\pgfsetstrokecolor{currentstroke}%
\pgfsetstrokeopacity{0.900000}%
\pgfsetdash{}{0pt}%
\pgfpathmoveto{\pgfqpoint{11.734417in}{7.818481in}}%
\pgfpathcurveto{\pgfqpoint{11.744671in}{7.818481in}}{\pgfqpoint{11.754507in}{7.822555in}}{\pgfqpoint{11.761757in}{7.829806in}}%
\pgfpathcurveto{\pgfqpoint{11.769008in}{7.837057in}}{\pgfqpoint{11.773082in}{7.846892in}}{\pgfqpoint{11.773082in}{7.857146in}}%
\pgfpathcurveto{\pgfqpoint{11.773082in}{7.867400in}}{\pgfqpoint{11.769008in}{7.877236in}}{\pgfqpoint{11.761757in}{7.884487in}}%
\pgfpathcurveto{\pgfqpoint{11.754507in}{7.891737in}}{\pgfqpoint{11.744671in}{7.895811in}}{\pgfqpoint{11.734417in}{7.895811in}}%
\pgfpathcurveto{\pgfqpoint{11.724163in}{7.895811in}}{\pgfqpoint{11.714327in}{7.891737in}}{\pgfqpoint{11.707077in}{7.884487in}}%
\pgfpathcurveto{\pgfqpoint{11.699826in}{7.877236in}}{\pgfqpoint{11.695752in}{7.867400in}}{\pgfqpoint{11.695752in}{7.857146in}}%
\pgfpathcurveto{\pgfqpoint{11.695752in}{7.846892in}}{\pgfqpoint{11.699826in}{7.837057in}}{\pgfqpoint{11.707077in}{7.829806in}}%
\pgfpathcurveto{\pgfqpoint{11.714327in}{7.822555in}}{\pgfqpoint{11.724163in}{7.818481in}}{\pgfqpoint{11.734417in}{7.818481in}}%
\pgfpathlineto{\pgfqpoint{11.734417in}{7.818481in}}%
\pgfpathclose%
\pgfusepath{stroke}%
\end{pgfscope}%
\begin{pgfscope}%
\pgfpathrectangle{\pgfqpoint{8.036091in}{4.050000in}}{\pgfqpoint{5.400000in}{5.400000in}}%
\pgfusepath{clip}%
\pgfsetbuttcap%
\pgfsetroundjoin%
\pgfsetlinewidth{1.003750pt}%
\definecolor{currentstroke}{rgb}{0.149020,0.274510,0.325490}%
\pgfsetstrokecolor{currentstroke}%
\pgfsetstrokeopacity{0.900000}%
\pgfsetdash{}{0pt}%
\pgfpathmoveto{\pgfqpoint{11.932084in}{8.063520in}}%
\pgfpathcurveto{\pgfqpoint{11.942338in}{8.063520in}}{\pgfqpoint{11.952174in}{8.067594in}}{\pgfqpoint{11.959424in}{8.074845in}}%
\pgfpathcurveto{\pgfqpoint{11.966675in}{8.082096in}}{\pgfqpoint{11.970749in}{8.091931in}}{\pgfqpoint{11.970749in}{8.102185in}}%
\pgfpathcurveto{\pgfqpoint{11.970749in}{8.112439in}}{\pgfqpoint{11.966675in}{8.122275in}}{\pgfqpoint{11.959424in}{8.129525in}}%
\pgfpathcurveto{\pgfqpoint{11.952174in}{8.136776in}}{\pgfqpoint{11.942338in}{8.140850in}}{\pgfqpoint{11.932084in}{8.140850in}}%
\pgfpathcurveto{\pgfqpoint{11.921830in}{8.140850in}}{\pgfqpoint{11.911994in}{8.136776in}}{\pgfqpoint{11.904744in}{8.129525in}}%
\pgfpathcurveto{\pgfqpoint{11.897493in}{8.122275in}}{\pgfqpoint{11.893419in}{8.112439in}}{\pgfqpoint{11.893419in}{8.102185in}}%
\pgfpathcurveto{\pgfqpoint{11.893419in}{8.091931in}}{\pgfqpoint{11.897493in}{8.082096in}}{\pgfqpoint{11.904744in}{8.074845in}}%
\pgfpathcurveto{\pgfqpoint{11.911994in}{8.067594in}}{\pgfqpoint{11.921830in}{8.063520in}}{\pgfqpoint{11.932084in}{8.063520in}}%
\pgfpathlineto{\pgfqpoint{11.932084in}{8.063520in}}%
\pgfpathclose%
\pgfusepath{stroke}%
\end{pgfscope}%
\begin{pgfscope}%
\pgfpathrectangle{\pgfqpoint{8.036091in}{4.050000in}}{\pgfqpoint{5.400000in}{5.400000in}}%
\pgfusepath{clip}%
\pgfsetbuttcap%
\pgfsetroundjoin%
\pgfsetlinewidth{1.003750pt}%
\definecolor{currentstroke}{rgb}{0.000000,0.619608,0.450980}%
\pgfsetstrokecolor{currentstroke}%
\pgfsetstrokeopacity{0.900000}%
\pgfsetdash{}{0pt}%
\pgfpathmoveto{\pgfqpoint{12.279605in}{8.212399in}}%
\pgfpathcurveto{\pgfqpoint{12.289859in}{8.212399in}}{\pgfqpoint{12.299695in}{8.216473in}}{\pgfqpoint{12.306945in}{8.223724in}}%
\pgfpathcurveto{\pgfqpoint{12.314196in}{8.230974in}}{\pgfqpoint{12.318270in}{8.240810in}}{\pgfqpoint{12.318270in}{8.251064in}}%
\pgfpathcurveto{\pgfqpoint{12.318270in}{8.261318in}}{\pgfqpoint{12.314196in}{8.271153in}}{\pgfqpoint{12.306945in}{8.278404in}}%
\pgfpathcurveto{\pgfqpoint{12.299695in}{8.285655in}}{\pgfqpoint{12.289859in}{8.289729in}}{\pgfqpoint{12.279605in}{8.289729in}}%
\pgfpathcurveto{\pgfqpoint{12.269351in}{8.289729in}}{\pgfqpoint{12.259515in}{8.285655in}}{\pgfqpoint{12.252265in}{8.278404in}}%
\pgfpathcurveto{\pgfqpoint{12.245014in}{8.271153in}}{\pgfqpoint{12.240940in}{8.261318in}}{\pgfqpoint{12.240940in}{8.251064in}}%
\pgfpathcurveto{\pgfqpoint{12.240940in}{8.240810in}}{\pgfqpoint{12.245014in}{8.230974in}}{\pgfqpoint{12.252265in}{8.223724in}}%
\pgfpathcurveto{\pgfqpoint{12.259515in}{8.216473in}}{\pgfqpoint{12.269351in}{8.212399in}}{\pgfqpoint{12.279605in}{8.212399in}}%
\pgfpathlineto{\pgfqpoint{12.279605in}{8.212399in}}%
\pgfpathclose%
\pgfusepath{stroke}%
\end{pgfscope}%
\begin{pgfscope}%
\pgfpathrectangle{\pgfqpoint{8.036091in}{4.050000in}}{\pgfqpoint{5.400000in}{5.400000in}}%
\pgfusepath{clip}%
\pgfsetbuttcap%
\pgfsetroundjoin%
\pgfsetlinewidth{1.003750pt}%
\definecolor{currentstroke}{rgb}{0.000000,0.619608,0.450980}%
\pgfsetstrokecolor{currentstroke}%
\pgfsetstrokeopacity{0.900000}%
\pgfsetdash{}{0pt}%
\pgfpathmoveto{\pgfqpoint{12.548102in}{8.680996in}}%
\pgfpathcurveto{\pgfqpoint{12.558356in}{8.680996in}}{\pgfqpoint{12.568191in}{8.685070in}}{\pgfqpoint{12.575442in}{8.692321in}}%
\pgfpathcurveto{\pgfqpoint{12.582693in}{8.699572in}}{\pgfqpoint{12.586767in}{8.709407in}}{\pgfqpoint{12.586767in}{8.719661in}}%
\pgfpathcurveto{\pgfqpoint{12.586767in}{8.729915in}}{\pgfqpoint{12.582693in}{8.739751in}}{\pgfqpoint{12.575442in}{8.747002in}}%
\pgfpathcurveto{\pgfqpoint{12.568191in}{8.754252in}}{\pgfqpoint{12.558356in}{8.758326in}}{\pgfqpoint{12.548102in}{8.758326in}}%
\pgfpathcurveto{\pgfqpoint{12.537848in}{8.758326in}}{\pgfqpoint{12.528012in}{8.754252in}}{\pgfqpoint{12.520762in}{8.747002in}}%
\pgfpathcurveto{\pgfqpoint{12.513511in}{8.739751in}}{\pgfqpoint{12.509437in}{8.729915in}}{\pgfqpoint{12.509437in}{8.719661in}}%
\pgfpathcurveto{\pgfqpoint{12.509437in}{8.709407in}}{\pgfqpoint{12.513511in}{8.699572in}}{\pgfqpoint{12.520762in}{8.692321in}}%
\pgfpathcurveto{\pgfqpoint{12.528012in}{8.685070in}}{\pgfqpoint{12.537848in}{8.680996in}}{\pgfqpoint{12.548102in}{8.680996in}}%
\pgfpathlineto{\pgfqpoint{12.548102in}{8.680996in}}%
\pgfpathclose%
\pgfusepath{stroke}%
\end{pgfscope}%
\begin{pgfscope}%
\pgfpathrectangle{\pgfqpoint{8.036091in}{4.050000in}}{\pgfqpoint{5.400000in}{5.400000in}}%
\pgfusepath{clip}%
\pgfsetbuttcap%
\pgfsetroundjoin%
\pgfsetlinewidth{1.003750pt}%
\definecolor{currentstroke}{rgb}{0.549020,0.337255,0.294118}%
\pgfsetstrokecolor{currentstroke}%
\pgfsetstrokeopacity{0.900000}%
\pgfsetdash{}{0pt}%
\pgfpathmoveto{\pgfqpoint{12.532061in}{6.650920in}}%
\pgfpathcurveto{\pgfqpoint{12.542316in}{6.650920in}}{\pgfqpoint{12.552151in}{6.654994in}}{\pgfqpoint{12.559402in}{6.662245in}}%
\pgfpathcurveto{\pgfqpoint{12.566652in}{6.669495in}}{\pgfqpoint{12.570726in}{6.679331in}}{\pgfqpoint{12.570726in}{6.689585in}}%
\pgfpathcurveto{\pgfqpoint{12.570726in}{6.699839in}}{\pgfqpoint{12.566652in}{6.709675in}}{\pgfqpoint{12.559402in}{6.716925in}}%
\pgfpathcurveto{\pgfqpoint{12.552151in}{6.724176in}}{\pgfqpoint{12.542316in}{6.728250in}}{\pgfqpoint{12.532061in}{6.728250in}}%
\pgfpathcurveto{\pgfqpoint{12.521807in}{6.728250in}}{\pgfqpoint{12.511972in}{6.724176in}}{\pgfqpoint{12.504721in}{6.716925in}}%
\pgfpathcurveto{\pgfqpoint{12.497470in}{6.709675in}}{\pgfqpoint{12.493396in}{6.699839in}}{\pgfqpoint{12.493396in}{6.689585in}}%
\pgfpathcurveto{\pgfqpoint{12.493396in}{6.679331in}}{\pgfqpoint{12.497470in}{6.669495in}}{\pgfqpoint{12.504721in}{6.662245in}}%
\pgfpathcurveto{\pgfqpoint{12.511972in}{6.654994in}}{\pgfqpoint{12.521807in}{6.650920in}}{\pgfqpoint{12.532061in}{6.650920in}}%
\pgfpathlineto{\pgfqpoint{12.532061in}{6.650920in}}%
\pgfpathclose%
\pgfusepath{stroke}%
\end{pgfscope}%
\begin{pgfscope}%
\pgfsetbuttcap%
\pgfsetroundjoin%
\definecolor{currentfill}{rgb}{0.352941,0.407843,0.450980}%
\pgfsetfillcolor{currentfill}%
\pgfsetlinewidth{0.000000pt}%
\definecolor{currentstroke}{rgb}{0.352941,0.407843,0.450980}%
\pgfsetstrokecolor{currentstroke}%
\pgfsetdash{}{0pt}%
\pgfpathmoveto{\pgfqpoint{8.061408in}{9.595947in}}%
\pgfpathlineto{\pgfqpoint{8.106826in}{9.595947in}}%
\pgfpathlineto{\pgfqpoint{8.106826in}{9.585000in}}%
\pgfpathlineto{\pgfqpoint{8.045760in}{9.585000in}}%
\pgfpathlineto{\pgfqpoint{8.045760in}{9.595947in}}%
\pgfpathquadraticcurveto{\pgfqpoint{8.053161in}{9.603617in}}{\pgfqpoint{8.065943in}{9.616522in}}%
\pgfpathquadraticcurveto{\pgfqpoint{8.078746in}{9.629449in}}{\pgfqpoint{8.082024in}{9.633201in}}%
\pgfpathquadraticcurveto{\pgfqpoint{8.088271in}{9.640211in}}{\pgfqpoint{8.090745in}{9.645076in}}%
\pgfpathquadraticcurveto{\pgfqpoint{8.093239in}{9.649941in}}{\pgfqpoint{8.093239in}{9.654642in}}%
\pgfpathquadraticcurveto{\pgfqpoint{8.093239in}{9.662311in}}{\pgfqpoint{8.087858in}{9.667135in}}%
\pgfpathquadraticcurveto{\pgfqpoint{8.082478in}{9.671980in}}{\pgfqpoint{8.073839in}{9.671980in}}%
\pgfpathquadraticcurveto{\pgfqpoint{8.067716in}{9.671980in}}{\pgfqpoint{8.060913in}{9.669857in}}%
\pgfpathquadraticcurveto{\pgfqpoint{8.054130in}{9.667733in}}{\pgfqpoint{8.046399in}{9.663404in}}%
\pgfpathlineto{\pgfqpoint{8.046399in}{9.676557in}}%
\pgfpathquadraticcurveto{\pgfqpoint{8.054254in}{9.679711in}}{\pgfqpoint{8.061078in}{9.681319in}}%
\pgfpathquadraticcurveto{\pgfqpoint{8.067923in}{9.682928in}}{\pgfqpoint{8.073592in}{9.682928in}}%
\pgfpathquadraticcurveto{\pgfqpoint{8.088539in}{9.682928in}}{\pgfqpoint{8.097424in}{9.675444in}}%
\pgfpathquadraticcurveto{\pgfqpoint{8.106310in}{9.667981in}}{\pgfqpoint{8.106310in}{9.655487in}}%
\pgfpathquadraticcurveto{\pgfqpoint{8.106310in}{9.649550in}}{\pgfqpoint{8.104084in}{9.644231in}}%
\pgfpathquadraticcurveto{\pgfqpoint{8.101878in}{9.638932in}}{\pgfqpoint{8.096002in}{9.631717in}}%
\pgfpathquadraticcurveto{\pgfqpoint{8.094394in}{9.629840in}}{\pgfqpoint{8.085756in}{9.620914in}}%
\pgfpathquadraticcurveto{\pgfqpoint{8.077138in}{9.611987in}}{\pgfqpoint{8.061408in}{9.595947in}}%
\pgfpathlineto{\pgfqpoint{8.061408in}{9.595947in}}%
\pgfpathclose%
\pgfpathmoveto{\pgfqpoint{8.161974in}{9.630686in}}%
\pgfpathquadraticcurveto{\pgfqpoint{8.152697in}{9.630686in}}{\pgfqpoint{8.147378in}{9.625717in}}%
\pgfpathquadraticcurveto{\pgfqpoint{8.142079in}{9.620749in}}{\pgfqpoint{8.142079in}{9.612069in}}%
\pgfpathquadraticcurveto{\pgfqpoint{8.142079in}{9.603369in}}{\pgfqpoint{8.147378in}{9.598401in}}%
\pgfpathquadraticcurveto{\pgfqpoint{8.152697in}{9.593432in}}{\pgfqpoint{8.161974in}{9.593432in}}%
\pgfpathquadraticcurveto{\pgfqpoint{8.171251in}{9.593432in}}{\pgfqpoint{8.176591in}{9.598421in}}%
\pgfpathquadraticcurveto{\pgfqpoint{8.181951in}{9.603431in}}{\pgfqpoint{8.181951in}{9.612069in}}%
\pgfpathquadraticcurveto{\pgfqpoint{8.181951in}{9.620749in}}{\pgfqpoint{8.176632in}{9.625717in}}%
\pgfpathquadraticcurveto{\pgfqpoint{8.171334in}{9.630686in}}{\pgfqpoint{8.161974in}{9.630686in}}%
\pgfpathlineto{\pgfqpoint{8.161974in}{9.630686in}}%
\pgfpathclose%
\pgfpathmoveto{\pgfqpoint{8.148965in}{9.636211in}}%
\pgfpathquadraticcurveto{\pgfqpoint{8.140595in}{9.638273in}}{\pgfqpoint{8.135915in}{9.644004in}}%
\pgfpathquadraticcurveto{\pgfqpoint{8.131256in}{9.649756in}}{\pgfqpoint{8.131256in}{9.658002in}}%
\pgfpathquadraticcurveto{\pgfqpoint{8.131256in}{9.669527in}}{\pgfqpoint{8.139461in}{9.676227in}}%
\pgfpathquadraticcurveto{\pgfqpoint{8.147687in}{9.682928in}}{\pgfqpoint{8.161974in}{9.682928in}}%
\pgfpathquadraticcurveto{\pgfqpoint{8.176344in}{9.682928in}}{\pgfqpoint{8.184528in}{9.676227in}}%
\pgfpathquadraticcurveto{\pgfqpoint{8.192713in}{9.669527in}}{\pgfqpoint{8.192713in}{9.658002in}}%
\pgfpathquadraticcurveto{\pgfqpoint{8.192713in}{9.649756in}}{\pgfqpoint{8.188033in}{9.644004in}}%
\pgfpathquadraticcurveto{\pgfqpoint{8.183374in}{9.638273in}}{\pgfqpoint{8.175066in}{9.636211in}}%
\pgfpathquadraticcurveto{\pgfqpoint{8.184467in}{9.634026in}}{\pgfqpoint{8.189703in}{9.627635in}}%
\pgfpathquadraticcurveto{\pgfqpoint{8.194960in}{9.621264in}}{\pgfqpoint{8.194960in}{9.612069in}}%
\pgfpathquadraticcurveto{\pgfqpoint{8.194960in}{9.598071in}}{\pgfqpoint{8.186425in}{9.590587in}}%
\pgfpathquadraticcurveto{\pgfqpoint{8.177890in}{9.583124in}}{\pgfqpoint{8.161974in}{9.583124in}}%
\pgfpathquadraticcurveto{\pgfqpoint{8.146079in}{9.583124in}}{\pgfqpoint{8.137523in}{9.590587in}}%
\pgfpathquadraticcurveto{\pgfqpoint{8.128988in}{9.598071in}}{\pgfqpoint{8.128988in}{9.612069in}}%
\pgfpathquadraticcurveto{\pgfqpoint{8.128988in}{9.621264in}}{\pgfqpoint{8.134266in}{9.627635in}}%
\pgfpathquadraticcurveto{\pgfqpoint{8.139564in}{9.634026in}}{\pgfqpoint{8.148965in}{9.636211in}}%
\pgfpathlineto{\pgfqpoint{8.148965in}{9.636211in}}%
\pgfpathclose%
\pgfpathmoveto{\pgfqpoint{8.144203in}{9.656765in}}%
\pgfpathquadraticcurveto{\pgfqpoint{8.144203in}{9.649302in}}{\pgfqpoint{8.148862in}{9.645117in}}%
\pgfpathquadraticcurveto{\pgfqpoint{8.153542in}{9.640932in}}{\pgfqpoint{8.161974in}{9.640932in}}%
\pgfpathquadraticcurveto{\pgfqpoint{8.170365in}{9.640932in}}{\pgfqpoint{8.175086in}{9.645117in}}%
\pgfpathquadraticcurveto{\pgfqpoint{8.179828in}{9.649302in}}{\pgfqpoint{8.179828in}{9.656765in}}%
\pgfpathquadraticcurveto{\pgfqpoint{8.179828in}{9.664249in}}{\pgfqpoint{8.175086in}{9.668434in}}%
\pgfpathquadraticcurveto{\pgfqpoint{8.170365in}{9.672619in}}{\pgfqpoint{8.161974in}{9.672619in}}%
\pgfpathquadraticcurveto{\pgfqpoint{8.153542in}{9.672619in}}{\pgfqpoint{8.148862in}{9.668434in}}%
\pgfpathquadraticcurveto{\pgfqpoint{8.144203in}{9.664249in}}{\pgfqpoint{8.144203in}{9.656765in}}%
\pgfpathlineto{\pgfqpoint{8.144203in}{9.656765in}}%
\pgfpathclose%
\pgfpathmoveto{\pgfqpoint{8.245924in}{9.630686in}}%
\pgfpathquadraticcurveto{\pgfqpoint{8.236646in}{9.630686in}}{\pgfqpoint{8.231327in}{9.625717in}}%
\pgfpathquadraticcurveto{\pgfqpoint{8.226029in}{9.620749in}}{\pgfqpoint{8.226029in}{9.612069in}}%
\pgfpathquadraticcurveto{\pgfqpoint{8.226029in}{9.603369in}}{\pgfqpoint{8.231327in}{9.598401in}}%
\pgfpathquadraticcurveto{\pgfqpoint{8.236646in}{9.593432in}}{\pgfqpoint{8.245924in}{9.593432in}}%
\pgfpathquadraticcurveto{\pgfqpoint{8.255201in}{9.593432in}}{\pgfqpoint{8.260541in}{9.598421in}}%
\pgfpathquadraticcurveto{\pgfqpoint{8.265901in}{9.603431in}}{\pgfqpoint{8.265901in}{9.612069in}}%
\pgfpathquadraticcurveto{\pgfqpoint{8.265901in}{9.620749in}}{\pgfqpoint{8.260582in}{9.625717in}}%
\pgfpathquadraticcurveto{\pgfqpoint{8.255284in}{9.630686in}}{\pgfqpoint{8.245924in}{9.630686in}}%
\pgfpathlineto{\pgfqpoint{8.245924in}{9.630686in}}%
\pgfpathclose%
\pgfpathmoveto{\pgfqpoint{8.232915in}{9.636211in}}%
\pgfpathquadraticcurveto{\pgfqpoint{8.224545in}{9.638273in}}{\pgfqpoint{8.219865in}{9.644004in}}%
\pgfpathquadraticcurveto{\pgfqpoint{8.215205in}{9.649756in}}{\pgfqpoint{8.215205in}{9.658002in}}%
\pgfpathquadraticcurveto{\pgfqpoint{8.215205in}{9.669527in}}{\pgfqpoint{8.223411in}{9.676227in}}%
\pgfpathquadraticcurveto{\pgfqpoint{8.231637in}{9.682928in}}{\pgfqpoint{8.245924in}{9.682928in}}%
\pgfpathquadraticcurveto{\pgfqpoint{8.260293in}{9.682928in}}{\pgfqpoint{8.268478in}{9.676227in}}%
\pgfpathquadraticcurveto{\pgfqpoint{8.276663in}{9.669527in}}{\pgfqpoint{8.276663in}{9.658002in}}%
\pgfpathquadraticcurveto{\pgfqpoint{8.276663in}{9.649756in}}{\pgfqpoint{8.271983in}{9.644004in}}%
\pgfpathquadraticcurveto{\pgfqpoint{8.267324in}{9.638273in}}{\pgfqpoint{8.259015in}{9.636211in}}%
\pgfpathquadraticcurveto{\pgfqpoint{8.268416in}{9.634026in}}{\pgfqpoint{8.273653in}{9.627635in}}%
\pgfpathquadraticcurveto{\pgfqpoint{8.278910in}{9.621264in}}{\pgfqpoint{8.278910in}{9.612069in}}%
\pgfpathquadraticcurveto{\pgfqpoint{8.278910in}{9.598071in}}{\pgfqpoint{8.270375in}{9.590587in}}%
\pgfpathquadraticcurveto{\pgfqpoint{8.261840in}{9.583124in}}{\pgfqpoint{8.245924in}{9.583124in}}%
\pgfpathquadraticcurveto{\pgfqpoint{8.230029in}{9.583124in}}{\pgfqpoint{8.221473in}{9.590587in}}%
\pgfpathquadraticcurveto{\pgfqpoint{8.212938in}{9.598071in}}{\pgfqpoint{8.212938in}{9.612069in}}%
\pgfpathquadraticcurveto{\pgfqpoint{8.212938in}{9.621264in}}{\pgfqpoint{8.218215in}{9.627635in}}%
\pgfpathquadraticcurveto{\pgfqpoint{8.223514in}{9.634026in}}{\pgfqpoint{8.232915in}{9.636211in}}%
\pgfpathlineto{\pgfqpoint{8.232915in}{9.636211in}}%
\pgfpathclose%
\pgfpathmoveto{\pgfqpoint{8.228153in}{9.656765in}}%
\pgfpathquadraticcurveto{\pgfqpoint{8.228153in}{9.649302in}}{\pgfqpoint{8.232812in}{9.645117in}}%
\pgfpathquadraticcurveto{\pgfqpoint{8.237492in}{9.640932in}}{\pgfqpoint{8.245924in}{9.640932in}}%
\pgfpathquadraticcurveto{\pgfqpoint{8.254315in}{9.640932in}}{\pgfqpoint{8.259036in}{9.645117in}}%
\pgfpathquadraticcurveto{\pgfqpoint{8.263778in}{9.649302in}}{\pgfqpoint{8.263778in}{9.656765in}}%
\pgfpathquadraticcurveto{\pgfqpoint{8.263778in}{9.664249in}}{\pgfqpoint{8.259036in}{9.668434in}}%
\pgfpathquadraticcurveto{\pgfqpoint{8.254315in}{9.672619in}}{\pgfqpoint{8.245924in}{9.672619in}}%
\pgfpathquadraticcurveto{\pgfqpoint{8.237492in}{9.672619in}}{\pgfqpoint{8.232812in}{9.668434in}}%
\pgfpathquadraticcurveto{\pgfqpoint{8.228153in}{9.664249in}}{\pgfqpoint{8.228153in}{9.656765in}}%
\pgfpathlineto{\pgfqpoint{8.228153in}{9.656765in}}%
\pgfpathclose%
\pgfpathmoveto{\pgfqpoint{8.384115in}{9.646086in}}%
\pgfpathquadraticcurveto{\pgfqpoint{8.382115in}{9.647241in}}{\pgfqpoint{8.379765in}{9.647777in}}%
\pgfpathquadraticcurveto{\pgfqpoint{8.377415in}{9.648333in}}{\pgfqpoint{8.374590in}{9.648333in}}%
\pgfpathquadraticcurveto{\pgfqpoint{8.364529in}{9.648333in}}{\pgfqpoint{8.359149in}{9.641798in}}%
\pgfpathquadraticcurveto{\pgfqpoint{8.353768in}{9.635263in}}{\pgfqpoint{8.353768in}{9.623016in}}%
\pgfpathlineto{\pgfqpoint{8.353768in}{9.585000in}}%
\pgfpathlineto{\pgfqpoint{8.341852in}{9.585000in}}%
\pgfpathlineto{\pgfqpoint{8.341852in}{9.657157in}}%
\pgfpathlineto{\pgfqpoint{8.353768in}{9.657157in}}%
\pgfpathlineto{\pgfqpoint{8.353768in}{9.645942in}}%
\pgfpathquadraticcurveto{\pgfqpoint{8.357520in}{9.652518in}}{\pgfqpoint{8.363499in}{9.655693in}}%
\pgfpathquadraticcurveto{\pgfqpoint{8.369498in}{9.658889in}}{\pgfqpoint{8.378074in}{9.658889in}}%
\pgfpathquadraticcurveto{\pgfqpoint{8.379291in}{9.658889in}}{\pgfqpoint{8.380775in}{9.658724in}}%
\pgfpathquadraticcurveto{\pgfqpoint{8.382260in}{9.658580in}}{\pgfqpoint{8.384053in}{9.658250in}}%
\pgfpathlineto{\pgfqpoint{8.384115in}{9.646086in}}%
\pgfpathlineto{\pgfqpoint{8.384115in}{9.646086in}}%
\pgfpathclose%
\pgfpathmoveto{\pgfqpoint{8.455365in}{9.624047in}}%
\pgfpathlineto{\pgfqpoint{8.455365in}{9.618254in}}%
\pgfpathlineto{\pgfqpoint{8.400855in}{9.618254in}}%
\pgfpathquadraticcurveto{\pgfqpoint{8.401639in}{9.606008in}}{\pgfqpoint{8.408236in}{9.599596in}}%
\pgfpathquadraticcurveto{\pgfqpoint{8.414833in}{9.593185in}}{\pgfqpoint{8.426626in}{9.593185in}}%
\pgfpathquadraticcurveto{\pgfqpoint{8.433450in}{9.593185in}}{\pgfqpoint{8.439862in}{9.594855in}}%
\pgfpathquadraticcurveto{\pgfqpoint{8.446273in}{9.596525in}}{\pgfqpoint{8.452602in}{9.599885in}}%
\pgfpathlineto{\pgfqpoint{8.452602in}{9.588670in}}%
\pgfpathquadraticcurveto{\pgfqpoint{8.446211in}{9.585969in}}{\pgfqpoint{8.439511in}{9.584546in}}%
\pgfpathquadraticcurveto{\pgfqpoint{8.432811in}{9.583124in}}{\pgfqpoint{8.425925in}{9.583124in}}%
\pgfpathquadraticcurveto{\pgfqpoint{8.408648in}{9.583124in}}{\pgfqpoint{8.398567in}{9.593164in}}%
\pgfpathquadraticcurveto{\pgfqpoint{8.388486in}{9.603225in}}{\pgfqpoint{8.388486in}{9.620378in}}%
\pgfpathquadraticcurveto{\pgfqpoint{8.388486in}{9.638087in}}{\pgfqpoint{8.398052in}{9.648478in}}%
\pgfpathquadraticcurveto{\pgfqpoint{8.407618in}{9.658889in}}{\pgfqpoint{8.423863in}{9.658889in}}%
\pgfpathquadraticcurveto{\pgfqpoint{8.438418in}{9.658889in}}{\pgfqpoint{8.446892in}{9.649508in}}%
\pgfpathquadraticcurveto{\pgfqpoint{8.455365in}{9.640149in}}{\pgfqpoint{8.455365in}{9.624047in}}%
\pgfpathlineto{\pgfqpoint{8.455365in}{9.624047in}}%
\pgfpathclose%
\pgfpathmoveto{\pgfqpoint{8.443511in}{9.627531in}}%
\pgfpathquadraticcurveto{\pgfqpoint{8.443387in}{9.637242in}}{\pgfqpoint{8.438068in}{9.643035in}}%
\pgfpathquadraticcurveto{\pgfqpoint{8.432749in}{9.648849in}}{\pgfqpoint{8.423987in}{9.648849in}}%
\pgfpathquadraticcurveto{\pgfqpoint{8.414071in}{9.648849in}}{\pgfqpoint{8.408112in}{9.643241in}}%
\pgfpathquadraticcurveto{\pgfqpoint{8.402154in}{9.637633in}}{\pgfqpoint{8.401247in}{9.627449in}}%
\pgfpathlineto{\pgfqpoint{8.443511in}{9.627531in}}%
\pgfpathlineto{\pgfqpoint{8.443511in}{9.627531in}}%
\pgfpathclose%
\pgfpathmoveto{\pgfqpoint{8.486557in}{9.677650in}}%
\pgfpathlineto{\pgfqpoint{8.486557in}{9.657157in}}%
\pgfpathlineto{\pgfqpoint{8.510967in}{9.657157in}}%
\pgfpathlineto{\pgfqpoint{8.510967in}{9.647942in}}%
\pgfpathlineto{\pgfqpoint{8.486557in}{9.647942in}}%
\pgfpathlineto{\pgfqpoint{8.486557in}{9.608771in}}%
\pgfpathquadraticcurveto{\pgfqpoint{8.486557in}{9.599947in}}{\pgfqpoint{8.488970in}{9.597432in}}%
\pgfpathquadraticcurveto{\pgfqpoint{8.491382in}{9.594916in}}{\pgfqpoint{8.498804in}{9.594916in}}%
\pgfpathlineto{\pgfqpoint{8.510967in}{9.594916in}}%
\pgfpathlineto{\pgfqpoint{8.510967in}{9.585000in}}%
\pgfpathlineto{\pgfqpoint{8.498804in}{9.585000in}}%
\pgfpathquadraticcurveto{\pgfqpoint{8.485073in}{9.585000in}}{\pgfqpoint{8.479857in}{9.590113in}}%
\pgfpathquadraticcurveto{\pgfqpoint{8.474641in}{9.595246in}}{\pgfqpoint{8.474641in}{9.608771in}}%
\pgfpathlineto{\pgfqpoint{8.474641in}{9.647942in}}%
\pgfpathlineto{\pgfqpoint{8.465941in}{9.647942in}}%
\pgfpathlineto{\pgfqpoint{8.465941in}{9.657157in}}%
\pgfpathlineto{\pgfqpoint{8.474641in}{9.657157in}}%
\pgfpathlineto{\pgfqpoint{8.474641in}{9.677650in}}%
\pgfpathlineto{\pgfqpoint{8.486557in}{9.677650in}}%
\pgfpathlineto{\pgfqpoint{8.486557in}{9.677650in}}%
\pgfpathclose%
\pgfpathmoveto{\pgfqpoint{8.559354in}{9.621264in}}%
\pgfpathquadraticcurveto{\pgfqpoint{8.544984in}{9.621264in}}{\pgfqpoint{8.539438in}{9.617986in}}%
\pgfpathquadraticcurveto{\pgfqpoint{8.533893in}{9.614708in}}{\pgfqpoint{8.533893in}{9.606771in}}%
\pgfpathquadraticcurveto{\pgfqpoint{8.533893in}{9.600462in}}{\pgfqpoint{8.538057in}{9.596751in}}%
\pgfpathquadraticcurveto{\pgfqpoint{8.542222in}{9.593061in}}{\pgfqpoint{8.549355in}{9.593061in}}%
\pgfpathquadraticcurveto{\pgfqpoint{8.559230in}{9.593061in}}{\pgfqpoint{8.565188in}{9.600050in}}%
\pgfpathquadraticcurveto{\pgfqpoint{8.571146in}{9.607039in}}{\pgfqpoint{8.571146in}{9.618625in}}%
\pgfpathlineto{\pgfqpoint{8.571146in}{9.621264in}}%
\pgfpathlineto{\pgfqpoint{8.559354in}{9.621264in}}%
\pgfpathlineto{\pgfqpoint{8.559354in}{9.621264in}}%
\pgfpathclose%
\pgfpathmoveto{\pgfqpoint{8.583001in}{9.626171in}}%
\pgfpathlineto{\pgfqpoint{8.583001in}{9.585000in}}%
\pgfpathlineto{\pgfqpoint{8.571146in}{9.585000in}}%
\pgfpathlineto{\pgfqpoint{8.571146in}{9.595947in}}%
\pgfpathquadraticcurveto{\pgfqpoint{8.567085in}{9.589391in}}{\pgfqpoint{8.561024in}{9.586258in}}%
\pgfpathquadraticcurveto{\pgfqpoint{8.554962in}{9.583124in}}{\pgfqpoint{8.546201in}{9.583124in}}%
\pgfpathquadraticcurveto{\pgfqpoint{8.535130in}{9.583124in}}{\pgfqpoint{8.528574in}{9.589350in}}%
\pgfpathquadraticcurveto{\pgfqpoint{8.522038in}{9.595576in}}{\pgfqpoint{8.522038in}{9.606008in}}%
\pgfpathquadraticcurveto{\pgfqpoint{8.522038in}{9.618172in}}{\pgfqpoint{8.530182in}{9.624357in}}%
\pgfpathquadraticcurveto{\pgfqpoint{8.538346in}{9.630541in}}{\pgfqpoint{8.554509in}{9.630541in}}%
\pgfpathlineto{\pgfqpoint{8.571146in}{9.630541in}}%
\pgfpathlineto{\pgfqpoint{8.571146in}{9.631717in}}%
\pgfpathquadraticcurveto{\pgfqpoint{8.571146in}{9.639901in}}{\pgfqpoint{8.565765in}{9.644375in}}%
\pgfpathquadraticcurveto{\pgfqpoint{8.560385in}{9.648849in}}{\pgfqpoint{8.550654in}{9.648849in}}%
\pgfpathquadraticcurveto{\pgfqpoint{8.544469in}{9.648849in}}{\pgfqpoint{8.538593in}{9.647364in}}%
\pgfpathquadraticcurveto{\pgfqpoint{8.532738in}{9.645880in}}{\pgfqpoint{8.527337in}{9.642911in}}%
\pgfpathlineto{\pgfqpoint{8.527337in}{9.653879in}}%
\pgfpathquadraticcurveto{\pgfqpoint{8.533831in}{9.656394in}}{\pgfqpoint{8.539954in}{9.657631in}}%
\pgfpathquadraticcurveto{\pgfqpoint{8.546077in}{9.658889in}}{\pgfqpoint{8.551870in}{9.658889in}}%
\pgfpathquadraticcurveto{\pgfqpoint{8.567538in}{9.658889in}}{\pgfqpoint{8.575270in}{9.650766in}}%
\pgfpathquadraticcurveto{\pgfqpoint{8.583001in}{9.642664in}}{\pgfqpoint{8.583001in}{9.626171in}}%
\pgfpathlineto{\pgfqpoint{8.583001in}{9.626171in}}%
\pgfpathclose%
\pgfpathmoveto{\pgfqpoint{8.607410in}{9.657157in}}%
\pgfpathlineto{\pgfqpoint{8.619265in}{9.657157in}}%
\pgfpathlineto{\pgfqpoint{8.619265in}{9.585000in}}%
\pgfpathlineto{\pgfqpoint{8.607410in}{9.585000in}}%
\pgfpathlineto{\pgfqpoint{8.607410in}{9.657157in}}%
\pgfpathlineto{\pgfqpoint{8.607410in}{9.657157in}}%
\pgfpathclose%
\pgfpathmoveto{\pgfqpoint{8.607410in}{9.685257in}}%
\pgfpathlineto{\pgfqpoint{8.619265in}{9.685257in}}%
\pgfpathlineto{\pgfqpoint{8.619265in}{9.670228in}}%
\pgfpathlineto{\pgfqpoint{8.607410in}{9.670228in}}%
\pgfpathlineto{\pgfqpoint{8.607410in}{9.685257in}}%
\pgfpathlineto{\pgfqpoint{8.607410in}{9.685257in}}%
\pgfpathclose%
\pgfpathmoveto{\pgfqpoint{8.704060in}{9.628562in}}%
\pgfpathlineto{\pgfqpoint{8.704060in}{9.585000in}}%
\pgfpathlineto{\pgfqpoint{8.692205in}{9.585000in}}%
\pgfpathlineto{\pgfqpoint{8.692205in}{9.628171in}}%
\pgfpathquadraticcurveto{\pgfqpoint{8.692205in}{9.638417in}}{\pgfqpoint{8.688206in}{9.643488in}}%
\pgfpathquadraticcurveto{\pgfqpoint{8.684206in}{9.648581in}}{\pgfqpoint{8.676228in}{9.648581in}}%
\pgfpathquadraticcurveto{\pgfqpoint{8.666620in}{9.648581in}}{\pgfqpoint{8.661075in}{9.642458in}}%
\pgfpathquadraticcurveto{\pgfqpoint{8.655529in}{9.636355in}}{\pgfqpoint{8.655529in}{9.625779in}}%
\pgfpathlineto{\pgfqpoint{8.655529in}{9.585000in}}%
\pgfpathlineto{\pgfqpoint{8.643613in}{9.585000in}}%
\pgfpathlineto{\pgfqpoint{8.643613in}{9.657157in}}%
\pgfpathlineto{\pgfqpoint{8.655529in}{9.657157in}}%
\pgfpathlineto{\pgfqpoint{8.655529in}{9.645942in}}%
\pgfpathquadraticcurveto{\pgfqpoint{8.659796in}{9.652457in}}{\pgfqpoint{8.665548in}{9.655673in}}%
\pgfpathquadraticcurveto{\pgfqpoint{8.671321in}{9.658889in}}{\pgfqpoint{8.678867in}{9.658889in}}%
\pgfpathquadraticcurveto{\pgfqpoint{8.691298in}{9.658889in}}{\pgfqpoint{8.697669in}{9.651199in}}%
\pgfpathquadraticcurveto{\pgfqpoint{8.704060in}{9.643509in}}{\pgfqpoint{8.704060in}{9.628562in}}%
\pgfpathlineto{\pgfqpoint{8.704060in}{9.628562in}}%
\pgfpathclose%
\pgfpathmoveto{\pgfqpoint{8.789411in}{9.624047in}}%
\pgfpathlineto{\pgfqpoint{8.789411in}{9.618254in}}%
\pgfpathlineto{\pgfqpoint{8.734902in}{9.618254in}}%
\pgfpathquadraticcurveto{\pgfqpoint{8.735685in}{9.606008in}}{\pgfqpoint{8.742282in}{9.599596in}}%
\pgfpathquadraticcurveto{\pgfqpoint{8.748880in}{9.593185in}}{\pgfqpoint{8.760672in}{9.593185in}}%
\pgfpathquadraticcurveto{\pgfqpoint{8.767496in}{9.593185in}}{\pgfqpoint{8.773908in}{9.594855in}}%
\pgfpathquadraticcurveto{\pgfqpoint{8.780319in}{9.596525in}}{\pgfqpoint{8.786649in}{9.599885in}}%
\pgfpathlineto{\pgfqpoint{8.786649in}{9.588670in}}%
\pgfpathquadraticcurveto{\pgfqpoint{8.780258in}{9.585969in}}{\pgfqpoint{8.773557in}{9.584546in}}%
\pgfpathquadraticcurveto{\pgfqpoint{8.766857in}{9.583124in}}{\pgfqpoint{8.759971in}{9.583124in}}%
\pgfpathquadraticcurveto{\pgfqpoint{8.742695in}{9.583124in}}{\pgfqpoint{8.732613in}{9.593164in}}%
\pgfpathquadraticcurveto{\pgfqpoint{8.722532in}{9.603225in}}{\pgfqpoint{8.722532in}{9.620378in}}%
\pgfpathquadraticcurveto{\pgfqpoint{8.722532in}{9.638087in}}{\pgfqpoint{8.732098in}{9.648478in}}%
\pgfpathquadraticcurveto{\pgfqpoint{8.741664in}{9.658889in}}{\pgfqpoint{8.757909in}{9.658889in}}%
\pgfpathquadraticcurveto{\pgfqpoint{8.772465in}{9.658889in}}{\pgfqpoint{8.780938in}{9.649508in}}%
\pgfpathquadraticcurveto{\pgfqpoint{8.789411in}{9.640149in}}{\pgfqpoint{8.789411in}{9.624047in}}%
\pgfpathlineto{\pgfqpoint{8.789411in}{9.624047in}}%
\pgfpathclose%
\pgfpathmoveto{\pgfqpoint{8.777557in}{9.627531in}}%
\pgfpathquadraticcurveto{\pgfqpoint{8.777433in}{9.637242in}}{\pgfqpoint{8.772114in}{9.643035in}}%
\pgfpathquadraticcurveto{\pgfqpoint{8.766795in}{9.648849in}}{\pgfqpoint{8.758033in}{9.648849in}}%
\pgfpathquadraticcurveto{\pgfqpoint{8.748117in}{9.648849in}}{\pgfqpoint{8.742159in}{9.643241in}}%
\pgfpathquadraticcurveto{\pgfqpoint{8.736201in}{9.637633in}}{\pgfqpoint{8.735293in}{9.627449in}}%
\pgfpathlineto{\pgfqpoint{8.777557in}{9.627531in}}%
\pgfpathlineto{\pgfqpoint{8.777557in}{9.627531in}}%
\pgfpathclose%
\pgfpathmoveto{\pgfqpoint{8.856352in}{9.646210in}}%
\pgfpathlineto{\pgfqpoint{8.856352in}{9.685257in}}%
\pgfpathlineto{\pgfqpoint{8.868207in}{9.685257in}}%
\pgfpathlineto{\pgfqpoint{8.868207in}{9.585000in}}%
\pgfpathlineto{\pgfqpoint{8.856352in}{9.585000in}}%
\pgfpathlineto{\pgfqpoint{8.856352in}{9.595824in}}%
\pgfpathquadraticcurveto{\pgfqpoint{8.852621in}{9.589391in}}{\pgfqpoint{8.846910in}{9.586258in}}%
\pgfpathquadraticcurveto{\pgfqpoint{8.841220in}{9.583124in}}{\pgfqpoint{8.833221in}{9.583124in}}%
\pgfpathquadraticcurveto{\pgfqpoint{8.820150in}{9.583124in}}{\pgfqpoint{8.811924in}{9.593556in}}%
\pgfpathquadraticcurveto{\pgfqpoint{8.803719in}{9.604008in}}{\pgfqpoint{8.803719in}{9.621017in}}%
\pgfpathquadraticcurveto{\pgfqpoint{8.803719in}{9.638025in}}{\pgfqpoint{8.811924in}{9.648457in}}%
\pgfpathquadraticcurveto{\pgfqpoint{8.820150in}{9.658889in}}{\pgfqpoint{8.833221in}{9.658889in}}%
\pgfpathquadraticcurveto{\pgfqpoint{8.841220in}{9.658889in}}{\pgfqpoint{8.846910in}{9.655755in}}%
\pgfpathquadraticcurveto{\pgfqpoint{8.852621in}{9.652642in}}{\pgfqpoint{8.856352in}{9.646210in}}%
\pgfpathlineto{\pgfqpoint{8.856352in}{9.646210in}}%
\pgfpathclose%
\pgfpathmoveto{\pgfqpoint{8.815965in}{9.621017in}}%
\pgfpathquadraticcurveto{\pgfqpoint{8.815965in}{9.607946in}}{\pgfqpoint{8.821346in}{9.600503in}}%
\pgfpathquadraticcurveto{\pgfqpoint{8.826727in}{9.593061in}}{\pgfqpoint{8.836128in}{9.593061in}}%
\pgfpathquadraticcurveto{\pgfqpoint{8.845529in}{9.593061in}}{\pgfqpoint{8.850930in}{9.600503in}}%
\pgfpathquadraticcurveto{\pgfqpoint{8.856352in}{9.607946in}}{\pgfqpoint{8.856352in}{9.621017in}}%
\pgfpathquadraticcurveto{\pgfqpoint{8.856352in}{9.634087in}}{\pgfqpoint{8.850930in}{9.641530in}}%
\pgfpathquadraticcurveto{\pgfqpoint{8.845529in}{9.648972in}}{\pgfqpoint{8.836128in}{9.648972in}}%
\pgfpathquadraticcurveto{\pgfqpoint{8.826727in}{9.648972in}}{\pgfqpoint{8.821346in}{9.641530in}}%
\pgfpathquadraticcurveto{\pgfqpoint{8.815965in}{9.634087in}}{\pgfqpoint{8.815965in}{9.621017in}}%
\pgfpathlineto{\pgfqpoint{8.815965in}{9.621017in}}%
\pgfpathclose%
\pgfpathmoveto{\pgfqpoint{8.955641in}{9.681196in}}%
\pgfpathlineto{\pgfqpoint{8.966588in}{9.681196in}}%
\pgfpathlineto{\pgfqpoint{8.933086in}{9.572754in}}%
\pgfpathlineto{\pgfqpoint{8.922139in}{9.572754in}}%
\pgfpathlineto{\pgfqpoint{8.955641in}{9.681196in}}%
\pgfpathlineto{\pgfqpoint{8.955641in}{9.681196in}}%
\pgfpathclose%
\pgfpathmoveto{\pgfqpoint{9.062062in}{9.636871in}}%
\pgfpathquadraticcurveto{\pgfqpoint{9.071401in}{9.634871in}}{\pgfqpoint{9.076638in}{9.628542in}}%
\pgfpathquadraticcurveto{\pgfqpoint{9.081895in}{9.622233in}}{\pgfqpoint{9.081895in}{9.612956in}}%
\pgfpathquadraticcurveto{\pgfqpoint{9.081895in}{9.598730in}}{\pgfqpoint{9.072102in}{9.590917in}}%
\pgfpathquadraticcurveto{\pgfqpoint{9.062309in}{9.583124in}}{\pgfqpoint{9.044270in}{9.583124in}}%
\pgfpathquadraticcurveto{\pgfqpoint{9.038230in}{9.583124in}}{\pgfqpoint{9.031818in}{9.584320in}}%
\pgfpathquadraticcurveto{\pgfqpoint{9.025406in}{9.585515in}}{\pgfqpoint{9.018582in}{9.587907in}}%
\pgfpathlineto{\pgfqpoint{9.018582in}{9.600462in}}%
\pgfpathquadraticcurveto{\pgfqpoint{9.023984in}{9.597308in}}{\pgfqpoint{9.030416in}{9.595700in}}%
\pgfpathquadraticcurveto{\pgfqpoint{9.036869in}{9.594092in}}{\pgfqpoint{9.043899in}{9.594092in}}%
\pgfpathquadraticcurveto{\pgfqpoint{9.056125in}{9.594092in}}{\pgfqpoint{9.062536in}{9.598916in}}%
\pgfpathquadraticcurveto{\pgfqpoint{9.068948in}{9.603740in}}{\pgfqpoint{9.068948in}{9.612956in}}%
\pgfpathquadraticcurveto{\pgfqpoint{9.068948in}{9.621470in}}{\pgfqpoint{9.062990in}{9.626253in}}%
\pgfpathquadraticcurveto{\pgfqpoint{9.057032in}{9.631057in}}{\pgfqpoint{9.046414in}{9.631057in}}%
\pgfpathlineto{\pgfqpoint{9.035199in}{9.631057in}}%
\pgfpathlineto{\pgfqpoint{9.035199in}{9.641757in}}%
\pgfpathlineto{\pgfqpoint{9.046930in}{9.641757in}}%
\pgfpathquadraticcurveto{\pgfqpoint{9.056516in}{9.641757in}}{\pgfqpoint{9.061608in}{9.645591in}}%
\pgfpathquadraticcurveto{\pgfqpoint{9.066701in}{9.649426in}}{\pgfqpoint{9.066701in}{9.656642in}}%
\pgfpathquadraticcurveto{\pgfqpoint{9.066701in}{9.664043in}}{\pgfqpoint{9.061444in}{9.668001in}}%
\pgfpathquadraticcurveto{\pgfqpoint{9.056207in}{9.671980in}}{\pgfqpoint{9.046414in}{9.671980in}}%
\pgfpathquadraticcurveto{\pgfqpoint{9.041054in}{9.671980in}}{\pgfqpoint{9.034931in}{9.670805in}}%
\pgfpathquadraticcurveto{\pgfqpoint{9.028808in}{9.669651in}}{\pgfqpoint{9.021469in}{9.667218in}}%
\pgfpathlineto{\pgfqpoint{9.021469in}{9.678804in}}%
\pgfpathquadraticcurveto{\pgfqpoint{9.028890in}{9.680866in}}{\pgfqpoint{9.035364in}{9.681897in}}%
\pgfpathquadraticcurveto{\pgfqpoint{9.041837in}{9.682928in}}{\pgfqpoint{9.047569in}{9.682928in}}%
\pgfpathquadraticcurveto{\pgfqpoint{9.062392in}{9.682928in}}{\pgfqpoint{9.071010in}{9.676186in}}%
\pgfpathquadraticcurveto{\pgfqpoint{9.079648in}{9.669465in}}{\pgfqpoint{9.079648in}{9.658002in}}%
\pgfpathquadraticcurveto{\pgfqpoint{9.079648in}{9.650003in}}{\pgfqpoint{9.075071in}{9.644499in}}%
\pgfpathquadraticcurveto{\pgfqpoint{9.070494in}{9.638994in}}{\pgfqpoint{9.062062in}{9.636871in}}%
\pgfpathlineto{\pgfqpoint{9.062062in}{9.636871in}}%
\pgfpathclose%
\pgfpathmoveto{\pgfqpoint{9.106717in}{9.681196in}}%
\pgfpathlineto{\pgfqpoint{9.157804in}{9.681196in}}%
\pgfpathlineto{\pgfqpoint{9.157804in}{9.670228in}}%
\pgfpathlineto{\pgfqpoint{9.118633in}{9.670228in}}%
\pgfpathlineto{\pgfqpoint{9.118633in}{9.646663in}}%
\pgfpathquadraticcurveto{\pgfqpoint{9.121458in}{9.647632in}}{\pgfqpoint{9.124282in}{9.648107in}}%
\pgfpathquadraticcurveto{\pgfqpoint{9.127127in}{9.648581in}}{\pgfqpoint{9.129972in}{9.648581in}}%
\pgfpathquadraticcurveto{\pgfqpoint{9.146074in}{9.648581in}}{\pgfqpoint{9.155475in}{9.639757in}}%
\pgfpathquadraticcurveto{\pgfqpoint{9.164896in}{9.630933in}}{\pgfqpoint{9.164896in}{9.615863in}}%
\pgfpathquadraticcurveto{\pgfqpoint{9.164896in}{9.600339in}}{\pgfqpoint{9.155227in}{9.591721in}}%
\pgfpathquadraticcurveto{\pgfqpoint{9.145558in}{9.583124in}}{\pgfqpoint{9.127972in}{9.583124in}}%
\pgfpathquadraticcurveto{\pgfqpoint{9.121911in}{9.583124in}}{\pgfqpoint{9.115623in}{9.584155in}}%
\pgfpathquadraticcurveto{\pgfqpoint{9.109356in}{9.585186in}}{\pgfqpoint{9.102656in}{9.587247in}}%
\pgfpathlineto{\pgfqpoint{9.102656in}{9.600339in}}%
\pgfpathquadraticcurveto{\pgfqpoint{9.108449in}{9.597184in}}{\pgfqpoint{9.114634in}{9.595638in}}%
\pgfpathquadraticcurveto{\pgfqpoint{9.120819in}{9.594092in}}{\pgfqpoint{9.127704in}{9.594092in}}%
\pgfpathquadraticcurveto{\pgfqpoint{9.138858in}{9.594092in}}{\pgfqpoint{9.145352in}{9.599947in}}%
\pgfpathquadraticcurveto{\pgfqpoint{9.151867in}{9.605802in}}{\pgfqpoint{9.151867in}{9.615863in}}%
\pgfpathquadraticcurveto{\pgfqpoint{9.151867in}{9.625903in}}{\pgfqpoint{9.145352in}{9.631758in}}%
\pgfpathquadraticcurveto{\pgfqpoint{9.138858in}{9.637633in}}{\pgfqpoint{9.127704in}{9.637633in}}%
\pgfpathquadraticcurveto{\pgfqpoint{9.122488in}{9.637633in}}{\pgfqpoint{9.117293in}{9.636479in}}%
\pgfpathquadraticcurveto{\pgfqpoint{9.112118in}{9.635324in}}{\pgfqpoint{9.106717in}{9.632871in}}%
\pgfpathlineto{\pgfqpoint{9.106717in}{9.681196in}}%
\pgfpathlineto{\pgfqpoint{9.106717in}{9.681196in}}%
\pgfpathclose%
\pgfpathmoveto{\pgfqpoint{9.187244in}{9.681196in}}%
\pgfpathlineto{\pgfqpoint{9.249093in}{9.681196in}}%
\pgfpathlineto{\pgfqpoint{9.249093in}{9.675650in}}%
\pgfpathlineto{\pgfqpoint{9.214169in}{9.585000in}}%
\pgfpathlineto{\pgfqpoint{9.200583in}{9.585000in}}%
\pgfpathlineto{\pgfqpoint{9.233446in}{9.670228in}}%
\pgfpathlineto{\pgfqpoint{9.187244in}{9.670228in}}%
\pgfpathlineto{\pgfqpoint{9.187244in}{9.681196in}}%
\pgfpathlineto{\pgfqpoint{9.187244in}{9.681196in}}%
\pgfpathclose%
\pgfpathmoveto{\pgfqpoint{9.366668in}{9.654395in}}%
\pgfpathlineto{\pgfqpoint{9.366668in}{9.643303in}}%
\pgfpathquadraticcurveto{\pgfqpoint{9.361638in}{9.646086in}}{\pgfqpoint{9.356587in}{9.647467in}}%
\pgfpathquadraticcurveto{\pgfqpoint{9.351536in}{9.648849in}}{\pgfqpoint{9.346382in}{9.648849in}}%
\pgfpathquadraticcurveto{\pgfqpoint{9.334837in}{9.648849in}}{\pgfqpoint{9.328446in}{9.641530in}}%
\pgfpathquadraticcurveto{\pgfqpoint{9.322075in}{9.634232in}}{\pgfqpoint{9.322075in}{9.621017in}}%
\pgfpathquadraticcurveto{\pgfqpoint{9.322075in}{9.607802in}}{\pgfqpoint{9.328446in}{9.600483in}}%
\pgfpathquadraticcurveto{\pgfqpoint{9.334837in}{9.593185in}}{\pgfqpoint{9.346382in}{9.593185in}}%
\pgfpathquadraticcurveto{\pgfqpoint{9.351536in}{9.593185in}}{\pgfqpoint{9.356587in}{9.594566in}}%
\pgfpathquadraticcurveto{\pgfqpoint{9.361638in}{9.595947in}}{\pgfqpoint{9.366668in}{9.598730in}}%
\pgfpathlineto{\pgfqpoint{9.366668in}{9.587763in}}%
\pgfpathquadraticcurveto{\pgfqpoint{9.361700in}{9.585454in}}{\pgfqpoint{9.356381in}{9.584299in}}%
\pgfpathquadraticcurveto{\pgfqpoint{9.351082in}{9.583124in}}{\pgfqpoint{9.345083in}{9.583124in}}%
\pgfpathquadraticcurveto{\pgfqpoint{9.328775in}{9.583124in}}{\pgfqpoint{9.319168in}{9.593370in}}%
\pgfpathquadraticcurveto{\pgfqpoint{9.309582in}{9.603617in}}{\pgfqpoint{9.309582in}{9.621017in}}%
\pgfpathquadraticcurveto{\pgfqpoint{9.309582in}{9.638664in}}{\pgfqpoint{9.319271in}{9.648766in}}%
\pgfpathquadraticcurveto{\pgfqpoint{9.328982in}{9.658889in}}{\pgfqpoint{9.345866in}{9.658889in}}%
\pgfpathquadraticcurveto{\pgfqpoint{9.351330in}{9.658889in}}{\pgfqpoint{9.356546in}{9.657755in}}%
\pgfpathquadraticcurveto{\pgfqpoint{9.361761in}{9.656642in}}{\pgfqpoint{9.366668in}{9.654395in}}%
\pgfpathlineto{\pgfqpoint{9.366668in}{9.654395in}}%
\pgfpathclose%
\pgfpathmoveto{\pgfqpoint{9.415240in}{9.648849in}}%
\pgfpathquadraticcurveto{\pgfqpoint{9.405715in}{9.648849in}}{\pgfqpoint{9.400170in}{9.641406in}}%
\pgfpathquadraticcurveto{\pgfqpoint{9.394624in}{9.633964in}}{\pgfqpoint{9.394624in}{9.621017in}}%
\pgfpathquadraticcurveto{\pgfqpoint{9.394624in}{9.608070in}}{\pgfqpoint{9.400128in}{9.600627in}}%
\pgfpathquadraticcurveto{\pgfqpoint{9.405654in}{9.593185in}}{\pgfqpoint{9.415240in}{9.593185in}}%
\pgfpathquadraticcurveto{\pgfqpoint{9.424724in}{9.593185in}}{\pgfqpoint{9.430249in}{9.600648in}}%
\pgfpathquadraticcurveto{\pgfqpoint{9.435795in}{9.608132in}}{\pgfqpoint{9.435795in}{9.621017in}}%
\pgfpathquadraticcurveto{\pgfqpoint{9.435795in}{9.633840in}}{\pgfqpoint{9.430249in}{9.641344in}}%
\pgfpathquadraticcurveto{\pgfqpoint{9.424724in}{9.648849in}}{\pgfqpoint{9.415240in}{9.648849in}}%
\pgfpathlineto{\pgfqpoint{9.415240in}{9.648849in}}%
\pgfpathclose%
\pgfpathmoveto{\pgfqpoint{9.415240in}{9.658889in}}%
\pgfpathquadraticcurveto{\pgfqpoint{9.430702in}{9.658889in}}{\pgfqpoint{9.439526in}{9.648828in}}%
\pgfpathquadraticcurveto{\pgfqpoint{9.448371in}{9.638788in}}{\pgfqpoint{9.448371in}{9.621017in}}%
\pgfpathquadraticcurveto{\pgfqpoint{9.448371in}{9.603307in}}{\pgfqpoint{9.439526in}{9.593205in}}%
\pgfpathquadraticcurveto{\pgfqpoint{9.430702in}{9.583124in}}{\pgfqpoint{9.415240in}{9.583124in}}%
\pgfpathquadraticcurveto{\pgfqpoint{9.399716in}{9.583124in}}{\pgfqpoint{9.390913in}{9.593205in}}%
\pgfpathquadraticcurveto{\pgfqpoint{9.382130in}{9.603307in}}{\pgfqpoint{9.382130in}{9.621017in}}%
\pgfpathquadraticcurveto{\pgfqpoint{9.382130in}{9.638788in}}{\pgfqpoint{9.390913in}{9.648828in}}%
\pgfpathquadraticcurveto{\pgfqpoint{9.399716in}{9.658889in}}{\pgfqpoint{9.415240in}{9.658889in}}%
\pgfpathlineto{\pgfqpoint{9.415240in}{9.658889in}}%
\pgfpathclose%
\pgfpathmoveto{\pgfqpoint{9.524197in}{9.643303in}}%
\pgfpathquadraticcurveto{\pgfqpoint{9.528651in}{9.651302in}}{\pgfqpoint{9.534835in}{9.655095in}}%
\pgfpathquadraticcurveto{\pgfqpoint{9.541020in}{9.658889in}}{\pgfqpoint{9.549391in}{9.658889in}}%
\pgfpathquadraticcurveto{\pgfqpoint{9.560668in}{9.658889in}}{\pgfqpoint{9.566791in}{9.650993in}}%
\pgfpathquadraticcurveto{\pgfqpoint{9.572914in}{9.643117in}}{\pgfqpoint{9.572914in}{9.628562in}}%
\pgfpathlineto{\pgfqpoint{9.572914in}{9.585000in}}%
\pgfpathlineto{\pgfqpoint{9.560998in}{9.585000in}}%
\pgfpathlineto{\pgfqpoint{9.560998in}{9.628171in}}%
\pgfpathquadraticcurveto{\pgfqpoint{9.560998in}{9.638541in}}{\pgfqpoint{9.557307in}{9.643550in}}%
\pgfpathquadraticcurveto{\pgfqpoint{9.553638in}{9.648581in}}{\pgfqpoint{9.546113in}{9.648581in}}%
\pgfpathquadraticcurveto{\pgfqpoint{9.536897in}{9.648581in}}{\pgfqpoint{9.531537in}{9.642458in}}%
\pgfpathquadraticcurveto{\pgfqpoint{9.526197in}{9.636355in}}{\pgfqpoint{9.526197in}{9.625779in}}%
\pgfpathlineto{\pgfqpoint{9.526197in}{9.585000in}}%
\pgfpathlineto{\pgfqpoint{9.514281in}{9.585000in}}%
\pgfpathlineto{\pgfqpoint{9.514281in}{9.628171in}}%
\pgfpathquadraticcurveto{\pgfqpoint{9.514281in}{9.638602in}}{\pgfqpoint{9.510611in}{9.643592in}}%
\pgfpathquadraticcurveto{\pgfqpoint{9.506942in}{9.648581in}}{\pgfqpoint{9.499272in}{9.648581in}}%
\pgfpathquadraticcurveto{\pgfqpoint{9.490181in}{9.648581in}}{\pgfqpoint{9.484820in}{9.642437in}}%
\pgfpathquadraticcurveto{\pgfqpoint{9.479481in}{9.636293in}}{\pgfqpoint{9.479481in}{9.625779in}}%
\pgfpathlineto{\pgfqpoint{9.479481in}{9.585000in}}%
\pgfpathlineto{\pgfqpoint{9.467564in}{9.585000in}}%
\pgfpathlineto{\pgfqpoint{9.467564in}{9.657157in}}%
\pgfpathlineto{\pgfqpoint{9.479481in}{9.657157in}}%
\pgfpathlineto{\pgfqpoint{9.479481in}{9.645942in}}%
\pgfpathquadraticcurveto{\pgfqpoint{9.483542in}{9.652580in}}{\pgfqpoint{9.489212in}{9.655735in}}%
\pgfpathquadraticcurveto{\pgfqpoint{9.494881in}{9.658889in}}{\pgfqpoint{9.502674in}{9.658889in}}%
\pgfpathquadraticcurveto{\pgfqpoint{9.510549in}{9.658889in}}{\pgfqpoint{9.516054in}{9.654889in}}%
\pgfpathquadraticcurveto{\pgfqpoint{9.521559in}{9.650910in}}{\pgfqpoint{9.524197in}{9.643303in}}%
\pgfpathlineto{\pgfqpoint{9.524197in}{9.643303in}}%
\pgfpathclose%
\pgfpathmoveto{\pgfqpoint{9.608003in}{9.595824in}}%
\pgfpathlineto{\pgfqpoint{9.608003in}{9.557560in}}%
\pgfpathlineto{\pgfqpoint{9.596087in}{9.557560in}}%
\pgfpathlineto{\pgfqpoint{9.596087in}{9.657157in}}%
\pgfpathlineto{\pgfqpoint{9.608003in}{9.657157in}}%
\pgfpathlineto{\pgfqpoint{9.608003in}{9.646210in}}%
\pgfpathquadraticcurveto{\pgfqpoint{9.611755in}{9.652642in}}{\pgfqpoint{9.617445in}{9.655755in}}%
\pgfpathquadraticcurveto{\pgfqpoint{9.623156in}{9.658889in}}{\pgfqpoint{9.631072in}{9.658889in}}%
\pgfpathquadraticcurveto{\pgfqpoint{9.644226in}{9.658889in}}{\pgfqpoint{9.652431in}{9.648457in}}%
\pgfpathquadraticcurveto{\pgfqpoint{9.660657in}{9.638025in}}{\pgfqpoint{9.660657in}{9.621017in}}%
\pgfpathquadraticcurveto{\pgfqpoint{9.660657in}{9.604008in}}{\pgfqpoint{9.652431in}{9.593556in}}%
\pgfpathquadraticcurveto{\pgfqpoint{9.644226in}{9.583124in}}{\pgfqpoint{9.631072in}{9.583124in}}%
\pgfpathquadraticcurveto{\pgfqpoint{9.623156in}{9.583124in}}{\pgfqpoint{9.617445in}{9.586258in}}%
\pgfpathquadraticcurveto{\pgfqpoint{9.611755in}{9.589391in}}{\pgfqpoint{9.608003in}{9.595824in}}%
\pgfpathlineto{\pgfqpoint{9.608003in}{9.595824in}}%
\pgfpathclose%
\pgfpathmoveto{\pgfqpoint{9.648349in}{9.621017in}}%
\pgfpathquadraticcurveto{\pgfqpoint{9.648349in}{9.634087in}}{\pgfqpoint{9.642968in}{9.641530in}}%
\pgfpathquadraticcurveto{\pgfqpoint{9.637587in}{9.648972in}}{\pgfqpoint{9.628186in}{9.648972in}}%
\pgfpathquadraticcurveto{\pgfqpoint{9.618765in}{9.648972in}}{\pgfqpoint{9.613384in}{9.641530in}}%
\pgfpathquadraticcurveto{\pgfqpoint{9.608003in}{9.634087in}}{\pgfqpoint{9.608003in}{9.621017in}}%
\pgfpathquadraticcurveto{\pgfqpoint{9.608003in}{9.607946in}}{\pgfqpoint{9.613384in}{9.600503in}}%
\pgfpathquadraticcurveto{\pgfqpoint{9.618765in}{9.593061in}}{\pgfqpoint{9.628186in}{9.593061in}}%
\pgfpathquadraticcurveto{\pgfqpoint{9.637587in}{9.593061in}}{\pgfqpoint{9.642968in}{9.600503in}}%
\pgfpathquadraticcurveto{\pgfqpoint{9.648349in}{9.607946in}}{\pgfqpoint{9.648349in}{9.621017in}}%
\pgfpathlineto{\pgfqpoint{9.648349in}{9.621017in}}%
\pgfpathclose%
\pgfpathmoveto{\pgfqpoint{9.713105in}{9.621264in}}%
\pgfpathquadraticcurveto{\pgfqpoint{9.698735in}{9.621264in}}{\pgfqpoint{9.693189in}{9.617986in}}%
\pgfpathquadraticcurveto{\pgfqpoint{9.687644in}{9.614708in}}{\pgfqpoint{9.687644in}{9.606771in}}%
\pgfpathquadraticcurveto{\pgfqpoint{9.687644in}{9.600462in}}{\pgfqpoint{9.691808in}{9.596751in}}%
\pgfpathquadraticcurveto{\pgfqpoint{9.695973in}{9.593061in}}{\pgfqpoint{9.703106in}{9.593061in}}%
\pgfpathquadraticcurveto{\pgfqpoint{9.712981in}{9.593061in}}{\pgfqpoint{9.718939in}{9.600050in}}%
\pgfpathquadraticcurveto{\pgfqpoint{9.724897in}{9.607039in}}{\pgfqpoint{9.724897in}{9.618625in}}%
\pgfpathlineto{\pgfqpoint{9.724897in}{9.621264in}}%
\pgfpathlineto{\pgfqpoint{9.713105in}{9.621264in}}%
\pgfpathlineto{\pgfqpoint{9.713105in}{9.621264in}}%
\pgfpathclose%
\pgfpathmoveto{\pgfqpoint{9.736752in}{9.626171in}}%
\pgfpathlineto{\pgfqpoint{9.736752in}{9.585000in}}%
\pgfpathlineto{\pgfqpoint{9.724897in}{9.585000in}}%
\pgfpathlineto{\pgfqpoint{9.724897in}{9.595947in}}%
\pgfpathquadraticcurveto{\pgfqpoint{9.720836in}{9.589391in}}{\pgfqpoint{9.714775in}{9.586258in}}%
\pgfpathquadraticcurveto{\pgfqpoint{9.708714in}{9.583124in}}{\pgfqpoint{9.699952in}{9.583124in}}%
\pgfpathquadraticcurveto{\pgfqpoint{9.688881in}{9.583124in}}{\pgfqpoint{9.682325in}{9.589350in}}%
\pgfpathquadraticcurveto{\pgfqpoint{9.675789in}{9.595576in}}{\pgfqpoint{9.675789in}{9.606008in}}%
\pgfpathquadraticcurveto{\pgfqpoint{9.675789in}{9.618172in}}{\pgfqpoint{9.683933in}{9.624357in}}%
\pgfpathquadraticcurveto{\pgfqpoint{9.692097in}{9.630541in}}{\pgfqpoint{9.708260in}{9.630541in}}%
\pgfpathlineto{\pgfqpoint{9.724897in}{9.630541in}}%
\pgfpathlineto{\pgfqpoint{9.724897in}{9.631717in}}%
\pgfpathquadraticcurveto{\pgfqpoint{9.724897in}{9.639901in}}{\pgfqpoint{9.719516in}{9.644375in}}%
\pgfpathquadraticcurveto{\pgfqpoint{9.714136in}{9.648849in}}{\pgfqpoint{9.704405in}{9.648849in}}%
\pgfpathquadraticcurveto{\pgfqpoint{9.698220in}{9.648849in}}{\pgfqpoint{9.692344in}{9.647364in}}%
\pgfpathquadraticcurveto{\pgfqpoint{9.686489in}{9.645880in}}{\pgfqpoint{9.681088in}{9.642911in}}%
\pgfpathlineto{\pgfqpoint{9.681088in}{9.653879in}}%
\pgfpathquadraticcurveto{\pgfqpoint{9.687582in}{9.656394in}}{\pgfqpoint{9.693705in}{9.657631in}}%
\pgfpathquadraticcurveto{\pgfqpoint{9.699828in}{9.658889in}}{\pgfqpoint{9.705621in}{9.658889in}}%
\pgfpathquadraticcurveto{\pgfqpoint{9.721289in}{9.658889in}}{\pgfqpoint{9.729021in}{9.650766in}}%
\pgfpathquadraticcurveto{\pgfqpoint{9.736752in}{9.642664in}}{\pgfqpoint{9.736752in}{9.626171in}}%
\pgfpathlineto{\pgfqpoint{9.736752in}{9.626171in}}%
\pgfpathclose%
\pgfpathmoveto{\pgfqpoint{9.802971in}{9.646086in}}%
\pgfpathquadraticcurveto{\pgfqpoint{9.800972in}{9.647241in}}{\pgfqpoint{9.798621in}{9.647777in}}%
\pgfpathquadraticcurveto{\pgfqpoint{9.796271in}{9.648333in}}{\pgfqpoint{9.793447in}{9.648333in}}%
\pgfpathquadraticcurveto{\pgfqpoint{9.783386in}{9.648333in}}{\pgfqpoint{9.778005in}{9.641798in}}%
\pgfpathquadraticcurveto{\pgfqpoint{9.772624in}{9.635263in}}{\pgfqpoint{9.772624in}{9.623016in}}%
\pgfpathlineto{\pgfqpoint{9.772624in}{9.585000in}}%
\pgfpathlineto{\pgfqpoint{9.760708in}{9.585000in}}%
\pgfpathlineto{\pgfqpoint{9.760708in}{9.657157in}}%
\pgfpathlineto{\pgfqpoint{9.772624in}{9.657157in}}%
\pgfpathlineto{\pgfqpoint{9.772624in}{9.645942in}}%
\pgfpathquadraticcurveto{\pgfqpoint{9.776376in}{9.652518in}}{\pgfqpoint{9.782355in}{9.655693in}}%
\pgfpathquadraticcurveto{\pgfqpoint{9.788354in}{9.658889in}}{\pgfqpoint{9.796931in}{9.658889in}}%
\pgfpathquadraticcurveto{\pgfqpoint{9.798147in}{9.658889in}}{\pgfqpoint{9.799631in}{9.658724in}}%
\pgfpathquadraticcurveto{\pgfqpoint{9.801116in}{9.658580in}}{\pgfqpoint{9.802909in}{9.658250in}}%
\pgfpathlineto{\pgfqpoint{9.802971in}{9.646086in}}%
\pgfpathlineto{\pgfqpoint{9.802971in}{9.646086in}}%
\pgfpathclose%
\pgfpathmoveto{\pgfqpoint{9.815403in}{9.657157in}}%
\pgfpathlineto{\pgfqpoint{9.827257in}{9.657157in}}%
\pgfpathlineto{\pgfqpoint{9.827257in}{9.585000in}}%
\pgfpathlineto{\pgfqpoint{9.815403in}{9.585000in}}%
\pgfpathlineto{\pgfqpoint{9.815403in}{9.657157in}}%
\pgfpathlineto{\pgfqpoint{9.815403in}{9.657157in}}%
\pgfpathclose%
\pgfpathmoveto{\pgfqpoint{9.815403in}{9.685257in}}%
\pgfpathlineto{\pgfqpoint{9.827257in}{9.685257in}}%
\pgfpathlineto{\pgfqpoint{9.827257in}{9.670228in}}%
\pgfpathlineto{\pgfqpoint{9.815403in}{9.670228in}}%
\pgfpathlineto{\pgfqpoint{9.815403in}{9.685257in}}%
\pgfpathlineto{\pgfqpoint{9.815403in}{9.685257in}}%
\pgfpathclose%
\pgfpathmoveto{\pgfqpoint{9.898054in}{9.655034in}}%
\pgfpathlineto{\pgfqpoint{9.898054in}{9.643818in}}%
\pgfpathquadraticcurveto{\pgfqpoint{9.893044in}{9.646395in}}{\pgfqpoint{9.887622in}{9.647674in}}%
\pgfpathquadraticcurveto{\pgfqpoint{9.882220in}{9.648972in}}{\pgfqpoint{9.876407in}{9.648972in}}%
\pgfpathquadraticcurveto{\pgfqpoint{9.867583in}{9.648972in}}{\pgfqpoint{9.863171in}{9.646272in}}%
\pgfpathquadraticcurveto{\pgfqpoint{9.858759in}{9.643571in}}{\pgfqpoint{9.858759in}{9.638149in}}%
\pgfpathquadraticcurveto{\pgfqpoint{9.858759in}{9.634026in}}{\pgfqpoint{9.861913in}{9.631675in}}%
\pgfpathquadraticcurveto{\pgfqpoint{9.865068in}{9.629325in}}{\pgfqpoint{9.874613in}{9.627202in}}%
\pgfpathlineto{\pgfqpoint{9.878674in}{9.626294in}}%
\pgfpathquadraticcurveto{\pgfqpoint{9.891292in}{9.623594in}}{\pgfqpoint{9.896611in}{9.618666in}}%
\pgfpathquadraticcurveto{\pgfqpoint{9.901930in}{9.613739in}}{\pgfqpoint{9.901930in}{9.604915in}}%
\pgfpathquadraticcurveto{\pgfqpoint{9.901930in}{9.594855in}}{\pgfqpoint{9.893972in}{9.588979in}}%
\pgfpathquadraticcurveto{\pgfqpoint{9.886014in}{9.583124in}}{\pgfqpoint{9.872098in}{9.583124in}}%
\pgfpathquadraticcurveto{\pgfqpoint{9.866305in}{9.583124in}}{\pgfqpoint{9.860017in}{9.584258in}}%
\pgfpathquadraticcurveto{\pgfqpoint{9.853729in}{9.585392in}}{\pgfqpoint{9.846781in}{9.587639in}}%
\pgfpathlineto{\pgfqpoint{9.846781in}{9.599885in}}%
\pgfpathquadraticcurveto{\pgfqpoint{9.853358in}{9.596463in}}{\pgfqpoint{9.859728in}{9.594752in}}%
\pgfpathquadraticcurveto{\pgfqpoint{9.866099in}{9.593061in}}{\pgfqpoint{9.872366in}{9.593061in}}%
\pgfpathquadraticcurveto{\pgfqpoint{9.880736in}{9.593061in}}{\pgfqpoint{9.885230in}{9.595927in}}%
\pgfpathquadraticcurveto{\pgfqpoint{9.889745in}{9.598792in}}{\pgfqpoint{9.889745in}{9.604008in}}%
\pgfpathquadraticcurveto{\pgfqpoint{9.889745in}{9.608832in}}{\pgfqpoint{9.886488in}{9.611410in}}%
\pgfpathquadraticcurveto{\pgfqpoint{9.883251in}{9.613987in}}{\pgfqpoint{9.872222in}{9.616378in}}%
\pgfpathlineto{\pgfqpoint{9.868098in}{9.617347in}}%
\pgfpathquadraticcurveto{\pgfqpoint{9.857089in}{9.619656in}}{\pgfqpoint{9.852182in}{9.624460in}}%
\pgfpathquadraticcurveto{\pgfqpoint{9.847296in}{9.629263in}}{\pgfqpoint{9.847296in}{9.637633in}}%
\pgfpathquadraticcurveto{\pgfqpoint{9.847296in}{9.647818in}}{\pgfqpoint{9.854512in}{9.653343in}}%
\pgfpathquadraticcurveto{\pgfqpoint{9.861728in}{9.658889in}}{\pgfqpoint{9.875005in}{9.658889in}}%
\pgfpathquadraticcurveto{\pgfqpoint{9.881561in}{9.658889in}}{\pgfqpoint{9.887354in}{9.657920in}}%
\pgfpathquadraticcurveto{\pgfqpoint{9.893168in}{9.656972in}}{\pgfqpoint{9.898054in}{9.655034in}}%
\pgfpathlineto{\pgfqpoint{9.898054in}{9.655034in}}%
\pgfpathclose%
\pgfpathmoveto{\pgfqpoint{9.948749in}{9.648849in}}%
\pgfpathquadraticcurveto{\pgfqpoint{9.939225in}{9.648849in}}{\pgfqpoint{9.933679in}{9.641406in}}%
\pgfpathquadraticcurveto{\pgfqpoint{9.928133in}{9.633964in}}{\pgfqpoint{9.928133in}{9.621017in}}%
\pgfpathquadraticcurveto{\pgfqpoint{9.928133in}{9.608070in}}{\pgfqpoint{9.933638in}{9.600627in}}%
\pgfpathquadraticcurveto{\pgfqpoint{9.939163in}{9.593185in}}{\pgfqpoint{9.948749in}{9.593185in}}%
\pgfpathquadraticcurveto{\pgfqpoint{9.958233in}{9.593185in}}{\pgfqpoint{9.963758in}{9.600648in}}%
\pgfpathquadraticcurveto{\pgfqpoint{9.969304in}{9.608132in}}{\pgfqpoint{9.969304in}{9.621017in}}%
\pgfpathquadraticcurveto{\pgfqpoint{9.969304in}{9.633840in}}{\pgfqpoint{9.963758in}{9.641344in}}%
\pgfpathquadraticcurveto{\pgfqpoint{9.958233in}{9.648849in}}{\pgfqpoint{9.948749in}{9.648849in}}%
\pgfpathlineto{\pgfqpoint{9.948749in}{9.648849in}}%
\pgfpathclose%
\pgfpathmoveto{\pgfqpoint{9.948749in}{9.658889in}}%
\pgfpathquadraticcurveto{\pgfqpoint{9.964212in}{9.658889in}}{\pgfqpoint{9.973035in}{9.648828in}}%
\pgfpathquadraticcurveto{\pgfqpoint{9.981880in}{9.638788in}}{\pgfqpoint{9.981880in}{9.621017in}}%
\pgfpathquadraticcurveto{\pgfqpoint{9.981880in}{9.603307in}}{\pgfqpoint{9.973035in}{9.593205in}}%
\pgfpathquadraticcurveto{\pgfqpoint{9.964212in}{9.583124in}}{\pgfqpoint{9.948749in}{9.583124in}}%
\pgfpathquadraticcurveto{\pgfqpoint{9.933225in}{9.583124in}}{\pgfqpoint{9.924422in}{9.593205in}}%
\pgfpathquadraticcurveto{\pgfqpoint{9.915640in}{9.603307in}}{\pgfqpoint{9.915640in}{9.621017in}}%
\pgfpathquadraticcurveto{\pgfqpoint{9.915640in}{9.638788in}}{\pgfqpoint{9.924422in}{9.648828in}}%
\pgfpathquadraticcurveto{\pgfqpoint{9.933225in}{9.658889in}}{\pgfqpoint{9.948749in}{9.658889in}}%
\pgfpathlineto{\pgfqpoint{9.948749in}{9.658889in}}%
\pgfpathclose%
\pgfpathmoveto{\pgfqpoint{10.061521in}{9.628562in}}%
\pgfpathlineto{\pgfqpoint{10.061521in}{9.585000in}}%
\pgfpathlineto{\pgfqpoint{10.049666in}{9.585000in}}%
\pgfpathlineto{\pgfqpoint{10.049666in}{9.628171in}}%
\pgfpathquadraticcurveto{\pgfqpoint{10.049666in}{9.638417in}}{\pgfqpoint{10.045667in}{9.643488in}}%
\pgfpathquadraticcurveto{\pgfqpoint{10.041667in}{9.648581in}}{\pgfqpoint{10.033689in}{9.648581in}}%
\pgfpathquadraticcurveto{\pgfqpoint{10.024081in}{9.648581in}}{\pgfqpoint{10.018536in}{9.642458in}}%
\pgfpathquadraticcurveto{\pgfqpoint{10.012990in}{9.636355in}}{\pgfqpoint{10.012990in}{9.625779in}}%
\pgfpathlineto{\pgfqpoint{10.012990in}{9.585000in}}%
\pgfpathlineto{\pgfqpoint{10.001074in}{9.585000in}}%
\pgfpathlineto{\pgfqpoint{10.001074in}{9.657157in}}%
\pgfpathlineto{\pgfqpoint{10.012990in}{9.657157in}}%
\pgfpathlineto{\pgfqpoint{10.012990in}{9.645942in}}%
\pgfpathquadraticcurveto{\pgfqpoint{10.017257in}{9.652457in}}{\pgfqpoint{10.023009in}{9.655673in}}%
\pgfpathquadraticcurveto{\pgfqpoint{10.028782in}{9.658889in}}{\pgfqpoint{10.036327in}{9.658889in}}%
\pgfpathquadraticcurveto{\pgfqpoint{10.048759in}{9.658889in}}{\pgfqpoint{10.055130in}{9.651199in}}%
\pgfpathquadraticcurveto{\pgfqpoint{10.061521in}{9.643509in}}{\pgfqpoint{10.061521in}{9.628562in}}%
\pgfpathlineto{\pgfqpoint{10.061521in}{9.628562in}}%
\pgfpathclose%
\pgfpathmoveto{\pgfqpoint{10.131142in}{9.655034in}}%
\pgfpathlineto{\pgfqpoint{10.131142in}{9.643818in}}%
\pgfpathquadraticcurveto{\pgfqpoint{10.126132in}{9.646395in}}{\pgfqpoint{10.120710in}{9.647674in}}%
\pgfpathquadraticcurveto{\pgfqpoint{10.115309in}{9.648972in}}{\pgfqpoint{10.109495in}{9.648972in}}%
\pgfpathquadraticcurveto{\pgfqpoint{10.100671in}{9.648972in}}{\pgfqpoint{10.096259in}{9.646272in}}%
\pgfpathquadraticcurveto{\pgfqpoint{10.091847in}{9.643571in}}{\pgfqpoint{10.091847in}{9.638149in}}%
\pgfpathquadraticcurveto{\pgfqpoint{10.091847in}{9.634026in}}{\pgfqpoint{10.095001in}{9.631675in}}%
\pgfpathquadraticcurveto{\pgfqpoint{10.098156in}{9.629325in}}{\pgfqpoint{10.107701in}{9.627202in}}%
\pgfpathlineto{\pgfqpoint{10.111763in}{9.626294in}}%
\pgfpathquadraticcurveto{\pgfqpoint{10.124380in}{9.623594in}}{\pgfqpoint{10.129699in}{9.618666in}}%
\pgfpathquadraticcurveto{\pgfqpoint{10.135018in}{9.613739in}}{\pgfqpoint{10.135018in}{9.604915in}}%
\pgfpathquadraticcurveto{\pgfqpoint{10.135018in}{9.594855in}}{\pgfqpoint{10.127060in}{9.588979in}}%
\pgfpathquadraticcurveto{\pgfqpoint{10.119102in}{9.583124in}}{\pgfqpoint{10.105186in}{9.583124in}}%
\pgfpathquadraticcurveto{\pgfqpoint{10.099393in}{9.583124in}}{\pgfqpoint{10.093105in}{9.584258in}}%
\pgfpathquadraticcurveto{\pgfqpoint{10.086817in}{9.585392in}}{\pgfqpoint{10.079869in}{9.587639in}}%
\pgfpathlineto{\pgfqpoint{10.079869in}{9.599885in}}%
\pgfpathquadraticcurveto{\pgfqpoint{10.086446in}{9.596463in}}{\pgfqpoint{10.092816in}{9.594752in}}%
\pgfpathquadraticcurveto{\pgfqpoint{10.099187in}{9.593061in}}{\pgfqpoint{10.105454in}{9.593061in}}%
\pgfpathquadraticcurveto{\pgfqpoint{10.113824in}{9.593061in}}{\pgfqpoint{10.118319in}{9.595927in}}%
\pgfpathquadraticcurveto{\pgfqpoint{10.122834in}{9.598792in}}{\pgfqpoint{10.122834in}{9.604008in}}%
\pgfpathquadraticcurveto{\pgfqpoint{10.122834in}{9.608832in}}{\pgfqpoint{10.119576in}{9.611410in}}%
\pgfpathquadraticcurveto{\pgfqpoint{10.116339in}{9.613987in}}{\pgfqpoint{10.105310in}{9.616378in}}%
\pgfpathlineto{\pgfqpoint{10.101186in}{9.617347in}}%
\pgfpathquadraticcurveto{\pgfqpoint{10.090177in}{9.619656in}}{\pgfqpoint{10.085271in}{9.624460in}}%
\pgfpathquadraticcurveto{\pgfqpoint{10.080385in}{9.629263in}}{\pgfqpoint{10.080385in}{9.637633in}}%
\pgfpathquadraticcurveto{\pgfqpoint{10.080385in}{9.647818in}}{\pgfqpoint{10.087600in}{9.653343in}}%
\pgfpathquadraticcurveto{\pgfqpoint{10.094816in}{9.658889in}}{\pgfqpoint{10.108093in}{9.658889in}}%
\pgfpathquadraticcurveto{\pgfqpoint{10.114649in}{9.658889in}}{\pgfqpoint{10.120442in}{9.657920in}}%
\pgfpathquadraticcurveto{\pgfqpoint{10.126256in}{9.656972in}}{\pgfqpoint{10.131142in}{9.655034in}}%
\pgfpathlineto{\pgfqpoint{10.131142in}{9.655034in}}%
\pgfpathclose%
\pgfpathmoveto{\pgfqpoint{10.156912in}{9.653219in}}%
\pgfpathlineto{\pgfqpoint{10.170498in}{9.653219in}}%
\pgfpathlineto{\pgfqpoint{10.170498in}{9.636871in}}%
\pgfpathlineto{\pgfqpoint{10.156912in}{9.636871in}}%
\pgfpathlineto{\pgfqpoint{10.156912in}{9.653219in}}%
\pgfpathlineto{\pgfqpoint{10.156912in}{9.653219in}}%
\pgfpathclose%
\pgfpathmoveto{\pgfqpoint{10.156912in}{9.601369in}}%
\pgfpathlineto{\pgfqpoint{10.170498in}{9.601369in}}%
\pgfpathlineto{\pgfqpoint{10.170498in}{9.590278in}}%
\pgfpathlineto{\pgfqpoint{10.159943in}{9.569661in}}%
\pgfpathlineto{\pgfqpoint{10.151635in}{9.569661in}}%
\pgfpathlineto{\pgfqpoint{10.156912in}{9.590278in}}%
\pgfpathlineto{\pgfqpoint{10.156912in}{9.601369in}}%
\pgfpathlineto{\pgfqpoint{10.156912in}{9.601369in}}%
\pgfpathclose%
\pgfpathmoveto{\pgfqpoint{10.277703in}{9.669857in}}%
\pgfpathlineto{\pgfqpoint{10.244841in}{9.618502in}}%
\pgfpathlineto{\pgfqpoint{10.277703in}{9.618502in}}%
\pgfpathlineto{\pgfqpoint{10.277703in}{9.669857in}}%
\pgfpathlineto{\pgfqpoint{10.277703in}{9.669857in}}%
\pgfpathclose%
\pgfpathmoveto{\pgfqpoint{10.274281in}{9.681196in}}%
\pgfpathlineto{\pgfqpoint{10.290650in}{9.681196in}}%
\pgfpathlineto{\pgfqpoint{10.290650in}{9.618502in}}%
\pgfpathlineto{\pgfqpoint{10.304381in}{9.618502in}}%
\pgfpathlineto{\pgfqpoint{10.304381in}{9.607678in}}%
\pgfpathlineto{\pgfqpoint{10.290650in}{9.607678in}}%
\pgfpathlineto{\pgfqpoint{10.290650in}{9.585000in}}%
\pgfpathlineto{\pgfqpoint{10.277703in}{9.585000in}}%
\pgfpathlineto{\pgfqpoint{10.277703in}{9.607678in}}%
\pgfpathlineto{\pgfqpoint{10.234285in}{9.607678in}}%
\pgfpathlineto{\pgfqpoint{10.234285in}{9.620233in}}%
\pgfpathlineto{\pgfqpoint{10.274281in}{9.681196in}}%
\pgfpathlineto{\pgfqpoint{10.274281in}{9.681196in}}%
\pgfpathclose%
\pgfpathmoveto{\pgfqpoint{10.365323in}{9.636871in}}%
\pgfpathquadraticcurveto{\pgfqpoint{10.374662in}{9.634871in}}{\pgfqpoint{10.379898in}{9.628542in}}%
\pgfpathquadraticcurveto{\pgfqpoint{10.385156in}{9.622233in}}{\pgfqpoint{10.385156in}{9.612956in}}%
\pgfpathquadraticcurveto{\pgfqpoint{10.385156in}{9.598730in}}{\pgfqpoint{10.375363in}{9.590917in}}%
\pgfpathquadraticcurveto{\pgfqpoint{10.365570in}{9.583124in}}{\pgfqpoint{10.347531in}{9.583124in}}%
\pgfpathquadraticcurveto{\pgfqpoint{10.341490in}{9.583124in}}{\pgfqpoint{10.335079in}{9.584320in}}%
\pgfpathquadraticcurveto{\pgfqpoint{10.328667in}{9.585515in}}{\pgfqpoint{10.321843in}{9.587907in}}%
\pgfpathlineto{\pgfqpoint{10.321843in}{9.600462in}}%
\pgfpathquadraticcurveto{\pgfqpoint{10.327244in}{9.597308in}}{\pgfqpoint{10.333677in}{9.595700in}}%
\pgfpathquadraticcurveto{\pgfqpoint{10.340130in}{9.594092in}}{\pgfqpoint{10.347160in}{9.594092in}}%
\pgfpathquadraticcurveto{\pgfqpoint{10.359385in}{9.594092in}}{\pgfqpoint{10.365797in}{9.598916in}}%
\pgfpathquadraticcurveto{\pgfqpoint{10.372209in}{9.603740in}}{\pgfqpoint{10.372209in}{9.612956in}}%
\pgfpathquadraticcurveto{\pgfqpoint{10.372209in}{9.621470in}}{\pgfqpoint{10.366250in}{9.626253in}}%
\pgfpathquadraticcurveto{\pgfqpoint{10.360292in}{9.631057in}}{\pgfqpoint{10.349675in}{9.631057in}}%
\pgfpathlineto{\pgfqpoint{10.338460in}{9.631057in}}%
\pgfpathlineto{\pgfqpoint{10.338460in}{9.641757in}}%
\pgfpathlineto{\pgfqpoint{10.350190in}{9.641757in}}%
\pgfpathquadraticcurveto{\pgfqpoint{10.359777in}{9.641757in}}{\pgfqpoint{10.364869in}{9.645591in}}%
\pgfpathquadraticcurveto{\pgfqpoint{10.369961in}{9.649426in}}{\pgfqpoint{10.369961in}{9.656642in}}%
\pgfpathquadraticcurveto{\pgfqpoint{10.369961in}{9.664043in}}{\pgfqpoint{10.364704in}{9.668001in}}%
\pgfpathquadraticcurveto{\pgfqpoint{10.359468in}{9.671980in}}{\pgfqpoint{10.349675in}{9.671980in}}%
\pgfpathquadraticcurveto{\pgfqpoint{10.344315in}{9.671980in}}{\pgfqpoint{10.338192in}{9.670805in}}%
\pgfpathquadraticcurveto{\pgfqpoint{10.332069in}{9.669651in}}{\pgfqpoint{10.324729in}{9.667218in}}%
\pgfpathlineto{\pgfqpoint{10.324729in}{9.678804in}}%
\pgfpathquadraticcurveto{\pgfqpoint{10.332151in}{9.680866in}}{\pgfqpoint{10.338625in}{9.681897in}}%
\pgfpathquadraticcurveto{\pgfqpoint{10.345098in}{9.682928in}}{\pgfqpoint{10.350829in}{9.682928in}}%
\pgfpathquadraticcurveto{\pgfqpoint{10.365653in}{9.682928in}}{\pgfqpoint{10.374270in}{9.676186in}}%
\pgfpathquadraticcurveto{\pgfqpoint{10.382908in}{9.669465in}}{\pgfqpoint{10.382908in}{9.658002in}}%
\pgfpathquadraticcurveto{\pgfqpoint{10.382908in}{9.650003in}}{\pgfqpoint{10.378332in}{9.644499in}}%
\pgfpathquadraticcurveto{\pgfqpoint{10.373755in}{9.638994in}}{\pgfqpoint{10.365323in}{9.636871in}}%
\pgfpathlineto{\pgfqpoint{10.365323in}{9.636871in}}%
\pgfpathclose%
\pgfpathmoveto{\pgfqpoint{10.502029in}{9.654395in}}%
\pgfpathlineto{\pgfqpoint{10.502029in}{9.643303in}}%
\pgfpathquadraticcurveto{\pgfqpoint{10.496999in}{9.646086in}}{\pgfqpoint{10.491948in}{9.647467in}}%
\pgfpathquadraticcurveto{\pgfqpoint{10.486897in}{9.648849in}}{\pgfqpoint{10.481743in}{9.648849in}}%
\pgfpathquadraticcurveto{\pgfqpoint{10.470198in}{9.648849in}}{\pgfqpoint{10.463807in}{9.641530in}}%
\pgfpathquadraticcurveto{\pgfqpoint{10.457436in}{9.634232in}}{\pgfqpoint{10.457436in}{9.621017in}}%
\pgfpathquadraticcurveto{\pgfqpoint{10.457436in}{9.607802in}}{\pgfqpoint{10.463807in}{9.600483in}}%
\pgfpathquadraticcurveto{\pgfqpoint{10.470198in}{9.593185in}}{\pgfqpoint{10.481743in}{9.593185in}}%
\pgfpathquadraticcurveto{\pgfqpoint{10.486897in}{9.593185in}}{\pgfqpoint{10.491948in}{9.594566in}}%
\pgfpathquadraticcurveto{\pgfqpoint{10.496999in}{9.595947in}}{\pgfqpoint{10.502029in}{9.598730in}}%
\pgfpathlineto{\pgfqpoint{10.502029in}{9.587763in}}%
\pgfpathquadraticcurveto{\pgfqpoint{10.497061in}{9.585454in}}{\pgfqpoint{10.491742in}{9.584299in}}%
\pgfpathquadraticcurveto{\pgfqpoint{10.486444in}{9.583124in}}{\pgfqpoint{10.480444in}{9.583124in}}%
\pgfpathquadraticcurveto{\pgfqpoint{10.464137in}{9.583124in}}{\pgfqpoint{10.454529in}{9.593370in}}%
\pgfpathquadraticcurveto{\pgfqpoint{10.444943in}{9.603617in}}{\pgfqpoint{10.444943in}{9.621017in}}%
\pgfpathquadraticcurveto{\pgfqpoint{10.444943in}{9.638664in}}{\pgfqpoint{10.454633in}{9.648766in}}%
\pgfpathquadraticcurveto{\pgfqpoint{10.464343in}{9.658889in}}{\pgfqpoint{10.481228in}{9.658889in}}%
\pgfpathquadraticcurveto{\pgfqpoint{10.486691in}{9.658889in}}{\pgfqpoint{10.491907in}{9.657755in}}%
\pgfpathquadraticcurveto{\pgfqpoint{10.497123in}{9.656642in}}{\pgfqpoint{10.502029in}{9.654395in}}%
\pgfpathlineto{\pgfqpoint{10.502029in}{9.654395in}}%
\pgfpathclose%
\pgfpathmoveto{\pgfqpoint{10.555446in}{9.621264in}}%
\pgfpathquadraticcurveto{\pgfqpoint{10.541077in}{9.621264in}}{\pgfqpoint{10.535531in}{9.617986in}}%
\pgfpathquadraticcurveto{\pgfqpoint{10.529985in}{9.614708in}}{\pgfqpoint{10.529985in}{9.606771in}}%
\pgfpathquadraticcurveto{\pgfqpoint{10.529985in}{9.600462in}}{\pgfqpoint{10.534150in}{9.596751in}}%
\pgfpathquadraticcurveto{\pgfqpoint{10.538314in}{9.593061in}}{\pgfqpoint{10.545447in}{9.593061in}}%
\pgfpathquadraticcurveto{\pgfqpoint{10.555323in}{9.593061in}}{\pgfqpoint{10.561281in}{9.600050in}}%
\pgfpathquadraticcurveto{\pgfqpoint{10.567239in}{9.607039in}}{\pgfqpoint{10.567239in}{9.618625in}}%
\pgfpathlineto{\pgfqpoint{10.567239in}{9.621264in}}%
\pgfpathlineto{\pgfqpoint{10.555446in}{9.621264in}}%
\pgfpathlineto{\pgfqpoint{10.555446in}{9.621264in}}%
\pgfpathclose%
\pgfpathmoveto{\pgfqpoint{10.579093in}{9.626171in}}%
\pgfpathlineto{\pgfqpoint{10.579093in}{9.585000in}}%
\pgfpathlineto{\pgfqpoint{10.567239in}{9.585000in}}%
\pgfpathlineto{\pgfqpoint{10.567239in}{9.595947in}}%
\pgfpathquadraticcurveto{\pgfqpoint{10.563177in}{9.589391in}}{\pgfqpoint{10.557116in}{9.586258in}}%
\pgfpathquadraticcurveto{\pgfqpoint{10.551055in}{9.583124in}}{\pgfqpoint{10.542293in}{9.583124in}}%
\pgfpathquadraticcurveto{\pgfqpoint{10.531222in}{9.583124in}}{\pgfqpoint{10.524666in}{9.589350in}}%
\pgfpathquadraticcurveto{\pgfqpoint{10.518131in}{9.595576in}}{\pgfqpoint{10.518131in}{9.606008in}}%
\pgfpathquadraticcurveto{\pgfqpoint{10.518131in}{9.618172in}}{\pgfqpoint{10.526274in}{9.624357in}}%
\pgfpathquadraticcurveto{\pgfqpoint{10.534438in}{9.630541in}}{\pgfqpoint{10.550602in}{9.630541in}}%
\pgfpathlineto{\pgfqpoint{10.567239in}{9.630541in}}%
\pgfpathlineto{\pgfqpoint{10.567239in}{9.631717in}}%
\pgfpathquadraticcurveto{\pgfqpoint{10.567239in}{9.639901in}}{\pgfqpoint{10.561858in}{9.644375in}}%
\pgfpathquadraticcurveto{\pgfqpoint{10.556477in}{9.648849in}}{\pgfqpoint{10.546746in}{9.648849in}}%
\pgfpathquadraticcurveto{\pgfqpoint{10.540561in}{9.648849in}}{\pgfqpoint{10.534686in}{9.647364in}}%
\pgfpathquadraticcurveto{\pgfqpoint{10.528831in}{9.645880in}}{\pgfqpoint{10.523429in}{9.642911in}}%
\pgfpathlineto{\pgfqpoint{10.523429in}{9.653879in}}%
\pgfpathquadraticcurveto{\pgfqpoint{10.529923in}{9.656394in}}{\pgfqpoint{10.536046in}{9.657631in}}%
\pgfpathquadraticcurveto{\pgfqpoint{10.542169in}{9.658889in}}{\pgfqpoint{10.547963in}{9.658889in}}%
\pgfpathquadraticcurveto{\pgfqpoint{10.563631in}{9.658889in}}{\pgfqpoint{10.571362in}{9.650766in}}%
\pgfpathquadraticcurveto{\pgfqpoint{10.579093in}{9.642664in}}{\pgfqpoint{10.579093in}{9.626171in}}%
\pgfpathlineto{\pgfqpoint{10.579093in}{9.626171in}}%
\pgfpathclose%
\pgfpathmoveto{\pgfqpoint{10.614966in}{9.595824in}}%
\pgfpathlineto{\pgfqpoint{10.614966in}{9.557560in}}%
\pgfpathlineto{\pgfqpoint{10.603049in}{9.557560in}}%
\pgfpathlineto{\pgfqpoint{10.603049in}{9.657157in}}%
\pgfpathlineto{\pgfqpoint{10.614966in}{9.657157in}}%
\pgfpathlineto{\pgfqpoint{10.614966in}{9.646210in}}%
\pgfpathquadraticcurveto{\pgfqpoint{10.618718in}{9.652642in}}{\pgfqpoint{10.624408in}{9.655755in}}%
\pgfpathquadraticcurveto{\pgfqpoint{10.630119in}{9.658889in}}{\pgfqpoint{10.638035in}{9.658889in}}%
\pgfpathquadraticcurveto{\pgfqpoint{10.651189in}{9.658889in}}{\pgfqpoint{10.659394in}{9.648457in}}%
\pgfpathquadraticcurveto{\pgfqpoint{10.667620in}{9.638025in}}{\pgfqpoint{10.667620in}{9.621017in}}%
\pgfpathquadraticcurveto{\pgfqpoint{10.667620in}{9.604008in}}{\pgfqpoint{10.659394in}{9.593556in}}%
\pgfpathquadraticcurveto{\pgfqpoint{10.651189in}{9.583124in}}{\pgfqpoint{10.638035in}{9.583124in}}%
\pgfpathquadraticcurveto{\pgfqpoint{10.630119in}{9.583124in}}{\pgfqpoint{10.624408in}{9.586258in}}%
\pgfpathquadraticcurveto{\pgfqpoint{10.618718in}{9.589391in}}{\pgfqpoint{10.614966in}{9.595824in}}%
\pgfpathlineto{\pgfqpoint{10.614966in}{9.595824in}}%
\pgfpathclose%
\pgfpathmoveto{\pgfqpoint{10.655312in}{9.621017in}}%
\pgfpathquadraticcurveto{\pgfqpoint{10.655312in}{9.634087in}}{\pgfqpoint{10.649931in}{9.641530in}}%
\pgfpathquadraticcurveto{\pgfqpoint{10.644550in}{9.648972in}}{\pgfqpoint{10.635149in}{9.648972in}}%
\pgfpathquadraticcurveto{\pgfqpoint{10.625727in}{9.648972in}}{\pgfqpoint{10.620347in}{9.641530in}}%
\pgfpathquadraticcurveto{\pgfqpoint{10.614966in}{9.634087in}}{\pgfqpoint{10.614966in}{9.621017in}}%
\pgfpathquadraticcurveto{\pgfqpoint{10.614966in}{9.607946in}}{\pgfqpoint{10.620347in}{9.600503in}}%
\pgfpathquadraticcurveto{\pgfqpoint{10.625727in}{9.593061in}}{\pgfqpoint{10.635149in}{9.593061in}}%
\pgfpathquadraticcurveto{\pgfqpoint{10.644550in}{9.593061in}}{\pgfqpoint{10.649931in}{9.600503in}}%
\pgfpathquadraticcurveto{\pgfqpoint{10.655312in}{9.607946in}}{\pgfqpoint{10.655312in}{9.621017in}}%
\pgfpathlineto{\pgfqpoint{10.655312in}{9.621017in}}%
\pgfpathclose%
\pgfpathmoveto{\pgfqpoint{10.698730in}{9.595824in}}%
\pgfpathlineto{\pgfqpoint{10.698730in}{9.557560in}}%
\pgfpathlineto{\pgfqpoint{10.686814in}{9.557560in}}%
\pgfpathlineto{\pgfqpoint{10.686814in}{9.657157in}}%
\pgfpathlineto{\pgfqpoint{10.698730in}{9.657157in}}%
\pgfpathlineto{\pgfqpoint{10.698730in}{9.646210in}}%
\pgfpathquadraticcurveto{\pgfqpoint{10.702482in}{9.652642in}}{\pgfqpoint{10.708172in}{9.655755in}}%
\pgfpathquadraticcurveto{\pgfqpoint{10.713883in}{9.658889in}}{\pgfqpoint{10.721799in}{9.658889in}}%
\pgfpathquadraticcurveto{\pgfqpoint{10.734953in}{9.658889in}}{\pgfqpoint{10.743158in}{9.648457in}}%
\pgfpathquadraticcurveto{\pgfqpoint{10.751384in}{9.638025in}}{\pgfqpoint{10.751384in}{9.621017in}}%
\pgfpathquadraticcurveto{\pgfqpoint{10.751384in}{9.604008in}}{\pgfqpoint{10.743158in}{9.593556in}}%
\pgfpathquadraticcurveto{\pgfqpoint{10.734953in}{9.583124in}}{\pgfqpoint{10.721799in}{9.583124in}}%
\pgfpathquadraticcurveto{\pgfqpoint{10.713883in}{9.583124in}}{\pgfqpoint{10.708172in}{9.586258in}}%
\pgfpathquadraticcurveto{\pgfqpoint{10.702482in}{9.589391in}}{\pgfqpoint{10.698730in}{9.595824in}}%
\pgfpathlineto{\pgfqpoint{10.698730in}{9.595824in}}%
\pgfpathclose%
\pgfpathmoveto{\pgfqpoint{10.739076in}{9.621017in}}%
\pgfpathquadraticcurveto{\pgfqpoint{10.739076in}{9.634087in}}{\pgfqpoint{10.733695in}{9.641530in}}%
\pgfpathquadraticcurveto{\pgfqpoint{10.728314in}{9.648972in}}{\pgfqpoint{10.718913in}{9.648972in}}%
\pgfpathquadraticcurveto{\pgfqpoint{10.709492in}{9.648972in}}{\pgfqpoint{10.704111in}{9.641530in}}%
\pgfpathquadraticcurveto{\pgfqpoint{10.698730in}{9.634087in}}{\pgfqpoint{10.698730in}{9.621017in}}%
\pgfpathquadraticcurveto{\pgfqpoint{10.698730in}{9.607946in}}{\pgfqpoint{10.704111in}{9.600503in}}%
\pgfpathquadraticcurveto{\pgfqpoint{10.709492in}{9.593061in}}{\pgfqpoint{10.718913in}{9.593061in}}%
\pgfpathquadraticcurveto{\pgfqpoint{10.728314in}{9.593061in}}{\pgfqpoint{10.733695in}{9.600503in}}%
\pgfpathquadraticcurveto{\pgfqpoint{10.739076in}{9.607946in}}{\pgfqpoint{10.739076in}{9.621017in}}%
\pgfpathlineto{\pgfqpoint{10.739076in}{9.621017in}}%
\pgfpathclose%
\pgfpathmoveto{\pgfqpoint{10.832756in}{9.624047in}}%
\pgfpathlineto{\pgfqpoint{10.832756in}{9.618254in}}%
\pgfpathlineto{\pgfqpoint{10.778247in}{9.618254in}}%
\pgfpathquadraticcurveto{\pgfqpoint{10.779030in}{9.606008in}}{\pgfqpoint{10.785628in}{9.599596in}}%
\pgfpathquadraticcurveto{\pgfqpoint{10.792225in}{9.593185in}}{\pgfqpoint{10.804017in}{9.593185in}}%
\pgfpathquadraticcurveto{\pgfqpoint{10.810841in}{9.593185in}}{\pgfqpoint{10.817253in}{9.594855in}}%
\pgfpathquadraticcurveto{\pgfqpoint{10.823665in}{9.596525in}}{\pgfqpoint{10.829994in}{9.599885in}}%
\pgfpathlineto{\pgfqpoint{10.829994in}{9.588670in}}%
\pgfpathquadraticcurveto{\pgfqpoint{10.823603in}{9.585969in}}{\pgfqpoint{10.816903in}{9.584546in}}%
\pgfpathquadraticcurveto{\pgfqpoint{10.810202in}{9.583124in}}{\pgfqpoint{10.803316in}{9.583124in}}%
\pgfpathquadraticcurveto{\pgfqpoint{10.786040in}{9.583124in}}{\pgfqpoint{10.775959in}{9.593164in}}%
\pgfpathquadraticcurveto{\pgfqpoint{10.765877in}{9.603225in}}{\pgfqpoint{10.765877in}{9.620378in}}%
\pgfpathquadraticcurveto{\pgfqpoint{10.765877in}{9.638087in}}{\pgfqpoint{10.775443in}{9.648478in}}%
\pgfpathquadraticcurveto{\pgfqpoint{10.785009in}{9.658889in}}{\pgfqpoint{10.801255in}{9.658889in}}%
\pgfpathquadraticcurveto{\pgfqpoint{10.815810in}{9.658889in}}{\pgfqpoint{10.824283in}{9.649508in}}%
\pgfpathquadraticcurveto{\pgfqpoint{10.832756in}{9.640149in}}{\pgfqpoint{10.832756in}{9.624047in}}%
\pgfpathlineto{\pgfqpoint{10.832756in}{9.624047in}}%
\pgfpathclose%
\pgfpathmoveto{\pgfqpoint{10.820902in}{9.627531in}}%
\pgfpathquadraticcurveto{\pgfqpoint{10.820778in}{9.637242in}}{\pgfqpoint{10.815459in}{9.643035in}}%
\pgfpathquadraticcurveto{\pgfqpoint{10.810140in}{9.648849in}}{\pgfqpoint{10.801378in}{9.648849in}}%
\pgfpathquadraticcurveto{\pgfqpoint{10.791462in}{9.648849in}}{\pgfqpoint{10.785504in}{9.643241in}}%
\pgfpathquadraticcurveto{\pgfqpoint{10.779546in}{9.637633in}}{\pgfqpoint{10.778639in}{9.627449in}}%
\pgfpathlineto{\pgfqpoint{10.820902in}{9.627531in}}%
\pgfpathlineto{\pgfqpoint{10.820902in}{9.627531in}}%
\pgfpathclose%
\pgfpathmoveto{\pgfqpoint{10.899698in}{9.646210in}}%
\pgfpathlineto{\pgfqpoint{10.899698in}{9.685257in}}%
\pgfpathlineto{\pgfqpoint{10.911552in}{9.685257in}}%
\pgfpathlineto{\pgfqpoint{10.911552in}{9.585000in}}%
\pgfpathlineto{\pgfqpoint{10.899698in}{9.585000in}}%
\pgfpathlineto{\pgfqpoint{10.899698in}{9.595824in}}%
\pgfpathquadraticcurveto{\pgfqpoint{10.895966in}{9.589391in}}{\pgfqpoint{10.890255in}{9.586258in}}%
\pgfpathquadraticcurveto{\pgfqpoint{10.884565in}{9.583124in}}{\pgfqpoint{10.876566in}{9.583124in}}%
\pgfpathquadraticcurveto{\pgfqpoint{10.863495in}{9.583124in}}{\pgfqpoint{10.855270in}{9.593556in}}%
\pgfpathquadraticcurveto{\pgfqpoint{10.847064in}{9.604008in}}{\pgfqpoint{10.847064in}{9.621017in}}%
\pgfpathquadraticcurveto{\pgfqpoint{10.847064in}{9.638025in}}{\pgfqpoint{10.855270in}{9.648457in}}%
\pgfpathquadraticcurveto{\pgfqpoint{10.863495in}{9.658889in}}{\pgfqpoint{10.876566in}{9.658889in}}%
\pgfpathquadraticcurveto{\pgfqpoint{10.884565in}{9.658889in}}{\pgfqpoint{10.890255in}{9.655755in}}%
\pgfpathquadraticcurveto{\pgfqpoint{10.895966in}{9.652642in}}{\pgfqpoint{10.899698in}{9.646210in}}%
\pgfpathlineto{\pgfqpoint{10.899698in}{9.646210in}}%
\pgfpathclose%
\pgfpathmoveto{\pgfqpoint{10.859310in}{9.621017in}}%
\pgfpathquadraticcurveto{\pgfqpoint{10.859310in}{9.607946in}}{\pgfqpoint{10.864691in}{9.600503in}}%
\pgfpathquadraticcurveto{\pgfqpoint{10.870072in}{9.593061in}}{\pgfqpoint{10.879473in}{9.593061in}}%
\pgfpathquadraticcurveto{\pgfqpoint{10.888874in}{9.593061in}}{\pgfqpoint{10.894276in}{9.600503in}}%
\pgfpathquadraticcurveto{\pgfqpoint{10.899698in}{9.607946in}}{\pgfqpoint{10.899698in}{9.621017in}}%
\pgfpathquadraticcurveto{\pgfqpoint{10.899698in}{9.634087in}}{\pgfqpoint{10.894276in}{9.641530in}}%
\pgfpathquadraticcurveto{\pgfqpoint{10.888874in}{9.648972in}}{\pgfqpoint{10.879473in}{9.648972in}}%
\pgfpathquadraticcurveto{\pgfqpoint{10.870072in}{9.648972in}}{\pgfqpoint{10.864691in}{9.641530in}}%
\pgfpathquadraticcurveto{\pgfqpoint{10.859310in}{9.634087in}}{\pgfqpoint{10.859310in}{9.621017in}}%
\pgfpathlineto{\pgfqpoint{10.859310in}{9.621017in}}%
\pgfpathclose%
\pgfpathmoveto{\pgfqpoint{11.006387in}{9.685113in}}%
\pgfpathquadraticcurveto{\pgfqpoint{10.997770in}{9.670310in}}{\pgfqpoint{10.993564in}{9.655796in}}%
\pgfpathquadraticcurveto{\pgfqpoint{10.989379in}{9.641303in}}{\pgfqpoint{10.989379in}{9.626418in}}%
\pgfpathquadraticcurveto{\pgfqpoint{10.989379in}{9.611554in}}{\pgfqpoint{10.993605in}{9.596957in}}%
\pgfpathquadraticcurveto{\pgfqpoint{10.997831in}{9.582361in}}{\pgfqpoint{11.006387in}{9.567600in}}%
\pgfpathlineto{\pgfqpoint{10.996079in}{9.567600in}}%
\pgfpathquadraticcurveto{\pgfqpoint{10.986431in}{9.582753in}}{\pgfqpoint{10.981627in}{9.597370in}}%
\pgfpathquadraticcurveto{\pgfqpoint{10.976823in}{9.611987in}}{\pgfqpoint{10.976823in}{9.626418in}}%
\pgfpathquadraticcurveto{\pgfqpoint{10.976823in}{9.640788in}}{\pgfqpoint{10.981586in}{9.655343in}}%
\pgfpathquadraticcurveto{\pgfqpoint{10.986369in}{9.669919in}}{\pgfqpoint{10.996079in}{9.685113in}}%
\pgfpathlineto{\pgfqpoint{11.006387in}{9.685113in}}%
\pgfpathlineto{\pgfqpoint{11.006387in}{9.685113in}}%
\pgfpathclose%
\pgfpathmoveto{\pgfqpoint{11.076874in}{9.646210in}}%
\pgfpathlineto{\pgfqpoint{11.076874in}{9.685257in}}%
\pgfpathlineto{\pgfqpoint{11.088729in}{9.685257in}}%
\pgfpathlineto{\pgfqpoint{11.088729in}{9.585000in}}%
\pgfpathlineto{\pgfqpoint{11.076874in}{9.585000in}}%
\pgfpathlineto{\pgfqpoint{11.076874in}{9.595824in}}%
\pgfpathquadraticcurveto{\pgfqpoint{11.073143in}{9.589391in}}{\pgfqpoint{11.067432in}{9.586258in}}%
\pgfpathquadraticcurveto{\pgfqpoint{11.061742in}{9.583124in}}{\pgfqpoint{11.053743in}{9.583124in}}%
\pgfpathquadraticcurveto{\pgfqpoint{11.040672in}{9.583124in}}{\pgfqpoint{11.032446in}{9.593556in}}%
\pgfpathquadraticcurveto{\pgfqpoint{11.024241in}{9.604008in}}{\pgfqpoint{11.024241in}{9.621017in}}%
\pgfpathquadraticcurveto{\pgfqpoint{11.024241in}{9.638025in}}{\pgfqpoint{11.032446in}{9.648457in}}%
\pgfpathquadraticcurveto{\pgfqpoint{11.040672in}{9.658889in}}{\pgfqpoint{11.053743in}{9.658889in}}%
\pgfpathquadraticcurveto{\pgfqpoint{11.061742in}{9.658889in}}{\pgfqpoint{11.067432in}{9.655755in}}%
\pgfpathquadraticcurveto{\pgfqpoint{11.073143in}{9.652642in}}{\pgfqpoint{11.076874in}{9.646210in}}%
\pgfpathlineto{\pgfqpoint{11.076874in}{9.646210in}}%
\pgfpathclose%
\pgfpathmoveto{\pgfqpoint{11.036487in}{9.621017in}}%
\pgfpathquadraticcurveto{\pgfqpoint{11.036487in}{9.607946in}}{\pgfqpoint{11.041868in}{9.600503in}}%
\pgfpathquadraticcurveto{\pgfqpoint{11.047249in}{9.593061in}}{\pgfqpoint{11.056650in}{9.593061in}}%
\pgfpathquadraticcurveto{\pgfqpoint{11.066051in}{9.593061in}}{\pgfqpoint{11.071452in}{9.600503in}}%
\pgfpathquadraticcurveto{\pgfqpoint{11.076874in}{9.607946in}}{\pgfqpoint{11.076874in}{9.621017in}}%
\pgfpathquadraticcurveto{\pgfqpoint{11.076874in}{9.634087in}}{\pgfqpoint{11.071452in}{9.641530in}}%
\pgfpathquadraticcurveto{\pgfqpoint{11.066051in}{9.648972in}}{\pgfqpoint{11.056650in}{9.648972in}}%
\pgfpathquadraticcurveto{\pgfqpoint{11.047249in}{9.648972in}}{\pgfqpoint{11.041868in}{9.641530in}}%
\pgfpathquadraticcurveto{\pgfqpoint{11.036487in}{9.634087in}}{\pgfqpoint{11.036487in}{9.621017in}}%
\pgfpathlineto{\pgfqpoint{11.036487in}{9.621017in}}%
\pgfpathclose%
\pgfpathmoveto{\pgfqpoint{11.113159in}{9.657157in}}%
\pgfpathlineto{\pgfqpoint{11.125013in}{9.657157in}}%
\pgfpathlineto{\pgfqpoint{11.125013in}{9.585000in}}%
\pgfpathlineto{\pgfqpoint{11.113159in}{9.585000in}}%
\pgfpathlineto{\pgfqpoint{11.113159in}{9.657157in}}%
\pgfpathlineto{\pgfqpoint{11.113159in}{9.657157in}}%
\pgfpathclose%
\pgfpathmoveto{\pgfqpoint{11.113159in}{9.685257in}}%
\pgfpathlineto{\pgfqpoint{11.125013in}{9.685257in}}%
\pgfpathlineto{\pgfqpoint{11.125013in}{9.670228in}}%
\pgfpathlineto{\pgfqpoint{11.113159in}{9.670228in}}%
\pgfpathlineto{\pgfqpoint{11.113159in}{9.685257in}}%
\pgfpathlineto{\pgfqpoint{11.113159in}{9.685257in}}%
\pgfpathclose%
\pgfpathmoveto{\pgfqpoint{11.182615in}{9.621264in}}%
\pgfpathquadraticcurveto{\pgfqpoint{11.168246in}{9.621264in}}{\pgfqpoint{11.162700in}{9.617986in}}%
\pgfpathquadraticcurveto{\pgfqpoint{11.157154in}{9.614708in}}{\pgfqpoint{11.157154in}{9.606771in}}%
\pgfpathquadraticcurveto{\pgfqpoint{11.157154in}{9.600462in}}{\pgfqpoint{11.161319in}{9.596751in}}%
\pgfpathquadraticcurveto{\pgfqpoint{11.165483in}{9.593061in}}{\pgfqpoint{11.172617in}{9.593061in}}%
\pgfpathquadraticcurveto{\pgfqpoint{11.182492in}{9.593061in}}{\pgfqpoint{11.188450in}{9.600050in}}%
\pgfpathquadraticcurveto{\pgfqpoint{11.194408in}{9.607039in}}{\pgfqpoint{11.194408in}{9.618625in}}%
\pgfpathlineto{\pgfqpoint{11.194408in}{9.621264in}}%
\pgfpathlineto{\pgfqpoint{11.182615in}{9.621264in}}%
\pgfpathlineto{\pgfqpoint{11.182615in}{9.621264in}}%
\pgfpathclose%
\pgfpathmoveto{\pgfqpoint{11.206262in}{9.626171in}}%
\pgfpathlineto{\pgfqpoint{11.206262in}{9.585000in}}%
\pgfpathlineto{\pgfqpoint{11.194408in}{9.585000in}}%
\pgfpathlineto{\pgfqpoint{11.194408in}{9.595947in}}%
\pgfpathquadraticcurveto{\pgfqpoint{11.190347in}{9.589391in}}{\pgfqpoint{11.184285in}{9.586258in}}%
\pgfpathquadraticcurveto{\pgfqpoint{11.178224in}{9.583124in}}{\pgfqpoint{11.169462in}{9.583124in}}%
\pgfpathquadraticcurveto{\pgfqpoint{11.158391in}{9.583124in}}{\pgfqpoint{11.151835in}{9.589350in}}%
\pgfpathquadraticcurveto{\pgfqpoint{11.145300in}{9.595576in}}{\pgfqpoint{11.145300in}{9.606008in}}%
\pgfpathquadraticcurveto{\pgfqpoint{11.145300in}{9.618172in}}{\pgfqpoint{11.153443in}{9.624357in}}%
\pgfpathquadraticcurveto{\pgfqpoint{11.161607in}{9.630541in}}{\pgfqpoint{11.177771in}{9.630541in}}%
\pgfpathlineto{\pgfqpoint{11.194408in}{9.630541in}}%
\pgfpathlineto{\pgfqpoint{11.194408in}{9.631717in}}%
\pgfpathquadraticcurveto{\pgfqpoint{11.194408in}{9.639901in}}{\pgfqpoint{11.189027in}{9.644375in}}%
\pgfpathquadraticcurveto{\pgfqpoint{11.183646in}{9.648849in}}{\pgfqpoint{11.173915in}{9.648849in}}%
\pgfpathquadraticcurveto{\pgfqpoint{11.167730in}{9.648849in}}{\pgfqpoint{11.161855in}{9.647364in}}%
\pgfpathquadraticcurveto{\pgfqpoint{11.156000in}{9.645880in}}{\pgfqpoint{11.150598in}{9.642911in}}%
\pgfpathlineto{\pgfqpoint{11.150598in}{9.653879in}}%
\pgfpathquadraticcurveto{\pgfqpoint{11.157092in}{9.656394in}}{\pgfqpoint{11.163215in}{9.657631in}}%
\pgfpathquadraticcurveto{\pgfqpoint{11.169339in}{9.658889in}}{\pgfqpoint{11.175132in}{9.658889in}}%
\pgfpathquadraticcurveto{\pgfqpoint{11.190800in}{9.658889in}}{\pgfqpoint{11.198531in}{9.650766in}}%
\pgfpathquadraticcurveto{\pgfqpoint{11.206262in}{9.642664in}}{\pgfqpoint{11.206262in}{9.626171in}}%
\pgfpathlineto{\pgfqpoint{11.206262in}{9.626171in}}%
\pgfpathclose%
\pgfpathmoveto{\pgfqpoint{11.286852in}{9.643303in}}%
\pgfpathquadraticcurveto{\pgfqpoint{11.291305in}{9.651302in}}{\pgfqpoint{11.297490in}{9.655095in}}%
\pgfpathquadraticcurveto{\pgfqpoint{11.303674in}{9.658889in}}{\pgfqpoint{11.312045in}{9.658889in}}%
\pgfpathquadraticcurveto{\pgfqpoint{11.323322in}{9.658889in}}{\pgfqpoint{11.329445in}{9.650993in}}%
\pgfpathquadraticcurveto{\pgfqpoint{11.335568in}{9.643117in}}{\pgfqpoint{11.335568in}{9.628562in}}%
\pgfpathlineto{\pgfqpoint{11.335568in}{9.585000in}}%
\pgfpathlineto{\pgfqpoint{11.323652in}{9.585000in}}%
\pgfpathlineto{\pgfqpoint{11.323652in}{9.628171in}}%
\pgfpathquadraticcurveto{\pgfqpoint{11.323652in}{9.638541in}}{\pgfqpoint{11.319961in}{9.643550in}}%
\pgfpathquadraticcurveto{\pgfqpoint{11.316292in}{9.648581in}}{\pgfqpoint{11.308767in}{9.648581in}}%
\pgfpathquadraticcurveto{\pgfqpoint{11.299551in}{9.648581in}}{\pgfqpoint{11.294191in}{9.642458in}}%
\pgfpathquadraticcurveto{\pgfqpoint{11.288851in}{9.636355in}}{\pgfqpoint{11.288851in}{9.625779in}}%
\pgfpathlineto{\pgfqpoint{11.288851in}{9.585000in}}%
\pgfpathlineto{\pgfqpoint{11.276935in}{9.585000in}}%
\pgfpathlineto{\pgfqpoint{11.276935in}{9.628171in}}%
\pgfpathquadraticcurveto{\pgfqpoint{11.276935in}{9.638602in}}{\pgfqpoint{11.273265in}{9.643592in}}%
\pgfpathquadraticcurveto{\pgfqpoint{11.269596in}{9.648581in}}{\pgfqpoint{11.261926in}{9.648581in}}%
\pgfpathquadraticcurveto{\pgfqpoint{11.252835in}{9.648581in}}{\pgfqpoint{11.247474in}{9.642437in}}%
\pgfpathquadraticcurveto{\pgfqpoint{11.242135in}{9.636293in}}{\pgfqpoint{11.242135in}{9.625779in}}%
\pgfpathlineto{\pgfqpoint{11.242135in}{9.585000in}}%
\pgfpathlineto{\pgfqpoint{11.230219in}{9.585000in}}%
\pgfpathlineto{\pgfqpoint{11.230219in}{9.657157in}}%
\pgfpathlineto{\pgfqpoint{11.242135in}{9.657157in}}%
\pgfpathlineto{\pgfqpoint{11.242135in}{9.645942in}}%
\pgfpathquadraticcurveto{\pgfqpoint{11.246196in}{9.652580in}}{\pgfqpoint{11.251866in}{9.655735in}}%
\pgfpathquadraticcurveto{\pgfqpoint{11.257535in}{9.658889in}}{\pgfqpoint{11.265328in}{9.658889in}}%
\pgfpathquadraticcurveto{\pgfqpoint{11.273204in}{9.658889in}}{\pgfqpoint{11.278708in}{9.654889in}}%
\pgfpathquadraticcurveto{\pgfqpoint{11.284213in}{9.650910in}}{\pgfqpoint{11.286852in}{9.643303in}}%
\pgfpathlineto{\pgfqpoint{11.286852in}{9.643303in}}%
\pgfpathclose%
\pgfpathmoveto{\pgfqpoint{11.387150in}{9.648849in}}%
\pgfpathquadraticcurveto{\pgfqpoint{11.377625in}{9.648849in}}{\pgfqpoint{11.372079in}{9.641406in}}%
\pgfpathquadraticcurveto{\pgfqpoint{11.366534in}{9.633964in}}{\pgfqpoint{11.366534in}{9.621017in}}%
\pgfpathquadraticcurveto{\pgfqpoint{11.366534in}{9.608070in}}{\pgfqpoint{11.372038in}{9.600627in}}%
\pgfpathquadraticcurveto{\pgfqpoint{11.377563in}{9.593185in}}{\pgfqpoint{11.387150in}{9.593185in}}%
\pgfpathquadraticcurveto{\pgfqpoint{11.396633in}{9.593185in}}{\pgfqpoint{11.402159in}{9.600648in}}%
\pgfpathquadraticcurveto{\pgfqpoint{11.407704in}{9.608132in}}{\pgfqpoint{11.407704in}{9.621017in}}%
\pgfpathquadraticcurveto{\pgfqpoint{11.407704in}{9.633840in}}{\pgfqpoint{11.402159in}{9.641344in}}%
\pgfpathquadraticcurveto{\pgfqpoint{11.396633in}{9.648849in}}{\pgfqpoint{11.387150in}{9.648849in}}%
\pgfpathlineto{\pgfqpoint{11.387150in}{9.648849in}}%
\pgfpathclose%
\pgfpathmoveto{\pgfqpoint{11.387150in}{9.658889in}}%
\pgfpathquadraticcurveto{\pgfqpoint{11.402612in}{9.658889in}}{\pgfqpoint{11.411436in}{9.648828in}}%
\pgfpathquadraticcurveto{\pgfqpoint{11.420280in}{9.638788in}}{\pgfqpoint{11.420280in}{9.621017in}}%
\pgfpathquadraticcurveto{\pgfqpoint{11.420280in}{9.603307in}}{\pgfqpoint{11.411436in}{9.593205in}}%
\pgfpathquadraticcurveto{\pgfqpoint{11.402612in}{9.583124in}}{\pgfqpoint{11.387150in}{9.583124in}}%
\pgfpathquadraticcurveto{\pgfqpoint{11.371626in}{9.583124in}}{\pgfqpoint{11.362823in}{9.593205in}}%
\pgfpathquadraticcurveto{\pgfqpoint{11.354040in}{9.603307in}}{\pgfqpoint{11.354040in}{9.621017in}}%
\pgfpathquadraticcurveto{\pgfqpoint{11.354040in}{9.638788in}}{\pgfqpoint{11.362823in}{9.648828in}}%
\pgfpathquadraticcurveto{\pgfqpoint{11.371626in}{9.658889in}}{\pgfqpoint{11.387150in}{9.658889in}}%
\pgfpathlineto{\pgfqpoint{11.387150in}{9.658889in}}%
\pgfpathclose%
\pgfpathmoveto{\pgfqpoint{11.499921in}{9.628562in}}%
\pgfpathlineto{\pgfqpoint{11.499921in}{9.585000in}}%
\pgfpathlineto{\pgfqpoint{11.488067in}{9.585000in}}%
\pgfpathlineto{\pgfqpoint{11.488067in}{9.628171in}}%
\pgfpathquadraticcurveto{\pgfqpoint{11.488067in}{9.638417in}}{\pgfqpoint{11.484067in}{9.643488in}}%
\pgfpathquadraticcurveto{\pgfqpoint{11.480068in}{9.648581in}}{\pgfqpoint{11.472089in}{9.648581in}}%
\pgfpathquadraticcurveto{\pgfqpoint{11.462482in}{9.648581in}}{\pgfqpoint{11.456936in}{9.642458in}}%
\pgfpathquadraticcurveto{\pgfqpoint{11.451390in}{9.636355in}}{\pgfqpoint{11.451390in}{9.625779in}}%
\pgfpathlineto{\pgfqpoint{11.451390in}{9.585000in}}%
\pgfpathlineto{\pgfqpoint{11.439474in}{9.585000in}}%
\pgfpathlineto{\pgfqpoint{11.439474in}{9.657157in}}%
\pgfpathlineto{\pgfqpoint{11.451390in}{9.657157in}}%
\pgfpathlineto{\pgfqpoint{11.451390in}{9.645942in}}%
\pgfpathquadraticcurveto{\pgfqpoint{11.455658in}{9.652457in}}{\pgfqpoint{11.461410in}{9.655673in}}%
\pgfpathquadraticcurveto{\pgfqpoint{11.467182in}{9.658889in}}{\pgfqpoint{11.474728in}{9.658889in}}%
\pgfpathquadraticcurveto{\pgfqpoint{11.487160in}{9.658889in}}{\pgfqpoint{11.493530in}{9.651199in}}%
\pgfpathquadraticcurveto{\pgfqpoint{11.499921in}{9.643509in}}{\pgfqpoint{11.499921in}{9.628562in}}%
\pgfpathlineto{\pgfqpoint{11.499921in}{9.628562in}}%
\pgfpathclose%
\pgfpathmoveto{\pgfqpoint{11.571027in}{9.646210in}}%
\pgfpathlineto{\pgfqpoint{11.571027in}{9.685257in}}%
\pgfpathlineto{\pgfqpoint{11.582881in}{9.685257in}}%
\pgfpathlineto{\pgfqpoint{11.582881in}{9.585000in}}%
\pgfpathlineto{\pgfqpoint{11.571027in}{9.585000in}}%
\pgfpathlineto{\pgfqpoint{11.571027in}{9.595824in}}%
\pgfpathquadraticcurveto{\pgfqpoint{11.567295in}{9.589391in}}{\pgfqpoint{11.561585in}{9.586258in}}%
\pgfpathquadraticcurveto{\pgfqpoint{11.555895in}{9.583124in}}{\pgfqpoint{11.547895in}{9.583124in}}%
\pgfpathquadraticcurveto{\pgfqpoint{11.534825in}{9.583124in}}{\pgfqpoint{11.526599in}{9.593556in}}%
\pgfpathquadraticcurveto{\pgfqpoint{11.518393in}{9.604008in}}{\pgfqpoint{11.518393in}{9.621017in}}%
\pgfpathquadraticcurveto{\pgfqpoint{11.518393in}{9.638025in}}{\pgfqpoint{11.526599in}{9.648457in}}%
\pgfpathquadraticcurveto{\pgfqpoint{11.534825in}{9.658889in}}{\pgfqpoint{11.547895in}{9.658889in}}%
\pgfpathquadraticcurveto{\pgfqpoint{11.555895in}{9.658889in}}{\pgfqpoint{11.561585in}{9.655755in}}%
\pgfpathquadraticcurveto{\pgfqpoint{11.567295in}{9.652642in}}{\pgfqpoint{11.571027in}{9.646210in}}%
\pgfpathlineto{\pgfqpoint{11.571027in}{9.646210in}}%
\pgfpathclose%
\pgfpathmoveto{\pgfqpoint{11.530640in}{9.621017in}}%
\pgfpathquadraticcurveto{\pgfqpoint{11.530640in}{9.607946in}}{\pgfqpoint{11.536020in}{9.600503in}}%
\pgfpathquadraticcurveto{\pgfqpoint{11.541401in}{9.593061in}}{\pgfqpoint{11.550802in}{9.593061in}}%
\pgfpathquadraticcurveto{\pgfqpoint{11.560203in}{9.593061in}}{\pgfqpoint{11.565605in}{9.600503in}}%
\pgfpathquadraticcurveto{\pgfqpoint{11.571027in}{9.607946in}}{\pgfqpoint{11.571027in}{9.621017in}}%
\pgfpathquadraticcurveto{\pgfqpoint{11.571027in}{9.634087in}}{\pgfqpoint{11.565605in}{9.641530in}}%
\pgfpathquadraticcurveto{\pgfqpoint{11.560203in}{9.648972in}}{\pgfqpoint{11.550802in}{9.648972in}}%
\pgfpathquadraticcurveto{\pgfqpoint{11.541401in}{9.648972in}}{\pgfqpoint{11.536020in}{9.641530in}}%
\pgfpathquadraticcurveto{\pgfqpoint{11.530640in}{9.634087in}}{\pgfqpoint{11.530640in}{9.621017in}}%
\pgfpathlineto{\pgfqpoint{11.530640in}{9.621017in}}%
\pgfpathclose%
\pgfpathmoveto{\pgfqpoint{11.653307in}{9.655034in}}%
\pgfpathlineto{\pgfqpoint{11.653307in}{9.643818in}}%
\pgfpathquadraticcurveto{\pgfqpoint{11.648297in}{9.646395in}}{\pgfqpoint{11.642875in}{9.647674in}}%
\pgfpathquadraticcurveto{\pgfqpoint{11.637473in}{9.648972in}}{\pgfqpoint{11.631659in}{9.648972in}}%
\pgfpathquadraticcurveto{\pgfqpoint{11.622836in}{9.648972in}}{\pgfqpoint{11.618424in}{9.646272in}}%
\pgfpathquadraticcurveto{\pgfqpoint{11.614012in}{9.643571in}}{\pgfqpoint{11.614012in}{9.638149in}}%
\pgfpathquadraticcurveto{\pgfqpoint{11.614012in}{9.634026in}}{\pgfqpoint{11.617166in}{9.631675in}}%
\pgfpathquadraticcurveto{\pgfqpoint{11.620321in}{9.629325in}}{\pgfqpoint{11.629866in}{9.627202in}}%
\pgfpathlineto{\pgfqpoint{11.633927in}{9.626294in}}%
\pgfpathquadraticcurveto{\pgfqpoint{11.646544in}{9.623594in}}{\pgfqpoint{11.651863in}{9.618666in}}%
\pgfpathquadraticcurveto{\pgfqpoint{11.657182in}{9.613739in}}{\pgfqpoint{11.657182in}{9.604915in}}%
\pgfpathquadraticcurveto{\pgfqpoint{11.657182in}{9.594855in}}{\pgfqpoint{11.649225in}{9.588979in}}%
\pgfpathquadraticcurveto{\pgfqpoint{11.641267in}{9.583124in}}{\pgfqpoint{11.627351in}{9.583124in}}%
\pgfpathquadraticcurveto{\pgfqpoint{11.621557in}{9.583124in}}{\pgfqpoint{11.615270in}{9.584258in}}%
\pgfpathquadraticcurveto{\pgfqpoint{11.608982in}{9.585392in}}{\pgfqpoint{11.602034in}{9.587639in}}%
\pgfpathlineto{\pgfqpoint{11.602034in}{9.599885in}}%
\pgfpathquadraticcurveto{\pgfqpoint{11.608610in}{9.596463in}}{\pgfqpoint{11.614981in}{9.594752in}}%
\pgfpathquadraticcurveto{\pgfqpoint{11.621351in}{9.593061in}}{\pgfqpoint{11.627619in}{9.593061in}}%
\pgfpathquadraticcurveto{\pgfqpoint{11.635989in}{9.593061in}}{\pgfqpoint{11.640483in}{9.595927in}}%
\pgfpathquadraticcurveto{\pgfqpoint{11.644998in}{9.598792in}}{\pgfqpoint{11.644998in}{9.604008in}}%
\pgfpathquadraticcurveto{\pgfqpoint{11.644998in}{9.608832in}}{\pgfqpoint{11.641741in}{9.611410in}}%
\pgfpathquadraticcurveto{\pgfqpoint{11.638504in}{9.613987in}}{\pgfqpoint{11.627474in}{9.616378in}}%
\pgfpathlineto{\pgfqpoint{11.623351in}{9.617347in}}%
\pgfpathquadraticcurveto{\pgfqpoint{11.612342in}{9.619656in}}{\pgfqpoint{11.607435in}{9.624460in}}%
\pgfpathquadraticcurveto{\pgfqpoint{11.602549in}{9.629263in}}{\pgfqpoint{11.602549in}{9.637633in}}%
\pgfpathquadraticcurveto{\pgfqpoint{11.602549in}{9.647818in}}{\pgfqpoint{11.609765in}{9.653343in}}%
\pgfpathquadraticcurveto{\pgfqpoint{11.616981in}{9.658889in}}{\pgfqpoint{11.630258in}{9.658889in}}%
\pgfpathquadraticcurveto{\pgfqpoint{11.636814in}{9.658889in}}{\pgfqpoint{11.642607in}{9.657920in}}%
\pgfpathquadraticcurveto{\pgfqpoint{11.648421in}{9.656972in}}{\pgfqpoint{11.653307in}{9.655034in}}%
\pgfpathlineto{\pgfqpoint{11.653307in}{9.655034in}}%
\pgfpathclose%
\pgfpathmoveto{\pgfqpoint{11.674191in}{9.685113in}}%
\pgfpathlineto{\pgfqpoint{11.684499in}{9.685113in}}%
\pgfpathquadraticcurveto{\pgfqpoint{11.694148in}{9.669919in}}{\pgfqpoint{11.698951in}{9.655343in}}%
\pgfpathquadraticcurveto{\pgfqpoint{11.703755in}{9.640788in}}{\pgfqpoint{11.703755in}{9.626418in}}%
\pgfpathquadraticcurveto{\pgfqpoint{11.703755in}{9.611987in}}{\pgfqpoint{11.698951in}{9.597370in}}%
\pgfpathquadraticcurveto{\pgfqpoint{11.694148in}{9.582753in}}{\pgfqpoint{11.684499in}{9.567600in}}%
\pgfpathlineto{\pgfqpoint{11.674191in}{9.567600in}}%
\pgfpathquadraticcurveto{\pgfqpoint{11.682747in}{9.582361in}}{\pgfqpoint{11.686973in}{9.596957in}}%
\pgfpathquadraticcurveto{\pgfqpoint{11.691199in}{9.611554in}}{\pgfqpoint{11.691199in}{9.626418in}}%
\pgfpathquadraticcurveto{\pgfqpoint{11.691199in}{9.641303in}}{\pgfqpoint{11.686973in}{9.655796in}}%
\pgfpathquadraticcurveto{\pgfqpoint{11.682747in}{9.670310in}}{\pgfqpoint{11.674191in}{9.685113in}}%
\pgfpathlineto{\pgfqpoint{11.674191in}{9.685113in}}%
\pgfpathclose%
\pgfusepath{fill}%
\end{pgfscope}%
\begin{pgfscope}%
\pgfsetbuttcap%
\pgfsetroundjoin%
\definecolor{currentfill}{rgb}{0.090196,0.168627,0.227451}%
\pgfsetfillcolor{currentfill}%
\pgfsetlinewidth{0.000000pt}%
\definecolor{currentstroke}{rgb}{0.090196,0.168627,0.227451}%
\pgfsetstrokecolor{currentstroke}%
\pgfsetdash{}{0pt}%
\pgfpathmoveto{\pgfqpoint{8.110709in}{9.925751in}}%
\pgfpathquadraticcurveto{\pgfqpoint{8.119368in}{9.925751in}}{\pgfqpoint{8.123804in}{9.929549in}}%
\pgfpathquadraticcurveto{\pgfqpoint{8.128270in}{9.933346in}}{\pgfqpoint{8.128270in}{9.940760in}}%
\pgfpathquadraticcurveto{\pgfqpoint{8.128270in}{9.948082in}}{\pgfqpoint{8.123804in}{9.951910in}}%
\pgfpathquadraticcurveto{\pgfqpoint{8.119368in}{9.955768in}}{\pgfqpoint{8.110709in}{9.955768in}}%
\pgfpathlineto{\pgfqpoint{8.090505in}{9.955768in}}%
\pgfpathlineto{\pgfqpoint{8.090505in}{9.925751in}}%
\pgfpathlineto{\pgfqpoint{8.110709in}{9.925751in}}%
\pgfpathlineto{\pgfqpoint{8.110709in}{9.925751in}}%
\pgfpathclose%
\pgfpathmoveto{\pgfqpoint{8.111955in}{9.863772in}}%
\pgfpathquadraticcurveto{\pgfqpoint{8.122953in}{9.863772in}}{\pgfqpoint{8.128513in}{9.868420in}}%
\pgfpathquadraticcurveto{\pgfqpoint{8.134073in}{9.873069in}}{\pgfqpoint{8.134073in}{9.882457in}}%
\pgfpathquadraticcurveto{\pgfqpoint{8.134073in}{9.891693in}}{\pgfqpoint{8.128574in}{9.896280in}}%
\pgfpathquadraticcurveto{\pgfqpoint{8.123074in}{9.900898in}}{\pgfqpoint{8.111955in}{9.900898in}}%
\pgfpathlineto{\pgfqpoint{8.090505in}{9.900898in}}%
\pgfpathlineto{\pgfqpoint{8.090505in}{9.863772in}}%
\pgfpathlineto{\pgfqpoint{8.111955in}{9.863772in}}%
\pgfpathlineto{\pgfqpoint{8.111955in}{9.863772in}}%
\pgfpathclose%
\pgfpathmoveto{\pgfqpoint{8.145952in}{9.914753in}}%
\pgfpathquadraticcurveto{\pgfqpoint{8.157710in}{9.911319in}}{\pgfqpoint{8.164151in}{9.902114in}}%
\pgfpathquadraticcurveto{\pgfqpoint{8.170622in}{9.892908in}}{\pgfqpoint{8.170622in}{9.879540in}}%
\pgfpathquadraticcurveto{\pgfqpoint{8.170622in}{9.859032in}}{\pgfqpoint{8.156768in}{9.848945in}}%
\pgfpathquadraticcurveto{\pgfqpoint{8.142914in}{9.838889in}}{\pgfqpoint{8.114598in}{9.838889in}}%
\pgfpathlineto{\pgfqpoint{8.053955in}{9.838889in}}%
\pgfpathlineto{\pgfqpoint{8.053955in}{9.980651in}}%
\pgfpathlineto{\pgfqpoint{8.108825in}{9.980651in}}%
\pgfpathquadraticcurveto{\pgfqpoint{8.138357in}{9.980651in}}{\pgfqpoint{8.151603in}{9.971719in}}%
\pgfpathquadraticcurveto{\pgfqpoint{8.164850in}{9.962786in}}{\pgfqpoint{8.164850in}{9.943129in}}%
\pgfpathquadraticcurveto{\pgfqpoint{8.164850in}{9.932799in}}{\pgfqpoint{8.159988in}{9.925538in}}%
\pgfpathquadraticcurveto{\pgfqpoint{8.155158in}{9.918277in}}{\pgfqpoint{8.145952in}{9.914753in}}%
\pgfpathlineto{\pgfqpoint{8.145952in}{9.914753in}}%
\pgfpathclose%
\pgfpathmoveto{\pgfqpoint{8.202159in}{9.980651in}}%
\pgfpathlineto{\pgfqpoint{8.238708in}{9.980651in}}%
\pgfpathlineto{\pgfqpoint{8.238708in}{9.852652in}}%
\pgfpathquadraticcurveto{\pgfqpoint{8.238708in}{9.826159in}}{\pgfqpoint{8.224307in}{9.813064in}}%
\pgfpathquadraticcurveto{\pgfqpoint{8.209936in}{9.799970in}}{\pgfqpoint{8.180770in}{9.799970in}}%
\pgfpathlineto{\pgfqpoint{8.173387in}{9.799970in}}%
\pgfpathlineto{\pgfqpoint{8.173387in}{9.827587in}}%
\pgfpathlineto{\pgfqpoint{8.179068in}{9.827587in}}%
\pgfpathquadraticcurveto{\pgfqpoint{8.190462in}{9.827587in}}{\pgfqpoint{8.196295in}{9.833967in}}%
\pgfpathquadraticcurveto{\pgfqpoint{8.202159in}{9.840317in}}{\pgfqpoint{8.202159in}{9.852652in}}%
\pgfpathlineto{\pgfqpoint{8.202159in}{9.980651in}}%
\pgfpathlineto{\pgfqpoint{8.202159in}{9.980651in}}%
\pgfpathclose%
\pgfpathmoveto{\pgfqpoint{8.373118in}{9.976185in}}%
\pgfpathlineto{\pgfqpoint{8.373118in}{9.946168in}}%
\pgfpathquadraticcurveto{\pgfqpoint{8.361451in}{9.951393in}}{\pgfqpoint{8.350331in}{9.954036in}}%
\pgfpathquadraticcurveto{\pgfqpoint{8.339242in}{9.956710in}}{\pgfqpoint{8.329368in}{9.956710in}}%
\pgfpathquadraticcurveto{\pgfqpoint{8.316273in}{9.956710in}}{\pgfqpoint{8.309984in}{9.953095in}}%
\pgfpathquadraticcurveto{\pgfqpoint{8.303725in}{9.949510in}}{\pgfqpoint{8.303725in}{9.941914in}}%
\pgfpathquadraticcurveto{\pgfqpoint{8.303725in}{9.936202in}}{\pgfqpoint{8.307949in}{9.933012in}}%
\pgfpathquadraticcurveto{\pgfqpoint{8.312172in}{9.929852in}}{\pgfqpoint{8.323291in}{9.927574in}}%
\pgfpathlineto{\pgfqpoint{8.338847in}{9.924444in}}%
\pgfpathquadraticcurveto{\pgfqpoint{8.362484in}{9.919674in}}{\pgfqpoint{8.372449in}{9.909983in}}%
\pgfpathquadraticcurveto{\pgfqpoint{8.382445in}{9.900321in}}{\pgfqpoint{8.382445in}{9.882457in}}%
\pgfpathquadraticcurveto{\pgfqpoint{8.382445in}{9.859032in}}{\pgfqpoint{8.368530in}{9.847578in}}%
\pgfpathquadraticcurveto{\pgfqpoint{8.354615in}{9.836124in}}{\pgfqpoint{8.326026in}{9.836124in}}%
\pgfpathquadraticcurveto{\pgfqpoint{8.312567in}{9.836124in}}{\pgfqpoint{8.298986in}{9.838707in}}%
\pgfpathquadraticcurveto{\pgfqpoint{8.285405in}{9.841259in}}{\pgfqpoint{8.271824in}{9.846302in}}%
\pgfpathlineto{\pgfqpoint{8.271824in}{9.877140in}}%
\pgfpathquadraticcurveto{\pgfqpoint{8.285405in}{9.869939in}}{\pgfqpoint{8.298074in}{9.866263in}}%
\pgfpathquadraticcurveto{\pgfqpoint{8.310744in}{9.862617in}}{\pgfqpoint{8.322532in}{9.862617in}}%
\pgfpathquadraticcurveto{\pgfqpoint{8.334502in}{9.862617in}}{\pgfqpoint{8.340852in}{9.866597in}}%
\pgfpathquadraticcurveto{\pgfqpoint{8.347202in}{9.870608in}}{\pgfqpoint{8.347202in}{9.878021in}}%
\pgfpathquadraticcurveto{\pgfqpoint{8.347202in}{9.884644in}}{\pgfqpoint{8.342888in}{9.888260in}}%
\pgfpathquadraticcurveto{\pgfqpoint{8.338574in}{9.891875in}}{\pgfqpoint{8.325661in}{9.894731in}}%
\pgfpathlineto{\pgfqpoint{8.311503in}{9.897860in}}%
\pgfpathquadraticcurveto{\pgfqpoint{8.290236in}{9.902418in}}{\pgfqpoint{8.280392in}{9.912383in}}%
\pgfpathquadraticcurveto{\pgfqpoint{8.270579in}{9.922348in}}{\pgfqpoint{8.270579in}{9.939240in}}%
\pgfpathquadraticcurveto{\pgfqpoint{8.270579in}{9.960417in}}{\pgfqpoint{8.284251in}{9.971810in}}%
\pgfpathquadraticcurveto{\pgfqpoint{8.297923in}{9.983203in}}{\pgfqpoint{8.323565in}{9.983203in}}%
\pgfpathquadraticcurveto{\pgfqpoint{8.335262in}{9.983203in}}{\pgfqpoint{8.347597in}{9.981441in}}%
\pgfpathquadraticcurveto{\pgfqpoint{8.359932in}{9.979679in}}{\pgfqpoint{8.373118in}{9.976185in}}%
\pgfpathlineto{\pgfqpoint{8.373118in}{9.976185in}}%
\pgfpathclose%
\pgfpathmoveto{\pgfqpoint{8.467302in}{9.945226in}}%
\pgfpathlineto{\pgfqpoint{8.501299in}{9.945226in}}%
\pgfpathlineto{\pgfqpoint{8.527792in}{9.871732in}}%
\pgfpathlineto{\pgfqpoint{8.554164in}{9.945226in}}%
\pgfpathlineto{\pgfqpoint{8.588252in}{9.945226in}}%
\pgfpathlineto{\pgfqpoint{8.546386in}{9.838889in}}%
\pgfpathlineto{\pgfqpoint{8.509077in}{9.838889in}}%
\pgfpathlineto{\pgfqpoint{8.467302in}{9.945226in}}%
\pgfpathlineto{\pgfqpoint{8.467302in}{9.945226in}}%
\pgfpathclose%
\pgfpathmoveto{\pgfqpoint{8.690518in}{9.941914in}}%
\pgfpathlineto{\pgfqpoint{8.690518in}{9.916089in}}%
\pgfpathquadraticcurveto{\pgfqpoint{8.679611in}{9.920647in}}{\pgfqpoint{8.669433in}{9.922925in}}%
\pgfpathquadraticcurveto{\pgfqpoint{8.659285in}{9.925204in}}{\pgfqpoint{8.650262in}{9.925204in}}%
\pgfpathquadraticcurveto{\pgfqpoint{8.640570in}{9.925204in}}{\pgfqpoint{8.635861in}{9.922773in}}%
\pgfpathquadraticcurveto{\pgfqpoint{8.631182in}{9.920343in}}{\pgfqpoint{8.631182in}{9.915330in}}%
\pgfpathquadraticcurveto{\pgfqpoint{8.631182in}{9.911228in}}{\pgfqpoint{8.634737in}{9.909041in}}%
\pgfpathquadraticcurveto{\pgfqpoint{8.638291in}{9.906884in}}{\pgfqpoint{8.647497in}{9.905820in}}%
\pgfpathlineto{\pgfqpoint{8.653482in}{9.904970in}}%
\pgfpathquadraticcurveto{\pgfqpoint{8.679611in}{9.901658in}}{\pgfqpoint{8.688604in}{9.894063in}}%
\pgfpathquadraticcurveto{\pgfqpoint{8.697627in}{9.886467in}}{\pgfqpoint{8.697627in}{9.870213in}}%
\pgfpathquadraticcurveto{\pgfqpoint{8.697627in}{9.853229in}}{\pgfqpoint{8.685080in}{9.844661in}}%
\pgfpathquadraticcurveto{\pgfqpoint{8.672562in}{9.836124in}}{\pgfqpoint{8.647710in}{9.836124in}}%
\pgfpathquadraticcurveto{\pgfqpoint{8.637167in}{9.836124in}}{\pgfqpoint{8.625896in}{9.837795in}}%
\pgfpathquadraticcurveto{\pgfqpoint{8.614654in}{9.839466in}}{\pgfqpoint{8.602775in}{9.842778in}}%
\pgfpathlineto{\pgfqpoint{8.602775in}{9.868602in}}%
\pgfpathquadraticcurveto{\pgfqpoint{8.612953in}{9.863681in}}{\pgfqpoint{8.623617in}{9.861189in}}%
\pgfpathquadraticcurveto{\pgfqpoint{8.634311in}{9.858728in}}{\pgfqpoint{8.645310in}{9.858728in}}%
\pgfpathquadraticcurveto{\pgfqpoint{8.655305in}{9.858728in}}{\pgfqpoint{8.660318in}{9.861463in}}%
\pgfpathquadraticcurveto{\pgfqpoint{8.665362in}{9.864227in}}{\pgfqpoint{8.665362in}{9.869666in}}%
\pgfpathquadraticcurveto{\pgfqpoint{8.665362in}{9.874223in}}{\pgfqpoint{8.661898in}{9.876441in}}%
\pgfpathquadraticcurveto{\pgfqpoint{8.658435in}{9.878659in}}{\pgfqpoint{8.648074in}{9.879905in}}%
\pgfpathlineto{\pgfqpoint{8.642089in}{9.880664in}}%
\pgfpathquadraticcurveto{\pgfqpoint{8.619394in}{9.883520in}}{\pgfqpoint{8.610279in}{9.891207in}}%
\pgfpathquadraticcurveto{\pgfqpoint{8.601165in}{9.898893in}}{\pgfqpoint{8.601165in}{9.914570in}}%
\pgfpathquadraticcurveto{\pgfqpoint{8.601165in}{9.931463in}}{\pgfqpoint{8.612740in}{9.939605in}}%
\pgfpathquadraticcurveto{\pgfqpoint{8.624346in}{9.947778in}}{\pgfqpoint{8.648257in}{9.947778in}}%
\pgfpathquadraticcurveto{\pgfqpoint{8.657675in}{9.947778in}}{\pgfqpoint{8.668005in}{9.946350in}}%
\pgfpathquadraticcurveto{\pgfqpoint{8.678365in}{9.944952in}}{\pgfqpoint{8.690518in}{9.941914in}}%
\pgfpathlineto{\pgfqpoint{8.690518in}{9.941914in}}%
\pgfpathclose%
\pgfpathmoveto{\pgfqpoint{8.775466in}{9.980651in}}%
\pgfpathlineto{\pgfqpoint{8.906108in}{9.980651in}}%
\pgfpathlineto{\pgfqpoint{8.906108in}{9.953003in}}%
\pgfpathlineto{\pgfqpoint{8.859107in}{9.953003in}}%
\pgfpathlineto{\pgfqpoint{8.859107in}{9.838889in}}%
\pgfpathlineto{\pgfqpoint{8.822558in}{9.838889in}}%
\pgfpathlineto{\pgfqpoint{8.822558in}{9.953003in}}%
\pgfpathlineto{\pgfqpoint{8.775466in}{9.953003in}}%
\pgfpathlineto{\pgfqpoint{8.775466in}{9.980651in}}%
\pgfpathlineto{\pgfqpoint{8.775466in}{9.980651in}}%
\pgfpathclose%
\pgfpathmoveto{\pgfqpoint{8.912975in}{9.980651in}}%
\pgfpathlineto{\pgfqpoint{8.948005in}{9.980651in}}%
\pgfpathlineto{\pgfqpoint{8.972493in}{9.877626in}}%
\pgfpathlineto{\pgfqpoint{8.996798in}{9.980651in}}%
\pgfpathlineto{\pgfqpoint{9.032011in}{9.980651in}}%
\pgfpathlineto{\pgfqpoint{9.056317in}{9.877626in}}%
\pgfpathlineto{\pgfqpoint{9.080835in}{9.980651in}}%
\pgfpathlineto{\pgfqpoint{9.115561in}{9.980651in}}%
\pgfpathlineto{\pgfqpoint{9.082141in}{9.838889in}}%
\pgfpathlineto{\pgfqpoint{9.040001in}{9.838889in}}%
\pgfpathlineto{\pgfqpoint{9.014268in}{9.946654in}}%
\pgfpathlineto{\pgfqpoint{8.988838in}{9.838889in}}%
\pgfpathlineto{\pgfqpoint{8.946668in}{9.838889in}}%
\pgfpathlineto{\pgfqpoint{8.912975in}{9.980651in}}%
\pgfpathlineto{\pgfqpoint{8.912975in}{9.980651in}}%
\pgfpathclose%
\pgfpathmoveto{\pgfqpoint{9.139502in}{9.980651in}}%
\pgfpathlineto{\pgfqpoint{9.238122in}{9.980651in}}%
\pgfpathlineto{\pgfqpoint{9.238122in}{9.953003in}}%
\pgfpathlineto{\pgfqpoint{9.176052in}{9.953003in}}%
\pgfpathlineto{\pgfqpoint{9.176052in}{9.926632in}}%
\pgfpathlineto{\pgfqpoint{9.234446in}{9.926632in}}%
\pgfpathlineto{\pgfqpoint{9.234446in}{9.898984in}}%
\pgfpathlineto{\pgfqpoint{9.176052in}{9.898984in}}%
\pgfpathlineto{\pgfqpoint{9.176052in}{9.838889in}}%
\pgfpathlineto{\pgfqpoint{9.139502in}{9.838889in}}%
\pgfpathlineto{\pgfqpoint{9.139502in}{9.980651in}}%
\pgfpathlineto{\pgfqpoint{9.139502in}{9.980651in}}%
\pgfpathclose%
\pgfpathmoveto{\pgfqpoint{9.272332in}{9.980651in}}%
\pgfpathlineto{\pgfqpoint{9.370952in}{9.980651in}}%
\pgfpathlineto{\pgfqpoint{9.370952in}{9.953003in}}%
\pgfpathlineto{\pgfqpoint{9.308882in}{9.953003in}}%
\pgfpathlineto{\pgfqpoint{9.308882in}{9.926632in}}%
\pgfpathlineto{\pgfqpoint{9.367276in}{9.926632in}}%
\pgfpathlineto{\pgfqpoint{9.367276in}{9.898984in}}%
\pgfpathlineto{\pgfqpoint{9.308882in}{9.898984in}}%
\pgfpathlineto{\pgfqpoint{9.308882in}{9.866506in}}%
\pgfpathlineto{\pgfqpoint{9.373048in}{9.866506in}}%
\pgfpathlineto{\pgfqpoint{9.373048in}{9.838889in}}%
\pgfpathlineto{\pgfqpoint{9.272332in}{9.838889in}}%
\pgfpathlineto{\pgfqpoint{9.272332in}{9.980651in}}%
\pgfpathlineto{\pgfqpoint{9.272332in}{9.980651in}}%
\pgfpathclose%
\pgfusepath{fill}%
\end{pgfscope}%
\begin{pgfscope}%
\pgfpathrectangle{\pgfqpoint{8.036091in}{4.050000in}}{\pgfqpoint{5.400000in}{5.400000in}}%
\pgfusepath{clip}%
\pgfsetbuttcap%
\pgfsetroundjoin%
\definecolor{currentfill}{rgb}{0.541176,0.380392,0.658824}%
\pgfsetfillcolor{currentfill}%
\pgfsetfillopacity{0.900000}%
\pgfsetlinewidth{0.803000pt}%
\definecolor{currentstroke}{rgb}{0.541176,0.380392,0.658824}%
\pgfsetstrokecolor{currentstroke}%
\pgfsetstrokeopacity{0.900000}%
\pgfsetdash{}{0pt}%
\pgfsys@defobject{currentmarker}{\pgfqpoint{-0.038665in}{-0.038665in}}{\pgfqpoint{0.038665in}{0.038665in}}{%
\pgfpathmoveto{\pgfqpoint{0.000000in}{-0.038665in}}%
\pgfpathcurveto{\pgfqpoint{0.010254in}{-0.038665in}}{\pgfqpoint{0.020090in}{-0.034591in}}{\pgfqpoint{0.027340in}{-0.027340in}}%
\pgfpathcurveto{\pgfqpoint{0.034591in}{-0.020090in}}{\pgfqpoint{0.038665in}{-0.010254in}}{\pgfqpoint{0.038665in}{0.000000in}}%
\pgfpathcurveto{\pgfqpoint{0.038665in}{0.010254in}}{\pgfqpoint{0.034591in}{0.020090in}}{\pgfqpoint{0.027340in}{0.027340in}}%
\pgfpathcurveto{\pgfqpoint{0.020090in}{0.034591in}}{\pgfqpoint{0.010254in}{0.038665in}}{\pgfqpoint{0.000000in}{0.038665in}}%
\pgfpathcurveto{\pgfqpoint{-0.010254in}{0.038665in}}{\pgfqpoint{-0.020090in}{0.034591in}}{\pgfqpoint{-0.027340in}{0.027340in}}%
\pgfpathcurveto{\pgfqpoint{-0.034591in}{0.020090in}}{\pgfqpoint{-0.038665in}{0.010254in}}{\pgfqpoint{-0.038665in}{0.000000in}}%
\pgfpathcurveto{\pgfqpoint{-0.038665in}{-0.010254in}}{\pgfqpoint{-0.034591in}{-0.020090in}}{\pgfqpoint{-0.027340in}{-0.027340in}}%
\pgfpathcurveto{\pgfqpoint{-0.020090in}{-0.034591in}}{\pgfqpoint{-0.010254in}{-0.038665in}}{\pgfqpoint{0.000000in}{-0.038665in}}%
\pgfpathlineto{\pgfqpoint{0.000000in}{-0.038665in}}%
\pgfpathclose%
\pgfusepath{stroke,fill}%
}%
\begin{pgfscope}%
\pgfsys@transformshift{10.876974in}{8.459270in}%
\pgfsys@useobject{currentmarker}{}%
\end{pgfscope}%
\end{pgfscope}%
\begin{pgfscope}%
\pgfsetbuttcap%
\pgfsetroundjoin%
\definecolor{currentfill}{rgb}{0.541176,0.380392,0.658824}%
\pgfsetfillcolor{currentfill}%
\pgfsetfillopacity{0.900000}%
\pgfsetlinewidth{0.803000pt}%
\definecolor{currentstroke}{rgb}{0.541176,0.380392,0.658824}%
\pgfsetstrokecolor{currentstroke}%
\pgfsetstrokeopacity{0.900000}%
\pgfsetdash{}{0pt}%
\pgfsys@defobject{currentmarker}{\pgfqpoint{-0.059738in}{-0.059738in}}{\pgfqpoint{0.059738in}{0.059738in}}{%
\pgfpathmoveto{\pgfqpoint{0.000000in}{-0.059738in}}%
\pgfpathlineto{\pgfqpoint{0.059738in}{0.000000in}}%
\pgfpathlineto{\pgfqpoint{0.000000in}{0.059738in}}%
\pgfpathlineto{\pgfqpoint{-0.059738in}{0.000000in}}%
\pgfpathlineto{\pgfqpoint{0.000000in}{-0.059738in}}%
\pgfpathclose%
\pgfusepath{stroke,fill}%
}%
\begin{pgfscope}%
\pgfsys@transformshift{12.372509in}{9.450000in}%
\pgfsys@useobject{currentmarker}{}%
\end{pgfscope}%
\end{pgfscope}%
\begin{pgfscope}%
\pgfpathrectangle{\pgfqpoint{8.036091in}{4.050000in}}{\pgfqpoint{5.400000in}{5.400000in}}%
\pgfusepath{clip}%
\pgfsetbuttcap%
\pgfsetroundjoin%
\definecolor{currentfill}{rgb}{0.321569,0.321569,0.321569}%
\pgfsetfillcolor{currentfill}%
\pgfsetfillopacity{0.900000}%
\pgfsetlinewidth{0.803000pt}%
\definecolor{currentstroke}{rgb}{0.321569,0.321569,0.321569}%
\pgfsetstrokecolor{currentstroke}%
\pgfsetstrokeopacity{0.900000}%
\pgfsetdash{}{0pt}%
\pgfsys@defobject{currentmarker}{\pgfqpoint{-0.038665in}{-0.038665in}}{\pgfqpoint{0.038665in}{0.038665in}}{%
\pgfpathmoveto{\pgfqpoint{0.000000in}{-0.038665in}}%
\pgfpathcurveto{\pgfqpoint{0.010254in}{-0.038665in}}{\pgfqpoint{0.020090in}{-0.034591in}}{\pgfqpoint{0.027340in}{-0.027340in}}%
\pgfpathcurveto{\pgfqpoint{0.034591in}{-0.020090in}}{\pgfqpoint{0.038665in}{-0.010254in}}{\pgfqpoint{0.038665in}{0.000000in}}%
\pgfpathcurveto{\pgfqpoint{0.038665in}{0.010254in}}{\pgfqpoint{0.034591in}{0.020090in}}{\pgfqpoint{0.027340in}{0.027340in}}%
\pgfpathcurveto{\pgfqpoint{0.020090in}{0.034591in}}{\pgfqpoint{0.010254in}{0.038665in}}{\pgfqpoint{0.000000in}{0.038665in}}%
\pgfpathcurveto{\pgfqpoint{-0.010254in}{0.038665in}}{\pgfqpoint{-0.020090in}{0.034591in}}{\pgfqpoint{-0.027340in}{0.027340in}}%
\pgfpathcurveto{\pgfqpoint{-0.034591in}{0.020090in}}{\pgfqpoint{-0.038665in}{0.010254in}}{\pgfqpoint{-0.038665in}{0.000000in}}%
\pgfpathcurveto{\pgfqpoint{-0.038665in}{-0.010254in}}{\pgfqpoint{-0.034591in}{-0.020090in}}{\pgfqpoint{-0.027340in}{-0.027340in}}%
\pgfpathcurveto{\pgfqpoint{-0.020090in}{-0.034591in}}{\pgfqpoint{-0.010254in}{-0.038665in}}{\pgfqpoint{0.000000in}{-0.038665in}}%
\pgfpathlineto{\pgfqpoint{0.000000in}{-0.038665in}}%
\pgfpathclose%
\pgfusepath{stroke,fill}%
}%
\begin{pgfscope}%
\pgfsys@transformshift{9.299179in}{5.017234in}%
\pgfsys@useobject{currentmarker}{}%
\end{pgfscope}%
\begin{pgfscope}%
\pgfsys@transformshift{9.484233in}{5.094659in}%
\pgfsys@useobject{currentmarker}{}%
\end{pgfscope}%
\begin{pgfscope}%
\pgfsys@transformshift{11.751002in}{8.291480in}%
\pgfsys@useobject{currentmarker}{}%
\end{pgfscope}%
\begin{pgfscope}%
\pgfsys@transformshift{11.563606in}{8.131639in}%
\pgfsys@useobject{currentmarker}{}%
\end{pgfscope}%
\begin{pgfscope}%
\pgfsys@transformshift{11.912496in}{7.437628in}%
\pgfsys@useobject{currentmarker}{}%
\end{pgfscope}%
\begin{pgfscope}%
\pgfsys@transformshift{12.053022in}{7.564920in}%
\pgfsys@useobject{currentmarker}{}%
\end{pgfscope}%
\begin{pgfscope}%
\pgfsys@transformshift{12.253052in}{7.810754in}%
\pgfsys@useobject{currentmarker}{}%
\end{pgfscope}%
\begin{pgfscope}%
\pgfsys@transformshift{11.926795in}{7.378902in}%
\pgfsys@useobject{currentmarker}{}%
\end{pgfscope}%
\begin{pgfscope}%
\pgfsys@transformshift{12.329971in}{8.872387in}%
\pgfsys@useobject{currentmarker}{}%
\end{pgfscope}%
\begin{pgfscope}%
\pgfsys@transformshift{12.010298in}{8.593765in}%
\pgfsys@useobject{currentmarker}{}%
\end{pgfscope}%
\begin{pgfscope}%
\pgfsys@transformshift{12.256767in}{8.820778in}%
\pgfsys@useobject{currentmarker}{}%
\end{pgfscope}%
\begin{pgfscope}%
\pgfsys@transformshift{9.554930in}{5.197704in}%
\pgfsys@useobject{currentmarker}{}%
\end{pgfscope}%
\begin{pgfscope}%
\pgfsys@transformshift{9.641620in}{5.279739in}%
\pgfsys@useobject{currentmarker}{}%
\end{pgfscope}%
\begin{pgfscope}%
\pgfsys@transformshift{9.240715in}{4.792708in}%
\pgfsys@useobject{currentmarker}{}%
\end{pgfscope}%
\begin{pgfscope}%
\pgfsys@transformshift{11.267728in}{7.841900in}%
\pgfsys@useobject{currentmarker}{}%
\end{pgfscope}%
\begin{pgfscope}%
\pgfsys@transformshift{11.576049in}{8.325511in}%
\pgfsys@useobject{currentmarker}{}%
\end{pgfscope}%
\begin{pgfscope}%
\pgfsys@transformshift{11.188625in}{7.740751in}%
\pgfsys@useobject{currentmarker}{}%
\end{pgfscope}%
\end{pgfscope}%
\begin{pgfscope}%
\pgfsetbuttcap%
\pgfsetroundjoin%
\definecolor{currentfill}{rgb}{0.321569,0.321569,0.321569}%
\pgfsetfillcolor{currentfill}%
\pgfsetfillopacity{0.900000}%
\pgfsetlinewidth{0.803000pt}%
\definecolor{currentstroke}{rgb}{0.321569,0.321569,0.321569}%
\pgfsetstrokecolor{currentstroke}%
\pgfsetstrokeopacity{0.900000}%
\pgfsetdash{}{0pt}%
\pgfsys@defobject{currentmarker}{\pgfqpoint{-0.059738in}{-0.059738in}}{\pgfqpoint{0.059738in}{0.059738in}}{%
\pgfpathmoveto{\pgfqpoint{0.000000in}{-0.059738in}}%
\pgfpathlineto{\pgfqpoint{0.059738in}{0.000000in}}%
\pgfpathlineto{\pgfqpoint{0.000000in}{0.059738in}}%
\pgfpathlineto{\pgfqpoint{-0.059738in}{0.000000in}}%
\pgfpathlineto{\pgfqpoint{0.000000in}{-0.059738in}}%
\pgfpathclose%
\pgfusepath{stroke,fill}%
}%
\begin{pgfscope}%
\pgfsys@transformshift{8.036091in}{4.050000in}%
\pgfsys@useobject{currentmarker}{}%
\end{pgfscope}%
\end{pgfscope}%
\begin{pgfscope}%
\pgfpathrectangle{\pgfqpoint{8.036091in}{4.050000in}}{\pgfqpoint{5.400000in}{5.400000in}}%
\pgfusepath{clip}%
\pgfsetbuttcap%
\pgfsetroundjoin%
\definecolor{currentfill}{rgb}{0.200000,0.133333,0.533333}%
\pgfsetfillcolor{currentfill}%
\pgfsetfillopacity{0.900000}%
\pgfsetlinewidth{0.803000pt}%
\definecolor{currentstroke}{rgb}{0.200000,0.133333,0.533333}%
\pgfsetstrokecolor{currentstroke}%
\pgfsetstrokeopacity{0.900000}%
\pgfsetdash{}{0pt}%
\pgfsys@defobject{currentmarker}{\pgfqpoint{-0.038665in}{-0.038665in}}{\pgfqpoint{0.038665in}{0.038665in}}{%
\pgfpathmoveto{\pgfqpoint{0.000000in}{-0.038665in}}%
\pgfpathcurveto{\pgfqpoint{0.010254in}{-0.038665in}}{\pgfqpoint{0.020090in}{-0.034591in}}{\pgfqpoint{0.027340in}{-0.027340in}}%
\pgfpathcurveto{\pgfqpoint{0.034591in}{-0.020090in}}{\pgfqpoint{0.038665in}{-0.010254in}}{\pgfqpoint{0.038665in}{0.000000in}}%
\pgfpathcurveto{\pgfqpoint{0.038665in}{0.010254in}}{\pgfqpoint{0.034591in}{0.020090in}}{\pgfqpoint{0.027340in}{0.027340in}}%
\pgfpathcurveto{\pgfqpoint{0.020090in}{0.034591in}}{\pgfqpoint{0.010254in}{0.038665in}}{\pgfqpoint{0.000000in}{0.038665in}}%
\pgfpathcurveto{\pgfqpoint{-0.010254in}{0.038665in}}{\pgfqpoint{-0.020090in}{0.034591in}}{\pgfqpoint{-0.027340in}{0.027340in}}%
\pgfpathcurveto{\pgfqpoint{-0.034591in}{0.020090in}}{\pgfqpoint{-0.038665in}{0.010254in}}{\pgfqpoint{-0.038665in}{0.000000in}}%
\pgfpathcurveto{\pgfqpoint{-0.038665in}{-0.010254in}}{\pgfqpoint{-0.034591in}{-0.020090in}}{\pgfqpoint{-0.027340in}{-0.027340in}}%
\pgfpathcurveto{\pgfqpoint{-0.020090in}{-0.034591in}}{\pgfqpoint{-0.010254in}{-0.038665in}}{\pgfqpoint{0.000000in}{-0.038665in}}%
\pgfpathlineto{\pgfqpoint{0.000000in}{-0.038665in}}%
\pgfpathclose%
\pgfusepath{stroke,fill}%
}%
\begin{pgfscope}%
\pgfsys@transformshift{9.212616in}{4.588697in}%
\pgfsys@useobject{currentmarker}{}%
\end{pgfscope}%
\begin{pgfscope}%
\pgfsys@transformshift{9.792876in}{5.084548in}%
\pgfsys@useobject{currentmarker}{}%
\end{pgfscope}%
\begin{pgfscope}%
\pgfsys@transformshift{9.033007in}{4.315329in}%
\pgfsys@useobject{currentmarker}{}%
\end{pgfscope}%
\begin{pgfscope}%
\pgfsys@transformshift{9.684518in}{4.951369in}%
\pgfsys@useobject{currentmarker}{}%
\end{pgfscope}%
\end{pgfscope}%
\begin{pgfscope}%
\pgfsetbuttcap%
\pgfsetroundjoin%
\definecolor{currentfill}{rgb}{0.176471,0.713725,0.643137}%
\pgfsetfillcolor{currentfill}%
\pgfsetfillopacity{0.900000}%
\pgfsetlinewidth{0.803000pt}%
\definecolor{currentstroke}{rgb}{0.176471,0.713725,0.643137}%
\pgfsetstrokecolor{currentstroke}%
\pgfsetstrokeopacity{0.900000}%
\pgfsetdash{}{0pt}%
\pgfsys@defobject{currentmarker}{\pgfqpoint{-0.059738in}{-0.059738in}}{\pgfqpoint{0.059738in}{0.059738in}}{%
\pgfpathmoveto{\pgfqpoint{0.000000in}{-0.059738in}}%
\pgfpathlineto{\pgfqpoint{0.059738in}{0.000000in}}%
\pgfpathlineto{\pgfqpoint{0.000000in}{0.059738in}}%
\pgfpathlineto{\pgfqpoint{-0.059738in}{0.000000in}}%
\pgfpathlineto{\pgfqpoint{0.000000in}{-0.059738in}}%
\pgfpathclose%
\pgfusepath{stroke,fill}%
}%
\begin{pgfscope}%
\pgfsys@transformshift{11.479441in}{9.450000in}%
\pgfsys@useobject{currentmarker}{}%
\end{pgfscope}%
\end{pgfscope}%
\begin{pgfscope}%
\pgfpathrectangle{\pgfqpoint{8.036091in}{4.050000in}}{\pgfqpoint{5.400000in}{5.400000in}}%
\pgfusepath{clip}%
\pgfsetbuttcap%
\pgfsetroundjoin%
\definecolor{currentfill}{rgb}{0.000000,0.447059,0.698039}%
\pgfsetfillcolor{currentfill}%
\pgfsetfillopacity{0.900000}%
\pgfsetlinewidth{0.803000pt}%
\definecolor{currentstroke}{rgb}{0.000000,0.447059,0.698039}%
\pgfsetstrokecolor{currentstroke}%
\pgfsetstrokeopacity{0.900000}%
\pgfsetdash{}{0pt}%
\pgfsys@defobject{currentmarker}{\pgfqpoint{-0.038665in}{-0.038665in}}{\pgfqpoint{0.038665in}{0.038665in}}{%
\pgfpathmoveto{\pgfqpoint{0.000000in}{-0.038665in}}%
\pgfpathcurveto{\pgfqpoint{0.010254in}{-0.038665in}}{\pgfqpoint{0.020090in}{-0.034591in}}{\pgfqpoint{0.027340in}{-0.027340in}}%
\pgfpathcurveto{\pgfqpoint{0.034591in}{-0.020090in}}{\pgfqpoint{0.038665in}{-0.010254in}}{\pgfqpoint{0.038665in}{0.000000in}}%
\pgfpathcurveto{\pgfqpoint{0.038665in}{0.010254in}}{\pgfqpoint{0.034591in}{0.020090in}}{\pgfqpoint{0.027340in}{0.027340in}}%
\pgfpathcurveto{\pgfqpoint{0.020090in}{0.034591in}}{\pgfqpoint{0.010254in}{0.038665in}}{\pgfqpoint{0.000000in}{0.038665in}}%
\pgfpathcurveto{\pgfqpoint{-0.010254in}{0.038665in}}{\pgfqpoint{-0.020090in}{0.034591in}}{\pgfqpoint{-0.027340in}{0.027340in}}%
\pgfpathcurveto{\pgfqpoint{-0.034591in}{0.020090in}}{\pgfqpoint{-0.038665in}{0.010254in}}{\pgfqpoint{-0.038665in}{0.000000in}}%
\pgfpathcurveto{\pgfqpoint{-0.038665in}{-0.010254in}}{\pgfqpoint{-0.034591in}{-0.020090in}}{\pgfqpoint{-0.027340in}{-0.027340in}}%
\pgfpathcurveto{\pgfqpoint{-0.020090in}{-0.034591in}}{\pgfqpoint{-0.010254in}{-0.038665in}}{\pgfqpoint{0.000000in}{-0.038665in}}%
\pgfpathlineto{\pgfqpoint{0.000000in}{-0.038665in}}%
\pgfpathclose%
\pgfusepath{stroke,fill}%
}%
\begin{pgfscope}%
\pgfsys@transformshift{10.754255in}{8.949705in}%
\pgfsys@useobject{currentmarker}{}%
\end{pgfscope}%
\end{pgfscope}%
\begin{pgfscope}%
\pgfsetbuttcap%
\pgfsetroundjoin%
\definecolor{currentfill}{rgb}{0.000000,0.447059,0.698039}%
\pgfsetfillcolor{currentfill}%
\pgfsetfillopacity{0.900000}%
\pgfsetlinewidth{0.803000pt}%
\definecolor{currentstroke}{rgb}{0.000000,0.447059,0.698039}%
\pgfsetstrokecolor{currentstroke}%
\pgfsetstrokeopacity{0.900000}%
\pgfsetdash{}{0pt}%
\pgfsys@defobject{currentmarker}{\pgfqpoint{-0.059738in}{-0.059738in}}{\pgfqpoint{0.059738in}{0.059738in}}{%
\pgfpathmoveto{\pgfqpoint{0.000000in}{-0.059738in}}%
\pgfpathlineto{\pgfqpoint{0.059738in}{0.000000in}}%
\pgfpathlineto{\pgfqpoint{0.000000in}{0.059738in}}%
\pgfpathlineto{\pgfqpoint{-0.059738in}{0.000000in}}%
\pgfpathlineto{\pgfqpoint{0.000000in}{-0.059738in}}%
\pgfpathclose%
\pgfusepath{stroke,fill}%
}%
\begin{pgfscope}%
\pgfsys@transformshift{9.261824in}{4.050000in}%
\pgfsys@useobject{currentmarker}{}%
\end{pgfscope}%
\begin{pgfscope}%
\pgfsys@transformshift{9.282566in}{4.050000in}%
\pgfsys@useobject{currentmarker}{}%
\end{pgfscope}%
\begin{pgfscope}%
\pgfsys@transformshift{9.719338in}{4.050000in}%
\pgfsys@useobject{currentmarker}{}%
\end{pgfscope}%
\begin{pgfscope}%
\pgfsys@transformshift{9.779679in}{4.050000in}%
\pgfsys@useobject{currentmarker}{}%
\end{pgfscope}%
\end{pgfscope}%
\begin{pgfscope}%
\pgfsetbuttcap%
\pgfsetroundjoin%
\definecolor{currentfill}{rgb}{0.337255,0.705882,0.913725}%
\pgfsetfillcolor{currentfill}%
\pgfsetfillopacity{0.900000}%
\pgfsetlinewidth{0.803000pt}%
\definecolor{currentstroke}{rgb}{0.337255,0.705882,0.913725}%
\pgfsetstrokecolor{currentstroke}%
\pgfsetstrokeopacity{0.900000}%
\pgfsetdash{}{0pt}%
\pgfsys@defobject{currentmarker}{\pgfqpoint{-0.059738in}{-0.059738in}}{\pgfqpoint{0.059738in}{0.059738in}}{%
\pgfpathmoveto{\pgfqpoint{0.000000in}{-0.059738in}}%
\pgfpathlineto{\pgfqpoint{0.059738in}{0.000000in}}%
\pgfpathlineto{\pgfqpoint{0.000000in}{0.059738in}}%
\pgfpathlineto{\pgfqpoint{-0.059738in}{0.000000in}}%
\pgfpathlineto{\pgfqpoint{0.000000in}{-0.059738in}}%
\pgfpathclose%
\pgfusepath{stroke,fill}%
}%
\begin{pgfscope}%
\pgfsys@transformshift{8.625030in}{4.050000in}%
\pgfsys@useobject{currentmarker}{}%
\end{pgfscope}%
\end{pgfscope}%
\begin{pgfscope}%
\pgfpathrectangle{\pgfqpoint{8.036091in}{4.050000in}}{\pgfqpoint{5.400000in}{5.400000in}}%
\pgfusepath{clip}%
\pgfsetbuttcap%
\pgfsetroundjoin%
\definecolor{currentfill}{rgb}{0.662745,0.368627,0.000000}%
\pgfsetfillcolor{currentfill}%
\pgfsetfillopacity{0.900000}%
\pgfsetlinewidth{0.803000pt}%
\definecolor{currentstroke}{rgb}{0.662745,0.368627,0.000000}%
\pgfsetstrokecolor{currentstroke}%
\pgfsetstrokeopacity{0.900000}%
\pgfsetdash{}{0pt}%
\pgfsys@defobject{currentmarker}{\pgfqpoint{-0.038665in}{-0.038665in}}{\pgfqpoint{0.038665in}{0.038665in}}{%
\pgfpathmoveto{\pgfqpoint{0.000000in}{-0.038665in}}%
\pgfpathcurveto{\pgfqpoint{0.010254in}{-0.038665in}}{\pgfqpoint{0.020090in}{-0.034591in}}{\pgfqpoint{0.027340in}{-0.027340in}}%
\pgfpathcurveto{\pgfqpoint{0.034591in}{-0.020090in}}{\pgfqpoint{0.038665in}{-0.010254in}}{\pgfqpoint{0.038665in}{0.000000in}}%
\pgfpathcurveto{\pgfqpoint{0.038665in}{0.010254in}}{\pgfqpoint{0.034591in}{0.020090in}}{\pgfqpoint{0.027340in}{0.027340in}}%
\pgfpathcurveto{\pgfqpoint{0.020090in}{0.034591in}}{\pgfqpoint{0.010254in}{0.038665in}}{\pgfqpoint{0.000000in}{0.038665in}}%
\pgfpathcurveto{\pgfqpoint{-0.010254in}{0.038665in}}{\pgfqpoint{-0.020090in}{0.034591in}}{\pgfqpoint{-0.027340in}{0.027340in}}%
\pgfpathcurveto{\pgfqpoint{-0.034591in}{0.020090in}}{\pgfqpoint{-0.038665in}{0.010254in}}{\pgfqpoint{-0.038665in}{0.000000in}}%
\pgfpathcurveto{\pgfqpoint{-0.038665in}{-0.010254in}}{\pgfqpoint{-0.034591in}{-0.020090in}}{\pgfqpoint{-0.027340in}{-0.027340in}}%
\pgfpathcurveto{\pgfqpoint{-0.020090in}{-0.034591in}}{\pgfqpoint{-0.010254in}{-0.038665in}}{\pgfqpoint{0.000000in}{-0.038665in}}%
\pgfpathlineto{\pgfqpoint{0.000000in}{-0.038665in}}%
\pgfpathclose%
\pgfusepath{stroke,fill}%
}%
\begin{pgfscope}%
\pgfsys@transformshift{11.877332in}{8.679035in}%
\pgfsys@useobject{currentmarker}{}%
\end{pgfscope}%
\end{pgfscope}%
\begin{pgfscope}%
\pgfpathrectangle{\pgfqpoint{8.036091in}{4.050000in}}{\pgfqpoint{5.400000in}{5.400000in}}%
\pgfusepath{clip}%
\pgfsetbuttcap%
\pgfsetroundjoin%
\definecolor{currentfill}{rgb}{0.901961,0.403922,0.619608}%
\pgfsetfillcolor{currentfill}%
\pgfsetfillopacity{0.900000}%
\pgfsetlinewidth{0.803000pt}%
\definecolor{currentstroke}{rgb}{0.901961,0.403922,0.619608}%
\pgfsetstrokecolor{currentstroke}%
\pgfsetstrokeopacity{0.900000}%
\pgfsetdash{}{0pt}%
\pgfsys@defobject{currentmarker}{\pgfqpoint{-0.038665in}{-0.038665in}}{\pgfqpoint{0.038665in}{0.038665in}}{%
\pgfpathmoveto{\pgfqpoint{0.000000in}{-0.038665in}}%
\pgfpathcurveto{\pgfqpoint{0.010254in}{-0.038665in}}{\pgfqpoint{0.020090in}{-0.034591in}}{\pgfqpoint{0.027340in}{-0.027340in}}%
\pgfpathcurveto{\pgfqpoint{0.034591in}{-0.020090in}}{\pgfqpoint{0.038665in}{-0.010254in}}{\pgfqpoint{0.038665in}{0.000000in}}%
\pgfpathcurveto{\pgfqpoint{0.038665in}{0.010254in}}{\pgfqpoint{0.034591in}{0.020090in}}{\pgfqpoint{0.027340in}{0.027340in}}%
\pgfpathcurveto{\pgfqpoint{0.020090in}{0.034591in}}{\pgfqpoint{0.010254in}{0.038665in}}{\pgfqpoint{0.000000in}{0.038665in}}%
\pgfpathcurveto{\pgfqpoint{-0.010254in}{0.038665in}}{\pgfqpoint{-0.020090in}{0.034591in}}{\pgfqpoint{-0.027340in}{0.027340in}}%
\pgfpathcurveto{\pgfqpoint{-0.034591in}{0.020090in}}{\pgfqpoint{-0.038665in}{0.010254in}}{\pgfqpoint{-0.038665in}{0.000000in}}%
\pgfpathcurveto{\pgfqpoint{-0.038665in}{-0.010254in}}{\pgfqpoint{-0.034591in}{-0.020090in}}{\pgfqpoint{-0.027340in}{-0.027340in}}%
\pgfpathcurveto{\pgfqpoint{-0.020090in}{-0.034591in}}{\pgfqpoint{-0.010254in}{-0.038665in}}{\pgfqpoint{0.000000in}{-0.038665in}}%
\pgfpathlineto{\pgfqpoint{0.000000in}{-0.038665in}}%
\pgfpathclose%
\pgfusepath{stroke,fill}%
}%
\begin{pgfscope}%
\pgfsys@transformshift{11.932279in}{6.603856in}%
\pgfsys@useobject{currentmarker}{}%
\end{pgfscope}%
\end{pgfscope}%
\begin{pgfscope}%
\pgfsetbuttcap%
\pgfsetroundjoin%
\definecolor{currentfill}{rgb}{0.941176,0.584314,0.337255}%
\pgfsetfillcolor{currentfill}%
\pgfsetfillopacity{0.900000}%
\pgfsetlinewidth{0.803000pt}%
\definecolor{currentstroke}{rgb}{0.941176,0.584314,0.337255}%
\pgfsetstrokecolor{currentstroke}%
\pgfsetstrokeopacity{0.900000}%
\pgfsetdash{}{0pt}%
\pgfsys@defobject{currentmarker}{\pgfqpoint{-0.059738in}{-0.059738in}}{\pgfqpoint{0.059738in}{0.059738in}}{%
\pgfpathmoveto{\pgfqpoint{0.000000in}{-0.059738in}}%
\pgfpathlineto{\pgfqpoint{0.059738in}{0.000000in}}%
\pgfpathlineto{\pgfqpoint{0.000000in}{0.059738in}}%
\pgfpathlineto{\pgfqpoint{-0.059738in}{0.000000in}}%
\pgfpathlineto{\pgfqpoint{0.000000in}{-0.059738in}}%
\pgfpathclose%
\pgfusepath{stroke,fill}%
}%
\begin{pgfscope}%
\pgfsys@transformshift{13.436091in}{9.450000in}%
\pgfsys@useobject{currentmarker}{}%
\end{pgfscope}%
\begin{pgfscope}%
\pgfsys@transformshift{13.143236in}{9.450000in}%
\pgfsys@useobject{currentmarker}{}%
\end{pgfscope}%
\begin{pgfscope}%
\pgfsys@transformshift{13.436091in}{9.450000in}%
\pgfsys@useobject{currentmarker}{}%
\end{pgfscope}%
\end{pgfscope}%
\begin{pgfscope}%
\pgfpathrectangle{\pgfqpoint{8.036091in}{4.050000in}}{\pgfqpoint{5.400000in}{5.400000in}}%
\pgfusepath{clip}%
\pgfsetbuttcap%
\pgfsetroundjoin%
\definecolor{currentfill}{rgb}{0.549020,0.160784,0.505882}%
\pgfsetfillcolor{currentfill}%
\pgfsetfillopacity{0.900000}%
\pgfsetlinewidth{0.803000pt}%
\definecolor{currentstroke}{rgb}{0.549020,0.160784,0.505882}%
\pgfsetstrokecolor{currentstroke}%
\pgfsetstrokeopacity{0.900000}%
\pgfsetdash{}{0pt}%
\pgfsys@defobject{currentmarker}{\pgfqpoint{-0.038665in}{-0.038665in}}{\pgfqpoint{0.038665in}{0.038665in}}{%
\pgfpathmoveto{\pgfqpoint{0.000000in}{-0.038665in}}%
\pgfpathcurveto{\pgfqpoint{0.010254in}{-0.038665in}}{\pgfqpoint{0.020090in}{-0.034591in}}{\pgfqpoint{0.027340in}{-0.027340in}}%
\pgfpathcurveto{\pgfqpoint{0.034591in}{-0.020090in}}{\pgfqpoint{0.038665in}{-0.010254in}}{\pgfqpoint{0.038665in}{0.000000in}}%
\pgfpathcurveto{\pgfqpoint{0.038665in}{0.010254in}}{\pgfqpoint{0.034591in}{0.020090in}}{\pgfqpoint{0.027340in}{0.027340in}}%
\pgfpathcurveto{\pgfqpoint{0.020090in}{0.034591in}}{\pgfqpoint{0.010254in}{0.038665in}}{\pgfqpoint{0.000000in}{0.038665in}}%
\pgfpathcurveto{\pgfqpoint{-0.010254in}{0.038665in}}{\pgfqpoint{-0.020090in}{0.034591in}}{\pgfqpoint{-0.027340in}{0.027340in}}%
\pgfpathcurveto{\pgfqpoint{-0.034591in}{0.020090in}}{\pgfqpoint{-0.038665in}{0.010254in}}{\pgfqpoint{-0.038665in}{0.000000in}}%
\pgfpathcurveto{\pgfqpoint{-0.038665in}{-0.010254in}}{\pgfqpoint{-0.034591in}{-0.020090in}}{\pgfqpoint{-0.027340in}{-0.027340in}}%
\pgfpathcurveto{\pgfqpoint{-0.020090in}{-0.034591in}}{\pgfqpoint{-0.010254in}{-0.038665in}}{\pgfqpoint{0.000000in}{-0.038665in}}%
\pgfpathlineto{\pgfqpoint{0.000000in}{-0.038665in}}%
\pgfpathclose%
\pgfusepath{stroke,fill}%
}%
\begin{pgfscope}%
\pgfsys@transformshift{12.346950in}{6.171195in}%
\pgfsys@useobject{currentmarker}{}%
\end{pgfscope}%
\end{pgfscope}%
\begin{pgfscope}%
\pgfsetbuttcap%
\pgfsetroundjoin%
\definecolor{currentfill}{rgb}{0.549020,0.160784,0.505882}%
\pgfsetfillcolor{currentfill}%
\pgfsetfillopacity{0.900000}%
\pgfsetlinewidth{0.803000pt}%
\definecolor{currentstroke}{rgb}{0.549020,0.160784,0.505882}%
\pgfsetstrokecolor{currentstroke}%
\pgfsetstrokeopacity{0.900000}%
\pgfsetdash{}{0pt}%
\pgfsys@defobject{currentmarker}{\pgfqpoint{-0.059738in}{-0.059738in}}{\pgfqpoint{0.059738in}{0.059738in}}{%
\pgfpathmoveto{\pgfqpoint{0.000000in}{-0.059738in}}%
\pgfpathlineto{\pgfqpoint{0.059738in}{0.000000in}}%
\pgfpathlineto{\pgfqpoint{0.000000in}{0.059738in}}%
\pgfpathlineto{\pgfqpoint{-0.059738in}{0.000000in}}%
\pgfpathlineto{\pgfqpoint{0.000000in}{-0.059738in}}%
\pgfpathclose%
\pgfusepath{stroke,fill}%
}%
\begin{pgfscope}%
\pgfsys@transformshift{13.436091in}{4.702196in}%
\pgfsys@useobject{currentmarker}{}%
\end{pgfscope}%
\begin{pgfscope}%
\pgfsys@transformshift{13.436091in}{9.450000in}%
\pgfsys@useobject{currentmarker}{}%
\end{pgfscope}%
\end{pgfscope}%
\begin{pgfscope}%
\pgfpathrectangle{\pgfqpoint{8.036091in}{4.050000in}}{\pgfqpoint{5.400000in}{5.400000in}}%
\pgfusepath{clip}%
\pgfsetbuttcap%
\pgfsetroundjoin%
\definecolor{currentfill}{rgb}{0.000000,0.560784,0.611765}%
\pgfsetfillcolor{currentfill}%
\pgfsetfillopacity{0.900000}%
\pgfsetlinewidth{0.803000pt}%
\definecolor{currentstroke}{rgb}{0.000000,0.560784,0.611765}%
\pgfsetstrokecolor{currentstroke}%
\pgfsetstrokeopacity{0.900000}%
\pgfsetdash{}{0pt}%
\pgfsys@defobject{currentmarker}{\pgfqpoint{-0.038665in}{-0.038665in}}{\pgfqpoint{0.038665in}{0.038665in}}{%
\pgfpathmoveto{\pgfqpoint{0.000000in}{-0.038665in}}%
\pgfpathcurveto{\pgfqpoint{0.010254in}{-0.038665in}}{\pgfqpoint{0.020090in}{-0.034591in}}{\pgfqpoint{0.027340in}{-0.027340in}}%
\pgfpathcurveto{\pgfqpoint{0.034591in}{-0.020090in}}{\pgfqpoint{0.038665in}{-0.010254in}}{\pgfqpoint{0.038665in}{0.000000in}}%
\pgfpathcurveto{\pgfqpoint{0.038665in}{0.010254in}}{\pgfqpoint{0.034591in}{0.020090in}}{\pgfqpoint{0.027340in}{0.027340in}}%
\pgfpathcurveto{\pgfqpoint{0.020090in}{0.034591in}}{\pgfqpoint{0.010254in}{0.038665in}}{\pgfqpoint{0.000000in}{0.038665in}}%
\pgfpathcurveto{\pgfqpoint{-0.010254in}{0.038665in}}{\pgfqpoint{-0.020090in}{0.034591in}}{\pgfqpoint{-0.027340in}{0.027340in}}%
\pgfpathcurveto{\pgfqpoint{-0.034591in}{0.020090in}}{\pgfqpoint{-0.038665in}{0.010254in}}{\pgfqpoint{-0.038665in}{0.000000in}}%
\pgfpathcurveto{\pgfqpoint{-0.038665in}{-0.010254in}}{\pgfqpoint{-0.034591in}{-0.020090in}}{\pgfqpoint{-0.027340in}{-0.027340in}}%
\pgfpathcurveto{\pgfqpoint{-0.020090in}{-0.034591in}}{\pgfqpoint{-0.010254in}{-0.038665in}}{\pgfqpoint{0.000000in}{-0.038665in}}%
\pgfpathlineto{\pgfqpoint{0.000000in}{-0.038665in}}%
\pgfpathclose%
\pgfusepath{stroke,fill}%
}%
\begin{pgfscope}%
\pgfsys@transformshift{8.715229in}{5.277759in}%
\pgfsys@useobject{currentmarker}{}%
\end{pgfscope}%
\end{pgfscope}%
\begin{pgfscope}%
\pgfsetbuttcap%
\pgfsetroundjoin%
\definecolor{currentfill}{rgb}{0.000000,0.560784,0.611765}%
\pgfsetfillcolor{currentfill}%
\pgfsetfillopacity{0.900000}%
\pgfsetlinewidth{0.803000pt}%
\definecolor{currentstroke}{rgb}{0.000000,0.560784,0.611765}%
\pgfsetstrokecolor{currentstroke}%
\pgfsetstrokeopacity{0.900000}%
\pgfsetdash{}{0pt}%
\pgfsys@defobject{currentmarker}{\pgfqpoint{-0.059738in}{-0.059738in}}{\pgfqpoint{0.059738in}{0.059738in}}{%
\pgfpathmoveto{\pgfqpoint{0.000000in}{-0.059738in}}%
\pgfpathlineto{\pgfqpoint{0.059738in}{0.000000in}}%
\pgfpathlineto{\pgfqpoint{0.000000in}{0.059738in}}%
\pgfpathlineto{\pgfqpoint{-0.059738in}{0.000000in}}%
\pgfpathlineto{\pgfqpoint{0.000000in}{-0.059738in}}%
\pgfpathclose%
\pgfusepath{stroke,fill}%
}%
\begin{pgfscope}%
\pgfsys@transformshift{8.036091in}{4.050000in}%
\pgfsys@useobject{currentmarker}{}%
\end{pgfscope}%
\end{pgfscope}%
\begin{pgfscope}%
\pgfsetbuttcap%
\pgfsetroundjoin%
\pgfsetlinewidth{0.903375pt}%
\definecolor{currentstroke}{rgb}{0.321569,0.321569,0.321569}%
\pgfsetstrokecolor{currentstroke}%
\pgfsetdash{}{0pt}%
\pgfpathmoveto{\pgfqpoint{8.036091in}{3.977166in}}%
\pgfpathlineto{\pgfqpoint{8.108925in}{4.050000in}}%
\pgfpathlineto{\pgfqpoint{8.036091in}{4.122834in}}%
\pgfpathlineto{\pgfqpoint{7.963257in}{4.050000in}}%
\pgfpathlineto{\pgfqpoint{8.036091in}{3.977166in}}%
\pgfpathclose%
\pgfusepath{stroke}%
\end{pgfscope}%
\begin{pgfscope}%
\pgfsetbuttcap%
\pgfsetroundjoin%
\pgfsetlinewidth{0.903375pt}%
\definecolor{currentstroke}{rgb}{0.321569,0.321569,0.321569}%
\pgfsetstrokecolor{currentstroke}%
\pgfsetdash{}{0pt}%
\pgfpathmoveto{\pgfqpoint{8.036091in}{3.977166in}}%
\pgfpathlineto{\pgfqpoint{8.108925in}{4.050000in}}%
\pgfpathlineto{\pgfqpoint{8.036091in}{4.122834in}}%
\pgfpathlineto{\pgfqpoint{7.963257in}{4.050000in}}%
\pgfpathlineto{\pgfqpoint{8.036091in}{3.977166in}}%
\pgfpathclose%
\pgfusepath{stroke}%
\end{pgfscope}%
\begin{pgfscope}%
\pgfsetbuttcap%
\pgfsetroundjoin%
\pgfsetlinewidth{0.903375pt}%
\definecolor{currentstroke}{rgb}{0.321569,0.321569,0.321569}%
\pgfsetstrokecolor{currentstroke}%
\pgfsetdash{}{0pt}%
\pgfpathmoveto{\pgfqpoint{8.036091in}{3.977166in}}%
\pgfpathlineto{\pgfqpoint{8.108925in}{4.050000in}}%
\pgfpathlineto{\pgfqpoint{8.036091in}{4.122834in}}%
\pgfpathlineto{\pgfqpoint{7.963257in}{4.050000in}}%
\pgfpathlineto{\pgfqpoint{8.036091in}{3.977166in}}%
\pgfpathclose%
\pgfusepath{stroke}%
\end{pgfscope}%
\begin{pgfscope}%
\pgfsetbuttcap%
\pgfsetroundjoin%
\pgfsetlinewidth{0.903375pt}%
\definecolor{currentstroke}{rgb}{0.321569,0.321569,0.321569}%
\pgfsetstrokecolor{currentstroke}%
\pgfsetdash{}{0pt}%
\pgfpathmoveto{\pgfqpoint{8.036091in}{3.977166in}}%
\pgfpathlineto{\pgfqpoint{8.108925in}{4.050000in}}%
\pgfpathlineto{\pgfqpoint{8.036091in}{4.122834in}}%
\pgfpathlineto{\pgfqpoint{7.963257in}{4.050000in}}%
\pgfpathlineto{\pgfqpoint{8.036091in}{3.977166in}}%
\pgfpathclose%
\pgfusepath{stroke}%
\end{pgfscope}%
\begin{pgfscope}%
\pgfpathrectangle{\pgfqpoint{8.036091in}{4.050000in}}{\pgfqpoint{5.400000in}{5.400000in}}%
\pgfusepath{clip}%
\pgfsetbuttcap%
\pgfsetroundjoin%
\pgfsetlinewidth{0.803000pt}%
\definecolor{currentstroke}{rgb}{0.768627,0.603922,0.000000}%
\pgfsetstrokecolor{currentstroke}%
\pgfsetdash{}{0pt}%
\pgfpathmoveto{\pgfqpoint{12.165740in}{7.904813in}}%
\pgfpathlineto{\pgfqpoint{12.258909in}{7.904813in}}%
\pgfpathlineto{\pgfqpoint{12.258909in}{7.997982in}}%
\pgfpathlineto{\pgfqpoint{12.165740in}{7.997982in}}%
\pgfpathlineto{\pgfqpoint{12.165740in}{7.904813in}}%
\pgfpathclose%
\pgfusepath{stroke}%
\end{pgfscope}%
\begin{pgfscope}%
\pgfpathrectangle{\pgfqpoint{8.036091in}{4.050000in}}{\pgfqpoint{5.400000in}{5.400000in}}%
\pgfusepath{clip}%
\pgfsetbuttcap%
\pgfsetroundjoin%
\pgfsetlinewidth{0.803000pt}%
\definecolor{currentstroke}{rgb}{0.768627,0.603922,0.000000}%
\pgfsetstrokecolor{currentstroke}%
\pgfsetdash{}{0pt}%
\pgfpathmoveto{\pgfqpoint{11.919235in}{7.921261in}}%
\pgfpathlineto{\pgfqpoint{12.012405in}{7.921261in}}%
\pgfpathlineto{\pgfqpoint{12.012405in}{8.014431in}}%
\pgfpathlineto{\pgfqpoint{11.919235in}{8.014431in}}%
\pgfpathlineto{\pgfqpoint{11.919235in}{7.921261in}}%
\pgfpathclose%
\pgfusepath{stroke}%
\end{pgfscope}%
\begin{pgfscope}%
\pgfpathrectangle{\pgfqpoint{8.036091in}{4.050000in}}{\pgfqpoint{5.400000in}{5.400000in}}%
\pgfusepath{clip}%
\pgfsetbuttcap%
\pgfsetroundjoin%
\pgfsetlinewidth{0.803000pt}%
\definecolor{currentstroke}{rgb}{0.768627,0.603922,0.000000}%
\pgfsetstrokecolor{currentstroke}%
\pgfsetdash{}{0pt}%
\pgfpathmoveto{\pgfqpoint{10.369257in}{6.028176in}}%
\pgfpathlineto{\pgfqpoint{10.462427in}{6.028176in}}%
\pgfpathlineto{\pgfqpoint{10.462427in}{6.121346in}}%
\pgfpathlineto{\pgfqpoint{10.369257in}{6.121346in}}%
\pgfpathlineto{\pgfqpoint{10.369257in}{6.028176in}}%
\pgfpathclose%
\pgfusepath{stroke}%
\end{pgfscope}%
\begin{pgfscope}%
\pgfpathrectangle{\pgfqpoint{8.036091in}{4.050000in}}{\pgfqpoint{5.400000in}{5.400000in}}%
\pgfusepath{clip}%
\pgfsetbuttcap%
\pgfsetroundjoin%
\pgfsetlinewidth{0.803000pt}%
\definecolor{currentstroke}{rgb}{0.768627,0.603922,0.000000}%
\pgfsetstrokecolor{currentstroke}%
\pgfsetdash{}{0pt}%
\pgfpathmoveto{\pgfqpoint{10.668108in}{6.610062in}}%
\pgfpathlineto{\pgfqpoint{10.761277in}{6.610062in}}%
\pgfpathlineto{\pgfqpoint{10.761277in}{6.703232in}}%
\pgfpathlineto{\pgfqpoint{10.668108in}{6.703232in}}%
\pgfpathlineto{\pgfqpoint{10.668108in}{6.610062in}}%
\pgfpathclose%
\pgfusepath{stroke}%
\end{pgfscope}%
\begin{pgfscope}%
\pgfsetbuttcap%
\pgfsetroundjoin%
\pgfsetlinewidth{0.903375pt}%
\definecolor{currentstroke}{rgb}{0.768627,0.603922,0.000000}%
\pgfsetstrokecolor{currentstroke}%
\pgfsetdash{}{0pt}%
\pgfpathmoveto{\pgfqpoint{8.036091in}{3.977166in}}%
\pgfpathlineto{\pgfqpoint{8.108925in}{4.050000in}}%
\pgfpathlineto{\pgfqpoint{8.036091in}{4.122834in}}%
\pgfpathlineto{\pgfqpoint{7.963257in}{4.050000in}}%
\pgfpathlineto{\pgfqpoint{8.036091in}{3.977166in}}%
\pgfpathclose%
\pgfusepath{stroke}%
\end{pgfscope}%
\begin{pgfscope}%
\pgfsetbuttcap%
\pgfsetroundjoin%
\pgfsetlinewidth{0.903375pt}%
\definecolor{currentstroke}{rgb}{0.768627,0.603922,0.000000}%
\pgfsetstrokecolor{currentstroke}%
\pgfsetdash{}{0pt}%
\pgfpathmoveto{\pgfqpoint{12.826361in}{9.377166in}}%
\pgfpathlineto{\pgfqpoint{12.899194in}{9.450000in}}%
\pgfpathlineto{\pgfqpoint{12.826361in}{9.522834in}}%
\pgfpathlineto{\pgfqpoint{12.753527in}{9.450000in}}%
\pgfpathlineto{\pgfqpoint{12.826361in}{9.377166in}}%
\pgfpathclose%
\pgfusepath{stroke}%
\end{pgfscope}%
\begin{pgfscope}%
\pgfsetbuttcap%
\pgfsetroundjoin%
\pgfsetlinewidth{0.903375pt}%
\definecolor{currentstroke}{rgb}{0.768627,0.603922,0.000000}%
\pgfsetstrokecolor{currentstroke}%
\pgfsetdash{}{0pt}%
\pgfpathmoveto{\pgfqpoint{12.779482in}{9.377166in}}%
\pgfpathlineto{\pgfqpoint{12.852316in}{9.450000in}}%
\pgfpathlineto{\pgfqpoint{12.779482in}{9.522834in}}%
\pgfpathlineto{\pgfqpoint{12.706648in}{9.450000in}}%
\pgfpathlineto{\pgfqpoint{12.779482in}{9.377166in}}%
\pgfpathclose%
\pgfusepath{stroke}%
\end{pgfscope}%
\begin{pgfscope}%
\pgfsetbuttcap%
\pgfsetroundjoin%
\pgfsetlinewidth{0.903375pt}%
\definecolor{currentstroke}{rgb}{0.768627,0.603922,0.000000}%
\pgfsetstrokecolor{currentstroke}%
\pgfsetdash{}{0pt}%
\pgfpathmoveto{\pgfqpoint{12.623336in}{9.377166in}}%
\pgfpathlineto{\pgfqpoint{12.696170in}{9.450000in}}%
\pgfpathlineto{\pgfqpoint{12.623336in}{9.522834in}}%
\pgfpathlineto{\pgfqpoint{12.550502in}{9.450000in}}%
\pgfpathlineto{\pgfqpoint{12.623336in}{9.377166in}}%
\pgfpathclose%
\pgfusepath{stroke}%
\end{pgfscope}%
\begin{pgfscope}%
\pgfsetbuttcap%
\pgfsetroundjoin%
\pgfsetlinewidth{0.903375pt}%
\definecolor{currentstroke}{rgb}{0.768627,0.603922,0.000000}%
\pgfsetstrokecolor{currentstroke}%
\pgfsetdash{}{0pt}%
\pgfpathmoveto{\pgfqpoint{8.036091in}{3.977166in}}%
\pgfpathlineto{\pgfqpoint{8.108925in}{4.050000in}}%
\pgfpathlineto{\pgfqpoint{8.036091in}{4.122834in}}%
\pgfpathlineto{\pgfqpoint{7.963257in}{4.050000in}}%
\pgfpathlineto{\pgfqpoint{8.036091in}{3.977166in}}%
\pgfpathclose%
\pgfusepath{stroke}%
\end{pgfscope}%
\begin{pgfscope}%
\pgfsetbuttcap%
\pgfsetroundjoin%
\pgfsetlinewidth{0.903375pt}%
\definecolor{currentstroke}{rgb}{0.709804,0.741176,0.207843}%
\pgfsetstrokecolor{currentstroke}%
\pgfsetdash{}{0pt}%
\pgfpathmoveto{\pgfqpoint{8.607902in}{3.977166in}}%
\pgfpathlineto{\pgfqpoint{8.680736in}{4.050000in}}%
\pgfpathlineto{\pgfqpoint{8.607902in}{4.122834in}}%
\pgfpathlineto{\pgfqpoint{8.535068in}{4.050000in}}%
\pgfpathlineto{\pgfqpoint{8.607902in}{3.977166in}}%
\pgfpathclose%
\pgfusepath{stroke}%
\end{pgfscope}%
\begin{pgfscope}%
\pgfsetbuttcap%
\pgfsetroundjoin%
\pgfsetlinewidth{0.903375pt}%
\definecolor{currentstroke}{rgb}{0.709804,0.741176,0.207843}%
\pgfsetstrokecolor{currentstroke}%
\pgfsetdash{}{0pt}%
\pgfpathmoveto{\pgfqpoint{8.126892in}{3.977166in}}%
\pgfpathlineto{\pgfqpoint{8.199726in}{4.050000in}}%
\pgfpathlineto{\pgfqpoint{8.126892in}{4.122834in}}%
\pgfpathlineto{\pgfqpoint{8.054058in}{4.050000in}}%
\pgfpathlineto{\pgfqpoint{8.126892in}{3.977166in}}%
\pgfpathclose%
\pgfusepath{stroke}%
\end{pgfscope}%
\begin{pgfscope}%
\pgfsetbuttcap%
\pgfsetroundjoin%
\pgfsetlinewidth{0.903375pt}%
\definecolor{currentstroke}{rgb}{0.709804,0.741176,0.207843}%
\pgfsetstrokecolor{currentstroke}%
\pgfsetdash{}{0pt}%
\pgfpathmoveto{\pgfqpoint{8.364839in}{3.977166in}}%
\pgfpathlineto{\pgfqpoint{8.437673in}{4.050000in}}%
\pgfpathlineto{\pgfqpoint{8.364839in}{4.122834in}}%
\pgfpathlineto{\pgfqpoint{8.292005in}{4.050000in}}%
\pgfpathlineto{\pgfqpoint{8.364839in}{3.977166in}}%
\pgfpathclose%
\pgfusepath{stroke}%
\end{pgfscope}%
\begin{pgfscope}%
\pgfsetbuttcap%
\pgfsetroundjoin%
\pgfsetlinewidth{0.903375pt}%
\definecolor{currentstroke}{rgb}{0.709804,0.741176,0.207843}%
\pgfsetstrokecolor{currentstroke}%
\pgfsetdash{}{0pt}%
\pgfpathmoveto{\pgfqpoint{8.036091in}{3.977166in}}%
\pgfpathlineto{\pgfqpoint{8.108925in}{4.050000in}}%
\pgfpathlineto{\pgfqpoint{8.036091in}{4.122834in}}%
\pgfpathlineto{\pgfqpoint{7.963257in}{4.050000in}}%
\pgfpathlineto{\pgfqpoint{8.036091in}{3.977166in}}%
\pgfpathclose%
\pgfusepath{stroke}%
\end{pgfscope}%
\begin{pgfscope}%
\pgfsetbuttcap%
\pgfsetroundjoin%
\pgfsetlinewidth{0.903375pt}%
\definecolor{currentstroke}{rgb}{0.709804,0.741176,0.207843}%
\pgfsetstrokecolor{currentstroke}%
\pgfsetdash{}{0pt}%
\pgfpathmoveto{\pgfqpoint{8.036091in}{3.977166in}}%
\pgfpathlineto{\pgfqpoint{8.108925in}{4.050000in}}%
\pgfpathlineto{\pgfqpoint{8.036091in}{4.122834in}}%
\pgfpathlineto{\pgfqpoint{7.963257in}{4.050000in}}%
\pgfpathlineto{\pgfqpoint{8.036091in}{3.977166in}}%
\pgfpathclose%
\pgfusepath{stroke}%
\end{pgfscope}%
\begin{pgfscope}%
\pgfpathrectangle{\pgfqpoint{8.036091in}{4.050000in}}{\pgfqpoint{5.400000in}{5.400000in}}%
\pgfusepath{clip}%
\pgfsetbuttcap%
\pgfsetroundjoin%
\pgfsetlinewidth{0.803000pt}%
\definecolor{currentstroke}{rgb}{0.200000,0.133333,0.533333}%
\pgfsetstrokecolor{currentstroke}%
\pgfsetdash{}{0pt}%
\pgfpathmoveto{\pgfqpoint{11.127061in}{7.323025in}}%
\pgfpathlineto{\pgfqpoint{11.220231in}{7.323025in}}%
\pgfpathlineto{\pgfqpoint{11.220231in}{7.416195in}}%
\pgfpathlineto{\pgfqpoint{11.127061in}{7.416195in}}%
\pgfpathlineto{\pgfqpoint{11.127061in}{7.323025in}}%
\pgfpathclose%
\pgfusepath{stroke}%
\end{pgfscope}%
\begin{pgfscope}%
\pgfpathrectangle{\pgfqpoint{8.036091in}{4.050000in}}{\pgfqpoint{5.400000in}{5.400000in}}%
\pgfusepath{clip}%
\pgfsetbuttcap%
\pgfsetroundjoin%
\pgfsetlinewidth{0.803000pt}%
\definecolor{currentstroke}{rgb}{0.200000,0.133333,0.533333}%
\pgfsetstrokecolor{currentstroke}%
\pgfsetdash{}{0pt}%
\pgfpathmoveto{\pgfqpoint{10.923260in}{6.997765in}}%
\pgfpathlineto{\pgfqpoint{11.016430in}{6.997765in}}%
\pgfpathlineto{\pgfqpoint{11.016430in}{7.090935in}}%
\pgfpathlineto{\pgfqpoint{10.923260in}{7.090935in}}%
\pgfpathlineto{\pgfqpoint{10.923260in}{6.997765in}}%
\pgfpathclose%
\pgfusepath{stroke}%
\end{pgfscope}%
\begin{pgfscope}%
\pgfpathrectangle{\pgfqpoint{8.036091in}{4.050000in}}{\pgfqpoint{5.400000in}{5.400000in}}%
\pgfusepath{clip}%
\pgfsetbuttcap%
\pgfsetroundjoin%
\pgfsetlinewidth{0.803000pt}%
\definecolor{currentstroke}{rgb}{0.200000,0.133333,0.533333}%
\pgfsetstrokecolor{currentstroke}%
\pgfsetdash{}{0pt}%
\pgfpathmoveto{\pgfqpoint{10.645079in}{6.642523in}}%
\pgfpathlineto{\pgfqpoint{10.738248in}{6.642523in}}%
\pgfpathlineto{\pgfqpoint{10.738248in}{6.735693in}}%
\pgfpathlineto{\pgfqpoint{10.645079in}{6.735693in}}%
\pgfpathlineto{\pgfqpoint{10.645079in}{6.642523in}}%
\pgfpathclose%
\pgfusepath{stroke}%
\end{pgfscope}%
\begin{pgfscope}%
\pgfpathrectangle{\pgfqpoint{8.036091in}{4.050000in}}{\pgfqpoint{5.400000in}{5.400000in}}%
\pgfusepath{clip}%
\pgfsetbuttcap%
\pgfsetroundjoin%
\pgfsetlinewidth{0.803000pt}%
\definecolor{currentstroke}{rgb}{0.901961,0.196078,0.156863}%
\pgfsetstrokecolor{currentstroke}%
\pgfsetdash{}{0pt}%
\pgfpathmoveto{\pgfqpoint{12.112602in}{7.720761in}}%
\pgfpathlineto{\pgfqpoint{12.205772in}{7.720761in}}%
\pgfpathlineto{\pgfqpoint{12.205772in}{7.813930in}}%
\pgfpathlineto{\pgfqpoint{12.112602in}{7.813930in}}%
\pgfpathlineto{\pgfqpoint{12.112602in}{7.720761in}}%
\pgfpathclose%
\pgfusepath{stroke}%
\end{pgfscope}%
\begin{pgfscope}%
\pgfpathrectangle{\pgfqpoint{8.036091in}{4.050000in}}{\pgfqpoint{5.400000in}{5.400000in}}%
\pgfusepath{clip}%
\pgfsetbuttcap%
\pgfsetroundjoin%
\pgfsetlinewidth{0.803000pt}%
\definecolor{currentstroke}{rgb}{0.901961,0.196078,0.156863}%
\pgfsetstrokecolor{currentstroke}%
\pgfsetdash{}{0pt}%
\pgfpathmoveto{\pgfqpoint{10.895721in}{6.567171in}}%
\pgfpathlineto{\pgfqpoint{10.988891in}{6.567171in}}%
\pgfpathlineto{\pgfqpoint{10.988891in}{6.660341in}}%
\pgfpathlineto{\pgfqpoint{10.895721in}{6.660341in}}%
\pgfpathlineto{\pgfqpoint{10.895721in}{6.567171in}}%
\pgfpathclose%
\pgfusepath{stroke}%
\end{pgfscope}%
\begin{pgfscope}%
\pgfsetbuttcap%
\pgfsetroundjoin%
\pgfsetlinewidth{0.803000pt}%
\definecolor{currentstroke}{rgb}{0.901961,0.196078,0.156863}%
\pgfsetstrokecolor{currentstroke}%
\pgfsetdash{}{0pt}%
\pgfpathmoveto{\pgfqpoint{13.436091in}{9.356314in}}%
\pgfpathlineto{\pgfqpoint{13.529777in}{9.450000in}}%
\pgfpathlineto{\pgfqpoint{13.436091in}{9.543686in}}%
\pgfpathlineto{\pgfqpoint{13.342405in}{9.450000in}}%
\pgfpathlineto{\pgfqpoint{13.436091in}{9.356314in}}%
\pgfpathclose%
\pgfusepath{stroke}%
\end{pgfscope}%
\begin{pgfscope}%
\pgfsetbuttcap%
\pgfsetroundjoin%
\pgfsetlinewidth{0.803000pt}%
\definecolor{currentstroke}{rgb}{0.901961,0.196078,0.156863}%
\pgfsetstrokecolor{currentstroke}%
\pgfsetdash{}{0pt}%
\pgfpathmoveto{\pgfqpoint{13.389506in}{9.403415in}}%
\pgfpathlineto{\pgfqpoint{13.482676in}{9.403415in}}%
\pgfpathlineto{\pgfqpoint{13.482676in}{9.496585in}}%
\pgfpathlineto{\pgfqpoint{13.389506in}{9.496585in}}%
\pgfpathlineto{\pgfqpoint{13.389506in}{9.403415in}}%
\pgfpathclose%
\pgfusepath{stroke}%
\end{pgfscope}%
\begin{pgfscope}%
\pgfsetbuttcap%
\pgfsetroundjoin%
\pgfsetlinewidth{0.803000pt}%
\definecolor{currentstroke}{rgb}{0.498039,0.560784,0.678431}%
\pgfsetstrokecolor{currentstroke}%
\pgfsetdash{}{0pt}%
\pgfpathmoveto{\pgfqpoint{8.036091in}{3.956314in}}%
\pgfpathlineto{\pgfqpoint{8.129777in}{4.050000in}}%
\pgfpathlineto{\pgfqpoint{8.036091in}{4.143686in}}%
\pgfpathlineto{\pgfqpoint{7.942405in}{4.050000in}}%
\pgfpathlineto{\pgfqpoint{8.036091in}{3.956314in}}%
\pgfpathclose%
\pgfusepath{stroke}%
\end{pgfscope}%
\begin{pgfscope}%
\pgfsetbuttcap%
\pgfsetroundjoin%
\pgfsetlinewidth{0.803000pt}%
\definecolor{currentstroke}{rgb}{0.709804,0.121569,0.270588}%
\pgfsetstrokecolor{currentstroke}%
\pgfsetdash{}{0pt}%
\pgfpathmoveto{\pgfqpoint{13.188796in}{9.356314in}}%
\pgfpathlineto{\pgfqpoint{13.282481in}{9.450000in}}%
\pgfpathlineto{\pgfqpoint{13.188796in}{9.543686in}}%
\pgfpathlineto{\pgfqpoint{13.095110in}{9.450000in}}%
\pgfpathlineto{\pgfqpoint{13.188796in}{9.356314in}}%
\pgfpathclose%
\pgfusepath{stroke}%
\end{pgfscope}%
\begin{pgfscope}%
\pgfpathrectangle{\pgfqpoint{8.036091in}{4.050000in}}{\pgfqpoint{5.400000in}{5.400000in}}%
\pgfusepath{clip}%
\pgfsetbuttcap%
\pgfsetroundjoin%
\definecolor{currentfill}{rgb}{0.541176,0.380392,0.658824}%
\pgfsetfillcolor{currentfill}%
\pgfsetlinewidth{1.254687pt}%
\definecolor{currentstroke}{rgb}{0.541176,0.380392,0.658824}%
\pgfsetstrokecolor{currentstroke}%
\pgfsetdash{}{0pt}%
\pgfsys@defobject{currentmarker}{\pgfqpoint{-0.043368in}{-0.043368in}}{\pgfqpoint{0.043368in}{0.043368in}}{%
\pgfpathmoveto{\pgfqpoint{-0.043368in}{-0.043368in}}%
\pgfpathlineto{\pgfqpoint{0.043368in}{0.043368in}}%
\pgfpathmoveto{\pgfqpoint{-0.043368in}{0.043368in}}%
\pgfpathlineto{\pgfqpoint{0.043368in}{-0.043368in}}%
\pgfusepath{stroke,fill}%
}%
\begin{pgfscope}%
\pgfsys@transformshift{11.130555in}{8.873485in}%
\pgfsys@useobject{currentmarker}{}%
\end{pgfscope}%
\begin{pgfscope}%
\pgfsys@transformshift{11.305180in}{8.780970in}%
\pgfsys@useobject{currentmarker}{}%
\end{pgfscope}%
\begin{pgfscope}%
\pgfsys@transformshift{11.182277in}{9.331092in}%
\pgfsys@useobject{currentmarker}{}%
\end{pgfscope}%
\begin{pgfscope}%
\pgfsys@transformshift{10.082481in}{6.710598in}%
\pgfsys@useobject{currentmarker}{}%
\end{pgfscope}%
\begin{pgfscope}%
\pgfsys@transformshift{10.725317in}{6.497209in}%
\pgfsys@useobject{currentmarker}{}%
\end{pgfscope}%
\begin{pgfscope}%
\pgfsys@transformshift{11.811615in}{5.828729in}%
\pgfsys@useobject{currentmarker}{}%
\end{pgfscope}%
\begin{pgfscope}%
\pgfsys@transformshift{10.178124in}{4.876134in}%
\pgfsys@useobject{currentmarker}{}%
\end{pgfscope}%
\begin{pgfscope}%
\pgfsys@transformshift{10.400591in}{6.754754in}%
\pgfsys@useobject{currentmarker}{}%
\end{pgfscope}%
\begin{pgfscope}%
\pgfsys@transformshift{10.427200in}{5.525686in}%
\pgfsys@useobject{currentmarker}{}%
\end{pgfscope}%
\begin{pgfscope}%
\pgfsys@transformshift{10.816222in}{7.121107in}%
\pgfsys@useobject{currentmarker}{}%
\end{pgfscope}%
\begin{pgfscope}%
\pgfsys@transformshift{10.868661in}{7.109337in}%
\pgfsys@useobject{currentmarker}{}%
\end{pgfscope}%
\end{pgfscope}%
\begin{pgfscope}%
\pgfpathrectangle{\pgfqpoint{8.036091in}{4.050000in}}{\pgfqpoint{5.400000in}{5.400000in}}%
\pgfusepath{clip}%
\pgfsetbuttcap%
\pgfsetroundjoin%
\definecolor{currentfill}{rgb}{0.321569,0.321569,0.321569}%
\pgfsetfillcolor{currentfill}%
\pgfsetlinewidth{1.254687pt}%
\definecolor{currentstroke}{rgb}{0.321569,0.321569,0.321569}%
\pgfsetstrokecolor{currentstroke}%
\pgfsetdash{}{0pt}%
\pgfsys@defobject{currentmarker}{\pgfqpoint{-0.043368in}{-0.043368in}}{\pgfqpoint{0.043368in}{0.043368in}}{%
\pgfpathmoveto{\pgfqpoint{-0.043368in}{-0.043368in}}%
\pgfpathlineto{\pgfqpoint{0.043368in}{0.043368in}}%
\pgfpathmoveto{\pgfqpoint{-0.043368in}{0.043368in}}%
\pgfpathlineto{\pgfqpoint{0.043368in}{-0.043368in}}%
\pgfusepath{stroke,fill}%
}%
\begin{pgfscope}%
\pgfsys@transformshift{9.822949in}{5.645904in}%
\pgfsys@useobject{currentmarker}{}%
\end{pgfscope}%
\begin{pgfscope}%
\pgfsys@transformshift{10.011486in}{5.927015in}%
\pgfsys@useobject{currentmarker}{}%
\end{pgfscope}%
\begin{pgfscope}%
\pgfsys@transformshift{9.627703in}{5.429145in}%
\pgfsys@useobject{currentmarker}{}%
\end{pgfscope}%
\begin{pgfscope}%
\pgfsys@transformshift{10.061926in}{5.868260in}%
\pgfsys@useobject{currentmarker}{}%
\end{pgfscope}%
\begin{pgfscope}%
\pgfsys@transformshift{9.489292in}{5.450176in}%
\pgfsys@useobject{currentmarker}{}%
\end{pgfscope}%
\begin{pgfscope}%
\pgfsys@transformshift{9.423983in}{5.541751in}%
\pgfsys@useobject{currentmarker}{}%
\end{pgfscope}%
\begin{pgfscope}%
\pgfsys@transformshift{9.797039in}{5.582538in}%
\pgfsys@useobject{currentmarker}{}%
\end{pgfscope}%
\begin{pgfscope}%
\pgfsys@transformshift{9.943351in}{5.770600in}%
\pgfsys@useobject{currentmarker}{}%
\end{pgfscope}%
\begin{pgfscope}%
\pgfsys@transformshift{9.451021in}{5.409478in}%
\pgfsys@useobject{currentmarker}{}%
\end{pgfscope}%
\begin{pgfscope}%
\pgfsys@transformshift{9.839137in}{5.588361in}%
\pgfsys@useobject{currentmarker}{}%
\end{pgfscope}%
\begin{pgfscope}%
\pgfsys@transformshift{10.079536in}{5.860073in}%
\pgfsys@useobject{currentmarker}{}%
\end{pgfscope}%
\begin{pgfscope}%
\pgfsys@transformshift{9.557294in}{5.444755in}%
\pgfsys@useobject{currentmarker}{}%
\end{pgfscope}%
\begin{pgfscope}%
\pgfsys@transformshift{9.888333in}{5.958769in}%
\pgfsys@useobject{currentmarker}{}%
\end{pgfscope}%
\begin{pgfscope}%
\pgfsys@transformshift{11.624514in}{7.890284in}%
\pgfsys@useobject{currentmarker}{}%
\end{pgfscope}%
\begin{pgfscope}%
\pgfsys@transformshift{9.313789in}{5.458422in}%
\pgfsys@useobject{currentmarker}{}%
\end{pgfscope}%
\begin{pgfscope}%
\pgfsys@transformshift{10.317242in}{6.372552in}%
\pgfsys@useobject{currentmarker}{}%
\end{pgfscope}%
\begin{pgfscope}%
\pgfsys@transformshift{9.100011in}{4.318944in}%
\pgfsys@useobject{currentmarker}{}%
\end{pgfscope}%
\begin{pgfscope}%
\pgfsys@transformshift{11.706989in}{6.803791in}%
\pgfsys@useobject{currentmarker}{}%
\end{pgfscope}%
\begin{pgfscope}%
\pgfsys@transformshift{13.284923in}{9.273225in}%
\pgfsys@useobject{currentmarker}{}%
\end{pgfscope}%
\end{pgfscope}%
\begin{pgfscope}%
\pgfsetbuttcap%
\pgfsetroundjoin%
\definecolor{currentfill}{rgb}{0.321569,0.321569,0.321569}%
\pgfsetfillcolor{currentfill}%
\pgfsetlinewidth{1.254687pt}%
\definecolor{currentstroke}{rgb}{0.321569,0.321569,0.321569}%
\pgfsetstrokecolor{currentstroke}%
\pgfsetdash{}{0pt}%
\pgfsys@defobject{currentmarker}{\pgfqpoint{-0.043368in}{-0.043368in}}{\pgfqpoint{0.043368in}{0.043368in}}{%
\pgfpathmoveto{\pgfqpoint{-0.043368in}{-0.043368in}}%
\pgfpathlineto{\pgfqpoint{0.043368in}{0.043368in}}%
\pgfpathmoveto{\pgfqpoint{-0.043368in}{0.043368in}}%
\pgfpathlineto{\pgfqpoint{0.043368in}{-0.043368in}}%
\pgfusepath{stroke,fill}%
}%
\begin{pgfscope}%
\pgfsys@transformshift{8.036091in}{4.050000in}%
\pgfsys@useobject{currentmarker}{}%
\end{pgfscope}%
\begin{pgfscope}%
\pgfsys@transformshift{8.036091in}{4.050000in}%
\pgfsys@useobject{currentmarker}{}%
\end{pgfscope}%
\begin{pgfscope}%
\pgfsys@transformshift{8.036091in}{4.050000in}%
\pgfsys@useobject{currentmarker}{}%
\end{pgfscope}%
\begin{pgfscope}%
\pgfsys@transformshift{8.036091in}{4.050000in}%
\pgfsys@useobject{currentmarker}{}%
\end{pgfscope}%
\end{pgfscope}%
\begin{pgfscope}%
\pgfpathrectangle{\pgfqpoint{8.036091in}{4.050000in}}{\pgfqpoint{5.400000in}{5.400000in}}%
\pgfusepath{clip}%
\pgfsetbuttcap%
\pgfsetroundjoin%
\definecolor{currentfill}{rgb}{0.768627,0.603922,0.000000}%
\pgfsetfillcolor{currentfill}%
\pgfsetlinewidth{1.254687pt}%
\definecolor{currentstroke}{rgb}{0.768627,0.603922,0.000000}%
\pgfsetstrokecolor{currentstroke}%
\pgfsetdash{}{0pt}%
\pgfsys@defobject{currentmarker}{\pgfqpoint{-0.043368in}{-0.043368in}}{\pgfqpoint{0.043368in}{0.043368in}}{%
\pgfpathmoveto{\pgfqpoint{-0.043368in}{-0.043368in}}%
\pgfpathlineto{\pgfqpoint{0.043368in}{0.043368in}}%
\pgfpathmoveto{\pgfqpoint{-0.043368in}{0.043368in}}%
\pgfpathlineto{\pgfqpoint{0.043368in}{-0.043368in}}%
\pgfusepath{stroke,fill}%
}%
\begin{pgfscope}%
\pgfsys@transformshift{12.212325in}{7.951398in}%
\pgfsys@useobject{currentmarker}{}%
\end{pgfscope}%
\begin{pgfscope}%
\pgfsys@transformshift{11.965820in}{7.967846in}%
\pgfsys@useobject{currentmarker}{}%
\end{pgfscope}%
\end{pgfscope}%
\begin{pgfscope}%
\pgfpathrectangle{\pgfqpoint{8.036091in}{4.050000in}}{\pgfqpoint{5.400000in}{5.400000in}}%
\pgfusepath{clip}%
\pgfsetbuttcap%
\pgfsetroundjoin%
\definecolor{currentfill}{rgb}{0.768627,0.603922,0.000000}%
\pgfsetfillcolor{currentfill}%
\pgfsetlinewidth{1.254687pt}%
\definecolor{currentstroke}{rgb}{0.768627,0.603922,0.000000}%
\pgfsetstrokecolor{currentstroke}%
\pgfsetdash{}{0pt}%
\pgfsys@defobject{currentmarker}{\pgfqpoint{-0.043368in}{-0.043368in}}{\pgfqpoint{0.043368in}{0.043368in}}{%
\pgfpathmoveto{\pgfqpoint{-0.043368in}{-0.043368in}}%
\pgfpathlineto{\pgfqpoint{0.043368in}{0.043368in}}%
\pgfpathmoveto{\pgfqpoint{-0.043368in}{0.043368in}}%
\pgfpathlineto{\pgfqpoint{0.043368in}{-0.043368in}}%
\pgfusepath{stroke,fill}%
}%
\begin{pgfscope}%
\pgfsys@transformshift{10.415842in}{6.074761in}%
\pgfsys@useobject{currentmarker}{}%
\end{pgfscope}%
\begin{pgfscope}%
\pgfsys@transformshift{10.714693in}{6.656647in}%
\pgfsys@useobject{currentmarker}{}%
\end{pgfscope}%
\end{pgfscope}%
\begin{pgfscope}%
\pgfpathrectangle{\pgfqpoint{8.036091in}{4.050000in}}{\pgfqpoint{5.400000in}{5.400000in}}%
\pgfusepath{clip}%
\pgfsetbuttcap%
\pgfsetroundjoin%
\definecolor{currentfill}{rgb}{0.768627,0.603922,0.000000}%
\pgfsetfillcolor{currentfill}%
\pgfsetlinewidth{1.254687pt}%
\definecolor{currentstroke}{rgb}{0.768627,0.603922,0.000000}%
\pgfsetstrokecolor{currentstroke}%
\pgfsetdash{}{0pt}%
\pgfsys@defobject{currentmarker}{\pgfqpoint{-0.043368in}{-0.043368in}}{\pgfqpoint{0.043368in}{0.043368in}}{%
\pgfpathmoveto{\pgfqpoint{-0.043368in}{-0.043368in}}%
\pgfpathlineto{\pgfqpoint{0.043368in}{0.043368in}}%
\pgfpathmoveto{\pgfqpoint{-0.043368in}{0.043368in}}%
\pgfpathlineto{\pgfqpoint{0.043368in}{-0.043368in}}%
\pgfusepath{stroke,fill}%
}%
\begin{pgfscope}%
\pgfsys@transformshift{11.976047in}{8.126731in}%
\pgfsys@useobject{currentmarker}{}%
\end{pgfscope}%
\begin{pgfscope}%
\pgfsys@transformshift{11.514899in}{8.268112in}%
\pgfsys@useobject{currentmarker}{}%
\end{pgfscope}%
\begin{pgfscope}%
\pgfsys@transformshift{8.552864in}{4.844266in}%
\pgfsys@useobject{currentmarker}{}%
\end{pgfscope}%
\begin{pgfscope}%
\pgfsys@transformshift{11.051655in}{6.628636in}%
\pgfsys@useobject{currentmarker}{}%
\end{pgfscope}%
\begin{pgfscope}%
\pgfsys@transformshift{12.347638in}{8.808224in}%
\pgfsys@useobject{currentmarker}{}%
\end{pgfscope}%
\end{pgfscope}%
\begin{pgfscope}%
\pgfsetbuttcap%
\pgfsetroundjoin%
\definecolor{currentfill}{rgb}{0.768627,0.603922,0.000000}%
\pgfsetfillcolor{currentfill}%
\pgfsetlinewidth{1.254687pt}%
\definecolor{currentstroke}{rgb}{0.768627,0.603922,0.000000}%
\pgfsetstrokecolor{currentstroke}%
\pgfsetdash{}{0pt}%
\pgfsys@defobject{currentmarker}{\pgfqpoint{-0.043368in}{-0.043368in}}{\pgfqpoint{0.043368in}{0.043368in}}{%
\pgfpathmoveto{\pgfqpoint{-0.043368in}{-0.043368in}}%
\pgfpathlineto{\pgfqpoint{0.043368in}{0.043368in}}%
\pgfpathmoveto{\pgfqpoint{-0.043368in}{0.043368in}}%
\pgfpathlineto{\pgfqpoint{0.043368in}{-0.043368in}}%
\pgfusepath{stroke,fill}%
}%
\begin{pgfscope}%
\pgfsys@transformshift{8.036091in}{4.050000in}%
\pgfsys@useobject{currentmarker}{}%
\end{pgfscope}%
\begin{pgfscope}%
\pgfsys@transformshift{12.826361in}{9.450000in}%
\pgfsys@useobject{currentmarker}{}%
\end{pgfscope}%
\begin{pgfscope}%
\pgfsys@transformshift{12.779482in}{9.450000in}%
\pgfsys@useobject{currentmarker}{}%
\end{pgfscope}%
\begin{pgfscope}%
\pgfsys@transformshift{12.623336in}{9.450000in}%
\pgfsys@useobject{currentmarker}{}%
\end{pgfscope}%
\begin{pgfscope}%
\pgfsys@transformshift{8.036091in}{4.050000in}%
\pgfsys@useobject{currentmarker}{}%
\end{pgfscope}%
\end{pgfscope}%
\begin{pgfscope}%
\pgfpathrectangle{\pgfqpoint{8.036091in}{4.050000in}}{\pgfqpoint{5.400000in}{5.400000in}}%
\pgfusepath{clip}%
\pgfsetbuttcap%
\pgfsetroundjoin%
\definecolor{currentfill}{rgb}{0.709804,0.741176,0.207843}%
\pgfsetfillcolor{currentfill}%
\pgfsetlinewidth{1.254687pt}%
\definecolor{currentstroke}{rgb}{0.709804,0.741176,0.207843}%
\pgfsetstrokecolor{currentstroke}%
\pgfsetdash{}{0pt}%
\pgfsys@defobject{currentmarker}{\pgfqpoint{-0.043368in}{-0.043368in}}{\pgfqpoint{0.043368in}{0.043368in}}{%
\pgfpathmoveto{\pgfqpoint{-0.043368in}{-0.043368in}}%
\pgfpathlineto{\pgfqpoint{0.043368in}{0.043368in}}%
\pgfpathmoveto{\pgfqpoint{-0.043368in}{0.043368in}}%
\pgfpathlineto{\pgfqpoint{0.043368in}{-0.043368in}}%
\pgfusepath{stroke,fill}%
}%
\begin{pgfscope}%
\pgfsys@transformshift{9.138089in}{4.806055in}%
\pgfsys@useobject{currentmarker}{}%
\end{pgfscope}%
\begin{pgfscope}%
\pgfsys@transformshift{9.426994in}{5.221143in}%
\pgfsys@useobject{currentmarker}{}%
\end{pgfscope}%
\begin{pgfscope}%
\pgfsys@transformshift{8.835872in}{4.329082in}%
\pgfsys@useobject{currentmarker}{}%
\end{pgfscope}%
\begin{pgfscope}%
\pgfsys@transformshift{8.972418in}{5.460856in}%
\pgfsys@useobject{currentmarker}{}%
\end{pgfscope}%
\begin{pgfscope}%
\pgfsys@transformshift{9.412740in}{5.348391in}%
\pgfsys@useobject{currentmarker}{}%
\end{pgfscope}%
\begin{pgfscope}%
\pgfsys@transformshift{9.298132in}{4.358153in}%
\pgfsys@useobject{currentmarker}{}%
\end{pgfscope}%
\begin{pgfscope}%
\pgfsys@transformshift{8.945493in}{5.485523in}%
\pgfsys@useobject{currentmarker}{}%
\end{pgfscope}%
\begin{pgfscope}%
\pgfsys@transformshift{9.519481in}{5.501668in}%
\pgfsys@useobject{currentmarker}{}%
\end{pgfscope}%
\begin{pgfscope}%
\pgfsys@transformshift{10.114215in}{5.644769in}%
\pgfsys@useobject{currentmarker}{}%
\end{pgfscope}%
\begin{pgfscope}%
\pgfsys@transformshift{9.757639in}{5.615270in}%
\pgfsys@useobject{currentmarker}{}%
\end{pgfscope}%
\begin{pgfscope}%
\pgfsys@transformshift{10.583743in}{6.300968in}%
\pgfsys@useobject{currentmarker}{}%
\end{pgfscope}%
\end{pgfscope}%
\begin{pgfscope}%
\pgfsetbuttcap%
\pgfsetroundjoin%
\definecolor{currentfill}{rgb}{0.709804,0.741176,0.207843}%
\pgfsetfillcolor{currentfill}%
\pgfsetlinewidth{1.254687pt}%
\definecolor{currentstroke}{rgb}{0.709804,0.741176,0.207843}%
\pgfsetstrokecolor{currentstroke}%
\pgfsetdash{}{0pt}%
\pgfsys@defobject{currentmarker}{\pgfqpoint{-0.043368in}{-0.043368in}}{\pgfqpoint{0.043368in}{0.043368in}}{%
\pgfpathmoveto{\pgfqpoint{-0.043368in}{-0.043368in}}%
\pgfpathlineto{\pgfqpoint{0.043368in}{0.043368in}}%
\pgfpathmoveto{\pgfqpoint{-0.043368in}{0.043368in}}%
\pgfpathlineto{\pgfqpoint{0.043368in}{-0.043368in}}%
\pgfusepath{stroke,fill}%
}%
\begin{pgfscope}%
\pgfsys@transformshift{8.607902in}{4.050000in}%
\pgfsys@useobject{currentmarker}{}%
\end{pgfscope}%
\begin{pgfscope}%
\pgfsys@transformshift{8.126892in}{4.050000in}%
\pgfsys@useobject{currentmarker}{}%
\end{pgfscope}%
\begin{pgfscope}%
\pgfsys@transformshift{8.364839in}{4.050000in}%
\pgfsys@useobject{currentmarker}{}%
\end{pgfscope}%
\begin{pgfscope}%
\pgfsys@transformshift{8.036091in}{4.050000in}%
\pgfsys@useobject{currentmarker}{}%
\end{pgfscope}%
\begin{pgfscope}%
\pgfsys@transformshift{8.036091in}{4.050000in}%
\pgfsys@useobject{currentmarker}{}%
\end{pgfscope}%
\end{pgfscope}%
\begin{pgfscope}%
\pgfpathrectangle{\pgfqpoint{8.036091in}{4.050000in}}{\pgfqpoint{5.400000in}{5.400000in}}%
\pgfusepath{clip}%
\pgfsetbuttcap%
\pgfsetroundjoin%
\definecolor{currentfill}{rgb}{0.200000,0.133333,0.533333}%
\pgfsetfillcolor{currentfill}%
\pgfsetlinewidth{1.254687pt}%
\definecolor{currentstroke}{rgb}{0.200000,0.133333,0.533333}%
\pgfsetstrokecolor{currentstroke}%
\pgfsetdash{}{0pt}%
\pgfsys@defobject{currentmarker}{\pgfqpoint{-0.043368in}{-0.043368in}}{\pgfqpoint{0.043368in}{0.043368in}}{%
\pgfpathmoveto{\pgfqpoint{-0.043368in}{-0.043368in}}%
\pgfpathlineto{\pgfqpoint{0.043368in}{0.043368in}}%
\pgfpathmoveto{\pgfqpoint{-0.043368in}{0.043368in}}%
\pgfpathlineto{\pgfqpoint{0.043368in}{-0.043368in}}%
\pgfusepath{stroke,fill}%
}%
\begin{pgfscope}%
\pgfsys@transformshift{11.173646in}{7.369610in}%
\pgfsys@useobject{currentmarker}{}%
\end{pgfscope}%
\begin{pgfscope}%
\pgfsys@transformshift{10.969845in}{7.044350in}%
\pgfsys@useobject{currentmarker}{}%
\end{pgfscope}%
\begin{pgfscope}%
\pgfsys@transformshift{10.691664in}{6.689108in}%
\pgfsys@useobject{currentmarker}{}%
\end{pgfscope}%
\end{pgfscope}%
\begin{pgfscope}%
\pgfpathrectangle{\pgfqpoint{8.036091in}{4.050000in}}{\pgfqpoint{5.400000in}{5.400000in}}%
\pgfusepath{clip}%
\pgfsetbuttcap%
\pgfsetroundjoin%
\definecolor{currentfill}{rgb}{0.200000,0.133333,0.533333}%
\pgfsetfillcolor{currentfill}%
\pgfsetlinewidth{1.254687pt}%
\definecolor{currentstroke}{rgb}{0.200000,0.133333,0.533333}%
\pgfsetstrokecolor{currentstroke}%
\pgfsetdash{}{0pt}%
\pgfsys@defobject{currentmarker}{\pgfqpoint{-0.043368in}{-0.043368in}}{\pgfqpoint{0.043368in}{0.043368in}}{%
\pgfpathmoveto{\pgfqpoint{-0.043368in}{-0.043368in}}%
\pgfpathlineto{\pgfqpoint{0.043368in}{0.043368in}}%
\pgfpathmoveto{\pgfqpoint{-0.043368in}{0.043368in}}%
\pgfpathlineto{\pgfqpoint{0.043368in}{-0.043368in}}%
\pgfusepath{stroke,fill}%
}%
\begin{pgfscope}%
\pgfsys@transformshift{8.982261in}{4.561158in}%
\pgfsys@useobject{currentmarker}{}%
\end{pgfscope}%
\begin{pgfscope}%
\pgfsys@transformshift{9.190754in}{4.755084in}%
\pgfsys@useobject{currentmarker}{}%
\end{pgfscope}%
\end{pgfscope}%
\begin{pgfscope}%
\pgfpathrectangle{\pgfqpoint{8.036091in}{4.050000in}}{\pgfqpoint{5.400000in}{5.400000in}}%
\pgfusepath{clip}%
\pgfsetbuttcap%
\pgfsetroundjoin%
\definecolor{currentfill}{rgb}{0.337255,0.705882,0.913725}%
\pgfsetfillcolor{currentfill}%
\pgfsetlinewidth{1.254687pt}%
\definecolor{currentstroke}{rgb}{0.337255,0.705882,0.913725}%
\pgfsetstrokecolor{currentstroke}%
\pgfsetdash{}{0pt}%
\pgfsys@defobject{currentmarker}{\pgfqpoint{-0.043368in}{-0.043368in}}{\pgfqpoint{0.043368in}{0.043368in}}{%
\pgfpathmoveto{\pgfqpoint{-0.043368in}{-0.043368in}}%
\pgfpathlineto{\pgfqpoint{0.043368in}{0.043368in}}%
\pgfpathmoveto{\pgfqpoint{-0.043368in}{0.043368in}}%
\pgfpathlineto{\pgfqpoint{0.043368in}{-0.043368in}}%
\pgfusepath{stroke,fill}%
}%
\begin{pgfscope}%
\pgfsys@transformshift{9.011593in}{5.563039in}%
\pgfsys@useobject{currentmarker}{}%
\end{pgfscope}%
\begin{pgfscope}%
\pgfsys@transformshift{9.121104in}{5.334784in}%
\pgfsys@useobject{currentmarker}{}%
\end{pgfscope}%
\end{pgfscope}%
\begin{pgfscope}%
\pgfpathrectangle{\pgfqpoint{8.036091in}{4.050000in}}{\pgfqpoint{5.400000in}{5.400000in}}%
\pgfusepath{clip}%
\pgfsetbuttcap%
\pgfsetroundjoin%
\definecolor{currentfill}{rgb}{0.901961,0.196078,0.156863}%
\pgfsetfillcolor{currentfill}%
\pgfsetlinewidth{1.254687pt}%
\definecolor{currentstroke}{rgb}{0.901961,0.196078,0.156863}%
\pgfsetstrokecolor{currentstroke}%
\pgfsetdash{}{0pt}%
\pgfsys@defobject{currentmarker}{\pgfqpoint{-0.043368in}{-0.043368in}}{\pgfqpoint{0.043368in}{0.043368in}}{%
\pgfpathmoveto{\pgfqpoint{-0.043368in}{-0.043368in}}%
\pgfpathlineto{\pgfqpoint{0.043368in}{0.043368in}}%
\pgfpathmoveto{\pgfqpoint{-0.043368in}{0.043368in}}%
\pgfpathlineto{\pgfqpoint{0.043368in}{-0.043368in}}%
\pgfusepath{stroke,fill}%
}%
\begin{pgfscope}%
\pgfsys@transformshift{12.159187in}{7.767346in}%
\pgfsys@useobject{currentmarker}{}%
\end{pgfscope}%
\begin{pgfscope}%
\pgfsys@transformshift{10.942306in}{6.613756in}%
\pgfsys@useobject{currentmarker}{}%
\end{pgfscope}%
\end{pgfscope}%
\begin{pgfscope}%
\pgfsetbuttcap%
\pgfsetroundjoin%
\definecolor{currentfill}{rgb}{0.901961,0.196078,0.156863}%
\pgfsetfillcolor{currentfill}%
\pgfsetlinewidth{1.254687pt}%
\definecolor{currentstroke}{rgb}{0.901961,0.196078,0.156863}%
\pgfsetstrokecolor{currentstroke}%
\pgfsetdash{}{0pt}%
\pgfsys@defobject{currentmarker}{\pgfqpoint{-0.043368in}{-0.043368in}}{\pgfqpoint{0.043368in}{0.043368in}}{%
\pgfpathmoveto{\pgfqpoint{-0.043368in}{-0.043368in}}%
\pgfpathlineto{\pgfqpoint{0.043368in}{0.043368in}}%
\pgfpathmoveto{\pgfqpoint{-0.043368in}{0.043368in}}%
\pgfpathlineto{\pgfqpoint{0.043368in}{-0.043368in}}%
\pgfusepath{stroke,fill}%
}%
\begin{pgfscope}%
\pgfsys@transformshift{13.436091in}{9.450000in}%
\pgfsys@useobject{currentmarker}{}%
\end{pgfscope}%
\end{pgfscope}%
\begin{pgfscope}%
\pgfpathrectangle{\pgfqpoint{8.036091in}{4.050000in}}{\pgfqpoint{5.400000in}{5.400000in}}%
\pgfusepath{clip}%
\pgfsetbuttcap%
\pgfsetroundjoin%
\definecolor{currentfill}{rgb}{0.835294,0.368627,0.000000}%
\pgfsetfillcolor{currentfill}%
\pgfsetlinewidth{1.254687pt}%
\definecolor{currentstroke}{rgb}{0.835294,0.368627,0.000000}%
\pgfsetstrokecolor{currentstroke}%
\pgfsetdash{}{0pt}%
\pgfsys@defobject{currentmarker}{\pgfqpoint{-0.043368in}{-0.043368in}}{\pgfqpoint{0.043368in}{0.043368in}}{%
\pgfpathmoveto{\pgfqpoint{-0.043368in}{-0.043368in}}%
\pgfpathlineto{\pgfqpoint{0.043368in}{0.043368in}}%
\pgfpathmoveto{\pgfqpoint{-0.043368in}{0.043368in}}%
\pgfpathlineto{\pgfqpoint{0.043368in}{-0.043368in}}%
\pgfusepath{stroke,fill}%
}%
\begin{pgfscope}%
\pgfsys@transformshift{11.850186in}{8.332332in}%
\pgfsys@useobject{currentmarker}{}%
\end{pgfscope}%
\end{pgfscope}%
\begin{pgfscope}%
\pgfpathrectangle{\pgfqpoint{8.036091in}{4.050000in}}{\pgfqpoint{5.400000in}{5.400000in}}%
\pgfusepath{clip}%
\pgfsetbuttcap%
\pgfsetroundjoin%
\pgfsetlinewidth{1.003750pt}%
\definecolor{currentstroke}{rgb}{0.541176,0.380392,0.658824}%
\pgfsetstrokecolor{currentstroke}%
\pgfsetstrokeopacity{0.900000}%
\pgfsetdash{}{0pt}%
\pgfpathmoveto{\pgfqpoint{11.915200in}{8.534658in}}%
\pgfpathlineto{\pgfqpoint{11.999683in}{8.534658in}}%
\pgfpathlineto{\pgfqpoint{11.999683in}{8.619141in}}%
\pgfpathlineto{\pgfqpoint{11.915200in}{8.619141in}}%
\pgfpathlineto{\pgfqpoint{11.915200in}{8.534658in}}%
\pgfpathclose%
\pgfusepath{stroke}%
\end{pgfscope}%
\begin{pgfscope}%
\pgfpathrectangle{\pgfqpoint{8.036091in}{4.050000in}}{\pgfqpoint{5.400000in}{5.400000in}}%
\pgfusepath{clip}%
\pgfsetbuttcap%
\pgfsetroundjoin%
\pgfsetlinewidth{1.003750pt}%
\definecolor{currentstroke}{rgb}{0.541176,0.380392,0.658824}%
\pgfsetstrokecolor{currentstroke}%
\pgfsetstrokeopacity{0.900000}%
\pgfsetdash{}{0pt}%
\pgfpathmoveto{\pgfqpoint{11.775883in}{7.566383in}}%
\pgfpathlineto{\pgfqpoint{11.860365in}{7.566383in}}%
\pgfpathlineto{\pgfqpoint{11.860365in}{7.650866in}}%
\pgfpathlineto{\pgfqpoint{11.775883in}{7.650866in}}%
\pgfpathlineto{\pgfqpoint{11.775883in}{7.566383in}}%
\pgfpathclose%
\pgfusepath{stroke}%
\end{pgfscope}%
\begin{pgfscope}%
\pgfpathrectangle{\pgfqpoint{8.036091in}{4.050000in}}{\pgfqpoint{5.400000in}{5.400000in}}%
\pgfusepath{clip}%
\pgfsetbuttcap%
\pgfsetroundjoin%
\pgfsetlinewidth{1.003750pt}%
\definecolor{currentstroke}{rgb}{0.541176,0.380392,0.658824}%
\pgfsetstrokecolor{currentstroke}%
\pgfsetstrokeopacity{0.900000}%
\pgfsetdash{}{0pt}%
\pgfpathmoveto{\pgfqpoint{12.084415in}{7.419689in}}%
\pgfpathlineto{\pgfqpoint{12.168898in}{7.419689in}}%
\pgfpathlineto{\pgfqpoint{12.168898in}{7.504172in}}%
\pgfpathlineto{\pgfqpoint{12.084415in}{7.504172in}}%
\pgfpathlineto{\pgfqpoint{12.084415in}{7.419689in}}%
\pgfpathclose%
\pgfusepath{stroke}%
\end{pgfscope}%
\begin{pgfscope}%
\pgfpathrectangle{\pgfqpoint{8.036091in}{4.050000in}}{\pgfqpoint{5.400000in}{5.400000in}}%
\pgfusepath{clip}%
\pgfsetbuttcap%
\pgfsetroundjoin%
\pgfsetlinewidth{1.003750pt}%
\definecolor{currentstroke}{rgb}{0.541176,0.380392,0.658824}%
\pgfsetstrokecolor{currentstroke}%
\pgfsetstrokeopacity{0.900000}%
\pgfsetdash{}{0pt}%
\pgfpathmoveto{\pgfqpoint{11.862615in}{8.692214in}}%
\pgfpathlineto{\pgfqpoint{11.947098in}{8.692214in}}%
\pgfpathlineto{\pgfqpoint{11.947098in}{8.776697in}}%
\pgfpathlineto{\pgfqpoint{11.862615in}{8.776697in}}%
\pgfpathlineto{\pgfqpoint{11.862615in}{8.692214in}}%
\pgfpathclose%
\pgfusepath{stroke}%
\end{pgfscope}%
\begin{pgfscope}%
\pgfpathrectangle{\pgfqpoint{8.036091in}{4.050000in}}{\pgfqpoint{5.400000in}{5.400000in}}%
\pgfusepath{clip}%
\pgfsetbuttcap%
\pgfsetroundjoin%
\pgfsetlinewidth{1.003750pt}%
\definecolor{currentstroke}{rgb}{0.541176,0.380392,0.658824}%
\pgfsetstrokecolor{currentstroke}%
\pgfsetstrokeopacity{0.900000}%
\pgfsetdash{}{0pt}%
\pgfpathmoveto{\pgfqpoint{11.406559in}{7.915907in}}%
\pgfpathlineto{\pgfqpoint{11.491042in}{7.915907in}}%
\pgfpathlineto{\pgfqpoint{11.491042in}{8.000390in}}%
\pgfpathlineto{\pgfqpoint{11.406559in}{8.000390in}}%
\pgfpathlineto{\pgfqpoint{11.406559in}{7.915907in}}%
\pgfpathclose%
\pgfusepath{stroke}%
\end{pgfscope}%
\begin{pgfscope}%
\pgfpathrectangle{\pgfqpoint{8.036091in}{4.050000in}}{\pgfqpoint{5.400000in}{5.400000in}}%
\pgfusepath{clip}%
\pgfsetbuttcap%
\pgfsetroundjoin%
\pgfsetlinewidth{1.003750pt}%
\definecolor{currentstroke}{rgb}{0.541176,0.380392,0.658824}%
\pgfsetstrokecolor{currentstroke}%
\pgfsetstrokeopacity{0.900000}%
\pgfsetdash{}{0pt}%
\pgfpathmoveto{\pgfqpoint{12.357139in}{8.160043in}}%
\pgfpathlineto{\pgfqpoint{12.441622in}{8.160043in}}%
\pgfpathlineto{\pgfqpoint{12.441622in}{8.244525in}}%
\pgfpathlineto{\pgfqpoint{12.357139in}{8.244525in}}%
\pgfpathlineto{\pgfqpoint{12.357139in}{8.160043in}}%
\pgfpathclose%
\pgfusepath{stroke}%
\end{pgfscope}%
\begin{pgfscope}%
\pgfpathrectangle{\pgfqpoint{8.036091in}{4.050000in}}{\pgfqpoint{5.400000in}{5.400000in}}%
\pgfusepath{clip}%
\pgfsetbuttcap%
\pgfsetroundjoin%
\pgfsetlinewidth{1.003750pt}%
\definecolor{currentstroke}{rgb}{0.541176,0.380392,0.658824}%
\pgfsetstrokecolor{currentstroke}%
\pgfsetstrokeopacity{0.900000}%
\pgfsetdash{}{0pt}%
\pgfpathmoveto{\pgfqpoint{11.025656in}{7.789639in}}%
\pgfpathlineto{\pgfqpoint{11.110139in}{7.789639in}}%
\pgfpathlineto{\pgfqpoint{11.110139in}{7.874122in}}%
\pgfpathlineto{\pgfqpoint{11.025656in}{7.874122in}}%
\pgfpathlineto{\pgfqpoint{11.025656in}{7.789639in}}%
\pgfpathclose%
\pgfusepath{stroke}%
\end{pgfscope}%
\begin{pgfscope}%
\pgfpathrectangle{\pgfqpoint{8.036091in}{4.050000in}}{\pgfqpoint{5.400000in}{5.400000in}}%
\pgfusepath{clip}%
\pgfsetbuttcap%
\pgfsetroundjoin%
\pgfsetlinewidth{1.003750pt}%
\definecolor{currentstroke}{rgb}{0.541176,0.380392,0.658824}%
\pgfsetstrokecolor{currentstroke}%
\pgfsetstrokeopacity{0.900000}%
\pgfsetdash{}{0pt}%
\pgfpathmoveto{\pgfqpoint{10.532962in}{6.928721in}}%
\pgfpathlineto{\pgfqpoint{10.617444in}{6.928721in}}%
\pgfpathlineto{\pgfqpoint{10.617444in}{7.013204in}}%
\pgfpathlineto{\pgfqpoint{10.532962in}{7.013204in}}%
\pgfpathlineto{\pgfqpoint{10.532962in}{6.928721in}}%
\pgfpathclose%
\pgfusepath{stroke}%
\end{pgfscope}%
\begin{pgfscope}%
\pgfpathrectangle{\pgfqpoint{8.036091in}{4.050000in}}{\pgfqpoint{5.400000in}{5.400000in}}%
\pgfusepath{clip}%
\pgfsetbuttcap%
\pgfsetroundjoin%
\pgfsetlinewidth{1.003750pt}%
\definecolor{currentstroke}{rgb}{0.541176,0.380392,0.658824}%
\pgfsetstrokecolor{currentstroke}%
\pgfsetstrokeopacity{0.900000}%
\pgfsetdash{}{0pt}%
\pgfpathmoveto{\pgfqpoint{9.407253in}{5.539818in}}%
\pgfpathlineto{\pgfqpoint{9.491736in}{5.539818in}}%
\pgfpathlineto{\pgfqpoint{9.491736in}{5.624301in}}%
\pgfpathlineto{\pgfqpoint{9.407253in}{5.624301in}}%
\pgfpathlineto{\pgfqpoint{9.407253in}{5.539818in}}%
\pgfpathclose%
\pgfusepath{stroke}%
\end{pgfscope}%
\begin{pgfscope}%
\pgfpathrectangle{\pgfqpoint{8.036091in}{4.050000in}}{\pgfqpoint{5.400000in}{5.400000in}}%
\pgfusepath{clip}%
\pgfsetbuttcap%
\pgfsetroundjoin%
\pgfsetlinewidth{1.003750pt}%
\definecolor{currentstroke}{rgb}{0.541176,0.380392,0.658824}%
\pgfsetstrokecolor{currentstroke}%
\pgfsetstrokeopacity{0.900000}%
\pgfsetdash{}{0pt}%
\pgfpathmoveto{\pgfqpoint{10.981574in}{6.942572in}}%
\pgfpathlineto{\pgfqpoint{11.066057in}{6.942572in}}%
\pgfpathlineto{\pgfqpoint{11.066057in}{7.027055in}}%
\pgfpathlineto{\pgfqpoint{10.981574in}{7.027055in}}%
\pgfpathlineto{\pgfqpoint{10.981574in}{6.942572in}}%
\pgfpathclose%
\pgfusepath{stroke}%
\end{pgfscope}%
\begin{pgfscope}%
\pgfpathrectangle{\pgfqpoint{8.036091in}{4.050000in}}{\pgfqpoint{5.400000in}{5.400000in}}%
\pgfusepath{clip}%
\pgfsetbuttcap%
\pgfsetroundjoin%
\pgfsetlinewidth{1.003750pt}%
\definecolor{currentstroke}{rgb}{0.541176,0.380392,0.658824}%
\pgfsetstrokecolor{currentstroke}%
\pgfsetstrokeopacity{0.900000}%
\pgfsetdash{}{0pt}%
\pgfpathmoveto{\pgfqpoint{11.964854in}{7.561281in}}%
\pgfpathlineto{\pgfqpoint{12.049336in}{7.561281in}}%
\pgfpathlineto{\pgfqpoint{12.049336in}{7.645764in}}%
\pgfpathlineto{\pgfqpoint{11.964854in}{7.645764in}}%
\pgfpathlineto{\pgfqpoint{11.964854in}{7.561281in}}%
\pgfpathclose%
\pgfusepath{stroke}%
\end{pgfscope}%
\begin{pgfscope}%
\pgfpathrectangle{\pgfqpoint{8.036091in}{4.050000in}}{\pgfqpoint{5.400000in}{5.400000in}}%
\pgfusepath{clip}%
\pgfsetbuttcap%
\pgfsetroundjoin%
\pgfsetlinewidth{1.003750pt}%
\definecolor{currentstroke}{rgb}{0.541176,0.380392,0.658824}%
\pgfsetstrokecolor{currentstroke}%
\pgfsetstrokeopacity{0.900000}%
\pgfsetdash{}{0pt}%
\pgfpathmoveto{\pgfqpoint{10.461453in}{5.829290in}}%
\pgfpathlineto{\pgfqpoint{10.545936in}{5.829290in}}%
\pgfpathlineto{\pgfqpoint{10.545936in}{5.913773in}}%
\pgfpathlineto{\pgfqpoint{10.461453in}{5.913773in}}%
\pgfpathlineto{\pgfqpoint{10.461453in}{5.829290in}}%
\pgfpathclose%
\pgfusepath{stroke}%
\end{pgfscope}%
\begin{pgfscope}%
\pgfpathrectangle{\pgfqpoint{8.036091in}{4.050000in}}{\pgfqpoint{5.400000in}{5.400000in}}%
\pgfusepath{clip}%
\pgfsetbuttcap%
\pgfsetroundjoin%
\pgfsetlinewidth{1.003750pt}%
\definecolor{currentstroke}{rgb}{0.541176,0.380392,0.658824}%
\pgfsetstrokecolor{currentstroke}%
\pgfsetstrokeopacity{0.900000}%
\pgfsetdash{}{0pt}%
\pgfpathmoveto{\pgfqpoint{10.428543in}{5.561768in}}%
\pgfpathlineto{\pgfqpoint{10.513025in}{5.561768in}}%
\pgfpathlineto{\pgfqpoint{10.513025in}{5.646251in}}%
\pgfpathlineto{\pgfqpoint{10.428543in}{5.646251in}}%
\pgfpathlineto{\pgfqpoint{10.428543in}{5.561768in}}%
\pgfpathclose%
\pgfusepath{stroke}%
\end{pgfscope}%
\begin{pgfscope}%
\pgfpathrectangle{\pgfqpoint{8.036091in}{4.050000in}}{\pgfqpoint{5.400000in}{5.400000in}}%
\pgfusepath{clip}%
\pgfsetbuttcap%
\pgfsetroundjoin%
\pgfsetlinewidth{1.003750pt}%
\definecolor{currentstroke}{rgb}{0.541176,0.380392,0.658824}%
\pgfsetstrokecolor{currentstroke}%
\pgfsetstrokeopacity{0.900000}%
\pgfsetdash{}{0pt}%
\pgfpathmoveto{\pgfqpoint{10.134321in}{5.805930in}}%
\pgfpathlineto{\pgfqpoint{10.218804in}{5.805930in}}%
\pgfpathlineto{\pgfqpoint{10.218804in}{5.890413in}}%
\pgfpathlineto{\pgfqpoint{10.134321in}{5.890413in}}%
\pgfpathlineto{\pgfqpoint{10.134321in}{5.805930in}}%
\pgfpathclose%
\pgfusepath{stroke}%
\end{pgfscope}%
\begin{pgfscope}%
\pgfpathrectangle{\pgfqpoint{8.036091in}{4.050000in}}{\pgfqpoint{5.400000in}{5.400000in}}%
\pgfusepath{clip}%
\pgfsetbuttcap%
\pgfsetroundjoin%
\pgfsetlinewidth{1.003750pt}%
\definecolor{currentstroke}{rgb}{0.541176,0.380392,0.658824}%
\pgfsetstrokecolor{currentstroke}%
\pgfsetstrokeopacity{0.900000}%
\pgfsetdash{}{0pt}%
\pgfpathmoveto{\pgfqpoint{9.536846in}{5.370125in}}%
\pgfpathlineto{\pgfqpoint{9.621329in}{5.370125in}}%
\pgfpathlineto{\pgfqpoint{9.621329in}{5.454608in}}%
\pgfpathlineto{\pgfqpoint{9.536846in}{5.454608in}}%
\pgfpathlineto{\pgfqpoint{9.536846in}{5.370125in}}%
\pgfpathclose%
\pgfusepath{stroke}%
\end{pgfscope}%
\begin{pgfscope}%
\pgfpathrectangle{\pgfqpoint{8.036091in}{4.050000in}}{\pgfqpoint{5.400000in}{5.400000in}}%
\pgfusepath{clip}%
\pgfsetbuttcap%
\pgfsetroundjoin%
\pgfsetlinewidth{1.003750pt}%
\definecolor{currentstroke}{rgb}{0.541176,0.380392,0.658824}%
\pgfsetstrokecolor{currentstroke}%
\pgfsetstrokeopacity{0.900000}%
\pgfsetdash{}{0pt}%
\pgfpathmoveto{\pgfqpoint{11.689358in}{7.788481in}}%
\pgfpathlineto{\pgfqpoint{11.773841in}{7.788481in}}%
\pgfpathlineto{\pgfqpoint{11.773841in}{7.872963in}}%
\pgfpathlineto{\pgfqpoint{11.689358in}{7.872963in}}%
\pgfpathlineto{\pgfqpoint{11.689358in}{7.788481in}}%
\pgfpathclose%
\pgfusepath{stroke}%
\end{pgfscope}%
\begin{pgfscope}%
\pgfpathrectangle{\pgfqpoint{8.036091in}{4.050000in}}{\pgfqpoint{5.400000in}{5.400000in}}%
\pgfusepath{clip}%
\pgfsetbuttcap%
\pgfsetroundjoin%
\pgfsetlinewidth{1.003750pt}%
\definecolor{currentstroke}{rgb}{0.541176,0.380392,0.658824}%
\pgfsetstrokecolor{currentstroke}%
\pgfsetstrokeopacity{0.900000}%
\pgfsetdash{}{0pt}%
\pgfpathmoveto{\pgfqpoint{11.279108in}{7.072169in}}%
\pgfpathlineto{\pgfqpoint{11.363591in}{7.072169in}}%
\pgfpathlineto{\pgfqpoint{11.363591in}{7.156652in}}%
\pgfpathlineto{\pgfqpoint{11.279108in}{7.156652in}}%
\pgfpathlineto{\pgfqpoint{11.279108in}{7.072169in}}%
\pgfpathclose%
\pgfusepath{stroke}%
\end{pgfscope}%
\begin{pgfscope}%
\pgfpathrectangle{\pgfqpoint{8.036091in}{4.050000in}}{\pgfqpoint{5.400000in}{5.400000in}}%
\pgfusepath{clip}%
\pgfsetbuttcap%
\pgfsetroundjoin%
\pgfsetlinewidth{1.003750pt}%
\definecolor{currentstroke}{rgb}{0.541176,0.380392,0.658824}%
\pgfsetstrokecolor{currentstroke}%
\pgfsetstrokeopacity{0.900000}%
\pgfsetdash{}{0pt}%
\pgfpathmoveto{\pgfqpoint{11.035566in}{8.400758in}}%
\pgfpathlineto{\pgfqpoint{11.120049in}{8.400758in}}%
\pgfpathlineto{\pgfqpoint{11.120049in}{8.485240in}}%
\pgfpathlineto{\pgfqpoint{11.035566in}{8.485240in}}%
\pgfpathlineto{\pgfqpoint{11.035566in}{8.400758in}}%
\pgfpathclose%
\pgfusepath{stroke}%
\end{pgfscope}%
\begin{pgfscope}%
\pgfpathrectangle{\pgfqpoint{8.036091in}{4.050000in}}{\pgfqpoint{5.400000in}{5.400000in}}%
\pgfusepath{clip}%
\pgfsetbuttcap%
\pgfsetroundjoin%
\pgfsetlinewidth{1.003750pt}%
\definecolor{currentstroke}{rgb}{0.541176,0.380392,0.658824}%
\pgfsetstrokecolor{currentstroke}%
\pgfsetstrokeopacity{0.900000}%
\pgfsetdash{}{0pt}%
\pgfpathmoveto{\pgfqpoint{10.723791in}{5.866316in}}%
\pgfpathlineto{\pgfqpoint{10.808274in}{5.866316in}}%
\pgfpathlineto{\pgfqpoint{10.808274in}{5.950799in}}%
\pgfpathlineto{\pgfqpoint{10.723791in}{5.950799in}}%
\pgfpathlineto{\pgfqpoint{10.723791in}{5.866316in}}%
\pgfpathclose%
\pgfusepath{stroke}%
\end{pgfscope}%
\begin{pgfscope}%
\pgfpathrectangle{\pgfqpoint{8.036091in}{4.050000in}}{\pgfqpoint{5.400000in}{5.400000in}}%
\pgfusepath{clip}%
\pgfsetbuttcap%
\pgfsetroundjoin%
\pgfsetlinewidth{1.003750pt}%
\definecolor{currentstroke}{rgb}{0.541176,0.380392,0.658824}%
\pgfsetstrokecolor{currentstroke}%
\pgfsetstrokeopacity{0.900000}%
\pgfsetdash{}{0pt}%
\pgfpathmoveto{\pgfqpoint{10.120554in}{4.835232in}}%
\pgfpathlineto{\pgfqpoint{10.205037in}{4.835232in}}%
\pgfpathlineto{\pgfqpoint{10.205037in}{4.919714in}}%
\pgfpathlineto{\pgfqpoint{10.120554in}{4.919714in}}%
\pgfpathlineto{\pgfqpoint{10.120554in}{4.835232in}}%
\pgfpathclose%
\pgfusepath{stroke}%
\end{pgfscope}%
\begin{pgfscope}%
\pgfpathrectangle{\pgfqpoint{8.036091in}{4.050000in}}{\pgfqpoint{5.400000in}{5.400000in}}%
\pgfusepath{clip}%
\pgfsetbuttcap%
\pgfsetroundjoin%
\pgfsetlinewidth{1.003750pt}%
\definecolor{currentstroke}{rgb}{0.541176,0.380392,0.658824}%
\pgfsetstrokecolor{currentstroke}%
\pgfsetstrokeopacity{0.900000}%
\pgfsetdash{}{0pt}%
\pgfpathmoveto{\pgfqpoint{10.485859in}{6.483275in}}%
\pgfpathlineto{\pgfqpoint{10.570342in}{6.483275in}}%
\pgfpathlineto{\pgfqpoint{10.570342in}{6.567758in}}%
\pgfpathlineto{\pgfqpoint{10.485859in}{6.567758in}}%
\pgfpathlineto{\pgfqpoint{10.485859in}{6.483275in}}%
\pgfpathclose%
\pgfusepath{stroke}%
\end{pgfscope}%
\begin{pgfscope}%
\pgfpathrectangle{\pgfqpoint{8.036091in}{4.050000in}}{\pgfqpoint{5.400000in}{5.400000in}}%
\pgfusepath{clip}%
\pgfsetbuttcap%
\pgfsetroundjoin%
\pgfsetlinewidth{1.003750pt}%
\definecolor{currentstroke}{rgb}{0.541176,0.380392,0.658824}%
\pgfsetstrokecolor{currentstroke}%
\pgfsetstrokeopacity{0.900000}%
\pgfsetdash{}{0pt}%
\pgfpathmoveto{\pgfqpoint{9.766245in}{5.048133in}}%
\pgfpathlineto{\pgfqpoint{9.850728in}{5.048133in}}%
\pgfpathlineto{\pgfqpoint{9.850728in}{5.132616in}}%
\pgfpathlineto{\pgfqpoint{9.766245in}{5.132616in}}%
\pgfpathlineto{\pgfqpoint{9.766245in}{5.048133in}}%
\pgfpathclose%
\pgfusepath{stroke}%
\end{pgfscope}%
\begin{pgfscope}%
\pgfpathrectangle{\pgfqpoint{8.036091in}{4.050000in}}{\pgfqpoint{5.400000in}{5.400000in}}%
\pgfusepath{clip}%
\pgfsetbuttcap%
\pgfsetroundjoin%
\pgfsetlinewidth{1.003750pt}%
\definecolor{currentstroke}{rgb}{0.541176,0.380392,0.658824}%
\pgfsetstrokecolor{currentstroke}%
\pgfsetstrokeopacity{0.900000}%
\pgfsetdash{}{0pt}%
\pgfpathmoveto{\pgfqpoint{10.310462in}{6.590923in}}%
\pgfpathlineto{\pgfqpoint{10.394945in}{6.590923in}}%
\pgfpathlineto{\pgfqpoint{10.394945in}{6.675406in}}%
\pgfpathlineto{\pgfqpoint{10.310462in}{6.675406in}}%
\pgfpathlineto{\pgfqpoint{10.310462in}{6.590923in}}%
\pgfpathclose%
\pgfusepath{stroke}%
\end{pgfscope}%
\begin{pgfscope}%
\pgfpathrectangle{\pgfqpoint{8.036091in}{4.050000in}}{\pgfqpoint{5.400000in}{5.400000in}}%
\pgfusepath{clip}%
\pgfsetbuttcap%
\pgfsetroundjoin%
\pgfsetlinewidth{1.003750pt}%
\definecolor{currentstroke}{rgb}{0.541176,0.380392,0.658824}%
\pgfsetstrokecolor{currentstroke}%
\pgfsetstrokeopacity{0.900000}%
\pgfsetdash{}{0pt}%
\pgfpathmoveto{\pgfqpoint{10.295862in}{6.765385in}}%
\pgfpathlineto{\pgfqpoint{10.380345in}{6.765385in}}%
\pgfpathlineto{\pgfqpoint{10.380345in}{6.849867in}}%
\pgfpathlineto{\pgfqpoint{10.295862in}{6.849867in}}%
\pgfpathlineto{\pgfqpoint{10.295862in}{6.765385in}}%
\pgfpathclose%
\pgfusepath{stroke}%
\end{pgfscope}%
\begin{pgfscope}%
\pgfpathrectangle{\pgfqpoint{8.036091in}{4.050000in}}{\pgfqpoint{5.400000in}{5.400000in}}%
\pgfusepath{clip}%
\pgfsetbuttcap%
\pgfsetroundjoin%
\pgfsetlinewidth{1.003750pt}%
\definecolor{currentstroke}{rgb}{0.768627,0.603922,0.000000}%
\pgfsetstrokecolor{currentstroke}%
\pgfsetstrokeopacity{0.900000}%
\pgfsetdash{}{0pt}%
\pgfpathmoveto{\pgfqpoint{12.207922in}{8.256200in}}%
\pgfpathlineto{\pgfqpoint{12.292405in}{8.256200in}}%
\pgfpathlineto{\pgfqpoint{12.292405in}{8.340682in}}%
\pgfpathlineto{\pgfqpoint{12.207922in}{8.340682in}}%
\pgfpathlineto{\pgfqpoint{12.207922in}{8.256200in}}%
\pgfpathclose%
\pgfusepath{stroke}%
\end{pgfscope}%
\begin{pgfscope}%
\pgfpathrectangle{\pgfqpoint{8.036091in}{4.050000in}}{\pgfqpoint{5.400000in}{5.400000in}}%
\pgfusepath{clip}%
\pgfsetbuttcap%
\pgfsetroundjoin%
\pgfsetlinewidth{1.003750pt}%
\definecolor{currentstroke}{rgb}{0.768627,0.603922,0.000000}%
\pgfsetstrokecolor{currentstroke}%
\pgfsetstrokeopacity{0.900000}%
\pgfsetdash{}{0pt}%
\pgfpathmoveto{\pgfqpoint{12.237441in}{8.338198in}}%
\pgfpathlineto{\pgfqpoint{12.321923in}{8.338198in}}%
\pgfpathlineto{\pgfqpoint{12.321923in}{8.422681in}}%
\pgfpathlineto{\pgfqpoint{12.237441in}{8.422681in}}%
\pgfpathlineto{\pgfqpoint{12.237441in}{8.338198in}}%
\pgfpathclose%
\pgfusepath{stroke}%
\end{pgfscope}%
\begin{pgfscope}%
\pgfpathrectangle{\pgfqpoint{8.036091in}{4.050000in}}{\pgfqpoint{5.400000in}{5.400000in}}%
\pgfusepath{clip}%
\pgfsetbuttcap%
\pgfsetroundjoin%
\pgfsetlinewidth{1.003750pt}%
\definecolor{currentstroke}{rgb}{0.768627,0.603922,0.000000}%
\pgfsetstrokecolor{currentstroke}%
\pgfsetstrokeopacity{0.900000}%
\pgfsetdash{}{0pt}%
\pgfpathmoveto{\pgfqpoint{10.453687in}{6.766780in}}%
\pgfpathlineto{\pgfqpoint{10.538170in}{6.766780in}}%
\pgfpathlineto{\pgfqpoint{10.538170in}{6.851263in}}%
\pgfpathlineto{\pgfqpoint{10.453687in}{6.851263in}}%
\pgfpathlineto{\pgfqpoint{10.453687in}{6.766780in}}%
\pgfpathclose%
\pgfusepath{stroke}%
\end{pgfscope}%
\begin{pgfscope}%
\pgfpathrectangle{\pgfqpoint{8.036091in}{4.050000in}}{\pgfqpoint{5.400000in}{5.400000in}}%
\pgfusepath{clip}%
\pgfsetbuttcap%
\pgfsetroundjoin%
\pgfsetlinewidth{1.003750pt}%
\definecolor{currentstroke}{rgb}{0.768627,0.603922,0.000000}%
\pgfsetstrokecolor{currentstroke}%
\pgfsetstrokeopacity{0.900000}%
\pgfsetdash{}{0pt}%
\pgfpathmoveto{\pgfqpoint{10.327966in}{6.764729in}}%
\pgfpathlineto{\pgfqpoint{10.412449in}{6.764729in}}%
\pgfpathlineto{\pgfqpoint{10.412449in}{6.849212in}}%
\pgfpathlineto{\pgfqpoint{10.327966in}{6.849212in}}%
\pgfpathlineto{\pgfqpoint{10.327966in}{6.764729in}}%
\pgfpathclose%
\pgfusepath{stroke}%
\end{pgfscope}%
\begin{pgfscope}%
\pgfpathrectangle{\pgfqpoint{8.036091in}{4.050000in}}{\pgfqpoint{5.400000in}{5.400000in}}%
\pgfusepath{clip}%
\pgfsetbuttcap%
\pgfsetroundjoin%
\pgfsetlinewidth{1.003750pt}%
\definecolor{currentstroke}{rgb}{0.200000,0.133333,0.533333}%
\pgfsetstrokecolor{currentstroke}%
\pgfsetstrokeopacity{0.900000}%
\pgfsetdash{}{0pt}%
\pgfpathmoveto{\pgfqpoint{9.568199in}{5.017436in}}%
\pgfpathlineto{\pgfqpoint{9.652682in}{5.017436in}}%
\pgfpathlineto{\pgfqpoint{9.652682in}{5.101919in}}%
\pgfpathlineto{\pgfqpoint{9.568199in}{5.101919in}}%
\pgfpathlineto{\pgfqpoint{9.568199in}{5.017436in}}%
\pgfpathclose%
\pgfusepath{stroke}%
\end{pgfscope}%
\begin{pgfscope}%
\pgfpathrectangle{\pgfqpoint{8.036091in}{4.050000in}}{\pgfqpoint{5.400000in}{5.400000in}}%
\pgfusepath{clip}%
\pgfsetbuttcap%
\pgfsetroundjoin%
\pgfsetlinewidth{1.003750pt}%
\definecolor{currentstroke}{rgb}{0.200000,0.133333,0.533333}%
\pgfsetstrokecolor{currentstroke}%
\pgfsetstrokeopacity{0.900000}%
\pgfsetdash{}{0pt}%
\pgfpathmoveto{\pgfqpoint{9.578626in}{5.079553in}}%
\pgfpathlineto{\pgfqpoint{9.663109in}{5.079553in}}%
\pgfpathlineto{\pgfqpoint{9.663109in}{5.164036in}}%
\pgfpathlineto{\pgfqpoint{9.578626in}{5.164036in}}%
\pgfpathlineto{\pgfqpoint{9.578626in}{5.079553in}}%
\pgfpathclose%
\pgfusepath{stroke}%
\end{pgfscope}%
\begin{pgfscope}%
\pgfpathrectangle{\pgfqpoint{8.036091in}{4.050000in}}{\pgfqpoint{5.400000in}{5.400000in}}%
\pgfusepath{clip}%
\pgfsetbuttcap%
\pgfsetroundjoin%
\pgfsetlinewidth{1.003750pt}%
\definecolor{currentstroke}{rgb}{0.200000,0.133333,0.533333}%
\pgfsetstrokecolor{currentstroke}%
\pgfsetstrokeopacity{0.900000}%
\pgfsetdash{}{0pt}%
\pgfpathmoveto{\pgfqpoint{10.003467in}{5.737513in}}%
\pgfpathlineto{\pgfqpoint{10.087950in}{5.737513in}}%
\pgfpathlineto{\pgfqpoint{10.087950in}{5.821996in}}%
\pgfpathlineto{\pgfqpoint{10.003467in}{5.821996in}}%
\pgfpathlineto{\pgfqpoint{10.003467in}{5.737513in}}%
\pgfpathclose%
\pgfusepath{stroke}%
\end{pgfscope}%
\begin{pgfscope}%
\pgfpathrectangle{\pgfqpoint{8.036091in}{4.050000in}}{\pgfqpoint{5.400000in}{5.400000in}}%
\pgfusepath{clip}%
\pgfsetbuttcap%
\pgfsetroundjoin%
\pgfsetlinewidth{1.003750pt}%
\definecolor{currentstroke}{rgb}{0.176471,0.713725,0.643137}%
\pgfsetstrokecolor{currentstroke}%
\pgfsetstrokeopacity{0.900000}%
\pgfsetdash{}{0pt}%
\pgfpathmoveto{\pgfqpoint{12.067656in}{8.175036in}}%
\pgfpathlineto{\pgfqpoint{12.152139in}{8.175036in}}%
\pgfpathlineto{\pgfqpoint{12.152139in}{8.259519in}}%
\pgfpathlineto{\pgfqpoint{12.067656in}{8.259519in}}%
\pgfpathlineto{\pgfqpoint{12.067656in}{8.175036in}}%
\pgfpathclose%
\pgfusepath{stroke}%
\end{pgfscope}%
\begin{pgfscope}%
\pgfpathrectangle{\pgfqpoint{8.036091in}{4.050000in}}{\pgfqpoint{5.400000in}{5.400000in}}%
\pgfusepath{clip}%
\pgfsetbuttcap%
\pgfsetroundjoin%
\pgfsetlinewidth{1.003750pt}%
\definecolor{currentstroke}{rgb}{0.176471,0.713725,0.643137}%
\pgfsetstrokecolor{currentstroke}%
\pgfsetstrokeopacity{0.900000}%
\pgfsetdash{}{0pt}%
\pgfpathmoveto{\pgfqpoint{12.322979in}{8.804231in}}%
\pgfpathlineto{\pgfqpoint{12.407462in}{8.804231in}}%
\pgfpathlineto{\pgfqpoint{12.407462in}{8.888714in}}%
\pgfpathlineto{\pgfqpoint{12.322979in}{8.888714in}}%
\pgfpathlineto{\pgfqpoint{12.322979in}{8.804231in}}%
\pgfpathclose%
\pgfusepath{stroke}%
\end{pgfscope}%
\begin{pgfscope}%
\pgfpathrectangle{\pgfqpoint{8.036091in}{4.050000in}}{\pgfqpoint{5.400000in}{5.400000in}}%
\pgfusepath{clip}%
\pgfsetbuttcap%
\pgfsetroundjoin%
\pgfsetlinewidth{1.003750pt}%
\definecolor{currentstroke}{rgb}{0.176471,0.713725,0.643137}%
\pgfsetstrokecolor{currentstroke}%
\pgfsetstrokeopacity{0.900000}%
\pgfsetdash{}{0pt}%
\pgfpathmoveto{\pgfqpoint{11.872333in}{8.007255in}}%
\pgfpathlineto{\pgfqpoint{11.956815in}{8.007255in}}%
\pgfpathlineto{\pgfqpoint{11.956815in}{8.091737in}}%
\pgfpathlineto{\pgfqpoint{11.872333in}{8.091737in}}%
\pgfpathlineto{\pgfqpoint{11.872333in}{8.007255in}}%
\pgfpathclose%
\pgfusepath{stroke}%
\end{pgfscope}%
\begin{pgfscope}%
\pgfpathrectangle{\pgfqpoint{8.036091in}{4.050000in}}{\pgfqpoint{5.400000in}{5.400000in}}%
\pgfusepath{clip}%
\pgfsetbuttcap%
\pgfsetroundjoin%
\pgfsetlinewidth{1.003750pt}%
\definecolor{currentstroke}{rgb}{0.176471,0.713725,0.643137}%
\pgfsetstrokecolor{currentstroke}%
\pgfsetstrokeopacity{0.900000}%
\pgfsetdash{}{0pt}%
\pgfpathmoveto{\pgfqpoint{12.419266in}{8.461537in}}%
\pgfpathlineto{\pgfqpoint{12.503748in}{8.461537in}}%
\pgfpathlineto{\pgfqpoint{12.503748in}{8.546020in}}%
\pgfpathlineto{\pgfqpoint{12.419266in}{8.546020in}}%
\pgfpathlineto{\pgfqpoint{12.419266in}{8.461537in}}%
\pgfpathclose%
\pgfusepath{stroke}%
\end{pgfscope}%
\begin{pgfscope}%
\pgfpathrectangle{\pgfqpoint{8.036091in}{4.050000in}}{\pgfqpoint{5.400000in}{5.400000in}}%
\pgfusepath{clip}%
\pgfsetbuttcap%
\pgfsetroundjoin%
\pgfsetlinewidth{1.003750pt}%
\definecolor{currentstroke}{rgb}{0.176471,0.713725,0.643137}%
\pgfsetstrokecolor{currentstroke}%
\pgfsetstrokeopacity{0.900000}%
\pgfsetdash{}{0pt}%
\pgfpathmoveto{\pgfqpoint{10.113315in}{6.360203in}}%
\pgfpathlineto{\pgfqpoint{10.197798in}{6.360203in}}%
\pgfpathlineto{\pgfqpoint{10.197798in}{6.444686in}}%
\pgfpathlineto{\pgfqpoint{10.113315in}{6.444686in}}%
\pgfpathlineto{\pgfqpoint{10.113315in}{6.360203in}}%
\pgfpathclose%
\pgfusepath{stroke}%
\end{pgfscope}%
\begin{pgfscope}%
\pgfpathrectangle{\pgfqpoint{8.036091in}{4.050000in}}{\pgfqpoint{5.400000in}{5.400000in}}%
\pgfusepath{clip}%
\pgfsetbuttcap%
\pgfsetroundjoin%
\pgfsetlinewidth{1.003750pt}%
\definecolor{currentstroke}{rgb}{0.901961,0.196078,0.156863}%
\pgfsetstrokecolor{currentstroke}%
\pgfsetstrokeopacity{0.900000}%
\pgfsetdash{}{0pt}%
\pgfpathmoveto{\pgfqpoint{13.031084in}{8.661520in}}%
\pgfpathlineto{\pgfqpoint{13.115566in}{8.661520in}}%
\pgfpathlineto{\pgfqpoint{13.115566in}{8.746002in}}%
\pgfpathlineto{\pgfqpoint{13.031084in}{8.746002in}}%
\pgfpathlineto{\pgfqpoint{13.031084in}{8.661520in}}%
\pgfpathclose%
\pgfusepath{stroke}%
\end{pgfscope}%
\begin{pgfscope}%
\pgfpathrectangle{\pgfqpoint{8.036091in}{4.050000in}}{\pgfqpoint{5.400000in}{5.400000in}}%
\pgfusepath{clip}%
\pgfsetbuttcap%
\pgfsetroundjoin%
\pgfsetlinewidth{1.003750pt}%
\definecolor{currentstroke}{rgb}{0.498039,0.560784,0.678431}%
\pgfsetstrokecolor{currentstroke}%
\pgfsetstrokeopacity{0.900000}%
\pgfsetdash{}{0pt}%
\pgfpathmoveto{\pgfqpoint{11.944154in}{7.364987in}}%
\pgfpathlineto{\pgfqpoint{12.028637in}{7.364987in}}%
\pgfpathlineto{\pgfqpoint{12.028637in}{7.449470in}}%
\pgfpathlineto{\pgfqpoint{11.944154in}{7.449470in}}%
\pgfpathlineto{\pgfqpoint{11.944154in}{7.364987in}}%
\pgfpathclose%
\pgfusepath{stroke}%
\end{pgfscope}%
\begin{pgfscope}%
\pgfsetbuttcap%
\pgfsetroundjoin%
\pgfsetlinewidth{1.003750pt}%
\definecolor{currentstroke}{rgb}{0.498039,0.560784,0.678431}%
\pgfsetstrokecolor{currentstroke}%
\pgfsetstrokeopacity{0.900000}%
\pgfsetdash{}{0pt}%
\pgfpathmoveto{\pgfqpoint{7.993850in}{4.007759in}}%
\pgfpathlineto{\pgfqpoint{8.078332in}{4.007759in}}%
\pgfpathlineto{\pgfqpoint{8.078332in}{4.092241in}}%
\pgfpathlineto{\pgfqpoint{7.993850in}{4.092241in}}%
\pgfpathlineto{\pgfqpoint{7.993850in}{4.007759in}}%
\pgfpathclose%
\pgfusepath{stroke}%
\end{pgfscope}%
\begin{pgfscope}%
\pgfsetbuttcap%
\pgfsetroundjoin%
\definecolor{currentfill}{rgb}{0.709804,0.121569,0.270588}%
\pgfsetfillcolor{currentfill}%
\pgfsetfillopacity{0.900000}%
\pgfsetlinewidth{0.803000pt}%
\definecolor{currentstroke}{rgb}{0.709804,0.121569,0.270588}%
\pgfsetstrokecolor{currentstroke}%
\pgfsetstrokeopacity{0.900000}%
\pgfsetdash{}{0pt}%
\pgfsys@defobject{currentmarker}{\pgfqpoint{-0.042241in}{-0.042241in}}{\pgfqpoint{0.042241in}{0.042241in}}{%
\pgfpathmoveto{\pgfqpoint{-0.042241in}{-0.042241in}}%
\pgfpathlineto{\pgfqpoint{0.042241in}{-0.042241in}}%
\pgfpathlineto{\pgfqpoint{0.042241in}{0.042241in}}%
\pgfpathlineto{\pgfqpoint{-0.042241in}{0.042241in}}%
\pgfpathlineto{\pgfqpoint{-0.042241in}{-0.042241in}}%
\pgfpathclose%
\pgfusepath{stroke,fill}%
}%
\begin{pgfscope}%
\pgfsys@transformshift{13.188796in}{9.450000in}%
\pgfsys@useobject{currentmarker}{}%
\end{pgfscope}%
\end{pgfscope}%
\begin{pgfscope}%
\pgfsetbuttcap%
\pgfsetmiterjoin%
\definecolor{currentfill}{rgb}{1.000000,1.000000,1.000000}%
\pgfsetfillcolor{currentfill}%
\pgfsetfillopacity{0.970000}%
\pgfsetlinewidth{0.702625pt}%
\definecolor{currentstroke}{rgb}{0.796078,0.827451,0.850980}%
\pgfsetstrokecolor{currentstroke}%
\pgfsetstrokeopacity{0.970000}%
\pgfsetdash{}{0pt}%
\pgfpathmoveto{\pgfqpoint{8.225091in}{7.579333in}}%
\pgfpathlineto{\pgfqpoint{10.341966in}{7.579333in}}%
\pgfpathquadraticcurveto{\pgfqpoint{10.419883in}{7.579333in}}{\pgfqpoint{10.419883in}{7.657250in}}%
\pgfpathlineto{\pgfqpoint{10.419883in}{8.316000in}}%
\pgfpathquadraticcurveto{\pgfqpoint{10.419883in}{8.393917in}}{\pgfqpoint{10.341966in}{8.393917in}}%
\pgfpathlineto{\pgfqpoint{8.225091in}{8.393917in}}%
\pgfpathquadraticcurveto{\pgfqpoint{8.147174in}{8.393917in}}{\pgfqpoint{8.147174in}{8.316000in}}%
\pgfpathlineto{\pgfqpoint{8.147174in}{7.657250in}}%
\pgfpathquadraticcurveto{\pgfqpoint{8.147174in}{7.579333in}}{\pgfqpoint{8.225091in}{7.579333in}}%
\pgfpathlineto{\pgfqpoint{8.225091in}{7.579333in}}%
\pgfpathclose%
\pgfusepath{stroke,fill}%
\end{pgfscope}%
\begin{pgfscope}%
\pgfsetbuttcap%
\pgfsetroundjoin%
\definecolor{currentfill}{rgb}{0.090196,0.168627,0.227451}%
\pgfsetfillcolor{currentfill}%
\pgfsetlinewidth{0.000000pt}%
\definecolor{currentstroke}{rgb}{0.090196,0.168627,0.227451}%
\pgfsetstrokecolor{currentstroke}%
\pgfsetdash{}{0pt}%
\pgfpathmoveto{\pgfqpoint{8.300905in}{8.266105in}}%
\pgfpathlineto{\pgfqpoint{8.300905in}{8.252470in}}%
\pgfpathquadraticcurveto{\pgfqpoint{8.292958in}{8.256277in}}{\pgfqpoint{8.285897in}{8.258137in}}%
\pgfpathquadraticcurveto{\pgfqpoint{8.278836in}{8.260018in}}{\pgfqpoint{8.272261in}{8.260018in}}%
\pgfpathquadraticcurveto{\pgfqpoint{8.260862in}{8.260018in}}{\pgfqpoint{8.254664in}{8.255591in}}%
\pgfpathquadraticcurveto{\pgfqpoint{8.248466in}{8.251164in}}{\pgfqpoint{8.248466in}{8.242996in}}%
\pgfpathquadraticcurveto{\pgfqpoint{8.248466in}{8.236134in}}{\pgfqpoint{8.252583in}{8.232637in}}%
\pgfpathquadraticcurveto{\pgfqpoint{8.256700in}{8.229161in}}{\pgfqpoint{8.268189in}{8.227014in}}%
\pgfpathlineto{\pgfqpoint{8.276622in}{8.225288in}}%
\pgfpathquadraticcurveto{\pgfqpoint{8.292250in}{8.222299in}}{\pgfqpoint{8.299687in}{8.214796in}}%
\pgfpathquadraticcurveto{\pgfqpoint{8.307125in}{8.207292in}}{\pgfqpoint{8.307125in}{8.194719in}}%
\pgfpathquadraticcurveto{\pgfqpoint{8.307125in}{8.179689in}}{\pgfqpoint{8.297053in}{8.171941in}}%
\pgfpathquadraticcurveto{\pgfqpoint{8.287004in}{8.164194in}}{\pgfqpoint{8.267569in}{8.164194in}}%
\pgfpathquadraticcurveto{\pgfqpoint{8.260242in}{8.164194in}}{\pgfqpoint{8.251963in}{8.165854in}}%
\pgfpathquadraticcurveto{\pgfqpoint{8.243707in}{8.167514in}}{\pgfqpoint{8.234853in}{8.170768in}}%
\pgfpathlineto{\pgfqpoint{8.234853in}{8.185156in}}%
\pgfpathquadraticcurveto{\pgfqpoint{8.243353in}{8.180397in}}{\pgfqpoint{8.251521in}{8.177962in}}%
\pgfpathquadraticcurveto{\pgfqpoint{8.259689in}{8.175549in}}{\pgfqpoint{8.267569in}{8.175549in}}%
\pgfpathquadraticcurveto{\pgfqpoint{8.279522in}{8.175549in}}{\pgfqpoint{8.286030in}{8.180242in}}%
\pgfpathquadraticcurveto{\pgfqpoint{8.292538in}{8.184957in}}{\pgfqpoint{8.292538in}{8.193678in}}%
\pgfpathquadraticcurveto{\pgfqpoint{8.292538in}{8.201271in}}{\pgfqpoint{8.287867in}{8.205565in}}%
\pgfpathquadraticcurveto{\pgfqpoint{8.283196in}{8.209859in}}{\pgfqpoint{8.272549in}{8.212007in}}%
\pgfpathlineto{\pgfqpoint{8.264027in}{8.213667in}}%
\pgfpathquadraticcurveto{\pgfqpoint{8.248400in}{8.216766in}}{\pgfqpoint{8.241405in}{8.223406in}}%
\pgfpathquadraticcurveto{\pgfqpoint{8.234432in}{8.230047in}}{\pgfqpoint{8.234432in}{8.241889in}}%
\pgfpathquadraticcurveto{\pgfqpoint{8.234432in}{8.255591in}}{\pgfqpoint{8.244083in}{8.263471in}}%
\pgfpathquadraticcurveto{\pgfqpoint{8.253734in}{8.271352in}}{\pgfqpoint{8.270668in}{8.271352in}}%
\pgfpathquadraticcurveto{\pgfqpoint{8.277950in}{8.271352in}}{\pgfqpoint{8.285476in}{8.270023in}}%
\pgfpathquadraticcurveto{\pgfqpoint{8.293025in}{8.268717in}}{\pgfqpoint{8.300905in}{8.266105in}}%
\pgfpathlineto{\pgfqpoint{8.300905in}{8.266105in}}%
\pgfpathclose%
\pgfpathmoveto{\pgfqpoint{8.328375in}{8.243682in}}%
\pgfpathlineto{\pgfqpoint{8.341103in}{8.243682in}}%
\pgfpathlineto{\pgfqpoint{8.341103in}{8.166208in}}%
\pgfpathlineto{\pgfqpoint{8.328375in}{8.166208in}}%
\pgfpathlineto{\pgfqpoint{8.328375in}{8.243682in}}%
\pgfpathlineto{\pgfqpoint{8.328375in}{8.243682in}}%
\pgfpathclose%
\pgfpathmoveto{\pgfqpoint{8.328375in}{8.273853in}}%
\pgfpathlineto{\pgfqpoint{8.341103in}{8.273853in}}%
\pgfpathlineto{\pgfqpoint{8.341103in}{8.257716in}}%
\pgfpathlineto{\pgfqpoint{8.328375in}{8.257716in}}%
\pgfpathlineto{\pgfqpoint{8.328375in}{8.273853in}}%
\pgfpathlineto{\pgfqpoint{8.328375in}{8.273853in}}%
\pgfpathclose%
\pgfpathmoveto{\pgfqpoint{8.418709in}{8.205853in}}%
\pgfpathquadraticcurveto{\pgfqpoint{8.418709in}{8.219688in}}{\pgfqpoint{8.412998in}{8.227280in}}%
\pgfpathquadraticcurveto{\pgfqpoint{8.407310in}{8.234895in}}{\pgfqpoint{8.396995in}{8.234895in}}%
\pgfpathquadraticcurveto{\pgfqpoint{8.386768in}{8.234895in}}{\pgfqpoint{8.381057in}{8.227280in}}%
\pgfpathquadraticcurveto{\pgfqpoint{8.375346in}{8.219688in}}{\pgfqpoint{8.375346in}{8.205853in}}%
\pgfpathquadraticcurveto{\pgfqpoint{8.375346in}{8.192085in}}{\pgfqpoint{8.381057in}{8.184470in}}%
\pgfpathquadraticcurveto{\pgfqpoint{8.386768in}{8.176855in}}{\pgfqpoint{8.396995in}{8.176855in}}%
\pgfpathquadraticcurveto{\pgfqpoint{8.407310in}{8.176855in}}{\pgfqpoint{8.412998in}{8.184470in}}%
\pgfpathquadraticcurveto{\pgfqpoint{8.418709in}{8.192085in}}{\pgfqpoint{8.418709in}{8.205853in}}%
\pgfpathlineto{\pgfqpoint{8.418709in}{8.205853in}}%
\pgfpathclose%
\pgfpathmoveto{\pgfqpoint{8.431437in}{8.175815in}}%
\pgfpathquadraticcurveto{\pgfqpoint{8.431437in}{8.156048in}}{\pgfqpoint{8.422650in}{8.146397in}}%
\pgfpathquadraticcurveto{\pgfqpoint{8.413884in}{8.136746in}}{\pgfqpoint{8.395755in}{8.136746in}}%
\pgfpathquadraticcurveto{\pgfqpoint{8.389048in}{8.136746in}}{\pgfqpoint{8.383094in}{8.137742in}}%
\pgfpathquadraticcurveto{\pgfqpoint{8.377139in}{8.138738in}}{\pgfqpoint{8.371539in}{8.140819in}}%
\pgfpathlineto{\pgfqpoint{8.371539in}{8.153193in}}%
\pgfpathquadraticcurveto{\pgfqpoint{8.377139in}{8.150160in}}{\pgfqpoint{8.382607in}{8.148721in}}%
\pgfpathquadraticcurveto{\pgfqpoint{8.388074in}{8.147260in}}{\pgfqpoint{8.393741in}{8.147260in}}%
\pgfpathquadraticcurveto{\pgfqpoint{8.406269in}{8.147260in}}{\pgfqpoint{8.412489in}{8.153790in}}%
\pgfpathquadraticcurveto{\pgfqpoint{8.418709in}{8.160320in}}{\pgfqpoint{8.418709in}{8.173535in}}%
\pgfpathlineto{\pgfqpoint{8.418709in}{8.179844in}}%
\pgfpathquadraticcurveto{\pgfqpoint{8.414769in}{8.172982in}}{\pgfqpoint{8.408616in}{8.169595in}}%
\pgfpathquadraticcurveto{\pgfqpoint{8.402462in}{8.166208in}}{\pgfqpoint{8.393873in}{8.166208in}}%
\pgfpathquadraticcurveto{\pgfqpoint{8.379640in}{8.166208in}}{\pgfqpoint{8.370919in}{8.177055in}}%
\pgfpathquadraticcurveto{\pgfqpoint{8.362198in}{8.187923in}}{\pgfqpoint{8.362198in}{8.205853in}}%
\pgfpathquadraticcurveto{\pgfqpoint{8.362198in}{8.223827in}}{\pgfqpoint{8.370919in}{8.234673in}}%
\pgfpathquadraticcurveto{\pgfqpoint{8.379640in}{8.245542in}}{\pgfqpoint{8.393873in}{8.245542in}}%
\pgfpathquadraticcurveto{\pgfqpoint{8.402462in}{8.245542in}}{\pgfqpoint{8.408616in}{8.242155in}}%
\pgfpathquadraticcurveto{\pgfqpoint{8.414769in}{8.238768in}}{\pgfqpoint{8.418709in}{8.231928in}}%
\pgfpathlineto{\pgfqpoint{8.418709in}{8.243682in}}%
\pgfpathlineto{\pgfqpoint{8.431437in}{8.243682in}}%
\pgfpathlineto{\pgfqpoint{8.431437in}{8.175815in}}%
\pgfpathlineto{\pgfqpoint{8.431437in}{8.175815in}}%
\pgfpathclose%
\pgfpathmoveto{\pgfqpoint{8.522082in}{8.212980in}}%
\pgfpathlineto{\pgfqpoint{8.522082in}{8.166208in}}%
\pgfpathlineto{\pgfqpoint{8.509354in}{8.166208in}}%
\pgfpathlineto{\pgfqpoint{8.509354in}{8.212560in}}%
\pgfpathquadraticcurveto{\pgfqpoint{8.509354in}{8.223561in}}{\pgfqpoint{8.505060in}{8.229007in}}%
\pgfpathquadraticcurveto{\pgfqpoint{8.500765in}{8.234474in}}{\pgfqpoint{8.492199in}{8.234474in}}%
\pgfpathquadraticcurveto{\pgfqpoint{8.481884in}{8.234474in}}{\pgfqpoint{8.475929in}{8.227900in}}%
\pgfpathquadraticcurveto{\pgfqpoint{8.469975in}{8.221348in}}{\pgfqpoint{8.469975in}{8.209992in}}%
\pgfpathlineto{\pgfqpoint{8.469975in}{8.166208in}}%
\pgfpathlineto{\pgfqpoint{8.457181in}{8.166208in}}%
\pgfpathlineto{\pgfqpoint{8.457181in}{8.243682in}}%
\pgfpathlineto{\pgfqpoint{8.469975in}{8.243682in}}%
\pgfpathlineto{\pgfqpoint{8.469975in}{8.231641in}}%
\pgfpathquadraticcurveto{\pgfqpoint{8.474557in}{8.238635in}}{\pgfqpoint{8.480733in}{8.242089in}}%
\pgfpathquadraticcurveto{\pgfqpoint{8.486931in}{8.245542in}}{\pgfqpoint{8.495032in}{8.245542in}}%
\pgfpathquadraticcurveto{\pgfqpoint{8.508380in}{8.245542in}}{\pgfqpoint{8.515220in}{8.237285in}}%
\pgfpathquadraticcurveto{\pgfqpoint{8.522082in}{8.229029in}}{\pgfqpoint{8.522082in}{8.212980in}}%
\pgfpathlineto{\pgfqpoint{8.522082in}{8.212980in}}%
\pgfpathclose%
\pgfpathmoveto{\pgfqpoint{8.541030in}{8.210678in}}%
\pgfpathlineto{\pgfqpoint{8.578306in}{8.210678in}}%
\pgfpathlineto{\pgfqpoint{8.578306in}{8.199345in}}%
\pgfpathlineto{\pgfqpoint{8.541030in}{8.199345in}}%
\pgfpathlineto{\pgfqpoint{8.541030in}{8.210678in}}%
\pgfpathlineto{\pgfqpoint{8.541030in}{8.210678in}}%
\pgfpathclose%
\pgfpathmoveto{\pgfqpoint{8.643450in}{8.231796in}}%
\pgfpathquadraticcurveto{\pgfqpoint{8.641303in}{8.233035in}}{\pgfqpoint{8.638780in}{8.233611in}}%
\pgfpathquadraticcurveto{\pgfqpoint{8.636256in}{8.234208in}}{\pgfqpoint{8.633224in}{8.234208in}}%
\pgfpathquadraticcurveto{\pgfqpoint{8.622422in}{8.234208in}}{\pgfqpoint{8.616644in}{8.227191in}}%
\pgfpathquadraticcurveto{\pgfqpoint{8.610867in}{8.220174in}}{\pgfqpoint{8.610867in}{8.207026in}}%
\pgfpathlineto{\pgfqpoint{8.610867in}{8.166208in}}%
\pgfpathlineto{\pgfqpoint{8.598073in}{8.166208in}}%
\pgfpathlineto{\pgfqpoint{8.598073in}{8.243682in}}%
\pgfpathlineto{\pgfqpoint{8.610867in}{8.243682in}}%
\pgfpathlineto{\pgfqpoint{8.610867in}{8.231641in}}%
\pgfpathquadraticcurveto{\pgfqpoint{8.614896in}{8.238702in}}{\pgfqpoint{8.621315in}{8.242111in}}%
\pgfpathquadraticcurveto{\pgfqpoint{8.627756in}{8.245542in}}{\pgfqpoint{8.636965in}{8.245542in}}%
\pgfpathquadraticcurveto{\pgfqpoint{8.638271in}{8.245542in}}{\pgfqpoint{8.639864in}{8.245365in}}%
\pgfpathquadraticcurveto{\pgfqpoint{8.641458in}{8.245210in}}{\pgfqpoint{8.643384in}{8.244855in}}%
\pgfpathlineto{\pgfqpoint{8.643450in}{8.231796in}}%
\pgfpathlineto{\pgfqpoint{8.643450in}{8.231796in}}%
\pgfpathclose%
\pgfpathmoveto{\pgfqpoint{8.719950in}{8.208133in}}%
\pgfpathlineto{\pgfqpoint{8.719950in}{8.201913in}}%
\pgfpathlineto{\pgfqpoint{8.661424in}{8.201913in}}%
\pgfpathquadraticcurveto{\pgfqpoint{8.662265in}{8.188764in}}{\pgfqpoint{8.669349in}{8.181880in}}%
\pgfpathquadraticcurveto{\pgfqpoint{8.676432in}{8.174996in}}{\pgfqpoint{8.689094in}{8.174996in}}%
\pgfpathquadraticcurveto{\pgfqpoint{8.696420in}{8.174996in}}{\pgfqpoint{8.703304in}{8.176789in}}%
\pgfpathquadraticcurveto{\pgfqpoint{8.710189in}{8.178582in}}{\pgfqpoint{8.716984in}{8.182190in}}%
\pgfpathlineto{\pgfqpoint{8.716984in}{8.170148in}}%
\pgfpathquadraticcurveto{\pgfqpoint{8.710122in}{8.167249in}}{\pgfqpoint{8.702928in}{8.165721in}}%
\pgfpathquadraticcurveto{\pgfqpoint{8.695734in}{8.164194in}}{\pgfqpoint{8.688341in}{8.164194in}}%
\pgfpathquadraticcurveto{\pgfqpoint{8.669791in}{8.164194in}}{\pgfqpoint{8.658967in}{8.174974in}}%
\pgfpathquadraticcurveto{\pgfqpoint{8.648143in}{8.185776in}}{\pgfqpoint{8.648143in}{8.204193in}}%
\pgfpathquadraticcurveto{\pgfqpoint{8.648143in}{8.223207in}}{\pgfqpoint{8.658414in}{8.234363in}}%
\pgfpathquadraticcurveto{\pgfqpoint{8.668685in}{8.245542in}}{\pgfqpoint{8.686127in}{8.245542in}}%
\pgfpathquadraticcurveto{\pgfqpoint{8.701755in}{8.245542in}}{\pgfqpoint{8.710853in}{8.235470in}}%
\pgfpathquadraticcurveto{\pgfqpoint{8.719950in}{8.225421in}}{\pgfqpoint{8.719950in}{8.208133in}}%
\pgfpathlineto{\pgfqpoint{8.719950in}{8.208133in}}%
\pgfpathclose%
\pgfpathmoveto{\pgfqpoint{8.707222in}{8.211874in}}%
\pgfpathquadraticcurveto{\pgfqpoint{8.707090in}{8.222299in}}{\pgfqpoint{8.701379in}{8.228520in}}%
\pgfpathquadraticcurveto{\pgfqpoint{8.695668in}{8.234762in}}{\pgfqpoint{8.686260in}{8.234762in}}%
\pgfpathquadraticcurveto{\pgfqpoint{8.675613in}{8.234762in}}{\pgfqpoint{8.669216in}{8.228741in}}%
\pgfpathquadraticcurveto{\pgfqpoint{8.662819in}{8.222720in}}{\pgfqpoint{8.661845in}{8.211785in}}%
\pgfpathlineto{\pgfqpoint{8.707222in}{8.211874in}}%
\pgfpathlineto{\pgfqpoint{8.707222in}{8.211874in}}%
\pgfpathclose%
\pgfpathmoveto{\pgfqpoint{8.731726in}{8.243682in}}%
\pgfpathlineto{\pgfqpoint{8.745207in}{8.243682in}}%
\pgfpathlineto{\pgfqpoint{8.769423in}{8.178671in}}%
\pgfpathlineto{\pgfqpoint{8.793639in}{8.243682in}}%
\pgfpathlineto{\pgfqpoint{8.807120in}{8.243682in}}%
\pgfpathlineto{\pgfqpoint{8.778056in}{8.166208in}}%
\pgfpathlineto{\pgfqpoint{8.760768in}{8.166208in}}%
\pgfpathlineto{\pgfqpoint{8.731726in}{8.243682in}}%
\pgfpathlineto{\pgfqpoint{8.731726in}{8.243682in}}%
\pgfpathclose%
\pgfpathmoveto{\pgfqpoint{8.890969in}{8.208133in}}%
\pgfpathlineto{\pgfqpoint{8.890969in}{8.201913in}}%
\pgfpathlineto{\pgfqpoint{8.832442in}{8.201913in}}%
\pgfpathquadraticcurveto{\pgfqpoint{8.833284in}{8.188764in}}{\pgfqpoint{8.840367in}{8.181880in}}%
\pgfpathquadraticcurveto{\pgfqpoint{8.847450in}{8.174996in}}{\pgfqpoint{8.860112in}{8.174996in}}%
\pgfpathquadraticcurveto{\pgfqpoint{8.867439in}{8.174996in}}{\pgfqpoint{8.874323in}{8.176789in}}%
\pgfpathquadraticcurveto{\pgfqpoint{8.881207in}{8.178582in}}{\pgfqpoint{8.888002in}{8.182190in}}%
\pgfpathlineto{\pgfqpoint{8.888002in}{8.170148in}}%
\pgfpathquadraticcurveto{\pgfqpoint{8.881140in}{8.167249in}}{\pgfqpoint{8.873946in}{8.165721in}}%
\pgfpathquadraticcurveto{\pgfqpoint{8.866752in}{8.164194in}}{\pgfqpoint{8.859359in}{8.164194in}}%
\pgfpathquadraticcurveto{\pgfqpoint{8.840810in}{8.164194in}}{\pgfqpoint{8.829985in}{8.174974in}}%
\pgfpathquadraticcurveto{\pgfqpoint{8.819161in}{8.185776in}}{\pgfqpoint{8.819161in}{8.204193in}}%
\pgfpathquadraticcurveto{\pgfqpoint{8.819161in}{8.223207in}}{\pgfqpoint{8.829432in}{8.234363in}}%
\pgfpathquadraticcurveto{\pgfqpoint{8.839703in}{8.245542in}}{\pgfqpoint{8.857146in}{8.245542in}}%
\pgfpathquadraticcurveto{\pgfqpoint{8.872773in}{8.245542in}}{\pgfqpoint{8.881871in}{8.235470in}}%
\pgfpathquadraticcurveto{\pgfqpoint{8.890969in}{8.225421in}}{\pgfqpoint{8.890969in}{8.208133in}}%
\pgfpathlineto{\pgfqpoint{8.890969in}{8.208133in}}%
\pgfpathclose%
\pgfpathmoveto{\pgfqpoint{8.878241in}{8.211874in}}%
\pgfpathquadraticcurveto{\pgfqpoint{8.878108in}{8.222299in}}{\pgfqpoint{8.872397in}{8.228520in}}%
\pgfpathquadraticcurveto{\pgfqpoint{8.866686in}{8.234762in}}{\pgfqpoint{8.857278in}{8.234762in}}%
\pgfpathquadraticcurveto{\pgfqpoint{8.846631in}{8.234762in}}{\pgfqpoint{8.840234in}{8.228741in}}%
\pgfpathquadraticcurveto{\pgfqpoint{8.833837in}{8.222720in}}{\pgfqpoint{8.832863in}{8.211785in}}%
\pgfpathlineto{\pgfqpoint{8.878241in}{8.211874in}}%
\pgfpathlineto{\pgfqpoint{8.878241in}{8.211874in}}%
\pgfpathclose%
\pgfpathmoveto{\pgfqpoint{8.956755in}{8.231796in}}%
\pgfpathquadraticcurveto{\pgfqpoint{8.954608in}{8.233035in}}{\pgfqpoint{8.952084in}{8.233611in}}%
\pgfpathquadraticcurveto{\pgfqpoint{8.949561in}{8.234208in}}{\pgfqpoint{8.946528in}{8.234208in}}%
\pgfpathquadraticcurveto{\pgfqpoint{8.935726in}{8.234208in}}{\pgfqpoint{8.929949in}{8.227191in}}%
\pgfpathquadraticcurveto{\pgfqpoint{8.924172in}{8.220174in}}{\pgfqpoint{8.924172in}{8.207026in}}%
\pgfpathlineto{\pgfqpoint{8.924172in}{8.166208in}}%
\pgfpathlineto{\pgfqpoint{8.911377in}{8.166208in}}%
\pgfpathlineto{\pgfqpoint{8.911377in}{8.243682in}}%
\pgfpathlineto{\pgfqpoint{8.924172in}{8.243682in}}%
\pgfpathlineto{\pgfqpoint{8.924172in}{8.231641in}}%
\pgfpathquadraticcurveto{\pgfqpoint{8.928200in}{8.238702in}}{\pgfqpoint{8.934620in}{8.242111in}}%
\pgfpathquadraticcurveto{\pgfqpoint{8.941061in}{8.245542in}}{\pgfqpoint{8.950269in}{8.245542in}}%
\pgfpathquadraticcurveto{\pgfqpoint{8.951575in}{8.245542in}}{\pgfqpoint{8.953169in}{8.245365in}}%
\pgfpathquadraticcurveto{\pgfqpoint{8.954763in}{8.245210in}}{\pgfqpoint{8.956689in}{8.244855in}}%
\pgfpathlineto{\pgfqpoint{8.956755in}{8.231796in}}%
\pgfpathlineto{\pgfqpoint{8.956755in}{8.231796in}}%
\pgfpathclose%
\pgfpathmoveto{\pgfqpoint{9.019487in}{8.241402in}}%
\pgfpathlineto{\pgfqpoint{9.019487in}{8.229361in}}%
\pgfpathquadraticcurveto{\pgfqpoint{9.014108in}{8.232128in}}{\pgfqpoint{9.008286in}{8.233500in}}%
\pgfpathquadraticcurveto{\pgfqpoint{9.002487in}{8.234895in}}{\pgfqpoint{8.996245in}{8.234895in}}%
\pgfpathquadraticcurveto{\pgfqpoint{8.986771in}{8.234895in}}{\pgfqpoint{8.982034in}{8.231995in}}%
\pgfpathquadraticcurveto{\pgfqpoint{8.977297in}{8.229095in}}{\pgfqpoint{8.977297in}{8.223273in}}%
\pgfpathquadraticcurveto{\pgfqpoint{8.977297in}{8.218846in}}{\pgfqpoint{8.980683in}{8.216323in}}%
\pgfpathquadraticcurveto{\pgfqpoint{8.984070in}{8.213799in}}{\pgfqpoint{8.994319in}{8.211520in}}%
\pgfpathlineto{\pgfqpoint{8.998679in}{8.210546in}}%
\pgfpathquadraticcurveto{\pgfqpoint{9.012226in}{8.207646in}}{\pgfqpoint{9.017937in}{8.202355in}}%
\pgfpathquadraticcurveto{\pgfqpoint{9.023648in}{8.197065in}}{\pgfqpoint{9.023648in}{8.187591in}}%
\pgfpathquadraticcurveto{\pgfqpoint{9.023648in}{8.176789in}}{\pgfqpoint{9.015104in}{8.170480in}}%
\pgfpathquadraticcurveto{\pgfqpoint{9.006560in}{8.164194in}}{\pgfqpoint{8.991618in}{8.164194in}}%
\pgfpathquadraticcurveto{\pgfqpoint{8.985398in}{8.164194in}}{\pgfqpoint{8.978647in}{8.165411in}}%
\pgfpathquadraticcurveto{\pgfqpoint{8.971896in}{8.166629in}}{\pgfqpoint{8.964436in}{8.169042in}}%
\pgfpathlineto{\pgfqpoint{8.964436in}{8.182190in}}%
\pgfpathquadraticcurveto{\pgfqpoint{8.971497in}{8.178516in}}{\pgfqpoint{8.978337in}{8.176678in}}%
\pgfpathquadraticcurveto{\pgfqpoint{8.985177in}{8.174863in}}{\pgfqpoint{8.991906in}{8.174863in}}%
\pgfpathquadraticcurveto{\pgfqpoint{9.000893in}{8.174863in}}{\pgfqpoint{9.005719in}{8.177940in}}%
\pgfpathquadraticcurveto{\pgfqpoint{9.010566in}{8.181017in}}{\pgfqpoint{9.010566in}{8.186617in}}%
\pgfpathquadraticcurveto{\pgfqpoint{9.010566in}{8.191797in}}{\pgfqpoint{9.007069in}{8.194564in}}%
\pgfpathquadraticcurveto{\pgfqpoint{9.003594in}{8.197331in}}{\pgfqpoint{8.991751in}{8.199898in}}%
\pgfpathlineto{\pgfqpoint{8.987324in}{8.200939in}}%
\pgfpathquadraticcurveto{\pgfqpoint{8.975504in}{8.203418in}}{\pgfqpoint{8.970235in}{8.208576in}}%
\pgfpathquadraticcurveto{\pgfqpoint{8.964989in}{8.213733in}}{\pgfqpoint{8.964989in}{8.222720in}}%
\pgfpathquadraticcurveto{\pgfqpoint{8.964989in}{8.233655in}}{\pgfqpoint{8.972737in}{8.239587in}}%
\pgfpathquadraticcurveto{\pgfqpoint{8.980484in}{8.245542in}}{\pgfqpoint{8.994739in}{8.245542in}}%
\pgfpathquadraticcurveto{\pgfqpoint{9.001778in}{8.245542in}}{\pgfqpoint{9.007998in}{8.244501in}}%
\pgfpathquadraticcurveto{\pgfqpoint{9.014241in}{8.243483in}}{\pgfqpoint{9.019487in}{8.241402in}}%
\pgfpathlineto{\pgfqpoint{9.019487in}{8.241402in}}%
\pgfpathclose%
\pgfpathmoveto{\pgfqpoint{9.079120in}{8.205145in}}%
\pgfpathquadraticcurveto{\pgfqpoint{9.063691in}{8.205145in}}{\pgfqpoint{9.057737in}{8.201625in}}%
\pgfpathquadraticcurveto{\pgfqpoint{9.051782in}{8.198105in}}{\pgfqpoint{9.051782in}{8.189583in}}%
\pgfpathquadraticcurveto{\pgfqpoint{9.051782in}{8.182810in}}{\pgfqpoint{9.056254in}{8.178826in}}%
\pgfpathquadraticcurveto{\pgfqpoint{9.060725in}{8.174863in}}{\pgfqpoint{9.068384in}{8.174863in}}%
\pgfpathquadraticcurveto{\pgfqpoint{9.078987in}{8.174863in}}{\pgfqpoint{9.085384in}{8.182367in}}%
\pgfpathquadraticcurveto{\pgfqpoint{9.091781in}{8.189871in}}{\pgfqpoint{9.091781in}{8.202311in}}%
\pgfpathlineto{\pgfqpoint{9.091781in}{8.205145in}}%
\pgfpathlineto{\pgfqpoint{9.079120in}{8.205145in}}%
\pgfpathlineto{\pgfqpoint{9.079120in}{8.205145in}}%
\pgfpathclose%
\pgfpathmoveto{\pgfqpoint{9.104509in}{8.210413in}}%
\pgfpathlineto{\pgfqpoint{9.104509in}{8.166208in}}%
\pgfpathlineto{\pgfqpoint{9.091781in}{8.166208in}}%
\pgfpathlineto{\pgfqpoint{9.091781in}{8.177962in}}%
\pgfpathquadraticcurveto{\pgfqpoint{9.087420in}{8.170923in}}{\pgfqpoint{9.080913in}{8.167559in}}%
\pgfpathquadraticcurveto{\pgfqpoint{9.074405in}{8.164194in}}{\pgfqpoint{9.064997in}{8.164194in}}%
\pgfpathquadraticcurveto{\pgfqpoint{9.053110in}{8.164194in}}{\pgfqpoint{9.046071in}{8.170879in}}%
\pgfpathquadraticcurveto{\pgfqpoint{9.039054in}{8.177564in}}{\pgfqpoint{9.039054in}{8.188764in}}%
\pgfpathquadraticcurveto{\pgfqpoint{9.039054in}{8.201824in}}{\pgfqpoint{9.047798in}{8.208465in}}%
\pgfpathquadraticcurveto{\pgfqpoint{9.056564in}{8.215105in}}{\pgfqpoint{9.073918in}{8.215105in}}%
\pgfpathlineto{\pgfqpoint{9.091781in}{8.215105in}}%
\pgfpathlineto{\pgfqpoint{9.091781in}{8.216367in}}%
\pgfpathquadraticcurveto{\pgfqpoint{9.091781in}{8.225155in}}{\pgfqpoint{9.086004in}{8.229958in}}%
\pgfpathquadraticcurveto{\pgfqpoint{9.080226in}{8.234762in}}{\pgfqpoint{9.069778in}{8.234762in}}%
\pgfpathquadraticcurveto{\pgfqpoint{9.063138in}{8.234762in}}{\pgfqpoint{9.056829in}{8.233168in}}%
\pgfpathquadraticcurveto{\pgfqpoint{9.050543in}{8.231574in}}{\pgfqpoint{9.044743in}{8.228387in}}%
\pgfpathlineto{\pgfqpoint{9.044743in}{8.240163in}}%
\pgfpathquadraticcurveto{\pgfqpoint{9.051716in}{8.242863in}}{\pgfqpoint{9.058290in}{8.244191in}}%
\pgfpathquadraticcurveto{\pgfqpoint{9.064864in}{8.245542in}}{\pgfqpoint{9.071084in}{8.245542in}}%
\pgfpathquadraticcurveto{\pgfqpoint{9.087907in}{8.245542in}}{\pgfqpoint{9.096208in}{8.236820in}}%
\pgfpathquadraticcurveto{\pgfqpoint{9.104509in}{8.228121in}}{\pgfqpoint{9.104509in}{8.210413in}}%
\pgfpathlineto{\pgfqpoint{9.104509in}{8.210413in}}%
\pgfpathclose%
\pgfpathmoveto{\pgfqpoint{9.130717in}{8.273853in}}%
\pgfpathlineto{\pgfqpoint{9.143445in}{8.273853in}}%
\pgfpathlineto{\pgfqpoint{9.143445in}{8.166208in}}%
\pgfpathlineto{\pgfqpoint{9.130717in}{8.166208in}}%
\pgfpathlineto{\pgfqpoint{9.130717in}{8.273853in}}%
\pgfpathlineto{\pgfqpoint{9.130717in}{8.273853in}}%
\pgfpathclose%
\pgfpathmoveto{\pgfqpoint{9.219458in}{8.241402in}}%
\pgfpathlineto{\pgfqpoint{9.219458in}{8.229361in}}%
\pgfpathquadraticcurveto{\pgfqpoint{9.214079in}{8.232128in}}{\pgfqpoint{9.208258in}{8.233500in}}%
\pgfpathquadraticcurveto{\pgfqpoint{9.202458in}{8.234895in}}{\pgfqpoint{9.196216in}{8.234895in}}%
\pgfpathquadraticcurveto{\pgfqpoint{9.186742in}{8.234895in}}{\pgfqpoint{9.182005in}{8.231995in}}%
\pgfpathquadraticcurveto{\pgfqpoint{9.177268in}{8.229095in}}{\pgfqpoint{9.177268in}{8.223273in}}%
\pgfpathquadraticcurveto{\pgfqpoint{9.177268in}{8.218846in}}{\pgfqpoint{9.180655in}{8.216323in}}%
\pgfpathquadraticcurveto{\pgfqpoint{9.184041in}{8.213799in}}{\pgfqpoint{9.194290in}{8.211520in}}%
\pgfpathlineto{\pgfqpoint{9.198651in}{8.210546in}}%
\pgfpathquadraticcurveto{\pgfqpoint{9.212198in}{8.207646in}}{\pgfqpoint{9.217909in}{8.202355in}}%
\pgfpathquadraticcurveto{\pgfqpoint{9.223620in}{8.197065in}}{\pgfqpoint{9.223620in}{8.187591in}}%
\pgfpathquadraticcurveto{\pgfqpoint{9.223620in}{8.176789in}}{\pgfqpoint{9.215075in}{8.170480in}}%
\pgfpathquadraticcurveto{\pgfqpoint{9.206531in}{8.164194in}}{\pgfqpoint{9.191590in}{8.164194in}}%
\pgfpathquadraticcurveto{\pgfqpoint{9.185370in}{8.164194in}}{\pgfqpoint{9.178618in}{8.165411in}}%
\pgfpathquadraticcurveto{\pgfqpoint{9.171867in}{8.166629in}}{\pgfqpoint{9.164407in}{8.169042in}}%
\pgfpathlineto{\pgfqpoint{9.164407in}{8.182190in}}%
\pgfpathquadraticcurveto{\pgfqpoint{9.171469in}{8.178516in}}{\pgfqpoint{9.178308in}{8.176678in}}%
\pgfpathquadraticcurveto{\pgfqpoint{9.185148in}{8.174863in}}{\pgfqpoint{9.191877in}{8.174863in}}%
\pgfpathquadraticcurveto{\pgfqpoint{9.200864in}{8.174863in}}{\pgfqpoint{9.205690in}{8.177940in}}%
\pgfpathquadraticcurveto{\pgfqpoint{9.210538in}{8.181017in}}{\pgfqpoint{9.210538in}{8.186617in}}%
\pgfpathquadraticcurveto{\pgfqpoint{9.210538in}{8.191797in}}{\pgfqpoint{9.207040in}{8.194564in}}%
\pgfpathquadraticcurveto{\pgfqpoint{9.203565in}{8.197331in}}{\pgfqpoint{9.191722in}{8.199898in}}%
\pgfpathlineto{\pgfqpoint{9.187295in}{8.200939in}}%
\pgfpathquadraticcurveto{\pgfqpoint{9.175475in}{8.203418in}}{\pgfqpoint{9.170207in}{8.208576in}}%
\pgfpathquadraticcurveto{\pgfqpoint{9.164961in}{8.213733in}}{\pgfqpoint{9.164961in}{8.222720in}}%
\pgfpathquadraticcurveto{\pgfqpoint{9.164961in}{8.233655in}}{\pgfqpoint{9.172708in}{8.239587in}}%
\pgfpathquadraticcurveto{\pgfqpoint{9.180455in}{8.245542in}}{\pgfqpoint{9.194711in}{8.245542in}}%
\pgfpathquadraticcurveto{\pgfqpoint{9.201750in}{8.245542in}}{\pgfqpoint{9.207970in}{8.244501in}}%
\pgfpathquadraticcurveto{\pgfqpoint{9.214212in}{8.243483in}}{\pgfqpoint{9.219458in}{8.241402in}}%
\pgfpathlineto{\pgfqpoint{9.219458in}{8.241402in}}%
\pgfpathclose%
\pgfpathmoveto{\pgfqpoint{9.247127in}{8.183784in}}%
\pgfpathlineto{\pgfqpoint{9.261715in}{8.183784in}}%
\pgfpathlineto{\pgfqpoint{9.261715in}{8.166208in}}%
\pgfpathlineto{\pgfqpoint{9.247127in}{8.166208in}}%
\pgfpathlineto{\pgfqpoint{9.247127in}{8.183784in}}%
\pgfpathlineto{\pgfqpoint{9.247127in}{8.183784in}}%
\pgfpathclose%
\pgfpathmoveto{\pgfqpoint{9.247127in}{8.239454in}}%
\pgfpathlineto{\pgfqpoint{9.261715in}{8.239454in}}%
\pgfpathlineto{\pgfqpoint{9.261715in}{8.221901in}}%
\pgfpathlineto{\pgfqpoint{9.247127in}{8.221901in}}%
\pgfpathlineto{\pgfqpoint{9.247127in}{8.239454in}}%
\pgfpathlineto{\pgfqpoint{9.247127in}{8.239454in}}%
\pgfpathclose%
\pgfpathmoveto{\pgfqpoint{9.370045in}{8.223406in}}%
\pgfpathquadraticcurveto{\pgfqpoint{9.360638in}{8.223406in}}{\pgfqpoint{9.355126in}{8.216965in}}%
\pgfpathquadraticcurveto{\pgfqpoint{9.349636in}{8.210546in}}{\pgfqpoint{9.349636in}{8.199345in}}%
\pgfpathquadraticcurveto{\pgfqpoint{9.349636in}{8.188211in}}{\pgfqpoint{9.355126in}{8.181725in}}%
\pgfpathquadraticcurveto{\pgfqpoint{9.360638in}{8.175262in}}{\pgfqpoint{9.370045in}{8.175262in}}%
\pgfpathquadraticcurveto{\pgfqpoint{9.379453in}{8.175262in}}{\pgfqpoint{9.384942in}{8.181725in}}%
\pgfpathquadraticcurveto{\pgfqpoint{9.390432in}{8.188211in}}{\pgfqpoint{9.390432in}{8.199345in}}%
\pgfpathquadraticcurveto{\pgfqpoint{9.390432in}{8.210546in}}{\pgfqpoint{9.384942in}{8.216965in}}%
\pgfpathquadraticcurveto{\pgfqpoint{9.379453in}{8.223406in}}{\pgfqpoint{9.370045in}{8.223406in}}%
\pgfpathlineto{\pgfqpoint{9.370045in}{8.223406in}}%
\pgfpathclose%
\pgfpathmoveto{\pgfqpoint{9.397781in}{8.267212in}}%
\pgfpathlineto{\pgfqpoint{9.397781in}{8.254484in}}%
\pgfpathquadraticcurveto{\pgfqpoint{9.392513in}{8.256964in}}{\pgfqpoint{9.387156in}{8.258270in}}%
\pgfpathquadraticcurveto{\pgfqpoint{9.381799in}{8.259598in}}{\pgfqpoint{9.376531in}{8.259598in}}%
\pgfpathquadraticcurveto{\pgfqpoint{9.362696in}{8.259598in}}{\pgfqpoint{9.355392in}{8.250257in}}%
\pgfpathquadraticcurveto{\pgfqpoint{9.348109in}{8.240915in}}{\pgfqpoint{9.347069in}{8.222034in}}%
\pgfpathquadraticcurveto{\pgfqpoint{9.351142in}{8.228055in}}{\pgfqpoint{9.357295in}{8.231264in}}%
\pgfpathquadraticcurveto{\pgfqpoint{9.363471in}{8.234474in}}{\pgfqpoint{9.370864in}{8.234474in}}%
\pgfpathquadraticcurveto{\pgfqpoint{9.386426in}{8.234474in}}{\pgfqpoint{9.395457in}{8.225022in}}%
\pgfpathquadraticcurveto{\pgfqpoint{9.404488in}{8.215592in}}{\pgfqpoint{9.404488in}{8.199345in}}%
\pgfpathquadraticcurveto{\pgfqpoint{9.404488in}{8.183430in}}{\pgfqpoint{9.395080in}{8.173801in}}%
\pgfpathquadraticcurveto{\pgfqpoint{9.385673in}{8.164194in}}{\pgfqpoint{9.370045in}{8.164194in}}%
\pgfpathquadraticcurveto{\pgfqpoint{9.352116in}{8.164194in}}{\pgfqpoint{9.342642in}{8.177918in}}%
\pgfpathquadraticcurveto{\pgfqpoint{9.333168in}{8.191664in}}{\pgfqpoint{9.333168in}{8.217740in}}%
\pgfpathquadraticcurveto{\pgfqpoint{9.333168in}{8.242221in}}{\pgfqpoint{9.344789in}{8.256786in}}%
\pgfpathquadraticcurveto{\pgfqpoint{9.356410in}{8.271352in}}{\pgfqpoint{9.375978in}{8.271352in}}%
\pgfpathquadraticcurveto{\pgfqpoint{9.381246in}{8.271352in}}{\pgfqpoint{9.386603in}{8.270311in}}%
\pgfpathquadraticcurveto{\pgfqpoint{9.391959in}{8.269271in}}{\pgfqpoint{9.397781in}{8.267212in}}%
\pgfpathlineto{\pgfqpoint{9.397781in}{8.267212in}}%
\pgfpathclose%
\pgfpathmoveto{\pgfqpoint{9.428549in}{8.183784in}}%
\pgfpathlineto{\pgfqpoint{9.443159in}{8.183784in}}%
\pgfpathlineto{\pgfqpoint{9.443159in}{8.166208in}}%
\pgfpathlineto{\pgfqpoint{9.428549in}{8.166208in}}%
\pgfpathlineto{\pgfqpoint{9.428549in}{8.183784in}}%
\pgfpathlineto{\pgfqpoint{9.428549in}{8.183784in}}%
\pgfpathclose%
\pgfpathmoveto{\pgfqpoint{9.473993in}{8.168355in}}%
\pgfpathlineto{\pgfqpoint{9.473993in}{8.181083in}}%
\pgfpathquadraticcurveto{\pgfqpoint{9.479261in}{8.178582in}}{\pgfqpoint{9.484640in}{8.177276in}}%
\pgfpathquadraticcurveto{\pgfqpoint{9.490041in}{8.175970in}}{\pgfqpoint{9.495243in}{8.175970in}}%
\pgfpathquadraticcurveto{\pgfqpoint{9.509078in}{8.175970in}}{\pgfqpoint{9.516360in}{8.185267in}}%
\pgfpathquadraticcurveto{\pgfqpoint{9.523665in}{8.194564in}}{\pgfqpoint{9.524705in}{8.213534in}}%
\pgfpathquadraticcurveto{\pgfqpoint{9.520699in}{8.207579in}}{\pgfqpoint{9.514523in}{8.204392in}}%
\pgfpathquadraticcurveto{\pgfqpoint{9.508370in}{8.201204in}}{\pgfqpoint{9.500910in}{8.201204in}}%
\pgfpathquadraticcurveto{\pgfqpoint{9.485415in}{8.201204in}}{\pgfqpoint{9.476384in}{8.210568in}}%
\pgfpathquadraticcurveto{\pgfqpoint{9.467353in}{8.219953in}}{\pgfqpoint{9.467353in}{8.236223in}}%
\pgfpathquadraticcurveto{\pgfqpoint{9.467353in}{8.252116in}}{\pgfqpoint{9.476760in}{8.261723in}}%
\pgfpathquadraticcurveto{\pgfqpoint{9.486168in}{8.271352in}}{\pgfqpoint{9.501795in}{8.271352in}}%
\pgfpathquadraticcurveto{\pgfqpoint{9.519725in}{8.271352in}}{\pgfqpoint{9.529155in}{8.257605in}}%
\pgfpathquadraticcurveto{\pgfqpoint{9.538607in}{8.243882in}}{\pgfqpoint{9.538607in}{8.217740in}}%
\pgfpathquadraticcurveto{\pgfqpoint{9.538607in}{8.193324in}}{\pgfqpoint{9.527008in}{8.178759in}}%
\pgfpathquadraticcurveto{\pgfqpoint{9.515431in}{8.164194in}}{\pgfqpoint{9.495863in}{8.164194in}}%
\pgfpathquadraticcurveto{\pgfqpoint{9.490595in}{8.164194in}}{\pgfqpoint{9.485194in}{8.165234in}}%
\pgfpathquadraticcurveto{\pgfqpoint{9.479815in}{8.166275in}}{\pgfqpoint{9.473993in}{8.168355in}}%
\pgfpathlineto{\pgfqpoint{9.473993in}{8.168355in}}%
\pgfpathclose%
\pgfpathmoveto{\pgfqpoint{9.501795in}{8.212139in}}%
\pgfpathquadraticcurveto{\pgfqpoint{9.511203in}{8.212139in}}{\pgfqpoint{9.516692in}{8.218559in}}%
\pgfpathquadraticcurveto{\pgfqpoint{9.522204in}{8.225000in}}{\pgfqpoint{9.522204in}{8.236223in}}%
\pgfpathquadraticcurveto{\pgfqpoint{9.522204in}{8.247357in}}{\pgfqpoint{9.516692in}{8.253820in}}%
\pgfpathquadraticcurveto{\pgfqpoint{9.511203in}{8.260284in}}{\pgfqpoint{9.501795in}{8.260284in}}%
\pgfpathquadraticcurveto{\pgfqpoint{9.492388in}{8.260284in}}{\pgfqpoint{9.486898in}{8.253820in}}%
\pgfpathquadraticcurveto{\pgfqpoint{9.481409in}{8.247357in}}{\pgfqpoint{9.481409in}{8.236223in}}%
\pgfpathquadraticcurveto{\pgfqpoint{9.481409in}{8.225000in}}{\pgfqpoint{9.486898in}{8.218559in}}%
\pgfpathquadraticcurveto{\pgfqpoint{9.492388in}{8.212139in}}{\pgfqpoint{9.501795in}{8.212139in}}%
\pgfpathlineto{\pgfqpoint{9.501795in}{8.212139in}}%
\pgfpathclose%
\pgfpathmoveto{\pgfqpoint{9.651564in}{8.211652in}}%
\pgfpathquadraticcurveto{\pgfqpoint{9.645543in}{8.211652in}}{\pgfqpoint{9.642112in}{8.206539in}}%
\pgfpathquadraticcurveto{\pgfqpoint{9.638703in}{8.201426in}}{\pgfqpoint{9.638703in}{8.192284in}}%
\pgfpathquadraticcurveto{\pgfqpoint{9.638703in}{8.183297in}}{\pgfqpoint{9.642112in}{8.178139in}}%
\pgfpathquadraticcurveto{\pgfqpoint{9.645543in}{8.172982in}}{\pgfqpoint{9.651564in}{8.172982in}}%
\pgfpathquadraticcurveto{\pgfqpoint{9.657452in}{8.172982in}}{\pgfqpoint{9.660860in}{8.178139in}}%
\pgfpathquadraticcurveto{\pgfqpoint{9.664291in}{8.183297in}}{\pgfqpoint{9.664291in}{8.192284in}}%
\pgfpathquadraticcurveto{\pgfqpoint{9.664291in}{8.201359in}}{\pgfqpoint{9.660860in}{8.206495in}}%
\pgfpathquadraticcurveto{\pgfqpoint{9.657452in}{8.211652in}}{\pgfqpoint{9.651564in}{8.211652in}}%
\pgfpathlineto{\pgfqpoint{9.651564in}{8.211652in}}%
\pgfpathclose%
\pgfpathmoveto{\pgfqpoint{9.651564in}{8.220440in}}%
\pgfpathquadraticcurveto{\pgfqpoint{9.662498in}{8.220440in}}{\pgfqpoint{9.668918in}{8.212826in}}%
\pgfpathquadraticcurveto{\pgfqpoint{9.675359in}{8.205233in}}{\pgfqpoint{9.675359in}{8.192284in}}%
\pgfpathquadraticcurveto{\pgfqpoint{9.675359in}{8.179357in}}{\pgfqpoint{9.668896in}{8.171764in}}%
\pgfpathquadraticcurveto{\pgfqpoint{9.662432in}{8.164194in}}{\pgfqpoint{9.651564in}{8.164194in}}%
\pgfpathquadraticcurveto{\pgfqpoint{9.640496in}{8.164194in}}{\pgfqpoint{9.634054in}{8.171764in}}%
\pgfpathquadraticcurveto{\pgfqpoint{9.627635in}{8.179357in}}{\pgfqpoint{9.627635in}{8.192284in}}%
\pgfpathquadraticcurveto{\pgfqpoint{9.627635in}{8.205299in}}{\pgfqpoint{9.634099in}{8.212870in}}%
\pgfpathquadraticcurveto{\pgfqpoint{9.640562in}{8.220440in}}{\pgfqpoint{9.651564in}{8.220440in}}%
\pgfpathlineto{\pgfqpoint{9.651564in}{8.220440in}}%
\pgfpathclose%
\pgfpathmoveto{\pgfqpoint{9.580177in}{8.262564in}}%
\pgfpathquadraticcurveto{\pgfqpoint{9.574222in}{8.262564in}}{\pgfqpoint{9.570791in}{8.257406in}}%
\pgfpathquadraticcurveto{\pgfqpoint{9.567383in}{8.252271in}}{\pgfqpoint{9.567383in}{8.243262in}}%
\pgfpathquadraticcurveto{\pgfqpoint{9.567383in}{8.234142in}}{\pgfqpoint{9.570769in}{8.229007in}}%
\pgfpathquadraticcurveto{\pgfqpoint{9.574156in}{8.223893in}}{\pgfqpoint{9.580177in}{8.223893in}}%
\pgfpathquadraticcurveto{\pgfqpoint{9.586198in}{8.223893in}}{\pgfqpoint{9.589607in}{8.229007in}}%
\pgfpathquadraticcurveto{\pgfqpoint{9.593038in}{8.234142in}}{\pgfqpoint{9.593038in}{8.243262in}}%
\pgfpathquadraticcurveto{\pgfqpoint{9.593038in}{8.252182in}}{\pgfqpoint{9.589584in}{8.257362in}}%
\pgfpathquadraticcurveto{\pgfqpoint{9.586131in}{8.262564in}}{\pgfqpoint{9.580177in}{8.262564in}}%
\pgfpathlineto{\pgfqpoint{9.580177in}{8.262564in}}%
\pgfpathclose%
\pgfpathmoveto{\pgfqpoint{9.642643in}{8.271352in}}%
\pgfpathlineto{\pgfqpoint{9.653711in}{8.271352in}}%
\pgfpathlineto{\pgfqpoint{9.589097in}{8.164194in}}%
\pgfpathlineto{\pgfqpoint{9.578030in}{8.164194in}}%
\pgfpathlineto{\pgfqpoint{9.642643in}{8.271352in}}%
\pgfpathlineto{\pgfqpoint{9.642643in}{8.271352in}}%
\pgfpathclose%
\pgfpathmoveto{\pgfqpoint{9.580177in}{8.271352in}}%
\pgfpathquadraticcurveto{\pgfqpoint{9.591112in}{8.271352in}}{\pgfqpoint{9.597597in}{8.263781in}}%
\pgfpathquadraticcurveto{\pgfqpoint{9.604105in}{8.256211in}}{\pgfqpoint{9.604105in}{8.243262in}}%
\pgfpathquadraticcurveto{\pgfqpoint{9.604105in}{8.230202in}}{\pgfqpoint{9.597642in}{8.222654in}}%
\pgfpathquadraticcurveto{\pgfqpoint{9.591178in}{8.215105in}}{\pgfqpoint{9.580177in}{8.215105in}}%
\pgfpathquadraticcurveto{\pgfqpoint{9.569176in}{8.215105in}}{\pgfqpoint{9.562778in}{8.222676in}}%
\pgfpathquadraticcurveto{\pgfqpoint{9.556381in}{8.230268in}}{\pgfqpoint{9.556381in}{8.243262in}}%
\pgfpathquadraticcurveto{\pgfqpoint{9.556381in}{8.256145in}}{\pgfqpoint{9.562801in}{8.263737in}}%
\pgfpathquadraticcurveto{\pgfqpoint{9.569242in}{8.271352in}}{\pgfqpoint{9.580177in}{8.271352in}}%
\pgfpathlineto{\pgfqpoint{9.580177in}{8.271352in}}%
\pgfpathclose%
\pgfusepath{fill}%
\end{pgfscope}%
\begin{pgfscope}%
\pgfsetbuttcap%
\pgfsetroundjoin%
\definecolor{currentfill}{rgb}{0.090196,0.168627,0.227451}%
\pgfsetfillcolor{currentfill}%
\pgfsetlinewidth{0.000000pt}%
\definecolor{currentstroke}{rgb}{0.090196,0.168627,0.227451}%
\pgfsetstrokecolor{currentstroke}%
\pgfsetdash{}{0pt}%
\pgfpathmoveto{\pgfqpoint{8.300905in}{8.046522in}}%
\pgfpathlineto{\pgfqpoint{8.300905in}{8.032887in}}%
\pgfpathquadraticcurveto{\pgfqpoint{8.292958in}{8.036694in}}{\pgfqpoint{8.285897in}{8.038553in}}%
\pgfpathquadraticcurveto{\pgfqpoint{8.278836in}{8.040435in}}{\pgfqpoint{8.272261in}{8.040435in}}%
\pgfpathquadraticcurveto{\pgfqpoint{8.260862in}{8.040435in}}{\pgfqpoint{8.254664in}{8.036008in}}%
\pgfpathquadraticcurveto{\pgfqpoint{8.248466in}{8.031581in}}{\pgfqpoint{8.248466in}{8.023413in}}%
\pgfpathquadraticcurveto{\pgfqpoint{8.248466in}{8.016551in}}{\pgfqpoint{8.252583in}{8.013053in}}%
\pgfpathquadraticcurveto{\pgfqpoint{8.256700in}{8.009578in}}{\pgfqpoint{8.268189in}{8.007431in}}%
\pgfpathlineto{\pgfqpoint{8.276622in}{8.005704in}}%
\pgfpathquadraticcurveto{\pgfqpoint{8.292250in}{8.002716in}}{\pgfqpoint{8.299687in}{7.995212in}}%
\pgfpathquadraticcurveto{\pgfqpoint{8.307125in}{7.987708in}}{\pgfqpoint{8.307125in}{7.975135in}}%
\pgfpathquadraticcurveto{\pgfqpoint{8.307125in}{7.960105in}}{\pgfqpoint{8.297053in}{7.952358in}}%
\pgfpathquadraticcurveto{\pgfqpoint{8.287004in}{7.944611in}}{\pgfqpoint{8.267569in}{7.944611in}}%
\pgfpathquadraticcurveto{\pgfqpoint{8.260242in}{7.944611in}}{\pgfqpoint{8.251963in}{7.946271in}}%
\pgfpathquadraticcurveto{\pgfqpoint{8.243707in}{7.947931in}}{\pgfqpoint{8.234853in}{7.951185in}}%
\pgfpathlineto{\pgfqpoint{8.234853in}{7.965573in}}%
\pgfpathquadraticcurveto{\pgfqpoint{8.243353in}{7.960814in}}{\pgfqpoint{8.251521in}{7.958379in}}%
\pgfpathquadraticcurveto{\pgfqpoint{8.259689in}{7.955966in}}{\pgfqpoint{8.267569in}{7.955966in}}%
\pgfpathquadraticcurveto{\pgfqpoint{8.279522in}{7.955966in}}{\pgfqpoint{8.286030in}{7.960659in}}%
\pgfpathquadraticcurveto{\pgfqpoint{8.292538in}{7.965374in}}{\pgfqpoint{8.292538in}{7.974095in}}%
\pgfpathquadraticcurveto{\pgfqpoint{8.292538in}{7.981688in}}{\pgfqpoint{8.287867in}{7.985982in}}%
\pgfpathquadraticcurveto{\pgfqpoint{8.283196in}{7.990276in}}{\pgfqpoint{8.272549in}{7.992423in}}%
\pgfpathlineto{\pgfqpoint{8.264027in}{7.994083in}}%
\pgfpathquadraticcurveto{\pgfqpoint{8.248400in}{7.997182in}}{\pgfqpoint{8.241405in}{8.003823in}}%
\pgfpathquadraticcurveto{\pgfqpoint{8.234432in}{8.010464in}}{\pgfqpoint{8.234432in}{8.022306in}}%
\pgfpathquadraticcurveto{\pgfqpoint{8.234432in}{8.036008in}}{\pgfqpoint{8.244083in}{8.043888in}}%
\pgfpathquadraticcurveto{\pgfqpoint{8.253734in}{8.051768in}}{\pgfqpoint{8.270668in}{8.051768in}}%
\pgfpathquadraticcurveto{\pgfqpoint{8.277950in}{8.051768in}}{\pgfqpoint{8.285476in}{8.050440in}}%
\pgfpathquadraticcurveto{\pgfqpoint{8.293025in}{8.049134in}}{\pgfqpoint{8.300905in}{8.046522in}}%
\pgfpathlineto{\pgfqpoint{8.300905in}{8.046522in}}%
\pgfpathclose%
\pgfpathmoveto{\pgfqpoint{8.328375in}{8.024099in}}%
\pgfpathlineto{\pgfqpoint{8.341103in}{8.024099in}}%
\pgfpathlineto{\pgfqpoint{8.341103in}{7.946625in}}%
\pgfpathlineto{\pgfqpoint{8.328375in}{7.946625in}}%
\pgfpathlineto{\pgfqpoint{8.328375in}{8.024099in}}%
\pgfpathlineto{\pgfqpoint{8.328375in}{8.024099in}}%
\pgfpathclose%
\pgfpathmoveto{\pgfqpoint{8.328375in}{8.054270in}}%
\pgfpathlineto{\pgfqpoint{8.341103in}{8.054270in}}%
\pgfpathlineto{\pgfqpoint{8.341103in}{8.038133in}}%
\pgfpathlineto{\pgfqpoint{8.328375in}{8.038133in}}%
\pgfpathlineto{\pgfqpoint{8.328375in}{8.054270in}}%
\pgfpathlineto{\pgfqpoint{8.328375in}{8.054270in}}%
\pgfpathclose%
\pgfpathmoveto{\pgfqpoint{8.418709in}{7.986270in}}%
\pgfpathquadraticcurveto{\pgfqpoint{8.418709in}{8.000104in}}{\pgfqpoint{8.412998in}{8.007697in}}%
\pgfpathquadraticcurveto{\pgfqpoint{8.407310in}{8.015311in}}{\pgfqpoint{8.396995in}{8.015311in}}%
\pgfpathquadraticcurveto{\pgfqpoint{8.386768in}{8.015311in}}{\pgfqpoint{8.381057in}{8.007697in}}%
\pgfpathquadraticcurveto{\pgfqpoint{8.375346in}{8.000104in}}{\pgfqpoint{8.375346in}{7.986270in}}%
\pgfpathquadraticcurveto{\pgfqpoint{8.375346in}{7.972501in}}{\pgfqpoint{8.381057in}{7.964887in}}%
\pgfpathquadraticcurveto{\pgfqpoint{8.386768in}{7.957272in}}{\pgfqpoint{8.396995in}{7.957272in}}%
\pgfpathquadraticcurveto{\pgfqpoint{8.407310in}{7.957272in}}{\pgfqpoint{8.412998in}{7.964887in}}%
\pgfpathquadraticcurveto{\pgfqpoint{8.418709in}{7.972501in}}{\pgfqpoint{8.418709in}{7.986270in}}%
\pgfpathlineto{\pgfqpoint{8.418709in}{7.986270in}}%
\pgfpathclose%
\pgfpathmoveto{\pgfqpoint{8.431437in}{7.956232in}}%
\pgfpathquadraticcurveto{\pgfqpoint{8.431437in}{7.936465in}}{\pgfqpoint{8.422650in}{7.926814in}}%
\pgfpathquadraticcurveto{\pgfqpoint{8.413884in}{7.917163in}}{\pgfqpoint{8.395755in}{7.917163in}}%
\pgfpathquadraticcurveto{\pgfqpoint{8.389048in}{7.917163in}}{\pgfqpoint{8.383094in}{7.918159in}}%
\pgfpathquadraticcurveto{\pgfqpoint{8.377139in}{7.919155in}}{\pgfqpoint{8.371539in}{7.921236in}}%
\pgfpathlineto{\pgfqpoint{8.371539in}{7.933609in}}%
\pgfpathquadraticcurveto{\pgfqpoint{8.377139in}{7.930577in}}{\pgfqpoint{8.382607in}{7.929138in}}%
\pgfpathquadraticcurveto{\pgfqpoint{8.388074in}{7.927677in}}{\pgfqpoint{8.393741in}{7.927677in}}%
\pgfpathquadraticcurveto{\pgfqpoint{8.406269in}{7.927677in}}{\pgfqpoint{8.412489in}{7.934207in}}%
\pgfpathquadraticcurveto{\pgfqpoint{8.418709in}{7.940737in}}{\pgfqpoint{8.418709in}{7.953952in}}%
\pgfpathlineto{\pgfqpoint{8.418709in}{7.960260in}}%
\pgfpathquadraticcurveto{\pgfqpoint{8.414769in}{7.953398in}}{\pgfqpoint{8.408616in}{7.950012in}}%
\pgfpathquadraticcurveto{\pgfqpoint{8.402462in}{7.946625in}}{\pgfqpoint{8.393873in}{7.946625in}}%
\pgfpathquadraticcurveto{\pgfqpoint{8.379640in}{7.946625in}}{\pgfqpoint{8.370919in}{7.957471in}}%
\pgfpathquadraticcurveto{\pgfqpoint{8.362198in}{7.968340in}}{\pgfqpoint{8.362198in}{7.986270in}}%
\pgfpathquadraticcurveto{\pgfqpoint{8.362198in}{8.004243in}}{\pgfqpoint{8.370919in}{8.015090in}}%
\pgfpathquadraticcurveto{\pgfqpoint{8.379640in}{8.025958in}}{\pgfqpoint{8.393873in}{8.025958in}}%
\pgfpathquadraticcurveto{\pgfqpoint{8.402462in}{8.025958in}}{\pgfqpoint{8.408616in}{8.022572in}}%
\pgfpathquadraticcurveto{\pgfqpoint{8.414769in}{8.019185in}}{\pgfqpoint{8.418709in}{8.012345in}}%
\pgfpathlineto{\pgfqpoint{8.418709in}{8.024099in}}%
\pgfpathlineto{\pgfqpoint{8.431437in}{8.024099in}}%
\pgfpathlineto{\pgfqpoint{8.431437in}{7.956232in}}%
\pgfpathlineto{\pgfqpoint{8.431437in}{7.956232in}}%
\pgfpathclose%
\pgfpathmoveto{\pgfqpoint{8.522082in}{7.993397in}}%
\pgfpathlineto{\pgfqpoint{8.522082in}{7.946625in}}%
\pgfpathlineto{\pgfqpoint{8.509354in}{7.946625in}}%
\pgfpathlineto{\pgfqpoint{8.509354in}{7.992977in}}%
\pgfpathquadraticcurveto{\pgfqpoint{8.509354in}{8.003978in}}{\pgfqpoint{8.505060in}{8.009423in}}%
\pgfpathquadraticcurveto{\pgfqpoint{8.500765in}{8.014891in}}{\pgfqpoint{8.492199in}{8.014891in}}%
\pgfpathquadraticcurveto{\pgfqpoint{8.481884in}{8.014891in}}{\pgfqpoint{8.475929in}{8.008316in}}%
\pgfpathquadraticcurveto{\pgfqpoint{8.469975in}{8.001764in}}{\pgfqpoint{8.469975in}{7.990409in}}%
\pgfpathlineto{\pgfqpoint{8.469975in}{7.946625in}}%
\pgfpathlineto{\pgfqpoint{8.457181in}{7.946625in}}%
\pgfpathlineto{\pgfqpoint{8.457181in}{8.024099in}}%
\pgfpathlineto{\pgfqpoint{8.469975in}{8.024099in}}%
\pgfpathlineto{\pgfqpoint{8.469975in}{8.012057in}}%
\pgfpathquadraticcurveto{\pgfqpoint{8.474557in}{8.019052in}}{\pgfqpoint{8.480733in}{8.022505in}}%
\pgfpathquadraticcurveto{\pgfqpoint{8.486931in}{8.025958in}}{\pgfqpoint{8.495032in}{8.025958in}}%
\pgfpathquadraticcurveto{\pgfqpoint{8.508380in}{8.025958in}}{\pgfqpoint{8.515220in}{8.017702in}}%
\pgfpathquadraticcurveto{\pgfqpoint{8.522082in}{8.009445in}}{\pgfqpoint{8.522082in}{7.993397in}}%
\pgfpathlineto{\pgfqpoint{8.522082in}{7.993397in}}%
\pgfpathclose%
\pgfpathmoveto{\pgfqpoint{8.547449in}{8.024099in}}%
\pgfpathlineto{\pgfqpoint{8.560177in}{8.024099in}}%
\pgfpathlineto{\pgfqpoint{8.560177in}{7.946625in}}%
\pgfpathlineto{\pgfqpoint{8.547449in}{7.946625in}}%
\pgfpathlineto{\pgfqpoint{8.547449in}{8.024099in}}%
\pgfpathlineto{\pgfqpoint{8.547449in}{8.024099in}}%
\pgfpathclose%
\pgfpathmoveto{\pgfqpoint{8.547449in}{8.054270in}}%
\pgfpathlineto{\pgfqpoint{8.560177in}{8.054270in}}%
\pgfpathlineto{\pgfqpoint{8.560177in}{8.038133in}}%
\pgfpathlineto{\pgfqpoint{8.547449in}{8.038133in}}%
\pgfpathlineto{\pgfqpoint{8.547449in}{8.054270in}}%
\pgfpathlineto{\pgfqpoint{8.547449in}{8.054270in}}%
\pgfpathclose%
\pgfpathmoveto{\pgfqpoint{8.649405in}{8.024099in}}%
\pgfpathlineto{\pgfqpoint{8.649405in}{7.946625in}}%
\pgfpathlineto{\pgfqpoint{8.636610in}{7.946625in}}%
\pgfpathlineto{\pgfqpoint{8.636610in}{8.014204in}}%
\pgfpathlineto{\pgfqpoint{8.601681in}{8.014204in}}%
\pgfpathlineto{\pgfqpoint{8.601681in}{7.946625in}}%
\pgfpathlineto{\pgfqpoint{8.588886in}{7.946625in}}%
\pgfpathlineto{\pgfqpoint{8.588886in}{8.014204in}}%
\pgfpathlineto{\pgfqpoint{8.576712in}{8.014204in}}%
\pgfpathlineto{\pgfqpoint{8.576712in}{8.024099in}}%
\pgfpathlineto{\pgfqpoint{8.588886in}{8.024099in}}%
\pgfpathlineto{\pgfqpoint{8.588886in}{8.029500in}}%
\pgfpathquadraticcurveto{\pgfqpoint{8.588886in}{8.042161in}}{\pgfqpoint{8.594863in}{8.048204in}}%
\pgfpathquadraticcurveto{\pgfqpoint{8.600862in}{8.054270in}}{\pgfqpoint{8.613235in}{8.054270in}}%
\pgfpathlineto{\pgfqpoint{8.626030in}{8.054270in}}%
\pgfpathlineto{\pgfqpoint{8.626030in}{8.043667in}}%
\pgfpathlineto{\pgfqpoint{8.613855in}{8.043667in}}%
\pgfpathquadraticcurveto{\pgfqpoint{8.607015in}{8.043667in}}{\pgfqpoint{8.604337in}{8.040900in}}%
\pgfpathquadraticcurveto{\pgfqpoint{8.601681in}{8.038133in}}{\pgfqpoint{8.601681in}{8.030939in}}%
\pgfpathlineto{\pgfqpoint{8.601681in}{8.024099in}}%
\pgfpathlineto{\pgfqpoint{8.649405in}{8.024099in}}%
\pgfpathlineto{\pgfqpoint{8.649405in}{8.024099in}}%
\pgfpathclose%
\pgfpathmoveto{\pgfqpoint{8.636610in}{8.054115in}}%
\pgfpathlineto{\pgfqpoint{8.649405in}{8.054115in}}%
\pgfpathlineto{\pgfqpoint{8.649405in}{8.038000in}}%
\pgfpathlineto{\pgfqpoint{8.636610in}{8.038000in}}%
\pgfpathlineto{\pgfqpoint{8.636610in}{8.054115in}}%
\pgfpathlineto{\pgfqpoint{8.636610in}{8.054115in}}%
\pgfpathclose%
\pgfpathmoveto{\pgfqpoint{8.731793in}{8.021133in}}%
\pgfpathlineto{\pgfqpoint{8.731793in}{8.009224in}}%
\pgfpathquadraticcurveto{\pgfqpoint{8.726392in}{8.012212in}}{\pgfqpoint{8.720969in}{8.013695in}}%
\pgfpathquadraticcurveto{\pgfqpoint{8.715545in}{8.015178in}}{\pgfqpoint{8.710011in}{8.015178in}}%
\pgfpathquadraticcurveto{\pgfqpoint{8.697616in}{8.015178in}}{\pgfqpoint{8.690754in}{8.007320in}}%
\pgfpathquadraticcurveto{\pgfqpoint{8.683914in}{7.999484in}}{\pgfqpoint{8.683914in}{7.985296in}}%
\pgfpathquadraticcurveto{\pgfqpoint{8.683914in}{7.971107in}}{\pgfqpoint{8.690754in}{7.963249in}}%
\pgfpathquadraticcurveto{\pgfqpoint{8.697616in}{7.955413in}}{\pgfqpoint{8.710011in}{7.955413in}}%
\pgfpathquadraticcurveto{\pgfqpoint{8.715545in}{7.955413in}}{\pgfqpoint{8.720969in}{7.956896in}}%
\pgfpathquadraticcurveto{\pgfqpoint{8.726392in}{7.958379in}}{\pgfqpoint{8.731793in}{7.961367in}}%
\pgfpathlineto{\pgfqpoint{8.731793in}{7.949591in}}%
\pgfpathquadraticcurveto{\pgfqpoint{8.726458in}{7.947112in}}{\pgfqpoint{8.720747in}{7.945872in}}%
\pgfpathquadraticcurveto{\pgfqpoint{8.715058in}{7.944611in}}{\pgfqpoint{8.708617in}{7.944611in}}%
\pgfpathquadraticcurveto{\pgfqpoint{8.691108in}{7.944611in}}{\pgfqpoint{8.680793in}{7.955612in}}%
\pgfpathquadraticcurveto{\pgfqpoint{8.670500in}{7.966613in}}{\pgfqpoint{8.670500in}{7.985296in}}%
\pgfpathquadraticcurveto{\pgfqpoint{8.670500in}{8.004243in}}{\pgfqpoint{8.680903in}{8.015090in}}%
\pgfpathquadraticcurveto{\pgfqpoint{8.691329in}{8.025958in}}{\pgfqpoint{8.709458in}{8.025958in}}%
\pgfpathquadraticcurveto{\pgfqpoint{8.715324in}{8.025958in}}{\pgfqpoint{8.720924in}{8.024741in}}%
\pgfpathquadraticcurveto{\pgfqpoint{8.726525in}{8.023546in}}{\pgfqpoint{8.731793in}{8.021133in}}%
\pgfpathlineto{\pgfqpoint{8.731793in}{8.021133in}}%
\pgfpathclose%
\pgfpathmoveto{\pgfqpoint{8.789146in}{7.985561in}}%
\pgfpathquadraticcurveto{\pgfqpoint{8.773717in}{7.985561in}}{\pgfqpoint{8.767763in}{7.982042in}}%
\pgfpathquadraticcurveto{\pgfqpoint{8.761808in}{7.978522in}}{\pgfqpoint{8.761808in}{7.970000in}}%
\pgfpathquadraticcurveto{\pgfqpoint{8.761808in}{7.963227in}}{\pgfqpoint{8.766280in}{7.959242in}}%
\pgfpathquadraticcurveto{\pgfqpoint{8.770751in}{7.955280in}}{\pgfqpoint{8.778410in}{7.955280in}}%
\pgfpathquadraticcurveto{\pgfqpoint{8.789013in}{7.955280in}}{\pgfqpoint{8.795410in}{7.962784in}}%
\pgfpathquadraticcurveto{\pgfqpoint{8.801807in}{7.970288in}}{\pgfqpoint{8.801807in}{7.982728in}}%
\pgfpathlineto{\pgfqpoint{8.801807in}{7.985561in}}%
\pgfpathlineto{\pgfqpoint{8.789146in}{7.985561in}}%
\pgfpathlineto{\pgfqpoint{8.789146in}{7.985561in}}%
\pgfpathclose%
\pgfpathmoveto{\pgfqpoint{8.814535in}{7.990829in}}%
\pgfpathlineto{\pgfqpoint{8.814535in}{7.946625in}}%
\pgfpathlineto{\pgfqpoint{8.801807in}{7.946625in}}%
\pgfpathlineto{\pgfqpoint{8.801807in}{7.958379in}}%
\pgfpathquadraticcurveto{\pgfqpoint{8.797446in}{7.951340in}}{\pgfqpoint{8.790939in}{7.947975in}}%
\pgfpathquadraticcurveto{\pgfqpoint{8.784431in}{7.944611in}}{\pgfqpoint{8.775023in}{7.944611in}}%
\pgfpathquadraticcurveto{\pgfqpoint{8.763136in}{7.944611in}}{\pgfqpoint{8.756097in}{7.951296in}}%
\pgfpathquadraticcurveto{\pgfqpoint{8.749080in}{7.957980in}}{\pgfqpoint{8.749080in}{7.969181in}}%
\pgfpathquadraticcurveto{\pgfqpoint{8.749080in}{7.982241in}}{\pgfqpoint{8.757824in}{7.988882in}}%
\pgfpathquadraticcurveto{\pgfqpoint{8.766590in}{7.995522in}}{\pgfqpoint{8.783944in}{7.995522in}}%
\pgfpathlineto{\pgfqpoint{8.801807in}{7.995522in}}%
\pgfpathlineto{\pgfqpoint{8.801807in}{7.996784in}}%
\pgfpathquadraticcurveto{\pgfqpoint{8.801807in}{8.005572in}}{\pgfqpoint{8.796030in}{8.010375in}}%
\pgfpathquadraticcurveto{\pgfqpoint{8.790252in}{8.015178in}}{\pgfqpoint{8.779804in}{8.015178in}}%
\pgfpathquadraticcurveto{\pgfqpoint{8.773164in}{8.015178in}}{\pgfqpoint{8.766855in}{8.013585in}}%
\pgfpathquadraticcurveto{\pgfqpoint{8.760569in}{8.011991in}}{\pgfqpoint{8.754769in}{8.008803in}}%
\pgfpathlineto{\pgfqpoint{8.754769in}{8.020579in}}%
\pgfpathquadraticcurveto{\pgfqpoint{8.761742in}{8.023280in}}{\pgfqpoint{8.768316in}{8.024608in}}%
\pgfpathquadraticcurveto{\pgfqpoint{8.774890in}{8.025958in}}{\pgfqpoint{8.781110in}{8.025958in}}%
\pgfpathquadraticcurveto{\pgfqpoint{8.797933in}{8.025958in}}{\pgfqpoint{8.806234in}{8.017237in}}%
\pgfpathquadraticcurveto{\pgfqpoint{8.814535in}{8.008538in}}{\pgfqpoint{8.814535in}{7.990829in}}%
\pgfpathlineto{\pgfqpoint{8.814535in}{7.990829in}}%
\pgfpathclose%
\pgfpathmoveto{\pgfqpoint{8.905157in}{7.993397in}}%
\pgfpathlineto{\pgfqpoint{8.905157in}{7.946625in}}%
\pgfpathlineto{\pgfqpoint{8.892429in}{7.946625in}}%
\pgfpathlineto{\pgfqpoint{8.892429in}{7.992977in}}%
\pgfpathquadraticcurveto{\pgfqpoint{8.892429in}{8.003978in}}{\pgfqpoint{8.888135in}{8.009423in}}%
\pgfpathquadraticcurveto{\pgfqpoint{8.883841in}{8.014891in}}{\pgfqpoint{8.875275in}{8.014891in}}%
\pgfpathquadraticcurveto{\pgfqpoint{8.864959in}{8.014891in}}{\pgfqpoint{8.859005in}{8.008316in}}%
\pgfpathquadraticcurveto{\pgfqpoint{8.853051in}{8.001764in}}{\pgfqpoint{8.853051in}{7.990409in}}%
\pgfpathlineto{\pgfqpoint{8.853051in}{7.946625in}}%
\pgfpathlineto{\pgfqpoint{8.840256in}{7.946625in}}%
\pgfpathlineto{\pgfqpoint{8.840256in}{8.024099in}}%
\pgfpathlineto{\pgfqpoint{8.853051in}{8.024099in}}%
\pgfpathlineto{\pgfqpoint{8.853051in}{8.012057in}}%
\pgfpathquadraticcurveto{\pgfqpoint{8.857633in}{8.019052in}}{\pgfqpoint{8.863808in}{8.022505in}}%
\pgfpathquadraticcurveto{\pgfqpoint{8.870006in}{8.025958in}}{\pgfqpoint{8.878108in}{8.025958in}}%
\pgfpathquadraticcurveto{\pgfqpoint{8.891455in}{8.025958in}}{\pgfqpoint{8.898295in}{8.017702in}}%
\pgfpathquadraticcurveto{\pgfqpoint{8.905157in}{8.009445in}}{\pgfqpoint{8.905157in}{7.993397in}}%
\pgfpathlineto{\pgfqpoint{8.905157in}{7.993397in}}%
\pgfpathclose%
\pgfpathmoveto{\pgfqpoint{8.943120in}{8.046102in}}%
\pgfpathlineto{\pgfqpoint{8.943120in}{8.024099in}}%
\pgfpathlineto{\pgfqpoint{8.969328in}{8.024099in}}%
\pgfpathlineto{\pgfqpoint{8.969328in}{8.014204in}}%
\pgfpathlineto{\pgfqpoint{8.943120in}{8.014204in}}%
\pgfpathlineto{\pgfqpoint{8.943120in}{7.972147in}}%
\pgfpathquadraticcurveto{\pgfqpoint{8.943120in}{7.962673in}}{\pgfqpoint{8.945709in}{7.959973in}}%
\pgfpathquadraticcurveto{\pgfqpoint{8.948299in}{7.957272in}}{\pgfqpoint{8.956268in}{7.957272in}}%
\pgfpathlineto{\pgfqpoint{8.969328in}{7.957272in}}%
\pgfpathlineto{\pgfqpoint{8.969328in}{7.946625in}}%
\pgfpathlineto{\pgfqpoint{8.956268in}{7.946625in}}%
\pgfpathquadraticcurveto{\pgfqpoint{8.941526in}{7.946625in}}{\pgfqpoint{8.935926in}{7.952115in}}%
\pgfpathquadraticcurveto{\pgfqpoint{8.930325in}{7.957626in}}{\pgfqpoint{8.930325in}{7.972147in}}%
\pgfpathlineto{\pgfqpoint{8.930325in}{8.014204in}}%
\pgfpathlineto{\pgfqpoint{8.920984in}{8.014204in}}%
\pgfpathlineto{\pgfqpoint{8.920984in}{8.024099in}}%
\pgfpathlineto{\pgfqpoint{8.930325in}{8.024099in}}%
\pgfpathlineto{\pgfqpoint{8.930325in}{8.046102in}}%
\pgfpathlineto{\pgfqpoint{8.943120in}{8.046102in}}%
\pgfpathlineto{\pgfqpoint{8.943120in}{8.046102in}}%
\pgfpathclose%
\pgfpathmoveto{\pgfqpoint{9.082064in}{8.012345in}}%
\pgfpathlineto{\pgfqpoint{9.082064in}{8.054270in}}%
\pgfpathlineto{\pgfqpoint{9.094791in}{8.054270in}}%
\pgfpathlineto{\pgfqpoint{9.094791in}{7.946625in}}%
\pgfpathlineto{\pgfqpoint{9.082064in}{7.946625in}}%
\pgfpathlineto{\pgfqpoint{9.082064in}{7.958246in}}%
\pgfpathquadraticcurveto{\pgfqpoint{9.078057in}{7.951340in}}{\pgfqpoint{9.071926in}{7.947975in}}%
\pgfpathquadraticcurveto{\pgfqpoint{9.065816in}{7.944611in}}{\pgfqpoint{9.057228in}{7.944611in}}%
\pgfpathquadraticcurveto{\pgfqpoint{9.043194in}{7.944611in}}{\pgfqpoint{9.034362in}{7.955811in}}%
\pgfpathquadraticcurveto{\pgfqpoint{9.025552in}{7.967034in}}{\pgfqpoint{9.025552in}{7.985296in}}%
\pgfpathquadraticcurveto{\pgfqpoint{9.025552in}{8.003557in}}{\pgfqpoint{9.034362in}{8.014758in}}%
\pgfpathquadraticcurveto{\pgfqpoint{9.043194in}{8.025958in}}{\pgfqpoint{9.057228in}{8.025958in}}%
\pgfpathquadraticcurveto{\pgfqpoint{9.065816in}{8.025958in}}{\pgfqpoint{9.071926in}{8.022594in}}%
\pgfpathquadraticcurveto{\pgfqpoint{9.078057in}{8.019251in}}{\pgfqpoint{9.082064in}{8.012345in}}%
\pgfpathlineto{\pgfqpoint{9.082064in}{8.012345in}}%
\pgfpathclose%
\pgfpathmoveto{\pgfqpoint{9.038700in}{7.985296in}}%
\pgfpathquadraticcurveto{\pgfqpoint{9.038700in}{7.971262in}}{\pgfqpoint{9.044478in}{7.963271in}}%
\pgfpathquadraticcurveto{\pgfqpoint{9.050255in}{7.955280in}}{\pgfqpoint{9.060349in}{7.955280in}}%
\pgfpathquadraticcurveto{\pgfqpoint{9.070442in}{7.955280in}}{\pgfqpoint{9.076242in}{7.963271in}}%
\pgfpathquadraticcurveto{\pgfqpoint{9.082064in}{7.971262in}}{\pgfqpoint{9.082064in}{7.985296in}}%
\pgfpathquadraticcurveto{\pgfqpoint{9.082064in}{7.999329in}}{\pgfqpoint{9.076242in}{8.007320in}}%
\pgfpathquadraticcurveto{\pgfqpoint{9.070442in}{8.015311in}}{\pgfqpoint{9.060349in}{8.015311in}}%
\pgfpathquadraticcurveto{\pgfqpoint{9.050255in}{8.015311in}}{\pgfqpoint{9.044478in}{8.007320in}}%
\pgfpathquadraticcurveto{\pgfqpoint{9.038700in}{7.999329in}}{\pgfqpoint{9.038700in}{7.985296in}}%
\pgfpathlineto{\pgfqpoint{9.038700in}{7.985296in}}%
\pgfpathclose%
\pgfpathmoveto{\pgfqpoint{9.121022in}{8.024099in}}%
\pgfpathlineto{\pgfqpoint{9.133750in}{8.024099in}}%
\pgfpathlineto{\pgfqpoint{9.133750in}{7.946625in}}%
\pgfpathlineto{\pgfqpoint{9.121022in}{7.946625in}}%
\pgfpathlineto{\pgfqpoint{9.121022in}{8.024099in}}%
\pgfpathlineto{\pgfqpoint{9.121022in}{8.024099in}}%
\pgfpathclose%
\pgfpathmoveto{\pgfqpoint{9.121022in}{8.054270in}}%
\pgfpathlineto{\pgfqpoint{9.133750in}{8.054270in}}%
\pgfpathlineto{\pgfqpoint{9.133750in}{8.038133in}}%
\pgfpathlineto{\pgfqpoint{9.121022in}{8.038133in}}%
\pgfpathlineto{\pgfqpoint{9.121022in}{8.054270in}}%
\pgfpathlineto{\pgfqpoint{9.121022in}{8.054270in}}%
\pgfpathclose%
\pgfpathmoveto{\pgfqpoint{9.247327in}{8.054270in}}%
\pgfpathlineto{\pgfqpoint{9.247327in}{8.043667in}}%
\pgfpathlineto{\pgfqpoint{9.235152in}{8.043667in}}%
\pgfpathquadraticcurveto{\pgfqpoint{9.228312in}{8.043667in}}{\pgfqpoint{9.225634in}{8.040900in}}%
\pgfpathquadraticcurveto{\pgfqpoint{9.222978in}{8.038133in}}{\pgfqpoint{9.222978in}{8.030939in}}%
\pgfpathlineto{\pgfqpoint{9.222978in}{8.024099in}}%
\pgfpathlineto{\pgfqpoint{9.243940in}{8.024099in}}%
\pgfpathlineto{\pgfqpoint{9.243940in}{8.014204in}}%
\pgfpathlineto{\pgfqpoint{9.222978in}{8.014204in}}%
\pgfpathlineto{\pgfqpoint{9.222978in}{7.946625in}}%
\pgfpathlineto{\pgfqpoint{9.210183in}{7.946625in}}%
\pgfpathlineto{\pgfqpoint{9.210183in}{8.014204in}}%
\pgfpathlineto{\pgfqpoint{9.175254in}{8.014204in}}%
\pgfpathlineto{\pgfqpoint{9.175254in}{7.946625in}}%
\pgfpathlineto{\pgfqpoint{9.162459in}{7.946625in}}%
\pgfpathlineto{\pgfqpoint{9.162459in}{8.014204in}}%
\pgfpathlineto{\pgfqpoint{9.150285in}{8.014204in}}%
\pgfpathlineto{\pgfqpoint{9.150285in}{8.024099in}}%
\pgfpathlineto{\pgfqpoint{9.162459in}{8.024099in}}%
\pgfpathlineto{\pgfqpoint{9.162459in}{8.029500in}}%
\pgfpathquadraticcurveto{\pgfqpoint{9.162459in}{8.042427in}}{\pgfqpoint{9.168480in}{8.048337in}}%
\pgfpathquadraticcurveto{\pgfqpoint{9.174501in}{8.054270in}}{\pgfqpoint{9.187561in}{8.054270in}}%
\pgfpathlineto{\pgfqpoint{9.199603in}{8.054270in}}%
\pgfpathlineto{\pgfqpoint{9.199603in}{8.043667in}}%
\pgfpathlineto{\pgfqpoint{9.187428in}{8.043667in}}%
\pgfpathquadraticcurveto{\pgfqpoint{9.180588in}{8.043667in}}{\pgfqpoint{9.177888in}{8.040900in}}%
\pgfpathquadraticcurveto{\pgfqpoint{9.175254in}{8.038133in}}{\pgfqpoint{9.175254in}{8.030939in}}%
\pgfpathlineto{\pgfqpoint{9.175254in}{8.024099in}}%
\pgfpathlineto{\pgfqpoint{9.210183in}{8.024099in}}%
\pgfpathlineto{\pgfqpoint{9.210183in}{8.029500in}}%
\pgfpathquadraticcurveto{\pgfqpoint{9.210183in}{8.042427in}}{\pgfqpoint{9.216204in}{8.048337in}}%
\pgfpathquadraticcurveto{\pgfqpoint{9.222225in}{8.054270in}}{\pgfqpoint{9.235307in}{8.054270in}}%
\pgfpathlineto{\pgfqpoint{9.247327in}{8.054270in}}%
\pgfpathlineto{\pgfqpoint{9.247327in}{8.054270in}}%
\pgfpathclose%
\pgfpathmoveto{\pgfqpoint{9.324247in}{7.988549in}}%
\pgfpathlineto{\pgfqpoint{9.324247in}{7.982329in}}%
\pgfpathlineto{\pgfqpoint{9.265721in}{7.982329in}}%
\pgfpathquadraticcurveto{\pgfqpoint{9.266562in}{7.969181in}}{\pgfqpoint{9.273646in}{7.962297in}}%
\pgfpathquadraticcurveto{\pgfqpoint{9.280729in}{7.955413in}}{\pgfqpoint{9.293390in}{7.955413in}}%
\pgfpathquadraticcurveto{\pgfqpoint{9.300717in}{7.955413in}}{\pgfqpoint{9.307601in}{7.957206in}}%
\pgfpathquadraticcurveto{\pgfqpoint{9.314485in}{7.958999in}}{\pgfqpoint{9.321281in}{7.962607in}}%
\pgfpathlineto{\pgfqpoint{9.321281in}{7.950565in}}%
\pgfpathquadraticcurveto{\pgfqpoint{9.314419in}{7.947665in}}{\pgfqpoint{9.307225in}{7.946138in}}%
\pgfpathquadraticcurveto{\pgfqpoint{9.300031in}{7.944611in}}{\pgfqpoint{9.292638in}{7.944611in}}%
\pgfpathquadraticcurveto{\pgfqpoint{9.274088in}{7.944611in}}{\pgfqpoint{9.263264in}{7.955391in}}%
\pgfpathquadraticcurveto{\pgfqpoint{9.252440in}{7.966193in}}{\pgfqpoint{9.252440in}{7.984609in}}%
\pgfpathquadraticcurveto{\pgfqpoint{9.252440in}{8.003624in}}{\pgfqpoint{9.262711in}{8.014780in}}%
\pgfpathquadraticcurveto{\pgfqpoint{9.272982in}{8.025958in}}{\pgfqpoint{9.290424in}{8.025958in}}%
\pgfpathquadraticcurveto{\pgfqpoint{9.306052in}{8.025958in}}{\pgfqpoint{9.315150in}{8.015887in}}%
\pgfpathquadraticcurveto{\pgfqpoint{9.324247in}{8.005837in}}{\pgfqpoint{9.324247in}{7.988549in}}%
\pgfpathlineto{\pgfqpoint{9.324247in}{7.988549in}}%
\pgfpathclose%
\pgfpathmoveto{\pgfqpoint{9.311519in}{7.992290in}}%
\pgfpathquadraticcurveto{\pgfqpoint{9.311386in}{8.002716in}}{\pgfqpoint{9.305676in}{8.008936in}}%
\pgfpathquadraticcurveto{\pgfqpoint{9.299965in}{8.015178in}}{\pgfqpoint{9.290557in}{8.015178in}}%
\pgfpathquadraticcurveto{\pgfqpoint{9.279910in}{8.015178in}}{\pgfqpoint{9.273513in}{8.009158in}}%
\pgfpathquadraticcurveto{\pgfqpoint{9.267116in}{8.003137in}}{\pgfqpoint{9.266142in}{7.992202in}}%
\pgfpathlineto{\pgfqpoint{9.311519in}{7.992290in}}%
\pgfpathlineto{\pgfqpoint{9.311519in}{7.992290in}}%
\pgfpathclose%
\pgfpathmoveto{\pgfqpoint{9.390034in}{8.012212in}}%
\pgfpathquadraticcurveto{\pgfqpoint{9.387886in}{8.013452in}}{\pgfqpoint{9.385363in}{8.014027in}}%
\pgfpathquadraticcurveto{\pgfqpoint{9.382840in}{8.014625in}}{\pgfqpoint{9.379807in}{8.014625in}}%
\pgfpathquadraticcurveto{\pgfqpoint{9.369005in}{8.014625in}}{\pgfqpoint{9.363228in}{8.007608in}}%
\pgfpathquadraticcurveto{\pgfqpoint{9.357450in}{8.000591in}}{\pgfqpoint{9.357450in}{7.987443in}}%
\pgfpathlineto{\pgfqpoint{9.357450in}{7.946625in}}%
\pgfpathlineto{\pgfqpoint{9.344656in}{7.946625in}}%
\pgfpathlineto{\pgfqpoint{9.344656in}{8.024099in}}%
\pgfpathlineto{\pgfqpoint{9.357450in}{8.024099in}}%
\pgfpathlineto{\pgfqpoint{9.357450in}{8.012057in}}%
\pgfpathquadraticcurveto{\pgfqpoint{9.361479in}{8.019118in}}{\pgfqpoint{9.367898in}{8.022527in}}%
\pgfpathquadraticcurveto{\pgfqpoint{9.374340in}{8.025958in}}{\pgfqpoint{9.383548in}{8.025958in}}%
\pgfpathquadraticcurveto{\pgfqpoint{9.384854in}{8.025958in}}{\pgfqpoint{9.386448in}{8.025781in}}%
\pgfpathquadraticcurveto{\pgfqpoint{9.388041in}{8.025626in}}{\pgfqpoint{9.389967in}{8.025272in}}%
\pgfpathlineto{\pgfqpoint{9.390034in}{8.012212in}}%
\pgfpathlineto{\pgfqpoint{9.390034in}{8.012212in}}%
\pgfpathclose%
\pgfpathmoveto{\pgfqpoint{9.466534in}{7.988549in}}%
\pgfpathlineto{\pgfqpoint{9.466534in}{7.982329in}}%
\pgfpathlineto{\pgfqpoint{9.408008in}{7.982329in}}%
\pgfpathquadraticcurveto{\pgfqpoint{9.408849in}{7.969181in}}{\pgfqpoint{9.415932in}{7.962297in}}%
\pgfpathquadraticcurveto{\pgfqpoint{9.423015in}{7.955413in}}{\pgfqpoint{9.435677in}{7.955413in}}%
\pgfpathquadraticcurveto{\pgfqpoint{9.443004in}{7.955413in}}{\pgfqpoint{9.449888in}{7.957206in}}%
\pgfpathquadraticcurveto{\pgfqpoint{9.456772in}{7.958999in}}{\pgfqpoint{9.463567in}{7.962607in}}%
\pgfpathlineto{\pgfqpoint{9.463567in}{7.950565in}}%
\pgfpathquadraticcurveto{\pgfqpoint{9.456705in}{7.947665in}}{\pgfqpoint{9.449511in}{7.946138in}}%
\pgfpathquadraticcurveto{\pgfqpoint{9.442317in}{7.944611in}}{\pgfqpoint{9.434924in}{7.944611in}}%
\pgfpathquadraticcurveto{\pgfqpoint{9.416375in}{7.944611in}}{\pgfqpoint{9.405551in}{7.955391in}}%
\pgfpathquadraticcurveto{\pgfqpoint{9.394726in}{7.966193in}}{\pgfqpoint{9.394726in}{7.984609in}}%
\pgfpathquadraticcurveto{\pgfqpoint{9.394726in}{8.003624in}}{\pgfqpoint{9.404997in}{8.014780in}}%
\pgfpathquadraticcurveto{\pgfqpoint{9.415268in}{8.025958in}}{\pgfqpoint{9.432711in}{8.025958in}}%
\pgfpathquadraticcurveto{\pgfqpoint{9.448338in}{8.025958in}}{\pgfqpoint{9.457436in}{8.015887in}}%
\pgfpathquadraticcurveto{\pgfqpoint{9.466534in}{8.005837in}}{\pgfqpoint{9.466534in}{7.988549in}}%
\pgfpathlineto{\pgfqpoint{9.466534in}{7.988549in}}%
\pgfpathclose%
\pgfpathmoveto{\pgfqpoint{9.453806in}{7.992290in}}%
\pgfpathquadraticcurveto{\pgfqpoint{9.453673in}{8.002716in}}{\pgfqpoint{9.447962in}{8.008936in}}%
\pgfpathquadraticcurveto{\pgfqpoint{9.442251in}{8.015178in}}{\pgfqpoint{9.432844in}{8.015178in}}%
\pgfpathquadraticcurveto{\pgfqpoint{9.422196in}{8.015178in}}{\pgfqpoint{9.415799in}{8.009158in}}%
\pgfpathquadraticcurveto{\pgfqpoint{9.409402in}{8.003137in}}{\pgfqpoint{9.408428in}{7.992202in}}%
\pgfpathlineto{\pgfqpoint{9.453806in}{7.992290in}}%
\pgfpathlineto{\pgfqpoint{9.453806in}{7.992290in}}%
\pgfpathclose%
\pgfpathmoveto{\pgfqpoint{9.551844in}{7.993397in}}%
\pgfpathlineto{\pgfqpoint{9.551844in}{7.946625in}}%
\pgfpathlineto{\pgfqpoint{9.539116in}{7.946625in}}%
\pgfpathlineto{\pgfqpoint{9.539116in}{7.992977in}}%
\pgfpathquadraticcurveto{\pgfqpoint{9.539116in}{8.003978in}}{\pgfqpoint{9.534821in}{8.009423in}}%
\pgfpathquadraticcurveto{\pgfqpoint{9.530527in}{8.014891in}}{\pgfqpoint{9.521961in}{8.014891in}}%
\pgfpathquadraticcurveto{\pgfqpoint{9.511646in}{8.014891in}}{\pgfqpoint{9.505691in}{8.008316in}}%
\pgfpathquadraticcurveto{\pgfqpoint{9.499737in}{8.001764in}}{\pgfqpoint{9.499737in}{7.990409in}}%
\pgfpathlineto{\pgfqpoint{9.499737in}{7.946625in}}%
\pgfpathlineto{\pgfqpoint{9.486942in}{7.946625in}}%
\pgfpathlineto{\pgfqpoint{9.486942in}{8.024099in}}%
\pgfpathlineto{\pgfqpoint{9.499737in}{8.024099in}}%
\pgfpathlineto{\pgfqpoint{9.499737in}{8.012057in}}%
\pgfpathquadraticcurveto{\pgfqpoint{9.504319in}{8.019052in}}{\pgfqpoint{9.510495in}{8.022505in}}%
\pgfpathquadraticcurveto{\pgfqpoint{9.516692in}{8.025958in}}{\pgfqpoint{9.524794in}{8.025958in}}%
\pgfpathquadraticcurveto{\pgfqpoint{9.538142in}{8.025958in}}{\pgfqpoint{9.544982in}{8.017702in}}%
\pgfpathquadraticcurveto{\pgfqpoint{9.551844in}{8.009445in}}{\pgfqpoint{9.551844in}{7.993397in}}%
\pgfpathlineto{\pgfqpoint{9.551844in}{7.993397in}}%
\pgfpathclose%
\pgfpathmoveto{\pgfqpoint{9.632970in}{8.021133in}}%
\pgfpathlineto{\pgfqpoint{9.632970in}{8.009224in}}%
\pgfpathquadraticcurveto{\pgfqpoint{9.627569in}{8.012212in}}{\pgfqpoint{9.622146in}{8.013695in}}%
\pgfpathquadraticcurveto{\pgfqpoint{9.616722in}{8.015178in}}{\pgfqpoint{9.611189in}{8.015178in}}%
\pgfpathquadraticcurveto{\pgfqpoint{9.598793in}{8.015178in}}{\pgfqpoint{9.591931in}{8.007320in}}%
\pgfpathquadraticcurveto{\pgfqpoint{9.585091in}{7.999484in}}{\pgfqpoint{9.585091in}{7.985296in}}%
\pgfpathquadraticcurveto{\pgfqpoint{9.585091in}{7.971107in}}{\pgfqpoint{9.591931in}{7.963249in}}%
\pgfpathquadraticcurveto{\pgfqpoint{9.598793in}{7.955413in}}{\pgfqpoint{9.611189in}{7.955413in}}%
\pgfpathquadraticcurveto{\pgfqpoint{9.616722in}{7.955413in}}{\pgfqpoint{9.622146in}{7.956896in}}%
\pgfpathquadraticcurveto{\pgfqpoint{9.627569in}{7.958379in}}{\pgfqpoint{9.632970in}{7.961367in}}%
\pgfpathlineto{\pgfqpoint{9.632970in}{7.949591in}}%
\pgfpathquadraticcurveto{\pgfqpoint{9.627635in}{7.947112in}}{\pgfqpoint{9.621924in}{7.945872in}}%
\pgfpathquadraticcurveto{\pgfqpoint{9.616235in}{7.944611in}}{\pgfqpoint{9.609794in}{7.944611in}}%
\pgfpathquadraticcurveto{\pgfqpoint{9.592285in}{7.944611in}}{\pgfqpoint{9.581970in}{7.955612in}}%
\pgfpathquadraticcurveto{\pgfqpoint{9.571677in}{7.966613in}}{\pgfqpoint{9.571677in}{7.985296in}}%
\pgfpathquadraticcurveto{\pgfqpoint{9.571677in}{8.004243in}}{\pgfqpoint{9.582080in}{8.015090in}}%
\pgfpathquadraticcurveto{\pgfqpoint{9.592506in}{8.025958in}}{\pgfqpoint{9.610635in}{8.025958in}}%
\pgfpathquadraticcurveto{\pgfqpoint{9.616501in}{8.025958in}}{\pgfqpoint{9.622101in}{8.024741in}}%
\pgfpathquadraticcurveto{\pgfqpoint{9.627702in}{8.023546in}}{\pgfqpoint{9.632970in}{8.021133in}}%
\pgfpathlineto{\pgfqpoint{9.632970in}{8.021133in}}%
\pgfpathclose%
\pgfpathmoveto{\pgfqpoint{9.721379in}{7.988549in}}%
\pgfpathlineto{\pgfqpoint{9.721379in}{7.982329in}}%
\pgfpathlineto{\pgfqpoint{9.662853in}{7.982329in}}%
\pgfpathquadraticcurveto{\pgfqpoint{9.663694in}{7.969181in}}{\pgfqpoint{9.670777in}{7.962297in}}%
\pgfpathquadraticcurveto{\pgfqpoint{9.677860in}{7.955413in}}{\pgfqpoint{9.690522in}{7.955413in}}%
\pgfpathquadraticcurveto{\pgfqpoint{9.697849in}{7.955413in}}{\pgfqpoint{9.704733in}{7.957206in}}%
\pgfpathquadraticcurveto{\pgfqpoint{9.711617in}{7.958999in}}{\pgfqpoint{9.718413in}{7.962607in}}%
\pgfpathlineto{\pgfqpoint{9.718413in}{7.950565in}}%
\pgfpathquadraticcurveto{\pgfqpoint{9.711551in}{7.947665in}}{\pgfqpoint{9.704357in}{7.946138in}}%
\pgfpathquadraticcurveto{\pgfqpoint{9.697163in}{7.944611in}}{\pgfqpoint{9.689769in}{7.944611in}}%
\pgfpathquadraticcurveto{\pgfqpoint{9.671220in}{7.944611in}}{\pgfqpoint{9.660396in}{7.955391in}}%
\pgfpathquadraticcurveto{\pgfqpoint{9.649571in}{7.966193in}}{\pgfqpoint{9.649571in}{7.984609in}}%
\pgfpathquadraticcurveto{\pgfqpoint{9.649571in}{8.003624in}}{\pgfqpoint{9.659842in}{8.014780in}}%
\pgfpathquadraticcurveto{\pgfqpoint{9.670113in}{8.025958in}}{\pgfqpoint{9.687556in}{8.025958in}}%
\pgfpathquadraticcurveto{\pgfqpoint{9.703183in}{8.025958in}}{\pgfqpoint{9.712281in}{8.015887in}}%
\pgfpathquadraticcurveto{\pgfqpoint{9.721379in}{8.005837in}}{\pgfqpoint{9.721379in}{7.988549in}}%
\pgfpathlineto{\pgfqpoint{9.721379in}{7.988549in}}%
\pgfpathclose%
\pgfpathmoveto{\pgfqpoint{9.708651in}{7.992290in}}%
\pgfpathquadraticcurveto{\pgfqpoint{9.708518in}{8.002716in}}{\pgfqpoint{9.702807in}{8.008936in}}%
\pgfpathquadraticcurveto{\pgfqpoint{9.697096in}{8.015178in}}{\pgfqpoint{9.687689in}{8.015178in}}%
\pgfpathquadraticcurveto{\pgfqpoint{9.677041in}{8.015178in}}{\pgfqpoint{9.670644in}{8.009158in}}%
\pgfpathquadraticcurveto{\pgfqpoint{9.664247in}{8.003137in}}{\pgfqpoint{9.663273in}{7.992202in}}%
\pgfpathlineto{\pgfqpoint{9.708651in}{7.992290in}}%
\pgfpathlineto{\pgfqpoint{9.708651in}{7.992290in}}%
\pgfpathclose%
\pgfpathmoveto{\pgfqpoint{9.791659in}{8.021819in}}%
\pgfpathlineto{\pgfqpoint{9.791659in}{8.009777in}}%
\pgfpathquadraticcurveto{\pgfqpoint{9.786280in}{8.012544in}}{\pgfqpoint{9.780458in}{8.013917in}}%
\pgfpathquadraticcurveto{\pgfqpoint{9.774659in}{8.015311in}}{\pgfqpoint{9.768416in}{8.015311in}}%
\pgfpathquadraticcurveto{\pgfqpoint{9.758942in}{8.015311in}}{\pgfqpoint{9.754205in}{8.012411in}}%
\pgfpathquadraticcurveto{\pgfqpoint{9.749469in}{8.009512in}}{\pgfqpoint{9.749469in}{8.003690in}}%
\pgfpathquadraticcurveto{\pgfqpoint{9.749469in}{7.999263in}}{\pgfqpoint{9.752855in}{7.996740in}}%
\pgfpathquadraticcurveto{\pgfqpoint{9.756242in}{7.994216in}}{\pgfqpoint{9.766491in}{7.991936in}}%
\pgfpathlineto{\pgfqpoint{9.770851in}{7.990962in}}%
\pgfpathquadraticcurveto{\pgfqpoint{9.784398in}{7.988062in}}{\pgfqpoint{9.790109in}{7.982772in}}%
\pgfpathquadraticcurveto{\pgfqpoint{9.795820in}{7.977482in}}{\pgfqpoint{9.795820in}{7.968008in}}%
\pgfpathquadraticcurveto{\pgfqpoint{9.795820in}{7.957206in}}{\pgfqpoint{9.787276in}{7.950897in}}%
\pgfpathquadraticcurveto{\pgfqpoint{9.778732in}{7.944611in}}{\pgfqpoint{9.763790in}{7.944611in}}%
\pgfpathquadraticcurveto{\pgfqpoint{9.757570in}{7.944611in}}{\pgfqpoint{9.750819in}{7.945828in}}%
\pgfpathquadraticcurveto{\pgfqpoint{9.744067in}{7.947046in}}{\pgfqpoint{9.736608in}{7.949458in}}%
\pgfpathlineto{\pgfqpoint{9.736608in}{7.962607in}}%
\pgfpathquadraticcurveto{\pgfqpoint{9.743669in}{7.958932in}}{\pgfqpoint{9.750509in}{7.957095in}}%
\pgfpathquadraticcurveto{\pgfqpoint{9.757349in}{7.955280in}}{\pgfqpoint{9.764078in}{7.955280in}}%
\pgfpathquadraticcurveto{\pgfqpoint{9.773065in}{7.955280in}}{\pgfqpoint{9.777890in}{7.958357in}}%
\pgfpathquadraticcurveto{\pgfqpoint{9.782738in}{7.961434in}}{\pgfqpoint{9.782738in}{7.967034in}}%
\pgfpathquadraticcurveto{\pgfqpoint{9.782738in}{7.972214in}}{\pgfqpoint{9.779241in}{7.974980in}}%
\pgfpathquadraticcurveto{\pgfqpoint{9.775765in}{7.977747in}}{\pgfqpoint{9.763923in}{7.980315in}}%
\pgfpathlineto{\pgfqpoint{9.759496in}{7.981355in}}%
\pgfpathquadraticcurveto{\pgfqpoint{9.747676in}{7.983835in}}{\pgfqpoint{9.742407in}{7.988992in}}%
\pgfpathquadraticcurveto{\pgfqpoint{9.737161in}{7.994150in}}{\pgfqpoint{9.737161in}{8.003137in}}%
\pgfpathquadraticcurveto{\pgfqpoint{9.737161in}{8.014072in}}{\pgfqpoint{9.744909in}{8.020004in}}%
\pgfpathquadraticcurveto{\pgfqpoint{9.752656in}{8.025958in}}{\pgfqpoint{9.766911in}{8.025958in}}%
\pgfpathquadraticcurveto{\pgfqpoint{9.773950in}{8.025958in}}{\pgfqpoint{9.780170in}{8.024918in}}%
\pgfpathquadraticcurveto{\pgfqpoint{9.786413in}{8.023900in}}{\pgfqpoint{9.791659in}{8.021819in}}%
\pgfpathlineto{\pgfqpoint{9.791659in}{8.021819in}}%
\pgfpathclose%
\pgfpathmoveto{\pgfqpoint{9.819328in}{7.964201in}}%
\pgfpathlineto{\pgfqpoint{9.833915in}{7.964201in}}%
\pgfpathlineto{\pgfqpoint{9.833915in}{7.946625in}}%
\pgfpathlineto{\pgfqpoint{9.819328in}{7.946625in}}%
\pgfpathlineto{\pgfqpoint{9.819328in}{7.964201in}}%
\pgfpathlineto{\pgfqpoint{9.819328in}{7.964201in}}%
\pgfpathclose%
\pgfpathmoveto{\pgfqpoint{9.819328in}{8.019871in}}%
\pgfpathlineto{\pgfqpoint{9.833915in}{8.019871in}}%
\pgfpathlineto{\pgfqpoint{9.833915in}{8.002318in}}%
\pgfpathlineto{\pgfqpoint{9.819328in}{8.002318in}}%
\pgfpathlineto{\pgfqpoint{9.819328in}{8.019871in}}%
\pgfpathlineto{\pgfqpoint{9.819328in}{8.019871in}}%
\pgfpathclose%
\pgfpathmoveto{\pgfqpoint{9.913049in}{7.958379in}}%
\pgfpathlineto{\pgfqpoint{9.935871in}{7.958379in}}%
\pgfpathlineto{\pgfqpoint{9.935871in}{8.037181in}}%
\pgfpathlineto{\pgfqpoint{9.911035in}{8.032201in}}%
\pgfpathlineto{\pgfqpoint{9.911035in}{8.044928in}}%
\pgfpathlineto{\pgfqpoint{9.935738in}{8.049909in}}%
\pgfpathlineto{\pgfqpoint{9.949705in}{8.049909in}}%
\pgfpathlineto{\pgfqpoint{9.949705in}{7.958379in}}%
\pgfpathlineto{\pgfqpoint{9.972527in}{7.958379in}}%
\pgfpathlineto{\pgfqpoint{9.972527in}{7.946625in}}%
\pgfpathlineto{\pgfqpoint{9.913049in}{7.946625in}}%
\pgfpathlineto{\pgfqpoint{9.913049in}{7.958379in}}%
\pgfpathlineto{\pgfqpoint{9.913049in}{7.958379in}}%
\pgfpathclose%
\pgfpathmoveto{\pgfqpoint{10.039155in}{8.037734in}}%
\pgfpathlineto{\pgfqpoint{10.003871in}{7.982595in}}%
\pgfpathlineto{\pgfqpoint{10.039155in}{7.982595in}}%
\pgfpathlineto{\pgfqpoint{10.039155in}{8.037734in}}%
\pgfpathlineto{\pgfqpoint{10.039155in}{8.037734in}}%
\pgfpathclose%
\pgfpathmoveto{\pgfqpoint{10.035480in}{8.049909in}}%
\pgfpathlineto{\pgfqpoint{10.053056in}{8.049909in}}%
\pgfpathlineto{\pgfqpoint{10.053056in}{7.982595in}}%
\pgfpathlineto{\pgfqpoint{10.067798in}{7.982595in}}%
\pgfpathlineto{\pgfqpoint{10.067798in}{7.970974in}}%
\pgfpathlineto{\pgfqpoint{10.053056in}{7.970974in}}%
\pgfpathlineto{\pgfqpoint{10.053056in}{7.946625in}}%
\pgfpathlineto{\pgfqpoint{10.039155in}{7.946625in}}%
\pgfpathlineto{\pgfqpoint{10.039155in}{7.970974in}}%
\pgfpathlineto{\pgfqpoint{9.992538in}{7.970974in}}%
\pgfpathlineto{\pgfqpoint{9.992538in}{7.984454in}}%
\pgfpathlineto{\pgfqpoint{10.035480in}{8.049909in}}%
\pgfpathlineto{\pgfqpoint{10.035480in}{8.049909in}}%
\pgfpathclose%
\pgfpathmoveto{\pgfqpoint{10.090885in}{7.964201in}}%
\pgfpathlineto{\pgfqpoint{10.105495in}{7.964201in}}%
\pgfpathlineto{\pgfqpoint{10.105495in}{7.946625in}}%
\pgfpathlineto{\pgfqpoint{10.090885in}{7.946625in}}%
\pgfpathlineto{\pgfqpoint{10.090885in}{7.964201in}}%
\pgfpathlineto{\pgfqpoint{10.090885in}{7.964201in}}%
\pgfpathclose%
\pgfpathmoveto{\pgfqpoint{10.167540in}{8.003823in}}%
\pgfpathquadraticcurveto{\pgfqpoint{10.158133in}{8.003823in}}{\pgfqpoint{10.152621in}{7.997382in}}%
\pgfpathquadraticcurveto{\pgfqpoint{10.147131in}{7.990962in}}{\pgfqpoint{10.147131in}{7.979762in}}%
\pgfpathquadraticcurveto{\pgfqpoint{10.147131in}{7.968628in}}{\pgfqpoint{10.152621in}{7.962142in}}%
\pgfpathquadraticcurveto{\pgfqpoint{10.158133in}{7.955678in}}{\pgfqpoint{10.167540in}{7.955678in}}%
\pgfpathquadraticcurveto{\pgfqpoint{10.176948in}{7.955678in}}{\pgfqpoint{10.182437in}{7.962142in}}%
\pgfpathquadraticcurveto{\pgfqpoint{10.187927in}{7.968628in}}{\pgfqpoint{10.187927in}{7.979762in}}%
\pgfpathquadraticcurveto{\pgfqpoint{10.187927in}{7.990962in}}{\pgfqpoint{10.182437in}{7.997382in}}%
\pgfpathquadraticcurveto{\pgfqpoint{10.176948in}{8.003823in}}{\pgfqpoint{10.167540in}{8.003823in}}%
\pgfpathlineto{\pgfqpoint{10.167540in}{8.003823in}}%
\pgfpathclose%
\pgfpathmoveto{\pgfqpoint{10.195276in}{8.047629in}}%
\pgfpathlineto{\pgfqpoint{10.195276in}{8.034901in}}%
\pgfpathquadraticcurveto{\pgfqpoint{10.190008in}{8.037380in}}{\pgfqpoint{10.184651in}{8.038686in}}%
\pgfpathquadraticcurveto{\pgfqpoint{10.179294in}{8.040014in}}{\pgfqpoint{10.174026in}{8.040014in}}%
\pgfpathquadraticcurveto{\pgfqpoint{10.160191in}{8.040014in}}{\pgfqpoint{10.152886in}{8.030673in}}%
\pgfpathquadraticcurveto{\pgfqpoint{10.145604in}{8.021332in}}{\pgfqpoint{10.144564in}{8.002451in}}%
\pgfpathquadraticcurveto{\pgfqpoint{10.148636in}{8.008471in}}{\pgfqpoint{10.154790in}{8.011681in}}%
\pgfpathquadraticcurveto{\pgfqpoint{10.160966in}{8.014891in}}{\pgfqpoint{10.168359in}{8.014891in}}%
\pgfpathquadraticcurveto{\pgfqpoint{10.183920in}{8.014891in}}{\pgfqpoint{10.192952in}{8.005439in}}%
\pgfpathquadraticcurveto{\pgfqpoint{10.201983in}{7.996009in}}{\pgfqpoint{10.201983in}{7.979762in}}%
\pgfpathquadraticcurveto{\pgfqpoint{10.201983in}{7.963846in}}{\pgfqpoint{10.192575in}{7.954217in}}%
\pgfpathquadraticcurveto{\pgfqpoint{10.183168in}{7.944611in}}{\pgfqpoint{10.167540in}{7.944611in}}%
\pgfpathquadraticcurveto{\pgfqpoint{10.149610in}{7.944611in}}{\pgfqpoint{10.140136in}{7.958335in}}%
\pgfpathquadraticcurveto{\pgfqpoint{10.130663in}{7.972081in}}{\pgfqpoint{10.130663in}{7.998156in}}%
\pgfpathquadraticcurveto{\pgfqpoint{10.130663in}{8.022638in}}{\pgfqpoint{10.142284in}{8.037203in}}%
\pgfpathquadraticcurveto{\pgfqpoint{10.153905in}{8.051768in}}{\pgfqpoint{10.173472in}{8.051768in}}%
\pgfpathquadraticcurveto{\pgfqpoint{10.178741in}{8.051768in}}{\pgfqpoint{10.184097in}{8.050728in}}%
\pgfpathquadraticcurveto{\pgfqpoint{10.189454in}{8.049688in}}{\pgfqpoint{10.195276in}{8.047629in}}%
\pgfpathlineto{\pgfqpoint{10.195276in}{8.047629in}}%
\pgfpathclose%
\pgfpathmoveto{\pgfqpoint{10.313900in}{7.992069in}}%
\pgfpathquadraticcurveto{\pgfqpoint{10.307879in}{7.992069in}}{\pgfqpoint{10.304448in}{7.986956in}}%
\pgfpathquadraticcurveto{\pgfqpoint{10.301039in}{7.981842in}}{\pgfqpoint{10.301039in}{7.972701in}}%
\pgfpathquadraticcurveto{\pgfqpoint{10.301039in}{7.963714in}}{\pgfqpoint{10.304448in}{7.958556in}}%
\pgfpathquadraticcurveto{\pgfqpoint{10.307879in}{7.953398in}}{\pgfqpoint{10.313900in}{7.953398in}}%
\pgfpathquadraticcurveto{\pgfqpoint{10.319788in}{7.953398in}}{\pgfqpoint{10.323196in}{7.958556in}}%
\pgfpathquadraticcurveto{\pgfqpoint{10.326627in}{7.963714in}}{\pgfqpoint{10.326627in}{7.972701in}}%
\pgfpathquadraticcurveto{\pgfqpoint{10.326627in}{7.981776in}}{\pgfqpoint{10.323196in}{7.986911in}}%
\pgfpathquadraticcurveto{\pgfqpoint{10.319788in}{7.992069in}}{\pgfqpoint{10.313900in}{7.992069in}}%
\pgfpathlineto{\pgfqpoint{10.313900in}{7.992069in}}%
\pgfpathclose%
\pgfpathmoveto{\pgfqpoint{10.313900in}{8.000857in}}%
\pgfpathquadraticcurveto{\pgfqpoint{10.324834in}{8.000857in}}{\pgfqpoint{10.331254in}{7.993242in}}%
\pgfpathquadraticcurveto{\pgfqpoint{10.337695in}{7.985650in}}{\pgfqpoint{10.337695in}{7.972701in}}%
\pgfpathquadraticcurveto{\pgfqpoint{10.337695in}{7.959773in}}{\pgfqpoint{10.331232in}{7.952181in}}%
\pgfpathquadraticcurveto{\pgfqpoint{10.324768in}{7.944611in}}{\pgfqpoint{10.313900in}{7.944611in}}%
\pgfpathquadraticcurveto{\pgfqpoint{10.302832in}{7.944611in}}{\pgfqpoint{10.296390in}{7.952181in}}%
\pgfpathquadraticcurveto{\pgfqpoint{10.289971in}{7.959773in}}{\pgfqpoint{10.289971in}{7.972701in}}%
\pgfpathquadraticcurveto{\pgfqpoint{10.289971in}{7.985716in}}{\pgfqpoint{10.296435in}{7.993286in}}%
\pgfpathquadraticcurveto{\pgfqpoint{10.302898in}{8.000857in}}{\pgfqpoint{10.313900in}{8.000857in}}%
\pgfpathlineto{\pgfqpoint{10.313900in}{8.000857in}}%
\pgfpathclose%
\pgfpathmoveto{\pgfqpoint{10.242513in}{8.042980in}}%
\pgfpathquadraticcurveto{\pgfqpoint{10.236558in}{8.042980in}}{\pgfqpoint{10.233127in}{8.037823in}}%
\pgfpathquadraticcurveto{\pgfqpoint{10.229719in}{8.032688in}}{\pgfqpoint{10.229719in}{8.023678in}}%
\pgfpathquadraticcurveto{\pgfqpoint{10.229719in}{8.014559in}}{\pgfqpoint{10.233105in}{8.009423in}}%
\pgfpathquadraticcurveto{\pgfqpoint{10.236492in}{8.004310in}}{\pgfqpoint{10.242513in}{8.004310in}}%
\pgfpathquadraticcurveto{\pgfqpoint{10.248534in}{8.004310in}}{\pgfqpoint{10.251942in}{8.009423in}}%
\pgfpathquadraticcurveto{\pgfqpoint{10.255373in}{8.014559in}}{\pgfqpoint{10.255373in}{8.023678in}}%
\pgfpathquadraticcurveto{\pgfqpoint{10.255373in}{8.032599in}}{\pgfqpoint{10.251920in}{8.037779in}}%
\pgfpathquadraticcurveto{\pgfqpoint{10.248467in}{8.042980in}}{\pgfqpoint{10.242513in}{8.042980in}}%
\pgfpathlineto{\pgfqpoint{10.242513in}{8.042980in}}%
\pgfpathclose%
\pgfpathmoveto{\pgfqpoint{10.304979in}{8.051768in}}%
\pgfpathlineto{\pgfqpoint{10.316047in}{8.051768in}}%
\pgfpathlineto{\pgfqpoint{10.251433in}{7.944611in}}%
\pgfpathlineto{\pgfqpoint{10.240366in}{7.944611in}}%
\pgfpathlineto{\pgfqpoint{10.304979in}{8.051768in}}%
\pgfpathlineto{\pgfqpoint{10.304979in}{8.051768in}}%
\pgfpathclose%
\pgfpathmoveto{\pgfqpoint{10.242513in}{8.051768in}}%
\pgfpathquadraticcurveto{\pgfqpoint{10.253448in}{8.051768in}}{\pgfqpoint{10.259933in}{8.044198in}}%
\pgfpathquadraticcurveto{\pgfqpoint{10.266441in}{8.036628in}}{\pgfqpoint{10.266441in}{8.023678in}}%
\pgfpathquadraticcurveto{\pgfqpoint{10.266441in}{8.010618in}}{\pgfqpoint{10.259978in}{8.003070in}}%
\pgfpathquadraticcurveto{\pgfqpoint{10.253514in}{7.995522in}}{\pgfqpoint{10.242513in}{7.995522in}}%
\pgfpathquadraticcurveto{\pgfqpoint{10.231511in}{7.995522in}}{\pgfqpoint{10.225114in}{8.003092in}}%
\pgfpathquadraticcurveto{\pgfqpoint{10.218717in}{8.010685in}}{\pgfqpoint{10.218717in}{8.023678in}}%
\pgfpathquadraticcurveto{\pgfqpoint{10.218717in}{8.036561in}}{\pgfqpoint{10.225136in}{8.044154in}}%
\pgfpathquadraticcurveto{\pgfqpoint{10.231578in}{8.051768in}}{\pgfqpoint{10.242513in}{8.051768in}}%
\pgfpathlineto{\pgfqpoint{10.242513in}{8.051768in}}%
\pgfpathclose%
\pgfusepath{fill}%
\end{pgfscope}%
\begin{pgfscope}%
\pgfsetbuttcap%
\pgfsetroundjoin%
\definecolor{currentfill}{rgb}{0.090196,0.168627,0.227451}%
\pgfsetfillcolor{currentfill}%
\pgfsetlinewidth{0.000000pt}%
\definecolor{currentstroke}{rgb}{0.090196,0.168627,0.227451}%
\pgfsetstrokecolor{currentstroke}%
\pgfsetdash{}{0pt}%
\pgfpathmoveto{\pgfqpoint{8.300905in}{7.826939in}}%
\pgfpathlineto{\pgfqpoint{8.300905in}{7.813303in}}%
\pgfpathquadraticcurveto{\pgfqpoint{8.292958in}{7.817111in}}{\pgfqpoint{8.285897in}{7.818970in}}%
\pgfpathquadraticcurveto{\pgfqpoint{8.278836in}{7.820852in}}{\pgfqpoint{8.272261in}{7.820852in}}%
\pgfpathquadraticcurveto{\pgfqpoint{8.260862in}{7.820852in}}{\pgfqpoint{8.254664in}{7.816424in}}%
\pgfpathquadraticcurveto{\pgfqpoint{8.248466in}{7.811997in}}{\pgfqpoint{8.248466in}{7.803829in}}%
\pgfpathquadraticcurveto{\pgfqpoint{8.248466in}{7.796967in}}{\pgfqpoint{8.252583in}{7.793470in}}%
\pgfpathquadraticcurveto{\pgfqpoint{8.256700in}{7.789995in}}{\pgfqpoint{8.268189in}{7.787848in}}%
\pgfpathlineto{\pgfqpoint{8.276622in}{7.786121in}}%
\pgfpathquadraticcurveto{\pgfqpoint{8.292250in}{7.783133in}}{\pgfqpoint{8.299687in}{7.775629in}}%
\pgfpathquadraticcurveto{\pgfqpoint{8.307125in}{7.768125in}}{\pgfqpoint{8.307125in}{7.755552in}}%
\pgfpathquadraticcurveto{\pgfqpoint{8.307125in}{7.740522in}}{\pgfqpoint{8.297053in}{7.732775in}}%
\pgfpathquadraticcurveto{\pgfqpoint{8.287004in}{7.725027in}}{\pgfqpoint{8.267569in}{7.725027in}}%
\pgfpathquadraticcurveto{\pgfqpoint{8.260242in}{7.725027in}}{\pgfqpoint{8.251963in}{7.726687in}}%
\pgfpathquadraticcurveto{\pgfqpoint{8.243707in}{7.728348in}}{\pgfqpoint{8.234853in}{7.731602in}}%
\pgfpathlineto{\pgfqpoint{8.234853in}{7.745990in}}%
\pgfpathquadraticcurveto{\pgfqpoint{8.243353in}{7.741230in}}{\pgfqpoint{8.251521in}{7.738796in}}%
\pgfpathquadraticcurveto{\pgfqpoint{8.259689in}{7.736383in}}{\pgfqpoint{8.267569in}{7.736383in}}%
\pgfpathquadraticcurveto{\pgfqpoint{8.279522in}{7.736383in}}{\pgfqpoint{8.286030in}{7.741076in}}%
\pgfpathquadraticcurveto{\pgfqpoint{8.292538in}{7.745790in}}{\pgfqpoint{8.292538in}{7.754512in}}%
\pgfpathquadraticcurveto{\pgfqpoint{8.292538in}{7.762104in}}{\pgfqpoint{8.287867in}{7.766398in}}%
\pgfpathquadraticcurveto{\pgfqpoint{8.283196in}{7.770693in}}{\pgfqpoint{8.272549in}{7.772840in}}%
\pgfpathlineto{\pgfqpoint{8.264027in}{7.774500in}}%
\pgfpathquadraticcurveto{\pgfqpoint{8.248400in}{7.777599in}}{\pgfqpoint{8.241405in}{7.784240in}}%
\pgfpathquadraticcurveto{\pgfqpoint{8.234432in}{7.790880in}}{\pgfqpoint{8.234432in}{7.802723in}}%
\pgfpathquadraticcurveto{\pgfqpoint{8.234432in}{7.816424in}}{\pgfqpoint{8.244083in}{7.824305in}}%
\pgfpathquadraticcurveto{\pgfqpoint{8.253734in}{7.832185in}}{\pgfqpoint{8.270668in}{7.832185in}}%
\pgfpathquadraticcurveto{\pgfqpoint{8.277950in}{7.832185in}}{\pgfqpoint{8.285476in}{7.830857in}}%
\pgfpathquadraticcurveto{\pgfqpoint{8.293025in}{7.829551in}}{\pgfqpoint{8.300905in}{7.826939in}}%
\pgfpathlineto{\pgfqpoint{8.300905in}{7.826939in}}%
\pgfpathclose%
\pgfpathmoveto{\pgfqpoint{8.328375in}{7.804516in}}%
\pgfpathlineto{\pgfqpoint{8.341103in}{7.804516in}}%
\pgfpathlineto{\pgfqpoint{8.341103in}{7.727042in}}%
\pgfpathlineto{\pgfqpoint{8.328375in}{7.727042in}}%
\pgfpathlineto{\pgfqpoint{8.328375in}{7.804516in}}%
\pgfpathlineto{\pgfqpoint{8.328375in}{7.804516in}}%
\pgfpathclose%
\pgfpathmoveto{\pgfqpoint{8.328375in}{7.834686in}}%
\pgfpathlineto{\pgfqpoint{8.341103in}{7.834686in}}%
\pgfpathlineto{\pgfqpoint{8.341103in}{7.818549in}}%
\pgfpathlineto{\pgfqpoint{8.328375in}{7.818549in}}%
\pgfpathlineto{\pgfqpoint{8.328375in}{7.834686in}}%
\pgfpathlineto{\pgfqpoint{8.328375in}{7.834686in}}%
\pgfpathclose%
\pgfpathmoveto{\pgfqpoint{8.418709in}{7.766686in}}%
\pgfpathquadraticcurveto{\pgfqpoint{8.418709in}{7.780521in}}{\pgfqpoint{8.412998in}{7.788113in}}%
\pgfpathquadraticcurveto{\pgfqpoint{8.407310in}{7.795728in}}{\pgfqpoint{8.396995in}{7.795728in}}%
\pgfpathquadraticcurveto{\pgfqpoint{8.386768in}{7.795728in}}{\pgfqpoint{8.381057in}{7.788113in}}%
\pgfpathquadraticcurveto{\pgfqpoint{8.375346in}{7.780521in}}{\pgfqpoint{8.375346in}{7.766686in}}%
\pgfpathquadraticcurveto{\pgfqpoint{8.375346in}{7.752918in}}{\pgfqpoint{8.381057in}{7.745303in}}%
\pgfpathquadraticcurveto{\pgfqpoint{8.386768in}{7.737689in}}{\pgfqpoint{8.396995in}{7.737689in}}%
\pgfpathquadraticcurveto{\pgfqpoint{8.407310in}{7.737689in}}{\pgfqpoint{8.412998in}{7.745303in}}%
\pgfpathquadraticcurveto{\pgfqpoint{8.418709in}{7.752918in}}{\pgfqpoint{8.418709in}{7.766686in}}%
\pgfpathlineto{\pgfqpoint{8.418709in}{7.766686in}}%
\pgfpathclose%
\pgfpathmoveto{\pgfqpoint{8.431437in}{7.736648in}}%
\pgfpathquadraticcurveto{\pgfqpoint{8.431437in}{7.716882in}}{\pgfqpoint{8.422650in}{7.707230in}}%
\pgfpathquadraticcurveto{\pgfqpoint{8.413884in}{7.697579in}}{\pgfqpoint{8.395755in}{7.697579in}}%
\pgfpathquadraticcurveto{\pgfqpoint{8.389048in}{7.697579in}}{\pgfqpoint{8.383094in}{7.698576in}}%
\pgfpathquadraticcurveto{\pgfqpoint{8.377139in}{7.699572in}}{\pgfqpoint{8.371539in}{7.701652in}}%
\pgfpathlineto{\pgfqpoint{8.371539in}{7.714026in}}%
\pgfpathquadraticcurveto{\pgfqpoint{8.377139in}{7.710993in}}{\pgfqpoint{8.382607in}{7.709555in}}%
\pgfpathquadraticcurveto{\pgfqpoint{8.388074in}{7.708094in}}{\pgfqpoint{8.393741in}{7.708094in}}%
\pgfpathquadraticcurveto{\pgfqpoint{8.406269in}{7.708094in}}{\pgfqpoint{8.412489in}{7.714624in}}%
\pgfpathquadraticcurveto{\pgfqpoint{8.418709in}{7.721154in}}{\pgfqpoint{8.418709in}{7.734368in}}%
\pgfpathlineto{\pgfqpoint{8.418709in}{7.740677in}}%
\pgfpathquadraticcurveto{\pgfqpoint{8.414769in}{7.733815in}}{\pgfqpoint{8.408616in}{7.730428in}}%
\pgfpathquadraticcurveto{\pgfqpoint{8.402462in}{7.727042in}}{\pgfqpoint{8.393873in}{7.727042in}}%
\pgfpathquadraticcurveto{\pgfqpoint{8.379640in}{7.727042in}}{\pgfqpoint{8.370919in}{7.737888in}}%
\pgfpathquadraticcurveto{\pgfqpoint{8.362198in}{7.748757in}}{\pgfqpoint{8.362198in}{7.766686in}}%
\pgfpathquadraticcurveto{\pgfqpoint{8.362198in}{7.784660in}}{\pgfqpoint{8.370919in}{7.795507in}}%
\pgfpathquadraticcurveto{\pgfqpoint{8.379640in}{7.806375in}}{\pgfqpoint{8.393873in}{7.806375in}}%
\pgfpathquadraticcurveto{\pgfqpoint{8.402462in}{7.806375in}}{\pgfqpoint{8.408616in}{7.802988in}}%
\pgfpathquadraticcurveto{\pgfqpoint{8.414769in}{7.799602in}}{\pgfqpoint{8.418709in}{7.792762in}}%
\pgfpathlineto{\pgfqpoint{8.418709in}{7.804516in}}%
\pgfpathlineto{\pgfqpoint{8.431437in}{7.804516in}}%
\pgfpathlineto{\pgfqpoint{8.431437in}{7.736648in}}%
\pgfpathlineto{\pgfqpoint{8.431437in}{7.736648in}}%
\pgfpathclose%
\pgfpathmoveto{\pgfqpoint{8.522082in}{7.773814in}}%
\pgfpathlineto{\pgfqpoint{8.522082in}{7.727042in}}%
\pgfpathlineto{\pgfqpoint{8.509354in}{7.727042in}}%
\pgfpathlineto{\pgfqpoint{8.509354in}{7.773393in}}%
\pgfpathquadraticcurveto{\pgfqpoint{8.509354in}{7.784395in}}{\pgfqpoint{8.505060in}{7.789840in}}%
\pgfpathquadraticcurveto{\pgfqpoint{8.500765in}{7.795307in}}{\pgfqpoint{8.492199in}{7.795307in}}%
\pgfpathquadraticcurveto{\pgfqpoint{8.481884in}{7.795307in}}{\pgfqpoint{8.475929in}{7.788733in}}%
\pgfpathquadraticcurveto{\pgfqpoint{8.469975in}{7.782181in}}{\pgfqpoint{8.469975in}{7.770826in}}%
\pgfpathlineto{\pgfqpoint{8.469975in}{7.727042in}}%
\pgfpathlineto{\pgfqpoint{8.457181in}{7.727042in}}%
\pgfpathlineto{\pgfqpoint{8.457181in}{7.804516in}}%
\pgfpathlineto{\pgfqpoint{8.469975in}{7.804516in}}%
\pgfpathlineto{\pgfqpoint{8.469975in}{7.792474in}}%
\pgfpathquadraticcurveto{\pgfqpoint{8.474557in}{7.799469in}}{\pgfqpoint{8.480733in}{7.802922in}}%
\pgfpathquadraticcurveto{\pgfqpoint{8.486931in}{7.806375in}}{\pgfqpoint{8.495032in}{7.806375in}}%
\pgfpathquadraticcurveto{\pgfqpoint{8.508380in}{7.806375in}}{\pgfqpoint{8.515220in}{7.798118in}}%
\pgfpathquadraticcurveto{\pgfqpoint{8.522082in}{7.789862in}}{\pgfqpoint{8.522082in}{7.773814in}}%
\pgfpathlineto{\pgfqpoint{8.522082in}{7.773814in}}%
\pgfpathclose%
\pgfpathmoveto{\pgfqpoint{8.547449in}{7.804516in}}%
\pgfpathlineto{\pgfqpoint{8.560177in}{7.804516in}}%
\pgfpathlineto{\pgfqpoint{8.560177in}{7.727042in}}%
\pgfpathlineto{\pgfqpoint{8.547449in}{7.727042in}}%
\pgfpathlineto{\pgfqpoint{8.547449in}{7.804516in}}%
\pgfpathlineto{\pgfqpoint{8.547449in}{7.804516in}}%
\pgfpathclose%
\pgfpathmoveto{\pgfqpoint{8.547449in}{7.834686in}}%
\pgfpathlineto{\pgfqpoint{8.560177in}{7.834686in}}%
\pgfpathlineto{\pgfqpoint{8.560177in}{7.818549in}}%
\pgfpathlineto{\pgfqpoint{8.547449in}{7.818549in}}%
\pgfpathlineto{\pgfqpoint{8.547449in}{7.834686in}}%
\pgfpathlineto{\pgfqpoint{8.547449in}{7.834686in}}%
\pgfpathclose%
\pgfpathmoveto{\pgfqpoint{8.649405in}{7.804516in}}%
\pgfpathlineto{\pgfqpoint{8.649405in}{7.727042in}}%
\pgfpathlineto{\pgfqpoint{8.636610in}{7.727042in}}%
\pgfpathlineto{\pgfqpoint{8.636610in}{7.794621in}}%
\pgfpathlineto{\pgfqpoint{8.601681in}{7.794621in}}%
\pgfpathlineto{\pgfqpoint{8.601681in}{7.727042in}}%
\pgfpathlineto{\pgfqpoint{8.588886in}{7.727042in}}%
\pgfpathlineto{\pgfqpoint{8.588886in}{7.794621in}}%
\pgfpathlineto{\pgfqpoint{8.576712in}{7.794621in}}%
\pgfpathlineto{\pgfqpoint{8.576712in}{7.804516in}}%
\pgfpathlineto{\pgfqpoint{8.588886in}{7.804516in}}%
\pgfpathlineto{\pgfqpoint{8.588886in}{7.809917in}}%
\pgfpathquadraticcurveto{\pgfqpoint{8.588886in}{7.822578in}}{\pgfqpoint{8.594863in}{7.828621in}}%
\pgfpathquadraticcurveto{\pgfqpoint{8.600862in}{7.834686in}}{\pgfqpoint{8.613235in}{7.834686in}}%
\pgfpathlineto{\pgfqpoint{8.626030in}{7.834686in}}%
\pgfpathlineto{\pgfqpoint{8.626030in}{7.824083in}}%
\pgfpathlineto{\pgfqpoint{8.613855in}{7.824083in}}%
\pgfpathquadraticcurveto{\pgfqpoint{8.607015in}{7.824083in}}{\pgfqpoint{8.604337in}{7.821316in}}%
\pgfpathquadraticcurveto{\pgfqpoint{8.601681in}{7.818549in}}{\pgfqpoint{8.601681in}{7.811355in}}%
\pgfpathlineto{\pgfqpoint{8.601681in}{7.804516in}}%
\pgfpathlineto{\pgfqpoint{8.649405in}{7.804516in}}%
\pgfpathlineto{\pgfqpoint{8.649405in}{7.804516in}}%
\pgfpathclose%
\pgfpathmoveto{\pgfqpoint{8.636610in}{7.834531in}}%
\pgfpathlineto{\pgfqpoint{8.649405in}{7.834531in}}%
\pgfpathlineto{\pgfqpoint{8.649405in}{7.818417in}}%
\pgfpathlineto{\pgfqpoint{8.636610in}{7.818417in}}%
\pgfpathlineto{\pgfqpoint{8.636610in}{7.834531in}}%
\pgfpathlineto{\pgfqpoint{8.636610in}{7.834531in}}%
\pgfpathclose%
\pgfpathmoveto{\pgfqpoint{8.731793in}{7.801549in}}%
\pgfpathlineto{\pgfqpoint{8.731793in}{7.789641in}}%
\pgfpathquadraticcurveto{\pgfqpoint{8.726392in}{7.792629in}}{\pgfqpoint{8.720969in}{7.794112in}}%
\pgfpathquadraticcurveto{\pgfqpoint{8.715545in}{7.795595in}}{\pgfqpoint{8.710011in}{7.795595in}}%
\pgfpathquadraticcurveto{\pgfqpoint{8.697616in}{7.795595in}}{\pgfqpoint{8.690754in}{7.787737in}}%
\pgfpathquadraticcurveto{\pgfqpoint{8.683914in}{7.779901in}}{\pgfqpoint{8.683914in}{7.765712in}}%
\pgfpathquadraticcurveto{\pgfqpoint{8.683914in}{7.751523in}}{\pgfqpoint{8.690754in}{7.743665in}}%
\pgfpathquadraticcurveto{\pgfqpoint{8.697616in}{7.735829in}}{\pgfqpoint{8.710011in}{7.735829in}}%
\pgfpathquadraticcurveto{\pgfqpoint{8.715545in}{7.735829in}}{\pgfqpoint{8.720969in}{7.737312in}}%
\pgfpathquadraticcurveto{\pgfqpoint{8.726392in}{7.738796in}}{\pgfqpoint{8.731793in}{7.741784in}}%
\pgfpathlineto{\pgfqpoint{8.731793in}{7.730008in}}%
\pgfpathquadraticcurveto{\pgfqpoint{8.726458in}{7.727529in}}{\pgfqpoint{8.720747in}{7.726289in}}%
\pgfpathquadraticcurveto{\pgfqpoint{8.715058in}{7.725027in}}{\pgfqpoint{8.708617in}{7.725027in}}%
\pgfpathquadraticcurveto{\pgfqpoint{8.691108in}{7.725027in}}{\pgfqpoint{8.680793in}{7.736029in}}%
\pgfpathquadraticcurveto{\pgfqpoint{8.670500in}{7.747030in}}{\pgfqpoint{8.670500in}{7.765712in}}%
\pgfpathquadraticcurveto{\pgfqpoint{8.670500in}{7.784660in}}{\pgfqpoint{8.680903in}{7.795507in}}%
\pgfpathquadraticcurveto{\pgfqpoint{8.691329in}{7.806375in}}{\pgfqpoint{8.709458in}{7.806375in}}%
\pgfpathquadraticcurveto{\pgfqpoint{8.715324in}{7.806375in}}{\pgfqpoint{8.720924in}{7.805158in}}%
\pgfpathquadraticcurveto{\pgfqpoint{8.726525in}{7.803962in}}{\pgfqpoint{8.731793in}{7.801549in}}%
\pgfpathlineto{\pgfqpoint{8.731793in}{7.801549in}}%
\pgfpathclose%
\pgfpathmoveto{\pgfqpoint{8.789146in}{7.765978in}}%
\pgfpathquadraticcurveto{\pgfqpoint{8.773717in}{7.765978in}}{\pgfqpoint{8.767763in}{7.762458in}}%
\pgfpathquadraticcurveto{\pgfqpoint{8.761808in}{7.758939in}}{\pgfqpoint{8.761808in}{7.750417in}}%
\pgfpathquadraticcurveto{\pgfqpoint{8.761808in}{7.743643in}}{\pgfqpoint{8.766280in}{7.739659in}}%
\pgfpathquadraticcurveto{\pgfqpoint{8.770751in}{7.735697in}}{\pgfqpoint{8.778410in}{7.735697in}}%
\pgfpathquadraticcurveto{\pgfqpoint{8.789013in}{7.735697in}}{\pgfqpoint{8.795410in}{7.743201in}}%
\pgfpathquadraticcurveto{\pgfqpoint{8.801807in}{7.750704in}}{\pgfqpoint{8.801807in}{7.763145in}}%
\pgfpathlineto{\pgfqpoint{8.801807in}{7.765978in}}%
\pgfpathlineto{\pgfqpoint{8.789146in}{7.765978in}}%
\pgfpathlineto{\pgfqpoint{8.789146in}{7.765978in}}%
\pgfpathclose%
\pgfpathmoveto{\pgfqpoint{8.814535in}{7.771246in}}%
\pgfpathlineto{\pgfqpoint{8.814535in}{7.727042in}}%
\pgfpathlineto{\pgfqpoint{8.801807in}{7.727042in}}%
\pgfpathlineto{\pgfqpoint{8.801807in}{7.738796in}}%
\pgfpathquadraticcurveto{\pgfqpoint{8.797446in}{7.731757in}}{\pgfqpoint{8.790939in}{7.728392in}}%
\pgfpathquadraticcurveto{\pgfqpoint{8.784431in}{7.725027in}}{\pgfqpoint{8.775023in}{7.725027in}}%
\pgfpathquadraticcurveto{\pgfqpoint{8.763136in}{7.725027in}}{\pgfqpoint{8.756097in}{7.731712in}}%
\pgfpathquadraticcurveto{\pgfqpoint{8.749080in}{7.738397in}}{\pgfqpoint{8.749080in}{7.749598in}}%
\pgfpathquadraticcurveto{\pgfqpoint{8.749080in}{7.762658in}}{\pgfqpoint{8.757824in}{7.769298in}}%
\pgfpathquadraticcurveto{\pgfqpoint{8.766590in}{7.775939in}}{\pgfqpoint{8.783944in}{7.775939in}}%
\pgfpathlineto{\pgfqpoint{8.801807in}{7.775939in}}%
\pgfpathlineto{\pgfqpoint{8.801807in}{7.777201in}}%
\pgfpathquadraticcurveto{\pgfqpoint{8.801807in}{7.785988in}}{\pgfqpoint{8.796030in}{7.790792in}}%
\pgfpathquadraticcurveto{\pgfqpoint{8.790252in}{7.795595in}}{\pgfqpoint{8.779804in}{7.795595in}}%
\pgfpathquadraticcurveto{\pgfqpoint{8.773164in}{7.795595in}}{\pgfqpoint{8.766855in}{7.794001in}}%
\pgfpathquadraticcurveto{\pgfqpoint{8.760569in}{7.792408in}}{\pgfqpoint{8.754769in}{7.789220in}}%
\pgfpathlineto{\pgfqpoint{8.754769in}{7.800996in}}%
\pgfpathquadraticcurveto{\pgfqpoint{8.761742in}{7.803697in}}{\pgfqpoint{8.768316in}{7.805025in}}%
\pgfpathquadraticcurveto{\pgfqpoint{8.774890in}{7.806375in}}{\pgfqpoint{8.781110in}{7.806375in}}%
\pgfpathquadraticcurveto{\pgfqpoint{8.797933in}{7.806375in}}{\pgfqpoint{8.806234in}{7.797654in}}%
\pgfpathquadraticcurveto{\pgfqpoint{8.814535in}{7.788954in}}{\pgfqpoint{8.814535in}{7.771246in}}%
\pgfpathlineto{\pgfqpoint{8.814535in}{7.771246in}}%
\pgfpathclose%
\pgfpathmoveto{\pgfqpoint{8.905157in}{7.773814in}}%
\pgfpathlineto{\pgfqpoint{8.905157in}{7.727042in}}%
\pgfpathlineto{\pgfqpoint{8.892429in}{7.727042in}}%
\pgfpathlineto{\pgfqpoint{8.892429in}{7.773393in}}%
\pgfpathquadraticcurveto{\pgfqpoint{8.892429in}{7.784395in}}{\pgfqpoint{8.888135in}{7.789840in}}%
\pgfpathquadraticcurveto{\pgfqpoint{8.883841in}{7.795307in}}{\pgfqpoint{8.875275in}{7.795307in}}%
\pgfpathquadraticcurveto{\pgfqpoint{8.864959in}{7.795307in}}{\pgfqpoint{8.859005in}{7.788733in}}%
\pgfpathquadraticcurveto{\pgfqpoint{8.853051in}{7.782181in}}{\pgfqpoint{8.853051in}{7.770826in}}%
\pgfpathlineto{\pgfqpoint{8.853051in}{7.727042in}}%
\pgfpathlineto{\pgfqpoint{8.840256in}{7.727042in}}%
\pgfpathlineto{\pgfqpoint{8.840256in}{7.804516in}}%
\pgfpathlineto{\pgfqpoint{8.853051in}{7.804516in}}%
\pgfpathlineto{\pgfqpoint{8.853051in}{7.792474in}}%
\pgfpathquadraticcurveto{\pgfqpoint{8.857633in}{7.799469in}}{\pgfqpoint{8.863808in}{7.802922in}}%
\pgfpathquadraticcurveto{\pgfqpoint{8.870006in}{7.806375in}}{\pgfqpoint{8.878108in}{7.806375in}}%
\pgfpathquadraticcurveto{\pgfqpoint{8.891455in}{7.806375in}}{\pgfqpoint{8.898295in}{7.798118in}}%
\pgfpathquadraticcurveto{\pgfqpoint{8.905157in}{7.789862in}}{\pgfqpoint{8.905157in}{7.773814in}}%
\pgfpathlineto{\pgfqpoint{8.905157in}{7.773814in}}%
\pgfpathclose%
\pgfpathmoveto{\pgfqpoint{8.943120in}{7.826518in}}%
\pgfpathlineto{\pgfqpoint{8.943120in}{7.804516in}}%
\pgfpathlineto{\pgfqpoint{8.969328in}{7.804516in}}%
\pgfpathlineto{\pgfqpoint{8.969328in}{7.794621in}}%
\pgfpathlineto{\pgfqpoint{8.943120in}{7.794621in}}%
\pgfpathlineto{\pgfqpoint{8.943120in}{7.752564in}}%
\pgfpathquadraticcurveto{\pgfqpoint{8.943120in}{7.743090in}}{\pgfqpoint{8.945709in}{7.740389in}}%
\pgfpathquadraticcurveto{\pgfqpoint{8.948299in}{7.737689in}}{\pgfqpoint{8.956268in}{7.737689in}}%
\pgfpathlineto{\pgfqpoint{8.969328in}{7.737689in}}%
\pgfpathlineto{\pgfqpoint{8.969328in}{7.727042in}}%
\pgfpathlineto{\pgfqpoint{8.956268in}{7.727042in}}%
\pgfpathquadraticcurveto{\pgfqpoint{8.941526in}{7.727042in}}{\pgfqpoint{8.935926in}{7.732531in}}%
\pgfpathquadraticcurveto{\pgfqpoint{8.930325in}{7.738043in}}{\pgfqpoint{8.930325in}{7.752564in}}%
\pgfpathlineto{\pgfqpoint{8.930325in}{7.794621in}}%
\pgfpathlineto{\pgfqpoint{8.920984in}{7.794621in}}%
\pgfpathlineto{\pgfqpoint{8.920984in}{7.804516in}}%
\pgfpathlineto{\pgfqpoint{8.930325in}{7.804516in}}%
\pgfpathlineto{\pgfqpoint{8.930325in}{7.826518in}}%
\pgfpathlineto{\pgfqpoint{8.943120in}{7.826518in}}%
\pgfpathlineto{\pgfqpoint{8.943120in}{7.826518in}}%
\pgfpathclose%
\pgfpathmoveto{\pgfqpoint{9.082905in}{7.815871in}}%
\pgfpathlineto{\pgfqpoint{9.082905in}{7.777333in}}%
\pgfpathlineto{\pgfqpoint{9.121420in}{7.777333in}}%
\pgfpathlineto{\pgfqpoint{9.121420in}{7.765579in}}%
\pgfpathlineto{\pgfqpoint{9.082905in}{7.765579in}}%
\pgfpathlineto{\pgfqpoint{9.082905in}{7.727042in}}%
\pgfpathlineto{\pgfqpoint{9.071284in}{7.727042in}}%
\pgfpathlineto{\pgfqpoint{9.071284in}{7.765579in}}%
\pgfpathlineto{\pgfqpoint{9.032746in}{7.765579in}}%
\pgfpathlineto{\pgfqpoint{9.032746in}{7.777333in}}%
\pgfpathlineto{\pgfqpoint{9.071284in}{7.777333in}}%
\pgfpathlineto{\pgfqpoint{9.071284in}{7.815871in}}%
\pgfpathlineto{\pgfqpoint{9.082905in}{7.815871in}}%
\pgfpathlineto{\pgfqpoint{9.082905in}{7.815871in}}%
\pgfpathclose%
\pgfpathmoveto{\pgfqpoint{9.244360in}{7.775474in}}%
\pgfpathquadraticcurveto{\pgfqpoint{9.248854in}{7.773947in}}{\pgfqpoint{9.253104in}{7.768966in}}%
\pgfpathquadraticcurveto{\pgfqpoint{9.257354in}{7.763986in}}{\pgfqpoint{9.261648in}{7.755264in}}%
\pgfpathlineto{\pgfqpoint{9.275837in}{7.727042in}}%
\pgfpathlineto{\pgfqpoint{9.260807in}{7.727042in}}%
\pgfpathlineto{\pgfqpoint{9.247614in}{7.753538in}}%
\pgfpathquadraticcurveto{\pgfqpoint{9.242479in}{7.763919in}}{\pgfqpoint{9.237676in}{7.767306in}}%
\pgfpathquadraticcurveto{\pgfqpoint{9.232872in}{7.770693in}}{\pgfqpoint{9.224571in}{7.770693in}}%
\pgfpathlineto{\pgfqpoint{9.209342in}{7.770693in}}%
\pgfpathlineto{\pgfqpoint{9.209342in}{7.727042in}}%
\pgfpathlineto{\pgfqpoint{9.195375in}{7.727042in}}%
\pgfpathlineto{\pgfqpoint{9.195375in}{7.830326in}}%
\pgfpathlineto{\pgfqpoint{9.226918in}{7.830326in}}%
\pgfpathquadraticcurveto{\pgfqpoint{9.244626in}{7.830326in}}{\pgfqpoint{9.253347in}{7.822910in}}%
\pgfpathquadraticcurveto{\pgfqpoint{9.262069in}{7.815517in}}{\pgfqpoint{9.262069in}{7.800576in}}%
\pgfpathquadraticcurveto{\pgfqpoint{9.262069in}{7.790814in}}{\pgfqpoint{9.257531in}{7.784372in}}%
\pgfpathquadraticcurveto{\pgfqpoint{9.252993in}{7.777953in}}{\pgfqpoint{9.244360in}{7.775474in}}%
\pgfpathlineto{\pgfqpoint{9.244360in}{7.775474in}}%
\pgfpathclose%
\pgfpathmoveto{\pgfqpoint{9.209342in}{7.818837in}}%
\pgfpathlineto{\pgfqpoint{9.209342in}{7.782181in}}%
\pgfpathlineto{\pgfqpoint{9.226918in}{7.782181in}}%
\pgfpathquadraticcurveto{\pgfqpoint{9.237011in}{7.782181in}}{\pgfqpoint{9.242169in}{7.786852in}}%
\pgfpathquadraticcurveto{\pgfqpoint{9.247327in}{7.791522in}}{\pgfqpoint{9.247327in}{7.800576in}}%
\pgfpathquadraticcurveto{\pgfqpoint{9.247327in}{7.809629in}}{\pgfqpoint{9.242169in}{7.814233in}}%
\pgfpathquadraticcurveto{\pgfqpoint{9.237011in}{7.818837in}}{\pgfqpoint{9.226918in}{7.818837in}}%
\pgfpathlineto{\pgfqpoint{9.209342in}{7.818837in}}%
\pgfpathlineto{\pgfqpoint{9.209342in}{7.818837in}}%
\pgfpathclose%
\pgfpathmoveto{\pgfqpoint{9.353178in}{7.768966in}}%
\pgfpathlineto{\pgfqpoint{9.353178in}{7.762746in}}%
\pgfpathlineto{\pgfqpoint{9.294652in}{7.762746in}}%
\pgfpathquadraticcurveto{\pgfqpoint{9.295493in}{7.749598in}}{\pgfqpoint{9.302577in}{7.742714in}}%
\pgfpathquadraticcurveto{\pgfqpoint{9.309660in}{7.735829in}}{\pgfqpoint{9.322321in}{7.735829in}}%
\pgfpathquadraticcurveto{\pgfqpoint{9.329648in}{7.735829in}}{\pgfqpoint{9.336532in}{7.737622in}}%
\pgfpathquadraticcurveto{\pgfqpoint{9.343416in}{7.739415in}}{\pgfqpoint{9.350212in}{7.743023in}}%
\pgfpathlineto{\pgfqpoint{9.350212in}{7.730982in}}%
\pgfpathquadraticcurveto{\pgfqpoint{9.343350in}{7.728082in}}{\pgfqpoint{9.336156in}{7.726555in}}%
\pgfpathquadraticcurveto{\pgfqpoint{9.328962in}{7.725027in}}{\pgfqpoint{9.321569in}{7.725027in}}%
\pgfpathquadraticcurveto{\pgfqpoint{9.303019in}{7.725027in}}{\pgfqpoint{9.292195in}{7.735807in}}%
\pgfpathquadraticcurveto{\pgfqpoint{9.281371in}{7.746609in}}{\pgfqpoint{9.281371in}{7.765026in}}%
\pgfpathquadraticcurveto{\pgfqpoint{9.281371in}{7.784040in}}{\pgfqpoint{9.291642in}{7.795197in}}%
\pgfpathquadraticcurveto{\pgfqpoint{9.301913in}{7.806375in}}{\pgfqpoint{9.319355in}{7.806375in}}%
\pgfpathquadraticcurveto{\pgfqpoint{9.334983in}{7.806375in}}{\pgfqpoint{9.344080in}{7.796303in}}%
\pgfpathquadraticcurveto{\pgfqpoint{9.353178in}{7.786254in}}{\pgfqpoint{9.353178in}{7.768966in}}%
\pgfpathlineto{\pgfqpoint{9.353178in}{7.768966in}}%
\pgfpathclose%
\pgfpathmoveto{\pgfqpoint{9.340450in}{7.772707in}}%
\pgfpathquadraticcurveto{\pgfqpoint{9.340317in}{7.783133in}}{\pgfqpoint{9.334607in}{7.789353in}}%
\pgfpathquadraticcurveto{\pgfqpoint{9.328896in}{7.795595in}}{\pgfqpoint{9.319488in}{7.795595in}}%
\pgfpathquadraticcurveto{\pgfqpoint{9.308841in}{7.795595in}}{\pgfqpoint{9.302444in}{7.789574in}}%
\pgfpathquadraticcurveto{\pgfqpoint{9.296047in}{7.783553in}}{\pgfqpoint{9.295073in}{7.772618in}}%
\pgfpathlineto{\pgfqpoint{9.340450in}{7.772707in}}%
\pgfpathlineto{\pgfqpoint{9.340450in}{7.772707in}}%
\pgfpathclose%
\pgfpathmoveto{\pgfqpoint{9.364954in}{7.804516in}}%
\pgfpathlineto{\pgfqpoint{9.378435in}{7.804516in}}%
\pgfpathlineto{\pgfqpoint{9.402651in}{7.739504in}}%
\pgfpathlineto{\pgfqpoint{9.426867in}{7.804516in}}%
\pgfpathlineto{\pgfqpoint{9.440347in}{7.804516in}}%
\pgfpathlineto{\pgfqpoint{9.411284in}{7.727042in}}%
\pgfpathlineto{\pgfqpoint{9.393996in}{7.727042in}}%
\pgfpathlineto{\pgfqpoint{9.364954in}{7.804516in}}%
\pgfpathlineto{\pgfqpoint{9.364954in}{7.804516in}}%
\pgfpathclose%
\pgfpathmoveto{\pgfqpoint{9.524196in}{7.768966in}}%
\pgfpathlineto{\pgfqpoint{9.524196in}{7.762746in}}%
\pgfpathlineto{\pgfqpoint{9.465670in}{7.762746in}}%
\pgfpathquadraticcurveto{\pgfqpoint{9.466511in}{7.749598in}}{\pgfqpoint{9.473595in}{7.742714in}}%
\pgfpathquadraticcurveto{\pgfqpoint{9.480678in}{7.735829in}}{\pgfqpoint{9.493340in}{7.735829in}}%
\pgfpathquadraticcurveto{\pgfqpoint{9.500666in}{7.735829in}}{\pgfqpoint{9.507551in}{7.737622in}}%
\pgfpathquadraticcurveto{\pgfqpoint{9.514435in}{7.739415in}}{\pgfqpoint{9.521230in}{7.743023in}}%
\pgfpathlineto{\pgfqpoint{9.521230in}{7.730982in}}%
\pgfpathquadraticcurveto{\pgfqpoint{9.514368in}{7.728082in}}{\pgfqpoint{9.507174in}{7.726555in}}%
\pgfpathquadraticcurveto{\pgfqpoint{9.499980in}{7.725027in}}{\pgfqpoint{9.492587in}{7.725027in}}%
\pgfpathquadraticcurveto{\pgfqpoint{9.474038in}{7.725027in}}{\pgfqpoint{9.463213in}{7.735807in}}%
\pgfpathquadraticcurveto{\pgfqpoint{9.452389in}{7.746609in}}{\pgfqpoint{9.452389in}{7.765026in}}%
\pgfpathquadraticcurveto{\pgfqpoint{9.452389in}{7.784040in}}{\pgfqpoint{9.462660in}{7.795197in}}%
\pgfpathquadraticcurveto{\pgfqpoint{9.472931in}{7.806375in}}{\pgfqpoint{9.490373in}{7.806375in}}%
\pgfpathquadraticcurveto{\pgfqpoint{9.506001in}{7.806375in}}{\pgfqpoint{9.515099in}{7.796303in}}%
\pgfpathquadraticcurveto{\pgfqpoint{9.524196in}{7.786254in}}{\pgfqpoint{9.524196in}{7.768966in}}%
\pgfpathlineto{\pgfqpoint{9.524196in}{7.768966in}}%
\pgfpathclose%
\pgfpathmoveto{\pgfqpoint{9.511469in}{7.772707in}}%
\pgfpathquadraticcurveto{\pgfqpoint{9.511336in}{7.783133in}}{\pgfqpoint{9.505625in}{7.789353in}}%
\pgfpathquadraticcurveto{\pgfqpoint{9.499914in}{7.795595in}}{\pgfqpoint{9.490506in}{7.795595in}}%
\pgfpathquadraticcurveto{\pgfqpoint{9.479859in}{7.795595in}}{\pgfqpoint{9.473462in}{7.789574in}}%
\pgfpathquadraticcurveto{\pgfqpoint{9.467065in}{7.783553in}}{\pgfqpoint{9.466091in}{7.772618in}}%
\pgfpathlineto{\pgfqpoint{9.511469in}{7.772707in}}%
\pgfpathlineto{\pgfqpoint{9.511469in}{7.772707in}}%
\pgfpathclose%
\pgfpathmoveto{\pgfqpoint{9.589983in}{7.792629in}}%
\pgfpathquadraticcurveto{\pgfqpoint{9.587836in}{7.793868in}}{\pgfqpoint{9.585312in}{7.794444in}}%
\pgfpathquadraticcurveto{\pgfqpoint{9.582789in}{7.795042in}}{\pgfqpoint{9.579756in}{7.795042in}}%
\pgfpathquadraticcurveto{\pgfqpoint{9.568954in}{7.795042in}}{\pgfqpoint{9.563177in}{7.788025in}}%
\pgfpathquadraticcurveto{\pgfqpoint{9.557400in}{7.781008in}}{\pgfqpoint{9.557400in}{7.767859in}}%
\pgfpathlineto{\pgfqpoint{9.557400in}{7.727042in}}%
\pgfpathlineto{\pgfqpoint{9.544605in}{7.727042in}}%
\pgfpathlineto{\pgfqpoint{9.544605in}{7.804516in}}%
\pgfpathlineto{\pgfqpoint{9.557400in}{7.804516in}}%
\pgfpathlineto{\pgfqpoint{9.557400in}{7.792474in}}%
\pgfpathquadraticcurveto{\pgfqpoint{9.561428in}{7.799535in}}{\pgfqpoint{9.567847in}{7.802944in}}%
\pgfpathquadraticcurveto{\pgfqpoint{9.574289in}{7.806375in}}{\pgfqpoint{9.583497in}{7.806375in}}%
\pgfpathquadraticcurveto{\pgfqpoint{9.584803in}{7.806375in}}{\pgfqpoint{9.586397in}{7.806198in}}%
\pgfpathquadraticcurveto{\pgfqpoint{9.587991in}{7.806043in}}{\pgfqpoint{9.589916in}{7.805689in}}%
\pgfpathlineto{\pgfqpoint{9.589983in}{7.792629in}}%
\pgfpathlineto{\pgfqpoint{9.589983in}{7.792629in}}%
\pgfpathclose%
\pgfpathmoveto{\pgfqpoint{9.652715in}{7.802236in}}%
\pgfpathlineto{\pgfqpoint{9.652715in}{7.790194in}}%
\pgfpathquadraticcurveto{\pgfqpoint{9.647336in}{7.792961in}}{\pgfqpoint{9.641514in}{7.794333in}}%
\pgfpathquadraticcurveto{\pgfqpoint{9.635715in}{7.795728in}}{\pgfqpoint{9.629472in}{7.795728in}}%
\pgfpathquadraticcurveto{\pgfqpoint{9.619998in}{7.795728in}}{\pgfqpoint{9.615261in}{7.792828in}}%
\pgfpathquadraticcurveto{\pgfqpoint{9.610525in}{7.789928in}}{\pgfqpoint{9.610525in}{7.784107in}}%
\pgfpathquadraticcurveto{\pgfqpoint{9.610525in}{7.779680in}}{\pgfqpoint{9.613911in}{7.777156in}}%
\pgfpathquadraticcurveto{\pgfqpoint{9.617298in}{7.774633in}}{\pgfqpoint{9.627547in}{7.772353in}}%
\pgfpathlineto{\pgfqpoint{9.631907in}{7.771379in}}%
\pgfpathquadraticcurveto{\pgfqpoint{9.645454in}{7.768479in}}{\pgfqpoint{9.651165in}{7.763189in}}%
\pgfpathquadraticcurveto{\pgfqpoint{9.656876in}{7.757898in}}{\pgfqpoint{9.656876in}{7.748424in}}%
\pgfpathquadraticcurveto{\pgfqpoint{9.656876in}{7.737622in}}{\pgfqpoint{9.648332in}{7.731314in}}%
\pgfpathquadraticcurveto{\pgfqpoint{9.639788in}{7.725027in}}{\pgfqpoint{9.624846in}{7.725027in}}%
\pgfpathquadraticcurveto{\pgfqpoint{9.618626in}{7.725027in}}{\pgfqpoint{9.611875in}{7.726245in}}%
\pgfpathquadraticcurveto{\pgfqpoint{9.605123in}{7.727462in}}{\pgfqpoint{9.597664in}{7.729875in}}%
\pgfpathlineto{\pgfqpoint{9.597664in}{7.743023in}}%
\pgfpathquadraticcurveto{\pgfqpoint{9.604725in}{7.739349in}}{\pgfqpoint{9.611565in}{7.737512in}}%
\pgfpathquadraticcurveto{\pgfqpoint{9.618405in}{7.735697in}}{\pgfqpoint{9.625134in}{7.735697in}}%
\pgfpathquadraticcurveto{\pgfqpoint{9.634121in}{7.735697in}}{\pgfqpoint{9.638946in}{7.738773in}}%
\pgfpathquadraticcurveto{\pgfqpoint{9.643794in}{7.741850in}}{\pgfqpoint{9.643794in}{7.747451in}}%
\pgfpathquadraticcurveto{\pgfqpoint{9.643794in}{7.752630in}}{\pgfqpoint{9.640297in}{7.755397in}}%
\pgfpathquadraticcurveto{\pgfqpoint{9.636821in}{7.758164in}}{\pgfqpoint{9.624979in}{7.760732in}}%
\pgfpathlineto{\pgfqpoint{9.620552in}{7.761772in}}%
\pgfpathquadraticcurveto{\pgfqpoint{9.608732in}{7.764251in}}{\pgfqpoint{9.603463in}{7.769409in}}%
\pgfpathquadraticcurveto{\pgfqpoint{9.598217in}{7.774566in}}{\pgfqpoint{9.598217in}{7.783553in}}%
\pgfpathquadraticcurveto{\pgfqpoint{9.598217in}{7.794488in}}{\pgfqpoint{9.605965in}{7.800421in}}%
\pgfpathquadraticcurveto{\pgfqpoint{9.613712in}{7.806375in}}{\pgfqpoint{9.627967in}{7.806375in}}%
\pgfpathquadraticcurveto{\pgfqpoint{9.635006in}{7.806375in}}{\pgfqpoint{9.641226in}{7.805335in}}%
\pgfpathquadraticcurveto{\pgfqpoint{9.647469in}{7.804316in}}{\pgfqpoint{9.652715in}{7.802236in}}%
\pgfpathlineto{\pgfqpoint{9.652715in}{7.802236in}}%
\pgfpathclose%
\pgfpathmoveto{\pgfqpoint{9.712347in}{7.765978in}}%
\pgfpathquadraticcurveto{\pgfqpoint{9.696919in}{7.765978in}}{\pgfqpoint{9.690965in}{7.762458in}}%
\pgfpathquadraticcurveto{\pgfqpoint{9.685010in}{7.758939in}}{\pgfqpoint{9.685010in}{7.750417in}}%
\pgfpathquadraticcurveto{\pgfqpoint{9.685010in}{7.743643in}}{\pgfqpoint{9.689482in}{7.739659in}}%
\pgfpathquadraticcurveto{\pgfqpoint{9.693953in}{7.735697in}}{\pgfqpoint{9.701612in}{7.735697in}}%
\pgfpathquadraticcurveto{\pgfqpoint{9.712215in}{7.735697in}}{\pgfqpoint{9.718612in}{7.743201in}}%
\pgfpathquadraticcurveto{\pgfqpoint{9.725009in}{7.750704in}}{\pgfqpoint{9.725009in}{7.763145in}}%
\pgfpathlineto{\pgfqpoint{9.725009in}{7.765978in}}%
\pgfpathlineto{\pgfqpoint{9.712347in}{7.765978in}}%
\pgfpathlineto{\pgfqpoint{9.712347in}{7.765978in}}%
\pgfpathclose%
\pgfpathmoveto{\pgfqpoint{9.737737in}{7.771246in}}%
\pgfpathlineto{\pgfqpoint{9.737737in}{7.727042in}}%
\pgfpathlineto{\pgfqpoint{9.725009in}{7.727042in}}%
\pgfpathlineto{\pgfqpoint{9.725009in}{7.738796in}}%
\pgfpathquadraticcurveto{\pgfqpoint{9.720648in}{7.731757in}}{\pgfqpoint{9.714140in}{7.728392in}}%
\pgfpathquadraticcurveto{\pgfqpoint{9.707633in}{7.725027in}}{\pgfqpoint{9.698225in}{7.725027in}}%
\pgfpathquadraticcurveto{\pgfqpoint{9.686338in}{7.725027in}}{\pgfqpoint{9.679299in}{7.731712in}}%
\pgfpathquadraticcurveto{\pgfqpoint{9.672282in}{7.738397in}}{\pgfqpoint{9.672282in}{7.749598in}}%
\pgfpathquadraticcurveto{\pgfqpoint{9.672282in}{7.762658in}}{\pgfqpoint{9.681026in}{7.769298in}}%
\pgfpathquadraticcurveto{\pgfqpoint{9.689791in}{7.775939in}}{\pgfqpoint{9.707146in}{7.775939in}}%
\pgfpathlineto{\pgfqpoint{9.725009in}{7.775939in}}%
\pgfpathlineto{\pgfqpoint{9.725009in}{7.777201in}}%
\pgfpathquadraticcurveto{\pgfqpoint{9.725009in}{7.785988in}}{\pgfqpoint{9.719232in}{7.790792in}}%
\pgfpathquadraticcurveto{\pgfqpoint{9.713454in}{7.795595in}}{\pgfqpoint{9.703006in}{7.795595in}}%
\pgfpathquadraticcurveto{\pgfqpoint{9.696366in}{7.795595in}}{\pgfqpoint{9.690057in}{7.794001in}}%
\pgfpathquadraticcurveto{\pgfqpoint{9.683771in}{7.792408in}}{\pgfqpoint{9.677971in}{7.789220in}}%
\pgfpathlineto{\pgfqpoint{9.677971in}{7.800996in}}%
\pgfpathquadraticcurveto{\pgfqpoint{9.684944in}{7.803697in}}{\pgfqpoint{9.691518in}{7.805025in}}%
\pgfpathquadraticcurveto{\pgfqpoint{9.698092in}{7.806375in}}{\pgfqpoint{9.704312in}{7.806375in}}%
\pgfpathquadraticcurveto{\pgfqpoint{9.721135in}{7.806375in}}{\pgfqpoint{9.729436in}{7.797654in}}%
\pgfpathquadraticcurveto{\pgfqpoint{9.737737in}{7.788954in}}{\pgfqpoint{9.737737in}{7.771246in}}%
\pgfpathlineto{\pgfqpoint{9.737737in}{7.771246in}}%
\pgfpathclose%
\pgfpathmoveto{\pgfqpoint{9.763945in}{7.834686in}}%
\pgfpathlineto{\pgfqpoint{9.776673in}{7.834686in}}%
\pgfpathlineto{\pgfqpoint{9.776673in}{7.727042in}}%
\pgfpathlineto{\pgfqpoint{9.763945in}{7.727042in}}%
\pgfpathlineto{\pgfqpoint{9.763945in}{7.834686in}}%
\pgfpathlineto{\pgfqpoint{9.763945in}{7.834686in}}%
\pgfpathclose%
\pgfpathmoveto{\pgfqpoint{9.806556in}{7.744617in}}%
\pgfpathlineto{\pgfqpoint{9.821143in}{7.744617in}}%
\pgfpathlineto{\pgfqpoint{9.821143in}{7.727042in}}%
\pgfpathlineto{\pgfqpoint{9.806556in}{7.727042in}}%
\pgfpathlineto{\pgfqpoint{9.806556in}{7.744617in}}%
\pgfpathlineto{\pgfqpoint{9.806556in}{7.744617in}}%
\pgfpathclose%
\pgfpathmoveto{\pgfqpoint{9.806556in}{7.800288in}}%
\pgfpathlineto{\pgfqpoint{9.821143in}{7.800288in}}%
\pgfpathlineto{\pgfqpoint{9.821143in}{7.782734in}}%
\pgfpathlineto{\pgfqpoint{9.806556in}{7.782734in}}%
\pgfpathlineto{\pgfqpoint{9.806556in}{7.800288in}}%
\pgfpathlineto{\pgfqpoint{9.806556in}{7.800288in}}%
\pgfpathclose%
\pgfpathmoveto{\pgfqpoint{9.900277in}{7.738796in}}%
\pgfpathlineto{\pgfqpoint{9.923099in}{7.738796in}}%
\pgfpathlineto{\pgfqpoint{9.923099in}{7.817598in}}%
\pgfpathlineto{\pgfqpoint{9.898263in}{7.812617in}}%
\pgfpathlineto{\pgfqpoint{9.898263in}{7.825345in}}%
\pgfpathlineto{\pgfqpoint{9.922966in}{7.830326in}}%
\pgfpathlineto{\pgfqpoint{9.936933in}{7.830326in}}%
\pgfpathlineto{\pgfqpoint{9.936933in}{7.738796in}}%
\pgfpathlineto{\pgfqpoint{9.959755in}{7.738796in}}%
\pgfpathlineto{\pgfqpoint{9.959755in}{7.727042in}}%
\pgfpathlineto{\pgfqpoint{9.900277in}{7.727042in}}%
\pgfpathlineto{\pgfqpoint{9.900277in}{7.738796in}}%
\pgfpathlineto{\pgfqpoint{9.900277in}{7.738796in}}%
\pgfpathclose%
\pgfpathmoveto{\pgfqpoint{9.987978in}{7.744617in}}%
\pgfpathlineto{\pgfqpoint{10.002587in}{7.744617in}}%
\pgfpathlineto{\pgfqpoint{10.002587in}{7.727042in}}%
\pgfpathlineto{\pgfqpoint{9.987978in}{7.727042in}}%
\pgfpathlineto{\pgfqpoint{9.987978in}{7.744617in}}%
\pgfpathlineto{\pgfqpoint{9.987978in}{7.744617in}}%
\pgfpathclose%
\pgfpathmoveto{\pgfqpoint{10.062884in}{7.821117in}}%
\pgfpathquadraticcurveto{\pgfqpoint{10.052104in}{7.821117in}}{\pgfqpoint{10.046659in}{7.810492in}}%
\pgfpathquadraticcurveto{\pgfqpoint{10.041235in}{7.799889in}}{\pgfqpoint{10.041235in}{7.778573in}}%
\pgfpathquadraticcurveto{\pgfqpoint{10.041235in}{7.757345in}}{\pgfqpoint{10.046659in}{7.746720in}}%
\pgfpathquadraticcurveto{\pgfqpoint{10.052104in}{7.736095in}}{\pgfqpoint{10.062884in}{7.736095in}}%
\pgfpathquadraticcurveto{\pgfqpoint{10.073752in}{7.736095in}}{\pgfqpoint{10.079176in}{7.746720in}}%
\pgfpathquadraticcurveto{\pgfqpoint{10.084621in}{7.757345in}}{\pgfqpoint{10.084621in}{7.778573in}}%
\pgfpathquadraticcurveto{\pgfqpoint{10.084621in}{7.799889in}}{\pgfqpoint{10.079176in}{7.810492in}}%
\pgfpathquadraticcurveto{\pgfqpoint{10.073752in}{7.821117in}}{\pgfqpoint{10.062884in}{7.821117in}}%
\pgfpathlineto{\pgfqpoint{10.062884in}{7.821117in}}%
\pgfpathclose%
\pgfpathmoveto{\pgfqpoint{10.062884in}{7.832185in}}%
\pgfpathquadraticcurveto{\pgfqpoint{10.080260in}{7.832185in}}{\pgfqpoint{10.089424in}{7.818439in}}%
\pgfpathquadraticcurveto{\pgfqpoint{10.098588in}{7.804715in}}{\pgfqpoint{10.098588in}{7.778573in}}%
\pgfpathquadraticcurveto{\pgfqpoint{10.098588in}{7.752497in}}{\pgfqpoint{10.089424in}{7.738751in}}%
\pgfpathquadraticcurveto{\pgfqpoint{10.080260in}{7.725027in}}{\pgfqpoint{10.062884in}{7.725027in}}%
\pgfpathquadraticcurveto{\pgfqpoint{10.045530in}{7.725027in}}{\pgfqpoint{10.036366in}{7.738751in}}%
\pgfpathquadraticcurveto{\pgfqpoint{10.027202in}{7.752497in}}{\pgfqpoint{10.027202in}{7.778573in}}%
\pgfpathquadraticcurveto{\pgfqpoint{10.027202in}{7.804715in}}{\pgfqpoint{10.036366in}{7.818439in}}%
\pgfpathquadraticcurveto{\pgfqpoint{10.045530in}{7.832185in}}{\pgfqpoint{10.062884in}{7.832185in}}%
\pgfpathlineto{\pgfqpoint{10.062884in}{7.832185in}}%
\pgfpathclose%
\pgfpathmoveto{\pgfqpoint{10.210992in}{7.772486in}}%
\pgfpathquadraticcurveto{\pgfqpoint{10.204971in}{7.772486in}}{\pgfqpoint{10.201540in}{7.767372in}}%
\pgfpathquadraticcurveto{\pgfqpoint{10.198131in}{7.762259in}}{\pgfqpoint{10.198131in}{7.753117in}}%
\pgfpathquadraticcurveto{\pgfqpoint{10.198131in}{7.744130in}}{\pgfqpoint{10.201540in}{7.738973in}}%
\pgfpathquadraticcurveto{\pgfqpoint{10.204971in}{7.733815in}}{\pgfqpoint{10.210992in}{7.733815in}}%
\pgfpathquadraticcurveto{\pgfqpoint{10.216880in}{7.733815in}}{\pgfqpoint{10.220289in}{7.738973in}}%
\pgfpathquadraticcurveto{\pgfqpoint{10.223720in}{7.744130in}}{\pgfqpoint{10.223720in}{7.753117in}}%
\pgfpathquadraticcurveto{\pgfqpoint{10.223720in}{7.762193in}}{\pgfqpoint{10.220289in}{7.767328in}}%
\pgfpathquadraticcurveto{\pgfqpoint{10.216880in}{7.772486in}}{\pgfqpoint{10.210992in}{7.772486in}}%
\pgfpathlineto{\pgfqpoint{10.210992in}{7.772486in}}%
\pgfpathclose%
\pgfpathmoveto{\pgfqpoint{10.210992in}{7.781273in}}%
\pgfpathquadraticcurveto{\pgfqpoint{10.221927in}{7.781273in}}{\pgfqpoint{10.228346in}{7.773659in}}%
\pgfpathquadraticcurveto{\pgfqpoint{10.234788in}{7.766066in}}{\pgfqpoint{10.234788in}{7.753117in}}%
\pgfpathquadraticcurveto{\pgfqpoint{10.234788in}{7.740190in}}{\pgfqpoint{10.228324in}{7.732598in}}%
\pgfpathquadraticcurveto{\pgfqpoint{10.221860in}{7.725027in}}{\pgfqpoint{10.210992in}{7.725027in}}%
\pgfpathquadraticcurveto{\pgfqpoint{10.199924in}{7.725027in}}{\pgfqpoint{10.193483in}{7.732598in}}%
\pgfpathquadraticcurveto{\pgfqpoint{10.187064in}{7.740190in}}{\pgfqpoint{10.187064in}{7.753117in}}%
\pgfpathquadraticcurveto{\pgfqpoint{10.187064in}{7.766133in}}{\pgfqpoint{10.193527in}{7.773703in}}%
\pgfpathquadraticcurveto{\pgfqpoint{10.199991in}{7.781273in}}{\pgfqpoint{10.210992in}{7.781273in}}%
\pgfpathlineto{\pgfqpoint{10.210992in}{7.781273in}}%
\pgfpathclose%
\pgfpathmoveto{\pgfqpoint{10.139605in}{7.823397in}}%
\pgfpathquadraticcurveto{\pgfqpoint{10.133651in}{7.823397in}}{\pgfqpoint{10.130220in}{7.818240in}}%
\pgfpathquadraticcurveto{\pgfqpoint{10.126811in}{7.813104in}}{\pgfqpoint{10.126811in}{7.804095in}}%
\pgfpathquadraticcurveto{\pgfqpoint{10.126811in}{7.794975in}}{\pgfqpoint{10.130198in}{7.789840in}}%
\pgfpathquadraticcurveto{\pgfqpoint{10.133584in}{7.784727in}}{\pgfqpoint{10.139605in}{7.784727in}}%
\pgfpathquadraticcurveto{\pgfqpoint{10.145626in}{7.784727in}}{\pgfqpoint{10.149035in}{7.789840in}}%
\pgfpathquadraticcurveto{\pgfqpoint{10.152466in}{7.794975in}}{\pgfqpoint{10.152466in}{7.804095in}}%
\pgfpathquadraticcurveto{\pgfqpoint{10.152466in}{7.813016in}}{\pgfqpoint{10.149013in}{7.818195in}}%
\pgfpathquadraticcurveto{\pgfqpoint{10.145560in}{7.823397in}}{\pgfqpoint{10.139605in}{7.823397in}}%
\pgfpathlineto{\pgfqpoint{10.139605in}{7.823397in}}%
\pgfpathclose%
\pgfpathmoveto{\pgfqpoint{10.202071in}{7.832185in}}%
\pgfpathlineto{\pgfqpoint{10.213139in}{7.832185in}}%
\pgfpathlineto{\pgfqpoint{10.148526in}{7.725027in}}%
\pgfpathlineto{\pgfqpoint{10.137458in}{7.725027in}}%
\pgfpathlineto{\pgfqpoint{10.202071in}{7.832185in}}%
\pgfpathlineto{\pgfqpoint{10.202071in}{7.832185in}}%
\pgfpathclose%
\pgfpathmoveto{\pgfqpoint{10.139605in}{7.832185in}}%
\pgfpathquadraticcurveto{\pgfqpoint{10.150540in}{7.832185in}}{\pgfqpoint{10.157026in}{7.824615in}}%
\pgfpathquadraticcurveto{\pgfqpoint{10.163534in}{7.817044in}}{\pgfqpoint{10.163534in}{7.804095in}}%
\pgfpathquadraticcurveto{\pgfqpoint{10.163534in}{7.791035in}}{\pgfqpoint{10.157070in}{7.783487in}}%
\pgfpathquadraticcurveto{\pgfqpoint{10.150607in}{7.775939in}}{\pgfqpoint{10.139605in}{7.775939in}}%
\pgfpathquadraticcurveto{\pgfqpoint{10.128604in}{7.775939in}}{\pgfqpoint{10.122207in}{7.783509in}}%
\pgfpathquadraticcurveto{\pgfqpoint{10.115810in}{7.791102in}}{\pgfqpoint{10.115810in}{7.804095in}}%
\pgfpathquadraticcurveto{\pgfqpoint{10.115810in}{7.816978in}}{\pgfqpoint{10.122229in}{7.824570in}}%
\pgfpathquadraticcurveto{\pgfqpoint{10.128670in}{7.832185in}}{\pgfqpoint{10.139605in}{7.832185in}}%
\pgfpathlineto{\pgfqpoint{10.139605in}{7.832185in}}%
\pgfpathclose%
\pgfusepath{fill}%
\end{pgfscope}%
\begin{pgfscope}%
\pgfsetbuttcap%
\pgfsetroundjoin%
\definecolor{currentfill}{rgb}{0.541176,0.380392,0.658824}%
\pgfsetfillcolor{currentfill}%
\pgfsetlinewidth{1.003750pt}%
\definecolor{currentstroke}{rgb}{0.000000,0.000000,0.000000}%
\pgfsetstrokecolor{currentstroke}%
\pgfsetstrokeopacity{0.000000}%
\pgfsetdash{}{0pt}%
\pgfsys@defobject{currentmarker}{\pgfqpoint{-0.038194in}{-0.038194in}}{\pgfqpoint{0.038194in}{0.038194in}}{%
\pgfpathmoveto{\pgfqpoint{0.000000in}{-0.038194in}}%
\pgfpathcurveto{\pgfqpoint{0.010129in}{-0.038194in}}{\pgfqpoint{0.019845in}{-0.034170in}}{\pgfqpoint{0.027008in}{-0.027008in}}%
\pgfpathcurveto{\pgfqpoint{0.034170in}{-0.019845in}}{\pgfqpoint{0.038194in}{-0.010129in}}{\pgfqpoint{0.038194in}{0.000000in}}%
\pgfpathcurveto{\pgfqpoint{0.038194in}{0.010129in}}{\pgfqpoint{0.034170in}{0.019845in}}{\pgfqpoint{0.027008in}{0.027008in}}%
\pgfpathcurveto{\pgfqpoint{0.019845in}{0.034170in}}{\pgfqpoint{0.010129in}{0.038194in}}{\pgfqpoint{0.000000in}{0.038194in}}%
\pgfpathcurveto{\pgfqpoint{-0.010129in}{0.038194in}}{\pgfqpoint{-0.019845in}{0.034170in}}{\pgfqpoint{-0.027008in}{0.027008in}}%
\pgfpathcurveto{\pgfqpoint{-0.034170in}{0.019845in}}{\pgfqpoint{-0.038194in}{0.010129in}}{\pgfqpoint{-0.038194in}{0.000000in}}%
\pgfpathcurveto{\pgfqpoint{-0.038194in}{-0.010129in}}{\pgfqpoint{-0.034170in}{-0.019845in}}{\pgfqpoint{-0.027008in}{-0.027008in}}%
\pgfpathcurveto{\pgfqpoint{-0.019845in}{-0.034170in}}{\pgfqpoint{-0.010129in}{-0.038194in}}{\pgfqpoint{0.000000in}{-0.038194in}}%
\pgfpathlineto{\pgfqpoint{0.000000in}{-0.038194in}}%
\pgfpathclose%
\pgfusepath{stroke,fill}%
}%
\begin{pgfscope}%
\pgfsys@transformshift{1.078000in}{3.153600in}%
\pgfsys@useobject{currentmarker}{}%
\end{pgfscope}%
\end{pgfscope}%
\begin{pgfscope}%
\pgfsetbuttcap%
\pgfsetroundjoin%
\definecolor{currentfill}{rgb}{0.709804,0.741176,0.207843}%
\pgfsetfillcolor{currentfill}%
\pgfsetlinewidth{1.003750pt}%
\definecolor{currentstroke}{rgb}{0.000000,0.000000,0.000000}%
\pgfsetstrokecolor{currentstroke}%
\pgfsetstrokeopacity{0.000000}%
\pgfsetdash{}{0pt}%
\pgfsys@defobject{currentmarker}{\pgfqpoint{-0.038194in}{-0.038194in}}{\pgfqpoint{0.038194in}{0.038194in}}{%
\pgfpathmoveto{\pgfqpoint{0.000000in}{-0.038194in}}%
\pgfpathcurveto{\pgfqpoint{0.010129in}{-0.038194in}}{\pgfqpoint{0.019845in}{-0.034170in}}{\pgfqpoint{0.027008in}{-0.027008in}}%
\pgfpathcurveto{\pgfqpoint{0.034170in}{-0.019845in}}{\pgfqpoint{0.038194in}{-0.010129in}}{\pgfqpoint{0.038194in}{0.000000in}}%
\pgfpathcurveto{\pgfqpoint{0.038194in}{0.010129in}}{\pgfqpoint{0.034170in}{0.019845in}}{\pgfqpoint{0.027008in}{0.027008in}}%
\pgfpathcurveto{\pgfqpoint{0.019845in}{0.034170in}}{\pgfqpoint{0.010129in}{0.038194in}}{\pgfqpoint{0.000000in}{0.038194in}}%
\pgfpathcurveto{\pgfqpoint{-0.010129in}{0.038194in}}{\pgfqpoint{-0.019845in}{0.034170in}}{\pgfqpoint{-0.027008in}{0.027008in}}%
\pgfpathcurveto{\pgfqpoint{-0.034170in}{0.019845in}}{\pgfqpoint{-0.038194in}{0.010129in}}{\pgfqpoint{-0.038194in}{0.000000in}}%
\pgfpathcurveto{\pgfqpoint{-0.038194in}{-0.010129in}}{\pgfqpoint{-0.034170in}{-0.019845in}}{\pgfqpoint{-0.027008in}{-0.027008in}}%
\pgfpathcurveto{\pgfqpoint{-0.019845in}{-0.034170in}}{\pgfqpoint{-0.010129in}{-0.038194in}}{\pgfqpoint{0.000000in}{-0.038194in}}%
\pgfpathlineto{\pgfqpoint{0.000000in}{-0.038194in}}%
\pgfpathclose%
\pgfusepath{stroke,fill}%
}%
\begin{pgfscope}%
\pgfsys@transformshift{1.078000in}{2.754000in}%
\pgfsys@useobject{currentmarker}{}%
\end{pgfscope}%
\end{pgfscope}%
\begin{pgfscope}%
\pgfsetbuttcap%
\pgfsetroundjoin%
\definecolor{currentfill}{rgb}{0.000000,0.447059,0.698039}%
\pgfsetfillcolor{currentfill}%
\pgfsetlinewidth{1.003750pt}%
\definecolor{currentstroke}{rgb}{0.000000,0.000000,0.000000}%
\pgfsetstrokecolor{currentstroke}%
\pgfsetstrokeopacity{0.000000}%
\pgfsetdash{}{0pt}%
\pgfsys@defobject{currentmarker}{\pgfqpoint{-0.038194in}{-0.038194in}}{\pgfqpoint{0.038194in}{0.038194in}}{%
\pgfpathmoveto{\pgfqpoint{0.000000in}{-0.038194in}}%
\pgfpathcurveto{\pgfqpoint{0.010129in}{-0.038194in}}{\pgfqpoint{0.019845in}{-0.034170in}}{\pgfqpoint{0.027008in}{-0.027008in}}%
\pgfpathcurveto{\pgfqpoint{0.034170in}{-0.019845in}}{\pgfqpoint{0.038194in}{-0.010129in}}{\pgfqpoint{0.038194in}{0.000000in}}%
\pgfpathcurveto{\pgfqpoint{0.038194in}{0.010129in}}{\pgfqpoint{0.034170in}{0.019845in}}{\pgfqpoint{0.027008in}{0.027008in}}%
\pgfpathcurveto{\pgfqpoint{0.019845in}{0.034170in}}{\pgfqpoint{0.010129in}{0.038194in}}{\pgfqpoint{0.000000in}{0.038194in}}%
\pgfpathcurveto{\pgfqpoint{-0.010129in}{0.038194in}}{\pgfqpoint{-0.019845in}{0.034170in}}{\pgfqpoint{-0.027008in}{0.027008in}}%
\pgfpathcurveto{\pgfqpoint{-0.034170in}{0.019845in}}{\pgfqpoint{-0.038194in}{0.010129in}}{\pgfqpoint{-0.038194in}{0.000000in}}%
\pgfpathcurveto{\pgfqpoint{-0.038194in}{-0.010129in}}{\pgfqpoint{-0.034170in}{-0.019845in}}{\pgfqpoint{-0.027008in}{-0.027008in}}%
\pgfpathcurveto{\pgfqpoint{-0.019845in}{-0.034170in}}{\pgfqpoint{-0.010129in}{-0.038194in}}{\pgfqpoint{0.000000in}{-0.038194in}}%
\pgfpathlineto{\pgfqpoint{0.000000in}{-0.038194in}}%
\pgfpathclose%
\pgfusepath{stroke,fill}%
}%
\begin{pgfscope}%
\pgfsys@transformshift{1.078000in}{2.354400in}%
\pgfsys@useobject{currentmarker}{}%
\end{pgfscope}%
\end{pgfscope}%
\begin{pgfscope}%
\pgfsetbuttcap%
\pgfsetroundjoin%
\definecolor{currentfill}{rgb}{0.149020,0.274510,0.325490}%
\pgfsetfillcolor{currentfill}%
\pgfsetlinewidth{1.003750pt}%
\definecolor{currentstroke}{rgb}{0.000000,0.000000,0.000000}%
\pgfsetstrokecolor{currentstroke}%
\pgfsetstrokeopacity{0.000000}%
\pgfsetdash{}{0pt}%
\pgfsys@defobject{currentmarker}{\pgfqpoint{-0.038194in}{-0.038194in}}{\pgfqpoint{0.038194in}{0.038194in}}{%
\pgfpathmoveto{\pgfqpoint{0.000000in}{-0.038194in}}%
\pgfpathcurveto{\pgfqpoint{0.010129in}{-0.038194in}}{\pgfqpoint{0.019845in}{-0.034170in}}{\pgfqpoint{0.027008in}{-0.027008in}}%
\pgfpathcurveto{\pgfqpoint{0.034170in}{-0.019845in}}{\pgfqpoint{0.038194in}{-0.010129in}}{\pgfqpoint{0.038194in}{0.000000in}}%
\pgfpathcurveto{\pgfqpoint{0.038194in}{0.010129in}}{\pgfqpoint{0.034170in}{0.019845in}}{\pgfqpoint{0.027008in}{0.027008in}}%
\pgfpathcurveto{\pgfqpoint{0.019845in}{0.034170in}}{\pgfqpoint{0.010129in}{0.038194in}}{\pgfqpoint{0.000000in}{0.038194in}}%
\pgfpathcurveto{\pgfqpoint{-0.010129in}{0.038194in}}{\pgfqpoint{-0.019845in}{0.034170in}}{\pgfqpoint{-0.027008in}{0.027008in}}%
\pgfpathcurveto{\pgfqpoint{-0.034170in}{0.019845in}}{\pgfqpoint{-0.038194in}{0.010129in}}{\pgfqpoint{-0.038194in}{0.000000in}}%
\pgfpathcurveto{\pgfqpoint{-0.038194in}{-0.010129in}}{\pgfqpoint{-0.034170in}{-0.019845in}}{\pgfqpoint{-0.027008in}{-0.027008in}}%
\pgfpathcurveto{\pgfqpoint{-0.019845in}{-0.034170in}}{\pgfqpoint{-0.010129in}{-0.038194in}}{\pgfqpoint{0.000000in}{-0.038194in}}%
\pgfpathlineto{\pgfqpoint{0.000000in}{-0.038194in}}%
\pgfpathclose%
\pgfusepath{stroke,fill}%
}%
\begin{pgfscope}%
\pgfsys@transformshift{1.078000in}{1.954800in}%
\pgfsys@useobject{currentmarker}{}%
\end{pgfscope}%
\end{pgfscope}%
\begin{pgfscope}%
\pgfsetbuttcap%
\pgfsetroundjoin%
\definecolor{currentfill}{rgb}{0.549020,0.337255,0.294118}%
\pgfsetfillcolor{currentfill}%
\pgfsetlinewidth{1.003750pt}%
\definecolor{currentstroke}{rgb}{0.000000,0.000000,0.000000}%
\pgfsetstrokecolor{currentstroke}%
\pgfsetstrokeopacity{0.000000}%
\pgfsetdash{}{0pt}%
\pgfsys@defobject{currentmarker}{\pgfqpoint{-0.038194in}{-0.038194in}}{\pgfqpoint{0.038194in}{0.038194in}}{%
\pgfpathmoveto{\pgfqpoint{0.000000in}{-0.038194in}}%
\pgfpathcurveto{\pgfqpoint{0.010129in}{-0.038194in}}{\pgfqpoint{0.019845in}{-0.034170in}}{\pgfqpoint{0.027008in}{-0.027008in}}%
\pgfpathcurveto{\pgfqpoint{0.034170in}{-0.019845in}}{\pgfqpoint{0.038194in}{-0.010129in}}{\pgfqpoint{0.038194in}{0.000000in}}%
\pgfpathcurveto{\pgfqpoint{0.038194in}{0.010129in}}{\pgfqpoint{0.034170in}{0.019845in}}{\pgfqpoint{0.027008in}{0.027008in}}%
\pgfpathcurveto{\pgfqpoint{0.019845in}{0.034170in}}{\pgfqpoint{0.010129in}{0.038194in}}{\pgfqpoint{0.000000in}{0.038194in}}%
\pgfpathcurveto{\pgfqpoint{-0.010129in}{0.038194in}}{\pgfqpoint{-0.019845in}{0.034170in}}{\pgfqpoint{-0.027008in}{0.027008in}}%
\pgfpathcurveto{\pgfqpoint{-0.034170in}{0.019845in}}{\pgfqpoint{-0.038194in}{0.010129in}}{\pgfqpoint{-0.038194in}{0.000000in}}%
\pgfpathcurveto{\pgfqpoint{-0.038194in}{-0.010129in}}{\pgfqpoint{-0.034170in}{-0.019845in}}{\pgfqpoint{-0.027008in}{-0.027008in}}%
\pgfpathcurveto{\pgfqpoint{-0.019845in}{-0.034170in}}{\pgfqpoint{-0.010129in}{-0.038194in}}{\pgfqpoint{0.000000in}{-0.038194in}}%
\pgfpathlineto{\pgfqpoint{0.000000in}{-0.038194in}}%
\pgfpathclose%
\pgfusepath{stroke,fill}%
}%
\begin{pgfscope}%
\pgfsys@transformshift{1.078000in}{1.555200in}%
\pgfsys@useobject{currentmarker}{}%
\end{pgfscope}%
\end{pgfscope}%
\begin{pgfscope}%
\pgfsetbuttcap%
\pgfsetroundjoin%
\definecolor{currentfill}{rgb}{0.835294,0.368627,0.000000}%
\pgfsetfillcolor{currentfill}%
\pgfsetlinewidth{1.003750pt}%
\definecolor{currentstroke}{rgb}{0.000000,0.000000,0.000000}%
\pgfsetstrokecolor{currentstroke}%
\pgfsetstrokeopacity{0.000000}%
\pgfsetdash{}{0pt}%
\pgfsys@defobject{currentmarker}{\pgfqpoint{-0.038194in}{-0.038194in}}{\pgfqpoint{0.038194in}{0.038194in}}{%
\pgfpathmoveto{\pgfqpoint{0.000000in}{-0.038194in}}%
\pgfpathcurveto{\pgfqpoint{0.010129in}{-0.038194in}}{\pgfqpoint{0.019845in}{-0.034170in}}{\pgfqpoint{0.027008in}{-0.027008in}}%
\pgfpathcurveto{\pgfqpoint{0.034170in}{-0.019845in}}{\pgfqpoint{0.038194in}{-0.010129in}}{\pgfqpoint{0.038194in}{0.000000in}}%
\pgfpathcurveto{\pgfqpoint{0.038194in}{0.010129in}}{\pgfqpoint{0.034170in}{0.019845in}}{\pgfqpoint{0.027008in}{0.027008in}}%
\pgfpathcurveto{\pgfqpoint{0.019845in}{0.034170in}}{\pgfqpoint{0.010129in}{0.038194in}}{\pgfqpoint{0.000000in}{0.038194in}}%
\pgfpathcurveto{\pgfqpoint{-0.010129in}{0.038194in}}{\pgfqpoint{-0.019845in}{0.034170in}}{\pgfqpoint{-0.027008in}{0.027008in}}%
\pgfpathcurveto{\pgfqpoint{-0.034170in}{0.019845in}}{\pgfqpoint{-0.038194in}{0.010129in}}{\pgfqpoint{-0.038194in}{0.000000in}}%
\pgfpathcurveto{\pgfqpoint{-0.038194in}{-0.010129in}}{\pgfqpoint{-0.034170in}{-0.019845in}}{\pgfqpoint{-0.027008in}{-0.027008in}}%
\pgfpathcurveto{\pgfqpoint{-0.019845in}{-0.034170in}}{\pgfqpoint{-0.010129in}{-0.038194in}}{\pgfqpoint{0.000000in}{-0.038194in}}%
\pgfpathlineto{\pgfqpoint{0.000000in}{-0.038194in}}%
\pgfpathclose%
\pgfusepath{stroke,fill}%
}%
\begin{pgfscope}%
\pgfsys@transformshift{4.312000in}{3.153600in}%
\pgfsys@useobject{currentmarker}{}%
\end{pgfscope}%
\end{pgfscope}%
\begin{pgfscope}%
\pgfsetbuttcap%
\pgfsetroundjoin%
\definecolor{currentfill}{rgb}{0.662745,0.368627,0.000000}%
\pgfsetfillcolor{currentfill}%
\pgfsetlinewidth{1.003750pt}%
\definecolor{currentstroke}{rgb}{0.000000,0.000000,0.000000}%
\pgfsetstrokecolor{currentstroke}%
\pgfsetstrokeopacity{0.000000}%
\pgfsetdash{}{0pt}%
\pgfsys@defobject{currentmarker}{\pgfqpoint{-0.038194in}{-0.038194in}}{\pgfqpoint{0.038194in}{0.038194in}}{%
\pgfpathmoveto{\pgfqpoint{0.000000in}{-0.038194in}}%
\pgfpathcurveto{\pgfqpoint{0.010129in}{-0.038194in}}{\pgfqpoint{0.019845in}{-0.034170in}}{\pgfqpoint{0.027008in}{-0.027008in}}%
\pgfpathcurveto{\pgfqpoint{0.034170in}{-0.019845in}}{\pgfqpoint{0.038194in}{-0.010129in}}{\pgfqpoint{0.038194in}{0.000000in}}%
\pgfpathcurveto{\pgfqpoint{0.038194in}{0.010129in}}{\pgfqpoint{0.034170in}{0.019845in}}{\pgfqpoint{0.027008in}{0.027008in}}%
\pgfpathcurveto{\pgfqpoint{0.019845in}{0.034170in}}{\pgfqpoint{0.010129in}{0.038194in}}{\pgfqpoint{0.000000in}{0.038194in}}%
\pgfpathcurveto{\pgfqpoint{-0.010129in}{0.038194in}}{\pgfqpoint{-0.019845in}{0.034170in}}{\pgfqpoint{-0.027008in}{0.027008in}}%
\pgfpathcurveto{\pgfqpoint{-0.034170in}{0.019845in}}{\pgfqpoint{-0.038194in}{0.010129in}}{\pgfqpoint{-0.038194in}{0.000000in}}%
\pgfpathcurveto{\pgfqpoint{-0.038194in}{-0.010129in}}{\pgfqpoint{-0.034170in}{-0.019845in}}{\pgfqpoint{-0.027008in}{-0.027008in}}%
\pgfpathcurveto{\pgfqpoint{-0.019845in}{-0.034170in}}{\pgfqpoint{-0.010129in}{-0.038194in}}{\pgfqpoint{0.000000in}{-0.038194in}}%
\pgfpathlineto{\pgfqpoint{0.000000in}{-0.038194in}}%
\pgfpathclose%
\pgfusepath{stroke,fill}%
}%
\begin{pgfscope}%
\pgfsys@transformshift{4.312000in}{2.754000in}%
\pgfsys@useobject{currentmarker}{}%
\end{pgfscope}%
\end{pgfscope}%
\begin{pgfscope}%
\pgfsetbuttcap%
\pgfsetroundjoin%
\definecolor{currentfill}{rgb}{0.337255,0.705882,0.913725}%
\pgfsetfillcolor{currentfill}%
\pgfsetlinewidth{1.003750pt}%
\definecolor{currentstroke}{rgb}{0.000000,0.000000,0.000000}%
\pgfsetstrokecolor{currentstroke}%
\pgfsetstrokeopacity{0.000000}%
\pgfsetdash{}{0pt}%
\pgfsys@defobject{currentmarker}{\pgfqpoint{-0.038194in}{-0.038194in}}{\pgfqpoint{0.038194in}{0.038194in}}{%
\pgfpathmoveto{\pgfqpoint{0.000000in}{-0.038194in}}%
\pgfpathcurveto{\pgfqpoint{0.010129in}{-0.038194in}}{\pgfqpoint{0.019845in}{-0.034170in}}{\pgfqpoint{0.027008in}{-0.027008in}}%
\pgfpathcurveto{\pgfqpoint{0.034170in}{-0.019845in}}{\pgfqpoint{0.038194in}{-0.010129in}}{\pgfqpoint{0.038194in}{0.000000in}}%
\pgfpathcurveto{\pgfqpoint{0.038194in}{0.010129in}}{\pgfqpoint{0.034170in}{0.019845in}}{\pgfqpoint{0.027008in}{0.027008in}}%
\pgfpathcurveto{\pgfqpoint{0.019845in}{0.034170in}}{\pgfqpoint{0.010129in}{0.038194in}}{\pgfqpoint{0.000000in}{0.038194in}}%
\pgfpathcurveto{\pgfqpoint{-0.010129in}{0.038194in}}{\pgfqpoint{-0.019845in}{0.034170in}}{\pgfqpoint{-0.027008in}{0.027008in}}%
\pgfpathcurveto{\pgfqpoint{-0.034170in}{0.019845in}}{\pgfqpoint{-0.038194in}{0.010129in}}{\pgfqpoint{-0.038194in}{0.000000in}}%
\pgfpathcurveto{\pgfqpoint{-0.038194in}{-0.010129in}}{\pgfqpoint{-0.034170in}{-0.019845in}}{\pgfqpoint{-0.027008in}{-0.027008in}}%
\pgfpathcurveto{\pgfqpoint{-0.019845in}{-0.034170in}}{\pgfqpoint{-0.010129in}{-0.038194in}}{\pgfqpoint{0.000000in}{-0.038194in}}%
\pgfpathlineto{\pgfqpoint{0.000000in}{-0.038194in}}%
\pgfpathclose%
\pgfusepath{stroke,fill}%
}%
\begin{pgfscope}%
\pgfsys@transformshift{4.312000in}{2.354400in}%
\pgfsys@useobject{currentmarker}{}%
\end{pgfscope}%
\end{pgfscope}%
\begin{pgfscope}%
\pgfsetbuttcap%
\pgfsetroundjoin%
\definecolor{currentfill}{rgb}{0.000000,0.619608,0.450980}%
\pgfsetfillcolor{currentfill}%
\pgfsetlinewidth{1.003750pt}%
\definecolor{currentstroke}{rgb}{0.000000,0.000000,0.000000}%
\pgfsetstrokecolor{currentstroke}%
\pgfsetstrokeopacity{0.000000}%
\pgfsetdash{}{0pt}%
\pgfsys@defobject{currentmarker}{\pgfqpoint{-0.038194in}{-0.038194in}}{\pgfqpoint{0.038194in}{0.038194in}}{%
\pgfpathmoveto{\pgfqpoint{0.000000in}{-0.038194in}}%
\pgfpathcurveto{\pgfqpoint{0.010129in}{-0.038194in}}{\pgfqpoint{0.019845in}{-0.034170in}}{\pgfqpoint{0.027008in}{-0.027008in}}%
\pgfpathcurveto{\pgfqpoint{0.034170in}{-0.019845in}}{\pgfqpoint{0.038194in}{-0.010129in}}{\pgfqpoint{0.038194in}{0.000000in}}%
\pgfpathcurveto{\pgfqpoint{0.038194in}{0.010129in}}{\pgfqpoint{0.034170in}{0.019845in}}{\pgfqpoint{0.027008in}{0.027008in}}%
\pgfpathcurveto{\pgfqpoint{0.019845in}{0.034170in}}{\pgfqpoint{0.010129in}{0.038194in}}{\pgfqpoint{0.000000in}{0.038194in}}%
\pgfpathcurveto{\pgfqpoint{-0.010129in}{0.038194in}}{\pgfqpoint{-0.019845in}{0.034170in}}{\pgfqpoint{-0.027008in}{0.027008in}}%
\pgfpathcurveto{\pgfqpoint{-0.034170in}{0.019845in}}{\pgfqpoint{-0.038194in}{0.010129in}}{\pgfqpoint{-0.038194in}{0.000000in}}%
\pgfpathcurveto{\pgfqpoint{-0.038194in}{-0.010129in}}{\pgfqpoint{-0.034170in}{-0.019845in}}{\pgfqpoint{-0.027008in}{-0.027008in}}%
\pgfpathcurveto{\pgfqpoint{-0.019845in}{-0.034170in}}{\pgfqpoint{-0.010129in}{-0.038194in}}{\pgfqpoint{0.000000in}{-0.038194in}}%
\pgfpathlineto{\pgfqpoint{0.000000in}{-0.038194in}}%
\pgfpathclose%
\pgfusepath{stroke,fill}%
}%
\begin{pgfscope}%
\pgfsys@transformshift{4.312000in}{1.954800in}%
\pgfsys@useobject{currentmarker}{}%
\end{pgfscope}%
\end{pgfscope}%
\begin{pgfscope}%
\pgfsetbuttcap%
\pgfsetroundjoin%
\definecolor{currentfill}{rgb}{0.709804,0.121569,0.270588}%
\pgfsetfillcolor{currentfill}%
\pgfsetlinewidth{1.003750pt}%
\definecolor{currentstroke}{rgb}{0.000000,0.000000,0.000000}%
\pgfsetstrokecolor{currentstroke}%
\pgfsetstrokeopacity{0.000000}%
\pgfsetdash{}{0pt}%
\pgfsys@defobject{currentmarker}{\pgfqpoint{-0.038194in}{-0.038194in}}{\pgfqpoint{0.038194in}{0.038194in}}{%
\pgfpathmoveto{\pgfqpoint{0.000000in}{-0.038194in}}%
\pgfpathcurveto{\pgfqpoint{0.010129in}{-0.038194in}}{\pgfqpoint{0.019845in}{-0.034170in}}{\pgfqpoint{0.027008in}{-0.027008in}}%
\pgfpathcurveto{\pgfqpoint{0.034170in}{-0.019845in}}{\pgfqpoint{0.038194in}{-0.010129in}}{\pgfqpoint{0.038194in}{0.000000in}}%
\pgfpathcurveto{\pgfqpoint{0.038194in}{0.010129in}}{\pgfqpoint{0.034170in}{0.019845in}}{\pgfqpoint{0.027008in}{0.027008in}}%
\pgfpathcurveto{\pgfqpoint{0.019845in}{0.034170in}}{\pgfqpoint{0.010129in}{0.038194in}}{\pgfqpoint{0.000000in}{0.038194in}}%
\pgfpathcurveto{\pgfqpoint{-0.010129in}{0.038194in}}{\pgfqpoint{-0.019845in}{0.034170in}}{\pgfqpoint{-0.027008in}{0.027008in}}%
\pgfpathcurveto{\pgfqpoint{-0.034170in}{0.019845in}}{\pgfqpoint{-0.038194in}{0.010129in}}{\pgfqpoint{-0.038194in}{0.000000in}}%
\pgfpathcurveto{\pgfqpoint{-0.038194in}{-0.010129in}}{\pgfqpoint{-0.034170in}{-0.019845in}}{\pgfqpoint{-0.027008in}{-0.027008in}}%
\pgfpathcurveto{\pgfqpoint{-0.019845in}{-0.034170in}}{\pgfqpoint{-0.010129in}{-0.038194in}}{\pgfqpoint{0.000000in}{-0.038194in}}%
\pgfpathlineto{\pgfqpoint{0.000000in}{-0.038194in}}%
\pgfpathclose%
\pgfusepath{stroke,fill}%
}%
\begin{pgfscope}%
\pgfsys@transformshift{4.312000in}{1.555200in}%
\pgfsys@useobject{currentmarker}{}%
\end{pgfscope}%
\end{pgfscope}%
\begin{pgfscope}%
\pgfsetbuttcap%
\pgfsetroundjoin%
\definecolor{currentfill}{rgb}{0.901961,0.403922,0.619608}%
\pgfsetfillcolor{currentfill}%
\pgfsetlinewidth{1.003750pt}%
\definecolor{currentstroke}{rgb}{0.000000,0.000000,0.000000}%
\pgfsetstrokecolor{currentstroke}%
\pgfsetstrokeopacity{0.000000}%
\pgfsetdash{}{0pt}%
\pgfsys@defobject{currentmarker}{\pgfqpoint{-0.038194in}{-0.038194in}}{\pgfqpoint{0.038194in}{0.038194in}}{%
\pgfpathmoveto{\pgfqpoint{0.000000in}{-0.038194in}}%
\pgfpathcurveto{\pgfqpoint{0.010129in}{-0.038194in}}{\pgfqpoint{0.019845in}{-0.034170in}}{\pgfqpoint{0.027008in}{-0.027008in}}%
\pgfpathcurveto{\pgfqpoint{0.034170in}{-0.019845in}}{\pgfqpoint{0.038194in}{-0.010129in}}{\pgfqpoint{0.038194in}{0.000000in}}%
\pgfpathcurveto{\pgfqpoint{0.038194in}{0.010129in}}{\pgfqpoint{0.034170in}{0.019845in}}{\pgfqpoint{0.027008in}{0.027008in}}%
\pgfpathcurveto{\pgfqpoint{0.019845in}{0.034170in}}{\pgfqpoint{0.010129in}{0.038194in}}{\pgfqpoint{0.000000in}{0.038194in}}%
\pgfpathcurveto{\pgfqpoint{-0.010129in}{0.038194in}}{\pgfqpoint{-0.019845in}{0.034170in}}{\pgfqpoint{-0.027008in}{0.027008in}}%
\pgfpathcurveto{\pgfqpoint{-0.034170in}{0.019845in}}{\pgfqpoint{-0.038194in}{0.010129in}}{\pgfqpoint{-0.038194in}{0.000000in}}%
\pgfpathcurveto{\pgfqpoint{-0.038194in}{-0.010129in}}{\pgfqpoint{-0.034170in}{-0.019845in}}{\pgfqpoint{-0.027008in}{-0.027008in}}%
\pgfpathcurveto{\pgfqpoint{-0.019845in}{-0.034170in}}{\pgfqpoint{-0.010129in}{-0.038194in}}{\pgfqpoint{0.000000in}{-0.038194in}}%
\pgfpathlineto{\pgfqpoint{0.000000in}{-0.038194in}}%
\pgfpathclose%
\pgfusepath{stroke,fill}%
}%
\begin{pgfscope}%
\pgfsys@transformshift{7.546000in}{3.153600in}%
\pgfsys@useobject{currentmarker}{}%
\end{pgfscope}%
\end{pgfscope}%
\begin{pgfscope}%
\pgfsetbuttcap%
\pgfsetroundjoin%
\definecolor{currentfill}{rgb}{0.549020,0.160784,0.505882}%
\pgfsetfillcolor{currentfill}%
\pgfsetlinewidth{1.003750pt}%
\definecolor{currentstroke}{rgb}{0.000000,0.000000,0.000000}%
\pgfsetstrokecolor{currentstroke}%
\pgfsetstrokeopacity{0.000000}%
\pgfsetdash{}{0pt}%
\pgfsys@defobject{currentmarker}{\pgfqpoint{-0.038194in}{-0.038194in}}{\pgfqpoint{0.038194in}{0.038194in}}{%
\pgfpathmoveto{\pgfqpoint{0.000000in}{-0.038194in}}%
\pgfpathcurveto{\pgfqpoint{0.010129in}{-0.038194in}}{\pgfqpoint{0.019845in}{-0.034170in}}{\pgfqpoint{0.027008in}{-0.027008in}}%
\pgfpathcurveto{\pgfqpoint{0.034170in}{-0.019845in}}{\pgfqpoint{0.038194in}{-0.010129in}}{\pgfqpoint{0.038194in}{0.000000in}}%
\pgfpathcurveto{\pgfqpoint{0.038194in}{0.010129in}}{\pgfqpoint{0.034170in}{0.019845in}}{\pgfqpoint{0.027008in}{0.027008in}}%
\pgfpathcurveto{\pgfqpoint{0.019845in}{0.034170in}}{\pgfqpoint{0.010129in}{0.038194in}}{\pgfqpoint{0.000000in}{0.038194in}}%
\pgfpathcurveto{\pgfqpoint{-0.010129in}{0.038194in}}{\pgfqpoint{-0.019845in}{0.034170in}}{\pgfqpoint{-0.027008in}{0.027008in}}%
\pgfpathcurveto{\pgfqpoint{-0.034170in}{0.019845in}}{\pgfqpoint{-0.038194in}{0.010129in}}{\pgfqpoint{-0.038194in}{0.000000in}}%
\pgfpathcurveto{\pgfqpoint{-0.038194in}{-0.010129in}}{\pgfqpoint{-0.034170in}{-0.019845in}}{\pgfqpoint{-0.027008in}{-0.027008in}}%
\pgfpathcurveto{\pgfqpoint{-0.019845in}{-0.034170in}}{\pgfqpoint{-0.010129in}{-0.038194in}}{\pgfqpoint{0.000000in}{-0.038194in}}%
\pgfpathlineto{\pgfqpoint{0.000000in}{-0.038194in}}%
\pgfpathclose%
\pgfusepath{stroke,fill}%
}%
\begin{pgfscope}%
\pgfsys@transformshift{7.546000in}{2.754000in}%
\pgfsys@useobject{currentmarker}{}%
\end{pgfscope}%
\end{pgfscope}%
\begin{pgfscope}%
\pgfsetbuttcap%
\pgfsetroundjoin%
\definecolor{currentfill}{rgb}{0.941176,0.584314,0.337255}%
\pgfsetfillcolor{currentfill}%
\pgfsetlinewidth{1.003750pt}%
\definecolor{currentstroke}{rgb}{0.000000,0.000000,0.000000}%
\pgfsetstrokecolor{currentstroke}%
\pgfsetstrokeopacity{0.000000}%
\pgfsetdash{}{0pt}%
\pgfsys@defobject{currentmarker}{\pgfqpoint{-0.038194in}{-0.038194in}}{\pgfqpoint{0.038194in}{0.038194in}}{%
\pgfpathmoveto{\pgfqpoint{0.000000in}{-0.038194in}}%
\pgfpathcurveto{\pgfqpoint{0.010129in}{-0.038194in}}{\pgfqpoint{0.019845in}{-0.034170in}}{\pgfqpoint{0.027008in}{-0.027008in}}%
\pgfpathcurveto{\pgfqpoint{0.034170in}{-0.019845in}}{\pgfqpoint{0.038194in}{-0.010129in}}{\pgfqpoint{0.038194in}{0.000000in}}%
\pgfpathcurveto{\pgfqpoint{0.038194in}{0.010129in}}{\pgfqpoint{0.034170in}{0.019845in}}{\pgfqpoint{0.027008in}{0.027008in}}%
\pgfpathcurveto{\pgfqpoint{0.019845in}{0.034170in}}{\pgfqpoint{0.010129in}{0.038194in}}{\pgfqpoint{0.000000in}{0.038194in}}%
\pgfpathcurveto{\pgfqpoint{-0.010129in}{0.038194in}}{\pgfqpoint{-0.019845in}{0.034170in}}{\pgfqpoint{-0.027008in}{0.027008in}}%
\pgfpathcurveto{\pgfqpoint{-0.034170in}{0.019845in}}{\pgfqpoint{-0.038194in}{0.010129in}}{\pgfqpoint{-0.038194in}{0.000000in}}%
\pgfpathcurveto{\pgfqpoint{-0.038194in}{-0.010129in}}{\pgfqpoint{-0.034170in}{-0.019845in}}{\pgfqpoint{-0.027008in}{-0.027008in}}%
\pgfpathcurveto{\pgfqpoint{-0.019845in}{-0.034170in}}{\pgfqpoint{-0.010129in}{-0.038194in}}{\pgfqpoint{0.000000in}{-0.038194in}}%
\pgfpathlineto{\pgfqpoint{0.000000in}{-0.038194in}}%
\pgfpathclose%
\pgfusepath{stroke,fill}%
}%
\begin{pgfscope}%
\pgfsys@transformshift{7.546000in}{2.354400in}%
\pgfsys@useobject{currentmarker}{}%
\end{pgfscope}%
\end{pgfscope}%
\begin{pgfscope}%
\pgfsetbuttcap%
\pgfsetroundjoin%
\definecolor{currentfill}{rgb}{0.000000,0.560784,0.611765}%
\pgfsetfillcolor{currentfill}%
\pgfsetlinewidth{1.003750pt}%
\definecolor{currentstroke}{rgb}{0.000000,0.000000,0.000000}%
\pgfsetstrokecolor{currentstroke}%
\pgfsetstrokeopacity{0.000000}%
\pgfsetdash{}{0pt}%
\pgfsys@defobject{currentmarker}{\pgfqpoint{-0.038194in}{-0.038194in}}{\pgfqpoint{0.038194in}{0.038194in}}{%
\pgfpathmoveto{\pgfqpoint{0.000000in}{-0.038194in}}%
\pgfpathcurveto{\pgfqpoint{0.010129in}{-0.038194in}}{\pgfqpoint{0.019845in}{-0.034170in}}{\pgfqpoint{0.027008in}{-0.027008in}}%
\pgfpathcurveto{\pgfqpoint{0.034170in}{-0.019845in}}{\pgfqpoint{0.038194in}{-0.010129in}}{\pgfqpoint{0.038194in}{0.000000in}}%
\pgfpathcurveto{\pgfqpoint{0.038194in}{0.010129in}}{\pgfqpoint{0.034170in}{0.019845in}}{\pgfqpoint{0.027008in}{0.027008in}}%
\pgfpathcurveto{\pgfqpoint{0.019845in}{0.034170in}}{\pgfqpoint{0.010129in}{0.038194in}}{\pgfqpoint{0.000000in}{0.038194in}}%
\pgfpathcurveto{\pgfqpoint{-0.010129in}{0.038194in}}{\pgfqpoint{-0.019845in}{0.034170in}}{\pgfqpoint{-0.027008in}{0.027008in}}%
\pgfpathcurveto{\pgfqpoint{-0.034170in}{0.019845in}}{\pgfqpoint{-0.038194in}{0.010129in}}{\pgfqpoint{-0.038194in}{0.000000in}}%
\pgfpathcurveto{\pgfqpoint{-0.038194in}{-0.010129in}}{\pgfqpoint{-0.034170in}{-0.019845in}}{\pgfqpoint{-0.027008in}{-0.027008in}}%
\pgfpathcurveto{\pgfqpoint{-0.019845in}{-0.034170in}}{\pgfqpoint{-0.010129in}{-0.038194in}}{\pgfqpoint{0.000000in}{-0.038194in}}%
\pgfpathlineto{\pgfqpoint{0.000000in}{-0.038194in}}%
\pgfpathclose%
\pgfusepath{stroke,fill}%
}%
\begin{pgfscope}%
\pgfsys@transformshift{7.546000in}{1.954800in}%
\pgfsys@useobject{currentmarker}{}%
\end{pgfscope}%
\end{pgfscope}%
\begin{pgfscope}%
\pgfsetbuttcap%
\pgfsetroundjoin%
\definecolor{currentfill}{rgb}{0.200000,0.133333,0.533333}%
\pgfsetfillcolor{currentfill}%
\pgfsetlinewidth{1.003750pt}%
\definecolor{currentstroke}{rgb}{0.000000,0.000000,0.000000}%
\pgfsetstrokecolor{currentstroke}%
\pgfsetstrokeopacity{0.000000}%
\pgfsetdash{}{0pt}%
\pgfsys@defobject{currentmarker}{\pgfqpoint{-0.038194in}{-0.038194in}}{\pgfqpoint{0.038194in}{0.038194in}}{%
\pgfpathmoveto{\pgfqpoint{0.000000in}{-0.038194in}}%
\pgfpathcurveto{\pgfqpoint{0.010129in}{-0.038194in}}{\pgfqpoint{0.019845in}{-0.034170in}}{\pgfqpoint{0.027008in}{-0.027008in}}%
\pgfpathcurveto{\pgfqpoint{0.034170in}{-0.019845in}}{\pgfqpoint{0.038194in}{-0.010129in}}{\pgfqpoint{0.038194in}{0.000000in}}%
\pgfpathcurveto{\pgfqpoint{0.038194in}{0.010129in}}{\pgfqpoint{0.034170in}{0.019845in}}{\pgfqpoint{0.027008in}{0.027008in}}%
\pgfpathcurveto{\pgfqpoint{0.019845in}{0.034170in}}{\pgfqpoint{0.010129in}{0.038194in}}{\pgfqpoint{0.000000in}{0.038194in}}%
\pgfpathcurveto{\pgfqpoint{-0.010129in}{0.038194in}}{\pgfqpoint{-0.019845in}{0.034170in}}{\pgfqpoint{-0.027008in}{0.027008in}}%
\pgfpathcurveto{\pgfqpoint{-0.034170in}{0.019845in}}{\pgfqpoint{-0.038194in}{0.010129in}}{\pgfqpoint{-0.038194in}{0.000000in}}%
\pgfpathcurveto{\pgfqpoint{-0.038194in}{-0.010129in}}{\pgfqpoint{-0.034170in}{-0.019845in}}{\pgfqpoint{-0.027008in}{-0.027008in}}%
\pgfpathcurveto{\pgfqpoint{-0.019845in}{-0.034170in}}{\pgfqpoint{-0.010129in}{-0.038194in}}{\pgfqpoint{0.000000in}{-0.038194in}}%
\pgfpathlineto{\pgfqpoint{0.000000in}{-0.038194in}}%
\pgfpathclose%
\pgfusepath{stroke,fill}%
}%
\begin{pgfscope}%
\pgfsys@transformshift{7.546000in}{1.555200in}%
\pgfsys@useobject{currentmarker}{}%
\end{pgfscope}%
\end{pgfscope}%
\begin{pgfscope}%
\pgfsetbuttcap%
\pgfsetroundjoin%
\definecolor{currentfill}{rgb}{0.901961,0.196078,0.156863}%
\pgfsetfillcolor{currentfill}%
\pgfsetlinewidth{1.003750pt}%
\definecolor{currentstroke}{rgb}{0.000000,0.000000,0.000000}%
\pgfsetstrokecolor{currentstroke}%
\pgfsetstrokeopacity{0.000000}%
\pgfsetdash{}{0pt}%
\pgfsys@defobject{currentmarker}{\pgfqpoint{-0.038194in}{-0.038194in}}{\pgfqpoint{0.038194in}{0.038194in}}{%
\pgfpathmoveto{\pgfqpoint{0.000000in}{-0.038194in}}%
\pgfpathcurveto{\pgfqpoint{0.010129in}{-0.038194in}}{\pgfqpoint{0.019845in}{-0.034170in}}{\pgfqpoint{0.027008in}{-0.027008in}}%
\pgfpathcurveto{\pgfqpoint{0.034170in}{-0.019845in}}{\pgfqpoint{0.038194in}{-0.010129in}}{\pgfqpoint{0.038194in}{0.000000in}}%
\pgfpathcurveto{\pgfqpoint{0.038194in}{0.010129in}}{\pgfqpoint{0.034170in}{0.019845in}}{\pgfqpoint{0.027008in}{0.027008in}}%
\pgfpathcurveto{\pgfqpoint{0.019845in}{0.034170in}}{\pgfqpoint{0.010129in}{0.038194in}}{\pgfqpoint{0.000000in}{0.038194in}}%
\pgfpathcurveto{\pgfqpoint{-0.010129in}{0.038194in}}{\pgfqpoint{-0.019845in}{0.034170in}}{\pgfqpoint{-0.027008in}{0.027008in}}%
\pgfpathcurveto{\pgfqpoint{-0.034170in}{0.019845in}}{\pgfqpoint{-0.038194in}{0.010129in}}{\pgfqpoint{-0.038194in}{0.000000in}}%
\pgfpathcurveto{\pgfqpoint{-0.038194in}{-0.010129in}}{\pgfqpoint{-0.034170in}{-0.019845in}}{\pgfqpoint{-0.027008in}{-0.027008in}}%
\pgfpathcurveto{\pgfqpoint{-0.019845in}{-0.034170in}}{\pgfqpoint{-0.010129in}{-0.038194in}}{\pgfqpoint{0.000000in}{-0.038194in}}%
\pgfpathlineto{\pgfqpoint{0.000000in}{-0.038194in}}%
\pgfpathclose%
\pgfusepath{stroke,fill}%
}%
\begin{pgfscope}%
\pgfsys@transformshift{10.780000in}{3.153600in}%
\pgfsys@useobject{currentmarker}{}%
\end{pgfscope}%
\end{pgfscope}%
\begin{pgfscope}%
\pgfsetbuttcap%
\pgfsetroundjoin%
\definecolor{currentfill}{rgb}{0.768627,0.603922,0.000000}%
\pgfsetfillcolor{currentfill}%
\pgfsetlinewidth{1.003750pt}%
\definecolor{currentstroke}{rgb}{0.000000,0.000000,0.000000}%
\pgfsetstrokecolor{currentstroke}%
\pgfsetstrokeopacity{0.000000}%
\pgfsetdash{}{0pt}%
\pgfsys@defobject{currentmarker}{\pgfqpoint{-0.038194in}{-0.038194in}}{\pgfqpoint{0.038194in}{0.038194in}}{%
\pgfpathmoveto{\pgfqpoint{0.000000in}{-0.038194in}}%
\pgfpathcurveto{\pgfqpoint{0.010129in}{-0.038194in}}{\pgfqpoint{0.019845in}{-0.034170in}}{\pgfqpoint{0.027008in}{-0.027008in}}%
\pgfpathcurveto{\pgfqpoint{0.034170in}{-0.019845in}}{\pgfqpoint{0.038194in}{-0.010129in}}{\pgfqpoint{0.038194in}{0.000000in}}%
\pgfpathcurveto{\pgfqpoint{0.038194in}{0.010129in}}{\pgfqpoint{0.034170in}{0.019845in}}{\pgfqpoint{0.027008in}{0.027008in}}%
\pgfpathcurveto{\pgfqpoint{0.019845in}{0.034170in}}{\pgfqpoint{0.010129in}{0.038194in}}{\pgfqpoint{0.000000in}{0.038194in}}%
\pgfpathcurveto{\pgfqpoint{-0.010129in}{0.038194in}}{\pgfqpoint{-0.019845in}{0.034170in}}{\pgfqpoint{-0.027008in}{0.027008in}}%
\pgfpathcurveto{\pgfqpoint{-0.034170in}{0.019845in}}{\pgfqpoint{-0.038194in}{0.010129in}}{\pgfqpoint{-0.038194in}{0.000000in}}%
\pgfpathcurveto{\pgfqpoint{-0.038194in}{-0.010129in}}{\pgfqpoint{-0.034170in}{-0.019845in}}{\pgfqpoint{-0.027008in}{-0.027008in}}%
\pgfpathcurveto{\pgfqpoint{-0.019845in}{-0.034170in}}{\pgfqpoint{-0.010129in}{-0.038194in}}{\pgfqpoint{0.000000in}{-0.038194in}}%
\pgfpathlineto{\pgfqpoint{0.000000in}{-0.038194in}}%
\pgfpathclose%
\pgfusepath{stroke,fill}%
}%
\begin{pgfscope}%
\pgfsys@transformshift{10.780000in}{2.754000in}%
\pgfsys@useobject{currentmarker}{}%
\end{pgfscope}%
\end{pgfscope}%
\begin{pgfscope}%
\pgfsetbuttcap%
\pgfsetroundjoin%
\definecolor{currentfill}{rgb}{0.498039,0.560784,0.678431}%
\pgfsetfillcolor{currentfill}%
\pgfsetlinewidth{1.003750pt}%
\definecolor{currentstroke}{rgb}{0.000000,0.000000,0.000000}%
\pgfsetstrokecolor{currentstroke}%
\pgfsetstrokeopacity{0.000000}%
\pgfsetdash{}{0pt}%
\pgfsys@defobject{currentmarker}{\pgfqpoint{-0.038194in}{-0.038194in}}{\pgfqpoint{0.038194in}{0.038194in}}{%
\pgfpathmoveto{\pgfqpoint{0.000000in}{-0.038194in}}%
\pgfpathcurveto{\pgfqpoint{0.010129in}{-0.038194in}}{\pgfqpoint{0.019845in}{-0.034170in}}{\pgfqpoint{0.027008in}{-0.027008in}}%
\pgfpathcurveto{\pgfqpoint{0.034170in}{-0.019845in}}{\pgfqpoint{0.038194in}{-0.010129in}}{\pgfqpoint{0.038194in}{0.000000in}}%
\pgfpathcurveto{\pgfqpoint{0.038194in}{0.010129in}}{\pgfqpoint{0.034170in}{0.019845in}}{\pgfqpoint{0.027008in}{0.027008in}}%
\pgfpathcurveto{\pgfqpoint{0.019845in}{0.034170in}}{\pgfqpoint{0.010129in}{0.038194in}}{\pgfqpoint{0.000000in}{0.038194in}}%
\pgfpathcurveto{\pgfqpoint{-0.010129in}{0.038194in}}{\pgfqpoint{-0.019845in}{0.034170in}}{\pgfqpoint{-0.027008in}{0.027008in}}%
\pgfpathcurveto{\pgfqpoint{-0.034170in}{0.019845in}}{\pgfqpoint{-0.038194in}{0.010129in}}{\pgfqpoint{-0.038194in}{0.000000in}}%
\pgfpathcurveto{\pgfqpoint{-0.038194in}{-0.010129in}}{\pgfqpoint{-0.034170in}{-0.019845in}}{\pgfqpoint{-0.027008in}{-0.027008in}}%
\pgfpathcurveto{\pgfqpoint{-0.019845in}{-0.034170in}}{\pgfqpoint{-0.010129in}{-0.038194in}}{\pgfqpoint{0.000000in}{-0.038194in}}%
\pgfpathlineto{\pgfqpoint{0.000000in}{-0.038194in}}%
\pgfpathclose%
\pgfusepath{stroke,fill}%
}%
\begin{pgfscope}%
\pgfsys@transformshift{10.780000in}{2.354400in}%
\pgfsys@useobject{currentmarker}{}%
\end{pgfscope}%
\end{pgfscope}%
\begin{pgfscope}%
\pgfsetbuttcap%
\pgfsetroundjoin%
\definecolor{currentfill}{rgb}{0.176471,0.713725,0.643137}%
\pgfsetfillcolor{currentfill}%
\pgfsetlinewidth{1.003750pt}%
\definecolor{currentstroke}{rgb}{0.000000,0.000000,0.000000}%
\pgfsetstrokecolor{currentstroke}%
\pgfsetstrokeopacity{0.000000}%
\pgfsetdash{}{0pt}%
\pgfsys@defobject{currentmarker}{\pgfqpoint{-0.038194in}{-0.038194in}}{\pgfqpoint{0.038194in}{0.038194in}}{%
\pgfpathmoveto{\pgfqpoint{0.000000in}{-0.038194in}}%
\pgfpathcurveto{\pgfqpoint{0.010129in}{-0.038194in}}{\pgfqpoint{0.019845in}{-0.034170in}}{\pgfqpoint{0.027008in}{-0.027008in}}%
\pgfpathcurveto{\pgfqpoint{0.034170in}{-0.019845in}}{\pgfqpoint{0.038194in}{-0.010129in}}{\pgfqpoint{0.038194in}{0.000000in}}%
\pgfpathcurveto{\pgfqpoint{0.038194in}{0.010129in}}{\pgfqpoint{0.034170in}{0.019845in}}{\pgfqpoint{0.027008in}{0.027008in}}%
\pgfpathcurveto{\pgfqpoint{0.019845in}{0.034170in}}{\pgfqpoint{0.010129in}{0.038194in}}{\pgfqpoint{0.000000in}{0.038194in}}%
\pgfpathcurveto{\pgfqpoint{-0.010129in}{0.038194in}}{\pgfqpoint{-0.019845in}{0.034170in}}{\pgfqpoint{-0.027008in}{0.027008in}}%
\pgfpathcurveto{\pgfqpoint{-0.034170in}{0.019845in}}{\pgfqpoint{-0.038194in}{0.010129in}}{\pgfqpoint{-0.038194in}{0.000000in}}%
\pgfpathcurveto{\pgfqpoint{-0.038194in}{-0.010129in}}{\pgfqpoint{-0.034170in}{-0.019845in}}{\pgfqpoint{-0.027008in}{-0.027008in}}%
\pgfpathcurveto{\pgfqpoint{-0.019845in}{-0.034170in}}{\pgfqpoint{-0.010129in}{-0.038194in}}{\pgfqpoint{0.000000in}{-0.038194in}}%
\pgfpathlineto{\pgfqpoint{0.000000in}{-0.038194in}}%
\pgfpathclose%
\pgfusepath{stroke,fill}%
}%
\begin{pgfscope}%
\pgfsys@transformshift{10.780000in}{1.954800in}%
\pgfsys@useobject{currentmarker}{}%
\end{pgfscope}%
\end{pgfscope}%
\begin{pgfscope}%
\pgfsetbuttcap%
\pgfsetroundjoin%
\definecolor{currentfill}{rgb}{0.321569,0.321569,0.321569}%
\pgfsetfillcolor{currentfill}%
\pgfsetlinewidth{1.003750pt}%
\definecolor{currentstroke}{rgb}{0.000000,0.000000,0.000000}%
\pgfsetstrokecolor{currentstroke}%
\pgfsetstrokeopacity{0.000000}%
\pgfsetdash{}{0pt}%
\pgfsys@defobject{currentmarker}{\pgfqpoint{-0.038194in}{-0.038194in}}{\pgfqpoint{0.038194in}{0.038194in}}{%
\pgfpathmoveto{\pgfqpoint{0.000000in}{-0.038194in}}%
\pgfpathcurveto{\pgfqpoint{0.010129in}{-0.038194in}}{\pgfqpoint{0.019845in}{-0.034170in}}{\pgfqpoint{0.027008in}{-0.027008in}}%
\pgfpathcurveto{\pgfqpoint{0.034170in}{-0.019845in}}{\pgfqpoint{0.038194in}{-0.010129in}}{\pgfqpoint{0.038194in}{0.000000in}}%
\pgfpathcurveto{\pgfqpoint{0.038194in}{0.010129in}}{\pgfqpoint{0.034170in}{0.019845in}}{\pgfqpoint{0.027008in}{0.027008in}}%
\pgfpathcurveto{\pgfqpoint{0.019845in}{0.034170in}}{\pgfqpoint{0.010129in}{0.038194in}}{\pgfqpoint{0.000000in}{0.038194in}}%
\pgfpathcurveto{\pgfqpoint{-0.010129in}{0.038194in}}{\pgfqpoint{-0.019845in}{0.034170in}}{\pgfqpoint{-0.027008in}{0.027008in}}%
\pgfpathcurveto{\pgfqpoint{-0.034170in}{0.019845in}}{\pgfqpoint{-0.038194in}{0.010129in}}{\pgfqpoint{-0.038194in}{0.000000in}}%
\pgfpathcurveto{\pgfqpoint{-0.038194in}{-0.010129in}}{\pgfqpoint{-0.034170in}{-0.019845in}}{\pgfqpoint{-0.027008in}{-0.027008in}}%
\pgfpathcurveto{\pgfqpoint{-0.019845in}{-0.034170in}}{\pgfqpoint{-0.010129in}{-0.038194in}}{\pgfqpoint{0.000000in}{-0.038194in}}%
\pgfpathlineto{\pgfqpoint{0.000000in}{-0.038194in}}%
\pgfpathclose%
\pgfusepath{stroke,fill}%
}%
\begin{pgfscope}%
\pgfsys@transformshift{10.780000in}{1.555200in}%
\pgfsys@useobject{currentmarker}{}%
\end{pgfscope}%
\end{pgfscope}%
\begin{pgfscope}%
\pgfsetbuttcap%
\pgfsetroundjoin%
\definecolor{currentfill}{rgb}{0.090196,0.168627,0.227451}%
\pgfsetfillcolor{currentfill}%
\pgfsetlinewidth{0.000000pt}%
\definecolor{currentstroke}{rgb}{0.090196,0.168627,0.227451}%
\pgfsetstrokecolor{currentstroke}%
\pgfsetdash{}{0pt}%
\pgfpathmoveto{\pgfqpoint{1.110111in}{10.351665in}}%
\pgfpathquadraticcurveto{\pgfqpoint{1.095701in}{10.344200in}}{\pgfqpoint{1.080120in}{10.340424in}}%
\pgfpathquadraticcurveto{\pgfqpoint{1.064538in}{10.336604in}}{\pgfqpoint{1.047568in}{10.336604in}}%
\pgfpathquadraticcurveto{\pgfqpoint{0.996960in}{10.336604in}}{\pgfqpoint{0.967403in}{10.364903in}}%
\pgfpathquadraticcurveto{\pgfqpoint{0.937845in}{10.393202in}}{\pgfqpoint{0.937845in}{10.441596in}}%
\pgfpathquadraticcurveto{\pgfqpoint{0.937845in}{10.490163in}}{\pgfqpoint{0.967403in}{10.518419in}}%
\pgfpathquadraticcurveto{\pgfqpoint{0.996960in}{10.546717in}}{\pgfqpoint{1.047568in}{10.546717in}}%
\pgfpathquadraticcurveto{\pgfqpoint{1.064538in}{10.546717in}}{\pgfqpoint{1.080120in}{10.542898in}}%
\pgfpathquadraticcurveto{\pgfqpoint{1.095701in}{10.539122in}}{\pgfqpoint{1.110111in}{10.531656in}}%
\pgfpathlineto{\pgfqpoint{1.110111in}{10.489773in}}%
\pgfpathquadraticcurveto{\pgfqpoint{1.095571in}{10.499669in}}{\pgfqpoint{1.081465in}{10.504269in}}%
\pgfpathquadraticcurveto{\pgfqpoint{1.067359in}{10.508870in}}{\pgfqpoint{1.051778in}{10.508870in}}%
\pgfpathquadraticcurveto{\pgfqpoint{1.023826in}{10.508870in}}{\pgfqpoint{1.007811in}{10.490945in}}%
\pgfpathquadraticcurveto{\pgfqpoint{0.991839in}{10.473063in}}{\pgfqpoint{0.991839in}{10.441596in}}%
\pgfpathquadraticcurveto{\pgfqpoint{0.991839in}{10.410259in}}{\pgfqpoint{1.007811in}{10.392334in}}%
\pgfpathquadraticcurveto{\pgfqpoint{1.023826in}{10.374452in}}{\pgfqpoint{1.051778in}{10.374452in}}%
\pgfpathquadraticcurveto{\pgfqpoint{1.067359in}{10.374452in}}{\pgfqpoint{1.081465in}{10.379052in}}%
\pgfpathquadraticcurveto{\pgfqpoint{1.095571in}{10.383696in}}{\pgfqpoint{1.110111in}{10.393592in}}%
\pgfpathlineto{\pgfqpoint{1.110111in}{10.351665in}}%
\pgfpathlineto{\pgfqpoint{1.110111in}{10.351665in}}%
\pgfpathclose%
\pgfpathmoveto{\pgfqpoint{1.299434in}{10.536691in}}%
\pgfpathlineto{\pgfqpoint{1.299434in}{10.493809in}}%
\pgfpathquadraticcurveto{\pgfqpoint{1.282767in}{10.501274in}}{\pgfqpoint{1.266882in}{10.505051in}}%
\pgfpathquadraticcurveto{\pgfqpoint{1.251040in}{10.508870in}}{\pgfqpoint{1.236934in}{10.508870in}}%
\pgfpathquadraticcurveto{\pgfqpoint{1.218227in}{10.508870in}}{\pgfqpoint{1.209243in}{10.503705in}}%
\pgfpathquadraticcurveto{\pgfqpoint{1.200302in}{10.498584in}}{\pgfqpoint{1.200302in}{10.487733in}}%
\pgfpathquadraticcurveto{\pgfqpoint{1.200302in}{10.479573in}}{\pgfqpoint{1.206335in}{10.475016in}}%
\pgfpathquadraticcurveto{\pgfqpoint{1.212368in}{10.470502in}}{\pgfqpoint{1.228253in}{10.467247in}}%
\pgfpathlineto{\pgfqpoint{1.250476in}{10.462776in}}%
\pgfpathquadraticcurveto{\pgfqpoint{1.284243in}{10.455962in}}{\pgfqpoint{1.298479in}{10.442116in}}%
\pgfpathquadraticcurveto{\pgfqpoint{1.312759in}{10.428314in}}{\pgfqpoint{1.312759in}{10.402794in}}%
\pgfpathquadraticcurveto{\pgfqpoint{1.312759in}{10.369330in}}{\pgfqpoint{1.292880in}{10.352967in}}%
\pgfpathquadraticcurveto{\pgfqpoint{1.273002in}{10.336604in}}{\pgfqpoint{1.232160in}{10.336604in}}%
\pgfpathquadraticcurveto{\pgfqpoint{1.212932in}{10.336604in}}{\pgfqpoint{1.193531in}{10.340294in}}%
\pgfpathquadraticcurveto{\pgfqpoint{1.174130in}{10.343939in}}{\pgfqpoint{1.154729in}{10.351144in}}%
\pgfpathlineto{\pgfqpoint{1.154729in}{10.395198in}}%
\pgfpathquadraticcurveto{\pgfqpoint{1.174130in}{10.384912in}}{\pgfqpoint{1.192229in}{10.379660in}}%
\pgfpathquadraticcurveto{\pgfqpoint{1.210328in}{10.374452in}}{\pgfqpoint{1.227168in}{10.374452in}}%
\pgfpathquadraticcurveto{\pgfqpoint{1.244269in}{10.374452in}}{\pgfqpoint{1.253340in}{10.380137in}}%
\pgfpathquadraticcurveto{\pgfqpoint{1.262411in}{10.385866in}}{\pgfqpoint{1.262411in}{10.396457in}}%
\pgfpathquadraticcurveto{\pgfqpoint{1.262411in}{10.405919in}}{\pgfqpoint{1.256248in}{10.411084in}}%
\pgfpathquadraticcurveto{\pgfqpoint{1.250085in}{10.416248in}}{\pgfqpoint{1.231639in}{10.420328in}}%
\pgfpathlineto{\pgfqpoint{1.211413in}{10.424799in}}%
\pgfpathquadraticcurveto{\pgfqpoint{1.181031in}{10.431309in}}{\pgfqpoint{1.166969in}{10.445545in}}%
\pgfpathquadraticcurveto{\pgfqpoint{1.152950in}{10.459781in}}{\pgfqpoint{1.152950in}{10.483913in}}%
\pgfpathquadraticcurveto{\pgfqpoint{1.152950in}{10.514165in}}{\pgfqpoint{1.172481in}{10.530441in}}%
\pgfpathquadraticcurveto{\pgfqpoint{1.192012in}{10.546717in}}{\pgfqpoint{1.228644in}{10.546717in}}%
\pgfpathquadraticcurveto{\pgfqpoint{1.245354in}{10.546717in}}{\pgfqpoint{1.262976in}{10.544200in}}%
\pgfpathquadraticcurveto{\pgfqpoint{1.280597in}{10.541682in}}{\pgfqpoint{1.299434in}{10.536691in}}%
\pgfpathlineto{\pgfqpoint{1.299434in}{10.536691in}}%
\pgfpathclose%
\pgfpathmoveto{\pgfqpoint{1.481465in}{10.377446in}}%
\pgfpathlineto{\pgfqpoint{1.399825in}{10.377446in}}%
\pgfpathlineto{\pgfqpoint{1.386934in}{10.340554in}}%
\pgfpathlineto{\pgfqpoint{1.334417in}{10.340554in}}%
\pgfpathlineto{\pgfqpoint{1.409417in}{10.543071in}}%
\pgfpathlineto{\pgfqpoint{1.471700in}{10.543071in}}%
\pgfpathlineto{\pgfqpoint{1.546700in}{10.340554in}}%
\pgfpathlineto{\pgfqpoint{1.494226in}{10.340554in}}%
\pgfpathlineto{\pgfqpoint{1.481465in}{10.377446in}}%
\pgfpathlineto{\pgfqpoint{1.481465in}{10.377446in}}%
\pgfpathclose%
\pgfpathmoveto{\pgfqpoint{1.412845in}{10.415033in}}%
\pgfpathlineto{\pgfqpoint{1.468314in}{10.415033in}}%
\pgfpathlineto{\pgfqpoint{1.440623in}{10.495589in}}%
\pgfpathlineto{\pgfqpoint{1.412845in}{10.415033in}}%
\pgfpathlineto{\pgfqpoint{1.412845in}{10.415033in}}%
\pgfpathclose%
\pgfpathmoveto{\pgfqpoint{1.736153in}{10.408913in}}%
\pgfpathquadraticcurveto{\pgfqpoint{1.720962in}{10.408913in}}{\pgfqpoint{1.713280in}{10.403748in}}%
\pgfpathquadraticcurveto{\pgfqpoint{1.705641in}{10.398627in}}{\pgfqpoint{1.705641in}{10.388557in}}%
\pgfpathquadraticcurveto{\pgfqpoint{1.705641in}{10.379356in}}{\pgfqpoint{1.711804in}{10.374104in}}%
\pgfpathquadraticcurveto{\pgfqpoint{1.718010in}{10.368896in}}{\pgfqpoint{1.728991in}{10.368896in}}%
\pgfpathquadraticcurveto{\pgfqpoint{1.742663in}{10.368896in}}{\pgfqpoint{1.752038in}{10.378705in}}%
\pgfpathquadraticcurveto{\pgfqpoint{1.761413in}{10.388557in}}{\pgfqpoint{1.761413in}{10.403358in}}%
\pgfpathlineto{\pgfqpoint{1.761413in}{10.408913in}}%
\pgfpathlineto{\pgfqpoint{1.736153in}{10.408913in}}%
\pgfpathlineto{\pgfqpoint{1.736153in}{10.408913in}}%
\pgfpathclose%
\pgfpathmoveto{\pgfqpoint{1.810372in}{10.427229in}}%
\pgfpathlineto{\pgfqpoint{1.810372in}{10.340554in}}%
\pgfpathlineto{\pgfqpoint{1.761413in}{10.340554in}}%
\pgfpathlineto{\pgfqpoint{1.761413in}{10.363080in}}%
\pgfpathquadraticcurveto{\pgfqpoint{1.751648in}{10.349235in}}{\pgfqpoint{1.739408in}{10.342898in}}%
\pgfpathquadraticcurveto{\pgfqpoint{1.727212in}{10.336604in}}{\pgfqpoint{1.709720in}{10.336604in}}%
\pgfpathquadraticcurveto{\pgfqpoint{1.686109in}{10.336604in}}{\pgfqpoint{1.671396in}{10.350363in}}%
\pgfpathquadraticcurveto{\pgfqpoint{1.656682in}{10.364165in}}{\pgfqpoint{1.656682in}{10.386127in}}%
\pgfpathquadraticcurveto{\pgfqpoint{1.656682in}{10.412863in}}{\pgfqpoint{1.675042in}{10.425320in}}%
\pgfpathquadraticcurveto{\pgfqpoint{1.693444in}{10.437820in}}{\pgfqpoint{1.732767in}{10.437820in}}%
\pgfpathlineto{\pgfqpoint{1.761413in}{10.437820in}}%
\pgfpathlineto{\pgfqpoint{1.761413in}{10.441596in}}%
\pgfpathquadraticcurveto{\pgfqpoint{1.761413in}{10.453141in}}{\pgfqpoint{1.752299in}{10.458479in}}%
\pgfpathquadraticcurveto{\pgfqpoint{1.743227in}{10.463861in}}{\pgfqpoint{1.723957in}{10.463861in}}%
\pgfpathquadraticcurveto{\pgfqpoint{1.708375in}{10.463861in}}{\pgfqpoint{1.694920in}{10.460736in}}%
\pgfpathquadraticcurveto{\pgfqpoint{1.681509in}{10.457611in}}{\pgfqpoint{1.669964in}{10.451361in}}%
\pgfpathlineto{\pgfqpoint{1.669964in}{10.488384in}}%
\pgfpathquadraticcurveto{\pgfqpoint{1.685589in}{10.492203in}}{\pgfqpoint{1.701300in}{10.494156in}}%
\pgfpathquadraticcurveto{\pgfqpoint{1.717056in}{10.496110in}}{\pgfqpoint{1.732767in}{10.496110in}}%
\pgfpathquadraticcurveto{\pgfqpoint{1.773870in}{10.496110in}}{\pgfqpoint{1.792099in}{10.479920in}}%
\pgfpathquadraticcurveto{\pgfqpoint{1.810372in}{10.463731in}}{\pgfqpoint{1.810372in}{10.427229in}}%
\pgfpathlineto{\pgfqpoint{1.810372in}{10.427229in}}%
\pgfpathclose%
\pgfpathmoveto{\pgfqpoint{2.008245in}{10.433045in}}%
\pgfpathlineto{\pgfqpoint{2.008245in}{10.340554in}}%
\pgfpathlineto{\pgfqpoint{1.959417in}{10.340554in}}%
\pgfpathlineto{\pgfqpoint{1.959417in}{10.355615in}}%
\pgfpathlineto{\pgfqpoint{1.959417in}{10.411344in}}%
\pgfpathquadraticcurveto{\pgfqpoint{1.959417in}{10.431005in}}{\pgfqpoint{1.958549in}{10.438471in}}%
\pgfpathquadraticcurveto{\pgfqpoint{1.957681in}{10.445936in}}{\pgfqpoint{1.955510in}{10.449452in}}%
\pgfpathquadraticcurveto{\pgfqpoint{1.952646in}{10.454226in}}{\pgfqpoint{1.947741in}{10.456873in}}%
\pgfpathquadraticcurveto{\pgfqpoint{1.942880in}{10.459521in}}{\pgfqpoint{1.936630in}{10.459521in}}%
\pgfpathquadraticcurveto{\pgfqpoint{1.921439in}{10.459521in}}{\pgfqpoint{1.912759in}{10.447759in}}%
\pgfpathquadraticcurveto{\pgfqpoint{1.904078in}{10.436040in}}{\pgfqpoint{1.904078in}{10.415294in}}%
\pgfpathlineto{\pgfqpoint{1.904078in}{10.340554in}}%
\pgfpathlineto{\pgfqpoint{1.855554in}{10.340554in}}%
\pgfpathlineto{\pgfqpoint{1.855554in}{10.492464in}}%
\pgfpathlineto{\pgfqpoint{1.904078in}{10.492464in}}%
\pgfpathlineto{\pgfqpoint{1.904078in}{10.470241in}}%
\pgfpathquadraticcurveto{\pgfqpoint{1.915059in}{10.483523in}}{\pgfqpoint{1.927385in}{10.489816in}}%
\pgfpathquadraticcurveto{\pgfqpoint{1.939755in}{10.496110in}}{\pgfqpoint{1.954686in}{10.496110in}}%
\pgfpathquadraticcurveto{\pgfqpoint{1.980988in}{10.496110in}}{\pgfqpoint{1.994616in}{10.479964in}}%
\pgfpathquadraticcurveto{\pgfqpoint{2.008245in}{10.463861in}}{\pgfqpoint{2.008245in}{10.433045in}}%
\pgfpathlineto{\pgfqpoint{2.008245in}{10.433045in}}%
\pgfpathclose%
\pgfpathmoveto{\pgfqpoint{2.156639in}{10.470241in}}%
\pgfpathlineto{\pgfqpoint{2.156639in}{10.551622in}}%
\pgfpathlineto{\pgfqpoint{2.205467in}{10.551622in}}%
\pgfpathlineto{\pgfqpoint{2.205467in}{10.340554in}}%
\pgfpathlineto{\pgfqpoint{2.156639in}{10.340554in}}%
\pgfpathlineto{\pgfqpoint{2.156639in}{10.362516in}}%
\pgfpathquadraticcurveto{\pgfqpoint{2.146613in}{10.349104in}}{\pgfqpoint{2.134503in}{10.342854in}}%
\pgfpathquadraticcurveto{\pgfqpoint{2.122438in}{10.336604in}}{\pgfqpoint{2.106595in}{10.336604in}}%
\pgfpathquadraticcurveto{\pgfqpoint{2.078514in}{10.336604in}}{\pgfqpoint{2.060458in}{10.358913in}}%
\pgfpathquadraticcurveto{\pgfqpoint{2.042446in}{10.381266in}}{\pgfqpoint{2.042446in}{10.416379in}}%
\pgfpathquadraticcurveto{\pgfqpoint{2.042446in}{10.451491in}}{\pgfqpoint{2.060458in}{10.473801in}}%
\pgfpathquadraticcurveto{\pgfqpoint{2.078514in}{10.496110in}}{\pgfqpoint{2.106595in}{10.496110in}}%
\pgfpathquadraticcurveto{\pgfqpoint{2.122307in}{10.496110in}}{\pgfqpoint{2.134460in}{10.489816in}}%
\pgfpathquadraticcurveto{\pgfqpoint{2.146613in}{10.483523in}}{\pgfqpoint{2.156639in}{10.470241in}}%
\pgfpathlineto{\pgfqpoint{2.156639in}{10.470241in}}%
\pgfpathclose%
\pgfpathmoveto{\pgfqpoint{2.124608in}{10.371891in}}%
\pgfpathquadraticcurveto{\pgfqpoint{2.140233in}{10.371891in}}{\pgfqpoint{2.148436in}{10.383262in}}%
\pgfpathquadraticcurveto{\pgfqpoint{2.156639in}{10.394677in}}{\pgfqpoint{2.156639in}{10.416379in}}%
\pgfpathquadraticcurveto{\pgfqpoint{2.156639in}{10.438080in}}{\pgfqpoint{2.148436in}{10.449452in}}%
\pgfpathquadraticcurveto{\pgfqpoint{2.140233in}{10.460866in}}{\pgfqpoint{2.124608in}{10.460866in}}%
\pgfpathquadraticcurveto{\pgfqpoint{2.109156in}{10.460866in}}{\pgfqpoint{2.100953in}{10.449452in}}%
\pgfpathquadraticcurveto{\pgfqpoint{2.092750in}{10.438080in}}{\pgfqpoint{2.092750in}{10.416379in}}%
\pgfpathquadraticcurveto{\pgfqpoint{2.092750in}{10.394677in}}{\pgfqpoint{2.100953in}{10.383262in}}%
\pgfpathquadraticcurveto{\pgfqpoint{2.109156in}{10.371891in}}{\pgfqpoint{2.124608in}{10.371891in}}%
\pgfpathlineto{\pgfqpoint{2.124608in}{10.371891in}}%
\pgfpathclose%
\pgfpathmoveto{\pgfqpoint{2.432073in}{10.464643in}}%
\pgfpathquadraticcurveto{\pgfqpoint{2.444443in}{10.464643in}}{\pgfqpoint{2.450780in}{10.470068in}}%
\pgfpathquadraticcurveto{\pgfqpoint{2.457160in}{10.475493in}}{\pgfqpoint{2.457160in}{10.486084in}}%
\pgfpathquadraticcurveto{\pgfqpoint{2.457160in}{10.496544in}}{\pgfqpoint{2.450780in}{10.502012in}}%
\pgfpathquadraticcurveto{\pgfqpoint{2.444443in}{10.507524in}}{\pgfqpoint{2.432073in}{10.507524in}}%
\pgfpathlineto{\pgfqpoint{2.403210in}{10.507524in}}%
\pgfpathlineto{\pgfqpoint{2.403210in}{10.464643in}}%
\pgfpathlineto{\pgfqpoint{2.432073in}{10.464643in}}%
\pgfpathlineto{\pgfqpoint{2.432073in}{10.464643in}}%
\pgfpathclose%
\pgfpathmoveto{\pgfqpoint{2.433852in}{10.376101in}}%
\pgfpathquadraticcurveto{\pgfqpoint{2.449564in}{10.376101in}}{\pgfqpoint{2.457507in}{10.382741in}}%
\pgfpathquadraticcurveto{\pgfqpoint{2.465450in}{10.389382in}}{\pgfqpoint{2.465450in}{10.402794in}}%
\pgfpathquadraticcurveto{\pgfqpoint{2.465450in}{10.415988in}}{\pgfqpoint{2.457594in}{10.422542in}}%
\pgfpathquadraticcurveto{\pgfqpoint{2.449738in}{10.429139in}}{\pgfqpoint{2.433852in}{10.429139in}}%
\pgfpathlineto{\pgfqpoint{2.403210in}{10.429139in}}%
\pgfpathlineto{\pgfqpoint{2.403210in}{10.376101in}}%
\pgfpathlineto{\pgfqpoint{2.433852in}{10.376101in}}%
\pgfpathlineto{\pgfqpoint{2.433852in}{10.376101in}}%
\pgfpathclose%
\pgfpathmoveto{\pgfqpoint{2.482420in}{10.448931in}}%
\pgfpathquadraticcurveto{\pgfqpoint{2.499217in}{10.444026in}}{\pgfqpoint{2.508418in}{10.430875in}}%
\pgfpathquadraticcurveto{\pgfqpoint{2.517663in}{10.417724in}}{\pgfqpoint{2.517663in}{10.398627in}}%
\pgfpathquadraticcurveto{\pgfqpoint{2.517663in}{10.369330in}}{\pgfqpoint{2.497872in}{10.354920in}}%
\pgfpathquadraticcurveto{\pgfqpoint{2.478080in}{10.340554in}}{\pgfqpoint{2.437628in}{10.340554in}}%
\pgfpathlineto{\pgfqpoint{2.350997in}{10.340554in}}%
\pgfpathlineto{\pgfqpoint{2.350997in}{10.543071in}}%
\pgfpathlineto{\pgfqpoint{2.429382in}{10.543071in}}%
\pgfpathquadraticcurveto{\pgfqpoint{2.471569in}{10.543071in}}{\pgfqpoint{2.490493in}{10.530311in}}%
\pgfpathquadraticcurveto{\pgfqpoint{2.509417in}{10.517551in}}{\pgfqpoint{2.509417in}{10.489469in}}%
\pgfpathquadraticcurveto{\pgfqpoint{2.509417in}{10.474712in}}{\pgfqpoint{2.502472in}{10.464339in}}%
\pgfpathquadraticcurveto{\pgfqpoint{2.495571in}{10.453965in}}{\pgfqpoint{2.482420in}{10.448931in}}%
\pgfpathlineto{\pgfqpoint{2.482420in}{10.448931in}}%
\pgfpathclose%
\pgfpathmoveto{\pgfqpoint{2.562715in}{10.543071in}}%
\pgfpathlineto{\pgfqpoint{2.614929in}{10.543071in}}%
\pgfpathlineto{\pgfqpoint{2.614929in}{10.360215in}}%
\pgfpathquadraticcurveto{\pgfqpoint{2.614929in}{10.322368in}}{\pgfqpoint{2.594356in}{10.303662in}}%
\pgfpathquadraticcurveto{\pgfqpoint{2.573826in}{10.284955in}}{\pgfqpoint{2.532160in}{10.284955in}}%
\pgfpathlineto{\pgfqpoint{2.521613in}{10.284955in}}%
\pgfpathlineto{\pgfqpoint{2.521613in}{10.324408in}}%
\pgfpathlineto{\pgfqpoint{2.529729in}{10.324408in}}%
\pgfpathquadraticcurveto{\pgfqpoint{2.546005in}{10.324408in}}{\pgfqpoint{2.554339in}{10.333523in}}%
\pgfpathquadraticcurveto{\pgfqpoint{2.562715in}{10.342594in}}{\pgfqpoint{2.562715in}{10.360215in}}%
\pgfpathlineto{\pgfqpoint{2.562715in}{10.543071in}}%
\pgfpathlineto{\pgfqpoint{2.562715in}{10.543071in}}%
\pgfpathclose%
\pgfpathmoveto{\pgfqpoint{2.806943in}{10.536691in}}%
\pgfpathlineto{\pgfqpoint{2.806943in}{10.493809in}}%
\pgfpathquadraticcurveto{\pgfqpoint{2.790276in}{10.501274in}}{\pgfqpoint{2.774391in}{10.505051in}}%
\pgfpathquadraticcurveto{\pgfqpoint{2.758549in}{10.508870in}}{\pgfqpoint{2.744443in}{10.508870in}}%
\pgfpathquadraticcurveto{\pgfqpoint{2.725736in}{10.508870in}}{\pgfqpoint{2.716752in}{10.503705in}}%
\pgfpathquadraticcurveto{\pgfqpoint{2.707811in}{10.498584in}}{\pgfqpoint{2.707811in}{10.487733in}}%
\pgfpathquadraticcurveto{\pgfqpoint{2.707811in}{10.479573in}}{\pgfqpoint{2.713844in}{10.475016in}}%
\pgfpathquadraticcurveto{\pgfqpoint{2.719877in}{10.470502in}}{\pgfqpoint{2.735762in}{10.467247in}}%
\pgfpathlineto{\pgfqpoint{2.757984in}{10.462776in}}%
\pgfpathquadraticcurveto{\pgfqpoint{2.791752in}{10.455962in}}{\pgfqpoint{2.805988in}{10.442116in}}%
\pgfpathquadraticcurveto{\pgfqpoint{2.820267in}{10.428314in}}{\pgfqpoint{2.820267in}{10.402794in}}%
\pgfpathquadraticcurveto{\pgfqpoint{2.820267in}{10.369330in}}{\pgfqpoint{2.800389in}{10.352967in}}%
\pgfpathquadraticcurveto{\pgfqpoint{2.780510in}{10.336604in}}{\pgfqpoint{2.739668in}{10.336604in}}%
\pgfpathquadraticcurveto{\pgfqpoint{2.720441in}{10.336604in}}{\pgfqpoint{2.701040in}{10.340294in}}%
\pgfpathquadraticcurveto{\pgfqpoint{2.681639in}{10.343939in}}{\pgfqpoint{2.662238in}{10.351144in}}%
\pgfpathlineto{\pgfqpoint{2.662238in}{10.395198in}}%
\pgfpathquadraticcurveto{\pgfqpoint{2.681639in}{10.384912in}}{\pgfqpoint{2.699738in}{10.379660in}}%
\pgfpathquadraticcurveto{\pgfqpoint{2.717837in}{10.374452in}}{\pgfqpoint{2.734677in}{10.374452in}}%
\pgfpathquadraticcurveto{\pgfqpoint{2.751778in}{10.374452in}}{\pgfqpoint{2.760849in}{10.380137in}}%
\pgfpathquadraticcurveto{\pgfqpoint{2.769920in}{10.385866in}}{\pgfqpoint{2.769920in}{10.396457in}}%
\pgfpathquadraticcurveto{\pgfqpoint{2.769920in}{10.405919in}}{\pgfqpoint{2.763757in}{10.411084in}}%
\pgfpathquadraticcurveto{\pgfqpoint{2.757594in}{10.416248in}}{\pgfqpoint{2.739148in}{10.420328in}}%
\pgfpathlineto{\pgfqpoint{2.718922in}{10.424799in}}%
\pgfpathquadraticcurveto{\pgfqpoint{2.688540in}{10.431309in}}{\pgfqpoint{2.674477in}{10.445545in}}%
\pgfpathquadraticcurveto{\pgfqpoint{2.660458in}{10.459781in}}{\pgfqpoint{2.660458in}{10.483913in}}%
\pgfpathquadraticcurveto{\pgfqpoint{2.660458in}{10.514165in}}{\pgfqpoint{2.679990in}{10.530441in}}%
\pgfpathquadraticcurveto{\pgfqpoint{2.699521in}{10.546717in}}{\pgfqpoint{2.736153in}{10.546717in}}%
\pgfpathquadraticcurveto{\pgfqpoint{2.752863in}{10.546717in}}{\pgfqpoint{2.770484in}{10.544200in}}%
\pgfpathquadraticcurveto{\pgfqpoint{2.788106in}{10.541682in}}{\pgfqpoint{2.806943in}{10.536691in}}%
\pgfpathlineto{\pgfqpoint{2.806943in}{10.536691in}}%
\pgfpathclose%
\pgfpathmoveto{\pgfqpoint{2.941491in}{10.492464in}}%
\pgfpathlineto{\pgfqpoint{2.990059in}{10.492464in}}%
\pgfpathlineto{\pgfqpoint{3.027906in}{10.387472in}}%
\pgfpathlineto{\pgfqpoint{3.065580in}{10.492464in}}%
\pgfpathlineto{\pgfqpoint{3.114278in}{10.492464in}}%
\pgfpathlineto{\pgfqpoint{3.054469in}{10.340554in}}%
\pgfpathlineto{\pgfqpoint{3.001170in}{10.340554in}}%
\pgfpathlineto{\pgfqpoint{2.941491in}{10.492464in}}%
\pgfpathlineto{\pgfqpoint{2.941491in}{10.492464in}}%
\pgfpathclose%
\pgfpathmoveto{\pgfqpoint{3.293314in}{10.416899in}}%
\pgfpathlineto{\pgfqpoint{3.293314in}{10.403097in}}%
\pgfpathlineto{\pgfqpoint{3.179816in}{10.403097in}}%
\pgfpathquadraticcurveto{\pgfqpoint{3.181552in}{10.385997in}}{\pgfqpoint{3.192142in}{10.377446in}}%
\pgfpathquadraticcurveto{\pgfqpoint{3.202733in}{10.368896in}}{\pgfqpoint{3.221700in}{10.368896in}}%
\pgfpathquadraticcurveto{\pgfqpoint{3.237021in}{10.368896in}}{\pgfqpoint{3.253080in}{10.373453in}}%
\pgfpathquadraticcurveto{\pgfqpoint{3.269182in}{10.378011in}}{\pgfqpoint{3.286153in}{10.387212in}}%
\pgfpathlineto{\pgfqpoint{3.286153in}{10.349799in}}%
\pgfpathquadraticcurveto{\pgfqpoint{3.268922in}{10.343288in}}{\pgfqpoint{3.251691in}{10.339946in}}%
\pgfpathquadraticcurveto{\pgfqpoint{3.234460in}{10.336604in}}{\pgfqpoint{3.217229in}{10.336604in}}%
\pgfpathquadraticcurveto{\pgfqpoint{3.175997in}{10.336604in}}{\pgfqpoint{3.153123in}{10.357568in}}%
\pgfpathquadraticcurveto{\pgfqpoint{3.130293in}{10.378531in}}{\pgfqpoint{3.130293in}{10.416379in}}%
\pgfpathquadraticcurveto{\pgfqpoint{3.130293in}{10.453531in}}{\pgfqpoint{3.152733in}{10.474799in}}%
\pgfpathquadraticcurveto{\pgfqpoint{3.175172in}{10.496110in}}{\pgfqpoint{3.214538in}{10.496110in}}%
\pgfpathquadraticcurveto{\pgfqpoint{3.250345in}{10.496110in}}{\pgfqpoint{3.271830in}{10.474538in}}%
\pgfpathquadraticcurveto{\pgfqpoint{3.293314in}{10.453011in}}{\pgfqpoint{3.293314in}{10.416899in}}%
\pgfpathlineto{\pgfqpoint{3.293314in}{10.416899in}}%
\pgfpathclose%
\pgfpathmoveto{\pgfqpoint{3.243401in}{10.433045in}}%
\pgfpathquadraticcurveto{\pgfqpoint{3.243401in}{10.446891in}}{\pgfqpoint{3.235328in}{10.455354in}}%
\pgfpathquadraticcurveto{\pgfqpoint{3.227255in}{10.463861in}}{\pgfqpoint{3.214234in}{10.463861in}}%
\pgfpathquadraticcurveto{\pgfqpoint{3.200128in}{10.463861in}}{\pgfqpoint{3.191318in}{10.455919in}}%
\pgfpathquadraticcurveto{\pgfqpoint{3.182507in}{10.447976in}}{\pgfqpoint{3.180337in}{10.433045in}}%
\pgfpathlineto{\pgfqpoint{3.243401in}{10.433045in}}%
\pgfpathlineto{\pgfqpoint{3.243401in}{10.433045in}}%
\pgfpathclose%
\pgfpathmoveto{\pgfqpoint{3.442967in}{10.451101in}}%
\pgfpathquadraticcurveto{\pgfqpoint{3.436587in}{10.454096in}}{\pgfqpoint{3.430250in}{10.455485in}}%
\pgfpathquadraticcurveto{\pgfqpoint{3.423957in}{10.456917in}}{\pgfqpoint{3.417576in}{10.456917in}}%
\pgfpathquadraticcurveto{\pgfqpoint{3.398870in}{10.456917in}}{\pgfqpoint{3.388757in}{10.444894in}}%
\pgfpathquadraticcurveto{\pgfqpoint{3.378644in}{10.432915in}}{\pgfqpoint{3.378644in}{10.410563in}}%
\pgfpathlineto{\pgfqpoint{3.378644in}{10.340554in}}%
\pgfpathlineto{\pgfqpoint{3.330120in}{10.340554in}}%
\pgfpathlineto{\pgfqpoint{3.330120in}{10.492464in}}%
\pgfpathlineto{\pgfqpoint{3.378644in}{10.492464in}}%
\pgfpathlineto{\pgfqpoint{3.378644in}{10.467507in}}%
\pgfpathquadraticcurveto{\pgfqpoint{3.388019in}{10.482438in}}{\pgfqpoint{3.400128in}{10.489252in}}%
\pgfpathquadraticcurveto{\pgfqpoint{3.412281in}{10.496110in}}{\pgfqpoint{3.429252in}{10.496110in}}%
\pgfpathquadraticcurveto{\pgfqpoint{3.431682in}{10.496110in}}{\pgfqpoint{3.434503in}{10.495893in}}%
\pgfpathquadraticcurveto{\pgfqpoint{3.437368in}{10.495719in}}{\pgfqpoint{3.442793in}{10.495024in}}%
\pgfpathlineto{\pgfqpoint{3.442967in}{10.451101in}}%
\pgfpathlineto{\pgfqpoint{3.442967in}{10.451101in}}%
\pgfpathclose%
\pgfpathmoveto{\pgfqpoint{3.585762in}{10.487733in}}%
\pgfpathlineto{\pgfqpoint{3.585762in}{10.450840in}}%
\pgfpathquadraticcurveto{\pgfqpoint{3.570181in}{10.457351in}}{\pgfqpoint{3.555641in}{10.460606in}}%
\pgfpathquadraticcurveto{\pgfqpoint{3.541144in}{10.463861in}}{\pgfqpoint{3.528253in}{10.463861in}}%
\pgfpathquadraticcurveto{\pgfqpoint{3.514408in}{10.463861in}}{\pgfqpoint{3.507681in}{10.460389in}}%
\pgfpathquadraticcurveto{\pgfqpoint{3.500997in}{10.456917in}}{\pgfqpoint{3.500997in}{10.449755in}}%
\pgfpathquadraticcurveto{\pgfqpoint{3.500997in}{10.443896in}}{\pgfqpoint{3.506075in}{10.440771in}}%
\pgfpathquadraticcurveto{\pgfqpoint{3.511153in}{10.437689in}}{\pgfqpoint{3.524304in}{10.436170in}}%
\pgfpathlineto{\pgfqpoint{3.532854in}{10.434955in}}%
\pgfpathquadraticcurveto{\pgfqpoint{3.570181in}{10.430224in}}{\pgfqpoint{3.583028in}{10.419373in}}%
\pgfpathquadraticcurveto{\pgfqpoint{3.595918in}{10.408523in}}{\pgfqpoint{3.595918in}{10.385302in}}%
\pgfpathquadraticcurveto{\pgfqpoint{3.595918in}{10.361040in}}{\pgfqpoint{3.577993in}{10.348801in}}%
\pgfpathquadraticcurveto{\pgfqpoint{3.560111in}{10.336604in}}{\pgfqpoint{3.524608in}{10.336604in}}%
\pgfpathquadraticcurveto{\pgfqpoint{3.509547in}{10.336604in}}{\pgfqpoint{3.493444in}{10.338991in}}%
\pgfpathquadraticcurveto{\pgfqpoint{3.477385in}{10.341379in}}{\pgfqpoint{3.460415in}{10.346110in}}%
\pgfpathlineto{\pgfqpoint{3.460415in}{10.383002in}}%
\pgfpathquadraticcurveto{\pgfqpoint{3.474955in}{10.375971in}}{\pgfqpoint{3.490189in}{10.372412in}}%
\pgfpathquadraticcurveto{\pgfqpoint{3.505467in}{10.368896in}}{\pgfqpoint{3.521179in}{10.368896in}}%
\pgfpathquadraticcurveto{\pgfqpoint{3.535458in}{10.368896in}}{\pgfqpoint{3.542620in}{10.372802in}}%
\pgfpathquadraticcurveto{\pgfqpoint{3.549825in}{10.376752in}}{\pgfqpoint{3.549825in}{10.384521in}}%
\pgfpathquadraticcurveto{\pgfqpoint{3.549825in}{10.391031in}}{\pgfqpoint{3.544877in}{10.394200in}}%
\pgfpathquadraticcurveto{\pgfqpoint{3.539929in}{10.397368in}}{\pgfqpoint{3.525128in}{10.399148in}}%
\pgfpathlineto{\pgfqpoint{3.516578in}{10.400233in}}%
\pgfpathquadraticcurveto{\pgfqpoint{3.484156in}{10.404313in}}{\pgfqpoint{3.471135in}{10.415294in}}%
\pgfpathquadraticcurveto{\pgfqpoint{3.458115in}{10.426274in}}{\pgfqpoint{3.458115in}{10.448670in}}%
\pgfpathquadraticcurveto{\pgfqpoint{3.458115in}{10.472802in}}{\pgfqpoint{3.474651in}{10.484434in}}%
\pgfpathquadraticcurveto{\pgfqpoint{3.491231in}{10.496110in}}{\pgfqpoint{3.525389in}{10.496110in}}%
\pgfpathquadraticcurveto{\pgfqpoint{3.538844in}{10.496110in}}{\pgfqpoint{3.553601in}{10.494070in}}%
\pgfpathquadraticcurveto{\pgfqpoint{3.568401in}{10.492073in}}{\pgfqpoint{3.585762in}{10.487733in}}%
\pgfpathlineto{\pgfqpoint{3.585762in}{10.487733in}}%
\pgfpathclose%
\pgfpathmoveto{\pgfqpoint{3.630771in}{10.399712in}}%
\pgfpathlineto{\pgfqpoint{3.630771in}{10.492464in}}%
\pgfpathlineto{\pgfqpoint{3.679599in}{10.492464in}}%
\pgfpathlineto{\pgfqpoint{3.679599in}{10.477273in}}%
\pgfpathquadraticcurveto{\pgfqpoint{3.679599in}{10.464946in}}{\pgfqpoint{3.679469in}{10.446283in}}%
\pgfpathquadraticcurveto{\pgfqpoint{3.679339in}{10.427620in}}{\pgfqpoint{3.679339in}{10.421413in}}%
\pgfpathquadraticcurveto{\pgfqpoint{3.679339in}{10.403097in}}{\pgfqpoint{3.680293in}{10.395024in}}%
\pgfpathquadraticcurveto{\pgfqpoint{3.681248in}{10.386952in}}{\pgfqpoint{3.683549in}{10.383262in}}%
\pgfpathquadraticcurveto{\pgfqpoint{3.686500in}{10.378531in}}{\pgfqpoint{3.691318in}{10.375927in}}%
\pgfpathquadraticcurveto{\pgfqpoint{3.696135in}{10.373366in}}{\pgfqpoint{3.702385in}{10.373366in}}%
\pgfpathquadraticcurveto{\pgfqpoint{3.717576in}{10.373366in}}{\pgfqpoint{3.726257in}{10.385042in}}%
\pgfpathquadraticcurveto{\pgfqpoint{3.734938in}{10.396717in}}{\pgfqpoint{3.734938in}{10.417464in}}%
\pgfpathlineto{\pgfqpoint{3.734938in}{10.492464in}}%
\pgfpathlineto{\pgfqpoint{3.783505in}{10.492464in}}%
\pgfpathlineto{\pgfqpoint{3.783505in}{10.340554in}}%
\pgfpathlineto{\pgfqpoint{3.734938in}{10.340554in}}%
\pgfpathlineto{\pgfqpoint{3.734938in}{10.362516in}}%
\pgfpathquadraticcurveto{\pgfqpoint{3.723957in}{10.349235in}}{\pgfqpoint{3.711674in}{10.342898in}}%
\pgfpathquadraticcurveto{\pgfqpoint{3.699391in}{10.336604in}}{\pgfqpoint{3.684634in}{10.336604in}}%
\pgfpathquadraticcurveto{\pgfqpoint{3.658288in}{10.336604in}}{\pgfqpoint{3.644530in}{10.352750in}}%
\pgfpathquadraticcurveto{\pgfqpoint{3.630771in}{10.368896in}}{\pgfqpoint{3.630771in}{10.399712in}}%
\pgfpathlineto{\pgfqpoint{3.630771in}{10.399712in}}%
\pgfpathclose%
\pgfpathmoveto{\pgfqpoint{3.948826in}{10.487733in}}%
\pgfpathlineto{\pgfqpoint{3.948826in}{10.450840in}}%
\pgfpathquadraticcurveto{\pgfqpoint{3.933245in}{10.457351in}}{\pgfqpoint{3.918705in}{10.460606in}}%
\pgfpathquadraticcurveto{\pgfqpoint{3.904208in}{10.463861in}}{\pgfqpoint{3.891318in}{10.463861in}}%
\pgfpathquadraticcurveto{\pgfqpoint{3.877472in}{10.463861in}}{\pgfqpoint{3.870745in}{10.460389in}}%
\pgfpathquadraticcurveto{\pgfqpoint{3.864061in}{10.456917in}}{\pgfqpoint{3.864061in}{10.449755in}}%
\pgfpathquadraticcurveto{\pgfqpoint{3.864061in}{10.443896in}}{\pgfqpoint{3.869139in}{10.440771in}}%
\pgfpathquadraticcurveto{\pgfqpoint{3.874217in}{10.437689in}}{\pgfqpoint{3.887368in}{10.436170in}}%
\pgfpathlineto{\pgfqpoint{3.895918in}{10.434955in}}%
\pgfpathquadraticcurveto{\pgfqpoint{3.933245in}{10.430224in}}{\pgfqpoint{3.946092in}{10.419373in}}%
\pgfpathquadraticcurveto{\pgfqpoint{3.958983in}{10.408523in}}{\pgfqpoint{3.958983in}{10.385302in}}%
\pgfpathquadraticcurveto{\pgfqpoint{3.958983in}{10.361040in}}{\pgfqpoint{3.941057in}{10.348801in}}%
\pgfpathquadraticcurveto{\pgfqpoint{3.923175in}{10.336604in}}{\pgfqpoint{3.887672in}{10.336604in}}%
\pgfpathquadraticcurveto{\pgfqpoint{3.872611in}{10.336604in}}{\pgfqpoint{3.856509in}{10.338991in}}%
\pgfpathquadraticcurveto{\pgfqpoint{3.840450in}{10.341379in}}{\pgfqpoint{3.823479in}{10.346110in}}%
\pgfpathlineto{\pgfqpoint{3.823479in}{10.383002in}}%
\pgfpathquadraticcurveto{\pgfqpoint{3.838019in}{10.375971in}}{\pgfqpoint{3.853253in}{10.372412in}}%
\pgfpathquadraticcurveto{\pgfqpoint{3.868531in}{10.368896in}}{\pgfqpoint{3.884243in}{10.368896in}}%
\pgfpathquadraticcurveto{\pgfqpoint{3.898523in}{10.368896in}}{\pgfqpoint{3.905684in}{10.372802in}}%
\pgfpathquadraticcurveto{\pgfqpoint{3.912889in}{10.376752in}}{\pgfqpoint{3.912889in}{10.384521in}}%
\pgfpathquadraticcurveto{\pgfqpoint{3.912889in}{10.391031in}}{\pgfqpoint{3.907941in}{10.394200in}}%
\pgfpathquadraticcurveto{\pgfqpoint{3.902993in}{10.397368in}}{\pgfqpoint{3.888193in}{10.399148in}}%
\pgfpathlineto{\pgfqpoint{3.879642in}{10.400233in}}%
\pgfpathquadraticcurveto{\pgfqpoint{3.847220in}{10.404313in}}{\pgfqpoint{3.834200in}{10.415294in}}%
\pgfpathquadraticcurveto{\pgfqpoint{3.821179in}{10.426274in}}{\pgfqpoint{3.821179in}{10.448670in}}%
\pgfpathquadraticcurveto{\pgfqpoint{3.821179in}{10.472802in}}{\pgfqpoint{3.837715in}{10.484434in}}%
\pgfpathquadraticcurveto{\pgfqpoint{3.854295in}{10.496110in}}{\pgfqpoint{3.888453in}{10.496110in}}%
\pgfpathquadraticcurveto{\pgfqpoint{3.901908in}{10.496110in}}{\pgfqpoint{3.916665in}{10.494070in}}%
\pgfpathquadraticcurveto{\pgfqpoint{3.931465in}{10.492073in}}{\pgfqpoint{3.948826in}{10.487733in}}%
\pgfpathlineto{\pgfqpoint{3.948826in}{10.487733in}}%
\pgfpathclose%
\pgfpathmoveto{\pgfqpoint{4.070181in}{10.543071in}}%
\pgfpathlineto{\pgfqpoint{4.256812in}{10.543071in}}%
\pgfpathlineto{\pgfqpoint{4.256812in}{10.503575in}}%
\pgfpathlineto{\pgfqpoint{4.189668in}{10.503575in}}%
\pgfpathlineto{\pgfqpoint{4.189668in}{10.340554in}}%
\pgfpathlineto{\pgfqpoint{4.137455in}{10.340554in}}%
\pgfpathlineto{\pgfqpoint{4.137455in}{10.503575in}}%
\pgfpathlineto{\pgfqpoint{4.070181in}{10.503575in}}%
\pgfpathlineto{\pgfqpoint{4.070181in}{10.543071in}}%
\pgfpathlineto{\pgfqpoint{4.070181in}{10.543071in}}%
\pgfpathclose%
\pgfpathmoveto{\pgfqpoint{4.266622in}{10.543071in}}%
\pgfpathlineto{\pgfqpoint{4.316665in}{10.543071in}}%
\pgfpathlineto{\pgfqpoint{4.351648in}{10.395893in}}%
\pgfpathlineto{\pgfqpoint{4.386370in}{10.543071in}}%
\pgfpathlineto{\pgfqpoint{4.436674in}{10.543071in}}%
\pgfpathlineto{\pgfqpoint{4.471396in}{10.395893in}}%
\pgfpathlineto{\pgfqpoint{4.506422in}{10.543071in}}%
\pgfpathlineto{\pgfqpoint{4.556031in}{10.543071in}}%
\pgfpathlineto{\pgfqpoint{4.508288in}{10.340554in}}%
\pgfpathlineto{\pgfqpoint{4.448089in}{10.340554in}}%
\pgfpathlineto{\pgfqpoint{4.411326in}{10.494504in}}%
\pgfpathlineto{\pgfqpoint{4.374998in}{10.340554in}}%
\pgfpathlineto{\pgfqpoint{4.314755in}{10.340554in}}%
\pgfpathlineto{\pgfqpoint{4.266622in}{10.543071in}}%
\pgfpathlineto{\pgfqpoint{4.266622in}{10.543071in}}%
\pgfpathclose%
\pgfpathmoveto{\pgfqpoint{4.590233in}{10.543071in}}%
\pgfpathlineto{\pgfqpoint{4.731118in}{10.543071in}}%
\pgfpathlineto{\pgfqpoint{4.731118in}{10.503575in}}%
\pgfpathlineto{\pgfqpoint{4.642446in}{10.503575in}}%
\pgfpathlineto{\pgfqpoint{4.642446in}{10.465901in}}%
\pgfpathlineto{\pgfqpoint{4.725866in}{10.465901in}}%
\pgfpathlineto{\pgfqpoint{4.725866in}{10.426405in}}%
\pgfpathlineto{\pgfqpoint{4.642446in}{10.426405in}}%
\pgfpathlineto{\pgfqpoint{4.642446in}{10.340554in}}%
\pgfpathlineto{\pgfqpoint{4.590233in}{10.340554in}}%
\pgfpathlineto{\pgfqpoint{4.590233in}{10.543071in}}%
\pgfpathlineto{\pgfqpoint{4.590233in}{10.543071in}}%
\pgfpathclose%
\pgfpathmoveto{\pgfqpoint{4.779990in}{10.543071in}}%
\pgfpathlineto{\pgfqpoint{4.920875in}{10.543071in}}%
\pgfpathlineto{\pgfqpoint{4.920875in}{10.503575in}}%
\pgfpathlineto{\pgfqpoint{4.832203in}{10.503575in}}%
\pgfpathlineto{\pgfqpoint{4.832203in}{10.465901in}}%
\pgfpathlineto{\pgfqpoint{4.915623in}{10.465901in}}%
\pgfpathlineto{\pgfqpoint{4.915623in}{10.426405in}}%
\pgfpathlineto{\pgfqpoint{4.832203in}{10.426405in}}%
\pgfpathlineto{\pgfqpoint{4.832203in}{10.380007in}}%
\pgfpathlineto{\pgfqpoint{4.923870in}{10.380007in}}%
\pgfpathlineto{\pgfqpoint{4.923870in}{10.340554in}}%
\pgfpathlineto{\pgfqpoint{4.779990in}{10.340554in}}%
\pgfpathlineto{\pgfqpoint{4.779990in}{10.543071in}}%
\pgfpathlineto{\pgfqpoint{4.779990in}{10.543071in}}%
\pgfpathclose%
\pgfusepath{fill}%
\end{pgfscope}%
\begin{pgfscope}%
\pgfsetbuttcap%
\pgfsetroundjoin%
\definecolor{currentfill}{rgb}{0.309804,0.372549,0.419608}%
\pgfsetfillcolor{currentfill}%
\pgfsetlinewidth{0.000000pt}%
\definecolor{currentstroke}{rgb}{0.309804,0.372549,0.419608}%
\pgfsetstrokecolor{currentstroke}%
\pgfsetdash{}{0pt}%
\pgfpathmoveto{\pgfqpoint{0.953314in}{10.153876in}}%
\pgfpathlineto{\pgfqpoint{1.005903in}{10.153876in}}%
\pgfpathlineto{\pgfqpoint{1.005903in}{10.141200in}}%
\pgfpathlineto{\pgfqpoint{0.935196in}{10.141200in}}%
\pgfpathlineto{\pgfqpoint{0.935196in}{10.153876in}}%
\pgfpathquadraticcurveto{\pgfqpoint{0.943766in}{10.162756in}}{\pgfqpoint{0.958566in}{10.177700in}}%
\pgfpathquadraticcurveto{\pgfqpoint{0.973390in}{10.192667in}}{\pgfqpoint{0.977186in}{10.197012in}}%
\pgfpathquadraticcurveto{\pgfqpoint{0.984419in}{10.205128in}}{\pgfqpoint{0.987283in}{10.210762in}}%
\pgfpathquadraticcurveto{\pgfqpoint{0.990172in}{10.216395in}}{\pgfqpoint{0.990172in}{10.221838in}}%
\pgfpathquadraticcurveto{\pgfqpoint{0.990172in}{10.230718in}}{\pgfqpoint{0.983941in}{10.236304in}}%
\pgfpathquadraticcurveto{\pgfqpoint{0.977711in}{10.241914in}}{\pgfqpoint{0.967709in}{10.241914in}}%
\pgfpathquadraticcurveto{\pgfqpoint{0.960619in}{10.241914in}}{\pgfqpoint{0.952741in}{10.239455in}}%
\pgfpathquadraticcurveto{\pgfqpoint{0.944888in}{10.236996in}}{\pgfqpoint{0.935936in}{10.231983in}}%
\pgfpathlineto{\pgfqpoint{0.935936in}{10.247213in}}%
\pgfpathquadraticcurveto{\pgfqpoint{0.945031in}{10.250866in}}{\pgfqpoint{0.952932in}{10.252728in}}%
\pgfpathquadraticcurveto{\pgfqpoint{0.960858in}{10.254590in}}{\pgfqpoint{0.967422in}{10.254590in}}%
\pgfpathquadraticcurveto{\pgfqpoint{0.984729in}{10.254590in}}{\pgfqpoint{0.995018in}{10.245924in}}%
\pgfpathquadraticcurveto{\pgfqpoint{1.005306in}{10.237283in}}{\pgfqpoint{1.005306in}{10.222817in}}%
\pgfpathquadraticcurveto{\pgfqpoint{1.005306in}{10.215942in}}{\pgfqpoint{1.002728in}{10.209783in}}%
\pgfpathquadraticcurveto{\pgfqpoint{1.000174in}{10.203648in}}{\pgfqpoint{0.993371in}{10.195293in}}%
\pgfpathquadraticcurveto{\pgfqpoint{0.991509in}{10.193121in}}{\pgfqpoint{0.981507in}{10.182784in}}%
\pgfpathquadraticcurveto{\pgfqpoint{0.971528in}{10.172448in}}{\pgfqpoint{0.953314in}{10.153876in}}%
\pgfpathlineto{\pgfqpoint{0.953314in}{10.153876in}}%
\pgfpathclose%
\pgfpathmoveto{\pgfqpoint{1.069760in}{10.242654in}}%
\pgfpathquadraticcurveto{\pgfqpoint{1.058134in}{10.242654in}}{\pgfqpoint{1.052262in}{10.231196in}}%
\pgfpathquadraticcurveto{\pgfqpoint{1.046413in}{10.219761in}}{\pgfqpoint{1.046413in}{10.196773in}}%
\pgfpathquadraticcurveto{\pgfqpoint{1.046413in}{10.173880in}}{\pgfqpoint{1.052262in}{10.162422in}}%
\pgfpathquadraticcurveto{\pgfqpoint{1.058134in}{10.150963in}}{\pgfqpoint{1.069760in}{10.150963in}}%
\pgfpathquadraticcurveto{\pgfqpoint{1.081480in}{10.150963in}}{\pgfqpoint{1.087329in}{10.162422in}}%
\pgfpathquadraticcurveto{\pgfqpoint{1.093201in}{10.173880in}}{\pgfqpoint{1.093201in}{10.196773in}}%
\pgfpathquadraticcurveto{\pgfqpoint{1.093201in}{10.219761in}}{\pgfqpoint{1.087329in}{10.231196in}}%
\pgfpathquadraticcurveto{\pgfqpoint{1.081480in}{10.242654in}}{\pgfqpoint{1.069760in}{10.242654in}}%
\pgfpathlineto{\pgfqpoint{1.069760in}{10.242654in}}%
\pgfpathclose%
\pgfpathmoveto{\pgfqpoint{1.069760in}{10.254590in}}%
\pgfpathquadraticcurveto{\pgfqpoint{1.088499in}{10.254590in}}{\pgfqpoint{1.098382in}{10.239766in}}%
\pgfpathquadraticcurveto{\pgfqpoint{1.108264in}{10.224965in}}{\pgfqpoint{1.108264in}{10.196773in}}%
\pgfpathquadraticcurveto{\pgfqpoint{1.108264in}{10.168652in}}{\pgfqpoint{1.098382in}{10.153828in}}%
\pgfpathquadraticcurveto{\pgfqpoint{1.088499in}{10.139028in}}{\pgfqpoint{1.069760in}{10.139028in}}%
\pgfpathquadraticcurveto{\pgfqpoint{1.051044in}{10.139028in}}{\pgfqpoint{1.041161in}{10.153828in}}%
\pgfpathquadraticcurveto{\pgfqpoint{1.031279in}{10.168652in}}{\pgfqpoint{1.031279in}{10.196773in}}%
\pgfpathquadraticcurveto{\pgfqpoint{1.031279in}{10.224965in}}{\pgfqpoint{1.041161in}{10.239766in}}%
\pgfpathquadraticcurveto{\pgfqpoint{1.051044in}{10.254590in}}{\pgfqpoint{1.069760in}{10.254590in}}%
\pgfpathlineto{\pgfqpoint{1.069760in}{10.254590in}}%
\pgfpathclose%
\pgfpathmoveto{\pgfqpoint{1.234616in}{10.222292in}}%
\pgfpathlineto{\pgfqpoint{1.234616in}{10.209305in}}%
\pgfpathquadraticcurveto{\pgfqpoint{1.228816in}{10.212289in}}{\pgfqpoint{1.222537in}{10.213769in}}%
\pgfpathquadraticcurveto{\pgfqpoint{1.216283in}{10.215273in}}{\pgfqpoint{1.209551in}{10.215273in}}%
\pgfpathquadraticcurveto{\pgfqpoint{1.199334in}{10.215273in}}{\pgfqpoint{1.194226in}{10.212146in}}%
\pgfpathquadraticcurveto{\pgfqpoint{1.189117in}{10.209019in}}{\pgfqpoint{1.189117in}{10.202741in}}%
\pgfpathquadraticcurveto{\pgfqpoint{1.189117in}{10.197966in}}{\pgfqpoint{1.192770in}{10.195245in}}%
\pgfpathquadraticcurveto{\pgfqpoint{1.196422in}{10.192524in}}{\pgfqpoint{1.207474in}{10.190065in}}%
\pgfpathlineto{\pgfqpoint{1.212177in}{10.189015in}}%
\pgfpathquadraticcurveto{\pgfqpoint{1.226786in}{10.185887in}}{\pgfqpoint{1.232945in}{10.180182in}}%
\pgfpathquadraticcurveto{\pgfqpoint{1.239104in}{10.174477in}}{\pgfqpoint{1.239104in}{10.164260in}}%
\pgfpathquadraticcurveto{\pgfqpoint{1.239104in}{10.152611in}}{\pgfqpoint{1.229890in}{10.145807in}}%
\pgfpathquadraticcurveto{\pgfqpoint{1.220675in}{10.139028in}}{\pgfqpoint{1.204562in}{10.139028in}}%
\pgfpathquadraticcurveto{\pgfqpoint{1.197854in}{10.139028in}}{\pgfqpoint{1.190573in}{10.140341in}}%
\pgfpathquadraticcurveto{\pgfqpoint{1.183293in}{10.141654in}}{\pgfqpoint{1.175248in}{10.144256in}}%
\pgfpathlineto{\pgfqpoint{1.175248in}{10.158435in}}%
\pgfpathquadraticcurveto{\pgfqpoint{1.182863in}{10.154473in}}{\pgfqpoint{1.190239in}{10.152491in}}%
\pgfpathquadraticcurveto{\pgfqpoint{1.197615in}{10.150534in}}{\pgfqpoint{1.204872in}{10.150534in}}%
\pgfpathquadraticcurveto{\pgfqpoint{1.214564in}{10.150534in}}{\pgfqpoint{1.219768in}{10.153852in}}%
\pgfpathquadraticcurveto{\pgfqpoint{1.224996in}{10.157170in}}{\pgfqpoint{1.224996in}{10.163210in}}%
\pgfpathquadraticcurveto{\pgfqpoint{1.224996in}{10.168795in}}{\pgfqpoint{1.221224in}{10.171779in}}%
\pgfpathquadraticcurveto{\pgfqpoint{1.217477in}{10.174763in}}{\pgfqpoint{1.204705in}{10.177532in}}%
\pgfpathlineto{\pgfqpoint{1.199931in}{10.178654in}}%
\pgfpathquadraticcurveto{\pgfqpoint{1.187184in}{10.181328in}}{\pgfqpoint{1.181502in}{10.186890in}}%
\pgfpathquadraticcurveto{\pgfqpoint{1.175845in}{10.192452in}}{\pgfqpoint{1.175845in}{10.202144in}}%
\pgfpathquadraticcurveto{\pgfqpoint{1.175845in}{10.213937in}}{\pgfqpoint{1.184200in}{10.220334in}}%
\pgfpathquadraticcurveto{\pgfqpoint{1.192555in}{10.226756in}}{\pgfqpoint{1.207928in}{10.226756in}}%
\pgfpathquadraticcurveto{\pgfqpoint{1.215519in}{10.226756in}}{\pgfqpoint{1.222227in}{10.225634in}}%
\pgfpathquadraticcurveto{\pgfqpoint{1.228959in}{10.224536in}}{\pgfqpoint{1.234616in}{10.222292in}}%
\pgfpathlineto{\pgfqpoint{1.234616in}{10.222292in}}%
\pgfpathclose%
\pgfpathmoveto{\pgfqpoint{1.274530in}{10.248479in}}%
\pgfpathlineto{\pgfqpoint{1.274530in}{10.224750in}}%
\pgfpathlineto{\pgfqpoint{1.302793in}{10.224750in}}%
\pgfpathlineto{\pgfqpoint{1.302793in}{10.214080in}}%
\pgfpathlineto{\pgfqpoint{1.274530in}{10.214080in}}%
\pgfpathlineto{\pgfqpoint{1.274530in}{10.168724in}}%
\pgfpathquadraticcurveto{\pgfqpoint{1.274530in}{10.158507in}}{\pgfqpoint{1.277322in}{10.155595in}}%
\pgfpathquadraticcurveto{\pgfqpoint{1.280115in}{10.152682in}}{\pgfqpoint{1.288709in}{10.152682in}}%
\pgfpathlineto{\pgfqpoint{1.302793in}{10.152682in}}%
\pgfpathlineto{\pgfqpoint{1.302793in}{10.141200in}}%
\pgfpathlineto{\pgfqpoint{1.288709in}{10.141200in}}%
\pgfpathquadraticcurveto{\pgfqpoint{1.272811in}{10.141200in}}{\pgfqpoint{1.266771in}{10.147120in}}%
\pgfpathquadraticcurveto{\pgfqpoint{1.260732in}{10.153064in}}{\pgfqpoint{1.260732in}{10.168724in}}%
\pgfpathlineto{\pgfqpoint{1.260732in}{10.214080in}}%
\pgfpathlineto{\pgfqpoint{1.250658in}{10.214080in}}%
\pgfpathlineto{\pgfqpoint{1.250658in}{10.224750in}}%
\pgfpathlineto{\pgfqpoint{1.260732in}{10.224750in}}%
\pgfpathlineto{\pgfqpoint{1.260732in}{10.248479in}}%
\pgfpathlineto{\pgfqpoint{1.274530in}{10.248479in}}%
\pgfpathlineto{\pgfqpoint{1.274530in}{10.248479in}}%
\pgfpathclose%
\pgfpathmoveto{\pgfqpoint{1.319432in}{10.174167in}}%
\pgfpathlineto{\pgfqpoint{1.319432in}{10.224750in}}%
\pgfpathlineto{\pgfqpoint{1.333158in}{10.224750in}}%
\pgfpathlineto{\pgfqpoint{1.333158in}{10.174692in}}%
\pgfpathquadraticcurveto{\pgfqpoint{1.333158in}{10.162828in}}{\pgfqpoint{1.337765in}{10.156884in}}%
\pgfpathquadraticcurveto{\pgfqpoint{1.342396in}{10.150963in}}{\pgfqpoint{1.351658in}{10.150963in}}%
\pgfpathquadraticcurveto{\pgfqpoint{1.362759in}{10.150963in}}{\pgfqpoint{1.369204in}{10.158053in}}%
\pgfpathquadraticcurveto{\pgfqpoint{1.375673in}{10.165143in}}{\pgfqpoint{1.375673in}{10.177389in}}%
\pgfpathlineto{\pgfqpoint{1.375673in}{10.224750in}}%
\pgfpathlineto{\pgfqpoint{1.389399in}{10.224750in}}%
\pgfpathlineto{\pgfqpoint{1.389399in}{10.141200in}}%
\pgfpathlineto{\pgfqpoint{1.375673in}{10.141200in}}%
\pgfpathlineto{\pgfqpoint{1.375673in}{10.154043in}}%
\pgfpathquadraticcurveto{\pgfqpoint{1.370684in}{10.146428in}}{\pgfqpoint{1.364072in}{10.142728in}}%
\pgfpathquadraticcurveto{\pgfqpoint{1.357483in}{10.139028in}}{\pgfqpoint{1.348746in}{10.139028in}}%
\pgfpathquadraticcurveto{\pgfqpoint{1.334352in}{10.139028in}}{\pgfqpoint{1.326880in}{10.147980in}}%
\pgfpathquadraticcurveto{\pgfqpoint{1.319432in}{10.156931in}}{\pgfqpoint{1.319432in}{10.174167in}}%
\pgfpathlineto{\pgfqpoint{1.319432in}{10.174167in}}%
\pgfpathclose%
\pgfpathmoveto{\pgfqpoint{1.353974in}{10.226756in}}%
\pgfpathlineto{\pgfqpoint{1.353974in}{10.226756in}}%
\pgfpathlineto{\pgfqpoint{1.353974in}{10.226756in}}%
\pgfpathclose%
\pgfpathmoveto{\pgfqpoint{1.472639in}{10.212075in}}%
\pgfpathlineto{\pgfqpoint{1.472639in}{10.257287in}}%
\pgfpathlineto{\pgfqpoint{1.486365in}{10.257287in}}%
\pgfpathlineto{\pgfqpoint{1.486365in}{10.141200in}}%
\pgfpathlineto{\pgfqpoint{1.472639in}{10.141200in}}%
\pgfpathlineto{\pgfqpoint{1.472639in}{10.153733in}}%
\pgfpathquadraticcurveto{\pgfqpoint{1.468319in}{10.146285in}}{\pgfqpoint{1.461706in}{10.142656in}}%
\pgfpathquadraticcurveto{\pgfqpoint{1.455118in}{10.139028in}}{\pgfqpoint{1.445855in}{10.139028in}}%
\pgfpathquadraticcurveto{\pgfqpoint{1.430721in}{10.139028in}}{\pgfqpoint{1.421196in}{10.151107in}}%
\pgfpathquadraticcurveto{\pgfqpoint{1.411695in}{10.163210in}}{\pgfqpoint{1.411695in}{10.182904in}}%
\pgfpathquadraticcurveto{\pgfqpoint{1.411695in}{10.202598in}}{\pgfqpoint{1.421196in}{10.214677in}}%
\pgfpathquadraticcurveto{\pgfqpoint{1.430721in}{10.226756in}}{\pgfqpoint{1.445855in}{10.226756in}}%
\pgfpathquadraticcurveto{\pgfqpoint{1.455118in}{10.226756in}}{\pgfqpoint{1.461706in}{10.223127in}}%
\pgfpathquadraticcurveto{\pgfqpoint{1.468319in}{10.219522in}}{\pgfqpoint{1.472639in}{10.212075in}}%
\pgfpathlineto{\pgfqpoint{1.472639in}{10.212075in}}%
\pgfpathclose%
\pgfpathmoveto{\pgfqpoint{1.425875in}{10.182904in}}%
\pgfpathquadraticcurveto{\pgfqpoint{1.425875in}{10.167769in}}{\pgfqpoint{1.432105in}{10.159151in}}%
\pgfpathquadraticcurveto{\pgfqpoint{1.438336in}{10.150534in}}{\pgfqpoint{1.449221in}{10.150534in}}%
\pgfpathquadraticcurveto{\pgfqpoint{1.460107in}{10.150534in}}{\pgfqpoint{1.466361in}{10.159151in}}%
\pgfpathquadraticcurveto{\pgfqpoint{1.472639in}{10.167769in}}{\pgfqpoint{1.472639in}{10.182904in}}%
\pgfpathquadraticcurveto{\pgfqpoint{1.472639in}{10.198038in}}{\pgfqpoint{1.466361in}{10.206656in}}%
\pgfpathquadraticcurveto{\pgfqpoint{1.460107in}{10.215273in}}{\pgfqpoint{1.449221in}{10.215273in}}%
\pgfpathquadraticcurveto{\pgfqpoint{1.438336in}{10.215273in}}{\pgfqpoint{1.432105in}{10.206656in}}%
\pgfpathquadraticcurveto{\pgfqpoint{1.425875in}{10.198038in}}{\pgfqpoint{1.425875in}{10.182904in}}%
\pgfpathlineto{\pgfqpoint{1.425875in}{10.182904in}}%
\pgfpathclose%
\pgfpathmoveto{\pgfqpoint{1.514653in}{10.224750in}}%
\pgfpathlineto{\pgfqpoint{1.528379in}{10.224750in}}%
\pgfpathlineto{\pgfqpoint{1.528379in}{10.141200in}}%
\pgfpathlineto{\pgfqpoint{1.514653in}{10.141200in}}%
\pgfpathlineto{\pgfqpoint{1.514653in}{10.224750in}}%
\pgfpathlineto{\pgfqpoint{1.514653in}{10.224750in}}%
\pgfpathclose%
\pgfpathmoveto{\pgfqpoint{1.514653in}{10.257287in}}%
\pgfpathlineto{\pgfqpoint{1.528379in}{10.257287in}}%
\pgfpathlineto{\pgfqpoint{1.528379in}{10.239885in}}%
\pgfpathlineto{\pgfqpoint{1.514653in}{10.239885in}}%
\pgfpathlineto{\pgfqpoint{1.514653in}{10.257287in}}%
\pgfpathlineto{\pgfqpoint{1.514653in}{10.257287in}}%
\pgfpathclose%
\pgfpathmoveto{\pgfqpoint{1.628568in}{10.186413in}}%
\pgfpathlineto{\pgfqpoint{1.628568in}{10.179705in}}%
\pgfpathlineto{\pgfqpoint{1.565452in}{10.179705in}}%
\pgfpathquadraticcurveto{\pgfqpoint{1.566359in}{10.165525in}}{\pgfqpoint{1.573998in}{10.158101in}}%
\pgfpathquadraticcurveto{\pgfqpoint{1.581637in}{10.150677in}}{\pgfqpoint{1.595291in}{10.150677in}}%
\pgfpathquadraticcurveto{\pgfqpoint{1.603193in}{10.150677in}}{\pgfqpoint{1.610617in}{10.152611in}}%
\pgfpathquadraticcurveto{\pgfqpoint{1.618041in}{10.154544in}}{\pgfqpoint{1.625369in}{10.158435in}}%
\pgfpathlineto{\pgfqpoint{1.625369in}{10.145449in}}%
\pgfpathquadraticcurveto{\pgfqpoint{1.617969in}{10.142322in}}{\pgfqpoint{1.610211in}{10.140675in}}%
\pgfpathquadraticcurveto{\pgfqpoint{1.602453in}{10.139028in}}{\pgfqpoint{1.594480in}{10.139028in}}%
\pgfpathquadraticcurveto{\pgfqpoint{1.574475in}{10.139028in}}{\pgfqpoint{1.562802in}{10.150653in}}%
\pgfpathquadraticcurveto{\pgfqpoint{1.551129in}{10.162302in}}{\pgfqpoint{1.551129in}{10.182164in}}%
\pgfpathquadraticcurveto{\pgfqpoint{1.551129in}{10.202669in}}{\pgfqpoint{1.562205in}{10.214700in}}%
\pgfpathquadraticcurveto{\pgfqpoint{1.573282in}{10.226756in}}{\pgfqpoint{1.592092in}{10.226756in}}%
\pgfpathquadraticcurveto{\pgfqpoint{1.608946in}{10.226756in}}{\pgfqpoint{1.618757in}{10.215894in}}%
\pgfpathquadraticcurveto{\pgfqpoint{1.628568in}{10.205056in}}{\pgfqpoint{1.628568in}{10.186413in}}%
\pgfpathlineto{\pgfqpoint{1.628568in}{10.186413in}}%
\pgfpathclose%
\pgfpathmoveto{\pgfqpoint{1.614842in}{10.190447in}}%
\pgfpathquadraticcurveto{\pgfqpoint{1.614699in}{10.201690in}}{\pgfqpoint{1.608540in}{10.208398in}}%
\pgfpathquadraticcurveto{\pgfqpoint{1.602381in}{10.215130in}}{\pgfqpoint{1.592236in}{10.215130in}}%
\pgfpathquadraticcurveto{\pgfqpoint{1.580753in}{10.215130in}}{\pgfqpoint{1.573855in}{10.208637in}}%
\pgfpathquadraticcurveto{\pgfqpoint{1.566956in}{10.202144in}}{\pgfqpoint{1.565905in}{10.190351in}}%
\pgfpathlineto{\pgfqpoint{1.614842in}{10.190447in}}%
\pgfpathlineto{\pgfqpoint{1.614842in}{10.190447in}}%
\pgfpathclose%
\pgfpathmoveto{\pgfqpoint{1.704360in}{10.222292in}}%
\pgfpathlineto{\pgfqpoint{1.704360in}{10.209305in}}%
\pgfpathquadraticcurveto{\pgfqpoint{1.698559in}{10.212289in}}{\pgfqpoint{1.692281in}{10.213769in}}%
\pgfpathquadraticcurveto{\pgfqpoint{1.686027in}{10.215273in}}{\pgfqpoint{1.679295in}{10.215273in}}%
\pgfpathquadraticcurveto{\pgfqpoint{1.669078in}{10.215273in}}{\pgfqpoint{1.663970in}{10.212146in}}%
\pgfpathquadraticcurveto{\pgfqpoint{1.658861in}{10.209019in}}{\pgfqpoint{1.658861in}{10.202741in}}%
\pgfpathquadraticcurveto{\pgfqpoint{1.658861in}{10.197966in}}{\pgfqpoint{1.662513in}{10.195245in}}%
\pgfpathquadraticcurveto{\pgfqpoint{1.666166in}{10.192524in}}{\pgfqpoint{1.677218in}{10.190065in}}%
\pgfpathlineto{\pgfqpoint{1.681921in}{10.189015in}}%
\pgfpathquadraticcurveto{\pgfqpoint{1.696530in}{10.185887in}}{\pgfqpoint{1.702689in}{10.180182in}}%
\pgfpathquadraticcurveto{\pgfqpoint{1.708848in}{10.174477in}}{\pgfqpoint{1.708848in}{10.164260in}}%
\pgfpathquadraticcurveto{\pgfqpoint{1.708848in}{10.152611in}}{\pgfqpoint{1.699634in}{10.145807in}}%
\pgfpathquadraticcurveto{\pgfqpoint{1.690419in}{10.139028in}}{\pgfqpoint{1.674306in}{10.139028in}}%
\pgfpathquadraticcurveto{\pgfqpoint{1.667598in}{10.139028in}}{\pgfqpoint{1.660317in}{10.140341in}}%
\pgfpathquadraticcurveto{\pgfqpoint{1.653036in}{10.141654in}}{\pgfqpoint{1.644992in}{10.144256in}}%
\pgfpathlineto{\pgfqpoint{1.644992in}{10.158435in}}%
\pgfpathquadraticcurveto{\pgfqpoint{1.652607in}{10.154473in}}{\pgfqpoint{1.659983in}{10.152491in}}%
\pgfpathquadraticcurveto{\pgfqpoint{1.667359in}{10.150534in}}{\pgfqpoint{1.674616in}{10.150534in}}%
\pgfpathquadraticcurveto{\pgfqpoint{1.684308in}{10.150534in}}{\pgfqpoint{1.689512in}{10.153852in}}%
\pgfpathquadraticcurveto{\pgfqpoint{1.694740in}{10.157170in}}{\pgfqpoint{1.694740in}{10.163210in}}%
\pgfpathquadraticcurveto{\pgfqpoint{1.694740in}{10.168795in}}{\pgfqpoint{1.690968in}{10.171779in}}%
\pgfpathquadraticcurveto{\pgfqpoint{1.687220in}{10.174763in}}{\pgfqpoint{1.674449in}{10.177532in}}%
\pgfpathlineto{\pgfqpoint{1.669675in}{10.178654in}}%
\pgfpathquadraticcurveto{\pgfqpoint{1.656928in}{10.181328in}}{\pgfqpoint{1.651246in}{10.186890in}}%
\pgfpathquadraticcurveto{\pgfqpoint{1.645589in}{10.192452in}}{\pgfqpoint{1.645589in}{10.202144in}}%
\pgfpathquadraticcurveto{\pgfqpoint{1.645589in}{10.213937in}}{\pgfqpoint{1.653944in}{10.220334in}}%
\pgfpathquadraticcurveto{\pgfqpoint{1.662299in}{10.226756in}}{\pgfqpoint{1.677672in}{10.226756in}}%
\pgfpathquadraticcurveto{\pgfqpoint{1.685263in}{10.226756in}}{\pgfqpoint{1.691971in}{10.225634in}}%
\pgfpathquadraticcurveto{\pgfqpoint{1.698703in}{10.224536in}}{\pgfqpoint{1.704360in}{10.222292in}}%
\pgfpathlineto{\pgfqpoint{1.704360in}{10.222292in}}%
\pgfpathclose%
\pgfpathmoveto{\pgfqpoint{1.781179in}{10.203720in}}%
\pgfpathlineto{\pgfqpoint{1.796934in}{10.203720in}}%
\pgfpathlineto{\pgfqpoint{1.796934in}{10.184766in}}%
\pgfpathlineto{\pgfqpoint{1.781179in}{10.184766in}}%
\pgfpathlineto{\pgfqpoint{1.781179in}{10.203720in}}%
\pgfpathlineto{\pgfqpoint{1.781179in}{10.203720in}}%
\pgfpathclose%
\pgfpathmoveto{\pgfqpoint{1.876951in}{10.252585in}}%
\pgfpathlineto{\pgfqpoint{1.940951in}{10.252585in}}%
\pgfpathlineto{\pgfqpoint{1.940951in}{10.239885in}}%
\pgfpathlineto{\pgfqpoint{1.892014in}{10.239885in}}%
\pgfpathlineto{\pgfqpoint{1.892014in}{10.207062in}}%
\pgfpathlineto{\pgfqpoint{1.936177in}{10.207062in}}%
\pgfpathlineto{\pgfqpoint{1.936177in}{10.194386in}}%
\pgfpathlineto{\pgfqpoint{1.892014in}{10.194386in}}%
\pgfpathlineto{\pgfqpoint{1.892014in}{10.141200in}}%
\pgfpathlineto{\pgfqpoint{1.876951in}{10.141200in}}%
\pgfpathlineto{\pgfqpoint{1.876951in}{10.252585in}}%
\pgfpathlineto{\pgfqpoint{1.876951in}{10.252585in}}%
\pgfpathclose%
\pgfpathmoveto{\pgfqpoint{1.953102in}{10.224750in}}%
\pgfpathlineto{\pgfqpoint{1.966828in}{10.224750in}}%
\pgfpathlineto{\pgfqpoint{1.966828in}{10.141200in}}%
\pgfpathlineto{\pgfqpoint{1.953102in}{10.141200in}}%
\pgfpathlineto{\pgfqpoint{1.953102in}{10.224750in}}%
\pgfpathlineto{\pgfqpoint{1.953102in}{10.224750in}}%
\pgfpathclose%
\pgfpathmoveto{\pgfqpoint{1.953102in}{10.257287in}}%
\pgfpathlineto{\pgfqpoint{1.966828in}{10.257287in}}%
\pgfpathlineto{\pgfqpoint{1.966828in}{10.239885in}}%
\pgfpathlineto{\pgfqpoint{1.953102in}{10.239885in}}%
\pgfpathlineto{\pgfqpoint{1.953102in}{10.257287in}}%
\pgfpathlineto{\pgfqpoint{1.953102in}{10.257287in}}%
\pgfpathclose%
\pgfpathmoveto{\pgfqpoint{1.995545in}{10.257287in}}%
\pgfpathlineto{\pgfqpoint{2.009271in}{10.257287in}}%
\pgfpathlineto{\pgfqpoint{2.009271in}{10.141200in}}%
\pgfpathlineto{\pgfqpoint{1.995545in}{10.141200in}}%
\pgfpathlineto{\pgfqpoint{1.995545in}{10.257287in}}%
\pgfpathlineto{\pgfqpoint{1.995545in}{10.257287in}}%
\pgfpathclose%
\pgfpathmoveto{\pgfqpoint{2.037989in}{10.257287in}}%
\pgfpathlineto{\pgfqpoint{2.051715in}{10.257287in}}%
\pgfpathlineto{\pgfqpoint{2.051715in}{10.141200in}}%
\pgfpathlineto{\pgfqpoint{2.037989in}{10.141200in}}%
\pgfpathlineto{\pgfqpoint{2.037989in}{10.257287in}}%
\pgfpathlineto{\pgfqpoint{2.037989in}{10.257287in}}%
\pgfpathclose%
\pgfpathmoveto{\pgfqpoint{2.151904in}{10.186413in}}%
\pgfpathlineto{\pgfqpoint{2.151904in}{10.179705in}}%
\pgfpathlineto{\pgfqpoint{2.088787in}{10.179705in}}%
\pgfpathquadraticcurveto{\pgfqpoint{2.089694in}{10.165525in}}{\pgfqpoint{2.097333in}{10.158101in}}%
\pgfpathquadraticcurveto{\pgfqpoint{2.104972in}{10.150677in}}{\pgfqpoint{2.118627in}{10.150677in}}%
\pgfpathquadraticcurveto{\pgfqpoint{2.126528in}{10.150677in}}{\pgfqpoint{2.133952in}{10.152611in}}%
\pgfpathquadraticcurveto{\pgfqpoint{2.141376in}{10.154544in}}{\pgfqpoint{2.148705in}{10.158435in}}%
\pgfpathlineto{\pgfqpoint{2.148705in}{10.145449in}}%
\pgfpathquadraticcurveto{\pgfqpoint{2.141305in}{10.142322in}}{\pgfqpoint{2.133546in}{10.140675in}}%
\pgfpathquadraticcurveto{\pgfqpoint{2.125788in}{10.139028in}}{\pgfqpoint{2.117815in}{10.139028in}}%
\pgfpathquadraticcurveto{\pgfqpoint{2.097811in}{10.139028in}}{\pgfqpoint{2.086138in}{10.150653in}}%
\pgfpathquadraticcurveto{\pgfqpoint{2.074464in}{10.162302in}}{\pgfqpoint{2.074464in}{10.182164in}}%
\pgfpathquadraticcurveto{\pgfqpoint{2.074464in}{10.202669in}}{\pgfqpoint{2.085541in}{10.214700in}}%
\pgfpathquadraticcurveto{\pgfqpoint{2.096617in}{10.226756in}}{\pgfqpoint{2.115428in}{10.226756in}}%
\pgfpathquadraticcurveto{\pgfqpoint{2.132281in}{10.226756in}}{\pgfqpoint{2.142092in}{10.215894in}}%
\pgfpathquadraticcurveto{\pgfqpoint{2.151904in}{10.205056in}}{\pgfqpoint{2.151904in}{10.186413in}}%
\pgfpathlineto{\pgfqpoint{2.151904in}{10.186413in}}%
\pgfpathclose%
\pgfpathmoveto{\pgfqpoint{2.138178in}{10.190447in}}%
\pgfpathquadraticcurveto{\pgfqpoint{2.138034in}{10.201690in}}{\pgfqpoint{2.131875in}{10.208398in}}%
\pgfpathquadraticcurveto{\pgfqpoint{2.125717in}{10.215130in}}{\pgfqpoint{2.115571in}{10.215130in}}%
\pgfpathquadraticcurveto{\pgfqpoint{2.104089in}{10.215130in}}{\pgfqpoint{2.097190in}{10.208637in}}%
\pgfpathquadraticcurveto{\pgfqpoint{2.090291in}{10.202144in}}{\pgfqpoint{2.089241in}{10.190351in}}%
\pgfpathlineto{\pgfqpoint{2.138178in}{10.190447in}}%
\pgfpathlineto{\pgfqpoint{2.138178in}{10.190447in}}%
\pgfpathclose%
\pgfpathmoveto{\pgfqpoint{2.229414in}{10.212075in}}%
\pgfpathlineto{\pgfqpoint{2.229414in}{10.257287in}}%
\pgfpathlineto{\pgfqpoint{2.243141in}{10.257287in}}%
\pgfpathlineto{\pgfqpoint{2.243141in}{10.141200in}}%
\pgfpathlineto{\pgfqpoint{2.229414in}{10.141200in}}%
\pgfpathlineto{\pgfqpoint{2.229414in}{10.153733in}}%
\pgfpathquadraticcurveto{\pgfqpoint{2.225094in}{10.146285in}}{\pgfqpoint{2.218481in}{10.142656in}}%
\pgfpathquadraticcurveto{\pgfqpoint{2.211893in}{10.139028in}}{\pgfqpoint{2.202631in}{10.139028in}}%
\pgfpathquadraticcurveto{\pgfqpoint{2.187496in}{10.139028in}}{\pgfqpoint{2.177971in}{10.151107in}}%
\pgfpathquadraticcurveto{\pgfqpoint{2.168470in}{10.163210in}}{\pgfqpoint{2.168470in}{10.182904in}}%
\pgfpathquadraticcurveto{\pgfqpoint{2.168470in}{10.202598in}}{\pgfqpoint{2.177971in}{10.214677in}}%
\pgfpathquadraticcurveto{\pgfqpoint{2.187496in}{10.226756in}}{\pgfqpoint{2.202631in}{10.226756in}}%
\pgfpathquadraticcurveto{\pgfqpoint{2.211893in}{10.226756in}}{\pgfqpoint{2.218481in}{10.223127in}}%
\pgfpathquadraticcurveto{\pgfqpoint{2.225094in}{10.219522in}}{\pgfqpoint{2.229414in}{10.212075in}}%
\pgfpathlineto{\pgfqpoint{2.229414in}{10.212075in}}%
\pgfpathclose%
\pgfpathmoveto{\pgfqpoint{2.182650in}{10.182904in}}%
\pgfpathquadraticcurveto{\pgfqpoint{2.182650in}{10.167769in}}{\pgfqpoint{2.188881in}{10.159151in}}%
\pgfpathquadraticcurveto{\pgfqpoint{2.195111in}{10.150534in}}{\pgfqpoint{2.205997in}{10.150534in}}%
\pgfpathquadraticcurveto{\pgfqpoint{2.216882in}{10.150534in}}{\pgfqpoint{2.223136in}{10.159151in}}%
\pgfpathquadraticcurveto{\pgfqpoint{2.229414in}{10.167769in}}{\pgfqpoint{2.229414in}{10.182904in}}%
\pgfpathquadraticcurveto{\pgfqpoint{2.229414in}{10.198038in}}{\pgfqpoint{2.223136in}{10.206656in}}%
\pgfpathquadraticcurveto{\pgfqpoint{2.216882in}{10.215273in}}{\pgfqpoint{2.205997in}{10.215273in}}%
\pgfpathquadraticcurveto{\pgfqpoint{2.195111in}{10.215273in}}{\pgfqpoint{2.188881in}{10.206656in}}%
\pgfpathquadraticcurveto{\pgfqpoint{2.182650in}{10.198038in}}{\pgfqpoint{2.182650in}{10.182904in}}%
\pgfpathlineto{\pgfqpoint{2.182650in}{10.182904in}}%
\pgfpathclose%
\pgfpathmoveto{\pgfqpoint{2.274938in}{10.160154in}}%
\pgfpathlineto{\pgfqpoint{2.290669in}{10.160154in}}%
\pgfpathlineto{\pgfqpoint{2.290669in}{10.141200in}}%
\pgfpathlineto{\pgfqpoint{2.274938in}{10.141200in}}%
\pgfpathlineto{\pgfqpoint{2.274938in}{10.160154in}}%
\pgfpathlineto{\pgfqpoint{2.274938in}{10.160154in}}%
\pgfpathclose%
\pgfpathmoveto{\pgfqpoint{2.274938in}{10.220191in}}%
\pgfpathlineto{\pgfqpoint{2.290669in}{10.220191in}}%
\pgfpathlineto{\pgfqpoint{2.290669in}{10.201261in}}%
\pgfpathlineto{\pgfqpoint{2.274938in}{10.201261in}}%
\pgfpathlineto{\pgfqpoint{2.274938in}{10.220191in}}%
\pgfpathlineto{\pgfqpoint{2.274938in}{10.220191in}}%
\pgfpathclose%
\pgfpathmoveto{\pgfqpoint{2.384723in}{10.153733in}}%
\pgfpathlineto{\pgfqpoint{2.384723in}{10.109427in}}%
\pgfpathlineto{\pgfqpoint{2.370925in}{10.109427in}}%
\pgfpathlineto{\pgfqpoint{2.370925in}{10.224750in}}%
\pgfpathlineto{\pgfqpoint{2.384723in}{10.224750in}}%
\pgfpathlineto{\pgfqpoint{2.384723in}{10.212075in}}%
\pgfpathquadraticcurveto{\pgfqpoint{2.389067in}{10.219522in}}{\pgfqpoint{2.395656in}{10.223127in}}%
\pgfpathquadraticcurveto{\pgfqpoint{2.402268in}{10.226756in}}{\pgfqpoint{2.411435in}{10.226756in}}%
\pgfpathquadraticcurveto{\pgfqpoint{2.426665in}{10.226756in}}{\pgfqpoint{2.436166in}{10.214677in}}%
\pgfpathquadraticcurveto{\pgfqpoint{2.445691in}{10.202598in}}{\pgfqpoint{2.445691in}{10.182904in}}%
\pgfpathquadraticcurveto{\pgfqpoint{2.445691in}{10.163210in}}{\pgfqpoint{2.436166in}{10.151107in}}%
\pgfpathquadraticcurveto{\pgfqpoint{2.426665in}{10.139028in}}{\pgfqpoint{2.411435in}{10.139028in}}%
\pgfpathquadraticcurveto{\pgfqpoint{2.402268in}{10.139028in}}{\pgfqpoint{2.395656in}{10.142656in}}%
\pgfpathquadraticcurveto{\pgfqpoint{2.389067in}{10.146285in}}{\pgfqpoint{2.384723in}{10.153733in}}%
\pgfpathlineto{\pgfqpoint{2.384723in}{10.153733in}}%
\pgfpathclose%
\pgfpathmoveto{\pgfqpoint{2.431439in}{10.182904in}}%
\pgfpathquadraticcurveto{\pgfqpoint{2.431439in}{10.198038in}}{\pgfqpoint{2.425209in}{10.206656in}}%
\pgfpathquadraticcurveto{\pgfqpoint{2.418978in}{10.215273in}}{\pgfqpoint{2.408093in}{10.215273in}}%
\pgfpathquadraticcurveto{\pgfqpoint{2.397184in}{10.215273in}}{\pgfqpoint{2.390953in}{10.206656in}}%
\pgfpathquadraticcurveto{\pgfqpoint{2.384723in}{10.198038in}}{\pgfqpoint{2.384723in}{10.182904in}}%
\pgfpathquadraticcurveto{\pgfqpoint{2.384723in}{10.167769in}}{\pgfqpoint{2.390953in}{10.159151in}}%
\pgfpathquadraticcurveto{\pgfqpoint{2.397184in}{10.150534in}}{\pgfqpoint{2.408093in}{10.150534in}}%
\pgfpathquadraticcurveto{\pgfqpoint{2.418978in}{10.150534in}}{\pgfqpoint{2.425209in}{10.159151in}}%
\pgfpathquadraticcurveto{\pgfqpoint{2.431439in}{10.167769in}}{\pgfqpoint{2.431439in}{10.182904in}}%
\pgfpathlineto{\pgfqpoint{2.431439in}{10.182904in}}%
\pgfpathclose%
\pgfpathmoveto{\pgfqpoint{2.506420in}{10.183190in}}%
\pgfpathquadraticcurveto{\pgfqpoint{2.489781in}{10.183190in}}{\pgfqpoint{2.483360in}{10.179394in}}%
\pgfpathquadraticcurveto{\pgfqpoint{2.476938in}{10.175599in}}{\pgfqpoint{2.476938in}{10.166408in}}%
\pgfpathquadraticcurveto{\pgfqpoint{2.476938in}{10.159104in}}{\pgfqpoint{2.481760in}{10.154807in}}%
\pgfpathquadraticcurveto{\pgfqpoint{2.486582in}{10.150534in}}{\pgfqpoint{2.494842in}{10.150534in}}%
\pgfpathquadraticcurveto{\pgfqpoint{2.506276in}{10.150534in}}{\pgfqpoint{2.513175in}{10.158626in}}%
\pgfpathquadraticcurveto{\pgfqpoint{2.520074in}{10.166719in}}{\pgfqpoint{2.520074in}{10.180134in}}%
\pgfpathlineto{\pgfqpoint{2.520074in}{10.183190in}}%
\pgfpathlineto{\pgfqpoint{2.506420in}{10.183190in}}%
\pgfpathlineto{\pgfqpoint{2.506420in}{10.183190in}}%
\pgfpathclose%
\pgfpathmoveto{\pgfqpoint{2.533800in}{10.188871in}}%
\pgfpathlineto{\pgfqpoint{2.533800in}{10.141200in}}%
\pgfpathlineto{\pgfqpoint{2.520074in}{10.141200in}}%
\pgfpathlineto{\pgfqpoint{2.520074in}{10.153876in}}%
\pgfpathquadraticcurveto{\pgfqpoint{2.515372in}{10.146285in}}{\pgfqpoint{2.508353in}{10.142656in}}%
\pgfpathquadraticcurveto{\pgfqpoint{2.501335in}{10.139028in}}{\pgfqpoint{2.491190in}{10.139028in}}%
\pgfpathquadraticcurveto{\pgfqpoint{2.478371in}{10.139028in}}{\pgfqpoint{2.470780in}{10.146237in}}%
\pgfpathquadraticcurveto{\pgfqpoint{2.463212in}{10.153446in}}{\pgfqpoint{2.463212in}{10.165525in}}%
\pgfpathquadraticcurveto{\pgfqpoint{2.463212in}{10.179609in}}{\pgfqpoint{2.472641in}{10.186771in}}%
\pgfpathquadraticcurveto{\pgfqpoint{2.482095in}{10.193932in}}{\pgfqpoint{2.500810in}{10.193932in}}%
\pgfpathlineto{\pgfqpoint{2.520074in}{10.193932in}}%
\pgfpathlineto{\pgfqpoint{2.520074in}{10.195293in}}%
\pgfpathquadraticcurveto{\pgfqpoint{2.520074in}{10.204770in}}{\pgfqpoint{2.513844in}{10.209950in}}%
\pgfpathquadraticcurveto{\pgfqpoint{2.507613in}{10.215130in}}{\pgfqpoint{2.496346in}{10.215130in}}%
\pgfpathquadraticcurveto{\pgfqpoint{2.489184in}{10.215130in}}{\pgfqpoint{2.482381in}{10.213411in}}%
\pgfpathquadraticcurveto{\pgfqpoint{2.475602in}{10.211693in}}{\pgfqpoint{2.469347in}{10.208255in}}%
\pgfpathlineto{\pgfqpoint{2.469347in}{10.220955in}}%
\pgfpathquadraticcurveto{\pgfqpoint{2.476867in}{10.223867in}}{\pgfqpoint{2.483957in}{10.225299in}}%
\pgfpathquadraticcurveto{\pgfqpoint{2.491046in}{10.226756in}}{\pgfqpoint{2.497754in}{10.226756in}}%
\pgfpathquadraticcurveto{\pgfqpoint{2.515897in}{10.226756in}}{\pgfqpoint{2.524849in}{10.217350in}}%
\pgfpathquadraticcurveto{\pgfqpoint{2.533800in}{10.207969in}}{\pgfqpoint{2.533800in}{10.188871in}}%
\pgfpathlineto{\pgfqpoint{2.533800in}{10.188871in}}%
\pgfpathclose%
\pgfpathmoveto{\pgfqpoint{2.562064in}{10.224750in}}%
\pgfpathlineto{\pgfqpoint{2.575790in}{10.224750in}}%
\pgfpathlineto{\pgfqpoint{2.575790in}{10.141200in}}%
\pgfpathlineto{\pgfqpoint{2.562064in}{10.141200in}}%
\pgfpathlineto{\pgfqpoint{2.562064in}{10.224750in}}%
\pgfpathlineto{\pgfqpoint{2.562064in}{10.224750in}}%
\pgfpathclose%
\pgfpathmoveto{\pgfqpoint{2.562064in}{10.257287in}}%
\pgfpathlineto{\pgfqpoint{2.575790in}{10.257287in}}%
\pgfpathlineto{\pgfqpoint{2.575790in}{10.239885in}}%
\pgfpathlineto{\pgfqpoint{2.562064in}{10.239885in}}%
\pgfpathlineto{\pgfqpoint{2.562064in}{10.257287in}}%
\pgfpathlineto{\pgfqpoint{2.562064in}{10.257287in}}%
\pgfpathclose%
\pgfpathmoveto{\pgfqpoint{2.652919in}{10.211931in}}%
\pgfpathquadraticcurveto{\pgfqpoint{2.650604in}{10.213268in}}{\pgfqpoint{2.647882in}{10.213889in}}%
\pgfpathquadraticcurveto{\pgfqpoint{2.645161in}{10.214533in}}{\pgfqpoint{2.641891in}{10.214533in}}%
\pgfpathquadraticcurveto{\pgfqpoint{2.630241in}{10.214533in}}{\pgfqpoint{2.624011in}{10.206966in}}%
\pgfpathquadraticcurveto{\pgfqpoint{2.617780in}{10.199399in}}{\pgfqpoint{2.617780in}{10.185219in}}%
\pgfpathlineto{\pgfqpoint{2.617780in}{10.141200in}}%
\pgfpathlineto{\pgfqpoint{2.603983in}{10.141200in}}%
\pgfpathlineto{\pgfqpoint{2.603983in}{10.224750in}}%
\pgfpathlineto{\pgfqpoint{2.617780in}{10.224750in}}%
\pgfpathlineto{\pgfqpoint{2.617780in}{10.211764in}}%
\pgfpathquadraticcurveto{\pgfqpoint{2.622125in}{10.219379in}}{\pgfqpoint{2.629048in}{10.223055in}}%
\pgfpathquadraticcurveto{\pgfqpoint{2.635994in}{10.226756in}}{\pgfqpoint{2.645925in}{10.226756in}}%
\pgfpathquadraticcurveto{\pgfqpoint{2.647333in}{10.226756in}}{\pgfqpoint{2.649052in}{10.226565in}}%
\pgfpathquadraticcurveto{\pgfqpoint{2.650771in}{10.226397in}}{\pgfqpoint{2.652848in}{10.226016in}}%
\pgfpathlineto{\pgfqpoint{2.652919in}{10.211931in}}%
\pgfpathlineto{\pgfqpoint{2.652919in}{10.211931in}}%
\pgfpathclose%
\pgfpathmoveto{\pgfqpoint{2.735419in}{10.186413in}}%
\pgfpathlineto{\pgfqpoint{2.735419in}{10.179705in}}%
\pgfpathlineto{\pgfqpoint{2.672303in}{10.179705in}}%
\pgfpathquadraticcurveto{\pgfqpoint{2.673210in}{10.165525in}}{\pgfqpoint{2.680849in}{10.158101in}}%
\pgfpathquadraticcurveto{\pgfqpoint{2.688488in}{10.150677in}}{\pgfqpoint{2.702142in}{10.150677in}}%
\pgfpathquadraticcurveto{\pgfqpoint{2.710044in}{10.150677in}}{\pgfqpoint{2.717468in}{10.152611in}}%
\pgfpathquadraticcurveto{\pgfqpoint{2.724892in}{10.154544in}}{\pgfqpoint{2.732220in}{10.158435in}}%
\pgfpathlineto{\pgfqpoint{2.732220in}{10.145449in}}%
\pgfpathquadraticcurveto{\pgfqpoint{2.724820in}{10.142322in}}{\pgfqpoint{2.717062in}{10.140675in}}%
\pgfpathquadraticcurveto{\pgfqpoint{2.709304in}{10.139028in}}{\pgfqpoint{2.701331in}{10.139028in}}%
\pgfpathquadraticcurveto{\pgfqpoint{2.681326in}{10.139028in}}{\pgfqpoint{2.669653in}{10.150653in}}%
\pgfpathquadraticcurveto{\pgfqpoint{2.657980in}{10.162302in}}{\pgfqpoint{2.657980in}{10.182164in}}%
\pgfpathquadraticcurveto{\pgfqpoint{2.657980in}{10.202669in}}{\pgfqpoint{2.669056in}{10.214700in}}%
\pgfpathquadraticcurveto{\pgfqpoint{2.680133in}{10.226756in}}{\pgfqpoint{2.698944in}{10.226756in}}%
\pgfpathquadraticcurveto{\pgfqpoint{2.715797in}{10.226756in}}{\pgfqpoint{2.725608in}{10.215894in}}%
\pgfpathquadraticcurveto{\pgfqpoint{2.735419in}{10.205056in}}{\pgfqpoint{2.735419in}{10.186413in}}%
\pgfpathlineto{\pgfqpoint{2.735419in}{10.186413in}}%
\pgfpathclose%
\pgfpathmoveto{\pgfqpoint{2.721693in}{10.190447in}}%
\pgfpathquadraticcurveto{\pgfqpoint{2.721550in}{10.201690in}}{\pgfqpoint{2.715391in}{10.208398in}}%
\pgfpathquadraticcurveto{\pgfqpoint{2.709232in}{10.215130in}}{\pgfqpoint{2.699087in}{10.215130in}}%
\pgfpathquadraticcurveto{\pgfqpoint{2.687605in}{10.215130in}}{\pgfqpoint{2.680706in}{10.208637in}}%
\pgfpathquadraticcurveto{\pgfqpoint{2.673807in}{10.202144in}}{\pgfqpoint{2.672757in}{10.190351in}}%
\pgfpathlineto{\pgfqpoint{2.721693in}{10.190447in}}%
\pgfpathlineto{\pgfqpoint{2.721693in}{10.190447in}}%
\pgfpathclose%
\pgfpathmoveto{\pgfqpoint{2.812930in}{10.212075in}}%
\pgfpathlineto{\pgfqpoint{2.812930in}{10.257287in}}%
\pgfpathlineto{\pgfqpoint{2.826656in}{10.257287in}}%
\pgfpathlineto{\pgfqpoint{2.826656in}{10.141200in}}%
\pgfpathlineto{\pgfqpoint{2.812930in}{10.141200in}}%
\pgfpathlineto{\pgfqpoint{2.812930in}{10.153733in}}%
\pgfpathquadraticcurveto{\pgfqpoint{2.808609in}{10.146285in}}{\pgfqpoint{2.801997in}{10.142656in}}%
\pgfpathquadraticcurveto{\pgfqpoint{2.795408in}{10.139028in}}{\pgfqpoint{2.786146in}{10.139028in}}%
\pgfpathquadraticcurveto{\pgfqpoint{2.771012in}{10.139028in}}{\pgfqpoint{2.761487in}{10.151107in}}%
\pgfpathquadraticcurveto{\pgfqpoint{2.751986in}{10.163210in}}{\pgfqpoint{2.751986in}{10.182904in}}%
\pgfpathquadraticcurveto{\pgfqpoint{2.751986in}{10.202598in}}{\pgfqpoint{2.761487in}{10.214677in}}%
\pgfpathquadraticcurveto{\pgfqpoint{2.771012in}{10.226756in}}{\pgfqpoint{2.786146in}{10.226756in}}%
\pgfpathquadraticcurveto{\pgfqpoint{2.795408in}{10.226756in}}{\pgfqpoint{2.801997in}{10.223127in}}%
\pgfpathquadraticcurveto{\pgfqpoint{2.808609in}{10.219522in}}{\pgfqpoint{2.812930in}{10.212075in}}%
\pgfpathlineto{\pgfqpoint{2.812930in}{10.212075in}}%
\pgfpathclose%
\pgfpathmoveto{\pgfqpoint{2.766166in}{10.182904in}}%
\pgfpathquadraticcurveto{\pgfqpoint{2.766166in}{10.167769in}}{\pgfqpoint{2.772396in}{10.159151in}}%
\pgfpathquadraticcurveto{\pgfqpoint{2.778627in}{10.150534in}}{\pgfqpoint{2.789512in}{10.150534in}}%
\pgfpathquadraticcurveto{\pgfqpoint{2.800398in}{10.150534in}}{\pgfqpoint{2.806652in}{10.159151in}}%
\pgfpathquadraticcurveto{\pgfqpoint{2.812930in}{10.167769in}}{\pgfqpoint{2.812930in}{10.182904in}}%
\pgfpathquadraticcurveto{\pgfqpoint{2.812930in}{10.198038in}}{\pgfqpoint{2.806652in}{10.206656in}}%
\pgfpathquadraticcurveto{\pgfqpoint{2.800398in}{10.215273in}}{\pgfqpoint{2.789512in}{10.215273in}}%
\pgfpathquadraticcurveto{\pgfqpoint{2.778627in}{10.215273in}}{\pgfqpoint{2.772396in}{10.206656in}}%
\pgfpathquadraticcurveto{\pgfqpoint{2.766166in}{10.198038in}}{\pgfqpoint{2.766166in}{10.182904in}}%
\pgfpathlineto{\pgfqpoint{2.766166in}{10.182904in}}%
\pgfpathclose%
\pgfpathmoveto{\pgfqpoint{2.921187in}{10.257956in}}%
\pgfpathlineto{\pgfqpoint{2.921187in}{10.105178in}}%
\pgfpathlineto{\pgfqpoint{2.908512in}{10.105178in}}%
\pgfpathlineto{\pgfqpoint{2.908512in}{10.257956in}}%
\pgfpathlineto{\pgfqpoint{2.921187in}{10.257956in}}%
\pgfpathlineto{\pgfqpoint{2.921187in}{10.257956in}}%
\pgfpathclose%
\pgfpathmoveto{\pgfqpoint{2.968549in}{10.248479in}}%
\pgfpathlineto{\pgfqpoint{2.968549in}{10.224750in}}%
\pgfpathlineto{\pgfqpoint{2.996812in}{10.224750in}}%
\pgfpathlineto{\pgfqpoint{2.996812in}{10.214080in}}%
\pgfpathlineto{\pgfqpoint{2.968549in}{10.214080in}}%
\pgfpathlineto{\pgfqpoint{2.968549in}{10.168724in}}%
\pgfpathquadraticcurveto{\pgfqpoint{2.968549in}{10.158507in}}{\pgfqpoint{2.971342in}{10.155595in}}%
\pgfpathquadraticcurveto{\pgfqpoint{2.974135in}{10.152682in}}{\pgfqpoint{2.982728in}{10.152682in}}%
\pgfpathlineto{\pgfqpoint{2.996812in}{10.152682in}}%
\pgfpathlineto{\pgfqpoint{2.996812in}{10.141200in}}%
\pgfpathlineto{\pgfqpoint{2.982728in}{10.141200in}}%
\pgfpathquadraticcurveto{\pgfqpoint{2.966830in}{10.141200in}}{\pgfqpoint{2.960790in}{10.147120in}}%
\pgfpathquadraticcurveto{\pgfqpoint{2.954751in}{10.153064in}}{\pgfqpoint{2.954751in}{10.168724in}}%
\pgfpathlineto{\pgfqpoint{2.954751in}{10.214080in}}%
\pgfpathlineto{\pgfqpoint{2.944677in}{10.214080in}}%
\pgfpathlineto{\pgfqpoint{2.944677in}{10.224750in}}%
\pgfpathlineto{\pgfqpoint{2.954751in}{10.224750in}}%
\pgfpathlineto{\pgfqpoint{2.954751in}{10.248479in}}%
\pgfpathlineto{\pgfqpoint{2.968549in}{10.248479in}}%
\pgfpathlineto{\pgfqpoint{2.968549in}{10.248479in}}%
\pgfpathclose%
\pgfpathmoveto{\pgfqpoint{3.032548in}{10.257956in}}%
\pgfpathlineto{\pgfqpoint{3.032548in}{10.105178in}}%
\pgfpathlineto{\pgfqpoint{3.019872in}{10.105178in}}%
\pgfpathlineto{\pgfqpoint{3.019872in}{10.257956in}}%
\pgfpathlineto{\pgfqpoint{3.032548in}{10.257956in}}%
\pgfpathlineto{\pgfqpoint{3.032548in}{10.257956in}}%
\pgfpathclose%
\pgfpathmoveto{\pgfqpoint{3.116671in}{10.216395in}}%
\pgfpathlineto{\pgfqpoint{3.116671in}{10.229978in}}%
\pgfpathlineto{\pgfqpoint{3.212301in}{10.195293in}}%
\pgfpathlineto{\pgfqpoint{3.212301in}{10.182904in}}%
\pgfpathlineto{\pgfqpoint{3.116671in}{10.148218in}}%
\pgfpathlineto{\pgfqpoint{3.116671in}{10.161801in}}%
\pgfpathlineto{\pgfqpoint{3.193514in}{10.189015in}}%
\pgfpathlineto{\pgfqpoint{3.116671in}{10.216395in}}%
\pgfpathlineto{\pgfqpoint{3.116671in}{10.216395in}}%
\pgfpathclose%
\pgfpathmoveto{\pgfqpoint{3.296018in}{10.153876in}}%
\pgfpathlineto{\pgfqpoint{3.320630in}{10.153876in}}%
\pgfpathlineto{\pgfqpoint{3.320630in}{10.238858in}}%
\pgfpathlineto{\pgfqpoint{3.293846in}{10.233487in}}%
\pgfpathlineto{\pgfqpoint{3.293846in}{10.247213in}}%
\pgfpathlineto{\pgfqpoint{3.320487in}{10.252585in}}%
\pgfpathlineto{\pgfqpoint{3.335549in}{10.252585in}}%
\pgfpathlineto{\pgfqpoint{3.335549in}{10.153876in}}%
\pgfpathlineto{\pgfqpoint{3.360161in}{10.153876in}}%
\pgfpathlineto{\pgfqpoint{3.360161in}{10.141200in}}%
\pgfpathlineto{\pgfqpoint{3.296018in}{10.141200in}}%
\pgfpathlineto{\pgfqpoint{3.296018in}{10.153876in}}%
\pgfpathlineto{\pgfqpoint{3.296018in}{10.153876in}}%
\pgfpathclose%
\pgfpathmoveto{\pgfqpoint{3.390597in}{10.160154in}}%
\pgfpathlineto{\pgfqpoint{3.406352in}{10.160154in}}%
\pgfpathlineto{\pgfqpoint{3.406352in}{10.141200in}}%
\pgfpathlineto{\pgfqpoint{3.390597in}{10.141200in}}%
\pgfpathlineto{\pgfqpoint{3.390597in}{10.160154in}}%
\pgfpathlineto{\pgfqpoint{3.390597in}{10.160154in}}%
\pgfpathclose%
\pgfpathmoveto{\pgfqpoint{3.439605in}{10.143516in}}%
\pgfpathlineto{\pgfqpoint{3.439605in}{10.157242in}}%
\pgfpathquadraticcurveto{\pgfqpoint{3.445287in}{10.154544in}}{\pgfqpoint{3.451088in}{10.153136in}}%
\pgfpathquadraticcurveto{\pgfqpoint{3.456912in}{10.151727in}}{\pgfqpoint{3.462522in}{10.151727in}}%
\pgfpathquadraticcurveto{\pgfqpoint{3.477442in}{10.151727in}}{\pgfqpoint{3.485296in}{10.161753in}}%
\pgfpathquadraticcurveto{\pgfqpoint{3.493173in}{10.171779in}}{\pgfqpoint{3.494295in}{10.192237in}}%
\pgfpathquadraticcurveto{\pgfqpoint{3.489974in}{10.185816in}}{\pgfqpoint{3.483314in}{10.182378in}}%
\pgfpathquadraticcurveto{\pgfqpoint{3.476678in}{10.178941in}}{\pgfqpoint{3.468633in}{10.178941in}}%
\pgfpathquadraticcurveto{\pgfqpoint{3.451923in}{10.178941in}}{\pgfqpoint{3.442184in}{10.189039in}}%
\pgfpathquadraticcurveto{\pgfqpoint{3.432444in}{10.199160in}}{\pgfqpoint{3.432444in}{10.216706in}}%
\pgfpathquadraticcurveto{\pgfqpoint{3.432444in}{10.233845in}}{\pgfqpoint{3.442589in}{10.244206in}}%
\pgfpathquadraticcurveto{\pgfqpoint{3.452735in}{10.254590in}}{\pgfqpoint{3.469588in}{10.254590in}}%
\pgfpathquadraticcurveto{\pgfqpoint{3.488924in}{10.254590in}}{\pgfqpoint{3.499093in}{10.239766in}}%
\pgfpathquadraticcurveto{\pgfqpoint{3.509286in}{10.224965in}}{\pgfqpoint{3.509286in}{10.196773in}}%
\pgfpathquadraticcurveto{\pgfqpoint{3.509286in}{10.170443in}}{\pgfqpoint{3.496778in}{10.154735in}}%
\pgfpathquadraticcurveto{\pgfqpoint{3.484293in}{10.139028in}}{\pgfqpoint{3.463191in}{10.139028in}}%
\pgfpathquadraticcurveto{\pgfqpoint{3.457509in}{10.139028in}}{\pgfqpoint{3.451684in}{10.140150in}}%
\pgfpathquadraticcurveto{\pgfqpoint{3.445884in}{10.141272in}}{\pgfqpoint{3.439605in}{10.143516in}}%
\pgfpathlineto{\pgfqpoint{3.439605in}{10.143516in}}%
\pgfpathclose%
\pgfpathmoveto{\pgfqpoint{3.469588in}{10.190733in}}%
\pgfpathquadraticcurveto{\pgfqpoint{3.479734in}{10.190733in}}{\pgfqpoint{3.485654in}{10.197656in}}%
\pgfpathquadraticcurveto{\pgfqpoint{3.491598in}{10.204603in}}{\pgfqpoint{3.491598in}{10.216706in}}%
\pgfpathquadraticcurveto{\pgfqpoint{3.491598in}{10.228713in}}{\pgfqpoint{3.485654in}{10.235684in}}%
\pgfpathquadraticcurveto{\pgfqpoint{3.479734in}{10.242654in}}{\pgfqpoint{3.469588in}{10.242654in}}%
\pgfpathquadraticcurveto{\pgfqpoint{3.459443in}{10.242654in}}{\pgfqpoint{3.453523in}{10.235684in}}%
\pgfpathquadraticcurveto{\pgfqpoint{3.447602in}{10.228713in}}{\pgfqpoint{3.447602in}{10.216706in}}%
\pgfpathquadraticcurveto{\pgfqpoint{3.447602in}{10.204603in}}{\pgfqpoint{3.453523in}{10.197656in}}%
\pgfpathquadraticcurveto{\pgfqpoint{3.459443in}{10.190733in}}{\pgfqpoint{3.469588in}{10.190733in}}%
\pgfpathlineto{\pgfqpoint{3.469588in}{10.190733in}}%
\pgfpathclose%
\pgfpathmoveto{\pgfqpoint{3.570469in}{10.202884in}}%
\pgfpathquadraticcurveto{\pgfqpoint{3.560324in}{10.202884in}}{\pgfqpoint{3.554380in}{10.195937in}}%
\pgfpathquadraticcurveto{\pgfqpoint{3.548460in}{10.189015in}}{\pgfqpoint{3.548460in}{10.176936in}}%
\pgfpathquadraticcurveto{\pgfqpoint{3.548460in}{10.164928in}}{\pgfqpoint{3.554380in}{10.157934in}}%
\pgfpathquadraticcurveto{\pgfqpoint{3.560324in}{10.150963in}}{\pgfqpoint{3.570469in}{10.150963in}}%
\pgfpathquadraticcurveto{\pgfqpoint{3.580615in}{10.150963in}}{\pgfqpoint{3.586535in}{10.157934in}}%
\pgfpathquadraticcurveto{\pgfqpoint{3.592455in}{10.164928in}}{\pgfqpoint{3.592455in}{10.176936in}}%
\pgfpathquadraticcurveto{\pgfqpoint{3.592455in}{10.189015in}}{\pgfqpoint{3.586535in}{10.195937in}}%
\pgfpathquadraticcurveto{\pgfqpoint{3.580615in}{10.202884in}}{\pgfqpoint{3.570469in}{10.202884in}}%
\pgfpathlineto{\pgfqpoint{3.570469in}{10.202884in}}%
\pgfpathclose%
\pgfpathmoveto{\pgfqpoint{3.600380in}{10.250126in}}%
\pgfpathlineto{\pgfqpoint{3.600380in}{10.236400in}}%
\pgfpathquadraticcurveto{\pgfqpoint{3.594699in}{10.239073in}}{\pgfqpoint{3.588922in}{10.240482in}}%
\pgfpathquadraticcurveto{\pgfqpoint{3.583145in}{10.241914in}}{\pgfqpoint{3.577464in}{10.241914in}}%
\pgfpathquadraticcurveto{\pgfqpoint{3.562544in}{10.241914in}}{\pgfqpoint{3.554666in}{10.231840in}}%
\pgfpathquadraticcurveto{\pgfqpoint{3.546812in}{10.221766in}}{\pgfqpoint{3.545691in}{10.201404in}}%
\pgfpathquadraticcurveto{\pgfqpoint{3.550083in}{10.207897in}}{\pgfqpoint{3.556719in}{10.211358in}}%
\pgfpathquadraticcurveto{\pgfqpoint{3.563379in}{10.214820in}}{\pgfqpoint{3.571352in}{10.214820in}}%
\pgfpathquadraticcurveto{\pgfqpoint{3.588134in}{10.214820in}}{\pgfqpoint{3.597874in}{10.204627in}}%
\pgfpathquadraticcurveto{\pgfqpoint{3.607613in}{10.194457in}}{\pgfqpoint{3.607613in}{10.176936in}}%
\pgfpathquadraticcurveto{\pgfqpoint{3.607613in}{10.159772in}}{\pgfqpoint{3.597468in}{10.149388in}}%
\pgfpathquadraticcurveto{\pgfqpoint{3.587322in}{10.139028in}}{\pgfqpoint{3.570469in}{10.139028in}}%
\pgfpathquadraticcurveto{\pgfqpoint{3.551133in}{10.139028in}}{\pgfqpoint{3.540916in}{10.153828in}}%
\pgfpathquadraticcurveto{\pgfqpoint{3.530699in}{10.168652in}}{\pgfqpoint{3.530699in}{10.196773in}}%
\pgfpathquadraticcurveto{\pgfqpoint{3.530699in}{10.223175in}}{\pgfqpoint{3.543232in}{10.238882in}}%
\pgfpathquadraticcurveto{\pgfqpoint{3.555764in}{10.254590in}}{\pgfqpoint{3.576867in}{10.254590in}}%
\pgfpathquadraticcurveto{\pgfqpoint{3.582548in}{10.254590in}}{\pgfqpoint{3.588325in}{10.253468in}}%
\pgfpathquadraticcurveto{\pgfqpoint{3.594102in}{10.252346in}}{\pgfqpoint{3.600380in}{10.250126in}}%
\pgfpathlineto{\pgfqpoint{3.600380in}{10.250126in}}%
\pgfpathclose%
\pgfpathmoveto{\pgfqpoint{3.682116in}{10.203720in}}%
\pgfpathlineto{\pgfqpoint{3.697872in}{10.203720in}}%
\pgfpathlineto{\pgfqpoint{3.697872in}{10.184766in}}%
\pgfpathlineto{\pgfqpoint{3.682116in}{10.184766in}}%
\pgfpathlineto{\pgfqpoint{3.682116in}{10.203720in}}%
\pgfpathlineto{\pgfqpoint{3.682116in}{10.203720in}}%
\pgfpathclose%
\pgfpathmoveto{\pgfqpoint{3.777889in}{10.252585in}}%
\pgfpathlineto{\pgfqpoint{3.792952in}{10.252585in}}%
\pgfpathlineto{\pgfqpoint{3.792952in}{10.206918in}}%
\pgfpathlineto{\pgfqpoint{3.847713in}{10.206918in}}%
\pgfpathlineto{\pgfqpoint{3.847713in}{10.252585in}}%
\pgfpathlineto{\pgfqpoint{3.862776in}{10.252585in}}%
\pgfpathlineto{\pgfqpoint{3.862776in}{10.141200in}}%
\pgfpathlineto{\pgfqpoint{3.847713in}{10.141200in}}%
\pgfpathlineto{\pgfqpoint{3.847713in}{10.194243in}}%
\pgfpathlineto{\pgfqpoint{3.792952in}{10.194243in}}%
\pgfpathlineto{\pgfqpoint{3.792952in}{10.141200in}}%
\pgfpathlineto{\pgfqpoint{3.777889in}{10.141200in}}%
\pgfpathlineto{\pgfqpoint{3.777889in}{10.252585in}}%
\pgfpathlineto{\pgfqpoint{3.777889in}{10.252585in}}%
\pgfpathclose%
\pgfpathmoveto{\pgfqpoint{3.924556in}{10.215130in}}%
\pgfpathquadraticcurveto{\pgfqpoint{3.913527in}{10.215130in}}{\pgfqpoint{3.907105in}{10.206512in}}%
\pgfpathquadraticcurveto{\pgfqpoint{3.900684in}{10.197895in}}{\pgfqpoint{3.900684in}{10.182904in}}%
\pgfpathquadraticcurveto{\pgfqpoint{3.900684in}{10.167912in}}{\pgfqpoint{3.907058in}{10.159295in}}%
\pgfpathquadraticcurveto{\pgfqpoint{3.913455in}{10.150677in}}{\pgfqpoint{3.924556in}{10.150677in}}%
\pgfpathquadraticcurveto{\pgfqpoint{3.935536in}{10.150677in}}{\pgfqpoint{3.941934in}{10.159318in}}%
\pgfpathquadraticcurveto{\pgfqpoint{3.948355in}{10.167984in}}{\pgfqpoint{3.948355in}{10.182904in}}%
\pgfpathquadraticcurveto{\pgfqpoint{3.948355in}{10.197752in}}{\pgfqpoint{3.941934in}{10.206441in}}%
\pgfpathquadraticcurveto{\pgfqpoint{3.935536in}{10.215130in}}{\pgfqpoint{3.924556in}{10.215130in}}%
\pgfpathlineto{\pgfqpoint{3.924556in}{10.215130in}}%
\pgfpathclose%
\pgfpathmoveto{\pgfqpoint{3.924556in}{10.226756in}}%
\pgfpathquadraticcurveto{\pgfqpoint{3.942459in}{10.226756in}}{\pgfqpoint{3.952676in}{10.215106in}}%
\pgfpathquadraticcurveto{\pgfqpoint{3.962917in}{10.203481in}}{\pgfqpoint{3.962917in}{10.182904in}}%
\pgfpathquadraticcurveto{\pgfqpoint{3.962917in}{10.162398in}}{\pgfqpoint{3.952676in}{10.150701in}}%
\pgfpathquadraticcurveto{\pgfqpoint{3.942459in}{10.139028in}}{\pgfqpoint{3.924556in}{10.139028in}}%
\pgfpathquadraticcurveto{\pgfqpoint{3.906580in}{10.139028in}}{\pgfqpoint{3.896387in}{10.150701in}}%
\pgfpathquadraticcurveto{\pgfqpoint{3.886218in}{10.162398in}}{\pgfqpoint{3.886218in}{10.182904in}}%
\pgfpathquadraticcurveto{\pgfqpoint{3.886218in}{10.203481in}}{\pgfqpoint{3.896387in}{10.215106in}}%
\pgfpathquadraticcurveto{\pgfqpoint{3.906580in}{10.226756in}}{\pgfqpoint{3.924556in}{10.226756in}}%
\pgfpathlineto{\pgfqpoint{3.924556in}{10.226756in}}%
\pgfpathclose%
\pgfpathmoveto{\pgfqpoint{3.985667in}{10.257287in}}%
\pgfpathlineto{\pgfqpoint{3.999393in}{10.257287in}}%
\pgfpathlineto{\pgfqpoint{3.999393in}{10.141200in}}%
\pgfpathlineto{\pgfqpoint{3.985667in}{10.141200in}}%
\pgfpathlineto{\pgfqpoint{3.985667in}{10.257287in}}%
\pgfpathlineto{\pgfqpoint{3.985667in}{10.257287in}}%
\pgfpathclose%
\pgfpathmoveto{\pgfqpoint{4.028110in}{10.257287in}}%
\pgfpathlineto{\pgfqpoint{4.041836in}{10.257287in}}%
\pgfpathlineto{\pgfqpoint{4.041836in}{10.141200in}}%
\pgfpathlineto{\pgfqpoint{4.028110in}{10.141200in}}%
\pgfpathlineto{\pgfqpoint{4.028110in}{10.257287in}}%
\pgfpathlineto{\pgfqpoint{4.028110in}{10.257287in}}%
\pgfpathclose%
\pgfpathmoveto{\pgfqpoint{4.102924in}{10.215130in}}%
\pgfpathquadraticcurveto{\pgfqpoint{4.091895in}{10.215130in}}{\pgfqpoint{4.085474in}{10.206512in}}%
\pgfpathquadraticcurveto{\pgfqpoint{4.079052in}{10.197895in}}{\pgfqpoint{4.079052in}{10.182904in}}%
\pgfpathquadraticcurveto{\pgfqpoint{4.079052in}{10.167912in}}{\pgfqpoint{4.085426in}{10.159295in}}%
\pgfpathquadraticcurveto{\pgfqpoint{4.091823in}{10.150677in}}{\pgfqpoint{4.102924in}{10.150677in}}%
\pgfpathquadraticcurveto{\pgfqpoint{4.113905in}{10.150677in}}{\pgfqpoint{4.120302in}{10.159318in}}%
\pgfpathquadraticcurveto{\pgfqpoint{4.126724in}{10.167984in}}{\pgfqpoint{4.126724in}{10.182904in}}%
\pgfpathquadraticcurveto{\pgfqpoint{4.126724in}{10.197752in}}{\pgfqpoint{4.120302in}{10.206441in}}%
\pgfpathquadraticcurveto{\pgfqpoint{4.113905in}{10.215130in}}{\pgfqpoint{4.102924in}{10.215130in}}%
\pgfpathlineto{\pgfqpoint{4.102924in}{10.215130in}}%
\pgfpathclose%
\pgfpathmoveto{\pgfqpoint{4.102924in}{10.226756in}}%
\pgfpathquadraticcurveto{\pgfqpoint{4.120827in}{10.226756in}}{\pgfqpoint{4.131044in}{10.215106in}}%
\pgfpathquadraticcurveto{\pgfqpoint{4.141285in}{10.203481in}}{\pgfqpoint{4.141285in}{10.182904in}}%
\pgfpathquadraticcurveto{\pgfqpoint{4.141285in}{10.162398in}}{\pgfqpoint{4.131044in}{10.150701in}}%
\pgfpathquadraticcurveto{\pgfqpoint{4.120827in}{10.139028in}}{\pgfqpoint{4.102924in}{10.139028in}}%
\pgfpathquadraticcurveto{\pgfqpoint{4.084948in}{10.139028in}}{\pgfqpoint{4.074755in}{10.150701in}}%
\pgfpathquadraticcurveto{\pgfqpoint{4.064586in}{10.162398in}}{\pgfqpoint{4.064586in}{10.182904in}}%
\pgfpathquadraticcurveto{\pgfqpoint{4.064586in}{10.203481in}}{\pgfqpoint{4.074755in}{10.215106in}}%
\pgfpathquadraticcurveto{\pgfqpoint{4.084948in}{10.226756in}}{\pgfqpoint{4.102924in}{10.226756in}}%
\pgfpathlineto{\pgfqpoint{4.102924in}{10.226756in}}%
\pgfpathclose%
\pgfpathmoveto{\pgfqpoint{4.156062in}{10.224750in}}%
\pgfpathlineto{\pgfqpoint{4.169788in}{10.224750in}}%
\pgfpathlineto{\pgfqpoint{4.186951in}{10.159557in}}%
\pgfpathlineto{\pgfqpoint{4.204020in}{10.224750in}}%
\pgfpathlineto{\pgfqpoint{4.220204in}{10.224750in}}%
\pgfpathlineto{\pgfqpoint{4.237368in}{10.159557in}}%
\pgfpathlineto{\pgfqpoint{4.254460in}{10.224750in}}%
\pgfpathlineto{\pgfqpoint{4.268186in}{10.224750in}}%
\pgfpathlineto{\pgfqpoint{4.246320in}{10.141200in}}%
\pgfpathlineto{\pgfqpoint{4.230135in}{10.141200in}}%
\pgfpathlineto{\pgfqpoint{4.212160in}{10.209687in}}%
\pgfpathlineto{\pgfqpoint{4.194113in}{10.141200in}}%
\pgfpathlineto{\pgfqpoint{4.177904in}{10.141200in}}%
\pgfpathlineto{\pgfqpoint{4.156062in}{10.224750in}}%
\pgfpathlineto{\pgfqpoint{4.156062in}{10.224750in}}%
\pgfpathclose%
\pgfpathmoveto{\pgfqpoint{4.284132in}{10.160154in}}%
\pgfpathlineto{\pgfqpoint{4.299864in}{10.160154in}}%
\pgfpathlineto{\pgfqpoint{4.299864in}{10.141200in}}%
\pgfpathlineto{\pgfqpoint{4.284132in}{10.141200in}}%
\pgfpathlineto{\pgfqpoint{4.284132in}{10.160154in}}%
\pgfpathlineto{\pgfqpoint{4.284132in}{10.160154in}}%
\pgfpathclose%
\pgfpathmoveto{\pgfqpoint{4.284132in}{10.220191in}}%
\pgfpathlineto{\pgfqpoint{4.299864in}{10.220191in}}%
\pgfpathlineto{\pgfqpoint{4.299864in}{10.201261in}}%
\pgfpathlineto{\pgfqpoint{4.284132in}{10.201261in}}%
\pgfpathlineto{\pgfqpoint{4.284132in}{10.220191in}}%
\pgfpathlineto{\pgfqpoint{4.284132in}{10.220191in}}%
\pgfpathclose%
\pgfpathmoveto{\pgfqpoint{4.413015in}{10.215130in}}%
\pgfpathquadraticcurveto{\pgfqpoint{4.401986in}{10.215130in}}{\pgfqpoint{4.395565in}{10.206512in}}%
\pgfpathquadraticcurveto{\pgfqpoint{4.389143in}{10.197895in}}{\pgfqpoint{4.389143in}{10.182904in}}%
\pgfpathquadraticcurveto{\pgfqpoint{4.389143in}{10.167912in}}{\pgfqpoint{4.395517in}{10.159295in}}%
\pgfpathquadraticcurveto{\pgfqpoint{4.401914in}{10.150677in}}{\pgfqpoint{4.413015in}{10.150677in}}%
\pgfpathquadraticcurveto{\pgfqpoint{4.423996in}{10.150677in}}{\pgfqpoint{4.430393in}{10.159318in}}%
\pgfpathquadraticcurveto{\pgfqpoint{4.436815in}{10.167984in}}{\pgfqpoint{4.436815in}{10.182904in}}%
\pgfpathquadraticcurveto{\pgfqpoint{4.436815in}{10.197752in}}{\pgfqpoint{4.430393in}{10.206441in}}%
\pgfpathquadraticcurveto{\pgfqpoint{4.423996in}{10.215130in}}{\pgfqpoint{4.413015in}{10.215130in}}%
\pgfpathlineto{\pgfqpoint{4.413015in}{10.215130in}}%
\pgfpathclose%
\pgfpathmoveto{\pgfqpoint{4.413015in}{10.226756in}}%
\pgfpathquadraticcurveto{\pgfqpoint{4.430918in}{10.226756in}}{\pgfqpoint{4.441135in}{10.215106in}}%
\pgfpathquadraticcurveto{\pgfqpoint{4.451376in}{10.203481in}}{\pgfqpoint{4.451376in}{10.182904in}}%
\pgfpathquadraticcurveto{\pgfqpoint{4.451376in}{10.162398in}}{\pgfqpoint{4.441135in}{10.150701in}}%
\pgfpathquadraticcurveto{\pgfqpoint{4.430918in}{10.139028in}}{\pgfqpoint{4.413015in}{10.139028in}}%
\pgfpathquadraticcurveto{\pgfqpoint{4.395039in}{10.139028in}}{\pgfqpoint{4.384846in}{10.150701in}}%
\pgfpathquadraticcurveto{\pgfqpoint{4.374677in}{10.162398in}}{\pgfqpoint{4.374677in}{10.182904in}}%
\pgfpathquadraticcurveto{\pgfqpoint{4.374677in}{10.203481in}}{\pgfqpoint{4.384846in}{10.215106in}}%
\pgfpathquadraticcurveto{\pgfqpoint{4.395039in}{10.226756in}}{\pgfqpoint{4.413015in}{10.226756in}}%
\pgfpathlineto{\pgfqpoint{4.413015in}{10.226756in}}%
\pgfpathclose%
\pgfpathmoveto{\pgfqpoint{4.487709in}{10.248479in}}%
\pgfpathlineto{\pgfqpoint{4.487709in}{10.224750in}}%
\pgfpathlineto{\pgfqpoint{4.515973in}{10.224750in}}%
\pgfpathlineto{\pgfqpoint{4.515973in}{10.214080in}}%
\pgfpathlineto{\pgfqpoint{4.487709in}{10.214080in}}%
\pgfpathlineto{\pgfqpoint{4.487709in}{10.168724in}}%
\pgfpathquadraticcurveto{\pgfqpoint{4.487709in}{10.158507in}}{\pgfqpoint{4.490502in}{10.155595in}}%
\pgfpathquadraticcurveto{\pgfqpoint{4.493295in}{10.152682in}}{\pgfqpoint{4.501888in}{10.152682in}}%
\pgfpathlineto{\pgfqpoint{4.515973in}{10.152682in}}%
\pgfpathlineto{\pgfqpoint{4.515973in}{10.141200in}}%
\pgfpathlineto{\pgfqpoint{4.501888in}{10.141200in}}%
\pgfpathquadraticcurveto{\pgfqpoint{4.485990in}{10.141200in}}{\pgfqpoint{4.479951in}{10.147120in}}%
\pgfpathquadraticcurveto{\pgfqpoint{4.473911in}{10.153064in}}{\pgfqpoint{4.473911in}{10.168724in}}%
\pgfpathlineto{\pgfqpoint{4.473911in}{10.214080in}}%
\pgfpathlineto{\pgfqpoint{4.463837in}{10.214080in}}%
\pgfpathlineto{\pgfqpoint{4.463837in}{10.224750in}}%
\pgfpathlineto{\pgfqpoint{4.473911in}{10.224750in}}%
\pgfpathlineto{\pgfqpoint{4.473911in}{10.248479in}}%
\pgfpathlineto{\pgfqpoint{4.487709in}{10.248479in}}%
\pgfpathlineto{\pgfqpoint{4.487709in}{10.248479in}}%
\pgfpathclose%
\pgfpathmoveto{\pgfqpoint{4.603486in}{10.191641in}}%
\pgfpathlineto{\pgfqpoint{4.603486in}{10.141200in}}%
\pgfpathlineto{\pgfqpoint{4.589760in}{10.141200in}}%
\pgfpathlineto{\pgfqpoint{4.589760in}{10.191187in}}%
\pgfpathquadraticcurveto{\pgfqpoint{4.589760in}{10.203051in}}{\pgfqpoint{4.585128in}{10.208924in}}%
\pgfpathquadraticcurveto{\pgfqpoint{4.580497in}{10.214820in}}{\pgfqpoint{4.571259in}{10.214820in}}%
\pgfpathquadraticcurveto{\pgfqpoint{4.560135in}{10.214820in}}{\pgfqpoint{4.553714in}{10.207730in}}%
\pgfpathquadraticcurveto{\pgfqpoint{4.547292in}{10.200664in}}{\pgfqpoint{4.547292in}{10.188418in}}%
\pgfpathlineto{\pgfqpoint{4.547292in}{10.141200in}}%
\pgfpathlineto{\pgfqpoint{4.533494in}{10.141200in}}%
\pgfpathlineto{\pgfqpoint{4.533494in}{10.257287in}}%
\pgfpathlineto{\pgfqpoint{4.547292in}{10.257287in}}%
\pgfpathlineto{\pgfqpoint{4.547292in}{10.211764in}}%
\pgfpathquadraticcurveto{\pgfqpoint{4.552234in}{10.219308in}}{\pgfqpoint{4.558894in}{10.223032in}}%
\pgfpathquadraticcurveto{\pgfqpoint{4.565578in}{10.226756in}}{\pgfqpoint{4.574315in}{10.226756in}}%
\pgfpathquadraticcurveto{\pgfqpoint{4.588709in}{10.226756in}}{\pgfqpoint{4.596086in}{10.217851in}}%
\pgfpathquadraticcurveto{\pgfqpoint{4.603486in}{10.208947in}}{\pgfqpoint{4.603486in}{10.191641in}}%
\pgfpathlineto{\pgfqpoint{4.603486in}{10.191641in}}%
\pgfpathclose%
\pgfpathmoveto{\pgfqpoint{4.702314in}{10.186413in}}%
\pgfpathlineto{\pgfqpoint{4.702314in}{10.179705in}}%
\pgfpathlineto{\pgfqpoint{4.639197in}{10.179705in}}%
\pgfpathquadraticcurveto{\pgfqpoint{4.640105in}{10.165525in}}{\pgfqpoint{4.647743in}{10.158101in}}%
\pgfpathquadraticcurveto{\pgfqpoint{4.655382in}{10.150677in}}{\pgfqpoint{4.669037in}{10.150677in}}%
\pgfpathquadraticcurveto{\pgfqpoint{4.676938in}{10.150677in}}{\pgfqpoint{4.684362in}{10.152611in}}%
\pgfpathquadraticcurveto{\pgfqpoint{4.691786in}{10.154544in}}{\pgfqpoint{4.699115in}{10.158435in}}%
\pgfpathlineto{\pgfqpoint{4.699115in}{10.145449in}}%
\pgfpathquadraticcurveto{\pgfqpoint{4.691715in}{10.142322in}}{\pgfqpoint{4.683957in}{10.140675in}}%
\pgfpathquadraticcurveto{\pgfqpoint{4.676198in}{10.139028in}}{\pgfqpoint{4.668225in}{10.139028in}}%
\pgfpathquadraticcurveto{\pgfqpoint{4.648221in}{10.139028in}}{\pgfqpoint{4.636548in}{10.150653in}}%
\pgfpathquadraticcurveto{\pgfqpoint{4.624875in}{10.162302in}}{\pgfqpoint{4.624875in}{10.182164in}}%
\pgfpathquadraticcurveto{\pgfqpoint{4.624875in}{10.202669in}}{\pgfqpoint{4.635951in}{10.214700in}}%
\pgfpathquadraticcurveto{\pgfqpoint{4.647027in}{10.226756in}}{\pgfqpoint{4.665838in}{10.226756in}}%
\pgfpathquadraticcurveto{\pgfqpoint{4.682691in}{10.226756in}}{\pgfqpoint{4.692503in}{10.215894in}}%
\pgfpathquadraticcurveto{\pgfqpoint{4.702314in}{10.205056in}}{\pgfqpoint{4.702314in}{10.186413in}}%
\pgfpathlineto{\pgfqpoint{4.702314in}{10.186413in}}%
\pgfpathclose%
\pgfpathmoveto{\pgfqpoint{4.688588in}{10.190447in}}%
\pgfpathquadraticcurveto{\pgfqpoint{4.688444in}{10.201690in}}{\pgfqpoint{4.682286in}{10.208398in}}%
\pgfpathquadraticcurveto{\pgfqpoint{4.676127in}{10.215130in}}{\pgfqpoint{4.665981in}{10.215130in}}%
\pgfpathquadraticcurveto{\pgfqpoint{4.654499in}{10.215130in}}{\pgfqpoint{4.647600in}{10.208637in}}%
\pgfpathquadraticcurveto{\pgfqpoint{4.640701in}{10.202144in}}{\pgfqpoint{4.639651in}{10.190351in}}%
\pgfpathlineto{\pgfqpoint{4.688588in}{10.190447in}}%
\pgfpathlineto{\pgfqpoint{4.688588in}{10.190447in}}%
\pgfpathclose%
\pgfpathmoveto{\pgfqpoint{4.773260in}{10.211931in}}%
\pgfpathquadraticcurveto{\pgfqpoint{4.770944in}{10.213268in}}{\pgfqpoint{4.768223in}{10.213889in}}%
\pgfpathquadraticcurveto{\pgfqpoint{4.765502in}{10.214533in}}{\pgfqpoint{4.762231in}{10.214533in}}%
\pgfpathquadraticcurveto{\pgfqpoint{4.750582in}{10.214533in}}{\pgfqpoint{4.744352in}{10.206966in}}%
\pgfpathquadraticcurveto{\pgfqpoint{4.738121in}{10.199399in}}{\pgfqpoint{4.738121in}{10.185219in}}%
\pgfpathlineto{\pgfqpoint{4.738121in}{10.141200in}}%
\pgfpathlineto{\pgfqpoint{4.724323in}{10.141200in}}%
\pgfpathlineto{\pgfqpoint{4.724323in}{10.224750in}}%
\pgfpathlineto{\pgfqpoint{4.738121in}{10.224750in}}%
\pgfpathlineto{\pgfqpoint{4.738121in}{10.211764in}}%
\pgfpathquadraticcurveto{\pgfqpoint{4.742466in}{10.219379in}}{\pgfqpoint{4.749388in}{10.223055in}}%
\pgfpathquadraticcurveto{\pgfqpoint{4.756335in}{10.226756in}}{\pgfqpoint{4.766266in}{10.226756in}}%
\pgfpathquadraticcurveto{\pgfqpoint{4.767674in}{10.226756in}}{\pgfqpoint{4.769393in}{10.226565in}}%
\pgfpathquadraticcurveto{\pgfqpoint{4.771112in}{10.226397in}}{\pgfqpoint{4.773188in}{10.226016in}}%
\pgfpathlineto{\pgfqpoint{4.773260in}{10.211931in}}%
\pgfpathlineto{\pgfqpoint{4.773260in}{10.211931in}}%
\pgfpathclose%
\pgfpathmoveto{\pgfqpoint{4.779681in}{10.224750in}}%
\pgfpathlineto{\pgfqpoint{4.793408in}{10.224750in}}%
\pgfpathlineto{\pgfqpoint{4.810571in}{10.159557in}}%
\pgfpathlineto{\pgfqpoint{4.827639in}{10.224750in}}%
\pgfpathlineto{\pgfqpoint{4.843824in}{10.224750in}}%
\pgfpathlineto{\pgfqpoint{4.860988in}{10.159557in}}%
\pgfpathlineto{\pgfqpoint{4.878080in}{10.224750in}}%
\pgfpathlineto{\pgfqpoint{4.891806in}{10.224750in}}%
\pgfpathlineto{\pgfqpoint{4.869940in}{10.141200in}}%
\pgfpathlineto{\pgfqpoint{4.853755in}{10.141200in}}%
\pgfpathlineto{\pgfqpoint{4.835780in}{10.209687in}}%
\pgfpathlineto{\pgfqpoint{4.817733in}{10.141200in}}%
\pgfpathlineto{\pgfqpoint{4.801524in}{10.141200in}}%
\pgfpathlineto{\pgfqpoint{4.779681in}{10.224750in}}%
\pgfpathlineto{\pgfqpoint{4.779681in}{10.224750in}}%
\pgfpathclose%
\pgfpathmoveto{\pgfqpoint{4.912598in}{10.224750in}}%
\pgfpathlineto{\pgfqpoint{4.926324in}{10.224750in}}%
\pgfpathlineto{\pgfqpoint{4.926324in}{10.141200in}}%
\pgfpathlineto{\pgfqpoint{4.912598in}{10.141200in}}%
\pgfpathlineto{\pgfqpoint{4.912598in}{10.224750in}}%
\pgfpathlineto{\pgfqpoint{4.912598in}{10.224750in}}%
\pgfpathclose%
\pgfpathmoveto{\pgfqpoint{4.912598in}{10.257287in}}%
\pgfpathlineto{\pgfqpoint{4.926324in}{10.257287in}}%
\pgfpathlineto{\pgfqpoint{4.926324in}{10.239885in}}%
\pgfpathlineto{\pgfqpoint{4.912598in}{10.239885in}}%
\pgfpathlineto{\pgfqpoint{4.912598in}{10.257287in}}%
\pgfpathlineto{\pgfqpoint{4.912598in}{10.257287in}}%
\pgfpathclose%
\pgfpathmoveto{\pgfqpoint{5.008299in}{10.222292in}}%
\pgfpathlineto{\pgfqpoint{5.008299in}{10.209305in}}%
\pgfpathquadraticcurveto{\pgfqpoint{5.002498in}{10.212289in}}{\pgfqpoint{4.996220in}{10.213769in}}%
\pgfpathquadraticcurveto{\pgfqpoint{4.989966in}{10.215273in}}{\pgfqpoint{4.983234in}{10.215273in}}%
\pgfpathquadraticcurveto{\pgfqpoint{4.973017in}{10.215273in}}{\pgfqpoint{4.967908in}{10.212146in}}%
\pgfpathquadraticcurveto{\pgfqpoint{4.962800in}{10.209019in}}{\pgfqpoint{4.962800in}{10.202741in}}%
\pgfpathquadraticcurveto{\pgfqpoint{4.962800in}{10.197966in}}{\pgfqpoint{4.966452in}{10.195245in}}%
\pgfpathquadraticcurveto{\pgfqpoint{4.970105in}{10.192524in}}{\pgfqpoint{4.981157in}{10.190065in}}%
\pgfpathlineto{\pgfqpoint{4.985860in}{10.189015in}}%
\pgfpathquadraticcurveto{\pgfqpoint{5.000469in}{10.185887in}}{\pgfqpoint{5.006628in}{10.180182in}}%
\pgfpathquadraticcurveto{\pgfqpoint{5.012787in}{10.174477in}}{\pgfqpoint{5.012787in}{10.164260in}}%
\pgfpathquadraticcurveto{\pgfqpoint{5.012787in}{10.152611in}}{\pgfqpoint{5.003572in}{10.145807in}}%
\pgfpathquadraticcurveto{\pgfqpoint{4.994358in}{10.139028in}}{\pgfqpoint{4.978245in}{10.139028in}}%
\pgfpathquadraticcurveto{\pgfqpoint{4.971537in}{10.139028in}}{\pgfqpoint{4.964256in}{10.140341in}}%
\pgfpathquadraticcurveto{\pgfqpoint{4.956975in}{10.141654in}}{\pgfqpoint{4.948931in}{10.144256in}}%
\pgfpathlineto{\pgfqpoint{4.948931in}{10.158435in}}%
\pgfpathquadraticcurveto{\pgfqpoint{4.956546in}{10.154473in}}{\pgfqpoint{4.963922in}{10.152491in}}%
\pgfpathquadraticcurveto{\pgfqpoint{4.971298in}{10.150534in}}{\pgfqpoint{4.978555in}{10.150534in}}%
\pgfpathquadraticcurveto{\pgfqpoint{4.988247in}{10.150534in}}{\pgfqpoint{4.993451in}{10.153852in}}%
\pgfpathquadraticcurveto{\pgfqpoint{4.998679in}{10.157170in}}{\pgfqpoint{4.998679in}{10.163210in}}%
\pgfpathquadraticcurveto{\pgfqpoint{4.998679in}{10.168795in}}{\pgfqpoint{4.994907in}{10.171779in}}%
\pgfpathquadraticcurveto{\pgfqpoint{4.991159in}{10.174763in}}{\pgfqpoint{4.978388in}{10.177532in}}%
\pgfpathlineto{\pgfqpoint{4.973614in}{10.178654in}}%
\pgfpathquadraticcurveto{\pgfqpoint{4.960866in}{10.181328in}}{\pgfqpoint{4.955185in}{10.186890in}}%
\pgfpathquadraticcurveto{\pgfqpoint{4.949527in}{10.192452in}}{\pgfqpoint{4.949527in}{10.202144in}}%
\pgfpathquadraticcurveto{\pgfqpoint{4.949527in}{10.213937in}}{\pgfqpoint{4.957882in}{10.220334in}}%
\pgfpathquadraticcurveto{\pgfqpoint{4.966237in}{10.226756in}}{\pgfqpoint{4.981611in}{10.226756in}}%
\pgfpathquadraticcurveto{\pgfqpoint{4.989202in}{10.226756in}}{\pgfqpoint{4.995910in}{10.225634in}}%
\pgfpathquadraticcurveto{\pgfqpoint{5.002641in}{10.224536in}}{\pgfqpoint{5.008299in}{10.222292in}}%
\pgfpathlineto{\pgfqpoint{5.008299in}{10.222292in}}%
\pgfpathclose%
\pgfpathmoveto{\pgfqpoint{5.106101in}{10.186413in}}%
\pgfpathlineto{\pgfqpoint{5.106101in}{10.179705in}}%
\pgfpathlineto{\pgfqpoint{5.042984in}{10.179705in}}%
\pgfpathquadraticcurveto{\pgfqpoint{5.043891in}{10.165525in}}{\pgfqpoint{5.051530in}{10.158101in}}%
\pgfpathquadraticcurveto{\pgfqpoint{5.059169in}{10.150677in}}{\pgfqpoint{5.072824in}{10.150677in}}%
\pgfpathquadraticcurveto{\pgfqpoint{5.080725in}{10.150677in}}{\pgfqpoint{5.088149in}{10.152611in}}%
\pgfpathquadraticcurveto{\pgfqpoint{5.095573in}{10.154544in}}{\pgfqpoint{5.102902in}{10.158435in}}%
\pgfpathlineto{\pgfqpoint{5.102902in}{10.145449in}}%
\pgfpathquadraticcurveto{\pgfqpoint{5.095502in}{10.142322in}}{\pgfqpoint{5.087743in}{10.140675in}}%
\pgfpathquadraticcurveto{\pgfqpoint{5.079985in}{10.139028in}}{\pgfqpoint{5.072012in}{10.139028in}}%
\pgfpathquadraticcurveto{\pgfqpoint{5.052008in}{10.139028in}}{\pgfqpoint{5.040335in}{10.150653in}}%
\pgfpathquadraticcurveto{\pgfqpoint{5.028661in}{10.162302in}}{\pgfqpoint{5.028661in}{10.182164in}}%
\pgfpathquadraticcurveto{\pgfqpoint{5.028661in}{10.202669in}}{\pgfqpoint{5.039738in}{10.214700in}}%
\pgfpathquadraticcurveto{\pgfqpoint{5.050814in}{10.226756in}}{\pgfqpoint{5.069625in}{10.226756in}}%
\pgfpathquadraticcurveto{\pgfqpoint{5.086478in}{10.226756in}}{\pgfqpoint{5.096289in}{10.215894in}}%
\pgfpathquadraticcurveto{\pgfqpoint{5.106101in}{10.205056in}}{\pgfqpoint{5.106101in}{10.186413in}}%
\pgfpathlineto{\pgfqpoint{5.106101in}{10.186413in}}%
\pgfpathclose%
\pgfpathmoveto{\pgfqpoint{5.092375in}{10.190447in}}%
\pgfpathquadraticcurveto{\pgfqpoint{5.092231in}{10.201690in}}{\pgfqpoint{5.086072in}{10.208398in}}%
\pgfpathquadraticcurveto{\pgfqpoint{5.079914in}{10.215130in}}{\pgfqpoint{5.069768in}{10.215130in}}%
\pgfpathquadraticcurveto{\pgfqpoint{5.058286in}{10.215130in}}{\pgfqpoint{5.051387in}{10.208637in}}%
\pgfpathquadraticcurveto{\pgfqpoint{5.044488in}{10.202144in}}{\pgfqpoint{5.043438in}{10.190351in}}%
\pgfpathlineto{\pgfqpoint{5.092375in}{10.190447in}}%
\pgfpathlineto{\pgfqpoint{5.092375in}{10.190447in}}%
\pgfpathclose%
\pgfpathmoveto{\pgfqpoint{5.179124in}{10.203720in}}%
\pgfpathlineto{\pgfqpoint{5.194879in}{10.203720in}}%
\pgfpathlineto{\pgfqpoint{5.194879in}{10.184766in}}%
\pgfpathlineto{\pgfqpoint{5.179124in}{10.184766in}}%
\pgfpathlineto{\pgfqpoint{5.179124in}{10.203720in}}%
\pgfpathlineto{\pgfqpoint{5.179124in}{10.203720in}}%
\pgfpathclose%
\pgfpathmoveto{\pgfqpoint{5.358303in}{10.243991in}}%
\pgfpathlineto{\pgfqpoint{5.358303in}{10.228116in}}%
\pgfpathquadraticcurveto{\pgfqpoint{5.350688in}{10.235206in}}{\pgfqpoint{5.342071in}{10.238691in}}%
\pgfpathquadraticcurveto{\pgfqpoint{5.333453in}{10.242200in}}{\pgfqpoint{5.323761in}{10.242200in}}%
\pgfpathquadraticcurveto{\pgfqpoint{5.304664in}{10.242200in}}{\pgfqpoint{5.294519in}{10.230527in}}%
\pgfpathquadraticcurveto{\pgfqpoint{5.284373in}{10.218854in}}{\pgfqpoint{5.284373in}{10.196773in}}%
\pgfpathquadraticcurveto{\pgfqpoint{5.284373in}{10.174763in}}{\pgfqpoint{5.294519in}{10.163090in}}%
\pgfpathquadraticcurveto{\pgfqpoint{5.304664in}{10.151417in}}{\pgfqpoint{5.323761in}{10.151417in}}%
\pgfpathquadraticcurveto{\pgfqpoint{5.333453in}{10.151417in}}{\pgfqpoint{5.342071in}{10.154926in}}%
\pgfpathquadraticcurveto{\pgfqpoint{5.350688in}{10.158435in}}{\pgfqpoint{5.358303in}{10.165525in}}%
\pgfpathlineto{\pgfqpoint{5.358303in}{10.149770in}}%
\pgfpathquadraticcurveto{\pgfqpoint{5.350402in}{10.144399in}}{\pgfqpoint{5.341546in}{10.141701in}}%
\pgfpathquadraticcurveto{\pgfqpoint{5.332713in}{10.139028in}}{\pgfqpoint{5.322878in}{10.139028in}}%
\pgfpathquadraticcurveto{\pgfqpoint{5.297574in}{10.139028in}}{\pgfqpoint{5.283013in}{10.154496in}}%
\pgfpathquadraticcurveto{\pgfqpoint{5.268475in}{10.169989in}}{\pgfqpoint{5.268475in}{10.196773in}}%
\pgfpathquadraticcurveto{\pgfqpoint{5.268475in}{10.223628in}}{\pgfqpoint{5.283013in}{10.239097in}}%
\pgfpathquadraticcurveto{\pgfqpoint{5.297574in}{10.254590in}}{\pgfqpoint{5.322878in}{10.254590in}}%
\pgfpathquadraticcurveto{\pgfqpoint{5.332856in}{10.254590in}}{\pgfqpoint{5.341689in}{10.251940in}}%
\pgfpathquadraticcurveto{\pgfqpoint{5.350545in}{10.249290in}}{\pgfqpoint{5.358303in}{10.243991in}}%
\pgfpathlineto{\pgfqpoint{5.358303in}{10.243991in}}%
\pgfpathclose%
\pgfpathmoveto{\pgfqpoint{5.429393in}{10.211931in}}%
\pgfpathquadraticcurveto{\pgfqpoint{5.427077in}{10.213268in}}{\pgfqpoint{5.424356in}{10.213889in}}%
\pgfpathquadraticcurveto{\pgfqpoint{5.421635in}{10.214533in}}{\pgfqpoint{5.418364in}{10.214533in}}%
\pgfpathquadraticcurveto{\pgfqpoint{5.406715in}{10.214533in}}{\pgfqpoint{5.400484in}{10.206966in}}%
\pgfpathquadraticcurveto{\pgfqpoint{5.394254in}{10.199399in}}{\pgfqpoint{5.394254in}{10.185219in}}%
\pgfpathlineto{\pgfqpoint{5.394254in}{10.141200in}}%
\pgfpathlineto{\pgfqpoint{5.380456in}{10.141200in}}%
\pgfpathlineto{\pgfqpoint{5.380456in}{10.224750in}}%
\pgfpathlineto{\pgfqpoint{5.394254in}{10.224750in}}%
\pgfpathlineto{\pgfqpoint{5.394254in}{10.211764in}}%
\pgfpathquadraticcurveto{\pgfqpoint{5.398599in}{10.219379in}}{\pgfqpoint{5.405521in}{10.223055in}}%
\pgfpathquadraticcurveto{\pgfqpoint{5.412468in}{10.226756in}}{\pgfqpoint{5.422398in}{10.226756in}}%
\pgfpathquadraticcurveto{\pgfqpoint{5.423807in}{10.226756in}}{\pgfqpoint{5.425526in}{10.226565in}}%
\pgfpathquadraticcurveto{\pgfqpoint{5.427244in}{10.226397in}}{\pgfqpoint{5.429321in}{10.226016in}}%
\pgfpathlineto{\pgfqpoint{5.429393in}{10.211931in}}%
\pgfpathlineto{\pgfqpoint{5.429393in}{10.211931in}}%
\pgfpathclose%
\pgfpathmoveto{\pgfqpoint{5.472791in}{10.215130in}}%
\pgfpathquadraticcurveto{\pgfqpoint{5.461763in}{10.215130in}}{\pgfqpoint{5.455341in}{10.206512in}}%
\pgfpathquadraticcurveto{\pgfqpoint{5.448920in}{10.197895in}}{\pgfqpoint{5.448920in}{10.182904in}}%
\pgfpathquadraticcurveto{\pgfqpoint{5.448920in}{10.167912in}}{\pgfqpoint{5.455293in}{10.159295in}}%
\pgfpathquadraticcurveto{\pgfqpoint{5.461691in}{10.150677in}}{\pgfqpoint{5.472791in}{10.150677in}}%
\pgfpathquadraticcurveto{\pgfqpoint{5.483772in}{10.150677in}}{\pgfqpoint{5.490170in}{10.159318in}}%
\pgfpathquadraticcurveto{\pgfqpoint{5.496591in}{10.167984in}}{\pgfqpoint{5.496591in}{10.182904in}}%
\pgfpathquadraticcurveto{\pgfqpoint{5.496591in}{10.197752in}}{\pgfqpoint{5.490170in}{10.206441in}}%
\pgfpathquadraticcurveto{\pgfqpoint{5.483772in}{10.215130in}}{\pgfqpoint{5.472791in}{10.215130in}}%
\pgfpathlineto{\pgfqpoint{5.472791in}{10.215130in}}%
\pgfpathclose%
\pgfpathmoveto{\pgfqpoint{5.472791in}{10.226756in}}%
\pgfpathquadraticcurveto{\pgfqpoint{5.490695in}{10.226756in}}{\pgfqpoint{5.500912in}{10.215106in}}%
\pgfpathquadraticcurveto{\pgfqpoint{5.511153in}{10.203481in}}{\pgfqpoint{5.511153in}{10.182904in}}%
\pgfpathquadraticcurveto{\pgfqpoint{5.511153in}{10.162398in}}{\pgfqpoint{5.500912in}{10.150701in}}%
\pgfpathquadraticcurveto{\pgfqpoint{5.490695in}{10.139028in}}{\pgfqpoint{5.472791in}{10.139028in}}%
\pgfpathquadraticcurveto{\pgfqpoint{5.454816in}{10.139028in}}{\pgfqpoint{5.444623in}{10.150701in}}%
\pgfpathquadraticcurveto{\pgfqpoint{5.434454in}{10.162398in}}{\pgfqpoint{5.434454in}{10.182904in}}%
\pgfpathquadraticcurveto{\pgfqpoint{5.434454in}{10.203481in}}{\pgfqpoint{5.444623in}{10.215106in}}%
\pgfpathquadraticcurveto{\pgfqpoint{5.454816in}{10.226756in}}{\pgfqpoint{5.472791in}{10.226756in}}%
\pgfpathlineto{\pgfqpoint{5.472791in}{10.226756in}}%
\pgfpathclose%
\pgfpathmoveto{\pgfqpoint{5.587160in}{10.222292in}}%
\pgfpathlineto{\pgfqpoint{5.587160in}{10.209305in}}%
\pgfpathquadraticcurveto{\pgfqpoint{5.581359in}{10.212289in}}{\pgfqpoint{5.575081in}{10.213769in}}%
\pgfpathquadraticcurveto{\pgfqpoint{5.568826in}{10.215273in}}{\pgfqpoint{5.562095in}{10.215273in}}%
\pgfpathquadraticcurveto{\pgfqpoint{5.551878in}{10.215273in}}{\pgfqpoint{5.546769in}{10.212146in}}%
\pgfpathquadraticcurveto{\pgfqpoint{5.541661in}{10.209019in}}{\pgfqpoint{5.541661in}{10.202741in}}%
\pgfpathquadraticcurveto{\pgfqpoint{5.541661in}{10.197966in}}{\pgfqpoint{5.545313in}{10.195245in}}%
\pgfpathquadraticcurveto{\pgfqpoint{5.548965in}{10.192524in}}{\pgfqpoint{5.560018in}{10.190065in}}%
\pgfpathlineto{\pgfqpoint{5.564720in}{10.189015in}}%
\pgfpathquadraticcurveto{\pgfqpoint{5.579330in}{10.185887in}}{\pgfqpoint{5.585489in}{10.180182in}}%
\pgfpathquadraticcurveto{\pgfqpoint{5.591648in}{10.174477in}}{\pgfqpoint{5.591648in}{10.164260in}}%
\pgfpathquadraticcurveto{\pgfqpoint{5.591648in}{10.152611in}}{\pgfqpoint{5.582433in}{10.145807in}}%
\pgfpathquadraticcurveto{\pgfqpoint{5.573219in}{10.139028in}}{\pgfqpoint{5.557105in}{10.139028in}}%
\pgfpathquadraticcurveto{\pgfqpoint{5.550398in}{10.139028in}}{\pgfqpoint{5.543117in}{10.140341in}}%
\pgfpathquadraticcurveto{\pgfqpoint{5.535836in}{10.141654in}}{\pgfqpoint{5.527791in}{10.144256in}}%
\pgfpathlineto{\pgfqpoint{5.527791in}{10.158435in}}%
\pgfpathquadraticcurveto{\pgfqpoint{5.535406in}{10.154473in}}{\pgfqpoint{5.542783in}{10.152491in}}%
\pgfpathquadraticcurveto{\pgfqpoint{5.550159in}{10.150534in}}{\pgfqpoint{5.557416in}{10.150534in}}%
\pgfpathquadraticcurveto{\pgfqpoint{5.567108in}{10.150534in}}{\pgfqpoint{5.572312in}{10.153852in}}%
\pgfpathquadraticcurveto{\pgfqpoint{5.577539in}{10.157170in}}{\pgfqpoint{5.577539in}{10.163210in}}%
\pgfpathquadraticcurveto{\pgfqpoint{5.577539in}{10.168795in}}{\pgfqpoint{5.573768in}{10.171779in}}%
\pgfpathquadraticcurveto{\pgfqpoint{5.570020in}{10.174763in}}{\pgfqpoint{5.557249in}{10.177532in}}%
\pgfpathlineto{\pgfqpoint{5.552474in}{10.178654in}}%
\pgfpathquadraticcurveto{\pgfqpoint{5.539727in}{10.181328in}}{\pgfqpoint{5.534046in}{10.186890in}}%
\pgfpathquadraticcurveto{\pgfqpoint{5.528388in}{10.192452in}}{\pgfqpoint{5.528388in}{10.202144in}}%
\pgfpathquadraticcurveto{\pgfqpoint{5.528388in}{10.213937in}}{\pgfqpoint{5.536743in}{10.220334in}}%
\pgfpathquadraticcurveto{\pgfqpoint{5.545098in}{10.226756in}}{\pgfqpoint{5.560471in}{10.226756in}}%
\pgfpathquadraticcurveto{\pgfqpoint{5.568062in}{10.226756in}}{\pgfqpoint{5.574770in}{10.225634in}}%
\pgfpathquadraticcurveto{\pgfqpoint{5.581502in}{10.224536in}}{\pgfqpoint{5.587160in}{10.222292in}}%
\pgfpathlineto{\pgfqpoint{5.587160in}{10.222292in}}%
\pgfpathclose%
\pgfpathmoveto{\pgfqpoint{5.666747in}{10.222292in}}%
\pgfpathlineto{\pgfqpoint{5.666747in}{10.209305in}}%
\pgfpathquadraticcurveto{\pgfqpoint{5.660947in}{10.212289in}}{\pgfqpoint{5.654668in}{10.213769in}}%
\pgfpathquadraticcurveto{\pgfqpoint{5.648414in}{10.215273in}}{\pgfqpoint{5.641682in}{10.215273in}}%
\pgfpathquadraticcurveto{\pgfqpoint{5.631465in}{10.215273in}}{\pgfqpoint{5.626357in}{10.212146in}}%
\pgfpathquadraticcurveto{\pgfqpoint{5.621248in}{10.209019in}}{\pgfqpoint{5.621248in}{10.202741in}}%
\pgfpathquadraticcurveto{\pgfqpoint{5.621248in}{10.197966in}}{\pgfqpoint{5.624901in}{10.195245in}}%
\pgfpathquadraticcurveto{\pgfqpoint{5.628553in}{10.192524in}}{\pgfqpoint{5.639605in}{10.190065in}}%
\pgfpathlineto{\pgfqpoint{5.644308in}{10.189015in}}%
\pgfpathquadraticcurveto{\pgfqpoint{5.658918in}{10.185887in}}{\pgfqpoint{5.665076in}{10.180182in}}%
\pgfpathquadraticcurveto{\pgfqpoint{5.671235in}{10.174477in}}{\pgfqpoint{5.671235in}{10.164260in}}%
\pgfpathquadraticcurveto{\pgfqpoint{5.671235in}{10.152611in}}{\pgfqpoint{5.662021in}{10.145807in}}%
\pgfpathquadraticcurveto{\pgfqpoint{5.652806in}{10.139028in}}{\pgfqpoint{5.636693in}{10.139028in}}%
\pgfpathquadraticcurveto{\pgfqpoint{5.629985in}{10.139028in}}{\pgfqpoint{5.622704in}{10.140341in}}%
\pgfpathquadraticcurveto{\pgfqpoint{5.615424in}{10.141654in}}{\pgfqpoint{5.607379in}{10.144256in}}%
\pgfpathlineto{\pgfqpoint{5.607379in}{10.158435in}}%
\pgfpathquadraticcurveto{\pgfqpoint{5.614994in}{10.154473in}}{\pgfqpoint{5.622370in}{10.152491in}}%
\pgfpathquadraticcurveto{\pgfqpoint{5.629747in}{10.150534in}}{\pgfqpoint{5.637003in}{10.150534in}}%
\pgfpathquadraticcurveto{\pgfqpoint{5.646695in}{10.150534in}}{\pgfqpoint{5.651899in}{10.153852in}}%
\pgfpathquadraticcurveto{\pgfqpoint{5.657127in}{10.157170in}}{\pgfqpoint{5.657127in}{10.163210in}}%
\pgfpathquadraticcurveto{\pgfqpoint{5.657127in}{10.168795in}}{\pgfqpoint{5.653355in}{10.171779in}}%
\pgfpathquadraticcurveto{\pgfqpoint{5.649608in}{10.174763in}}{\pgfqpoint{5.636836in}{10.177532in}}%
\pgfpathlineto{\pgfqpoint{5.632062in}{10.178654in}}%
\pgfpathquadraticcurveto{\pgfqpoint{5.619315in}{10.181328in}}{\pgfqpoint{5.613633in}{10.186890in}}%
\pgfpathquadraticcurveto{\pgfqpoint{5.607976in}{10.192452in}}{\pgfqpoint{5.607976in}{10.202144in}}%
\pgfpathquadraticcurveto{\pgfqpoint{5.607976in}{10.213937in}}{\pgfqpoint{5.616331in}{10.220334in}}%
\pgfpathquadraticcurveto{\pgfqpoint{5.624686in}{10.226756in}}{\pgfqpoint{5.640059in}{10.226756in}}%
\pgfpathquadraticcurveto{\pgfqpoint{5.647650in}{10.226756in}}{\pgfqpoint{5.654358in}{10.225634in}}%
\pgfpathquadraticcurveto{\pgfqpoint{5.661090in}{10.224536in}}{\pgfqpoint{5.666747in}{10.222292in}}%
\pgfpathlineto{\pgfqpoint{5.666747in}{10.222292in}}%
\pgfpathclose%
\pgfpathmoveto{\pgfqpoint{5.696587in}{10.160154in}}%
\pgfpathlineto{\pgfqpoint{5.712318in}{10.160154in}}%
\pgfpathlineto{\pgfqpoint{5.712318in}{10.141200in}}%
\pgfpathlineto{\pgfqpoint{5.696587in}{10.141200in}}%
\pgfpathlineto{\pgfqpoint{5.696587in}{10.160154in}}%
\pgfpathlineto{\pgfqpoint{5.696587in}{10.160154in}}%
\pgfpathclose%
\pgfpathmoveto{\pgfqpoint{5.696587in}{10.220191in}}%
\pgfpathlineto{\pgfqpoint{5.712318in}{10.220191in}}%
\pgfpathlineto{\pgfqpoint{5.712318in}{10.201261in}}%
\pgfpathlineto{\pgfqpoint{5.696587in}{10.201261in}}%
\pgfpathlineto{\pgfqpoint{5.696587in}{10.220191in}}%
\pgfpathlineto{\pgfqpoint{5.696587in}{10.220191in}}%
\pgfpathclose%
\pgfpathmoveto{\pgfqpoint{5.835400in}{10.257287in}}%
\pgfpathlineto{\pgfqpoint{5.835400in}{10.245853in}}%
\pgfpathlineto{\pgfqpoint{5.822270in}{10.245853in}}%
\pgfpathquadraticcurveto{\pgfqpoint{5.814894in}{10.245853in}}{\pgfqpoint{5.812006in}{10.242869in}}%
\pgfpathquadraticcurveto{\pgfqpoint{5.809141in}{10.239885in}}{\pgfqpoint{5.809141in}{10.232127in}}%
\pgfpathlineto{\pgfqpoint{5.809141in}{10.224750in}}%
\pgfpathlineto{\pgfqpoint{5.831747in}{10.224750in}}%
\pgfpathlineto{\pgfqpoint{5.831747in}{10.214080in}}%
\pgfpathlineto{\pgfqpoint{5.809141in}{10.214080in}}%
\pgfpathlineto{\pgfqpoint{5.809141in}{10.141200in}}%
\pgfpathlineto{\pgfqpoint{5.795343in}{10.141200in}}%
\pgfpathlineto{\pgfqpoint{5.795343in}{10.214080in}}%
\pgfpathlineto{\pgfqpoint{5.782214in}{10.214080in}}%
\pgfpathlineto{\pgfqpoint{5.782214in}{10.224750in}}%
\pgfpathlineto{\pgfqpoint{5.795343in}{10.224750in}}%
\pgfpathlineto{\pgfqpoint{5.795343in}{10.230575in}}%
\pgfpathquadraticcurveto{\pgfqpoint{5.795343in}{10.244516in}}{\pgfqpoint{5.801836in}{10.250890in}}%
\pgfpathquadraticcurveto{\pgfqpoint{5.808329in}{10.257287in}}{\pgfqpoint{5.822414in}{10.257287in}}%
\pgfpathlineto{\pgfqpoint{5.835400in}{10.257287in}}%
\pgfpathlineto{\pgfqpoint{5.835400in}{10.257287in}}%
\pgfpathclose%
\pgfpathmoveto{\pgfqpoint{5.895293in}{10.211931in}}%
\pgfpathquadraticcurveto{\pgfqpoint{5.892978in}{10.213268in}}{\pgfqpoint{5.890257in}{10.213889in}}%
\pgfpathquadraticcurveto{\pgfqpoint{5.887535in}{10.214533in}}{\pgfqpoint{5.884265in}{10.214533in}}%
\pgfpathquadraticcurveto{\pgfqpoint{5.872615in}{10.214533in}}{\pgfqpoint{5.866385in}{10.206966in}}%
\pgfpathquadraticcurveto{\pgfqpoint{5.860155in}{10.199399in}}{\pgfqpoint{5.860155in}{10.185219in}}%
\pgfpathlineto{\pgfqpoint{5.860155in}{10.141200in}}%
\pgfpathlineto{\pgfqpoint{5.846357in}{10.141200in}}%
\pgfpathlineto{\pgfqpoint{5.846357in}{10.224750in}}%
\pgfpathlineto{\pgfqpoint{5.860155in}{10.224750in}}%
\pgfpathlineto{\pgfqpoint{5.860155in}{10.211764in}}%
\pgfpathquadraticcurveto{\pgfqpoint{5.864499in}{10.219379in}}{\pgfqpoint{5.871422in}{10.223055in}}%
\pgfpathquadraticcurveto{\pgfqpoint{5.878368in}{10.226756in}}{\pgfqpoint{5.888299in}{10.226756in}}%
\pgfpathquadraticcurveto{\pgfqpoint{5.889707in}{10.226756in}}{\pgfqpoint{5.891426in}{10.226565in}}%
\pgfpathquadraticcurveto{\pgfqpoint{5.893145in}{10.226397in}}{\pgfqpoint{5.895222in}{10.226016in}}%
\pgfpathlineto{\pgfqpoint{5.895293in}{10.211931in}}%
\pgfpathlineto{\pgfqpoint{5.895293in}{10.211931in}}%
\pgfpathclose%
\pgfpathmoveto{\pgfqpoint{5.947668in}{10.183190in}}%
\pgfpathquadraticcurveto{\pgfqpoint{5.931029in}{10.183190in}}{\pgfqpoint{5.924608in}{10.179394in}}%
\pgfpathquadraticcurveto{\pgfqpoint{5.918186in}{10.175599in}}{\pgfqpoint{5.918186in}{10.166408in}}%
\pgfpathquadraticcurveto{\pgfqpoint{5.918186in}{10.159104in}}{\pgfqpoint{5.923008in}{10.154807in}}%
\pgfpathquadraticcurveto{\pgfqpoint{5.927830in}{10.150534in}}{\pgfqpoint{5.936090in}{10.150534in}}%
\pgfpathquadraticcurveto{\pgfqpoint{5.947524in}{10.150534in}}{\pgfqpoint{5.954423in}{10.158626in}}%
\pgfpathquadraticcurveto{\pgfqpoint{5.961322in}{10.166719in}}{\pgfqpoint{5.961322in}{10.180134in}}%
\pgfpathlineto{\pgfqpoint{5.961322in}{10.183190in}}%
\pgfpathlineto{\pgfqpoint{5.947668in}{10.183190in}}%
\pgfpathlineto{\pgfqpoint{5.947668in}{10.183190in}}%
\pgfpathclose%
\pgfpathmoveto{\pgfqpoint{5.975048in}{10.188871in}}%
\pgfpathlineto{\pgfqpoint{5.975048in}{10.141200in}}%
\pgfpathlineto{\pgfqpoint{5.961322in}{10.141200in}}%
\pgfpathlineto{\pgfqpoint{5.961322in}{10.153876in}}%
\pgfpathquadraticcurveto{\pgfqpoint{5.956619in}{10.146285in}}{\pgfqpoint{5.949601in}{10.142656in}}%
\pgfpathquadraticcurveto{\pgfqpoint{5.942583in}{10.139028in}}{\pgfqpoint{5.932437in}{10.139028in}}%
\pgfpathquadraticcurveto{\pgfqpoint{5.919618in}{10.139028in}}{\pgfqpoint{5.912027in}{10.146237in}}%
\pgfpathquadraticcurveto{\pgfqpoint{5.904460in}{10.153446in}}{\pgfqpoint{5.904460in}{10.165525in}}%
\pgfpathquadraticcurveto{\pgfqpoint{5.904460in}{10.179609in}}{\pgfqpoint{5.913889in}{10.186771in}}%
\pgfpathquadraticcurveto{\pgfqpoint{5.923342in}{10.193932in}}{\pgfqpoint{5.942058in}{10.193932in}}%
\pgfpathlineto{\pgfqpoint{5.961322in}{10.193932in}}%
\pgfpathlineto{\pgfqpoint{5.961322in}{10.195293in}}%
\pgfpathquadraticcurveto{\pgfqpoint{5.961322in}{10.204770in}}{\pgfqpoint{5.955092in}{10.209950in}}%
\pgfpathquadraticcurveto{\pgfqpoint{5.948861in}{10.215130in}}{\pgfqpoint{5.937594in}{10.215130in}}%
\pgfpathquadraticcurveto{\pgfqpoint{5.930432in}{10.215130in}}{\pgfqpoint{5.923629in}{10.213411in}}%
\pgfpathquadraticcurveto{\pgfqpoint{5.916849in}{10.211693in}}{\pgfqpoint{5.910595in}{10.208255in}}%
\pgfpathlineto{\pgfqpoint{5.910595in}{10.220955in}}%
\pgfpathquadraticcurveto{\pgfqpoint{5.918115in}{10.223867in}}{\pgfqpoint{5.925204in}{10.225299in}}%
\pgfpathquadraticcurveto{\pgfqpoint{5.932294in}{10.226756in}}{\pgfqpoint{5.939002in}{10.226756in}}%
\pgfpathquadraticcurveto{\pgfqpoint{5.957145in}{10.226756in}}{\pgfqpoint{5.966096in}{10.217350in}}%
\pgfpathquadraticcurveto{\pgfqpoint{5.975048in}{10.207969in}}{\pgfqpoint{5.975048in}{10.188871in}}%
\pgfpathlineto{\pgfqpoint{5.975048in}{10.188871in}}%
\pgfpathclose%
\pgfpathmoveto{\pgfqpoint{6.058288in}{10.183954in}}%
\pgfpathquadraticcurveto{\pgfqpoint{6.058288in}{10.198874in}}{\pgfqpoint{6.052129in}{10.207062in}}%
\pgfpathquadraticcurveto{\pgfqpoint{6.045994in}{10.215273in}}{\pgfqpoint{6.034870in}{10.215273in}}%
\pgfpathquadraticcurveto{\pgfqpoint{6.023842in}{10.215273in}}{\pgfqpoint{6.017683in}{10.207062in}}%
\pgfpathquadraticcurveto{\pgfqpoint{6.011524in}{10.198874in}}{\pgfqpoint{6.011524in}{10.183954in}}%
\pgfpathquadraticcurveto{\pgfqpoint{6.011524in}{10.169106in}}{\pgfqpoint{6.017683in}{10.160894in}}%
\pgfpathquadraticcurveto{\pgfqpoint{6.023842in}{10.152682in}}{\pgfqpoint{6.034870in}{10.152682in}}%
\pgfpathquadraticcurveto{\pgfqpoint{6.045994in}{10.152682in}}{\pgfqpoint{6.052129in}{10.160894in}}%
\pgfpathquadraticcurveto{\pgfqpoint{6.058288in}{10.169106in}}{\pgfqpoint{6.058288in}{10.183954in}}%
\pgfpathlineto{\pgfqpoint{6.058288in}{10.183954in}}%
\pgfpathclose%
\pgfpathmoveto{\pgfqpoint{6.072014in}{10.151560in}}%
\pgfpathquadraticcurveto{\pgfqpoint{6.072014in}{10.130243in}}{\pgfqpoint{6.062537in}{10.119835in}}%
\pgfpathquadraticcurveto{\pgfqpoint{6.053084in}{10.109427in}}{\pgfqpoint{6.033533in}{10.109427in}}%
\pgfpathquadraticcurveto{\pgfqpoint{6.026300in}{10.109427in}}{\pgfqpoint{6.019879in}{10.110501in}}%
\pgfpathquadraticcurveto{\pgfqpoint{6.013457in}{10.111575in}}{\pgfqpoint{6.007418in}{10.113819in}}%
\pgfpathlineto{\pgfqpoint{6.007418in}{10.127164in}}%
\pgfpathquadraticcurveto{\pgfqpoint{6.013457in}{10.123893in}}{\pgfqpoint{6.019354in}{10.122341in}}%
\pgfpathquadraticcurveto{\pgfqpoint{6.025250in}{10.120766in}}{\pgfqpoint{6.031361in}{10.120766in}}%
\pgfpathquadraticcurveto{\pgfqpoint{6.044872in}{10.120766in}}{\pgfqpoint{6.051580in}{10.127808in}}%
\pgfpathquadraticcurveto{\pgfqpoint{6.058288in}{10.134850in}}{\pgfqpoint{6.058288in}{10.149101in}}%
\pgfpathlineto{\pgfqpoint{6.058288in}{10.155905in}}%
\pgfpathquadraticcurveto{\pgfqpoint{6.054039in}{10.148505in}}{\pgfqpoint{6.047403in}{10.144852in}}%
\pgfpathquadraticcurveto{\pgfqpoint{6.040766in}{10.141200in}}{\pgfqpoint{6.031504in}{10.141200in}}%
\pgfpathquadraticcurveto{\pgfqpoint{6.016155in}{10.141200in}}{\pgfqpoint{6.006750in}{10.152897in}}%
\pgfpathquadraticcurveto{\pgfqpoint{5.997344in}{10.164618in}}{\pgfqpoint{5.997344in}{10.183954in}}%
\pgfpathquadraticcurveto{\pgfqpoint{5.997344in}{10.203338in}}{\pgfqpoint{6.006750in}{10.215035in}}%
\pgfpathquadraticcurveto{\pgfqpoint{6.016155in}{10.226756in}}{\pgfqpoint{6.031504in}{10.226756in}}%
\pgfpathquadraticcurveto{\pgfqpoint{6.040766in}{10.226756in}}{\pgfqpoint{6.047403in}{10.223103in}}%
\pgfpathquadraticcurveto{\pgfqpoint{6.054039in}{10.219451in}}{\pgfqpoint{6.058288in}{10.212075in}}%
\pgfpathlineto{\pgfqpoint{6.058288in}{10.224750in}}%
\pgfpathlineto{\pgfqpoint{6.072014in}{10.224750in}}%
\pgfpathlineto{\pgfqpoint{6.072014in}{10.151560in}}%
\pgfpathlineto{\pgfqpoint{6.072014in}{10.151560in}}%
\pgfpathclose%
\pgfpathmoveto{\pgfqpoint{6.100302in}{10.224750in}}%
\pgfpathlineto{\pgfqpoint{6.114028in}{10.224750in}}%
\pgfpathlineto{\pgfqpoint{6.114028in}{10.141200in}}%
\pgfpathlineto{\pgfqpoint{6.100302in}{10.141200in}}%
\pgfpathlineto{\pgfqpoint{6.100302in}{10.224750in}}%
\pgfpathlineto{\pgfqpoint{6.100302in}{10.224750in}}%
\pgfpathclose%
\pgfpathmoveto{\pgfqpoint{6.100302in}{10.257287in}}%
\pgfpathlineto{\pgfqpoint{6.114028in}{10.257287in}}%
\pgfpathlineto{\pgfqpoint{6.114028in}{10.239885in}}%
\pgfpathlineto{\pgfqpoint{6.100302in}{10.239885in}}%
\pgfpathlineto{\pgfqpoint{6.100302in}{10.257287in}}%
\pgfpathlineto{\pgfqpoint{6.100302in}{10.257287in}}%
\pgfpathclose%
\pgfpathmoveto{\pgfqpoint{6.142746in}{10.257287in}}%
\pgfpathlineto{\pgfqpoint{6.156472in}{10.257287in}}%
\pgfpathlineto{\pgfqpoint{6.156472in}{10.141200in}}%
\pgfpathlineto{\pgfqpoint{6.142746in}{10.141200in}}%
\pgfpathlineto{\pgfqpoint{6.142746in}{10.257287in}}%
\pgfpathlineto{\pgfqpoint{6.142746in}{10.257287in}}%
\pgfpathclose%
\pgfpathmoveto{\pgfqpoint{6.256661in}{10.186413in}}%
\pgfpathlineto{\pgfqpoint{6.256661in}{10.179705in}}%
\pgfpathlineto{\pgfqpoint{6.193544in}{10.179705in}}%
\pgfpathquadraticcurveto{\pgfqpoint{6.194451in}{10.165525in}}{\pgfqpoint{6.202090in}{10.158101in}}%
\pgfpathquadraticcurveto{\pgfqpoint{6.209729in}{10.150677in}}{\pgfqpoint{6.223384in}{10.150677in}}%
\pgfpathquadraticcurveto{\pgfqpoint{6.231285in}{10.150677in}}{\pgfqpoint{6.238709in}{10.152611in}}%
\pgfpathquadraticcurveto{\pgfqpoint{6.246133in}{10.154544in}}{\pgfqpoint{6.253462in}{10.158435in}}%
\pgfpathlineto{\pgfqpoint{6.253462in}{10.145449in}}%
\pgfpathquadraticcurveto{\pgfqpoint{6.246062in}{10.142322in}}{\pgfqpoint{6.238303in}{10.140675in}}%
\pgfpathquadraticcurveto{\pgfqpoint{6.230545in}{10.139028in}}{\pgfqpoint{6.222572in}{10.139028in}}%
\pgfpathquadraticcurveto{\pgfqpoint{6.202568in}{10.139028in}}{\pgfqpoint{6.190895in}{10.150653in}}%
\pgfpathquadraticcurveto{\pgfqpoint{6.179221in}{10.162302in}}{\pgfqpoint{6.179221in}{10.182164in}}%
\pgfpathquadraticcurveto{\pgfqpoint{6.179221in}{10.202669in}}{\pgfqpoint{6.190298in}{10.214700in}}%
\pgfpathquadraticcurveto{\pgfqpoint{6.201374in}{10.226756in}}{\pgfqpoint{6.220185in}{10.226756in}}%
\pgfpathquadraticcurveto{\pgfqpoint{6.237038in}{10.226756in}}{\pgfqpoint{6.246849in}{10.215894in}}%
\pgfpathquadraticcurveto{\pgfqpoint{6.256661in}{10.205056in}}{\pgfqpoint{6.256661in}{10.186413in}}%
\pgfpathlineto{\pgfqpoint{6.256661in}{10.186413in}}%
\pgfpathclose%
\pgfpathmoveto{\pgfqpoint{6.242934in}{10.190447in}}%
\pgfpathquadraticcurveto{\pgfqpoint{6.242791in}{10.201690in}}{\pgfqpoint{6.236632in}{10.208398in}}%
\pgfpathquadraticcurveto{\pgfqpoint{6.230474in}{10.215130in}}{\pgfqpoint{6.220328in}{10.215130in}}%
\pgfpathquadraticcurveto{\pgfqpoint{6.208846in}{10.215130in}}{\pgfqpoint{6.201947in}{10.208637in}}%
\pgfpathquadraticcurveto{\pgfqpoint{6.195048in}{10.202144in}}{\pgfqpoint{6.193998in}{10.190351in}}%
\pgfpathlineto{\pgfqpoint{6.242934in}{10.190447in}}%
\pgfpathlineto{\pgfqpoint{6.242934in}{10.190447in}}%
\pgfpathclose%
\pgfpathmoveto{\pgfqpoint{6.327750in}{10.224750in}}%
\pgfpathlineto{\pgfqpoint{6.341476in}{10.224750in}}%
\pgfpathlineto{\pgfqpoint{6.341476in}{10.141200in}}%
\pgfpathlineto{\pgfqpoint{6.327750in}{10.141200in}}%
\pgfpathlineto{\pgfqpoint{6.327750in}{10.224750in}}%
\pgfpathlineto{\pgfqpoint{6.327750in}{10.224750in}}%
\pgfpathclose%
\pgfpathmoveto{\pgfqpoint{6.327750in}{10.257287in}}%
\pgfpathlineto{\pgfqpoint{6.341476in}{10.257287in}}%
\pgfpathlineto{\pgfqpoint{6.341476in}{10.239885in}}%
\pgfpathlineto{\pgfqpoint{6.327750in}{10.239885in}}%
\pgfpathlineto{\pgfqpoint{6.327750in}{10.257287in}}%
\pgfpathlineto{\pgfqpoint{6.327750in}{10.257287in}}%
\pgfpathclose%
\pgfpathmoveto{\pgfqpoint{6.439660in}{10.191641in}}%
\pgfpathlineto{\pgfqpoint{6.439660in}{10.141200in}}%
\pgfpathlineto{\pgfqpoint{6.425934in}{10.141200in}}%
\pgfpathlineto{\pgfqpoint{6.425934in}{10.191187in}}%
\pgfpathquadraticcurveto{\pgfqpoint{6.425934in}{10.203051in}}{\pgfqpoint{6.421303in}{10.208924in}}%
\pgfpathquadraticcurveto{\pgfqpoint{6.416671in}{10.214820in}}{\pgfqpoint{6.407433in}{10.214820in}}%
\pgfpathquadraticcurveto{\pgfqpoint{6.396309in}{10.214820in}}{\pgfqpoint{6.389888in}{10.207730in}}%
\pgfpathquadraticcurveto{\pgfqpoint{6.383466in}{10.200664in}}{\pgfqpoint{6.383466in}{10.188418in}}%
\pgfpathlineto{\pgfqpoint{6.383466in}{10.141200in}}%
\pgfpathlineto{\pgfqpoint{6.369668in}{10.141200in}}%
\pgfpathlineto{\pgfqpoint{6.369668in}{10.224750in}}%
\pgfpathlineto{\pgfqpoint{6.383466in}{10.224750in}}%
\pgfpathlineto{\pgfqpoint{6.383466in}{10.211764in}}%
\pgfpathquadraticcurveto{\pgfqpoint{6.388408in}{10.219308in}}{\pgfqpoint{6.395068in}{10.223032in}}%
\pgfpathquadraticcurveto{\pgfqpoint{6.401752in}{10.226756in}}{\pgfqpoint{6.410489in}{10.226756in}}%
\pgfpathquadraticcurveto{\pgfqpoint{6.424883in}{10.226756in}}{\pgfqpoint{6.432260in}{10.217851in}}%
\pgfpathquadraticcurveto{\pgfqpoint{6.439660in}{10.208947in}}{\pgfqpoint{6.439660in}{10.191641in}}%
\pgfpathlineto{\pgfqpoint{6.439660in}{10.191641in}}%
\pgfpathclose%
\pgfpathmoveto{\pgfqpoint{6.509317in}{10.257287in}}%
\pgfpathlineto{\pgfqpoint{6.509317in}{10.245853in}}%
\pgfpathlineto{\pgfqpoint{6.496187in}{10.245853in}}%
\pgfpathquadraticcurveto{\pgfqpoint{6.488811in}{10.245853in}}{\pgfqpoint{6.485923in}{10.242869in}}%
\pgfpathquadraticcurveto{\pgfqpoint{6.483058in}{10.239885in}}{\pgfqpoint{6.483058in}{10.232127in}}%
\pgfpathlineto{\pgfqpoint{6.483058in}{10.224750in}}%
\pgfpathlineto{\pgfqpoint{6.505664in}{10.224750in}}%
\pgfpathlineto{\pgfqpoint{6.505664in}{10.214080in}}%
\pgfpathlineto{\pgfqpoint{6.483058in}{10.214080in}}%
\pgfpathlineto{\pgfqpoint{6.483058in}{10.141200in}}%
\pgfpathlineto{\pgfqpoint{6.469260in}{10.141200in}}%
\pgfpathlineto{\pgfqpoint{6.469260in}{10.214080in}}%
\pgfpathlineto{\pgfqpoint{6.456131in}{10.214080in}}%
\pgfpathlineto{\pgfqpoint{6.456131in}{10.224750in}}%
\pgfpathlineto{\pgfqpoint{6.469260in}{10.224750in}}%
\pgfpathlineto{\pgfqpoint{6.469260in}{10.230575in}}%
\pgfpathquadraticcurveto{\pgfqpoint{6.469260in}{10.244516in}}{\pgfqpoint{6.475753in}{10.250890in}}%
\pgfpathquadraticcurveto{\pgfqpoint{6.482247in}{10.257287in}}{\pgfqpoint{6.496331in}{10.257287in}}%
\pgfpathlineto{\pgfqpoint{6.509317in}{10.257287in}}%
\pgfpathlineto{\pgfqpoint{6.509317in}{10.257287in}}%
\pgfpathclose%
\pgfpathmoveto{\pgfqpoint{6.592270in}{10.186413in}}%
\pgfpathlineto{\pgfqpoint{6.592270in}{10.179705in}}%
\pgfpathlineto{\pgfqpoint{6.529154in}{10.179705in}}%
\pgfpathquadraticcurveto{\pgfqpoint{6.530061in}{10.165525in}}{\pgfqpoint{6.537700in}{10.158101in}}%
\pgfpathquadraticcurveto{\pgfqpoint{6.545339in}{10.150677in}}{\pgfqpoint{6.558993in}{10.150677in}}%
\pgfpathquadraticcurveto{\pgfqpoint{6.566895in}{10.150677in}}{\pgfqpoint{6.574319in}{10.152611in}}%
\pgfpathquadraticcurveto{\pgfqpoint{6.581743in}{10.154544in}}{\pgfqpoint{6.589072in}{10.158435in}}%
\pgfpathlineto{\pgfqpoint{6.589072in}{10.145449in}}%
\pgfpathquadraticcurveto{\pgfqpoint{6.581671in}{10.142322in}}{\pgfqpoint{6.573913in}{10.140675in}}%
\pgfpathquadraticcurveto{\pgfqpoint{6.566155in}{10.139028in}}{\pgfqpoint{6.558182in}{10.139028in}}%
\pgfpathquadraticcurveto{\pgfqpoint{6.538178in}{10.139028in}}{\pgfqpoint{6.526504in}{10.150653in}}%
\pgfpathquadraticcurveto{\pgfqpoint{6.514831in}{10.162302in}}{\pgfqpoint{6.514831in}{10.182164in}}%
\pgfpathquadraticcurveto{\pgfqpoint{6.514831in}{10.202669in}}{\pgfqpoint{6.525908in}{10.214700in}}%
\pgfpathquadraticcurveto{\pgfqpoint{6.536984in}{10.226756in}}{\pgfqpoint{6.555795in}{10.226756in}}%
\pgfpathquadraticcurveto{\pgfqpoint{6.572648in}{10.226756in}}{\pgfqpoint{6.582459in}{10.215894in}}%
\pgfpathquadraticcurveto{\pgfqpoint{6.592270in}{10.205056in}}{\pgfqpoint{6.592270in}{10.186413in}}%
\pgfpathlineto{\pgfqpoint{6.592270in}{10.186413in}}%
\pgfpathclose%
\pgfpathmoveto{\pgfqpoint{6.578544in}{10.190447in}}%
\pgfpathquadraticcurveto{\pgfqpoint{6.578401in}{10.201690in}}{\pgfqpoint{6.572242in}{10.208398in}}%
\pgfpathquadraticcurveto{\pgfqpoint{6.566083in}{10.215130in}}{\pgfqpoint{6.555938in}{10.215130in}}%
\pgfpathquadraticcurveto{\pgfqpoint{6.544456in}{10.215130in}}{\pgfqpoint{6.537557in}{10.208637in}}%
\pgfpathquadraticcurveto{\pgfqpoint{6.530658in}{10.202144in}}{\pgfqpoint{6.529608in}{10.190351in}}%
\pgfpathlineto{\pgfqpoint{6.578544in}{10.190447in}}%
\pgfpathlineto{\pgfqpoint{6.578544in}{10.190447in}}%
\pgfpathclose%
\pgfpathmoveto{\pgfqpoint{6.663217in}{10.211931in}}%
\pgfpathquadraticcurveto{\pgfqpoint{6.660901in}{10.213268in}}{\pgfqpoint{6.658180in}{10.213889in}}%
\pgfpathquadraticcurveto{\pgfqpoint{6.655458in}{10.214533in}}{\pgfqpoint{6.652188in}{10.214533in}}%
\pgfpathquadraticcurveto{\pgfqpoint{6.640539in}{10.214533in}}{\pgfqpoint{6.634308in}{10.206966in}}%
\pgfpathquadraticcurveto{\pgfqpoint{6.628078in}{10.199399in}}{\pgfqpoint{6.628078in}{10.185219in}}%
\pgfpathlineto{\pgfqpoint{6.628078in}{10.141200in}}%
\pgfpathlineto{\pgfqpoint{6.614280in}{10.141200in}}%
\pgfpathlineto{\pgfqpoint{6.614280in}{10.224750in}}%
\pgfpathlineto{\pgfqpoint{6.628078in}{10.224750in}}%
\pgfpathlineto{\pgfqpoint{6.628078in}{10.211764in}}%
\pgfpathquadraticcurveto{\pgfqpoint{6.632422in}{10.219379in}}{\pgfqpoint{6.639345in}{10.223055in}}%
\pgfpathquadraticcurveto{\pgfqpoint{6.646292in}{10.226756in}}{\pgfqpoint{6.656222in}{10.226756in}}%
\pgfpathquadraticcurveto{\pgfqpoint{6.657631in}{10.226756in}}{\pgfqpoint{6.659349in}{10.226565in}}%
\pgfpathquadraticcurveto{\pgfqpoint{6.661068in}{10.226397in}}{\pgfqpoint{6.663145in}{10.226016in}}%
\pgfpathlineto{\pgfqpoint{6.663217in}{10.211931in}}%
\pgfpathlineto{\pgfqpoint{6.663217in}{10.211931in}}%
\pgfpathclose%
\pgfpathmoveto{\pgfqpoint{6.745717in}{10.186413in}}%
\pgfpathlineto{\pgfqpoint{6.745717in}{10.179705in}}%
\pgfpathlineto{\pgfqpoint{6.682600in}{10.179705in}}%
\pgfpathquadraticcurveto{\pgfqpoint{6.683507in}{10.165525in}}{\pgfqpoint{6.691146in}{10.158101in}}%
\pgfpathquadraticcurveto{\pgfqpoint{6.698785in}{10.150677in}}{\pgfqpoint{6.712440in}{10.150677in}}%
\pgfpathquadraticcurveto{\pgfqpoint{6.720341in}{10.150677in}}{\pgfqpoint{6.727765in}{10.152611in}}%
\pgfpathquadraticcurveto{\pgfqpoint{6.735189in}{10.154544in}}{\pgfqpoint{6.742518in}{10.158435in}}%
\pgfpathlineto{\pgfqpoint{6.742518in}{10.145449in}}%
\pgfpathquadraticcurveto{\pgfqpoint{6.735118in}{10.142322in}}{\pgfqpoint{6.727359in}{10.140675in}}%
\pgfpathquadraticcurveto{\pgfqpoint{6.719601in}{10.139028in}}{\pgfqpoint{6.711628in}{10.139028in}}%
\pgfpathquadraticcurveto{\pgfqpoint{6.691624in}{10.139028in}}{\pgfqpoint{6.679951in}{10.150653in}}%
\pgfpathquadraticcurveto{\pgfqpoint{6.668277in}{10.162302in}}{\pgfqpoint{6.668277in}{10.182164in}}%
\pgfpathquadraticcurveto{\pgfqpoint{6.668277in}{10.202669in}}{\pgfqpoint{6.679354in}{10.214700in}}%
\pgfpathquadraticcurveto{\pgfqpoint{6.690430in}{10.226756in}}{\pgfqpoint{6.709241in}{10.226756in}}%
\pgfpathquadraticcurveto{\pgfqpoint{6.726094in}{10.226756in}}{\pgfqpoint{6.735905in}{10.215894in}}%
\pgfpathquadraticcurveto{\pgfqpoint{6.745717in}{10.205056in}}{\pgfqpoint{6.745717in}{10.186413in}}%
\pgfpathlineto{\pgfqpoint{6.745717in}{10.186413in}}%
\pgfpathclose%
\pgfpathmoveto{\pgfqpoint{6.731990in}{10.190447in}}%
\pgfpathquadraticcurveto{\pgfqpoint{6.731847in}{10.201690in}}{\pgfqpoint{6.725688in}{10.208398in}}%
\pgfpathquadraticcurveto{\pgfqpoint{6.719530in}{10.215130in}}{\pgfqpoint{6.709384in}{10.215130in}}%
\pgfpathquadraticcurveto{\pgfqpoint{6.697902in}{10.215130in}}{\pgfqpoint{6.691003in}{10.208637in}}%
\pgfpathquadraticcurveto{\pgfqpoint{6.684104in}{10.202144in}}{\pgfqpoint{6.683054in}{10.190351in}}%
\pgfpathlineto{\pgfqpoint{6.731990in}{10.190447in}}%
\pgfpathlineto{\pgfqpoint{6.731990in}{10.190447in}}%
\pgfpathclose%
\pgfpathmoveto{\pgfqpoint{6.837717in}{10.191641in}}%
\pgfpathlineto{\pgfqpoint{6.837717in}{10.141200in}}%
\pgfpathlineto{\pgfqpoint{6.823991in}{10.141200in}}%
\pgfpathlineto{\pgfqpoint{6.823991in}{10.191187in}}%
\pgfpathquadraticcurveto{\pgfqpoint{6.823991in}{10.203051in}}{\pgfqpoint{6.819360in}{10.208924in}}%
\pgfpathquadraticcurveto{\pgfqpoint{6.814729in}{10.214820in}}{\pgfqpoint{6.805491in}{10.214820in}}%
\pgfpathquadraticcurveto{\pgfqpoint{6.794367in}{10.214820in}}{\pgfqpoint{6.787945in}{10.207730in}}%
\pgfpathquadraticcurveto{\pgfqpoint{6.781524in}{10.200664in}}{\pgfqpoint{6.781524in}{10.188418in}}%
\pgfpathlineto{\pgfqpoint{6.781524in}{10.141200in}}%
\pgfpathlineto{\pgfqpoint{6.767726in}{10.141200in}}%
\pgfpathlineto{\pgfqpoint{6.767726in}{10.224750in}}%
\pgfpathlineto{\pgfqpoint{6.781524in}{10.224750in}}%
\pgfpathlineto{\pgfqpoint{6.781524in}{10.211764in}}%
\pgfpathquadraticcurveto{\pgfqpoint{6.786465in}{10.219308in}}{\pgfqpoint{6.793125in}{10.223032in}}%
\pgfpathquadraticcurveto{\pgfqpoint{6.799809in}{10.226756in}}{\pgfqpoint{6.808546in}{10.226756in}}%
\pgfpathquadraticcurveto{\pgfqpoint{6.822941in}{10.226756in}}{\pgfqpoint{6.830317in}{10.217851in}}%
\pgfpathquadraticcurveto{\pgfqpoint{6.837717in}{10.208947in}}{\pgfqpoint{6.837717in}{10.191641in}}%
\pgfpathlineto{\pgfqpoint{6.837717in}{10.191641in}}%
\pgfpathclose%
\pgfpathmoveto{\pgfqpoint{6.925207in}{10.221552in}}%
\pgfpathlineto{\pgfqpoint{6.925207in}{10.208709in}}%
\pgfpathquadraticcurveto{\pgfqpoint{6.919382in}{10.211931in}}{\pgfqpoint{6.913533in}{10.213531in}}%
\pgfpathquadraticcurveto{\pgfqpoint{6.907685in}{10.215130in}}{\pgfqpoint{6.901717in}{10.215130in}}%
\pgfpathquadraticcurveto{\pgfqpoint{6.888349in}{10.215130in}}{\pgfqpoint{6.880949in}{10.206656in}}%
\pgfpathquadraticcurveto{\pgfqpoint{6.873572in}{10.198205in}}{\pgfqpoint{6.873572in}{10.182904in}}%
\pgfpathquadraticcurveto{\pgfqpoint{6.873572in}{10.167602in}}{\pgfqpoint{6.880949in}{10.159128in}}%
\pgfpathquadraticcurveto{\pgfqpoint{6.888349in}{10.150677in}}{\pgfqpoint{6.901717in}{10.150677in}}%
\pgfpathquadraticcurveto{\pgfqpoint{6.907685in}{10.150677in}}{\pgfqpoint{6.913533in}{10.152276in}}%
\pgfpathquadraticcurveto{\pgfqpoint{6.919382in}{10.153876in}}{\pgfqpoint{6.925207in}{10.157098in}}%
\pgfpathlineto{\pgfqpoint{6.925207in}{10.144399in}}%
\pgfpathquadraticcurveto{\pgfqpoint{6.919454in}{10.141725in}}{\pgfqpoint{6.913295in}{10.140388in}}%
\pgfpathquadraticcurveto{\pgfqpoint{6.907160in}{10.139028in}}{\pgfqpoint{6.900213in}{10.139028in}}%
\pgfpathquadraticcurveto{\pgfqpoint{6.881331in}{10.139028in}}{\pgfqpoint{6.870207in}{10.150892in}}%
\pgfpathquadraticcurveto{\pgfqpoint{6.859106in}{10.162756in}}{\pgfqpoint{6.859106in}{10.182904in}}%
\pgfpathquadraticcurveto{\pgfqpoint{6.859106in}{10.203338in}}{\pgfqpoint{6.870326in}{10.215035in}}%
\pgfpathquadraticcurveto{\pgfqpoint{6.881569in}{10.226756in}}{\pgfqpoint{6.901120in}{10.226756in}}%
\pgfpathquadraticcurveto{\pgfqpoint{6.907446in}{10.226756in}}{\pgfqpoint{6.913486in}{10.225443in}}%
\pgfpathquadraticcurveto{\pgfqpoint{6.919525in}{10.224154in}}{\pgfqpoint{6.925207in}{10.221552in}}%
\pgfpathlineto{\pgfqpoint{6.925207in}{10.221552in}}%
\pgfpathclose%
\pgfpathmoveto{\pgfqpoint{7.020549in}{10.186413in}}%
\pgfpathlineto{\pgfqpoint{7.020549in}{10.179705in}}%
\pgfpathlineto{\pgfqpoint{6.957433in}{10.179705in}}%
\pgfpathquadraticcurveto{\pgfqpoint{6.958340in}{10.165525in}}{\pgfqpoint{6.965979in}{10.158101in}}%
\pgfpathquadraticcurveto{\pgfqpoint{6.973618in}{10.150677in}}{\pgfqpoint{6.987273in}{10.150677in}}%
\pgfpathquadraticcurveto{\pgfqpoint{6.995174in}{10.150677in}}{\pgfqpoint{7.002598in}{10.152611in}}%
\pgfpathquadraticcurveto{\pgfqpoint{7.010022in}{10.154544in}}{\pgfqpoint{7.017351in}{10.158435in}}%
\pgfpathlineto{\pgfqpoint{7.017351in}{10.145449in}}%
\pgfpathquadraticcurveto{\pgfqpoint{7.009951in}{10.142322in}}{\pgfqpoint{7.002192in}{10.140675in}}%
\pgfpathquadraticcurveto{\pgfqpoint{6.994434in}{10.139028in}}{\pgfqpoint{6.986461in}{10.139028in}}%
\pgfpathquadraticcurveto{\pgfqpoint{6.966457in}{10.139028in}}{\pgfqpoint{6.954783in}{10.150653in}}%
\pgfpathquadraticcurveto{\pgfqpoint{6.943110in}{10.162302in}}{\pgfqpoint{6.943110in}{10.182164in}}%
\pgfpathquadraticcurveto{\pgfqpoint{6.943110in}{10.202669in}}{\pgfqpoint{6.954187in}{10.214700in}}%
\pgfpathquadraticcurveto{\pgfqpoint{6.965263in}{10.226756in}}{\pgfqpoint{6.984074in}{10.226756in}}%
\pgfpathquadraticcurveto{\pgfqpoint{7.000927in}{10.226756in}}{\pgfqpoint{7.010738in}{10.215894in}}%
\pgfpathquadraticcurveto{\pgfqpoint{7.020549in}{10.205056in}}{\pgfqpoint{7.020549in}{10.186413in}}%
\pgfpathlineto{\pgfqpoint{7.020549in}{10.186413in}}%
\pgfpathclose%
\pgfpathmoveto{\pgfqpoint{7.006823in}{10.190447in}}%
\pgfpathquadraticcurveto{\pgfqpoint{7.006680in}{10.201690in}}{\pgfqpoint{7.000521in}{10.208398in}}%
\pgfpathquadraticcurveto{\pgfqpoint{6.994362in}{10.215130in}}{\pgfqpoint{6.984217in}{10.215130in}}%
\pgfpathquadraticcurveto{\pgfqpoint{6.972735in}{10.215130in}}{\pgfqpoint{6.965836in}{10.208637in}}%
\pgfpathquadraticcurveto{\pgfqpoint{6.958937in}{10.202144in}}{\pgfqpoint{6.957887in}{10.190351in}}%
\pgfpathlineto{\pgfqpoint{7.006823in}{10.190447in}}%
\pgfpathlineto{\pgfqpoint{7.006823in}{10.190447in}}%
\pgfpathclose%
\pgfusepath{fill}%
\end{pgfscope}%
\begin{pgfscope}%
\pgfsetbuttcap%
\pgfsetroundjoin%
\definecolor{currentfill}{rgb}{0.090196,0.168627,0.227451}%
\pgfsetfillcolor{currentfill}%
\pgfsetlinewidth{0.000000pt}%
\definecolor{currentstroke}{rgb}{0.090196,0.168627,0.227451}%
\pgfsetstrokecolor{currentstroke}%
\pgfsetdash{}{0pt}%
\pgfpathmoveto{\pgfqpoint{1.316611in}{3.175090in}}%
\pgfpathlineto{\pgfqpoint{1.316611in}{3.162391in}}%
\pgfpathquadraticcurveto{\pgfqpoint{1.309210in}{3.165937in}}{\pgfqpoint{1.302633in}{3.167668in}}%
\pgfpathquadraticcurveto{\pgfqpoint{1.296056in}{3.169421in}}{\pgfqpoint{1.289933in}{3.169421in}}%
\pgfpathquadraticcurveto{\pgfqpoint{1.279316in}{3.169421in}}{\pgfqpoint{1.273543in}{3.165297in}}%
\pgfpathquadraticcurveto{\pgfqpoint{1.267771in}{3.161174in}}{\pgfqpoint{1.267771in}{3.153567in}}%
\pgfpathquadraticcurveto{\pgfqpoint{1.267771in}{3.147176in}}{\pgfqpoint{1.271605in}{3.143918in}}%
\pgfpathquadraticcurveto{\pgfqpoint{1.275440in}{3.140682in}}{\pgfqpoint{1.286140in}{3.138682in}}%
\pgfpathlineto{\pgfqpoint{1.293995in}{3.137074in}}%
\pgfpathquadraticcurveto{\pgfqpoint{1.308550in}{3.134291in}}{\pgfqpoint{1.315477in}{3.127302in}}%
\pgfpathquadraticcurveto{\pgfqpoint{1.322404in}{3.120313in}}{\pgfqpoint{1.322404in}{3.108603in}}%
\pgfpathquadraticcurveto{\pgfqpoint{1.322404in}{3.094604in}}{\pgfqpoint{1.313024in}{3.087388in}}%
\pgfpathquadraticcurveto{\pgfqpoint{1.303664in}{3.080173in}}{\pgfqpoint{1.285563in}{3.080173in}}%
\pgfpathquadraticcurveto{\pgfqpoint{1.278739in}{3.080173in}}{\pgfqpoint{1.271028in}{3.081719in}}%
\pgfpathquadraticcurveto{\pgfqpoint{1.263338in}{3.083265in}}{\pgfqpoint{1.255092in}{3.086296in}}%
\pgfpathlineto{\pgfqpoint{1.255092in}{3.099696in}}%
\pgfpathquadraticcurveto{\pgfqpoint{1.263008in}{3.095264in}}{\pgfqpoint{1.270616in}{3.092996in}}%
\pgfpathquadraticcurveto{\pgfqpoint{1.278223in}{3.090749in}}{\pgfqpoint{1.285563in}{3.090749in}}%
\pgfpathquadraticcurveto{\pgfqpoint{1.296696in}{3.090749in}}{\pgfqpoint{1.302757in}{3.095120in}}%
\pgfpathquadraticcurveto{\pgfqpoint{1.308818in}{3.099511in}}{\pgfqpoint{1.308818in}{3.107634in}}%
\pgfpathquadraticcurveto{\pgfqpoint{1.308818in}{3.114705in}}{\pgfqpoint{1.304468in}{3.118705in}}%
\pgfpathquadraticcurveto{\pgfqpoint{1.300118in}{3.122704in}}{\pgfqpoint{1.290201in}{3.124704in}}%
\pgfpathlineto{\pgfqpoint{1.282264in}{3.126250in}}%
\pgfpathquadraticcurveto{\pgfqpoint{1.267709in}{3.129136in}}{\pgfqpoint{1.261194in}{3.135321in}}%
\pgfpathquadraticcurveto{\pgfqpoint{1.254700in}{3.141506in}}{\pgfqpoint{1.254700in}{3.152536in}}%
\pgfpathquadraticcurveto{\pgfqpoint{1.254700in}{3.165297in}}{\pgfqpoint{1.263689in}{3.172637in}}%
\pgfpathquadraticcurveto{\pgfqpoint{1.272678in}{3.179976in}}{\pgfqpoint{1.288449in}{3.179976in}}%
\pgfpathquadraticcurveto{\pgfqpoint{1.295232in}{3.179976in}}{\pgfqpoint{1.302241in}{3.178739in}}%
\pgfpathquadraticcurveto{\pgfqpoint{1.309271in}{3.177523in}}{\pgfqpoint{1.316611in}{3.175090in}}%
\pgfpathlineto{\pgfqpoint{1.316611in}{3.175090in}}%
\pgfpathclose%
\pgfpathmoveto{\pgfqpoint{1.353926in}{3.174699in}}%
\pgfpathlineto{\pgfqpoint{1.353926in}{3.154206in}}%
\pgfpathlineto{\pgfqpoint{1.378336in}{3.154206in}}%
\pgfpathlineto{\pgfqpoint{1.378336in}{3.144990in}}%
\pgfpathlineto{\pgfqpoint{1.353926in}{3.144990in}}%
\pgfpathlineto{\pgfqpoint{1.353926in}{3.105819in}}%
\pgfpathquadraticcurveto{\pgfqpoint{1.353926in}{3.096996in}}{\pgfqpoint{1.356339in}{3.094480in}}%
\pgfpathquadraticcurveto{\pgfqpoint{1.358751in}{3.091965in}}{\pgfqpoint{1.366173in}{3.091965in}}%
\pgfpathlineto{\pgfqpoint{1.378336in}{3.091965in}}%
\pgfpathlineto{\pgfqpoint{1.378336in}{3.082049in}}%
\pgfpathlineto{\pgfqpoint{1.366173in}{3.082049in}}%
\pgfpathquadraticcurveto{\pgfqpoint{1.352442in}{3.082049in}}{\pgfqpoint{1.347226in}{3.087162in}}%
\pgfpathquadraticcurveto{\pgfqpoint{1.342010in}{3.092295in}}{\pgfqpoint{1.342010in}{3.105819in}}%
\pgfpathlineto{\pgfqpoint{1.342010in}{3.144990in}}%
\pgfpathlineto{\pgfqpoint{1.333310in}{3.144990in}}%
\pgfpathlineto{\pgfqpoint{1.333310in}{3.154206in}}%
\pgfpathlineto{\pgfqpoint{1.342010in}{3.154206in}}%
\pgfpathlineto{\pgfqpoint{1.342010in}{3.174699in}}%
\pgfpathlineto{\pgfqpoint{1.353926in}{3.174699in}}%
\pgfpathlineto{\pgfqpoint{1.353926in}{3.174699in}}%
\pgfpathclose%
\pgfpathmoveto{\pgfqpoint{1.392706in}{3.110520in}}%
\pgfpathlineto{\pgfqpoint{1.392706in}{3.154206in}}%
\pgfpathlineto{\pgfqpoint{1.404560in}{3.154206in}}%
\pgfpathlineto{\pgfqpoint{1.404560in}{3.110973in}}%
\pgfpathquadraticcurveto{\pgfqpoint{1.404560in}{3.100727in}}{\pgfqpoint{1.408539in}{3.095594in}}%
\pgfpathquadraticcurveto{\pgfqpoint{1.412539in}{3.090481in}}{\pgfqpoint{1.420538in}{3.090481in}}%
\pgfpathquadraticcurveto{\pgfqpoint{1.430124in}{3.090481in}}{\pgfqpoint{1.435691in}{3.096604in}}%
\pgfpathquadraticcurveto{\pgfqpoint{1.441278in}{3.102727in}}{\pgfqpoint{1.441278in}{3.113303in}}%
\pgfpathlineto{\pgfqpoint{1.441278in}{3.154206in}}%
\pgfpathlineto{\pgfqpoint{1.453132in}{3.154206in}}%
\pgfpathlineto{\pgfqpoint{1.453132in}{3.082049in}}%
\pgfpathlineto{\pgfqpoint{1.441278in}{3.082049in}}%
\pgfpathlineto{\pgfqpoint{1.441278in}{3.093140in}}%
\pgfpathquadraticcurveto{\pgfqpoint{1.436969in}{3.086564in}}{\pgfqpoint{1.431258in}{3.083368in}}%
\pgfpathquadraticcurveto{\pgfqpoint{1.425568in}{3.080173in}}{\pgfqpoint{1.418023in}{3.080173in}}%
\pgfpathquadraticcurveto{\pgfqpoint{1.405591in}{3.080173in}}{\pgfqpoint{1.399138in}{3.087904in}}%
\pgfpathquadraticcurveto{\pgfqpoint{1.392706in}{3.095635in}}{\pgfqpoint{1.392706in}{3.110520in}}%
\pgfpathlineto{\pgfqpoint{1.392706in}{3.110520in}}%
\pgfpathclose%
\pgfpathmoveto{\pgfqpoint{1.422538in}{3.155938in}}%
\pgfpathlineto{\pgfqpoint{1.422538in}{3.155938in}}%
\pgfpathlineto{\pgfqpoint{1.422538in}{3.155938in}}%
\pgfpathclose%
\pgfpathmoveto{\pgfqpoint{1.525021in}{3.143259in}}%
\pgfpathlineto{\pgfqpoint{1.525021in}{3.182306in}}%
\pgfpathlineto{\pgfqpoint{1.536876in}{3.182306in}}%
\pgfpathlineto{\pgfqpoint{1.536876in}{3.082049in}}%
\pgfpathlineto{\pgfqpoint{1.525021in}{3.082049in}}%
\pgfpathlineto{\pgfqpoint{1.525021in}{3.092872in}}%
\pgfpathquadraticcurveto{\pgfqpoint{1.521290in}{3.086440in}}{\pgfqpoint{1.515579in}{3.083306in}}%
\pgfpathquadraticcurveto{\pgfqpoint{1.509889in}{3.080173in}}{\pgfqpoint{1.501890in}{3.080173in}}%
\pgfpathquadraticcurveto{\pgfqpoint{1.488819in}{3.080173in}}{\pgfqpoint{1.480593in}{3.090605in}}%
\pgfpathquadraticcurveto{\pgfqpoint{1.472388in}{3.101057in}}{\pgfqpoint{1.472388in}{3.118066in}}%
\pgfpathquadraticcurveto{\pgfqpoint{1.472388in}{3.135074in}}{\pgfqpoint{1.480593in}{3.145506in}}%
\pgfpathquadraticcurveto{\pgfqpoint{1.488819in}{3.155938in}}{\pgfqpoint{1.501890in}{3.155938in}}%
\pgfpathquadraticcurveto{\pgfqpoint{1.509889in}{3.155938in}}{\pgfqpoint{1.515579in}{3.152804in}}%
\pgfpathquadraticcurveto{\pgfqpoint{1.521290in}{3.149691in}}{\pgfqpoint{1.525021in}{3.143259in}}%
\pgfpathlineto{\pgfqpoint{1.525021in}{3.143259in}}%
\pgfpathclose%
\pgfpathmoveto{\pgfqpoint{1.484634in}{3.118066in}}%
\pgfpathquadraticcurveto{\pgfqpoint{1.484634in}{3.104995in}}{\pgfqpoint{1.490015in}{3.097552in}}%
\pgfpathquadraticcurveto{\pgfqpoint{1.495396in}{3.090110in}}{\pgfqpoint{1.504797in}{3.090110in}}%
\pgfpathquadraticcurveto{\pgfqpoint{1.514198in}{3.090110in}}{\pgfqpoint{1.519599in}{3.097552in}}%
\pgfpathquadraticcurveto{\pgfqpoint{1.525021in}{3.104995in}}{\pgfqpoint{1.525021in}{3.118066in}}%
\pgfpathquadraticcurveto{\pgfqpoint{1.525021in}{3.131136in}}{\pgfqpoint{1.519599in}{3.138579in}}%
\pgfpathquadraticcurveto{\pgfqpoint{1.514198in}{3.146021in}}{\pgfqpoint{1.504797in}{3.146021in}}%
\pgfpathquadraticcurveto{\pgfqpoint{1.495396in}{3.146021in}}{\pgfqpoint{1.490015in}{3.138579in}}%
\pgfpathquadraticcurveto{\pgfqpoint{1.484634in}{3.131136in}}{\pgfqpoint{1.484634in}{3.118066in}}%
\pgfpathlineto{\pgfqpoint{1.484634in}{3.118066in}}%
\pgfpathclose%
\pgfpathmoveto{\pgfqpoint{1.591323in}{3.075348in}}%
\pgfpathquadraticcurveto{\pgfqpoint{1.586314in}{3.062463in}}{\pgfqpoint{1.581531in}{3.058546in}}%
\pgfpathquadraticcurveto{\pgfqpoint{1.576768in}{3.054608in}}{\pgfqpoint{1.568790in}{3.054608in}}%
\pgfpathlineto{\pgfqpoint{1.559306in}{3.054608in}}%
\pgfpathlineto{\pgfqpoint{1.559306in}{3.064525in}}%
\pgfpathlineto{\pgfqpoint{1.566275in}{3.064525in}}%
\pgfpathquadraticcurveto{\pgfqpoint{1.571161in}{3.064525in}}{\pgfqpoint{1.573861in}{3.066855in}}%
\pgfpathquadraticcurveto{\pgfqpoint{1.576583in}{3.069164in}}{\pgfqpoint{1.579861in}{3.077802in}}%
\pgfpathlineto{\pgfqpoint{1.581984in}{3.083203in}}%
\pgfpathlineto{\pgfqpoint{1.552812in}{3.154206in}}%
\pgfpathlineto{\pgfqpoint{1.565367in}{3.154206in}}%
\pgfpathlineto{\pgfqpoint{1.587922in}{3.097779in}}%
\pgfpathlineto{\pgfqpoint{1.610476in}{3.154206in}}%
\pgfpathlineto{\pgfqpoint{1.623031in}{3.154206in}}%
\pgfpathlineto{\pgfqpoint{1.591323in}{3.075348in}}%
\pgfpathlineto{\pgfqpoint{1.591323in}{3.075348in}}%
\pgfpathclose%
\pgfpathmoveto{\pgfqpoint{1.685272in}{3.092996in}}%
\pgfpathlineto{\pgfqpoint{1.706527in}{3.092996in}}%
\pgfpathlineto{\pgfqpoint{1.706527in}{3.166390in}}%
\pgfpathlineto{\pgfqpoint{1.683396in}{3.161751in}}%
\pgfpathlineto{\pgfqpoint{1.683396in}{3.173606in}}%
\pgfpathlineto{\pgfqpoint{1.706404in}{3.178245in}}%
\pgfpathlineto{\pgfqpoint{1.719413in}{3.178245in}}%
\pgfpathlineto{\pgfqpoint{1.719413in}{3.092996in}}%
\pgfpathlineto{\pgfqpoint{1.740668in}{3.092996in}}%
\pgfpathlineto{\pgfqpoint{1.740668in}{3.082049in}}%
\pgfpathlineto{\pgfqpoint{1.685272in}{3.082049in}}%
\pgfpathlineto{\pgfqpoint{1.685272in}{3.092996in}}%
\pgfpathlineto{\pgfqpoint{1.685272in}{3.092996in}}%
\pgfpathclose%
\pgfpathmoveto{\pgfqpoint{1.835689in}{3.182162in}}%
\pgfpathquadraticcurveto{\pgfqpoint{1.827071in}{3.167359in}}{\pgfqpoint{1.822865in}{3.152845in}}%
\pgfpathquadraticcurveto{\pgfqpoint{1.818680in}{3.138352in}}{\pgfqpoint{1.818680in}{3.123467in}}%
\pgfpathquadraticcurveto{\pgfqpoint{1.818680in}{3.108603in}}{\pgfqpoint{1.822906in}{3.094006in}}%
\pgfpathquadraticcurveto{\pgfqpoint{1.827133in}{3.079410in}}{\pgfqpoint{1.835689in}{3.064649in}}%
\pgfpathlineto{\pgfqpoint{1.825380in}{3.064649in}}%
\pgfpathquadraticcurveto{\pgfqpoint{1.815732in}{3.079802in}}{\pgfqpoint{1.810928in}{3.094419in}}%
\pgfpathquadraticcurveto{\pgfqpoint{1.806125in}{3.109036in}}{\pgfqpoint{1.806125in}{3.123467in}}%
\pgfpathquadraticcurveto{\pgfqpoint{1.806125in}{3.137837in}}{\pgfqpoint{1.810887in}{3.152392in}}%
\pgfpathquadraticcurveto{\pgfqpoint{1.815670in}{3.166967in}}{\pgfqpoint{1.825380in}{3.182162in}}%
\pgfpathlineto{\pgfqpoint{1.835689in}{3.182162in}}%
\pgfpathlineto{\pgfqpoint{1.835689in}{3.182162in}}%
\pgfpathclose%
\pgfpathmoveto{\pgfqpoint{1.871582in}{3.092996in}}%
\pgfpathlineto{\pgfqpoint{1.916999in}{3.092996in}}%
\pgfpathlineto{\pgfqpoint{1.916999in}{3.082049in}}%
\pgfpathlineto{\pgfqpoint{1.855934in}{3.082049in}}%
\pgfpathlineto{\pgfqpoint{1.855934in}{3.092996in}}%
\pgfpathquadraticcurveto{\pgfqpoint{1.863335in}{3.100665in}}{\pgfqpoint{1.876117in}{3.113571in}}%
\pgfpathquadraticcurveto{\pgfqpoint{1.888920in}{3.126498in}}{\pgfqpoint{1.892198in}{3.130250in}}%
\pgfpathquadraticcurveto{\pgfqpoint{1.898445in}{3.137259in}}{\pgfqpoint{1.900919in}{3.142125in}}%
\pgfpathquadraticcurveto{\pgfqpoint{1.903413in}{3.146990in}}{\pgfqpoint{1.903413in}{3.151691in}}%
\pgfpathquadraticcurveto{\pgfqpoint{1.903413in}{3.159360in}}{\pgfqpoint{1.898032in}{3.164184in}}%
\pgfpathquadraticcurveto{\pgfqpoint{1.892651in}{3.169029in}}{\pgfqpoint{1.884013in}{3.169029in}}%
\pgfpathquadraticcurveto{\pgfqpoint{1.877890in}{3.169029in}}{\pgfqpoint{1.871087in}{3.166906in}}%
\pgfpathquadraticcurveto{\pgfqpoint{1.864304in}{3.164782in}}{\pgfqpoint{1.856573in}{3.160453in}}%
\pgfpathlineto{\pgfqpoint{1.856573in}{3.173606in}}%
\pgfpathquadraticcurveto{\pgfqpoint{1.864428in}{3.176760in}}{\pgfqpoint{1.871252in}{3.178368in}}%
\pgfpathquadraticcurveto{\pgfqpoint{1.878096in}{3.179976in}}{\pgfqpoint{1.883766in}{3.179976in}}%
\pgfpathquadraticcurveto{\pgfqpoint{1.898713in}{3.179976in}}{\pgfqpoint{1.907598in}{3.172493in}}%
\pgfpathquadraticcurveto{\pgfqpoint{1.916484in}{3.165029in}}{\pgfqpoint{1.916484in}{3.152536in}}%
\pgfpathquadraticcurveto{\pgfqpoint{1.916484in}{3.146598in}}{\pgfqpoint{1.914257in}{3.141279in}}%
\pgfpathquadraticcurveto{\pgfqpoint{1.912051in}{3.135981in}}{\pgfqpoint{1.906176in}{3.128765in}}%
\pgfpathquadraticcurveto{\pgfqpoint{1.904568in}{3.126889in}}{\pgfqpoint{1.895929in}{3.117962in}}%
\pgfpathquadraticcurveto{\pgfqpoint{1.887312in}{3.109036in}}{\pgfqpoint{1.871582in}{3.092996in}}%
\pgfpathlineto{\pgfqpoint{1.871582in}{3.092996in}}%
\pgfpathclose%
\pgfpathmoveto{\pgfqpoint{1.972148in}{3.169668in}}%
\pgfpathquadraticcurveto{\pgfqpoint{1.962108in}{3.169668in}}{\pgfqpoint{1.957036in}{3.159772in}}%
\pgfpathquadraticcurveto{\pgfqpoint{1.951985in}{3.149897in}}{\pgfqpoint{1.951985in}{3.130044in}}%
\pgfpathquadraticcurveto{\pgfqpoint{1.951985in}{3.110273in}}{\pgfqpoint{1.957036in}{3.100377in}}%
\pgfpathquadraticcurveto{\pgfqpoint{1.962108in}{3.090481in}}{\pgfqpoint{1.972148in}{3.090481in}}%
\pgfpathquadraticcurveto{\pgfqpoint{1.982271in}{3.090481in}}{\pgfqpoint{1.987322in}{3.100377in}}%
\pgfpathquadraticcurveto{\pgfqpoint{1.992393in}{3.110273in}}{\pgfqpoint{1.992393in}{3.130044in}}%
\pgfpathquadraticcurveto{\pgfqpoint{1.992393in}{3.149897in}}{\pgfqpoint{1.987322in}{3.159772in}}%
\pgfpathquadraticcurveto{\pgfqpoint{1.982271in}{3.169668in}}{\pgfqpoint{1.972148in}{3.169668in}}%
\pgfpathlineto{\pgfqpoint{1.972148in}{3.169668in}}%
\pgfpathclose%
\pgfpathmoveto{\pgfqpoint{1.972148in}{3.179976in}}%
\pgfpathquadraticcurveto{\pgfqpoint{1.988332in}{3.179976in}}{\pgfqpoint{1.996867in}{3.167174in}}%
\pgfpathquadraticcurveto{\pgfqpoint{2.005402in}{3.154391in}}{\pgfqpoint{2.005402in}{3.130044in}}%
\pgfpathquadraticcurveto{\pgfqpoint{2.005402in}{3.105758in}}{\pgfqpoint{1.996867in}{3.092955in}}%
\pgfpathquadraticcurveto{\pgfqpoint{1.988332in}{3.080173in}}{\pgfqpoint{1.972148in}{3.080173in}}%
\pgfpathquadraticcurveto{\pgfqpoint{1.955985in}{3.080173in}}{\pgfqpoint{1.947450in}{3.092955in}}%
\pgfpathquadraticcurveto{\pgfqpoint{1.938914in}{3.105758in}}{\pgfqpoint{1.938914in}{3.130044in}}%
\pgfpathquadraticcurveto{\pgfqpoint{1.938914in}{3.154391in}}{\pgfqpoint{1.947450in}{3.167174in}}%
\pgfpathquadraticcurveto{\pgfqpoint{1.955985in}{3.179976in}}{\pgfqpoint{1.972148in}{3.179976in}}%
\pgfpathlineto{\pgfqpoint{1.972148in}{3.179976in}}%
\pgfpathclose%
\pgfpathmoveto{\pgfqpoint{2.030533in}{3.092996in}}%
\pgfpathlineto{\pgfqpoint{2.051789in}{3.092996in}}%
\pgfpathlineto{\pgfqpoint{2.051789in}{3.166390in}}%
\pgfpathlineto{\pgfqpoint{2.028657in}{3.161751in}}%
\pgfpathlineto{\pgfqpoint{2.028657in}{3.173606in}}%
\pgfpathlineto{\pgfqpoint{2.051665in}{3.178245in}}%
\pgfpathlineto{\pgfqpoint{2.064674in}{3.178245in}}%
\pgfpathlineto{\pgfqpoint{2.064674in}{3.092996in}}%
\pgfpathlineto{\pgfqpoint{2.085929in}{3.092996in}}%
\pgfpathlineto{\pgfqpoint{2.085929in}{3.082049in}}%
\pgfpathlineto{\pgfqpoint{2.030533in}{3.082049in}}%
\pgfpathlineto{\pgfqpoint{2.030533in}{3.092996in}}%
\pgfpathlineto{\pgfqpoint{2.030533in}{3.092996in}}%
\pgfpathclose%
\pgfpathmoveto{\pgfqpoint{2.108937in}{3.178245in}}%
\pgfpathlineto{\pgfqpoint{2.170786in}{3.178245in}}%
\pgfpathlineto{\pgfqpoint{2.170786in}{3.172699in}}%
\pgfpathlineto{\pgfqpoint{2.135862in}{3.082049in}}%
\pgfpathlineto{\pgfqpoint{2.122276in}{3.082049in}}%
\pgfpathlineto{\pgfqpoint{2.155138in}{3.167277in}}%
\pgfpathlineto{\pgfqpoint{2.108937in}{3.167277in}}%
\pgfpathlineto{\pgfqpoint{2.108937in}{3.178245in}}%
\pgfpathlineto{\pgfqpoint{2.108937in}{3.178245in}}%
\pgfpathclose%
\pgfpathmoveto{\pgfqpoint{2.192640in}{3.182162in}}%
\pgfpathlineto{\pgfqpoint{2.202948in}{3.182162in}}%
\pgfpathquadraticcurveto{\pgfqpoint{2.212596in}{3.166967in}}{\pgfqpoint{2.217400in}{3.152392in}}%
\pgfpathquadraticcurveto{\pgfqpoint{2.222203in}{3.137837in}}{\pgfqpoint{2.222203in}{3.123467in}}%
\pgfpathquadraticcurveto{\pgfqpoint{2.222203in}{3.109036in}}{\pgfqpoint{2.217400in}{3.094419in}}%
\pgfpathquadraticcurveto{\pgfqpoint{2.212596in}{3.079802in}}{\pgfqpoint{2.202948in}{3.064649in}}%
\pgfpathlineto{\pgfqpoint{2.192640in}{3.064649in}}%
\pgfpathquadraticcurveto{\pgfqpoint{2.201195in}{3.079410in}}{\pgfqpoint{2.205422in}{3.094006in}}%
\pgfpathquadraticcurveto{\pgfqpoint{2.209648in}{3.108603in}}{\pgfqpoint{2.209648in}{3.123467in}}%
\pgfpathquadraticcurveto{\pgfqpoint{2.209648in}{3.138352in}}{\pgfqpoint{2.205422in}{3.152845in}}%
\pgfpathquadraticcurveto{\pgfqpoint{2.201195in}{3.167359in}}{\pgfqpoint{2.192640in}{3.182162in}}%
\pgfpathlineto{\pgfqpoint{2.192640in}{3.182162in}}%
\pgfpathclose%
\pgfusepath{fill}%
\end{pgfscope}%
\begin{pgfscope}%
\pgfsetbuttcap%
\pgfsetroundjoin%
\definecolor{currentfill}{rgb}{0.090196,0.168627,0.227451}%
\pgfsetfillcolor{currentfill}%
\pgfsetlinewidth{0.000000pt}%
\definecolor{currentstroke}{rgb}{0.090196,0.168627,0.227451}%
\pgfsetstrokecolor{currentstroke}%
\pgfsetdash{}{0pt}%
\pgfpathmoveto{\pgfqpoint{1.316611in}{2.775490in}}%
\pgfpathlineto{\pgfqpoint{1.316611in}{2.762791in}}%
\pgfpathquadraticcurveto{\pgfqpoint{1.309210in}{2.766337in}}{\pgfqpoint{1.302633in}{2.768068in}}%
\pgfpathquadraticcurveto{\pgfqpoint{1.296056in}{2.769821in}}{\pgfqpoint{1.289933in}{2.769821in}}%
\pgfpathquadraticcurveto{\pgfqpoint{1.279316in}{2.769821in}}{\pgfqpoint{1.273543in}{2.765697in}}%
\pgfpathquadraticcurveto{\pgfqpoint{1.267771in}{2.761574in}}{\pgfqpoint{1.267771in}{2.753967in}}%
\pgfpathquadraticcurveto{\pgfqpoint{1.267771in}{2.747576in}}{\pgfqpoint{1.271605in}{2.744318in}}%
\pgfpathquadraticcurveto{\pgfqpoint{1.275440in}{2.741082in}}{\pgfqpoint{1.286140in}{2.739082in}}%
\pgfpathlineto{\pgfqpoint{1.293995in}{2.737474in}}%
\pgfpathquadraticcurveto{\pgfqpoint{1.308550in}{2.734691in}}{\pgfqpoint{1.315477in}{2.727702in}}%
\pgfpathquadraticcurveto{\pgfqpoint{1.322404in}{2.720713in}}{\pgfqpoint{1.322404in}{2.709003in}}%
\pgfpathquadraticcurveto{\pgfqpoint{1.322404in}{2.695004in}}{\pgfqpoint{1.313024in}{2.687788in}}%
\pgfpathquadraticcurveto{\pgfqpoint{1.303664in}{2.680573in}}{\pgfqpoint{1.285563in}{2.680573in}}%
\pgfpathquadraticcurveto{\pgfqpoint{1.278739in}{2.680573in}}{\pgfqpoint{1.271028in}{2.682119in}}%
\pgfpathquadraticcurveto{\pgfqpoint{1.263338in}{2.683665in}}{\pgfqpoint{1.255092in}{2.686696in}}%
\pgfpathlineto{\pgfqpoint{1.255092in}{2.700096in}}%
\pgfpathquadraticcurveto{\pgfqpoint{1.263008in}{2.695664in}}{\pgfqpoint{1.270616in}{2.693396in}}%
\pgfpathquadraticcurveto{\pgfqpoint{1.278223in}{2.691149in}}{\pgfqpoint{1.285563in}{2.691149in}}%
\pgfpathquadraticcurveto{\pgfqpoint{1.296696in}{2.691149in}}{\pgfqpoint{1.302757in}{2.695520in}}%
\pgfpathquadraticcurveto{\pgfqpoint{1.308818in}{2.699911in}}{\pgfqpoint{1.308818in}{2.708034in}}%
\pgfpathquadraticcurveto{\pgfqpoint{1.308818in}{2.715105in}}{\pgfqpoint{1.304468in}{2.719105in}}%
\pgfpathquadraticcurveto{\pgfqpoint{1.300118in}{2.723104in}}{\pgfqpoint{1.290201in}{2.725104in}}%
\pgfpathlineto{\pgfqpoint{1.282264in}{2.726650in}}%
\pgfpathquadraticcurveto{\pgfqpoint{1.267709in}{2.729536in}}{\pgfqpoint{1.261194in}{2.735721in}}%
\pgfpathquadraticcurveto{\pgfqpoint{1.254700in}{2.741906in}}{\pgfqpoint{1.254700in}{2.752936in}}%
\pgfpathquadraticcurveto{\pgfqpoint{1.254700in}{2.765697in}}{\pgfqpoint{1.263689in}{2.773037in}}%
\pgfpathquadraticcurveto{\pgfqpoint{1.272678in}{2.780376in}}{\pgfqpoint{1.288449in}{2.780376in}}%
\pgfpathquadraticcurveto{\pgfqpoint{1.295232in}{2.780376in}}{\pgfqpoint{1.302241in}{2.779139in}}%
\pgfpathquadraticcurveto{\pgfqpoint{1.309271in}{2.777923in}}{\pgfqpoint{1.316611in}{2.775490in}}%
\pgfpathlineto{\pgfqpoint{1.316611in}{2.775490in}}%
\pgfpathclose%
\pgfpathmoveto{\pgfqpoint{1.353926in}{2.775099in}}%
\pgfpathlineto{\pgfqpoint{1.353926in}{2.754606in}}%
\pgfpathlineto{\pgfqpoint{1.378336in}{2.754606in}}%
\pgfpathlineto{\pgfqpoint{1.378336in}{2.745390in}}%
\pgfpathlineto{\pgfqpoint{1.353926in}{2.745390in}}%
\pgfpathlineto{\pgfqpoint{1.353926in}{2.706219in}}%
\pgfpathquadraticcurveto{\pgfqpoint{1.353926in}{2.697396in}}{\pgfqpoint{1.356339in}{2.694880in}}%
\pgfpathquadraticcurveto{\pgfqpoint{1.358751in}{2.692365in}}{\pgfqpoint{1.366173in}{2.692365in}}%
\pgfpathlineto{\pgfqpoint{1.378336in}{2.692365in}}%
\pgfpathlineto{\pgfqpoint{1.378336in}{2.682449in}}%
\pgfpathlineto{\pgfqpoint{1.366173in}{2.682449in}}%
\pgfpathquadraticcurveto{\pgfqpoint{1.352442in}{2.682449in}}{\pgfqpoint{1.347226in}{2.687562in}}%
\pgfpathquadraticcurveto{\pgfqpoint{1.342010in}{2.692695in}}{\pgfqpoint{1.342010in}{2.706219in}}%
\pgfpathlineto{\pgfqpoint{1.342010in}{2.745390in}}%
\pgfpathlineto{\pgfqpoint{1.333310in}{2.745390in}}%
\pgfpathlineto{\pgfqpoint{1.333310in}{2.754606in}}%
\pgfpathlineto{\pgfqpoint{1.342010in}{2.754606in}}%
\pgfpathlineto{\pgfqpoint{1.342010in}{2.775099in}}%
\pgfpathlineto{\pgfqpoint{1.353926in}{2.775099in}}%
\pgfpathlineto{\pgfqpoint{1.353926in}{2.775099in}}%
\pgfpathclose%
\pgfpathmoveto{\pgfqpoint{1.392706in}{2.710920in}}%
\pgfpathlineto{\pgfqpoint{1.392706in}{2.754606in}}%
\pgfpathlineto{\pgfqpoint{1.404560in}{2.754606in}}%
\pgfpathlineto{\pgfqpoint{1.404560in}{2.711373in}}%
\pgfpathquadraticcurveto{\pgfqpoint{1.404560in}{2.701127in}}{\pgfqpoint{1.408539in}{2.695994in}}%
\pgfpathquadraticcurveto{\pgfqpoint{1.412539in}{2.690881in}}{\pgfqpoint{1.420538in}{2.690881in}}%
\pgfpathquadraticcurveto{\pgfqpoint{1.430124in}{2.690881in}}{\pgfqpoint{1.435691in}{2.697004in}}%
\pgfpathquadraticcurveto{\pgfqpoint{1.441278in}{2.703127in}}{\pgfqpoint{1.441278in}{2.713703in}}%
\pgfpathlineto{\pgfqpoint{1.441278in}{2.754606in}}%
\pgfpathlineto{\pgfqpoint{1.453132in}{2.754606in}}%
\pgfpathlineto{\pgfqpoint{1.453132in}{2.682449in}}%
\pgfpathlineto{\pgfqpoint{1.441278in}{2.682449in}}%
\pgfpathlineto{\pgfqpoint{1.441278in}{2.693540in}}%
\pgfpathquadraticcurveto{\pgfqpoint{1.436969in}{2.686964in}}{\pgfqpoint{1.431258in}{2.683768in}}%
\pgfpathquadraticcurveto{\pgfqpoint{1.425568in}{2.680573in}}{\pgfqpoint{1.418023in}{2.680573in}}%
\pgfpathquadraticcurveto{\pgfqpoint{1.405591in}{2.680573in}}{\pgfqpoint{1.399138in}{2.688304in}}%
\pgfpathquadraticcurveto{\pgfqpoint{1.392706in}{2.696035in}}{\pgfqpoint{1.392706in}{2.710920in}}%
\pgfpathlineto{\pgfqpoint{1.392706in}{2.710920in}}%
\pgfpathclose%
\pgfpathmoveto{\pgfqpoint{1.422538in}{2.756338in}}%
\pgfpathlineto{\pgfqpoint{1.422538in}{2.756338in}}%
\pgfpathlineto{\pgfqpoint{1.422538in}{2.756338in}}%
\pgfpathclose%
\pgfpathmoveto{\pgfqpoint{1.525021in}{2.743659in}}%
\pgfpathlineto{\pgfqpoint{1.525021in}{2.782706in}}%
\pgfpathlineto{\pgfqpoint{1.536876in}{2.782706in}}%
\pgfpathlineto{\pgfqpoint{1.536876in}{2.682449in}}%
\pgfpathlineto{\pgfqpoint{1.525021in}{2.682449in}}%
\pgfpathlineto{\pgfqpoint{1.525021in}{2.693272in}}%
\pgfpathquadraticcurveto{\pgfqpoint{1.521290in}{2.686840in}}{\pgfqpoint{1.515579in}{2.683706in}}%
\pgfpathquadraticcurveto{\pgfqpoint{1.509889in}{2.680573in}}{\pgfqpoint{1.501890in}{2.680573in}}%
\pgfpathquadraticcurveto{\pgfqpoint{1.488819in}{2.680573in}}{\pgfqpoint{1.480593in}{2.691005in}}%
\pgfpathquadraticcurveto{\pgfqpoint{1.472388in}{2.701457in}}{\pgfqpoint{1.472388in}{2.718466in}}%
\pgfpathquadraticcurveto{\pgfqpoint{1.472388in}{2.735474in}}{\pgfqpoint{1.480593in}{2.745906in}}%
\pgfpathquadraticcurveto{\pgfqpoint{1.488819in}{2.756338in}}{\pgfqpoint{1.501890in}{2.756338in}}%
\pgfpathquadraticcurveto{\pgfqpoint{1.509889in}{2.756338in}}{\pgfqpoint{1.515579in}{2.753204in}}%
\pgfpathquadraticcurveto{\pgfqpoint{1.521290in}{2.750091in}}{\pgfqpoint{1.525021in}{2.743659in}}%
\pgfpathlineto{\pgfqpoint{1.525021in}{2.743659in}}%
\pgfpathclose%
\pgfpathmoveto{\pgfqpoint{1.484634in}{2.718466in}}%
\pgfpathquadraticcurveto{\pgfqpoint{1.484634in}{2.705395in}}{\pgfqpoint{1.490015in}{2.697952in}}%
\pgfpathquadraticcurveto{\pgfqpoint{1.495396in}{2.690510in}}{\pgfqpoint{1.504797in}{2.690510in}}%
\pgfpathquadraticcurveto{\pgfqpoint{1.514198in}{2.690510in}}{\pgfqpoint{1.519599in}{2.697952in}}%
\pgfpathquadraticcurveto{\pgfqpoint{1.525021in}{2.705395in}}{\pgfqpoint{1.525021in}{2.718466in}}%
\pgfpathquadraticcurveto{\pgfqpoint{1.525021in}{2.731536in}}{\pgfqpoint{1.519599in}{2.738979in}}%
\pgfpathquadraticcurveto{\pgfqpoint{1.514198in}{2.746421in}}{\pgfqpoint{1.504797in}{2.746421in}}%
\pgfpathquadraticcurveto{\pgfqpoint{1.495396in}{2.746421in}}{\pgfqpoint{1.490015in}{2.738979in}}%
\pgfpathquadraticcurveto{\pgfqpoint{1.484634in}{2.731536in}}{\pgfqpoint{1.484634in}{2.718466in}}%
\pgfpathlineto{\pgfqpoint{1.484634in}{2.718466in}}%
\pgfpathclose%
\pgfpathmoveto{\pgfqpoint{1.591323in}{2.675748in}}%
\pgfpathquadraticcurveto{\pgfqpoint{1.586314in}{2.662863in}}{\pgfqpoint{1.581531in}{2.658946in}}%
\pgfpathquadraticcurveto{\pgfqpoint{1.576768in}{2.655008in}}{\pgfqpoint{1.568790in}{2.655008in}}%
\pgfpathlineto{\pgfqpoint{1.559306in}{2.655008in}}%
\pgfpathlineto{\pgfqpoint{1.559306in}{2.664925in}}%
\pgfpathlineto{\pgfqpoint{1.566275in}{2.664925in}}%
\pgfpathquadraticcurveto{\pgfqpoint{1.571161in}{2.664925in}}{\pgfqpoint{1.573861in}{2.667255in}}%
\pgfpathquadraticcurveto{\pgfqpoint{1.576583in}{2.669564in}}{\pgfqpoint{1.579861in}{2.678202in}}%
\pgfpathlineto{\pgfqpoint{1.581984in}{2.683603in}}%
\pgfpathlineto{\pgfqpoint{1.552812in}{2.754606in}}%
\pgfpathlineto{\pgfqpoint{1.565367in}{2.754606in}}%
\pgfpathlineto{\pgfqpoint{1.587922in}{2.698179in}}%
\pgfpathlineto{\pgfqpoint{1.610476in}{2.754606in}}%
\pgfpathlineto{\pgfqpoint{1.623031in}{2.754606in}}%
\pgfpathlineto{\pgfqpoint{1.591323in}{2.675748in}}%
\pgfpathlineto{\pgfqpoint{1.591323in}{2.675748in}}%
\pgfpathclose%
\pgfpathmoveto{\pgfqpoint{1.694219in}{2.693396in}}%
\pgfpathlineto{\pgfqpoint{1.739637in}{2.693396in}}%
\pgfpathlineto{\pgfqpoint{1.739637in}{2.682449in}}%
\pgfpathlineto{\pgfqpoint{1.678572in}{2.682449in}}%
\pgfpathlineto{\pgfqpoint{1.678572in}{2.693396in}}%
\pgfpathquadraticcurveto{\pgfqpoint{1.685973in}{2.701065in}}{\pgfqpoint{1.698755in}{2.713971in}}%
\pgfpathquadraticcurveto{\pgfqpoint{1.711558in}{2.726898in}}{\pgfqpoint{1.714836in}{2.730650in}}%
\pgfpathquadraticcurveto{\pgfqpoint{1.721082in}{2.737659in}}{\pgfqpoint{1.723556in}{2.742525in}}%
\pgfpathquadraticcurveto{\pgfqpoint{1.726051in}{2.747390in}}{\pgfqpoint{1.726051in}{2.752091in}}%
\pgfpathquadraticcurveto{\pgfqpoint{1.726051in}{2.759760in}}{\pgfqpoint{1.720670in}{2.764584in}}%
\pgfpathquadraticcurveto{\pgfqpoint{1.715289in}{2.769429in}}{\pgfqpoint{1.706651in}{2.769429in}}%
\pgfpathquadraticcurveto{\pgfqpoint{1.700528in}{2.769429in}}{\pgfqpoint{1.693725in}{2.767306in}}%
\pgfpathquadraticcurveto{\pgfqpoint{1.686942in}{2.765182in}}{\pgfqpoint{1.679211in}{2.760853in}}%
\pgfpathlineto{\pgfqpoint{1.679211in}{2.774006in}}%
\pgfpathquadraticcurveto{\pgfqpoint{1.687066in}{2.777160in}}{\pgfqpoint{1.693890in}{2.778768in}}%
\pgfpathquadraticcurveto{\pgfqpoint{1.700734in}{2.780376in}}{\pgfqpoint{1.706404in}{2.780376in}}%
\pgfpathquadraticcurveto{\pgfqpoint{1.721350in}{2.780376in}}{\pgfqpoint{1.730236in}{2.772893in}}%
\pgfpathquadraticcurveto{\pgfqpoint{1.739122in}{2.765429in}}{\pgfqpoint{1.739122in}{2.752936in}}%
\pgfpathquadraticcurveto{\pgfqpoint{1.739122in}{2.746998in}}{\pgfqpoint{1.736895in}{2.741679in}}%
\pgfpathquadraticcurveto{\pgfqpoint{1.734689in}{2.736381in}}{\pgfqpoint{1.728814in}{2.729165in}}%
\pgfpathquadraticcurveto{\pgfqpoint{1.727206in}{2.727289in}}{\pgfqpoint{1.718567in}{2.718362in}}%
\pgfpathquadraticcurveto{\pgfqpoint{1.709950in}{2.709436in}}{\pgfqpoint{1.694219in}{2.693396in}}%
\pgfpathlineto{\pgfqpoint{1.694219in}{2.693396in}}%
\pgfpathclose%
\pgfpathmoveto{\pgfqpoint{1.835689in}{2.782562in}}%
\pgfpathquadraticcurveto{\pgfqpoint{1.827071in}{2.767759in}}{\pgfqpoint{1.822865in}{2.753245in}}%
\pgfpathquadraticcurveto{\pgfqpoint{1.818680in}{2.738752in}}{\pgfqpoint{1.818680in}{2.723867in}}%
\pgfpathquadraticcurveto{\pgfqpoint{1.818680in}{2.709003in}}{\pgfqpoint{1.822906in}{2.694406in}}%
\pgfpathquadraticcurveto{\pgfqpoint{1.827133in}{2.679810in}}{\pgfqpoint{1.835689in}{2.665049in}}%
\pgfpathlineto{\pgfqpoint{1.825380in}{2.665049in}}%
\pgfpathquadraticcurveto{\pgfqpoint{1.815732in}{2.680202in}}{\pgfqpoint{1.810928in}{2.694819in}}%
\pgfpathquadraticcurveto{\pgfqpoint{1.806125in}{2.709436in}}{\pgfqpoint{1.806125in}{2.723867in}}%
\pgfpathquadraticcurveto{\pgfqpoint{1.806125in}{2.738237in}}{\pgfqpoint{1.810887in}{2.752792in}}%
\pgfpathquadraticcurveto{\pgfqpoint{1.815670in}{2.767367in}}{\pgfqpoint{1.825380in}{2.782562in}}%
\pgfpathlineto{\pgfqpoint{1.835689in}{2.782562in}}%
\pgfpathlineto{\pgfqpoint{1.835689in}{2.782562in}}%
\pgfpathclose%
\pgfpathmoveto{\pgfqpoint{1.871582in}{2.693396in}}%
\pgfpathlineto{\pgfqpoint{1.916999in}{2.693396in}}%
\pgfpathlineto{\pgfqpoint{1.916999in}{2.682449in}}%
\pgfpathlineto{\pgfqpoint{1.855934in}{2.682449in}}%
\pgfpathlineto{\pgfqpoint{1.855934in}{2.693396in}}%
\pgfpathquadraticcurveto{\pgfqpoint{1.863335in}{2.701065in}}{\pgfqpoint{1.876117in}{2.713971in}}%
\pgfpathquadraticcurveto{\pgfqpoint{1.888920in}{2.726898in}}{\pgfqpoint{1.892198in}{2.730650in}}%
\pgfpathquadraticcurveto{\pgfqpoint{1.898445in}{2.737659in}}{\pgfqpoint{1.900919in}{2.742525in}}%
\pgfpathquadraticcurveto{\pgfqpoint{1.903413in}{2.747390in}}{\pgfqpoint{1.903413in}{2.752091in}}%
\pgfpathquadraticcurveto{\pgfqpoint{1.903413in}{2.759760in}}{\pgfqpoint{1.898032in}{2.764584in}}%
\pgfpathquadraticcurveto{\pgfqpoint{1.892651in}{2.769429in}}{\pgfqpoint{1.884013in}{2.769429in}}%
\pgfpathquadraticcurveto{\pgfqpoint{1.877890in}{2.769429in}}{\pgfqpoint{1.871087in}{2.767306in}}%
\pgfpathquadraticcurveto{\pgfqpoint{1.864304in}{2.765182in}}{\pgfqpoint{1.856573in}{2.760853in}}%
\pgfpathlineto{\pgfqpoint{1.856573in}{2.774006in}}%
\pgfpathquadraticcurveto{\pgfqpoint{1.864428in}{2.777160in}}{\pgfqpoint{1.871252in}{2.778768in}}%
\pgfpathquadraticcurveto{\pgfqpoint{1.878096in}{2.780376in}}{\pgfqpoint{1.883766in}{2.780376in}}%
\pgfpathquadraticcurveto{\pgfqpoint{1.898713in}{2.780376in}}{\pgfqpoint{1.907598in}{2.772893in}}%
\pgfpathquadraticcurveto{\pgfqpoint{1.916484in}{2.765429in}}{\pgfqpoint{1.916484in}{2.752936in}}%
\pgfpathquadraticcurveto{\pgfqpoint{1.916484in}{2.746998in}}{\pgfqpoint{1.914257in}{2.741679in}}%
\pgfpathquadraticcurveto{\pgfqpoint{1.912051in}{2.736381in}}{\pgfqpoint{1.906176in}{2.729165in}}%
\pgfpathquadraticcurveto{\pgfqpoint{1.904568in}{2.727289in}}{\pgfqpoint{1.895929in}{2.718362in}}%
\pgfpathquadraticcurveto{\pgfqpoint{1.887312in}{2.709436in}}{\pgfqpoint{1.871582in}{2.693396in}}%
\pgfpathlineto{\pgfqpoint{1.871582in}{2.693396in}}%
\pgfpathclose%
\pgfpathmoveto{\pgfqpoint{1.972148in}{2.770068in}}%
\pgfpathquadraticcurveto{\pgfqpoint{1.962108in}{2.770068in}}{\pgfqpoint{1.957036in}{2.760172in}}%
\pgfpathquadraticcurveto{\pgfqpoint{1.951985in}{2.750297in}}{\pgfqpoint{1.951985in}{2.730444in}}%
\pgfpathquadraticcurveto{\pgfqpoint{1.951985in}{2.710673in}}{\pgfqpoint{1.957036in}{2.700777in}}%
\pgfpathquadraticcurveto{\pgfqpoint{1.962108in}{2.690881in}}{\pgfqpoint{1.972148in}{2.690881in}}%
\pgfpathquadraticcurveto{\pgfqpoint{1.982271in}{2.690881in}}{\pgfqpoint{1.987322in}{2.700777in}}%
\pgfpathquadraticcurveto{\pgfqpoint{1.992393in}{2.710673in}}{\pgfqpoint{1.992393in}{2.730444in}}%
\pgfpathquadraticcurveto{\pgfqpoint{1.992393in}{2.750297in}}{\pgfqpoint{1.987322in}{2.760172in}}%
\pgfpathquadraticcurveto{\pgfqpoint{1.982271in}{2.770068in}}{\pgfqpoint{1.972148in}{2.770068in}}%
\pgfpathlineto{\pgfqpoint{1.972148in}{2.770068in}}%
\pgfpathclose%
\pgfpathmoveto{\pgfqpoint{1.972148in}{2.780376in}}%
\pgfpathquadraticcurveto{\pgfqpoint{1.988332in}{2.780376in}}{\pgfqpoint{1.996867in}{2.767574in}}%
\pgfpathquadraticcurveto{\pgfqpoint{2.005402in}{2.754791in}}{\pgfqpoint{2.005402in}{2.730444in}}%
\pgfpathquadraticcurveto{\pgfqpoint{2.005402in}{2.706158in}}{\pgfqpoint{1.996867in}{2.693355in}}%
\pgfpathquadraticcurveto{\pgfqpoint{1.988332in}{2.680573in}}{\pgfqpoint{1.972148in}{2.680573in}}%
\pgfpathquadraticcurveto{\pgfqpoint{1.955985in}{2.680573in}}{\pgfqpoint{1.947450in}{2.693355in}}%
\pgfpathquadraticcurveto{\pgfqpoint{1.938914in}{2.706158in}}{\pgfqpoint{1.938914in}{2.730444in}}%
\pgfpathquadraticcurveto{\pgfqpoint{1.938914in}{2.754791in}}{\pgfqpoint{1.947450in}{2.767574in}}%
\pgfpathquadraticcurveto{\pgfqpoint{1.955985in}{2.780376in}}{\pgfqpoint{1.972148in}{2.780376in}}%
\pgfpathlineto{\pgfqpoint{1.972148in}{2.780376in}}%
\pgfpathclose%
\pgfpathmoveto{\pgfqpoint{2.030533in}{2.693396in}}%
\pgfpathlineto{\pgfqpoint{2.051789in}{2.693396in}}%
\pgfpathlineto{\pgfqpoint{2.051789in}{2.766790in}}%
\pgfpathlineto{\pgfqpoint{2.028657in}{2.762151in}}%
\pgfpathlineto{\pgfqpoint{2.028657in}{2.774006in}}%
\pgfpathlineto{\pgfqpoint{2.051665in}{2.778645in}}%
\pgfpathlineto{\pgfqpoint{2.064674in}{2.778645in}}%
\pgfpathlineto{\pgfqpoint{2.064674in}{2.693396in}}%
\pgfpathlineto{\pgfqpoint{2.085929in}{2.693396in}}%
\pgfpathlineto{\pgfqpoint{2.085929in}{2.682449in}}%
\pgfpathlineto{\pgfqpoint{2.030533in}{2.682449in}}%
\pgfpathlineto{\pgfqpoint{2.030533in}{2.693396in}}%
\pgfpathlineto{\pgfqpoint{2.030533in}{2.693396in}}%
\pgfpathclose%
\pgfpathmoveto{\pgfqpoint{2.108937in}{2.778645in}}%
\pgfpathlineto{\pgfqpoint{2.170786in}{2.778645in}}%
\pgfpathlineto{\pgfqpoint{2.170786in}{2.773099in}}%
\pgfpathlineto{\pgfqpoint{2.135862in}{2.682449in}}%
\pgfpathlineto{\pgfqpoint{2.122276in}{2.682449in}}%
\pgfpathlineto{\pgfqpoint{2.155138in}{2.767677in}}%
\pgfpathlineto{\pgfqpoint{2.108937in}{2.767677in}}%
\pgfpathlineto{\pgfqpoint{2.108937in}{2.778645in}}%
\pgfpathlineto{\pgfqpoint{2.108937in}{2.778645in}}%
\pgfpathclose%
\pgfpathmoveto{\pgfqpoint{2.192640in}{2.782562in}}%
\pgfpathlineto{\pgfqpoint{2.202948in}{2.782562in}}%
\pgfpathquadraticcurveto{\pgfqpoint{2.212596in}{2.767367in}}{\pgfqpoint{2.217400in}{2.752792in}}%
\pgfpathquadraticcurveto{\pgfqpoint{2.222203in}{2.738237in}}{\pgfqpoint{2.222203in}{2.723867in}}%
\pgfpathquadraticcurveto{\pgfqpoint{2.222203in}{2.709436in}}{\pgfqpoint{2.217400in}{2.694819in}}%
\pgfpathquadraticcurveto{\pgfqpoint{2.212596in}{2.680202in}}{\pgfqpoint{2.202948in}{2.665049in}}%
\pgfpathlineto{\pgfqpoint{2.192640in}{2.665049in}}%
\pgfpathquadraticcurveto{\pgfqpoint{2.201195in}{2.679810in}}{\pgfqpoint{2.205422in}{2.694406in}}%
\pgfpathquadraticcurveto{\pgfqpoint{2.209648in}{2.709003in}}{\pgfqpoint{2.209648in}{2.723867in}}%
\pgfpathquadraticcurveto{\pgfqpoint{2.209648in}{2.738752in}}{\pgfqpoint{2.205422in}{2.753245in}}%
\pgfpathquadraticcurveto{\pgfqpoint{2.201195in}{2.767759in}}{\pgfqpoint{2.192640in}{2.782562in}}%
\pgfpathlineto{\pgfqpoint{2.192640in}{2.782562in}}%
\pgfpathclose%
\pgfusepath{fill}%
\end{pgfscope}%
\begin{pgfscope}%
\pgfsetbuttcap%
\pgfsetroundjoin%
\definecolor{currentfill}{rgb}{0.090196,0.168627,0.227451}%
\pgfsetfillcolor{currentfill}%
\pgfsetlinewidth{0.000000pt}%
\definecolor{currentstroke}{rgb}{0.090196,0.168627,0.227451}%
\pgfsetstrokecolor{currentstroke}%
\pgfsetdash{}{0pt}%
\pgfpathmoveto{\pgfqpoint{1.316611in}{2.375890in}}%
\pgfpathlineto{\pgfqpoint{1.316611in}{2.363191in}}%
\pgfpathquadraticcurveto{\pgfqpoint{1.309210in}{2.366737in}}{\pgfqpoint{1.302633in}{2.368468in}}%
\pgfpathquadraticcurveto{\pgfqpoint{1.296056in}{2.370221in}}{\pgfqpoint{1.289933in}{2.370221in}}%
\pgfpathquadraticcurveto{\pgfqpoint{1.279316in}{2.370221in}}{\pgfqpoint{1.273543in}{2.366097in}}%
\pgfpathquadraticcurveto{\pgfqpoint{1.267771in}{2.361974in}}{\pgfqpoint{1.267771in}{2.354367in}}%
\pgfpathquadraticcurveto{\pgfqpoint{1.267771in}{2.347976in}}{\pgfqpoint{1.271605in}{2.344718in}}%
\pgfpathquadraticcurveto{\pgfqpoint{1.275440in}{2.341482in}}{\pgfqpoint{1.286140in}{2.339482in}}%
\pgfpathlineto{\pgfqpoint{1.293995in}{2.337874in}}%
\pgfpathquadraticcurveto{\pgfqpoint{1.308550in}{2.335091in}}{\pgfqpoint{1.315477in}{2.328102in}}%
\pgfpathquadraticcurveto{\pgfqpoint{1.322404in}{2.321113in}}{\pgfqpoint{1.322404in}{2.309403in}}%
\pgfpathquadraticcurveto{\pgfqpoint{1.322404in}{2.295404in}}{\pgfqpoint{1.313024in}{2.288188in}}%
\pgfpathquadraticcurveto{\pgfqpoint{1.303664in}{2.280973in}}{\pgfqpoint{1.285563in}{2.280973in}}%
\pgfpathquadraticcurveto{\pgfqpoint{1.278739in}{2.280973in}}{\pgfqpoint{1.271028in}{2.282519in}}%
\pgfpathquadraticcurveto{\pgfqpoint{1.263338in}{2.284065in}}{\pgfqpoint{1.255092in}{2.287096in}}%
\pgfpathlineto{\pgfqpoint{1.255092in}{2.300496in}}%
\pgfpathquadraticcurveto{\pgfqpoint{1.263008in}{2.296064in}}{\pgfqpoint{1.270616in}{2.293796in}}%
\pgfpathquadraticcurveto{\pgfqpoint{1.278223in}{2.291549in}}{\pgfqpoint{1.285563in}{2.291549in}}%
\pgfpathquadraticcurveto{\pgfqpoint{1.296696in}{2.291549in}}{\pgfqpoint{1.302757in}{2.295920in}}%
\pgfpathquadraticcurveto{\pgfqpoint{1.308818in}{2.300311in}}{\pgfqpoint{1.308818in}{2.308434in}}%
\pgfpathquadraticcurveto{\pgfqpoint{1.308818in}{2.315505in}}{\pgfqpoint{1.304468in}{2.319505in}}%
\pgfpathquadraticcurveto{\pgfqpoint{1.300118in}{2.323504in}}{\pgfqpoint{1.290201in}{2.325504in}}%
\pgfpathlineto{\pgfqpoint{1.282264in}{2.327050in}}%
\pgfpathquadraticcurveto{\pgfqpoint{1.267709in}{2.329936in}}{\pgfqpoint{1.261194in}{2.336121in}}%
\pgfpathquadraticcurveto{\pgfqpoint{1.254700in}{2.342306in}}{\pgfqpoint{1.254700in}{2.353336in}}%
\pgfpathquadraticcurveto{\pgfqpoint{1.254700in}{2.366097in}}{\pgfqpoint{1.263689in}{2.373437in}}%
\pgfpathquadraticcurveto{\pgfqpoint{1.272678in}{2.380776in}}{\pgfqpoint{1.288449in}{2.380776in}}%
\pgfpathquadraticcurveto{\pgfqpoint{1.295232in}{2.380776in}}{\pgfqpoint{1.302241in}{2.379539in}}%
\pgfpathquadraticcurveto{\pgfqpoint{1.309271in}{2.378323in}}{\pgfqpoint{1.316611in}{2.375890in}}%
\pgfpathlineto{\pgfqpoint{1.316611in}{2.375890in}}%
\pgfpathclose%
\pgfpathmoveto{\pgfqpoint{1.353926in}{2.375499in}}%
\pgfpathlineto{\pgfqpoint{1.353926in}{2.355006in}}%
\pgfpathlineto{\pgfqpoint{1.378336in}{2.355006in}}%
\pgfpathlineto{\pgfqpoint{1.378336in}{2.345790in}}%
\pgfpathlineto{\pgfqpoint{1.353926in}{2.345790in}}%
\pgfpathlineto{\pgfqpoint{1.353926in}{2.306619in}}%
\pgfpathquadraticcurveto{\pgfqpoint{1.353926in}{2.297796in}}{\pgfqpoint{1.356339in}{2.295280in}}%
\pgfpathquadraticcurveto{\pgfqpoint{1.358751in}{2.292765in}}{\pgfqpoint{1.366173in}{2.292765in}}%
\pgfpathlineto{\pgfqpoint{1.378336in}{2.292765in}}%
\pgfpathlineto{\pgfqpoint{1.378336in}{2.282849in}}%
\pgfpathlineto{\pgfqpoint{1.366173in}{2.282849in}}%
\pgfpathquadraticcurveto{\pgfqpoint{1.352442in}{2.282849in}}{\pgfqpoint{1.347226in}{2.287962in}}%
\pgfpathquadraticcurveto{\pgfqpoint{1.342010in}{2.293095in}}{\pgfqpoint{1.342010in}{2.306619in}}%
\pgfpathlineto{\pgfqpoint{1.342010in}{2.345790in}}%
\pgfpathlineto{\pgfqpoint{1.333310in}{2.345790in}}%
\pgfpathlineto{\pgfqpoint{1.333310in}{2.355006in}}%
\pgfpathlineto{\pgfqpoint{1.342010in}{2.355006in}}%
\pgfpathlineto{\pgfqpoint{1.342010in}{2.375499in}}%
\pgfpathlineto{\pgfqpoint{1.353926in}{2.375499in}}%
\pgfpathlineto{\pgfqpoint{1.353926in}{2.375499in}}%
\pgfpathclose%
\pgfpathmoveto{\pgfqpoint{1.392706in}{2.311320in}}%
\pgfpathlineto{\pgfqpoint{1.392706in}{2.355006in}}%
\pgfpathlineto{\pgfqpoint{1.404560in}{2.355006in}}%
\pgfpathlineto{\pgfqpoint{1.404560in}{2.311773in}}%
\pgfpathquadraticcurveto{\pgfqpoint{1.404560in}{2.301527in}}{\pgfqpoint{1.408539in}{2.296394in}}%
\pgfpathquadraticcurveto{\pgfqpoint{1.412539in}{2.291281in}}{\pgfqpoint{1.420538in}{2.291281in}}%
\pgfpathquadraticcurveto{\pgfqpoint{1.430124in}{2.291281in}}{\pgfqpoint{1.435691in}{2.297404in}}%
\pgfpathquadraticcurveto{\pgfqpoint{1.441278in}{2.303527in}}{\pgfqpoint{1.441278in}{2.314103in}}%
\pgfpathlineto{\pgfqpoint{1.441278in}{2.355006in}}%
\pgfpathlineto{\pgfqpoint{1.453132in}{2.355006in}}%
\pgfpathlineto{\pgfqpoint{1.453132in}{2.282849in}}%
\pgfpathlineto{\pgfqpoint{1.441278in}{2.282849in}}%
\pgfpathlineto{\pgfqpoint{1.441278in}{2.293940in}}%
\pgfpathquadraticcurveto{\pgfqpoint{1.436969in}{2.287364in}}{\pgfqpoint{1.431258in}{2.284168in}}%
\pgfpathquadraticcurveto{\pgfqpoint{1.425568in}{2.280973in}}{\pgfqpoint{1.418023in}{2.280973in}}%
\pgfpathquadraticcurveto{\pgfqpoint{1.405591in}{2.280973in}}{\pgfqpoint{1.399138in}{2.288704in}}%
\pgfpathquadraticcurveto{\pgfqpoint{1.392706in}{2.296435in}}{\pgfqpoint{1.392706in}{2.311320in}}%
\pgfpathlineto{\pgfqpoint{1.392706in}{2.311320in}}%
\pgfpathclose%
\pgfpathmoveto{\pgfqpoint{1.422538in}{2.356738in}}%
\pgfpathlineto{\pgfqpoint{1.422538in}{2.356738in}}%
\pgfpathlineto{\pgfqpoint{1.422538in}{2.356738in}}%
\pgfpathclose%
\pgfpathmoveto{\pgfqpoint{1.525021in}{2.344059in}}%
\pgfpathlineto{\pgfqpoint{1.525021in}{2.383106in}}%
\pgfpathlineto{\pgfqpoint{1.536876in}{2.383106in}}%
\pgfpathlineto{\pgfqpoint{1.536876in}{2.282849in}}%
\pgfpathlineto{\pgfqpoint{1.525021in}{2.282849in}}%
\pgfpathlineto{\pgfqpoint{1.525021in}{2.293672in}}%
\pgfpathquadraticcurveto{\pgfqpoint{1.521290in}{2.287240in}}{\pgfqpoint{1.515579in}{2.284106in}}%
\pgfpathquadraticcurveto{\pgfqpoint{1.509889in}{2.280973in}}{\pgfqpoint{1.501890in}{2.280973in}}%
\pgfpathquadraticcurveto{\pgfqpoint{1.488819in}{2.280973in}}{\pgfqpoint{1.480593in}{2.291405in}}%
\pgfpathquadraticcurveto{\pgfqpoint{1.472388in}{2.301857in}}{\pgfqpoint{1.472388in}{2.318866in}}%
\pgfpathquadraticcurveto{\pgfqpoint{1.472388in}{2.335874in}}{\pgfqpoint{1.480593in}{2.346306in}}%
\pgfpathquadraticcurveto{\pgfqpoint{1.488819in}{2.356738in}}{\pgfqpoint{1.501890in}{2.356738in}}%
\pgfpathquadraticcurveto{\pgfqpoint{1.509889in}{2.356738in}}{\pgfqpoint{1.515579in}{2.353604in}}%
\pgfpathquadraticcurveto{\pgfqpoint{1.521290in}{2.350491in}}{\pgfqpoint{1.525021in}{2.344059in}}%
\pgfpathlineto{\pgfqpoint{1.525021in}{2.344059in}}%
\pgfpathclose%
\pgfpathmoveto{\pgfqpoint{1.484634in}{2.318866in}}%
\pgfpathquadraticcurveto{\pgfqpoint{1.484634in}{2.305795in}}{\pgfqpoint{1.490015in}{2.298352in}}%
\pgfpathquadraticcurveto{\pgfqpoint{1.495396in}{2.290910in}}{\pgfqpoint{1.504797in}{2.290910in}}%
\pgfpathquadraticcurveto{\pgfqpoint{1.514198in}{2.290910in}}{\pgfqpoint{1.519599in}{2.298352in}}%
\pgfpathquadraticcurveto{\pgfqpoint{1.525021in}{2.305795in}}{\pgfqpoint{1.525021in}{2.318866in}}%
\pgfpathquadraticcurveto{\pgfqpoint{1.525021in}{2.331936in}}{\pgfqpoint{1.519599in}{2.339379in}}%
\pgfpathquadraticcurveto{\pgfqpoint{1.514198in}{2.346821in}}{\pgfqpoint{1.504797in}{2.346821in}}%
\pgfpathquadraticcurveto{\pgfqpoint{1.495396in}{2.346821in}}{\pgfqpoint{1.490015in}{2.339379in}}%
\pgfpathquadraticcurveto{\pgfqpoint{1.484634in}{2.331936in}}{\pgfqpoint{1.484634in}{2.318866in}}%
\pgfpathlineto{\pgfqpoint{1.484634in}{2.318866in}}%
\pgfpathclose%
\pgfpathmoveto{\pgfqpoint{1.591323in}{2.276148in}}%
\pgfpathquadraticcurveto{\pgfqpoint{1.586314in}{2.263263in}}{\pgfqpoint{1.581531in}{2.259346in}}%
\pgfpathquadraticcurveto{\pgfqpoint{1.576768in}{2.255408in}}{\pgfqpoint{1.568790in}{2.255408in}}%
\pgfpathlineto{\pgfqpoint{1.559306in}{2.255408in}}%
\pgfpathlineto{\pgfqpoint{1.559306in}{2.265325in}}%
\pgfpathlineto{\pgfqpoint{1.566275in}{2.265325in}}%
\pgfpathquadraticcurveto{\pgfqpoint{1.571161in}{2.265325in}}{\pgfqpoint{1.573861in}{2.267655in}}%
\pgfpathquadraticcurveto{\pgfqpoint{1.576583in}{2.269964in}}{\pgfqpoint{1.579861in}{2.278602in}}%
\pgfpathlineto{\pgfqpoint{1.581984in}{2.284003in}}%
\pgfpathlineto{\pgfqpoint{1.552812in}{2.355006in}}%
\pgfpathlineto{\pgfqpoint{1.565367in}{2.355006in}}%
\pgfpathlineto{\pgfqpoint{1.587922in}{2.298579in}}%
\pgfpathlineto{\pgfqpoint{1.610476in}{2.355006in}}%
\pgfpathlineto{\pgfqpoint{1.623031in}{2.355006in}}%
\pgfpathlineto{\pgfqpoint{1.591323in}{2.276148in}}%
\pgfpathlineto{\pgfqpoint{1.591323in}{2.276148in}}%
\pgfpathclose%
\pgfpathmoveto{\pgfqpoint{1.718773in}{2.367706in}}%
\pgfpathlineto{\pgfqpoint{1.685911in}{2.316350in}}%
\pgfpathlineto{\pgfqpoint{1.718773in}{2.316350in}}%
\pgfpathlineto{\pgfqpoint{1.718773in}{2.367706in}}%
\pgfpathlineto{\pgfqpoint{1.718773in}{2.367706in}}%
\pgfpathclose%
\pgfpathmoveto{\pgfqpoint{1.715351in}{2.379045in}}%
\pgfpathlineto{\pgfqpoint{1.731720in}{2.379045in}}%
\pgfpathlineto{\pgfqpoint{1.731720in}{2.316350in}}%
\pgfpathlineto{\pgfqpoint{1.745451in}{2.316350in}}%
\pgfpathlineto{\pgfqpoint{1.745451in}{2.305527in}}%
\pgfpathlineto{\pgfqpoint{1.731720in}{2.305527in}}%
\pgfpathlineto{\pgfqpoint{1.731720in}{2.282849in}}%
\pgfpathlineto{\pgfqpoint{1.718773in}{2.282849in}}%
\pgfpathlineto{\pgfqpoint{1.718773in}{2.305527in}}%
\pgfpathlineto{\pgfqpoint{1.675355in}{2.305527in}}%
\pgfpathlineto{\pgfqpoint{1.675355in}{2.318082in}}%
\pgfpathlineto{\pgfqpoint{1.715351in}{2.379045in}}%
\pgfpathlineto{\pgfqpoint{1.715351in}{2.379045in}}%
\pgfpathclose%
\pgfpathmoveto{\pgfqpoint{1.835689in}{2.382962in}}%
\pgfpathquadraticcurveto{\pgfqpoint{1.827071in}{2.368159in}}{\pgfqpoint{1.822865in}{2.353645in}}%
\pgfpathquadraticcurveto{\pgfqpoint{1.818680in}{2.339152in}}{\pgfqpoint{1.818680in}{2.324267in}}%
\pgfpathquadraticcurveto{\pgfqpoint{1.818680in}{2.309403in}}{\pgfqpoint{1.822906in}{2.294806in}}%
\pgfpathquadraticcurveto{\pgfqpoint{1.827133in}{2.280210in}}{\pgfqpoint{1.835689in}{2.265449in}}%
\pgfpathlineto{\pgfqpoint{1.825380in}{2.265449in}}%
\pgfpathquadraticcurveto{\pgfqpoint{1.815732in}{2.280602in}}{\pgfqpoint{1.810928in}{2.295219in}}%
\pgfpathquadraticcurveto{\pgfqpoint{1.806125in}{2.309836in}}{\pgfqpoint{1.806125in}{2.324267in}}%
\pgfpathquadraticcurveto{\pgfqpoint{1.806125in}{2.338637in}}{\pgfqpoint{1.810887in}{2.353192in}}%
\pgfpathquadraticcurveto{\pgfqpoint{1.815670in}{2.367767in}}{\pgfqpoint{1.825380in}{2.382962in}}%
\pgfpathlineto{\pgfqpoint{1.835689in}{2.382962in}}%
\pgfpathlineto{\pgfqpoint{1.835689in}{2.382962in}}%
\pgfpathclose%
\pgfpathmoveto{\pgfqpoint{1.871582in}{2.293796in}}%
\pgfpathlineto{\pgfqpoint{1.916999in}{2.293796in}}%
\pgfpathlineto{\pgfqpoint{1.916999in}{2.282849in}}%
\pgfpathlineto{\pgfqpoint{1.855934in}{2.282849in}}%
\pgfpathlineto{\pgfqpoint{1.855934in}{2.293796in}}%
\pgfpathquadraticcurveto{\pgfqpoint{1.863335in}{2.301465in}}{\pgfqpoint{1.876117in}{2.314371in}}%
\pgfpathquadraticcurveto{\pgfqpoint{1.888920in}{2.327298in}}{\pgfqpoint{1.892198in}{2.331050in}}%
\pgfpathquadraticcurveto{\pgfqpoint{1.898445in}{2.338059in}}{\pgfqpoint{1.900919in}{2.342925in}}%
\pgfpathquadraticcurveto{\pgfqpoint{1.903413in}{2.347790in}}{\pgfqpoint{1.903413in}{2.352491in}}%
\pgfpathquadraticcurveto{\pgfqpoint{1.903413in}{2.360160in}}{\pgfqpoint{1.898032in}{2.364984in}}%
\pgfpathquadraticcurveto{\pgfqpoint{1.892651in}{2.369829in}}{\pgfqpoint{1.884013in}{2.369829in}}%
\pgfpathquadraticcurveto{\pgfqpoint{1.877890in}{2.369829in}}{\pgfqpoint{1.871087in}{2.367706in}}%
\pgfpathquadraticcurveto{\pgfqpoint{1.864304in}{2.365582in}}{\pgfqpoint{1.856573in}{2.361253in}}%
\pgfpathlineto{\pgfqpoint{1.856573in}{2.374406in}}%
\pgfpathquadraticcurveto{\pgfqpoint{1.864428in}{2.377560in}}{\pgfqpoint{1.871252in}{2.379168in}}%
\pgfpathquadraticcurveto{\pgfqpoint{1.878096in}{2.380776in}}{\pgfqpoint{1.883766in}{2.380776in}}%
\pgfpathquadraticcurveto{\pgfqpoint{1.898713in}{2.380776in}}{\pgfqpoint{1.907598in}{2.373293in}}%
\pgfpathquadraticcurveto{\pgfqpoint{1.916484in}{2.365829in}}{\pgfqpoint{1.916484in}{2.353336in}}%
\pgfpathquadraticcurveto{\pgfqpoint{1.916484in}{2.347398in}}{\pgfqpoint{1.914257in}{2.342079in}}%
\pgfpathquadraticcurveto{\pgfqpoint{1.912051in}{2.336781in}}{\pgfqpoint{1.906176in}{2.329565in}}%
\pgfpathquadraticcurveto{\pgfqpoint{1.904568in}{2.327689in}}{\pgfqpoint{1.895929in}{2.318762in}}%
\pgfpathquadraticcurveto{\pgfqpoint{1.887312in}{2.309836in}}{\pgfqpoint{1.871582in}{2.293796in}}%
\pgfpathlineto{\pgfqpoint{1.871582in}{2.293796in}}%
\pgfpathclose%
\pgfpathmoveto{\pgfqpoint{1.972148in}{2.370468in}}%
\pgfpathquadraticcurveto{\pgfqpoint{1.962108in}{2.370468in}}{\pgfqpoint{1.957036in}{2.360572in}}%
\pgfpathquadraticcurveto{\pgfqpoint{1.951985in}{2.350697in}}{\pgfqpoint{1.951985in}{2.330844in}}%
\pgfpathquadraticcurveto{\pgfqpoint{1.951985in}{2.311073in}}{\pgfqpoint{1.957036in}{2.301177in}}%
\pgfpathquadraticcurveto{\pgfqpoint{1.962108in}{2.291281in}}{\pgfqpoint{1.972148in}{2.291281in}}%
\pgfpathquadraticcurveto{\pgfqpoint{1.982271in}{2.291281in}}{\pgfqpoint{1.987322in}{2.301177in}}%
\pgfpathquadraticcurveto{\pgfqpoint{1.992393in}{2.311073in}}{\pgfqpoint{1.992393in}{2.330844in}}%
\pgfpathquadraticcurveto{\pgfqpoint{1.992393in}{2.350697in}}{\pgfqpoint{1.987322in}{2.360572in}}%
\pgfpathquadraticcurveto{\pgfqpoint{1.982271in}{2.370468in}}{\pgfqpoint{1.972148in}{2.370468in}}%
\pgfpathlineto{\pgfqpoint{1.972148in}{2.370468in}}%
\pgfpathclose%
\pgfpathmoveto{\pgfqpoint{1.972148in}{2.380776in}}%
\pgfpathquadraticcurveto{\pgfqpoint{1.988332in}{2.380776in}}{\pgfqpoint{1.996867in}{2.367974in}}%
\pgfpathquadraticcurveto{\pgfqpoint{2.005402in}{2.355191in}}{\pgfqpoint{2.005402in}{2.330844in}}%
\pgfpathquadraticcurveto{\pgfqpoint{2.005402in}{2.306558in}}{\pgfqpoint{1.996867in}{2.293755in}}%
\pgfpathquadraticcurveto{\pgfqpoint{1.988332in}{2.280973in}}{\pgfqpoint{1.972148in}{2.280973in}}%
\pgfpathquadraticcurveto{\pgfqpoint{1.955985in}{2.280973in}}{\pgfqpoint{1.947450in}{2.293755in}}%
\pgfpathquadraticcurveto{\pgfqpoint{1.938914in}{2.306558in}}{\pgfqpoint{1.938914in}{2.330844in}}%
\pgfpathquadraticcurveto{\pgfqpoint{1.938914in}{2.355191in}}{\pgfqpoint{1.947450in}{2.367974in}}%
\pgfpathquadraticcurveto{\pgfqpoint{1.955985in}{2.380776in}}{\pgfqpoint{1.972148in}{2.380776in}}%
\pgfpathlineto{\pgfqpoint{1.972148in}{2.380776in}}%
\pgfpathclose%
\pgfpathmoveto{\pgfqpoint{2.030533in}{2.293796in}}%
\pgfpathlineto{\pgfqpoint{2.051789in}{2.293796in}}%
\pgfpathlineto{\pgfqpoint{2.051789in}{2.367190in}}%
\pgfpathlineto{\pgfqpoint{2.028657in}{2.362551in}}%
\pgfpathlineto{\pgfqpoint{2.028657in}{2.374406in}}%
\pgfpathlineto{\pgfqpoint{2.051665in}{2.379045in}}%
\pgfpathlineto{\pgfqpoint{2.064674in}{2.379045in}}%
\pgfpathlineto{\pgfqpoint{2.064674in}{2.293796in}}%
\pgfpathlineto{\pgfqpoint{2.085929in}{2.293796in}}%
\pgfpathlineto{\pgfqpoint{2.085929in}{2.282849in}}%
\pgfpathlineto{\pgfqpoint{2.030533in}{2.282849in}}%
\pgfpathlineto{\pgfqpoint{2.030533in}{2.293796in}}%
\pgfpathlineto{\pgfqpoint{2.030533in}{2.293796in}}%
\pgfpathclose%
\pgfpathmoveto{\pgfqpoint{2.140047in}{2.328535in}}%
\pgfpathquadraticcurveto{\pgfqpoint{2.130770in}{2.328535in}}{\pgfqpoint{2.125451in}{2.323566in}}%
\pgfpathquadraticcurveto{\pgfqpoint{2.120153in}{2.318597in}}{\pgfqpoint{2.120153in}{2.309918in}}%
\pgfpathquadraticcurveto{\pgfqpoint{2.120153in}{2.301218in}}{\pgfqpoint{2.125451in}{2.296249in}}%
\pgfpathquadraticcurveto{\pgfqpoint{2.130770in}{2.291281in}}{\pgfqpoint{2.140047in}{2.291281in}}%
\pgfpathquadraticcurveto{\pgfqpoint{2.149325in}{2.291281in}}{\pgfqpoint{2.154664in}{2.296270in}}%
\pgfpathquadraticcurveto{\pgfqpoint{2.160025in}{2.301280in}}{\pgfqpoint{2.160025in}{2.309918in}}%
\pgfpathquadraticcurveto{\pgfqpoint{2.160025in}{2.318597in}}{\pgfqpoint{2.154706in}{2.323566in}}%
\pgfpathquadraticcurveto{\pgfqpoint{2.149407in}{2.328535in}}{\pgfqpoint{2.140047in}{2.328535in}}%
\pgfpathlineto{\pgfqpoint{2.140047in}{2.328535in}}%
\pgfpathclose%
\pgfpathmoveto{\pgfqpoint{2.127038in}{2.334060in}}%
\pgfpathquadraticcurveto{\pgfqpoint{2.118668in}{2.336121in}}{\pgfqpoint{2.113988in}{2.341853in}}%
\pgfpathquadraticcurveto{\pgfqpoint{2.109329in}{2.347605in}}{\pgfqpoint{2.109329in}{2.355851in}}%
\pgfpathquadraticcurveto{\pgfqpoint{2.109329in}{2.367376in}}{\pgfqpoint{2.117534in}{2.374076in}}%
\pgfpathquadraticcurveto{\pgfqpoint{2.125760in}{2.380776in}}{\pgfqpoint{2.140047in}{2.380776in}}%
\pgfpathquadraticcurveto{\pgfqpoint{2.154417in}{2.380776in}}{\pgfqpoint{2.162602in}{2.374076in}}%
\pgfpathquadraticcurveto{\pgfqpoint{2.170786in}{2.367376in}}{\pgfqpoint{2.170786in}{2.355851in}}%
\pgfpathquadraticcurveto{\pgfqpoint{2.170786in}{2.347605in}}{\pgfqpoint{2.166106in}{2.341853in}}%
\pgfpathquadraticcurveto{\pgfqpoint{2.161447in}{2.336121in}}{\pgfqpoint{2.153139in}{2.334060in}}%
\pgfpathquadraticcurveto{\pgfqpoint{2.162540in}{2.331874in}}{\pgfqpoint{2.167776in}{2.325483in}}%
\pgfpathquadraticcurveto{\pgfqpoint{2.173033in}{2.319113in}}{\pgfqpoint{2.173033in}{2.309918in}}%
\pgfpathquadraticcurveto{\pgfqpoint{2.173033in}{2.295920in}}{\pgfqpoint{2.164498in}{2.288436in}}%
\pgfpathquadraticcurveto{\pgfqpoint{2.155963in}{2.280973in}}{\pgfqpoint{2.140047in}{2.280973in}}%
\pgfpathquadraticcurveto{\pgfqpoint{2.124152in}{2.280973in}}{\pgfqpoint{2.115596in}{2.288436in}}%
\pgfpathquadraticcurveto{\pgfqpoint{2.107061in}{2.295920in}}{\pgfqpoint{2.107061in}{2.309918in}}%
\pgfpathquadraticcurveto{\pgfqpoint{2.107061in}{2.319113in}}{\pgfqpoint{2.112339in}{2.325483in}}%
\pgfpathquadraticcurveto{\pgfqpoint{2.117637in}{2.331874in}}{\pgfqpoint{2.127038in}{2.334060in}}%
\pgfpathlineto{\pgfqpoint{2.127038in}{2.334060in}}%
\pgfpathclose%
\pgfpathmoveto{\pgfqpoint{2.122276in}{2.354614in}}%
\pgfpathquadraticcurveto{\pgfqpoint{2.122276in}{2.347151in}}{\pgfqpoint{2.126935in}{2.342966in}}%
\pgfpathquadraticcurveto{\pgfqpoint{2.131615in}{2.338781in}}{\pgfqpoint{2.140047in}{2.338781in}}%
\pgfpathquadraticcurveto{\pgfqpoint{2.148438in}{2.338781in}}{\pgfqpoint{2.153159in}{2.342966in}}%
\pgfpathquadraticcurveto{\pgfqpoint{2.157901in}{2.347151in}}{\pgfqpoint{2.157901in}{2.354614in}}%
\pgfpathquadraticcurveto{\pgfqpoint{2.157901in}{2.362098in}}{\pgfqpoint{2.153159in}{2.366283in}}%
\pgfpathquadraticcurveto{\pgfqpoint{2.148438in}{2.370468in}}{\pgfqpoint{2.140047in}{2.370468in}}%
\pgfpathquadraticcurveto{\pgfqpoint{2.131615in}{2.370468in}}{\pgfqpoint{2.126935in}{2.366283in}}%
\pgfpathquadraticcurveto{\pgfqpoint{2.122276in}{2.362098in}}{\pgfqpoint{2.122276in}{2.354614in}}%
\pgfpathlineto{\pgfqpoint{2.122276in}{2.354614in}}%
\pgfpathclose%
\pgfpathmoveto{\pgfqpoint{2.192640in}{2.382962in}}%
\pgfpathlineto{\pgfqpoint{2.202948in}{2.382962in}}%
\pgfpathquadraticcurveto{\pgfqpoint{2.212596in}{2.367767in}}{\pgfqpoint{2.217400in}{2.353192in}}%
\pgfpathquadraticcurveto{\pgfqpoint{2.222203in}{2.338637in}}{\pgfqpoint{2.222203in}{2.324267in}}%
\pgfpathquadraticcurveto{\pgfqpoint{2.222203in}{2.309836in}}{\pgfqpoint{2.217400in}{2.295219in}}%
\pgfpathquadraticcurveto{\pgfqpoint{2.212596in}{2.280602in}}{\pgfqpoint{2.202948in}{2.265449in}}%
\pgfpathlineto{\pgfqpoint{2.192640in}{2.265449in}}%
\pgfpathquadraticcurveto{\pgfqpoint{2.201195in}{2.280210in}}{\pgfqpoint{2.205422in}{2.294806in}}%
\pgfpathquadraticcurveto{\pgfqpoint{2.209648in}{2.309403in}}{\pgfqpoint{2.209648in}{2.324267in}}%
\pgfpathquadraticcurveto{\pgfqpoint{2.209648in}{2.339152in}}{\pgfqpoint{2.205422in}{2.353645in}}%
\pgfpathquadraticcurveto{\pgfqpoint{2.201195in}{2.368159in}}{\pgfqpoint{2.192640in}{2.382962in}}%
\pgfpathlineto{\pgfqpoint{2.192640in}{2.382962in}}%
\pgfpathclose%
\pgfusepath{fill}%
\end{pgfscope}%
\begin{pgfscope}%
\pgfsetbuttcap%
\pgfsetroundjoin%
\definecolor{currentfill}{rgb}{0.090196,0.168627,0.227451}%
\pgfsetfillcolor{currentfill}%
\pgfsetlinewidth{0.000000pt}%
\definecolor{currentstroke}{rgb}{0.090196,0.168627,0.227451}%
\pgfsetstrokecolor{currentstroke}%
\pgfsetdash{}{0pt}%
\pgfpathmoveto{\pgfqpoint{1.316611in}{1.976290in}}%
\pgfpathlineto{\pgfqpoint{1.316611in}{1.963591in}}%
\pgfpathquadraticcurveto{\pgfqpoint{1.309210in}{1.967137in}}{\pgfqpoint{1.302633in}{1.968868in}}%
\pgfpathquadraticcurveto{\pgfqpoint{1.296056in}{1.970621in}}{\pgfqpoint{1.289933in}{1.970621in}}%
\pgfpathquadraticcurveto{\pgfqpoint{1.279316in}{1.970621in}}{\pgfqpoint{1.273543in}{1.966497in}}%
\pgfpathquadraticcurveto{\pgfqpoint{1.267771in}{1.962374in}}{\pgfqpoint{1.267771in}{1.954767in}}%
\pgfpathquadraticcurveto{\pgfqpoint{1.267771in}{1.948376in}}{\pgfqpoint{1.271605in}{1.945118in}}%
\pgfpathquadraticcurveto{\pgfqpoint{1.275440in}{1.941882in}}{\pgfqpoint{1.286140in}{1.939882in}}%
\pgfpathlineto{\pgfqpoint{1.293995in}{1.938274in}}%
\pgfpathquadraticcurveto{\pgfqpoint{1.308550in}{1.935491in}}{\pgfqpoint{1.315477in}{1.928502in}}%
\pgfpathquadraticcurveto{\pgfqpoint{1.322404in}{1.921513in}}{\pgfqpoint{1.322404in}{1.909803in}}%
\pgfpathquadraticcurveto{\pgfqpoint{1.322404in}{1.895804in}}{\pgfqpoint{1.313024in}{1.888588in}}%
\pgfpathquadraticcurveto{\pgfqpoint{1.303664in}{1.881373in}}{\pgfqpoint{1.285563in}{1.881373in}}%
\pgfpathquadraticcurveto{\pgfqpoint{1.278739in}{1.881373in}}{\pgfqpoint{1.271028in}{1.882919in}}%
\pgfpathquadraticcurveto{\pgfqpoint{1.263338in}{1.884465in}}{\pgfqpoint{1.255092in}{1.887496in}}%
\pgfpathlineto{\pgfqpoint{1.255092in}{1.900896in}}%
\pgfpathquadraticcurveto{\pgfqpoint{1.263008in}{1.896464in}}{\pgfqpoint{1.270616in}{1.894196in}}%
\pgfpathquadraticcurveto{\pgfqpoint{1.278223in}{1.891949in}}{\pgfqpoint{1.285563in}{1.891949in}}%
\pgfpathquadraticcurveto{\pgfqpoint{1.296696in}{1.891949in}}{\pgfqpoint{1.302757in}{1.896320in}}%
\pgfpathquadraticcurveto{\pgfqpoint{1.308818in}{1.900711in}}{\pgfqpoint{1.308818in}{1.908834in}}%
\pgfpathquadraticcurveto{\pgfqpoint{1.308818in}{1.915905in}}{\pgfqpoint{1.304468in}{1.919905in}}%
\pgfpathquadraticcurveto{\pgfqpoint{1.300118in}{1.923904in}}{\pgfqpoint{1.290201in}{1.925904in}}%
\pgfpathlineto{\pgfqpoint{1.282264in}{1.927450in}}%
\pgfpathquadraticcurveto{\pgfqpoint{1.267709in}{1.930336in}}{\pgfqpoint{1.261194in}{1.936521in}}%
\pgfpathquadraticcurveto{\pgfqpoint{1.254700in}{1.942706in}}{\pgfqpoint{1.254700in}{1.953736in}}%
\pgfpathquadraticcurveto{\pgfqpoint{1.254700in}{1.966497in}}{\pgfqpoint{1.263689in}{1.973837in}}%
\pgfpathquadraticcurveto{\pgfqpoint{1.272678in}{1.981176in}}{\pgfqpoint{1.288449in}{1.981176in}}%
\pgfpathquadraticcurveto{\pgfqpoint{1.295232in}{1.981176in}}{\pgfqpoint{1.302241in}{1.979939in}}%
\pgfpathquadraticcurveto{\pgfqpoint{1.309271in}{1.978723in}}{\pgfqpoint{1.316611in}{1.976290in}}%
\pgfpathlineto{\pgfqpoint{1.316611in}{1.976290in}}%
\pgfpathclose%
\pgfpathmoveto{\pgfqpoint{1.353926in}{1.975899in}}%
\pgfpathlineto{\pgfqpoint{1.353926in}{1.955406in}}%
\pgfpathlineto{\pgfqpoint{1.378336in}{1.955406in}}%
\pgfpathlineto{\pgfqpoint{1.378336in}{1.946190in}}%
\pgfpathlineto{\pgfqpoint{1.353926in}{1.946190in}}%
\pgfpathlineto{\pgfqpoint{1.353926in}{1.907019in}}%
\pgfpathquadraticcurveto{\pgfqpoint{1.353926in}{1.898196in}}{\pgfqpoint{1.356339in}{1.895680in}}%
\pgfpathquadraticcurveto{\pgfqpoint{1.358751in}{1.893165in}}{\pgfqpoint{1.366173in}{1.893165in}}%
\pgfpathlineto{\pgfqpoint{1.378336in}{1.893165in}}%
\pgfpathlineto{\pgfqpoint{1.378336in}{1.883249in}}%
\pgfpathlineto{\pgfqpoint{1.366173in}{1.883249in}}%
\pgfpathquadraticcurveto{\pgfqpoint{1.352442in}{1.883249in}}{\pgfqpoint{1.347226in}{1.888362in}}%
\pgfpathquadraticcurveto{\pgfqpoint{1.342010in}{1.893495in}}{\pgfqpoint{1.342010in}{1.907019in}}%
\pgfpathlineto{\pgfqpoint{1.342010in}{1.946190in}}%
\pgfpathlineto{\pgfqpoint{1.333310in}{1.946190in}}%
\pgfpathlineto{\pgfqpoint{1.333310in}{1.955406in}}%
\pgfpathlineto{\pgfqpoint{1.342010in}{1.955406in}}%
\pgfpathlineto{\pgfqpoint{1.342010in}{1.975899in}}%
\pgfpathlineto{\pgfqpoint{1.353926in}{1.975899in}}%
\pgfpathlineto{\pgfqpoint{1.353926in}{1.975899in}}%
\pgfpathclose%
\pgfpathmoveto{\pgfqpoint{1.392706in}{1.911720in}}%
\pgfpathlineto{\pgfqpoint{1.392706in}{1.955406in}}%
\pgfpathlineto{\pgfqpoint{1.404560in}{1.955406in}}%
\pgfpathlineto{\pgfqpoint{1.404560in}{1.912173in}}%
\pgfpathquadraticcurveto{\pgfqpoint{1.404560in}{1.901927in}}{\pgfqpoint{1.408539in}{1.896794in}}%
\pgfpathquadraticcurveto{\pgfqpoint{1.412539in}{1.891681in}}{\pgfqpoint{1.420538in}{1.891681in}}%
\pgfpathquadraticcurveto{\pgfqpoint{1.430124in}{1.891681in}}{\pgfqpoint{1.435691in}{1.897804in}}%
\pgfpathquadraticcurveto{\pgfqpoint{1.441278in}{1.903927in}}{\pgfqpoint{1.441278in}{1.914503in}}%
\pgfpathlineto{\pgfqpoint{1.441278in}{1.955406in}}%
\pgfpathlineto{\pgfqpoint{1.453132in}{1.955406in}}%
\pgfpathlineto{\pgfqpoint{1.453132in}{1.883249in}}%
\pgfpathlineto{\pgfqpoint{1.441278in}{1.883249in}}%
\pgfpathlineto{\pgfqpoint{1.441278in}{1.894340in}}%
\pgfpathquadraticcurveto{\pgfqpoint{1.436969in}{1.887764in}}{\pgfqpoint{1.431258in}{1.884568in}}%
\pgfpathquadraticcurveto{\pgfqpoint{1.425568in}{1.881373in}}{\pgfqpoint{1.418023in}{1.881373in}}%
\pgfpathquadraticcurveto{\pgfqpoint{1.405591in}{1.881373in}}{\pgfqpoint{1.399138in}{1.889104in}}%
\pgfpathquadraticcurveto{\pgfqpoint{1.392706in}{1.896835in}}{\pgfqpoint{1.392706in}{1.911720in}}%
\pgfpathlineto{\pgfqpoint{1.392706in}{1.911720in}}%
\pgfpathclose%
\pgfpathmoveto{\pgfqpoint{1.422538in}{1.957138in}}%
\pgfpathlineto{\pgfqpoint{1.422538in}{1.957138in}}%
\pgfpathlineto{\pgfqpoint{1.422538in}{1.957138in}}%
\pgfpathclose%
\pgfpathmoveto{\pgfqpoint{1.525021in}{1.944459in}}%
\pgfpathlineto{\pgfqpoint{1.525021in}{1.983506in}}%
\pgfpathlineto{\pgfqpoint{1.536876in}{1.983506in}}%
\pgfpathlineto{\pgfqpoint{1.536876in}{1.883249in}}%
\pgfpathlineto{\pgfqpoint{1.525021in}{1.883249in}}%
\pgfpathlineto{\pgfqpoint{1.525021in}{1.894072in}}%
\pgfpathquadraticcurveto{\pgfqpoint{1.521290in}{1.887640in}}{\pgfqpoint{1.515579in}{1.884506in}}%
\pgfpathquadraticcurveto{\pgfqpoint{1.509889in}{1.881373in}}{\pgfqpoint{1.501890in}{1.881373in}}%
\pgfpathquadraticcurveto{\pgfqpoint{1.488819in}{1.881373in}}{\pgfqpoint{1.480593in}{1.891805in}}%
\pgfpathquadraticcurveto{\pgfqpoint{1.472388in}{1.902257in}}{\pgfqpoint{1.472388in}{1.919266in}}%
\pgfpathquadraticcurveto{\pgfqpoint{1.472388in}{1.936274in}}{\pgfqpoint{1.480593in}{1.946706in}}%
\pgfpathquadraticcurveto{\pgfqpoint{1.488819in}{1.957138in}}{\pgfqpoint{1.501890in}{1.957138in}}%
\pgfpathquadraticcurveto{\pgfqpoint{1.509889in}{1.957138in}}{\pgfqpoint{1.515579in}{1.954004in}}%
\pgfpathquadraticcurveto{\pgfqpoint{1.521290in}{1.950891in}}{\pgfqpoint{1.525021in}{1.944459in}}%
\pgfpathlineto{\pgfqpoint{1.525021in}{1.944459in}}%
\pgfpathclose%
\pgfpathmoveto{\pgfqpoint{1.484634in}{1.919266in}}%
\pgfpathquadraticcurveto{\pgfqpoint{1.484634in}{1.906195in}}{\pgfqpoint{1.490015in}{1.898752in}}%
\pgfpathquadraticcurveto{\pgfqpoint{1.495396in}{1.891310in}}{\pgfqpoint{1.504797in}{1.891310in}}%
\pgfpathquadraticcurveto{\pgfqpoint{1.514198in}{1.891310in}}{\pgfqpoint{1.519599in}{1.898752in}}%
\pgfpathquadraticcurveto{\pgfqpoint{1.525021in}{1.906195in}}{\pgfqpoint{1.525021in}{1.919266in}}%
\pgfpathquadraticcurveto{\pgfqpoint{1.525021in}{1.932336in}}{\pgfqpoint{1.519599in}{1.939779in}}%
\pgfpathquadraticcurveto{\pgfqpoint{1.514198in}{1.947221in}}{\pgfqpoint{1.504797in}{1.947221in}}%
\pgfpathquadraticcurveto{\pgfqpoint{1.495396in}{1.947221in}}{\pgfqpoint{1.490015in}{1.939779in}}%
\pgfpathquadraticcurveto{\pgfqpoint{1.484634in}{1.932336in}}{\pgfqpoint{1.484634in}{1.919266in}}%
\pgfpathlineto{\pgfqpoint{1.484634in}{1.919266in}}%
\pgfpathclose%
\pgfpathmoveto{\pgfqpoint{1.591323in}{1.876548in}}%
\pgfpathquadraticcurveto{\pgfqpoint{1.586314in}{1.863663in}}{\pgfqpoint{1.581531in}{1.859746in}}%
\pgfpathquadraticcurveto{\pgfqpoint{1.576768in}{1.855808in}}{\pgfqpoint{1.568790in}{1.855808in}}%
\pgfpathlineto{\pgfqpoint{1.559306in}{1.855808in}}%
\pgfpathlineto{\pgfqpoint{1.559306in}{1.865725in}}%
\pgfpathlineto{\pgfqpoint{1.566275in}{1.865725in}}%
\pgfpathquadraticcurveto{\pgfqpoint{1.571161in}{1.865725in}}{\pgfqpoint{1.573861in}{1.868055in}}%
\pgfpathquadraticcurveto{\pgfqpoint{1.576583in}{1.870364in}}{\pgfqpoint{1.579861in}{1.879002in}}%
\pgfpathlineto{\pgfqpoint{1.581984in}{1.884403in}}%
\pgfpathlineto{\pgfqpoint{1.552812in}{1.955406in}}%
\pgfpathlineto{\pgfqpoint{1.565367in}{1.955406in}}%
\pgfpathlineto{\pgfqpoint{1.587922in}{1.898979in}}%
\pgfpathlineto{\pgfqpoint{1.610476in}{1.955406in}}%
\pgfpathlineto{\pgfqpoint{1.623031in}{1.955406in}}%
\pgfpathlineto{\pgfqpoint{1.591323in}{1.876548in}}%
\pgfpathlineto{\pgfqpoint{1.591323in}{1.876548in}}%
\pgfpathclose%
\pgfpathmoveto{\pgfqpoint{1.683148in}{1.979445in}}%
\pgfpathlineto{\pgfqpoint{1.734236in}{1.979445in}}%
\pgfpathlineto{\pgfqpoint{1.734236in}{1.968477in}}%
\pgfpathlineto{\pgfqpoint{1.695065in}{1.968477in}}%
\pgfpathlineto{\pgfqpoint{1.695065in}{1.944912in}}%
\pgfpathquadraticcurveto{\pgfqpoint{1.697889in}{1.945881in}}{\pgfqpoint{1.700714in}{1.946355in}}%
\pgfpathquadraticcurveto{\pgfqpoint{1.703559in}{1.946830in}}{\pgfqpoint{1.706404in}{1.946830in}}%
\pgfpathquadraticcurveto{\pgfqpoint{1.722505in}{1.946830in}}{\pgfqpoint{1.731906in}{1.938006in}}%
\pgfpathquadraticcurveto{\pgfqpoint{1.741328in}{1.929182in}}{\pgfqpoint{1.741328in}{1.914111in}}%
\pgfpathquadraticcurveto{\pgfqpoint{1.741328in}{1.898587in}}{\pgfqpoint{1.731659in}{1.889970in}}%
\pgfpathquadraticcurveto{\pgfqpoint{1.721990in}{1.881373in}}{\pgfqpoint{1.704404in}{1.881373in}}%
\pgfpathquadraticcurveto{\pgfqpoint{1.698343in}{1.881373in}}{\pgfqpoint{1.692055in}{1.882404in}}%
\pgfpathquadraticcurveto{\pgfqpoint{1.685787in}{1.883434in}}{\pgfqpoint{1.679087in}{1.885496in}}%
\pgfpathlineto{\pgfqpoint{1.679087in}{1.898587in}}%
\pgfpathquadraticcurveto{\pgfqpoint{1.684880in}{1.895433in}}{\pgfqpoint{1.691065in}{1.893887in}}%
\pgfpathquadraticcurveto{\pgfqpoint{1.697250in}{1.892341in}}{\pgfqpoint{1.704136in}{1.892341in}}%
\pgfpathquadraticcurveto{\pgfqpoint{1.715289in}{1.892341in}}{\pgfqpoint{1.721783in}{1.898196in}}%
\pgfpathquadraticcurveto{\pgfqpoint{1.728298in}{1.904051in}}{\pgfqpoint{1.728298in}{1.914111in}}%
\pgfpathquadraticcurveto{\pgfqpoint{1.728298in}{1.924152in}}{\pgfqpoint{1.721783in}{1.930007in}}%
\pgfpathquadraticcurveto{\pgfqpoint{1.715289in}{1.935882in}}{\pgfqpoint{1.704136in}{1.935882in}}%
\pgfpathquadraticcurveto{\pgfqpoint{1.698920in}{1.935882in}}{\pgfqpoint{1.693725in}{1.934728in}}%
\pgfpathquadraticcurveto{\pgfqpoint{1.688550in}{1.933573in}}{\pgfqpoint{1.683148in}{1.931120in}}%
\pgfpathlineto{\pgfqpoint{1.683148in}{1.979445in}}%
\pgfpathlineto{\pgfqpoint{1.683148in}{1.979445in}}%
\pgfpathclose%
\pgfpathmoveto{\pgfqpoint{1.835689in}{1.983362in}}%
\pgfpathquadraticcurveto{\pgfqpoint{1.827071in}{1.968559in}}{\pgfqpoint{1.822865in}{1.954045in}}%
\pgfpathquadraticcurveto{\pgfqpoint{1.818680in}{1.939552in}}{\pgfqpoint{1.818680in}{1.924667in}}%
\pgfpathquadraticcurveto{\pgfqpoint{1.818680in}{1.909803in}}{\pgfqpoint{1.822906in}{1.895206in}}%
\pgfpathquadraticcurveto{\pgfqpoint{1.827133in}{1.880610in}}{\pgfqpoint{1.835689in}{1.865849in}}%
\pgfpathlineto{\pgfqpoint{1.825380in}{1.865849in}}%
\pgfpathquadraticcurveto{\pgfqpoint{1.815732in}{1.881002in}}{\pgfqpoint{1.810928in}{1.895619in}}%
\pgfpathquadraticcurveto{\pgfqpoint{1.806125in}{1.910236in}}{\pgfqpoint{1.806125in}{1.924667in}}%
\pgfpathquadraticcurveto{\pgfqpoint{1.806125in}{1.939037in}}{\pgfqpoint{1.810887in}{1.953592in}}%
\pgfpathquadraticcurveto{\pgfqpoint{1.815670in}{1.968167in}}{\pgfqpoint{1.825380in}{1.983362in}}%
\pgfpathlineto{\pgfqpoint{1.835689in}{1.983362in}}%
\pgfpathlineto{\pgfqpoint{1.835689in}{1.983362in}}%
\pgfpathclose%
\pgfpathmoveto{\pgfqpoint{1.871582in}{1.894196in}}%
\pgfpathlineto{\pgfqpoint{1.916999in}{1.894196in}}%
\pgfpathlineto{\pgfqpoint{1.916999in}{1.883249in}}%
\pgfpathlineto{\pgfqpoint{1.855934in}{1.883249in}}%
\pgfpathlineto{\pgfqpoint{1.855934in}{1.894196in}}%
\pgfpathquadraticcurveto{\pgfqpoint{1.863335in}{1.901865in}}{\pgfqpoint{1.876117in}{1.914771in}}%
\pgfpathquadraticcurveto{\pgfqpoint{1.888920in}{1.927698in}}{\pgfqpoint{1.892198in}{1.931450in}}%
\pgfpathquadraticcurveto{\pgfqpoint{1.898445in}{1.938459in}}{\pgfqpoint{1.900919in}{1.943325in}}%
\pgfpathquadraticcurveto{\pgfqpoint{1.903413in}{1.948190in}}{\pgfqpoint{1.903413in}{1.952891in}}%
\pgfpathquadraticcurveto{\pgfqpoint{1.903413in}{1.960560in}}{\pgfqpoint{1.898032in}{1.965384in}}%
\pgfpathquadraticcurveto{\pgfqpoint{1.892651in}{1.970229in}}{\pgfqpoint{1.884013in}{1.970229in}}%
\pgfpathquadraticcurveto{\pgfqpoint{1.877890in}{1.970229in}}{\pgfqpoint{1.871087in}{1.968106in}}%
\pgfpathquadraticcurveto{\pgfqpoint{1.864304in}{1.965982in}}{\pgfqpoint{1.856573in}{1.961653in}}%
\pgfpathlineto{\pgfqpoint{1.856573in}{1.974806in}}%
\pgfpathquadraticcurveto{\pgfqpoint{1.864428in}{1.977960in}}{\pgfqpoint{1.871252in}{1.979568in}}%
\pgfpathquadraticcurveto{\pgfqpoint{1.878096in}{1.981176in}}{\pgfqpoint{1.883766in}{1.981176in}}%
\pgfpathquadraticcurveto{\pgfqpoint{1.898713in}{1.981176in}}{\pgfqpoint{1.907598in}{1.973693in}}%
\pgfpathquadraticcurveto{\pgfqpoint{1.916484in}{1.966229in}}{\pgfqpoint{1.916484in}{1.953736in}}%
\pgfpathquadraticcurveto{\pgfqpoint{1.916484in}{1.947798in}}{\pgfqpoint{1.914257in}{1.942479in}}%
\pgfpathquadraticcurveto{\pgfqpoint{1.912051in}{1.937181in}}{\pgfqpoint{1.906176in}{1.929965in}}%
\pgfpathquadraticcurveto{\pgfqpoint{1.904568in}{1.928089in}}{\pgfqpoint{1.895929in}{1.919162in}}%
\pgfpathquadraticcurveto{\pgfqpoint{1.887312in}{1.910236in}}{\pgfqpoint{1.871582in}{1.894196in}}%
\pgfpathlineto{\pgfqpoint{1.871582in}{1.894196in}}%
\pgfpathclose%
\pgfpathmoveto{\pgfqpoint{1.972148in}{1.970868in}}%
\pgfpathquadraticcurveto{\pgfqpoint{1.962108in}{1.970868in}}{\pgfqpoint{1.957036in}{1.960972in}}%
\pgfpathquadraticcurveto{\pgfqpoint{1.951985in}{1.951097in}}{\pgfqpoint{1.951985in}{1.931244in}}%
\pgfpathquadraticcurveto{\pgfqpoint{1.951985in}{1.911473in}}{\pgfqpoint{1.957036in}{1.901577in}}%
\pgfpathquadraticcurveto{\pgfqpoint{1.962108in}{1.891681in}}{\pgfqpoint{1.972148in}{1.891681in}}%
\pgfpathquadraticcurveto{\pgfqpoint{1.982271in}{1.891681in}}{\pgfqpoint{1.987322in}{1.901577in}}%
\pgfpathquadraticcurveto{\pgfqpoint{1.992393in}{1.911473in}}{\pgfqpoint{1.992393in}{1.931244in}}%
\pgfpathquadraticcurveto{\pgfqpoint{1.992393in}{1.951097in}}{\pgfqpoint{1.987322in}{1.960972in}}%
\pgfpathquadraticcurveto{\pgfqpoint{1.982271in}{1.970868in}}{\pgfqpoint{1.972148in}{1.970868in}}%
\pgfpathlineto{\pgfqpoint{1.972148in}{1.970868in}}%
\pgfpathclose%
\pgfpathmoveto{\pgfqpoint{1.972148in}{1.981176in}}%
\pgfpathquadraticcurveto{\pgfqpoint{1.988332in}{1.981176in}}{\pgfqpoint{1.996867in}{1.968374in}}%
\pgfpathquadraticcurveto{\pgfqpoint{2.005402in}{1.955591in}}{\pgfqpoint{2.005402in}{1.931244in}}%
\pgfpathquadraticcurveto{\pgfqpoint{2.005402in}{1.906958in}}{\pgfqpoint{1.996867in}{1.894155in}}%
\pgfpathquadraticcurveto{\pgfqpoint{1.988332in}{1.881373in}}{\pgfqpoint{1.972148in}{1.881373in}}%
\pgfpathquadraticcurveto{\pgfqpoint{1.955985in}{1.881373in}}{\pgfqpoint{1.947450in}{1.894155in}}%
\pgfpathquadraticcurveto{\pgfqpoint{1.938914in}{1.906958in}}{\pgfqpoint{1.938914in}{1.931244in}}%
\pgfpathquadraticcurveto{\pgfqpoint{1.938914in}{1.955591in}}{\pgfqpoint{1.947450in}{1.968374in}}%
\pgfpathquadraticcurveto{\pgfqpoint{1.955985in}{1.981176in}}{\pgfqpoint{1.972148in}{1.981176in}}%
\pgfpathlineto{\pgfqpoint{1.972148in}{1.981176in}}%
\pgfpathclose%
\pgfpathmoveto{\pgfqpoint{2.030533in}{1.894196in}}%
\pgfpathlineto{\pgfqpoint{2.051789in}{1.894196in}}%
\pgfpathlineto{\pgfqpoint{2.051789in}{1.967590in}}%
\pgfpathlineto{\pgfqpoint{2.028657in}{1.962951in}}%
\pgfpathlineto{\pgfqpoint{2.028657in}{1.974806in}}%
\pgfpathlineto{\pgfqpoint{2.051665in}{1.979445in}}%
\pgfpathlineto{\pgfqpoint{2.064674in}{1.979445in}}%
\pgfpathlineto{\pgfqpoint{2.064674in}{1.894196in}}%
\pgfpathlineto{\pgfqpoint{2.085929in}{1.894196in}}%
\pgfpathlineto{\pgfqpoint{2.085929in}{1.883249in}}%
\pgfpathlineto{\pgfqpoint{2.030533in}{1.883249in}}%
\pgfpathlineto{\pgfqpoint{2.030533in}{1.894196in}}%
\pgfpathlineto{\pgfqpoint{2.030533in}{1.894196in}}%
\pgfpathclose%
\pgfpathmoveto{\pgfqpoint{2.140047in}{1.928935in}}%
\pgfpathquadraticcurveto{\pgfqpoint{2.130770in}{1.928935in}}{\pgfqpoint{2.125451in}{1.923966in}}%
\pgfpathquadraticcurveto{\pgfqpoint{2.120153in}{1.918997in}}{\pgfqpoint{2.120153in}{1.910318in}}%
\pgfpathquadraticcurveto{\pgfqpoint{2.120153in}{1.901618in}}{\pgfqpoint{2.125451in}{1.896649in}}%
\pgfpathquadraticcurveto{\pgfqpoint{2.130770in}{1.891681in}}{\pgfqpoint{2.140047in}{1.891681in}}%
\pgfpathquadraticcurveto{\pgfqpoint{2.149325in}{1.891681in}}{\pgfqpoint{2.154664in}{1.896670in}}%
\pgfpathquadraticcurveto{\pgfqpoint{2.160025in}{1.901680in}}{\pgfqpoint{2.160025in}{1.910318in}}%
\pgfpathquadraticcurveto{\pgfqpoint{2.160025in}{1.918997in}}{\pgfqpoint{2.154706in}{1.923966in}}%
\pgfpathquadraticcurveto{\pgfqpoint{2.149407in}{1.928935in}}{\pgfqpoint{2.140047in}{1.928935in}}%
\pgfpathlineto{\pgfqpoint{2.140047in}{1.928935in}}%
\pgfpathclose%
\pgfpathmoveto{\pgfqpoint{2.127038in}{1.934460in}}%
\pgfpathquadraticcurveto{\pgfqpoint{2.118668in}{1.936521in}}{\pgfqpoint{2.113988in}{1.942253in}}%
\pgfpathquadraticcurveto{\pgfqpoint{2.109329in}{1.948005in}}{\pgfqpoint{2.109329in}{1.956251in}}%
\pgfpathquadraticcurveto{\pgfqpoint{2.109329in}{1.967776in}}{\pgfqpoint{2.117534in}{1.974476in}}%
\pgfpathquadraticcurveto{\pgfqpoint{2.125760in}{1.981176in}}{\pgfqpoint{2.140047in}{1.981176in}}%
\pgfpathquadraticcurveto{\pgfqpoint{2.154417in}{1.981176in}}{\pgfqpoint{2.162602in}{1.974476in}}%
\pgfpathquadraticcurveto{\pgfqpoint{2.170786in}{1.967776in}}{\pgfqpoint{2.170786in}{1.956251in}}%
\pgfpathquadraticcurveto{\pgfqpoint{2.170786in}{1.948005in}}{\pgfqpoint{2.166106in}{1.942253in}}%
\pgfpathquadraticcurveto{\pgfqpoint{2.161447in}{1.936521in}}{\pgfqpoint{2.153139in}{1.934460in}}%
\pgfpathquadraticcurveto{\pgfqpoint{2.162540in}{1.932274in}}{\pgfqpoint{2.167776in}{1.925883in}}%
\pgfpathquadraticcurveto{\pgfqpoint{2.173033in}{1.919513in}}{\pgfqpoint{2.173033in}{1.910318in}}%
\pgfpathquadraticcurveto{\pgfqpoint{2.173033in}{1.896320in}}{\pgfqpoint{2.164498in}{1.888836in}}%
\pgfpathquadraticcurveto{\pgfqpoint{2.155963in}{1.881373in}}{\pgfqpoint{2.140047in}{1.881373in}}%
\pgfpathquadraticcurveto{\pgfqpoint{2.124152in}{1.881373in}}{\pgfqpoint{2.115596in}{1.888836in}}%
\pgfpathquadraticcurveto{\pgfqpoint{2.107061in}{1.896320in}}{\pgfqpoint{2.107061in}{1.910318in}}%
\pgfpathquadraticcurveto{\pgfqpoint{2.107061in}{1.919513in}}{\pgfqpoint{2.112339in}{1.925883in}}%
\pgfpathquadraticcurveto{\pgfqpoint{2.117637in}{1.932274in}}{\pgfqpoint{2.127038in}{1.934460in}}%
\pgfpathlineto{\pgfqpoint{2.127038in}{1.934460in}}%
\pgfpathclose%
\pgfpathmoveto{\pgfqpoint{2.122276in}{1.955014in}}%
\pgfpathquadraticcurveto{\pgfqpoint{2.122276in}{1.947551in}}{\pgfqpoint{2.126935in}{1.943366in}}%
\pgfpathquadraticcurveto{\pgfqpoint{2.131615in}{1.939181in}}{\pgfqpoint{2.140047in}{1.939181in}}%
\pgfpathquadraticcurveto{\pgfqpoint{2.148438in}{1.939181in}}{\pgfqpoint{2.153159in}{1.943366in}}%
\pgfpathquadraticcurveto{\pgfqpoint{2.157901in}{1.947551in}}{\pgfqpoint{2.157901in}{1.955014in}}%
\pgfpathquadraticcurveto{\pgfqpoint{2.157901in}{1.962498in}}{\pgfqpoint{2.153159in}{1.966683in}}%
\pgfpathquadraticcurveto{\pgfqpoint{2.148438in}{1.970868in}}{\pgfqpoint{2.140047in}{1.970868in}}%
\pgfpathquadraticcurveto{\pgfqpoint{2.131615in}{1.970868in}}{\pgfqpoint{2.126935in}{1.966683in}}%
\pgfpathquadraticcurveto{\pgfqpoint{2.122276in}{1.962498in}}{\pgfqpoint{2.122276in}{1.955014in}}%
\pgfpathlineto{\pgfqpoint{2.122276in}{1.955014in}}%
\pgfpathclose%
\pgfpathmoveto{\pgfqpoint{2.192640in}{1.983362in}}%
\pgfpathlineto{\pgfqpoint{2.202948in}{1.983362in}}%
\pgfpathquadraticcurveto{\pgfqpoint{2.212596in}{1.968167in}}{\pgfqpoint{2.217400in}{1.953592in}}%
\pgfpathquadraticcurveto{\pgfqpoint{2.222203in}{1.939037in}}{\pgfqpoint{2.222203in}{1.924667in}}%
\pgfpathquadraticcurveto{\pgfqpoint{2.222203in}{1.910236in}}{\pgfqpoint{2.217400in}{1.895619in}}%
\pgfpathquadraticcurveto{\pgfqpoint{2.212596in}{1.881002in}}{\pgfqpoint{2.202948in}{1.865849in}}%
\pgfpathlineto{\pgfqpoint{2.192640in}{1.865849in}}%
\pgfpathquadraticcurveto{\pgfqpoint{2.201195in}{1.880610in}}{\pgfqpoint{2.205422in}{1.895206in}}%
\pgfpathquadraticcurveto{\pgfqpoint{2.209648in}{1.909803in}}{\pgfqpoint{2.209648in}{1.924667in}}%
\pgfpathquadraticcurveto{\pgfqpoint{2.209648in}{1.939552in}}{\pgfqpoint{2.205422in}{1.954045in}}%
\pgfpathquadraticcurveto{\pgfqpoint{2.201195in}{1.968559in}}{\pgfqpoint{2.192640in}{1.983362in}}%
\pgfpathlineto{\pgfqpoint{2.192640in}{1.983362in}}%
\pgfpathclose%
\pgfusepath{fill}%
\end{pgfscope}%
\begin{pgfscope}%
\pgfsetbuttcap%
\pgfsetroundjoin%
\definecolor{currentfill}{rgb}{0.090196,0.168627,0.227451}%
\pgfsetfillcolor{currentfill}%
\pgfsetlinewidth{0.000000pt}%
\definecolor{currentstroke}{rgb}{0.090196,0.168627,0.227451}%
\pgfsetstrokecolor{currentstroke}%
\pgfsetdash{}{0pt}%
\pgfpathmoveto{\pgfqpoint{1.316611in}{1.576690in}}%
\pgfpathlineto{\pgfqpoint{1.316611in}{1.563991in}}%
\pgfpathquadraticcurveto{\pgfqpoint{1.309210in}{1.567537in}}{\pgfqpoint{1.302633in}{1.569268in}}%
\pgfpathquadraticcurveto{\pgfqpoint{1.296056in}{1.571021in}}{\pgfqpoint{1.289933in}{1.571021in}}%
\pgfpathquadraticcurveto{\pgfqpoint{1.279316in}{1.571021in}}{\pgfqpoint{1.273543in}{1.566897in}}%
\pgfpathquadraticcurveto{\pgfqpoint{1.267771in}{1.562774in}}{\pgfqpoint{1.267771in}{1.555167in}}%
\pgfpathquadraticcurveto{\pgfqpoint{1.267771in}{1.548776in}}{\pgfqpoint{1.271605in}{1.545518in}}%
\pgfpathquadraticcurveto{\pgfqpoint{1.275440in}{1.542282in}}{\pgfqpoint{1.286140in}{1.540282in}}%
\pgfpathlineto{\pgfqpoint{1.293995in}{1.538674in}}%
\pgfpathquadraticcurveto{\pgfqpoint{1.308550in}{1.535891in}}{\pgfqpoint{1.315477in}{1.528902in}}%
\pgfpathquadraticcurveto{\pgfqpoint{1.322404in}{1.521913in}}{\pgfqpoint{1.322404in}{1.510203in}}%
\pgfpathquadraticcurveto{\pgfqpoint{1.322404in}{1.496204in}}{\pgfqpoint{1.313024in}{1.488988in}}%
\pgfpathquadraticcurveto{\pgfqpoint{1.303664in}{1.481773in}}{\pgfqpoint{1.285563in}{1.481773in}}%
\pgfpathquadraticcurveto{\pgfqpoint{1.278739in}{1.481773in}}{\pgfqpoint{1.271028in}{1.483319in}}%
\pgfpathquadraticcurveto{\pgfqpoint{1.263338in}{1.484865in}}{\pgfqpoint{1.255092in}{1.487896in}}%
\pgfpathlineto{\pgfqpoint{1.255092in}{1.501296in}}%
\pgfpathquadraticcurveto{\pgfqpoint{1.263008in}{1.496864in}}{\pgfqpoint{1.270616in}{1.494596in}}%
\pgfpathquadraticcurveto{\pgfqpoint{1.278223in}{1.492349in}}{\pgfqpoint{1.285563in}{1.492349in}}%
\pgfpathquadraticcurveto{\pgfqpoint{1.296696in}{1.492349in}}{\pgfqpoint{1.302757in}{1.496720in}}%
\pgfpathquadraticcurveto{\pgfqpoint{1.308818in}{1.501111in}}{\pgfqpoint{1.308818in}{1.509234in}}%
\pgfpathquadraticcurveto{\pgfqpoint{1.308818in}{1.516305in}}{\pgfqpoint{1.304468in}{1.520305in}}%
\pgfpathquadraticcurveto{\pgfqpoint{1.300118in}{1.524304in}}{\pgfqpoint{1.290201in}{1.526304in}}%
\pgfpathlineto{\pgfqpoint{1.282264in}{1.527850in}}%
\pgfpathquadraticcurveto{\pgfqpoint{1.267709in}{1.530736in}}{\pgfqpoint{1.261194in}{1.536921in}}%
\pgfpathquadraticcurveto{\pgfqpoint{1.254700in}{1.543106in}}{\pgfqpoint{1.254700in}{1.554136in}}%
\pgfpathquadraticcurveto{\pgfqpoint{1.254700in}{1.566897in}}{\pgfqpoint{1.263689in}{1.574237in}}%
\pgfpathquadraticcurveto{\pgfqpoint{1.272678in}{1.581576in}}{\pgfqpoint{1.288449in}{1.581576in}}%
\pgfpathquadraticcurveto{\pgfqpoint{1.295232in}{1.581576in}}{\pgfqpoint{1.302241in}{1.580339in}}%
\pgfpathquadraticcurveto{\pgfqpoint{1.309271in}{1.579123in}}{\pgfqpoint{1.316611in}{1.576690in}}%
\pgfpathlineto{\pgfqpoint{1.316611in}{1.576690in}}%
\pgfpathclose%
\pgfpathmoveto{\pgfqpoint{1.353926in}{1.576299in}}%
\pgfpathlineto{\pgfqpoint{1.353926in}{1.555806in}}%
\pgfpathlineto{\pgfqpoint{1.378336in}{1.555806in}}%
\pgfpathlineto{\pgfqpoint{1.378336in}{1.546590in}}%
\pgfpathlineto{\pgfqpoint{1.353926in}{1.546590in}}%
\pgfpathlineto{\pgfqpoint{1.353926in}{1.507419in}}%
\pgfpathquadraticcurveto{\pgfqpoint{1.353926in}{1.498596in}}{\pgfqpoint{1.356339in}{1.496080in}}%
\pgfpathquadraticcurveto{\pgfqpoint{1.358751in}{1.493565in}}{\pgfqpoint{1.366173in}{1.493565in}}%
\pgfpathlineto{\pgfqpoint{1.378336in}{1.493565in}}%
\pgfpathlineto{\pgfqpoint{1.378336in}{1.483649in}}%
\pgfpathlineto{\pgfqpoint{1.366173in}{1.483649in}}%
\pgfpathquadraticcurveto{\pgfqpoint{1.352442in}{1.483649in}}{\pgfqpoint{1.347226in}{1.488762in}}%
\pgfpathquadraticcurveto{\pgfqpoint{1.342010in}{1.493895in}}{\pgfqpoint{1.342010in}{1.507419in}}%
\pgfpathlineto{\pgfqpoint{1.342010in}{1.546590in}}%
\pgfpathlineto{\pgfqpoint{1.333310in}{1.546590in}}%
\pgfpathlineto{\pgfqpoint{1.333310in}{1.555806in}}%
\pgfpathlineto{\pgfqpoint{1.342010in}{1.555806in}}%
\pgfpathlineto{\pgfqpoint{1.342010in}{1.576299in}}%
\pgfpathlineto{\pgfqpoint{1.353926in}{1.576299in}}%
\pgfpathlineto{\pgfqpoint{1.353926in}{1.576299in}}%
\pgfpathclose%
\pgfpathmoveto{\pgfqpoint{1.392706in}{1.512120in}}%
\pgfpathlineto{\pgfqpoint{1.392706in}{1.555806in}}%
\pgfpathlineto{\pgfqpoint{1.404560in}{1.555806in}}%
\pgfpathlineto{\pgfqpoint{1.404560in}{1.512573in}}%
\pgfpathquadraticcurveto{\pgfqpoint{1.404560in}{1.502327in}}{\pgfqpoint{1.408539in}{1.497194in}}%
\pgfpathquadraticcurveto{\pgfqpoint{1.412539in}{1.492081in}}{\pgfqpoint{1.420538in}{1.492081in}}%
\pgfpathquadraticcurveto{\pgfqpoint{1.430124in}{1.492081in}}{\pgfqpoint{1.435691in}{1.498204in}}%
\pgfpathquadraticcurveto{\pgfqpoint{1.441278in}{1.504327in}}{\pgfqpoint{1.441278in}{1.514903in}}%
\pgfpathlineto{\pgfqpoint{1.441278in}{1.555806in}}%
\pgfpathlineto{\pgfqpoint{1.453132in}{1.555806in}}%
\pgfpathlineto{\pgfqpoint{1.453132in}{1.483649in}}%
\pgfpathlineto{\pgfqpoint{1.441278in}{1.483649in}}%
\pgfpathlineto{\pgfqpoint{1.441278in}{1.494740in}}%
\pgfpathquadraticcurveto{\pgfqpoint{1.436969in}{1.488164in}}{\pgfqpoint{1.431258in}{1.484968in}}%
\pgfpathquadraticcurveto{\pgfqpoint{1.425568in}{1.481773in}}{\pgfqpoint{1.418023in}{1.481773in}}%
\pgfpathquadraticcurveto{\pgfqpoint{1.405591in}{1.481773in}}{\pgfqpoint{1.399138in}{1.489504in}}%
\pgfpathquadraticcurveto{\pgfqpoint{1.392706in}{1.497235in}}{\pgfqpoint{1.392706in}{1.512120in}}%
\pgfpathlineto{\pgfqpoint{1.392706in}{1.512120in}}%
\pgfpathclose%
\pgfpathmoveto{\pgfqpoint{1.422538in}{1.557538in}}%
\pgfpathlineto{\pgfqpoint{1.422538in}{1.557538in}}%
\pgfpathlineto{\pgfqpoint{1.422538in}{1.557538in}}%
\pgfpathclose%
\pgfpathmoveto{\pgfqpoint{1.525021in}{1.544859in}}%
\pgfpathlineto{\pgfqpoint{1.525021in}{1.583906in}}%
\pgfpathlineto{\pgfqpoint{1.536876in}{1.583906in}}%
\pgfpathlineto{\pgfqpoint{1.536876in}{1.483649in}}%
\pgfpathlineto{\pgfqpoint{1.525021in}{1.483649in}}%
\pgfpathlineto{\pgfqpoint{1.525021in}{1.494472in}}%
\pgfpathquadraticcurveto{\pgfqpoint{1.521290in}{1.488040in}}{\pgfqpoint{1.515579in}{1.484906in}}%
\pgfpathquadraticcurveto{\pgfqpoint{1.509889in}{1.481773in}}{\pgfqpoint{1.501890in}{1.481773in}}%
\pgfpathquadraticcurveto{\pgfqpoint{1.488819in}{1.481773in}}{\pgfqpoint{1.480593in}{1.492205in}}%
\pgfpathquadraticcurveto{\pgfqpoint{1.472388in}{1.502657in}}{\pgfqpoint{1.472388in}{1.519666in}}%
\pgfpathquadraticcurveto{\pgfqpoint{1.472388in}{1.536674in}}{\pgfqpoint{1.480593in}{1.547106in}}%
\pgfpathquadraticcurveto{\pgfqpoint{1.488819in}{1.557538in}}{\pgfqpoint{1.501890in}{1.557538in}}%
\pgfpathquadraticcurveto{\pgfqpoint{1.509889in}{1.557538in}}{\pgfqpoint{1.515579in}{1.554404in}}%
\pgfpathquadraticcurveto{\pgfqpoint{1.521290in}{1.551291in}}{\pgfqpoint{1.525021in}{1.544859in}}%
\pgfpathlineto{\pgfqpoint{1.525021in}{1.544859in}}%
\pgfpathclose%
\pgfpathmoveto{\pgfqpoint{1.484634in}{1.519666in}}%
\pgfpathquadraticcurveto{\pgfqpoint{1.484634in}{1.506595in}}{\pgfqpoint{1.490015in}{1.499152in}}%
\pgfpathquadraticcurveto{\pgfqpoint{1.495396in}{1.491710in}}{\pgfqpoint{1.504797in}{1.491710in}}%
\pgfpathquadraticcurveto{\pgfqpoint{1.514198in}{1.491710in}}{\pgfqpoint{1.519599in}{1.499152in}}%
\pgfpathquadraticcurveto{\pgfqpoint{1.525021in}{1.506595in}}{\pgfqpoint{1.525021in}{1.519666in}}%
\pgfpathquadraticcurveto{\pgfqpoint{1.525021in}{1.532736in}}{\pgfqpoint{1.519599in}{1.540179in}}%
\pgfpathquadraticcurveto{\pgfqpoint{1.514198in}{1.547621in}}{\pgfqpoint{1.504797in}{1.547621in}}%
\pgfpathquadraticcurveto{\pgfqpoint{1.495396in}{1.547621in}}{\pgfqpoint{1.490015in}{1.540179in}}%
\pgfpathquadraticcurveto{\pgfqpoint{1.484634in}{1.532736in}}{\pgfqpoint{1.484634in}{1.519666in}}%
\pgfpathlineto{\pgfqpoint{1.484634in}{1.519666in}}%
\pgfpathclose%
\pgfpathmoveto{\pgfqpoint{1.591323in}{1.476948in}}%
\pgfpathquadraticcurveto{\pgfqpoint{1.586314in}{1.464063in}}{\pgfqpoint{1.581531in}{1.460146in}}%
\pgfpathquadraticcurveto{\pgfqpoint{1.576768in}{1.456208in}}{\pgfqpoint{1.568790in}{1.456208in}}%
\pgfpathlineto{\pgfqpoint{1.559306in}{1.456208in}}%
\pgfpathlineto{\pgfqpoint{1.559306in}{1.466125in}}%
\pgfpathlineto{\pgfqpoint{1.566275in}{1.466125in}}%
\pgfpathquadraticcurveto{\pgfqpoint{1.571161in}{1.466125in}}{\pgfqpoint{1.573861in}{1.468455in}}%
\pgfpathquadraticcurveto{\pgfqpoint{1.576583in}{1.470764in}}{\pgfqpoint{1.579861in}{1.479402in}}%
\pgfpathlineto{\pgfqpoint{1.581984in}{1.484803in}}%
\pgfpathlineto{\pgfqpoint{1.552812in}{1.555806in}}%
\pgfpathlineto{\pgfqpoint{1.565367in}{1.555806in}}%
\pgfpathlineto{\pgfqpoint{1.587922in}{1.499379in}}%
\pgfpathlineto{\pgfqpoint{1.610476in}{1.555806in}}%
\pgfpathlineto{\pgfqpoint{1.623031in}{1.555806in}}%
\pgfpathlineto{\pgfqpoint{1.591323in}{1.476948in}}%
\pgfpathlineto{\pgfqpoint{1.591323in}{1.476948in}}%
\pgfpathclose%
\pgfpathmoveto{\pgfqpoint{1.712465in}{1.536921in}}%
\pgfpathquadraticcurveto{\pgfqpoint{1.703703in}{1.536921in}}{\pgfqpoint{1.698569in}{1.530922in}}%
\pgfpathquadraticcurveto{\pgfqpoint{1.693457in}{1.524943in}}{\pgfqpoint{1.693457in}{1.514511in}}%
\pgfpathquadraticcurveto{\pgfqpoint{1.693457in}{1.504141in}}{\pgfqpoint{1.698569in}{1.498101in}}%
\pgfpathquadraticcurveto{\pgfqpoint{1.703703in}{1.492081in}}{\pgfqpoint{1.712465in}{1.492081in}}%
\pgfpathquadraticcurveto{\pgfqpoint{1.721227in}{1.492081in}}{\pgfqpoint{1.726340in}{1.498101in}}%
\pgfpathquadraticcurveto{\pgfqpoint{1.731452in}{1.504141in}}{\pgfqpoint{1.731452in}{1.514511in}}%
\pgfpathquadraticcurveto{\pgfqpoint{1.731452in}{1.524943in}}{\pgfqpoint{1.726340in}{1.530922in}}%
\pgfpathquadraticcurveto{\pgfqpoint{1.721227in}{1.536921in}}{\pgfqpoint{1.712465in}{1.536921in}}%
\pgfpathlineto{\pgfqpoint{1.712465in}{1.536921in}}%
\pgfpathclose%
\pgfpathmoveto{\pgfqpoint{1.738297in}{1.577721in}}%
\pgfpathlineto{\pgfqpoint{1.738297in}{1.565867in}}%
\pgfpathquadraticcurveto{\pgfqpoint{1.733390in}{1.568176in}}{\pgfqpoint{1.728401in}{1.569392in}}%
\pgfpathquadraticcurveto{\pgfqpoint{1.723412in}{1.570629in}}{\pgfqpoint{1.718505in}{1.570629in}}%
\pgfpathquadraticcurveto{\pgfqpoint{1.705620in}{1.570629in}}{\pgfqpoint{1.698817in}{1.561929in}}%
\pgfpathquadraticcurveto{\pgfqpoint{1.692034in}{1.553229in}}{\pgfqpoint{1.691065in}{1.535643in}}%
\pgfpathquadraticcurveto{\pgfqpoint{1.694859in}{1.541251in}}{\pgfqpoint{1.700590in}{1.544240in}}%
\pgfpathquadraticcurveto{\pgfqpoint{1.706342in}{1.547230in}}{\pgfqpoint{1.713228in}{1.547230in}}%
\pgfpathquadraticcurveto{\pgfqpoint{1.727721in}{1.547230in}}{\pgfqpoint{1.736132in}{1.538426in}}%
\pgfpathquadraticcurveto{\pgfqpoint{1.744544in}{1.529644in}}{\pgfqpoint{1.744544in}{1.514511in}}%
\pgfpathquadraticcurveto{\pgfqpoint{1.744544in}{1.499688in}}{\pgfqpoint{1.735782in}{1.490720in}}%
\pgfpathquadraticcurveto{\pgfqpoint{1.727020in}{1.481773in}}{\pgfqpoint{1.712465in}{1.481773in}}%
\pgfpathquadraticcurveto{\pgfqpoint{1.695766in}{1.481773in}}{\pgfqpoint{1.686942in}{1.494555in}}%
\pgfpathquadraticcurveto{\pgfqpoint{1.678118in}{1.507358in}}{\pgfqpoint{1.678118in}{1.531644in}}%
\pgfpathquadraticcurveto{\pgfqpoint{1.678118in}{1.554445in}}{\pgfqpoint{1.688942in}{1.568011in}}%
\pgfpathquadraticcurveto{\pgfqpoint{1.699765in}{1.581576in}}{\pgfqpoint{1.717990in}{1.581576in}}%
\pgfpathquadraticcurveto{\pgfqpoint{1.722897in}{1.581576in}}{\pgfqpoint{1.727886in}{1.580607in}}%
\pgfpathquadraticcurveto{\pgfqpoint{1.732875in}{1.579638in}}{\pgfqpoint{1.738297in}{1.577721in}}%
\pgfpathlineto{\pgfqpoint{1.738297in}{1.577721in}}%
\pgfpathclose%
\pgfpathmoveto{\pgfqpoint{1.835689in}{1.583762in}}%
\pgfpathquadraticcurveto{\pgfqpoint{1.827071in}{1.568959in}}{\pgfqpoint{1.822865in}{1.554445in}}%
\pgfpathquadraticcurveto{\pgfqpoint{1.818680in}{1.539952in}}{\pgfqpoint{1.818680in}{1.525067in}}%
\pgfpathquadraticcurveto{\pgfqpoint{1.818680in}{1.510203in}}{\pgfqpoint{1.822906in}{1.495606in}}%
\pgfpathquadraticcurveto{\pgfqpoint{1.827133in}{1.481010in}}{\pgfqpoint{1.835689in}{1.466249in}}%
\pgfpathlineto{\pgfqpoint{1.825380in}{1.466249in}}%
\pgfpathquadraticcurveto{\pgfqpoint{1.815732in}{1.481402in}}{\pgfqpoint{1.810928in}{1.496019in}}%
\pgfpathquadraticcurveto{\pgfqpoint{1.806125in}{1.510636in}}{\pgfqpoint{1.806125in}{1.525067in}}%
\pgfpathquadraticcurveto{\pgfqpoint{1.806125in}{1.539437in}}{\pgfqpoint{1.810887in}{1.553992in}}%
\pgfpathquadraticcurveto{\pgfqpoint{1.815670in}{1.568567in}}{\pgfqpoint{1.825380in}{1.583762in}}%
\pgfpathlineto{\pgfqpoint{1.835689in}{1.583762in}}%
\pgfpathlineto{\pgfqpoint{1.835689in}{1.583762in}}%
\pgfpathclose%
\pgfpathmoveto{\pgfqpoint{1.871582in}{1.494596in}}%
\pgfpathlineto{\pgfqpoint{1.916999in}{1.494596in}}%
\pgfpathlineto{\pgfqpoint{1.916999in}{1.483649in}}%
\pgfpathlineto{\pgfqpoint{1.855934in}{1.483649in}}%
\pgfpathlineto{\pgfqpoint{1.855934in}{1.494596in}}%
\pgfpathquadraticcurveto{\pgfqpoint{1.863335in}{1.502265in}}{\pgfqpoint{1.876117in}{1.515171in}}%
\pgfpathquadraticcurveto{\pgfqpoint{1.888920in}{1.528098in}}{\pgfqpoint{1.892198in}{1.531850in}}%
\pgfpathquadraticcurveto{\pgfqpoint{1.898445in}{1.538859in}}{\pgfqpoint{1.900919in}{1.543725in}}%
\pgfpathquadraticcurveto{\pgfqpoint{1.903413in}{1.548590in}}{\pgfqpoint{1.903413in}{1.553291in}}%
\pgfpathquadraticcurveto{\pgfqpoint{1.903413in}{1.560960in}}{\pgfqpoint{1.898032in}{1.565784in}}%
\pgfpathquadraticcurveto{\pgfqpoint{1.892651in}{1.570629in}}{\pgfqpoint{1.884013in}{1.570629in}}%
\pgfpathquadraticcurveto{\pgfqpoint{1.877890in}{1.570629in}}{\pgfqpoint{1.871087in}{1.568506in}}%
\pgfpathquadraticcurveto{\pgfqpoint{1.864304in}{1.566382in}}{\pgfqpoint{1.856573in}{1.562053in}}%
\pgfpathlineto{\pgfqpoint{1.856573in}{1.575206in}}%
\pgfpathquadraticcurveto{\pgfqpoint{1.864428in}{1.578360in}}{\pgfqpoint{1.871252in}{1.579968in}}%
\pgfpathquadraticcurveto{\pgfqpoint{1.878096in}{1.581576in}}{\pgfqpoint{1.883766in}{1.581576in}}%
\pgfpathquadraticcurveto{\pgfqpoint{1.898713in}{1.581576in}}{\pgfqpoint{1.907598in}{1.574093in}}%
\pgfpathquadraticcurveto{\pgfqpoint{1.916484in}{1.566629in}}{\pgfqpoint{1.916484in}{1.554136in}}%
\pgfpathquadraticcurveto{\pgfqpoint{1.916484in}{1.548198in}}{\pgfqpoint{1.914257in}{1.542879in}}%
\pgfpathquadraticcurveto{\pgfqpoint{1.912051in}{1.537581in}}{\pgfqpoint{1.906176in}{1.530365in}}%
\pgfpathquadraticcurveto{\pgfqpoint{1.904568in}{1.528489in}}{\pgfqpoint{1.895929in}{1.519562in}}%
\pgfpathquadraticcurveto{\pgfqpoint{1.887312in}{1.510636in}}{\pgfqpoint{1.871582in}{1.494596in}}%
\pgfpathlineto{\pgfqpoint{1.871582in}{1.494596in}}%
\pgfpathclose%
\pgfpathmoveto{\pgfqpoint{1.972148in}{1.571268in}}%
\pgfpathquadraticcurveto{\pgfqpoint{1.962108in}{1.571268in}}{\pgfqpoint{1.957036in}{1.561372in}}%
\pgfpathquadraticcurveto{\pgfqpoint{1.951985in}{1.551497in}}{\pgfqpoint{1.951985in}{1.531644in}}%
\pgfpathquadraticcurveto{\pgfqpoint{1.951985in}{1.511873in}}{\pgfqpoint{1.957036in}{1.501977in}}%
\pgfpathquadraticcurveto{\pgfqpoint{1.962108in}{1.492081in}}{\pgfqpoint{1.972148in}{1.492081in}}%
\pgfpathquadraticcurveto{\pgfqpoint{1.982271in}{1.492081in}}{\pgfqpoint{1.987322in}{1.501977in}}%
\pgfpathquadraticcurveto{\pgfqpoint{1.992393in}{1.511873in}}{\pgfqpoint{1.992393in}{1.531644in}}%
\pgfpathquadraticcurveto{\pgfqpoint{1.992393in}{1.551497in}}{\pgfqpoint{1.987322in}{1.561372in}}%
\pgfpathquadraticcurveto{\pgfqpoint{1.982271in}{1.571268in}}{\pgfqpoint{1.972148in}{1.571268in}}%
\pgfpathlineto{\pgfqpoint{1.972148in}{1.571268in}}%
\pgfpathclose%
\pgfpathmoveto{\pgfqpoint{1.972148in}{1.581576in}}%
\pgfpathquadraticcurveto{\pgfqpoint{1.988332in}{1.581576in}}{\pgfqpoint{1.996867in}{1.568774in}}%
\pgfpathquadraticcurveto{\pgfqpoint{2.005402in}{1.555991in}}{\pgfqpoint{2.005402in}{1.531644in}}%
\pgfpathquadraticcurveto{\pgfqpoint{2.005402in}{1.507358in}}{\pgfqpoint{1.996867in}{1.494555in}}%
\pgfpathquadraticcurveto{\pgfqpoint{1.988332in}{1.481773in}}{\pgfqpoint{1.972148in}{1.481773in}}%
\pgfpathquadraticcurveto{\pgfqpoint{1.955985in}{1.481773in}}{\pgfqpoint{1.947450in}{1.494555in}}%
\pgfpathquadraticcurveto{\pgfqpoint{1.938914in}{1.507358in}}{\pgfqpoint{1.938914in}{1.531644in}}%
\pgfpathquadraticcurveto{\pgfqpoint{1.938914in}{1.555991in}}{\pgfqpoint{1.947450in}{1.568774in}}%
\pgfpathquadraticcurveto{\pgfqpoint{1.955985in}{1.581576in}}{\pgfqpoint{1.972148in}{1.581576in}}%
\pgfpathlineto{\pgfqpoint{1.972148in}{1.581576in}}%
\pgfpathclose%
\pgfpathmoveto{\pgfqpoint{2.030533in}{1.494596in}}%
\pgfpathlineto{\pgfqpoint{2.051789in}{1.494596in}}%
\pgfpathlineto{\pgfqpoint{2.051789in}{1.567990in}}%
\pgfpathlineto{\pgfqpoint{2.028657in}{1.563351in}}%
\pgfpathlineto{\pgfqpoint{2.028657in}{1.575206in}}%
\pgfpathlineto{\pgfqpoint{2.051665in}{1.579845in}}%
\pgfpathlineto{\pgfqpoint{2.064674in}{1.579845in}}%
\pgfpathlineto{\pgfqpoint{2.064674in}{1.494596in}}%
\pgfpathlineto{\pgfqpoint{2.085929in}{1.494596in}}%
\pgfpathlineto{\pgfqpoint{2.085929in}{1.483649in}}%
\pgfpathlineto{\pgfqpoint{2.030533in}{1.483649in}}%
\pgfpathlineto{\pgfqpoint{2.030533in}{1.494596in}}%
\pgfpathlineto{\pgfqpoint{2.030533in}{1.494596in}}%
\pgfpathclose%
\pgfpathmoveto{\pgfqpoint{2.140047in}{1.529335in}}%
\pgfpathquadraticcurveto{\pgfqpoint{2.130770in}{1.529335in}}{\pgfqpoint{2.125451in}{1.524366in}}%
\pgfpathquadraticcurveto{\pgfqpoint{2.120153in}{1.519397in}}{\pgfqpoint{2.120153in}{1.510718in}}%
\pgfpathquadraticcurveto{\pgfqpoint{2.120153in}{1.502018in}}{\pgfqpoint{2.125451in}{1.497049in}}%
\pgfpathquadraticcurveto{\pgfqpoint{2.130770in}{1.492081in}}{\pgfqpoint{2.140047in}{1.492081in}}%
\pgfpathquadraticcurveto{\pgfqpoint{2.149325in}{1.492081in}}{\pgfqpoint{2.154664in}{1.497070in}}%
\pgfpathquadraticcurveto{\pgfqpoint{2.160025in}{1.502080in}}{\pgfqpoint{2.160025in}{1.510718in}}%
\pgfpathquadraticcurveto{\pgfqpoint{2.160025in}{1.519397in}}{\pgfqpoint{2.154706in}{1.524366in}}%
\pgfpathquadraticcurveto{\pgfqpoint{2.149407in}{1.529335in}}{\pgfqpoint{2.140047in}{1.529335in}}%
\pgfpathlineto{\pgfqpoint{2.140047in}{1.529335in}}%
\pgfpathclose%
\pgfpathmoveto{\pgfqpoint{2.127038in}{1.534860in}}%
\pgfpathquadraticcurveto{\pgfqpoint{2.118668in}{1.536921in}}{\pgfqpoint{2.113988in}{1.542653in}}%
\pgfpathquadraticcurveto{\pgfqpoint{2.109329in}{1.548405in}}{\pgfqpoint{2.109329in}{1.556651in}}%
\pgfpathquadraticcurveto{\pgfqpoint{2.109329in}{1.568176in}}{\pgfqpoint{2.117534in}{1.574876in}}%
\pgfpathquadraticcurveto{\pgfqpoint{2.125760in}{1.581576in}}{\pgfqpoint{2.140047in}{1.581576in}}%
\pgfpathquadraticcurveto{\pgfqpoint{2.154417in}{1.581576in}}{\pgfqpoint{2.162602in}{1.574876in}}%
\pgfpathquadraticcurveto{\pgfqpoint{2.170786in}{1.568176in}}{\pgfqpoint{2.170786in}{1.556651in}}%
\pgfpathquadraticcurveto{\pgfqpoint{2.170786in}{1.548405in}}{\pgfqpoint{2.166106in}{1.542653in}}%
\pgfpathquadraticcurveto{\pgfqpoint{2.161447in}{1.536921in}}{\pgfqpoint{2.153139in}{1.534860in}}%
\pgfpathquadraticcurveto{\pgfqpoint{2.162540in}{1.532674in}}{\pgfqpoint{2.167776in}{1.526283in}}%
\pgfpathquadraticcurveto{\pgfqpoint{2.173033in}{1.519913in}}{\pgfqpoint{2.173033in}{1.510718in}}%
\pgfpathquadraticcurveto{\pgfqpoint{2.173033in}{1.496720in}}{\pgfqpoint{2.164498in}{1.489236in}}%
\pgfpathquadraticcurveto{\pgfqpoint{2.155963in}{1.481773in}}{\pgfqpoint{2.140047in}{1.481773in}}%
\pgfpathquadraticcurveto{\pgfqpoint{2.124152in}{1.481773in}}{\pgfqpoint{2.115596in}{1.489236in}}%
\pgfpathquadraticcurveto{\pgfqpoint{2.107061in}{1.496720in}}{\pgfqpoint{2.107061in}{1.510718in}}%
\pgfpathquadraticcurveto{\pgfqpoint{2.107061in}{1.519913in}}{\pgfqpoint{2.112339in}{1.526283in}}%
\pgfpathquadraticcurveto{\pgfqpoint{2.117637in}{1.532674in}}{\pgfqpoint{2.127038in}{1.534860in}}%
\pgfpathlineto{\pgfqpoint{2.127038in}{1.534860in}}%
\pgfpathclose%
\pgfpathmoveto{\pgfqpoint{2.122276in}{1.555414in}}%
\pgfpathquadraticcurveto{\pgfqpoint{2.122276in}{1.547951in}}{\pgfqpoint{2.126935in}{1.543766in}}%
\pgfpathquadraticcurveto{\pgfqpoint{2.131615in}{1.539581in}}{\pgfqpoint{2.140047in}{1.539581in}}%
\pgfpathquadraticcurveto{\pgfqpoint{2.148438in}{1.539581in}}{\pgfqpoint{2.153159in}{1.543766in}}%
\pgfpathquadraticcurveto{\pgfqpoint{2.157901in}{1.547951in}}{\pgfqpoint{2.157901in}{1.555414in}}%
\pgfpathquadraticcurveto{\pgfqpoint{2.157901in}{1.562898in}}{\pgfqpoint{2.153159in}{1.567083in}}%
\pgfpathquadraticcurveto{\pgfqpoint{2.148438in}{1.571268in}}{\pgfqpoint{2.140047in}{1.571268in}}%
\pgfpathquadraticcurveto{\pgfqpoint{2.131615in}{1.571268in}}{\pgfqpoint{2.126935in}{1.567083in}}%
\pgfpathquadraticcurveto{\pgfqpoint{2.122276in}{1.562898in}}{\pgfqpoint{2.122276in}{1.555414in}}%
\pgfpathlineto{\pgfqpoint{2.122276in}{1.555414in}}%
\pgfpathclose%
\pgfpathmoveto{\pgfqpoint{2.192640in}{1.583762in}}%
\pgfpathlineto{\pgfqpoint{2.202948in}{1.583762in}}%
\pgfpathquadraticcurveto{\pgfqpoint{2.212596in}{1.568567in}}{\pgfqpoint{2.217400in}{1.553992in}}%
\pgfpathquadraticcurveto{\pgfqpoint{2.222203in}{1.539437in}}{\pgfqpoint{2.222203in}{1.525067in}}%
\pgfpathquadraticcurveto{\pgfqpoint{2.222203in}{1.510636in}}{\pgfqpoint{2.217400in}{1.496019in}}%
\pgfpathquadraticcurveto{\pgfqpoint{2.212596in}{1.481402in}}{\pgfqpoint{2.202948in}{1.466249in}}%
\pgfpathlineto{\pgfqpoint{2.192640in}{1.466249in}}%
\pgfpathquadraticcurveto{\pgfqpoint{2.201195in}{1.481010in}}{\pgfqpoint{2.205422in}{1.495606in}}%
\pgfpathquadraticcurveto{\pgfqpoint{2.209648in}{1.510203in}}{\pgfqpoint{2.209648in}{1.525067in}}%
\pgfpathquadraticcurveto{\pgfqpoint{2.209648in}{1.539952in}}{\pgfqpoint{2.205422in}{1.554445in}}%
\pgfpathquadraticcurveto{\pgfqpoint{2.201195in}{1.568959in}}{\pgfqpoint{2.192640in}{1.583762in}}%
\pgfpathlineto{\pgfqpoint{2.192640in}{1.583762in}}%
\pgfpathclose%
\pgfusepath{fill}%
\end{pgfscope}%
\begin{pgfscope}%
\pgfsetbuttcap%
\pgfsetroundjoin%
\definecolor{currentfill}{rgb}{0.090196,0.168627,0.227451}%
\pgfsetfillcolor{currentfill}%
\pgfsetlinewidth{0.000000pt}%
\definecolor{currentstroke}{rgb}{0.090196,0.168627,0.227451}%
\pgfsetstrokecolor{currentstroke}%
\pgfsetdash{}{0pt}%
\pgfpathmoveto{\pgfqpoint{4.550611in}{3.175090in}}%
\pgfpathlineto{\pgfqpoint{4.550611in}{3.162391in}}%
\pgfpathquadraticcurveto{\pgfqpoint{4.543210in}{3.165937in}}{\pgfqpoint{4.536633in}{3.167668in}}%
\pgfpathquadraticcurveto{\pgfqpoint{4.530056in}{3.169421in}}{\pgfqpoint{4.523933in}{3.169421in}}%
\pgfpathquadraticcurveto{\pgfqpoint{4.513316in}{3.169421in}}{\pgfqpoint{4.507543in}{3.165297in}}%
\pgfpathquadraticcurveto{\pgfqpoint{4.501771in}{3.161174in}}{\pgfqpoint{4.501771in}{3.153567in}}%
\pgfpathquadraticcurveto{\pgfqpoint{4.501771in}{3.147176in}}{\pgfqpoint{4.505605in}{3.143918in}}%
\pgfpathquadraticcurveto{\pgfqpoint{4.509440in}{3.140682in}}{\pgfqpoint{4.520140in}{3.138682in}}%
\pgfpathlineto{\pgfqpoint{4.527995in}{3.137074in}}%
\pgfpathquadraticcurveto{\pgfqpoint{4.542550in}{3.134291in}}{\pgfqpoint{4.549477in}{3.127302in}}%
\pgfpathquadraticcurveto{\pgfqpoint{4.556404in}{3.120313in}}{\pgfqpoint{4.556404in}{3.108603in}}%
\pgfpathquadraticcurveto{\pgfqpoint{4.556404in}{3.094604in}}{\pgfqpoint{4.547024in}{3.087388in}}%
\pgfpathquadraticcurveto{\pgfqpoint{4.537664in}{3.080173in}}{\pgfqpoint{4.519563in}{3.080173in}}%
\pgfpathquadraticcurveto{\pgfqpoint{4.512739in}{3.080173in}}{\pgfqpoint{4.505028in}{3.081719in}}%
\pgfpathquadraticcurveto{\pgfqpoint{4.497338in}{3.083265in}}{\pgfqpoint{4.489092in}{3.086296in}}%
\pgfpathlineto{\pgfqpoint{4.489092in}{3.099696in}}%
\pgfpathquadraticcurveto{\pgfqpoint{4.497008in}{3.095264in}}{\pgfqpoint{4.504616in}{3.092996in}}%
\pgfpathquadraticcurveto{\pgfqpoint{4.512223in}{3.090749in}}{\pgfqpoint{4.519563in}{3.090749in}}%
\pgfpathquadraticcurveto{\pgfqpoint{4.530696in}{3.090749in}}{\pgfqpoint{4.536757in}{3.095120in}}%
\pgfpathquadraticcurveto{\pgfqpoint{4.542818in}{3.099511in}}{\pgfqpoint{4.542818in}{3.107634in}}%
\pgfpathquadraticcurveto{\pgfqpoint{4.542818in}{3.114705in}}{\pgfqpoint{4.538468in}{3.118705in}}%
\pgfpathquadraticcurveto{\pgfqpoint{4.534118in}{3.122704in}}{\pgfqpoint{4.524201in}{3.124704in}}%
\pgfpathlineto{\pgfqpoint{4.516264in}{3.126250in}}%
\pgfpathquadraticcurveto{\pgfqpoint{4.501709in}{3.129136in}}{\pgfqpoint{4.495194in}{3.135321in}}%
\pgfpathquadraticcurveto{\pgfqpoint{4.488700in}{3.141506in}}{\pgfqpoint{4.488700in}{3.152536in}}%
\pgfpathquadraticcurveto{\pgfqpoint{4.488700in}{3.165297in}}{\pgfqpoint{4.497689in}{3.172637in}}%
\pgfpathquadraticcurveto{\pgfqpoint{4.506678in}{3.179976in}}{\pgfqpoint{4.522449in}{3.179976in}}%
\pgfpathquadraticcurveto{\pgfqpoint{4.529232in}{3.179976in}}{\pgfqpoint{4.536241in}{3.178739in}}%
\pgfpathquadraticcurveto{\pgfqpoint{4.543271in}{3.177523in}}{\pgfqpoint{4.550611in}{3.175090in}}%
\pgfpathlineto{\pgfqpoint{4.550611in}{3.175090in}}%
\pgfpathclose%
\pgfpathmoveto{\pgfqpoint{4.587926in}{3.174699in}}%
\pgfpathlineto{\pgfqpoint{4.587926in}{3.154206in}}%
\pgfpathlineto{\pgfqpoint{4.612336in}{3.154206in}}%
\pgfpathlineto{\pgfqpoint{4.612336in}{3.144990in}}%
\pgfpathlineto{\pgfqpoint{4.587926in}{3.144990in}}%
\pgfpathlineto{\pgfqpoint{4.587926in}{3.105819in}}%
\pgfpathquadraticcurveto{\pgfqpoint{4.587926in}{3.096996in}}{\pgfqpoint{4.590339in}{3.094480in}}%
\pgfpathquadraticcurveto{\pgfqpoint{4.592751in}{3.091965in}}{\pgfqpoint{4.600173in}{3.091965in}}%
\pgfpathlineto{\pgfqpoint{4.612336in}{3.091965in}}%
\pgfpathlineto{\pgfqpoint{4.612336in}{3.082049in}}%
\pgfpathlineto{\pgfqpoint{4.600173in}{3.082049in}}%
\pgfpathquadraticcurveto{\pgfqpoint{4.586442in}{3.082049in}}{\pgfqpoint{4.581226in}{3.087162in}}%
\pgfpathquadraticcurveto{\pgfqpoint{4.576010in}{3.092295in}}{\pgfqpoint{4.576010in}{3.105819in}}%
\pgfpathlineto{\pgfqpoint{4.576010in}{3.144990in}}%
\pgfpathlineto{\pgfqpoint{4.567310in}{3.144990in}}%
\pgfpathlineto{\pgfqpoint{4.567310in}{3.154206in}}%
\pgfpathlineto{\pgfqpoint{4.576010in}{3.154206in}}%
\pgfpathlineto{\pgfqpoint{4.576010in}{3.174699in}}%
\pgfpathlineto{\pgfqpoint{4.587926in}{3.174699in}}%
\pgfpathlineto{\pgfqpoint{4.587926in}{3.174699in}}%
\pgfpathclose%
\pgfpathmoveto{\pgfqpoint{4.626706in}{3.110520in}}%
\pgfpathlineto{\pgfqpoint{4.626706in}{3.154206in}}%
\pgfpathlineto{\pgfqpoint{4.638560in}{3.154206in}}%
\pgfpathlineto{\pgfqpoint{4.638560in}{3.110973in}}%
\pgfpathquadraticcurveto{\pgfqpoint{4.638560in}{3.100727in}}{\pgfqpoint{4.642539in}{3.095594in}}%
\pgfpathquadraticcurveto{\pgfqpoint{4.646539in}{3.090481in}}{\pgfqpoint{4.654538in}{3.090481in}}%
\pgfpathquadraticcurveto{\pgfqpoint{4.664124in}{3.090481in}}{\pgfqpoint{4.669691in}{3.096604in}}%
\pgfpathquadraticcurveto{\pgfqpoint{4.675278in}{3.102727in}}{\pgfqpoint{4.675278in}{3.113303in}}%
\pgfpathlineto{\pgfqpoint{4.675278in}{3.154206in}}%
\pgfpathlineto{\pgfqpoint{4.687132in}{3.154206in}}%
\pgfpathlineto{\pgfqpoint{4.687132in}{3.082049in}}%
\pgfpathlineto{\pgfqpoint{4.675278in}{3.082049in}}%
\pgfpathlineto{\pgfqpoint{4.675278in}{3.093140in}}%
\pgfpathquadraticcurveto{\pgfqpoint{4.670969in}{3.086564in}}{\pgfqpoint{4.665258in}{3.083368in}}%
\pgfpathquadraticcurveto{\pgfqpoint{4.659568in}{3.080173in}}{\pgfqpoint{4.652023in}{3.080173in}}%
\pgfpathquadraticcurveto{\pgfqpoint{4.639591in}{3.080173in}}{\pgfqpoint{4.633138in}{3.087904in}}%
\pgfpathquadraticcurveto{\pgfqpoint{4.626706in}{3.095635in}}{\pgfqpoint{4.626706in}{3.110520in}}%
\pgfpathlineto{\pgfqpoint{4.626706in}{3.110520in}}%
\pgfpathclose%
\pgfpathmoveto{\pgfqpoint{4.656538in}{3.155938in}}%
\pgfpathlineto{\pgfqpoint{4.656538in}{3.155938in}}%
\pgfpathlineto{\pgfqpoint{4.656538in}{3.155938in}}%
\pgfpathclose%
\pgfpathmoveto{\pgfqpoint{4.759021in}{3.143259in}}%
\pgfpathlineto{\pgfqpoint{4.759021in}{3.182306in}}%
\pgfpathlineto{\pgfqpoint{4.770876in}{3.182306in}}%
\pgfpathlineto{\pgfqpoint{4.770876in}{3.082049in}}%
\pgfpathlineto{\pgfqpoint{4.759021in}{3.082049in}}%
\pgfpathlineto{\pgfqpoint{4.759021in}{3.092872in}}%
\pgfpathquadraticcurveto{\pgfqpoint{4.755290in}{3.086440in}}{\pgfqpoint{4.749579in}{3.083306in}}%
\pgfpathquadraticcurveto{\pgfqpoint{4.743889in}{3.080173in}}{\pgfqpoint{4.735890in}{3.080173in}}%
\pgfpathquadraticcurveto{\pgfqpoint{4.722819in}{3.080173in}}{\pgfqpoint{4.714593in}{3.090605in}}%
\pgfpathquadraticcurveto{\pgfqpoint{4.706388in}{3.101057in}}{\pgfqpoint{4.706388in}{3.118066in}}%
\pgfpathquadraticcurveto{\pgfqpoint{4.706388in}{3.135074in}}{\pgfqpoint{4.714593in}{3.145506in}}%
\pgfpathquadraticcurveto{\pgfqpoint{4.722819in}{3.155938in}}{\pgfqpoint{4.735890in}{3.155938in}}%
\pgfpathquadraticcurveto{\pgfqpoint{4.743889in}{3.155938in}}{\pgfqpoint{4.749579in}{3.152804in}}%
\pgfpathquadraticcurveto{\pgfqpoint{4.755290in}{3.149691in}}{\pgfqpoint{4.759021in}{3.143259in}}%
\pgfpathlineto{\pgfqpoint{4.759021in}{3.143259in}}%
\pgfpathclose%
\pgfpathmoveto{\pgfqpoint{4.718634in}{3.118066in}}%
\pgfpathquadraticcurveto{\pgfqpoint{4.718634in}{3.104995in}}{\pgfqpoint{4.724015in}{3.097552in}}%
\pgfpathquadraticcurveto{\pgfqpoint{4.729396in}{3.090110in}}{\pgfqpoint{4.738797in}{3.090110in}}%
\pgfpathquadraticcurveto{\pgfqpoint{4.748198in}{3.090110in}}{\pgfqpoint{4.753599in}{3.097552in}}%
\pgfpathquadraticcurveto{\pgfqpoint{4.759021in}{3.104995in}}{\pgfqpoint{4.759021in}{3.118066in}}%
\pgfpathquadraticcurveto{\pgfqpoint{4.759021in}{3.131136in}}{\pgfqpoint{4.753599in}{3.138579in}}%
\pgfpathquadraticcurveto{\pgfqpoint{4.748198in}{3.146021in}}{\pgfqpoint{4.738797in}{3.146021in}}%
\pgfpathquadraticcurveto{\pgfqpoint{4.729396in}{3.146021in}}{\pgfqpoint{4.724015in}{3.138579in}}%
\pgfpathquadraticcurveto{\pgfqpoint{4.718634in}{3.131136in}}{\pgfqpoint{4.718634in}{3.118066in}}%
\pgfpathlineto{\pgfqpoint{4.718634in}{3.118066in}}%
\pgfpathclose%
\pgfpathmoveto{\pgfqpoint{4.825323in}{3.075348in}}%
\pgfpathquadraticcurveto{\pgfqpoint{4.820314in}{3.062463in}}{\pgfqpoint{4.815531in}{3.058546in}}%
\pgfpathquadraticcurveto{\pgfqpoint{4.810768in}{3.054608in}}{\pgfqpoint{4.802790in}{3.054608in}}%
\pgfpathlineto{\pgfqpoint{4.793306in}{3.054608in}}%
\pgfpathlineto{\pgfqpoint{4.793306in}{3.064525in}}%
\pgfpathlineto{\pgfqpoint{4.800275in}{3.064525in}}%
\pgfpathquadraticcurveto{\pgfqpoint{4.805161in}{3.064525in}}{\pgfqpoint{4.807861in}{3.066855in}}%
\pgfpathquadraticcurveto{\pgfqpoint{4.810583in}{3.069164in}}{\pgfqpoint{4.813861in}{3.077802in}}%
\pgfpathlineto{\pgfqpoint{4.815984in}{3.083203in}}%
\pgfpathlineto{\pgfqpoint{4.786812in}{3.154206in}}%
\pgfpathlineto{\pgfqpoint{4.799367in}{3.154206in}}%
\pgfpathlineto{\pgfqpoint{4.821922in}{3.097779in}}%
\pgfpathlineto{\pgfqpoint{4.844476in}{3.154206in}}%
\pgfpathlineto{\pgfqpoint{4.857031in}{3.154206in}}%
\pgfpathlineto{\pgfqpoint{4.825323in}{3.075348in}}%
\pgfpathlineto{\pgfqpoint{4.825323in}{3.075348in}}%
\pgfpathclose%
\pgfpathmoveto{\pgfqpoint{4.913726in}{3.178245in}}%
\pgfpathlineto{\pgfqpoint{4.975575in}{3.178245in}}%
\pgfpathlineto{\pgfqpoint{4.975575in}{3.172699in}}%
\pgfpathlineto{\pgfqpoint{4.940651in}{3.082049in}}%
\pgfpathlineto{\pgfqpoint{4.927065in}{3.082049in}}%
\pgfpathlineto{\pgfqpoint{4.959927in}{3.167277in}}%
\pgfpathlineto{\pgfqpoint{4.913726in}{3.167277in}}%
\pgfpathlineto{\pgfqpoint{4.913726in}{3.178245in}}%
\pgfpathlineto{\pgfqpoint{4.913726in}{3.178245in}}%
\pgfpathclose%
\pgfpathmoveto{\pgfqpoint{5.069689in}{3.182162in}}%
\pgfpathquadraticcurveto{\pgfqpoint{5.061071in}{3.167359in}}{\pgfqpoint{5.056865in}{3.152845in}}%
\pgfpathquadraticcurveto{\pgfqpoint{5.052680in}{3.138352in}}{\pgfqpoint{5.052680in}{3.123467in}}%
\pgfpathquadraticcurveto{\pgfqpoint{5.052680in}{3.108603in}}{\pgfqpoint{5.056906in}{3.094006in}}%
\pgfpathquadraticcurveto{\pgfqpoint{5.061133in}{3.079410in}}{\pgfqpoint{5.069689in}{3.064649in}}%
\pgfpathlineto{\pgfqpoint{5.059380in}{3.064649in}}%
\pgfpathquadraticcurveto{\pgfqpoint{5.049732in}{3.079802in}}{\pgfqpoint{5.044928in}{3.094419in}}%
\pgfpathquadraticcurveto{\pgfqpoint{5.040125in}{3.109036in}}{\pgfqpoint{5.040125in}{3.123467in}}%
\pgfpathquadraticcurveto{\pgfqpoint{5.040125in}{3.137837in}}{\pgfqpoint{5.044887in}{3.152392in}}%
\pgfpathquadraticcurveto{\pgfqpoint{5.049670in}{3.166967in}}{\pgfqpoint{5.059380in}{3.182162in}}%
\pgfpathlineto{\pgfqpoint{5.069689in}{3.182162in}}%
\pgfpathlineto{\pgfqpoint{5.069689in}{3.182162in}}%
\pgfpathclose%
\pgfpathmoveto{\pgfqpoint{5.105582in}{3.092996in}}%
\pgfpathlineto{\pgfqpoint{5.150999in}{3.092996in}}%
\pgfpathlineto{\pgfqpoint{5.150999in}{3.082049in}}%
\pgfpathlineto{\pgfqpoint{5.089934in}{3.082049in}}%
\pgfpathlineto{\pgfqpoint{5.089934in}{3.092996in}}%
\pgfpathquadraticcurveto{\pgfqpoint{5.097335in}{3.100665in}}{\pgfqpoint{5.110117in}{3.113571in}}%
\pgfpathquadraticcurveto{\pgfqpoint{5.122920in}{3.126498in}}{\pgfqpoint{5.126198in}{3.130250in}}%
\pgfpathquadraticcurveto{\pgfqpoint{5.132445in}{3.137259in}}{\pgfqpoint{5.134919in}{3.142125in}}%
\pgfpathquadraticcurveto{\pgfqpoint{5.137413in}{3.146990in}}{\pgfqpoint{5.137413in}{3.151691in}}%
\pgfpathquadraticcurveto{\pgfqpoint{5.137413in}{3.159360in}}{\pgfqpoint{5.132032in}{3.164184in}}%
\pgfpathquadraticcurveto{\pgfqpoint{5.126651in}{3.169029in}}{\pgfqpoint{5.118013in}{3.169029in}}%
\pgfpathquadraticcurveto{\pgfqpoint{5.111890in}{3.169029in}}{\pgfqpoint{5.105087in}{3.166906in}}%
\pgfpathquadraticcurveto{\pgfqpoint{5.098304in}{3.164782in}}{\pgfqpoint{5.090573in}{3.160453in}}%
\pgfpathlineto{\pgfqpoint{5.090573in}{3.173606in}}%
\pgfpathquadraticcurveto{\pgfqpoint{5.098428in}{3.176760in}}{\pgfqpoint{5.105252in}{3.178368in}}%
\pgfpathquadraticcurveto{\pgfqpoint{5.112096in}{3.179976in}}{\pgfqpoint{5.117766in}{3.179976in}}%
\pgfpathquadraticcurveto{\pgfqpoint{5.132713in}{3.179976in}}{\pgfqpoint{5.141598in}{3.172493in}}%
\pgfpathquadraticcurveto{\pgfqpoint{5.150484in}{3.165029in}}{\pgfqpoint{5.150484in}{3.152536in}}%
\pgfpathquadraticcurveto{\pgfqpoint{5.150484in}{3.146598in}}{\pgfqpoint{5.148257in}{3.141279in}}%
\pgfpathquadraticcurveto{\pgfqpoint{5.146051in}{3.135981in}}{\pgfqpoint{5.140176in}{3.128765in}}%
\pgfpathquadraticcurveto{\pgfqpoint{5.138568in}{3.126889in}}{\pgfqpoint{5.129929in}{3.117962in}}%
\pgfpathquadraticcurveto{\pgfqpoint{5.121312in}{3.109036in}}{\pgfqpoint{5.105582in}{3.092996in}}%
\pgfpathlineto{\pgfqpoint{5.105582in}{3.092996in}}%
\pgfpathclose%
\pgfpathmoveto{\pgfqpoint{5.206148in}{3.169668in}}%
\pgfpathquadraticcurveto{\pgfqpoint{5.196108in}{3.169668in}}{\pgfqpoint{5.191036in}{3.159772in}}%
\pgfpathquadraticcurveto{\pgfqpoint{5.185985in}{3.149897in}}{\pgfqpoint{5.185985in}{3.130044in}}%
\pgfpathquadraticcurveto{\pgfqpoint{5.185985in}{3.110273in}}{\pgfqpoint{5.191036in}{3.100377in}}%
\pgfpathquadraticcurveto{\pgfqpoint{5.196108in}{3.090481in}}{\pgfqpoint{5.206148in}{3.090481in}}%
\pgfpathquadraticcurveto{\pgfqpoint{5.216271in}{3.090481in}}{\pgfqpoint{5.221322in}{3.100377in}}%
\pgfpathquadraticcurveto{\pgfqpoint{5.226393in}{3.110273in}}{\pgfqpoint{5.226393in}{3.130044in}}%
\pgfpathquadraticcurveto{\pgfqpoint{5.226393in}{3.149897in}}{\pgfqpoint{5.221322in}{3.159772in}}%
\pgfpathquadraticcurveto{\pgfqpoint{5.216271in}{3.169668in}}{\pgfqpoint{5.206148in}{3.169668in}}%
\pgfpathlineto{\pgfqpoint{5.206148in}{3.169668in}}%
\pgfpathclose%
\pgfpathmoveto{\pgfqpoint{5.206148in}{3.179976in}}%
\pgfpathquadraticcurveto{\pgfqpoint{5.222332in}{3.179976in}}{\pgfqpoint{5.230867in}{3.167174in}}%
\pgfpathquadraticcurveto{\pgfqpoint{5.239402in}{3.154391in}}{\pgfqpoint{5.239402in}{3.130044in}}%
\pgfpathquadraticcurveto{\pgfqpoint{5.239402in}{3.105758in}}{\pgfqpoint{5.230867in}{3.092955in}}%
\pgfpathquadraticcurveto{\pgfqpoint{5.222332in}{3.080173in}}{\pgfqpoint{5.206148in}{3.080173in}}%
\pgfpathquadraticcurveto{\pgfqpoint{5.189985in}{3.080173in}}{\pgfqpoint{5.181450in}{3.092955in}}%
\pgfpathquadraticcurveto{\pgfqpoint{5.172914in}{3.105758in}}{\pgfqpoint{5.172914in}{3.130044in}}%
\pgfpathquadraticcurveto{\pgfqpoint{5.172914in}{3.154391in}}{\pgfqpoint{5.181450in}{3.167174in}}%
\pgfpathquadraticcurveto{\pgfqpoint{5.189985in}{3.179976in}}{\pgfqpoint{5.206148in}{3.179976in}}%
\pgfpathlineto{\pgfqpoint{5.206148in}{3.179976in}}%
\pgfpathclose%
\pgfpathmoveto{\pgfqpoint{5.264533in}{3.092996in}}%
\pgfpathlineto{\pgfqpoint{5.285789in}{3.092996in}}%
\pgfpathlineto{\pgfqpoint{5.285789in}{3.166390in}}%
\pgfpathlineto{\pgfqpoint{5.262657in}{3.161751in}}%
\pgfpathlineto{\pgfqpoint{5.262657in}{3.173606in}}%
\pgfpathlineto{\pgfqpoint{5.285665in}{3.178245in}}%
\pgfpathlineto{\pgfqpoint{5.298674in}{3.178245in}}%
\pgfpathlineto{\pgfqpoint{5.298674in}{3.092996in}}%
\pgfpathlineto{\pgfqpoint{5.319929in}{3.092996in}}%
\pgfpathlineto{\pgfqpoint{5.319929in}{3.082049in}}%
\pgfpathlineto{\pgfqpoint{5.264533in}{3.082049in}}%
\pgfpathlineto{\pgfqpoint{5.264533in}{3.092996in}}%
\pgfpathlineto{\pgfqpoint{5.264533in}{3.092996in}}%
\pgfpathclose%
\pgfpathmoveto{\pgfqpoint{5.346607in}{3.084049in}}%
\pgfpathlineto{\pgfqpoint{5.346607in}{3.095903in}}%
\pgfpathquadraticcurveto{\pgfqpoint{5.351514in}{3.093573in}}{\pgfqpoint{5.356523in}{3.092357in}}%
\pgfpathquadraticcurveto{\pgfqpoint{5.361554in}{3.091141in}}{\pgfqpoint{5.366399in}{3.091141in}}%
\pgfpathquadraticcurveto{\pgfqpoint{5.379284in}{3.091141in}}{\pgfqpoint{5.386067in}{3.099799in}}%
\pgfpathquadraticcurveto{\pgfqpoint{5.392870in}{3.108458in}}{\pgfqpoint{5.393839in}{3.126126in}}%
\pgfpathquadraticcurveto{\pgfqpoint{5.390107in}{3.120581in}}{\pgfqpoint{5.384355in}{3.117612in}}%
\pgfpathquadraticcurveto{\pgfqpoint{5.378624in}{3.114643in}}{\pgfqpoint{5.371676in}{3.114643in}}%
\pgfpathquadraticcurveto{\pgfqpoint{5.357245in}{3.114643in}}{\pgfqpoint{5.348834in}{3.123364in}}%
\pgfpathquadraticcurveto{\pgfqpoint{5.340422in}{3.132105in}}{\pgfqpoint{5.340422in}{3.147258in}}%
\pgfpathquadraticcurveto{\pgfqpoint{5.340422in}{3.162061in}}{\pgfqpoint{5.349184in}{3.171008in}}%
\pgfpathquadraticcurveto{\pgfqpoint{5.357946in}{3.179976in}}{\pgfqpoint{5.372501in}{3.179976in}}%
\pgfpathquadraticcurveto{\pgfqpoint{5.389200in}{3.179976in}}{\pgfqpoint{5.397983in}{3.167174in}}%
\pgfpathquadraticcurveto{\pgfqpoint{5.406786in}{3.154391in}}{\pgfqpoint{5.406786in}{3.130044in}}%
\pgfpathquadraticcurveto{\pgfqpoint{5.406786in}{3.107304in}}{\pgfqpoint{5.395983in}{3.093738in}}%
\pgfpathquadraticcurveto{\pgfqpoint{5.385201in}{3.080173in}}{\pgfqpoint{5.366976in}{3.080173in}}%
\pgfpathquadraticcurveto{\pgfqpoint{5.362069in}{3.080173in}}{\pgfqpoint{5.357039in}{3.081142in}}%
\pgfpathquadraticcurveto{\pgfqpoint{5.352029in}{3.082111in}}{\pgfqpoint{5.346607in}{3.084049in}}%
\pgfpathlineto{\pgfqpoint{5.346607in}{3.084049in}}%
\pgfpathclose%
\pgfpathmoveto{\pgfqpoint{5.372501in}{3.124828in}}%
\pgfpathquadraticcurveto{\pgfqpoint{5.381263in}{3.124828in}}{\pgfqpoint{5.386376in}{3.130806in}}%
\pgfpathquadraticcurveto{\pgfqpoint{5.391509in}{3.136806in}}{\pgfqpoint{5.391509in}{3.147258in}}%
\pgfpathquadraticcurveto{\pgfqpoint{5.391509in}{3.157628in}}{\pgfqpoint{5.386376in}{3.163648in}}%
\pgfpathquadraticcurveto{\pgfqpoint{5.381263in}{3.169668in}}{\pgfqpoint{5.372501in}{3.169668in}}%
\pgfpathquadraticcurveto{\pgfqpoint{5.363739in}{3.169668in}}{\pgfqpoint{5.358626in}{3.163648in}}%
\pgfpathquadraticcurveto{\pgfqpoint{5.353513in}{3.157628in}}{\pgfqpoint{5.353513in}{3.147258in}}%
\pgfpathquadraticcurveto{\pgfqpoint{5.353513in}{3.136806in}}{\pgfqpoint{5.358626in}{3.130806in}}%
\pgfpathquadraticcurveto{\pgfqpoint{5.363739in}{3.124828in}}{\pgfqpoint{5.372501in}{3.124828in}}%
\pgfpathlineto{\pgfqpoint{5.372501in}{3.124828in}}%
\pgfpathclose%
\pgfpathmoveto{\pgfqpoint{5.426640in}{3.182162in}}%
\pgfpathlineto{\pgfqpoint{5.436948in}{3.182162in}}%
\pgfpathquadraticcurveto{\pgfqpoint{5.446596in}{3.166967in}}{\pgfqpoint{5.451400in}{3.152392in}}%
\pgfpathquadraticcurveto{\pgfqpoint{5.456203in}{3.137837in}}{\pgfqpoint{5.456203in}{3.123467in}}%
\pgfpathquadraticcurveto{\pgfqpoint{5.456203in}{3.109036in}}{\pgfqpoint{5.451400in}{3.094419in}}%
\pgfpathquadraticcurveto{\pgfqpoint{5.446596in}{3.079802in}}{\pgfqpoint{5.436948in}{3.064649in}}%
\pgfpathlineto{\pgfqpoint{5.426640in}{3.064649in}}%
\pgfpathquadraticcurveto{\pgfqpoint{5.435195in}{3.079410in}}{\pgfqpoint{5.439422in}{3.094006in}}%
\pgfpathquadraticcurveto{\pgfqpoint{5.443648in}{3.108603in}}{\pgfqpoint{5.443648in}{3.123467in}}%
\pgfpathquadraticcurveto{\pgfqpoint{5.443648in}{3.138352in}}{\pgfqpoint{5.439422in}{3.152845in}}%
\pgfpathquadraticcurveto{\pgfqpoint{5.435195in}{3.167359in}}{\pgfqpoint{5.426640in}{3.182162in}}%
\pgfpathlineto{\pgfqpoint{5.426640in}{3.182162in}}%
\pgfpathclose%
\pgfusepath{fill}%
\end{pgfscope}%
\begin{pgfscope}%
\pgfsetbuttcap%
\pgfsetroundjoin%
\definecolor{currentfill}{rgb}{0.090196,0.168627,0.227451}%
\pgfsetfillcolor{currentfill}%
\pgfsetlinewidth{0.000000pt}%
\definecolor{currentstroke}{rgb}{0.090196,0.168627,0.227451}%
\pgfsetstrokecolor{currentstroke}%
\pgfsetdash{}{0pt}%
\pgfpathmoveto{\pgfqpoint{4.550611in}{2.775490in}}%
\pgfpathlineto{\pgfqpoint{4.550611in}{2.762791in}}%
\pgfpathquadraticcurveto{\pgfqpoint{4.543210in}{2.766337in}}{\pgfqpoint{4.536633in}{2.768068in}}%
\pgfpathquadraticcurveto{\pgfqpoint{4.530056in}{2.769821in}}{\pgfqpoint{4.523933in}{2.769821in}}%
\pgfpathquadraticcurveto{\pgfqpoint{4.513316in}{2.769821in}}{\pgfqpoint{4.507543in}{2.765697in}}%
\pgfpathquadraticcurveto{\pgfqpoint{4.501771in}{2.761574in}}{\pgfqpoint{4.501771in}{2.753967in}}%
\pgfpathquadraticcurveto{\pgfqpoint{4.501771in}{2.747576in}}{\pgfqpoint{4.505605in}{2.744318in}}%
\pgfpathquadraticcurveto{\pgfqpoint{4.509440in}{2.741082in}}{\pgfqpoint{4.520140in}{2.739082in}}%
\pgfpathlineto{\pgfqpoint{4.527995in}{2.737474in}}%
\pgfpathquadraticcurveto{\pgfqpoint{4.542550in}{2.734691in}}{\pgfqpoint{4.549477in}{2.727702in}}%
\pgfpathquadraticcurveto{\pgfqpoint{4.556404in}{2.720713in}}{\pgfqpoint{4.556404in}{2.709003in}}%
\pgfpathquadraticcurveto{\pgfqpoint{4.556404in}{2.695004in}}{\pgfqpoint{4.547024in}{2.687788in}}%
\pgfpathquadraticcurveto{\pgfqpoint{4.537664in}{2.680573in}}{\pgfqpoint{4.519563in}{2.680573in}}%
\pgfpathquadraticcurveto{\pgfqpoint{4.512739in}{2.680573in}}{\pgfqpoint{4.505028in}{2.682119in}}%
\pgfpathquadraticcurveto{\pgfqpoint{4.497338in}{2.683665in}}{\pgfqpoint{4.489092in}{2.686696in}}%
\pgfpathlineto{\pgfqpoint{4.489092in}{2.700096in}}%
\pgfpathquadraticcurveto{\pgfqpoint{4.497008in}{2.695664in}}{\pgfqpoint{4.504616in}{2.693396in}}%
\pgfpathquadraticcurveto{\pgfqpoint{4.512223in}{2.691149in}}{\pgfqpoint{4.519563in}{2.691149in}}%
\pgfpathquadraticcurveto{\pgfqpoint{4.530696in}{2.691149in}}{\pgfqpoint{4.536757in}{2.695520in}}%
\pgfpathquadraticcurveto{\pgfqpoint{4.542818in}{2.699911in}}{\pgfqpoint{4.542818in}{2.708034in}}%
\pgfpathquadraticcurveto{\pgfqpoint{4.542818in}{2.715105in}}{\pgfqpoint{4.538468in}{2.719105in}}%
\pgfpathquadraticcurveto{\pgfqpoint{4.534118in}{2.723104in}}{\pgfqpoint{4.524201in}{2.725104in}}%
\pgfpathlineto{\pgfqpoint{4.516264in}{2.726650in}}%
\pgfpathquadraticcurveto{\pgfqpoint{4.501709in}{2.729536in}}{\pgfqpoint{4.495194in}{2.735721in}}%
\pgfpathquadraticcurveto{\pgfqpoint{4.488700in}{2.741906in}}{\pgfqpoint{4.488700in}{2.752936in}}%
\pgfpathquadraticcurveto{\pgfqpoint{4.488700in}{2.765697in}}{\pgfqpoint{4.497689in}{2.773037in}}%
\pgfpathquadraticcurveto{\pgfqpoint{4.506678in}{2.780376in}}{\pgfqpoint{4.522449in}{2.780376in}}%
\pgfpathquadraticcurveto{\pgfqpoint{4.529232in}{2.780376in}}{\pgfqpoint{4.536241in}{2.779139in}}%
\pgfpathquadraticcurveto{\pgfqpoint{4.543271in}{2.777923in}}{\pgfqpoint{4.550611in}{2.775490in}}%
\pgfpathlineto{\pgfqpoint{4.550611in}{2.775490in}}%
\pgfpathclose%
\pgfpathmoveto{\pgfqpoint{4.587926in}{2.775099in}}%
\pgfpathlineto{\pgfqpoint{4.587926in}{2.754606in}}%
\pgfpathlineto{\pgfqpoint{4.612336in}{2.754606in}}%
\pgfpathlineto{\pgfqpoint{4.612336in}{2.745390in}}%
\pgfpathlineto{\pgfqpoint{4.587926in}{2.745390in}}%
\pgfpathlineto{\pgfqpoint{4.587926in}{2.706219in}}%
\pgfpathquadraticcurveto{\pgfqpoint{4.587926in}{2.697396in}}{\pgfqpoint{4.590339in}{2.694880in}}%
\pgfpathquadraticcurveto{\pgfqpoint{4.592751in}{2.692365in}}{\pgfqpoint{4.600173in}{2.692365in}}%
\pgfpathlineto{\pgfqpoint{4.612336in}{2.692365in}}%
\pgfpathlineto{\pgfqpoint{4.612336in}{2.682449in}}%
\pgfpathlineto{\pgfqpoint{4.600173in}{2.682449in}}%
\pgfpathquadraticcurveto{\pgfqpoint{4.586442in}{2.682449in}}{\pgfqpoint{4.581226in}{2.687562in}}%
\pgfpathquadraticcurveto{\pgfqpoint{4.576010in}{2.692695in}}{\pgfqpoint{4.576010in}{2.706219in}}%
\pgfpathlineto{\pgfqpoint{4.576010in}{2.745390in}}%
\pgfpathlineto{\pgfqpoint{4.567310in}{2.745390in}}%
\pgfpathlineto{\pgfqpoint{4.567310in}{2.754606in}}%
\pgfpathlineto{\pgfqpoint{4.576010in}{2.754606in}}%
\pgfpathlineto{\pgfqpoint{4.576010in}{2.775099in}}%
\pgfpathlineto{\pgfqpoint{4.587926in}{2.775099in}}%
\pgfpathlineto{\pgfqpoint{4.587926in}{2.775099in}}%
\pgfpathclose%
\pgfpathmoveto{\pgfqpoint{4.626706in}{2.710920in}}%
\pgfpathlineto{\pgfqpoint{4.626706in}{2.754606in}}%
\pgfpathlineto{\pgfqpoint{4.638560in}{2.754606in}}%
\pgfpathlineto{\pgfqpoint{4.638560in}{2.711373in}}%
\pgfpathquadraticcurveto{\pgfqpoint{4.638560in}{2.701127in}}{\pgfqpoint{4.642539in}{2.695994in}}%
\pgfpathquadraticcurveto{\pgfqpoint{4.646539in}{2.690881in}}{\pgfqpoint{4.654538in}{2.690881in}}%
\pgfpathquadraticcurveto{\pgfqpoint{4.664124in}{2.690881in}}{\pgfqpoint{4.669691in}{2.697004in}}%
\pgfpathquadraticcurveto{\pgfqpoint{4.675278in}{2.703127in}}{\pgfqpoint{4.675278in}{2.713703in}}%
\pgfpathlineto{\pgfqpoint{4.675278in}{2.754606in}}%
\pgfpathlineto{\pgfqpoint{4.687132in}{2.754606in}}%
\pgfpathlineto{\pgfqpoint{4.687132in}{2.682449in}}%
\pgfpathlineto{\pgfqpoint{4.675278in}{2.682449in}}%
\pgfpathlineto{\pgfqpoint{4.675278in}{2.693540in}}%
\pgfpathquadraticcurveto{\pgfqpoint{4.670969in}{2.686964in}}{\pgfqpoint{4.665258in}{2.683768in}}%
\pgfpathquadraticcurveto{\pgfqpoint{4.659568in}{2.680573in}}{\pgfqpoint{4.652023in}{2.680573in}}%
\pgfpathquadraticcurveto{\pgfqpoint{4.639591in}{2.680573in}}{\pgfqpoint{4.633138in}{2.688304in}}%
\pgfpathquadraticcurveto{\pgfqpoint{4.626706in}{2.696035in}}{\pgfqpoint{4.626706in}{2.710920in}}%
\pgfpathlineto{\pgfqpoint{4.626706in}{2.710920in}}%
\pgfpathclose%
\pgfpathmoveto{\pgfqpoint{4.656538in}{2.756338in}}%
\pgfpathlineto{\pgfqpoint{4.656538in}{2.756338in}}%
\pgfpathlineto{\pgfqpoint{4.656538in}{2.756338in}}%
\pgfpathclose%
\pgfpathmoveto{\pgfqpoint{4.759021in}{2.743659in}}%
\pgfpathlineto{\pgfqpoint{4.759021in}{2.782706in}}%
\pgfpathlineto{\pgfqpoint{4.770876in}{2.782706in}}%
\pgfpathlineto{\pgfqpoint{4.770876in}{2.682449in}}%
\pgfpathlineto{\pgfqpoint{4.759021in}{2.682449in}}%
\pgfpathlineto{\pgfqpoint{4.759021in}{2.693272in}}%
\pgfpathquadraticcurveto{\pgfqpoint{4.755290in}{2.686840in}}{\pgfqpoint{4.749579in}{2.683706in}}%
\pgfpathquadraticcurveto{\pgfqpoint{4.743889in}{2.680573in}}{\pgfqpoint{4.735890in}{2.680573in}}%
\pgfpathquadraticcurveto{\pgfqpoint{4.722819in}{2.680573in}}{\pgfqpoint{4.714593in}{2.691005in}}%
\pgfpathquadraticcurveto{\pgfqpoint{4.706388in}{2.701457in}}{\pgfqpoint{4.706388in}{2.718466in}}%
\pgfpathquadraticcurveto{\pgfqpoint{4.706388in}{2.735474in}}{\pgfqpoint{4.714593in}{2.745906in}}%
\pgfpathquadraticcurveto{\pgfqpoint{4.722819in}{2.756338in}}{\pgfqpoint{4.735890in}{2.756338in}}%
\pgfpathquadraticcurveto{\pgfqpoint{4.743889in}{2.756338in}}{\pgfqpoint{4.749579in}{2.753204in}}%
\pgfpathquadraticcurveto{\pgfqpoint{4.755290in}{2.750091in}}{\pgfqpoint{4.759021in}{2.743659in}}%
\pgfpathlineto{\pgfqpoint{4.759021in}{2.743659in}}%
\pgfpathclose%
\pgfpathmoveto{\pgfqpoint{4.718634in}{2.718466in}}%
\pgfpathquadraticcurveto{\pgfqpoint{4.718634in}{2.705395in}}{\pgfqpoint{4.724015in}{2.697952in}}%
\pgfpathquadraticcurveto{\pgfqpoint{4.729396in}{2.690510in}}{\pgfqpoint{4.738797in}{2.690510in}}%
\pgfpathquadraticcurveto{\pgfqpoint{4.748198in}{2.690510in}}{\pgfqpoint{4.753599in}{2.697952in}}%
\pgfpathquadraticcurveto{\pgfqpoint{4.759021in}{2.705395in}}{\pgfqpoint{4.759021in}{2.718466in}}%
\pgfpathquadraticcurveto{\pgfqpoint{4.759021in}{2.731536in}}{\pgfqpoint{4.753599in}{2.738979in}}%
\pgfpathquadraticcurveto{\pgfqpoint{4.748198in}{2.746421in}}{\pgfqpoint{4.738797in}{2.746421in}}%
\pgfpathquadraticcurveto{\pgfqpoint{4.729396in}{2.746421in}}{\pgfqpoint{4.724015in}{2.738979in}}%
\pgfpathquadraticcurveto{\pgfqpoint{4.718634in}{2.731536in}}{\pgfqpoint{4.718634in}{2.718466in}}%
\pgfpathlineto{\pgfqpoint{4.718634in}{2.718466in}}%
\pgfpathclose%
\pgfpathmoveto{\pgfqpoint{4.825323in}{2.675748in}}%
\pgfpathquadraticcurveto{\pgfqpoint{4.820314in}{2.662863in}}{\pgfqpoint{4.815531in}{2.658946in}}%
\pgfpathquadraticcurveto{\pgfqpoint{4.810768in}{2.655008in}}{\pgfqpoint{4.802790in}{2.655008in}}%
\pgfpathlineto{\pgfqpoint{4.793306in}{2.655008in}}%
\pgfpathlineto{\pgfqpoint{4.793306in}{2.664925in}}%
\pgfpathlineto{\pgfqpoint{4.800275in}{2.664925in}}%
\pgfpathquadraticcurveto{\pgfqpoint{4.805161in}{2.664925in}}{\pgfqpoint{4.807861in}{2.667255in}}%
\pgfpathquadraticcurveto{\pgfqpoint{4.810583in}{2.669564in}}{\pgfqpoint{4.813861in}{2.678202in}}%
\pgfpathlineto{\pgfqpoint{4.815984in}{2.683603in}}%
\pgfpathlineto{\pgfqpoint{4.786812in}{2.754606in}}%
\pgfpathlineto{\pgfqpoint{4.799367in}{2.754606in}}%
\pgfpathlineto{\pgfqpoint{4.821922in}{2.698179in}}%
\pgfpathlineto{\pgfqpoint{4.844476in}{2.754606in}}%
\pgfpathlineto{\pgfqpoint{4.857031in}{2.754606in}}%
\pgfpathlineto{\pgfqpoint{4.825323in}{2.675748in}}%
\pgfpathlineto{\pgfqpoint{4.825323in}{2.675748in}}%
\pgfpathclose%
\pgfpathmoveto{\pgfqpoint{4.944836in}{2.728135in}}%
\pgfpathquadraticcurveto{\pgfqpoint{4.935559in}{2.728135in}}{\pgfqpoint{4.930240in}{2.723166in}}%
\pgfpathquadraticcurveto{\pgfqpoint{4.924941in}{2.718197in}}{\pgfqpoint{4.924941in}{2.709518in}}%
\pgfpathquadraticcurveto{\pgfqpoint{4.924941in}{2.700818in}}{\pgfqpoint{4.930240in}{2.695849in}}%
\pgfpathquadraticcurveto{\pgfqpoint{4.935559in}{2.690881in}}{\pgfqpoint{4.944836in}{2.690881in}}%
\pgfpathquadraticcurveto{\pgfqpoint{4.954113in}{2.690881in}}{\pgfqpoint{4.959453in}{2.695870in}}%
\pgfpathquadraticcurveto{\pgfqpoint{4.964813in}{2.700880in}}{\pgfqpoint{4.964813in}{2.709518in}}%
\pgfpathquadraticcurveto{\pgfqpoint{4.964813in}{2.718197in}}{\pgfqpoint{4.959494in}{2.723166in}}%
\pgfpathquadraticcurveto{\pgfqpoint{4.954196in}{2.728135in}}{\pgfqpoint{4.944836in}{2.728135in}}%
\pgfpathlineto{\pgfqpoint{4.944836in}{2.728135in}}%
\pgfpathclose%
\pgfpathmoveto{\pgfqpoint{4.931827in}{2.733660in}}%
\pgfpathquadraticcurveto{\pgfqpoint{4.923457in}{2.735721in}}{\pgfqpoint{4.918777in}{2.741453in}}%
\pgfpathquadraticcurveto{\pgfqpoint{4.914118in}{2.747205in}}{\pgfqpoint{4.914118in}{2.755451in}}%
\pgfpathquadraticcurveto{\pgfqpoint{4.914118in}{2.766976in}}{\pgfqpoint{4.922323in}{2.773676in}}%
\pgfpathquadraticcurveto{\pgfqpoint{4.930549in}{2.780376in}}{\pgfqpoint{4.944836in}{2.780376in}}%
\pgfpathquadraticcurveto{\pgfqpoint{4.959206in}{2.780376in}}{\pgfqpoint{4.967390in}{2.773676in}}%
\pgfpathquadraticcurveto{\pgfqpoint{4.975575in}{2.766976in}}{\pgfqpoint{4.975575in}{2.755451in}}%
\pgfpathquadraticcurveto{\pgfqpoint{4.975575in}{2.747205in}}{\pgfqpoint{4.970895in}{2.741453in}}%
\pgfpathquadraticcurveto{\pgfqpoint{4.966236in}{2.735721in}}{\pgfqpoint{4.957928in}{2.733660in}}%
\pgfpathquadraticcurveto{\pgfqpoint{4.967329in}{2.731474in}}{\pgfqpoint{4.972565in}{2.725083in}}%
\pgfpathquadraticcurveto{\pgfqpoint{4.977822in}{2.718713in}}{\pgfqpoint{4.977822in}{2.709518in}}%
\pgfpathquadraticcurveto{\pgfqpoint{4.977822in}{2.695520in}}{\pgfqpoint{4.969287in}{2.688036in}}%
\pgfpathquadraticcurveto{\pgfqpoint{4.960752in}{2.680573in}}{\pgfqpoint{4.944836in}{2.680573in}}%
\pgfpathquadraticcurveto{\pgfqpoint{4.928941in}{2.680573in}}{\pgfqpoint{4.920385in}{2.688036in}}%
\pgfpathquadraticcurveto{\pgfqpoint{4.911850in}{2.695520in}}{\pgfqpoint{4.911850in}{2.709518in}}%
\pgfpathquadraticcurveto{\pgfqpoint{4.911850in}{2.718713in}}{\pgfqpoint{4.917128in}{2.725083in}}%
\pgfpathquadraticcurveto{\pgfqpoint{4.922426in}{2.731474in}}{\pgfqpoint{4.931827in}{2.733660in}}%
\pgfpathlineto{\pgfqpoint{4.931827in}{2.733660in}}%
\pgfpathclose%
\pgfpathmoveto{\pgfqpoint{4.927065in}{2.754214in}}%
\pgfpathquadraticcurveto{\pgfqpoint{4.927065in}{2.746751in}}{\pgfqpoint{4.931724in}{2.742566in}}%
\pgfpathquadraticcurveto{\pgfqpoint{4.936404in}{2.738381in}}{\pgfqpoint{4.944836in}{2.738381in}}%
\pgfpathquadraticcurveto{\pgfqpoint{4.953227in}{2.738381in}}{\pgfqpoint{4.957948in}{2.742566in}}%
\pgfpathquadraticcurveto{\pgfqpoint{4.962690in}{2.746751in}}{\pgfqpoint{4.962690in}{2.754214in}}%
\pgfpathquadraticcurveto{\pgfqpoint{4.962690in}{2.761698in}}{\pgfqpoint{4.957948in}{2.765883in}}%
\pgfpathquadraticcurveto{\pgfqpoint{4.953227in}{2.770068in}}{\pgfqpoint{4.944836in}{2.770068in}}%
\pgfpathquadraticcurveto{\pgfqpoint{4.936404in}{2.770068in}}{\pgfqpoint{4.931724in}{2.765883in}}%
\pgfpathquadraticcurveto{\pgfqpoint{4.927065in}{2.761698in}}{\pgfqpoint{4.927065in}{2.754214in}}%
\pgfpathlineto{\pgfqpoint{4.927065in}{2.754214in}}%
\pgfpathclose%
\pgfpathmoveto{\pgfqpoint{5.069689in}{2.782562in}}%
\pgfpathquadraticcurveto{\pgfqpoint{5.061071in}{2.767759in}}{\pgfqpoint{5.056865in}{2.753245in}}%
\pgfpathquadraticcurveto{\pgfqpoint{5.052680in}{2.738752in}}{\pgfqpoint{5.052680in}{2.723867in}}%
\pgfpathquadraticcurveto{\pgfqpoint{5.052680in}{2.709003in}}{\pgfqpoint{5.056906in}{2.694406in}}%
\pgfpathquadraticcurveto{\pgfqpoint{5.061133in}{2.679810in}}{\pgfqpoint{5.069689in}{2.665049in}}%
\pgfpathlineto{\pgfqpoint{5.059380in}{2.665049in}}%
\pgfpathquadraticcurveto{\pgfqpoint{5.049732in}{2.680202in}}{\pgfqpoint{5.044928in}{2.694819in}}%
\pgfpathquadraticcurveto{\pgfqpoint{5.040125in}{2.709436in}}{\pgfqpoint{5.040125in}{2.723867in}}%
\pgfpathquadraticcurveto{\pgfqpoint{5.040125in}{2.738237in}}{\pgfqpoint{5.044887in}{2.752792in}}%
\pgfpathquadraticcurveto{\pgfqpoint{5.049670in}{2.767367in}}{\pgfqpoint{5.059380in}{2.782562in}}%
\pgfpathlineto{\pgfqpoint{5.069689in}{2.782562in}}%
\pgfpathlineto{\pgfqpoint{5.069689in}{2.782562in}}%
\pgfpathclose%
\pgfpathmoveto{\pgfqpoint{5.105582in}{2.693396in}}%
\pgfpathlineto{\pgfqpoint{5.150999in}{2.693396in}}%
\pgfpathlineto{\pgfqpoint{5.150999in}{2.682449in}}%
\pgfpathlineto{\pgfqpoint{5.089934in}{2.682449in}}%
\pgfpathlineto{\pgfqpoint{5.089934in}{2.693396in}}%
\pgfpathquadraticcurveto{\pgfqpoint{5.097335in}{2.701065in}}{\pgfqpoint{5.110117in}{2.713971in}}%
\pgfpathquadraticcurveto{\pgfqpoint{5.122920in}{2.726898in}}{\pgfqpoint{5.126198in}{2.730650in}}%
\pgfpathquadraticcurveto{\pgfqpoint{5.132445in}{2.737659in}}{\pgfqpoint{5.134919in}{2.742525in}}%
\pgfpathquadraticcurveto{\pgfqpoint{5.137413in}{2.747390in}}{\pgfqpoint{5.137413in}{2.752091in}}%
\pgfpathquadraticcurveto{\pgfqpoint{5.137413in}{2.759760in}}{\pgfqpoint{5.132032in}{2.764584in}}%
\pgfpathquadraticcurveto{\pgfqpoint{5.126651in}{2.769429in}}{\pgfqpoint{5.118013in}{2.769429in}}%
\pgfpathquadraticcurveto{\pgfqpoint{5.111890in}{2.769429in}}{\pgfqpoint{5.105087in}{2.767306in}}%
\pgfpathquadraticcurveto{\pgfqpoint{5.098304in}{2.765182in}}{\pgfqpoint{5.090573in}{2.760853in}}%
\pgfpathlineto{\pgfqpoint{5.090573in}{2.774006in}}%
\pgfpathquadraticcurveto{\pgfqpoint{5.098428in}{2.777160in}}{\pgfqpoint{5.105252in}{2.778768in}}%
\pgfpathquadraticcurveto{\pgfqpoint{5.112096in}{2.780376in}}{\pgfqpoint{5.117766in}{2.780376in}}%
\pgfpathquadraticcurveto{\pgfqpoint{5.132713in}{2.780376in}}{\pgfqpoint{5.141598in}{2.772893in}}%
\pgfpathquadraticcurveto{\pgfqpoint{5.150484in}{2.765429in}}{\pgfqpoint{5.150484in}{2.752936in}}%
\pgfpathquadraticcurveto{\pgfqpoint{5.150484in}{2.746998in}}{\pgfqpoint{5.148257in}{2.741679in}}%
\pgfpathquadraticcurveto{\pgfqpoint{5.146051in}{2.736381in}}{\pgfqpoint{5.140176in}{2.729165in}}%
\pgfpathquadraticcurveto{\pgfqpoint{5.138568in}{2.727289in}}{\pgfqpoint{5.129929in}{2.718362in}}%
\pgfpathquadraticcurveto{\pgfqpoint{5.121312in}{2.709436in}}{\pgfqpoint{5.105582in}{2.693396in}}%
\pgfpathlineto{\pgfqpoint{5.105582in}{2.693396in}}%
\pgfpathclose%
\pgfpathmoveto{\pgfqpoint{5.206148in}{2.770068in}}%
\pgfpathquadraticcurveto{\pgfqpoint{5.196108in}{2.770068in}}{\pgfqpoint{5.191036in}{2.760172in}}%
\pgfpathquadraticcurveto{\pgfqpoint{5.185985in}{2.750297in}}{\pgfqpoint{5.185985in}{2.730444in}}%
\pgfpathquadraticcurveto{\pgfqpoint{5.185985in}{2.710673in}}{\pgfqpoint{5.191036in}{2.700777in}}%
\pgfpathquadraticcurveto{\pgfqpoint{5.196108in}{2.690881in}}{\pgfqpoint{5.206148in}{2.690881in}}%
\pgfpathquadraticcurveto{\pgfqpoint{5.216271in}{2.690881in}}{\pgfqpoint{5.221322in}{2.700777in}}%
\pgfpathquadraticcurveto{\pgfqpoint{5.226393in}{2.710673in}}{\pgfqpoint{5.226393in}{2.730444in}}%
\pgfpathquadraticcurveto{\pgfqpoint{5.226393in}{2.750297in}}{\pgfqpoint{5.221322in}{2.760172in}}%
\pgfpathquadraticcurveto{\pgfqpoint{5.216271in}{2.770068in}}{\pgfqpoint{5.206148in}{2.770068in}}%
\pgfpathlineto{\pgfqpoint{5.206148in}{2.770068in}}%
\pgfpathclose%
\pgfpathmoveto{\pgfqpoint{5.206148in}{2.780376in}}%
\pgfpathquadraticcurveto{\pgfqpoint{5.222332in}{2.780376in}}{\pgfqpoint{5.230867in}{2.767574in}}%
\pgfpathquadraticcurveto{\pgfqpoint{5.239402in}{2.754791in}}{\pgfqpoint{5.239402in}{2.730444in}}%
\pgfpathquadraticcurveto{\pgfqpoint{5.239402in}{2.706158in}}{\pgfqpoint{5.230867in}{2.693355in}}%
\pgfpathquadraticcurveto{\pgfqpoint{5.222332in}{2.680573in}}{\pgfqpoint{5.206148in}{2.680573in}}%
\pgfpathquadraticcurveto{\pgfqpoint{5.189985in}{2.680573in}}{\pgfqpoint{5.181450in}{2.693355in}}%
\pgfpathquadraticcurveto{\pgfqpoint{5.172914in}{2.706158in}}{\pgfqpoint{5.172914in}{2.730444in}}%
\pgfpathquadraticcurveto{\pgfqpoint{5.172914in}{2.754791in}}{\pgfqpoint{5.181450in}{2.767574in}}%
\pgfpathquadraticcurveto{\pgfqpoint{5.189985in}{2.780376in}}{\pgfqpoint{5.206148in}{2.780376in}}%
\pgfpathlineto{\pgfqpoint{5.206148in}{2.780376in}}%
\pgfpathclose%
\pgfpathmoveto{\pgfqpoint{5.273481in}{2.693396in}}%
\pgfpathlineto{\pgfqpoint{5.318899in}{2.693396in}}%
\pgfpathlineto{\pgfqpoint{5.318899in}{2.682449in}}%
\pgfpathlineto{\pgfqpoint{5.257833in}{2.682449in}}%
\pgfpathlineto{\pgfqpoint{5.257833in}{2.693396in}}%
\pgfpathquadraticcurveto{\pgfqpoint{5.265234in}{2.701065in}}{\pgfqpoint{5.278016in}{2.713971in}}%
\pgfpathquadraticcurveto{\pgfqpoint{5.290819in}{2.726898in}}{\pgfqpoint{5.294097in}{2.730650in}}%
\pgfpathquadraticcurveto{\pgfqpoint{5.300344in}{2.737659in}}{\pgfqpoint{5.302818in}{2.742525in}}%
\pgfpathquadraticcurveto{\pgfqpoint{5.305313in}{2.747390in}}{\pgfqpoint{5.305313in}{2.752091in}}%
\pgfpathquadraticcurveto{\pgfqpoint{5.305313in}{2.759760in}}{\pgfqpoint{5.299932in}{2.764584in}}%
\pgfpathquadraticcurveto{\pgfqpoint{5.294551in}{2.769429in}}{\pgfqpoint{5.285913in}{2.769429in}}%
\pgfpathquadraticcurveto{\pgfqpoint{5.279789in}{2.769429in}}{\pgfqpoint{5.272986in}{2.767306in}}%
\pgfpathquadraticcurveto{\pgfqpoint{5.266203in}{2.765182in}}{\pgfqpoint{5.258472in}{2.760853in}}%
\pgfpathlineto{\pgfqpoint{5.258472in}{2.774006in}}%
\pgfpathquadraticcurveto{\pgfqpoint{5.266327in}{2.777160in}}{\pgfqpoint{5.273151in}{2.778768in}}%
\pgfpathquadraticcurveto{\pgfqpoint{5.279996in}{2.780376in}}{\pgfqpoint{5.285665in}{2.780376in}}%
\pgfpathquadraticcurveto{\pgfqpoint{5.300612in}{2.780376in}}{\pgfqpoint{5.309498in}{2.772893in}}%
\pgfpathquadraticcurveto{\pgfqpoint{5.318383in}{2.765429in}}{\pgfqpoint{5.318383in}{2.752936in}}%
\pgfpathquadraticcurveto{\pgfqpoint{5.318383in}{2.746998in}}{\pgfqpoint{5.316157in}{2.741679in}}%
\pgfpathquadraticcurveto{\pgfqpoint{5.313951in}{2.736381in}}{\pgfqpoint{5.308075in}{2.729165in}}%
\pgfpathquadraticcurveto{\pgfqpoint{5.306467in}{2.727289in}}{\pgfqpoint{5.297829in}{2.718362in}}%
\pgfpathquadraticcurveto{\pgfqpoint{5.289211in}{2.709436in}}{\pgfqpoint{5.273481in}{2.693396in}}%
\pgfpathlineto{\pgfqpoint{5.273481in}{2.693396in}}%
\pgfpathclose%
\pgfpathmoveto{\pgfqpoint{5.374047in}{2.770068in}}%
\pgfpathquadraticcurveto{\pgfqpoint{5.364007in}{2.770068in}}{\pgfqpoint{5.358936in}{2.760172in}}%
\pgfpathquadraticcurveto{\pgfqpoint{5.353885in}{2.750297in}}{\pgfqpoint{5.353885in}{2.730444in}}%
\pgfpathquadraticcurveto{\pgfqpoint{5.353885in}{2.710673in}}{\pgfqpoint{5.358936in}{2.700777in}}%
\pgfpathquadraticcurveto{\pgfqpoint{5.364007in}{2.690881in}}{\pgfqpoint{5.374047in}{2.690881in}}%
\pgfpathquadraticcurveto{\pgfqpoint{5.384170in}{2.690881in}}{\pgfqpoint{5.389221in}{2.700777in}}%
\pgfpathquadraticcurveto{\pgfqpoint{5.394293in}{2.710673in}}{\pgfqpoint{5.394293in}{2.730444in}}%
\pgfpathquadraticcurveto{\pgfqpoint{5.394293in}{2.750297in}}{\pgfqpoint{5.389221in}{2.760172in}}%
\pgfpathquadraticcurveto{\pgfqpoint{5.384170in}{2.770068in}}{\pgfqpoint{5.374047in}{2.770068in}}%
\pgfpathlineto{\pgfqpoint{5.374047in}{2.770068in}}%
\pgfpathclose%
\pgfpathmoveto{\pgfqpoint{5.374047in}{2.780376in}}%
\pgfpathquadraticcurveto{\pgfqpoint{5.390231in}{2.780376in}}{\pgfqpoint{5.398766in}{2.767574in}}%
\pgfpathquadraticcurveto{\pgfqpoint{5.407301in}{2.754791in}}{\pgfqpoint{5.407301in}{2.730444in}}%
\pgfpathquadraticcurveto{\pgfqpoint{5.407301in}{2.706158in}}{\pgfqpoint{5.398766in}{2.693355in}}%
\pgfpathquadraticcurveto{\pgfqpoint{5.390231in}{2.680573in}}{\pgfqpoint{5.374047in}{2.680573in}}%
\pgfpathquadraticcurveto{\pgfqpoint{5.357884in}{2.680573in}}{\pgfqpoint{5.349349in}{2.693355in}}%
\pgfpathquadraticcurveto{\pgfqpoint{5.340814in}{2.706158in}}{\pgfqpoint{5.340814in}{2.730444in}}%
\pgfpathquadraticcurveto{\pgfqpoint{5.340814in}{2.754791in}}{\pgfqpoint{5.349349in}{2.767574in}}%
\pgfpathquadraticcurveto{\pgfqpoint{5.357884in}{2.780376in}}{\pgfqpoint{5.374047in}{2.780376in}}%
\pgfpathlineto{\pgfqpoint{5.374047in}{2.780376in}}%
\pgfpathclose%
\pgfpathmoveto{\pgfqpoint{5.426640in}{2.782562in}}%
\pgfpathlineto{\pgfqpoint{5.436948in}{2.782562in}}%
\pgfpathquadraticcurveto{\pgfqpoint{5.446596in}{2.767367in}}{\pgfqpoint{5.451400in}{2.752792in}}%
\pgfpathquadraticcurveto{\pgfqpoint{5.456203in}{2.738237in}}{\pgfqpoint{5.456203in}{2.723867in}}%
\pgfpathquadraticcurveto{\pgfqpoint{5.456203in}{2.709436in}}{\pgfqpoint{5.451400in}{2.694819in}}%
\pgfpathquadraticcurveto{\pgfqpoint{5.446596in}{2.680202in}}{\pgfqpoint{5.436948in}{2.665049in}}%
\pgfpathlineto{\pgfqpoint{5.426640in}{2.665049in}}%
\pgfpathquadraticcurveto{\pgfqpoint{5.435195in}{2.679810in}}{\pgfqpoint{5.439422in}{2.694406in}}%
\pgfpathquadraticcurveto{\pgfqpoint{5.443648in}{2.709003in}}{\pgfqpoint{5.443648in}{2.723867in}}%
\pgfpathquadraticcurveto{\pgfqpoint{5.443648in}{2.738752in}}{\pgfqpoint{5.439422in}{2.753245in}}%
\pgfpathquadraticcurveto{\pgfqpoint{5.435195in}{2.767759in}}{\pgfqpoint{5.426640in}{2.782562in}}%
\pgfpathlineto{\pgfqpoint{5.426640in}{2.782562in}}%
\pgfpathclose%
\pgfusepath{fill}%
\end{pgfscope}%
\begin{pgfscope}%
\pgfsetbuttcap%
\pgfsetroundjoin%
\definecolor{currentfill}{rgb}{0.090196,0.168627,0.227451}%
\pgfsetfillcolor{currentfill}%
\pgfsetlinewidth{0.000000pt}%
\definecolor{currentstroke}{rgb}{0.090196,0.168627,0.227451}%
\pgfsetstrokecolor{currentstroke}%
\pgfsetdash{}{0pt}%
\pgfpathmoveto{\pgfqpoint{4.550611in}{2.375890in}}%
\pgfpathlineto{\pgfqpoint{4.550611in}{2.363191in}}%
\pgfpathquadraticcurveto{\pgfqpoint{4.543210in}{2.366737in}}{\pgfqpoint{4.536633in}{2.368468in}}%
\pgfpathquadraticcurveto{\pgfqpoint{4.530056in}{2.370221in}}{\pgfqpoint{4.523933in}{2.370221in}}%
\pgfpathquadraticcurveto{\pgfqpoint{4.513316in}{2.370221in}}{\pgfqpoint{4.507543in}{2.366097in}}%
\pgfpathquadraticcurveto{\pgfqpoint{4.501771in}{2.361974in}}{\pgfqpoint{4.501771in}{2.354367in}}%
\pgfpathquadraticcurveto{\pgfqpoint{4.501771in}{2.347976in}}{\pgfqpoint{4.505605in}{2.344718in}}%
\pgfpathquadraticcurveto{\pgfqpoint{4.509440in}{2.341482in}}{\pgfqpoint{4.520140in}{2.339482in}}%
\pgfpathlineto{\pgfqpoint{4.527995in}{2.337874in}}%
\pgfpathquadraticcurveto{\pgfqpoint{4.542550in}{2.335091in}}{\pgfqpoint{4.549477in}{2.328102in}}%
\pgfpathquadraticcurveto{\pgfqpoint{4.556404in}{2.321113in}}{\pgfqpoint{4.556404in}{2.309403in}}%
\pgfpathquadraticcurveto{\pgfqpoint{4.556404in}{2.295404in}}{\pgfqpoint{4.547024in}{2.288188in}}%
\pgfpathquadraticcurveto{\pgfqpoint{4.537664in}{2.280973in}}{\pgfqpoint{4.519563in}{2.280973in}}%
\pgfpathquadraticcurveto{\pgfqpoint{4.512739in}{2.280973in}}{\pgfqpoint{4.505028in}{2.282519in}}%
\pgfpathquadraticcurveto{\pgfqpoint{4.497338in}{2.284065in}}{\pgfqpoint{4.489092in}{2.287096in}}%
\pgfpathlineto{\pgfqpoint{4.489092in}{2.300496in}}%
\pgfpathquadraticcurveto{\pgfqpoint{4.497008in}{2.296064in}}{\pgfqpoint{4.504616in}{2.293796in}}%
\pgfpathquadraticcurveto{\pgfqpoint{4.512223in}{2.291549in}}{\pgfqpoint{4.519563in}{2.291549in}}%
\pgfpathquadraticcurveto{\pgfqpoint{4.530696in}{2.291549in}}{\pgfqpoint{4.536757in}{2.295920in}}%
\pgfpathquadraticcurveto{\pgfqpoint{4.542818in}{2.300311in}}{\pgfqpoint{4.542818in}{2.308434in}}%
\pgfpathquadraticcurveto{\pgfqpoint{4.542818in}{2.315505in}}{\pgfqpoint{4.538468in}{2.319505in}}%
\pgfpathquadraticcurveto{\pgfqpoint{4.534118in}{2.323504in}}{\pgfqpoint{4.524201in}{2.325504in}}%
\pgfpathlineto{\pgfqpoint{4.516264in}{2.327050in}}%
\pgfpathquadraticcurveto{\pgfqpoint{4.501709in}{2.329936in}}{\pgfqpoint{4.495194in}{2.336121in}}%
\pgfpathquadraticcurveto{\pgfqpoint{4.488700in}{2.342306in}}{\pgfqpoint{4.488700in}{2.353336in}}%
\pgfpathquadraticcurveto{\pgfqpoint{4.488700in}{2.366097in}}{\pgfqpoint{4.497689in}{2.373437in}}%
\pgfpathquadraticcurveto{\pgfqpoint{4.506678in}{2.380776in}}{\pgfqpoint{4.522449in}{2.380776in}}%
\pgfpathquadraticcurveto{\pgfqpoint{4.529232in}{2.380776in}}{\pgfqpoint{4.536241in}{2.379539in}}%
\pgfpathquadraticcurveto{\pgfqpoint{4.543271in}{2.378323in}}{\pgfqpoint{4.550611in}{2.375890in}}%
\pgfpathlineto{\pgfqpoint{4.550611in}{2.375890in}}%
\pgfpathclose%
\pgfpathmoveto{\pgfqpoint{4.587926in}{2.375499in}}%
\pgfpathlineto{\pgfqpoint{4.587926in}{2.355006in}}%
\pgfpathlineto{\pgfqpoint{4.612336in}{2.355006in}}%
\pgfpathlineto{\pgfqpoint{4.612336in}{2.345790in}}%
\pgfpathlineto{\pgfqpoint{4.587926in}{2.345790in}}%
\pgfpathlineto{\pgfqpoint{4.587926in}{2.306619in}}%
\pgfpathquadraticcurveto{\pgfqpoint{4.587926in}{2.297796in}}{\pgfqpoint{4.590339in}{2.295280in}}%
\pgfpathquadraticcurveto{\pgfqpoint{4.592751in}{2.292765in}}{\pgfqpoint{4.600173in}{2.292765in}}%
\pgfpathlineto{\pgfqpoint{4.612336in}{2.292765in}}%
\pgfpathlineto{\pgfqpoint{4.612336in}{2.282849in}}%
\pgfpathlineto{\pgfqpoint{4.600173in}{2.282849in}}%
\pgfpathquadraticcurveto{\pgfqpoint{4.586442in}{2.282849in}}{\pgfqpoint{4.581226in}{2.287962in}}%
\pgfpathquadraticcurveto{\pgfqpoint{4.576010in}{2.293095in}}{\pgfqpoint{4.576010in}{2.306619in}}%
\pgfpathlineto{\pgfqpoint{4.576010in}{2.345790in}}%
\pgfpathlineto{\pgfqpoint{4.567310in}{2.345790in}}%
\pgfpathlineto{\pgfqpoint{4.567310in}{2.355006in}}%
\pgfpathlineto{\pgfqpoint{4.576010in}{2.355006in}}%
\pgfpathlineto{\pgfqpoint{4.576010in}{2.375499in}}%
\pgfpathlineto{\pgfqpoint{4.587926in}{2.375499in}}%
\pgfpathlineto{\pgfqpoint{4.587926in}{2.375499in}}%
\pgfpathclose%
\pgfpathmoveto{\pgfqpoint{4.626706in}{2.311320in}}%
\pgfpathlineto{\pgfqpoint{4.626706in}{2.355006in}}%
\pgfpathlineto{\pgfqpoint{4.638560in}{2.355006in}}%
\pgfpathlineto{\pgfqpoint{4.638560in}{2.311773in}}%
\pgfpathquadraticcurveto{\pgfqpoint{4.638560in}{2.301527in}}{\pgfqpoint{4.642539in}{2.296394in}}%
\pgfpathquadraticcurveto{\pgfqpoint{4.646539in}{2.291281in}}{\pgfqpoint{4.654538in}{2.291281in}}%
\pgfpathquadraticcurveto{\pgfqpoint{4.664124in}{2.291281in}}{\pgfqpoint{4.669691in}{2.297404in}}%
\pgfpathquadraticcurveto{\pgfqpoint{4.675278in}{2.303527in}}{\pgfqpoint{4.675278in}{2.314103in}}%
\pgfpathlineto{\pgfqpoint{4.675278in}{2.355006in}}%
\pgfpathlineto{\pgfqpoint{4.687132in}{2.355006in}}%
\pgfpathlineto{\pgfqpoint{4.687132in}{2.282849in}}%
\pgfpathlineto{\pgfqpoint{4.675278in}{2.282849in}}%
\pgfpathlineto{\pgfqpoint{4.675278in}{2.293940in}}%
\pgfpathquadraticcurveto{\pgfqpoint{4.670969in}{2.287364in}}{\pgfqpoint{4.665258in}{2.284168in}}%
\pgfpathquadraticcurveto{\pgfqpoint{4.659568in}{2.280973in}}{\pgfqpoint{4.652023in}{2.280973in}}%
\pgfpathquadraticcurveto{\pgfqpoint{4.639591in}{2.280973in}}{\pgfqpoint{4.633138in}{2.288704in}}%
\pgfpathquadraticcurveto{\pgfqpoint{4.626706in}{2.296435in}}{\pgfqpoint{4.626706in}{2.311320in}}%
\pgfpathlineto{\pgfqpoint{4.626706in}{2.311320in}}%
\pgfpathclose%
\pgfpathmoveto{\pgfqpoint{4.656538in}{2.356738in}}%
\pgfpathlineto{\pgfqpoint{4.656538in}{2.356738in}}%
\pgfpathlineto{\pgfqpoint{4.656538in}{2.356738in}}%
\pgfpathclose%
\pgfpathmoveto{\pgfqpoint{4.759021in}{2.344059in}}%
\pgfpathlineto{\pgfqpoint{4.759021in}{2.383106in}}%
\pgfpathlineto{\pgfqpoint{4.770876in}{2.383106in}}%
\pgfpathlineto{\pgfqpoint{4.770876in}{2.282849in}}%
\pgfpathlineto{\pgfqpoint{4.759021in}{2.282849in}}%
\pgfpathlineto{\pgfqpoint{4.759021in}{2.293672in}}%
\pgfpathquadraticcurveto{\pgfqpoint{4.755290in}{2.287240in}}{\pgfqpoint{4.749579in}{2.284106in}}%
\pgfpathquadraticcurveto{\pgfqpoint{4.743889in}{2.280973in}}{\pgfqpoint{4.735890in}{2.280973in}}%
\pgfpathquadraticcurveto{\pgfqpoint{4.722819in}{2.280973in}}{\pgfqpoint{4.714593in}{2.291405in}}%
\pgfpathquadraticcurveto{\pgfqpoint{4.706388in}{2.301857in}}{\pgfqpoint{4.706388in}{2.318866in}}%
\pgfpathquadraticcurveto{\pgfqpoint{4.706388in}{2.335874in}}{\pgfqpoint{4.714593in}{2.346306in}}%
\pgfpathquadraticcurveto{\pgfqpoint{4.722819in}{2.356738in}}{\pgfqpoint{4.735890in}{2.356738in}}%
\pgfpathquadraticcurveto{\pgfqpoint{4.743889in}{2.356738in}}{\pgfqpoint{4.749579in}{2.353604in}}%
\pgfpathquadraticcurveto{\pgfqpoint{4.755290in}{2.350491in}}{\pgfqpoint{4.759021in}{2.344059in}}%
\pgfpathlineto{\pgfqpoint{4.759021in}{2.344059in}}%
\pgfpathclose%
\pgfpathmoveto{\pgfqpoint{4.718634in}{2.318866in}}%
\pgfpathquadraticcurveto{\pgfqpoint{4.718634in}{2.305795in}}{\pgfqpoint{4.724015in}{2.298352in}}%
\pgfpathquadraticcurveto{\pgfqpoint{4.729396in}{2.290910in}}{\pgfqpoint{4.738797in}{2.290910in}}%
\pgfpathquadraticcurveto{\pgfqpoint{4.748198in}{2.290910in}}{\pgfqpoint{4.753599in}{2.298352in}}%
\pgfpathquadraticcurveto{\pgfqpoint{4.759021in}{2.305795in}}{\pgfqpoint{4.759021in}{2.318866in}}%
\pgfpathquadraticcurveto{\pgfqpoint{4.759021in}{2.331936in}}{\pgfqpoint{4.753599in}{2.339379in}}%
\pgfpathquadraticcurveto{\pgfqpoint{4.748198in}{2.346821in}}{\pgfqpoint{4.738797in}{2.346821in}}%
\pgfpathquadraticcurveto{\pgfqpoint{4.729396in}{2.346821in}}{\pgfqpoint{4.724015in}{2.339379in}}%
\pgfpathquadraticcurveto{\pgfqpoint{4.718634in}{2.331936in}}{\pgfqpoint{4.718634in}{2.318866in}}%
\pgfpathlineto{\pgfqpoint{4.718634in}{2.318866in}}%
\pgfpathclose%
\pgfpathmoveto{\pgfqpoint{4.825323in}{2.276148in}}%
\pgfpathquadraticcurveto{\pgfqpoint{4.820314in}{2.263263in}}{\pgfqpoint{4.815531in}{2.259346in}}%
\pgfpathquadraticcurveto{\pgfqpoint{4.810768in}{2.255408in}}{\pgfqpoint{4.802790in}{2.255408in}}%
\pgfpathlineto{\pgfqpoint{4.793306in}{2.255408in}}%
\pgfpathlineto{\pgfqpoint{4.793306in}{2.265325in}}%
\pgfpathlineto{\pgfqpoint{4.800275in}{2.265325in}}%
\pgfpathquadraticcurveto{\pgfqpoint{4.805161in}{2.265325in}}{\pgfqpoint{4.807861in}{2.267655in}}%
\pgfpathquadraticcurveto{\pgfqpoint{4.810583in}{2.269964in}}{\pgfqpoint{4.813861in}{2.278602in}}%
\pgfpathlineto{\pgfqpoint{4.815984in}{2.284003in}}%
\pgfpathlineto{\pgfqpoint{4.786812in}{2.355006in}}%
\pgfpathlineto{\pgfqpoint{4.799367in}{2.355006in}}%
\pgfpathlineto{\pgfqpoint{4.821922in}{2.298579in}}%
\pgfpathlineto{\pgfqpoint{4.844476in}{2.355006in}}%
\pgfpathlineto{\pgfqpoint{4.857031in}{2.355006in}}%
\pgfpathlineto{\pgfqpoint{4.825323in}{2.276148in}}%
\pgfpathlineto{\pgfqpoint{4.825323in}{2.276148in}}%
\pgfpathclose%
\pgfpathmoveto{\pgfqpoint{4.917396in}{2.284849in}}%
\pgfpathlineto{\pgfqpoint{4.917396in}{2.296703in}}%
\pgfpathquadraticcurveto{\pgfqpoint{4.922303in}{2.294373in}}{\pgfqpoint{4.927312in}{2.293157in}}%
\pgfpathquadraticcurveto{\pgfqpoint{4.932343in}{2.291941in}}{\pgfqpoint{4.937188in}{2.291941in}}%
\pgfpathquadraticcurveto{\pgfqpoint{4.950073in}{2.291941in}}{\pgfqpoint{4.956855in}{2.300599in}}%
\pgfpathquadraticcurveto{\pgfqpoint{4.963659in}{2.309258in}}{\pgfqpoint{4.964628in}{2.326926in}}%
\pgfpathquadraticcurveto{\pgfqpoint{4.960896in}{2.321381in}}{\pgfqpoint{4.955144in}{2.318412in}}%
\pgfpathquadraticcurveto{\pgfqpoint{4.949413in}{2.315443in}}{\pgfqpoint{4.942465in}{2.315443in}}%
\pgfpathquadraticcurveto{\pgfqpoint{4.928034in}{2.315443in}}{\pgfqpoint{4.919622in}{2.324164in}}%
\pgfpathquadraticcurveto{\pgfqpoint{4.911211in}{2.332905in}}{\pgfqpoint{4.911211in}{2.348058in}}%
\pgfpathquadraticcurveto{\pgfqpoint{4.911211in}{2.362861in}}{\pgfqpoint{4.919973in}{2.371808in}}%
\pgfpathquadraticcurveto{\pgfqpoint{4.928735in}{2.380776in}}{\pgfqpoint{4.943290in}{2.380776in}}%
\pgfpathquadraticcurveto{\pgfqpoint{4.959989in}{2.380776in}}{\pgfqpoint{4.968772in}{2.367974in}}%
\pgfpathquadraticcurveto{\pgfqpoint{4.977575in}{2.355191in}}{\pgfqpoint{4.977575in}{2.330844in}}%
\pgfpathquadraticcurveto{\pgfqpoint{4.977575in}{2.308104in}}{\pgfqpoint{4.966772in}{2.294538in}}%
\pgfpathquadraticcurveto{\pgfqpoint{4.955990in}{2.280973in}}{\pgfqpoint{4.937765in}{2.280973in}}%
\pgfpathquadraticcurveto{\pgfqpoint{4.932858in}{2.280973in}}{\pgfqpoint{4.927828in}{2.281942in}}%
\pgfpathquadraticcurveto{\pgfqpoint{4.922818in}{2.282911in}}{\pgfqpoint{4.917396in}{2.284849in}}%
\pgfpathlineto{\pgfqpoint{4.917396in}{2.284849in}}%
\pgfpathclose%
\pgfpathmoveto{\pgfqpoint{4.943290in}{2.325628in}}%
\pgfpathquadraticcurveto{\pgfqpoint{4.952052in}{2.325628in}}{\pgfqpoint{4.957165in}{2.331606in}}%
\pgfpathquadraticcurveto{\pgfqpoint{4.962298in}{2.337606in}}{\pgfqpoint{4.962298in}{2.348058in}}%
\pgfpathquadraticcurveto{\pgfqpoint{4.962298in}{2.358428in}}{\pgfqpoint{4.957165in}{2.364448in}}%
\pgfpathquadraticcurveto{\pgfqpoint{4.952052in}{2.370468in}}{\pgfqpoint{4.943290in}{2.370468in}}%
\pgfpathquadraticcurveto{\pgfqpoint{4.934528in}{2.370468in}}{\pgfqpoint{4.929415in}{2.364448in}}%
\pgfpathquadraticcurveto{\pgfqpoint{4.924302in}{2.358428in}}{\pgfqpoint{4.924302in}{2.348058in}}%
\pgfpathquadraticcurveto{\pgfqpoint{4.924302in}{2.337606in}}{\pgfqpoint{4.929415in}{2.331606in}}%
\pgfpathquadraticcurveto{\pgfqpoint{4.934528in}{2.325628in}}{\pgfqpoint{4.943290in}{2.325628in}}%
\pgfpathlineto{\pgfqpoint{4.943290in}{2.325628in}}%
\pgfpathclose%
\pgfpathmoveto{\pgfqpoint{5.069689in}{2.382962in}}%
\pgfpathquadraticcurveto{\pgfqpoint{5.061071in}{2.368159in}}{\pgfqpoint{5.056865in}{2.353645in}}%
\pgfpathquadraticcurveto{\pgfqpoint{5.052680in}{2.339152in}}{\pgfqpoint{5.052680in}{2.324267in}}%
\pgfpathquadraticcurveto{\pgfqpoint{5.052680in}{2.309403in}}{\pgfqpoint{5.056906in}{2.294806in}}%
\pgfpathquadraticcurveto{\pgfqpoint{5.061133in}{2.280210in}}{\pgfqpoint{5.069689in}{2.265449in}}%
\pgfpathlineto{\pgfqpoint{5.059380in}{2.265449in}}%
\pgfpathquadraticcurveto{\pgfqpoint{5.049732in}{2.280602in}}{\pgfqpoint{5.044928in}{2.295219in}}%
\pgfpathquadraticcurveto{\pgfqpoint{5.040125in}{2.309836in}}{\pgfqpoint{5.040125in}{2.324267in}}%
\pgfpathquadraticcurveto{\pgfqpoint{5.040125in}{2.338637in}}{\pgfqpoint{5.044887in}{2.353192in}}%
\pgfpathquadraticcurveto{\pgfqpoint{5.049670in}{2.367767in}}{\pgfqpoint{5.059380in}{2.382962in}}%
\pgfpathlineto{\pgfqpoint{5.069689in}{2.382962in}}%
\pgfpathlineto{\pgfqpoint{5.069689in}{2.382962in}}%
\pgfpathclose%
\pgfpathmoveto{\pgfqpoint{5.105582in}{2.293796in}}%
\pgfpathlineto{\pgfqpoint{5.150999in}{2.293796in}}%
\pgfpathlineto{\pgfqpoint{5.150999in}{2.282849in}}%
\pgfpathlineto{\pgfqpoint{5.089934in}{2.282849in}}%
\pgfpathlineto{\pgfqpoint{5.089934in}{2.293796in}}%
\pgfpathquadraticcurveto{\pgfqpoint{5.097335in}{2.301465in}}{\pgfqpoint{5.110117in}{2.314371in}}%
\pgfpathquadraticcurveto{\pgfqpoint{5.122920in}{2.327298in}}{\pgfqpoint{5.126198in}{2.331050in}}%
\pgfpathquadraticcurveto{\pgfqpoint{5.132445in}{2.338059in}}{\pgfqpoint{5.134919in}{2.342925in}}%
\pgfpathquadraticcurveto{\pgfqpoint{5.137413in}{2.347790in}}{\pgfqpoint{5.137413in}{2.352491in}}%
\pgfpathquadraticcurveto{\pgfqpoint{5.137413in}{2.360160in}}{\pgfqpoint{5.132032in}{2.364984in}}%
\pgfpathquadraticcurveto{\pgfqpoint{5.126651in}{2.369829in}}{\pgfqpoint{5.118013in}{2.369829in}}%
\pgfpathquadraticcurveto{\pgfqpoint{5.111890in}{2.369829in}}{\pgfqpoint{5.105087in}{2.367706in}}%
\pgfpathquadraticcurveto{\pgfqpoint{5.098304in}{2.365582in}}{\pgfqpoint{5.090573in}{2.361253in}}%
\pgfpathlineto{\pgfqpoint{5.090573in}{2.374406in}}%
\pgfpathquadraticcurveto{\pgfqpoint{5.098428in}{2.377560in}}{\pgfqpoint{5.105252in}{2.379168in}}%
\pgfpathquadraticcurveto{\pgfqpoint{5.112096in}{2.380776in}}{\pgfqpoint{5.117766in}{2.380776in}}%
\pgfpathquadraticcurveto{\pgfqpoint{5.132713in}{2.380776in}}{\pgfqpoint{5.141598in}{2.373293in}}%
\pgfpathquadraticcurveto{\pgfqpoint{5.150484in}{2.365829in}}{\pgfqpoint{5.150484in}{2.353336in}}%
\pgfpathquadraticcurveto{\pgfqpoint{5.150484in}{2.347398in}}{\pgfqpoint{5.148257in}{2.342079in}}%
\pgfpathquadraticcurveto{\pgfqpoint{5.146051in}{2.336781in}}{\pgfqpoint{5.140176in}{2.329565in}}%
\pgfpathquadraticcurveto{\pgfqpoint{5.138568in}{2.327689in}}{\pgfqpoint{5.129929in}{2.318762in}}%
\pgfpathquadraticcurveto{\pgfqpoint{5.121312in}{2.309836in}}{\pgfqpoint{5.105582in}{2.293796in}}%
\pgfpathlineto{\pgfqpoint{5.105582in}{2.293796in}}%
\pgfpathclose%
\pgfpathmoveto{\pgfqpoint{5.206148in}{2.370468in}}%
\pgfpathquadraticcurveto{\pgfqpoint{5.196108in}{2.370468in}}{\pgfqpoint{5.191036in}{2.360572in}}%
\pgfpathquadraticcurveto{\pgfqpoint{5.185985in}{2.350697in}}{\pgfqpoint{5.185985in}{2.330844in}}%
\pgfpathquadraticcurveto{\pgfqpoint{5.185985in}{2.311073in}}{\pgfqpoint{5.191036in}{2.301177in}}%
\pgfpathquadraticcurveto{\pgfqpoint{5.196108in}{2.291281in}}{\pgfqpoint{5.206148in}{2.291281in}}%
\pgfpathquadraticcurveto{\pgfqpoint{5.216271in}{2.291281in}}{\pgfqpoint{5.221322in}{2.301177in}}%
\pgfpathquadraticcurveto{\pgfqpoint{5.226393in}{2.311073in}}{\pgfqpoint{5.226393in}{2.330844in}}%
\pgfpathquadraticcurveto{\pgfqpoint{5.226393in}{2.350697in}}{\pgfqpoint{5.221322in}{2.360572in}}%
\pgfpathquadraticcurveto{\pgfqpoint{5.216271in}{2.370468in}}{\pgfqpoint{5.206148in}{2.370468in}}%
\pgfpathlineto{\pgfqpoint{5.206148in}{2.370468in}}%
\pgfpathclose%
\pgfpathmoveto{\pgfqpoint{5.206148in}{2.380776in}}%
\pgfpathquadraticcurveto{\pgfqpoint{5.222332in}{2.380776in}}{\pgfqpoint{5.230867in}{2.367974in}}%
\pgfpathquadraticcurveto{\pgfqpoint{5.239402in}{2.355191in}}{\pgfqpoint{5.239402in}{2.330844in}}%
\pgfpathquadraticcurveto{\pgfqpoint{5.239402in}{2.306558in}}{\pgfqpoint{5.230867in}{2.293755in}}%
\pgfpathquadraticcurveto{\pgfqpoint{5.222332in}{2.280973in}}{\pgfqpoint{5.206148in}{2.280973in}}%
\pgfpathquadraticcurveto{\pgfqpoint{5.189985in}{2.280973in}}{\pgfqpoint{5.181450in}{2.293755in}}%
\pgfpathquadraticcurveto{\pgfqpoint{5.172914in}{2.306558in}}{\pgfqpoint{5.172914in}{2.330844in}}%
\pgfpathquadraticcurveto{\pgfqpoint{5.172914in}{2.355191in}}{\pgfqpoint{5.181450in}{2.367974in}}%
\pgfpathquadraticcurveto{\pgfqpoint{5.189985in}{2.380776in}}{\pgfqpoint{5.206148in}{2.380776in}}%
\pgfpathlineto{\pgfqpoint{5.206148in}{2.380776in}}%
\pgfpathclose%
\pgfpathmoveto{\pgfqpoint{5.273481in}{2.293796in}}%
\pgfpathlineto{\pgfqpoint{5.318899in}{2.293796in}}%
\pgfpathlineto{\pgfqpoint{5.318899in}{2.282849in}}%
\pgfpathlineto{\pgfqpoint{5.257833in}{2.282849in}}%
\pgfpathlineto{\pgfqpoint{5.257833in}{2.293796in}}%
\pgfpathquadraticcurveto{\pgfqpoint{5.265234in}{2.301465in}}{\pgfqpoint{5.278016in}{2.314371in}}%
\pgfpathquadraticcurveto{\pgfqpoint{5.290819in}{2.327298in}}{\pgfqpoint{5.294097in}{2.331050in}}%
\pgfpathquadraticcurveto{\pgfqpoint{5.300344in}{2.338059in}}{\pgfqpoint{5.302818in}{2.342925in}}%
\pgfpathquadraticcurveto{\pgfqpoint{5.305313in}{2.347790in}}{\pgfqpoint{5.305313in}{2.352491in}}%
\pgfpathquadraticcurveto{\pgfqpoint{5.305313in}{2.360160in}}{\pgfqpoint{5.299932in}{2.364984in}}%
\pgfpathquadraticcurveto{\pgfqpoint{5.294551in}{2.369829in}}{\pgfqpoint{5.285913in}{2.369829in}}%
\pgfpathquadraticcurveto{\pgfqpoint{5.279789in}{2.369829in}}{\pgfqpoint{5.272986in}{2.367706in}}%
\pgfpathquadraticcurveto{\pgfqpoint{5.266203in}{2.365582in}}{\pgfqpoint{5.258472in}{2.361253in}}%
\pgfpathlineto{\pgfqpoint{5.258472in}{2.374406in}}%
\pgfpathquadraticcurveto{\pgfqpoint{5.266327in}{2.377560in}}{\pgfqpoint{5.273151in}{2.379168in}}%
\pgfpathquadraticcurveto{\pgfqpoint{5.279996in}{2.380776in}}{\pgfqpoint{5.285665in}{2.380776in}}%
\pgfpathquadraticcurveto{\pgfqpoint{5.300612in}{2.380776in}}{\pgfqpoint{5.309498in}{2.373293in}}%
\pgfpathquadraticcurveto{\pgfqpoint{5.318383in}{2.365829in}}{\pgfqpoint{5.318383in}{2.353336in}}%
\pgfpathquadraticcurveto{\pgfqpoint{5.318383in}{2.347398in}}{\pgfqpoint{5.316157in}{2.342079in}}%
\pgfpathquadraticcurveto{\pgfqpoint{5.313951in}{2.336781in}}{\pgfqpoint{5.308075in}{2.329565in}}%
\pgfpathquadraticcurveto{\pgfqpoint{5.306467in}{2.327689in}}{\pgfqpoint{5.297829in}{2.318762in}}%
\pgfpathquadraticcurveto{\pgfqpoint{5.289211in}{2.309836in}}{\pgfqpoint{5.273481in}{2.293796in}}%
\pgfpathlineto{\pgfqpoint{5.273481in}{2.293796in}}%
\pgfpathclose%
\pgfpathmoveto{\pgfqpoint{5.374047in}{2.370468in}}%
\pgfpathquadraticcurveto{\pgfqpoint{5.364007in}{2.370468in}}{\pgfqpoint{5.358936in}{2.360572in}}%
\pgfpathquadraticcurveto{\pgfqpoint{5.353885in}{2.350697in}}{\pgfqpoint{5.353885in}{2.330844in}}%
\pgfpathquadraticcurveto{\pgfqpoint{5.353885in}{2.311073in}}{\pgfqpoint{5.358936in}{2.301177in}}%
\pgfpathquadraticcurveto{\pgfqpoint{5.364007in}{2.291281in}}{\pgfqpoint{5.374047in}{2.291281in}}%
\pgfpathquadraticcurveto{\pgfqpoint{5.384170in}{2.291281in}}{\pgfqpoint{5.389221in}{2.301177in}}%
\pgfpathquadraticcurveto{\pgfqpoint{5.394293in}{2.311073in}}{\pgfqpoint{5.394293in}{2.330844in}}%
\pgfpathquadraticcurveto{\pgfqpoint{5.394293in}{2.350697in}}{\pgfqpoint{5.389221in}{2.360572in}}%
\pgfpathquadraticcurveto{\pgfqpoint{5.384170in}{2.370468in}}{\pgfqpoint{5.374047in}{2.370468in}}%
\pgfpathlineto{\pgfqpoint{5.374047in}{2.370468in}}%
\pgfpathclose%
\pgfpathmoveto{\pgfqpoint{5.374047in}{2.380776in}}%
\pgfpathquadraticcurveto{\pgfqpoint{5.390231in}{2.380776in}}{\pgfqpoint{5.398766in}{2.367974in}}%
\pgfpathquadraticcurveto{\pgfqpoint{5.407301in}{2.355191in}}{\pgfqpoint{5.407301in}{2.330844in}}%
\pgfpathquadraticcurveto{\pgfqpoint{5.407301in}{2.306558in}}{\pgfqpoint{5.398766in}{2.293755in}}%
\pgfpathquadraticcurveto{\pgfqpoint{5.390231in}{2.280973in}}{\pgfqpoint{5.374047in}{2.280973in}}%
\pgfpathquadraticcurveto{\pgfqpoint{5.357884in}{2.280973in}}{\pgfqpoint{5.349349in}{2.293755in}}%
\pgfpathquadraticcurveto{\pgfqpoint{5.340814in}{2.306558in}}{\pgfqpoint{5.340814in}{2.330844in}}%
\pgfpathquadraticcurveto{\pgfqpoint{5.340814in}{2.355191in}}{\pgfqpoint{5.349349in}{2.367974in}}%
\pgfpathquadraticcurveto{\pgfqpoint{5.357884in}{2.380776in}}{\pgfqpoint{5.374047in}{2.380776in}}%
\pgfpathlineto{\pgfqpoint{5.374047in}{2.380776in}}%
\pgfpathclose%
\pgfpathmoveto{\pgfqpoint{5.426640in}{2.382962in}}%
\pgfpathlineto{\pgfqpoint{5.436948in}{2.382962in}}%
\pgfpathquadraticcurveto{\pgfqpoint{5.446596in}{2.367767in}}{\pgfqpoint{5.451400in}{2.353192in}}%
\pgfpathquadraticcurveto{\pgfqpoint{5.456203in}{2.338637in}}{\pgfqpoint{5.456203in}{2.324267in}}%
\pgfpathquadraticcurveto{\pgfqpoint{5.456203in}{2.309836in}}{\pgfqpoint{5.451400in}{2.295219in}}%
\pgfpathquadraticcurveto{\pgfqpoint{5.446596in}{2.280602in}}{\pgfqpoint{5.436948in}{2.265449in}}%
\pgfpathlineto{\pgfqpoint{5.426640in}{2.265449in}}%
\pgfpathquadraticcurveto{\pgfqpoint{5.435195in}{2.280210in}}{\pgfqpoint{5.439422in}{2.294806in}}%
\pgfpathquadraticcurveto{\pgfqpoint{5.443648in}{2.309403in}}{\pgfqpoint{5.443648in}{2.324267in}}%
\pgfpathquadraticcurveto{\pgfqpoint{5.443648in}{2.339152in}}{\pgfqpoint{5.439422in}{2.353645in}}%
\pgfpathquadraticcurveto{\pgfqpoint{5.435195in}{2.368159in}}{\pgfqpoint{5.426640in}{2.382962in}}%
\pgfpathlineto{\pgfqpoint{5.426640in}{2.382962in}}%
\pgfpathclose%
\pgfusepath{fill}%
\end{pgfscope}%
\begin{pgfscope}%
\pgfsetbuttcap%
\pgfsetroundjoin%
\definecolor{currentfill}{rgb}{0.090196,0.168627,0.227451}%
\pgfsetfillcolor{currentfill}%
\pgfsetlinewidth{0.000000pt}%
\definecolor{currentstroke}{rgb}{0.090196,0.168627,0.227451}%
\pgfsetstrokecolor{currentstroke}%
\pgfsetdash{}{0pt}%
\pgfpathmoveto{\pgfqpoint{4.550611in}{1.976290in}}%
\pgfpathlineto{\pgfqpoint{4.550611in}{1.963591in}}%
\pgfpathquadraticcurveto{\pgfqpoint{4.543210in}{1.967137in}}{\pgfqpoint{4.536633in}{1.968868in}}%
\pgfpathquadraticcurveto{\pgfqpoint{4.530056in}{1.970621in}}{\pgfqpoint{4.523933in}{1.970621in}}%
\pgfpathquadraticcurveto{\pgfqpoint{4.513316in}{1.970621in}}{\pgfqpoint{4.507543in}{1.966497in}}%
\pgfpathquadraticcurveto{\pgfqpoint{4.501771in}{1.962374in}}{\pgfqpoint{4.501771in}{1.954767in}}%
\pgfpathquadraticcurveto{\pgfqpoint{4.501771in}{1.948376in}}{\pgfqpoint{4.505605in}{1.945118in}}%
\pgfpathquadraticcurveto{\pgfqpoint{4.509440in}{1.941882in}}{\pgfqpoint{4.520140in}{1.939882in}}%
\pgfpathlineto{\pgfqpoint{4.527995in}{1.938274in}}%
\pgfpathquadraticcurveto{\pgfqpoint{4.542550in}{1.935491in}}{\pgfqpoint{4.549477in}{1.928502in}}%
\pgfpathquadraticcurveto{\pgfqpoint{4.556404in}{1.921513in}}{\pgfqpoint{4.556404in}{1.909803in}}%
\pgfpathquadraticcurveto{\pgfqpoint{4.556404in}{1.895804in}}{\pgfqpoint{4.547024in}{1.888588in}}%
\pgfpathquadraticcurveto{\pgfqpoint{4.537664in}{1.881373in}}{\pgfqpoint{4.519563in}{1.881373in}}%
\pgfpathquadraticcurveto{\pgfqpoint{4.512739in}{1.881373in}}{\pgfqpoint{4.505028in}{1.882919in}}%
\pgfpathquadraticcurveto{\pgfqpoint{4.497338in}{1.884465in}}{\pgfqpoint{4.489092in}{1.887496in}}%
\pgfpathlineto{\pgfqpoint{4.489092in}{1.900896in}}%
\pgfpathquadraticcurveto{\pgfqpoint{4.497008in}{1.896464in}}{\pgfqpoint{4.504616in}{1.894196in}}%
\pgfpathquadraticcurveto{\pgfqpoint{4.512223in}{1.891949in}}{\pgfqpoint{4.519563in}{1.891949in}}%
\pgfpathquadraticcurveto{\pgfqpoint{4.530696in}{1.891949in}}{\pgfqpoint{4.536757in}{1.896320in}}%
\pgfpathquadraticcurveto{\pgfqpoint{4.542818in}{1.900711in}}{\pgfqpoint{4.542818in}{1.908834in}}%
\pgfpathquadraticcurveto{\pgfqpoint{4.542818in}{1.915905in}}{\pgfqpoint{4.538468in}{1.919905in}}%
\pgfpathquadraticcurveto{\pgfqpoint{4.534118in}{1.923904in}}{\pgfqpoint{4.524201in}{1.925904in}}%
\pgfpathlineto{\pgfqpoint{4.516264in}{1.927450in}}%
\pgfpathquadraticcurveto{\pgfqpoint{4.501709in}{1.930336in}}{\pgfqpoint{4.495194in}{1.936521in}}%
\pgfpathquadraticcurveto{\pgfqpoint{4.488700in}{1.942706in}}{\pgfqpoint{4.488700in}{1.953736in}}%
\pgfpathquadraticcurveto{\pgfqpoint{4.488700in}{1.966497in}}{\pgfqpoint{4.497689in}{1.973837in}}%
\pgfpathquadraticcurveto{\pgfqpoint{4.506678in}{1.981176in}}{\pgfqpoint{4.522449in}{1.981176in}}%
\pgfpathquadraticcurveto{\pgfqpoint{4.529232in}{1.981176in}}{\pgfqpoint{4.536241in}{1.979939in}}%
\pgfpathquadraticcurveto{\pgfqpoint{4.543271in}{1.978723in}}{\pgfqpoint{4.550611in}{1.976290in}}%
\pgfpathlineto{\pgfqpoint{4.550611in}{1.976290in}}%
\pgfpathclose%
\pgfpathmoveto{\pgfqpoint{4.587926in}{1.975899in}}%
\pgfpathlineto{\pgfqpoint{4.587926in}{1.955406in}}%
\pgfpathlineto{\pgfqpoint{4.612336in}{1.955406in}}%
\pgfpathlineto{\pgfqpoint{4.612336in}{1.946190in}}%
\pgfpathlineto{\pgfqpoint{4.587926in}{1.946190in}}%
\pgfpathlineto{\pgfqpoint{4.587926in}{1.907019in}}%
\pgfpathquadraticcurveto{\pgfqpoint{4.587926in}{1.898196in}}{\pgfqpoint{4.590339in}{1.895680in}}%
\pgfpathquadraticcurveto{\pgfqpoint{4.592751in}{1.893165in}}{\pgfqpoint{4.600173in}{1.893165in}}%
\pgfpathlineto{\pgfqpoint{4.612336in}{1.893165in}}%
\pgfpathlineto{\pgfqpoint{4.612336in}{1.883249in}}%
\pgfpathlineto{\pgfqpoint{4.600173in}{1.883249in}}%
\pgfpathquadraticcurveto{\pgfqpoint{4.586442in}{1.883249in}}{\pgfqpoint{4.581226in}{1.888362in}}%
\pgfpathquadraticcurveto{\pgfqpoint{4.576010in}{1.893495in}}{\pgfqpoint{4.576010in}{1.907019in}}%
\pgfpathlineto{\pgfqpoint{4.576010in}{1.946190in}}%
\pgfpathlineto{\pgfqpoint{4.567310in}{1.946190in}}%
\pgfpathlineto{\pgfqpoint{4.567310in}{1.955406in}}%
\pgfpathlineto{\pgfqpoint{4.576010in}{1.955406in}}%
\pgfpathlineto{\pgfqpoint{4.576010in}{1.975899in}}%
\pgfpathlineto{\pgfqpoint{4.587926in}{1.975899in}}%
\pgfpathlineto{\pgfqpoint{4.587926in}{1.975899in}}%
\pgfpathclose%
\pgfpathmoveto{\pgfqpoint{4.626706in}{1.911720in}}%
\pgfpathlineto{\pgfqpoint{4.626706in}{1.955406in}}%
\pgfpathlineto{\pgfqpoint{4.638560in}{1.955406in}}%
\pgfpathlineto{\pgfqpoint{4.638560in}{1.912173in}}%
\pgfpathquadraticcurveto{\pgfqpoint{4.638560in}{1.901927in}}{\pgfqpoint{4.642539in}{1.896794in}}%
\pgfpathquadraticcurveto{\pgfqpoint{4.646539in}{1.891681in}}{\pgfqpoint{4.654538in}{1.891681in}}%
\pgfpathquadraticcurveto{\pgfqpoint{4.664124in}{1.891681in}}{\pgfqpoint{4.669691in}{1.897804in}}%
\pgfpathquadraticcurveto{\pgfqpoint{4.675278in}{1.903927in}}{\pgfqpoint{4.675278in}{1.914503in}}%
\pgfpathlineto{\pgfqpoint{4.675278in}{1.955406in}}%
\pgfpathlineto{\pgfqpoint{4.687132in}{1.955406in}}%
\pgfpathlineto{\pgfqpoint{4.687132in}{1.883249in}}%
\pgfpathlineto{\pgfqpoint{4.675278in}{1.883249in}}%
\pgfpathlineto{\pgfqpoint{4.675278in}{1.894340in}}%
\pgfpathquadraticcurveto{\pgfqpoint{4.670969in}{1.887764in}}{\pgfqpoint{4.665258in}{1.884568in}}%
\pgfpathquadraticcurveto{\pgfqpoint{4.659568in}{1.881373in}}{\pgfqpoint{4.652023in}{1.881373in}}%
\pgfpathquadraticcurveto{\pgfqpoint{4.639591in}{1.881373in}}{\pgfqpoint{4.633138in}{1.889104in}}%
\pgfpathquadraticcurveto{\pgfqpoint{4.626706in}{1.896835in}}{\pgfqpoint{4.626706in}{1.911720in}}%
\pgfpathlineto{\pgfqpoint{4.626706in}{1.911720in}}%
\pgfpathclose%
\pgfpathmoveto{\pgfqpoint{4.656538in}{1.957138in}}%
\pgfpathlineto{\pgfqpoint{4.656538in}{1.957138in}}%
\pgfpathlineto{\pgfqpoint{4.656538in}{1.957138in}}%
\pgfpathclose%
\pgfpathmoveto{\pgfqpoint{4.759021in}{1.944459in}}%
\pgfpathlineto{\pgfqpoint{4.759021in}{1.983506in}}%
\pgfpathlineto{\pgfqpoint{4.770876in}{1.983506in}}%
\pgfpathlineto{\pgfqpoint{4.770876in}{1.883249in}}%
\pgfpathlineto{\pgfqpoint{4.759021in}{1.883249in}}%
\pgfpathlineto{\pgfqpoint{4.759021in}{1.894072in}}%
\pgfpathquadraticcurveto{\pgfqpoint{4.755290in}{1.887640in}}{\pgfqpoint{4.749579in}{1.884506in}}%
\pgfpathquadraticcurveto{\pgfqpoint{4.743889in}{1.881373in}}{\pgfqpoint{4.735890in}{1.881373in}}%
\pgfpathquadraticcurveto{\pgfqpoint{4.722819in}{1.881373in}}{\pgfqpoint{4.714593in}{1.891805in}}%
\pgfpathquadraticcurveto{\pgfqpoint{4.706388in}{1.902257in}}{\pgfqpoint{4.706388in}{1.919266in}}%
\pgfpathquadraticcurveto{\pgfqpoint{4.706388in}{1.936274in}}{\pgfqpoint{4.714593in}{1.946706in}}%
\pgfpathquadraticcurveto{\pgfqpoint{4.722819in}{1.957138in}}{\pgfqpoint{4.735890in}{1.957138in}}%
\pgfpathquadraticcurveto{\pgfqpoint{4.743889in}{1.957138in}}{\pgfqpoint{4.749579in}{1.954004in}}%
\pgfpathquadraticcurveto{\pgfqpoint{4.755290in}{1.950891in}}{\pgfqpoint{4.759021in}{1.944459in}}%
\pgfpathlineto{\pgfqpoint{4.759021in}{1.944459in}}%
\pgfpathclose%
\pgfpathmoveto{\pgfqpoint{4.718634in}{1.919266in}}%
\pgfpathquadraticcurveto{\pgfqpoint{4.718634in}{1.906195in}}{\pgfqpoint{4.724015in}{1.898752in}}%
\pgfpathquadraticcurveto{\pgfqpoint{4.729396in}{1.891310in}}{\pgfqpoint{4.738797in}{1.891310in}}%
\pgfpathquadraticcurveto{\pgfqpoint{4.748198in}{1.891310in}}{\pgfqpoint{4.753599in}{1.898752in}}%
\pgfpathquadraticcurveto{\pgfqpoint{4.759021in}{1.906195in}}{\pgfqpoint{4.759021in}{1.919266in}}%
\pgfpathquadraticcurveto{\pgfqpoint{4.759021in}{1.932336in}}{\pgfqpoint{4.753599in}{1.939779in}}%
\pgfpathquadraticcurveto{\pgfqpoint{4.748198in}{1.947221in}}{\pgfqpoint{4.738797in}{1.947221in}}%
\pgfpathquadraticcurveto{\pgfqpoint{4.729396in}{1.947221in}}{\pgfqpoint{4.724015in}{1.939779in}}%
\pgfpathquadraticcurveto{\pgfqpoint{4.718634in}{1.932336in}}{\pgfqpoint{4.718634in}{1.919266in}}%
\pgfpathlineto{\pgfqpoint{4.718634in}{1.919266in}}%
\pgfpathclose%
\pgfpathmoveto{\pgfqpoint{4.825323in}{1.876548in}}%
\pgfpathquadraticcurveto{\pgfqpoint{4.820314in}{1.863663in}}{\pgfqpoint{4.815531in}{1.859746in}}%
\pgfpathquadraticcurveto{\pgfqpoint{4.810768in}{1.855808in}}{\pgfqpoint{4.802790in}{1.855808in}}%
\pgfpathlineto{\pgfqpoint{4.793306in}{1.855808in}}%
\pgfpathlineto{\pgfqpoint{4.793306in}{1.865725in}}%
\pgfpathlineto{\pgfqpoint{4.800275in}{1.865725in}}%
\pgfpathquadraticcurveto{\pgfqpoint{4.805161in}{1.865725in}}{\pgfqpoint{4.807861in}{1.868055in}}%
\pgfpathquadraticcurveto{\pgfqpoint{4.810583in}{1.870364in}}{\pgfqpoint{4.813861in}{1.879002in}}%
\pgfpathlineto{\pgfqpoint{4.815984in}{1.884403in}}%
\pgfpathlineto{\pgfqpoint{4.786812in}{1.955406in}}%
\pgfpathlineto{\pgfqpoint{4.799367in}{1.955406in}}%
\pgfpathlineto{\pgfqpoint{4.821922in}{1.898979in}}%
\pgfpathlineto{\pgfqpoint{4.844476in}{1.955406in}}%
\pgfpathlineto{\pgfqpoint{4.857031in}{1.955406in}}%
\pgfpathlineto{\pgfqpoint{4.825323in}{1.876548in}}%
\pgfpathlineto{\pgfqpoint{4.825323in}{1.876548in}}%
\pgfpathclose%
\pgfpathmoveto{\pgfqpoint{4.919272in}{1.894196in}}%
\pgfpathlineto{\pgfqpoint{4.940527in}{1.894196in}}%
\pgfpathlineto{\pgfqpoint{4.940527in}{1.967590in}}%
\pgfpathlineto{\pgfqpoint{4.917396in}{1.962951in}}%
\pgfpathlineto{\pgfqpoint{4.917396in}{1.974806in}}%
\pgfpathlineto{\pgfqpoint{4.940404in}{1.979445in}}%
\pgfpathlineto{\pgfqpoint{4.953413in}{1.979445in}}%
\pgfpathlineto{\pgfqpoint{4.953413in}{1.894196in}}%
\pgfpathlineto{\pgfqpoint{4.974668in}{1.894196in}}%
\pgfpathlineto{\pgfqpoint{4.974668in}{1.883249in}}%
\pgfpathlineto{\pgfqpoint{4.919272in}{1.883249in}}%
\pgfpathlineto{\pgfqpoint{4.919272in}{1.894196in}}%
\pgfpathlineto{\pgfqpoint{4.919272in}{1.894196in}}%
\pgfpathclose%
\pgfpathmoveto{\pgfqpoint{5.028786in}{1.970868in}}%
\pgfpathquadraticcurveto{\pgfqpoint{5.018746in}{1.970868in}}{\pgfqpoint{5.013674in}{1.960972in}}%
\pgfpathquadraticcurveto{\pgfqpoint{5.008623in}{1.951097in}}{\pgfqpoint{5.008623in}{1.931244in}}%
\pgfpathquadraticcurveto{\pgfqpoint{5.008623in}{1.911473in}}{\pgfqpoint{5.013674in}{1.901577in}}%
\pgfpathquadraticcurveto{\pgfqpoint{5.018746in}{1.891681in}}{\pgfqpoint{5.028786in}{1.891681in}}%
\pgfpathquadraticcurveto{\pgfqpoint{5.038908in}{1.891681in}}{\pgfqpoint{5.043959in}{1.901577in}}%
\pgfpathquadraticcurveto{\pgfqpoint{5.049031in}{1.911473in}}{\pgfqpoint{5.049031in}{1.931244in}}%
\pgfpathquadraticcurveto{\pgfqpoint{5.049031in}{1.951097in}}{\pgfqpoint{5.043959in}{1.960972in}}%
\pgfpathquadraticcurveto{\pgfqpoint{5.038908in}{1.970868in}}{\pgfqpoint{5.028786in}{1.970868in}}%
\pgfpathlineto{\pgfqpoint{5.028786in}{1.970868in}}%
\pgfpathclose%
\pgfpathmoveto{\pgfqpoint{5.028786in}{1.981176in}}%
\pgfpathquadraticcurveto{\pgfqpoint{5.044970in}{1.981176in}}{\pgfqpoint{5.053505in}{1.968374in}}%
\pgfpathquadraticcurveto{\pgfqpoint{5.062040in}{1.955591in}}{\pgfqpoint{5.062040in}{1.931244in}}%
\pgfpathquadraticcurveto{\pgfqpoint{5.062040in}{1.906958in}}{\pgfqpoint{5.053505in}{1.894155in}}%
\pgfpathquadraticcurveto{\pgfqpoint{5.044970in}{1.881373in}}{\pgfqpoint{5.028786in}{1.881373in}}%
\pgfpathquadraticcurveto{\pgfqpoint{5.012623in}{1.881373in}}{\pgfqpoint{5.004087in}{1.894155in}}%
\pgfpathquadraticcurveto{\pgfqpoint{4.995552in}{1.906958in}}{\pgfqpoint{4.995552in}{1.931244in}}%
\pgfpathquadraticcurveto{\pgfqpoint{4.995552in}{1.955591in}}{\pgfqpoint{5.004087in}{1.968374in}}%
\pgfpathquadraticcurveto{\pgfqpoint{5.012623in}{1.981176in}}{\pgfqpoint{5.028786in}{1.981176in}}%
\pgfpathlineto{\pgfqpoint{5.028786in}{1.981176in}}%
\pgfpathclose%
\pgfpathmoveto{\pgfqpoint{5.153638in}{1.983362in}}%
\pgfpathquadraticcurveto{\pgfqpoint{5.145021in}{1.968559in}}{\pgfqpoint{5.140815in}{1.954045in}}%
\pgfpathquadraticcurveto{\pgfqpoint{5.136630in}{1.939552in}}{\pgfqpoint{5.136630in}{1.924667in}}%
\pgfpathquadraticcurveto{\pgfqpoint{5.136630in}{1.909803in}}{\pgfqpoint{5.140856in}{1.895206in}}%
\pgfpathquadraticcurveto{\pgfqpoint{5.145082in}{1.880610in}}{\pgfqpoint{5.153638in}{1.865849in}}%
\pgfpathlineto{\pgfqpoint{5.143330in}{1.865849in}}%
\pgfpathquadraticcurveto{\pgfqpoint{5.133682in}{1.881002in}}{\pgfqpoint{5.128878in}{1.895619in}}%
\pgfpathquadraticcurveto{\pgfqpoint{5.124074in}{1.910236in}}{\pgfqpoint{5.124074in}{1.924667in}}%
\pgfpathquadraticcurveto{\pgfqpoint{5.124074in}{1.939037in}}{\pgfqpoint{5.128837in}{1.953592in}}%
\pgfpathquadraticcurveto{\pgfqpoint{5.133620in}{1.968167in}}{\pgfqpoint{5.143330in}{1.983362in}}%
\pgfpathlineto{\pgfqpoint{5.153638in}{1.983362in}}%
\pgfpathlineto{\pgfqpoint{5.153638in}{1.983362in}}%
\pgfpathclose%
\pgfpathmoveto{\pgfqpoint{5.189531in}{1.894196in}}%
\pgfpathlineto{\pgfqpoint{5.234949in}{1.894196in}}%
\pgfpathlineto{\pgfqpoint{5.234949in}{1.883249in}}%
\pgfpathlineto{\pgfqpoint{5.173883in}{1.883249in}}%
\pgfpathlineto{\pgfqpoint{5.173883in}{1.894196in}}%
\pgfpathquadraticcurveto{\pgfqpoint{5.181285in}{1.901865in}}{\pgfqpoint{5.194067in}{1.914771in}}%
\pgfpathquadraticcurveto{\pgfqpoint{5.206870in}{1.927698in}}{\pgfqpoint{5.210148in}{1.931450in}}%
\pgfpathquadraticcurveto{\pgfqpoint{5.216394in}{1.938459in}}{\pgfqpoint{5.218868in}{1.943325in}}%
\pgfpathquadraticcurveto{\pgfqpoint{5.221363in}{1.948190in}}{\pgfqpoint{5.221363in}{1.952891in}}%
\pgfpathquadraticcurveto{\pgfqpoint{5.221363in}{1.960560in}}{\pgfqpoint{5.215982in}{1.965384in}}%
\pgfpathquadraticcurveto{\pgfqpoint{5.210601in}{1.970229in}}{\pgfqpoint{5.201963in}{1.970229in}}%
\pgfpathquadraticcurveto{\pgfqpoint{5.195840in}{1.970229in}}{\pgfqpoint{5.189036in}{1.968106in}}%
\pgfpathquadraticcurveto{\pgfqpoint{5.182254in}{1.965982in}}{\pgfqpoint{5.174523in}{1.961653in}}%
\pgfpathlineto{\pgfqpoint{5.174523in}{1.974806in}}%
\pgfpathquadraticcurveto{\pgfqpoint{5.182377in}{1.977960in}}{\pgfqpoint{5.189201in}{1.979568in}}%
\pgfpathquadraticcurveto{\pgfqpoint{5.196046in}{1.981176in}}{\pgfqpoint{5.201715in}{1.981176in}}%
\pgfpathquadraticcurveto{\pgfqpoint{5.216662in}{1.981176in}}{\pgfqpoint{5.225548in}{1.973693in}}%
\pgfpathquadraticcurveto{\pgfqpoint{5.234434in}{1.966229in}}{\pgfqpoint{5.234434in}{1.953736in}}%
\pgfpathquadraticcurveto{\pgfqpoint{5.234434in}{1.947798in}}{\pgfqpoint{5.232207in}{1.942479in}}%
\pgfpathquadraticcurveto{\pgfqpoint{5.230001in}{1.937181in}}{\pgfqpoint{5.224125in}{1.929965in}}%
\pgfpathquadraticcurveto{\pgfqpoint{5.222517in}{1.928089in}}{\pgfqpoint{5.213879in}{1.919162in}}%
\pgfpathquadraticcurveto{\pgfqpoint{5.205262in}{1.910236in}}{\pgfqpoint{5.189531in}{1.894196in}}%
\pgfpathlineto{\pgfqpoint{5.189531in}{1.894196in}}%
\pgfpathclose%
\pgfpathmoveto{\pgfqpoint{5.290098in}{1.970868in}}%
\pgfpathquadraticcurveto{\pgfqpoint{5.280058in}{1.970868in}}{\pgfqpoint{5.274986in}{1.960972in}}%
\pgfpathquadraticcurveto{\pgfqpoint{5.269935in}{1.951097in}}{\pgfqpoint{5.269935in}{1.931244in}}%
\pgfpathquadraticcurveto{\pgfqpoint{5.269935in}{1.911473in}}{\pgfqpoint{5.274986in}{1.901577in}}%
\pgfpathquadraticcurveto{\pgfqpoint{5.280058in}{1.891681in}}{\pgfqpoint{5.290098in}{1.891681in}}%
\pgfpathquadraticcurveto{\pgfqpoint{5.300220in}{1.891681in}}{\pgfqpoint{5.305271in}{1.901577in}}%
\pgfpathquadraticcurveto{\pgfqpoint{5.310343in}{1.911473in}}{\pgfqpoint{5.310343in}{1.931244in}}%
\pgfpathquadraticcurveto{\pgfqpoint{5.310343in}{1.951097in}}{\pgfqpoint{5.305271in}{1.960972in}}%
\pgfpathquadraticcurveto{\pgfqpoint{5.300220in}{1.970868in}}{\pgfqpoint{5.290098in}{1.970868in}}%
\pgfpathlineto{\pgfqpoint{5.290098in}{1.970868in}}%
\pgfpathclose%
\pgfpathmoveto{\pgfqpoint{5.290098in}{1.981176in}}%
\pgfpathquadraticcurveto{\pgfqpoint{5.306281in}{1.981176in}}{\pgfqpoint{5.314817in}{1.968374in}}%
\pgfpathquadraticcurveto{\pgfqpoint{5.323352in}{1.955591in}}{\pgfqpoint{5.323352in}{1.931244in}}%
\pgfpathquadraticcurveto{\pgfqpoint{5.323352in}{1.906958in}}{\pgfqpoint{5.314817in}{1.894155in}}%
\pgfpathquadraticcurveto{\pgfqpoint{5.306281in}{1.881373in}}{\pgfqpoint{5.290098in}{1.881373in}}%
\pgfpathquadraticcurveto{\pgfqpoint{5.273934in}{1.881373in}}{\pgfqpoint{5.265399in}{1.894155in}}%
\pgfpathquadraticcurveto{\pgfqpoint{5.256864in}{1.906958in}}{\pgfqpoint{5.256864in}{1.931244in}}%
\pgfpathquadraticcurveto{\pgfqpoint{5.256864in}{1.955591in}}{\pgfqpoint{5.265399in}{1.968374in}}%
\pgfpathquadraticcurveto{\pgfqpoint{5.273934in}{1.981176in}}{\pgfqpoint{5.290098in}{1.981176in}}%
\pgfpathlineto{\pgfqpoint{5.290098in}{1.981176in}}%
\pgfpathclose%
\pgfpathmoveto{\pgfqpoint{5.357431in}{1.894196in}}%
\pgfpathlineto{\pgfqpoint{5.402848in}{1.894196in}}%
\pgfpathlineto{\pgfqpoint{5.402848in}{1.883249in}}%
\pgfpathlineto{\pgfqpoint{5.341783in}{1.883249in}}%
\pgfpathlineto{\pgfqpoint{5.341783in}{1.894196in}}%
\pgfpathquadraticcurveto{\pgfqpoint{5.349184in}{1.901865in}}{\pgfqpoint{5.361966in}{1.914771in}}%
\pgfpathquadraticcurveto{\pgfqpoint{5.374769in}{1.927698in}}{\pgfqpoint{5.378047in}{1.931450in}}%
\pgfpathquadraticcurveto{\pgfqpoint{5.384294in}{1.938459in}}{\pgfqpoint{5.386768in}{1.943325in}}%
\pgfpathquadraticcurveto{\pgfqpoint{5.389262in}{1.948190in}}{\pgfqpoint{5.389262in}{1.952891in}}%
\pgfpathquadraticcurveto{\pgfqpoint{5.389262in}{1.960560in}}{\pgfqpoint{5.383881in}{1.965384in}}%
\pgfpathquadraticcurveto{\pgfqpoint{5.378500in}{1.970229in}}{\pgfqpoint{5.369862in}{1.970229in}}%
\pgfpathquadraticcurveto{\pgfqpoint{5.363739in}{1.970229in}}{\pgfqpoint{5.356936in}{1.968106in}}%
\pgfpathquadraticcurveto{\pgfqpoint{5.350153in}{1.965982in}}{\pgfqpoint{5.342422in}{1.961653in}}%
\pgfpathlineto{\pgfqpoint{5.342422in}{1.974806in}}%
\pgfpathquadraticcurveto{\pgfqpoint{5.350277in}{1.977960in}}{\pgfqpoint{5.357101in}{1.979568in}}%
\pgfpathquadraticcurveto{\pgfqpoint{5.363945in}{1.981176in}}{\pgfqpoint{5.369615in}{1.981176in}}%
\pgfpathquadraticcurveto{\pgfqpoint{5.384562in}{1.981176in}}{\pgfqpoint{5.393447in}{1.973693in}}%
\pgfpathquadraticcurveto{\pgfqpoint{5.402333in}{1.966229in}}{\pgfqpoint{5.402333in}{1.953736in}}%
\pgfpathquadraticcurveto{\pgfqpoint{5.402333in}{1.947798in}}{\pgfqpoint{5.400106in}{1.942479in}}%
\pgfpathquadraticcurveto{\pgfqpoint{5.397900in}{1.937181in}}{\pgfqpoint{5.392025in}{1.929965in}}%
\pgfpathquadraticcurveto{\pgfqpoint{5.390417in}{1.928089in}}{\pgfqpoint{5.381778in}{1.919162in}}%
\pgfpathquadraticcurveto{\pgfqpoint{5.373161in}{1.910236in}}{\pgfqpoint{5.357431in}{1.894196in}}%
\pgfpathlineto{\pgfqpoint{5.357431in}{1.894196in}}%
\pgfpathclose%
\pgfpathmoveto{\pgfqpoint{5.432433in}{1.894196in}}%
\pgfpathlineto{\pgfqpoint{5.453688in}{1.894196in}}%
\pgfpathlineto{\pgfqpoint{5.453688in}{1.967590in}}%
\pgfpathlineto{\pgfqpoint{5.430557in}{1.962951in}}%
\pgfpathlineto{\pgfqpoint{5.430557in}{1.974806in}}%
\pgfpathlineto{\pgfqpoint{5.453564in}{1.979445in}}%
\pgfpathlineto{\pgfqpoint{5.466573in}{1.979445in}}%
\pgfpathlineto{\pgfqpoint{5.466573in}{1.894196in}}%
\pgfpathlineto{\pgfqpoint{5.487829in}{1.894196in}}%
\pgfpathlineto{\pgfqpoint{5.487829in}{1.883249in}}%
\pgfpathlineto{\pgfqpoint{5.432433in}{1.883249in}}%
\pgfpathlineto{\pgfqpoint{5.432433in}{1.894196in}}%
\pgfpathlineto{\pgfqpoint{5.432433in}{1.894196in}}%
\pgfpathclose%
\pgfpathmoveto{\pgfqpoint{5.510589in}{1.983362in}}%
\pgfpathlineto{\pgfqpoint{5.520897in}{1.983362in}}%
\pgfpathquadraticcurveto{\pgfqpoint{5.530546in}{1.968167in}}{\pgfqpoint{5.535349in}{1.953592in}}%
\pgfpathquadraticcurveto{\pgfqpoint{5.540153in}{1.939037in}}{\pgfqpoint{5.540153in}{1.924667in}}%
\pgfpathquadraticcurveto{\pgfqpoint{5.540153in}{1.910236in}}{\pgfqpoint{5.535349in}{1.895619in}}%
\pgfpathquadraticcurveto{\pgfqpoint{5.530546in}{1.881002in}}{\pgfqpoint{5.520897in}{1.865849in}}%
\pgfpathlineto{\pgfqpoint{5.510589in}{1.865849in}}%
\pgfpathquadraticcurveto{\pgfqpoint{5.519145in}{1.880610in}}{\pgfqpoint{5.523371in}{1.895206in}}%
\pgfpathquadraticcurveto{\pgfqpoint{5.527598in}{1.909803in}}{\pgfqpoint{5.527598in}{1.924667in}}%
\pgfpathquadraticcurveto{\pgfqpoint{5.527598in}{1.939552in}}{\pgfqpoint{5.523371in}{1.954045in}}%
\pgfpathquadraticcurveto{\pgfqpoint{5.519145in}{1.968559in}}{\pgfqpoint{5.510589in}{1.983362in}}%
\pgfpathlineto{\pgfqpoint{5.510589in}{1.983362in}}%
\pgfpathclose%
\pgfusepath{fill}%
\end{pgfscope}%
\begin{pgfscope}%
\pgfsetbuttcap%
\pgfsetroundjoin%
\definecolor{currentfill}{rgb}{0.090196,0.168627,0.227451}%
\pgfsetfillcolor{currentfill}%
\pgfsetlinewidth{0.000000pt}%
\definecolor{currentstroke}{rgb}{0.090196,0.168627,0.227451}%
\pgfsetstrokecolor{currentstroke}%
\pgfsetdash{}{0pt}%
\pgfpathmoveto{\pgfqpoint{4.550611in}{1.576690in}}%
\pgfpathlineto{\pgfqpoint{4.550611in}{1.563991in}}%
\pgfpathquadraticcurveto{\pgfqpoint{4.543210in}{1.567537in}}{\pgfqpoint{4.536633in}{1.569268in}}%
\pgfpathquadraticcurveto{\pgfqpoint{4.530056in}{1.571021in}}{\pgfqpoint{4.523933in}{1.571021in}}%
\pgfpathquadraticcurveto{\pgfqpoint{4.513316in}{1.571021in}}{\pgfqpoint{4.507543in}{1.566897in}}%
\pgfpathquadraticcurveto{\pgfqpoint{4.501771in}{1.562774in}}{\pgfqpoint{4.501771in}{1.555167in}}%
\pgfpathquadraticcurveto{\pgfqpoint{4.501771in}{1.548776in}}{\pgfqpoint{4.505605in}{1.545518in}}%
\pgfpathquadraticcurveto{\pgfqpoint{4.509440in}{1.542282in}}{\pgfqpoint{4.520140in}{1.540282in}}%
\pgfpathlineto{\pgfqpoint{4.527995in}{1.538674in}}%
\pgfpathquadraticcurveto{\pgfqpoint{4.542550in}{1.535891in}}{\pgfqpoint{4.549477in}{1.528902in}}%
\pgfpathquadraticcurveto{\pgfqpoint{4.556404in}{1.521913in}}{\pgfqpoint{4.556404in}{1.510203in}}%
\pgfpathquadraticcurveto{\pgfqpoint{4.556404in}{1.496204in}}{\pgfqpoint{4.547024in}{1.488988in}}%
\pgfpathquadraticcurveto{\pgfqpoint{4.537664in}{1.481773in}}{\pgfqpoint{4.519563in}{1.481773in}}%
\pgfpathquadraticcurveto{\pgfqpoint{4.512739in}{1.481773in}}{\pgfqpoint{4.505028in}{1.483319in}}%
\pgfpathquadraticcurveto{\pgfqpoint{4.497338in}{1.484865in}}{\pgfqpoint{4.489092in}{1.487896in}}%
\pgfpathlineto{\pgfqpoint{4.489092in}{1.501296in}}%
\pgfpathquadraticcurveto{\pgfqpoint{4.497008in}{1.496864in}}{\pgfqpoint{4.504616in}{1.494596in}}%
\pgfpathquadraticcurveto{\pgfqpoint{4.512223in}{1.492349in}}{\pgfqpoint{4.519563in}{1.492349in}}%
\pgfpathquadraticcurveto{\pgfqpoint{4.530696in}{1.492349in}}{\pgfqpoint{4.536757in}{1.496720in}}%
\pgfpathquadraticcurveto{\pgfqpoint{4.542818in}{1.501111in}}{\pgfqpoint{4.542818in}{1.509234in}}%
\pgfpathquadraticcurveto{\pgfqpoint{4.542818in}{1.516305in}}{\pgfqpoint{4.538468in}{1.520305in}}%
\pgfpathquadraticcurveto{\pgfqpoint{4.534118in}{1.524304in}}{\pgfqpoint{4.524201in}{1.526304in}}%
\pgfpathlineto{\pgfqpoint{4.516264in}{1.527850in}}%
\pgfpathquadraticcurveto{\pgfqpoint{4.501709in}{1.530736in}}{\pgfqpoint{4.495194in}{1.536921in}}%
\pgfpathquadraticcurveto{\pgfqpoint{4.488700in}{1.543106in}}{\pgfqpoint{4.488700in}{1.554136in}}%
\pgfpathquadraticcurveto{\pgfqpoint{4.488700in}{1.566897in}}{\pgfqpoint{4.497689in}{1.574237in}}%
\pgfpathquadraticcurveto{\pgfqpoint{4.506678in}{1.581576in}}{\pgfqpoint{4.522449in}{1.581576in}}%
\pgfpathquadraticcurveto{\pgfqpoint{4.529232in}{1.581576in}}{\pgfqpoint{4.536241in}{1.580339in}}%
\pgfpathquadraticcurveto{\pgfqpoint{4.543271in}{1.579123in}}{\pgfqpoint{4.550611in}{1.576690in}}%
\pgfpathlineto{\pgfqpoint{4.550611in}{1.576690in}}%
\pgfpathclose%
\pgfpathmoveto{\pgfqpoint{4.587926in}{1.576299in}}%
\pgfpathlineto{\pgfqpoint{4.587926in}{1.555806in}}%
\pgfpathlineto{\pgfqpoint{4.612336in}{1.555806in}}%
\pgfpathlineto{\pgfqpoint{4.612336in}{1.546590in}}%
\pgfpathlineto{\pgfqpoint{4.587926in}{1.546590in}}%
\pgfpathlineto{\pgfqpoint{4.587926in}{1.507419in}}%
\pgfpathquadraticcurveto{\pgfqpoint{4.587926in}{1.498596in}}{\pgfqpoint{4.590339in}{1.496080in}}%
\pgfpathquadraticcurveto{\pgfqpoint{4.592751in}{1.493565in}}{\pgfqpoint{4.600173in}{1.493565in}}%
\pgfpathlineto{\pgfqpoint{4.612336in}{1.493565in}}%
\pgfpathlineto{\pgfqpoint{4.612336in}{1.483649in}}%
\pgfpathlineto{\pgfqpoint{4.600173in}{1.483649in}}%
\pgfpathquadraticcurveto{\pgfqpoint{4.586442in}{1.483649in}}{\pgfqpoint{4.581226in}{1.488762in}}%
\pgfpathquadraticcurveto{\pgfqpoint{4.576010in}{1.493895in}}{\pgfqpoint{4.576010in}{1.507419in}}%
\pgfpathlineto{\pgfqpoint{4.576010in}{1.546590in}}%
\pgfpathlineto{\pgfqpoint{4.567310in}{1.546590in}}%
\pgfpathlineto{\pgfqpoint{4.567310in}{1.555806in}}%
\pgfpathlineto{\pgfqpoint{4.576010in}{1.555806in}}%
\pgfpathlineto{\pgfqpoint{4.576010in}{1.576299in}}%
\pgfpathlineto{\pgfqpoint{4.587926in}{1.576299in}}%
\pgfpathlineto{\pgfqpoint{4.587926in}{1.576299in}}%
\pgfpathclose%
\pgfpathmoveto{\pgfqpoint{4.626706in}{1.512120in}}%
\pgfpathlineto{\pgfqpoint{4.626706in}{1.555806in}}%
\pgfpathlineto{\pgfqpoint{4.638560in}{1.555806in}}%
\pgfpathlineto{\pgfqpoint{4.638560in}{1.512573in}}%
\pgfpathquadraticcurveto{\pgfqpoint{4.638560in}{1.502327in}}{\pgfqpoint{4.642539in}{1.497194in}}%
\pgfpathquadraticcurveto{\pgfqpoint{4.646539in}{1.492081in}}{\pgfqpoint{4.654538in}{1.492081in}}%
\pgfpathquadraticcurveto{\pgfqpoint{4.664124in}{1.492081in}}{\pgfqpoint{4.669691in}{1.498204in}}%
\pgfpathquadraticcurveto{\pgfqpoint{4.675278in}{1.504327in}}{\pgfqpoint{4.675278in}{1.514903in}}%
\pgfpathlineto{\pgfqpoint{4.675278in}{1.555806in}}%
\pgfpathlineto{\pgfqpoint{4.687132in}{1.555806in}}%
\pgfpathlineto{\pgfqpoint{4.687132in}{1.483649in}}%
\pgfpathlineto{\pgfqpoint{4.675278in}{1.483649in}}%
\pgfpathlineto{\pgfqpoint{4.675278in}{1.494740in}}%
\pgfpathquadraticcurveto{\pgfqpoint{4.670969in}{1.488164in}}{\pgfqpoint{4.665258in}{1.484968in}}%
\pgfpathquadraticcurveto{\pgfqpoint{4.659568in}{1.481773in}}{\pgfqpoint{4.652023in}{1.481773in}}%
\pgfpathquadraticcurveto{\pgfqpoint{4.639591in}{1.481773in}}{\pgfqpoint{4.633138in}{1.489504in}}%
\pgfpathquadraticcurveto{\pgfqpoint{4.626706in}{1.497235in}}{\pgfqpoint{4.626706in}{1.512120in}}%
\pgfpathlineto{\pgfqpoint{4.626706in}{1.512120in}}%
\pgfpathclose%
\pgfpathmoveto{\pgfqpoint{4.656538in}{1.557538in}}%
\pgfpathlineto{\pgfqpoint{4.656538in}{1.557538in}}%
\pgfpathlineto{\pgfqpoint{4.656538in}{1.557538in}}%
\pgfpathclose%
\pgfpathmoveto{\pgfqpoint{4.759021in}{1.544859in}}%
\pgfpathlineto{\pgfqpoint{4.759021in}{1.583906in}}%
\pgfpathlineto{\pgfqpoint{4.770876in}{1.583906in}}%
\pgfpathlineto{\pgfqpoint{4.770876in}{1.483649in}}%
\pgfpathlineto{\pgfqpoint{4.759021in}{1.483649in}}%
\pgfpathlineto{\pgfqpoint{4.759021in}{1.494472in}}%
\pgfpathquadraticcurveto{\pgfqpoint{4.755290in}{1.488040in}}{\pgfqpoint{4.749579in}{1.484906in}}%
\pgfpathquadraticcurveto{\pgfqpoint{4.743889in}{1.481773in}}{\pgfqpoint{4.735890in}{1.481773in}}%
\pgfpathquadraticcurveto{\pgfqpoint{4.722819in}{1.481773in}}{\pgfqpoint{4.714593in}{1.492205in}}%
\pgfpathquadraticcurveto{\pgfqpoint{4.706388in}{1.502657in}}{\pgfqpoint{4.706388in}{1.519666in}}%
\pgfpathquadraticcurveto{\pgfqpoint{4.706388in}{1.536674in}}{\pgfqpoint{4.714593in}{1.547106in}}%
\pgfpathquadraticcurveto{\pgfqpoint{4.722819in}{1.557538in}}{\pgfqpoint{4.735890in}{1.557538in}}%
\pgfpathquadraticcurveto{\pgfqpoint{4.743889in}{1.557538in}}{\pgfqpoint{4.749579in}{1.554404in}}%
\pgfpathquadraticcurveto{\pgfqpoint{4.755290in}{1.551291in}}{\pgfqpoint{4.759021in}{1.544859in}}%
\pgfpathlineto{\pgfqpoint{4.759021in}{1.544859in}}%
\pgfpathclose%
\pgfpathmoveto{\pgfqpoint{4.718634in}{1.519666in}}%
\pgfpathquadraticcurveto{\pgfqpoint{4.718634in}{1.506595in}}{\pgfqpoint{4.724015in}{1.499152in}}%
\pgfpathquadraticcurveto{\pgfqpoint{4.729396in}{1.491710in}}{\pgfqpoint{4.738797in}{1.491710in}}%
\pgfpathquadraticcurveto{\pgfqpoint{4.748198in}{1.491710in}}{\pgfqpoint{4.753599in}{1.499152in}}%
\pgfpathquadraticcurveto{\pgfqpoint{4.759021in}{1.506595in}}{\pgfqpoint{4.759021in}{1.519666in}}%
\pgfpathquadraticcurveto{\pgfqpoint{4.759021in}{1.532736in}}{\pgfqpoint{4.753599in}{1.540179in}}%
\pgfpathquadraticcurveto{\pgfqpoint{4.748198in}{1.547621in}}{\pgfqpoint{4.738797in}{1.547621in}}%
\pgfpathquadraticcurveto{\pgfqpoint{4.729396in}{1.547621in}}{\pgfqpoint{4.724015in}{1.540179in}}%
\pgfpathquadraticcurveto{\pgfqpoint{4.718634in}{1.532736in}}{\pgfqpoint{4.718634in}{1.519666in}}%
\pgfpathlineto{\pgfqpoint{4.718634in}{1.519666in}}%
\pgfpathclose%
\pgfpathmoveto{\pgfqpoint{4.825323in}{1.476948in}}%
\pgfpathquadraticcurveto{\pgfqpoint{4.820314in}{1.464063in}}{\pgfqpoint{4.815531in}{1.460146in}}%
\pgfpathquadraticcurveto{\pgfqpoint{4.810768in}{1.456208in}}{\pgfqpoint{4.802790in}{1.456208in}}%
\pgfpathlineto{\pgfqpoint{4.793306in}{1.456208in}}%
\pgfpathlineto{\pgfqpoint{4.793306in}{1.466125in}}%
\pgfpathlineto{\pgfqpoint{4.800275in}{1.466125in}}%
\pgfpathquadraticcurveto{\pgfqpoint{4.805161in}{1.466125in}}{\pgfqpoint{4.807861in}{1.468455in}}%
\pgfpathquadraticcurveto{\pgfqpoint{4.810583in}{1.470764in}}{\pgfqpoint{4.813861in}{1.479402in}}%
\pgfpathlineto{\pgfqpoint{4.815984in}{1.484803in}}%
\pgfpathlineto{\pgfqpoint{4.786812in}{1.555806in}}%
\pgfpathlineto{\pgfqpoint{4.799367in}{1.555806in}}%
\pgfpathlineto{\pgfqpoint{4.821922in}{1.499379in}}%
\pgfpathlineto{\pgfqpoint{4.844476in}{1.555806in}}%
\pgfpathlineto{\pgfqpoint{4.857031in}{1.555806in}}%
\pgfpathlineto{\pgfqpoint{4.825323in}{1.476948in}}%
\pgfpathlineto{\pgfqpoint{4.825323in}{1.476948in}}%
\pgfpathclose%
\pgfpathmoveto{\pgfqpoint{4.919272in}{1.494596in}}%
\pgfpathlineto{\pgfqpoint{4.940527in}{1.494596in}}%
\pgfpathlineto{\pgfqpoint{4.940527in}{1.567990in}}%
\pgfpathlineto{\pgfqpoint{4.917396in}{1.563351in}}%
\pgfpathlineto{\pgfqpoint{4.917396in}{1.575206in}}%
\pgfpathlineto{\pgfqpoint{4.940404in}{1.579845in}}%
\pgfpathlineto{\pgfqpoint{4.953413in}{1.579845in}}%
\pgfpathlineto{\pgfqpoint{4.953413in}{1.494596in}}%
\pgfpathlineto{\pgfqpoint{4.974668in}{1.494596in}}%
\pgfpathlineto{\pgfqpoint{4.974668in}{1.483649in}}%
\pgfpathlineto{\pgfqpoint{4.919272in}{1.483649in}}%
\pgfpathlineto{\pgfqpoint{4.919272in}{1.494596in}}%
\pgfpathlineto{\pgfqpoint{4.919272in}{1.494596in}}%
\pgfpathclose%
\pgfpathmoveto{\pgfqpoint{5.003222in}{1.494596in}}%
\pgfpathlineto{\pgfqpoint{5.024477in}{1.494596in}}%
\pgfpathlineto{\pgfqpoint{5.024477in}{1.567990in}}%
\pgfpathlineto{\pgfqpoint{5.001345in}{1.563351in}}%
\pgfpathlineto{\pgfqpoint{5.001345in}{1.575206in}}%
\pgfpathlineto{\pgfqpoint{5.024353in}{1.579845in}}%
\pgfpathlineto{\pgfqpoint{5.037362in}{1.579845in}}%
\pgfpathlineto{\pgfqpoint{5.037362in}{1.494596in}}%
\pgfpathlineto{\pgfqpoint{5.058618in}{1.494596in}}%
\pgfpathlineto{\pgfqpoint{5.058618in}{1.483649in}}%
\pgfpathlineto{\pgfqpoint{5.003222in}{1.483649in}}%
\pgfpathlineto{\pgfqpoint{5.003222in}{1.494596in}}%
\pgfpathlineto{\pgfqpoint{5.003222in}{1.494596in}}%
\pgfpathclose%
\pgfpathmoveto{\pgfqpoint{5.153638in}{1.583762in}}%
\pgfpathquadraticcurveto{\pgfqpoint{5.145021in}{1.568959in}}{\pgfqpoint{5.140815in}{1.554445in}}%
\pgfpathquadraticcurveto{\pgfqpoint{5.136630in}{1.539952in}}{\pgfqpoint{5.136630in}{1.525067in}}%
\pgfpathquadraticcurveto{\pgfqpoint{5.136630in}{1.510203in}}{\pgfqpoint{5.140856in}{1.495606in}}%
\pgfpathquadraticcurveto{\pgfqpoint{5.145082in}{1.481010in}}{\pgfqpoint{5.153638in}{1.466249in}}%
\pgfpathlineto{\pgfqpoint{5.143330in}{1.466249in}}%
\pgfpathquadraticcurveto{\pgfqpoint{5.133682in}{1.481402in}}{\pgfqpoint{5.128878in}{1.496019in}}%
\pgfpathquadraticcurveto{\pgfqpoint{5.124074in}{1.510636in}}{\pgfqpoint{5.124074in}{1.525067in}}%
\pgfpathquadraticcurveto{\pgfqpoint{5.124074in}{1.539437in}}{\pgfqpoint{5.128837in}{1.553992in}}%
\pgfpathquadraticcurveto{\pgfqpoint{5.133620in}{1.568567in}}{\pgfqpoint{5.143330in}{1.583762in}}%
\pgfpathlineto{\pgfqpoint{5.153638in}{1.583762in}}%
\pgfpathlineto{\pgfqpoint{5.153638in}{1.583762in}}%
\pgfpathclose%
\pgfpathmoveto{\pgfqpoint{5.189531in}{1.494596in}}%
\pgfpathlineto{\pgfqpoint{5.234949in}{1.494596in}}%
\pgfpathlineto{\pgfqpoint{5.234949in}{1.483649in}}%
\pgfpathlineto{\pgfqpoint{5.173883in}{1.483649in}}%
\pgfpathlineto{\pgfqpoint{5.173883in}{1.494596in}}%
\pgfpathquadraticcurveto{\pgfqpoint{5.181285in}{1.502265in}}{\pgfqpoint{5.194067in}{1.515171in}}%
\pgfpathquadraticcurveto{\pgfqpoint{5.206870in}{1.528098in}}{\pgfqpoint{5.210148in}{1.531850in}}%
\pgfpathquadraticcurveto{\pgfqpoint{5.216394in}{1.538859in}}{\pgfqpoint{5.218868in}{1.543725in}}%
\pgfpathquadraticcurveto{\pgfqpoint{5.221363in}{1.548590in}}{\pgfqpoint{5.221363in}{1.553291in}}%
\pgfpathquadraticcurveto{\pgfqpoint{5.221363in}{1.560960in}}{\pgfqpoint{5.215982in}{1.565784in}}%
\pgfpathquadraticcurveto{\pgfqpoint{5.210601in}{1.570629in}}{\pgfqpoint{5.201963in}{1.570629in}}%
\pgfpathquadraticcurveto{\pgfqpoint{5.195840in}{1.570629in}}{\pgfqpoint{5.189036in}{1.568506in}}%
\pgfpathquadraticcurveto{\pgfqpoint{5.182254in}{1.566382in}}{\pgfqpoint{5.174523in}{1.562053in}}%
\pgfpathlineto{\pgfqpoint{5.174523in}{1.575206in}}%
\pgfpathquadraticcurveto{\pgfqpoint{5.182377in}{1.578360in}}{\pgfqpoint{5.189201in}{1.579968in}}%
\pgfpathquadraticcurveto{\pgfqpoint{5.196046in}{1.581576in}}{\pgfqpoint{5.201715in}{1.581576in}}%
\pgfpathquadraticcurveto{\pgfqpoint{5.216662in}{1.581576in}}{\pgfqpoint{5.225548in}{1.574093in}}%
\pgfpathquadraticcurveto{\pgfqpoint{5.234434in}{1.566629in}}{\pgfqpoint{5.234434in}{1.554136in}}%
\pgfpathquadraticcurveto{\pgfqpoint{5.234434in}{1.548198in}}{\pgfqpoint{5.232207in}{1.542879in}}%
\pgfpathquadraticcurveto{\pgfqpoint{5.230001in}{1.537581in}}{\pgfqpoint{5.224125in}{1.530365in}}%
\pgfpathquadraticcurveto{\pgfqpoint{5.222517in}{1.528489in}}{\pgfqpoint{5.213879in}{1.519562in}}%
\pgfpathquadraticcurveto{\pgfqpoint{5.205262in}{1.510636in}}{\pgfqpoint{5.189531in}{1.494596in}}%
\pgfpathlineto{\pgfqpoint{5.189531in}{1.494596in}}%
\pgfpathclose%
\pgfpathmoveto{\pgfqpoint{5.290098in}{1.571268in}}%
\pgfpathquadraticcurveto{\pgfqpoint{5.280058in}{1.571268in}}{\pgfqpoint{5.274986in}{1.561372in}}%
\pgfpathquadraticcurveto{\pgfqpoint{5.269935in}{1.551497in}}{\pgfqpoint{5.269935in}{1.531644in}}%
\pgfpathquadraticcurveto{\pgfqpoint{5.269935in}{1.511873in}}{\pgfqpoint{5.274986in}{1.501977in}}%
\pgfpathquadraticcurveto{\pgfqpoint{5.280058in}{1.492081in}}{\pgfqpoint{5.290098in}{1.492081in}}%
\pgfpathquadraticcurveto{\pgfqpoint{5.300220in}{1.492081in}}{\pgfqpoint{5.305271in}{1.501977in}}%
\pgfpathquadraticcurveto{\pgfqpoint{5.310343in}{1.511873in}}{\pgfqpoint{5.310343in}{1.531644in}}%
\pgfpathquadraticcurveto{\pgfqpoint{5.310343in}{1.551497in}}{\pgfqpoint{5.305271in}{1.561372in}}%
\pgfpathquadraticcurveto{\pgfqpoint{5.300220in}{1.571268in}}{\pgfqpoint{5.290098in}{1.571268in}}%
\pgfpathlineto{\pgfqpoint{5.290098in}{1.571268in}}%
\pgfpathclose%
\pgfpathmoveto{\pgfqpoint{5.290098in}{1.581576in}}%
\pgfpathquadraticcurveto{\pgfqpoint{5.306281in}{1.581576in}}{\pgfqpoint{5.314817in}{1.568774in}}%
\pgfpathquadraticcurveto{\pgfqpoint{5.323352in}{1.555991in}}{\pgfqpoint{5.323352in}{1.531644in}}%
\pgfpathquadraticcurveto{\pgfqpoint{5.323352in}{1.507358in}}{\pgfqpoint{5.314817in}{1.494555in}}%
\pgfpathquadraticcurveto{\pgfqpoint{5.306281in}{1.481773in}}{\pgfqpoint{5.290098in}{1.481773in}}%
\pgfpathquadraticcurveto{\pgfqpoint{5.273934in}{1.481773in}}{\pgfqpoint{5.265399in}{1.494555in}}%
\pgfpathquadraticcurveto{\pgfqpoint{5.256864in}{1.507358in}}{\pgfqpoint{5.256864in}{1.531644in}}%
\pgfpathquadraticcurveto{\pgfqpoint{5.256864in}{1.555991in}}{\pgfqpoint{5.265399in}{1.568774in}}%
\pgfpathquadraticcurveto{\pgfqpoint{5.273934in}{1.581576in}}{\pgfqpoint{5.290098in}{1.581576in}}%
\pgfpathlineto{\pgfqpoint{5.290098in}{1.581576in}}%
\pgfpathclose%
\pgfpathmoveto{\pgfqpoint{5.357431in}{1.494596in}}%
\pgfpathlineto{\pgfqpoint{5.402848in}{1.494596in}}%
\pgfpathlineto{\pgfqpoint{5.402848in}{1.483649in}}%
\pgfpathlineto{\pgfqpoint{5.341783in}{1.483649in}}%
\pgfpathlineto{\pgfqpoint{5.341783in}{1.494596in}}%
\pgfpathquadraticcurveto{\pgfqpoint{5.349184in}{1.502265in}}{\pgfqpoint{5.361966in}{1.515171in}}%
\pgfpathquadraticcurveto{\pgfqpoint{5.374769in}{1.528098in}}{\pgfqpoint{5.378047in}{1.531850in}}%
\pgfpathquadraticcurveto{\pgfqpoint{5.384294in}{1.538859in}}{\pgfqpoint{5.386768in}{1.543725in}}%
\pgfpathquadraticcurveto{\pgfqpoint{5.389262in}{1.548590in}}{\pgfqpoint{5.389262in}{1.553291in}}%
\pgfpathquadraticcurveto{\pgfqpoint{5.389262in}{1.560960in}}{\pgfqpoint{5.383881in}{1.565784in}}%
\pgfpathquadraticcurveto{\pgfqpoint{5.378500in}{1.570629in}}{\pgfqpoint{5.369862in}{1.570629in}}%
\pgfpathquadraticcurveto{\pgfqpoint{5.363739in}{1.570629in}}{\pgfqpoint{5.356936in}{1.568506in}}%
\pgfpathquadraticcurveto{\pgfqpoint{5.350153in}{1.566382in}}{\pgfqpoint{5.342422in}{1.562053in}}%
\pgfpathlineto{\pgfqpoint{5.342422in}{1.575206in}}%
\pgfpathquadraticcurveto{\pgfqpoint{5.350277in}{1.578360in}}{\pgfqpoint{5.357101in}{1.579968in}}%
\pgfpathquadraticcurveto{\pgfqpoint{5.363945in}{1.581576in}}{\pgfqpoint{5.369615in}{1.581576in}}%
\pgfpathquadraticcurveto{\pgfqpoint{5.384562in}{1.581576in}}{\pgfqpoint{5.393447in}{1.574093in}}%
\pgfpathquadraticcurveto{\pgfqpoint{5.402333in}{1.566629in}}{\pgfqpoint{5.402333in}{1.554136in}}%
\pgfpathquadraticcurveto{\pgfqpoint{5.402333in}{1.548198in}}{\pgfqpoint{5.400106in}{1.542879in}}%
\pgfpathquadraticcurveto{\pgfqpoint{5.397900in}{1.537581in}}{\pgfqpoint{5.392025in}{1.530365in}}%
\pgfpathquadraticcurveto{\pgfqpoint{5.390417in}{1.528489in}}{\pgfqpoint{5.381778in}{1.519562in}}%
\pgfpathquadraticcurveto{\pgfqpoint{5.373161in}{1.510636in}}{\pgfqpoint{5.357431in}{1.494596in}}%
\pgfpathlineto{\pgfqpoint{5.357431in}{1.494596in}}%
\pgfpathclose%
\pgfpathmoveto{\pgfqpoint{5.432433in}{1.494596in}}%
\pgfpathlineto{\pgfqpoint{5.453688in}{1.494596in}}%
\pgfpathlineto{\pgfqpoint{5.453688in}{1.567990in}}%
\pgfpathlineto{\pgfqpoint{5.430557in}{1.563351in}}%
\pgfpathlineto{\pgfqpoint{5.430557in}{1.575206in}}%
\pgfpathlineto{\pgfqpoint{5.453564in}{1.579845in}}%
\pgfpathlineto{\pgfqpoint{5.466573in}{1.579845in}}%
\pgfpathlineto{\pgfqpoint{5.466573in}{1.494596in}}%
\pgfpathlineto{\pgfqpoint{5.487829in}{1.494596in}}%
\pgfpathlineto{\pgfqpoint{5.487829in}{1.483649in}}%
\pgfpathlineto{\pgfqpoint{5.432433in}{1.483649in}}%
\pgfpathlineto{\pgfqpoint{5.432433in}{1.494596in}}%
\pgfpathlineto{\pgfqpoint{5.432433in}{1.494596in}}%
\pgfpathclose%
\pgfpathmoveto{\pgfqpoint{5.510589in}{1.583762in}}%
\pgfpathlineto{\pgfqpoint{5.520897in}{1.583762in}}%
\pgfpathquadraticcurveto{\pgfqpoint{5.530546in}{1.568567in}}{\pgfqpoint{5.535349in}{1.553992in}}%
\pgfpathquadraticcurveto{\pgfqpoint{5.540153in}{1.539437in}}{\pgfqpoint{5.540153in}{1.525067in}}%
\pgfpathquadraticcurveto{\pgfqpoint{5.540153in}{1.510636in}}{\pgfqpoint{5.535349in}{1.496019in}}%
\pgfpathquadraticcurveto{\pgfqpoint{5.530546in}{1.481402in}}{\pgfqpoint{5.520897in}{1.466249in}}%
\pgfpathlineto{\pgfqpoint{5.510589in}{1.466249in}}%
\pgfpathquadraticcurveto{\pgfqpoint{5.519145in}{1.481010in}}{\pgfqpoint{5.523371in}{1.495606in}}%
\pgfpathquadraticcurveto{\pgfqpoint{5.527598in}{1.510203in}}{\pgfqpoint{5.527598in}{1.525067in}}%
\pgfpathquadraticcurveto{\pgfqpoint{5.527598in}{1.539952in}}{\pgfqpoint{5.523371in}{1.554445in}}%
\pgfpathquadraticcurveto{\pgfqpoint{5.519145in}{1.568959in}}{\pgfqpoint{5.510589in}{1.583762in}}%
\pgfpathlineto{\pgfqpoint{5.510589in}{1.583762in}}%
\pgfpathclose%
\pgfusepath{fill}%
\end{pgfscope}%
\begin{pgfscope}%
\pgfsetbuttcap%
\pgfsetroundjoin%
\definecolor{currentfill}{rgb}{0.090196,0.168627,0.227451}%
\pgfsetfillcolor{currentfill}%
\pgfsetlinewidth{0.000000pt}%
\definecolor{currentstroke}{rgb}{0.090196,0.168627,0.227451}%
\pgfsetstrokecolor{currentstroke}%
\pgfsetdash{}{0pt}%
\pgfpathmoveto{\pgfqpoint{7.784611in}{3.175090in}}%
\pgfpathlineto{\pgfqpoint{7.784611in}{3.162391in}}%
\pgfpathquadraticcurveto{\pgfqpoint{7.777210in}{3.165937in}}{\pgfqpoint{7.770633in}{3.167668in}}%
\pgfpathquadraticcurveto{\pgfqpoint{7.764056in}{3.169421in}}{\pgfqpoint{7.757933in}{3.169421in}}%
\pgfpathquadraticcurveto{\pgfqpoint{7.747316in}{3.169421in}}{\pgfqpoint{7.741543in}{3.165297in}}%
\pgfpathquadraticcurveto{\pgfqpoint{7.735771in}{3.161174in}}{\pgfqpoint{7.735771in}{3.153567in}}%
\pgfpathquadraticcurveto{\pgfqpoint{7.735771in}{3.147176in}}{\pgfqpoint{7.739605in}{3.143918in}}%
\pgfpathquadraticcurveto{\pgfqpoint{7.743440in}{3.140682in}}{\pgfqpoint{7.754140in}{3.138682in}}%
\pgfpathlineto{\pgfqpoint{7.761995in}{3.137074in}}%
\pgfpathquadraticcurveto{\pgfqpoint{7.776550in}{3.134291in}}{\pgfqpoint{7.783477in}{3.127302in}}%
\pgfpathquadraticcurveto{\pgfqpoint{7.790404in}{3.120313in}}{\pgfqpoint{7.790404in}{3.108603in}}%
\pgfpathquadraticcurveto{\pgfqpoint{7.790404in}{3.094604in}}{\pgfqpoint{7.781024in}{3.087388in}}%
\pgfpathquadraticcurveto{\pgfqpoint{7.771664in}{3.080173in}}{\pgfqpoint{7.753563in}{3.080173in}}%
\pgfpathquadraticcurveto{\pgfqpoint{7.746739in}{3.080173in}}{\pgfqpoint{7.739028in}{3.081719in}}%
\pgfpathquadraticcurveto{\pgfqpoint{7.731338in}{3.083265in}}{\pgfqpoint{7.723092in}{3.086296in}}%
\pgfpathlineto{\pgfqpoint{7.723092in}{3.099696in}}%
\pgfpathquadraticcurveto{\pgfqpoint{7.731008in}{3.095264in}}{\pgfqpoint{7.738616in}{3.092996in}}%
\pgfpathquadraticcurveto{\pgfqpoint{7.746223in}{3.090749in}}{\pgfqpoint{7.753563in}{3.090749in}}%
\pgfpathquadraticcurveto{\pgfqpoint{7.764696in}{3.090749in}}{\pgfqpoint{7.770757in}{3.095120in}}%
\pgfpathquadraticcurveto{\pgfqpoint{7.776818in}{3.099511in}}{\pgfqpoint{7.776818in}{3.107634in}}%
\pgfpathquadraticcurveto{\pgfqpoint{7.776818in}{3.114705in}}{\pgfqpoint{7.772468in}{3.118705in}}%
\pgfpathquadraticcurveto{\pgfqpoint{7.768118in}{3.122704in}}{\pgfqpoint{7.758201in}{3.124704in}}%
\pgfpathlineto{\pgfqpoint{7.750264in}{3.126250in}}%
\pgfpathquadraticcurveto{\pgfqpoint{7.735709in}{3.129136in}}{\pgfqpoint{7.729194in}{3.135321in}}%
\pgfpathquadraticcurveto{\pgfqpoint{7.722700in}{3.141506in}}{\pgfqpoint{7.722700in}{3.152536in}}%
\pgfpathquadraticcurveto{\pgfqpoint{7.722700in}{3.165297in}}{\pgfqpoint{7.731689in}{3.172637in}}%
\pgfpathquadraticcurveto{\pgfqpoint{7.740678in}{3.179976in}}{\pgfqpoint{7.756449in}{3.179976in}}%
\pgfpathquadraticcurveto{\pgfqpoint{7.763232in}{3.179976in}}{\pgfqpoint{7.770241in}{3.178739in}}%
\pgfpathquadraticcurveto{\pgfqpoint{7.777271in}{3.177523in}}{\pgfqpoint{7.784611in}{3.175090in}}%
\pgfpathlineto{\pgfqpoint{7.784611in}{3.175090in}}%
\pgfpathclose%
\pgfpathmoveto{\pgfqpoint{7.821926in}{3.174699in}}%
\pgfpathlineto{\pgfqpoint{7.821926in}{3.154206in}}%
\pgfpathlineto{\pgfqpoint{7.846336in}{3.154206in}}%
\pgfpathlineto{\pgfqpoint{7.846336in}{3.144990in}}%
\pgfpathlineto{\pgfqpoint{7.821926in}{3.144990in}}%
\pgfpathlineto{\pgfqpoint{7.821926in}{3.105819in}}%
\pgfpathquadraticcurveto{\pgfqpoint{7.821926in}{3.096996in}}{\pgfqpoint{7.824339in}{3.094480in}}%
\pgfpathquadraticcurveto{\pgfqpoint{7.826751in}{3.091965in}}{\pgfqpoint{7.834173in}{3.091965in}}%
\pgfpathlineto{\pgfqpoint{7.846336in}{3.091965in}}%
\pgfpathlineto{\pgfqpoint{7.846336in}{3.082049in}}%
\pgfpathlineto{\pgfqpoint{7.834173in}{3.082049in}}%
\pgfpathquadraticcurveto{\pgfqpoint{7.820442in}{3.082049in}}{\pgfqpoint{7.815226in}{3.087162in}}%
\pgfpathquadraticcurveto{\pgfqpoint{7.810010in}{3.092295in}}{\pgfqpoint{7.810010in}{3.105819in}}%
\pgfpathlineto{\pgfqpoint{7.810010in}{3.144990in}}%
\pgfpathlineto{\pgfqpoint{7.801310in}{3.144990in}}%
\pgfpathlineto{\pgfqpoint{7.801310in}{3.154206in}}%
\pgfpathlineto{\pgfqpoint{7.810010in}{3.154206in}}%
\pgfpathlineto{\pgfqpoint{7.810010in}{3.174699in}}%
\pgfpathlineto{\pgfqpoint{7.821926in}{3.174699in}}%
\pgfpathlineto{\pgfqpoint{7.821926in}{3.174699in}}%
\pgfpathclose%
\pgfpathmoveto{\pgfqpoint{7.860706in}{3.110520in}}%
\pgfpathlineto{\pgfqpoint{7.860706in}{3.154206in}}%
\pgfpathlineto{\pgfqpoint{7.872560in}{3.154206in}}%
\pgfpathlineto{\pgfqpoint{7.872560in}{3.110973in}}%
\pgfpathquadraticcurveto{\pgfqpoint{7.872560in}{3.100727in}}{\pgfqpoint{7.876539in}{3.095594in}}%
\pgfpathquadraticcurveto{\pgfqpoint{7.880539in}{3.090481in}}{\pgfqpoint{7.888538in}{3.090481in}}%
\pgfpathquadraticcurveto{\pgfqpoint{7.898124in}{3.090481in}}{\pgfqpoint{7.903691in}{3.096604in}}%
\pgfpathquadraticcurveto{\pgfqpoint{7.909278in}{3.102727in}}{\pgfqpoint{7.909278in}{3.113303in}}%
\pgfpathlineto{\pgfqpoint{7.909278in}{3.154206in}}%
\pgfpathlineto{\pgfqpoint{7.921132in}{3.154206in}}%
\pgfpathlineto{\pgfqpoint{7.921132in}{3.082049in}}%
\pgfpathlineto{\pgfqpoint{7.909278in}{3.082049in}}%
\pgfpathlineto{\pgfqpoint{7.909278in}{3.093140in}}%
\pgfpathquadraticcurveto{\pgfqpoint{7.904969in}{3.086564in}}{\pgfqpoint{7.899258in}{3.083368in}}%
\pgfpathquadraticcurveto{\pgfqpoint{7.893568in}{3.080173in}}{\pgfqpoint{7.886023in}{3.080173in}}%
\pgfpathquadraticcurveto{\pgfqpoint{7.873591in}{3.080173in}}{\pgfqpoint{7.867138in}{3.087904in}}%
\pgfpathquadraticcurveto{\pgfqpoint{7.860706in}{3.095635in}}{\pgfqpoint{7.860706in}{3.110520in}}%
\pgfpathlineto{\pgfqpoint{7.860706in}{3.110520in}}%
\pgfpathclose%
\pgfpathmoveto{\pgfqpoint{7.890538in}{3.155938in}}%
\pgfpathlineto{\pgfqpoint{7.890538in}{3.155938in}}%
\pgfpathlineto{\pgfqpoint{7.890538in}{3.155938in}}%
\pgfpathclose%
\pgfpathmoveto{\pgfqpoint{7.993021in}{3.143259in}}%
\pgfpathlineto{\pgfqpoint{7.993021in}{3.182306in}}%
\pgfpathlineto{\pgfqpoint{8.004876in}{3.182306in}}%
\pgfpathlineto{\pgfqpoint{8.004876in}{3.082049in}}%
\pgfpathlineto{\pgfqpoint{7.993021in}{3.082049in}}%
\pgfpathlineto{\pgfqpoint{7.993021in}{3.092872in}}%
\pgfpathquadraticcurveto{\pgfqpoint{7.989290in}{3.086440in}}{\pgfqpoint{7.983579in}{3.083306in}}%
\pgfpathquadraticcurveto{\pgfqpoint{7.977889in}{3.080173in}}{\pgfqpoint{7.969890in}{3.080173in}}%
\pgfpathquadraticcurveto{\pgfqpoint{7.956819in}{3.080173in}}{\pgfqpoint{7.948593in}{3.090605in}}%
\pgfpathquadraticcurveto{\pgfqpoint{7.940388in}{3.101057in}}{\pgfqpoint{7.940388in}{3.118066in}}%
\pgfpathquadraticcurveto{\pgfqpoint{7.940388in}{3.135074in}}{\pgfqpoint{7.948593in}{3.145506in}}%
\pgfpathquadraticcurveto{\pgfqpoint{7.956819in}{3.155938in}}{\pgfqpoint{7.969890in}{3.155938in}}%
\pgfpathquadraticcurveto{\pgfqpoint{7.977889in}{3.155938in}}{\pgfqpoint{7.983579in}{3.152804in}}%
\pgfpathquadraticcurveto{\pgfqpoint{7.989290in}{3.149691in}}{\pgfqpoint{7.993021in}{3.143259in}}%
\pgfpathlineto{\pgfqpoint{7.993021in}{3.143259in}}%
\pgfpathclose%
\pgfpathmoveto{\pgfqpoint{7.952634in}{3.118066in}}%
\pgfpathquadraticcurveto{\pgfqpoint{7.952634in}{3.104995in}}{\pgfqpoint{7.958015in}{3.097552in}}%
\pgfpathquadraticcurveto{\pgfqpoint{7.963396in}{3.090110in}}{\pgfqpoint{7.972797in}{3.090110in}}%
\pgfpathquadraticcurveto{\pgfqpoint{7.982198in}{3.090110in}}{\pgfqpoint{7.987599in}{3.097552in}}%
\pgfpathquadraticcurveto{\pgfqpoint{7.993021in}{3.104995in}}{\pgfqpoint{7.993021in}{3.118066in}}%
\pgfpathquadraticcurveto{\pgfqpoint{7.993021in}{3.131136in}}{\pgfqpoint{7.987599in}{3.138579in}}%
\pgfpathquadraticcurveto{\pgfqpoint{7.982198in}{3.146021in}}{\pgfqpoint{7.972797in}{3.146021in}}%
\pgfpathquadraticcurveto{\pgfqpoint{7.963396in}{3.146021in}}{\pgfqpoint{7.958015in}{3.138579in}}%
\pgfpathquadraticcurveto{\pgfqpoint{7.952634in}{3.131136in}}{\pgfqpoint{7.952634in}{3.118066in}}%
\pgfpathlineto{\pgfqpoint{7.952634in}{3.118066in}}%
\pgfpathclose%
\pgfpathmoveto{\pgfqpoint{8.059323in}{3.075348in}}%
\pgfpathquadraticcurveto{\pgfqpoint{8.054314in}{3.062463in}}{\pgfqpoint{8.049531in}{3.058546in}}%
\pgfpathquadraticcurveto{\pgfqpoint{8.044768in}{3.054608in}}{\pgfqpoint{8.036790in}{3.054608in}}%
\pgfpathlineto{\pgfqpoint{8.027306in}{3.054608in}}%
\pgfpathlineto{\pgfqpoint{8.027306in}{3.064525in}}%
\pgfpathlineto{\pgfqpoint{8.034275in}{3.064525in}}%
\pgfpathquadraticcurveto{\pgfqpoint{8.039161in}{3.064525in}}{\pgfqpoint{8.041861in}{3.066855in}}%
\pgfpathquadraticcurveto{\pgfqpoint{8.044583in}{3.069164in}}{\pgfqpoint{8.047861in}{3.077802in}}%
\pgfpathlineto{\pgfqpoint{8.049984in}{3.083203in}}%
\pgfpathlineto{\pgfqpoint{8.020812in}{3.154206in}}%
\pgfpathlineto{\pgfqpoint{8.033367in}{3.154206in}}%
\pgfpathlineto{\pgfqpoint{8.055922in}{3.097779in}}%
\pgfpathlineto{\pgfqpoint{8.078476in}{3.154206in}}%
\pgfpathlineto{\pgfqpoint{8.091031in}{3.154206in}}%
\pgfpathlineto{\pgfqpoint{8.059323in}{3.075348in}}%
\pgfpathlineto{\pgfqpoint{8.059323in}{3.075348in}}%
\pgfpathclose%
\pgfpathmoveto{\pgfqpoint{8.153272in}{3.092996in}}%
\pgfpathlineto{\pgfqpoint{8.174527in}{3.092996in}}%
\pgfpathlineto{\pgfqpoint{8.174527in}{3.166390in}}%
\pgfpathlineto{\pgfqpoint{8.151396in}{3.161751in}}%
\pgfpathlineto{\pgfqpoint{8.151396in}{3.173606in}}%
\pgfpathlineto{\pgfqpoint{8.174404in}{3.178245in}}%
\pgfpathlineto{\pgfqpoint{8.187413in}{3.178245in}}%
\pgfpathlineto{\pgfqpoint{8.187413in}{3.092996in}}%
\pgfpathlineto{\pgfqpoint{8.208668in}{3.092996in}}%
\pgfpathlineto{\pgfqpoint{8.208668in}{3.082049in}}%
\pgfpathlineto{\pgfqpoint{8.153272in}{3.082049in}}%
\pgfpathlineto{\pgfqpoint{8.153272in}{3.092996in}}%
\pgfpathlineto{\pgfqpoint{8.153272in}{3.092996in}}%
\pgfpathclose%
\pgfpathmoveto{\pgfqpoint{8.246169in}{3.092996in}}%
\pgfpathlineto{\pgfqpoint{8.291587in}{3.092996in}}%
\pgfpathlineto{\pgfqpoint{8.291587in}{3.082049in}}%
\pgfpathlineto{\pgfqpoint{8.230521in}{3.082049in}}%
\pgfpathlineto{\pgfqpoint{8.230521in}{3.092996in}}%
\pgfpathquadraticcurveto{\pgfqpoint{8.237923in}{3.100665in}}{\pgfqpoint{8.250705in}{3.113571in}}%
\pgfpathquadraticcurveto{\pgfqpoint{8.263507in}{3.126498in}}{\pgfqpoint{8.266785in}{3.130250in}}%
\pgfpathquadraticcurveto{\pgfqpoint{8.273032in}{3.137259in}}{\pgfqpoint{8.275506in}{3.142125in}}%
\pgfpathquadraticcurveto{\pgfqpoint{8.278001in}{3.146990in}}{\pgfqpoint{8.278001in}{3.151691in}}%
\pgfpathquadraticcurveto{\pgfqpoint{8.278001in}{3.159360in}}{\pgfqpoint{8.272620in}{3.164184in}}%
\pgfpathquadraticcurveto{\pgfqpoint{8.267239in}{3.169029in}}{\pgfqpoint{8.258601in}{3.169029in}}%
\pgfpathquadraticcurveto{\pgfqpoint{8.252478in}{3.169029in}}{\pgfqpoint{8.245674in}{3.166906in}}%
\pgfpathquadraticcurveto{\pgfqpoint{8.238891in}{3.164782in}}{\pgfqpoint{8.231160in}{3.160453in}}%
\pgfpathlineto{\pgfqpoint{8.231160in}{3.173606in}}%
\pgfpathquadraticcurveto{\pgfqpoint{8.239015in}{3.176760in}}{\pgfqpoint{8.245839in}{3.178368in}}%
\pgfpathquadraticcurveto{\pgfqpoint{8.252684in}{3.179976in}}{\pgfqpoint{8.258353in}{3.179976in}}%
\pgfpathquadraticcurveto{\pgfqpoint{8.273300in}{3.179976in}}{\pgfqpoint{8.282186in}{3.172493in}}%
\pgfpathquadraticcurveto{\pgfqpoint{8.291071in}{3.165029in}}{\pgfqpoint{8.291071in}{3.152536in}}%
\pgfpathquadraticcurveto{\pgfqpoint{8.291071in}{3.146598in}}{\pgfqpoint{8.288845in}{3.141279in}}%
\pgfpathquadraticcurveto{\pgfqpoint{8.286639in}{3.135981in}}{\pgfqpoint{8.280763in}{3.128765in}}%
\pgfpathquadraticcurveto{\pgfqpoint{8.279155in}{3.126889in}}{\pgfqpoint{8.270517in}{3.117962in}}%
\pgfpathquadraticcurveto{\pgfqpoint{8.261899in}{3.109036in}}{\pgfqpoint{8.246169in}{3.092996in}}%
\pgfpathlineto{\pgfqpoint{8.246169in}{3.092996in}}%
\pgfpathclose%
\pgfpathmoveto{\pgfqpoint{8.387638in}{3.182162in}}%
\pgfpathquadraticcurveto{\pgfqpoint{8.379021in}{3.167359in}}{\pgfqpoint{8.374815in}{3.152845in}}%
\pgfpathquadraticcurveto{\pgfqpoint{8.370630in}{3.138352in}}{\pgfqpoint{8.370630in}{3.123467in}}%
\pgfpathquadraticcurveto{\pgfqpoint{8.370630in}{3.108603in}}{\pgfqpoint{8.374856in}{3.094006in}}%
\pgfpathquadraticcurveto{\pgfqpoint{8.379082in}{3.079410in}}{\pgfqpoint{8.387638in}{3.064649in}}%
\pgfpathlineto{\pgfqpoint{8.377330in}{3.064649in}}%
\pgfpathquadraticcurveto{\pgfqpoint{8.367682in}{3.079802in}}{\pgfqpoint{8.362878in}{3.094419in}}%
\pgfpathquadraticcurveto{\pgfqpoint{8.358074in}{3.109036in}}{\pgfqpoint{8.358074in}{3.123467in}}%
\pgfpathquadraticcurveto{\pgfqpoint{8.358074in}{3.137837in}}{\pgfqpoint{8.362837in}{3.152392in}}%
\pgfpathquadraticcurveto{\pgfqpoint{8.367620in}{3.166967in}}{\pgfqpoint{8.377330in}{3.182162in}}%
\pgfpathlineto{\pgfqpoint{8.387638in}{3.182162in}}%
\pgfpathlineto{\pgfqpoint{8.387638in}{3.182162in}}%
\pgfpathclose%
\pgfpathmoveto{\pgfqpoint{8.423531in}{3.092996in}}%
\pgfpathlineto{\pgfqpoint{8.468949in}{3.092996in}}%
\pgfpathlineto{\pgfqpoint{8.468949in}{3.082049in}}%
\pgfpathlineto{\pgfqpoint{8.407883in}{3.082049in}}%
\pgfpathlineto{\pgfqpoint{8.407883in}{3.092996in}}%
\pgfpathquadraticcurveto{\pgfqpoint{8.415285in}{3.100665in}}{\pgfqpoint{8.428067in}{3.113571in}}%
\pgfpathquadraticcurveto{\pgfqpoint{8.440870in}{3.126498in}}{\pgfqpoint{8.444148in}{3.130250in}}%
\pgfpathquadraticcurveto{\pgfqpoint{8.450394in}{3.137259in}}{\pgfqpoint{8.452868in}{3.142125in}}%
\pgfpathquadraticcurveto{\pgfqpoint{8.455363in}{3.146990in}}{\pgfqpoint{8.455363in}{3.151691in}}%
\pgfpathquadraticcurveto{\pgfqpoint{8.455363in}{3.159360in}}{\pgfqpoint{8.449982in}{3.164184in}}%
\pgfpathquadraticcurveto{\pgfqpoint{8.444601in}{3.169029in}}{\pgfqpoint{8.435963in}{3.169029in}}%
\pgfpathquadraticcurveto{\pgfqpoint{8.429840in}{3.169029in}}{\pgfqpoint{8.423036in}{3.166906in}}%
\pgfpathquadraticcurveto{\pgfqpoint{8.416254in}{3.164782in}}{\pgfqpoint{8.408523in}{3.160453in}}%
\pgfpathlineto{\pgfqpoint{8.408523in}{3.173606in}}%
\pgfpathquadraticcurveto{\pgfqpoint{8.416377in}{3.176760in}}{\pgfqpoint{8.423201in}{3.178368in}}%
\pgfpathquadraticcurveto{\pgfqpoint{8.430046in}{3.179976in}}{\pgfqpoint{8.435715in}{3.179976in}}%
\pgfpathquadraticcurveto{\pgfqpoint{8.450662in}{3.179976in}}{\pgfqpoint{8.459548in}{3.172493in}}%
\pgfpathquadraticcurveto{\pgfqpoint{8.468434in}{3.165029in}}{\pgfqpoint{8.468434in}{3.152536in}}%
\pgfpathquadraticcurveto{\pgfqpoint{8.468434in}{3.146598in}}{\pgfqpoint{8.466207in}{3.141279in}}%
\pgfpathquadraticcurveto{\pgfqpoint{8.464001in}{3.135981in}}{\pgfqpoint{8.458125in}{3.128765in}}%
\pgfpathquadraticcurveto{\pgfqpoint{8.456517in}{3.126889in}}{\pgfqpoint{8.447879in}{3.117962in}}%
\pgfpathquadraticcurveto{\pgfqpoint{8.439262in}{3.109036in}}{\pgfqpoint{8.423531in}{3.092996in}}%
\pgfpathlineto{\pgfqpoint{8.423531in}{3.092996in}}%
\pgfpathclose%
\pgfpathmoveto{\pgfqpoint{8.524098in}{3.169668in}}%
\pgfpathquadraticcurveto{\pgfqpoint{8.514058in}{3.169668in}}{\pgfqpoint{8.508986in}{3.159772in}}%
\pgfpathquadraticcurveto{\pgfqpoint{8.503935in}{3.149897in}}{\pgfqpoint{8.503935in}{3.130044in}}%
\pgfpathquadraticcurveto{\pgfqpoint{8.503935in}{3.110273in}}{\pgfqpoint{8.508986in}{3.100377in}}%
\pgfpathquadraticcurveto{\pgfqpoint{8.514058in}{3.090481in}}{\pgfqpoint{8.524098in}{3.090481in}}%
\pgfpathquadraticcurveto{\pgfqpoint{8.534220in}{3.090481in}}{\pgfqpoint{8.539271in}{3.100377in}}%
\pgfpathquadraticcurveto{\pgfqpoint{8.544343in}{3.110273in}}{\pgfqpoint{8.544343in}{3.130044in}}%
\pgfpathquadraticcurveto{\pgfqpoint{8.544343in}{3.149897in}}{\pgfqpoint{8.539271in}{3.159772in}}%
\pgfpathquadraticcurveto{\pgfqpoint{8.534220in}{3.169668in}}{\pgfqpoint{8.524098in}{3.169668in}}%
\pgfpathlineto{\pgfqpoint{8.524098in}{3.169668in}}%
\pgfpathclose%
\pgfpathmoveto{\pgfqpoint{8.524098in}{3.179976in}}%
\pgfpathquadraticcurveto{\pgfqpoint{8.540281in}{3.179976in}}{\pgfqpoint{8.548817in}{3.167174in}}%
\pgfpathquadraticcurveto{\pgfqpoint{8.557352in}{3.154391in}}{\pgfqpoint{8.557352in}{3.130044in}}%
\pgfpathquadraticcurveto{\pgfqpoint{8.557352in}{3.105758in}}{\pgfqpoint{8.548817in}{3.092955in}}%
\pgfpathquadraticcurveto{\pgfqpoint{8.540281in}{3.080173in}}{\pgfqpoint{8.524098in}{3.080173in}}%
\pgfpathquadraticcurveto{\pgfqpoint{8.507934in}{3.080173in}}{\pgfqpoint{8.499399in}{3.092955in}}%
\pgfpathquadraticcurveto{\pgfqpoint{8.490864in}{3.105758in}}{\pgfqpoint{8.490864in}{3.130044in}}%
\pgfpathquadraticcurveto{\pgfqpoint{8.490864in}{3.154391in}}{\pgfqpoint{8.499399in}{3.167174in}}%
\pgfpathquadraticcurveto{\pgfqpoint{8.507934in}{3.179976in}}{\pgfqpoint{8.524098in}{3.179976in}}%
\pgfpathlineto{\pgfqpoint{8.524098in}{3.179976in}}%
\pgfpathclose%
\pgfpathmoveto{\pgfqpoint{8.591431in}{3.092996in}}%
\pgfpathlineto{\pgfqpoint{8.636848in}{3.092996in}}%
\pgfpathlineto{\pgfqpoint{8.636848in}{3.082049in}}%
\pgfpathlineto{\pgfqpoint{8.575783in}{3.082049in}}%
\pgfpathlineto{\pgfqpoint{8.575783in}{3.092996in}}%
\pgfpathquadraticcurveto{\pgfqpoint{8.583184in}{3.100665in}}{\pgfqpoint{8.595966in}{3.113571in}}%
\pgfpathquadraticcurveto{\pgfqpoint{8.608769in}{3.126498in}}{\pgfqpoint{8.612047in}{3.130250in}}%
\pgfpathquadraticcurveto{\pgfqpoint{8.618294in}{3.137259in}}{\pgfqpoint{8.620768in}{3.142125in}}%
\pgfpathquadraticcurveto{\pgfqpoint{8.623262in}{3.146990in}}{\pgfqpoint{8.623262in}{3.151691in}}%
\pgfpathquadraticcurveto{\pgfqpoint{8.623262in}{3.159360in}}{\pgfqpoint{8.617881in}{3.164184in}}%
\pgfpathquadraticcurveto{\pgfqpoint{8.612500in}{3.169029in}}{\pgfqpoint{8.603862in}{3.169029in}}%
\pgfpathquadraticcurveto{\pgfqpoint{8.597739in}{3.169029in}}{\pgfqpoint{8.590936in}{3.166906in}}%
\pgfpathquadraticcurveto{\pgfqpoint{8.584153in}{3.164782in}}{\pgfqpoint{8.576422in}{3.160453in}}%
\pgfpathlineto{\pgfqpoint{8.576422in}{3.173606in}}%
\pgfpathquadraticcurveto{\pgfqpoint{8.584277in}{3.176760in}}{\pgfqpoint{8.591101in}{3.178368in}}%
\pgfpathquadraticcurveto{\pgfqpoint{8.597945in}{3.179976in}}{\pgfqpoint{8.603615in}{3.179976in}}%
\pgfpathquadraticcurveto{\pgfqpoint{8.618562in}{3.179976in}}{\pgfqpoint{8.627447in}{3.172493in}}%
\pgfpathquadraticcurveto{\pgfqpoint{8.636333in}{3.165029in}}{\pgfqpoint{8.636333in}{3.152536in}}%
\pgfpathquadraticcurveto{\pgfqpoint{8.636333in}{3.146598in}}{\pgfqpoint{8.634106in}{3.141279in}}%
\pgfpathquadraticcurveto{\pgfqpoint{8.631900in}{3.135981in}}{\pgfqpoint{8.626025in}{3.128765in}}%
\pgfpathquadraticcurveto{\pgfqpoint{8.624417in}{3.126889in}}{\pgfqpoint{8.615778in}{3.117962in}}%
\pgfpathquadraticcurveto{\pgfqpoint{8.607161in}{3.109036in}}{\pgfqpoint{8.591431in}{3.092996in}}%
\pgfpathlineto{\pgfqpoint{8.591431in}{3.092996in}}%
\pgfpathclose%
\pgfpathmoveto{\pgfqpoint{8.675380in}{3.092996in}}%
\pgfpathlineto{\pgfqpoint{8.720798in}{3.092996in}}%
\pgfpathlineto{\pgfqpoint{8.720798in}{3.082049in}}%
\pgfpathlineto{\pgfqpoint{8.659732in}{3.082049in}}%
\pgfpathlineto{\pgfqpoint{8.659732in}{3.092996in}}%
\pgfpathquadraticcurveto{\pgfqpoint{8.667134in}{3.100665in}}{\pgfqpoint{8.679916in}{3.113571in}}%
\pgfpathquadraticcurveto{\pgfqpoint{8.692719in}{3.126498in}}{\pgfqpoint{8.695997in}{3.130250in}}%
\pgfpathquadraticcurveto{\pgfqpoint{8.702243in}{3.137259in}}{\pgfqpoint{8.704717in}{3.142125in}}%
\pgfpathquadraticcurveto{\pgfqpoint{8.707212in}{3.146990in}}{\pgfqpoint{8.707212in}{3.151691in}}%
\pgfpathquadraticcurveto{\pgfqpoint{8.707212in}{3.159360in}}{\pgfqpoint{8.701831in}{3.164184in}}%
\pgfpathquadraticcurveto{\pgfqpoint{8.696450in}{3.169029in}}{\pgfqpoint{8.687812in}{3.169029in}}%
\pgfpathquadraticcurveto{\pgfqpoint{8.681689in}{3.169029in}}{\pgfqpoint{8.674885in}{3.166906in}}%
\pgfpathquadraticcurveto{\pgfqpoint{8.668103in}{3.164782in}}{\pgfqpoint{8.660372in}{3.160453in}}%
\pgfpathlineto{\pgfqpoint{8.660372in}{3.173606in}}%
\pgfpathquadraticcurveto{\pgfqpoint{8.668226in}{3.176760in}}{\pgfqpoint{8.675050in}{3.178368in}}%
\pgfpathquadraticcurveto{\pgfqpoint{8.681895in}{3.179976in}}{\pgfqpoint{8.687564in}{3.179976in}}%
\pgfpathquadraticcurveto{\pgfqpoint{8.702511in}{3.179976in}}{\pgfqpoint{8.711397in}{3.172493in}}%
\pgfpathquadraticcurveto{\pgfqpoint{8.720283in}{3.165029in}}{\pgfqpoint{8.720283in}{3.152536in}}%
\pgfpathquadraticcurveto{\pgfqpoint{8.720283in}{3.146598in}}{\pgfqpoint{8.718056in}{3.141279in}}%
\pgfpathquadraticcurveto{\pgfqpoint{8.715850in}{3.135981in}}{\pgfqpoint{8.709974in}{3.128765in}}%
\pgfpathquadraticcurveto{\pgfqpoint{8.708366in}{3.126889in}}{\pgfqpoint{8.699728in}{3.117962in}}%
\pgfpathquadraticcurveto{\pgfqpoint{8.691110in}{3.109036in}}{\pgfqpoint{8.675380in}{3.092996in}}%
\pgfpathlineto{\pgfqpoint{8.675380in}{3.092996in}}%
\pgfpathclose%
\pgfpathmoveto{\pgfqpoint{8.744589in}{3.182162in}}%
\pgfpathlineto{\pgfqpoint{8.754897in}{3.182162in}}%
\pgfpathquadraticcurveto{\pgfqpoint{8.764546in}{3.166967in}}{\pgfqpoint{8.769349in}{3.152392in}}%
\pgfpathquadraticcurveto{\pgfqpoint{8.774153in}{3.137837in}}{\pgfqpoint{8.774153in}{3.123467in}}%
\pgfpathquadraticcurveto{\pgfqpoint{8.774153in}{3.109036in}}{\pgfqpoint{8.769349in}{3.094419in}}%
\pgfpathquadraticcurveto{\pgfqpoint{8.764546in}{3.079802in}}{\pgfqpoint{8.754897in}{3.064649in}}%
\pgfpathlineto{\pgfqpoint{8.744589in}{3.064649in}}%
\pgfpathquadraticcurveto{\pgfqpoint{8.753145in}{3.079410in}}{\pgfqpoint{8.757371in}{3.094006in}}%
\pgfpathquadraticcurveto{\pgfqpoint{8.761598in}{3.108603in}}{\pgfqpoint{8.761598in}{3.123467in}}%
\pgfpathquadraticcurveto{\pgfqpoint{8.761598in}{3.138352in}}{\pgfqpoint{8.757371in}{3.152845in}}%
\pgfpathquadraticcurveto{\pgfqpoint{8.753145in}{3.167359in}}{\pgfqpoint{8.744589in}{3.182162in}}%
\pgfpathlineto{\pgfqpoint{8.744589in}{3.182162in}}%
\pgfpathclose%
\pgfusepath{fill}%
\end{pgfscope}%
\begin{pgfscope}%
\pgfsetbuttcap%
\pgfsetroundjoin%
\definecolor{currentfill}{rgb}{0.090196,0.168627,0.227451}%
\pgfsetfillcolor{currentfill}%
\pgfsetlinewidth{0.000000pt}%
\definecolor{currentstroke}{rgb}{0.090196,0.168627,0.227451}%
\pgfsetstrokecolor{currentstroke}%
\pgfsetdash{}{0pt}%
\pgfpathmoveto{\pgfqpoint{7.784611in}{2.775490in}}%
\pgfpathlineto{\pgfqpoint{7.784611in}{2.762791in}}%
\pgfpathquadraticcurveto{\pgfqpoint{7.777210in}{2.766337in}}{\pgfqpoint{7.770633in}{2.768068in}}%
\pgfpathquadraticcurveto{\pgfqpoint{7.764056in}{2.769821in}}{\pgfqpoint{7.757933in}{2.769821in}}%
\pgfpathquadraticcurveto{\pgfqpoint{7.747316in}{2.769821in}}{\pgfqpoint{7.741543in}{2.765697in}}%
\pgfpathquadraticcurveto{\pgfqpoint{7.735771in}{2.761574in}}{\pgfqpoint{7.735771in}{2.753967in}}%
\pgfpathquadraticcurveto{\pgfqpoint{7.735771in}{2.747576in}}{\pgfqpoint{7.739605in}{2.744318in}}%
\pgfpathquadraticcurveto{\pgfqpoint{7.743440in}{2.741082in}}{\pgfqpoint{7.754140in}{2.739082in}}%
\pgfpathlineto{\pgfqpoint{7.761995in}{2.737474in}}%
\pgfpathquadraticcurveto{\pgfqpoint{7.776550in}{2.734691in}}{\pgfqpoint{7.783477in}{2.727702in}}%
\pgfpathquadraticcurveto{\pgfqpoint{7.790404in}{2.720713in}}{\pgfqpoint{7.790404in}{2.709003in}}%
\pgfpathquadraticcurveto{\pgfqpoint{7.790404in}{2.695004in}}{\pgfqpoint{7.781024in}{2.687788in}}%
\pgfpathquadraticcurveto{\pgfqpoint{7.771664in}{2.680573in}}{\pgfqpoint{7.753563in}{2.680573in}}%
\pgfpathquadraticcurveto{\pgfqpoint{7.746739in}{2.680573in}}{\pgfqpoint{7.739028in}{2.682119in}}%
\pgfpathquadraticcurveto{\pgfqpoint{7.731338in}{2.683665in}}{\pgfqpoint{7.723092in}{2.686696in}}%
\pgfpathlineto{\pgfqpoint{7.723092in}{2.700096in}}%
\pgfpathquadraticcurveto{\pgfqpoint{7.731008in}{2.695664in}}{\pgfqpoint{7.738616in}{2.693396in}}%
\pgfpathquadraticcurveto{\pgfqpoint{7.746223in}{2.691149in}}{\pgfqpoint{7.753563in}{2.691149in}}%
\pgfpathquadraticcurveto{\pgfqpoint{7.764696in}{2.691149in}}{\pgfqpoint{7.770757in}{2.695520in}}%
\pgfpathquadraticcurveto{\pgfqpoint{7.776818in}{2.699911in}}{\pgfqpoint{7.776818in}{2.708034in}}%
\pgfpathquadraticcurveto{\pgfqpoint{7.776818in}{2.715105in}}{\pgfqpoint{7.772468in}{2.719105in}}%
\pgfpathquadraticcurveto{\pgfqpoint{7.768118in}{2.723104in}}{\pgfqpoint{7.758201in}{2.725104in}}%
\pgfpathlineto{\pgfqpoint{7.750264in}{2.726650in}}%
\pgfpathquadraticcurveto{\pgfqpoint{7.735709in}{2.729536in}}{\pgfqpoint{7.729194in}{2.735721in}}%
\pgfpathquadraticcurveto{\pgfqpoint{7.722700in}{2.741906in}}{\pgfqpoint{7.722700in}{2.752936in}}%
\pgfpathquadraticcurveto{\pgfqpoint{7.722700in}{2.765697in}}{\pgfqpoint{7.731689in}{2.773037in}}%
\pgfpathquadraticcurveto{\pgfqpoint{7.740678in}{2.780376in}}{\pgfqpoint{7.756449in}{2.780376in}}%
\pgfpathquadraticcurveto{\pgfqpoint{7.763232in}{2.780376in}}{\pgfqpoint{7.770241in}{2.779139in}}%
\pgfpathquadraticcurveto{\pgfqpoint{7.777271in}{2.777923in}}{\pgfqpoint{7.784611in}{2.775490in}}%
\pgfpathlineto{\pgfqpoint{7.784611in}{2.775490in}}%
\pgfpathclose%
\pgfpathmoveto{\pgfqpoint{7.821926in}{2.775099in}}%
\pgfpathlineto{\pgfqpoint{7.821926in}{2.754606in}}%
\pgfpathlineto{\pgfqpoint{7.846336in}{2.754606in}}%
\pgfpathlineto{\pgfqpoint{7.846336in}{2.745390in}}%
\pgfpathlineto{\pgfqpoint{7.821926in}{2.745390in}}%
\pgfpathlineto{\pgfqpoint{7.821926in}{2.706219in}}%
\pgfpathquadraticcurveto{\pgfqpoint{7.821926in}{2.697396in}}{\pgfqpoint{7.824339in}{2.694880in}}%
\pgfpathquadraticcurveto{\pgfqpoint{7.826751in}{2.692365in}}{\pgfqpoint{7.834173in}{2.692365in}}%
\pgfpathlineto{\pgfqpoint{7.846336in}{2.692365in}}%
\pgfpathlineto{\pgfqpoint{7.846336in}{2.682449in}}%
\pgfpathlineto{\pgfqpoint{7.834173in}{2.682449in}}%
\pgfpathquadraticcurveto{\pgfqpoint{7.820442in}{2.682449in}}{\pgfqpoint{7.815226in}{2.687562in}}%
\pgfpathquadraticcurveto{\pgfqpoint{7.810010in}{2.692695in}}{\pgfqpoint{7.810010in}{2.706219in}}%
\pgfpathlineto{\pgfqpoint{7.810010in}{2.745390in}}%
\pgfpathlineto{\pgfqpoint{7.801310in}{2.745390in}}%
\pgfpathlineto{\pgfqpoint{7.801310in}{2.754606in}}%
\pgfpathlineto{\pgfqpoint{7.810010in}{2.754606in}}%
\pgfpathlineto{\pgfqpoint{7.810010in}{2.775099in}}%
\pgfpathlineto{\pgfqpoint{7.821926in}{2.775099in}}%
\pgfpathlineto{\pgfqpoint{7.821926in}{2.775099in}}%
\pgfpathclose%
\pgfpathmoveto{\pgfqpoint{7.860706in}{2.710920in}}%
\pgfpathlineto{\pgfqpoint{7.860706in}{2.754606in}}%
\pgfpathlineto{\pgfqpoint{7.872560in}{2.754606in}}%
\pgfpathlineto{\pgfqpoint{7.872560in}{2.711373in}}%
\pgfpathquadraticcurveto{\pgfqpoint{7.872560in}{2.701127in}}{\pgfqpoint{7.876539in}{2.695994in}}%
\pgfpathquadraticcurveto{\pgfqpoint{7.880539in}{2.690881in}}{\pgfqpoint{7.888538in}{2.690881in}}%
\pgfpathquadraticcurveto{\pgfqpoint{7.898124in}{2.690881in}}{\pgfqpoint{7.903691in}{2.697004in}}%
\pgfpathquadraticcurveto{\pgfqpoint{7.909278in}{2.703127in}}{\pgfqpoint{7.909278in}{2.713703in}}%
\pgfpathlineto{\pgfqpoint{7.909278in}{2.754606in}}%
\pgfpathlineto{\pgfqpoint{7.921132in}{2.754606in}}%
\pgfpathlineto{\pgfqpoint{7.921132in}{2.682449in}}%
\pgfpathlineto{\pgfqpoint{7.909278in}{2.682449in}}%
\pgfpathlineto{\pgfqpoint{7.909278in}{2.693540in}}%
\pgfpathquadraticcurveto{\pgfqpoint{7.904969in}{2.686964in}}{\pgfqpoint{7.899258in}{2.683768in}}%
\pgfpathquadraticcurveto{\pgfqpoint{7.893568in}{2.680573in}}{\pgfqpoint{7.886023in}{2.680573in}}%
\pgfpathquadraticcurveto{\pgfqpoint{7.873591in}{2.680573in}}{\pgfqpoint{7.867138in}{2.688304in}}%
\pgfpathquadraticcurveto{\pgfqpoint{7.860706in}{2.696035in}}{\pgfqpoint{7.860706in}{2.710920in}}%
\pgfpathlineto{\pgfqpoint{7.860706in}{2.710920in}}%
\pgfpathclose%
\pgfpathmoveto{\pgfqpoint{7.890538in}{2.756338in}}%
\pgfpathlineto{\pgfqpoint{7.890538in}{2.756338in}}%
\pgfpathlineto{\pgfqpoint{7.890538in}{2.756338in}}%
\pgfpathclose%
\pgfpathmoveto{\pgfqpoint{7.993021in}{2.743659in}}%
\pgfpathlineto{\pgfqpoint{7.993021in}{2.782706in}}%
\pgfpathlineto{\pgfqpoint{8.004876in}{2.782706in}}%
\pgfpathlineto{\pgfqpoint{8.004876in}{2.682449in}}%
\pgfpathlineto{\pgfqpoint{7.993021in}{2.682449in}}%
\pgfpathlineto{\pgfqpoint{7.993021in}{2.693272in}}%
\pgfpathquadraticcurveto{\pgfqpoint{7.989290in}{2.686840in}}{\pgfqpoint{7.983579in}{2.683706in}}%
\pgfpathquadraticcurveto{\pgfqpoint{7.977889in}{2.680573in}}{\pgfqpoint{7.969890in}{2.680573in}}%
\pgfpathquadraticcurveto{\pgfqpoint{7.956819in}{2.680573in}}{\pgfqpoint{7.948593in}{2.691005in}}%
\pgfpathquadraticcurveto{\pgfqpoint{7.940388in}{2.701457in}}{\pgfqpoint{7.940388in}{2.718466in}}%
\pgfpathquadraticcurveto{\pgfqpoint{7.940388in}{2.735474in}}{\pgfqpoint{7.948593in}{2.745906in}}%
\pgfpathquadraticcurveto{\pgfqpoint{7.956819in}{2.756338in}}{\pgfqpoint{7.969890in}{2.756338in}}%
\pgfpathquadraticcurveto{\pgfqpoint{7.977889in}{2.756338in}}{\pgfqpoint{7.983579in}{2.753204in}}%
\pgfpathquadraticcurveto{\pgfqpoint{7.989290in}{2.750091in}}{\pgfqpoint{7.993021in}{2.743659in}}%
\pgfpathlineto{\pgfqpoint{7.993021in}{2.743659in}}%
\pgfpathclose%
\pgfpathmoveto{\pgfqpoint{7.952634in}{2.718466in}}%
\pgfpathquadraticcurveto{\pgfqpoint{7.952634in}{2.705395in}}{\pgfqpoint{7.958015in}{2.697952in}}%
\pgfpathquadraticcurveto{\pgfqpoint{7.963396in}{2.690510in}}{\pgfqpoint{7.972797in}{2.690510in}}%
\pgfpathquadraticcurveto{\pgfqpoint{7.982198in}{2.690510in}}{\pgfqpoint{7.987599in}{2.697952in}}%
\pgfpathquadraticcurveto{\pgfqpoint{7.993021in}{2.705395in}}{\pgfqpoint{7.993021in}{2.718466in}}%
\pgfpathquadraticcurveto{\pgfqpoint{7.993021in}{2.731536in}}{\pgfqpoint{7.987599in}{2.738979in}}%
\pgfpathquadraticcurveto{\pgfqpoint{7.982198in}{2.746421in}}{\pgfqpoint{7.972797in}{2.746421in}}%
\pgfpathquadraticcurveto{\pgfqpoint{7.963396in}{2.746421in}}{\pgfqpoint{7.958015in}{2.738979in}}%
\pgfpathquadraticcurveto{\pgfqpoint{7.952634in}{2.731536in}}{\pgfqpoint{7.952634in}{2.718466in}}%
\pgfpathlineto{\pgfqpoint{7.952634in}{2.718466in}}%
\pgfpathclose%
\pgfpathmoveto{\pgfqpoint{8.059323in}{2.675748in}}%
\pgfpathquadraticcurveto{\pgfqpoint{8.054314in}{2.662863in}}{\pgfqpoint{8.049531in}{2.658946in}}%
\pgfpathquadraticcurveto{\pgfqpoint{8.044768in}{2.655008in}}{\pgfqpoint{8.036790in}{2.655008in}}%
\pgfpathlineto{\pgfqpoint{8.027306in}{2.655008in}}%
\pgfpathlineto{\pgfqpoint{8.027306in}{2.664925in}}%
\pgfpathlineto{\pgfqpoint{8.034275in}{2.664925in}}%
\pgfpathquadraticcurveto{\pgfqpoint{8.039161in}{2.664925in}}{\pgfqpoint{8.041861in}{2.667255in}}%
\pgfpathquadraticcurveto{\pgfqpoint{8.044583in}{2.669564in}}{\pgfqpoint{8.047861in}{2.678202in}}%
\pgfpathlineto{\pgfqpoint{8.049984in}{2.683603in}}%
\pgfpathlineto{\pgfqpoint{8.020812in}{2.754606in}}%
\pgfpathlineto{\pgfqpoint{8.033367in}{2.754606in}}%
\pgfpathlineto{\pgfqpoint{8.055922in}{2.698179in}}%
\pgfpathlineto{\pgfqpoint{8.078476in}{2.754606in}}%
\pgfpathlineto{\pgfqpoint{8.091031in}{2.754606in}}%
\pgfpathlineto{\pgfqpoint{8.059323in}{2.675748in}}%
\pgfpathlineto{\pgfqpoint{8.059323in}{2.675748in}}%
\pgfpathclose%
\pgfpathmoveto{\pgfqpoint{8.153272in}{2.693396in}}%
\pgfpathlineto{\pgfqpoint{8.174527in}{2.693396in}}%
\pgfpathlineto{\pgfqpoint{8.174527in}{2.766790in}}%
\pgfpathlineto{\pgfqpoint{8.151396in}{2.762151in}}%
\pgfpathlineto{\pgfqpoint{8.151396in}{2.774006in}}%
\pgfpathlineto{\pgfqpoint{8.174404in}{2.778645in}}%
\pgfpathlineto{\pgfqpoint{8.187413in}{2.778645in}}%
\pgfpathlineto{\pgfqpoint{8.187413in}{2.693396in}}%
\pgfpathlineto{\pgfqpoint{8.208668in}{2.693396in}}%
\pgfpathlineto{\pgfqpoint{8.208668in}{2.682449in}}%
\pgfpathlineto{\pgfqpoint{8.153272in}{2.682449in}}%
\pgfpathlineto{\pgfqpoint{8.153272in}{2.693396in}}%
\pgfpathlineto{\pgfqpoint{8.153272in}{2.693396in}}%
\pgfpathclose%
\pgfpathmoveto{\pgfqpoint{8.274393in}{2.734319in}}%
\pgfpathquadraticcurveto{\pgfqpoint{8.283732in}{2.732320in}}{\pgfqpoint{8.288969in}{2.725990in}}%
\pgfpathquadraticcurveto{\pgfqpoint{8.294226in}{2.719682in}}{\pgfqpoint{8.294226in}{2.710405in}}%
\pgfpathquadraticcurveto{\pgfqpoint{8.294226in}{2.696179in}}{\pgfqpoint{8.284433in}{2.688366in}}%
\pgfpathquadraticcurveto{\pgfqpoint{8.274640in}{2.680573in}}{\pgfqpoint{8.256601in}{2.680573in}}%
\pgfpathquadraticcurveto{\pgfqpoint{8.250560in}{2.680573in}}{\pgfqpoint{8.244149in}{2.681768in}}%
\pgfpathquadraticcurveto{\pgfqpoint{8.237737in}{2.682964in}}{\pgfqpoint{8.230913in}{2.685356in}}%
\pgfpathlineto{\pgfqpoint{8.230913in}{2.697911in}}%
\pgfpathquadraticcurveto{\pgfqpoint{8.236314in}{2.694757in}}{\pgfqpoint{8.242747in}{2.693149in}}%
\pgfpathquadraticcurveto{\pgfqpoint{8.249200in}{2.691541in}}{\pgfqpoint{8.256230in}{2.691541in}}%
\pgfpathquadraticcurveto{\pgfqpoint{8.268455in}{2.691541in}}{\pgfqpoint{8.274867in}{2.696365in}}%
\pgfpathquadraticcurveto{\pgfqpoint{8.281279in}{2.701189in}}{\pgfqpoint{8.281279in}{2.710405in}}%
\pgfpathquadraticcurveto{\pgfqpoint{8.281279in}{2.718919in}}{\pgfqpoint{8.275321in}{2.723702in}}%
\pgfpathquadraticcurveto{\pgfqpoint{8.269362in}{2.728506in}}{\pgfqpoint{8.258745in}{2.728506in}}%
\pgfpathlineto{\pgfqpoint{8.247530in}{2.728506in}}%
\pgfpathlineto{\pgfqpoint{8.247530in}{2.739206in}}%
\pgfpathlineto{\pgfqpoint{8.259260in}{2.739206in}}%
\pgfpathquadraticcurveto{\pgfqpoint{8.268847in}{2.739206in}}{\pgfqpoint{8.273939in}{2.743040in}}%
\pgfpathquadraticcurveto{\pgfqpoint{8.279031in}{2.746875in}}{\pgfqpoint{8.279031in}{2.754091in}}%
\pgfpathquadraticcurveto{\pgfqpoint{8.279031in}{2.761492in}}{\pgfqpoint{8.273774in}{2.765450in}}%
\pgfpathquadraticcurveto{\pgfqpoint{8.268538in}{2.769429in}}{\pgfqpoint{8.258745in}{2.769429in}}%
\pgfpathquadraticcurveto{\pgfqpoint{8.253385in}{2.769429in}}{\pgfqpoint{8.247262in}{2.768254in}}%
\pgfpathquadraticcurveto{\pgfqpoint{8.241139in}{2.767099in}}{\pgfqpoint{8.233799in}{2.764667in}}%
\pgfpathlineto{\pgfqpoint{8.233799in}{2.776253in}}%
\pgfpathquadraticcurveto{\pgfqpoint{8.241221in}{2.778315in}}{\pgfqpoint{8.247695in}{2.779345in}}%
\pgfpathquadraticcurveto{\pgfqpoint{8.254168in}{2.780376in}}{\pgfqpoint{8.259900in}{2.780376in}}%
\pgfpathquadraticcurveto{\pgfqpoint{8.274723in}{2.780376in}}{\pgfqpoint{8.283340in}{2.773635in}}%
\pgfpathquadraticcurveto{\pgfqpoint{8.291979in}{2.766914in}}{\pgfqpoint{8.291979in}{2.755451in}}%
\pgfpathquadraticcurveto{\pgfqpoint{8.291979in}{2.747452in}}{\pgfqpoint{8.287402in}{2.741947in}}%
\pgfpathquadraticcurveto{\pgfqpoint{8.282825in}{2.736443in}}{\pgfqpoint{8.274393in}{2.734319in}}%
\pgfpathlineto{\pgfqpoint{8.274393in}{2.734319in}}%
\pgfpathclose%
\pgfpathmoveto{\pgfqpoint{8.387638in}{2.782562in}}%
\pgfpathquadraticcurveto{\pgfqpoint{8.379021in}{2.767759in}}{\pgfqpoint{8.374815in}{2.753245in}}%
\pgfpathquadraticcurveto{\pgfqpoint{8.370630in}{2.738752in}}{\pgfqpoint{8.370630in}{2.723867in}}%
\pgfpathquadraticcurveto{\pgfqpoint{8.370630in}{2.709003in}}{\pgfqpoint{8.374856in}{2.694406in}}%
\pgfpathquadraticcurveto{\pgfqpoint{8.379082in}{2.679810in}}{\pgfqpoint{8.387638in}{2.665049in}}%
\pgfpathlineto{\pgfqpoint{8.377330in}{2.665049in}}%
\pgfpathquadraticcurveto{\pgfqpoint{8.367682in}{2.680202in}}{\pgfqpoint{8.362878in}{2.694819in}}%
\pgfpathquadraticcurveto{\pgfqpoint{8.358074in}{2.709436in}}{\pgfqpoint{8.358074in}{2.723867in}}%
\pgfpathquadraticcurveto{\pgfqpoint{8.358074in}{2.738237in}}{\pgfqpoint{8.362837in}{2.752792in}}%
\pgfpathquadraticcurveto{\pgfqpoint{8.367620in}{2.767367in}}{\pgfqpoint{8.377330in}{2.782562in}}%
\pgfpathlineto{\pgfqpoint{8.387638in}{2.782562in}}%
\pgfpathlineto{\pgfqpoint{8.387638in}{2.782562in}}%
\pgfpathclose%
\pgfpathmoveto{\pgfqpoint{8.423531in}{2.693396in}}%
\pgfpathlineto{\pgfqpoint{8.468949in}{2.693396in}}%
\pgfpathlineto{\pgfqpoint{8.468949in}{2.682449in}}%
\pgfpathlineto{\pgfqpoint{8.407883in}{2.682449in}}%
\pgfpathlineto{\pgfqpoint{8.407883in}{2.693396in}}%
\pgfpathquadraticcurveto{\pgfqpoint{8.415285in}{2.701065in}}{\pgfqpoint{8.428067in}{2.713971in}}%
\pgfpathquadraticcurveto{\pgfqpoint{8.440870in}{2.726898in}}{\pgfqpoint{8.444148in}{2.730650in}}%
\pgfpathquadraticcurveto{\pgfqpoint{8.450394in}{2.737659in}}{\pgfqpoint{8.452868in}{2.742525in}}%
\pgfpathquadraticcurveto{\pgfqpoint{8.455363in}{2.747390in}}{\pgfqpoint{8.455363in}{2.752091in}}%
\pgfpathquadraticcurveto{\pgfqpoint{8.455363in}{2.759760in}}{\pgfqpoint{8.449982in}{2.764584in}}%
\pgfpathquadraticcurveto{\pgfqpoint{8.444601in}{2.769429in}}{\pgfqpoint{8.435963in}{2.769429in}}%
\pgfpathquadraticcurveto{\pgfqpoint{8.429840in}{2.769429in}}{\pgfqpoint{8.423036in}{2.767306in}}%
\pgfpathquadraticcurveto{\pgfqpoint{8.416254in}{2.765182in}}{\pgfqpoint{8.408523in}{2.760853in}}%
\pgfpathlineto{\pgfqpoint{8.408523in}{2.774006in}}%
\pgfpathquadraticcurveto{\pgfqpoint{8.416377in}{2.777160in}}{\pgfqpoint{8.423201in}{2.778768in}}%
\pgfpathquadraticcurveto{\pgfqpoint{8.430046in}{2.780376in}}{\pgfqpoint{8.435715in}{2.780376in}}%
\pgfpathquadraticcurveto{\pgfqpoint{8.450662in}{2.780376in}}{\pgfqpoint{8.459548in}{2.772893in}}%
\pgfpathquadraticcurveto{\pgfqpoint{8.468434in}{2.765429in}}{\pgfqpoint{8.468434in}{2.752936in}}%
\pgfpathquadraticcurveto{\pgfqpoint{8.468434in}{2.746998in}}{\pgfqpoint{8.466207in}{2.741679in}}%
\pgfpathquadraticcurveto{\pgfqpoint{8.464001in}{2.736381in}}{\pgfqpoint{8.458125in}{2.729165in}}%
\pgfpathquadraticcurveto{\pgfqpoint{8.456517in}{2.727289in}}{\pgfqpoint{8.447879in}{2.718362in}}%
\pgfpathquadraticcurveto{\pgfqpoint{8.439262in}{2.709436in}}{\pgfqpoint{8.423531in}{2.693396in}}%
\pgfpathlineto{\pgfqpoint{8.423531in}{2.693396in}}%
\pgfpathclose%
\pgfpathmoveto{\pgfqpoint{8.524098in}{2.770068in}}%
\pgfpathquadraticcurveto{\pgfqpoint{8.514058in}{2.770068in}}{\pgfqpoint{8.508986in}{2.760172in}}%
\pgfpathquadraticcurveto{\pgfqpoint{8.503935in}{2.750297in}}{\pgfqpoint{8.503935in}{2.730444in}}%
\pgfpathquadraticcurveto{\pgfqpoint{8.503935in}{2.710673in}}{\pgfqpoint{8.508986in}{2.700777in}}%
\pgfpathquadraticcurveto{\pgfqpoint{8.514058in}{2.690881in}}{\pgfqpoint{8.524098in}{2.690881in}}%
\pgfpathquadraticcurveto{\pgfqpoint{8.534220in}{2.690881in}}{\pgfqpoint{8.539271in}{2.700777in}}%
\pgfpathquadraticcurveto{\pgfqpoint{8.544343in}{2.710673in}}{\pgfqpoint{8.544343in}{2.730444in}}%
\pgfpathquadraticcurveto{\pgfqpoint{8.544343in}{2.750297in}}{\pgfqpoint{8.539271in}{2.760172in}}%
\pgfpathquadraticcurveto{\pgfqpoint{8.534220in}{2.770068in}}{\pgfqpoint{8.524098in}{2.770068in}}%
\pgfpathlineto{\pgfqpoint{8.524098in}{2.770068in}}%
\pgfpathclose%
\pgfpathmoveto{\pgfqpoint{8.524098in}{2.780376in}}%
\pgfpathquadraticcurveto{\pgfqpoint{8.540281in}{2.780376in}}{\pgfqpoint{8.548817in}{2.767574in}}%
\pgfpathquadraticcurveto{\pgfqpoint{8.557352in}{2.754791in}}{\pgfqpoint{8.557352in}{2.730444in}}%
\pgfpathquadraticcurveto{\pgfqpoint{8.557352in}{2.706158in}}{\pgfqpoint{8.548817in}{2.693355in}}%
\pgfpathquadraticcurveto{\pgfqpoint{8.540281in}{2.680573in}}{\pgfqpoint{8.524098in}{2.680573in}}%
\pgfpathquadraticcurveto{\pgfqpoint{8.507934in}{2.680573in}}{\pgfqpoint{8.499399in}{2.693355in}}%
\pgfpathquadraticcurveto{\pgfqpoint{8.490864in}{2.706158in}}{\pgfqpoint{8.490864in}{2.730444in}}%
\pgfpathquadraticcurveto{\pgfqpoint{8.490864in}{2.754791in}}{\pgfqpoint{8.499399in}{2.767574in}}%
\pgfpathquadraticcurveto{\pgfqpoint{8.507934in}{2.780376in}}{\pgfqpoint{8.524098in}{2.780376in}}%
\pgfpathlineto{\pgfqpoint{8.524098in}{2.780376in}}%
\pgfpathclose%
\pgfpathmoveto{\pgfqpoint{8.591431in}{2.693396in}}%
\pgfpathlineto{\pgfqpoint{8.636848in}{2.693396in}}%
\pgfpathlineto{\pgfqpoint{8.636848in}{2.682449in}}%
\pgfpathlineto{\pgfqpoint{8.575783in}{2.682449in}}%
\pgfpathlineto{\pgfqpoint{8.575783in}{2.693396in}}%
\pgfpathquadraticcurveto{\pgfqpoint{8.583184in}{2.701065in}}{\pgfqpoint{8.595966in}{2.713971in}}%
\pgfpathquadraticcurveto{\pgfqpoint{8.608769in}{2.726898in}}{\pgfqpoint{8.612047in}{2.730650in}}%
\pgfpathquadraticcurveto{\pgfqpoint{8.618294in}{2.737659in}}{\pgfqpoint{8.620768in}{2.742525in}}%
\pgfpathquadraticcurveto{\pgfqpoint{8.623262in}{2.747390in}}{\pgfqpoint{8.623262in}{2.752091in}}%
\pgfpathquadraticcurveto{\pgfqpoint{8.623262in}{2.759760in}}{\pgfqpoint{8.617881in}{2.764584in}}%
\pgfpathquadraticcurveto{\pgfqpoint{8.612500in}{2.769429in}}{\pgfqpoint{8.603862in}{2.769429in}}%
\pgfpathquadraticcurveto{\pgfqpoint{8.597739in}{2.769429in}}{\pgfqpoint{8.590936in}{2.767306in}}%
\pgfpathquadraticcurveto{\pgfqpoint{8.584153in}{2.765182in}}{\pgfqpoint{8.576422in}{2.760853in}}%
\pgfpathlineto{\pgfqpoint{8.576422in}{2.774006in}}%
\pgfpathquadraticcurveto{\pgfqpoint{8.584277in}{2.777160in}}{\pgfqpoint{8.591101in}{2.778768in}}%
\pgfpathquadraticcurveto{\pgfqpoint{8.597945in}{2.780376in}}{\pgfqpoint{8.603615in}{2.780376in}}%
\pgfpathquadraticcurveto{\pgfqpoint{8.618562in}{2.780376in}}{\pgfqpoint{8.627447in}{2.772893in}}%
\pgfpathquadraticcurveto{\pgfqpoint{8.636333in}{2.765429in}}{\pgfqpoint{8.636333in}{2.752936in}}%
\pgfpathquadraticcurveto{\pgfqpoint{8.636333in}{2.746998in}}{\pgfqpoint{8.634106in}{2.741679in}}%
\pgfpathquadraticcurveto{\pgfqpoint{8.631900in}{2.736381in}}{\pgfqpoint{8.626025in}{2.729165in}}%
\pgfpathquadraticcurveto{\pgfqpoint{8.624417in}{2.727289in}}{\pgfqpoint{8.615778in}{2.718362in}}%
\pgfpathquadraticcurveto{\pgfqpoint{8.607161in}{2.709436in}}{\pgfqpoint{8.591431in}{2.693396in}}%
\pgfpathlineto{\pgfqpoint{8.591431in}{2.693396in}}%
\pgfpathclose%
\pgfpathmoveto{\pgfqpoint{8.703604in}{2.734319in}}%
\pgfpathquadraticcurveto{\pgfqpoint{8.712943in}{2.732320in}}{\pgfqpoint{8.718180in}{2.725990in}}%
\pgfpathquadraticcurveto{\pgfqpoint{8.723437in}{2.719682in}}{\pgfqpoint{8.723437in}{2.710405in}}%
\pgfpathquadraticcurveto{\pgfqpoint{8.723437in}{2.696179in}}{\pgfqpoint{8.713644in}{2.688366in}}%
\pgfpathquadraticcurveto{\pgfqpoint{8.703851in}{2.680573in}}{\pgfqpoint{8.685812in}{2.680573in}}%
\pgfpathquadraticcurveto{\pgfqpoint{8.679771in}{2.680573in}}{\pgfqpoint{8.673360in}{2.681768in}}%
\pgfpathquadraticcurveto{\pgfqpoint{8.666948in}{2.682964in}}{\pgfqpoint{8.660124in}{2.685356in}}%
\pgfpathlineto{\pgfqpoint{8.660124in}{2.697911in}}%
\pgfpathquadraticcurveto{\pgfqpoint{8.665526in}{2.694757in}}{\pgfqpoint{8.671958in}{2.693149in}}%
\pgfpathquadraticcurveto{\pgfqpoint{8.678411in}{2.691541in}}{\pgfqpoint{8.685441in}{2.691541in}}%
\pgfpathquadraticcurveto{\pgfqpoint{8.697666in}{2.691541in}}{\pgfqpoint{8.704078in}{2.696365in}}%
\pgfpathquadraticcurveto{\pgfqpoint{8.710490in}{2.701189in}}{\pgfqpoint{8.710490in}{2.710405in}}%
\pgfpathquadraticcurveto{\pgfqpoint{8.710490in}{2.718919in}}{\pgfqpoint{8.704532in}{2.723702in}}%
\pgfpathquadraticcurveto{\pgfqpoint{8.698574in}{2.728506in}}{\pgfqpoint{8.687956in}{2.728506in}}%
\pgfpathlineto{\pgfqpoint{8.676741in}{2.728506in}}%
\pgfpathlineto{\pgfqpoint{8.676741in}{2.739206in}}%
\pgfpathlineto{\pgfqpoint{8.688472in}{2.739206in}}%
\pgfpathquadraticcurveto{\pgfqpoint{8.698058in}{2.739206in}}{\pgfqpoint{8.703150in}{2.743040in}}%
\pgfpathquadraticcurveto{\pgfqpoint{8.708243in}{2.746875in}}{\pgfqpoint{8.708243in}{2.754091in}}%
\pgfpathquadraticcurveto{\pgfqpoint{8.708243in}{2.761492in}}{\pgfqpoint{8.702985in}{2.765450in}}%
\pgfpathquadraticcurveto{\pgfqpoint{8.697749in}{2.769429in}}{\pgfqpoint{8.687956in}{2.769429in}}%
\pgfpathquadraticcurveto{\pgfqpoint{8.682596in}{2.769429in}}{\pgfqpoint{8.676473in}{2.768254in}}%
\pgfpathquadraticcurveto{\pgfqpoint{8.670350in}{2.767099in}}{\pgfqpoint{8.663010in}{2.764667in}}%
\pgfpathlineto{\pgfqpoint{8.663010in}{2.776253in}}%
\pgfpathquadraticcurveto{\pgfqpoint{8.670432in}{2.778315in}}{\pgfqpoint{8.676906in}{2.779345in}}%
\pgfpathquadraticcurveto{\pgfqpoint{8.683379in}{2.780376in}}{\pgfqpoint{8.689111in}{2.780376in}}%
\pgfpathquadraticcurveto{\pgfqpoint{8.703934in}{2.780376in}}{\pgfqpoint{8.712551in}{2.773635in}}%
\pgfpathquadraticcurveto{\pgfqpoint{8.721190in}{2.766914in}}{\pgfqpoint{8.721190in}{2.755451in}}%
\pgfpathquadraticcurveto{\pgfqpoint{8.721190in}{2.747452in}}{\pgfqpoint{8.716613in}{2.741947in}}%
\pgfpathquadraticcurveto{\pgfqpoint{8.712036in}{2.736443in}}{\pgfqpoint{8.703604in}{2.734319in}}%
\pgfpathlineto{\pgfqpoint{8.703604in}{2.734319in}}%
\pgfpathclose%
\pgfpathmoveto{\pgfqpoint{8.744589in}{2.782562in}}%
\pgfpathlineto{\pgfqpoint{8.754897in}{2.782562in}}%
\pgfpathquadraticcurveto{\pgfqpoint{8.764546in}{2.767367in}}{\pgfqpoint{8.769349in}{2.752792in}}%
\pgfpathquadraticcurveto{\pgfqpoint{8.774153in}{2.738237in}}{\pgfqpoint{8.774153in}{2.723867in}}%
\pgfpathquadraticcurveto{\pgfqpoint{8.774153in}{2.709436in}}{\pgfqpoint{8.769349in}{2.694819in}}%
\pgfpathquadraticcurveto{\pgfqpoint{8.764546in}{2.680202in}}{\pgfqpoint{8.754897in}{2.665049in}}%
\pgfpathlineto{\pgfqpoint{8.744589in}{2.665049in}}%
\pgfpathquadraticcurveto{\pgfqpoint{8.753145in}{2.679810in}}{\pgfqpoint{8.757371in}{2.694406in}}%
\pgfpathquadraticcurveto{\pgfqpoint{8.761598in}{2.709003in}}{\pgfqpoint{8.761598in}{2.723867in}}%
\pgfpathquadraticcurveto{\pgfqpoint{8.761598in}{2.738752in}}{\pgfqpoint{8.757371in}{2.753245in}}%
\pgfpathquadraticcurveto{\pgfqpoint{8.753145in}{2.767759in}}{\pgfqpoint{8.744589in}{2.782562in}}%
\pgfpathlineto{\pgfqpoint{8.744589in}{2.782562in}}%
\pgfpathclose%
\pgfusepath{fill}%
\end{pgfscope}%
\begin{pgfscope}%
\pgfsetbuttcap%
\pgfsetroundjoin%
\definecolor{currentfill}{rgb}{0.090196,0.168627,0.227451}%
\pgfsetfillcolor{currentfill}%
\pgfsetlinewidth{0.000000pt}%
\definecolor{currentstroke}{rgb}{0.090196,0.168627,0.227451}%
\pgfsetstrokecolor{currentstroke}%
\pgfsetdash{}{0pt}%
\pgfpathmoveto{\pgfqpoint{7.784611in}{2.375890in}}%
\pgfpathlineto{\pgfqpoint{7.784611in}{2.363191in}}%
\pgfpathquadraticcurveto{\pgfqpoint{7.777210in}{2.366737in}}{\pgfqpoint{7.770633in}{2.368468in}}%
\pgfpathquadraticcurveto{\pgfqpoint{7.764056in}{2.370221in}}{\pgfqpoint{7.757933in}{2.370221in}}%
\pgfpathquadraticcurveto{\pgfqpoint{7.747316in}{2.370221in}}{\pgfqpoint{7.741543in}{2.366097in}}%
\pgfpathquadraticcurveto{\pgfqpoint{7.735771in}{2.361974in}}{\pgfqpoint{7.735771in}{2.354367in}}%
\pgfpathquadraticcurveto{\pgfqpoint{7.735771in}{2.347976in}}{\pgfqpoint{7.739605in}{2.344718in}}%
\pgfpathquadraticcurveto{\pgfqpoint{7.743440in}{2.341482in}}{\pgfqpoint{7.754140in}{2.339482in}}%
\pgfpathlineto{\pgfqpoint{7.761995in}{2.337874in}}%
\pgfpathquadraticcurveto{\pgfqpoint{7.776550in}{2.335091in}}{\pgfqpoint{7.783477in}{2.328102in}}%
\pgfpathquadraticcurveto{\pgfqpoint{7.790404in}{2.321113in}}{\pgfqpoint{7.790404in}{2.309403in}}%
\pgfpathquadraticcurveto{\pgfqpoint{7.790404in}{2.295404in}}{\pgfqpoint{7.781024in}{2.288188in}}%
\pgfpathquadraticcurveto{\pgfqpoint{7.771664in}{2.280973in}}{\pgfqpoint{7.753563in}{2.280973in}}%
\pgfpathquadraticcurveto{\pgfqpoint{7.746739in}{2.280973in}}{\pgfqpoint{7.739028in}{2.282519in}}%
\pgfpathquadraticcurveto{\pgfqpoint{7.731338in}{2.284065in}}{\pgfqpoint{7.723092in}{2.287096in}}%
\pgfpathlineto{\pgfqpoint{7.723092in}{2.300496in}}%
\pgfpathquadraticcurveto{\pgfqpoint{7.731008in}{2.296064in}}{\pgfqpoint{7.738616in}{2.293796in}}%
\pgfpathquadraticcurveto{\pgfqpoint{7.746223in}{2.291549in}}{\pgfqpoint{7.753563in}{2.291549in}}%
\pgfpathquadraticcurveto{\pgfqpoint{7.764696in}{2.291549in}}{\pgfqpoint{7.770757in}{2.295920in}}%
\pgfpathquadraticcurveto{\pgfqpoint{7.776818in}{2.300311in}}{\pgfqpoint{7.776818in}{2.308434in}}%
\pgfpathquadraticcurveto{\pgfqpoint{7.776818in}{2.315505in}}{\pgfqpoint{7.772468in}{2.319505in}}%
\pgfpathquadraticcurveto{\pgfqpoint{7.768118in}{2.323504in}}{\pgfqpoint{7.758201in}{2.325504in}}%
\pgfpathlineto{\pgfqpoint{7.750264in}{2.327050in}}%
\pgfpathquadraticcurveto{\pgfqpoint{7.735709in}{2.329936in}}{\pgfqpoint{7.729194in}{2.336121in}}%
\pgfpathquadraticcurveto{\pgfqpoint{7.722700in}{2.342306in}}{\pgfqpoint{7.722700in}{2.353336in}}%
\pgfpathquadraticcurveto{\pgfqpoint{7.722700in}{2.366097in}}{\pgfqpoint{7.731689in}{2.373437in}}%
\pgfpathquadraticcurveto{\pgfqpoint{7.740678in}{2.380776in}}{\pgfqpoint{7.756449in}{2.380776in}}%
\pgfpathquadraticcurveto{\pgfqpoint{7.763232in}{2.380776in}}{\pgfqpoint{7.770241in}{2.379539in}}%
\pgfpathquadraticcurveto{\pgfqpoint{7.777271in}{2.378323in}}{\pgfqpoint{7.784611in}{2.375890in}}%
\pgfpathlineto{\pgfqpoint{7.784611in}{2.375890in}}%
\pgfpathclose%
\pgfpathmoveto{\pgfqpoint{7.821926in}{2.375499in}}%
\pgfpathlineto{\pgfqpoint{7.821926in}{2.355006in}}%
\pgfpathlineto{\pgfqpoint{7.846336in}{2.355006in}}%
\pgfpathlineto{\pgfqpoint{7.846336in}{2.345790in}}%
\pgfpathlineto{\pgfqpoint{7.821926in}{2.345790in}}%
\pgfpathlineto{\pgfqpoint{7.821926in}{2.306619in}}%
\pgfpathquadraticcurveto{\pgfqpoint{7.821926in}{2.297796in}}{\pgfqpoint{7.824339in}{2.295280in}}%
\pgfpathquadraticcurveto{\pgfqpoint{7.826751in}{2.292765in}}{\pgfqpoint{7.834173in}{2.292765in}}%
\pgfpathlineto{\pgfqpoint{7.846336in}{2.292765in}}%
\pgfpathlineto{\pgfqpoint{7.846336in}{2.282849in}}%
\pgfpathlineto{\pgfqpoint{7.834173in}{2.282849in}}%
\pgfpathquadraticcurveto{\pgfqpoint{7.820442in}{2.282849in}}{\pgfqpoint{7.815226in}{2.287962in}}%
\pgfpathquadraticcurveto{\pgfqpoint{7.810010in}{2.293095in}}{\pgfqpoint{7.810010in}{2.306619in}}%
\pgfpathlineto{\pgfqpoint{7.810010in}{2.345790in}}%
\pgfpathlineto{\pgfqpoint{7.801310in}{2.345790in}}%
\pgfpathlineto{\pgfqpoint{7.801310in}{2.355006in}}%
\pgfpathlineto{\pgfqpoint{7.810010in}{2.355006in}}%
\pgfpathlineto{\pgfqpoint{7.810010in}{2.375499in}}%
\pgfpathlineto{\pgfqpoint{7.821926in}{2.375499in}}%
\pgfpathlineto{\pgfqpoint{7.821926in}{2.375499in}}%
\pgfpathclose%
\pgfpathmoveto{\pgfqpoint{7.860706in}{2.311320in}}%
\pgfpathlineto{\pgfqpoint{7.860706in}{2.355006in}}%
\pgfpathlineto{\pgfqpoint{7.872560in}{2.355006in}}%
\pgfpathlineto{\pgfqpoint{7.872560in}{2.311773in}}%
\pgfpathquadraticcurveto{\pgfqpoint{7.872560in}{2.301527in}}{\pgfqpoint{7.876539in}{2.296394in}}%
\pgfpathquadraticcurveto{\pgfqpoint{7.880539in}{2.291281in}}{\pgfqpoint{7.888538in}{2.291281in}}%
\pgfpathquadraticcurveto{\pgfqpoint{7.898124in}{2.291281in}}{\pgfqpoint{7.903691in}{2.297404in}}%
\pgfpathquadraticcurveto{\pgfqpoint{7.909278in}{2.303527in}}{\pgfqpoint{7.909278in}{2.314103in}}%
\pgfpathlineto{\pgfqpoint{7.909278in}{2.355006in}}%
\pgfpathlineto{\pgfqpoint{7.921132in}{2.355006in}}%
\pgfpathlineto{\pgfqpoint{7.921132in}{2.282849in}}%
\pgfpathlineto{\pgfqpoint{7.909278in}{2.282849in}}%
\pgfpathlineto{\pgfqpoint{7.909278in}{2.293940in}}%
\pgfpathquadraticcurveto{\pgfqpoint{7.904969in}{2.287364in}}{\pgfqpoint{7.899258in}{2.284168in}}%
\pgfpathquadraticcurveto{\pgfqpoint{7.893568in}{2.280973in}}{\pgfqpoint{7.886023in}{2.280973in}}%
\pgfpathquadraticcurveto{\pgfqpoint{7.873591in}{2.280973in}}{\pgfqpoint{7.867138in}{2.288704in}}%
\pgfpathquadraticcurveto{\pgfqpoint{7.860706in}{2.296435in}}{\pgfqpoint{7.860706in}{2.311320in}}%
\pgfpathlineto{\pgfqpoint{7.860706in}{2.311320in}}%
\pgfpathclose%
\pgfpathmoveto{\pgfqpoint{7.890538in}{2.356738in}}%
\pgfpathlineto{\pgfqpoint{7.890538in}{2.356738in}}%
\pgfpathlineto{\pgfqpoint{7.890538in}{2.356738in}}%
\pgfpathclose%
\pgfpathmoveto{\pgfqpoint{7.993021in}{2.344059in}}%
\pgfpathlineto{\pgfqpoint{7.993021in}{2.383106in}}%
\pgfpathlineto{\pgfqpoint{8.004876in}{2.383106in}}%
\pgfpathlineto{\pgfqpoint{8.004876in}{2.282849in}}%
\pgfpathlineto{\pgfqpoint{7.993021in}{2.282849in}}%
\pgfpathlineto{\pgfqpoint{7.993021in}{2.293672in}}%
\pgfpathquadraticcurveto{\pgfqpoint{7.989290in}{2.287240in}}{\pgfqpoint{7.983579in}{2.284106in}}%
\pgfpathquadraticcurveto{\pgfqpoint{7.977889in}{2.280973in}}{\pgfqpoint{7.969890in}{2.280973in}}%
\pgfpathquadraticcurveto{\pgfqpoint{7.956819in}{2.280973in}}{\pgfqpoint{7.948593in}{2.291405in}}%
\pgfpathquadraticcurveto{\pgfqpoint{7.940388in}{2.301857in}}{\pgfqpoint{7.940388in}{2.318866in}}%
\pgfpathquadraticcurveto{\pgfqpoint{7.940388in}{2.335874in}}{\pgfqpoint{7.948593in}{2.346306in}}%
\pgfpathquadraticcurveto{\pgfqpoint{7.956819in}{2.356738in}}{\pgfqpoint{7.969890in}{2.356738in}}%
\pgfpathquadraticcurveto{\pgfqpoint{7.977889in}{2.356738in}}{\pgfqpoint{7.983579in}{2.353604in}}%
\pgfpathquadraticcurveto{\pgfqpoint{7.989290in}{2.350491in}}{\pgfqpoint{7.993021in}{2.344059in}}%
\pgfpathlineto{\pgfqpoint{7.993021in}{2.344059in}}%
\pgfpathclose%
\pgfpathmoveto{\pgfqpoint{7.952634in}{2.318866in}}%
\pgfpathquadraticcurveto{\pgfqpoint{7.952634in}{2.305795in}}{\pgfqpoint{7.958015in}{2.298352in}}%
\pgfpathquadraticcurveto{\pgfqpoint{7.963396in}{2.290910in}}{\pgfqpoint{7.972797in}{2.290910in}}%
\pgfpathquadraticcurveto{\pgfqpoint{7.982198in}{2.290910in}}{\pgfqpoint{7.987599in}{2.298352in}}%
\pgfpathquadraticcurveto{\pgfqpoint{7.993021in}{2.305795in}}{\pgfqpoint{7.993021in}{2.318866in}}%
\pgfpathquadraticcurveto{\pgfqpoint{7.993021in}{2.331936in}}{\pgfqpoint{7.987599in}{2.339379in}}%
\pgfpathquadraticcurveto{\pgfqpoint{7.982198in}{2.346821in}}{\pgfqpoint{7.972797in}{2.346821in}}%
\pgfpathquadraticcurveto{\pgfqpoint{7.963396in}{2.346821in}}{\pgfqpoint{7.958015in}{2.339379in}}%
\pgfpathquadraticcurveto{\pgfqpoint{7.952634in}{2.331936in}}{\pgfqpoint{7.952634in}{2.318866in}}%
\pgfpathlineto{\pgfqpoint{7.952634in}{2.318866in}}%
\pgfpathclose%
\pgfpathmoveto{\pgfqpoint{8.059323in}{2.276148in}}%
\pgfpathquadraticcurveto{\pgfqpoint{8.054314in}{2.263263in}}{\pgfqpoint{8.049531in}{2.259346in}}%
\pgfpathquadraticcurveto{\pgfqpoint{8.044768in}{2.255408in}}{\pgfqpoint{8.036790in}{2.255408in}}%
\pgfpathlineto{\pgfqpoint{8.027306in}{2.255408in}}%
\pgfpathlineto{\pgfqpoint{8.027306in}{2.265325in}}%
\pgfpathlineto{\pgfqpoint{8.034275in}{2.265325in}}%
\pgfpathquadraticcurveto{\pgfqpoint{8.039161in}{2.265325in}}{\pgfqpoint{8.041861in}{2.267655in}}%
\pgfpathquadraticcurveto{\pgfqpoint{8.044583in}{2.269964in}}{\pgfqpoint{8.047861in}{2.278602in}}%
\pgfpathlineto{\pgfqpoint{8.049984in}{2.284003in}}%
\pgfpathlineto{\pgfqpoint{8.020812in}{2.355006in}}%
\pgfpathlineto{\pgfqpoint{8.033367in}{2.355006in}}%
\pgfpathlineto{\pgfqpoint{8.055922in}{2.298579in}}%
\pgfpathlineto{\pgfqpoint{8.078476in}{2.355006in}}%
\pgfpathlineto{\pgfqpoint{8.091031in}{2.355006in}}%
\pgfpathlineto{\pgfqpoint{8.059323in}{2.276148in}}%
\pgfpathlineto{\pgfqpoint{8.059323in}{2.276148in}}%
\pgfpathclose%
\pgfpathmoveto{\pgfqpoint{8.153272in}{2.293796in}}%
\pgfpathlineto{\pgfqpoint{8.174527in}{2.293796in}}%
\pgfpathlineto{\pgfqpoint{8.174527in}{2.367190in}}%
\pgfpathlineto{\pgfqpoint{8.151396in}{2.362551in}}%
\pgfpathlineto{\pgfqpoint{8.151396in}{2.374406in}}%
\pgfpathlineto{\pgfqpoint{8.174404in}{2.379045in}}%
\pgfpathlineto{\pgfqpoint{8.187413in}{2.379045in}}%
\pgfpathlineto{\pgfqpoint{8.187413in}{2.293796in}}%
\pgfpathlineto{\pgfqpoint{8.208668in}{2.293796in}}%
\pgfpathlineto{\pgfqpoint{8.208668in}{2.282849in}}%
\pgfpathlineto{\pgfqpoint{8.153272in}{2.282849in}}%
\pgfpathlineto{\pgfqpoint{8.153272in}{2.293796in}}%
\pgfpathlineto{\pgfqpoint{8.153272in}{2.293796in}}%
\pgfpathclose%
\pgfpathmoveto{\pgfqpoint{8.270723in}{2.367706in}}%
\pgfpathlineto{\pgfqpoint{8.237861in}{2.316350in}}%
\pgfpathlineto{\pgfqpoint{8.270723in}{2.316350in}}%
\pgfpathlineto{\pgfqpoint{8.270723in}{2.367706in}}%
\pgfpathlineto{\pgfqpoint{8.270723in}{2.367706in}}%
\pgfpathclose%
\pgfpathmoveto{\pgfqpoint{8.267301in}{2.379045in}}%
\pgfpathlineto{\pgfqpoint{8.283670in}{2.379045in}}%
\pgfpathlineto{\pgfqpoint{8.283670in}{2.316350in}}%
\pgfpathlineto{\pgfqpoint{8.297401in}{2.316350in}}%
\pgfpathlineto{\pgfqpoint{8.297401in}{2.305527in}}%
\pgfpathlineto{\pgfqpoint{8.283670in}{2.305527in}}%
\pgfpathlineto{\pgfqpoint{8.283670in}{2.282849in}}%
\pgfpathlineto{\pgfqpoint{8.270723in}{2.282849in}}%
\pgfpathlineto{\pgfqpoint{8.270723in}{2.305527in}}%
\pgfpathlineto{\pgfqpoint{8.227305in}{2.305527in}}%
\pgfpathlineto{\pgfqpoint{8.227305in}{2.318082in}}%
\pgfpathlineto{\pgfqpoint{8.267301in}{2.379045in}}%
\pgfpathlineto{\pgfqpoint{8.267301in}{2.379045in}}%
\pgfpathclose%
\pgfpathmoveto{\pgfqpoint{8.387638in}{2.382962in}}%
\pgfpathquadraticcurveto{\pgfqpoint{8.379021in}{2.368159in}}{\pgfqpoint{8.374815in}{2.353645in}}%
\pgfpathquadraticcurveto{\pgfqpoint{8.370630in}{2.339152in}}{\pgfqpoint{8.370630in}{2.324267in}}%
\pgfpathquadraticcurveto{\pgfqpoint{8.370630in}{2.309403in}}{\pgfqpoint{8.374856in}{2.294806in}}%
\pgfpathquadraticcurveto{\pgfqpoint{8.379082in}{2.280210in}}{\pgfqpoint{8.387638in}{2.265449in}}%
\pgfpathlineto{\pgfqpoint{8.377330in}{2.265449in}}%
\pgfpathquadraticcurveto{\pgfqpoint{8.367682in}{2.280602in}}{\pgfqpoint{8.362878in}{2.295219in}}%
\pgfpathquadraticcurveto{\pgfqpoint{8.358074in}{2.309836in}}{\pgfqpoint{8.358074in}{2.324267in}}%
\pgfpathquadraticcurveto{\pgfqpoint{8.358074in}{2.338637in}}{\pgfqpoint{8.362837in}{2.353192in}}%
\pgfpathquadraticcurveto{\pgfqpoint{8.367620in}{2.367767in}}{\pgfqpoint{8.377330in}{2.382962in}}%
\pgfpathlineto{\pgfqpoint{8.387638in}{2.382962in}}%
\pgfpathlineto{\pgfqpoint{8.387638in}{2.382962in}}%
\pgfpathclose%
\pgfpathmoveto{\pgfqpoint{8.423531in}{2.293796in}}%
\pgfpathlineto{\pgfqpoint{8.468949in}{2.293796in}}%
\pgfpathlineto{\pgfqpoint{8.468949in}{2.282849in}}%
\pgfpathlineto{\pgfqpoint{8.407883in}{2.282849in}}%
\pgfpathlineto{\pgfqpoint{8.407883in}{2.293796in}}%
\pgfpathquadraticcurveto{\pgfqpoint{8.415285in}{2.301465in}}{\pgfqpoint{8.428067in}{2.314371in}}%
\pgfpathquadraticcurveto{\pgfqpoint{8.440870in}{2.327298in}}{\pgfqpoint{8.444148in}{2.331050in}}%
\pgfpathquadraticcurveto{\pgfqpoint{8.450394in}{2.338059in}}{\pgfqpoint{8.452868in}{2.342925in}}%
\pgfpathquadraticcurveto{\pgfqpoint{8.455363in}{2.347790in}}{\pgfqpoint{8.455363in}{2.352491in}}%
\pgfpathquadraticcurveto{\pgfqpoint{8.455363in}{2.360160in}}{\pgfqpoint{8.449982in}{2.364984in}}%
\pgfpathquadraticcurveto{\pgfqpoint{8.444601in}{2.369829in}}{\pgfqpoint{8.435963in}{2.369829in}}%
\pgfpathquadraticcurveto{\pgfqpoint{8.429840in}{2.369829in}}{\pgfqpoint{8.423036in}{2.367706in}}%
\pgfpathquadraticcurveto{\pgfqpoint{8.416254in}{2.365582in}}{\pgfqpoint{8.408523in}{2.361253in}}%
\pgfpathlineto{\pgfqpoint{8.408523in}{2.374406in}}%
\pgfpathquadraticcurveto{\pgfqpoint{8.416377in}{2.377560in}}{\pgfqpoint{8.423201in}{2.379168in}}%
\pgfpathquadraticcurveto{\pgfqpoint{8.430046in}{2.380776in}}{\pgfqpoint{8.435715in}{2.380776in}}%
\pgfpathquadraticcurveto{\pgfqpoint{8.450662in}{2.380776in}}{\pgfqpoint{8.459548in}{2.373293in}}%
\pgfpathquadraticcurveto{\pgfqpoint{8.468434in}{2.365829in}}{\pgfqpoint{8.468434in}{2.353336in}}%
\pgfpathquadraticcurveto{\pgfqpoint{8.468434in}{2.347398in}}{\pgfqpoint{8.466207in}{2.342079in}}%
\pgfpathquadraticcurveto{\pgfqpoint{8.464001in}{2.336781in}}{\pgfqpoint{8.458125in}{2.329565in}}%
\pgfpathquadraticcurveto{\pgfqpoint{8.456517in}{2.327689in}}{\pgfqpoint{8.447879in}{2.318762in}}%
\pgfpathquadraticcurveto{\pgfqpoint{8.439262in}{2.309836in}}{\pgfqpoint{8.423531in}{2.293796in}}%
\pgfpathlineto{\pgfqpoint{8.423531in}{2.293796in}}%
\pgfpathclose%
\pgfpathmoveto{\pgfqpoint{8.524098in}{2.370468in}}%
\pgfpathquadraticcurveto{\pgfqpoint{8.514058in}{2.370468in}}{\pgfqpoint{8.508986in}{2.360572in}}%
\pgfpathquadraticcurveto{\pgfqpoint{8.503935in}{2.350697in}}{\pgfqpoint{8.503935in}{2.330844in}}%
\pgfpathquadraticcurveto{\pgfqpoint{8.503935in}{2.311073in}}{\pgfqpoint{8.508986in}{2.301177in}}%
\pgfpathquadraticcurveto{\pgfqpoint{8.514058in}{2.291281in}}{\pgfqpoint{8.524098in}{2.291281in}}%
\pgfpathquadraticcurveto{\pgfqpoint{8.534220in}{2.291281in}}{\pgfqpoint{8.539271in}{2.301177in}}%
\pgfpathquadraticcurveto{\pgfqpoint{8.544343in}{2.311073in}}{\pgfqpoint{8.544343in}{2.330844in}}%
\pgfpathquadraticcurveto{\pgfqpoint{8.544343in}{2.350697in}}{\pgfqpoint{8.539271in}{2.360572in}}%
\pgfpathquadraticcurveto{\pgfqpoint{8.534220in}{2.370468in}}{\pgfqpoint{8.524098in}{2.370468in}}%
\pgfpathlineto{\pgfqpoint{8.524098in}{2.370468in}}%
\pgfpathclose%
\pgfpathmoveto{\pgfqpoint{8.524098in}{2.380776in}}%
\pgfpathquadraticcurveto{\pgfqpoint{8.540281in}{2.380776in}}{\pgfqpoint{8.548817in}{2.367974in}}%
\pgfpathquadraticcurveto{\pgfqpoint{8.557352in}{2.355191in}}{\pgfqpoint{8.557352in}{2.330844in}}%
\pgfpathquadraticcurveto{\pgfqpoint{8.557352in}{2.306558in}}{\pgfqpoint{8.548817in}{2.293755in}}%
\pgfpathquadraticcurveto{\pgfqpoint{8.540281in}{2.280973in}}{\pgfqpoint{8.524098in}{2.280973in}}%
\pgfpathquadraticcurveto{\pgfqpoint{8.507934in}{2.280973in}}{\pgfqpoint{8.499399in}{2.293755in}}%
\pgfpathquadraticcurveto{\pgfqpoint{8.490864in}{2.306558in}}{\pgfqpoint{8.490864in}{2.330844in}}%
\pgfpathquadraticcurveto{\pgfqpoint{8.490864in}{2.355191in}}{\pgfqpoint{8.499399in}{2.367974in}}%
\pgfpathquadraticcurveto{\pgfqpoint{8.507934in}{2.380776in}}{\pgfqpoint{8.524098in}{2.380776in}}%
\pgfpathlineto{\pgfqpoint{8.524098in}{2.380776in}}%
\pgfpathclose%
\pgfpathmoveto{\pgfqpoint{8.591431in}{2.293796in}}%
\pgfpathlineto{\pgfqpoint{8.636848in}{2.293796in}}%
\pgfpathlineto{\pgfqpoint{8.636848in}{2.282849in}}%
\pgfpathlineto{\pgfqpoint{8.575783in}{2.282849in}}%
\pgfpathlineto{\pgfqpoint{8.575783in}{2.293796in}}%
\pgfpathquadraticcurveto{\pgfqpoint{8.583184in}{2.301465in}}{\pgfqpoint{8.595966in}{2.314371in}}%
\pgfpathquadraticcurveto{\pgfqpoint{8.608769in}{2.327298in}}{\pgfqpoint{8.612047in}{2.331050in}}%
\pgfpathquadraticcurveto{\pgfqpoint{8.618294in}{2.338059in}}{\pgfqpoint{8.620768in}{2.342925in}}%
\pgfpathquadraticcurveto{\pgfqpoint{8.623262in}{2.347790in}}{\pgfqpoint{8.623262in}{2.352491in}}%
\pgfpathquadraticcurveto{\pgfqpoint{8.623262in}{2.360160in}}{\pgfqpoint{8.617881in}{2.364984in}}%
\pgfpathquadraticcurveto{\pgfqpoint{8.612500in}{2.369829in}}{\pgfqpoint{8.603862in}{2.369829in}}%
\pgfpathquadraticcurveto{\pgfqpoint{8.597739in}{2.369829in}}{\pgfqpoint{8.590936in}{2.367706in}}%
\pgfpathquadraticcurveto{\pgfqpoint{8.584153in}{2.365582in}}{\pgfqpoint{8.576422in}{2.361253in}}%
\pgfpathlineto{\pgfqpoint{8.576422in}{2.374406in}}%
\pgfpathquadraticcurveto{\pgfqpoint{8.584277in}{2.377560in}}{\pgfqpoint{8.591101in}{2.379168in}}%
\pgfpathquadraticcurveto{\pgfqpoint{8.597945in}{2.380776in}}{\pgfqpoint{8.603615in}{2.380776in}}%
\pgfpathquadraticcurveto{\pgfqpoint{8.618562in}{2.380776in}}{\pgfqpoint{8.627447in}{2.373293in}}%
\pgfpathquadraticcurveto{\pgfqpoint{8.636333in}{2.365829in}}{\pgfqpoint{8.636333in}{2.353336in}}%
\pgfpathquadraticcurveto{\pgfqpoint{8.636333in}{2.347398in}}{\pgfqpoint{8.634106in}{2.342079in}}%
\pgfpathquadraticcurveto{\pgfqpoint{8.631900in}{2.336781in}}{\pgfqpoint{8.626025in}{2.329565in}}%
\pgfpathquadraticcurveto{\pgfqpoint{8.624417in}{2.327689in}}{\pgfqpoint{8.615778in}{2.318762in}}%
\pgfpathquadraticcurveto{\pgfqpoint{8.607161in}{2.309836in}}{\pgfqpoint{8.591431in}{2.293796in}}%
\pgfpathlineto{\pgfqpoint{8.591431in}{2.293796in}}%
\pgfpathclose%
\pgfpathmoveto{\pgfqpoint{8.703604in}{2.334719in}}%
\pgfpathquadraticcurveto{\pgfqpoint{8.712943in}{2.332720in}}{\pgfqpoint{8.718180in}{2.326390in}}%
\pgfpathquadraticcurveto{\pgfqpoint{8.723437in}{2.320082in}}{\pgfqpoint{8.723437in}{2.310805in}}%
\pgfpathquadraticcurveto{\pgfqpoint{8.723437in}{2.296579in}}{\pgfqpoint{8.713644in}{2.288766in}}%
\pgfpathquadraticcurveto{\pgfqpoint{8.703851in}{2.280973in}}{\pgfqpoint{8.685812in}{2.280973in}}%
\pgfpathquadraticcurveto{\pgfqpoint{8.679771in}{2.280973in}}{\pgfqpoint{8.673360in}{2.282168in}}%
\pgfpathquadraticcurveto{\pgfqpoint{8.666948in}{2.283364in}}{\pgfqpoint{8.660124in}{2.285756in}}%
\pgfpathlineto{\pgfqpoint{8.660124in}{2.298311in}}%
\pgfpathquadraticcurveto{\pgfqpoint{8.665526in}{2.295157in}}{\pgfqpoint{8.671958in}{2.293549in}}%
\pgfpathquadraticcurveto{\pgfqpoint{8.678411in}{2.291941in}}{\pgfqpoint{8.685441in}{2.291941in}}%
\pgfpathquadraticcurveto{\pgfqpoint{8.697666in}{2.291941in}}{\pgfqpoint{8.704078in}{2.296765in}}%
\pgfpathquadraticcurveto{\pgfqpoint{8.710490in}{2.301589in}}{\pgfqpoint{8.710490in}{2.310805in}}%
\pgfpathquadraticcurveto{\pgfqpoint{8.710490in}{2.319319in}}{\pgfqpoint{8.704532in}{2.324102in}}%
\pgfpathquadraticcurveto{\pgfqpoint{8.698574in}{2.328906in}}{\pgfqpoint{8.687956in}{2.328906in}}%
\pgfpathlineto{\pgfqpoint{8.676741in}{2.328906in}}%
\pgfpathlineto{\pgfqpoint{8.676741in}{2.339606in}}%
\pgfpathlineto{\pgfqpoint{8.688472in}{2.339606in}}%
\pgfpathquadraticcurveto{\pgfqpoint{8.698058in}{2.339606in}}{\pgfqpoint{8.703150in}{2.343440in}}%
\pgfpathquadraticcurveto{\pgfqpoint{8.708243in}{2.347275in}}{\pgfqpoint{8.708243in}{2.354491in}}%
\pgfpathquadraticcurveto{\pgfqpoint{8.708243in}{2.361892in}}{\pgfqpoint{8.702985in}{2.365850in}}%
\pgfpathquadraticcurveto{\pgfqpoint{8.697749in}{2.369829in}}{\pgfqpoint{8.687956in}{2.369829in}}%
\pgfpathquadraticcurveto{\pgfqpoint{8.682596in}{2.369829in}}{\pgfqpoint{8.676473in}{2.368654in}}%
\pgfpathquadraticcurveto{\pgfqpoint{8.670350in}{2.367499in}}{\pgfqpoint{8.663010in}{2.365067in}}%
\pgfpathlineto{\pgfqpoint{8.663010in}{2.376653in}}%
\pgfpathquadraticcurveto{\pgfqpoint{8.670432in}{2.378715in}}{\pgfqpoint{8.676906in}{2.379745in}}%
\pgfpathquadraticcurveto{\pgfqpoint{8.683379in}{2.380776in}}{\pgfqpoint{8.689111in}{2.380776in}}%
\pgfpathquadraticcurveto{\pgfqpoint{8.703934in}{2.380776in}}{\pgfqpoint{8.712551in}{2.374035in}}%
\pgfpathquadraticcurveto{\pgfqpoint{8.721190in}{2.367314in}}{\pgfqpoint{8.721190in}{2.355851in}}%
\pgfpathquadraticcurveto{\pgfqpoint{8.721190in}{2.347852in}}{\pgfqpoint{8.716613in}{2.342347in}}%
\pgfpathquadraticcurveto{\pgfqpoint{8.712036in}{2.336843in}}{\pgfqpoint{8.703604in}{2.334719in}}%
\pgfpathlineto{\pgfqpoint{8.703604in}{2.334719in}}%
\pgfpathclose%
\pgfpathmoveto{\pgfqpoint{8.744589in}{2.382962in}}%
\pgfpathlineto{\pgfqpoint{8.754897in}{2.382962in}}%
\pgfpathquadraticcurveto{\pgfqpoint{8.764546in}{2.367767in}}{\pgfqpoint{8.769349in}{2.353192in}}%
\pgfpathquadraticcurveto{\pgfqpoint{8.774153in}{2.338637in}}{\pgfqpoint{8.774153in}{2.324267in}}%
\pgfpathquadraticcurveto{\pgfqpoint{8.774153in}{2.309836in}}{\pgfqpoint{8.769349in}{2.295219in}}%
\pgfpathquadraticcurveto{\pgfqpoint{8.764546in}{2.280602in}}{\pgfqpoint{8.754897in}{2.265449in}}%
\pgfpathlineto{\pgfqpoint{8.744589in}{2.265449in}}%
\pgfpathquadraticcurveto{\pgfqpoint{8.753145in}{2.280210in}}{\pgfqpoint{8.757371in}{2.294806in}}%
\pgfpathquadraticcurveto{\pgfqpoint{8.761598in}{2.309403in}}{\pgfqpoint{8.761598in}{2.324267in}}%
\pgfpathquadraticcurveto{\pgfqpoint{8.761598in}{2.339152in}}{\pgfqpoint{8.757371in}{2.353645in}}%
\pgfpathquadraticcurveto{\pgfqpoint{8.753145in}{2.368159in}}{\pgfqpoint{8.744589in}{2.382962in}}%
\pgfpathlineto{\pgfqpoint{8.744589in}{2.382962in}}%
\pgfpathclose%
\pgfusepath{fill}%
\end{pgfscope}%
\begin{pgfscope}%
\pgfsetbuttcap%
\pgfsetroundjoin%
\definecolor{currentfill}{rgb}{0.090196,0.168627,0.227451}%
\pgfsetfillcolor{currentfill}%
\pgfsetlinewidth{0.000000pt}%
\definecolor{currentstroke}{rgb}{0.090196,0.168627,0.227451}%
\pgfsetstrokecolor{currentstroke}%
\pgfsetdash{}{0pt}%
\pgfpathmoveto{\pgfqpoint{7.784611in}{1.976290in}}%
\pgfpathlineto{\pgfqpoint{7.784611in}{1.963591in}}%
\pgfpathquadraticcurveto{\pgfqpoint{7.777210in}{1.967137in}}{\pgfqpoint{7.770633in}{1.968868in}}%
\pgfpathquadraticcurveto{\pgfqpoint{7.764056in}{1.970621in}}{\pgfqpoint{7.757933in}{1.970621in}}%
\pgfpathquadraticcurveto{\pgfqpoint{7.747316in}{1.970621in}}{\pgfqpoint{7.741543in}{1.966497in}}%
\pgfpathquadraticcurveto{\pgfqpoint{7.735771in}{1.962374in}}{\pgfqpoint{7.735771in}{1.954767in}}%
\pgfpathquadraticcurveto{\pgfqpoint{7.735771in}{1.948376in}}{\pgfqpoint{7.739605in}{1.945118in}}%
\pgfpathquadraticcurveto{\pgfqpoint{7.743440in}{1.941882in}}{\pgfqpoint{7.754140in}{1.939882in}}%
\pgfpathlineto{\pgfqpoint{7.761995in}{1.938274in}}%
\pgfpathquadraticcurveto{\pgfqpoint{7.776550in}{1.935491in}}{\pgfqpoint{7.783477in}{1.928502in}}%
\pgfpathquadraticcurveto{\pgfqpoint{7.790404in}{1.921513in}}{\pgfqpoint{7.790404in}{1.909803in}}%
\pgfpathquadraticcurveto{\pgfqpoint{7.790404in}{1.895804in}}{\pgfqpoint{7.781024in}{1.888588in}}%
\pgfpathquadraticcurveto{\pgfqpoint{7.771664in}{1.881373in}}{\pgfqpoint{7.753563in}{1.881373in}}%
\pgfpathquadraticcurveto{\pgfqpoint{7.746739in}{1.881373in}}{\pgfqpoint{7.739028in}{1.882919in}}%
\pgfpathquadraticcurveto{\pgfqpoint{7.731338in}{1.884465in}}{\pgfqpoint{7.723092in}{1.887496in}}%
\pgfpathlineto{\pgfqpoint{7.723092in}{1.900896in}}%
\pgfpathquadraticcurveto{\pgfqpoint{7.731008in}{1.896464in}}{\pgfqpoint{7.738616in}{1.894196in}}%
\pgfpathquadraticcurveto{\pgfqpoint{7.746223in}{1.891949in}}{\pgfqpoint{7.753563in}{1.891949in}}%
\pgfpathquadraticcurveto{\pgfqpoint{7.764696in}{1.891949in}}{\pgfqpoint{7.770757in}{1.896320in}}%
\pgfpathquadraticcurveto{\pgfqpoint{7.776818in}{1.900711in}}{\pgfqpoint{7.776818in}{1.908834in}}%
\pgfpathquadraticcurveto{\pgfqpoint{7.776818in}{1.915905in}}{\pgfqpoint{7.772468in}{1.919905in}}%
\pgfpathquadraticcurveto{\pgfqpoint{7.768118in}{1.923904in}}{\pgfqpoint{7.758201in}{1.925904in}}%
\pgfpathlineto{\pgfqpoint{7.750264in}{1.927450in}}%
\pgfpathquadraticcurveto{\pgfqpoint{7.735709in}{1.930336in}}{\pgfqpoint{7.729194in}{1.936521in}}%
\pgfpathquadraticcurveto{\pgfqpoint{7.722700in}{1.942706in}}{\pgfqpoint{7.722700in}{1.953736in}}%
\pgfpathquadraticcurveto{\pgfqpoint{7.722700in}{1.966497in}}{\pgfqpoint{7.731689in}{1.973837in}}%
\pgfpathquadraticcurveto{\pgfqpoint{7.740678in}{1.981176in}}{\pgfqpoint{7.756449in}{1.981176in}}%
\pgfpathquadraticcurveto{\pgfqpoint{7.763232in}{1.981176in}}{\pgfqpoint{7.770241in}{1.979939in}}%
\pgfpathquadraticcurveto{\pgfqpoint{7.777271in}{1.978723in}}{\pgfqpoint{7.784611in}{1.976290in}}%
\pgfpathlineto{\pgfqpoint{7.784611in}{1.976290in}}%
\pgfpathclose%
\pgfpathmoveto{\pgfqpoint{7.821926in}{1.975899in}}%
\pgfpathlineto{\pgfqpoint{7.821926in}{1.955406in}}%
\pgfpathlineto{\pgfqpoint{7.846336in}{1.955406in}}%
\pgfpathlineto{\pgfqpoint{7.846336in}{1.946190in}}%
\pgfpathlineto{\pgfqpoint{7.821926in}{1.946190in}}%
\pgfpathlineto{\pgfqpoint{7.821926in}{1.907019in}}%
\pgfpathquadraticcurveto{\pgfqpoint{7.821926in}{1.898196in}}{\pgfqpoint{7.824339in}{1.895680in}}%
\pgfpathquadraticcurveto{\pgfqpoint{7.826751in}{1.893165in}}{\pgfqpoint{7.834173in}{1.893165in}}%
\pgfpathlineto{\pgfqpoint{7.846336in}{1.893165in}}%
\pgfpathlineto{\pgfqpoint{7.846336in}{1.883249in}}%
\pgfpathlineto{\pgfqpoint{7.834173in}{1.883249in}}%
\pgfpathquadraticcurveto{\pgfqpoint{7.820442in}{1.883249in}}{\pgfqpoint{7.815226in}{1.888362in}}%
\pgfpathquadraticcurveto{\pgfqpoint{7.810010in}{1.893495in}}{\pgfqpoint{7.810010in}{1.907019in}}%
\pgfpathlineto{\pgfqpoint{7.810010in}{1.946190in}}%
\pgfpathlineto{\pgfqpoint{7.801310in}{1.946190in}}%
\pgfpathlineto{\pgfqpoint{7.801310in}{1.955406in}}%
\pgfpathlineto{\pgfqpoint{7.810010in}{1.955406in}}%
\pgfpathlineto{\pgfqpoint{7.810010in}{1.975899in}}%
\pgfpathlineto{\pgfqpoint{7.821926in}{1.975899in}}%
\pgfpathlineto{\pgfqpoint{7.821926in}{1.975899in}}%
\pgfpathclose%
\pgfpathmoveto{\pgfqpoint{7.860706in}{1.911720in}}%
\pgfpathlineto{\pgfqpoint{7.860706in}{1.955406in}}%
\pgfpathlineto{\pgfqpoint{7.872560in}{1.955406in}}%
\pgfpathlineto{\pgfqpoint{7.872560in}{1.912173in}}%
\pgfpathquadraticcurveto{\pgfqpoint{7.872560in}{1.901927in}}{\pgfqpoint{7.876539in}{1.896794in}}%
\pgfpathquadraticcurveto{\pgfqpoint{7.880539in}{1.891681in}}{\pgfqpoint{7.888538in}{1.891681in}}%
\pgfpathquadraticcurveto{\pgfqpoint{7.898124in}{1.891681in}}{\pgfqpoint{7.903691in}{1.897804in}}%
\pgfpathquadraticcurveto{\pgfqpoint{7.909278in}{1.903927in}}{\pgfqpoint{7.909278in}{1.914503in}}%
\pgfpathlineto{\pgfqpoint{7.909278in}{1.955406in}}%
\pgfpathlineto{\pgfqpoint{7.921132in}{1.955406in}}%
\pgfpathlineto{\pgfqpoint{7.921132in}{1.883249in}}%
\pgfpathlineto{\pgfqpoint{7.909278in}{1.883249in}}%
\pgfpathlineto{\pgfqpoint{7.909278in}{1.894340in}}%
\pgfpathquadraticcurveto{\pgfqpoint{7.904969in}{1.887764in}}{\pgfqpoint{7.899258in}{1.884568in}}%
\pgfpathquadraticcurveto{\pgfqpoint{7.893568in}{1.881373in}}{\pgfqpoint{7.886023in}{1.881373in}}%
\pgfpathquadraticcurveto{\pgfqpoint{7.873591in}{1.881373in}}{\pgfqpoint{7.867138in}{1.889104in}}%
\pgfpathquadraticcurveto{\pgfqpoint{7.860706in}{1.896835in}}{\pgfqpoint{7.860706in}{1.911720in}}%
\pgfpathlineto{\pgfqpoint{7.860706in}{1.911720in}}%
\pgfpathclose%
\pgfpathmoveto{\pgfqpoint{7.890538in}{1.957138in}}%
\pgfpathlineto{\pgfqpoint{7.890538in}{1.957138in}}%
\pgfpathlineto{\pgfqpoint{7.890538in}{1.957138in}}%
\pgfpathclose%
\pgfpathmoveto{\pgfqpoint{7.993021in}{1.944459in}}%
\pgfpathlineto{\pgfqpoint{7.993021in}{1.983506in}}%
\pgfpathlineto{\pgfqpoint{8.004876in}{1.983506in}}%
\pgfpathlineto{\pgfqpoint{8.004876in}{1.883249in}}%
\pgfpathlineto{\pgfqpoint{7.993021in}{1.883249in}}%
\pgfpathlineto{\pgfqpoint{7.993021in}{1.894072in}}%
\pgfpathquadraticcurveto{\pgfqpoint{7.989290in}{1.887640in}}{\pgfqpoint{7.983579in}{1.884506in}}%
\pgfpathquadraticcurveto{\pgfqpoint{7.977889in}{1.881373in}}{\pgfqpoint{7.969890in}{1.881373in}}%
\pgfpathquadraticcurveto{\pgfqpoint{7.956819in}{1.881373in}}{\pgfqpoint{7.948593in}{1.891805in}}%
\pgfpathquadraticcurveto{\pgfqpoint{7.940388in}{1.902257in}}{\pgfqpoint{7.940388in}{1.919266in}}%
\pgfpathquadraticcurveto{\pgfqpoint{7.940388in}{1.936274in}}{\pgfqpoint{7.948593in}{1.946706in}}%
\pgfpathquadraticcurveto{\pgfqpoint{7.956819in}{1.957138in}}{\pgfqpoint{7.969890in}{1.957138in}}%
\pgfpathquadraticcurveto{\pgfqpoint{7.977889in}{1.957138in}}{\pgfqpoint{7.983579in}{1.954004in}}%
\pgfpathquadraticcurveto{\pgfqpoint{7.989290in}{1.950891in}}{\pgfqpoint{7.993021in}{1.944459in}}%
\pgfpathlineto{\pgfqpoint{7.993021in}{1.944459in}}%
\pgfpathclose%
\pgfpathmoveto{\pgfqpoint{7.952634in}{1.919266in}}%
\pgfpathquadraticcurveto{\pgfqpoint{7.952634in}{1.906195in}}{\pgfqpoint{7.958015in}{1.898752in}}%
\pgfpathquadraticcurveto{\pgfqpoint{7.963396in}{1.891310in}}{\pgfqpoint{7.972797in}{1.891310in}}%
\pgfpathquadraticcurveto{\pgfqpoint{7.982198in}{1.891310in}}{\pgfqpoint{7.987599in}{1.898752in}}%
\pgfpathquadraticcurveto{\pgfqpoint{7.993021in}{1.906195in}}{\pgfqpoint{7.993021in}{1.919266in}}%
\pgfpathquadraticcurveto{\pgfqpoint{7.993021in}{1.932336in}}{\pgfqpoint{7.987599in}{1.939779in}}%
\pgfpathquadraticcurveto{\pgfqpoint{7.982198in}{1.947221in}}{\pgfqpoint{7.972797in}{1.947221in}}%
\pgfpathquadraticcurveto{\pgfqpoint{7.963396in}{1.947221in}}{\pgfqpoint{7.958015in}{1.939779in}}%
\pgfpathquadraticcurveto{\pgfqpoint{7.952634in}{1.932336in}}{\pgfqpoint{7.952634in}{1.919266in}}%
\pgfpathlineto{\pgfqpoint{7.952634in}{1.919266in}}%
\pgfpathclose%
\pgfpathmoveto{\pgfqpoint{8.059323in}{1.876548in}}%
\pgfpathquadraticcurveto{\pgfqpoint{8.054314in}{1.863663in}}{\pgfqpoint{8.049531in}{1.859746in}}%
\pgfpathquadraticcurveto{\pgfqpoint{8.044768in}{1.855808in}}{\pgfqpoint{8.036790in}{1.855808in}}%
\pgfpathlineto{\pgfqpoint{8.027306in}{1.855808in}}%
\pgfpathlineto{\pgfqpoint{8.027306in}{1.865725in}}%
\pgfpathlineto{\pgfqpoint{8.034275in}{1.865725in}}%
\pgfpathquadraticcurveto{\pgfqpoint{8.039161in}{1.865725in}}{\pgfqpoint{8.041861in}{1.868055in}}%
\pgfpathquadraticcurveto{\pgfqpoint{8.044583in}{1.870364in}}{\pgfqpoint{8.047861in}{1.879002in}}%
\pgfpathlineto{\pgfqpoint{8.049984in}{1.884403in}}%
\pgfpathlineto{\pgfqpoint{8.020812in}{1.955406in}}%
\pgfpathlineto{\pgfqpoint{8.033367in}{1.955406in}}%
\pgfpathlineto{\pgfqpoint{8.055922in}{1.898979in}}%
\pgfpathlineto{\pgfqpoint{8.078476in}{1.955406in}}%
\pgfpathlineto{\pgfqpoint{8.091031in}{1.955406in}}%
\pgfpathlineto{\pgfqpoint{8.059323in}{1.876548in}}%
\pgfpathlineto{\pgfqpoint{8.059323in}{1.876548in}}%
\pgfpathclose%
\pgfpathmoveto{\pgfqpoint{8.153272in}{1.894196in}}%
\pgfpathlineto{\pgfqpoint{8.174527in}{1.894196in}}%
\pgfpathlineto{\pgfqpoint{8.174527in}{1.967590in}}%
\pgfpathlineto{\pgfqpoint{8.151396in}{1.962951in}}%
\pgfpathlineto{\pgfqpoint{8.151396in}{1.974806in}}%
\pgfpathlineto{\pgfqpoint{8.174404in}{1.979445in}}%
\pgfpathlineto{\pgfqpoint{8.187413in}{1.979445in}}%
\pgfpathlineto{\pgfqpoint{8.187413in}{1.894196in}}%
\pgfpathlineto{\pgfqpoint{8.208668in}{1.894196in}}%
\pgfpathlineto{\pgfqpoint{8.208668in}{1.883249in}}%
\pgfpathlineto{\pgfqpoint{8.153272in}{1.883249in}}%
\pgfpathlineto{\pgfqpoint{8.153272in}{1.894196in}}%
\pgfpathlineto{\pgfqpoint{8.153272in}{1.894196in}}%
\pgfpathclose%
\pgfpathmoveto{\pgfqpoint{8.235098in}{1.979445in}}%
\pgfpathlineto{\pgfqpoint{8.286185in}{1.979445in}}%
\pgfpathlineto{\pgfqpoint{8.286185in}{1.968477in}}%
\pgfpathlineto{\pgfqpoint{8.247014in}{1.968477in}}%
\pgfpathlineto{\pgfqpoint{8.247014in}{1.944912in}}%
\pgfpathquadraticcurveto{\pgfqpoint{8.249839in}{1.945881in}}{\pgfqpoint{8.252663in}{1.946355in}}%
\pgfpathquadraticcurveto{\pgfqpoint{8.255508in}{1.946830in}}{\pgfqpoint{8.258353in}{1.946830in}}%
\pgfpathquadraticcurveto{\pgfqpoint{8.274455in}{1.946830in}}{\pgfqpoint{8.283856in}{1.938006in}}%
\pgfpathquadraticcurveto{\pgfqpoint{8.293277in}{1.929182in}}{\pgfqpoint{8.293277in}{1.914111in}}%
\pgfpathquadraticcurveto{\pgfqpoint{8.293277in}{1.898587in}}{\pgfqpoint{8.283608in}{1.889970in}}%
\pgfpathquadraticcurveto{\pgfqpoint{8.273939in}{1.881373in}}{\pgfqpoint{8.256354in}{1.881373in}}%
\pgfpathquadraticcurveto{\pgfqpoint{8.250292in}{1.881373in}}{\pgfqpoint{8.244004in}{1.882404in}}%
\pgfpathquadraticcurveto{\pgfqpoint{8.237737in}{1.883434in}}{\pgfqpoint{8.231037in}{1.885496in}}%
\pgfpathlineto{\pgfqpoint{8.231037in}{1.898587in}}%
\pgfpathquadraticcurveto{\pgfqpoint{8.236830in}{1.895433in}}{\pgfqpoint{8.243015in}{1.893887in}}%
\pgfpathquadraticcurveto{\pgfqpoint{8.249200in}{1.892341in}}{\pgfqpoint{8.256086in}{1.892341in}}%
\pgfpathquadraticcurveto{\pgfqpoint{8.267239in}{1.892341in}}{\pgfqpoint{8.273733in}{1.898196in}}%
\pgfpathquadraticcurveto{\pgfqpoint{8.280248in}{1.904051in}}{\pgfqpoint{8.280248in}{1.914111in}}%
\pgfpathquadraticcurveto{\pgfqpoint{8.280248in}{1.924152in}}{\pgfqpoint{8.273733in}{1.930007in}}%
\pgfpathquadraticcurveto{\pgfqpoint{8.267239in}{1.935882in}}{\pgfqpoint{8.256086in}{1.935882in}}%
\pgfpathquadraticcurveto{\pgfqpoint{8.250870in}{1.935882in}}{\pgfqpoint{8.245674in}{1.934728in}}%
\pgfpathquadraticcurveto{\pgfqpoint{8.240500in}{1.933573in}}{\pgfqpoint{8.235098in}{1.931120in}}%
\pgfpathlineto{\pgfqpoint{8.235098in}{1.979445in}}%
\pgfpathlineto{\pgfqpoint{8.235098in}{1.979445in}}%
\pgfpathclose%
\pgfpathmoveto{\pgfqpoint{8.387638in}{1.983362in}}%
\pgfpathquadraticcurveto{\pgfqpoint{8.379021in}{1.968559in}}{\pgfqpoint{8.374815in}{1.954045in}}%
\pgfpathquadraticcurveto{\pgfqpoint{8.370630in}{1.939552in}}{\pgfqpoint{8.370630in}{1.924667in}}%
\pgfpathquadraticcurveto{\pgfqpoint{8.370630in}{1.909803in}}{\pgfqpoint{8.374856in}{1.895206in}}%
\pgfpathquadraticcurveto{\pgfqpoint{8.379082in}{1.880610in}}{\pgfqpoint{8.387638in}{1.865849in}}%
\pgfpathlineto{\pgfqpoint{8.377330in}{1.865849in}}%
\pgfpathquadraticcurveto{\pgfqpoint{8.367682in}{1.881002in}}{\pgfqpoint{8.362878in}{1.895619in}}%
\pgfpathquadraticcurveto{\pgfqpoint{8.358074in}{1.910236in}}{\pgfqpoint{8.358074in}{1.924667in}}%
\pgfpathquadraticcurveto{\pgfqpoint{8.358074in}{1.939037in}}{\pgfqpoint{8.362837in}{1.953592in}}%
\pgfpathquadraticcurveto{\pgfqpoint{8.367620in}{1.968167in}}{\pgfqpoint{8.377330in}{1.983362in}}%
\pgfpathlineto{\pgfqpoint{8.387638in}{1.983362in}}%
\pgfpathlineto{\pgfqpoint{8.387638in}{1.983362in}}%
\pgfpathclose%
\pgfpathmoveto{\pgfqpoint{8.423531in}{1.894196in}}%
\pgfpathlineto{\pgfqpoint{8.468949in}{1.894196in}}%
\pgfpathlineto{\pgfqpoint{8.468949in}{1.883249in}}%
\pgfpathlineto{\pgfqpoint{8.407883in}{1.883249in}}%
\pgfpathlineto{\pgfqpoint{8.407883in}{1.894196in}}%
\pgfpathquadraticcurveto{\pgfqpoint{8.415285in}{1.901865in}}{\pgfqpoint{8.428067in}{1.914771in}}%
\pgfpathquadraticcurveto{\pgfqpoint{8.440870in}{1.927698in}}{\pgfqpoint{8.444148in}{1.931450in}}%
\pgfpathquadraticcurveto{\pgfqpoint{8.450394in}{1.938459in}}{\pgfqpoint{8.452868in}{1.943325in}}%
\pgfpathquadraticcurveto{\pgfqpoint{8.455363in}{1.948190in}}{\pgfqpoint{8.455363in}{1.952891in}}%
\pgfpathquadraticcurveto{\pgfqpoint{8.455363in}{1.960560in}}{\pgfqpoint{8.449982in}{1.965384in}}%
\pgfpathquadraticcurveto{\pgfqpoint{8.444601in}{1.970229in}}{\pgfqpoint{8.435963in}{1.970229in}}%
\pgfpathquadraticcurveto{\pgfqpoint{8.429840in}{1.970229in}}{\pgfqpoint{8.423036in}{1.968106in}}%
\pgfpathquadraticcurveto{\pgfqpoint{8.416254in}{1.965982in}}{\pgfqpoint{8.408523in}{1.961653in}}%
\pgfpathlineto{\pgfqpoint{8.408523in}{1.974806in}}%
\pgfpathquadraticcurveto{\pgfqpoint{8.416377in}{1.977960in}}{\pgfqpoint{8.423201in}{1.979568in}}%
\pgfpathquadraticcurveto{\pgfqpoint{8.430046in}{1.981176in}}{\pgfqpoint{8.435715in}{1.981176in}}%
\pgfpathquadraticcurveto{\pgfqpoint{8.450662in}{1.981176in}}{\pgfqpoint{8.459548in}{1.973693in}}%
\pgfpathquadraticcurveto{\pgfqpoint{8.468434in}{1.966229in}}{\pgfqpoint{8.468434in}{1.953736in}}%
\pgfpathquadraticcurveto{\pgfqpoint{8.468434in}{1.947798in}}{\pgfqpoint{8.466207in}{1.942479in}}%
\pgfpathquadraticcurveto{\pgfqpoint{8.464001in}{1.937181in}}{\pgfqpoint{8.458125in}{1.929965in}}%
\pgfpathquadraticcurveto{\pgfqpoint{8.456517in}{1.928089in}}{\pgfqpoint{8.447879in}{1.919162in}}%
\pgfpathquadraticcurveto{\pgfqpoint{8.439262in}{1.910236in}}{\pgfqpoint{8.423531in}{1.894196in}}%
\pgfpathlineto{\pgfqpoint{8.423531in}{1.894196in}}%
\pgfpathclose%
\pgfpathmoveto{\pgfqpoint{8.524098in}{1.970868in}}%
\pgfpathquadraticcurveto{\pgfqpoint{8.514058in}{1.970868in}}{\pgfqpoint{8.508986in}{1.960972in}}%
\pgfpathquadraticcurveto{\pgfqpoint{8.503935in}{1.951097in}}{\pgfqpoint{8.503935in}{1.931244in}}%
\pgfpathquadraticcurveto{\pgfqpoint{8.503935in}{1.911473in}}{\pgfqpoint{8.508986in}{1.901577in}}%
\pgfpathquadraticcurveto{\pgfqpoint{8.514058in}{1.891681in}}{\pgfqpoint{8.524098in}{1.891681in}}%
\pgfpathquadraticcurveto{\pgfqpoint{8.534220in}{1.891681in}}{\pgfqpoint{8.539271in}{1.901577in}}%
\pgfpathquadraticcurveto{\pgfqpoint{8.544343in}{1.911473in}}{\pgfqpoint{8.544343in}{1.931244in}}%
\pgfpathquadraticcurveto{\pgfqpoint{8.544343in}{1.951097in}}{\pgfqpoint{8.539271in}{1.960972in}}%
\pgfpathquadraticcurveto{\pgfqpoint{8.534220in}{1.970868in}}{\pgfqpoint{8.524098in}{1.970868in}}%
\pgfpathlineto{\pgfqpoint{8.524098in}{1.970868in}}%
\pgfpathclose%
\pgfpathmoveto{\pgfqpoint{8.524098in}{1.981176in}}%
\pgfpathquadraticcurveto{\pgfqpoint{8.540281in}{1.981176in}}{\pgfqpoint{8.548817in}{1.968374in}}%
\pgfpathquadraticcurveto{\pgfqpoint{8.557352in}{1.955591in}}{\pgfqpoint{8.557352in}{1.931244in}}%
\pgfpathquadraticcurveto{\pgfqpoint{8.557352in}{1.906958in}}{\pgfqpoint{8.548817in}{1.894155in}}%
\pgfpathquadraticcurveto{\pgfqpoint{8.540281in}{1.881373in}}{\pgfqpoint{8.524098in}{1.881373in}}%
\pgfpathquadraticcurveto{\pgfqpoint{8.507934in}{1.881373in}}{\pgfqpoint{8.499399in}{1.894155in}}%
\pgfpathquadraticcurveto{\pgfqpoint{8.490864in}{1.906958in}}{\pgfqpoint{8.490864in}{1.931244in}}%
\pgfpathquadraticcurveto{\pgfqpoint{8.490864in}{1.955591in}}{\pgfqpoint{8.499399in}{1.968374in}}%
\pgfpathquadraticcurveto{\pgfqpoint{8.507934in}{1.981176in}}{\pgfqpoint{8.524098in}{1.981176in}}%
\pgfpathlineto{\pgfqpoint{8.524098in}{1.981176in}}%
\pgfpathclose%
\pgfpathmoveto{\pgfqpoint{8.591431in}{1.894196in}}%
\pgfpathlineto{\pgfqpoint{8.636848in}{1.894196in}}%
\pgfpathlineto{\pgfqpoint{8.636848in}{1.883249in}}%
\pgfpathlineto{\pgfqpoint{8.575783in}{1.883249in}}%
\pgfpathlineto{\pgfqpoint{8.575783in}{1.894196in}}%
\pgfpathquadraticcurveto{\pgfqpoint{8.583184in}{1.901865in}}{\pgfqpoint{8.595966in}{1.914771in}}%
\pgfpathquadraticcurveto{\pgfqpoint{8.608769in}{1.927698in}}{\pgfqpoint{8.612047in}{1.931450in}}%
\pgfpathquadraticcurveto{\pgfqpoint{8.618294in}{1.938459in}}{\pgfqpoint{8.620768in}{1.943325in}}%
\pgfpathquadraticcurveto{\pgfqpoint{8.623262in}{1.948190in}}{\pgfqpoint{8.623262in}{1.952891in}}%
\pgfpathquadraticcurveto{\pgfqpoint{8.623262in}{1.960560in}}{\pgfqpoint{8.617881in}{1.965384in}}%
\pgfpathquadraticcurveto{\pgfqpoint{8.612500in}{1.970229in}}{\pgfqpoint{8.603862in}{1.970229in}}%
\pgfpathquadraticcurveto{\pgfqpoint{8.597739in}{1.970229in}}{\pgfqpoint{8.590936in}{1.968106in}}%
\pgfpathquadraticcurveto{\pgfqpoint{8.584153in}{1.965982in}}{\pgfqpoint{8.576422in}{1.961653in}}%
\pgfpathlineto{\pgfqpoint{8.576422in}{1.974806in}}%
\pgfpathquadraticcurveto{\pgfqpoint{8.584277in}{1.977960in}}{\pgfqpoint{8.591101in}{1.979568in}}%
\pgfpathquadraticcurveto{\pgfqpoint{8.597945in}{1.981176in}}{\pgfqpoint{8.603615in}{1.981176in}}%
\pgfpathquadraticcurveto{\pgfqpoint{8.618562in}{1.981176in}}{\pgfqpoint{8.627447in}{1.973693in}}%
\pgfpathquadraticcurveto{\pgfqpoint{8.636333in}{1.966229in}}{\pgfqpoint{8.636333in}{1.953736in}}%
\pgfpathquadraticcurveto{\pgfqpoint{8.636333in}{1.947798in}}{\pgfqpoint{8.634106in}{1.942479in}}%
\pgfpathquadraticcurveto{\pgfqpoint{8.631900in}{1.937181in}}{\pgfqpoint{8.626025in}{1.929965in}}%
\pgfpathquadraticcurveto{\pgfqpoint{8.624417in}{1.928089in}}{\pgfqpoint{8.615778in}{1.919162in}}%
\pgfpathquadraticcurveto{\pgfqpoint{8.607161in}{1.910236in}}{\pgfqpoint{8.591431in}{1.894196in}}%
\pgfpathlineto{\pgfqpoint{8.591431in}{1.894196in}}%
\pgfpathclose%
\pgfpathmoveto{\pgfqpoint{8.703604in}{1.935119in}}%
\pgfpathquadraticcurveto{\pgfqpoint{8.712943in}{1.933120in}}{\pgfqpoint{8.718180in}{1.926790in}}%
\pgfpathquadraticcurveto{\pgfqpoint{8.723437in}{1.920482in}}{\pgfqpoint{8.723437in}{1.911205in}}%
\pgfpathquadraticcurveto{\pgfqpoint{8.723437in}{1.896979in}}{\pgfqpoint{8.713644in}{1.889166in}}%
\pgfpathquadraticcurveto{\pgfqpoint{8.703851in}{1.881373in}}{\pgfqpoint{8.685812in}{1.881373in}}%
\pgfpathquadraticcurveto{\pgfqpoint{8.679771in}{1.881373in}}{\pgfqpoint{8.673360in}{1.882568in}}%
\pgfpathquadraticcurveto{\pgfqpoint{8.666948in}{1.883764in}}{\pgfqpoint{8.660124in}{1.886156in}}%
\pgfpathlineto{\pgfqpoint{8.660124in}{1.898711in}}%
\pgfpathquadraticcurveto{\pgfqpoint{8.665526in}{1.895557in}}{\pgfqpoint{8.671958in}{1.893949in}}%
\pgfpathquadraticcurveto{\pgfqpoint{8.678411in}{1.892341in}}{\pgfqpoint{8.685441in}{1.892341in}}%
\pgfpathquadraticcurveto{\pgfqpoint{8.697666in}{1.892341in}}{\pgfqpoint{8.704078in}{1.897165in}}%
\pgfpathquadraticcurveto{\pgfqpoint{8.710490in}{1.901989in}}{\pgfqpoint{8.710490in}{1.911205in}}%
\pgfpathquadraticcurveto{\pgfqpoint{8.710490in}{1.919719in}}{\pgfqpoint{8.704532in}{1.924502in}}%
\pgfpathquadraticcurveto{\pgfqpoint{8.698574in}{1.929306in}}{\pgfqpoint{8.687956in}{1.929306in}}%
\pgfpathlineto{\pgfqpoint{8.676741in}{1.929306in}}%
\pgfpathlineto{\pgfqpoint{8.676741in}{1.940006in}}%
\pgfpathlineto{\pgfqpoint{8.688472in}{1.940006in}}%
\pgfpathquadraticcurveto{\pgfqpoint{8.698058in}{1.940006in}}{\pgfqpoint{8.703150in}{1.943840in}}%
\pgfpathquadraticcurveto{\pgfqpoint{8.708243in}{1.947675in}}{\pgfqpoint{8.708243in}{1.954891in}}%
\pgfpathquadraticcurveto{\pgfqpoint{8.708243in}{1.962292in}}{\pgfqpoint{8.702985in}{1.966250in}}%
\pgfpathquadraticcurveto{\pgfqpoint{8.697749in}{1.970229in}}{\pgfqpoint{8.687956in}{1.970229in}}%
\pgfpathquadraticcurveto{\pgfqpoint{8.682596in}{1.970229in}}{\pgfqpoint{8.676473in}{1.969054in}}%
\pgfpathquadraticcurveto{\pgfqpoint{8.670350in}{1.967899in}}{\pgfqpoint{8.663010in}{1.965467in}}%
\pgfpathlineto{\pgfqpoint{8.663010in}{1.977053in}}%
\pgfpathquadraticcurveto{\pgfqpoint{8.670432in}{1.979115in}}{\pgfqpoint{8.676906in}{1.980145in}}%
\pgfpathquadraticcurveto{\pgfqpoint{8.683379in}{1.981176in}}{\pgfqpoint{8.689111in}{1.981176in}}%
\pgfpathquadraticcurveto{\pgfqpoint{8.703934in}{1.981176in}}{\pgfqpoint{8.712551in}{1.974435in}}%
\pgfpathquadraticcurveto{\pgfqpoint{8.721190in}{1.967714in}}{\pgfqpoint{8.721190in}{1.956251in}}%
\pgfpathquadraticcurveto{\pgfqpoint{8.721190in}{1.948252in}}{\pgfqpoint{8.716613in}{1.942747in}}%
\pgfpathquadraticcurveto{\pgfqpoint{8.712036in}{1.937243in}}{\pgfqpoint{8.703604in}{1.935119in}}%
\pgfpathlineto{\pgfqpoint{8.703604in}{1.935119in}}%
\pgfpathclose%
\pgfpathmoveto{\pgfqpoint{8.744589in}{1.983362in}}%
\pgfpathlineto{\pgfqpoint{8.754897in}{1.983362in}}%
\pgfpathquadraticcurveto{\pgfqpoint{8.764546in}{1.968167in}}{\pgfqpoint{8.769349in}{1.953592in}}%
\pgfpathquadraticcurveto{\pgfqpoint{8.774153in}{1.939037in}}{\pgfqpoint{8.774153in}{1.924667in}}%
\pgfpathquadraticcurveto{\pgfqpoint{8.774153in}{1.910236in}}{\pgfqpoint{8.769349in}{1.895619in}}%
\pgfpathquadraticcurveto{\pgfqpoint{8.764546in}{1.881002in}}{\pgfqpoint{8.754897in}{1.865849in}}%
\pgfpathlineto{\pgfqpoint{8.744589in}{1.865849in}}%
\pgfpathquadraticcurveto{\pgfqpoint{8.753145in}{1.880610in}}{\pgfqpoint{8.757371in}{1.895206in}}%
\pgfpathquadraticcurveto{\pgfqpoint{8.761598in}{1.909803in}}{\pgfqpoint{8.761598in}{1.924667in}}%
\pgfpathquadraticcurveto{\pgfqpoint{8.761598in}{1.939552in}}{\pgfqpoint{8.757371in}{1.954045in}}%
\pgfpathquadraticcurveto{\pgfqpoint{8.753145in}{1.968559in}}{\pgfqpoint{8.744589in}{1.983362in}}%
\pgfpathlineto{\pgfqpoint{8.744589in}{1.983362in}}%
\pgfpathclose%
\pgfusepath{fill}%
\end{pgfscope}%
\begin{pgfscope}%
\pgfsetbuttcap%
\pgfsetroundjoin%
\definecolor{currentfill}{rgb}{0.090196,0.168627,0.227451}%
\pgfsetfillcolor{currentfill}%
\pgfsetlinewidth{0.000000pt}%
\definecolor{currentstroke}{rgb}{0.090196,0.168627,0.227451}%
\pgfsetstrokecolor{currentstroke}%
\pgfsetdash{}{0pt}%
\pgfpathmoveto{\pgfqpoint{7.784611in}{1.576690in}}%
\pgfpathlineto{\pgfqpoint{7.784611in}{1.563991in}}%
\pgfpathquadraticcurveto{\pgfqpoint{7.777210in}{1.567537in}}{\pgfqpoint{7.770633in}{1.569268in}}%
\pgfpathquadraticcurveto{\pgfqpoint{7.764056in}{1.571021in}}{\pgfqpoint{7.757933in}{1.571021in}}%
\pgfpathquadraticcurveto{\pgfqpoint{7.747316in}{1.571021in}}{\pgfqpoint{7.741543in}{1.566897in}}%
\pgfpathquadraticcurveto{\pgfqpoint{7.735771in}{1.562774in}}{\pgfqpoint{7.735771in}{1.555167in}}%
\pgfpathquadraticcurveto{\pgfqpoint{7.735771in}{1.548776in}}{\pgfqpoint{7.739605in}{1.545518in}}%
\pgfpathquadraticcurveto{\pgfqpoint{7.743440in}{1.542282in}}{\pgfqpoint{7.754140in}{1.540282in}}%
\pgfpathlineto{\pgfqpoint{7.761995in}{1.538674in}}%
\pgfpathquadraticcurveto{\pgfqpoint{7.776550in}{1.535891in}}{\pgfqpoint{7.783477in}{1.528902in}}%
\pgfpathquadraticcurveto{\pgfqpoint{7.790404in}{1.521913in}}{\pgfqpoint{7.790404in}{1.510203in}}%
\pgfpathquadraticcurveto{\pgfqpoint{7.790404in}{1.496204in}}{\pgfqpoint{7.781024in}{1.488988in}}%
\pgfpathquadraticcurveto{\pgfqpoint{7.771664in}{1.481773in}}{\pgfqpoint{7.753563in}{1.481773in}}%
\pgfpathquadraticcurveto{\pgfqpoint{7.746739in}{1.481773in}}{\pgfqpoint{7.739028in}{1.483319in}}%
\pgfpathquadraticcurveto{\pgfqpoint{7.731338in}{1.484865in}}{\pgfqpoint{7.723092in}{1.487896in}}%
\pgfpathlineto{\pgfqpoint{7.723092in}{1.501296in}}%
\pgfpathquadraticcurveto{\pgfqpoint{7.731008in}{1.496864in}}{\pgfqpoint{7.738616in}{1.494596in}}%
\pgfpathquadraticcurveto{\pgfqpoint{7.746223in}{1.492349in}}{\pgfqpoint{7.753563in}{1.492349in}}%
\pgfpathquadraticcurveto{\pgfqpoint{7.764696in}{1.492349in}}{\pgfqpoint{7.770757in}{1.496720in}}%
\pgfpathquadraticcurveto{\pgfqpoint{7.776818in}{1.501111in}}{\pgfqpoint{7.776818in}{1.509234in}}%
\pgfpathquadraticcurveto{\pgfqpoint{7.776818in}{1.516305in}}{\pgfqpoint{7.772468in}{1.520305in}}%
\pgfpathquadraticcurveto{\pgfqpoint{7.768118in}{1.524304in}}{\pgfqpoint{7.758201in}{1.526304in}}%
\pgfpathlineto{\pgfqpoint{7.750264in}{1.527850in}}%
\pgfpathquadraticcurveto{\pgfqpoint{7.735709in}{1.530736in}}{\pgfqpoint{7.729194in}{1.536921in}}%
\pgfpathquadraticcurveto{\pgfqpoint{7.722700in}{1.543106in}}{\pgfqpoint{7.722700in}{1.554136in}}%
\pgfpathquadraticcurveto{\pgfqpoint{7.722700in}{1.566897in}}{\pgfqpoint{7.731689in}{1.574237in}}%
\pgfpathquadraticcurveto{\pgfqpoint{7.740678in}{1.581576in}}{\pgfqpoint{7.756449in}{1.581576in}}%
\pgfpathquadraticcurveto{\pgfqpoint{7.763232in}{1.581576in}}{\pgfqpoint{7.770241in}{1.580339in}}%
\pgfpathquadraticcurveto{\pgfqpoint{7.777271in}{1.579123in}}{\pgfqpoint{7.784611in}{1.576690in}}%
\pgfpathlineto{\pgfqpoint{7.784611in}{1.576690in}}%
\pgfpathclose%
\pgfpathmoveto{\pgfqpoint{7.821926in}{1.576299in}}%
\pgfpathlineto{\pgfqpoint{7.821926in}{1.555806in}}%
\pgfpathlineto{\pgfqpoint{7.846336in}{1.555806in}}%
\pgfpathlineto{\pgfqpoint{7.846336in}{1.546590in}}%
\pgfpathlineto{\pgfqpoint{7.821926in}{1.546590in}}%
\pgfpathlineto{\pgfqpoint{7.821926in}{1.507419in}}%
\pgfpathquadraticcurveto{\pgfqpoint{7.821926in}{1.498596in}}{\pgfqpoint{7.824339in}{1.496080in}}%
\pgfpathquadraticcurveto{\pgfqpoint{7.826751in}{1.493565in}}{\pgfqpoint{7.834173in}{1.493565in}}%
\pgfpathlineto{\pgfqpoint{7.846336in}{1.493565in}}%
\pgfpathlineto{\pgfqpoint{7.846336in}{1.483649in}}%
\pgfpathlineto{\pgfqpoint{7.834173in}{1.483649in}}%
\pgfpathquadraticcurveto{\pgfqpoint{7.820442in}{1.483649in}}{\pgfqpoint{7.815226in}{1.488762in}}%
\pgfpathquadraticcurveto{\pgfqpoint{7.810010in}{1.493895in}}{\pgfqpoint{7.810010in}{1.507419in}}%
\pgfpathlineto{\pgfqpoint{7.810010in}{1.546590in}}%
\pgfpathlineto{\pgfqpoint{7.801310in}{1.546590in}}%
\pgfpathlineto{\pgfqpoint{7.801310in}{1.555806in}}%
\pgfpathlineto{\pgfqpoint{7.810010in}{1.555806in}}%
\pgfpathlineto{\pgfqpoint{7.810010in}{1.576299in}}%
\pgfpathlineto{\pgfqpoint{7.821926in}{1.576299in}}%
\pgfpathlineto{\pgfqpoint{7.821926in}{1.576299in}}%
\pgfpathclose%
\pgfpathmoveto{\pgfqpoint{7.860706in}{1.512120in}}%
\pgfpathlineto{\pgfqpoint{7.860706in}{1.555806in}}%
\pgfpathlineto{\pgfqpoint{7.872560in}{1.555806in}}%
\pgfpathlineto{\pgfqpoint{7.872560in}{1.512573in}}%
\pgfpathquadraticcurveto{\pgfqpoint{7.872560in}{1.502327in}}{\pgfqpoint{7.876539in}{1.497194in}}%
\pgfpathquadraticcurveto{\pgfqpoint{7.880539in}{1.492081in}}{\pgfqpoint{7.888538in}{1.492081in}}%
\pgfpathquadraticcurveto{\pgfqpoint{7.898124in}{1.492081in}}{\pgfqpoint{7.903691in}{1.498204in}}%
\pgfpathquadraticcurveto{\pgfqpoint{7.909278in}{1.504327in}}{\pgfqpoint{7.909278in}{1.514903in}}%
\pgfpathlineto{\pgfqpoint{7.909278in}{1.555806in}}%
\pgfpathlineto{\pgfqpoint{7.921132in}{1.555806in}}%
\pgfpathlineto{\pgfqpoint{7.921132in}{1.483649in}}%
\pgfpathlineto{\pgfqpoint{7.909278in}{1.483649in}}%
\pgfpathlineto{\pgfqpoint{7.909278in}{1.494740in}}%
\pgfpathquadraticcurveto{\pgfqpoint{7.904969in}{1.488164in}}{\pgfqpoint{7.899258in}{1.484968in}}%
\pgfpathquadraticcurveto{\pgfqpoint{7.893568in}{1.481773in}}{\pgfqpoint{7.886023in}{1.481773in}}%
\pgfpathquadraticcurveto{\pgfqpoint{7.873591in}{1.481773in}}{\pgfqpoint{7.867138in}{1.489504in}}%
\pgfpathquadraticcurveto{\pgfqpoint{7.860706in}{1.497235in}}{\pgfqpoint{7.860706in}{1.512120in}}%
\pgfpathlineto{\pgfqpoint{7.860706in}{1.512120in}}%
\pgfpathclose%
\pgfpathmoveto{\pgfqpoint{7.890538in}{1.557538in}}%
\pgfpathlineto{\pgfqpoint{7.890538in}{1.557538in}}%
\pgfpathlineto{\pgfqpoint{7.890538in}{1.557538in}}%
\pgfpathclose%
\pgfpathmoveto{\pgfqpoint{7.993021in}{1.544859in}}%
\pgfpathlineto{\pgfqpoint{7.993021in}{1.583906in}}%
\pgfpathlineto{\pgfqpoint{8.004876in}{1.583906in}}%
\pgfpathlineto{\pgfqpoint{8.004876in}{1.483649in}}%
\pgfpathlineto{\pgfqpoint{7.993021in}{1.483649in}}%
\pgfpathlineto{\pgfqpoint{7.993021in}{1.494472in}}%
\pgfpathquadraticcurveto{\pgfqpoint{7.989290in}{1.488040in}}{\pgfqpoint{7.983579in}{1.484906in}}%
\pgfpathquadraticcurveto{\pgfqpoint{7.977889in}{1.481773in}}{\pgfqpoint{7.969890in}{1.481773in}}%
\pgfpathquadraticcurveto{\pgfqpoint{7.956819in}{1.481773in}}{\pgfqpoint{7.948593in}{1.492205in}}%
\pgfpathquadraticcurveto{\pgfqpoint{7.940388in}{1.502657in}}{\pgfqpoint{7.940388in}{1.519666in}}%
\pgfpathquadraticcurveto{\pgfqpoint{7.940388in}{1.536674in}}{\pgfqpoint{7.948593in}{1.547106in}}%
\pgfpathquadraticcurveto{\pgfqpoint{7.956819in}{1.557538in}}{\pgfqpoint{7.969890in}{1.557538in}}%
\pgfpathquadraticcurveto{\pgfqpoint{7.977889in}{1.557538in}}{\pgfqpoint{7.983579in}{1.554404in}}%
\pgfpathquadraticcurveto{\pgfqpoint{7.989290in}{1.551291in}}{\pgfqpoint{7.993021in}{1.544859in}}%
\pgfpathlineto{\pgfqpoint{7.993021in}{1.544859in}}%
\pgfpathclose%
\pgfpathmoveto{\pgfqpoint{7.952634in}{1.519666in}}%
\pgfpathquadraticcurveto{\pgfqpoint{7.952634in}{1.506595in}}{\pgfqpoint{7.958015in}{1.499152in}}%
\pgfpathquadraticcurveto{\pgfqpoint{7.963396in}{1.491710in}}{\pgfqpoint{7.972797in}{1.491710in}}%
\pgfpathquadraticcurveto{\pgfqpoint{7.982198in}{1.491710in}}{\pgfqpoint{7.987599in}{1.499152in}}%
\pgfpathquadraticcurveto{\pgfqpoint{7.993021in}{1.506595in}}{\pgfqpoint{7.993021in}{1.519666in}}%
\pgfpathquadraticcurveto{\pgfqpoint{7.993021in}{1.532736in}}{\pgfqpoint{7.987599in}{1.540179in}}%
\pgfpathquadraticcurveto{\pgfqpoint{7.982198in}{1.547621in}}{\pgfqpoint{7.972797in}{1.547621in}}%
\pgfpathquadraticcurveto{\pgfqpoint{7.963396in}{1.547621in}}{\pgfqpoint{7.958015in}{1.540179in}}%
\pgfpathquadraticcurveto{\pgfqpoint{7.952634in}{1.532736in}}{\pgfqpoint{7.952634in}{1.519666in}}%
\pgfpathlineto{\pgfqpoint{7.952634in}{1.519666in}}%
\pgfpathclose%
\pgfpathmoveto{\pgfqpoint{8.059323in}{1.476948in}}%
\pgfpathquadraticcurveto{\pgfqpoint{8.054314in}{1.464063in}}{\pgfqpoint{8.049531in}{1.460146in}}%
\pgfpathquadraticcurveto{\pgfqpoint{8.044768in}{1.456208in}}{\pgfqpoint{8.036790in}{1.456208in}}%
\pgfpathlineto{\pgfqpoint{8.027306in}{1.456208in}}%
\pgfpathlineto{\pgfqpoint{8.027306in}{1.466125in}}%
\pgfpathlineto{\pgfqpoint{8.034275in}{1.466125in}}%
\pgfpathquadraticcurveto{\pgfqpoint{8.039161in}{1.466125in}}{\pgfqpoint{8.041861in}{1.468455in}}%
\pgfpathquadraticcurveto{\pgfqpoint{8.044583in}{1.470764in}}{\pgfqpoint{8.047861in}{1.479402in}}%
\pgfpathlineto{\pgfqpoint{8.049984in}{1.484803in}}%
\pgfpathlineto{\pgfqpoint{8.020812in}{1.555806in}}%
\pgfpathlineto{\pgfqpoint{8.033367in}{1.555806in}}%
\pgfpathlineto{\pgfqpoint{8.055922in}{1.499379in}}%
\pgfpathlineto{\pgfqpoint{8.078476in}{1.555806in}}%
\pgfpathlineto{\pgfqpoint{8.091031in}{1.555806in}}%
\pgfpathlineto{\pgfqpoint{8.059323in}{1.476948in}}%
\pgfpathlineto{\pgfqpoint{8.059323in}{1.476948in}}%
\pgfpathclose%
\pgfpathmoveto{\pgfqpoint{8.153272in}{1.494596in}}%
\pgfpathlineto{\pgfqpoint{8.174527in}{1.494596in}}%
\pgfpathlineto{\pgfqpoint{8.174527in}{1.567990in}}%
\pgfpathlineto{\pgfqpoint{8.151396in}{1.563351in}}%
\pgfpathlineto{\pgfqpoint{8.151396in}{1.575206in}}%
\pgfpathlineto{\pgfqpoint{8.174404in}{1.579845in}}%
\pgfpathlineto{\pgfqpoint{8.187413in}{1.579845in}}%
\pgfpathlineto{\pgfqpoint{8.187413in}{1.494596in}}%
\pgfpathlineto{\pgfqpoint{8.208668in}{1.494596in}}%
\pgfpathlineto{\pgfqpoint{8.208668in}{1.483649in}}%
\pgfpathlineto{\pgfqpoint{8.153272in}{1.483649in}}%
\pgfpathlineto{\pgfqpoint{8.153272in}{1.494596in}}%
\pgfpathlineto{\pgfqpoint{8.153272in}{1.494596in}}%
\pgfpathclose%
\pgfpathmoveto{\pgfqpoint{8.264414in}{1.536921in}}%
\pgfpathquadraticcurveto{\pgfqpoint{8.255653in}{1.536921in}}{\pgfqpoint{8.250519in}{1.530922in}}%
\pgfpathquadraticcurveto{\pgfqpoint{8.245406in}{1.524943in}}{\pgfqpoint{8.245406in}{1.514511in}}%
\pgfpathquadraticcurveto{\pgfqpoint{8.245406in}{1.504141in}}{\pgfqpoint{8.250519in}{1.498101in}}%
\pgfpathquadraticcurveto{\pgfqpoint{8.255653in}{1.492081in}}{\pgfqpoint{8.264414in}{1.492081in}}%
\pgfpathquadraticcurveto{\pgfqpoint{8.273176in}{1.492081in}}{\pgfqpoint{8.278289in}{1.498101in}}%
\pgfpathquadraticcurveto{\pgfqpoint{8.283402in}{1.504141in}}{\pgfqpoint{8.283402in}{1.514511in}}%
\pgfpathquadraticcurveto{\pgfqpoint{8.283402in}{1.524943in}}{\pgfqpoint{8.278289in}{1.530922in}}%
\pgfpathquadraticcurveto{\pgfqpoint{8.273176in}{1.536921in}}{\pgfqpoint{8.264414in}{1.536921in}}%
\pgfpathlineto{\pgfqpoint{8.264414in}{1.536921in}}%
\pgfpathclose%
\pgfpathmoveto{\pgfqpoint{8.290247in}{1.577721in}}%
\pgfpathlineto{\pgfqpoint{8.290247in}{1.565867in}}%
\pgfpathquadraticcurveto{\pgfqpoint{8.285340in}{1.568176in}}{\pgfqpoint{8.280351in}{1.569392in}}%
\pgfpathquadraticcurveto{\pgfqpoint{8.275362in}{1.570629in}}{\pgfqpoint{8.270455in}{1.570629in}}%
\pgfpathquadraticcurveto{\pgfqpoint{8.257570in}{1.570629in}}{\pgfqpoint{8.250766in}{1.561929in}}%
\pgfpathquadraticcurveto{\pgfqpoint{8.243984in}{1.553229in}}{\pgfqpoint{8.243015in}{1.535643in}}%
\pgfpathquadraticcurveto{\pgfqpoint{8.246808in}{1.541251in}}{\pgfqpoint{8.252539in}{1.544240in}}%
\pgfpathquadraticcurveto{\pgfqpoint{8.258291in}{1.547230in}}{\pgfqpoint{8.265177in}{1.547230in}}%
\pgfpathquadraticcurveto{\pgfqpoint{8.279671in}{1.547230in}}{\pgfqpoint{8.288082in}{1.538426in}}%
\pgfpathquadraticcurveto{\pgfqpoint{8.296493in}{1.529644in}}{\pgfqpoint{8.296493in}{1.514511in}}%
\pgfpathquadraticcurveto{\pgfqpoint{8.296493in}{1.499688in}}{\pgfqpoint{8.287732in}{1.490720in}}%
\pgfpathquadraticcurveto{\pgfqpoint{8.278970in}{1.481773in}}{\pgfqpoint{8.264414in}{1.481773in}}%
\pgfpathquadraticcurveto{\pgfqpoint{8.247715in}{1.481773in}}{\pgfqpoint{8.238891in}{1.494555in}}%
\pgfpathquadraticcurveto{\pgfqpoint{8.230068in}{1.507358in}}{\pgfqpoint{8.230068in}{1.531644in}}%
\pgfpathquadraticcurveto{\pgfqpoint{8.230068in}{1.554445in}}{\pgfqpoint{8.240891in}{1.568011in}}%
\pgfpathquadraticcurveto{\pgfqpoint{8.251715in}{1.581576in}}{\pgfqpoint{8.269940in}{1.581576in}}%
\pgfpathquadraticcurveto{\pgfqpoint{8.274846in}{1.581576in}}{\pgfqpoint{8.279836in}{1.580607in}}%
\pgfpathquadraticcurveto{\pgfqpoint{8.284825in}{1.579638in}}{\pgfqpoint{8.290247in}{1.577721in}}%
\pgfpathlineto{\pgfqpoint{8.290247in}{1.577721in}}%
\pgfpathclose%
\pgfpathmoveto{\pgfqpoint{8.387638in}{1.583762in}}%
\pgfpathquadraticcurveto{\pgfqpoint{8.379021in}{1.568959in}}{\pgfqpoint{8.374815in}{1.554445in}}%
\pgfpathquadraticcurveto{\pgfqpoint{8.370630in}{1.539952in}}{\pgfqpoint{8.370630in}{1.525067in}}%
\pgfpathquadraticcurveto{\pgfqpoint{8.370630in}{1.510203in}}{\pgfqpoint{8.374856in}{1.495606in}}%
\pgfpathquadraticcurveto{\pgfqpoint{8.379082in}{1.481010in}}{\pgfqpoint{8.387638in}{1.466249in}}%
\pgfpathlineto{\pgfqpoint{8.377330in}{1.466249in}}%
\pgfpathquadraticcurveto{\pgfqpoint{8.367682in}{1.481402in}}{\pgfqpoint{8.362878in}{1.496019in}}%
\pgfpathquadraticcurveto{\pgfqpoint{8.358074in}{1.510636in}}{\pgfqpoint{8.358074in}{1.525067in}}%
\pgfpathquadraticcurveto{\pgfqpoint{8.358074in}{1.539437in}}{\pgfqpoint{8.362837in}{1.553992in}}%
\pgfpathquadraticcurveto{\pgfqpoint{8.367620in}{1.568567in}}{\pgfqpoint{8.377330in}{1.583762in}}%
\pgfpathlineto{\pgfqpoint{8.387638in}{1.583762in}}%
\pgfpathlineto{\pgfqpoint{8.387638in}{1.583762in}}%
\pgfpathclose%
\pgfpathmoveto{\pgfqpoint{8.423531in}{1.494596in}}%
\pgfpathlineto{\pgfqpoint{8.468949in}{1.494596in}}%
\pgfpathlineto{\pgfqpoint{8.468949in}{1.483649in}}%
\pgfpathlineto{\pgfqpoint{8.407883in}{1.483649in}}%
\pgfpathlineto{\pgfqpoint{8.407883in}{1.494596in}}%
\pgfpathquadraticcurveto{\pgfqpoint{8.415285in}{1.502265in}}{\pgfqpoint{8.428067in}{1.515171in}}%
\pgfpathquadraticcurveto{\pgfqpoint{8.440870in}{1.528098in}}{\pgfqpoint{8.444148in}{1.531850in}}%
\pgfpathquadraticcurveto{\pgfqpoint{8.450394in}{1.538859in}}{\pgfqpoint{8.452868in}{1.543725in}}%
\pgfpathquadraticcurveto{\pgfqpoint{8.455363in}{1.548590in}}{\pgfqpoint{8.455363in}{1.553291in}}%
\pgfpathquadraticcurveto{\pgfqpoint{8.455363in}{1.560960in}}{\pgfqpoint{8.449982in}{1.565784in}}%
\pgfpathquadraticcurveto{\pgfqpoint{8.444601in}{1.570629in}}{\pgfqpoint{8.435963in}{1.570629in}}%
\pgfpathquadraticcurveto{\pgfqpoint{8.429840in}{1.570629in}}{\pgfqpoint{8.423036in}{1.568506in}}%
\pgfpathquadraticcurveto{\pgfqpoint{8.416254in}{1.566382in}}{\pgfqpoint{8.408523in}{1.562053in}}%
\pgfpathlineto{\pgfqpoint{8.408523in}{1.575206in}}%
\pgfpathquadraticcurveto{\pgfqpoint{8.416377in}{1.578360in}}{\pgfqpoint{8.423201in}{1.579968in}}%
\pgfpathquadraticcurveto{\pgfqpoint{8.430046in}{1.581576in}}{\pgfqpoint{8.435715in}{1.581576in}}%
\pgfpathquadraticcurveto{\pgfqpoint{8.450662in}{1.581576in}}{\pgfqpoint{8.459548in}{1.574093in}}%
\pgfpathquadraticcurveto{\pgfqpoint{8.468434in}{1.566629in}}{\pgfqpoint{8.468434in}{1.554136in}}%
\pgfpathquadraticcurveto{\pgfqpoint{8.468434in}{1.548198in}}{\pgfqpoint{8.466207in}{1.542879in}}%
\pgfpathquadraticcurveto{\pgfqpoint{8.464001in}{1.537581in}}{\pgfqpoint{8.458125in}{1.530365in}}%
\pgfpathquadraticcurveto{\pgfqpoint{8.456517in}{1.528489in}}{\pgfqpoint{8.447879in}{1.519562in}}%
\pgfpathquadraticcurveto{\pgfqpoint{8.439262in}{1.510636in}}{\pgfqpoint{8.423531in}{1.494596in}}%
\pgfpathlineto{\pgfqpoint{8.423531in}{1.494596in}}%
\pgfpathclose%
\pgfpathmoveto{\pgfqpoint{8.524098in}{1.571268in}}%
\pgfpathquadraticcurveto{\pgfqpoint{8.514058in}{1.571268in}}{\pgfqpoint{8.508986in}{1.561372in}}%
\pgfpathquadraticcurveto{\pgfqpoint{8.503935in}{1.551497in}}{\pgfqpoint{8.503935in}{1.531644in}}%
\pgfpathquadraticcurveto{\pgfqpoint{8.503935in}{1.511873in}}{\pgfqpoint{8.508986in}{1.501977in}}%
\pgfpathquadraticcurveto{\pgfqpoint{8.514058in}{1.492081in}}{\pgfqpoint{8.524098in}{1.492081in}}%
\pgfpathquadraticcurveto{\pgfqpoint{8.534220in}{1.492081in}}{\pgfqpoint{8.539271in}{1.501977in}}%
\pgfpathquadraticcurveto{\pgfqpoint{8.544343in}{1.511873in}}{\pgfqpoint{8.544343in}{1.531644in}}%
\pgfpathquadraticcurveto{\pgfqpoint{8.544343in}{1.551497in}}{\pgfqpoint{8.539271in}{1.561372in}}%
\pgfpathquadraticcurveto{\pgfqpoint{8.534220in}{1.571268in}}{\pgfqpoint{8.524098in}{1.571268in}}%
\pgfpathlineto{\pgfqpoint{8.524098in}{1.571268in}}%
\pgfpathclose%
\pgfpathmoveto{\pgfqpoint{8.524098in}{1.581576in}}%
\pgfpathquadraticcurveto{\pgfqpoint{8.540281in}{1.581576in}}{\pgfqpoint{8.548817in}{1.568774in}}%
\pgfpathquadraticcurveto{\pgfqpoint{8.557352in}{1.555991in}}{\pgfqpoint{8.557352in}{1.531644in}}%
\pgfpathquadraticcurveto{\pgfqpoint{8.557352in}{1.507358in}}{\pgfqpoint{8.548817in}{1.494555in}}%
\pgfpathquadraticcurveto{\pgfqpoint{8.540281in}{1.481773in}}{\pgfqpoint{8.524098in}{1.481773in}}%
\pgfpathquadraticcurveto{\pgfqpoint{8.507934in}{1.481773in}}{\pgfqpoint{8.499399in}{1.494555in}}%
\pgfpathquadraticcurveto{\pgfqpoint{8.490864in}{1.507358in}}{\pgfqpoint{8.490864in}{1.531644in}}%
\pgfpathquadraticcurveto{\pgfqpoint{8.490864in}{1.555991in}}{\pgfqpoint{8.499399in}{1.568774in}}%
\pgfpathquadraticcurveto{\pgfqpoint{8.507934in}{1.581576in}}{\pgfqpoint{8.524098in}{1.581576in}}%
\pgfpathlineto{\pgfqpoint{8.524098in}{1.581576in}}%
\pgfpathclose%
\pgfpathmoveto{\pgfqpoint{8.591431in}{1.494596in}}%
\pgfpathlineto{\pgfqpoint{8.636848in}{1.494596in}}%
\pgfpathlineto{\pgfqpoint{8.636848in}{1.483649in}}%
\pgfpathlineto{\pgfqpoint{8.575783in}{1.483649in}}%
\pgfpathlineto{\pgfqpoint{8.575783in}{1.494596in}}%
\pgfpathquadraticcurveto{\pgfqpoint{8.583184in}{1.502265in}}{\pgfqpoint{8.595966in}{1.515171in}}%
\pgfpathquadraticcurveto{\pgfqpoint{8.608769in}{1.528098in}}{\pgfqpoint{8.612047in}{1.531850in}}%
\pgfpathquadraticcurveto{\pgfqpoint{8.618294in}{1.538859in}}{\pgfqpoint{8.620768in}{1.543725in}}%
\pgfpathquadraticcurveto{\pgfqpoint{8.623262in}{1.548590in}}{\pgfqpoint{8.623262in}{1.553291in}}%
\pgfpathquadraticcurveto{\pgfqpoint{8.623262in}{1.560960in}}{\pgfqpoint{8.617881in}{1.565784in}}%
\pgfpathquadraticcurveto{\pgfqpoint{8.612500in}{1.570629in}}{\pgfqpoint{8.603862in}{1.570629in}}%
\pgfpathquadraticcurveto{\pgfqpoint{8.597739in}{1.570629in}}{\pgfqpoint{8.590936in}{1.568506in}}%
\pgfpathquadraticcurveto{\pgfqpoint{8.584153in}{1.566382in}}{\pgfqpoint{8.576422in}{1.562053in}}%
\pgfpathlineto{\pgfqpoint{8.576422in}{1.575206in}}%
\pgfpathquadraticcurveto{\pgfqpoint{8.584277in}{1.578360in}}{\pgfqpoint{8.591101in}{1.579968in}}%
\pgfpathquadraticcurveto{\pgfqpoint{8.597945in}{1.581576in}}{\pgfqpoint{8.603615in}{1.581576in}}%
\pgfpathquadraticcurveto{\pgfqpoint{8.618562in}{1.581576in}}{\pgfqpoint{8.627447in}{1.574093in}}%
\pgfpathquadraticcurveto{\pgfqpoint{8.636333in}{1.566629in}}{\pgfqpoint{8.636333in}{1.554136in}}%
\pgfpathquadraticcurveto{\pgfqpoint{8.636333in}{1.548198in}}{\pgfqpoint{8.634106in}{1.542879in}}%
\pgfpathquadraticcurveto{\pgfqpoint{8.631900in}{1.537581in}}{\pgfqpoint{8.626025in}{1.530365in}}%
\pgfpathquadraticcurveto{\pgfqpoint{8.624417in}{1.528489in}}{\pgfqpoint{8.615778in}{1.519562in}}%
\pgfpathquadraticcurveto{\pgfqpoint{8.607161in}{1.510636in}}{\pgfqpoint{8.591431in}{1.494596in}}%
\pgfpathlineto{\pgfqpoint{8.591431in}{1.494596in}}%
\pgfpathclose%
\pgfpathmoveto{\pgfqpoint{8.699934in}{1.568506in}}%
\pgfpathlineto{\pgfqpoint{8.667072in}{1.517150in}}%
\pgfpathlineto{\pgfqpoint{8.699934in}{1.517150in}}%
\pgfpathlineto{\pgfqpoint{8.699934in}{1.568506in}}%
\pgfpathlineto{\pgfqpoint{8.699934in}{1.568506in}}%
\pgfpathclose%
\pgfpathmoveto{\pgfqpoint{8.696512in}{1.579845in}}%
\pgfpathlineto{\pgfqpoint{8.712881in}{1.579845in}}%
\pgfpathlineto{\pgfqpoint{8.712881in}{1.517150in}}%
\pgfpathlineto{\pgfqpoint{8.726612in}{1.517150in}}%
\pgfpathlineto{\pgfqpoint{8.726612in}{1.506327in}}%
\pgfpathlineto{\pgfqpoint{8.712881in}{1.506327in}}%
\pgfpathlineto{\pgfqpoint{8.712881in}{1.483649in}}%
\pgfpathlineto{\pgfqpoint{8.699934in}{1.483649in}}%
\pgfpathlineto{\pgfqpoint{8.699934in}{1.506327in}}%
\pgfpathlineto{\pgfqpoint{8.656516in}{1.506327in}}%
\pgfpathlineto{\pgfqpoint{8.656516in}{1.518882in}}%
\pgfpathlineto{\pgfqpoint{8.696512in}{1.579845in}}%
\pgfpathlineto{\pgfqpoint{8.696512in}{1.579845in}}%
\pgfpathclose%
\pgfpathmoveto{\pgfqpoint{8.744589in}{1.583762in}}%
\pgfpathlineto{\pgfqpoint{8.754897in}{1.583762in}}%
\pgfpathquadraticcurveto{\pgfqpoint{8.764546in}{1.568567in}}{\pgfqpoint{8.769349in}{1.553992in}}%
\pgfpathquadraticcurveto{\pgfqpoint{8.774153in}{1.539437in}}{\pgfqpoint{8.774153in}{1.525067in}}%
\pgfpathquadraticcurveto{\pgfqpoint{8.774153in}{1.510636in}}{\pgfqpoint{8.769349in}{1.496019in}}%
\pgfpathquadraticcurveto{\pgfqpoint{8.764546in}{1.481402in}}{\pgfqpoint{8.754897in}{1.466249in}}%
\pgfpathlineto{\pgfqpoint{8.744589in}{1.466249in}}%
\pgfpathquadraticcurveto{\pgfqpoint{8.753145in}{1.481010in}}{\pgfqpoint{8.757371in}{1.495606in}}%
\pgfpathquadraticcurveto{\pgfqpoint{8.761598in}{1.510203in}}{\pgfqpoint{8.761598in}{1.525067in}}%
\pgfpathquadraticcurveto{\pgfqpoint{8.761598in}{1.539952in}}{\pgfqpoint{8.757371in}{1.554445in}}%
\pgfpathquadraticcurveto{\pgfqpoint{8.753145in}{1.568959in}}{\pgfqpoint{8.744589in}{1.583762in}}%
\pgfpathlineto{\pgfqpoint{8.744589in}{1.583762in}}%
\pgfpathclose%
\pgfusepath{fill}%
\end{pgfscope}%
\begin{pgfscope}%
\pgfsetbuttcap%
\pgfsetroundjoin%
\definecolor{currentfill}{rgb}{0.090196,0.168627,0.227451}%
\pgfsetfillcolor{currentfill}%
\pgfsetlinewidth{0.000000pt}%
\definecolor{currentstroke}{rgb}{0.090196,0.168627,0.227451}%
\pgfsetstrokecolor{currentstroke}%
\pgfsetdash{}{0pt}%
\pgfpathmoveto{\pgfqpoint{11.018611in}{3.175090in}}%
\pgfpathlineto{\pgfqpoint{11.018611in}{3.162391in}}%
\pgfpathquadraticcurveto{\pgfqpoint{11.011210in}{3.165937in}}{\pgfqpoint{11.004633in}{3.167668in}}%
\pgfpathquadraticcurveto{\pgfqpoint{10.998056in}{3.169421in}}{\pgfqpoint{10.991933in}{3.169421in}}%
\pgfpathquadraticcurveto{\pgfqpoint{10.981316in}{3.169421in}}{\pgfqpoint{10.975543in}{3.165297in}}%
\pgfpathquadraticcurveto{\pgfqpoint{10.969771in}{3.161174in}}{\pgfqpoint{10.969771in}{3.153567in}}%
\pgfpathquadraticcurveto{\pgfqpoint{10.969771in}{3.147176in}}{\pgfqpoint{10.973605in}{3.143918in}}%
\pgfpathquadraticcurveto{\pgfqpoint{10.977440in}{3.140682in}}{\pgfqpoint{10.988140in}{3.138682in}}%
\pgfpathlineto{\pgfqpoint{10.995995in}{3.137074in}}%
\pgfpathquadraticcurveto{\pgfqpoint{11.010550in}{3.134291in}}{\pgfqpoint{11.017477in}{3.127302in}}%
\pgfpathquadraticcurveto{\pgfqpoint{11.024404in}{3.120313in}}{\pgfqpoint{11.024404in}{3.108603in}}%
\pgfpathquadraticcurveto{\pgfqpoint{11.024404in}{3.094604in}}{\pgfqpoint{11.015024in}{3.087388in}}%
\pgfpathquadraticcurveto{\pgfqpoint{11.005664in}{3.080173in}}{\pgfqpoint{10.987563in}{3.080173in}}%
\pgfpathquadraticcurveto{\pgfqpoint{10.980739in}{3.080173in}}{\pgfqpoint{10.973028in}{3.081719in}}%
\pgfpathquadraticcurveto{\pgfqpoint{10.965338in}{3.083265in}}{\pgfqpoint{10.957092in}{3.086296in}}%
\pgfpathlineto{\pgfqpoint{10.957092in}{3.099696in}}%
\pgfpathquadraticcurveto{\pgfqpoint{10.965008in}{3.095264in}}{\pgfqpoint{10.972616in}{3.092996in}}%
\pgfpathquadraticcurveto{\pgfqpoint{10.980223in}{3.090749in}}{\pgfqpoint{10.987563in}{3.090749in}}%
\pgfpathquadraticcurveto{\pgfqpoint{10.998696in}{3.090749in}}{\pgfqpoint{11.004757in}{3.095120in}}%
\pgfpathquadraticcurveto{\pgfqpoint{11.010818in}{3.099511in}}{\pgfqpoint{11.010818in}{3.107634in}}%
\pgfpathquadraticcurveto{\pgfqpoint{11.010818in}{3.114705in}}{\pgfqpoint{11.006468in}{3.118705in}}%
\pgfpathquadraticcurveto{\pgfqpoint{11.002118in}{3.122704in}}{\pgfqpoint{10.992201in}{3.124704in}}%
\pgfpathlineto{\pgfqpoint{10.984264in}{3.126250in}}%
\pgfpathquadraticcurveto{\pgfqpoint{10.969709in}{3.129136in}}{\pgfqpoint{10.963194in}{3.135321in}}%
\pgfpathquadraticcurveto{\pgfqpoint{10.956700in}{3.141506in}}{\pgfqpoint{10.956700in}{3.152536in}}%
\pgfpathquadraticcurveto{\pgfqpoint{10.956700in}{3.165297in}}{\pgfqpoint{10.965689in}{3.172637in}}%
\pgfpathquadraticcurveto{\pgfqpoint{10.974678in}{3.179976in}}{\pgfqpoint{10.990449in}{3.179976in}}%
\pgfpathquadraticcurveto{\pgfqpoint{10.997232in}{3.179976in}}{\pgfqpoint{11.004241in}{3.178739in}}%
\pgfpathquadraticcurveto{\pgfqpoint{11.011271in}{3.177523in}}{\pgfqpoint{11.018611in}{3.175090in}}%
\pgfpathlineto{\pgfqpoint{11.018611in}{3.175090in}}%
\pgfpathclose%
\pgfpathmoveto{\pgfqpoint{11.055926in}{3.174699in}}%
\pgfpathlineto{\pgfqpoint{11.055926in}{3.154206in}}%
\pgfpathlineto{\pgfqpoint{11.080336in}{3.154206in}}%
\pgfpathlineto{\pgfqpoint{11.080336in}{3.144990in}}%
\pgfpathlineto{\pgfqpoint{11.055926in}{3.144990in}}%
\pgfpathlineto{\pgfqpoint{11.055926in}{3.105819in}}%
\pgfpathquadraticcurveto{\pgfqpoint{11.055926in}{3.096996in}}{\pgfqpoint{11.058339in}{3.094480in}}%
\pgfpathquadraticcurveto{\pgfqpoint{11.060751in}{3.091965in}}{\pgfqpoint{11.068173in}{3.091965in}}%
\pgfpathlineto{\pgfqpoint{11.080336in}{3.091965in}}%
\pgfpathlineto{\pgfqpoint{11.080336in}{3.082049in}}%
\pgfpathlineto{\pgfqpoint{11.068173in}{3.082049in}}%
\pgfpathquadraticcurveto{\pgfqpoint{11.054442in}{3.082049in}}{\pgfqpoint{11.049226in}{3.087162in}}%
\pgfpathquadraticcurveto{\pgfqpoint{11.044010in}{3.092295in}}{\pgfqpoint{11.044010in}{3.105819in}}%
\pgfpathlineto{\pgfqpoint{11.044010in}{3.144990in}}%
\pgfpathlineto{\pgfqpoint{11.035310in}{3.144990in}}%
\pgfpathlineto{\pgfqpoint{11.035310in}{3.154206in}}%
\pgfpathlineto{\pgfqpoint{11.044010in}{3.154206in}}%
\pgfpathlineto{\pgfqpoint{11.044010in}{3.174699in}}%
\pgfpathlineto{\pgfqpoint{11.055926in}{3.174699in}}%
\pgfpathlineto{\pgfqpoint{11.055926in}{3.174699in}}%
\pgfpathclose%
\pgfpathmoveto{\pgfqpoint{11.094706in}{3.110520in}}%
\pgfpathlineto{\pgfqpoint{11.094706in}{3.154206in}}%
\pgfpathlineto{\pgfqpoint{11.106560in}{3.154206in}}%
\pgfpathlineto{\pgfqpoint{11.106560in}{3.110973in}}%
\pgfpathquadraticcurveto{\pgfqpoint{11.106560in}{3.100727in}}{\pgfqpoint{11.110539in}{3.095594in}}%
\pgfpathquadraticcurveto{\pgfqpoint{11.114539in}{3.090481in}}{\pgfqpoint{11.122538in}{3.090481in}}%
\pgfpathquadraticcurveto{\pgfqpoint{11.132124in}{3.090481in}}{\pgfqpoint{11.137691in}{3.096604in}}%
\pgfpathquadraticcurveto{\pgfqpoint{11.143278in}{3.102727in}}{\pgfqpoint{11.143278in}{3.113303in}}%
\pgfpathlineto{\pgfqpoint{11.143278in}{3.154206in}}%
\pgfpathlineto{\pgfqpoint{11.155132in}{3.154206in}}%
\pgfpathlineto{\pgfqpoint{11.155132in}{3.082049in}}%
\pgfpathlineto{\pgfqpoint{11.143278in}{3.082049in}}%
\pgfpathlineto{\pgfqpoint{11.143278in}{3.093140in}}%
\pgfpathquadraticcurveto{\pgfqpoint{11.138969in}{3.086564in}}{\pgfqpoint{11.133258in}{3.083368in}}%
\pgfpathquadraticcurveto{\pgfqpoint{11.127568in}{3.080173in}}{\pgfqpoint{11.120023in}{3.080173in}}%
\pgfpathquadraticcurveto{\pgfqpoint{11.107591in}{3.080173in}}{\pgfqpoint{11.101138in}{3.087904in}}%
\pgfpathquadraticcurveto{\pgfqpoint{11.094706in}{3.095635in}}{\pgfqpoint{11.094706in}{3.110520in}}%
\pgfpathlineto{\pgfqpoint{11.094706in}{3.110520in}}%
\pgfpathclose%
\pgfpathmoveto{\pgfqpoint{11.124538in}{3.155938in}}%
\pgfpathlineto{\pgfqpoint{11.124538in}{3.155938in}}%
\pgfpathlineto{\pgfqpoint{11.124538in}{3.155938in}}%
\pgfpathclose%
\pgfpathmoveto{\pgfqpoint{11.227021in}{3.143259in}}%
\pgfpathlineto{\pgfqpoint{11.227021in}{3.182306in}}%
\pgfpathlineto{\pgfqpoint{11.238876in}{3.182306in}}%
\pgfpathlineto{\pgfqpoint{11.238876in}{3.082049in}}%
\pgfpathlineto{\pgfqpoint{11.227021in}{3.082049in}}%
\pgfpathlineto{\pgfqpoint{11.227021in}{3.092872in}}%
\pgfpathquadraticcurveto{\pgfqpoint{11.223290in}{3.086440in}}{\pgfqpoint{11.217579in}{3.083306in}}%
\pgfpathquadraticcurveto{\pgfqpoint{11.211889in}{3.080173in}}{\pgfqpoint{11.203890in}{3.080173in}}%
\pgfpathquadraticcurveto{\pgfqpoint{11.190819in}{3.080173in}}{\pgfqpoint{11.182593in}{3.090605in}}%
\pgfpathquadraticcurveto{\pgfqpoint{11.174388in}{3.101057in}}{\pgfqpoint{11.174388in}{3.118066in}}%
\pgfpathquadraticcurveto{\pgfqpoint{11.174388in}{3.135074in}}{\pgfqpoint{11.182593in}{3.145506in}}%
\pgfpathquadraticcurveto{\pgfqpoint{11.190819in}{3.155938in}}{\pgfqpoint{11.203890in}{3.155938in}}%
\pgfpathquadraticcurveto{\pgfqpoint{11.211889in}{3.155938in}}{\pgfqpoint{11.217579in}{3.152804in}}%
\pgfpathquadraticcurveto{\pgfqpoint{11.223290in}{3.149691in}}{\pgfqpoint{11.227021in}{3.143259in}}%
\pgfpathlineto{\pgfqpoint{11.227021in}{3.143259in}}%
\pgfpathclose%
\pgfpathmoveto{\pgfqpoint{11.186634in}{3.118066in}}%
\pgfpathquadraticcurveto{\pgfqpoint{11.186634in}{3.104995in}}{\pgfqpoint{11.192015in}{3.097552in}}%
\pgfpathquadraticcurveto{\pgfqpoint{11.197396in}{3.090110in}}{\pgfqpoint{11.206797in}{3.090110in}}%
\pgfpathquadraticcurveto{\pgfqpoint{11.216198in}{3.090110in}}{\pgfqpoint{11.221599in}{3.097552in}}%
\pgfpathquadraticcurveto{\pgfqpoint{11.227021in}{3.104995in}}{\pgfqpoint{11.227021in}{3.118066in}}%
\pgfpathquadraticcurveto{\pgfqpoint{11.227021in}{3.131136in}}{\pgfqpoint{11.221599in}{3.138579in}}%
\pgfpathquadraticcurveto{\pgfqpoint{11.216198in}{3.146021in}}{\pgfqpoint{11.206797in}{3.146021in}}%
\pgfpathquadraticcurveto{\pgfqpoint{11.197396in}{3.146021in}}{\pgfqpoint{11.192015in}{3.138579in}}%
\pgfpathquadraticcurveto{\pgfqpoint{11.186634in}{3.131136in}}{\pgfqpoint{11.186634in}{3.118066in}}%
\pgfpathlineto{\pgfqpoint{11.186634in}{3.118066in}}%
\pgfpathclose%
\pgfpathmoveto{\pgfqpoint{11.293323in}{3.075348in}}%
\pgfpathquadraticcurveto{\pgfqpoint{11.288314in}{3.062463in}}{\pgfqpoint{11.283531in}{3.058546in}}%
\pgfpathquadraticcurveto{\pgfqpoint{11.278768in}{3.054608in}}{\pgfqpoint{11.270790in}{3.054608in}}%
\pgfpathlineto{\pgfqpoint{11.261306in}{3.054608in}}%
\pgfpathlineto{\pgfqpoint{11.261306in}{3.064525in}}%
\pgfpathlineto{\pgfqpoint{11.268275in}{3.064525in}}%
\pgfpathquadraticcurveto{\pgfqpoint{11.273161in}{3.064525in}}{\pgfqpoint{11.275861in}{3.066855in}}%
\pgfpathquadraticcurveto{\pgfqpoint{11.278583in}{3.069164in}}{\pgfqpoint{11.281861in}{3.077802in}}%
\pgfpathlineto{\pgfqpoint{11.283984in}{3.083203in}}%
\pgfpathlineto{\pgfqpoint{11.254812in}{3.154206in}}%
\pgfpathlineto{\pgfqpoint{11.267367in}{3.154206in}}%
\pgfpathlineto{\pgfqpoint{11.289922in}{3.097779in}}%
\pgfpathlineto{\pgfqpoint{11.312476in}{3.154206in}}%
\pgfpathlineto{\pgfqpoint{11.325031in}{3.154206in}}%
\pgfpathlineto{\pgfqpoint{11.293323in}{3.075348in}}%
\pgfpathlineto{\pgfqpoint{11.293323in}{3.075348in}}%
\pgfpathclose%
\pgfpathmoveto{\pgfqpoint{11.387272in}{3.092996in}}%
\pgfpathlineto{\pgfqpoint{11.408527in}{3.092996in}}%
\pgfpathlineto{\pgfqpoint{11.408527in}{3.166390in}}%
\pgfpathlineto{\pgfqpoint{11.385396in}{3.161751in}}%
\pgfpathlineto{\pgfqpoint{11.385396in}{3.173606in}}%
\pgfpathlineto{\pgfqpoint{11.408404in}{3.178245in}}%
\pgfpathlineto{\pgfqpoint{11.421413in}{3.178245in}}%
\pgfpathlineto{\pgfqpoint{11.421413in}{3.092996in}}%
\pgfpathlineto{\pgfqpoint{11.442668in}{3.092996in}}%
\pgfpathlineto{\pgfqpoint{11.442668in}{3.082049in}}%
\pgfpathlineto{\pgfqpoint{11.387272in}{3.082049in}}%
\pgfpathlineto{\pgfqpoint{11.387272in}{3.092996in}}%
\pgfpathlineto{\pgfqpoint{11.387272in}{3.092996in}}%
\pgfpathclose%
\pgfpathmoveto{\pgfqpoint{11.465676in}{3.178245in}}%
\pgfpathlineto{\pgfqpoint{11.527525in}{3.178245in}}%
\pgfpathlineto{\pgfqpoint{11.527525in}{3.172699in}}%
\pgfpathlineto{\pgfqpoint{11.492601in}{3.082049in}}%
\pgfpathlineto{\pgfqpoint{11.479015in}{3.082049in}}%
\pgfpathlineto{\pgfqpoint{11.511877in}{3.167277in}}%
\pgfpathlineto{\pgfqpoint{11.465676in}{3.167277in}}%
\pgfpathlineto{\pgfqpoint{11.465676in}{3.178245in}}%
\pgfpathlineto{\pgfqpoint{11.465676in}{3.178245in}}%
\pgfpathclose%
\pgfpathmoveto{\pgfqpoint{11.621638in}{3.182162in}}%
\pgfpathquadraticcurveto{\pgfqpoint{11.613021in}{3.167359in}}{\pgfqpoint{11.608815in}{3.152845in}}%
\pgfpathquadraticcurveto{\pgfqpoint{11.604630in}{3.138352in}}{\pgfqpoint{11.604630in}{3.123467in}}%
\pgfpathquadraticcurveto{\pgfqpoint{11.604630in}{3.108603in}}{\pgfqpoint{11.608856in}{3.094006in}}%
\pgfpathquadraticcurveto{\pgfqpoint{11.613082in}{3.079410in}}{\pgfqpoint{11.621638in}{3.064649in}}%
\pgfpathlineto{\pgfqpoint{11.611330in}{3.064649in}}%
\pgfpathquadraticcurveto{\pgfqpoint{11.601682in}{3.079802in}}{\pgfqpoint{11.596878in}{3.094419in}}%
\pgfpathquadraticcurveto{\pgfqpoint{11.592074in}{3.109036in}}{\pgfqpoint{11.592074in}{3.123467in}}%
\pgfpathquadraticcurveto{\pgfqpoint{11.592074in}{3.137837in}}{\pgfqpoint{11.596837in}{3.152392in}}%
\pgfpathquadraticcurveto{\pgfqpoint{11.601620in}{3.166967in}}{\pgfqpoint{11.611330in}{3.182162in}}%
\pgfpathlineto{\pgfqpoint{11.621638in}{3.182162in}}%
\pgfpathlineto{\pgfqpoint{11.621638in}{3.182162in}}%
\pgfpathclose%
\pgfpathmoveto{\pgfqpoint{11.657531in}{3.092996in}}%
\pgfpathlineto{\pgfqpoint{11.702949in}{3.092996in}}%
\pgfpathlineto{\pgfqpoint{11.702949in}{3.082049in}}%
\pgfpathlineto{\pgfqpoint{11.641883in}{3.082049in}}%
\pgfpathlineto{\pgfqpoint{11.641883in}{3.092996in}}%
\pgfpathquadraticcurveto{\pgfqpoint{11.649285in}{3.100665in}}{\pgfqpoint{11.662067in}{3.113571in}}%
\pgfpathquadraticcurveto{\pgfqpoint{11.674870in}{3.126498in}}{\pgfqpoint{11.678148in}{3.130250in}}%
\pgfpathquadraticcurveto{\pgfqpoint{11.684394in}{3.137259in}}{\pgfqpoint{11.686868in}{3.142125in}}%
\pgfpathquadraticcurveto{\pgfqpoint{11.689363in}{3.146990in}}{\pgfqpoint{11.689363in}{3.151691in}}%
\pgfpathquadraticcurveto{\pgfqpoint{11.689363in}{3.159360in}}{\pgfqpoint{11.683982in}{3.164184in}}%
\pgfpathquadraticcurveto{\pgfqpoint{11.678601in}{3.169029in}}{\pgfqpoint{11.669963in}{3.169029in}}%
\pgfpathquadraticcurveto{\pgfqpoint{11.663840in}{3.169029in}}{\pgfqpoint{11.657036in}{3.166906in}}%
\pgfpathquadraticcurveto{\pgfqpoint{11.650254in}{3.164782in}}{\pgfqpoint{11.642523in}{3.160453in}}%
\pgfpathlineto{\pgfqpoint{11.642523in}{3.173606in}}%
\pgfpathquadraticcurveto{\pgfqpoint{11.650377in}{3.176760in}}{\pgfqpoint{11.657201in}{3.178368in}}%
\pgfpathquadraticcurveto{\pgfqpoint{11.664046in}{3.179976in}}{\pgfqpoint{11.669715in}{3.179976in}}%
\pgfpathquadraticcurveto{\pgfqpoint{11.684662in}{3.179976in}}{\pgfqpoint{11.693548in}{3.172493in}}%
\pgfpathquadraticcurveto{\pgfqpoint{11.702434in}{3.165029in}}{\pgfqpoint{11.702434in}{3.152536in}}%
\pgfpathquadraticcurveto{\pgfqpoint{11.702434in}{3.146598in}}{\pgfqpoint{11.700207in}{3.141279in}}%
\pgfpathquadraticcurveto{\pgfqpoint{11.698001in}{3.135981in}}{\pgfqpoint{11.692125in}{3.128765in}}%
\pgfpathquadraticcurveto{\pgfqpoint{11.690517in}{3.126889in}}{\pgfqpoint{11.681879in}{3.117962in}}%
\pgfpathquadraticcurveto{\pgfqpoint{11.673262in}{3.109036in}}{\pgfqpoint{11.657531in}{3.092996in}}%
\pgfpathlineto{\pgfqpoint{11.657531in}{3.092996in}}%
\pgfpathclose%
\pgfpathmoveto{\pgfqpoint{11.758098in}{3.169668in}}%
\pgfpathquadraticcurveto{\pgfqpoint{11.748058in}{3.169668in}}{\pgfqpoint{11.742986in}{3.159772in}}%
\pgfpathquadraticcurveto{\pgfqpoint{11.737935in}{3.149897in}}{\pgfqpoint{11.737935in}{3.130044in}}%
\pgfpathquadraticcurveto{\pgfqpoint{11.737935in}{3.110273in}}{\pgfqpoint{11.742986in}{3.100377in}}%
\pgfpathquadraticcurveto{\pgfqpoint{11.748058in}{3.090481in}}{\pgfqpoint{11.758098in}{3.090481in}}%
\pgfpathquadraticcurveto{\pgfqpoint{11.768220in}{3.090481in}}{\pgfqpoint{11.773271in}{3.100377in}}%
\pgfpathquadraticcurveto{\pgfqpoint{11.778343in}{3.110273in}}{\pgfqpoint{11.778343in}{3.130044in}}%
\pgfpathquadraticcurveto{\pgfqpoint{11.778343in}{3.149897in}}{\pgfqpoint{11.773271in}{3.159772in}}%
\pgfpathquadraticcurveto{\pgfqpoint{11.768220in}{3.169668in}}{\pgfqpoint{11.758098in}{3.169668in}}%
\pgfpathlineto{\pgfqpoint{11.758098in}{3.169668in}}%
\pgfpathclose%
\pgfpathmoveto{\pgfqpoint{11.758098in}{3.179976in}}%
\pgfpathquadraticcurveto{\pgfqpoint{11.774281in}{3.179976in}}{\pgfqpoint{11.782817in}{3.167174in}}%
\pgfpathquadraticcurveto{\pgfqpoint{11.791352in}{3.154391in}}{\pgfqpoint{11.791352in}{3.130044in}}%
\pgfpathquadraticcurveto{\pgfqpoint{11.791352in}{3.105758in}}{\pgfqpoint{11.782817in}{3.092955in}}%
\pgfpathquadraticcurveto{\pgfqpoint{11.774281in}{3.080173in}}{\pgfqpoint{11.758098in}{3.080173in}}%
\pgfpathquadraticcurveto{\pgfqpoint{11.741934in}{3.080173in}}{\pgfqpoint{11.733399in}{3.092955in}}%
\pgfpathquadraticcurveto{\pgfqpoint{11.724864in}{3.105758in}}{\pgfqpoint{11.724864in}{3.130044in}}%
\pgfpathquadraticcurveto{\pgfqpoint{11.724864in}{3.154391in}}{\pgfqpoint{11.733399in}{3.167174in}}%
\pgfpathquadraticcurveto{\pgfqpoint{11.741934in}{3.179976in}}{\pgfqpoint{11.758098in}{3.179976in}}%
\pgfpathlineto{\pgfqpoint{11.758098in}{3.179976in}}%
\pgfpathclose%
\pgfpathmoveto{\pgfqpoint{11.825431in}{3.092996in}}%
\pgfpathlineto{\pgfqpoint{11.870848in}{3.092996in}}%
\pgfpathlineto{\pgfqpoint{11.870848in}{3.082049in}}%
\pgfpathlineto{\pgfqpoint{11.809783in}{3.082049in}}%
\pgfpathlineto{\pgfqpoint{11.809783in}{3.092996in}}%
\pgfpathquadraticcurveto{\pgfqpoint{11.817184in}{3.100665in}}{\pgfqpoint{11.829966in}{3.113571in}}%
\pgfpathquadraticcurveto{\pgfqpoint{11.842769in}{3.126498in}}{\pgfqpoint{11.846047in}{3.130250in}}%
\pgfpathquadraticcurveto{\pgfqpoint{11.852294in}{3.137259in}}{\pgfqpoint{11.854768in}{3.142125in}}%
\pgfpathquadraticcurveto{\pgfqpoint{11.857262in}{3.146990in}}{\pgfqpoint{11.857262in}{3.151691in}}%
\pgfpathquadraticcurveto{\pgfqpoint{11.857262in}{3.159360in}}{\pgfqpoint{11.851881in}{3.164184in}}%
\pgfpathquadraticcurveto{\pgfqpoint{11.846500in}{3.169029in}}{\pgfqpoint{11.837862in}{3.169029in}}%
\pgfpathquadraticcurveto{\pgfqpoint{11.831739in}{3.169029in}}{\pgfqpoint{11.824936in}{3.166906in}}%
\pgfpathquadraticcurveto{\pgfqpoint{11.818153in}{3.164782in}}{\pgfqpoint{11.810422in}{3.160453in}}%
\pgfpathlineto{\pgfqpoint{11.810422in}{3.173606in}}%
\pgfpathquadraticcurveto{\pgfqpoint{11.818277in}{3.176760in}}{\pgfqpoint{11.825101in}{3.178368in}}%
\pgfpathquadraticcurveto{\pgfqpoint{11.831945in}{3.179976in}}{\pgfqpoint{11.837615in}{3.179976in}}%
\pgfpathquadraticcurveto{\pgfqpoint{11.852562in}{3.179976in}}{\pgfqpoint{11.861447in}{3.172493in}}%
\pgfpathquadraticcurveto{\pgfqpoint{11.870333in}{3.165029in}}{\pgfqpoint{11.870333in}{3.152536in}}%
\pgfpathquadraticcurveto{\pgfqpoint{11.870333in}{3.146598in}}{\pgfqpoint{11.868106in}{3.141279in}}%
\pgfpathquadraticcurveto{\pgfqpoint{11.865900in}{3.135981in}}{\pgfqpoint{11.860025in}{3.128765in}}%
\pgfpathquadraticcurveto{\pgfqpoint{11.858417in}{3.126889in}}{\pgfqpoint{11.849778in}{3.117962in}}%
\pgfpathquadraticcurveto{\pgfqpoint{11.841161in}{3.109036in}}{\pgfqpoint{11.825431in}{3.092996in}}%
\pgfpathlineto{\pgfqpoint{11.825431in}{3.092996in}}%
\pgfpathclose%
\pgfpathmoveto{\pgfqpoint{11.933934in}{3.166906in}}%
\pgfpathlineto{\pgfqpoint{11.901072in}{3.115550in}}%
\pgfpathlineto{\pgfqpoint{11.933934in}{3.115550in}}%
\pgfpathlineto{\pgfqpoint{11.933934in}{3.166906in}}%
\pgfpathlineto{\pgfqpoint{11.933934in}{3.166906in}}%
\pgfpathclose%
\pgfpathmoveto{\pgfqpoint{11.930512in}{3.178245in}}%
\pgfpathlineto{\pgfqpoint{11.946881in}{3.178245in}}%
\pgfpathlineto{\pgfqpoint{11.946881in}{3.115550in}}%
\pgfpathlineto{\pgfqpoint{11.960612in}{3.115550in}}%
\pgfpathlineto{\pgfqpoint{11.960612in}{3.104727in}}%
\pgfpathlineto{\pgfqpoint{11.946881in}{3.104727in}}%
\pgfpathlineto{\pgfqpoint{11.946881in}{3.082049in}}%
\pgfpathlineto{\pgfqpoint{11.933934in}{3.082049in}}%
\pgfpathlineto{\pgfqpoint{11.933934in}{3.104727in}}%
\pgfpathlineto{\pgfqpoint{11.890516in}{3.104727in}}%
\pgfpathlineto{\pgfqpoint{11.890516in}{3.117282in}}%
\pgfpathlineto{\pgfqpoint{11.930512in}{3.178245in}}%
\pgfpathlineto{\pgfqpoint{11.930512in}{3.178245in}}%
\pgfpathclose%
\pgfpathmoveto{\pgfqpoint{11.978589in}{3.182162in}}%
\pgfpathlineto{\pgfqpoint{11.988897in}{3.182162in}}%
\pgfpathquadraticcurveto{\pgfqpoint{11.998546in}{3.166967in}}{\pgfqpoint{12.003349in}{3.152392in}}%
\pgfpathquadraticcurveto{\pgfqpoint{12.008153in}{3.137837in}}{\pgfqpoint{12.008153in}{3.123467in}}%
\pgfpathquadraticcurveto{\pgfqpoint{12.008153in}{3.109036in}}{\pgfqpoint{12.003349in}{3.094419in}}%
\pgfpathquadraticcurveto{\pgfqpoint{11.998546in}{3.079802in}}{\pgfqpoint{11.988897in}{3.064649in}}%
\pgfpathlineto{\pgfqpoint{11.978589in}{3.064649in}}%
\pgfpathquadraticcurveto{\pgfqpoint{11.987145in}{3.079410in}}{\pgfqpoint{11.991371in}{3.094006in}}%
\pgfpathquadraticcurveto{\pgfqpoint{11.995598in}{3.108603in}}{\pgfqpoint{11.995598in}{3.123467in}}%
\pgfpathquadraticcurveto{\pgfqpoint{11.995598in}{3.138352in}}{\pgfqpoint{11.991371in}{3.152845in}}%
\pgfpathquadraticcurveto{\pgfqpoint{11.987145in}{3.167359in}}{\pgfqpoint{11.978589in}{3.182162in}}%
\pgfpathlineto{\pgfqpoint{11.978589in}{3.182162in}}%
\pgfpathclose%
\pgfusepath{fill}%
\end{pgfscope}%
\begin{pgfscope}%
\pgfsetbuttcap%
\pgfsetroundjoin%
\definecolor{currentfill}{rgb}{0.090196,0.168627,0.227451}%
\pgfsetfillcolor{currentfill}%
\pgfsetlinewidth{0.000000pt}%
\definecolor{currentstroke}{rgb}{0.090196,0.168627,0.227451}%
\pgfsetstrokecolor{currentstroke}%
\pgfsetdash{}{0pt}%
\pgfpathmoveto{\pgfqpoint{11.018611in}{2.775490in}}%
\pgfpathlineto{\pgfqpoint{11.018611in}{2.762791in}}%
\pgfpathquadraticcurveto{\pgfqpoint{11.011210in}{2.766337in}}{\pgfqpoint{11.004633in}{2.768068in}}%
\pgfpathquadraticcurveto{\pgfqpoint{10.998056in}{2.769821in}}{\pgfqpoint{10.991933in}{2.769821in}}%
\pgfpathquadraticcurveto{\pgfqpoint{10.981316in}{2.769821in}}{\pgfqpoint{10.975543in}{2.765697in}}%
\pgfpathquadraticcurveto{\pgfqpoint{10.969771in}{2.761574in}}{\pgfqpoint{10.969771in}{2.753967in}}%
\pgfpathquadraticcurveto{\pgfqpoint{10.969771in}{2.747576in}}{\pgfqpoint{10.973605in}{2.744318in}}%
\pgfpathquadraticcurveto{\pgfqpoint{10.977440in}{2.741082in}}{\pgfqpoint{10.988140in}{2.739082in}}%
\pgfpathlineto{\pgfqpoint{10.995995in}{2.737474in}}%
\pgfpathquadraticcurveto{\pgfqpoint{11.010550in}{2.734691in}}{\pgfqpoint{11.017477in}{2.727702in}}%
\pgfpathquadraticcurveto{\pgfqpoint{11.024404in}{2.720713in}}{\pgfqpoint{11.024404in}{2.709003in}}%
\pgfpathquadraticcurveto{\pgfqpoint{11.024404in}{2.695004in}}{\pgfqpoint{11.015024in}{2.687788in}}%
\pgfpathquadraticcurveto{\pgfqpoint{11.005664in}{2.680573in}}{\pgfqpoint{10.987563in}{2.680573in}}%
\pgfpathquadraticcurveto{\pgfqpoint{10.980739in}{2.680573in}}{\pgfqpoint{10.973028in}{2.682119in}}%
\pgfpathquadraticcurveto{\pgfqpoint{10.965338in}{2.683665in}}{\pgfqpoint{10.957092in}{2.686696in}}%
\pgfpathlineto{\pgfqpoint{10.957092in}{2.700096in}}%
\pgfpathquadraticcurveto{\pgfqpoint{10.965008in}{2.695664in}}{\pgfqpoint{10.972616in}{2.693396in}}%
\pgfpathquadraticcurveto{\pgfqpoint{10.980223in}{2.691149in}}{\pgfqpoint{10.987563in}{2.691149in}}%
\pgfpathquadraticcurveto{\pgfqpoint{10.998696in}{2.691149in}}{\pgfqpoint{11.004757in}{2.695520in}}%
\pgfpathquadraticcurveto{\pgfqpoint{11.010818in}{2.699911in}}{\pgfqpoint{11.010818in}{2.708034in}}%
\pgfpathquadraticcurveto{\pgfqpoint{11.010818in}{2.715105in}}{\pgfqpoint{11.006468in}{2.719105in}}%
\pgfpathquadraticcurveto{\pgfqpoint{11.002118in}{2.723104in}}{\pgfqpoint{10.992201in}{2.725104in}}%
\pgfpathlineto{\pgfqpoint{10.984264in}{2.726650in}}%
\pgfpathquadraticcurveto{\pgfqpoint{10.969709in}{2.729536in}}{\pgfqpoint{10.963194in}{2.735721in}}%
\pgfpathquadraticcurveto{\pgfqpoint{10.956700in}{2.741906in}}{\pgfqpoint{10.956700in}{2.752936in}}%
\pgfpathquadraticcurveto{\pgfqpoint{10.956700in}{2.765697in}}{\pgfqpoint{10.965689in}{2.773037in}}%
\pgfpathquadraticcurveto{\pgfqpoint{10.974678in}{2.780376in}}{\pgfqpoint{10.990449in}{2.780376in}}%
\pgfpathquadraticcurveto{\pgfqpoint{10.997232in}{2.780376in}}{\pgfqpoint{11.004241in}{2.779139in}}%
\pgfpathquadraticcurveto{\pgfqpoint{11.011271in}{2.777923in}}{\pgfqpoint{11.018611in}{2.775490in}}%
\pgfpathlineto{\pgfqpoint{11.018611in}{2.775490in}}%
\pgfpathclose%
\pgfpathmoveto{\pgfqpoint{11.055926in}{2.775099in}}%
\pgfpathlineto{\pgfqpoint{11.055926in}{2.754606in}}%
\pgfpathlineto{\pgfqpoint{11.080336in}{2.754606in}}%
\pgfpathlineto{\pgfqpoint{11.080336in}{2.745390in}}%
\pgfpathlineto{\pgfqpoint{11.055926in}{2.745390in}}%
\pgfpathlineto{\pgfqpoint{11.055926in}{2.706219in}}%
\pgfpathquadraticcurveto{\pgfqpoint{11.055926in}{2.697396in}}{\pgfqpoint{11.058339in}{2.694880in}}%
\pgfpathquadraticcurveto{\pgfqpoint{11.060751in}{2.692365in}}{\pgfqpoint{11.068173in}{2.692365in}}%
\pgfpathlineto{\pgfqpoint{11.080336in}{2.692365in}}%
\pgfpathlineto{\pgfqpoint{11.080336in}{2.682449in}}%
\pgfpathlineto{\pgfqpoint{11.068173in}{2.682449in}}%
\pgfpathquadraticcurveto{\pgfqpoint{11.054442in}{2.682449in}}{\pgfqpoint{11.049226in}{2.687562in}}%
\pgfpathquadraticcurveto{\pgfqpoint{11.044010in}{2.692695in}}{\pgfqpoint{11.044010in}{2.706219in}}%
\pgfpathlineto{\pgfqpoint{11.044010in}{2.745390in}}%
\pgfpathlineto{\pgfqpoint{11.035310in}{2.745390in}}%
\pgfpathlineto{\pgfqpoint{11.035310in}{2.754606in}}%
\pgfpathlineto{\pgfqpoint{11.044010in}{2.754606in}}%
\pgfpathlineto{\pgfqpoint{11.044010in}{2.775099in}}%
\pgfpathlineto{\pgfqpoint{11.055926in}{2.775099in}}%
\pgfpathlineto{\pgfqpoint{11.055926in}{2.775099in}}%
\pgfpathclose%
\pgfpathmoveto{\pgfqpoint{11.094706in}{2.710920in}}%
\pgfpathlineto{\pgfqpoint{11.094706in}{2.754606in}}%
\pgfpathlineto{\pgfqpoint{11.106560in}{2.754606in}}%
\pgfpathlineto{\pgfqpoint{11.106560in}{2.711373in}}%
\pgfpathquadraticcurveto{\pgfqpoint{11.106560in}{2.701127in}}{\pgfqpoint{11.110539in}{2.695994in}}%
\pgfpathquadraticcurveto{\pgfqpoint{11.114539in}{2.690881in}}{\pgfqpoint{11.122538in}{2.690881in}}%
\pgfpathquadraticcurveto{\pgfqpoint{11.132124in}{2.690881in}}{\pgfqpoint{11.137691in}{2.697004in}}%
\pgfpathquadraticcurveto{\pgfqpoint{11.143278in}{2.703127in}}{\pgfqpoint{11.143278in}{2.713703in}}%
\pgfpathlineto{\pgfqpoint{11.143278in}{2.754606in}}%
\pgfpathlineto{\pgfqpoint{11.155132in}{2.754606in}}%
\pgfpathlineto{\pgfqpoint{11.155132in}{2.682449in}}%
\pgfpathlineto{\pgfqpoint{11.143278in}{2.682449in}}%
\pgfpathlineto{\pgfqpoint{11.143278in}{2.693540in}}%
\pgfpathquadraticcurveto{\pgfqpoint{11.138969in}{2.686964in}}{\pgfqpoint{11.133258in}{2.683768in}}%
\pgfpathquadraticcurveto{\pgfqpoint{11.127568in}{2.680573in}}{\pgfqpoint{11.120023in}{2.680573in}}%
\pgfpathquadraticcurveto{\pgfqpoint{11.107591in}{2.680573in}}{\pgfqpoint{11.101138in}{2.688304in}}%
\pgfpathquadraticcurveto{\pgfqpoint{11.094706in}{2.696035in}}{\pgfqpoint{11.094706in}{2.710920in}}%
\pgfpathlineto{\pgfqpoint{11.094706in}{2.710920in}}%
\pgfpathclose%
\pgfpathmoveto{\pgfqpoint{11.124538in}{2.756338in}}%
\pgfpathlineto{\pgfqpoint{11.124538in}{2.756338in}}%
\pgfpathlineto{\pgfqpoint{11.124538in}{2.756338in}}%
\pgfpathclose%
\pgfpathmoveto{\pgfqpoint{11.227021in}{2.743659in}}%
\pgfpathlineto{\pgfqpoint{11.227021in}{2.782706in}}%
\pgfpathlineto{\pgfqpoint{11.238876in}{2.782706in}}%
\pgfpathlineto{\pgfqpoint{11.238876in}{2.682449in}}%
\pgfpathlineto{\pgfqpoint{11.227021in}{2.682449in}}%
\pgfpathlineto{\pgfqpoint{11.227021in}{2.693272in}}%
\pgfpathquadraticcurveto{\pgfqpoint{11.223290in}{2.686840in}}{\pgfqpoint{11.217579in}{2.683706in}}%
\pgfpathquadraticcurveto{\pgfqpoint{11.211889in}{2.680573in}}{\pgfqpoint{11.203890in}{2.680573in}}%
\pgfpathquadraticcurveto{\pgfqpoint{11.190819in}{2.680573in}}{\pgfqpoint{11.182593in}{2.691005in}}%
\pgfpathquadraticcurveto{\pgfqpoint{11.174388in}{2.701457in}}{\pgfqpoint{11.174388in}{2.718466in}}%
\pgfpathquadraticcurveto{\pgfqpoint{11.174388in}{2.735474in}}{\pgfqpoint{11.182593in}{2.745906in}}%
\pgfpathquadraticcurveto{\pgfqpoint{11.190819in}{2.756338in}}{\pgfqpoint{11.203890in}{2.756338in}}%
\pgfpathquadraticcurveto{\pgfqpoint{11.211889in}{2.756338in}}{\pgfqpoint{11.217579in}{2.753204in}}%
\pgfpathquadraticcurveto{\pgfqpoint{11.223290in}{2.750091in}}{\pgfqpoint{11.227021in}{2.743659in}}%
\pgfpathlineto{\pgfqpoint{11.227021in}{2.743659in}}%
\pgfpathclose%
\pgfpathmoveto{\pgfqpoint{11.186634in}{2.718466in}}%
\pgfpathquadraticcurveto{\pgfqpoint{11.186634in}{2.705395in}}{\pgfqpoint{11.192015in}{2.697952in}}%
\pgfpathquadraticcurveto{\pgfqpoint{11.197396in}{2.690510in}}{\pgfqpoint{11.206797in}{2.690510in}}%
\pgfpathquadraticcurveto{\pgfqpoint{11.216198in}{2.690510in}}{\pgfqpoint{11.221599in}{2.697952in}}%
\pgfpathquadraticcurveto{\pgfqpoint{11.227021in}{2.705395in}}{\pgfqpoint{11.227021in}{2.718466in}}%
\pgfpathquadraticcurveto{\pgfqpoint{11.227021in}{2.731536in}}{\pgfqpoint{11.221599in}{2.738979in}}%
\pgfpathquadraticcurveto{\pgfqpoint{11.216198in}{2.746421in}}{\pgfqpoint{11.206797in}{2.746421in}}%
\pgfpathquadraticcurveto{\pgfqpoint{11.197396in}{2.746421in}}{\pgfqpoint{11.192015in}{2.738979in}}%
\pgfpathquadraticcurveto{\pgfqpoint{11.186634in}{2.731536in}}{\pgfqpoint{11.186634in}{2.718466in}}%
\pgfpathlineto{\pgfqpoint{11.186634in}{2.718466in}}%
\pgfpathclose%
\pgfpathmoveto{\pgfqpoint{11.293323in}{2.675748in}}%
\pgfpathquadraticcurveto{\pgfqpoint{11.288314in}{2.662863in}}{\pgfqpoint{11.283531in}{2.658946in}}%
\pgfpathquadraticcurveto{\pgfqpoint{11.278768in}{2.655008in}}{\pgfqpoint{11.270790in}{2.655008in}}%
\pgfpathlineto{\pgfqpoint{11.261306in}{2.655008in}}%
\pgfpathlineto{\pgfqpoint{11.261306in}{2.664925in}}%
\pgfpathlineto{\pgfqpoint{11.268275in}{2.664925in}}%
\pgfpathquadraticcurveto{\pgfqpoint{11.273161in}{2.664925in}}{\pgfqpoint{11.275861in}{2.667255in}}%
\pgfpathquadraticcurveto{\pgfqpoint{11.278583in}{2.669564in}}{\pgfqpoint{11.281861in}{2.678202in}}%
\pgfpathlineto{\pgfqpoint{11.283984in}{2.683603in}}%
\pgfpathlineto{\pgfqpoint{11.254812in}{2.754606in}}%
\pgfpathlineto{\pgfqpoint{11.267367in}{2.754606in}}%
\pgfpathlineto{\pgfqpoint{11.289922in}{2.698179in}}%
\pgfpathlineto{\pgfqpoint{11.312476in}{2.754606in}}%
\pgfpathlineto{\pgfqpoint{11.325031in}{2.754606in}}%
\pgfpathlineto{\pgfqpoint{11.293323in}{2.675748in}}%
\pgfpathlineto{\pgfqpoint{11.293323in}{2.675748in}}%
\pgfpathclose%
\pgfpathmoveto{\pgfqpoint{11.387272in}{2.693396in}}%
\pgfpathlineto{\pgfqpoint{11.408527in}{2.693396in}}%
\pgfpathlineto{\pgfqpoint{11.408527in}{2.766790in}}%
\pgfpathlineto{\pgfqpoint{11.385396in}{2.762151in}}%
\pgfpathlineto{\pgfqpoint{11.385396in}{2.774006in}}%
\pgfpathlineto{\pgfqpoint{11.408404in}{2.778645in}}%
\pgfpathlineto{\pgfqpoint{11.421413in}{2.778645in}}%
\pgfpathlineto{\pgfqpoint{11.421413in}{2.693396in}}%
\pgfpathlineto{\pgfqpoint{11.442668in}{2.693396in}}%
\pgfpathlineto{\pgfqpoint{11.442668in}{2.682449in}}%
\pgfpathlineto{\pgfqpoint{11.387272in}{2.682449in}}%
\pgfpathlineto{\pgfqpoint{11.387272in}{2.693396in}}%
\pgfpathlineto{\pgfqpoint{11.387272in}{2.693396in}}%
\pgfpathclose%
\pgfpathmoveto{\pgfqpoint{11.496786in}{2.728135in}}%
\pgfpathquadraticcurveto{\pgfqpoint{11.487508in}{2.728135in}}{\pgfqpoint{11.482189in}{2.723166in}}%
\pgfpathquadraticcurveto{\pgfqpoint{11.476891in}{2.718197in}}{\pgfqpoint{11.476891in}{2.709518in}}%
\pgfpathquadraticcurveto{\pgfqpoint{11.476891in}{2.700818in}}{\pgfqpoint{11.482189in}{2.695849in}}%
\pgfpathquadraticcurveto{\pgfqpoint{11.487508in}{2.690881in}}{\pgfqpoint{11.496786in}{2.690881in}}%
\pgfpathquadraticcurveto{\pgfqpoint{11.506063in}{2.690881in}}{\pgfqpoint{11.511403in}{2.695870in}}%
\pgfpathquadraticcurveto{\pgfqpoint{11.516763in}{2.700880in}}{\pgfqpoint{11.516763in}{2.709518in}}%
\pgfpathquadraticcurveto{\pgfqpoint{11.516763in}{2.718197in}}{\pgfqpoint{11.511444in}{2.723166in}}%
\pgfpathquadraticcurveto{\pgfqpoint{11.506146in}{2.728135in}}{\pgfqpoint{11.496786in}{2.728135in}}%
\pgfpathlineto{\pgfqpoint{11.496786in}{2.728135in}}%
\pgfpathclose%
\pgfpathmoveto{\pgfqpoint{11.483777in}{2.733660in}}%
\pgfpathquadraticcurveto{\pgfqpoint{11.475407in}{2.735721in}}{\pgfqpoint{11.470727in}{2.741453in}}%
\pgfpathquadraticcurveto{\pgfqpoint{11.466067in}{2.747205in}}{\pgfqpoint{11.466067in}{2.755451in}}%
\pgfpathquadraticcurveto{\pgfqpoint{11.466067in}{2.766976in}}{\pgfqpoint{11.474273in}{2.773676in}}%
\pgfpathquadraticcurveto{\pgfqpoint{11.482499in}{2.780376in}}{\pgfqpoint{11.496786in}{2.780376in}}%
\pgfpathquadraticcurveto{\pgfqpoint{11.511155in}{2.780376in}}{\pgfqpoint{11.519340in}{2.773676in}}%
\pgfpathquadraticcurveto{\pgfqpoint{11.527525in}{2.766976in}}{\pgfqpoint{11.527525in}{2.755451in}}%
\pgfpathquadraticcurveto{\pgfqpoint{11.527525in}{2.747205in}}{\pgfqpoint{11.522845in}{2.741453in}}%
\pgfpathquadraticcurveto{\pgfqpoint{11.518186in}{2.735721in}}{\pgfqpoint{11.509877in}{2.733660in}}%
\pgfpathquadraticcurveto{\pgfqpoint{11.519278in}{2.731474in}}{\pgfqpoint{11.524515in}{2.725083in}}%
\pgfpathquadraticcurveto{\pgfqpoint{11.529772in}{2.718713in}}{\pgfqpoint{11.529772in}{2.709518in}}%
\pgfpathquadraticcurveto{\pgfqpoint{11.529772in}{2.695520in}}{\pgfqpoint{11.521237in}{2.688036in}}%
\pgfpathquadraticcurveto{\pgfqpoint{11.512702in}{2.680573in}}{\pgfqpoint{11.496786in}{2.680573in}}%
\pgfpathquadraticcurveto{\pgfqpoint{11.480891in}{2.680573in}}{\pgfqpoint{11.472335in}{2.688036in}}%
\pgfpathquadraticcurveto{\pgfqpoint{11.463800in}{2.695520in}}{\pgfqpoint{11.463800in}{2.709518in}}%
\pgfpathquadraticcurveto{\pgfqpoint{11.463800in}{2.718713in}}{\pgfqpoint{11.469077in}{2.725083in}}%
\pgfpathquadraticcurveto{\pgfqpoint{11.474376in}{2.731474in}}{\pgfqpoint{11.483777in}{2.733660in}}%
\pgfpathlineto{\pgfqpoint{11.483777in}{2.733660in}}%
\pgfpathclose%
\pgfpathmoveto{\pgfqpoint{11.479015in}{2.754214in}}%
\pgfpathquadraticcurveto{\pgfqpoint{11.479015in}{2.746751in}}{\pgfqpoint{11.483674in}{2.742566in}}%
\pgfpathquadraticcurveto{\pgfqpoint{11.488354in}{2.738381in}}{\pgfqpoint{11.496786in}{2.738381in}}%
\pgfpathquadraticcurveto{\pgfqpoint{11.505177in}{2.738381in}}{\pgfqpoint{11.509898in}{2.742566in}}%
\pgfpathquadraticcurveto{\pgfqpoint{11.514640in}{2.746751in}}{\pgfqpoint{11.514640in}{2.754214in}}%
\pgfpathquadraticcurveto{\pgfqpoint{11.514640in}{2.761698in}}{\pgfqpoint{11.509898in}{2.765883in}}%
\pgfpathquadraticcurveto{\pgfqpoint{11.505177in}{2.770068in}}{\pgfqpoint{11.496786in}{2.770068in}}%
\pgfpathquadraticcurveto{\pgfqpoint{11.488354in}{2.770068in}}{\pgfqpoint{11.483674in}{2.765883in}}%
\pgfpathquadraticcurveto{\pgfqpoint{11.479015in}{2.761698in}}{\pgfqpoint{11.479015in}{2.754214in}}%
\pgfpathlineto{\pgfqpoint{11.479015in}{2.754214in}}%
\pgfpathclose%
\pgfpathmoveto{\pgfqpoint{11.621638in}{2.782562in}}%
\pgfpathquadraticcurveto{\pgfqpoint{11.613021in}{2.767759in}}{\pgfqpoint{11.608815in}{2.753245in}}%
\pgfpathquadraticcurveto{\pgfqpoint{11.604630in}{2.738752in}}{\pgfqpoint{11.604630in}{2.723867in}}%
\pgfpathquadraticcurveto{\pgfqpoint{11.604630in}{2.709003in}}{\pgfqpoint{11.608856in}{2.694406in}}%
\pgfpathquadraticcurveto{\pgfqpoint{11.613082in}{2.679810in}}{\pgfqpoint{11.621638in}{2.665049in}}%
\pgfpathlineto{\pgfqpoint{11.611330in}{2.665049in}}%
\pgfpathquadraticcurveto{\pgfqpoint{11.601682in}{2.680202in}}{\pgfqpoint{11.596878in}{2.694819in}}%
\pgfpathquadraticcurveto{\pgfqpoint{11.592074in}{2.709436in}}{\pgfqpoint{11.592074in}{2.723867in}}%
\pgfpathquadraticcurveto{\pgfqpoint{11.592074in}{2.738237in}}{\pgfqpoint{11.596837in}{2.752792in}}%
\pgfpathquadraticcurveto{\pgfqpoint{11.601620in}{2.767367in}}{\pgfqpoint{11.611330in}{2.782562in}}%
\pgfpathlineto{\pgfqpoint{11.621638in}{2.782562in}}%
\pgfpathlineto{\pgfqpoint{11.621638in}{2.782562in}}%
\pgfpathclose%
\pgfpathmoveto{\pgfqpoint{11.657531in}{2.693396in}}%
\pgfpathlineto{\pgfqpoint{11.702949in}{2.693396in}}%
\pgfpathlineto{\pgfqpoint{11.702949in}{2.682449in}}%
\pgfpathlineto{\pgfqpoint{11.641883in}{2.682449in}}%
\pgfpathlineto{\pgfqpoint{11.641883in}{2.693396in}}%
\pgfpathquadraticcurveto{\pgfqpoint{11.649285in}{2.701065in}}{\pgfqpoint{11.662067in}{2.713971in}}%
\pgfpathquadraticcurveto{\pgfqpoint{11.674870in}{2.726898in}}{\pgfqpoint{11.678148in}{2.730650in}}%
\pgfpathquadraticcurveto{\pgfqpoint{11.684394in}{2.737659in}}{\pgfqpoint{11.686868in}{2.742525in}}%
\pgfpathquadraticcurveto{\pgfqpoint{11.689363in}{2.747390in}}{\pgfqpoint{11.689363in}{2.752091in}}%
\pgfpathquadraticcurveto{\pgfqpoint{11.689363in}{2.759760in}}{\pgfqpoint{11.683982in}{2.764584in}}%
\pgfpathquadraticcurveto{\pgfqpoint{11.678601in}{2.769429in}}{\pgfqpoint{11.669963in}{2.769429in}}%
\pgfpathquadraticcurveto{\pgfqpoint{11.663840in}{2.769429in}}{\pgfqpoint{11.657036in}{2.767306in}}%
\pgfpathquadraticcurveto{\pgfqpoint{11.650254in}{2.765182in}}{\pgfqpoint{11.642523in}{2.760853in}}%
\pgfpathlineto{\pgfqpoint{11.642523in}{2.774006in}}%
\pgfpathquadraticcurveto{\pgfqpoint{11.650377in}{2.777160in}}{\pgfqpoint{11.657201in}{2.778768in}}%
\pgfpathquadraticcurveto{\pgfqpoint{11.664046in}{2.780376in}}{\pgfqpoint{11.669715in}{2.780376in}}%
\pgfpathquadraticcurveto{\pgfqpoint{11.684662in}{2.780376in}}{\pgfqpoint{11.693548in}{2.772893in}}%
\pgfpathquadraticcurveto{\pgfqpoint{11.702434in}{2.765429in}}{\pgfqpoint{11.702434in}{2.752936in}}%
\pgfpathquadraticcurveto{\pgfqpoint{11.702434in}{2.746998in}}{\pgfqpoint{11.700207in}{2.741679in}}%
\pgfpathquadraticcurveto{\pgfqpoint{11.698001in}{2.736381in}}{\pgfqpoint{11.692125in}{2.729165in}}%
\pgfpathquadraticcurveto{\pgfqpoint{11.690517in}{2.727289in}}{\pgfqpoint{11.681879in}{2.718362in}}%
\pgfpathquadraticcurveto{\pgfqpoint{11.673262in}{2.709436in}}{\pgfqpoint{11.657531in}{2.693396in}}%
\pgfpathlineto{\pgfqpoint{11.657531in}{2.693396in}}%
\pgfpathclose%
\pgfpathmoveto{\pgfqpoint{11.758098in}{2.770068in}}%
\pgfpathquadraticcurveto{\pgfqpoint{11.748058in}{2.770068in}}{\pgfqpoint{11.742986in}{2.760172in}}%
\pgfpathquadraticcurveto{\pgfqpoint{11.737935in}{2.750297in}}{\pgfqpoint{11.737935in}{2.730444in}}%
\pgfpathquadraticcurveto{\pgfqpoint{11.737935in}{2.710673in}}{\pgfqpoint{11.742986in}{2.700777in}}%
\pgfpathquadraticcurveto{\pgfqpoint{11.748058in}{2.690881in}}{\pgfqpoint{11.758098in}{2.690881in}}%
\pgfpathquadraticcurveto{\pgfqpoint{11.768220in}{2.690881in}}{\pgfqpoint{11.773271in}{2.700777in}}%
\pgfpathquadraticcurveto{\pgfqpoint{11.778343in}{2.710673in}}{\pgfqpoint{11.778343in}{2.730444in}}%
\pgfpathquadraticcurveto{\pgfqpoint{11.778343in}{2.750297in}}{\pgfqpoint{11.773271in}{2.760172in}}%
\pgfpathquadraticcurveto{\pgfqpoint{11.768220in}{2.770068in}}{\pgfqpoint{11.758098in}{2.770068in}}%
\pgfpathlineto{\pgfqpoint{11.758098in}{2.770068in}}%
\pgfpathclose%
\pgfpathmoveto{\pgfqpoint{11.758098in}{2.780376in}}%
\pgfpathquadraticcurveto{\pgfqpoint{11.774281in}{2.780376in}}{\pgfqpoint{11.782817in}{2.767574in}}%
\pgfpathquadraticcurveto{\pgfqpoint{11.791352in}{2.754791in}}{\pgfqpoint{11.791352in}{2.730444in}}%
\pgfpathquadraticcurveto{\pgfqpoint{11.791352in}{2.706158in}}{\pgfqpoint{11.782817in}{2.693355in}}%
\pgfpathquadraticcurveto{\pgfqpoint{11.774281in}{2.680573in}}{\pgfqpoint{11.758098in}{2.680573in}}%
\pgfpathquadraticcurveto{\pgfqpoint{11.741934in}{2.680573in}}{\pgfqpoint{11.733399in}{2.693355in}}%
\pgfpathquadraticcurveto{\pgfqpoint{11.724864in}{2.706158in}}{\pgfqpoint{11.724864in}{2.730444in}}%
\pgfpathquadraticcurveto{\pgfqpoint{11.724864in}{2.754791in}}{\pgfqpoint{11.733399in}{2.767574in}}%
\pgfpathquadraticcurveto{\pgfqpoint{11.741934in}{2.780376in}}{\pgfqpoint{11.758098in}{2.780376in}}%
\pgfpathlineto{\pgfqpoint{11.758098in}{2.780376in}}%
\pgfpathclose%
\pgfpathmoveto{\pgfqpoint{11.825431in}{2.693396in}}%
\pgfpathlineto{\pgfqpoint{11.870848in}{2.693396in}}%
\pgfpathlineto{\pgfqpoint{11.870848in}{2.682449in}}%
\pgfpathlineto{\pgfqpoint{11.809783in}{2.682449in}}%
\pgfpathlineto{\pgfqpoint{11.809783in}{2.693396in}}%
\pgfpathquadraticcurveto{\pgfqpoint{11.817184in}{2.701065in}}{\pgfqpoint{11.829966in}{2.713971in}}%
\pgfpathquadraticcurveto{\pgfqpoint{11.842769in}{2.726898in}}{\pgfqpoint{11.846047in}{2.730650in}}%
\pgfpathquadraticcurveto{\pgfqpoint{11.852294in}{2.737659in}}{\pgfqpoint{11.854768in}{2.742525in}}%
\pgfpathquadraticcurveto{\pgfqpoint{11.857262in}{2.747390in}}{\pgfqpoint{11.857262in}{2.752091in}}%
\pgfpathquadraticcurveto{\pgfqpoint{11.857262in}{2.759760in}}{\pgfqpoint{11.851881in}{2.764584in}}%
\pgfpathquadraticcurveto{\pgfqpoint{11.846500in}{2.769429in}}{\pgfqpoint{11.837862in}{2.769429in}}%
\pgfpathquadraticcurveto{\pgfqpoint{11.831739in}{2.769429in}}{\pgfqpoint{11.824936in}{2.767306in}}%
\pgfpathquadraticcurveto{\pgfqpoint{11.818153in}{2.765182in}}{\pgfqpoint{11.810422in}{2.760853in}}%
\pgfpathlineto{\pgfqpoint{11.810422in}{2.774006in}}%
\pgfpathquadraticcurveto{\pgfqpoint{11.818277in}{2.777160in}}{\pgfqpoint{11.825101in}{2.778768in}}%
\pgfpathquadraticcurveto{\pgfqpoint{11.831945in}{2.780376in}}{\pgfqpoint{11.837615in}{2.780376in}}%
\pgfpathquadraticcurveto{\pgfqpoint{11.852562in}{2.780376in}}{\pgfqpoint{11.861447in}{2.772893in}}%
\pgfpathquadraticcurveto{\pgfqpoint{11.870333in}{2.765429in}}{\pgfqpoint{11.870333in}{2.752936in}}%
\pgfpathquadraticcurveto{\pgfqpoint{11.870333in}{2.746998in}}{\pgfqpoint{11.868106in}{2.741679in}}%
\pgfpathquadraticcurveto{\pgfqpoint{11.865900in}{2.736381in}}{\pgfqpoint{11.860025in}{2.729165in}}%
\pgfpathquadraticcurveto{\pgfqpoint{11.858417in}{2.727289in}}{\pgfqpoint{11.849778in}{2.718362in}}%
\pgfpathquadraticcurveto{\pgfqpoint{11.841161in}{2.709436in}}{\pgfqpoint{11.825431in}{2.693396in}}%
\pgfpathlineto{\pgfqpoint{11.825431in}{2.693396in}}%
\pgfpathclose%
\pgfpathmoveto{\pgfqpoint{11.933934in}{2.767306in}}%
\pgfpathlineto{\pgfqpoint{11.901072in}{2.715950in}}%
\pgfpathlineto{\pgfqpoint{11.933934in}{2.715950in}}%
\pgfpathlineto{\pgfqpoint{11.933934in}{2.767306in}}%
\pgfpathlineto{\pgfqpoint{11.933934in}{2.767306in}}%
\pgfpathclose%
\pgfpathmoveto{\pgfqpoint{11.930512in}{2.778645in}}%
\pgfpathlineto{\pgfqpoint{11.946881in}{2.778645in}}%
\pgfpathlineto{\pgfqpoint{11.946881in}{2.715950in}}%
\pgfpathlineto{\pgfqpoint{11.960612in}{2.715950in}}%
\pgfpathlineto{\pgfqpoint{11.960612in}{2.705127in}}%
\pgfpathlineto{\pgfqpoint{11.946881in}{2.705127in}}%
\pgfpathlineto{\pgfqpoint{11.946881in}{2.682449in}}%
\pgfpathlineto{\pgfqpoint{11.933934in}{2.682449in}}%
\pgfpathlineto{\pgfqpoint{11.933934in}{2.705127in}}%
\pgfpathlineto{\pgfqpoint{11.890516in}{2.705127in}}%
\pgfpathlineto{\pgfqpoint{11.890516in}{2.717682in}}%
\pgfpathlineto{\pgfqpoint{11.930512in}{2.778645in}}%
\pgfpathlineto{\pgfqpoint{11.930512in}{2.778645in}}%
\pgfpathclose%
\pgfpathmoveto{\pgfqpoint{11.978589in}{2.782562in}}%
\pgfpathlineto{\pgfqpoint{11.988897in}{2.782562in}}%
\pgfpathquadraticcurveto{\pgfqpoint{11.998546in}{2.767367in}}{\pgfqpoint{12.003349in}{2.752792in}}%
\pgfpathquadraticcurveto{\pgfqpoint{12.008153in}{2.738237in}}{\pgfqpoint{12.008153in}{2.723867in}}%
\pgfpathquadraticcurveto{\pgfqpoint{12.008153in}{2.709436in}}{\pgfqpoint{12.003349in}{2.694819in}}%
\pgfpathquadraticcurveto{\pgfqpoint{11.998546in}{2.680202in}}{\pgfqpoint{11.988897in}{2.665049in}}%
\pgfpathlineto{\pgfqpoint{11.978589in}{2.665049in}}%
\pgfpathquadraticcurveto{\pgfqpoint{11.987145in}{2.679810in}}{\pgfqpoint{11.991371in}{2.694406in}}%
\pgfpathquadraticcurveto{\pgfqpoint{11.995598in}{2.709003in}}{\pgfqpoint{11.995598in}{2.723867in}}%
\pgfpathquadraticcurveto{\pgfqpoint{11.995598in}{2.738752in}}{\pgfqpoint{11.991371in}{2.753245in}}%
\pgfpathquadraticcurveto{\pgfqpoint{11.987145in}{2.767759in}}{\pgfqpoint{11.978589in}{2.782562in}}%
\pgfpathlineto{\pgfqpoint{11.978589in}{2.782562in}}%
\pgfpathclose%
\pgfusepath{fill}%
\end{pgfscope}%
\begin{pgfscope}%
\pgfsetbuttcap%
\pgfsetroundjoin%
\definecolor{currentfill}{rgb}{0.090196,0.168627,0.227451}%
\pgfsetfillcolor{currentfill}%
\pgfsetlinewidth{0.000000pt}%
\definecolor{currentstroke}{rgb}{0.090196,0.168627,0.227451}%
\pgfsetstrokecolor{currentstroke}%
\pgfsetdash{}{0pt}%
\pgfpathmoveto{\pgfqpoint{11.018611in}{2.375890in}}%
\pgfpathlineto{\pgfqpoint{11.018611in}{2.363191in}}%
\pgfpathquadraticcurveto{\pgfqpoint{11.011210in}{2.366737in}}{\pgfqpoint{11.004633in}{2.368468in}}%
\pgfpathquadraticcurveto{\pgfqpoint{10.998056in}{2.370221in}}{\pgfqpoint{10.991933in}{2.370221in}}%
\pgfpathquadraticcurveto{\pgfqpoint{10.981316in}{2.370221in}}{\pgfqpoint{10.975543in}{2.366097in}}%
\pgfpathquadraticcurveto{\pgfqpoint{10.969771in}{2.361974in}}{\pgfqpoint{10.969771in}{2.354367in}}%
\pgfpathquadraticcurveto{\pgfqpoint{10.969771in}{2.347976in}}{\pgfqpoint{10.973605in}{2.344718in}}%
\pgfpathquadraticcurveto{\pgfqpoint{10.977440in}{2.341482in}}{\pgfqpoint{10.988140in}{2.339482in}}%
\pgfpathlineto{\pgfqpoint{10.995995in}{2.337874in}}%
\pgfpathquadraticcurveto{\pgfqpoint{11.010550in}{2.335091in}}{\pgfqpoint{11.017477in}{2.328102in}}%
\pgfpathquadraticcurveto{\pgfqpoint{11.024404in}{2.321113in}}{\pgfqpoint{11.024404in}{2.309403in}}%
\pgfpathquadraticcurveto{\pgfqpoint{11.024404in}{2.295404in}}{\pgfqpoint{11.015024in}{2.288188in}}%
\pgfpathquadraticcurveto{\pgfqpoint{11.005664in}{2.280973in}}{\pgfqpoint{10.987563in}{2.280973in}}%
\pgfpathquadraticcurveto{\pgfqpoint{10.980739in}{2.280973in}}{\pgfqpoint{10.973028in}{2.282519in}}%
\pgfpathquadraticcurveto{\pgfqpoint{10.965338in}{2.284065in}}{\pgfqpoint{10.957092in}{2.287096in}}%
\pgfpathlineto{\pgfqpoint{10.957092in}{2.300496in}}%
\pgfpathquadraticcurveto{\pgfqpoint{10.965008in}{2.296064in}}{\pgfqpoint{10.972616in}{2.293796in}}%
\pgfpathquadraticcurveto{\pgfqpoint{10.980223in}{2.291549in}}{\pgfqpoint{10.987563in}{2.291549in}}%
\pgfpathquadraticcurveto{\pgfqpoint{10.998696in}{2.291549in}}{\pgfqpoint{11.004757in}{2.295920in}}%
\pgfpathquadraticcurveto{\pgfqpoint{11.010818in}{2.300311in}}{\pgfqpoint{11.010818in}{2.308434in}}%
\pgfpathquadraticcurveto{\pgfqpoint{11.010818in}{2.315505in}}{\pgfqpoint{11.006468in}{2.319505in}}%
\pgfpathquadraticcurveto{\pgfqpoint{11.002118in}{2.323504in}}{\pgfqpoint{10.992201in}{2.325504in}}%
\pgfpathlineto{\pgfqpoint{10.984264in}{2.327050in}}%
\pgfpathquadraticcurveto{\pgfqpoint{10.969709in}{2.329936in}}{\pgfqpoint{10.963194in}{2.336121in}}%
\pgfpathquadraticcurveto{\pgfqpoint{10.956700in}{2.342306in}}{\pgfqpoint{10.956700in}{2.353336in}}%
\pgfpathquadraticcurveto{\pgfqpoint{10.956700in}{2.366097in}}{\pgfqpoint{10.965689in}{2.373437in}}%
\pgfpathquadraticcurveto{\pgfqpoint{10.974678in}{2.380776in}}{\pgfqpoint{10.990449in}{2.380776in}}%
\pgfpathquadraticcurveto{\pgfqpoint{10.997232in}{2.380776in}}{\pgfqpoint{11.004241in}{2.379539in}}%
\pgfpathquadraticcurveto{\pgfqpoint{11.011271in}{2.378323in}}{\pgfqpoint{11.018611in}{2.375890in}}%
\pgfpathlineto{\pgfqpoint{11.018611in}{2.375890in}}%
\pgfpathclose%
\pgfpathmoveto{\pgfqpoint{11.055926in}{2.375499in}}%
\pgfpathlineto{\pgfqpoint{11.055926in}{2.355006in}}%
\pgfpathlineto{\pgfqpoint{11.080336in}{2.355006in}}%
\pgfpathlineto{\pgfqpoint{11.080336in}{2.345790in}}%
\pgfpathlineto{\pgfqpoint{11.055926in}{2.345790in}}%
\pgfpathlineto{\pgfqpoint{11.055926in}{2.306619in}}%
\pgfpathquadraticcurveto{\pgfqpoint{11.055926in}{2.297796in}}{\pgfqpoint{11.058339in}{2.295280in}}%
\pgfpathquadraticcurveto{\pgfqpoint{11.060751in}{2.292765in}}{\pgfqpoint{11.068173in}{2.292765in}}%
\pgfpathlineto{\pgfqpoint{11.080336in}{2.292765in}}%
\pgfpathlineto{\pgfqpoint{11.080336in}{2.282849in}}%
\pgfpathlineto{\pgfqpoint{11.068173in}{2.282849in}}%
\pgfpathquadraticcurveto{\pgfqpoint{11.054442in}{2.282849in}}{\pgfqpoint{11.049226in}{2.287962in}}%
\pgfpathquadraticcurveto{\pgfqpoint{11.044010in}{2.293095in}}{\pgfqpoint{11.044010in}{2.306619in}}%
\pgfpathlineto{\pgfqpoint{11.044010in}{2.345790in}}%
\pgfpathlineto{\pgfqpoint{11.035310in}{2.345790in}}%
\pgfpathlineto{\pgfqpoint{11.035310in}{2.355006in}}%
\pgfpathlineto{\pgfqpoint{11.044010in}{2.355006in}}%
\pgfpathlineto{\pgfqpoint{11.044010in}{2.375499in}}%
\pgfpathlineto{\pgfqpoint{11.055926in}{2.375499in}}%
\pgfpathlineto{\pgfqpoint{11.055926in}{2.375499in}}%
\pgfpathclose%
\pgfpathmoveto{\pgfqpoint{11.094706in}{2.311320in}}%
\pgfpathlineto{\pgfqpoint{11.094706in}{2.355006in}}%
\pgfpathlineto{\pgfqpoint{11.106560in}{2.355006in}}%
\pgfpathlineto{\pgfqpoint{11.106560in}{2.311773in}}%
\pgfpathquadraticcurveto{\pgfqpoint{11.106560in}{2.301527in}}{\pgfqpoint{11.110539in}{2.296394in}}%
\pgfpathquadraticcurveto{\pgfqpoint{11.114539in}{2.291281in}}{\pgfqpoint{11.122538in}{2.291281in}}%
\pgfpathquadraticcurveto{\pgfqpoint{11.132124in}{2.291281in}}{\pgfqpoint{11.137691in}{2.297404in}}%
\pgfpathquadraticcurveto{\pgfqpoint{11.143278in}{2.303527in}}{\pgfqpoint{11.143278in}{2.314103in}}%
\pgfpathlineto{\pgfqpoint{11.143278in}{2.355006in}}%
\pgfpathlineto{\pgfqpoint{11.155132in}{2.355006in}}%
\pgfpathlineto{\pgfqpoint{11.155132in}{2.282849in}}%
\pgfpathlineto{\pgfqpoint{11.143278in}{2.282849in}}%
\pgfpathlineto{\pgfqpoint{11.143278in}{2.293940in}}%
\pgfpathquadraticcurveto{\pgfqpoint{11.138969in}{2.287364in}}{\pgfqpoint{11.133258in}{2.284168in}}%
\pgfpathquadraticcurveto{\pgfqpoint{11.127568in}{2.280973in}}{\pgfqpoint{11.120023in}{2.280973in}}%
\pgfpathquadraticcurveto{\pgfqpoint{11.107591in}{2.280973in}}{\pgfqpoint{11.101138in}{2.288704in}}%
\pgfpathquadraticcurveto{\pgfqpoint{11.094706in}{2.296435in}}{\pgfqpoint{11.094706in}{2.311320in}}%
\pgfpathlineto{\pgfqpoint{11.094706in}{2.311320in}}%
\pgfpathclose%
\pgfpathmoveto{\pgfqpoint{11.124538in}{2.356738in}}%
\pgfpathlineto{\pgfqpoint{11.124538in}{2.356738in}}%
\pgfpathlineto{\pgfqpoint{11.124538in}{2.356738in}}%
\pgfpathclose%
\pgfpathmoveto{\pgfqpoint{11.227021in}{2.344059in}}%
\pgfpathlineto{\pgfqpoint{11.227021in}{2.383106in}}%
\pgfpathlineto{\pgfqpoint{11.238876in}{2.383106in}}%
\pgfpathlineto{\pgfqpoint{11.238876in}{2.282849in}}%
\pgfpathlineto{\pgfqpoint{11.227021in}{2.282849in}}%
\pgfpathlineto{\pgfqpoint{11.227021in}{2.293672in}}%
\pgfpathquadraticcurveto{\pgfqpoint{11.223290in}{2.287240in}}{\pgfqpoint{11.217579in}{2.284106in}}%
\pgfpathquadraticcurveto{\pgfqpoint{11.211889in}{2.280973in}}{\pgfqpoint{11.203890in}{2.280973in}}%
\pgfpathquadraticcurveto{\pgfqpoint{11.190819in}{2.280973in}}{\pgfqpoint{11.182593in}{2.291405in}}%
\pgfpathquadraticcurveto{\pgfqpoint{11.174388in}{2.301857in}}{\pgfqpoint{11.174388in}{2.318866in}}%
\pgfpathquadraticcurveto{\pgfqpoint{11.174388in}{2.335874in}}{\pgfqpoint{11.182593in}{2.346306in}}%
\pgfpathquadraticcurveto{\pgfqpoint{11.190819in}{2.356738in}}{\pgfqpoint{11.203890in}{2.356738in}}%
\pgfpathquadraticcurveto{\pgfqpoint{11.211889in}{2.356738in}}{\pgfqpoint{11.217579in}{2.353604in}}%
\pgfpathquadraticcurveto{\pgfqpoint{11.223290in}{2.350491in}}{\pgfqpoint{11.227021in}{2.344059in}}%
\pgfpathlineto{\pgfqpoint{11.227021in}{2.344059in}}%
\pgfpathclose%
\pgfpathmoveto{\pgfqpoint{11.186634in}{2.318866in}}%
\pgfpathquadraticcurveto{\pgfqpoint{11.186634in}{2.305795in}}{\pgfqpoint{11.192015in}{2.298352in}}%
\pgfpathquadraticcurveto{\pgfqpoint{11.197396in}{2.290910in}}{\pgfqpoint{11.206797in}{2.290910in}}%
\pgfpathquadraticcurveto{\pgfqpoint{11.216198in}{2.290910in}}{\pgfqpoint{11.221599in}{2.298352in}}%
\pgfpathquadraticcurveto{\pgfqpoint{11.227021in}{2.305795in}}{\pgfqpoint{11.227021in}{2.318866in}}%
\pgfpathquadraticcurveto{\pgfqpoint{11.227021in}{2.331936in}}{\pgfqpoint{11.221599in}{2.339379in}}%
\pgfpathquadraticcurveto{\pgfqpoint{11.216198in}{2.346821in}}{\pgfqpoint{11.206797in}{2.346821in}}%
\pgfpathquadraticcurveto{\pgfqpoint{11.197396in}{2.346821in}}{\pgfqpoint{11.192015in}{2.339379in}}%
\pgfpathquadraticcurveto{\pgfqpoint{11.186634in}{2.331936in}}{\pgfqpoint{11.186634in}{2.318866in}}%
\pgfpathlineto{\pgfqpoint{11.186634in}{2.318866in}}%
\pgfpathclose%
\pgfpathmoveto{\pgfqpoint{11.293323in}{2.276148in}}%
\pgfpathquadraticcurveto{\pgfqpoint{11.288314in}{2.263263in}}{\pgfqpoint{11.283531in}{2.259346in}}%
\pgfpathquadraticcurveto{\pgfqpoint{11.278768in}{2.255408in}}{\pgfqpoint{11.270790in}{2.255408in}}%
\pgfpathlineto{\pgfqpoint{11.261306in}{2.255408in}}%
\pgfpathlineto{\pgfqpoint{11.261306in}{2.265325in}}%
\pgfpathlineto{\pgfqpoint{11.268275in}{2.265325in}}%
\pgfpathquadraticcurveto{\pgfqpoint{11.273161in}{2.265325in}}{\pgfqpoint{11.275861in}{2.267655in}}%
\pgfpathquadraticcurveto{\pgfqpoint{11.278583in}{2.269964in}}{\pgfqpoint{11.281861in}{2.278602in}}%
\pgfpathlineto{\pgfqpoint{11.283984in}{2.284003in}}%
\pgfpathlineto{\pgfqpoint{11.254812in}{2.355006in}}%
\pgfpathlineto{\pgfqpoint{11.267367in}{2.355006in}}%
\pgfpathlineto{\pgfqpoint{11.289922in}{2.298579in}}%
\pgfpathlineto{\pgfqpoint{11.312476in}{2.355006in}}%
\pgfpathlineto{\pgfqpoint{11.325031in}{2.355006in}}%
\pgfpathlineto{\pgfqpoint{11.293323in}{2.276148in}}%
\pgfpathlineto{\pgfqpoint{11.293323in}{2.276148in}}%
\pgfpathclose%
\pgfpathmoveto{\pgfqpoint{11.396219in}{2.293796in}}%
\pgfpathlineto{\pgfqpoint{11.441637in}{2.293796in}}%
\pgfpathlineto{\pgfqpoint{11.441637in}{2.282849in}}%
\pgfpathlineto{\pgfqpoint{11.380572in}{2.282849in}}%
\pgfpathlineto{\pgfqpoint{11.380572in}{2.293796in}}%
\pgfpathquadraticcurveto{\pgfqpoint{11.387973in}{2.301465in}}{\pgfqpoint{11.400755in}{2.314371in}}%
\pgfpathquadraticcurveto{\pgfqpoint{11.413558in}{2.327298in}}{\pgfqpoint{11.416836in}{2.331050in}}%
\pgfpathquadraticcurveto{\pgfqpoint{11.423082in}{2.338059in}}{\pgfqpoint{11.425556in}{2.342925in}}%
\pgfpathquadraticcurveto{\pgfqpoint{11.428051in}{2.347790in}}{\pgfqpoint{11.428051in}{2.352491in}}%
\pgfpathquadraticcurveto{\pgfqpoint{11.428051in}{2.360160in}}{\pgfqpoint{11.422670in}{2.364984in}}%
\pgfpathquadraticcurveto{\pgfqpoint{11.417289in}{2.369829in}}{\pgfqpoint{11.408651in}{2.369829in}}%
\pgfpathquadraticcurveto{\pgfqpoint{11.402528in}{2.369829in}}{\pgfqpoint{11.395725in}{2.367706in}}%
\pgfpathquadraticcurveto{\pgfqpoint{11.388942in}{2.365582in}}{\pgfqpoint{11.381211in}{2.361253in}}%
\pgfpathlineto{\pgfqpoint{11.381211in}{2.374406in}}%
\pgfpathquadraticcurveto{\pgfqpoint{11.389066in}{2.377560in}}{\pgfqpoint{11.395890in}{2.379168in}}%
\pgfpathquadraticcurveto{\pgfqpoint{11.402734in}{2.380776in}}{\pgfqpoint{11.408404in}{2.380776in}}%
\pgfpathquadraticcurveto{\pgfqpoint{11.423350in}{2.380776in}}{\pgfqpoint{11.432236in}{2.373293in}}%
\pgfpathquadraticcurveto{\pgfqpoint{11.441122in}{2.365829in}}{\pgfqpoint{11.441122in}{2.353336in}}%
\pgfpathquadraticcurveto{\pgfqpoint{11.441122in}{2.347398in}}{\pgfqpoint{11.438895in}{2.342079in}}%
\pgfpathquadraticcurveto{\pgfqpoint{11.436689in}{2.336781in}}{\pgfqpoint{11.430814in}{2.329565in}}%
\pgfpathquadraticcurveto{\pgfqpoint{11.429206in}{2.327689in}}{\pgfqpoint{11.420567in}{2.318762in}}%
\pgfpathquadraticcurveto{\pgfqpoint{11.411950in}{2.309836in}}{\pgfqpoint{11.396219in}{2.293796in}}%
\pgfpathlineto{\pgfqpoint{11.396219in}{2.293796in}}%
\pgfpathclose%
\pgfpathmoveto{\pgfqpoint{11.496786in}{2.370468in}}%
\pgfpathquadraticcurveto{\pgfqpoint{11.486746in}{2.370468in}}{\pgfqpoint{11.481674in}{2.360572in}}%
\pgfpathquadraticcurveto{\pgfqpoint{11.476623in}{2.350697in}}{\pgfqpoint{11.476623in}{2.330844in}}%
\pgfpathquadraticcurveto{\pgfqpoint{11.476623in}{2.311073in}}{\pgfqpoint{11.481674in}{2.301177in}}%
\pgfpathquadraticcurveto{\pgfqpoint{11.486746in}{2.291281in}}{\pgfqpoint{11.496786in}{2.291281in}}%
\pgfpathquadraticcurveto{\pgfqpoint{11.506908in}{2.291281in}}{\pgfqpoint{11.511959in}{2.301177in}}%
\pgfpathquadraticcurveto{\pgfqpoint{11.517031in}{2.311073in}}{\pgfqpoint{11.517031in}{2.330844in}}%
\pgfpathquadraticcurveto{\pgfqpoint{11.517031in}{2.350697in}}{\pgfqpoint{11.511959in}{2.360572in}}%
\pgfpathquadraticcurveto{\pgfqpoint{11.506908in}{2.370468in}}{\pgfqpoint{11.496786in}{2.370468in}}%
\pgfpathlineto{\pgfqpoint{11.496786in}{2.370468in}}%
\pgfpathclose%
\pgfpathmoveto{\pgfqpoint{11.496786in}{2.380776in}}%
\pgfpathquadraticcurveto{\pgfqpoint{11.512970in}{2.380776in}}{\pgfqpoint{11.521505in}{2.367974in}}%
\pgfpathquadraticcurveto{\pgfqpoint{11.530040in}{2.355191in}}{\pgfqpoint{11.530040in}{2.330844in}}%
\pgfpathquadraticcurveto{\pgfqpoint{11.530040in}{2.306558in}}{\pgfqpoint{11.521505in}{2.293755in}}%
\pgfpathquadraticcurveto{\pgfqpoint{11.512970in}{2.280973in}}{\pgfqpoint{11.496786in}{2.280973in}}%
\pgfpathquadraticcurveto{\pgfqpoint{11.480623in}{2.280973in}}{\pgfqpoint{11.472087in}{2.293755in}}%
\pgfpathquadraticcurveto{\pgfqpoint{11.463552in}{2.306558in}}{\pgfqpoint{11.463552in}{2.330844in}}%
\pgfpathquadraticcurveto{\pgfqpoint{11.463552in}{2.355191in}}{\pgfqpoint{11.472087in}{2.367974in}}%
\pgfpathquadraticcurveto{\pgfqpoint{11.480623in}{2.380776in}}{\pgfqpoint{11.496786in}{2.380776in}}%
\pgfpathlineto{\pgfqpoint{11.496786in}{2.380776in}}%
\pgfpathclose%
\pgfpathmoveto{\pgfqpoint{11.621638in}{2.382962in}}%
\pgfpathquadraticcurveto{\pgfqpoint{11.613021in}{2.368159in}}{\pgfqpoint{11.608815in}{2.353645in}}%
\pgfpathquadraticcurveto{\pgfqpoint{11.604630in}{2.339152in}}{\pgfqpoint{11.604630in}{2.324267in}}%
\pgfpathquadraticcurveto{\pgfqpoint{11.604630in}{2.309403in}}{\pgfqpoint{11.608856in}{2.294806in}}%
\pgfpathquadraticcurveto{\pgfqpoint{11.613082in}{2.280210in}}{\pgfqpoint{11.621638in}{2.265449in}}%
\pgfpathlineto{\pgfqpoint{11.611330in}{2.265449in}}%
\pgfpathquadraticcurveto{\pgfqpoint{11.601682in}{2.280602in}}{\pgfqpoint{11.596878in}{2.295219in}}%
\pgfpathquadraticcurveto{\pgfqpoint{11.592074in}{2.309836in}}{\pgfqpoint{11.592074in}{2.324267in}}%
\pgfpathquadraticcurveto{\pgfqpoint{11.592074in}{2.338637in}}{\pgfqpoint{11.596837in}{2.353192in}}%
\pgfpathquadraticcurveto{\pgfqpoint{11.601620in}{2.367767in}}{\pgfqpoint{11.611330in}{2.382962in}}%
\pgfpathlineto{\pgfqpoint{11.621638in}{2.382962in}}%
\pgfpathlineto{\pgfqpoint{11.621638in}{2.382962in}}%
\pgfpathclose%
\pgfpathmoveto{\pgfqpoint{11.657531in}{2.293796in}}%
\pgfpathlineto{\pgfqpoint{11.702949in}{2.293796in}}%
\pgfpathlineto{\pgfqpoint{11.702949in}{2.282849in}}%
\pgfpathlineto{\pgfqpoint{11.641883in}{2.282849in}}%
\pgfpathlineto{\pgfqpoint{11.641883in}{2.293796in}}%
\pgfpathquadraticcurveto{\pgfqpoint{11.649285in}{2.301465in}}{\pgfqpoint{11.662067in}{2.314371in}}%
\pgfpathquadraticcurveto{\pgfqpoint{11.674870in}{2.327298in}}{\pgfqpoint{11.678148in}{2.331050in}}%
\pgfpathquadraticcurveto{\pgfqpoint{11.684394in}{2.338059in}}{\pgfqpoint{11.686868in}{2.342925in}}%
\pgfpathquadraticcurveto{\pgfqpoint{11.689363in}{2.347790in}}{\pgfqpoint{11.689363in}{2.352491in}}%
\pgfpathquadraticcurveto{\pgfqpoint{11.689363in}{2.360160in}}{\pgfqpoint{11.683982in}{2.364984in}}%
\pgfpathquadraticcurveto{\pgfqpoint{11.678601in}{2.369829in}}{\pgfqpoint{11.669963in}{2.369829in}}%
\pgfpathquadraticcurveto{\pgfqpoint{11.663840in}{2.369829in}}{\pgfqpoint{11.657036in}{2.367706in}}%
\pgfpathquadraticcurveto{\pgfqpoint{11.650254in}{2.365582in}}{\pgfqpoint{11.642523in}{2.361253in}}%
\pgfpathlineto{\pgfqpoint{11.642523in}{2.374406in}}%
\pgfpathquadraticcurveto{\pgfqpoint{11.650377in}{2.377560in}}{\pgfqpoint{11.657201in}{2.379168in}}%
\pgfpathquadraticcurveto{\pgfqpoint{11.664046in}{2.380776in}}{\pgfqpoint{11.669715in}{2.380776in}}%
\pgfpathquadraticcurveto{\pgfqpoint{11.684662in}{2.380776in}}{\pgfqpoint{11.693548in}{2.373293in}}%
\pgfpathquadraticcurveto{\pgfqpoint{11.702434in}{2.365829in}}{\pgfqpoint{11.702434in}{2.353336in}}%
\pgfpathquadraticcurveto{\pgfqpoint{11.702434in}{2.347398in}}{\pgfqpoint{11.700207in}{2.342079in}}%
\pgfpathquadraticcurveto{\pgfqpoint{11.698001in}{2.336781in}}{\pgfqpoint{11.692125in}{2.329565in}}%
\pgfpathquadraticcurveto{\pgfqpoint{11.690517in}{2.327689in}}{\pgfqpoint{11.681879in}{2.318762in}}%
\pgfpathquadraticcurveto{\pgfqpoint{11.673262in}{2.309836in}}{\pgfqpoint{11.657531in}{2.293796in}}%
\pgfpathlineto{\pgfqpoint{11.657531in}{2.293796in}}%
\pgfpathclose%
\pgfpathmoveto{\pgfqpoint{11.758098in}{2.370468in}}%
\pgfpathquadraticcurveto{\pgfqpoint{11.748058in}{2.370468in}}{\pgfqpoint{11.742986in}{2.360572in}}%
\pgfpathquadraticcurveto{\pgfqpoint{11.737935in}{2.350697in}}{\pgfqpoint{11.737935in}{2.330844in}}%
\pgfpathquadraticcurveto{\pgfqpoint{11.737935in}{2.311073in}}{\pgfqpoint{11.742986in}{2.301177in}}%
\pgfpathquadraticcurveto{\pgfqpoint{11.748058in}{2.291281in}}{\pgfqpoint{11.758098in}{2.291281in}}%
\pgfpathquadraticcurveto{\pgfqpoint{11.768220in}{2.291281in}}{\pgfqpoint{11.773271in}{2.301177in}}%
\pgfpathquadraticcurveto{\pgfqpoint{11.778343in}{2.311073in}}{\pgfqpoint{11.778343in}{2.330844in}}%
\pgfpathquadraticcurveto{\pgfqpoint{11.778343in}{2.350697in}}{\pgfqpoint{11.773271in}{2.360572in}}%
\pgfpathquadraticcurveto{\pgfqpoint{11.768220in}{2.370468in}}{\pgfqpoint{11.758098in}{2.370468in}}%
\pgfpathlineto{\pgfqpoint{11.758098in}{2.370468in}}%
\pgfpathclose%
\pgfpathmoveto{\pgfqpoint{11.758098in}{2.380776in}}%
\pgfpathquadraticcurveto{\pgfqpoint{11.774281in}{2.380776in}}{\pgfqpoint{11.782817in}{2.367974in}}%
\pgfpathquadraticcurveto{\pgfqpoint{11.791352in}{2.355191in}}{\pgfqpoint{11.791352in}{2.330844in}}%
\pgfpathquadraticcurveto{\pgfqpoint{11.791352in}{2.306558in}}{\pgfqpoint{11.782817in}{2.293755in}}%
\pgfpathquadraticcurveto{\pgfqpoint{11.774281in}{2.280973in}}{\pgfqpoint{11.758098in}{2.280973in}}%
\pgfpathquadraticcurveto{\pgfqpoint{11.741934in}{2.280973in}}{\pgfqpoint{11.733399in}{2.293755in}}%
\pgfpathquadraticcurveto{\pgfqpoint{11.724864in}{2.306558in}}{\pgfqpoint{11.724864in}{2.330844in}}%
\pgfpathquadraticcurveto{\pgfqpoint{11.724864in}{2.355191in}}{\pgfqpoint{11.733399in}{2.367974in}}%
\pgfpathquadraticcurveto{\pgfqpoint{11.741934in}{2.380776in}}{\pgfqpoint{11.758098in}{2.380776in}}%
\pgfpathlineto{\pgfqpoint{11.758098in}{2.380776in}}%
\pgfpathclose%
\pgfpathmoveto{\pgfqpoint{11.825431in}{2.293796in}}%
\pgfpathlineto{\pgfqpoint{11.870848in}{2.293796in}}%
\pgfpathlineto{\pgfqpoint{11.870848in}{2.282849in}}%
\pgfpathlineto{\pgfqpoint{11.809783in}{2.282849in}}%
\pgfpathlineto{\pgfqpoint{11.809783in}{2.293796in}}%
\pgfpathquadraticcurveto{\pgfqpoint{11.817184in}{2.301465in}}{\pgfqpoint{11.829966in}{2.314371in}}%
\pgfpathquadraticcurveto{\pgfqpoint{11.842769in}{2.327298in}}{\pgfqpoint{11.846047in}{2.331050in}}%
\pgfpathquadraticcurveto{\pgfqpoint{11.852294in}{2.338059in}}{\pgfqpoint{11.854768in}{2.342925in}}%
\pgfpathquadraticcurveto{\pgfqpoint{11.857262in}{2.347790in}}{\pgfqpoint{11.857262in}{2.352491in}}%
\pgfpathquadraticcurveto{\pgfqpoint{11.857262in}{2.360160in}}{\pgfqpoint{11.851881in}{2.364984in}}%
\pgfpathquadraticcurveto{\pgfqpoint{11.846500in}{2.369829in}}{\pgfqpoint{11.837862in}{2.369829in}}%
\pgfpathquadraticcurveto{\pgfqpoint{11.831739in}{2.369829in}}{\pgfqpoint{11.824936in}{2.367706in}}%
\pgfpathquadraticcurveto{\pgfqpoint{11.818153in}{2.365582in}}{\pgfqpoint{11.810422in}{2.361253in}}%
\pgfpathlineto{\pgfqpoint{11.810422in}{2.374406in}}%
\pgfpathquadraticcurveto{\pgfqpoint{11.818277in}{2.377560in}}{\pgfqpoint{11.825101in}{2.379168in}}%
\pgfpathquadraticcurveto{\pgfqpoint{11.831945in}{2.380776in}}{\pgfqpoint{11.837615in}{2.380776in}}%
\pgfpathquadraticcurveto{\pgfqpoint{11.852562in}{2.380776in}}{\pgfqpoint{11.861447in}{2.373293in}}%
\pgfpathquadraticcurveto{\pgfqpoint{11.870333in}{2.365829in}}{\pgfqpoint{11.870333in}{2.353336in}}%
\pgfpathquadraticcurveto{\pgfqpoint{11.870333in}{2.347398in}}{\pgfqpoint{11.868106in}{2.342079in}}%
\pgfpathquadraticcurveto{\pgfqpoint{11.865900in}{2.336781in}}{\pgfqpoint{11.860025in}{2.329565in}}%
\pgfpathquadraticcurveto{\pgfqpoint{11.858417in}{2.327689in}}{\pgfqpoint{11.849778in}{2.318762in}}%
\pgfpathquadraticcurveto{\pgfqpoint{11.841161in}{2.309836in}}{\pgfqpoint{11.825431in}{2.293796in}}%
\pgfpathlineto{\pgfqpoint{11.825431in}{2.293796in}}%
\pgfpathclose%
\pgfpathmoveto{\pgfqpoint{11.933934in}{2.367706in}}%
\pgfpathlineto{\pgfqpoint{11.901072in}{2.316350in}}%
\pgfpathlineto{\pgfqpoint{11.933934in}{2.316350in}}%
\pgfpathlineto{\pgfqpoint{11.933934in}{2.367706in}}%
\pgfpathlineto{\pgfqpoint{11.933934in}{2.367706in}}%
\pgfpathclose%
\pgfpathmoveto{\pgfqpoint{11.930512in}{2.379045in}}%
\pgfpathlineto{\pgfqpoint{11.946881in}{2.379045in}}%
\pgfpathlineto{\pgfqpoint{11.946881in}{2.316350in}}%
\pgfpathlineto{\pgfqpoint{11.960612in}{2.316350in}}%
\pgfpathlineto{\pgfqpoint{11.960612in}{2.305527in}}%
\pgfpathlineto{\pgfqpoint{11.946881in}{2.305527in}}%
\pgfpathlineto{\pgfqpoint{11.946881in}{2.282849in}}%
\pgfpathlineto{\pgfqpoint{11.933934in}{2.282849in}}%
\pgfpathlineto{\pgfqpoint{11.933934in}{2.305527in}}%
\pgfpathlineto{\pgfqpoint{11.890516in}{2.305527in}}%
\pgfpathlineto{\pgfqpoint{11.890516in}{2.318082in}}%
\pgfpathlineto{\pgfqpoint{11.930512in}{2.379045in}}%
\pgfpathlineto{\pgfqpoint{11.930512in}{2.379045in}}%
\pgfpathclose%
\pgfpathmoveto{\pgfqpoint{11.978589in}{2.382962in}}%
\pgfpathlineto{\pgfqpoint{11.988897in}{2.382962in}}%
\pgfpathquadraticcurveto{\pgfqpoint{11.998546in}{2.367767in}}{\pgfqpoint{12.003349in}{2.353192in}}%
\pgfpathquadraticcurveto{\pgfqpoint{12.008153in}{2.338637in}}{\pgfqpoint{12.008153in}{2.324267in}}%
\pgfpathquadraticcurveto{\pgfqpoint{12.008153in}{2.309836in}}{\pgfqpoint{12.003349in}{2.295219in}}%
\pgfpathquadraticcurveto{\pgfqpoint{11.998546in}{2.280602in}}{\pgfqpoint{11.988897in}{2.265449in}}%
\pgfpathlineto{\pgfqpoint{11.978589in}{2.265449in}}%
\pgfpathquadraticcurveto{\pgfqpoint{11.987145in}{2.280210in}}{\pgfqpoint{11.991371in}{2.294806in}}%
\pgfpathquadraticcurveto{\pgfqpoint{11.995598in}{2.309403in}}{\pgfqpoint{11.995598in}{2.324267in}}%
\pgfpathquadraticcurveto{\pgfqpoint{11.995598in}{2.339152in}}{\pgfqpoint{11.991371in}{2.353645in}}%
\pgfpathquadraticcurveto{\pgfqpoint{11.987145in}{2.368159in}}{\pgfqpoint{11.978589in}{2.382962in}}%
\pgfpathlineto{\pgfqpoint{11.978589in}{2.382962in}}%
\pgfpathclose%
\pgfusepath{fill}%
\end{pgfscope}%
\begin{pgfscope}%
\pgfsetbuttcap%
\pgfsetroundjoin%
\definecolor{currentfill}{rgb}{0.090196,0.168627,0.227451}%
\pgfsetfillcolor{currentfill}%
\pgfsetlinewidth{0.000000pt}%
\definecolor{currentstroke}{rgb}{0.090196,0.168627,0.227451}%
\pgfsetstrokecolor{currentstroke}%
\pgfsetdash{}{0pt}%
\pgfpathmoveto{\pgfqpoint{11.018611in}{1.976290in}}%
\pgfpathlineto{\pgfqpoint{11.018611in}{1.963591in}}%
\pgfpathquadraticcurveto{\pgfqpoint{11.011210in}{1.967137in}}{\pgfqpoint{11.004633in}{1.968868in}}%
\pgfpathquadraticcurveto{\pgfqpoint{10.998056in}{1.970621in}}{\pgfqpoint{10.991933in}{1.970621in}}%
\pgfpathquadraticcurveto{\pgfqpoint{10.981316in}{1.970621in}}{\pgfqpoint{10.975543in}{1.966497in}}%
\pgfpathquadraticcurveto{\pgfqpoint{10.969771in}{1.962374in}}{\pgfqpoint{10.969771in}{1.954767in}}%
\pgfpathquadraticcurveto{\pgfqpoint{10.969771in}{1.948376in}}{\pgfqpoint{10.973605in}{1.945118in}}%
\pgfpathquadraticcurveto{\pgfqpoint{10.977440in}{1.941882in}}{\pgfqpoint{10.988140in}{1.939882in}}%
\pgfpathlineto{\pgfqpoint{10.995995in}{1.938274in}}%
\pgfpathquadraticcurveto{\pgfqpoint{11.010550in}{1.935491in}}{\pgfqpoint{11.017477in}{1.928502in}}%
\pgfpathquadraticcurveto{\pgfqpoint{11.024404in}{1.921513in}}{\pgfqpoint{11.024404in}{1.909803in}}%
\pgfpathquadraticcurveto{\pgfqpoint{11.024404in}{1.895804in}}{\pgfqpoint{11.015024in}{1.888588in}}%
\pgfpathquadraticcurveto{\pgfqpoint{11.005664in}{1.881373in}}{\pgfqpoint{10.987563in}{1.881373in}}%
\pgfpathquadraticcurveto{\pgfqpoint{10.980739in}{1.881373in}}{\pgfqpoint{10.973028in}{1.882919in}}%
\pgfpathquadraticcurveto{\pgfqpoint{10.965338in}{1.884465in}}{\pgfqpoint{10.957092in}{1.887496in}}%
\pgfpathlineto{\pgfqpoint{10.957092in}{1.900896in}}%
\pgfpathquadraticcurveto{\pgfqpoint{10.965008in}{1.896464in}}{\pgfqpoint{10.972616in}{1.894196in}}%
\pgfpathquadraticcurveto{\pgfqpoint{10.980223in}{1.891949in}}{\pgfqpoint{10.987563in}{1.891949in}}%
\pgfpathquadraticcurveto{\pgfqpoint{10.998696in}{1.891949in}}{\pgfqpoint{11.004757in}{1.896320in}}%
\pgfpathquadraticcurveto{\pgfqpoint{11.010818in}{1.900711in}}{\pgfqpoint{11.010818in}{1.908834in}}%
\pgfpathquadraticcurveto{\pgfqpoint{11.010818in}{1.915905in}}{\pgfqpoint{11.006468in}{1.919905in}}%
\pgfpathquadraticcurveto{\pgfqpoint{11.002118in}{1.923904in}}{\pgfqpoint{10.992201in}{1.925904in}}%
\pgfpathlineto{\pgfqpoint{10.984264in}{1.927450in}}%
\pgfpathquadraticcurveto{\pgfqpoint{10.969709in}{1.930336in}}{\pgfqpoint{10.963194in}{1.936521in}}%
\pgfpathquadraticcurveto{\pgfqpoint{10.956700in}{1.942706in}}{\pgfqpoint{10.956700in}{1.953736in}}%
\pgfpathquadraticcurveto{\pgfqpoint{10.956700in}{1.966497in}}{\pgfqpoint{10.965689in}{1.973837in}}%
\pgfpathquadraticcurveto{\pgfqpoint{10.974678in}{1.981176in}}{\pgfqpoint{10.990449in}{1.981176in}}%
\pgfpathquadraticcurveto{\pgfqpoint{10.997232in}{1.981176in}}{\pgfqpoint{11.004241in}{1.979939in}}%
\pgfpathquadraticcurveto{\pgfqpoint{11.011271in}{1.978723in}}{\pgfqpoint{11.018611in}{1.976290in}}%
\pgfpathlineto{\pgfqpoint{11.018611in}{1.976290in}}%
\pgfpathclose%
\pgfpathmoveto{\pgfqpoint{11.055926in}{1.975899in}}%
\pgfpathlineto{\pgfqpoint{11.055926in}{1.955406in}}%
\pgfpathlineto{\pgfqpoint{11.080336in}{1.955406in}}%
\pgfpathlineto{\pgfqpoint{11.080336in}{1.946190in}}%
\pgfpathlineto{\pgfqpoint{11.055926in}{1.946190in}}%
\pgfpathlineto{\pgfqpoint{11.055926in}{1.907019in}}%
\pgfpathquadraticcurveto{\pgfqpoint{11.055926in}{1.898196in}}{\pgfqpoint{11.058339in}{1.895680in}}%
\pgfpathquadraticcurveto{\pgfqpoint{11.060751in}{1.893165in}}{\pgfqpoint{11.068173in}{1.893165in}}%
\pgfpathlineto{\pgfqpoint{11.080336in}{1.893165in}}%
\pgfpathlineto{\pgfqpoint{11.080336in}{1.883249in}}%
\pgfpathlineto{\pgfqpoint{11.068173in}{1.883249in}}%
\pgfpathquadraticcurveto{\pgfqpoint{11.054442in}{1.883249in}}{\pgfqpoint{11.049226in}{1.888362in}}%
\pgfpathquadraticcurveto{\pgfqpoint{11.044010in}{1.893495in}}{\pgfqpoint{11.044010in}{1.907019in}}%
\pgfpathlineto{\pgfqpoint{11.044010in}{1.946190in}}%
\pgfpathlineto{\pgfqpoint{11.035310in}{1.946190in}}%
\pgfpathlineto{\pgfqpoint{11.035310in}{1.955406in}}%
\pgfpathlineto{\pgfqpoint{11.044010in}{1.955406in}}%
\pgfpathlineto{\pgfqpoint{11.044010in}{1.975899in}}%
\pgfpathlineto{\pgfqpoint{11.055926in}{1.975899in}}%
\pgfpathlineto{\pgfqpoint{11.055926in}{1.975899in}}%
\pgfpathclose%
\pgfpathmoveto{\pgfqpoint{11.094706in}{1.911720in}}%
\pgfpathlineto{\pgfqpoint{11.094706in}{1.955406in}}%
\pgfpathlineto{\pgfqpoint{11.106560in}{1.955406in}}%
\pgfpathlineto{\pgfqpoint{11.106560in}{1.912173in}}%
\pgfpathquadraticcurveto{\pgfqpoint{11.106560in}{1.901927in}}{\pgfqpoint{11.110539in}{1.896794in}}%
\pgfpathquadraticcurveto{\pgfqpoint{11.114539in}{1.891681in}}{\pgfqpoint{11.122538in}{1.891681in}}%
\pgfpathquadraticcurveto{\pgfqpoint{11.132124in}{1.891681in}}{\pgfqpoint{11.137691in}{1.897804in}}%
\pgfpathquadraticcurveto{\pgfqpoint{11.143278in}{1.903927in}}{\pgfqpoint{11.143278in}{1.914503in}}%
\pgfpathlineto{\pgfqpoint{11.143278in}{1.955406in}}%
\pgfpathlineto{\pgfqpoint{11.155132in}{1.955406in}}%
\pgfpathlineto{\pgfqpoint{11.155132in}{1.883249in}}%
\pgfpathlineto{\pgfqpoint{11.143278in}{1.883249in}}%
\pgfpathlineto{\pgfqpoint{11.143278in}{1.894340in}}%
\pgfpathquadraticcurveto{\pgfqpoint{11.138969in}{1.887764in}}{\pgfqpoint{11.133258in}{1.884568in}}%
\pgfpathquadraticcurveto{\pgfqpoint{11.127568in}{1.881373in}}{\pgfqpoint{11.120023in}{1.881373in}}%
\pgfpathquadraticcurveto{\pgfqpoint{11.107591in}{1.881373in}}{\pgfqpoint{11.101138in}{1.889104in}}%
\pgfpathquadraticcurveto{\pgfqpoint{11.094706in}{1.896835in}}{\pgfqpoint{11.094706in}{1.911720in}}%
\pgfpathlineto{\pgfqpoint{11.094706in}{1.911720in}}%
\pgfpathclose%
\pgfpathmoveto{\pgfqpoint{11.124538in}{1.957138in}}%
\pgfpathlineto{\pgfqpoint{11.124538in}{1.957138in}}%
\pgfpathlineto{\pgfqpoint{11.124538in}{1.957138in}}%
\pgfpathclose%
\pgfpathmoveto{\pgfqpoint{11.227021in}{1.944459in}}%
\pgfpathlineto{\pgfqpoint{11.227021in}{1.983506in}}%
\pgfpathlineto{\pgfqpoint{11.238876in}{1.983506in}}%
\pgfpathlineto{\pgfqpoint{11.238876in}{1.883249in}}%
\pgfpathlineto{\pgfqpoint{11.227021in}{1.883249in}}%
\pgfpathlineto{\pgfqpoint{11.227021in}{1.894072in}}%
\pgfpathquadraticcurveto{\pgfqpoint{11.223290in}{1.887640in}}{\pgfqpoint{11.217579in}{1.884506in}}%
\pgfpathquadraticcurveto{\pgfqpoint{11.211889in}{1.881373in}}{\pgfqpoint{11.203890in}{1.881373in}}%
\pgfpathquadraticcurveto{\pgfqpoint{11.190819in}{1.881373in}}{\pgfqpoint{11.182593in}{1.891805in}}%
\pgfpathquadraticcurveto{\pgfqpoint{11.174388in}{1.902257in}}{\pgfqpoint{11.174388in}{1.919266in}}%
\pgfpathquadraticcurveto{\pgfqpoint{11.174388in}{1.936274in}}{\pgfqpoint{11.182593in}{1.946706in}}%
\pgfpathquadraticcurveto{\pgfqpoint{11.190819in}{1.957138in}}{\pgfqpoint{11.203890in}{1.957138in}}%
\pgfpathquadraticcurveto{\pgfqpoint{11.211889in}{1.957138in}}{\pgfqpoint{11.217579in}{1.954004in}}%
\pgfpathquadraticcurveto{\pgfqpoint{11.223290in}{1.950891in}}{\pgfqpoint{11.227021in}{1.944459in}}%
\pgfpathlineto{\pgfqpoint{11.227021in}{1.944459in}}%
\pgfpathclose%
\pgfpathmoveto{\pgfqpoint{11.186634in}{1.919266in}}%
\pgfpathquadraticcurveto{\pgfqpoint{11.186634in}{1.906195in}}{\pgfqpoint{11.192015in}{1.898752in}}%
\pgfpathquadraticcurveto{\pgfqpoint{11.197396in}{1.891310in}}{\pgfqpoint{11.206797in}{1.891310in}}%
\pgfpathquadraticcurveto{\pgfqpoint{11.216198in}{1.891310in}}{\pgfqpoint{11.221599in}{1.898752in}}%
\pgfpathquadraticcurveto{\pgfqpoint{11.227021in}{1.906195in}}{\pgfqpoint{11.227021in}{1.919266in}}%
\pgfpathquadraticcurveto{\pgfqpoint{11.227021in}{1.932336in}}{\pgfqpoint{11.221599in}{1.939779in}}%
\pgfpathquadraticcurveto{\pgfqpoint{11.216198in}{1.947221in}}{\pgfqpoint{11.206797in}{1.947221in}}%
\pgfpathquadraticcurveto{\pgfqpoint{11.197396in}{1.947221in}}{\pgfqpoint{11.192015in}{1.939779in}}%
\pgfpathquadraticcurveto{\pgfqpoint{11.186634in}{1.932336in}}{\pgfqpoint{11.186634in}{1.919266in}}%
\pgfpathlineto{\pgfqpoint{11.186634in}{1.919266in}}%
\pgfpathclose%
\pgfpathmoveto{\pgfqpoint{11.293323in}{1.876548in}}%
\pgfpathquadraticcurveto{\pgfqpoint{11.288314in}{1.863663in}}{\pgfqpoint{11.283531in}{1.859746in}}%
\pgfpathquadraticcurveto{\pgfqpoint{11.278768in}{1.855808in}}{\pgfqpoint{11.270790in}{1.855808in}}%
\pgfpathlineto{\pgfqpoint{11.261306in}{1.855808in}}%
\pgfpathlineto{\pgfqpoint{11.261306in}{1.865725in}}%
\pgfpathlineto{\pgfqpoint{11.268275in}{1.865725in}}%
\pgfpathquadraticcurveto{\pgfqpoint{11.273161in}{1.865725in}}{\pgfqpoint{11.275861in}{1.868055in}}%
\pgfpathquadraticcurveto{\pgfqpoint{11.278583in}{1.870364in}}{\pgfqpoint{11.281861in}{1.879002in}}%
\pgfpathlineto{\pgfqpoint{11.283984in}{1.884403in}}%
\pgfpathlineto{\pgfqpoint{11.254812in}{1.955406in}}%
\pgfpathlineto{\pgfqpoint{11.267367in}{1.955406in}}%
\pgfpathlineto{\pgfqpoint{11.289922in}{1.898979in}}%
\pgfpathlineto{\pgfqpoint{11.312476in}{1.955406in}}%
\pgfpathlineto{\pgfqpoint{11.325031in}{1.955406in}}%
\pgfpathlineto{\pgfqpoint{11.293323in}{1.876548in}}%
\pgfpathlineto{\pgfqpoint{11.293323in}{1.876548in}}%
\pgfpathclose%
\pgfpathmoveto{\pgfqpoint{11.396219in}{1.894196in}}%
\pgfpathlineto{\pgfqpoint{11.441637in}{1.894196in}}%
\pgfpathlineto{\pgfqpoint{11.441637in}{1.883249in}}%
\pgfpathlineto{\pgfqpoint{11.380572in}{1.883249in}}%
\pgfpathlineto{\pgfqpoint{11.380572in}{1.894196in}}%
\pgfpathquadraticcurveto{\pgfqpoint{11.387973in}{1.901865in}}{\pgfqpoint{11.400755in}{1.914771in}}%
\pgfpathquadraticcurveto{\pgfqpoint{11.413558in}{1.927698in}}{\pgfqpoint{11.416836in}{1.931450in}}%
\pgfpathquadraticcurveto{\pgfqpoint{11.423082in}{1.938459in}}{\pgfqpoint{11.425556in}{1.943325in}}%
\pgfpathquadraticcurveto{\pgfqpoint{11.428051in}{1.948190in}}{\pgfqpoint{11.428051in}{1.952891in}}%
\pgfpathquadraticcurveto{\pgfqpoint{11.428051in}{1.960560in}}{\pgfqpoint{11.422670in}{1.965384in}}%
\pgfpathquadraticcurveto{\pgfqpoint{11.417289in}{1.970229in}}{\pgfqpoint{11.408651in}{1.970229in}}%
\pgfpathquadraticcurveto{\pgfqpoint{11.402528in}{1.970229in}}{\pgfqpoint{11.395725in}{1.968106in}}%
\pgfpathquadraticcurveto{\pgfqpoint{11.388942in}{1.965982in}}{\pgfqpoint{11.381211in}{1.961653in}}%
\pgfpathlineto{\pgfqpoint{11.381211in}{1.974806in}}%
\pgfpathquadraticcurveto{\pgfqpoint{11.389066in}{1.977960in}}{\pgfqpoint{11.395890in}{1.979568in}}%
\pgfpathquadraticcurveto{\pgfqpoint{11.402734in}{1.981176in}}{\pgfqpoint{11.408404in}{1.981176in}}%
\pgfpathquadraticcurveto{\pgfqpoint{11.423350in}{1.981176in}}{\pgfqpoint{11.432236in}{1.973693in}}%
\pgfpathquadraticcurveto{\pgfqpoint{11.441122in}{1.966229in}}{\pgfqpoint{11.441122in}{1.953736in}}%
\pgfpathquadraticcurveto{\pgfqpoint{11.441122in}{1.947798in}}{\pgfqpoint{11.438895in}{1.942479in}}%
\pgfpathquadraticcurveto{\pgfqpoint{11.436689in}{1.937181in}}{\pgfqpoint{11.430814in}{1.929965in}}%
\pgfpathquadraticcurveto{\pgfqpoint{11.429206in}{1.928089in}}{\pgfqpoint{11.420567in}{1.919162in}}%
\pgfpathquadraticcurveto{\pgfqpoint{11.411950in}{1.910236in}}{\pgfqpoint{11.396219in}{1.894196in}}%
\pgfpathlineto{\pgfqpoint{11.396219in}{1.894196in}}%
\pgfpathclose%
\pgfpathmoveto{\pgfqpoint{11.471222in}{1.894196in}}%
\pgfpathlineto{\pgfqpoint{11.492477in}{1.894196in}}%
\pgfpathlineto{\pgfqpoint{11.492477in}{1.967590in}}%
\pgfpathlineto{\pgfqpoint{11.469345in}{1.962951in}}%
\pgfpathlineto{\pgfqpoint{11.469345in}{1.974806in}}%
\pgfpathlineto{\pgfqpoint{11.492353in}{1.979445in}}%
\pgfpathlineto{\pgfqpoint{11.505362in}{1.979445in}}%
\pgfpathlineto{\pgfqpoint{11.505362in}{1.894196in}}%
\pgfpathlineto{\pgfqpoint{11.526618in}{1.894196in}}%
\pgfpathlineto{\pgfqpoint{11.526618in}{1.883249in}}%
\pgfpathlineto{\pgfqpoint{11.471222in}{1.883249in}}%
\pgfpathlineto{\pgfqpoint{11.471222in}{1.894196in}}%
\pgfpathlineto{\pgfqpoint{11.471222in}{1.894196in}}%
\pgfpathclose%
\pgfpathmoveto{\pgfqpoint{11.621638in}{1.983362in}}%
\pgfpathquadraticcurveto{\pgfqpoint{11.613021in}{1.968559in}}{\pgfqpoint{11.608815in}{1.954045in}}%
\pgfpathquadraticcurveto{\pgfqpoint{11.604630in}{1.939552in}}{\pgfqpoint{11.604630in}{1.924667in}}%
\pgfpathquadraticcurveto{\pgfqpoint{11.604630in}{1.909803in}}{\pgfqpoint{11.608856in}{1.895206in}}%
\pgfpathquadraticcurveto{\pgfqpoint{11.613082in}{1.880610in}}{\pgfqpoint{11.621638in}{1.865849in}}%
\pgfpathlineto{\pgfqpoint{11.611330in}{1.865849in}}%
\pgfpathquadraticcurveto{\pgfqpoint{11.601682in}{1.881002in}}{\pgfqpoint{11.596878in}{1.895619in}}%
\pgfpathquadraticcurveto{\pgfqpoint{11.592074in}{1.910236in}}{\pgfqpoint{11.592074in}{1.924667in}}%
\pgfpathquadraticcurveto{\pgfqpoint{11.592074in}{1.939037in}}{\pgfqpoint{11.596837in}{1.953592in}}%
\pgfpathquadraticcurveto{\pgfqpoint{11.601620in}{1.968167in}}{\pgfqpoint{11.611330in}{1.983362in}}%
\pgfpathlineto{\pgfqpoint{11.621638in}{1.983362in}}%
\pgfpathlineto{\pgfqpoint{11.621638in}{1.983362in}}%
\pgfpathclose%
\pgfpathmoveto{\pgfqpoint{11.657531in}{1.894196in}}%
\pgfpathlineto{\pgfqpoint{11.702949in}{1.894196in}}%
\pgfpathlineto{\pgfqpoint{11.702949in}{1.883249in}}%
\pgfpathlineto{\pgfqpoint{11.641883in}{1.883249in}}%
\pgfpathlineto{\pgfqpoint{11.641883in}{1.894196in}}%
\pgfpathquadraticcurveto{\pgfqpoint{11.649285in}{1.901865in}}{\pgfqpoint{11.662067in}{1.914771in}}%
\pgfpathquadraticcurveto{\pgfqpoint{11.674870in}{1.927698in}}{\pgfqpoint{11.678148in}{1.931450in}}%
\pgfpathquadraticcurveto{\pgfqpoint{11.684394in}{1.938459in}}{\pgfqpoint{11.686868in}{1.943325in}}%
\pgfpathquadraticcurveto{\pgfqpoint{11.689363in}{1.948190in}}{\pgfqpoint{11.689363in}{1.952891in}}%
\pgfpathquadraticcurveto{\pgfqpoint{11.689363in}{1.960560in}}{\pgfqpoint{11.683982in}{1.965384in}}%
\pgfpathquadraticcurveto{\pgfqpoint{11.678601in}{1.970229in}}{\pgfqpoint{11.669963in}{1.970229in}}%
\pgfpathquadraticcurveto{\pgfqpoint{11.663840in}{1.970229in}}{\pgfqpoint{11.657036in}{1.968106in}}%
\pgfpathquadraticcurveto{\pgfqpoint{11.650254in}{1.965982in}}{\pgfqpoint{11.642523in}{1.961653in}}%
\pgfpathlineto{\pgfqpoint{11.642523in}{1.974806in}}%
\pgfpathquadraticcurveto{\pgfqpoint{11.650377in}{1.977960in}}{\pgfqpoint{11.657201in}{1.979568in}}%
\pgfpathquadraticcurveto{\pgfqpoint{11.664046in}{1.981176in}}{\pgfqpoint{11.669715in}{1.981176in}}%
\pgfpathquadraticcurveto{\pgfqpoint{11.684662in}{1.981176in}}{\pgfqpoint{11.693548in}{1.973693in}}%
\pgfpathquadraticcurveto{\pgfqpoint{11.702434in}{1.966229in}}{\pgfqpoint{11.702434in}{1.953736in}}%
\pgfpathquadraticcurveto{\pgfqpoint{11.702434in}{1.947798in}}{\pgfqpoint{11.700207in}{1.942479in}}%
\pgfpathquadraticcurveto{\pgfqpoint{11.698001in}{1.937181in}}{\pgfqpoint{11.692125in}{1.929965in}}%
\pgfpathquadraticcurveto{\pgfqpoint{11.690517in}{1.928089in}}{\pgfqpoint{11.681879in}{1.919162in}}%
\pgfpathquadraticcurveto{\pgfqpoint{11.673262in}{1.910236in}}{\pgfqpoint{11.657531in}{1.894196in}}%
\pgfpathlineto{\pgfqpoint{11.657531in}{1.894196in}}%
\pgfpathclose%
\pgfpathmoveto{\pgfqpoint{11.758098in}{1.970868in}}%
\pgfpathquadraticcurveto{\pgfqpoint{11.748058in}{1.970868in}}{\pgfqpoint{11.742986in}{1.960972in}}%
\pgfpathquadraticcurveto{\pgfqpoint{11.737935in}{1.951097in}}{\pgfqpoint{11.737935in}{1.931244in}}%
\pgfpathquadraticcurveto{\pgfqpoint{11.737935in}{1.911473in}}{\pgfqpoint{11.742986in}{1.901577in}}%
\pgfpathquadraticcurveto{\pgfqpoint{11.748058in}{1.891681in}}{\pgfqpoint{11.758098in}{1.891681in}}%
\pgfpathquadraticcurveto{\pgfqpoint{11.768220in}{1.891681in}}{\pgfqpoint{11.773271in}{1.901577in}}%
\pgfpathquadraticcurveto{\pgfqpoint{11.778343in}{1.911473in}}{\pgfqpoint{11.778343in}{1.931244in}}%
\pgfpathquadraticcurveto{\pgfqpoint{11.778343in}{1.951097in}}{\pgfqpoint{11.773271in}{1.960972in}}%
\pgfpathquadraticcurveto{\pgfqpoint{11.768220in}{1.970868in}}{\pgfqpoint{11.758098in}{1.970868in}}%
\pgfpathlineto{\pgfqpoint{11.758098in}{1.970868in}}%
\pgfpathclose%
\pgfpathmoveto{\pgfqpoint{11.758098in}{1.981176in}}%
\pgfpathquadraticcurveto{\pgfqpoint{11.774281in}{1.981176in}}{\pgfqpoint{11.782817in}{1.968374in}}%
\pgfpathquadraticcurveto{\pgfqpoint{11.791352in}{1.955591in}}{\pgfqpoint{11.791352in}{1.931244in}}%
\pgfpathquadraticcurveto{\pgfqpoint{11.791352in}{1.906958in}}{\pgfqpoint{11.782817in}{1.894155in}}%
\pgfpathquadraticcurveto{\pgfqpoint{11.774281in}{1.881373in}}{\pgfqpoint{11.758098in}{1.881373in}}%
\pgfpathquadraticcurveto{\pgfqpoint{11.741934in}{1.881373in}}{\pgfqpoint{11.733399in}{1.894155in}}%
\pgfpathquadraticcurveto{\pgfqpoint{11.724864in}{1.906958in}}{\pgfqpoint{11.724864in}{1.931244in}}%
\pgfpathquadraticcurveto{\pgfqpoint{11.724864in}{1.955591in}}{\pgfqpoint{11.733399in}{1.968374in}}%
\pgfpathquadraticcurveto{\pgfqpoint{11.741934in}{1.981176in}}{\pgfqpoint{11.758098in}{1.981176in}}%
\pgfpathlineto{\pgfqpoint{11.758098in}{1.981176in}}%
\pgfpathclose%
\pgfpathmoveto{\pgfqpoint{11.825431in}{1.894196in}}%
\pgfpathlineto{\pgfqpoint{11.870848in}{1.894196in}}%
\pgfpathlineto{\pgfqpoint{11.870848in}{1.883249in}}%
\pgfpathlineto{\pgfqpoint{11.809783in}{1.883249in}}%
\pgfpathlineto{\pgfqpoint{11.809783in}{1.894196in}}%
\pgfpathquadraticcurveto{\pgfqpoint{11.817184in}{1.901865in}}{\pgfqpoint{11.829966in}{1.914771in}}%
\pgfpathquadraticcurveto{\pgfqpoint{11.842769in}{1.927698in}}{\pgfqpoint{11.846047in}{1.931450in}}%
\pgfpathquadraticcurveto{\pgfqpoint{11.852294in}{1.938459in}}{\pgfqpoint{11.854768in}{1.943325in}}%
\pgfpathquadraticcurveto{\pgfqpoint{11.857262in}{1.948190in}}{\pgfqpoint{11.857262in}{1.952891in}}%
\pgfpathquadraticcurveto{\pgfqpoint{11.857262in}{1.960560in}}{\pgfqpoint{11.851881in}{1.965384in}}%
\pgfpathquadraticcurveto{\pgfqpoint{11.846500in}{1.970229in}}{\pgfqpoint{11.837862in}{1.970229in}}%
\pgfpathquadraticcurveto{\pgfqpoint{11.831739in}{1.970229in}}{\pgfqpoint{11.824936in}{1.968106in}}%
\pgfpathquadraticcurveto{\pgfqpoint{11.818153in}{1.965982in}}{\pgfqpoint{11.810422in}{1.961653in}}%
\pgfpathlineto{\pgfqpoint{11.810422in}{1.974806in}}%
\pgfpathquadraticcurveto{\pgfqpoint{11.818277in}{1.977960in}}{\pgfqpoint{11.825101in}{1.979568in}}%
\pgfpathquadraticcurveto{\pgfqpoint{11.831945in}{1.981176in}}{\pgfqpoint{11.837615in}{1.981176in}}%
\pgfpathquadraticcurveto{\pgfqpoint{11.852562in}{1.981176in}}{\pgfqpoint{11.861447in}{1.973693in}}%
\pgfpathquadraticcurveto{\pgfqpoint{11.870333in}{1.966229in}}{\pgfqpoint{11.870333in}{1.953736in}}%
\pgfpathquadraticcurveto{\pgfqpoint{11.870333in}{1.947798in}}{\pgfqpoint{11.868106in}{1.942479in}}%
\pgfpathquadraticcurveto{\pgfqpoint{11.865900in}{1.937181in}}{\pgfqpoint{11.860025in}{1.929965in}}%
\pgfpathquadraticcurveto{\pgfqpoint{11.858417in}{1.928089in}}{\pgfqpoint{11.849778in}{1.919162in}}%
\pgfpathquadraticcurveto{\pgfqpoint{11.841161in}{1.910236in}}{\pgfqpoint{11.825431in}{1.894196in}}%
\pgfpathlineto{\pgfqpoint{11.825431in}{1.894196in}}%
\pgfpathclose%
\pgfpathmoveto{\pgfqpoint{11.898309in}{1.979445in}}%
\pgfpathlineto{\pgfqpoint{11.949396in}{1.979445in}}%
\pgfpathlineto{\pgfqpoint{11.949396in}{1.968477in}}%
\pgfpathlineto{\pgfqpoint{11.910225in}{1.968477in}}%
\pgfpathlineto{\pgfqpoint{11.910225in}{1.944912in}}%
\pgfpathquadraticcurveto{\pgfqpoint{11.913050in}{1.945881in}}{\pgfqpoint{11.915874in}{1.946355in}}%
\pgfpathquadraticcurveto{\pgfqpoint{11.918719in}{1.946830in}}{\pgfqpoint{11.921564in}{1.946830in}}%
\pgfpathquadraticcurveto{\pgfqpoint{11.937666in}{1.946830in}}{\pgfqpoint{11.947067in}{1.938006in}}%
\pgfpathquadraticcurveto{\pgfqpoint{11.956488in}{1.929182in}}{\pgfqpoint{11.956488in}{1.914111in}}%
\pgfpathquadraticcurveto{\pgfqpoint{11.956488in}{1.898587in}}{\pgfqpoint{11.946819in}{1.889970in}}%
\pgfpathquadraticcurveto{\pgfqpoint{11.937150in}{1.881373in}}{\pgfqpoint{11.919565in}{1.881373in}}%
\pgfpathquadraticcurveto{\pgfqpoint{11.913503in}{1.881373in}}{\pgfqpoint{11.907215in}{1.882404in}}%
\pgfpathquadraticcurveto{\pgfqpoint{11.900948in}{1.883434in}}{\pgfqpoint{11.894248in}{1.885496in}}%
\pgfpathlineto{\pgfqpoint{11.894248in}{1.898587in}}%
\pgfpathquadraticcurveto{\pgfqpoint{11.900041in}{1.895433in}}{\pgfqpoint{11.906226in}{1.893887in}}%
\pgfpathquadraticcurveto{\pgfqpoint{11.912411in}{1.892341in}}{\pgfqpoint{11.919297in}{1.892341in}}%
\pgfpathquadraticcurveto{\pgfqpoint{11.930450in}{1.892341in}}{\pgfqpoint{11.936944in}{1.898196in}}%
\pgfpathquadraticcurveto{\pgfqpoint{11.943459in}{1.904051in}}{\pgfqpoint{11.943459in}{1.914111in}}%
\pgfpathquadraticcurveto{\pgfqpoint{11.943459in}{1.924152in}}{\pgfqpoint{11.936944in}{1.930007in}}%
\pgfpathquadraticcurveto{\pgfqpoint{11.930450in}{1.935882in}}{\pgfqpoint{11.919297in}{1.935882in}}%
\pgfpathquadraticcurveto{\pgfqpoint{11.914081in}{1.935882in}}{\pgfqpoint{11.908885in}{1.934728in}}%
\pgfpathquadraticcurveto{\pgfqpoint{11.903711in}{1.933573in}}{\pgfqpoint{11.898309in}{1.931120in}}%
\pgfpathlineto{\pgfqpoint{11.898309in}{1.979445in}}%
\pgfpathlineto{\pgfqpoint{11.898309in}{1.979445in}}%
\pgfpathclose%
\pgfpathmoveto{\pgfqpoint{11.978589in}{1.983362in}}%
\pgfpathlineto{\pgfqpoint{11.988897in}{1.983362in}}%
\pgfpathquadraticcurveto{\pgfqpoint{11.998546in}{1.968167in}}{\pgfqpoint{12.003349in}{1.953592in}}%
\pgfpathquadraticcurveto{\pgfqpoint{12.008153in}{1.939037in}}{\pgfqpoint{12.008153in}{1.924667in}}%
\pgfpathquadraticcurveto{\pgfqpoint{12.008153in}{1.910236in}}{\pgfqpoint{12.003349in}{1.895619in}}%
\pgfpathquadraticcurveto{\pgfqpoint{11.998546in}{1.881002in}}{\pgfqpoint{11.988897in}{1.865849in}}%
\pgfpathlineto{\pgfqpoint{11.978589in}{1.865849in}}%
\pgfpathquadraticcurveto{\pgfqpoint{11.987145in}{1.880610in}}{\pgfqpoint{11.991371in}{1.895206in}}%
\pgfpathquadraticcurveto{\pgfqpoint{11.995598in}{1.909803in}}{\pgfqpoint{11.995598in}{1.924667in}}%
\pgfpathquadraticcurveto{\pgfqpoint{11.995598in}{1.939552in}}{\pgfqpoint{11.991371in}{1.954045in}}%
\pgfpathquadraticcurveto{\pgfqpoint{11.987145in}{1.968559in}}{\pgfqpoint{11.978589in}{1.983362in}}%
\pgfpathlineto{\pgfqpoint{11.978589in}{1.983362in}}%
\pgfpathclose%
\pgfusepath{fill}%
\end{pgfscope}%
\begin{pgfscope}%
\pgfsetbuttcap%
\pgfsetroundjoin%
\definecolor{currentfill}{rgb}{0.090196,0.168627,0.227451}%
\pgfsetfillcolor{currentfill}%
\pgfsetlinewidth{0.000000pt}%
\definecolor{currentstroke}{rgb}{0.090196,0.168627,0.227451}%
\pgfsetstrokecolor{currentstroke}%
\pgfsetdash{}{0pt}%
\pgfpathmoveto{\pgfqpoint{11.018611in}{1.576690in}}%
\pgfpathlineto{\pgfqpoint{11.018611in}{1.563991in}}%
\pgfpathquadraticcurveto{\pgfqpoint{11.011210in}{1.567537in}}{\pgfqpoint{11.004633in}{1.569268in}}%
\pgfpathquadraticcurveto{\pgfqpoint{10.998056in}{1.571021in}}{\pgfqpoint{10.991933in}{1.571021in}}%
\pgfpathquadraticcurveto{\pgfqpoint{10.981316in}{1.571021in}}{\pgfqpoint{10.975543in}{1.566897in}}%
\pgfpathquadraticcurveto{\pgfqpoint{10.969771in}{1.562774in}}{\pgfqpoint{10.969771in}{1.555167in}}%
\pgfpathquadraticcurveto{\pgfqpoint{10.969771in}{1.548776in}}{\pgfqpoint{10.973605in}{1.545518in}}%
\pgfpathquadraticcurveto{\pgfqpoint{10.977440in}{1.542282in}}{\pgfqpoint{10.988140in}{1.540282in}}%
\pgfpathlineto{\pgfqpoint{10.995995in}{1.538674in}}%
\pgfpathquadraticcurveto{\pgfqpoint{11.010550in}{1.535891in}}{\pgfqpoint{11.017477in}{1.528902in}}%
\pgfpathquadraticcurveto{\pgfqpoint{11.024404in}{1.521913in}}{\pgfqpoint{11.024404in}{1.510203in}}%
\pgfpathquadraticcurveto{\pgfqpoint{11.024404in}{1.496204in}}{\pgfqpoint{11.015024in}{1.488988in}}%
\pgfpathquadraticcurveto{\pgfqpoint{11.005664in}{1.481773in}}{\pgfqpoint{10.987563in}{1.481773in}}%
\pgfpathquadraticcurveto{\pgfqpoint{10.980739in}{1.481773in}}{\pgfqpoint{10.973028in}{1.483319in}}%
\pgfpathquadraticcurveto{\pgfqpoint{10.965338in}{1.484865in}}{\pgfqpoint{10.957092in}{1.487896in}}%
\pgfpathlineto{\pgfqpoint{10.957092in}{1.501296in}}%
\pgfpathquadraticcurveto{\pgfqpoint{10.965008in}{1.496864in}}{\pgfqpoint{10.972616in}{1.494596in}}%
\pgfpathquadraticcurveto{\pgfqpoint{10.980223in}{1.492349in}}{\pgfqpoint{10.987563in}{1.492349in}}%
\pgfpathquadraticcurveto{\pgfqpoint{10.998696in}{1.492349in}}{\pgfqpoint{11.004757in}{1.496720in}}%
\pgfpathquadraticcurveto{\pgfqpoint{11.010818in}{1.501111in}}{\pgfqpoint{11.010818in}{1.509234in}}%
\pgfpathquadraticcurveto{\pgfqpoint{11.010818in}{1.516305in}}{\pgfqpoint{11.006468in}{1.520305in}}%
\pgfpathquadraticcurveto{\pgfqpoint{11.002118in}{1.524304in}}{\pgfqpoint{10.992201in}{1.526304in}}%
\pgfpathlineto{\pgfqpoint{10.984264in}{1.527850in}}%
\pgfpathquadraticcurveto{\pgfqpoint{10.969709in}{1.530736in}}{\pgfqpoint{10.963194in}{1.536921in}}%
\pgfpathquadraticcurveto{\pgfqpoint{10.956700in}{1.543106in}}{\pgfqpoint{10.956700in}{1.554136in}}%
\pgfpathquadraticcurveto{\pgfqpoint{10.956700in}{1.566897in}}{\pgfqpoint{10.965689in}{1.574237in}}%
\pgfpathquadraticcurveto{\pgfqpoint{10.974678in}{1.581576in}}{\pgfqpoint{10.990449in}{1.581576in}}%
\pgfpathquadraticcurveto{\pgfqpoint{10.997232in}{1.581576in}}{\pgfqpoint{11.004241in}{1.580339in}}%
\pgfpathquadraticcurveto{\pgfqpoint{11.011271in}{1.579123in}}{\pgfqpoint{11.018611in}{1.576690in}}%
\pgfpathlineto{\pgfqpoint{11.018611in}{1.576690in}}%
\pgfpathclose%
\pgfpathmoveto{\pgfqpoint{11.055926in}{1.576299in}}%
\pgfpathlineto{\pgfqpoint{11.055926in}{1.555806in}}%
\pgfpathlineto{\pgfqpoint{11.080336in}{1.555806in}}%
\pgfpathlineto{\pgfqpoint{11.080336in}{1.546590in}}%
\pgfpathlineto{\pgfqpoint{11.055926in}{1.546590in}}%
\pgfpathlineto{\pgfqpoint{11.055926in}{1.507419in}}%
\pgfpathquadraticcurveto{\pgfqpoint{11.055926in}{1.498596in}}{\pgfqpoint{11.058339in}{1.496080in}}%
\pgfpathquadraticcurveto{\pgfqpoint{11.060751in}{1.493565in}}{\pgfqpoint{11.068173in}{1.493565in}}%
\pgfpathlineto{\pgfqpoint{11.080336in}{1.493565in}}%
\pgfpathlineto{\pgfqpoint{11.080336in}{1.483649in}}%
\pgfpathlineto{\pgfqpoint{11.068173in}{1.483649in}}%
\pgfpathquadraticcurveto{\pgfqpoint{11.054442in}{1.483649in}}{\pgfqpoint{11.049226in}{1.488762in}}%
\pgfpathquadraticcurveto{\pgfqpoint{11.044010in}{1.493895in}}{\pgfqpoint{11.044010in}{1.507419in}}%
\pgfpathlineto{\pgfqpoint{11.044010in}{1.546590in}}%
\pgfpathlineto{\pgfqpoint{11.035310in}{1.546590in}}%
\pgfpathlineto{\pgfqpoint{11.035310in}{1.555806in}}%
\pgfpathlineto{\pgfqpoint{11.044010in}{1.555806in}}%
\pgfpathlineto{\pgfqpoint{11.044010in}{1.576299in}}%
\pgfpathlineto{\pgfqpoint{11.055926in}{1.576299in}}%
\pgfpathlineto{\pgfqpoint{11.055926in}{1.576299in}}%
\pgfpathclose%
\pgfpathmoveto{\pgfqpoint{11.094706in}{1.512120in}}%
\pgfpathlineto{\pgfqpoint{11.094706in}{1.555806in}}%
\pgfpathlineto{\pgfqpoint{11.106560in}{1.555806in}}%
\pgfpathlineto{\pgfqpoint{11.106560in}{1.512573in}}%
\pgfpathquadraticcurveto{\pgfqpoint{11.106560in}{1.502327in}}{\pgfqpoint{11.110539in}{1.497194in}}%
\pgfpathquadraticcurveto{\pgfqpoint{11.114539in}{1.492081in}}{\pgfqpoint{11.122538in}{1.492081in}}%
\pgfpathquadraticcurveto{\pgfqpoint{11.132124in}{1.492081in}}{\pgfqpoint{11.137691in}{1.498204in}}%
\pgfpathquadraticcurveto{\pgfqpoint{11.143278in}{1.504327in}}{\pgfqpoint{11.143278in}{1.514903in}}%
\pgfpathlineto{\pgfqpoint{11.143278in}{1.555806in}}%
\pgfpathlineto{\pgfqpoint{11.155132in}{1.555806in}}%
\pgfpathlineto{\pgfqpoint{11.155132in}{1.483649in}}%
\pgfpathlineto{\pgfqpoint{11.143278in}{1.483649in}}%
\pgfpathlineto{\pgfqpoint{11.143278in}{1.494740in}}%
\pgfpathquadraticcurveto{\pgfqpoint{11.138969in}{1.488164in}}{\pgfqpoint{11.133258in}{1.484968in}}%
\pgfpathquadraticcurveto{\pgfqpoint{11.127568in}{1.481773in}}{\pgfqpoint{11.120023in}{1.481773in}}%
\pgfpathquadraticcurveto{\pgfqpoint{11.107591in}{1.481773in}}{\pgfqpoint{11.101138in}{1.489504in}}%
\pgfpathquadraticcurveto{\pgfqpoint{11.094706in}{1.497235in}}{\pgfqpoint{11.094706in}{1.512120in}}%
\pgfpathlineto{\pgfqpoint{11.094706in}{1.512120in}}%
\pgfpathclose%
\pgfpathmoveto{\pgfqpoint{11.124538in}{1.557538in}}%
\pgfpathlineto{\pgfqpoint{11.124538in}{1.557538in}}%
\pgfpathlineto{\pgfqpoint{11.124538in}{1.557538in}}%
\pgfpathclose%
\pgfpathmoveto{\pgfqpoint{11.227021in}{1.544859in}}%
\pgfpathlineto{\pgfqpoint{11.227021in}{1.583906in}}%
\pgfpathlineto{\pgfqpoint{11.238876in}{1.583906in}}%
\pgfpathlineto{\pgfqpoint{11.238876in}{1.483649in}}%
\pgfpathlineto{\pgfqpoint{11.227021in}{1.483649in}}%
\pgfpathlineto{\pgfqpoint{11.227021in}{1.494472in}}%
\pgfpathquadraticcurveto{\pgfqpoint{11.223290in}{1.488040in}}{\pgfqpoint{11.217579in}{1.484906in}}%
\pgfpathquadraticcurveto{\pgfqpoint{11.211889in}{1.481773in}}{\pgfqpoint{11.203890in}{1.481773in}}%
\pgfpathquadraticcurveto{\pgfqpoint{11.190819in}{1.481773in}}{\pgfqpoint{11.182593in}{1.492205in}}%
\pgfpathquadraticcurveto{\pgfqpoint{11.174388in}{1.502657in}}{\pgfqpoint{11.174388in}{1.519666in}}%
\pgfpathquadraticcurveto{\pgfqpoint{11.174388in}{1.536674in}}{\pgfqpoint{11.182593in}{1.547106in}}%
\pgfpathquadraticcurveto{\pgfqpoint{11.190819in}{1.557538in}}{\pgfqpoint{11.203890in}{1.557538in}}%
\pgfpathquadraticcurveto{\pgfqpoint{11.211889in}{1.557538in}}{\pgfqpoint{11.217579in}{1.554404in}}%
\pgfpathquadraticcurveto{\pgfqpoint{11.223290in}{1.551291in}}{\pgfqpoint{11.227021in}{1.544859in}}%
\pgfpathlineto{\pgfqpoint{11.227021in}{1.544859in}}%
\pgfpathclose%
\pgfpathmoveto{\pgfqpoint{11.186634in}{1.519666in}}%
\pgfpathquadraticcurveto{\pgfqpoint{11.186634in}{1.506595in}}{\pgfqpoint{11.192015in}{1.499152in}}%
\pgfpathquadraticcurveto{\pgfqpoint{11.197396in}{1.491710in}}{\pgfqpoint{11.206797in}{1.491710in}}%
\pgfpathquadraticcurveto{\pgfqpoint{11.216198in}{1.491710in}}{\pgfqpoint{11.221599in}{1.499152in}}%
\pgfpathquadraticcurveto{\pgfqpoint{11.227021in}{1.506595in}}{\pgfqpoint{11.227021in}{1.519666in}}%
\pgfpathquadraticcurveto{\pgfqpoint{11.227021in}{1.532736in}}{\pgfqpoint{11.221599in}{1.540179in}}%
\pgfpathquadraticcurveto{\pgfqpoint{11.216198in}{1.547621in}}{\pgfqpoint{11.206797in}{1.547621in}}%
\pgfpathquadraticcurveto{\pgfqpoint{11.197396in}{1.547621in}}{\pgfqpoint{11.192015in}{1.540179in}}%
\pgfpathquadraticcurveto{\pgfqpoint{11.186634in}{1.532736in}}{\pgfqpoint{11.186634in}{1.519666in}}%
\pgfpathlineto{\pgfqpoint{11.186634in}{1.519666in}}%
\pgfpathclose%
\pgfpathmoveto{\pgfqpoint{11.293323in}{1.476948in}}%
\pgfpathquadraticcurveto{\pgfqpoint{11.288314in}{1.464063in}}{\pgfqpoint{11.283531in}{1.460146in}}%
\pgfpathquadraticcurveto{\pgfqpoint{11.278768in}{1.456208in}}{\pgfqpoint{11.270790in}{1.456208in}}%
\pgfpathlineto{\pgfqpoint{11.261306in}{1.456208in}}%
\pgfpathlineto{\pgfqpoint{11.261306in}{1.466125in}}%
\pgfpathlineto{\pgfqpoint{11.268275in}{1.466125in}}%
\pgfpathquadraticcurveto{\pgfqpoint{11.273161in}{1.466125in}}{\pgfqpoint{11.275861in}{1.468455in}}%
\pgfpathquadraticcurveto{\pgfqpoint{11.278583in}{1.470764in}}{\pgfqpoint{11.281861in}{1.479402in}}%
\pgfpathlineto{\pgfqpoint{11.283984in}{1.484803in}}%
\pgfpathlineto{\pgfqpoint{11.254812in}{1.555806in}}%
\pgfpathlineto{\pgfqpoint{11.267367in}{1.555806in}}%
\pgfpathlineto{\pgfqpoint{11.289922in}{1.499379in}}%
\pgfpathlineto{\pgfqpoint{11.312476in}{1.555806in}}%
\pgfpathlineto{\pgfqpoint{11.325031in}{1.555806in}}%
\pgfpathlineto{\pgfqpoint{11.293323in}{1.476948in}}%
\pgfpathlineto{\pgfqpoint{11.293323in}{1.476948in}}%
\pgfpathclose%
\pgfpathmoveto{\pgfqpoint{11.396219in}{1.494596in}}%
\pgfpathlineto{\pgfqpoint{11.441637in}{1.494596in}}%
\pgfpathlineto{\pgfqpoint{11.441637in}{1.483649in}}%
\pgfpathlineto{\pgfqpoint{11.380572in}{1.483649in}}%
\pgfpathlineto{\pgfqpoint{11.380572in}{1.494596in}}%
\pgfpathquadraticcurveto{\pgfqpoint{11.387973in}{1.502265in}}{\pgfqpoint{11.400755in}{1.515171in}}%
\pgfpathquadraticcurveto{\pgfqpoint{11.413558in}{1.528098in}}{\pgfqpoint{11.416836in}{1.531850in}}%
\pgfpathquadraticcurveto{\pgfqpoint{11.423082in}{1.538859in}}{\pgfqpoint{11.425556in}{1.543725in}}%
\pgfpathquadraticcurveto{\pgfqpoint{11.428051in}{1.548590in}}{\pgfqpoint{11.428051in}{1.553291in}}%
\pgfpathquadraticcurveto{\pgfqpoint{11.428051in}{1.560960in}}{\pgfqpoint{11.422670in}{1.565784in}}%
\pgfpathquadraticcurveto{\pgfqpoint{11.417289in}{1.570629in}}{\pgfqpoint{11.408651in}{1.570629in}}%
\pgfpathquadraticcurveto{\pgfqpoint{11.402528in}{1.570629in}}{\pgfqpoint{11.395725in}{1.568506in}}%
\pgfpathquadraticcurveto{\pgfqpoint{11.388942in}{1.566382in}}{\pgfqpoint{11.381211in}{1.562053in}}%
\pgfpathlineto{\pgfqpoint{11.381211in}{1.575206in}}%
\pgfpathquadraticcurveto{\pgfqpoint{11.389066in}{1.578360in}}{\pgfqpoint{11.395890in}{1.579968in}}%
\pgfpathquadraticcurveto{\pgfqpoint{11.402734in}{1.581576in}}{\pgfqpoint{11.408404in}{1.581576in}}%
\pgfpathquadraticcurveto{\pgfqpoint{11.423350in}{1.581576in}}{\pgfqpoint{11.432236in}{1.574093in}}%
\pgfpathquadraticcurveto{\pgfqpoint{11.441122in}{1.566629in}}{\pgfqpoint{11.441122in}{1.554136in}}%
\pgfpathquadraticcurveto{\pgfqpoint{11.441122in}{1.548198in}}{\pgfqpoint{11.438895in}{1.542879in}}%
\pgfpathquadraticcurveto{\pgfqpoint{11.436689in}{1.537581in}}{\pgfqpoint{11.430814in}{1.530365in}}%
\pgfpathquadraticcurveto{\pgfqpoint{11.429206in}{1.528489in}}{\pgfqpoint{11.420567in}{1.519562in}}%
\pgfpathquadraticcurveto{\pgfqpoint{11.411950in}{1.510636in}}{\pgfqpoint{11.396219in}{1.494596in}}%
\pgfpathlineto{\pgfqpoint{11.396219in}{1.494596in}}%
\pgfpathclose%
\pgfpathmoveto{\pgfqpoint{11.480169in}{1.494596in}}%
\pgfpathlineto{\pgfqpoint{11.525587in}{1.494596in}}%
\pgfpathlineto{\pgfqpoint{11.525587in}{1.483649in}}%
\pgfpathlineto{\pgfqpoint{11.464521in}{1.483649in}}%
\pgfpathlineto{\pgfqpoint{11.464521in}{1.494596in}}%
\pgfpathquadraticcurveto{\pgfqpoint{11.471923in}{1.502265in}}{\pgfqpoint{11.484705in}{1.515171in}}%
\pgfpathquadraticcurveto{\pgfqpoint{11.497507in}{1.528098in}}{\pgfqpoint{11.500785in}{1.531850in}}%
\pgfpathquadraticcurveto{\pgfqpoint{11.507032in}{1.538859in}}{\pgfqpoint{11.509506in}{1.543725in}}%
\pgfpathquadraticcurveto{\pgfqpoint{11.512001in}{1.548590in}}{\pgfqpoint{11.512001in}{1.553291in}}%
\pgfpathquadraticcurveto{\pgfqpoint{11.512001in}{1.560960in}}{\pgfqpoint{11.506620in}{1.565784in}}%
\pgfpathquadraticcurveto{\pgfqpoint{11.501239in}{1.570629in}}{\pgfqpoint{11.492601in}{1.570629in}}%
\pgfpathquadraticcurveto{\pgfqpoint{11.486478in}{1.570629in}}{\pgfqpoint{11.479674in}{1.568506in}}%
\pgfpathquadraticcurveto{\pgfqpoint{11.472891in}{1.566382in}}{\pgfqpoint{11.465160in}{1.562053in}}%
\pgfpathlineto{\pgfqpoint{11.465160in}{1.575206in}}%
\pgfpathquadraticcurveto{\pgfqpoint{11.473015in}{1.578360in}}{\pgfqpoint{11.479839in}{1.579968in}}%
\pgfpathquadraticcurveto{\pgfqpoint{11.486684in}{1.581576in}}{\pgfqpoint{11.492353in}{1.581576in}}%
\pgfpathquadraticcurveto{\pgfqpoint{11.507300in}{1.581576in}}{\pgfqpoint{11.516186in}{1.574093in}}%
\pgfpathquadraticcurveto{\pgfqpoint{11.525071in}{1.566629in}}{\pgfqpoint{11.525071in}{1.554136in}}%
\pgfpathquadraticcurveto{\pgfqpoint{11.525071in}{1.548198in}}{\pgfqpoint{11.522845in}{1.542879in}}%
\pgfpathquadraticcurveto{\pgfqpoint{11.520639in}{1.537581in}}{\pgfqpoint{11.514763in}{1.530365in}}%
\pgfpathquadraticcurveto{\pgfqpoint{11.513155in}{1.528489in}}{\pgfqpoint{11.504517in}{1.519562in}}%
\pgfpathquadraticcurveto{\pgfqpoint{11.495899in}{1.510636in}}{\pgfqpoint{11.480169in}{1.494596in}}%
\pgfpathlineto{\pgfqpoint{11.480169in}{1.494596in}}%
\pgfpathclose%
\pgfpathmoveto{\pgfqpoint{11.621638in}{1.583762in}}%
\pgfpathquadraticcurveto{\pgfqpoint{11.613021in}{1.568959in}}{\pgfqpoint{11.608815in}{1.554445in}}%
\pgfpathquadraticcurveto{\pgfqpoint{11.604630in}{1.539952in}}{\pgfqpoint{11.604630in}{1.525067in}}%
\pgfpathquadraticcurveto{\pgfqpoint{11.604630in}{1.510203in}}{\pgfqpoint{11.608856in}{1.495606in}}%
\pgfpathquadraticcurveto{\pgfqpoint{11.613082in}{1.481010in}}{\pgfqpoint{11.621638in}{1.466249in}}%
\pgfpathlineto{\pgfqpoint{11.611330in}{1.466249in}}%
\pgfpathquadraticcurveto{\pgfqpoint{11.601682in}{1.481402in}}{\pgfqpoint{11.596878in}{1.496019in}}%
\pgfpathquadraticcurveto{\pgfqpoint{11.592074in}{1.510636in}}{\pgfqpoint{11.592074in}{1.525067in}}%
\pgfpathquadraticcurveto{\pgfqpoint{11.592074in}{1.539437in}}{\pgfqpoint{11.596837in}{1.553992in}}%
\pgfpathquadraticcurveto{\pgfqpoint{11.601620in}{1.568567in}}{\pgfqpoint{11.611330in}{1.583762in}}%
\pgfpathlineto{\pgfqpoint{11.621638in}{1.583762in}}%
\pgfpathlineto{\pgfqpoint{11.621638in}{1.583762in}}%
\pgfpathclose%
\pgfpathmoveto{\pgfqpoint{11.657531in}{1.494596in}}%
\pgfpathlineto{\pgfqpoint{11.702949in}{1.494596in}}%
\pgfpathlineto{\pgfqpoint{11.702949in}{1.483649in}}%
\pgfpathlineto{\pgfqpoint{11.641883in}{1.483649in}}%
\pgfpathlineto{\pgfqpoint{11.641883in}{1.494596in}}%
\pgfpathquadraticcurveto{\pgfqpoint{11.649285in}{1.502265in}}{\pgfqpoint{11.662067in}{1.515171in}}%
\pgfpathquadraticcurveto{\pgfqpoint{11.674870in}{1.528098in}}{\pgfqpoint{11.678148in}{1.531850in}}%
\pgfpathquadraticcurveto{\pgfqpoint{11.684394in}{1.538859in}}{\pgfqpoint{11.686868in}{1.543725in}}%
\pgfpathquadraticcurveto{\pgfqpoint{11.689363in}{1.548590in}}{\pgfqpoint{11.689363in}{1.553291in}}%
\pgfpathquadraticcurveto{\pgfqpoint{11.689363in}{1.560960in}}{\pgfqpoint{11.683982in}{1.565784in}}%
\pgfpathquadraticcurveto{\pgfqpoint{11.678601in}{1.570629in}}{\pgfqpoint{11.669963in}{1.570629in}}%
\pgfpathquadraticcurveto{\pgfqpoint{11.663840in}{1.570629in}}{\pgfqpoint{11.657036in}{1.568506in}}%
\pgfpathquadraticcurveto{\pgfqpoint{11.650254in}{1.566382in}}{\pgfqpoint{11.642523in}{1.562053in}}%
\pgfpathlineto{\pgfqpoint{11.642523in}{1.575206in}}%
\pgfpathquadraticcurveto{\pgfqpoint{11.650377in}{1.578360in}}{\pgfqpoint{11.657201in}{1.579968in}}%
\pgfpathquadraticcurveto{\pgfqpoint{11.664046in}{1.581576in}}{\pgfqpoint{11.669715in}{1.581576in}}%
\pgfpathquadraticcurveto{\pgfqpoint{11.684662in}{1.581576in}}{\pgfqpoint{11.693548in}{1.574093in}}%
\pgfpathquadraticcurveto{\pgfqpoint{11.702434in}{1.566629in}}{\pgfqpoint{11.702434in}{1.554136in}}%
\pgfpathquadraticcurveto{\pgfqpoint{11.702434in}{1.548198in}}{\pgfqpoint{11.700207in}{1.542879in}}%
\pgfpathquadraticcurveto{\pgfqpoint{11.698001in}{1.537581in}}{\pgfqpoint{11.692125in}{1.530365in}}%
\pgfpathquadraticcurveto{\pgfqpoint{11.690517in}{1.528489in}}{\pgfqpoint{11.681879in}{1.519562in}}%
\pgfpathquadraticcurveto{\pgfqpoint{11.673262in}{1.510636in}}{\pgfqpoint{11.657531in}{1.494596in}}%
\pgfpathlineto{\pgfqpoint{11.657531in}{1.494596in}}%
\pgfpathclose%
\pgfpathmoveto{\pgfqpoint{11.758098in}{1.571268in}}%
\pgfpathquadraticcurveto{\pgfqpoint{11.748058in}{1.571268in}}{\pgfqpoint{11.742986in}{1.561372in}}%
\pgfpathquadraticcurveto{\pgfqpoint{11.737935in}{1.551497in}}{\pgfqpoint{11.737935in}{1.531644in}}%
\pgfpathquadraticcurveto{\pgfqpoint{11.737935in}{1.511873in}}{\pgfqpoint{11.742986in}{1.501977in}}%
\pgfpathquadraticcurveto{\pgfqpoint{11.748058in}{1.492081in}}{\pgfqpoint{11.758098in}{1.492081in}}%
\pgfpathquadraticcurveto{\pgfqpoint{11.768220in}{1.492081in}}{\pgfqpoint{11.773271in}{1.501977in}}%
\pgfpathquadraticcurveto{\pgfqpoint{11.778343in}{1.511873in}}{\pgfqpoint{11.778343in}{1.531644in}}%
\pgfpathquadraticcurveto{\pgfqpoint{11.778343in}{1.551497in}}{\pgfqpoint{11.773271in}{1.561372in}}%
\pgfpathquadraticcurveto{\pgfqpoint{11.768220in}{1.571268in}}{\pgfqpoint{11.758098in}{1.571268in}}%
\pgfpathlineto{\pgfqpoint{11.758098in}{1.571268in}}%
\pgfpathclose%
\pgfpathmoveto{\pgfqpoint{11.758098in}{1.581576in}}%
\pgfpathquadraticcurveto{\pgfqpoint{11.774281in}{1.581576in}}{\pgfqpoint{11.782817in}{1.568774in}}%
\pgfpathquadraticcurveto{\pgfqpoint{11.791352in}{1.555991in}}{\pgfqpoint{11.791352in}{1.531644in}}%
\pgfpathquadraticcurveto{\pgfqpoint{11.791352in}{1.507358in}}{\pgfqpoint{11.782817in}{1.494555in}}%
\pgfpathquadraticcurveto{\pgfqpoint{11.774281in}{1.481773in}}{\pgfqpoint{11.758098in}{1.481773in}}%
\pgfpathquadraticcurveto{\pgfqpoint{11.741934in}{1.481773in}}{\pgfqpoint{11.733399in}{1.494555in}}%
\pgfpathquadraticcurveto{\pgfqpoint{11.724864in}{1.507358in}}{\pgfqpoint{11.724864in}{1.531644in}}%
\pgfpathquadraticcurveto{\pgfqpoint{11.724864in}{1.555991in}}{\pgfqpoint{11.733399in}{1.568774in}}%
\pgfpathquadraticcurveto{\pgfqpoint{11.741934in}{1.581576in}}{\pgfqpoint{11.758098in}{1.581576in}}%
\pgfpathlineto{\pgfqpoint{11.758098in}{1.581576in}}%
\pgfpathclose%
\pgfpathmoveto{\pgfqpoint{11.825431in}{1.494596in}}%
\pgfpathlineto{\pgfqpoint{11.870848in}{1.494596in}}%
\pgfpathlineto{\pgfqpoint{11.870848in}{1.483649in}}%
\pgfpathlineto{\pgfqpoint{11.809783in}{1.483649in}}%
\pgfpathlineto{\pgfqpoint{11.809783in}{1.494596in}}%
\pgfpathquadraticcurveto{\pgfqpoint{11.817184in}{1.502265in}}{\pgfqpoint{11.829966in}{1.515171in}}%
\pgfpathquadraticcurveto{\pgfqpoint{11.842769in}{1.528098in}}{\pgfqpoint{11.846047in}{1.531850in}}%
\pgfpathquadraticcurveto{\pgfqpoint{11.852294in}{1.538859in}}{\pgfqpoint{11.854768in}{1.543725in}}%
\pgfpathquadraticcurveto{\pgfqpoint{11.857262in}{1.548590in}}{\pgfqpoint{11.857262in}{1.553291in}}%
\pgfpathquadraticcurveto{\pgfqpoint{11.857262in}{1.560960in}}{\pgfqpoint{11.851881in}{1.565784in}}%
\pgfpathquadraticcurveto{\pgfqpoint{11.846500in}{1.570629in}}{\pgfqpoint{11.837862in}{1.570629in}}%
\pgfpathquadraticcurveto{\pgfqpoint{11.831739in}{1.570629in}}{\pgfqpoint{11.824936in}{1.568506in}}%
\pgfpathquadraticcurveto{\pgfqpoint{11.818153in}{1.566382in}}{\pgfqpoint{11.810422in}{1.562053in}}%
\pgfpathlineto{\pgfqpoint{11.810422in}{1.575206in}}%
\pgfpathquadraticcurveto{\pgfqpoint{11.818277in}{1.578360in}}{\pgfqpoint{11.825101in}{1.579968in}}%
\pgfpathquadraticcurveto{\pgfqpoint{11.831945in}{1.581576in}}{\pgfqpoint{11.837615in}{1.581576in}}%
\pgfpathquadraticcurveto{\pgfqpoint{11.852562in}{1.581576in}}{\pgfqpoint{11.861447in}{1.574093in}}%
\pgfpathquadraticcurveto{\pgfqpoint{11.870333in}{1.566629in}}{\pgfqpoint{11.870333in}{1.554136in}}%
\pgfpathquadraticcurveto{\pgfqpoint{11.870333in}{1.548198in}}{\pgfqpoint{11.868106in}{1.542879in}}%
\pgfpathquadraticcurveto{\pgfqpoint{11.865900in}{1.537581in}}{\pgfqpoint{11.860025in}{1.530365in}}%
\pgfpathquadraticcurveto{\pgfqpoint{11.858417in}{1.528489in}}{\pgfqpoint{11.849778in}{1.519562in}}%
\pgfpathquadraticcurveto{\pgfqpoint{11.841161in}{1.510636in}}{\pgfqpoint{11.825431in}{1.494596in}}%
\pgfpathlineto{\pgfqpoint{11.825431in}{1.494596in}}%
\pgfpathclose%
\pgfpathmoveto{\pgfqpoint{11.898309in}{1.579845in}}%
\pgfpathlineto{\pgfqpoint{11.949396in}{1.579845in}}%
\pgfpathlineto{\pgfqpoint{11.949396in}{1.568877in}}%
\pgfpathlineto{\pgfqpoint{11.910225in}{1.568877in}}%
\pgfpathlineto{\pgfqpoint{11.910225in}{1.545312in}}%
\pgfpathquadraticcurveto{\pgfqpoint{11.913050in}{1.546281in}}{\pgfqpoint{11.915874in}{1.546755in}}%
\pgfpathquadraticcurveto{\pgfqpoint{11.918719in}{1.547230in}}{\pgfqpoint{11.921564in}{1.547230in}}%
\pgfpathquadraticcurveto{\pgfqpoint{11.937666in}{1.547230in}}{\pgfqpoint{11.947067in}{1.538406in}}%
\pgfpathquadraticcurveto{\pgfqpoint{11.956488in}{1.529582in}}{\pgfqpoint{11.956488in}{1.514511in}}%
\pgfpathquadraticcurveto{\pgfqpoint{11.956488in}{1.498987in}}{\pgfqpoint{11.946819in}{1.490370in}}%
\pgfpathquadraticcurveto{\pgfqpoint{11.937150in}{1.481773in}}{\pgfqpoint{11.919565in}{1.481773in}}%
\pgfpathquadraticcurveto{\pgfqpoint{11.913503in}{1.481773in}}{\pgfqpoint{11.907215in}{1.482804in}}%
\pgfpathquadraticcurveto{\pgfqpoint{11.900948in}{1.483834in}}{\pgfqpoint{11.894248in}{1.485896in}}%
\pgfpathlineto{\pgfqpoint{11.894248in}{1.498987in}}%
\pgfpathquadraticcurveto{\pgfqpoint{11.900041in}{1.495833in}}{\pgfqpoint{11.906226in}{1.494287in}}%
\pgfpathquadraticcurveto{\pgfqpoint{11.912411in}{1.492741in}}{\pgfqpoint{11.919297in}{1.492741in}}%
\pgfpathquadraticcurveto{\pgfqpoint{11.930450in}{1.492741in}}{\pgfqpoint{11.936944in}{1.498596in}}%
\pgfpathquadraticcurveto{\pgfqpoint{11.943459in}{1.504451in}}{\pgfqpoint{11.943459in}{1.514511in}}%
\pgfpathquadraticcurveto{\pgfqpoint{11.943459in}{1.524552in}}{\pgfqpoint{11.936944in}{1.530407in}}%
\pgfpathquadraticcurveto{\pgfqpoint{11.930450in}{1.536282in}}{\pgfqpoint{11.919297in}{1.536282in}}%
\pgfpathquadraticcurveto{\pgfqpoint{11.914081in}{1.536282in}}{\pgfqpoint{11.908885in}{1.535128in}}%
\pgfpathquadraticcurveto{\pgfqpoint{11.903711in}{1.533973in}}{\pgfqpoint{11.898309in}{1.531520in}}%
\pgfpathlineto{\pgfqpoint{11.898309in}{1.579845in}}%
\pgfpathlineto{\pgfqpoint{11.898309in}{1.579845in}}%
\pgfpathclose%
\pgfpathmoveto{\pgfqpoint{11.978589in}{1.583762in}}%
\pgfpathlineto{\pgfqpoint{11.988897in}{1.583762in}}%
\pgfpathquadraticcurveto{\pgfqpoint{11.998546in}{1.568567in}}{\pgfqpoint{12.003349in}{1.553992in}}%
\pgfpathquadraticcurveto{\pgfqpoint{12.008153in}{1.539437in}}{\pgfqpoint{12.008153in}{1.525067in}}%
\pgfpathquadraticcurveto{\pgfqpoint{12.008153in}{1.510636in}}{\pgfqpoint{12.003349in}{1.496019in}}%
\pgfpathquadraticcurveto{\pgfqpoint{11.998546in}{1.481402in}}{\pgfqpoint{11.988897in}{1.466249in}}%
\pgfpathlineto{\pgfqpoint{11.978589in}{1.466249in}}%
\pgfpathquadraticcurveto{\pgfqpoint{11.987145in}{1.481010in}}{\pgfqpoint{11.991371in}{1.495606in}}%
\pgfpathquadraticcurveto{\pgfqpoint{11.995598in}{1.510203in}}{\pgfqpoint{11.995598in}{1.525067in}}%
\pgfpathquadraticcurveto{\pgfqpoint{11.995598in}{1.539952in}}{\pgfqpoint{11.991371in}{1.554445in}}%
\pgfpathquadraticcurveto{\pgfqpoint{11.987145in}{1.568959in}}{\pgfqpoint{11.978589in}{1.583762in}}%
\pgfpathlineto{\pgfqpoint{11.978589in}{1.583762in}}%
\pgfpathclose%
\pgfusepath{fill}%
\end{pgfscope}%
\begin{pgfscope}%
\pgfsetbuttcap%
\pgfsetroundjoin%
\definecolor{currentfill}{rgb}{0.309804,0.372549,0.419608}%
\pgfsetfillcolor{currentfill}%
\pgfsetlinewidth{0.000000pt}%
\definecolor{currentstroke}{rgb}{0.309804,0.372549,0.419608}%
\pgfsetstrokecolor{currentstroke}%
\pgfsetdash{}{0pt}%
\pgfpathmoveto{\pgfqpoint{0.998246in}{1.341160in}}%
\pgfpathlineto{\pgfqpoint{0.998246in}{1.329182in}}%
\pgfpathquadraticcurveto{\pgfqpoint{0.992500in}{1.334532in}}{\pgfqpoint{0.985998in}{1.337162in}}%
\pgfpathquadraticcurveto{\pgfqpoint{0.979495in}{1.339809in}}{\pgfqpoint{0.972183in}{1.339809in}}%
\pgfpathquadraticcurveto{\pgfqpoint{0.957773in}{1.339809in}}{\pgfqpoint{0.950118in}{1.331001in}}%
\pgfpathquadraticcurveto{\pgfqpoint{0.942462in}{1.322194in}}{\pgfqpoint{0.942462in}{1.305532in}}%
\pgfpathquadraticcurveto{\pgfqpoint{0.942462in}{1.288925in}}{\pgfqpoint{0.950118in}{1.280117in}}%
\pgfpathquadraticcurveto{\pgfqpoint{0.957773in}{1.271309in}}{\pgfqpoint{0.972183in}{1.271309in}}%
\pgfpathquadraticcurveto{\pgfqpoint{0.979495in}{1.271309in}}{\pgfqpoint{0.985998in}{1.273957in}}%
\pgfpathquadraticcurveto{\pgfqpoint{0.992500in}{1.276605in}}{\pgfqpoint{0.998246in}{1.281954in}}%
\pgfpathlineto{\pgfqpoint{0.998246in}{1.270066in}}%
\pgfpathquadraticcurveto{\pgfqpoint{0.992284in}{1.266014in}}{\pgfqpoint{0.985602in}{1.263978in}}%
\pgfpathquadraticcurveto{\pgfqpoint{0.978937in}{1.261961in}}{\pgfqpoint{0.971516in}{1.261961in}}%
\pgfpathquadraticcurveto{\pgfqpoint{0.952423in}{1.261961in}}{\pgfqpoint{0.941436in}{1.273633in}}%
\pgfpathquadraticcurveto{\pgfqpoint{0.930466in}{1.285323in}}{\pgfqpoint{0.930466in}{1.305532in}}%
\pgfpathquadraticcurveto{\pgfqpoint{0.930466in}{1.325796in}}{\pgfqpoint{0.941436in}{1.337468in}}%
\pgfpathquadraticcurveto{\pgfqpoint{0.952423in}{1.349158in}}{\pgfqpoint{0.971516in}{1.349158in}}%
\pgfpathquadraticcurveto{\pgfqpoint{0.979045in}{1.349158in}}{\pgfqpoint{0.985710in}{1.347158in}}%
\pgfpathquadraticcurveto{\pgfqpoint{0.992392in}{1.345159in}}{\pgfqpoint{0.998246in}{1.341160in}}%
\pgfpathlineto{\pgfqpoint{0.998246in}{1.341160in}}%
\pgfpathclose%
\pgfpathmoveto{\pgfqpoint{1.015358in}{1.326643in}}%
\pgfpathlineto{\pgfqpoint{1.025715in}{1.326643in}}%
\pgfpathlineto{\pgfqpoint{1.025715in}{1.263600in}}%
\pgfpathlineto{\pgfqpoint{1.015358in}{1.263600in}}%
\pgfpathlineto{\pgfqpoint{1.015358in}{1.326643in}}%
\pgfpathlineto{\pgfqpoint{1.015358in}{1.326643in}}%
\pgfpathclose%
\pgfpathmoveto{\pgfqpoint{1.015358in}{1.351193in}}%
\pgfpathlineto{\pgfqpoint{1.025715in}{1.351193in}}%
\pgfpathlineto{\pgfqpoint{1.025715in}{1.338062in}}%
\pgfpathlineto{\pgfqpoint{1.015358in}{1.338062in}}%
\pgfpathlineto{\pgfqpoint{1.015358in}{1.351193in}}%
\pgfpathlineto{\pgfqpoint{1.015358in}{1.351193in}}%
\pgfpathclose%
\pgfpathmoveto{\pgfqpoint{1.083912in}{1.316970in}}%
\pgfpathquadraticcurveto{\pgfqpoint{1.082165in}{1.317979in}}{\pgfqpoint{1.080111in}{1.318447in}}%
\pgfpathquadraticcurveto{\pgfqpoint{1.078058in}{1.318933in}}{\pgfqpoint{1.075590in}{1.318933in}}%
\pgfpathquadraticcurveto{\pgfqpoint{1.066800in}{1.318933in}}{\pgfqpoint{1.062099in}{1.313223in}}%
\pgfpathquadraticcurveto{\pgfqpoint{1.057398in}{1.307514in}}{\pgfqpoint{1.057398in}{1.296814in}}%
\pgfpathlineto{\pgfqpoint{1.057398in}{1.263600in}}%
\pgfpathlineto{\pgfqpoint{1.046987in}{1.263600in}}%
\pgfpathlineto{\pgfqpoint{1.046987in}{1.326643in}}%
\pgfpathlineto{\pgfqpoint{1.057398in}{1.326643in}}%
\pgfpathlineto{\pgfqpoint{1.057398in}{1.316844in}}%
\pgfpathquadraticcurveto{\pgfqpoint{1.060676in}{1.322590in}}{\pgfqpoint{1.065900in}{1.325364in}}%
\pgfpathquadraticcurveto{\pgfqpoint{1.071141in}{1.328156in}}{\pgfqpoint{1.078634in}{1.328156in}}%
\pgfpathquadraticcurveto{\pgfqpoint{1.079697in}{1.328156in}}{\pgfqpoint{1.080994in}{1.328011in}}%
\pgfpathquadraticcurveto{\pgfqpoint{1.082291in}{1.327885in}}{\pgfqpoint{1.083858in}{1.327597in}}%
\pgfpathlineto{\pgfqpoint{1.083912in}{1.316970in}}%
\pgfpathlineto{\pgfqpoint{1.083912in}{1.316970in}}%
\pgfpathclose%
\pgfpathmoveto{\pgfqpoint{1.137606in}{1.324229in}}%
\pgfpathlineto{\pgfqpoint{1.137606in}{1.314538in}}%
\pgfpathquadraticcurveto{\pgfqpoint{1.133211in}{1.316970in}}{\pgfqpoint{1.128798in}{1.318177in}}%
\pgfpathquadraticcurveto{\pgfqpoint{1.124385in}{1.319384in}}{\pgfqpoint{1.119882in}{1.319384in}}%
\pgfpathquadraticcurveto{\pgfqpoint{1.109795in}{1.319384in}}{\pgfqpoint{1.104212in}{1.312989in}}%
\pgfpathquadraticcurveto{\pgfqpoint{1.098646in}{1.306613in}}{\pgfqpoint{1.098646in}{1.295067in}}%
\pgfpathquadraticcurveto{\pgfqpoint{1.098646in}{1.283521in}}{\pgfqpoint{1.104212in}{1.277127in}}%
\pgfpathquadraticcurveto{\pgfqpoint{1.109795in}{1.270751in}}{\pgfqpoint{1.119882in}{1.270751in}}%
\pgfpathquadraticcurveto{\pgfqpoint{1.124385in}{1.270751in}}{\pgfqpoint{1.128798in}{1.271958in}}%
\pgfpathquadraticcurveto{\pgfqpoint{1.133211in}{1.273164in}}{\pgfqpoint{1.137606in}{1.275596in}}%
\pgfpathlineto{\pgfqpoint{1.137606in}{1.266014in}}%
\pgfpathquadraticcurveto{\pgfqpoint{1.133265in}{1.263996in}}{\pgfqpoint{1.128618in}{1.262988in}}%
\pgfpathquadraticcurveto{\pgfqpoint{1.123989in}{1.261961in}}{\pgfqpoint{1.118747in}{1.261961in}}%
\pgfpathquadraticcurveto{\pgfqpoint{1.104500in}{1.261961in}}{\pgfqpoint{1.096106in}{1.270913in}}%
\pgfpathquadraticcurveto{\pgfqpoint{1.087730in}{1.279865in}}{\pgfqpoint{1.087730in}{1.295067in}}%
\pgfpathquadraticcurveto{\pgfqpoint{1.087730in}{1.310486in}}{\pgfqpoint{1.096196in}{1.319312in}}%
\pgfpathquadraticcurveto{\pgfqpoint{1.104680in}{1.328156in}}{\pgfqpoint{1.119432in}{1.328156in}}%
\pgfpathquadraticcurveto{\pgfqpoint{1.124205in}{1.328156in}}{\pgfqpoint{1.128762in}{1.327165in}}%
\pgfpathquadraticcurveto{\pgfqpoint{1.133319in}{1.326192in}}{\pgfqpoint{1.137606in}{1.324229in}}%
\pgfpathlineto{\pgfqpoint{1.137606in}{1.324229in}}%
\pgfpathclose%
\pgfpathmoveto{\pgfqpoint{1.155618in}{1.351193in}}%
\pgfpathlineto{\pgfqpoint{1.165975in}{1.351193in}}%
\pgfpathlineto{\pgfqpoint{1.165975in}{1.263600in}}%
\pgfpathlineto{\pgfqpoint{1.155618in}{1.263600in}}%
\pgfpathlineto{\pgfqpoint{1.155618in}{1.351193in}}%
\pgfpathlineto{\pgfqpoint{1.155618in}{1.351193in}}%
\pgfpathclose%
\pgfpathmoveto{\pgfqpoint{1.241572in}{1.297715in}}%
\pgfpathlineto{\pgfqpoint{1.241572in}{1.292654in}}%
\pgfpathlineto{\pgfqpoint{1.193948in}{1.292654in}}%
\pgfpathquadraticcurveto{\pgfqpoint{1.194633in}{1.281954in}}{\pgfqpoint{1.200396in}{1.276353in}}%
\pgfpathquadraticcurveto{\pgfqpoint{1.206160in}{1.270751in}}{\pgfqpoint{1.216463in}{1.270751in}}%
\pgfpathquadraticcurveto{\pgfqpoint{1.222425in}{1.270751in}}{\pgfqpoint{1.228027in}{1.272210in}}%
\pgfpathquadraticcurveto{\pgfqpoint{1.233629in}{1.273669in}}{\pgfqpoint{1.239159in}{1.276605in}}%
\pgfpathlineto{\pgfqpoint{1.239159in}{1.266806in}}%
\pgfpathquadraticcurveto{\pgfqpoint{1.233575in}{1.264447in}}{\pgfqpoint{1.227721in}{1.263204in}}%
\pgfpathquadraticcurveto{\pgfqpoint{1.221867in}{1.261961in}}{\pgfqpoint{1.215851in}{1.261961in}}%
\pgfpathquadraticcurveto{\pgfqpoint{1.200757in}{1.261961in}}{\pgfqpoint{1.191949in}{1.270733in}}%
\pgfpathquadraticcurveto{\pgfqpoint{1.183141in}{1.279523in}}{\pgfqpoint{1.183141in}{1.294509in}}%
\pgfpathquadraticcurveto{\pgfqpoint{1.183141in}{1.309981in}}{\pgfqpoint{1.191498in}{1.319059in}}%
\pgfpathquadraticcurveto{\pgfqpoint{1.199856in}{1.328156in}}{\pgfqpoint{1.214050in}{1.328156in}}%
\pgfpathquadraticcurveto{\pgfqpoint{1.226766in}{1.328156in}}{\pgfqpoint{1.234169in}{1.319960in}}%
\pgfpathquadraticcurveto{\pgfqpoint{1.241572in}{1.311783in}}{\pgfqpoint{1.241572in}{1.297715in}}%
\pgfpathlineto{\pgfqpoint{1.241572in}{1.297715in}}%
\pgfpathclose%
\pgfpathmoveto{\pgfqpoint{1.231215in}{1.300759in}}%
\pgfpathquadraticcurveto{\pgfqpoint{1.231107in}{1.309243in}}{\pgfqpoint{1.226460in}{1.314304in}}%
\pgfpathquadraticcurveto{\pgfqpoint{1.221813in}{1.319384in}}{\pgfqpoint{1.214158in}{1.319384in}}%
\pgfpathquadraticcurveto{\pgfqpoint{1.205494in}{1.319384in}}{\pgfqpoint{1.200288in}{1.314484in}}%
\pgfpathquadraticcurveto{\pgfqpoint{1.195083in}{1.309585in}}{\pgfqpoint{1.194290in}{1.300687in}}%
\pgfpathlineto{\pgfqpoint{1.231215in}{1.300759in}}%
\pgfpathlineto{\pgfqpoint{1.231215in}{1.300759in}}%
\pgfpathclose%
\pgfpathmoveto{\pgfqpoint{1.298761in}{1.324787in}}%
\pgfpathlineto{\pgfqpoint{1.298761in}{1.314989in}}%
\pgfpathquadraticcurveto{\pgfqpoint{1.294384in}{1.317240in}}{\pgfqpoint{1.289647in}{1.318357in}}%
\pgfpathquadraticcurveto{\pgfqpoint{1.284928in}{1.319492in}}{\pgfqpoint{1.279848in}{1.319492in}}%
\pgfpathquadraticcurveto{\pgfqpoint{1.272139in}{1.319492in}}{\pgfqpoint{1.268284in}{1.317132in}}%
\pgfpathquadraticcurveto{\pgfqpoint{1.264430in}{1.314773in}}{\pgfqpoint{1.264430in}{1.310035in}}%
\pgfpathquadraticcurveto{\pgfqpoint{1.264430in}{1.306433in}}{\pgfqpoint{1.267186in}{1.304380in}}%
\pgfpathquadraticcurveto{\pgfqpoint{1.269941in}{1.302326in}}{\pgfqpoint{1.278281in}{1.300471in}}%
\pgfpathlineto{\pgfqpoint{1.281829in}{1.299678in}}%
\pgfpathquadraticcurveto{\pgfqpoint{1.292853in}{1.297319in}}{\pgfqpoint{1.297500in}{1.293014in}}%
\pgfpathquadraticcurveto{\pgfqpoint{1.302147in}{1.288709in}}{\pgfqpoint{1.302147in}{1.281000in}}%
\pgfpathquadraticcurveto{\pgfqpoint{1.302147in}{1.272210in}}{\pgfqpoint{1.295194in}{1.267076in}}%
\pgfpathquadraticcurveto{\pgfqpoint{1.288242in}{1.261961in}}{\pgfqpoint{1.276084in}{1.261961in}}%
\pgfpathquadraticcurveto{\pgfqpoint{1.271022in}{1.261961in}}{\pgfqpoint{1.265528in}{1.262952in}}%
\pgfpathquadraticcurveto{\pgfqpoint{1.260035in}{1.263942in}}{\pgfqpoint{1.253965in}{1.265906in}}%
\pgfpathlineto{\pgfqpoint{1.253965in}{1.276605in}}%
\pgfpathquadraticcurveto{\pgfqpoint{1.259711in}{1.273615in}}{\pgfqpoint{1.265276in}{1.272120in}}%
\pgfpathquadraticcurveto{\pgfqpoint{1.270842in}{1.270643in}}{\pgfqpoint{1.276318in}{1.270643in}}%
\pgfpathquadraticcurveto{\pgfqpoint{1.283631in}{1.270643in}}{\pgfqpoint{1.287557in}{1.273146in}}%
\pgfpathquadraticcurveto{\pgfqpoint{1.291502in}{1.275650in}}{\pgfqpoint{1.291502in}{1.280207in}}%
\pgfpathquadraticcurveto{\pgfqpoint{1.291502in}{1.284422in}}{\pgfqpoint{1.288656in}{1.286674in}}%
\pgfpathquadraticcurveto{\pgfqpoint{1.285828in}{1.288925in}}{\pgfqpoint{1.276192in}{1.291014in}}%
\pgfpathlineto{\pgfqpoint{1.272589in}{1.291861in}}%
\pgfpathquadraticcurveto{\pgfqpoint{1.262971in}{1.293878in}}{\pgfqpoint{1.258684in}{1.298075in}}%
\pgfpathquadraticcurveto{\pgfqpoint{1.254415in}{1.302272in}}{\pgfqpoint{1.254415in}{1.309585in}}%
\pgfpathquadraticcurveto{\pgfqpoint{1.254415in}{1.318483in}}{\pgfqpoint{1.260719in}{1.323310in}}%
\pgfpathquadraticcurveto{\pgfqpoint{1.267023in}{1.328156in}}{\pgfqpoint{1.278623in}{1.328156in}}%
\pgfpathquadraticcurveto{\pgfqpoint{1.284351in}{1.328156in}}{\pgfqpoint{1.289413in}{1.327309in}}%
\pgfpathquadraticcurveto{\pgfqpoint{1.294492in}{1.326480in}}{\pgfqpoint{1.298761in}{1.324787in}}%
\pgfpathlineto{\pgfqpoint{1.298761in}{1.324787in}}%
\pgfpathclose%
\pgfpathmoveto{\pgfqpoint{1.321276in}{1.277902in}}%
\pgfpathlineto{\pgfqpoint{1.333146in}{1.277902in}}%
\pgfpathlineto{\pgfqpoint{1.333146in}{1.263600in}}%
\pgfpathlineto{\pgfqpoint{1.321276in}{1.263600in}}%
\pgfpathlineto{\pgfqpoint{1.321276in}{1.277902in}}%
\pgfpathlineto{\pgfqpoint{1.321276in}{1.277902in}}%
\pgfpathclose%
\pgfpathmoveto{\pgfqpoint{1.321276in}{1.323202in}}%
\pgfpathlineto{\pgfqpoint{1.333146in}{1.323202in}}%
\pgfpathlineto{\pgfqpoint{1.333146in}{1.308919in}}%
\pgfpathlineto{\pgfqpoint{1.321276in}{1.308919in}}%
\pgfpathlineto{\pgfqpoint{1.321276in}{1.323202in}}%
\pgfpathlineto{\pgfqpoint{1.321276in}{1.323202in}}%
\pgfpathclose%
\pgfusepath{fill}%
\end{pgfscope}%
\begin{pgfscope}%
\pgfsetbuttcap%
\pgfsetroundjoin%
\definecolor{currentfill}{rgb}{0.309804,0.372549,0.419608}%
\pgfsetfillcolor{currentfill}%
\pgfsetlinewidth{0.000000pt}%
\definecolor{currentstroke}{rgb}{0.309804,0.372549,0.419608}%
\pgfsetstrokecolor{currentstroke}%
\pgfsetdash{}{0pt}%
\pgfpathmoveto{\pgfqpoint{1.458231in}{1.324787in}}%
\pgfpathlineto{\pgfqpoint{1.458231in}{1.314989in}}%
\pgfpathquadraticcurveto{\pgfqpoint{1.453854in}{1.317240in}}{\pgfqpoint{1.449117in}{1.318357in}}%
\pgfpathquadraticcurveto{\pgfqpoint{1.444398in}{1.319492in}}{\pgfqpoint{1.439319in}{1.319492in}}%
\pgfpathquadraticcurveto{\pgfqpoint{1.431609in}{1.319492in}}{\pgfqpoint{1.427755in}{1.317132in}}%
\pgfpathquadraticcurveto{\pgfqpoint{1.423900in}{1.314773in}}{\pgfqpoint{1.423900in}{1.310035in}}%
\pgfpathquadraticcurveto{\pgfqpoint{1.423900in}{1.306433in}}{\pgfqpoint{1.426656in}{1.304380in}}%
\pgfpathquadraticcurveto{\pgfqpoint{1.429412in}{1.302326in}}{\pgfqpoint{1.437752in}{1.300471in}}%
\pgfpathlineto{\pgfqpoint{1.441300in}{1.299678in}}%
\pgfpathquadraticcurveto{\pgfqpoint{1.452323in}{1.297319in}}{\pgfqpoint{1.456970in}{1.293014in}}%
\pgfpathquadraticcurveto{\pgfqpoint{1.461618in}{1.288709in}}{\pgfqpoint{1.461618in}{1.281000in}}%
\pgfpathquadraticcurveto{\pgfqpoint{1.461618in}{1.272210in}}{\pgfqpoint{1.454665in}{1.267076in}}%
\pgfpathquadraticcurveto{\pgfqpoint{1.447712in}{1.261961in}}{\pgfqpoint{1.435554in}{1.261961in}}%
\pgfpathquadraticcurveto{\pgfqpoint{1.430493in}{1.261961in}}{\pgfqpoint{1.424999in}{1.262952in}}%
\pgfpathquadraticcurveto{\pgfqpoint{1.419505in}{1.263942in}}{\pgfqpoint{1.413435in}{1.265906in}}%
\pgfpathlineto{\pgfqpoint{1.413435in}{1.276605in}}%
\pgfpathquadraticcurveto{\pgfqpoint{1.419181in}{1.273615in}}{\pgfqpoint{1.424747in}{1.272120in}}%
\pgfpathquadraticcurveto{\pgfqpoint{1.430313in}{1.270643in}}{\pgfqpoint{1.435788in}{1.270643in}}%
\pgfpathquadraticcurveto{\pgfqpoint{1.443101in}{1.270643in}}{\pgfqpoint{1.447028in}{1.273146in}}%
\pgfpathquadraticcurveto{\pgfqpoint{1.450972in}{1.275650in}}{\pgfqpoint{1.450972in}{1.280207in}}%
\pgfpathquadraticcurveto{\pgfqpoint{1.450972in}{1.284422in}}{\pgfqpoint{1.448127in}{1.286674in}}%
\pgfpathquadraticcurveto{\pgfqpoint{1.445299in}{1.288925in}}{\pgfqpoint{1.435662in}{1.291014in}}%
\pgfpathlineto{\pgfqpoint{1.432060in}{1.291861in}}%
\pgfpathquadraticcurveto{\pgfqpoint{1.422441in}{1.293878in}}{\pgfqpoint{1.418154in}{1.298075in}}%
\pgfpathquadraticcurveto{\pgfqpoint{1.413885in}{1.302272in}}{\pgfqpoint{1.413885in}{1.309585in}}%
\pgfpathquadraticcurveto{\pgfqpoint{1.413885in}{1.318483in}}{\pgfqpoint{1.420190in}{1.323310in}}%
\pgfpathquadraticcurveto{\pgfqpoint{1.426494in}{1.328156in}}{\pgfqpoint{1.438094in}{1.328156in}}%
\pgfpathquadraticcurveto{\pgfqpoint{1.443822in}{1.328156in}}{\pgfqpoint{1.448883in}{1.327309in}}%
\pgfpathquadraticcurveto{\pgfqpoint{1.453962in}{1.326480in}}{\pgfqpoint{1.458231in}{1.324787in}}%
\pgfpathlineto{\pgfqpoint{1.458231in}{1.324787in}}%
\pgfpathclose%
\pgfpathmoveto{\pgfqpoint{1.488348in}{1.344547in}}%
\pgfpathlineto{\pgfqpoint{1.488348in}{1.326643in}}%
\pgfpathlineto{\pgfqpoint{1.509674in}{1.326643in}}%
\pgfpathlineto{\pgfqpoint{1.509674in}{1.318591in}}%
\pgfpathlineto{\pgfqpoint{1.488348in}{1.318591in}}%
\pgfpathlineto{\pgfqpoint{1.488348in}{1.284368in}}%
\pgfpathquadraticcurveto{\pgfqpoint{1.488348in}{1.276659in}}{\pgfqpoint{1.490455in}{1.274461in}}%
\pgfpathquadraticcurveto{\pgfqpoint{1.492563in}{1.272264in}}{\pgfqpoint{1.499047in}{1.272264in}}%
\pgfpathlineto{\pgfqpoint{1.509674in}{1.272264in}}%
\pgfpathlineto{\pgfqpoint{1.509674in}{1.263600in}}%
\pgfpathlineto{\pgfqpoint{1.499047in}{1.263600in}}%
\pgfpathquadraticcurveto{\pgfqpoint{1.487051in}{1.263600in}}{\pgfqpoint{1.482494in}{1.268067in}}%
\pgfpathquadraticcurveto{\pgfqpoint{1.477937in}{1.272552in}}{\pgfqpoint{1.477937in}{1.284368in}}%
\pgfpathlineto{\pgfqpoint{1.477937in}{1.318591in}}%
\pgfpathlineto{\pgfqpoint{1.470336in}{1.318591in}}%
\pgfpathlineto{\pgfqpoint{1.470336in}{1.326643in}}%
\pgfpathlineto{\pgfqpoint{1.477937in}{1.326643in}}%
\pgfpathlineto{\pgfqpoint{1.477937in}{1.344547in}}%
\pgfpathlineto{\pgfqpoint{1.488348in}{1.344547in}}%
\pgfpathlineto{\pgfqpoint{1.488348in}{1.344547in}}%
\pgfpathclose%
\pgfpathmoveto{\pgfqpoint{1.551949in}{1.295283in}}%
\pgfpathquadraticcurveto{\pgfqpoint{1.539394in}{1.295283in}}{\pgfqpoint{1.534549in}{1.292419in}}%
\pgfpathquadraticcurveto{\pgfqpoint{1.529704in}{1.289556in}}{\pgfqpoint{1.529704in}{1.282621in}}%
\pgfpathquadraticcurveto{\pgfqpoint{1.529704in}{1.277109in}}{\pgfqpoint{1.533342in}{1.273867in}}%
\pgfpathquadraticcurveto{\pgfqpoint{1.536980in}{1.270643in}}{\pgfqpoint{1.543213in}{1.270643in}}%
\pgfpathquadraticcurveto{\pgfqpoint{1.551840in}{1.270643in}}{\pgfqpoint{1.557046in}{1.276749in}}%
\pgfpathquadraticcurveto{\pgfqpoint{1.562252in}{1.282855in}}{\pgfqpoint{1.562252in}{1.292978in}}%
\pgfpathlineto{\pgfqpoint{1.562252in}{1.295283in}}%
\pgfpathlineto{\pgfqpoint{1.551949in}{1.295283in}}%
\pgfpathlineto{\pgfqpoint{1.551949in}{1.295283in}}%
\pgfpathclose%
\pgfpathmoveto{\pgfqpoint{1.572609in}{1.299570in}}%
\pgfpathlineto{\pgfqpoint{1.572609in}{1.263600in}}%
\pgfpathlineto{\pgfqpoint{1.562252in}{1.263600in}}%
\pgfpathlineto{\pgfqpoint{1.562252in}{1.273164in}}%
\pgfpathquadraticcurveto{\pgfqpoint{1.558703in}{1.267437in}}{\pgfqpoint{1.553408in}{1.264699in}}%
\pgfpathquadraticcurveto{\pgfqpoint{1.548112in}{1.261961in}}{\pgfqpoint{1.540457in}{1.261961in}}%
\pgfpathquadraticcurveto{\pgfqpoint{1.530784in}{1.261961in}}{\pgfqpoint{1.525056in}{1.267401in}}%
\pgfpathquadraticcurveto{\pgfqpoint{1.519347in}{1.272840in}}{\pgfqpoint{1.519347in}{1.281954in}}%
\pgfpathquadraticcurveto{\pgfqpoint{1.519347in}{1.292582in}}{\pgfqpoint{1.526461in}{1.297985in}}%
\pgfpathquadraticcurveto{\pgfqpoint{1.533594in}{1.303389in}}{\pgfqpoint{1.547716in}{1.303389in}}%
\pgfpathlineto{\pgfqpoint{1.562252in}{1.303389in}}%
\pgfpathlineto{\pgfqpoint{1.562252in}{1.304416in}}%
\pgfpathquadraticcurveto{\pgfqpoint{1.562252in}{1.311566in}}{\pgfqpoint{1.557550in}{1.315475in}}%
\pgfpathquadraticcurveto{\pgfqpoint{1.552849in}{1.319384in}}{\pgfqpoint{1.544347in}{1.319384in}}%
\pgfpathquadraticcurveto{\pgfqpoint{1.538944in}{1.319384in}}{\pgfqpoint{1.533810in}{1.318087in}}%
\pgfpathquadraticcurveto{\pgfqpoint{1.528695in}{1.316790in}}{\pgfqpoint{1.523976in}{1.314196in}}%
\pgfpathlineto{\pgfqpoint{1.523976in}{1.323779in}}%
\pgfpathquadraticcurveto{\pgfqpoint{1.529650in}{1.325976in}}{\pgfqpoint{1.534999in}{1.327057in}}%
\pgfpathquadraticcurveto{\pgfqpoint{1.540349in}{1.328156in}}{\pgfqpoint{1.545410in}{1.328156in}}%
\pgfpathquadraticcurveto{\pgfqpoint{1.559099in}{1.328156in}}{\pgfqpoint{1.565854in}{1.321059in}}%
\pgfpathquadraticcurveto{\pgfqpoint{1.572609in}{1.313980in}}{\pgfqpoint{1.572609in}{1.299570in}}%
\pgfpathlineto{\pgfqpoint{1.572609in}{1.299570in}}%
\pgfpathclose%
\pgfpathmoveto{\pgfqpoint{1.604184in}{1.344547in}}%
\pgfpathlineto{\pgfqpoint{1.604184in}{1.326643in}}%
\pgfpathlineto{\pgfqpoint{1.625510in}{1.326643in}}%
\pgfpathlineto{\pgfqpoint{1.625510in}{1.318591in}}%
\pgfpathlineto{\pgfqpoint{1.604184in}{1.318591in}}%
\pgfpathlineto{\pgfqpoint{1.604184in}{1.284368in}}%
\pgfpathquadraticcurveto{\pgfqpoint{1.604184in}{1.276659in}}{\pgfqpoint{1.606291in}{1.274461in}}%
\pgfpathquadraticcurveto{\pgfqpoint{1.608399in}{1.272264in}}{\pgfqpoint{1.614883in}{1.272264in}}%
\pgfpathlineto{\pgfqpoint{1.625510in}{1.272264in}}%
\pgfpathlineto{\pgfqpoint{1.625510in}{1.263600in}}%
\pgfpathlineto{\pgfqpoint{1.614883in}{1.263600in}}%
\pgfpathquadraticcurveto{\pgfqpoint{1.602887in}{1.263600in}}{\pgfqpoint{1.598330in}{1.268067in}}%
\pgfpathquadraticcurveto{\pgfqpoint{1.593773in}{1.272552in}}{\pgfqpoint{1.593773in}{1.284368in}}%
\pgfpathlineto{\pgfqpoint{1.593773in}{1.318591in}}%
\pgfpathlineto{\pgfqpoint{1.586172in}{1.318591in}}%
\pgfpathlineto{\pgfqpoint{1.586172in}{1.326643in}}%
\pgfpathlineto{\pgfqpoint{1.593773in}{1.326643in}}%
\pgfpathlineto{\pgfqpoint{1.593773in}{1.344547in}}%
\pgfpathlineto{\pgfqpoint{1.604184in}{1.344547in}}%
\pgfpathlineto{\pgfqpoint{1.604184in}{1.344547in}}%
\pgfpathclose%
\pgfpathmoveto{\pgfqpoint{1.639127in}{1.326643in}}%
\pgfpathlineto{\pgfqpoint{1.649484in}{1.326643in}}%
\pgfpathlineto{\pgfqpoint{1.649484in}{1.263600in}}%
\pgfpathlineto{\pgfqpoint{1.639127in}{1.263600in}}%
\pgfpathlineto{\pgfqpoint{1.639127in}{1.326643in}}%
\pgfpathlineto{\pgfqpoint{1.639127in}{1.326643in}}%
\pgfpathclose%
\pgfpathmoveto{\pgfqpoint{1.639127in}{1.351193in}}%
\pgfpathlineto{\pgfqpoint{1.649484in}{1.351193in}}%
\pgfpathlineto{\pgfqpoint{1.649484in}{1.338062in}}%
\pgfpathlineto{\pgfqpoint{1.639127in}{1.338062in}}%
\pgfpathlineto{\pgfqpoint{1.639127in}{1.351193in}}%
\pgfpathlineto{\pgfqpoint{1.639127in}{1.351193in}}%
\pgfpathclose%
\pgfpathmoveto{\pgfqpoint{1.716526in}{1.324229in}}%
\pgfpathlineto{\pgfqpoint{1.716526in}{1.314538in}}%
\pgfpathquadraticcurveto{\pgfqpoint{1.712131in}{1.316970in}}{\pgfqpoint{1.707718in}{1.318177in}}%
\pgfpathquadraticcurveto{\pgfqpoint{1.703305in}{1.319384in}}{\pgfqpoint{1.698802in}{1.319384in}}%
\pgfpathquadraticcurveto{\pgfqpoint{1.688715in}{1.319384in}}{\pgfqpoint{1.683131in}{1.312989in}}%
\pgfpathquadraticcurveto{\pgfqpoint{1.677565in}{1.306613in}}{\pgfqpoint{1.677565in}{1.295067in}}%
\pgfpathquadraticcurveto{\pgfqpoint{1.677565in}{1.283521in}}{\pgfqpoint{1.683131in}{1.277127in}}%
\pgfpathquadraticcurveto{\pgfqpoint{1.688715in}{1.270751in}}{\pgfqpoint{1.698802in}{1.270751in}}%
\pgfpathquadraticcurveto{\pgfqpoint{1.703305in}{1.270751in}}{\pgfqpoint{1.707718in}{1.271958in}}%
\pgfpathquadraticcurveto{\pgfqpoint{1.712131in}{1.273164in}}{\pgfqpoint{1.716526in}{1.275596in}}%
\pgfpathlineto{\pgfqpoint{1.716526in}{1.266014in}}%
\pgfpathquadraticcurveto{\pgfqpoint{1.712185in}{1.263996in}}{\pgfqpoint{1.707538in}{1.262988in}}%
\pgfpathquadraticcurveto{\pgfqpoint{1.702908in}{1.261961in}}{\pgfqpoint{1.697667in}{1.261961in}}%
\pgfpathquadraticcurveto{\pgfqpoint{1.683419in}{1.261961in}}{\pgfqpoint{1.675026in}{1.270913in}}%
\pgfpathquadraticcurveto{\pgfqpoint{1.666650in}{1.279865in}}{\pgfqpoint{1.666650in}{1.295067in}}%
\pgfpathquadraticcurveto{\pgfqpoint{1.666650in}{1.310486in}}{\pgfqpoint{1.675116in}{1.319312in}}%
\pgfpathquadraticcurveto{\pgfqpoint{1.683599in}{1.328156in}}{\pgfqpoint{1.698351in}{1.328156in}}%
\pgfpathquadraticcurveto{\pgfqpoint{1.703125in}{1.328156in}}{\pgfqpoint{1.707682in}{1.327165in}}%
\pgfpathquadraticcurveto{\pgfqpoint{1.712239in}{1.326192in}}{\pgfqpoint{1.716526in}{1.324229in}}%
\pgfpathlineto{\pgfqpoint{1.716526in}{1.324229in}}%
\pgfpathclose%
\pgfusepath{fill}%
\end{pgfscope}%
\begin{pgfscope}%
\pgfsetbuttcap%
\pgfsetroundjoin%
\definecolor{currentfill}{rgb}{0.309804,0.372549,0.419608}%
\pgfsetfillcolor{currentfill}%
\pgfsetlinewidth{0.000000pt}%
\definecolor{currentstroke}{rgb}{0.309804,0.372549,0.419608}%
\pgfsetstrokecolor{currentstroke}%
\pgfsetdash{}{0pt}%
\pgfpathmoveto{\pgfqpoint{1.843916in}{1.324787in}}%
\pgfpathlineto{\pgfqpoint{1.843916in}{1.314989in}}%
\pgfpathquadraticcurveto{\pgfqpoint{1.839539in}{1.317240in}}{\pgfqpoint{1.834802in}{1.318357in}}%
\pgfpathquadraticcurveto{\pgfqpoint{1.830083in}{1.319492in}}{\pgfqpoint{1.825003in}{1.319492in}}%
\pgfpathquadraticcurveto{\pgfqpoint{1.817294in}{1.319492in}}{\pgfqpoint{1.813440in}{1.317132in}}%
\pgfpathquadraticcurveto{\pgfqpoint{1.809585in}{1.314773in}}{\pgfqpoint{1.809585in}{1.310035in}}%
\pgfpathquadraticcurveto{\pgfqpoint{1.809585in}{1.306433in}}{\pgfqpoint{1.812341in}{1.304380in}}%
\pgfpathquadraticcurveto{\pgfqpoint{1.815097in}{1.302326in}}{\pgfqpoint{1.823436in}{1.300471in}}%
\pgfpathlineto{\pgfqpoint{1.826985in}{1.299678in}}%
\pgfpathquadraticcurveto{\pgfqpoint{1.838008in}{1.297319in}}{\pgfqpoint{1.842655in}{1.293014in}}%
\pgfpathquadraticcurveto{\pgfqpoint{1.847303in}{1.288709in}}{\pgfqpoint{1.847303in}{1.281000in}}%
\pgfpathquadraticcurveto{\pgfqpoint{1.847303in}{1.272210in}}{\pgfqpoint{1.840350in}{1.267076in}}%
\pgfpathquadraticcurveto{\pgfqpoint{1.833397in}{1.261961in}}{\pgfqpoint{1.821239in}{1.261961in}}%
\pgfpathquadraticcurveto{\pgfqpoint{1.816178in}{1.261961in}}{\pgfqpoint{1.810684in}{1.262952in}}%
\pgfpathquadraticcurveto{\pgfqpoint{1.805190in}{1.263942in}}{\pgfqpoint{1.799120in}{1.265906in}}%
\pgfpathlineto{\pgfqpoint{1.799120in}{1.276605in}}%
\pgfpathquadraticcurveto{\pgfqpoint{1.804866in}{1.273615in}}{\pgfqpoint{1.810432in}{1.272120in}}%
\pgfpathquadraticcurveto{\pgfqpoint{1.815997in}{1.270643in}}{\pgfqpoint{1.821473in}{1.270643in}}%
\pgfpathquadraticcurveto{\pgfqpoint{1.828786in}{1.270643in}}{\pgfqpoint{1.832713in}{1.273146in}}%
\pgfpathquadraticcurveto{\pgfqpoint{1.836657in}{1.275650in}}{\pgfqpoint{1.836657in}{1.280207in}}%
\pgfpathquadraticcurveto{\pgfqpoint{1.836657in}{1.284422in}}{\pgfqpoint{1.833811in}{1.286674in}}%
\pgfpathquadraticcurveto{\pgfqpoint{1.830984in}{1.288925in}}{\pgfqpoint{1.821347in}{1.291014in}}%
\pgfpathlineto{\pgfqpoint{1.817745in}{1.291861in}}%
\pgfpathquadraticcurveto{\pgfqpoint{1.808126in}{1.293878in}}{\pgfqpoint{1.803839in}{1.298075in}}%
\pgfpathquadraticcurveto{\pgfqpoint{1.799570in}{1.302272in}}{\pgfqpoint{1.799570in}{1.309585in}}%
\pgfpathquadraticcurveto{\pgfqpoint{1.799570in}{1.318483in}}{\pgfqpoint{1.805875in}{1.323310in}}%
\pgfpathquadraticcurveto{\pgfqpoint{1.812179in}{1.328156in}}{\pgfqpoint{1.823779in}{1.328156in}}%
\pgfpathquadraticcurveto{\pgfqpoint{1.829507in}{1.328156in}}{\pgfqpoint{1.834568in}{1.327309in}}%
\pgfpathquadraticcurveto{\pgfqpoint{1.839647in}{1.326480in}}{\pgfqpoint{1.843916in}{1.324787in}}%
\pgfpathlineto{\pgfqpoint{1.843916in}{1.324787in}}%
\pgfpathclose%
\pgfpathmoveto{\pgfqpoint{1.873798in}{1.273056in}}%
\pgfpathlineto{\pgfqpoint{1.873798in}{1.239626in}}%
\pgfpathlineto{\pgfqpoint{1.863387in}{1.239626in}}%
\pgfpathlineto{\pgfqpoint{1.863387in}{1.326643in}}%
\pgfpathlineto{\pgfqpoint{1.873798in}{1.326643in}}%
\pgfpathlineto{\pgfqpoint{1.873798in}{1.317078in}}%
\pgfpathquadraticcurveto{\pgfqpoint{1.877077in}{1.322698in}}{\pgfqpoint{1.882048in}{1.325418in}}%
\pgfpathquadraticcurveto{\pgfqpoint{1.887037in}{1.328156in}}{\pgfqpoint{1.893954in}{1.328156in}}%
\pgfpathquadraticcurveto{\pgfqpoint{1.905446in}{1.328156in}}{\pgfqpoint{1.912615in}{1.319041in}}%
\pgfpathquadraticcurveto{\pgfqpoint{1.919801in}{1.309927in}}{\pgfqpoint{1.919801in}{1.295067in}}%
\pgfpathquadraticcurveto{\pgfqpoint{1.919801in}{1.280207in}}{\pgfqpoint{1.912615in}{1.271075in}}%
\pgfpathquadraticcurveto{\pgfqpoint{1.905446in}{1.261961in}}{\pgfqpoint{1.893954in}{1.261961in}}%
\pgfpathquadraticcurveto{\pgfqpoint{1.887037in}{1.261961in}}{\pgfqpoint{1.882048in}{1.264699in}}%
\pgfpathquadraticcurveto{\pgfqpoint{1.877077in}{1.267437in}}{\pgfqpoint{1.873798in}{1.273056in}}%
\pgfpathlineto{\pgfqpoint{1.873798in}{1.273056in}}%
\pgfpathclose%
\pgfpathmoveto{\pgfqpoint{1.909048in}{1.295067in}}%
\pgfpathquadraticcurveto{\pgfqpoint{1.909048in}{1.306487in}}{\pgfqpoint{1.904347in}{1.312989in}}%
\pgfpathquadraticcurveto{\pgfqpoint{1.899646in}{1.319492in}}{\pgfqpoint{1.891432in}{1.319492in}}%
\pgfpathquadraticcurveto{\pgfqpoint{1.883201in}{1.319492in}}{\pgfqpoint{1.878500in}{1.312989in}}%
\pgfpathquadraticcurveto{\pgfqpoint{1.873798in}{1.306487in}}{\pgfqpoint{1.873798in}{1.295067in}}%
\pgfpathquadraticcurveto{\pgfqpoint{1.873798in}{1.283648in}}{\pgfqpoint{1.878500in}{1.277145in}}%
\pgfpathquadraticcurveto{\pgfqpoint{1.883201in}{1.270643in}}{\pgfqpoint{1.891432in}{1.270643in}}%
\pgfpathquadraticcurveto{\pgfqpoint{1.899646in}{1.270643in}}{\pgfqpoint{1.904347in}{1.277145in}}%
\pgfpathquadraticcurveto{\pgfqpoint{1.909048in}{1.283648in}}{\pgfqpoint{1.909048in}{1.295067in}}%
\pgfpathlineto{\pgfqpoint{1.909048in}{1.295067in}}%
\pgfpathclose%
\pgfpathmoveto{\pgfqpoint{1.990895in}{1.297715in}}%
\pgfpathlineto{\pgfqpoint{1.990895in}{1.292654in}}%
\pgfpathlineto{\pgfqpoint{1.943271in}{1.292654in}}%
\pgfpathquadraticcurveto{\pgfqpoint{1.943956in}{1.281954in}}{\pgfqpoint{1.949720in}{1.276353in}}%
\pgfpathquadraticcurveto{\pgfqpoint{1.955484in}{1.270751in}}{\pgfqpoint{1.965786in}{1.270751in}}%
\pgfpathquadraticcurveto{\pgfqpoint{1.971748in}{1.270751in}}{\pgfqpoint{1.977350in}{1.272210in}}%
\pgfpathquadraticcurveto{\pgfqpoint{1.982952in}{1.273669in}}{\pgfqpoint{1.988482in}{1.276605in}}%
\pgfpathlineto{\pgfqpoint{1.988482in}{1.266806in}}%
\pgfpathquadraticcurveto{\pgfqpoint{1.982898in}{1.264447in}}{\pgfqpoint{1.977044in}{1.263204in}}%
\pgfpathquadraticcurveto{\pgfqpoint{1.971190in}{1.261961in}}{\pgfqpoint{1.965174in}{1.261961in}}%
\pgfpathquadraticcurveto{\pgfqpoint{1.950080in}{1.261961in}}{\pgfqpoint{1.941272in}{1.270733in}}%
\pgfpathquadraticcurveto{\pgfqpoint{1.932464in}{1.279523in}}{\pgfqpoint{1.932464in}{1.294509in}}%
\pgfpathquadraticcurveto{\pgfqpoint{1.932464in}{1.309981in}}{\pgfqpoint{1.940822in}{1.319059in}}%
\pgfpathquadraticcurveto{\pgfqpoint{1.949179in}{1.328156in}}{\pgfqpoint{1.963373in}{1.328156in}}%
\pgfpathquadraticcurveto{\pgfqpoint{1.976089in}{1.328156in}}{\pgfqpoint{1.983492in}{1.319960in}}%
\pgfpathquadraticcurveto{\pgfqpoint{1.990895in}{1.311783in}}{\pgfqpoint{1.990895in}{1.297715in}}%
\pgfpathlineto{\pgfqpoint{1.990895in}{1.297715in}}%
\pgfpathclose%
\pgfpathmoveto{\pgfqpoint{1.980538in}{1.300759in}}%
\pgfpathquadraticcurveto{\pgfqpoint{1.980430in}{1.309243in}}{\pgfqpoint{1.975783in}{1.314304in}}%
\pgfpathquadraticcurveto{\pgfqpoint{1.971136in}{1.319384in}}{\pgfqpoint{1.963481in}{1.319384in}}%
\pgfpathquadraticcurveto{\pgfqpoint{1.954817in}{1.319384in}}{\pgfqpoint{1.949612in}{1.314484in}}%
\pgfpathquadraticcurveto{\pgfqpoint{1.944406in}{1.309585in}}{\pgfqpoint{1.943613in}{1.300687in}}%
\pgfpathlineto{\pgfqpoint{1.980538in}{1.300759in}}%
\pgfpathlineto{\pgfqpoint{1.980538in}{1.300759in}}%
\pgfpathclose%
\pgfpathmoveto{\pgfqpoint{2.053271in}{1.324229in}}%
\pgfpathlineto{\pgfqpoint{2.053271in}{1.314538in}}%
\pgfpathquadraticcurveto{\pgfqpoint{2.048877in}{1.316970in}}{\pgfqpoint{2.044464in}{1.318177in}}%
\pgfpathquadraticcurveto{\pgfqpoint{2.040051in}{1.319384in}}{\pgfqpoint{2.035548in}{1.319384in}}%
\pgfpathquadraticcurveto{\pgfqpoint{2.025461in}{1.319384in}}{\pgfqpoint{2.019877in}{1.312989in}}%
\pgfpathquadraticcurveto{\pgfqpoint{2.014311in}{1.306613in}}{\pgfqpoint{2.014311in}{1.295067in}}%
\pgfpathquadraticcurveto{\pgfqpoint{2.014311in}{1.283521in}}{\pgfqpoint{2.019877in}{1.277127in}}%
\pgfpathquadraticcurveto{\pgfqpoint{2.025461in}{1.270751in}}{\pgfqpoint{2.035548in}{1.270751in}}%
\pgfpathquadraticcurveto{\pgfqpoint{2.040051in}{1.270751in}}{\pgfqpoint{2.044464in}{1.271958in}}%
\pgfpathquadraticcurveto{\pgfqpoint{2.048877in}{1.273164in}}{\pgfqpoint{2.053271in}{1.275596in}}%
\pgfpathlineto{\pgfqpoint{2.053271in}{1.266014in}}%
\pgfpathquadraticcurveto{\pgfqpoint{2.048931in}{1.263996in}}{\pgfqpoint{2.044283in}{1.262988in}}%
\pgfpathquadraticcurveto{\pgfqpoint{2.039654in}{1.261961in}}{\pgfqpoint{2.034413in}{1.261961in}}%
\pgfpathquadraticcurveto{\pgfqpoint{2.020165in}{1.261961in}}{\pgfqpoint{2.011771in}{1.270913in}}%
\pgfpathquadraticcurveto{\pgfqpoint{2.003396in}{1.279865in}}{\pgfqpoint{2.003396in}{1.295067in}}%
\pgfpathquadraticcurveto{\pgfqpoint{2.003396in}{1.310486in}}{\pgfqpoint{2.011862in}{1.319312in}}%
\pgfpathquadraticcurveto{\pgfqpoint{2.020345in}{1.328156in}}{\pgfqpoint{2.035097in}{1.328156in}}%
\pgfpathquadraticcurveto{\pgfqpoint{2.039870in}{1.328156in}}{\pgfqpoint{2.044428in}{1.327165in}}%
\pgfpathquadraticcurveto{\pgfqpoint{2.048985in}{1.326192in}}{\pgfqpoint{2.053271in}{1.324229in}}%
\pgfpathlineto{\pgfqpoint{2.053271in}{1.324229in}}%
\pgfpathclose%
\pgfpathmoveto{\pgfqpoint{2.071284in}{1.326643in}}%
\pgfpathlineto{\pgfqpoint{2.081641in}{1.326643in}}%
\pgfpathlineto{\pgfqpoint{2.081641in}{1.263600in}}%
\pgfpathlineto{\pgfqpoint{2.071284in}{1.263600in}}%
\pgfpathlineto{\pgfqpoint{2.071284in}{1.326643in}}%
\pgfpathlineto{\pgfqpoint{2.071284in}{1.326643in}}%
\pgfpathclose%
\pgfpathmoveto{\pgfqpoint{2.071284in}{1.351193in}}%
\pgfpathlineto{\pgfqpoint{2.081641in}{1.351193in}}%
\pgfpathlineto{\pgfqpoint{2.081641in}{1.338062in}}%
\pgfpathlineto{\pgfqpoint{2.071284in}{1.338062in}}%
\pgfpathlineto{\pgfqpoint{2.071284in}{1.351193in}}%
\pgfpathlineto{\pgfqpoint{2.071284in}{1.351193in}}%
\pgfpathclose%
\pgfpathmoveto{\pgfqpoint{2.154248in}{1.326643in}}%
\pgfpathlineto{\pgfqpoint{2.154248in}{1.263600in}}%
\pgfpathlineto{\pgfqpoint{2.143837in}{1.263600in}}%
\pgfpathlineto{\pgfqpoint{2.143837in}{1.318591in}}%
\pgfpathlineto{\pgfqpoint{2.115413in}{1.318591in}}%
\pgfpathlineto{\pgfqpoint{2.115413in}{1.263600in}}%
\pgfpathlineto{\pgfqpoint{2.105002in}{1.263600in}}%
\pgfpathlineto{\pgfqpoint{2.105002in}{1.318591in}}%
\pgfpathlineto{\pgfqpoint{2.095096in}{1.318591in}}%
\pgfpathlineto{\pgfqpoint{2.095096in}{1.326643in}}%
\pgfpathlineto{\pgfqpoint{2.105002in}{1.326643in}}%
\pgfpathlineto{\pgfqpoint{2.105002in}{1.331038in}}%
\pgfpathquadraticcurveto{\pgfqpoint{2.105002in}{1.341340in}}{\pgfqpoint{2.109866in}{1.346258in}}%
\pgfpathquadraticcurveto{\pgfqpoint{2.114747in}{1.351193in}}{\pgfqpoint{2.124816in}{1.351193in}}%
\pgfpathlineto{\pgfqpoint{2.135227in}{1.351193in}}%
\pgfpathlineto{\pgfqpoint{2.135227in}{1.342565in}}%
\pgfpathlineto{\pgfqpoint{2.125320in}{1.342565in}}%
\pgfpathquadraticcurveto{\pgfqpoint{2.119754in}{1.342565in}}{\pgfqpoint{2.117575in}{1.340314in}}%
\pgfpathquadraticcurveto{\pgfqpoint{2.115413in}{1.338062in}}{\pgfqpoint{2.115413in}{1.332208in}}%
\pgfpathlineto{\pgfqpoint{2.115413in}{1.326643in}}%
\pgfpathlineto{\pgfqpoint{2.154248in}{1.326643in}}%
\pgfpathlineto{\pgfqpoint{2.154248in}{1.326643in}}%
\pgfpathclose%
\pgfpathmoveto{\pgfqpoint{2.143837in}{1.351067in}}%
\pgfpathlineto{\pgfqpoint{2.154248in}{1.351067in}}%
\pgfpathlineto{\pgfqpoint{2.154248in}{1.337954in}}%
\pgfpathlineto{\pgfqpoint{2.143837in}{1.337954in}}%
\pgfpathlineto{\pgfqpoint{2.143837in}{1.351067in}}%
\pgfpathlineto{\pgfqpoint{2.143837in}{1.351067in}}%
\pgfpathclose%
\pgfpathmoveto{\pgfqpoint{2.221289in}{1.324229in}}%
\pgfpathlineto{\pgfqpoint{2.221289in}{1.314538in}}%
\pgfpathquadraticcurveto{\pgfqpoint{2.216894in}{1.316970in}}{\pgfqpoint{2.212481in}{1.318177in}}%
\pgfpathquadraticcurveto{\pgfqpoint{2.208068in}{1.319384in}}{\pgfqpoint{2.203565in}{1.319384in}}%
\pgfpathquadraticcurveto{\pgfqpoint{2.193478in}{1.319384in}}{\pgfqpoint{2.187894in}{1.312989in}}%
\pgfpathquadraticcurveto{\pgfqpoint{2.182329in}{1.306613in}}{\pgfqpoint{2.182329in}{1.295067in}}%
\pgfpathquadraticcurveto{\pgfqpoint{2.182329in}{1.283521in}}{\pgfqpoint{2.187894in}{1.277127in}}%
\pgfpathquadraticcurveto{\pgfqpoint{2.193478in}{1.270751in}}{\pgfqpoint{2.203565in}{1.270751in}}%
\pgfpathquadraticcurveto{\pgfqpoint{2.208068in}{1.270751in}}{\pgfqpoint{2.212481in}{1.271958in}}%
\pgfpathquadraticcurveto{\pgfqpoint{2.216894in}{1.273164in}}{\pgfqpoint{2.221289in}{1.275596in}}%
\pgfpathlineto{\pgfqpoint{2.221289in}{1.266014in}}%
\pgfpathquadraticcurveto{\pgfqpoint{2.216948in}{1.263996in}}{\pgfqpoint{2.212301in}{1.262988in}}%
\pgfpathquadraticcurveto{\pgfqpoint{2.207672in}{1.261961in}}{\pgfqpoint{2.202430in}{1.261961in}}%
\pgfpathquadraticcurveto{\pgfqpoint{2.188183in}{1.261961in}}{\pgfqpoint{2.179789in}{1.270913in}}%
\pgfpathquadraticcurveto{\pgfqpoint{2.171413in}{1.279865in}}{\pgfqpoint{2.171413in}{1.295067in}}%
\pgfpathquadraticcurveto{\pgfqpoint{2.171413in}{1.310486in}}{\pgfqpoint{2.179879in}{1.319312in}}%
\pgfpathquadraticcurveto{\pgfqpoint{2.188363in}{1.328156in}}{\pgfqpoint{2.203115in}{1.328156in}}%
\pgfpathquadraticcurveto{\pgfqpoint{2.207888in}{1.328156in}}{\pgfqpoint{2.212445in}{1.327165in}}%
\pgfpathquadraticcurveto{\pgfqpoint{2.217002in}{1.326192in}}{\pgfqpoint{2.221289in}{1.324229in}}%
\pgfpathlineto{\pgfqpoint{2.221289in}{1.324229in}}%
\pgfpathclose%
\pgfpathmoveto{\pgfqpoint{2.267958in}{1.295283in}}%
\pgfpathquadraticcurveto{\pgfqpoint{2.255404in}{1.295283in}}{\pgfqpoint{2.250559in}{1.292419in}}%
\pgfpathquadraticcurveto{\pgfqpoint{2.245713in}{1.289556in}}{\pgfqpoint{2.245713in}{1.282621in}}%
\pgfpathquadraticcurveto{\pgfqpoint{2.245713in}{1.277109in}}{\pgfqpoint{2.249352in}{1.273867in}}%
\pgfpathquadraticcurveto{\pgfqpoint{2.252990in}{1.270643in}}{\pgfqpoint{2.259222in}{1.270643in}}%
\pgfpathquadraticcurveto{\pgfqpoint{2.267850in}{1.270643in}}{\pgfqpoint{2.273056in}{1.276749in}}%
\pgfpathquadraticcurveto{\pgfqpoint{2.278261in}{1.282855in}}{\pgfqpoint{2.278261in}{1.292978in}}%
\pgfpathlineto{\pgfqpoint{2.278261in}{1.295283in}}%
\pgfpathlineto{\pgfqpoint{2.267958in}{1.295283in}}%
\pgfpathlineto{\pgfqpoint{2.267958in}{1.295283in}}%
\pgfpathclose%
\pgfpathmoveto{\pgfqpoint{2.288618in}{1.299570in}}%
\pgfpathlineto{\pgfqpoint{2.288618in}{1.263600in}}%
\pgfpathlineto{\pgfqpoint{2.278261in}{1.263600in}}%
\pgfpathlineto{\pgfqpoint{2.278261in}{1.273164in}}%
\pgfpathquadraticcurveto{\pgfqpoint{2.274713in}{1.267437in}}{\pgfqpoint{2.269417in}{1.264699in}}%
\pgfpathquadraticcurveto{\pgfqpoint{2.264122in}{1.261961in}}{\pgfqpoint{2.256467in}{1.261961in}}%
\pgfpathquadraticcurveto{\pgfqpoint{2.246794in}{1.261961in}}{\pgfqpoint{2.241066in}{1.267401in}}%
\pgfpathquadraticcurveto{\pgfqpoint{2.235356in}{1.272840in}}{\pgfqpoint{2.235356in}{1.281954in}}%
\pgfpathquadraticcurveto{\pgfqpoint{2.235356in}{1.292582in}}{\pgfqpoint{2.242471in}{1.297985in}}%
\pgfpathquadraticcurveto{\pgfqpoint{2.249604in}{1.303389in}}{\pgfqpoint{2.263725in}{1.303389in}}%
\pgfpathlineto{\pgfqpoint{2.278261in}{1.303389in}}%
\pgfpathlineto{\pgfqpoint{2.278261in}{1.304416in}}%
\pgfpathquadraticcurveto{\pgfqpoint{2.278261in}{1.311566in}}{\pgfqpoint{2.273560in}{1.315475in}}%
\pgfpathquadraticcurveto{\pgfqpoint{2.268859in}{1.319384in}}{\pgfqpoint{2.260357in}{1.319384in}}%
\pgfpathquadraticcurveto{\pgfqpoint{2.254954in}{1.319384in}}{\pgfqpoint{2.249820in}{1.318087in}}%
\pgfpathquadraticcurveto{\pgfqpoint{2.244705in}{1.316790in}}{\pgfqpoint{2.239985in}{1.314196in}}%
\pgfpathlineto{\pgfqpoint{2.239985in}{1.323779in}}%
\pgfpathquadraticcurveto{\pgfqpoint{2.245659in}{1.325976in}}{\pgfqpoint{2.251009in}{1.327057in}}%
\pgfpathquadraticcurveto{\pgfqpoint{2.256359in}{1.328156in}}{\pgfqpoint{2.261420in}{1.328156in}}%
\pgfpathquadraticcurveto{\pgfqpoint{2.275109in}{1.328156in}}{\pgfqpoint{2.281864in}{1.321059in}}%
\pgfpathquadraticcurveto{\pgfqpoint{2.288618in}{1.313980in}}{\pgfqpoint{2.288618in}{1.299570in}}%
\pgfpathlineto{\pgfqpoint{2.288618in}{1.299570in}}%
\pgfpathclose%
\pgfpathmoveto{\pgfqpoint{2.320194in}{1.344547in}}%
\pgfpathlineto{\pgfqpoint{2.320194in}{1.326643in}}%
\pgfpathlineto{\pgfqpoint{2.341520in}{1.326643in}}%
\pgfpathlineto{\pgfqpoint{2.341520in}{1.318591in}}%
\pgfpathlineto{\pgfqpoint{2.320194in}{1.318591in}}%
\pgfpathlineto{\pgfqpoint{2.320194in}{1.284368in}}%
\pgfpathquadraticcurveto{\pgfqpoint{2.320194in}{1.276659in}}{\pgfqpoint{2.322301in}{1.274461in}}%
\pgfpathquadraticcurveto{\pgfqpoint{2.324408in}{1.272264in}}{\pgfqpoint{2.330893in}{1.272264in}}%
\pgfpathlineto{\pgfqpoint{2.341520in}{1.272264in}}%
\pgfpathlineto{\pgfqpoint{2.341520in}{1.263600in}}%
\pgfpathlineto{\pgfqpoint{2.330893in}{1.263600in}}%
\pgfpathquadraticcurveto{\pgfqpoint{2.318897in}{1.263600in}}{\pgfqpoint{2.314340in}{1.268067in}}%
\pgfpathquadraticcurveto{\pgfqpoint{2.309783in}{1.272552in}}{\pgfqpoint{2.309783in}{1.284368in}}%
\pgfpathlineto{\pgfqpoint{2.309783in}{1.318591in}}%
\pgfpathlineto{\pgfqpoint{2.302181in}{1.318591in}}%
\pgfpathlineto{\pgfqpoint{2.302181in}{1.326643in}}%
\pgfpathlineto{\pgfqpoint{2.309783in}{1.326643in}}%
\pgfpathlineto{\pgfqpoint{2.309783in}{1.344547in}}%
\pgfpathlineto{\pgfqpoint{2.320194in}{1.344547in}}%
\pgfpathlineto{\pgfqpoint{2.320194in}{1.344547in}}%
\pgfpathclose%
\pgfpathmoveto{\pgfqpoint{2.355137in}{1.326643in}}%
\pgfpathlineto{\pgfqpoint{2.365494in}{1.326643in}}%
\pgfpathlineto{\pgfqpoint{2.365494in}{1.263600in}}%
\pgfpathlineto{\pgfqpoint{2.355137in}{1.263600in}}%
\pgfpathlineto{\pgfqpoint{2.355137in}{1.326643in}}%
\pgfpathlineto{\pgfqpoint{2.355137in}{1.326643in}}%
\pgfpathclose%
\pgfpathmoveto{\pgfqpoint{2.355137in}{1.351193in}}%
\pgfpathlineto{\pgfqpoint{2.365494in}{1.351193in}}%
\pgfpathlineto{\pgfqpoint{2.365494in}{1.338062in}}%
\pgfpathlineto{\pgfqpoint{2.355137in}{1.338062in}}%
\pgfpathlineto{\pgfqpoint{2.355137in}{1.351193in}}%
\pgfpathlineto{\pgfqpoint{2.355137in}{1.351193in}}%
\pgfpathclose%
\pgfpathmoveto{\pgfqpoint{2.411587in}{1.319384in}}%
\pgfpathquadraticcurveto{\pgfqpoint{2.403266in}{1.319384in}}{\pgfqpoint{2.398420in}{1.312881in}}%
\pgfpathquadraticcurveto{\pgfqpoint{2.393575in}{1.306379in}}{\pgfqpoint{2.393575in}{1.295067in}}%
\pgfpathquadraticcurveto{\pgfqpoint{2.393575in}{1.283756in}}{\pgfqpoint{2.398384in}{1.277253in}}%
\pgfpathquadraticcurveto{\pgfqpoint{2.403212in}{1.270751in}}{\pgfqpoint{2.411587in}{1.270751in}}%
\pgfpathquadraticcurveto{\pgfqpoint{2.419873in}{1.270751in}}{\pgfqpoint{2.424700in}{1.277271in}}%
\pgfpathquadraticcurveto{\pgfqpoint{2.429545in}{1.283810in}}{\pgfqpoint{2.429545in}{1.295067in}}%
\pgfpathquadraticcurveto{\pgfqpoint{2.429545in}{1.306271in}}{\pgfqpoint{2.424700in}{1.312827in}}%
\pgfpathquadraticcurveto{\pgfqpoint{2.419873in}{1.319384in}}{\pgfqpoint{2.411587in}{1.319384in}}%
\pgfpathlineto{\pgfqpoint{2.411587in}{1.319384in}}%
\pgfpathclose%
\pgfpathmoveto{\pgfqpoint{2.411587in}{1.328156in}}%
\pgfpathquadraticcurveto{\pgfqpoint{2.425096in}{1.328156in}}{\pgfqpoint{2.432806in}{1.319366in}}%
\pgfpathquadraticcurveto{\pgfqpoint{2.440533in}{1.310594in}}{\pgfqpoint{2.440533in}{1.295067in}}%
\pgfpathquadraticcurveto{\pgfqpoint{2.440533in}{1.279595in}}{\pgfqpoint{2.432806in}{1.270769in}}%
\pgfpathquadraticcurveto{\pgfqpoint{2.425096in}{1.261961in}}{\pgfqpoint{2.411587in}{1.261961in}}%
\pgfpathquadraticcurveto{\pgfqpoint{2.398024in}{1.261961in}}{\pgfqpoint{2.390333in}{1.270769in}}%
\pgfpathquadraticcurveto{\pgfqpoint{2.382660in}{1.279595in}}{\pgfqpoint{2.382660in}{1.295067in}}%
\pgfpathquadraticcurveto{\pgfqpoint{2.382660in}{1.310594in}}{\pgfqpoint{2.390333in}{1.319366in}}%
\pgfpathquadraticcurveto{\pgfqpoint{2.398024in}{1.328156in}}{\pgfqpoint{2.411587in}{1.328156in}}%
\pgfpathlineto{\pgfqpoint{2.411587in}{1.328156in}}%
\pgfpathclose%
\pgfpathmoveto{\pgfqpoint{2.510114in}{1.301660in}}%
\pgfpathlineto{\pgfqpoint{2.510114in}{1.263600in}}%
\pgfpathlineto{\pgfqpoint{2.499757in}{1.263600in}}%
\pgfpathlineto{\pgfqpoint{2.499757in}{1.301317in}}%
\pgfpathquadraticcurveto{\pgfqpoint{2.499757in}{1.310269in}}{\pgfqpoint{2.496262in}{1.314700in}}%
\pgfpathquadraticcurveto{\pgfqpoint{2.492768in}{1.319149in}}{\pgfqpoint{2.485797in}{1.319149in}}%
\pgfpathquadraticcurveto{\pgfqpoint{2.477404in}{1.319149in}}{\pgfqpoint{2.472558in}{1.313800in}}%
\pgfpathquadraticcurveto{\pgfqpoint{2.467713in}{1.308468in}}{\pgfqpoint{2.467713in}{1.299228in}}%
\pgfpathlineto{\pgfqpoint{2.467713in}{1.263600in}}%
\pgfpathlineto{\pgfqpoint{2.457302in}{1.263600in}}%
\pgfpathlineto{\pgfqpoint{2.457302in}{1.326643in}}%
\pgfpathlineto{\pgfqpoint{2.467713in}{1.326643in}}%
\pgfpathlineto{\pgfqpoint{2.467713in}{1.316844in}}%
\pgfpathquadraticcurveto{\pgfqpoint{2.471442in}{1.322536in}}{\pgfqpoint{2.476467in}{1.325346in}}%
\pgfpathquadraticcurveto{\pgfqpoint{2.481510in}{1.328156in}}{\pgfqpoint{2.488103in}{1.328156in}}%
\pgfpathquadraticcurveto{\pgfqpoint{2.498964in}{1.328156in}}{\pgfqpoint{2.504530in}{1.321437in}}%
\pgfpathquadraticcurveto{\pgfqpoint{2.510114in}{1.314718in}}{\pgfqpoint{2.510114in}{1.301660in}}%
\pgfpathlineto{\pgfqpoint{2.510114in}{1.301660in}}%
\pgfpathclose%
\pgfpathmoveto{\pgfqpoint{2.570941in}{1.324787in}}%
\pgfpathlineto{\pgfqpoint{2.570941in}{1.314989in}}%
\pgfpathquadraticcurveto{\pgfqpoint{2.566564in}{1.317240in}}{\pgfqpoint{2.561827in}{1.318357in}}%
\pgfpathquadraticcurveto{\pgfqpoint{2.557107in}{1.319492in}}{\pgfqpoint{2.552028in}{1.319492in}}%
\pgfpathquadraticcurveto{\pgfqpoint{2.544319in}{1.319492in}}{\pgfqpoint{2.540464in}{1.317132in}}%
\pgfpathquadraticcurveto{\pgfqpoint{2.536610in}{1.314773in}}{\pgfqpoint{2.536610in}{1.310035in}}%
\pgfpathquadraticcurveto{\pgfqpoint{2.536610in}{1.306433in}}{\pgfqpoint{2.539365in}{1.304380in}}%
\pgfpathquadraticcurveto{\pgfqpoint{2.542121in}{1.302326in}}{\pgfqpoint{2.550461in}{1.300471in}}%
\pgfpathlineto{\pgfqpoint{2.554009in}{1.299678in}}%
\pgfpathquadraticcurveto{\pgfqpoint{2.565033in}{1.297319in}}{\pgfqpoint{2.569680in}{1.293014in}}%
\pgfpathquadraticcurveto{\pgfqpoint{2.574327in}{1.288709in}}{\pgfqpoint{2.574327in}{1.281000in}}%
\pgfpathquadraticcurveto{\pgfqpoint{2.574327in}{1.272210in}}{\pgfqpoint{2.567374in}{1.267076in}}%
\pgfpathquadraticcurveto{\pgfqpoint{2.560422in}{1.261961in}}{\pgfqpoint{2.548263in}{1.261961in}}%
\pgfpathquadraticcurveto{\pgfqpoint{2.543202in}{1.261961in}}{\pgfqpoint{2.537708in}{1.262952in}}%
\pgfpathquadraticcurveto{\pgfqpoint{2.532215in}{1.263942in}}{\pgfqpoint{2.526145in}{1.265906in}}%
\pgfpathlineto{\pgfqpoint{2.526145in}{1.276605in}}%
\pgfpathquadraticcurveto{\pgfqpoint{2.531890in}{1.273615in}}{\pgfqpoint{2.537456in}{1.272120in}}%
\pgfpathquadraticcurveto{\pgfqpoint{2.543022in}{1.270643in}}{\pgfqpoint{2.548498in}{1.270643in}}%
\pgfpathquadraticcurveto{\pgfqpoint{2.555811in}{1.270643in}}{\pgfqpoint{2.559737in}{1.273146in}}%
\pgfpathquadraticcurveto{\pgfqpoint{2.563682in}{1.275650in}}{\pgfqpoint{2.563682in}{1.280207in}}%
\pgfpathquadraticcurveto{\pgfqpoint{2.563682in}{1.284422in}}{\pgfqpoint{2.560836in}{1.286674in}}%
\pgfpathquadraticcurveto{\pgfqpoint{2.558008in}{1.288925in}}{\pgfqpoint{2.548372in}{1.291014in}}%
\pgfpathlineto{\pgfqpoint{2.544769in}{1.291861in}}%
\pgfpathquadraticcurveto{\pgfqpoint{2.535151in}{1.293878in}}{\pgfqpoint{2.530864in}{1.298075in}}%
\pgfpathquadraticcurveto{\pgfqpoint{2.526595in}{1.302272in}}{\pgfqpoint{2.526595in}{1.309585in}}%
\pgfpathquadraticcurveto{\pgfqpoint{2.526595in}{1.318483in}}{\pgfqpoint{2.532899in}{1.323310in}}%
\pgfpathquadraticcurveto{\pgfqpoint{2.539203in}{1.328156in}}{\pgfqpoint{2.550803in}{1.328156in}}%
\pgfpathquadraticcurveto{\pgfqpoint{2.556531in}{1.328156in}}{\pgfqpoint{2.561592in}{1.327309in}}%
\pgfpathquadraticcurveto{\pgfqpoint{2.566672in}{1.326480in}}{\pgfqpoint{2.570941in}{1.324787in}}%
\pgfpathlineto{\pgfqpoint{2.570941in}{1.324787in}}%
\pgfpathclose%
\pgfpathmoveto{\pgfqpoint{2.592267in}{1.277902in}}%
\pgfpathlineto{\pgfqpoint{2.604155in}{1.277902in}}%
\pgfpathlineto{\pgfqpoint{2.604155in}{1.263600in}}%
\pgfpathlineto{\pgfqpoint{2.592267in}{1.263600in}}%
\pgfpathlineto{\pgfqpoint{2.592267in}{1.277902in}}%
\pgfpathlineto{\pgfqpoint{2.592267in}{1.277902in}}%
\pgfpathclose%
\pgfusepath{fill}%
\end{pgfscope}%
\begin{pgfscope}%
\pgfsetbuttcap%
\pgfsetroundjoin%
\definecolor{currentfill}{rgb}{0.309804,0.372549,0.419608}%
\pgfsetfillcolor{currentfill}%
\pgfsetlinewidth{0.000000pt}%
\definecolor{currentstroke}{rgb}{0.309804,0.372549,0.419608}%
\pgfsetstrokecolor{currentstroke}%
\pgfsetdash{}{0pt}%
\pgfpathmoveto{\pgfqpoint{2.770246in}{1.344889in}}%
\pgfpathlineto{\pgfqpoint{2.770246in}{1.333793in}}%
\pgfpathquadraticcurveto{\pgfqpoint{2.763780in}{1.336891in}}{\pgfqpoint{2.758034in}{1.338404in}}%
\pgfpathquadraticcurveto{\pgfqpoint{2.752288in}{1.339936in}}{\pgfqpoint{2.746939in}{1.339936in}}%
\pgfpathquadraticcurveto{\pgfqpoint{2.737662in}{1.339936in}}{\pgfqpoint{2.732619in}{1.336333in}}%
\pgfpathquadraticcurveto{\pgfqpoint{2.727576in}{1.332731in}}{\pgfqpoint{2.727576in}{1.326084in}}%
\pgfpathquadraticcurveto{\pgfqpoint{2.727576in}{1.320500in}}{\pgfqpoint{2.730926in}{1.317654in}}%
\pgfpathquadraticcurveto{\pgfqpoint{2.734276in}{1.314827in}}{\pgfqpoint{2.743624in}{1.313079in}}%
\pgfpathlineto{\pgfqpoint{2.750487in}{1.311674in}}%
\pgfpathquadraticcurveto{\pgfqpoint{2.763204in}{1.309243in}}{\pgfqpoint{2.769256in}{1.303137in}}%
\pgfpathquadraticcurveto{\pgfqpoint{2.775308in}{1.297031in}}{\pgfqpoint{2.775308in}{1.286800in}}%
\pgfpathquadraticcurveto{\pgfqpoint{2.775308in}{1.274569in}}{\pgfqpoint{2.767112in}{1.268265in}}%
\pgfpathquadraticcurveto{\pgfqpoint{2.758935in}{1.261961in}}{\pgfqpoint{2.743120in}{1.261961in}}%
\pgfpathquadraticcurveto{\pgfqpoint{2.737158in}{1.261961in}}{\pgfqpoint{2.730421in}{1.263312in}}%
\pgfpathquadraticcurveto{\pgfqpoint{2.723703in}{1.264663in}}{\pgfqpoint{2.716498in}{1.267311in}}%
\pgfpathlineto{\pgfqpoint{2.716498in}{1.279018in}}%
\pgfpathquadraticcurveto{\pgfqpoint{2.723415in}{1.275146in}}{\pgfqpoint{2.730061in}{1.273164in}}%
\pgfpathquadraticcurveto{\pgfqpoint{2.736708in}{1.271201in}}{\pgfqpoint{2.743120in}{1.271201in}}%
\pgfpathquadraticcurveto{\pgfqpoint{2.752847in}{1.271201in}}{\pgfqpoint{2.758142in}{1.275020in}}%
\pgfpathquadraticcurveto{\pgfqpoint{2.763438in}{1.278856in}}{\pgfqpoint{2.763438in}{1.285953in}}%
\pgfpathquadraticcurveto{\pgfqpoint{2.763438in}{1.292131in}}{\pgfqpoint{2.759637in}{1.295626in}}%
\pgfpathquadraticcurveto{\pgfqpoint{2.755837in}{1.299120in}}{\pgfqpoint{2.747173in}{1.300867in}}%
\pgfpathlineto{\pgfqpoint{2.740238in}{1.302218in}}%
\pgfpathquadraticcurveto{\pgfqpoint{2.727521in}{1.304740in}}{\pgfqpoint{2.721830in}{1.310143in}}%
\pgfpathquadraticcurveto{\pgfqpoint{2.716156in}{1.315547in}}{\pgfqpoint{2.716156in}{1.325184in}}%
\pgfpathquadraticcurveto{\pgfqpoint{2.716156in}{1.336333in}}{\pgfqpoint{2.724009in}{1.342745in}}%
\pgfpathquadraticcurveto{\pgfqpoint{2.731862in}{1.349158in}}{\pgfqpoint{2.745642in}{1.349158in}}%
\pgfpathquadraticcurveto{\pgfqpoint{2.751568in}{1.349158in}}{\pgfqpoint{2.757692in}{1.348077in}}%
\pgfpathquadraticcurveto{\pgfqpoint{2.763834in}{1.347014in}}{\pgfqpoint{2.770246in}{1.344889in}}%
\pgfpathlineto{\pgfqpoint{2.770246in}{1.344889in}}%
\pgfpathclose%
\pgfpathmoveto{\pgfqpoint{2.798796in}{1.295067in}}%
\pgfpathquadraticcurveto{\pgfqpoint{2.798796in}{1.283648in}}{\pgfqpoint{2.803497in}{1.277145in}}%
\pgfpathquadraticcurveto{\pgfqpoint{2.808198in}{1.270643in}}{\pgfqpoint{2.816411in}{1.270643in}}%
\pgfpathquadraticcurveto{\pgfqpoint{2.824625in}{1.270643in}}{\pgfqpoint{2.829344in}{1.277145in}}%
\pgfpathquadraticcurveto{\pgfqpoint{2.834081in}{1.283648in}}{\pgfqpoint{2.834081in}{1.295067in}}%
\pgfpathquadraticcurveto{\pgfqpoint{2.834081in}{1.306487in}}{\pgfqpoint{2.829344in}{1.312989in}}%
\pgfpathquadraticcurveto{\pgfqpoint{2.824625in}{1.319492in}}{\pgfqpoint{2.816411in}{1.319492in}}%
\pgfpathquadraticcurveto{\pgfqpoint{2.808198in}{1.319492in}}{\pgfqpoint{2.803497in}{1.312989in}}%
\pgfpathquadraticcurveto{\pgfqpoint{2.798796in}{1.306487in}}{\pgfqpoint{2.798796in}{1.295067in}}%
\pgfpathlineto{\pgfqpoint{2.798796in}{1.295067in}}%
\pgfpathclose%
\pgfpathmoveto{\pgfqpoint{2.834081in}{1.273056in}}%
\pgfpathquadraticcurveto{\pgfqpoint{2.830821in}{1.267437in}}{\pgfqpoint{2.825832in}{1.264699in}}%
\pgfpathquadraticcurveto{\pgfqpoint{2.820860in}{1.261961in}}{\pgfqpoint{2.813872in}{1.261961in}}%
\pgfpathquadraticcurveto{\pgfqpoint{2.802452in}{1.261961in}}{\pgfqpoint{2.795265in}{1.271075in}}%
\pgfpathquadraticcurveto{\pgfqpoint{2.788096in}{1.280207in}}{\pgfqpoint{2.788096in}{1.295067in}}%
\pgfpathquadraticcurveto{\pgfqpoint{2.788096in}{1.309927in}}{\pgfqpoint{2.795265in}{1.319041in}}%
\pgfpathquadraticcurveto{\pgfqpoint{2.802452in}{1.328156in}}{\pgfqpoint{2.813872in}{1.328156in}}%
\pgfpathquadraticcurveto{\pgfqpoint{2.820860in}{1.328156in}}{\pgfqpoint{2.825832in}{1.325418in}}%
\pgfpathquadraticcurveto{\pgfqpoint{2.830821in}{1.322698in}}{\pgfqpoint{2.834081in}{1.317078in}}%
\pgfpathlineto{\pgfqpoint{2.834081in}{1.326643in}}%
\pgfpathlineto{\pgfqpoint{2.844438in}{1.326643in}}%
\pgfpathlineto{\pgfqpoint{2.844438in}{1.239626in}}%
\pgfpathlineto{\pgfqpoint{2.834081in}{1.239626in}}%
\pgfpathlineto{\pgfqpoint{2.834081in}{1.273056in}}%
\pgfpathlineto{\pgfqpoint{2.834081in}{1.273056in}}%
\pgfpathclose%
\pgfpathmoveto{\pgfqpoint{2.864720in}{1.288475in}}%
\pgfpathlineto{\pgfqpoint{2.864720in}{1.326643in}}%
\pgfpathlineto{\pgfqpoint{2.875077in}{1.326643in}}%
\pgfpathlineto{\pgfqpoint{2.875077in}{1.288871in}}%
\pgfpathquadraticcurveto{\pgfqpoint{2.875077in}{1.279919in}}{\pgfqpoint{2.878553in}{1.275434in}}%
\pgfpathquadraticcurveto{\pgfqpoint{2.882048in}{1.270967in}}{\pgfqpoint{2.889036in}{1.270967in}}%
\pgfpathquadraticcurveto{\pgfqpoint{2.897412in}{1.270967in}}{\pgfqpoint{2.902275in}{1.276317in}}%
\pgfpathquadraticcurveto{\pgfqpoint{2.907157in}{1.281666in}}{\pgfqpoint{2.907157in}{1.290906in}}%
\pgfpathlineto{\pgfqpoint{2.907157in}{1.326643in}}%
\pgfpathlineto{\pgfqpoint{2.917514in}{1.326643in}}%
\pgfpathlineto{\pgfqpoint{2.917514in}{1.263600in}}%
\pgfpathlineto{\pgfqpoint{2.907157in}{1.263600in}}%
\pgfpathlineto{\pgfqpoint{2.907157in}{1.273291in}}%
\pgfpathquadraticcurveto{\pgfqpoint{2.903392in}{1.267545in}}{\pgfqpoint{2.898403in}{1.264753in}}%
\pgfpathquadraticcurveto{\pgfqpoint{2.893431in}{1.261961in}}{\pgfqpoint{2.886839in}{1.261961in}}%
\pgfpathquadraticcurveto{\pgfqpoint{2.875978in}{1.261961in}}{\pgfqpoint{2.870340in}{1.268715in}}%
\pgfpathquadraticcurveto{\pgfqpoint{2.864720in}{1.275470in}}{\pgfqpoint{2.864720in}{1.288475in}}%
\pgfpathlineto{\pgfqpoint{2.864720in}{1.288475in}}%
\pgfpathclose%
\pgfpathmoveto{\pgfqpoint{2.890784in}{1.328156in}}%
\pgfpathlineto{\pgfqpoint{2.890784in}{1.328156in}}%
\pgfpathlineto{\pgfqpoint{2.890784in}{1.328156in}}%
\pgfpathclose%
\pgfpathmoveto{\pgfqpoint{2.967497in}{1.295283in}}%
\pgfpathquadraticcurveto{\pgfqpoint{2.954943in}{1.295283in}}{\pgfqpoint{2.950098in}{1.292419in}}%
\pgfpathquadraticcurveto{\pgfqpoint{2.945252in}{1.289556in}}{\pgfqpoint{2.945252in}{1.282621in}}%
\pgfpathquadraticcurveto{\pgfqpoint{2.945252in}{1.277109in}}{\pgfqpoint{2.948891in}{1.273867in}}%
\pgfpathquadraticcurveto{\pgfqpoint{2.952529in}{1.270643in}}{\pgfqpoint{2.958762in}{1.270643in}}%
\pgfpathquadraticcurveto{\pgfqpoint{2.967389in}{1.270643in}}{\pgfqpoint{2.972595in}{1.276749in}}%
\pgfpathquadraticcurveto{\pgfqpoint{2.977800in}{1.282855in}}{\pgfqpoint{2.977800in}{1.292978in}}%
\pgfpathlineto{\pgfqpoint{2.977800in}{1.295283in}}%
\pgfpathlineto{\pgfqpoint{2.967497in}{1.295283in}}%
\pgfpathlineto{\pgfqpoint{2.967497in}{1.295283in}}%
\pgfpathclose%
\pgfpathmoveto{\pgfqpoint{2.988157in}{1.299570in}}%
\pgfpathlineto{\pgfqpoint{2.988157in}{1.263600in}}%
\pgfpathlineto{\pgfqpoint{2.977800in}{1.263600in}}%
\pgfpathlineto{\pgfqpoint{2.977800in}{1.273164in}}%
\pgfpathquadraticcurveto{\pgfqpoint{2.974252in}{1.267437in}}{\pgfqpoint{2.968956in}{1.264699in}}%
\pgfpathquadraticcurveto{\pgfqpoint{2.963661in}{1.261961in}}{\pgfqpoint{2.956006in}{1.261961in}}%
\pgfpathquadraticcurveto{\pgfqpoint{2.946333in}{1.261961in}}{\pgfqpoint{2.940605in}{1.267401in}}%
\pgfpathquadraticcurveto{\pgfqpoint{2.934895in}{1.272840in}}{\pgfqpoint{2.934895in}{1.281954in}}%
\pgfpathquadraticcurveto{\pgfqpoint{2.934895in}{1.292582in}}{\pgfqpoint{2.942010in}{1.297985in}}%
\pgfpathquadraticcurveto{\pgfqpoint{2.949143in}{1.303389in}}{\pgfqpoint{2.963265in}{1.303389in}}%
\pgfpathlineto{\pgfqpoint{2.977800in}{1.303389in}}%
\pgfpathlineto{\pgfqpoint{2.977800in}{1.304416in}}%
\pgfpathquadraticcurveto{\pgfqpoint{2.977800in}{1.311566in}}{\pgfqpoint{2.973099in}{1.315475in}}%
\pgfpathquadraticcurveto{\pgfqpoint{2.968398in}{1.319384in}}{\pgfqpoint{2.959896in}{1.319384in}}%
\pgfpathquadraticcurveto{\pgfqpoint{2.954493in}{1.319384in}}{\pgfqpoint{2.949359in}{1.318087in}}%
\pgfpathquadraticcurveto{\pgfqpoint{2.944244in}{1.316790in}}{\pgfqpoint{2.939525in}{1.314196in}}%
\pgfpathlineto{\pgfqpoint{2.939525in}{1.323779in}}%
\pgfpathquadraticcurveto{\pgfqpoint{2.945198in}{1.325976in}}{\pgfqpoint{2.950548in}{1.327057in}}%
\pgfpathquadraticcurveto{\pgfqpoint{2.955898in}{1.328156in}}{\pgfqpoint{2.960959in}{1.328156in}}%
\pgfpathquadraticcurveto{\pgfqpoint{2.974648in}{1.328156in}}{\pgfqpoint{2.981403in}{1.321059in}}%
\pgfpathquadraticcurveto{\pgfqpoint{2.988157in}{1.313980in}}{\pgfqpoint{2.988157in}{1.299570in}}%
\pgfpathlineto{\pgfqpoint{2.988157in}{1.299570in}}%
\pgfpathclose%
\pgfpathmoveto{\pgfqpoint{3.046012in}{1.316970in}}%
\pgfpathquadraticcurveto{\pgfqpoint{3.044265in}{1.317979in}}{\pgfqpoint{3.042212in}{1.318447in}}%
\pgfpathquadraticcurveto{\pgfqpoint{3.040158in}{1.318933in}}{\pgfqpoint{3.037691in}{1.318933in}}%
\pgfpathquadraticcurveto{\pgfqpoint{3.028901in}{1.318933in}}{\pgfqpoint{3.024200in}{1.313223in}}%
\pgfpathquadraticcurveto{\pgfqpoint{3.019498in}{1.307514in}}{\pgfqpoint{3.019498in}{1.296814in}}%
\pgfpathlineto{\pgfqpoint{3.019498in}{1.263600in}}%
\pgfpathlineto{\pgfqpoint{3.009087in}{1.263600in}}%
\pgfpathlineto{\pgfqpoint{3.009087in}{1.326643in}}%
\pgfpathlineto{\pgfqpoint{3.019498in}{1.326643in}}%
\pgfpathlineto{\pgfqpoint{3.019498in}{1.316844in}}%
\pgfpathquadraticcurveto{\pgfqpoint{3.022777in}{1.322590in}}{\pgfqpoint{3.028000in}{1.325364in}}%
\pgfpathquadraticcurveto{\pgfqpoint{3.033242in}{1.328156in}}{\pgfqpoint{3.040735in}{1.328156in}}%
\pgfpathquadraticcurveto{\pgfqpoint{3.041798in}{1.328156in}}{\pgfqpoint{3.043094in}{1.328011in}}%
\pgfpathquadraticcurveto{\pgfqpoint{3.044391in}{1.327885in}}{\pgfqpoint{3.045958in}{1.327597in}}%
\pgfpathlineto{\pgfqpoint{3.046012in}{1.316970in}}%
\pgfpathlineto{\pgfqpoint{3.046012in}{1.316970in}}%
\pgfpathclose%
\pgfpathmoveto{\pgfqpoint{3.108262in}{1.297715in}}%
\pgfpathlineto{\pgfqpoint{3.108262in}{1.292654in}}%
\pgfpathlineto{\pgfqpoint{3.060638in}{1.292654in}}%
\pgfpathquadraticcurveto{\pgfqpoint{3.061323in}{1.281954in}}{\pgfqpoint{3.067087in}{1.276353in}}%
\pgfpathquadraticcurveto{\pgfqpoint{3.072850in}{1.270751in}}{\pgfqpoint{3.083153in}{1.270751in}}%
\pgfpathquadraticcurveto{\pgfqpoint{3.089115in}{1.270751in}}{\pgfqpoint{3.094717in}{1.272210in}}%
\pgfpathquadraticcurveto{\pgfqpoint{3.100319in}{1.273669in}}{\pgfqpoint{3.105849in}{1.276605in}}%
\pgfpathlineto{\pgfqpoint{3.105849in}{1.266806in}}%
\pgfpathquadraticcurveto{\pgfqpoint{3.100265in}{1.264447in}}{\pgfqpoint{3.094411in}{1.263204in}}%
\pgfpathquadraticcurveto{\pgfqpoint{3.088557in}{1.261961in}}{\pgfqpoint{3.082541in}{1.261961in}}%
\pgfpathquadraticcurveto{\pgfqpoint{3.067447in}{1.261961in}}{\pgfqpoint{3.058639in}{1.270733in}}%
\pgfpathquadraticcurveto{\pgfqpoint{3.049831in}{1.279523in}}{\pgfqpoint{3.049831in}{1.294509in}}%
\pgfpathquadraticcurveto{\pgfqpoint{3.049831in}{1.309981in}}{\pgfqpoint{3.058189in}{1.319059in}}%
\pgfpathquadraticcurveto{\pgfqpoint{3.066546in}{1.328156in}}{\pgfqpoint{3.080740in}{1.328156in}}%
\pgfpathquadraticcurveto{\pgfqpoint{3.093456in}{1.328156in}}{\pgfqpoint{3.100859in}{1.319960in}}%
\pgfpathquadraticcurveto{\pgfqpoint{3.108262in}{1.311783in}}{\pgfqpoint{3.108262in}{1.297715in}}%
\pgfpathlineto{\pgfqpoint{3.108262in}{1.297715in}}%
\pgfpathclose%
\pgfpathmoveto{\pgfqpoint{3.097905in}{1.300759in}}%
\pgfpathquadraticcurveto{\pgfqpoint{3.097797in}{1.309243in}}{\pgfqpoint{3.093150in}{1.314304in}}%
\pgfpathquadraticcurveto{\pgfqpoint{3.088503in}{1.319384in}}{\pgfqpoint{3.080848in}{1.319384in}}%
\pgfpathquadraticcurveto{\pgfqpoint{3.072184in}{1.319384in}}{\pgfqpoint{3.066979in}{1.314484in}}%
\pgfpathquadraticcurveto{\pgfqpoint{3.061773in}{1.309585in}}{\pgfqpoint{3.060980in}{1.300687in}}%
\pgfpathlineto{\pgfqpoint{3.097905in}{1.300759in}}%
\pgfpathlineto{\pgfqpoint{3.097905in}{1.300759in}}%
\pgfpathclose%
\pgfpathmoveto{\pgfqpoint{3.165451in}{1.324787in}}%
\pgfpathlineto{\pgfqpoint{3.165451in}{1.314989in}}%
\pgfpathquadraticcurveto{\pgfqpoint{3.161074in}{1.317240in}}{\pgfqpoint{3.156337in}{1.318357in}}%
\pgfpathquadraticcurveto{\pgfqpoint{3.151618in}{1.319492in}}{\pgfqpoint{3.146538in}{1.319492in}}%
\pgfpathquadraticcurveto{\pgfqpoint{3.138829in}{1.319492in}}{\pgfqpoint{3.134974in}{1.317132in}}%
\pgfpathquadraticcurveto{\pgfqpoint{3.131120in}{1.314773in}}{\pgfqpoint{3.131120in}{1.310035in}}%
\pgfpathquadraticcurveto{\pgfqpoint{3.131120in}{1.306433in}}{\pgfqpoint{3.133876in}{1.304380in}}%
\pgfpathquadraticcurveto{\pgfqpoint{3.136632in}{1.302326in}}{\pgfqpoint{3.144971in}{1.300471in}}%
\pgfpathlineto{\pgfqpoint{3.148520in}{1.299678in}}%
\pgfpathquadraticcurveto{\pgfqpoint{3.159543in}{1.297319in}}{\pgfqpoint{3.164190in}{1.293014in}}%
\pgfpathquadraticcurveto{\pgfqpoint{3.168837in}{1.288709in}}{\pgfqpoint{3.168837in}{1.281000in}}%
\pgfpathquadraticcurveto{\pgfqpoint{3.168837in}{1.272210in}}{\pgfqpoint{3.161885in}{1.267076in}}%
\pgfpathquadraticcurveto{\pgfqpoint{3.154932in}{1.261961in}}{\pgfqpoint{3.142774in}{1.261961in}}%
\pgfpathquadraticcurveto{\pgfqpoint{3.137712in}{1.261961in}}{\pgfqpoint{3.132219in}{1.262952in}}%
\pgfpathquadraticcurveto{\pgfqpoint{3.126725in}{1.263942in}}{\pgfqpoint{3.120655in}{1.265906in}}%
\pgfpathlineto{\pgfqpoint{3.120655in}{1.276605in}}%
\pgfpathquadraticcurveto{\pgfqpoint{3.126401in}{1.273615in}}{\pgfqpoint{3.131966in}{1.272120in}}%
\pgfpathquadraticcurveto{\pgfqpoint{3.137532in}{1.270643in}}{\pgfqpoint{3.143008in}{1.270643in}}%
\pgfpathquadraticcurveto{\pgfqpoint{3.150321in}{1.270643in}}{\pgfqpoint{3.154247in}{1.273146in}}%
\pgfpathquadraticcurveto{\pgfqpoint{3.158192in}{1.275650in}}{\pgfqpoint{3.158192in}{1.280207in}}%
\pgfpathquadraticcurveto{\pgfqpoint{3.158192in}{1.284422in}}{\pgfqpoint{3.155346in}{1.286674in}}%
\pgfpathquadraticcurveto{\pgfqpoint{3.152518in}{1.288925in}}{\pgfqpoint{3.142882in}{1.291014in}}%
\pgfpathlineto{\pgfqpoint{3.139279in}{1.291861in}}%
\pgfpathquadraticcurveto{\pgfqpoint{3.129661in}{1.293878in}}{\pgfqpoint{3.125374in}{1.298075in}}%
\pgfpathquadraticcurveto{\pgfqpoint{3.121105in}{1.302272in}}{\pgfqpoint{3.121105in}{1.309585in}}%
\pgfpathquadraticcurveto{\pgfqpoint{3.121105in}{1.318483in}}{\pgfqpoint{3.127409in}{1.323310in}}%
\pgfpathquadraticcurveto{\pgfqpoint{3.133714in}{1.328156in}}{\pgfqpoint{3.145313in}{1.328156in}}%
\pgfpathquadraticcurveto{\pgfqpoint{3.151041in}{1.328156in}}{\pgfqpoint{3.156103in}{1.327309in}}%
\pgfpathquadraticcurveto{\pgfqpoint{3.161182in}{1.326480in}}{\pgfqpoint{3.165451in}{1.324787in}}%
\pgfpathlineto{\pgfqpoint{3.165451in}{1.324787in}}%
\pgfpathclose%
\pgfpathmoveto{\pgfqpoint{3.187966in}{1.277902in}}%
\pgfpathlineto{\pgfqpoint{3.199836in}{1.277902in}}%
\pgfpathlineto{\pgfqpoint{3.199836in}{1.263600in}}%
\pgfpathlineto{\pgfqpoint{3.187966in}{1.263600in}}%
\pgfpathlineto{\pgfqpoint{3.187966in}{1.277902in}}%
\pgfpathlineto{\pgfqpoint{3.187966in}{1.277902in}}%
\pgfpathclose%
\pgfpathmoveto{\pgfqpoint{3.187966in}{1.323202in}}%
\pgfpathlineto{\pgfqpoint{3.199836in}{1.323202in}}%
\pgfpathlineto{\pgfqpoint{3.199836in}{1.308919in}}%
\pgfpathlineto{\pgfqpoint{3.187966in}{1.308919in}}%
\pgfpathlineto{\pgfqpoint{3.187966in}{1.323202in}}%
\pgfpathlineto{\pgfqpoint{3.187966in}{1.323202in}}%
\pgfpathclose%
\pgfusepath{fill}%
\end{pgfscope}%
\begin{pgfscope}%
\pgfsetbuttcap%
\pgfsetroundjoin%
\definecolor{currentfill}{rgb}{0.309804,0.372549,0.419608}%
\pgfsetfillcolor{currentfill}%
\pgfsetlinewidth{0.000000pt}%
\definecolor{currentstroke}{rgb}{0.309804,0.372549,0.419608}%
\pgfsetstrokecolor{currentstroke}%
\pgfsetdash{}{0pt}%
\pgfpathmoveto{\pgfqpoint{3.356529in}{1.297715in}}%
\pgfpathlineto{\pgfqpoint{3.356529in}{1.292654in}}%
\pgfpathlineto{\pgfqpoint{3.308905in}{1.292654in}}%
\pgfpathquadraticcurveto{\pgfqpoint{3.309590in}{1.281954in}}{\pgfqpoint{3.315354in}{1.276353in}}%
\pgfpathquadraticcurveto{\pgfqpoint{3.321117in}{1.270751in}}{\pgfqpoint{3.331420in}{1.270751in}}%
\pgfpathquadraticcurveto{\pgfqpoint{3.337382in}{1.270751in}}{\pgfqpoint{3.342984in}{1.272210in}}%
\pgfpathquadraticcurveto{\pgfqpoint{3.348586in}{1.273669in}}{\pgfqpoint{3.354116in}{1.276605in}}%
\pgfpathlineto{\pgfqpoint{3.354116in}{1.266806in}}%
\pgfpathquadraticcurveto{\pgfqpoint{3.348532in}{1.264447in}}{\pgfqpoint{3.342678in}{1.263204in}}%
\pgfpathquadraticcurveto{\pgfqpoint{3.336824in}{1.261961in}}{\pgfqpoint{3.330808in}{1.261961in}}%
\pgfpathquadraticcurveto{\pgfqpoint{3.315714in}{1.261961in}}{\pgfqpoint{3.306906in}{1.270733in}}%
\pgfpathquadraticcurveto{\pgfqpoint{3.298098in}{1.279523in}}{\pgfqpoint{3.298098in}{1.294509in}}%
\pgfpathquadraticcurveto{\pgfqpoint{3.298098in}{1.309981in}}{\pgfqpoint{3.306456in}{1.319059in}}%
\pgfpathquadraticcurveto{\pgfqpoint{3.314813in}{1.328156in}}{\pgfqpoint{3.329007in}{1.328156in}}%
\pgfpathquadraticcurveto{\pgfqpoint{3.341723in}{1.328156in}}{\pgfqpoint{3.349126in}{1.319960in}}%
\pgfpathquadraticcurveto{\pgfqpoint{3.356529in}{1.311783in}}{\pgfqpoint{3.356529in}{1.297715in}}%
\pgfpathlineto{\pgfqpoint{3.356529in}{1.297715in}}%
\pgfpathclose%
\pgfpathmoveto{\pgfqpoint{3.346172in}{1.300759in}}%
\pgfpathquadraticcurveto{\pgfqpoint{3.346064in}{1.309243in}}{\pgfqpoint{3.341417in}{1.314304in}}%
\pgfpathquadraticcurveto{\pgfqpoint{3.336770in}{1.319384in}}{\pgfqpoint{3.329115in}{1.319384in}}%
\pgfpathquadraticcurveto{\pgfqpoint{3.320451in}{1.319384in}}{\pgfqpoint{3.315245in}{1.314484in}}%
\pgfpathquadraticcurveto{\pgfqpoint{3.310040in}{1.309585in}}{\pgfqpoint{3.309247in}{1.300687in}}%
\pgfpathlineto{\pgfqpoint{3.346172in}{1.300759in}}%
\pgfpathlineto{\pgfqpoint{3.346172in}{1.300759in}}%
\pgfpathclose%
\pgfpathmoveto{\pgfqpoint{3.366112in}{1.326643in}}%
\pgfpathlineto{\pgfqpoint{3.377081in}{1.326643in}}%
\pgfpathlineto{\pgfqpoint{3.396786in}{1.273741in}}%
\pgfpathlineto{\pgfqpoint{3.416492in}{1.326643in}}%
\pgfpathlineto{\pgfqpoint{3.427461in}{1.326643in}}%
\pgfpathlineto{\pgfqpoint{3.403811in}{1.263600in}}%
\pgfpathlineto{\pgfqpoint{3.389744in}{1.263600in}}%
\pgfpathlineto{\pgfqpoint{3.366112in}{1.326643in}}%
\pgfpathlineto{\pgfqpoint{3.366112in}{1.326643in}}%
\pgfpathclose%
\pgfpathmoveto{\pgfqpoint{3.495691in}{1.297715in}}%
\pgfpathlineto{\pgfqpoint{3.495691in}{1.292654in}}%
\pgfpathlineto{\pgfqpoint{3.448067in}{1.292654in}}%
\pgfpathquadraticcurveto{\pgfqpoint{3.448752in}{1.281954in}}{\pgfqpoint{3.454515in}{1.276353in}}%
\pgfpathquadraticcurveto{\pgfqpoint{3.460279in}{1.270751in}}{\pgfqpoint{3.470582in}{1.270751in}}%
\pgfpathquadraticcurveto{\pgfqpoint{3.476544in}{1.270751in}}{\pgfqpoint{3.482146in}{1.272210in}}%
\pgfpathquadraticcurveto{\pgfqpoint{3.487748in}{1.273669in}}{\pgfqpoint{3.493278in}{1.276605in}}%
\pgfpathlineto{\pgfqpoint{3.493278in}{1.266806in}}%
\pgfpathquadraticcurveto{\pgfqpoint{3.487694in}{1.264447in}}{\pgfqpoint{3.481840in}{1.263204in}}%
\pgfpathquadraticcurveto{\pgfqpoint{3.475986in}{1.261961in}}{\pgfqpoint{3.469970in}{1.261961in}}%
\pgfpathquadraticcurveto{\pgfqpoint{3.454876in}{1.261961in}}{\pgfqpoint{3.446068in}{1.270733in}}%
\pgfpathquadraticcurveto{\pgfqpoint{3.437260in}{1.279523in}}{\pgfqpoint{3.437260in}{1.294509in}}%
\pgfpathquadraticcurveto{\pgfqpoint{3.437260in}{1.309981in}}{\pgfqpoint{3.445617in}{1.319059in}}%
\pgfpathquadraticcurveto{\pgfqpoint{3.453975in}{1.328156in}}{\pgfqpoint{3.468169in}{1.328156in}}%
\pgfpathquadraticcurveto{\pgfqpoint{3.480885in}{1.328156in}}{\pgfqpoint{3.488288in}{1.319960in}}%
\pgfpathquadraticcurveto{\pgfqpoint{3.495691in}{1.311783in}}{\pgfqpoint{3.495691in}{1.297715in}}%
\pgfpathlineto{\pgfqpoint{3.495691in}{1.297715in}}%
\pgfpathclose%
\pgfpathmoveto{\pgfqpoint{3.485334in}{1.300759in}}%
\pgfpathquadraticcurveto{\pgfqpoint{3.485226in}{1.309243in}}{\pgfqpoint{3.480579in}{1.314304in}}%
\pgfpathquadraticcurveto{\pgfqpoint{3.475932in}{1.319384in}}{\pgfqpoint{3.468277in}{1.319384in}}%
\pgfpathquadraticcurveto{\pgfqpoint{3.459613in}{1.319384in}}{\pgfqpoint{3.454407in}{1.314484in}}%
\pgfpathquadraticcurveto{\pgfqpoint{3.449202in}{1.309585in}}{\pgfqpoint{3.448409in}{1.300687in}}%
\pgfpathlineto{\pgfqpoint{3.485334in}{1.300759in}}%
\pgfpathlineto{\pgfqpoint{3.485334in}{1.300759in}}%
\pgfpathclose%
\pgfpathmoveto{\pgfqpoint{3.565110in}{1.301660in}}%
\pgfpathlineto{\pgfqpoint{3.565110in}{1.263600in}}%
\pgfpathlineto{\pgfqpoint{3.554753in}{1.263600in}}%
\pgfpathlineto{\pgfqpoint{3.554753in}{1.301317in}}%
\pgfpathquadraticcurveto{\pgfqpoint{3.554753in}{1.310269in}}{\pgfqpoint{3.551259in}{1.314700in}}%
\pgfpathquadraticcurveto{\pgfqpoint{3.547764in}{1.319149in}}{\pgfqpoint{3.540794in}{1.319149in}}%
\pgfpathquadraticcurveto{\pgfqpoint{3.532400in}{1.319149in}}{\pgfqpoint{3.527555in}{1.313800in}}%
\pgfpathquadraticcurveto{\pgfqpoint{3.522709in}{1.308468in}}{\pgfqpoint{3.522709in}{1.299228in}}%
\pgfpathlineto{\pgfqpoint{3.522709in}{1.263600in}}%
\pgfpathlineto{\pgfqpoint{3.512298in}{1.263600in}}%
\pgfpathlineto{\pgfqpoint{3.512298in}{1.326643in}}%
\pgfpathlineto{\pgfqpoint{3.522709in}{1.326643in}}%
\pgfpathlineto{\pgfqpoint{3.522709in}{1.316844in}}%
\pgfpathquadraticcurveto{\pgfqpoint{3.526438in}{1.322536in}}{\pgfqpoint{3.531463in}{1.325346in}}%
\pgfpathquadraticcurveto{\pgfqpoint{3.536507in}{1.328156in}}{\pgfqpoint{3.543099in}{1.328156in}}%
\pgfpathquadraticcurveto{\pgfqpoint{3.553961in}{1.328156in}}{\pgfqpoint{3.559526in}{1.321437in}}%
\pgfpathquadraticcurveto{\pgfqpoint{3.565110in}{1.314718in}}{\pgfqpoint{3.565110in}{1.301660in}}%
\pgfpathlineto{\pgfqpoint{3.565110in}{1.301660in}}%
\pgfpathclose%
\pgfpathmoveto{\pgfqpoint{3.596001in}{1.344547in}}%
\pgfpathlineto{\pgfqpoint{3.596001in}{1.326643in}}%
\pgfpathlineto{\pgfqpoint{3.617327in}{1.326643in}}%
\pgfpathlineto{\pgfqpoint{3.617327in}{1.318591in}}%
\pgfpathlineto{\pgfqpoint{3.596001in}{1.318591in}}%
\pgfpathlineto{\pgfqpoint{3.596001in}{1.284368in}}%
\pgfpathquadraticcurveto{\pgfqpoint{3.596001in}{1.276659in}}{\pgfqpoint{3.598108in}{1.274461in}}%
\pgfpathquadraticcurveto{\pgfqpoint{3.600216in}{1.272264in}}{\pgfqpoint{3.606700in}{1.272264in}}%
\pgfpathlineto{\pgfqpoint{3.617327in}{1.272264in}}%
\pgfpathlineto{\pgfqpoint{3.617327in}{1.263600in}}%
\pgfpathlineto{\pgfqpoint{3.606700in}{1.263600in}}%
\pgfpathquadraticcurveto{\pgfqpoint{3.594704in}{1.263600in}}{\pgfqpoint{3.590147in}{1.268067in}}%
\pgfpathquadraticcurveto{\pgfqpoint{3.585590in}{1.272552in}}{\pgfqpoint{3.585590in}{1.284368in}}%
\pgfpathlineto{\pgfqpoint{3.585590in}{1.318591in}}%
\pgfpathlineto{\pgfqpoint{3.577989in}{1.318591in}}%
\pgfpathlineto{\pgfqpoint{3.577989in}{1.326643in}}%
\pgfpathlineto{\pgfqpoint{3.585590in}{1.326643in}}%
\pgfpathlineto{\pgfqpoint{3.585590in}{1.344547in}}%
\pgfpathlineto{\pgfqpoint{3.596001in}{1.344547in}}%
\pgfpathlineto{\pgfqpoint{3.596001in}{1.344547in}}%
\pgfpathclose%
\pgfpathmoveto{\pgfqpoint{3.625721in}{1.299786in}}%
\pgfpathlineto{\pgfqpoint{3.656053in}{1.299786in}}%
\pgfpathlineto{\pgfqpoint{3.656053in}{1.290564in}}%
\pgfpathlineto{\pgfqpoint{3.625721in}{1.290564in}}%
\pgfpathlineto{\pgfqpoint{3.625721in}{1.299786in}}%
\pgfpathlineto{\pgfqpoint{3.625721in}{1.299786in}}%
\pgfpathclose%
\pgfpathmoveto{\pgfqpoint{3.712720in}{1.324787in}}%
\pgfpathlineto{\pgfqpoint{3.712720in}{1.314989in}}%
\pgfpathquadraticcurveto{\pgfqpoint{3.708343in}{1.317240in}}{\pgfqpoint{3.703605in}{1.318357in}}%
\pgfpathquadraticcurveto{\pgfqpoint{3.698886in}{1.319492in}}{\pgfqpoint{3.693807in}{1.319492in}}%
\pgfpathquadraticcurveto{\pgfqpoint{3.686098in}{1.319492in}}{\pgfqpoint{3.682243in}{1.317132in}}%
\pgfpathquadraticcurveto{\pgfqpoint{3.678388in}{1.314773in}}{\pgfqpoint{3.678388in}{1.310035in}}%
\pgfpathquadraticcurveto{\pgfqpoint{3.678388in}{1.306433in}}{\pgfqpoint{3.681144in}{1.304380in}}%
\pgfpathquadraticcurveto{\pgfqpoint{3.683900in}{1.302326in}}{\pgfqpoint{3.692240in}{1.300471in}}%
\pgfpathlineto{\pgfqpoint{3.695788in}{1.299678in}}%
\pgfpathquadraticcurveto{\pgfqpoint{3.706812in}{1.297319in}}{\pgfqpoint{3.711459in}{1.293014in}}%
\pgfpathquadraticcurveto{\pgfqpoint{3.716106in}{1.288709in}}{\pgfqpoint{3.716106in}{1.281000in}}%
\pgfpathquadraticcurveto{\pgfqpoint{3.716106in}{1.272210in}}{\pgfqpoint{3.709153in}{1.267076in}}%
\pgfpathquadraticcurveto{\pgfqpoint{3.702201in}{1.261961in}}{\pgfqpoint{3.690042in}{1.261961in}}%
\pgfpathquadraticcurveto{\pgfqpoint{3.684981in}{1.261961in}}{\pgfqpoint{3.679487in}{1.262952in}}%
\pgfpathquadraticcurveto{\pgfqpoint{3.673993in}{1.263942in}}{\pgfqpoint{3.667923in}{1.265906in}}%
\pgfpathlineto{\pgfqpoint{3.667923in}{1.276605in}}%
\pgfpathquadraticcurveto{\pgfqpoint{3.673669in}{1.273615in}}{\pgfqpoint{3.679235in}{1.272120in}}%
\pgfpathquadraticcurveto{\pgfqpoint{3.684801in}{1.270643in}}{\pgfqpoint{3.690276in}{1.270643in}}%
\pgfpathquadraticcurveto{\pgfqpoint{3.697589in}{1.270643in}}{\pgfqpoint{3.701516in}{1.273146in}}%
\pgfpathquadraticcurveto{\pgfqpoint{3.705461in}{1.275650in}}{\pgfqpoint{3.705461in}{1.280207in}}%
\pgfpathquadraticcurveto{\pgfqpoint{3.705461in}{1.284422in}}{\pgfqpoint{3.702615in}{1.286674in}}%
\pgfpathquadraticcurveto{\pgfqpoint{3.699787in}{1.288925in}}{\pgfqpoint{3.690150in}{1.291014in}}%
\pgfpathlineto{\pgfqpoint{3.686548in}{1.291861in}}%
\pgfpathquadraticcurveto{\pgfqpoint{3.676929in}{1.293878in}}{\pgfqpoint{3.672643in}{1.298075in}}%
\pgfpathquadraticcurveto{\pgfqpoint{3.668374in}{1.302272in}}{\pgfqpoint{3.668374in}{1.309585in}}%
\pgfpathquadraticcurveto{\pgfqpoint{3.668374in}{1.318483in}}{\pgfqpoint{3.674678in}{1.323310in}}%
\pgfpathquadraticcurveto{\pgfqpoint{3.680982in}{1.328156in}}{\pgfqpoint{3.692582in}{1.328156in}}%
\pgfpathquadraticcurveto{\pgfqpoint{3.698310in}{1.328156in}}{\pgfqpoint{3.703371in}{1.327309in}}%
\pgfpathquadraticcurveto{\pgfqpoint{3.708451in}{1.326480in}}{\pgfqpoint{3.712720in}{1.324787in}}%
\pgfpathlineto{\pgfqpoint{3.712720in}{1.324787in}}%
\pgfpathclose%
\pgfpathmoveto{\pgfqpoint{3.742836in}{1.344547in}}%
\pgfpathlineto{\pgfqpoint{3.742836in}{1.326643in}}%
\pgfpathlineto{\pgfqpoint{3.764162in}{1.326643in}}%
\pgfpathlineto{\pgfqpoint{3.764162in}{1.318591in}}%
\pgfpathlineto{\pgfqpoint{3.742836in}{1.318591in}}%
\pgfpathlineto{\pgfqpoint{3.742836in}{1.284368in}}%
\pgfpathquadraticcurveto{\pgfqpoint{3.742836in}{1.276659in}}{\pgfqpoint{3.744943in}{1.274461in}}%
\pgfpathquadraticcurveto{\pgfqpoint{3.747051in}{1.272264in}}{\pgfqpoint{3.753535in}{1.272264in}}%
\pgfpathlineto{\pgfqpoint{3.764162in}{1.272264in}}%
\pgfpathlineto{\pgfqpoint{3.764162in}{1.263600in}}%
\pgfpathlineto{\pgfqpoint{3.753535in}{1.263600in}}%
\pgfpathquadraticcurveto{\pgfqpoint{3.741539in}{1.263600in}}{\pgfqpoint{3.736982in}{1.268067in}}%
\pgfpathquadraticcurveto{\pgfqpoint{3.732425in}{1.272552in}}{\pgfqpoint{3.732425in}{1.284368in}}%
\pgfpathlineto{\pgfqpoint{3.732425in}{1.318591in}}%
\pgfpathlineto{\pgfqpoint{3.724824in}{1.318591in}}%
\pgfpathlineto{\pgfqpoint{3.724824in}{1.326643in}}%
\pgfpathlineto{\pgfqpoint{3.732425in}{1.326643in}}%
\pgfpathlineto{\pgfqpoint{3.732425in}{1.344547in}}%
\pgfpathlineto{\pgfqpoint{3.742836in}{1.344547in}}%
\pgfpathlineto{\pgfqpoint{3.742836in}{1.344547in}}%
\pgfpathclose%
\pgfpathmoveto{\pgfqpoint{3.776717in}{1.288475in}}%
\pgfpathlineto{\pgfqpoint{3.776717in}{1.326643in}}%
\pgfpathlineto{\pgfqpoint{3.787074in}{1.326643in}}%
\pgfpathlineto{\pgfqpoint{3.787074in}{1.288871in}}%
\pgfpathquadraticcurveto{\pgfqpoint{3.787074in}{1.279919in}}{\pgfqpoint{3.790550in}{1.275434in}}%
\pgfpathquadraticcurveto{\pgfqpoint{3.794044in}{1.270967in}}{\pgfqpoint{3.801033in}{1.270967in}}%
\pgfpathquadraticcurveto{\pgfqpoint{3.809409in}{1.270967in}}{\pgfqpoint{3.814272in}{1.276317in}}%
\pgfpathquadraticcurveto{\pgfqpoint{3.819153in}{1.281666in}}{\pgfqpoint{3.819153in}{1.290906in}}%
\pgfpathlineto{\pgfqpoint{3.819153in}{1.326643in}}%
\pgfpathlineto{\pgfqpoint{3.829510in}{1.326643in}}%
\pgfpathlineto{\pgfqpoint{3.829510in}{1.263600in}}%
\pgfpathlineto{\pgfqpoint{3.819153in}{1.263600in}}%
\pgfpathlineto{\pgfqpoint{3.819153in}{1.273291in}}%
\pgfpathquadraticcurveto{\pgfqpoint{3.815389in}{1.267545in}}{\pgfqpoint{3.810400in}{1.264753in}}%
\pgfpathquadraticcurveto{\pgfqpoint{3.805428in}{1.261961in}}{\pgfqpoint{3.798836in}{1.261961in}}%
\pgfpathquadraticcurveto{\pgfqpoint{3.787974in}{1.261961in}}{\pgfqpoint{3.782337in}{1.268715in}}%
\pgfpathquadraticcurveto{\pgfqpoint{3.776717in}{1.275470in}}{\pgfqpoint{3.776717in}{1.288475in}}%
\pgfpathlineto{\pgfqpoint{3.776717in}{1.288475in}}%
\pgfpathclose%
\pgfpathmoveto{\pgfqpoint{3.802780in}{1.328156in}}%
\pgfpathlineto{\pgfqpoint{3.802780in}{1.328156in}}%
\pgfpathlineto{\pgfqpoint{3.802780in}{1.328156in}}%
\pgfpathclose%
\pgfpathmoveto{\pgfqpoint{3.892319in}{1.317078in}}%
\pgfpathlineto{\pgfqpoint{3.892319in}{1.351193in}}%
\pgfpathlineto{\pgfqpoint{3.902676in}{1.351193in}}%
\pgfpathlineto{\pgfqpoint{3.902676in}{1.263600in}}%
\pgfpathlineto{\pgfqpoint{3.892319in}{1.263600in}}%
\pgfpathlineto{\pgfqpoint{3.892319in}{1.273056in}}%
\pgfpathquadraticcurveto{\pgfqpoint{3.889059in}{1.267437in}}{\pgfqpoint{3.884069in}{1.264699in}}%
\pgfpathquadraticcurveto{\pgfqpoint{3.879098in}{1.261961in}}{\pgfqpoint{3.872109in}{1.261961in}}%
\pgfpathquadraticcurveto{\pgfqpoint{3.860689in}{1.261961in}}{\pgfqpoint{3.853503in}{1.271075in}}%
\pgfpathquadraticcurveto{\pgfqpoint{3.846334in}{1.280207in}}{\pgfqpoint{3.846334in}{1.295067in}}%
\pgfpathquadraticcurveto{\pgfqpoint{3.846334in}{1.309927in}}{\pgfqpoint{3.853503in}{1.319041in}}%
\pgfpathquadraticcurveto{\pgfqpoint{3.860689in}{1.328156in}}{\pgfqpoint{3.872109in}{1.328156in}}%
\pgfpathquadraticcurveto{\pgfqpoint{3.879098in}{1.328156in}}{\pgfqpoint{3.884069in}{1.325418in}}%
\pgfpathquadraticcurveto{\pgfqpoint{3.889059in}{1.322698in}}{\pgfqpoint{3.892319in}{1.317078in}}%
\pgfpathlineto{\pgfqpoint{3.892319in}{1.317078in}}%
\pgfpathclose%
\pgfpathmoveto{\pgfqpoint{3.857033in}{1.295067in}}%
\pgfpathquadraticcurveto{\pgfqpoint{3.857033in}{1.283648in}}{\pgfqpoint{3.861734in}{1.277145in}}%
\pgfpathquadraticcurveto{\pgfqpoint{3.866435in}{1.270643in}}{\pgfqpoint{3.874649in}{1.270643in}}%
\pgfpathquadraticcurveto{\pgfqpoint{3.882862in}{1.270643in}}{\pgfqpoint{3.887582in}{1.277145in}}%
\pgfpathquadraticcurveto{\pgfqpoint{3.892319in}{1.283648in}}{\pgfqpoint{3.892319in}{1.295067in}}%
\pgfpathquadraticcurveto{\pgfqpoint{3.892319in}{1.306487in}}{\pgfqpoint{3.887582in}{1.312989in}}%
\pgfpathquadraticcurveto{\pgfqpoint{3.882862in}{1.319492in}}{\pgfqpoint{3.874649in}{1.319492in}}%
\pgfpathquadraticcurveto{\pgfqpoint{3.866435in}{1.319492in}}{\pgfqpoint{3.861734in}{1.312989in}}%
\pgfpathquadraticcurveto{\pgfqpoint{3.857033in}{1.306487in}}{\pgfqpoint{3.857033in}{1.295067in}}%
\pgfpathlineto{\pgfqpoint{3.857033in}{1.295067in}}%
\pgfpathclose%
\pgfpathmoveto{\pgfqpoint{3.950246in}{1.257746in}}%
\pgfpathquadraticcurveto{\pgfqpoint{3.945869in}{1.246488in}}{\pgfqpoint{3.941690in}{1.243066in}}%
\pgfpathquadraticcurveto{\pgfqpoint{3.937529in}{1.239626in}}{\pgfqpoint{3.930559in}{1.239626in}}%
\pgfpathlineto{\pgfqpoint{3.922273in}{1.239626in}}%
\pgfpathlineto{\pgfqpoint{3.922273in}{1.248290in}}%
\pgfpathlineto{\pgfqpoint{3.928361in}{1.248290in}}%
\pgfpathquadraticcurveto{\pgfqpoint{3.932630in}{1.248290in}}{\pgfqpoint{3.934990in}{1.250325in}}%
\pgfpathquadraticcurveto{\pgfqpoint{3.937367in}{1.252342in}}{\pgfqpoint{3.940231in}{1.259889in}}%
\pgfpathlineto{\pgfqpoint{3.942086in}{1.264609in}}%
\pgfpathlineto{\pgfqpoint{3.916599in}{1.326643in}}%
\pgfpathlineto{\pgfqpoint{3.927569in}{1.326643in}}%
\pgfpathlineto{\pgfqpoint{3.947274in}{1.277343in}}%
\pgfpathlineto{\pgfqpoint{3.966979in}{1.326643in}}%
\pgfpathlineto{\pgfqpoint{3.977949in}{1.326643in}}%
\pgfpathlineto{\pgfqpoint{3.950246in}{1.257746in}}%
\pgfpathlineto{\pgfqpoint{3.950246in}{1.257746in}}%
\pgfpathclose%
\pgfusepath{fill}%
\end{pgfscope}%
\begin{pgfscope}%
\pgfsetbuttcap%
\pgfsetroundjoin%
\definecolor{currentfill}{rgb}{0.309804,0.372549,0.419608}%
\pgfsetfillcolor{currentfill}%
\pgfsetlinewidth{0.000000pt}%
\definecolor{currentstroke}{rgb}{0.309804,0.372549,0.419608}%
\pgfsetstrokecolor{currentstroke}%
\pgfsetdash{}{0pt}%
\pgfpathmoveto{\pgfqpoint{4.106943in}{1.295283in}}%
\pgfpathquadraticcurveto{\pgfqpoint{4.094389in}{1.295283in}}{\pgfqpoint{4.089543in}{1.292419in}}%
\pgfpathquadraticcurveto{\pgfqpoint{4.084698in}{1.289556in}}{\pgfqpoint{4.084698in}{1.282621in}}%
\pgfpathquadraticcurveto{\pgfqpoint{4.084698in}{1.277109in}}{\pgfqpoint{4.088337in}{1.273867in}}%
\pgfpathquadraticcurveto{\pgfqpoint{4.091975in}{1.270643in}}{\pgfqpoint{4.098207in}{1.270643in}}%
\pgfpathquadraticcurveto{\pgfqpoint{4.106835in}{1.270643in}}{\pgfqpoint{4.112041in}{1.276749in}}%
\pgfpathquadraticcurveto{\pgfqpoint{4.117246in}{1.282855in}}{\pgfqpoint{4.117246in}{1.292978in}}%
\pgfpathlineto{\pgfqpoint{4.117246in}{1.295283in}}%
\pgfpathlineto{\pgfqpoint{4.106943in}{1.295283in}}%
\pgfpathlineto{\pgfqpoint{4.106943in}{1.295283in}}%
\pgfpathclose%
\pgfpathmoveto{\pgfqpoint{4.127603in}{1.299570in}}%
\pgfpathlineto{\pgfqpoint{4.127603in}{1.263600in}}%
\pgfpathlineto{\pgfqpoint{4.117246in}{1.263600in}}%
\pgfpathlineto{\pgfqpoint{4.117246in}{1.273164in}}%
\pgfpathquadraticcurveto{\pgfqpoint{4.113698in}{1.267437in}}{\pgfqpoint{4.108402in}{1.264699in}}%
\pgfpathquadraticcurveto{\pgfqpoint{4.103107in}{1.261961in}}{\pgfqpoint{4.095451in}{1.261961in}}%
\pgfpathquadraticcurveto{\pgfqpoint{4.085779in}{1.261961in}}{\pgfqpoint{4.080051in}{1.267401in}}%
\pgfpathquadraticcurveto{\pgfqpoint{4.074341in}{1.272840in}}{\pgfqpoint{4.074341in}{1.281954in}}%
\pgfpathquadraticcurveto{\pgfqpoint{4.074341in}{1.292582in}}{\pgfqpoint{4.081456in}{1.297985in}}%
\pgfpathquadraticcurveto{\pgfqpoint{4.088589in}{1.303389in}}{\pgfqpoint{4.102710in}{1.303389in}}%
\pgfpathlineto{\pgfqpoint{4.117246in}{1.303389in}}%
\pgfpathlineto{\pgfqpoint{4.117246in}{1.304416in}}%
\pgfpathquadraticcurveto{\pgfqpoint{4.117246in}{1.311566in}}{\pgfqpoint{4.112545in}{1.315475in}}%
\pgfpathquadraticcurveto{\pgfqpoint{4.107844in}{1.319384in}}{\pgfqpoint{4.099342in}{1.319384in}}%
\pgfpathquadraticcurveto{\pgfqpoint{4.093938in}{1.319384in}}{\pgfqpoint{4.088805in}{1.318087in}}%
\pgfpathquadraticcurveto{\pgfqpoint{4.083689in}{1.316790in}}{\pgfqpoint{4.078970in}{1.314196in}}%
\pgfpathlineto{\pgfqpoint{4.078970in}{1.323779in}}%
\pgfpathquadraticcurveto{\pgfqpoint{4.084644in}{1.325976in}}{\pgfqpoint{4.089994in}{1.327057in}}%
\pgfpathquadraticcurveto{\pgfqpoint{4.095343in}{1.328156in}}{\pgfqpoint{4.100405in}{1.328156in}}%
\pgfpathquadraticcurveto{\pgfqpoint{4.114094in}{1.328156in}}{\pgfqpoint{4.120849in}{1.321059in}}%
\pgfpathquadraticcurveto{\pgfqpoint{4.127603in}{1.313980in}}{\pgfqpoint{4.127603in}{1.299570in}}%
\pgfpathlineto{\pgfqpoint{4.127603in}{1.299570in}}%
\pgfpathclose%
\pgfpathmoveto{\pgfqpoint{4.141508in}{1.326643in}}%
\pgfpathlineto{\pgfqpoint{4.152478in}{1.326643in}}%
\pgfpathlineto{\pgfqpoint{4.172183in}{1.273741in}}%
\pgfpathlineto{\pgfqpoint{4.191888in}{1.326643in}}%
\pgfpathlineto{\pgfqpoint{4.202858in}{1.326643in}}%
\pgfpathlineto{\pgfqpoint{4.179208in}{1.263600in}}%
\pgfpathlineto{\pgfqpoint{4.165140in}{1.263600in}}%
\pgfpathlineto{\pgfqpoint{4.141508in}{1.326643in}}%
\pgfpathlineto{\pgfqpoint{4.141508in}{1.326643in}}%
\pgfpathclose%
\pgfpathmoveto{\pgfqpoint{4.271088in}{1.297715in}}%
\pgfpathlineto{\pgfqpoint{4.271088in}{1.292654in}}%
\pgfpathlineto{\pgfqpoint{4.223464in}{1.292654in}}%
\pgfpathquadraticcurveto{\pgfqpoint{4.224148in}{1.281954in}}{\pgfqpoint{4.229912in}{1.276353in}}%
\pgfpathquadraticcurveto{\pgfqpoint{4.235676in}{1.270751in}}{\pgfqpoint{4.245979in}{1.270751in}}%
\pgfpathquadraticcurveto{\pgfqpoint{4.251941in}{1.270751in}}{\pgfqpoint{4.257543in}{1.272210in}}%
\pgfpathquadraticcurveto{\pgfqpoint{4.263145in}{1.273669in}}{\pgfqpoint{4.268674in}{1.276605in}}%
\pgfpathlineto{\pgfqpoint{4.268674in}{1.266806in}}%
\pgfpathquadraticcurveto{\pgfqpoint{4.263090in}{1.264447in}}{\pgfqpoint{4.257237in}{1.263204in}}%
\pgfpathquadraticcurveto{\pgfqpoint{4.251383in}{1.261961in}}{\pgfqpoint{4.245367in}{1.261961in}}%
\pgfpathquadraticcurveto{\pgfqpoint{4.230272in}{1.261961in}}{\pgfqpoint{4.221464in}{1.270733in}}%
\pgfpathquadraticcurveto{\pgfqpoint{4.212656in}{1.279523in}}{\pgfqpoint{4.212656in}{1.294509in}}%
\pgfpathquadraticcurveto{\pgfqpoint{4.212656in}{1.309981in}}{\pgfqpoint{4.221014in}{1.319059in}}%
\pgfpathquadraticcurveto{\pgfqpoint{4.229372in}{1.328156in}}{\pgfqpoint{4.243565in}{1.328156in}}%
\pgfpathquadraticcurveto{\pgfqpoint{4.256282in}{1.328156in}}{\pgfqpoint{4.263685in}{1.319960in}}%
\pgfpathquadraticcurveto{\pgfqpoint{4.271088in}{1.311783in}}{\pgfqpoint{4.271088in}{1.297715in}}%
\pgfpathlineto{\pgfqpoint{4.271088in}{1.297715in}}%
\pgfpathclose%
\pgfpathmoveto{\pgfqpoint{4.260731in}{1.300759in}}%
\pgfpathquadraticcurveto{\pgfqpoint{4.260623in}{1.309243in}}{\pgfqpoint{4.255976in}{1.314304in}}%
\pgfpathquadraticcurveto{\pgfqpoint{4.251329in}{1.319384in}}{\pgfqpoint{4.243673in}{1.319384in}}%
\pgfpathquadraticcurveto{\pgfqpoint{4.235010in}{1.319384in}}{\pgfqpoint{4.229804in}{1.314484in}}%
\pgfpathquadraticcurveto{\pgfqpoint{4.224599in}{1.309585in}}{\pgfqpoint{4.223806in}{1.300687in}}%
\pgfpathlineto{\pgfqpoint{4.260731in}{1.300759in}}%
\pgfpathlineto{\pgfqpoint{4.260731in}{1.300759in}}%
\pgfpathclose%
\pgfpathmoveto{\pgfqpoint{4.324620in}{1.316970in}}%
\pgfpathquadraticcurveto{\pgfqpoint{4.322873in}{1.317979in}}{\pgfqpoint{4.320819in}{1.318447in}}%
\pgfpathquadraticcurveto{\pgfqpoint{4.318766in}{1.318933in}}{\pgfqpoint{4.316298in}{1.318933in}}%
\pgfpathquadraticcurveto{\pgfqpoint{4.307508in}{1.318933in}}{\pgfqpoint{4.302807in}{1.313223in}}%
\pgfpathquadraticcurveto{\pgfqpoint{4.298106in}{1.307514in}}{\pgfqpoint{4.298106in}{1.296814in}}%
\pgfpathlineto{\pgfqpoint{4.298106in}{1.263600in}}%
\pgfpathlineto{\pgfqpoint{4.287695in}{1.263600in}}%
\pgfpathlineto{\pgfqpoint{4.287695in}{1.326643in}}%
\pgfpathlineto{\pgfqpoint{4.298106in}{1.326643in}}%
\pgfpathlineto{\pgfqpoint{4.298106in}{1.316844in}}%
\pgfpathquadraticcurveto{\pgfqpoint{4.301384in}{1.322590in}}{\pgfqpoint{4.306608in}{1.325364in}}%
\pgfpathquadraticcurveto{\pgfqpoint{4.311849in}{1.328156in}}{\pgfqpoint{4.319342in}{1.328156in}}%
\pgfpathquadraticcurveto{\pgfqpoint{4.320405in}{1.328156in}}{\pgfqpoint{4.321702in}{1.328011in}}%
\pgfpathquadraticcurveto{\pgfqpoint{4.322999in}{1.327885in}}{\pgfqpoint{4.324566in}{1.327597in}}%
\pgfpathlineto{\pgfqpoint{4.324620in}{1.316970in}}%
\pgfpathlineto{\pgfqpoint{4.324620in}{1.316970in}}%
\pgfpathclose%
\pgfpathmoveto{\pgfqpoint{4.364139in}{1.295283in}}%
\pgfpathquadraticcurveto{\pgfqpoint{4.351584in}{1.295283in}}{\pgfqpoint{4.346739in}{1.292419in}}%
\pgfpathquadraticcurveto{\pgfqpoint{4.341894in}{1.289556in}}{\pgfqpoint{4.341894in}{1.282621in}}%
\pgfpathquadraticcurveto{\pgfqpoint{4.341894in}{1.277109in}}{\pgfqpoint{4.345532in}{1.273867in}}%
\pgfpathquadraticcurveto{\pgfqpoint{4.349171in}{1.270643in}}{\pgfqpoint{4.355403in}{1.270643in}}%
\pgfpathquadraticcurveto{\pgfqpoint{4.364031in}{1.270643in}}{\pgfqpoint{4.369236in}{1.276749in}}%
\pgfpathquadraticcurveto{\pgfqpoint{4.374442in}{1.282855in}}{\pgfqpoint{4.374442in}{1.292978in}}%
\pgfpathlineto{\pgfqpoint{4.374442in}{1.295283in}}%
\pgfpathlineto{\pgfqpoint{4.364139in}{1.295283in}}%
\pgfpathlineto{\pgfqpoint{4.364139in}{1.295283in}}%
\pgfpathclose%
\pgfpathmoveto{\pgfqpoint{4.384799in}{1.299570in}}%
\pgfpathlineto{\pgfqpoint{4.384799in}{1.263600in}}%
\pgfpathlineto{\pgfqpoint{4.374442in}{1.263600in}}%
\pgfpathlineto{\pgfqpoint{4.374442in}{1.273164in}}%
\pgfpathquadraticcurveto{\pgfqpoint{4.370893in}{1.267437in}}{\pgfqpoint{4.365598in}{1.264699in}}%
\pgfpathquadraticcurveto{\pgfqpoint{4.360302in}{1.261961in}}{\pgfqpoint{4.352647in}{1.261961in}}%
\pgfpathquadraticcurveto{\pgfqpoint{4.342974in}{1.261961in}}{\pgfqpoint{4.337247in}{1.267401in}}%
\pgfpathquadraticcurveto{\pgfqpoint{4.331537in}{1.272840in}}{\pgfqpoint{4.331537in}{1.281954in}}%
\pgfpathquadraticcurveto{\pgfqpoint{4.331537in}{1.292582in}}{\pgfqpoint{4.338651in}{1.297985in}}%
\pgfpathquadraticcurveto{\pgfqpoint{4.345784in}{1.303389in}}{\pgfqpoint{4.359906in}{1.303389in}}%
\pgfpathlineto{\pgfqpoint{4.374442in}{1.303389in}}%
\pgfpathlineto{\pgfqpoint{4.374442in}{1.304416in}}%
\pgfpathquadraticcurveto{\pgfqpoint{4.374442in}{1.311566in}}{\pgfqpoint{4.369740in}{1.315475in}}%
\pgfpathquadraticcurveto{\pgfqpoint{4.365039in}{1.319384in}}{\pgfqpoint{4.356538in}{1.319384in}}%
\pgfpathquadraticcurveto{\pgfqpoint{4.351134in}{1.319384in}}{\pgfqpoint{4.346000in}{1.318087in}}%
\pgfpathquadraticcurveto{\pgfqpoint{4.340885in}{1.316790in}}{\pgfqpoint{4.336166in}{1.314196in}}%
\pgfpathlineto{\pgfqpoint{4.336166in}{1.323779in}}%
\pgfpathquadraticcurveto{\pgfqpoint{4.341840in}{1.325976in}}{\pgfqpoint{4.347189in}{1.327057in}}%
\pgfpathquadraticcurveto{\pgfqpoint{4.352539in}{1.328156in}}{\pgfqpoint{4.357600in}{1.328156in}}%
\pgfpathquadraticcurveto{\pgfqpoint{4.371289in}{1.328156in}}{\pgfqpoint{4.378044in}{1.321059in}}%
\pgfpathquadraticcurveto{\pgfqpoint{4.384799in}{1.313980in}}{\pgfqpoint{4.384799in}{1.299570in}}%
\pgfpathlineto{\pgfqpoint{4.384799in}{1.299570in}}%
\pgfpathclose%
\pgfpathmoveto{\pgfqpoint{4.447607in}{1.295860in}}%
\pgfpathquadraticcurveto{\pgfqpoint{4.447607in}{1.307117in}}{\pgfqpoint{4.442960in}{1.313296in}}%
\pgfpathquadraticcurveto{\pgfqpoint{4.438331in}{1.319492in}}{\pgfqpoint{4.429937in}{1.319492in}}%
\pgfpathquadraticcurveto{\pgfqpoint{4.421615in}{1.319492in}}{\pgfqpoint{4.416968in}{1.313296in}}%
\pgfpathquadraticcurveto{\pgfqpoint{4.412321in}{1.307117in}}{\pgfqpoint{4.412321in}{1.295860in}}%
\pgfpathquadraticcurveto{\pgfqpoint{4.412321in}{1.284656in}}{\pgfqpoint{4.416968in}{1.278460in}}%
\pgfpathquadraticcurveto{\pgfqpoint{4.421615in}{1.272264in}}{\pgfqpoint{4.429937in}{1.272264in}}%
\pgfpathquadraticcurveto{\pgfqpoint{4.438331in}{1.272264in}}{\pgfqpoint{4.442960in}{1.278460in}}%
\pgfpathquadraticcurveto{\pgfqpoint{4.447607in}{1.284656in}}{\pgfqpoint{4.447607in}{1.295860in}}%
\pgfpathlineto{\pgfqpoint{4.447607in}{1.295860in}}%
\pgfpathclose%
\pgfpathmoveto{\pgfqpoint{4.457964in}{1.271417in}}%
\pgfpathquadraticcurveto{\pgfqpoint{4.457964in}{1.255332in}}{\pgfqpoint{4.450813in}{1.247479in}}%
\pgfpathquadraticcurveto{\pgfqpoint{4.443680in}{1.239626in}}{\pgfqpoint{4.428928in}{1.239626in}}%
\pgfpathquadraticcurveto{\pgfqpoint{4.423471in}{1.239626in}}{\pgfqpoint{4.418625in}{1.240436in}}%
\pgfpathquadraticcurveto{\pgfqpoint{4.413780in}{1.241247in}}{\pgfqpoint{4.409223in}{1.242940in}}%
\pgfpathlineto{\pgfqpoint{4.409223in}{1.253009in}}%
\pgfpathquadraticcurveto{\pgfqpoint{4.413780in}{1.250541in}}{\pgfqpoint{4.418229in}{1.249370in}}%
\pgfpathquadraticcurveto{\pgfqpoint{4.422678in}{1.248182in}}{\pgfqpoint{4.427289in}{1.248182in}}%
\pgfpathquadraticcurveto{\pgfqpoint{4.437484in}{1.248182in}}{\pgfqpoint{4.442546in}{1.253495in}}%
\pgfpathquadraticcurveto{\pgfqpoint{4.447607in}{1.258809in}}{\pgfqpoint{4.447607in}{1.269562in}}%
\pgfpathlineto{\pgfqpoint{4.447607in}{1.274695in}}%
\pgfpathquadraticcurveto{\pgfqpoint{4.444401in}{1.269112in}}{\pgfqpoint{4.439393in}{1.266356in}}%
\pgfpathquadraticcurveto{\pgfqpoint{4.434386in}{1.263600in}}{\pgfqpoint{4.427397in}{1.263600in}}%
\pgfpathquadraticcurveto{\pgfqpoint{4.415816in}{1.263600in}}{\pgfqpoint{4.408719in}{1.272426in}}%
\pgfpathquadraticcurveto{\pgfqpoint{4.401622in}{1.281270in}}{\pgfqpoint{4.401622in}{1.295860in}}%
\pgfpathquadraticcurveto{\pgfqpoint{4.401622in}{1.310486in}}{\pgfqpoint{4.408719in}{1.319312in}}%
\pgfpathquadraticcurveto{\pgfqpoint{4.415816in}{1.328156in}}{\pgfqpoint{4.427397in}{1.328156in}}%
\pgfpathquadraticcurveto{\pgfqpoint{4.434386in}{1.328156in}}{\pgfqpoint{4.439393in}{1.325400in}}%
\pgfpathquadraticcurveto{\pgfqpoint{4.444401in}{1.322644in}}{\pgfqpoint{4.447607in}{1.317078in}}%
\pgfpathlineto{\pgfqpoint{4.447607in}{1.326643in}}%
\pgfpathlineto{\pgfqpoint{4.457964in}{1.326643in}}%
\pgfpathlineto{\pgfqpoint{4.457964in}{1.271417in}}%
\pgfpathlineto{\pgfqpoint{4.457964in}{1.271417in}}%
\pgfpathclose%
\pgfpathmoveto{\pgfqpoint{4.533237in}{1.297715in}}%
\pgfpathlineto{\pgfqpoint{4.533237in}{1.292654in}}%
\pgfpathlineto{\pgfqpoint{4.485613in}{1.292654in}}%
\pgfpathquadraticcurveto{\pgfqpoint{4.486297in}{1.281954in}}{\pgfqpoint{4.492061in}{1.276353in}}%
\pgfpathquadraticcurveto{\pgfqpoint{4.497825in}{1.270751in}}{\pgfqpoint{4.508128in}{1.270751in}}%
\pgfpathquadraticcurveto{\pgfqpoint{4.514090in}{1.270751in}}{\pgfqpoint{4.519692in}{1.272210in}}%
\pgfpathquadraticcurveto{\pgfqpoint{4.525293in}{1.273669in}}{\pgfqpoint{4.530823in}{1.276605in}}%
\pgfpathlineto{\pgfqpoint{4.530823in}{1.266806in}}%
\pgfpathquadraticcurveto{\pgfqpoint{4.525239in}{1.264447in}}{\pgfqpoint{4.519385in}{1.263204in}}%
\pgfpathquadraticcurveto{\pgfqpoint{4.513531in}{1.261961in}}{\pgfqpoint{4.507515in}{1.261961in}}%
\pgfpathquadraticcurveto{\pgfqpoint{4.492421in}{1.261961in}}{\pgfqpoint{4.483613in}{1.270733in}}%
\pgfpathquadraticcurveto{\pgfqpoint{4.474805in}{1.279523in}}{\pgfqpoint{4.474805in}{1.294509in}}%
\pgfpathquadraticcurveto{\pgfqpoint{4.474805in}{1.309981in}}{\pgfqpoint{4.483163in}{1.319059in}}%
\pgfpathquadraticcurveto{\pgfqpoint{4.491521in}{1.328156in}}{\pgfqpoint{4.505714in}{1.328156in}}%
\pgfpathquadraticcurveto{\pgfqpoint{4.518431in}{1.328156in}}{\pgfqpoint{4.525834in}{1.319960in}}%
\pgfpathquadraticcurveto{\pgfqpoint{4.533237in}{1.311783in}}{\pgfqpoint{4.533237in}{1.297715in}}%
\pgfpathlineto{\pgfqpoint{4.533237in}{1.297715in}}%
\pgfpathclose%
\pgfpathmoveto{\pgfqpoint{4.522880in}{1.300759in}}%
\pgfpathquadraticcurveto{\pgfqpoint{4.522772in}{1.309243in}}{\pgfqpoint{4.518125in}{1.314304in}}%
\pgfpathquadraticcurveto{\pgfqpoint{4.513477in}{1.319384in}}{\pgfqpoint{4.505822in}{1.319384in}}%
\pgfpathquadraticcurveto{\pgfqpoint{4.497158in}{1.319384in}}{\pgfqpoint{4.491953in}{1.314484in}}%
\pgfpathquadraticcurveto{\pgfqpoint{4.486747in}{1.309585in}}{\pgfqpoint{4.485955in}{1.300687in}}%
\pgfpathlineto{\pgfqpoint{4.522880in}{1.300759in}}%
\pgfpathlineto{\pgfqpoint{4.522880in}{1.300759in}}%
\pgfpathclose%
\pgfpathmoveto{\pgfqpoint{4.590425in}{1.324787in}}%
\pgfpathlineto{\pgfqpoint{4.590425in}{1.314989in}}%
\pgfpathquadraticcurveto{\pgfqpoint{4.586048in}{1.317240in}}{\pgfqpoint{4.581311in}{1.318357in}}%
\pgfpathquadraticcurveto{\pgfqpoint{4.576592in}{1.319492in}}{\pgfqpoint{4.571513in}{1.319492in}}%
\pgfpathquadraticcurveto{\pgfqpoint{4.563803in}{1.319492in}}{\pgfqpoint{4.559949in}{1.317132in}}%
\pgfpathquadraticcurveto{\pgfqpoint{4.556094in}{1.314773in}}{\pgfqpoint{4.556094in}{1.310035in}}%
\pgfpathquadraticcurveto{\pgfqpoint{4.556094in}{1.306433in}}{\pgfqpoint{4.558850in}{1.304380in}}%
\pgfpathquadraticcurveto{\pgfqpoint{4.561606in}{1.302326in}}{\pgfqpoint{4.569946in}{1.300471in}}%
\pgfpathlineto{\pgfqpoint{4.573494in}{1.299678in}}%
\pgfpathquadraticcurveto{\pgfqpoint{4.584517in}{1.297319in}}{\pgfqpoint{4.589164in}{1.293014in}}%
\pgfpathquadraticcurveto{\pgfqpoint{4.593812in}{1.288709in}}{\pgfqpoint{4.593812in}{1.281000in}}%
\pgfpathquadraticcurveto{\pgfqpoint{4.593812in}{1.272210in}}{\pgfqpoint{4.586859in}{1.267076in}}%
\pgfpathquadraticcurveto{\pgfqpoint{4.579906in}{1.261961in}}{\pgfqpoint{4.567748in}{1.261961in}}%
\pgfpathquadraticcurveto{\pgfqpoint{4.562687in}{1.261961in}}{\pgfqpoint{4.557193in}{1.262952in}}%
\pgfpathquadraticcurveto{\pgfqpoint{4.551699in}{1.263942in}}{\pgfqpoint{4.545629in}{1.265906in}}%
\pgfpathlineto{\pgfqpoint{4.545629in}{1.276605in}}%
\pgfpathquadraticcurveto{\pgfqpoint{4.551375in}{1.273615in}}{\pgfqpoint{4.556941in}{1.272120in}}%
\pgfpathquadraticcurveto{\pgfqpoint{4.562507in}{1.270643in}}{\pgfqpoint{4.567982in}{1.270643in}}%
\pgfpathquadraticcurveto{\pgfqpoint{4.575295in}{1.270643in}}{\pgfqpoint{4.579222in}{1.273146in}}%
\pgfpathquadraticcurveto{\pgfqpoint{4.583166in}{1.275650in}}{\pgfqpoint{4.583166in}{1.280207in}}%
\pgfpathquadraticcurveto{\pgfqpoint{4.583166in}{1.284422in}}{\pgfqpoint{4.580321in}{1.286674in}}%
\pgfpathquadraticcurveto{\pgfqpoint{4.577493in}{1.288925in}}{\pgfqpoint{4.567856in}{1.291014in}}%
\pgfpathlineto{\pgfqpoint{4.564254in}{1.291861in}}%
\pgfpathquadraticcurveto{\pgfqpoint{4.554635in}{1.293878in}}{\pgfqpoint{4.550348in}{1.298075in}}%
\pgfpathquadraticcurveto{\pgfqpoint{4.546079in}{1.302272in}}{\pgfqpoint{4.546079in}{1.309585in}}%
\pgfpathquadraticcurveto{\pgfqpoint{4.546079in}{1.318483in}}{\pgfqpoint{4.552384in}{1.323310in}}%
\pgfpathquadraticcurveto{\pgfqpoint{4.558688in}{1.328156in}}{\pgfqpoint{4.570288in}{1.328156in}}%
\pgfpathquadraticcurveto{\pgfqpoint{4.576016in}{1.328156in}}{\pgfqpoint{4.581077in}{1.327309in}}%
\pgfpathquadraticcurveto{\pgfqpoint{4.586156in}{1.326480in}}{\pgfqpoint{4.590425in}{1.324787in}}%
\pgfpathlineto{\pgfqpoint{4.590425in}{1.324787in}}%
\pgfpathclose%
\pgfpathmoveto{\pgfqpoint{4.612941in}{1.277902in}}%
\pgfpathlineto{\pgfqpoint{4.624811in}{1.277902in}}%
\pgfpathlineto{\pgfqpoint{4.624811in}{1.268211in}}%
\pgfpathlineto{\pgfqpoint{4.615588in}{1.250199in}}%
\pgfpathlineto{\pgfqpoint{4.608329in}{1.250199in}}%
\pgfpathlineto{\pgfqpoint{4.612941in}{1.268211in}}%
\pgfpathlineto{\pgfqpoint{4.612941in}{1.277902in}}%
\pgfpathlineto{\pgfqpoint{4.612941in}{1.277902in}}%
\pgfpathclose%
\pgfusepath{fill}%
\end{pgfscope}%
\begin{pgfscope}%
\pgfsetbuttcap%
\pgfsetroundjoin%
\definecolor{currentfill}{rgb}{0.309804,0.372549,0.419608}%
\pgfsetfillcolor{currentfill}%
\pgfsetlinewidth{0.000000pt}%
\definecolor{currentstroke}{rgb}{0.309804,0.372549,0.419608}%
\pgfsetstrokecolor{currentstroke}%
\pgfsetdash{}{0pt}%
\pgfpathmoveto{\pgfqpoint{4.707955in}{1.326643in}}%
\pgfpathlineto{\pgfqpoint{4.718312in}{1.326643in}}%
\pgfpathlineto{\pgfqpoint{4.731262in}{1.277451in}}%
\pgfpathlineto{\pgfqpoint{4.744141in}{1.326643in}}%
\pgfpathlineto{\pgfqpoint{4.756353in}{1.326643in}}%
\pgfpathlineto{\pgfqpoint{4.769304in}{1.277451in}}%
\pgfpathlineto{\pgfqpoint{4.782201in}{1.326643in}}%
\pgfpathlineto{\pgfqpoint{4.792558in}{1.326643in}}%
\pgfpathlineto{\pgfqpoint{4.776059in}{1.263600in}}%
\pgfpathlineto{\pgfqpoint{4.763846in}{1.263600in}}%
\pgfpathlineto{\pgfqpoint{4.750283in}{1.315277in}}%
\pgfpathlineto{\pgfqpoint{4.736666in}{1.263600in}}%
\pgfpathlineto{\pgfqpoint{4.724436in}{1.263600in}}%
\pgfpathlineto{\pgfqpoint{4.707955in}{1.326643in}}%
\pgfpathlineto{\pgfqpoint{4.707955in}{1.326643in}}%
\pgfpathclose%
\pgfpathmoveto{\pgfqpoint{4.808246in}{1.326643in}}%
\pgfpathlineto{\pgfqpoint{4.818603in}{1.326643in}}%
\pgfpathlineto{\pgfqpoint{4.818603in}{1.263600in}}%
\pgfpathlineto{\pgfqpoint{4.808246in}{1.263600in}}%
\pgfpathlineto{\pgfqpoint{4.808246in}{1.326643in}}%
\pgfpathlineto{\pgfqpoint{4.808246in}{1.326643in}}%
\pgfpathclose%
\pgfpathmoveto{\pgfqpoint{4.808246in}{1.351193in}}%
\pgfpathlineto{\pgfqpoint{4.818603in}{1.351193in}}%
\pgfpathlineto{\pgfqpoint{4.818603in}{1.338062in}}%
\pgfpathlineto{\pgfqpoint{4.808246in}{1.338062in}}%
\pgfpathlineto{\pgfqpoint{4.808246in}{1.351193in}}%
\pgfpathlineto{\pgfqpoint{4.808246in}{1.351193in}}%
\pgfpathclose%
\pgfpathmoveto{\pgfqpoint{4.850521in}{1.344547in}}%
\pgfpathlineto{\pgfqpoint{4.850521in}{1.326643in}}%
\pgfpathlineto{\pgfqpoint{4.871847in}{1.326643in}}%
\pgfpathlineto{\pgfqpoint{4.871847in}{1.318591in}}%
\pgfpathlineto{\pgfqpoint{4.850521in}{1.318591in}}%
\pgfpathlineto{\pgfqpoint{4.850521in}{1.284368in}}%
\pgfpathquadraticcurveto{\pgfqpoint{4.850521in}{1.276659in}}{\pgfqpoint{4.852628in}{1.274461in}}%
\pgfpathquadraticcurveto{\pgfqpoint{4.854736in}{1.272264in}}{\pgfqpoint{4.861220in}{1.272264in}}%
\pgfpathlineto{\pgfqpoint{4.871847in}{1.272264in}}%
\pgfpathlineto{\pgfqpoint{4.871847in}{1.263600in}}%
\pgfpathlineto{\pgfqpoint{4.861220in}{1.263600in}}%
\pgfpathquadraticcurveto{\pgfqpoint{4.849224in}{1.263600in}}{\pgfqpoint{4.844667in}{1.268067in}}%
\pgfpathquadraticcurveto{\pgfqpoint{4.840110in}{1.272552in}}{\pgfqpoint{4.840110in}{1.284368in}}%
\pgfpathlineto{\pgfqpoint{4.840110in}{1.318591in}}%
\pgfpathlineto{\pgfqpoint{4.832509in}{1.318591in}}%
\pgfpathlineto{\pgfqpoint{4.832509in}{1.326643in}}%
\pgfpathlineto{\pgfqpoint{4.840110in}{1.326643in}}%
\pgfpathlineto{\pgfqpoint{4.840110in}{1.344547in}}%
\pgfpathlineto{\pgfqpoint{4.850521in}{1.344547in}}%
\pgfpathlineto{\pgfqpoint{4.850521in}{1.344547in}}%
\pgfpathclose%
\pgfpathmoveto{\pgfqpoint{4.937880in}{1.301660in}}%
\pgfpathlineto{\pgfqpoint{4.937880in}{1.263600in}}%
\pgfpathlineto{\pgfqpoint{4.927523in}{1.263600in}}%
\pgfpathlineto{\pgfqpoint{4.927523in}{1.301317in}}%
\pgfpathquadraticcurveto{\pgfqpoint{4.927523in}{1.310269in}}{\pgfqpoint{4.924028in}{1.314700in}}%
\pgfpathquadraticcurveto{\pgfqpoint{4.920534in}{1.319149in}}{\pgfqpoint{4.913563in}{1.319149in}}%
\pgfpathquadraticcurveto{\pgfqpoint{4.905170in}{1.319149in}}{\pgfqpoint{4.900324in}{1.313800in}}%
\pgfpathquadraticcurveto{\pgfqpoint{4.895479in}{1.308468in}}{\pgfqpoint{4.895479in}{1.299228in}}%
\pgfpathlineto{\pgfqpoint{4.895479in}{1.263600in}}%
\pgfpathlineto{\pgfqpoint{4.885068in}{1.263600in}}%
\pgfpathlineto{\pgfqpoint{4.885068in}{1.351193in}}%
\pgfpathlineto{\pgfqpoint{4.895479in}{1.351193in}}%
\pgfpathlineto{\pgfqpoint{4.895479in}{1.316844in}}%
\pgfpathquadraticcurveto{\pgfqpoint{4.899208in}{1.322536in}}{\pgfqpoint{4.904233in}{1.325346in}}%
\pgfpathquadraticcurveto{\pgfqpoint{4.909276in}{1.328156in}}{\pgfqpoint{4.915869in}{1.328156in}}%
\pgfpathquadraticcurveto{\pgfqpoint{4.926730in}{1.328156in}}{\pgfqpoint{4.932296in}{1.321437in}}%
\pgfpathquadraticcurveto{\pgfqpoint{4.937880in}{1.314718in}}{\pgfqpoint{4.937880in}{1.301660in}}%
\pgfpathlineto{\pgfqpoint{4.937880in}{1.301660in}}%
\pgfpathclose%
\pgfusepath{fill}%
\end{pgfscope}%
\begin{pgfscope}%
\pgfsetbuttcap%
\pgfsetroundjoin%
\definecolor{currentfill}{rgb}{0.309804,0.372549,0.419608}%
\pgfsetfillcolor{currentfill}%
\pgfsetlinewidth{0.000000pt}%
\definecolor{currentstroke}{rgb}{0.309804,0.372549,0.419608}%
\pgfsetstrokecolor{currentstroke}%
\pgfsetdash{}{0pt}%
\pgfpathmoveto{\pgfqpoint{5.039670in}{1.273056in}}%
\pgfpathlineto{\pgfqpoint{5.039670in}{1.239626in}}%
\pgfpathlineto{\pgfqpoint{5.029259in}{1.239626in}}%
\pgfpathlineto{\pgfqpoint{5.029259in}{1.326643in}}%
\pgfpathlineto{\pgfqpoint{5.039670in}{1.326643in}}%
\pgfpathlineto{\pgfqpoint{5.039670in}{1.317078in}}%
\pgfpathquadraticcurveto{\pgfqpoint{5.042949in}{1.322698in}}{\pgfqpoint{5.047920in}{1.325418in}}%
\pgfpathquadraticcurveto{\pgfqpoint{5.052909in}{1.328156in}}{\pgfqpoint{5.059826in}{1.328156in}}%
\pgfpathquadraticcurveto{\pgfqpoint{5.071318in}{1.328156in}}{\pgfqpoint{5.078487in}{1.319041in}}%
\pgfpathquadraticcurveto{\pgfqpoint{5.085673in}{1.309927in}}{\pgfqpoint{5.085673in}{1.295067in}}%
\pgfpathquadraticcurveto{\pgfqpoint{5.085673in}{1.280207in}}{\pgfqpoint{5.078487in}{1.271075in}}%
\pgfpathquadraticcurveto{\pgfqpoint{5.071318in}{1.261961in}}{\pgfqpoint{5.059826in}{1.261961in}}%
\pgfpathquadraticcurveto{\pgfqpoint{5.052909in}{1.261961in}}{\pgfqpoint{5.047920in}{1.264699in}}%
\pgfpathquadraticcurveto{\pgfqpoint{5.042949in}{1.267437in}}{\pgfqpoint{5.039670in}{1.273056in}}%
\pgfpathlineto{\pgfqpoint{5.039670in}{1.273056in}}%
\pgfpathclose%
\pgfpathmoveto{\pgfqpoint{5.074920in}{1.295067in}}%
\pgfpathquadraticcurveto{\pgfqpoint{5.074920in}{1.306487in}}{\pgfqpoint{5.070219in}{1.312989in}}%
\pgfpathquadraticcurveto{\pgfqpoint{5.065518in}{1.319492in}}{\pgfqpoint{5.057304in}{1.319492in}}%
\pgfpathquadraticcurveto{\pgfqpoint{5.049073in}{1.319492in}}{\pgfqpoint{5.044372in}{1.312989in}}%
\pgfpathquadraticcurveto{\pgfqpoint{5.039670in}{1.306487in}}{\pgfqpoint{5.039670in}{1.295067in}}%
\pgfpathquadraticcurveto{\pgfqpoint{5.039670in}{1.283648in}}{\pgfqpoint{5.044372in}{1.277145in}}%
\pgfpathquadraticcurveto{\pgfqpoint{5.049073in}{1.270643in}}{\pgfqpoint{5.057304in}{1.270643in}}%
\pgfpathquadraticcurveto{\pgfqpoint{5.065518in}{1.270643in}}{\pgfqpoint{5.070219in}{1.277145in}}%
\pgfpathquadraticcurveto{\pgfqpoint{5.074920in}{1.283648in}}{\pgfqpoint{5.074920in}{1.295067in}}%
\pgfpathlineto{\pgfqpoint{5.074920in}{1.295067in}}%
\pgfpathclose%
\pgfpathmoveto{\pgfqpoint{5.139368in}{1.316970in}}%
\pgfpathquadraticcurveto{\pgfqpoint{5.137620in}{1.317979in}}{\pgfqpoint{5.135567in}{1.318447in}}%
\pgfpathquadraticcurveto{\pgfqpoint{5.133514in}{1.318933in}}{\pgfqpoint{5.131046in}{1.318933in}}%
\pgfpathquadraticcurveto{\pgfqpoint{5.122256in}{1.318933in}}{\pgfqpoint{5.117555in}{1.313223in}}%
\pgfpathquadraticcurveto{\pgfqpoint{5.112854in}{1.307514in}}{\pgfqpoint{5.112854in}{1.296814in}}%
\pgfpathlineto{\pgfqpoint{5.112854in}{1.263600in}}%
\pgfpathlineto{\pgfqpoint{5.102443in}{1.263600in}}%
\pgfpathlineto{\pgfqpoint{5.102443in}{1.326643in}}%
\pgfpathlineto{\pgfqpoint{5.112854in}{1.326643in}}%
\pgfpathlineto{\pgfqpoint{5.112854in}{1.316844in}}%
\pgfpathquadraticcurveto{\pgfqpoint{5.116132in}{1.322590in}}{\pgfqpoint{5.121355in}{1.325364in}}%
\pgfpathquadraticcurveto{\pgfqpoint{5.126597in}{1.328156in}}{\pgfqpoint{5.134090in}{1.328156in}}%
\pgfpathquadraticcurveto{\pgfqpoint{5.135153in}{1.328156in}}{\pgfqpoint{5.136450in}{1.328011in}}%
\pgfpathquadraticcurveto{\pgfqpoint{5.137747in}{1.327885in}}{\pgfqpoint{5.139314in}{1.327597in}}%
\pgfpathlineto{\pgfqpoint{5.139368in}{1.316970in}}%
\pgfpathlineto{\pgfqpoint{5.139368in}{1.316970in}}%
\pgfpathclose%
\pgfpathmoveto{\pgfqpoint{5.201618in}{1.297715in}}%
\pgfpathlineto{\pgfqpoint{5.201618in}{1.292654in}}%
\pgfpathlineto{\pgfqpoint{5.153993in}{1.292654in}}%
\pgfpathquadraticcurveto{\pgfqpoint{5.154678in}{1.281954in}}{\pgfqpoint{5.160442in}{1.276353in}}%
\pgfpathquadraticcurveto{\pgfqpoint{5.166206in}{1.270751in}}{\pgfqpoint{5.176509in}{1.270751in}}%
\pgfpathquadraticcurveto{\pgfqpoint{5.182471in}{1.270751in}}{\pgfqpoint{5.188072in}{1.272210in}}%
\pgfpathquadraticcurveto{\pgfqpoint{5.193674in}{1.273669in}}{\pgfqpoint{5.199204in}{1.276605in}}%
\pgfpathlineto{\pgfqpoint{5.199204in}{1.266806in}}%
\pgfpathquadraticcurveto{\pgfqpoint{5.193620in}{1.264447in}}{\pgfqpoint{5.187766in}{1.263204in}}%
\pgfpathquadraticcurveto{\pgfqpoint{5.181912in}{1.261961in}}{\pgfqpoint{5.175896in}{1.261961in}}%
\pgfpathquadraticcurveto{\pgfqpoint{5.160802in}{1.261961in}}{\pgfqpoint{5.151994in}{1.270733in}}%
\pgfpathquadraticcurveto{\pgfqpoint{5.143186in}{1.279523in}}{\pgfqpoint{5.143186in}{1.294509in}}%
\pgfpathquadraticcurveto{\pgfqpoint{5.143186in}{1.309981in}}{\pgfqpoint{5.151544in}{1.319059in}}%
\pgfpathquadraticcurveto{\pgfqpoint{5.159901in}{1.328156in}}{\pgfqpoint{5.174095in}{1.328156in}}%
\pgfpathquadraticcurveto{\pgfqpoint{5.186812in}{1.328156in}}{\pgfqpoint{5.194215in}{1.319960in}}%
\pgfpathquadraticcurveto{\pgfqpoint{5.201618in}{1.311783in}}{\pgfqpoint{5.201618in}{1.297715in}}%
\pgfpathlineto{\pgfqpoint{5.201618in}{1.297715in}}%
\pgfpathclose%
\pgfpathmoveto{\pgfqpoint{5.191261in}{1.300759in}}%
\pgfpathquadraticcurveto{\pgfqpoint{5.191153in}{1.309243in}}{\pgfqpoint{5.186505in}{1.314304in}}%
\pgfpathquadraticcurveto{\pgfqpoint{5.181858in}{1.319384in}}{\pgfqpoint{5.174203in}{1.319384in}}%
\pgfpathquadraticcurveto{\pgfqpoint{5.165539in}{1.319384in}}{\pgfqpoint{5.160334in}{1.314484in}}%
\pgfpathquadraticcurveto{\pgfqpoint{5.155128in}{1.309585in}}{\pgfqpoint{5.154336in}{1.300687in}}%
\pgfpathlineto{\pgfqpoint{5.191261in}{1.300759in}}%
\pgfpathlineto{\pgfqpoint{5.191261in}{1.300759in}}%
\pgfpathclose%
\pgfusepath{fill}%
\end{pgfscope}%
\begin{pgfscope}%
\pgfsetbuttcap%
\pgfsetroundjoin%
\definecolor{currentfill}{rgb}{0.309804,0.372549,0.419608}%
\pgfsetfillcolor{currentfill}%
\pgfsetlinewidth{0.000000pt}%
\definecolor{currentstroke}{rgb}{0.309804,0.372549,0.419608}%
\pgfsetstrokecolor{currentstroke}%
\pgfsetdash{}{0pt}%
\pgfpathmoveto{\pgfqpoint{5.311498in}{1.295283in}}%
\pgfpathquadraticcurveto{\pgfqpoint{5.298943in}{1.295283in}}{\pgfqpoint{5.294098in}{1.292419in}}%
\pgfpathquadraticcurveto{\pgfqpoint{5.289253in}{1.289556in}}{\pgfqpoint{5.289253in}{1.282621in}}%
\pgfpathquadraticcurveto{\pgfqpoint{5.289253in}{1.277109in}}{\pgfqpoint{5.292891in}{1.273867in}}%
\pgfpathquadraticcurveto{\pgfqpoint{5.296530in}{1.270643in}}{\pgfqpoint{5.302762in}{1.270643in}}%
\pgfpathquadraticcurveto{\pgfqpoint{5.311390in}{1.270643in}}{\pgfqpoint{5.316595in}{1.276749in}}%
\pgfpathquadraticcurveto{\pgfqpoint{5.321801in}{1.282855in}}{\pgfqpoint{5.321801in}{1.292978in}}%
\pgfpathlineto{\pgfqpoint{5.321801in}{1.295283in}}%
\pgfpathlineto{\pgfqpoint{5.311498in}{1.295283in}}%
\pgfpathlineto{\pgfqpoint{5.311498in}{1.295283in}}%
\pgfpathclose%
\pgfpathmoveto{\pgfqpoint{5.332158in}{1.299570in}}%
\pgfpathlineto{\pgfqpoint{5.332158in}{1.263600in}}%
\pgfpathlineto{\pgfqpoint{5.321801in}{1.263600in}}%
\pgfpathlineto{\pgfqpoint{5.321801in}{1.273164in}}%
\pgfpathquadraticcurveto{\pgfqpoint{5.318252in}{1.267437in}}{\pgfqpoint{5.312957in}{1.264699in}}%
\pgfpathquadraticcurveto{\pgfqpoint{5.307661in}{1.261961in}}{\pgfqpoint{5.300006in}{1.261961in}}%
\pgfpathquadraticcurveto{\pgfqpoint{5.290334in}{1.261961in}}{\pgfqpoint{5.284606in}{1.267401in}}%
\pgfpathquadraticcurveto{\pgfqpoint{5.278896in}{1.272840in}}{\pgfqpoint{5.278896in}{1.281954in}}%
\pgfpathquadraticcurveto{\pgfqpoint{5.278896in}{1.292582in}}{\pgfqpoint{5.286011in}{1.297985in}}%
\pgfpathquadraticcurveto{\pgfqpoint{5.293143in}{1.303389in}}{\pgfqpoint{5.307265in}{1.303389in}}%
\pgfpathlineto{\pgfqpoint{5.321801in}{1.303389in}}%
\pgfpathlineto{\pgfqpoint{5.321801in}{1.304416in}}%
\pgfpathquadraticcurveto{\pgfqpoint{5.321801in}{1.311566in}}{\pgfqpoint{5.317100in}{1.315475in}}%
\pgfpathquadraticcurveto{\pgfqpoint{5.312398in}{1.319384in}}{\pgfqpoint{5.303897in}{1.319384in}}%
\pgfpathquadraticcurveto{\pgfqpoint{5.298493in}{1.319384in}}{\pgfqpoint{5.293360in}{1.318087in}}%
\pgfpathquadraticcurveto{\pgfqpoint{5.288244in}{1.316790in}}{\pgfqpoint{5.283525in}{1.314196in}}%
\pgfpathlineto{\pgfqpoint{5.283525in}{1.323779in}}%
\pgfpathquadraticcurveto{\pgfqpoint{5.289199in}{1.325976in}}{\pgfqpoint{5.294548in}{1.327057in}}%
\pgfpathquadraticcurveto{\pgfqpoint{5.299898in}{1.328156in}}{\pgfqpoint{5.304959in}{1.328156in}}%
\pgfpathquadraticcurveto{\pgfqpoint{5.318649in}{1.328156in}}{\pgfqpoint{5.325403in}{1.321059in}}%
\pgfpathquadraticcurveto{\pgfqpoint{5.332158in}{1.313980in}}{\pgfqpoint{5.332158in}{1.299570in}}%
\pgfpathlineto{\pgfqpoint{5.332158in}{1.299570in}}%
\pgfpathclose%
\pgfpathmoveto{\pgfqpoint{5.405900in}{1.301660in}}%
\pgfpathlineto{\pgfqpoint{5.405900in}{1.263600in}}%
\pgfpathlineto{\pgfqpoint{5.395543in}{1.263600in}}%
\pgfpathlineto{\pgfqpoint{5.395543in}{1.301317in}}%
\pgfpathquadraticcurveto{\pgfqpoint{5.395543in}{1.310269in}}{\pgfqpoint{5.392048in}{1.314700in}}%
\pgfpathquadraticcurveto{\pgfqpoint{5.388554in}{1.319149in}}{\pgfqpoint{5.381583in}{1.319149in}}%
\pgfpathquadraticcurveto{\pgfqpoint{5.373189in}{1.319149in}}{\pgfqpoint{5.368344in}{1.313800in}}%
\pgfpathquadraticcurveto{\pgfqpoint{5.363499in}{1.308468in}}{\pgfqpoint{5.363499in}{1.299228in}}%
\pgfpathlineto{\pgfqpoint{5.363499in}{1.263600in}}%
\pgfpathlineto{\pgfqpoint{5.353088in}{1.263600in}}%
\pgfpathlineto{\pgfqpoint{5.353088in}{1.326643in}}%
\pgfpathlineto{\pgfqpoint{5.363499in}{1.326643in}}%
\pgfpathlineto{\pgfqpoint{5.363499in}{1.316844in}}%
\pgfpathquadraticcurveto{\pgfqpoint{5.367227in}{1.322536in}}{\pgfqpoint{5.372253in}{1.325346in}}%
\pgfpathquadraticcurveto{\pgfqpoint{5.377296in}{1.328156in}}{\pgfqpoint{5.383889in}{1.328156in}}%
\pgfpathquadraticcurveto{\pgfqpoint{5.394750in}{1.328156in}}{\pgfqpoint{5.400316in}{1.321437in}}%
\pgfpathquadraticcurveto{\pgfqpoint{5.405900in}{1.314718in}}{\pgfqpoint{5.405900in}{1.301660in}}%
\pgfpathlineto{\pgfqpoint{5.405900in}{1.301660in}}%
\pgfpathclose%
\pgfpathmoveto{\pgfqpoint{5.468023in}{1.317078in}}%
\pgfpathlineto{\pgfqpoint{5.468023in}{1.351193in}}%
\pgfpathlineto{\pgfqpoint{5.478380in}{1.351193in}}%
\pgfpathlineto{\pgfqpoint{5.478380in}{1.263600in}}%
\pgfpathlineto{\pgfqpoint{5.468023in}{1.263600in}}%
\pgfpathlineto{\pgfqpoint{5.468023in}{1.273056in}}%
\pgfpathquadraticcurveto{\pgfqpoint{5.464763in}{1.267437in}}{\pgfqpoint{5.459774in}{1.264699in}}%
\pgfpathquadraticcurveto{\pgfqpoint{5.454803in}{1.261961in}}{\pgfqpoint{5.447814in}{1.261961in}}%
\pgfpathquadraticcurveto{\pgfqpoint{5.436394in}{1.261961in}}{\pgfqpoint{5.429207in}{1.271075in}}%
\pgfpathquadraticcurveto{\pgfqpoint{5.422038in}{1.280207in}}{\pgfqpoint{5.422038in}{1.295067in}}%
\pgfpathquadraticcurveto{\pgfqpoint{5.422038in}{1.309927in}}{\pgfqpoint{5.429207in}{1.319041in}}%
\pgfpathquadraticcurveto{\pgfqpoint{5.436394in}{1.328156in}}{\pgfqpoint{5.447814in}{1.328156in}}%
\pgfpathquadraticcurveto{\pgfqpoint{5.454803in}{1.328156in}}{\pgfqpoint{5.459774in}{1.325418in}}%
\pgfpathquadraticcurveto{\pgfqpoint{5.464763in}{1.322698in}}{\pgfqpoint{5.468023in}{1.317078in}}%
\pgfpathlineto{\pgfqpoint{5.468023in}{1.317078in}}%
\pgfpathclose%
\pgfpathmoveto{\pgfqpoint{5.432738in}{1.295067in}}%
\pgfpathquadraticcurveto{\pgfqpoint{5.432738in}{1.283648in}}{\pgfqpoint{5.437439in}{1.277145in}}%
\pgfpathquadraticcurveto{\pgfqpoint{5.442140in}{1.270643in}}{\pgfqpoint{5.450354in}{1.270643in}}%
\pgfpathquadraticcurveto{\pgfqpoint{5.458567in}{1.270643in}}{\pgfqpoint{5.463286in}{1.277145in}}%
\pgfpathquadraticcurveto{\pgfqpoint{5.468023in}{1.283648in}}{\pgfqpoint{5.468023in}{1.295067in}}%
\pgfpathquadraticcurveto{\pgfqpoint{5.468023in}{1.306487in}}{\pgfqpoint{5.463286in}{1.312989in}}%
\pgfpathquadraticcurveto{\pgfqpoint{5.458567in}{1.319492in}}{\pgfqpoint{5.450354in}{1.319492in}}%
\pgfpathquadraticcurveto{\pgfqpoint{5.442140in}{1.319492in}}{\pgfqpoint{5.437439in}{1.312989in}}%
\pgfpathquadraticcurveto{\pgfqpoint{5.432738in}{1.306487in}}{\pgfqpoint{5.432738in}{1.295067in}}%
\pgfpathlineto{\pgfqpoint{5.432738in}{1.295067in}}%
\pgfpathclose%
\pgfusepath{fill}%
\end{pgfscope}%
\begin{pgfscope}%
\pgfsetbuttcap%
\pgfsetroundjoin%
\definecolor{currentfill}{rgb}{0.309804,0.372549,0.419608}%
\pgfsetfillcolor{currentfill}%
\pgfsetlinewidth{0.000000pt}%
\definecolor{currentstroke}{rgb}{0.309804,0.372549,0.419608}%
\pgfsetstrokecolor{currentstroke}%
\pgfsetdash{}{0pt}%
\pgfpathmoveto{\pgfqpoint{5.578540in}{1.273056in}}%
\pgfpathlineto{\pgfqpoint{5.578540in}{1.239626in}}%
\pgfpathlineto{\pgfqpoint{5.568129in}{1.239626in}}%
\pgfpathlineto{\pgfqpoint{5.568129in}{1.326643in}}%
\pgfpathlineto{\pgfqpoint{5.578540in}{1.326643in}}%
\pgfpathlineto{\pgfqpoint{5.578540in}{1.317078in}}%
\pgfpathquadraticcurveto{\pgfqpoint{5.581818in}{1.322698in}}{\pgfqpoint{5.586790in}{1.325418in}}%
\pgfpathquadraticcurveto{\pgfqpoint{5.591779in}{1.328156in}}{\pgfqpoint{5.598696in}{1.328156in}}%
\pgfpathquadraticcurveto{\pgfqpoint{5.610187in}{1.328156in}}{\pgfqpoint{5.617356in}{1.319041in}}%
\pgfpathquadraticcurveto{\pgfqpoint{5.624543in}{1.309927in}}{\pgfqpoint{5.624543in}{1.295067in}}%
\pgfpathquadraticcurveto{\pgfqpoint{5.624543in}{1.280207in}}{\pgfqpoint{5.617356in}{1.271075in}}%
\pgfpathquadraticcurveto{\pgfqpoint{5.610187in}{1.261961in}}{\pgfqpoint{5.598696in}{1.261961in}}%
\pgfpathquadraticcurveto{\pgfqpoint{5.591779in}{1.261961in}}{\pgfqpoint{5.586790in}{1.264699in}}%
\pgfpathquadraticcurveto{\pgfqpoint{5.581818in}{1.267437in}}{\pgfqpoint{5.578540in}{1.273056in}}%
\pgfpathlineto{\pgfqpoint{5.578540in}{1.273056in}}%
\pgfpathclose%
\pgfpathmoveto{\pgfqpoint{5.613790in}{1.295067in}}%
\pgfpathquadraticcurveto{\pgfqpoint{5.613790in}{1.306487in}}{\pgfqpoint{5.609089in}{1.312989in}}%
\pgfpathquadraticcurveto{\pgfqpoint{5.604388in}{1.319492in}}{\pgfqpoint{5.596174in}{1.319492in}}%
\pgfpathquadraticcurveto{\pgfqpoint{5.587942in}{1.319492in}}{\pgfqpoint{5.583241in}{1.312989in}}%
\pgfpathquadraticcurveto{\pgfqpoint{5.578540in}{1.306487in}}{\pgfqpoint{5.578540in}{1.295067in}}%
\pgfpathquadraticcurveto{\pgfqpoint{5.578540in}{1.283648in}}{\pgfqpoint{5.583241in}{1.277145in}}%
\pgfpathquadraticcurveto{\pgfqpoint{5.587942in}{1.270643in}}{\pgfqpoint{5.596174in}{1.270643in}}%
\pgfpathquadraticcurveto{\pgfqpoint{5.604388in}{1.270643in}}{\pgfqpoint{5.609089in}{1.277145in}}%
\pgfpathquadraticcurveto{\pgfqpoint{5.613790in}{1.283648in}}{\pgfqpoint{5.613790in}{1.295067in}}%
\pgfpathlineto{\pgfqpoint{5.613790in}{1.295067in}}%
\pgfpathclose%
\pgfpathmoveto{\pgfqpoint{5.666133in}{1.319384in}}%
\pgfpathquadraticcurveto{\pgfqpoint{5.657812in}{1.319384in}}{\pgfqpoint{5.652966in}{1.312881in}}%
\pgfpathquadraticcurveto{\pgfqpoint{5.648121in}{1.306379in}}{\pgfqpoint{5.648121in}{1.295067in}}%
\pgfpathquadraticcurveto{\pgfqpoint{5.648121in}{1.283756in}}{\pgfqpoint{5.652930in}{1.277253in}}%
\pgfpathquadraticcurveto{\pgfqpoint{5.657758in}{1.270751in}}{\pgfqpoint{5.666133in}{1.270751in}}%
\pgfpathquadraticcurveto{\pgfqpoint{5.674419in}{1.270751in}}{\pgfqpoint{5.679246in}{1.277271in}}%
\pgfpathquadraticcurveto{\pgfqpoint{5.684091in}{1.283810in}}{\pgfqpoint{5.684091in}{1.295067in}}%
\pgfpathquadraticcurveto{\pgfqpoint{5.684091in}{1.306271in}}{\pgfqpoint{5.679246in}{1.312827in}}%
\pgfpathquadraticcurveto{\pgfqpoint{5.674419in}{1.319384in}}{\pgfqpoint{5.666133in}{1.319384in}}%
\pgfpathlineto{\pgfqpoint{5.666133in}{1.319384in}}%
\pgfpathclose%
\pgfpathmoveto{\pgfqpoint{5.666133in}{1.328156in}}%
\pgfpathquadraticcurveto{\pgfqpoint{5.679642in}{1.328156in}}{\pgfqpoint{5.687352in}{1.319366in}}%
\pgfpathquadraticcurveto{\pgfqpoint{5.695079in}{1.310594in}}{\pgfqpoint{5.695079in}{1.295067in}}%
\pgfpathquadraticcurveto{\pgfqpoint{5.695079in}{1.279595in}}{\pgfqpoint{5.687352in}{1.270769in}}%
\pgfpathquadraticcurveto{\pgfqpoint{5.679642in}{1.261961in}}{\pgfqpoint{5.666133in}{1.261961in}}%
\pgfpathquadraticcurveto{\pgfqpoint{5.652570in}{1.261961in}}{\pgfqpoint{5.644879in}{1.270769in}}%
\pgfpathquadraticcurveto{\pgfqpoint{5.637206in}{1.279595in}}{\pgfqpoint{5.637206in}{1.295067in}}%
\pgfpathquadraticcurveto{\pgfqpoint{5.637206in}{1.310594in}}{\pgfqpoint{5.644879in}{1.319366in}}%
\pgfpathquadraticcurveto{\pgfqpoint{5.652570in}{1.328156in}}{\pgfqpoint{5.666133in}{1.328156in}}%
\pgfpathlineto{\pgfqpoint{5.666133in}{1.328156in}}%
\pgfpathclose%
\pgfpathmoveto{\pgfqpoint{5.752429in}{1.324787in}}%
\pgfpathlineto{\pgfqpoint{5.752429in}{1.314989in}}%
\pgfpathquadraticcurveto{\pgfqpoint{5.748053in}{1.317240in}}{\pgfqpoint{5.743315in}{1.318357in}}%
\pgfpathquadraticcurveto{\pgfqpoint{5.738596in}{1.319492in}}{\pgfqpoint{5.733517in}{1.319492in}}%
\pgfpathquadraticcurveto{\pgfqpoint{5.725808in}{1.319492in}}{\pgfqpoint{5.721953in}{1.317132in}}%
\pgfpathquadraticcurveto{\pgfqpoint{5.718098in}{1.314773in}}{\pgfqpoint{5.718098in}{1.310035in}}%
\pgfpathquadraticcurveto{\pgfqpoint{5.718098in}{1.306433in}}{\pgfqpoint{5.720854in}{1.304380in}}%
\pgfpathquadraticcurveto{\pgfqpoint{5.723610in}{1.302326in}}{\pgfqpoint{5.731950in}{1.300471in}}%
\pgfpathlineto{\pgfqpoint{5.735498in}{1.299678in}}%
\pgfpathquadraticcurveto{\pgfqpoint{5.746521in}{1.297319in}}{\pgfqpoint{5.751169in}{1.293014in}}%
\pgfpathquadraticcurveto{\pgfqpoint{5.755816in}{1.288709in}}{\pgfqpoint{5.755816in}{1.281000in}}%
\pgfpathquadraticcurveto{\pgfqpoint{5.755816in}{1.272210in}}{\pgfqpoint{5.748863in}{1.267076in}}%
\pgfpathquadraticcurveto{\pgfqpoint{5.741910in}{1.261961in}}{\pgfqpoint{5.729752in}{1.261961in}}%
\pgfpathquadraticcurveto{\pgfqpoint{5.724691in}{1.261961in}}{\pgfqpoint{5.719197in}{1.262952in}}%
\pgfpathquadraticcurveto{\pgfqpoint{5.713703in}{1.263942in}}{\pgfqpoint{5.707633in}{1.265906in}}%
\pgfpathlineto{\pgfqpoint{5.707633in}{1.276605in}}%
\pgfpathquadraticcurveto{\pgfqpoint{5.713379in}{1.273615in}}{\pgfqpoint{5.718945in}{1.272120in}}%
\pgfpathquadraticcurveto{\pgfqpoint{5.724511in}{1.270643in}}{\pgfqpoint{5.729986in}{1.270643in}}%
\pgfpathquadraticcurveto{\pgfqpoint{5.737299in}{1.270643in}}{\pgfqpoint{5.741226in}{1.273146in}}%
\pgfpathquadraticcurveto{\pgfqpoint{5.745171in}{1.275650in}}{\pgfqpoint{5.745171in}{1.280207in}}%
\pgfpathquadraticcurveto{\pgfqpoint{5.745171in}{1.284422in}}{\pgfqpoint{5.742325in}{1.286674in}}%
\pgfpathquadraticcurveto{\pgfqpoint{5.739497in}{1.288925in}}{\pgfqpoint{5.729860in}{1.291014in}}%
\pgfpathlineto{\pgfqpoint{5.726258in}{1.291861in}}%
\pgfpathquadraticcurveto{\pgfqpoint{5.716639in}{1.293878in}}{\pgfqpoint{5.712352in}{1.298075in}}%
\pgfpathquadraticcurveto{\pgfqpoint{5.708084in}{1.302272in}}{\pgfqpoint{5.708084in}{1.309585in}}%
\pgfpathquadraticcurveto{\pgfqpoint{5.708084in}{1.318483in}}{\pgfqpoint{5.714388in}{1.323310in}}%
\pgfpathquadraticcurveto{\pgfqpoint{5.720692in}{1.328156in}}{\pgfqpoint{5.732292in}{1.328156in}}%
\pgfpathquadraticcurveto{\pgfqpoint{5.738020in}{1.328156in}}{\pgfqpoint{5.743081in}{1.327309in}}%
\pgfpathquadraticcurveto{\pgfqpoint{5.748161in}{1.326480in}}{\pgfqpoint{5.752429in}{1.324787in}}%
\pgfpathlineto{\pgfqpoint{5.752429in}{1.324787in}}%
\pgfpathclose%
\pgfpathmoveto{\pgfqpoint{5.782546in}{1.344547in}}%
\pgfpathlineto{\pgfqpoint{5.782546in}{1.326643in}}%
\pgfpathlineto{\pgfqpoint{5.803872in}{1.326643in}}%
\pgfpathlineto{\pgfqpoint{5.803872in}{1.318591in}}%
\pgfpathlineto{\pgfqpoint{5.782546in}{1.318591in}}%
\pgfpathlineto{\pgfqpoint{5.782546in}{1.284368in}}%
\pgfpathquadraticcurveto{\pgfqpoint{5.782546in}{1.276659in}}{\pgfqpoint{5.784653in}{1.274461in}}%
\pgfpathquadraticcurveto{\pgfqpoint{5.786761in}{1.272264in}}{\pgfqpoint{5.793245in}{1.272264in}}%
\pgfpathlineto{\pgfqpoint{5.803872in}{1.272264in}}%
\pgfpathlineto{\pgfqpoint{5.803872in}{1.263600in}}%
\pgfpathlineto{\pgfqpoint{5.793245in}{1.263600in}}%
\pgfpathquadraticcurveto{\pgfqpoint{5.781249in}{1.263600in}}{\pgfqpoint{5.776692in}{1.268067in}}%
\pgfpathquadraticcurveto{\pgfqpoint{5.772135in}{1.272552in}}{\pgfqpoint{5.772135in}{1.284368in}}%
\pgfpathlineto{\pgfqpoint{5.772135in}{1.318591in}}%
\pgfpathlineto{\pgfqpoint{5.764534in}{1.318591in}}%
\pgfpathlineto{\pgfqpoint{5.764534in}{1.326643in}}%
\pgfpathlineto{\pgfqpoint{5.772135in}{1.326643in}}%
\pgfpathlineto{\pgfqpoint{5.772135in}{1.344547in}}%
\pgfpathlineto{\pgfqpoint{5.782546in}{1.344547in}}%
\pgfpathlineto{\pgfqpoint{5.782546in}{1.344547in}}%
\pgfpathclose%
\pgfusepath{fill}%
\end{pgfscope}%
\begin{pgfscope}%
\pgfsetbuttcap%
\pgfsetroundjoin%
\definecolor{currentfill}{rgb}{0.309804,0.372549,0.419608}%
\pgfsetfillcolor{currentfill}%
\pgfsetlinewidth{0.000000pt}%
\definecolor{currentstroke}{rgb}{0.309804,0.372549,0.419608}%
\pgfsetstrokecolor{currentstroke}%
\pgfsetdash{}{0pt}%
\pgfpathmoveto{\pgfqpoint{5.938139in}{1.297715in}}%
\pgfpathlineto{\pgfqpoint{5.938139in}{1.292654in}}%
\pgfpathlineto{\pgfqpoint{5.890515in}{1.292654in}}%
\pgfpathquadraticcurveto{\pgfqpoint{5.891199in}{1.281954in}}{\pgfqpoint{5.896963in}{1.276353in}}%
\pgfpathquadraticcurveto{\pgfqpoint{5.902727in}{1.270751in}}{\pgfqpoint{5.913030in}{1.270751in}}%
\pgfpathquadraticcurveto{\pgfqpoint{5.918992in}{1.270751in}}{\pgfqpoint{5.924594in}{1.272210in}}%
\pgfpathquadraticcurveto{\pgfqpoint{5.930195in}{1.273669in}}{\pgfqpoint{5.935725in}{1.276605in}}%
\pgfpathlineto{\pgfqpoint{5.935725in}{1.266806in}}%
\pgfpathquadraticcurveto{\pgfqpoint{5.930141in}{1.264447in}}{\pgfqpoint{5.924287in}{1.263204in}}%
\pgfpathquadraticcurveto{\pgfqpoint{5.918433in}{1.261961in}}{\pgfqpoint{5.912417in}{1.261961in}}%
\pgfpathquadraticcurveto{\pgfqpoint{5.897323in}{1.261961in}}{\pgfqpoint{5.888515in}{1.270733in}}%
\pgfpathquadraticcurveto{\pgfqpoint{5.879707in}{1.279523in}}{\pgfqpoint{5.879707in}{1.294509in}}%
\pgfpathquadraticcurveto{\pgfqpoint{5.879707in}{1.309981in}}{\pgfqpoint{5.888065in}{1.319059in}}%
\pgfpathquadraticcurveto{\pgfqpoint{5.896423in}{1.328156in}}{\pgfqpoint{5.910616in}{1.328156in}}%
\pgfpathquadraticcurveto{\pgfqpoint{5.923333in}{1.328156in}}{\pgfqpoint{5.930736in}{1.319960in}}%
\pgfpathquadraticcurveto{\pgfqpoint{5.938139in}{1.311783in}}{\pgfqpoint{5.938139in}{1.297715in}}%
\pgfpathlineto{\pgfqpoint{5.938139in}{1.297715in}}%
\pgfpathclose%
\pgfpathmoveto{\pgfqpoint{5.927782in}{1.300759in}}%
\pgfpathquadraticcurveto{\pgfqpoint{5.927674in}{1.309243in}}{\pgfqpoint{5.923026in}{1.314304in}}%
\pgfpathquadraticcurveto{\pgfqpoint{5.918379in}{1.319384in}}{\pgfqpoint{5.910724in}{1.319384in}}%
\pgfpathquadraticcurveto{\pgfqpoint{5.902060in}{1.319384in}}{\pgfqpoint{5.896855in}{1.314484in}}%
\pgfpathquadraticcurveto{\pgfqpoint{5.891649in}{1.309585in}}{\pgfqpoint{5.890857in}{1.300687in}}%
\pgfpathlineto{\pgfqpoint{5.927782in}{1.300759in}}%
\pgfpathlineto{\pgfqpoint{5.927782in}{1.300759in}}%
\pgfpathclose%
\pgfpathmoveto{\pgfqpoint{6.007558in}{1.301660in}}%
\pgfpathlineto{\pgfqpoint{6.007558in}{1.263600in}}%
\pgfpathlineto{\pgfqpoint{5.997201in}{1.263600in}}%
\pgfpathlineto{\pgfqpoint{5.997201in}{1.301317in}}%
\pgfpathquadraticcurveto{\pgfqpoint{5.997201in}{1.310269in}}{\pgfqpoint{5.993706in}{1.314700in}}%
\pgfpathquadraticcurveto{\pgfqpoint{5.990212in}{1.319149in}}{\pgfqpoint{5.983241in}{1.319149in}}%
\pgfpathquadraticcurveto{\pgfqpoint{5.974847in}{1.319149in}}{\pgfqpoint{5.970002in}{1.313800in}}%
\pgfpathquadraticcurveto{\pgfqpoint{5.965157in}{1.308468in}}{\pgfqpoint{5.965157in}{1.299228in}}%
\pgfpathlineto{\pgfqpoint{5.965157in}{1.263600in}}%
\pgfpathlineto{\pgfqpoint{5.954746in}{1.263600in}}%
\pgfpathlineto{\pgfqpoint{5.954746in}{1.326643in}}%
\pgfpathlineto{\pgfqpoint{5.965157in}{1.326643in}}%
\pgfpathlineto{\pgfqpoint{5.965157in}{1.316844in}}%
\pgfpathquadraticcurveto{\pgfqpoint{5.968885in}{1.322536in}}{\pgfqpoint{5.973911in}{1.325346in}}%
\pgfpathquadraticcurveto{\pgfqpoint{5.978954in}{1.328156in}}{\pgfqpoint{5.985547in}{1.328156in}}%
\pgfpathquadraticcurveto{\pgfqpoint{5.996408in}{1.328156in}}{\pgfqpoint{6.001974in}{1.321437in}}%
\pgfpathquadraticcurveto{\pgfqpoint{6.007558in}{1.314718in}}{\pgfqpoint{6.007558in}{1.301660in}}%
\pgfpathlineto{\pgfqpoint{6.007558in}{1.301660in}}%
\pgfpathclose%
\pgfpathmoveto{\pgfqpoint{6.038448in}{1.344547in}}%
\pgfpathlineto{\pgfqpoint{6.038448in}{1.326643in}}%
\pgfpathlineto{\pgfqpoint{6.059775in}{1.326643in}}%
\pgfpathlineto{\pgfqpoint{6.059775in}{1.318591in}}%
\pgfpathlineto{\pgfqpoint{6.038448in}{1.318591in}}%
\pgfpathlineto{\pgfqpoint{6.038448in}{1.284368in}}%
\pgfpathquadraticcurveto{\pgfqpoint{6.038448in}{1.276659in}}{\pgfqpoint{6.040556in}{1.274461in}}%
\pgfpathquadraticcurveto{\pgfqpoint{6.042663in}{1.272264in}}{\pgfqpoint{6.049148in}{1.272264in}}%
\pgfpathlineto{\pgfqpoint{6.059775in}{1.272264in}}%
\pgfpathlineto{\pgfqpoint{6.059775in}{1.263600in}}%
\pgfpathlineto{\pgfqpoint{6.049148in}{1.263600in}}%
\pgfpathquadraticcurveto{\pgfqpoint{6.037151in}{1.263600in}}{\pgfqpoint{6.032594in}{1.268067in}}%
\pgfpathquadraticcurveto{\pgfqpoint{6.028037in}{1.272552in}}{\pgfqpoint{6.028037in}{1.284368in}}%
\pgfpathlineto{\pgfqpoint{6.028037in}{1.318591in}}%
\pgfpathlineto{\pgfqpoint{6.020436in}{1.318591in}}%
\pgfpathlineto{\pgfqpoint{6.020436in}{1.326643in}}%
\pgfpathlineto{\pgfqpoint{6.028037in}{1.326643in}}%
\pgfpathlineto{\pgfqpoint{6.028037in}{1.344547in}}%
\pgfpathlineto{\pgfqpoint{6.038448in}{1.344547in}}%
\pgfpathlineto{\pgfqpoint{6.038448in}{1.344547in}}%
\pgfpathclose%
\pgfpathmoveto{\pgfqpoint{6.127320in}{1.297715in}}%
\pgfpathlineto{\pgfqpoint{6.127320in}{1.292654in}}%
\pgfpathlineto{\pgfqpoint{6.079696in}{1.292654in}}%
\pgfpathquadraticcurveto{\pgfqpoint{6.080381in}{1.281954in}}{\pgfqpoint{6.086145in}{1.276353in}}%
\pgfpathquadraticcurveto{\pgfqpoint{6.091908in}{1.270751in}}{\pgfqpoint{6.102211in}{1.270751in}}%
\pgfpathquadraticcurveto{\pgfqpoint{6.108173in}{1.270751in}}{\pgfqpoint{6.113775in}{1.272210in}}%
\pgfpathquadraticcurveto{\pgfqpoint{6.119377in}{1.273669in}}{\pgfqpoint{6.124907in}{1.276605in}}%
\pgfpathlineto{\pgfqpoint{6.124907in}{1.266806in}}%
\pgfpathquadraticcurveto{\pgfqpoint{6.119323in}{1.264447in}}{\pgfqpoint{6.113469in}{1.263204in}}%
\pgfpathquadraticcurveto{\pgfqpoint{6.107615in}{1.261961in}}{\pgfqpoint{6.101599in}{1.261961in}}%
\pgfpathquadraticcurveto{\pgfqpoint{6.086505in}{1.261961in}}{\pgfqpoint{6.077697in}{1.270733in}}%
\pgfpathquadraticcurveto{\pgfqpoint{6.068889in}{1.279523in}}{\pgfqpoint{6.068889in}{1.294509in}}%
\pgfpathquadraticcurveto{\pgfqpoint{6.068889in}{1.309981in}}{\pgfqpoint{6.077247in}{1.319059in}}%
\pgfpathquadraticcurveto{\pgfqpoint{6.085604in}{1.328156in}}{\pgfqpoint{6.099798in}{1.328156in}}%
\pgfpathquadraticcurveto{\pgfqpoint{6.112514in}{1.328156in}}{\pgfqpoint{6.119917in}{1.319960in}}%
\pgfpathquadraticcurveto{\pgfqpoint{6.127320in}{1.311783in}}{\pgfqpoint{6.127320in}{1.297715in}}%
\pgfpathlineto{\pgfqpoint{6.127320in}{1.297715in}}%
\pgfpathclose%
\pgfpathmoveto{\pgfqpoint{6.116963in}{1.300759in}}%
\pgfpathquadraticcurveto{\pgfqpoint{6.116855in}{1.309243in}}{\pgfqpoint{6.112208in}{1.314304in}}%
\pgfpathquadraticcurveto{\pgfqpoint{6.107561in}{1.319384in}}{\pgfqpoint{6.099906in}{1.319384in}}%
\pgfpathquadraticcurveto{\pgfqpoint{6.091242in}{1.319384in}}{\pgfqpoint{6.086036in}{1.314484in}}%
\pgfpathquadraticcurveto{\pgfqpoint{6.080831in}{1.309585in}}{\pgfqpoint{6.080038in}{1.300687in}}%
\pgfpathlineto{\pgfqpoint{6.116963in}{1.300759in}}%
\pgfpathlineto{\pgfqpoint{6.116963in}{1.300759in}}%
\pgfpathclose%
\pgfpathmoveto{\pgfqpoint{6.180852in}{1.316970in}}%
\pgfpathquadraticcurveto{\pgfqpoint{6.179105in}{1.317979in}}{\pgfqpoint{6.177052in}{1.318447in}}%
\pgfpathquadraticcurveto{\pgfqpoint{6.174998in}{1.318933in}}{\pgfqpoint{6.172531in}{1.318933in}}%
\pgfpathquadraticcurveto{\pgfqpoint{6.163741in}{1.318933in}}{\pgfqpoint{6.159040in}{1.313223in}}%
\pgfpathquadraticcurveto{\pgfqpoint{6.154339in}{1.307514in}}{\pgfqpoint{6.154339in}{1.296814in}}%
\pgfpathlineto{\pgfqpoint{6.154339in}{1.263600in}}%
\pgfpathlineto{\pgfqpoint{6.143928in}{1.263600in}}%
\pgfpathlineto{\pgfqpoint{6.143928in}{1.326643in}}%
\pgfpathlineto{\pgfqpoint{6.154339in}{1.326643in}}%
\pgfpathlineto{\pgfqpoint{6.154339in}{1.316844in}}%
\pgfpathquadraticcurveto{\pgfqpoint{6.157617in}{1.322590in}}{\pgfqpoint{6.162840in}{1.325364in}}%
\pgfpathquadraticcurveto{\pgfqpoint{6.168082in}{1.328156in}}{\pgfqpoint{6.175575in}{1.328156in}}%
\pgfpathquadraticcurveto{\pgfqpoint{6.176638in}{1.328156in}}{\pgfqpoint{6.177934in}{1.328011in}}%
\pgfpathquadraticcurveto{\pgfqpoint{6.179231in}{1.327885in}}{\pgfqpoint{6.180798in}{1.327597in}}%
\pgfpathlineto{\pgfqpoint{6.180852in}{1.316970in}}%
\pgfpathlineto{\pgfqpoint{6.180852in}{1.316970in}}%
\pgfpathclose%
\pgfpathmoveto{\pgfqpoint{6.191714in}{1.326643in}}%
\pgfpathlineto{\pgfqpoint{6.202071in}{1.326643in}}%
\pgfpathlineto{\pgfqpoint{6.202071in}{1.263600in}}%
\pgfpathlineto{\pgfqpoint{6.191714in}{1.263600in}}%
\pgfpathlineto{\pgfqpoint{6.191714in}{1.326643in}}%
\pgfpathlineto{\pgfqpoint{6.191714in}{1.326643in}}%
\pgfpathclose%
\pgfpathmoveto{\pgfqpoint{6.191714in}{1.351193in}}%
\pgfpathlineto{\pgfqpoint{6.202071in}{1.351193in}}%
\pgfpathlineto{\pgfqpoint{6.202071in}{1.338062in}}%
\pgfpathlineto{\pgfqpoint{6.191714in}{1.338062in}}%
\pgfpathlineto{\pgfqpoint{6.191714in}{1.351193in}}%
\pgfpathlineto{\pgfqpoint{6.191714in}{1.351193in}}%
\pgfpathclose%
\pgfpathmoveto{\pgfqpoint{6.276155in}{1.301660in}}%
\pgfpathlineto{\pgfqpoint{6.276155in}{1.263600in}}%
\pgfpathlineto{\pgfqpoint{6.265798in}{1.263600in}}%
\pgfpathlineto{\pgfqpoint{6.265798in}{1.301317in}}%
\pgfpathquadraticcurveto{\pgfqpoint{6.265798in}{1.310269in}}{\pgfqpoint{6.262303in}{1.314700in}}%
\pgfpathquadraticcurveto{\pgfqpoint{6.258809in}{1.319149in}}{\pgfqpoint{6.251838in}{1.319149in}}%
\pgfpathquadraticcurveto{\pgfqpoint{6.243445in}{1.319149in}}{\pgfqpoint{6.238599in}{1.313800in}}%
\pgfpathquadraticcurveto{\pgfqpoint{6.233754in}{1.308468in}}{\pgfqpoint{6.233754in}{1.299228in}}%
\pgfpathlineto{\pgfqpoint{6.233754in}{1.263600in}}%
\pgfpathlineto{\pgfqpoint{6.223343in}{1.263600in}}%
\pgfpathlineto{\pgfqpoint{6.223343in}{1.326643in}}%
\pgfpathlineto{\pgfqpoint{6.233754in}{1.326643in}}%
\pgfpathlineto{\pgfqpoint{6.233754in}{1.316844in}}%
\pgfpathquadraticcurveto{\pgfqpoint{6.237483in}{1.322536in}}{\pgfqpoint{6.242508in}{1.325346in}}%
\pgfpathquadraticcurveto{\pgfqpoint{6.247551in}{1.328156in}}{\pgfqpoint{6.254144in}{1.328156in}}%
\pgfpathquadraticcurveto{\pgfqpoint{6.265005in}{1.328156in}}{\pgfqpoint{6.270571in}{1.321437in}}%
\pgfpathquadraticcurveto{\pgfqpoint{6.276155in}{1.314718in}}{\pgfqpoint{6.276155in}{1.301660in}}%
\pgfpathlineto{\pgfqpoint{6.276155in}{1.301660in}}%
\pgfpathclose%
\pgfpathmoveto{\pgfqpoint{6.338279in}{1.295860in}}%
\pgfpathquadraticcurveto{\pgfqpoint{6.338279in}{1.307117in}}{\pgfqpoint{6.333632in}{1.313296in}}%
\pgfpathquadraticcurveto{\pgfqpoint{6.329002in}{1.319492in}}{\pgfqpoint{6.320609in}{1.319492in}}%
\pgfpathquadraticcurveto{\pgfqpoint{6.312287in}{1.319492in}}{\pgfqpoint{6.307640in}{1.313296in}}%
\pgfpathquadraticcurveto{\pgfqpoint{6.302993in}{1.307117in}}{\pgfqpoint{6.302993in}{1.295860in}}%
\pgfpathquadraticcurveto{\pgfqpoint{6.302993in}{1.284656in}}{\pgfqpoint{6.307640in}{1.278460in}}%
\pgfpathquadraticcurveto{\pgfqpoint{6.312287in}{1.272264in}}{\pgfqpoint{6.320609in}{1.272264in}}%
\pgfpathquadraticcurveto{\pgfqpoint{6.329002in}{1.272264in}}{\pgfqpoint{6.333632in}{1.278460in}}%
\pgfpathquadraticcurveto{\pgfqpoint{6.338279in}{1.284656in}}{\pgfqpoint{6.338279in}{1.295860in}}%
\pgfpathlineto{\pgfqpoint{6.338279in}{1.295860in}}%
\pgfpathclose%
\pgfpathmoveto{\pgfqpoint{6.348636in}{1.271417in}}%
\pgfpathquadraticcurveto{\pgfqpoint{6.348636in}{1.255332in}}{\pgfqpoint{6.341485in}{1.247479in}}%
\pgfpathquadraticcurveto{\pgfqpoint{6.334352in}{1.239626in}}{\pgfqpoint{6.319600in}{1.239626in}}%
\pgfpathquadraticcurveto{\pgfqpoint{6.314142in}{1.239626in}}{\pgfqpoint{6.309297in}{1.240436in}}%
\pgfpathquadraticcurveto{\pgfqpoint{6.304452in}{1.241247in}}{\pgfqpoint{6.299895in}{1.242940in}}%
\pgfpathlineto{\pgfqpoint{6.299895in}{1.253009in}}%
\pgfpathquadraticcurveto{\pgfqpoint{6.304452in}{1.250541in}}{\pgfqpoint{6.308901in}{1.249370in}}%
\pgfpathquadraticcurveto{\pgfqpoint{6.313350in}{1.248182in}}{\pgfqpoint{6.317961in}{1.248182in}}%
\pgfpathquadraticcurveto{\pgfqpoint{6.328156in}{1.248182in}}{\pgfqpoint{6.333217in}{1.253495in}}%
\pgfpathquadraticcurveto{\pgfqpoint{6.338279in}{1.258809in}}{\pgfqpoint{6.338279in}{1.269562in}}%
\pgfpathlineto{\pgfqpoint{6.338279in}{1.274695in}}%
\pgfpathquadraticcurveto{\pgfqpoint{6.335072in}{1.269112in}}{\pgfqpoint{6.330065in}{1.266356in}}%
\pgfpathquadraticcurveto{\pgfqpoint{6.325058in}{1.263600in}}{\pgfqpoint{6.318069in}{1.263600in}}%
\pgfpathquadraticcurveto{\pgfqpoint{6.306487in}{1.263600in}}{\pgfqpoint{6.299390in}{1.272426in}}%
\pgfpathquadraticcurveto{\pgfqpoint{6.292294in}{1.281270in}}{\pgfqpoint{6.292294in}{1.295860in}}%
\pgfpathquadraticcurveto{\pgfqpoint{6.292294in}{1.310486in}}{\pgfqpoint{6.299390in}{1.319312in}}%
\pgfpathquadraticcurveto{\pgfqpoint{6.306487in}{1.328156in}}{\pgfqpoint{6.318069in}{1.328156in}}%
\pgfpathquadraticcurveto{\pgfqpoint{6.325058in}{1.328156in}}{\pgfqpoint{6.330065in}{1.325400in}}%
\pgfpathquadraticcurveto{\pgfqpoint{6.335072in}{1.322644in}}{\pgfqpoint{6.338279in}{1.317078in}}%
\pgfpathlineto{\pgfqpoint{6.338279in}{1.326643in}}%
\pgfpathlineto{\pgfqpoint{6.348636in}{1.326643in}}%
\pgfpathlineto{\pgfqpoint{6.348636in}{1.271417in}}%
\pgfpathlineto{\pgfqpoint{6.348636in}{1.271417in}}%
\pgfpathclose%
\pgfusepath{fill}%
\end{pgfscope}%
\begin{pgfscope}%
\pgfsetbuttcap%
\pgfsetroundjoin%
\definecolor{currentfill}{rgb}{0.309804,0.372549,0.419608}%
\pgfsetfillcolor{currentfill}%
\pgfsetlinewidth{0.000000pt}%
\definecolor{currentstroke}{rgb}{0.309804,0.372549,0.419608}%
\pgfsetstrokecolor{currentstroke}%
\pgfsetdash{}{0pt}%
\pgfpathmoveto{\pgfqpoint{6.490080in}{1.324787in}}%
\pgfpathlineto{\pgfqpoint{6.490080in}{1.314989in}}%
\pgfpathquadraticcurveto{\pgfqpoint{6.485703in}{1.317240in}}{\pgfqpoint{6.480966in}{1.318357in}}%
\pgfpathquadraticcurveto{\pgfqpoint{6.476247in}{1.319492in}}{\pgfqpoint{6.471168in}{1.319492in}}%
\pgfpathquadraticcurveto{\pgfqpoint{6.463458in}{1.319492in}}{\pgfqpoint{6.459604in}{1.317132in}}%
\pgfpathquadraticcurveto{\pgfqpoint{6.455749in}{1.314773in}}{\pgfqpoint{6.455749in}{1.310035in}}%
\pgfpathquadraticcurveto{\pgfqpoint{6.455749in}{1.306433in}}{\pgfqpoint{6.458505in}{1.304380in}}%
\pgfpathquadraticcurveto{\pgfqpoint{6.461261in}{1.302326in}}{\pgfqpoint{6.469600in}{1.300471in}}%
\pgfpathlineto{\pgfqpoint{6.473149in}{1.299678in}}%
\pgfpathquadraticcurveto{\pgfqpoint{6.484172in}{1.297319in}}{\pgfqpoint{6.488819in}{1.293014in}}%
\pgfpathquadraticcurveto{\pgfqpoint{6.493467in}{1.288709in}}{\pgfqpoint{6.493467in}{1.281000in}}%
\pgfpathquadraticcurveto{\pgfqpoint{6.493467in}{1.272210in}}{\pgfqpoint{6.486514in}{1.267076in}}%
\pgfpathquadraticcurveto{\pgfqpoint{6.479561in}{1.261961in}}{\pgfqpoint{6.467403in}{1.261961in}}%
\pgfpathquadraticcurveto{\pgfqpoint{6.462342in}{1.261961in}}{\pgfqpoint{6.456848in}{1.262952in}}%
\pgfpathquadraticcurveto{\pgfqpoint{6.451354in}{1.263942in}}{\pgfqpoint{6.445284in}{1.265906in}}%
\pgfpathlineto{\pgfqpoint{6.445284in}{1.276605in}}%
\pgfpathquadraticcurveto{\pgfqpoint{6.451030in}{1.273615in}}{\pgfqpoint{6.456596in}{1.272120in}}%
\pgfpathquadraticcurveto{\pgfqpoint{6.462161in}{1.270643in}}{\pgfqpoint{6.467637in}{1.270643in}}%
\pgfpathquadraticcurveto{\pgfqpoint{6.474950in}{1.270643in}}{\pgfqpoint{6.478877in}{1.273146in}}%
\pgfpathquadraticcurveto{\pgfqpoint{6.482821in}{1.275650in}}{\pgfqpoint{6.482821in}{1.280207in}}%
\pgfpathquadraticcurveto{\pgfqpoint{6.482821in}{1.284422in}}{\pgfqpoint{6.479975in}{1.286674in}}%
\pgfpathquadraticcurveto{\pgfqpoint{6.477148in}{1.288925in}}{\pgfqpoint{6.467511in}{1.291014in}}%
\pgfpathlineto{\pgfqpoint{6.463909in}{1.291861in}}%
\pgfpathquadraticcurveto{\pgfqpoint{6.454290in}{1.293878in}}{\pgfqpoint{6.450003in}{1.298075in}}%
\pgfpathquadraticcurveto{\pgfqpoint{6.445734in}{1.302272in}}{\pgfqpoint{6.445734in}{1.309585in}}%
\pgfpathquadraticcurveto{\pgfqpoint{6.445734in}{1.318483in}}{\pgfqpoint{6.452039in}{1.323310in}}%
\pgfpathquadraticcurveto{\pgfqpoint{6.458343in}{1.328156in}}{\pgfqpoint{6.469943in}{1.328156in}}%
\pgfpathquadraticcurveto{\pgfqpoint{6.475671in}{1.328156in}}{\pgfqpoint{6.480732in}{1.327309in}}%
\pgfpathquadraticcurveto{\pgfqpoint{6.485811in}{1.326480in}}{\pgfqpoint{6.490080in}{1.324787in}}%
\pgfpathlineto{\pgfqpoint{6.490080in}{1.324787in}}%
\pgfpathclose%
\pgfpathmoveto{\pgfqpoint{6.563876in}{1.297715in}}%
\pgfpathlineto{\pgfqpoint{6.563876in}{1.292654in}}%
\pgfpathlineto{\pgfqpoint{6.516252in}{1.292654in}}%
\pgfpathquadraticcurveto{\pgfqpoint{6.516936in}{1.281954in}}{\pgfqpoint{6.522700in}{1.276353in}}%
\pgfpathquadraticcurveto{\pgfqpoint{6.528464in}{1.270751in}}{\pgfqpoint{6.538767in}{1.270751in}}%
\pgfpathquadraticcurveto{\pgfqpoint{6.544729in}{1.270751in}}{\pgfqpoint{6.550331in}{1.272210in}}%
\pgfpathquadraticcurveto{\pgfqpoint{6.555933in}{1.273669in}}{\pgfqpoint{6.561462in}{1.276605in}}%
\pgfpathlineto{\pgfqpoint{6.561462in}{1.266806in}}%
\pgfpathquadraticcurveto{\pgfqpoint{6.555879in}{1.264447in}}{\pgfqpoint{6.550025in}{1.263204in}}%
\pgfpathquadraticcurveto{\pgfqpoint{6.544171in}{1.261961in}}{\pgfqpoint{6.538155in}{1.261961in}}%
\pgfpathquadraticcurveto{\pgfqpoint{6.523061in}{1.261961in}}{\pgfqpoint{6.514253in}{1.270733in}}%
\pgfpathquadraticcurveto{\pgfqpoint{6.505445in}{1.279523in}}{\pgfqpoint{6.505445in}{1.294509in}}%
\pgfpathquadraticcurveto{\pgfqpoint{6.505445in}{1.309981in}}{\pgfqpoint{6.513802in}{1.319059in}}%
\pgfpathquadraticcurveto{\pgfqpoint{6.522160in}{1.328156in}}{\pgfqpoint{6.536354in}{1.328156in}}%
\pgfpathquadraticcurveto{\pgfqpoint{6.549070in}{1.328156in}}{\pgfqpoint{6.556473in}{1.319960in}}%
\pgfpathquadraticcurveto{\pgfqpoint{6.563876in}{1.311783in}}{\pgfqpoint{6.563876in}{1.297715in}}%
\pgfpathlineto{\pgfqpoint{6.563876in}{1.297715in}}%
\pgfpathclose%
\pgfpathmoveto{\pgfqpoint{6.553519in}{1.300759in}}%
\pgfpathquadraticcurveto{\pgfqpoint{6.553411in}{1.309243in}}{\pgfqpoint{6.548764in}{1.314304in}}%
\pgfpathquadraticcurveto{\pgfqpoint{6.544117in}{1.319384in}}{\pgfqpoint{6.536462in}{1.319384in}}%
\pgfpathquadraticcurveto{\pgfqpoint{6.527798in}{1.319384in}}{\pgfqpoint{6.522592in}{1.314484in}}%
\pgfpathquadraticcurveto{\pgfqpoint{6.517387in}{1.309585in}}{\pgfqpoint{6.516594in}{1.300687in}}%
\pgfpathlineto{\pgfqpoint{6.553519in}{1.300759in}}%
\pgfpathlineto{\pgfqpoint{6.553519in}{1.300759in}}%
\pgfpathclose%
\pgfpathmoveto{\pgfqpoint{6.590894in}{1.273056in}}%
\pgfpathlineto{\pgfqpoint{6.590894in}{1.239626in}}%
\pgfpathlineto{\pgfqpoint{6.580483in}{1.239626in}}%
\pgfpathlineto{\pgfqpoint{6.580483in}{1.326643in}}%
\pgfpathlineto{\pgfqpoint{6.590894in}{1.326643in}}%
\pgfpathlineto{\pgfqpoint{6.590894in}{1.317078in}}%
\pgfpathquadraticcurveto{\pgfqpoint{6.594173in}{1.322698in}}{\pgfqpoint{6.599144in}{1.325418in}}%
\pgfpathquadraticcurveto{\pgfqpoint{6.604133in}{1.328156in}}{\pgfqpoint{6.611050in}{1.328156in}}%
\pgfpathquadraticcurveto{\pgfqpoint{6.622542in}{1.328156in}}{\pgfqpoint{6.629711in}{1.319041in}}%
\pgfpathquadraticcurveto{\pgfqpoint{6.636897in}{1.309927in}}{\pgfqpoint{6.636897in}{1.295067in}}%
\pgfpathquadraticcurveto{\pgfqpoint{6.636897in}{1.280207in}}{\pgfqpoint{6.629711in}{1.271075in}}%
\pgfpathquadraticcurveto{\pgfqpoint{6.622542in}{1.261961in}}{\pgfqpoint{6.611050in}{1.261961in}}%
\pgfpathquadraticcurveto{\pgfqpoint{6.604133in}{1.261961in}}{\pgfqpoint{6.599144in}{1.264699in}}%
\pgfpathquadraticcurveto{\pgfqpoint{6.594173in}{1.267437in}}{\pgfqpoint{6.590894in}{1.273056in}}%
\pgfpathlineto{\pgfqpoint{6.590894in}{1.273056in}}%
\pgfpathclose%
\pgfpathmoveto{\pgfqpoint{6.626144in}{1.295067in}}%
\pgfpathquadraticcurveto{\pgfqpoint{6.626144in}{1.306487in}}{\pgfqpoint{6.621443in}{1.312989in}}%
\pgfpathquadraticcurveto{\pgfqpoint{6.616742in}{1.319492in}}{\pgfqpoint{6.608528in}{1.319492in}}%
\pgfpathquadraticcurveto{\pgfqpoint{6.600297in}{1.319492in}}{\pgfqpoint{6.595595in}{1.312989in}}%
\pgfpathquadraticcurveto{\pgfqpoint{6.590894in}{1.306487in}}{\pgfqpoint{6.590894in}{1.295067in}}%
\pgfpathquadraticcurveto{\pgfqpoint{6.590894in}{1.283648in}}{\pgfqpoint{6.595595in}{1.277145in}}%
\pgfpathquadraticcurveto{\pgfqpoint{6.600297in}{1.270643in}}{\pgfqpoint{6.608528in}{1.270643in}}%
\pgfpathquadraticcurveto{\pgfqpoint{6.616742in}{1.270643in}}{\pgfqpoint{6.621443in}{1.277145in}}%
\pgfpathquadraticcurveto{\pgfqpoint{6.626144in}{1.283648in}}{\pgfqpoint{6.626144in}{1.295067in}}%
\pgfpathlineto{\pgfqpoint{6.626144in}{1.295067in}}%
\pgfpathclose%
\pgfpathmoveto{\pgfqpoint{6.682720in}{1.295283in}}%
\pgfpathquadraticcurveto{\pgfqpoint{6.670166in}{1.295283in}}{\pgfqpoint{6.665321in}{1.292419in}}%
\pgfpathquadraticcurveto{\pgfqpoint{6.660475in}{1.289556in}}{\pgfqpoint{6.660475in}{1.282621in}}%
\pgfpathquadraticcurveto{\pgfqpoint{6.660475in}{1.277109in}}{\pgfqpoint{6.664114in}{1.273867in}}%
\pgfpathquadraticcurveto{\pgfqpoint{6.667752in}{1.270643in}}{\pgfqpoint{6.673984in}{1.270643in}}%
\pgfpathquadraticcurveto{\pgfqpoint{6.682612in}{1.270643in}}{\pgfqpoint{6.687818in}{1.276749in}}%
\pgfpathquadraticcurveto{\pgfqpoint{6.693023in}{1.282855in}}{\pgfqpoint{6.693023in}{1.292978in}}%
\pgfpathlineto{\pgfqpoint{6.693023in}{1.295283in}}%
\pgfpathlineto{\pgfqpoint{6.682720in}{1.295283in}}%
\pgfpathlineto{\pgfqpoint{6.682720in}{1.295283in}}%
\pgfpathclose%
\pgfpathmoveto{\pgfqpoint{6.703380in}{1.299570in}}%
\pgfpathlineto{\pgfqpoint{6.703380in}{1.263600in}}%
\pgfpathlineto{\pgfqpoint{6.693023in}{1.263600in}}%
\pgfpathlineto{\pgfqpoint{6.693023in}{1.273164in}}%
\pgfpathquadraticcurveto{\pgfqpoint{6.689475in}{1.267437in}}{\pgfqpoint{6.684179in}{1.264699in}}%
\pgfpathquadraticcurveto{\pgfqpoint{6.678884in}{1.261961in}}{\pgfqpoint{6.671229in}{1.261961in}}%
\pgfpathquadraticcurveto{\pgfqpoint{6.661556in}{1.261961in}}{\pgfqpoint{6.655828in}{1.267401in}}%
\pgfpathquadraticcurveto{\pgfqpoint{6.650118in}{1.272840in}}{\pgfqpoint{6.650118in}{1.281954in}}%
\pgfpathquadraticcurveto{\pgfqpoint{6.650118in}{1.292582in}}{\pgfqpoint{6.657233in}{1.297985in}}%
\pgfpathquadraticcurveto{\pgfqpoint{6.664366in}{1.303389in}}{\pgfqpoint{6.678487in}{1.303389in}}%
\pgfpathlineto{\pgfqpoint{6.693023in}{1.303389in}}%
\pgfpathlineto{\pgfqpoint{6.693023in}{1.304416in}}%
\pgfpathquadraticcurveto{\pgfqpoint{6.693023in}{1.311566in}}{\pgfqpoint{6.688322in}{1.315475in}}%
\pgfpathquadraticcurveto{\pgfqpoint{6.683621in}{1.319384in}}{\pgfqpoint{6.675119in}{1.319384in}}%
\pgfpathquadraticcurveto{\pgfqpoint{6.669715in}{1.319384in}}{\pgfqpoint{6.664582in}{1.318087in}}%
\pgfpathquadraticcurveto{\pgfqpoint{6.659467in}{1.316790in}}{\pgfqpoint{6.654747in}{1.314196in}}%
\pgfpathlineto{\pgfqpoint{6.654747in}{1.323779in}}%
\pgfpathquadraticcurveto{\pgfqpoint{6.660421in}{1.325976in}}{\pgfqpoint{6.665771in}{1.327057in}}%
\pgfpathquadraticcurveto{\pgfqpoint{6.671120in}{1.328156in}}{\pgfqpoint{6.676182in}{1.328156in}}%
\pgfpathquadraticcurveto{\pgfqpoint{6.689871in}{1.328156in}}{\pgfqpoint{6.696626in}{1.321059in}}%
\pgfpathquadraticcurveto{\pgfqpoint{6.703380in}{1.313980in}}{\pgfqpoint{6.703380in}{1.299570in}}%
\pgfpathlineto{\pgfqpoint{6.703380in}{1.299570in}}%
\pgfpathclose%
\pgfpathmoveto{\pgfqpoint{6.761235in}{1.316970in}}%
\pgfpathquadraticcurveto{\pgfqpoint{6.759488in}{1.317979in}}{\pgfqpoint{6.757435in}{1.318447in}}%
\pgfpathquadraticcurveto{\pgfqpoint{6.755381in}{1.318933in}}{\pgfqpoint{6.752914in}{1.318933in}}%
\pgfpathquadraticcurveto{\pgfqpoint{6.744124in}{1.318933in}}{\pgfqpoint{6.739423in}{1.313223in}}%
\pgfpathquadraticcurveto{\pgfqpoint{6.734721in}{1.307514in}}{\pgfqpoint{6.734721in}{1.296814in}}%
\pgfpathlineto{\pgfqpoint{6.734721in}{1.263600in}}%
\pgfpathlineto{\pgfqpoint{6.724310in}{1.263600in}}%
\pgfpathlineto{\pgfqpoint{6.724310in}{1.326643in}}%
\pgfpathlineto{\pgfqpoint{6.734721in}{1.326643in}}%
\pgfpathlineto{\pgfqpoint{6.734721in}{1.316844in}}%
\pgfpathquadraticcurveto{\pgfqpoint{6.738000in}{1.322590in}}{\pgfqpoint{6.743223in}{1.325364in}}%
\pgfpathquadraticcurveto{\pgfqpoint{6.748465in}{1.328156in}}{\pgfqpoint{6.755958in}{1.328156in}}%
\pgfpathquadraticcurveto{\pgfqpoint{6.757020in}{1.328156in}}{\pgfqpoint{6.758317in}{1.328011in}}%
\pgfpathquadraticcurveto{\pgfqpoint{6.759614in}{1.327885in}}{\pgfqpoint{6.761181in}{1.327597in}}%
\pgfpathlineto{\pgfqpoint{6.761235in}{1.316970in}}%
\pgfpathlineto{\pgfqpoint{6.761235in}{1.316970in}}%
\pgfpathclose%
\pgfpathmoveto{\pgfqpoint{6.800754in}{1.295283in}}%
\pgfpathquadraticcurveto{\pgfqpoint{6.788199in}{1.295283in}}{\pgfqpoint{6.783354in}{1.292419in}}%
\pgfpathquadraticcurveto{\pgfqpoint{6.778509in}{1.289556in}}{\pgfqpoint{6.778509in}{1.282621in}}%
\pgfpathquadraticcurveto{\pgfqpoint{6.778509in}{1.277109in}}{\pgfqpoint{6.782147in}{1.273867in}}%
\pgfpathquadraticcurveto{\pgfqpoint{6.785786in}{1.270643in}}{\pgfqpoint{6.792018in}{1.270643in}}%
\pgfpathquadraticcurveto{\pgfqpoint{6.800646in}{1.270643in}}{\pgfqpoint{6.805851in}{1.276749in}}%
\pgfpathquadraticcurveto{\pgfqpoint{6.811057in}{1.282855in}}{\pgfqpoint{6.811057in}{1.292978in}}%
\pgfpathlineto{\pgfqpoint{6.811057in}{1.295283in}}%
\pgfpathlineto{\pgfqpoint{6.800754in}{1.295283in}}%
\pgfpathlineto{\pgfqpoint{6.800754in}{1.295283in}}%
\pgfpathclose%
\pgfpathmoveto{\pgfqpoint{6.821414in}{1.299570in}}%
\pgfpathlineto{\pgfqpoint{6.821414in}{1.263600in}}%
\pgfpathlineto{\pgfqpoint{6.811057in}{1.263600in}}%
\pgfpathlineto{\pgfqpoint{6.811057in}{1.273164in}}%
\pgfpathquadraticcurveto{\pgfqpoint{6.807508in}{1.267437in}}{\pgfqpoint{6.802213in}{1.264699in}}%
\pgfpathquadraticcurveto{\pgfqpoint{6.796917in}{1.261961in}}{\pgfqpoint{6.789262in}{1.261961in}}%
\pgfpathquadraticcurveto{\pgfqpoint{6.779590in}{1.261961in}}{\pgfqpoint{6.773862in}{1.267401in}}%
\pgfpathquadraticcurveto{\pgfqpoint{6.768152in}{1.272840in}}{\pgfqpoint{6.768152in}{1.281954in}}%
\pgfpathquadraticcurveto{\pgfqpoint{6.768152in}{1.292582in}}{\pgfqpoint{6.775267in}{1.297985in}}%
\pgfpathquadraticcurveto{\pgfqpoint{6.782400in}{1.303389in}}{\pgfqpoint{6.796521in}{1.303389in}}%
\pgfpathlineto{\pgfqpoint{6.811057in}{1.303389in}}%
\pgfpathlineto{\pgfqpoint{6.811057in}{1.304416in}}%
\pgfpathquadraticcurveto{\pgfqpoint{6.811057in}{1.311566in}}{\pgfqpoint{6.806356in}{1.315475in}}%
\pgfpathquadraticcurveto{\pgfqpoint{6.801655in}{1.319384in}}{\pgfqpoint{6.793153in}{1.319384in}}%
\pgfpathquadraticcurveto{\pgfqpoint{6.787749in}{1.319384in}}{\pgfqpoint{6.782616in}{1.318087in}}%
\pgfpathquadraticcurveto{\pgfqpoint{6.777500in}{1.316790in}}{\pgfqpoint{6.772781in}{1.314196in}}%
\pgfpathlineto{\pgfqpoint{6.772781in}{1.323779in}}%
\pgfpathquadraticcurveto{\pgfqpoint{6.778455in}{1.325976in}}{\pgfqpoint{6.783804in}{1.327057in}}%
\pgfpathquadraticcurveto{\pgfqpoint{6.789154in}{1.328156in}}{\pgfqpoint{6.794215in}{1.328156in}}%
\pgfpathquadraticcurveto{\pgfqpoint{6.807905in}{1.328156in}}{\pgfqpoint{6.814659in}{1.321059in}}%
\pgfpathquadraticcurveto{\pgfqpoint{6.821414in}{1.313980in}}{\pgfqpoint{6.821414in}{1.299570in}}%
\pgfpathlineto{\pgfqpoint{6.821414in}{1.299570in}}%
\pgfpathclose%
\pgfpathmoveto{\pgfqpoint{6.852989in}{1.344547in}}%
\pgfpathlineto{\pgfqpoint{6.852989in}{1.326643in}}%
\pgfpathlineto{\pgfqpoint{6.874316in}{1.326643in}}%
\pgfpathlineto{\pgfqpoint{6.874316in}{1.318591in}}%
\pgfpathlineto{\pgfqpoint{6.852989in}{1.318591in}}%
\pgfpathlineto{\pgfqpoint{6.852989in}{1.284368in}}%
\pgfpathquadraticcurveto{\pgfqpoint{6.852989in}{1.276659in}}{\pgfqpoint{6.855097in}{1.274461in}}%
\pgfpathquadraticcurveto{\pgfqpoint{6.857204in}{1.272264in}}{\pgfqpoint{6.863688in}{1.272264in}}%
\pgfpathlineto{\pgfqpoint{6.874316in}{1.272264in}}%
\pgfpathlineto{\pgfqpoint{6.874316in}{1.263600in}}%
\pgfpathlineto{\pgfqpoint{6.863688in}{1.263600in}}%
\pgfpathquadraticcurveto{\pgfqpoint{6.851692in}{1.263600in}}{\pgfqpoint{6.847135in}{1.268067in}}%
\pgfpathquadraticcurveto{\pgfqpoint{6.842578in}{1.272552in}}{\pgfqpoint{6.842578in}{1.284368in}}%
\pgfpathlineto{\pgfqpoint{6.842578in}{1.318591in}}%
\pgfpathlineto{\pgfqpoint{6.834977in}{1.318591in}}%
\pgfpathlineto{\pgfqpoint{6.834977in}{1.326643in}}%
\pgfpathlineto{\pgfqpoint{6.842578in}{1.326643in}}%
\pgfpathlineto{\pgfqpoint{6.842578in}{1.344547in}}%
\pgfpathlineto{\pgfqpoint{6.852989in}{1.344547in}}%
\pgfpathlineto{\pgfqpoint{6.852989in}{1.344547in}}%
\pgfpathclose%
\pgfpathmoveto{\pgfqpoint{6.941861in}{1.297715in}}%
\pgfpathlineto{\pgfqpoint{6.941861in}{1.292654in}}%
\pgfpathlineto{\pgfqpoint{6.894237in}{1.292654in}}%
\pgfpathquadraticcurveto{\pgfqpoint{6.894921in}{1.281954in}}{\pgfqpoint{6.900685in}{1.276353in}}%
\pgfpathquadraticcurveto{\pgfqpoint{6.906449in}{1.270751in}}{\pgfqpoint{6.916752in}{1.270751in}}%
\pgfpathquadraticcurveto{\pgfqpoint{6.922714in}{1.270751in}}{\pgfqpoint{6.928316in}{1.272210in}}%
\pgfpathquadraticcurveto{\pgfqpoint{6.933918in}{1.273669in}}{\pgfqpoint{6.939447in}{1.276605in}}%
\pgfpathlineto{\pgfqpoint{6.939447in}{1.266806in}}%
\pgfpathquadraticcurveto{\pgfqpoint{6.933864in}{1.264447in}}{\pgfqpoint{6.928010in}{1.263204in}}%
\pgfpathquadraticcurveto{\pgfqpoint{6.922156in}{1.261961in}}{\pgfqpoint{6.916140in}{1.261961in}}%
\pgfpathquadraticcurveto{\pgfqpoint{6.901046in}{1.261961in}}{\pgfqpoint{6.892238in}{1.270733in}}%
\pgfpathquadraticcurveto{\pgfqpoint{6.883430in}{1.279523in}}{\pgfqpoint{6.883430in}{1.294509in}}%
\pgfpathquadraticcurveto{\pgfqpoint{6.883430in}{1.309981in}}{\pgfqpoint{6.891787in}{1.319059in}}%
\pgfpathquadraticcurveto{\pgfqpoint{6.900145in}{1.328156in}}{\pgfqpoint{6.914339in}{1.328156in}}%
\pgfpathquadraticcurveto{\pgfqpoint{6.927055in}{1.328156in}}{\pgfqpoint{6.934458in}{1.319960in}}%
\pgfpathquadraticcurveto{\pgfqpoint{6.941861in}{1.311783in}}{\pgfqpoint{6.941861in}{1.297715in}}%
\pgfpathlineto{\pgfqpoint{6.941861in}{1.297715in}}%
\pgfpathclose%
\pgfpathmoveto{\pgfqpoint{6.931504in}{1.300759in}}%
\pgfpathquadraticcurveto{\pgfqpoint{6.931396in}{1.309243in}}{\pgfqpoint{6.926749in}{1.314304in}}%
\pgfpathquadraticcurveto{\pgfqpoint{6.922102in}{1.319384in}}{\pgfqpoint{6.914447in}{1.319384in}}%
\pgfpathquadraticcurveto{\pgfqpoint{6.905783in}{1.319384in}}{\pgfqpoint{6.900577in}{1.314484in}}%
\pgfpathquadraticcurveto{\pgfqpoint{6.895372in}{1.309585in}}{\pgfqpoint{6.894579in}{1.300687in}}%
\pgfpathlineto{\pgfqpoint{6.931504in}{1.300759in}}%
\pgfpathlineto{\pgfqpoint{6.931504in}{1.300759in}}%
\pgfpathclose%
\pgfpathmoveto{\pgfqpoint{6.958865in}{1.351193in}}%
\pgfpathlineto{\pgfqpoint{6.969222in}{1.351193in}}%
\pgfpathlineto{\pgfqpoint{6.969222in}{1.263600in}}%
\pgfpathlineto{\pgfqpoint{6.958865in}{1.263600in}}%
\pgfpathlineto{\pgfqpoint{6.958865in}{1.351193in}}%
\pgfpathlineto{\pgfqpoint{6.958865in}{1.351193in}}%
\pgfpathclose%
\pgfpathmoveto{\pgfqpoint{7.017116in}{1.257746in}}%
\pgfpathquadraticcurveto{\pgfqpoint{7.012739in}{1.246488in}}{\pgfqpoint{7.008560in}{1.243066in}}%
\pgfpathquadraticcurveto{\pgfqpoint{7.004399in}{1.239626in}}{\pgfqpoint{6.997429in}{1.239626in}}%
\pgfpathlineto{\pgfqpoint{6.989143in}{1.239626in}}%
\pgfpathlineto{\pgfqpoint{6.989143in}{1.248290in}}%
\pgfpathlineto{\pgfqpoint{6.995231in}{1.248290in}}%
\pgfpathquadraticcurveto{\pgfqpoint{6.999500in}{1.248290in}}{\pgfqpoint{7.001860in}{1.250325in}}%
\pgfpathquadraticcurveto{\pgfqpoint{7.004237in}{1.252342in}}{\pgfqpoint{7.007101in}{1.259889in}}%
\pgfpathlineto{\pgfqpoint{7.008956in}{1.264609in}}%
\pgfpathlineto{\pgfqpoint{6.983469in}{1.326643in}}%
\pgfpathlineto{\pgfqpoint{6.994439in}{1.326643in}}%
\pgfpathlineto{\pgfqpoint{7.014144in}{1.277343in}}%
\pgfpathlineto{\pgfqpoint{7.033849in}{1.326643in}}%
\pgfpathlineto{\pgfqpoint{7.044819in}{1.326643in}}%
\pgfpathlineto{\pgfqpoint{7.017116in}{1.257746in}}%
\pgfpathlineto{\pgfqpoint{7.017116in}{1.257746in}}%
\pgfpathclose%
\pgfpathmoveto{\pgfqpoint{7.044152in}{1.277902in}}%
\pgfpathlineto{\pgfqpoint{7.056040in}{1.277902in}}%
\pgfpathlineto{\pgfqpoint{7.056040in}{1.263600in}}%
\pgfpathlineto{\pgfqpoint{7.044152in}{1.263600in}}%
\pgfpathlineto{\pgfqpoint{7.044152in}{1.277902in}}%
\pgfpathlineto{\pgfqpoint{7.044152in}{1.277902in}}%
\pgfpathclose%
\pgfusepath{fill}%
\end{pgfscope}%
\begin{pgfscope}%
\pgfsetbuttcap%
\pgfsetroundjoin%
\definecolor{currentfill}{rgb}{0.309804,0.372549,0.419608}%
\pgfsetfillcolor{currentfill}%
\pgfsetlinewidth{0.000000pt}%
\definecolor{currentstroke}{rgb}{0.309804,0.372549,0.419608}%
\pgfsetstrokecolor{currentstroke}%
\pgfsetdash{}{0pt}%
\pgfpathmoveto{\pgfqpoint{7.160990in}{1.338296in}}%
\pgfpathlineto{\pgfqpoint{7.160990in}{1.306721in}}%
\pgfpathlineto{\pgfqpoint{7.175291in}{1.306721in}}%
\pgfpathquadraticcurveto{\pgfqpoint{7.183235in}{1.306721in}}{\pgfqpoint{7.187558in}{1.310828in}}%
\pgfpathquadraticcurveto{\pgfqpoint{7.191899in}{1.314935in}}{\pgfqpoint{7.191899in}{1.322536in}}%
\pgfpathquadraticcurveto{\pgfqpoint{7.191899in}{1.330083in}}{\pgfqpoint{7.187558in}{1.334190in}}%
\pgfpathquadraticcurveto{\pgfqpoint{7.183235in}{1.338296in}}{\pgfqpoint{7.175291in}{1.338296in}}%
\pgfpathlineto{\pgfqpoint{7.160990in}{1.338296in}}%
\pgfpathlineto{\pgfqpoint{7.160990in}{1.338296in}}%
\pgfpathclose%
\pgfpathmoveto{\pgfqpoint{7.149624in}{1.347645in}}%
\pgfpathlineto{\pgfqpoint{7.175291in}{1.347645in}}%
\pgfpathquadraticcurveto{\pgfqpoint{7.189431in}{1.347645in}}{\pgfqpoint{7.196654in}{1.341250in}}%
\pgfpathquadraticcurveto{\pgfqpoint{7.203895in}{1.334856in}}{\pgfqpoint{7.203895in}{1.322536in}}%
\pgfpathquadraticcurveto{\pgfqpoint{7.203895in}{1.310089in}}{\pgfqpoint{7.196654in}{1.303731in}}%
\pgfpathquadraticcurveto{\pgfqpoint{7.189431in}{1.297373in}}{\pgfqpoint{7.175291in}{1.297373in}}%
\pgfpathlineto{\pgfqpoint{7.160990in}{1.297373in}}%
\pgfpathlineto{\pgfqpoint{7.160990in}{1.263600in}}%
\pgfpathlineto{\pgfqpoint{7.149624in}{1.263600in}}%
\pgfpathlineto{\pgfqpoint{7.149624in}{1.347645in}}%
\pgfpathlineto{\pgfqpoint{7.149624in}{1.347645in}}%
\pgfpathclose%
\pgfpathmoveto{\pgfqpoint{7.253194in}{1.316970in}}%
\pgfpathquadraticcurveto{\pgfqpoint{7.251447in}{1.317979in}}{\pgfqpoint{7.249393in}{1.318447in}}%
\pgfpathquadraticcurveto{\pgfqpoint{7.247340in}{1.318933in}}{\pgfqpoint{7.244872in}{1.318933in}}%
\pgfpathquadraticcurveto{\pgfqpoint{7.236082in}{1.318933in}}{\pgfqpoint{7.231381in}{1.313223in}}%
\pgfpathquadraticcurveto{\pgfqpoint{7.226680in}{1.307514in}}{\pgfqpoint{7.226680in}{1.296814in}}%
\pgfpathlineto{\pgfqpoint{7.226680in}{1.263600in}}%
\pgfpathlineto{\pgfqpoint{7.216269in}{1.263600in}}%
\pgfpathlineto{\pgfqpoint{7.216269in}{1.326643in}}%
\pgfpathlineto{\pgfqpoint{7.226680in}{1.326643in}}%
\pgfpathlineto{\pgfqpoint{7.226680in}{1.316844in}}%
\pgfpathquadraticcurveto{\pgfqpoint{7.229958in}{1.322590in}}{\pgfqpoint{7.235182in}{1.325364in}}%
\pgfpathquadraticcurveto{\pgfqpoint{7.240423in}{1.328156in}}{\pgfqpoint{7.247916in}{1.328156in}}%
\pgfpathquadraticcurveto{\pgfqpoint{7.248979in}{1.328156in}}{\pgfqpoint{7.250276in}{1.328011in}}%
\pgfpathquadraticcurveto{\pgfqpoint{7.251573in}{1.327885in}}{\pgfqpoint{7.253140in}{1.327597in}}%
\pgfpathlineto{\pgfqpoint{7.253194in}{1.316970in}}%
\pgfpathlineto{\pgfqpoint{7.253194in}{1.316970in}}%
\pgfpathclose%
\pgfpathmoveto{\pgfqpoint{7.315444in}{1.297715in}}%
\pgfpathlineto{\pgfqpoint{7.315444in}{1.292654in}}%
\pgfpathlineto{\pgfqpoint{7.267820in}{1.292654in}}%
\pgfpathquadraticcurveto{\pgfqpoint{7.268504in}{1.281954in}}{\pgfqpoint{7.274268in}{1.276353in}}%
\pgfpathquadraticcurveto{\pgfqpoint{7.280032in}{1.270751in}}{\pgfqpoint{7.290335in}{1.270751in}}%
\pgfpathquadraticcurveto{\pgfqpoint{7.296297in}{1.270751in}}{\pgfqpoint{7.301899in}{1.272210in}}%
\pgfpathquadraticcurveto{\pgfqpoint{7.307501in}{1.273669in}}{\pgfqpoint{7.313030in}{1.276605in}}%
\pgfpathlineto{\pgfqpoint{7.313030in}{1.266806in}}%
\pgfpathquadraticcurveto{\pgfqpoint{7.307447in}{1.264447in}}{\pgfqpoint{7.301593in}{1.263204in}}%
\pgfpathquadraticcurveto{\pgfqpoint{7.295739in}{1.261961in}}{\pgfqpoint{7.289723in}{1.261961in}}%
\pgfpathquadraticcurveto{\pgfqpoint{7.274628in}{1.261961in}}{\pgfqpoint{7.265821in}{1.270733in}}%
\pgfpathquadraticcurveto{\pgfqpoint{7.257013in}{1.279523in}}{\pgfqpoint{7.257013in}{1.294509in}}%
\pgfpathquadraticcurveto{\pgfqpoint{7.257013in}{1.309981in}}{\pgfqpoint{7.265370in}{1.319059in}}%
\pgfpathquadraticcurveto{\pgfqpoint{7.273728in}{1.328156in}}{\pgfqpoint{7.287921in}{1.328156in}}%
\pgfpathquadraticcurveto{\pgfqpoint{7.300638in}{1.328156in}}{\pgfqpoint{7.308041in}{1.319960in}}%
\pgfpathquadraticcurveto{\pgfqpoint{7.315444in}{1.311783in}}{\pgfqpoint{7.315444in}{1.297715in}}%
\pgfpathlineto{\pgfqpoint{7.315444in}{1.297715in}}%
\pgfpathclose%
\pgfpathmoveto{\pgfqpoint{7.305087in}{1.300759in}}%
\pgfpathquadraticcurveto{\pgfqpoint{7.304979in}{1.309243in}}{\pgfqpoint{7.300332in}{1.314304in}}%
\pgfpathquadraticcurveto{\pgfqpoint{7.295685in}{1.319384in}}{\pgfqpoint{7.288030in}{1.319384in}}%
\pgfpathquadraticcurveto{\pgfqpoint{7.279366in}{1.319384in}}{\pgfqpoint{7.274160in}{1.314484in}}%
\pgfpathquadraticcurveto{\pgfqpoint{7.268955in}{1.309585in}}{\pgfqpoint{7.268162in}{1.300687in}}%
\pgfpathlineto{\pgfqpoint{7.305087in}{1.300759in}}%
\pgfpathlineto{\pgfqpoint{7.305087in}{1.300759in}}%
\pgfpathclose%
\pgfusepath{fill}%
\end{pgfscope}%
\begin{pgfscope}%
\pgfsetbuttcap%
\pgfsetroundjoin%
\definecolor{currentfill}{rgb}{0.309804,0.372549,0.419608}%
\pgfsetfillcolor{currentfill}%
\pgfsetlinewidth{0.000000pt}%
\definecolor{currentstroke}{rgb}{0.309804,0.372549,0.419608}%
\pgfsetstrokecolor{currentstroke}%
\pgfsetdash{}{0pt}%
\pgfpathmoveto{\pgfqpoint{7.418985in}{1.295283in}}%
\pgfpathquadraticcurveto{\pgfqpoint{7.406430in}{1.295283in}}{\pgfqpoint{7.401585in}{1.292419in}}%
\pgfpathquadraticcurveto{\pgfqpoint{7.396740in}{1.289556in}}{\pgfqpoint{7.396740in}{1.282621in}}%
\pgfpathquadraticcurveto{\pgfqpoint{7.396740in}{1.277109in}}{\pgfqpoint{7.400378in}{1.273867in}}%
\pgfpathquadraticcurveto{\pgfqpoint{7.404017in}{1.270643in}}{\pgfqpoint{7.410249in}{1.270643in}}%
\pgfpathquadraticcurveto{\pgfqpoint{7.418877in}{1.270643in}}{\pgfqpoint{7.424082in}{1.276749in}}%
\pgfpathquadraticcurveto{\pgfqpoint{7.429288in}{1.282855in}}{\pgfqpoint{7.429288in}{1.292978in}}%
\pgfpathlineto{\pgfqpoint{7.429288in}{1.295283in}}%
\pgfpathlineto{\pgfqpoint{7.418985in}{1.295283in}}%
\pgfpathlineto{\pgfqpoint{7.418985in}{1.295283in}}%
\pgfpathclose%
\pgfpathmoveto{\pgfqpoint{7.439645in}{1.299570in}}%
\pgfpathlineto{\pgfqpoint{7.439645in}{1.263600in}}%
\pgfpathlineto{\pgfqpoint{7.429288in}{1.263600in}}%
\pgfpathlineto{\pgfqpoint{7.429288in}{1.273164in}}%
\pgfpathquadraticcurveto{\pgfqpoint{7.425739in}{1.267437in}}{\pgfqpoint{7.420444in}{1.264699in}}%
\pgfpathquadraticcurveto{\pgfqpoint{7.415148in}{1.261961in}}{\pgfqpoint{7.407493in}{1.261961in}}%
\pgfpathquadraticcurveto{\pgfqpoint{7.397821in}{1.261961in}}{\pgfqpoint{7.392093in}{1.267401in}}%
\pgfpathquadraticcurveto{\pgfqpoint{7.386383in}{1.272840in}}{\pgfqpoint{7.386383in}{1.281954in}}%
\pgfpathquadraticcurveto{\pgfqpoint{7.386383in}{1.292582in}}{\pgfqpoint{7.393498in}{1.297985in}}%
\pgfpathquadraticcurveto{\pgfqpoint{7.400630in}{1.303389in}}{\pgfqpoint{7.414752in}{1.303389in}}%
\pgfpathlineto{\pgfqpoint{7.429288in}{1.303389in}}%
\pgfpathlineto{\pgfqpoint{7.429288in}{1.304416in}}%
\pgfpathquadraticcurveto{\pgfqpoint{7.429288in}{1.311566in}}{\pgfqpoint{7.424587in}{1.315475in}}%
\pgfpathquadraticcurveto{\pgfqpoint{7.419885in}{1.319384in}}{\pgfqpoint{7.411384in}{1.319384in}}%
\pgfpathquadraticcurveto{\pgfqpoint{7.405980in}{1.319384in}}{\pgfqpoint{7.400847in}{1.318087in}}%
\pgfpathquadraticcurveto{\pgfqpoint{7.395731in}{1.316790in}}{\pgfqpoint{7.391012in}{1.314196in}}%
\pgfpathlineto{\pgfqpoint{7.391012in}{1.323779in}}%
\pgfpathquadraticcurveto{\pgfqpoint{7.396686in}{1.325976in}}{\pgfqpoint{7.402035in}{1.327057in}}%
\pgfpathquadraticcurveto{\pgfqpoint{7.407385in}{1.328156in}}{\pgfqpoint{7.412446in}{1.328156in}}%
\pgfpathquadraticcurveto{\pgfqpoint{7.426136in}{1.328156in}}{\pgfqpoint{7.432890in}{1.321059in}}%
\pgfpathquadraticcurveto{\pgfqpoint{7.439645in}{1.313980in}}{\pgfqpoint{7.439645in}{1.299570in}}%
\pgfpathlineto{\pgfqpoint{7.439645in}{1.299570in}}%
\pgfpathclose%
\pgfpathmoveto{\pgfqpoint{7.453550in}{1.326643in}}%
\pgfpathlineto{\pgfqpoint{7.464520in}{1.326643in}}%
\pgfpathlineto{\pgfqpoint{7.484225in}{1.273741in}}%
\pgfpathlineto{\pgfqpoint{7.503930in}{1.326643in}}%
\pgfpathlineto{\pgfqpoint{7.514900in}{1.326643in}}%
\pgfpathlineto{\pgfqpoint{7.491250in}{1.263600in}}%
\pgfpathlineto{\pgfqpoint{7.477182in}{1.263600in}}%
\pgfpathlineto{\pgfqpoint{7.453550in}{1.326643in}}%
\pgfpathlineto{\pgfqpoint{7.453550in}{1.326643in}}%
\pgfpathclose%
\pgfpathmoveto{\pgfqpoint{7.583130in}{1.297715in}}%
\pgfpathlineto{\pgfqpoint{7.583130in}{1.292654in}}%
\pgfpathlineto{\pgfqpoint{7.535505in}{1.292654in}}%
\pgfpathquadraticcurveto{\pgfqpoint{7.536190in}{1.281954in}}{\pgfqpoint{7.541954in}{1.276353in}}%
\pgfpathquadraticcurveto{\pgfqpoint{7.547718in}{1.270751in}}{\pgfqpoint{7.558021in}{1.270751in}}%
\pgfpathquadraticcurveto{\pgfqpoint{7.563983in}{1.270751in}}{\pgfqpoint{7.569584in}{1.272210in}}%
\pgfpathquadraticcurveto{\pgfqpoint{7.575186in}{1.273669in}}{\pgfqpoint{7.580716in}{1.276605in}}%
\pgfpathlineto{\pgfqpoint{7.580716in}{1.266806in}}%
\pgfpathquadraticcurveto{\pgfqpoint{7.575132in}{1.264447in}}{\pgfqpoint{7.569278in}{1.263204in}}%
\pgfpathquadraticcurveto{\pgfqpoint{7.563424in}{1.261961in}}{\pgfqpoint{7.557408in}{1.261961in}}%
\pgfpathquadraticcurveto{\pgfqpoint{7.542314in}{1.261961in}}{\pgfqpoint{7.533506in}{1.270733in}}%
\pgfpathquadraticcurveto{\pgfqpoint{7.524698in}{1.279523in}}{\pgfqpoint{7.524698in}{1.294509in}}%
\pgfpathquadraticcurveto{\pgfqpoint{7.524698in}{1.309981in}}{\pgfqpoint{7.533056in}{1.319059in}}%
\pgfpathquadraticcurveto{\pgfqpoint{7.541413in}{1.328156in}}{\pgfqpoint{7.555607in}{1.328156in}}%
\pgfpathquadraticcurveto{\pgfqpoint{7.568324in}{1.328156in}}{\pgfqpoint{7.575727in}{1.319960in}}%
\pgfpathquadraticcurveto{\pgfqpoint{7.583130in}{1.311783in}}{\pgfqpoint{7.583130in}{1.297715in}}%
\pgfpathlineto{\pgfqpoint{7.583130in}{1.297715in}}%
\pgfpathclose%
\pgfpathmoveto{\pgfqpoint{7.572773in}{1.300759in}}%
\pgfpathquadraticcurveto{\pgfqpoint{7.572664in}{1.309243in}}{\pgfqpoint{7.568017in}{1.314304in}}%
\pgfpathquadraticcurveto{\pgfqpoint{7.563370in}{1.319384in}}{\pgfqpoint{7.555715in}{1.319384in}}%
\pgfpathquadraticcurveto{\pgfqpoint{7.547051in}{1.319384in}}{\pgfqpoint{7.541846in}{1.314484in}}%
\pgfpathquadraticcurveto{\pgfqpoint{7.536640in}{1.309585in}}{\pgfqpoint{7.535848in}{1.300687in}}%
\pgfpathlineto{\pgfqpoint{7.572773in}{1.300759in}}%
\pgfpathlineto{\pgfqpoint{7.572773in}{1.300759in}}%
\pgfpathclose%
\pgfpathmoveto{\pgfqpoint{7.636662in}{1.316970in}}%
\pgfpathquadraticcurveto{\pgfqpoint{7.634914in}{1.317979in}}{\pgfqpoint{7.632861in}{1.318447in}}%
\pgfpathquadraticcurveto{\pgfqpoint{7.630808in}{1.318933in}}{\pgfqpoint{7.628340in}{1.318933in}}%
\pgfpathquadraticcurveto{\pgfqpoint{7.619550in}{1.318933in}}{\pgfqpoint{7.614849in}{1.313223in}}%
\pgfpathquadraticcurveto{\pgfqpoint{7.610148in}{1.307514in}}{\pgfqpoint{7.610148in}{1.296814in}}%
\pgfpathlineto{\pgfqpoint{7.610148in}{1.263600in}}%
\pgfpathlineto{\pgfqpoint{7.599737in}{1.263600in}}%
\pgfpathlineto{\pgfqpoint{7.599737in}{1.326643in}}%
\pgfpathlineto{\pgfqpoint{7.610148in}{1.326643in}}%
\pgfpathlineto{\pgfqpoint{7.610148in}{1.316844in}}%
\pgfpathquadraticcurveto{\pgfqpoint{7.613426in}{1.322590in}}{\pgfqpoint{7.618650in}{1.325364in}}%
\pgfpathquadraticcurveto{\pgfqpoint{7.623891in}{1.328156in}}{\pgfqpoint{7.631384in}{1.328156in}}%
\pgfpathquadraticcurveto{\pgfqpoint{7.632447in}{1.328156in}}{\pgfqpoint{7.633744in}{1.328011in}}%
\pgfpathquadraticcurveto{\pgfqpoint{7.635041in}{1.327885in}}{\pgfqpoint{7.636608in}{1.327597in}}%
\pgfpathlineto{\pgfqpoint{7.636662in}{1.316970in}}%
\pgfpathlineto{\pgfqpoint{7.636662in}{1.316970in}}%
\pgfpathclose%
\pgfpathmoveto{\pgfqpoint{7.676180in}{1.295283in}}%
\pgfpathquadraticcurveto{\pgfqpoint{7.663626in}{1.295283in}}{\pgfqpoint{7.658781in}{1.292419in}}%
\pgfpathquadraticcurveto{\pgfqpoint{7.653935in}{1.289556in}}{\pgfqpoint{7.653935in}{1.282621in}}%
\pgfpathquadraticcurveto{\pgfqpoint{7.653935in}{1.277109in}}{\pgfqpoint{7.657574in}{1.273867in}}%
\pgfpathquadraticcurveto{\pgfqpoint{7.661212in}{1.270643in}}{\pgfqpoint{7.667444in}{1.270643in}}%
\pgfpathquadraticcurveto{\pgfqpoint{7.676072in}{1.270643in}}{\pgfqpoint{7.681278in}{1.276749in}}%
\pgfpathquadraticcurveto{\pgfqpoint{7.686483in}{1.282855in}}{\pgfqpoint{7.686483in}{1.292978in}}%
\pgfpathlineto{\pgfqpoint{7.686483in}{1.295283in}}%
\pgfpathlineto{\pgfqpoint{7.676180in}{1.295283in}}%
\pgfpathlineto{\pgfqpoint{7.676180in}{1.295283in}}%
\pgfpathclose%
\pgfpathmoveto{\pgfqpoint{7.696840in}{1.299570in}}%
\pgfpathlineto{\pgfqpoint{7.696840in}{1.263600in}}%
\pgfpathlineto{\pgfqpoint{7.686483in}{1.263600in}}%
\pgfpathlineto{\pgfqpoint{7.686483in}{1.273164in}}%
\pgfpathquadraticcurveto{\pgfqpoint{7.682935in}{1.267437in}}{\pgfqpoint{7.677639in}{1.264699in}}%
\pgfpathquadraticcurveto{\pgfqpoint{7.672344in}{1.261961in}}{\pgfqpoint{7.664689in}{1.261961in}}%
\pgfpathquadraticcurveto{\pgfqpoint{7.655016in}{1.261961in}}{\pgfqpoint{7.649288in}{1.267401in}}%
\pgfpathquadraticcurveto{\pgfqpoint{7.643578in}{1.272840in}}{\pgfqpoint{7.643578in}{1.281954in}}%
\pgfpathquadraticcurveto{\pgfqpoint{7.643578in}{1.292582in}}{\pgfqpoint{7.650693in}{1.297985in}}%
\pgfpathquadraticcurveto{\pgfqpoint{7.657826in}{1.303389in}}{\pgfqpoint{7.671947in}{1.303389in}}%
\pgfpathlineto{\pgfqpoint{7.686483in}{1.303389in}}%
\pgfpathlineto{\pgfqpoint{7.686483in}{1.304416in}}%
\pgfpathquadraticcurveto{\pgfqpoint{7.686483in}{1.311566in}}{\pgfqpoint{7.681782in}{1.315475in}}%
\pgfpathquadraticcurveto{\pgfqpoint{7.677081in}{1.319384in}}{\pgfqpoint{7.668579in}{1.319384in}}%
\pgfpathquadraticcurveto{\pgfqpoint{7.663176in}{1.319384in}}{\pgfqpoint{7.658042in}{1.318087in}}%
\pgfpathquadraticcurveto{\pgfqpoint{7.652927in}{1.316790in}}{\pgfqpoint{7.648207in}{1.314196in}}%
\pgfpathlineto{\pgfqpoint{7.648207in}{1.323779in}}%
\pgfpathquadraticcurveto{\pgfqpoint{7.653881in}{1.325976in}}{\pgfqpoint{7.659231in}{1.327057in}}%
\pgfpathquadraticcurveto{\pgfqpoint{7.664581in}{1.328156in}}{\pgfqpoint{7.669642in}{1.328156in}}%
\pgfpathquadraticcurveto{\pgfqpoint{7.683331in}{1.328156in}}{\pgfqpoint{7.690086in}{1.321059in}}%
\pgfpathquadraticcurveto{\pgfqpoint{7.696840in}{1.313980in}}{\pgfqpoint{7.696840in}{1.299570in}}%
\pgfpathlineto{\pgfqpoint{7.696840in}{1.299570in}}%
\pgfpathclose%
\pgfpathmoveto{\pgfqpoint{7.759649in}{1.295860in}}%
\pgfpathquadraticcurveto{\pgfqpoint{7.759649in}{1.307117in}}{\pgfqpoint{7.755002in}{1.313296in}}%
\pgfpathquadraticcurveto{\pgfqpoint{7.750372in}{1.319492in}}{\pgfqpoint{7.741979in}{1.319492in}}%
\pgfpathquadraticcurveto{\pgfqpoint{7.733657in}{1.319492in}}{\pgfqpoint{7.729010in}{1.313296in}}%
\pgfpathquadraticcurveto{\pgfqpoint{7.724363in}{1.307117in}}{\pgfqpoint{7.724363in}{1.295860in}}%
\pgfpathquadraticcurveto{\pgfqpoint{7.724363in}{1.284656in}}{\pgfqpoint{7.729010in}{1.278460in}}%
\pgfpathquadraticcurveto{\pgfqpoint{7.733657in}{1.272264in}}{\pgfqpoint{7.741979in}{1.272264in}}%
\pgfpathquadraticcurveto{\pgfqpoint{7.750372in}{1.272264in}}{\pgfqpoint{7.755002in}{1.278460in}}%
\pgfpathquadraticcurveto{\pgfqpoint{7.759649in}{1.284656in}}{\pgfqpoint{7.759649in}{1.295860in}}%
\pgfpathlineto{\pgfqpoint{7.759649in}{1.295860in}}%
\pgfpathclose%
\pgfpathmoveto{\pgfqpoint{7.770006in}{1.271417in}}%
\pgfpathquadraticcurveto{\pgfqpoint{7.770006in}{1.255332in}}{\pgfqpoint{7.762855in}{1.247479in}}%
\pgfpathquadraticcurveto{\pgfqpoint{7.755722in}{1.239626in}}{\pgfqpoint{7.740970in}{1.239626in}}%
\pgfpathquadraticcurveto{\pgfqpoint{7.735512in}{1.239626in}}{\pgfqpoint{7.730667in}{1.240436in}}%
\pgfpathquadraticcurveto{\pgfqpoint{7.725822in}{1.241247in}}{\pgfqpoint{7.721265in}{1.242940in}}%
\pgfpathlineto{\pgfqpoint{7.721265in}{1.253009in}}%
\pgfpathquadraticcurveto{\pgfqpoint{7.725822in}{1.250541in}}{\pgfqpoint{7.730271in}{1.249370in}}%
\pgfpathquadraticcurveto{\pgfqpoint{7.734720in}{1.248182in}}{\pgfqpoint{7.739331in}{1.248182in}}%
\pgfpathquadraticcurveto{\pgfqpoint{7.749526in}{1.248182in}}{\pgfqpoint{7.754587in}{1.253495in}}%
\pgfpathquadraticcurveto{\pgfqpoint{7.759649in}{1.258809in}}{\pgfqpoint{7.759649in}{1.269562in}}%
\pgfpathlineto{\pgfqpoint{7.759649in}{1.274695in}}%
\pgfpathquadraticcurveto{\pgfqpoint{7.756442in}{1.269112in}}{\pgfqpoint{7.751435in}{1.266356in}}%
\pgfpathquadraticcurveto{\pgfqpoint{7.746428in}{1.263600in}}{\pgfqpoint{7.739439in}{1.263600in}}%
\pgfpathquadraticcurveto{\pgfqpoint{7.727857in}{1.263600in}}{\pgfqpoint{7.720760in}{1.272426in}}%
\pgfpathquadraticcurveto{\pgfqpoint{7.713664in}{1.281270in}}{\pgfqpoint{7.713664in}{1.295860in}}%
\pgfpathquadraticcurveto{\pgfqpoint{7.713664in}{1.310486in}}{\pgfqpoint{7.720760in}{1.319312in}}%
\pgfpathquadraticcurveto{\pgfqpoint{7.727857in}{1.328156in}}{\pgfqpoint{7.739439in}{1.328156in}}%
\pgfpathquadraticcurveto{\pgfqpoint{7.746428in}{1.328156in}}{\pgfqpoint{7.751435in}{1.325400in}}%
\pgfpathquadraticcurveto{\pgfqpoint{7.756442in}{1.322644in}}{\pgfqpoint{7.759649in}{1.317078in}}%
\pgfpathlineto{\pgfqpoint{7.759649in}{1.326643in}}%
\pgfpathlineto{\pgfqpoint{7.770006in}{1.326643in}}%
\pgfpathlineto{\pgfqpoint{7.770006in}{1.271417in}}%
\pgfpathlineto{\pgfqpoint{7.770006in}{1.271417in}}%
\pgfpathclose%
\pgfpathmoveto{\pgfqpoint{7.845278in}{1.297715in}}%
\pgfpathlineto{\pgfqpoint{7.845278in}{1.292654in}}%
\pgfpathlineto{\pgfqpoint{7.797654in}{1.292654in}}%
\pgfpathquadraticcurveto{\pgfqpoint{7.798339in}{1.281954in}}{\pgfqpoint{7.804103in}{1.276353in}}%
\pgfpathquadraticcurveto{\pgfqpoint{7.809867in}{1.270751in}}{\pgfqpoint{7.820169in}{1.270751in}}%
\pgfpathquadraticcurveto{\pgfqpoint{7.826132in}{1.270751in}}{\pgfqpoint{7.831733in}{1.272210in}}%
\pgfpathquadraticcurveto{\pgfqpoint{7.837335in}{1.273669in}}{\pgfqpoint{7.842865in}{1.276605in}}%
\pgfpathlineto{\pgfqpoint{7.842865in}{1.266806in}}%
\pgfpathquadraticcurveto{\pgfqpoint{7.837281in}{1.264447in}}{\pgfqpoint{7.831427in}{1.263204in}}%
\pgfpathquadraticcurveto{\pgfqpoint{7.825573in}{1.261961in}}{\pgfqpoint{7.819557in}{1.261961in}}%
\pgfpathquadraticcurveto{\pgfqpoint{7.804463in}{1.261961in}}{\pgfqpoint{7.795655in}{1.270733in}}%
\pgfpathquadraticcurveto{\pgfqpoint{7.786847in}{1.279523in}}{\pgfqpoint{7.786847in}{1.294509in}}%
\pgfpathquadraticcurveto{\pgfqpoint{7.786847in}{1.309981in}}{\pgfqpoint{7.795205in}{1.319059in}}%
\pgfpathquadraticcurveto{\pgfqpoint{7.803562in}{1.328156in}}{\pgfqpoint{7.817756in}{1.328156in}}%
\pgfpathquadraticcurveto{\pgfqpoint{7.830472in}{1.328156in}}{\pgfqpoint{7.837875in}{1.319960in}}%
\pgfpathquadraticcurveto{\pgfqpoint{7.845278in}{1.311783in}}{\pgfqpoint{7.845278in}{1.297715in}}%
\pgfpathlineto{\pgfqpoint{7.845278in}{1.297715in}}%
\pgfpathclose%
\pgfpathmoveto{\pgfqpoint{7.834921in}{1.300759in}}%
\pgfpathquadraticcurveto{\pgfqpoint{7.834813in}{1.309243in}}{\pgfqpoint{7.830166in}{1.314304in}}%
\pgfpathquadraticcurveto{\pgfqpoint{7.825519in}{1.319384in}}{\pgfqpoint{7.817864in}{1.319384in}}%
\pgfpathquadraticcurveto{\pgfqpoint{7.809200in}{1.319384in}}{\pgfqpoint{7.803995in}{1.314484in}}%
\pgfpathquadraticcurveto{\pgfqpoint{7.798789in}{1.309585in}}{\pgfqpoint{7.797997in}{1.300687in}}%
\pgfpathlineto{\pgfqpoint{7.834921in}{1.300759in}}%
\pgfpathlineto{\pgfqpoint{7.834921in}{1.300759in}}%
\pgfpathclose%
\pgfpathmoveto{\pgfqpoint{7.902467in}{1.324787in}}%
\pgfpathlineto{\pgfqpoint{7.902467in}{1.314989in}}%
\pgfpathquadraticcurveto{\pgfqpoint{7.898090in}{1.317240in}}{\pgfqpoint{7.893353in}{1.318357in}}%
\pgfpathquadraticcurveto{\pgfqpoint{7.888634in}{1.319492in}}{\pgfqpoint{7.883554in}{1.319492in}}%
\pgfpathquadraticcurveto{\pgfqpoint{7.875845in}{1.319492in}}{\pgfqpoint{7.871990in}{1.317132in}}%
\pgfpathquadraticcurveto{\pgfqpoint{7.868136in}{1.314773in}}{\pgfqpoint{7.868136in}{1.310035in}}%
\pgfpathquadraticcurveto{\pgfqpoint{7.868136in}{1.306433in}}{\pgfqpoint{7.870892in}{1.304380in}}%
\pgfpathquadraticcurveto{\pgfqpoint{7.873648in}{1.302326in}}{\pgfqpoint{7.881987in}{1.300471in}}%
\pgfpathlineto{\pgfqpoint{7.885536in}{1.299678in}}%
\pgfpathquadraticcurveto{\pgfqpoint{7.896559in}{1.297319in}}{\pgfqpoint{7.901206in}{1.293014in}}%
\pgfpathquadraticcurveto{\pgfqpoint{7.905853in}{1.288709in}}{\pgfqpoint{7.905853in}{1.281000in}}%
\pgfpathquadraticcurveto{\pgfqpoint{7.905853in}{1.272210in}}{\pgfqpoint{7.898901in}{1.267076in}}%
\pgfpathquadraticcurveto{\pgfqpoint{7.891948in}{1.261961in}}{\pgfqpoint{7.879790in}{1.261961in}}%
\pgfpathquadraticcurveto{\pgfqpoint{7.874728in}{1.261961in}}{\pgfqpoint{7.869235in}{1.262952in}}%
\pgfpathquadraticcurveto{\pgfqpoint{7.863741in}{1.263942in}}{\pgfqpoint{7.857671in}{1.265906in}}%
\pgfpathlineto{\pgfqpoint{7.857671in}{1.276605in}}%
\pgfpathquadraticcurveto{\pgfqpoint{7.863417in}{1.273615in}}{\pgfqpoint{7.868982in}{1.272120in}}%
\pgfpathquadraticcurveto{\pgfqpoint{7.874548in}{1.270643in}}{\pgfqpoint{7.880024in}{1.270643in}}%
\pgfpathquadraticcurveto{\pgfqpoint{7.887337in}{1.270643in}}{\pgfqpoint{7.891263in}{1.273146in}}%
\pgfpathquadraticcurveto{\pgfqpoint{7.895208in}{1.275650in}}{\pgfqpoint{7.895208in}{1.280207in}}%
\pgfpathquadraticcurveto{\pgfqpoint{7.895208in}{1.284422in}}{\pgfqpoint{7.892362in}{1.286674in}}%
\pgfpathquadraticcurveto{\pgfqpoint{7.889534in}{1.288925in}}{\pgfqpoint{7.879898in}{1.291014in}}%
\pgfpathlineto{\pgfqpoint{7.876295in}{1.291861in}}%
\pgfpathquadraticcurveto{\pgfqpoint{7.866677in}{1.293878in}}{\pgfqpoint{7.862390in}{1.298075in}}%
\pgfpathquadraticcurveto{\pgfqpoint{7.858121in}{1.302272in}}{\pgfqpoint{7.858121in}{1.309585in}}%
\pgfpathquadraticcurveto{\pgfqpoint{7.858121in}{1.318483in}}{\pgfqpoint{7.864425in}{1.323310in}}%
\pgfpathquadraticcurveto{\pgfqpoint{7.870730in}{1.328156in}}{\pgfqpoint{7.882329in}{1.328156in}}%
\pgfpathquadraticcurveto{\pgfqpoint{7.888057in}{1.328156in}}{\pgfqpoint{7.893119in}{1.327309in}}%
\pgfpathquadraticcurveto{\pgfqpoint{7.898198in}{1.326480in}}{\pgfqpoint{7.902467in}{1.324787in}}%
\pgfpathlineto{\pgfqpoint{7.902467in}{1.324787in}}%
\pgfpathclose%
\pgfusepath{fill}%
\end{pgfscope}%
\begin{pgfscope}%
\pgfsetbuttcap%
\pgfsetroundjoin%
\definecolor{currentfill}{rgb}{0.309804,0.372549,0.419608}%
\pgfsetfillcolor{currentfill}%
\pgfsetlinewidth{0.000000pt}%
\definecolor{currentstroke}{rgb}{0.309804,0.372549,0.419608}%
\pgfsetstrokecolor{currentstroke}%
\pgfsetdash{}{0pt}%
\pgfpathmoveto{\pgfqpoint{8.014670in}{1.295283in}}%
\pgfpathquadraticcurveto{\pgfqpoint{8.002115in}{1.295283in}}{\pgfqpoint{7.997270in}{1.292419in}}%
\pgfpathquadraticcurveto{\pgfqpoint{7.992425in}{1.289556in}}{\pgfqpoint{7.992425in}{1.282621in}}%
\pgfpathquadraticcurveto{\pgfqpoint{7.992425in}{1.277109in}}{\pgfqpoint{7.996063in}{1.273867in}}%
\pgfpathquadraticcurveto{\pgfqpoint{7.999702in}{1.270643in}}{\pgfqpoint{8.005934in}{1.270643in}}%
\pgfpathquadraticcurveto{\pgfqpoint{8.014562in}{1.270643in}}{\pgfqpoint{8.019767in}{1.276749in}}%
\pgfpathquadraticcurveto{\pgfqpoint{8.024973in}{1.282855in}}{\pgfqpoint{8.024973in}{1.292978in}}%
\pgfpathlineto{\pgfqpoint{8.024973in}{1.295283in}}%
\pgfpathlineto{\pgfqpoint{8.014670in}{1.295283in}}%
\pgfpathlineto{\pgfqpoint{8.014670in}{1.295283in}}%
\pgfpathclose%
\pgfpathmoveto{\pgfqpoint{8.035330in}{1.299570in}}%
\pgfpathlineto{\pgfqpoint{8.035330in}{1.263600in}}%
\pgfpathlineto{\pgfqpoint{8.024973in}{1.263600in}}%
\pgfpathlineto{\pgfqpoint{8.024973in}{1.273164in}}%
\pgfpathquadraticcurveto{\pgfqpoint{8.021424in}{1.267437in}}{\pgfqpoint{8.016129in}{1.264699in}}%
\pgfpathquadraticcurveto{\pgfqpoint{8.010833in}{1.261961in}}{\pgfqpoint{8.003178in}{1.261961in}}%
\pgfpathquadraticcurveto{\pgfqpoint{7.993505in}{1.261961in}}{\pgfqpoint{7.987778in}{1.267401in}}%
\pgfpathquadraticcurveto{\pgfqpoint{7.982068in}{1.272840in}}{\pgfqpoint{7.982068in}{1.281954in}}%
\pgfpathquadraticcurveto{\pgfqpoint{7.982068in}{1.292582in}}{\pgfqpoint{7.989183in}{1.297985in}}%
\pgfpathquadraticcurveto{\pgfqpoint{7.996315in}{1.303389in}}{\pgfqpoint{8.010437in}{1.303389in}}%
\pgfpathlineto{\pgfqpoint{8.024973in}{1.303389in}}%
\pgfpathlineto{\pgfqpoint{8.024973in}{1.304416in}}%
\pgfpathquadraticcurveto{\pgfqpoint{8.024973in}{1.311566in}}{\pgfqpoint{8.020271in}{1.315475in}}%
\pgfpathquadraticcurveto{\pgfqpoint{8.015570in}{1.319384in}}{\pgfqpoint{8.007069in}{1.319384in}}%
\pgfpathquadraticcurveto{\pgfqpoint{8.001665in}{1.319384in}}{\pgfqpoint{7.996531in}{1.318087in}}%
\pgfpathquadraticcurveto{\pgfqpoint{7.991416in}{1.316790in}}{\pgfqpoint{7.986697in}{1.314196in}}%
\pgfpathlineto{\pgfqpoint{7.986697in}{1.323779in}}%
\pgfpathquadraticcurveto{\pgfqpoint{7.992371in}{1.325976in}}{\pgfqpoint{7.997720in}{1.327057in}}%
\pgfpathquadraticcurveto{\pgfqpoint{8.003070in}{1.328156in}}{\pgfqpoint{8.008131in}{1.328156in}}%
\pgfpathquadraticcurveto{\pgfqpoint{8.021821in}{1.328156in}}{\pgfqpoint{8.028575in}{1.321059in}}%
\pgfpathquadraticcurveto{\pgfqpoint{8.035330in}{1.313980in}}{\pgfqpoint{8.035330in}{1.299570in}}%
\pgfpathlineto{\pgfqpoint{8.035330in}{1.299570in}}%
\pgfpathclose%
\pgfpathmoveto{\pgfqpoint{8.093185in}{1.316970in}}%
\pgfpathquadraticcurveto{\pgfqpoint{8.091437in}{1.317979in}}{\pgfqpoint{8.089384in}{1.318447in}}%
\pgfpathquadraticcurveto{\pgfqpoint{8.087331in}{1.318933in}}{\pgfqpoint{8.084863in}{1.318933in}}%
\pgfpathquadraticcurveto{\pgfqpoint{8.076073in}{1.318933in}}{\pgfqpoint{8.071372in}{1.313223in}}%
\pgfpathquadraticcurveto{\pgfqpoint{8.066671in}{1.307514in}}{\pgfqpoint{8.066671in}{1.296814in}}%
\pgfpathlineto{\pgfqpoint{8.066671in}{1.263600in}}%
\pgfpathlineto{\pgfqpoint{8.056260in}{1.263600in}}%
\pgfpathlineto{\pgfqpoint{8.056260in}{1.326643in}}%
\pgfpathlineto{\pgfqpoint{8.066671in}{1.326643in}}%
\pgfpathlineto{\pgfqpoint{8.066671in}{1.316844in}}%
\pgfpathquadraticcurveto{\pgfqpoint{8.069949in}{1.322590in}}{\pgfqpoint{8.075173in}{1.325364in}}%
\pgfpathquadraticcurveto{\pgfqpoint{8.080414in}{1.328156in}}{\pgfqpoint{8.087907in}{1.328156in}}%
\pgfpathquadraticcurveto{\pgfqpoint{8.088970in}{1.328156in}}{\pgfqpoint{8.090267in}{1.328011in}}%
\pgfpathquadraticcurveto{\pgfqpoint{8.091564in}{1.327885in}}{\pgfqpoint{8.093131in}{1.327597in}}%
\pgfpathlineto{\pgfqpoint{8.093185in}{1.316970in}}%
\pgfpathlineto{\pgfqpoint{8.093185in}{1.316970in}}%
\pgfpathclose%
\pgfpathmoveto{\pgfqpoint{8.155435in}{1.297715in}}%
\pgfpathlineto{\pgfqpoint{8.155435in}{1.292654in}}%
\pgfpathlineto{\pgfqpoint{8.107811in}{1.292654in}}%
\pgfpathquadraticcurveto{\pgfqpoint{8.108495in}{1.281954in}}{\pgfqpoint{8.114259in}{1.276353in}}%
\pgfpathquadraticcurveto{\pgfqpoint{8.120023in}{1.270751in}}{\pgfqpoint{8.130326in}{1.270751in}}%
\pgfpathquadraticcurveto{\pgfqpoint{8.136288in}{1.270751in}}{\pgfqpoint{8.141890in}{1.272210in}}%
\pgfpathquadraticcurveto{\pgfqpoint{8.147491in}{1.273669in}}{\pgfqpoint{8.153021in}{1.276605in}}%
\pgfpathlineto{\pgfqpoint{8.153021in}{1.266806in}}%
\pgfpathquadraticcurveto{\pgfqpoint{8.147437in}{1.264447in}}{\pgfqpoint{8.141583in}{1.263204in}}%
\pgfpathquadraticcurveto{\pgfqpoint{8.135729in}{1.261961in}}{\pgfqpoint{8.129713in}{1.261961in}}%
\pgfpathquadraticcurveto{\pgfqpoint{8.114619in}{1.261961in}}{\pgfqpoint{8.105811in}{1.270733in}}%
\pgfpathquadraticcurveto{\pgfqpoint{8.097003in}{1.279523in}}{\pgfqpoint{8.097003in}{1.294509in}}%
\pgfpathquadraticcurveto{\pgfqpoint{8.097003in}{1.309981in}}{\pgfqpoint{8.105361in}{1.319059in}}%
\pgfpathquadraticcurveto{\pgfqpoint{8.113719in}{1.328156in}}{\pgfqpoint{8.127912in}{1.328156in}}%
\pgfpathquadraticcurveto{\pgfqpoint{8.140629in}{1.328156in}}{\pgfqpoint{8.148032in}{1.319960in}}%
\pgfpathquadraticcurveto{\pgfqpoint{8.155435in}{1.311783in}}{\pgfqpoint{8.155435in}{1.297715in}}%
\pgfpathlineto{\pgfqpoint{8.155435in}{1.297715in}}%
\pgfpathclose%
\pgfpathmoveto{\pgfqpoint{8.145078in}{1.300759in}}%
\pgfpathquadraticcurveto{\pgfqpoint{8.144970in}{1.309243in}}{\pgfqpoint{8.140322in}{1.314304in}}%
\pgfpathquadraticcurveto{\pgfqpoint{8.135675in}{1.319384in}}{\pgfqpoint{8.128020in}{1.319384in}}%
\pgfpathquadraticcurveto{\pgfqpoint{8.119356in}{1.319384in}}{\pgfqpoint{8.114151in}{1.314484in}}%
\pgfpathquadraticcurveto{\pgfqpoint{8.108945in}{1.309585in}}{\pgfqpoint{8.108153in}{1.300687in}}%
\pgfpathlineto{\pgfqpoint{8.145078in}{1.300759in}}%
\pgfpathlineto{\pgfqpoint{8.145078in}{1.300759in}}%
\pgfpathclose%
\pgfusepath{fill}%
\end{pgfscope}%
\begin{pgfscope}%
\pgfsetbuttcap%
\pgfsetroundjoin%
\definecolor{currentfill}{rgb}{0.309804,0.372549,0.419608}%
\pgfsetfillcolor{currentfill}%
\pgfsetlinewidth{0.000000pt}%
\definecolor{currentstroke}{rgb}{0.309804,0.372549,0.419608}%
\pgfsetstrokecolor{currentstroke}%
\pgfsetdash{}{0pt}%
\pgfpathmoveto{\pgfqpoint{8.270679in}{1.317078in}}%
\pgfpathlineto{\pgfqpoint{8.270679in}{1.351193in}}%
\pgfpathlineto{\pgfqpoint{8.281036in}{1.351193in}}%
\pgfpathlineto{\pgfqpoint{8.281036in}{1.263600in}}%
\pgfpathlineto{\pgfqpoint{8.270679in}{1.263600in}}%
\pgfpathlineto{\pgfqpoint{8.270679in}{1.273056in}}%
\pgfpathquadraticcurveto{\pgfqpoint{8.267419in}{1.267437in}}{\pgfqpoint{8.262430in}{1.264699in}}%
\pgfpathquadraticcurveto{\pgfqpoint{8.257458in}{1.261961in}}{\pgfqpoint{8.250470in}{1.261961in}}%
\pgfpathquadraticcurveto{\pgfqpoint{8.239050in}{1.261961in}}{\pgfqpoint{8.231863in}{1.271075in}}%
\pgfpathquadraticcurveto{\pgfqpoint{8.224694in}{1.280207in}}{\pgfqpoint{8.224694in}{1.295067in}}%
\pgfpathquadraticcurveto{\pgfqpoint{8.224694in}{1.309927in}}{\pgfqpoint{8.231863in}{1.319041in}}%
\pgfpathquadraticcurveto{\pgfqpoint{8.239050in}{1.328156in}}{\pgfqpoint{8.250470in}{1.328156in}}%
\pgfpathquadraticcurveto{\pgfqpoint{8.257458in}{1.328156in}}{\pgfqpoint{8.262430in}{1.325418in}}%
\pgfpathquadraticcurveto{\pgfqpoint{8.267419in}{1.322698in}}{\pgfqpoint{8.270679in}{1.317078in}}%
\pgfpathlineto{\pgfqpoint{8.270679in}{1.317078in}}%
\pgfpathclose%
\pgfpathmoveto{\pgfqpoint{8.235393in}{1.295067in}}%
\pgfpathquadraticcurveto{\pgfqpoint{8.235393in}{1.283648in}}{\pgfqpoint{8.240095in}{1.277145in}}%
\pgfpathquadraticcurveto{\pgfqpoint{8.244796in}{1.270643in}}{\pgfqpoint{8.253009in}{1.270643in}}%
\pgfpathquadraticcurveto{\pgfqpoint{8.261223in}{1.270643in}}{\pgfqpoint{8.265942in}{1.277145in}}%
\pgfpathquadraticcurveto{\pgfqpoint{8.270679in}{1.283648in}}{\pgfqpoint{8.270679in}{1.295067in}}%
\pgfpathquadraticcurveto{\pgfqpoint{8.270679in}{1.306487in}}{\pgfqpoint{8.265942in}{1.312989in}}%
\pgfpathquadraticcurveto{\pgfqpoint{8.261223in}{1.319492in}}{\pgfqpoint{8.253009in}{1.319492in}}%
\pgfpathquadraticcurveto{\pgfqpoint{8.244796in}{1.319492in}}{\pgfqpoint{8.240095in}{1.312989in}}%
\pgfpathquadraticcurveto{\pgfqpoint{8.235393in}{1.306487in}}{\pgfqpoint{8.235393in}{1.295067in}}%
\pgfpathlineto{\pgfqpoint{8.235393in}{1.295067in}}%
\pgfpathclose%
\pgfpathmoveto{\pgfqpoint{8.302381in}{1.326643in}}%
\pgfpathlineto{\pgfqpoint{8.312738in}{1.326643in}}%
\pgfpathlineto{\pgfqpoint{8.312738in}{1.263600in}}%
\pgfpathlineto{\pgfqpoint{8.302381in}{1.263600in}}%
\pgfpathlineto{\pgfqpoint{8.302381in}{1.326643in}}%
\pgfpathlineto{\pgfqpoint{8.302381in}{1.326643in}}%
\pgfpathclose%
\pgfpathmoveto{\pgfqpoint{8.302381in}{1.351193in}}%
\pgfpathlineto{\pgfqpoint{8.312738in}{1.351193in}}%
\pgfpathlineto{\pgfqpoint{8.312738in}{1.338062in}}%
\pgfpathlineto{\pgfqpoint{8.302381in}{1.338062in}}%
\pgfpathlineto{\pgfqpoint{8.302381in}{1.351193in}}%
\pgfpathlineto{\pgfqpoint{8.302381in}{1.351193in}}%
\pgfpathclose%
\pgfpathmoveto{\pgfqpoint{8.363064in}{1.295283in}}%
\pgfpathquadraticcurveto{\pgfqpoint{8.350509in}{1.295283in}}{\pgfqpoint{8.345664in}{1.292419in}}%
\pgfpathquadraticcurveto{\pgfqpoint{8.340819in}{1.289556in}}{\pgfqpoint{8.340819in}{1.282621in}}%
\pgfpathquadraticcurveto{\pgfqpoint{8.340819in}{1.277109in}}{\pgfqpoint{8.344457in}{1.273867in}}%
\pgfpathquadraticcurveto{\pgfqpoint{8.348095in}{1.270643in}}{\pgfqpoint{8.354328in}{1.270643in}}%
\pgfpathquadraticcurveto{\pgfqpoint{8.362956in}{1.270643in}}{\pgfqpoint{8.368161in}{1.276749in}}%
\pgfpathquadraticcurveto{\pgfqpoint{8.373367in}{1.282855in}}{\pgfqpoint{8.373367in}{1.292978in}}%
\pgfpathlineto{\pgfqpoint{8.373367in}{1.295283in}}%
\pgfpathlineto{\pgfqpoint{8.363064in}{1.295283in}}%
\pgfpathlineto{\pgfqpoint{8.363064in}{1.295283in}}%
\pgfpathclose%
\pgfpathmoveto{\pgfqpoint{8.383724in}{1.299570in}}%
\pgfpathlineto{\pgfqpoint{8.383724in}{1.263600in}}%
\pgfpathlineto{\pgfqpoint{8.373367in}{1.263600in}}%
\pgfpathlineto{\pgfqpoint{8.373367in}{1.273164in}}%
\pgfpathquadraticcurveto{\pgfqpoint{8.369818in}{1.267437in}}{\pgfqpoint{8.364523in}{1.264699in}}%
\pgfpathquadraticcurveto{\pgfqpoint{8.359227in}{1.261961in}}{\pgfqpoint{8.351572in}{1.261961in}}%
\pgfpathquadraticcurveto{\pgfqpoint{8.341899in}{1.261961in}}{\pgfqpoint{8.336171in}{1.267401in}}%
\pgfpathquadraticcurveto{\pgfqpoint{8.330462in}{1.272840in}}{\pgfqpoint{8.330462in}{1.281954in}}%
\pgfpathquadraticcurveto{\pgfqpoint{8.330462in}{1.292582in}}{\pgfqpoint{8.337576in}{1.297985in}}%
\pgfpathquadraticcurveto{\pgfqpoint{8.344709in}{1.303389in}}{\pgfqpoint{8.358831in}{1.303389in}}%
\pgfpathlineto{\pgfqpoint{8.373367in}{1.303389in}}%
\pgfpathlineto{\pgfqpoint{8.373367in}{1.304416in}}%
\pgfpathquadraticcurveto{\pgfqpoint{8.373367in}{1.311566in}}{\pgfqpoint{8.368665in}{1.315475in}}%
\pgfpathquadraticcurveto{\pgfqpoint{8.363964in}{1.319384in}}{\pgfqpoint{8.355462in}{1.319384in}}%
\pgfpathquadraticcurveto{\pgfqpoint{8.350059in}{1.319384in}}{\pgfqpoint{8.344925in}{1.318087in}}%
\pgfpathquadraticcurveto{\pgfqpoint{8.339810in}{1.316790in}}{\pgfqpoint{8.335091in}{1.314196in}}%
\pgfpathlineto{\pgfqpoint{8.335091in}{1.323779in}}%
\pgfpathquadraticcurveto{\pgfqpoint{8.340765in}{1.325976in}}{\pgfqpoint{8.346114in}{1.327057in}}%
\pgfpathquadraticcurveto{\pgfqpoint{8.351464in}{1.328156in}}{\pgfqpoint{8.356525in}{1.328156in}}%
\pgfpathquadraticcurveto{\pgfqpoint{8.370214in}{1.328156in}}{\pgfqpoint{8.376969in}{1.321059in}}%
\pgfpathquadraticcurveto{\pgfqpoint{8.383724in}{1.313980in}}{\pgfqpoint{8.383724in}{1.299570in}}%
\pgfpathlineto{\pgfqpoint{8.383724in}{1.299570in}}%
\pgfpathclose%
\pgfpathmoveto{\pgfqpoint{8.446532in}{1.295860in}}%
\pgfpathquadraticcurveto{\pgfqpoint{8.446532in}{1.307117in}}{\pgfqpoint{8.441885in}{1.313296in}}%
\pgfpathquadraticcurveto{\pgfqpoint{8.437256in}{1.319492in}}{\pgfqpoint{8.428862in}{1.319492in}}%
\pgfpathquadraticcurveto{\pgfqpoint{8.420540in}{1.319492in}}{\pgfqpoint{8.415893in}{1.313296in}}%
\pgfpathquadraticcurveto{\pgfqpoint{8.411246in}{1.307117in}}{\pgfqpoint{8.411246in}{1.295860in}}%
\pgfpathquadraticcurveto{\pgfqpoint{8.411246in}{1.284656in}}{\pgfqpoint{8.415893in}{1.278460in}}%
\pgfpathquadraticcurveto{\pgfqpoint{8.420540in}{1.272264in}}{\pgfqpoint{8.428862in}{1.272264in}}%
\pgfpathquadraticcurveto{\pgfqpoint{8.437256in}{1.272264in}}{\pgfqpoint{8.441885in}{1.278460in}}%
\pgfpathquadraticcurveto{\pgfqpoint{8.446532in}{1.284656in}}{\pgfqpoint{8.446532in}{1.295860in}}%
\pgfpathlineto{\pgfqpoint{8.446532in}{1.295860in}}%
\pgfpathclose%
\pgfpathmoveto{\pgfqpoint{8.456889in}{1.271417in}}%
\pgfpathquadraticcurveto{\pgfqpoint{8.456889in}{1.255332in}}{\pgfqpoint{8.449738in}{1.247479in}}%
\pgfpathquadraticcurveto{\pgfqpoint{8.442605in}{1.239626in}}{\pgfqpoint{8.427853in}{1.239626in}}%
\pgfpathquadraticcurveto{\pgfqpoint{8.422396in}{1.239626in}}{\pgfqpoint{8.417550in}{1.240436in}}%
\pgfpathquadraticcurveto{\pgfqpoint{8.412705in}{1.241247in}}{\pgfqpoint{8.408148in}{1.242940in}}%
\pgfpathlineto{\pgfqpoint{8.408148in}{1.253009in}}%
\pgfpathquadraticcurveto{\pgfqpoint{8.412705in}{1.250541in}}{\pgfqpoint{8.417154in}{1.249370in}}%
\pgfpathquadraticcurveto{\pgfqpoint{8.421603in}{1.248182in}}{\pgfqpoint{8.426214in}{1.248182in}}%
\pgfpathquadraticcurveto{\pgfqpoint{8.436409in}{1.248182in}}{\pgfqpoint{8.441470in}{1.253495in}}%
\pgfpathquadraticcurveto{\pgfqpoint{8.446532in}{1.258809in}}{\pgfqpoint{8.446532in}{1.269562in}}%
\pgfpathlineto{\pgfqpoint{8.446532in}{1.274695in}}%
\pgfpathquadraticcurveto{\pgfqpoint{8.443326in}{1.269112in}}{\pgfqpoint{8.438318in}{1.266356in}}%
\pgfpathquadraticcurveto{\pgfqpoint{8.433311in}{1.263600in}}{\pgfqpoint{8.426322in}{1.263600in}}%
\pgfpathquadraticcurveto{\pgfqpoint{8.414740in}{1.263600in}}{\pgfqpoint{8.407644in}{1.272426in}}%
\pgfpathquadraticcurveto{\pgfqpoint{8.400547in}{1.281270in}}{\pgfqpoint{8.400547in}{1.295860in}}%
\pgfpathquadraticcurveto{\pgfqpoint{8.400547in}{1.310486in}}{\pgfqpoint{8.407644in}{1.319312in}}%
\pgfpathquadraticcurveto{\pgfqpoint{8.414740in}{1.328156in}}{\pgfqpoint{8.426322in}{1.328156in}}%
\pgfpathquadraticcurveto{\pgfqpoint{8.433311in}{1.328156in}}{\pgfqpoint{8.438318in}{1.325400in}}%
\pgfpathquadraticcurveto{\pgfqpoint{8.443326in}{1.322644in}}{\pgfqpoint{8.446532in}{1.317078in}}%
\pgfpathlineto{\pgfqpoint{8.446532in}{1.326643in}}%
\pgfpathlineto{\pgfqpoint{8.456889in}{1.326643in}}%
\pgfpathlineto{\pgfqpoint{8.456889in}{1.271417in}}%
\pgfpathlineto{\pgfqpoint{8.456889in}{1.271417in}}%
\pgfpathclose%
\pgfpathmoveto{\pgfqpoint{8.530649in}{1.301660in}}%
\pgfpathlineto{\pgfqpoint{8.530649in}{1.263600in}}%
\pgfpathlineto{\pgfqpoint{8.520292in}{1.263600in}}%
\pgfpathlineto{\pgfqpoint{8.520292in}{1.301317in}}%
\pgfpathquadraticcurveto{\pgfqpoint{8.520292in}{1.310269in}}{\pgfqpoint{8.516797in}{1.314700in}}%
\pgfpathquadraticcurveto{\pgfqpoint{8.513303in}{1.319149in}}{\pgfqpoint{8.506332in}{1.319149in}}%
\pgfpathquadraticcurveto{\pgfqpoint{8.497939in}{1.319149in}}{\pgfqpoint{8.493093in}{1.313800in}}%
\pgfpathquadraticcurveto{\pgfqpoint{8.488248in}{1.308468in}}{\pgfqpoint{8.488248in}{1.299228in}}%
\pgfpathlineto{\pgfqpoint{8.488248in}{1.263600in}}%
\pgfpathlineto{\pgfqpoint{8.477837in}{1.263600in}}%
\pgfpathlineto{\pgfqpoint{8.477837in}{1.326643in}}%
\pgfpathlineto{\pgfqpoint{8.488248in}{1.326643in}}%
\pgfpathlineto{\pgfqpoint{8.488248in}{1.316844in}}%
\pgfpathquadraticcurveto{\pgfqpoint{8.491977in}{1.322536in}}{\pgfqpoint{8.497002in}{1.325346in}}%
\pgfpathquadraticcurveto{\pgfqpoint{8.502045in}{1.328156in}}{\pgfqpoint{8.508638in}{1.328156in}}%
\pgfpathquadraticcurveto{\pgfqpoint{8.519499in}{1.328156in}}{\pgfqpoint{8.525065in}{1.321437in}}%
\pgfpathquadraticcurveto{\pgfqpoint{8.530649in}{1.314718in}}{\pgfqpoint{8.530649in}{1.301660in}}%
\pgfpathlineto{\pgfqpoint{8.530649in}{1.301660in}}%
\pgfpathclose%
\pgfpathmoveto{\pgfqpoint{8.575715in}{1.319384in}}%
\pgfpathquadraticcurveto{\pgfqpoint{8.567393in}{1.319384in}}{\pgfqpoint{8.562548in}{1.312881in}}%
\pgfpathquadraticcurveto{\pgfqpoint{8.557703in}{1.306379in}}{\pgfqpoint{8.557703in}{1.295067in}}%
\pgfpathquadraticcurveto{\pgfqpoint{8.557703in}{1.283756in}}{\pgfqpoint{8.562512in}{1.277253in}}%
\pgfpathquadraticcurveto{\pgfqpoint{8.567339in}{1.270751in}}{\pgfqpoint{8.575715in}{1.270751in}}%
\pgfpathquadraticcurveto{\pgfqpoint{8.584001in}{1.270751in}}{\pgfqpoint{8.588828in}{1.277271in}}%
\pgfpathquadraticcurveto{\pgfqpoint{8.593673in}{1.283810in}}{\pgfqpoint{8.593673in}{1.295067in}}%
\pgfpathquadraticcurveto{\pgfqpoint{8.593673in}{1.306271in}}{\pgfqpoint{8.588828in}{1.312827in}}%
\pgfpathquadraticcurveto{\pgfqpoint{8.584001in}{1.319384in}}{\pgfqpoint{8.575715in}{1.319384in}}%
\pgfpathlineto{\pgfqpoint{8.575715in}{1.319384in}}%
\pgfpathclose%
\pgfpathmoveto{\pgfqpoint{8.575715in}{1.328156in}}%
\pgfpathquadraticcurveto{\pgfqpoint{8.589224in}{1.328156in}}{\pgfqpoint{8.596933in}{1.319366in}}%
\pgfpathquadraticcurveto{\pgfqpoint{8.604661in}{1.310594in}}{\pgfqpoint{8.604661in}{1.295067in}}%
\pgfpathquadraticcurveto{\pgfqpoint{8.604661in}{1.279595in}}{\pgfqpoint{8.596933in}{1.270769in}}%
\pgfpathquadraticcurveto{\pgfqpoint{8.589224in}{1.261961in}}{\pgfqpoint{8.575715in}{1.261961in}}%
\pgfpathquadraticcurveto{\pgfqpoint{8.562152in}{1.261961in}}{\pgfqpoint{8.554461in}{1.270769in}}%
\pgfpathquadraticcurveto{\pgfqpoint{8.546788in}{1.279595in}}{\pgfqpoint{8.546788in}{1.295067in}}%
\pgfpathquadraticcurveto{\pgfqpoint{8.546788in}{1.310594in}}{\pgfqpoint{8.554461in}{1.319366in}}%
\pgfpathquadraticcurveto{\pgfqpoint{8.562152in}{1.328156in}}{\pgfqpoint{8.575715in}{1.328156in}}%
\pgfpathlineto{\pgfqpoint{8.575715in}{1.328156in}}%
\pgfpathclose%
\pgfpathmoveto{\pgfqpoint{8.662011in}{1.324787in}}%
\pgfpathlineto{\pgfqpoint{8.662011in}{1.314989in}}%
\pgfpathquadraticcurveto{\pgfqpoint{8.657634in}{1.317240in}}{\pgfqpoint{8.652897in}{1.318357in}}%
\pgfpathquadraticcurveto{\pgfqpoint{8.648178in}{1.319492in}}{\pgfqpoint{8.643099in}{1.319492in}}%
\pgfpathquadraticcurveto{\pgfqpoint{8.635389in}{1.319492in}}{\pgfqpoint{8.631535in}{1.317132in}}%
\pgfpathquadraticcurveto{\pgfqpoint{8.627680in}{1.314773in}}{\pgfqpoint{8.627680in}{1.310035in}}%
\pgfpathquadraticcurveto{\pgfqpoint{8.627680in}{1.306433in}}{\pgfqpoint{8.630436in}{1.304380in}}%
\pgfpathquadraticcurveto{\pgfqpoint{8.633192in}{1.302326in}}{\pgfqpoint{8.641531in}{1.300471in}}%
\pgfpathlineto{\pgfqpoint{8.645080in}{1.299678in}}%
\pgfpathquadraticcurveto{\pgfqpoint{8.656103in}{1.297319in}}{\pgfqpoint{8.660750in}{1.293014in}}%
\pgfpathquadraticcurveto{\pgfqpoint{8.665398in}{1.288709in}}{\pgfqpoint{8.665398in}{1.281000in}}%
\pgfpathquadraticcurveto{\pgfqpoint{8.665398in}{1.272210in}}{\pgfqpoint{8.658445in}{1.267076in}}%
\pgfpathquadraticcurveto{\pgfqpoint{8.651492in}{1.261961in}}{\pgfqpoint{8.639334in}{1.261961in}}%
\pgfpathquadraticcurveto{\pgfqpoint{8.634273in}{1.261961in}}{\pgfqpoint{8.628779in}{1.262952in}}%
\pgfpathquadraticcurveto{\pgfqpoint{8.623285in}{1.263942in}}{\pgfqpoint{8.617215in}{1.265906in}}%
\pgfpathlineto{\pgfqpoint{8.617215in}{1.276605in}}%
\pgfpathquadraticcurveto{\pgfqpoint{8.622961in}{1.273615in}}{\pgfqpoint{8.628527in}{1.272120in}}%
\pgfpathquadraticcurveto{\pgfqpoint{8.634092in}{1.270643in}}{\pgfqpoint{8.639568in}{1.270643in}}%
\pgfpathquadraticcurveto{\pgfqpoint{8.646881in}{1.270643in}}{\pgfqpoint{8.650808in}{1.273146in}}%
\pgfpathquadraticcurveto{\pgfqpoint{8.654752in}{1.275650in}}{\pgfqpoint{8.654752in}{1.280207in}}%
\pgfpathquadraticcurveto{\pgfqpoint{8.654752in}{1.284422in}}{\pgfqpoint{8.651906in}{1.286674in}}%
\pgfpathquadraticcurveto{\pgfqpoint{8.649079in}{1.288925in}}{\pgfqpoint{8.639442in}{1.291014in}}%
\pgfpathlineto{\pgfqpoint{8.635840in}{1.291861in}}%
\pgfpathquadraticcurveto{\pgfqpoint{8.626221in}{1.293878in}}{\pgfqpoint{8.621934in}{1.298075in}}%
\pgfpathquadraticcurveto{\pgfqpoint{8.617665in}{1.302272in}}{\pgfqpoint{8.617665in}{1.309585in}}%
\pgfpathquadraticcurveto{\pgfqpoint{8.617665in}{1.318483in}}{\pgfqpoint{8.623970in}{1.323310in}}%
\pgfpathquadraticcurveto{\pgfqpoint{8.630274in}{1.328156in}}{\pgfqpoint{8.641874in}{1.328156in}}%
\pgfpathquadraticcurveto{\pgfqpoint{8.647602in}{1.328156in}}{\pgfqpoint{8.652663in}{1.327309in}}%
\pgfpathquadraticcurveto{\pgfqpoint{8.657742in}{1.326480in}}{\pgfqpoint{8.662011in}{1.324787in}}%
\pgfpathlineto{\pgfqpoint{8.662011in}{1.324787in}}%
\pgfpathclose%
\pgfpathmoveto{\pgfqpoint{8.692128in}{1.344547in}}%
\pgfpathlineto{\pgfqpoint{8.692128in}{1.326643in}}%
\pgfpathlineto{\pgfqpoint{8.713454in}{1.326643in}}%
\pgfpathlineto{\pgfqpoint{8.713454in}{1.318591in}}%
\pgfpathlineto{\pgfqpoint{8.692128in}{1.318591in}}%
\pgfpathlineto{\pgfqpoint{8.692128in}{1.284368in}}%
\pgfpathquadraticcurveto{\pgfqpoint{8.692128in}{1.276659in}}{\pgfqpoint{8.694235in}{1.274461in}}%
\pgfpathquadraticcurveto{\pgfqpoint{8.696342in}{1.272264in}}{\pgfqpoint{8.702827in}{1.272264in}}%
\pgfpathlineto{\pgfqpoint{8.713454in}{1.272264in}}%
\pgfpathlineto{\pgfqpoint{8.713454in}{1.263600in}}%
\pgfpathlineto{\pgfqpoint{8.702827in}{1.263600in}}%
\pgfpathquadraticcurveto{\pgfqpoint{8.690831in}{1.263600in}}{\pgfqpoint{8.686274in}{1.268067in}}%
\pgfpathquadraticcurveto{\pgfqpoint{8.681717in}{1.272552in}}{\pgfqpoint{8.681717in}{1.284368in}}%
\pgfpathlineto{\pgfqpoint{8.681717in}{1.318591in}}%
\pgfpathlineto{\pgfqpoint{8.674115in}{1.318591in}}%
\pgfpathlineto{\pgfqpoint{8.674115in}{1.326643in}}%
\pgfpathlineto{\pgfqpoint{8.681717in}{1.326643in}}%
\pgfpathlineto{\pgfqpoint{8.681717in}{1.344547in}}%
\pgfpathlineto{\pgfqpoint{8.692128in}{1.344547in}}%
\pgfpathlineto{\pgfqpoint{8.692128in}{1.344547in}}%
\pgfpathclose%
\pgfpathmoveto{\pgfqpoint{8.727071in}{1.326643in}}%
\pgfpathlineto{\pgfqpoint{8.737428in}{1.326643in}}%
\pgfpathlineto{\pgfqpoint{8.737428in}{1.263600in}}%
\pgfpathlineto{\pgfqpoint{8.727071in}{1.263600in}}%
\pgfpathlineto{\pgfqpoint{8.727071in}{1.326643in}}%
\pgfpathlineto{\pgfqpoint{8.727071in}{1.326643in}}%
\pgfpathclose%
\pgfpathmoveto{\pgfqpoint{8.727071in}{1.351193in}}%
\pgfpathlineto{\pgfqpoint{8.737428in}{1.351193in}}%
\pgfpathlineto{\pgfqpoint{8.737428in}{1.338062in}}%
\pgfpathlineto{\pgfqpoint{8.727071in}{1.338062in}}%
\pgfpathlineto{\pgfqpoint{8.727071in}{1.351193in}}%
\pgfpathlineto{\pgfqpoint{8.727071in}{1.351193in}}%
\pgfpathclose%
\pgfpathmoveto{\pgfqpoint{8.804469in}{1.324229in}}%
\pgfpathlineto{\pgfqpoint{8.804469in}{1.314538in}}%
\pgfpathquadraticcurveto{\pgfqpoint{8.800074in}{1.316970in}}{\pgfqpoint{8.795661in}{1.318177in}}%
\pgfpathquadraticcurveto{\pgfqpoint{8.791248in}{1.319384in}}{\pgfqpoint{8.786745in}{1.319384in}}%
\pgfpathquadraticcurveto{\pgfqpoint{8.776659in}{1.319384in}}{\pgfqpoint{8.771075in}{1.312989in}}%
\pgfpathquadraticcurveto{\pgfqpoint{8.765509in}{1.306613in}}{\pgfqpoint{8.765509in}{1.295067in}}%
\pgfpathquadraticcurveto{\pgfqpoint{8.765509in}{1.283521in}}{\pgfqpoint{8.771075in}{1.277127in}}%
\pgfpathquadraticcurveto{\pgfqpoint{8.776659in}{1.270751in}}{\pgfqpoint{8.786745in}{1.270751in}}%
\pgfpathquadraticcurveto{\pgfqpoint{8.791248in}{1.270751in}}{\pgfqpoint{8.795661in}{1.271958in}}%
\pgfpathquadraticcurveto{\pgfqpoint{8.800074in}{1.273164in}}{\pgfqpoint{8.804469in}{1.275596in}}%
\pgfpathlineto{\pgfqpoint{8.804469in}{1.266014in}}%
\pgfpathquadraticcurveto{\pgfqpoint{8.800128in}{1.263996in}}{\pgfqpoint{8.795481in}{1.262988in}}%
\pgfpathquadraticcurveto{\pgfqpoint{8.790852in}{1.261961in}}{\pgfqpoint{8.785611in}{1.261961in}}%
\pgfpathquadraticcurveto{\pgfqpoint{8.771363in}{1.261961in}}{\pgfqpoint{8.762969in}{1.270913in}}%
\pgfpathquadraticcurveto{\pgfqpoint{8.754594in}{1.279865in}}{\pgfqpoint{8.754594in}{1.295067in}}%
\pgfpathquadraticcurveto{\pgfqpoint{8.754594in}{1.310486in}}{\pgfqpoint{8.763059in}{1.319312in}}%
\pgfpathquadraticcurveto{\pgfqpoint{8.771543in}{1.328156in}}{\pgfqpoint{8.786295in}{1.328156in}}%
\pgfpathquadraticcurveto{\pgfqpoint{8.791068in}{1.328156in}}{\pgfqpoint{8.795625in}{1.327165in}}%
\pgfpathquadraticcurveto{\pgfqpoint{8.800183in}{1.326192in}}{\pgfqpoint{8.804469in}{1.324229in}}%
\pgfpathlineto{\pgfqpoint{8.804469in}{1.324229in}}%
\pgfpathclose%
\pgfpathmoveto{\pgfqpoint{8.862667in}{1.324787in}}%
\pgfpathlineto{\pgfqpoint{8.862667in}{1.314989in}}%
\pgfpathquadraticcurveto{\pgfqpoint{8.858290in}{1.317240in}}{\pgfqpoint{8.853553in}{1.318357in}}%
\pgfpathquadraticcurveto{\pgfqpoint{8.848833in}{1.319492in}}{\pgfqpoint{8.843754in}{1.319492in}}%
\pgfpathquadraticcurveto{\pgfqpoint{8.836045in}{1.319492in}}{\pgfqpoint{8.832190in}{1.317132in}}%
\pgfpathquadraticcurveto{\pgfqpoint{8.828336in}{1.314773in}}{\pgfqpoint{8.828336in}{1.310035in}}%
\pgfpathquadraticcurveto{\pgfqpoint{8.828336in}{1.306433in}}{\pgfqpoint{8.831091in}{1.304380in}}%
\pgfpathquadraticcurveto{\pgfqpoint{8.833847in}{1.302326in}}{\pgfqpoint{8.842187in}{1.300471in}}%
\pgfpathlineto{\pgfqpoint{8.845735in}{1.299678in}}%
\pgfpathquadraticcurveto{\pgfqpoint{8.856759in}{1.297319in}}{\pgfqpoint{8.861406in}{1.293014in}}%
\pgfpathquadraticcurveto{\pgfqpoint{8.866053in}{1.288709in}}{\pgfqpoint{8.866053in}{1.281000in}}%
\pgfpathquadraticcurveto{\pgfqpoint{8.866053in}{1.272210in}}{\pgfqpoint{8.859100in}{1.267076in}}%
\pgfpathquadraticcurveto{\pgfqpoint{8.852148in}{1.261961in}}{\pgfqpoint{8.839989in}{1.261961in}}%
\pgfpathquadraticcurveto{\pgfqpoint{8.834928in}{1.261961in}}{\pgfqpoint{8.829434in}{1.262952in}}%
\pgfpathquadraticcurveto{\pgfqpoint{8.823941in}{1.263942in}}{\pgfqpoint{8.817870in}{1.265906in}}%
\pgfpathlineto{\pgfqpoint{8.817870in}{1.276605in}}%
\pgfpathquadraticcurveto{\pgfqpoint{8.823616in}{1.273615in}}{\pgfqpoint{8.829182in}{1.272120in}}%
\pgfpathquadraticcurveto{\pgfqpoint{8.834748in}{1.270643in}}{\pgfqpoint{8.840224in}{1.270643in}}%
\pgfpathquadraticcurveto{\pgfqpoint{8.847536in}{1.270643in}}{\pgfqpoint{8.851463in}{1.273146in}}%
\pgfpathquadraticcurveto{\pgfqpoint{8.855408in}{1.275650in}}{\pgfqpoint{8.855408in}{1.280207in}}%
\pgfpathquadraticcurveto{\pgfqpoint{8.855408in}{1.284422in}}{\pgfqpoint{8.852562in}{1.286674in}}%
\pgfpathquadraticcurveto{\pgfqpoint{8.849734in}{1.288925in}}{\pgfqpoint{8.840097in}{1.291014in}}%
\pgfpathlineto{\pgfqpoint{8.836495in}{1.291861in}}%
\pgfpathquadraticcurveto{\pgfqpoint{8.826877in}{1.293878in}}{\pgfqpoint{8.822590in}{1.298075in}}%
\pgfpathquadraticcurveto{\pgfqpoint{8.818321in}{1.302272in}}{\pgfqpoint{8.818321in}{1.309585in}}%
\pgfpathquadraticcurveto{\pgfqpoint{8.818321in}{1.318483in}}{\pgfqpoint{8.824625in}{1.323310in}}%
\pgfpathquadraticcurveto{\pgfqpoint{8.830929in}{1.328156in}}{\pgfqpoint{8.842529in}{1.328156in}}%
\pgfpathquadraticcurveto{\pgfqpoint{8.848257in}{1.328156in}}{\pgfqpoint{8.853318in}{1.327309in}}%
\pgfpathquadraticcurveto{\pgfqpoint{8.858398in}{1.326480in}}{\pgfqpoint{8.862667in}{1.324787in}}%
\pgfpathlineto{\pgfqpoint{8.862667in}{1.324787in}}%
\pgfpathclose%
\pgfpathmoveto{\pgfqpoint{8.885182in}{1.277902in}}%
\pgfpathlineto{\pgfqpoint{8.897052in}{1.277902in}}%
\pgfpathlineto{\pgfqpoint{8.897052in}{1.268211in}}%
\pgfpathlineto{\pgfqpoint{8.887830in}{1.250199in}}%
\pgfpathlineto{\pgfqpoint{8.880571in}{1.250199in}}%
\pgfpathlineto{\pgfqpoint{8.885182in}{1.268211in}}%
\pgfpathlineto{\pgfqpoint{8.885182in}{1.277902in}}%
\pgfpathlineto{\pgfqpoint{8.885182in}{1.277902in}}%
\pgfpathclose%
\pgfusepath{fill}%
\end{pgfscope}%
\begin{pgfscope}%
\pgfsetbuttcap%
\pgfsetroundjoin%
\definecolor{currentfill}{rgb}{0.309804,0.372549,0.419608}%
\pgfsetfillcolor{currentfill}%
\pgfsetlinewidth{0.000000pt}%
\definecolor{currentstroke}{rgb}{0.309804,0.372549,0.419608}%
\pgfsetstrokecolor{currentstroke}%
\pgfsetdash{}{0pt}%
\pgfpathmoveto{\pgfqpoint{9.057298in}{1.301660in}}%
\pgfpathlineto{\pgfqpoint{9.057298in}{1.263600in}}%
\pgfpathlineto{\pgfqpoint{9.046941in}{1.263600in}}%
\pgfpathlineto{\pgfqpoint{9.046941in}{1.301317in}}%
\pgfpathquadraticcurveto{\pgfqpoint{9.046941in}{1.310269in}}{\pgfqpoint{9.043446in}{1.314700in}}%
\pgfpathquadraticcurveto{\pgfqpoint{9.039952in}{1.319149in}}{\pgfqpoint{9.032981in}{1.319149in}}%
\pgfpathquadraticcurveto{\pgfqpoint{9.024587in}{1.319149in}}{\pgfqpoint{9.019742in}{1.313800in}}%
\pgfpathquadraticcurveto{\pgfqpoint{9.014897in}{1.308468in}}{\pgfqpoint{9.014897in}{1.299228in}}%
\pgfpathlineto{\pgfqpoint{9.014897in}{1.263600in}}%
\pgfpathlineto{\pgfqpoint{9.004486in}{1.263600in}}%
\pgfpathlineto{\pgfqpoint{9.004486in}{1.326643in}}%
\pgfpathlineto{\pgfqpoint{9.014897in}{1.326643in}}%
\pgfpathlineto{\pgfqpoint{9.014897in}{1.316844in}}%
\pgfpathquadraticcurveto{\pgfqpoint{9.018625in}{1.322536in}}{\pgfqpoint{9.023651in}{1.325346in}}%
\pgfpathquadraticcurveto{\pgfqpoint{9.028694in}{1.328156in}}{\pgfqpoint{9.035287in}{1.328156in}}%
\pgfpathquadraticcurveto{\pgfqpoint{9.046148in}{1.328156in}}{\pgfqpoint{9.051714in}{1.321437in}}%
\pgfpathquadraticcurveto{\pgfqpoint{9.057298in}{1.314718in}}{\pgfqpoint{9.057298in}{1.301660in}}%
\pgfpathlineto{\pgfqpoint{9.057298in}{1.301660in}}%
\pgfpathclose%
\pgfpathmoveto{\pgfqpoint{9.102364in}{1.319384in}}%
\pgfpathquadraticcurveto{\pgfqpoint{9.094042in}{1.319384in}}{\pgfqpoint{9.089197in}{1.312881in}}%
\pgfpathquadraticcurveto{\pgfqpoint{9.084352in}{1.306379in}}{\pgfqpoint{9.084352in}{1.295067in}}%
\pgfpathquadraticcurveto{\pgfqpoint{9.084352in}{1.283756in}}{\pgfqpoint{9.089161in}{1.277253in}}%
\pgfpathquadraticcurveto{\pgfqpoint{9.093988in}{1.270751in}}{\pgfqpoint{9.102364in}{1.270751in}}%
\pgfpathquadraticcurveto{\pgfqpoint{9.110650in}{1.270751in}}{\pgfqpoint{9.115477in}{1.277271in}}%
\pgfpathquadraticcurveto{\pgfqpoint{9.120322in}{1.283810in}}{\pgfqpoint{9.120322in}{1.295067in}}%
\pgfpathquadraticcurveto{\pgfqpoint{9.120322in}{1.306271in}}{\pgfqpoint{9.115477in}{1.312827in}}%
\pgfpathquadraticcurveto{\pgfqpoint{9.110650in}{1.319384in}}{\pgfqpoint{9.102364in}{1.319384in}}%
\pgfpathlineto{\pgfqpoint{9.102364in}{1.319384in}}%
\pgfpathclose%
\pgfpathmoveto{\pgfqpoint{9.102364in}{1.328156in}}%
\pgfpathquadraticcurveto{\pgfqpoint{9.115873in}{1.328156in}}{\pgfqpoint{9.123582in}{1.319366in}}%
\pgfpathquadraticcurveto{\pgfqpoint{9.131309in}{1.310594in}}{\pgfqpoint{9.131309in}{1.295067in}}%
\pgfpathquadraticcurveto{\pgfqpoint{9.131309in}{1.279595in}}{\pgfqpoint{9.123582in}{1.270769in}}%
\pgfpathquadraticcurveto{\pgfqpoint{9.115873in}{1.261961in}}{\pgfqpoint{9.102364in}{1.261961in}}%
\pgfpathquadraticcurveto{\pgfqpoint{9.088801in}{1.261961in}}{\pgfqpoint{9.081110in}{1.270769in}}%
\pgfpathquadraticcurveto{\pgfqpoint{9.073436in}{1.279595in}}{\pgfqpoint{9.073436in}{1.295067in}}%
\pgfpathquadraticcurveto{\pgfqpoint{9.073436in}{1.310594in}}{\pgfqpoint{9.081110in}{1.319366in}}%
\pgfpathquadraticcurveto{\pgfqpoint{9.088801in}{1.328156in}}{\pgfqpoint{9.102364in}{1.328156in}}%
\pgfpathlineto{\pgfqpoint{9.102364in}{1.328156in}}%
\pgfpathclose%
\pgfpathmoveto{\pgfqpoint{9.158724in}{1.344547in}}%
\pgfpathlineto{\pgfqpoint{9.158724in}{1.326643in}}%
\pgfpathlineto{\pgfqpoint{9.180050in}{1.326643in}}%
\pgfpathlineto{\pgfqpoint{9.180050in}{1.318591in}}%
\pgfpathlineto{\pgfqpoint{9.158724in}{1.318591in}}%
\pgfpathlineto{\pgfqpoint{9.158724in}{1.284368in}}%
\pgfpathquadraticcurveto{\pgfqpoint{9.158724in}{1.276659in}}{\pgfqpoint{9.160831in}{1.274461in}}%
\pgfpathquadraticcurveto{\pgfqpoint{9.162939in}{1.272264in}}{\pgfqpoint{9.169423in}{1.272264in}}%
\pgfpathlineto{\pgfqpoint{9.180050in}{1.272264in}}%
\pgfpathlineto{\pgfqpoint{9.180050in}{1.263600in}}%
\pgfpathlineto{\pgfqpoint{9.169423in}{1.263600in}}%
\pgfpathquadraticcurveto{\pgfqpoint{9.157427in}{1.263600in}}{\pgfqpoint{9.152870in}{1.268067in}}%
\pgfpathquadraticcurveto{\pgfqpoint{9.148313in}{1.272552in}}{\pgfqpoint{9.148313in}{1.284368in}}%
\pgfpathlineto{\pgfqpoint{9.148313in}{1.318591in}}%
\pgfpathlineto{\pgfqpoint{9.140712in}{1.318591in}}%
\pgfpathlineto{\pgfqpoint{9.140712in}{1.326643in}}%
\pgfpathlineto{\pgfqpoint{9.148313in}{1.326643in}}%
\pgfpathlineto{\pgfqpoint{9.148313in}{1.344547in}}%
\pgfpathlineto{\pgfqpoint{9.158724in}{1.344547in}}%
\pgfpathlineto{\pgfqpoint{9.158724in}{1.344547in}}%
\pgfpathclose%
\pgfusepath{fill}%
\end{pgfscope}%
\begin{pgfscope}%
\pgfsetbuttcap%
\pgfsetroundjoin%
\definecolor{currentfill}{rgb}{0.309804,0.372549,0.419608}%
\pgfsetfillcolor{currentfill}%
\pgfsetlinewidth{0.000000pt}%
\definecolor{currentstroke}{rgb}{0.309804,0.372549,0.419608}%
\pgfsetstrokecolor{currentstroke}%
\pgfsetdash{}{0pt}%
\pgfpathmoveto{\pgfqpoint{9.270816in}{1.344547in}}%
\pgfpathlineto{\pgfqpoint{9.270816in}{1.326643in}}%
\pgfpathlineto{\pgfqpoint{9.292142in}{1.326643in}}%
\pgfpathlineto{\pgfqpoint{9.292142in}{1.318591in}}%
\pgfpathlineto{\pgfqpoint{9.270816in}{1.318591in}}%
\pgfpathlineto{\pgfqpoint{9.270816in}{1.284368in}}%
\pgfpathquadraticcurveto{\pgfqpoint{9.270816in}{1.276659in}}{\pgfqpoint{9.272923in}{1.274461in}}%
\pgfpathquadraticcurveto{\pgfqpoint{9.275031in}{1.272264in}}{\pgfqpoint{9.281515in}{1.272264in}}%
\pgfpathlineto{\pgfqpoint{9.292142in}{1.272264in}}%
\pgfpathlineto{\pgfqpoint{9.292142in}{1.263600in}}%
\pgfpathlineto{\pgfqpoint{9.281515in}{1.263600in}}%
\pgfpathquadraticcurveto{\pgfqpoint{9.269519in}{1.263600in}}{\pgfqpoint{9.264962in}{1.268067in}}%
\pgfpathquadraticcurveto{\pgfqpoint{9.260405in}{1.272552in}}{\pgfqpoint{9.260405in}{1.284368in}}%
\pgfpathlineto{\pgfqpoint{9.260405in}{1.318591in}}%
\pgfpathlineto{\pgfqpoint{9.252804in}{1.318591in}}%
\pgfpathlineto{\pgfqpoint{9.252804in}{1.326643in}}%
\pgfpathlineto{\pgfqpoint{9.260405in}{1.326643in}}%
\pgfpathlineto{\pgfqpoint{9.260405in}{1.344547in}}%
\pgfpathlineto{\pgfqpoint{9.270816in}{1.344547in}}%
\pgfpathlineto{\pgfqpoint{9.270816in}{1.344547in}}%
\pgfpathclose%
\pgfpathmoveto{\pgfqpoint{9.342288in}{1.316970in}}%
\pgfpathquadraticcurveto{\pgfqpoint{9.340541in}{1.317979in}}{\pgfqpoint{9.338488in}{1.318447in}}%
\pgfpathquadraticcurveto{\pgfqpoint{9.336434in}{1.318933in}}{\pgfqpoint{9.333967in}{1.318933in}}%
\pgfpathquadraticcurveto{\pgfqpoint{9.325177in}{1.318933in}}{\pgfqpoint{9.320475in}{1.313223in}}%
\pgfpathquadraticcurveto{\pgfqpoint{9.315774in}{1.307514in}}{\pgfqpoint{9.315774in}{1.296814in}}%
\pgfpathlineto{\pgfqpoint{9.315774in}{1.263600in}}%
\pgfpathlineto{\pgfqpoint{9.305363in}{1.263600in}}%
\pgfpathlineto{\pgfqpoint{9.305363in}{1.326643in}}%
\pgfpathlineto{\pgfqpoint{9.315774in}{1.326643in}}%
\pgfpathlineto{\pgfqpoint{9.315774in}{1.316844in}}%
\pgfpathquadraticcurveto{\pgfqpoint{9.319053in}{1.322590in}}{\pgfqpoint{9.324276in}{1.325364in}}%
\pgfpathquadraticcurveto{\pgfqpoint{9.329518in}{1.328156in}}{\pgfqpoint{9.337011in}{1.328156in}}%
\pgfpathquadraticcurveto{\pgfqpoint{9.338073in}{1.328156in}}{\pgfqpoint{9.339370in}{1.328011in}}%
\pgfpathquadraticcurveto{\pgfqpoint{9.340667in}{1.327885in}}{\pgfqpoint{9.342234in}{1.327597in}}%
\pgfpathlineto{\pgfqpoint{9.342288in}{1.316970in}}%
\pgfpathlineto{\pgfqpoint{9.342288in}{1.316970in}}%
\pgfpathclose%
\pgfpathmoveto{\pgfqpoint{9.404538in}{1.297715in}}%
\pgfpathlineto{\pgfqpoint{9.404538in}{1.292654in}}%
\pgfpathlineto{\pgfqpoint{9.356914in}{1.292654in}}%
\pgfpathquadraticcurveto{\pgfqpoint{9.357599in}{1.281954in}}{\pgfqpoint{9.363362in}{1.276353in}}%
\pgfpathquadraticcurveto{\pgfqpoint{9.369126in}{1.270751in}}{\pgfqpoint{9.379429in}{1.270751in}}%
\pgfpathquadraticcurveto{\pgfqpoint{9.385391in}{1.270751in}}{\pgfqpoint{9.390993in}{1.272210in}}%
\pgfpathquadraticcurveto{\pgfqpoint{9.396595in}{1.273669in}}{\pgfqpoint{9.402125in}{1.276605in}}%
\pgfpathlineto{\pgfqpoint{9.402125in}{1.266806in}}%
\pgfpathquadraticcurveto{\pgfqpoint{9.396541in}{1.264447in}}{\pgfqpoint{9.390687in}{1.263204in}}%
\pgfpathquadraticcurveto{\pgfqpoint{9.384833in}{1.261961in}}{\pgfqpoint{9.378817in}{1.261961in}}%
\pgfpathquadraticcurveto{\pgfqpoint{9.363723in}{1.261961in}}{\pgfqpoint{9.354915in}{1.270733in}}%
\pgfpathquadraticcurveto{\pgfqpoint{9.346107in}{1.279523in}}{\pgfqpoint{9.346107in}{1.294509in}}%
\pgfpathquadraticcurveto{\pgfqpoint{9.346107in}{1.309981in}}{\pgfqpoint{9.354464in}{1.319059in}}%
\pgfpathquadraticcurveto{\pgfqpoint{9.362822in}{1.328156in}}{\pgfqpoint{9.377016in}{1.328156in}}%
\pgfpathquadraticcurveto{\pgfqpoint{9.389732in}{1.328156in}}{\pgfqpoint{9.397135in}{1.319960in}}%
\pgfpathquadraticcurveto{\pgfqpoint{9.404538in}{1.311783in}}{\pgfqpoint{9.404538in}{1.297715in}}%
\pgfpathlineto{\pgfqpoint{9.404538in}{1.297715in}}%
\pgfpathclose%
\pgfpathmoveto{\pgfqpoint{9.394181in}{1.300759in}}%
\pgfpathquadraticcurveto{\pgfqpoint{9.394073in}{1.309243in}}{\pgfqpoint{9.389426in}{1.314304in}}%
\pgfpathquadraticcurveto{\pgfqpoint{9.384779in}{1.319384in}}{\pgfqpoint{9.377124in}{1.319384in}}%
\pgfpathquadraticcurveto{\pgfqpoint{9.368460in}{1.319384in}}{\pgfqpoint{9.363254in}{1.314484in}}%
\pgfpathquadraticcurveto{\pgfqpoint{9.358049in}{1.309585in}}{\pgfqpoint{9.357256in}{1.300687in}}%
\pgfpathlineto{\pgfqpoint{9.394181in}{1.300759in}}%
\pgfpathlineto{\pgfqpoint{9.394181in}{1.300759in}}%
\pgfpathclose%
\pgfpathmoveto{\pgfqpoint{9.450199in}{1.295283in}}%
\pgfpathquadraticcurveto{\pgfqpoint{9.437645in}{1.295283in}}{\pgfqpoint{9.432799in}{1.292419in}}%
\pgfpathquadraticcurveto{\pgfqpoint{9.427954in}{1.289556in}}{\pgfqpoint{9.427954in}{1.282621in}}%
\pgfpathquadraticcurveto{\pgfqpoint{9.427954in}{1.277109in}}{\pgfqpoint{9.431592in}{1.273867in}}%
\pgfpathquadraticcurveto{\pgfqpoint{9.435231in}{1.270643in}}{\pgfqpoint{9.441463in}{1.270643in}}%
\pgfpathquadraticcurveto{\pgfqpoint{9.450091in}{1.270643in}}{\pgfqpoint{9.455296in}{1.276749in}}%
\pgfpathquadraticcurveto{\pgfqpoint{9.460502in}{1.282855in}}{\pgfqpoint{9.460502in}{1.292978in}}%
\pgfpathlineto{\pgfqpoint{9.460502in}{1.295283in}}%
\pgfpathlineto{\pgfqpoint{9.450199in}{1.295283in}}%
\pgfpathlineto{\pgfqpoint{9.450199in}{1.295283in}}%
\pgfpathclose%
\pgfpathmoveto{\pgfqpoint{9.470859in}{1.299570in}}%
\pgfpathlineto{\pgfqpoint{9.470859in}{1.263600in}}%
\pgfpathlineto{\pgfqpoint{9.460502in}{1.263600in}}%
\pgfpathlineto{\pgfqpoint{9.460502in}{1.273164in}}%
\pgfpathquadraticcurveto{\pgfqpoint{9.456954in}{1.267437in}}{\pgfqpoint{9.451658in}{1.264699in}}%
\pgfpathquadraticcurveto{\pgfqpoint{9.446362in}{1.261961in}}{\pgfqpoint{9.438707in}{1.261961in}}%
\pgfpathquadraticcurveto{\pgfqpoint{9.429035in}{1.261961in}}{\pgfqpoint{9.423307in}{1.267401in}}%
\pgfpathquadraticcurveto{\pgfqpoint{9.417597in}{1.272840in}}{\pgfqpoint{9.417597in}{1.281954in}}%
\pgfpathquadraticcurveto{\pgfqpoint{9.417597in}{1.292582in}}{\pgfqpoint{9.424712in}{1.297985in}}%
\pgfpathquadraticcurveto{\pgfqpoint{9.431845in}{1.303389in}}{\pgfqpoint{9.445966in}{1.303389in}}%
\pgfpathlineto{\pgfqpoint{9.460502in}{1.303389in}}%
\pgfpathlineto{\pgfqpoint{9.460502in}{1.304416in}}%
\pgfpathquadraticcurveto{\pgfqpoint{9.460502in}{1.311566in}}{\pgfqpoint{9.455801in}{1.315475in}}%
\pgfpathquadraticcurveto{\pgfqpoint{9.451100in}{1.319384in}}{\pgfqpoint{9.442598in}{1.319384in}}%
\pgfpathquadraticcurveto{\pgfqpoint{9.437194in}{1.319384in}}{\pgfqpoint{9.432061in}{1.318087in}}%
\pgfpathquadraticcurveto{\pgfqpoint{9.426945in}{1.316790in}}{\pgfqpoint{9.422226in}{1.314196in}}%
\pgfpathlineto{\pgfqpoint{9.422226in}{1.323779in}}%
\pgfpathquadraticcurveto{\pgfqpoint{9.427900in}{1.325976in}}{\pgfqpoint{9.433250in}{1.327057in}}%
\pgfpathquadraticcurveto{\pgfqpoint{9.438599in}{1.328156in}}{\pgfqpoint{9.443661in}{1.328156in}}%
\pgfpathquadraticcurveto{\pgfqpoint{9.457350in}{1.328156in}}{\pgfqpoint{9.464104in}{1.321059in}}%
\pgfpathquadraticcurveto{\pgfqpoint{9.470859in}{1.313980in}}{\pgfqpoint{9.470859in}{1.299570in}}%
\pgfpathlineto{\pgfqpoint{9.470859in}{1.299570in}}%
\pgfpathclose%
\pgfpathmoveto{\pgfqpoint{9.502434in}{1.344547in}}%
\pgfpathlineto{\pgfqpoint{9.502434in}{1.326643in}}%
\pgfpathlineto{\pgfqpoint{9.523761in}{1.326643in}}%
\pgfpathlineto{\pgfqpoint{9.523761in}{1.318591in}}%
\pgfpathlineto{\pgfqpoint{9.502434in}{1.318591in}}%
\pgfpathlineto{\pgfqpoint{9.502434in}{1.284368in}}%
\pgfpathquadraticcurveto{\pgfqpoint{9.502434in}{1.276659in}}{\pgfqpoint{9.504542in}{1.274461in}}%
\pgfpathquadraticcurveto{\pgfqpoint{9.506649in}{1.272264in}}{\pgfqpoint{9.513133in}{1.272264in}}%
\pgfpathlineto{\pgfqpoint{9.523761in}{1.272264in}}%
\pgfpathlineto{\pgfqpoint{9.523761in}{1.263600in}}%
\pgfpathlineto{\pgfqpoint{9.513133in}{1.263600in}}%
\pgfpathquadraticcurveto{\pgfqpoint{9.501137in}{1.263600in}}{\pgfqpoint{9.496580in}{1.268067in}}%
\pgfpathquadraticcurveto{\pgfqpoint{9.492023in}{1.272552in}}{\pgfqpoint{9.492023in}{1.284368in}}%
\pgfpathlineto{\pgfqpoint{9.492023in}{1.318591in}}%
\pgfpathlineto{\pgfqpoint{9.484422in}{1.318591in}}%
\pgfpathlineto{\pgfqpoint{9.484422in}{1.326643in}}%
\pgfpathlineto{\pgfqpoint{9.492023in}{1.326643in}}%
\pgfpathlineto{\pgfqpoint{9.492023in}{1.344547in}}%
\pgfpathlineto{\pgfqpoint{9.502434in}{1.344547in}}%
\pgfpathlineto{\pgfqpoint{9.502434in}{1.344547in}}%
\pgfpathclose%
\pgfpathmoveto{\pgfqpoint{9.586461in}{1.314538in}}%
\pgfpathquadraticcurveto{\pgfqpoint{9.590352in}{1.321527in}}{\pgfqpoint{9.595755in}{1.324841in}}%
\pgfpathquadraticcurveto{\pgfqpoint{9.601159in}{1.328156in}}{\pgfqpoint{9.608472in}{1.328156in}}%
\pgfpathquadraticcurveto{\pgfqpoint{9.618324in}{1.328156in}}{\pgfqpoint{9.623674in}{1.321257in}}%
\pgfpathquadraticcurveto{\pgfqpoint{9.629024in}{1.314376in}}{\pgfqpoint{9.629024in}{1.301660in}}%
\pgfpathlineto{\pgfqpoint{9.629024in}{1.263600in}}%
\pgfpathlineto{\pgfqpoint{9.618613in}{1.263600in}}%
\pgfpathlineto{\pgfqpoint{9.618613in}{1.301317in}}%
\pgfpathquadraticcurveto{\pgfqpoint{9.618613in}{1.310378in}}{\pgfqpoint{9.615388in}{1.314755in}}%
\pgfpathquadraticcurveto{\pgfqpoint{9.612182in}{1.319149in}}{\pgfqpoint{9.605608in}{1.319149in}}%
\pgfpathquadraticcurveto{\pgfqpoint{9.597556in}{1.319149in}}{\pgfqpoint{9.592873in}{1.313800in}}%
\pgfpathquadraticcurveto{\pgfqpoint{9.588208in}{1.308468in}}{\pgfqpoint{9.588208in}{1.299228in}}%
\pgfpathlineto{\pgfqpoint{9.588208in}{1.263600in}}%
\pgfpathlineto{\pgfqpoint{9.577797in}{1.263600in}}%
\pgfpathlineto{\pgfqpoint{9.577797in}{1.301317in}}%
\pgfpathquadraticcurveto{\pgfqpoint{9.577797in}{1.310432in}}{\pgfqpoint{9.574591in}{1.314791in}}%
\pgfpathquadraticcurveto{\pgfqpoint{9.571385in}{1.319149in}}{\pgfqpoint{9.564684in}{1.319149in}}%
\pgfpathquadraticcurveto{\pgfqpoint{9.556741in}{1.319149in}}{\pgfqpoint{9.552058in}{1.313782in}}%
\pgfpathquadraticcurveto{\pgfqpoint{9.547393in}{1.308414in}}{\pgfqpoint{9.547393in}{1.299228in}}%
\pgfpathlineto{\pgfqpoint{9.547393in}{1.263600in}}%
\pgfpathlineto{\pgfqpoint{9.536982in}{1.263600in}}%
\pgfpathlineto{\pgfqpoint{9.536982in}{1.326643in}}%
\pgfpathlineto{\pgfqpoint{9.547393in}{1.326643in}}%
\pgfpathlineto{\pgfqpoint{9.547393in}{1.316844in}}%
\pgfpathquadraticcurveto{\pgfqpoint{9.550941in}{1.322644in}}{\pgfqpoint{9.555894in}{1.325400in}}%
\pgfpathquadraticcurveto{\pgfqpoint{9.560848in}{1.328156in}}{\pgfqpoint{9.567656in}{1.328156in}}%
\pgfpathquadraticcurveto{\pgfqpoint{9.574537in}{1.328156in}}{\pgfqpoint{9.579346in}{1.324661in}}%
\pgfpathquadraticcurveto{\pgfqpoint{9.584155in}{1.321185in}}{\pgfqpoint{9.586461in}{1.314538in}}%
\pgfpathlineto{\pgfqpoint{9.586461in}{1.314538in}}%
\pgfpathclose%
\pgfpathmoveto{\pgfqpoint{9.703594in}{1.297715in}}%
\pgfpathlineto{\pgfqpoint{9.703594in}{1.292654in}}%
\pgfpathlineto{\pgfqpoint{9.655970in}{1.292654in}}%
\pgfpathquadraticcurveto{\pgfqpoint{9.656654in}{1.281954in}}{\pgfqpoint{9.662418in}{1.276353in}}%
\pgfpathquadraticcurveto{\pgfqpoint{9.668182in}{1.270751in}}{\pgfqpoint{9.678485in}{1.270751in}}%
\pgfpathquadraticcurveto{\pgfqpoint{9.684447in}{1.270751in}}{\pgfqpoint{9.690049in}{1.272210in}}%
\pgfpathquadraticcurveto{\pgfqpoint{9.695651in}{1.273669in}}{\pgfqpoint{9.701180in}{1.276605in}}%
\pgfpathlineto{\pgfqpoint{9.701180in}{1.266806in}}%
\pgfpathquadraticcurveto{\pgfqpoint{9.695597in}{1.264447in}}{\pgfqpoint{9.689743in}{1.263204in}}%
\pgfpathquadraticcurveto{\pgfqpoint{9.683889in}{1.261961in}}{\pgfqpoint{9.677873in}{1.261961in}}%
\pgfpathquadraticcurveto{\pgfqpoint{9.662778in}{1.261961in}}{\pgfqpoint{9.653970in}{1.270733in}}%
\pgfpathquadraticcurveto{\pgfqpoint{9.645163in}{1.279523in}}{\pgfqpoint{9.645163in}{1.294509in}}%
\pgfpathquadraticcurveto{\pgfqpoint{9.645163in}{1.309981in}}{\pgfqpoint{9.653520in}{1.319059in}}%
\pgfpathquadraticcurveto{\pgfqpoint{9.661878in}{1.328156in}}{\pgfqpoint{9.676071in}{1.328156in}}%
\pgfpathquadraticcurveto{\pgfqpoint{9.688788in}{1.328156in}}{\pgfqpoint{9.696191in}{1.319960in}}%
\pgfpathquadraticcurveto{\pgfqpoint{9.703594in}{1.311783in}}{\pgfqpoint{9.703594in}{1.297715in}}%
\pgfpathlineto{\pgfqpoint{9.703594in}{1.297715in}}%
\pgfpathclose%
\pgfpathmoveto{\pgfqpoint{9.693237in}{1.300759in}}%
\pgfpathquadraticcurveto{\pgfqpoint{9.693129in}{1.309243in}}{\pgfqpoint{9.688482in}{1.314304in}}%
\pgfpathquadraticcurveto{\pgfqpoint{9.683835in}{1.319384in}}{\pgfqpoint{9.676179in}{1.319384in}}%
\pgfpathquadraticcurveto{\pgfqpoint{9.667516in}{1.319384in}}{\pgfqpoint{9.662310in}{1.314484in}}%
\pgfpathquadraticcurveto{\pgfqpoint{9.657105in}{1.309585in}}{\pgfqpoint{9.656312in}{1.300687in}}%
\pgfpathlineto{\pgfqpoint{9.693237in}{1.300759in}}%
\pgfpathlineto{\pgfqpoint{9.693237in}{1.300759in}}%
\pgfpathclose%
\pgfpathmoveto{\pgfqpoint{9.773013in}{1.301660in}}%
\pgfpathlineto{\pgfqpoint{9.773013in}{1.263600in}}%
\pgfpathlineto{\pgfqpoint{9.762656in}{1.263600in}}%
\pgfpathlineto{\pgfqpoint{9.762656in}{1.301317in}}%
\pgfpathquadraticcurveto{\pgfqpoint{9.762656in}{1.310269in}}{\pgfqpoint{9.759161in}{1.314700in}}%
\pgfpathquadraticcurveto{\pgfqpoint{9.755667in}{1.319149in}}{\pgfqpoint{9.748696in}{1.319149in}}%
\pgfpathquadraticcurveto{\pgfqpoint{9.740303in}{1.319149in}}{\pgfqpoint{9.735457in}{1.313800in}}%
\pgfpathquadraticcurveto{\pgfqpoint{9.730612in}{1.308468in}}{\pgfqpoint{9.730612in}{1.299228in}}%
\pgfpathlineto{\pgfqpoint{9.730612in}{1.263600in}}%
\pgfpathlineto{\pgfqpoint{9.720201in}{1.263600in}}%
\pgfpathlineto{\pgfqpoint{9.720201in}{1.326643in}}%
\pgfpathlineto{\pgfqpoint{9.730612in}{1.326643in}}%
\pgfpathlineto{\pgfqpoint{9.730612in}{1.316844in}}%
\pgfpathquadraticcurveto{\pgfqpoint{9.734341in}{1.322536in}}{\pgfqpoint{9.739366in}{1.325346in}}%
\pgfpathquadraticcurveto{\pgfqpoint{9.744410in}{1.328156in}}{\pgfqpoint{9.751002in}{1.328156in}}%
\pgfpathquadraticcurveto{\pgfqpoint{9.761863in}{1.328156in}}{\pgfqpoint{9.767429in}{1.321437in}}%
\pgfpathquadraticcurveto{\pgfqpoint{9.773013in}{1.314718in}}{\pgfqpoint{9.773013in}{1.301660in}}%
\pgfpathlineto{\pgfqpoint{9.773013in}{1.301660in}}%
\pgfpathclose%
\pgfpathmoveto{\pgfqpoint{9.803904in}{1.344547in}}%
\pgfpathlineto{\pgfqpoint{9.803904in}{1.326643in}}%
\pgfpathlineto{\pgfqpoint{9.825230in}{1.326643in}}%
\pgfpathlineto{\pgfqpoint{9.825230in}{1.318591in}}%
\pgfpathlineto{\pgfqpoint{9.803904in}{1.318591in}}%
\pgfpathlineto{\pgfqpoint{9.803904in}{1.284368in}}%
\pgfpathquadraticcurveto{\pgfqpoint{9.803904in}{1.276659in}}{\pgfqpoint{9.806011in}{1.274461in}}%
\pgfpathquadraticcurveto{\pgfqpoint{9.808118in}{1.272264in}}{\pgfqpoint{9.814603in}{1.272264in}}%
\pgfpathlineto{\pgfqpoint{9.825230in}{1.272264in}}%
\pgfpathlineto{\pgfqpoint{9.825230in}{1.263600in}}%
\pgfpathlineto{\pgfqpoint{9.814603in}{1.263600in}}%
\pgfpathquadraticcurveto{\pgfqpoint{9.802607in}{1.263600in}}{\pgfqpoint{9.798050in}{1.268067in}}%
\pgfpathquadraticcurveto{\pgfqpoint{9.793493in}{1.272552in}}{\pgfqpoint{9.793493in}{1.284368in}}%
\pgfpathlineto{\pgfqpoint{9.793493in}{1.318591in}}%
\pgfpathlineto{\pgfqpoint{9.785891in}{1.318591in}}%
\pgfpathlineto{\pgfqpoint{9.785891in}{1.326643in}}%
\pgfpathlineto{\pgfqpoint{9.793493in}{1.326643in}}%
\pgfpathlineto{\pgfqpoint{9.793493in}{1.344547in}}%
\pgfpathlineto{\pgfqpoint{9.803904in}{1.344547in}}%
\pgfpathlineto{\pgfqpoint{9.803904in}{1.344547in}}%
\pgfpathclose%
\pgfusepath{fill}%
\end{pgfscope}%
\begin{pgfscope}%
\pgfsetbuttcap%
\pgfsetroundjoin%
\definecolor{currentfill}{rgb}{0.309804,0.372549,0.419608}%
\pgfsetfillcolor{currentfill}%
\pgfsetlinewidth{0.000000pt}%
\definecolor{currentstroke}{rgb}{0.309804,0.372549,0.419608}%
\pgfsetstrokecolor{currentstroke}%
\pgfsetdash{}{0pt}%
\pgfpathmoveto{\pgfqpoint{9.977680in}{1.297715in}}%
\pgfpathlineto{\pgfqpoint{9.977680in}{1.292654in}}%
\pgfpathlineto{\pgfqpoint{9.930056in}{1.292654in}}%
\pgfpathquadraticcurveto{\pgfqpoint{9.930741in}{1.281954in}}{\pgfqpoint{9.936505in}{1.276353in}}%
\pgfpathquadraticcurveto{\pgfqpoint{9.942268in}{1.270751in}}{\pgfqpoint{9.952571in}{1.270751in}}%
\pgfpathquadraticcurveto{\pgfqpoint{9.958533in}{1.270751in}}{\pgfqpoint{9.964135in}{1.272210in}}%
\pgfpathquadraticcurveto{\pgfqpoint{9.969737in}{1.273669in}}{\pgfqpoint{9.975267in}{1.276605in}}%
\pgfpathlineto{\pgfqpoint{9.975267in}{1.266806in}}%
\pgfpathquadraticcurveto{\pgfqpoint{9.969683in}{1.264447in}}{\pgfqpoint{9.963829in}{1.263204in}}%
\pgfpathquadraticcurveto{\pgfqpoint{9.957975in}{1.261961in}}{\pgfqpoint{9.951959in}{1.261961in}}%
\pgfpathquadraticcurveto{\pgfqpoint{9.936865in}{1.261961in}}{\pgfqpoint{9.928057in}{1.270733in}}%
\pgfpathquadraticcurveto{\pgfqpoint{9.919249in}{1.279523in}}{\pgfqpoint{9.919249in}{1.294509in}}%
\pgfpathquadraticcurveto{\pgfqpoint{9.919249in}{1.309981in}}{\pgfqpoint{9.927607in}{1.319059in}}%
\pgfpathquadraticcurveto{\pgfqpoint{9.935964in}{1.328156in}}{\pgfqpoint{9.950158in}{1.328156in}}%
\pgfpathquadraticcurveto{\pgfqpoint{9.962874in}{1.328156in}}{\pgfqpoint{9.970277in}{1.319960in}}%
\pgfpathquadraticcurveto{\pgfqpoint{9.977680in}{1.311783in}}{\pgfqpoint{9.977680in}{1.297715in}}%
\pgfpathlineto{\pgfqpoint{9.977680in}{1.297715in}}%
\pgfpathclose%
\pgfpathmoveto{\pgfqpoint{9.967323in}{1.300759in}}%
\pgfpathquadraticcurveto{\pgfqpoint{9.967215in}{1.309243in}}{\pgfqpoint{9.962568in}{1.314304in}}%
\pgfpathquadraticcurveto{\pgfqpoint{9.957921in}{1.319384in}}{\pgfqpoint{9.950266in}{1.319384in}}%
\pgfpathquadraticcurveto{\pgfqpoint{9.941602in}{1.319384in}}{\pgfqpoint{9.936396in}{1.314484in}}%
\pgfpathquadraticcurveto{\pgfqpoint{9.931191in}{1.309585in}}{\pgfqpoint{9.930398in}{1.300687in}}%
\pgfpathlineto{\pgfqpoint{9.967323in}{1.300759in}}%
\pgfpathlineto{\pgfqpoint{9.967323in}{1.300759in}}%
\pgfpathclose%
\pgfpathmoveto{\pgfqpoint{10.065436in}{1.351193in}}%
\pgfpathlineto{\pgfqpoint{10.065436in}{1.342565in}}%
\pgfpathlineto{\pgfqpoint{10.055529in}{1.342565in}}%
\pgfpathquadraticcurveto{\pgfqpoint{10.049963in}{1.342565in}}{\pgfqpoint{10.047784in}{1.340314in}}%
\pgfpathquadraticcurveto{\pgfqpoint{10.045622in}{1.338062in}}{\pgfqpoint{10.045622in}{1.332208in}}%
\pgfpathlineto{\pgfqpoint{10.045622in}{1.326643in}}%
\pgfpathlineto{\pgfqpoint{10.062680in}{1.326643in}}%
\pgfpathlineto{\pgfqpoint{10.062680in}{1.318591in}}%
\pgfpathlineto{\pgfqpoint{10.045622in}{1.318591in}}%
\pgfpathlineto{\pgfqpoint{10.045622in}{1.263600in}}%
\pgfpathlineto{\pgfqpoint{10.035211in}{1.263600in}}%
\pgfpathlineto{\pgfqpoint{10.035211in}{1.318591in}}%
\pgfpathlineto{\pgfqpoint{10.006788in}{1.318591in}}%
\pgfpathlineto{\pgfqpoint{10.006788in}{1.263600in}}%
\pgfpathlineto{\pgfqpoint{9.996377in}{1.263600in}}%
\pgfpathlineto{\pgfqpoint{9.996377in}{1.318591in}}%
\pgfpathlineto{\pgfqpoint{9.986470in}{1.318591in}}%
\pgfpathlineto{\pgfqpoint{9.986470in}{1.326643in}}%
\pgfpathlineto{\pgfqpoint{9.996377in}{1.326643in}}%
\pgfpathlineto{\pgfqpoint{9.996377in}{1.331038in}}%
\pgfpathquadraticcurveto{\pgfqpoint{9.996377in}{1.341557in}}{\pgfqpoint{10.001276in}{1.346366in}}%
\pgfpathquadraticcurveto{\pgfqpoint{10.006176in}{1.351193in}}{\pgfqpoint{10.016803in}{1.351193in}}%
\pgfpathlineto{\pgfqpoint{10.026601in}{1.351193in}}%
\pgfpathlineto{\pgfqpoint{10.026601in}{1.342565in}}%
\pgfpathlineto{\pgfqpoint{10.016695in}{1.342565in}}%
\pgfpathquadraticcurveto{\pgfqpoint{10.011129in}{1.342565in}}{\pgfqpoint{10.008931in}{1.340314in}}%
\pgfpathquadraticcurveto{\pgfqpoint{10.006788in}{1.338062in}}{\pgfqpoint{10.006788in}{1.332208in}}%
\pgfpathlineto{\pgfqpoint{10.006788in}{1.326643in}}%
\pgfpathlineto{\pgfqpoint{10.035211in}{1.326643in}}%
\pgfpathlineto{\pgfqpoint{10.035211in}{1.331038in}}%
\pgfpathquadraticcurveto{\pgfqpoint{10.035211in}{1.341557in}}{\pgfqpoint{10.040110in}{1.346366in}}%
\pgfpathquadraticcurveto{\pgfqpoint{10.045010in}{1.351193in}}{\pgfqpoint{10.055655in}{1.351193in}}%
\pgfpathlineto{\pgfqpoint{10.065436in}{1.351193in}}%
\pgfpathlineto{\pgfqpoint{10.065436in}{1.351193in}}%
\pgfpathclose%
\pgfpathmoveto{\pgfqpoint{10.128028in}{1.297715in}}%
\pgfpathlineto{\pgfqpoint{10.128028in}{1.292654in}}%
\pgfpathlineto{\pgfqpoint{10.080404in}{1.292654in}}%
\pgfpathquadraticcurveto{\pgfqpoint{10.081088in}{1.281954in}}{\pgfqpoint{10.086852in}{1.276353in}}%
\pgfpathquadraticcurveto{\pgfqpoint{10.092616in}{1.270751in}}{\pgfqpoint{10.102919in}{1.270751in}}%
\pgfpathquadraticcurveto{\pgfqpoint{10.108881in}{1.270751in}}{\pgfqpoint{10.114483in}{1.272210in}}%
\pgfpathquadraticcurveto{\pgfqpoint{10.120084in}{1.273669in}}{\pgfqpoint{10.125614in}{1.276605in}}%
\pgfpathlineto{\pgfqpoint{10.125614in}{1.266806in}}%
\pgfpathquadraticcurveto{\pgfqpoint{10.120030in}{1.264447in}}{\pgfqpoint{10.114176in}{1.263204in}}%
\pgfpathquadraticcurveto{\pgfqpoint{10.108322in}{1.261961in}}{\pgfqpoint{10.102306in}{1.261961in}}%
\pgfpathquadraticcurveto{\pgfqpoint{10.087212in}{1.261961in}}{\pgfqpoint{10.078404in}{1.270733in}}%
\pgfpathquadraticcurveto{\pgfqpoint{10.069596in}{1.279523in}}{\pgfqpoint{10.069596in}{1.294509in}}%
\pgfpathquadraticcurveto{\pgfqpoint{10.069596in}{1.309981in}}{\pgfqpoint{10.077954in}{1.319059in}}%
\pgfpathquadraticcurveto{\pgfqpoint{10.086312in}{1.328156in}}{\pgfqpoint{10.100505in}{1.328156in}}%
\pgfpathquadraticcurveto{\pgfqpoint{10.113222in}{1.328156in}}{\pgfqpoint{10.120625in}{1.319960in}}%
\pgfpathquadraticcurveto{\pgfqpoint{10.128028in}{1.311783in}}{\pgfqpoint{10.128028in}{1.297715in}}%
\pgfpathlineto{\pgfqpoint{10.128028in}{1.297715in}}%
\pgfpathclose%
\pgfpathmoveto{\pgfqpoint{10.117671in}{1.300759in}}%
\pgfpathquadraticcurveto{\pgfqpoint{10.117563in}{1.309243in}}{\pgfqpoint{10.112916in}{1.314304in}}%
\pgfpathquadraticcurveto{\pgfqpoint{10.108268in}{1.319384in}}{\pgfqpoint{10.100613in}{1.319384in}}%
\pgfpathquadraticcurveto{\pgfqpoint{10.091949in}{1.319384in}}{\pgfqpoint{10.086744in}{1.314484in}}%
\pgfpathquadraticcurveto{\pgfqpoint{10.081538in}{1.309585in}}{\pgfqpoint{10.080746in}{1.300687in}}%
\pgfpathlineto{\pgfqpoint{10.117671in}{1.300759in}}%
\pgfpathlineto{\pgfqpoint{10.117671in}{1.300759in}}%
\pgfpathclose%
\pgfpathmoveto{\pgfqpoint{10.190404in}{1.324229in}}%
\pgfpathlineto{\pgfqpoint{10.190404in}{1.314538in}}%
\pgfpathquadraticcurveto{\pgfqpoint{10.186009in}{1.316970in}}{\pgfqpoint{10.181596in}{1.318177in}}%
\pgfpathquadraticcurveto{\pgfqpoint{10.177183in}{1.319384in}}{\pgfqpoint{10.172680in}{1.319384in}}%
\pgfpathquadraticcurveto{\pgfqpoint{10.162593in}{1.319384in}}{\pgfqpoint{10.157009in}{1.312989in}}%
\pgfpathquadraticcurveto{\pgfqpoint{10.151444in}{1.306613in}}{\pgfqpoint{10.151444in}{1.295067in}}%
\pgfpathquadraticcurveto{\pgfqpoint{10.151444in}{1.283521in}}{\pgfqpoint{10.157009in}{1.277127in}}%
\pgfpathquadraticcurveto{\pgfqpoint{10.162593in}{1.270751in}}{\pgfqpoint{10.172680in}{1.270751in}}%
\pgfpathquadraticcurveto{\pgfqpoint{10.177183in}{1.270751in}}{\pgfqpoint{10.181596in}{1.271958in}}%
\pgfpathquadraticcurveto{\pgfqpoint{10.186009in}{1.273164in}}{\pgfqpoint{10.190404in}{1.275596in}}%
\pgfpathlineto{\pgfqpoint{10.190404in}{1.266014in}}%
\pgfpathquadraticcurveto{\pgfqpoint{10.186063in}{1.263996in}}{\pgfqpoint{10.181416in}{1.262988in}}%
\pgfpathquadraticcurveto{\pgfqpoint{10.176787in}{1.261961in}}{\pgfqpoint{10.171545in}{1.261961in}}%
\pgfpathquadraticcurveto{\pgfqpoint{10.157298in}{1.261961in}}{\pgfqpoint{10.148904in}{1.270913in}}%
\pgfpathquadraticcurveto{\pgfqpoint{10.140528in}{1.279865in}}{\pgfqpoint{10.140528in}{1.295067in}}%
\pgfpathquadraticcurveto{\pgfqpoint{10.140528in}{1.310486in}}{\pgfqpoint{10.148994in}{1.319312in}}%
\pgfpathquadraticcurveto{\pgfqpoint{10.157478in}{1.328156in}}{\pgfqpoint{10.172230in}{1.328156in}}%
\pgfpathquadraticcurveto{\pgfqpoint{10.177003in}{1.328156in}}{\pgfqpoint{10.181560in}{1.327165in}}%
\pgfpathquadraticcurveto{\pgfqpoint{10.186117in}{1.326192in}}{\pgfqpoint{10.190404in}{1.324229in}}%
\pgfpathlineto{\pgfqpoint{10.190404in}{1.324229in}}%
\pgfpathclose%
\pgfpathmoveto{\pgfqpoint{10.218665in}{1.344547in}}%
\pgfpathlineto{\pgfqpoint{10.218665in}{1.326643in}}%
\pgfpathlineto{\pgfqpoint{10.239991in}{1.326643in}}%
\pgfpathlineto{\pgfqpoint{10.239991in}{1.318591in}}%
\pgfpathlineto{\pgfqpoint{10.218665in}{1.318591in}}%
\pgfpathlineto{\pgfqpoint{10.218665in}{1.284368in}}%
\pgfpathquadraticcurveto{\pgfqpoint{10.218665in}{1.276659in}}{\pgfqpoint{10.220772in}{1.274461in}}%
\pgfpathquadraticcurveto{\pgfqpoint{10.222880in}{1.272264in}}{\pgfqpoint{10.229364in}{1.272264in}}%
\pgfpathlineto{\pgfqpoint{10.239991in}{1.272264in}}%
\pgfpathlineto{\pgfqpoint{10.239991in}{1.263600in}}%
\pgfpathlineto{\pgfqpoint{10.229364in}{1.263600in}}%
\pgfpathquadraticcurveto{\pgfqpoint{10.217368in}{1.263600in}}{\pgfqpoint{10.212811in}{1.268067in}}%
\pgfpathquadraticcurveto{\pgfqpoint{10.208254in}{1.272552in}}{\pgfqpoint{10.208254in}{1.284368in}}%
\pgfpathlineto{\pgfqpoint{10.208254in}{1.318591in}}%
\pgfpathlineto{\pgfqpoint{10.200653in}{1.318591in}}%
\pgfpathlineto{\pgfqpoint{10.200653in}{1.326643in}}%
\pgfpathlineto{\pgfqpoint{10.208254in}{1.326643in}}%
\pgfpathlineto{\pgfqpoint{10.208254in}{1.344547in}}%
\pgfpathlineto{\pgfqpoint{10.218665in}{1.344547in}}%
\pgfpathlineto{\pgfqpoint{10.218665in}{1.344547in}}%
\pgfpathclose%
\pgfpathmoveto{\pgfqpoint{10.293794in}{1.324787in}}%
\pgfpathlineto{\pgfqpoint{10.293794in}{1.314989in}}%
\pgfpathquadraticcurveto{\pgfqpoint{10.289417in}{1.317240in}}{\pgfqpoint{10.284679in}{1.318357in}}%
\pgfpathquadraticcurveto{\pgfqpoint{10.279960in}{1.319492in}}{\pgfqpoint{10.274881in}{1.319492in}}%
\pgfpathquadraticcurveto{\pgfqpoint{10.267172in}{1.319492in}}{\pgfqpoint{10.263317in}{1.317132in}}%
\pgfpathquadraticcurveto{\pgfqpoint{10.259462in}{1.314773in}}{\pgfqpoint{10.259462in}{1.310035in}}%
\pgfpathquadraticcurveto{\pgfqpoint{10.259462in}{1.306433in}}{\pgfqpoint{10.262218in}{1.304380in}}%
\pgfpathquadraticcurveto{\pgfqpoint{10.264974in}{1.302326in}}{\pgfqpoint{10.273314in}{1.300471in}}%
\pgfpathlineto{\pgfqpoint{10.276862in}{1.299678in}}%
\pgfpathquadraticcurveto{\pgfqpoint{10.287886in}{1.297319in}}{\pgfqpoint{10.292533in}{1.293014in}}%
\pgfpathquadraticcurveto{\pgfqpoint{10.297180in}{1.288709in}}{\pgfqpoint{10.297180in}{1.281000in}}%
\pgfpathquadraticcurveto{\pgfqpoint{10.297180in}{1.272210in}}{\pgfqpoint{10.290227in}{1.267076in}}%
\pgfpathquadraticcurveto{\pgfqpoint{10.283275in}{1.261961in}}{\pgfqpoint{10.271116in}{1.261961in}}%
\pgfpathquadraticcurveto{\pgfqpoint{10.266055in}{1.261961in}}{\pgfqpoint{10.260561in}{1.262952in}}%
\pgfpathquadraticcurveto{\pgfqpoint{10.255067in}{1.263942in}}{\pgfqpoint{10.248997in}{1.265906in}}%
\pgfpathlineto{\pgfqpoint{10.248997in}{1.276605in}}%
\pgfpathquadraticcurveto{\pgfqpoint{10.254743in}{1.273615in}}{\pgfqpoint{10.260309in}{1.272120in}}%
\pgfpathquadraticcurveto{\pgfqpoint{10.265875in}{1.270643in}}{\pgfqpoint{10.271350in}{1.270643in}}%
\pgfpathquadraticcurveto{\pgfqpoint{10.278663in}{1.270643in}}{\pgfqpoint{10.282590in}{1.273146in}}%
\pgfpathquadraticcurveto{\pgfqpoint{10.286535in}{1.275650in}}{\pgfqpoint{10.286535in}{1.280207in}}%
\pgfpathquadraticcurveto{\pgfqpoint{10.286535in}{1.284422in}}{\pgfqpoint{10.283689in}{1.286674in}}%
\pgfpathquadraticcurveto{\pgfqpoint{10.280861in}{1.288925in}}{\pgfqpoint{10.271224in}{1.291014in}}%
\pgfpathlineto{\pgfqpoint{10.267622in}{1.291861in}}%
\pgfpathquadraticcurveto{\pgfqpoint{10.258003in}{1.293878in}}{\pgfqpoint{10.253717in}{1.298075in}}%
\pgfpathquadraticcurveto{\pgfqpoint{10.249448in}{1.302272in}}{\pgfqpoint{10.249448in}{1.309585in}}%
\pgfpathquadraticcurveto{\pgfqpoint{10.249448in}{1.318483in}}{\pgfqpoint{10.255752in}{1.323310in}}%
\pgfpathquadraticcurveto{\pgfqpoint{10.262056in}{1.328156in}}{\pgfqpoint{10.273656in}{1.328156in}}%
\pgfpathquadraticcurveto{\pgfqpoint{10.279384in}{1.328156in}}{\pgfqpoint{10.284445in}{1.327309in}}%
\pgfpathquadraticcurveto{\pgfqpoint{10.289525in}{1.326480in}}{\pgfqpoint{10.293794in}{1.324787in}}%
\pgfpathlineto{\pgfqpoint{10.293794in}{1.324787in}}%
\pgfpathclose%
\pgfpathmoveto{\pgfqpoint{10.315120in}{1.277902in}}%
\pgfpathlineto{\pgfqpoint{10.327008in}{1.277902in}}%
\pgfpathlineto{\pgfqpoint{10.327008in}{1.263600in}}%
\pgfpathlineto{\pgfqpoint{10.315120in}{1.263600in}}%
\pgfpathlineto{\pgfqpoint{10.315120in}{1.277902in}}%
\pgfpathlineto{\pgfqpoint{10.315120in}{1.277902in}}%
\pgfpathclose%
\pgfusepath{fill}%
\end{pgfscope}%
\begin{pgfscope}%
\pgfsetbuttcap%
\pgfsetroundjoin%
\definecolor{currentfill}{rgb}{0.309804,0.372549,0.419608}%
\pgfsetfillcolor{currentfill}%
\pgfsetlinewidth{0.000000pt}%
\definecolor{currentstroke}{rgb}{0.309804,0.372549,0.419608}%
\pgfsetstrokecolor{currentstroke}%
\pgfsetdash{}{0pt}%
\pgfpathmoveto{\pgfqpoint{10.470267in}{1.344889in}}%
\pgfpathlineto{\pgfqpoint{10.470267in}{1.333793in}}%
\pgfpathquadraticcurveto{\pgfqpoint{10.463801in}{1.336891in}}{\pgfqpoint{10.458055in}{1.338404in}}%
\pgfpathquadraticcurveto{\pgfqpoint{10.452309in}{1.339936in}}{\pgfqpoint{10.446959in}{1.339936in}}%
\pgfpathquadraticcurveto{\pgfqpoint{10.437683in}{1.339936in}}{\pgfqpoint{10.432640in}{1.336333in}}%
\pgfpathquadraticcurveto{\pgfqpoint{10.427596in}{1.332731in}}{\pgfqpoint{10.427596in}{1.326084in}}%
\pgfpathquadraticcurveto{\pgfqpoint{10.427596in}{1.320500in}}{\pgfqpoint{10.430947in}{1.317654in}}%
\pgfpathquadraticcurveto{\pgfqpoint{10.434297in}{1.314827in}}{\pgfqpoint{10.443645in}{1.313079in}}%
\pgfpathlineto{\pgfqpoint{10.450508in}{1.311674in}}%
\pgfpathquadraticcurveto{\pgfqpoint{10.463224in}{1.309243in}}{\pgfqpoint{10.469276in}{1.303137in}}%
\pgfpathquadraticcurveto{\pgfqpoint{10.475329in}{1.297031in}}{\pgfqpoint{10.475329in}{1.286800in}}%
\pgfpathquadraticcurveto{\pgfqpoint{10.475329in}{1.274569in}}{\pgfqpoint{10.467133in}{1.268265in}}%
\pgfpathquadraticcurveto{\pgfqpoint{10.458956in}{1.261961in}}{\pgfqpoint{10.443141in}{1.261961in}}%
\pgfpathquadraticcurveto{\pgfqpoint{10.437179in}{1.261961in}}{\pgfqpoint{10.430442in}{1.263312in}}%
\pgfpathquadraticcurveto{\pgfqpoint{10.423724in}{1.264663in}}{\pgfqpoint{10.416519in}{1.267311in}}%
\pgfpathlineto{\pgfqpoint{10.416519in}{1.279018in}}%
\pgfpathquadraticcurveto{\pgfqpoint{10.423436in}{1.275146in}}{\pgfqpoint{10.430082in}{1.273164in}}%
\pgfpathquadraticcurveto{\pgfqpoint{10.436729in}{1.271201in}}{\pgfqpoint{10.443141in}{1.271201in}}%
\pgfpathquadraticcurveto{\pgfqpoint{10.452867in}{1.271201in}}{\pgfqpoint{10.458163in}{1.275020in}}%
\pgfpathquadraticcurveto{\pgfqpoint{10.463459in}{1.278856in}}{\pgfqpoint{10.463459in}{1.285953in}}%
\pgfpathquadraticcurveto{\pgfqpoint{10.463459in}{1.292131in}}{\pgfqpoint{10.459658in}{1.295626in}}%
\pgfpathquadraticcurveto{\pgfqpoint{10.455857in}{1.299120in}}{\pgfqpoint{10.447194in}{1.300867in}}%
\pgfpathlineto{\pgfqpoint{10.440259in}{1.302218in}}%
\pgfpathquadraticcurveto{\pgfqpoint{10.427542in}{1.304740in}}{\pgfqpoint{10.421850in}{1.310143in}}%
\pgfpathquadraticcurveto{\pgfqpoint{10.416177in}{1.315547in}}{\pgfqpoint{10.416177in}{1.325184in}}%
\pgfpathquadraticcurveto{\pgfqpoint{10.416177in}{1.336333in}}{\pgfqpoint{10.424030in}{1.342745in}}%
\pgfpathquadraticcurveto{\pgfqpoint{10.431883in}{1.349158in}}{\pgfqpoint{10.445663in}{1.349158in}}%
\pgfpathquadraticcurveto{\pgfqpoint{10.451589in}{1.349158in}}{\pgfqpoint{10.457713in}{1.348077in}}%
\pgfpathquadraticcurveto{\pgfqpoint{10.463855in}{1.347014in}}{\pgfqpoint{10.470267in}{1.344889in}}%
\pgfpathlineto{\pgfqpoint{10.470267in}{1.344889in}}%
\pgfpathclose%
\pgfpathmoveto{\pgfqpoint{10.545036in}{1.301660in}}%
\pgfpathlineto{\pgfqpoint{10.545036in}{1.263600in}}%
\pgfpathlineto{\pgfqpoint{10.534679in}{1.263600in}}%
\pgfpathlineto{\pgfqpoint{10.534679in}{1.301317in}}%
\pgfpathquadraticcurveto{\pgfqpoint{10.534679in}{1.310269in}}{\pgfqpoint{10.531184in}{1.314700in}}%
\pgfpathquadraticcurveto{\pgfqpoint{10.527690in}{1.319149in}}{\pgfqpoint{10.520719in}{1.319149in}}%
\pgfpathquadraticcurveto{\pgfqpoint{10.512326in}{1.319149in}}{\pgfqpoint{10.507480in}{1.313800in}}%
\pgfpathquadraticcurveto{\pgfqpoint{10.502635in}{1.308468in}}{\pgfqpoint{10.502635in}{1.299228in}}%
\pgfpathlineto{\pgfqpoint{10.502635in}{1.263600in}}%
\pgfpathlineto{\pgfqpoint{10.492224in}{1.263600in}}%
\pgfpathlineto{\pgfqpoint{10.492224in}{1.351193in}}%
\pgfpathlineto{\pgfqpoint{10.502635in}{1.351193in}}%
\pgfpathlineto{\pgfqpoint{10.502635in}{1.316844in}}%
\pgfpathquadraticcurveto{\pgfqpoint{10.506363in}{1.322536in}}{\pgfqpoint{10.511389in}{1.325346in}}%
\pgfpathquadraticcurveto{\pgfqpoint{10.516432in}{1.328156in}}{\pgfqpoint{10.523025in}{1.328156in}}%
\pgfpathquadraticcurveto{\pgfqpoint{10.533886in}{1.328156in}}{\pgfqpoint{10.539452in}{1.321437in}}%
\pgfpathquadraticcurveto{\pgfqpoint{10.545036in}{1.314718in}}{\pgfqpoint{10.545036in}{1.301660in}}%
\pgfpathlineto{\pgfqpoint{10.545036in}{1.301660in}}%
\pgfpathclose%
\pgfpathmoveto{\pgfqpoint{10.594335in}{1.295283in}}%
\pgfpathquadraticcurveto{\pgfqpoint{10.581780in}{1.295283in}}{\pgfqpoint{10.576935in}{1.292419in}}%
\pgfpathquadraticcurveto{\pgfqpoint{10.572090in}{1.289556in}}{\pgfqpoint{10.572090in}{1.282621in}}%
\pgfpathquadraticcurveto{\pgfqpoint{10.572090in}{1.277109in}}{\pgfqpoint{10.575728in}{1.273867in}}%
\pgfpathquadraticcurveto{\pgfqpoint{10.579367in}{1.270643in}}{\pgfqpoint{10.585599in}{1.270643in}}%
\pgfpathquadraticcurveto{\pgfqpoint{10.594227in}{1.270643in}}{\pgfqpoint{10.599432in}{1.276749in}}%
\pgfpathquadraticcurveto{\pgfqpoint{10.604638in}{1.282855in}}{\pgfqpoint{10.604638in}{1.292978in}}%
\pgfpathlineto{\pgfqpoint{10.604638in}{1.295283in}}%
\pgfpathlineto{\pgfqpoint{10.594335in}{1.295283in}}%
\pgfpathlineto{\pgfqpoint{10.594335in}{1.295283in}}%
\pgfpathclose%
\pgfpathmoveto{\pgfqpoint{10.614995in}{1.299570in}}%
\pgfpathlineto{\pgfqpoint{10.614995in}{1.263600in}}%
\pgfpathlineto{\pgfqpoint{10.604638in}{1.263600in}}%
\pgfpathlineto{\pgfqpoint{10.604638in}{1.273164in}}%
\pgfpathquadraticcurveto{\pgfqpoint{10.601089in}{1.267437in}}{\pgfqpoint{10.595794in}{1.264699in}}%
\pgfpathquadraticcurveto{\pgfqpoint{10.590498in}{1.261961in}}{\pgfqpoint{10.582843in}{1.261961in}}%
\pgfpathquadraticcurveto{\pgfqpoint{10.573171in}{1.261961in}}{\pgfqpoint{10.567443in}{1.267401in}}%
\pgfpathquadraticcurveto{\pgfqpoint{10.561733in}{1.272840in}}{\pgfqpoint{10.561733in}{1.281954in}}%
\pgfpathquadraticcurveto{\pgfqpoint{10.561733in}{1.292582in}}{\pgfqpoint{10.568848in}{1.297985in}}%
\pgfpathquadraticcurveto{\pgfqpoint{10.575980in}{1.303389in}}{\pgfqpoint{10.590102in}{1.303389in}}%
\pgfpathlineto{\pgfqpoint{10.604638in}{1.303389in}}%
\pgfpathlineto{\pgfqpoint{10.604638in}{1.304416in}}%
\pgfpathquadraticcurveto{\pgfqpoint{10.604638in}{1.311566in}}{\pgfqpoint{10.599937in}{1.315475in}}%
\pgfpathquadraticcurveto{\pgfqpoint{10.595235in}{1.319384in}}{\pgfqpoint{10.586734in}{1.319384in}}%
\pgfpathquadraticcurveto{\pgfqpoint{10.581330in}{1.319384in}}{\pgfqpoint{10.576197in}{1.318087in}}%
\pgfpathquadraticcurveto{\pgfqpoint{10.571081in}{1.316790in}}{\pgfqpoint{10.566362in}{1.314196in}}%
\pgfpathlineto{\pgfqpoint{10.566362in}{1.323779in}}%
\pgfpathquadraticcurveto{\pgfqpoint{10.572036in}{1.325976in}}{\pgfqpoint{10.577385in}{1.327057in}}%
\pgfpathquadraticcurveto{\pgfqpoint{10.582735in}{1.328156in}}{\pgfqpoint{10.587796in}{1.328156in}}%
\pgfpathquadraticcurveto{\pgfqpoint{10.601486in}{1.328156in}}{\pgfqpoint{10.608240in}{1.321059in}}%
\pgfpathquadraticcurveto{\pgfqpoint{10.614995in}{1.313980in}}{\pgfqpoint{10.614995in}{1.299570in}}%
\pgfpathlineto{\pgfqpoint{10.614995in}{1.299570in}}%
\pgfpathclose%
\pgfpathmoveto{\pgfqpoint{10.672850in}{1.316970in}}%
\pgfpathquadraticcurveto{\pgfqpoint{10.671103in}{1.317979in}}{\pgfqpoint{10.669049in}{1.318447in}}%
\pgfpathquadraticcurveto{\pgfqpoint{10.666996in}{1.318933in}}{\pgfqpoint{10.664528in}{1.318933in}}%
\pgfpathquadraticcurveto{\pgfqpoint{10.655738in}{1.318933in}}{\pgfqpoint{10.651037in}{1.313223in}}%
\pgfpathquadraticcurveto{\pgfqpoint{10.646336in}{1.307514in}}{\pgfqpoint{10.646336in}{1.296814in}}%
\pgfpathlineto{\pgfqpoint{10.646336in}{1.263600in}}%
\pgfpathlineto{\pgfqpoint{10.635925in}{1.263600in}}%
\pgfpathlineto{\pgfqpoint{10.635925in}{1.326643in}}%
\pgfpathlineto{\pgfqpoint{10.646336in}{1.326643in}}%
\pgfpathlineto{\pgfqpoint{10.646336in}{1.316844in}}%
\pgfpathquadraticcurveto{\pgfqpoint{10.649614in}{1.322590in}}{\pgfqpoint{10.654838in}{1.325364in}}%
\pgfpathquadraticcurveto{\pgfqpoint{10.660079in}{1.328156in}}{\pgfqpoint{10.667572in}{1.328156in}}%
\pgfpathquadraticcurveto{\pgfqpoint{10.668635in}{1.328156in}}{\pgfqpoint{10.669932in}{1.328011in}}%
\pgfpathquadraticcurveto{\pgfqpoint{10.671229in}{1.327885in}}{\pgfqpoint{10.672796in}{1.327597in}}%
\pgfpathlineto{\pgfqpoint{10.672850in}{1.316970in}}%
\pgfpathlineto{\pgfqpoint{10.672850in}{1.316970in}}%
\pgfpathclose%
\pgfpathmoveto{\pgfqpoint{10.735100in}{1.297715in}}%
\pgfpathlineto{\pgfqpoint{10.735100in}{1.292654in}}%
\pgfpathlineto{\pgfqpoint{10.687476in}{1.292654in}}%
\pgfpathquadraticcurveto{\pgfqpoint{10.688160in}{1.281954in}}{\pgfqpoint{10.693924in}{1.276353in}}%
\pgfpathquadraticcurveto{\pgfqpoint{10.699688in}{1.270751in}}{\pgfqpoint{10.709991in}{1.270751in}}%
\pgfpathquadraticcurveto{\pgfqpoint{10.715953in}{1.270751in}}{\pgfqpoint{10.721555in}{1.272210in}}%
\pgfpathquadraticcurveto{\pgfqpoint{10.727156in}{1.273669in}}{\pgfqpoint{10.732686in}{1.276605in}}%
\pgfpathlineto{\pgfqpoint{10.732686in}{1.266806in}}%
\pgfpathquadraticcurveto{\pgfqpoint{10.727102in}{1.264447in}}{\pgfqpoint{10.721248in}{1.263204in}}%
\pgfpathquadraticcurveto{\pgfqpoint{10.715395in}{1.261961in}}{\pgfqpoint{10.709378in}{1.261961in}}%
\pgfpathquadraticcurveto{\pgfqpoint{10.694284in}{1.261961in}}{\pgfqpoint{10.685476in}{1.270733in}}%
\pgfpathquadraticcurveto{\pgfqpoint{10.676668in}{1.279523in}}{\pgfqpoint{10.676668in}{1.294509in}}%
\pgfpathquadraticcurveto{\pgfqpoint{10.676668in}{1.309981in}}{\pgfqpoint{10.685026in}{1.319059in}}%
\pgfpathquadraticcurveto{\pgfqpoint{10.693384in}{1.328156in}}{\pgfqpoint{10.707577in}{1.328156in}}%
\pgfpathquadraticcurveto{\pgfqpoint{10.720294in}{1.328156in}}{\pgfqpoint{10.727697in}{1.319960in}}%
\pgfpathquadraticcurveto{\pgfqpoint{10.735100in}{1.311783in}}{\pgfqpoint{10.735100in}{1.297715in}}%
\pgfpathlineto{\pgfqpoint{10.735100in}{1.297715in}}%
\pgfpathclose%
\pgfpathmoveto{\pgfqpoint{10.724743in}{1.300759in}}%
\pgfpathquadraticcurveto{\pgfqpoint{10.724635in}{1.309243in}}{\pgfqpoint{10.719988in}{1.314304in}}%
\pgfpathquadraticcurveto{\pgfqpoint{10.715340in}{1.319384in}}{\pgfqpoint{10.707685in}{1.319384in}}%
\pgfpathquadraticcurveto{\pgfqpoint{10.699021in}{1.319384in}}{\pgfqpoint{10.693816in}{1.314484in}}%
\pgfpathquadraticcurveto{\pgfqpoint{10.688610in}{1.309585in}}{\pgfqpoint{10.687818in}{1.300687in}}%
\pgfpathlineto{\pgfqpoint{10.724743in}{1.300759in}}%
\pgfpathlineto{\pgfqpoint{10.724743in}{1.300759in}}%
\pgfpathclose%
\pgfpathmoveto{\pgfqpoint{10.792288in}{1.324787in}}%
\pgfpathlineto{\pgfqpoint{10.792288in}{1.314989in}}%
\pgfpathquadraticcurveto{\pgfqpoint{10.787911in}{1.317240in}}{\pgfqpoint{10.783174in}{1.318357in}}%
\pgfpathquadraticcurveto{\pgfqpoint{10.778455in}{1.319492in}}{\pgfqpoint{10.773376in}{1.319492in}}%
\pgfpathquadraticcurveto{\pgfqpoint{10.765666in}{1.319492in}}{\pgfqpoint{10.761812in}{1.317132in}}%
\pgfpathquadraticcurveto{\pgfqpoint{10.757957in}{1.314773in}}{\pgfqpoint{10.757957in}{1.310035in}}%
\pgfpathquadraticcurveto{\pgfqpoint{10.757957in}{1.306433in}}{\pgfqpoint{10.760713in}{1.304380in}}%
\pgfpathquadraticcurveto{\pgfqpoint{10.763469in}{1.302326in}}{\pgfqpoint{10.771809in}{1.300471in}}%
\pgfpathlineto{\pgfqpoint{10.775357in}{1.299678in}}%
\pgfpathquadraticcurveto{\pgfqpoint{10.786380in}{1.297319in}}{\pgfqpoint{10.791028in}{1.293014in}}%
\pgfpathquadraticcurveto{\pgfqpoint{10.795675in}{1.288709in}}{\pgfqpoint{10.795675in}{1.281000in}}%
\pgfpathquadraticcurveto{\pgfqpoint{10.795675in}{1.272210in}}{\pgfqpoint{10.788722in}{1.267076in}}%
\pgfpathquadraticcurveto{\pgfqpoint{10.781769in}{1.261961in}}{\pgfqpoint{10.769611in}{1.261961in}}%
\pgfpathquadraticcurveto{\pgfqpoint{10.764550in}{1.261961in}}{\pgfqpoint{10.759056in}{1.262952in}}%
\pgfpathquadraticcurveto{\pgfqpoint{10.753562in}{1.263942in}}{\pgfqpoint{10.747492in}{1.265906in}}%
\pgfpathlineto{\pgfqpoint{10.747492in}{1.276605in}}%
\pgfpathquadraticcurveto{\pgfqpoint{10.753238in}{1.273615in}}{\pgfqpoint{10.758804in}{1.272120in}}%
\pgfpathquadraticcurveto{\pgfqpoint{10.764370in}{1.270643in}}{\pgfqpoint{10.769845in}{1.270643in}}%
\pgfpathquadraticcurveto{\pgfqpoint{10.777158in}{1.270643in}}{\pgfqpoint{10.781085in}{1.273146in}}%
\pgfpathquadraticcurveto{\pgfqpoint{10.785030in}{1.275650in}}{\pgfqpoint{10.785030in}{1.280207in}}%
\pgfpathquadraticcurveto{\pgfqpoint{10.785030in}{1.284422in}}{\pgfqpoint{10.782184in}{1.286674in}}%
\pgfpathquadraticcurveto{\pgfqpoint{10.779356in}{1.288925in}}{\pgfqpoint{10.769719in}{1.291014in}}%
\pgfpathlineto{\pgfqpoint{10.766117in}{1.291861in}}%
\pgfpathquadraticcurveto{\pgfqpoint{10.756498in}{1.293878in}}{\pgfqpoint{10.752211in}{1.298075in}}%
\pgfpathquadraticcurveto{\pgfqpoint{10.747942in}{1.302272in}}{\pgfqpoint{10.747942in}{1.309585in}}%
\pgfpathquadraticcurveto{\pgfqpoint{10.747942in}{1.318483in}}{\pgfqpoint{10.754247in}{1.323310in}}%
\pgfpathquadraticcurveto{\pgfqpoint{10.760551in}{1.328156in}}{\pgfqpoint{10.772151in}{1.328156in}}%
\pgfpathquadraticcurveto{\pgfqpoint{10.777879in}{1.328156in}}{\pgfqpoint{10.782940in}{1.327309in}}%
\pgfpathquadraticcurveto{\pgfqpoint{10.788020in}{1.326480in}}{\pgfqpoint{10.792288in}{1.324787in}}%
\pgfpathlineto{\pgfqpoint{10.792288in}{1.324787in}}%
\pgfpathclose%
\pgfusepath{fill}%
\end{pgfscope}%
\begin{pgfscope}%
\pgfsetbuttcap%
\pgfsetroundjoin%
\definecolor{currentfill}{rgb}{0.309804,0.372549,0.419608}%
\pgfsetfillcolor{currentfill}%
\pgfsetlinewidth{0.000000pt}%
\definecolor{currentstroke}{rgb}{0.309804,0.372549,0.419608}%
\pgfsetstrokecolor{currentstroke}%
\pgfsetdash{}{0pt}%
\pgfpathmoveto{\pgfqpoint{10.881559in}{1.288475in}}%
\pgfpathlineto{\pgfqpoint{10.881559in}{1.326643in}}%
\pgfpathlineto{\pgfqpoint{10.891916in}{1.326643in}}%
\pgfpathlineto{\pgfqpoint{10.891916in}{1.288871in}}%
\pgfpathquadraticcurveto{\pgfqpoint{10.891916in}{1.279919in}}{\pgfqpoint{10.895392in}{1.275434in}}%
\pgfpathquadraticcurveto{\pgfqpoint{10.898887in}{1.270967in}}{\pgfqpoint{10.905875in}{1.270967in}}%
\pgfpathquadraticcurveto{\pgfqpoint{10.914251in}{1.270967in}}{\pgfqpoint{10.919114in}{1.276317in}}%
\pgfpathquadraticcurveto{\pgfqpoint{10.923996in}{1.281666in}}{\pgfqpoint{10.923996in}{1.290906in}}%
\pgfpathlineto{\pgfqpoint{10.923996in}{1.326643in}}%
\pgfpathlineto{\pgfqpoint{10.934353in}{1.326643in}}%
\pgfpathlineto{\pgfqpoint{10.934353in}{1.263600in}}%
\pgfpathlineto{\pgfqpoint{10.923996in}{1.263600in}}%
\pgfpathlineto{\pgfqpoint{10.923996in}{1.273291in}}%
\pgfpathquadraticcurveto{\pgfqpoint{10.920231in}{1.267545in}}{\pgfqpoint{10.915242in}{1.264753in}}%
\pgfpathquadraticcurveto{\pgfqpoint{10.910270in}{1.261961in}}{\pgfqpoint{10.903678in}{1.261961in}}%
\pgfpathquadraticcurveto{\pgfqpoint{10.892817in}{1.261961in}}{\pgfqpoint{10.887179in}{1.268715in}}%
\pgfpathquadraticcurveto{\pgfqpoint{10.881559in}{1.275470in}}{\pgfqpoint{10.881559in}{1.288475in}}%
\pgfpathlineto{\pgfqpoint{10.881559in}{1.288475in}}%
\pgfpathclose%
\pgfpathmoveto{\pgfqpoint{10.907623in}{1.328156in}}%
\pgfpathlineto{\pgfqpoint{10.907623in}{1.328156in}}%
\pgfpathlineto{\pgfqpoint{10.907623in}{1.328156in}}%
\pgfpathclose%
\pgfpathmoveto{\pgfqpoint{10.995864in}{1.324787in}}%
\pgfpathlineto{\pgfqpoint{10.995864in}{1.314989in}}%
\pgfpathquadraticcurveto{\pgfqpoint{10.991487in}{1.317240in}}{\pgfqpoint{10.986750in}{1.318357in}}%
\pgfpathquadraticcurveto{\pgfqpoint{10.982031in}{1.319492in}}{\pgfqpoint{10.976951in}{1.319492in}}%
\pgfpathquadraticcurveto{\pgfqpoint{10.969242in}{1.319492in}}{\pgfqpoint{10.965388in}{1.317132in}}%
\pgfpathquadraticcurveto{\pgfqpoint{10.961533in}{1.314773in}}{\pgfqpoint{10.961533in}{1.310035in}}%
\pgfpathquadraticcurveto{\pgfqpoint{10.961533in}{1.306433in}}{\pgfqpoint{10.964289in}{1.304380in}}%
\pgfpathquadraticcurveto{\pgfqpoint{10.967045in}{1.302326in}}{\pgfqpoint{10.975384in}{1.300471in}}%
\pgfpathlineto{\pgfqpoint{10.978933in}{1.299678in}}%
\pgfpathquadraticcurveto{\pgfqpoint{10.989956in}{1.297319in}}{\pgfqpoint{10.994603in}{1.293014in}}%
\pgfpathquadraticcurveto{\pgfqpoint{10.999250in}{1.288709in}}{\pgfqpoint{10.999250in}{1.281000in}}%
\pgfpathquadraticcurveto{\pgfqpoint{10.999250in}{1.272210in}}{\pgfqpoint{10.992298in}{1.267076in}}%
\pgfpathquadraticcurveto{\pgfqpoint{10.985345in}{1.261961in}}{\pgfqpoint{10.973187in}{1.261961in}}%
\pgfpathquadraticcurveto{\pgfqpoint{10.968125in}{1.261961in}}{\pgfqpoint{10.962632in}{1.262952in}}%
\pgfpathquadraticcurveto{\pgfqpoint{10.957138in}{1.263942in}}{\pgfqpoint{10.951068in}{1.265906in}}%
\pgfpathlineto{\pgfqpoint{10.951068in}{1.276605in}}%
\pgfpathquadraticcurveto{\pgfqpoint{10.956814in}{1.273615in}}{\pgfqpoint{10.962380in}{1.272120in}}%
\pgfpathquadraticcurveto{\pgfqpoint{10.967945in}{1.270643in}}{\pgfqpoint{10.973421in}{1.270643in}}%
\pgfpathquadraticcurveto{\pgfqpoint{10.980734in}{1.270643in}}{\pgfqpoint{10.984661in}{1.273146in}}%
\pgfpathquadraticcurveto{\pgfqpoint{10.988605in}{1.275650in}}{\pgfqpoint{10.988605in}{1.280207in}}%
\pgfpathquadraticcurveto{\pgfqpoint{10.988605in}{1.284422in}}{\pgfqpoint{10.985759in}{1.286674in}}%
\pgfpathquadraticcurveto{\pgfqpoint{10.982931in}{1.288925in}}{\pgfqpoint{10.973295in}{1.291014in}}%
\pgfpathlineto{\pgfqpoint{10.969692in}{1.291861in}}%
\pgfpathquadraticcurveto{\pgfqpoint{10.960074in}{1.293878in}}{\pgfqpoint{10.955787in}{1.298075in}}%
\pgfpathquadraticcurveto{\pgfqpoint{10.951518in}{1.302272in}}{\pgfqpoint{10.951518in}{1.309585in}}%
\pgfpathquadraticcurveto{\pgfqpoint{10.951518in}{1.318483in}}{\pgfqpoint{10.957822in}{1.323310in}}%
\pgfpathquadraticcurveto{\pgfqpoint{10.964127in}{1.328156in}}{\pgfqpoint{10.975727in}{1.328156in}}%
\pgfpathquadraticcurveto{\pgfqpoint{10.981454in}{1.328156in}}{\pgfqpoint{10.986516in}{1.327309in}}%
\pgfpathquadraticcurveto{\pgfqpoint{10.991595in}{1.326480in}}{\pgfqpoint{10.995864in}{1.324787in}}%
\pgfpathlineto{\pgfqpoint{10.995864in}{1.324787in}}%
\pgfpathclose%
\pgfpathmoveto{\pgfqpoint{11.069660in}{1.297715in}}%
\pgfpathlineto{\pgfqpoint{11.069660in}{1.292654in}}%
\pgfpathlineto{\pgfqpoint{11.022036in}{1.292654in}}%
\pgfpathquadraticcurveto{\pgfqpoint{11.022720in}{1.281954in}}{\pgfqpoint{11.028484in}{1.276353in}}%
\pgfpathquadraticcurveto{\pgfqpoint{11.034248in}{1.270751in}}{\pgfqpoint{11.044551in}{1.270751in}}%
\pgfpathquadraticcurveto{\pgfqpoint{11.050513in}{1.270751in}}{\pgfqpoint{11.056115in}{1.272210in}}%
\pgfpathquadraticcurveto{\pgfqpoint{11.061717in}{1.273669in}}{\pgfqpoint{11.067246in}{1.276605in}}%
\pgfpathlineto{\pgfqpoint{11.067246in}{1.266806in}}%
\pgfpathquadraticcurveto{\pgfqpoint{11.061663in}{1.264447in}}{\pgfqpoint{11.055809in}{1.263204in}}%
\pgfpathquadraticcurveto{\pgfqpoint{11.049955in}{1.261961in}}{\pgfqpoint{11.043939in}{1.261961in}}%
\pgfpathquadraticcurveto{\pgfqpoint{11.028844in}{1.261961in}}{\pgfqpoint{11.020036in}{1.270733in}}%
\pgfpathquadraticcurveto{\pgfqpoint{11.011229in}{1.279523in}}{\pgfqpoint{11.011229in}{1.294509in}}%
\pgfpathquadraticcurveto{\pgfqpoint{11.011229in}{1.309981in}}{\pgfqpoint{11.019586in}{1.319059in}}%
\pgfpathquadraticcurveto{\pgfqpoint{11.027944in}{1.328156in}}{\pgfqpoint{11.042137in}{1.328156in}}%
\pgfpathquadraticcurveto{\pgfqpoint{11.054854in}{1.328156in}}{\pgfqpoint{11.062257in}{1.319960in}}%
\pgfpathquadraticcurveto{\pgfqpoint{11.069660in}{1.311783in}}{\pgfqpoint{11.069660in}{1.297715in}}%
\pgfpathlineto{\pgfqpoint{11.069660in}{1.297715in}}%
\pgfpathclose%
\pgfpathmoveto{\pgfqpoint{11.059303in}{1.300759in}}%
\pgfpathquadraticcurveto{\pgfqpoint{11.059195in}{1.309243in}}{\pgfqpoint{11.054548in}{1.314304in}}%
\pgfpathquadraticcurveto{\pgfqpoint{11.049901in}{1.319384in}}{\pgfqpoint{11.042245in}{1.319384in}}%
\pgfpathquadraticcurveto{\pgfqpoint{11.033582in}{1.319384in}}{\pgfqpoint{11.028376in}{1.314484in}}%
\pgfpathquadraticcurveto{\pgfqpoint{11.023171in}{1.309585in}}{\pgfqpoint{11.022378in}{1.300687in}}%
\pgfpathlineto{\pgfqpoint{11.059303in}{1.300759in}}%
\pgfpathlineto{\pgfqpoint{11.059303in}{1.300759in}}%
\pgfpathclose%
\pgfusepath{fill}%
\end{pgfscope}%
\begin{pgfscope}%
\pgfsetbuttcap%
\pgfsetroundjoin%
\definecolor{currentfill}{rgb}{0.309804,0.372549,0.419608}%
\pgfsetfillcolor{currentfill}%
\pgfsetlinewidth{0.000000pt}%
\definecolor{currentstroke}{rgb}{0.309804,0.372549,0.419608}%
\pgfsetstrokecolor{currentstroke}%
\pgfsetdash{}{0pt}%
\pgfpathmoveto{\pgfqpoint{11.184835in}{1.316970in}}%
\pgfpathquadraticcurveto{\pgfqpoint{11.183088in}{1.317979in}}{\pgfqpoint{11.181035in}{1.318447in}}%
\pgfpathquadraticcurveto{\pgfqpoint{11.178981in}{1.318933in}}{\pgfqpoint{11.176514in}{1.318933in}}%
\pgfpathquadraticcurveto{\pgfqpoint{11.167724in}{1.318933in}}{\pgfqpoint{11.163023in}{1.313223in}}%
\pgfpathquadraticcurveto{\pgfqpoint{11.158321in}{1.307514in}}{\pgfqpoint{11.158321in}{1.296814in}}%
\pgfpathlineto{\pgfqpoint{11.158321in}{1.263600in}}%
\pgfpathlineto{\pgfqpoint{11.147910in}{1.263600in}}%
\pgfpathlineto{\pgfqpoint{11.147910in}{1.326643in}}%
\pgfpathlineto{\pgfqpoint{11.158321in}{1.326643in}}%
\pgfpathlineto{\pgfqpoint{11.158321in}{1.316844in}}%
\pgfpathquadraticcurveto{\pgfqpoint{11.161600in}{1.322590in}}{\pgfqpoint{11.166823in}{1.325364in}}%
\pgfpathquadraticcurveto{\pgfqpoint{11.172065in}{1.328156in}}{\pgfqpoint{11.179558in}{1.328156in}}%
\pgfpathquadraticcurveto{\pgfqpoint{11.180620in}{1.328156in}}{\pgfqpoint{11.181917in}{1.328011in}}%
\pgfpathquadraticcurveto{\pgfqpoint{11.183214in}{1.327885in}}{\pgfqpoint{11.184781in}{1.327597in}}%
\pgfpathlineto{\pgfqpoint{11.184835in}{1.316970in}}%
\pgfpathlineto{\pgfqpoint{11.184835in}{1.316970in}}%
\pgfpathclose%
\pgfpathmoveto{\pgfqpoint{11.247085in}{1.297715in}}%
\pgfpathlineto{\pgfqpoint{11.247085in}{1.292654in}}%
\pgfpathlineto{\pgfqpoint{11.199461in}{1.292654in}}%
\pgfpathquadraticcurveto{\pgfqpoint{11.200146in}{1.281954in}}{\pgfqpoint{11.205910in}{1.276353in}}%
\pgfpathquadraticcurveto{\pgfqpoint{11.211673in}{1.270751in}}{\pgfqpoint{11.221976in}{1.270751in}}%
\pgfpathquadraticcurveto{\pgfqpoint{11.227938in}{1.270751in}}{\pgfqpoint{11.233540in}{1.272210in}}%
\pgfpathquadraticcurveto{\pgfqpoint{11.239142in}{1.273669in}}{\pgfqpoint{11.244672in}{1.276605in}}%
\pgfpathlineto{\pgfqpoint{11.244672in}{1.266806in}}%
\pgfpathquadraticcurveto{\pgfqpoint{11.239088in}{1.264447in}}{\pgfqpoint{11.233234in}{1.263204in}}%
\pgfpathquadraticcurveto{\pgfqpoint{11.227380in}{1.261961in}}{\pgfqpoint{11.221364in}{1.261961in}}%
\pgfpathquadraticcurveto{\pgfqpoint{11.206270in}{1.261961in}}{\pgfqpoint{11.197462in}{1.270733in}}%
\pgfpathquadraticcurveto{\pgfqpoint{11.188654in}{1.279523in}}{\pgfqpoint{11.188654in}{1.294509in}}%
\pgfpathquadraticcurveto{\pgfqpoint{11.188654in}{1.309981in}}{\pgfqpoint{11.197012in}{1.319059in}}%
\pgfpathquadraticcurveto{\pgfqpoint{11.205369in}{1.328156in}}{\pgfqpoint{11.219563in}{1.328156in}}%
\pgfpathquadraticcurveto{\pgfqpoint{11.232279in}{1.328156in}}{\pgfqpoint{11.239682in}{1.319960in}}%
\pgfpathquadraticcurveto{\pgfqpoint{11.247085in}{1.311783in}}{\pgfqpoint{11.247085in}{1.297715in}}%
\pgfpathlineto{\pgfqpoint{11.247085in}{1.297715in}}%
\pgfpathclose%
\pgfpathmoveto{\pgfqpoint{11.236728in}{1.300759in}}%
\pgfpathquadraticcurveto{\pgfqpoint{11.236620in}{1.309243in}}{\pgfqpoint{11.231973in}{1.314304in}}%
\pgfpathquadraticcurveto{\pgfqpoint{11.227326in}{1.319384in}}{\pgfqpoint{11.219671in}{1.319384in}}%
\pgfpathquadraticcurveto{\pgfqpoint{11.211007in}{1.319384in}}{\pgfqpoint{11.205801in}{1.314484in}}%
\pgfpathquadraticcurveto{\pgfqpoint{11.200596in}{1.309585in}}{\pgfqpoint{11.199803in}{1.300687in}}%
\pgfpathlineto{\pgfqpoint{11.236728in}{1.300759in}}%
\pgfpathlineto{\pgfqpoint{11.236728in}{1.300759in}}%
\pgfpathclose%
\pgfpathmoveto{\pgfqpoint{11.274338in}{1.344547in}}%
\pgfpathlineto{\pgfqpoint{11.274338in}{1.326643in}}%
\pgfpathlineto{\pgfqpoint{11.295664in}{1.326643in}}%
\pgfpathlineto{\pgfqpoint{11.295664in}{1.318591in}}%
\pgfpathlineto{\pgfqpoint{11.274338in}{1.318591in}}%
\pgfpathlineto{\pgfqpoint{11.274338in}{1.284368in}}%
\pgfpathquadraticcurveto{\pgfqpoint{11.274338in}{1.276659in}}{\pgfqpoint{11.276445in}{1.274461in}}%
\pgfpathquadraticcurveto{\pgfqpoint{11.278553in}{1.272264in}}{\pgfqpoint{11.285037in}{1.272264in}}%
\pgfpathlineto{\pgfqpoint{11.295664in}{1.272264in}}%
\pgfpathlineto{\pgfqpoint{11.295664in}{1.263600in}}%
\pgfpathlineto{\pgfqpoint{11.285037in}{1.263600in}}%
\pgfpathquadraticcurveto{\pgfqpoint{11.273041in}{1.263600in}}{\pgfqpoint{11.268484in}{1.268067in}}%
\pgfpathquadraticcurveto{\pgfqpoint{11.263927in}{1.272552in}}{\pgfqpoint{11.263927in}{1.284368in}}%
\pgfpathlineto{\pgfqpoint{11.263927in}{1.318591in}}%
\pgfpathlineto{\pgfqpoint{11.256326in}{1.318591in}}%
\pgfpathlineto{\pgfqpoint{11.256326in}{1.326643in}}%
\pgfpathlineto{\pgfqpoint{11.263927in}{1.326643in}}%
\pgfpathlineto{\pgfqpoint{11.263927in}{1.344547in}}%
\pgfpathlineto{\pgfqpoint{11.274338in}{1.344547in}}%
\pgfpathlineto{\pgfqpoint{11.274338in}{1.344547in}}%
\pgfpathclose%
\pgfpathmoveto{\pgfqpoint{11.337939in}{1.295283in}}%
\pgfpathquadraticcurveto{\pgfqpoint{11.325384in}{1.295283in}}{\pgfqpoint{11.320539in}{1.292419in}}%
\pgfpathquadraticcurveto{\pgfqpoint{11.315694in}{1.289556in}}{\pgfqpoint{11.315694in}{1.282621in}}%
\pgfpathquadraticcurveto{\pgfqpoint{11.315694in}{1.277109in}}{\pgfqpoint{11.319332in}{1.273867in}}%
\pgfpathquadraticcurveto{\pgfqpoint{11.322970in}{1.270643in}}{\pgfqpoint{11.329203in}{1.270643in}}%
\pgfpathquadraticcurveto{\pgfqpoint{11.337831in}{1.270643in}}{\pgfqpoint{11.343036in}{1.276749in}}%
\pgfpathquadraticcurveto{\pgfqpoint{11.348242in}{1.282855in}}{\pgfqpoint{11.348242in}{1.292978in}}%
\pgfpathlineto{\pgfqpoint{11.348242in}{1.295283in}}%
\pgfpathlineto{\pgfqpoint{11.337939in}{1.295283in}}%
\pgfpathlineto{\pgfqpoint{11.337939in}{1.295283in}}%
\pgfpathclose%
\pgfpathmoveto{\pgfqpoint{11.358599in}{1.299570in}}%
\pgfpathlineto{\pgfqpoint{11.358599in}{1.263600in}}%
\pgfpathlineto{\pgfqpoint{11.348242in}{1.263600in}}%
\pgfpathlineto{\pgfqpoint{11.348242in}{1.273164in}}%
\pgfpathquadraticcurveto{\pgfqpoint{11.344693in}{1.267437in}}{\pgfqpoint{11.339398in}{1.264699in}}%
\pgfpathquadraticcurveto{\pgfqpoint{11.334102in}{1.261961in}}{\pgfqpoint{11.326447in}{1.261961in}}%
\pgfpathquadraticcurveto{\pgfqpoint{11.316774in}{1.261961in}}{\pgfqpoint{11.311046in}{1.267401in}}%
\pgfpathquadraticcurveto{\pgfqpoint{11.305337in}{1.272840in}}{\pgfqpoint{11.305337in}{1.281954in}}%
\pgfpathquadraticcurveto{\pgfqpoint{11.305337in}{1.292582in}}{\pgfqpoint{11.312451in}{1.297985in}}%
\pgfpathquadraticcurveto{\pgfqpoint{11.319584in}{1.303389in}}{\pgfqpoint{11.333706in}{1.303389in}}%
\pgfpathlineto{\pgfqpoint{11.348242in}{1.303389in}}%
\pgfpathlineto{\pgfqpoint{11.348242in}{1.304416in}}%
\pgfpathquadraticcurveto{\pgfqpoint{11.348242in}{1.311566in}}{\pgfqpoint{11.343540in}{1.315475in}}%
\pgfpathquadraticcurveto{\pgfqpoint{11.338839in}{1.319384in}}{\pgfqpoint{11.330337in}{1.319384in}}%
\pgfpathquadraticcurveto{\pgfqpoint{11.324934in}{1.319384in}}{\pgfqpoint{11.319800in}{1.318087in}}%
\pgfpathquadraticcurveto{\pgfqpoint{11.314685in}{1.316790in}}{\pgfqpoint{11.309966in}{1.314196in}}%
\pgfpathlineto{\pgfqpoint{11.309966in}{1.323779in}}%
\pgfpathquadraticcurveto{\pgfqpoint{11.315640in}{1.325976in}}{\pgfqpoint{11.320989in}{1.327057in}}%
\pgfpathquadraticcurveto{\pgfqpoint{11.326339in}{1.328156in}}{\pgfqpoint{11.331400in}{1.328156in}}%
\pgfpathquadraticcurveto{\pgfqpoint{11.345089in}{1.328156in}}{\pgfqpoint{11.351844in}{1.321059in}}%
\pgfpathquadraticcurveto{\pgfqpoint{11.358599in}{1.313980in}}{\pgfqpoint{11.358599in}{1.299570in}}%
\pgfpathlineto{\pgfqpoint{11.358599in}{1.299570in}}%
\pgfpathclose%
\pgfpathmoveto{\pgfqpoint{11.379925in}{1.326643in}}%
\pgfpathlineto{\pgfqpoint{11.390282in}{1.326643in}}%
\pgfpathlineto{\pgfqpoint{11.390282in}{1.263600in}}%
\pgfpathlineto{\pgfqpoint{11.379925in}{1.263600in}}%
\pgfpathlineto{\pgfqpoint{11.379925in}{1.326643in}}%
\pgfpathlineto{\pgfqpoint{11.379925in}{1.326643in}}%
\pgfpathclose%
\pgfpathmoveto{\pgfqpoint{11.379925in}{1.351193in}}%
\pgfpathlineto{\pgfqpoint{11.390282in}{1.351193in}}%
\pgfpathlineto{\pgfqpoint{11.390282in}{1.338062in}}%
\pgfpathlineto{\pgfqpoint{11.379925in}{1.338062in}}%
\pgfpathlineto{\pgfqpoint{11.379925in}{1.351193in}}%
\pgfpathlineto{\pgfqpoint{11.379925in}{1.351193in}}%
\pgfpathclose%
\pgfpathmoveto{\pgfqpoint{11.464366in}{1.301660in}}%
\pgfpathlineto{\pgfqpoint{11.464366in}{1.263600in}}%
\pgfpathlineto{\pgfqpoint{11.454009in}{1.263600in}}%
\pgfpathlineto{\pgfqpoint{11.454009in}{1.301317in}}%
\pgfpathquadraticcurveto{\pgfqpoint{11.454009in}{1.310269in}}{\pgfqpoint{11.450515in}{1.314700in}}%
\pgfpathquadraticcurveto{\pgfqpoint{11.447020in}{1.319149in}}{\pgfqpoint{11.440049in}{1.319149in}}%
\pgfpathquadraticcurveto{\pgfqpoint{11.431656in}{1.319149in}}{\pgfqpoint{11.426811in}{1.313800in}}%
\pgfpathquadraticcurveto{\pgfqpoint{11.421965in}{1.308468in}}{\pgfqpoint{11.421965in}{1.299228in}}%
\pgfpathlineto{\pgfqpoint{11.421965in}{1.263600in}}%
\pgfpathlineto{\pgfqpoint{11.411554in}{1.263600in}}%
\pgfpathlineto{\pgfqpoint{11.411554in}{1.326643in}}%
\pgfpathlineto{\pgfqpoint{11.421965in}{1.326643in}}%
\pgfpathlineto{\pgfqpoint{11.421965in}{1.316844in}}%
\pgfpathquadraticcurveto{\pgfqpoint{11.425694in}{1.322536in}}{\pgfqpoint{11.430719in}{1.325346in}}%
\pgfpathquadraticcurveto{\pgfqpoint{11.435763in}{1.328156in}}{\pgfqpoint{11.442355in}{1.328156in}}%
\pgfpathquadraticcurveto{\pgfqpoint{11.453216in}{1.328156in}}{\pgfqpoint{11.458782in}{1.321437in}}%
\pgfpathquadraticcurveto{\pgfqpoint{11.464366in}{1.314718in}}{\pgfqpoint{11.464366in}{1.301660in}}%
\pgfpathlineto{\pgfqpoint{11.464366in}{1.301660in}}%
\pgfpathclose%
\pgfpathmoveto{\pgfqpoint{11.538936in}{1.297715in}}%
\pgfpathlineto{\pgfqpoint{11.538936in}{1.292654in}}%
\pgfpathlineto{\pgfqpoint{11.491312in}{1.292654in}}%
\pgfpathquadraticcurveto{\pgfqpoint{11.491997in}{1.281954in}}{\pgfqpoint{11.497760in}{1.276353in}}%
\pgfpathquadraticcurveto{\pgfqpoint{11.503524in}{1.270751in}}{\pgfqpoint{11.513827in}{1.270751in}}%
\pgfpathquadraticcurveto{\pgfqpoint{11.519789in}{1.270751in}}{\pgfqpoint{11.525391in}{1.272210in}}%
\pgfpathquadraticcurveto{\pgfqpoint{11.530993in}{1.273669in}}{\pgfqpoint{11.536523in}{1.276605in}}%
\pgfpathlineto{\pgfqpoint{11.536523in}{1.266806in}}%
\pgfpathquadraticcurveto{\pgfqpoint{11.530939in}{1.264447in}}{\pgfqpoint{11.525085in}{1.263204in}}%
\pgfpathquadraticcurveto{\pgfqpoint{11.519231in}{1.261961in}}{\pgfqpoint{11.513215in}{1.261961in}}%
\pgfpathquadraticcurveto{\pgfqpoint{11.498121in}{1.261961in}}{\pgfqpoint{11.489313in}{1.270733in}}%
\pgfpathquadraticcurveto{\pgfqpoint{11.480505in}{1.279523in}}{\pgfqpoint{11.480505in}{1.294509in}}%
\pgfpathquadraticcurveto{\pgfqpoint{11.480505in}{1.309981in}}{\pgfqpoint{11.488862in}{1.319059in}}%
\pgfpathquadraticcurveto{\pgfqpoint{11.497220in}{1.328156in}}{\pgfqpoint{11.511414in}{1.328156in}}%
\pgfpathquadraticcurveto{\pgfqpoint{11.524130in}{1.328156in}}{\pgfqpoint{11.531533in}{1.319960in}}%
\pgfpathquadraticcurveto{\pgfqpoint{11.538936in}{1.311783in}}{\pgfqpoint{11.538936in}{1.297715in}}%
\pgfpathlineto{\pgfqpoint{11.538936in}{1.297715in}}%
\pgfpathclose%
\pgfpathmoveto{\pgfqpoint{11.528579in}{1.300759in}}%
\pgfpathquadraticcurveto{\pgfqpoint{11.528471in}{1.309243in}}{\pgfqpoint{11.523824in}{1.314304in}}%
\pgfpathquadraticcurveto{\pgfqpoint{11.519177in}{1.319384in}}{\pgfqpoint{11.511522in}{1.319384in}}%
\pgfpathquadraticcurveto{\pgfqpoint{11.502858in}{1.319384in}}{\pgfqpoint{11.497652in}{1.314484in}}%
\pgfpathquadraticcurveto{\pgfqpoint{11.492447in}{1.309585in}}{\pgfqpoint{11.491654in}{1.300687in}}%
\pgfpathlineto{\pgfqpoint{11.528579in}{1.300759in}}%
\pgfpathlineto{\pgfqpoint{11.528579in}{1.300759in}}%
\pgfpathclose%
\pgfpathmoveto{\pgfqpoint{11.597422in}{1.317078in}}%
\pgfpathlineto{\pgfqpoint{11.597422in}{1.351193in}}%
\pgfpathlineto{\pgfqpoint{11.607779in}{1.351193in}}%
\pgfpathlineto{\pgfqpoint{11.607779in}{1.263600in}}%
\pgfpathlineto{\pgfqpoint{11.597422in}{1.263600in}}%
\pgfpathlineto{\pgfqpoint{11.597422in}{1.273056in}}%
\pgfpathquadraticcurveto{\pgfqpoint{11.594161in}{1.267437in}}{\pgfqpoint{11.589172in}{1.264699in}}%
\pgfpathquadraticcurveto{\pgfqpoint{11.584201in}{1.261961in}}{\pgfqpoint{11.577212in}{1.261961in}}%
\pgfpathquadraticcurveto{\pgfqpoint{11.565792in}{1.261961in}}{\pgfqpoint{11.558605in}{1.271075in}}%
\pgfpathquadraticcurveto{\pgfqpoint{11.551437in}{1.280207in}}{\pgfqpoint{11.551437in}{1.295067in}}%
\pgfpathquadraticcurveto{\pgfqpoint{11.551437in}{1.309927in}}{\pgfqpoint{11.558605in}{1.319041in}}%
\pgfpathquadraticcurveto{\pgfqpoint{11.565792in}{1.328156in}}{\pgfqpoint{11.577212in}{1.328156in}}%
\pgfpathquadraticcurveto{\pgfqpoint{11.584201in}{1.328156in}}{\pgfqpoint{11.589172in}{1.325418in}}%
\pgfpathquadraticcurveto{\pgfqpoint{11.594161in}{1.322698in}}{\pgfqpoint{11.597422in}{1.317078in}}%
\pgfpathlineto{\pgfqpoint{11.597422in}{1.317078in}}%
\pgfpathclose%
\pgfpathmoveto{\pgfqpoint{11.562136in}{1.295067in}}%
\pgfpathquadraticcurveto{\pgfqpoint{11.562136in}{1.283648in}}{\pgfqpoint{11.566837in}{1.277145in}}%
\pgfpathquadraticcurveto{\pgfqpoint{11.571538in}{1.270643in}}{\pgfqpoint{11.579752in}{1.270643in}}%
\pgfpathquadraticcurveto{\pgfqpoint{11.587965in}{1.270643in}}{\pgfqpoint{11.592684in}{1.277145in}}%
\pgfpathquadraticcurveto{\pgfqpoint{11.597422in}{1.283648in}}{\pgfqpoint{11.597422in}{1.295067in}}%
\pgfpathquadraticcurveto{\pgfqpoint{11.597422in}{1.306487in}}{\pgfqpoint{11.592684in}{1.312989in}}%
\pgfpathquadraticcurveto{\pgfqpoint{11.587965in}{1.319492in}}{\pgfqpoint{11.579752in}{1.319492in}}%
\pgfpathquadraticcurveto{\pgfqpoint{11.571538in}{1.319492in}}{\pgfqpoint{11.566837in}{1.312989in}}%
\pgfpathquadraticcurveto{\pgfqpoint{11.562136in}{1.306487in}}{\pgfqpoint{11.562136in}{1.295067in}}%
\pgfpathlineto{\pgfqpoint{11.562136in}{1.295067in}}%
\pgfpathclose%
\pgfusepath{fill}%
\end{pgfscope}%
\begin{pgfscope}%
\pgfsetbuttcap%
\pgfsetroundjoin%
\definecolor{currentfill}{rgb}{0.309804,0.372549,0.419608}%
\pgfsetfillcolor{currentfill}%
\pgfsetlinewidth{0.000000pt}%
\definecolor{currentstroke}{rgb}{0.309804,0.372549,0.419608}%
\pgfsetstrokecolor{currentstroke}%
\pgfsetdash{}{0pt}%
\pgfpathmoveto{\pgfqpoint{11.746864in}{1.324229in}}%
\pgfpathlineto{\pgfqpoint{11.746864in}{1.314538in}}%
\pgfpathquadraticcurveto{\pgfqpoint{11.742469in}{1.316970in}}{\pgfqpoint{11.738056in}{1.318177in}}%
\pgfpathquadraticcurveto{\pgfqpoint{11.733643in}{1.319384in}}{\pgfqpoint{11.729140in}{1.319384in}}%
\pgfpathquadraticcurveto{\pgfqpoint{11.719053in}{1.319384in}}{\pgfqpoint{11.713470in}{1.312989in}}%
\pgfpathquadraticcurveto{\pgfqpoint{11.707904in}{1.306613in}}{\pgfqpoint{11.707904in}{1.295067in}}%
\pgfpathquadraticcurveto{\pgfqpoint{11.707904in}{1.283521in}}{\pgfqpoint{11.713470in}{1.277127in}}%
\pgfpathquadraticcurveto{\pgfqpoint{11.719053in}{1.270751in}}{\pgfqpoint{11.729140in}{1.270751in}}%
\pgfpathquadraticcurveto{\pgfqpoint{11.733643in}{1.270751in}}{\pgfqpoint{11.738056in}{1.271958in}}%
\pgfpathquadraticcurveto{\pgfqpoint{11.742469in}{1.273164in}}{\pgfqpoint{11.746864in}{1.275596in}}%
\pgfpathlineto{\pgfqpoint{11.746864in}{1.266014in}}%
\pgfpathquadraticcurveto{\pgfqpoint{11.742523in}{1.263996in}}{\pgfqpoint{11.737876in}{1.262988in}}%
\pgfpathquadraticcurveto{\pgfqpoint{11.733247in}{1.261961in}}{\pgfqpoint{11.728005in}{1.261961in}}%
\pgfpathquadraticcurveto{\pgfqpoint{11.713758in}{1.261961in}}{\pgfqpoint{11.705364in}{1.270913in}}%
\pgfpathquadraticcurveto{\pgfqpoint{11.696988in}{1.279865in}}{\pgfqpoint{11.696988in}{1.295067in}}%
\pgfpathquadraticcurveto{\pgfqpoint{11.696988in}{1.310486in}}{\pgfqpoint{11.705454in}{1.319312in}}%
\pgfpathquadraticcurveto{\pgfqpoint{11.713938in}{1.328156in}}{\pgfqpoint{11.728690in}{1.328156in}}%
\pgfpathquadraticcurveto{\pgfqpoint{11.733463in}{1.328156in}}{\pgfqpoint{11.738020in}{1.327165in}}%
\pgfpathquadraticcurveto{\pgfqpoint{11.742577in}{1.326192in}}{\pgfqpoint{11.746864in}{1.324229in}}%
\pgfpathlineto{\pgfqpoint{11.746864in}{1.324229in}}%
\pgfpathclose%
\pgfpathmoveto{\pgfqpoint{11.789301in}{1.319384in}}%
\pgfpathquadraticcurveto{\pgfqpoint{11.780979in}{1.319384in}}{\pgfqpoint{11.776134in}{1.312881in}}%
\pgfpathquadraticcurveto{\pgfqpoint{11.771289in}{1.306379in}}{\pgfqpoint{11.771289in}{1.295067in}}%
\pgfpathquadraticcurveto{\pgfqpoint{11.771289in}{1.283756in}}{\pgfqpoint{11.776098in}{1.277253in}}%
\pgfpathquadraticcurveto{\pgfqpoint{11.780925in}{1.270751in}}{\pgfqpoint{11.789301in}{1.270751in}}%
\pgfpathquadraticcurveto{\pgfqpoint{11.797586in}{1.270751in}}{\pgfqpoint{11.802414in}{1.277271in}}%
\pgfpathquadraticcurveto{\pgfqpoint{11.807259in}{1.283810in}}{\pgfqpoint{11.807259in}{1.295067in}}%
\pgfpathquadraticcurveto{\pgfqpoint{11.807259in}{1.306271in}}{\pgfqpoint{11.802414in}{1.312827in}}%
\pgfpathquadraticcurveto{\pgfqpoint{11.797586in}{1.319384in}}{\pgfqpoint{11.789301in}{1.319384in}}%
\pgfpathlineto{\pgfqpoint{11.789301in}{1.319384in}}%
\pgfpathclose%
\pgfpathmoveto{\pgfqpoint{11.789301in}{1.328156in}}%
\pgfpathquadraticcurveto{\pgfqpoint{11.802810in}{1.328156in}}{\pgfqpoint{11.810519in}{1.319366in}}%
\pgfpathquadraticcurveto{\pgfqpoint{11.818246in}{1.310594in}}{\pgfqpoint{11.818246in}{1.295067in}}%
\pgfpathquadraticcurveto{\pgfqpoint{11.818246in}{1.279595in}}{\pgfqpoint{11.810519in}{1.270769in}}%
\pgfpathquadraticcurveto{\pgfqpoint{11.802810in}{1.261961in}}{\pgfqpoint{11.789301in}{1.261961in}}%
\pgfpathquadraticcurveto{\pgfqpoint{11.775738in}{1.261961in}}{\pgfqpoint{11.768046in}{1.270769in}}%
\pgfpathquadraticcurveto{\pgfqpoint{11.760373in}{1.279595in}}{\pgfqpoint{11.760373in}{1.295067in}}%
\pgfpathquadraticcurveto{\pgfqpoint{11.760373in}{1.310594in}}{\pgfqpoint{11.768046in}{1.319366in}}%
\pgfpathquadraticcurveto{\pgfqpoint{11.775738in}{1.328156in}}{\pgfqpoint{11.789301in}{1.328156in}}%
\pgfpathlineto{\pgfqpoint{11.789301in}{1.328156in}}%
\pgfpathclose%
\pgfpathmoveto{\pgfqpoint{11.884495in}{1.314538in}}%
\pgfpathquadraticcurveto{\pgfqpoint{11.888386in}{1.321527in}}{\pgfqpoint{11.893789in}{1.324841in}}%
\pgfpathquadraticcurveto{\pgfqpoint{11.899193in}{1.328156in}}{\pgfqpoint{11.906506in}{1.328156in}}%
\pgfpathquadraticcurveto{\pgfqpoint{11.916359in}{1.328156in}}{\pgfqpoint{11.921708in}{1.321257in}}%
\pgfpathquadraticcurveto{\pgfqpoint{11.927058in}{1.314376in}}{\pgfqpoint{11.927058in}{1.301660in}}%
\pgfpathlineto{\pgfqpoint{11.927058in}{1.263600in}}%
\pgfpathlineto{\pgfqpoint{11.916647in}{1.263600in}}%
\pgfpathlineto{\pgfqpoint{11.916647in}{1.301317in}}%
\pgfpathquadraticcurveto{\pgfqpoint{11.916647in}{1.310378in}}{\pgfqpoint{11.913423in}{1.314755in}}%
\pgfpathquadraticcurveto{\pgfqpoint{11.910216in}{1.319149in}}{\pgfqpoint{11.903642in}{1.319149in}}%
\pgfpathquadraticcurveto{\pgfqpoint{11.895590in}{1.319149in}}{\pgfqpoint{11.890907in}{1.313800in}}%
\pgfpathquadraticcurveto{\pgfqpoint{11.886242in}{1.308468in}}{\pgfqpoint{11.886242in}{1.299228in}}%
\pgfpathlineto{\pgfqpoint{11.886242in}{1.263600in}}%
\pgfpathlineto{\pgfqpoint{11.875831in}{1.263600in}}%
\pgfpathlineto{\pgfqpoint{11.875831in}{1.301317in}}%
\pgfpathquadraticcurveto{\pgfqpoint{11.875831in}{1.310432in}}{\pgfqpoint{11.872625in}{1.314791in}}%
\pgfpathquadraticcurveto{\pgfqpoint{11.869419in}{1.319149in}}{\pgfqpoint{11.862718in}{1.319149in}}%
\pgfpathquadraticcurveto{\pgfqpoint{11.854775in}{1.319149in}}{\pgfqpoint{11.850092in}{1.313782in}}%
\pgfpathquadraticcurveto{\pgfqpoint{11.845427in}{1.308414in}}{\pgfqpoint{11.845427in}{1.299228in}}%
\pgfpathlineto{\pgfqpoint{11.845427in}{1.263600in}}%
\pgfpathlineto{\pgfqpoint{11.835016in}{1.263600in}}%
\pgfpathlineto{\pgfqpoint{11.835016in}{1.326643in}}%
\pgfpathlineto{\pgfqpoint{11.845427in}{1.326643in}}%
\pgfpathlineto{\pgfqpoint{11.845427in}{1.316844in}}%
\pgfpathquadraticcurveto{\pgfqpoint{11.848975in}{1.322644in}}{\pgfqpoint{11.853928in}{1.325400in}}%
\pgfpathquadraticcurveto{\pgfqpoint{11.858882in}{1.328156in}}{\pgfqpoint{11.865690in}{1.328156in}}%
\pgfpathquadraticcurveto{\pgfqpoint{11.872571in}{1.328156in}}{\pgfqpoint{11.877380in}{1.324661in}}%
\pgfpathquadraticcurveto{\pgfqpoint{11.882189in}{1.321185in}}{\pgfqpoint{11.884495in}{1.314538in}}%
\pgfpathlineto{\pgfqpoint{11.884495in}{1.314538in}}%
\pgfpathclose%
\pgfpathmoveto{\pgfqpoint{11.957714in}{1.273056in}}%
\pgfpathlineto{\pgfqpoint{11.957714in}{1.239626in}}%
\pgfpathlineto{\pgfqpoint{11.947303in}{1.239626in}}%
\pgfpathlineto{\pgfqpoint{11.947303in}{1.326643in}}%
\pgfpathlineto{\pgfqpoint{11.957714in}{1.326643in}}%
\pgfpathlineto{\pgfqpoint{11.957714in}{1.317078in}}%
\pgfpathquadraticcurveto{\pgfqpoint{11.960993in}{1.322698in}}{\pgfqpoint{11.965964in}{1.325418in}}%
\pgfpathquadraticcurveto{\pgfqpoint{11.970953in}{1.328156in}}{\pgfqpoint{11.977870in}{1.328156in}}%
\pgfpathquadraticcurveto{\pgfqpoint{11.989362in}{1.328156in}}{\pgfqpoint{11.996531in}{1.319041in}}%
\pgfpathquadraticcurveto{\pgfqpoint{12.003717in}{1.309927in}}{\pgfqpoint{12.003717in}{1.295067in}}%
\pgfpathquadraticcurveto{\pgfqpoint{12.003717in}{1.280207in}}{\pgfqpoint{11.996531in}{1.271075in}}%
\pgfpathquadraticcurveto{\pgfqpoint{11.989362in}{1.261961in}}{\pgfqpoint{11.977870in}{1.261961in}}%
\pgfpathquadraticcurveto{\pgfqpoint{11.970953in}{1.261961in}}{\pgfqpoint{11.965964in}{1.264699in}}%
\pgfpathquadraticcurveto{\pgfqpoint{11.960993in}{1.267437in}}{\pgfqpoint{11.957714in}{1.273056in}}%
\pgfpathlineto{\pgfqpoint{11.957714in}{1.273056in}}%
\pgfpathclose%
\pgfpathmoveto{\pgfqpoint{11.992964in}{1.295067in}}%
\pgfpathquadraticcurveto{\pgfqpoint{11.992964in}{1.306487in}}{\pgfqpoint{11.988263in}{1.312989in}}%
\pgfpathquadraticcurveto{\pgfqpoint{11.983562in}{1.319492in}}{\pgfqpoint{11.975348in}{1.319492in}}%
\pgfpathquadraticcurveto{\pgfqpoint{11.967117in}{1.319492in}}{\pgfqpoint{11.962416in}{1.312989in}}%
\pgfpathquadraticcurveto{\pgfqpoint{11.957714in}{1.306487in}}{\pgfqpoint{11.957714in}{1.295067in}}%
\pgfpathquadraticcurveto{\pgfqpoint{11.957714in}{1.283648in}}{\pgfqpoint{11.962416in}{1.277145in}}%
\pgfpathquadraticcurveto{\pgfqpoint{11.967117in}{1.270643in}}{\pgfqpoint{11.975348in}{1.270643in}}%
\pgfpathquadraticcurveto{\pgfqpoint{11.983562in}{1.270643in}}{\pgfqpoint{11.988263in}{1.277145in}}%
\pgfpathquadraticcurveto{\pgfqpoint{11.992964in}{1.283648in}}{\pgfqpoint{11.992964in}{1.295067in}}%
\pgfpathlineto{\pgfqpoint{11.992964in}{1.295067in}}%
\pgfpathclose%
\pgfpathmoveto{\pgfqpoint{12.049540in}{1.295283in}}%
\pgfpathquadraticcurveto{\pgfqpoint{12.036986in}{1.295283in}}{\pgfqpoint{12.032141in}{1.292419in}}%
\pgfpathquadraticcurveto{\pgfqpoint{12.027295in}{1.289556in}}{\pgfqpoint{12.027295in}{1.282621in}}%
\pgfpathquadraticcurveto{\pgfqpoint{12.027295in}{1.277109in}}{\pgfqpoint{12.030934in}{1.273867in}}%
\pgfpathquadraticcurveto{\pgfqpoint{12.034572in}{1.270643in}}{\pgfqpoint{12.040804in}{1.270643in}}%
\pgfpathquadraticcurveto{\pgfqpoint{12.049432in}{1.270643in}}{\pgfqpoint{12.054638in}{1.276749in}}%
\pgfpathquadraticcurveto{\pgfqpoint{12.059843in}{1.282855in}}{\pgfqpoint{12.059843in}{1.292978in}}%
\pgfpathlineto{\pgfqpoint{12.059843in}{1.295283in}}%
\pgfpathlineto{\pgfqpoint{12.049540in}{1.295283in}}%
\pgfpathlineto{\pgfqpoint{12.049540in}{1.295283in}}%
\pgfpathclose%
\pgfpathmoveto{\pgfqpoint{12.070200in}{1.299570in}}%
\pgfpathlineto{\pgfqpoint{12.070200in}{1.263600in}}%
\pgfpathlineto{\pgfqpoint{12.059843in}{1.263600in}}%
\pgfpathlineto{\pgfqpoint{12.059843in}{1.273164in}}%
\pgfpathquadraticcurveto{\pgfqpoint{12.056295in}{1.267437in}}{\pgfqpoint{12.050999in}{1.264699in}}%
\pgfpathquadraticcurveto{\pgfqpoint{12.045704in}{1.261961in}}{\pgfqpoint{12.038049in}{1.261961in}}%
\pgfpathquadraticcurveto{\pgfqpoint{12.028376in}{1.261961in}}{\pgfqpoint{12.022648in}{1.267401in}}%
\pgfpathquadraticcurveto{\pgfqpoint{12.016938in}{1.272840in}}{\pgfqpoint{12.016938in}{1.281954in}}%
\pgfpathquadraticcurveto{\pgfqpoint{12.016938in}{1.292582in}}{\pgfqpoint{12.024053in}{1.297985in}}%
\pgfpathquadraticcurveto{\pgfqpoint{12.031186in}{1.303389in}}{\pgfqpoint{12.045308in}{1.303389in}}%
\pgfpathlineto{\pgfqpoint{12.059843in}{1.303389in}}%
\pgfpathlineto{\pgfqpoint{12.059843in}{1.304416in}}%
\pgfpathquadraticcurveto{\pgfqpoint{12.059843in}{1.311566in}}{\pgfqpoint{12.055142in}{1.315475in}}%
\pgfpathquadraticcurveto{\pgfqpoint{12.050441in}{1.319384in}}{\pgfqpoint{12.041939in}{1.319384in}}%
\pgfpathquadraticcurveto{\pgfqpoint{12.036536in}{1.319384in}}{\pgfqpoint{12.031402in}{1.318087in}}%
\pgfpathquadraticcurveto{\pgfqpoint{12.026287in}{1.316790in}}{\pgfqpoint{12.021567in}{1.314196in}}%
\pgfpathlineto{\pgfqpoint{12.021567in}{1.323779in}}%
\pgfpathquadraticcurveto{\pgfqpoint{12.027241in}{1.325976in}}{\pgfqpoint{12.032591in}{1.327057in}}%
\pgfpathquadraticcurveto{\pgfqpoint{12.037941in}{1.328156in}}{\pgfqpoint{12.043002in}{1.328156in}}%
\pgfpathquadraticcurveto{\pgfqpoint{12.056691in}{1.328156in}}{\pgfqpoint{12.063446in}{1.321059in}}%
\pgfpathquadraticcurveto{\pgfqpoint{12.070200in}{1.313980in}}{\pgfqpoint{12.070200in}{1.299570in}}%
\pgfpathlineto{\pgfqpoint{12.070200in}{1.299570in}}%
\pgfpathclose%
\pgfpathmoveto{\pgfqpoint{12.128055in}{1.316970in}}%
\pgfpathquadraticcurveto{\pgfqpoint{12.126308in}{1.317979in}}{\pgfqpoint{12.124255in}{1.318447in}}%
\pgfpathquadraticcurveto{\pgfqpoint{12.122201in}{1.318933in}}{\pgfqpoint{12.119734in}{1.318933in}}%
\pgfpathquadraticcurveto{\pgfqpoint{12.110944in}{1.318933in}}{\pgfqpoint{12.106243in}{1.313223in}}%
\pgfpathquadraticcurveto{\pgfqpoint{12.101541in}{1.307514in}}{\pgfqpoint{12.101541in}{1.296814in}}%
\pgfpathlineto{\pgfqpoint{12.101541in}{1.263600in}}%
\pgfpathlineto{\pgfqpoint{12.091130in}{1.263600in}}%
\pgfpathlineto{\pgfqpoint{12.091130in}{1.326643in}}%
\pgfpathlineto{\pgfqpoint{12.101541in}{1.326643in}}%
\pgfpathlineto{\pgfqpoint{12.101541in}{1.316844in}}%
\pgfpathquadraticcurveto{\pgfqpoint{12.104820in}{1.322590in}}{\pgfqpoint{12.110043in}{1.325364in}}%
\pgfpathquadraticcurveto{\pgfqpoint{12.115285in}{1.328156in}}{\pgfqpoint{12.122778in}{1.328156in}}%
\pgfpathquadraticcurveto{\pgfqpoint{12.123840in}{1.328156in}}{\pgfqpoint{12.125137in}{1.328011in}}%
\pgfpathquadraticcurveto{\pgfqpoint{12.126434in}{1.327885in}}{\pgfqpoint{12.128001in}{1.327597in}}%
\pgfpathlineto{\pgfqpoint{12.128055in}{1.316970in}}%
\pgfpathlineto{\pgfqpoint{12.128055in}{1.316970in}}%
\pgfpathclose%
\pgfpathmoveto{\pgfqpoint{12.138917in}{1.326643in}}%
\pgfpathlineto{\pgfqpoint{12.149274in}{1.326643in}}%
\pgfpathlineto{\pgfqpoint{12.149274in}{1.263600in}}%
\pgfpathlineto{\pgfqpoint{12.138917in}{1.263600in}}%
\pgfpathlineto{\pgfqpoint{12.138917in}{1.326643in}}%
\pgfpathlineto{\pgfqpoint{12.138917in}{1.326643in}}%
\pgfpathclose%
\pgfpathmoveto{\pgfqpoint{12.138917in}{1.351193in}}%
\pgfpathlineto{\pgfqpoint{12.149274in}{1.351193in}}%
\pgfpathlineto{\pgfqpoint{12.149274in}{1.338062in}}%
\pgfpathlineto{\pgfqpoint{12.138917in}{1.338062in}}%
\pgfpathlineto{\pgfqpoint{12.138917in}{1.351193in}}%
\pgfpathlineto{\pgfqpoint{12.138917in}{1.351193in}}%
\pgfpathclose%
\pgfpathmoveto{\pgfqpoint{12.211127in}{1.324787in}}%
\pgfpathlineto{\pgfqpoint{12.211127in}{1.314989in}}%
\pgfpathquadraticcurveto{\pgfqpoint{12.206750in}{1.317240in}}{\pgfqpoint{12.202013in}{1.318357in}}%
\pgfpathquadraticcurveto{\pgfqpoint{12.197294in}{1.319492in}}{\pgfqpoint{12.192215in}{1.319492in}}%
\pgfpathquadraticcurveto{\pgfqpoint{12.184505in}{1.319492in}}{\pgfqpoint{12.180651in}{1.317132in}}%
\pgfpathquadraticcurveto{\pgfqpoint{12.176796in}{1.314773in}}{\pgfqpoint{12.176796in}{1.310035in}}%
\pgfpathquadraticcurveto{\pgfqpoint{12.176796in}{1.306433in}}{\pgfqpoint{12.179552in}{1.304380in}}%
\pgfpathquadraticcurveto{\pgfqpoint{12.182308in}{1.302326in}}{\pgfqpoint{12.190648in}{1.300471in}}%
\pgfpathlineto{\pgfqpoint{12.194196in}{1.299678in}}%
\pgfpathquadraticcurveto{\pgfqpoint{12.205219in}{1.297319in}}{\pgfqpoint{12.209867in}{1.293014in}}%
\pgfpathquadraticcurveto{\pgfqpoint{12.214514in}{1.288709in}}{\pgfqpoint{12.214514in}{1.281000in}}%
\pgfpathquadraticcurveto{\pgfqpoint{12.214514in}{1.272210in}}{\pgfqpoint{12.207561in}{1.267076in}}%
\pgfpathquadraticcurveto{\pgfqpoint{12.200608in}{1.261961in}}{\pgfqpoint{12.188450in}{1.261961in}}%
\pgfpathquadraticcurveto{\pgfqpoint{12.183389in}{1.261961in}}{\pgfqpoint{12.177895in}{1.262952in}}%
\pgfpathquadraticcurveto{\pgfqpoint{12.172401in}{1.263942in}}{\pgfqpoint{12.166331in}{1.265906in}}%
\pgfpathlineto{\pgfqpoint{12.166331in}{1.276605in}}%
\pgfpathquadraticcurveto{\pgfqpoint{12.172077in}{1.273615in}}{\pgfqpoint{12.177643in}{1.272120in}}%
\pgfpathquadraticcurveto{\pgfqpoint{12.183209in}{1.270643in}}{\pgfqpoint{12.188684in}{1.270643in}}%
\pgfpathquadraticcurveto{\pgfqpoint{12.195997in}{1.270643in}}{\pgfqpoint{12.199924in}{1.273146in}}%
\pgfpathquadraticcurveto{\pgfqpoint{12.203868in}{1.275650in}}{\pgfqpoint{12.203868in}{1.280207in}}%
\pgfpathquadraticcurveto{\pgfqpoint{12.203868in}{1.284422in}}{\pgfqpoint{12.201023in}{1.286674in}}%
\pgfpathquadraticcurveto{\pgfqpoint{12.198195in}{1.288925in}}{\pgfqpoint{12.188558in}{1.291014in}}%
\pgfpathlineto{\pgfqpoint{12.184956in}{1.291861in}}%
\pgfpathquadraticcurveto{\pgfqpoint{12.175337in}{1.293878in}}{\pgfqpoint{12.171050in}{1.298075in}}%
\pgfpathquadraticcurveto{\pgfqpoint{12.166781in}{1.302272in}}{\pgfqpoint{12.166781in}{1.309585in}}%
\pgfpathquadraticcurveto{\pgfqpoint{12.166781in}{1.318483in}}{\pgfqpoint{12.173086in}{1.323310in}}%
\pgfpathquadraticcurveto{\pgfqpoint{12.179390in}{1.328156in}}{\pgfqpoint{12.190990in}{1.328156in}}%
\pgfpathquadraticcurveto{\pgfqpoint{12.196718in}{1.328156in}}{\pgfqpoint{12.201779in}{1.327309in}}%
\pgfpathquadraticcurveto{\pgfqpoint{12.206859in}{1.326480in}}{\pgfqpoint{12.211127in}{1.324787in}}%
\pgfpathlineto{\pgfqpoint{12.211127in}{1.324787in}}%
\pgfpathclose%
\pgfpathmoveto{\pgfqpoint{12.255419in}{1.319384in}}%
\pgfpathquadraticcurveto{\pgfqpoint{12.247098in}{1.319384in}}{\pgfqpoint{12.242252in}{1.312881in}}%
\pgfpathquadraticcurveto{\pgfqpoint{12.237407in}{1.306379in}}{\pgfqpoint{12.237407in}{1.295067in}}%
\pgfpathquadraticcurveto{\pgfqpoint{12.237407in}{1.283756in}}{\pgfqpoint{12.242216in}{1.277253in}}%
\pgfpathquadraticcurveto{\pgfqpoint{12.247044in}{1.270751in}}{\pgfqpoint{12.255419in}{1.270751in}}%
\pgfpathquadraticcurveto{\pgfqpoint{12.263705in}{1.270751in}}{\pgfqpoint{12.268532in}{1.277271in}}%
\pgfpathquadraticcurveto{\pgfqpoint{12.273377in}{1.283810in}}{\pgfqpoint{12.273377in}{1.295067in}}%
\pgfpathquadraticcurveto{\pgfqpoint{12.273377in}{1.306271in}}{\pgfqpoint{12.268532in}{1.312827in}}%
\pgfpathquadraticcurveto{\pgfqpoint{12.263705in}{1.319384in}}{\pgfqpoint{12.255419in}{1.319384in}}%
\pgfpathlineto{\pgfqpoint{12.255419in}{1.319384in}}%
\pgfpathclose%
\pgfpathmoveto{\pgfqpoint{12.255419in}{1.328156in}}%
\pgfpathquadraticcurveto{\pgfqpoint{12.268928in}{1.328156in}}{\pgfqpoint{12.276638in}{1.319366in}}%
\pgfpathquadraticcurveto{\pgfqpoint{12.284365in}{1.310594in}}{\pgfqpoint{12.284365in}{1.295067in}}%
\pgfpathquadraticcurveto{\pgfqpoint{12.284365in}{1.279595in}}{\pgfqpoint{12.276638in}{1.270769in}}%
\pgfpathquadraticcurveto{\pgfqpoint{12.268928in}{1.261961in}}{\pgfqpoint{12.255419in}{1.261961in}}%
\pgfpathquadraticcurveto{\pgfqpoint{12.241856in}{1.261961in}}{\pgfqpoint{12.234165in}{1.270769in}}%
\pgfpathquadraticcurveto{\pgfqpoint{12.226492in}{1.279595in}}{\pgfqpoint{12.226492in}{1.295067in}}%
\pgfpathquadraticcurveto{\pgfqpoint{12.226492in}{1.310594in}}{\pgfqpoint{12.234165in}{1.319366in}}%
\pgfpathquadraticcurveto{\pgfqpoint{12.241856in}{1.328156in}}{\pgfqpoint{12.255419in}{1.328156in}}%
\pgfpathlineto{\pgfqpoint{12.255419in}{1.328156in}}%
\pgfpathclose%
\pgfpathmoveto{\pgfqpoint{12.353946in}{1.301660in}}%
\pgfpathlineto{\pgfqpoint{12.353946in}{1.263600in}}%
\pgfpathlineto{\pgfqpoint{12.343589in}{1.263600in}}%
\pgfpathlineto{\pgfqpoint{12.343589in}{1.301317in}}%
\pgfpathquadraticcurveto{\pgfqpoint{12.343589in}{1.310269in}}{\pgfqpoint{12.340094in}{1.314700in}}%
\pgfpathquadraticcurveto{\pgfqpoint{12.336600in}{1.319149in}}{\pgfqpoint{12.329629in}{1.319149in}}%
\pgfpathquadraticcurveto{\pgfqpoint{12.321236in}{1.319149in}}{\pgfqpoint{12.316390in}{1.313800in}}%
\pgfpathquadraticcurveto{\pgfqpoint{12.311545in}{1.308468in}}{\pgfqpoint{12.311545in}{1.299228in}}%
\pgfpathlineto{\pgfqpoint{12.311545in}{1.263600in}}%
\pgfpathlineto{\pgfqpoint{12.301134in}{1.263600in}}%
\pgfpathlineto{\pgfqpoint{12.301134in}{1.326643in}}%
\pgfpathlineto{\pgfqpoint{12.311545in}{1.326643in}}%
\pgfpathlineto{\pgfqpoint{12.311545in}{1.316844in}}%
\pgfpathquadraticcurveto{\pgfqpoint{12.315274in}{1.322536in}}{\pgfqpoint{12.320299in}{1.325346in}}%
\pgfpathquadraticcurveto{\pgfqpoint{12.325342in}{1.328156in}}{\pgfqpoint{12.331935in}{1.328156in}}%
\pgfpathquadraticcurveto{\pgfqpoint{12.342796in}{1.328156in}}{\pgfqpoint{12.348362in}{1.321437in}}%
\pgfpathquadraticcurveto{\pgfqpoint{12.353946in}{1.314718in}}{\pgfqpoint{12.353946in}{1.301660in}}%
\pgfpathlineto{\pgfqpoint{12.353946in}{1.301660in}}%
\pgfpathclose%
\pgfpathmoveto{\pgfqpoint{12.414773in}{1.324787in}}%
\pgfpathlineto{\pgfqpoint{12.414773in}{1.314989in}}%
\pgfpathquadraticcurveto{\pgfqpoint{12.410396in}{1.317240in}}{\pgfqpoint{12.405659in}{1.318357in}}%
\pgfpathquadraticcurveto{\pgfqpoint{12.400939in}{1.319492in}}{\pgfqpoint{12.395860in}{1.319492in}}%
\pgfpathquadraticcurveto{\pgfqpoint{12.388151in}{1.319492in}}{\pgfqpoint{12.384296in}{1.317132in}}%
\pgfpathquadraticcurveto{\pgfqpoint{12.380442in}{1.314773in}}{\pgfqpoint{12.380442in}{1.310035in}}%
\pgfpathquadraticcurveto{\pgfqpoint{12.380442in}{1.306433in}}{\pgfqpoint{12.383197in}{1.304380in}}%
\pgfpathquadraticcurveto{\pgfqpoint{12.385953in}{1.302326in}}{\pgfqpoint{12.394293in}{1.300471in}}%
\pgfpathlineto{\pgfqpoint{12.397841in}{1.299678in}}%
\pgfpathquadraticcurveto{\pgfqpoint{12.408865in}{1.297319in}}{\pgfqpoint{12.413512in}{1.293014in}}%
\pgfpathquadraticcurveto{\pgfqpoint{12.418159in}{1.288709in}}{\pgfqpoint{12.418159in}{1.281000in}}%
\pgfpathquadraticcurveto{\pgfqpoint{12.418159in}{1.272210in}}{\pgfqpoint{12.411206in}{1.267076in}}%
\pgfpathquadraticcurveto{\pgfqpoint{12.404254in}{1.261961in}}{\pgfqpoint{12.392095in}{1.261961in}}%
\pgfpathquadraticcurveto{\pgfqpoint{12.387034in}{1.261961in}}{\pgfqpoint{12.381540in}{1.262952in}}%
\pgfpathquadraticcurveto{\pgfqpoint{12.376047in}{1.263942in}}{\pgfqpoint{12.369977in}{1.265906in}}%
\pgfpathlineto{\pgfqpoint{12.369977in}{1.276605in}}%
\pgfpathquadraticcurveto{\pgfqpoint{12.375722in}{1.273615in}}{\pgfqpoint{12.381288in}{1.272120in}}%
\pgfpathquadraticcurveto{\pgfqpoint{12.386854in}{1.270643in}}{\pgfqpoint{12.392330in}{1.270643in}}%
\pgfpathquadraticcurveto{\pgfqpoint{12.399643in}{1.270643in}}{\pgfqpoint{12.403569in}{1.273146in}}%
\pgfpathquadraticcurveto{\pgfqpoint{12.407514in}{1.275650in}}{\pgfqpoint{12.407514in}{1.280207in}}%
\pgfpathquadraticcurveto{\pgfqpoint{12.407514in}{1.284422in}}{\pgfqpoint{12.404668in}{1.286674in}}%
\pgfpathquadraticcurveto{\pgfqpoint{12.401840in}{1.288925in}}{\pgfqpoint{12.392204in}{1.291014in}}%
\pgfpathlineto{\pgfqpoint{12.388601in}{1.291861in}}%
\pgfpathquadraticcurveto{\pgfqpoint{12.378983in}{1.293878in}}{\pgfqpoint{12.374696in}{1.298075in}}%
\pgfpathquadraticcurveto{\pgfqpoint{12.370427in}{1.302272in}}{\pgfqpoint{12.370427in}{1.309585in}}%
\pgfpathquadraticcurveto{\pgfqpoint{12.370427in}{1.318483in}}{\pgfqpoint{12.376731in}{1.323310in}}%
\pgfpathquadraticcurveto{\pgfqpoint{12.383035in}{1.328156in}}{\pgfqpoint{12.394635in}{1.328156in}}%
\pgfpathquadraticcurveto{\pgfqpoint{12.400363in}{1.328156in}}{\pgfqpoint{12.405424in}{1.327309in}}%
\pgfpathquadraticcurveto{\pgfqpoint{12.410504in}{1.326480in}}{\pgfqpoint{12.414773in}{1.324787in}}%
\pgfpathlineto{\pgfqpoint{12.414773in}{1.324787in}}%
\pgfpathclose%
\pgfpathmoveto{\pgfqpoint{12.437288in}{1.277902in}}%
\pgfpathlineto{\pgfqpoint{12.449158in}{1.277902in}}%
\pgfpathlineto{\pgfqpoint{12.449158in}{1.268211in}}%
\pgfpathlineto{\pgfqpoint{12.439936in}{1.250199in}}%
\pgfpathlineto{\pgfqpoint{12.432677in}{1.250199in}}%
\pgfpathlineto{\pgfqpoint{12.437288in}{1.268211in}}%
\pgfpathlineto{\pgfqpoint{12.437288in}{1.277902in}}%
\pgfpathlineto{\pgfqpoint{12.437288in}{1.277902in}}%
\pgfpathclose%
\pgfusepath{fill}%
\end{pgfscope}%
\begin{pgfscope}%
\pgfsetbuttcap%
\pgfsetroundjoin%
\definecolor{currentfill}{rgb}{0.309804,0.372549,0.419608}%
\pgfsetfillcolor{currentfill}%
\pgfsetlinewidth{0.000000pt}%
\definecolor{currentstroke}{rgb}{0.309804,0.372549,0.419608}%
\pgfsetstrokecolor{currentstroke}%
\pgfsetdash{}{0pt}%
\pgfpathmoveto{\pgfqpoint{12.557176in}{1.326643in}}%
\pgfpathlineto{\pgfqpoint{12.567533in}{1.326643in}}%
\pgfpathlineto{\pgfqpoint{12.567533in}{1.263600in}}%
\pgfpathlineto{\pgfqpoint{12.557176in}{1.263600in}}%
\pgfpathlineto{\pgfqpoint{12.557176in}{1.326643in}}%
\pgfpathlineto{\pgfqpoint{12.557176in}{1.326643in}}%
\pgfpathclose%
\pgfpathmoveto{\pgfqpoint{12.557176in}{1.351193in}}%
\pgfpathlineto{\pgfqpoint{12.567533in}{1.351193in}}%
\pgfpathlineto{\pgfqpoint{12.567533in}{1.338062in}}%
\pgfpathlineto{\pgfqpoint{12.557176in}{1.338062in}}%
\pgfpathlineto{\pgfqpoint{12.557176in}{1.351193in}}%
\pgfpathlineto{\pgfqpoint{12.557176in}{1.351193in}}%
\pgfpathclose%
\pgfpathmoveto{\pgfqpoint{12.641617in}{1.301660in}}%
\pgfpathlineto{\pgfqpoint{12.641617in}{1.263600in}}%
\pgfpathlineto{\pgfqpoint{12.631260in}{1.263600in}}%
\pgfpathlineto{\pgfqpoint{12.631260in}{1.301317in}}%
\pgfpathquadraticcurveto{\pgfqpoint{12.631260in}{1.310269in}}{\pgfqpoint{12.627766in}{1.314700in}}%
\pgfpathquadraticcurveto{\pgfqpoint{12.624272in}{1.319149in}}{\pgfqpoint{12.617301in}{1.319149in}}%
\pgfpathquadraticcurveto{\pgfqpoint{12.608907in}{1.319149in}}{\pgfqpoint{12.604062in}{1.313800in}}%
\pgfpathquadraticcurveto{\pgfqpoint{12.599217in}{1.308468in}}{\pgfqpoint{12.599217in}{1.299228in}}%
\pgfpathlineto{\pgfqpoint{12.599217in}{1.263600in}}%
\pgfpathlineto{\pgfqpoint{12.588806in}{1.263600in}}%
\pgfpathlineto{\pgfqpoint{12.588806in}{1.326643in}}%
\pgfpathlineto{\pgfqpoint{12.599217in}{1.326643in}}%
\pgfpathlineto{\pgfqpoint{12.599217in}{1.316844in}}%
\pgfpathquadraticcurveto{\pgfqpoint{12.602945in}{1.322536in}}{\pgfqpoint{12.607971in}{1.325346in}}%
\pgfpathquadraticcurveto{\pgfqpoint{12.613014in}{1.328156in}}{\pgfqpoint{12.619607in}{1.328156in}}%
\pgfpathquadraticcurveto{\pgfqpoint{12.630468in}{1.328156in}}{\pgfqpoint{12.636034in}{1.321437in}}%
\pgfpathquadraticcurveto{\pgfqpoint{12.641617in}{1.314718in}}{\pgfqpoint{12.641617in}{1.301660in}}%
\pgfpathlineto{\pgfqpoint{12.641617in}{1.301660in}}%
\pgfpathclose%
\pgfpathmoveto{\pgfqpoint{12.707632in}{1.324229in}}%
\pgfpathlineto{\pgfqpoint{12.707632in}{1.314538in}}%
\pgfpathquadraticcurveto{\pgfqpoint{12.703237in}{1.316970in}}{\pgfqpoint{12.698824in}{1.318177in}}%
\pgfpathquadraticcurveto{\pgfqpoint{12.694411in}{1.319384in}}{\pgfqpoint{12.689908in}{1.319384in}}%
\pgfpathquadraticcurveto{\pgfqpoint{12.679821in}{1.319384in}}{\pgfqpoint{12.674237in}{1.312989in}}%
\pgfpathquadraticcurveto{\pgfqpoint{12.668672in}{1.306613in}}{\pgfqpoint{12.668672in}{1.295067in}}%
\pgfpathquadraticcurveto{\pgfqpoint{12.668672in}{1.283521in}}{\pgfqpoint{12.674237in}{1.277127in}}%
\pgfpathquadraticcurveto{\pgfqpoint{12.679821in}{1.270751in}}{\pgfqpoint{12.689908in}{1.270751in}}%
\pgfpathquadraticcurveto{\pgfqpoint{12.694411in}{1.270751in}}{\pgfqpoint{12.698824in}{1.271958in}}%
\pgfpathquadraticcurveto{\pgfqpoint{12.703237in}{1.273164in}}{\pgfqpoint{12.707632in}{1.275596in}}%
\pgfpathlineto{\pgfqpoint{12.707632in}{1.266014in}}%
\pgfpathquadraticcurveto{\pgfqpoint{12.703291in}{1.263996in}}{\pgfqpoint{12.698644in}{1.262988in}}%
\pgfpathquadraticcurveto{\pgfqpoint{12.694015in}{1.261961in}}{\pgfqpoint{12.688773in}{1.261961in}}%
\pgfpathquadraticcurveto{\pgfqpoint{12.674526in}{1.261961in}}{\pgfqpoint{12.666132in}{1.270913in}}%
\pgfpathquadraticcurveto{\pgfqpoint{12.657756in}{1.279865in}}{\pgfqpoint{12.657756in}{1.295067in}}%
\pgfpathquadraticcurveto{\pgfqpoint{12.657756in}{1.310486in}}{\pgfqpoint{12.666222in}{1.319312in}}%
\pgfpathquadraticcurveto{\pgfqpoint{12.674706in}{1.328156in}}{\pgfqpoint{12.689458in}{1.328156in}}%
\pgfpathquadraticcurveto{\pgfqpoint{12.694231in}{1.328156in}}{\pgfqpoint{12.698788in}{1.327165in}}%
\pgfpathquadraticcurveto{\pgfqpoint{12.703345in}{1.326192in}}{\pgfqpoint{12.707632in}{1.324229in}}%
\pgfpathlineto{\pgfqpoint{12.707632in}{1.324229in}}%
\pgfpathclose%
\pgfpathmoveto{\pgfqpoint{12.725644in}{1.351193in}}%
\pgfpathlineto{\pgfqpoint{12.736001in}{1.351193in}}%
\pgfpathlineto{\pgfqpoint{12.736001in}{1.263600in}}%
\pgfpathlineto{\pgfqpoint{12.725644in}{1.263600in}}%
\pgfpathlineto{\pgfqpoint{12.725644in}{1.351193in}}%
\pgfpathlineto{\pgfqpoint{12.725644in}{1.351193in}}%
\pgfpathclose%
\pgfpathmoveto{\pgfqpoint{12.756607in}{1.288475in}}%
\pgfpathlineto{\pgfqpoint{12.756607in}{1.326643in}}%
\pgfpathlineto{\pgfqpoint{12.766964in}{1.326643in}}%
\pgfpathlineto{\pgfqpoint{12.766964in}{1.288871in}}%
\pgfpathquadraticcurveto{\pgfqpoint{12.766964in}{1.279919in}}{\pgfqpoint{12.770440in}{1.275434in}}%
\pgfpathquadraticcurveto{\pgfqpoint{12.773935in}{1.270967in}}{\pgfqpoint{12.780923in}{1.270967in}}%
\pgfpathquadraticcurveto{\pgfqpoint{12.789299in}{1.270967in}}{\pgfqpoint{12.794162in}{1.276317in}}%
\pgfpathquadraticcurveto{\pgfqpoint{12.799044in}{1.281666in}}{\pgfqpoint{12.799044in}{1.290906in}}%
\pgfpathlineto{\pgfqpoint{12.799044in}{1.326643in}}%
\pgfpathlineto{\pgfqpoint{12.809401in}{1.326643in}}%
\pgfpathlineto{\pgfqpoint{12.809401in}{1.263600in}}%
\pgfpathlineto{\pgfqpoint{12.799044in}{1.263600in}}%
\pgfpathlineto{\pgfqpoint{12.799044in}{1.273291in}}%
\pgfpathquadraticcurveto{\pgfqpoint{12.795279in}{1.267545in}}{\pgfqpoint{12.790290in}{1.264753in}}%
\pgfpathquadraticcurveto{\pgfqpoint{12.785318in}{1.261961in}}{\pgfqpoint{12.778726in}{1.261961in}}%
\pgfpathquadraticcurveto{\pgfqpoint{12.767865in}{1.261961in}}{\pgfqpoint{12.762227in}{1.268715in}}%
\pgfpathquadraticcurveto{\pgfqpoint{12.756607in}{1.275470in}}{\pgfqpoint{12.756607in}{1.288475in}}%
\pgfpathlineto{\pgfqpoint{12.756607in}{1.288475in}}%
\pgfpathclose%
\pgfpathmoveto{\pgfqpoint{12.782671in}{1.328156in}}%
\pgfpathlineto{\pgfqpoint{12.782671in}{1.328156in}}%
\pgfpathlineto{\pgfqpoint{12.782671in}{1.328156in}}%
\pgfpathclose%
\pgfpathmoveto{\pgfqpoint{12.872209in}{1.317078in}}%
\pgfpathlineto{\pgfqpoint{12.872209in}{1.351193in}}%
\pgfpathlineto{\pgfqpoint{12.882566in}{1.351193in}}%
\pgfpathlineto{\pgfqpoint{12.882566in}{1.263600in}}%
\pgfpathlineto{\pgfqpoint{12.872209in}{1.263600in}}%
\pgfpathlineto{\pgfqpoint{12.872209in}{1.273056in}}%
\pgfpathquadraticcurveto{\pgfqpoint{12.868949in}{1.267437in}}{\pgfqpoint{12.863959in}{1.264699in}}%
\pgfpathquadraticcurveto{\pgfqpoint{12.858988in}{1.261961in}}{\pgfqpoint{12.851999in}{1.261961in}}%
\pgfpathquadraticcurveto{\pgfqpoint{12.840580in}{1.261961in}}{\pgfqpoint{12.833393in}{1.271075in}}%
\pgfpathquadraticcurveto{\pgfqpoint{12.826224in}{1.280207in}}{\pgfqpoint{12.826224in}{1.295067in}}%
\pgfpathquadraticcurveto{\pgfqpoint{12.826224in}{1.309927in}}{\pgfqpoint{12.833393in}{1.319041in}}%
\pgfpathquadraticcurveto{\pgfqpoint{12.840580in}{1.328156in}}{\pgfqpoint{12.851999in}{1.328156in}}%
\pgfpathquadraticcurveto{\pgfqpoint{12.858988in}{1.328156in}}{\pgfqpoint{12.863959in}{1.325418in}}%
\pgfpathquadraticcurveto{\pgfqpoint{12.868949in}{1.322698in}}{\pgfqpoint{12.872209in}{1.317078in}}%
\pgfpathlineto{\pgfqpoint{12.872209in}{1.317078in}}%
\pgfpathclose%
\pgfpathmoveto{\pgfqpoint{12.836923in}{1.295067in}}%
\pgfpathquadraticcurveto{\pgfqpoint{12.836923in}{1.283648in}}{\pgfqpoint{12.841624in}{1.277145in}}%
\pgfpathquadraticcurveto{\pgfqpoint{12.846326in}{1.270643in}}{\pgfqpoint{12.854539in}{1.270643in}}%
\pgfpathquadraticcurveto{\pgfqpoint{12.862753in}{1.270643in}}{\pgfqpoint{12.867472in}{1.277145in}}%
\pgfpathquadraticcurveto{\pgfqpoint{12.872209in}{1.283648in}}{\pgfqpoint{12.872209in}{1.295067in}}%
\pgfpathquadraticcurveto{\pgfqpoint{12.872209in}{1.306487in}}{\pgfqpoint{12.867472in}{1.312989in}}%
\pgfpathquadraticcurveto{\pgfqpoint{12.862753in}{1.319492in}}{\pgfqpoint{12.854539in}{1.319492in}}%
\pgfpathquadraticcurveto{\pgfqpoint{12.846326in}{1.319492in}}{\pgfqpoint{12.841624in}{1.312989in}}%
\pgfpathquadraticcurveto{\pgfqpoint{12.836923in}{1.306487in}}{\pgfqpoint{12.836923in}{1.295067in}}%
\pgfpathlineto{\pgfqpoint{12.836923in}{1.295067in}}%
\pgfpathclose%
\pgfpathmoveto{\pgfqpoint{12.903910in}{1.326643in}}%
\pgfpathlineto{\pgfqpoint{12.914267in}{1.326643in}}%
\pgfpathlineto{\pgfqpoint{12.914267in}{1.263600in}}%
\pgfpathlineto{\pgfqpoint{12.903910in}{1.263600in}}%
\pgfpathlineto{\pgfqpoint{12.903910in}{1.326643in}}%
\pgfpathlineto{\pgfqpoint{12.903910in}{1.326643in}}%
\pgfpathclose%
\pgfpathmoveto{\pgfqpoint{12.903910in}{1.351193in}}%
\pgfpathlineto{\pgfqpoint{12.914267in}{1.351193in}}%
\pgfpathlineto{\pgfqpoint{12.914267in}{1.338062in}}%
\pgfpathlineto{\pgfqpoint{12.903910in}{1.338062in}}%
\pgfpathlineto{\pgfqpoint{12.903910in}{1.351193in}}%
\pgfpathlineto{\pgfqpoint{12.903910in}{1.351193in}}%
\pgfpathclose%
\pgfpathmoveto{\pgfqpoint{12.988351in}{1.301660in}}%
\pgfpathlineto{\pgfqpoint{12.988351in}{1.263600in}}%
\pgfpathlineto{\pgfqpoint{12.977994in}{1.263600in}}%
\pgfpathlineto{\pgfqpoint{12.977994in}{1.301317in}}%
\pgfpathquadraticcurveto{\pgfqpoint{12.977994in}{1.310269in}}{\pgfqpoint{12.974500in}{1.314700in}}%
\pgfpathquadraticcurveto{\pgfqpoint{12.971006in}{1.319149in}}{\pgfqpoint{12.964035in}{1.319149in}}%
\pgfpathquadraticcurveto{\pgfqpoint{12.955641in}{1.319149in}}{\pgfqpoint{12.950796in}{1.313800in}}%
\pgfpathquadraticcurveto{\pgfqpoint{12.945951in}{1.308468in}}{\pgfqpoint{12.945951in}{1.299228in}}%
\pgfpathlineto{\pgfqpoint{12.945951in}{1.263600in}}%
\pgfpathlineto{\pgfqpoint{12.935540in}{1.263600in}}%
\pgfpathlineto{\pgfqpoint{12.935540in}{1.326643in}}%
\pgfpathlineto{\pgfqpoint{12.945951in}{1.326643in}}%
\pgfpathlineto{\pgfqpoint{12.945951in}{1.316844in}}%
\pgfpathquadraticcurveto{\pgfqpoint{12.949679in}{1.322536in}}{\pgfqpoint{12.954705in}{1.325346in}}%
\pgfpathquadraticcurveto{\pgfqpoint{12.959748in}{1.328156in}}{\pgfqpoint{12.966340in}{1.328156in}}%
\pgfpathquadraticcurveto{\pgfqpoint{12.977202in}{1.328156in}}{\pgfqpoint{12.982768in}{1.321437in}}%
\pgfpathquadraticcurveto{\pgfqpoint{12.988351in}{1.314718in}}{\pgfqpoint{12.988351in}{1.301660in}}%
\pgfpathlineto{\pgfqpoint{12.988351in}{1.301660in}}%
\pgfpathclose%
\pgfpathmoveto{\pgfqpoint{13.050475in}{1.295860in}}%
\pgfpathquadraticcurveto{\pgfqpoint{13.050475in}{1.307117in}}{\pgfqpoint{13.045828in}{1.313296in}}%
\pgfpathquadraticcurveto{\pgfqpoint{13.041199in}{1.319492in}}{\pgfqpoint{13.032805in}{1.319492in}}%
\pgfpathquadraticcurveto{\pgfqpoint{13.024484in}{1.319492in}}{\pgfqpoint{13.019837in}{1.313296in}}%
\pgfpathquadraticcurveto{\pgfqpoint{13.015189in}{1.307117in}}{\pgfqpoint{13.015189in}{1.295860in}}%
\pgfpathquadraticcurveto{\pgfqpoint{13.015189in}{1.284656in}}{\pgfqpoint{13.019837in}{1.278460in}}%
\pgfpathquadraticcurveto{\pgfqpoint{13.024484in}{1.272264in}}{\pgfqpoint{13.032805in}{1.272264in}}%
\pgfpathquadraticcurveto{\pgfqpoint{13.041199in}{1.272264in}}{\pgfqpoint{13.045828in}{1.278460in}}%
\pgfpathquadraticcurveto{\pgfqpoint{13.050475in}{1.284656in}}{\pgfqpoint{13.050475in}{1.295860in}}%
\pgfpathlineto{\pgfqpoint{13.050475in}{1.295860in}}%
\pgfpathclose%
\pgfpathmoveto{\pgfqpoint{13.060832in}{1.271417in}}%
\pgfpathquadraticcurveto{\pgfqpoint{13.060832in}{1.255332in}}{\pgfqpoint{13.053681in}{1.247479in}}%
\pgfpathquadraticcurveto{\pgfqpoint{13.046549in}{1.239626in}}{\pgfqpoint{13.031797in}{1.239626in}}%
\pgfpathquadraticcurveto{\pgfqpoint{13.026339in}{1.239626in}}{\pgfqpoint{13.021494in}{1.240436in}}%
\pgfpathquadraticcurveto{\pgfqpoint{13.016648in}{1.241247in}}{\pgfqpoint{13.012091in}{1.242940in}}%
\pgfpathlineto{\pgfqpoint{13.012091in}{1.253009in}}%
\pgfpathquadraticcurveto{\pgfqpoint{13.016648in}{1.250541in}}{\pgfqpoint{13.021097in}{1.249370in}}%
\pgfpathquadraticcurveto{\pgfqpoint{13.025546in}{1.248182in}}{\pgfqpoint{13.030158in}{1.248182in}}%
\pgfpathquadraticcurveto{\pgfqpoint{13.040352in}{1.248182in}}{\pgfqpoint{13.045414in}{1.253495in}}%
\pgfpathquadraticcurveto{\pgfqpoint{13.050475in}{1.258809in}}{\pgfqpoint{13.050475in}{1.269562in}}%
\pgfpathlineto{\pgfqpoint{13.050475in}{1.274695in}}%
\pgfpathquadraticcurveto{\pgfqpoint{13.047269in}{1.269112in}}{\pgfqpoint{13.042262in}{1.266356in}}%
\pgfpathquadraticcurveto{\pgfqpoint{13.037254in}{1.263600in}}{\pgfqpoint{13.030266in}{1.263600in}}%
\pgfpathquadraticcurveto{\pgfqpoint{13.018684in}{1.263600in}}{\pgfqpoint{13.011587in}{1.272426in}}%
\pgfpathquadraticcurveto{\pgfqpoint{13.004490in}{1.281270in}}{\pgfqpoint{13.004490in}{1.295860in}}%
\pgfpathquadraticcurveto{\pgfqpoint{13.004490in}{1.310486in}}{\pgfqpoint{13.011587in}{1.319312in}}%
\pgfpathquadraticcurveto{\pgfqpoint{13.018684in}{1.328156in}}{\pgfqpoint{13.030266in}{1.328156in}}%
\pgfpathquadraticcurveto{\pgfqpoint{13.037254in}{1.328156in}}{\pgfqpoint{13.042262in}{1.325400in}}%
\pgfpathquadraticcurveto{\pgfqpoint{13.047269in}{1.322644in}}{\pgfqpoint{13.050475in}{1.317078in}}%
\pgfpathlineto{\pgfqpoint{13.050475in}{1.326643in}}%
\pgfpathlineto{\pgfqpoint{13.060832in}{1.326643in}}%
\pgfpathlineto{\pgfqpoint{13.060832in}{1.271417in}}%
\pgfpathlineto{\pgfqpoint{13.060832in}{1.271417in}}%
\pgfpathclose%
\pgfusepath{fill}%
\end{pgfscope}%
\begin{pgfscope}%
\pgfsetbuttcap%
\pgfsetroundjoin%
\definecolor{currentfill}{rgb}{0.309804,0.372549,0.419608}%
\pgfsetfillcolor{currentfill}%
\pgfsetlinewidth{0.000000pt}%
\definecolor{currentstroke}{rgb}{0.309804,0.372549,0.419608}%
\pgfsetstrokecolor{currentstroke}%
\pgfsetdash{}{0pt}%
\pgfpathmoveto{\pgfqpoint{13.194980in}{1.351193in}}%
\pgfpathlineto{\pgfqpoint{13.223800in}{1.351193in}}%
\pgfpathlineto{\pgfqpoint{13.223800in}{1.263600in}}%
\pgfpathlineto{\pgfqpoint{13.213389in}{1.263600in}}%
\pgfpathlineto{\pgfqpoint{13.213389in}{1.342565in}}%
\pgfpathlineto{\pgfqpoint{13.194872in}{1.342565in}}%
\pgfpathquadraticcurveto{\pgfqpoint{13.189306in}{1.342565in}}{\pgfqpoint{13.187127in}{1.340314in}}%
\pgfpathquadraticcurveto{\pgfqpoint{13.184965in}{1.338062in}}{\pgfqpoint{13.184965in}{1.332208in}}%
\pgfpathlineto{\pgfqpoint{13.184965in}{1.326643in}}%
\pgfpathlineto{\pgfqpoint{13.202023in}{1.326643in}}%
\pgfpathlineto{\pgfqpoint{13.202023in}{1.318591in}}%
\pgfpathlineto{\pgfqpoint{13.184965in}{1.318591in}}%
\pgfpathlineto{\pgfqpoint{13.184965in}{1.263600in}}%
\pgfpathlineto{\pgfqpoint{13.174554in}{1.263600in}}%
\pgfpathlineto{\pgfqpoint{13.174554in}{1.318591in}}%
\pgfpathlineto{\pgfqpoint{13.164648in}{1.318591in}}%
\pgfpathlineto{\pgfqpoint{13.164648in}{1.326643in}}%
\pgfpathlineto{\pgfqpoint{13.174554in}{1.326643in}}%
\pgfpathlineto{\pgfqpoint{13.174554in}{1.331038in}}%
\pgfpathquadraticcurveto{\pgfqpoint{13.174554in}{1.341557in}}{\pgfqpoint{13.179454in}{1.346366in}}%
\pgfpathquadraticcurveto{\pgfqpoint{13.184353in}{1.351193in}}{\pgfqpoint{13.194980in}{1.351193in}}%
\pgfpathlineto{\pgfqpoint{13.194980in}{1.351193in}}%
\pgfpathclose%
\pgfpathmoveto{\pgfqpoint{13.274126in}{1.295283in}}%
\pgfpathquadraticcurveto{\pgfqpoint{13.261571in}{1.295283in}}{\pgfqpoint{13.256726in}{1.292419in}}%
\pgfpathquadraticcurveto{\pgfqpoint{13.251881in}{1.289556in}}{\pgfqpoint{13.251881in}{1.282621in}}%
\pgfpathquadraticcurveto{\pgfqpoint{13.251881in}{1.277109in}}{\pgfqpoint{13.255519in}{1.273867in}}%
\pgfpathquadraticcurveto{\pgfqpoint{13.259158in}{1.270643in}}{\pgfqpoint{13.265390in}{1.270643in}}%
\pgfpathquadraticcurveto{\pgfqpoint{13.274018in}{1.270643in}}{\pgfqpoint{13.279223in}{1.276749in}}%
\pgfpathquadraticcurveto{\pgfqpoint{13.284429in}{1.282855in}}{\pgfqpoint{13.284429in}{1.292978in}}%
\pgfpathlineto{\pgfqpoint{13.284429in}{1.295283in}}%
\pgfpathlineto{\pgfqpoint{13.274126in}{1.295283in}}%
\pgfpathlineto{\pgfqpoint{13.274126in}{1.295283in}}%
\pgfpathclose%
\pgfpathmoveto{\pgfqpoint{13.294786in}{1.299570in}}%
\pgfpathlineto{\pgfqpoint{13.294786in}{1.263600in}}%
\pgfpathlineto{\pgfqpoint{13.284429in}{1.263600in}}%
\pgfpathlineto{\pgfqpoint{13.284429in}{1.273164in}}%
\pgfpathquadraticcurveto{\pgfqpoint{13.280880in}{1.267437in}}{\pgfqpoint{13.275585in}{1.264699in}}%
\pgfpathquadraticcurveto{\pgfqpoint{13.270289in}{1.261961in}}{\pgfqpoint{13.262634in}{1.261961in}}%
\pgfpathquadraticcurveto{\pgfqpoint{13.252961in}{1.261961in}}{\pgfqpoint{13.247234in}{1.267401in}}%
\pgfpathquadraticcurveto{\pgfqpoint{13.241524in}{1.272840in}}{\pgfqpoint{13.241524in}{1.281954in}}%
\pgfpathquadraticcurveto{\pgfqpoint{13.241524in}{1.292582in}}{\pgfqpoint{13.248638in}{1.297985in}}%
\pgfpathquadraticcurveto{\pgfqpoint{13.255771in}{1.303389in}}{\pgfqpoint{13.269893in}{1.303389in}}%
\pgfpathlineto{\pgfqpoint{13.284429in}{1.303389in}}%
\pgfpathlineto{\pgfqpoint{13.284429in}{1.304416in}}%
\pgfpathquadraticcurveto{\pgfqpoint{13.284429in}{1.311566in}}{\pgfqpoint{13.279727in}{1.315475in}}%
\pgfpathquadraticcurveto{\pgfqpoint{13.275026in}{1.319384in}}{\pgfqpoint{13.266525in}{1.319384in}}%
\pgfpathquadraticcurveto{\pgfqpoint{13.261121in}{1.319384in}}{\pgfqpoint{13.255987in}{1.318087in}}%
\pgfpathquadraticcurveto{\pgfqpoint{13.250872in}{1.316790in}}{\pgfqpoint{13.246153in}{1.314196in}}%
\pgfpathlineto{\pgfqpoint{13.246153in}{1.323779in}}%
\pgfpathquadraticcurveto{\pgfqpoint{13.251827in}{1.325976in}}{\pgfqpoint{13.257176in}{1.327057in}}%
\pgfpathquadraticcurveto{\pgfqpoint{13.262526in}{1.328156in}}{\pgfqpoint{13.267587in}{1.328156in}}%
\pgfpathquadraticcurveto{\pgfqpoint{13.281276in}{1.328156in}}{\pgfqpoint{13.288031in}{1.321059in}}%
\pgfpathquadraticcurveto{\pgfqpoint{13.294786in}{1.313980in}}{\pgfqpoint{13.294786in}{1.299570in}}%
\pgfpathlineto{\pgfqpoint{13.294786in}{1.299570in}}%
\pgfpathclose%
\pgfpathmoveto{\pgfqpoint{13.357594in}{1.295860in}}%
\pgfpathquadraticcurveto{\pgfqpoint{13.357594in}{1.307117in}}{\pgfqpoint{13.352947in}{1.313296in}}%
\pgfpathquadraticcurveto{\pgfqpoint{13.348318in}{1.319492in}}{\pgfqpoint{13.339924in}{1.319492in}}%
\pgfpathquadraticcurveto{\pgfqpoint{13.331602in}{1.319492in}}{\pgfqpoint{13.326955in}{1.313296in}}%
\pgfpathquadraticcurveto{\pgfqpoint{13.322308in}{1.307117in}}{\pgfqpoint{13.322308in}{1.295860in}}%
\pgfpathquadraticcurveto{\pgfqpoint{13.322308in}{1.284656in}}{\pgfqpoint{13.326955in}{1.278460in}}%
\pgfpathquadraticcurveto{\pgfqpoint{13.331602in}{1.272264in}}{\pgfqpoint{13.339924in}{1.272264in}}%
\pgfpathquadraticcurveto{\pgfqpoint{13.348318in}{1.272264in}}{\pgfqpoint{13.352947in}{1.278460in}}%
\pgfpathquadraticcurveto{\pgfqpoint{13.357594in}{1.284656in}}{\pgfqpoint{13.357594in}{1.295860in}}%
\pgfpathlineto{\pgfqpoint{13.357594in}{1.295860in}}%
\pgfpathclose%
\pgfpathmoveto{\pgfqpoint{13.367951in}{1.271417in}}%
\pgfpathquadraticcurveto{\pgfqpoint{13.367951in}{1.255332in}}{\pgfqpoint{13.360800in}{1.247479in}}%
\pgfpathquadraticcurveto{\pgfqpoint{13.353667in}{1.239626in}}{\pgfqpoint{13.338915in}{1.239626in}}%
\pgfpathquadraticcurveto{\pgfqpoint{13.333458in}{1.239626in}}{\pgfqpoint{13.328612in}{1.240436in}}%
\pgfpathquadraticcurveto{\pgfqpoint{13.323767in}{1.241247in}}{\pgfqpoint{13.319210in}{1.242940in}}%
\pgfpathlineto{\pgfqpoint{13.319210in}{1.253009in}}%
\pgfpathquadraticcurveto{\pgfqpoint{13.323767in}{1.250541in}}{\pgfqpoint{13.328216in}{1.249370in}}%
\pgfpathquadraticcurveto{\pgfqpoint{13.332665in}{1.248182in}}{\pgfqpoint{13.337276in}{1.248182in}}%
\pgfpathquadraticcurveto{\pgfqpoint{13.347471in}{1.248182in}}{\pgfqpoint{13.352533in}{1.253495in}}%
\pgfpathquadraticcurveto{\pgfqpoint{13.357594in}{1.258809in}}{\pgfqpoint{13.357594in}{1.269562in}}%
\pgfpathlineto{\pgfqpoint{13.357594in}{1.274695in}}%
\pgfpathquadraticcurveto{\pgfqpoint{13.354388in}{1.269112in}}{\pgfqpoint{13.349380in}{1.266356in}}%
\pgfpathquadraticcurveto{\pgfqpoint{13.344373in}{1.263600in}}{\pgfqpoint{13.337384in}{1.263600in}}%
\pgfpathquadraticcurveto{\pgfqpoint{13.325803in}{1.263600in}}{\pgfqpoint{13.318706in}{1.272426in}}%
\pgfpathquadraticcurveto{\pgfqpoint{13.311609in}{1.281270in}}{\pgfqpoint{13.311609in}{1.295860in}}%
\pgfpathquadraticcurveto{\pgfqpoint{13.311609in}{1.310486in}}{\pgfqpoint{13.318706in}{1.319312in}}%
\pgfpathquadraticcurveto{\pgfqpoint{13.325803in}{1.328156in}}{\pgfqpoint{13.337384in}{1.328156in}}%
\pgfpathquadraticcurveto{\pgfqpoint{13.344373in}{1.328156in}}{\pgfqpoint{13.349380in}{1.325400in}}%
\pgfpathquadraticcurveto{\pgfqpoint{13.354388in}{1.322644in}}{\pgfqpoint{13.357594in}{1.317078in}}%
\pgfpathlineto{\pgfqpoint{13.357594in}{1.326643in}}%
\pgfpathlineto{\pgfqpoint{13.367951in}{1.326643in}}%
\pgfpathlineto{\pgfqpoint{13.367951in}{1.271417in}}%
\pgfpathlineto{\pgfqpoint{13.367951in}{1.271417in}}%
\pgfpathclose%
\pgfpathmoveto{\pgfqpoint{13.430777in}{1.295860in}}%
\pgfpathquadraticcurveto{\pgfqpoint{13.430777in}{1.307117in}}{\pgfqpoint{13.426130in}{1.313296in}}%
\pgfpathquadraticcurveto{\pgfqpoint{13.421501in}{1.319492in}}{\pgfqpoint{13.413107in}{1.319492in}}%
\pgfpathquadraticcurveto{\pgfqpoint{13.404786in}{1.319492in}}{\pgfqpoint{13.400139in}{1.313296in}}%
\pgfpathquadraticcurveto{\pgfqpoint{13.395492in}{1.307117in}}{\pgfqpoint{13.395492in}{1.295860in}}%
\pgfpathquadraticcurveto{\pgfqpoint{13.395492in}{1.284656in}}{\pgfqpoint{13.400139in}{1.278460in}}%
\pgfpathquadraticcurveto{\pgfqpoint{13.404786in}{1.272264in}}{\pgfqpoint{13.413107in}{1.272264in}}%
\pgfpathquadraticcurveto{\pgfqpoint{13.421501in}{1.272264in}}{\pgfqpoint{13.426130in}{1.278460in}}%
\pgfpathquadraticcurveto{\pgfqpoint{13.430777in}{1.284656in}}{\pgfqpoint{13.430777in}{1.295860in}}%
\pgfpathlineto{\pgfqpoint{13.430777in}{1.295860in}}%
\pgfpathclose%
\pgfpathmoveto{\pgfqpoint{13.441134in}{1.271417in}}%
\pgfpathquadraticcurveto{\pgfqpoint{13.441134in}{1.255332in}}{\pgfqpoint{13.433984in}{1.247479in}}%
\pgfpathquadraticcurveto{\pgfqpoint{13.426851in}{1.239626in}}{\pgfqpoint{13.412099in}{1.239626in}}%
\pgfpathquadraticcurveto{\pgfqpoint{13.406641in}{1.239626in}}{\pgfqpoint{13.401796in}{1.240436in}}%
\pgfpathquadraticcurveto{\pgfqpoint{13.396951in}{1.241247in}}{\pgfqpoint{13.392393in}{1.242940in}}%
\pgfpathlineto{\pgfqpoint{13.392393in}{1.253009in}}%
\pgfpathquadraticcurveto{\pgfqpoint{13.396951in}{1.250541in}}{\pgfqpoint{13.401400in}{1.249370in}}%
\pgfpathquadraticcurveto{\pgfqpoint{13.405849in}{1.248182in}}{\pgfqpoint{13.410460in}{1.248182in}}%
\pgfpathquadraticcurveto{\pgfqpoint{13.420655in}{1.248182in}}{\pgfqpoint{13.425716in}{1.253495in}}%
\pgfpathquadraticcurveto{\pgfqpoint{13.430777in}{1.258809in}}{\pgfqpoint{13.430777in}{1.269562in}}%
\pgfpathlineto{\pgfqpoint{13.430777in}{1.274695in}}%
\pgfpathquadraticcurveto{\pgfqpoint{13.427571in}{1.269112in}}{\pgfqpoint{13.422564in}{1.266356in}}%
\pgfpathquadraticcurveto{\pgfqpoint{13.417556in}{1.263600in}}{\pgfqpoint{13.410568in}{1.263600in}}%
\pgfpathquadraticcurveto{\pgfqpoint{13.398986in}{1.263600in}}{\pgfqpoint{13.391889in}{1.272426in}}%
\pgfpathquadraticcurveto{\pgfqpoint{13.384792in}{1.281270in}}{\pgfqpoint{13.384792in}{1.295860in}}%
\pgfpathquadraticcurveto{\pgfqpoint{13.384792in}{1.310486in}}{\pgfqpoint{13.391889in}{1.319312in}}%
\pgfpathquadraticcurveto{\pgfqpoint{13.398986in}{1.328156in}}{\pgfqpoint{13.410568in}{1.328156in}}%
\pgfpathquadraticcurveto{\pgfqpoint{13.417556in}{1.328156in}}{\pgfqpoint{13.422564in}{1.325400in}}%
\pgfpathquadraticcurveto{\pgfqpoint{13.427571in}{1.322644in}}{\pgfqpoint{13.430777in}{1.317078in}}%
\pgfpathlineto{\pgfqpoint{13.430777in}{1.326643in}}%
\pgfpathlineto{\pgfqpoint{13.441134in}{1.326643in}}%
\pgfpathlineto{\pgfqpoint{13.441134in}{1.271417in}}%
\pgfpathlineto{\pgfqpoint{13.441134in}{1.271417in}}%
\pgfpathclose%
\pgfpathmoveto{\pgfqpoint{13.516407in}{1.297715in}}%
\pgfpathlineto{\pgfqpoint{13.516407in}{1.292654in}}%
\pgfpathlineto{\pgfqpoint{13.468783in}{1.292654in}}%
\pgfpathquadraticcurveto{\pgfqpoint{13.469467in}{1.281954in}}{\pgfqpoint{13.475231in}{1.276353in}}%
\pgfpathquadraticcurveto{\pgfqpoint{13.480995in}{1.270751in}}{\pgfqpoint{13.491298in}{1.270751in}}%
\pgfpathquadraticcurveto{\pgfqpoint{13.497260in}{1.270751in}}{\pgfqpoint{13.502862in}{1.272210in}}%
\pgfpathquadraticcurveto{\pgfqpoint{13.508464in}{1.273669in}}{\pgfqpoint{13.513993in}{1.276605in}}%
\pgfpathlineto{\pgfqpoint{13.513993in}{1.266806in}}%
\pgfpathquadraticcurveto{\pgfqpoint{13.508410in}{1.264447in}}{\pgfqpoint{13.502556in}{1.263204in}}%
\pgfpathquadraticcurveto{\pgfqpoint{13.496702in}{1.261961in}}{\pgfqpoint{13.490686in}{1.261961in}}%
\pgfpathquadraticcurveto{\pgfqpoint{13.475592in}{1.261961in}}{\pgfqpoint{13.466784in}{1.270733in}}%
\pgfpathquadraticcurveto{\pgfqpoint{13.457976in}{1.279523in}}{\pgfqpoint{13.457976in}{1.294509in}}%
\pgfpathquadraticcurveto{\pgfqpoint{13.457976in}{1.309981in}}{\pgfqpoint{13.466333in}{1.319059in}}%
\pgfpathquadraticcurveto{\pgfqpoint{13.474691in}{1.328156in}}{\pgfqpoint{13.488885in}{1.328156in}}%
\pgfpathquadraticcurveto{\pgfqpoint{13.501601in}{1.328156in}}{\pgfqpoint{13.509004in}{1.319960in}}%
\pgfpathquadraticcurveto{\pgfqpoint{13.516407in}{1.311783in}}{\pgfqpoint{13.516407in}{1.297715in}}%
\pgfpathlineto{\pgfqpoint{13.516407in}{1.297715in}}%
\pgfpathclose%
\pgfpathmoveto{\pgfqpoint{13.506050in}{1.300759in}}%
\pgfpathquadraticcurveto{\pgfqpoint{13.505942in}{1.309243in}}{\pgfqpoint{13.501295in}{1.314304in}}%
\pgfpathquadraticcurveto{\pgfqpoint{13.496648in}{1.319384in}}{\pgfqpoint{13.488993in}{1.319384in}}%
\pgfpathquadraticcurveto{\pgfqpoint{13.480329in}{1.319384in}}{\pgfqpoint{13.475123in}{1.314484in}}%
\pgfpathquadraticcurveto{\pgfqpoint{13.469918in}{1.309585in}}{\pgfqpoint{13.469125in}{1.300687in}}%
\pgfpathlineto{\pgfqpoint{13.506050in}{1.300759in}}%
\pgfpathlineto{\pgfqpoint{13.506050in}{1.300759in}}%
\pgfpathclose%
\pgfpathmoveto{\pgfqpoint{13.574893in}{1.317078in}}%
\pgfpathlineto{\pgfqpoint{13.574893in}{1.351193in}}%
\pgfpathlineto{\pgfqpoint{13.585250in}{1.351193in}}%
\pgfpathlineto{\pgfqpoint{13.585250in}{1.263600in}}%
\pgfpathlineto{\pgfqpoint{13.574893in}{1.263600in}}%
\pgfpathlineto{\pgfqpoint{13.574893in}{1.273056in}}%
\pgfpathquadraticcurveto{\pgfqpoint{13.571632in}{1.267437in}}{\pgfqpoint{13.566643in}{1.264699in}}%
\pgfpathquadraticcurveto{\pgfqpoint{13.561672in}{1.261961in}}{\pgfqpoint{13.554683in}{1.261961in}}%
\pgfpathquadraticcurveto{\pgfqpoint{13.543263in}{1.261961in}}{\pgfqpoint{13.536076in}{1.271075in}}%
\pgfpathquadraticcurveto{\pgfqpoint{13.528908in}{1.280207in}}{\pgfqpoint{13.528908in}{1.295067in}}%
\pgfpathquadraticcurveto{\pgfqpoint{13.528908in}{1.309927in}}{\pgfqpoint{13.536076in}{1.319041in}}%
\pgfpathquadraticcurveto{\pgfqpoint{13.543263in}{1.328156in}}{\pgfqpoint{13.554683in}{1.328156in}}%
\pgfpathquadraticcurveto{\pgfqpoint{13.561672in}{1.328156in}}{\pgfqpoint{13.566643in}{1.325418in}}%
\pgfpathquadraticcurveto{\pgfqpoint{13.571632in}{1.322698in}}{\pgfqpoint{13.574893in}{1.317078in}}%
\pgfpathlineto{\pgfqpoint{13.574893in}{1.317078in}}%
\pgfpathclose%
\pgfpathmoveto{\pgfqpoint{13.539607in}{1.295067in}}%
\pgfpathquadraticcurveto{\pgfqpoint{13.539607in}{1.283648in}}{\pgfqpoint{13.544308in}{1.277145in}}%
\pgfpathquadraticcurveto{\pgfqpoint{13.549009in}{1.270643in}}{\pgfqpoint{13.557223in}{1.270643in}}%
\pgfpathquadraticcurveto{\pgfqpoint{13.565436in}{1.270643in}}{\pgfqpoint{13.570155in}{1.277145in}}%
\pgfpathquadraticcurveto{\pgfqpoint{13.574893in}{1.283648in}}{\pgfqpoint{13.574893in}{1.295067in}}%
\pgfpathquadraticcurveto{\pgfqpoint{13.574893in}{1.306487in}}{\pgfqpoint{13.570155in}{1.312989in}}%
\pgfpathquadraticcurveto{\pgfqpoint{13.565436in}{1.319492in}}{\pgfqpoint{13.557223in}{1.319492in}}%
\pgfpathquadraticcurveto{\pgfqpoint{13.549009in}{1.319492in}}{\pgfqpoint{13.544308in}{1.312989in}}%
\pgfpathquadraticcurveto{\pgfqpoint{13.539607in}{1.306487in}}{\pgfqpoint{13.539607in}{1.295067in}}%
\pgfpathlineto{\pgfqpoint{13.539607in}{1.295067in}}%
\pgfpathclose%
\pgfusepath{fill}%
\end{pgfscope}%
\begin{pgfscope}%
\pgfsetbuttcap%
\pgfsetroundjoin%
\definecolor{currentfill}{rgb}{0.309804,0.372549,0.419608}%
\pgfsetfillcolor{currentfill}%
\pgfsetlinewidth{0.000000pt}%
\definecolor{currentstroke}{rgb}{0.309804,0.372549,0.419608}%
\pgfsetstrokecolor{currentstroke}%
\pgfsetdash{}{0pt}%
\pgfpathmoveto{\pgfqpoint{0.944876in}{1.089456in}}%
\pgfpathlineto{\pgfqpoint{0.944876in}{1.056026in}}%
\pgfpathlineto{\pgfqpoint{0.934465in}{1.056026in}}%
\pgfpathlineto{\pgfqpoint{0.934465in}{1.143043in}}%
\pgfpathlineto{\pgfqpoint{0.944876in}{1.143043in}}%
\pgfpathlineto{\pgfqpoint{0.944876in}{1.133478in}}%
\pgfpathquadraticcurveto{\pgfqpoint{0.948154in}{1.139098in}}{\pgfqpoint{0.953126in}{1.141818in}}%
\pgfpathquadraticcurveto{\pgfqpoint{0.958115in}{1.144556in}}{\pgfqpoint{0.965032in}{1.144556in}}%
\pgfpathquadraticcurveto{\pgfqpoint{0.976523in}{1.144556in}}{\pgfqpoint{0.983692in}{1.135441in}}%
\pgfpathquadraticcurveto{\pgfqpoint{0.990879in}{1.126327in}}{\pgfqpoint{0.990879in}{1.111467in}}%
\pgfpathquadraticcurveto{\pgfqpoint{0.990879in}{1.096607in}}{\pgfqpoint{0.983692in}{1.087475in}}%
\pgfpathquadraticcurveto{\pgfqpoint{0.976523in}{1.078361in}}{\pgfqpoint{0.965032in}{1.078361in}}%
\pgfpathquadraticcurveto{\pgfqpoint{0.958115in}{1.078361in}}{\pgfqpoint{0.953126in}{1.081099in}}%
\pgfpathquadraticcurveto{\pgfqpoint{0.948154in}{1.083837in}}{\pgfqpoint{0.944876in}{1.089456in}}%
\pgfpathlineto{\pgfqpoint{0.944876in}{1.089456in}}%
\pgfpathclose%
\pgfpathmoveto{\pgfqpoint{0.980126in}{1.111467in}}%
\pgfpathquadraticcurveto{\pgfqpoint{0.980126in}{1.122887in}}{\pgfqpoint{0.975425in}{1.129389in}}%
\pgfpathquadraticcurveto{\pgfqpoint{0.970724in}{1.135892in}}{\pgfqpoint{0.962510in}{1.135892in}}%
\pgfpathquadraticcurveto{\pgfqpoint{0.954278in}{1.135892in}}{\pgfqpoint{0.949577in}{1.129389in}}%
\pgfpathquadraticcurveto{\pgfqpoint{0.944876in}{1.122887in}}{\pgfqpoint{0.944876in}{1.111467in}}%
\pgfpathquadraticcurveto{\pgfqpoint{0.944876in}{1.100048in}}{\pgfqpoint{0.949577in}{1.093545in}}%
\pgfpathquadraticcurveto{\pgfqpoint{0.954278in}{1.087043in}}{\pgfqpoint{0.962510in}{1.087043in}}%
\pgfpathquadraticcurveto{\pgfqpoint{0.970724in}{1.087043in}}{\pgfqpoint{0.975425in}{1.093545in}}%
\pgfpathquadraticcurveto{\pgfqpoint{0.980126in}{1.100048in}}{\pgfqpoint{0.980126in}{1.111467in}}%
\pgfpathlineto{\pgfqpoint{0.980126in}{1.111467in}}%
\pgfpathclose%
\pgfpathmoveto{\pgfqpoint{1.032469in}{1.135784in}}%
\pgfpathquadraticcurveto{\pgfqpoint{1.024148in}{1.135784in}}{\pgfqpoint{1.019302in}{1.129281in}}%
\pgfpathquadraticcurveto{\pgfqpoint{1.014457in}{1.122779in}}{\pgfqpoint{1.014457in}{1.111467in}}%
\pgfpathquadraticcurveto{\pgfqpoint{1.014457in}{1.100156in}}{\pgfqpoint{1.019266in}{1.093653in}}%
\pgfpathquadraticcurveto{\pgfqpoint{1.024094in}{1.087151in}}{\pgfqpoint{1.032469in}{1.087151in}}%
\pgfpathquadraticcurveto{\pgfqpoint{1.040755in}{1.087151in}}{\pgfqpoint{1.045582in}{1.093671in}}%
\pgfpathquadraticcurveto{\pgfqpoint{1.050427in}{1.100210in}}{\pgfqpoint{1.050427in}{1.111467in}}%
\pgfpathquadraticcurveto{\pgfqpoint{1.050427in}{1.122671in}}{\pgfqpoint{1.045582in}{1.129227in}}%
\pgfpathquadraticcurveto{\pgfqpoint{1.040755in}{1.135784in}}{\pgfqpoint{1.032469in}{1.135784in}}%
\pgfpathlineto{\pgfqpoint{1.032469in}{1.135784in}}%
\pgfpathclose%
\pgfpathmoveto{\pgfqpoint{1.032469in}{1.144556in}}%
\pgfpathquadraticcurveto{\pgfqpoint{1.045978in}{1.144556in}}{\pgfqpoint{1.053688in}{1.135766in}}%
\pgfpathquadraticcurveto{\pgfqpoint{1.061415in}{1.126994in}}{\pgfqpoint{1.061415in}{1.111467in}}%
\pgfpathquadraticcurveto{\pgfqpoint{1.061415in}{1.095995in}}{\pgfqpoint{1.053688in}{1.087169in}}%
\pgfpathquadraticcurveto{\pgfqpoint{1.045978in}{1.078361in}}{\pgfqpoint{1.032469in}{1.078361in}}%
\pgfpathquadraticcurveto{\pgfqpoint{1.018906in}{1.078361in}}{\pgfqpoint{1.011215in}{1.087169in}}%
\pgfpathquadraticcurveto{\pgfqpoint{1.003542in}{1.095995in}}{\pgfqpoint{1.003542in}{1.111467in}}%
\pgfpathquadraticcurveto{\pgfqpoint{1.003542in}{1.126994in}}{\pgfqpoint{1.011215in}{1.135766in}}%
\pgfpathquadraticcurveto{\pgfqpoint{1.018906in}{1.144556in}}{\pgfqpoint{1.032469in}{1.144556in}}%
\pgfpathlineto{\pgfqpoint{1.032469in}{1.144556in}}%
\pgfpathclose%
\pgfpathmoveto{\pgfqpoint{1.078580in}{1.143043in}}%
\pgfpathlineto{\pgfqpoint{1.088937in}{1.143043in}}%
\pgfpathlineto{\pgfqpoint{1.088937in}{1.080000in}}%
\pgfpathlineto{\pgfqpoint{1.078580in}{1.080000in}}%
\pgfpathlineto{\pgfqpoint{1.078580in}{1.143043in}}%
\pgfpathlineto{\pgfqpoint{1.078580in}{1.143043in}}%
\pgfpathclose%
\pgfpathmoveto{\pgfqpoint{1.078580in}{1.167593in}}%
\pgfpathlineto{\pgfqpoint{1.088937in}{1.167593in}}%
\pgfpathlineto{\pgfqpoint{1.088937in}{1.154462in}}%
\pgfpathlineto{\pgfqpoint{1.078580in}{1.154462in}}%
\pgfpathlineto{\pgfqpoint{1.078580in}{1.167593in}}%
\pgfpathlineto{\pgfqpoint{1.078580in}{1.167593in}}%
\pgfpathclose%
\pgfpathmoveto{\pgfqpoint{1.163021in}{1.118060in}}%
\pgfpathlineto{\pgfqpoint{1.163021in}{1.080000in}}%
\pgfpathlineto{\pgfqpoint{1.152664in}{1.080000in}}%
\pgfpathlineto{\pgfqpoint{1.152664in}{1.117717in}}%
\pgfpathquadraticcurveto{\pgfqpoint{1.152664in}{1.126669in}}{\pgfqpoint{1.149170in}{1.131100in}}%
\pgfpathquadraticcurveto{\pgfqpoint{1.145676in}{1.135549in}}{\pgfqpoint{1.138705in}{1.135549in}}%
\pgfpathquadraticcurveto{\pgfqpoint{1.130311in}{1.135549in}}{\pgfqpoint{1.125466in}{1.130200in}}%
\pgfpathquadraticcurveto{\pgfqpoint{1.120621in}{1.124868in}}{\pgfqpoint{1.120621in}{1.115628in}}%
\pgfpathlineto{\pgfqpoint{1.120621in}{1.080000in}}%
\pgfpathlineto{\pgfqpoint{1.110210in}{1.080000in}}%
\pgfpathlineto{\pgfqpoint{1.110210in}{1.143043in}}%
\pgfpathlineto{\pgfqpoint{1.120621in}{1.143043in}}%
\pgfpathlineto{\pgfqpoint{1.120621in}{1.133244in}}%
\pgfpathquadraticcurveto{\pgfqpoint{1.124349in}{1.138936in}}{\pgfqpoint{1.129375in}{1.141746in}}%
\pgfpathquadraticcurveto{\pgfqpoint{1.134418in}{1.144556in}}{\pgfqpoint{1.141010in}{1.144556in}}%
\pgfpathquadraticcurveto{\pgfqpoint{1.151872in}{1.144556in}}{\pgfqpoint{1.157438in}{1.137837in}}%
\pgfpathquadraticcurveto{\pgfqpoint{1.163021in}{1.131118in}}{\pgfqpoint{1.163021in}{1.118060in}}%
\pgfpathlineto{\pgfqpoint{1.163021in}{1.118060in}}%
\pgfpathclose%
\pgfpathmoveto{\pgfqpoint{1.193912in}{1.160947in}}%
\pgfpathlineto{\pgfqpoint{1.193912in}{1.143043in}}%
\pgfpathlineto{\pgfqpoint{1.215238in}{1.143043in}}%
\pgfpathlineto{\pgfqpoint{1.215238in}{1.134991in}}%
\pgfpathlineto{\pgfqpoint{1.193912in}{1.134991in}}%
\pgfpathlineto{\pgfqpoint{1.193912in}{1.100768in}}%
\pgfpathquadraticcurveto{\pgfqpoint{1.193912in}{1.093059in}}{\pgfqpoint{1.196020in}{1.090861in}}%
\pgfpathquadraticcurveto{\pgfqpoint{1.198127in}{1.088664in}}{\pgfqpoint{1.204611in}{1.088664in}}%
\pgfpathlineto{\pgfqpoint{1.215238in}{1.088664in}}%
\pgfpathlineto{\pgfqpoint{1.215238in}{1.080000in}}%
\pgfpathlineto{\pgfqpoint{1.204611in}{1.080000in}}%
\pgfpathquadraticcurveto{\pgfqpoint{1.192615in}{1.080000in}}{\pgfqpoint{1.188058in}{1.084467in}}%
\pgfpathquadraticcurveto{\pgfqpoint{1.183501in}{1.088952in}}{\pgfqpoint{1.183501in}{1.100768in}}%
\pgfpathlineto{\pgfqpoint{1.183501in}{1.134991in}}%
\pgfpathlineto{\pgfqpoint{1.175900in}{1.134991in}}%
\pgfpathlineto{\pgfqpoint{1.175900in}{1.143043in}}%
\pgfpathlineto{\pgfqpoint{1.183501in}{1.143043in}}%
\pgfpathlineto{\pgfqpoint{1.183501in}{1.160947in}}%
\pgfpathlineto{\pgfqpoint{1.193912in}{1.160947in}}%
\pgfpathlineto{\pgfqpoint{1.193912in}{1.160947in}}%
\pgfpathclose%
\pgfpathmoveto{\pgfqpoint{1.269041in}{1.141187in}}%
\pgfpathlineto{\pgfqpoint{1.269041in}{1.131389in}}%
\pgfpathquadraticcurveto{\pgfqpoint{1.264664in}{1.133640in}}{\pgfqpoint{1.259927in}{1.134757in}}%
\pgfpathquadraticcurveto{\pgfqpoint{1.255207in}{1.135892in}}{\pgfqpoint{1.250128in}{1.135892in}}%
\pgfpathquadraticcurveto{\pgfqpoint{1.242419in}{1.135892in}}{\pgfqpoint{1.238564in}{1.133532in}}%
\pgfpathquadraticcurveto{\pgfqpoint{1.234710in}{1.131173in}}{\pgfqpoint{1.234710in}{1.126435in}}%
\pgfpathquadraticcurveto{\pgfqpoint{1.234710in}{1.122833in}}{\pgfqpoint{1.237465in}{1.120780in}}%
\pgfpathquadraticcurveto{\pgfqpoint{1.240221in}{1.118726in}}{\pgfqpoint{1.248561in}{1.116871in}}%
\pgfpathlineto{\pgfqpoint{1.252109in}{1.116078in}}%
\pgfpathquadraticcurveto{\pgfqpoint{1.263133in}{1.113719in}}{\pgfqpoint{1.267780in}{1.109414in}}%
\pgfpathquadraticcurveto{\pgfqpoint{1.272427in}{1.105109in}}{\pgfqpoint{1.272427in}{1.097400in}}%
\pgfpathquadraticcurveto{\pgfqpoint{1.272427in}{1.088610in}}{\pgfqpoint{1.265474in}{1.083476in}}%
\pgfpathquadraticcurveto{\pgfqpoint{1.258522in}{1.078361in}}{\pgfqpoint{1.246363in}{1.078361in}}%
\pgfpathquadraticcurveto{\pgfqpoint{1.241302in}{1.078361in}}{\pgfqpoint{1.235808in}{1.079352in}}%
\pgfpathquadraticcurveto{\pgfqpoint{1.230315in}{1.080342in}}{\pgfqpoint{1.224245in}{1.082306in}}%
\pgfpathlineto{\pgfqpoint{1.224245in}{1.093005in}}%
\pgfpathquadraticcurveto{\pgfqpoint{1.229990in}{1.090015in}}{\pgfqpoint{1.235556in}{1.088520in}}%
\pgfpathquadraticcurveto{\pgfqpoint{1.241122in}{1.087043in}}{\pgfqpoint{1.246598in}{1.087043in}}%
\pgfpathquadraticcurveto{\pgfqpoint{1.253911in}{1.087043in}}{\pgfqpoint{1.257837in}{1.089546in}}%
\pgfpathquadraticcurveto{\pgfqpoint{1.261782in}{1.092050in}}{\pgfqpoint{1.261782in}{1.096607in}}%
\pgfpathquadraticcurveto{\pgfqpoint{1.261782in}{1.100822in}}{\pgfqpoint{1.258936in}{1.103074in}}%
\pgfpathquadraticcurveto{\pgfqpoint{1.256108in}{1.105325in}}{\pgfqpoint{1.246472in}{1.107414in}}%
\pgfpathlineto{\pgfqpoint{1.242869in}{1.108261in}}%
\pgfpathquadraticcurveto{\pgfqpoint{1.233251in}{1.110278in}}{\pgfqpoint{1.228964in}{1.114475in}}%
\pgfpathquadraticcurveto{\pgfqpoint{1.224695in}{1.118672in}}{\pgfqpoint{1.224695in}{1.125985in}}%
\pgfpathquadraticcurveto{\pgfqpoint{1.224695in}{1.134883in}}{\pgfqpoint{1.230999in}{1.139710in}}%
\pgfpathquadraticcurveto{\pgfqpoint{1.237303in}{1.144556in}}{\pgfqpoint{1.248903in}{1.144556in}}%
\pgfpathquadraticcurveto{\pgfqpoint{1.254631in}{1.144556in}}{\pgfqpoint{1.259692in}{1.143709in}}%
\pgfpathquadraticcurveto{\pgfqpoint{1.264772in}{1.142880in}}{\pgfqpoint{1.269041in}{1.141187in}}%
\pgfpathlineto{\pgfqpoint{1.269041in}{1.141187in}}%
\pgfpathclose%
\pgfpathmoveto{\pgfqpoint{1.290367in}{1.094302in}}%
\pgfpathlineto{\pgfqpoint{1.302255in}{1.094302in}}%
\pgfpathlineto{\pgfqpoint{1.302255in}{1.080000in}}%
\pgfpathlineto{\pgfqpoint{1.290367in}{1.080000in}}%
\pgfpathlineto{\pgfqpoint{1.290367in}{1.094302in}}%
\pgfpathlineto{\pgfqpoint{1.290367in}{1.094302in}}%
\pgfpathclose%
\pgfusepath{fill}%
\end{pgfscope}%
\begin{pgfscope}%
\pgfsetbuttcap%
\pgfsetroundjoin%
\definecolor{currentfill}{rgb}{0.309804,0.372549,0.419608}%
\pgfsetfillcolor{currentfill}%
\pgfsetlinewidth{0.000000pt}%
\definecolor{currentstroke}{rgb}{0.309804,0.372549,0.419608}%
\pgfsetstrokecolor{currentstroke}%
\pgfsetdash{}{0pt}%
\pgfpathmoveto{\pgfqpoint{1.437588in}{1.161289in}}%
\pgfpathlineto{\pgfqpoint{1.437588in}{1.150193in}}%
\pgfpathquadraticcurveto{\pgfqpoint{1.431122in}{1.153291in}}{\pgfqpoint{1.425376in}{1.154804in}}%
\pgfpathquadraticcurveto{\pgfqpoint{1.419630in}{1.156336in}}{\pgfqpoint{1.414281in}{1.156336in}}%
\pgfpathquadraticcurveto{\pgfqpoint{1.405005in}{1.156336in}}{\pgfqpoint{1.399961in}{1.152733in}}%
\pgfpathquadraticcurveto{\pgfqpoint{1.394918in}{1.149131in}}{\pgfqpoint{1.394918in}{1.142484in}}%
\pgfpathquadraticcurveto{\pgfqpoint{1.394918in}{1.136900in}}{\pgfqpoint{1.398268in}{1.134054in}}%
\pgfpathquadraticcurveto{\pgfqpoint{1.401618in}{1.131227in}}{\pgfqpoint{1.410967in}{1.129479in}}%
\pgfpathlineto{\pgfqpoint{1.417829in}{1.128074in}}%
\pgfpathquadraticcurveto{\pgfqpoint{1.430546in}{1.125643in}}{\pgfqpoint{1.436598in}{1.119537in}}%
\pgfpathquadraticcurveto{\pgfqpoint{1.442650in}{1.113431in}}{\pgfqpoint{1.442650in}{1.103200in}}%
\pgfpathquadraticcurveto{\pgfqpoint{1.442650in}{1.090969in}}{\pgfqpoint{1.434454in}{1.084665in}}%
\pgfpathquadraticcurveto{\pgfqpoint{1.426277in}{1.078361in}}{\pgfqpoint{1.410462in}{1.078361in}}%
\pgfpathquadraticcurveto{\pgfqpoint{1.404500in}{1.078361in}}{\pgfqpoint{1.397764in}{1.079712in}}%
\pgfpathquadraticcurveto{\pgfqpoint{1.391045in}{1.081063in}}{\pgfqpoint{1.383840in}{1.083711in}}%
\pgfpathlineto{\pgfqpoint{1.383840in}{1.095418in}}%
\pgfpathquadraticcurveto{\pgfqpoint{1.390757in}{1.091546in}}{\pgfqpoint{1.397403in}{1.089564in}}%
\pgfpathquadraticcurveto{\pgfqpoint{1.404050in}{1.087601in}}{\pgfqpoint{1.410462in}{1.087601in}}%
\pgfpathquadraticcurveto{\pgfqpoint{1.420189in}{1.087601in}}{\pgfqpoint{1.425484in}{1.091420in}}%
\pgfpathquadraticcurveto{\pgfqpoint{1.430780in}{1.095256in}}{\pgfqpoint{1.430780in}{1.102353in}}%
\pgfpathquadraticcurveto{\pgfqpoint{1.430780in}{1.108531in}}{\pgfqpoint{1.426979in}{1.112026in}}%
\pgfpathquadraticcurveto{\pgfqpoint{1.423179in}{1.115520in}}{\pgfqpoint{1.414515in}{1.117267in}}%
\pgfpathlineto{\pgfqpoint{1.407580in}{1.118618in}}%
\pgfpathquadraticcurveto{\pgfqpoint{1.394864in}{1.121140in}}{\pgfqpoint{1.389172in}{1.126543in}}%
\pgfpathquadraticcurveto{\pgfqpoint{1.383498in}{1.131947in}}{\pgfqpoint{1.383498in}{1.141584in}}%
\pgfpathquadraticcurveto{\pgfqpoint{1.383498in}{1.152733in}}{\pgfqpoint{1.391351in}{1.159145in}}%
\pgfpathquadraticcurveto{\pgfqpoint{1.399205in}{1.165558in}}{\pgfqpoint{1.412984in}{1.165558in}}%
\pgfpathquadraticcurveto{\pgfqpoint{1.418910in}{1.165558in}}{\pgfqpoint{1.425034in}{1.164477in}}%
\pgfpathquadraticcurveto{\pgfqpoint{1.431176in}{1.163414in}}{\pgfqpoint{1.437588in}{1.161289in}}%
\pgfpathlineto{\pgfqpoint{1.437588in}{1.161289in}}%
\pgfpathclose%
\pgfpathmoveto{\pgfqpoint{1.459942in}{1.143043in}}%
\pgfpathlineto{\pgfqpoint{1.470299in}{1.143043in}}%
\pgfpathlineto{\pgfqpoint{1.470299in}{1.080000in}}%
\pgfpathlineto{\pgfqpoint{1.459942in}{1.080000in}}%
\pgfpathlineto{\pgfqpoint{1.459942in}{1.143043in}}%
\pgfpathlineto{\pgfqpoint{1.459942in}{1.143043in}}%
\pgfpathclose%
\pgfpathmoveto{\pgfqpoint{1.459942in}{1.167593in}}%
\pgfpathlineto{\pgfqpoint{1.470299in}{1.167593in}}%
\pgfpathlineto{\pgfqpoint{1.470299in}{1.154462in}}%
\pgfpathlineto{\pgfqpoint{1.459942in}{1.154462in}}%
\pgfpathlineto{\pgfqpoint{1.459942in}{1.167593in}}%
\pgfpathlineto{\pgfqpoint{1.459942in}{1.167593in}}%
\pgfpathclose%
\pgfpathmoveto{\pgfqpoint{1.533449in}{1.112260in}}%
\pgfpathquadraticcurveto{\pgfqpoint{1.533449in}{1.123517in}}{\pgfqpoint{1.528802in}{1.129696in}}%
\pgfpathquadraticcurveto{\pgfqpoint{1.524173in}{1.135892in}}{\pgfqpoint{1.515779in}{1.135892in}}%
\pgfpathquadraticcurveto{\pgfqpoint{1.507458in}{1.135892in}}{\pgfqpoint{1.502811in}{1.129696in}}%
\pgfpathquadraticcurveto{\pgfqpoint{1.498163in}{1.123517in}}{\pgfqpoint{1.498163in}{1.112260in}}%
\pgfpathquadraticcurveto{\pgfqpoint{1.498163in}{1.101056in}}{\pgfqpoint{1.502811in}{1.094860in}}%
\pgfpathquadraticcurveto{\pgfqpoint{1.507458in}{1.088664in}}{\pgfqpoint{1.515779in}{1.088664in}}%
\pgfpathquadraticcurveto{\pgfqpoint{1.524173in}{1.088664in}}{\pgfqpoint{1.528802in}{1.094860in}}%
\pgfpathquadraticcurveto{\pgfqpoint{1.533449in}{1.101056in}}{\pgfqpoint{1.533449in}{1.112260in}}%
\pgfpathlineto{\pgfqpoint{1.533449in}{1.112260in}}%
\pgfpathclose%
\pgfpathmoveto{\pgfqpoint{1.543806in}{1.087817in}}%
\pgfpathquadraticcurveto{\pgfqpoint{1.543806in}{1.071732in}}{\pgfqpoint{1.536655in}{1.063879in}}%
\pgfpathquadraticcurveto{\pgfqpoint{1.529523in}{1.056026in}}{\pgfqpoint{1.514771in}{1.056026in}}%
\pgfpathquadraticcurveto{\pgfqpoint{1.509313in}{1.056026in}}{\pgfqpoint{1.504468in}{1.056836in}}%
\pgfpathquadraticcurveto{\pgfqpoint{1.499622in}{1.057647in}}{\pgfqpoint{1.495065in}{1.059340in}}%
\pgfpathlineto{\pgfqpoint{1.495065in}{1.069409in}}%
\pgfpathquadraticcurveto{\pgfqpoint{1.499622in}{1.066941in}}{\pgfqpoint{1.504071in}{1.065770in}}%
\pgfpathquadraticcurveto{\pgfqpoint{1.508520in}{1.064582in}}{\pgfqpoint{1.513131in}{1.064582in}}%
\pgfpathquadraticcurveto{\pgfqpoint{1.523326in}{1.064582in}}{\pgfqpoint{1.528388in}{1.069895in}}%
\pgfpathquadraticcurveto{\pgfqpoint{1.533449in}{1.075209in}}{\pgfqpoint{1.533449in}{1.085962in}}%
\pgfpathlineto{\pgfqpoint{1.533449in}{1.091095in}}%
\pgfpathquadraticcurveto{\pgfqpoint{1.530243in}{1.085512in}}{\pgfqpoint{1.525236in}{1.082756in}}%
\pgfpathquadraticcurveto{\pgfqpoint{1.520228in}{1.080000in}}{\pgfqpoint{1.513240in}{1.080000in}}%
\pgfpathquadraticcurveto{\pgfqpoint{1.501658in}{1.080000in}}{\pgfqpoint{1.494561in}{1.088826in}}%
\pgfpathquadraticcurveto{\pgfqpoint{1.487464in}{1.097670in}}{\pgfqpoint{1.487464in}{1.112260in}}%
\pgfpathquadraticcurveto{\pgfqpoint{1.487464in}{1.126886in}}{\pgfqpoint{1.494561in}{1.135712in}}%
\pgfpathquadraticcurveto{\pgfqpoint{1.501658in}{1.144556in}}{\pgfqpoint{1.513240in}{1.144556in}}%
\pgfpathquadraticcurveto{\pgfqpoint{1.520228in}{1.144556in}}{\pgfqpoint{1.525236in}{1.141800in}}%
\pgfpathquadraticcurveto{\pgfqpoint{1.530243in}{1.139044in}}{\pgfqpoint{1.533449in}{1.133478in}}%
\pgfpathlineto{\pgfqpoint{1.533449in}{1.143043in}}%
\pgfpathlineto{\pgfqpoint{1.543806in}{1.143043in}}%
\pgfpathlineto{\pgfqpoint{1.543806in}{1.087817in}}%
\pgfpathlineto{\pgfqpoint{1.543806in}{1.087817in}}%
\pgfpathclose%
\pgfpathmoveto{\pgfqpoint{1.617566in}{1.118060in}}%
\pgfpathlineto{\pgfqpoint{1.617566in}{1.080000in}}%
\pgfpathlineto{\pgfqpoint{1.607209in}{1.080000in}}%
\pgfpathlineto{\pgfqpoint{1.607209in}{1.117717in}}%
\pgfpathquadraticcurveto{\pgfqpoint{1.607209in}{1.126669in}}{\pgfqpoint{1.603715in}{1.131100in}}%
\pgfpathquadraticcurveto{\pgfqpoint{1.600220in}{1.135549in}}{\pgfqpoint{1.593250in}{1.135549in}}%
\pgfpathquadraticcurveto{\pgfqpoint{1.584856in}{1.135549in}}{\pgfqpoint{1.580011in}{1.130200in}}%
\pgfpathquadraticcurveto{\pgfqpoint{1.575165in}{1.124868in}}{\pgfqpoint{1.575165in}{1.115628in}}%
\pgfpathlineto{\pgfqpoint{1.575165in}{1.080000in}}%
\pgfpathlineto{\pgfqpoint{1.564754in}{1.080000in}}%
\pgfpathlineto{\pgfqpoint{1.564754in}{1.143043in}}%
\pgfpathlineto{\pgfqpoint{1.575165in}{1.143043in}}%
\pgfpathlineto{\pgfqpoint{1.575165in}{1.133244in}}%
\pgfpathquadraticcurveto{\pgfqpoint{1.578894in}{1.138936in}}{\pgfqpoint{1.583919in}{1.141746in}}%
\pgfpathquadraticcurveto{\pgfqpoint{1.588963in}{1.144556in}}{\pgfqpoint{1.595555in}{1.144556in}}%
\pgfpathquadraticcurveto{\pgfqpoint{1.606416in}{1.144556in}}{\pgfqpoint{1.611982in}{1.137837in}}%
\pgfpathquadraticcurveto{\pgfqpoint{1.617566in}{1.131118in}}{\pgfqpoint{1.617566in}{1.118060in}}%
\pgfpathlineto{\pgfqpoint{1.617566in}{1.118060in}}%
\pgfpathclose%
\pgfpathmoveto{\pgfqpoint{1.638208in}{1.143043in}}%
\pgfpathlineto{\pgfqpoint{1.648565in}{1.143043in}}%
\pgfpathlineto{\pgfqpoint{1.648565in}{1.080000in}}%
\pgfpathlineto{\pgfqpoint{1.638208in}{1.080000in}}%
\pgfpathlineto{\pgfqpoint{1.638208in}{1.143043in}}%
\pgfpathlineto{\pgfqpoint{1.638208in}{1.143043in}}%
\pgfpathclose%
\pgfpathmoveto{\pgfqpoint{1.638208in}{1.167593in}}%
\pgfpathlineto{\pgfqpoint{1.648565in}{1.167593in}}%
\pgfpathlineto{\pgfqpoint{1.648565in}{1.154462in}}%
\pgfpathlineto{\pgfqpoint{1.638208in}{1.154462in}}%
\pgfpathlineto{\pgfqpoint{1.638208in}{1.167593in}}%
\pgfpathlineto{\pgfqpoint{1.638208in}{1.167593in}}%
\pgfpathclose%
\pgfpathmoveto{\pgfqpoint{1.721172in}{1.143043in}}%
\pgfpathlineto{\pgfqpoint{1.721172in}{1.080000in}}%
\pgfpathlineto{\pgfqpoint{1.710761in}{1.080000in}}%
\pgfpathlineto{\pgfqpoint{1.710761in}{1.134991in}}%
\pgfpathlineto{\pgfqpoint{1.682338in}{1.134991in}}%
\pgfpathlineto{\pgfqpoint{1.682338in}{1.080000in}}%
\pgfpathlineto{\pgfqpoint{1.671927in}{1.080000in}}%
\pgfpathlineto{\pgfqpoint{1.671927in}{1.134991in}}%
\pgfpathlineto{\pgfqpoint{1.662020in}{1.134991in}}%
\pgfpathlineto{\pgfqpoint{1.662020in}{1.143043in}}%
\pgfpathlineto{\pgfqpoint{1.671927in}{1.143043in}}%
\pgfpathlineto{\pgfqpoint{1.671927in}{1.147438in}}%
\pgfpathquadraticcurveto{\pgfqpoint{1.671927in}{1.157740in}}{\pgfqpoint{1.676790in}{1.162658in}}%
\pgfpathquadraticcurveto{\pgfqpoint{1.681671in}{1.167593in}}{\pgfqpoint{1.691740in}{1.167593in}}%
\pgfpathlineto{\pgfqpoint{1.702151in}{1.167593in}}%
\pgfpathlineto{\pgfqpoint{1.702151in}{1.158965in}}%
\pgfpathlineto{\pgfqpoint{1.692244in}{1.158965in}}%
\pgfpathquadraticcurveto{\pgfqpoint{1.686679in}{1.158965in}}{\pgfqpoint{1.684499in}{1.156714in}}%
\pgfpathquadraticcurveto{\pgfqpoint{1.682338in}{1.154462in}}{\pgfqpoint{1.682338in}{1.148608in}}%
\pgfpathlineto{\pgfqpoint{1.682338in}{1.143043in}}%
\pgfpathlineto{\pgfqpoint{1.721172in}{1.143043in}}%
\pgfpathlineto{\pgfqpoint{1.721172in}{1.143043in}}%
\pgfpathclose%
\pgfpathmoveto{\pgfqpoint{1.710761in}{1.167467in}}%
\pgfpathlineto{\pgfqpoint{1.721172in}{1.167467in}}%
\pgfpathlineto{\pgfqpoint{1.721172in}{1.154354in}}%
\pgfpathlineto{\pgfqpoint{1.710761in}{1.154354in}}%
\pgfpathlineto{\pgfqpoint{1.710761in}{1.167467in}}%
\pgfpathlineto{\pgfqpoint{1.710761in}{1.167467in}}%
\pgfpathclose%
\pgfpathmoveto{\pgfqpoint{1.788213in}{1.140629in}}%
\pgfpathlineto{\pgfqpoint{1.788213in}{1.130938in}}%
\pgfpathquadraticcurveto{\pgfqpoint{1.783818in}{1.133370in}}{\pgfqpoint{1.779405in}{1.134577in}}%
\pgfpathquadraticcurveto{\pgfqpoint{1.774992in}{1.135784in}}{\pgfqpoint{1.770489in}{1.135784in}}%
\pgfpathquadraticcurveto{\pgfqpoint{1.760402in}{1.135784in}}{\pgfqpoint{1.754819in}{1.129389in}}%
\pgfpathquadraticcurveto{\pgfqpoint{1.749253in}{1.123013in}}{\pgfqpoint{1.749253in}{1.111467in}}%
\pgfpathquadraticcurveto{\pgfqpoint{1.749253in}{1.099921in}}{\pgfqpoint{1.754819in}{1.093527in}}%
\pgfpathquadraticcurveto{\pgfqpoint{1.760402in}{1.087151in}}{\pgfqpoint{1.770489in}{1.087151in}}%
\pgfpathquadraticcurveto{\pgfqpoint{1.774992in}{1.087151in}}{\pgfqpoint{1.779405in}{1.088358in}}%
\pgfpathquadraticcurveto{\pgfqpoint{1.783818in}{1.089564in}}{\pgfqpoint{1.788213in}{1.091996in}}%
\pgfpathlineto{\pgfqpoint{1.788213in}{1.082414in}}%
\pgfpathquadraticcurveto{\pgfqpoint{1.783872in}{1.080396in}}{\pgfqpoint{1.779225in}{1.079388in}}%
\pgfpathquadraticcurveto{\pgfqpoint{1.774596in}{1.078361in}}{\pgfqpoint{1.769354in}{1.078361in}}%
\pgfpathquadraticcurveto{\pgfqpoint{1.755107in}{1.078361in}}{\pgfqpoint{1.746713in}{1.087313in}}%
\pgfpathquadraticcurveto{\pgfqpoint{1.738337in}{1.096265in}}{\pgfqpoint{1.738337in}{1.111467in}}%
\pgfpathquadraticcurveto{\pgfqpoint{1.738337in}{1.126886in}}{\pgfqpoint{1.746803in}{1.135712in}}%
\pgfpathquadraticcurveto{\pgfqpoint{1.755287in}{1.144556in}}{\pgfqpoint{1.770039in}{1.144556in}}%
\pgfpathquadraticcurveto{\pgfqpoint{1.774812in}{1.144556in}}{\pgfqpoint{1.779369in}{1.143565in}}%
\pgfpathquadraticcurveto{\pgfqpoint{1.783926in}{1.142592in}}{\pgfqpoint{1.788213in}{1.140629in}}%
\pgfpathlineto{\pgfqpoint{1.788213in}{1.140629in}}%
\pgfpathclose%
\pgfpathmoveto{\pgfqpoint{1.834883in}{1.111683in}}%
\pgfpathquadraticcurveto{\pgfqpoint{1.822328in}{1.111683in}}{\pgfqpoint{1.817483in}{1.108819in}}%
\pgfpathquadraticcurveto{\pgfqpoint{1.812638in}{1.105956in}}{\pgfqpoint{1.812638in}{1.099021in}}%
\pgfpathquadraticcurveto{\pgfqpoint{1.812638in}{1.093509in}}{\pgfqpoint{1.816276in}{1.090267in}}%
\pgfpathquadraticcurveto{\pgfqpoint{1.819914in}{1.087043in}}{\pgfqpoint{1.826147in}{1.087043in}}%
\pgfpathquadraticcurveto{\pgfqpoint{1.834774in}{1.087043in}}{\pgfqpoint{1.839980in}{1.093149in}}%
\pgfpathquadraticcurveto{\pgfqpoint{1.845186in}{1.099255in}}{\pgfqpoint{1.845186in}{1.109378in}}%
\pgfpathlineto{\pgfqpoint{1.845186in}{1.111683in}}%
\pgfpathlineto{\pgfqpoint{1.834883in}{1.111683in}}%
\pgfpathlineto{\pgfqpoint{1.834883in}{1.111683in}}%
\pgfpathclose%
\pgfpathmoveto{\pgfqpoint{1.855542in}{1.115970in}}%
\pgfpathlineto{\pgfqpoint{1.855542in}{1.080000in}}%
\pgfpathlineto{\pgfqpoint{1.845186in}{1.080000in}}%
\pgfpathlineto{\pgfqpoint{1.845186in}{1.089564in}}%
\pgfpathquadraticcurveto{\pgfqpoint{1.841637in}{1.083837in}}{\pgfqpoint{1.836342in}{1.081099in}}%
\pgfpathquadraticcurveto{\pgfqpoint{1.831046in}{1.078361in}}{\pgfqpoint{1.823391in}{1.078361in}}%
\pgfpathquadraticcurveto{\pgfqpoint{1.813718in}{1.078361in}}{\pgfqpoint{1.807990in}{1.083801in}}%
\pgfpathquadraticcurveto{\pgfqpoint{1.802281in}{1.089240in}}{\pgfqpoint{1.802281in}{1.098354in}}%
\pgfpathquadraticcurveto{\pgfqpoint{1.802281in}{1.108982in}}{\pgfqpoint{1.809395in}{1.114385in}}%
\pgfpathquadraticcurveto{\pgfqpoint{1.816528in}{1.119789in}}{\pgfqpoint{1.830650in}{1.119789in}}%
\pgfpathlineto{\pgfqpoint{1.845186in}{1.119789in}}%
\pgfpathlineto{\pgfqpoint{1.845186in}{1.120816in}}%
\pgfpathquadraticcurveto{\pgfqpoint{1.845186in}{1.127966in}}{\pgfqpoint{1.840484in}{1.131875in}}%
\pgfpathquadraticcurveto{\pgfqpoint{1.835783in}{1.135784in}}{\pgfqpoint{1.827281in}{1.135784in}}%
\pgfpathquadraticcurveto{\pgfqpoint{1.821878in}{1.135784in}}{\pgfqpoint{1.816744in}{1.134487in}}%
\pgfpathquadraticcurveto{\pgfqpoint{1.811629in}{1.133190in}}{\pgfqpoint{1.806910in}{1.130596in}}%
\pgfpathlineto{\pgfqpoint{1.806910in}{1.140179in}}%
\pgfpathquadraticcurveto{\pgfqpoint{1.812584in}{1.142376in}}{\pgfqpoint{1.817933in}{1.143457in}}%
\pgfpathquadraticcurveto{\pgfqpoint{1.823283in}{1.144556in}}{\pgfqpoint{1.828344in}{1.144556in}}%
\pgfpathquadraticcurveto{\pgfqpoint{1.842033in}{1.144556in}}{\pgfqpoint{1.848788in}{1.137459in}}%
\pgfpathquadraticcurveto{\pgfqpoint{1.855542in}{1.130380in}}{\pgfqpoint{1.855542in}{1.115970in}}%
\pgfpathlineto{\pgfqpoint{1.855542in}{1.115970in}}%
\pgfpathclose%
\pgfpathmoveto{\pgfqpoint{1.929284in}{1.118060in}}%
\pgfpathlineto{\pgfqpoint{1.929284in}{1.080000in}}%
\pgfpathlineto{\pgfqpoint{1.918927in}{1.080000in}}%
\pgfpathlineto{\pgfqpoint{1.918927in}{1.117717in}}%
\pgfpathquadraticcurveto{\pgfqpoint{1.918927in}{1.126669in}}{\pgfqpoint{1.915433in}{1.131100in}}%
\pgfpathquadraticcurveto{\pgfqpoint{1.911939in}{1.135549in}}{\pgfqpoint{1.904968in}{1.135549in}}%
\pgfpathquadraticcurveto{\pgfqpoint{1.896574in}{1.135549in}}{\pgfqpoint{1.891729in}{1.130200in}}%
\pgfpathquadraticcurveto{\pgfqpoint{1.886884in}{1.124868in}}{\pgfqpoint{1.886884in}{1.115628in}}%
\pgfpathlineto{\pgfqpoint{1.886884in}{1.080000in}}%
\pgfpathlineto{\pgfqpoint{1.876473in}{1.080000in}}%
\pgfpathlineto{\pgfqpoint{1.876473in}{1.143043in}}%
\pgfpathlineto{\pgfqpoint{1.886884in}{1.143043in}}%
\pgfpathlineto{\pgfqpoint{1.886884in}{1.133244in}}%
\pgfpathquadraticcurveto{\pgfqpoint{1.890612in}{1.138936in}}{\pgfqpoint{1.895638in}{1.141746in}}%
\pgfpathquadraticcurveto{\pgfqpoint{1.900681in}{1.144556in}}{\pgfqpoint{1.907273in}{1.144556in}}%
\pgfpathquadraticcurveto{\pgfqpoint{1.918135in}{1.144556in}}{\pgfqpoint{1.923700in}{1.137837in}}%
\pgfpathquadraticcurveto{\pgfqpoint{1.929284in}{1.131118in}}{\pgfqpoint{1.929284in}{1.118060in}}%
\pgfpathlineto{\pgfqpoint{1.929284in}{1.118060in}}%
\pgfpathclose%
\pgfpathmoveto{\pgfqpoint{1.960175in}{1.160947in}}%
\pgfpathlineto{\pgfqpoint{1.960175in}{1.143043in}}%
\pgfpathlineto{\pgfqpoint{1.981501in}{1.143043in}}%
\pgfpathlineto{\pgfqpoint{1.981501in}{1.134991in}}%
\pgfpathlineto{\pgfqpoint{1.960175in}{1.134991in}}%
\pgfpathlineto{\pgfqpoint{1.960175in}{1.100768in}}%
\pgfpathquadraticcurveto{\pgfqpoint{1.960175in}{1.093059in}}{\pgfqpoint{1.962283in}{1.090861in}}%
\pgfpathquadraticcurveto{\pgfqpoint{1.964390in}{1.088664in}}{\pgfqpoint{1.970874in}{1.088664in}}%
\pgfpathlineto{\pgfqpoint{1.981501in}{1.088664in}}%
\pgfpathlineto{\pgfqpoint{1.981501in}{1.080000in}}%
\pgfpathlineto{\pgfqpoint{1.970874in}{1.080000in}}%
\pgfpathquadraticcurveto{\pgfqpoint{1.958878in}{1.080000in}}{\pgfqpoint{1.954321in}{1.084467in}}%
\pgfpathquadraticcurveto{\pgfqpoint{1.949764in}{1.088952in}}{\pgfqpoint{1.949764in}{1.100768in}}%
\pgfpathlineto{\pgfqpoint{1.949764in}{1.134991in}}%
\pgfpathlineto{\pgfqpoint{1.942163in}{1.134991in}}%
\pgfpathlineto{\pgfqpoint{1.942163in}{1.143043in}}%
\pgfpathlineto{\pgfqpoint{1.949764in}{1.143043in}}%
\pgfpathlineto{\pgfqpoint{1.949764in}{1.160947in}}%
\pgfpathlineto{\pgfqpoint{1.960175in}{1.160947in}}%
\pgfpathlineto{\pgfqpoint{1.960175in}{1.160947in}}%
\pgfpathclose%
\pgfusepath{fill}%
\end{pgfscope}%
\begin{pgfscope}%
\pgfsetbuttcap%
\pgfsetroundjoin%
\definecolor{currentfill}{rgb}{0.309804,0.372549,0.419608}%
\pgfsetfillcolor{currentfill}%
\pgfsetlinewidth{0.000000pt}%
\definecolor{currentstroke}{rgb}{0.309804,0.372549,0.419608}%
\pgfsetstrokecolor{currentstroke}%
\pgfsetdash{}{0pt}%
\pgfpathmoveto{\pgfqpoint{2.120822in}{1.152283in}}%
\pgfpathlineto{\pgfqpoint{2.120822in}{1.120924in}}%
\pgfpathlineto{\pgfqpoint{2.152163in}{1.120924in}}%
\pgfpathlineto{\pgfqpoint{2.152163in}{1.111359in}}%
\pgfpathlineto{\pgfqpoint{2.120822in}{1.111359in}}%
\pgfpathlineto{\pgfqpoint{2.120822in}{1.080000in}}%
\pgfpathlineto{\pgfqpoint{2.111365in}{1.080000in}}%
\pgfpathlineto{\pgfqpoint{2.111365in}{1.111359in}}%
\pgfpathlineto{\pgfqpoint{2.080006in}{1.111359in}}%
\pgfpathlineto{\pgfqpoint{2.080006in}{1.120924in}}%
\pgfpathlineto{\pgfqpoint{2.111365in}{1.120924in}}%
\pgfpathlineto{\pgfqpoint{2.111365in}{1.152283in}}%
\pgfpathlineto{\pgfqpoint{2.120822in}{1.152283in}}%
\pgfpathlineto{\pgfqpoint{2.120822in}{1.152283in}}%
\pgfpathclose%
\pgfusepath{fill}%
\end{pgfscope}%
\begin{pgfscope}%
\pgfsetbuttcap%
\pgfsetroundjoin%
\definecolor{currentfill}{rgb}{0.309804,0.372549,0.419608}%
\pgfsetfillcolor{currentfill}%
\pgfsetlinewidth{0.000000pt}%
\definecolor{currentstroke}{rgb}{0.309804,0.372549,0.419608}%
\pgfsetstrokecolor{currentstroke}%
\pgfsetdash{}{0pt}%
\pgfpathmoveto{\pgfqpoint{2.260863in}{1.119411in}}%
\pgfpathquadraticcurveto{\pgfqpoint{2.264520in}{1.118168in}}{\pgfqpoint{2.267978in}{1.114115in}}%
\pgfpathquadraticcurveto{\pgfqpoint{2.271436in}{1.110062in}}{\pgfqpoint{2.274931in}{1.102965in}}%
\pgfpathlineto{\pgfqpoint{2.286476in}{1.080000in}}%
\pgfpathlineto{\pgfqpoint{2.274246in}{1.080000in}}%
\pgfpathlineto{\pgfqpoint{2.263511in}{1.101561in}}%
\pgfpathquadraticcurveto{\pgfqpoint{2.259332in}{1.110008in}}{\pgfqpoint{2.255423in}{1.112764in}}%
\pgfpathquadraticcurveto{\pgfqpoint{2.251515in}{1.115520in}}{\pgfqpoint{2.244760in}{1.115520in}}%
\pgfpathlineto{\pgfqpoint{2.232368in}{1.115520in}}%
\pgfpathlineto{\pgfqpoint{2.232368in}{1.080000in}}%
\pgfpathlineto{\pgfqpoint{2.221002in}{1.080000in}}%
\pgfpathlineto{\pgfqpoint{2.221002in}{1.164045in}}%
\pgfpathlineto{\pgfqpoint{2.246670in}{1.164045in}}%
\pgfpathquadraticcurveto{\pgfqpoint{2.261079in}{1.164045in}}{\pgfqpoint{2.268176in}{1.158011in}}%
\pgfpathquadraticcurveto{\pgfqpoint{2.275273in}{1.151995in}}{\pgfqpoint{2.275273in}{1.139836in}}%
\pgfpathquadraticcurveto{\pgfqpoint{2.275273in}{1.131893in}}{\pgfqpoint{2.271580in}{1.126651in}}%
\pgfpathquadraticcurveto{\pgfqpoint{2.267888in}{1.121428in}}{\pgfqpoint{2.260863in}{1.119411in}}%
\pgfpathlineto{\pgfqpoint{2.260863in}{1.119411in}}%
\pgfpathclose%
\pgfpathmoveto{\pgfqpoint{2.232368in}{1.154696in}}%
\pgfpathlineto{\pgfqpoint{2.232368in}{1.124868in}}%
\pgfpathlineto{\pgfqpoint{2.246670in}{1.124868in}}%
\pgfpathquadraticcurveto{\pgfqpoint{2.254883in}{1.124868in}}{\pgfqpoint{2.259080in}{1.128669in}}%
\pgfpathquadraticcurveto{\pgfqpoint{2.263277in}{1.132469in}}{\pgfqpoint{2.263277in}{1.139836in}}%
\pgfpathquadraticcurveto{\pgfqpoint{2.263277in}{1.147203in}}{\pgfqpoint{2.259080in}{1.150950in}}%
\pgfpathquadraticcurveto{\pgfqpoint{2.254883in}{1.154696in}}{\pgfqpoint{2.246670in}{1.154696in}}%
\pgfpathlineto{\pgfqpoint{2.232368in}{1.154696in}}%
\pgfpathlineto{\pgfqpoint{2.232368in}{1.154696in}}%
\pgfpathclose%
\pgfpathmoveto{\pgfqpoint{2.349411in}{1.114115in}}%
\pgfpathlineto{\pgfqpoint{2.349411in}{1.109054in}}%
\pgfpathlineto{\pgfqpoint{2.301787in}{1.109054in}}%
\pgfpathquadraticcurveto{\pgfqpoint{2.302471in}{1.098354in}}{\pgfqpoint{2.308235in}{1.092753in}}%
\pgfpathquadraticcurveto{\pgfqpoint{2.313999in}{1.087151in}}{\pgfqpoint{2.324302in}{1.087151in}}%
\pgfpathquadraticcurveto{\pgfqpoint{2.330264in}{1.087151in}}{\pgfqpoint{2.335866in}{1.088610in}}%
\pgfpathquadraticcurveto{\pgfqpoint{2.341468in}{1.090069in}}{\pgfqpoint{2.346997in}{1.093005in}}%
\pgfpathlineto{\pgfqpoint{2.346997in}{1.083206in}}%
\pgfpathquadraticcurveto{\pgfqpoint{2.341413in}{1.080847in}}{\pgfqpoint{2.335560in}{1.079604in}}%
\pgfpathquadraticcurveto{\pgfqpoint{2.329706in}{1.078361in}}{\pgfqpoint{2.323690in}{1.078361in}}%
\pgfpathquadraticcurveto{\pgfqpoint{2.308595in}{1.078361in}}{\pgfqpoint{2.299787in}{1.087133in}}%
\pgfpathquadraticcurveto{\pgfqpoint{2.290979in}{1.095923in}}{\pgfqpoint{2.290979in}{1.110909in}}%
\pgfpathquadraticcurveto{\pgfqpoint{2.290979in}{1.126381in}}{\pgfqpoint{2.299337in}{1.135459in}}%
\pgfpathquadraticcurveto{\pgfqpoint{2.307695in}{1.144556in}}{\pgfqpoint{2.321888in}{1.144556in}}%
\pgfpathquadraticcurveto{\pgfqpoint{2.334605in}{1.144556in}}{\pgfqpoint{2.342008in}{1.136360in}}%
\pgfpathquadraticcurveto{\pgfqpoint{2.349411in}{1.128183in}}{\pgfqpoint{2.349411in}{1.114115in}}%
\pgfpathlineto{\pgfqpoint{2.349411in}{1.114115in}}%
\pgfpathclose%
\pgfpathmoveto{\pgfqpoint{2.339054in}{1.117159in}}%
\pgfpathquadraticcurveto{\pgfqpoint{2.338946in}{1.125643in}}{\pgfqpoint{2.334299in}{1.130704in}}%
\pgfpathquadraticcurveto{\pgfqpoint{2.329652in}{1.135784in}}{\pgfqpoint{2.321996in}{1.135784in}}%
\pgfpathquadraticcurveto{\pgfqpoint{2.313333in}{1.135784in}}{\pgfqpoint{2.308127in}{1.130884in}}%
\pgfpathquadraticcurveto{\pgfqpoint{2.302922in}{1.125985in}}{\pgfqpoint{2.302129in}{1.117087in}}%
\pgfpathlineto{\pgfqpoint{2.339054in}{1.117159in}}%
\pgfpathlineto{\pgfqpoint{2.339054in}{1.117159in}}%
\pgfpathclose%
\pgfpathmoveto{\pgfqpoint{2.358993in}{1.143043in}}%
\pgfpathlineto{\pgfqpoint{2.369963in}{1.143043in}}%
\pgfpathlineto{\pgfqpoint{2.389668in}{1.090141in}}%
\pgfpathlineto{\pgfqpoint{2.409373in}{1.143043in}}%
\pgfpathlineto{\pgfqpoint{2.420343in}{1.143043in}}%
\pgfpathlineto{\pgfqpoint{2.396693in}{1.080000in}}%
\pgfpathlineto{\pgfqpoint{2.382625in}{1.080000in}}%
\pgfpathlineto{\pgfqpoint{2.358993in}{1.143043in}}%
\pgfpathlineto{\pgfqpoint{2.358993in}{1.143043in}}%
\pgfpathclose%
\pgfpathmoveto{\pgfqpoint{2.488573in}{1.114115in}}%
\pgfpathlineto{\pgfqpoint{2.488573in}{1.109054in}}%
\pgfpathlineto{\pgfqpoint{2.440949in}{1.109054in}}%
\pgfpathquadraticcurveto{\pgfqpoint{2.441633in}{1.098354in}}{\pgfqpoint{2.447397in}{1.092753in}}%
\pgfpathquadraticcurveto{\pgfqpoint{2.453161in}{1.087151in}}{\pgfqpoint{2.463464in}{1.087151in}}%
\pgfpathquadraticcurveto{\pgfqpoint{2.469426in}{1.087151in}}{\pgfqpoint{2.475028in}{1.088610in}}%
\pgfpathquadraticcurveto{\pgfqpoint{2.480629in}{1.090069in}}{\pgfqpoint{2.486159in}{1.093005in}}%
\pgfpathlineto{\pgfqpoint{2.486159in}{1.083206in}}%
\pgfpathquadraticcurveto{\pgfqpoint{2.480575in}{1.080847in}}{\pgfqpoint{2.474721in}{1.079604in}}%
\pgfpathquadraticcurveto{\pgfqpoint{2.468867in}{1.078361in}}{\pgfqpoint{2.462851in}{1.078361in}}%
\pgfpathquadraticcurveto{\pgfqpoint{2.447757in}{1.078361in}}{\pgfqpoint{2.438949in}{1.087133in}}%
\pgfpathquadraticcurveto{\pgfqpoint{2.430141in}{1.095923in}}{\pgfqpoint{2.430141in}{1.110909in}}%
\pgfpathquadraticcurveto{\pgfqpoint{2.430141in}{1.126381in}}{\pgfqpoint{2.438499in}{1.135459in}}%
\pgfpathquadraticcurveto{\pgfqpoint{2.446857in}{1.144556in}}{\pgfqpoint{2.461050in}{1.144556in}}%
\pgfpathquadraticcurveto{\pgfqpoint{2.473767in}{1.144556in}}{\pgfqpoint{2.481170in}{1.136360in}}%
\pgfpathquadraticcurveto{\pgfqpoint{2.488573in}{1.128183in}}{\pgfqpoint{2.488573in}{1.114115in}}%
\pgfpathlineto{\pgfqpoint{2.488573in}{1.114115in}}%
\pgfpathclose%
\pgfpathmoveto{\pgfqpoint{2.478216in}{1.117159in}}%
\pgfpathquadraticcurveto{\pgfqpoint{2.478108in}{1.125643in}}{\pgfqpoint{2.473461in}{1.130704in}}%
\pgfpathquadraticcurveto{\pgfqpoint{2.468813in}{1.135784in}}{\pgfqpoint{2.461158in}{1.135784in}}%
\pgfpathquadraticcurveto{\pgfqpoint{2.452494in}{1.135784in}}{\pgfqpoint{2.447289in}{1.130884in}}%
\pgfpathquadraticcurveto{\pgfqpoint{2.442083in}{1.125985in}}{\pgfqpoint{2.441291in}{1.117087in}}%
\pgfpathlineto{\pgfqpoint{2.478216in}{1.117159in}}%
\pgfpathlineto{\pgfqpoint{2.478216in}{1.117159in}}%
\pgfpathclose%
\pgfpathmoveto{\pgfqpoint{2.542105in}{1.133370in}}%
\pgfpathquadraticcurveto{\pgfqpoint{2.540358in}{1.134379in}}{\pgfqpoint{2.538304in}{1.134847in}}%
\pgfpathquadraticcurveto{\pgfqpoint{2.536251in}{1.135333in}}{\pgfqpoint{2.533783in}{1.135333in}}%
\pgfpathquadraticcurveto{\pgfqpoint{2.524993in}{1.135333in}}{\pgfqpoint{2.520292in}{1.129623in}}%
\pgfpathquadraticcurveto{\pgfqpoint{2.515591in}{1.123914in}}{\pgfqpoint{2.515591in}{1.113214in}}%
\pgfpathlineto{\pgfqpoint{2.515591in}{1.080000in}}%
\pgfpathlineto{\pgfqpoint{2.505180in}{1.080000in}}%
\pgfpathlineto{\pgfqpoint{2.505180in}{1.143043in}}%
\pgfpathlineto{\pgfqpoint{2.515591in}{1.143043in}}%
\pgfpathlineto{\pgfqpoint{2.515591in}{1.133244in}}%
\pgfpathquadraticcurveto{\pgfqpoint{2.518869in}{1.138990in}}{\pgfqpoint{2.524093in}{1.141764in}}%
\pgfpathquadraticcurveto{\pgfqpoint{2.529334in}{1.144556in}}{\pgfqpoint{2.536827in}{1.144556in}}%
\pgfpathquadraticcurveto{\pgfqpoint{2.537890in}{1.144556in}}{\pgfqpoint{2.539187in}{1.144411in}}%
\pgfpathquadraticcurveto{\pgfqpoint{2.540484in}{1.144285in}}{\pgfqpoint{2.542051in}{1.143997in}}%
\pgfpathlineto{\pgfqpoint{2.542105in}{1.133370in}}%
\pgfpathlineto{\pgfqpoint{2.542105in}{1.133370in}}%
\pgfpathclose%
\pgfpathmoveto{\pgfqpoint{2.593151in}{1.141187in}}%
\pgfpathlineto{\pgfqpoint{2.593151in}{1.131389in}}%
\pgfpathquadraticcurveto{\pgfqpoint{2.588774in}{1.133640in}}{\pgfqpoint{2.584037in}{1.134757in}}%
\pgfpathquadraticcurveto{\pgfqpoint{2.579318in}{1.135892in}}{\pgfqpoint{2.574239in}{1.135892in}}%
\pgfpathquadraticcurveto{\pgfqpoint{2.566529in}{1.135892in}}{\pgfqpoint{2.562675in}{1.133532in}}%
\pgfpathquadraticcurveto{\pgfqpoint{2.558820in}{1.131173in}}{\pgfqpoint{2.558820in}{1.126435in}}%
\pgfpathquadraticcurveto{\pgfqpoint{2.558820in}{1.122833in}}{\pgfqpoint{2.561576in}{1.120780in}}%
\pgfpathquadraticcurveto{\pgfqpoint{2.564332in}{1.118726in}}{\pgfqpoint{2.572672in}{1.116871in}}%
\pgfpathlineto{\pgfqpoint{2.576220in}{1.116078in}}%
\pgfpathquadraticcurveto{\pgfqpoint{2.587243in}{1.113719in}}{\pgfqpoint{2.591890in}{1.109414in}}%
\pgfpathquadraticcurveto{\pgfqpoint{2.596538in}{1.105109in}}{\pgfqpoint{2.596538in}{1.097400in}}%
\pgfpathquadraticcurveto{\pgfqpoint{2.596538in}{1.088610in}}{\pgfqpoint{2.589585in}{1.083476in}}%
\pgfpathquadraticcurveto{\pgfqpoint{2.582632in}{1.078361in}}{\pgfqpoint{2.570474in}{1.078361in}}%
\pgfpathquadraticcurveto{\pgfqpoint{2.565413in}{1.078361in}}{\pgfqpoint{2.559919in}{1.079352in}}%
\pgfpathquadraticcurveto{\pgfqpoint{2.554425in}{1.080342in}}{\pgfqpoint{2.548355in}{1.082306in}}%
\pgfpathlineto{\pgfqpoint{2.548355in}{1.093005in}}%
\pgfpathquadraticcurveto{\pgfqpoint{2.554101in}{1.090015in}}{\pgfqpoint{2.559667in}{1.088520in}}%
\pgfpathquadraticcurveto{\pgfqpoint{2.565233in}{1.087043in}}{\pgfqpoint{2.570708in}{1.087043in}}%
\pgfpathquadraticcurveto{\pgfqpoint{2.578021in}{1.087043in}}{\pgfqpoint{2.581948in}{1.089546in}}%
\pgfpathquadraticcurveto{\pgfqpoint{2.585892in}{1.092050in}}{\pgfqpoint{2.585892in}{1.096607in}}%
\pgfpathquadraticcurveto{\pgfqpoint{2.585892in}{1.100822in}}{\pgfqpoint{2.583047in}{1.103074in}}%
\pgfpathquadraticcurveto{\pgfqpoint{2.580219in}{1.105325in}}{\pgfqpoint{2.570582in}{1.107414in}}%
\pgfpathlineto{\pgfqpoint{2.566980in}{1.108261in}}%
\pgfpathquadraticcurveto{\pgfqpoint{2.557361in}{1.110278in}}{\pgfqpoint{2.553074in}{1.114475in}}%
\pgfpathquadraticcurveto{\pgfqpoint{2.548805in}{1.118672in}}{\pgfqpoint{2.548805in}{1.125985in}}%
\pgfpathquadraticcurveto{\pgfqpoint{2.548805in}{1.134883in}}{\pgfqpoint{2.555110in}{1.139710in}}%
\pgfpathquadraticcurveto{\pgfqpoint{2.561414in}{1.144556in}}{\pgfqpoint{2.573014in}{1.144556in}}%
\pgfpathquadraticcurveto{\pgfqpoint{2.578742in}{1.144556in}}{\pgfqpoint{2.583803in}{1.143709in}}%
\pgfpathquadraticcurveto{\pgfqpoint{2.588882in}{1.142880in}}{\pgfqpoint{2.593151in}{1.141187in}}%
\pgfpathlineto{\pgfqpoint{2.593151in}{1.141187in}}%
\pgfpathclose%
\pgfpathmoveto{\pgfqpoint{2.641676in}{1.111683in}}%
\pgfpathquadraticcurveto{\pgfqpoint{2.629122in}{1.111683in}}{\pgfqpoint{2.624276in}{1.108819in}}%
\pgfpathquadraticcurveto{\pgfqpoint{2.619431in}{1.105956in}}{\pgfqpoint{2.619431in}{1.099021in}}%
\pgfpathquadraticcurveto{\pgfqpoint{2.619431in}{1.093509in}}{\pgfqpoint{2.623070in}{1.090267in}}%
\pgfpathquadraticcurveto{\pgfqpoint{2.626708in}{1.087043in}}{\pgfqpoint{2.632940in}{1.087043in}}%
\pgfpathquadraticcurveto{\pgfqpoint{2.641568in}{1.087043in}}{\pgfqpoint{2.646774in}{1.093149in}}%
\pgfpathquadraticcurveto{\pgfqpoint{2.651979in}{1.099255in}}{\pgfqpoint{2.651979in}{1.109378in}}%
\pgfpathlineto{\pgfqpoint{2.651979in}{1.111683in}}%
\pgfpathlineto{\pgfqpoint{2.641676in}{1.111683in}}%
\pgfpathlineto{\pgfqpoint{2.641676in}{1.111683in}}%
\pgfpathclose%
\pgfpathmoveto{\pgfqpoint{2.662336in}{1.115970in}}%
\pgfpathlineto{\pgfqpoint{2.662336in}{1.080000in}}%
\pgfpathlineto{\pgfqpoint{2.651979in}{1.080000in}}%
\pgfpathlineto{\pgfqpoint{2.651979in}{1.089564in}}%
\pgfpathquadraticcurveto{\pgfqpoint{2.648431in}{1.083837in}}{\pgfqpoint{2.643135in}{1.081099in}}%
\pgfpathquadraticcurveto{\pgfqpoint{2.637839in}{1.078361in}}{\pgfqpoint{2.630184in}{1.078361in}}%
\pgfpathquadraticcurveto{\pgfqpoint{2.620512in}{1.078361in}}{\pgfqpoint{2.614784in}{1.083801in}}%
\pgfpathquadraticcurveto{\pgfqpoint{2.609074in}{1.089240in}}{\pgfqpoint{2.609074in}{1.098354in}}%
\pgfpathquadraticcurveto{\pgfqpoint{2.609074in}{1.108982in}}{\pgfqpoint{2.616189in}{1.114385in}}%
\pgfpathquadraticcurveto{\pgfqpoint{2.623322in}{1.119789in}}{\pgfqpoint{2.637443in}{1.119789in}}%
\pgfpathlineto{\pgfqpoint{2.651979in}{1.119789in}}%
\pgfpathlineto{\pgfqpoint{2.651979in}{1.120816in}}%
\pgfpathquadraticcurveto{\pgfqpoint{2.651979in}{1.127966in}}{\pgfqpoint{2.647278in}{1.131875in}}%
\pgfpathquadraticcurveto{\pgfqpoint{2.642577in}{1.135784in}}{\pgfqpoint{2.634075in}{1.135784in}}%
\pgfpathquadraticcurveto{\pgfqpoint{2.628671in}{1.135784in}}{\pgfqpoint{2.623538in}{1.134487in}}%
\pgfpathquadraticcurveto{\pgfqpoint{2.618422in}{1.133190in}}{\pgfqpoint{2.613703in}{1.130596in}}%
\pgfpathlineto{\pgfqpoint{2.613703in}{1.140179in}}%
\pgfpathquadraticcurveto{\pgfqpoint{2.619377in}{1.142376in}}{\pgfqpoint{2.624727in}{1.143457in}}%
\pgfpathquadraticcurveto{\pgfqpoint{2.630076in}{1.144556in}}{\pgfqpoint{2.635138in}{1.144556in}}%
\pgfpathquadraticcurveto{\pgfqpoint{2.648827in}{1.144556in}}{\pgfqpoint{2.655581in}{1.137459in}}%
\pgfpathquadraticcurveto{\pgfqpoint{2.662336in}{1.130380in}}{\pgfqpoint{2.662336in}{1.115970in}}%
\pgfpathlineto{\pgfqpoint{2.662336in}{1.115970in}}%
\pgfpathclose%
\pgfpathmoveto{\pgfqpoint{2.683662in}{1.167593in}}%
\pgfpathlineto{\pgfqpoint{2.694019in}{1.167593in}}%
\pgfpathlineto{\pgfqpoint{2.694019in}{1.080000in}}%
\pgfpathlineto{\pgfqpoint{2.683662in}{1.080000in}}%
\pgfpathlineto{\pgfqpoint{2.683662in}{1.167593in}}%
\pgfpathlineto{\pgfqpoint{2.683662in}{1.167593in}}%
\pgfpathclose%
\pgfpathmoveto{\pgfqpoint{2.718336in}{1.094302in}}%
\pgfpathlineto{\pgfqpoint{2.730206in}{1.094302in}}%
\pgfpathlineto{\pgfqpoint{2.730206in}{1.080000in}}%
\pgfpathlineto{\pgfqpoint{2.718336in}{1.080000in}}%
\pgfpathlineto{\pgfqpoint{2.718336in}{1.094302in}}%
\pgfpathlineto{\pgfqpoint{2.718336in}{1.094302in}}%
\pgfpathclose%
\pgfpathmoveto{\pgfqpoint{2.718336in}{1.139602in}}%
\pgfpathlineto{\pgfqpoint{2.730206in}{1.139602in}}%
\pgfpathlineto{\pgfqpoint{2.730206in}{1.125319in}}%
\pgfpathlineto{\pgfqpoint{2.718336in}{1.125319in}}%
\pgfpathlineto{\pgfqpoint{2.718336in}{1.139602in}}%
\pgfpathlineto{\pgfqpoint{2.718336in}{1.139602in}}%
\pgfpathclose%
\pgfusepath{fill}%
\end{pgfscope}%
\begin{pgfscope}%
\pgfsetbuttcap%
\pgfsetroundjoin%
\definecolor{currentfill}{rgb}{0.309804,0.372549,0.419608}%
\pgfsetfillcolor{currentfill}%
\pgfsetlinewidth{0.000000pt}%
\definecolor{currentstroke}{rgb}{0.309804,0.372549,0.419608}%
\pgfsetstrokecolor{currentstroke}%
\pgfsetdash{}{0pt}%
\pgfpathmoveto{\pgfqpoint{2.821717in}{1.135784in}}%
\pgfpathquadraticcurveto{\pgfqpoint{2.813395in}{1.135784in}}{\pgfqpoint{2.808550in}{1.129281in}}%
\pgfpathquadraticcurveto{\pgfqpoint{2.803705in}{1.122779in}}{\pgfqpoint{2.803705in}{1.111467in}}%
\pgfpathquadraticcurveto{\pgfqpoint{2.803705in}{1.100156in}}{\pgfqpoint{2.808514in}{1.093653in}}%
\pgfpathquadraticcurveto{\pgfqpoint{2.813341in}{1.087151in}}{\pgfqpoint{2.821717in}{1.087151in}}%
\pgfpathquadraticcurveto{\pgfqpoint{2.830003in}{1.087151in}}{\pgfqpoint{2.834830in}{1.093671in}}%
\pgfpathquadraticcurveto{\pgfqpoint{2.839675in}{1.100210in}}{\pgfqpoint{2.839675in}{1.111467in}}%
\pgfpathquadraticcurveto{\pgfqpoint{2.839675in}{1.122671in}}{\pgfqpoint{2.834830in}{1.129227in}}%
\pgfpathquadraticcurveto{\pgfqpoint{2.830003in}{1.135784in}}{\pgfqpoint{2.821717in}{1.135784in}}%
\pgfpathlineto{\pgfqpoint{2.821717in}{1.135784in}}%
\pgfpathclose%
\pgfpathmoveto{\pgfqpoint{2.821717in}{1.144556in}}%
\pgfpathquadraticcurveto{\pgfqpoint{2.835226in}{1.144556in}}{\pgfqpoint{2.842935in}{1.135766in}}%
\pgfpathquadraticcurveto{\pgfqpoint{2.850663in}{1.126994in}}{\pgfqpoint{2.850663in}{1.111467in}}%
\pgfpathquadraticcurveto{\pgfqpoint{2.850663in}{1.095995in}}{\pgfqpoint{2.842935in}{1.087169in}}%
\pgfpathquadraticcurveto{\pgfqpoint{2.835226in}{1.078361in}}{\pgfqpoint{2.821717in}{1.078361in}}%
\pgfpathquadraticcurveto{\pgfqpoint{2.808154in}{1.078361in}}{\pgfqpoint{2.800463in}{1.087169in}}%
\pgfpathquadraticcurveto{\pgfqpoint{2.792790in}{1.095995in}}{\pgfqpoint{2.792790in}{1.111467in}}%
\pgfpathquadraticcurveto{\pgfqpoint{2.792790in}{1.126994in}}{\pgfqpoint{2.800463in}{1.135766in}}%
\pgfpathquadraticcurveto{\pgfqpoint{2.808154in}{1.144556in}}{\pgfqpoint{2.821717in}{1.144556in}}%
\pgfpathlineto{\pgfqpoint{2.821717in}{1.144556in}}%
\pgfpathclose%
\pgfpathmoveto{\pgfqpoint{2.877843in}{1.089456in}}%
\pgfpathlineto{\pgfqpoint{2.877843in}{1.056026in}}%
\pgfpathlineto{\pgfqpoint{2.867432in}{1.056026in}}%
\pgfpathlineto{\pgfqpoint{2.867432in}{1.143043in}}%
\pgfpathlineto{\pgfqpoint{2.877843in}{1.143043in}}%
\pgfpathlineto{\pgfqpoint{2.877843in}{1.133478in}}%
\pgfpathquadraticcurveto{\pgfqpoint{2.881121in}{1.139098in}}{\pgfqpoint{2.886092in}{1.141818in}}%
\pgfpathquadraticcurveto{\pgfqpoint{2.891082in}{1.144556in}}{\pgfqpoint{2.897999in}{1.144556in}}%
\pgfpathquadraticcurveto{\pgfqpoint{2.909490in}{1.144556in}}{\pgfqpoint{2.916659in}{1.135441in}}%
\pgfpathquadraticcurveto{\pgfqpoint{2.923846in}{1.126327in}}{\pgfqpoint{2.923846in}{1.111467in}}%
\pgfpathquadraticcurveto{\pgfqpoint{2.923846in}{1.096607in}}{\pgfqpoint{2.916659in}{1.087475in}}%
\pgfpathquadraticcurveto{\pgfqpoint{2.909490in}{1.078361in}}{\pgfqpoint{2.897999in}{1.078361in}}%
\pgfpathquadraticcurveto{\pgfqpoint{2.891082in}{1.078361in}}{\pgfqpoint{2.886092in}{1.081099in}}%
\pgfpathquadraticcurveto{\pgfqpoint{2.881121in}{1.083837in}}{\pgfqpoint{2.877843in}{1.089456in}}%
\pgfpathlineto{\pgfqpoint{2.877843in}{1.089456in}}%
\pgfpathclose%
\pgfpathmoveto{\pgfqpoint{2.913093in}{1.111467in}}%
\pgfpathquadraticcurveto{\pgfqpoint{2.913093in}{1.122887in}}{\pgfqpoint{2.908392in}{1.129389in}}%
\pgfpathquadraticcurveto{\pgfqpoint{2.903690in}{1.135892in}}{\pgfqpoint{2.895477in}{1.135892in}}%
\pgfpathquadraticcurveto{\pgfqpoint{2.887245in}{1.135892in}}{\pgfqpoint{2.882544in}{1.129389in}}%
\pgfpathquadraticcurveto{\pgfqpoint{2.877843in}{1.122887in}}{\pgfqpoint{2.877843in}{1.111467in}}%
\pgfpathquadraticcurveto{\pgfqpoint{2.877843in}{1.100048in}}{\pgfqpoint{2.882544in}{1.093545in}}%
\pgfpathquadraticcurveto{\pgfqpoint{2.887245in}{1.087043in}}{\pgfqpoint{2.895477in}{1.087043in}}%
\pgfpathquadraticcurveto{\pgfqpoint{2.903690in}{1.087043in}}{\pgfqpoint{2.908392in}{1.093545in}}%
\pgfpathquadraticcurveto{\pgfqpoint{2.913093in}{1.100048in}}{\pgfqpoint{2.913093in}{1.111467in}}%
\pgfpathlineto{\pgfqpoint{2.913093in}{1.111467in}}%
\pgfpathclose%
\pgfpathmoveto{\pgfqpoint{2.951026in}{1.089456in}}%
\pgfpathlineto{\pgfqpoint{2.951026in}{1.056026in}}%
\pgfpathlineto{\pgfqpoint{2.940615in}{1.056026in}}%
\pgfpathlineto{\pgfqpoint{2.940615in}{1.143043in}}%
\pgfpathlineto{\pgfqpoint{2.951026in}{1.143043in}}%
\pgfpathlineto{\pgfqpoint{2.951026in}{1.133478in}}%
\pgfpathquadraticcurveto{\pgfqpoint{2.954305in}{1.139098in}}{\pgfqpoint{2.959276in}{1.141818in}}%
\pgfpathquadraticcurveto{\pgfqpoint{2.964265in}{1.144556in}}{\pgfqpoint{2.971182in}{1.144556in}}%
\pgfpathquadraticcurveto{\pgfqpoint{2.982674in}{1.144556in}}{\pgfqpoint{2.989842in}{1.135441in}}%
\pgfpathquadraticcurveto{\pgfqpoint{2.997029in}{1.126327in}}{\pgfqpoint{2.997029in}{1.111467in}}%
\pgfpathquadraticcurveto{\pgfqpoint{2.997029in}{1.096607in}}{\pgfqpoint{2.989842in}{1.087475in}}%
\pgfpathquadraticcurveto{\pgfqpoint{2.982674in}{1.078361in}}{\pgfqpoint{2.971182in}{1.078361in}}%
\pgfpathquadraticcurveto{\pgfqpoint{2.964265in}{1.078361in}}{\pgfqpoint{2.959276in}{1.081099in}}%
\pgfpathquadraticcurveto{\pgfqpoint{2.954305in}{1.083837in}}{\pgfqpoint{2.951026in}{1.089456in}}%
\pgfpathlineto{\pgfqpoint{2.951026in}{1.089456in}}%
\pgfpathclose%
\pgfpathmoveto{\pgfqpoint{2.986276in}{1.111467in}}%
\pgfpathquadraticcurveto{\pgfqpoint{2.986276in}{1.122887in}}{\pgfqpoint{2.981575in}{1.129389in}}%
\pgfpathquadraticcurveto{\pgfqpoint{2.976874in}{1.135892in}}{\pgfqpoint{2.968660in}{1.135892in}}%
\pgfpathquadraticcurveto{\pgfqpoint{2.960429in}{1.135892in}}{\pgfqpoint{2.955727in}{1.129389in}}%
\pgfpathquadraticcurveto{\pgfqpoint{2.951026in}{1.122887in}}{\pgfqpoint{2.951026in}{1.111467in}}%
\pgfpathquadraticcurveto{\pgfqpoint{2.951026in}{1.100048in}}{\pgfqpoint{2.955727in}{1.093545in}}%
\pgfpathquadraticcurveto{\pgfqpoint{2.960429in}{1.087043in}}{\pgfqpoint{2.968660in}{1.087043in}}%
\pgfpathquadraticcurveto{\pgfqpoint{2.976874in}{1.087043in}}{\pgfqpoint{2.981575in}{1.093545in}}%
\pgfpathquadraticcurveto{\pgfqpoint{2.986276in}{1.100048in}}{\pgfqpoint{2.986276in}{1.111467in}}%
\pgfpathlineto{\pgfqpoint{2.986276in}{1.111467in}}%
\pgfpathclose%
\pgfpathmoveto{\pgfqpoint{3.038619in}{1.135784in}}%
\pgfpathquadraticcurveto{\pgfqpoint{3.030298in}{1.135784in}}{\pgfqpoint{3.025453in}{1.129281in}}%
\pgfpathquadraticcurveto{\pgfqpoint{3.020607in}{1.122779in}}{\pgfqpoint{3.020607in}{1.111467in}}%
\pgfpathquadraticcurveto{\pgfqpoint{3.020607in}{1.100156in}}{\pgfqpoint{3.025416in}{1.093653in}}%
\pgfpathquadraticcurveto{\pgfqpoint{3.030244in}{1.087151in}}{\pgfqpoint{3.038619in}{1.087151in}}%
\pgfpathquadraticcurveto{\pgfqpoint{3.046905in}{1.087151in}}{\pgfqpoint{3.051732in}{1.093671in}}%
\pgfpathquadraticcurveto{\pgfqpoint{3.056578in}{1.100210in}}{\pgfqpoint{3.056578in}{1.111467in}}%
\pgfpathquadraticcurveto{\pgfqpoint{3.056578in}{1.122671in}}{\pgfqpoint{3.051732in}{1.129227in}}%
\pgfpathquadraticcurveto{\pgfqpoint{3.046905in}{1.135784in}}{\pgfqpoint{3.038619in}{1.135784in}}%
\pgfpathlineto{\pgfqpoint{3.038619in}{1.135784in}}%
\pgfpathclose%
\pgfpathmoveto{\pgfqpoint{3.038619in}{1.144556in}}%
\pgfpathquadraticcurveto{\pgfqpoint{3.052129in}{1.144556in}}{\pgfqpoint{3.059838in}{1.135766in}}%
\pgfpathquadraticcurveto{\pgfqpoint{3.067565in}{1.126994in}}{\pgfqpoint{3.067565in}{1.111467in}}%
\pgfpathquadraticcurveto{\pgfqpoint{3.067565in}{1.095995in}}{\pgfqpoint{3.059838in}{1.087169in}}%
\pgfpathquadraticcurveto{\pgfqpoint{3.052129in}{1.078361in}}{\pgfqpoint{3.038619in}{1.078361in}}%
\pgfpathquadraticcurveto{\pgfqpoint{3.025056in}{1.078361in}}{\pgfqpoint{3.017365in}{1.087169in}}%
\pgfpathquadraticcurveto{\pgfqpoint{3.009692in}{1.095995in}}{\pgfqpoint{3.009692in}{1.111467in}}%
\pgfpathquadraticcurveto{\pgfqpoint{3.009692in}{1.126994in}}{\pgfqpoint{3.017365in}{1.135766in}}%
\pgfpathquadraticcurveto{\pgfqpoint{3.025056in}{1.144556in}}{\pgfqpoint{3.038619in}{1.144556in}}%
\pgfpathlineto{\pgfqpoint{3.038619in}{1.144556in}}%
\pgfpathclose%
\pgfpathmoveto{\pgfqpoint{3.124916in}{1.141187in}}%
\pgfpathlineto{\pgfqpoint{3.124916in}{1.131389in}}%
\pgfpathquadraticcurveto{\pgfqpoint{3.120539in}{1.133640in}}{\pgfqpoint{3.115801in}{1.134757in}}%
\pgfpathquadraticcurveto{\pgfqpoint{3.111082in}{1.135892in}}{\pgfqpoint{3.106003in}{1.135892in}}%
\pgfpathquadraticcurveto{\pgfqpoint{3.098294in}{1.135892in}}{\pgfqpoint{3.094439in}{1.133532in}}%
\pgfpathquadraticcurveto{\pgfqpoint{3.090584in}{1.131173in}}{\pgfqpoint{3.090584in}{1.126435in}}%
\pgfpathquadraticcurveto{\pgfqpoint{3.090584in}{1.122833in}}{\pgfqpoint{3.093340in}{1.120780in}}%
\pgfpathquadraticcurveto{\pgfqpoint{3.096096in}{1.118726in}}{\pgfqpoint{3.104436in}{1.116871in}}%
\pgfpathlineto{\pgfqpoint{3.107984in}{1.116078in}}%
\pgfpathquadraticcurveto{\pgfqpoint{3.119008in}{1.113719in}}{\pgfqpoint{3.123655in}{1.109414in}}%
\pgfpathquadraticcurveto{\pgfqpoint{3.128302in}{1.105109in}}{\pgfqpoint{3.128302in}{1.097400in}}%
\pgfpathquadraticcurveto{\pgfqpoint{3.128302in}{1.088610in}}{\pgfqpoint{3.121349in}{1.083476in}}%
\pgfpathquadraticcurveto{\pgfqpoint{3.114397in}{1.078361in}}{\pgfqpoint{3.102238in}{1.078361in}}%
\pgfpathquadraticcurveto{\pgfqpoint{3.097177in}{1.078361in}}{\pgfqpoint{3.091683in}{1.079352in}}%
\pgfpathquadraticcurveto{\pgfqpoint{3.086189in}{1.080342in}}{\pgfqpoint{3.080119in}{1.082306in}}%
\pgfpathlineto{\pgfqpoint{3.080119in}{1.093005in}}%
\pgfpathquadraticcurveto{\pgfqpoint{3.085865in}{1.090015in}}{\pgfqpoint{3.091431in}{1.088520in}}%
\pgfpathquadraticcurveto{\pgfqpoint{3.096997in}{1.087043in}}{\pgfqpoint{3.102472in}{1.087043in}}%
\pgfpathquadraticcurveto{\pgfqpoint{3.109785in}{1.087043in}}{\pgfqpoint{3.113712in}{1.089546in}}%
\pgfpathquadraticcurveto{\pgfqpoint{3.117657in}{1.092050in}}{\pgfqpoint{3.117657in}{1.096607in}}%
\pgfpathquadraticcurveto{\pgfqpoint{3.117657in}{1.100822in}}{\pgfqpoint{3.114811in}{1.103074in}}%
\pgfpathquadraticcurveto{\pgfqpoint{3.111983in}{1.105325in}}{\pgfqpoint{3.102346in}{1.107414in}}%
\pgfpathlineto{\pgfqpoint{3.098744in}{1.108261in}}%
\pgfpathquadraticcurveto{\pgfqpoint{3.089125in}{1.110278in}}{\pgfqpoint{3.084839in}{1.114475in}}%
\pgfpathquadraticcurveto{\pgfqpoint{3.080570in}{1.118672in}}{\pgfqpoint{3.080570in}{1.125985in}}%
\pgfpathquadraticcurveto{\pgfqpoint{3.080570in}{1.134883in}}{\pgfqpoint{3.086874in}{1.139710in}}%
\pgfpathquadraticcurveto{\pgfqpoint{3.093178in}{1.144556in}}{\pgfqpoint{3.104778in}{1.144556in}}%
\pgfpathquadraticcurveto{\pgfqpoint{3.110506in}{1.144556in}}{\pgfqpoint{3.115567in}{1.143709in}}%
\pgfpathquadraticcurveto{\pgfqpoint{3.120647in}{1.142880in}}{\pgfqpoint{3.124916in}{1.141187in}}%
\pgfpathlineto{\pgfqpoint{3.124916in}{1.141187in}}%
\pgfpathclose%
\pgfpathmoveto{\pgfqpoint{3.144783in}{1.143043in}}%
\pgfpathlineto{\pgfqpoint{3.155140in}{1.143043in}}%
\pgfpathlineto{\pgfqpoint{3.155140in}{1.080000in}}%
\pgfpathlineto{\pgfqpoint{3.144783in}{1.080000in}}%
\pgfpathlineto{\pgfqpoint{3.144783in}{1.143043in}}%
\pgfpathlineto{\pgfqpoint{3.144783in}{1.143043in}}%
\pgfpathclose%
\pgfpathmoveto{\pgfqpoint{3.144783in}{1.167593in}}%
\pgfpathlineto{\pgfqpoint{3.155140in}{1.167593in}}%
\pgfpathlineto{\pgfqpoint{3.155140in}{1.154462in}}%
\pgfpathlineto{\pgfqpoint{3.144783in}{1.154462in}}%
\pgfpathlineto{\pgfqpoint{3.144783in}{1.167593in}}%
\pgfpathlineto{\pgfqpoint{3.144783in}{1.167593in}}%
\pgfpathclose%
\pgfpathmoveto{\pgfqpoint{3.187058in}{1.160947in}}%
\pgfpathlineto{\pgfqpoint{3.187058in}{1.143043in}}%
\pgfpathlineto{\pgfqpoint{3.208384in}{1.143043in}}%
\pgfpathlineto{\pgfqpoint{3.208384in}{1.134991in}}%
\pgfpathlineto{\pgfqpoint{3.187058in}{1.134991in}}%
\pgfpathlineto{\pgfqpoint{3.187058in}{1.100768in}}%
\pgfpathquadraticcurveto{\pgfqpoint{3.187058in}{1.093059in}}{\pgfqpoint{3.189165in}{1.090861in}}%
\pgfpathquadraticcurveto{\pgfqpoint{3.191272in}{1.088664in}}{\pgfqpoint{3.197757in}{1.088664in}}%
\pgfpathlineto{\pgfqpoint{3.208384in}{1.088664in}}%
\pgfpathlineto{\pgfqpoint{3.208384in}{1.080000in}}%
\pgfpathlineto{\pgfqpoint{3.197757in}{1.080000in}}%
\pgfpathquadraticcurveto{\pgfqpoint{3.185761in}{1.080000in}}{\pgfqpoint{3.181204in}{1.084467in}}%
\pgfpathquadraticcurveto{\pgfqpoint{3.176647in}{1.088952in}}{\pgfqpoint{3.176647in}{1.100768in}}%
\pgfpathlineto{\pgfqpoint{3.176647in}{1.134991in}}%
\pgfpathlineto{\pgfqpoint{3.169045in}{1.134991in}}%
\pgfpathlineto{\pgfqpoint{3.169045in}{1.143043in}}%
\pgfpathlineto{\pgfqpoint{3.176647in}{1.143043in}}%
\pgfpathlineto{\pgfqpoint{3.176647in}{1.160947in}}%
\pgfpathlineto{\pgfqpoint{3.187058in}{1.160947in}}%
\pgfpathlineto{\pgfqpoint{3.187058in}{1.160947in}}%
\pgfpathclose%
\pgfpathmoveto{\pgfqpoint{3.275930in}{1.114115in}}%
\pgfpathlineto{\pgfqpoint{3.275930in}{1.109054in}}%
\pgfpathlineto{\pgfqpoint{3.228305in}{1.109054in}}%
\pgfpathquadraticcurveto{\pgfqpoint{3.228990in}{1.098354in}}{\pgfqpoint{3.234754in}{1.092753in}}%
\pgfpathquadraticcurveto{\pgfqpoint{3.240518in}{1.087151in}}{\pgfqpoint{3.250821in}{1.087151in}}%
\pgfpathquadraticcurveto{\pgfqpoint{3.256783in}{1.087151in}}{\pgfqpoint{3.262384in}{1.088610in}}%
\pgfpathquadraticcurveto{\pgfqpoint{3.267986in}{1.090069in}}{\pgfqpoint{3.273516in}{1.093005in}}%
\pgfpathlineto{\pgfqpoint{3.273516in}{1.083206in}}%
\pgfpathquadraticcurveto{\pgfqpoint{3.267932in}{1.080847in}}{\pgfqpoint{3.262078in}{1.079604in}}%
\pgfpathquadraticcurveto{\pgfqpoint{3.256224in}{1.078361in}}{\pgfqpoint{3.250208in}{1.078361in}}%
\pgfpathquadraticcurveto{\pgfqpoint{3.235114in}{1.078361in}}{\pgfqpoint{3.226306in}{1.087133in}}%
\pgfpathquadraticcurveto{\pgfqpoint{3.217498in}{1.095923in}}{\pgfqpoint{3.217498in}{1.110909in}}%
\pgfpathquadraticcurveto{\pgfqpoint{3.217498in}{1.126381in}}{\pgfqpoint{3.225856in}{1.135459in}}%
\pgfpathquadraticcurveto{\pgfqpoint{3.234213in}{1.144556in}}{\pgfqpoint{3.248407in}{1.144556in}}%
\pgfpathquadraticcurveto{\pgfqpoint{3.261124in}{1.144556in}}{\pgfqpoint{3.268527in}{1.136360in}}%
\pgfpathquadraticcurveto{\pgfqpoint{3.275930in}{1.128183in}}{\pgfqpoint{3.275930in}{1.114115in}}%
\pgfpathlineto{\pgfqpoint{3.275930in}{1.114115in}}%
\pgfpathclose%
\pgfpathmoveto{\pgfqpoint{3.265573in}{1.117159in}}%
\pgfpathquadraticcurveto{\pgfqpoint{3.265464in}{1.125643in}}{\pgfqpoint{3.260817in}{1.130704in}}%
\pgfpathquadraticcurveto{\pgfqpoint{3.256170in}{1.135784in}}{\pgfqpoint{3.248515in}{1.135784in}}%
\pgfpathquadraticcurveto{\pgfqpoint{3.239851in}{1.135784in}}{\pgfqpoint{3.234646in}{1.130884in}}%
\pgfpathquadraticcurveto{\pgfqpoint{3.229440in}{1.125985in}}{\pgfqpoint{3.228648in}{1.117087in}}%
\pgfpathlineto{\pgfqpoint{3.265573in}{1.117159in}}%
\pgfpathlineto{\pgfqpoint{3.265573in}{1.117159in}}%
\pgfpathclose%
\pgfusepath{fill}%
\end{pgfscope}%
\begin{pgfscope}%
\pgfsetbuttcap%
\pgfsetroundjoin%
\definecolor{currentfill}{rgb}{0.309804,0.372549,0.419608}%
\pgfsetfillcolor{currentfill}%
\pgfsetlinewidth{0.000000pt}%
\definecolor{currentstroke}{rgb}{0.309804,0.372549,0.419608}%
\pgfsetstrokecolor{currentstroke}%
\pgfsetdash{}{0pt}%
\pgfpathmoveto{\pgfqpoint{3.401605in}{1.118060in}}%
\pgfpathlineto{\pgfqpoint{3.401605in}{1.080000in}}%
\pgfpathlineto{\pgfqpoint{3.391248in}{1.080000in}}%
\pgfpathlineto{\pgfqpoint{3.391248in}{1.117717in}}%
\pgfpathquadraticcurveto{\pgfqpoint{3.391248in}{1.126669in}}{\pgfqpoint{3.387753in}{1.131100in}}%
\pgfpathquadraticcurveto{\pgfqpoint{3.384259in}{1.135549in}}{\pgfqpoint{3.377288in}{1.135549in}}%
\pgfpathquadraticcurveto{\pgfqpoint{3.368895in}{1.135549in}}{\pgfqpoint{3.364049in}{1.130200in}}%
\pgfpathquadraticcurveto{\pgfqpoint{3.359204in}{1.124868in}}{\pgfqpoint{3.359204in}{1.115628in}}%
\pgfpathlineto{\pgfqpoint{3.359204in}{1.080000in}}%
\pgfpathlineto{\pgfqpoint{3.348793in}{1.080000in}}%
\pgfpathlineto{\pgfqpoint{3.348793in}{1.143043in}}%
\pgfpathlineto{\pgfqpoint{3.359204in}{1.143043in}}%
\pgfpathlineto{\pgfqpoint{3.359204in}{1.133244in}}%
\pgfpathquadraticcurveto{\pgfqpoint{3.362933in}{1.138936in}}{\pgfqpoint{3.367958in}{1.141746in}}%
\pgfpathquadraticcurveto{\pgfqpoint{3.373002in}{1.144556in}}{\pgfqpoint{3.379594in}{1.144556in}}%
\pgfpathquadraticcurveto{\pgfqpoint{3.390455in}{1.144556in}}{\pgfqpoint{3.396021in}{1.137837in}}%
\pgfpathquadraticcurveto{\pgfqpoint{3.401605in}{1.131118in}}{\pgfqpoint{3.401605in}{1.118060in}}%
\pgfpathlineto{\pgfqpoint{3.401605in}{1.118060in}}%
\pgfpathclose%
\pgfpathmoveto{\pgfqpoint{3.446671in}{1.135784in}}%
\pgfpathquadraticcurveto{\pgfqpoint{3.438350in}{1.135784in}}{\pgfqpoint{3.433504in}{1.129281in}}%
\pgfpathquadraticcurveto{\pgfqpoint{3.428659in}{1.122779in}}{\pgfqpoint{3.428659in}{1.111467in}}%
\pgfpathquadraticcurveto{\pgfqpoint{3.428659in}{1.100156in}}{\pgfqpoint{3.433468in}{1.093653in}}%
\pgfpathquadraticcurveto{\pgfqpoint{3.438296in}{1.087151in}}{\pgfqpoint{3.446671in}{1.087151in}}%
\pgfpathquadraticcurveto{\pgfqpoint{3.454957in}{1.087151in}}{\pgfqpoint{3.459784in}{1.093671in}}%
\pgfpathquadraticcurveto{\pgfqpoint{3.464629in}{1.100210in}}{\pgfqpoint{3.464629in}{1.111467in}}%
\pgfpathquadraticcurveto{\pgfqpoint{3.464629in}{1.122671in}}{\pgfqpoint{3.459784in}{1.129227in}}%
\pgfpathquadraticcurveto{\pgfqpoint{3.454957in}{1.135784in}}{\pgfqpoint{3.446671in}{1.135784in}}%
\pgfpathlineto{\pgfqpoint{3.446671in}{1.135784in}}%
\pgfpathclose%
\pgfpathmoveto{\pgfqpoint{3.446671in}{1.144556in}}%
\pgfpathquadraticcurveto{\pgfqpoint{3.460180in}{1.144556in}}{\pgfqpoint{3.467890in}{1.135766in}}%
\pgfpathquadraticcurveto{\pgfqpoint{3.475617in}{1.126994in}}{\pgfqpoint{3.475617in}{1.111467in}}%
\pgfpathquadraticcurveto{\pgfqpoint{3.475617in}{1.095995in}}{\pgfqpoint{3.467890in}{1.087169in}}%
\pgfpathquadraticcurveto{\pgfqpoint{3.460180in}{1.078361in}}{\pgfqpoint{3.446671in}{1.078361in}}%
\pgfpathquadraticcurveto{\pgfqpoint{3.433108in}{1.078361in}}{\pgfqpoint{3.425417in}{1.087169in}}%
\pgfpathquadraticcurveto{\pgfqpoint{3.417744in}{1.095995in}}{\pgfqpoint{3.417744in}{1.111467in}}%
\pgfpathquadraticcurveto{\pgfqpoint{3.417744in}{1.126994in}}{\pgfqpoint{3.425417in}{1.135766in}}%
\pgfpathquadraticcurveto{\pgfqpoint{3.433108in}{1.144556in}}{\pgfqpoint{3.446671in}{1.144556in}}%
\pgfpathlineto{\pgfqpoint{3.446671in}{1.144556in}}%
\pgfpathclose%
\pgfpathmoveto{\pgfqpoint{3.545198in}{1.118060in}}%
\pgfpathlineto{\pgfqpoint{3.545198in}{1.080000in}}%
\pgfpathlineto{\pgfqpoint{3.534841in}{1.080000in}}%
\pgfpathlineto{\pgfqpoint{3.534841in}{1.117717in}}%
\pgfpathquadraticcurveto{\pgfqpoint{3.534841in}{1.126669in}}{\pgfqpoint{3.531346in}{1.131100in}}%
\pgfpathquadraticcurveto{\pgfqpoint{3.527852in}{1.135549in}}{\pgfqpoint{3.520881in}{1.135549in}}%
\pgfpathquadraticcurveto{\pgfqpoint{3.512488in}{1.135549in}}{\pgfqpoint{3.507642in}{1.130200in}}%
\pgfpathquadraticcurveto{\pgfqpoint{3.502797in}{1.124868in}}{\pgfqpoint{3.502797in}{1.115628in}}%
\pgfpathlineto{\pgfqpoint{3.502797in}{1.080000in}}%
\pgfpathlineto{\pgfqpoint{3.492386in}{1.080000in}}%
\pgfpathlineto{\pgfqpoint{3.492386in}{1.143043in}}%
\pgfpathlineto{\pgfqpoint{3.502797in}{1.143043in}}%
\pgfpathlineto{\pgfqpoint{3.502797in}{1.133244in}}%
\pgfpathquadraticcurveto{\pgfqpoint{3.506526in}{1.138936in}}{\pgfqpoint{3.511551in}{1.141746in}}%
\pgfpathquadraticcurveto{\pgfqpoint{3.516594in}{1.144556in}}{\pgfqpoint{3.523187in}{1.144556in}}%
\pgfpathquadraticcurveto{\pgfqpoint{3.534048in}{1.144556in}}{\pgfqpoint{3.539614in}{1.137837in}}%
\pgfpathquadraticcurveto{\pgfqpoint{3.545198in}{1.131118in}}{\pgfqpoint{3.545198in}{1.118060in}}%
\pgfpathlineto{\pgfqpoint{3.545198in}{1.118060in}}%
\pgfpathclose%
\pgfpathmoveto{\pgfqpoint{3.561337in}{1.143043in}}%
\pgfpathlineto{\pgfqpoint{3.610528in}{1.143043in}}%
\pgfpathlineto{\pgfqpoint{3.610528in}{1.133586in}}%
\pgfpathlineto{\pgfqpoint{3.571586in}{1.088268in}}%
\pgfpathlineto{\pgfqpoint{3.610528in}{1.088268in}}%
\pgfpathlineto{\pgfqpoint{3.610528in}{1.080000in}}%
\pgfpathlineto{\pgfqpoint{3.559932in}{1.080000in}}%
\pgfpathlineto{\pgfqpoint{3.559932in}{1.089456in}}%
\pgfpathlineto{\pgfqpoint{3.598892in}{1.134775in}}%
\pgfpathlineto{\pgfqpoint{3.561337in}{1.134775in}}%
\pgfpathlineto{\pgfqpoint{3.561337in}{1.143043in}}%
\pgfpathlineto{\pgfqpoint{3.561337in}{1.143043in}}%
\pgfpathclose%
\pgfpathmoveto{\pgfqpoint{3.680271in}{1.114115in}}%
\pgfpathlineto{\pgfqpoint{3.680271in}{1.109054in}}%
\pgfpathlineto{\pgfqpoint{3.632647in}{1.109054in}}%
\pgfpathquadraticcurveto{\pgfqpoint{3.633331in}{1.098354in}}{\pgfqpoint{3.639095in}{1.092753in}}%
\pgfpathquadraticcurveto{\pgfqpoint{3.644859in}{1.087151in}}{\pgfqpoint{3.655162in}{1.087151in}}%
\pgfpathquadraticcurveto{\pgfqpoint{3.661124in}{1.087151in}}{\pgfqpoint{3.666726in}{1.088610in}}%
\pgfpathquadraticcurveto{\pgfqpoint{3.672327in}{1.090069in}}{\pgfqpoint{3.677857in}{1.093005in}}%
\pgfpathlineto{\pgfqpoint{3.677857in}{1.083206in}}%
\pgfpathquadraticcurveto{\pgfqpoint{3.672273in}{1.080847in}}{\pgfqpoint{3.666419in}{1.079604in}}%
\pgfpathquadraticcurveto{\pgfqpoint{3.660566in}{1.078361in}}{\pgfqpoint{3.654549in}{1.078361in}}%
\pgfpathquadraticcurveto{\pgfqpoint{3.639455in}{1.078361in}}{\pgfqpoint{3.630647in}{1.087133in}}%
\pgfpathquadraticcurveto{\pgfqpoint{3.621839in}{1.095923in}}{\pgfqpoint{3.621839in}{1.110909in}}%
\pgfpathquadraticcurveto{\pgfqpoint{3.621839in}{1.126381in}}{\pgfqpoint{3.630197in}{1.135459in}}%
\pgfpathquadraticcurveto{\pgfqpoint{3.638555in}{1.144556in}}{\pgfqpoint{3.652748in}{1.144556in}}%
\pgfpathquadraticcurveto{\pgfqpoint{3.665465in}{1.144556in}}{\pgfqpoint{3.672868in}{1.136360in}}%
\pgfpathquadraticcurveto{\pgfqpoint{3.680271in}{1.128183in}}{\pgfqpoint{3.680271in}{1.114115in}}%
\pgfpathlineto{\pgfqpoint{3.680271in}{1.114115in}}%
\pgfpathclose%
\pgfpathmoveto{\pgfqpoint{3.669914in}{1.117159in}}%
\pgfpathquadraticcurveto{\pgfqpoint{3.669806in}{1.125643in}}{\pgfqpoint{3.665159in}{1.130704in}}%
\pgfpathquadraticcurveto{\pgfqpoint{3.660512in}{1.135784in}}{\pgfqpoint{3.652856in}{1.135784in}}%
\pgfpathquadraticcurveto{\pgfqpoint{3.644192in}{1.135784in}}{\pgfqpoint{3.638987in}{1.130884in}}%
\pgfpathquadraticcurveto{\pgfqpoint{3.633781in}{1.125985in}}{\pgfqpoint{3.632989in}{1.117087in}}%
\pgfpathlineto{\pgfqpoint{3.669914in}{1.117159in}}%
\pgfpathlineto{\pgfqpoint{3.669914in}{1.117159in}}%
\pgfpathclose%
\pgfpathmoveto{\pgfqpoint{3.733803in}{1.133370in}}%
\pgfpathquadraticcurveto{\pgfqpoint{3.732056in}{1.134379in}}{\pgfqpoint{3.730002in}{1.134847in}}%
\pgfpathquadraticcurveto{\pgfqpoint{3.727949in}{1.135333in}}{\pgfqpoint{3.725481in}{1.135333in}}%
\pgfpathquadraticcurveto{\pgfqpoint{3.716691in}{1.135333in}}{\pgfqpoint{3.711990in}{1.129623in}}%
\pgfpathquadraticcurveto{\pgfqpoint{3.707289in}{1.123914in}}{\pgfqpoint{3.707289in}{1.113214in}}%
\pgfpathlineto{\pgfqpoint{3.707289in}{1.080000in}}%
\pgfpathlineto{\pgfqpoint{3.696878in}{1.080000in}}%
\pgfpathlineto{\pgfqpoint{3.696878in}{1.143043in}}%
\pgfpathlineto{\pgfqpoint{3.707289in}{1.143043in}}%
\pgfpathlineto{\pgfqpoint{3.707289in}{1.133244in}}%
\pgfpathquadraticcurveto{\pgfqpoint{3.710567in}{1.138990in}}{\pgfqpoint{3.715791in}{1.141764in}}%
\pgfpathquadraticcurveto{\pgfqpoint{3.721032in}{1.144556in}}{\pgfqpoint{3.728525in}{1.144556in}}%
\pgfpathquadraticcurveto{\pgfqpoint{3.729588in}{1.144556in}}{\pgfqpoint{3.730885in}{1.144411in}}%
\pgfpathquadraticcurveto{\pgfqpoint{3.732182in}{1.144285in}}{\pgfqpoint{3.733749in}{1.143997in}}%
\pgfpathlineto{\pgfqpoint{3.733803in}{1.133370in}}%
\pgfpathlineto{\pgfqpoint{3.733803in}{1.133370in}}%
\pgfpathclose%
\pgfpathmoveto{\pgfqpoint{3.766549in}{1.135784in}}%
\pgfpathquadraticcurveto{\pgfqpoint{3.758227in}{1.135784in}}{\pgfqpoint{3.753382in}{1.129281in}}%
\pgfpathquadraticcurveto{\pgfqpoint{3.748537in}{1.122779in}}{\pgfqpoint{3.748537in}{1.111467in}}%
\pgfpathquadraticcurveto{\pgfqpoint{3.748537in}{1.100156in}}{\pgfqpoint{3.753346in}{1.093653in}}%
\pgfpathquadraticcurveto{\pgfqpoint{3.758173in}{1.087151in}}{\pgfqpoint{3.766549in}{1.087151in}}%
\pgfpathquadraticcurveto{\pgfqpoint{3.774835in}{1.087151in}}{\pgfqpoint{3.779662in}{1.093671in}}%
\pgfpathquadraticcurveto{\pgfqpoint{3.784507in}{1.100210in}}{\pgfqpoint{3.784507in}{1.111467in}}%
\pgfpathquadraticcurveto{\pgfqpoint{3.784507in}{1.122671in}}{\pgfqpoint{3.779662in}{1.129227in}}%
\pgfpathquadraticcurveto{\pgfqpoint{3.774835in}{1.135784in}}{\pgfqpoint{3.766549in}{1.135784in}}%
\pgfpathlineto{\pgfqpoint{3.766549in}{1.135784in}}%
\pgfpathclose%
\pgfpathmoveto{\pgfqpoint{3.766549in}{1.144556in}}%
\pgfpathquadraticcurveto{\pgfqpoint{3.780058in}{1.144556in}}{\pgfqpoint{3.787767in}{1.135766in}}%
\pgfpathquadraticcurveto{\pgfqpoint{3.795495in}{1.126994in}}{\pgfqpoint{3.795495in}{1.111467in}}%
\pgfpathquadraticcurveto{\pgfqpoint{3.795495in}{1.095995in}}{\pgfqpoint{3.787767in}{1.087169in}}%
\pgfpathquadraticcurveto{\pgfqpoint{3.780058in}{1.078361in}}{\pgfqpoint{3.766549in}{1.078361in}}%
\pgfpathquadraticcurveto{\pgfqpoint{3.752986in}{1.078361in}}{\pgfqpoint{3.745295in}{1.087169in}}%
\pgfpathquadraticcurveto{\pgfqpoint{3.737622in}{1.095995in}}{\pgfqpoint{3.737622in}{1.111467in}}%
\pgfpathquadraticcurveto{\pgfqpoint{3.737622in}{1.126994in}}{\pgfqpoint{3.745295in}{1.135766in}}%
\pgfpathquadraticcurveto{\pgfqpoint{3.752986in}{1.144556in}}{\pgfqpoint{3.766549in}{1.144556in}}%
\pgfpathlineto{\pgfqpoint{3.766549in}{1.144556in}}%
\pgfpathclose%
\pgfusepath{fill}%
\end{pgfscope}%
\begin{pgfscope}%
\pgfsetbuttcap%
\pgfsetroundjoin%
\definecolor{currentfill}{rgb}{0.309804,0.372549,0.419608}%
\pgfsetfillcolor{currentfill}%
\pgfsetlinewidth{0.000000pt}%
\definecolor{currentstroke}{rgb}{0.309804,0.372549,0.419608}%
\pgfsetstrokecolor{currentstroke}%
\pgfsetdash{}{0pt}%
\pgfpathmoveto{\pgfqpoint{3.908771in}{1.141187in}}%
\pgfpathlineto{\pgfqpoint{3.908771in}{1.131389in}}%
\pgfpathquadraticcurveto{\pgfqpoint{3.904394in}{1.133640in}}{\pgfqpoint{3.899657in}{1.134757in}}%
\pgfpathquadraticcurveto{\pgfqpoint{3.894938in}{1.135892in}}{\pgfqpoint{3.889859in}{1.135892in}}%
\pgfpathquadraticcurveto{\pgfqpoint{3.882149in}{1.135892in}}{\pgfqpoint{3.878295in}{1.133532in}}%
\pgfpathquadraticcurveto{\pgfqpoint{3.874440in}{1.131173in}}{\pgfqpoint{3.874440in}{1.126435in}}%
\pgfpathquadraticcurveto{\pgfqpoint{3.874440in}{1.122833in}}{\pgfqpoint{3.877196in}{1.120780in}}%
\pgfpathquadraticcurveto{\pgfqpoint{3.879952in}{1.118726in}}{\pgfqpoint{3.888292in}{1.116871in}}%
\pgfpathlineto{\pgfqpoint{3.891840in}{1.116078in}}%
\pgfpathquadraticcurveto{\pgfqpoint{3.902863in}{1.113719in}}{\pgfqpoint{3.907511in}{1.109414in}}%
\pgfpathquadraticcurveto{\pgfqpoint{3.912158in}{1.105109in}}{\pgfqpoint{3.912158in}{1.097400in}}%
\pgfpathquadraticcurveto{\pgfqpoint{3.912158in}{1.088610in}}{\pgfqpoint{3.905205in}{1.083476in}}%
\pgfpathquadraticcurveto{\pgfqpoint{3.898252in}{1.078361in}}{\pgfqpoint{3.886094in}{1.078361in}}%
\pgfpathquadraticcurveto{\pgfqpoint{3.881033in}{1.078361in}}{\pgfqpoint{3.875539in}{1.079352in}}%
\pgfpathquadraticcurveto{\pgfqpoint{3.870045in}{1.080342in}}{\pgfqpoint{3.863975in}{1.082306in}}%
\pgfpathlineto{\pgfqpoint{3.863975in}{1.093005in}}%
\pgfpathquadraticcurveto{\pgfqpoint{3.869721in}{1.090015in}}{\pgfqpoint{3.875287in}{1.088520in}}%
\pgfpathquadraticcurveto{\pgfqpoint{3.880853in}{1.087043in}}{\pgfqpoint{3.886328in}{1.087043in}}%
\pgfpathquadraticcurveto{\pgfqpoint{3.893641in}{1.087043in}}{\pgfqpoint{3.897568in}{1.089546in}}%
\pgfpathquadraticcurveto{\pgfqpoint{3.901513in}{1.092050in}}{\pgfqpoint{3.901513in}{1.096607in}}%
\pgfpathquadraticcurveto{\pgfqpoint{3.901513in}{1.100822in}}{\pgfqpoint{3.898667in}{1.103074in}}%
\pgfpathquadraticcurveto{\pgfqpoint{3.895839in}{1.105325in}}{\pgfqpoint{3.886202in}{1.107414in}}%
\pgfpathlineto{\pgfqpoint{3.882600in}{1.108261in}}%
\pgfpathquadraticcurveto{\pgfqpoint{3.872981in}{1.110278in}}{\pgfqpoint{3.868694in}{1.114475in}}%
\pgfpathquadraticcurveto{\pgfqpoint{3.864426in}{1.118672in}}{\pgfqpoint{3.864426in}{1.125985in}}%
\pgfpathquadraticcurveto{\pgfqpoint{3.864426in}{1.134883in}}{\pgfqpoint{3.870730in}{1.139710in}}%
\pgfpathquadraticcurveto{\pgfqpoint{3.877034in}{1.144556in}}{\pgfqpoint{3.888634in}{1.144556in}}%
\pgfpathquadraticcurveto{\pgfqpoint{3.894362in}{1.144556in}}{\pgfqpoint{3.899423in}{1.143709in}}%
\pgfpathquadraticcurveto{\pgfqpoint{3.904503in}{1.142880in}}{\pgfqpoint{3.908771in}{1.141187in}}%
\pgfpathlineto{\pgfqpoint{3.908771in}{1.141187in}}%
\pgfpathclose%
\pgfpathmoveto{\pgfqpoint{3.928639in}{1.143043in}}%
\pgfpathlineto{\pgfqpoint{3.938996in}{1.143043in}}%
\pgfpathlineto{\pgfqpoint{3.938996in}{1.080000in}}%
\pgfpathlineto{\pgfqpoint{3.928639in}{1.080000in}}%
\pgfpathlineto{\pgfqpoint{3.928639in}{1.143043in}}%
\pgfpathlineto{\pgfqpoint{3.928639in}{1.143043in}}%
\pgfpathclose%
\pgfpathmoveto{\pgfqpoint{3.928639in}{1.167593in}}%
\pgfpathlineto{\pgfqpoint{3.938996in}{1.167593in}}%
\pgfpathlineto{\pgfqpoint{3.938996in}{1.154462in}}%
\pgfpathlineto{\pgfqpoint{3.928639in}{1.154462in}}%
\pgfpathlineto{\pgfqpoint{3.928639in}{1.167593in}}%
\pgfpathlineto{\pgfqpoint{3.928639in}{1.167593in}}%
\pgfpathclose%
\pgfpathmoveto{\pgfqpoint{4.002146in}{1.112260in}}%
\pgfpathquadraticcurveto{\pgfqpoint{4.002146in}{1.123517in}}{\pgfqpoint{3.997499in}{1.129696in}}%
\pgfpathquadraticcurveto{\pgfqpoint{3.992870in}{1.135892in}}{\pgfqpoint{3.984477in}{1.135892in}}%
\pgfpathquadraticcurveto{\pgfqpoint{3.976155in}{1.135892in}}{\pgfqpoint{3.971508in}{1.129696in}}%
\pgfpathquadraticcurveto{\pgfqpoint{3.966861in}{1.123517in}}{\pgfqpoint{3.966861in}{1.112260in}}%
\pgfpathquadraticcurveto{\pgfqpoint{3.966861in}{1.101056in}}{\pgfqpoint{3.971508in}{1.094860in}}%
\pgfpathquadraticcurveto{\pgfqpoint{3.976155in}{1.088664in}}{\pgfqpoint{3.984477in}{1.088664in}}%
\pgfpathquadraticcurveto{\pgfqpoint{3.992870in}{1.088664in}}{\pgfqpoint{3.997499in}{1.094860in}}%
\pgfpathquadraticcurveto{\pgfqpoint{4.002146in}{1.101056in}}{\pgfqpoint{4.002146in}{1.112260in}}%
\pgfpathlineto{\pgfqpoint{4.002146in}{1.112260in}}%
\pgfpathclose%
\pgfpathmoveto{\pgfqpoint{4.012503in}{1.087817in}}%
\pgfpathquadraticcurveto{\pgfqpoint{4.012503in}{1.071732in}}{\pgfqpoint{4.005353in}{1.063879in}}%
\pgfpathquadraticcurveto{\pgfqpoint{3.998220in}{1.056026in}}{\pgfqpoint{3.983468in}{1.056026in}}%
\pgfpathquadraticcurveto{\pgfqpoint{3.978010in}{1.056026in}}{\pgfqpoint{3.973165in}{1.056836in}}%
\pgfpathquadraticcurveto{\pgfqpoint{3.968320in}{1.057647in}}{\pgfqpoint{3.963763in}{1.059340in}}%
\pgfpathlineto{\pgfqpoint{3.963763in}{1.069409in}}%
\pgfpathquadraticcurveto{\pgfqpoint{3.968320in}{1.066941in}}{\pgfqpoint{3.972769in}{1.065770in}}%
\pgfpathquadraticcurveto{\pgfqpoint{3.977218in}{1.064582in}}{\pgfqpoint{3.981829in}{1.064582in}}%
\pgfpathquadraticcurveto{\pgfqpoint{3.992024in}{1.064582in}}{\pgfqpoint{3.997085in}{1.069895in}}%
\pgfpathquadraticcurveto{\pgfqpoint{4.002146in}{1.075209in}}{\pgfqpoint{4.002146in}{1.085962in}}%
\pgfpathlineto{\pgfqpoint{4.002146in}{1.091095in}}%
\pgfpathquadraticcurveto{\pgfqpoint{3.998940in}{1.085512in}}{\pgfqpoint{3.993933in}{1.082756in}}%
\pgfpathquadraticcurveto{\pgfqpoint{3.988926in}{1.080000in}}{\pgfqpoint{3.981937in}{1.080000in}}%
\pgfpathquadraticcurveto{\pgfqpoint{3.970355in}{1.080000in}}{\pgfqpoint{3.963258in}{1.088826in}}%
\pgfpathquadraticcurveto{\pgfqpoint{3.956161in}{1.097670in}}{\pgfqpoint{3.956161in}{1.112260in}}%
\pgfpathquadraticcurveto{\pgfqpoint{3.956161in}{1.126886in}}{\pgfqpoint{3.963258in}{1.135712in}}%
\pgfpathquadraticcurveto{\pgfqpoint{3.970355in}{1.144556in}}{\pgfqpoint{3.981937in}{1.144556in}}%
\pgfpathquadraticcurveto{\pgfqpoint{3.988926in}{1.144556in}}{\pgfqpoint{3.993933in}{1.141800in}}%
\pgfpathquadraticcurveto{\pgfqpoint{3.998940in}{1.139044in}}{\pgfqpoint{4.002146in}{1.133478in}}%
\pgfpathlineto{\pgfqpoint{4.002146in}{1.143043in}}%
\pgfpathlineto{\pgfqpoint{4.012503in}{1.143043in}}%
\pgfpathlineto{\pgfqpoint{4.012503in}{1.087817in}}%
\pgfpathlineto{\pgfqpoint{4.012503in}{1.087817in}}%
\pgfpathclose%
\pgfpathmoveto{\pgfqpoint{4.086263in}{1.118060in}}%
\pgfpathlineto{\pgfqpoint{4.086263in}{1.080000in}}%
\pgfpathlineto{\pgfqpoint{4.075906in}{1.080000in}}%
\pgfpathlineto{\pgfqpoint{4.075906in}{1.117717in}}%
\pgfpathquadraticcurveto{\pgfqpoint{4.075906in}{1.126669in}}{\pgfqpoint{4.072412in}{1.131100in}}%
\pgfpathquadraticcurveto{\pgfqpoint{4.068917in}{1.135549in}}{\pgfqpoint{4.061947in}{1.135549in}}%
\pgfpathquadraticcurveto{\pgfqpoint{4.053553in}{1.135549in}}{\pgfqpoint{4.048708in}{1.130200in}}%
\pgfpathquadraticcurveto{\pgfqpoint{4.043863in}{1.124868in}}{\pgfqpoint{4.043863in}{1.115628in}}%
\pgfpathlineto{\pgfqpoint{4.043863in}{1.080000in}}%
\pgfpathlineto{\pgfqpoint{4.033452in}{1.080000in}}%
\pgfpathlineto{\pgfqpoint{4.033452in}{1.143043in}}%
\pgfpathlineto{\pgfqpoint{4.043863in}{1.143043in}}%
\pgfpathlineto{\pgfqpoint{4.043863in}{1.133244in}}%
\pgfpathquadraticcurveto{\pgfqpoint{4.047591in}{1.138936in}}{\pgfqpoint{4.052616in}{1.141746in}}%
\pgfpathquadraticcurveto{\pgfqpoint{4.057660in}{1.144556in}}{\pgfqpoint{4.064252in}{1.144556in}}%
\pgfpathquadraticcurveto{\pgfqpoint{4.075114in}{1.144556in}}{\pgfqpoint{4.080679in}{1.137837in}}%
\pgfpathquadraticcurveto{\pgfqpoint{4.086263in}{1.131118in}}{\pgfqpoint{4.086263in}{1.118060in}}%
\pgfpathlineto{\pgfqpoint{4.086263in}{1.118060in}}%
\pgfpathclose%
\pgfpathmoveto{\pgfqpoint{4.147090in}{1.141187in}}%
\pgfpathlineto{\pgfqpoint{4.147090in}{1.131389in}}%
\pgfpathquadraticcurveto{\pgfqpoint{4.142713in}{1.133640in}}{\pgfqpoint{4.137976in}{1.134757in}}%
\pgfpathquadraticcurveto{\pgfqpoint{4.133257in}{1.135892in}}{\pgfqpoint{4.128177in}{1.135892in}}%
\pgfpathquadraticcurveto{\pgfqpoint{4.120468in}{1.135892in}}{\pgfqpoint{4.116614in}{1.133532in}}%
\pgfpathquadraticcurveto{\pgfqpoint{4.112759in}{1.131173in}}{\pgfqpoint{4.112759in}{1.126435in}}%
\pgfpathquadraticcurveto{\pgfqpoint{4.112759in}{1.122833in}}{\pgfqpoint{4.115515in}{1.120780in}}%
\pgfpathquadraticcurveto{\pgfqpoint{4.118271in}{1.118726in}}{\pgfqpoint{4.126610in}{1.116871in}}%
\pgfpathlineto{\pgfqpoint{4.130159in}{1.116078in}}%
\pgfpathquadraticcurveto{\pgfqpoint{4.141182in}{1.113719in}}{\pgfqpoint{4.145829in}{1.109414in}}%
\pgfpathquadraticcurveto{\pgfqpoint{4.150477in}{1.105109in}}{\pgfqpoint{4.150477in}{1.097400in}}%
\pgfpathquadraticcurveto{\pgfqpoint{4.150477in}{1.088610in}}{\pgfqpoint{4.143524in}{1.083476in}}%
\pgfpathquadraticcurveto{\pgfqpoint{4.136571in}{1.078361in}}{\pgfqpoint{4.124413in}{1.078361in}}%
\pgfpathquadraticcurveto{\pgfqpoint{4.119352in}{1.078361in}}{\pgfqpoint{4.113858in}{1.079352in}}%
\pgfpathquadraticcurveto{\pgfqpoint{4.108364in}{1.080342in}}{\pgfqpoint{4.102294in}{1.082306in}}%
\pgfpathlineto{\pgfqpoint{4.102294in}{1.093005in}}%
\pgfpathquadraticcurveto{\pgfqpoint{4.108040in}{1.090015in}}{\pgfqpoint{4.113606in}{1.088520in}}%
\pgfpathquadraticcurveto{\pgfqpoint{4.119171in}{1.087043in}}{\pgfqpoint{4.124647in}{1.087043in}}%
\pgfpathquadraticcurveto{\pgfqpoint{4.131960in}{1.087043in}}{\pgfqpoint{4.135887in}{1.089546in}}%
\pgfpathquadraticcurveto{\pgfqpoint{4.139831in}{1.092050in}}{\pgfqpoint{4.139831in}{1.096607in}}%
\pgfpathquadraticcurveto{\pgfqpoint{4.139831in}{1.100822in}}{\pgfqpoint{4.136985in}{1.103074in}}%
\pgfpathquadraticcurveto{\pgfqpoint{4.134158in}{1.105325in}}{\pgfqpoint{4.124521in}{1.107414in}}%
\pgfpathlineto{\pgfqpoint{4.120919in}{1.108261in}}%
\pgfpathquadraticcurveto{\pgfqpoint{4.111300in}{1.110278in}}{\pgfqpoint{4.107013in}{1.114475in}}%
\pgfpathquadraticcurveto{\pgfqpoint{4.102744in}{1.118672in}}{\pgfqpoint{4.102744in}{1.125985in}}%
\pgfpathquadraticcurveto{\pgfqpoint{4.102744in}{1.134883in}}{\pgfqpoint{4.109049in}{1.139710in}}%
\pgfpathquadraticcurveto{\pgfqpoint{4.115353in}{1.144556in}}{\pgfqpoint{4.126953in}{1.144556in}}%
\pgfpathquadraticcurveto{\pgfqpoint{4.132681in}{1.144556in}}{\pgfqpoint{4.137742in}{1.143709in}}%
\pgfpathquadraticcurveto{\pgfqpoint{4.142821in}{1.142880in}}{\pgfqpoint{4.147090in}{1.141187in}}%
\pgfpathlineto{\pgfqpoint{4.147090in}{1.141187in}}%
\pgfpathclose%
\pgfusepath{fill}%
\end{pgfscope}%
\begin{pgfscope}%
\pgfsetbuttcap%
\pgfsetroundjoin%
\definecolor{currentfill}{rgb}{0.309804,0.372549,0.419608}%
\pgfsetfillcolor{currentfill}%
\pgfsetlinewidth{0.000000pt}%
\definecolor{currentstroke}{rgb}{0.309804,0.372549,0.419608}%
\pgfsetstrokecolor{currentstroke}%
\pgfsetdash{}{0pt}%
\pgfpathmoveto{\pgfqpoint{4.249141in}{1.111683in}}%
\pgfpathquadraticcurveto{\pgfqpoint{4.236586in}{1.111683in}}{\pgfqpoint{4.231741in}{1.108819in}}%
\pgfpathquadraticcurveto{\pgfqpoint{4.226896in}{1.105956in}}{\pgfqpoint{4.226896in}{1.099021in}}%
\pgfpathquadraticcurveto{\pgfqpoint{4.226896in}{1.093509in}}{\pgfqpoint{4.230534in}{1.090267in}}%
\pgfpathquadraticcurveto{\pgfqpoint{4.234172in}{1.087043in}}{\pgfqpoint{4.240405in}{1.087043in}}%
\pgfpathquadraticcurveto{\pgfqpoint{4.249032in}{1.087043in}}{\pgfqpoint{4.254238in}{1.093149in}}%
\pgfpathquadraticcurveto{\pgfqpoint{4.259443in}{1.099255in}}{\pgfqpoint{4.259443in}{1.109378in}}%
\pgfpathlineto{\pgfqpoint{4.259443in}{1.111683in}}%
\pgfpathlineto{\pgfqpoint{4.249141in}{1.111683in}}%
\pgfpathlineto{\pgfqpoint{4.249141in}{1.111683in}}%
\pgfpathclose%
\pgfpathmoveto{\pgfqpoint{4.269800in}{1.115970in}}%
\pgfpathlineto{\pgfqpoint{4.269800in}{1.080000in}}%
\pgfpathlineto{\pgfqpoint{4.259443in}{1.080000in}}%
\pgfpathlineto{\pgfqpoint{4.259443in}{1.089564in}}%
\pgfpathquadraticcurveto{\pgfqpoint{4.255895in}{1.083837in}}{\pgfqpoint{4.250600in}{1.081099in}}%
\pgfpathquadraticcurveto{\pgfqpoint{4.245304in}{1.078361in}}{\pgfqpoint{4.237649in}{1.078361in}}%
\pgfpathquadraticcurveto{\pgfqpoint{4.227976in}{1.078361in}}{\pgfqpoint{4.222248in}{1.083801in}}%
\pgfpathquadraticcurveto{\pgfqpoint{4.216539in}{1.089240in}}{\pgfqpoint{4.216539in}{1.098354in}}%
\pgfpathquadraticcurveto{\pgfqpoint{4.216539in}{1.108982in}}{\pgfqpoint{4.223653in}{1.114385in}}%
\pgfpathquadraticcurveto{\pgfqpoint{4.230786in}{1.119789in}}{\pgfqpoint{4.244908in}{1.119789in}}%
\pgfpathlineto{\pgfqpoint{4.259443in}{1.119789in}}%
\pgfpathlineto{\pgfqpoint{4.259443in}{1.120816in}}%
\pgfpathquadraticcurveto{\pgfqpoint{4.259443in}{1.127966in}}{\pgfqpoint{4.254742in}{1.131875in}}%
\pgfpathquadraticcurveto{\pgfqpoint{4.250041in}{1.135784in}}{\pgfqpoint{4.241539in}{1.135784in}}%
\pgfpathquadraticcurveto{\pgfqpoint{4.236136in}{1.135784in}}{\pgfqpoint{4.231002in}{1.134487in}}%
\pgfpathquadraticcurveto{\pgfqpoint{4.225887in}{1.133190in}}{\pgfqpoint{4.221168in}{1.130596in}}%
\pgfpathlineto{\pgfqpoint{4.221168in}{1.140179in}}%
\pgfpathquadraticcurveto{\pgfqpoint{4.226841in}{1.142376in}}{\pgfqpoint{4.232191in}{1.143457in}}%
\pgfpathquadraticcurveto{\pgfqpoint{4.237541in}{1.144556in}}{\pgfqpoint{4.242602in}{1.144556in}}%
\pgfpathquadraticcurveto{\pgfqpoint{4.256291in}{1.144556in}}{\pgfqpoint{4.263046in}{1.137459in}}%
\pgfpathquadraticcurveto{\pgfqpoint{4.269800in}{1.130380in}}{\pgfqpoint{4.269800in}{1.115970in}}%
\pgfpathlineto{\pgfqpoint{4.269800in}{1.115970in}}%
\pgfpathclose%
\pgfpathmoveto{\pgfqpoint{4.343542in}{1.118060in}}%
\pgfpathlineto{\pgfqpoint{4.343542in}{1.080000in}}%
\pgfpathlineto{\pgfqpoint{4.333185in}{1.080000in}}%
\pgfpathlineto{\pgfqpoint{4.333185in}{1.117717in}}%
\pgfpathquadraticcurveto{\pgfqpoint{4.333185in}{1.126669in}}{\pgfqpoint{4.329691in}{1.131100in}}%
\pgfpathquadraticcurveto{\pgfqpoint{4.326197in}{1.135549in}}{\pgfqpoint{4.319226in}{1.135549in}}%
\pgfpathquadraticcurveto{\pgfqpoint{4.310832in}{1.135549in}}{\pgfqpoint{4.305987in}{1.130200in}}%
\pgfpathquadraticcurveto{\pgfqpoint{4.301142in}{1.124868in}}{\pgfqpoint{4.301142in}{1.115628in}}%
\pgfpathlineto{\pgfqpoint{4.301142in}{1.080000in}}%
\pgfpathlineto{\pgfqpoint{4.290731in}{1.080000in}}%
\pgfpathlineto{\pgfqpoint{4.290731in}{1.143043in}}%
\pgfpathlineto{\pgfqpoint{4.301142in}{1.143043in}}%
\pgfpathlineto{\pgfqpoint{4.301142in}{1.133244in}}%
\pgfpathquadraticcurveto{\pgfqpoint{4.304870in}{1.138936in}}{\pgfqpoint{4.309896in}{1.141746in}}%
\pgfpathquadraticcurveto{\pgfqpoint{4.314939in}{1.144556in}}{\pgfqpoint{4.321531in}{1.144556in}}%
\pgfpathquadraticcurveto{\pgfqpoint{4.332393in}{1.144556in}}{\pgfqpoint{4.337958in}{1.137837in}}%
\pgfpathquadraticcurveto{\pgfqpoint{4.343542in}{1.131118in}}{\pgfqpoint{4.343542in}{1.118060in}}%
\pgfpathlineto{\pgfqpoint{4.343542in}{1.118060in}}%
\pgfpathclose%
\pgfpathmoveto{\pgfqpoint{4.405666in}{1.133478in}}%
\pgfpathlineto{\pgfqpoint{4.405666in}{1.167593in}}%
\pgfpathlineto{\pgfqpoint{4.416023in}{1.167593in}}%
\pgfpathlineto{\pgfqpoint{4.416023in}{1.080000in}}%
\pgfpathlineto{\pgfqpoint{4.405666in}{1.080000in}}%
\pgfpathlineto{\pgfqpoint{4.405666in}{1.089456in}}%
\pgfpathquadraticcurveto{\pgfqpoint{4.402406in}{1.083837in}}{\pgfqpoint{4.397417in}{1.081099in}}%
\pgfpathquadraticcurveto{\pgfqpoint{4.392445in}{1.078361in}}{\pgfqpoint{4.385457in}{1.078361in}}%
\pgfpathquadraticcurveto{\pgfqpoint{4.374037in}{1.078361in}}{\pgfqpoint{4.366850in}{1.087475in}}%
\pgfpathquadraticcurveto{\pgfqpoint{4.359681in}{1.096607in}}{\pgfqpoint{4.359681in}{1.111467in}}%
\pgfpathquadraticcurveto{\pgfqpoint{4.359681in}{1.126327in}}{\pgfqpoint{4.366850in}{1.135441in}}%
\pgfpathquadraticcurveto{\pgfqpoint{4.374037in}{1.144556in}}{\pgfqpoint{4.385457in}{1.144556in}}%
\pgfpathquadraticcurveto{\pgfqpoint{4.392445in}{1.144556in}}{\pgfqpoint{4.397417in}{1.141818in}}%
\pgfpathquadraticcurveto{\pgfqpoint{4.402406in}{1.139098in}}{\pgfqpoint{4.405666in}{1.133478in}}%
\pgfpathlineto{\pgfqpoint{4.405666in}{1.133478in}}%
\pgfpathclose%
\pgfpathmoveto{\pgfqpoint{4.370380in}{1.111467in}}%
\pgfpathquadraticcurveto{\pgfqpoint{4.370380in}{1.100048in}}{\pgfqpoint{4.375082in}{1.093545in}}%
\pgfpathquadraticcurveto{\pgfqpoint{4.379783in}{1.087043in}}{\pgfqpoint{4.387996in}{1.087043in}}%
\pgfpathquadraticcurveto{\pgfqpoint{4.396210in}{1.087043in}}{\pgfqpoint{4.400929in}{1.093545in}}%
\pgfpathquadraticcurveto{\pgfqpoint{4.405666in}{1.100048in}}{\pgfqpoint{4.405666in}{1.111467in}}%
\pgfpathquadraticcurveto{\pgfqpoint{4.405666in}{1.122887in}}{\pgfqpoint{4.400929in}{1.129389in}}%
\pgfpathquadraticcurveto{\pgfqpoint{4.396210in}{1.135892in}}{\pgfqpoint{4.387996in}{1.135892in}}%
\pgfpathquadraticcurveto{\pgfqpoint{4.379783in}{1.135892in}}{\pgfqpoint{4.375082in}{1.129389in}}%
\pgfpathquadraticcurveto{\pgfqpoint{4.370380in}{1.122887in}}{\pgfqpoint{4.370380in}{1.111467in}}%
\pgfpathlineto{\pgfqpoint{4.370380in}{1.111467in}}%
\pgfpathclose%
\pgfusepath{fill}%
\end{pgfscope}%
\begin{pgfscope}%
\pgfsetbuttcap%
\pgfsetroundjoin%
\definecolor{currentfill}{rgb}{0.309804,0.372549,0.419608}%
\pgfsetfillcolor{currentfill}%
\pgfsetlinewidth{0.000000pt}%
\definecolor{currentstroke}{rgb}{0.309804,0.372549,0.419608}%
\pgfsetstrokecolor{currentstroke}%
\pgfsetdash{}{0pt}%
\pgfpathmoveto{\pgfqpoint{4.521037in}{1.111683in}}%
\pgfpathquadraticcurveto{\pgfqpoint{4.508483in}{1.111683in}}{\pgfqpoint{4.503638in}{1.108819in}}%
\pgfpathquadraticcurveto{\pgfqpoint{4.498792in}{1.105956in}}{\pgfqpoint{4.498792in}{1.099021in}}%
\pgfpathquadraticcurveto{\pgfqpoint{4.498792in}{1.093509in}}{\pgfqpoint{4.502431in}{1.090267in}}%
\pgfpathquadraticcurveto{\pgfqpoint{4.506069in}{1.087043in}}{\pgfqpoint{4.512302in}{1.087043in}}%
\pgfpathquadraticcurveto{\pgfqpoint{4.520929in}{1.087043in}}{\pgfqpoint{4.526135in}{1.093149in}}%
\pgfpathquadraticcurveto{\pgfqpoint{4.531340in}{1.099255in}}{\pgfqpoint{4.531340in}{1.109378in}}%
\pgfpathlineto{\pgfqpoint{4.531340in}{1.111683in}}%
\pgfpathlineto{\pgfqpoint{4.521037in}{1.111683in}}%
\pgfpathlineto{\pgfqpoint{4.521037in}{1.111683in}}%
\pgfpathclose%
\pgfpathmoveto{\pgfqpoint{4.541697in}{1.115970in}}%
\pgfpathlineto{\pgfqpoint{4.541697in}{1.080000in}}%
\pgfpathlineto{\pgfqpoint{4.531340in}{1.080000in}}%
\pgfpathlineto{\pgfqpoint{4.531340in}{1.089564in}}%
\pgfpathquadraticcurveto{\pgfqpoint{4.527792in}{1.083837in}}{\pgfqpoint{4.522496in}{1.081099in}}%
\pgfpathquadraticcurveto{\pgfqpoint{4.517201in}{1.078361in}}{\pgfqpoint{4.509546in}{1.078361in}}%
\pgfpathquadraticcurveto{\pgfqpoint{4.499873in}{1.078361in}}{\pgfqpoint{4.494145in}{1.083801in}}%
\pgfpathquadraticcurveto{\pgfqpoint{4.488435in}{1.089240in}}{\pgfqpoint{4.488435in}{1.098354in}}%
\pgfpathquadraticcurveto{\pgfqpoint{4.488435in}{1.108982in}}{\pgfqpoint{4.495550in}{1.114385in}}%
\pgfpathquadraticcurveto{\pgfqpoint{4.502683in}{1.119789in}}{\pgfqpoint{4.516805in}{1.119789in}}%
\pgfpathlineto{\pgfqpoint{4.531340in}{1.119789in}}%
\pgfpathlineto{\pgfqpoint{4.531340in}{1.120816in}}%
\pgfpathquadraticcurveto{\pgfqpoint{4.531340in}{1.127966in}}{\pgfqpoint{4.526639in}{1.131875in}}%
\pgfpathquadraticcurveto{\pgfqpoint{4.521938in}{1.135784in}}{\pgfqpoint{4.513436in}{1.135784in}}%
\pgfpathquadraticcurveto{\pgfqpoint{4.508033in}{1.135784in}}{\pgfqpoint{4.502899in}{1.134487in}}%
\pgfpathquadraticcurveto{\pgfqpoint{4.497784in}{1.133190in}}{\pgfqpoint{4.493065in}{1.130596in}}%
\pgfpathlineto{\pgfqpoint{4.493065in}{1.140179in}}%
\pgfpathquadraticcurveto{\pgfqpoint{4.498738in}{1.142376in}}{\pgfqpoint{4.504088in}{1.143457in}}%
\pgfpathquadraticcurveto{\pgfqpoint{4.509438in}{1.144556in}}{\pgfqpoint{4.514499in}{1.144556in}}%
\pgfpathquadraticcurveto{\pgfqpoint{4.528188in}{1.144556in}}{\pgfqpoint{4.534943in}{1.137459in}}%
\pgfpathquadraticcurveto{\pgfqpoint{4.541697in}{1.130380in}}{\pgfqpoint{4.541697in}{1.115970in}}%
\pgfpathlineto{\pgfqpoint{4.541697in}{1.115970in}}%
\pgfpathclose%
\pgfusepath{fill}%
\end{pgfscope}%
\begin{pgfscope}%
\pgfsetbuttcap%
\pgfsetroundjoin%
\definecolor{currentfill}{rgb}{0.309804,0.372549,0.419608}%
\pgfsetfillcolor{currentfill}%
\pgfsetlinewidth{0.000000pt}%
\definecolor{currentstroke}{rgb}{0.309804,0.372549,0.419608}%
\pgfsetstrokecolor{currentstroke}%
\pgfsetdash{}{0pt}%
\pgfpathmoveto{\pgfqpoint{4.655215in}{1.143043in}}%
\pgfpathlineto{\pgfqpoint{4.655215in}{1.080000in}}%
\pgfpathlineto{\pgfqpoint{4.644804in}{1.080000in}}%
\pgfpathlineto{\pgfqpoint{4.644804in}{1.134991in}}%
\pgfpathlineto{\pgfqpoint{4.616381in}{1.134991in}}%
\pgfpathlineto{\pgfqpoint{4.616381in}{1.080000in}}%
\pgfpathlineto{\pgfqpoint{4.605970in}{1.080000in}}%
\pgfpathlineto{\pgfqpoint{4.605970in}{1.134991in}}%
\pgfpathlineto{\pgfqpoint{4.596063in}{1.134991in}}%
\pgfpathlineto{\pgfqpoint{4.596063in}{1.143043in}}%
\pgfpathlineto{\pgfqpoint{4.605970in}{1.143043in}}%
\pgfpathlineto{\pgfqpoint{4.605970in}{1.147438in}}%
\pgfpathquadraticcurveto{\pgfqpoint{4.605970in}{1.157740in}}{\pgfqpoint{4.610833in}{1.162658in}}%
\pgfpathquadraticcurveto{\pgfqpoint{4.615715in}{1.167593in}}{\pgfqpoint{4.625783in}{1.167593in}}%
\pgfpathlineto{\pgfqpoint{4.636194in}{1.167593in}}%
\pgfpathlineto{\pgfqpoint{4.636194in}{1.158965in}}%
\pgfpathlineto{\pgfqpoint{4.626288in}{1.158965in}}%
\pgfpathquadraticcurveto{\pgfqpoint{4.620722in}{1.158965in}}{\pgfqpoint{4.618543in}{1.156714in}}%
\pgfpathquadraticcurveto{\pgfqpoint{4.616381in}{1.154462in}}{\pgfqpoint{4.616381in}{1.148608in}}%
\pgfpathlineto{\pgfqpoint{4.616381in}{1.143043in}}%
\pgfpathlineto{\pgfqpoint{4.655215in}{1.143043in}}%
\pgfpathlineto{\pgfqpoint{4.655215in}{1.143043in}}%
\pgfpathclose%
\pgfpathmoveto{\pgfqpoint{4.644804in}{1.167467in}}%
\pgfpathlineto{\pgfqpoint{4.655215in}{1.167467in}}%
\pgfpathlineto{\pgfqpoint{4.655215in}{1.154354in}}%
\pgfpathlineto{\pgfqpoint{4.644804in}{1.154354in}}%
\pgfpathlineto{\pgfqpoint{4.644804in}{1.167467in}}%
\pgfpathlineto{\pgfqpoint{4.644804in}{1.167467in}}%
\pgfpathclose%
\pgfpathmoveto{\pgfqpoint{4.676884in}{1.167593in}}%
\pgfpathlineto{\pgfqpoint{4.687241in}{1.167593in}}%
\pgfpathlineto{\pgfqpoint{4.687241in}{1.080000in}}%
\pgfpathlineto{\pgfqpoint{4.676884in}{1.080000in}}%
\pgfpathlineto{\pgfqpoint{4.676884in}{1.167593in}}%
\pgfpathlineto{\pgfqpoint{4.676884in}{1.167593in}}%
\pgfpathclose%
\pgfpathmoveto{\pgfqpoint{4.708910in}{1.167593in}}%
\pgfpathlineto{\pgfqpoint{4.719267in}{1.167593in}}%
\pgfpathlineto{\pgfqpoint{4.719267in}{1.080000in}}%
\pgfpathlineto{\pgfqpoint{4.708910in}{1.080000in}}%
\pgfpathlineto{\pgfqpoint{4.708910in}{1.167593in}}%
\pgfpathlineto{\pgfqpoint{4.708910in}{1.167593in}}%
\pgfpathclose%
\pgfpathmoveto{\pgfqpoint{4.794864in}{1.114115in}}%
\pgfpathlineto{\pgfqpoint{4.794864in}{1.109054in}}%
\pgfpathlineto{\pgfqpoint{4.747239in}{1.109054in}}%
\pgfpathquadraticcurveto{\pgfqpoint{4.747924in}{1.098354in}}{\pgfqpoint{4.753688in}{1.092753in}}%
\pgfpathquadraticcurveto{\pgfqpoint{4.759452in}{1.087151in}}{\pgfqpoint{4.769755in}{1.087151in}}%
\pgfpathquadraticcurveto{\pgfqpoint{4.775717in}{1.087151in}}{\pgfqpoint{4.781318in}{1.088610in}}%
\pgfpathquadraticcurveto{\pgfqpoint{4.786920in}{1.090069in}}{\pgfqpoint{4.792450in}{1.093005in}}%
\pgfpathlineto{\pgfqpoint{4.792450in}{1.083206in}}%
\pgfpathquadraticcurveto{\pgfqpoint{4.786866in}{1.080847in}}{\pgfqpoint{4.781012in}{1.079604in}}%
\pgfpathquadraticcurveto{\pgfqpoint{4.775158in}{1.078361in}}{\pgfqpoint{4.769142in}{1.078361in}}%
\pgfpathquadraticcurveto{\pgfqpoint{4.754048in}{1.078361in}}{\pgfqpoint{4.745240in}{1.087133in}}%
\pgfpathquadraticcurveto{\pgfqpoint{4.736432in}{1.095923in}}{\pgfqpoint{4.736432in}{1.110909in}}%
\pgfpathquadraticcurveto{\pgfqpoint{4.736432in}{1.126381in}}{\pgfqpoint{4.744790in}{1.135459in}}%
\pgfpathquadraticcurveto{\pgfqpoint{4.753147in}{1.144556in}}{\pgfqpoint{4.767341in}{1.144556in}}%
\pgfpathquadraticcurveto{\pgfqpoint{4.780058in}{1.144556in}}{\pgfqpoint{4.787461in}{1.136360in}}%
\pgfpathquadraticcurveto{\pgfqpoint{4.794864in}{1.128183in}}{\pgfqpoint{4.794864in}{1.114115in}}%
\pgfpathlineto{\pgfqpoint{4.794864in}{1.114115in}}%
\pgfpathclose%
\pgfpathmoveto{\pgfqpoint{4.784507in}{1.117159in}}%
\pgfpathquadraticcurveto{\pgfqpoint{4.784398in}{1.125643in}}{\pgfqpoint{4.779751in}{1.130704in}}%
\pgfpathquadraticcurveto{\pgfqpoint{4.775104in}{1.135784in}}{\pgfqpoint{4.767449in}{1.135784in}}%
\pgfpathquadraticcurveto{\pgfqpoint{4.758785in}{1.135784in}}{\pgfqpoint{4.753580in}{1.130884in}}%
\pgfpathquadraticcurveto{\pgfqpoint{4.748374in}{1.125985in}}{\pgfqpoint{4.747582in}{1.117087in}}%
\pgfpathlineto{\pgfqpoint{4.784507in}{1.117159in}}%
\pgfpathlineto{\pgfqpoint{4.784507in}{1.117159in}}%
\pgfpathclose%
\pgfpathmoveto{\pgfqpoint{4.853349in}{1.133478in}}%
\pgfpathlineto{\pgfqpoint{4.853349in}{1.167593in}}%
\pgfpathlineto{\pgfqpoint{4.863706in}{1.167593in}}%
\pgfpathlineto{\pgfqpoint{4.863706in}{1.080000in}}%
\pgfpathlineto{\pgfqpoint{4.853349in}{1.080000in}}%
\pgfpathlineto{\pgfqpoint{4.853349in}{1.089456in}}%
\pgfpathquadraticcurveto{\pgfqpoint{4.850089in}{1.083837in}}{\pgfqpoint{4.845099in}{1.081099in}}%
\pgfpathquadraticcurveto{\pgfqpoint{4.840128in}{1.078361in}}{\pgfqpoint{4.833139in}{1.078361in}}%
\pgfpathquadraticcurveto{\pgfqpoint{4.821720in}{1.078361in}}{\pgfqpoint{4.814533in}{1.087475in}}%
\pgfpathquadraticcurveto{\pgfqpoint{4.807364in}{1.096607in}}{\pgfqpoint{4.807364in}{1.111467in}}%
\pgfpathquadraticcurveto{\pgfqpoint{4.807364in}{1.126327in}}{\pgfqpoint{4.814533in}{1.135441in}}%
\pgfpathquadraticcurveto{\pgfqpoint{4.821720in}{1.144556in}}{\pgfqpoint{4.833139in}{1.144556in}}%
\pgfpathquadraticcurveto{\pgfqpoint{4.840128in}{1.144556in}}{\pgfqpoint{4.845099in}{1.141818in}}%
\pgfpathquadraticcurveto{\pgfqpoint{4.850089in}{1.139098in}}{\pgfqpoint{4.853349in}{1.133478in}}%
\pgfpathlineto{\pgfqpoint{4.853349in}{1.133478in}}%
\pgfpathclose%
\pgfpathmoveto{\pgfqpoint{4.818063in}{1.111467in}}%
\pgfpathquadraticcurveto{\pgfqpoint{4.818063in}{1.100048in}}{\pgfqpoint{4.822764in}{1.093545in}}%
\pgfpathquadraticcurveto{\pgfqpoint{4.827466in}{1.087043in}}{\pgfqpoint{4.835679in}{1.087043in}}%
\pgfpathquadraticcurveto{\pgfqpoint{4.843893in}{1.087043in}}{\pgfqpoint{4.848612in}{1.093545in}}%
\pgfpathquadraticcurveto{\pgfqpoint{4.853349in}{1.100048in}}{\pgfqpoint{4.853349in}{1.111467in}}%
\pgfpathquadraticcurveto{\pgfqpoint{4.853349in}{1.122887in}}{\pgfqpoint{4.848612in}{1.129389in}}%
\pgfpathquadraticcurveto{\pgfqpoint{4.843893in}{1.135892in}}{\pgfqpoint{4.835679in}{1.135892in}}%
\pgfpathquadraticcurveto{\pgfqpoint{4.827466in}{1.135892in}}{\pgfqpoint{4.822764in}{1.129389in}}%
\pgfpathquadraticcurveto{\pgfqpoint{4.818063in}{1.122887in}}{\pgfqpoint{4.818063in}{1.111467in}}%
\pgfpathlineto{\pgfqpoint{4.818063in}{1.111467in}}%
\pgfpathclose%
\pgfusepath{fill}%
\end{pgfscope}%
\begin{pgfscope}%
\pgfsetbuttcap%
\pgfsetroundjoin%
\definecolor{currentfill}{rgb}{0.309804,0.372549,0.419608}%
\pgfsetfillcolor{currentfill}%
\pgfsetlinewidth{0.000000pt}%
\definecolor{currentstroke}{rgb}{0.309804,0.372549,0.419608}%
\pgfsetstrokecolor{currentstroke}%
\pgfsetdash{}{0pt}%
\pgfpathmoveto{\pgfqpoint{4.985257in}{1.130938in}}%
\pgfpathquadraticcurveto{\pgfqpoint{4.989148in}{1.137927in}}{\pgfqpoint{4.994551in}{1.141241in}}%
\pgfpathquadraticcurveto{\pgfqpoint{4.999955in}{1.144556in}}{\pgfqpoint{5.007268in}{1.144556in}}%
\pgfpathquadraticcurveto{\pgfqpoint{5.017120in}{1.144556in}}{\pgfqpoint{5.022470in}{1.137657in}}%
\pgfpathquadraticcurveto{\pgfqpoint{5.027820in}{1.130776in}}{\pgfqpoint{5.027820in}{1.118060in}}%
\pgfpathlineto{\pgfqpoint{5.027820in}{1.080000in}}%
\pgfpathlineto{\pgfqpoint{5.017409in}{1.080000in}}%
\pgfpathlineto{\pgfqpoint{5.017409in}{1.117717in}}%
\pgfpathquadraticcurveto{\pgfqpoint{5.017409in}{1.126778in}}{\pgfqpoint{5.014184in}{1.131155in}}%
\pgfpathquadraticcurveto{\pgfqpoint{5.010978in}{1.135549in}}{\pgfqpoint{5.004404in}{1.135549in}}%
\pgfpathquadraticcurveto{\pgfqpoint{4.996352in}{1.135549in}}{\pgfqpoint{4.991669in}{1.130200in}}%
\pgfpathquadraticcurveto{\pgfqpoint{4.987004in}{1.124868in}}{\pgfqpoint{4.987004in}{1.115628in}}%
\pgfpathlineto{\pgfqpoint{4.987004in}{1.080000in}}%
\pgfpathlineto{\pgfqpoint{4.976593in}{1.080000in}}%
\pgfpathlineto{\pgfqpoint{4.976593in}{1.117717in}}%
\pgfpathquadraticcurveto{\pgfqpoint{4.976593in}{1.126832in}}{\pgfqpoint{4.973387in}{1.131191in}}%
\pgfpathquadraticcurveto{\pgfqpoint{4.970181in}{1.135549in}}{\pgfqpoint{4.963480in}{1.135549in}}%
\pgfpathquadraticcurveto{\pgfqpoint{4.955537in}{1.135549in}}{\pgfqpoint{4.950854in}{1.130182in}}%
\pgfpathquadraticcurveto{\pgfqpoint{4.946189in}{1.124814in}}{\pgfqpoint{4.946189in}{1.115628in}}%
\pgfpathlineto{\pgfqpoint{4.946189in}{1.080000in}}%
\pgfpathlineto{\pgfqpoint{4.935778in}{1.080000in}}%
\pgfpathlineto{\pgfqpoint{4.935778in}{1.143043in}}%
\pgfpathlineto{\pgfqpoint{4.946189in}{1.143043in}}%
\pgfpathlineto{\pgfqpoint{4.946189in}{1.133244in}}%
\pgfpathquadraticcurveto{\pgfqpoint{4.949737in}{1.139044in}}{\pgfqpoint{4.954690in}{1.141800in}}%
\pgfpathquadraticcurveto{\pgfqpoint{4.959644in}{1.144556in}}{\pgfqpoint{4.966452in}{1.144556in}}%
\pgfpathquadraticcurveto{\pgfqpoint{4.973333in}{1.144556in}}{\pgfqpoint{4.978142in}{1.141061in}}%
\pgfpathquadraticcurveto{\pgfqpoint{4.982951in}{1.137585in}}{\pgfqpoint{4.985257in}{1.130938in}}%
\pgfpathlineto{\pgfqpoint{4.985257in}{1.130938in}}%
\pgfpathclose%
\pgfpathmoveto{\pgfqpoint{5.077119in}{1.111683in}}%
\pgfpathquadraticcurveto{\pgfqpoint{5.064564in}{1.111683in}}{\pgfqpoint{5.059719in}{1.108819in}}%
\pgfpathquadraticcurveto{\pgfqpoint{5.054874in}{1.105956in}}{\pgfqpoint{5.054874in}{1.099021in}}%
\pgfpathquadraticcurveto{\pgfqpoint{5.054874in}{1.093509in}}{\pgfqpoint{5.058512in}{1.090267in}}%
\pgfpathquadraticcurveto{\pgfqpoint{5.062151in}{1.087043in}}{\pgfqpoint{5.068383in}{1.087043in}}%
\pgfpathquadraticcurveto{\pgfqpoint{5.077011in}{1.087043in}}{\pgfqpoint{5.082216in}{1.093149in}}%
\pgfpathquadraticcurveto{\pgfqpoint{5.087422in}{1.099255in}}{\pgfqpoint{5.087422in}{1.109378in}}%
\pgfpathlineto{\pgfqpoint{5.087422in}{1.111683in}}%
\pgfpathlineto{\pgfqpoint{5.077119in}{1.111683in}}%
\pgfpathlineto{\pgfqpoint{5.077119in}{1.111683in}}%
\pgfpathclose%
\pgfpathmoveto{\pgfqpoint{5.097779in}{1.115970in}}%
\pgfpathlineto{\pgfqpoint{5.097779in}{1.080000in}}%
\pgfpathlineto{\pgfqpoint{5.087422in}{1.080000in}}%
\pgfpathlineto{\pgfqpoint{5.087422in}{1.089564in}}%
\pgfpathquadraticcurveto{\pgfqpoint{5.083873in}{1.083837in}}{\pgfqpoint{5.078578in}{1.081099in}}%
\pgfpathquadraticcurveto{\pgfqpoint{5.073282in}{1.078361in}}{\pgfqpoint{5.065627in}{1.078361in}}%
\pgfpathquadraticcurveto{\pgfqpoint{5.055955in}{1.078361in}}{\pgfqpoint{5.050227in}{1.083801in}}%
\pgfpathquadraticcurveto{\pgfqpoint{5.044517in}{1.089240in}}{\pgfqpoint{5.044517in}{1.098354in}}%
\pgfpathquadraticcurveto{\pgfqpoint{5.044517in}{1.108982in}}{\pgfqpoint{5.051632in}{1.114385in}}%
\pgfpathquadraticcurveto{\pgfqpoint{5.058765in}{1.119789in}}{\pgfqpoint{5.072886in}{1.119789in}}%
\pgfpathlineto{\pgfqpoint{5.087422in}{1.119789in}}%
\pgfpathlineto{\pgfqpoint{5.087422in}{1.120816in}}%
\pgfpathquadraticcurveto{\pgfqpoint{5.087422in}{1.127966in}}{\pgfqpoint{5.082721in}{1.131875in}}%
\pgfpathquadraticcurveto{\pgfqpoint{5.078020in}{1.135784in}}{\pgfqpoint{5.069518in}{1.135784in}}%
\pgfpathquadraticcurveto{\pgfqpoint{5.064114in}{1.135784in}}{\pgfqpoint{5.058981in}{1.134487in}}%
\pgfpathquadraticcurveto{\pgfqpoint{5.053865in}{1.133190in}}{\pgfqpoint{5.049146in}{1.130596in}}%
\pgfpathlineto{\pgfqpoint{5.049146in}{1.140179in}}%
\pgfpathquadraticcurveto{\pgfqpoint{5.054820in}{1.142376in}}{\pgfqpoint{5.060169in}{1.143457in}}%
\pgfpathquadraticcurveto{\pgfqpoint{5.065519in}{1.144556in}}{\pgfqpoint{5.070581in}{1.144556in}}%
\pgfpathquadraticcurveto{\pgfqpoint{5.084270in}{1.144556in}}{\pgfqpoint{5.091024in}{1.137459in}}%
\pgfpathquadraticcurveto{\pgfqpoint{5.097779in}{1.130380in}}{\pgfqpoint{5.097779in}{1.115970in}}%
\pgfpathlineto{\pgfqpoint{5.097779in}{1.115970in}}%
\pgfpathclose%
\pgfpathmoveto{\pgfqpoint{5.155634in}{1.133370in}}%
\pgfpathquadraticcurveto{\pgfqpoint{5.153887in}{1.134379in}}{\pgfqpoint{5.151833in}{1.134847in}}%
\pgfpathquadraticcurveto{\pgfqpoint{5.149780in}{1.135333in}}{\pgfqpoint{5.147312in}{1.135333in}}%
\pgfpathquadraticcurveto{\pgfqpoint{5.138522in}{1.135333in}}{\pgfqpoint{5.133821in}{1.129623in}}%
\pgfpathquadraticcurveto{\pgfqpoint{5.129120in}{1.123914in}}{\pgfqpoint{5.129120in}{1.113214in}}%
\pgfpathlineto{\pgfqpoint{5.129120in}{1.080000in}}%
\pgfpathlineto{\pgfqpoint{5.118709in}{1.080000in}}%
\pgfpathlineto{\pgfqpoint{5.118709in}{1.143043in}}%
\pgfpathlineto{\pgfqpoint{5.129120in}{1.143043in}}%
\pgfpathlineto{\pgfqpoint{5.129120in}{1.133244in}}%
\pgfpathquadraticcurveto{\pgfqpoint{5.132398in}{1.138990in}}{\pgfqpoint{5.137622in}{1.141764in}}%
\pgfpathquadraticcurveto{\pgfqpoint{5.142863in}{1.144556in}}{\pgfqpoint{5.150356in}{1.144556in}}%
\pgfpathquadraticcurveto{\pgfqpoint{5.151419in}{1.144556in}}{\pgfqpoint{5.152716in}{1.144411in}}%
\pgfpathquadraticcurveto{\pgfqpoint{5.154013in}{1.144285in}}{\pgfqpoint{5.155580in}{1.143997in}}%
\pgfpathlineto{\pgfqpoint{5.155634in}{1.133370in}}%
\pgfpathlineto{\pgfqpoint{5.155634in}{1.133370in}}%
\pgfpathclose%
\pgfpathmoveto{\pgfqpoint{5.166099in}{1.167593in}}%
\pgfpathlineto{\pgfqpoint{5.176510in}{1.167593in}}%
\pgfpathlineto{\pgfqpoint{5.176510in}{1.115862in}}%
\pgfpathlineto{\pgfqpoint{5.207419in}{1.143043in}}%
\pgfpathlineto{\pgfqpoint{5.220640in}{1.143043in}}%
\pgfpathlineto{\pgfqpoint{5.187209in}{1.113557in}}%
\pgfpathlineto{\pgfqpoint{5.222063in}{1.080000in}}%
\pgfpathlineto{\pgfqpoint{5.208554in}{1.080000in}}%
\pgfpathlineto{\pgfqpoint{5.176510in}{1.110783in}}%
\pgfpathlineto{\pgfqpoint{5.176510in}{1.080000in}}%
\pgfpathlineto{\pgfqpoint{5.166099in}{1.080000in}}%
\pgfpathlineto{\pgfqpoint{5.166099in}{1.167593in}}%
\pgfpathlineto{\pgfqpoint{5.166099in}{1.167593in}}%
\pgfpathclose%
\pgfpathmoveto{\pgfqpoint{5.283070in}{1.114115in}}%
\pgfpathlineto{\pgfqpoint{5.283070in}{1.109054in}}%
\pgfpathlineto{\pgfqpoint{5.235446in}{1.109054in}}%
\pgfpathquadraticcurveto{\pgfqpoint{5.236130in}{1.098354in}}{\pgfqpoint{5.241894in}{1.092753in}}%
\pgfpathquadraticcurveto{\pgfqpoint{5.247658in}{1.087151in}}{\pgfqpoint{5.257961in}{1.087151in}}%
\pgfpathquadraticcurveto{\pgfqpoint{5.263923in}{1.087151in}}{\pgfqpoint{5.269525in}{1.088610in}}%
\pgfpathquadraticcurveto{\pgfqpoint{5.275127in}{1.090069in}}{\pgfqpoint{5.280656in}{1.093005in}}%
\pgfpathlineto{\pgfqpoint{5.280656in}{1.083206in}}%
\pgfpathquadraticcurveto{\pgfqpoint{5.275072in}{1.080847in}}{\pgfqpoint{5.269219in}{1.079604in}}%
\pgfpathquadraticcurveto{\pgfqpoint{5.263365in}{1.078361in}}{\pgfqpoint{5.257349in}{1.078361in}}%
\pgfpathquadraticcurveto{\pgfqpoint{5.242254in}{1.078361in}}{\pgfqpoint{5.233446in}{1.087133in}}%
\pgfpathquadraticcurveto{\pgfqpoint{5.224638in}{1.095923in}}{\pgfqpoint{5.224638in}{1.110909in}}%
\pgfpathquadraticcurveto{\pgfqpoint{5.224638in}{1.126381in}}{\pgfqpoint{5.232996in}{1.135459in}}%
\pgfpathquadraticcurveto{\pgfqpoint{5.241354in}{1.144556in}}{\pgfqpoint{5.255547in}{1.144556in}}%
\pgfpathquadraticcurveto{\pgfqpoint{5.268264in}{1.144556in}}{\pgfqpoint{5.275667in}{1.136360in}}%
\pgfpathquadraticcurveto{\pgfqpoint{5.283070in}{1.128183in}}{\pgfqpoint{5.283070in}{1.114115in}}%
\pgfpathlineto{\pgfqpoint{5.283070in}{1.114115in}}%
\pgfpathclose%
\pgfpathmoveto{\pgfqpoint{5.272713in}{1.117159in}}%
\pgfpathquadraticcurveto{\pgfqpoint{5.272605in}{1.125643in}}{\pgfqpoint{5.267958in}{1.130704in}}%
\pgfpathquadraticcurveto{\pgfqpoint{5.263311in}{1.135784in}}{\pgfqpoint{5.255655in}{1.135784in}}%
\pgfpathquadraticcurveto{\pgfqpoint{5.246992in}{1.135784in}}{\pgfqpoint{5.241786in}{1.130884in}}%
\pgfpathquadraticcurveto{\pgfqpoint{5.236581in}{1.125985in}}{\pgfqpoint{5.235788in}{1.117087in}}%
\pgfpathlineto{\pgfqpoint{5.272713in}{1.117159in}}%
\pgfpathlineto{\pgfqpoint{5.272713in}{1.117159in}}%
\pgfpathclose%
\pgfpathmoveto{\pgfqpoint{5.336602in}{1.133370in}}%
\pgfpathquadraticcurveto{\pgfqpoint{5.334855in}{1.134379in}}{\pgfqpoint{5.332801in}{1.134847in}}%
\pgfpathquadraticcurveto{\pgfqpoint{5.330748in}{1.135333in}}{\pgfqpoint{5.328280in}{1.135333in}}%
\pgfpathquadraticcurveto{\pgfqpoint{5.319490in}{1.135333in}}{\pgfqpoint{5.314789in}{1.129623in}}%
\pgfpathquadraticcurveto{\pgfqpoint{5.310088in}{1.123914in}}{\pgfqpoint{5.310088in}{1.113214in}}%
\pgfpathlineto{\pgfqpoint{5.310088in}{1.080000in}}%
\pgfpathlineto{\pgfqpoint{5.299677in}{1.080000in}}%
\pgfpathlineto{\pgfqpoint{5.299677in}{1.143043in}}%
\pgfpathlineto{\pgfqpoint{5.310088in}{1.143043in}}%
\pgfpathlineto{\pgfqpoint{5.310088in}{1.133244in}}%
\pgfpathquadraticcurveto{\pgfqpoint{5.313366in}{1.138990in}}{\pgfqpoint{5.318590in}{1.141764in}}%
\pgfpathquadraticcurveto{\pgfqpoint{5.323831in}{1.144556in}}{\pgfqpoint{5.331324in}{1.144556in}}%
\pgfpathquadraticcurveto{\pgfqpoint{5.332387in}{1.144556in}}{\pgfqpoint{5.333684in}{1.144411in}}%
\pgfpathquadraticcurveto{\pgfqpoint{5.334981in}{1.144285in}}{\pgfqpoint{5.336548in}{1.143997in}}%
\pgfpathlineto{\pgfqpoint{5.336602in}{1.133370in}}%
\pgfpathlineto{\pgfqpoint{5.336602in}{1.133370in}}%
\pgfpathclose%
\pgfpathmoveto{\pgfqpoint{5.338349in}{1.094302in}}%
\pgfpathlineto{\pgfqpoint{5.350237in}{1.094302in}}%
\pgfpathlineto{\pgfqpoint{5.350237in}{1.080000in}}%
\pgfpathlineto{\pgfqpoint{5.338349in}{1.080000in}}%
\pgfpathlineto{\pgfqpoint{5.338349in}{1.094302in}}%
\pgfpathlineto{\pgfqpoint{5.338349in}{1.094302in}}%
\pgfpathclose%
\pgfusepath{fill}%
\end{pgfscope}%
\begin{pgfscope}%
\pgfsetbuttcap%
\pgfsetroundjoin%
\definecolor{currentfill}{rgb}{0.309804,0.372549,0.419608}%
\pgfsetfillcolor{currentfill}%
\pgfsetlinewidth{0.000000pt}%
\definecolor{currentstroke}{rgb}{0.309804,0.372549,0.419608}%
\pgfsetstrokecolor{currentstroke}%
\pgfsetdash{}{0pt}%
\pgfpathmoveto{\pgfqpoint{5.445199in}{1.154696in}}%
\pgfpathlineto{\pgfqpoint{5.445199in}{1.089348in}}%
\pgfpathlineto{\pgfqpoint{5.458942in}{1.089348in}}%
\pgfpathquadraticcurveto{\pgfqpoint{5.476342in}{1.089348in}}{\pgfqpoint{5.484412in}{1.097220in}}%
\pgfpathquadraticcurveto{\pgfqpoint{5.492481in}{1.105109in}}{\pgfqpoint{5.492481in}{1.122112in}}%
\pgfpathquadraticcurveto{\pgfqpoint{5.492481in}{1.138990in}}{\pgfqpoint{5.484412in}{1.146843in}}%
\pgfpathquadraticcurveto{\pgfqpoint{5.476342in}{1.154696in}}{\pgfqpoint{5.458942in}{1.154696in}}%
\pgfpathlineto{\pgfqpoint{5.445199in}{1.154696in}}%
\pgfpathlineto{\pgfqpoint{5.445199in}{1.154696in}}%
\pgfpathclose%
\pgfpathmoveto{\pgfqpoint{5.433834in}{1.164045in}}%
\pgfpathlineto{\pgfqpoint{5.457195in}{1.164045in}}%
\pgfpathquadraticcurveto{\pgfqpoint{5.481620in}{1.164045in}}{\pgfqpoint{5.493039in}{1.153886in}}%
\pgfpathquadraticcurveto{\pgfqpoint{5.504477in}{1.143727in}}{\pgfqpoint{5.504477in}{1.122112in}}%
\pgfpathquadraticcurveto{\pgfqpoint{5.504477in}{1.100372in}}{\pgfqpoint{5.492985in}{1.090177in}}%
\pgfpathquadraticcurveto{\pgfqpoint{5.481512in}{1.080000in}}{\pgfqpoint{5.457195in}{1.080000in}}%
\pgfpathlineto{\pgfqpoint{5.433834in}{1.080000in}}%
\pgfpathlineto{\pgfqpoint{5.433834in}{1.164045in}}%
\pgfpathlineto{\pgfqpoint{5.433834in}{1.164045in}}%
\pgfpathclose%
\pgfpathmoveto{\pgfqpoint{5.522147in}{1.143043in}}%
\pgfpathlineto{\pgfqpoint{5.532504in}{1.143043in}}%
\pgfpathlineto{\pgfqpoint{5.532504in}{1.080000in}}%
\pgfpathlineto{\pgfqpoint{5.522147in}{1.080000in}}%
\pgfpathlineto{\pgfqpoint{5.522147in}{1.143043in}}%
\pgfpathlineto{\pgfqpoint{5.522147in}{1.143043in}}%
\pgfpathclose%
\pgfpathmoveto{\pgfqpoint{5.522147in}{1.167593in}}%
\pgfpathlineto{\pgfqpoint{5.532504in}{1.167593in}}%
\pgfpathlineto{\pgfqpoint{5.532504in}{1.154462in}}%
\pgfpathlineto{\pgfqpoint{5.522147in}{1.154462in}}%
\pgfpathlineto{\pgfqpoint{5.522147in}{1.167593in}}%
\pgfpathlineto{\pgfqpoint{5.522147in}{1.167593in}}%
\pgfpathclose%
\pgfpathmoveto{\pgfqpoint{5.582830in}{1.111683in}}%
\pgfpathquadraticcurveto{\pgfqpoint{5.570276in}{1.111683in}}{\pgfqpoint{5.565430in}{1.108819in}}%
\pgfpathquadraticcurveto{\pgfqpoint{5.560585in}{1.105956in}}{\pgfqpoint{5.560585in}{1.099021in}}%
\pgfpathquadraticcurveto{\pgfqpoint{5.560585in}{1.093509in}}{\pgfqpoint{5.564223in}{1.090267in}}%
\pgfpathquadraticcurveto{\pgfqpoint{5.567862in}{1.087043in}}{\pgfqpoint{5.574094in}{1.087043in}}%
\pgfpathquadraticcurveto{\pgfqpoint{5.582722in}{1.087043in}}{\pgfqpoint{5.587927in}{1.093149in}}%
\pgfpathquadraticcurveto{\pgfqpoint{5.593133in}{1.099255in}}{\pgfqpoint{5.593133in}{1.109378in}}%
\pgfpathlineto{\pgfqpoint{5.593133in}{1.111683in}}%
\pgfpathlineto{\pgfqpoint{5.582830in}{1.111683in}}%
\pgfpathlineto{\pgfqpoint{5.582830in}{1.111683in}}%
\pgfpathclose%
\pgfpathmoveto{\pgfqpoint{5.603490in}{1.115970in}}%
\pgfpathlineto{\pgfqpoint{5.603490in}{1.080000in}}%
\pgfpathlineto{\pgfqpoint{5.593133in}{1.080000in}}%
\pgfpathlineto{\pgfqpoint{5.593133in}{1.089564in}}%
\pgfpathquadraticcurveto{\pgfqpoint{5.589585in}{1.083837in}}{\pgfqpoint{5.584289in}{1.081099in}}%
\pgfpathquadraticcurveto{\pgfqpoint{5.578993in}{1.078361in}}{\pgfqpoint{5.571338in}{1.078361in}}%
\pgfpathquadraticcurveto{\pgfqpoint{5.561666in}{1.078361in}}{\pgfqpoint{5.555938in}{1.083801in}}%
\pgfpathquadraticcurveto{\pgfqpoint{5.550228in}{1.089240in}}{\pgfqpoint{5.550228in}{1.098354in}}%
\pgfpathquadraticcurveto{\pgfqpoint{5.550228in}{1.108982in}}{\pgfqpoint{5.557343in}{1.114385in}}%
\pgfpathquadraticcurveto{\pgfqpoint{5.564476in}{1.119789in}}{\pgfqpoint{5.578597in}{1.119789in}}%
\pgfpathlineto{\pgfqpoint{5.593133in}{1.119789in}}%
\pgfpathlineto{\pgfqpoint{5.593133in}{1.120816in}}%
\pgfpathquadraticcurveto{\pgfqpoint{5.593133in}{1.127966in}}{\pgfqpoint{5.588432in}{1.131875in}}%
\pgfpathquadraticcurveto{\pgfqpoint{5.583731in}{1.135784in}}{\pgfqpoint{5.575229in}{1.135784in}}%
\pgfpathquadraticcurveto{\pgfqpoint{5.569825in}{1.135784in}}{\pgfqpoint{5.564692in}{1.134487in}}%
\pgfpathquadraticcurveto{\pgfqpoint{5.559576in}{1.133190in}}{\pgfqpoint{5.554857in}{1.130596in}}%
\pgfpathlineto{\pgfqpoint{5.554857in}{1.140179in}}%
\pgfpathquadraticcurveto{\pgfqpoint{5.560531in}{1.142376in}}{\pgfqpoint{5.565881in}{1.143457in}}%
\pgfpathquadraticcurveto{\pgfqpoint{5.571230in}{1.144556in}}{\pgfqpoint{5.576292in}{1.144556in}}%
\pgfpathquadraticcurveto{\pgfqpoint{5.589981in}{1.144556in}}{\pgfqpoint{5.596735in}{1.137459in}}%
\pgfpathquadraticcurveto{\pgfqpoint{5.603490in}{1.130380in}}{\pgfqpoint{5.603490in}{1.115970in}}%
\pgfpathlineto{\pgfqpoint{5.603490in}{1.115970in}}%
\pgfpathclose%
\pgfpathmoveto{\pgfqpoint{5.673899in}{1.130938in}}%
\pgfpathquadraticcurveto{\pgfqpoint{5.677790in}{1.137927in}}{\pgfqpoint{5.683194in}{1.141241in}}%
\pgfpathquadraticcurveto{\pgfqpoint{5.688597in}{1.144556in}}{\pgfqpoint{5.695910in}{1.144556in}}%
\pgfpathquadraticcurveto{\pgfqpoint{5.705763in}{1.144556in}}{\pgfqpoint{5.711113in}{1.137657in}}%
\pgfpathquadraticcurveto{\pgfqpoint{5.716462in}{1.130776in}}{\pgfqpoint{5.716462in}{1.118060in}}%
\pgfpathlineto{\pgfqpoint{5.716462in}{1.080000in}}%
\pgfpathlineto{\pgfqpoint{5.706051in}{1.080000in}}%
\pgfpathlineto{\pgfqpoint{5.706051in}{1.117717in}}%
\pgfpathquadraticcurveto{\pgfqpoint{5.706051in}{1.126778in}}{\pgfqpoint{5.702827in}{1.131155in}}%
\pgfpathquadraticcurveto{\pgfqpoint{5.699621in}{1.135549in}}{\pgfqpoint{5.693046in}{1.135549in}}%
\pgfpathquadraticcurveto{\pgfqpoint{5.684995in}{1.135549in}}{\pgfqpoint{5.680312in}{1.130200in}}%
\pgfpathquadraticcurveto{\pgfqpoint{5.675647in}{1.124868in}}{\pgfqpoint{5.675647in}{1.115628in}}%
\pgfpathlineto{\pgfqpoint{5.675647in}{1.080000in}}%
\pgfpathlineto{\pgfqpoint{5.665236in}{1.080000in}}%
\pgfpathlineto{\pgfqpoint{5.665236in}{1.117717in}}%
\pgfpathquadraticcurveto{\pgfqpoint{5.665236in}{1.126832in}}{\pgfqpoint{5.662029in}{1.131191in}}%
\pgfpathquadraticcurveto{\pgfqpoint{5.658823in}{1.135549in}}{\pgfqpoint{5.652123in}{1.135549in}}%
\pgfpathquadraticcurveto{\pgfqpoint{5.644179in}{1.135549in}}{\pgfqpoint{5.639496in}{1.130182in}}%
\pgfpathquadraticcurveto{\pgfqpoint{5.634831in}{1.124814in}}{\pgfqpoint{5.634831in}{1.115628in}}%
\pgfpathlineto{\pgfqpoint{5.634831in}{1.080000in}}%
\pgfpathlineto{\pgfqpoint{5.624420in}{1.080000in}}%
\pgfpathlineto{\pgfqpoint{5.624420in}{1.143043in}}%
\pgfpathlineto{\pgfqpoint{5.634831in}{1.143043in}}%
\pgfpathlineto{\pgfqpoint{5.634831in}{1.133244in}}%
\pgfpathquadraticcurveto{\pgfqpoint{5.638380in}{1.139044in}}{\pgfqpoint{5.643333in}{1.141800in}}%
\pgfpathquadraticcurveto{\pgfqpoint{5.648286in}{1.144556in}}{\pgfqpoint{5.655095in}{1.144556in}}%
\pgfpathquadraticcurveto{\pgfqpoint{5.661975in}{1.144556in}}{\pgfqpoint{5.666785in}{1.141061in}}%
\pgfpathquadraticcurveto{\pgfqpoint{5.671594in}{1.137585in}}{\pgfqpoint{5.673899in}{1.130938in}}%
\pgfpathlineto{\pgfqpoint{5.673899in}{1.130938in}}%
\pgfpathclose%
\pgfpathmoveto{\pgfqpoint{5.761529in}{1.135784in}}%
\pgfpathquadraticcurveto{\pgfqpoint{5.753207in}{1.135784in}}{\pgfqpoint{5.748362in}{1.129281in}}%
\pgfpathquadraticcurveto{\pgfqpoint{5.743516in}{1.122779in}}{\pgfqpoint{5.743516in}{1.111467in}}%
\pgfpathquadraticcurveto{\pgfqpoint{5.743516in}{1.100156in}}{\pgfqpoint{5.748326in}{1.093653in}}%
\pgfpathquadraticcurveto{\pgfqpoint{5.753153in}{1.087151in}}{\pgfqpoint{5.761529in}{1.087151in}}%
\pgfpathquadraticcurveto{\pgfqpoint{5.769814in}{1.087151in}}{\pgfqpoint{5.774641in}{1.093671in}}%
\pgfpathquadraticcurveto{\pgfqpoint{5.779487in}{1.100210in}}{\pgfqpoint{5.779487in}{1.111467in}}%
\pgfpathquadraticcurveto{\pgfqpoint{5.779487in}{1.122671in}}{\pgfqpoint{5.774641in}{1.129227in}}%
\pgfpathquadraticcurveto{\pgfqpoint{5.769814in}{1.135784in}}{\pgfqpoint{5.761529in}{1.135784in}}%
\pgfpathlineto{\pgfqpoint{5.761529in}{1.135784in}}%
\pgfpathclose%
\pgfpathmoveto{\pgfqpoint{5.761529in}{1.144556in}}%
\pgfpathquadraticcurveto{\pgfqpoint{5.775038in}{1.144556in}}{\pgfqpoint{5.782747in}{1.135766in}}%
\pgfpathquadraticcurveto{\pgfqpoint{5.790474in}{1.126994in}}{\pgfqpoint{5.790474in}{1.111467in}}%
\pgfpathquadraticcurveto{\pgfqpoint{5.790474in}{1.095995in}}{\pgfqpoint{5.782747in}{1.087169in}}%
\pgfpathquadraticcurveto{\pgfqpoint{5.775038in}{1.078361in}}{\pgfqpoint{5.761529in}{1.078361in}}%
\pgfpathquadraticcurveto{\pgfqpoint{5.747965in}{1.078361in}}{\pgfqpoint{5.740274in}{1.087169in}}%
\pgfpathquadraticcurveto{\pgfqpoint{5.732601in}{1.095995in}}{\pgfqpoint{5.732601in}{1.111467in}}%
\pgfpathquadraticcurveto{\pgfqpoint{5.732601in}{1.126994in}}{\pgfqpoint{5.740274in}{1.135766in}}%
\pgfpathquadraticcurveto{\pgfqpoint{5.747965in}{1.144556in}}{\pgfqpoint{5.761529in}{1.144556in}}%
\pgfpathlineto{\pgfqpoint{5.761529in}{1.144556in}}%
\pgfpathclose%
\pgfpathmoveto{\pgfqpoint{5.860055in}{1.118060in}}%
\pgfpathlineto{\pgfqpoint{5.860055in}{1.080000in}}%
\pgfpathlineto{\pgfqpoint{5.849698in}{1.080000in}}%
\pgfpathlineto{\pgfqpoint{5.849698in}{1.117717in}}%
\pgfpathquadraticcurveto{\pgfqpoint{5.849698in}{1.126669in}}{\pgfqpoint{5.846204in}{1.131100in}}%
\pgfpathquadraticcurveto{\pgfqpoint{5.842709in}{1.135549in}}{\pgfqpoint{5.835739in}{1.135549in}}%
\pgfpathquadraticcurveto{\pgfqpoint{5.827345in}{1.135549in}}{\pgfqpoint{5.822500in}{1.130200in}}%
\pgfpathquadraticcurveto{\pgfqpoint{5.817654in}{1.124868in}}{\pgfqpoint{5.817654in}{1.115628in}}%
\pgfpathlineto{\pgfqpoint{5.817654in}{1.080000in}}%
\pgfpathlineto{\pgfqpoint{5.807243in}{1.080000in}}%
\pgfpathlineto{\pgfqpoint{5.807243in}{1.143043in}}%
\pgfpathlineto{\pgfqpoint{5.817654in}{1.143043in}}%
\pgfpathlineto{\pgfqpoint{5.817654in}{1.133244in}}%
\pgfpathquadraticcurveto{\pgfqpoint{5.821383in}{1.138936in}}{\pgfqpoint{5.826408in}{1.141746in}}%
\pgfpathquadraticcurveto{\pgfqpoint{5.831452in}{1.144556in}}{\pgfqpoint{5.838044in}{1.144556in}}%
\pgfpathquadraticcurveto{\pgfqpoint{5.848906in}{1.144556in}}{\pgfqpoint{5.854471in}{1.137837in}}%
\pgfpathquadraticcurveto{\pgfqpoint{5.860055in}{1.131118in}}{\pgfqpoint{5.860055in}{1.118060in}}%
\pgfpathlineto{\pgfqpoint{5.860055in}{1.118060in}}%
\pgfpathclose%
\pgfpathmoveto{\pgfqpoint{5.922179in}{1.133478in}}%
\pgfpathlineto{\pgfqpoint{5.922179in}{1.167593in}}%
\pgfpathlineto{\pgfqpoint{5.932536in}{1.167593in}}%
\pgfpathlineto{\pgfqpoint{5.932536in}{1.080000in}}%
\pgfpathlineto{\pgfqpoint{5.922179in}{1.080000in}}%
\pgfpathlineto{\pgfqpoint{5.922179in}{1.089456in}}%
\pgfpathquadraticcurveto{\pgfqpoint{5.918919in}{1.083837in}}{\pgfqpoint{5.913929in}{1.081099in}}%
\pgfpathquadraticcurveto{\pgfqpoint{5.908958in}{1.078361in}}{\pgfqpoint{5.901969in}{1.078361in}}%
\pgfpathquadraticcurveto{\pgfqpoint{5.890550in}{1.078361in}}{\pgfqpoint{5.883363in}{1.087475in}}%
\pgfpathquadraticcurveto{\pgfqpoint{5.876194in}{1.096607in}}{\pgfqpoint{5.876194in}{1.111467in}}%
\pgfpathquadraticcurveto{\pgfqpoint{5.876194in}{1.126327in}}{\pgfqpoint{5.883363in}{1.135441in}}%
\pgfpathquadraticcurveto{\pgfqpoint{5.890550in}{1.144556in}}{\pgfqpoint{5.901969in}{1.144556in}}%
\pgfpathquadraticcurveto{\pgfqpoint{5.908958in}{1.144556in}}{\pgfqpoint{5.913929in}{1.141818in}}%
\pgfpathquadraticcurveto{\pgfqpoint{5.918919in}{1.139098in}}{\pgfqpoint{5.922179in}{1.133478in}}%
\pgfpathlineto{\pgfqpoint{5.922179in}{1.133478in}}%
\pgfpathclose%
\pgfpathmoveto{\pgfqpoint{5.886893in}{1.111467in}}%
\pgfpathquadraticcurveto{\pgfqpoint{5.886893in}{1.100048in}}{\pgfqpoint{5.891594in}{1.093545in}}%
\pgfpathquadraticcurveto{\pgfqpoint{5.896296in}{1.087043in}}{\pgfqpoint{5.904509in}{1.087043in}}%
\pgfpathquadraticcurveto{\pgfqpoint{5.912723in}{1.087043in}}{\pgfqpoint{5.917442in}{1.093545in}}%
\pgfpathquadraticcurveto{\pgfqpoint{5.922179in}{1.100048in}}{\pgfqpoint{5.922179in}{1.111467in}}%
\pgfpathquadraticcurveto{\pgfqpoint{5.922179in}{1.122887in}}{\pgfqpoint{5.917442in}{1.129389in}}%
\pgfpathquadraticcurveto{\pgfqpoint{5.912723in}{1.135892in}}{\pgfqpoint{5.904509in}{1.135892in}}%
\pgfpathquadraticcurveto{\pgfqpoint{5.896296in}{1.135892in}}{\pgfqpoint{5.891594in}{1.129389in}}%
\pgfpathquadraticcurveto{\pgfqpoint{5.886893in}{1.122887in}}{\pgfqpoint{5.886893in}{1.111467in}}%
\pgfpathlineto{\pgfqpoint{5.886893in}{1.111467in}}%
\pgfpathclose%
\pgfpathmoveto{\pgfqpoint{5.994065in}{1.141187in}}%
\pgfpathlineto{\pgfqpoint{5.994065in}{1.131389in}}%
\pgfpathquadraticcurveto{\pgfqpoint{5.989689in}{1.133640in}}{\pgfqpoint{5.984951in}{1.134757in}}%
\pgfpathquadraticcurveto{\pgfqpoint{5.980232in}{1.135892in}}{\pgfqpoint{5.975153in}{1.135892in}}%
\pgfpathquadraticcurveto{\pgfqpoint{5.967444in}{1.135892in}}{\pgfqpoint{5.963589in}{1.133532in}}%
\pgfpathquadraticcurveto{\pgfqpoint{5.959734in}{1.131173in}}{\pgfqpoint{5.959734in}{1.126435in}}%
\pgfpathquadraticcurveto{\pgfqpoint{5.959734in}{1.122833in}}{\pgfqpoint{5.962490in}{1.120780in}}%
\pgfpathquadraticcurveto{\pgfqpoint{5.965246in}{1.118726in}}{\pgfqpoint{5.973586in}{1.116871in}}%
\pgfpathlineto{\pgfqpoint{5.977134in}{1.116078in}}%
\pgfpathquadraticcurveto{\pgfqpoint{5.988158in}{1.113719in}}{\pgfqpoint{5.992805in}{1.109414in}}%
\pgfpathquadraticcurveto{\pgfqpoint{5.997452in}{1.105109in}}{\pgfqpoint{5.997452in}{1.097400in}}%
\pgfpathquadraticcurveto{\pgfqpoint{5.997452in}{1.088610in}}{\pgfqpoint{5.990499in}{1.083476in}}%
\pgfpathquadraticcurveto{\pgfqpoint{5.983546in}{1.078361in}}{\pgfqpoint{5.971388in}{1.078361in}}%
\pgfpathquadraticcurveto{\pgfqpoint{5.966327in}{1.078361in}}{\pgfqpoint{5.960833in}{1.079352in}}%
\pgfpathquadraticcurveto{\pgfqpoint{5.955339in}{1.080342in}}{\pgfqpoint{5.949269in}{1.082306in}}%
\pgfpathlineto{\pgfqpoint{5.949269in}{1.093005in}}%
\pgfpathquadraticcurveto{\pgfqpoint{5.955015in}{1.090015in}}{\pgfqpoint{5.960581in}{1.088520in}}%
\pgfpathquadraticcurveto{\pgfqpoint{5.966147in}{1.087043in}}{\pgfqpoint{5.971622in}{1.087043in}}%
\pgfpathquadraticcurveto{\pgfqpoint{5.978935in}{1.087043in}}{\pgfqpoint{5.982862in}{1.089546in}}%
\pgfpathquadraticcurveto{\pgfqpoint{5.986807in}{1.092050in}}{\pgfqpoint{5.986807in}{1.096607in}}%
\pgfpathquadraticcurveto{\pgfqpoint{5.986807in}{1.100822in}}{\pgfqpoint{5.983961in}{1.103074in}}%
\pgfpathquadraticcurveto{\pgfqpoint{5.981133in}{1.105325in}}{\pgfqpoint{5.971496in}{1.107414in}}%
\pgfpathlineto{\pgfqpoint{5.967894in}{1.108261in}}%
\pgfpathquadraticcurveto{\pgfqpoint{5.958275in}{1.110278in}}{\pgfqpoint{5.953988in}{1.114475in}}%
\pgfpathquadraticcurveto{\pgfqpoint{5.949720in}{1.118672in}}{\pgfqpoint{5.949720in}{1.125985in}}%
\pgfpathquadraticcurveto{\pgfqpoint{5.949720in}{1.134883in}}{\pgfqpoint{5.956024in}{1.139710in}}%
\pgfpathquadraticcurveto{\pgfqpoint{5.962328in}{1.144556in}}{\pgfqpoint{5.973928in}{1.144556in}}%
\pgfpathquadraticcurveto{\pgfqpoint{5.979656in}{1.144556in}}{\pgfqpoint{5.984717in}{1.143709in}}%
\pgfpathquadraticcurveto{\pgfqpoint{5.989797in}{1.142880in}}{\pgfqpoint{5.994065in}{1.141187in}}%
\pgfpathlineto{\pgfqpoint{5.994065in}{1.141187in}}%
\pgfpathclose%
\pgfpathmoveto{\pgfqpoint{6.016581in}{1.094302in}}%
\pgfpathlineto{\pgfqpoint{6.028451in}{1.094302in}}%
\pgfpathlineto{\pgfqpoint{6.028451in}{1.080000in}}%
\pgfpathlineto{\pgfqpoint{6.016581in}{1.080000in}}%
\pgfpathlineto{\pgfqpoint{6.016581in}{1.094302in}}%
\pgfpathlineto{\pgfqpoint{6.016581in}{1.094302in}}%
\pgfpathclose%
\pgfpathmoveto{\pgfqpoint{6.016581in}{1.139602in}}%
\pgfpathlineto{\pgfqpoint{6.028451in}{1.139602in}}%
\pgfpathlineto{\pgfqpoint{6.028451in}{1.125319in}}%
\pgfpathlineto{\pgfqpoint{6.016581in}{1.125319in}}%
\pgfpathlineto{\pgfqpoint{6.016581in}{1.139602in}}%
\pgfpathlineto{\pgfqpoint{6.016581in}{1.139602in}}%
\pgfpathclose%
\pgfusepath{fill}%
\end{pgfscope}%
\begin{pgfscope}%
\pgfsetbuttcap%
\pgfsetroundjoin%
\definecolor{currentfill}{rgb}{0.309804,0.372549,0.419608}%
\pgfsetfillcolor{currentfill}%
\pgfsetlinewidth{0.000000pt}%
\definecolor{currentstroke}{rgb}{0.309804,0.372549,0.419608}%
\pgfsetstrokecolor{currentstroke}%
\pgfsetdash{}{0pt}%
\pgfpathmoveto{\pgfqpoint{6.160653in}{1.140629in}}%
\pgfpathlineto{\pgfqpoint{6.160653in}{1.130938in}}%
\pgfpathquadraticcurveto{\pgfqpoint{6.156258in}{1.133370in}}{\pgfqpoint{6.151845in}{1.134577in}}%
\pgfpathquadraticcurveto{\pgfqpoint{6.147432in}{1.135784in}}{\pgfqpoint{6.142929in}{1.135784in}}%
\pgfpathquadraticcurveto{\pgfqpoint{6.132842in}{1.135784in}}{\pgfqpoint{6.127258in}{1.129389in}}%
\pgfpathquadraticcurveto{\pgfqpoint{6.121692in}{1.123013in}}{\pgfqpoint{6.121692in}{1.111467in}}%
\pgfpathquadraticcurveto{\pgfqpoint{6.121692in}{1.099921in}}{\pgfqpoint{6.127258in}{1.093527in}}%
\pgfpathquadraticcurveto{\pgfqpoint{6.132842in}{1.087151in}}{\pgfqpoint{6.142929in}{1.087151in}}%
\pgfpathquadraticcurveto{\pgfqpoint{6.147432in}{1.087151in}}{\pgfqpoint{6.151845in}{1.088358in}}%
\pgfpathquadraticcurveto{\pgfqpoint{6.156258in}{1.089564in}}{\pgfqpoint{6.160653in}{1.091996in}}%
\pgfpathlineto{\pgfqpoint{6.160653in}{1.082414in}}%
\pgfpathquadraticcurveto{\pgfqpoint{6.156312in}{1.080396in}}{\pgfqpoint{6.151665in}{1.079388in}}%
\pgfpathquadraticcurveto{\pgfqpoint{6.147036in}{1.078361in}}{\pgfqpoint{6.141794in}{1.078361in}}%
\pgfpathquadraticcurveto{\pgfqpoint{6.127546in}{1.078361in}}{\pgfqpoint{6.119153in}{1.087313in}}%
\pgfpathquadraticcurveto{\pgfqpoint{6.110777in}{1.096265in}}{\pgfqpoint{6.110777in}{1.111467in}}%
\pgfpathquadraticcurveto{\pgfqpoint{6.110777in}{1.126886in}}{\pgfqpoint{6.119243in}{1.135712in}}%
\pgfpathquadraticcurveto{\pgfqpoint{6.127726in}{1.144556in}}{\pgfqpoint{6.142478in}{1.144556in}}%
\pgfpathquadraticcurveto{\pgfqpoint{6.147252in}{1.144556in}}{\pgfqpoint{6.151809in}{1.143565in}}%
\pgfpathquadraticcurveto{\pgfqpoint{6.156366in}{1.142592in}}{\pgfqpoint{6.160653in}{1.140629in}}%
\pgfpathlineto{\pgfqpoint{6.160653in}{1.140629in}}%
\pgfpathclose%
\pgfpathmoveto{\pgfqpoint{6.203089in}{1.135784in}}%
\pgfpathquadraticcurveto{\pgfqpoint{6.194768in}{1.135784in}}{\pgfqpoint{6.189922in}{1.129281in}}%
\pgfpathquadraticcurveto{\pgfqpoint{6.185077in}{1.122779in}}{\pgfqpoint{6.185077in}{1.111467in}}%
\pgfpathquadraticcurveto{\pgfqpoint{6.185077in}{1.100156in}}{\pgfqpoint{6.189886in}{1.093653in}}%
\pgfpathquadraticcurveto{\pgfqpoint{6.194714in}{1.087151in}}{\pgfqpoint{6.203089in}{1.087151in}}%
\pgfpathquadraticcurveto{\pgfqpoint{6.211375in}{1.087151in}}{\pgfqpoint{6.216202in}{1.093671in}}%
\pgfpathquadraticcurveto{\pgfqpoint{6.221047in}{1.100210in}}{\pgfqpoint{6.221047in}{1.111467in}}%
\pgfpathquadraticcurveto{\pgfqpoint{6.221047in}{1.122671in}}{\pgfqpoint{6.216202in}{1.129227in}}%
\pgfpathquadraticcurveto{\pgfqpoint{6.211375in}{1.135784in}}{\pgfqpoint{6.203089in}{1.135784in}}%
\pgfpathlineto{\pgfqpoint{6.203089in}{1.135784in}}%
\pgfpathclose%
\pgfpathmoveto{\pgfqpoint{6.203089in}{1.144556in}}%
\pgfpathquadraticcurveto{\pgfqpoint{6.216598in}{1.144556in}}{\pgfqpoint{6.224308in}{1.135766in}}%
\pgfpathquadraticcurveto{\pgfqpoint{6.232035in}{1.126994in}}{\pgfqpoint{6.232035in}{1.111467in}}%
\pgfpathquadraticcurveto{\pgfqpoint{6.232035in}{1.095995in}}{\pgfqpoint{6.224308in}{1.087169in}}%
\pgfpathquadraticcurveto{\pgfqpoint{6.216598in}{1.078361in}}{\pgfqpoint{6.203089in}{1.078361in}}%
\pgfpathquadraticcurveto{\pgfqpoint{6.189526in}{1.078361in}}{\pgfqpoint{6.181835in}{1.087169in}}%
\pgfpathquadraticcurveto{\pgfqpoint{6.174162in}{1.095995in}}{\pgfqpoint{6.174162in}{1.111467in}}%
\pgfpathquadraticcurveto{\pgfqpoint{6.174162in}{1.126994in}}{\pgfqpoint{6.181835in}{1.135766in}}%
\pgfpathquadraticcurveto{\pgfqpoint{6.189526in}{1.144556in}}{\pgfqpoint{6.203089in}{1.144556in}}%
\pgfpathlineto{\pgfqpoint{6.203089in}{1.144556in}}%
\pgfpathclose%
\pgfpathmoveto{\pgfqpoint{6.273625in}{1.135784in}}%
\pgfpathquadraticcurveto{\pgfqpoint{6.265303in}{1.135784in}}{\pgfqpoint{6.260458in}{1.129281in}}%
\pgfpathquadraticcurveto{\pgfqpoint{6.255613in}{1.122779in}}{\pgfqpoint{6.255613in}{1.111467in}}%
\pgfpathquadraticcurveto{\pgfqpoint{6.255613in}{1.100156in}}{\pgfqpoint{6.260422in}{1.093653in}}%
\pgfpathquadraticcurveto{\pgfqpoint{6.265249in}{1.087151in}}{\pgfqpoint{6.273625in}{1.087151in}}%
\pgfpathquadraticcurveto{\pgfqpoint{6.281911in}{1.087151in}}{\pgfqpoint{6.286738in}{1.093671in}}%
\pgfpathquadraticcurveto{\pgfqpoint{6.291583in}{1.100210in}}{\pgfqpoint{6.291583in}{1.111467in}}%
\pgfpathquadraticcurveto{\pgfqpoint{6.291583in}{1.122671in}}{\pgfqpoint{6.286738in}{1.129227in}}%
\pgfpathquadraticcurveto{\pgfqpoint{6.281911in}{1.135784in}}{\pgfqpoint{6.273625in}{1.135784in}}%
\pgfpathlineto{\pgfqpoint{6.273625in}{1.135784in}}%
\pgfpathclose%
\pgfpathmoveto{\pgfqpoint{6.273625in}{1.144556in}}%
\pgfpathquadraticcurveto{\pgfqpoint{6.287134in}{1.144556in}}{\pgfqpoint{6.294843in}{1.135766in}}%
\pgfpathquadraticcurveto{\pgfqpoint{6.302570in}{1.126994in}}{\pgfqpoint{6.302570in}{1.111467in}}%
\pgfpathquadraticcurveto{\pgfqpoint{6.302570in}{1.095995in}}{\pgfqpoint{6.294843in}{1.087169in}}%
\pgfpathquadraticcurveto{\pgfqpoint{6.287134in}{1.078361in}}{\pgfqpoint{6.273625in}{1.078361in}}%
\pgfpathquadraticcurveto{\pgfqpoint{6.260062in}{1.078361in}}{\pgfqpoint{6.252371in}{1.087169in}}%
\pgfpathquadraticcurveto{\pgfqpoint{6.244697in}{1.095995in}}{\pgfqpoint{6.244697in}{1.111467in}}%
\pgfpathquadraticcurveto{\pgfqpoint{6.244697in}{1.126994in}}{\pgfqpoint{6.252371in}{1.135766in}}%
\pgfpathquadraticcurveto{\pgfqpoint{6.260062in}{1.144556in}}{\pgfqpoint{6.273625in}{1.144556in}}%
\pgfpathlineto{\pgfqpoint{6.273625in}{1.144556in}}%
\pgfpathclose%
\pgfpathmoveto{\pgfqpoint{6.356265in}{1.133370in}}%
\pgfpathquadraticcurveto{\pgfqpoint{6.354517in}{1.134379in}}{\pgfqpoint{6.352464in}{1.134847in}}%
\pgfpathquadraticcurveto{\pgfqpoint{6.350411in}{1.135333in}}{\pgfqpoint{6.347943in}{1.135333in}}%
\pgfpathquadraticcurveto{\pgfqpoint{6.339153in}{1.135333in}}{\pgfqpoint{6.334452in}{1.129623in}}%
\pgfpathquadraticcurveto{\pgfqpoint{6.329751in}{1.123914in}}{\pgfqpoint{6.329751in}{1.113214in}}%
\pgfpathlineto{\pgfqpoint{6.329751in}{1.080000in}}%
\pgfpathlineto{\pgfqpoint{6.319340in}{1.080000in}}%
\pgfpathlineto{\pgfqpoint{6.319340in}{1.143043in}}%
\pgfpathlineto{\pgfqpoint{6.329751in}{1.143043in}}%
\pgfpathlineto{\pgfqpoint{6.329751in}{1.133244in}}%
\pgfpathquadraticcurveto{\pgfqpoint{6.333029in}{1.138990in}}{\pgfqpoint{6.338253in}{1.141764in}}%
\pgfpathquadraticcurveto{\pgfqpoint{6.343494in}{1.144556in}}{\pgfqpoint{6.350987in}{1.144556in}}%
\pgfpathquadraticcurveto{\pgfqpoint{6.352050in}{1.144556in}}{\pgfqpoint{6.353347in}{1.144411in}}%
\pgfpathquadraticcurveto{\pgfqpoint{6.354644in}{1.144285in}}{\pgfqpoint{6.356211in}{1.143997in}}%
\pgfpathlineto{\pgfqpoint{6.356265in}{1.133370in}}%
\pgfpathlineto{\pgfqpoint{6.356265in}{1.133370in}}%
\pgfpathclose%
\pgfpathmoveto{\pgfqpoint{6.406591in}{1.133478in}}%
\pgfpathlineto{\pgfqpoint{6.406591in}{1.167593in}}%
\pgfpathlineto{\pgfqpoint{6.416948in}{1.167593in}}%
\pgfpathlineto{\pgfqpoint{6.416948in}{1.080000in}}%
\pgfpathlineto{\pgfqpoint{6.406591in}{1.080000in}}%
\pgfpathlineto{\pgfqpoint{6.406591in}{1.089456in}}%
\pgfpathquadraticcurveto{\pgfqpoint{6.403330in}{1.083837in}}{\pgfqpoint{6.398341in}{1.081099in}}%
\pgfpathquadraticcurveto{\pgfqpoint{6.393370in}{1.078361in}}{\pgfqpoint{6.386381in}{1.078361in}}%
\pgfpathquadraticcurveto{\pgfqpoint{6.374961in}{1.078361in}}{\pgfqpoint{6.367774in}{1.087475in}}%
\pgfpathquadraticcurveto{\pgfqpoint{6.360606in}{1.096607in}}{\pgfqpoint{6.360606in}{1.111467in}}%
\pgfpathquadraticcurveto{\pgfqpoint{6.360606in}{1.126327in}}{\pgfqpoint{6.367774in}{1.135441in}}%
\pgfpathquadraticcurveto{\pgfqpoint{6.374961in}{1.144556in}}{\pgfqpoint{6.386381in}{1.144556in}}%
\pgfpathquadraticcurveto{\pgfqpoint{6.393370in}{1.144556in}}{\pgfqpoint{6.398341in}{1.141818in}}%
\pgfpathquadraticcurveto{\pgfqpoint{6.403330in}{1.139098in}}{\pgfqpoint{6.406591in}{1.133478in}}%
\pgfpathlineto{\pgfqpoint{6.406591in}{1.133478in}}%
\pgfpathclose%
\pgfpathmoveto{\pgfqpoint{6.371305in}{1.111467in}}%
\pgfpathquadraticcurveto{\pgfqpoint{6.371305in}{1.100048in}}{\pgfqpoint{6.376006in}{1.093545in}}%
\pgfpathquadraticcurveto{\pgfqpoint{6.380707in}{1.087043in}}{\pgfqpoint{6.388921in}{1.087043in}}%
\pgfpathquadraticcurveto{\pgfqpoint{6.397134in}{1.087043in}}{\pgfqpoint{6.401853in}{1.093545in}}%
\pgfpathquadraticcurveto{\pgfqpoint{6.406591in}{1.100048in}}{\pgfqpoint{6.406591in}{1.111467in}}%
\pgfpathquadraticcurveto{\pgfqpoint{6.406591in}{1.122887in}}{\pgfqpoint{6.401853in}{1.129389in}}%
\pgfpathquadraticcurveto{\pgfqpoint{6.397134in}{1.135892in}}{\pgfqpoint{6.388921in}{1.135892in}}%
\pgfpathquadraticcurveto{\pgfqpoint{6.380707in}{1.135892in}}{\pgfqpoint{6.376006in}{1.129389in}}%
\pgfpathquadraticcurveto{\pgfqpoint{6.371305in}{1.122887in}}{\pgfqpoint{6.371305in}{1.111467in}}%
\pgfpathlineto{\pgfqpoint{6.371305in}{1.111467in}}%
\pgfpathclose%
\pgfpathmoveto{\pgfqpoint{6.438292in}{1.143043in}}%
\pgfpathlineto{\pgfqpoint{6.448649in}{1.143043in}}%
\pgfpathlineto{\pgfqpoint{6.448649in}{1.080000in}}%
\pgfpathlineto{\pgfqpoint{6.438292in}{1.080000in}}%
\pgfpathlineto{\pgfqpoint{6.438292in}{1.143043in}}%
\pgfpathlineto{\pgfqpoint{6.438292in}{1.143043in}}%
\pgfpathclose%
\pgfpathmoveto{\pgfqpoint{6.438292in}{1.167593in}}%
\pgfpathlineto{\pgfqpoint{6.448649in}{1.167593in}}%
\pgfpathlineto{\pgfqpoint{6.448649in}{1.154462in}}%
\pgfpathlineto{\pgfqpoint{6.438292in}{1.154462in}}%
\pgfpathlineto{\pgfqpoint{6.438292in}{1.167593in}}%
\pgfpathlineto{\pgfqpoint{6.438292in}{1.167593in}}%
\pgfpathclose%
\pgfpathmoveto{\pgfqpoint{6.522733in}{1.118060in}}%
\pgfpathlineto{\pgfqpoint{6.522733in}{1.080000in}}%
\pgfpathlineto{\pgfqpoint{6.512376in}{1.080000in}}%
\pgfpathlineto{\pgfqpoint{6.512376in}{1.117717in}}%
\pgfpathquadraticcurveto{\pgfqpoint{6.512376in}{1.126669in}}{\pgfqpoint{6.508882in}{1.131100in}}%
\pgfpathquadraticcurveto{\pgfqpoint{6.505387in}{1.135549in}}{\pgfqpoint{6.498417in}{1.135549in}}%
\pgfpathquadraticcurveto{\pgfqpoint{6.490023in}{1.135549in}}{\pgfqpoint{6.485178in}{1.130200in}}%
\pgfpathquadraticcurveto{\pgfqpoint{6.480332in}{1.124868in}}{\pgfqpoint{6.480332in}{1.115628in}}%
\pgfpathlineto{\pgfqpoint{6.480332in}{1.080000in}}%
\pgfpathlineto{\pgfqpoint{6.469921in}{1.080000in}}%
\pgfpathlineto{\pgfqpoint{6.469921in}{1.143043in}}%
\pgfpathlineto{\pgfqpoint{6.480332in}{1.143043in}}%
\pgfpathlineto{\pgfqpoint{6.480332in}{1.133244in}}%
\pgfpathquadraticcurveto{\pgfqpoint{6.484061in}{1.138936in}}{\pgfqpoint{6.489086in}{1.141746in}}%
\pgfpathquadraticcurveto{\pgfqpoint{6.494130in}{1.144556in}}{\pgfqpoint{6.500722in}{1.144556in}}%
\pgfpathquadraticcurveto{\pgfqpoint{6.511583in}{1.144556in}}{\pgfqpoint{6.517149in}{1.137837in}}%
\pgfpathquadraticcurveto{\pgfqpoint{6.522733in}{1.131118in}}{\pgfqpoint{6.522733in}{1.118060in}}%
\pgfpathlineto{\pgfqpoint{6.522733in}{1.118060in}}%
\pgfpathclose%
\pgfpathmoveto{\pgfqpoint{6.572032in}{1.111683in}}%
\pgfpathquadraticcurveto{\pgfqpoint{6.559478in}{1.111683in}}{\pgfqpoint{6.554633in}{1.108819in}}%
\pgfpathquadraticcurveto{\pgfqpoint{6.549787in}{1.105956in}}{\pgfqpoint{6.549787in}{1.099021in}}%
\pgfpathquadraticcurveto{\pgfqpoint{6.549787in}{1.093509in}}{\pgfqpoint{6.553426in}{1.090267in}}%
\pgfpathquadraticcurveto{\pgfqpoint{6.557064in}{1.087043in}}{\pgfqpoint{6.563296in}{1.087043in}}%
\pgfpathquadraticcurveto{\pgfqpoint{6.571924in}{1.087043in}}{\pgfqpoint{6.577130in}{1.093149in}}%
\pgfpathquadraticcurveto{\pgfqpoint{6.582335in}{1.099255in}}{\pgfqpoint{6.582335in}{1.109378in}}%
\pgfpathlineto{\pgfqpoint{6.582335in}{1.111683in}}%
\pgfpathlineto{\pgfqpoint{6.572032in}{1.111683in}}%
\pgfpathlineto{\pgfqpoint{6.572032in}{1.111683in}}%
\pgfpathclose%
\pgfpathmoveto{\pgfqpoint{6.592692in}{1.115970in}}%
\pgfpathlineto{\pgfqpoint{6.592692in}{1.080000in}}%
\pgfpathlineto{\pgfqpoint{6.582335in}{1.080000in}}%
\pgfpathlineto{\pgfqpoint{6.582335in}{1.089564in}}%
\pgfpathquadraticcurveto{\pgfqpoint{6.578787in}{1.083837in}}{\pgfqpoint{6.573491in}{1.081099in}}%
\pgfpathquadraticcurveto{\pgfqpoint{6.568196in}{1.078361in}}{\pgfqpoint{6.560540in}{1.078361in}}%
\pgfpathquadraticcurveto{\pgfqpoint{6.550868in}{1.078361in}}{\pgfqpoint{6.545140in}{1.083801in}}%
\pgfpathquadraticcurveto{\pgfqpoint{6.539430in}{1.089240in}}{\pgfqpoint{6.539430in}{1.098354in}}%
\pgfpathquadraticcurveto{\pgfqpoint{6.539430in}{1.108982in}}{\pgfqpoint{6.546545in}{1.114385in}}%
\pgfpathquadraticcurveto{\pgfqpoint{6.553678in}{1.119789in}}{\pgfqpoint{6.567799in}{1.119789in}}%
\pgfpathlineto{\pgfqpoint{6.582335in}{1.119789in}}%
\pgfpathlineto{\pgfqpoint{6.582335in}{1.120816in}}%
\pgfpathquadraticcurveto{\pgfqpoint{6.582335in}{1.127966in}}{\pgfqpoint{6.577634in}{1.131875in}}%
\pgfpathquadraticcurveto{\pgfqpoint{6.572933in}{1.135784in}}{\pgfqpoint{6.564431in}{1.135784in}}%
\pgfpathquadraticcurveto{\pgfqpoint{6.559027in}{1.135784in}}{\pgfqpoint{6.553894in}{1.134487in}}%
\pgfpathquadraticcurveto{\pgfqpoint{6.548779in}{1.133190in}}{\pgfqpoint{6.544059in}{1.130596in}}%
\pgfpathlineto{\pgfqpoint{6.544059in}{1.140179in}}%
\pgfpathquadraticcurveto{\pgfqpoint{6.549733in}{1.142376in}}{\pgfqpoint{6.555083in}{1.143457in}}%
\pgfpathquadraticcurveto{\pgfqpoint{6.560432in}{1.144556in}}{\pgfqpoint{6.565494in}{1.144556in}}%
\pgfpathquadraticcurveto{\pgfqpoint{6.579183in}{1.144556in}}{\pgfqpoint{6.585938in}{1.137459in}}%
\pgfpathquadraticcurveto{\pgfqpoint{6.592692in}{1.130380in}}{\pgfqpoint{6.592692in}{1.115970in}}%
\pgfpathlineto{\pgfqpoint{6.592692in}{1.115970in}}%
\pgfpathclose%
\pgfpathmoveto{\pgfqpoint{6.624267in}{1.160947in}}%
\pgfpathlineto{\pgfqpoint{6.624267in}{1.143043in}}%
\pgfpathlineto{\pgfqpoint{6.645594in}{1.143043in}}%
\pgfpathlineto{\pgfqpoint{6.645594in}{1.134991in}}%
\pgfpathlineto{\pgfqpoint{6.624267in}{1.134991in}}%
\pgfpathlineto{\pgfqpoint{6.624267in}{1.100768in}}%
\pgfpathquadraticcurveto{\pgfqpoint{6.624267in}{1.093059in}}{\pgfqpoint{6.626375in}{1.090861in}}%
\pgfpathquadraticcurveto{\pgfqpoint{6.628482in}{1.088664in}}{\pgfqpoint{6.634967in}{1.088664in}}%
\pgfpathlineto{\pgfqpoint{6.645594in}{1.088664in}}%
\pgfpathlineto{\pgfqpoint{6.645594in}{1.080000in}}%
\pgfpathlineto{\pgfqpoint{6.634967in}{1.080000in}}%
\pgfpathquadraticcurveto{\pgfqpoint{6.622971in}{1.080000in}}{\pgfqpoint{6.618414in}{1.084467in}}%
\pgfpathquadraticcurveto{\pgfqpoint{6.613856in}{1.088952in}}{\pgfqpoint{6.613856in}{1.100768in}}%
\pgfpathlineto{\pgfqpoint{6.613856in}{1.134991in}}%
\pgfpathlineto{\pgfqpoint{6.606255in}{1.134991in}}%
\pgfpathlineto{\pgfqpoint{6.606255in}{1.143043in}}%
\pgfpathlineto{\pgfqpoint{6.613856in}{1.143043in}}%
\pgfpathlineto{\pgfqpoint{6.613856in}{1.160947in}}%
\pgfpathlineto{\pgfqpoint{6.624267in}{1.160947in}}%
\pgfpathlineto{\pgfqpoint{6.624267in}{1.160947in}}%
\pgfpathclose%
\pgfpathmoveto{\pgfqpoint{6.713139in}{1.114115in}}%
\pgfpathlineto{\pgfqpoint{6.713139in}{1.109054in}}%
\pgfpathlineto{\pgfqpoint{6.665515in}{1.109054in}}%
\pgfpathquadraticcurveto{\pgfqpoint{6.666200in}{1.098354in}}{\pgfqpoint{6.671964in}{1.092753in}}%
\pgfpathquadraticcurveto{\pgfqpoint{6.677728in}{1.087151in}}{\pgfqpoint{6.688031in}{1.087151in}}%
\pgfpathquadraticcurveto{\pgfqpoint{6.693993in}{1.087151in}}{\pgfqpoint{6.699594in}{1.088610in}}%
\pgfpathquadraticcurveto{\pgfqpoint{6.705196in}{1.090069in}}{\pgfqpoint{6.710726in}{1.093005in}}%
\pgfpathlineto{\pgfqpoint{6.710726in}{1.083206in}}%
\pgfpathquadraticcurveto{\pgfqpoint{6.705142in}{1.080847in}}{\pgfqpoint{6.699288in}{1.079604in}}%
\pgfpathquadraticcurveto{\pgfqpoint{6.693434in}{1.078361in}}{\pgfqpoint{6.687418in}{1.078361in}}%
\pgfpathquadraticcurveto{\pgfqpoint{6.672324in}{1.078361in}}{\pgfqpoint{6.663516in}{1.087133in}}%
\pgfpathquadraticcurveto{\pgfqpoint{6.654708in}{1.095923in}}{\pgfqpoint{6.654708in}{1.110909in}}%
\pgfpathquadraticcurveto{\pgfqpoint{6.654708in}{1.126381in}}{\pgfqpoint{6.663066in}{1.135459in}}%
\pgfpathquadraticcurveto{\pgfqpoint{6.671423in}{1.144556in}}{\pgfqpoint{6.685617in}{1.144556in}}%
\pgfpathquadraticcurveto{\pgfqpoint{6.698333in}{1.144556in}}{\pgfqpoint{6.705736in}{1.136360in}}%
\pgfpathquadraticcurveto{\pgfqpoint{6.713139in}{1.128183in}}{\pgfqpoint{6.713139in}{1.114115in}}%
\pgfpathlineto{\pgfqpoint{6.713139in}{1.114115in}}%
\pgfpathclose%
\pgfpathmoveto{\pgfqpoint{6.702782in}{1.117159in}}%
\pgfpathquadraticcurveto{\pgfqpoint{6.702674in}{1.125643in}}{\pgfqpoint{6.698027in}{1.130704in}}%
\pgfpathquadraticcurveto{\pgfqpoint{6.693380in}{1.135784in}}{\pgfqpoint{6.685725in}{1.135784in}}%
\pgfpathquadraticcurveto{\pgfqpoint{6.677061in}{1.135784in}}{\pgfqpoint{6.671856in}{1.130884in}}%
\pgfpathquadraticcurveto{\pgfqpoint{6.666650in}{1.125985in}}{\pgfqpoint{6.665858in}{1.117087in}}%
\pgfpathlineto{\pgfqpoint{6.702782in}{1.117159in}}%
\pgfpathlineto{\pgfqpoint{6.702782in}{1.117159in}}%
\pgfpathclose%
\pgfpathmoveto{\pgfqpoint{6.770328in}{1.141187in}}%
\pgfpathlineto{\pgfqpoint{6.770328in}{1.131389in}}%
\pgfpathquadraticcurveto{\pgfqpoint{6.765951in}{1.133640in}}{\pgfqpoint{6.761214in}{1.134757in}}%
\pgfpathquadraticcurveto{\pgfqpoint{6.756495in}{1.135892in}}{\pgfqpoint{6.751415in}{1.135892in}}%
\pgfpathquadraticcurveto{\pgfqpoint{6.743706in}{1.135892in}}{\pgfqpoint{6.739851in}{1.133532in}}%
\pgfpathquadraticcurveto{\pgfqpoint{6.735997in}{1.131173in}}{\pgfqpoint{6.735997in}{1.126435in}}%
\pgfpathquadraticcurveto{\pgfqpoint{6.735997in}{1.122833in}}{\pgfqpoint{6.738753in}{1.120780in}}%
\pgfpathquadraticcurveto{\pgfqpoint{6.741509in}{1.118726in}}{\pgfqpoint{6.749848in}{1.116871in}}%
\pgfpathlineto{\pgfqpoint{6.753397in}{1.116078in}}%
\pgfpathquadraticcurveto{\pgfqpoint{6.764420in}{1.113719in}}{\pgfqpoint{6.769067in}{1.109414in}}%
\pgfpathquadraticcurveto{\pgfqpoint{6.773714in}{1.105109in}}{\pgfqpoint{6.773714in}{1.097400in}}%
\pgfpathquadraticcurveto{\pgfqpoint{6.773714in}{1.088610in}}{\pgfqpoint{6.766762in}{1.083476in}}%
\pgfpathquadraticcurveto{\pgfqpoint{6.759809in}{1.078361in}}{\pgfqpoint{6.747651in}{1.078361in}}%
\pgfpathquadraticcurveto{\pgfqpoint{6.742589in}{1.078361in}}{\pgfqpoint{6.737096in}{1.079352in}}%
\pgfpathquadraticcurveto{\pgfqpoint{6.731602in}{1.080342in}}{\pgfqpoint{6.725532in}{1.082306in}}%
\pgfpathlineto{\pgfqpoint{6.725532in}{1.093005in}}%
\pgfpathquadraticcurveto{\pgfqpoint{6.731278in}{1.090015in}}{\pgfqpoint{6.736843in}{1.088520in}}%
\pgfpathquadraticcurveto{\pgfqpoint{6.742409in}{1.087043in}}{\pgfqpoint{6.747885in}{1.087043in}}%
\pgfpathquadraticcurveto{\pgfqpoint{6.755198in}{1.087043in}}{\pgfqpoint{6.759124in}{1.089546in}}%
\pgfpathquadraticcurveto{\pgfqpoint{6.763069in}{1.092050in}}{\pgfqpoint{6.763069in}{1.096607in}}%
\pgfpathquadraticcurveto{\pgfqpoint{6.763069in}{1.100822in}}{\pgfqpoint{6.760223in}{1.103074in}}%
\pgfpathquadraticcurveto{\pgfqpoint{6.757395in}{1.105325in}}{\pgfqpoint{6.747759in}{1.107414in}}%
\pgfpathlineto{\pgfqpoint{6.744156in}{1.108261in}}%
\pgfpathquadraticcurveto{\pgfqpoint{6.734538in}{1.110278in}}{\pgfqpoint{6.730251in}{1.114475in}}%
\pgfpathquadraticcurveto{\pgfqpoint{6.725982in}{1.118672in}}{\pgfqpoint{6.725982in}{1.125985in}}%
\pgfpathquadraticcurveto{\pgfqpoint{6.725982in}{1.134883in}}{\pgfqpoint{6.732286in}{1.139710in}}%
\pgfpathquadraticcurveto{\pgfqpoint{6.738591in}{1.144556in}}{\pgfqpoint{6.750190in}{1.144556in}}%
\pgfpathquadraticcurveto{\pgfqpoint{6.755918in}{1.144556in}}{\pgfqpoint{6.760980in}{1.143709in}}%
\pgfpathquadraticcurveto{\pgfqpoint{6.766059in}{1.142880in}}{\pgfqpoint{6.770328in}{1.141187in}}%
\pgfpathlineto{\pgfqpoint{6.770328in}{1.141187in}}%
\pgfpathclose%
\pgfusepath{fill}%
\end{pgfscope}%
\begin{pgfscope}%
\pgfsetbuttcap%
\pgfsetroundjoin%
\definecolor{currentfill}{rgb}{0.309804,0.372549,0.419608}%
\pgfsetfillcolor{currentfill}%
\pgfsetlinewidth{0.000000pt}%
\definecolor{currentstroke}{rgb}{0.309804,0.372549,0.419608}%
\pgfsetstrokecolor{currentstroke}%
\pgfsetdash{}{0pt}%
\pgfpathmoveto{\pgfqpoint{6.900518in}{1.140629in}}%
\pgfpathlineto{\pgfqpoint{6.900518in}{1.130938in}}%
\pgfpathquadraticcurveto{\pgfqpoint{6.896123in}{1.133370in}}{\pgfqpoint{6.891710in}{1.134577in}}%
\pgfpathquadraticcurveto{\pgfqpoint{6.887297in}{1.135784in}}{\pgfqpoint{6.882794in}{1.135784in}}%
\pgfpathquadraticcurveto{\pgfqpoint{6.872708in}{1.135784in}}{\pgfqpoint{6.867124in}{1.129389in}}%
\pgfpathquadraticcurveto{\pgfqpoint{6.861558in}{1.123013in}}{\pgfqpoint{6.861558in}{1.111467in}}%
\pgfpathquadraticcurveto{\pgfqpoint{6.861558in}{1.099921in}}{\pgfqpoint{6.867124in}{1.093527in}}%
\pgfpathquadraticcurveto{\pgfqpoint{6.872708in}{1.087151in}}{\pgfqpoint{6.882794in}{1.087151in}}%
\pgfpathquadraticcurveto{\pgfqpoint{6.887297in}{1.087151in}}{\pgfqpoint{6.891710in}{1.088358in}}%
\pgfpathquadraticcurveto{\pgfqpoint{6.896123in}{1.089564in}}{\pgfqpoint{6.900518in}{1.091996in}}%
\pgfpathlineto{\pgfqpoint{6.900518in}{1.082414in}}%
\pgfpathquadraticcurveto{\pgfqpoint{6.896177in}{1.080396in}}{\pgfqpoint{6.891530in}{1.079388in}}%
\pgfpathquadraticcurveto{\pgfqpoint{6.886901in}{1.078361in}}{\pgfqpoint{6.881660in}{1.078361in}}%
\pgfpathquadraticcurveto{\pgfqpoint{6.867412in}{1.078361in}}{\pgfqpoint{6.859018in}{1.087313in}}%
\pgfpathquadraticcurveto{\pgfqpoint{6.850643in}{1.096265in}}{\pgfqpoint{6.850643in}{1.111467in}}%
\pgfpathquadraticcurveto{\pgfqpoint{6.850643in}{1.126886in}}{\pgfqpoint{6.859108in}{1.135712in}}%
\pgfpathquadraticcurveto{\pgfqpoint{6.867592in}{1.144556in}}{\pgfqpoint{6.882344in}{1.144556in}}%
\pgfpathquadraticcurveto{\pgfqpoint{6.887117in}{1.144556in}}{\pgfqpoint{6.891674in}{1.143565in}}%
\pgfpathquadraticcurveto{\pgfqpoint{6.896231in}{1.142592in}}{\pgfqpoint{6.900518in}{1.140629in}}%
\pgfpathlineto{\pgfqpoint{6.900518in}{1.140629in}}%
\pgfpathclose%
\pgfpathmoveto{\pgfqpoint{6.947188in}{1.111683in}}%
\pgfpathquadraticcurveto{\pgfqpoint{6.934633in}{1.111683in}}{\pgfqpoint{6.929788in}{1.108819in}}%
\pgfpathquadraticcurveto{\pgfqpoint{6.924943in}{1.105956in}}{\pgfqpoint{6.924943in}{1.099021in}}%
\pgfpathquadraticcurveto{\pgfqpoint{6.924943in}{1.093509in}}{\pgfqpoint{6.928581in}{1.090267in}}%
\pgfpathquadraticcurveto{\pgfqpoint{6.932220in}{1.087043in}}{\pgfqpoint{6.938452in}{1.087043in}}%
\pgfpathquadraticcurveto{\pgfqpoint{6.947080in}{1.087043in}}{\pgfqpoint{6.952285in}{1.093149in}}%
\pgfpathquadraticcurveto{\pgfqpoint{6.957491in}{1.099255in}}{\pgfqpoint{6.957491in}{1.109378in}}%
\pgfpathlineto{\pgfqpoint{6.957491in}{1.111683in}}%
\pgfpathlineto{\pgfqpoint{6.947188in}{1.111683in}}%
\pgfpathlineto{\pgfqpoint{6.947188in}{1.111683in}}%
\pgfpathclose%
\pgfpathmoveto{\pgfqpoint{6.967848in}{1.115970in}}%
\pgfpathlineto{\pgfqpoint{6.967848in}{1.080000in}}%
\pgfpathlineto{\pgfqpoint{6.957491in}{1.080000in}}%
\pgfpathlineto{\pgfqpoint{6.957491in}{1.089564in}}%
\pgfpathquadraticcurveto{\pgfqpoint{6.953942in}{1.083837in}}{\pgfqpoint{6.948647in}{1.081099in}}%
\pgfpathquadraticcurveto{\pgfqpoint{6.943351in}{1.078361in}}{\pgfqpoint{6.935696in}{1.078361in}}%
\pgfpathquadraticcurveto{\pgfqpoint{6.926024in}{1.078361in}}{\pgfqpoint{6.920296in}{1.083801in}}%
\pgfpathquadraticcurveto{\pgfqpoint{6.914586in}{1.089240in}}{\pgfqpoint{6.914586in}{1.098354in}}%
\pgfpathquadraticcurveto{\pgfqpoint{6.914586in}{1.108982in}}{\pgfqpoint{6.921701in}{1.114385in}}%
\pgfpathquadraticcurveto{\pgfqpoint{6.928833in}{1.119789in}}{\pgfqpoint{6.942955in}{1.119789in}}%
\pgfpathlineto{\pgfqpoint{6.957491in}{1.119789in}}%
\pgfpathlineto{\pgfqpoint{6.957491in}{1.120816in}}%
\pgfpathquadraticcurveto{\pgfqpoint{6.957491in}{1.127966in}}{\pgfqpoint{6.952790in}{1.131875in}}%
\pgfpathquadraticcurveto{\pgfqpoint{6.948088in}{1.135784in}}{\pgfqpoint{6.939587in}{1.135784in}}%
\pgfpathquadraticcurveto{\pgfqpoint{6.934183in}{1.135784in}}{\pgfqpoint{6.929050in}{1.134487in}}%
\pgfpathquadraticcurveto{\pgfqpoint{6.923934in}{1.133190in}}{\pgfqpoint{6.919215in}{1.130596in}}%
\pgfpathlineto{\pgfqpoint{6.919215in}{1.140179in}}%
\pgfpathquadraticcurveto{\pgfqpoint{6.924889in}{1.142376in}}{\pgfqpoint{6.930238in}{1.143457in}}%
\pgfpathquadraticcurveto{\pgfqpoint{6.935588in}{1.144556in}}{\pgfqpoint{6.940649in}{1.144556in}}%
\pgfpathquadraticcurveto{\pgfqpoint{6.954339in}{1.144556in}}{\pgfqpoint{6.961093in}{1.137459in}}%
\pgfpathquadraticcurveto{\pgfqpoint{6.967848in}{1.130380in}}{\pgfqpoint{6.967848in}{1.115970in}}%
\pgfpathlineto{\pgfqpoint{6.967848in}{1.115970in}}%
\pgfpathclose%
\pgfpathmoveto{\pgfqpoint{6.999189in}{1.089456in}}%
\pgfpathlineto{\pgfqpoint{6.999189in}{1.056026in}}%
\pgfpathlineto{\pgfqpoint{6.988778in}{1.056026in}}%
\pgfpathlineto{\pgfqpoint{6.988778in}{1.143043in}}%
\pgfpathlineto{\pgfqpoint{6.999189in}{1.143043in}}%
\pgfpathlineto{\pgfqpoint{6.999189in}{1.133478in}}%
\pgfpathquadraticcurveto{\pgfqpoint{7.002467in}{1.139098in}}{\pgfqpoint{7.007438in}{1.141818in}}%
\pgfpathquadraticcurveto{\pgfqpoint{7.012428in}{1.144556in}}{\pgfqpoint{7.019344in}{1.144556in}}%
\pgfpathquadraticcurveto{\pgfqpoint{7.030836in}{1.144556in}}{\pgfqpoint{7.038005in}{1.135441in}}%
\pgfpathquadraticcurveto{\pgfqpoint{7.045192in}{1.126327in}}{\pgfqpoint{7.045192in}{1.111467in}}%
\pgfpathquadraticcurveto{\pgfqpoint{7.045192in}{1.096607in}}{\pgfqpoint{7.038005in}{1.087475in}}%
\pgfpathquadraticcurveto{\pgfqpoint{7.030836in}{1.078361in}}{\pgfqpoint{7.019344in}{1.078361in}}%
\pgfpathquadraticcurveto{\pgfqpoint{7.012428in}{1.078361in}}{\pgfqpoint{7.007438in}{1.081099in}}%
\pgfpathquadraticcurveto{\pgfqpoint{7.002467in}{1.083837in}}{\pgfqpoint{6.999189in}{1.089456in}}%
\pgfpathlineto{\pgfqpoint{6.999189in}{1.089456in}}%
\pgfpathclose%
\pgfpathmoveto{\pgfqpoint{7.034439in}{1.111467in}}%
\pgfpathquadraticcurveto{\pgfqpoint{7.034439in}{1.122887in}}{\pgfqpoint{7.029738in}{1.129389in}}%
\pgfpathquadraticcurveto{\pgfqpoint{7.025036in}{1.135892in}}{\pgfqpoint{7.016823in}{1.135892in}}%
\pgfpathquadraticcurveto{\pgfqpoint{7.008591in}{1.135892in}}{\pgfqpoint{7.003890in}{1.129389in}}%
\pgfpathquadraticcurveto{\pgfqpoint{6.999189in}{1.122887in}}{\pgfqpoint{6.999189in}{1.111467in}}%
\pgfpathquadraticcurveto{\pgfqpoint{6.999189in}{1.100048in}}{\pgfqpoint{7.003890in}{1.093545in}}%
\pgfpathquadraticcurveto{\pgfqpoint{7.008591in}{1.087043in}}{\pgfqpoint{7.016823in}{1.087043in}}%
\pgfpathquadraticcurveto{\pgfqpoint{7.025036in}{1.087043in}}{\pgfqpoint{7.029738in}{1.093545in}}%
\pgfpathquadraticcurveto{\pgfqpoint{7.034439in}{1.100048in}}{\pgfqpoint{7.034439in}{1.111467in}}%
\pgfpathlineto{\pgfqpoint{7.034439in}{1.111467in}}%
\pgfpathclose%
\pgfpathmoveto{\pgfqpoint{7.072372in}{1.089456in}}%
\pgfpathlineto{\pgfqpoint{7.072372in}{1.056026in}}%
\pgfpathlineto{\pgfqpoint{7.061961in}{1.056026in}}%
\pgfpathlineto{\pgfqpoint{7.061961in}{1.143043in}}%
\pgfpathlineto{\pgfqpoint{7.072372in}{1.143043in}}%
\pgfpathlineto{\pgfqpoint{7.072372in}{1.133478in}}%
\pgfpathquadraticcurveto{\pgfqpoint{7.075650in}{1.139098in}}{\pgfqpoint{7.080622in}{1.141818in}}%
\pgfpathquadraticcurveto{\pgfqpoint{7.085611in}{1.144556in}}{\pgfqpoint{7.092528in}{1.144556in}}%
\pgfpathquadraticcurveto{\pgfqpoint{7.104020in}{1.144556in}}{\pgfqpoint{7.111188in}{1.135441in}}%
\pgfpathquadraticcurveto{\pgfqpoint{7.118375in}{1.126327in}}{\pgfqpoint{7.118375in}{1.111467in}}%
\pgfpathquadraticcurveto{\pgfqpoint{7.118375in}{1.096607in}}{\pgfqpoint{7.111188in}{1.087475in}}%
\pgfpathquadraticcurveto{\pgfqpoint{7.104020in}{1.078361in}}{\pgfqpoint{7.092528in}{1.078361in}}%
\pgfpathquadraticcurveto{\pgfqpoint{7.085611in}{1.078361in}}{\pgfqpoint{7.080622in}{1.081099in}}%
\pgfpathquadraticcurveto{\pgfqpoint{7.075650in}{1.083837in}}{\pgfqpoint{7.072372in}{1.089456in}}%
\pgfpathlineto{\pgfqpoint{7.072372in}{1.089456in}}%
\pgfpathclose%
\pgfpathmoveto{\pgfqpoint{7.107622in}{1.111467in}}%
\pgfpathquadraticcurveto{\pgfqpoint{7.107622in}{1.122887in}}{\pgfqpoint{7.102921in}{1.129389in}}%
\pgfpathquadraticcurveto{\pgfqpoint{7.098220in}{1.135892in}}{\pgfqpoint{7.090006in}{1.135892in}}%
\pgfpathquadraticcurveto{\pgfqpoint{7.081775in}{1.135892in}}{\pgfqpoint{7.077073in}{1.129389in}}%
\pgfpathquadraticcurveto{\pgfqpoint{7.072372in}{1.122887in}}{\pgfqpoint{7.072372in}{1.111467in}}%
\pgfpathquadraticcurveto{\pgfqpoint{7.072372in}{1.100048in}}{\pgfqpoint{7.077073in}{1.093545in}}%
\pgfpathquadraticcurveto{\pgfqpoint{7.081775in}{1.087043in}}{\pgfqpoint{7.090006in}{1.087043in}}%
\pgfpathquadraticcurveto{\pgfqpoint{7.098220in}{1.087043in}}{\pgfqpoint{7.102921in}{1.093545in}}%
\pgfpathquadraticcurveto{\pgfqpoint{7.107622in}{1.100048in}}{\pgfqpoint{7.107622in}{1.111467in}}%
\pgfpathlineto{\pgfqpoint{7.107622in}{1.111467in}}%
\pgfpathclose%
\pgfpathmoveto{\pgfqpoint{7.189469in}{1.114115in}}%
\pgfpathlineto{\pgfqpoint{7.189469in}{1.109054in}}%
\pgfpathlineto{\pgfqpoint{7.141845in}{1.109054in}}%
\pgfpathquadraticcurveto{\pgfqpoint{7.142530in}{1.098354in}}{\pgfqpoint{7.148293in}{1.092753in}}%
\pgfpathquadraticcurveto{\pgfqpoint{7.154057in}{1.087151in}}{\pgfqpoint{7.164360in}{1.087151in}}%
\pgfpathquadraticcurveto{\pgfqpoint{7.170322in}{1.087151in}}{\pgfqpoint{7.175924in}{1.088610in}}%
\pgfpathquadraticcurveto{\pgfqpoint{7.181526in}{1.090069in}}{\pgfqpoint{7.187056in}{1.093005in}}%
\pgfpathlineto{\pgfqpoint{7.187056in}{1.083206in}}%
\pgfpathquadraticcurveto{\pgfqpoint{7.181472in}{1.080847in}}{\pgfqpoint{7.175618in}{1.079604in}}%
\pgfpathquadraticcurveto{\pgfqpoint{7.169764in}{1.078361in}}{\pgfqpoint{7.163748in}{1.078361in}}%
\pgfpathquadraticcurveto{\pgfqpoint{7.148654in}{1.078361in}}{\pgfqpoint{7.139846in}{1.087133in}}%
\pgfpathquadraticcurveto{\pgfqpoint{7.131038in}{1.095923in}}{\pgfqpoint{7.131038in}{1.110909in}}%
\pgfpathquadraticcurveto{\pgfqpoint{7.131038in}{1.126381in}}{\pgfqpoint{7.139395in}{1.135459in}}%
\pgfpathquadraticcurveto{\pgfqpoint{7.147753in}{1.144556in}}{\pgfqpoint{7.161947in}{1.144556in}}%
\pgfpathquadraticcurveto{\pgfqpoint{7.174663in}{1.144556in}}{\pgfqpoint{7.182066in}{1.136360in}}%
\pgfpathquadraticcurveto{\pgfqpoint{7.189469in}{1.128183in}}{\pgfqpoint{7.189469in}{1.114115in}}%
\pgfpathlineto{\pgfqpoint{7.189469in}{1.114115in}}%
\pgfpathclose%
\pgfpathmoveto{\pgfqpoint{7.179112in}{1.117159in}}%
\pgfpathquadraticcurveto{\pgfqpoint{7.179004in}{1.125643in}}{\pgfqpoint{7.174357in}{1.130704in}}%
\pgfpathquadraticcurveto{\pgfqpoint{7.169710in}{1.135784in}}{\pgfqpoint{7.162055in}{1.135784in}}%
\pgfpathquadraticcurveto{\pgfqpoint{7.153391in}{1.135784in}}{\pgfqpoint{7.148185in}{1.130884in}}%
\pgfpathquadraticcurveto{\pgfqpoint{7.142980in}{1.125985in}}{\pgfqpoint{7.142187in}{1.117087in}}%
\pgfpathlineto{\pgfqpoint{7.179112in}{1.117159in}}%
\pgfpathlineto{\pgfqpoint{7.179112in}{1.117159in}}%
\pgfpathclose%
\pgfpathmoveto{\pgfqpoint{7.247955in}{1.133478in}}%
\pgfpathlineto{\pgfqpoint{7.247955in}{1.167593in}}%
\pgfpathlineto{\pgfqpoint{7.258312in}{1.167593in}}%
\pgfpathlineto{\pgfqpoint{7.258312in}{1.080000in}}%
\pgfpathlineto{\pgfqpoint{7.247955in}{1.080000in}}%
\pgfpathlineto{\pgfqpoint{7.247955in}{1.089456in}}%
\pgfpathquadraticcurveto{\pgfqpoint{7.244695in}{1.083837in}}{\pgfqpoint{7.239705in}{1.081099in}}%
\pgfpathquadraticcurveto{\pgfqpoint{7.234734in}{1.078361in}}{\pgfqpoint{7.227745in}{1.078361in}}%
\pgfpathquadraticcurveto{\pgfqpoint{7.216325in}{1.078361in}}{\pgfqpoint{7.209139in}{1.087475in}}%
\pgfpathquadraticcurveto{\pgfqpoint{7.201970in}{1.096607in}}{\pgfqpoint{7.201970in}{1.111467in}}%
\pgfpathquadraticcurveto{\pgfqpoint{7.201970in}{1.126327in}}{\pgfqpoint{7.209139in}{1.135441in}}%
\pgfpathquadraticcurveto{\pgfqpoint{7.216325in}{1.144556in}}{\pgfqpoint{7.227745in}{1.144556in}}%
\pgfpathquadraticcurveto{\pgfqpoint{7.234734in}{1.144556in}}{\pgfqpoint{7.239705in}{1.141818in}}%
\pgfpathquadraticcurveto{\pgfqpoint{7.244695in}{1.139098in}}{\pgfqpoint{7.247955in}{1.133478in}}%
\pgfpathlineto{\pgfqpoint{7.247955in}{1.133478in}}%
\pgfpathclose%
\pgfpathmoveto{\pgfqpoint{7.212669in}{1.111467in}}%
\pgfpathquadraticcurveto{\pgfqpoint{7.212669in}{1.100048in}}{\pgfqpoint{7.217370in}{1.093545in}}%
\pgfpathquadraticcurveto{\pgfqpoint{7.222071in}{1.087043in}}{\pgfqpoint{7.230285in}{1.087043in}}%
\pgfpathquadraticcurveto{\pgfqpoint{7.238498in}{1.087043in}}{\pgfqpoint{7.243218in}{1.093545in}}%
\pgfpathquadraticcurveto{\pgfqpoint{7.247955in}{1.100048in}}{\pgfqpoint{7.247955in}{1.111467in}}%
\pgfpathquadraticcurveto{\pgfqpoint{7.247955in}{1.122887in}}{\pgfqpoint{7.243218in}{1.129389in}}%
\pgfpathquadraticcurveto{\pgfqpoint{7.238498in}{1.135892in}}{\pgfqpoint{7.230285in}{1.135892in}}%
\pgfpathquadraticcurveto{\pgfqpoint{7.222071in}{1.135892in}}{\pgfqpoint{7.217370in}{1.129389in}}%
\pgfpathquadraticcurveto{\pgfqpoint{7.212669in}{1.122887in}}{\pgfqpoint{7.212669in}{1.111467in}}%
\pgfpathlineto{\pgfqpoint{7.212669in}{1.111467in}}%
\pgfpathclose%
\pgfusepath{fill}%
\end{pgfscope}%
\begin{pgfscope}%
\pgfsetbuttcap%
\pgfsetroundjoin%
\definecolor{currentfill}{rgb}{0.309804,0.372549,0.419608}%
\pgfsetfillcolor{currentfill}%
\pgfsetlinewidth{0.000000pt}%
\definecolor{currentstroke}{rgb}{0.309804,0.372549,0.419608}%
\pgfsetstrokecolor{currentstroke}%
\pgfsetdash{}{0pt}%
\pgfpathmoveto{\pgfqpoint{7.375700in}{1.111683in}}%
\pgfpathquadraticcurveto{\pgfqpoint{7.363145in}{1.111683in}}{\pgfqpoint{7.358300in}{1.108819in}}%
\pgfpathquadraticcurveto{\pgfqpoint{7.353455in}{1.105956in}}{\pgfqpoint{7.353455in}{1.099021in}}%
\pgfpathquadraticcurveto{\pgfqpoint{7.353455in}{1.093509in}}{\pgfqpoint{7.357093in}{1.090267in}}%
\pgfpathquadraticcurveto{\pgfqpoint{7.360732in}{1.087043in}}{\pgfqpoint{7.366964in}{1.087043in}}%
\pgfpathquadraticcurveto{\pgfqpoint{7.375592in}{1.087043in}}{\pgfqpoint{7.380797in}{1.093149in}}%
\pgfpathquadraticcurveto{\pgfqpoint{7.386003in}{1.099255in}}{\pgfqpoint{7.386003in}{1.109378in}}%
\pgfpathlineto{\pgfqpoint{7.386003in}{1.111683in}}%
\pgfpathlineto{\pgfqpoint{7.375700in}{1.111683in}}%
\pgfpathlineto{\pgfqpoint{7.375700in}{1.111683in}}%
\pgfpathclose%
\pgfpathmoveto{\pgfqpoint{7.396360in}{1.115970in}}%
\pgfpathlineto{\pgfqpoint{7.396360in}{1.080000in}}%
\pgfpathlineto{\pgfqpoint{7.386003in}{1.080000in}}%
\pgfpathlineto{\pgfqpoint{7.386003in}{1.089564in}}%
\pgfpathquadraticcurveto{\pgfqpoint{7.382454in}{1.083837in}}{\pgfqpoint{7.377159in}{1.081099in}}%
\pgfpathquadraticcurveto{\pgfqpoint{7.371863in}{1.078361in}}{\pgfqpoint{7.364208in}{1.078361in}}%
\pgfpathquadraticcurveto{\pgfqpoint{7.354536in}{1.078361in}}{\pgfqpoint{7.348808in}{1.083801in}}%
\pgfpathquadraticcurveto{\pgfqpoint{7.343098in}{1.089240in}}{\pgfqpoint{7.343098in}{1.098354in}}%
\pgfpathquadraticcurveto{\pgfqpoint{7.343098in}{1.108982in}}{\pgfqpoint{7.350213in}{1.114385in}}%
\pgfpathquadraticcurveto{\pgfqpoint{7.357346in}{1.119789in}}{\pgfqpoint{7.371467in}{1.119789in}}%
\pgfpathlineto{\pgfqpoint{7.386003in}{1.119789in}}%
\pgfpathlineto{\pgfqpoint{7.386003in}{1.120816in}}%
\pgfpathquadraticcurveto{\pgfqpoint{7.386003in}{1.127966in}}{\pgfqpoint{7.381302in}{1.131875in}}%
\pgfpathquadraticcurveto{\pgfqpoint{7.376601in}{1.135784in}}{\pgfqpoint{7.368099in}{1.135784in}}%
\pgfpathquadraticcurveto{\pgfqpoint{7.362695in}{1.135784in}}{\pgfqpoint{7.357562in}{1.134487in}}%
\pgfpathquadraticcurveto{\pgfqpoint{7.352446in}{1.133190in}}{\pgfqpoint{7.347727in}{1.130596in}}%
\pgfpathlineto{\pgfqpoint{7.347727in}{1.140179in}}%
\pgfpathquadraticcurveto{\pgfqpoint{7.353401in}{1.142376in}}{\pgfqpoint{7.358750in}{1.143457in}}%
\pgfpathquadraticcurveto{\pgfqpoint{7.364100in}{1.144556in}}{\pgfqpoint{7.369162in}{1.144556in}}%
\pgfpathquadraticcurveto{\pgfqpoint{7.382851in}{1.144556in}}{\pgfqpoint{7.389605in}{1.137459in}}%
\pgfpathquadraticcurveto{\pgfqpoint{7.396360in}{1.130380in}}{\pgfqpoint{7.396360in}{1.115970in}}%
\pgfpathlineto{\pgfqpoint{7.396360in}{1.115970in}}%
\pgfpathclose%
\pgfpathmoveto{\pgfqpoint{7.427935in}{1.160947in}}%
\pgfpathlineto{\pgfqpoint{7.427935in}{1.143043in}}%
\pgfpathlineto{\pgfqpoint{7.449262in}{1.143043in}}%
\pgfpathlineto{\pgfqpoint{7.449262in}{1.134991in}}%
\pgfpathlineto{\pgfqpoint{7.427935in}{1.134991in}}%
\pgfpathlineto{\pgfqpoint{7.427935in}{1.100768in}}%
\pgfpathquadraticcurveto{\pgfqpoint{7.427935in}{1.093059in}}{\pgfqpoint{7.430043in}{1.090861in}}%
\pgfpathquadraticcurveto{\pgfqpoint{7.432150in}{1.088664in}}{\pgfqpoint{7.438634in}{1.088664in}}%
\pgfpathlineto{\pgfqpoint{7.449262in}{1.088664in}}%
\pgfpathlineto{\pgfqpoint{7.449262in}{1.080000in}}%
\pgfpathlineto{\pgfqpoint{7.438634in}{1.080000in}}%
\pgfpathquadraticcurveto{\pgfqpoint{7.426638in}{1.080000in}}{\pgfqpoint{7.422081in}{1.084467in}}%
\pgfpathquadraticcurveto{\pgfqpoint{7.417524in}{1.088952in}}{\pgfqpoint{7.417524in}{1.100768in}}%
\pgfpathlineto{\pgfqpoint{7.417524in}{1.134991in}}%
\pgfpathlineto{\pgfqpoint{7.409923in}{1.134991in}}%
\pgfpathlineto{\pgfqpoint{7.409923in}{1.143043in}}%
\pgfpathlineto{\pgfqpoint{7.417524in}{1.143043in}}%
\pgfpathlineto{\pgfqpoint{7.417524in}{1.160947in}}%
\pgfpathlineto{\pgfqpoint{7.427935in}{1.160947in}}%
\pgfpathlineto{\pgfqpoint{7.427935in}{1.160947in}}%
\pgfpathclose%
\pgfusepath{fill}%
\end{pgfscope}%
\begin{pgfscope}%
\pgfsetbuttcap%
\pgfsetroundjoin%
\definecolor{currentfill}{rgb}{0.309804,0.372549,0.419608}%
\pgfsetfillcolor{currentfill}%
\pgfsetlinewidth{0.000000pt}%
\definecolor{currentstroke}{rgb}{0.309804,0.372549,0.419608}%
\pgfsetstrokecolor{currentstroke}%
\pgfsetdash{}{0pt}%
\pgfpathmoveto{\pgfqpoint{7.551106in}{1.152283in}}%
\pgfpathlineto{\pgfqpoint{7.551106in}{1.130488in}}%
\pgfpathlineto{\pgfqpoint{7.582447in}{1.130488in}}%
\pgfpathlineto{\pgfqpoint{7.582447in}{1.120924in}}%
\pgfpathlineto{\pgfqpoint{7.551106in}{1.120924in}}%
\pgfpathlineto{\pgfqpoint{7.551106in}{1.099147in}}%
\pgfpathlineto{\pgfqpoint{7.541650in}{1.099147in}}%
\pgfpathlineto{\pgfqpoint{7.541650in}{1.120924in}}%
\pgfpathlineto{\pgfqpoint{7.510290in}{1.120924in}}%
\pgfpathlineto{\pgfqpoint{7.510290in}{1.130488in}}%
\pgfpathlineto{\pgfqpoint{7.541650in}{1.130488in}}%
\pgfpathlineto{\pgfqpoint{7.541650in}{1.152283in}}%
\pgfpathlineto{\pgfqpoint{7.551106in}{1.152283in}}%
\pgfpathlineto{\pgfqpoint{7.551106in}{1.152283in}}%
\pgfpathclose%
\pgfpathmoveto{\pgfqpoint{7.510290in}{1.089564in}}%
\pgfpathlineto{\pgfqpoint{7.582447in}{1.089564in}}%
\pgfpathlineto{\pgfqpoint{7.582447in}{1.080000in}}%
\pgfpathlineto{\pgfqpoint{7.510290in}{1.080000in}}%
\pgfpathlineto{\pgfqpoint{7.510290in}{1.089564in}}%
\pgfpathlineto{\pgfqpoint{7.510290in}{1.089564in}}%
\pgfpathclose%
\pgfpathmoveto{\pgfqpoint{7.607124in}{1.164045in}}%
\pgfpathlineto{\pgfqpoint{7.651758in}{1.164045in}}%
\pgfpathlineto{\pgfqpoint{7.651758in}{1.154462in}}%
\pgfpathlineto{\pgfqpoint{7.617535in}{1.154462in}}%
\pgfpathlineto{\pgfqpoint{7.617535in}{1.133874in}}%
\pgfpathquadraticcurveto{\pgfqpoint{7.620002in}{1.134721in}}{\pgfqpoint{7.622470in}{1.135135in}}%
\pgfpathquadraticcurveto{\pgfqpoint{7.624956in}{1.135549in}}{\pgfqpoint{7.627441in}{1.135549in}}%
\pgfpathquadraticcurveto{\pgfqpoint{7.641509in}{1.135549in}}{\pgfqpoint{7.649722in}{1.127840in}}%
\pgfpathquadraticcurveto{\pgfqpoint{7.657954in}{1.120131in}}{\pgfqpoint{7.657954in}{1.106964in}}%
\pgfpathquadraticcurveto{\pgfqpoint{7.657954in}{1.093401in}}{\pgfqpoint{7.649506in}{1.085872in}}%
\pgfpathquadraticcurveto{\pgfqpoint{7.641059in}{1.078361in}}{\pgfqpoint{7.625694in}{1.078361in}}%
\pgfpathquadraticcurveto{\pgfqpoint{7.620399in}{1.078361in}}{\pgfqpoint{7.614905in}{1.079262in}}%
\pgfpathquadraticcurveto{\pgfqpoint{7.609429in}{1.080162in}}{\pgfqpoint{7.603575in}{1.081963in}}%
\pgfpathlineto{\pgfqpoint{7.603575in}{1.093401in}}%
\pgfpathquadraticcurveto{\pgfqpoint{7.608637in}{1.090645in}}{\pgfqpoint{7.614040in}{1.089294in}}%
\pgfpathquadraticcurveto{\pgfqpoint{7.619444in}{1.087943in}}{\pgfqpoint{7.625460in}{1.087943in}}%
\pgfpathquadraticcurveto{\pgfqpoint{7.635205in}{1.087943in}}{\pgfqpoint{7.640878in}{1.093059in}}%
\pgfpathquadraticcurveto{\pgfqpoint{7.646570in}{1.098174in}}{\pgfqpoint{7.646570in}{1.106964in}}%
\pgfpathquadraticcurveto{\pgfqpoint{7.646570in}{1.115736in}}{\pgfqpoint{7.640878in}{1.120852in}}%
\pgfpathquadraticcurveto{\pgfqpoint{7.635205in}{1.125985in}}{\pgfqpoint{7.625460in}{1.125985in}}%
\pgfpathquadraticcurveto{\pgfqpoint{7.620903in}{1.125985in}}{\pgfqpoint{7.616364in}{1.124976in}}%
\pgfpathquadraticcurveto{\pgfqpoint{7.611843in}{1.123968in}}{\pgfqpoint{7.607124in}{1.121824in}}%
\pgfpathlineto{\pgfqpoint{7.607124in}{1.164045in}}%
\pgfpathlineto{\pgfqpoint{7.607124in}{1.164045in}}%
\pgfpathclose%
\pgfpathmoveto{\pgfqpoint{7.681532in}{1.139602in}}%
\pgfpathlineto{\pgfqpoint{7.693402in}{1.139602in}}%
\pgfpathlineto{\pgfqpoint{7.693402in}{1.125319in}}%
\pgfpathlineto{\pgfqpoint{7.681532in}{1.125319in}}%
\pgfpathlineto{\pgfqpoint{7.681532in}{1.139602in}}%
\pgfpathlineto{\pgfqpoint{7.681532in}{1.139602in}}%
\pgfpathclose%
\pgfpathmoveto{\pgfqpoint{7.681532in}{1.094302in}}%
\pgfpathlineto{\pgfqpoint{7.693402in}{1.094302in}}%
\pgfpathlineto{\pgfqpoint{7.693402in}{1.084611in}}%
\pgfpathlineto{\pgfqpoint{7.684180in}{1.066599in}}%
\pgfpathlineto{\pgfqpoint{7.676921in}{1.066599in}}%
\pgfpathlineto{\pgfqpoint{7.681532in}{1.084611in}}%
\pgfpathlineto{\pgfqpoint{7.681532in}{1.094302in}}%
\pgfpathlineto{\pgfqpoint{7.681532in}{1.094302in}}%
\pgfpathclose%
\pgfusepath{fill}%
\end{pgfscope}%
\begin{pgfscope}%
\pgfsetbuttcap%
\pgfsetroundjoin%
\definecolor{currentfill}{rgb}{0.309804,0.372549,0.419608}%
\pgfsetfillcolor{currentfill}%
\pgfsetlinewidth{0.000000pt}%
\definecolor{currentstroke}{rgb}{0.309804,0.372549,0.419608}%
\pgfsetstrokecolor{currentstroke}%
\pgfsetdash{}{0pt}%
\pgfpathmoveto{\pgfqpoint{7.812318in}{1.133478in}}%
\pgfpathlineto{\pgfqpoint{7.812318in}{1.167593in}}%
\pgfpathlineto{\pgfqpoint{7.822675in}{1.167593in}}%
\pgfpathlineto{\pgfqpoint{7.822675in}{1.080000in}}%
\pgfpathlineto{\pgfqpoint{7.812318in}{1.080000in}}%
\pgfpathlineto{\pgfqpoint{7.812318in}{1.089456in}}%
\pgfpathquadraticcurveto{\pgfqpoint{7.809058in}{1.083837in}}{\pgfqpoint{7.804069in}{1.081099in}}%
\pgfpathquadraticcurveto{\pgfqpoint{7.799097in}{1.078361in}}{\pgfqpoint{7.792109in}{1.078361in}}%
\pgfpathquadraticcurveto{\pgfqpoint{7.780689in}{1.078361in}}{\pgfqpoint{7.773502in}{1.087475in}}%
\pgfpathquadraticcurveto{\pgfqpoint{7.766333in}{1.096607in}}{\pgfqpoint{7.766333in}{1.111467in}}%
\pgfpathquadraticcurveto{\pgfqpoint{7.766333in}{1.126327in}}{\pgfqpoint{7.773502in}{1.135441in}}%
\pgfpathquadraticcurveto{\pgfqpoint{7.780689in}{1.144556in}}{\pgfqpoint{7.792109in}{1.144556in}}%
\pgfpathquadraticcurveto{\pgfqpoint{7.799097in}{1.144556in}}{\pgfqpoint{7.804069in}{1.141818in}}%
\pgfpathquadraticcurveto{\pgfqpoint{7.809058in}{1.139098in}}{\pgfqpoint{7.812318in}{1.133478in}}%
\pgfpathlineto{\pgfqpoint{7.812318in}{1.133478in}}%
\pgfpathclose%
\pgfpathmoveto{\pgfqpoint{7.777033in}{1.111467in}}%
\pgfpathquadraticcurveto{\pgfqpoint{7.777033in}{1.100048in}}{\pgfqpoint{7.781734in}{1.093545in}}%
\pgfpathquadraticcurveto{\pgfqpoint{7.786435in}{1.087043in}}{\pgfqpoint{7.794648in}{1.087043in}}%
\pgfpathquadraticcurveto{\pgfqpoint{7.802862in}{1.087043in}}{\pgfqpoint{7.807581in}{1.093545in}}%
\pgfpathquadraticcurveto{\pgfqpoint{7.812318in}{1.100048in}}{\pgfqpoint{7.812318in}{1.111467in}}%
\pgfpathquadraticcurveto{\pgfqpoint{7.812318in}{1.122887in}}{\pgfqpoint{7.807581in}{1.129389in}}%
\pgfpathquadraticcurveto{\pgfqpoint{7.802862in}{1.135892in}}{\pgfqpoint{7.794648in}{1.135892in}}%
\pgfpathquadraticcurveto{\pgfqpoint{7.786435in}{1.135892in}}{\pgfqpoint{7.781734in}{1.129389in}}%
\pgfpathquadraticcurveto{\pgfqpoint{7.777033in}{1.122887in}}{\pgfqpoint{7.777033in}{1.111467in}}%
\pgfpathlineto{\pgfqpoint{7.777033in}{1.111467in}}%
\pgfpathclose%
\pgfpathmoveto{\pgfqpoint{7.844020in}{1.143043in}}%
\pgfpathlineto{\pgfqpoint{7.854377in}{1.143043in}}%
\pgfpathlineto{\pgfqpoint{7.854377in}{1.080000in}}%
\pgfpathlineto{\pgfqpoint{7.844020in}{1.080000in}}%
\pgfpathlineto{\pgfqpoint{7.844020in}{1.143043in}}%
\pgfpathlineto{\pgfqpoint{7.844020in}{1.143043in}}%
\pgfpathclose%
\pgfpathmoveto{\pgfqpoint{7.844020in}{1.167593in}}%
\pgfpathlineto{\pgfqpoint{7.854377in}{1.167593in}}%
\pgfpathlineto{\pgfqpoint{7.854377in}{1.154462in}}%
\pgfpathlineto{\pgfqpoint{7.844020in}{1.154462in}}%
\pgfpathlineto{\pgfqpoint{7.844020in}{1.167593in}}%
\pgfpathlineto{\pgfqpoint{7.844020in}{1.167593in}}%
\pgfpathclose%
\pgfpathmoveto{\pgfqpoint{7.916230in}{1.141187in}}%
\pgfpathlineto{\pgfqpoint{7.916230in}{1.131389in}}%
\pgfpathquadraticcurveto{\pgfqpoint{7.911853in}{1.133640in}}{\pgfqpoint{7.907116in}{1.134757in}}%
\pgfpathquadraticcurveto{\pgfqpoint{7.902397in}{1.135892in}}{\pgfqpoint{7.897318in}{1.135892in}}%
\pgfpathquadraticcurveto{\pgfqpoint{7.889608in}{1.135892in}}{\pgfqpoint{7.885754in}{1.133532in}}%
\pgfpathquadraticcurveto{\pgfqpoint{7.881899in}{1.131173in}}{\pgfqpoint{7.881899in}{1.126435in}}%
\pgfpathquadraticcurveto{\pgfqpoint{7.881899in}{1.122833in}}{\pgfqpoint{7.884655in}{1.120780in}}%
\pgfpathquadraticcurveto{\pgfqpoint{7.887411in}{1.118726in}}{\pgfqpoint{7.895751in}{1.116871in}}%
\pgfpathlineto{\pgfqpoint{7.899299in}{1.116078in}}%
\pgfpathquadraticcurveto{\pgfqpoint{7.910322in}{1.113719in}}{\pgfqpoint{7.914970in}{1.109414in}}%
\pgfpathquadraticcurveto{\pgfqpoint{7.919617in}{1.105109in}}{\pgfqpoint{7.919617in}{1.097400in}}%
\pgfpathquadraticcurveto{\pgfqpoint{7.919617in}{1.088610in}}{\pgfqpoint{7.912664in}{1.083476in}}%
\pgfpathquadraticcurveto{\pgfqpoint{7.905711in}{1.078361in}}{\pgfqpoint{7.893553in}{1.078361in}}%
\pgfpathquadraticcurveto{\pgfqpoint{7.888492in}{1.078361in}}{\pgfqpoint{7.882998in}{1.079352in}}%
\pgfpathquadraticcurveto{\pgfqpoint{7.877504in}{1.080342in}}{\pgfqpoint{7.871434in}{1.082306in}}%
\pgfpathlineto{\pgfqpoint{7.871434in}{1.093005in}}%
\pgfpathquadraticcurveto{\pgfqpoint{7.877180in}{1.090015in}}{\pgfqpoint{7.882746in}{1.088520in}}%
\pgfpathquadraticcurveto{\pgfqpoint{7.888312in}{1.087043in}}{\pgfqpoint{7.893787in}{1.087043in}}%
\pgfpathquadraticcurveto{\pgfqpoint{7.901100in}{1.087043in}}{\pgfqpoint{7.905027in}{1.089546in}}%
\pgfpathquadraticcurveto{\pgfqpoint{7.908972in}{1.092050in}}{\pgfqpoint{7.908972in}{1.096607in}}%
\pgfpathquadraticcurveto{\pgfqpoint{7.908972in}{1.100822in}}{\pgfqpoint{7.906126in}{1.103074in}}%
\pgfpathquadraticcurveto{\pgfqpoint{7.903298in}{1.105325in}}{\pgfqpoint{7.893661in}{1.107414in}}%
\pgfpathlineto{\pgfqpoint{7.890059in}{1.108261in}}%
\pgfpathquadraticcurveto{\pgfqpoint{7.880440in}{1.110278in}}{\pgfqpoint{7.876153in}{1.114475in}}%
\pgfpathquadraticcurveto{\pgfqpoint{7.871885in}{1.118672in}}{\pgfqpoint{7.871885in}{1.125985in}}%
\pgfpathquadraticcurveto{\pgfqpoint{7.871885in}{1.134883in}}{\pgfqpoint{7.878189in}{1.139710in}}%
\pgfpathquadraticcurveto{\pgfqpoint{7.884493in}{1.144556in}}{\pgfqpoint{7.896093in}{1.144556in}}%
\pgfpathquadraticcurveto{\pgfqpoint{7.901821in}{1.144556in}}{\pgfqpoint{7.906882in}{1.143709in}}%
\pgfpathquadraticcurveto{\pgfqpoint{7.911962in}{1.142880in}}{\pgfqpoint{7.916230in}{1.141187in}}%
\pgfpathlineto{\pgfqpoint{7.916230in}{1.141187in}}%
\pgfpathclose%
\pgfpathmoveto{\pgfqpoint{7.946113in}{1.089456in}}%
\pgfpathlineto{\pgfqpoint{7.946113in}{1.056026in}}%
\pgfpathlineto{\pgfqpoint{7.935702in}{1.056026in}}%
\pgfpathlineto{\pgfqpoint{7.935702in}{1.143043in}}%
\pgfpathlineto{\pgfqpoint{7.946113in}{1.143043in}}%
\pgfpathlineto{\pgfqpoint{7.946113in}{1.133478in}}%
\pgfpathquadraticcurveto{\pgfqpoint{7.949391in}{1.139098in}}{\pgfqpoint{7.954362in}{1.141818in}}%
\pgfpathquadraticcurveto{\pgfqpoint{7.959352in}{1.144556in}}{\pgfqpoint{7.966268in}{1.144556in}}%
\pgfpathquadraticcurveto{\pgfqpoint{7.977760in}{1.144556in}}{\pgfqpoint{7.984929in}{1.135441in}}%
\pgfpathquadraticcurveto{\pgfqpoint{7.992116in}{1.126327in}}{\pgfqpoint{7.992116in}{1.111467in}}%
\pgfpathquadraticcurveto{\pgfqpoint{7.992116in}{1.096607in}}{\pgfqpoint{7.984929in}{1.087475in}}%
\pgfpathquadraticcurveto{\pgfqpoint{7.977760in}{1.078361in}}{\pgfqpoint{7.966268in}{1.078361in}}%
\pgfpathquadraticcurveto{\pgfqpoint{7.959352in}{1.078361in}}{\pgfqpoint{7.954362in}{1.081099in}}%
\pgfpathquadraticcurveto{\pgfqpoint{7.949391in}{1.083837in}}{\pgfqpoint{7.946113in}{1.089456in}}%
\pgfpathlineto{\pgfqpoint{7.946113in}{1.089456in}}%
\pgfpathclose%
\pgfpathmoveto{\pgfqpoint{7.981362in}{1.111467in}}%
\pgfpathquadraticcurveto{\pgfqpoint{7.981362in}{1.122887in}}{\pgfqpoint{7.976661in}{1.129389in}}%
\pgfpathquadraticcurveto{\pgfqpoint{7.971960in}{1.135892in}}{\pgfqpoint{7.963746in}{1.135892in}}%
\pgfpathquadraticcurveto{\pgfqpoint{7.955515in}{1.135892in}}{\pgfqpoint{7.950814in}{1.129389in}}%
\pgfpathquadraticcurveto{\pgfqpoint{7.946113in}{1.122887in}}{\pgfqpoint{7.946113in}{1.111467in}}%
\pgfpathquadraticcurveto{\pgfqpoint{7.946113in}{1.100048in}}{\pgfqpoint{7.950814in}{1.093545in}}%
\pgfpathquadraticcurveto{\pgfqpoint{7.955515in}{1.087043in}}{\pgfqpoint{7.963746in}{1.087043in}}%
\pgfpathquadraticcurveto{\pgfqpoint{7.971960in}{1.087043in}}{\pgfqpoint{7.976661in}{1.093545in}}%
\pgfpathquadraticcurveto{\pgfqpoint{7.981362in}{1.100048in}}{\pgfqpoint{7.981362in}{1.111467in}}%
\pgfpathlineto{\pgfqpoint{7.981362in}{1.111467in}}%
\pgfpathclose%
\pgfpathmoveto{\pgfqpoint{8.009281in}{1.167593in}}%
\pgfpathlineto{\pgfqpoint{8.019638in}{1.167593in}}%
\pgfpathlineto{\pgfqpoint{8.019638in}{1.080000in}}%
\pgfpathlineto{\pgfqpoint{8.009281in}{1.080000in}}%
\pgfpathlineto{\pgfqpoint{8.009281in}{1.167593in}}%
\pgfpathlineto{\pgfqpoint{8.009281in}{1.167593in}}%
\pgfpathclose%
\pgfpathmoveto{\pgfqpoint{8.069964in}{1.111683in}}%
\pgfpathquadraticcurveto{\pgfqpoint{8.057410in}{1.111683in}}{\pgfqpoint{8.052564in}{1.108819in}}%
\pgfpathquadraticcurveto{\pgfqpoint{8.047719in}{1.105956in}}{\pgfqpoint{8.047719in}{1.099021in}}%
\pgfpathquadraticcurveto{\pgfqpoint{8.047719in}{1.093509in}}{\pgfqpoint{8.051358in}{1.090267in}}%
\pgfpathquadraticcurveto{\pgfqpoint{8.054996in}{1.087043in}}{\pgfqpoint{8.061228in}{1.087043in}}%
\pgfpathquadraticcurveto{\pgfqpoint{8.069856in}{1.087043in}}{\pgfqpoint{8.075062in}{1.093149in}}%
\pgfpathquadraticcurveto{\pgfqpoint{8.080267in}{1.099255in}}{\pgfqpoint{8.080267in}{1.109378in}}%
\pgfpathlineto{\pgfqpoint{8.080267in}{1.111683in}}%
\pgfpathlineto{\pgfqpoint{8.069964in}{1.111683in}}%
\pgfpathlineto{\pgfqpoint{8.069964in}{1.111683in}}%
\pgfpathclose%
\pgfpathmoveto{\pgfqpoint{8.090624in}{1.115970in}}%
\pgfpathlineto{\pgfqpoint{8.090624in}{1.080000in}}%
\pgfpathlineto{\pgfqpoint{8.080267in}{1.080000in}}%
\pgfpathlineto{\pgfqpoint{8.080267in}{1.089564in}}%
\pgfpathquadraticcurveto{\pgfqpoint{8.076719in}{1.083837in}}{\pgfqpoint{8.071423in}{1.081099in}}%
\pgfpathquadraticcurveto{\pgfqpoint{8.066128in}{1.078361in}}{\pgfqpoint{8.058472in}{1.078361in}}%
\pgfpathquadraticcurveto{\pgfqpoint{8.048800in}{1.078361in}}{\pgfqpoint{8.043072in}{1.083801in}}%
\pgfpathquadraticcurveto{\pgfqpoint{8.037362in}{1.089240in}}{\pgfqpoint{8.037362in}{1.098354in}}%
\pgfpathquadraticcurveto{\pgfqpoint{8.037362in}{1.108982in}}{\pgfqpoint{8.044477in}{1.114385in}}%
\pgfpathquadraticcurveto{\pgfqpoint{8.051610in}{1.119789in}}{\pgfqpoint{8.065731in}{1.119789in}}%
\pgfpathlineto{\pgfqpoint{8.080267in}{1.119789in}}%
\pgfpathlineto{\pgfqpoint{8.080267in}{1.120816in}}%
\pgfpathquadraticcurveto{\pgfqpoint{8.080267in}{1.127966in}}{\pgfqpoint{8.075566in}{1.131875in}}%
\pgfpathquadraticcurveto{\pgfqpoint{8.070865in}{1.135784in}}{\pgfqpoint{8.062363in}{1.135784in}}%
\pgfpathquadraticcurveto{\pgfqpoint{8.056959in}{1.135784in}}{\pgfqpoint{8.051826in}{1.134487in}}%
\pgfpathquadraticcurveto{\pgfqpoint{8.046710in}{1.133190in}}{\pgfqpoint{8.041991in}{1.130596in}}%
\pgfpathlineto{\pgfqpoint{8.041991in}{1.140179in}}%
\pgfpathquadraticcurveto{\pgfqpoint{8.047665in}{1.142376in}}{\pgfqpoint{8.053015in}{1.143457in}}%
\pgfpathquadraticcurveto{\pgfqpoint{8.058364in}{1.144556in}}{\pgfqpoint{8.063426in}{1.144556in}}%
\pgfpathquadraticcurveto{\pgfqpoint{8.077115in}{1.144556in}}{\pgfqpoint{8.083870in}{1.137459in}}%
\pgfpathquadraticcurveto{\pgfqpoint{8.090624in}{1.130380in}}{\pgfqpoint{8.090624in}{1.115970in}}%
\pgfpathlineto{\pgfqpoint{8.090624in}{1.115970in}}%
\pgfpathclose%
\pgfpathmoveto{\pgfqpoint{8.138176in}{1.074146in}}%
\pgfpathquadraticcurveto{\pgfqpoint{8.133799in}{1.062888in}}{\pgfqpoint{8.129620in}{1.059466in}}%
\pgfpathquadraticcurveto{\pgfqpoint{8.125460in}{1.056026in}}{\pgfqpoint{8.118489in}{1.056026in}}%
\pgfpathlineto{\pgfqpoint{8.110203in}{1.056026in}}%
\pgfpathlineto{\pgfqpoint{8.110203in}{1.064690in}}%
\pgfpathlineto{\pgfqpoint{8.116291in}{1.064690in}}%
\pgfpathquadraticcurveto{\pgfqpoint{8.120560in}{1.064690in}}{\pgfqpoint{8.122920in}{1.066725in}}%
\pgfpathquadraticcurveto{\pgfqpoint{8.125297in}{1.068742in}}{\pgfqpoint{8.128161in}{1.076289in}}%
\pgfpathlineto{\pgfqpoint{8.130017in}{1.081009in}}%
\pgfpathlineto{\pgfqpoint{8.104529in}{1.143043in}}%
\pgfpathlineto{\pgfqpoint{8.115499in}{1.143043in}}%
\pgfpathlineto{\pgfqpoint{8.135204in}{1.093743in}}%
\pgfpathlineto{\pgfqpoint{8.154909in}{1.143043in}}%
\pgfpathlineto{\pgfqpoint{8.165879in}{1.143043in}}%
\pgfpathlineto{\pgfqpoint{8.138176in}{1.074146in}}%
\pgfpathlineto{\pgfqpoint{8.138176in}{1.074146in}}%
\pgfpathclose%
\pgfpathmoveto{\pgfqpoint{8.234109in}{1.114115in}}%
\pgfpathlineto{\pgfqpoint{8.234109in}{1.109054in}}%
\pgfpathlineto{\pgfqpoint{8.186485in}{1.109054in}}%
\pgfpathquadraticcurveto{\pgfqpoint{8.187169in}{1.098354in}}{\pgfqpoint{8.192933in}{1.092753in}}%
\pgfpathquadraticcurveto{\pgfqpoint{8.198697in}{1.087151in}}{\pgfqpoint{8.209000in}{1.087151in}}%
\pgfpathquadraticcurveto{\pgfqpoint{8.214962in}{1.087151in}}{\pgfqpoint{8.220564in}{1.088610in}}%
\pgfpathquadraticcurveto{\pgfqpoint{8.226166in}{1.090069in}}{\pgfqpoint{8.231695in}{1.093005in}}%
\pgfpathlineto{\pgfqpoint{8.231695in}{1.083206in}}%
\pgfpathquadraticcurveto{\pgfqpoint{8.226112in}{1.080847in}}{\pgfqpoint{8.220258in}{1.079604in}}%
\pgfpathquadraticcurveto{\pgfqpoint{8.214404in}{1.078361in}}{\pgfqpoint{8.208388in}{1.078361in}}%
\pgfpathquadraticcurveto{\pgfqpoint{8.193293in}{1.078361in}}{\pgfqpoint{8.184485in}{1.087133in}}%
\pgfpathquadraticcurveto{\pgfqpoint{8.175677in}{1.095923in}}{\pgfqpoint{8.175677in}{1.110909in}}%
\pgfpathquadraticcurveto{\pgfqpoint{8.175677in}{1.126381in}}{\pgfqpoint{8.184035in}{1.135459in}}%
\pgfpathquadraticcurveto{\pgfqpoint{8.192393in}{1.144556in}}{\pgfqpoint{8.206586in}{1.144556in}}%
\pgfpathquadraticcurveto{\pgfqpoint{8.219303in}{1.144556in}}{\pgfqpoint{8.226706in}{1.136360in}}%
\pgfpathquadraticcurveto{\pgfqpoint{8.234109in}{1.128183in}}{\pgfqpoint{8.234109in}{1.114115in}}%
\pgfpathlineto{\pgfqpoint{8.234109in}{1.114115in}}%
\pgfpathclose%
\pgfpathmoveto{\pgfqpoint{8.223752in}{1.117159in}}%
\pgfpathquadraticcurveto{\pgfqpoint{8.223644in}{1.125643in}}{\pgfqpoint{8.218997in}{1.130704in}}%
\pgfpathquadraticcurveto{\pgfqpoint{8.214350in}{1.135784in}}{\pgfqpoint{8.206694in}{1.135784in}}%
\pgfpathquadraticcurveto{\pgfqpoint{8.198031in}{1.135784in}}{\pgfqpoint{8.192825in}{1.130884in}}%
\pgfpathquadraticcurveto{\pgfqpoint{8.187620in}{1.125985in}}{\pgfqpoint{8.186827in}{1.117087in}}%
\pgfpathlineto{\pgfqpoint{8.223752in}{1.117159in}}%
\pgfpathlineto{\pgfqpoint{8.223752in}{1.117159in}}%
\pgfpathclose%
\pgfpathmoveto{\pgfqpoint{8.292594in}{1.133478in}}%
\pgfpathlineto{\pgfqpoint{8.292594in}{1.167593in}}%
\pgfpathlineto{\pgfqpoint{8.302951in}{1.167593in}}%
\pgfpathlineto{\pgfqpoint{8.302951in}{1.080000in}}%
\pgfpathlineto{\pgfqpoint{8.292594in}{1.080000in}}%
\pgfpathlineto{\pgfqpoint{8.292594in}{1.089456in}}%
\pgfpathquadraticcurveto{\pgfqpoint{8.289334in}{1.083837in}}{\pgfqpoint{8.284345in}{1.081099in}}%
\pgfpathquadraticcurveto{\pgfqpoint{8.279373in}{1.078361in}}{\pgfqpoint{8.272385in}{1.078361in}}%
\pgfpathquadraticcurveto{\pgfqpoint{8.260965in}{1.078361in}}{\pgfqpoint{8.253778in}{1.087475in}}%
\pgfpathquadraticcurveto{\pgfqpoint{8.246609in}{1.096607in}}{\pgfqpoint{8.246609in}{1.111467in}}%
\pgfpathquadraticcurveto{\pgfqpoint{8.246609in}{1.126327in}}{\pgfqpoint{8.253778in}{1.135441in}}%
\pgfpathquadraticcurveto{\pgfqpoint{8.260965in}{1.144556in}}{\pgfqpoint{8.272385in}{1.144556in}}%
\pgfpathquadraticcurveto{\pgfqpoint{8.279373in}{1.144556in}}{\pgfqpoint{8.284345in}{1.141818in}}%
\pgfpathquadraticcurveto{\pgfqpoint{8.289334in}{1.139098in}}{\pgfqpoint{8.292594in}{1.133478in}}%
\pgfpathlineto{\pgfqpoint{8.292594in}{1.133478in}}%
\pgfpathclose%
\pgfpathmoveto{\pgfqpoint{8.257309in}{1.111467in}}%
\pgfpathquadraticcurveto{\pgfqpoint{8.257309in}{1.100048in}}{\pgfqpoint{8.262010in}{1.093545in}}%
\pgfpathquadraticcurveto{\pgfqpoint{8.266711in}{1.087043in}}{\pgfqpoint{8.274924in}{1.087043in}}%
\pgfpathquadraticcurveto{\pgfqpoint{8.283138in}{1.087043in}}{\pgfqpoint{8.287857in}{1.093545in}}%
\pgfpathquadraticcurveto{\pgfqpoint{8.292594in}{1.100048in}}{\pgfqpoint{8.292594in}{1.111467in}}%
\pgfpathquadraticcurveto{\pgfqpoint{8.292594in}{1.122887in}}{\pgfqpoint{8.287857in}{1.129389in}}%
\pgfpathquadraticcurveto{\pgfqpoint{8.283138in}{1.135892in}}{\pgfqpoint{8.274924in}{1.135892in}}%
\pgfpathquadraticcurveto{\pgfqpoint{8.266711in}{1.135892in}}{\pgfqpoint{8.262010in}{1.129389in}}%
\pgfpathquadraticcurveto{\pgfqpoint{8.257309in}{1.122887in}}{\pgfqpoint{8.257309in}{1.111467in}}%
\pgfpathlineto{\pgfqpoint{8.257309in}{1.111467in}}%
\pgfpathclose%
\pgfusepath{fill}%
\end{pgfscope}%
\begin{pgfscope}%
\pgfsetbuttcap%
\pgfsetroundjoin%
\definecolor{currentfill}{rgb}{0.309804,0.372549,0.419608}%
\pgfsetfillcolor{currentfill}%
\pgfsetlinewidth{0.000000pt}%
\definecolor{currentstroke}{rgb}{0.309804,0.372549,0.419608}%
\pgfsetstrokecolor{currentstroke}%
\pgfsetdash{}{0pt}%
\pgfpathmoveto{\pgfqpoint{8.424215in}{1.112260in}}%
\pgfpathquadraticcurveto{\pgfqpoint{8.424215in}{1.123517in}}{\pgfqpoint{8.419568in}{1.129696in}}%
\pgfpathquadraticcurveto{\pgfqpoint{8.414939in}{1.135892in}}{\pgfqpoint{8.406545in}{1.135892in}}%
\pgfpathquadraticcurveto{\pgfqpoint{8.398224in}{1.135892in}}{\pgfqpoint{8.393577in}{1.129696in}}%
\pgfpathquadraticcurveto{\pgfqpoint{8.388929in}{1.123517in}}{\pgfqpoint{8.388929in}{1.112260in}}%
\pgfpathquadraticcurveto{\pgfqpoint{8.388929in}{1.101056in}}{\pgfqpoint{8.393577in}{1.094860in}}%
\pgfpathquadraticcurveto{\pgfqpoint{8.398224in}{1.088664in}}{\pgfqpoint{8.406545in}{1.088664in}}%
\pgfpathquadraticcurveto{\pgfqpoint{8.414939in}{1.088664in}}{\pgfqpoint{8.419568in}{1.094860in}}%
\pgfpathquadraticcurveto{\pgfqpoint{8.424215in}{1.101056in}}{\pgfqpoint{8.424215in}{1.112260in}}%
\pgfpathlineto{\pgfqpoint{8.424215in}{1.112260in}}%
\pgfpathclose%
\pgfpathmoveto{\pgfqpoint{8.434572in}{1.087817in}}%
\pgfpathquadraticcurveto{\pgfqpoint{8.434572in}{1.071732in}}{\pgfqpoint{8.427421in}{1.063879in}}%
\pgfpathquadraticcurveto{\pgfqpoint{8.420289in}{1.056026in}}{\pgfqpoint{8.405537in}{1.056026in}}%
\pgfpathquadraticcurveto{\pgfqpoint{8.400079in}{1.056026in}}{\pgfqpoint{8.395234in}{1.056836in}}%
\pgfpathquadraticcurveto{\pgfqpoint{8.390388in}{1.057647in}}{\pgfqpoint{8.385831in}{1.059340in}}%
\pgfpathlineto{\pgfqpoint{8.385831in}{1.069409in}}%
\pgfpathquadraticcurveto{\pgfqpoint{8.390388in}{1.066941in}}{\pgfqpoint{8.394837in}{1.065770in}}%
\pgfpathquadraticcurveto{\pgfqpoint{8.399286in}{1.064582in}}{\pgfqpoint{8.403897in}{1.064582in}}%
\pgfpathquadraticcurveto{\pgfqpoint{8.414092in}{1.064582in}}{\pgfqpoint{8.419154in}{1.069895in}}%
\pgfpathquadraticcurveto{\pgfqpoint{8.424215in}{1.075209in}}{\pgfqpoint{8.424215in}{1.085962in}}%
\pgfpathlineto{\pgfqpoint{8.424215in}{1.091095in}}%
\pgfpathquadraticcurveto{\pgfqpoint{8.421009in}{1.085512in}}{\pgfqpoint{8.416002in}{1.082756in}}%
\pgfpathquadraticcurveto{\pgfqpoint{8.410994in}{1.080000in}}{\pgfqpoint{8.404006in}{1.080000in}}%
\pgfpathquadraticcurveto{\pgfqpoint{8.392424in}{1.080000in}}{\pgfqpoint{8.385327in}{1.088826in}}%
\pgfpathquadraticcurveto{\pgfqpoint{8.378230in}{1.097670in}}{\pgfqpoint{8.378230in}{1.112260in}}%
\pgfpathquadraticcurveto{\pgfqpoint{8.378230in}{1.126886in}}{\pgfqpoint{8.385327in}{1.135712in}}%
\pgfpathquadraticcurveto{\pgfqpoint{8.392424in}{1.144556in}}{\pgfqpoint{8.404006in}{1.144556in}}%
\pgfpathquadraticcurveto{\pgfqpoint{8.410994in}{1.144556in}}{\pgfqpoint{8.416002in}{1.141800in}}%
\pgfpathquadraticcurveto{\pgfqpoint{8.421009in}{1.139044in}}{\pgfqpoint{8.424215in}{1.133478in}}%
\pgfpathlineto{\pgfqpoint{8.424215in}{1.143043in}}%
\pgfpathlineto{\pgfqpoint{8.434572in}{1.143043in}}%
\pgfpathlineto{\pgfqpoint{8.434572in}{1.087817in}}%
\pgfpathlineto{\pgfqpoint{8.434572in}{1.087817in}}%
\pgfpathclose%
\pgfpathmoveto{\pgfqpoint{8.484574in}{1.111683in}}%
\pgfpathquadraticcurveto{\pgfqpoint{8.472019in}{1.111683in}}{\pgfqpoint{8.467174in}{1.108819in}}%
\pgfpathquadraticcurveto{\pgfqpoint{8.462329in}{1.105956in}}{\pgfqpoint{8.462329in}{1.099021in}}%
\pgfpathquadraticcurveto{\pgfqpoint{8.462329in}{1.093509in}}{\pgfqpoint{8.465967in}{1.090267in}}%
\pgfpathquadraticcurveto{\pgfqpoint{8.469606in}{1.087043in}}{\pgfqpoint{8.475838in}{1.087043in}}%
\pgfpathquadraticcurveto{\pgfqpoint{8.484466in}{1.087043in}}{\pgfqpoint{8.489671in}{1.093149in}}%
\pgfpathquadraticcurveto{\pgfqpoint{8.494877in}{1.099255in}}{\pgfqpoint{8.494877in}{1.109378in}}%
\pgfpathlineto{\pgfqpoint{8.494877in}{1.111683in}}%
\pgfpathlineto{\pgfqpoint{8.484574in}{1.111683in}}%
\pgfpathlineto{\pgfqpoint{8.484574in}{1.111683in}}%
\pgfpathclose%
\pgfpathmoveto{\pgfqpoint{8.505234in}{1.115970in}}%
\pgfpathlineto{\pgfqpoint{8.505234in}{1.080000in}}%
\pgfpathlineto{\pgfqpoint{8.494877in}{1.080000in}}%
\pgfpathlineto{\pgfqpoint{8.494877in}{1.089564in}}%
\pgfpathquadraticcurveto{\pgfqpoint{8.491328in}{1.083837in}}{\pgfqpoint{8.486033in}{1.081099in}}%
\pgfpathquadraticcurveto{\pgfqpoint{8.480737in}{1.078361in}}{\pgfqpoint{8.473082in}{1.078361in}}%
\pgfpathquadraticcurveto{\pgfqpoint{8.463410in}{1.078361in}}{\pgfqpoint{8.457682in}{1.083801in}}%
\pgfpathquadraticcurveto{\pgfqpoint{8.451972in}{1.089240in}}{\pgfqpoint{8.451972in}{1.098354in}}%
\pgfpathquadraticcurveto{\pgfqpoint{8.451972in}{1.108982in}}{\pgfqpoint{8.459087in}{1.114385in}}%
\pgfpathquadraticcurveto{\pgfqpoint{8.466220in}{1.119789in}}{\pgfqpoint{8.480341in}{1.119789in}}%
\pgfpathlineto{\pgfqpoint{8.494877in}{1.119789in}}%
\pgfpathlineto{\pgfqpoint{8.494877in}{1.120816in}}%
\pgfpathquadraticcurveto{\pgfqpoint{8.494877in}{1.127966in}}{\pgfqpoint{8.490176in}{1.131875in}}%
\pgfpathquadraticcurveto{\pgfqpoint{8.485475in}{1.135784in}}{\pgfqpoint{8.476973in}{1.135784in}}%
\pgfpathquadraticcurveto{\pgfqpoint{8.471569in}{1.135784in}}{\pgfqpoint{8.466436in}{1.134487in}}%
\pgfpathquadraticcurveto{\pgfqpoint{8.461320in}{1.133190in}}{\pgfqpoint{8.456601in}{1.130596in}}%
\pgfpathlineto{\pgfqpoint{8.456601in}{1.140179in}}%
\pgfpathquadraticcurveto{\pgfqpoint{8.462275in}{1.142376in}}{\pgfqpoint{8.467624in}{1.143457in}}%
\pgfpathquadraticcurveto{\pgfqpoint{8.472974in}{1.144556in}}{\pgfqpoint{8.478036in}{1.144556in}}%
\pgfpathquadraticcurveto{\pgfqpoint{8.491725in}{1.144556in}}{\pgfqpoint{8.498479in}{1.137459in}}%
\pgfpathquadraticcurveto{\pgfqpoint{8.505234in}{1.130380in}}{\pgfqpoint{8.505234in}{1.115970in}}%
\pgfpathlineto{\pgfqpoint{8.505234in}{1.115970in}}%
\pgfpathclose%
\pgfpathmoveto{\pgfqpoint{8.536575in}{1.089456in}}%
\pgfpathlineto{\pgfqpoint{8.536575in}{1.056026in}}%
\pgfpathlineto{\pgfqpoint{8.526164in}{1.056026in}}%
\pgfpathlineto{\pgfqpoint{8.526164in}{1.143043in}}%
\pgfpathlineto{\pgfqpoint{8.536575in}{1.143043in}}%
\pgfpathlineto{\pgfqpoint{8.536575in}{1.133478in}}%
\pgfpathquadraticcurveto{\pgfqpoint{8.539853in}{1.139098in}}{\pgfqpoint{8.544825in}{1.141818in}}%
\pgfpathquadraticcurveto{\pgfqpoint{8.549814in}{1.144556in}}{\pgfqpoint{8.556731in}{1.144556in}}%
\pgfpathquadraticcurveto{\pgfqpoint{8.568222in}{1.144556in}}{\pgfqpoint{8.575391in}{1.135441in}}%
\pgfpathquadraticcurveto{\pgfqpoint{8.582578in}{1.126327in}}{\pgfqpoint{8.582578in}{1.111467in}}%
\pgfpathquadraticcurveto{\pgfqpoint{8.582578in}{1.096607in}}{\pgfqpoint{8.575391in}{1.087475in}}%
\pgfpathquadraticcurveto{\pgfqpoint{8.568222in}{1.078361in}}{\pgfqpoint{8.556731in}{1.078361in}}%
\pgfpathquadraticcurveto{\pgfqpoint{8.549814in}{1.078361in}}{\pgfqpoint{8.544825in}{1.081099in}}%
\pgfpathquadraticcurveto{\pgfqpoint{8.539853in}{1.083837in}}{\pgfqpoint{8.536575in}{1.089456in}}%
\pgfpathlineto{\pgfqpoint{8.536575in}{1.089456in}}%
\pgfpathclose%
\pgfpathmoveto{\pgfqpoint{8.571825in}{1.111467in}}%
\pgfpathquadraticcurveto{\pgfqpoint{8.571825in}{1.122887in}}{\pgfqpoint{8.567124in}{1.129389in}}%
\pgfpathquadraticcurveto{\pgfqpoint{8.562422in}{1.135892in}}{\pgfqpoint{8.554209in}{1.135892in}}%
\pgfpathquadraticcurveto{\pgfqpoint{8.545977in}{1.135892in}}{\pgfqpoint{8.541276in}{1.129389in}}%
\pgfpathquadraticcurveto{\pgfqpoint{8.536575in}{1.122887in}}{\pgfqpoint{8.536575in}{1.111467in}}%
\pgfpathquadraticcurveto{\pgfqpoint{8.536575in}{1.100048in}}{\pgfqpoint{8.541276in}{1.093545in}}%
\pgfpathquadraticcurveto{\pgfqpoint{8.545977in}{1.087043in}}{\pgfqpoint{8.554209in}{1.087043in}}%
\pgfpathquadraticcurveto{\pgfqpoint{8.562422in}{1.087043in}}{\pgfqpoint{8.567124in}{1.093545in}}%
\pgfpathquadraticcurveto{\pgfqpoint{8.571825in}{1.100048in}}{\pgfqpoint{8.571825in}{1.111467in}}%
\pgfpathlineto{\pgfqpoint{8.571825in}{1.111467in}}%
\pgfpathclose%
\pgfpathmoveto{\pgfqpoint{8.639929in}{1.141187in}}%
\pgfpathlineto{\pgfqpoint{8.639929in}{1.131389in}}%
\pgfpathquadraticcurveto{\pgfqpoint{8.635552in}{1.133640in}}{\pgfqpoint{8.630815in}{1.134757in}}%
\pgfpathquadraticcurveto{\pgfqpoint{8.626095in}{1.135892in}}{\pgfqpoint{8.621016in}{1.135892in}}%
\pgfpathquadraticcurveto{\pgfqpoint{8.613307in}{1.135892in}}{\pgfqpoint{8.609452in}{1.133532in}}%
\pgfpathquadraticcurveto{\pgfqpoint{8.605598in}{1.131173in}}{\pgfqpoint{8.605598in}{1.126435in}}%
\pgfpathquadraticcurveto{\pgfqpoint{8.605598in}{1.122833in}}{\pgfqpoint{8.608353in}{1.120780in}}%
\pgfpathquadraticcurveto{\pgfqpoint{8.611109in}{1.118726in}}{\pgfqpoint{8.619449in}{1.116871in}}%
\pgfpathlineto{\pgfqpoint{8.622997in}{1.116078in}}%
\pgfpathquadraticcurveto{\pgfqpoint{8.634021in}{1.113719in}}{\pgfqpoint{8.638668in}{1.109414in}}%
\pgfpathquadraticcurveto{\pgfqpoint{8.643315in}{1.105109in}}{\pgfqpoint{8.643315in}{1.097400in}}%
\pgfpathquadraticcurveto{\pgfqpoint{8.643315in}{1.088610in}}{\pgfqpoint{8.636362in}{1.083476in}}%
\pgfpathquadraticcurveto{\pgfqpoint{8.629410in}{1.078361in}}{\pgfqpoint{8.617251in}{1.078361in}}%
\pgfpathquadraticcurveto{\pgfqpoint{8.612190in}{1.078361in}}{\pgfqpoint{8.606696in}{1.079352in}}%
\pgfpathquadraticcurveto{\pgfqpoint{8.601203in}{1.080342in}}{\pgfqpoint{8.595133in}{1.082306in}}%
\pgfpathlineto{\pgfqpoint{8.595133in}{1.093005in}}%
\pgfpathquadraticcurveto{\pgfqpoint{8.600878in}{1.090015in}}{\pgfqpoint{8.606444in}{1.088520in}}%
\pgfpathquadraticcurveto{\pgfqpoint{8.612010in}{1.087043in}}{\pgfqpoint{8.617486in}{1.087043in}}%
\pgfpathquadraticcurveto{\pgfqpoint{8.624799in}{1.087043in}}{\pgfqpoint{8.628725in}{1.089546in}}%
\pgfpathquadraticcurveto{\pgfqpoint{8.632670in}{1.092050in}}{\pgfqpoint{8.632670in}{1.096607in}}%
\pgfpathquadraticcurveto{\pgfqpoint{8.632670in}{1.100822in}}{\pgfqpoint{8.629824in}{1.103074in}}%
\pgfpathquadraticcurveto{\pgfqpoint{8.626996in}{1.105325in}}{\pgfqpoint{8.617360in}{1.107414in}}%
\pgfpathlineto{\pgfqpoint{8.613757in}{1.108261in}}%
\pgfpathquadraticcurveto{\pgfqpoint{8.604139in}{1.110278in}}{\pgfqpoint{8.599852in}{1.114475in}}%
\pgfpathquadraticcurveto{\pgfqpoint{8.595583in}{1.118672in}}{\pgfqpoint{8.595583in}{1.125985in}}%
\pgfpathquadraticcurveto{\pgfqpoint{8.595583in}{1.134883in}}{\pgfqpoint{8.601887in}{1.139710in}}%
\pgfpathquadraticcurveto{\pgfqpoint{8.608191in}{1.144556in}}{\pgfqpoint{8.619791in}{1.144556in}}%
\pgfpathquadraticcurveto{\pgfqpoint{8.625519in}{1.144556in}}{\pgfqpoint{8.630580in}{1.143709in}}%
\pgfpathquadraticcurveto{\pgfqpoint{8.635660in}{1.142880in}}{\pgfqpoint{8.639929in}{1.141187in}}%
\pgfpathlineto{\pgfqpoint{8.639929in}{1.141187in}}%
\pgfpathclose%
\pgfusepath{fill}%
\end{pgfscope}%
\begin{pgfscope}%
\pgfsetbuttcap%
\pgfsetroundjoin%
\definecolor{currentfill}{rgb}{0.309804,0.372549,0.419608}%
\pgfsetfillcolor{currentfill}%
\pgfsetlinewidth{0.000000pt}%
\definecolor{currentstroke}{rgb}{0.309804,0.372549,0.419608}%
\pgfsetstrokecolor{currentstroke}%
\pgfsetdash{}{0pt}%
\pgfpathmoveto{\pgfqpoint{8.763713in}{1.130938in}}%
\pgfpathquadraticcurveto{\pgfqpoint{8.767604in}{1.137927in}}{\pgfqpoint{8.773007in}{1.141241in}}%
\pgfpathquadraticcurveto{\pgfqpoint{8.778411in}{1.144556in}}{\pgfqpoint{8.785724in}{1.144556in}}%
\pgfpathquadraticcurveto{\pgfqpoint{8.795577in}{1.144556in}}{\pgfqpoint{8.800926in}{1.137657in}}%
\pgfpathquadraticcurveto{\pgfqpoint{8.806276in}{1.130776in}}{\pgfqpoint{8.806276in}{1.118060in}}%
\pgfpathlineto{\pgfqpoint{8.806276in}{1.080000in}}%
\pgfpathlineto{\pgfqpoint{8.795865in}{1.080000in}}%
\pgfpathlineto{\pgfqpoint{8.795865in}{1.117717in}}%
\pgfpathquadraticcurveto{\pgfqpoint{8.795865in}{1.126778in}}{\pgfqpoint{8.792641in}{1.131155in}}%
\pgfpathquadraticcurveto{\pgfqpoint{8.789435in}{1.135549in}}{\pgfqpoint{8.782860in}{1.135549in}}%
\pgfpathquadraticcurveto{\pgfqpoint{8.774809in}{1.135549in}}{\pgfqpoint{8.770126in}{1.130200in}}%
\pgfpathquadraticcurveto{\pgfqpoint{8.765460in}{1.124868in}}{\pgfqpoint{8.765460in}{1.115628in}}%
\pgfpathlineto{\pgfqpoint{8.765460in}{1.080000in}}%
\pgfpathlineto{\pgfqpoint{8.755049in}{1.080000in}}%
\pgfpathlineto{\pgfqpoint{8.755049in}{1.117717in}}%
\pgfpathquadraticcurveto{\pgfqpoint{8.755049in}{1.126832in}}{\pgfqpoint{8.751843in}{1.131191in}}%
\pgfpathquadraticcurveto{\pgfqpoint{8.748637in}{1.135549in}}{\pgfqpoint{8.741937in}{1.135549in}}%
\pgfpathquadraticcurveto{\pgfqpoint{8.733993in}{1.135549in}}{\pgfqpoint{8.729310in}{1.130182in}}%
\pgfpathquadraticcurveto{\pgfqpoint{8.724645in}{1.124814in}}{\pgfqpoint{8.724645in}{1.115628in}}%
\pgfpathlineto{\pgfqpoint{8.724645in}{1.080000in}}%
\pgfpathlineto{\pgfqpoint{8.714234in}{1.080000in}}%
\pgfpathlineto{\pgfqpoint{8.714234in}{1.143043in}}%
\pgfpathlineto{\pgfqpoint{8.724645in}{1.143043in}}%
\pgfpathlineto{\pgfqpoint{8.724645in}{1.133244in}}%
\pgfpathquadraticcurveto{\pgfqpoint{8.728193in}{1.139044in}}{\pgfqpoint{8.733147in}{1.141800in}}%
\pgfpathquadraticcurveto{\pgfqpoint{8.738100in}{1.144556in}}{\pgfqpoint{8.744909in}{1.144556in}}%
\pgfpathquadraticcurveto{\pgfqpoint{8.751789in}{1.144556in}}{\pgfqpoint{8.756598in}{1.141061in}}%
\pgfpathquadraticcurveto{\pgfqpoint{8.761408in}{1.137585in}}{\pgfqpoint{8.763713in}{1.130938in}}%
\pgfpathlineto{\pgfqpoint{8.763713in}{1.130938in}}%
\pgfpathclose%
\pgfpathmoveto{\pgfqpoint{8.855575in}{1.111683in}}%
\pgfpathquadraticcurveto{\pgfqpoint{8.843021in}{1.111683in}}{\pgfqpoint{8.838175in}{1.108819in}}%
\pgfpathquadraticcurveto{\pgfqpoint{8.833330in}{1.105956in}}{\pgfqpoint{8.833330in}{1.099021in}}%
\pgfpathquadraticcurveto{\pgfqpoint{8.833330in}{1.093509in}}{\pgfqpoint{8.836969in}{1.090267in}}%
\pgfpathquadraticcurveto{\pgfqpoint{8.840607in}{1.087043in}}{\pgfqpoint{8.846839in}{1.087043in}}%
\pgfpathquadraticcurveto{\pgfqpoint{8.855467in}{1.087043in}}{\pgfqpoint{8.860673in}{1.093149in}}%
\pgfpathquadraticcurveto{\pgfqpoint{8.865878in}{1.099255in}}{\pgfqpoint{8.865878in}{1.109378in}}%
\pgfpathlineto{\pgfqpoint{8.865878in}{1.111683in}}%
\pgfpathlineto{\pgfqpoint{8.855575in}{1.111683in}}%
\pgfpathlineto{\pgfqpoint{8.855575in}{1.111683in}}%
\pgfpathclose%
\pgfpathmoveto{\pgfqpoint{8.876235in}{1.115970in}}%
\pgfpathlineto{\pgfqpoint{8.876235in}{1.080000in}}%
\pgfpathlineto{\pgfqpoint{8.865878in}{1.080000in}}%
\pgfpathlineto{\pgfqpoint{8.865878in}{1.089564in}}%
\pgfpathquadraticcurveto{\pgfqpoint{8.862330in}{1.083837in}}{\pgfqpoint{8.857034in}{1.081099in}}%
\pgfpathquadraticcurveto{\pgfqpoint{8.851739in}{1.078361in}}{\pgfqpoint{8.844083in}{1.078361in}}%
\pgfpathquadraticcurveto{\pgfqpoint{8.834411in}{1.078361in}}{\pgfqpoint{8.828683in}{1.083801in}}%
\pgfpathquadraticcurveto{\pgfqpoint{8.822973in}{1.089240in}}{\pgfqpoint{8.822973in}{1.098354in}}%
\pgfpathquadraticcurveto{\pgfqpoint{8.822973in}{1.108982in}}{\pgfqpoint{8.830088in}{1.114385in}}%
\pgfpathquadraticcurveto{\pgfqpoint{8.837221in}{1.119789in}}{\pgfqpoint{8.851342in}{1.119789in}}%
\pgfpathlineto{\pgfqpoint{8.865878in}{1.119789in}}%
\pgfpathlineto{\pgfqpoint{8.865878in}{1.120816in}}%
\pgfpathquadraticcurveto{\pgfqpoint{8.865878in}{1.127966in}}{\pgfqpoint{8.861177in}{1.131875in}}%
\pgfpathquadraticcurveto{\pgfqpoint{8.856476in}{1.135784in}}{\pgfqpoint{8.847974in}{1.135784in}}%
\pgfpathquadraticcurveto{\pgfqpoint{8.842570in}{1.135784in}}{\pgfqpoint{8.837437in}{1.134487in}}%
\pgfpathquadraticcurveto{\pgfqpoint{8.832321in}{1.133190in}}{\pgfqpoint{8.827602in}{1.130596in}}%
\pgfpathlineto{\pgfqpoint{8.827602in}{1.140179in}}%
\pgfpathquadraticcurveto{\pgfqpoint{8.833276in}{1.142376in}}{\pgfqpoint{8.838626in}{1.143457in}}%
\pgfpathquadraticcurveto{\pgfqpoint{8.843975in}{1.144556in}}{\pgfqpoint{8.849037in}{1.144556in}}%
\pgfpathquadraticcurveto{\pgfqpoint{8.862726in}{1.144556in}}{\pgfqpoint{8.869481in}{1.137459in}}%
\pgfpathquadraticcurveto{\pgfqpoint{8.876235in}{1.130380in}}{\pgfqpoint{8.876235in}{1.115970in}}%
\pgfpathlineto{\pgfqpoint{8.876235in}{1.115970in}}%
\pgfpathclose%
\pgfpathmoveto{\pgfqpoint{8.923787in}{1.074146in}}%
\pgfpathquadraticcurveto{\pgfqpoint{8.919410in}{1.062888in}}{\pgfqpoint{8.915231in}{1.059466in}}%
\pgfpathquadraticcurveto{\pgfqpoint{8.911071in}{1.056026in}}{\pgfqpoint{8.904100in}{1.056026in}}%
\pgfpathlineto{\pgfqpoint{8.895814in}{1.056026in}}%
\pgfpathlineto{\pgfqpoint{8.895814in}{1.064690in}}%
\pgfpathlineto{\pgfqpoint{8.901902in}{1.064690in}}%
\pgfpathquadraticcurveto{\pgfqpoint{8.906171in}{1.064690in}}{\pgfqpoint{8.908531in}{1.066725in}}%
\pgfpathquadraticcurveto{\pgfqpoint{8.910909in}{1.068742in}}{\pgfqpoint{8.913772in}{1.076289in}}%
\pgfpathlineto{\pgfqpoint{8.915628in}{1.081009in}}%
\pgfpathlineto{\pgfqpoint{8.890140in}{1.143043in}}%
\pgfpathlineto{\pgfqpoint{8.901110in}{1.143043in}}%
\pgfpathlineto{\pgfqpoint{8.920815in}{1.093743in}}%
\pgfpathlineto{\pgfqpoint{8.940520in}{1.143043in}}%
\pgfpathlineto{\pgfqpoint{8.951490in}{1.143043in}}%
\pgfpathlineto{\pgfqpoint{8.923787in}{1.074146in}}%
\pgfpathlineto{\pgfqpoint{8.923787in}{1.074146in}}%
\pgfpathclose%
\pgfusepath{fill}%
\end{pgfscope}%
\begin{pgfscope}%
\pgfsetbuttcap%
\pgfsetroundjoin%
\definecolor{currentfill}{rgb}{0.309804,0.372549,0.419608}%
\pgfsetfillcolor{currentfill}%
\pgfsetlinewidth{0.000000pt}%
\definecolor{currentstroke}{rgb}{0.309804,0.372549,0.419608}%
\pgfsetstrokecolor{currentstroke}%
\pgfsetdash{}{0pt}%
\pgfpathmoveto{\pgfqpoint{9.058009in}{1.133478in}}%
\pgfpathlineto{\pgfqpoint{9.058009in}{1.167593in}}%
\pgfpathlineto{\pgfqpoint{9.068366in}{1.167593in}}%
\pgfpathlineto{\pgfqpoint{9.068366in}{1.080000in}}%
\pgfpathlineto{\pgfqpoint{9.058009in}{1.080000in}}%
\pgfpathlineto{\pgfqpoint{9.058009in}{1.089456in}}%
\pgfpathquadraticcurveto{\pgfqpoint{9.054749in}{1.083837in}}{\pgfqpoint{9.049759in}{1.081099in}}%
\pgfpathquadraticcurveto{\pgfqpoint{9.044788in}{1.078361in}}{\pgfqpoint{9.037799in}{1.078361in}}%
\pgfpathquadraticcurveto{\pgfqpoint{9.026380in}{1.078361in}}{\pgfqpoint{9.019193in}{1.087475in}}%
\pgfpathquadraticcurveto{\pgfqpoint{9.012024in}{1.096607in}}{\pgfqpoint{9.012024in}{1.111467in}}%
\pgfpathquadraticcurveto{\pgfqpoint{9.012024in}{1.126327in}}{\pgfqpoint{9.019193in}{1.135441in}}%
\pgfpathquadraticcurveto{\pgfqpoint{9.026380in}{1.144556in}}{\pgfqpoint{9.037799in}{1.144556in}}%
\pgfpathquadraticcurveto{\pgfqpoint{9.044788in}{1.144556in}}{\pgfqpoint{9.049759in}{1.141818in}}%
\pgfpathquadraticcurveto{\pgfqpoint{9.054749in}{1.139098in}}{\pgfqpoint{9.058009in}{1.133478in}}%
\pgfpathlineto{\pgfqpoint{9.058009in}{1.133478in}}%
\pgfpathclose%
\pgfpathmoveto{\pgfqpoint{9.022723in}{1.111467in}}%
\pgfpathquadraticcurveto{\pgfqpoint{9.022723in}{1.100048in}}{\pgfqpoint{9.027424in}{1.093545in}}%
\pgfpathquadraticcurveto{\pgfqpoint{9.032125in}{1.087043in}}{\pgfqpoint{9.040339in}{1.087043in}}%
\pgfpathquadraticcurveto{\pgfqpoint{9.048553in}{1.087043in}}{\pgfqpoint{9.053272in}{1.093545in}}%
\pgfpathquadraticcurveto{\pgfqpoint{9.058009in}{1.100048in}}{\pgfqpoint{9.058009in}{1.111467in}}%
\pgfpathquadraticcurveto{\pgfqpoint{9.058009in}{1.122887in}}{\pgfqpoint{9.053272in}{1.129389in}}%
\pgfpathquadraticcurveto{\pgfqpoint{9.048553in}{1.135892in}}{\pgfqpoint{9.040339in}{1.135892in}}%
\pgfpathquadraticcurveto{\pgfqpoint{9.032125in}{1.135892in}}{\pgfqpoint{9.027424in}{1.129389in}}%
\pgfpathquadraticcurveto{\pgfqpoint{9.022723in}{1.122887in}}{\pgfqpoint{9.022723in}{1.111467in}}%
\pgfpathlineto{\pgfqpoint{9.022723in}{1.111467in}}%
\pgfpathclose%
\pgfpathmoveto{\pgfqpoint{9.089710in}{1.143043in}}%
\pgfpathlineto{\pgfqpoint{9.100067in}{1.143043in}}%
\pgfpathlineto{\pgfqpoint{9.100067in}{1.080000in}}%
\pgfpathlineto{\pgfqpoint{9.089710in}{1.080000in}}%
\pgfpathlineto{\pgfqpoint{9.089710in}{1.143043in}}%
\pgfpathlineto{\pgfqpoint{9.089710in}{1.143043in}}%
\pgfpathclose%
\pgfpathmoveto{\pgfqpoint{9.089710in}{1.167593in}}%
\pgfpathlineto{\pgfqpoint{9.100067in}{1.167593in}}%
\pgfpathlineto{\pgfqpoint{9.100067in}{1.154462in}}%
\pgfpathlineto{\pgfqpoint{9.089710in}{1.154462in}}%
\pgfpathlineto{\pgfqpoint{9.089710in}{1.167593in}}%
\pgfpathlineto{\pgfqpoint{9.089710in}{1.167593in}}%
\pgfpathclose%
\pgfpathmoveto{\pgfqpoint{9.192488in}{1.167593in}}%
\pgfpathlineto{\pgfqpoint{9.192488in}{1.158965in}}%
\pgfpathlineto{\pgfqpoint{9.182581in}{1.158965in}}%
\pgfpathquadraticcurveto{\pgfqpoint{9.177015in}{1.158965in}}{\pgfqpoint{9.174836in}{1.156714in}}%
\pgfpathquadraticcurveto{\pgfqpoint{9.172674in}{1.154462in}}{\pgfqpoint{9.172674in}{1.148608in}}%
\pgfpathlineto{\pgfqpoint{9.172674in}{1.143043in}}%
\pgfpathlineto{\pgfqpoint{9.189732in}{1.143043in}}%
\pgfpathlineto{\pgfqpoint{9.189732in}{1.134991in}}%
\pgfpathlineto{\pgfqpoint{9.172674in}{1.134991in}}%
\pgfpathlineto{\pgfqpoint{9.172674in}{1.080000in}}%
\pgfpathlineto{\pgfqpoint{9.162263in}{1.080000in}}%
\pgfpathlineto{\pgfqpoint{9.162263in}{1.134991in}}%
\pgfpathlineto{\pgfqpoint{9.133840in}{1.134991in}}%
\pgfpathlineto{\pgfqpoint{9.133840in}{1.080000in}}%
\pgfpathlineto{\pgfqpoint{9.123429in}{1.080000in}}%
\pgfpathlineto{\pgfqpoint{9.123429in}{1.134991in}}%
\pgfpathlineto{\pgfqpoint{9.113522in}{1.134991in}}%
\pgfpathlineto{\pgfqpoint{9.113522in}{1.143043in}}%
\pgfpathlineto{\pgfqpoint{9.123429in}{1.143043in}}%
\pgfpathlineto{\pgfqpoint{9.123429in}{1.147438in}}%
\pgfpathquadraticcurveto{\pgfqpoint{9.123429in}{1.157957in}}{\pgfqpoint{9.128328in}{1.162766in}}%
\pgfpathquadraticcurveto{\pgfqpoint{9.133228in}{1.167593in}}{\pgfqpoint{9.143855in}{1.167593in}}%
\pgfpathlineto{\pgfqpoint{9.153653in}{1.167593in}}%
\pgfpathlineto{\pgfqpoint{9.153653in}{1.158965in}}%
\pgfpathlineto{\pgfqpoint{9.143747in}{1.158965in}}%
\pgfpathquadraticcurveto{\pgfqpoint{9.138181in}{1.158965in}}{\pgfqpoint{9.135984in}{1.156714in}}%
\pgfpathquadraticcurveto{\pgfqpoint{9.133840in}{1.154462in}}{\pgfqpoint{9.133840in}{1.148608in}}%
\pgfpathlineto{\pgfqpoint{9.133840in}{1.143043in}}%
\pgfpathlineto{\pgfqpoint{9.162263in}{1.143043in}}%
\pgfpathlineto{\pgfqpoint{9.162263in}{1.147438in}}%
\pgfpathquadraticcurveto{\pgfqpoint{9.162263in}{1.157957in}}{\pgfqpoint{9.167163in}{1.162766in}}%
\pgfpathquadraticcurveto{\pgfqpoint{9.172062in}{1.167593in}}{\pgfqpoint{9.182707in}{1.167593in}}%
\pgfpathlineto{\pgfqpoint{9.192488in}{1.167593in}}%
\pgfpathlineto{\pgfqpoint{9.192488in}{1.167593in}}%
\pgfpathclose%
\pgfpathmoveto{\pgfqpoint{9.255080in}{1.114115in}}%
\pgfpathlineto{\pgfqpoint{9.255080in}{1.109054in}}%
\pgfpathlineto{\pgfqpoint{9.207456in}{1.109054in}}%
\pgfpathquadraticcurveto{\pgfqpoint{9.208140in}{1.098354in}}{\pgfqpoint{9.213904in}{1.092753in}}%
\pgfpathquadraticcurveto{\pgfqpoint{9.219668in}{1.087151in}}{\pgfqpoint{9.229971in}{1.087151in}}%
\pgfpathquadraticcurveto{\pgfqpoint{9.235933in}{1.087151in}}{\pgfqpoint{9.241535in}{1.088610in}}%
\pgfpathquadraticcurveto{\pgfqpoint{9.247137in}{1.090069in}}{\pgfqpoint{9.252666in}{1.093005in}}%
\pgfpathlineto{\pgfqpoint{9.252666in}{1.083206in}}%
\pgfpathquadraticcurveto{\pgfqpoint{9.247083in}{1.080847in}}{\pgfqpoint{9.241229in}{1.079604in}}%
\pgfpathquadraticcurveto{\pgfqpoint{9.235375in}{1.078361in}}{\pgfqpoint{9.229359in}{1.078361in}}%
\pgfpathquadraticcurveto{\pgfqpoint{9.214264in}{1.078361in}}{\pgfqpoint{9.205456in}{1.087133in}}%
\pgfpathquadraticcurveto{\pgfqpoint{9.196648in}{1.095923in}}{\pgfqpoint{9.196648in}{1.110909in}}%
\pgfpathquadraticcurveto{\pgfqpoint{9.196648in}{1.126381in}}{\pgfqpoint{9.205006in}{1.135459in}}%
\pgfpathquadraticcurveto{\pgfqpoint{9.213364in}{1.144556in}}{\pgfqpoint{9.227557in}{1.144556in}}%
\pgfpathquadraticcurveto{\pgfqpoint{9.240274in}{1.144556in}}{\pgfqpoint{9.247677in}{1.136360in}}%
\pgfpathquadraticcurveto{\pgfqpoint{9.255080in}{1.128183in}}{\pgfqpoint{9.255080in}{1.114115in}}%
\pgfpathlineto{\pgfqpoint{9.255080in}{1.114115in}}%
\pgfpathclose%
\pgfpathmoveto{\pgfqpoint{9.244723in}{1.117159in}}%
\pgfpathquadraticcurveto{\pgfqpoint{9.244615in}{1.125643in}}{\pgfqpoint{9.239968in}{1.130704in}}%
\pgfpathquadraticcurveto{\pgfqpoint{9.235321in}{1.135784in}}{\pgfqpoint{9.227665in}{1.135784in}}%
\pgfpathquadraticcurveto{\pgfqpoint{9.219002in}{1.135784in}}{\pgfqpoint{9.213796in}{1.130884in}}%
\pgfpathquadraticcurveto{\pgfqpoint{9.208591in}{1.125985in}}{\pgfqpoint{9.207798in}{1.117087in}}%
\pgfpathlineto{\pgfqpoint{9.244723in}{1.117159in}}%
\pgfpathlineto{\pgfqpoint{9.244723in}{1.117159in}}%
\pgfpathclose%
\pgfpathmoveto{\pgfqpoint{9.308612in}{1.133370in}}%
\pgfpathquadraticcurveto{\pgfqpoint{9.306865in}{1.134379in}}{\pgfqpoint{9.304811in}{1.134847in}}%
\pgfpathquadraticcurveto{\pgfqpoint{9.302758in}{1.135333in}}{\pgfqpoint{9.300290in}{1.135333in}}%
\pgfpathquadraticcurveto{\pgfqpoint{9.291500in}{1.135333in}}{\pgfqpoint{9.286799in}{1.129623in}}%
\pgfpathquadraticcurveto{\pgfqpoint{9.282098in}{1.123914in}}{\pgfqpoint{9.282098in}{1.113214in}}%
\pgfpathlineto{\pgfqpoint{9.282098in}{1.080000in}}%
\pgfpathlineto{\pgfqpoint{9.271687in}{1.080000in}}%
\pgfpathlineto{\pgfqpoint{9.271687in}{1.143043in}}%
\pgfpathlineto{\pgfqpoint{9.282098in}{1.143043in}}%
\pgfpathlineto{\pgfqpoint{9.282098in}{1.133244in}}%
\pgfpathquadraticcurveto{\pgfqpoint{9.285376in}{1.138990in}}{\pgfqpoint{9.290600in}{1.141764in}}%
\pgfpathquadraticcurveto{\pgfqpoint{9.295841in}{1.144556in}}{\pgfqpoint{9.303334in}{1.144556in}}%
\pgfpathquadraticcurveto{\pgfqpoint{9.304397in}{1.144556in}}{\pgfqpoint{9.305694in}{1.144411in}}%
\pgfpathquadraticcurveto{\pgfqpoint{9.306991in}{1.144285in}}{\pgfqpoint{9.308558in}{1.143997in}}%
\pgfpathlineto{\pgfqpoint{9.308612in}{1.133370in}}%
\pgfpathlineto{\pgfqpoint{9.308612in}{1.133370in}}%
\pgfpathclose%
\pgfusepath{fill}%
\end{pgfscope}%
\begin{pgfscope}%
\pgfsetbuttcap%
\pgfsetroundjoin%
\definecolor{currentfill}{rgb}{0.309804,0.372549,0.419608}%
\pgfsetfillcolor{currentfill}%
\pgfsetlinewidth{0.000000pt}%
\definecolor{currentstroke}{rgb}{0.309804,0.372549,0.419608}%
\pgfsetstrokecolor{currentstroke}%
\pgfsetdash{}{0pt}%
\pgfpathmoveto{\pgfqpoint{9.400341in}{1.167593in}}%
\pgfpathlineto{\pgfqpoint{9.400341in}{1.158965in}}%
\pgfpathlineto{\pgfqpoint{9.390435in}{1.158965in}}%
\pgfpathquadraticcurveto{\pgfqpoint{9.384869in}{1.158965in}}{\pgfqpoint{9.382689in}{1.156714in}}%
\pgfpathquadraticcurveto{\pgfqpoint{9.380528in}{1.154462in}}{\pgfqpoint{9.380528in}{1.148608in}}%
\pgfpathlineto{\pgfqpoint{9.380528in}{1.143043in}}%
\pgfpathlineto{\pgfqpoint{9.397586in}{1.143043in}}%
\pgfpathlineto{\pgfqpoint{9.397586in}{1.134991in}}%
\pgfpathlineto{\pgfqpoint{9.380528in}{1.134991in}}%
\pgfpathlineto{\pgfqpoint{9.380528in}{1.080000in}}%
\pgfpathlineto{\pgfqpoint{9.370117in}{1.080000in}}%
\pgfpathlineto{\pgfqpoint{9.370117in}{1.134991in}}%
\pgfpathlineto{\pgfqpoint{9.360210in}{1.134991in}}%
\pgfpathlineto{\pgfqpoint{9.360210in}{1.143043in}}%
\pgfpathlineto{\pgfqpoint{9.370117in}{1.143043in}}%
\pgfpathlineto{\pgfqpoint{9.370117in}{1.147438in}}%
\pgfpathquadraticcurveto{\pgfqpoint{9.370117in}{1.157957in}}{\pgfqpoint{9.375016in}{1.162766in}}%
\pgfpathquadraticcurveto{\pgfqpoint{9.379916in}{1.167593in}}{\pgfqpoint{9.390543in}{1.167593in}}%
\pgfpathlineto{\pgfqpoint{9.400341in}{1.167593in}}%
\pgfpathlineto{\pgfqpoint{9.400341in}{1.167593in}}%
\pgfpathclose%
\pgfpathmoveto{\pgfqpoint{9.445534in}{1.133370in}}%
\pgfpathquadraticcurveto{\pgfqpoint{9.443787in}{1.134379in}}{\pgfqpoint{9.441733in}{1.134847in}}%
\pgfpathquadraticcurveto{\pgfqpoint{9.439680in}{1.135333in}}{\pgfqpoint{9.437212in}{1.135333in}}%
\pgfpathquadraticcurveto{\pgfqpoint{9.428422in}{1.135333in}}{\pgfqpoint{9.423721in}{1.129623in}}%
\pgfpathquadraticcurveto{\pgfqpoint{9.419020in}{1.123914in}}{\pgfqpoint{9.419020in}{1.113214in}}%
\pgfpathlineto{\pgfqpoint{9.419020in}{1.080000in}}%
\pgfpathlineto{\pgfqpoint{9.408609in}{1.080000in}}%
\pgfpathlineto{\pgfqpoint{9.408609in}{1.143043in}}%
\pgfpathlineto{\pgfqpoint{9.419020in}{1.143043in}}%
\pgfpathlineto{\pgfqpoint{9.419020in}{1.133244in}}%
\pgfpathquadraticcurveto{\pgfqpoint{9.422298in}{1.138990in}}{\pgfqpoint{9.427522in}{1.141764in}}%
\pgfpathquadraticcurveto{\pgfqpoint{9.432763in}{1.144556in}}{\pgfqpoint{9.440256in}{1.144556in}}%
\pgfpathquadraticcurveto{\pgfqpoint{9.441319in}{1.144556in}}{\pgfqpoint{9.442616in}{1.144411in}}%
\pgfpathquadraticcurveto{\pgfqpoint{9.443913in}{1.144285in}}{\pgfqpoint{9.445480in}{1.143997in}}%
\pgfpathlineto{\pgfqpoint{9.445534in}{1.133370in}}%
\pgfpathlineto{\pgfqpoint{9.445534in}{1.133370in}}%
\pgfpathclose%
\pgfpathmoveto{\pgfqpoint{9.478280in}{1.135784in}}%
\pgfpathquadraticcurveto{\pgfqpoint{9.469958in}{1.135784in}}{\pgfqpoint{9.465113in}{1.129281in}}%
\pgfpathquadraticcurveto{\pgfqpoint{9.460268in}{1.122779in}}{\pgfqpoint{9.460268in}{1.111467in}}%
\pgfpathquadraticcurveto{\pgfqpoint{9.460268in}{1.100156in}}{\pgfqpoint{9.465077in}{1.093653in}}%
\pgfpathquadraticcurveto{\pgfqpoint{9.469904in}{1.087151in}}{\pgfqpoint{9.478280in}{1.087151in}}%
\pgfpathquadraticcurveto{\pgfqpoint{9.486566in}{1.087151in}}{\pgfqpoint{9.491393in}{1.093671in}}%
\pgfpathquadraticcurveto{\pgfqpoint{9.496238in}{1.100210in}}{\pgfqpoint{9.496238in}{1.111467in}}%
\pgfpathquadraticcurveto{\pgfqpoint{9.496238in}{1.122671in}}{\pgfqpoint{9.491393in}{1.129227in}}%
\pgfpathquadraticcurveto{\pgfqpoint{9.486566in}{1.135784in}}{\pgfqpoint{9.478280in}{1.135784in}}%
\pgfpathlineto{\pgfqpoint{9.478280in}{1.135784in}}%
\pgfpathclose%
\pgfpathmoveto{\pgfqpoint{9.478280in}{1.144556in}}%
\pgfpathquadraticcurveto{\pgfqpoint{9.491789in}{1.144556in}}{\pgfqpoint{9.499498in}{1.135766in}}%
\pgfpathquadraticcurveto{\pgfqpoint{9.507225in}{1.126994in}}{\pgfqpoint{9.507225in}{1.111467in}}%
\pgfpathquadraticcurveto{\pgfqpoint{9.507225in}{1.095995in}}{\pgfqpoint{9.499498in}{1.087169in}}%
\pgfpathquadraticcurveto{\pgfqpoint{9.491789in}{1.078361in}}{\pgfqpoint{9.478280in}{1.078361in}}%
\pgfpathquadraticcurveto{\pgfqpoint{9.464717in}{1.078361in}}{\pgfqpoint{9.457026in}{1.087169in}}%
\pgfpathquadraticcurveto{\pgfqpoint{9.449352in}{1.095995in}}{\pgfqpoint{9.449352in}{1.111467in}}%
\pgfpathquadraticcurveto{\pgfqpoint{9.449352in}{1.126994in}}{\pgfqpoint{9.457026in}{1.135766in}}%
\pgfpathquadraticcurveto{\pgfqpoint{9.464717in}{1.144556in}}{\pgfqpoint{9.478280in}{1.144556in}}%
\pgfpathlineto{\pgfqpoint{9.478280in}{1.144556in}}%
\pgfpathclose%
\pgfpathmoveto{\pgfqpoint{9.573474in}{1.130938in}}%
\pgfpathquadraticcurveto{\pgfqpoint{9.577365in}{1.137927in}}{\pgfqpoint{9.582768in}{1.141241in}}%
\pgfpathquadraticcurveto{\pgfqpoint{9.588172in}{1.144556in}}{\pgfqpoint{9.595485in}{1.144556in}}%
\pgfpathquadraticcurveto{\pgfqpoint{9.605338in}{1.144556in}}{\pgfqpoint{9.610687in}{1.137657in}}%
\pgfpathquadraticcurveto{\pgfqpoint{9.616037in}{1.130776in}}{\pgfqpoint{9.616037in}{1.118060in}}%
\pgfpathlineto{\pgfqpoint{9.616037in}{1.080000in}}%
\pgfpathlineto{\pgfqpoint{9.605626in}{1.080000in}}%
\pgfpathlineto{\pgfqpoint{9.605626in}{1.117717in}}%
\pgfpathquadraticcurveto{\pgfqpoint{9.605626in}{1.126778in}}{\pgfqpoint{9.602402in}{1.131155in}}%
\pgfpathquadraticcurveto{\pgfqpoint{9.599196in}{1.135549in}}{\pgfqpoint{9.592621in}{1.135549in}}%
\pgfpathquadraticcurveto{\pgfqpoint{9.584570in}{1.135549in}}{\pgfqpoint{9.579887in}{1.130200in}}%
\pgfpathquadraticcurveto{\pgfqpoint{9.575221in}{1.124868in}}{\pgfqpoint{9.575221in}{1.115628in}}%
\pgfpathlineto{\pgfqpoint{9.575221in}{1.080000in}}%
\pgfpathlineto{\pgfqpoint{9.564810in}{1.080000in}}%
\pgfpathlineto{\pgfqpoint{9.564810in}{1.117717in}}%
\pgfpathquadraticcurveto{\pgfqpoint{9.564810in}{1.126832in}}{\pgfqpoint{9.561604in}{1.131191in}}%
\pgfpathquadraticcurveto{\pgfqpoint{9.558398in}{1.135549in}}{\pgfqpoint{9.551697in}{1.135549in}}%
\pgfpathquadraticcurveto{\pgfqpoint{9.543754in}{1.135549in}}{\pgfqpoint{9.539071in}{1.130182in}}%
\pgfpathquadraticcurveto{\pgfqpoint{9.534406in}{1.124814in}}{\pgfqpoint{9.534406in}{1.115628in}}%
\pgfpathlineto{\pgfqpoint{9.534406in}{1.080000in}}%
\pgfpathlineto{\pgfqpoint{9.523995in}{1.080000in}}%
\pgfpathlineto{\pgfqpoint{9.523995in}{1.143043in}}%
\pgfpathlineto{\pgfqpoint{9.534406in}{1.143043in}}%
\pgfpathlineto{\pgfqpoint{9.534406in}{1.133244in}}%
\pgfpathquadraticcurveto{\pgfqpoint{9.537954in}{1.139044in}}{\pgfqpoint{9.542908in}{1.141800in}}%
\pgfpathquadraticcurveto{\pgfqpoint{9.547861in}{1.144556in}}{\pgfqpoint{9.554669in}{1.144556in}}%
\pgfpathquadraticcurveto{\pgfqpoint{9.561550in}{1.144556in}}{\pgfqpoint{9.566359in}{1.141061in}}%
\pgfpathquadraticcurveto{\pgfqpoint{9.571169in}{1.137585in}}{\pgfqpoint{9.573474in}{1.130938in}}%
\pgfpathlineto{\pgfqpoint{9.573474in}{1.130938in}}%
\pgfpathclose%
\pgfusepath{fill}%
\end{pgfscope}%
\begin{pgfscope}%
\pgfsetbuttcap%
\pgfsetroundjoin%
\definecolor{currentfill}{rgb}{0.309804,0.372549,0.419608}%
\pgfsetfillcolor{currentfill}%
\pgfsetlinewidth{0.000000pt}%
\definecolor{currentstroke}{rgb}{0.309804,0.372549,0.419608}%
\pgfsetstrokecolor{currentstroke}%
\pgfsetdash{}{0pt}%
\pgfpathmoveto{\pgfqpoint{9.716478in}{1.111683in}}%
\pgfpathquadraticcurveto{\pgfqpoint{9.703924in}{1.111683in}}{\pgfqpoint{9.699078in}{1.108819in}}%
\pgfpathquadraticcurveto{\pgfqpoint{9.694233in}{1.105956in}}{\pgfqpoint{9.694233in}{1.099021in}}%
\pgfpathquadraticcurveto{\pgfqpoint{9.694233in}{1.093509in}}{\pgfqpoint{9.697871in}{1.090267in}}%
\pgfpathquadraticcurveto{\pgfqpoint{9.701510in}{1.087043in}}{\pgfqpoint{9.707742in}{1.087043in}}%
\pgfpathquadraticcurveto{\pgfqpoint{9.716370in}{1.087043in}}{\pgfqpoint{9.721575in}{1.093149in}}%
\pgfpathquadraticcurveto{\pgfqpoint{9.726781in}{1.099255in}}{\pgfqpoint{9.726781in}{1.109378in}}%
\pgfpathlineto{\pgfqpoint{9.726781in}{1.111683in}}%
\pgfpathlineto{\pgfqpoint{9.716478in}{1.111683in}}%
\pgfpathlineto{\pgfqpoint{9.716478in}{1.111683in}}%
\pgfpathclose%
\pgfpathmoveto{\pgfqpoint{9.737138in}{1.115970in}}%
\pgfpathlineto{\pgfqpoint{9.737138in}{1.080000in}}%
\pgfpathlineto{\pgfqpoint{9.726781in}{1.080000in}}%
\pgfpathlineto{\pgfqpoint{9.726781in}{1.089564in}}%
\pgfpathquadraticcurveto{\pgfqpoint{9.723233in}{1.083837in}}{\pgfqpoint{9.717937in}{1.081099in}}%
\pgfpathquadraticcurveto{\pgfqpoint{9.712641in}{1.078361in}}{\pgfqpoint{9.704986in}{1.078361in}}%
\pgfpathquadraticcurveto{\pgfqpoint{9.695314in}{1.078361in}}{\pgfqpoint{9.689586in}{1.083801in}}%
\pgfpathquadraticcurveto{\pgfqpoint{9.683876in}{1.089240in}}{\pgfqpoint{9.683876in}{1.098354in}}%
\pgfpathquadraticcurveto{\pgfqpoint{9.683876in}{1.108982in}}{\pgfqpoint{9.690991in}{1.114385in}}%
\pgfpathquadraticcurveto{\pgfqpoint{9.698124in}{1.119789in}}{\pgfqpoint{9.712245in}{1.119789in}}%
\pgfpathlineto{\pgfqpoint{9.726781in}{1.119789in}}%
\pgfpathlineto{\pgfqpoint{9.726781in}{1.120816in}}%
\pgfpathquadraticcurveto{\pgfqpoint{9.726781in}{1.127966in}}{\pgfqpoint{9.722080in}{1.131875in}}%
\pgfpathquadraticcurveto{\pgfqpoint{9.717379in}{1.135784in}}{\pgfqpoint{9.708877in}{1.135784in}}%
\pgfpathquadraticcurveto{\pgfqpoint{9.703473in}{1.135784in}}{\pgfqpoint{9.698340in}{1.134487in}}%
\pgfpathquadraticcurveto{\pgfqpoint{9.693224in}{1.133190in}}{\pgfqpoint{9.688505in}{1.130596in}}%
\pgfpathlineto{\pgfqpoint{9.688505in}{1.140179in}}%
\pgfpathquadraticcurveto{\pgfqpoint{9.694179in}{1.142376in}}{\pgfqpoint{9.699529in}{1.143457in}}%
\pgfpathquadraticcurveto{\pgfqpoint{9.704878in}{1.144556in}}{\pgfqpoint{9.709940in}{1.144556in}}%
\pgfpathquadraticcurveto{\pgfqpoint{9.723629in}{1.144556in}}{\pgfqpoint{9.730383in}{1.137459in}}%
\pgfpathquadraticcurveto{\pgfqpoint{9.737138in}{1.130380in}}{\pgfqpoint{9.737138in}{1.115970in}}%
\pgfpathlineto{\pgfqpoint{9.737138in}{1.115970in}}%
\pgfpathclose%
\pgfpathmoveto{\pgfqpoint{9.803837in}{1.140629in}}%
\pgfpathlineto{\pgfqpoint{9.803837in}{1.130938in}}%
\pgfpathquadraticcurveto{\pgfqpoint{9.799442in}{1.133370in}}{\pgfqpoint{9.795029in}{1.134577in}}%
\pgfpathquadraticcurveto{\pgfqpoint{9.790616in}{1.135784in}}{\pgfqpoint{9.786113in}{1.135784in}}%
\pgfpathquadraticcurveto{\pgfqpoint{9.776026in}{1.135784in}}{\pgfqpoint{9.770442in}{1.129389in}}%
\pgfpathquadraticcurveto{\pgfqpoint{9.764877in}{1.123013in}}{\pgfqpoint{9.764877in}{1.111467in}}%
\pgfpathquadraticcurveto{\pgfqpoint{9.764877in}{1.099921in}}{\pgfqpoint{9.770442in}{1.093527in}}%
\pgfpathquadraticcurveto{\pgfqpoint{9.776026in}{1.087151in}}{\pgfqpoint{9.786113in}{1.087151in}}%
\pgfpathquadraticcurveto{\pgfqpoint{9.790616in}{1.087151in}}{\pgfqpoint{9.795029in}{1.088358in}}%
\pgfpathquadraticcurveto{\pgfqpoint{9.799442in}{1.089564in}}{\pgfqpoint{9.803837in}{1.091996in}}%
\pgfpathlineto{\pgfqpoint{9.803837in}{1.082414in}}%
\pgfpathquadraticcurveto{\pgfqpoint{9.799496in}{1.080396in}}{\pgfqpoint{9.794849in}{1.079388in}}%
\pgfpathquadraticcurveto{\pgfqpoint{9.790220in}{1.078361in}}{\pgfqpoint{9.784978in}{1.078361in}}%
\pgfpathquadraticcurveto{\pgfqpoint{9.770731in}{1.078361in}}{\pgfqpoint{9.762337in}{1.087313in}}%
\pgfpathquadraticcurveto{\pgfqpoint{9.753961in}{1.096265in}}{\pgfqpoint{9.753961in}{1.111467in}}%
\pgfpathquadraticcurveto{\pgfqpoint{9.753961in}{1.126886in}}{\pgfqpoint{9.762427in}{1.135712in}}%
\pgfpathquadraticcurveto{\pgfqpoint{9.770911in}{1.144556in}}{\pgfqpoint{9.785663in}{1.144556in}}%
\pgfpathquadraticcurveto{\pgfqpoint{9.790436in}{1.144556in}}{\pgfqpoint{9.794993in}{1.143565in}}%
\pgfpathquadraticcurveto{\pgfqpoint{9.799550in}{1.142592in}}{\pgfqpoint{9.803837in}{1.140629in}}%
\pgfpathlineto{\pgfqpoint{9.803837in}{1.140629in}}%
\pgfpathclose%
\pgfpathmoveto{\pgfqpoint{9.832098in}{1.160947in}}%
\pgfpathlineto{\pgfqpoint{9.832098in}{1.143043in}}%
\pgfpathlineto{\pgfqpoint{9.853424in}{1.143043in}}%
\pgfpathlineto{\pgfqpoint{9.853424in}{1.134991in}}%
\pgfpathlineto{\pgfqpoint{9.832098in}{1.134991in}}%
\pgfpathlineto{\pgfqpoint{9.832098in}{1.100768in}}%
\pgfpathquadraticcurveto{\pgfqpoint{9.832098in}{1.093059in}}{\pgfqpoint{9.834205in}{1.090861in}}%
\pgfpathquadraticcurveto{\pgfqpoint{9.836313in}{1.088664in}}{\pgfqpoint{9.842797in}{1.088664in}}%
\pgfpathlineto{\pgfqpoint{9.853424in}{1.088664in}}%
\pgfpathlineto{\pgfqpoint{9.853424in}{1.080000in}}%
\pgfpathlineto{\pgfqpoint{9.842797in}{1.080000in}}%
\pgfpathquadraticcurveto{\pgfqpoint{9.830801in}{1.080000in}}{\pgfqpoint{9.826244in}{1.084467in}}%
\pgfpathquadraticcurveto{\pgfqpoint{9.821687in}{1.088952in}}{\pgfqpoint{9.821687in}{1.100768in}}%
\pgfpathlineto{\pgfqpoint{9.821687in}{1.134991in}}%
\pgfpathlineto{\pgfqpoint{9.814086in}{1.134991in}}%
\pgfpathlineto{\pgfqpoint{9.814086in}{1.143043in}}%
\pgfpathlineto{\pgfqpoint{9.821687in}{1.143043in}}%
\pgfpathlineto{\pgfqpoint{9.821687in}{1.160947in}}%
\pgfpathlineto{\pgfqpoint{9.832098in}{1.160947in}}%
\pgfpathlineto{\pgfqpoint{9.832098in}{1.160947in}}%
\pgfpathclose%
\pgfpathmoveto{\pgfqpoint{9.865979in}{1.104875in}}%
\pgfpathlineto{\pgfqpoint{9.865979in}{1.143043in}}%
\pgfpathlineto{\pgfqpoint{9.876336in}{1.143043in}}%
\pgfpathlineto{\pgfqpoint{9.876336in}{1.105271in}}%
\pgfpathquadraticcurveto{\pgfqpoint{9.876336in}{1.096319in}}{\pgfqpoint{9.879812in}{1.091834in}}%
\pgfpathquadraticcurveto{\pgfqpoint{9.883307in}{1.087367in}}{\pgfqpoint{9.890295in}{1.087367in}}%
\pgfpathquadraticcurveto{\pgfqpoint{9.898671in}{1.087367in}}{\pgfqpoint{9.903534in}{1.092717in}}%
\pgfpathquadraticcurveto{\pgfqpoint{9.908416in}{1.098066in}}{\pgfqpoint{9.908416in}{1.107306in}}%
\pgfpathlineto{\pgfqpoint{9.908416in}{1.143043in}}%
\pgfpathlineto{\pgfqpoint{9.918773in}{1.143043in}}%
\pgfpathlineto{\pgfqpoint{9.918773in}{1.080000in}}%
\pgfpathlineto{\pgfqpoint{9.908416in}{1.080000in}}%
\pgfpathlineto{\pgfqpoint{9.908416in}{1.089691in}}%
\pgfpathquadraticcurveto{\pgfqpoint{9.904651in}{1.083945in}}{\pgfqpoint{9.899662in}{1.081153in}}%
\pgfpathquadraticcurveto{\pgfqpoint{9.894690in}{1.078361in}}{\pgfqpoint{9.888098in}{1.078361in}}%
\pgfpathquadraticcurveto{\pgfqpoint{9.877237in}{1.078361in}}{\pgfqpoint{9.871599in}{1.085115in}}%
\pgfpathquadraticcurveto{\pgfqpoint{9.865979in}{1.091870in}}{\pgfqpoint{9.865979in}{1.104875in}}%
\pgfpathlineto{\pgfqpoint{9.865979in}{1.104875in}}%
\pgfpathclose%
\pgfpathmoveto{\pgfqpoint{9.892042in}{1.144556in}}%
\pgfpathlineto{\pgfqpoint{9.892042in}{1.144556in}}%
\pgfpathlineto{\pgfqpoint{9.892042in}{1.144556in}}%
\pgfpathclose%
\pgfpathmoveto{\pgfqpoint{9.968756in}{1.111683in}}%
\pgfpathquadraticcurveto{\pgfqpoint{9.956202in}{1.111683in}}{\pgfqpoint{9.951357in}{1.108819in}}%
\pgfpathquadraticcurveto{\pgfqpoint{9.946511in}{1.105956in}}{\pgfqpoint{9.946511in}{1.099021in}}%
\pgfpathquadraticcurveto{\pgfqpoint{9.946511in}{1.093509in}}{\pgfqpoint{9.950150in}{1.090267in}}%
\pgfpathquadraticcurveto{\pgfqpoint{9.953788in}{1.087043in}}{\pgfqpoint{9.960020in}{1.087043in}}%
\pgfpathquadraticcurveto{\pgfqpoint{9.968648in}{1.087043in}}{\pgfqpoint{9.973854in}{1.093149in}}%
\pgfpathquadraticcurveto{\pgfqpoint{9.979059in}{1.099255in}}{\pgfqpoint{9.979059in}{1.109378in}}%
\pgfpathlineto{\pgfqpoint{9.979059in}{1.111683in}}%
\pgfpathlineto{\pgfqpoint{9.968756in}{1.111683in}}%
\pgfpathlineto{\pgfqpoint{9.968756in}{1.111683in}}%
\pgfpathclose%
\pgfpathmoveto{\pgfqpoint{9.989416in}{1.115970in}}%
\pgfpathlineto{\pgfqpoint{9.989416in}{1.080000in}}%
\pgfpathlineto{\pgfqpoint{9.979059in}{1.080000in}}%
\pgfpathlineto{\pgfqpoint{9.979059in}{1.089564in}}%
\pgfpathquadraticcurveto{\pgfqpoint{9.975511in}{1.083837in}}{\pgfqpoint{9.970215in}{1.081099in}}%
\pgfpathquadraticcurveto{\pgfqpoint{9.964920in}{1.078361in}}{\pgfqpoint{9.957264in}{1.078361in}}%
\pgfpathquadraticcurveto{\pgfqpoint{9.947592in}{1.078361in}}{\pgfqpoint{9.941864in}{1.083801in}}%
\pgfpathquadraticcurveto{\pgfqpoint{9.936154in}{1.089240in}}{\pgfqpoint{9.936154in}{1.098354in}}%
\pgfpathquadraticcurveto{\pgfqpoint{9.936154in}{1.108982in}}{\pgfqpoint{9.943269in}{1.114385in}}%
\pgfpathquadraticcurveto{\pgfqpoint{9.950402in}{1.119789in}}{\pgfqpoint{9.964523in}{1.119789in}}%
\pgfpathlineto{\pgfqpoint{9.979059in}{1.119789in}}%
\pgfpathlineto{\pgfqpoint{9.979059in}{1.120816in}}%
\pgfpathquadraticcurveto{\pgfqpoint{9.979059in}{1.127966in}}{\pgfqpoint{9.974358in}{1.131875in}}%
\pgfpathquadraticcurveto{\pgfqpoint{9.969657in}{1.135784in}}{\pgfqpoint{9.961155in}{1.135784in}}%
\pgfpathquadraticcurveto{\pgfqpoint{9.955751in}{1.135784in}}{\pgfqpoint{9.950618in}{1.134487in}}%
\pgfpathquadraticcurveto{\pgfqpoint{9.945503in}{1.133190in}}{\pgfqpoint{9.940783in}{1.130596in}}%
\pgfpathlineto{\pgfqpoint{9.940783in}{1.140179in}}%
\pgfpathquadraticcurveto{\pgfqpoint{9.946457in}{1.142376in}}{\pgfqpoint{9.951807in}{1.143457in}}%
\pgfpathquadraticcurveto{\pgfqpoint{9.957156in}{1.144556in}}{\pgfqpoint{9.962218in}{1.144556in}}%
\pgfpathquadraticcurveto{\pgfqpoint{9.975907in}{1.144556in}}{\pgfqpoint{9.982662in}{1.137459in}}%
\pgfpathquadraticcurveto{\pgfqpoint{9.989416in}{1.130380in}}{\pgfqpoint{9.989416in}{1.115970in}}%
\pgfpathlineto{\pgfqpoint{9.989416in}{1.115970in}}%
\pgfpathclose%
\pgfpathmoveto{\pgfqpoint{10.010743in}{1.167593in}}%
\pgfpathlineto{\pgfqpoint{10.021100in}{1.167593in}}%
\pgfpathlineto{\pgfqpoint{10.021100in}{1.080000in}}%
\pgfpathlineto{\pgfqpoint{10.010743in}{1.080000in}}%
\pgfpathlineto{\pgfqpoint{10.010743in}{1.167593in}}%
\pgfpathlineto{\pgfqpoint{10.010743in}{1.167593in}}%
\pgfpathclose%
\pgfusepath{fill}%
\end{pgfscope}%
\begin{pgfscope}%
\pgfsetbuttcap%
\pgfsetroundjoin%
\definecolor{currentfill}{rgb}{0.309804,0.372549,0.419608}%
\pgfsetfillcolor{currentfill}%
\pgfsetlinewidth{0.000000pt}%
\definecolor{currentstroke}{rgb}{0.309804,0.372549,0.419608}%
\pgfsetstrokecolor{currentstroke}%
\pgfsetdash{}{0pt}%
\pgfpathmoveto{\pgfqpoint{10.141200in}{1.112260in}}%
\pgfpathquadraticcurveto{\pgfqpoint{10.141200in}{1.123517in}}{\pgfqpoint{10.136552in}{1.129696in}}%
\pgfpathquadraticcurveto{\pgfqpoint{10.131923in}{1.135892in}}{\pgfqpoint{10.123530in}{1.135892in}}%
\pgfpathquadraticcurveto{\pgfqpoint{10.115208in}{1.135892in}}{\pgfqpoint{10.110561in}{1.129696in}}%
\pgfpathquadraticcurveto{\pgfqpoint{10.105914in}{1.123517in}}{\pgfqpoint{10.105914in}{1.112260in}}%
\pgfpathquadraticcurveto{\pgfqpoint{10.105914in}{1.101056in}}{\pgfqpoint{10.110561in}{1.094860in}}%
\pgfpathquadraticcurveto{\pgfqpoint{10.115208in}{1.088664in}}{\pgfqpoint{10.123530in}{1.088664in}}%
\pgfpathquadraticcurveto{\pgfqpoint{10.131923in}{1.088664in}}{\pgfqpoint{10.136552in}{1.094860in}}%
\pgfpathquadraticcurveto{\pgfqpoint{10.141200in}{1.101056in}}{\pgfqpoint{10.141200in}{1.112260in}}%
\pgfpathlineto{\pgfqpoint{10.141200in}{1.112260in}}%
\pgfpathclose%
\pgfpathmoveto{\pgfqpoint{10.151557in}{1.087817in}}%
\pgfpathquadraticcurveto{\pgfqpoint{10.151557in}{1.071732in}}{\pgfqpoint{10.144406in}{1.063879in}}%
\pgfpathquadraticcurveto{\pgfqpoint{10.137273in}{1.056026in}}{\pgfqpoint{10.122521in}{1.056026in}}%
\pgfpathquadraticcurveto{\pgfqpoint{10.117063in}{1.056026in}}{\pgfqpoint{10.112218in}{1.056836in}}%
\pgfpathquadraticcurveto{\pgfqpoint{10.107373in}{1.057647in}}{\pgfqpoint{10.102816in}{1.059340in}}%
\pgfpathlineto{\pgfqpoint{10.102816in}{1.069409in}}%
\pgfpathquadraticcurveto{\pgfqpoint{10.107373in}{1.066941in}}{\pgfqpoint{10.111822in}{1.065770in}}%
\pgfpathquadraticcurveto{\pgfqpoint{10.116271in}{1.064582in}}{\pgfqpoint{10.120882in}{1.064582in}}%
\pgfpathquadraticcurveto{\pgfqpoint{10.131077in}{1.064582in}}{\pgfqpoint{10.136138in}{1.069895in}}%
\pgfpathquadraticcurveto{\pgfqpoint{10.141200in}{1.075209in}}{\pgfqpoint{10.141200in}{1.085962in}}%
\pgfpathlineto{\pgfqpoint{10.141200in}{1.091095in}}%
\pgfpathquadraticcurveto{\pgfqpoint{10.137993in}{1.085512in}}{\pgfqpoint{10.132986in}{1.082756in}}%
\pgfpathquadraticcurveto{\pgfqpoint{10.127979in}{1.080000in}}{\pgfqpoint{10.120990in}{1.080000in}}%
\pgfpathquadraticcurveto{\pgfqpoint{10.109408in}{1.080000in}}{\pgfqpoint{10.102311in}{1.088826in}}%
\pgfpathquadraticcurveto{\pgfqpoint{10.095215in}{1.097670in}}{\pgfqpoint{10.095215in}{1.112260in}}%
\pgfpathquadraticcurveto{\pgfqpoint{10.095215in}{1.126886in}}{\pgfqpoint{10.102311in}{1.135712in}}%
\pgfpathquadraticcurveto{\pgfqpoint{10.109408in}{1.144556in}}{\pgfqpoint{10.120990in}{1.144556in}}%
\pgfpathquadraticcurveto{\pgfqpoint{10.127979in}{1.144556in}}{\pgfqpoint{10.132986in}{1.141800in}}%
\pgfpathquadraticcurveto{\pgfqpoint{10.137993in}{1.139044in}}{\pgfqpoint{10.141200in}{1.133478in}}%
\pgfpathlineto{\pgfqpoint{10.141200in}{1.143043in}}%
\pgfpathlineto{\pgfqpoint{10.151557in}{1.143043in}}%
\pgfpathlineto{\pgfqpoint{10.151557in}{1.087817in}}%
\pgfpathlineto{\pgfqpoint{10.151557in}{1.087817in}}%
\pgfpathclose%
\pgfpathmoveto{\pgfqpoint{10.201558in}{1.111683in}}%
\pgfpathquadraticcurveto{\pgfqpoint{10.189004in}{1.111683in}}{\pgfqpoint{10.184159in}{1.108819in}}%
\pgfpathquadraticcurveto{\pgfqpoint{10.179313in}{1.105956in}}{\pgfqpoint{10.179313in}{1.099021in}}%
\pgfpathquadraticcurveto{\pgfqpoint{10.179313in}{1.093509in}}{\pgfqpoint{10.182952in}{1.090267in}}%
\pgfpathquadraticcurveto{\pgfqpoint{10.186590in}{1.087043in}}{\pgfqpoint{10.192822in}{1.087043in}}%
\pgfpathquadraticcurveto{\pgfqpoint{10.201450in}{1.087043in}}{\pgfqpoint{10.206656in}{1.093149in}}%
\pgfpathquadraticcurveto{\pgfqpoint{10.211861in}{1.099255in}}{\pgfqpoint{10.211861in}{1.109378in}}%
\pgfpathlineto{\pgfqpoint{10.211861in}{1.111683in}}%
\pgfpathlineto{\pgfqpoint{10.201558in}{1.111683in}}%
\pgfpathlineto{\pgfqpoint{10.201558in}{1.111683in}}%
\pgfpathclose%
\pgfpathmoveto{\pgfqpoint{10.222218in}{1.115970in}}%
\pgfpathlineto{\pgfqpoint{10.222218in}{1.080000in}}%
\pgfpathlineto{\pgfqpoint{10.211861in}{1.080000in}}%
\pgfpathlineto{\pgfqpoint{10.211861in}{1.089564in}}%
\pgfpathquadraticcurveto{\pgfqpoint{10.208313in}{1.083837in}}{\pgfqpoint{10.203017in}{1.081099in}}%
\pgfpathquadraticcurveto{\pgfqpoint{10.197722in}{1.078361in}}{\pgfqpoint{10.190067in}{1.078361in}}%
\pgfpathquadraticcurveto{\pgfqpoint{10.180394in}{1.078361in}}{\pgfqpoint{10.174666in}{1.083801in}}%
\pgfpathquadraticcurveto{\pgfqpoint{10.168956in}{1.089240in}}{\pgfqpoint{10.168956in}{1.098354in}}%
\pgfpathquadraticcurveto{\pgfqpoint{10.168956in}{1.108982in}}{\pgfqpoint{10.176071in}{1.114385in}}%
\pgfpathquadraticcurveto{\pgfqpoint{10.183204in}{1.119789in}}{\pgfqpoint{10.197325in}{1.119789in}}%
\pgfpathlineto{\pgfqpoint{10.211861in}{1.119789in}}%
\pgfpathlineto{\pgfqpoint{10.211861in}{1.120816in}}%
\pgfpathquadraticcurveto{\pgfqpoint{10.211861in}{1.127966in}}{\pgfqpoint{10.207160in}{1.131875in}}%
\pgfpathquadraticcurveto{\pgfqpoint{10.202459in}{1.135784in}}{\pgfqpoint{10.193957in}{1.135784in}}%
\pgfpathquadraticcurveto{\pgfqpoint{10.188554in}{1.135784in}}{\pgfqpoint{10.183420in}{1.134487in}}%
\pgfpathquadraticcurveto{\pgfqpoint{10.178305in}{1.133190in}}{\pgfqpoint{10.173585in}{1.130596in}}%
\pgfpathlineto{\pgfqpoint{10.173585in}{1.140179in}}%
\pgfpathquadraticcurveto{\pgfqpoint{10.179259in}{1.142376in}}{\pgfqpoint{10.184609in}{1.143457in}}%
\pgfpathquadraticcurveto{\pgfqpoint{10.189958in}{1.144556in}}{\pgfqpoint{10.195020in}{1.144556in}}%
\pgfpathquadraticcurveto{\pgfqpoint{10.208709in}{1.144556in}}{\pgfqpoint{10.215464in}{1.137459in}}%
\pgfpathquadraticcurveto{\pgfqpoint{10.222218in}{1.130380in}}{\pgfqpoint{10.222218in}{1.115970in}}%
\pgfpathlineto{\pgfqpoint{10.222218in}{1.115970in}}%
\pgfpathclose%
\pgfpathmoveto{\pgfqpoint{10.253559in}{1.089456in}}%
\pgfpathlineto{\pgfqpoint{10.253559in}{1.056026in}}%
\pgfpathlineto{\pgfqpoint{10.243148in}{1.056026in}}%
\pgfpathlineto{\pgfqpoint{10.243148in}{1.143043in}}%
\pgfpathlineto{\pgfqpoint{10.253559in}{1.143043in}}%
\pgfpathlineto{\pgfqpoint{10.253559in}{1.133478in}}%
\pgfpathquadraticcurveto{\pgfqpoint{10.256838in}{1.139098in}}{\pgfqpoint{10.261809in}{1.141818in}}%
\pgfpathquadraticcurveto{\pgfqpoint{10.266798in}{1.144556in}}{\pgfqpoint{10.273715in}{1.144556in}}%
\pgfpathquadraticcurveto{\pgfqpoint{10.285207in}{1.144556in}}{\pgfqpoint{10.292376in}{1.135441in}}%
\pgfpathquadraticcurveto{\pgfqpoint{10.299562in}{1.126327in}}{\pgfqpoint{10.299562in}{1.111467in}}%
\pgfpathquadraticcurveto{\pgfqpoint{10.299562in}{1.096607in}}{\pgfqpoint{10.292376in}{1.087475in}}%
\pgfpathquadraticcurveto{\pgfqpoint{10.285207in}{1.078361in}}{\pgfqpoint{10.273715in}{1.078361in}}%
\pgfpathquadraticcurveto{\pgfqpoint{10.266798in}{1.078361in}}{\pgfqpoint{10.261809in}{1.081099in}}%
\pgfpathquadraticcurveto{\pgfqpoint{10.256838in}{1.083837in}}{\pgfqpoint{10.253559in}{1.089456in}}%
\pgfpathlineto{\pgfqpoint{10.253559in}{1.089456in}}%
\pgfpathclose%
\pgfpathmoveto{\pgfqpoint{10.288809in}{1.111467in}}%
\pgfpathquadraticcurveto{\pgfqpoint{10.288809in}{1.122887in}}{\pgfqpoint{10.284108in}{1.129389in}}%
\pgfpathquadraticcurveto{\pgfqpoint{10.279407in}{1.135892in}}{\pgfqpoint{10.271193in}{1.135892in}}%
\pgfpathquadraticcurveto{\pgfqpoint{10.262962in}{1.135892in}}{\pgfqpoint{10.258261in}{1.129389in}}%
\pgfpathquadraticcurveto{\pgfqpoint{10.253559in}{1.122887in}}{\pgfqpoint{10.253559in}{1.111467in}}%
\pgfpathquadraticcurveto{\pgfqpoint{10.253559in}{1.100048in}}{\pgfqpoint{10.258261in}{1.093545in}}%
\pgfpathquadraticcurveto{\pgfqpoint{10.262962in}{1.087043in}}{\pgfqpoint{10.271193in}{1.087043in}}%
\pgfpathquadraticcurveto{\pgfqpoint{10.279407in}{1.087043in}}{\pgfqpoint{10.284108in}{1.093545in}}%
\pgfpathquadraticcurveto{\pgfqpoint{10.288809in}{1.100048in}}{\pgfqpoint{10.288809in}{1.111467in}}%
\pgfpathlineto{\pgfqpoint{10.288809in}{1.111467in}}%
\pgfpathclose%
\pgfpathmoveto{\pgfqpoint{10.356913in}{1.141187in}}%
\pgfpathlineto{\pgfqpoint{10.356913in}{1.131389in}}%
\pgfpathquadraticcurveto{\pgfqpoint{10.352536in}{1.133640in}}{\pgfqpoint{10.347799in}{1.134757in}}%
\pgfpathquadraticcurveto{\pgfqpoint{10.343080in}{1.135892in}}{\pgfqpoint{10.338000in}{1.135892in}}%
\pgfpathquadraticcurveto{\pgfqpoint{10.330291in}{1.135892in}}{\pgfqpoint{10.326437in}{1.133532in}}%
\pgfpathquadraticcurveto{\pgfqpoint{10.322582in}{1.131173in}}{\pgfqpoint{10.322582in}{1.126435in}}%
\pgfpathquadraticcurveto{\pgfqpoint{10.322582in}{1.122833in}}{\pgfqpoint{10.325338in}{1.120780in}}%
\pgfpathquadraticcurveto{\pgfqpoint{10.328094in}{1.118726in}}{\pgfqpoint{10.336433in}{1.116871in}}%
\pgfpathlineto{\pgfqpoint{10.339982in}{1.116078in}}%
\pgfpathquadraticcurveto{\pgfqpoint{10.351005in}{1.113719in}}{\pgfqpoint{10.355652in}{1.109414in}}%
\pgfpathquadraticcurveto{\pgfqpoint{10.360299in}{1.105109in}}{\pgfqpoint{10.360299in}{1.097400in}}%
\pgfpathquadraticcurveto{\pgfqpoint{10.360299in}{1.088610in}}{\pgfqpoint{10.353347in}{1.083476in}}%
\pgfpathquadraticcurveto{\pgfqpoint{10.346394in}{1.078361in}}{\pgfqpoint{10.334236in}{1.078361in}}%
\pgfpathquadraticcurveto{\pgfqpoint{10.329174in}{1.078361in}}{\pgfqpoint{10.323681in}{1.079352in}}%
\pgfpathquadraticcurveto{\pgfqpoint{10.318187in}{1.080342in}}{\pgfqpoint{10.312117in}{1.082306in}}%
\pgfpathlineto{\pgfqpoint{10.312117in}{1.093005in}}%
\pgfpathquadraticcurveto{\pgfqpoint{10.317863in}{1.090015in}}{\pgfqpoint{10.323429in}{1.088520in}}%
\pgfpathquadraticcurveto{\pgfqpoint{10.328994in}{1.087043in}}{\pgfqpoint{10.334470in}{1.087043in}}%
\pgfpathquadraticcurveto{\pgfqpoint{10.341783in}{1.087043in}}{\pgfqpoint{10.345710in}{1.089546in}}%
\pgfpathquadraticcurveto{\pgfqpoint{10.349654in}{1.092050in}}{\pgfqpoint{10.349654in}{1.096607in}}%
\pgfpathquadraticcurveto{\pgfqpoint{10.349654in}{1.100822in}}{\pgfqpoint{10.346808in}{1.103074in}}%
\pgfpathquadraticcurveto{\pgfqpoint{10.343980in}{1.105325in}}{\pgfqpoint{10.334344in}{1.107414in}}%
\pgfpathlineto{\pgfqpoint{10.330741in}{1.108261in}}%
\pgfpathquadraticcurveto{\pgfqpoint{10.321123in}{1.110278in}}{\pgfqpoint{10.316836in}{1.114475in}}%
\pgfpathquadraticcurveto{\pgfqpoint{10.312567in}{1.118672in}}{\pgfqpoint{10.312567in}{1.125985in}}%
\pgfpathquadraticcurveto{\pgfqpoint{10.312567in}{1.134883in}}{\pgfqpoint{10.318871in}{1.139710in}}%
\pgfpathquadraticcurveto{\pgfqpoint{10.325176in}{1.144556in}}{\pgfqpoint{10.336776in}{1.144556in}}%
\pgfpathquadraticcurveto{\pgfqpoint{10.342503in}{1.144556in}}{\pgfqpoint{10.347565in}{1.143709in}}%
\pgfpathquadraticcurveto{\pgfqpoint{10.352644in}{1.142880in}}{\pgfqpoint{10.356913in}{1.141187in}}%
\pgfpathlineto{\pgfqpoint{10.356913in}{1.141187in}}%
\pgfpathclose%
\pgfpathmoveto{\pgfqpoint{10.378239in}{1.094302in}}%
\pgfpathlineto{\pgfqpoint{10.390128in}{1.094302in}}%
\pgfpathlineto{\pgfqpoint{10.390128in}{1.080000in}}%
\pgfpathlineto{\pgfqpoint{10.378239in}{1.080000in}}%
\pgfpathlineto{\pgfqpoint{10.378239in}{1.094302in}}%
\pgfpathlineto{\pgfqpoint{10.378239in}{1.094302in}}%
\pgfpathclose%
\pgfusepath{fill}%
\end{pgfscope}%
\begin{pgfscope}%
\pgfsetbuttcap%
\pgfsetroundjoin%
\definecolor{currentfill}{rgb}{0.309804,0.372549,0.419608}%
\pgfsetfillcolor{currentfill}%
\pgfsetlinewidth{0.000000pt}%
\definecolor{currentstroke}{rgb}{0.309804,0.372549,0.419608}%
\pgfsetstrokecolor{currentstroke}%
\pgfsetdash{}{0pt}%
\pgfpathmoveto{\pgfqpoint{10.472065in}{1.164045in}}%
\pgfpathlineto{\pgfqpoint{10.525201in}{1.164045in}}%
\pgfpathlineto{\pgfqpoint{10.525201in}{1.154462in}}%
\pgfpathlineto{\pgfqpoint{10.483430in}{1.154462in}}%
\pgfpathlineto{\pgfqpoint{10.483430in}{1.129587in}}%
\pgfpathlineto{\pgfqpoint{10.523453in}{1.129587in}}%
\pgfpathlineto{\pgfqpoint{10.523453in}{1.120023in}}%
\pgfpathlineto{\pgfqpoint{10.483430in}{1.120023in}}%
\pgfpathlineto{\pgfqpoint{10.483430in}{1.089564in}}%
\pgfpathlineto{\pgfqpoint{10.526209in}{1.089564in}}%
\pgfpathlineto{\pgfqpoint{10.526209in}{1.080000in}}%
\pgfpathlineto{\pgfqpoint{10.472065in}{1.080000in}}%
\pgfpathlineto{\pgfqpoint{10.472065in}{1.164045in}}%
\pgfpathlineto{\pgfqpoint{10.472065in}{1.164045in}}%
\pgfpathclose%
\pgfpathmoveto{\pgfqpoint{10.596871in}{1.143043in}}%
\pgfpathlineto{\pgfqpoint{10.574068in}{1.112368in}}%
\pgfpathlineto{\pgfqpoint{10.598042in}{1.080000in}}%
\pgfpathlineto{\pgfqpoint{10.585830in}{1.080000in}}%
\pgfpathlineto{\pgfqpoint{10.567475in}{1.104767in}}%
\pgfpathlineto{\pgfqpoint{10.549139in}{1.080000in}}%
\pgfpathlineto{\pgfqpoint{10.536909in}{1.080000in}}%
\pgfpathlineto{\pgfqpoint{10.561405in}{1.112980in}}%
\pgfpathlineto{\pgfqpoint{10.538998in}{1.143043in}}%
\pgfpathlineto{\pgfqpoint{10.551210in}{1.143043in}}%
\pgfpathlineto{\pgfqpoint{10.567925in}{1.120581in}}%
\pgfpathlineto{\pgfqpoint{10.584641in}{1.143043in}}%
\pgfpathlineto{\pgfqpoint{10.596871in}{1.143043in}}%
\pgfpathlineto{\pgfqpoint{10.596871in}{1.143043in}}%
\pgfpathclose%
\pgfpathmoveto{\pgfqpoint{10.656041in}{1.140629in}}%
\pgfpathlineto{\pgfqpoint{10.656041in}{1.130938in}}%
\pgfpathquadraticcurveto{\pgfqpoint{10.651646in}{1.133370in}}{\pgfqpoint{10.647233in}{1.134577in}}%
\pgfpathquadraticcurveto{\pgfqpoint{10.642820in}{1.135784in}}{\pgfqpoint{10.638317in}{1.135784in}}%
\pgfpathquadraticcurveto{\pgfqpoint{10.628230in}{1.135784in}}{\pgfqpoint{10.622646in}{1.129389in}}%
\pgfpathquadraticcurveto{\pgfqpoint{10.617081in}{1.123013in}}{\pgfqpoint{10.617081in}{1.111467in}}%
\pgfpathquadraticcurveto{\pgfqpoint{10.617081in}{1.099921in}}{\pgfqpoint{10.622646in}{1.093527in}}%
\pgfpathquadraticcurveto{\pgfqpoint{10.628230in}{1.087151in}}{\pgfqpoint{10.638317in}{1.087151in}}%
\pgfpathquadraticcurveto{\pgfqpoint{10.642820in}{1.087151in}}{\pgfqpoint{10.647233in}{1.088358in}}%
\pgfpathquadraticcurveto{\pgfqpoint{10.651646in}{1.089564in}}{\pgfqpoint{10.656041in}{1.091996in}}%
\pgfpathlineto{\pgfqpoint{10.656041in}{1.082414in}}%
\pgfpathquadraticcurveto{\pgfqpoint{10.651700in}{1.080396in}}{\pgfqpoint{10.647053in}{1.079388in}}%
\pgfpathquadraticcurveto{\pgfqpoint{10.642424in}{1.078361in}}{\pgfqpoint{10.637182in}{1.078361in}}%
\pgfpathquadraticcurveto{\pgfqpoint{10.622935in}{1.078361in}}{\pgfqpoint{10.614541in}{1.087313in}}%
\pgfpathquadraticcurveto{\pgfqpoint{10.606165in}{1.096265in}}{\pgfqpoint{10.606165in}{1.111467in}}%
\pgfpathquadraticcurveto{\pgfqpoint{10.606165in}{1.126886in}}{\pgfqpoint{10.614631in}{1.135712in}}%
\pgfpathquadraticcurveto{\pgfqpoint{10.623115in}{1.144556in}}{\pgfqpoint{10.637867in}{1.144556in}}%
\pgfpathquadraticcurveto{\pgfqpoint{10.642640in}{1.144556in}}{\pgfqpoint{10.647197in}{1.143565in}}%
\pgfpathquadraticcurveto{\pgfqpoint{10.651754in}{1.142592in}}{\pgfqpoint{10.656041in}{1.140629in}}%
\pgfpathlineto{\pgfqpoint{10.656041in}{1.140629in}}%
\pgfpathclose%
\pgfpathmoveto{\pgfqpoint{10.674053in}{1.167593in}}%
\pgfpathlineto{\pgfqpoint{10.684410in}{1.167593in}}%
\pgfpathlineto{\pgfqpoint{10.684410in}{1.080000in}}%
\pgfpathlineto{\pgfqpoint{10.674053in}{1.080000in}}%
\pgfpathlineto{\pgfqpoint{10.674053in}{1.167593in}}%
\pgfpathlineto{\pgfqpoint{10.674053in}{1.167593in}}%
\pgfpathclose%
\pgfpathmoveto{\pgfqpoint{10.705016in}{1.104875in}}%
\pgfpathlineto{\pgfqpoint{10.705016in}{1.143043in}}%
\pgfpathlineto{\pgfqpoint{10.715373in}{1.143043in}}%
\pgfpathlineto{\pgfqpoint{10.715373in}{1.105271in}}%
\pgfpathquadraticcurveto{\pgfqpoint{10.715373in}{1.096319in}}{\pgfqpoint{10.718849in}{1.091834in}}%
\pgfpathquadraticcurveto{\pgfqpoint{10.722344in}{1.087367in}}{\pgfqpoint{10.729332in}{1.087367in}}%
\pgfpathquadraticcurveto{\pgfqpoint{10.737708in}{1.087367in}}{\pgfqpoint{10.742571in}{1.092717in}}%
\pgfpathquadraticcurveto{\pgfqpoint{10.747453in}{1.098066in}}{\pgfqpoint{10.747453in}{1.107306in}}%
\pgfpathlineto{\pgfqpoint{10.747453in}{1.143043in}}%
\pgfpathlineto{\pgfqpoint{10.757810in}{1.143043in}}%
\pgfpathlineto{\pgfqpoint{10.757810in}{1.080000in}}%
\pgfpathlineto{\pgfqpoint{10.747453in}{1.080000in}}%
\pgfpathlineto{\pgfqpoint{10.747453in}{1.089691in}}%
\pgfpathquadraticcurveto{\pgfqpoint{10.743688in}{1.083945in}}{\pgfqpoint{10.738699in}{1.081153in}}%
\pgfpathquadraticcurveto{\pgfqpoint{10.733727in}{1.078361in}}{\pgfqpoint{10.727135in}{1.078361in}}%
\pgfpathquadraticcurveto{\pgfqpoint{10.716274in}{1.078361in}}{\pgfqpoint{10.710636in}{1.085115in}}%
\pgfpathquadraticcurveto{\pgfqpoint{10.705016in}{1.091870in}}{\pgfqpoint{10.705016in}{1.104875in}}%
\pgfpathlineto{\pgfqpoint{10.705016in}{1.104875in}}%
\pgfpathclose%
\pgfpathmoveto{\pgfqpoint{10.731080in}{1.144556in}}%
\pgfpathlineto{\pgfqpoint{10.731080in}{1.144556in}}%
\pgfpathlineto{\pgfqpoint{10.731080in}{1.144556in}}%
\pgfpathclose%
\pgfpathmoveto{\pgfqpoint{10.820618in}{1.133478in}}%
\pgfpathlineto{\pgfqpoint{10.820618in}{1.167593in}}%
\pgfpathlineto{\pgfqpoint{10.830975in}{1.167593in}}%
\pgfpathlineto{\pgfqpoint{10.830975in}{1.080000in}}%
\pgfpathlineto{\pgfqpoint{10.820618in}{1.080000in}}%
\pgfpathlineto{\pgfqpoint{10.820618in}{1.089456in}}%
\pgfpathquadraticcurveto{\pgfqpoint{10.817358in}{1.083837in}}{\pgfqpoint{10.812368in}{1.081099in}}%
\pgfpathquadraticcurveto{\pgfqpoint{10.807397in}{1.078361in}}{\pgfqpoint{10.800408in}{1.078361in}}%
\pgfpathquadraticcurveto{\pgfqpoint{10.788989in}{1.078361in}}{\pgfqpoint{10.781802in}{1.087475in}}%
\pgfpathquadraticcurveto{\pgfqpoint{10.774633in}{1.096607in}}{\pgfqpoint{10.774633in}{1.111467in}}%
\pgfpathquadraticcurveto{\pgfqpoint{10.774633in}{1.126327in}}{\pgfqpoint{10.781802in}{1.135441in}}%
\pgfpathquadraticcurveto{\pgfqpoint{10.788989in}{1.144556in}}{\pgfqpoint{10.800408in}{1.144556in}}%
\pgfpathquadraticcurveto{\pgfqpoint{10.807397in}{1.144556in}}{\pgfqpoint{10.812368in}{1.141818in}}%
\pgfpathquadraticcurveto{\pgfqpoint{10.817358in}{1.139098in}}{\pgfqpoint{10.820618in}{1.133478in}}%
\pgfpathlineto{\pgfqpoint{10.820618in}{1.133478in}}%
\pgfpathclose%
\pgfpathmoveto{\pgfqpoint{10.785332in}{1.111467in}}%
\pgfpathquadraticcurveto{\pgfqpoint{10.785332in}{1.100048in}}{\pgfqpoint{10.790033in}{1.093545in}}%
\pgfpathquadraticcurveto{\pgfqpoint{10.794734in}{1.087043in}}{\pgfqpoint{10.802948in}{1.087043in}}%
\pgfpathquadraticcurveto{\pgfqpoint{10.811162in}{1.087043in}}{\pgfqpoint{10.815881in}{1.093545in}}%
\pgfpathquadraticcurveto{\pgfqpoint{10.820618in}{1.100048in}}{\pgfqpoint{10.820618in}{1.111467in}}%
\pgfpathquadraticcurveto{\pgfqpoint{10.820618in}{1.122887in}}{\pgfqpoint{10.815881in}{1.129389in}}%
\pgfpathquadraticcurveto{\pgfqpoint{10.811162in}{1.135892in}}{\pgfqpoint{10.802948in}{1.135892in}}%
\pgfpathquadraticcurveto{\pgfqpoint{10.794734in}{1.135892in}}{\pgfqpoint{10.790033in}{1.129389in}}%
\pgfpathquadraticcurveto{\pgfqpoint{10.785332in}{1.122887in}}{\pgfqpoint{10.785332in}{1.111467in}}%
\pgfpathlineto{\pgfqpoint{10.785332in}{1.111467in}}%
\pgfpathclose%
\pgfpathmoveto{\pgfqpoint{10.906248in}{1.114115in}}%
\pgfpathlineto{\pgfqpoint{10.906248in}{1.109054in}}%
\pgfpathlineto{\pgfqpoint{10.858624in}{1.109054in}}%
\pgfpathquadraticcurveto{\pgfqpoint{10.859308in}{1.098354in}}{\pgfqpoint{10.865072in}{1.092753in}}%
\pgfpathquadraticcurveto{\pgfqpoint{10.870836in}{1.087151in}}{\pgfqpoint{10.881139in}{1.087151in}}%
\pgfpathquadraticcurveto{\pgfqpoint{10.887101in}{1.087151in}}{\pgfqpoint{10.892703in}{1.088610in}}%
\pgfpathquadraticcurveto{\pgfqpoint{10.898304in}{1.090069in}}{\pgfqpoint{10.903834in}{1.093005in}}%
\pgfpathlineto{\pgfqpoint{10.903834in}{1.083206in}}%
\pgfpathquadraticcurveto{\pgfqpoint{10.898250in}{1.080847in}}{\pgfqpoint{10.892396in}{1.079604in}}%
\pgfpathquadraticcurveto{\pgfqpoint{10.886542in}{1.078361in}}{\pgfqpoint{10.880526in}{1.078361in}}%
\pgfpathquadraticcurveto{\pgfqpoint{10.865432in}{1.078361in}}{\pgfqpoint{10.856624in}{1.087133in}}%
\pgfpathquadraticcurveto{\pgfqpoint{10.847816in}{1.095923in}}{\pgfqpoint{10.847816in}{1.110909in}}%
\pgfpathquadraticcurveto{\pgfqpoint{10.847816in}{1.126381in}}{\pgfqpoint{10.856174in}{1.135459in}}%
\pgfpathquadraticcurveto{\pgfqpoint{10.864532in}{1.144556in}}{\pgfqpoint{10.878725in}{1.144556in}}%
\pgfpathquadraticcurveto{\pgfqpoint{10.891442in}{1.144556in}}{\pgfqpoint{10.898845in}{1.136360in}}%
\pgfpathquadraticcurveto{\pgfqpoint{10.906248in}{1.128183in}}{\pgfqpoint{10.906248in}{1.114115in}}%
\pgfpathlineto{\pgfqpoint{10.906248in}{1.114115in}}%
\pgfpathclose%
\pgfpathmoveto{\pgfqpoint{10.895891in}{1.117159in}}%
\pgfpathquadraticcurveto{\pgfqpoint{10.895783in}{1.125643in}}{\pgfqpoint{10.891136in}{1.130704in}}%
\pgfpathquadraticcurveto{\pgfqpoint{10.886488in}{1.135784in}}{\pgfqpoint{10.878833in}{1.135784in}}%
\pgfpathquadraticcurveto{\pgfqpoint{10.870169in}{1.135784in}}{\pgfqpoint{10.864964in}{1.130884in}}%
\pgfpathquadraticcurveto{\pgfqpoint{10.859758in}{1.125985in}}{\pgfqpoint{10.858966in}{1.117087in}}%
\pgfpathlineto{\pgfqpoint{10.895891in}{1.117159in}}%
\pgfpathlineto{\pgfqpoint{10.895891in}{1.117159in}}%
\pgfpathclose%
\pgfpathmoveto{\pgfqpoint{10.964733in}{1.133478in}}%
\pgfpathlineto{\pgfqpoint{10.964733in}{1.167593in}}%
\pgfpathlineto{\pgfqpoint{10.975090in}{1.167593in}}%
\pgfpathlineto{\pgfqpoint{10.975090in}{1.080000in}}%
\pgfpathlineto{\pgfqpoint{10.964733in}{1.080000in}}%
\pgfpathlineto{\pgfqpoint{10.964733in}{1.089456in}}%
\pgfpathquadraticcurveto{\pgfqpoint{10.961473in}{1.083837in}}{\pgfqpoint{10.956484in}{1.081099in}}%
\pgfpathquadraticcurveto{\pgfqpoint{10.951512in}{1.078361in}}{\pgfqpoint{10.944524in}{1.078361in}}%
\pgfpathquadraticcurveto{\pgfqpoint{10.933104in}{1.078361in}}{\pgfqpoint{10.925917in}{1.087475in}}%
\pgfpathquadraticcurveto{\pgfqpoint{10.918748in}{1.096607in}}{\pgfqpoint{10.918748in}{1.111467in}}%
\pgfpathquadraticcurveto{\pgfqpoint{10.918748in}{1.126327in}}{\pgfqpoint{10.925917in}{1.135441in}}%
\pgfpathquadraticcurveto{\pgfqpoint{10.933104in}{1.144556in}}{\pgfqpoint{10.944524in}{1.144556in}}%
\pgfpathquadraticcurveto{\pgfqpoint{10.951512in}{1.144556in}}{\pgfqpoint{10.956484in}{1.141818in}}%
\pgfpathquadraticcurveto{\pgfqpoint{10.961473in}{1.139098in}}{\pgfqpoint{10.964733in}{1.133478in}}%
\pgfpathlineto{\pgfqpoint{10.964733in}{1.133478in}}%
\pgfpathclose%
\pgfpathmoveto{\pgfqpoint{10.929447in}{1.111467in}}%
\pgfpathquadraticcurveto{\pgfqpoint{10.929447in}{1.100048in}}{\pgfqpoint{10.934149in}{1.093545in}}%
\pgfpathquadraticcurveto{\pgfqpoint{10.938850in}{1.087043in}}{\pgfqpoint{10.947063in}{1.087043in}}%
\pgfpathquadraticcurveto{\pgfqpoint{10.955277in}{1.087043in}}{\pgfqpoint{10.959996in}{1.093545in}}%
\pgfpathquadraticcurveto{\pgfqpoint{10.964733in}{1.100048in}}{\pgfqpoint{10.964733in}{1.111467in}}%
\pgfpathquadraticcurveto{\pgfqpoint{10.964733in}{1.122887in}}{\pgfqpoint{10.959996in}{1.129389in}}%
\pgfpathquadraticcurveto{\pgfqpoint{10.955277in}{1.135892in}}{\pgfqpoint{10.947063in}{1.135892in}}%
\pgfpathquadraticcurveto{\pgfqpoint{10.938850in}{1.135892in}}{\pgfqpoint{10.934149in}{1.129389in}}%
\pgfpathquadraticcurveto{\pgfqpoint{10.929447in}{1.122887in}}{\pgfqpoint{10.929447in}{1.111467in}}%
\pgfpathlineto{\pgfqpoint{10.929447in}{1.111467in}}%
\pgfpathclose%
\pgfusepath{fill}%
\end{pgfscope}%
\begin{pgfscope}%
\pgfsetbuttcap%
\pgfsetroundjoin%
\definecolor{currentfill}{rgb}{0.309804,0.372549,0.419608}%
\pgfsetfillcolor{currentfill}%
\pgfsetlinewidth{0.000000pt}%
\definecolor{currentstroke}{rgb}{0.309804,0.372549,0.419608}%
\pgfsetstrokecolor{currentstroke}%
\pgfsetdash{}{0pt}%
\pgfpathmoveto{\pgfqpoint{11.071480in}{1.143043in}}%
\pgfpathlineto{\pgfqpoint{11.081837in}{1.143043in}}%
\pgfpathlineto{\pgfqpoint{11.081837in}{1.080000in}}%
\pgfpathlineto{\pgfqpoint{11.071480in}{1.080000in}}%
\pgfpathlineto{\pgfqpoint{11.071480in}{1.143043in}}%
\pgfpathlineto{\pgfqpoint{11.071480in}{1.143043in}}%
\pgfpathclose%
\pgfpathmoveto{\pgfqpoint{11.071480in}{1.167593in}}%
\pgfpathlineto{\pgfqpoint{11.081837in}{1.167593in}}%
\pgfpathlineto{\pgfqpoint{11.081837in}{1.154462in}}%
\pgfpathlineto{\pgfqpoint{11.071480in}{1.154462in}}%
\pgfpathlineto{\pgfqpoint{11.071480in}{1.167593in}}%
\pgfpathlineto{\pgfqpoint{11.071480in}{1.167593in}}%
\pgfpathclose%
\pgfpathmoveto{\pgfqpoint{11.135423in}{1.167593in}}%
\pgfpathlineto{\pgfqpoint{11.135423in}{1.158965in}}%
\pgfpathlineto{\pgfqpoint{11.125517in}{1.158965in}}%
\pgfpathquadraticcurveto{\pgfqpoint{11.119951in}{1.158965in}}{\pgfqpoint{11.117771in}{1.156714in}}%
\pgfpathquadraticcurveto{\pgfqpoint{11.115610in}{1.154462in}}{\pgfqpoint{11.115610in}{1.148608in}}%
\pgfpathlineto{\pgfqpoint{11.115610in}{1.143043in}}%
\pgfpathlineto{\pgfqpoint{11.132667in}{1.143043in}}%
\pgfpathlineto{\pgfqpoint{11.132667in}{1.134991in}}%
\pgfpathlineto{\pgfqpoint{11.115610in}{1.134991in}}%
\pgfpathlineto{\pgfqpoint{11.115610in}{1.080000in}}%
\pgfpathlineto{\pgfqpoint{11.105199in}{1.080000in}}%
\pgfpathlineto{\pgfqpoint{11.105199in}{1.134991in}}%
\pgfpathlineto{\pgfqpoint{11.095292in}{1.134991in}}%
\pgfpathlineto{\pgfqpoint{11.095292in}{1.143043in}}%
\pgfpathlineto{\pgfqpoint{11.105199in}{1.143043in}}%
\pgfpathlineto{\pgfqpoint{11.105199in}{1.147438in}}%
\pgfpathquadraticcurveto{\pgfqpoint{11.105199in}{1.157957in}}{\pgfqpoint{11.110098in}{1.162766in}}%
\pgfpathquadraticcurveto{\pgfqpoint{11.114997in}{1.167593in}}{\pgfqpoint{11.125625in}{1.167593in}}%
\pgfpathlineto{\pgfqpoint{11.135423in}{1.167593in}}%
\pgfpathlineto{\pgfqpoint{11.135423in}{1.167593in}}%
\pgfpathclose%
\pgfusepath{fill}%
\end{pgfscope}%
\begin{pgfscope}%
\pgfsetbuttcap%
\pgfsetroundjoin%
\definecolor{currentfill}{rgb}{0.309804,0.372549,0.419608}%
\pgfsetfillcolor{currentfill}%
\pgfsetlinewidth{0.000000pt}%
\definecolor{currentstroke}{rgb}{0.309804,0.372549,0.419608}%
\pgfsetstrokecolor{currentstroke}%
\pgfsetdash{}{0pt}%
\pgfpathmoveto{\pgfqpoint{11.237305in}{1.114115in}}%
\pgfpathlineto{\pgfqpoint{11.237305in}{1.109054in}}%
\pgfpathlineto{\pgfqpoint{11.189681in}{1.109054in}}%
\pgfpathquadraticcurveto{\pgfqpoint{11.190366in}{1.098354in}}{\pgfqpoint{11.196130in}{1.092753in}}%
\pgfpathquadraticcurveto{\pgfqpoint{11.201893in}{1.087151in}}{\pgfqpoint{11.212196in}{1.087151in}}%
\pgfpathquadraticcurveto{\pgfqpoint{11.218158in}{1.087151in}}{\pgfqpoint{11.223760in}{1.088610in}}%
\pgfpathquadraticcurveto{\pgfqpoint{11.229362in}{1.090069in}}{\pgfqpoint{11.234892in}{1.093005in}}%
\pgfpathlineto{\pgfqpoint{11.234892in}{1.083206in}}%
\pgfpathquadraticcurveto{\pgfqpoint{11.229308in}{1.080847in}}{\pgfqpoint{11.223454in}{1.079604in}}%
\pgfpathquadraticcurveto{\pgfqpoint{11.217600in}{1.078361in}}{\pgfqpoint{11.211584in}{1.078361in}}%
\pgfpathquadraticcurveto{\pgfqpoint{11.196490in}{1.078361in}}{\pgfqpoint{11.187682in}{1.087133in}}%
\pgfpathquadraticcurveto{\pgfqpoint{11.178874in}{1.095923in}}{\pgfqpoint{11.178874in}{1.110909in}}%
\pgfpathquadraticcurveto{\pgfqpoint{11.178874in}{1.126381in}}{\pgfqpoint{11.187232in}{1.135459in}}%
\pgfpathquadraticcurveto{\pgfqpoint{11.195589in}{1.144556in}}{\pgfqpoint{11.209783in}{1.144556in}}%
\pgfpathquadraticcurveto{\pgfqpoint{11.222499in}{1.144556in}}{\pgfqpoint{11.229902in}{1.136360in}}%
\pgfpathquadraticcurveto{\pgfqpoint{11.237305in}{1.128183in}}{\pgfqpoint{11.237305in}{1.114115in}}%
\pgfpathlineto{\pgfqpoint{11.237305in}{1.114115in}}%
\pgfpathclose%
\pgfpathmoveto{\pgfqpoint{11.226948in}{1.117159in}}%
\pgfpathquadraticcurveto{\pgfqpoint{11.226840in}{1.125643in}}{\pgfqpoint{11.222193in}{1.130704in}}%
\pgfpathquadraticcurveto{\pgfqpoint{11.217546in}{1.135784in}}{\pgfqpoint{11.209891in}{1.135784in}}%
\pgfpathquadraticcurveto{\pgfqpoint{11.201227in}{1.135784in}}{\pgfqpoint{11.196021in}{1.130884in}}%
\pgfpathquadraticcurveto{\pgfqpoint{11.190816in}{1.125985in}}{\pgfqpoint{11.190023in}{1.117087in}}%
\pgfpathlineto{\pgfqpoint{11.226948in}{1.117159in}}%
\pgfpathlineto{\pgfqpoint{11.226948in}{1.117159in}}%
\pgfpathclose%
\pgfpathmoveto{\pgfqpoint{11.254309in}{1.143043in}}%
\pgfpathlineto{\pgfqpoint{11.264666in}{1.143043in}}%
\pgfpathlineto{\pgfqpoint{11.264666in}{1.080000in}}%
\pgfpathlineto{\pgfqpoint{11.254309in}{1.080000in}}%
\pgfpathlineto{\pgfqpoint{11.254309in}{1.143043in}}%
\pgfpathlineto{\pgfqpoint{11.254309in}{1.143043in}}%
\pgfpathclose%
\pgfpathmoveto{\pgfqpoint{11.254309in}{1.167593in}}%
\pgfpathlineto{\pgfqpoint{11.264666in}{1.167593in}}%
\pgfpathlineto{\pgfqpoint{11.264666in}{1.154462in}}%
\pgfpathlineto{\pgfqpoint{11.254309in}{1.154462in}}%
\pgfpathlineto{\pgfqpoint{11.254309in}{1.167593in}}%
\pgfpathlineto{\pgfqpoint{11.254309in}{1.167593in}}%
\pgfpathclose%
\pgfpathmoveto{\pgfqpoint{11.296583in}{1.160947in}}%
\pgfpathlineto{\pgfqpoint{11.296583in}{1.143043in}}%
\pgfpathlineto{\pgfqpoint{11.317910in}{1.143043in}}%
\pgfpathlineto{\pgfqpoint{11.317910in}{1.134991in}}%
\pgfpathlineto{\pgfqpoint{11.296583in}{1.134991in}}%
\pgfpathlineto{\pgfqpoint{11.296583in}{1.100768in}}%
\pgfpathquadraticcurveto{\pgfqpoint{11.296583in}{1.093059in}}{\pgfqpoint{11.298691in}{1.090861in}}%
\pgfpathquadraticcurveto{\pgfqpoint{11.300798in}{1.088664in}}{\pgfqpoint{11.307283in}{1.088664in}}%
\pgfpathlineto{\pgfqpoint{11.317910in}{1.088664in}}%
\pgfpathlineto{\pgfqpoint{11.317910in}{1.080000in}}%
\pgfpathlineto{\pgfqpoint{11.307283in}{1.080000in}}%
\pgfpathquadraticcurveto{\pgfqpoint{11.295286in}{1.080000in}}{\pgfqpoint{11.290729in}{1.084467in}}%
\pgfpathquadraticcurveto{\pgfqpoint{11.286172in}{1.088952in}}{\pgfqpoint{11.286172in}{1.100768in}}%
\pgfpathlineto{\pgfqpoint{11.286172in}{1.134991in}}%
\pgfpathlineto{\pgfqpoint{11.278571in}{1.134991in}}%
\pgfpathlineto{\pgfqpoint{11.278571in}{1.143043in}}%
\pgfpathlineto{\pgfqpoint{11.286172in}{1.143043in}}%
\pgfpathlineto{\pgfqpoint{11.286172in}{1.160947in}}%
\pgfpathlineto{\pgfqpoint{11.296583in}{1.160947in}}%
\pgfpathlineto{\pgfqpoint{11.296583in}{1.160947in}}%
\pgfpathclose%
\pgfpathmoveto{\pgfqpoint{11.383942in}{1.118060in}}%
\pgfpathlineto{\pgfqpoint{11.383942in}{1.080000in}}%
\pgfpathlineto{\pgfqpoint{11.373585in}{1.080000in}}%
\pgfpathlineto{\pgfqpoint{11.373585in}{1.117717in}}%
\pgfpathquadraticcurveto{\pgfqpoint{11.373585in}{1.126669in}}{\pgfqpoint{11.370091in}{1.131100in}}%
\pgfpathquadraticcurveto{\pgfqpoint{11.366597in}{1.135549in}}{\pgfqpoint{11.359626in}{1.135549in}}%
\pgfpathquadraticcurveto{\pgfqpoint{11.351232in}{1.135549in}}{\pgfqpoint{11.346387in}{1.130200in}}%
\pgfpathquadraticcurveto{\pgfqpoint{11.341542in}{1.124868in}}{\pgfqpoint{11.341542in}{1.115628in}}%
\pgfpathlineto{\pgfqpoint{11.341542in}{1.080000in}}%
\pgfpathlineto{\pgfqpoint{11.331131in}{1.080000in}}%
\pgfpathlineto{\pgfqpoint{11.331131in}{1.167593in}}%
\pgfpathlineto{\pgfqpoint{11.341542in}{1.167593in}}%
\pgfpathlineto{\pgfqpoint{11.341542in}{1.133244in}}%
\pgfpathquadraticcurveto{\pgfqpoint{11.345270in}{1.138936in}}{\pgfqpoint{11.350296in}{1.141746in}}%
\pgfpathquadraticcurveto{\pgfqpoint{11.355339in}{1.144556in}}{\pgfqpoint{11.361931in}{1.144556in}}%
\pgfpathquadraticcurveto{\pgfqpoint{11.372793in}{1.144556in}}{\pgfqpoint{11.378359in}{1.137837in}}%
\pgfpathquadraticcurveto{\pgfqpoint{11.383942in}{1.131118in}}{\pgfqpoint{11.383942in}{1.118060in}}%
\pgfpathlineto{\pgfqpoint{11.383942in}{1.118060in}}%
\pgfpathclose%
\pgfpathmoveto{\pgfqpoint{11.458513in}{1.114115in}}%
\pgfpathlineto{\pgfqpoint{11.458513in}{1.109054in}}%
\pgfpathlineto{\pgfqpoint{11.410888in}{1.109054in}}%
\pgfpathquadraticcurveto{\pgfqpoint{11.411573in}{1.098354in}}{\pgfqpoint{11.417337in}{1.092753in}}%
\pgfpathquadraticcurveto{\pgfqpoint{11.423101in}{1.087151in}}{\pgfqpoint{11.433404in}{1.087151in}}%
\pgfpathquadraticcurveto{\pgfqpoint{11.439366in}{1.087151in}}{\pgfqpoint{11.444967in}{1.088610in}}%
\pgfpathquadraticcurveto{\pgfqpoint{11.450569in}{1.090069in}}{\pgfqpoint{11.456099in}{1.093005in}}%
\pgfpathlineto{\pgfqpoint{11.456099in}{1.083206in}}%
\pgfpathquadraticcurveto{\pgfqpoint{11.450515in}{1.080847in}}{\pgfqpoint{11.444661in}{1.079604in}}%
\pgfpathquadraticcurveto{\pgfqpoint{11.438807in}{1.078361in}}{\pgfqpoint{11.432791in}{1.078361in}}%
\pgfpathquadraticcurveto{\pgfqpoint{11.417697in}{1.078361in}}{\pgfqpoint{11.408889in}{1.087133in}}%
\pgfpathquadraticcurveto{\pgfqpoint{11.400081in}{1.095923in}}{\pgfqpoint{11.400081in}{1.110909in}}%
\pgfpathquadraticcurveto{\pgfqpoint{11.400081in}{1.126381in}}{\pgfqpoint{11.408439in}{1.135459in}}%
\pgfpathquadraticcurveto{\pgfqpoint{11.416796in}{1.144556in}}{\pgfqpoint{11.430990in}{1.144556in}}%
\pgfpathquadraticcurveto{\pgfqpoint{11.443707in}{1.144556in}}{\pgfqpoint{11.451110in}{1.136360in}}%
\pgfpathquadraticcurveto{\pgfqpoint{11.458513in}{1.128183in}}{\pgfqpoint{11.458513in}{1.114115in}}%
\pgfpathlineto{\pgfqpoint{11.458513in}{1.114115in}}%
\pgfpathclose%
\pgfpathmoveto{\pgfqpoint{11.448156in}{1.117159in}}%
\pgfpathquadraticcurveto{\pgfqpoint{11.448048in}{1.125643in}}{\pgfqpoint{11.443400in}{1.130704in}}%
\pgfpathquadraticcurveto{\pgfqpoint{11.438753in}{1.135784in}}{\pgfqpoint{11.431098in}{1.135784in}}%
\pgfpathquadraticcurveto{\pgfqpoint{11.422434in}{1.135784in}}{\pgfqpoint{11.417229in}{1.130884in}}%
\pgfpathquadraticcurveto{\pgfqpoint{11.412023in}{1.125985in}}{\pgfqpoint{11.411231in}{1.117087in}}%
\pgfpathlineto{\pgfqpoint{11.448156in}{1.117159in}}%
\pgfpathlineto{\pgfqpoint{11.448156in}{1.117159in}}%
\pgfpathclose%
\pgfpathmoveto{\pgfqpoint{11.512045in}{1.133370in}}%
\pgfpathquadraticcurveto{\pgfqpoint{11.510298in}{1.134379in}}{\pgfqpoint{11.508244in}{1.134847in}}%
\pgfpathquadraticcurveto{\pgfqpoint{11.506191in}{1.135333in}}{\pgfqpoint{11.503723in}{1.135333in}}%
\pgfpathquadraticcurveto{\pgfqpoint{11.494933in}{1.135333in}}{\pgfqpoint{11.490232in}{1.129623in}}%
\pgfpathquadraticcurveto{\pgfqpoint{11.485531in}{1.123914in}}{\pgfqpoint{11.485531in}{1.113214in}}%
\pgfpathlineto{\pgfqpoint{11.485531in}{1.080000in}}%
\pgfpathlineto{\pgfqpoint{11.475120in}{1.080000in}}%
\pgfpathlineto{\pgfqpoint{11.475120in}{1.143043in}}%
\pgfpathlineto{\pgfqpoint{11.485531in}{1.143043in}}%
\pgfpathlineto{\pgfqpoint{11.485531in}{1.133244in}}%
\pgfpathquadraticcurveto{\pgfqpoint{11.488809in}{1.138990in}}{\pgfqpoint{11.494033in}{1.141764in}}%
\pgfpathquadraticcurveto{\pgfqpoint{11.499274in}{1.144556in}}{\pgfqpoint{11.506767in}{1.144556in}}%
\pgfpathquadraticcurveto{\pgfqpoint{11.507830in}{1.144556in}}{\pgfqpoint{11.509127in}{1.144411in}}%
\pgfpathquadraticcurveto{\pgfqpoint{11.510424in}{1.144285in}}{\pgfqpoint{11.511991in}{1.143997in}}%
\pgfpathlineto{\pgfqpoint{11.512045in}{1.133370in}}%
\pgfpathlineto{\pgfqpoint{11.512045in}{1.133370in}}%
\pgfpathclose%
\pgfusepath{fill}%
\end{pgfscope}%
\begin{pgfscope}%
\pgfsetbuttcap%
\pgfsetroundjoin%
\definecolor{currentfill}{rgb}{0.309804,0.372549,0.419608}%
\pgfsetfillcolor{currentfill}%
\pgfsetlinewidth{0.000000pt}%
\definecolor{currentstroke}{rgb}{0.309804,0.372549,0.419608}%
\pgfsetstrokecolor{currentstroke}%
\pgfsetdash{}{0pt}%
\pgfpathmoveto{\pgfqpoint{11.629202in}{1.114115in}}%
\pgfpathlineto{\pgfqpoint{11.629202in}{1.109054in}}%
\pgfpathlineto{\pgfqpoint{11.581578in}{1.109054in}}%
\pgfpathquadraticcurveto{\pgfqpoint{11.582263in}{1.098354in}}{\pgfqpoint{11.588026in}{1.092753in}}%
\pgfpathquadraticcurveto{\pgfqpoint{11.593790in}{1.087151in}}{\pgfqpoint{11.604093in}{1.087151in}}%
\pgfpathquadraticcurveto{\pgfqpoint{11.610055in}{1.087151in}}{\pgfqpoint{11.615657in}{1.088610in}}%
\pgfpathquadraticcurveto{\pgfqpoint{11.621259in}{1.090069in}}{\pgfqpoint{11.626789in}{1.093005in}}%
\pgfpathlineto{\pgfqpoint{11.626789in}{1.083206in}}%
\pgfpathquadraticcurveto{\pgfqpoint{11.621205in}{1.080847in}}{\pgfqpoint{11.615351in}{1.079604in}}%
\pgfpathquadraticcurveto{\pgfqpoint{11.609497in}{1.078361in}}{\pgfqpoint{11.603481in}{1.078361in}}%
\pgfpathquadraticcurveto{\pgfqpoint{11.588387in}{1.078361in}}{\pgfqpoint{11.579579in}{1.087133in}}%
\pgfpathquadraticcurveto{\pgfqpoint{11.570771in}{1.095923in}}{\pgfqpoint{11.570771in}{1.110909in}}%
\pgfpathquadraticcurveto{\pgfqpoint{11.570771in}{1.126381in}}{\pgfqpoint{11.579128in}{1.135459in}}%
\pgfpathquadraticcurveto{\pgfqpoint{11.587486in}{1.144556in}}{\pgfqpoint{11.601680in}{1.144556in}}%
\pgfpathquadraticcurveto{\pgfqpoint{11.614396in}{1.144556in}}{\pgfqpoint{11.621799in}{1.136360in}}%
\pgfpathquadraticcurveto{\pgfqpoint{11.629202in}{1.128183in}}{\pgfqpoint{11.629202in}{1.114115in}}%
\pgfpathlineto{\pgfqpoint{11.629202in}{1.114115in}}%
\pgfpathclose%
\pgfpathmoveto{\pgfqpoint{11.618845in}{1.117159in}}%
\pgfpathquadraticcurveto{\pgfqpoint{11.618737in}{1.125643in}}{\pgfqpoint{11.614090in}{1.130704in}}%
\pgfpathquadraticcurveto{\pgfqpoint{11.609443in}{1.135784in}}{\pgfqpoint{11.601788in}{1.135784in}}%
\pgfpathquadraticcurveto{\pgfqpoint{11.593124in}{1.135784in}}{\pgfqpoint{11.587918in}{1.130884in}}%
\pgfpathquadraticcurveto{\pgfqpoint{11.582713in}{1.125985in}}{\pgfqpoint{11.581920in}{1.117087in}}%
\pgfpathlineto{\pgfqpoint{11.618845in}{1.117159in}}%
\pgfpathlineto{\pgfqpoint{11.618845in}{1.117159in}}%
\pgfpathclose%
\pgfpathmoveto{\pgfqpoint{11.686391in}{1.141187in}}%
\pgfpathlineto{\pgfqpoint{11.686391in}{1.131389in}}%
\pgfpathquadraticcurveto{\pgfqpoint{11.682014in}{1.133640in}}{\pgfqpoint{11.677277in}{1.134757in}}%
\pgfpathquadraticcurveto{\pgfqpoint{11.672557in}{1.135892in}}{\pgfqpoint{11.667478in}{1.135892in}}%
\pgfpathquadraticcurveto{\pgfqpoint{11.659769in}{1.135892in}}{\pgfqpoint{11.655914in}{1.133532in}}%
\pgfpathquadraticcurveto{\pgfqpoint{11.652060in}{1.131173in}}{\pgfqpoint{11.652060in}{1.126435in}}%
\pgfpathquadraticcurveto{\pgfqpoint{11.652060in}{1.122833in}}{\pgfqpoint{11.654815in}{1.120780in}}%
\pgfpathquadraticcurveto{\pgfqpoint{11.657571in}{1.118726in}}{\pgfqpoint{11.665911in}{1.116871in}}%
\pgfpathlineto{\pgfqpoint{11.669459in}{1.116078in}}%
\pgfpathquadraticcurveto{\pgfqpoint{11.680483in}{1.113719in}}{\pgfqpoint{11.685130in}{1.109414in}}%
\pgfpathquadraticcurveto{\pgfqpoint{11.689777in}{1.105109in}}{\pgfqpoint{11.689777in}{1.097400in}}%
\pgfpathquadraticcurveto{\pgfqpoint{11.689777in}{1.088610in}}{\pgfqpoint{11.682824in}{1.083476in}}%
\pgfpathquadraticcurveto{\pgfqpoint{11.675872in}{1.078361in}}{\pgfqpoint{11.663713in}{1.078361in}}%
\pgfpathquadraticcurveto{\pgfqpoint{11.658652in}{1.078361in}}{\pgfqpoint{11.653158in}{1.079352in}}%
\pgfpathquadraticcurveto{\pgfqpoint{11.647665in}{1.080342in}}{\pgfqpoint{11.641595in}{1.082306in}}%
\pgfpathlineto{\pgfqpoint{11.641595in}{1.093005in}}%
\pgfpathquadraticcurveto{\pgfqpoint{11.647340in}{1.090015in}}{\pgfqpoint{11.652906in}{1.088520in}}%
\pgfpathquadraticcurveto{\pgfqpoint{11.658472in}{1.087043in}}{\pgfqpoint{11.663948in}{1.087043in}}%
\pgfpathquadraticcurveto{\pgfqpoint{11.671261in}{1.087043in}}{\pgfqpoint{11.675187in}{1.089546in}}%
\pgfpathquadraticcurveto{\pgfqpoint{11.679132in}{1.092050in}}{\pgfqpoint{11.679132in}{1.096607in}}%
\pgfpathquadraticcurveto{\pgfqpoint{11.679132in}{1.100822in}}{\pgfqpoint{11.676286in}{1.103074in}}%
\pgfpathquadraticcurveto{\pgfqpoint{11.673458in}{1.105325in}}{\pgfqpoint{11.663822in}{1.107414in}}%
\pgfpathlineto{\pgfqpoint{11.660219in}{1.108261in}}%
\pgfpathquadraticcurveto{\pgfqpoint{11.650601in}{1.110278in}}{\pgfqpoint{11.646314in}{1.114475in}}%
\pgfpathquadraticcurveto{\pgfqpoint{11.642045in}{1.118672in}}{\pgfqpoint{11.642045in}{1.125985in}}%
\pgfpathquadraticcurveto{\pgfqpoint{11.642045in}{1.134883in}}{\pgfqpoint{11.648349in}{1.139710in}}%
\pgfpathquadraticcurveto{\pgfqpoint{11.654653in}{1.144556in}}{\pgfqpoint{11.666253in}{1.144556in}}%
\pgfpathquadraticcurveto{\pgfqpoint{11.671981in}{1.144556in}}{\pgfqpoint{11.677042in}{1.143709in}}%
\pgfpathquadraticcurveto{\pgfqpoint{11.682122in}{1.142880in}}{\pgfqpoint{11.686391in}{1.141187in}}%
\pgfpathlineto{\pgfqpoint{11.686391in}{1.141187in}}%
\pgfpathclose%
\pgfpathmoveto{\pgfqpoint{11.716507in}{1.160947in}}%
\pgfpathlineto{\pgfqpoint{11.716507in}{1.143043in}}%
\pgfpathlineto{\pgfqpoint{11.737834in}{1.143043in}}%
\pgfpathlineto{\pgfqpoint{11.737834in}{1.134991in}}%
\pgfpathlineto{\pgfqpoint{11.716507in}{1.134991in}}%
\pgfpathlineto{\pgfqpoint{11.716507in}{1.100768in}}%
\pgfpathquadraticcurveto{\pgfqpoint{11.716507in}{1.093059in}}{\pgfqpoint{11.718615in}{1.090861in}}%
\pgfpathquadraticcurveto{\pgfqpoint{11.720722in}{1.088664in}}{\pgfqpoint{11.727206in}{1.088664in}}%
\pgfpathlineto{\pgfqpoint{11.737834in}{1.088664in}}%
\pgfpathlineto{\pgfqpoint{11.737834in}{1.080000in}}%
\pgfpathlineto{\pgfqpoint{11.727206in}{1.080000in}}%
\pgfpathquadraticcurveto{\pgfqpoint{11.715210in}{1.080000in}}{\pgfqpoint{11.710653in}{1.084467in}}%
\pgfpathquadraticcurveto{\pgfqpoint{11.706096in}{1.088952in}}{\pgfqpoint{11.706096in}{1.100768in}}%
\pgfpathlineto{\pgfqpoint{11.706096in}{1.134991in}}%
\pgfpathlineto{\pgfqpoint{11.698495in}{1.134991in}}%
\pgfpathlineto{\pgfqpoint{11.698495in}{1.143043in}}%
\pgfpathlineto{\pgfqpoint{11.706096in}{1.143043in}}%
\pgfpathlineto{\pgfqpoint{11.706096in}{1.160947in}}%
\pgfpathlineto{\pgfqpoint{11.716507in}{1.160947in}}%
\pgfpathlineto{\pgfqpoint{11.716507in}{1.160947in}}%
\pgfpathclose%
\pgfpathmoveto{\pgfqpoint{11.751451in}{1.143043in}}%
\pgfpathlineto{\pgfqpoint{11.761808in}{1.143043in}}%
\pgfpathlineto{\pgfqpoint{11.761808in}{1.080000in}}%
\pgfpathlineto{\pgfqpoint{11.751451in}{1.080000in}}%
\pgfpathlineto{\pgfqpoint{11.751451in}{1.143043in}}%
\pgfpathlineto{\pgfqpoint{11.751451in}{1.143043in}}%
\pgfpathclose%
\pgfpathmoveto{\pgfqpoint{11.751451in}{1.167593in}}%
\pgfpathlineto{\pgfqpoint{11.761808in}{1.167593in}}%
\pgfpathlineto{\pgfqpoint{11.761808in}{1.154462in}}%
\pgfpathlineto{\pgfqpoint{11.751451in}{1.154462in}}%
\pgfpathlineto{\pgfqpoint{11.751451in}{1.167593in}}%
\pgfpathlineto{\pgfqpoint{11.751451in}{1.167593in}}%
\pgfpathclose%
\pgfpathmoveto{\pgfqpoint{11.832559in}{1.130938in}}%
\pgfpathquadraticcurveto{\pgfqpoint{11.836450in}{1.137927in}}{\pgfqpoint{11.841854in}{1.141241in}}%
\pgfpathquadraticcurveto{\pgfqpoint{11.847257in}{1.144556in}}{\pgfqpoint{11.854570in}{1.144556in}}%
\pgfpathquadraticcurveto{\pgfqpoint{11.864423in}{1.144556in}}{\pgfqpoint{11.869773in}{1.137657in}}%
\pgfpathquadraticcurveto{\pgfqpoint{11.875122in}{1.130776in}}{\pgfqpoint{11.875122in}{1.118060in}}%
\pgfpathlineto{\pgfqpoint{11.875122in}{1.080000in}}%
\pgfpathlineto{\pgfqpoint{11.864711in}{1.080000in}}%
\pgfpathlineto{\pgfqpoint{11.864711in}{1.117717in}}%
\pgfpathquadraticcurveto{\pgfqpoint{11.864711in}{1.126778in}}{\pgfqpoint{11.861487in}{1.131155in}}%
\pgfpathquadraticcurveto{\pgfqpoint{11.858281in}{1.135549in}}{\pgfqpoint{11.851706in}{1.135549in}}%
\pgfpathquadraticcurveto{\pgfqpoint{11.843655in}{1.135549in}}{\pgfqpoint{11.838972in}{1.130200in}}%
\pgfpathquadraticcurveto{\pgfqpoint{11.834307in}{1.124868in}}{\pgfqpoint{11.834307in}{1.115628in}}%
\pgfpathlineto{\pgfqpoint{11.834307in}{1.080000in}}%
\pgfpathlineto{\pgfqpoint{11.823896in}{1.080000in}}%
\pgfpathlineto{\pgfqpoint{11.823896in}{1.117717in}}%
\pgfpathquadraticcurveto{\pgfqpoint{11.823896in}{1.126832in}}{\pgfqpoint{11.820689in}{1.131191in}}%
\pgfpathquadraticcurveto{\pgfqpoint{11.817483in}{1.135549in}}{\pgfqpoint{11.810783in}{1.135549in}}%
\pgfpathquadraticcurveto{\pgfqpoint{11.802839in}{1.135549in}}{\pgfqpoint{11.798156in}{1.130182in}}%
\pgfpathquadraticcurveto{\pgfqpoint{11.793491in}{1.124814in}}{\pgfqpoint{11.793491in}{1.115628in}}%
\pgfpathlineto{\pgfqpoint{11.793491in}{1.080000in}}%
\pgfpathlineto{\pgfqpoint{11.783080in}{1.080000in}}%
\pgfpathlineto{\pgfqpoint{11.783080in}{1.143043in}}%
\pgfpathlineto{\pgfqpoint{11.793491in}{1.143043in}}%
\pgfpathlineto{\pgfqpoint{11.793491in}{1.133244in}}%
\pgfpathquadraticcurveto{\pgfqpoint{11.797039in}{1.139044in}}{\pgfqpoint{11.801993in}{1.141800in}}%
\pgfpathquadraticcurveto{\pgfqpoint{11.806946in}{1.144556in}}{\pgfqpoint{11.813755in}{1.144556in}}%
\pgfpathquadraticcurveto{\pgfqpoint{11.820635in}{1.144556in}}{\pgfqpoint{11.825445in}{1.141061in}}%
\pgfpathquadraticcurveto{\pgfqpoint{11.830254in}{1.137585in}}{\pgfqpoint{11.832559in}{1.130938in}}%
\pgfpathlineto{\pgfqpoint{11.832559in}{1.130938in}}%
\pgfpathclose%
\pgfpathmoveto{\pgfqpoint{11.924421in}{1.111683in}}%
\pgfpathquadraticcurveto{\pgfqpoint{11.911867in}{1.111683in}}{\pgfqpoint{11.907022in}{1.108819in}}%
\pgfpathquadraticcurveto{\pgfqpoint{11.902176in}{1.105956in}}{\pgfqpoint{11.902176in}{1.099021in}}%
\pgfpathquadraticcurveto{\pgfqpoint{11.902176in}{1.093509in}}{\pgfqpoint{11.905815in}{1.090267in}}%
\pgfpathquadraticcurveto{\pgfqpoint{11.909453in}{1.087043in}}{\pgfqpoint{11.915686in}{1.087043in}}%
\pgfpathquadraticcurveto{\pgfqpoint{11.924313in}{1.087043in}}{\pgfqpoint{11.929519in}{1.093149in}}%
\pgfpathquadraticcurveto{\pgfqpoint{11.934724in}{1.099255in}}{\pgfqpoint{11.934724in}{1.109378in}}%
\pgfpathlineto{\pgfqpoint{11.934724in}{1.111683in}}%
\pgfpathlineto{\pgfqpoint{11.924421in}{1.111683in}}%
\pgfpathlineto{\pgfqpoint{11.924421in}{1.111683in}}%
\pgfpathclose%
\pgfpathmoveto{\pgfqpoint{11.945081in}{1.115970in}}%
\pgfpathlineto{\pgfqpoint{11.945081in}{1.080000in}}%
\pgfpathlineto{\pgfqpoint{11.934724in}{1.080000in}}%
\pgfpathlineto{\pgfqpoint{11.934724in}{1.089564in}}%
\pgfpathquadraticcurveto{\pgfqpoint{11.931176in}{1.083837in}}{\pgfqpoint{11.925880in}{1.081099in}}%
\pgfpathquadraticcurveto{\pgfqpoint{11.920585in}{1.078361in}}{\pgfqpoint{11.912930in}{1.078361in}}%
\pgfpathquadraticcurveto{\pgfqpoint{11.903257in}{1.078361in}}{\pgfqpoint{11.897529in}{1.083801in}}%
\pgfpathquadraticcurveto{\pgfqpoint{11.891819in}{1.089240in}}{\pgfqpoint{11.891819in}{1.098354in}}%
\pgfpathquadraticcurveto{\pgfqpoint{11.891819in}{1.108982in}}{\pgfqpoint{11.898934in}{1.114385in}}%
\pgfpathquadraticcurveto{\pgfqpoint{11.906067in}{1.119789in}}{\pgfqpoint{11.920189in}{1.119789in}}%
\pgfpathlineto{\pgfqpoint{11.934724in}{1.119789in}}%
\pgfpathlineto{\pgfqpoint{11.934724in}{1.120816in}}%
\pgfpathquadraticcurveto{\pgfqpoint{11.934724in}{1.127966in}}{\pgfqpoint{11.930023in}{1.131875in}}%
\pgfpathquadraticcurveto{\pgfqpoint{11.925322in}{1.135784in}}{\pgfqpoint{11.916820in}{1.135784in}}%
\pgfpathquadraticcurveto{\pgfqpoint{11.911417in}{1.135784in}}{\pgfqpoint{11.906283in}{1.134487in}}%
\pgfpathquadraticcurveto{\pgfqpoint{11.901168in}{1.133190in}}{\pgfqpoint{11.896449in}{1.130596in}}%
\pgfpathlineto{\pgfqpoint{11.896449in}{1.140179in}}%
\pgfpathquadraticcurveto{\pgfqpoint{11.902122in}{1.142376in}}{\pgfqpoint{11.907472in}{1.143457in}}%
\pgfpathquadraticcurveto{\pgfqpoint{11.912822in}{1.144556in}}{\pgfqpoint{11.917883in}{1.144556in}}%
\pgfpathquadraticcurveto{\pgfqpoint{11.931572in}{1.144556in}}{\pgfqpoint{11.938327in}{1.137459in}}%
\pgfpathquadraticcurveto{\pgfqpoint{11.945081in}{1.130380in}}{\pgfqpoint{11.945081in}{1.115970in}}%
\pgfpathlineto{\pgfqpoint{11.945081in}{1.115970in}}%
\pgfpathclose%
\pgfpathmoveto{\pgfqpoint{11.976657in}{1.160947in}}%
\pgfpathlineto{\pgfqpoint{11.976657in}{1.143043in}}%
\pgfpathlineto{\pgfqpoint{11.997983in}{1.143043in}}%
\pgfpathlineto{\pgfqpoint{11.997983in}{1.134991in}}%
\pgfpathlineto{\pgfqpoint{11.976657in}{1.134991in}}%
\pgfpathlineto{\pgfqpoint{11.976657in}{1.100768in}}%
\pgfpathquadraticcurveto{\pgfqpoint{11.976657in}{1.093059in}}{\pgfqpoint{11.978764in}{1.090861in}}%
\pgfpathquadraticcurveto{\pgfqpoint{11.980871in}{1.088664in}}{\pgfqpoint{11.987356in}{1.088664in}}%
\pgfpathlineto{\pgfqpoint{11.997983in}{1.088664in}}%
\pgfpathlineto{\pgfqpoint{11.997983in}{1.080000in}}%
\pgfpathlineto{\pgfqpoint{11.987356in}{1.080000in}}%
\pgfpathquadraticcurveto{\pgfqpoint{11.975360in}{1.080000in}}{\pgfqpoint{11.970803in}{1.084467in}}%
\pgfpathquadraticcurveto{\pgfqpoint{11.966246in}{1.088952in}}{\pgfqpoint{11.966246in}{1.100768in}}%
\pgfpathlineto{\pgfqpoint{11.966246in}{1.134991in}}%
\pgfpathlineto{\pgfqpoint{11.958644in}{1.134991in}}%
\pgfpathlineto{\pgfqpoint{11.958644in}{1.143043in}}%
\pgfpathlineto{\pgfqpoint{11.966246in}{1.143043in}}%
\pgfpathlineto{\pgfqpoint{11.966246in}{1.160947in}}%
\pgfpathlineto{\pgfqpoint{11.976657in}{1.160947in}}%
\pgfpathlineto{\pgfqpoint{11.976657in}{1.160947in}}%
\pgfpathclose%
\pgfpathmoveto{\pgfqpoint{12.065529in}{1.114115in}}%
\pgfpathlineto{\pgfqpoint{12.065529in}{1.109054in}}%
\pgfpathlineto{\pgfqpoint{12.017904in}{1.109054in}}%
\pgfpathquadraticcurveto{\pgfqpoint{12.018589in}{1.098354in}}{\pgfqpoint{12.024353in}{1.092753in}}%
\pgfpathquadraticcurveto{\pgfqpoint{12.030117in}{1.087151in}}{\pgfqpoint{12.040420in}{1.087151in}}%
\pgfpathquadraticcurveto{\pgfqpoint{12.046382in}{1.087151in}}{\pgfqpoint{12.051983in}{1.088610in}}%
\pgfpathquadraticcurveto{\pgfqpoint{12.057585in}{1.090069in}}{\pgfqpoint{12.063115in}{1.093005in}}%
\pgfpathlineto{\pgfqpoint{12.063115in}{1.083206in}}%
\pgfpathquadraticcurveto{\pgfqpoint{12.057531in}{1.080847in}}{\pgfqpoint{12.051677in}{1.079604in}}%
\pgfpathquadraticcurveto{\pgfqpoint{12.045823in}{1.078361in}}{\pgfqpoint{12.039807in}{1.078361in}}%
\pgfpathquadraticcurveto{\pgfqpoint{12.024713in}{1.078361in}}{\pgfqpoint{12.015905in}{1.087133in}}%
\pgfpathquadraticcurveto{\pgfqpoint{12.007097in}{1.095923in}}{\pgfqpoint{12.007097in}{1.110909in}}%
\pgfpathquadraticcurveto{\pgfqpoint{12.007097in}{1.126381in}}{\pgfqpoint{12.015455in}{1.135459in}}%
\pgfpathquadraticcurveto{\pgfqpoint{12.023812in}{1.144556in}}{\pgfqpoint{12.038006in}{1.144556in}}%
\pgfpathquadraticcurveto{\pgfqpoint{12.050723in}{1.144556in}}{\pgfqpoint{12.058126in}{1.136360in}}%
\pgfpathquadraticcurveto{\pgfqpoint{12.065529in}{1.128183in}}{\pgfqpoint{12.065529in}{1.114115in}}%
\pgfpathlineto{\pgfqpoint{12.065529in}{1.114115in}}%
\pgfpathclose%
\pgfpathmoveto{\pgfqpoint{12.055172in}{1.117159in}}%
\pgfpathquadraticcurveto{\pgfqpoint{12.055064in}{1.125643in}}{\pgfqpoint{12.050416in}{1.130704in}}%
\pgfpathquadraticcurveto{\pgfqpoint{12.045769in}{1.135784in}}{\pgfqpoint{12.038114in}{1.135784in}}%
\pgfpathquadraticcurveto{\pgfqpoint{12.029450in}{1.135784in}}{\pgfqpoint{12.024245in}{1.130884in}}%
\pgfpathquadraticcurveto{\pgfqpoint{12.019039in}{1.125985in}}{\pgfqpoint{12.018247in}{1.117087in}}%
\pgfpathlineto{\pgfqpoint{12.055172in}{1.117159in}}%
\pgfpathlineto{\pgfqpoint{12.055172in}{1.117159in}}%
\pgfpathclose%
\pgfusepath{fill}%
\end{pgfscope}%
\begin{pgfscope}%
\pgfsetbuttcap%
\pgfsetroundjoin%
\definecolor{currentfill}{rgb}{0.309804,0.372549,0.419608}%
\pgfsetfillcolor{currentfill}%
\pgfsetlinewidth{0.000000pt}%
\definecolor{currentstroke}{rgb}{0.309804,0.372549,0.419608}%
\pgfsetstrokecolor{currentstroke}%
\pgfsetdash{}{0pt}%
\pgfpathmoveto{\pgfqpoint{12.189586in}{1.118060in}}%
\pgfpathlineto{\pgfqpoint{12.189586in}{1.080000in}}%
\pgfpathlineto{\pgfqpoint{12.179229in}{1.080000in}}%
\pgfpathlineto{\pgfqpoint{12.179229in}{1.117717in}}%
\pgfpathquadraticcurveto{\pgfqpoint{12.179229in}{1.126669in}}{\pgfqpoint{12.175735in}{1.131100in}}%
\pgfpathquadraticcurveto{\pgfqpoint{12.172240in}{1.135549in}}{\pgfqpoint{12.165270in}{1.135549in}}%
\pgfpathquadraticcurveto{\pgfqpoint{12.156876in}{1.135549in}}{\pgfqpoint{12.152031in}{1.130200in}}%
\pgfpathquadraticcurveto{\pgfqpoint{12.147185in}{1.124868in}}{\pgfqpoint{12.147185in}{1.115628in}}%
\pgfpathlineto{\pgfqpoint{12.147185in}{1.080000in}}%
\pgfpathlineto{\pgfqpoint{12.136774in}{1.080000in}}%
\pgfpathlineto{\pgfqpoint{12.136774in}{1.167593in}}%
\pgfpathlineto{\pgfqpoint{12.147185in}{1.167593in}}%
\pgfpathlineto{\pgfqpoint{12.147185in}{1.133244in}}%
\pgfpathquadraticcurveto{\pgfqpoint{12.150914in}{1.138936in}}{\pgfqpoint{12.155939in}{1.141746in}}%
\pgfpathquadraticcurveto{\pgfqpoint{12.160983in}{1.144556in}}{\pgfqpoint{12.167575in}{1.144556in}}%
\pgfpathquadraticcurveto{\pgfqpoint{12.178437in}{1.144556in}}{\pgfqpoint{12.184002in}{1.137837in}}%
\pgfpathquadraticcurveto{\pgfqpoint{12.189586in}{1.131118in}}{\pgfqpoint{12.189586in}{1.118060in}}%
\pgfpathlineto{\pgfqpoint{12.189586in}{1.118060in}}%
\pgfpathclose%
\pgfpathmoveto{\pgfqpoint{12.238885in}{1.111683in}}%
\pgfpathquadraticcurveto{\pgfqpoint{12.226331in}{1.111683in}}{\pgfqpoint{12.221486in}{1.108819in}}%
\pgfpathquadraticcurveto{\pgfqpoint{12.216640in}{1.105956in}}{\pgfqpoint{12.216640in}{1.099021in}}%
\pgfpathquadraticcurveto{\pgfqpoint{12.216640in}{1.093509in}}{\pgfqpoint{12.220279in}{1.090267in}}%
\pgfpathquadraticcurveto{\pgfqpoint{12.223917in}{1.087043in}}{\pgfqpoint{12.230149in}{1.087043in}}%
\pgfpathquadraticcurveto{\pgfqpoint{12.238777in}{1.087043in}}{\pgfqpoint{12.243983in}{1.093149in}}%
\pgfpathquadraticcurveto{\pgfqpoint{12.249188in}{1.099255in}}{\pgfqpoint{12.249188in}{1.109378in}}%
\pgfpathlineto{\pgfqpoint{12.249188in}{1.111683in}}%
\pgfpathlineto{\pgfqpoint{12.238885in}{1.111683in}}%
\pgfpathlineto{\pgfqpoint{12.238885in}{1.111683in}}%
\pgfpathclose%
\pgfpathmoveto{\pgfqpoint{12.259545in}{1.115970in}}%
\pgfpathlineto{\pgfqpoint{12.259545in}{1.080000in}}%
\pgfpathlineto{\pgfqpoint{12.249188in}{1.080000in}}%
\pgfpathlineto{\pgfqpoint{12.249188in}{1.089564in}}%
\pgfpathquadraticcurveto{\pgfqpoint{12.245640in}{1.083837in}}{\pgfqpoint{12.240344in}{1.081099in}}%
\pgfpathquadraticcurveto{\pgfqpoint{12.235049in}{1.078361in}}{\pgfqpoint{12.227394in}{1.078361in}}%
\pgfpathquadraticcurveto{\pgfqpoint{12.217721in}{1.078361in}}{\pgfqpoint{12.211993in}{1.083801in}}%
\pgfpathquadraticcurveto{\pgfqpoint{12.206283in}{1.089240in}}{\pgfqpoint{12.206283in}{1.098354in}}%
\pgfpathquadraticcurveto{\pgfqpoint{12.206283in}{1.108982in}}{\pgfqpoint{12.213398in}{1.114385in}}%
\pgfpathquadraticcurveto{\pgfqpoint{12.220531in}{1.119789in}}{\pgfqpoint{12.234652in}{1.119789in}}%
\pgfpathlineto{\pgfqpoint{12.249188in}{1.119789in}}%
\pgfpathlineto{\pgfqpoint{12.249188in}{1.120816in}}%
\pgfpathquadraticcurveto{\pgfqpoint{12.249188in}{1.127966in}}{\pgfqpoint{12.244487in}{1.131875in}}%
\pgfpathquadraticcurveto{\pgfqpoint{12.239786in}{1.135784in}}{\pgfqpoint{12.231284in}{1.135784in}}%
\pgfpathquadraticcurveto{\pgfqpoint{12.225881in}{1.135784in}}{\pgfqpoint{12.220747in}{1.134487in}}%
\pgfpathquadraticcurveto{\pgfqpoint{12.215632in}{1.133190in}}{\pgfqpoint{12.210912in}{1.130596in}}%
\pgfpathlineto{\pgfqpoint{12.210912in}{1.140179in}}%
\pgfpathquadraticcurveto{\pgfqpoint{12.216586in}{1.142376in}}{\pgfqpoint{12.221936in}{1.143457in}}%
\pgfpathquadraticcurveto{\pgfqpoint{12.227286in}{1.144556in}}{\pgfqpoint{12.232347in}{1.144556in}}%
\pgfpathquadraticcurveto{\pgfqpoint{12.246036in}{1.144556in}}{\pgfqpoint{12.252791in}{1.137459in}}%
\pgfpathquadraticcurveto{\pgfqpoint{12.259545in}{1.130380in}}{\pgfqpoint{12.259545in}{1.115970in}}%
\pgfpathlineto{\pgfqpoint{12.259545in}{1.115970in}}%
\pgfpathclose%
\pgfpathmoveto{\pgfqpoint{12.321057in}{1.141187in}}%
\pgfpathlineto{\pgfqpoint{12.321057in}{1.131389in}}%
\pgfpathquadraticcurveto{\pgfqpoint{12.316680in}{1.133640in}}{\pgfqpoint{12.311943in}{1.134757in}}%
\pgfpathquadraticcurveto{\pgfqpoint{12.307223in}{1.135892in}}{\pgfqpoint{12.302144in}{1.135892in}}%
\pgfpathquadraticcurveto{\pgfqpoint{12.294435in}{1.135892in}}{\pgfqpoint{12.290580in}{1.133532in}}%
\pgfpathquadraticcurveto{\pgfqpoint{12.286726in}{1.131173in}}{\pgfqpoint{12.286726in}{1.126435in}}%
\pgfpathquadraticcurveto{\pgfqpoint{12.286726in}{1.122833in}}{\pgfqpoint{12.289481in}{1.120780in}}%
\pgfpathquadraticcurveto{\pgfqpoint{12.292237in}{1.118726in}}{\pgfqpoint{12.300577in}{1.116871in}}%
\pgfpathlineto{\pgfqpoint{12.304125in}{1.116078in}}%
\pgfpathquadraticcurveto{\pgfqpoint{12.315149in}{1.113719in}}{\pgfqpoint{12.319796in}{1.109414in}}%
\pgfpathquadraticcurveto{\pgfqpoint{12.324443in}{1.105109in}}{\pgfqpoint{12.324443in}{1.097400in}}%
\pgfpathquadraticcurveto{\pgfqpoint{12.324443in}{1.088610in}}{\pgfqpoint{12.317490in}{1.083476in}}%
\pgfpathquadraticcurveto{\pgfqpoint{12.310538in}{1.078361in}}{\pgfqpoint{12.298379in}{1.078361in}}%
\pgfpathquadraticcurveto{\pgfqpoint{12.293318in}{1.078361in}}{\pgfqpoint{12.287824in}{1.079352in}}%
\pgfpathquadraticcurveto{\pgfqpoint{12.282331in}{1.080342in}}{\pgfqpoint{12.276261in}{1.082306in}}%
\pgfpathlineto{\pgfqpoint{12.276261in}{1.093005in}}%
\pgfpathquadraticcurveto{\pgfqpoint{12.282006in}{1.090015in}}{\pgfqpoint{12.287572in}{1.088520in}}%
\pgfpathquadraticcurveto{\pgfqpoint{12.293138in}{1.087043in}}{\pgfqpoint{12.298614in}{1.087043in}}%
\pgfpathquadraticcurveto{\pgfqpoint{12.305927in}{1.087043in}}{\pgfqpoint{12.309853in}{1.089546in}}%
\pgfpathquadraticcurveto{\pgfqpoint{12.313798in}{1.092050in}}{\pgfqpoint{12.313798in}{1.096607in}}%
\pgfpathquadraticcurveto{\pgfqpoint{12.313798in}{1.100822in}}{\pgfqpoint{12.310952in}{1.103074in}}%
\pgfpathquadraticcurveto{\pgfqpoint{12.308124in}{1.105325in}}{\pgfqpoint{12.298488in}{1.107414in}}%
\pgfpathlineto{\pgfqpoint{12.294885in}{1.108261in}}%
\pgfpathquadraticcurveto{\pgfqpoint{12.285267in}{1.110278in}}{\pgfqpoint{12.280980in}{1.114475in}}%
\pgfpathquadraticcurveto{\pgfqpoint{12.276711in}{1.118672in}}{\pgfqpoint{12.276711in}{1.125985in}}%
\pgfpathquadraticcurveto{\pgfqpoint{12.276711in}{1.134883in}}{\pgfqpoint{12.283015in}{1.139710in}}%
\pgfpathquadraticcurveto{\pgfqpoint{12.289319in}{1.144556in}}{\pgfqpoint{12.300919in}{1.144556in}}%
\pgfpathquadraticcurveto{\pgfqpoint{12.306647in}{1.144556in}}{\pgfqpoint{12.311708in}{1.143709in}}%
\pgfpathquadraticcurveto{\pgfqpoint{12.316788in}{1.142880in}}{\pgfqpoint{12.321057in}{1.141187in}}%
\pgfpathlineto{\pgfqpoint{12.321057in}{1.141187in}}%
\pgfpathclose%
\pgfusepath{fill}%
\end{pgfscope}%
\begin{pgfscope}%
\pgfsetbuttcap%
\pgfsetroundjoin%
\definecolor{currentfill}{rgb}{0.309804,0.372549,0.419608}%
\pgfsetfillcolor{currentfill}%
\pgfsetlinewidth{0.000000pt}%
\definecolor{currentstroke}{rgb}{0.309804,0.372549,0.419608}%
\pgfsetstrokecolor{currentstroke}%
\pgfsetdash{}{0pt}%
\pgfpathmoveto{\pgfqpoint{12.442996in}{1.114115in}}%
\pgfpathlineto{\pgfqpoint{12.442996in}{1.109054in}}%
\pgfpathlineto{\pgfqpoint{12.395372in}{1.109054in}}%
\pgfpathquadraticcurveto{\pgfqpoint{12.396056in}{1.098354in}}{\pgfqpoint{12.401820in}{1.092753in}}%
\pgfpathquadraticcurveto{\pgfqpoint{12.407584in}{1.087151in}}{\pgfqpoint{12.417887in}{1.087151in}}%
\pgfpathquadraticcurveto{\pgfqpoint{12.423849in}{1.087151in}}{\pgfqpoint{12.429451in}{1.088610in}}%
\pgfpathquadraticcurveto{\pgfqpoint{12.435053in}{1.090069in}}{\pgfqpoint{12.440582in}{1.093005in}}%
\pgfpathlineto{\pgfqpoint{12.440582in}{1.083206in}}%
\pgfpathquadraticcurveto{\pgfqpoint{12.434999in}{1.080847in}}{\pgfqpoint{12.429145in}{1.079604in}}%
\pgfpathquadraticcurveto{\pgfqpoint{12.423291in}{1.078361in}}{\pgfqpoint{12.417275in}{1.078361in}}%
\pgfpathquadraticcurveto{\pgfqpoint{12.402180in}{1.078361in}}{\pgfqpoint{12.393372in}{1.087133in}}%
\pgfpathquadraticcurveto{\pgfqpoint{12.384565in}{1.095923in}}{\pgfqpoint{12.384565in}{1.110909in}}%
\pgfpathquadraticcurveto{\pgfqpoint{12.384565in}{1.126381in}}{\pgfqpoint{12.392922in}{1.135459in}}%
\pgfpathquadraticcurveto{\pgfqpoint{12.401280in}{1.144556in}}{\pgfqpoint{12.415473in}{1.144556in}}%
\pgfpathquadraticcurveto{\pgfqpoint{12.428190in}{1.144556in}}{\pgfqpoint{12.435593in}{1.136360in}}%
\pgfpathquadraticcurveto{\pgfqpoint{12.442996in}{1.128183in}}{\pgfqpoint{12.442996in}{1.114115in}}%
\pgfpathlineto{\pgfqpoint{12.442996in}{1.114115in}}%
\pgfpathclose%
\pgfpathmoveto{\pgfqpoint{12.432639in}{1.117159in}}%
\pgfpathquadraticcurveto{\pgfqpoint{12.432531in}{1.125643in}}{\pgfqpoint{12.427884in}{1.130704in}}%
\pgfpathquadraticcurveto{\pgfqpoint{12.423237in}{1.135784in}}{\pgfqpoint{12.415581in}{1.135784in}}%
\pgfpathquadraticcurveto{\pgfqpoint{12.406918in}{1.135784in}}{\pgfqpoint{12.401712in}{1.130884in}}%
\pgfpathquadraticcurveto{\pgfqpoint{12.396507in}{1.125985in}}{\pgfqpoint{12.395714in}{1.117087in}}%
\pgfpathlineto{\pgfqpoint{12.432639in}{1.117159in}}%
\pgfpathlineto{\pgfqpoint{12.432639in}{1.117159in}}%
\pgfpathclose%
\pgfpathmoveto{\pgfqpoint{12.530751in}{1.167593in}}%
\pgfpathlineto{\pgfqpoint{12.530751in}{1.158965in}}%
\pgfpathlineto{\pgfqpoint{12.520844in}{1.158965in}}%
\pgfpathquadraticcurveto{\pgfqpoint{12.515279in}{1.158965in}}{\pgfqpoint{12.513099in}{1.156714in}}%
\pgfpathquadraticcurveto{\pgfqpoint{12.510938in}{1.154462in}}{\pgfqpoint{12.510938in}{1.148608in}}%
\pgfpathlineto{\pgfqpoint{12.510938in}{1.143043in}}%
\pgfpathlineto{\pgfqpoint{12.527995in}{1.143043in}}%
\pgfpathlineto{\pgfqpoint{12.527995in}{1.134991in}}%
\pgfpathlineto{\pgfqpoint{12.510938in}{1.134991in}}%
\pgfpathlineto{\pgfqpoint{12.510938in}{1.080000in}}%
\pgfpathlineto{\pgfqpoint{12.500527in}{1.080000in}}%
\pgfpathlineto{\pgfqpoint{12.500527in}{1.134991in}}%
\pgfpathlineto{\pgfqpoint{12.472104in}{1.134991in}}%
\pgfpathlineto{\pgfqpoint{12.472104in}{1.080000in}}%
\pgfpathlineto{\pgfqpoint{12.461693in}{1.080000in}}%
\pgfpathlineto{\pgfqpoint{12.461693in}{1.134991in}}%
\pgfpathlineto{\pgfqpoint{12.451786in}{1.134991in}}%
\pgfpathlineto{\pgfqpoint{12.451786in}{1.143043in}}%
\pgfpathlineto{\pgfqpoint{12.461693in}{1.143043in}}%
\pgfpathlineto{\pgfqpoint{12.461693in}{1.147438in}}%
\pgfpathquadraticcurveto{\pgfqpoint{12.461693in}{1.157957in}}{\pgfqpoint{12.466592in}{1.162766in}}%
\pgfpathquadraticcurveto{\pgfqpoint{12.471491in}{1.167593in}}{\pgfqpoint{12.482118in}{1.167593in}}%
\pgfpathlineto{\pgfqpoint{12.491917in}{1.167593in}}%
\pgfpathlineto{\pgfqpoint{12.491917in}{1.158965in}}%
\pgfpathlineto{\pgfqpoint{12.482010in}{1.158965in}}%
\pgfpathquadraticcurveto{\pgfqpoint{12.476445in}{1.158965in}}{\pgfqpoint{12.474247in}{1.156714in}}%
\pgfpathquadraticcurveto{\pgfqpoint{12.472104in}{1.154462in}}{\pgfqpoint{12.472104in}{1.148608in}}%
\pgfpathlineto{\pgfqpoint{12.472104in}{1.143043in}}%
\pgfpathlineto{\pgfqpoint{12.500527in}{1.143043in}}%
\pgfpathlineto{\pgfqpoint{12.500527in}{1.147438in}}%
\pgfpathquadraticcurveto{\pgfqpoint{12.500527in}{1.157957in}}{\pgfqpoint{12.505426in}{1.162766in}}%
\pgfpathquadraticcurveto{\pgfqpoint{12.510325in}{1.167593in}}{\pgfqpoint{12.520971in}{1.167593in}}%
\pgfpathlineto{\pgfqpoint{12.530751in}{1.167593in}}%
\pgfpathlineto{\pgfqpoint{12.530751in}{1.167593in}}%
\pgfpathclose%
\pgfpathmoveto{\pgfqpoint{12.593343in}{1.114115in}}%
\pgfpathlineto{\pgfqpoint{12.593343in}{1.109054in}}%
\pgfpathlineto{\pgfqpoint{12.545719in}{1.109054in}}%
\pgfpathquadraticcurveto{\pgfqpoint{12.546404in}{1.098354in}}{\pgfqpoint{12.552168in}{1.092753in}}%
\pgfpathquadraticcurveto{\pgfqpoint{12.557932in}{1.087151in}}{\pgfqpoint{12.568234in}{1.087151in}}%
\pgfpathquadraticcurveto{\pgfqpoint{12.574196in}{1.087151in}}{\pgfqpoint{12.579798in}{1.088610in}}%
\pgfpathquadraticcurveto{\pgfqpoint{12.585400in}{1.090069in}}{\pgfqpoint{12.590930in}{1.093005in}}%
\pgfpathlineto{\pgfqpoint{12.590930in}{1.083206in}}%
\pgfpathquadraticcurveto{\pgfqpoint{12.585346in}{1.080847in}}{\pgfqpoint{12.579492in}{1.079604in}}%
\pgfpathquadraticcurveto{\pgfqpoint{12.573638in}{1.078361in}}{\pgfqpoint{12.567622in}{1.078361in}}%
\pgfpathquadraticcurveto{\pgfqpoint{12.552528in}{1.078361in}}{\pgfqpoint{12.543720in}{1.087133in}}%
\pgfpathquadraticcurveto{\pgfqpoint{12.534912in}{1.095923in}}{\pgfqpoint{12.534912in}{1.110909in}}%
\pgfpathquadraticcurveto{\pgfqpoint{12.534912in}{1.126381in}}{\pgfqpoint{12.543270in}{1.135459in}}%
\pgfpathquadraticcurveto{\pgfqpoint{12.551627in}{1.144556in}}{\pgfqpoint{12.565821in}{1.144556in}}%
\pgfpathquadraticcurveto{\pgfqpoint{12.578537in}{1.144556in}}{\pgfqpoint{12.585940in}{1.136360in}}%
\pgfpathquadraticcurveto{\pgfqpoint{12.593343in}{1.128183in}}{\pgfqpoint{12.593343in}{1.114115in}}%
\pgfpathlineto{\pgfqpoint{12.593343in}{1.114115in}}%
\pgfpathclose%
\pgfpathmoveto{\pgfqpoint{12.582986in}{1.117159in}}%
\pgfpathquadraticcurveto{\pgfqpoint{12.582878in}{1.125643in}}{\pgfqpoint{12.578231in}{1.130704in}}%
\pgfpathquadraticcurveto{\pgfqpoint{12.573584in}{1.135784in}}{\pgfqpoint{12.565929in}{1.135784in}}%
\pgfpathquadraticcurveto{\pgfqpoint{12.557265in}{1.135784in}}{\pgfqpoint{12.552060in}{1.130884in}}%
\pgfpathquadraticcurveto{\pgfqpoint{12.546854in}{1.125985in}}{\pgfqpoint{12.546062in}{1.117087in}}%
\pgfpathlineto{\pgfqpoint{12.582986in}{1.117159in}}%
\pgfpathlineto{\pgfqpoint{12.582986in}{1.117159in}}%
\pgfpathclose%
\pgfpathmoveto{\pgfqpoint{12.655719in}{1.140629in}}%
\pgfpathlineto{\pgfqpoint{12.655719in}{1.130938in}}%
\pgfpathquadraticcurveto{\pgfqpoint{12.651325in}{1.133370in}}{\pgfqpoint{12.646912in}{1.134577in}}%
\pgfpathquadraticcurveto{\pgfqpoint{12.642499in}{1.135784in}}{\pgfqpoint{12.637996in}{1.135784in}}%
\pgfpathquadraticcurveto{\pgfqpoint{12.627909in}{1.135784in}}{\pgfqpoint{12.622325in}{1.129389in}}%
\pgfpathquadraticcurveto{\pgfqpoint{12.616759in}{1.123013in}}{\pgfqpoint{12.616759in}{1.111467in}}%
\pgfpathquadraticcurveto{\pgfqpoint{12.616759in}{1.099921in}}{\pgfqpoint{12.622325in}{1.093527in}}%
\pgfpathquadraticcurveto{\pgfqpoint{12.627909in}{1.087151in}}{\pgfqpoint{12.637996in}{1.087151in}}%
\pgfpathquadraticcurveto{\pgfqpoint{12.642499in}{1.087151in}}{\pgfqpoint{12.646912in}{1.088358in}}%
\pgfpathquadraticcurveto{\pgfqpoint{12.651325in}{1.089564in}}{\pgfqpoint{12.655719in}{1.091996in}}%
\pgfpathlineto{\pgfqpoint{12.655719in}{1.082414in}}%
\pgfpathquadraticcurveto{\pgfqpoint{12.651379in}{1.080396in}}{\pgfqpoint{12.646731in}{1.079388in}}%
\pgfpathquadraticcurveto{\pgfqpoint{12.642102in}{1.078361in}}{\pgfqpoint{12.636861in}{1.078361in}}%
\pgfpathquadraticcurveto{\pgfqpoint{12.622613in}{1.078361in}}{\pgfqpoint{12.614219in}{1.087313in}}%
\pgfpathquadraticcurveto{\pgfqpoint{12.605844in}{1.096265in}}{\pgfqpoint{12.605844in}{1.111467in}}%
\pgfpathquadraticcurveto{\pgfqpoint{12.605844in}{1.126886in}}{\pgfqpoint{12.614310in}{1.135712in}}%
\pgfpathquadraticcurveto{\pgfqpoint{12.622793in}{1.144556in}}{\pgfqpoint{12.637545in}{1.144556in}}%
\pgfpathquadraticcurveto{\pgfqpoint{12.642318in}{1.144556in}}{\pgfqpoint{12.646876in}{1.143565in}}%
\pgfpathquadraticcurveto{\pgfqpoint{12.651433in}{1.142592in}}{\pgfqpoint{12.655719in}{1.140629in}}%
\pgfpathlineto{\pgfqpoint{12.655719in}{1.140629in}}%
\pgfpathclose%
\pgfpathmoveto{\pgfqpoint{12.683981in}{1.160947in}}%
\pgfpathlineto{\pgfqpoint{12.683981in}{1.143043in}}%
\pgfpathlineto{\pgfqpoint{12.705307in}{1.143043in}}%
\pgfpathlineto{\pgfqpoint{12.705307in}{1.134991in}}%
\pgfpathlineto{\pgfqpoint{12.683981in}{1.134991in}}%
\pgfpathlineto{\pgfqpoint{12.683981in}{1.100768in}}%
\pgfpathquadraticcurveto{\pgfqpoint{12.683981in}{1.093059in}}{\pgfqpoint{12.686088in}{1.090861in}}%
\pgfpathquadraticcurveto{\pgfqpoint{12.688195in}{1.088664in}}{\pgfqpoint{12.694680in}{1.088664in}}%
\pgfpathlineto{\pgfqpoint{12.705307in}{1.088664in}}%
\pgfpathlineto{\pgfqpoint{12.705307in}{1.080000in}}%
\pgfpathlineto{\pgfqpoint{12.694680in}{1.080000in}}%
\pgfpathquadraticcurveto{\pgfqpoint{12.682684in}{1.080000in}}{\pgfqpoint{12.678127in}{1.084467in}}%
\pgfpathquadraticcurveto{\pgfqpoint{12.673570in}{1.088952in}}{\pgfqpoint{12.673570in}{1.100768in}}%
\pgfpathlineto{\pgfqpoint{12.673570in}{1.134991in}}%
\pgfpathlineto{\pgfqpoint{12.665968in}{1.134991in}}%
\pgfpathlineto{\pgfqpoint{12.665968in}{1.143043in}}%
\pgfpathlineto{\pgfqpoint{12.673570in}{1.143043in}}%
\pgfpathlineto{\pgfqpoint{12.673570in}{1.160947in}}%
\pgfpathlineto{\pgfqpoint{12.683981in}{1.160947in}}%
\pgfpathlineto{\pgfqpoint{12.683981in}{1.160947in}}%
\pgfpathclose%
\pgfpathmoveto{\pgfqpoint{12.718924in}{1.143043in}}%
\pgfpathlineto{\pgfqpoint{12.729281in}{1.143043in}}%
\pgfpathlineto{\pgfqpoint{12.729281in}{1.080000in}}%
\pgfpathlineto{\pgfqpoint{12.718924in}{1.080000in}}%
\pgfpathlineto{\pgfqpoint{12.718924in}{1.143043in}}%
\pgfpathlineto{\pgfqpoint{12.718924in}{1.143043in}}%
\pgfpathclose%
\pgfpathmoveto{\pgfqpoint{12.718924in}{1.167593in}}%
\pgfpathlineto{\pgfqpoint{12.729281in}{1.167593in}}%
\pgfpathlineto{\pgfqpoint{12.729281in}{1.154462in}}%
\pgfpathlineto{\pgfqpoint{12.718924in}{1.154462in}}%
\pgfpathlineto{\pgfqpoint{12.718924in}{1.167593in}}%
\pgfpathlineto{\pgfqpoint{12.718924in}{1.167593in}}%
\pgfpathclose%
\pgfpathmoveto{\pgfqpoint{12.743529in}{1.143043in}}%
\pgfpathlineto{\pgfqpoint{12.754498in}{1.143043in}}%
\pgfpathlineto{\pgfqpoint{12.774203in}{1.090141in}}%
\pgfpathlineto{\pgfqpoint{12.793909in}{1.143043in}}%
\pgfpathlineto{\pgfqpoint{12.804878in}{1.143043in}}%
\pgfpathlineto{\pgfqpoint{12.781228in}{1.080000in}}%
\pgfpathlineto{\pgfqpoint{12.767161in}{1.080000in}}%
\pgfpathlineto{\pgfqpoint{12.743529in}{1.143043in}}%
\pgfpathlineto{\pgfqpoint{12.743529in}{1.143043in}}%
\pgfpathclose%
\pgfpathmoveto{\pgfqpoint{12.873108in}{1.114115in}}%
\pgfpathlineto{\pgfqpoint{12.873108in}{1.109054in}}%
\pgfpathlineto{\pgfqpoint{12.825484in}{1.109054in}}%
\pgfpathquadraticcurveto{\pgfqpoint{12.826168in}{1.098354in}}{\pgfqpoint{12.831932in}{1.092753in}}%
\pgfpathquadraticcurveto{\pgfqpoint{12.837696in}{1.087151in}}{\pgfqpoint{12.847999in}{1.087151in}}%
\pgfpathquadraticcurveto{\pgfqpoint{12.853961in}{1.087151in}}{\pgfqpoint{12.859563in}{1.088610in}}%
\pgfpathquadraticcurveto{\pgfqpoint{12.865165in}{1.090069in}}{\pgfqpoint{12.870695in}{1.093005in}}%
\pgfpathlineto{\pgfqpoint{12.870695in}{1.083206in}}%
\pgfpathquadraticcurveto{\pgfqpoint{12.865111in}{1.080847in}}{\pgfqpoint{12.859257in}{1.079604in}}%
\pgfpathquadraticcurveto{\pgfqpoint{12.853403in}{1.078361in}}{\pgfqpoint{12.847387in}{1.078361in}}%
\pgfpathquadraticcurveto{\pgfqpoint{12.832293in}{1.078361in}}{\pgfqpoint{12.823485in}{1.087133in}}%
\pgfpathquadraticcurveto{\pgfqpoint{12.814677in}{1.095923in}}{\pgfqpoint{12.814677in}{1.110909in}}%
\pgfpathquadraticcurveto{\pgfqpoint{12.814677in}{1.126381in}}{\pgfqpoint{12.823034in}{1.135459in}}%
\pgfpathquadraticcurveto{\pgfqpoint{12.831392in}{1.144556in}}{\pgfqpoint{12.845586in}{1.144556in}}%
\pgfpathquadraticcurveto{\pgfqpoint{12.858302in}{1.144556in}}{\pgfqpoint{12.865705in}{1.136360in}}%
\pgfpathquadraticcurveto{\pgfqpoint{12.873108in}{1.128183in}}{\pgfqpoint{12.873108in}{1.114115in}}%
\pgfpathlineto{\pgfqpoint{12.873108in}{1.114115in}}%
\pgfpathclose%
\pgfpathmoveto{\pgfqpoint{12.862751in}{1.117159in}}%
\pgfpathquadraticcurveto{\pgfqpoint{12.862643in}{1.125643in}}{\pgfqpoint{12.857996in}{1.130704in}}%
\pgfpathquadraticcurveto{\pgfqpoint{12.853349in}{1.135784in}}{\pgfqpoint{12.845694in}{1.135784in}}%
\pgfpathquadraticcurveto{\pgfqpoint{12.837030in}{1.135784in}}{\pgfqpoint{12.831824in}{1.130884in}}%
\pgfpathquadraticcurveto{\pgfqpoint{12.826619in}{1.125985in}}{\pgfqpoint{12.825826in}{1.117087in}}%
\pgfpathlineto{\pgfqpoint{12.862751in}{1.117159in}}%
\pgfpathlineto{\pgfqpoint{12.862751in}{1.117159in}}%
\pgfpathclose%
\pgfusepath{fill}%
\end{pgfscope}%
\begin{pgfscope}%
\pgfsetbuttcap%
\pgfsetroundjoin%
\definecolor{currentfill}{rgb}{0.309804,0.372549,0.419608}%
\pgfsetfillcolor{currentfill}%
\pgfsetlinewidth{0.000000pt}%
\definecolor{currentstroke}{rgb}{0.309804,0.372549,0.419608}%
\pgfsetstrokecolor{currentstroke}%
\pgfsetdash{}{0pt}%
\pgfpathmoveto{\pgfqpoint{12.986337in}{1.140629in}}%
\pgfpathlineto{\pgfqpoint{12.986337in}{1.130938in}}%
\pgfpathquadraticcurveto{\pgfqpoint{12.981942in}{1.133370in}}{\pgfqpoint{12.977529in}{1.134577in}}%
\pgfpathquadraticcurveto{\pgfqpoint{12.973116in}{1.135784in}}{\pgfqpoint{12.968613in}{1.135784in}}%
\pgfpathquadraticcurveto{\pgfqpoint{12.958526in}{1.135784in}}{\pgfqpoint{12.952943in}{1.129389in}}%
\pgfpathquadraticcurveto{\pgfqpoint{12.947377in}{1.123013in}}{\pgfqpoint{12.947377in}{1.111467in}}%
\pgfpathquadraticcurveto{\pgfqpoint{12.947377in}{1.099921in}}{\pgfqpoint{12.952943in}{1.093527in}}%
\pgfpathquadraticcurveto{\pgfqpoint{12.958526in}{1.087151in}}{\pgfqpoint{12.968613in}{1.087151in}}%
\pgfpathquadraticcurveto{\pgfqpoint{12.973116in}{1.087151in}}{\pgfqpoint{12.977529in}{1.088358in}}%
\pgfpathquadraticcurveto{\pgfqpoint{12.981942in}{1.089564in}}{\pgfqpoint{12.986337in}{1.091996in}}%
\pgfpathlineto{\pgfqpoint{12.986337in}{1.082414in}}%
\pgfpathquadraticcurveto{\pgfqpoint{12.981996in}{1.080396in}}{\pgfqpoint{12.977349in}{1.079388in}}%
\pgfpathquadraticcurveto{\pgfqpoint{12.972720in}{1.078361in}}{\pgfqpoint{12.967478in}{1.078361in}}%
\pgfpathquadraticcurveto{\pgfqpoint{12.953231in}{1.078361in}}{\pgfqpoint{12.944837in}{1.087313in}}%
\pgfpathquadraticcurveto{\pgfqpoint{12.936461in}{1.096265in}}{\pgfqpoint{12.936461in}{1.111467in}}%
\pgfpathquadraticcurveto{\pgfqpoint{12.936461in}{1.126886in}}{\pgfqpoint{12.944927in}{1.135712in}}%
\pgfpathquadraticcurveto{\pgfqpoint{12.953411in}{1.144556in}}{\pgfqpoint{12.968163in}{1.144556in}}%
\pgfpathquadraticcurveto{\pgfqpoint{12.972936in}{1.144556in}}{\pgfqpoint{12.977493in}{1.143565in}}%
\pgfpathquadraticcurveto{\pgfqpoint{12.982050in}{1.142592in}}{\pgfqpoint{12.986337in}{1.140629in}}%
\pgfpathlineto{\pgfqpoint{12.986337in}{1.140629in}}%
\pgfpathclose%
\pgfpathmoveto{\pgfqpoint{13.004349in}{1.167593in}}%
\pgfpathlineto{\pgfqpoint{13.014706in}{1.167593in}}%
\pgfpathlineto{\pgfqpoint{13.014706in}{1.080000in}}%
\pgfpathlineto{\pgfqpoint{13.004349in}{1.080000in}}%
\pgfpathlineto{\pgfqpoint{13.004349in}{1.167593in}}%
\pgfpathlineto{\pgfqpoint{13.004349in}{1.167593in}}%
\pgfpathclose%
\pgfpathmoveto{\pgfqpoint{13.035312in}{1.104875in}}%
\pgfpathlineto{\pgfqpoint{13.035312in}{1.143043in}}%
\pgfpathlineto{\pgfqpoint{13.045669in}{1.143043in}}%
\pgfpathlineto{\pgfqpoint{13.045669in}{1.105271in}}%
\pgfpathquadraticcurveto{\pgfqpoint{13.045669in}{1.096319in}}{\pgfqpoint{13.049145in}{1.091834in}}%
\pgfpathquadraticcurveto{\pgfqpoint{13.052640in}{1.087367in}}{\pgfqpoint{13.059629in}{1.087367in}}%
\pgfpathquadraticcurveto{\pgfqpoint{13.068004in}{1.087367in}}{\pgfqpoint{13.072867in}{1.092717in}}%
\pgfpathquadraticcurveto{\pgfqpoint{13.077749in}{1.098066in}}{\pgfqpoint{13.077749in}{1.107306in}}%
\pgfpathlineto{\pgfqpoint{13.077749in}{1.143043in}}%
\pgfpathlineto{\pgfqpoint{13.088106in}{1.143043in}}%
\pgfpathlineto{\pgfqpoint{13.088106in}{1.080000in}}%
\pgfpathlineto{\pgfqpoint{13.077749in}{1.080000in}}%
\pgfpathlineto{\pgfqpoint{13.077749in}{1.089691in}}%
\pgfpathquadraticcurveto{\pgfqpoint{13.073984in}{1.083945in}}{\pgfqpoint{13.068995in}{1.081153in}}%
\pgfpathquadraticcurveto{\pgfqpoint{13.064023in}{1.078361in}}{\pgfqpoint{13.057431in}{1.078361in}}%
\pgfpathquadraticcurveto{\pgfqpoint{13.046570in}{1.078361in}}{\pgfqpoint{13.040932in}{1.085115in}}%
\pgfpathquadraticcurveto{\pgfqpoint{13.035312in}{1.091870in}}{\pgfqpoint{13.035312in}{1.104875in}}%
\pgfpathlineto{\pgfqpoint{13.035312in}{1.104875in}}%
\pgfpathclose%
\pgfpathmoveto{\pgfqpoint{13.061376in}{1.144556in}}%
\pgfpathlineto{\pgfqpoint{13.061376in}{1.144556in}}%
\pgfpathlineto{\pgfqpoint{13.061376in}{1.144556in}}%
\pgfpathclose%
\pgfpathmoveto{\pgfqpoint{13.149617in}{1.141187in}}%
\pgfpathlineto{\pgfqpoint{13.149617in}{1.131389in}}%
\pgfpathquadraticcurveto{\pgfqpoint{13.145240in}{1.133640in}}{\pgfqpoint{13.140503in}{1.134757in}}%
\pgfpathquadraticcurveto{\pgfqpoint{13.135784in}{1.135892in}}{\pgfqpoint{13.130704in}{1.135892in}}%
\pgfpathquadraticcurveto{\pgfqpoint{13.122995in}{1.135892in}}{\pgfqpoint{13.119141in}{1.133532in}}%
\pgfpathquadraticcurveto{\pgfqpoint{13.115286in}{1.131173in}}{\pgfqpoint{13.115286in}{1.126435in}}%
\pgfpathquadraticcurveto{\pgfqpoint{13.115286in}{1.122833in}}{\pgfqpoint{13.118042in}{1.120780in}}%
\pgfpathquadraticcurveto{\pgfqpoint{13.120798in}{1.118726in}}{\pgfqpoint{13.129137in}{1.116871in}}%
\pgfpathlineto{\pgfqpoint{13.132686in}{1.116078in}}%
\pgfpathquadraticcurveto{\pgfqpoint{13.143709in}{1.113719in}}{\pgfqpoint{13.148356in}{1.109414in}}%
\pgfpathquadraticcurveto{\pgfqpoint{13.153004in}{1.105109in}}{\pgfqpoint{13.153004in}{1.097400in}}%
\pgfpathquadraticcurveto{\pgfqpoint{13.153004in}{1.088610in}}{\pgfqpoint{13.146051in}{1.083476in}}%
\pgfpathquadraticcurveto{\pgfqpoint{13.139098in}{1.078361in}}{\pgfqpoint{13.126940in}{1.078361in}}%
\pgfpathquadraticcurveto{\pgfqpoint{13.121879in}{1.078361in}}{\pgfqpoint{13.116385in}{1.079352in}}%
\pgfpathquadraticcurveto{\pgfqpoint{13.110891in}{1.080342in}}{\pgfqpoint{13.104821in}{1.082306in}}%
\pgfpathlineto{\pgfqpoint{13.104821in}{1.093005in}}%
\pgfpathquadraticcurveto{\pgfqpoint{13.110567in}{1.090015in}}{\pgfqpoint{13.116133in}{1.088520in}}%
\pgfpathquadraticcurveto{\pgfqpoint{13.121698in}{1.087043in}}{\pgfqpoint{13.127174in}{1.087043in}}%
\pgfpathquadraticcurveto{\pgfqpoint{13.134487in}{1.087043in}}{\pgfqpoint{13.138414in}{1.089546in}}%
\pgfpathquadraticcurveto{\pgfqpoint{13.142358in}{1.092050in}}{\pgfqpoint{13.142358in}{1.096607in}}%
\pgfpathquadraticcurveto{\pgfqpoint{13.142358in}{1.100822in}}{\pgfqpoint{13.139512in}{1.103074in}}%
\pgfpathquadraticcurveto{\pgfqpoint{13.136685in}{1.105325in}}{\pgfqpoint{13.127048in}{1.107414in}}%
\pgfpathlineto{\pgfqpoint{13.123446in}{1.108261in}}%
\pgfpathquadraticcurveto{\pgfqpoint{13.113827in}{1.110278in}}{\pgfqpoint{13.109540in}{1.114475in}}%
\pgfpathquadraticcurveto{\pgfqpoint{13.105271in}{1.118672in}}{\pgfqpoint{13.105271in}{1.125985in}}%
\pgfpathquadraticcurveto{\pgfqpoint{13.105271in}{1.134883in}}{\pgfqpoint{13.111576in}{1.139710in}}%
\pgfpathquadraticcurveto{\pgfqpoint{13.117880in}{1.144556in}}{\pgfqpoint{13.129480in}{1.144556in}}%
\pgfpathquadraticcurveto{\pgfqpoint{13.135208in}{1.144556in}}{\pgfqpoint{13.140269in}{1.143709in}}%
\pgfpathquadraticcurveto{\pgfqpoint{13.145348in}{1.142880in}}{\pgfqpoint{13.149617in}{1.141187in}}%
\pgfpathlineto{\pgfqpoint{13.149617in}{1.141187in}}%
\pgfpathclose%
\pgfpathmoveto{\pgfqpoint{13.179734in}{1.160947in}}%
\pgfpathlineto{\pgfqpoint{13.179734in}{1.143043in}}%
\pgfpathlineto{\pgfqpoint{13.201060in}{1.143043in}}%
\pgfpathlineto{\pgfqpoint{13.201060in}{1.134991in}}%
\pgfpathlineto{\pgfqpoint{13.179734in}{1.134991in}}%
\pgfpathlineto{\pgfqpoint{13.179734in}{1.100768in}}%
\pgfpathquadraticcurveto{\pgfqpoint{13.179734in}{1.093059in}}{\pgfqpoint{13.181841in}{1.090861in}}%
\pgfpathquadraticcurveto{\pgfqpoint{13.183948in}{1.088664in}}{\pgfqpoint{13.190433in}{1.088664in}}%
\pgfpathlineto{\pgfqpoint{13.201060in}{1.088664in}}%
\pgfpathlineto{\pgfqpoint{13.201060in}{1.080000in}}%
\pgfpathlineto{\pgfqpoint{13.190433in}{1.080000in}}%
\pgfpathquadraticcurveto{\pgfqpoint{13.178437in}{1.080000in}}{\pgfqpoint{13.173880in}{1.084467in}}%
\pgfpathquadraticcurveto{\pgfqpoint{13.169323in}{1.088952in}}{\pgfqpoint{13.169323in}{1.100768in}}%
\pgfpathlineto{\pgfqpoint{13.169323in}{1.134991in}}%
\pgfpathlineto{\pgfqpoint{13.161721in}{1.134991in}}%
\pgfpathlineto{\pgfqpoint{13.161721in}{1.143043in}}%
\pgfpathlineto{\pgfqpoint{13.169323in}{1.143043in}}%
\pgfpathlineto{\pgfqpoint{13.169323in}{1.160947in}}%
\pgfpathlineto{\pgfqpoint{13.179734in}{1.160947in}}%
\pgfpathlineto{\pgfqpoint{13.179734in}{1.160947in}}%
\pgfpathclose%
\pgfpathmoveto{\pgfqpoint{13.268606in}{1.114115in}}%
\pgfpathlineto{\pgfqpoint{13.268606in}{1.109054in}}%
\pgfpathlineto{\pgfqpoint{13.220981in}{1.109054in}}%
\pgfpathquadraticcurveto{\pgfqpoint{13.221666in}{1.098354in}}{\pgfqpoint{13.227430in}{1.092753in}}%
\pgfpathquadraticcurveto{\pgfqpoint{13.233194in}{1.087151in}}{\pgfqpoint{13.243497in}{1.087151in}}%
\pgfpathquadraticcurveto{\pgfqpoint{13.249459in}{1.087151in}}{\pgfqpoint{13.255060in}{1.088610in}}%
\pgfpathquadraticcurveto{\pgfqpoint{13.260662in}{1.090069in}}{\pgfqpoint{13.266192in}{1.093005in}}%
\pgfpathlineto{\pgfqpoint{13.266192in}{1.083206in}}%
\pgfpathquadraticcurveto{\pgfqpoint{13.260608in}{1.080847in}}{\pgfqpoint{13.254754in}{1.079604in}}%
\pgfpathquadraticcurveto{\pgfqpoint{13.248900in}{1.078361in}}{\pgfqpoint{13.242884in}{1.078361in}}%
\pgfpathquadraticcurveto{\pgfqpoint{13.227790in}{1.078361in}}{\pgfqpoint{13.218982in}{1.087133in}}%
\pgfpathquadraticcurveto{\pgfqpoint{13.210174in}{1.095923in}}{\pgfqpoint{13.210174in}{1.110909in}}%
\pgfpathquadraticcurveto{\pgfqpoint{13.210174in}{1.126381in}}{\pgfqpoint{13.218532in}{1.135459in}}%
\pgfpathquadraticcurveto{\pgfqpoint{13.226889in}{1.144556in}}{\pgfqpoint{13.241083in}{1.144556in}}%
\pgfpathquadraticcurveto{\pgfqpoint{13.253800in}{1.144556in}}{\pgfqpoint{13.261203in}{1.136360in}}%
\pgfpathquadraticcurveto{\pgfqpoint{13.268606in}{1.128183in}}{\pgfqpoint{13.268606in}{1.114115in}}%
\pgfpathlineto{\pgfqpoint{13.268606in}{1.114115in}}%
\pgfpathclose%
\pgfpathmoveto{\pgfqpoint{13.258249in}{1.117159in}}%
\pgfpathquadraticcurveto{\pgfqpoint{13.258140in}{1.125643in}}{\pgfqpoint{13.253493in}{1.130704in}}%
\pgfpathquadraticcurveto{\pgfqpoint{13.248846in}{1.135784in}}{\pgfqpoint{13.241191in}{1.135784in}}%
\pgfpathquadraticcurveto{\pgfqpoint{13.232527in}{1.135784in}}{\pgfqpoint{13.227322in}{1.130884in}}%
\pgfpathquadraticcurveto{\pgfqpoint{13.222116in}{1.125985in}}{\pgfqpoint{13.221324in}{1.117087in}}%
\pgfpathlineto{\pgfqpoint{13.258249in}{1.117159in}}%
\pgfpathlineto{\pgfqpoint{13.258249in}{1.117159in}}%
\pgfpathclose%
\pgfpathmoveto{\pgfqpoint{13.322138in}{1.133370in}}%
\pgfpathquadraticcurveto{\pgfqpoint{13.320390in}{1.134379in}}{\pgfqpoint{13.318337in}{1.134847in}}%
\pgfpathquadraticcurveto{\pgfqpoint{13.316284in}{1.135333in}}{\pgfqpoint{13.313816in}{1.135333in}}%
\pgfpathquadraticcurveto{\pgfqpoint{13.305026in}{1.135333in}}{\pgfqpoint{13.300325in}{1.129623in}}%
\pgfpathquadraticcurveto{\pgfqpoint{13.295624in}{1.123914in}}{\pgfqpoint{13.295624in}{1.113214in}}%
\pgfpathlineto{\pgfqpoint{13.295624in}{1.080000in}}%
\pgfpathlineto{\pgfqpoint{13.285213in}{1.080000in}}%
\pgfpathlineto{\pgfqpoint{13.285213in}{1.143043in}}%
\pgfpathlineto{\pgfqpoint{13.295624in}{1.143043in}}%
\pgfpathlineto{\pgfqpoint{13.295624in}{1.133244in}}%
\pgfpathquadraticcurveto{\pgfqpoint{13.298902in}{1.138990in}}{\pgfqpoint{13.304125in}{1.141764in}}%
\pgfpathquadraticcurveto{\pgfqpoint{13.309367in}{1.144556in}}{\pgfqpoint{13.316860in}{1.144556in}}%
\pgfpathquadraticcurveto{\pgfqpoint{13.317923in}{1.144556in}}{\pgfqpoint{13.319220in}{1.144411in}}%
\pgfpathquadraticcurveto{\pgfqpoint{13.320517in}{1.144285in}}{\pgfqpoint{13.322084in}{1.143997in}}%
\pgfpathlineto{\pgfqpoint{13.322138in}{1.133370in}}%
\pgfpathlineto{\pgfqpoint{13.322138in}{1.133370in}}%
\pgfpathclose%
\pgfpathmoveto{\pgfqpoint{13.373184in}{1.141187in}}%
\pgfpathlineto{\pgfqpoint{13.373184in}{1.131389in}}%
\pgfpathquadraticcurveto{\pgfqpoint{13.368807in}{1.133640in}}{\pgfqpoint{13.364070in}{1.134757in}}%
\pgfpathquadraticcurveto{\pgfqpoint{13.359351in}{1.135892in}}{\pgfqpoint{13.354271in}{1.135892in}}%
\pgfpathquadraticcurveto{\pgfqpoint{13.346562in}{1.135892in}}{\pgfqpoint{13.342708in}{1.133532in}}%
\pgfpathquadraticcurveto{\pgfqpoint{13.338853in}{1.131173in}}{\pgfqpoint{13.338853in}{1.126435in}}%
\pgfpathquadraticcurveto{\pgfqpoint{13.338853in}{1.122833in}}{\pgfqpoint{13.341609in}{1.120780in}}%
\pgfpathquadraticcurveto{\pgfqpoint{13.344365in}{1.118726in}}{\pgfqpoint{13.352704in}{1.116871in}}%
\pgfpathlineto{\pgfqpoint{13.356253in}{1.116078in}}%
\pgfpathquadraticcurveto{\pgfqpoint{13.367276in}{1.113719in}}{\pgfqpoint{13.371923in}{1.109414in}}%
\pgfpathquadraticcurveto{\pgfqpoint{13.376570in}{1.105109in}}{\pgfqpoint{13.376570in}{1.097400in}}%
\pgfpathquadraticcurveto{\pgfqpoint{13.376570in}{1.088610in}}{\pgfqpoint{13.369618in}{1.083476in}}%
\pgfpathquadraticcurveto{\pgfqpoint{13.362665in}{1.078361in}}{\pgfqpoint{13.350507in}{1.078361in}}%
\pgfpathquadraticcurveto{\pgfqpoint{13.345445in}{1.078361in}}{\pgfqpoint{13.339952in}{1.079352in}}%
\pgfpathquadraticcurveto{\pgfqpoint{13.334458in}{1.080342in}}{\pgfqpoint{13.328388in}{1.082306in}}%
\pgfpathlineto{\pgfqpoint{13.328388in}{1.093005in}}%
\pgfpathquadraticcurveto{\pgfqpoint{13.334134in}{1.090015in}}{\pgfqpoint{13.339699in}{1.088520in}}%
\pgfpathquadraticcurveto{\pgfqpoint{13.345265in}{1.087043in}}{\pgfqpoint{13.350741in}{1.087043in}}%
\pgfpathquadraticcurveto{\pgfqpoint{13.358054in}{1.087043in}}{\pgfqpoint{13.361981in}{1.089546in}}%
\pgfpathquadraticcurveto{\pgfqpoint{13.365925in}{1.092050in}}{\pgfqpoint{13.365925in}{1.096607in}}%
\pgfpathquadraticcurveto{\pgfqpoint{13.365925in}{1.100822in}}{\pgfqpoint{13.363079in}{1.103074in}}%
\pgfpathquadraticcurveto{\pgfqpoint{13.360251in}{1.105325in}}{\pgfqpoint{13.350615in}{1.107414in}}%
\pgfpathlineto{\pgfqpoint{13.347012in}{1.108261in}}%
\pgfpathquadraticcurveto{\pgfqpoint{13.337394in}{1.110278in}}{\pgfqpoint{13.333107in}{1.114475in}}%
\pgfpathquadraticcurveto{\pgfqpoint{13.328838in}{1.118672in}}{\pgfqpoint{13.328838in}{1.125985in}}%
\pgfpathquadraticcurveto{\pgfqpoint{13.328838in}{1.134883in}}{\pgfqpoint{13.335142in}{1.139710in}}%
\pgfpathquadraticcurveto{\pgfqpoint{13.341447in}{1.144556in}}{\pgfqpoint{13.353046in}{1.144556in}}%
\pgfpathquadraticcurveto{\pgfqpoint{13.358774in}{1.144556in}}{\pgfqpoint{13.363836in}{1.143709in}}%
\pgfpathquadraticcurveto{\pgfqpoint{13.368915in}{1.142880in}}{\pgfqpoint{13.373184in}{1.141187in}}%
\pgfpathlineto{\pgfqpoint{13.373184in}{1.141187in}}%
\pgfpathclose%
\pgfusepath{fill}%
\end{pgfscope}%
\begin{pgfscope}%
\pgfsetbuttcap%
\pgfsetroundjoin%
\definecolor{currentfill}{rgb}{0.309804,0.372549,0.419608}%
\pgfsetfillcolor{currentfill}%
\pgfsetlinewidth{0.000000pt}%
\definecolor{currentstroke}{rgb}{0.309804,0.372549,0.419608}%
\pgfsetstrokecolor{currentstroke}%
\pgfsetdash{}{0pt}%
\pgfpathmoveto{\pgfqpoint{13.526369in}{1.136738in}}%
\pgfpathlineto{\pgfqpoint{13.468280in}{1.116078in}}%
\pgfpathlineto{\pgfqpoint{13.526369in}{1.095544in}}%
\pgfpathlineto{\pgfqpoint{13.526369in}{1.085296in}}%
\pgfpathlineto{\pgfqpoint{13.454212in}{1.111467in}}%
\pgfpathlineto{\pgfqpoint{13.454212in}{1.120816in}}%
\pgfpathlineto{\pgfqpoint{13.526369in}{1.146987in}}%
\pgfpathlineto{\pgfqpoint{13.526369in}{1.136738in}}%
\pgfpathlineto{\pgfqpoint{13.526369in}{1.136738in}}%
\pgfpathclose%
\pgfpathmoveto{\pgfqpoint{13.551046in}{1.164045in}}%
\pgfpathlineto{\pgfqpoint{13.595680in}{1.164045in}}%
\pgfpathlineto{\pgfqpoint{13.595680in}{1.154462in}}%
\pgfpathlineto{\pgfqpoint{13.561457in}{1.154462in}}%
\pgfpathlineto{\pgfqpoint{13.561457in}{1.133874in}}%
\pgfpathquadraticcurveto{\pgfqpoint{13.563924in}{1.134721in}}{\pgfqpoint{13.566392in}{1.135135in}}%
\pgfpathquadraticcurveto{\pgfqpoint{13.568878in}{1.135549in}}{\pgfqpoint{13.571363in}{1.135549in}}%
\pgfpathquadraticcurveto{\pgfqpoint{13.585431in}{1.135549in}}{\pgfqpoint{13.593644in}{1.127840in}}%
\pgfpathquadraticcurveto{\pgfqpoint{13.601876in}{1.120131in}}{\pgfqpoint{13.601876in}{1.106964in}}%
\pgfpathquadraticcurveto{\pgfqpoint{13.601876in}{1.093401in}}{\pgfqpoint{13.593428in}{1.085872in}}%
\pgfpathquadraticcurveto{\pgfqpoint{13.584980in}{1.078361in}}{\pgfqpoint{13.569616in}{1.078361in}}%
\pgfpathquadraticcurveto{\pgfqpoint{13.564321in}{1.078361in}}{\pgfqpoint{13.558827in}{1.079262in}}%
\pgfpathquadraticcurveto{\pgfqpoint{13.553351in}{1.080162in}}{\pgfqpoint{13.547497in}{1.081963in}}%
\pgfpathlineto{\pgfqpoint{13.547497in}{1.093401in}}%
\pgfpathquadraticcurveto{\pgfqpoint{13.552559in}{1.090645in}}{\pgfqpoint{13.557962in}{1.089294in}}%
\pgfpathquadraticcurveto{\pgfqpoint{13.563366in}{1.087943in}}{\pgfqpoint{13.569382in}{1.087943in}}%
\pgfpathquadraticcurveto{\pgfqpoint{13.579127in}{1.087943in}}{\pgfqpoint{13.584800in}{1.093059in}}%
\pgfpathquadraticcurveto{\pgfqpoint{13.590492in}{1.098174in}}{\pgfqpoint{13.590492in}{1.106964in}}%
\pgfpathquadraticcurveto{\pgfqpoint{13.590492in}{1.115736in}}{\pgfqpoint{13.584800in}{1.120852in}}%
\pgfpathquadraticcurveto{\pgfqpoint{13.579127in}{1.125985in}}{\pgfqpoint{13.569382in}{1.125985in}}%
\pgfpathquadraticcurveto{\pgfqpoint{13.564825in}{1.125985in}}{\pgfqpoint{13.560286in}{1.124976in}}%
\pgfpathquadraticcurveto{\pgfqpoint{13.555765in}{1.123968in}}{\pgfqpoint{13.551046in}{1.121824in}}%
\pgfpathlineto{\pgfqpoint{13.551046in}{1.164045in}}%
\pgfpathlineto{\pgfqpoint{13.551046in}{1.164045in}}%
\pgfpathclose%
\pgfusepath{fill}%
\end{pgfscope}%
\begin{pgfscope}%
\pgfsetbuttcap%
\pgfsetroundjoin%
\definecolor{currentfill}{rgb}{0.309804,0.372549,0.419608}%
\pgfsetfillcolor{currentfill}%
\pgfsetlinewidth{0.000000pt}%
\definecolor{currentstroke}{rgb}{0.309804,0.372549,0.419608}%
\pgfsetstrokecolor{currentstroke}%
\pgfsetdash{}{0pt}%
\pgfpathmoveto{\pgfqpoint{0.959286in}{0.952184in}}%
\pgfpathquadraticcurveto{\pgfqpoint{0.950964in}{0.952184in}}{\pgfqpoint{0.946119in}{0.945681in}}%
\pgfpathquadraticcurveto{\pgfqpoint{0.941274in}{0.939179in}}{\pgfqpoint{0.941274in}{0.927867in}}%
\pgfpathquadraticcurveto{\pgfqpoint{0.941274in}{0.916556in}}{\pgfqpoint{0.946083in}{0.910053in}}%
\pgfpathquadraticcurveto{\pgfqpoint{0.950910in}{0.903551in}}{\pgfqpoint{0.959286in}{0.903551in}}%
\pgfpathquadraticcurveto{\pgfqpoint{0.967571in}{0.903551in}}{\pgfqpoint{0.972399in}{0.910071in}}%
\pgfpathquadraticcurveto{\pgfqpoint{0.977244in}{0.916610in}}{\pgfqpoint{0.977244in}{0.927867in}}%
\pgfpathquadraticcurveto{\pgfqpoint{0.977244in}{0.939071in}}{\pgfqpoint{0.972399in}{0.945627in}}%
\pgfpathquadraticcurveto{\pgfqpoint{0.967571in}{0.952184in}}{\pgfqpoint{0.959286in}{0.952184in}}%
\pgfpathlineto{\pgfqpoint{0.959286in}{0.952184in}}%
\pgfpathclose%
\pgfpathmoveto{\pgfqpoint{0.959286in}{0.960956in}}%
\pgfpathquadraticcurveto{\pgfqpoint{0.972795in}{0.960956in}}{\pgfqpoint{0.980504in}{0.952166in}}%
\pgfpathquadraticcurveto{\pgfqpoint{0.988231in}{0.943394in}}{\pgfqpoint{0.988231in}{0.927867in}}%
\pgfpathquadraticcurveto{\pgfqpoint{0.988231in}{0.912395in}}{\pgfqpoint{0.980504in}{0.903569in}}%
\pgfpathquadraticcurveto{\pgfqpoint{0.972795in}{0.894761in}}{\pgfqpoint{0.959286in}{0.894761in}}%
\pgfpathquadraticcurveto{\pgfqpoint{0.945723in}{0.894761in}}{\pgfqpoint{0.938031in}{0.903569in}}%
\pgfpathquadraticcurveto{\pgfqpoint{0.930358in}{0.912395in}}{\pgfqpoint{0.930358in}{0.927867in}}%
\pgfpathquadraticcurveto{\pgfqpoint{0.930358in}{0.943394in}}{\pgfqpoint{0.938031in}{0.952166in}}%
\pgfpathquadraticcurveto{\pgfqpoint{0.945723in}{0.960956in}}{\pgfqpoint{0.959286in}{0.960956in}}%
\pgfpathlineto{\pgfqpoint{0.959286in}{0.960956in}}%
\pgfpathclose%
\pgfpathmoveto{\pgfqpoint{1.041926in}{0.949770in}}%
\pgfpathquadraticcurveto{\pgfqpoint{1.040178in}{0.950779in}}{\pgfqpoint{1.038125in}{0.951247in}}%
\pgfpathquadraticcurveto{\pgfqpoint{1.036072in}{0.951733in}}{\pgfqpoint{1.033604in}{0.951733in}}%
\pgfpathquadraticcurveto{\pgfqpoint{1.024814in}{0.951733in}}{\pgfqpoint{1.020113in}{0.946023in}}%
\pgfpathquadraticcurveto{\pgfqpoint{1.015412in}{0.940314in}}{\pgfqpoint{1.015412in}{0.929614in}}%
\pgfpathlineto{\pgfqpoint{1.015412in}{0.896400in}}%
\pgfpathlineto{\pgfqpoint{1.005001in}{0.896400in}}%
\pgfpathlineto{\pgfqpoint{1.005001in}{0.959443in}}%
\pgfpathlineto{\pgfqpoint{1.015412in}{0.959443in}}%
\pgfpathlineto{\pgfqpoint{1.015412in}{0.949644in}}%
\pgfpathquadraticcurveto{\pgfqpoint{1.018690in}{0.955390in}}{\pgfqpoint{1.023913in}{0.958164in}}%
\pgfpathquadraticcurveto{\pgfqpoint{1.029155in}{0.960956in}}{\pgfqpoint{1.036648in}{0.960956in}}%
\pgfpathquadraticcurveto{\pgfqpoint{1.037711in}{0.960956in}}{\pgfqpoint{1.039008in}{0.960811in}}%
\pgfpathquadraticcurveto{\pgfqpoint{1.040304in}{0.960685in}}{\pgfqpoint{1.041872in}{0.960397in}}%
\pgfpathlineto{\pgfqpoint{1.041926in}{0.949770in}}%
\pgfpathlineto{\pgfqpoint{1.041926in}{0.949770in}}%
\pgfpathclose%
\pgfusepath{fill}%
\end{pgfscope}%
\begin{pgfscope}%
\pgfsetbuttcap%
\pgfsetroundjoin%
\definecolor{currentfill}{rgb}{0.309804,0.372549,0.419608}%
\pgfsetfillcolor{currentfill}%
\pgfsetlinewidth{0.000000pt}%
\definecolor{currentstroke}{rgb}{0.309804,0.372549,0.419608}%
\pgfsetstrokecolor{currentstroke}%
\pgfsetdash{}{0pt}%
\pgfpathmoveto{\pgfqpoint{1.151428in}{0.947338in}}%
\pgfpathquadraticcurveto{\pgfqpoint{1.155319in}{0.954327in}}{\pgfqpoint{1.160723in}{0.957641in}}%
\pgfpathquadraticcurveto{\pgfqpoint{1.166126in}{0.960956in}}{\pgfqpoint{1.173439in}{0.960956in}}%
\pgfpathquadraticcurveto{\pgfqpoint{1.183292in}{0.960956in}}{\pgfqpoint{1.188641in}{0.954057in}}%
\pgfpathquadraticcurveto{\pgfqpoint{1.193991in}{0.947176in}}{\pgfqpoint{1.193991in}{0.934460in}}%
\pgfpathlineto{\pgfqpoint{1.193991in}{0.896400in}}%
\pgfpathlineto{\pgfqpoint{1.183580in}{0.896400in}}%
\pgfpathlineto{\pgfqpoint{1.183580in}{0.934117in}}%
\pgfpathquadraticcurveto{\pgfqpoint{1.183580in}{0.943178in}}{\pgfqpoint{1.180356in}{0.947555in}}%
\pgfpathquadraticcurveto{\pgfqpoint{1.177150in}{0.951949in}}{\pgfqpoint{1.170575in}{0.951949in}}%
\pgfpathquadraticcurveto{\pgfqpoint{1.162524in}{0.951949in}}{\pgfqpoint{1.157841in}{0.946600in}}%
\pgfpathquadraticcurveto{\pgfqpoint{1.153175in}{0.941268in}}{\pgfqpoint{1.153175in}{0.932028in}}%
\pgfpathlineto{\pgfqpoint{1.153175in}{0.896400in}}%
\pgfpathlineto{\pgfqpoint{1.142764in}{0.896400in}}%
\pgfpathlineto{\pgfqpoint{1.142764in}{0.934117in}}%
\pgfpathquadraticcurveto{\pgfqpoint{1.142764in}{0.943232in}}{\pgfqpoint{1.139558in}{0.947591in}}%
\pgfpathquadraticcurveto{\pgfqpoint{1.136352in}{0.951949in}}{\pgfqpoint{1.129652in}{0.951949in}}%
\pgfpathquadraticcurveto{\pgfqpoint{1.121708in}{0.951949in}}{\pgfqpoint{1.117025in}{0.946582in}}%
\pgfpathquadraticcurveto{\pgfqpoint{1.112360in}{0.941214in}}{\pgfqpoint{1.112360in}{0.932028in}}%
\pgfpathlineto{\pgfqpoint{1.112360in}{0.896400in}}%
\pgfpathlineto{\pgfqpoint{1.101949in}{0.896400in}}%
\pgfpathlineto{\pgfqpoint{1.101949in}{0.959443in}}%
\pgfpathlineto{\pgfqpoint{1.112360in}{0.959443in}}%
\pgfpathlineto{\pgfqpoint{1.112360in}{0.949644in}}%
\pgfpathquadraticcurveto{\pgfqpoint{1.115908in}{0.955444in}}{\pgfqpoint{1.120862in}{0.958200in}}%
\pgfpathquadraticcurveto{\pgfqpoint{1.125815in}{0.960956in}}{\pgfqpoint{1.132624in}{0.960956in}}%
\pgfpathquadraticcurveto{\pgfqpoint{1.139504in}{0.960956in}}{\pgfqpoint{1.144313in}{0.957461in}}%
\pgfpathquadraticcurveto{\pgfqpoint{1.149123in}{0.953985in}}{\pgfqpoint{1.151428in}{0.947338in}}%
\pgfpathlineto{\pgfqpoint{1.151428in}{0.947338in}}%
\pgfpathclose%
\pgfpathmoveto{\pgfqpoint{1.243290in}{0.928083in}}%
\pgfpathquadraticcurveto{\pgfqpoint{1.230736in}{0.928083in}}{\pgfqpoint{1.225890in}{0.925219in}}%
\pgfpathquadraticcurveto{\pgfqpoint{1.221045in}{0.922356in}}{\pgfqpoint{1.221045in}{0.915421in}}%
\pgfpathquadraticcurveto{\pgfqpoint{1.221045in}{0.909909in}}{\pgfqpoint{1.224684in}{0.906667in}}%
\pgfpathquadraticcurveto{\pgfqpoint{1.228322in}{0.903443in}}{\pgfqpoint{1.234554in}{0.903443in}}%
\pgfpathquadraticcurveto{\pgfqpoint{1.243182in}{0.903443in}}{\pgfqpoint{1.248388in}{0.909549in}}%
\pgfpathquadraticcurveto{\pgfqpoint{1.253593in}{0.915655in}}{\pgfqpoint{1.253593in}{0.925778in}}%
\pgfpathlineto{\pgfqpoint{1.253593in}{0.928083in}}%
\pgfpathlineto{\pgfqpoint{1.243290in}{0.928083in}}%
\pgfpathlineto{\pgfqpoint{1.243290in}{0.928083in}}%
\pgfpathclose%
\pgfpathmoveto{\pgfqpoint{1.263950in}{0.932370in}}%
\pgfpathlineto{\pgfqpoint{1.263950in}{0.896400in}}%
\pgfpathlineto{\pgfqpoint{1.253593in}{0.896400in}}%
\pgfpathlineto{\pgfqpoint{1.253593in}{0.905964in}}%
\pgfpathquadraticcurveto{\pgfqpoint{1.250045in}{0.900237in}}{\pgfqpoint{1.244749in}{0.897499in}}%
\pgfpathquadraticcurveto{\pgfqpoint{1.239454in}{0.894761in}}{\pgfqpoint{1.231798in}{0.894761in}}%
\pgfpathquadraticcurveto{\pgfqpoint{1.222126in}{0.894761in}}{\pgfqpoint{1.216398in}{0.900201in}}%
\pgfpathquadraticcurveto{\pgfqpoint{1.210688in}{0.905640in}}{\pgfqpoint{1.210688in}{0.914754in}}%
\pgfpathquadraticcurveto{\pgfqpoint{1.210688in}{0.925382in}}{\pgfqpoint{1.217803in}{0.930785in}}%
\pgfpathquadraticcurveto{\pgfqpoint{1.224936in}{0.936189in}}{\pgfqpoint{1.239057in}{0.936189in}}%
\pgfpathlineto{\pgfqpoint{1.253593in}{0.936189in}}%
\pgfpathlineto{\pgfqpoint{1.253593in}{0.937216in}}%
\pgfpathquadraticcurveto{\pgfqpoint{1.253593in}{0.944366in}}{\pgfqpoint{1.248892in}{0.948275in}}%
\pgfpathquadraticcurveto{\pgfqpoint{1.244191in}{0.952184in}}{\pgfqpoint{1.235689in}{0.952184in}}%
\pgfpathquadraticcurveto{\pgfqpoint{1.230285in}{0.952184in}}{\pgfqpoint{1.225152in}{0.950887in}}%
\pgfpathquadraticcurveto{\pgfqpoint{1.220037in}{0.949590in}}{\pgfqpoint{1.215317in}{0.946996in}}%
\pgfpathlineto{\pgfqpoint{1.215317in}{0.956579in}}%
\pgfpathquadraticcurveto{\pgfqpoint{1.220991in}{0.958776in}}{\pgfqpoint{1.226341in}{0.959857in}}%
\pgfpathquadraticcurveto{\pgfqpoint{1.231690in}{0.960956in}}{\pgfqpoint{1.236752in}{0.960956in}}%
\pgfpathquadraticcurveto{\pgfqpoint{1.250441in}{0.960956in}}{\pgfqpoint{1.257196in}{0.953859in}}%
\pgfpathquadraticcurveto{\pgfqpoint{1.263950in}{0.946780in}}{\pgfqpoint{1.263950in}{0.932370in}}%
\pgfpathlineto{\pgfqpoint{1.263950in}{0.932370in}}%
\pgfpathclose%
\pgfpathmoveto{\pgfqpoint{1.337692in}{0.959443in}}%
\pgfpathlineto{\pgfqpoint{1.314889in}{0.928768in}}%
\pgfpathlineto{\pgfqpoint{1.338863in}{0.896400in}}%
\pgfpathlineto{\pgfqpoint{1.326650in}{0.896400in}}%
\pgfpathlineto{\pgfqpoint{1.308296in}{0.921167in}}%
\pgfpathlineto{\pgfqpoint{1.289960in}{0.896400in}}%
\pgfpathlineto{\pgfqpoint{1.277729in}{0.896400in}}%
\pgfpathlineto{\pgfqpoint{1.302226in}{0.929380in}}%
\pgfpathlineto{\pgfqpoint{1.279819in}{0.959443in}}%
\pgfpathlineto{\pgfqpoint{1.292031in}{0.959443in}}%
\pgfpathlineto{\pgfqpoint{1.308746in}{0.936981in}}%
\pgfpathlineto{\pgfqpoint{1.325462in}{0.959443in}}%
\pgfpathlineto{\pgfqpoint{1.337692in}{0.959443in}}%
\pgfpathlineto{\pgfqpoint{1.337692in}{0.959443in}}%
\pgfpathclose%
\pgfpathmoveto{\pgfqpoint{1.353507in}{0.959443in}}%
\pgfpathlineto{\pgfqpoint{1.363864in}{0.959443in}}%
\pgfpathlineto{\pgfqpoint{1.363864in}{0.896400in}}%
\pgfpathlineto{\pgfqpoint{1.353507in}{0.896400in}}%
\pgfpathlineto{\pgfqpoint{1.353507in}{0.959443in}}%
\pgfpathlineto{\pgfqpoint{1.353507in}{0.959443in}}%
\pgfpathclose%
\pgfpathmoveto{\pgfqpoint{1.353507in}{0.983993in}}%
\pgfpathlineto{\pgfqpoint{1.363864in}{0.983993in}}%
\pgfpathlineto{\pgfqpoint{1.363864in}{0.970862in}}%
\pgfpathlineto{\pgfqpoint{1.353507in}{0.970862in}}%
\pgfpathlineto{\pgfqpoint{1.353507in}{0.983993in}}%
\pgfpathlineto{\pgfqpoint{1.353507in}{0.983993in}}%
\pgfpathclose%
\pgfpathmoveto{\pgfqpoint{1.434615in}{0.947338in}}%
\pgfpathquadraticcurveto{\pgfqpoint{1.438506in}{0.954327in}}{\pgfqpoint{1.443910in}{0.957641in}}%
\pgfpathquadraticcurveto{\pgfqpoint{1.449313in}{0.960956in}}{\pgfqpoint{1.456626in}{0.960956in}}%
\pgfpathquadraticcurveto{\pgfqpoint{1.466479in}{0.960956in}}{\pgfqpoint{1.471828in}{0.954057in}}%
\pgfpathquadraticcurveto{\pgfqpoint{1.477178in}{0.947176in}}{\pgfqpoint{1.477178in}{0.934460in}}%
\pgfpathlineto{\pgfqpoint{1.477178in}{0.896400in}}%
\pgfpathlineto{\pgfqpoint{1.466767in}{0.896400in}}%
\pgfpathlineto{\pgfqpoint{1.466767in}{0.934117in}}%
\pgfpathquadraticcurveto{\pgfqpoint{1.466767in}{0.943178in}}{\pgfqpoint{1.463543in}{0.947555in}}%
\pgfpathquadraticcurveto{\pgfqpoint{1.460337in}{0.951949in}}{\pgfqpoint{1.453762in}{0.951949in}}%
\pgfpathquadraticcurveto{\pgfqpoint{1.445711in}{0.951949in}}{\pgfqpoint{1.441028in}{0.946600in}}%
\pgfpathquadraticcurveto{\pgfqpoint{1.436362in}{0.941268in}}{\pgfqpoint{1.436362in}{0.932028in}}%
\pgfpathlineto{\pgfqpoint{1.436362in}{0.896400in}}%
\pgfpathlineto{\pgfqpoint{1.425951in}{0.896400in}}%
\pgfpathlineto{\pgfqpoint{1.425951in}{0.934117in}}%
\pgfpathquadraticcurveto{\pgfqpoint{1.425951in}{0.943232in}}{\pgfqpoint{1.422745in}{0.947591in}}%
\pgfpathquadraticcurveto{\pgfqpoint{1.419539in}{0.951949in}}{\pgfqpoint{1.412839in}{0.951949in}}%
\pgfpathquadraticcurveto{\pgfqpoint{1.404895in}{0.951949in}}{\pgfqpoint{1.400212in}{0.946582in}}%
\pgfpathquadraticcurveto{\pgfqpoint{1.395547in}{0.941214in}}{\pgfqpoint{1.395547in}{0.932028in}}%
\pgfpathlineto{\pgfqpoint{1.395547in}{0.896400in}}%
\pgfpathlineto{\pgfqpoint{1.385136in}{0.896400in}}%
\pgfpathlineto{\pgfqpoint{1.385136in}{0.959443in}}%
\pgfpathlineto{\pgfqpoint{1.395547in}{0.959443in}}%
\pgfpathlineto{\pgfqpoint{1.395547in}{0.949644in}}%
\pgfpathquadraticcurveto{\pgfqpoint{1.399095in}{0.955444in}}{\pgfqpoint{1.404049in}{0.958200in}}%
\pgfpathquadraticcurveto{\pgfqpoint{1.409002in}{0.960956in}}{\pgfqpoint{1.415811in}{0.960956in}}%
\pgfpathquadraticcurveto{\pgfqpoint{1.422691in}{0.960956in}}{\pgfqpoint{1.427501in}{0.957461in}}%
\pgfpathquadraticcurveto{\pgfqpoint{1.432310in}{0.953985in}}{\pgfqpoint{1.434615in}{0.947338in}}%
\pgfpathlineto{\pgfqpoint{1.434615in}{0.947338in}}%
\pgfpathclose%
\pgfpathmoveto{\pgfqpoint{1.496757in}{0.921275in}}%
\pgfpathlineto{\pgfqpoint{1.496757in}{0.959443in}}%
\pgfpathlineto{\pgfqpoint{1.507114in}{0.959443in}}%
\pgfpathlineto{\pgfqpoint{1.507114in}{0.921671in}}%
\pgfpathquadraticcurveto{\pgfqpoint{1.507114in}{0.912719in}}{\pgfqpoint{1.510591in}{0.908234in}}%
\pgfpathquadraticcurveto{\pgfqpoint{1.514085in}{0.903767in}}{\pgfqpoint{1.521074in}{0.903767in}}%
\pgfpathquadraticcurveto{\pgfqpoint{1.529449in}{0.903767in}}{\pgfqpoint{1.534313in}{0.909117in}}%
\pgfpathquadraticcurveto{\pgfqpoint{1.539194in}{0.914466in}}{\pgfqpoint{1.539194in}{0.923706in}}%
\pgfpathlineto{\pgfqpoint{1.539194in}{0.959443in}}%
\pgfpathlineto{\pgfqpoint{1.549551in}{0.959443in}}%
\pgfpathlineto{\pgfqpoint{1.549551in}{0.896400in}}%
\pgfpathlineto{\pgfqpoint{1.539194in}{0.896400in}}%
\pgfpathlineto{\pgfqpoint{1.539194in}{0.906091in}}%
\pgfpathquadraticcurveto{\pgfqpoint{1.535429in}{0.900345in}}{\pgfqpoint{1.530440in}{0.897553in}}%
\pgfpathquadraticcurveto{\pgfqpoint{1.525469in}{0.894761in}}{\pgfqpoint{1.518876in}{0.894761in}}%
\pgfpathquadraticcurveto{\pgfqpoint{1.508015in}{0.894761in}}{\pgfqpoint{1.502377in}{0.901515in}}%
\pgfpathquadraticcurveto{\pgfqpoint{1.496757in}{0.908270in}}{\pgfqpoint{1.496757in}{0.921275in}}%
\pgfpathlineto{\pgfqpoint{1.496757in}{0.921275in}}%
\pgfpathclose%
\pgfpathmoveto{\pgfqpoint{1.522821in}{0.960956in}}%
\pgfpathlineto{\pgfqpoint{1.522821in}{0.960956in}}%
\pgfpathlineto{\pgfqpoint{1.522821in}{0.960956in}}%
\pgfpathclose%
\pgfpathmoveto{\pgfqpoint{1.619960in}{0.947338in}}%
\pgfpathquadraticcurveto{\pgfqpoint{1.623851in}{0.954327in}}{\pgfqpoint{1.629255in}{0.957641in}}%
\pgfpathquadraticcurveto{\pgfqpoint{1.634658in}{0.960956in}}{\pgfqpoint{1.641971in}{0.960956in}}%
\pgfpathquadraticcurveto{\pgfqpoint{1.651824in}{0.960956in}}{\pgfqpoint{1.657173in}{0.954057in}}%
\pgfpathquadraticcurveto{\pgfqpoint{1.662523in}{0.947176in}}{\pgfqpoint{1.662523in}{0.934460in}}%
\pgfpathlineto{\pgfqpoint{1.662523in}{0.896400in}}%
\pgfpathlineto{\pgfqpoint{1.652112in}{0.896400in}}%
\pgfpathlineto{\pgfqpoint{1.652112in}{0.934117in}}%
\pgfpathquadraticcurveto{\pgfqpoint{1.652112in}{0.943178in}}{\pgfqpoint{1.648888in}{0.947555in}}%
\pgfpathquadraticcurveto{\pgfqpoint{1.645682in}{0.951949in}}{\pgfqpoint{1.639107in}{0.951949in}}%
\pgfpathquadraticcurveto{\pgfqpoint{1.631056in}{0.951949in}}{\pgfqpoint{1.626373in}{0.946600in}}%
\pgfpathquadraticcurveto{\pgfqpoint{1.621708in}{0.941268in}}{\pgfqpoint{1.621708in}{0.932028in}}%
\pgfpathlineto{\pgfqpoint{1.621708in}{0.896400in}}%
\pgfpathlineto{\pgfqpoint{1.611297in}{0.896400in}}%
\pgfpathlineto{\pgfqpoint{1.611297in}{0.934117in}}%
\pgfpathquadraticcurveto{\pgfqpoint{1.611297in}{0.943232in}}{\pgfqpoint{1.608090in}{0.947591in}}%
\pgfpathquadraticcurveto{\pgfqpoint{1.604884in}{0.951949in}}{\pgfqpoint{1.598184in}{0.951949in}}%
\pgfpathquadraticcurveto{\pgfqpoint{1.590240in}{0.951949in}}{\pgfqpoint{1.585557in}{0.946582in}}%
\pgfpathquadraticcurveto{\pgfqpoint{1.580892in}{0.941214in}}{\pgfqpoint{1.580892in}{0.932028in}}%
\pgfpathlineto{\pgfqpoint{1.580892in}{0.896400in}}%
\pgfpathlineto{\pgfqpoint{1.570481in}{0.896400in}}%
\pgfpathlineto{\pgfqpoint{1.570481in}{0.959443in}}%
\pgfpathlineto{\pgfqpoint{1.580892in}{0.959443in}}%
\pgfpathlineto{\pgfqpoint{1.580892in}{0.949644in}}%
\pgfpathquadraticcurveto{\pgfqpoint{1.584440in}{0.955444in}}{\pgfqpoint{1.589394in}{0.958200in}}%
\pgfpathquadraticcurveto{\pgfqpoint{1.594347in}{0.960956in}}{\pgfqpoint{1.601156in}{0.960956in}}%
\pgfpathquadraticcurveto{\pgfqpoint{1.608036in}{0.960956in}}{\pgfqpoint{1.612846in}{0.957461in}}%
\pgfpathquadraticcurveto{\pgfqpoint{1.617655in}{0.953985in}}{\pgfqpoint{1.619960in}{0.947338in}}%
\pgfpathlineto{\pgfqpoint{1.619960in}{0.947338in}}%
\pgfpathclose%
\pgfusepath{fill}%
\end{pgfscope}%
\begin{pgfscope}%
\pgfsetbuttcap%
\pgfsetroundjoin%
\definecolor{currentfill}{rgb}{0.309804,0.372549,0.419608}%
\pgfsetfillcolor{currentfill}%
\pgfsetlinewidth{0.000000pt}%
\definecolor{currentstroke}{rgb}{0.309804,0.372549,0.419608}%
\pgfsetstrokecolor{currentstroke}%
\pgfsetdash{}{0pt}%
\pgfpathmoveto{\pgfqpoint{1.805202in}{0.957029in}}%
\pgfpathlineto{\pgfqpoint{1.805202in}{0.947338in}}%
\pgfpathquadraticcurveto{\pgfqpoint{1.800807in}{0.949770in}}{\pgfqpoint{1.796394in}{0.950977in}}%
\pgfpathquadraticcurveto{\pgfqpoint{1.791981in}{0.952184in}}{\pgfqpoint{1.787478in}{0.952184in}}%
\pgfpathquadraticcurveto{\pgfqpoint{1.777391in}{0.952184in}}{\pgfqpoint{1.771807in}{0.945789in}}%
\pgfpathquadraticcurveto{\pgfqpoint{1.766241in}{0.939413in}}{\pgfqpoint{1.766241in}{0.927867in}}%
\pgfpathquadraticcurveto{\pgfqpoint{1.766241in}{0.916321in}}{\pgfqpoint{1.771807in}{0.909927in}}%
\pgfpathquadraticcurveto{\pgfqpoint{1.777391in}{0.903551in}}{\pgfqpoint{1.787478in}{0.903551in}}%
\pgfpathquadraticcurveto{\pgfqpoint{1.791981in}{0.903551in}}{\pgfqpoint{1.796394in}{0.904758in}}%
\pgfpathquadraticcurveto{\pgfqpoint{1.800807in}{0.905964in}}{\pgfqpoint{1.805202in}{0.908396in}}%
\pgfpathlineto{\pgfqpoint{1.805202in}{0.898814in}}%
\pgfpathquadraticcurveto{\pgfqpoint{1.800861in}{0.896796in}}{\pgfqpoint{1.796213in}{0.895788in}}%
\pgfpathquadraticcurveto{\pgfqpoint{1.791584in}{0.894761in}}{\pgfqpoint{1.786343in}{0.894761in}}%
\pgfpathquadraticcurveto{\pgfqpoint{1.772095in}{0.894761in}}{\pgfqpoint{1.763702in}{0.903713in}}%
\pgfpathquadraticcurveto{\pgfqpoint{1.755326in}{0.912665in}}{\pgfqpoint{1.755326in}{0.927867in}}%
\pgfpathquadraticcurveto{\pgfqpoint{1.755326in}{0.943286in}}{\pgfqpoint{1.763792in}{0.952112in}}%
\pgfpathquadraticcurveto{\pgfqpoint{1.772275in}{0.960956in}}{\pgfqpoint{1.787027in}{0.960956in}}%
\pgfpathquadraticcurveto{\pgfqpoint{1.791800in}{0.960956in}}{\pgfqpoint{1.796358in}{0.959965in}}%
\pgfpathquadraticcurveto{\pgfqpoint{1.800915in}{0.958992in}}{\pgfqpoint{1.805202in}{0.957029in}}%
\pgfpathlineto{\pgfqpoint{1.805202in}{0.957029in}}%
\pgfpathclose%
\pgfpathmoveto{\pgfqpoint{1.823214in}{0.983993in}}%
\pgfpathlineto{\pgfqpoint{1.833571in}{0.983993in}}%
\pgfpathlineto{\pgfqpoint{1.833571in}{0.896400in}}%
\pgfpathlineto{\pgfqpoint{1.823214in}{0.896400in}}%
\pgfpathlineto{\pgfqpoint{1.823214in}{0.983993in}}%
\pgfpathlineto{\pgfqpoint{1.823214in}{0.983993in}}%
\pgfpathclose%
\pgfpathmoveto{\pgfqpoint{1.854177in}{0.921275in}}%
\pgfpathlineto{\pgfqpoint{1.854177in}{0.959443in}}%
\pgfpathlineto{\pgfqpoint{1.864534in}{0.959443in}}%
\pgfpathlineto{\pgfqpoint{1.864534in}{0.921671in}}%
\pgfpathquadraticcurveto{\pgfqpoint{1.864534in}{0.912719in}}{\pgfqpoint{1.868010in}{0.908234in}}%
\pgfpathquadraticcurveto{\pgfqpoint{1.871504in}{0.903767in}}{\pgfqpoint{1.878493in}{0.903767in}}%
\pgfpathquadraticcurveto{\pgfqpoint{1.886869in}{0.903767in}}{\pgfqpoint{1.891732in}{0.909117in}}%
\pgfpathquadraticcurveto{\pgfqpoint{1.896613in}{0.914466in}}{\pgfqpoint{1.896613in}{0.923706in}}%
\pgfpathlineto{\pgfqpoint{1.896613in}{0.959443in}}%
\pgfpathlineto{\pgfqpoint{1.906970in}{0.959443in}}%
\pgfpathlineto{\pgfqpoint{1.906970in}{0.896400in}}%
\pgfpathlineto{\pgfqpoint{1.896613in}{0.896400in}}%
\pgfpathlineto{\pgfqpoint{1.896613in}{0.906091in}}%
\pgfpathquadraticcurveto{\pgfqpoint{1.892849in}{0.900345in}}{\pgfqpoint{1.887859in}{0.897553in}}%
\pgfpathquadraticcurveto{\pgfqpoint{1.882888in}{0.894761in}}{\pgfqpoint{1.876296in}{0.894761in}}%
\pgfpathquadraticcurveto{\pgfqpoint{1.865434in}{0.894761in}}{\pgfqpoint{1.859796in}{0.901515in}}%
\pgfpathquadraticcurveto{\pgfqpoint{1.854177in}{0.908270in}}{\pgfqpoint{1.854177in}{0.921275in}}%
\pgfpathlineto{\pgfqpoint{1.854177in}{0.921275in}}%
\pgfpathclose%
\pgfpathmoveto{\pgfqpoint{1.880240in}{0.960956in}}%
\pgfpathlineto{\pgfqpoint{1.880240in}{0.960956in}}%
\pgfpathlineto{\pgfqpoint{1.880240in}{0.960956in}}%
\pgfpathclose%
\pgfpathmoveto{\pgfqpoint{1.968482in}{0.957587in}}%
\pgfpathlineto{\pgfqpoint{1.968482in}{0.947789in}}%
\pgfpathquadraticcurveto{\pgfqpoint{1.964105in}{0.950040in}}{\pgfqpoint{1.959368in}{0.951157in}}%
\pgfpathquadraticcurveto{\pgfqpoint{1.954648in}{0.952292in}}{\pgfqpoint{1.949569in}{0.952292in}}%
\pgfpathquadraticcurveto{\pgfqpoint{1.941860in}{0.952292in}}{\pgfqpoint{1.938005in}{0.949932in}}%
\pgfpathquadraticcurveto{\pgfqpoint{1.934151in}{0.947573in}}{\pgfqpoint{1.934151in}{0.942835in}}%
\pgfpathquadraticcurveto{\pgfqpoint{1.934151in}{0.939233in}}{\pgfqpoint{1.936906in}{0.937180in}}%
\pgfpathquadraticcurveto{\pgfqpoint{1.939662in}{0.935126in}}{\pgfqpoint{1.948002in}{0.933271in}}%
\pgfpathlineto{\pgfqpoint{1.951550in}{0.932478in}}%
\pgfpathquadraticcurveto{\pgfqpoint{1.962574in}{0.930119in}}{\pgfqpoint{1.967221in}{0.925814in}}%
\pgfpathquadraticcurveto{\pgfqpoint{1.971868in}{0.921509in}}{\pgfqpoint{1.971868in}{0.913800in}}%
\pgfpathquadraticcurveto{\pgfqpoint{1.971868in}{0.905010in}}{\pgfqpoint{1.964915in}{0.899876in}}%
\pgfpathquadraticcurveto{\pgfqpoint{1.957963in}{0.894761in}}{\pgfqpoint{1.945804in}{0.894761in}}%
\pgfpathquadraticcurveto{\pgfqpoint{1.940743in}{0.894761in}}{\pgfqpoint{1.935249in}{0.895752in}}%
\pgfpathquadraticcurveto{\pgfqpoint{1.929756in}{0.896742in}}{\pgfqpoint{1.923685in}{0.898706in}}%
\pgfpathlineto{\pgfqpoint{1.923685in}{0.909405in}}%
\pgfpathquadraticcurveto{\pgfqpoint{1.929431in}{0.906415in}}{\pgfqpoint{1.934997in}{0.904920in}}%
\pgfpathquadraticcurveto{\pgfqpoint{1.940563in}{0.903443in}}{\pgfqpoint{1.946039in}{0.903443in}}%
\pgfpathquadraticcurveto{\pgfqpoint{1.953351in}{0.903443in}}{\pgfqpoint{1.957278in}{0.905946in}}%
\pgfpathquadraticcurveto{\pgfqpoint{1.961223in}{0.908450in}}{\pgfqpoint{1.961223in}{0.913007in}}%
\pgfpathquadraticcurveto{\pgfqpoint{1.961223in}{0.917222in}}{\pgfqpoint{1.958377in}{0.919474in}}%
\pgfpathquadraticcurveto{\pgfqpoint{1.955549in}{0.921725in}}{\pgfqpoint{1.945912in}{0.923814in}}%
\pgfpathlineto{\pgfqpoint{1.942310in}{0.924661in}}%
\pgfpathquadraticcurveto{\pgfqpoint{1.932692in}{0.926678in}}{\pgfqpoint{1.928405in}{0.930875in}}%
\pgfpathquadraticcurveto{\pgfqpoint{1.924136in}{0.935072in}}{\pgfqpoint{1.924136in}{0.942385in}}%
\pgfpathquadraticcurveto{\pgfqpoint{1.924136in}{0.951283in}}{\pgfqpoint{1.930440in}{0.956110in}}%
\pgfpathquadraticcurveto{\pgfqpoint{1.936744in}{0.960956in}}{\pgfqpoint{1.948344in}{0.960956in}}%
\pgfpathquadraticcurveto{\pgfqpoint{1.954072in}{0.960956in}}{\pgfqpoint{1.959133in}{0.960109in}}%
\pgfpathquadraticcurveto{\pgfqpoint{1.964213in}{0.959280in}}{\pgfqpoint{1.968482in}{0.957587in}}%
\pgfpathlineto{\pgfqpoint{1.968482in}{0.957587in}}%
\pgfpathclose%
\pgfpathmoveto{\pgfqpoint{1.998598in}{0.977347in}}%
\pgfpathlineto{\pgfqpoint{1.998598in}{0.959443in}}%
\pgfpathlineto{\pgfqpoint{2.019924in}{0.959443in}}%
\pgfpathlineto{\pgfqpoint{2.019924in}{0.951391in}}%
\pgfpathlineto{\pgfqpoint{1.998598in}{0.951391in}}%
\pgfpathlineto{\pgfqpoint{1.998598in}{0.917168in}}%
\pgfpathquadraticcurveto{\pgfqpoint{1.998598in}{0.909459in}}{\pgfqpoint{2.000705in}{0.907261in}}%
\pgfpathquadraticcurveto{\pgfqpoint{2.002813in}{0.905064in}}{\pgfqpoint{2.009297in}{0.905064in}}%
\pgfpathlineto{\pgfqpoint{2.019924in}{0.905064in}}%
\pgfpathlineto{\pgfqpoint{2.019924in}{0.896400in}}%
\pgfpathlineto{\pgfqpoint{2.009297in}{0.896400in}}%
\pgfpathquadraticcurveto{\pgfqpoint{1.997301in}{0.896400in}}{\pgfqpoint{1.992744in}{0.900867in}}%
\pgfpathquadraticcurveto{\pgfqpoint{1.988187in}{0.905352in}}{\pgfqpoint{1.988187in}{0.917168in}}%
\pgfpathlineto{\pgfqpoint{1.988187in}{0.951391in}}%
\pgfpathlineto{\pgfqpoint{1.980586in}{0.951391in}}%
\pgfpathlineto{\pgfqpoint{1.980586in}{0.959443in}}%
\pgfpathlineto{\pgfqpoint{1.988187in}{0.959443in}}%
\pgfpathlineto{\pgfqpoint{1.988187in}{0.977347in}}%
\pgfpathlineto{\pgfqpoint{1.998598in}{0.977347in}}%
\pgfpathlineto{\pgfqpoint{1.998598in}{0.977347in}}%
\pgfpathclose%
\pgfpathmoveto{\pgfqpoint{2.087470in}{0.930515in}}%
\pgfpathlineto{\pgfqpoint{2.087470in}{0.925454in}}%
\pgfpathlineto{\pgfqpoint{2.039846in}{0.925454in}}%
\pgfpathquadraticcurveto{\pgfqpoint{2.040530in}{0.914754in}}{\pgfqpoint{2.046294in}{0.909153in}}%
\pgfpathquadraticcurveto{\pgfqpoint{2.052058in}{0.903551in}}{\pgfqpoint{2.062361in}{0.903551in}}%
\pgfpathquadraticcurveto{\pgfqpoint{2.068323in}{0.903551in}}{\pgfqpoint{2.073925in}{0.905010in}}%
\pgfpathquadraticcurveto{\pgfqpoint{2.079527in}{0.906469in}}{\pgfqpoint{2.085056in}{0.909405in}}%
\pgfpathlineto{\pgfqpoint{2.085056in}{0.899606in}}%
\pgfpathquadraticcurveto{\pgfqpoint{2.079473in}{0.897247in}}{\pgfqpoint{2.073619in}{0.896004in}}%
\pgfpathquadraticcurveto{\pgfqpoint{2.067765in}{0.894761in}}{\pgfqpoint{2.061749in}{0.894761in}}%
\pgfpathquadraticcurveto{\pgfqpoint{2.046654in}{0.894761in}}{\pgfqpoint{2.037846in}{0.903533in}}%
\pgfpathquadraticcurveto{\pgfqpoint{2.029039in}{0.912323in}}{\pgfqpoint{2.029039in}{0.927309in}}%
\pgfpathquadraticcurveto{\pgfqpoint{2.029039in}{0.942781in}}{\pgfqpoint{2.037396in}{0.951859in}}%
\pgfpathquadraticcurveto{\pgfqpoint{2.045754in}{0.960956in}}{\pgfqpoint{2.059947in}{0.960956in}}%
\pgfpathquadraticcurveto{\pgfqpoint{2.072664in}{0.960956in}}{\pgfqpoint{2.080067in}{0.952760in}}%
\pgfpathquadraticcurveto{\pgfqpoint{2.087470in}{0.944583in}}{\pgfqpoint{2.087470in}{0.930515in}}%
\pgfpathlineto{\pgfqpoint{2.087470in}{0.930515in}}%
\pgfpathclose%
\pgfpathmoveto{\pgfqpoint{2.077113in}{0.933559in}}%
\pgfpathquadraticcurveto{\pgfqpoint{2.077005in}{0.942043in}}{\pgfqpoint{2.072358in}{0.947104in}}%
\pgfpathquadraticcurveto{\pgfqpoint{2.067711in}{0.952184in}}{\pgfqpoint{2.060055in}{0.952184in}}%
\pgfpathquadraticcurveto{\pgfqpoint{2.051392in}{0.952184in}}{\pgfqpoint{2.046186in}{0.947284in}}%
\pgfpathquadraticcurveto{\pgfqpoint{2.040981in}{0.942385in}}{\pgfqpoint{2.040188in}{0.933487in}}%
\pgfpathlineto{\pgfqpoint{2.077113in}{0.933559in}}%
\pgfpathlineto{\pgfqpoint{2.077113in}{0.933559in}}%
\pgfpathclose%
\pgfpathmoveto{\pgfqpoint{2.141002in}{0.949770in}}%
\pgfpathquadraticcurveto{\pgfqpoint{2.139255in}{0.950779in}}{\pgfqpoint{2.137202in}{0.951247in}}%
\pgfpathquadraticcurveto{\pgfqpoint{2.135148in}{0.951733in}}{\pgfqpoint{2.132680in}{0.951733in}}%
\pgfpathquadraticcurveto{\pgfqpoint{2.123891in}{0.951733in}}{\pgfqpoint{2.119189in}{0.946023in}}%
\pgfpathquadraticcurveto{\pgfqpoint{2.114488in}{0.940314in}}{\pgfqpoint{2.114488in}{0.929614in}}%
\pgfpathlineto{\pgfqpoint{2.114488in}{0.896400in}}%
\pgfpathlineto{\pgfqpoint{2.104077in}{0.896400in}}%
\pgfpathlineto{\pgfqpoint{2.104077in}{0.959443in}}%
\pgfpathlineto{\pgfqpoint{2.114488in}{0.959443in}}%
\pgfpathlineto{\pgfqpoint{2.114488in}{0.949644in}}%
\pgfpathquadraticcurveto{\pgfqpoint{2.117766in}{0.955390in}}{\pgfqpoint{2.122990in}{0.958164in}}%
\pgfpathquadraticcurveto{\pgfqpoint{2.128231in}{0.960956in}}{\pgfqpoint{2.135725in}{0.960956in}}%
\pgfpathquadraticcurveto{\pgfqpoint{2.136787in}{0.960956in}}{\pgfqpoint{2.138084in}{0.960811in}}%
\pgfpathquadraticcurveto{\pgfqpoint{2.139381in}{0.960685in}}{\pgfqpoint{2.140948in}{0.960397in}}%
\pgfpathlineto{\pgfqpoint{2.141002in}{0.949770in}}%
\pgfpathlineto{\pgfqpoint{2.141002in}{0.949770in}}%
\pgfpathclose%
\pgfusepath{fill}%
\end{pgfscope}%
\begin{pgfscope}%
\pgfsetbuttcap%
\pgfsetroundjoin%
\definecolor{currentfill}{rgb}{0.309804,0.372549,0.419608}%
\pgfsetfillcolor{currentfill}%
\pgfsetlinewidth{0.000000pt}%
\definecolor{currentstroke}{rgb}{0.309804,0.372549,0.419608}%
\pgfsetstrokecolor{currentstroke}%
\pgfsetdash{}{0pt}%
\pgfpathmoveto{\pgfqpoint{2.217313in}{0.983993in}}%
\pgfpathlineto{\pgfqpoint{2.227670in}{0.983993in}}%
\pgfpathlineto{\pgfqpoint{2.227670in}{0.896400in}}%
\pgfpathlineto{\pgfqpoint{2.217313in}{0.896400in}}%
\pgfpathlineto{\pgfqpoint{2.217313in}{0.983993in}}%
\pgfpathlineto{\pgfqpoint{2.217313in}{0.983993in}}%
\pgfpathclose%
\pgfpathmoveto{\pgfqpoint{2.303267in}{0.930515in}}%
\pgfpathlineto{\pgfqpoint{2.303267in}{0.925454in}}%
\pgfpathlineto{\pgfqpoint{2.255643in}{0.925454in}}%
\pgfpathquadraticcurveto{\pgfqpoint{2.256327in}{0.914754in}}{\pgfqpoint{2.262091in}{0.909153in}}%
\pgfpathquadraticcurveto{\pgfqpoint{2.267855in}{0.903551in}}{\pgfqpoint{2.278158in}{0.903551in}}%
\pgfpathquadraticcurveto{\pgfqpoint{2.284120in}{0.903551in}}{\pgfqpoint{2.289722in}{0.905010in}}%
\pgfpathquadraticcurveto{\pgfqpoint{2.295323in}{0.906469in}}{\pgfqpoint{2.300853in}{0.909405in}}%
\pgfpathlineto{\pgfqpoint{2.300853in}{0.899606in}}%
\pgfpathquadraticcurveto{\pgfqpoint{2.295269in}{0.897247in}}{\pgfqpoint{2.289415in}{0.896004in}}%
\pgfpathquadraticcurveto{\pgfqpoint{2.283561in}{0.894761in}}{\pgfqpoint{2.277545in}{0.894761in}}%
\pgfpathquadraticcurveto{\pgfqpoint{2.262451in}{0.894761in}}{\pgfqpoint{2.253643in}{0.903533in}}%
\pgfpathquadraticcurveto{\pgfqpoint{2.244835in}{0.912323in}}{\pgfqpoint{2.244835in}{0.927309in}}%
\pgfpathquadraticcurveto{\pgfqpoint{2.244835in}{0.942781in}}{\pgfqpoint{2.253193in}{0.951859in}}%
\pgfpathquadraticcurveto{\pgfqpoint{2.261551in}{0.960956in}}{\pgfqpoint{2.275744in}{0.960956in}}%
\pgfpathquadraticcurveto{\pgfqpoint{2.288461in}{0.960956in}}{\pgfqpoint{2.295864in}{0.952760in}}%
\pgfpathquadraticcurveto{\pgfqpoint{2.303267in}{0.944583in}}{\pgfqpoint{2.303267in}{0.930515in}}%
\pgfpathlineto{\pgfqpoint{2.303267in}{0.930515in}}%
\pgfpathclose%
\pgfpathmoveto{\pgfqpoint{2.292910in}{0.933559in}}%
\pgfpathquadraticcurveto{\pgfqpoint{2.292802in}{0.942043in}}{\pgfqpoint{2.288155in}{0.947104in}}%
\pgfpathquadraticcurveto{\pgfqpoint{2.283507in}{0.952184in}}{\pgfqpoint{2.275852in}{0.952184in}}%
\pgfpathquadraticcurveto{\pgfqpoint{2.267188in}{0.952184in}}{\pgfqpoint{2.261983in}{0.947284in}}%
\pgfpathquadraticcurveto{\pgfqpoint{2.256777in}{0.942385in}}{\pgfqpoint{2.255985in}{0.933487in}}%
\pgfpathlineto{\pgfqpoint{2.292910in}{0.933559in}}%
\pgfpathlineto{\pgfqpoint{2.292910in}{0.933559in}}%
\pgfpathclose%
\pgfpathmoveto{\pgfqpoint{2.312849in}{0.959443in}}%
\pgfpathlineto{\pgfqpoint{2.323819in}{0.959443in}}%
\pgfpathlineto{\pgfqpoint{2.343524in}{0.906541in}}%
\pgfpathlineto{\pgfqpoint{2.363229in}{0.959443in}}%
\pgfpathlineto{\pgfqpoint{2.374199in}{0.959443in}}%
\pgfpathlineto{\pgfqpoint{2.350549in}{0.896400in}}%
\pgfpathlineto{\pgfqpoint{2.336481in}{0.896400in}}%
\pgfpathlineto{\pgfqpoint{2.312849in}{0.959443in}}%
\pgfpathlineto{\pgfqpoint{2.312849in}{0.959443in}}%
\pgfpathclose%
\pgfpathmoveto{\pgfqpoint{2.442429in}{0.930515in}}%
\pgfpathlineto{\pgfqpoint{2.442429in}{0.925454in}}%
\pgfpathlineto{\pgfqpoint{2.394804in}{0.925454in}}%
\pgfpathquadraticcurveto{\pgfqpoint{2.395489in}{0.914754in}}{\pgfqpoint{2.401253in}{0.909153in}}%
\pgfpathquadraticcurveto{\pgfqpoint{2.407017in}{0.903551in}}{\pgfqpoint{2.417320in}{0.903551in}}%
\pgfpathquadraticcurveto{\pgfqpoint{2.423282in}{0.903551in}}{\pgfqpoint{2.428883in}{0.905010in}}%
\pgfpathquadraticcurveto{\pgfqpoint{2.434485in}{0.906469in}}{\pgfqpoint{2.440015in}{0.909405in}}%
\pgfpathlineto{\pgfqpoint{2.440015in}{0.899606in}}%
\pgfpathquadraticcurveto{\pgfqpoint{2.434431in}{0.897247in}}{\pgfqpoint{2.428577in}{0.896004in}}%
\pgfpathquadraticcurveto{\pgfqpoint{2.422723in}{0.894761in}}{\pgfqpoint{2.416707in}{0.894761in}}%
\pgfpathquadraticcurveto{\pgfqpoint{2.401613in}{0.894761in}}{\pgfqpoint{2.392805in}{0.903533in}}%
\pgfpathquadraticcurveto{\pgfqpoint{2.383997in}{0.912323in}}{\pgfqpoint{2.383997in}{0.927309in}}%
\pgfpathquadraticcurveto{\pgfqpoint{2.383997in}{0.942781in}}{\pgfqpoint{2.392355in}{0.951859in}}%
\pgfpathquadraticcurveto{\pgfqpoint{2.400712in}{0.960956in}}{\pgfqpoint{2.414906in}{0.960956in}}%
\pgfpathquadraticcurveto{\pgfqpoint{2.427623in}{0.960956in}}{\pgfqpoint{2.435026in}{0.952760in}}%
\pgfpathquadraticcurveto{\pgfqpoint{2.442429in}{0.944583in}}{\pgfqpoint{2.442429in}{0.930515in}}%
\pgfpathlineto{\pgfqpoint{2.442429in}{0.930515in}}%
\pgfpathclose%
\pgfpathmoveto{\pgfqpoint{2.432072in}{0.933559in}}%
\pgfpathquadraticcurveto{\pgfqpoint{2.431964in}{0.942043in}}{\pgfqpoint{2.427316in}{0.947104in}}%
\pgfpathquadraticcurveto{\pgfqpoint{2.422669in}{0.952184in}}{\pgfqpoint{2.415014in}{0.952184in}}%
\pgfpathquadraticcurveto{\pgfqpoint{2.406350in}{0.952184in}}{\pgfqpoint{2.401145in}{0.947284in}}%
\pgfpathquadraticcurveto{\pgfqpoint{2.395939in}{0.942385in}}{\pgfqpoint{2.395147in}{0.933487in}}%
\pgfpathlineto{\pgfqpoint{2.432072in}{0.933559in}}%
\pgfpathlineto{\pgfqpoint{2.432072in}{0.933559in}}%
\pgfpathclose%
\pgfpathmoveto{\pgfqpoint{2.495961in}{0.949770in}}%
\pgfpathquadraticcurveto{\pgfqpoint{2.494214in}{0.950779in}}{\pgfqpoint{2.492160in}{0.951247in}}%
\pgfpathquadraticcurveto{\pgfqpoint{2.490107in}{0.951733in}}{\pgfqpoint{2.487639in}{0.951733in}}%
\pgfpathquadraticcurveto{\pgfqpoint{2.478849in}{0.951733in}}{\pgfqpoint{2.474148in}{0.946023in}}%
\pgfpathquadraticcurveto{\pgfqpoint{2.469447in}{0.940314in}}{\pgfqpoint{2.469447in}{0.929614in}}%
\pgfpathlineto{\pgfqpoint{2.469447in}{0.896400in}}%
\pgfpathlineto{\pgfqpoint{2.459036in}{0.896400in}}%
\pgfpathlineto{\pgfqpoint{2.459036in}{0.959443in}}%
\pgfpathlineto{\pgfqpoint{2.469447in}{0.959443in}}%
\pgfpathlineto{\pgfqpoint{2.469447in}{0.949644in}}%
\pgfpathquadraticcurveto{\pgfqpoint{2.472725in}{0.955390in}}{\pgfqpoint{2.477949in}{0.958164in}}%
\pgfpathquadraticcurveto{\pgfqpoint{2.483190in}{0.960956in}}{\pgfqpoint{2.490683in}{0.960956in}}%
\pgfpathquadraticcurveto{\pgfqpoint{2.491746in}{0.960956in}}{\pgfqpoint{2.493043in}{0.960811in}}%
\pgfpathquadraticcurveto{\pgfqpoint{2.494340in}{0.960685in}}{\pgfqpoint{2.495907in}{0.960397in}}%
\pgfpathlineto{\pgfqpoint{2.495961in}{0.949770in}}%
\pgfpathlineto{\pgfqpoint{2.495961in}{0.949770in}}%
\pgfpathclose%
\pgfpathmoveto{\pgfqpoint{2.535479in}{0.928083in}}%
\pgfpathquadraticcurveto{\pgfqpoint{2.522925in}{0.928083in}}{\pgfqpoint{2.518080in}{0.925219in}}%
\pgfpathquadraticcurveto{\pgfqpoint{2.513234in}{0.922356in}}{\pgfqpoint{2.513234in}{0.915421in}}%
\pgfpathquadraticcurveto{\pgfqpoint{2.513234in}{0.909909in}}{\pgfqpoint{2.516873in}{0.906667in}}%
\pgfpathquadraticcurveto{\pgfqpoint{2.520511in}{0.903443in}}{\pgfqpoint{2.526743in}{0.903443in}}%
\pgfpathquadraticcurveto{\pgfqpoint{2.535371in}{0.903443in}}{\pgfqpoint{2.540577in}{0.909549in}}%
\pgfpathquadraticcurveto{\pgfqpoint{2.545782in}{0.915655in}}{\pgfqpoint{2.545782in}{0.925778in}}%
\pgfpathlineto{\pgfqpoint{2.545782in}{0.928083in}}%
\pgfpathlineto{\pgfqpoint{2.535479in}{0.928083in}}%
\pgfpathlineto{\pgfqpoint{2.535479in}{0.928083in}}%
\pgfpathclose%
\pgfpathmoveto{\pgfqpoint{2.556139in}{0.932370in}}%
\pgfpathlineto{\pgfqpoint{2.556139in}{0.896400in}}%
\pgfpathlineto{\pgfqpoint{2.545782in}{0.896400in}}%
\pgfpathlineto{\pgfqpoint{2.545782in}{0.905964in}}%
\pgfpathquadraticcurveto{\pgfqpoint{2.542234in}{0.900237in}}{\pgfqpoint{2.536938in}{0.897499in}}%
\pgfpathquadraticcurveto{\pgfqpoint{2.531643in}{0.894761in}}{\pgfqpoint{2.523988in}{0.894761in}}%
\pgfpathquadraticcurveto{\pgfqpoint{2.514315in}{0.894761in}}{\pgfqpoint{2.508587in}{0.900201in}}%
\pgfpathquadraticcurveto{\pgfqpoint{2.502877in}{0.905640in}}{\pgfqpoint{2.502877in}{0.914754in}}%
\pgfpathquadraticcurveto{\pgfqpoint{2.502877in}{0.925382in}}{\pgfqpoint{2.509992in}{0.930785in}}%
\pgfpathquadraticcurveto{\pgfqpoint{2.517125in}{0.936189in}}{\pgfqpoint{2.531247in}{0.936189in}}%
\pgfpathlineto{\pgfqpoint{2.545782in}{0.936189in}}%
\pgfpathlineto{\pgfqpoint{2.545782in}{0.937216in}}%
\pgfpathquadraticcurveto{\pgfqpoint{2.545782in}{0.944366in}}{\pgfqpoint{2.541081in}{0.948275in}}%
\pgfpathquadraticcurveto{\pgfqpoint{2.536380in}{0.952184in}}{\pgfqpoint{2.527878in}{0.952184in}}%
\pgfpathquadraticcurveto{\pgfqpoint{2.522475in}{0.952184in}}{\pgfqpoint{2.517341in}{0.950887in}}%
\pgfpathquadraticcurveto{\pgfqpoint{2.512226in}{0.949590in}}{\pgfqpoint{2.507507in}{0.946996in}}%
\pgfpathlineto{\pgfqpoint{2.507507in}{0.956579in}}%
\pgfpathquadraticcurveto{\pgfqpoint{2.513180in}{0.958776in}}{\pgfqpoint{2.518530in}{0.959857in}}%
\pgfpathquadraticcurveto{\pgfqpoint{2.523880in}{0.960956in}}{\pgfqpoint{2.528941in}{0.960956in}}%
\pgfpathquadraticcurveto{\pgfqpoint{2.542630in}{0.960956in}}{\pgfqpoint{2.549385in}{0.953859in}}%
\pgfpathquadraticcurveto{\pgfqpoint{2.556139in}{0.946780in}}{\pgfqpoint{2.556139in}{0.932370in}}%
\pgfpathlineto{\pgfqpoint{2.556139in}{0.932370in}}%
\pgfpathclose%
\pgfpathmoveto{\pgfqpoint{2.618948in}{0.928660in}}%
\pgfpathquadraticcurveto{\pgfqpoint{2.618948in}{0.939917in}}{\pgfqpoint{2.614301in}{0.946096in}}%
\pgfpathquadraticcurveto{\pgfqpoint{2.609671in}{0.952292in}}{\pgfqpoint{2.601278in}{0.952292in}}%
\pgfpathquadraticcurveto{\pgfqpoint{2.592956in}{0.952292in}}{\pgfqpoint{2.588309in}{0.946096in}}%
\pgfpathquadraticcurveto{\pgfqpoint{2.583662in}{0.939917in}}{\pgfqpoint{2.583662in}{0.928660in}}%
\pgfpathquadraticcurveto{\pgfqpoint{2.583662in}{0.917456in}}{\pgfqpoint{2.588309in}{0.911260in}}%
\pgfpathquadraticcurveto{\pgfqpoint{2.592956in}{0.905064in}}{\pgfqpoint{2.601278in}{0.905064in}}%
\pgfpathquadraticcurveto{\pgfqpoint{2.609671in}{0.905064in}}{\pgfqpoint{2.614301in}{0.911260in}}%
\pgfpathquadraticcurveto{\pgfqpoint{2.618948in}{0.917456in}}{\pgfqpoint{2.618948in}{0.928660in}}%
\pgfpathlineto{\pgfqpoint{2.618948in}{0.928660in}}%
\pgfpathclose%
\pgfpathmoveto{\pgfqpoint{2.629305in}{0.904217in}}%
\pgfpathquadraticcurveto{\pgfqpoint{2.629305in}{0.888132in}}{\pgfqpoint{2.622154in}{0.880279in}}%
\pgfpathquadraticcurveto{\pgfqpoint{2.615021in}{0.872426in}}{\pgfqpoint{2.600269in}{0.872426in}}%
\pgfpathquadraticcurveto{\pgfqpoint{2.594811in}{0.872426in}}{\pgfqpoint{2.589966in}{0.873236in}}%
\pgfpathquadraticcurveto{\pgfqpoint{2.585121in}{0.874047in}}{\pgfqpoint{2.580564in}{0.875740in}}%
\pgfpathlineto{\pgfqpoint{2.580564in}{0.885809in}}%
\pgfpathquadraticcurveto{\pgfqpoint{2.585121in}{0.883341in}}{\pgfqpoint{2.589570in}{0.882170in}}%
\pgfpathquadraticcurveto{\pgfqpoint{2.594019in}{0.880982in}}{\pgfqpoint{2.598630in}{0.880982in}}%
\pgfpathquadraticcurveto{\pgfqpoint{2.608825in}{0.880982in}}{\pgfqpoint{2.613886in}{0.886295in}}%
\pgfpathquadraticcurveto{\pgfqpoint{2.618948in}{0.891609in}}{\pgfqpoint{2.618948in}{0.902362in}}%
\pgfpathlineto{\pgfqpoint{2.618948in}{0.907495in}}%
\pgfpathquadraticcurveto{\pgfqpoint{2.615742in}{0.901912in}}{\pgfqpoint{2.610734in}{0.899156in}}%
\pgfpathquadraticcurveto{\pgfqpoint{2.605727in}{0.896400in}}{\pgfqpoint{2.598738in}{0.896400in}}%
\pgfpathquadraticcurveto{\pgfqpoint{2.587156in}{0.896400in}}{\pgfqpoint{2.580059in}{0.905226in}}%
\pgfpathquadraticcurveto{\pgfqpoint{2.572963in}{0.914070in}}{\pgfqpoint{2.572963in}{0.928660in}}%
\pgfpathquadraticcurveto{\pgfqpoint{2.572963in}{0.943286in}}{\pgfqpoint{2.580059in}{0.952112in}}%
\pgfpathquadraticcurveto{\pgfqpoint{2.587156in}{0.960956in}}{\pgfqpoint{2.598738in}{0.960956in}}%
\pgfpathquadraticcurveto{\pgfqpoint{2.605727in}{0.960956in}}{\pgfqpoint{2.610734in}{0.958200in}}%
\pgfpathquadraticcurveto{\pgfqpoint{2.615742in}{0.955444in}}{\pgfqpoint{2.618948in}{0.949878in}}%
\pgfpathlineto{\pgfqpoint{2.618948in}{0.959443in}}%
\pgfpathlineto{\pgfqpoint{2.629305in}{0.959443in}}%
\pgfpathlineto{\pgfqpoint{2.629305in}{0.904217in}}%
\pgfpathlineto{\pgfqpoint{2.629305in}{0.904217in}}%
\pgfpathclose%
\pgfpathmoveto{\pgfqpoint{2.704577in}{0.930515in}}%
\pgfpathlineto{\pgfqpoint{2.704577in}{0.925454in}}%
\pgfpathlineto{\pgfqpoint{2.656953in}{0.925454in}}%
\pgfpathquadraticcurveto{\pgfqpoint{2.657638in}{0.914754in}}{\pgfqpoint{2.663402in}{0.909153in}}%
\pgfpathquadraticcurveto{\pgfqpoint{2.669166in}{0.903551in}}{\pgfqpoint{2.679469in}{0.903551in}}%
\pgfpathquadraticcurveto{\pgfqpoint{2.685431in}{0.903551in}}{\pgfqpoint{2.691032in}{0.905010in}}%
\pgfpathquadraticcurveto{\pgfqpoint{2.696634in}{0.906469in}}{\pgfqpoint{2.702164in}{0.909405in}}%
\pgfpathlineto{\pgfqpoint{2.702164in}{0.899606in}}%
\pgfpathquadraticcurveto{\pgfqpoint{2.696580in}{0.897247in}}{\pgfqpoint{2.690726in}{0.896004in}}%
\pgfpathquadraticcurveto{\pgfqpoint{2.684872in}{0.894761in}}{\pgfqpoint{2.678856in}{0.894761in}}%
\pgfpathquadraticcurveto{\pgfqpoint{2.663762in}{0.894761in}}{\pgfqpoint{2.654954in}{0.903533in}}%
\pgfpathquadraticcurveto{\pgfqpoint{2.646146in}{0.912323in}}{\pgfqpoint{2.646146in}{0.927309in}}%
\pgfpathquadraticcurveto{\pgfqpoint{2.646146in}{0.942781in}}{\pgfqpoint{2.654504in}{0.951859in}}%
\pgfpathquadraticcurveto{\pgfqpoint{2.662861in}{0.960956in}}{\pgfqpoint{2.677055in}{0.960956in}}%
\pgfpathquadraticcurveto{\pgfqpoint{2.689771in}{0.960956in}}{\pgfqpoint{2.697174in}{0.952760in}}%
\pgfpathquadraticcurveto{\pgfqpoint{2.704577in}{0.944583in}}{\pgfqpoint{2.704577in}{0.930515in}}%
\pgfpathlineto{\pgfqpoint{2.704577in}{0.930515in}}%
\pgfpathclose%
\pgfpathmoveto{\pgfqpoint{2.694220in}{0.933559in}}%
\pgfpathquadraticcurveto{\pgfqpoint{2.694112in}{0.942043in}}{\pgfqpoint{2.689465in}{0.947104in}}%
\pgfpathquadraticcurveto{\pgfqpoint{2.684818in}{0.952184in}}{\pgfqpoint{2.677163in}{0.952184in}}%
\pgfpathquadraticcurveto{\pgfqpoint{2.668499in}{0.952184in}}{\pgfqpoint{2.663294in}{0.947284in}}%
\pgfpathquadraticcurveto{\pgfqpoint{2.658088in}{0.942385in}}{\pgfqpoint{2.657296in}{0.933487in}}%
\pgfpathlineto{\pgfqpoint{2.694220in}{0.933559in}}%
\pgfpathlineto{\pgfqpoint{2.694220in}{0.933559in}}%
\pgfpathclose%
\pgfusepath{fill}%
\end{pgfscope}%
\begin{pgfscope}%
\pgfsetbuttcap%
\pgfsetroundjoin%
\definecolor{currentfill}{rgb}{0.309804,0.372549,0.419608}%
\pgfsetfillcolor{currentfill}%
\pgfsetlinewidth{0.000000pt}%
\definecolor{currentstroke}{rgb}{0.309804,0.372549,0.419608}%
\pgfsetstrokecolor{currentstroke}%
\pgfsetdash{}{0pt}%
\pgfpathmoveto{\pgfqpoint{2.776147in}{0.953138in}}%
\pgfpathlineto{\pgfqpoint{2.776147in}{0.963387in}}%
\pgfpathlineto{\pgfqpoint{2.848304in}{0.937216in}}%
\pgfpathlineto{\pgfqpoint{2.848304in}{0.927867in}}%
\pgfpathlineto{\pgfqpoint{2.776147in}{0.901696in}}%
\pgfpathlineto{\pgfqpoint{2.776147in}{0.911944in}}%
\pgfpathlineto{\pgfqpoint{2.834129in}{0.932478in}}%
\pgfpathlineto{\pgfqpoint{2.776147in}{0.953138in}}%
\pgfpathlineto{\pgfqpoint{2.776147in}{0.953138in}}%
\pgfpathclose%
\pgfpathmoveto{\pgfqpoint{2.872981in}{0.980445in}}%
\pgfpathlineto{\pgfqpoint{2.917615in}{0.980445in}}%
\pgfpathlineto{\pgfqpoint{2.917615in}{0.970862in}}%
\pgfpathlineto{\pgfqpoint{2.883392in}{0.970862in}}%
\pgfpathlineto{\pgfqpoint{2.883392in}{0.950274in}}%
\pgfpathquadraticcurveto{\pgfqpoint{2.885859in}{0.951121in}}{\pgfqpoint{2.888327in}{0.951535in}}%
\pgfpathquadraticcurveto{\pgfqpoint{2.890813in}{0.951949in}}{\pgfqpoint{2.893298in}{0.951949in}}%
\pgfpathquadraticcurveto{\pgfqpoint{2.907366in}{0.951949in}}{\pgfqpoint{2.915579in}{0.944240in}}%
\pgfpathquadraticcurveto{\pgfqpoint{2.923811in}{0.936531in}}{\pgfqpoint{2.923811in}{0.923364in}}%
\pgfpathquadraticcurveto{\pgfqpoint{2.923811in}{0.909801in}}{\pgfqpoint{2.915363in}{0.902272in}}%
\pgfpathquadraticcurveto{\pgfqpoint{2.906916in}{0.894761in}}{\pgfqpoint{2.891551in}{0.894761in}}%
\pgfpathquadraticcurveto{\pgfqpoint{2.886256in}{0.894761in}}{\pgfqpoint{2.880762in}{0.895662in}}%
\pgfpathquadraticcurveto{\pgfqpoint{2.875286in}{0.896562in}}{\pgfqpoint{2.869432in}{0.898363in}}%
\pgfpathlineto{\pgfqpoint{2.869432in}{0.909801in}}%
\pgfpathquadraticcurveto{\pgfqpoint{2.874494in}{0.907045in}}{\pgfqpoint{2.879897in}{0.905694in}}%
\pgfpathquadraticcurveto{\pgfqpoint{2.885301in}{0.904343in}}{\pgfqpoint{2.891317in}{0.904343in}}%
\pgfpathquadraticcurveto{\pgfqpoint{2.901062in}{0.904343in}}{\pgfqpoint{2.906736in}{0.909459in}}%
\pgfpathquadraticcurveto{\pgfqpoint{2.912427in}{0.914574in}}{\pgfqpoint{2.912427in}{0.923364in}}%
\pgfpathquadraticcurveto{\pgfqpoint{2.912427in}{0.932136in}}{\pgfqpoint{2.906736in}{0.937252in}}%
\pgfpathquadraticcurveto{\pgfqpoint{2.901062in}{0.942385in}}{\pgfqpoint{2.891317in}{0.942385in}}%
\pgfpathquadraticcurveto{\pgfqpoint{2.886760in}{0.942385in}}{\pgfqpoint{2.882221in}{0.941376in}}%
\pgfpathquadraticcurveto{\pgfqpoint{2.877700in}{0.940368in}}{\pgfqpoint{2.872981in}{0.938224in}}%
\pgfpathlineto{\pgfqpoint{2.872981in}{0.980445in}}%
\pgfpathlineto{\pgfqpoint{2.872981in}{0.980445in}}%
\pgfpathclose%
\pgfpathmoveto{\pgfqpoint{2.970517in}{0.972952in}}%
\pgfpathquadraticcurveto{\pgfqpoint{2.961745in}{0.972952in}}{\pgfqpoint{2.957314in}{0.964306in}}%
\pgfpathquadraticcurveto{\pgfqpoint{2.952901in}{0.955678in}}{\pgfqpoint{2.952901in}{0.938332in}}%
\pgfpathquadraticcurveto{\pgfqpoint{2.952901in}{0.921059in}}{\pgfqpoint{2.957314in}{0.912413in}}%
\pgfpathquadraticcurveto{\pgfqpoint{2.961745in}{0.903767in}}{\pgfqpoint{2.970517in}{0.903767in}}%
\pgfpathquadraticcurveto{\pgfqpoint{2.979361in}{0.903767in}}{\pgfqpoint{2.983774in}{0.912413in}}%
\pgfpathquadraticcurveto{\pgfqpoint{2.988204in}{0.921059in}}{\pgfqpoint{2.988204in}{0.938332in}}%
\pgfpathquadraticcurveto{\pgfqpoint{2.988204in}{0.955678in}}{\pgfqpoint{2.983774in}{0.964306in}}%
\pgfpathquadraticcurveto{\pgfqpoint{2.979361in}{0.972952in}}{\pgfqpoint{2.970517in}{0.972952in}}%
\pgfpathlineto{\pgfqpoint{2.970517in}{0.972952in}}%
\pgfpathclose%
\pgfpathmoveto{\pgfqpoint{2.970517in}{0.981958in}}%
\pgfpathquadraticcurveto{\pgfqpoint{2.984656in}{0.981958in}}{\pgfqpoint{2.992113in}{0.970772in}}%
\pgfpathquadraticcurveto{\pgfqpoint{2.999570in}{0.959605in}}{\pgfqpoint{2.999570in}{0.938332in}}%
\pgfpathquadraticcurveto{\pgfqpoint{2.999570in}{0.917114in}}{\pgfqpoint{2.992113in}{0.905928in}}%
\pgfpathquadraticcurveto{\pgfqpoint{2.984656in}{0.894761in}}{\pgfqpoint{2.970517in}{0.894761in}}%
\pgfpathquadraticcurveto{\pgfqpoint{2.956395in}{0.894761in}}{\pgfqpoint{2.948938in}{0.905928in}}%
\pgfpathquadraticcurveto{\pgfqpoint{2.941481in}{0.917114in}}{\pgfqpoint{2.941481in}{0.938332in}}%
\pgfpathquadraticcurveto{\pgfqpoint{2.941481in}{0.959605in}}{\pgfqpoint{2.948938in}{0.970772in}}%
\pgfpathquadraticcurveto{\pgfqpoint{2.956395in}{0.981958in}}{\pgfqpoint{2.970517in}{0.981958in}}%
\pgfpathlineto{\pgfqpoint{2.970517in}{0.981958in}}%
\pgfpathclose%
\pgfpathmoveto{\pgfqpoint{3.091036in}{0.933379in}}%
\pgfpathquadraticcurveto{\pgfqpoint{3.086137in}{0.933379in}}{\pgfqpoint{3.083345in}{0.929218in}}%
\pgfpathquadraticcurveto{\pgfqpoint{3.080571in}{0.925057in}}{\pgfqpoint{3.080571in}{0.917618in}}%
\pgfpathquadraticcurveto{\pgfqpoint{3.080571in}{0.910305in}}{\pgfqpoint{3.083345in}{0.906109in}}%
\pgfpathquadraticcurveto{\pgfqpoint{3.086137in}{0.901912in}}{\pgfqpoint{3.091036in}{0.901912in}}%
\pgfpathquadraticcurveto{\pgfqpoint{3.095827in}{0.901912in}}{\pgfqpoint{3.098601in}{0.906109in}}%
\pgfpathquadraticcurveto{\pgfqpoint{3.101393in}{0.910305in}}{\pgfqpoint{3.101393in}{0.917618in}}%
\pgfpathquadraticcurveto{\pgfqpoint{3.101393in}{0.925003in}}{\pgfqpoint{3.098601in}{0.929182in}}%
\pgfpathquadraticcurveto{\pgfqpoint{3.095827in}{0.933379in}}{\pgfqpoint{3.091036in}{0.933379in}}%
\pgfpathlineto{\pgfqpoint{3.091036in}{0.933379in}}%
\pgfpathclose%
\pgfpathmoveto{\pgfqpoint{3.091036in}{0.940530in}}%
\pgfpathquadraticcurveto{\pgfqpoint{3.099934in}{0.940530in}}{\pgfqpoint{3.105157in}{0.934334in}}%
\pgfpathquadraticcurveto{\pgfqpoint{3.110399in}{0.928155in}}{\pgfqpoint{3.110399in}{0.917618in}}%
\pgfpathquadraticcurveto{\pgfqpoint{3.110399in}{0.907099in}}{\pgfqpoint{3.105139in}{0.900921in}}%
\pgfpathquadraticcurveto{\pgfqpoint{3.099880in}{0.894761in}}{\pgfqpoint{3.091036in}{0.894761in}}%
\pgfpathquadraticcurveto{\pgfqpoint{3.082030in}{0.894761in}}{\pgfqpoint{3.076788in}{0.900921in}}%
\pgfpathquadraticcurveto{\pgfqpoint{3.071565in}{0.907099in}}{\pgfqpoint{3.071565in}{0.917618in}}%
\pgfpathquadraticcurveto{\pgfqpoint{3.071565in}{0.928209in}}{\pgfqpoint{3.076824in}{0.934370in}}%
\pgfpathquadraticcurveto{\pgfqpoint{3.082084in}{0.940530in}}{\pgfqpoint{3.091036in}{0.940530in}}%
\pgfpathlineto{\pgfqpoint{3.091036in}{0.940530in}}%
\pgfpathclose%
\pgfpathmoveto{\pgfqpoint{3.032947in}{0.974807in}}%
\pgfpathquadraticcurveto{\pgfqpoint{3.028101in}{0.974807in}}{\pgfqpoint{3.025310in}{0.970610in}}%
\pgfpathquadraticcurveto{\pgfqpoint{3.022536in}{0.966431in}}{\pgfqpoint{3.022536in}{0.959100in}}%
\pgfpathquadraticcurveto{\pgfqpoint{3.022536in}{0.951679in}}{\pgfqpoint{3.025292in}{0.947500in}}%
\pgfpathquadraticcurveto{\pgfqpoint{3.028047in}{0.943340in}}{\pgfqpoint{3.032947in}{0.943340in}}%
\pgfpathquadraticcurveto{\pgfqpoint{3.037846in}{0.943340in}}{\pgfqpoint{3.040620in}{0.947500in}}%
\pgfpathquadraticcurveto{\pgfqpoint{3.043412in}{0.951679in}}{\pgfqpoint{3.043412in}{0.959100in}}%
\pgfpathquadraticcurveto{\pgfqpoint{3.043412in}{0.966359in}}{\pgfqpoint{3.040602in}{0.970574in}}%
\pgfpathquadraticcurveto{\pgfqpoint{3.037792in}{0.974807in}}{\pgfqpoint{3.032947in}{0.974807in}}%
\pgfpathlineto{\pgfqpoint{3.032947in}{0.974807in}}%
\pgfpathclose%
\pgfpathmoveto{\pgfqpoint{3.083777in}{0.981958in}}%
\pgfpathlineto{\pgfqpoint{3.092783in}{0.981958in}}%
\pgfpathlineto{\pgfqpoint{3.040206in}{0.894761in}}%
\pgfpathlineto{\pgfqpoint{3.031200in}{0.894761in}}%
\pgfpathlineto{\pgfqpoint{3.083777in}{0.981958in}}%
\pgfpathlineto{\pgfqpoint{3.083777in}{0.981958in}}%
\pgfpathclose%
\pgfpathmoveto{\pgfqpoint{3.032947in}{0.981958in}}%
\pgfpathquadraticcurveto{\pgfqpoint{3.041845in}{0.981958in}}{\pgfqpoint{3.047122in}{0.975798in}}%
\pgfpathquadraticcurveto{\pgfqpoint{3.052418in}{0.969637in}}{\pgfqpoint{3.052418in}{0.959100in}}%
\pgfpathquadraticcurveto{\pgfqpoint{3.052418in}{0.948473in}}{\pgfqpoint{3.047158in}{0.942331in}}%
\pgfpathquadraticcurveto{\pgfqpoint{3.041899in}{0.936189in}}{\pgfqpoint{3.032947in}{0.936189in}}%
\pgfpathquadraticcurveto{\pgfqpoint{3.023995in}{0.936189in}}{\pgfqpoint{3.018789in}{0.942349in}}%
\pgfpathquadraticcurveto{\pgfqpoint{3.013584in}{0.948527in}}{\pgfqpoint{3.013584in}{0.959100in}}%
\pgfpathquadraticcurveto{\pgfqpoint{3.013584in}{0.969583in}}{\pgfqpoint{3.018807in}{0.975762in}}%
\pgfpathquadraticcurveto{\pgfqpoint{3.024049in}{0.981958in}}{\pgfqpoint{3.032947in}{0.981958in}}%
\pgfpathlineto{\pgfqpoint{3.032947in}{0.981958in}}%
\pgfpathclose%
\pgfpathmoveto{\pgfqpoint{3.129078in}{0.910702in}}%
\pgfpathlineto{\pgfqpoint{3.140966in}{0.910702in}}%
\pgfpathlineto{\pgfqpoint{3.140966in}{0.896400in}}%
\pgfpathlineto{\pgfqpoint{3.129078in}{0.896400in}}%
\pgfpathlineto{\pgfqpoint{3.129078in}{0.910702in}}%
\pgfpathlineto{\pgfqpoint{3.129078in}{0.910702in}}%
\pgfpathclose%
\pgfusepath{fill}%
\end{pgfscope}%
\begin{pgfscope}%
\pgfsetbuttcap%
\pgfsetroundjoin%
\definecolor{currentfill}{rgb}{0.309804,0.372549,0.419608}%
\pgfsetfillcolor{currentfill}%
\pgfsetlinewidth{0.000000pt}%
\definecolor{currentstroke}{rgb}{0.309804,0.372549,0.419608}%
\pgfsetstrokecolor{currentstroke}%
\pgfsetdash{}{0pt}%
\pgfpathmoveto{\pgfqpoint{3.295665in}{0.973960in}}%
\pgfpathlineto{\pgfqpoint{3.295665in}{0.961982in}}%
\pgfpathquadraticcurveto{\pgfqpoint{3.289919in}{0.967332in}}{\pgfqpoint{3.283417in}{0.969962in}}%
\pgfpathquadraticcurveto{\pgfqpoint{3.276914in}{0.972609in}}{\pgfqpoint{3.269601in}{0.972609in}}%
\pgfpathquadraticcurveto{\pgfqpoint{3.255192in}{0.972609in}}{\pgfqpoint{3.247537in}{0.963801in}}%
\pgfpathquadraticcurveto{\pgfqpoint{3.239881in}{0.954994in}}{\pgfqpoint{3.239881in}{0.938332in}}%
\pgfpathquadraticcurveto{\pgfqpoint{3.239881in}{0.921725in}}{\pgfqpoint{3.247537in}{0.912917in}}%
\pgfpathquadraticcurveto{\pgfqpoint{3.255192in}{0.904109in}}{\pgfqpoint{3.269601in}{0.904109in}}%
\pgfpathquadraticcurveto{\pgfqpoint{3.276914in}{0.904109in}}{\pgfqpoint{3.283417in}{0.906757in}}%
\pgfpathquadraticcurveto{\pgfqpoint{3.289919in}{0.909405in}}{\pgfqpoint{3.295665in}{0.914754in}}%
\pgfpathlineto{\pgfqpoint{3.295665in}{0.902866in}}%
\pgfpathquadraticcurveto{\pgfqpoint{3.289703in}{0.898814in}}{\pgfqpoint{3.283021in}{0.896778in}}%
\pgfpathquadraticcurveto{\pgfqpoint{3.276356in}{0.894761in}}{\pgfqpoint{3.268935in}{0.894761in}}%
\pgfpathquadraticcurveto{\pgfqpoint{3.249842in}{0.894761in}}{\pgfqpoint{3.238855in}{0.906433in}}%
\pgfpathquadraticcurveto{\pgfqpoint{3.227885in}{0.918123in}}{\pgfqpoint{3.227885in}{0.938332in}}%
\pgfpathquadraticcurveto{\pgfqpoint{3.227885in}{0.958596in}}{\pgfqpoint{3.238855in}{0.970268in}}%
\pgfpathquadraticcurveto{\pgfqpoint{3.249842in}{0.981958in}}{\pgfqpoint{3.268935in}{0.981958in}}%
\pgfpathquadraticcurveto{\pgfqpoint{3.276464in}{0.981958in}}{\pgfqpoint{3.283129in}{0.979958in}}%
\pgfpathquadraticcurveto{\pgfqpoint{3.289811in}{0.977959in}}{\pgfqpoint{3.295665in}{0.973960in}}%
\pgfpathlineto{\pgfqpoint{3.295665in}{0.973960in}}%
\pgfpathclose%
\pgfpathmoveto{\pgfqpoint{3.349305in}{0.949770in}}%
\pgfpathquadraticcurveto{\pgfqpoint{3.347558in}{0.950779in}}{\pgfqpoint{3.345505in}{0.951247in}}%
\pgfpathquadraticcurveto{\pgfqpoint{3.343451in}{0.951733in}}{\pgfqpoint{3.340984in}{0.951733in}}%
\pgfpathquadraticcurveto{\pgfqpoint{3.332194in}{0.951733in}}{\pgfqpoint{3.327493in}{0.946023in}}%
\pgfpathquadraticcurveto{\pgfqpoint{3.322791in}{0.940314in}}{\pgfqpoint{3.322791in}{0.929614in}}%
\pgfpathlineto{\pgfqpoint{3.322791in}{0.896400in}}%
\pgfpathlineto{\pgfqpoint{3.312380in}{0.896400in}}%
\pgfpathlineto{\pgfqpoint{3.312380in}{0.959443in}}%
\pgfpathlineto{\pgfqpoint{3.322791in}{0.959443in}}%
\pgfpathlineto{\pgfqpoint{3.322791in}{0.949644in}}%
\pgfpathquadraticcurveto{\pgfqpoint{3.326070in}{0.955390in}}{\pgfqpoint{3.331293in}{0.958164in}}%
\pgfpathquadraticcurveto{\pgfqpoint{3.336535in}{0.960956in}}{\pgfqpoint{3.344028in}{0.960956in}}%
\pgfpathquadraticcurveto{\pgfqpoint{3.345090in}{0.960956in}}{\pgfqpoint{3.346387in}{0.960811in}}%
\pgfpathquadraticcurveto{\pgfqpoint{3.347684in}{0.960685in}}{\pgfqpoint{3.349251in}{0.960397in}}%
\pgfpathlineto{\pgfqpoint{3.349305in}{0.949770in}}%
\pgfpathlineto{\pgfqpoint{3.349305in}{0.949770in}}%
\pgfpathclose%
\pgfpathmoveto{\pgfqpoint{3.382051in}{0.952184in}}%
\pgfpathquadraticcurveto{\pgfqpoint{3.373730in}{0.952184in}}{\pgfqpoint{3.368884in}{0.945681in}}%
\pgfpathquadraticcurveto{\pgfqpoint{3.364039in}{0.939179in}}{\pgfqpoint{3.364039in}{0.927867in}}%
\pgfpathquadraticcurveto{\pgfqpoint{3.364039in}{0.916556in}}{\pgfqpoint{3.368848in}{0.910053in}}%
\pgfpathquadraticcurveto{\pgfqpoint{3.373676in}{0.903551in}}{\pgfqpoint{3.382051in}{0.903551in}}%
\pgfpathquadraticcurveto{\pgfqpoint{3.390337in}{0.903551in}}{\pgfqpoint{3.395164in}{0.910071in}}%
\pgfpathquadraticcurveto{\pgfqpoint{3.400009in}{0.916610in}}{\pgfqpoint{3.400009in}{0.927867in}}%
\pgfpathquadraticcurveto{\pgfqpoint{3.400009in}{0.939071in}}{\pgfqpoint{3.395164in}{0.945627in}}%
\pgfpathquadraticcurveto{\pgfqpoint{3.390337in}{0.952184in}}{\pgfqpoint{3.382051in}{0.952184in}}%
\pgfpathlineto{\pgfqpoint{3.382051in}{0.952184in}}%
\pgfpathclose%
\pgfpathmoveto{\pgfqpoint{3.382051in}{0.960956in}}%
\pgfpathquadraticcurveto{\pgfqpoint{3.395560in}{0.960956in}}{\pgfqpoint{3.403270in}{0.952166in}}%
\pgfpathquadraticcurveto{\pgfqpoint{3.410997in}{0.943394in}}{\pgfqpoint{3.410997in}{0.927867in}}%
\pgfpathquadraticcurveto{\pgfqpoint{3.410997in}{0.912395in}}{\pgfqpoint{3.403270in}{0.903569in}}%
\pgfpathquadraticcurveto{\pgfqpoint{3.395560in}{0.894761in}}{\pgfqpoint{3.382051in}{0.894761in}}%
\pgfpathquadraticcurveto{\pgfqpoint{3.368488in}{0.894761in}}{\pgfqpoint{3.360797in}{0.903569in}}%
\pgfpathquadraticcurveto{\pgfqpoint{3.353124in}{0.912395in}}{\pgfqpoint{3.353124in}{0.927867in}}%
\pgfpathquadraticcurveto{\pgfqpoint{3.353124in}{0.943394in}}{\pgfqpoint{3.360797in}{0.952166in}}%
\pgfpathquadraticcurveto{\pgfqpoint{3.368488in}{0.960956in}}{\pgfqpoint{3.382051in}{0.960956in}}%
\pgfpathlineto{\pgfqpoint{3.382051in}{0.960956in}}%
\pgfpathclose%
\pgfpathmoveto{\pgfqpoint{3.468348in}{0.957587in}}%
\pgfpathlineto{\pgfqpoint{3.468348in}{0.947789in}}%
\pgfpathquadraticcurveto{\pgfqpoint{3.463971in}{0.950040in}}{\pgfqpoint{3.459233in}{0.951157in}}%
\pgfpathquadraticcurveto{\pgfqpoint{3.454514in}{0.952292in}}{\pgfqpoint{3.449435in}{0.952292in}}%
\pgfpathquadraticcurveto{\pgfqpoint{3.441726in}{0.952292in}}{\pgfqpoint{3.437871in}{0.949932in}}%
\pgfpathquadraticcurveto{\pgfqpoint{3.434016in}{0.947573in}}{\pgfqpoint{3.434016in}{0.942835in}}%
\pgfpathquadraticcurveto{\pgfqpoint{3.434016in}{0.939233in}}{\pgfqpoint{3.436772in}{0.937180in}}%
\pgfpathquadraticcurveto{\pgfqpoint{3.439528in}{0.935126in}}{\pgfqpoint{3.447868in}{0.933271in}}%
\pgfpathlineto{\pgfqpoint{3.451416in}{0.932478in}}%
\pgfpathquadraticcurveto{\pgfqpoint{3.462440in}{0.930119in}}{\pgfqpoint{3.467087in}{0.925814in}}%
\pgfpathquadraticcurveto{\pgfqpoint{3.471734in}{0.921509in}}{\pgfqpoint{3.471734in}{0.913800in}}%
\pgfpathquadraticcurveto{\pgfqpoint{3.471734in}{0.905010in}}{\pgfqpoint{3.464781in}{0.899876in}}%
\pgfpathquadraticcurveto{\pgfqpoint{3.457828in}{0.894761in}}{\pgfqpoint{3.445670in}{0.894761in}}%
\pgfpathquadraticcurveto{\pgfqpoint{3.440609in}{0.894761in}}{\pgfqpoint{3.435115in}{0.895752in}}%
\pgfpathquadraticcurveto{\pgfqpoint{3.429621in}{0.896742in}}{\pgfqpoint{3.423551in}{0.898706in}}%
\pgfpathlineto{\pgfqpoint{3.423551in}{0.909405in}}%
\pgfpathquadraticcurveto{\pgfqpoint{3.429297in}{0.906415in}}{\pgfqpoint{3.434863in}{0.904920in}}%
\pgfpathquadraticcurveto{\pgfqpoint{3.440429in}{0.903443in}}{\pgfqpoint{3.445904in}{0.903443in}}%
\pgfpathquadraticcurveto{\pgfqpoint{3.453217in}{0.903443in}}{\pgfqpoint{3.457144in}{0.905946in}}%
\pgfpathquadraticcurveto{\pgfqpoint{3.461089in}{0.908450in}}{\pgfqpoint{3.461089in}{0.913007in}}%
\pgfpathquadraticcurveto{\pgfqpoint{3.461089in}{0.917222in}}{\pgfqpoint{3.458243in}{0.919474in}}%
\pgfpathquadraticcurveto{\pgfqpoint{3.455415in}{0.921725in}}{\pgfqpoint{3.445778in}{0.923814in}}%
\pgfpathlineto{\pgfqpoint{3.442176in}{0.924661in}}%
\pgfpathquadraticcurveto{\pgfqpoint{3.432557in}{0.926678in}}{\pgfqpoint{3.428271in}{0.930875in}}%
\pgfpathquadraticcurveto{\pgfqpoint{3.424002in}{0.935072in}}{\pgfqpoint{3.424002in}{0.942385in}}%
\pgfpathquadraticcurveto{\pgfqpoint{3.424002in}{0.951283in}}{\pgfqpoint{3.430306in}{0.956110in}}%
\pgfpathquadraticcurveto{\pgfqpoint{3.436610in}{0.960956in}}{\pgfqpoint{3.448210in}{0.960956in}}%
\pgfpathquadraticcurveto{\pgfqpoint{3.453938in}{0.960956in}}{\pgfqpoint{3.458999in}{0.960109in}}%
\pgfpathquadraticcurveto{\pgfqpoint{3.464079in}{0.959280in}}{\pgfqpoint{3.468348in}{0.957587in}}%
\pgfpathlineto{\pgfqpoint{3.468348in}{0.957587in}}%
\pgfpathclose%
\pgfpathmoveto{\pgfqpoint{3.528400in}{0.957587in}}%
\pgfpathlineto{\pgfqpoint{3.528400in}{0.947789in}}%
\pgfpathquadraticcurveto{\pgfqpoint{3.524023in}{0.950040in}}{\pgfqpoint{3.519286in}{0.951157in}}%
\pgfpathquadraticcurveto{\pgfqpoint{3.514567in}{0.952292in}}{\pgfqpoint{3.509487in}{0.952292in}}%
\pgfpathquadraticcurveto{\pgfqpoint{3.501778in}{0.952292in}}{\pgfqpoint{3.497924in}{0.949932in}}%
\pgfpathquadraticcurveto{\pgfqpoint{3.494069in}{0.947573in}}{\pgfqpoint{3.494069in}{0.942835in}}%
\pgfpathquadraticcurveto{\pgfqpoint{3.494069in}{0.939233in}}{\pgfqpoint{3.496825in}{0.937180in}}%
\pgfpathquadraticcurveto{\pgfqpoint{3.499581in}{0.935126in}}{\pgfqpoint{3.507920in}{0.933271in}}%
\pgfpathlineto{\pgfqpoint{3.511469in}{0.932478in}}%
\pgfpathquadraticcurveto{\pgfqpoint{3.522492in}{0.930119in}}{\pgfqpoint{3.527139in}{0.925814in}}%
\pgfpathquadraticcurveto{\pgfqpoint{3.531786in}{0.921509in}}{\pgfqpoint{3.531786in}{0.913800in}}%
\pgfpathquadraticcurveto{\pgfqpoint{3.531786in}{0.905010in}}{\pgfqpoint{3.524834in}{0.899876in}}%
\pgfpathquadraticcurveto{\pgfqpoint{3.517881in}{0.894761in}}{\pgfqpoint{3.505723in}{0.894761in}}%
\pgfpathquadraticcurveto{\pgfqpoint{3.500661in}{0.894761in}}{\pgfqpoint{3.495168in}{0.895752in}}%
\pgfpathquadraticcurveto{\pgfqpoint{3.489674in}{0.896742in}}{\pgfqpoint{3.483604in}{0.898706in}}%
\pgfpathlineto{\pgfqpoint{3.483604in}{0.909405in}}%
\pgfpathquadraticcurveto{\pgfqpoint{3.489350in}{0.906415in}}{\pgfqpoint{3.494916in}{0.904920in}}%
\pgfpathquadraticcurveto{\pgfqpoint{3.500481in}{0.903443in}}{\pgfqpoint{3.505957in}{0.903443in}}%
\pgfpathquadraticcurveto{\pgfqpoint{3.513270in}{0.903443in}}{\pgfqpoint{3.517197in}{0.905946in}}%
\pgfpathquadraticcurveto{\pgfqpoint{3.521141in}{0.908450in}}{\pgfqpoint{3.521141in}{0.913007in}}%
\pgfpathquadraticcurveto{\pgfqpoint{3.521141in}{0.917222in}}{\pgfqpoint{3.518295in}{0.919474in}}%
\pgfpathquadraticcurveto{\pgfqpoint{3.515467in}{0.921725in}}{\pgfqpoint{3.505831in}{0.923814in}}%
\pgfpathlineto{\pgfqpoint{3.502228in}{0.924661in}}%
\pgfpathquadraticcurveto{\pgfqpoint{3.492610in}{0.926678in}}{\pgfqpoint{3.488323in}{0.930875in}}%
\pgfpathquadraticcurveto{\pgfqpoint{3.484054in}{0.935072in}}{\pgfqpoint{3.484054in}{0.942385in}}%
\pgfpathquadraticcurveto{\pgfqpoint{3.484054in}{0.951283in}}{\pgfqpoint{3.490358in}{0.956110in}}%
\pgfpathquadraticcurveto{\pgfqpoint{3.496663in}{0.960956in}}{\pgfqpoint{3.508263in}{0.960956in}}%
\pgfpathquadraticcurveto{\pgfqpoint{3.513990in}{0.960956in}}{\pgfqpoint{3.519052in}{0.960109in}}%
\pgfpathquadraticcurveto{\pgfqpoint{3.524131in}{0.959280in}}{\pgfqpoint{3.528400in}{0.957587in}}%
\pgfpathlineto{\pgfqpoint{3.528400in}{0.957587in}}%
\pgfpathclose%
\pgfpathmoveto{\pgfqpoint{3.602196in}{0.930515in}}%
\pgfpathlineto{\pgfqpoint{3.602196in}{0.925454in}}%
\pgfpathlineto{\pgfqpoint{3.554572in}{0.925454in}}%
\pgfpathquadraticcurveto{\pgfqpoint{3.555256in}{0.914754in}}{\pgfqpoint{3.561020in}{0.909153in}}%
\pgfpathquadraticcurveto{\pgfqpoint{3.566784in}{0.903551in}}{\pgfqpoint{3.577087in}{0.903551in}}%
\pgfpathquadraticcurveto{\pgfqpoint{3.583049in}{0.903551in}}{\pgfqpoint{3.588651in}{0.905010in}}%
\pgfpathquadraticcurveto{\pgfqpoint{3.594253in}{0.906469in}}{\pgfqpoint{3.599782in}{0.909405in}}%
\pgfpathlineto{\pgfqpoint{3.599782in}{0.899606in}}%
\pgfpathquadraticcurveto{\pgfqpoint{3.594198in}{0.897247in}}{\pgfqpoint{3.588345in}{0.896004in}}%
\pgfpathquadraticcurveto{\pgfqpoint{3.582491in}{0.894761in}}{\pgfqpoint{3.576475in}{0.894761in}}%
\pgfpathquadraticcurveto{\pgfqpoint{3.561380in}{0.894761in}}{\pgfqpoint{3.552572in}{0.903533in}}%
\pgfpathquadraticcurveto{\pgfqpoint{3.543764in}{0.912323in}}{\pgfqpoint{3.543764in}{0.927309in}}%
\pgfpathquadraticcurveto{\pgfqpoint{3.543764in}{0.942781in}}{\pgfqpoint{3.552122in}{0.951859in}}%
\pgfpathquadraticcurveto{\pgfqpoint{3.560480in}{0.960956in}}{\pgfqpoint{3.574673in}{0.960956in}}%
\pgfpathquadraticcurveto{\pgfqpoint{3.587390in}{0.960956in}}{\pgfqpoint{3.594793in}{0.952760in}}%
\pgfpathquadraticcurveto{\pgfqpoint{3.602196in}{0.944583in}}{\pgfqpoint{3.602196in}{0.930515in}}%
\pgfpathlineto{\pgfqpoint{3.602196in}{0.930515in}}%
\pgfpathclose%
\pgfpathmoveto{\pgfqpoint{3.591839in}{0.933559in}}%
\pgfpathquadraticcurveto{\pgfqpoint{3.591731in}{0.942043in}}{\pgfqpoint{3.587084in}{0.947104in}}%
\pgfpathquadraticcurveto{\pgfqpoint{3.582437in}{0.952184in}}{\pgfqpoint{3.574781in}{0.952184in}}%
\pgfpathquadraticcurveto{\pgfqpoint{3.566118in}{0.952184in}}{\pgfqpoint{3.560912in}{0.947284in}}%
\pgfpathquadraticcurveto{\pgfqpoint{3.555707in}{0.942385in}}{\pgfqpoint{3.554914in}{0.933487in}}%
\pgfpathlineto{\pgfqpoint{3.591839in}{0.933559in}}%
\pgfpathlineto{\pgfqpoint{3.591839in}{0.933559in}}%
\pgfpathclose%
\pgfpathmoveto{\pgfqpoint{3.659384in}{0.957587in}}%
\pgfpathlineto{\pgfqpoint{3.659384in}{0.947789in}}%
\pgfpathquadraticcurveto{\pgfqpoint{3.655008in}{0.950040in}}{\pgfqpoint{3.650270in}{0.951157in}}%
\pgfpathquadraticcurveto{\pgfqpoint{3.645551in}{0.952292in}}{\pgfqpoint{3.640472in}{0.952292in}}%
\pgfpathquadraticcurveto{\pgfqpoint{3.632763in}{0.952292in}}{\pgfqpoint{3.628908in}{0.949932in}}%
\pgfpathquadraticcurveto{\pgfqpoint{3.625053in}{0.947573in}}{\pgfqpoint{3.625053in}{0.942835in}}%
\pgfpathquadraticcurveto{\pgfqpoint{3.625053in}{0.939233in}}{\pgfqpoint{3.627809in}{0.937180in}}%
\pgfpathquadraticcurveto{\pgfqpoint{3.630565in}{0.935126in}}{\pgfqpoint{3.638905in}{0.933271in}}%
\pgfpathlineto{\pgfqpoint{3.642453in}{0.932478in}}%
\pgfpathquadraticcurveto{\pgfqpoint{3.653476in}{0.930119in}}{\pgfqpoint{3.658124in}{0.925814in}}%
\pgfpathquadraticcurveto{\pgfqpoint{3.662771in}{0.921509in}}{\pgfqpoint{3.662771in}{0.913800in}}%
\pgfpathquadraticcurveto{\pgfqpoint{3.662771in}{0.905010in}}{\pgfqpoint{3.655818in}{0.899876in}}%
\pgfpathquadraticcurveto{\pgfqpoint{3.648865in}{0.894761in}}{\pgfqpoint{3.636707in}{0.894761in}}%
\pgfpathquadraticcurveto{\pgfqpoint{3.631646in}{0.894761in}}{\pgfqpoint{3.626152in}{0.895752in}}%
\pgfpathquadraticcurveto{\pgfqpoint{3.620658in}{0.896742in}}{\pgfqpoint{3.614588in}{0.898706in}}%
\pgfpathlineto{\pgfqpoint{3.614588in}{0.909405in}}%
\pgfpathquadraticcurveto{\pgfqpoint{3.620334in}{0.906415in}}{\pgfqpoint{3.625900in}{0.904920in}}%
\pgfpathquadraticcurveto{\pgfqpoint{3.631466in}{0.903443in}}{\pgfqpoint{3.636941in}{0.903443in}}%
\pgfpathquadraticcurveto{\pgfqpoint{3.644254in}{0.903443in}}{\pgfqpoint{3.648181in}{0.905946in}}%
\pgfpathquadraticcurveto{\pgfqpoint{3.652126in}{0.908450in}}{\pgfqpoint{3.652126in}{0.913007in}}%
\pgfpathquadraticcurveto{\pgfqpoint{3.652126in}{0.917222in}}{\pgfqpoint{3.649280in}{0.919474in}}%
\pgfpathquadraticcurveto{\pgfqpoint{3.646452in}{0.921725in}}{\pgfqpoint{3.636815in}{0.923814in}}%
\pgfpathlineto{\pgfqpoint{3.633213in}{0.924661in}}%
\pgfpathquadraticcurveto{\pgfqpoint{3.623594in}{0.926678in}}{\pgfqpoint{3.619307in}{0.930875in}}%
\pgfpathquadraticcurveto{\pgfqpoint{3.615039in}{0.935072in}}{\pgfqpoint{3.615039in}{0.942385in}}%
\pgfpathquadraticcurveto{\pgfqpoint{3.615039in}{0.951283in}}{\pgfqpoint{3.621343in}{0.956110in}}%
\pgfpathquadraticcurveto{\pgfqpoint{3.627647in}{0.960956in}}{\pgfqpoint{3.639247in}{0.960956in}}%
\pgfpathquadraticcurveto{\pgfqpoint{3.644975in}{0.960956in}}{\pgfqpoint{3.650036in}{0.960109in}}%
\pgfpathquadraticcurveto{\pgfqpoint{3.655116in}{0.959280in}}{\pgfqpoint{3.659384in}{0.957587in}}%
\pgfpathlineto{\pgfqpoint{3.659384in}{0.957587in}}%
\pgfpathclose%
\pgfpathmoveto{\pgfqpoint{3.681900in}{0.910702in}}%
\pgfpathlineto{\pgfqpoint{3.693770in}{0.910702in}}%
\pgfpathlineto{\pgfqpoint{3.693770in}{0.896400in}}%
\pgfpathlineto{\pgfqpoint{3.681900in}{0.896400in}}%
\pgfpathlineto{\pgfqpoint{3.681900in}{0.910702in}}%
\pgfpathlineto{\pgfqpoint{3.681900in}{0.910702in}}%
\pgfpathclose%
\pgfpathmoveto{\pgfqpoint{3.681900in}{0.956002in}}%
\pgfpathlineto{\pgfqpoint{3.693770in}{0.956002in}}%
\pgfpathlineto{\pgfqpoint{3.693770in}{0.941719in}}%
\pgfpathlineto{\pgfqpoint{3.681900in}{0.941719in}}%
\pgfpathlineto{\pgfqpoint{3.681900in}{0.956002in}}%
\pgfpathlineto{\pgfqpoint{3.681900in}{0.956002in}}%
\pgfpathclose%
\pgfusepath{fill}%
\end{pgfscope}%
\begin{pgfscope}%
\pgfsetbuttcap%
\pgfsetroundjoin%
\definecolor{currentfill}{rgb}{0.309804,0.372549,0.419608}%
\pgfsetfillcolor{currentfill}%
\pgfsetlinewidth{0.000000pt}%
\definecolor{currentstroke}{rgb}{0.309804,0.372549,0.419608}%
\pgfsetstrokecolor{currentstroke}%
\pgfsetdash{}{0pt}%
\pgfpathmoveto{\pgfqpoint{3.821192in}{0.930515in}}%
\pgfpathlineto{\pgfqpoint{3.821192in}{0.925454in}}%
\pgfpathlineto{\pgfqpoint{3.773568in}{0.925454in}}%
\pgfpathquadraticcurveto{\pgfqpoint{3.774253in}{0.914754in}}{\pgfqpoint{3.780017in}{0.909153in}}%
\pgfpathquadraticcurveto{\pgfqpoint{3.785781in}{0.903551in}}{\pgfqpoint{3.796084in}{0.903551in}}%
\pgfpathquadraticcurveto{\pgfqpoint{3.802046in}{0.903551in}}{\pgfqpoint{3.807647in}{0.905010in}}%
\pgfpathquadraticcurveto{\pgfqpoint{3.813249in}{0.906469in}}{\pgfqpoint{3.818779in}{0.909405in}}%
\pgfpathlineto{\pgfqpoint{3.818779in}{0.899606in}}%
\pgfpathquadraticcurveto{\pgfqpoint{3.813195in}{0.897247in}}{\pgfqpoint{3.807341in}{0.896004in}}%
\pgfpathquadraticcurveto{\pgfqpoint{3.801487in}{0.894761in}}{\pgfqpoint{3.795471in}{0.894761in}}%
\pgfpathquadraticcurveto{\pgfqpoint{3.780377in}{0.894761in}}{\pgfqpoint{3.771569in}{0.903533in}}%
\pgfpathquadraticcurveto{\pgfqpoint{3.762761in}{0.912323in}}{\pgfqpoint{3.762761in}{0.927309in}}%
\pgfpathquadraticcurveto{\pgfqpoint{3.762761in}{0.942781in}}{\pgfqpoint{3.771119in}{0.951859in}}%
\pgfpathquadraticcurveto{\pgfqpoint{3.779476in}{0.960956in}}{\pgfqpoint{3.793670in}{0.960956in}}%
\pgfpathquadraticcurveto{\pgfqpoint{3.806387in}{0.960956in}}{\pgfqpoint{3.813789in}{0.952760in}}%
\pgfpathquadraticcurveto{\pgfqpoint{3.821192in}{0.944583in}}{\pgfqpoint{3.821192in}{0.930515in}}%
\pgfpathlineto{\pgfqpoint{3.821192in}{0.930515in}}%
\pgfpathclose%
\pgfpathmoveto{\pgfqpoint{3.810836in}{0.933559in}}%
\pgfpathquadraticcurveto{\pgfqpoint{3.810727in}{0.942043in}}{\pgfqpoint{3.806080in}{0.947104in}}%
\pgfpathquadraticcurveto{\pgfqpoint{3.801433in}{0.952184in}}{\pgfqpoint{3.793778in}{0.952184in}}%
\pgfpathquadraticcurveto{\pgfqpoint{3.785114in}{0.952184in}}{\pgfqpoint{3.779909in}{0.947284in}}%
\pgfpathquadraticcurveto{\pgfqpoint{3.774703in}{0.942385in}}{\pgfqpoint{3.773911in}{0.933487in}}%
\pgfpathlineto{\pgfqpoint{3.810836in}{0.933559in}}%
\pgfpathlineto{\pgfqpoint{3.810836in}{0.933559in}}%
\pgfpathclose%
\pgfpathmoveto{\pgfqpoint{3.908948in}{0.983993in}}%
\pgfpathlineto{\pgfqpoint{3.908948in}{0.975365in}}%
\pgfpathlineto{\pgfqpoint{3.899041in}{0.975365in}}%
\pgfpathquadraticcurveto{\pgfqpoint{3.893475in}{0.975365in}}{\pgfqpoint{3.891296in}{0.973114in}}%
\pgfpathquadraticcurveto{\pgfqpoint{3.889134in}{0.970862in}}{\pgfqpoint{3.889134in}{0.965008in}}%
\pgfpathlineto{\pgfqpoint{3.889134in}{0.959443in}}%
\pgfpathlineto{\pgfqpoint{3.906192in}{0.959443in}}%
\pgfpathlineto{\pgfqpoint{3.906192in}{0.951391in}}%
\pgfpathlineto{\pgfqpoint{3.889134in}{0.951391in}}%
\pgfpathlineto{\pgfqpoint{3.889134in}{0.896400in}}%
\pgfpathlineto{\pgfqpoint{3.878723in}{0.896400in}}%
\pgfpathlineto{\pgfqpoint{3.878723in}{0.951391in}}%
\pgfpathlineto{\pgfqpoint{3.850300in}{0.951391in}}%
\pgfpathlineto{\pgfqpoint{3.850300in}{0.896400in}}%
\pgfpathlineto{\pgfqpoint{3.839889in}{0.896400in}}%
\pgfpathlineto{\pgfqpoint{3.839889in}{0.951391in}}%
\pgfpathlineto{\pgfqpoint{3.829982in}{0.951391in}}%
\pgfpathlineto{\pgfqpoint{3.829982in}{0.959443in}}%
\pgfpathlineto{\pgfqpoint{3.839889in}{0.959443in}}%
\pgfpathlineto{\pgfqpoint{3.839889in}{0.963838in}}%
\pgfpathquadraticcurveto{\pgfqpoint{3.839889in}{0.974357in}}{\pgfqpoint{3.844788in}{0.979166in}}%
\pgfpathquadraticcurveto{\pgfqpoint{3.849688in}{0.983993in}}{\pgfqpoint{3.860315in}{0.983993in}}%
\pgfpathlineto{\pgfqpoint{3.870113in}{0.983993in}}%
\pgfpathlineto{\pgfqpoint{3.870113in}{0.975365in}}%
\pgfpathlineto{\pgfqpoint{3.860207in}{0.975365in}}%
\pgfpathquadraticcurveto{\pgfqpoint{3.854641in}{0.975365in}}{\pgfqpoint{3.852444in}{0.973114in}}%
\pgfpathquadraticcurveto{\pgfqpoint{3.850300in}{0.970862in}}{\pgfqpoint{3.850300in}{0.965008in}}%
\pgfpathlineto{\pgfqpoint{3.850300in}{0.959443in}}%
\pgfpathlineto{\pgfqpoint{3.878723in}{0.959443in}}%
\pgfpathlineto{\pgfqpoint{3.878723in}{0.963838in}}%
\pgfpathquadraticcurveto{\pgfqpoint{3.878723in}{0.974357in}}{\pgfqpoint{3.883623in}{0.979166in}}%
\pgfpathquadraticcurveto{\pgfqpoint{3.888522in}{0.983993in}}{\pgfqpoint{3.899167in}{0.983993in}}%
\pgfpathlineto{\pgfqpoint{3.908948in}{0.983993in}}%
\pgfpathlineto{\pgfqpoint{3.908948in}{0.983993in}}%
\pgfpathclose%
\pgfpathmoveto{\pgfqpoint{3.971540in}{0.930515in}}%
\pgfpathlineto{\pgfqpoint{3.971540in}{0.925454in}}%
\pgfpathlineto{\pgfqpoint{3.923916in}{0.925454in}}%
\pgfpathquadraticcurveto{\pgfqpoint{3.924600in}{0.914754in}}{\pgfqpoint{3.930364in}{0.909153in}}%
\pgfpathquadraticcurveto{\pgfqpoint{3.936128in}{0.903551in}}{\pgfqpoint{3.946431in}{0.903551in}}%
\pgfpathquadraticcurveto{\pgfqpoint{3.952393in}{0.903551in}}{\pgfqpoint{3.957995in}{0.905010in}}%
\pgfpathquadraticcurveto{\pgfqpoint{3.963597in}{0.906469in}}{\pgfqpoint{3.969126in}{0.909405in}}%
\pgfpathlineto{\pgfqpoint{3.969126in}{0.899606in}}%
\pgfpathquadraticcurveto{\pgfqpoint{3.963543in}{0.897247in}}{\pgfqpoint{3.957689in}{0.896004in}}%
\pgfpathquadraticcurveto{\pgfqpoint{3.951835in}{0.894761in}}{\pgfqpoint{3.945819in}{0.894761in}}%
\pgfpathquadraticcurveto{\pgfqpoint{3.930724in}{0.894761in}}{\pgfqpoint{3.921916in}{0.903533in}}%
\pgfpathquadraticcurveto{\pgfqpoint{3.913109in}{0.912323in}}{\pgfqpoint{3.913109in}{0.927309in}}%
\pgfpathquadraticcurveto{\pgfqpoint{3.913109in}{0.942781in}}{\pgfqpoint{3.921466in}{0.951859in}}%
\pgfpathquadraticcurveto{\pgfqpoint{3.929824in}{0.960956in}}{\pgfqpoint{3.944017in}{0.960956in}}%
\pgfpathquadraticcurveto{\pgfqpoint{3.956734in}{0.960956in}}{\pgfqpoint{3.964137in}{0.952760in}}%
\pgfpathquadraticcurveto{\pgfqpoint{3.971540in}{0.944583in}}{\pgfqpoint{3.971540in}{0.930515in}}%
\pgfpathlineto{\pgfqpoint{3.971540in}{0.930515in}}%
\pgfpathclose%
\pgfpathmoveto{\pgfqpoint{3.961183in}{0.933559in}}%
\pgfpathquadraticcurveto{\pgfqpoint{3.961075in}{0.942043in}}{\pgfqpoint{3.956428in}{0.947104in}}%
\pgfpathquadraticcurveto{\pgfqpoint{3.951781in}{0.952184in}}{\pgfqpoint{3.944125in}{0.952184in}}%
\pgfpathquadraticcurveto{\pgfqpoint{3.935462in}{0.952184in}}{\pgfqpoint{3.930256in}{0.947284in}}%
\pgfpathquadraticcurveto{\pgfqpoint{3.925051in}{0.942385in}}{\pgfqpoint{3.924258in}{0.933487in}}%
\pgfpathlineto{\pgfqpoint{3.961183in}{0.933559in}}%
\pgfpathlineto{\pgfqpoint{3.961183in}{0.933559in}}%
\pgfpathclose%
\pgfpathmoveto{\pgfqpoint{4.033916in}{0.957029in}}%
\pgfpathlineto{\pgfqpoint{4.033916in}{0.947338in}}%
\pgfpathquadraticcurveto{\pgfqpoint{4.029521in}{0.949770in}}{\pgfqpoint{4.025108in}{0.950977in}}%
\pgfpathquadraticcurveto{\pgfqpoint{4.020695in}{0.952184in}}{\pgfqpoint{4.016192in}{0.952184in}}%
\pgfpathquadraticcurveto{\pgfqpoint{4.006105in}{0.952184in}}{\pgfqpoint{4.000521in}{0.945789in}}%
\pgfpathquadraticcurveto{\pgfqpoint{3.994956in}{0.939413in}}{\pgfqpoint{3.994956in}{0.927867in}}%
\pgfpathquadraticcurveto{\pgfqpoint{3.994956in}{0.916321in}}{\pgfqpoint{4.000521in}{0.909927in}}%
\pgfpathquadraticcurveto{\pgfqpoint{4.006105in}{0.903551in}}{\pgfqpoint{4.016192in}{0.903551in}}%
\pgfpathquadraticcurveto{\pgfqpoint{4.020695in}{0.903551in}}{\pgfqpoint{4.025108in}{0.904758in}}%
\pgfpathquadraticcurveto{\pgfqpoint{4.029521in}{0.905964in}}{\pgfqpoint{4.033916in}{0.908396in}}%
\pgfpathlineto{\pgfqpoint{4.033916in}{0.898814in}}%
\pgfpathquadraticcurveto{\pgfqpoint{4.029575in}{0.896796in}}{\pgfqpoint{4.024928in}{0.895788in}}%
\pgfpathquadraticcurveto{\pgfqpoint{4.020299in}{0.894761in}}{\pgfqpoint{4.015057in}{0.894761in}}%
\pgfpathquadraticcurveto{\pgfqpoint{4.000810in}{0.894761in}}{\pgfqpoint{3.992416in}{0.903713in}}%
\pgfpathquadraticcurveto{\pgfqpoint{3.984040in}{0.912665in}}{\pgfqpoint{3.984040in}{0.927867in}}%
\pgfpathquadraticcurveto{\pgfqpoint{3.984040in}{0.943286in}}{\pgfqpoint{3.992506in}{0.952112in}}%
\pgfpathquadraticcurveto{\pgfqpoint{4.000990in}{0.960956in}}{\pgfqpoint{4.015742in}{0.960956in}}%
\pgfpathquadraticcurveto{\pgfqpoint{4.020515in}{0.960956in}}{\pgfqpoint{4.025072in}{0.959965in}}%
\pgfpathquadraticcurveto{\pgfqpoint{4.029629in}{0.958992in}}{\pgfqpoint{4.033916in}{0.957029in}}%
\pgfpathlineto{\pgfqpoint{4.033916in}{0.957029in}}%
\pgfpathclose%
\pgfpathmoveto{\pgfqpoint{4.062177in}{0.977347in}}%
\pgfpathlineto{\pgfqpoint{4.062177in}{0.959443in}}%
\pgfpathlineto{\pgfqpoint{4.083503in}{0.959443in}}%
\pgfpathlineto{\pgfqpoint{4.083503in}{0.951391in}}%
\pgfpathlineto{\pgfqpoint{4.062177in}{0.951391in}}%
\pgfpathlineto{\pgfqpoint{4.062177in}{0.917168in}}%
\pgfpathquadraticcurveto{\pgfqpoint{4.062177in}{0.909459in}}{\pgfqpoint{4.064285in}{0.907261in}}%
\pgfpathquadraticcurveto{\pgfqpoint{4.066392in}{0.905064in}}{\pgfqpoint{4.072876in}{0.905064in}}%
\pgfpathlineto{\pgfqpoint{4.083503in}{0.905064in}}%
\pgfpathlineto{\pgfqpoint{4.083503in}{0.896400in}}%
\pgfpathlineto{\pgfqpoint{4.072876in}{0.896400in}}%
\pgfpathquadraticcurveto{\pgfqpoint{4.060880in}{0.896400in}}{\pgfqpoint{4.056323in}{0.900867in}}%
\pgfpathquadraticcurveto{\pgfqpoint{4.051766in}{0.905352in}}{\pgfqpoint{4.051766in}{0.917168in}}%
\pgfpathlineto{\pgfqpoint{4.051766in}{0.951391in}}%
\pgfpathlineto{\pgfqpoint{4.044165in}{0.951391in}}%
\pgfpathlineto{\pgfqpoint{4.044165in}{0.959443in}}%
\pgfpathlineto{\pgfqpoint{4.051766in}{0.959443in}}%
\pgfpathlineto{\pgfqpoint{4.051766in}{0.977347in}}%
\pgfpathlineto{\pgfqpoint{4.062177in}{0.977347in}}%
\pgfpathlineto{\pgfqpoint{4.062177in}{0.977347in}}%
\pgfpathclose%
\pgfpathmoveto{\pgfqpoint{4.097121in}{0.959443in}}%
\pgfpathlineto{\pgfqpoint{4.107478in}{0.959443in}}%
\pgfpathlineto{\pgfqpoint{4.107478in}{0.896400in}}%
\pgfpathlineto{\pgfqpoint{4.097121in}{0.896400in}}%
\pgfpathlineto{\pgfqpoint{4.097121in}{0.959443in}}%
\pgfpathlineto{\pgfqpoint{4.097121in}{0.959443in}}%
\pgfpathclose%
\pgfpathmoveto{\pgfqpoint{4.097121in}{0.983993in}}%
\pgfpathlineto{\pgfqpoint{4.107478in}{0.983993in}}%
\pgfpathlineto{\pgfqpoint{4.107478in}{0.970862in}}%
\pgfpathlineto{\pgfqpoint{4.097121in}{0.970862in}}%
\pgfpathlineto{\pgfqpoint{4.097121in}{0.983993in}}%
\pgfpathlineto{\pgfqpoint{4.097121in}{0.983993in}}%
\pgfpathclose%
\pgfpathmoveto{\pgfqpoint{4.121725in}{0.959443in}}%
\pgfpathlineto{\pgfqpoint{4.132695in}{0.959443in}}%
\pgfpathlineto{\pgfqpoint{4.152400in}{0.906541in}}%
\pgfpathlineto{\pgfqpoint{4.172105in}{0.959443in}}%
\pgfpathlineto{\pgfqpoint{4.183075in}{0.959443in}}%
\pgfpathlineto{\pgfqpoint{4.159425in}{0.896400in}}%
\pgfpathlineto{\pgfqpoint{4.145357in}{0.896400in}}%
\pgfpathlineto{\pgfqpoint{4.121725in}{0.959443in}}%
\pgfpathlineto{\pgfqpoint{4.121725in}{0.959443in}}%
\pgfpathclose%
\pgfpathmoveto{\pgfqpoint{4.251305in}{0.930515in}}%
\pgfpathlineto{\pgfqpoint{4.251305in}{0.925454in}}%
\pgfpathlineto{\pgfqpoint{4.203681in}{0.925454in}}%
\pgfpathquadraticcurveto{\pgfqpoint{4.204365in}{0.914754in}}{\pgfqpoint{4.210129in}{0.909153in}}%
\pgfpathquadraticcurveto{\pgfqpoint{4.215893in}{0.903551in}}{\pgfqpoint{4.226196in}{0.903551in}}%
\pgfpathquadraticcurveto{\pgfqpoint{4.232158in}{0.903551in}}{\pgfqpoint{4.237760in}{0.905010in}}%
\pgfpathquadraticcurveto{\pgfqpoint{4.243361in}{0.906469in}}{\pgfqpoint{4.248891in}{0.909405in}}%
\pgfpathlineto{\pgfqpoint{4.248891in}{0.899606in}}%
\pgfpathquadraticcurveto{\pgfqpoint{4.243307in}{0.897247in}}{\pgfqpoint{4.237453in}{0.896004in}}%
\pgfpathquadraticcurveto{\pgfqpoint{4.231599in}{0.894761in}}{\pgfqpoint{4.225583in}{0.894761in}}%
\pgfpathquadraticcurveto{\pgfqpoint{4.210489in}{0.894761in}}{\pgfqpoint{4.201681in}{0.903533in}}%
\pgfpathquadraticcurveto{\pgfqpoint{4.192873in}{0.912323in}}{\pgfqpoint{4.192873in}{0.927309in}}%
\pgfpathquadraticcurveto{\pgfqpoint{4.192873in}{0.942781in}}{\pgfqpoint{4.201231in}{0.951859in}}%
\pgfpathquadraticcurveto{\pgfqpoint{4.209589in}{0.960956in}}{\pgfqpoint{4.223782in}{0.960956in}}%
\pgfpathquadraticcurveto{\pgfqpoint{4.236499in}{0.960956in}}{\pgfqpoint{4.243902in}{0.952760in}}%
\pgfpathquadraticcurveto{\pgfqpoint{4.251305in}{0.944583in}}{\pgfqpoint{4.251305in}{0.930515in}}%
\pgfpathlineto{\pgfqpoint{4.251305in}{0.930515in}}%
\pgfpathclose%
\pgfpathmoveto{\pgfqpoint{4.240948in}{0.933559in}}%
\pgfpathquadraticcurveto{\pgfqpoint{4.240840in}{0.942043in}}{\pgfqpoint{4.236192in}{0.947104in}}%
\pgfpathquadraticcurveto{\pgfqpoint{4.231545in}{0.952184in}}{\pgfqpoint{4.223890in}{0.952184in}}%
\pgfpathquadraticcurveto{\pgfqpoint{4.215226in}{0.952184in}}{\pgfqpoint{4.210021in}{0.947284in}}%
\pgfpathquadraticcurveto{\pgfqpoint{4.204815in}{0.942385in}}{\pgfqpoint{4.204023in}{0.933487in}}%
\pgfpathlineto{\pgfqpoint{4.240948in}{0.933559in}}%
\pgfpathlineto{\pgfqpoint{4.240948in}{0.933559in}}%
\pgfpathclose%
\pgfusepath{fill}%
\end{pgfscope}%
\begin{pgfscope}%
\pgfsetbuttcap%
\pgfsetroundjoin%
\definecolor{currentfill}{rgb}{0.309804,0.372549,0.419608}%
\pgfsetfillcolor{currentfill}%
\pgfsetlinewidth{0.000000pt}%
\definecolor{currentstroke}{rgb}{0.309804,0.372549,0.419608}%
\pgfsetstrokecolor{currentstroke}%
\pgfsetdash{}{0pt}%
\pgfpathmoveto{\pgfqpoint{4.370121in}{0.957029in}}%
\pgfpathlineto{\pgfqpoint{4.370121in}{0.947338in}}%
\pgfpathquadraticcurveto{\pgfqpoint{4.365726in}{0.949770in}}{\pgfqpoint{4.361313in}{0.950977in}}%
\pgfpathquadraticcurveto{\pgfqpoint{4.356900in}{0.952184in}}{\pgfqpoint{4.352397in}{0.952184in}}%
\pgfpathquadraticcurveto{\pgfqpoint{4.342310in}{0.952184in}}{\pgfqpoint{4.336726in}{0.945789in}}%
\pgfpathquadraticcurveto{\pgfqpoint{4.331160in}{0.939413in}}{\pgfqpoint{4.331160in}{0.927867in}}%
\pgfpathquadraticcurveto{\pgfqpoint{4.331160in}{0.916321in}}{\pgfqpoint{4.336726in}{0.909927in}}%
\pgfpathquadraticcurveto{\pgfqpoint{4.342310in}{0.903551in}}{\pgfqpoint{4.352397in}{0.903551in}}%
\pgfpathquadraticcurveto{\pgfqpoint{4.356900in}{0.903551in}}{\pgfqpoint{4.361313in}{0.904758in}}%
\pgfpathquadraticcurveto{\pgfqpoint{4.365726in}{0.905964in}}{\pgfqpoint{4.370121in}{0.908396in}}%
\pgfpathlineto{\pgfqpoint{4.370121in}{0.898814in}}%
\pgfpathquadraticcurveto{\pgfqpoint{4.365780in}{0.896796in}}{\pgfqpoint{4.361132in}{0.895788in}}%
\pgfpathquadraticcurveto{\pgfqpoint{4.356503in}{0.894761in}}{\pgfqpoint{4.351262in}{0.894761in}}%
\pgfpathquadraticcurveto{\pgfqpoint{4.337014in}{0.894761in}}{\pgfqpoint{4.328621in}{0.903713in}}%
\pgfpathquadraticcurveto{\pgfqpoint{4.320245in}{0.912665in}}{\pgfqpoint{4.320245in}{0.927867in}}%
\pgfpathquadraticcurveto{\pgfqpoint{4.320245in}{0.943286in}}{\pgfqpoint{4.328711in}{0.952112in}}%
\pgfpathquadraticcurveto{\pgfqpoint{4.337194in}{0.960956in}}{\pgfqpoint{4.351946in}{0.960956in}}%
\pgfpathquadraticcurveto{\pgfqpoint{4.356719in}{0.960956in}}{\pgfqpoint{4.361277in}{0.959965in}}%
\pgfpathquadraticcurveto{\pgfqpoint{4.365834in}{0.958992in}}{\pgfqpoint{4.370121in}{0.957029in}}%
\pgfpathlineto{\pgfqpoint{4.370121in}{0.957029in}}%
\pgfpathclose%
\pgfpathmoveto{\pgfqpoint{4.388133in}{0.983993in}}%
\pgfpathlineto{\pgfqpoint{4.398490in}{0.983993in}}%
\pgfpathlineto{\pgfqpoint{4.398490in}{0.896400in}}%
\pgfpathlineto{\pgfqpoint{4.388133in}{0.896400in}}%
\pgfpathlineto{\pgfqpoint{4.388133in}{0.983993in}}%
\pgfpathlineto{\pgfqpoint{4.388133in}{0.983993in}}%
\pgfpathclose%
\pgfpathmoveto{\pgfqpoint{4.419096in}{0.921275in}}%
\pgfpathlineto{\pgfqpoint{4.419096in}{0.959443in}}%
\pgfpathlineto{\pgfqpoint{4.429453in}{0.959443in}}%
\pgfpathlineto{\pgfqpoint{4.429453in}{0.921671in}}%
\pgfpathquadraticcurveto{\pgfqpoint{4.429453in}{0.912719in}}{\pgfqpoint{4.432929in}{0.908234in}}%
\pgfpathquadraticcurveto{\pgfqpoint{4.436423in}{0.903767in}}{\pgfqpoint{4.443412in}{0.903767in}}%
\pgfpathquadraticcurveto{\pgfqpoint{4.451788in}{0.903767in}}{\pgfqpoint{4.456651in}{0.909117in}}%
\pgfpathquadraticcurveto{\pgfqpoint{4.461532in}{0.914466in}}{\pgfqpoint{4.461532in}{0.923706in}}%
\pgfpathlineto{\pgfqpoint{4.461532in}{0.959443in}}%
\pgfpathlineto{\pgfqpoint{4.471889in}{0.959443in}}%
\pgfpathlineto{\pgfqpoint{4.471889in}{0.896400in}}%
\pgfpathlineto{\pgfqpoint{4.461532in}{0.896400in}}%
\pgfpathlineto{\pgfqpoint{4.461532in}{0.906091in}}%
\pgfpathquadraticcurveto{\pgfqpoint{4.457768in}{0.900345in}}{\pgfqpoint{4.452778in}{0.897553in}}%
\pgfpathquadraticcurveto{\pgfqpoint{4.447807in}{0.894761in}}{\pgfqpoint{4.441214in}{0.894761in}}%
\pgfpathquadraticcurveto{\pgfqpoint{4.430353in}{0.894761in}}{\pgfqpoint{4.424715in}{0.901515in}}%
\pgfpathquadraticcurveto{\pgfqpoint{4.419096in}{0.908270in}}{\pgfqpoint{4.419096in}{0.921275in}}%
\pgfpathlineto{\pgfqpoint{4.419096in}{0.921275in}}%
\pgfpathclose%
\pgfpathmoveto{\pgfqpoint{4.445159in}{0.960956in}}%
\pgfpathlineto{\pgfqpoint{4.445159in}{0.960956in}}%
\pgfpathlineto{\pgfqpoint{4.445159in}{0.960956in}}%
\pgfpathclose%
\pgfpathmoveto{\pgfqpoint{4.533401in}{0.957587in}}%
\pgfpathlineto{\pgfqpoint{4.533401in}{0.947789in}}%
\pgfpathquadraticcurveto{\pgfqpoint{4.529024in}{0.950040in}}{\pgfqpoint{4.524287in}{0.951157in}}%
\pgfpathquadraticcurveto{\pgfqpoint{4.519567in}{0.952292in}}{\pgfqpoint{4.514488in}{0.952292in}}%
\pgfpathquadraticcurveto{\pgfqpoint{4.506779in}{0.952292in}}{\pgfqpoint{4.502924in}{0.949932in}}%
\pgfpathquadraticcurveto{\pgfqpoint{4.499070in}{0.947573in}}{\pgfqpoint{4.499070in}{0.942835in}}%
\pgfpathquadraticcurveto{\pgfqpoint{4.499070in}{0.939233in}}{\pgfqpoint{4.501825in}{0.937180in}}%
\pgfpathquadraticcurveto{\pgfqpoint{4.504581in}{0.935126in}}{\pgfqpoint{4.512921in}{0.933271in}}%
\pgfpathlineto{\pgfqpoint{4.516469in}{0.932478in}}%
\pgfpathquadraticcurveto{\pgfqpoint{4.527493in}{0.930119in}}{\pgfqpoint{4.532140in}{0.925814in}}%
\pgfpathquadraticcurveto{\pgfqpoint{4.536787in}{0.921509in}}{\pgfqpoint{4.536787in}{0.913800in}}%
\pgfpathquadraticcurveto{\pgfqpoint{4.536787in}{0.905010in}}{\pgfqpoint{4.529834in}{0.899876in}}%
\pgfpathquadraticcurveto{\pgfqpoint{4.522882in}{0.894761in}}{\pgfqpoint{4.510723in}{0.894761in}}%
\pgfpathquadraticcurveto{\pgfqpoint{4.505662in}{0.894761in}}{\pgfqpoint{4.500168in}{0.895752in}}%
\pgfpathquadraticcurveto{\pgfqpoint{4.494675in}{0.896742in}}{\pgfqpoint{4.488604in}{0.898706in}}%
\pgfpathlineto{\pgfqpoint{4.488604in}{0.909405in}}%
\pgfpathquadraticcurveto{\pgfqpoint{4.494350in}{0.906415in}}{\pgfqpoint{4.499916in}{0.904920in}}%
\pgfpathquadraticcurveto{\pgfqpoint{4.505482in}{0.903443in}}{\pgfqpoint{4.510958in}{0.903443in}}%
\pgfpathquadraticcurveto{\pgfqpoint{4.518270in}{0.903443in}}{\pgfqpoint{4.522197in}{0.905946in}}%
\pgfpathquadraticcurveto{\pgfqpoint{4.526142in}{0.908450in}}{\pgfqpoint{4.526142in}{0.913007in}}%
\pgfpathquadraticcurveto{\pgfqpoint{4.526142in}{0.917222in}}{\pgfqpoint{4.523296in}{0.919474in}}%
\pgfpathquadraticcurveto{\pgfqpoint{4.520468in}{0.921725in}}{\pgfqpoint{4.510831in}{0.923814in}}%
\pgfpathlineto{\pgfqpoint{4.507229in}{0.924661in}}%
\pgfpathquadraticcurveto{\pgfqpoint{4.497611in}{0.926678in}}{\pgfqpoint{4.493324in}{0.930875in}}%
\pgfpathquadraticcurveto{\pgfqpoint{4.489055in}{0.935072in}}{\pgfqpoint{4.489055in}{0.942385in}}%
\pgfpathquadraticcurveto{\pgfqpoint{4.489055in}{0.951283in}}{\pgfqpoint{4.495359in}{0.956110in}}%
\pgfpathquadraticcurveto{\pgfqpoint{4.501663in}{0.960956in}}{\pgfqpoint{4.513263in}{0.960956in}}%
\pgfpathquadraticcurveto{\pgfqpoint{4.518991in}{0.960956in}}{\pgfqpoint{4.524052in}{0.960109in}}%
\pgfpathquadraticcurveto{\pgfqpoint{4.529132in}{0.959280in}}{\pgfqpoint{4.533401in}{0.957587in}}%
\pgfpathlineto{\pgfqpoint{4.533401in}{0.957587in}}%
\pgfpathclose%
\pgfpathmoveto{\pgfqpoint{4.563517in}{0.977347in}}%
\pgfpathlineto{\pgfqpoint{4.563517in}{0.959443in}}%
\pgfpathlineto{\pgfqpoint{4.584843in}{0.959443in}}%
\pgfpathlineto{\pgfqpoint{4.584843in}{0.951391in}}%
\pgfpathlineto{\pgfqpoint{4.563517in}{0.951391in}}%
\pgfpathlineto{\pgfqpoint{4.563517in}{0.917168in}}%
\pgfpathquadraticcurveto{\pgfqpoint{4.563517in}{0.909459in}}{\pgfqpoint{4.565624in}{0.907261in}}%
\pgfpathquadraticcurveto{\pgfqpoint{4.567732in}{0.905064in}}{\pgfqpoint{4.574216in}{0.905064in}}%
\pgfpathlineto{\pgfqpoint{4.584843in}{0.905064in}}%
\pgfpathlineto{\pgfqpoint{4.584843in}{0.896400in}}%
\pgfpathlineto{\pgfqpoint{4.574216in}{0.896400in}}%
\pgfpathquadraticcurveto{\pgfqpoint{4.562220in}{0.896400in}}{\pgfqpoint{4.557663in}{0.900867in}}%
\pgfpathquadraticcurveto{\pgfqpoint{4.553106in}{0.905352in}}{\pgfqpoint{4.553106in}{0.917168in}}%
\pgfpathlineto{\pgfqpoint{4.553106in}{0.951391in}}%
\pgfpathlineto{\pgfqpoint{4.545505in}{0.951391in}}%
\pgfpathlineto{\pgfqpoint{4.545505in}{0.959443in}}%
\pgfpathlineto{\pgfqpoint{4.553106in}{0.959443in}}%
\pgfpathlineto{\pgfqpoint{4.553106in}{0.977347in}}%
\pgfpathlineto{\pgfqpoint{4.563517in}{0.977347in}}%
\pgfpathlineto{\pgfqpoint{4.563517in}{0.977347in}}%
\pgfpathclose%
\pgfpathmoveto{\pgfqpoint{4.652389in}{0.930515in}}%
\pgfpathlineto{\pgfqpoint{4.652389in}{0.925454in}}%
\pgfpathlineto{\pgfqpoint{4.604765in}{0.925454in}}%
\pgfpathquadraticcurveto{\pgfqpoint{4.605449in}{0.914754in}}{\pgfqpoint{4.611213in}{0.909153in}}%
\pgfpathquadraticcurveto{\pgfqpoint{4.616977in}{0.903551in}}{\pgfqpoint{4.627280in}{0.903551in}}%
\pgfpathquadraticcurveto{\pgfqpoint{4.633242in}{0.903551in}}{\pgfqpoint{4.638844in}{0.905010in}}%
\pgfpathquadraticcurveto{\pgfqpoint{4.644446in}{0.906469in}}{\pgfqpoint{4.649975in}{0.909405in}}%
\pgfpathlineto{\pgfqpoint{4.649975in}{0.899606in}}%
\pgfpathquadraticcurveto{\pgfqpoint{4.644392in}{0.897247in}}{\pgfqpoint{4.638538in}{0.896004in}}%
\pgfpathquadraticcurveto{\pgfqpoint{4.632684in}{0.894761in}}{\pgfqpoint{4.626668in}{0.894761in}}%
\pgfpathquadraticcurveto{\pgfqpoint{4.611573in}{0.894761in}}{\pgfqpoint{4.602765in}{0.903533in}}%
\pgfpathquadraticcurveto{\pgfqpoint{4.593958in}{0.912323in}}{\pgfqpoint{4.593958in}{0.927309in}}%
\pgfpathquadraticcurveto{\pgfqpoint{4.593958in}{0.942781in}}{\pgfqpoint{4.602315in}{0.951859in}}%
\pgfpathquadraticcurveto{\pgfqpoint{4.610673in}{0.960956in}}{\pgfqpoint{4.624866in}{0.960956in}}%
\pgfpathquadraticcurveto{\pgfqpoint{4.637583in}{0.960956in}}{\pgfqpoint{4.644986in}{0.952760in}}%
\pgfpathquadraticcurveto{\pgfqpoint{4.652389in}{0.944583in}}{\pgfqpoint{4.652389in}{0.930515in}}%
\pgfpathlineto{\pgfqpoint{4.652389in}{0.930515in}}%
\pgfpathclose%
\pgfpathmoveto{\pgfqpoint{4.642032in}{0.933559in}}%
\pgfpathquadraticcurveto{\pgfqpoint{4.641924in}{0.942043in}}{\pgfqpoint{4.637277in}{0.947104in}}%
\pgfpathquadraticcurveto{\pgfqpoint{4.632630in}{0.952184in}}{\pgfqpoint{4.624974in}{0.952184in}}%
\pgfpathquadraticcurveto{\pgfqpoint{4.616311in}{0.952184in}}{\pgfqpoint{4.611105in}{0.947284in}}%
\pgfpathquadraticcurveto{\pgfqpoint{4.605900in}{0.942385in}}{\pgfqpoint{4.605107in}{0.933487in}}%
\pgfpathlineto{\pgfqpoint{4.642032in}{0.933559in}}%
\pgfpathlineto{\pgfqpoint{4.642032in}{0.933559in}}%
\pgfpathclose%
\pgfpathmoveto{\pgfqpoint{4.705921in}{0.949770in}}%
\pgfpathquadraticcurveto{\pgfqpoint{4.704174in}{0.950779in}}{\pgfqpoint{4.702121in}{0.951247in}}%
\pgfpathquadraticcurveto{\pgfqpoint{4.700067in}{0.951733in}}{\pgfqpoint{4.697599in}{0.951733in}}%
\pgfpathquadraticcurveto{\pgfqpoint{4.688810in}{0.951733in}}{\pgfqpoint{4.684108in}{0.946023in}}%
\pgfpathquadraticcurveto{\pgfqpoint{4.679407in}{0.940314in}}{\pgfqpoint{4.679407in}{0.929614in}}%
\pgfpathlineto{\pgfqpoint{4.679407in}{0.896400in}}%
\pgfpathlineto{\pgfqpoint{4.668996in}{0.896400in}}%
\pgfpathlineto{\pgfqpoint{4.668996in}{0.959443in}}%
\pgfpathlineto{\pgfqpoint{4.679407in}{0.959443in}}%
\pgfpathlineto{\pgfqpoint{4.679407in}{0.949644in}}%
\pgfpathquadraticcurveto{\pgfqpoint{4.682685in}{0.955390in}}{\pgfqpoint{4.687909in}{0.958164in}}%
\pgfpathquadraticcurveto{\pgfqpoint{4.693150in}{0.960956in}}{\pgfqpoint{4.700644in}{0.960956in}}%
\pgfpathquadraticcurveto{\pgfqpoint{4.701706in}{0.960956in}}{\pgfqpoint{4.703003in}{0.960811in}}%
\pgfpathquadraticcurveto{\pgfqpoint{4.704300in}{0.960685in}}{\pgfqpoint{4.705867in}{0.960397in}}%
\pgfpathlineto{\pgfqpoint{4.705921in}{0.949770in}}%
\pgfpathlineto{\pgfqpoint{4.705921in}{0.949770in}}%
\pgfpathclose%
\pgfpathmoveto{\pgfqpoint{4.756968in}{0.957587in}}%
\pgfpathlineto{\pgfqpoint{4.756968in}{0.947789in}}%
\pgfpathquadraticcurveto{\pgfqpoint{4.752591in}{0.950040in}}{\pgfqpoint{4.747853in}{0.951157in}}%
\pgfpathquadraticcurveto{\pgfqpoint{4.743134in}{0.952292in}}{\pgfqpoint{4.738055in}{0.952292in}}%
\pgfpathquadraticcurveto{\pgfqpoint{4.730346in}{0.952292in}}{\pgfqpoint{4.726491in}{0.949932in}}%
\pgfpathquadraticcurveto{\pgfqpoint{4.722636in}{0.947573in}}{\pgfqpoint{4.722636in}{0.942835in}}%
\pgfpathquadraticcurveto{\pgfqpoint{4.722636in}{0.939233in}}{\pgfqpoint{4.725392in}{0.937180in}}%
\pgfpathquadraticcurveto{\pgfqpoint{4.728148in}{0.935126in}}{\pgfqpoint{4.736488in}{0.933271in}}%
\pgfpathlineto{\pgfqpoint{4.740036in}{0.932478in}}%
\pgfpathquadraticcurveto{\pgfqpoint{4.751060in}{0.930119in}}{\pgfqpoint{4.755707in}{0.925814in}}%
\pgfpathquadraticcurveto{\pgfqpoint{4.760354in}{0.921509in}}{\pgfqpoint{4.760354in}{0.913800in}}%
\pgfpathquadraticcurveto{\pgfqpoint{4.760354in}{0.905010in}}{\pgfqpoint{4.753401in}{0.899876in}}%
\pgfpathquadraticcurveto{\pgfqpoint{4.746448in}{0.894761in}}{\pgfqpoint{4.734290in}{0.894761in}}%
\pgfpathquadraticcurveto{\pgfqpoint{4.729229in}{0.894761in}}{\pgfqpoint{4.723735in}{0.895752in}}%
\pgfpathquadraticcurveto{\pgfqpoint{4.718241in}{0.896742in}}{\pgfqpoint{4.712171in}{0.898706in}}%
\pgfpathlineto{\pgfqpoint{4.712171in}{0.909405in}}%
\pgfpathquadraticcurveto{\pgfqpoint{4.717917in}{0.906415in}}{\pgfqpoint{4.723483in}{0.904920in}}%
\pgfpathquadraticcurveto{\pgfqpoint{4.729049in}{0.903443in}}{\pgfqpoint{4.734524in}{0.903443in}}%
\pgfpathquadraticcurveto{\pgfqpoint{4.741837in}{0.903443in}}{\pgfqpoint{4.745764in}{0.905946in}}%
\pgfpathquadraticcurveto{\pgfqpoint{4.749709in}{0.908450in}}{\pgfqpoint{4.749709in}{0.913007in}}%
\pgfpathquadraticcurveto{\pgfqpoint{4.749709in}{0.917222in}}{\pgfqpoint{4.746863in}{0.919474in}}%
\pgfpathquadraticcurveto{\pgfqpoint{4.744035in}{0.921725in}}{\pgfqpoint{4.734398in}{0.923814in}}%
\pgfpathlineto{\pgfqpoint{4.730796in}{0.924661in}}%
\pgfpathquadraticcurveto{\pgfqpoint{4.721177in}{0.926678in}}{\pgfqpoint{4.716890in}{0.930875in}}%
\pgfpathquadraticcurveto{\pgfqpoint{4.712622in}{0.935072in}}{\pgfqpoint{4.712622in}{0.942385in}}%
\pgfpathquadraticcurveto{\pgfqpoint{4.712622in}{0.951283in}}{\pgfqpoint{4.718926in}{0.956110in}}%
\pgfpathquadraticcurveto{\pgfqpoint{4.725230in}{0.960956in}}{\pgfqpoint{4.736830in}{0.960956in}}%
\pgfpathquadraticcurveto{\pgfqpoint{4.742558in}{0.960956in}}{\pgfqpoint{4.747619in}{0.960109in}}%
\pgfpathquadraticcurveto{\pgfqpoint{4.752699in}{0.959280in}}{\pgfqpoint{4.756968in}{0.957587in}}%
\pgfpathlineto{\pgfqpoint{4.756968in}{0.957587in}}%
\pgfpathclose%
\pgfusepath{fill}%
\end{pgfscope}%
\begin{pgfscope}%
\pgfsetbuttcap%
\pgfsetroundjoin%
\definecolor{currentfill}{rgb}{0.309804,0.372549,0.419608}%
\pgfsetfillcolor{currentfill}%
\pgfsetlinewidth{0.000000pt}%
\definecolor{currentstroke}{rgb}{0.309804,0.372549,0.419608}%
\pgfsetstrokecolor{currentstroke}%
\pgfsetdash{}{0pt}%
\pgfpathmoveto{\pgfqpoint{4.915739in}{0.953138in}}%
\pgfpathlineto{\pgfqpoint{4.857650in}{0.932478in}}%
\pgfpathlineto{\pgfqpoint{4.915739in}{0.911944in}}%
\pgfpathlineto{\pgfqpoint{4.915739in}{0.901696in}}%
\pgfpathlineto{\pgfqpoint{4.843583in}{0.927867in}}%
\pgfpathlineto{\pgfqpoint{4.843583in}{0.937216in}}%
\pgfpathlineto{\pgfqpoint{4.915739in}{0.963387in}}%
\pgfpathlineto{\pgfqpoint{4.915739in}{0.953138in}}%
\pgfpathlineto{\pgfqpoint{4.915739in}{0.953138in}}%
\pgfpathclose%
\pgfpathmoveto{\pgfqpoint{4.942271in}{0.905964in}}%
\pgfpathlineto{\pgfqpoint{4.960842in}{0.905964in}}%
\pgfpathlineto{\pgfqpoint{4.960842in}{0.970088in}}%
\pgfpathlineto{\pgfqpoint{4.940632in}{0.966035in}}%
\pgfpathlineto{\pgfqpoint{4.940632in}{0.976392in}}%
\pgfpathlineto{\pgfqpoint{4.960734in}{0.980445in}}%
\pgfpathlineto{\pgfqpoint{4.972099in}{0.980445in}}%
\pgfpathlineto{\pgfqpoint{4.972099in}{0.905964in}}%
\pgfpathlineto{\pgfqpoint{4.990670in}{0.905964in}}%
\pgfpathlineto{\pgfqpoint{4.990670in}{0.896400in}}%
\pgfpathlineto{\pgfqpoint{4.942271in}{0.896400in}}%
\pgfpathlineto{\pgfqpoint{4.942271in}{0.905964in}}%
\pgfpathlineto{\pgfqpoint{4.942271in}{0.905964in}}%
\pgfpathclose%
\pgfpathmoveto{\pgfqpoint{5.037952in}{0.972952in}}%
\pgfpathquadraticcurveto{\pgfqpoint{5.029180in}{0.972952in}}{\pgfqpoint{5.024749in}{0.964306in}}%
\pgfpathquadraticcurveto{\pgfqpoint{5.020336in}{0.955678in}}{\pgfqpoint{5.020336in}{0.938332in}}%
\pgfpathquadraticcurveto{\pgfqpoint{5.020336in}{0.921059in}}{\pgfqpoint{5.024749in}{0.912413in}}%
\pgfpathquadraticcurveto{\pgfqpoint{5.029180in}{0.903767in}}{\pgfqpoint{5.037952in}{0.903767in}}%
\pgfpathquadraticcurveto{\pgfqpoint{5.046796in}{0.903767in}}{\pgfqpoint{5.051209in}{0.912413in}}%
\pgfpathquadraticcurveto{\pgfqpoint{5.055640in}{0.921059in}}{\pgfqpoint{5.055640in}{0.938332in}}%
\pgfpathquadraticcurveto{\pgfqpoint{5.055640in}{0.955678in}}{\pgfqpoint{5.051209in}{0.964306in}}%
\pgfpathquadraticcurveto{\pgfqpoint{5.046796in}{0.972952in}}{\pgfqpoint{5.037952in}{0.972952in}}%
\pgfpathlineto{\pgfqpoint{5.037952in}{0.972952in}}%
\pgfpathclose%
\pgfpathmoveto{\pgfqpoint{5.037952in}{0.981958in}}%
\pgfpathquadraticcurveto{\pgfqpoint{5.052091in}{0.981958in}}{\pgfqpoint{5.059548in}{0.970772in}}%
\pgfpathquadraticcurveto{\pgfqpoint{5.067005in}{0.959605in}}{\pgfqpoint{5.067005in}{0.938332in}}%
\pgfpathquadraticcurveto{\pgfqpoint{5.067005in}{0.917114in}}{\pgfqpoint{5.059548in}{0.905928in}}%
\pgfpathquadraticcurveto{\pgfqpoint{5.052091in}{0.894761in}}{\pgfqpoint{5.037952in}{0.894761in}}%
\pgfpathquadraticcurveto{\pgfqpoint{5.023830in}{0.894761in}}{\pgfqpoint{5.016373in}{0.905928in}}%
\pgfpathquadraticcurveto{\pgfqpoint{5.008916in}{0.917114in}}{\pgfqpoint{5.008916in}{0.938332in}}%
\pgfpathquadraticcurveto{\pgfqpoint{5.008916in}{0.959605in}}{\pgfqpoint{5.016373in}{0.970772in}}%
\pgfpathquadraticcurveto{\pgfqpoint{5.023830in}{0.981958in}}{\pgfqpoint{5.037952in}{0.981958in}}%
\pgfpathlineto{\pgfqpoint{5.037952in}{0.981958in}}%
\pgfpathclose%
\pgfusepath{fill}%
\end{pgfscope}%
\begin{pgfscope}%
\pgfsetbuttcap%
\pgfsetroundjoin%
\definecolor{currentfill}{rgb}{0.309804,0.372549,0.419608}%
\pgfsetfillcolor{currentfill}%
\pgfsetlinewidth{0.000000pt}%
\definecolor{currentstroke}{rgb}{0.309804,0.372549,0.419608}%
\pgfsetstrokecolor{currentstroke}%
\pgfsetdash{}{0pt}%
\pgfpathmoveto{\pgfqpoint{5.174140in}{0.952184in}}%
\pgfpathquadraticcurveto{\pgfqpoint{5.165818in}{0.952184in}}{\pgfqpoint{5.160973in}{0.945681in}}%
\pgfpathquadraticcurveto{\pgfqpoint{5.156128in}{0.939179in}}{\pgfqpoint{5.156128in}{0.927867in}}%
\pgfpathquadraticcurveto{\pgfqpoint{5.156128in}{0.916556in}}{\pgfqpoint{5.160937in}{0.910053in}}%
\pgfpathquadraticcurveto{\pgfqpoint{5.165764in}{0.903551in}}{\pgfqpoint{5.174140in}{0.903551in}}%
\pgfpathquadraticcurveto{\pgfqpoint{5.182426in}{0.903551in}}{\pgfqpoint{5.187253in}{0.910071in}}%
\pgfpathquadraticcurveto{\pgfqpoint{5.192098in}{0.916610in}}{\pgfqpoint{5.192098in}{0.927867in}}%
\pgfpathquadraticcurveto{\pgfqpoint{5.192098in}{0.939071in}}{\pgfqpoint{5.187253in}{0.945627in}}%
\pgfpathquadraticcurveto{\pgfqpoint{5.182426in}{0.952184in}}{\pgfqpoint{5.174140in}{0.952184in}}%
\pgfpathlineto{\pgfqpoint{5.174140in}{0.952184in}}%
\pgfpathclose%
\pgfpathmoveto{\pgfqpoint{5.174140in}{0.960956in}}%
\pgfpathquadraticcurveto{\pgfqpoint{5.187649in}{0.960956in}}{\pgfqpoint{5.195358in}{0.952166in}}%
\pgfpathquadraticcurveto{\pgfqpoint{5.203086in}{0.943394in}}{\pgfqpoint{5.203086in}{0.927867in}}%
\pgfpathquadraticcurveto{\pgfqpoint{5.203086in}{0.912395in}}{\pgfqpoint{5.195358in}{0.903569in}}%
\pgfpathquadraticcurveto{\pgfqpoint{5.187649in}{0.894761in}}{\pgfqpoint{5.174140in}{0.894761in}}%
\pgfpathquadraticcurveto{\pgfqpoint{5.160577in}{0.894761in}}{\pgfqpoint{5.152886in}{0.903569in}}%
\pgfpathquadraticcurveto{\pgfqpoint{5.145212in}{0.912395in}}{\pgfqpoint{5.145212in}{0.927867in}}%
\pgfpathquadraticcurveto{\pgfqpoint{5.145212in}{0.943394in}}{\pgfqpoint{5.152886in}{0.952166in}}%
\pgfpathquadraticcurveto{\pgfqpoint{5.160577in}{0.960956in}}{\pgfqpoint{5.174140in}{0.960956in}}%
\pgfpathlineto{\pgfqpoint{5.174140in}{0.960956in}}%
\pgfpathclose%
\pgfpathmoveto{\pgfqpoint{5.256780in}{0.949770in}}%
\pgfpathquadraticcurveto{\pgfqpoint{5.255033in}{0.950779in}}{\pgfqpoint{5.252979in}{0.951247in}}%
\pgfpathquadraticcurveto{\pgfqpoint{5.250926in}{0.951733in}}{\pgfqpoint{5.248458in}{0.951733in}}%
\pgfpathquadraticcurveto{\pgfqpoint{5.239668in}{0.951733in}}{\pgfqpoint{5.234967in}{0.946023in}}%
\pgfpathquadraticcurveto{\pgfqpoint{5.230266in}{0.940314in}}{\pgfqpoint{5.230266in}{0.929614in}}%
\pgfpathlineto{\pgfqpoint{5.230266in}{0.896400in}}%
\pgfpathlineto{\pgfqpoint{5.219855in}{0.896400in}}%
\pgfpathlineto{\pgfqpoint{5.219855in}{0.959443in}}%
\pgfpathlineto{\pgfqpoint{5.230266in}{0.959443in}}%
\pgfpathlineto{\pgfqpoint{5.230266in}{0.949644in}}%
\pgfpathquadraticcurveto{\pgfqpoint{5.233544in}{0.955390in}}{\pgfqpoint{5.238768in}{0.958164in}}%
\pgfpathquadraticcurveto{\pgfqpoint{5.244009in}{0.960956in}}{\pgfqpoint{5.251502in}{0.960956in}}%
\pgfpathquadraticcurveto{\pgfqpoint{5.252565in}{0.960956in}}{\pgfqpoint{5.253862in}{0.960811in}}%
\pgfpathquadraticcurveto{\pgfqpoint{5.255159in}{0.960685in}}{\pgfqpoint{5.256726in}{0.960397in}}%
\pgfpathlineto{\pgfqpoint{5.256780in}{0.949770in}}%
\pgfpathlineto{\pgfqpoint{5.256780in}{0.949770in}}%
\pgfpathclose%
\pgfusepath{fill}%
\end{pgfscope}%
\begin{pgfscope}%
\pgfsetbuttcap%
\pgfsetroundjoin%
\definecolor{currentfill}{rgb}{0.309804,0.372549,0.419608}%
\pgfsetfillcolor{currentfill}%
\pgfsetlinewidth{0.000000pt}%
\definecolor{currentstroke}{rgb}{0.309804,0.372549,0.419608}%
\pgfsetstrokecolor{currentstroke}%
\pgfsetdash{}{0pt}%
\pgfpathmoveto{\pgfqpoint{5.366282in}{0.947338in}}%
\pgfpathquadraticcurveto{\pgfqpoint{5.370173in}{0.954327in}}{\pgfqpoint{5.375577in}{0.957641in}}%
\pgfpathquadraticcurveto{\pgfqpoint{5.380980in}{0.960956in}}{\pgfqpoint{5.388293in}{0.960956in}}%
\pgfpathquadraticcurveto{\pgfqpoint{5.398146in}{0.960956in}}{\pgfqpoint{5.403496in}{0.954057in}}%
\pgfpathquadraticcurveto{\pgfqpoint{5.408845in}{0.947176in}}{\pgfqpoint{5.408845in}{0.934460in}}%
\pgfpathlineto{\pgfqpoint{5.408845in}{0.896400in}}%
\pgfpathlineto{\pgfqpoint{5.398434in}{0.896400in}}%
\pgfpathlineto{\pgfqpoint{5.398434in}{0.934117in}}%
\pgfpathquadraticcurveto{\pgfqpoint{5.398434in}{0.943178in}}{\pgfqpoint{5.395210in}{0.947555in}}%
\pgfpathquadraticcurveto{\pgfqpoint{5.392004in}{0.951949in}}{\pgfqpoint{5.385429in}{0.951949in}}%
\pgfpathquadraticcurveto{\pgfqpoint{5.377378in}{0.951949in}}{\pgfqpoint{5.372695in}{0.946600in}}%
\pgfpathquadraticcurveto{\pgfqpoint{5.368030in}{0.941268in}}{\pgfqpoint{5.368030in}{0.932028in}}%
\pgfpathlineto{\pgfqpoint{5.368030in}{0.896400in}}%
\pgfpathlineto{\pgfqpoint{5.357619in}{0.896400in}}%
\pgfpathlineto{\pgfqpoint{5.357619in}{0.934117in}}%
\pgfpathquadraticcurveto{\pgfqpoint{5.357619in}{0.943232in}}{\pgfqpoint{5.354412in}{0.947591in}}%
\pgfpathquadraticcurveto{\pgfqpoint{5.351206in}{0.951949in}}{\pgfqpoint{5.344506in}{0.951949in}}%
\pgfpathquadraticcurveto{\pgfqpoint{5.336562in}{0.951949in}}{\pgfqpoint{5.331879in}{0.946582in}}%
\pgfpathquadraticcurveto{\pgfqpoint{5.327214in}{0.941214in}}{\pgfqpoint{5.327214in}{0.932028in}}%
\pgfpathlineto{\pgfqpoint{5.327214in}{0.896400in}}%
\pgfpathlineto{\pgfqpoint{5.316803in}{0.896400in}}%
\pgfpathlineto{\pgfqpoint{5.316803in}{0.959443in}}%
\pgfpathlineto{\pgfqpoint{5.327214in}{0.959443in}}%
\pgfpathlineto{\pgfqpoint{5.327214in}{0.949644in}}%
\pgfpathquadraticcurveto{\pgfqpoint{5.330762in}{0.955444in}}{\pgfqpoint{5.335716in}{0.958200in}}%
\pgfpathquadraticcurveto{\pgfqpoint{5.340669in}{0.960956in}}{\pgfqpoint{5.347478in}{0.960956in}}%
\pgfpathquadraticcurveto{\pgfqpoint{5.354358in}{0.960956in}}{\pgfqpoint{5.359168in}{0.957461in}}%
\pgfpathquadraticcurveto{\pgfqpoint{5.363977in}{0.953985in}}{\pgfqpoint{5.366282in}{0.947338in}}%
\pgfpathlineto{\pgfqpoint{5.366282in}{0.947338in}}%
\pgfpathclose%
\pgfpathmoveto{\pgfqpoint{5.458144in}{0.928083in}}%
\pgfpathquadraticcurveto{\pgfqpoint{5.445590in}{0.928083in}}{\pgfqpoint{5.440745in}{0.925219in}}%
\pgfpathquadraticcurveto{\pgfqpoint{5.435899in}{0.922356in}}{\pgfqpoint{5.435899in}{0.915421in}}%
\pgfpathquadraticcurveto{\pgfqpoint{5.435899in}{0.909909in}}{\pgfqpoint{5.439538in}{0.906667in}}%
\pgfpathquadraticcurveto{\pgfqpoint{5.443176in}{0.903443in}}{\pgfqpoint{5.449408in}{0.903443in}}%
\pgfpathquadraticcurveto{\pgfqpoint{5.458036in}{0.903443in}}{\pgfqpoint{5.463242in}{0.909549in}}%
\pgfpathquadraticcurveto{\pgfqpoint{5.468447in}{0.915655in}}{\pgfqpoint{5.468447in}{0.925778in}}%
\pgfpathlineto{\pgfqpoint{5.468447in}{0.928083in}}%
\pgfpathlineto{\pgfqpoint{5.458144in}{0.928083in}}%
\pgfpathlineto{\pgfqpoint{5.458144in}{0.928083in}}%
\pgfpathclose%
\pgfpathmoveto{\pgfqpoint{5.478804in}{0.932370in}}%
\pgfpathlineto{\pgfqpoint{5.478804in}{0.896400in}}%
\pgfpathlineto{\pgfqpoint{5.468447in}{0.896400in}}%
\pgfpathlineto{\pgfqpoint{5.468447in}{0.905964in}}%
\pgfpathquadraticcurveto{\pgfqpoint{5.464899in}{0.900237in}}{\pgfqpoint{5.459603in}{0.897499in}}%
\pgfpathquadraticcurveto{\pgfqpoint{5.454308in}{0.894761in}}{\pgfqpoint{5.446653in}{0.894761in}}%
\pgfpathquadraticcurveto{\pgfqpoint{5.436980in}{0.894761in}}{\pgfqpoint{5.431252in}{0.900201in}}%
\pgfpathquadraticcurveto{\pgfqpoint{5.425542in}{0.905640in}}{\pgfqpoint{5.425542in}{0.914754in}}%
\pgfpathquadraticcurveto{\pgfqpoint{5.425542in}{0.925382in}}{\pgfqpoint{5.432657in}{0.930785in}}%
\pgfpathquadraticcurveto{\pgfqpoint{5.439790in}{0.936189in}}{\pgfqpoint{5.453912in}{0.936189in}}%
\pgfpathlineto{\pgfqpoint{5.468447in}{0.936189in}}%
\pgfpathlineto{\pgfqpoint{5.468447in}{0.937216in}}%
\pgfpathquadraticcurveto{\pgfqpoint{5.468447in}{0.944366in}}{\pgfqpoint{5.463746in}{0.948275in}}%
\pgfpathquadraticcurveto{\pgfqpoint{5.459045in}{0.952184in}}{\pgfqpoint{5.450543in}{0.952184in}}%
\pgfpathquadraticcurveto{\pgfqpoint{5.445140in}{0.952184in}}{\pgfqpoint{5.440006in}{0.950887in}}%
\pgfpathquadraticcurveto{\pgfqpoint{5.434891in}{0.949590in}}{\pgfqpoint{5.430172in}{0.946996in}}%
\pgfpathlineto{\pgfqpoint{5.430172in}{0.956579in}}%
\pgfpathquadraticcurveto{\pgfqpoint{5.435845in}{0.958776in}}{\pgfqpoint{5.441195in}{0.959857in}}%
\pgfpathquadraticcurveto{\pgfqpoint{5.446545in}{0.960956in}}{\pgfqpoint{5.451606in}{0.960956in}}%
\pgfpathquadraticcurveto{\pgfqpoint{5.465295in}{0.960956in}}{\pgfqpoint{5.472050in}{0.953859in}}%
\pgfpathquadraticcurveto{\pgfqpoint{5.478804in}{0.946780in}}{\pgfqpoint{5.478804in}{0.932370in}}%
\pgfpathlineto{\pgfqpoint{5.478804in}{0.932370in}}%
\pgfpathclose%
\pgfpathmoveto{\pgfqpoint{5.552546in}{0.959443in}}%
\pgfpathlineto{\pgfqpoint{5.529743in}{0.928768in}}%
\pgfpathlineto{\pgfqpoint{5.553717in}{0.896400in}}%
\pgfpathlineto{\pgfqpoint{5.541505in}{0.896400in}}%
\pgfpathlineto{\pgfqpoint{5.523150in}{0.921167in}}%
\pgfpathlineto{\pgfqpoint{5.504814in}{0.896400in}}%
\pgfpathlineto{\pgfqpoint{5.492584in}{0.896400in}}%
\pgfpathlineto{\pgfqpoint{5.517080in}{0.929380in}}%
\pgfpathlineto{\pgfqpoint{5.494673in}{0.959443in}}%
\pgfpathlineto{\pgfqpoint{5.506885in}{0.959443in}}%
\pgfpathlineto{\pgfqpoint{5.523601in}{0.936981in}}%
\pgfpathlineto{\pgfqpoint{5.540316in}{0.959443in}}%
\pgfpathlineto{\pgfqpoint{5.552546in}{0.959443in}}%
\pgfpathlineto{\pgfqpoint{5.552546in}{0.959443in}}%
\pgfpathclose%
\pgfpathmoveto{\pgfqpoint{5.568361in}{0.959443in}}%
\pgfpathlineto{\pgfqpoint{5.578718in}{0.959443in}}%
\pgfpathlineto{\pgfqpoint{5.578718in}{0.896400in}}%
\pgfpathlineto{\pgfqpoint{5.568361in}{0.896400in}}%
\pgfpathlineto{\pgfqpoint{5.568361in}{0.959443in}}%
\pgfpathlineto{\pgfqpoint{5.568361in}{0.959443in}}%
\pgfpathclose%
\pgfpathmoveto{\pgfqpoint{5.568361in}{0.983993in}}%
\pgfpathlineto{\pgfqpoint{5.578718in}{0.983993in}}%
\pgfpathlineto{\pgfqpoint{5.578718in}{0.970862in}}%
\pgfpathlineto{\pgfqpoint{5.568361in}{0.970862in}}%
\pgfpathlineto{\pgfqpoint{5.568361in}{0.983993in}}%
\pgfpathlineto{\pgfqpoint{5.568361in}{0.983993in}}%
\pgfpathclose%
\pgfpathmoveto{\pgfqpoint{5.649469in}{0.947338in}}%
\pgfpathquadraticcurveto{\pgfqpoint{5.653360in}{0.954327in}}{\pgfqpoint{5.658764in}{0.957641in}}%
\pgfpathquadraticcurveto{\pgfqpoint{5.664167in}{0.960956in}}{\pgfqpoint{5.671480in}{0.960956in}}%
\pgfpathquadraticcurveto{\pgfqpoint{5.681333in}{0.960956in}}{\pgfqpoint{5.686683in}{0.954057in}}%
\pgfpathquadraticcurveto{\pgfqpoint{5.692032in}{0.947176in}}{\pgfqpoint{5.692032in}{0.934460in}}%
\pgfpathlineto{\pgfqpoint{5.692032in}{0.896400in}}%
\pgfpathlineto{\pgfqpoint{5.681621in}{0.896400in}}%
\pgfpathlineto{\pgfqpoint{5.681621in}{0.934117in}}%
\pgfpathquadraticcurveto{\pgfqpoint{5.681621in}{0.943178in}}{\pgfqpoint{5.678397in}{0.947555in}}%
\pgfpathquadraticcurveto{\pgfqpoint{5.675191in}{0.951949in}}{\pgfqpoint{5.668616in}{0.951949in}}%
\pgfpathquadraticcurveto{\pgfqpoint{5.660565in}{0.951949in}}{\pgfqpoint{5.655882in}{0.946600in}}%
\pgfpathquadraticcurveto{\pgfqpoint{5.651217in}{0.941268in}}{\pgfqpoint{5.651217in}{0.932028in}}%
\pgfpathlineto{\pgfqpoint{5.651217in}{0.896400in}}%
\pgfpathlineto{\pgfqpoint{5.640806in}{0.896400in}}%
\pgfpathlineto{\pgfqpoint{5.640806in}{0.934117in}}%
\pgfpathquadraticcurveto{\pgfqpoint{5.640806in}{0.943232in}}{\pgfqpoint{5.637599in}{0.947591in}}%
\pgfpathquadraticcurveto{\pgfqpoint{5.634393in}{0.951949in}}{\pgfqpoint{5.627693in}{0.951949in}}%
\pgfpathquadraticcurveto{\pgfqpoint{5.619749in}{0.951949in}}{\pgfqpoint{5.615066in}{0.946582in}}%
\pgfpathquadraticcurveto{\pgfqpoint{5.610401in}{0.941214in}}{\pgfqpoint{5.610401in}{0.932028in}}%
\pgfpathlineto{\pgfqpoint{5.610401in}{0.896400in}}%
\pgfpathlineto{\pgfqpoint{5.599990in}{0.896400in}}%
\pgfpathlineto{\pgfqpoint{5.599990in}{0.959443in}}%
\pgfpathlineto{\pgfqpoint{5.610401in}{0.959443in}}%
\pgfpathlineto{\pgfqpoint{5.610401in}{0.949644in}}%
\pgfpathquadraticcurveto{\pgfqpoint{5.613950in}{0.955444in}}{\pgfqpoint{5.618903in}{0.958200in}}%
\pgfpathquadraticcurveto{\pgfqpoint{5.623856in}{0.960956in}}{\pgfqpoint{5.630665in}{0.960956in}}%
\pgfpathquadraticcurveto{\pgfqpoint{5.637545in}{0.960956in}}{\pgfqpoint{5.642355in}{0.957461in}}%
\pgfpathquadraticcurveto{\pgfqpoint{5.647164in}{0.953985in}}{\pgfqpoint{5.649469in}{0.947338in}}%
\pgfpathlineto{\pgfqpoint{5.649469in}{0.947338in}}%
\pgfpathclose%
\pgfpathmoveto{\pgfqpoint{5.711611in}{0.921275in}}%
\pgfpathlineto{\pgfqpoint{5.711611in}{0.959443in}}%
\pgfpathlineto{\pgfqpoint{5.721968in}{0.959443in}}%
\pgfpathlineto{\pgfqpoint{5.721968in}{0.921671in}}%
\pgfpathquadraticcurveto{\pgfqpoint{5.721968in}{0.912719in}}{\pgfqpoint{5.725445in}{0.908234in}}%
\pgfpathquadraticcurveto{\pgfqpoint{5.728939in}{0.903767in}}{\pgfqpoint{5.735928in}{0.903767in}}%
\pgfpathquadraticcurveto{\pgfqpoint{5.744303in}{0.903767in}}{\pgfqpoint{5.749167in}{0.909117in}}%
\pgfpathquadraticcurveto{\pgfqpoint{5.754048in}{0.914466in}}{\pgfqpoint{5.754048in}{0.923706in}}%
\pgfpathlineto{\pgfqpoint{5.754048in}{0.959443in}}%
\pgfpathlineto{\pgfqpoint{5.764405in}{0.959443in}}%
\pgfpathlineto{\pgfqpoint{5.764405in}{0.896400in}}%
\pgfpathlineto{\pgfqpoint{5.754048in}{0.896400in}}%
\pgfpathlineto{\pgfqpoint{5.754048in}{0.906091in}}%
\pgfpathquadraticcurveto{\pgfqpoint{5.750283in}{0.900345in}}{\pgfqpoint{5.745294in}{0.897553in}}%
\pgfpathquadraticcurveto{\pgfqpoint{5.740323in}{0.894761in}}{\pgfqpoint{5.733730in}{0.894761in}}%
\pgfpathquadraticcurveto{\pgfqpoint{5.722869in}{0.894761in}}{\pgfqpoint{5.717231in}{0.901515in}}%
\pgfpathquadraticcurveto{\pgfqpoint{5.711611in}{0.908270in}}{\pgfqpoint{5.711611in}{0.921275in}}%
\pgfpathlineto{\pgfqpoint{5.711611in}{0.921275in}}%
\pgfpathclose%
\pgfpathmoveto{\pgfqpoint{5.737675in}{0.960956in}}%
\pgfpathlineto{\pgfqpoint{5.737675in}{0.960956in}}%
\pgfpathlineto{\pgfqpoint{5.737675in}{0.960956in}}%
\pgfpathclose%
\pgfpathmoveto{\pgfqpoint{5.834815in}{0.947338in}}%
\pgfpathquadraticcurveto{\pgfqpoint{5.838705in}{0.954327in}}{\pgfqpoint{5.844109in}{0.957641in}}%
\pgfpathquadraticcurveto{\pgfqpoint{5.849512in}{0.960956in}}{\pgfqpoint{5.856825in}{0.960956in}}%
\pgfpathquadraticcurveto{\pgfqpoint{5.866678in}{0.960956in}}{\pgfqpoint{5.872028in}{0.954057in}}%
\pgfpathquadraticcurveto{\pgfqpoint{5.877377in}{0.947176in}}{\pgfqpoint{5.877377in}{0.934460in}}%
\pgfpathlineto{\pgfqpoint{5.877377in}{0.896400in}}%
\pgfpathlineto{\pgfqpoint{5.866966in}{0.896400in}}%
\pgfpathlineto{\pgfqpoint{5.866966in}{0.934117in}}%
\pgfpathquadraticcurveto{\pgfqpoint{5.866966in}{0.943178in}}{\pgfqpoint{5.863742in}{0.947555in}}%
\pgfpathquadraticcurveto{\pgfqpoint{5.860536in}{0.951949in}}{\pgfqpoint{5.853961in}{0.951949in}}%
\pgfpathquadraticcurveto{\pgfqpoint{5.845910in}{0.951949in}}{\pgfqpoint{5.841227in}{0.946600in}}%
\pgfpathquadraticcurveto{\pgfqpoint{5.836562in}{0.941268in}}{\pgfqpoint{5.836562in}{0.932028in}}%
\pgfpathlineto{\pgfqpoint{5.836562in}{0.896400in}}%
\pgfpathlineto{\pgfqpoint{5.826151in}{0.896400in}}%
\pgfpathlineto{\pgfqpoint{5.826151in}{0.934117in}}%
\pgfpathquadraticcurveto{\pgfqpoint{5.826151in}{0.943232in}}{\pgfqpoint{5.822945in}{0.947591in}}%
\pgfpathquadraticcurveto{\pgfqpoint{5.819738in}{0.951949in}}{\pgfqpoint{5.813038in}{0.951949in}}%
\pgfpathquadraticcurveto{\pgfqpoint{5.805094in}{0.951949in}}{\pgfqpoint{5.800411in}{0.946582in}}%
\pgfpathquadraticcurveto{\pgfqpoint{5.795746in}{0.941214in}}{\pgfqpoint{5.795746in}{0.932028in}}%
\pgfpathlineto{\pgfqpoint{5.795746in}{0.896400in}}%
\pgfpathlineto{\pgfqpoint{5.785335in}{0.896400in}}%
\pgfpathlineto{\pgfqpoint{5.785335in}{0.959443in}}%
\pgfpathlineto{\pgfqpoint{5.795746in}{0.959443in}}%
\pgfpathlineto{\pgfqpoint{5.795746in}{0.949644in}}%
\pgfpathquadraticcurveto{\pgfqpoint{5.799295in}{0.955444in}}{\pgfqpoint{5.804248in}{0.958200in}}%
\pgfpathquadraticcurveto{\pgfqpoint{5.809201in}{0.960956in}}{\pgfqpoint{5.816010in}{0.960956in}}%
\pgfpathquadraticcurveto{\pgfqpoint{5.822890in}{0.960956in}}{\pgfqpoint{5.827700in}{0.957461in}}%
\pgfpathquadraticcurveto{\pgfqpoint{5.832509in}{0.953985in}}{\pgfqpoint{5.834815in}{0.947338in}}%
\pgfpathlineto{\pgfqpoint{5.834815in}{0.947338in}}%
\pgfpathclose%
\pgfusepath{fill}%
\end{pgfscope}%
\begin{pgfscope}%
\pgfsetbuttcap%
\pgfsetroundjoin%
\definecolor{currentfill}{rgb}{0.309804,0.372549,0.419608}%
\pgfsetfillcolor{currentfill}%
\pgfsetlinewidth{0.000000pt}%
\definecolor{currentstroke}{rgb}{0.309804,0.372549,0.419608}%
\pgfsetstrokecolor{currentstroke}%
\pgfsetdash{}{0pt}%
\pgfpathmoveto{\pgfqpoint{5.974683in}{0.983993in}}%
\pgfpathlineto{\pgfqpoint{5.985040in}{0.983993in}}%
\pgfpathlineto{\pgfqpoint{5.985040in}{0.896400in}}%
\pgfpathlineto{\pgfqpoint{5.974683in}{0.896400in}}%
\pgfpathlineto{\pgfqpoint{5.974683in}{0.983993in}}%
\pgfpathlineto{\pgfqpoint{5.974683in}{0.983993in}}%
\pgfpathclose%
\pgfpathmoveto{\pgfqpoint{6.060637in}{0.930515in}}%
\pgfpathlineto{\pgfqpoint{6.060637in}{0.925454in}}%
\pgfpathlineto{\pgfqpoint{6.013013in}{0.925454in}}%
\pgfpathquadraticcurveto{\pgfqpoint{6.013697in}{0.914754in}}{\pgfqpoint{6.019461in}{0.909153in}}%
\pgfpathquadraticcurveto{\pgfqpoint{6.025225in}{0.903551in}}{\pgfqpoint{6.035528in}{0.903551in}}%
\pgfpathquadraticcurveto{\pgfqpoint{6.041490in}{0.903551in}}{\pgfqpoint{6.047092in}{0.905010in}}%
\pgfpathquadraticcurveto{\pgfqpoint{6.052694in}{0.906469in}}{\pgfqpoint{6.058223in}{0.909405in}}%
\pgfpathlineto{\pgfqpoint{6.058223in}{0.899606in}}%
\pgfpathquadraticcurveto{\pgfqpoint{6.052640in}{0.897247in}}{\pgfqpoint{6.046786in}{0.896004in}}%
\pgfpathquadraticcurveto{\pgfqpoint{6.040932in}{0.894761in}}{\pgfqpoint{6.034916in}{0.894761in}}%
\pgfpathquadraticcurveto{\pgfqpoint{6.019822in}{0.894761in}}{\pgfqpoint{6.011014in}{0.903533in}}%
\pgfpathquadraticcurveto{\pgfqpoint{6.002206in}{0.912323in}}{\pgfqpoint{6.002206in}{0.927309in}}%
\pgfpathquadraticcurveto{\pgfqpoint{6.002206in}{0.942781in}}{\pgfqpoint{6.010563in}{0.951859in}}%
\pgfpathquadraticcurveto{\pgfqpoint{6.018921in}{0.960956in}}{\pgfqpoint{6.033115in}{0.960956in}}%
\pgfpathquadraticcurveto{\pgfqpoint{6.045831in}{0.960956in}}{\pgfqpoint{6.053234in}{0.952760in}}%
\pgfpathquadraticcurveto{\pgfqpoint{6.060637in}{0.944583in}}{\pgfqpoint{6.060637in}{0.930515in}}%
\pgfpathlineto{\pgfqpoint{6.060637in}{0.930515in}}%
\pgfpathclose%
\pgfpathmoveto{\pgfqpoint{6.050280in}{0.933559in}}%
\pgfpathquadraticcurveto{\pgfqpoint{6.050172in}{0.942043in}}{\pgfqpoint{6.045525in}{0.947104in}}%
\pgfpathquadraticcurveto{\pgfqpoint{6.040878in}{0.952184in}}{\pgfqpoint{6.033223in}{0.952184in}}%
\pgfpathquadraticcurveto{\pgfqpoint{6.024559in}{0.952184in}}{\pgfqpoint{6.019353in}{0.947284in}}%
\pgfpathquadraticcurveto{\pgfqpoint{6.014148in}{0.942385in}}{\pgfqpoint{6.013355in}{0.933487in}}%
\pgfpathlineto{\pgfqpoint{6.050280in}{0.933559in}}%
\pgfpathlineto{\pgfqpoint{6.050280in}{0.933559in}}%
\pgfpathclose%
\pgfpathmoveto{\pgfqpoint{6.070220in}{0.959443in}}%
\pgfpathlineto{\pgfqpoint{6.081189in}{0.959443in}}%
\pgfpathlineto{\pgfqpoint{6.100894in}{0.906541in}}%
\pgfpathlineto{\pgfqpoint{6.120600in}{0.959443in}}%
\pgfpathlineto{\pgfqpoint{6.131569in}{0.959443in}}%
\pgfpathlineto{\pgfqpoint{6.107919in}{0.896400in}}%
\pgfpathlineto{\pgfqpoint{6.093851in}{0.896400in}}%
\pgfpathlineto{\pgfqpoint{6.070220in}{0.959443in}}%
\pgfpathlineto{\pgfqpoint{6.070220in}{0.959443in}}%
\pgfpathclose%
\pgfpathmoveto{\pgfqpoint{6.199799in}{0.930515in}}%
\pgfpathlineto{\pgfqpoint{6.199799in}{0.925454in}}%
\pgfpathlineto{\pgfqpoint{6.152175in}{0.925454in}}%
\pgfpathquadraticcurveto{\pgfqpoint{6.152859in}{0.914754in}}{\pgfqpoint{6.158623in}{0.909153in}}%
\pgfpathquadraticcurveto{\pgfqpoint{6.164387in}{0.903551in}}{\pgfqpoint{6.174690in}{0.903551in}}%
\pgfpathquadraticcurveto{\pgfqpoint{6.180652in}{0.903551in}}{\pgfqpoint{6.186254in}{0.905010in}}%
\pgfpathquadraticcurveto{\pgfqpoint{6.191856in}{0.906469in}}{\pgfqpoint{6.197385in}{0.909405in}}%
\pgfpathlineto{\pgfqpoint{6.197385in}{0.899606in}}%
\pgfpathquadraticcurveto{\pgfqpoint{6.191802in}{0.897247in}}{\pgfqpoint{6.185948in}{0.896004in}}%
\pgfpathquadraticcurveto{\pgfqpoint{6.180094in}{0.894761in}}{\pgfqpoint{6.174078in}{0.894761in}}%
\pgfpathquadraticcurveto{\pgfqpoint{6.158983in}{0.894761in}}{\pgfqpoint{6.150175in}{0.903533in}}%
\pgfpathquadraticcurveto{\pgfqpoint{6.141368in}{0.912323in}}{\pgfqpoint{6.141368in}{0.927309in}}%
\pgfpathquadraticcurveto{\pgfqpoint{6.141368in}{0.942781in}}{\pgfqpoint{6.149725in}{0.951859in}}%
\pgfpathquadraticcurveto{\pgfqpoint{6.158083in}{0.960956in}}{\pgfqpoint{6.172276in}{0.960956in}}%
\pgfpathquadraticcurveto{\pgfqpoint{6.184993in}{0.960956in}}{\pgfqpoint{6.192396in}{0.952760in}}%
\pgfpathquadraticcurveto{\pgfqpoint{6.199799in}{0.944583in}}{\pgfqpoint{6.199799in}{0.930515in}}%
\pgfpathlineto{\pgfqpoint{6.199799in}{0.930515in}}%
\pgfpathclose%
\pgfpathmoveto{\pgfqpoint{6.189442in}{0.933559in}}%
\pgfpathquadraticcurveto{\pgfqpoint{6.189334in}{0.942043in}}{\pgfqpoint{6.184687in}{0.947104in}}%
\pgfpathquadraticcurveto{\pgfqpoint{6.180040in}{0.952184in}}{\pgfqpoint{6.172384in}{0.952184in}}%
\pgfpathquadraticcurveto{\pgfqpoint{6.163721in}{0.952184in}}{\pgfqpoint{6.158515in}{0.947284in}}%
\pgfpathquadraticcurveto{\pgfqpoint{6.153310in}{0.942385in}}{\pgfqpoint{6.152517in}{0.933487in}}%
\pgfpathlineto{\pgfqpoint{6.189442in}{0.933559in}}%
\pgfpathlineto{\pgfqpoint{6.189442in}{0.933559in}}%
\pgfpathclose%
\pgfpathmoveto{\pgfqpoint{6.253331in}{0.949770in}}%
\pgfpathquadraticcurveto{\pgfqpoint{6.251584in}{0.950779in}}{\pgfqpoint{6.249531in}{0.951247in}}%
\pgfpathquadraticcurveto{\pgfqpoint{6.247477in}{0.951733in}}{\pgfqpoint{6.245009in}{0.951733in}}%
\pgfpathquadraticcurveto{\pgfqpoint{6.236220in}{0.951733in}}{\pgfqpoint{6.231518in}{0.946023in}}%
\pgfpathquadraticcurveto{\pgfqpoint{6.226817in}{0.940314in}}{\pgfqpoint{6.226817in}{0.929614in}}%
\pgfpathlineto{\pgfqpoint{6.226817in}{0.896400in}}%
\pgfpathlineto{\pgfqpoint{6.216406in}{0.896400in}}%
\pgfpathlineto{\pgfqpoint{6.216406in}{0.959443in}}%
\pgfpathlineto{\pgfqpoint{6.226817in}{0.959443in}}%
\pgfpathlineto{\pgfqpoint{6.226817in}{0.949644in}}%
\pgfpathquadraticcurveto{\pgfqpoint{6.230095in}{0.955390in}}{\pgfqpoint{6.235319in}{0.958164in}}%
\pgfpathquadraticcurveto{\pgfqpoint{6.240560in}{0.960956in}}{\pgfqpoint{6.248054in}{0.960956in}}%
\pgfpathquadraticcurveto{\pgfqpoint{6.249116in}{0.960956in}}{\pgfqpoint{6.250413in}{0.960811in}}%
\pgfpathquadraticcurveto{\pgfqpoint{6.251710in}{0.960685in}}{\pgfqpoint{6.253277in}{0.960397in}}%
\pgfpathlineto{\pgfqpoint{6.253331in}{0.949770in}}%
\pgfpathlineto{\pgfqpoint{6.253331in}{0.949770in}}%
\pgfpathclose%
\pgfpathmoveto{\pgfqpoint{6.292850in}{0.928083in}}%
\pgfpathquadraticcurveto{\pgfqpoint{6.280295in}{0.928083in}}{\pgfqpoint{6.275450in}{0.925219in}}%
\pgfpathquadraticcurveto{\pgfqpoint{6.270605in}{0.922356in}}{\pgfqpoint{6.270605in}{0.915421in}}%
\pgfpathquadraticcurveto{\pgfqpoint{6.270605in}{0.909909in}}{\pgfqpoint{6.274243in}{0.906667in}}%
\pgfpathquadraticcurveto{\pgfqpoint{6.277882in}{0.903443in}}{\pgfqpoint{6.284114in}{0.903443in}}%
\pgfpathquadraticcurveto{\pgfqpoint{6.292742in}{0.903443in}}{\pgfqpoint{6.297947in}{0.909549in}}%
\pgfpathquadraticcurveto{\pgfqpoint{6.303153in}{0.915655in}}{\pgfqpoint{6.303153in}{0.925778in}}%
\pgfpathlineto{\pgfqpoint{6.303153in}{0.928083in}}%
\pgfpathlineto{\pgfqpoint{6.292850in}{0.928083in}}%
\pgfpathlineto{\pgfqpoint{6.292850in}{0.928083in}}%
\pgfpathclose%
\pgfpathmoveto{\pgfqpoint{6.313510in}{0.932370in}}%
\pgfpathlineto{\pgfqpoint{6.313510in}{0.896400in}}%
\pgfpathlineto{\pgfqpoint{6.303153in}{0.896400in}}%
\pgfpathlineto{\pgfqpoint{6.303153in}{0.905964in}}%
\pgfpathquadraticcurveto{\pgfqpoint{6.299604in}{0.900237in}}{\pgfqpoint{6.294309in}{0.897499in}}%
\pgfpathquadraticcurveto{\pgfqpoint{6.289013in}{0.894761in}}{\pgfqpoint{6.281358in}{0.894761in}}%
\pgfpathquadraticcurveto{\pgfqpoint{6.271685in}{0.894761in}}{\pgfqpoint{6.265958in}{0.900201in}}%
\pgfpathquadraticcurveto{\pgfqpoint{6.260248in}{0.905640in}}{\pgfqpoint{6.260248in}{0.914754in}}%
\pgfpathquadraticcurveto{\pgfqpoint{6.260248in}{0.925382in}}{\pgfqpoint{6.267363in}{0.930785in}}%
\pgfpathquadraticcurveto{\pgfqpoint{6.274495in}{0.936189in}}{\pgfqpoint{6.288617in}{0.936189in}}%
\pgfpathlineto{\pgfqpoint{6.303153in}{0.936189in}}%
\pgfpathlineto{\pgfqpoint{6.303153in}{0.937216in}}%
\pgfpathquadraticcurveto{\pgfqpoint{6.303153in}{0.944366in}}{\pgfqpoint{6.298452in}{0.948275in}}%
\pgfpathquadraticcurveto{\pgfqpoint{6.293750in}{0.952184in}}{\pgfqpoint{6.285249in}{0.952184in}}%
\pgfpathquadraticcurveto{\pgfqpoint{6.279845in}{0.952184in}}{\pgfqpoint{6.274712in}{0.950887in}}%
\pgfpathquadraticcurveto{\pgfqpoint{6.269596in}{0.949590in}}{\pgfqpoint{6.264877in}{0.946996in}}%
\pgfpathlineto{\pgfqpoint{6.264877in}{0.956579in}}%
\pgfpathquadraticcurveto{\pgfqpoint{6.270551in}{0.958776in}}{\pgfqpoint{6.275900in}{0.959857in}}%
\pgfpathquadraticcurveto{\pgfqpoint{6.281250in}{0.960956in}}{\pgfqpoint{6.286311in}{0.960956in}}%
\pgfpathquadraticcurveto{\pgfqpoint{6.300001in}{0.960956in}}{\pgfqpoint{6.306755in}{0.953859in}}%
\pgfpathquadraticcurveto{\pgfqpoint{6.313510in}{0.946780in}}{\pgfqpoint{6.313510in}{0.932370in}}%
\pgfpathlineto{\pgfqpoint{6.313510in}{0.932370in}}%
\pgfpathclose%
\pgfpathmoveto{\pgfqpoint{6.376318in}{0.928660in}}%
\pgfpathquadraticcurveto{\pgfqpoint{6.376318in}{0.939917in}}{\pgfqpoint{6.371671in}{0.946096in}}%
\pgfpathquadraticcurveto{\pgfqpoint{6.367042in}{0.952292in}}{\pgfqpoint{6.358648in}{0.952292in}}%
\pgfpathquadraticcurveto{\pgfqpoint{6.350327in}{0.952292in}}{\pgfqpoint{6.345679in}{0.946096in}}%
\pgfpathquadraticcurveto{\pgfqpoint{6.341032in}{0.939917in}}{\pgfqpoint{6.341032in}{0.928660in}}%
\pgfpathquadraticcurveto{\pgfqpoint{6.341032in}{0.917456in}}{\pgfqpoint{6.345679in}{0.911260in}}%
\pgfpathquadraticcurveto{\pgfqpoint{6.350327in}{0.905064in}}{\pgfqpoint{6.358648in}{0.905064in}}%
\pgfpathquadraticcurveto{\pgfqpoint{6.367042in}{0.905064in}}{\pgfqpoint{6.371671in}{0.911260in}}%
\pgfpathquadraticcurveto{\pgfqpoint{6.376318in}{0.917456in}}{\pgfqpoint{6.376318in}{0.928660in}}%
\pgfpathlineto{\pgfqpoint{6.376318in}{0.928660in}}%
\pgfpathclose%
\pgfpathmoveto{\pgfqpoint{6.386675in}{0.904217in}}%
\pgfpathquadraticcurveto{\pgfqpoint{6.386675in}{0.888132in}}{\pgfqpoint{6.379524in}{0.880279in}}%
\pgfpathquadraticcurveto{\pgfqpoint{6.372391in}{0.872426in}}{\pgfqpoint{6.357639in}{0.872426in}}%
\pgfpathquadraticcurveto{\pgfqpoint{6.352182in}{0.872426in}}{\pgfqpoint{6.347337in}{0.873236in}}%
\pgfpathquadraticcurveto{\pgfqpoint{6.342491in}{0.874047in}}{\pgfqpoint{6.337934in}{0.875740in}}%
\pgfpathlineto{\pgfqpoint{6.337934in}{0.885809in}}%
\pgfpathquadraticcurveto{\pgfqpoint{6.342491in}{0.883341in}}{\pgfqpoint{6.346940in}{0.882170in}}%
\pgfpathquadraticcurveto{\pgfqpoint{6.351389in}{0.880982in}}{\pgfqpoint{6.356000in}{0.880982in}}%
\pgfpathquadraticcurveto{\pgfqpoint{6.366195in}{0.880982in}}{\pgfqpoint{6.371257in}{0.886295in}}%
\pgfpathquadraticcurveto{\pgfqpoint{6.376318in}{0.891609in}}{\pgfqpoint{6.376318in}{0.902362in}}%
\pgfpathlineto{\pgfqpoint{6.376318in}{0.907495in}}%
\pgfpathquadraticcurveto{\pgfqpoint{6.373112in}{0.901912in}}{\pgfqpoint{6.368105in}{0.899156in}}%
\pgfpathquadraticcurveto{\pgfqpoint{6.363097in}{0.896400in}}{\pgfqpoint{6.356108in}{0.896400in}}%
\pgfpathquadraticcurveto{\pgfqpoint{6.344527in}{0.896400in}}{\pgfqpoint{6.337430in}{0.905226in}}%
\pgfpathquadraticcurveto{\pgfqpoint{6.330333in}{0.914070in}}{\pgfqpoint{6.330333in}{0.928660in}}%
\pgfpathquadraticcurveto{\pgfqpoint{6.330333in}{0.943286in}}{\pgfqpoint{6.337430in}{0.952112in}}%
\pgfpathquadraticcurveto{\pgfqpoint{6.344527in}{0.960956in}}{\pgfqpoint{6.356108in}{0.960956in}}%
\pgfpathquadraticcurveto{\pgfqpoint{6.363097in}{0.960956in}}{\pgfqpoint{6.368105in}{0.958200in}}%
\pgfpathquadraticcurveto{\pgfqpoint{6.373112in}{0.955444in}}{\pgfqpoint{6.376318in}{0.949878in}}%
\pgfpathlineto{\pgfqpoint{6.376318in}{0.959443in}}%
\pgfpathlineto{\pgfqpoint{6.386675in}{0.959443in}}%
\pgfpathlineto{\pgfqpoint{6.386675in}{0.904217in}}%
\pgfpathlineto{\pgfqpoint{6.386675in}{0.904217in}}%
\pgfpathclose%
\pgfpathmoveto{\pgfqpoint{6.461948in}{0.930515in}}%
\pgfpathlineto{\pgfqpoint{6.461948in}{0.925454in}}%
\pgfpathlineto{\pgfqpoint{6.414324in}{0.925454in}}%
\pgfpathquadraticcurveto{\pgfqpoint{6.415008in}{0.914754in}}{\pgfqpoint{6.420772in}{0.909153in}}%
\pgfpathquadraticcurveto{\pgfqpoint{6.426536in}{0.903551in}}{\pgfqpoint{6.436839in}{0.903551in}}%
\pgfpathquadraticcurveto{\pgfqpoint{6.442801in}{0.903551in}}{\pgfqpoint{6.448403in}{0.905010in}}%
\pgfpathquadraticcurveto{\pgfqpoint{6.454004in}{0.906469in}}{\pgfqpoint{6.459534in}{0.909405in}}%
\pgfpathlineto{\pgfqpoint{6.459534in}{0.899606in}}%
\pgfpathquadraticcurveto{\pgfqpoint{6.453950in}{0.897247in}}{\pgfqpoint{6.448096in}{0.896004in}}%
\pgfpathquadraticcurveto{\pgfqpoint{6.442243in}{0.894761in}}{\pgfqpoint{6.436226in}{0.894761in}}%
\pgfpathquadraticcurveto{\pgfqpoint{6.421132in}{0.894761in}}{\pgfqpoint{6.412324in}{0.903533in}}%
\pgfpathquadraticcurveto{\pgfqpoint{6.403516in}{0.912323in}}{\pgfqpoint{6.403516in}{0.927309in}}%
\pgfpathquadraticcurveto{\pgfqpoint{6.403516in}{0.942781in}}{\pgfqpoint{6.411874in}{0.951859in}}%
\pgfpathquadraticcurveto{\pgfqpoint{6.420232in}{0.960956in}}{\pgfqpoint{6.434425in}{0.960956in}}%
\pgfpathquadraticcurveto{\pgfqpoint{6.447142in}{0.960956in}}{\pgfqpoint{6.454545in}{0.952760in}}%
\pgfpathquadraticcurveto{\pgfqpoint{6.461948in}{0.944583in}}{\pgfqpoint{6.461948in}{0.930515in}}%
\pgfpathlineto{\pgfqpoint{6.461948in}{0.930515in}}%
\pgfpathclose%
\pgfpathmoveto{\pgfqpoint{6.451591in}{0.933559in}}%
\pgfpathquadraticcurveto{\pgfqpoint{6.451483in}{0.942043in}}{\pgfqpoint{6.446836in}{0.947104in}}%
\pgfpathquadraticcurveto{\pgfqpoint{6.442189in}{0.952184in}}{\pgfqpoint{6.434533in}{0.952184in}}%
\pgfpathquadraticcurveto{\pgfqpoint{6.425870in}{0.952184in}}{\pgfqpoint{6.420664in}{0.947284in}}%
\pgfpathquadraticcurveto{\pgfqpoint{6.415458in}{0.942385in}}{\pgfqpoint{6.414666in}{0.933487in}}%
\pgfpathlineto{\pgfqpoint{6.451591in}{0.933559in}}%
\pgfpathlineto{\pgfqpoint{6.451591in}{0.933559in}}%
\pgfpathclose%
\pgfusepath{fill}%
\end{pgfscope}%
\begin{pgfscope}%
\pgfsetbuttcap%
\pgfsetroundjoin%
\definecolor{currentfill}{rgb}{0.309804,0.372549,0.419608}%
\pgfsetfillcolor{currentfill}%
\pgfsetlinewidth{0.000000pt}%
\definecolor{currentstroke}{rgb}{0.309804,0.372549,0.419608}%
\pgfsetstrokecolor{currentstroke}%
\pgfsetdash{}{0pt}%
\pgfpathmoveto{\pgfqpoint{6.533518in}{0.953138in}}%
\pgfpathlineto{\pgfqpoint{6.533518in}{0.963387in}}%
\pgfpathlineto{\pgfqpoint{6.605674in}{0.937216in}}%
\pgfpathlineto{\pgfqpoint{6.605674in}{0.927867in}}%
\pgfpathlineto{\pgfqpoint{6.533518in}{0.901696in}}%
\pgfpathlineto{\pgfqpoint{6.533518in}{0.911944in}}%
\pgfpathlineto{\pgfqpoint{6.591499in}{0.932478in}}%
\pgfpathlineto{\pgfqpoint{6.533518in}{0.953138in}}%
\pgfpathlineto{\pgfqpoint{6.533518in}{0.953138in}}%
\pgfpathclose%
\pgfpathmoveto{\pgfqpoint{6.640024in}{0.905964in}}%
\pgfpathlineto{\pgfqpoint{6.679704in}{0.905964in}}%
\pgfpathlineto{\pgfqpoint{6.679704in}{0.896400in}}%
\pgfpathlineto{\pgfqpoint{6.626352in}{0.896400in}}%
\pgfpathlineto{\pgfqpoint{6.626352in}{0.905964in}}%
\pgfpathquadraticcurveto{\pgfqpoint{6.632819in}{0.912665in}}{\pgfqpoint{6.643986in}{0.923941in}}%
\pgfpathquadraticcurveto{\pgfqpoint{6.655172in}{0.935234in}}{\pgfqpoint{6.658036in}{0.938512in}}%
\pgfpathquadraticcurveto{\pgfqpoint{6.663493in}{0.944637in}}{\pgfqpoint{6.665655in}{0.948887in}}%
\pgfpathquadraticcurveto{\pgfqpoint{6.667834in}{0.953138in}}{\pgfqpoint{6.667834in}{0.957245in}}%
\pgfpathquadraticcurveto{\pgfqpoint{6.667834in}{0.963946in}}{\pgfqpoint{6.663133in}{0.968160in}}%
\pgfpathquadraticcurveto{\pgfqpoint{6.658432in}{0.972393in}}{\pgfqpoint{6.650885in}{0.972393in}}%
\pgfpathquadraticcurveto{\pgfqpoint{6.645535in}{0.972393in}}{\pgfqpoint{6.639591in}{0.970538in}}%
\pgfpathquadraticcurveto{\pgfqpoint{6.633665in}{0.968683in}}{\pgfqpoint{6.626911in}{0.964900in}}%
\pgfpathlineto{\pgfqpoint{6.626911in}{0.976392in}}%
\pgfpathquadraticcurveto{\pgfqpoint{6.633773in}{0.979148in}}{\pgfqpoint{6.639735in}{0.980553in}}%
\pgfpathquadraticcurveto{\pgfqpoint{6.645715in}{0.981958in}}{\pgfqpoint{6.650669in}{0.981958in}}%
\pgfpathquadraticcurveto{\pgfqpoint{6.663728in}{0.981958in}}{\pgfqpoint{6.671491in}{0.975419in}}%
\pgfpathquadraticcurveto{\pgfqpoint{6.679254in}{0.968899in}}{\pgfqpoint{6.679254in}{0.957984in}}%
\pgfpathquadraticcurveto{\pgfqpoint{6.679254in}{0.952796in}}{\pgfqpoint{6.677309in}{0.948149in}}%
\pgfpathquadraticcurveto{\pgfqpoint{6.675382in}{0.943520in}}{\pgfqpoint{6.670248in}{0.937216in}}%
\pgfpathquadraticcurveto{\pgfqpoint{6.668843in}{0.935576in}}{\pgfqpoint{6.661296in}{0.927777in}}%
\pgfpathquadraticcurveto{\pgfqpoint{6.653767in}{0.919978in}}{\pgfqpoint{6.640024in}{0.905964in}}%
\pgfpathlineto{\pgfqpoint{6.640024in}{0.905964in}}%
\pgfpathclose%
\pgfpathmoveto{\pgfqpoint{6.703697in}{0.980445in}}%
\pgfpathlineto{\pgfqpoint{6.748331in}{0.980445in}}%
\pgfpathlineto{\pgfqpoint{6.748331in}{0.970862in}}%
\pgfpathlineto{\pgfqpoint{6.714108in}{0.970862in}}%
\pgfpathlineto{\pgfqpoint{6.714108in}{0.950274in}}%
\pgfpathquadraticcurveto{\pgfqpoint{6.716575in}{0.951121in}}{\pgfqpoint{6.719043in}{0.951535in}}%
\pgfpathquadraticcurveto{\pgfqpoint{6.721529in}{0.951949in}}{\pgfqpoint{6.724014in}{0.951949in}}%
\pgfpathquadraticcurveto{\pgfqpoint{6.738082in}{0.951949in}}{\pgfqpoint{6.746295in}{0.944240in}}%
\pgfpathquadraticcurveto{\pgfqpoint{6.754527in}{0.936531in}}{\pgfqpoint{6.754527in}{0.923364in}}%
\pgfpathquadraticcurveto{\pgfqpoint{6.754527in}{0.909801in}}{\pgfqpoint{6.746079in}{0.902272in}}%
\pgfpathquadraticcurveto{\pgfqpoint{6.737632in}{0.894761in}}{\pgfqpoint{6.722267in}{0.894761in}}%
\pgfpathquadraticcurveto{\pgfqpoint{6.716972in}{0.894761in}}{\pgfqpoint{6.711478in}{0.895662in}}%
\pgfpathquadraticcurveto{\pgfqpoint{6.706002in}{0.896562in}}{\pgfqpoint{6.700148in}{0.898363in}}%
\pgfpathlineto{\pgfqpoint{6.700148in}{0.909801in}}%
\pgfpathquadraticcurveto{\pgfqpoint{6.705210in}{0.907045in}}{\pgfqpoint{6.710613in}{0.905694in}}%
\pgfpathquadraticcurveto{\pgfqpoint{6.716017in}{0.904343in}}{\pgfqpoint{6.722033in}{0.904343in}}%
\pgfpathquadraticcurveto{\pgfqpoint{6.731778in}{0.904343in}}{\pgfqpoint{6.737451in}{0.909459in}}%
\pgfpathquadraticcurveto{\pgfqpoint{6.743143in}{0.914574in}}{\pgfqpoint{6.743143in}{0.923364in}}%
\pgfpathquadraticcurveto{\pgfqpoint{6.743143in}{0.932136in}}{\pgfqpoint{6.737451in}{0.937252in}}%
\pgfpathquadraticcurveto{\pgfqpoint{6.731778in}{0.942385in}}{\pgfqpoint{6.722033in}{0.942385in}}%
\pgfpathquadraticcurveto{\pgfqpoint{6.717476in}{0.942385in}}{\pgfqpoint{6.712937in}{0.941376in}}%
\pgfpathquadraticcurveto{\pgfqpoint{6.708416in}{0.940368in}}{\pgfqpoint{6.703697in}{0.938224in}}%
\pgfpathlineto{\pgfqpoint{6.703697in}{0.980445in}}%
\pgfpathlineto{\pgfqpoint{6.703697in}{0.980445in}}%
\pgfpathclose%
\pgfpathmoveto{\pgfqpoint{6.848406in}{0.933379in}}%
\pgfpathquadraticcurveto{\pgfqpoint{6.843507in}{0.933379in}}{\pgfqpoint{6.840715in}{0.929218in}}%
\pgfpathquadraticcurveto{\pgfqpoint{6.837941in}{0.925057in}}{\pgfqpoint{6.837941in}{0.917618in}}%
\pgfpathquadraticcurveto{\pgfqpoint{6.837941in}{0.910305in}}{\pgfqpoint{6.840715in}{0.906109in}}%
\pgfpathquadraticcurveto{\pgfqpoint{6.843507in}{0.901912in}}{\pgfqpoint{6.848406in}{0.901912in}}%
\pgfpathquadraticcurveto{\pgfqpoint{6.853197in}{0.901912in}}{\pgfqpoint{6.855971in}{0.906109in}}%
\pgfpathquadraticcurveto{\pgfqpoint{6.858763in}{0.910305in}}{\pgfqpoint{6.858763in}{0.917618in}}%
\pgfpathquadraticcurveto{\pgfqpoint{6.858763in}{0.925003in}}{\pgfqpoint{6.855971in}{0.929182in}}%
\pgfpathquadraticcurveto{\pgfqpoint{6.853197in}{0.933379in}}{\pgfqpoint{6.848406in}{0.933379in}}%
\pgfpathlineto{\pgfqpoint{6.848406in}{0.933379in}}%
\pgfpathclose%
\pgfpathmoveto{\pgfqpoint{6.848406in}{0.940530in}}%
\pgfpathquadraticcurveto{\pgfqpoint{6.857304in}{0.940530in}}{\pgfqpoint{6.862528in}{0.934334in}}%
\pgfpathquadraticcurveto{\pgfqpoint{6.867769in}{0.928155in}}{\pgfqpoint{6.867769in}{0.917618in}}%
\pgfpathquadraticcurveto{\pgfqpoint{6.867769in}{0.907099in}}{\pgfqpoint{6.862510in}{0.900921in}}%
\pgfpathquadraticcurveto{\pgfqpoint{6.857250in}{0.894761in}}{\pgfqpoint{6.848406in}{0.894761in}}%
\pgfpathquadraticcurveto{\pgfqpoint{6.839400in}{0.894761in}}{\pgfqpoint{6.834159in}{0.900921in}}%
\pgfpathquadraticcurveto{\pgfqpoint{6.828935in}{0.907099in}}{\pgfqpoint{6.828935in}{0.917618in}}%
\pgfpathquadraticcurveto{\pgfqpoint{6.828935in}{0.928209in}}{\pgfqpoint{6.834195in}{0.934370in}}%
\pgfpathquadraticcurveto{\pgfqpoint{6.839454in}{0.940530in}}{\pgfqpoint{6.848406in}{0.940530in}}%
\pgfpathlineto{\pgfqpoint{6.848406in}{0.940530in}}%
\pgfpathclose%
\pgfpathmoveto{\pgfqpoint{6.790317in}{0.974807in}}%
\pgfpathquadraticcurveto{\pgfqpoint{6.785472in}{0.974807in}}{\pgfqpoint{6.782680in}{0.970610in}}%
\pgfpathquadraticcurveto{\pgfqpoint{6.779906in}{0.966431in}}{\pgfqpoint{6.779906in}{0.959100in}}%
\pgfpathquadraticcurveto{\pgfqpoint{6.779906in}{0.951679in}}{\pgfqpoint{6.782662in}{0.947500in}}%
\pgfpathquadraticcurveto{\pgfqpoint{6.785418in}{0.943340in}}{\pgfqpoint{6.790317in}{0.943340in}}%
\pgfpathquadraticcurveto{\pgfqpoint{6.795216in}{0.943340in}}{\pgfqpoint{6.797990in}{0.947500in}}%
\pgfpathquadraticcurveto{\pgfqpoint{6.800782in}{0.951679in}}{\pgfqpoint{6.800782in}{0.959100in}}%
\pgfpathquadraticcurveto{\pgfqpoint{6.800782in}{0.966359in}}{\pgfqpoint{6.797972in}{0.970574in}}%
\pgfpathquadraticcurveto{\pgfqpoint{6.795162in}{0.974807in}}{\pgfqpoint{6.790317in}{0.974807in}}%
\pgfpathlineto{\pgfqpoint{6.790317in}{0.974807in}}%
\pgfpathclose%
\pgfpathmoveto{\pgfqpoint{6.841147in}{0.981958in}}%
\pgfpathlineto{\pgfqpoint{6.850153in}{0.981958in}}%
\pgfpathlineto{\pgfqpoint{6.797576in}{0.894761in}}%
\pgfpathlineto{\pgfqpoint{6.788570in}{0.894761in}}%
\pgfpathlineto{\pgfqpoint{6.841147in}{0.981958in}}%
\pgfpathlineto{\pgfqpoint{6.841147in}{0.981958in}}%
\pgfpathclose%
\pgfpathmoveto{\pgfqpoint{6.790317in}{0.981958in}}%
\pgfpathquadraticcurveto{\pgfqpoint{6.799215in}{0.981958in}}{\pgfqpoint{6.804493in}{0.975798in}}%
\pgfpathquadraticcurveto{\pgfqpoint{6.809788in}{0.969637in}}{\pgfqpoint{6.809788in}{0.959100in}}%
\pgfpathquadraticcurveto{\pgfqpoint{6.809788in}{0.948473in}}{\pgfqpoint{6.804529in}{0.942331in}}%
\pgfpathquadraticcurveto{\pgfqpoint{6.799269in}{0.936189in}}{\pgfqpoint{6.790317in}{0.936189in}}%
\pgfpathquadraticcurveto{\pgfqpoint{6.781365in}{0.936189in}}{\pgfqpoint{6.776160in}{0.942349in}}%
\pgfpathquadraticcurveto{\pgfqpoint{6.770954in}{0.948527in}}{\pgfqpoint{6.770954in}{0.959100in}}%
\pgfpathquadraticcurveto{\pgfqpoint{6.770954in}{0.969583in}}{\pgfqpoint{6.776178in}{0.975762in}}%
\pgfpathquadraticcurveto{\pgfqpoint{6.781419in}{0.981958in}}{\pgfqpoint{6.790317in}{0.981958in}}%
\pgfpathlineto{\pgfqpoint{6.790317in}{0.981958in}}%
\pgfpathclose%
\pgfusepath{fill}%
\end{pgfscope}%
\begin{pgfscope}%
\pgfsetbuttcap%
\pgfsetroundjoin%
\definecolor{currentfill}{rgb}{0.309804,0.372549,0.419608}%
\pgfsetfillcolor{currentfill}%
\pgfsetlinewidth{0.000000pt}%
\definecolor{currentstroke}{rgb}{0.309804,0.372549,0.419608}%
\pgfsetstrokecolor{currentstroke}%
\pgfsetdash{}{0pt}%
\pgfpathmoveto{\pgfqpoint{6.981568in}{0.983993in}}%
\pgfpathlineto{\pgfqpoint{6.981568in}{0.975365in}}%
\pgfpathlineto{\pgfqpoint{6.971662in}{0.975365in}}%
\pgfpathquadraticcurveto{\pgfqpoint{6.966096in}{0.975365in}}{\pgfqpoint{6.963916in}{0.973114in}}%
\pgfpathquadraticcurveto{\pgfqpoint{6.961755in}{0.970862in}}{\pgfqpoint{6.961755in}{0.965008in}}%
\pgfpathlineto{\pgfqpoint{6.961755in}{0.959443in}}%
\pgfpathlineto{\pgfqpoint{6.978812in}{0.959443in}}%
\pgfpathlineto{\pgfqpoint{6.978812in}{0.951391in}}%
\pgfpathlineto{\pgfqpoint{6.961755in}{0.951391in}}%
\pgfpathlineto{\pgfqpoint{6.961755in}{0.896400in}}%
\pgfpathlineto{\pgfqpoint{6.951344in}{0.896400in}}%
\pgfpathlineto{\pgfqpoint{6.951344in}{0.951391in}}%
\pgfpathlineto{\pgfqpoint{6.941437in}{0.951391in}}%
\pgfpathlineto{\pgfqpoint{6.941437in}{0.959443in}}%
\pgfpathlineto{\pgfqpoint{6.951344in}{0.959443in}}%
\pgfpathlineto{\pgfqpoint{6.951344in}{0.963838in}}%
\pgfpathquadraticcurveto{\pgfqpoint{6.951344in}{0.974357in}}{\pgfqpoint{6.956243in}{0.979166in}}%
\pgfpathquadraticcurveto{\pgfqpoint{6.961142in}{0.983993in}}{\pgfqpoint{6.971770in}{0.983993in}}%
\pgfpathlineto{\pgfqpoint{6.981568in}{0.983993in}}%
\pgfpathlineto{\pgfqpoint{6.981568in}{0.983993in}}%
\pgfpathclose%
\pgfpathmoveto{\pgfqpoint{7.014657in}{0.952184in}}%
\pgfpathquadraticcurveto{\pgfqpoint{7.006335in}{0.952184in}}{\pgfqpoint{7.001490in}{0.945681in}}%
\pgfpathquadraticcurveto{\pgfqpoint{6.996644in}{0.939179in}}{\pgfqpoint{6.996644in}{0.927867in}}%
\pgfpathquadraticcurveto{\pgfqpoint{6.996644in}{0.916556in}}{\pgfqpoint{7.001454in}{0.910053in}}%
\pgfpathquadraticcurveto{\pgfqpoint{7.006281in}{0.903551in}}{\pgfqpoint{7.014657in}{0.903551in}}%
\pgfpathquadraticcurveto{\pgfqpoint{7.022942in}{0.903551in}}{\pgfqpoint{7.027769in}{0.910071in}}%
\pgfpathquadraticcurveto{\pgfqpoint{7.032615in}{0.916610in}}{\pgfqpoint{7.032615in}{0.927867in}}%
\pgfpathquadraticcurveto{\pgfqpoint{7.032615in}{0.939071in}}{\pgfqpoint{7.027769in}{0.945627in}}%
\pgfpathquadraticcurveto{\pgfqpoint{7.022942in}{0.952184in}}{\pgfqpoint{7.014657in}{0.952184in}}%
\pgfpathlineto{\pgfqpoint{7.014657in}{0.952184in}}%
\pgfpathclose%
\pgfpathmoveto{\pgfqpoint{7.014657in}{0.960956in}}%
\pgfpathquadraticcurveto{\pgfqpoint{7.028166in}{0.960956in}}{\pgfqpoint{7.035875in}{0.952166in}}%
\pgfpathquadraticcurveto{\pgfqpoint{7.043602in}{0.943394in}}{\pgfqpoint{7.043602in}{0.927867in}}%
\pgfpathquadraticcurveto{\pgfqpoint{7.043602in}{0.912395in}}{\pgfqpoint{7.035875in}{0.903569in}}%
\pgfpathquadraticcurveto{\pgfqpoint{7.028166in}{0.894761in}}{\pgfqpoint{7.014657in}{0.894761in}}%
\pgfpathquadraticcurveto{\pgfqpoint{7.001093in}{0.894761in}}{\pgfqpoint{6.993402in}{0.903569in}}%
\pgfpathquadraticcurveto{\pgfqpoint{6.985729in}{0.912395in}}{\pgfqpoint{6.985729in}{0.927867in}}%
\pgfpathquadraticcurveto{\pgfqpoint{6.985729in}{0.943394in}}{\pgfqpoint{6.993402in}{0.952166in}}%
\pgfpathquadraticcurveto{\pgfqpoint{7.001093in}{0.960956in}}{\pgfqpoint{7.014657in}{0.960956in}}%
\pgfpathlineto{\pgfqpoint{7.014657in}{0.960956in}}%
\pgfpathclose%
\pgfpathmoveto{\pgfqpoint{7.097296in}{0.949770in}}%
\pgfpathquadraticcurveto{\pgfqpoint{7.095549in}{0.950779in}}{\pgfqpoint{7.093496in}{0.951247in}}%
\pgfpathquadraticcurveto{\pgfqpoint{7.091442in}{0.951733in}}{\pgfqpoint{7.088975in}{0.951733in}}%
\pgfpathquadraticcurveto{\pgfqpoint{7.080185in}{0.951733in}}{\pgfqpoint{7.075484in}{0.946023in}}%
\pgfpathquadraticcurveto{\pgfqpoint{7.070782in}{0.940314in}}{\pgfqpoint{7.070782in}{0.929614in}}%
\pgfpathlineto{\pgfqpoint{7.070782in}{0.896400in}}%
\pgfpathlineto{\pgfqpoint{7.060371in}{0.896400in}}%
\pgfpathlineto{\pgfqpoint{7.060371in}{0.959443in}}%
\pgfpathlineto{\pgfqpoint{7.070782in}{0.959443in}}%
\pgfpathlineto{\pgfqpoint{7.070782in}{0.949644in}}%
\pgfpathquadraticcurveto{\pgfqpoint{7.074061in}{0.955390in}}{\pgfqpoint{7.079284in}{0.958164in}}%
\pgfpathquadraticcurveto{\pgfqpoint{7.084526in}{0.960956in}}{\pgfqpoint{7.092019in}{0.960956in}}%
\pgfpathquadraticcurveto{\pgfqpoint{7.093081in}{0.960956in}}{\pgfqpoint{7.094378in}{0.960811in}}%
\pgfpathquadraticcurveto{\pgfqpoint{7.095675in}{0.960685in}}{\pgfqpoint{7.097242in}{0.960397in}}%
\pgfpathlineto{\pgfqpoint{7.097296in}{0.949770in}}%
\pgfpathlineto{\pgfqpoint{7.097296in}{0.949770in}}%
\pgfpathclose%
\pgfusepath{fill}%
\end{pgfscope}%
\begin{pgfscope}%
\pgfsetbuttcap%
\pgfsetroundjoin%
\definecolor{currentfill}{rgb}{0.309804,0.372549,0.419608}%
\pgfsetfillcolor{currentfill}%
\pgfsetlinewidth{0.000000pt}%
\definecolor{currentstroke}{rgb}{0.309804,0.372549,0.419608}%
\pgfsetstrokecolor{currentstroke}%
\pgfsetdash{}{0pt}%
\pgfpathmoveto{\pgfqpoint{7.211063in}{0.930515in}}%
\pgfpathlineto{\pgfqpoint{7.211063in}{0.925454in}}%
\pgfpathlineto{\pgfqpoint{7.163439in}{0.925454in}}%
\pgfpathquadraticcurveto{\pgfqpoint{7.164123in}{0.914754in}}{\pgfqpoint{7.169887in}{0.909153in}}%
\pgfpathquadraticcurveto{\pgfqpoint{7.175651in}{0.903551in}}{\pgfqpoint{7.185954in}{0.903551in}}%
\pgfpathquadraticcurveto{\pgfqpoint{7.191916in}{0.903551in}}{\pgfqpoint{7.197518in}{0.905010in}}%
\pgfpathquadraticcurveto{\pgfqpoint{7.203120in}{0.906469in}}{\pgfqpoint{7.208649in}{0.909405in}}%
\pgfpathlineto{\pgfqpoint{7.208649in}{0.899606in}}%
\pgfpathquadraticcurveto{\pgfqpoint{7.203065in}{0.897247in}}{\pgfqpoint{7.197212in}{0.896004in}}%
\pgfpathquadraticcurveto{\pgfqpoint{7.191358in}{0.894761in}}{\pgfqpoint{7.185342in}{0.894761in}}%
\pgfpathquadraticcurveto{\pgfqpoint{7.170247in}{0.894761in}}{\pgfqpoint{7.161439in}{0.903533in}}%
\pgfpathquadraticcurveto{\pgfqpoint{7.152631in}{0.912323in}}{\pgfqpoint{7.152631in}{0.927309in}}%
\pgfpathquadraticcurveto{\pgfqpoint{7.152631in}{0.942781in}}{\pgfqpoint{7.160989in}{0.951859in}}%
\pgfpathquadraticcurveto{\pgfqpoint{7.169347in}{0.960956in}}{\pgfqpoint{7.183540in}{0.960956in}}%
\pgfpathquadraticcurveto{\pgfqpoint{7.196257in}{0.960956in}}{\pgfqpoint{7.203660in}{0.952760in}}%
\pgfpathquadraticcurveto{\pgfqpoint{7.211063in}{0.944583in}}{\pgfqpoint{7.211063in}{0.930515in}}%
\pgfpathlineto{\pgfqpoint{7.211063in}{0.930515in}}%
\pgfpathclose%
\pgfpathmoveto{\pgfqpoint{7.200706in}{0.933559in}}%
\pgfpathquadraticcurveto{\pgfqpoint{7.200598in}{0.942043in}}{\pgfqpoint{7.195951in}{0.947104in}}%
\pgfpathquadraticcurveto{\pgfqpoint{7.191304in}{0.952184in}}{\pgfqpoint{7.183648in}{0.952184in}}%
\pgfpathquadraticcurveto{\pgfqpoint{7.174985in}{0.952184in}}{\pgfqpoint{7.169779in}{0.947284in}}%
\pgfpathquadraticcurveto{\pgfqpoint{7.164573in}{0.942385in}}{\pgfqpoint{7.163781in}{0.933487in}}%
\pgfpathlineto{\pgfqpoint{7.200706in}{0.933559in}}%
\pgfpathlineto{\pgfqpoint{7.200706in}{0.933559in}}%
\pgfpathclose%
\pgfpathmoveto{\pgfqpoint{7.228066in}{0.959443in}}%
\pgfpathlineto{\pgfqpoint{7.238423in}{0.959443in}}%
\pgfpathlineto{\pgfqpoint{7.238423in}{0.896400in}}%
\pgfpathlineto{\pgfqpoint{7.228066in}{0.896400in}}%
\pgfpathlineto{\pgfqpoint{7.228066in}{0.959443in}}%
\pgfpathlineto{\pgfqpoint{7.228066in}{0.959443in}}%
\pgfpathclose%
\pgfpathmoveto{\pgfqpoint{7.228066in}{0.983993in}}%
\pgfpathlineto{\pgfqpoint{7.238423in}{0.983993in}}%
\pgfpathlineto{\pgfqpoint{7.238423in}{0.970862in}}%
\pgfpathlineto{\pgfqpoint{7.228066in}{0.970862in}}%
\pgfpathlineto{\pgfqpoint{7.228066in}{0.983993in}}%
\pgfpathlineto{\pgfqpoint{7.228066in}{0.983993in}}%
\pgfpathclose%
\pgfpathmoveto{\pgfqpoint{7.270341in}{0.977347in}}%
\pgfpathlineto{\pgfqpoint{7.270341in}{0.959443in}}%
\pgfpathlineto{\pgfqpoint{7.291667in}{0.959443in}}%
\pgfpathlineto{\pgfqpoint{7.291667in}{0.951391in}}%
\pgfpathlineto{\pgfqpoint{7.270341in}{0.951391in}}%
\pgfpathlineto{\pgfqpoint{7.270341in}{0.917168in}}%
\pgfpathquadraticcurveto{\pgfqpoint{7.270341in}{0.909459in}}{\pgfqpoint{7.272448in}{0.907261in}}%
\pgfpathquadraticcurveto{\pgfqpoint{7.274556in}{0.905064in}}{\pgfqpoint{7.281040in}{0.905064in}}%
\pgfpathlineto{\pgfqpoint{7.291667in}{0.905064in}}%
\pgfpathlineto{\pgfqpoint{7.291667in}{0.896400in}}%
\pgfpathlineto{\pgfqpoint{7.281040in}{0.896400in}}%
\pgfpathquadraticcurveto{\pgfqpoint{7.269044in}{0.896400in}}{\pgfqpoint{7.264487in}{0.900867in}}%
\pgfpathquadraticcurveto{\pgfqpoint{7.259930in}{0.905352in}}{\pgfqpoint{7.259930in}{0.917168in}}%
\pgfpathlineto{\pgfqpoint{7.259930in}{0.951391in}}%
\pgfpathlineto{\pgfqpoint{7.252329in}{0.951391in}}%
\pgfpathlineto{\pgfqpoint{7.252329in}{0.959443in}}%
\pgfpathlineto{\pgfqpoint{7.259930in}{0.959443in}}%
\pgfpathlineto{\pgfqpoint{7.259930in}{0.977347in}}%
\pgfpathlineto{\pgfqpoint{7.270341in}{0.977347in}}%
\pgfpathlineto{\pgfqpoint{7.270341in}{0.977347in}}%
\pgfpathclose%
\pgfpathmoveto{\pgfqpoint{7.357700in}{0.934460in}}%
\pgfpathlineto{\pgfqpoint{7.357700in}{0.896400in}}%
\pgfpathlineto{\pgfqpoint{7.347343in}{0.896400in}}%
\pgfpathlineto{\pgfqpoint{7.347343in}{0.934117in}}%
\pgfpathquadraticcurveto{\pgfqpoint{7.347343in}{0.943069in}}{\pgfqpoint{7.343848in}{0.947500in}}%
\pgfpathquadraticcurveto{\pgfqpoint{7.340354in}{0.951949in}}{\pgfqpoint{7.333383in}{0.951949in}}%
\pgfpathquadraticcurveto{\pgfqpoint{7.324990in}{0.951949in}}{\pgfqpoint{7.320144in}{0.946600in}}%
\pgfpathquadraticcurveto{\pgfqpoint{7.315299in}{0.941268in}}{\pgfqpoint{7.315299in}{0.932028in}}%
\pgfpathlineto{\pgfqpoint{7.315299in}{0.896400in}}%
\pgfpathlineto{\pgfqpoint{7.304888in}{0.896400in}}%
\pgfpathlineto{\pgfqpoint{7.304888in}{0.983993in}}%
\pgfpathlineto{\pgfqpoint{7.315299in}{0.983993in}}%
\pgfpathlineto{\pgfqpoint{7.315299in}{0.949644in}}%
\pgfpathquadraticcurveto{\pgfqpoint{7.319028in}{0.955336in}}{\pgfqpoint{7.324053in}{0.958146in}}%
\pgfpathquadraticcurveto{\pgfqpoint{7.329096in}{0.960956in}}{\pgfqpoint{7.335689in}{0.960956in}}%
\pgfpathquadraticcurveto{\pgfqpoint{7.346550in}{0.960956in}}{\pgfqpoint{7.352116in}{0.954237in}}%
\pgfpathquadraticcurveto{\pgfqpoint{7.357700in}{0.947518in}}{\pgfqpoint{7.357700in}{0.934460in}}%
\pgfpathlineto{\pgfqpoint{7.357700in}{0.934460in}}%
\pgfpathclose%
\pgfpathmoveto{\pgfqpoint{7.432270in}{0.930515in}}%
\pgfpathlineto{\pgfqpoint{7.432270in}{0.925454in}}%
\pgfpathlineto{\pgfqpoint{7.384646in}{0.925454in}}%
\pgfpathquadraticcurveto{\pgfqpoint{7.385330in}{0.914754in}}{\pgfqpoint{7.391094in}{0.909153in}}%
\pgfpathquadraticcurveto{\pgfqpoint{7.396858in}{0.903551in}}{\pgfqpoint{7.407161in}{0.903551in}}%
\pgfpathquadraticcurveto{\pgfqpoint{7.413123in}{0.903551in}}{\pgfqpoint{7.418725in}{0.905010in}}%
\pgfpathquadraticcurveto{\pgfqpoint{7.424327in}{0.906469in}}{\pgfqpoint{7.429856in}{0.909405in}}%
\pgfpathlineto{\pgfqpoint{7.429856in}{0.899606in}}%
\pgfpathquadraticcurveto{\pgfqpoint{7.424273in}{0.897247in}}{\pgfqpoint{7.418419in}{0.896004in}}%
\pgfpathquadraticcurveto{\pgfqpoint{7.412565in}{0.894761in}}{\pgfqpoint{7.406549in}{0.894761in}}%
\pgfpathquadraticcurveto{\pgfqpoint{7.391455in}{0.894761in}}{\pgfqpoint{7.382647in}{0.903533in}}%
\pgfpathquadraticcurveto{\pgfqpoint{7.373839in}{0.912323in}}{\pgfqpoint{7.373839in}{0.927309in}}%
\pgfpathquadraticcurveto{\pgfqpoint{7.373839in}{0.942781in}}{\pgfqpoint{7.382196in}{0.951859in}}%
\pgfpathquadraticcurveto{\pgfqpoint{7.390554in}{0.960956in}}{\pgfqpoint{7.404748in}{0.960956in}}%
\pgfpathquadraticcurveto{\pgfqpoint{7.417464in}{0.960956in}}{\pgfqpoint{7.424867in}{0.952760in}}%
\pgfpathquadraticcurveto{\pgfqpoint{7.432270in}{0.944583in}}{\pgfqpoint{7.432270in}{0.930515in}}%
\pgfpathlineto{\pgfqpoint{7.432270in}{0.930515in}}%
\pgfpathclose%
\pgfpathmoveto{\pgfqpoint{7.421913in}{0.933559in}}%
\pgfpathquadraticcurveto{\pgfqpoint{7.421805in}{0.942043in}}{\pgfqpoint{7.417158in}{0.947104in}}%
\pgfpathquadraticcurveto{\pgfqpoint{7.412511in}{0.952184in}}{\pgfqpoint{7.404856in}{0.952184in}}%
\pgfpathquadraticcurveto{\pgfqpoint{7.396192in}{0.952184in}}{\pgfqpoint{7.390986in}{0.947284in}}%
\pgfpathquadraticcurveto{\pgfqpoint{7.385781in}{0.942385in}}{\pgfqpoint{7.384988in}{0.933487in}}%
\pgfpathlineto{\pgfqpoint{7.421913in}{0.933559in}}%
\pgfpathlineto{\pgfqpoint{7.421913in}{0.933559in}}%
\pgfpathclose%
\pgfpathmoveto{\pgfqpoint{7.485802in}{0.949770in}}%
\pgfpathquadraticcurveto{\pgfqpoint{7.484055in}{0.950779in}}{\pgfqpoint{7.482002in}{0.951247in}}%
\pgfpathquadraticcurveto{\pgfqpoint{7.479948in}{0.951733in}}{\pgfqpoint{7.477481in}{0.951733in}}%
\pgfpathquadraticcurveto{\pgfqpoint{7.468691in}{0.951733in}}{\pgfqpoint{7.463990in}{0.946023in}}%
\pgfpathquadraticcurveto{\pgfqpoint{7.459288in}{0.940314in}}{\pgfqpoint{7.459288in}{0.929614in}}%
\pgfpathlineto{\pgfqpoint{7.459288in}{0.896400in}}%
\pgfpathlineto{\pgfqpoint{7.448877in}{0.896400in}}%
\pgfpathlineto{\pgfqpoint{7.448877in}{0.959443in}}%
\pgfpathlineto{\pgfqpoint{7.459288in}{0.959443in}}%
\pgfpathlineto{\pgfqpoint{7.459288in}{0.949644in}}%
\pgfpathquadraticcurveto{\pgfqpoint{7.462567in}{0.955390in}}{\pgfqpoint{7.467790in}{0.958164in}}%
\pgfpathquadraticcurveto{\pgfqpoint{7.473032in}{0.960956in}}{\pgfqpoint{7.480525in}{0.960956in}}%
\pgfpathquadraticcurveto{\pgfqpoint{7.481587in}{0.960956in}}{\pgfqpoint{7.482884in}{0.960811in}}%
\pgfpathquadraticcurveto{\pgfqpoint{7.484181in}{0.960685in}}{\pgfqpoint{7.485748in}{0.960397in}}%
\pgfpathlineto{\pgfqpoint{7.485802in}{0.949770in}}%
\pgfpathlineto{\pgfqpoint{7.485802in}{0.949770in}}%
\pgfpathclose%
\pgfusepath{fill}%
\end{pgfscope}%
\begin{pgfscope}%
\pgfsetbuttcap%
\pgfsetroundjoin%
\definecolor{currentfill}{rgb}{0.309804,0.372549,0.419608}%
\pgfsetfillcolor{currentfill}%
\pgfsetlinewidth{0.000000pt}%
\definecolor{currentstroke}{rgb}{0.309804,0.372549,0.419608}%
\pgfsetstrokecolor{currentstroke}%
\pgfsetdash{}{0pt}%
\pgfpathmoveto{\pgfqpoint{7.608547in}{0.930515in}}%
\pgfpathlineto{\pgfqpoint{7.608547in}{0.925454in}}%
\pgfpathlineto{\pgfqpoint{7.560923in}{0.925454in}}%
\pgfpathquadraticcurveto{\pgfqpoint{7.561607in}{0.914754in}}{\pgfqpoint{7.567371in}{0.909153in}}%
\pgfpathquadraticcurveto{\pgfqpoint{7.573135in}{0.903551in}}{\pgfqpoint{7.583438in}{0.903551in}}%
\pgfpathquadraticcurveto{\pgfqpoint{7.589400in}{0.903551in}}{\pgfqpoint{7.595002in}{0.905010in}}%
\pgfpathquadraticcurveto{\pgfqpoint{7.600603in}{0.906469in}}{\pgfqpoint{7.606133in}{0.909405in}}%
\pgfpathlineto{\pgfqpoint{7.606133in}{0.899606in}}%
\pgfpathquadraticcurveto{\pgfqpoint{7.600549in}{0.897247in}}{\pgfqpoint{7.594695in}{0.896004in}}%
\pgfpathquadraticcurveto{\pgfqpoint{7.588841in}{0.894761in}}{\pgfqpoint{7.582825in}{0.894761in}}%
\pgfpathquadraticcurveto{\pgfqpoint{7.567731in}{0.894761in}}{\pgfqpoint{7.558923in}{0.903533in}}%
\pgfpathquadraticcurveto{\pgfqpoint{7.550115in}{0.912323in}}{\pgfqpoint{7.550115in}{0.927309in}}%
\pgfpathquadraticcurveto{\pgfqpoint{7.550115in}{0.942781in}}{\pgfqpoint{7.558473in}{0.951859in}}%
\pgfpathquadraticcurveto{\pgfqpoint{7.566831in}{0.960956in}}{\pgfqpoint{7.581024in}{0.960956in}}%
\pgfpathquadraticcurveto{\pgfqpoint{7.593741in}{0.960956in}}{\pgfqpoint{7.601144in}{0.952760in}}%
\pgfpathquadraticcurveto{\pgfqpoint{7.608547in}{0.944583in}}{\pgfqpoint{7.608547in}{0.930515in}}%
\pgfpathlineto{\pgfqpoint{7.608547in}{0.930515in}}%
\pgfpathclose%
\pgfpathmoveto{\pgfqpoint{7.598190in}{0.933559in}}%
\pgfpathquadraticcurveto{\pgfqpoint{7.598082in}{0.942043in}}{\pgfqpoint{7.593434in}{0.947104in}}%
\pgfpathquadraticcurveto{\pgfqpoint{7.588787in}{0.952184in}}{\pgfqpoint{7.581132in}{0.952184in}}%
\pgfpathquadraticcurveto{\pgfqpoint{7.572468in}{0.952184in}}{\pgfqpoint{7.567263in}{0.947284in}}%
\pgfpathquadraticcurveto{\pgfqpoint{7.562057in}{0.942385in}}{\pgfqpoint{7.561265in}{0.933487in}}%
\pgfpathlineto{\pgfqpoint{7.598190in}{0.933559in}}%
\pgfpathlineto{\pgfqpoint{7.598190in}{0.933559in}}%
\pgfpathclose%
\pgfpathmoveto{\pgfqpoint{7.665735in}{0.957587in}}%
\pgfpathlineto{\pgfqpoint{7.665735in}{0.947789in}}%
\pgfpathquadraticcurveto{\pgfqpoint{7.661358in}{0.950040in}}{\pgfqpoint{7.656621in}{0.951157in}}%
\pgfpathquadraticcurveto{\pgfqpoint{7.651902in}{0.952292in}}{\pgfqpoint{7.646822in}{0.952292in}}%
\pgfpathquadraticcurveto{\pgfqpoint{7.639113in}{0.952292in}}{\pgfqpoint{7.635259in}{0.949932in}}%
\pgfpathquadraticcurveto{\pgfqpoint{7.631404in}{0.947573in}}{\pgfqpoint{7.631404in}{0.942835in}}%
\pgfpathquadraticcurveto{\pgfqpoint{7.631404in}{0.939233in}}{\pgfqpoint{7.634160in}{0.937180in}}%
\pgfpathquadraticcurveto{\pgfqpoint{7.636916in}{0.935126in}}{\pgfqpoint{7.645255in}{0.933271in}}%
\pgfpathlineto{\pgfqpoint{7.648804in}{0.932478in}}%
\pgfpathquadraticcurveto{\pgfqpoint{7.659827in}{0.930119in}}{\pgfqpoint{7.664474in}{0.925814in}}%
\pgfpathquadraticcurveto{\pgfqpoint{7.669122in}{0.921509in}}{\pgfqpoint{7.669122in}{0.913800in}}%
\pgfpathquadraticcurveto{\pgfqpoint{7.669122in}{0.905010in}}{\pgfqpoint{7.662169in}{0.899876in}}%
\pgfpathquadraticcurveto{\pgfqpoint{7.655216in}{0.894761in}}{\pgfqpoint{7.643058in}{0.894761in}}%
\pgfpathquadraticcurveto{\pgfqpoint{7.637997in}{0.894761in}}{\pgfqpoint{7.632503in}{0.895752in}}%
\pgfpathquadraticcurveto{\pgfqpoint{7.627009in}{0.896742in}}{\pgfqpoint{7.620939in}{0.898706in}}%
\pgfpathlineto{\pgfqpoint{7.620939in}{0.909405in}}%
\pgfpathquadraticcurveto{\pgfqpoint{7.626685in}{0.906415in}}{\pgfqpoint{7.632251in}{0.904920in}}%
\pgfpathquadraticcurveto{\pgfqpoint{7.637816in}{0.903443in}}{\pgfqpoint{7.643292in}{0.903443in}}%
\pgfpathquadraticcurveto{\pgfqpoint{7.650605in}{0.903443in}}{\pgfqpoint{7.654532in}{0.905946in}}%
\pgfpathquadraticcurveto{\pgfqpoint{7.658476in}{0.908450in}}{\pgfqpoint{7.658476in}{0.913007in}}%
\pgfpathquadraticcurveto{\pgfqpoint{7.658476in}{0.917222in}}{\pgfqpoint{7.655630in}{0.919474in}}%
\pgfpathquadraticcurveto{\pgfqpoint{7.652803in}{0.921725in}}{\pgfqpoint{7.643166in}{0.923814in}}%
\pgfpathlineto{\pgfqpoint{7.639564in}{0.924661in}}%
\pgfpathquadraticcurveto{\pgfqpoint{7.629945in}{0.926678in}}{\pgfqpoint{7.625658in}{0.930875in}}%
\pgfpathquadraticcurveto{\pgfqpoint{7.621389in}{0.935072in}}{\pgfqpoint{7.621389in}{0.942385in}}%
\pgfpathquadraticcurveto{\pgfqpoint{7.621389in}{0.951283in}}{\pgfqpoint{7.627694in}{0.956110in}}%
\pgfpathquadraticcurveto{\pgfqpoint{7.633998in}{0.960956in}}{\pgfqpoint{7.645598in}{0.960956in}}%
\pgfpathquadraticcurveto{\pgfqpoint{7.651326in}{0.960956in}}{\pgfqpoint{7.656387in}{0.960109in}}%
\pgfpathquadraticcurveto{\pgfqpoint{7.661466in}{0.959280in}}{\pgfqpoint{7.665735in}{0.957587in}}%
\pgfpathlineto{\pgfqpoint{7.665735in}{0.957587in}}%
\pgfpathclose%
\pgfpathmoveto{\pgfqpoint{7.695852in}{0.977347in}}%
\pgfpathlineto{\pgfqpoint{7.695852in}{0.959443in}}%
\pgfpathlineto{\pgfqpoint{7.717178in}{0.959443in}}%
\pgfpathlineto{\pgfqpoint{7.717178in}{0.951391in}}%
\pgfpathlineto{\pgfqpoint{7.695852in}{0.951391in}}%
\pgfpathlineto{\pgfqpoint{7.695852in}{0.917168in}}%
\pgfpathquadraticcurveto{\pgfqpoint{7.695852in}{0.909459in}}{\pgfqpoint{7.697959in}{0.907261in}}%
\pgfpathquadraticcurveto{\pgfqpoint{7.700066in}{0.905064in}}{\pgfqpoint{7.706551in}{0.905064in}}%
\pgfpathlineto{\pgfqpoint{7.717178in}{0.905064in}}%
\pgfpathlineto{\pgfqpoint{7.717178in}{0.896400in}}%
\pgfpathlineto{\pgfqpoint{7.706551in}{0.896400in}}%
\pgfpathquadraticcurveto{\pgfqpoint{7.694555in}{0.896400in}}{\pgfqpoint{7.689998in}{0.900867in}}%
\pgfpathquadraticcurveto{\pgfqpoint{7.685441in}{0.905352in}}{\pgfqpoint{7.685441in}{0.917168in}}%
\pgfpathlineto{\pgfqpoint{7.685441in}{0.951391in}}%
\pgfpathlineto{\pgfqpoint{7.677839in}{0.951391in}}%
\pgfpathlineto{\pgfqpoint{7.677839in}{0.959443in}}%
\pgfpathlineto{\pgfqpoint{7.685441in}{0.959443in}}%
\pgfpathlineto{\pgfqpoint{7.685441in}{0.977347in}}%
\pgfpathlineto{\pgfqpoint{7.695852in}{0.977347in}}%
\pgfpathlineto{\pgfqpoint{7.695852in}{0.977347in}}%
\pgfpathclose%
\pgfpathmoveto{\pgfqpoint{7.730795in}{0.959443in}}%
\pgfpathlineto{\pgfqpoint{7.741152in}{0.959443in}}%
\pgfpathlineto{\pgfqpoint{7.741152in}{0.896400in}}%
\pgfpathlineto{\pgfqpoint{7.730795in}{0.896400in}}%
\pgfpathlineto{\pgfqpoint{7.730795in}{0.959443in}}%
\pgfpathlineto{\pgfqpoint{7.730795in}{0.959443in}}%
\pgfpathclose%
\pgfpathmoveto{\pgfqpoint{7.730795in}{0.983993in}}%
\pgfpathlineto{\pgfqpoint{7.741152in}{0.983993in}}%
\pgfpathlineto{\pgfqpoint{7.741152in}{0.970862in}}%
\pgfpathlineto{\pgfqpoint{7.730795in}{0.970862in}}%
\pgfpathlineto{\pgfqpoint{7.730795in}{0.983993in}}%
\pgfpathlineto{\pgfqpoint{7.730795in}{0.983993in}}%
\pgfpathclose%
\pgfpathmoveto{\pgfqpoint{7.811904in}{0.947338in}}%
\pgfpathquadraticcurveto{\pgfqpoint{7.815794in}{0.954327in}}{\pgfqpoint{7.821198in}{0.957641in}}%
\pgfpathquadraticcurveto{\pgfqpoint{7.826602in}{0.960956in}}{\pgfqpoint{7.833915in}{0.960956in}}%
\pgfpathquadraticcurveto{\pgfqpoint{7.843767in}{0.960956in}}{\pgfqpoint{7.849117in}{0.954057in}}%
\pgfpathquadraticcurveto{\pgfqpoint{7.854467in}{0.947176in}}{\pgfqpoint{7.854467in}{0.934460in}}%
\pgfpathlineto{\pgfqpoint{7.854467in}{0.896400in}}%
\pgfpathlineto{\pgfqpoint{7.844056in}{0.896400in}}%
\pgfpathlineto{\pgfqpoint{7.844056in}{0.934117in}}%
\pgfpathquadraticcurveto{\pgfqpoint{7.844056in}{0.943178in}}{\pgfqpoint{7.840831in}{0.947555in}}%
\pgfpathquadraticcurveto{\pgfqpoint{7.837625in}{0.951949in}}{\pgfqpoint{7.831051in}{0.951949in}}%
\pgfpathquadraticcurveto{\pgfqpoint{7.822999in}{0.951949in}}{\pgfqpoint{7.818316in}{0.946600in}}%
\pgfpathquadraticcurveto{\pgfqpoint{7.813651in}{0.941268in}}{\pgfqpoint{7.813651in}{0.932028in}}%
\pgfpathlineto{\pgfqpoint{7.813651in}{0.896400in}}%
\pgfpathlineto{\pgfqpoint{7.803240in}{0.896400in}}%
\pgfpathlineto{\pgfqpoint{7.803240in}{0.934117in}}%
\pgfpathquadraticcurveto{\pgfqpoint{7.803240in}{0.943232in}}{\pgfqpoint{7.800034in}{0.947591in}}%
\pgfpathquadraticcurveto{\pgfqpoint{7.796828in}{0.951949in}}{\pgfqpoint{7.790127in}{0.951949in}}%
\pgfpathquadraticcurveto{\pgfqpoint{7.782184in}{0.951949in}}{\pgfqpoint{7.777501in}{0.946582in}}%
\pgfpathquadraticcurveto{\pgfqpoint{7.772836in}{0.941214in}}{\pgfqpoint{7.772836in}{0.932028in}}%
\pgfpathlineto{\pgfqpoint{7.772836in}{0.896400in}}%
\pgfpathlineto{\pgfqpoint{7.762424in}{0.896400in}}%
\pgfpathlineto{\pgfqpoint{7.762424in}{0.959443in}}%
\pgfpathlineto{\pgfqpoint{7.772836in}{0.959443in}}%
\pgfpathlineto{\pgfqpoint{7.772836in}{0.949644in}}%
\pgfpathquadraticcurveto{\pgfqpoint{7.776384in}{0.955444in}}{\pgfqpoint{7.781337in}{0.958200in}}%
\pgfpathquadraticcurveto{\pgfqpoint{7.786291in}{0.960956in}}{\pgfqpoint{7.793099in}{0.960956in}}%
\pgfpathquadraticcurveto{\pgfqpoint{7.799980in}{0.960956in}}{\pgfqpoint{7.804789in}{0.957461in}}%
\pgfpathquadraticcurveto{\pgfqpoint{7.809598in}{0.953985in}}{\pgfqpoint{7.811904in}{0.947338in}}%
\pgfpathlineto{\pgfqpoint{7.811904in}{0.947338in}}%
\pgfpathclose%
\pgfpathmoveto{\pgfqpoint{7.903766in}{0.928083in}}%
\pgfpathquadraticcurveto{\pgfqpoint{7.891211in}{0.928083in}}{\pgfqpoint{7.886366in}{0.925219in}}%
\pgfpathquadraticcurveto{\pgfqpoint{7.881521in}{0.922356in}}{\pgfqpoint{7.881521in}{0.915421in}}%
\pgfpathquadraticcurveto{\pgfqpoint{7.881521in}{0.909909in}}{\pgfqpoint{7.885159in}{0.906667in}}%
\pgfpathquadraticcurveto{\pgfqpoint{7.888798in}{0.903443in}}{\pgfqpoint{7.895030in}{0.903443in}}%
\pgfpathquadraticcurveto{\pgfqpoint{7.903658in}{0.903443in}}{\pgfqpoint{7.908863in}{0.909549in}}%
\pgfpathquadraticcurveto{\pgfqpoint{7.914069in}{0.915655in}}{\pgfqpoint{7.914069in}{0.925778in}}%
\pgfpathlineto{\pgfqpoint{7.914069in}{0.928083in}}%
\pgfpathlineto{\pgfqpoint{7.903766in}{0.928083in}}%
\pgfpathlineto{\pgfqpoint{7.903766in}{0.928083in}}%
\pgfpathclose%
\pgfpathmoveto{\pgfqpoint{7.924426in}{0.932370in}}%
\pgfpathlineto{\pgfqpoint{7.924426in}{0.896400in}}%
\pgfpathlineto{\pgfqpoint{7.914069in}{0.896400in}}%
\pgfpathlineto{\pgfqpoint{7.914069in}{0.905964in}}%
\pgfpathquadraticcurveto{\pgfqpoint{7.910520in}{0.900237in}}{\pgfqpoint{7.905225in}{0.897499in}}%
\pgfpathquadraticcurveto{\pgfqpoint{7.899929in}{0.894761in}}{\pgfqpoint{7.892274in}{0.894761in}}%
\pgfpathquadraticcurveto{\pgfqpoint{7.882602in}{0.894761in}}{\pgfqpoint{7.876874in}{0.900201in}}%
\pgfpathquadraticcurveto{\pgfqpoint{7.871164in}{0.905640in}}{\pgfqpoint{7.871164in}{0.914754in}}%
\pgfpathquadraticcurveto{\pgfqpoint{7.871164in}{0.925382in}}{\pgfqpoint{7.878279in}{0.930785in}}%
\pgfpathquadraticcurveto{\pgfqpoint{7.885411in}{0.936189in}}{\pgfqpoint{7.899533in}{0.936189in}}%
\pgfpathlineto{\pgfqpoint{7.914069in}{0.936189in}}%
\pgfpathlineto{\pgfqpoint{7.914069in}{0.937216in}}%
\pgfpathquadraticcurveto{\pgfqpoint{7.914069in}{0.944366in}}{\pgfqpoint{7.909368in}{0.948275in}}%
\pgfpathquadraticcurveto{\pgfqpoint{7.904666in}{0.952184in}}{\pgfqpoint{7.896165in}{0.952184in}}%
\pgfpathquadraticcurveto{\pgfqpoint{7.890761in}{0.952184in}}{\pgfqpoint{7.885628in}{0.950887in}}%
\pgfpathquadraticcurveto{\pgfqpoint{7.880512in}{0.949590in}}{\pgfqpoint{7.875793in}{0.946996in}}%
\pgfpathlineto{\pgfqpoint{7.875793in}{0.956579in}}%
\pgfpathquadraticcurveto{\pgfqpoint{7.881467in}{0.958776in}}{\pgfqpoint{7.886816in}{0.959857in}}%
\pgfpathquadraticcurveto{\pgfqpoint{7.892166in}{0.960956in}}{\pgfqpoint{7.897227in}{0.960956in}}%
\pgfpathquadraticcurveto{\pgfqpoint{7.910917in}{0.960956in}}{\pgfqpoint{7.917671in}{0.953859in}}%
\pgfpathquadraticcurveto{\pgfqpoint{7.924426in}{0.946780in}}{\pgfqpoint{7.924426in}{0.932370in}}%
\pgfpathlineto{\pgfqpoint{7.924426in}{0.932370in}}%
\pgfpathclose%
\pgfpathmoveto{\pgfqpoint{7.956001in}{0.977347in}}%
\pgfpathlineto{\pgfqpoint{7.956001in}{0.959443in}}%
\pgfpathlineto{\pgfqpoint{7.977327in}{0.959443in}}%
\pgfpathlineto{\pgfqpoint{7.977327in}{0.951391in}}%
\pgfpathlineto{\pgfqpoint{7.956001in}{0.951391in}}%
\pgfpathlineto{\pgfqpoint{7.956001in}{0.917168in}}%
\pgfpathquadraticcurveto{\pgfqpoint{7.956001in}{0.909459in}}{\pgfqpoint{7.958109in}{0.907261in}}%
\pgfpathquadraticcurveto{\pgfqpoint{7.960216in}{0.905064in}}{\pgfqpoint{7.966700in}{0.905064in}}%
\pgfpathlineto{\pgfqpoint{7.977327in}{0.905064in}}%
\pgfpathlineto{\pgfqpoint{7.977327in}{0.896400in}}%
\pgfpathlineto{\pgfqpoint{7.966700in}{0.896400in}}%
\pgfpathquadraticcurveto{\pgfqpoint{7.954704in}{0.896400in}}{\pgfqpoint{7.950147in}{0.900867in}}%
\pgfpathquadraticcurveto{\pgfqpoint{7.945590in}{0.905352in}}{\pgfqpoint{7.945590in}{0.917168in}}%
\pgfpathlineto{\pgfqpoint{7.945590in}{0.951391in}}%
\pgfpathlineto{\pgfqpoint{7.937989in}{0.951391in}}%
\pgfpathlineto{\pgfqpoint{7.937989in}{0.959443in}}%
\pgfpathlineto{\pgfqpoint{7.945590in}{0.959443in}}%
\pgfpathlineto{\pgfqpoint{7.945590in}{0.977347in}}%
\pgfpathlineto{\pgfqpoint{7.956001in}{0.977347in}}%
\pgfpathlineto{\pgfqpoint{7.956001in}{0.977347in}}%
\pgfpathclose%
\pgfpathmoveto{\pgfqpoint{8.044873in}{0.930515in}}%
\pgfpathlineto{\pgfqpoint{8.044873in}{0.925454in}}%
\pgfpathlineto{\pgfqpoint{7.997249in}{0.925454in}}%
\pgfpathquadraticcurveto{\pgfqpoint{7.997933in}{0.914754in}}{\pgfqpoint{8.003697in}{0.909153in}}%
\pgfpathquadraticcurveto{\pgfqpoint{8.009461in}{0.903551in}}{\pgfqpoint{8.019764in}{0.903551in}}%
\pgfpathquadraticcurveto{\pgfqpoint{8.025726in}{0.903551in}}{\pgfqpoint{8.031328in}{0.905010in}}%
\pgfpathquadraticcurveto{\pgfqpoint{8.036930in}{0.906469in}}{\pgfqpoint{8.042459in}{0.909405in}}%
\pgfpathlineto{\pgfqpoint{8.042459in}{0.899606in}}%
\pgfpathquadraticcurveto{\pgfqpoint{8.036876in}{0.897247in}}{\pgfqpoint{8.031022in}{0.896004in}}%
\pgfpathquadraticcurveto{\pgfqpoint{8.025168in}{0.894761in}}{\pgfqpoint{8.019152in}{0.894761in}}%
\pgfpathquadraticcurveto{\pgfqpoint{8.004058in}{0.894761in}}{\pgfqpoint{7.995250in}{0.903533in}}%
\pgfpathquadraticcurveto{\pgfqpoint{7.986442in}{0.912323in}}{\pgfqpoint{7.986442in}{0.927309in}}%
\pgfpathquadraticcurveto{\pgfqpoint{7.986442in}{0.942781in}}{\pgfqpoint{7.994799in}{0.951859in}}%
\pgfpathquadraticcurveto{\pgfqpoint{8.003157in}{0.960956in}}{\pgfqpoint{8.017350in}{0.960956in}}%
\pgfpathquadraticcurveto{\pgfqpoint{8.030067in}{0.960956in}}{\pgfqpoint{8.037470in}{0.952760in}}%
\pgfpathquadraticcurveto{\pgfqpoint{8.044873in}{0.944583in}}{\pgfqpoint{8.044873in}{0.930515in}}%
\pgfpathlineto{\pgfqpoint{8.044873in}{0.930515in}}%
\pgfpathclose%
\pgfpathmoveto{\pgfqpoint{8.034516in}{0.933559in}}%
\pgfpathquadraticcurveto{\pgfqpoint{8.034408in}{0.942043in}}{\pgfqpoint{8.029761in}{0.947104in}}%
\pgfpathquadraticcurveto{\pgfqpoint{8.025114in}{0.952184in}}{\pgfqpoint{8.017459in}{0.952184in}}%
\pgfpathquadraticcurveto{\pgfqpoint{8.008795in}{0.952184in}}{\pgfqpoint{8.003589in}{0.947284in}}%
\pgfpathquadraticcurveto{\pgfqpoint{7.998384in}{0.942385in}}{\pgfqpoint{7.997591in}{0.933487in}}%
\pgfpathlineto{\pgfqpoint{8.034516in}{0.933559in}}%
\pgfpathlineto{\pgfqpoint{8.034516in}{0.933559in}}%
\pgfpathclose%
\pgfusepath{fill}%
\end{pgfscope}%
\begin{pgfscope}%
\pgfsetbuttcap%
\pgfsetroundjoin%
\definecolor{currentfill}{rgb}{0.309804,0.372549,0.419608}%
\pgfsetfillcolor{currentfill}%
\pgfsetlinewidth{0.000000pt}%
\definecolor{currentstroke}{rgb}{0.309804,0.372549,0.419608}%
\pgfsetstrokecolor{currentstroke}%
\pgfsetdash{}{0pt}%
\pgfpathmoveto{\pgfqpoint{8.146527in}{0.952184in}}%
\pgfpathquadraticcurveto{\pgfqpoint{8.138205in}{0.952184in}}{\pgfqpoint{8.133360in}{0.945681in}}%
\pgfpathquadraticcurveto{\pgfqpoint{8.128514in}{0.939179in}}{\pgfqpoint{8.128514in}{0.927867in}}%
\pgfpathquadraticcurveto{\pgfqpoint{8.128514in}{0.916556in}}{\pgfqpoint{8.133324in}{0.910053in}}%
\pgfpathquadraticcurveto{\pgfqpoint{8.138151in}{0.903551in}}{\pgfqpoint{8.146527in}{0.903551in}}%
\pgfpathquadraticcurveto{\pgfqpoint{8.154812in}{0.903551in}}{\pgfqpoint{8.159639in}{0.910071in}}%
\pgfpathquadraticcurveto{\pgfqpoint{8.164485in}{0.916610in}}{\pgfqpoint{8.164485in}{0.927867in}}%
\pgfpathquadraticcurveto{\pgfqpoint{8.164485in}{0.939071in}}{\pgfqpoint{8.159639in}{0.945627in}}%
\pgfpathquadraticcurveto{\pgfqpoint{8.154812in}{0.952184in}}{\pgfqpoint{8.146527in}{0.952184in}}%
\pgfpathlineto{\pgfqpoint{8.146527in}{0.952184in}}%
\pgfpathclose%
\pgfpathmoveto{\pgfqpoint{8.146527in}{0.960956in}}%
\pgfpathquadraticcurveto{\pgfqpoint{8.160036in}{0.960956in}}{\pgfqpoint{8.167745in}{0.952166in}}%
\pgfpathquadraticcurveto{\pgfqpoint{8.175472in}{0.943394in}}{\pgfqpoint{8.175472in}{0.927867in}}%
\pgfpathquadraticcurveto{\pgfqpoint{8.175472in}{0.912395in}}{\pgfqpoint{8.167745in}{0.903569in}}%
\pgfpathquadraticcurveto{\pgfqpoint{8.160036in}{0.894761in}}{\pgfqpoint{8.146527in}{0.894761in}}%
\pgfpathquadraticcurveto{\pgfqpoint{8.132963in}{0.894761in}}{\pgfqpoint{8.125272in}{0.903569in}}%
\pgfpathquadraticcurveto{\pgfqpoint{8.117599in}{0.912395in}}{\pgfqpoint{8.117599in}{0.927867in}}%
\pgfpathquadraticcurveto{\pgfqpoint{8.117599in}{0.943394in}}{\pgfqpoint{8.125272in}{0.952166in}}%
\pgfpathquadraticcurveto{\pgfqpoint{8.132963in}{0.960956in}}{\pgfqpoint{8.146527in}{0.960956in}}%
\pgfpathlineto{\pgfqpoint{8.146527in}{0.960956in}}%
\pgfpathclose%
\pgfpathmoveto{\pgfqpoint{8.229166in}{0.949770in}}%
\pgfpathquadraticcurveto{\pgfqpoint{8.227419in}{0.950779in}}{\pgfqpoint{8.225366in}{0.951247in}}%
\pgfpathquadraticcurveto{\pgfqpoint{8.223312in}{0.951733in}}{\pgfqpoint{8.220845in}{0.951733in}}%
\pgfpathquadraticcurveto{\pgfqpoint{8.212055in}{0.951733in}}{\pgfqpoint{8.207354in}{0.946023in}}%
\pgfpathquadraticcurveto{\pgfqpoint{8.202652in}{0.940314in}}{\pgfqpoint{8.202652in}{0.929614in}}%
\pgfpathlineto{\pgfqpoint{8.202652in}{0.896400in}}%
\pgfpathlineto{\pgfqpoint{8.192241in}{0.896400in}}%
\pgfpathlineto{\pgfqpoint{8.192241in}{0.959443in}}%
\pgfpathlineto{\pgfqpoint{8.202652in}{0.959443in}}%
\pgfpathlineto{\pgfqpoint{8.202652in}{0.949644in}}%
\pgfpathquadraticcurveto{\pgfqpoint{8.205931in}{0.955390in}}{\pgfqpoint{8.211154in}{0.958164in}}%
\pgfpathquadraticcurveto{\pgfqpoint{8.216396in}{0.960956in}}{\pgfqpoint{8.223889in}{0.960956in}}%
\pgfpathquadraticcurveto{\pgfqpoint{8.224951in}{0.960956in}}{\pgfqpoint{8.226248in}{0.960811in}}%
\pgfpathquadraticcurveto{\pgfqpoint{8.227545in}{0.960685in}}{\pgfqpoint{8.229112in}{0.960397in}}%
\pgfpathlineto{\pgfqpoint{8.229166in}{0.949770in}}%
\pgfpathlineto{\pgfqpoint{8.229166in}{0.949770in}}%
\pgfpathclose%
\pgfusepath{fill}%
\end{pgfscope}%
\begin{pgfscope}%
\pgfsetbuttcap%
\pgfsetroundjoin%
\definecolor{currentfill}{rgb}{0.309804,0.372549,0.419608}%
\pgfsetfillcolor{currentfill}%
\pgfsetlinewidth{0.000000pt}%
\definecolor{currentstroke}{rgb}{0.309804,0.372549,0.419608}%
\pgfsetstrokecolor{currentstroke}%
\pgfsetdash{}{0pt}%
\pgfpathmoveto{\pgfqpoint{8.299835in}{0.977347in}}%
\pgfpathlineto{\pgfqpoint{8.299835in}{0.959443in}}%
\pgfpathlineto{\pgfqpoint{8.321161in}{0.959443in}}%
\pgfpathlineto{\pgfqpoint{8.321161in}{0.951391in}}%
\pgfpathlineto{\pgfqpoint{8.299835in}{0.951391in}}%
\pgfpathlineto{\pgfqpoint{8.299835in}{0.917168in}}%
\pgfpathquadraticcurveto{\pgfqpoint{8.299835in}{0.909459in}}{\pgfqpoint{8.301942in}{0.907261in}}%
\pgfpathquadraticcurveto{\pgfqpoint{8.304050in}{0.905064in}}{\pgfqpoint{8.310534in}{0.905064in}}%
\pgfpathlineto{\pgfqpoint{8.321161in}{0.905064in}}%
\pgfpathlineto{\pgfqpoint{8.321161in}{0.896400in}}%
\pgfpathlineto{\pgfqpoint{8.310534in}{0.896400in}}%
\pgfpathquadraticcurveto{\pgfqpoint{8.298538in}{0.896400in}}{\pgfqpoint{8.293981in}{0.900867in}}%
\pgfpathquadraticcurveto{\pgfqpoint{8.289424in}{0.905352in}}{\pgfqpoint{8.289424in}{0.917168in}}%
\pgfpathlineto{\pgfqpoint{8.289424in}{0.951391in}}%
\pgfpathlineto{\pgfqpoint{8.281823in}{0.951391in}}%
\pgfpathlineto{\pgfqpoint{8.281823in}{0.959443in}}%
\pgfpathlineto{\pgfqpoint{8.289424in}{0.959443in}}%
\pgfpathlineto{\pgfqpoint{8.289424in}{0.977347in}}%
\pgfpathlineto{\pgfqpoint{8.299835in}{0.977347in}}%
\pgfpathlineto{\pgfqpoint{8.299835in}{0.977347in}}%
\pgfpathclose%
\pgfpathmoveto{\pgfqpoint{8.387194in}{0.934460in}}%
\pgfpathlineto{\pgfqpoint{8.387194in}{0.896400in}}%
\pgfpathlineto{\pgfqpoint{8.376837in}{0.896400in}}%
\pgfpathlineto{\pgfqpoint{8.376837in}{0.934117in}}%
\pgfpathquadraticcurveto{\pgfqpoint{8.376837in}{0.943069in}}{\pgfqpoint{8.373342in}{0.947500in}}%
\pgfpathquadraticcurveto{\pgfqpoint{8.369848in}{0.951949in}}{\pgfqpoint{8.362877in}{0.951949in}}%
\pgfpathquadraticcurveto{\pgfqpoint{8.354484in}{0.951949in}}{\pgfqpoint{8.349638in}{0.946600in}}%
\pgfpathquadraticcurveto{\pgfqpoint{8.344793in}{0.941268in}}{\pgfqpoint{8.344793in}{0.932028in}}%
\pgfpathlineto{\pgfqpoint{8.344793in}{0.896400in}}%
\pgfpathlineto{\pgfqpoint{8.334382in}{0.896400in}}%
\pgfpathlineto{\pgfqpoint{8.334382in}{0.983993in}}%
\pgfpathlineto{\pgfqpoint{8.344793in}{0.983993in}}%
\pgfpathlineto{\pgfqpoint{8.344793in}{0.949644in}}%
\pgfpathquadraticcurveto{\pgfqpoint{8.348522in}{0.955336in}}{\pgfqpoint{8.353547in}{0.958146in}}%
\pgfpathquadraticcurveto{\pgfqpoint{8.358590in}{0.960956in}}{\pgfqpoint{8.365183in}{0.960956in}}%
\pgfpathquadraticcurveto{\pgfqpoint{8.376044in}{0.960956in}}{\pgfqpoint{8.381610in}{0.954237in}}%
\pgfpathquadraticcurveto{\pgfqpoint{8.387194in}{0.947518in}}{\pgfqpoint{8.387194in}{0.934460in}}%
\pgfpathlineto{\pgfqpoint{8.387194in}{0.934460in}}%
\pgfpathclose%
\pgfpathmoveto{\pgfqpoint{8.461764in}{0.930515in}}%
\pgfpathlineto{\pgfqpoint{8.461764in}{0.925454in}}%
\pgfpathlineto{\pgfqpoint{8.414140in}{0.925454in}}%
\pgfpathquadraticcurveto{\pgfqpoint{8.414824in}{0.914754in}}{\pgfqpoint{8.420588in}{0.909153in}}%
\pgfpathquadraticcurveto{\pgfqpoint{8.426352in}{0.903551in}}{\pgfqpoint{8.436655in}{0.903551in}}%
\pgfpathquadraticcurveto{\pgfqpoint{8.442617in}{0.903551in}}{\pgfqpoint{8.448219in}{0.905010in}}%
\pgfpathquadraticcurveto{\pgfqpoint{8.453821in}{0.906469in}}{\pgfqpoint{8.459350in}{0.909405in}}%
\pgfpathlineto{\pgfqpoint{8.459350in}{0.899606in}}%
\pgfpathquadraticcurveto{\pgfqpoint{8.453767in}{0.897247in}}{\pgfqpoint{8.447913in}{0.896004in}}%
\pgfpathquadraticcurveto{\pgfqpoint{8.442059in}{0.894761in}}{\pgfqpoint{8.436043in}{0.894761in}}%
\pgfpathquadraticcurveto{\pgfqpoint{8.420948in}{0.894761in}}{\pgfqpoint{8.412141in}{0.903533in}}%
\pgfpathquadraticcurveto{\pgfqpoint{8.403333in}{0.912323in}}{\pgfqpoint{8.403333in}{0.927309in}}%
\pgfpathquadraticcurveto{\pgfqpoint{8.403333in}{0.942781in}}{\pgfqpoint{8.411690in}{0.951859in}}%
\pgfpathquadraticcurveto{\pgfqpoint{8.420048in}{0.960956in}}{\pgfqpoint{8.434241in}{0.960956in}}%
\pgfpathquadraticcurveto{\pgfqpoint{8.446958in}{0.960956in}}{\pgfqpoint{8.454361in}{0.952760in}}%
\pgfpathquadraticcurveto{\pgfqpoint{8.461764in}{0.944583in}}{\pgfqpoint{8.461764in}{0.930515in}}%
\pgfpathlineto{\pgfqpoint{8.461764in}{0.930515in}}%
\pgfpathclose%
\pgfpathmoveto{\pgfqpoint{8.451407in}{0.933559in}}%
\pgfpathquadraticcurveto{\pgfqpoint{8.451299in}{0.942043in}}{\pgfqpoint{8.446652in}{0.947104in}}%
\pgfpathquadraticcurveto{\pgfqpoint{8.442005in}{0.952184in}}{\pgfqpoint{8.434350in}{0.952184in}}%
\pgfpathquadraticcurveto{\pgfqpoint{8.425686in}{0.952184in}}{\pgfqpoint{8.420480in}{0.947284in}}%
\pgfpathquadraticcurveto{\pgfqpoint{8.415275in}{0.942385in}}{\pgfqpoint{8.414482in}{0.933487in}}%
\pgfpathlineto{\pgfqpoint{8.451407in}{0.933559in}}%
\pgfpathlineto{\pgfqpoint{8.451407in}{0.933559in}}%
\pgfpathclose%
\pgfpathmoveto{\pgfqpoint{8.478768in}{0.959443in}}%
\pgfpathlineto{\pgfqpoint{8.489124in}{0.959443in}}%
\pgfpathlineto{\pgfqpoint{8.489124in}{0.896400in}}%
\pgfpathlineto{\pgfqpoint{8.478768in}{0.896400in}}%
\pgfpathlineto{\pgfqpoint{8.478768in}{0.959443in}}%
\pgfpathlineto{\pgfqpoint{8.478768in}{0.959443in}}%
\pgfpathclose%
\pgfpathmoveto{\pgfqpoint{8.478768in}{0.983993in}}%
\pgfpathlineto{\pgfqpoint{8.489124in}{0.983993in}}%
\pgfpathlineto{\pgfqpoint{8.489124in}{0.970862in}}%
\pgfpathlineto{\pgfqpoint{8.478768in}{0.970862in}}%
\pgfpathlineto{\pgfqpoint{8.478768in}{0.983993in}}%
\pgfpathlineto{\pgfqpoint{8.478768in}{0.983993in}}%
\pgfpathclose%
\pgfpathmoveto{\pgfqpoint{8.547322in}{0.949770in}}%
\pgfpathquadraticcurveto{\pgfqpoint{8.545575in}{0.950779in}}{\pgfqpoint{8.543521in}{0.951247in}}%
\pgfpathquadraticcurveto{\pgfqpoint{8.541468in}{0.951733in}}{\pgfqpoint{8.539000in}{0.951733in}}%
\pgfpathquadraticcurveto{\pgfqpoint{8.530210in}{0.951733in}}{\pgfqpoint{8.525509in}{0.946023in}}%
\pgfpathquadraticcurveto{\pgfqpoint{8.520808in}{0.940314in}}{\pgfqpoint{8.520808in}{0.929614in}}%
\pgfpathlineto{\pgfqpoint{8.520808in}{0.896400in}}%
\pgfpathlineto{\pgfqpoint{8.510397in}{0.896400in}}%
\pgfpathlineto{\pgfqpoint{8.510397in}{0.959443in}}%
\pgfpathlineto{\pgfqpoint{8.520808in}{0.959443in}}%
\pgfpathlineto{\pgfqpoint{8.520808in}{0.949644in}}%
\pgfpathquadraticcurveto{\pgfqpoint{8.524086in}{0.955390in}}{\pgfqpoint{8.529310in}{0.958164in}}%
\pgfpathquadraticcurveto{\pgfqpoint{8.534551in}{0.960956in}}{\pgfqpoint{8.542044in}{0.960956in}}%
\pgfpathquadraticcurveto{\pgfqpoint{8.543107in}{0.960956in}}{\pgfqpoint{8.544404in}{0.960811in}}%
\pgfpathquadraticcurveto{\pgfqpoint{8.545701in}{0.960685in}}{\pgfqpoint{8.547268in}{0.960397in}}%
\pgfpathlineto{\pgfqpoint{8.547322in}{0.949770in}}%
\pgfpathlineto{\pgfqpoint{8.547322in}{0.949770in}}%
\pgfpathclose%
\pgfusepath{fill}%
\end{pgfscope}%
\begin{pgfscope}%
\pgfsetbuttcap%
\pgfsetroundjoin%
\definecolor{currentfill}{rgb}{0.309804,0.372549,0.419608}%
\pgfsetfillcolor{currentfill}%
\pgfsetlinewidth{0.000000pt}%
\definecolor{currentstroke}{rgb}{0.309804,0.372549,0.419608}%
\pgfsetstrokecolor{currentstroke}%
\pgfsetdash{}{0pt}%
\pgfpathmoveto{\pgfqpoint{8.658552in}{0.949878in}}%
\pgfpathlineto{\pgfqpoint{8.658552in}{0.983993in}}%
\pgfpathlineto{\pgfqpoint{8.668909in}{0.983993in}}%
\pgfpathlineto{\pgfqpoint{8.668909in}{0.896400in}}%
\pgfpathlineto{\pgfqpoint{8.658552in}{0.896400in}}%
\pgfpathlineto{\pgfqpoint{8.658552in}{0.905856in}}%
\pgfpathquadraticcurveto{\pgfqpoint{8.655291in}{0.900237in}}{\pgfqpoint{8.650302in}{0.897499in}}%
\pgfpathquadraticcurveto{\pgfqpoint{8.645331in}{0.894761in}}{\pgfqpoint{8.638342in}{0.894761in}}%
\pgfpathquadraticcurveto{\pgfqpoint{8.626922in}{0.894761in}}{\pgfqpoint{8.619735in}{0.903875in}}%
\pgfpathquadraticcurveto{\pgfqpoint{8.612567in}{0.913007in}}{\pgfqpoint{8.612567in}{0.927867in}}%
\pgfpathquadraticcurveto{\pgfqpoint{8.612567in}{0.942727in}}{\pgfqpoint{8.619735in}{0.951841in}}%
\pgfpathquadraticcurveto{\pgfqpoint{8.626922in}{0.960956in}}{\pgfqpoint{8.638342in}{0.960956in}}%
\pgfpathquadraticcurveto{\pgfqpoint{8.645331in}{0.960956in}}{\pgfqpoint{8.650302in}{0.958218in}}%
\pgfpathquadraticcurveto{\pgfqpoint{8.655291in}{0.955498in}}{\pgfqpoint{8.658552in}{0.949878in}}%
\pgfpathlineto{\pgfqpoint{8.658552in}{0.949878in}}%
\pgfpathclose%
\pgfpathmoveto{\pgfqpoint{8.623266in}{0.927867in}}%
\pgfpathquadraticcurveto{\pgfqpoint{8.623266in}{0.916448in}}{\pgfqpoint{8.627967in}{0.909945in}}%
\pgfpathquadraticcurveto{\pgfqpoint{8.632668in}{0.903443in}}{\pgfqpoint{8.640882in}{0.903443in}}%
\pgfpathquadraticcurveto{\pgfqpoint{8.649095in}{0.903443in}}{\pgfqpoint{8.653814in}{0.909945in}}%
\pgfpathquadraticcurveto{\pgfqpoint{8.658552in}{0.916448in}}{\pgfqpoint{8.658552in}{0.927867in}}%
\pgfpathquadraticcurveto{\pgfqpoint{8.658552in}{0.939287in}}{\pgfqpoint{8.653814in}{0.945789in}}%
\pgfpathquadraticcurveto{\pgfqpoint{8.649095in}{0.952292in}}{\pgfqpoint{8.640882in}{0.952292in}}%
\pgfpathquadraticcurveto{\pgfqpoint{8.632668in}{0.952292in}}{\pgfqpoint{8.627967in}{0.945789in}}%
\pgfpathquadraticcurveto{\pgfqpoint{8.623266in}{0.939287in}}{\pgfqpoint{8.623266in}{0.927867in}}%
\pgfpathlineto{\pgfqpoint{8.623266in}{0.927867in}}%
\pgfpathclose%
\pgfpathmoveto{\pgfqpoint{8.690253in}{0.959443in}}%
\pgfpathlineto{\pgfqpoint{8.700610in}{0.959443in}}%
\pgfpathlineto{\pgfqpoint{8.700610in}{0.896400in}}%
\pgfpathlineto{\pgfqpoint{8.690253in}{0.896400in}}%
\pgfpathlineto{\pgfqpoint{8.690253in}{0.959443in}}%
\pgfpathlineto{\pgfqpoint{8.690253in}{0.959443in}}%
\pgfpathclose%
\pgfpathmoveto{\pgfqpoint{8.690253in}{0.983993in}}%
\pgfpathlineto{\pgfqpoint{8.700610in}{0.983993in}}%
\pgfpathlineto{\pgfqpoint{8.700610in}{0.970862in}}%
\pgfpathlineto{\pgfqpoint{8.690253in}{0.970862in}}%
\pgfpathlineto{\pgfqpoint{8.690253in}{0.983993in}}%
\pgfpathlineto{\pgfqpoint{8.690253in}{0.983993in}}%
\pgfpathclose%
\pgfpathmoveto{\pgfqpoint{8.793030in}{0.983993in}}%
\pgfpathlineto{\pgfqpoint{8.793030in}{0.975365in}}%
\pgfpathlineto{\pgfqpoint{8.783124in}{0.975365in}}%
\pgfpathquadraticcurveto{\pgfqpoint{8.777558in}{0.975365in}}{\pgfqpoint{8.775378in}{0.973114in}}%
\pgfpathquadraticcurveto{\pgfqpoint{8.773217in}{0.970862in}}{\pgfqpoint{8.773217in}{0.965008in}}%
\pgfpathlineto{\pgfqpoint{8.773217in}{0.959443in}}%
\pgfpathlineto{\pgfqpoint{8.790275in}{0.959443in}}%
\pgfpathlineto{\pgfqpoint{8.790275in}{0.951391in}}%
\pgfpathlineto{\pgfqpoint{8.773217in}{0.951391in}}%
\pgfpathlineto{\pgfqpoint{8.773217in}{0.896400in}}%
\pgfpathlineto{\pgfqpoint{8.762806in}{0.896400in}}%
\pgfpathlineto{\pgfqpoint{8.762806in}{0.951391in}}%
\pgfpathlineto{\pgfqpoint{8.734383in}{0.951391in}}%
\pgfpathlineto{\pgfqpoint{8.734383in}{0.896400in}}%
\pgfpathlineto{\pgfqpoint{8.723972in}{0.896400in}}%
\pgfpathlineto{\pgfqpoint{8.723972in}{0.951391in}}%
\pgfpathlineto{\pgfqpoint{8.714065in}{0.951391in}}%
\pgfpathlineto{\pgfqpoint{8.714065in}{0.959443in}}%
\pgfpathlineto{\pgfqpoint{8.723972in}{0.959443in}}%
\pgfpathlineto{\pgfqpoint{8.723972in}{0.963838in}}%
\pgfpathquadraticcurveto{\pgfqpoint{8.723972in}{0.974357in}}{\pgfqpoint{8.728871in}{0.979166in}}%
\pgfpathquadraticcurveto{\pgfqpoint{8.733770in}{0.983993in}}{\pgfqpoint{8.744398in}{0.983993in}}%
\pgfpathlineto{\pgfqpoint{8.754196in}{0.983993in}}%
\pgfpathlineto{\pgfqpoint{8.754196in}{0.975365in}}%
\pgfpathlineto{\pgfqpoint{8.744289in}{0.975365in}}%
\pgfpathquadraticcurveto{\pgfqpoint{8.738724in}{0.975365in}}{\pgfqpoint{8.736526in}{0.973114in}}%
\pgfpathquadraticcurveto{\pgfqpoint{8.734383in}{0.970862in}}{\pgfqpoint{8.734383in}{0.965008in}}%
\pgfpathlineto{\pgfqpoint{8.734383in}{0.959443in}}%
\pgfpathlineto{\pgfqpoint{8.762806in}{0.959443in}}%
\pgfpathlineto{\pgfqpoint{8.762806in}{0.963838in}}%
\pgfpathquadraticcurveto{\pgfqpoint{8.762806in}{0.974357in}}{\pgfqpoint{8.767705in}{0.979166in}}%
\pgfpathquadraticcurveto{\pgfqpoint{8.772605in}{0.983993in}}{\pgfqpoint{8.783250in}{0.983993in}}%
\pgfpathlineto{\pgfqpoint{8.793030in}{0.983993in}}%
\pgfpathlineto{\pgfqpoint{8.793030in}{0.983993in}}%
\pgfpathclose%
\pgfpathmoveto{\pgfqpoint{8.855623in}{0.930515in}}%
\pgfpathlineto{\pgfqpoint{8.855623in}{0.925454in}}%
\pgfpathlineto{\pgfqpoint{8.807998in}{0.925454in}}%
\pgfpathquadraticcurveto{\pgfqpoint{8.808683in}{0.914754in}}{\pgfqpoint{8.814447in}{0.909153in}}%
\pgfpathquadraticcurveto{\pgfqpoint{8.820211in}{0.903551in}}{\pgfqpoint{8.830514in}{0.903551in}}%
\pgfpathquadraticcurveto{\pgfqpoint{8.836476in}{0.903551in}}{\pgfqpoint{8.842077in}{0.905010in}}%
\pgfpathquadraticcurveto{\pgfqpoint{8.847679in}{0.906469in}}{\pgfqpoint{8.853209in}{0.909405in}}%
\pgfpathlineto{\pgfqpoint{8.853209in}{0.899606in}}%
\pgfpathquadraticcurveto{\pgfqpoint{8.847625in}{0.897247in}}{\pgfqpoint{8.841771in}{0.896004in}}%
\pgfpathquadraticcurveto{\pgfqpoint{8.835917in}{0.894761in}}{\pgfqpoint{8.829901in}{0.894761in}}%
\pgfpathquadraticcurveto{\pgfqpoint{8.814807in}{0.894761in}}{\pgfqpoint{8.805999in}{0.903533in}}%
\pgfpathquadraticcurveto{\pgfqpoint{8.797191in}{0.912323in}}{\pgfqpoint{8.797191in}{0.927309in}}%
\pgfpathquadraticcurveto{\pgfqpoint{8.797191in}{0.942781in}}{\pgfqpoint{8.805549in}{0.951859in}}%
\pgfpathquadraticcurveto{\pgfqpoint{8.813906in}{0.960956in}}{\pgfqpoint{8.828100in}{0.960956in}}%
\pgfpathquadraticcurveto{\pgfqpoint{8.840817in}{0.960956in}}{\pgfqpoint{8.848220in}{0.952760in}}%
\pgfpathquadraticcurveto{\pgfqpoint{8.855623in}{0.944583in}}{\pgfqpoint{8.855623in}{0.930515in}}%
\pgfpathlineto{\pgfqpoint{8.855623in}{0.930515in}}%
\pgfpathclose%
\pgfpathmoveto{\pgfqpoint{8.845266in}{0.933559in}}%
\pgfpathquadraticcurveto{\pgfqpoint{8.845158in}{0.942043in}}{\pgfqpoint{8.840510in}{0.947104in}}%
\pgfpathquadraticcurveto{\pgfqpoint{8.835863in}{0.952184in}}{\pgfqpoint{8.828208in}{0.952184in}}%
\pgfpathquadraticcurveto{\pgfqpoint{8.819544in}{0.952184in}}{\pgfqpoint{8.814339in}{0.947284in}}%
\pgfpathquadraticcurveto{\pgfqpoint{8.809133in}{0.942385in}}{\pgfqpoint{8.808341in}{0.933487in}}%
\pgfpathlineto{\pgfqpoint{8.845266in}{0.933559in}}%
\pgfpathlineto{\pgfqpoint{8.845266in}{0.933559in}}%
\pgfpathclose%
\pgfpathmoveto{\pgfqpoint{8.909155in}{0.949770in}}%
\pgfpathquadraticcurveto{\pgfqpoint{8.907408in}{0.950779in}}{\pgfqpoint{8.905354in}{0.951247in}}%
\pgfpathquadraticcurveto{\pgfqpoint{8.903301in}{0.951733in}}{\pgfqpoint{8.900833in}{0.951733in}}%
\pgfpathquadraticcurveto{\pgfqpoint{8.892043in}{0.951733in}}{\pgfqpoint{8.887342in}{0.946023in}}%
\pgfpathquadraticcurveto{\pgfqpoint{8.882641in}{0.940314in}}{\pgfqpoint{8.882641in}{0.929614in}}%
\pgfpathlineto{\pgfqpoint{8.882641in}{0.896400in}}%
\pgfpathlineto{\pgfqpoint{8.872230in}{0.896400in}}%
\pgfpathlineto{\pgfqpoint{8.872230in}{0.959443in}}%
\pgfpathlineto{\pgfqpoint{8.882641in}{0.959443in}}%
\pgfpathlineto{\pgfqpoint{8.882641in}{0.949644in}}%
\pgfpathquadraticcurveto{\pgfqpoint{8.885919in}{0.955390in}}{\pgfqpoint{8.891143in}{0.958164in}}%
\pgfpathquadraticcurveto{\pgfqpoint{8.896384in}{0.960956in}}{\pgfqpoint{8.903877in}{0.960956in}}%
\pgfpathquadraticcurveto{\pgfqpoint{8.904940in}{0.960956in}}{\pgfqpoint{8.906237in}{0.960811in}}%
\pgfpathquadraticcurveto{\pgfqpoint{8.907534in}{0.960685in}}{\pgfqpoint{8.909101in}{0.960397in}}%
\pgfpathlineto{\pgfqpoint{8.909155in}{0.949770in}}%
\pgfpathlineto{\pgfqpoint{8.909155in}{0.949770in}}%
\pgfpathclose%
\pgfpathmoveto{\pgfqpoint{8.971405in}{0.930515in}}%
\pgfpathlineto{\pgfqpoint{8.971405in}{0.925454in}}%
\pgfpathlineto{\pgfqpoint{8.923781in}{0.925454in}}%
\pgfpathquadraticcurveto{\pgfqpoint{8.924465in}{0.914754in}}{\pgfqpoint{8.930229in}{0.909153in}}%
\pgfpathquadraticcurveto{\pgfqpoint{8.935993in}{0.903551in}}{\pgfqpoint{8.946296in}{0.903551in}}%
\pgfpathquadraticcurveto{\pgfqpoint{8.952258in}{0.903551in}}{\pgfqpoint{8.957860in}{0.905010in}}%
\pgfpathquadraticcurveto{\pgfqpoint{8.963461in}{0.906469in}}{\pgfqpoint{8.968991in}{0.909405in}}%
\pgfpathlineto{\pgfqpoint{8.968991in}{0.899606in}}%
\pgfpathquadraticcurveto{\pgfqpoint{8.963407in}{0.897247in}}{\pgfqpoint{8.957553in}{0.896004in}}%
\pgfpathquadraticcurveto{\pgfqpoint{8.951699in}{0.894761in}}{\pgfqpoint{8.945683in}{0.894761in}}%
\pgfpathquadraticcurveto{\pgfqpoint{8.930589in}{0.894761in}}{\pgfqpoint{8.921781in}{0.903533in}}%
\pgfpathquadraticcurveto{\pgfqpoint{8.912973in}{0.912323in}}{\pgfqpoint{8.912973in}{0.927309in}}%
\pgfpathquadraticcurveto{\pgfqpoint{8.912973in}{0.942781in}}{\pgfqpoint{8.921331in}{0.951859in}}%
\pgfpathquadraticcurveto{\pgfqpoint{8.929689in}{0.960956in}}{\pgfqpoint{8.943882in}{0.960956in}}%
\pgfpathquadraticcurveto{\pgfqpoint{8.956599in}{0.960956in}}{\pgfqpoint{8.964002in}{0.952760in}}%
\pgfpathquadraticcurveto{\pgfqpoint{8.971405in}{0.944583in}}{\pgfqpoint{8.971405in}{0.930515in}}%
\pgfpathlineto{\pgfqpoint{8.971405in}{0.930515in}}%
\pgfpathclose%
\pgfpathmoveto{\pgfqpoint{8.961048in}{0.933559in}}%
\pgfpathquadraticcurveto{\pgfqpoint{8.960940in}{0.942043in}}{\pgfqpoint{8.956293in}{0.947104in}}%
\pgfpathquadraticcurveto{\pgfqpoint{8.951645in}{0.952184in}}{\pgfqpoint{8.943990in}{0.952184in}}%
\pgfpathquadraticcurveto{\pgfqpoint{8.935326in}{0.952184in}}{\pgfqpoint{8.930121in}{0.947284in}}%
\pgfpathquadraticcurveto{\pgfqpoint{8.924915in}{0.942385in}}{\pgfqpoint{8.924123in}{0.933487in}}%
\pgfpathlineto{\pgfqpoint{8.961048in}{0.933559in}}%
\pgfpathlineto{\pgfqpoint{8.961048in}{0.933559in}}%
\pgfpathclose%
\pgfpathmoveto{\pgfqpoint{9.040824in}{0.934460in}}%
\pgfpathlineto{\pgfqpoint{9.040824in}{0.896400in}}%
\pgfpathlineto{\pgfqpoint{9.030467in}{0.896400in}}%
\pgfpathlineto{\pgfqpoint{9.030467in}{0.934117in}}%
\pgfpathquadraticcurveto{\pgfqpoint{9.030467in}{0.943069in}}{\pgfqpoint{9.026972in}{0.947500in}}%
\pgfpathquadraticcurveto{\pgfqpoint{9.023478in}{0.951949in}}{\pgfqpoint{9.016507in}{0.951949in}}%
\pgfpathquadraticcurveto{\pgfqpoint{9.008113in}{0.951949in}}{\pgfqpoint{9.003268in}{0.946600in}}%
\pgfpathquadraticcurveto{\pgfqpoint{8.998423in}{0.941268in}}{\pgfqpoint{8.998423in}{0.932028in}}%
\pgfpathlineto{\pgfqpoint{8.998423in}{0.896400in}}%
\pgfpathlineto{\pgfqpoint{8.988012in}{0.896400in}}%
\pgfpathlineto{\pgfqpoint{8.988012in}{0.959443in}}%
\pgfpathlineto{\pgfqpoint{8.998423in}{0.959443in}}%
\pgfpathlineto{\pgfqpoint{8.998423in}{0.949644in}}%
\pgfpathquadraticcurveto{\pgfqpoint{9.002151in}{0.955336in}}{\pgfqpoint{9.007177in}{0.958146in}}%
\pgfpathquadraticcurveto{\pgfqpoint{9.012220in}{0.960956in}}{\pgfqpoint{9.018813in}{0.960956in}}%
\pgfpathquadraticcurveto{\pgfqpoint{9.029674in}{0.960956in}}{\pgfqpoint{9.035240in}{0.954237in}}%
\pgfpathquadraticcurveto{\pgfqpoint{9.040824in}{0.947518in}}{\pgfqpoint{9.040824in}{0.934460in}}%
\pgfpathlineto{\pgfqpoint{9.040824in}{0.934460in}}%
\pgfpathclose%
\pgfpathmoveto{\pgfqpoint{9.106838in}{0.957029in}}%
\pgfpathlineto{\pgfqpoint{9.106838in}{0.947338in}}%
\pgfpathquadraticcurveto{\pgfqpoint{9.102443in}{0.949770in}}{\pgfqpoint{9.098030in}{0.950977in}}%
\pgfpathquadraticcurveto{\pgfqpoint{9.093617in}{0.952184in}}{\pgfqpoint{9.089114in}{0.952184in}}%
\pgfpathquadraticcurveto{\pgfqpoint{9.079027in}{0.952184in}}{\pgfqpoint{9.073444in}{0.945789in}}%
\pgfpathquadraticcurveto{\pgfqpoint{9.067878in}{0.939413in}}{\pgfqpoint{9.067878in}{0.927867in}}%
\pgfpathquadraticcurveto{\pgfqpoint{9.067878in}{0.916321in}}{\pgfqpoint{9.073444in}{0.909927in}}%
\pgfpathquadraticcurveto{\pgfqpoint{9.079027in}{0.903551in}}{\pgfqpoint{9.089114in}{0.903551in}}%
\pgfpathquadraticcurveto{\pgfqpoint{9.093617in}{0.903551in}}{\pgfqpoint{9.098030in}{0.904758in}}%
\pgfpathquadraticcurveto{\pgfqpoint{9.102443in}{0.905964in}}{\pgfqpoint{9.106838in}{0.908396in}}%
\pgfpathlineto{\pgfqpoint{9.106838in}{0.898814in}}%
\pgfpathquadraticcurveto{\pgfqpoint{9.102497in}{0.896796in}}{\pgfqpoint{9.097850in}{0.895788in}}%
\pgfpathquadraticcurveto{\pgfqpoint{9.093221in}{0.894761in}}{\pgfqpoint{9.087979in}{0.894761in}}%
\pgfpathquadraticcurveto{\pgfqpoint{9.073732in}{0.894761in}}{\pgfqpoint{9.065338in}{0.903713in}}%
\pgfpathquadraticcurveto{\pgfqpoint{9.056962in}{0.912665in}}{\pgfqpoint{9.056962in}{0.927867in}}%
\pgfpathquadraticcurveto{\pgfqpoint{9.056962in}{0.943286in}}{\pgfqpoint{9.065428in}{0.952112in}}%
\pgfpathquadraticcurveto{\pgfqpoint{9.073912in}{0.960956in}}{\pgfqpoint{9.088664in}{0.960956in}}%
\pgfpathquadraticcurveto{\pgfqpoint{9.093437in}{0.960956in}}{\pgfqpoint{9.097994in}{0.959965in}}%
\pgfpathquadraticcurveto{\pgfqpoint{9.102551in}{0.958992in}}{\pgfqpoint{9.106838in}{0.957029in}}%
\pgfpathlineto{\pgfqpoint{9.106838in}{0.957029in}}%
\pgfpathclose%
\pgfpathmoveto{\pgfqpoint{9.178779in}{0.930515in}}%
\pgfpathlineto{\pgfqpoint{9.178779in}{0.925454in}}%
\pgfpathlineto{\pgfqpoint{9.131155in}{0.925454in}}%
\pgfpathquadraticcurveto{\pgfqpoint{9.131839in}{0.914754in}}{\pgfqpoint{9.137603in}{0.909153in}}%
\pgfpathquadraticcurveto{\pgfqpoint{9.143367in}{0.903551in}}{\pgfqpoint{9.153670in}{0.903551in}}%
\pgfpathquadraticcurveto{\pgfqpoint{9.159632in}{0.903551in}}{\pgfqpoint{9.165234in}{0.905010in}}%
\pgfpathquadraticcurveto{\pgfqpoint{9.170835in}{0.906469in}}{\pgfqpoint{9.176365in}{0.909405in}}%
\pgfpathlineto{\pgfqpoint{9.176365in}{0.899606in}}%
\pgfpathquadraticcurveto{\pgfqpoint{9.170781in}{0.897247in}}{\pgfqpoint{9.164927in}{0.896004in}}%
\pgfpathquadraticcurveto{\pgfqpoint{9.159073in}{0.894761in}}{\pgfqpoint{9.153057in}{0.894761in}}%
\pgfpathquadraticcurveto{\pgfqpoint{9.137963in}{0.894761in}}{\pgfqpoint{9.129155in}{0.903533in}}%
\pgfpathquadraticcurveto{\pgfqpoint{9.120347in}{0.912323in}}{\pgfqpoint{9.120347in}{0.927309in}}%
\pgfpathquadraticcurveto{\pgfqpoint{9.120347in}{0.942781in}}{\pgfqpoint{9.128705in}{0.951859in}}%
\pgfpathquadraticcurveto{\pgfqpoint{9.137063in}{0.960956in}}{\pgfqpoint{9.151256in}{0.960956in}}%
\pgfpathquadraticcurveto{\pgfqpoint{9.163973in}{0.960956in}}{\pgfqpoint{9.171376in}{0.952760in}}%
\pgfpathquadraticcurveto{\pgfqpoint{9.178779in}{0.944583in}}{\pgfqpoint{9.178779in}{0.930515in}}%
\pgfpathlineto{\pgfqpoint{9.178779in}{0.930515in}}%
\pgfpathclose%
\pgfpathmoveto{\pgfqpoint{9.168422in}{0.933559in}}%
\pgfpathquadraticcurveto{\pgfqpoint{9.168314in}{0.942043in}}{\pgfqpoint{9.163666in}{0.947104in}}%
\pgfpathquadraticcurveto{\pgfqpoint{9.159019in}{0.952184in}}{\pgfqpoint{9.151364in}{0.952184in}}%
\pgfpathquadraticcurveto{\pgfqpoint{9.142700in}{0.952184in}}{\pgfqpoint{9.137495in}{0.947284in}}%
\pgfpathquadraticcurveto{\pgfqpoint{9.132289in}{0.942385in}}{\pgfqpoint{9.131497in}{0.933487in}}%
\pgfpathlineto{\pgfqpoint{9.168422in}{0.933559in}}%
\pgfpathlineto{\pgfqpoint{9.168422in}{0.933559in}}%
\pgfpathclose%
\pgfpathmoveto{\pgfqpoint{9.198430in}{0.910702in}}%
\pgfpathlineto{\pgfqpoint{9.210300in}{0.910702in}}%
\pgfpathlineto{\pgfqpoint{9.210300in}{0.901011in}}%
\pgfpathlineto{\pgfqpoint{9.201078in}{0.882999in}}%
\pgfpathlineto{\pgfqpoint{9.193819in}{0.882999in}}%
\pgfpathlineto{\pgfqpoint{9.198430in}{0.901011in}}%
\pgfpathlineto{\pgfqpoint{9.198430in}{0.910702in}}%
\pgfpathlineto{\pgfqpoint{9.198430in}{0.910702in}}%
\pgfpathclose%
\pgfusepath{fill}%
\end{pgfscope}%
\begin{pgfscope}%
\pgfsetbuttcap%
\pgfsetroundjoin%
\definecolor{currentfill}{rgb}{0.309804,0.372549,0.419608}%
\pgfsetfillcolor{currentfill}%
\pgfsetlinewidth{0.000000pt}%
\definecolor{currentstroke}{rgb}{0.309804,0.372549,0.419608}%
\pgfsetstrokecolor{currentstroke}%
\pgfsetdash{}{0pt}%
\pgfpathmoveto{\pgfqpoint{9.326478in}{0.952184in}}%
\pgfpathquadraticcurveto{\pgfqpoint{9.318156in}{0.952184in}}{\pgfqpoint{9.313311in}{0.945681in}}%
\pgfpathquadraticcurveto{\pgfqpoint{9.308466in}{0.939179in}}{\pgfqpoint{9.308466in}{0.927867in}}%
\pgfpathquadraticcurveto{\pgfqpoint{9.308466in}{0.916556in}}{\pgfqpoint{9.313275in}{0.910053in}}%
\pgfpathquadraticcurveto{\pgfqpoint{9.318102in}{0.903551in}}{\pgfqpoint{9.326478in}{0.903551in}}%
\pgfpathquadraticcurveto{\pgfqpoint{9.334764in}{0.903551in}}{\pgfqpoint{9.339591in}{0.910071in}}%
\pgfpathquadraticcurveto{\pgfqpoint{9.344436in}{0.916610in}}{\pgfqpoint{9.344436in}{0.927867in}}%
\pgfpathquadraticcurveto{\pgfqpoint{9.344436in}{0.939071in}}{\pgfqpoint{9.339591in}{0.945627in}}%
\pgfpathquadraticcurveto{\pgfqpoint{9.334764in}{0.952184in}}{\pgfqpoint{9.326478in}{0.952184in}}%
\pgfpathlineto{\pgfqpoint{9.326478in}{0.952184in}}%
\pgfpathclose%
\pgfpathmoveto{\pgfqpoint{9.326478in}{0.960956in}}%
\pgfpathquadraticcurveto{\pgfqpoint{9.339987in}{0.960956in}}{\pgfqpoint{9.347696in}{0.952166in}}%
\pgfpathquadraticcurveto{\pgfqpoint{9.355423in}{0.943394in}}{\pgfqpoint{9.355423in}{0.927867in}}%
\pgfpathquadraticcurveto{\pgfqpoint{9.355423in}{0.912395in}}{\pgfqpoint{9.347696in}{0.903569in}}%
\pgfpathquadraticcurveto{\pgfqpoint{9.339987in}{0.894761in}}{\pgfqpoint{9.326478in}{0.894761in}}%
\pgfpathquadraticcurveto{\pgfqpoint{9.312915in}{0.894761in}}{\pgfqpoint{9.305224in}{0.903569in}}%
\pgfpathquadraticcurveto{\pgfqpoint{9.297550in}{0.912395in}}{\pgfqpoint{9.297550in}{0.927867in}}%
\pgfpathquadraticcurveto{\pgfqpoint{9.297550in}{0.943394in}}{\pgfqpoint{9.305224in}{0.952166in}}%
\pgfpathquadraticcurveto{\pgfqpoint{9.312915in}{0.960956in}}{\pgfqpoint{9.326478in}{0.960956in}}%
\pgfpathlineto{\pgfqpoint{9.326478in}{0.960956in}}%
\pgfpathclose%
\pgfpathmoveto{\pgfqpoint{9.409118in}{0.949770in}}%
\pgfpathquadraticcurveto{\pgfqpoint{9.407371in}{0.950779in}}{\pgfqpoint{9.405317in}{0.951247in}}%
\pgfpathquadraticcurveto{\pgfqpoint{9.403264in}{0.951733in}}{\pgfqpoint{9.400796in}{0.951733in}}%
\pgfpathquadraticcurveto{\pgfqpoint{9.392006in}{0.951733in}}{\pgfqpoint{9.387305in}{0.946023in}}%
\pgfpathquadraticcurveto{\pgfqpoint{9.382604in}{0.940314in}}{\pgfqpoint{9.382604in}{0.929614in}}%
\pgfpathlineto{\pgfqpoint{9.382604in}{0.896400in}}%
\pgfpathlineto{\pgfqpoint{9.372193in}{0.896400in}}%
\pgfpathlineto{\pgfqpoint{9.372193in}{0.959443in}}%
\pgfpathlineto{\pgfqpoint{9.382604in}{0.959443in}}%
\pgfpathlineto{\pgfqpoint{9.382604in}{0.949644in}}%
\pgfpathquadraticcurveto{\pgfqpoint{9.385882in}{0.955390in}}{\pgfqpoint{9.391106in}{0.958164in}}%
\pgfpathquadraticcurveto{\pgfqpoint{9.396347in}{0.960956in}}{\pgfqpoint{9.403840in}{0.960956in}}%
\pgfpathquadraticcurveto{\pgfqpoint{9.404903in}{0.960956in}}{\pgfqpoint{9.406200in}{0.960811in}}%
\pgfpathquadraticcurveto{\pgfqpoint{9.407497in}{0.960685in}}{\pgfqpoint{9.409064in}{0.960397in}}%
\pgfpathlineto{\pgfqpoint{9.409118in}{0.949770in}}%
\pgfpathlineto{\pgfqpoint{9.409118in}{0.949770in}}%
\pgfpathclose%
\pgfusepath{fill}%
\end{pgfscope}%
\begin{pgfscope}%
\pgfsetbuttcap%
\pgfsetroundjoin%
\definecolor{currentfill}{rgb}{0.309804,0.372549,0.419608}%
\pgfsetfillcolor{currentfill}%
\pgfsetlinewidth{0.000000pt}%
\definecolor{currentstroke}{rgb}{0.309804,0.372549,0.419608}%
\pgfsetstrokecolor{currentstroke}%
\pgfsetdash{}{0pt}%
\pgfpathmoveto{\pgfqpoint{9.479552in}{0.905856in}}%
\pgfpathlineto{\pgfqpoint{9.479552in}{0.872426in}}%
\pgfpathlineto{\pgfqpoint{9.469141in}{0.872426in}}%
\pgfpathlineto{\pgfqpoint{9.469141in}{0.959443in}}%
\pgfpathlineto{\pgfqpoint{9.479552in}{0.959443in}}%
\pgfpathlineto{\pgfqpoint{9.479552in}{0.949878in}}%
\pgfpathquadraticcurveto{\pgfqpoint{9.482830in}{0.955498in}}{\pgfqpoint{9.487802in}{0.958218in}}%
\pgfpathquadraticcurveto{\pgfqpoint{9.492791in}{0.960956in}}{\pgfqpoint{9.499708in}{0.960956in}}%
\pgfpathquadraticcurveto{\pgfqpoint{9.511199in}{0.960956in}}{\pgfqpoint{9.518368in}{0.951841in}}%
\pgfpathquadraticcurveto{\pgfqpoint{9.525555in}{0.942727in}}{\pgfqpoint{9.525555in}{0.927867in}}%
\pgfpathquadraticcurveto{\pgfqpoint{9.525555in}{0.913007in}}{\pgfqpoint{9.518368in}{0.903875in}}%
\pgfpathquadraticcurveto{\pgfqpoint{9.511199in}{0.894761in}}{\pgfqpoint{9.499708in}{0.894761in}}%
\pgfpathquadraticcurveto{\pgfqpoint{9.492791in}{0.894761in}}{\pgfqpoint{9.487802in}{0.897499in}}%
\pgfpathquadraticcurveto{\pgfqpoint{9.482830in}{0.900237in}}{\pgfqpoint{9.479552in}{0.905856in}}%
\pgfpathlineto{\pgfqpoint{9.479552in}{0.905856in}}%
\pgfpathclose%
\pgfpathmoveto{\pgfqpoint{9.514802in}{0.927867in}}%
\pgfpathquadraticcurveto{\pgfqpoint{9.514802in}{0.939287in}}{\pgfqpoint{9.510101in}{0.945789in}}%
\pgfpathquadraticcurveto{\pgfqpoint{9.505399in}{0.952292in}}{\pgfqpoint{9.497186in}{0.952292in}}%
\pgfpathquadraticcurveto{\pgfqpoint{9.488954in}{0.952292in}}{\pgfqpoint{9.484253in}{0.945789in}}%
\pgfpathquadraticcurveto{\pgfqpoint{9.479552in}{0.939287in}}{\pgfqpoint{9.479552in}{0.927867in}}%
\pgfpathquadraticcurveto{\pgfqpoint{9.479552in}{0.916448in}}{\pgfqpoint{9.484253in}{0.909945in}}%
\pgfpathquadraticcurveto{\pgfqpoint{9.488954in}{0.903443in}}{\pgfqpoint{9.497186in}{0.903443in}}%
\pgfpathquadraticcurveto{\pgfqpoint{9.505399in}{0.903443in}}{\pgfqpoint{9.510101in}{0.909945in}}%
\pgfpathquadraticcurveto{\pgfqpoint{9.514802in}{0.916448in}}{\pgfqpoint{9.514802in}{0.927867in}}%
\pgfpathlineto{\pgfqpoint{9.514802in}{0.927867in}}%
\pgfpathclose%
\pgfpathmoveto{\pgfqpoint{9.579249in}{0.949770in}}%
\pgfpathquadraticcurveto{\pgfqpoint{9.577502in}{0.950779in}}{\pgfqpoint{9.575449in}{0.951247in}}%
\pgfpathquadraticcurveto{\pgfqpoint{9.573395in}{0.951733in}}{\pgfqpoint{9.570928in}{0.951733in}}%
\pgfpathquadraticcurveto{\pgfqpoint{9.562138in}{0.951733in}}{\pgfqpoint{9.557437in}{0.946023in}}%
\pgfpathquadraticcurveto{\pgfqpoint{9.552735in}{0.940314in}}{\pgfqpoint{9.552735in}{0.929614in}}%
\pgfpathlineto{\pgfqpoint{9.552735in}{0.896400in}}%
\pgfpathlineto{\pgfqpoint{9.542324in}{0.896400in}}%
\pgfpathlineto{\pgfqpoint{9.542324in}{0.959443in}}%
\pgfpathlineto{\pgfqpoint{9.552735in}{0.959443in}}%
\pgfpathlineto{\pgfqpoint{9.552735in}{0.949644in}}%
\pgfpathquadraticcurveto{\pgfqpoint{9.556014in}{0.955390in}}{\pgfqpoint{9.561237in}{0.958164in}}%
\pgfpathquadraticcurveto{\pgfqpoint{9.566479in}{0.960956in}}{\pgfqpoint{9.573972in}{0.960956in}}%
\pgfpathquadraticcurveto{\pgfqpoint{9.575034in}{0.960956in}}{\pgfqpoint{9.576331in}{0.960811in}}%
\pgfpathquadraticcurveto{\pgfqpoint{9.577628in}{0.960685in}}{\pgfqpoint{9.579195in}{0.960397in}}%
\pgfpathlineto{\pgfqpoint{9.579249in}{0.949770in}}%
\pgfpathlineto{\pgfqpoint{9.579249in}{0.949770in}}%
\pgfpathclose%
\pgfpathmoveto{\pgfqpoint{9.590111in}{0.959443in}}%
\pgfpathlineto{\pgfqpoint{9.600468in}{0.959443in}}%
\pgfpathlineto{\pgfqpoint{9.600468in}{0.896400in}}%
\pgfpathlineto{\pgfqpoint{9.590111in}{0.896400in}}%
\pgfpathlineto{\pgfqpoint{9.590111in}{0.959443in}}%
\pgfpathlineto{\pgfqpoint{9.590111in}{0.959443in}}%
\pgfpathclose%
\pgfpathmoveto{\pgfqpoint{9.590111in}{0.983993in}}%
\pgfpathlineto{\pgfqpoint{9.600468in}{0.983993in}}%
\pgfpathlineto{\pgfqpoint{9.600468in}{0.970862in}}%
\pgfpathlineto{\pgfqpoint{9.590111in}{0.970862in}}%
\pgfpathlineto{\pgfqpoint{9.590111in}{0.983993in}}%
\pgfpathlineto{\pgfqpoint{9.590111in}{0.983993in}}%
\pgfpathclose%
\pgfpathmoveto{\pgfqpoint{9.646561in}{0.952184in}}%
\pgfpathquadraticcurveto{\pgfqpoint{9.638239in}{0.952184in}}{\pgfqpoint{9.633394in}{0.945681in}}%
\pgfpathquadraticcurveto{\pgfqpoint{9.628549in}{0.939179in}}{\pgfqpoint{9.628549in}{0.927867in}}%
\pgfpathquadraticcurveto{\pgfqpoint{9.628549in}{0.916556in}}{\pgfqpoint{9.633358in}{0.910053in}}%
\pgfpathquadraticcurveto{\pgfqpoint{9.638185in}{0.903551in}}{\pgfqpoint{9.646561in}{0.903551in}}%
\pgfpathquadraticcurveto{\pgfqpoint{9.654846in}{0.903551in}}{\pgfqpoint{9.659674in}{0.910071in}}%
\pgfpathquadraticcurveto{\pgfqpoint{9.664519in}{0.916610in}}{\pgfqpoint{9.664519in}{0.927867in}}%
\pgfpathquadraticcurveto{\pgfqpoint{9.664519in}{0.939071in}}{\pgfqpoint{9.659674in}{0.945627in}}%
\pgfpathquadraticcurveto{\pgfqpoint{9.654846in}{0.952184in}}{\pgfqpoint{9.646561in}{0.952184in}}%
\pgfpathlineto{\pgfqpoint{9.646561in}{0.952184in}}%
\pgfpathclose%
\pgfpathmoveto{\pgfqpoint{9.646561in}{0.960956in}}%
\pgfpathquadraticcurveto{\pgfqpoint{9.660070in}{0.960956in}}{\pgfqpoint{9.667779in}{0.952166in}}%
\pgfpathquadraticcurveto{\pgfqpoint{9.675506in}{0.943394in}}{\pgfqpoint{9.675506in}{0.927867in}}%
\pgfpathquadraticcurveto{\pgfqpoint{9.675506in}{0.912395in}}{\pgfqpoint{9.667779in}{0.903569in}}%
\pgfpathquadraticcurveto{\pgfqpoint{9.660070in}{0.894761in}}{\pgfqpoint{9.646561in}{0.894761in}}%
\pgfpathquadraticcurveto{\pgfqpoint{9.632998in}{0.894761in}}{\pgfqpoint{9.625306in}{0.903569in}}%
\pgfpathquadraticcurveto{\pgfqpoint{9.617633in}{0.912395in}}{\pgfqpoint{9.617633in}{0.927867in}}%
\pgfpathquadraticcurveto{\pgfqpoint{9.617633in}{0.943394in}}{\pgfqpoint{9.625306in}{0.952166in}}%
\pgfpathquadraticcurveto{\pgfqpoint{9.632998in}{0.960956in}}{\pgfqpoint{9.646561in}{0.960956in}}%
\pgfpathlineto{\pgfqpoint{9.646561in}{0.960956in}}%
\pgfpathclose%
\pgfpathmoveto{\pgfqpoint{9.729200in}{0.949770in}}%
\pgfpathquadraticcurveto{\pgfqpoint{9.727453in}{0.950779in}}{\pgfqpoint{9.725400in}{0.951247in}}%
\pgfpathquadraticcurveto{\pgfqpoint{9.723346in}{0.951733in}}{\pgfqpoint{9.720879in}{0.951733in}}%
\pgfpathquadraticcurveto{\pgfqpoint{9.712089in}{0.951733in}}{\pgfqpoint{9.707388in}{0.946023in}}%
\pgfpathquadraticcurveto{\pgfqpoint{9.702687in}{0.940314in}}{\pgfqpoint{9.702687in}{0.929614in}}%
\pgfpathlineto{\pgfqpoint{9.702687in}{0.896400in}}%
\pgfpathlineto{\pgfqpoint{9.692276in}{0.896400in}}%
\pgfpathlineto{\pgfqpoint{9.692276in}{0.959443in}}%
\pgfpathlineto{\pgfqpoint{9.702687in}{0.959443in}}%
\pgfpathlineto{\pgfqpoint{9.702687in}{0.949644in}}%
\pgfpathquadraticcurveto{\pgfqpoint{9.705965in}{0.955390in}}{\pgfqpoint{9.711188in}{0.958164in}}%
\pgfpathquadraticcurveto{\pgfqpoint{9.716430in}{0.960956in}}{\pgfqpoint{9.723923in}{0.960956in}}%
\pgfpathquadraticcurveto{\pgfqpoint{9.724986in}{0.960956in}}{\pgfqpoint{9.726282in}{0.960811in}}%
\pgfpathquadraticcurveto{\pgfqpoint{9.727579in}{0.960685in}}{\pgfqpoint{9.729146in}{0.960397in}}%
\pgfpathlineto{\pgfqpoint{9.729200in}{0.949770in}}%
\pgfpathlineto{\pgfqpoint{9.729200in}{0.949770in}}%
\pgfpathclose%
\pgfusepath{fill}%
\end{pgfscope}%
\begin{pgfscope}%
\pgfsetbuttcap%
\pgfsetroundjoin%
\definecolor{currentfill}{rgb}{0.309804,0.372549,0.419608}%
\pgfsetfillcolor{currentfill}%
\pgfsetlinewidth{0.000000pt}%
\definecolor{currentstroke}{rgb}{0.309804,0.372549,0.419608}%
\pgfsetstrokecolor{currentstroke}%
\pgfsetdash{}{0pt}%
\pgfpathmoveto{\pgfqpoint{9.797021in}{0.959443in}}%
\pgfpathlineto{\pgfqpoint{9.807378in}{0.959443in}}%
\pgfpathlineto{\pgfqpoint{9.807378in}{0.896400in}}%
\pgfpathlineto{\pgfqpoint{9.797021in}{0.896400in}}%
\pgfpathlineto{\pgfqpoint{9.797021in}{0.959443in}}%
\pgfpathlineto{\pgfqpoint{9.797021in}{0.959443in}}%
\pgfpathclose%
\pgfpathmoveto{\pgfqpoint{9.797021in}{0.983993in}}%
\pgfpathlineto{\pgfqpoint{9.807378in}{0.983993in}}%
\pgfpathlineto{\pgfqpoint{9.807378in}{0.970862in}}%
\pgfpathlineto{\pgfqpoint{9.797021in}{0.970862in}}%
\pgfpathlineto{\pgfqpoint{9.797021in}{0.983993in}}%
\pgfpathlineto{\pgfqpoint{9.797021in}{0.983993in}}%
\pgfpathclose%
\pgfpathmoveto{\pgfqpoint{9.881462in}{0.934460in}}%
\pgfpathlineto{\pgfqpoint{9.881462in}{0.896400in}}%
\pgfpathlineto{\pgfqpoint{9.871105in}{0.896400in}}%
\pgfpathlineto{\pgfqpoint{9.871105in}{0.934117in}}%
\pgfpathquadraticcurveto{\pgfqpoint{9.871105in}{0.943069in}}{\pgfqpoint{9.867611in}{0.947500in}}%
\pgfpathquadraticcurveto{\pgfqpoint{9.864116in}{0.951949in}}{\pgfqpoint{9.857146in}{0.951949in}}%
\pgfpathquadraticcurveto{\pgfqpoint{9.848752in}{0.951949in}}{\pgfqpoint{9.843907in}{0.946600in}}%
\pgfpathquadraticcurveto{\pgfqpoint{9.839061in}{0.941268in}}{\pgfqpoint{9.839061in}{0.932028in}}%
\pgfpathlineto{\pgfqpoint{9.839061in}{0.896400in}}%
\pgfpathlineto{\pgfqpoint{9.828650in}{0.896400in}}%
\pgfpathlineto{\pgfqpoint{9.828650in}{0.959443in}}%
\pgfpathlineto{\pgfqpoint{9.839061in}{0.959443in}}%
\pgfpathlineto{\pgfqpoint{9.839061in}{0.949644in}}%
\pgfpathquadraticcurveto{\pgfqpoint{9.842790in}{0.955336in}}{\pgfqpoint{9.847815in}{0.958146in}}%
\pgfpathquadraticcurveto{\pgfqpoint{9.852859in}{0.960956in}}{\pgfqpoint{9.859451in}{0.960956in}}%
\pgfpathquadraticcurveto{\pgfqpoint{9.870313in}{0.960956in}}{\pgfqpoint{9.875878in}{0.954237in}}%
\pgfpathquadraticcurveto{\pgfqpoint{9.881462in}{0.947518in}}{\pgfqpoint{9.881462in}{0.934460in}}%
\pgfpathlineto{\pgfqpoint{9.881462in}{0.934460in}}%
\pgfpathclose%
\pgfpathmoveto{\pgfqpoint{9.934021in}{0.983993in}}%
\pgfpathlineto{\pgfqpoint{9.934021in}{0.975365in}}%
\pgfpathlineto{\pgfqpoint{9.924115in}{0.975365in}}%
\pgfpathquadraticcurveto{\pgfqpoint{9.918549in}{0.975365in}}{\pgfqpoint{9.916370in}{0.973114in}}%
\pgfpathquadraticcurveto{\pgfqpoint{9.914208in}{0.970862in}}{\pgfqpoint{9.914208in}{0.965008in}}%
\pgfpathlineto{\pgfqpoint{9.914208in}{0.959443in}}%
\pgfpathlineto{\pgfqpoint{9.931266in}{0.959443in}}%
\pgfpathlineto{\pgfqpoint{9.931266in}{0.951391in}}%
\pgfpathlineto{\pgfqpoint{9.914208in}{0.951391in}}%
\pgfpathlineto{\pgfqpoint{9.914208in}{0.896400in}}%
\pgfpathlineto{\pgfqpoint{9.903797in}{0.896400in}}%
\pgfpathlineto{\pgfqpoint{9.903797in}{0.951391in}}%
\pgfpathlineto{\pgfqpoint{9.893890in}{0.951391in}}%
\pgfpathlineto{\pgfqpoint{9.893890in}{0.959443in}}%
\pgfpathlineto{\pgfqpoint{9.903797in}{0.959443in}}%
\pgfpathlineto{\pgfqpoint{9.903797in}{0.963838in}}%
\pgfpathquadraticcurveto{\pgfqpoint{9.903797in}{0.974357in}}{\pgfqpoint{9.908696in}{0.979166in}}%
\pgfpathquadraticcurveto{\pgfqpoint{9.913596in}{0.983993in}}{\pgfqpoint{9.924223in}{0.983993in}}%
\pgfpathlineto{\pgfqpoint{9.934021in}{0.983993in}}%
\pgfpathlineto{\pgfqpoint{9.934021in}{0.983993in}}%
\pgfpathclose%
\pgfpathmoveto{\pgfqpoint{9.996614in}{0.930515in}}%
\pgfpathlineto{\pgfqpoint{9.996614in}{0.925454in}}%
\pgfpathlineto{\pgfqpoint{9.948990in}{0.925454in}}%
\pgfpathquadraticcurveto{\pgfqpoint{9.949674in}{0.914754in}}{\pgfqpoint{9.955438in}{0.909153in}}%
\pgfpathquadraticcurveto{\pgfqpoint{9.961202in}{0.903551in}}{\pgfqpoint{9.971505in}{0.903551in}}%
\pgfpathquadraticcurveto{\pgfqpoint{9.977467in}{0.903551in}}{\pgfqpoint{9.983069in}{0.905010in}}%
\pgfpathquadraticcurveto{\pgfqpoint{9.988670in}{0.906469in}}{\pgfqpoint{9.994200in}{0.909405in}}%
\pgfpathlineto{\pgfqpoint{9.994200in}{0.899606in}}%
\pgfpathquadraticcurveto{\pgfqpoint{9.988616in}{0.897247in}}{\pgfqpoint{9.982762in}{0.896004in}}%
\pgfpathquadraticcurveto{\pgfqpoint{9.976908in}{0.894761in}}{\pgfqpoint{9.970892in}{0.894761in}}%
\pgfpathquadraticcurveto{\pgfqpoint{9.955798in}{0.894761in}}{\pgfqpoint{9.946990in}{0.903533in}}%
\pgfpathquadraticcurveto{\pgfqpoint{9.938182in}{0.912323in}}{\pgfqpoint{9.938182in}{0.927309in}}%
\pgfpathquadraticcurveto{\pgfqpoint{9.938182in}{0.942781in}}{\pgfqpoint{9.946540in}{0.951859in}}%
\pgfpathquadraticcurveto{\pgfqpoint{9.954898in}{0.960956in}}{\pgfqpoint{9.969091in}{0.960956in}}%
\pgfpathquadraticcurveto{\pgfqpoint{9.981808in}{0.960956in}}{\pgfqpoint{9.989211in}{0.952760in}}%
\pgfpathquadraticcurveto{\pgfqpoint{9.996614in}{0.944583in}}{\pgfqpoint{9.996614in}{0.930515in}}%
\pgfpathlineto{\pgfqpoint{9.996614in}{0.930515in}}%
\pgfpathclose%
\pgfpathmoveto{\pgfqpoint{9.986257in}{0.933559in}}%
\pgfpathquadraticcurveto{\pgfqpoint{9.986149in}{0.942043in}}{\pgfqpoint{9.981502in}{0.947104in}}%
\pgfpathquadraticcurveto{\pgfqpoint{9.976854in}{0.952184in}}{\pgfqpoint{9.969199in}{0.952184in}}%
\pgfpathquadraticcurveto{\pgfqpoint{9.960535in}{0.952184in}}{\pgfqpoint{9.955330in}{0.947284in}}%
\pgfpathquadraticcurveto{\pgfqpoint{9.950124in}{0.942385in}}{\pgfqpoint{9.949332in}{0.933487in}}%
\pgfpathlineto{\pgfqpoint{9.986257in}{0.933559in}}%
\pgfpathlineto{\pgfqpoint{9.986257in}{0.933559in}}%
\pgfpathclose%
\pgfpathmoveto{\pgfqpoint{10.050146in}{0.949770in}}%
\pgfpathquadraticcurveto{\pgfqpoint{10.048399in}{0.950779in}}{\pgfqpoint{10.046345in}{0.951247in}}%
\pgfpathquadraticcurveto{\pgfqpoint{10.044292in}{0.951733in}}{\pgfqpoint{10.041824in}{0.951733in}}%
\pgfpathquadraticcurveto{\pgfqpoint{10.033034in}{0.951733in}}{\pgfqpoint{10.028333in}{0.946023in}}%
\pgfpathquadraticcurveto{\pgfqpoint{10.023632in}{0.940314in}}{\pgfqpoint{10.023632in}{0.929614in}}%
\pgfpathlineto{\pgfqpoint{10.023632in}{0.896400in}}%
\pgfpathlineto{\pgfqpoint{10.013221in}{0.896400in}}%
\pgfpathlineto{\pgfqpoint{10.013221in}{0.959443in}}%
\pgfpathlineto{\pgfqpoint{10.023632in}{0.959443in}}%
\pgfpathlineto{\pgfqpoint{10.023632in}{0.949644in}}%
\pgfpathquadraticcurveto{\pgfqpoint{10.026910in}{0.955390in}}{\pgfqpoint{10.032134in}{0.958164in}}%
\pgfpathquadraticcurveto{\pgfqpoint{10.037375in}{0.960956in}}{\pgfqpoint{10.044868in}{0.960956in}}%
\pgfpathquadraticcurveto{\pgfqpoint{10.045931in}{0.960956in}}{\pgfqpoint{10.047228in}{0.960811in}}%
\pgfpathquadraticcurveto{\pgfqpoint{10.048525in}{0.960685in}}{\pgfqpoint{10.050092in}{0.960397in}}%
\pgfpathlineto{\pgfqpoint{10.050146in}{0.949770in}}%
\pgfpathlineto{\pgfqpoint{10.050146in}{0.949770in}}%
\pgfpathclose%
\pgfpathmoveto{\pgfqpoint{10.112396in}{0.930515in}}%
\pgfpathlineto{\pgfqpoint{10.112396in}{0.925454in}}%
\pgfpathlineto{\pgfqpoint{10.064772in}{0.925454in}}%
\pgfpathquadraticcurveto{\pgfqpoint{10.065456in}{0.914754in}}{\pgfqpoint{10.071220in}{0.909153in}}%
\pgfpathquadraticcurveto{\pgfqpoint{10.076984in}{0.903551in}}{\pgfqpoint{10.087287in}{0.903551in}}%
\pgfpathquadraticcurveto{\pgfqpoint{10.093249in}{0.903551in}}{\pgfqpoint{10.098851in}{0.905010in}}%
\pgfpathquadraticcurveto{\pgfqpoint{10.104452in}{0.906469in}}{\pgfqpoint{10.109982in}{0.909405in}}%
\pgfpathlineto{\pgfqpoint{10.109982in}{0.899606in}}%
\pgfpathquadraticcurveto{\pgfqpoint{10.104398in}{0.897247in}}{\pgfqpoint{10.098544in}{0.896004in}}%
\pgfpathquadraticcurveto{\pgfqpoint{10.092691in}{0.894761in}}{\pgfqpoint{10.086674in}{0.894761in}}%
\pgfpathquadraticcurveto{\pgfqpoint{10.071580in}{0.894761in}}{\pgfqpoint{10.062772in}{0.903533in}}%
\pgfpathquadraticcurveto{\pgfqpoint{10.053964in}{0.912323in}}{\pgfqpoint{10.053964in}{0.927309in}}%
\pgfpathquadraticcurveto{\pgfqpoint{10.053964in}{0.942781in}}{\pgfqpoint{10.062322in}{0.951859in}}%
\pgfpathquadraticcurveto{\pgfqpoint{10.070680in}{0.960956in}}{\pgfqpoint{10.084873in}{0.960956in}}%
\pgfpathquadraticcurveto{\pgfqpoint{10.097590in}{0.960956in}}{\pgfqpoint{10.104993in}{0.952760in}}%
\pgfpathquadraticcurveto{\pgfqpoint{10.112396in}{0.944583in}}{\pgfqpoint{10.112396in}{0.930515in}}%
\pgfpathlineto{\pgfqpoint{10.112396in}{0.930515in}}%
\pgfpathclose%
\pgfpathmoveto{\pgfqpoint{10.102039in}{0.933559in}}%
\pgfpathquadraticcurveto{\pgfqpoint{10.101931in}{0.942043in}}{\pgfqpoint{10.097284in}{0.947104in}}%
\pgfpathquadraticcurveto{\pgfqpoint{10.092637in}{0.952184in}}{\pgfqpoint{10.084981in}{0.952184in}}%
\pgfpathquadraticcurveto{\pgfqpoint{10.076317in}{0.952184in}}{\pgfqpoint{10.071112in}{0.947284in}}%
\pgfpathquadraticcurveto{\pgfqpoint{10.065906in}{0.942385in}}{\pgfqpoint{10.065114in}{0.933487in}}%
\pgfpathlineto{\pgfqpoint{10.102039in}{0.933559in}}%
\pgfpathlineto{\pgfqpoint{10.102039in}{0.933559in}}%
\pgfpathclose%
\pgfpathmoveto{\pgfqpoint{10.181815in}{0.934460in}}%
\pgfpathlineto{\pgfqpoint{10.181815in}{0.896400in}}%
\pgfpathlineto{\pgfqpoint{10.171458in}{0.896400in}}%
\pgfpathlineto{\pgfqpoint{10.171458in}{0.934117in}}%
\pgfpathquadraticcurveto{\pgfqpoint{10.171458in}{0.943069in}}{\pgfqpoint{10.167963in}{0.947500in}}%
\pgfpathquadraticcurveto{\pgfqpoint{10.164469in}{0.951949in}}{\pgfqpoint{10.157498in}{0.951949in}}%
\pgfpathquadraticcurveto{\pgfqpoint{10.149105in}{0.951949in}}{\pgfqpoint{10.144259in}{0.946600in}}%
\pgfpathquadraticcurveto{\pgfqpoint{10.139414in}{0.941268in}}{\pgfqpoint{10.139414in}{0.932028in}}%
\pgfpathlineto{\pgfqpoint{10.139414in}{0.896400in}}%
\pgfpathlineto{\pgfqpoint{10.129003in}{0.896400in}}%
\pgfpathlineto{\pgfqpoint{10.129003in}{0.959443in}}%
\pgfpathlineto{\pgfqpoint{10.139414in}{0.959443in}}%
\pgfpathlineto{\pgfqpoint{10.139414in}{0.949644in}}%
\pgfpathquadraticcurveto{\pgfqpoint{10.143143in}{0.955336in}}{\pgfqpoint{10.148168in}{0.958146in}}%
\pgfpathquadraticcurveto{\pgfqpoint{10.153211in}{0.960956in}}{\pgfqpoint{10.159804in}{0.960956in}}%
\pgfpathquadraticcurveto{\pgfqpoint{10.170665in}{0.960956in}}{\pgfqpoint{10.176231in}{0.954237in}}%
\pgfpathquadraticcurveto{\pgfqpoint{10.181815in}{0.947518in}}{\pgfqpoint{10.181815in}{0.934460in}}%
\pgfpathlineto{\pgfqpoint{10.181815in}{0.934460in}}%
\pgfpathclose%
\pgfpathmoveto{\pgfqpoint{10.247829in}{0.957029in}}%
\pgfpathlineto{\pgfqpoint{10.247829in}{0.947338in}}%
\pgfpathquadraticcurveto{\pgfqpoint{10.243434in}{0.949770in}}{\pgfqpoint{10.239021in}{0.950977in}}%
\pgfpathquadraticcurveto{\pgfqpoint{10.234608in}{0.952184in}}{\pgfqpoint{10.230105in}{0.952184in}}%
\pgfpathquadraticcurveto{\pgfqpoint{10.220018in}{0.952184in}}{\pgfqpoint{10.214435in}{0.945789in}}%
\pgfpathquadraticcurveto{\pgfqpoint{10.208869in}{0.939413in}}{\pgfqpoint{10.208869in}{0.927867in}}%
\pgfpathquadraticcurveto{\pgfqpoint{10.208869in}{0.916321in}}{\pgfqpoint{10.214435in}{0.909927in}}%
\pgfpathquadraticcurveto{\pgfqpoint{10.220018in}{0.903551in}}{\pgfqpoint{10.230105in}{0.903551in}}%
\pgfpathquadraticcurveto{\pgfqpoint{10.234608in}{0.903551in}}{\pgfqpoint{10.239021in}{0.904758in}}%
\pgfpathquadraticcurveto{\pgfqpoint{10.243434in}{0.905964in}}{\pgfqpoint{10.247829in}{0.908396in}}%
\pgfpathlineto{\pgfqpoint{10.247829in}{0.898814in}}%
\pgfpathquadraticcurveto{\pgfqpoint{10.243488in}{0.896796in}}{\pgfqpoint{10.238841in}{0.895788in}}%
\pgfpathquadraticcurveto{\pgfqpoint{10.234212in}{0.894761in}}{\pgfqpoint{10.228970in}{0.894761in}}%
\pgfpathquadraticcurveto{\pgfqpoint{10.214723in}{0.894761in}}{\pgfqpoint{10.206329in}{0.903713in}}%
\pgfpathquadraticcurveto{\pgfqpoint{10.197954in}{0.912665in}}{\pgfqpoint{10.197954in}{0.927867in}}%
\pgfpathquadraticcurveto{\pgfqpoint{10.197954in}{0.943286in}}{\pgfqpoint{10.206419in}{0.952112in}}%
\pgfpathquadraticcurveto{\pgfqpoint{10.214903in}{0.960956in}}{\pgfqpoint{10.229655in}{0.960956in}}%
\pgfpathquadraticcurveto{\pgfqpoint{10.234428in}{0.960956in}}{\pgfqpoint{10.238985in}{0.959965in}}%
\pgfpathquadraticcurveto{\pgfqpoint{10.243542in}{0.958992in}}{\pgfqpoint{10.247829in}{0.957029in}}%
\pgfpathlineto{\pgfqpoint{10.247829in}{0.957029in}}%
\pgfpathclose%
\pgfpathmoveto{\pgfqpoint{10.319770in}{0.930515in}}%
\pgfpathlineto{\pgfqpoint{10.319770in}{0.925454in}}%
\pgfpathlineto{\pgfqpoint{10.272146in}{0.925454in}}%
\pgfpathquadraticcurveto{\pgfqpoint{10.272830in}{0.914754in}}{\pgfqpoint{10.278594in}{0.909153in}}%
\pgfpathquadraticcurveto{\pgfqpoint{10.284358in}{0.903551in}}{\pgfqpoint{10.294661in}{0.903551in}}%
\pgfpathquadraticcurveto{\pgfqpoint{10.300623in}{0.903551in}}{\pgfqpoint{10.306225in}{0.905010in}}%
\pgfpathquadraticcurveto{\pgfqpoint{10.311826in}{0.906469in}}{\pgfqpoint{10.317356in}{0.909405in}}%
\pgfpathlineto{\pgfqpoint{10.317356in}{0.899606in}}%
\pgfpathquadraticcurveto{\pgfqpoint{10.311772in}{0.897247in}}{\pgfqpoint{10.305918in}{0.896004in}}%
\pgfpathquadraticcurveto{\pgfqpoint{10.300064in}{0.894761in}}{\pgfqpoint{10.294048in}{0.894761in}}%
\pgfpathquadraticcurveto{\pgfqpoint{10.278954in}{0.894761in}}{\pgfqpoint{10.270146in}{0.903533in}}%
\pgfpathquadraticcurveto{\pgfqpoint{10.261338in}{0.912323in}}{\pgfqpoint{10.261338in}{0.927309in}}%
\pgfpathquadraticcurveto{\pgfqpoint{10.261338in}{0.942781in}}{\pgfqpoint{10.269696in}{0.951859in}}%
\pgfpathquadraticcurveto{\pgfqpoint{10.278054in}{0.960956in}}{\pgfqpoint{10.292247in}{0.960956in}}%
\pgfpathquadraticcurveto{\pgfqpoint{10.304964in}{0.960956in}}{\pgfqpoint{10.312367in}{0.952760in}}%
\pgfpathquadraticcurveto{\pgfqpoint{10.319770in}{0.944583in}}{\pgfqpoint{10.319770in}{0.930515in}}%
\pgfpathlineto{\pgfqpoint{10.319770in}{0.930515in}}%
\pgfpathclose%
\pgfpathmoveto{\pgfqpoint{10.309413in}{0.933559in}}%
\pgfpathquadraticcurveto{\pgfqpoint{10.309305in}{0.942043in}}{\pgfqpoint{10.304658in}{0.947104in}}%
\pgfpathquadraticcurveto{\pgfqpoint{10.300010in}{0.952184in}}{\pgfqpoint{10.292355in}{0.952184in}}%
\pgfpathquadraticcurveto{\pgfqpoint{10.283691in}{0.952184in}}{\pgfqpoint{10.278486in}{0.947284in}}%
\pgfpathquadraticcurveto{\pgfqpoint{10.273280in}{0.942385in}}{\pgfqpoint{10.272488in}{0.933487in}}%
\pgfpathlineto{\pgfqpoint{10.309413in}{0.933559in}}%
\pgfpathlineto{\pgfqpoint{10.309413in}{0.933559in}}%
\pgfpathclose%
\pgfusepath{fill}%
\end{pgfscope}%
\begin{pgfscope}%
\pgfsetbuttcap%
\pgfsetroundjoin%
\definecolor{currentfill}{rgb}{0.309804,0.372549,0.419608}%
\pgfsetfillcolor{currentfill}%
\pgfsetlinewidth{0.000000pt}%
\definecolor{currentstroke}{rgb}{0.309804,0.372549,0.419608}%
\pgfsetstrokecolor{currentstroke}%
\pgfsetdash{}{0pt}%
\pgfpathmoveto{\pgfqpoint{10.395989in}{0.959443in}}%
\pgfpathlineto{\pgfqpoint{10.406346in}{0.959443in}}%
\pgfpathlineto{\pgfqpoint{10.419297in}{0.910251in}}%
\pgfpathlineto{\pgfqpoint{10.432175in}{0.959443in}}%
\pgfpathlineto{\pgfqpoint{10.444387in}{0.959443in}}%
\pgfpathlineto{\pgfqpoint{10.457338in}{0.910251in}}%
\pgfpathlineto{\pgfqpoint{10.470235in}{0.959443in}}%
\pgfpathlineto{\pgfqpoint{10.480592in}{0.959443in}}%
\pgfpathlineto{\pgfqpoint{10.464093in}{0.896400in}}%
\pgfpathlineto{\pgfqpoint{10.451880in}{0.896400in}}%
\pgfpathlineto{\pgfqpoint{10.438317in}{0.948077in}}%
\pgfpathlineto{\pgfqpoint{10.424700in}{0.896400in}}%
\pgfpathlineto{\pgfqpoint{10.412470in}{0.896400in}}%
\pgfpathlineto{\pgfqpoint{10.395989in}{0.959443in}}%
\pgfpathlineto{\pgfqpoint{10.395989in}{0.959443in}}%
\pgfpathclose%
\pgfpathmoveto{\pgfqpoint{10.524938in}{0.928083in}}%
\pgfpathquadraticcurveto{\pgfqpoint{10.512383in}{0.928083in}}{\pgfqpoint{10.507538in}{0.925219in}}%
\pgfpathquadraticcurveto{\pgfqpoint{10.502693in}{0.922356in}}{\pgfqpoint{10.502693in}{0.915421in}}%
\pgfpathquadraticcurveto{\pgfqpoint{10.502693in}{0.909909in}}{\pgfqpoint{10.506331in}{0.906667in}}%
\pgfpathquadraticcurveto{\pgfqpoint{10.509970in}{0.903443in}}{\pgfqpoint{10.516202in}{0.903443in}}%
\pgfpathquadraticcurveto{\pgfqpoint{10.524830in}{0.903443in}}{\pgfqpoint{10.530035in}{0.909549in}}%
\pgfpathquadraticcurveto{\pgfqpoint{10.535241in}{0.915655in}}{\pgfqpoint{10.535241in}{0.925778in}}%
\pgfpathlineto{\pgfqpoint{10.535241in}{0.928083in}}%
\pgfpathlineto{\pgfqpoint{10.524938in}{0.928083in}}%
\pgfpathlineto{\pgfqpoint{10.524938in}{0.928083in}}%
\pgfpathclose%
\pgfpathmoveto{\pgfqpoint{10.545598in}{0.932370in}}%
\pgfpathlineto{\pgfqpoint{10.545598in}{0.896400in}}%
\pgfpathlineto{\pgfqpoint{10.535241in}{0.896400in}}%
\pgfpathlineto{\pgfqpoint{10.535241in}{0.905964in}}%
\pgfpathquadraticcurveto{\pgfqpoint{10.531692in}{0.900237in}}{\pgfqpoint{10.526397in}{0.897499in}}%
\pgfpathquadraticcurveto{\pgfqpoint{10.521101in}{0.894761in}}{\pgfqpoint{10.513446in}{0.894761in}}%
\pgfpathquadraticcurveto{\pgfqpoint{10.503774in}{0.894761in}}{\pgfqpoint{10.498046in}{0.900201in}}%
\pgfpathquadraticcurveto{\pgfqpoint{10.492336in}{0.905640in}}{\pgfqpoint{10.492336in}{0.914754in}}%
\pgfpathquadraticcurveto{\pgfqpoint{10.492336in}{0.925382in}}{\pgfqpoint{10.499451in}{0.930785in}}%
\pgfpathquadraticcurveto{\pgfqpoint{10.506583in}{0.936189in}}{\pgfqpoint{10.520705in}{0.936189in}}%
\pgfpathlineto{\pgfqpoint{10.535241in}{0.936189in}}%
\pgfpathlineto{\pgfqpoint{10.535241in}{0.937216in}}%
\pgfpathquadraticcurveto{\pgfqpoint{10.535241in}{0.944366in}}{\pgfqpoint{10.530540in}{0.948275in}}%
\pgfpathquadraticcurveto{\pgfqpoint{10.525838in}{0.952184in}}{\pgfqpoint{10.517337in}{0.952184in}}%
\pgfpathquadraticcurveto{\pgfqpoint{10.511933in}{0.952184in}}{\pgfqpoint{10.506800in}{0.950887in}}%
\pgfpathquadraticcurveto{\pgfqpoint{10.501684in}{0.949590in}}{\pgfqpoint{10.496965in}{0.946996in}}%
\pgfpathlineto{\pgfqpoint{10.496965in}{0.956579in}}%
\pgfpathquadraticcurveto{\pgfqpoint{10.502639in}{0.958776in}}{\pgfqpoint{10.507988in}{0.959857in}}%
\pgfpathquadraticcurveto{\pgfqpoint{10.513338in}{0.960956in}}{\pgfqpoint{10.518399in}{0.960956in}}%
\pgfpathquadraticcurveto{\pgfqpoint{10.532089in}{0.960956in}}{\pgfqpoint{10.538843in}{0.953859in}}%
\pgfpathquadraticcurveto{\pgfqpoint{10.545598in}{0.946780in}}{\pgfqpoint{10.545598in}{0.932370in}}%
\pgfpathlineto{\pgfqpoint{10.545598in}{0.932370in}}%
\pgfpathclose%
\pgfpathmoveto{\pgfqpoint{10.603453in}{0.949770in}}%
\pgfpathquadraticcurveto{\pgfqpoint{10.601706in}{0.950779in}}{\pgfqpoint{10.599652in}{0.951247in}}%
\pgfpathquadraticcurveto{\pgfqpoint{10.597599in}{0.951733in}}{\pgfqpoint{10.595131in}{0.951733in}}%
\pgfpathquadraticcurveto{\pgfqpoint{10.586341in}{0.951733in}}{\pgfqpoint{10.581640in}{0.946023in}}%
\pgfpathquadraticcurveto{\pgfqpoint{10.576939in}{0.940314in}}{\pgfqpoint{10.576939in}{0.929614in}}%
\pgfpathlineto{\pgfqpoint{10.576939in}{0.896400in}}%
\pgfpathlineto{\pgfqpoint{10.566528in}{0.896400in}}%
\pgfpathlineto{\pgfqpoint{10.566528in}{0.959443in}}%
\pgfpathlineto{\pgfqpoint{10.576939in}{0.959443in}}%
\pgfpathlineto{\pgfqpoint{10.576939in}{0.949644in}}%
\pgfpathquadraticcurveto{\pgfqpoint{10.580217in}{0.955390in}}{\pgfqpoint{10.585441in}{0.958164in}}%
\pgfpathquadraticcurveto{\pgfqpoint{10.590682in}{0.960956in}}{\pgfqpoint{10.598175in}{0.960956in}}%
\pgfpathquadraticcurveto{\pgfqpoint{10.599238in}{0.960956in}}{\pgfqpoint{10.600535in}{0.960811in}}%
\pgfpathquadraticcurveto{\pgfqpoint{10.601832in}{0.960685in}}{\pgfqpoint{10.603399in}{0.960397in}}%
\pgfpathlineto{\pgfqpoint{10.603453in}{0.949770in}}%
\pgfpathlineto{\pgfqpoint{10.603453in}{0.949770in}}%
\pgfpathclose%
\pgfpathmoveto{\pgfqpoint{10.664712in}{0.934460in}}%
\pgfpathlineto{\pgfqpoint{10.664712in}{0.896400in}}%
\pgfpathlineto{\pgfqpoint{10.654355in}{0.896400in}}%
\pgfpathlineto{\pgfqpoint{10.654355in}{0.934117in}}%
\pgfpathquadraticcurveto{\pgfqpoint{10.654355in}{0.943069in}}{\pgfqpoint{10.650861in}{0.947500in}}%
\pgfpathquadraticcurveto{\pgfqpoint{10.647366in}{0.951949in}}{\pgfqpoint{10.640396in}{0.951949in}}%
\pgfpathquadraticcurveto{\pgfqpoint{10.632002in}{0.951949in}}{\pgfqpoint{10.627157in}{0.946600in}}%
\pgfpathquadraticcurveto{\pgfqpoint{10.622311in}{0.941268in}}{\pgfqpoint{10.622311in}{0.932028in}}%
\pgfpathlineto{\pgfqpoint{10.622311in}{0.896400in}}%
\pgfpathlineto{\pgfqpoint{10.611900in}{0.896400in}}%
\pgfpathlineto{\pgfqpoint{10.611900in}{0.959443in}}%
\pgfpathlineto{\pgfqpoint{10.622311in}{0.959443in}}%
\pgfpathlineto{\pgfqpoint{10.622311in}{0.949644in}}%
\pgfpathquadraticcurveto{\pgfqpoint{10.626040in}{0.955336in}}{\pgfqpoint{10.631065in}{0.958146in}}%
\pgfpathquadraticcurveto{\pgfqpoint{10.636109in}{0.960956in}}{\pgfqpoint{10.642701in}{0.960956in}}%
\pgfpathquadraticcurveto{\pgfqpoint{10.653563in}{0.960956in}}{\pgfqpoint{10.659128in}{0.954237in}}%
\pgfpathquadraticcurveto{\pgfqpoint{10.664712in}{0.947518in}}{\pgfqpoint{10.664712in}{0.934460in}}%
\pgfpathlineto{\pgfqpoint{10.664712in}{0.934460in}}%
\pgfpathclose%
\pgfpathmoveto{\pgfqpoint{10.685354in}{0.959443in}}%
\pgfpathlineto{\pgfqpoint{10.695711in}{0.959443in}}%
\pgfpathlineto{\pgfqpoint{10.695711in}{0.896400in}}%
\pgfpathlineto{\pgfqpoint{10.685354in}{0.896400in}}%
\pgfpathlineto{\pgfqpoint{10.685354in}{0.959443in}}%
\pgfpathlineto{\pgfqpoint{10.685354in}{0.959443in}}%
\pgfpathclose%
\pgfpathmoveto{\pgfqpoint{10.685354in}{0.983993in}}%
\pgfpathlineto{\pgfqpoint{10.695711in}{0.983993in}}%
\pgfpathlineto{\pgfqpoint{10.695711in}{0.970862in}}%
\pgfpathlineto{\pgfqpoint{10.685354in}{0.970862in}}%
\pgfpathlineto{\pgfqpoint{10.685354in}{0.983993in}}%
\pgfpathlineto{\pgfqpoint{10.685354in}{0.983993in}}%
\pgfpathclose%
\pgfpathmoveto{\pgfqpoint{10.769795in}{0.934460in}}%
\pgfpathlineto{\pgfqpoint{10.769795in}{0.896400in}}%
\pgfpathlineto{\pgfqpoint{10.759438in}{0.896400in}}%
\pgfpathlineto{\pgfqpoint{10.759438in}{0.934117in}}%
\pgfpathquadraticcurveto{\pgfqpoint{10.759438in}{0.943069in}}{\pgfqpoint{10.755944in}{0.947500in}}%
\pgfpathquadraticcurveto{\pgfqpoint{10.752449in}{0.951949in}}{\pgfqpoint{10.745479in}{0.951949in}}%
\pgfpathquadraticcurveto{\pgfqpoint{10.737085in}{0.951949in}}{\pgfqpoint{10.732240in}{0.946600in}}%
\pgfpathquadraticcurveto{\pgfqpoint{10.727394in}{0.941268in}}{\pgfqpoint{10.727394in}{0.932028in}}%
\pgfpathlineto{\pgfqpoint{10.727394in}{0.896400in}}%
\pgfpathlineto{\pgfqpoint{10.716983in}{0.896400in}}%
\pgfpathlineto{\pgfqpoint{10.716983in}{0.959443in}}%
\pgfpathlineto{\pgfqpoint{10.727394in}{0.959443in}}%
\pgfpathlineto{\pgfqpoint{10.727394in}{0.949644in}}%
\pgfpathquadraticcurveto{\pgfqpoint{10.731123in}{0.955336in}}{\pgfqpoint{10.736148in}{0.958146in}}%
\pgfpathquadraticcurveto{\pgfqpoint{10.741192in}{0.960956in}}{\pgfqpoint{10.747784in}{0.960956in}}%
\pgfpathquadraticcurveto{\pgfqpoint{10.758645in}{0.960956in}}{\pgfqpoint{10.764211in}{0.954237in}}%
\pgfpathquadraticcurveto{\pgfqpoint{10.769795in}{0.947518in}}{\pgfqpoint{10.769795in}{0.934460in}}%
\pgfpathlineto{\pgfqpoint{10.769795in}{0.934460in}}%
\pgfpathclose%
\pgfpathmoveto{\pgfqpoint{10.831919in}{0.928660in}}%
\pgfpathquadraticcurveto{\pgfqpoint{10.831919in}{0.939917in}}{\pgfqpoint{10.827272in}{0.946096in}}%
\pgfpathquadraticcurveto{\pgfqpoint{10.822643in}{0.952292in}}{\pgfqpoint{10.814249in}{0.952292in}}%
\pgfpathquadraticcurveto{\pgfqpoint{10.805927in}{0.952292in}}{\pgfqpoint{10.801280in}{0.946096in}}%
\pgfpathquadraticcurveto{\pgfqpoint{10.796633in}{0.939917in}}{\pgfqpoint{10.796633in}{0.928660in}}%
\pgfpathquadraticcurveto{\pgfqpoint{10.796633in}{0.917456in}}{\pgfqpoint{10.801280in}{0.911260in}}%
\pgfpathquadraticcurveto{\pgfqpoint{10.805927in}{0.905064in}}{\pgfqpoint{10.814249in}{0.905064in}}%
\pgfpathquadraticcurveto{\pgfqpoint{10.822643in}{0.905064in}}{\pgfqpoint{10.827272in}{0.911260in}}%
\pgfpathquadraticcurveto{\pgfqpoint{10.831919in}{0.917456in}}{\pgfqpoint{10.831919in}{0.928660in}}%
\pgfpathlineto{\pgfqpoint{10.831919in}{0.928660in}}%
\pgfpathclose%
\pgfpathmoveto{\pgfqpoint{10.842276in}{0.904217in}}%
\pgfpathquadraticcurveto{\pgfqpoint{10.842276in}{0.888132in}}{\pgfqpoint{10.835125in}{0.880279in}}%
\pgfpathquadraticcurveto{\pgfqpoint{10.827992in}{0.872426in}}{\pgfqpoint{10.813240in}{0.872426in}}%
\pgfpathquadraticcurveto{\pgfqpoint{10.807783in}{0.872426in}}{\pgfqpoint{10.802937in}{0.873236in}}%
\pgfpathquadraticcurveto{\pgfqpoint{10.798092in}{0.874047in}}{\pgfqpoint{10.793535in}{0.875740in}}%
\pgfpathlineto{\pgfqpoint{10.793535in}{0.885809in}}%
\pgfpathquadraticcurveto{\pgfqpoint{10.798092in}{0.883341in}}{\pgfqpoint{10.802541in}{0.882170in}}%
\pgfpathquadraticcurveto{\pgfqpoint{10.806990in}{0.880982in}}{\pgfqpoint{10.811601in}{0.880982in}}%
\pgfpathquadraticcurveto{\pgfqpoint{10.821796in}{0.880982in}}{\pgfqpoint{10.826857in}{0.886295in}}%
\pgfpathquadraticcurveto{\pgfqpoint{10.831919in}{0.891609in}}{\pgfqpoint{10.831919in}{0.902362in}}%
\pgfpathlineto{\pgfqpoint{10.831919in}{0.907495in}}%
\pgfpathquadraticcurveto{\pgfqpoint{10.828713in}{0.901912in}}{\pgfqpoint{10.823705in}{0.899156in}}%
\pgfpathquadraticcurveto{\pgfqpoint{10.818698in}{0.896400in}}{\pgfqpoint{10.811709in}{0.896400in}}%
\pgfpathquadraticcurveto{\pgfqpoint{10.800127in}{0.896400in}}{\pgfqpoint{10.793031in}{0.905226in}}%
\pgfpathquadraticcurveto{\pgfqpoint{10.785934in}{0.914070in}}{\pgfqpoint{10.785934in}{0.928660in}}%
\pgfpathquadraticcurveto{\pgfqpoint{10.785934in}{0.943286in}}{\pgfqpoint{10.793031in}{0.952112in}}%
\pgfpathquadraticcurveto{\pgfqpoint{10.800127in}{0.960956in}}{\pgfqpoint{10.811709in}{0.960956in}}%
\pgfpathquadraticcurveto{\pgfqpoint{10.818698in}{0.960956in}}{\pgfqpoint{10.823705in}{0.958200in}}%
\pgfpathquadraticcurveto{\pgfqpoint{10.828713in}{0.955444in}}{\pgfqpoint{10.831919in}{0.949878in}}%
\pgfpathlineto{\pgfqpoint{10.831919in}{0.959443in}}%
\pgfpathlineto{\pgfqpoint{10.842276in}{0.959443in}}%
\pgfpathlineto{\pgfqpoint{10.842276in}{0.904217in}}%
\pgfpathlineto{\pgfqpoint{10.842276in}{0.904217in}}%
\pgfpathclose%
\pgfpathmoveto{\pgfqpoint{10.903805in}{0.957587in}}%
\pgfpathlineto{\pgfqpoint{10.903805in}{0.947789in}}%
\pgfpathquadraticcurveto{\pgfqpoint{10.899428in}{0.950040in}}{\pgfqpoint{10.894691in}{0.951157in}}%
\pgfpathquadraticcurveto{\pgfqpoint{10.889972in}{0.952292in}}{\pgfqpoint{10.884893in}{0.952292in}}%
\pgfpathquadraticcurveto{\pgfqpoint{10.877183in}{0.952292in}}{\pgfqpoint{10.873329in}{0.949932in}}%
\pgfpathquadraticcurveto{\pgfqpoint{10.869474in}{0.947573in}}{\pgfqpoint{10.869474in}{0.942835in}}%
\pgfpathquadraticcurveto{\pgfqpoint{10.869474in}{0.939233in}}{\pgfqpoint{10.872230in}{0.937180in}}%
\pgfpathquadraticcurveto{\pgfqpoint{10.874986in}{0.935126in}}{\pgfqpoint{10.883326in}{0.933271in}}%
\pgfpathlineto{\pgfqpoint{10.886874in}{0.932478in}}%
\pgfpathquadraticcurveto{\pgfqpoint{10.897897in}{0.930119in}}{\pgfqpoint{10.902545in}{0.925814in}}%
\pgfpathquadraticcurveto{\pgfqpoint{10.907192in}{0.921509in}}{\pgfqpoint{10.907192in}{0.913800in}}%
\pgfpathquadraticcurveto{\pgfqpoint{10.907192in}{0.905010in}}{\pgfqpoint{10.900239in}{0.899876in}}%
\pgfpathquadraticcurveto{\pgfqpoint{10.893286in}{0.894761in}}{\pgfqpoint{10.881128in}{0.894761in}}%
\pgfpathquadraticcurveto{\pgfqpoint{10.876067in}{0.894761in}}{\pgfqpoint{10.870573in}{0.895752in}}%
\pgfpathquadraticcurveto{\pgfqpoint{10.865079in}{0.896742in}}{\pgfqpoint{10.859009in}{0.898706in}}%
\pgfpathlineto{\pgfqpoint{10.859009in}{0.909405in}}%
\pgfpathquadraticcurveto{\pgfqpoint{10.864755in}{0.906415in}}{\pgfqpoint{10.870321in}{0.904920in}}%
\pgfpathquadraticcurveto{\pgfqpoint{10.875887in}{0.903443in}}{\pgfqpoint{10.881362in}{0.903443in}}%
\pgfpathquadraticcurveto{\pgfqpoint{10.888675in}{0.903443in}}{\pgfqpoint{10.892602in}{0.905946in}}%
\pgfpathquadraticcurveto{\pgfqpoint{10.896547in}{0.908450in}}{\pgfqpoint{10.896547in}{0.913007in}}%
\pgfpathquadraticcurveto{\pgfqpoint{10.896547in}{0.917222in}}{\pgfqpoint{10.893701in}{0.919474in}}%
\pgfpathquadraticcurveto{\pgfqpoint{10.890873in}{0.921725in}}{\pgfqpoint{10.881236in}{0.923814in}}%
\pgfpathlineto{\pgfqpoint{10.877634in}{0.924661in}}%
\pgfpathquadraticcurveto{\pgfqpoint{10.868015in}{0.926678in}}{\pgfqpoint{10.863728in}{0.930875in}}%
\pgfpathquadraticcurveto{\pgfqpoint{10.859459in}{0.935072in}}{\pgfqpoint{10.859459in}{0.942385in}}%
\pgfpathquadraticcurveto{\pgfqpoint{10.859459in}{0.951283in}}{\pgfqpoint{10.865764in}{0.956110in}}%
\pgfpathquadraticcurveto{\pgfqpoint{10.872068in}{0.960956in}}{\pgfqpoint{10.883668in}{0.960956in}}%
\pgfpathquadraticcurveto{\pgfqpoint{10.889396in}{0.960956in}}{\pgfqpoint{10.894457in}{0.960109in}}%
\pgfpathquadraticcurveto{\pgfqpoint{10.899537in}{0.959280in}}{\pgfqpoint{10.903805in}{0.957587in}}%
\pgfpathlineto{\pgfqpoint{10.903805in}{0.957587in}}%
\pgfpathclose%
\pgfpathmoveto{\pgfqpoint{10.926321in}{0.956002in}}%
\pgfpathlineto{\pgfqpoint{10.938191in}{0.956002in}}%
\pgfpathlineto{\pgfqpoint{10.938191in}{0.941719in}}%
\pgfpathlineto{\pgfqpoint{10.926321in}{0.941719in}}%
\pgfpathlineto{\pgfqpoint{10.926321in}{0.956002in}}%
\pgfpathlineto{\pgfqpoint{10.926321in}{0.956002in}}%
\pgfpathclose%
\pgfpathmoveto{\pgfqpoint{10.926321in}{0.910702in}}%
\pgfpathlineto{\pgfqpoint{10.938191in}{0.910702in}}%
\pgfpathlineto{\pgfqpoint{10.938191in}{0.901011in}}%
\pgfpathlineto{\pgfqpoint{10.928968in}{0.882999in}}%
\pgfpathlineto{\pgfqpoint{10.921709in}{0.882999in}}%
\pgfpathlineto{\pgfqpoint{10.926321in}{0.901011in}}%
\pgfpathlineto{\pgfqpoint{10.926321in}{0.910702in}}%
\pgfpathlineto{\pgfqpoint{10.926321in}{0.910702in}}%
\pgfpathclose%
\pgfusepath{fill}%
\end{pgfscope}%
\begin{pgfscope}%
\pgfsetbuttcap%
\pgfsetroundjoin%
\definecolor{currentfill}{rgb}{0.309804,0.372549,0.419608}%
\pgfsetfillcolor{currentfill}%
\pgfsetlinewidth{0.000000pt}%
\definecolor{currentstroke}{rgb}{0.309804,0.372549,0.419608}%
\pgfsetstrokecolor{currentstroke}%
\pgfsetdash{}{0pt}%
\pgfpathmoveto{\pgfqpoint{11.077643in}{0.957587in}}%
\pgfpathlineto{\pgfqpoint{11.077643in}{0.947789in}}%
\pgfpathquadraticcurveto{\pgfqpoint{11.073266in}{0.950040in}}{\pgfqpoint{11.068528in}{0.951157in}}%
\pgfpathquadraticcurveto{\pgfqpoint{11.063809in}{0.952292in}}{\pgfqpoint{11.058730in}{0.952292in}}%
\pgfpathquadraticcurveto{\pgfqpoint{11.051021in}{0.952292in}}{\pgfqpoint{11.047166in}{0.949932in}}%
\pgfpathquadraticcurveto{\pgfqpoint{11.043311in}{0.947573in}}{\pgfqpoint{11.043311in}{0.942835in}}%
\pgfpathquadraticcurveto{\pgfqpoint{11.043311in}{0.939233in}}{\pgfqpoint{11.046067in}{0.937180in}}%
\pgfpathquadraticcurveto{\pgfqpoint{11.048823in}{0.935126in}}{\pgfqpoint{11.057163in}{0.933271in}}%
\pgfpathlineto{\pgfqpoint{11.060711in}{0.932478in}}%
\pgfpathquadraticcurveto{\pgfqpoint{11.071735in}{0.930119in}}{\pgfqpoint{11.076382in}{0.925814in}}%
\pgfpathquadraticcurveto{\pgfqpoint{11.081029in}{0.921509in}}{\pgfqpoint{11.081029in}{0.913800in}}%
\pgfpathquadraticcurveto{\pgfqpoint{11.081029in}{0.905010in}}{\pgfqpoint{11.074076in}{0.899876in}}%
\pgfpathquadraticcurveto{\pgfqpoint{11.067123in}{0.894761in}}{\pgfqpoint{11.054965in}{0.894761in}}%
\pgfpathquadraticcurveto{\pgfqpoint{11.049904in}{0.894761in}}{\pgfqpoint{11.044410in}{0.895752in}}%
\pgfpathquadraticcurveto{\pgfqpoint{11.038916in}{0.896742in}}{\pgfqpoint{11.032846in}{0.898706in}}%
\pgfpathlineto{\pgfqpoint{11.032846in}{0.909405in}}%
\pgfpathquadraticcurveto{\pgfqpoint{11.038592in}{0.906415in}}{\pgfqpoint{11.044158in}{0.904920in}}%
\pgfpathquadraticcurveto{\pgfqpoint{11.049724in}{0.903443in}}{\pgfqpoint{11.055199in}{0.903443in}}%
\pgfpathquadraticcurveto{\pgfqpoint{11.062512in}{0.903443in}}{\pgfqpoint{11.066439in}{0.905946in}}%
\pgfpathquadraticcurveto{\pgfqpoint{11.070384in}{0.908450in}}{\pgfqpoint{11.070384in}{0.913007in}}%
\pgfpathquadraticcurveto{\pgfqpoint{11.070384in}{0.917222in}}{\pgfqpoint{11.067538in}{0.919474in}}%
\pgfpathquadraticcurveto{\pgfqpoint{11.064710in}{0.921725in}}{\pgfqpoint{11.055073in}{0.923814in}}%
\pgfpathlineto{\pgfqpoint{11.051471in}{0.924661in}}%
\pgfpathquadraticcurveto{\pgfqpoint{11.041852in}{0.926678in}}{\pgfqpoint{11.037565in}{0.930875in}}%
\pgfpathquadraticcurveto{\pgfqpoint{11.033297in}{0.935072in}}{\pgfqpoint{11.033297in}{0.942385in}}%
\pgfpathquadraticcurveto{\pgfqpoint{11.033297in}{0.951283in}}{\pgfqpoint{11.039601in}{0.956110in}}%
\pgfpathquadraticcurveto{\pgfqpoint{11.045905in}{0.960956in}}{\pgfqpoint{11.057505in}{0.960956in}}%
\pgfpathquadraticcurveto{\pgfqpoint{11.063233in}{0.960956in}}{\pgfqpoint{11.068294in}{0.960109in}}%
\pgfpathquadraticcurveto{\pgfqpoint{11.073374in}{0.959280in}}{\pgfqpoint{11.077643in}{0.957587in}}%
\pgfpathlineto{\pgfqpoint{11.077643in}{0.957587in}}%
\pgfpathclose%
\pgfpathmoveto{\pgfqpoint{11.097510in}{0.959443in}}%
\pgfpathlineto{\pgfqpoint{11.107867in}{0.959443in}}%
\pgfpathlineto{\pgfqpoint{11.107867in}{0.896400in}}%
\pgfpathlineto{\pgfqpoint{11.097510in}{0.896400in}}%
\pgfpathlineto{\pgfqpoint{11.097510in}{0.959443in}}%
\pgfpathlineto{\pgfqpoint{11.097510in}{0.959443in}}%
\pgfpathclose%
\pgfpathmoveto{\pgfqpoint{11.097510in}{0.983993in}}%
\pgfpathlineto{\pgfqpoint{11.107867in}{0.983993in}}%
\pgfpathlineto{\pgfqpoint{11.107867in}{0.970862in}}%
\pgfpathlineto{\pgfqpoint{11.097510in}{0.970862in}}%
\pgfpathlineto{\pgfqpoint{11.097510in}{0.983993in}}%
\pgfpathlineto{\pgfqpoint{11.097510in}{0.983993in}}%
\pgfpathclose%
\pgfpathmoveto{\pgfqpoint{11.171018in}{0.928660in}}%
\pgfpathquadraticcurveto{\pgfqpoint{11.171018in}{0.939917in}}{\pgfqpoint{11.166370in}{0.946096in}}%
\pgfpathquadraticcurveto{\pgfqpoint{11.161741in}{0.952292in}}{\pgfqpoint{11.153348in}{0.952292in}}%
\pgfpathquadraticcurveto{\pgfqpoint{11.145026in}{0.952292in}}{\pgfqpoint{11.140379in}{0.946096in}}%
\pgfpathquadraticcurveto{\pgfqpoint{11.135732in}{0.939917in}}{\pgfqpoint{11.135732in}{0.928660in}}%
\pgfpathquadraticcurveto{\pgfqpoint{11.135732in}{0.917456in}}{\pgfqpoint{11.140379in}{0.911260in}}%
\pgfpathquadraticcurveto{\pgfqpoint{11.145026in}{0.905064in}}{\pgfqpoint{11.153348in}{0.905064in}}%
\pgfpathquadraticcurveto{\pgfqpoint{11.161741in}{0.905064in}}{\pgfqpoint{11.166370in}{0.911260in}}%
\pgfpathquadraticcurveto{\pgfqpoint{11.171018in}{0.917456in}}{\pgfqpoint{11.171018in}{0.928660in}}%
\pgfpathlineto{\pgfqpoint{11.171018in}{0.928660in}}%
\pgfpathclose%
\pgfpathmoveto{\pgfqpoint{11.181374in}{0.904217in}}%
\pgfpathquadraticcurveto{\pgfqpoint{11.181374in}{0.888132in}}{\pgfqpoint{11.174224in}{0.880279in}}%
\pgfpathquadraticcurveto{\pgfqpoint{11.167091in}{0.872426in}}{\pgfqpoint{11.152339in}{0.872426in}}%
\pgfpathquadraticcurveto{\pgfqpoint{11.146881in}{0.872426in}}{\pgfqpoint{11.142036in}{0.873236in}}%
\pgfpathquadraticcurveto{\pgfqpoint{11.137191in}{0.874047in}}{\pgfqpoint{11.132634in}{0.875740in}}%
\pgfpathlineto{\pgfqpoint{11.132634in}{0.885809in}}%
\pgfpathquadraticcurveto{\pgfqpoint{11.137191in}{0.883341in}}{\pgfqpoint{11.141640in}{0.882170in}}%
\pgfpathquadraticcurveto{\pgfqpoint{11.146089in}{0.880982in}}{\pgfqpoint{11.150700in}{0.880982in}}%
\pgfpathquadraticcurveto{\pgfqpoint{11.160895in}{0.880982in}}{\pgfqpoint{11.165956in}{0.886295in}}%
\pgfpathquadraticcurveto{\pgfqpoint{11.171018in}{0.891609in}}{\pgfqpoint{11.171018in}{0.902362in}}%
\pgfpathlineto{\pgfqpoint{11.171018in}{0.907495in}}%
\pgfpathquadraticcurveto{\pgfqpoint{11.167811in}{0.901912in}}{\pgfqpoint{11.162804in}{0.899156in}}%
\pgfpathquadraticcurveto{\pgfqpoint{11.157797in}{0.896400in}}{\pgfqpoint{11.150808in}{0.896400in}}%
\pgfpathquadraticcurveto{\pgfqpoint{11.139226in}{0.896400in}}{\pgfqpoint{11.132129in}{0.905226in}}%
\pgfpathquadraticcurveto{\pgfqpoint{11.125032in}{0.914070in}}{\pgfqpoint{11.125032in}{0.928660in}}%
\pgfpathquadraticcurveto{\pgfqpoint{11.125032in}{0.943286in}}{\pgfqpoint{11.132129in}{0.952112in}}%
\pgfpathquadraticcurveto{\pgfqpoint{11.139226in}{0.960956in}}{\pgfqpoint{11.150808in}{0.960956in}}%
\pgfpathquadraticcurveto{\pgfqpoint{11.157797in}{0.960956in}}{\pgfqpoint{11.162804in}{0.958200in}}%
\pgfpathquadraticcurveto{\pgfqpoint{11.167811in}{0.955444in}}{\pgfqpoint{11.171018in}{0.949878in}}%
\pgfpathlineto{\pgfqpoint{11.171018in}{0.959443in}}%
\pgfpathlineto{\pgfqpoint{11.181374in}{0.959443in}}%
\pgfpathlineto{\pgfqpoint{11.181374in}{0.904217in}}%
\pgfpathlineto{\pgfqpoint{11.181374in}{0.904217in}}%
\pgfpathclose%
\pgfpathmoveto{\pgfqpoint{11.255134in}{0.934460in}}%
\pgfpathlineto{\pgfqpoint{11.255134in}{0.896400in}}%
\pgfpathlineto{\pgfqpoint{11.244777in}{0.896400in}}%
\pgfpathlineto{\pgfqpoint{11.244777in}{0.934117in}}%
\pgfpathquadraticcurveto{\pgfqpoint{11.244777in}{0.943069in}}{\pgfqpoint{11.241283in}{0.947500in}}%
\pgfpathquadraticcurveto{\pgfqpoint{11.237789in}{0.951949in}}{\pgfqpoint{11.230818in}{0.951949in}}%
\pgfpathquadraticcurveto{\pgfqpoint{11.222424in}{0.951949in}}{\pgfqpoint{11.217579in}{0.946600in}}%
\pgfpathquadraticcurveto{\pgfqpoint{11.212734in}{0.941268in}}{\pgfqpoint{11.212734in}{0.932028in}}%
\pgfpathlineto{\pgfqpoint{11.212734in}{0.896400in}}%
\pgfpathlineto{\pgfqpoint{11.202323in}{0.896400in}}%
\pgfpathlineto{\pgfqpoint{11.202323in}{0.959443in}}%
\pgfpathlineto{\pgfqpoint{11.212734in}{0.959443in}}%
\pgfpathlineto{\pgfqpoint{11.212734in}{0.949644in}}%
\pgfpathquadraticcurveto{\pgfqpoint{11.216462in}{0.955336in}}{\pgfqpoint{11.221488in}{0.958146in}}%
\pgfpathquadraticcurveto{\pgfqpoint{11.226531in}{0.960956in}}{\pgfqpoint{11.233123in}{0.960956in}}%
\pgfpathquadraticcurveto{\pgfqpoint{11.243985in}{0.960956in}}{\pgfqpoint{11.249550in}{0.954237in}}%
\pgfpathquadraticcurveto{\pgfqpoint{11.255134in}{0.947518in}}{\pgfqpoint{11.255134in}{0.934460in}}%
\pgfpathlineto{\pgfqpoint{11.255134in}{0.934460in}}%
\pgfpathclose%
\pgfpathmoveto{\pgfqpoint{11.275776in}{0.959443in}}%
\pgfpathlineto{\pgfqpoint{11.286133in}{0.959443in}}%
\pgfpathlineto{\pgfqpoint{11.286133in}{0.896400in}}%
\pgfpathlineto{\pgfqpoint{11.275776in}{0.896400in}}%
\pgfpathlineto{\pgfqpoint{11.275776in}{0.959443in}}%
\pgfpathlineto{\pgfqpoint{11.275776in}{0.959443in}}%
\pgfpathclose%
\pgfpathmoveto{\pgfqpoint{11.275776in}{0.983993in}}%
\pgfpathlineto{\pgfqpoint{11.286133in}{0.983993in}}%
\pgfpathlineto{\pgfqpoint{11.286133in}{0.970862in}}%
\pgfpathlineto{\pgfqpoint{11.275776in}{0.970862in}}%
\pgfpathlineto{\pgfqpoint{11.275776in}{0.983993in}}%
\pgfpathlineto{\pgfqpoint{11.275776in}{0.983993in}}%
\pgfpathclose%
\pgfpathmoveto{\pgfqpoint{11.358740in}{0.959443in}}%
\pgfpathlineto{\pgfqpoint{11.358740in}{0.896400in}}%
\pgfpathlineto{\pgfqpoint{11.348329in}{0.896400in}}%
\pgfpathlineto{\pgfqpoint{11.348329in}{0.951391in}}%
\pgfpathlineto{\pgfqpoint{11.319906in}{0.951391in}}%
\pgfpathlineto{\pgfqpoint{11.319906in}{0.896400in}}%
\pgfpathlineto{\pgfqpoint{11.309495in}{0.896400in}}%
\pgfpathlineto{\pgfqpoint{11.309495in}{0.951391in}}%
\pgfpathlineto{\pgfqpoint{11.299588in}{0.951391in}}%
\pgfpathlineto{\pgfqpoint{11.299588in}{0.959443in}}%
\pgfpathlineto{\pgfqpoint{11.309495in}{0.959443in}}%
\pgfpathlineto{\pgfqpoint{11.309495in}{0.963838in}}%
\pgfpathquadraticcurveto{\pgfqpoint{11.309495in}{0.974140in}}{\pgfqpoint{11.314358in}{0.979058in}}%
\pgfpathquadraticcurveto{\pgfqpoint{11.319240in}{0.983993in}}{\pgfqpoint{11.329308in}{0.983993in}}%
\pgfpathlineto{\pgfqpoint{11.339719in}{0.983993in}}%
\pgfpathlineto{\pgfqpoint{11.339719in}{0.975365in}}%
\pgfpathlineto{\pgfqpoint{11.329813in}{0.975365in}}%
\pgfpathquadraticcurveto{\pgfqpoint{11.324247in}{0.975365in}}{\pgfqpoint{11.322067in}{0.973114in}}%
\pgfpathquadraticcurveto{\pgfqpoint{11.319906in}{0.970862in}}{\pgfqpoint{11.319906in}{0.965008in}}%
\pgfpathlineto{\pgfqpoint{11.319906in}{0.959443in}}%
\pgfpathlineto{\pgfqpoint{11.358740in}{0.959443in}}%
\pgfpathlineto{\pgfqpoint{11.358740in}{0.959443in}}%
\pgfpathclose%
\pgfpathmoveto{\pgfqpoint{11.348329in}{0.983867in}}%
\pgfpathlineto{\pgfqpoint{11.358740in}{0.983867in}}%
\pgfpathlineto{\pgfqpoint{11.358740in}{0.970754in}}%
\pgfpathlineto{\pgfqpoint{11.348329in}{0.970754in}}%
\pgfpathlineto{\pgfqpoint{11.348329in}{0.983867in}}%
\pgfpathlineto{\pgfqpoint{11.348329in}{0.983867in}}%
\pgfpathclose%
\pgfpathmoveto{\pgfqpoint{11.425781in}{0.957029in}}%
\pgfpathlineto{\pgfqpoint{11.425781in}{0.947338in}}%
\pgfpathquadraticcurveto{\pgfqpoint{11.421386in}{0.949770in}}{\pgfqpoint{11.416973in}{0.950977in}}%
\pgfpathquadraticcurveto{\pgfqpoint{11.412560in}{0.952184in}}{\pgfqpoint{11.408057in}{0.952184in}}%
\pgfpathquadraticcurveto{\pgfqpoint{11.397971in}{0.952184in}}{\pgfqpoint{11.392387in}{0.945789in}}%
\pgfpathquadraticcurveto{\pgfqpoint{11.386821in}{0.939413in}}{\pgfqpoint{11.386821in}{0.927867in}}%
\pgfpathquadraticcurveto{\pgfqpoint{11.386821in}{0.916321in}}{\pgfqpoint{11.392387in}{0.909927in}}%
\pgfpathquadraticcurveto{\pgfqpoint{11.397971in}{0.903551in}}{\pgfqpoint{11.408057in}{0.903551in}}%
\pgfpathquadraticcurveto{\pgfqpoint{11.412560in}{0.903551in}}{\pgfqpoint{11.416973in}{0.904758in}}%
\pgfpathquadraticcurveto{\pgfqpoint{11.421386in}{0.905964in}}{\pgfqpoint{11.425781in}{0.908396in}}%
\pgfpathlineto{\pgfqpoint{11.425781in}{0.898814in}}%
\pgfpathquadraticcurveto{\pgfqpoint{11.421440in}{0.896796in}}{\pgfqpoint{11.416793in}{0.895788in}}%
\pgfpathquadraticcurveto{\pgfqpoint{11.412164in}{0.894761in}}{\pgfqpoint{11.406923in}{0.894761in}}%
\pgfpathquadraticcurveto{\pgfqpoint{11.392675in}{0.894761in}}{\pgfqpoint{11.384281in}{0.903713in}}%
\pgfpathquadraticcurveto{\pgfqpoint{11.375906in}{0.912665in}}{\pgfqpoint{11.375906in}{0.927867in}}%
\pgfpathquadraticcurveto{\pgfqpoint{11.375906in}{0.943286in}}{\pgfqpoint{11.384371in}{0.952112in}}%
\pgfpathquadraticcurveto{\pgfqpoint{11.392855in}{0.960956in}}{\pgfqpoint{11.407607in}{0.960956in}}%
\pgfpathquadraticcurveto{\pgfqpoint{11.412380in}{0.960956in}}{\pgfqpoint{11.416937in}{0.959965in}}%
\pgfpathquadraticcurveto{\pgfqpoint{11.421495in}{0.958992in}}{\pgfqpoint{11.425781in}{0.957029in}}%
\pgfpathlineto{\pgfqpoint{11.425781in}{0.957029in}}%
\pgfpathclose%
\pgfpathmoveto{\pgfqpoint{11.472451in}{0.928083in}}%
\pgfpathquadraticcurveto{\pgfqpoint{11.459896in}{0.928083in}}{\pgfqpoint{11.455051in}{0.925219in}}%
\pgfpathquadraticcurveto{\pgfqpoint{11.450206in}{0.922356in}}{\pgfqpoint{11.450206in}{0.915421in}}%
\pgfpathquadraticcurveto{\pgfqpoint{11.450206in}{0.909909in}}{\pgfqpoint{11.453844in}{0.906667in}}%
\pgfpathquadraticcurveto{\pgfqpoint{11.457483in}{0.903443in}}{\pgfqpoint{11.463715in}{0.903443in}}%
\pgfpathquadraticcurveto{\pgfqpoint{11.472343in}{0.903443in}}{\pgfqpoint{11.477548in}{0.909549in}}%
\pgfpathquadraticcurveto{\pgfqpoint{11.482754in}{0.915655in}}{\pgfqpoint{11.482754in}{0.925778in}}%
\pgfpathlineto{\pgfqpoint{11.482754in}{0.928083in}}%
\pgfpathlineto{\pgfqpoint{11.472451in}{0.928083in}}%
\pgfpathlineto{\pgfqpoint{11.472451in}{0.928083in}}%
\pgfpathclose%
\pgfpathmoveto{\pgfqpoint{11.493111in}{0.932370in}}%
\pgfpathlineto{\pgfqpoint{11.493111in}{0.896400in}}%
\pgfpathlineto{\pgfqpoint{11.482754in}{0.896400in}}%
\pgfpathlineto{\pgfqpoint{11.482754in}{0.905964in}}%
\pgfpathquadraticcurveto{\pgfqpoint{11.479205in}{0.900237in}}{\pgfqpoint{11.473910in}{0.897499in}}%
\pgfpathquadraticcurveto{\pgfqpoint{11.468614in}{0.894761in}}{\pgfqpoint{11.460959in}{0.894761in}}%
\pgfpathquadraticcurveto{\pgfqpoint{11.451287in}{0.894761in}}{\pgfqpoint{11.445559in}{0.900201in}}%
\pgfpathquadraticcurveto{\pgfqpoint{11.439849in}{0.905640in}}{\pgfqpoint{11.439849in}{0.914754in}}%
\pgfpathquadraticcurveto{\pgfqpoint{11.439849in}{0.925382in}}{\pgfqpoint{11.446964in}{0.930785in}}%
\pgfpathquadraticcurveto{\pgfqpoint{11.454096in}{0.936189in}}{\pgfqpoint{11.468218in}{0.936189in}}%
\pgfpathlineto{\pgfqpoint{11.482754in}{0.936189in}}%
\pgfpathlineto{\pgfqpoint{11.482754in}{0.937216in}}%
\pgfpathquadraticcurveto{\pgfqpoint{11.482754in}{0.944366in}}{\pgfqpoint{11.478053in}{0.948275in}}%
\pgfpathquadraticcurveto{\pgfqpoint{11.473351in}{0.952184in}}{\pgfqpoint{11.464850in}{0.952184in}}%
\pgfpathquadraticcurveto{\pgfqpoint{11.459446in}{0.952184in}}{\pgfqpoint{11.454313in}{0.950887in}}%
\pgfpathquadraticcurveto{\pgfqpoint{11.449197in}{0.949590in}}{\pgfqpoint{11.444478in}{0.946996in}}%
\pgfpathlineto{\pgfqpoint{11.444478in}{0.956579in}}%
\pgfpathquadraticcurveto{\pgfqpoint{11.450152in}{0.958776in}}{\pgfqpoint{11.455501in}{0.959857in}}%
\pgfpathquadraticcurveto{\pgfqpoint{11.460851in}{0.960956in}}{\pgfqpoint{11.465912in}{0.960956in}}%
\pgfpathquadraticcurveto{\pgfqpoint{11.479602in}{0.960956in}}{\pgfqpoint{11.486356in}{0.953859in}}%
\pgfpathquadraticcurveto{\pgfqpoint{11.493111in}{0.946780in}}{\pgfqpoint{11.493111in}{0.932370in}}%
\pgfpathlineto{\pgfqpoint{11.493111in}{0.932370in}}%
\pgfpathclose%
\pgfpathmoveto{\pgfqpoint{11.566853in}{0.934460in}}%
\pgfpathlineto{\pgfqpoint{11.566853in}{0.896400in}}%
\pgfpathlineto{\pgfqpoint{11.556496in}{0.896400in}}%
\pgfpathlineto{\pgfqpoint{11.556496in}{0.934117in}}%
\pgfpathquadraticcurveto{\pgfqpoint{11.556496in}{0.943069in}}{\pgfqpoint{11.553001in}{0.947500in}}%
\pgfpathquadraticcurveto{\pgfqpoint{11.549507in}{0.951949in}}{\pgfqpoint{11.542536in}{0.951949in}}%
\pgfpathquadraticcurveto{\pgfqpoint{11.534143in}{0.951949in}}{\pgfqpoint{11.529297in}{0.946600in}}%
\pgfpathquadraticcurveto{\pgfqpoint{11.524452in}{0.941268in}}{\pgfqpoint{11.524452in}{0.932028in}}%
\pgfpathlineto{\pgfqpoint{11.524452in}{0.896400in}}%
\pgfpathlineto{\pgfqpoint{11.514041in}{0.896400in}}%
\pgfpathlineto{\pgfqpoint{11.514041in}{0.959443in}}%
\pgfpathlineto{\pgfqpoint{11.524452in}{0.959443in}}%
\pgfpathlineto{\pgfqpoint{11.524452in}{0.949644in}}%
\pgfpathquadraticcurveto{\pgfqpoint{11.528180in}{0.955336in}}{\pgfqpoint{11.533206in}{0.958146in}}%
\pgfpathquadraticcurveto{\pgfqpoint{11.538249in}{0.960956in}}{\pgfqpoint{11.544842in}{0.960956in}}%
\pgfpathquadraticcurveto{\pgfqpoint{11.555703in}{0.960956in}}{\pgfqpoint{11.561269in}{0.954237in}}%
\pgfpathquadraticcurveto{\pgfqpoint{11.566853in}{0.947518in}}{\pgfqpoint{11.566853in}{0.934460in}}%
\pgfpathlineto{\pgfqpoint{11.566853in}{0.934460in}}%
\pgfpathclose%
\pgfpathmoveto{\pgfqpoint{11.632867in}{0.957029in}}%
\pgfpathlineto{\pgfqpoint{11.632867in}{0.947338in}}%
\pgfpathquadraticcurveto{\pgfqpoint{11.628472in}{0.949770in}}{\pgfqpoint{11.624059in}{0.950977in}}%
\pgfpathquadraticcurveto{\pgfqpoint{11.619646in}{0.952184in}}{\pgfqpoint{11.615143in}{0.952184in}}%
\pgfpathquadraticcurveto{\pgfqpoint{11.605056in}{0.952184in}}{\pgfqpoint{11.599473in}{0.945789in}}%
\pgfpathquadraticcurveto{\pgfqpoint{11.593907in}{0.939413in}}{\pgfqpoint{11.593907in}{0.927867in}}%
\pgfpathquadraticcurveto{\pgfqpoint{11.593907in}{0.916321in}}{\pgfqpoint{11.599473in}{0.909927in}}%
\pgfpathquadraticcurveto{\pgfqpoint{11.605056in}{0.903551in}}{\pgfqpoint{11.615143in}{0.903551in}}%
\pgfpathquadraticcurveto{\pgfqpoint{11.619646in}{0.903551in}}{\pgfqpoint{11.624059in}{0.904758in}}%
\pgfpathquadraticcurveto{\pgfqpoint{11.628472in}{0.905964in}}{\pgfqpoint{11.632867in}{0.908396in}}%
\pgfpathlineto{\pgfqpoint{11.632867in}{0.898814in}}%
\pgfpathquadraticcurveto{\pgfqpoint{11.628526in}{0.896796in}}{\pgfqpoint{11.623879in}{0.895788in}}%
\pgfpathquadraticcurveto{\pgfqpoint{11.619250in}{0.894761in}}{\pgfqpoint{11.614008in}{0.894761in}}%
\pgfpathquadraticcurveto{\pgfqpoint{11.599761in}{0.894761in}}{\pgfqpoint{11.591367in}{0.903713in}}%
\pgfpathquadraticcurveto{\pgfqpoint{11.582991in}{0.912665in}}{\pgfqpoint{11.582991in}{0.927867in}}%
\pgfpathquadraticcurveto{\pgfqpoint{11.582991in}{0.943286in}}{\pgfqpoint{11.591457in}{0.952112in}}%
\pgfpathquadraticcurveto{\pgfqpoint{11.599941in}{0.960956in}}{\pgfqpoint{11.614693in}{0.960956in}}%
\pgfpathquadraticcurveto{\pgfqpoint{11.619466in}{0.960956in}}{\pgfqpoint{11.624023in}{0.959965in}}%
\pgfpathquadraticcurveto{\pgfqpoint{11.628580in}{0.958992in}}{\pgfqpoint{11.632867in}{0.957029in}}%
\pgfpathlineto{\pgfqpoint{11.632867in}{0.957029in}}%
\pgfpathclose%
\pgfpathmoveto{\pgfqpoint{11.704808in}{0.930515in}}%
\pgfpathlineto{\pgfqpoint{11.704808in}{0.925454in}}%
\pgfpathlineto{\pgfqpoint{11.657184in}{0.925454in}}%
\pgfpathquadraticcurveto{\pgfqpoint{11.657868in}{0.914754in}}{\pgfqpoint{11.663632in}{0.909153in}}%
\pgfpathquadraticcurveto{\pgfqpoint{11.669396in}{0.903551in}}{\pgfqpoint{11.679699in}{0.903551in}}%
\pgfpathquadraticcurveto{\pgfqpoint{11.685661in}{0.903551in}}{\pgfqpoint{11.691263in}{0.905010in}}%
\pgfpathquadraticcurveto{\pgfqpoint{11.696864in}{0.906469in}}{\pgfqpoint{11.702394in}{0.909405in}}%
\pgfpathlineto{\pgfqpoint{11.702394in}{0.899606in}}%
\pgfpathquadraticcurveto{\pgfqpoint{11.696810in}{0.897247in}}{\pgfqpoint{11.690956in}{0.896004in}}%
\pgfpathquadraticcurveto{\pgfqpoint{11.685102in}{0.894761in}}{\pgfqpoint{11.679086in}{0.894761in}}%
\pgfpathquadraticcurveto{\pgfqpoint{11.663992in}{0.894761in}}{\pgfqpoint{11.655184in}{0.903533in}}%
\pgfpathquadraticcurveto{\pgfqpoint{11.646376in}{0.912323in}}{\pgfqpoint{11.646376in}{0.927309in}}%
\pgfpathquadraticcurveto{\pgfqpoint{11.646376in}{0.942781in}}{\pgfqpoint{11.654734in}{0.951859in}}%
\pgfpathquadraticcurveto{\pgfqpoint{11.663092in}{0.960956in}}{\pgfqpoint{11.677285in}{0.960956in}}%
\pgfpathquadraticcurveto{\pgfqpoint{11.690002in}{0.960956in}}{\pgfqpoint{11.697405in}{0.952760in}}%
\pgfpathquadraticcurveto{\pgfqpoint{11.704808in}{0.944583in}}{\pgfqpoint{11.704808in}{0.930515in}}%
\pgfpathlineto{\pgfqpoint{11.704808in}{0.930515in}}%
\pgfpathclose%
\pgfpathmoveto{\pgfqpoint{11.694451in}{0.933559in}}%
\pgfpathquadraticcurveto{\pgfqpoint{11.694343in}{0.942043in}}{\pgfqpoint{11.689695in}{0.947104in}}%
\pgfpathquadraticcurveto{\pgfqpoint{11.685048in}{0.952184in}}{\pgfqpoint{11.677393in}{0.952184in}}%
\pgfpathquadraticcurveto{\pgfqpoint{11.668729in}{0.952184in}}{\pgfqpoint{11.663524in}{0.947284in}}%
\pgfpathquadraticcurveto{\pgfqpoint{11.658318in}{0.942385in}}{\pgfqpoint{11.657526in}{0.933487in}}%
\pgfpathlineto{\pgfqpoint{11.694451in}{0.933559in}}%
\pgfpathlineto{\pgfqpoint{11.694451in}{0.933559in}}%
\pgfpathclose%
\pgfusepath{fill}%
\end{pgfscope}%
\begin{pgfscope}%
\pgfsetbuttcap%
\pgfsetroundjoin%
\definecolor{currentfill}{rgb}{0.309804,0.372549,0.419608}%
\pgfsetfillcolor{currentfill}%
\pgfsetlinewidth{0.000000pt}%
\definecolor{currentstroke}{rgb}{0.309804,0.372549,0.419608}%
\pgfsetstrokecolor{currentstroke}%
\pgfsetdash{}{0pt}%
\pgfpathmoveto{\pgfqpoint{11.855880in}{0.959443in}}%
\pgfpathlineto{\pgfqpoint{11.855880in}{0.896400in}}%
\pgfpathlineto{\pgfqpoint{11.845469in}{0.896400in}}%
\pgfpathlineto{\pgfqpoint{11.845469in}{0.951391in}}%
\pgfpathlineto{\pgfqpoint{11.817045in}{0.951391in}}%
\pgfpathlineto{\pgfqpoint{11.817045in}{0.896400in}}%
\pgfpathlineto{\pgfqpoint{11.806634in}{0.896400in}}%
\pgfpathlineto{\pgfqpoint{11.806634in}{0.951391in}}%
\pgfpathlineto{\pgfqpoint{11.796728in}{0.951391in}}%
\pgfpathlineto{\pgfqpoint{11.796728in}{0.959443in}}%
\pgfpathlineto{\pgfqpoint{11.806634in}{0.959443in}}%
\pgfpathlineto{\pgfqpoint{11.806634in}{0.963838in}}%
\pgfpathquadraticcurveto{\pgfqpoint{11.806634in}{0.974140in}}{\pgfqpoint{11.811498in}{0.979058in}}%
\pgfpathquadraticcurveto{\pgfqpoint{11.816379in}{0.983993in}}{\pgfqpoint{11.826448in}{0.983993in}}%
\pgfpathlineto{\pgfqpoint{11.836859in}{0.983993in}}%
\pgfpathlineto{\pgfqpoint{11.836859in}{0.975365in}}%
\pgfpathlineto{\pgfqpoint{11.826952in}{0.975365in}}%
\pgfpathquadraticcurveto{\pgfqpoint{11.821386in}{0.975365in}}{\pgfqpoint{11.819207in}{0.973114in}}%
\pgfpathquadraticcurveto{\pgfqpoint{11.817045in}{0.970862in}}{\pgfqpoint{11.817045in}{0.965008in}}%
\pgfpathlineto{\pgfqpoint{11.817045in}{0.959443in}}%
\pgfpathlineto{\pgfqpoint{11.855880in}{0.959443in}}%
\pgfpathlineto{\pgfqpoint{11.855880in}{0.959443in}}%
\pgfpathclose%
\pgfpathmoveto{\pgfqpoint{11.845469in}{0.983867in}}%
\pgfpathlineto{\pgfqpoint{11.855880in}{0.983867in}}%
\pgfpathlineto{\pgfqpoint{11.855880in}{0.970754in}}%
\pgfpathlineto{\pgfqpoint{11.845469in}{0.970754in}}%
\pgfpathlineto{\pgfqpoint{11.845469in}{0.983867in}}%
\pgfpathlineto{\pgfqpoint{11.845469in}{0.983867in}}%
\pgfpathclose%
\pgfpathmoveto{\pgfqpoint{11.877548in}{0.983993in}}%
\pgfpathlineto{\pgfqpoint{11.887905in}{0.983993in}}%
\pgfpathlineto{\pgfqpoint{11.887905in}{0.896400in}}%
\pgfpathlineto{\pgfqpoint{11.877548in}{0.896400in}}%
\pgfpathlineto{\pgfqpoint{11.877548in}{0.983993in}}%
\pgfpathlineto{\pgfqpoint{11.877548in}{0.983993in}}%
\pgfpathclose%
\pgfpathmoveto{\pgfqpoint{11.909574in}{0.983993in}}%
\pgfpathlineto{\pgfqpoint{11.919931in}{0.983993in}}%
\pgfpathlineto{\pgfqpoint{11.919931in}{0.896400in}}%
\pgfpathlineto{\pgfqpoint{11.909574in}{0.896400in}}%
\pgfpathlineto{\pgfqpoint{11.909574in}{0.983993in}}%
\pgfpathlineto{\pgfqpoint{11.909574in}{0.983993in}}%
\pgfpathclose%
\pgfusepath{fill}%
\end{pgfscope}%
\begin{pgfscope}%
\pgfsetbuttcap%
\pgfsetroundjoin%
\definecolor{currentfill}{rgb}{0.309804,0.372549,0.419608}%
\pgfsetfillcolor{currentfill}%
\pgfsetlinewidth{0.000000pt}%
\definecolor{currentstroke}{rgb}{0.309804,0.372549,0.419608}%
\pgfsetstrokecolor{currentstroke}%
\pgfsetdash{}{0pt}%
\pgfpathmoveto{\pgfqpoint{12.032610in}{0.957587in}}%
\pgfpathlineto{\pgfqpoint{12.032610in}{0.947789in}}%
\pgfpathquadraticcurveto{\pgfqpoint{12.028233in}{0.950040in}}{\pgfqpoint{12.023496in}{0.951157in}}%
\pgfpathquadraticcurveto{\pgfqpoint{12.018777in}{0.952292in}}{\pgfqpoint{12.013697in}{0.952292in}}%
\pgfpathquadraticcurveto{\pgfqpoint{12.005988in}{0.952292in}}{\pgfqpoint{12.002134in}{0.949932in}}%
\pgfpathquadraticcurveto{\pgfqpoint{11.998279in}{0.947573in}}{\pgfqpoint{11.998279in}{0.942835in}}%
\pgfpathquadraticcurveto{\pgfqpoint{11.998279in}{0.939233in}}{\pgfqpoint{12.001035in}{0.937180in}}%
\pgfpathquadraticcurveto{\pgfqpoint{12.003791in}{0.935126in}}{\pgfqpoint{12.012130in}{0.933271in}}%
\pgfpathlineto{\pgfqpoint{12.015679in}{0.932478in}}%
\pgfpathquadraticcurveto{\pgfqpoint{12.026702in}{0.930119in}}{\pgfqpoint{12.031349in}{0.925814in}}%
\pgfpathquadraticcurveto{\pgfqpoint{12.035996in}{0.921509in}}{\pgfqpoint{12.035996in}{0.913800in}}%
\pgfpathquadraticcurveto{\pgfqpoint{12.035996in}{0.905010in}}{\pgfqpoint{12.029044in}{0.899876in}}%
\pgfpathquadraticcurveto{\pgfqpoint{12.022091in}{0.894761in}}{\pgfqpoint{12.009933in}{0.894761in}}%
\pgfpathquadraticcurveto{\pgfqpoint{12.004871in}{0.894761in}}{\pgfqpoint{11.999378in}{0.895752in}}%
\pgfpathquadraticcurveto{\pgfqpoint{11.993884in}{0.896742in}}{\pgfqpoint{11.987814in}{0.898706in}}%
\pgfpathlineto{\pgfqpoint{11.987814in}{0.909405in}}%
\pgfpathquadraticcurveto{\pgfqpoint{11.993560in}{0.906415in}}{\pgfqpoint{11.999126in}{0.904920in}}%
\pgfpathquadraticcurveto{\pgfqpoint{12.004691in}{0.903443in}}{\pgfqpoint{12.010167in}{0.903443in}}%
\pgfpathquadraticcurveto{\pgfqpoint{12.017480in}{0.903443in}}{\pgfqpoint{12.021407in}{0.905946in}}%
\pgfpathquadraticcurveto{\pgfqpoint{12.025351in}{0.908450in}}{\pgfqpoint{12.025351in}{0.913007in}}%
\pgfpathquadraticcurveto{\pgfqpoint{12.025351in}{0.917222in}}{\pgfqpoint{12.022505in}{0.919474in}}%
\pgfpathquadraticcurveto{\pgfqpoint{12.019677in}{0.921725in}}{\pgfqpoint{12.010041in}{0.923814in}}%
\pgfpathlineto{\pgfqpoint{12.006438in}{0.924661in}}%
\pgfpathquadraticcurveto{\pgfqpoint{11.996820in}{0.926678in}}{\pgfqpoint{11.992533in}{0.930875in}}%
\pgfpathquadraticcurveto{\pgfqpoint{11.988264in}{0.935072in}}{\pgfqpoint{11.988264in}{0.942385in}}%
\pgfpathquadraticcurveto{\pgfqpoint{11.988264in}{0.951283in}}{\pgfqpoint{11.994568in}{0.956110in}}%
\pgfpathquadraticcurveto{\pgfqpoint{12.000873in}{0.960956in}}{\pgfqpoint{12.012473in}{0.960956in}}%
\pgfpathquadraticcurveto{\pgfqpoint{12.018200in}{0.960956in}}{\pgfqpoint{12.023262in}{0.960109in}}%
\pgfpathquadraticcurveto{\pgfqpoint{12.028341in}{0.959280in}}{\pgfqpoint{12.032610in}{0.957587in}}%
\pgfpathlineto{\pgfqpoint{12.032610in}{0.957587in}}%
\pgfpathclose%
\pgfpathmoveto{\pgfqpoint{12.051415in}{0.921275in}}%
\pgfpathlineto{\pgfqpoint{12.051415in}{0.959443in}}%
\pgfpathlineto{\pgfqpoint{12.061772in}{0.959443in}}%
\pgfpathlineto{\pgfqpoint{12.061772in}{0.921671in}}%
\pgfpathquadraticcurveto{\pgfqpoint{12.061772in}{0.912719in}}{\pgfqpoint{12.065248in}{0.908234in}}%
\pgfpathquadraticcurveto{\pgfqpoint{12.068742in}{0.903767in}}{\pgfqpoint{12.075731in}{0.903767in}}%
\pgfpathquadraticcurveto{\pgfqpoint{12.084107in}{0.903767in}}{\pgfqpoint{12.088970in}{0.909117in}}%
\pgfpathquadraticcurveto{\pgfqpoint{12.093851in}{0.914466in}}{\pgfqpoint{12.093851in}{0.923706in}}%
\pgfpathlineto{\pgfqpoint{12.093851in}{0.959443in}}%
\pgfpathlineto{\pgfqpoint{12.104208in}{0.959443in}}%
\pgfpathlineto{\pgfqpoint{12.104208in}{0.896400in}}%
\pgfpathlineto{\pgfqpoint{12.093851in}{0.896400in}}%
\pgfpathlineto{\pgfqpoint{12.093851in}{0.906091in}}%
\pgfpathquadraticcurveto{\pgfqpoint{12.090087in}{0.900345in}}{\pgfqpoint{12.085098in}{0.897553in}}%
\pgfpathquadraticcurveto{\pgfqpoint{12.080126in}{0.894761in}}{\pgfqpoint{12.073534in}{0.894761in}}%
\pgfpathquadraticcurveto{\pgfqpoint{12.062672in}{0.894761in}}{\pgfqpoint{12.057035in}{0.901515in}}%
\pgfpathquadraticcurveto{\pgfqpoint{12.051415in}{0.908270in}}{\pgfqpoint{12.051415in}{0.921275in}}%
\pgfpathlineto{\pgfqpoint{12.051415in}{0.921275in}}%
\pgfpathclose%
\pgfpathmoveto{\pgfqpoint{12.077478in}{0.960956in}}%
\pgfpathlineto{\pgfqpoint{12.077478in}{0.960956in}}%
\pgfpathlineto{\pgfqpoint{12.077478in}{0.960956in}}%
\pgfpathclose%
\pgfpathmoveto{\pgfqpoint{12.135550in}{0.905856in}}%
\pgfpathlineto{\pgfqpoint{12.135550in}{0.872426in}}%
\pgfpathlineto{\pgfqpoint{12.125139in}{0.872426in}}%
\pgfpathlineto{\pgfqpoint{12.125139in}{0.959443in}}%
\pgfpathlineto{\pgfqpoint{12.135550in}{0.959443in}}%
\pgfpathlineto{\pgfqpoint{12.135550in}{0.949878in}}%
\pgfpathquadraticcurveto{\pgfqpoint{12.138828in}{0.955498in}}{\pgfqpoint{12.143799in}{0.958218in}}%
\pgfpathquadraticcurveto{\pgfqpoint{12.148788in}{0.960956in}}{\pgfqpoint{12.155705in}{0.960956in}}%
\pgfpathquadraticcurveto{\pgfqpoint{12.167197in}{0.960956in}}{\pgfqpoint{12.174366in}{0.951841in}}%
\pgfpathquadraticcurveto{\pgfqpoint{12.181553in}{0.942727in}}{\pgfqpoint{12.181553in}{0.927867in}}%
\pgfpathquadraticcurveto{\pgfqpoint{12.181553in}{0.913007in}}{\pgfqpoint{12.174366in}{0.903875in}}%
\pgfpathquadraticcurveto{\pgfqpoint{12.167197in}{0.894761in}}{\pgfqpoint{12.155705in}{0.894761in}}%
\pgfpathquadraticcurveto{\pgfqpoint{12.148788in}{0.894761in}}{\pgfqpoint{12.143799in}{0.897499in}}%
\pgfpathquadraticcurveto{\pgfqpoint{12.138828in}{0.900237in}}{\pgfqpoint{12.135550in}{0.905856in}}%
\pgfpathlineto{\pgfqpoint{12.135550in}{0.905856in}}%
\pgfpathclose%
\pgfpathmoveto{\pgfqpoint{12.170799in}{0.927867in}}%
\pgfpathquadraticcurveto{\pgfqpoint{12.170799in}{0.939287in}}{\pgfqpoint{12.166098in}{0.945789in}}%
\pgfpathquadraticcurveto{\pgfqpoint{12.161397in}{0.952292in}}{\pgfqpoint{12.153183in}{0.952292in}}%
\pgfpathquadraticcurveto{\pgfqpoint{12.144952in}{0.952292in}}{\pgfqpoint{12.140251in}{0.945789in}}%
\pgfpathquadraticcurveto{\pgfqpoint{12.135550in}{0.939287in}}{\pgfqpoint{12.135550in}{0.927867in}}%
\pgfpathquadraticcurveto{\pgfqpoint{12.135550in}{0.916448in}}{\pgfqpoint{12.140251in}{0.909945in}}%
\pgfpathquadraticcurveto{\pgfqpoint{12.144952in}{0.903443in}}{\pgfqpoint{12.153183in}{0.903443in}}%
\pgfpathquadraticcurveto{\pgfqpoint{12.161397in}{0.903443in}}{\pgfqpoint{12.166098in}{0.909945in}}%
\pgfpathquadraticcurveto{\pgfqpoint{12.170799in}{0.916448in}}{\pgfqpoint{12.170799in}{0.927867in}}%
\pgfpathlineto{\pgfqpoint{12.170799in}{0.927867in}}%
\pgfpathclose%
\pgfpathmoveto{\pgfqpoint{12.208733in}{0.905856in}}%
\pgfpathlineto{\pgfqpoint{12.208733in}{0.872426in}}%
\pgfpathlineto{\pgfqpoint{12.198322in}{0.872426in}}%
\pgfpathlineto{\pgfqpoint{12.198322in}{0.959443in}}%
\pgfpathlineto{\pgfqpoint{12.208733in}{0.959443in}}%
\pgfpathlineto{\pgfqpoint{12.208733in}{0.949878in}}%
\pgfpathquadraticcurveto{\pgfqpoint{12.212011in}{0.955498in}}{\pgfqpoint{12.216982in}{0.958218in}}%
\pgfpathquadraticcurveto{\pgfqpoint{12.221972in}{0.960956in}}{\pgfqpoint{12.228889in}{0.960956in}}%
\pgfpathquadraticcurveto{\pgfqpoint{12.240380in}{0.960956in}}{\pgfqpoint{12.247549in}{0.951841in}}%
\pgfpathquadraticcurveto{\pgfqpoint{12.254736in}{0.942727in}}{\pgfqpoint{12.254736in}{0.927867in}}%
\pgfpathquadraticcurveto{\pgfqpoint{12.254736in}{0.913007in}}{\pgfqpoint{12.247549in}{0.903875in}}%
\pgfpathquadraticcurveto{\pgfqpoint{12.240380in}{0.894761in}}{\pgfqpoint{12.228889in}{0.894761in}}%
\pgfpathquadraticcurveto{\pgfqpoint{12.221972in}{0.894761in}}{\pgfqpoint{12.216982in}{0.897499in}}%
\pgfpathquadraticcurveto{\pgfqpoint{12.212011in}{0.900237in}}{\pgfqpoint{12.208733in}{0.905856in}}%
\pgfpathlineto{\pgfqpoint{12.208733in}{0.905856in}}%
\pgfpathclose%
\pgfpathmoveto{\pgfqpoint{12.243983in}{0.927867in}}%
\pgfpathquadraticcurveto{\pgfqpoint{12.243983in}{0.939287in}}{\pgfqpoint{12.239282in}{0.945789in}}%
\pgfpathquadraticcurveto{\pgfqpoint{12.234580in}{0.952292in}}{\pgfqpoint{12.226367in}{0.952292in}}%
\pgfpathquadraticcurveto{\pgfqpoint{12.218135in}{0.952292in}}{\pgfqpoint{12.213434in}{0.945789in}}%
\pgfpathquadraticcurveto{\pgfqpoint{12.208733in}{0.939287in}}{\pgfqpoint{12.208733in}{0.927867in}}%
\pgfpathquadraticcurveto{\pgfqpoint{12.208733in}{0.916448in}}{\pgfqpoint{12.213434in}{0.909945in}}%
\pgfpathquadraticcurveto{\pgfqpoint{12.218135in}{0.903443in}}{\pgfqpoint{12.226367in}{0.903443in}}%
\pgfpathquadraticcurveto{\pgfqpoint{12.234580in}{0.903443in}}{\pgfqpoint{12.239282in}{0.909945in}}%
\pgfpathquadraticcurveto{\pgfqpoint{12.243983in}{0.916448in}}{\pgfqpoint{12.243983in}{0.927867in}}%
\pgfpathlineto{\pgfqpoint{12.243983in}{0.927867in}}%
\pgfpathclose%
\pgfpathmoveto{\pgfqpoint{12.308430in}{0.949770in}}%
\pgfpathquadraticcurveto{\pgfqpoint{12.306683in}{0.950779in}}{\pgfqpoint{12.304630in}{0.951247in}}%
\pgfpathquadraticcurveto{\pgfqpoint{12.302576in}{0.951733in}}{\pgfqpoint{12.300109in}{0.951733in}}%
\pgfpathquadraticcurveto{\pgfqpoint{12.291319in}{0.951733in}}{\pgfqpoint{12.286617in}{0.946023in}}%
\pgfpathquadraticcurveto{\pgfqpoint{12.281916in}{0.940314in}}{\pgfqpoint{12.281916in}{0.929614in}}%
\pgfpathlineto{\pgfqpoint{12.281916in}{0.896400in}}%
\pgfpathlineto{\pgfqpoint{12.271505in}{0.896400in}}%
\pgfpathlineto{\pgfqpoint{12.271505in}{0.959443in}}%
\pgfpathlineto{\pgfqpoint{12.281916in}{0.959443in}}%
\pgfpathlineto{\pgfqpoint{12.281916in}{0.949644in}}%
\pgfpathquadraticcurveto{\pgfqpoint{12.285195in}{0.955390in}}{\pgfqpoint{12.290418in}{0.958164in}}%
\pgfpathquadraticcurveto{\pgfqpoint{12.295660in}{0.960956in}}{\pgfqpoint{12.303153in}{0.960956in}}%
\pgfpathquadraticcurveto{\pgfqpoint{12.304215in}{0.960956in}}{\pgfqpoint{12.305512in}{0.960811in}}%
\pgfpathquadraticcurveto{\pgfqpoint{12.306809in}{0.960685in}}{\pgfqpoint{12.308376in}{0.960397in}}%
\pgfpathlineto{\pgfqpoint{12.308430in}{0.949770in}}%
\pgfpathlineto{\pgfqpoint{12.308430in}{0.949770in}}%
\pgfpathclose%
\pgfpathmoveto{\pgfqpoint{12.370680in}{0.930515in}}%
\pgfpathlineto{\pgfqpoint{12.370680in}{0.925454in}}%
\pgfpathlineto{\pgfqpoint{12.323056in}{0.925454in}}%
\pgfpathquadraticcurveto{\pgfqpoint{12.323741in}{0.914754in}}{\pgfqpoint{12.329504in}{0.909153in}}%
\pgfpathquadraticcurveto{\pgfqpoint{12.335268in}{0.903551in}}{\pgfqpoint{12.345571in}{0.903551in}}%
\pgfpathquadraticcurveto{\pgfqpoint{12.351533in}{0.903551in}}{\pgfqpoint{12.357135in}{0.905010in}}%
\pgfpathquadraticcurveto{\pgfqpoint{12.362737in}{0.906469in}}{\pgfqpoint{12.368267in}{0.909405in}}%
\pgfpathlineto{\pgfqpoint{12.368267in}{0.899606in}}%
\pgfpathquadraticcurveto{\pgfqpoint{12.362683in}{0.897247in}}{\pgfqpoint{12.356829in}{0.896004in}}%
\pgfpathquadraticcurveto{\pgfqpoint{12.350975in}{0.894761in}}{\pgfqpoint{12.344959in}{0.894761in}}%
\pgfpathquadraticcurveto{\pgfqpoint{12.329865in}{0.894761in}}{\pgfqpoint{12.321057in}{0.903533in}}%
\pgfpathquadraticcurveto{\pgfqpoint{12.312249in}{0.912323in}}{\pgfqpoint{12.312249in}{0.927309in}}%
\pgfpathquadraticcurveto{\pgfqpoint{12.312249in}{0.942781in}}{\pgfqpoint{12.320606in}{0.951859in}}%
\pgfpathquadraticcurveto{\pgfqpoint{12.328964in}{0.960956in}}{\pgfqpoint{12.343158in}{0.960956in}}%
\pgfpathquadraticcurveto{\pgfqpoint{12.355874in}{0.960956in}}{\pgfqpoint{12.363277in}{0.952760in}}%
\pgfpathquadraticcurveto{\pgfqpoint{12.370680in}{0.944583in}}{\pgfqpoint{12.370680in}{0.930515in}}%
\pgfpathlineto{\pgfqpoint{12.370680in}{0.930515in}}%
\pgfpathclose%
\pgfpathmoveto{\pgfqpoint{12.360323in}{0.933559in}}%
\pgfpathquadraticcurveto{\pgfqpoint{12.360215in}{0.942043in}}{\pgfqpoint{12.355568in}{0.947104in}}%
\pgfpathquadraticcurveto{\pgfqpoint{12.350921in}{0.952184in}}{\pgfqpoint{12.343266in}{0.952184in}}%
\pgfpathquadraticcurveto{\pgfqpoint{12.334602in}{0.952184in}}{\pgfqpoint{12.329396in}{0.947284in}}%
\pgfpathquadraticcurveto{\pgfqpoint{12.324191in}{0.942385in}}{\pgfqpoint{12.323398in}{0.933487in}}%
\pgfpathlineto{\pgfqpoint{12.360323in}{0.933559in}}%
\pgfpathlineto{\pgfqpoint{12.360323in}{0.933559in}}%
\pgfpathclose%
\pgfpathmoveto{\pgfqpoint{12.427869in}{0.957587in}}%
\pgfpathlineto{\pgfqpoint{12.427869in}{0.947789in}}%
\pgfpathquadraticcurveto{\pgfqpoint{12.423492in}{0.950040in}}{\pgfqpoint{12.418755in}{0.951157in}}%
\pgfpathquadraticcurveto{\pgfqpoint{12.414035in}{0.952292in}}{\pgfqpoint{12.408956in}{0.952292in}}%
\pgfpathquadraticcurveto{\pgfqpoint{12.401247in}{0.952292in}}{\pgfqpoint{12.397392in}{0.949932in}}%
\pgfpathquadraticcurveto{\pgfqpoint{12.393538in}{0.947573in}}{\pgfqpoint{12.393538in}{0.942835in}}%
\pgfpathquadraticcurveto{\pgfqpoint{12.393538in}{0.939233in}}{\pgfqpoint{12.396293in}{0.937180in}}%
\pgfpathquadraticcurveto{\pgfqpoint{12.399049in}{0.935126in}}{\pgfqpoint{12.407389in}{0.933271in}}%
\pgfpathlineto{\pgfqpoint{12.410937in}{0.932478in}}%
\pgfpathquadraticcurveto{\pgfqpoint{12.421961in}{0.930119in}}{\pgfqpoint{12.426608in}{0.925814in}}%
\pgfpathquadraticcurveto{\pgfqpoint{12.431255in}{0.921509in}}{\pgfqpoint{12.431255in}{0.913800in}}%
\pgfpathquadraticcurveto{\pgfqpoint{12.431255in}{0.905010in}}{\pgfqpoint{12.424302in}{0.899876in}}%
\pgfpathquadraticcurveto{\pgfqpoint{12.417350in}{0.894761in}}{\pgfqpoint{12.405191in}{0.894761in}}%
\pgfpathquadraticcurveto{\pgfqpoint{12.400130in}{0.894761in}}{\pgfqpoint{12.394636in}{0.895752in}}%
\pgfpathquadraticcurveto{\pgfqpoint{12.389143in}{0.896742in}}{\pgfqpoint{12.383073in}{0.898706in}}%
\pgfpathlineto{\pgfqpoint{12.383073in}{0.909405in}}%
\pgfpathquadraticcurveto{\pgfqpoint{12.388818in}{0.906415in}}{\pgfqpoint{12.394384in}{0.904920in}}%
\pgfpathquadraticcurveto{\pgfqpoint{12.399950in}{0.903443in}}{\pgfqpoint{12.405426in}{0.903443in}}%
\pgfpathquadraticcurveto{\pgfqpoint{12.412739in}{0.903443in}}{\pgfqpoint{12.416665in}{0.905946in}}%
\pgfpathquadraticcurveto{\pgfqpoint{12.420610in}{0.908450in}}{\pgfqpoint{12.420610in}{0.913007in}}%
\pgfpathquadraticcurveto{\pgfqpoint{12.420610in}{0.917222in}}{\pgfqpoint{12.417764in}{0.919474in}}%
\pgfpathquadraticcurveto{\pgfqpoint{12.414936in}{0.921725in}}{\pgfqpoint{12.405300in}{0.923814in}}%
\pgfpathlineto{\pgfqpoint{12.401697in}{0.924661in}}%
\pgfpathquadraticcurveto{\pgfqpoint{12.392079in}{0.926678in}}{\pgfqpoint{12.387792in}{0.930875in}}%
\pgfpathquadraticcurveto{\pgfqpoint{12.383523in}{0.935072in}}{\pgfqpoint{12.383523in}{0.942385in}}%
\pgfpathquadraticcurveto{\pgfqpoint{12.383523in}{0.951283in}}{\pgfqpoint{12.389827in}{0.956110in}}%
\pgfpathquadraticcurveto{\pgfqpoint{12.396131in}{0.960956in}}{\pgfqpoint{12.407731in}{0.960956in}}%
\pgfpathquadraticcurveto{\pgfqpoint{12.413459in}{0.960956in}}{\pgfqpoint{12.418520in}{0.960109in}}%
\pgfpathquadraticcurveto{\pgfqpoint{12.423600in}{0.959280in}}{\pgfqpoint{12.427869in}{0.957587in}}%
\pgfpathlineto{\pgfqpoint{12.427869in}{0.957587in}}%
\pgfpathclose%
\pgfpathmoveto{\pgfqpoint{12.487921in}{0.957587in}}%
\pgfpathlineto{\pgfqpoint{12.487921in}{0.947789in}}%
\pgfpathquadraticcurveto{\pgfqpoint{12.483544in}{0.950040in}}{\pgfqpoint{12.478807in}{0.951157in}}%
\pgfpathquadraticcurveto{\pgfqpoint{12.474088in}{0.952292in}}{\pgfqpoint{12.469009in}{0.952292in}}%
\pgfpathquadraticcurveto{\pgfqpoint{12.461299in}{0.952292in}}{\pgfqpoint{12.457445in}{0.949932in}}%
\pgfpathquadraticcurveto{\pgfqpoint{12.453590in}{0.947573in}}{\pgfqpoint{12.453590in}{0.942835in}}%
\pgfpathquadraticcurveto{\pgfqpoint{12.453590in}{0.939233in}}{\pgfqpoint{12.456346in}{0.937180in}}%
\pgfpathquadraticcurveto{\pgfqpoint{12.459102in}{0.935126in}}{\pgfqpoint{12.467441in}{0.933271in}}%
\pgfpathlineto{\pgfqpoint{12.470990in}{0.932478in}}%
\pgfpathquadraticcurveto{\pgfqpoint{12.482013in}{0.930119in}}{\pgfqpoint{12.486660in}{0.925814in}}%
\pgfpathquadraticcurveto{\pgfqpoint{12.491308in}{0.921509in}}{\pgfqpoint{12.491308in}{0.913800in}}%
\pgfpathquadraticcurveto{\pgfqpoint{12.491308in}{0.905010in}}{\pgfqpoint{12.484355in}{0.899876in}}%
\pgfpathquadraticcurveto{\pgfqpoint{12.477402in}{0.894761in}}{\pgfqpoint{12.465244in}{0.894761in}}%
\pgfpathquadraticcurveto{\pgfqpoint{12.460183in}{0.894761in}}{\pgfqpoint{12.454689in}{0.895752in}}%
\pgfpathquadraticcurveto{\pgfqpoint{12.449195in}{0.896742in}}{\pgfqpoint{12.443125in}{0.898706in}}%
\pgfpathlineto{\pgfqpoint{12.443125in}{0.909405in}}%
\pgfpathquadraticcurveto{\pgfqpoint{12.448871in}{0.906415in}}{\pgfqpoint{12.454437in}{0.904920in}}%
\pgfpathquadraticcurveto{\pgfqpoint{12.460002in}{0.903443in}}{\pgfqpoint{12.465478in}{0.903443in}}%
\pgfpathquadraticcurveto{\pgfqpoint{12.472791in}{0.903443in}}{\pgfqpoint{12.476718in}{0.905946in}}%
\pgfpathquadraticcurveto{\pgfqpoint{12.480662in}{0.908450in}}{\pgfqpoint{12.480662in}{0.913007in}}%
\pgfpathquadraticcurveto{\pgfqpoint{12.480662in}{0.917222in}}{\pgfqpoint{12.477816in}{0.919474in}}%
\pgfpathquadraticcurveto{\pgfqpoint{12.474989in}{0.921725in}}{\pgfqpoint{12.465352in}{0.923814in}}%
\pgfpathlineto{\pgfqpoint{12.461750in}{0.924661in}}%
\pgfpathquadraticcurveto{\pgfqpoint{12.452131in}{0.926678in}}{\pgfqpoint{12.447844in}{0.930875in}}%
\pgfpathquadraticcurveto{\pgfqpoint{12.443575in}{0.935072in}}{\pgfqpoint{12.443575in}{0.942385in}}%
\pgfpathquadraticcurveto{\pgfqpoint{12.443575in}{0.951283in}}{\pgfqpoint{12.449880in}{0.956110in}}%
\pgfpathquadraticcurveto{\pgfqpoint{12.456184in}{0.960956in}}{\pgfqpoint{12.467784in}{0.960956in}}%
\pgfpathquadraticcurveto{\pgfqpoint{12.473512in}{0.960956in}}{\pgfqpoint{12.478573in}{0.960109in}}%
\pgfpathquadraticcurveto{\pgfqpoint{12.483652in}{0.959280in}}{\pgfqpoint{12.487921in}{0.957587in}}%
\pgfpathlineto{\pgfqpoint{12.487921in}{0.957587in}}%
\pgfpathclose%
\pgfpathmoveto{\pgfqpoint{12.561717in}{0.930515in}}%
\pgfpathlineto{\pgfqpoint{12.561717in}{0.925454in}}%
\pgfpathlineto{\pgfqpoint{12.514093in}{0.925454in}}%
\pgfpathquadraticcurveto{\pgfqpoint{12.514777in}{0.914754in}}{\pgfqpoint{12.520541in}{0.909153in}}%
\pgfpathquadraticcurveto{\pgfqpoint{12.526305in}{0.903551in}}{\pgfqpoint{12.536608in}{0.903551in}}%
\pgfpathquadraticcurveto{\pgfqpoint{12.542570in}{0.903551in}}{\pgfqpoint{12.548172in}{0.905010in}}%
\pgfpathquadraticcurveto{\pgfqpoint{12.553774in}{0.906469in}}{\pgfqpoint{12.559303in}{0.909405in}}%
\pgfpathlineto{\pgfqpoint{12.559303in}{0.899606in}}%
\pgfpathquadraticcurveto{\pgfqpoint{12.553720in}{0.897247in}}{\pgfqpoint{12.547866in}{0.896004in}}%
\pgfpathquadraticcurveto{\pgfqpoint{12.542012in}{0.894761in}}{\pgfqpoint{12.535996in}{0.894761in}}%
\pgfpathquadraticcurveto{\pgfqpoint{12.520902in}{0.894761in}}{\pgfqpoint{12.512094in}{0.903533in}}%
\pgfpathquadraticcurveto{\pgfqpoint{12.503286in}{0.912323in}}{\pgfqpoint{12.503286in}{0.927309in}}%
\pgfpathquadraticcurveto{\pgfqpoint{12.503286in}{0.942781in}}{\pgfqpoint{12.511643in}{0.951859in}}%
\pgfpathquadraticcurveto{\pgfqpoint{12.520001in}{0.960956in}}{\pgfqpoint{12.534195in}{0.960956in}}%
\pgfpathquadraticcurveto{\pgfqpoint{12.546911in}{0.960956in}}{\pgfqpoint{12.554314in}{0.952760in}}%
\pgfpathquadraticcurveto{\pgfqpoint{12.561717in}{0.944583in}}{\pgfqpoint{12.561717in}{0.930515in}}%
\pgfpathlineto{\pgfqpoint{12.561717in}{0.930515in}}%
\pgfpathclose%
\pgfpathmoveto{\pgfqpoint{12.551360in}{0.933559in}}%
\pgfpathquadraticcurveto{\pgfqpoint{12.551252in}{0.942043in}}{\pgfqpoint{12.546605in}{0.947104in}}%
\pgfpathquadraticcurveto{\pgfqpoint{12.541958in}{0.952184in}}{\pgfqpoint{12.534303in}{0.952184in}}%
\pgfpathquadraticcurveto{\pgfqpoint{12.525639in}{0.952184in}}{\pgfqpoint{12.520433in}{0.947284in}}%
\pgfpathquadraticcurveto{\pgfqpoint{12.515228in}{0.942385in}}{\pgfqpoint{12.514435in}{0.933487in}}%
\pgfpathlineto{\pgfqpoint{12.551360in}{0.933559in}}%
\pgfpathlineto{\pgfqpoint{12.551360in}{0.933559in}}%
\pgfpathclose%
\pgfpathmoveto{\pgfqpoint{12.620203in}{0.949878in}}%
\pgfpathlineto{\pgfqpoint{12.620203in}{0.983993in}}%
\pgfpathlineto{\pgfqpoint{12.630560in}{0.983993in}}%
\pgfpathlineto{\pgfqpoint{12.630560in}{0.896400in}}%
\pgfpathlineto{\pgfqpoint{12.620203in}{0.896400in}}%
\pgfpathlineto{\pgfqpoint{12.620203in}{0.905856in}}%
\pgfpathquadraticcurveto{\pgfqpoint{12.616942in}{0.900237in}}{\pgfqpoint{12.611953in}{0.897499in}}%
\pgfpathquadraticcurveto{\pgfqpoint{12.606982in}{0.894761in}}{\pgfqpoint{12.599993in}{0.894761in}}%
\pgfpathquadraticcurveto{\pgfqpoint{12.588573in}{0.894761in}}{\pgfqpoint{12.581386in}{0.903875in}}%
\pgfpathquadraticcurveto{\pgfqpoint{12.574218in}{0.913007in}}{\pgfqpoint{12.574218in}{0.927867in}}%
\pgfpathquadraticcurveto{\pgfqpoint{12.574218in}{0.942727in}}{\pgfqpoint{12.581386in}{0.951841in}}%
\pgfpathquadraticcurveto{\pgfqpoint{12.588573in}{0.960956in}}{\pgfqpoint{12.599993in}{0.960956in}}%
\pgfpathquadraticcurveto{\pgfqpoint{12.606982in}{0.960956in}}{\pgfqpoint{12.611953in}{0.958218in}}%
\pgfpathquadraticcurveto{\pgfqpoint{12.616942in}{0.955498in}}{\pgfqpoint{12.620203in}{0.949878in}}%
\pgfpathlineto{\pgfqpoint{12.620203in}{0.949878in}}%
\pgfpathclose%
\pgfpathmoveto{\pgfqpoint{12.584917in}{0.927867in}}%
\pgfpathquadraticcurveto{\pgfqpoint{12.584917in}{0.916448in}}{\pgfqpoint{12.589618in}{0.909945in}}%
\pgfpathquadraticcurveto{\pgfqpoint{12.594319in}{0.903443in}}{\pgfqpoint{12.602533in}{0.903443in}}%
\pgfpathquadraticcurveto{\pgfqpoint{12.610746in}{0.903443in}}{\pgfqpoint{12.615465in}{0.909945in}}%
\pgfpathquadraticcurveto{\pgfqpoint{12.620203in}{0.916448in}}{\pgfqpoint{12.620203in}{0.927867in}}%
\pgfpathquadraticcurveto{\pgfqpoint{12.620203in}{0.939287in}}{\pgfqpoint{12.615465in}{0.945789in}}%
\pgfpathquadraticcurveto{\pgfqpoint{12.610746in}{0.952292in}}{\pgfqpoint{12.602533in}{0.952292in}}%
\pgfpathquadraticcurveto{\pgfqpoint{12.594319in}{0.952292in}}{\pgfqpoint{12.589618in}{0.945789in}}%
\pgfpathquadraticcurveto{\pgfqpoint{12.584917in}{0.939287in}}{\pgfqpoint{12.584917in}{0.927867in}}%
\pgfpathlineto{\pgfqpoint{12.584917in}{0.927867in}}%
\pgfpathclose%
\pgfpathmoveto{\pgfqpoint{12.653363in}{0.910702in}}%
\pgfpathlineto{\pgfqpoint{12.665251in}{0.910702in}}%
\pgfpathlineto{\pgfqpoint{12.665251in}{0.896400in}}%
\pgfpathlineto{\pgfqpoint{12.653363in}{0.896400in}}%
\pgfpathlineto{\pgfqpoint{12.653363in}{0.910702in}}%
\pgfpathlineto{\pgfqpoint{12.653363in}{0.910702in}}%
\pgfpathclose%
\pgfusepath{fill}%
\end{pgfscope}%
\begin{pgfscope}%
\pgfsetbuttcap%
\pgfsetroundjoin%
\definecolor{currentfill}{rgb}{0.309804,0.372549,0.419608}%
\pgfsetfillcolor{currentfill}%
\pgfsetlinewidth{0.000000pt}%
\definecolor{currentstroke}{rgb}{0.309804,0.372549,0.419608}%
\pgfsetstrokecolor{currentstroke}%
\pgfsetdash{}{0pt}%
\pgfpathmoveto{\pgfqpoint{12.767859in}{0.980445in}}%
\pgfpathlineto{\pgfqpoint{12.779225in}{0.980445in}}%
\pgfpathlineto{\pgfqpoint{12.779225in}{0.905964in}}%
\pgfpathlineto{\pgfqpoint{12.820148in}{0.905964in}}%
\pgfpathlineto{\pgfqpoint{12.820148in}{0.896400in}}%
\pgfpathlineto{\pgfqpoint{12.767859in}{0.896400in}}%
\pgfpathlineto{\pgfqpoint{12.767859in}{0.980445in}}%
\pgfpathlineto{\pgfqpoint{12.767859in}{0.980445in}}%
\pgfpathclose%
\pgfpathmoveto{\pgfqpoint{12.883551in}{0.930515in}}%
\pgfpathlineto{\pgfqpoint{12.883551in}{0.925454in}}%
\pgfpathlineto{\pgfqpoint{12.835927in}{0.925454in}}%
\pgfpathquadraticcurveto{\pgfqpoint{12.836611in}{0.914754in}}{\pgfqpoint{12.842375in}{0.909153in}}%
\pgfpathquadraticcurveto{\pgfqpoint{12.848139in}{0.903551in}}{\pgfqpoint{12.858442in}{0.903551in}}%
\pgfpathquadraticcurveto{\pgfqpoint{12.864404in}{0.903551in}}{\pgfqpoint{12.870006in}{0.905010in}}%
\pgfpathquadraticcurveto{\pgfqpoint{12.875608in}{0.906469in}}{\pgfqpoint{12.881138in}{0.909405in}}%
\pgfpathlineto{\pgfqpoint{12.881138in}{0.899606in}}%
\pgfpathquadraticcurveto{\pgfqpoint{12.875554in}{0.897247in}}{\pgfqpoint{12.869700in}{0.896004in}}%
\pgfpathquadraticcurveto{\pgfqpoint{12.863846in}{0.894761in}}{\pgfqpoint{12.857830in}{0.894761in}}%
\pgfpathquadraticcurveto{\pgfqpoint{12.842736in}{0.894761in}}{\pgfqpoint{12.833928in}{0.903533in}}%
\pgfpathquadraticcurveto{\pgfqpoint{12.825120in}{0.912323in}}{\pgfqpoint{12.825120in}{0.927309in}}%
\pgfpathquadraticcurveto{\pgfqpoint{12.825120in}{0.942781in}}{\pgfqpoint{12.833477in}{0.951859in}}%
\pgfpathquadraticcurveto{\pgfqpoint{12.841835in}{0.960956in}}{\pgfqpoint{12.856029in}{0.960956in}}%
\pgfpathquadraticcurveto{\pgfqpoint{12.868745in}{0.960956in}}{\pgfqpoint{12.876148in}{0.952760in}}%
\pgfpathquadraticcurveto{\pgfqpoint{12.883551in}{0.944583in}}{\pgfqpoint{12.883551in}{0.930515in}}%
\pgfpathlineto{\pgfqpoint{12.883551in}{0.930515in}}%
\pgfpathclose%
\pgfpathmoveto{\pgfqpoint{12.873194in}{0.933559in}}%
\pgfpathquadraticcurveto{\pgfqpoint{12.873086in}{0.942043in}}{\pgfqpoint{12.868439in}{0.947104in}}%
\pgfpathquadraticcurveto{\pgfqpoint{12.863792in}{0.952184in}}{\pgfqpoint{12.856137in}{0.952184in}}%
\pgfpathquadraticcurveto{\pgfqpoint{12.847473in}{0.952184in}}{\pgfqpoint{12.842267in}{0.947284in}}%
\pgfpathquadraticcurveto{\pgfqpoint{12.837062in}{0.942385in}}{\pgfqpoint{12.836269in}{0.933487in}}%
\pgfpathlineto{\pgfqpoint{12.873194in}{0.933559in}}%
\pgfpathlineto{\pgfqpoint{12.873194in}{0.933559in}}%
\pgfpathclose%
\pgfpathmoveto{\pgfqpoint{12.893134in}{0.959443in}}%
\pgfpathlineto{\pgfqpoint{12.904103in}{0.959443in}}%
\pgfpathlineto{\pgfqpoint{12.923808in}{0.906541in}}%
\pgfpathlineto{\pgfqpoint{12.943514in}{0.959443in}}%
\pgfpathlineto{\pgfqpoint{12.954483in}{0.959443in}}%
\pgfpathlineto{\pgfqpoint{12.930833in}{0.896400in}}%
\pgfpathlineto{\pgfqpoint{12.916766in}{0.896400in}}%
\pgfpathlineto{\pgfqpoint{12.893134in}{0.959443in}}%
\pgfpathlineto{\pgfqpoint{12.893134in}{0.959443in}}%
\pgfpathclose%
\pgfpathmoveto{\pgfqpoint{13.022713in}{0.930515in}}%
\pgfpathlineto{\pgfqpoint{13.022713in}{0.925454in}}%
\pgfpathlineto{\pgfqpoint{12.975089in}{0.925454in}}%
\pgfpathquadraticcurveto{\pgfqpoint{12.975773in}{0.914754in}}{\pgfqpoint{12.981537in}{0.909153in}}%
\pgfpathquadraticcurveto{\pgfqpoint{12.987301in}{0.903551in}}{\pgfqpoint{12.997604in}{0.903551in}}%
\pgfpathquadraticcurveto{\pgfqpoint{13.003566in}{0.903551in}}{\pgfqpoint{13.009168in}{0.905010in}}%
\pgfpathquadraticcurveto{\pgfqpoint{13.014770in}{0.906469in}}{\pgfqpoint{13.020299in}{0.909405in}}%
\pgfpathlineto{\pgfqpoint{13.020299in}{0.899606in}}%
\pgfpathquadraticcurveto{\pgfqpoint{13.014716in}{0.897247in}}{\pgfqpoint{13.008862in}{0.896004in}}%
\pgfpathquadraticcurveto{\pgfqpoint{13.003008in}{0.894761in}}{\pgfqpoint{12.996992in}{0.894761in}}%
\pgfpathquadraticcurveto{\pgfqpoint{12.981897in}{0.894761in}}{\pgfqpoint{12.973090in}{0.903533in}}%
\pgfpathquadraticcurveto{\pgfqpoint{12.964282in}{0.912323in}}{\pgfqpoint{12.964282in}{0.927309in}}%
\pgfpathquadraticcurveto{\pgfqpoint{12.964282in}{0.942781in}}{\pgfqpoint{12.972639in}{0.951859in}}%
\pgfpathquadraticcurveto{\pgfqpoint{12.980997in}{0.960956in}}{\pgfqpoint{12.995190in}{0.960956in}}%
\pgfpathquadraticcurveto{\pgfqpoint{13.007907in}{0.960956in}}{\pgfqpoint{13.015310in}{0.952760in}}%
\pgfpathquadraticcurveto{\pgfqpoint{13.022713in}{0.944583in}}{\pgfqpoint{13.022713in}{0.930515in}}%
\pgfpathlineto{\pgfqpoint{13.022713in}{0.930515in}}%
\pgfpathclose%
\pgfpathmoveto{\pgfqpoint{13.012356in}{0.933559in}}%
\pgfpathquadraticcurveto{\pgfqpoint{13.012248in}{0.942043in}}{\pgfqpoint{13.007601in}{0.947104in}}%
\pgfpathquadraticcurveto{\pgfqpoint{13.002954in}{0.952184in}}{\pgfqpoint{12.995299in}{0.952184in}}%
\pgfpathquadraticcurveto{\pgfqpoint{12.986635in}{0.952184in}}{\pgfqpoint{12.981429in}{0.947284in}}%
\pgfpathquadraticcurveto{\pgfqpoint{12.976224in}{0.942385in}}{\pgfqpoint{12.975431in}{0.933487in}}%
\pgfpathlineto{\pgfqpoint{13.012356in}{0.933559in}}%
\pgfpathlineto{\pgfqpoint{13.012356in}{0.933559in}}%
\pgfpathclose%
\pgfpathmoveto{\pgfqpoint{13.076245in}{0.949770in}}%
\pgfpathquadraticcurveto{\pgfqpoint{13.074498in}{0.950779in}}{\pgfqpoint{13.072445in}{0.951247in}}%
\pgfpathquadraticcurveto{\pgfqpoint{13.070391in}{0.951733in}}{\pgfqpoint{13.067924in}{0.951733in}}%
\pgfpathquadraticcurveto{\pgfqpoint{13.059134in}{0.951733in}}{\pgfqpoint{13.054432in}{0.946023in}}%
\pgfpathquadraticcurveto{\pgfqpoint{13.049731in}{0.940314in}}{\pgfqpoint{13.049731in}{0.929614in}}%
\pgfpathlineto{\pgfqpoint{13.049731in}{0.896400in}}%
\pgfpathlineto{\pgfqpoint{13.039320in}{0.896400in}}%
\pgfpathlineto{\pgfqpoint{13.039320in}{0.959443in}}%
\pgfpathlineto{\pgfqpoint{13.049731in}{0.959443in}}%
\pgfpathlineto{\pgfqpoint{13.049731in}{0.949644in}}%
\pgfpathquadraticcurveto{\pgfqpoint{13.053009in}{0.955390in}}{\pgfqpoint{13.058233in}{0.958164in}}%
\pgfpathquadraticcurveto{\pgfqpoint{13.063475in}{0.960956in}}{\pgfqpoint{13.070968in}{0.960956in}}%
\pgfpathquadraticcurveto{\pgfqpoint{13.072030in}{0.960956in}}{\pgfqpoint{13.073327in}{0.960811in}}%
\pgfpathquadraticcurveto{\pgfqpoint{13.074624in}{0.960685in}}{\pgfqpoint{13.076191in}{0.960397in}}%
\pgfpathlineto{\pgfqpoint{13.076245in}{0.949770in}}%
\pgfpathlineto{\pgfqpoint{13.076245in}{0.949770in}}%
\pgfpathclose%
\pgfpathmoveto{\pgfqpoint{13.115764in}{0.928083in}}%
\pgfpathquadraticcurveto{\pgfqpoint{13.103209in}{0.928083in}}{\pgfqpoint{13.098364in}{0.925219in}}%
\pgfpathquadraticcurveto{\pgfqpoint{13.093519in}{0.922356in}}{\pgfqpoint{13.093519in}{0.915421in}}%
\pgfpathquadraticcurveto{\pgfqpoint{13.093519in}{0.909909in}}{\pgfqpoint{13.097157in}{0.906667in}}%
\pgfpathquadraticcurveto{\pgfqpoint{13.100796in}{0.903443in}}{\pgfqpoint{13.107028in}{0.903443in}}%
\pgfpathquadraticcurveto{\pgfqpoint{13.115656in}{0.903443in}}{\pgfqpoint{13.120861in}{0.909549in}}%
\pgfpathquadraticcurveto{\pgfqpoint{13.126067in}{0.915655in}}{\pgfqpoint{13.126067in}{0.925778in}}%
\pgfpathlineto{\pgfqpoint{13.126067in}{0.928083in}}%
\pgfpathlineto{\pgfqpoint{13.115764in}{0.928083in}}%
\pgfpathlineto{\pgfqpoint{13.115764in}{0.928083in}}%
\pgfpathclose%
\pgfpathmoveto{\pgfqpoint{13.136424in}{0.932370in}}%
\pgfpathlineto{\pgfqpoint{13.136424in}{0.896400in}}%
\pgfpathlineto{\pgfqpoint{13.126067in}{0.896400in}}%
\pgfpathlineto{\pgfqpoint{13.126067in}{0.905964in}}%
\pgfpathquadraticcurveto{\pgfqpoint{13.122518in}{0.900237in}}{\pgfqpoint{13.117223in}{0.897499in}}%
\pgfpathquadraticcurveto{\pgfqpoint{13.111927in}{0.894761in}}{\pgfqpoint{13.104272in}{0.894761in}}%
\pgfpathquadraticcurveto{\pgfqpoint{13.094600in}{0.894761in}}{\pgfqpoint{13.088872in}{0.900201in}}%
\pgfpathquadraticcurveto{\pgfqpoint{13.083162in}{0.905640in}}{\pgfqpoint{13.083162in}{0.914754in}}%
\pgfpathquadraticcurveto{\pgfqpoint{13.083162in}{0.925382in}}{\pgfqpoint{13.090277in}{0.930785in}}%
\pgfpathquadraticcurveto{\pgfqpoint{13.097409in}{0.936189in}}{\pgfqpoint{13.111531in}{0.936189in}}%
\pgfpathlineto{\pgfqpoint{13.126067in}{0.936189in}}%
\pgfpathlineto{\pgfqpoint{13.126067in}{0.937216in}}%
\pgfpathquadraticcurveto{\pgfqpoint{13.126067in}{0.944366in}}{\pgfqpoint{13.121366in}{0.948275in}}%
\pgfpathquadraticcurveto{\pgfqpoint{13.116664in}{0.952184in}}{\pgfqpoint{13.108163in}{0.952184in}}%
\pgfpathquadraticcurveto{\pgfqpoint{13.102759in}{0.952184in}}{\pgfqpoint{13.097626in}{0.950887in}}%
\pgfpathquadraticcurveto{\pgfqpoint{13.092510in}{0.949590in}}{\pgfqpoint{13.087791in}{0.946996in}}%
\pgfpathlineto{\pgfqpoint{13.087791in}{0.956579in}}%
\pgfpathquadraticcurveto{\pgfqpoint{13.093465in}{0.958776in}}{\pgfqpoint{13.098814in}{0.959857in}}%
\pgfpathquadraticcurveto{\pgfqpoint{13.104164in}{0.960956in}}{\pgfqpoint{13.109225in}{0.960956in}}%
\pgfpathquadraticcurveto{\pgfqpoint{13.122915in}{0.960956in}}{\pgfqpoint{13.129669in}{0.953859in}}%
\pgfpathquadraticcurveto{\pgfqpoint{13.136424in}{0.946780in}}{\pgfqpoint{13.136424in}{0.932370in}}%
\pgfpathlineto{\pgfqpoint{13.136424in}{0.932370in}}%
\pgfpathclose%
\pgfpathmoveto{\pgfqpoint{13.199232in}{0.928660in}}%
\pgfpathquadraticcurveto{\pgfqpoint{13.199232in}{0.939917in}}{\pgfqpoint{13.194585in}{0.946096in}}%
\pgfpathquadraticcurveto{\pgfqpoint{13.189956in}{0.952292in}}{\pgfqpoint{13.181562in}{0.952292in}}%
\pgfpathquadraticcurveto{\pgfqpoint{13.173241in}{0.952292in}}{\pgfqpoint{13.168593in}{0.946096in}}%
\pgfpathquadraticcurveto{\pgfqpoint{13.163946in}{0.939917in}}{\pgfqpoint{13.163946in}{0.928660in}}%
\pgfpathquadraticcurveto{\pgfqpoint{13.163946in}{0.917456in}}{\pgfqpoint{13.168593in}{0.911260in}}%
\pgfpathquadraticcurveto{\pgfqpoint{13.173241in}{0.905064in}}{\pgfqpoint{13.181562in}{0.905064in}}%
\pgfpathquadraticcurveto{\pgfqpoint{13.189956in}{0.905064in}}{\pgfqpoint{13.194585in}{0.911260in}}%
\pgfpathquadraticcurveto{\pgfqpoint{13.199232in}{0.917456in}}{\pgfqpoint{13.199232in}{0.928660in}}%
\pgfpathlineto{\pgfqpoint{13.199232in}{0.928660in}}%
\pgfpathclose%
\pgfpathmoveto{\pgfqpoint{13.209589in}{0.904217in}}%
\pgfpathquadraticcurveto{\pgfqpoint{13.209589in}{0.888132in}}{\pgfqpoint{13.202438in}{0.880279in}}%
\pgfpathquadraticcurveto{\pgfqpoint{13.195305in}{0.872426in}}{\pgfqpoint{13.180554in}{0.872426in}}%
\pgfpathquadraticcurveto{\pgfqpoint{13.175096in}{0.872426in}}{\pgfqpoint{13.170251in}{0.873236in}}%
\pgfpathquadraticcurveto{\pgfqpoint{13.165405in}{0.874047in}}{\pgfqpoint{13.160848in}{0.875740in}}%
\pgfpathlineto{\pgfqpoint{13.160848in}{0.885809in}}%
\pgfpathquadraticcurveto{\pgfqpoint{13.165405in}{0.883341in}}{\pgfqpoint{13.169854in}{0.882170in}}%
\pgfpathquadraticcurveto{\pgfqpoint{13.174303in}{0.880982in}}{\pgfqpoint{13.178914in}{0.880982in}}%
\pgfpathquadraticcurveto{\pgfqpoint{13.189109in}{0.880982in}}{\pgfqpoint{13.194171in}{0.886295in}}%
\pgfpathquadraticcurveto{\pgfqpoint{13.199232in}{0.891609in}}{\pgfqpoint{13.199232in}{0.902362in}}%
\pgfpathlineto{\pgfqpoint{13.199232in}{0.907495in}}%
\pgfpathquadraticcurveto{\pgfqpoint{13.196026in}{0.901912in}}{\pgfqpoint{13.191019in}{0.899156in}}%
\pgfpathquadraticcurveto{\pgfqpoint{13.186011in}{0.896400in}}{\pgfqpoint{13.179022in}{0.896400in}}%
\pgfpathquadraticcurveto{\pgfqpoint{13.167441in}{0.896400in}}{\pgfqpoint{13.160344in}{0.905226in}}%
\pgfpathquadraticcurveto{\pgfqpoint{13.153247in}{0.914070in}}{\pgfqpoint{13.153247in}{0.928660in}}%
\pgfpathquadraticcurveto{\pgfqpoint{13.153247in}{0.943286in}}{\pgfqpoint{13.160344in}{0.952112in}}%
\pgfpathquadraticcurveto{\pgfqpoint{13.167441in}{0.960956in}}{\pgfqpoint{13.179022in}{0.960956in}}%
\pgfpathquadraticcurveto{\pgfqpoint{13.186011in}{0.960956in}}{\pgfqpoint{13.191019in}{0.958200in}}%
\pgfpathquadraticcurveto{\pgfqpoint{13.196026in}{0.955444in}}{\pgfqpoint{13.199232in}{0.949878in}}%
\pgfpathlineto{\pgfqpoint{13.199232in}{0.959443in}}%
\pgfpathlineto{\pgfqpoint{13.209589in}{0.959443in}}%
\pgfpathlineto{\pgfqpoint{13.209589in}{0.904217in}}%
\pgfpathlineto{\pgfqpoint{13.209589in}{0.904217in}}%
\pgfpathclose%
\pgfpathmoveto{\pgfqpoint{13.284862in}{0.930515in}}%
\pgfpathlineto{\pgfqpoint{13.284862in}{0.925454in}}%
\pgfpathlineto{\pgfqpoint{13.237238in}{0.925454in}}%
\pgfpathquadraticcurveto{\pgfqpoint{13.237922in}{0.914754in}}{\pgfqpoint{13.243686in}{0.909153in}}%
\pgfpathquadraticcurveto{\pgfqpoint{13.249450in}{0.903551in}}{\pgfqpoint{13.259753in}{0.903551in}}%
\pgfpathquadraticcurveto{\pgfqpoint{13.265715in}{0.903551in}}{\pgfqpoint{13.271317in}{0.905010in}}%
\pgfpathquadraticcurveto{\pgfqpoint{13.276919in}{0.906469in}}{\pgfqpoint{13.282448in}{0.909405in}}%
\pgfpathlineto{\pgfqpoint{13.282448in}{0.899606in}}%
\pgfpathquadraticcurveto{\pgfqpoint{13.276865in}{0.897247in}}{\pgfqpoint{13.271011in}{0.896004in}}%
\pgfpathquadraticcurveto{\pgfqpoint{13.265157in}{0.894761in}}{\pgfqpoint{13.259141in}{0.894761in}}%
\pgfpathquadraticcurveto{\pgfqpoint{13.244046in}{0.894761in}}{\pgfqpoint{13.235238in}{0.903533in}}%
\pgfpathquadraticcurveto{\pgfqpoint{13.226430in}{0.912323in}}{\pgfqpoint{13.226430in}{0.927309in}}%
\pgfpathquadraticcurveto{\pgfqpoint{13.226430in}{0.942781in}}{\pgfqpoint{13.234788in}{0.951859in}}%
\pgfpathquadraticcurveto{\pgfqpoint{13.243146in}{0.960956in}}{\pgfqpoint{13.257339in}{0.960956in}}%
\pgfpathquadraticcurveto{\pgfqpoint{13.270056in}{0.960956in}}{\pgfqpoint{13.277459in}{0.952760in}}%
\pgfpathquadraticcurveto{\pgfqpoint{13.284862in}{0.944583in}}{\pgfqpoint{13.284862in}{0.930515in}}%
\pgfpathlineto{\pgfqpoint{13.284862in}{0.930515in}}%
\pgfpathclose%
\pgfpathmoveto{\pgfqpoint{13.274505in}{0.933559in}}%
\pgfpathquadraticcurveto{\pgfqpoint{13.274397in}{0.942043in}}{\pgfqpoint{13.269750in}{0.947104in}}%
\pgfpathquadraticcurveto{\pgfqpoint{13.265103in}{0.952184in}}{\pgfqpoint{13.257447in}{0.952184in}}%
\pgfpathquadraticcurveto{\pgfqpoint{13.248784in}{0.952184in}}{\pgfqpoint{13.243578in}{0.947284in}}%
\pgfpathquadraticcurveto{\pgfqpoint{13.238373in}{0.942385in}}{\pgfqpoint{13.237580in}{0.933487in}}%
\pgfpathlineto{\pgfqpoint{13.274505in}{0.933559in}}%
\pgfpathlineto{\pgfqpoint{13.274505in}{0.933559in}}%
\pgfpathclose%
\pgfusepath{fill}%
\end{pgfscope}%
\begin{pgfscope}%
\pgfsetbuttcap%
\pgfsetroundjoin%
\definecolor{currentfill}{rgb}{0.309804,0.372549,0.419608}%
\pgfsetfillcolor{currentfill}%
\pgfsetlinewidth{0.000000pt}%
\definecolor{currentstroke}{rgb}{0.309804,0.372549,0.419608}%
\pgfsetstrokecolor{currentstroke}%
\pgfsetdash{}{0pt}%
\pgfpathmoveto{\pgfqpoint{13.361799in}{0.921275in}}%
\pgfpathlineto{\pgfqpoint{13.361799in}{0.959443in}}%
\pgfpathlineto{\pgfqpoint{13.372156in}{0.959443in}}%
\pgfpathlineto{\pgfqpoint{13.372156in}{0.921671in}}%
\pgfpathquadraticcurveto{\pgfqpoint{13.372156in}{0.912719in}}{\pgfqpoint{13.375632in}{0.908234in}}%
\pgfpathquadraticcurveto{\pgfqpoint{13.379126in}{0.903767in}}{\pgfqpoint{13.386115in}{0.903767in}}%
\pgfpathquadraticcurveto{\pgfqpoint{13.394491in}{0.903767in}}{\pgfqpoint{13.399354in}{0.909117in}}%
\pgfpathquadraticcurveto{\pgfqpoint{13.404235in}{0.914466in}}{\pgfqpoint{13.404235in}{0.923706in}}%
\pgfpathlineto{\pgfqpoint{13.404235in}{0.959443in}}%
\pgfpathlineto{\pgfqpoint{13.414592in}{0.959443in}}%
\pgfpathlineto{\pgfqpoint{13.414592in}{0.896400in}}%
\pgfpathlineto{\pgfqpoint{13.404235in}{0.896400in}}%
\pgfpathlineto{\pgfqpoint{13.404235in}{0.906091in}}%
\pgfpathquadraticcurveto{\pgfqpoint{13.400471in}{0.900345in}}{\pgfqpoint{13.395481in}{0.897553in}}%
\pgfpathquadraticcurveto{\pgfqpoint{13.390510in}{0.894761in}}{\pgfqpoint{13.383918in}{0.894761in}}%
\pgfpathquadraticcurveto{\pgfqpoint{13.373056in}{0.894761in}}{\pgfqpoint{13.367418in}{0.901515in}}%
\pgfpathquadraticcurveto{\pgfqpoint{13.361799in}{0.908270in}}{\pgfqpoint{13.361799in}{0.921275in}}%
\pgfpathlineto{\pgfqpoint{13.361799in}{0.921275in}}%
\pgfpathclose%
\pgfpathmoveto{\pgfqpoint{13.387862in}{0.960956in}}%
\pgfpathlineto{\pgfqpoint{13.387862in}{0.960956in}}%
\pgfpathlineto{\pgfqpoint{13.387862in}{0.960956in}}%
\pgfpathclose%
\pgfpathmoveto{\pgfqpoint{13.476104in}{0.957587in}}%
\pgfpathlineto{\pgfqpoint{13.476104in}{0.947789in}}%
\pgfpathquadraticcurveto{\pgfqpoint{13.471727in}{0.950040in}}{\pgfqpoint{13.466990in}{0.951157in}}%
\pgfpathquadraticcurveto{\pgfqpoint{13.462270in}{0.952292in}}{\pgfqpoint{13.457191in}{0.952292in}}%
\pgfpathquadraticcurveto{\pgfqpoint{13.449482in}{0.952292in}}{\pgfqpoint{13.445627in}{0.949932in}}%
\pgfpathquadraticcurveto{\pgfqpoint{13.441773in}{0.947573in}}{\pgfqpoint{13.441773in}{0.942835in}}%
\pgfpathquadraticcurveto{\pgfqpoint{13.441773in}{0.939233in}}{\pgfqpoint{13.444528in}{0.937180in}}%
\pgfpathquadraticcurveto{\pgfqpoint{13.447284in}{0.935126in}}{\pgfqpoint{13.455624in}{0.933271in}}%
\pgfpathlineto{\pgfqpoint{13.459172in}{0.932478in}}%
\pgfpathquadraticcurveto{\pgfqpoint{13.470196in}{0.930119in}}{\pgfqpoint{13.474843in}{0.925814in}}%
\pgfpathquadraticcurveto{\pgfqpoint{13.479490in}{0.921509in}}{\pgfqpoint{13.479490in}{0.913800in}}%
\pgfpathquadraticcurveto{\pgfqpoint{13.479490in}{0.905010in}}{\pgfqpoint{13.472537in}{0.899876in}}%
\pgfpathquadraticcurveto{\pgfqpoint{13.465585in}{0.894761in}}{\pgfqpoint{13.453426in}{0.894761in}}%
\pgfpathquadraticcurveto{\pgfqpoint{13.448365in}{0.894761in}}{\pgfqpoint{13.442871in}{0.895752in}}%
\pgfpathquadraticcurveto{\pgfqpoint{13.437378in}{0.896742in}}{\pgfqpoint{13.431308in}{0.898706in}}%
\pgfpathlineto{\pgfqpoint{13.431308in}{0.909405in}}%
\pgfpathquadraticcurveto{\pgfqpoint{13.437053in}{0.906415in}}{\pgfqpoint{13.442619in}{0.904920in}}%
\pgfpathquadraticcurveto{\pgfqpoint{13.448185in}{0.903443in}}{\pgfqpoint{13.453661in}{0.903443in}}%
\pgfpathquadraticcurveto{\pgfqpoint{13.460974in}{0.903443in}}{\pgfqpoint{13.464900in}{0.905946in}}%
\pgfpathquadraticcurveto{\pgfqpoint{13.468845in}{0.908450in}}{\pgfqpoint{13.468845in}{0.913007in}}%
\pgfpathquadraticcurveto{\pgfqpoint{13.468845in}{0.917222in}}{\pgfqpoint{13.465999in}{0.919474in}}%
\pgfpathquadraticcurveto{\pgfqpoint{13.463171in}{0.921725in}}{\pgfqpoint{13.453535in}{0.923814in}}%
\pgfpathlineto{\pgfqpoint{13.449932in}{0.924661in}}%
\pgfpathquadraticcurveto{\pgfqpoint{13.440314in}{0.926678in}}{\pgfqpoint{13.436027in}{0.930875in}}%
\pgfpathquadraticcurveto{\pgfqpoint{13.431758in}{0.935072in}}{\pgfqpoint{13.431758in}{0.942385in}}%
\pgfpathquadraticcurveto{\pgfqpoint{13.431758in}{0.951283in}}{\pgfqpoint{13.438062in}{0.956110in}}%
\pgfpathquadraticcurveto{\pgfqpoint{13.444366in}{0.960956in}}{\pgfqpoint{13.455966in}{0.960956in}}%
\pgfpathquadraticcurveto{\pgfqpoint{13.461694in}{0.960956in}}{\pgfqpoint{13.466755in}{0.960109in}}%
\pgfpathquadraticcurveto{\pgfqpoint{13.471835in}{0.959280in}}{\pgfqpoint{13.476104in}{0.957587in}}%
\pgfpathlineto{\pgfqpoint{13.476104in}{0.957587in}}%
\pgfpathclose%
\pgfpathmoveto{\pgfqpoint{13.549900in}{0.930515in}}%
\pgfpathlineto{\pgfqpoint{13.549900in}{0.925454in}}%
\pgfpathlineto{\pgfqpoint{13.502275in}{0.925454in}}%
\pgfpathquadraticcurveto{\pgfqpoint{13.502960in}{0.914754in}}{\pgfqpoint{13.508724in}{0.909153in}}%
\pgfpathquadraticcurveto{\pgfqpoint{13.514488in}{0.903551in}}{\pgfqpoint{13.524791in}{0.903551in}}%
\pgfpathquadraticcurveto{\pgfqpoint{13.530753in}{0.903551in}}{\pgfqpoint{13.536354in}{0.905010in}}%
\pgfpathquadraticcurveto{\pgfqpoint{13.541956in}{0.906469in}}{\pgfqpoint{13.547486in}{0.909405in}}%
\pgfpathlineto{\pgfqpoint{13.547486in}{0.899606in}}%
\pgfpathquadraticcurveto{\pgfqpoint{13.541902in}{0.897247in}}{\pgfqpoint{13.536048in}{0.896004in}}%
\pgfpathquadraticcurveto{\pgfqpoint{13.530194in}{0.894761in}}{\pgfqpoint{13.524178in}{0.894761in}}%
\pgfpathquadraticcurveto{\pgfqpoint{13.509084in}{0.894761in}}{\pgfqpoint{13.500276in}{0.903533in}}%
\pgfpathquadraticcurveto{\pgfqpoint{13.491468in}{0.912323in}}{\pgfqpoint{13.491468in}{0.927309in}}%
\pgfpathquadraticcurveto{\pgfqpoint{13.491468in}{0.942781in}}{\pgfqpoint{13.499826in}{0.951859in}}%
\pgfpathquadraticcurveto{\pgfqpoint{13.508183in}{0.960956in}}{\pgfqpoint{13.522377in}{0.960956in}}%
\pgfpathquadraticcurveto{\pgfqpoint{13.535094in}{0.960956in}}{\pgfqpoint{13.542497in}{0.952760in}}%
\pgfpathquadraticcurveto{\pgfqpoint{13.549900in}{0.944583in}}{\pgfqpoint{13.549900in}{0.930515in}}%
\pgfpathlineto{\pgfqpoint{13.549900in}{0.930515in}}%
\pgfpathclose%
\pgfpathmoveto{\pgfqpoint{13.539543in}{0.933559in}}%
\pgfpathquadraticcurveto{\pgfqpoint{13.539434in}{0.942043in}}{\pgfqpoint{13.534787in}{0.947104in}}%
\pgfpathquadraticcurveto{\pgfqpoint{13.530140in}{0.952184in}}{\pgfqpoint{13.522485in}{0.952184in}}%
\pgfpathquadraticcurveto{\pgfqpoint{13.513821in}{0.952184in}}{\pgfqpoint{13.508616in}{0.947284in}}%
\pgfpathquadraticcurveto{\pgfqpoint{13.503410in}{0.942385in}}{\pgfqpoint{13.502618in}{0.933487in}}%
\pgfpathlineto{\pgfqpoint{13.539543in}{0.933559in}}%
\pgfpathlineto{\pgfqpoint{13.539543in}{0.933559in}}%
\pgfpathclose%
\pgfpathmoveto{\pgfqpoint{13.607088in}{0.957587in}}%
\pgfpathlineto{\pgfqpoint{13.607088in}{0.947789in}}%
\pgfpathquadraticcurveto{\pgfqpoint{13.602711in}{0.950040in}}{\pgfqpoint{13.597974in}{0.951157in}}%
\pgfpathquadraticcurveto{\pgfqpoint{13.593255in}{0.952292in}}{\pgfqpoint{13.588175in}{0.952292in}}%
\pgfpathquadraticcurveto{\pgfqpoint{13.580466in}{0.952292in}}{\pgfqpoint{13.576612in}{0.949932in}}%
\pgfpathquadraticcurveto{\pgfqpoint{13.572757in}{0.947573in}}{\pgfqpoint{13.572757in}{0.942835in}}%
\pgfpathquadraticcurveto{\pgfqpoint{13.572757in}{0.939233in}}{\pgfqpoint{13.575513in}{0.937180in}}%
\pgfpathquadraticcurveto{\pgfqpoint{13.578269in}{0.935126in}}{\pgfqpoint{13.586608in}{0.933271in}}%
\pgfpathlineto{\pgfqpoint{13.590157in}{0.932478in}}%
\pgfpathquadraticcurveto{\pgfqpoint{13.601180in}{0.930119in}}{\pgfqpoint{13.605827in}{0.925814in}}%
\pgfpathquadraticcurveto{\pgfqpoint{13.610474in}{0.921509in}}{\pgfqpoint{13.610474in}{0.913800in}}%
\pgfpathquadraticcurveto{\pgfqpoint{13.610474in}{0.905010in}}{\pgfqpoint{13.603522in}{0.899876in}}%
\pgfpathquadraticcurveto{\pgfqpoint{13.596569in}{0.894761in}}{\pgfqpoint{13.584411in}{0.894761in}}%
\pgfpathquadraticcurveto{\pgfqpoint{13.579349in}{0.894761in}}{\pgfqpoint{13.573856in}{0.895752in}}%
\pgfpathquadraticcurveto{\pgfqpoint{13.568362in}{0.896742in}}{\pgfqpoint{13.562292in}{0.898706in}}%
\pgfpathlineto{\pgfqpoint{13.562292in}{0.909405in}}%
\pgfpathquadraticcurveto{\pgfqpoint{13.568038in}{0.906415in}}{\pgfqpoint{13.573604in}{0.904920in}}%
\pgfpathquadraticcurveto{\pgfqpoint{13.579169in}{0.903443in}}{\pgfqpoint{13.584645in}{0.903443in}}%
\pgfpathquadraticcurveto{\pgfqpoint{13.591958in}{0.903443in}}{\pgfqpoint{13.595885in}{0.905946in}}%
\pgfpathquadraticcurveto{\pgfqpoint{13.599829in}{0.908450in}}{\pgfqpoint{13.599829in}{0.913007in}}%
\pgfpathquadraticcurveto{\pgfqpoint{13.599829in}{0.917222in}}{\pgfqpoint{13.596983in}{0.919474in}}%
\pgfpathquadraticcurveto{\pgfqpoint{13.594155in}{0.921725in}}{\pgfqpoint{13.584519in}{0.923814in}}%
\pgfpathlineto{\pgfqpoint{13.580916in}{0.924661in}}%
\pgfpathquadraticcurveto{\pgfqpoint{13.571298in}{0.926678in}}{\pgfqpoint{13.567011in}{0.930875in}}%
\pgfpathquadraticcurveto{\pgfqpoint{13.562742in}{0.935072in}}{\pgfqpoint{13.562742in}{0.942385in}}%
\pgfpathquadraticcurveto{\pgfqpoint{13.562742in}{0.951283in}}{\pgfqpoint{13.569046in}{0.956110in}}%
\pgfpathquadraticcurveto{\pgfqpoint{13.575351in}{0.960956in}}{\pgfqpoint{13.586951in}{0.960956in}}%
\pgfpathquadraticcurveto{\pgfqpoint{13.592678in}{0.960956in}}{\pgfqpoint{13.597740in}{0.960109in}}%
\pgfpathquadraticcurveto{\pgfqpoint{13.602819in}{0.959280in}}{\pgfqpoint{13.607088in}{0.957587in}}%
\pgfpathlineto{\pgfqpoint{13.607088in}{0.957587in}}%
\pgfpathclose%
\pgfusepath{fill}%
\end{pgfscope}%
\begin{pgfscope}%
\pgfsetbuttcap%
\pgfsetroundjoin%
\definecolor{currentfill}{rgb}{0.309804,0.372549,0.419608}%
\pgfsetfillcolor{currentfill}%
\pgfsetlinewidth{0.000000pt}%
\definecolor{currentstroke}{rgb}{0.309804,0.372549,0.419608}%
\pgfsetstrokecolor{currentstroke}%
\pgfsetdash{}{0pt}%
\pgfpathmoveto{\pgfqpoint{0.975046in}{0.773987in}}%
\pgfpathlineto{\pgfqpoint{0.975046in}{0.764189in}}%
\pgfpathquadraticcurveto{\pgfqpoint{0.970669in}{0.766440in}}{\pgfqpoint{0.965932in}{0.767557in}}%
\pgfpathquadraticcurveto{\pgfqpoint{0.961213in}{0.768692in}}{\pgfqpoint{0.956134in}{0.768692in}}%
\pgfpathquadraticcurveto{\pgfqpoint{0.948424in}{0.768692in}}{\pgfqpoint{0.944570in}{0.766332in}}%
\pgfpathquadraticcurveto{\pgfqpoint{0.940715in}{0.763973in}}{\pgfqpoint{0.940715in}{0.759235in}}%
\pgfpathquadraticcurveto{\pgfqpoint{0.940715in}{0.755633in}}{\pgfqpoint{0.943471in}{0.753580in}}%
\pgfpathquadraticcurveto{\pgfqpoint{0.946227in}{0.751526in}}{\pgfqpoint{0.954567in}{0.749671in}}%
\pgfpathlineto{\pgfqpoint{0.958115in}{0.748878in}}%
\pgfpathquadraticcurveto{\pgfqpoint{0.969138in}{0.746519in}}{\pgfqpoint{0.973786in}{0.742214in}}%
\pgfpathquadraticcurveto{\pgfqpoint{0.978433in}{0.737909in}}{\pgfqpoint{0.978433in}{0.730200in}}%
\pgfpathquadraticcurveto{\pgfqpoint{0.978433in}{0.721410in}}{\pgfqpoint{0.971480in}{0.716276in}}%
\pgfpathquadraticcurveto{\pgfqpoint{0.964527in}{0.711161in}}{\pgfqpoint{0.952369in}{0.711161in}}%
\pgfpathquadraticcurveto{\pgfqpoint{0.947308in}{0.711161in}}{\pgfqpoint{0.941814in}{0.712152in}}%
\pgfpathquadraticcurveto{\pgfqpoint{0.936320in}{0.713142in}}{\pgfqpoint{0.930250in}{0.715106in}}%
\pgfpathlineto{\pgfqpoint{0.930250in}{0.725805in}}%
\pgfpathquadraticcurveto{\pgfqpoint{0.935996in}{0.722815in}}{\pgfqpoint{0.941562in}{0.721320in}}%
\pgfpathquadraticcurveto{\pgfqpoint{0.947128in}{0.719843in}}{\pgfqpoint{0.952603in}{0.719843in}}%
\pgfpathquadraticcurveto{\pgfqpoint{0.959916in}{0.719843in}}{\pgfqpoint{0.963843in}{0.722346in}}%
\pgfpathquadraticcurveto{\pgfqpoint{0.967788in}{0.724850in}}{\pgfqpoint{0.967788in}{0.729407in}}%
\pgfpathquadraticcurveto{\pgfqpoint{0.967788in}{0.733622in}}{\pgfqpoint{0.964942in}{0.735874in}}%
\pgfpathquadraticcurveto{\pgfqpoint{0.962114in}{0.738125in}}{\pgfqpoint{0.952477in}{0.740214in}}%
\pgfpathlineto{\pgfqpoint{0.948875in}{0.741061in}}%
\pgfpathquadraticcurveto{\pgfqpoint{0.939256in}{0.743078in}}{\pgfqpoint{0.934969in}{0.747275in}}%
\pgfpathquadraticcurveto{\pgfqpoint{0.930701in}{0.751472in}}{\pgfqpoint{0.930701in}{0.758785in}}%
\pgfpathquadraticcurveto{\pgfqpoint{0.930701in}{0.767683in}}{\pgfqpoint{0.937005in}{0.772510in}}%
\pgfpathquadraticcurveto{\pgfqpoint{0.943309in}{0.777356in}}{\pgfqpoint{0.954909in}{0.777356in}}%
\pgfpathquadraticcurveto{\pgfqpoint{0.960637in}{0.777356in}}{\pgfqpoint{0.965698in}{0.776509in}}%
\pgfpathquadraticcurveto{\pgfqpoint{0.970778in}{0.775680in}}{\pgfqpoint{0.975046in}{0.773987in}}%
\pgfpathlineto{\pgfqpoint{0.975046in}{0.773987in}}%
\pgfpathclose%
\pgfpathmoveto{\pgfqpoint{1.001110in}{0.744267in}}%
\pgfpathquadraticcurveto{\pgfqpoint{1.001110in}{0.732848in}}{\pgfqpoint{1.005811in}{0.726345in}}%
\pgfpathquadraticcurveto{\pgfqpoint{1.010512in}{0.719843in}}{\pgfqpoint{1.018726in}{0.719843in}}%
\pgfpathquadraticcurveto{\pgfqpoint{1.026939in}{0.719843in}}{\pgfqpoint{1.031659in}{0.726345in}}%
\pgfpathquadraticcurveto{\pgfqpoint{1.036396in}{0.732848in}}{\pgfqpoint{1.036396in}{0.744267in}}%
\pgfpathquadraticcurveto{\pgfqpoint{1.036396in}{0.755687in}}{\pgfqpoint{1.031659in}{0.762189in}}%
\pgfpathquadraticcurveto{\pgfqpoint{1.026939in}{0.768692in}}{\pgfqpoint{1.018726in}{0.768692in}}%
\pgfpathquadraticcurveto{\pgfqpoint{1.010512in}{0.768692in}}{\pgfqpoint{1.005811in}{0.762189in}}%
\pgfpathquadraticcurveto{\pgfqpoint{1.001110in}{0.755687in}}{\pgfqpoint{1.001110in}{0.744267in}}%
\pgfpathlineto{\pgfqpoint{1.001110in}{0.744267in}}%
\pgfpathclose%
\pgfpathmoveto{\pgfqpoint{1.036396in}{0.722256in}}%
\pgfpathquadraticcurveto{\pgfqpoint{1.033136in}{0.716637in}}{\pgfqpoint{1.028146in}{0.713899in}}%
\pgfpathquadraticcurveto{\pgfqpoint{1.023175in}{0.711161in}}{\pgfqpoint{1.016186in}{0.711161in}}%
\pgfpathquadraticcurveto{\pgfqpoint{1.004766in}{0.711161in}}{\pgfqpoint{0.997580in}{0.720275in}}%
\pgfpathquadraticcurveto{\pgfqpoint{0.990411in}{0.729407in}}{\pgfqpoint{0.990411in}{0.744267in}}%
\pgfpathquadraticcurveto{\pgfqpoint{0.990411in}{0.759127in}}{\pgfqpoint{0.997580in}{0.768241in}}%
\pgfpathquadraticcurveto{\pgfqpoint{1.004766in}{0.777356in}}{\pgfqpoint{1.016186in}{0.777356in}}%
\pgfpathquadraticcurveto{\pgfqpoint{1.023175in}{0.777356in}}{\pgfqpoint{1.028146in}{0.774618in}}%
\pgfpathquadraticcurveto{\pgfqpoint{1.033136in}{0.771898in}}{\pgfqpoint{1.036396in}{0.766278in}}%
\pgfpathlineto{\pgfqpoint{1.036396in}{0.775843in}}%
\pgfpathlineto{\pgfqpoint{1.046753in}{0.775843in}}%
\pgfpathlineto{\pgfqpoint{1.046753in}{0.688826in}}%
\pgfpathlineto{\pgfqpoint{1.036396in}{0.688826in}}%
\pgfpathlineto{\pgfqpoint{1.036396in}{0.722256in}}%
\pgfpathlineto{\pgfqpoint{1.036396in}{0.722256in}}%
\pgfpathclose%
\pgfpathmoveto{\pgfqpoint{1.067035in}{0.737675in}}%
\pgfpathlineto{\pgfqpoint{1.067035in}{0.775843in}}%
\pgfpathlineto{\pgfqpoint{1.077391in}{0.775843in}}%
\pgfpathlineto{\pgfqpoint{1.077391in}{0.738071in}}%
\pgfpathquadraticcurveto{\pgfqpoint{1.077391in}{0.729119in}}{\pgfqpoint{1.080868in}{0.724634in}}%
\pgfpathquadraticcurveto{\pgfqpoint{1.084362in}{0.720167in}}{\pgfqpoint{1.091351in}{0.720167in}}%
\pgfpathquadraticcurveto{\pgfqpoint{1.099727in}{0.720167in}}{\pgfqpoint{1.104590in}{0.725517in}}%
\pgfpathquadraticcurveto{\pgfqpoint{1.109471in}{0.730866in}}{\pgfqpoint{1.109471in}{0.740106in}}%
\pgfpathlineto{\pgfqpoint{1.109471in}{0.775843in}}%
\pgfpathlineto{\pgfqpoint{1.119828in}{0.775843in}}%
\pgfpathlineto{\pgfqpoint{1.119828in}{0.712800in}}%
\pgfpathlineto{\pgfqpoint{1.109471in}{0.712800in}}%
\pgfpathlineto{\pgfqpoint{1.109471in}{0.722491in}}%
\pgfpathquadraticcurveto{\pgfqpoint{1.105707in}{0.716745in}}{\pgfqpoint{1.100717in}{0.713953in}}%
\pgfpathquadraticcurveto{\pgfqpoint{1.095746in}{0.711161in}}{\pgfqpoint{1.089153in}{0.711161in}}%
\pgfpathquadraticcurveto{\pgfqpoint{1.078292in}{0.711161in}}{\pgfqpoint{1.072654in}{0.717915in}}%
\pgfpathquadraticcurveto{\pgfqpoint{1.067035in}{0.724670in}}{\pgfqpoint{1.067035in}{0.737675in}}%
\pgfpathlineto{\pgfqpoint{1.067035in}{0.737675in}}%
\pgfpathclose%
\pgfpathmoveto{\pgfqpoint{1.093098in}{0.777356in}}%
\pgfpathlineto{\pgfqpoint{1.093098in}{0.777356in}}%
\pgfpathlineto{\pgfqpoint{1.093098in}{0.777356in}}%
\pgfpathclose%
\pgfpathmoveto{\pgfqpoint{1.169812in}{0.744483in}}%
\pgfpathquadraticcurveto{\pgfqpoint{1.157257in}{0.744483in}}{\pgfqpoint{1.152412in}{0.741619in}}%
\pgfpathquadraticcurveto{\pgfqpoint{1.147567in}{0.738756in}}{\pgfqpoint{1.147567in}{0.731821in}}%
\pgfpathquadraticcurveto{\pgfqpoint{1.147567in}{0.726309in}}{\pgfqpoint{1.151205in}{0.723067in}}%
\pgfpathquadraticcurveto{\pgfqpoint{1.154844in}{0.719843in}}{\pgfqpoint{1.161076in}{0.719843in}}%
\pgfpathquadraticcurveto{\pgfqpoint{1.169704in}{0.719843in}}{\pgfqpoint{1.174909in}{0.725949in}}%
\pgfpathquadraticcurveto{\pgfqpoint{1.180115in}{0.732055in}}{\pgfqpoint{1.180115in}{0.742178in}}%
\pgfpathlineto{\pgfqpoint{1.180115in}{0.744483in}}%
\pgfpathlineto{\pgfqpoint{1.169812in}{0.744483in}}%
\pgfpathlineto{\pgfqpoint{1.169812in}{0.744483in}}%
\pgfpathclose%
\pgfpathmoveto{\pgfqpoint{1.190472in}{0.748770in}}%
\pgfpathlineto{\pgfqpoint{1.190472in}{0.712800in}}%
\pgfpathlineto{\pgfqpoint{1.180115in}{0.712800in}}%
\pgfpathlineto{\pgfqpoint{1.180115in}{0.722364in}}%
\pgfpathquadraticcurveto{\pgfqpoint{1.176566in}{0.716637in}}{\pgfqpoint{1.171271in}{0.713899in}}%
\pgfpathquadraticcurveto{\pgfqpoint{1.165975in}{0.711161in}}{\pgfqpoint{1.158320in}{0.711161in}}%
\pgfpathquadraticcurveto{\pgfqpoint{1.148648in}{0.711161in}}{\pgfqpoint{1.142920in}{0.716601in}}%
\pgfpathquadraticcurveto{\pgfqpoint{1.137210in}{0.722040in}}{\pgfqpoint{1.137210in}{0.731154in}}%
\pgfpathquadraticcurveto{\pgfqpoint{1.137210in}{0.741782in}}{\pgfqpoint{1.144325in}{0.747185in}}%
\pgfpathquadraticcurveto{\pgfqpoint{1.151457in}{0.752589in}}{\pgfqpoint{1.165579in}{0.752589in}}%
\pgfpathlineto{\pgfqpoint{1.180115in}{0.752589in}}%
\pgfpathlineto{\pgfqpoint{1.180115in}{0.753616in}}%
\pgfpathquadraticcurveto{\pgfqpoint{1.180115in}{0.760766in}}{\pgfqpoint{1.175414in}{0.764675in}}%
\pgfpathquadraticcurveto{\pgfqpoint{1.170712in}{0.768584in}}{\pgfqpoint{1.162211in}{0.768584in}}%
\pgfpathquadraticcurveto{\pgfqpoint{1.156807in}{0.768584in}}{\pgfqpoint{1.151674in}{0.767287in}}%
\pgfpathquadraticcurveto{\pgfqpoint{1.146558in}{0.765990in}}{\pgfqpoint{1.141839in}{0.763396in}}%
\pgfpathlineto{\pgfqpoint{1.141839in}{0.772979in}}%
\pgfpathquadraticcurveto{\pgfqpoint{1.147513in}{0.775176in}}{\pgfqpoint{1.152862in}{0.776257in}}%
\pgfpathquadraticcurveto{\pgfqpoint{1.158212in}{0.777356in}}{\pgfqpoint{1.163273in}{0.777356in}}%
\pgfpathquadraticcurveto{\pgfqpoint{1.176963in}{0.777356in}}{\pgfqpoint{1.183717in}{0.770259in}}%
\pgfpathquadraticcurveto{\pgfqpoint{1.190472in}{0.763180in}}{\pgfqpoint{1.190472in}{0.748770in}}%
\pgfpathlineto{\pgfqpoint{1.190472in}{0.748770in}}%
\pgfpathclose%
\pgfpathmoveto{\pgfqpoint{1.248327in}{0.766170in}}%
\pgfpathquadraticcurveto{\pgfqpoint{1.246580in}{0.767179in}}{\pgfqpoint{1.244526in}{0.767647in}}%
\pgfpathquadraticcurveto{\pgfqpoint{1.242473in}{0.768133in}}{\pgfqpoint{1.240005in}{0.768133in}}%
\pgfpathquadraticcurveto{\pgfqpoint{1.231215in}{0.768133in}}{\pgfqpoint{1.226514in}{0.762423in}}%
\pgfpathquadraticcurveto{\pgfqpoint{1.221813in}{0.756714in}}{\pgfqpoint{1.221813in}{0.746014in}}%
\pgfpathlineto{\pgfqpoint{1.221813in}{0.712800in}}%
\pgfpathlineto{\pgfqpoint{1.211402in}{0.712800in}}%
\pgfpathlineto{\pgfqpoint{1.211402in}{0.775843in}}%
\pgfpathlineto{\pgfqpoint{1.221813in}{0.775843in}}%
\pgfpathlineto{\pgfqpoint{1.221813in}{0.766044in}}%
\pgfpathquadraticcurveto{\pgfqpoint{1.225091in}{0.771790in}}{\pgfqpoint{1.230315in}{0.774564in}}%
\pgfpathquadraticcurveto{\pgfqpoint{1.235556in}{0.777356in}}{\pgfqpoint{1.243049in}{0.777356in}}%
\pgfpathquadraticcurveto{\pgfqpoint{1.244112in}{0.777356in}}{\pgfqpoint{1.245409in}{0.777211in}}%
\pgfpathquadraticcurveto{\pgfqpoint{1.246706in}{0.777085in}}{\pgfqpoint{1.248273in}{0.776797in}}%
\pgfpathlineto{\pgfqpoint{1.248327in}{0.766170in}}%
\pgfpathlineto{\pgfqpoint{1.248327in}{0.766170in}}%
\pgfpathclose%
\pgfpathmoveto{\pgfqpoint{1.310577in}{0.746915in}}%
\pgfpathlineto{\pgfqpoint{1.310577in}{0.741854in}}%
\pgfpathlineto{\pgfqpoint{1.262953in}{0.741854in}}%
\pgfpathquadraticcurveto{\pgfqpoint{1.263637in}{0.731154in}}{\pgfqpoint{1.269401in}{0.725553in}}%
\pgfpathquadraticcurveto{\pgfqpoint{1.275165in}{0.719951in}}{\pgfqpoint{1.285468in}{0.719951in}}%
\pgfpathquadraticcurveto{\pgfqpoint{1.291430in}{0.719951in}}{\pgfqpoint{1.297032in}{0.721410in}}%
\pgfpathquadraticcurveto{\pgfqpoint{1.302633in}{0.722869in}}{\pgfqpoint{1.308163in}{0.725805in}}%
\pgfpathlineto{\pgfqpoint{1.308163in}{0.716006in}}%
\pgfpathquadraticcurveto{\pgfqpoint{1.302579in}{0.713647in}}{\pgfqpoint{1.296725in}{0.712404in}}%
\pgfpathquadraticcurveto{\pgfqpoint{1.290872in}{0.711161in}}{\pgfqpoint{1.284855in}{0.711161in}}%
\pgfpathquadraticcurveto{\pgfqpoint{1.269761in}{0.711161in}}{\pgfqpoint{1.260953in}{0.719933in}}%
\pgfpathquadraticcurveto{\pgfqpoint{1.252145in}{0.728723in}}{\pgfqpoint{1.252145in}{0.743709in}}%
\pgfpathquadraticcurveto{\pgfqpoint{1.252145in}{0.759181in}}{\pgfqpoint{1.260503in}{0.768259in}}%
\pgfpathquadraticcurveto{\pgfqpoint{1.268861in}{0.777356in}}{\pgfqpoint{1.283054in}{0.777356in}}%
\pgfpathquadraticcurveto{\pgfqpoint{1.295771in}{0.777356in}}{\pgfqpoint{1.303174in}{0.769160in}}%
\pgfpathquadraticcurveto{\pgfqpoint{1.310577in}{0.760983in}}{\pgfqpoint{1.310577in}{0.746915in}}%
\pgfpathlineto{\pgfqpoint{1.310577in}{0.746915in}}%
\pgfpathclose%
\pgfpathmoveto{\pgfqpoint{1.300220in}{0.749959in}}%
\pgfpathquadraticcurveto{\pgfqpoint{1.300112in}{0.758443in}}{\pgfqpoint{1.295465in}{0.763504in}}%
\pgfpathquadraticcurveto{\pgfqpoint{1.290817in}{0.768584in}}{\pgfqpoint{1.283162in}{0.768584in}}%
\pgfpathquadraticcurveto{\pgfqpoint{1.274498in}{0.768584in}}{\pgfqpoint{1.269293in}{0.763684in}}%
\pgfpathquadraticcurveto{\pgfqpoint{1.264087in}{0.758785in}}{\pgfqpoint{1.263295in}{0.749887in}}%
\pgfpathlineto{\pgfqpoint{1.300220in}{0.749959in}}%
\pgfpathlineto{\pgfqpoint{1.300220in}{0.749959in}}%
\pgfpathclose%
\pgfpathmoveto{\pgfqpoint{1.369062in}{0.766278in}}%
\pgfpathlineto{\pgfqpoint{1.369062in}{0.800393in}}%
\pgfpathlineto{\pgfqpoint{1.379419in}{0.800393in}}%
\pgfpathlineto{\pgfqpoint{1.379419in}{0.712800in}}%
\pgfpathlineto{\pgfqpoint{1.369062in}{0.712800in}}%
\pgfpathlineto{\pgfqpoint{1.369062in}{0.722256in}}%
\pgfpathquadraticcurveto{\pgfqpoint{1.365802in}{0.716637in}}{\pgfqpoint{1.360813in}{0.713899in}}%
\pgfpathquadraticcurveto{\pgfqpoint{1.355841in}{0.711161in}}{\pgfqpoint{1.348853in}{0.711161in}}%
\pgfpathquadraticcurveto{\pgfqpoint{1.337433in}{0.711161in}}{\pgfqpoint{1.330246in}{0.720275in}}%
\pgfpathquadraticcurveto{\pgfqpoint{1.323077in}{0.729407in}}{\pgfqpoint{1.323077in}{0.744267in}}%
\pgfpathquadraticcurveto{\pgfqpoint{1.323077in}{0.759127in}}{\pgfqpoint{1.330246in}{0.768241in}}%
\pgfpathquadraticcurveto{\pgfqpoint{1.337433in}{0.777356in}}{\pgfqpoint{1.348853in}{0.777356in}}%
\pgfpathquadraticcurveto{\pgfqpoint{1.355841in}{0.777356in}}{\pgfqpoint{1.360813in}{0.774618in}}%
\pgfpathquadraticcurveto{\pgfqpoint{1.365802in}{0.771898in}}{\pgfqpoint{1.369062in}{0.766278in}}%
\pgfpathlineto{\pgfqpoint{1.369062in}{0.766278in}}%
\pgfpathclose%
\pgfpathmoveto{\pgfqpoint{1.333776in}{0.744267in}}%
\pgfpathquadraticcurveto{\pgfqpoint{1.333776in}{0.732848in}}{\pgfqpoint{1.338478in}{0.726345in}}%
\pgfpathquadraticcurveto{\pgfqpoint{1.343179in}{0.719843in}}{\pgfqpoint{1.351392in}{0.719843in}}%
\pgfpathquadraticcurveto{\pgfqpoint{1.359606in}{0.719843in}}{\pgfqpoint{1.364325in}{0.726345in}}%
\pgfpathquadraticcurveto{\pgfqpoint{1.369062in}{0.732848in}}{\pgfqpoint{1.369062in}{0.744267in}}%
\pgfpathquadraticcurveto{\pgfqpoint{1.369062in}{0.755687in}}{\pgfqpoint{1.364325in}{0.762189in}}%
\pgfpathquadraticcurveto{\pgfqpoint{1.359606in}{0.768692in}}{\pgfqpoint{1.351392in}{0.768692in}}%
\pgfpathquadraticcurveto{\pgfqpoint{1.343179in}{0.768692in}}{\pgfqpoint{1.338478in}{0.762189in}}%
\pgfpathquadraticcurveto{\pgfqpoint{1.333776in}{0.755687in}}{\pgfqpoint{1.333776in}{0.744267in}}%
\pgfpathlineto{\pgfqpoint{1.333776in}{0.744267in}}%
\pgfpathclose%
\pgfusepath{fill}%
\end{pgfscope}%
\begin{pgfscope}%
\pgfsetbuttcap%
\pgfsetroundjoin%
\definecolor{currentfill}{rgb}{0.309804,0.372549,0.419608}%
\pgfsetfillcolor{currentfill}%
\pgfsetlinewidth{0.000000pt}%
\definecolor{currentstroke}{rgb}{0.309804,0.372549,0.419608}%
\pgfsetstrokecolor{currentstroke}%
\pgfsetdash{}{0pt}%
\pgfpathmoveto{\pgfqpoint{1.488479in}{0.768584in}}%
\pgfpathquadraticcurveto{\pgfqpoint{1.480157in}{0.768584in}}{\pgfqpoint{1.475312in}{0.762081in}}%
\pgfpathquadraticcurveto{\pgfqpoint{1.470467in}{0.755579in}}{\pgfqpoint{1.470467in}{0.744267in}}%
\pgfpathquadraticcurveto{\pgfqpoint{1.470467in}{0.732956in}}{\pgfqpoint{1.475276in}{0.726453in}}%
\pgfpathquadraticcurveto{\pgfqpoint{1.480103in}{0.719951in}}{\pgfqpoint{1.488479in}{0.719951in}}%
\pgfpathquadraticcurveto{\pgfqpoint{1.496764in}{0.719951in}}{\pgfqpoint{1.501592in}{0.726471in}}%
\pgfpathquadraticcurveto{\pgfqpoint{1.506437in}{0.733010in}}{\pgfqpoint{1.506437in}{0.744267in}}%
\pgfpathquadraticcurveto{\pgfqpoint{1.506437in}{0.755471in}}{\pgfqpoint{1.501592in}{0.762027in}}%
\pgfpathquadraticcurveto{\pgfqpoint{1.496764in}{0.768584in}}{\pgfqpoint{1.488479in}{0.768584in}}%
\pgfpathlineto{\pgfqpoint{1.488479in}{0.768584in}}%
\pgfpathclose%
\pgfpathmoveto{\pgfqpoint{1.488479in}{0.777356in}}%
\pgfpathquadraticcurveto{\pgfqpoint{1.501988in}{0.777356in}}{\pgfqpoint{1.509697in}{0.768566in}}%
\pgfpathquadraticcurveto{\pgfqpoint{1.517424in}{0.759794in}}{\pgfqpoint{1.517424in}{0.744267in}}%
\pgfpathquadraticcurveto{\pgfqpoint{1.517424in}{0.728795in}}{\pgfqpoint{1.509697in}{0.719969in}}%
\pgfpathquadraticcurveto{\pgfqpoint{1.501988in}{0.711161in}}{\pgfqpoint{1.488479in}{0.711161in}}%
\pgfpathquadraticcurveto{\pgfqpoint{1.474916in}{0.711161in}}{\pgfqpoint{1.467224in}{0.719969in}}%
\pgfpathquadraticcurveto{\pgfqpoint{1.459551in}{0.728795in}}{\pgfqpoint{1.459551in}{0.744267in}}%
\pgfpathquadraticcurveto{\pgfqpoint{1.459551in}{0.759794in}}{\pgfqpoint{1.467224in}{0.768566in}}%
\pgfpathquadraticcurveto{\pgfqpoint{1.474916in}{0.777356in}}{\pgfqpoint{1.488479in}{0.777356in}}%
\pgfpathlineto{\pgfqpoint{1.488479in}{0.777356in}}%
\pgfpathclose%
\pgfpathmoveto{\pgfqpoint{1.533527in}{0.737675in}}%
\pgfpathlineto{\pgfqpoint{1.533527in}{0.775843in}}%
\pgfpathlineto{\pgfqpoint{1.543884in}{0.775843in}}%
\pgfpathlineto{\pgfqpoint{1.543884in}{0.738071in}}%
\pgfpathquadraticcurveto{\pgfqpoint{1.543884in}{0.729119in}}{\pgfqpoint{1.547360in}{0.724634in}}%
\pgfpathquadraticcurveto{\pgfqpoint{1.550855in}{0.720167in}}{\pgfqpoint{1.557843in}{0.720167in}}%
\pgfpathquadraticcurveto{\pgfqpoint{1.566219in}{0.720167in}}{\pgfqpoint{1.571082in}{0.725517in}}%
\pgfpathquadraticcurveto{\pgfqpoint{1.575964in}{0.730866in}}{\pgfqpoint{1.575964in}{0.740106in}}%
\pgfpathlineto{\pgfqpoint{1.575964in}{0.775843in}}%
\pgfpathlineto{\pgfqpoint{1.586321in}{0.775843in}}%
\pgfpathlineto{\pgfqpoint{1.586321in}{0.712800in}}%
\pgfpathlineto{\pgfqpoint{1.575964in}{0.712800in}}%
\pgfpathlineto{\pgfqpoint{1.575964in}{0.722491in}}%
\pgfpathquadraticcurveto{\pgfqpoint{1.572199in}{0.716745in}}{\pgfqpoint{1.567210in}{0.713953in}}%
\pgfpathquadraticcurveto{\pgfqpoint{1.562238in}{0.711161in}}{\pgfqpoint{1.555646in}{0.711161in}}%
\pgfpathquadraticcurveto{\pgfqpoint{1.544785in}{0.711161in}}{\pgfqpoint{1.539147in}{0.717915in}}%
\pgfpathquadraticcurveto{\pgfqpoint{1.533527in}{0.724670in}}{\pgfqpoint{1.533527in}{0.737675in}}%
\pgfpathlineto{\pgfqpoint{1.533527in}{0.737675in}}%
\pgfpathclose%
\pgfpathmoveto{\pgfqpoint{1.559591in}{0.777356in}}%
\pgfpathlineto{\pgfqpoint{1.559591in}{0.777356in}}%
\pgfpathlineto{\pgfqpoint{1.559591in}{0.777356in}}%
\pgfpathclose%
\pgfpathmoveto{\pgfqpoint{1.617896in}{0.793747in}}%
\pgfpathlineto{\pgfqpoint{1.617896in}{0.775843in}}%
\pgfpathlineto{\pgfqpoint{1.639222in}{0.775843in}}%
\pgfpathlineto{\pgfqpoint{1.639222in}{0.767791in}}%
\pgfpathlineto{\pgfqpoint{1.617896in}{0.767791in}}%
\pgfpathlineto{\pgfqpoint{1.617896in}{0.733568in}}%
\pgfpathquadraticcurveto{\pgfqpoint{1.617896in}{0.725859in}}{\pgfqpoint{1.620003in}{0.723661in}}%
\pgfpathquadraticcurveto{\pgfqpoint{1.622111in}{0.721464in}}{\pgfqpoint{1.628595in}{0.721464in}}%
\pgfpathlineto{\pgfqpoint{1.639222in}{0.721464in}}%
\pgfpathlineto{\pgfqpoint{1.639222in}{0.712800in}}%
\pgfpathlineto{\pgfqpoint{1.628595in}{0.712800in}}%
\pgfpathquadraticcurveto{\pgfqpoint{1.616599in}{0.712800in}}{\pgfqpoint{1.612042in}{0.717267in}}%
\pgfpathquadraticcurveto{\pgfqpoint{1.607485in}{0.721752in}}{\pgfqpoint{1.607485in}{0.733568in}}%
\pgfpathlineto{\pgfqpoint{1.607485in}{0.767791in}}%
\pgfpathlineto{\pgfqpoint{1.599884in}{0.767791in}}%
\pgfpathlineto{\pgfqpoint{1.599884in}{0.775843in}}%
\pgfpathlineto{\pgfqpoint{1.607485in}{0.775843in}}%
\pgfpathlineto{\pgfqpoint{1.607485in}{0.793747in}}%
\pgfpathlineto{\pgfqpoint{1.617896in}{0.793747in}}%
\pgfpathlineto{\pgfqpoint{1.617896in}{0.793747in}}%
\pgfpathclose%
\pgfpathmoveto{\pgfqpoint{1.698212in}{0.773429in}}%
\pgfpathlineto{\pgfqpoint{1.698212in}{0.763738in}}%
\pgfpathquadraticcurveto{\pgfqpoint{1.693817in}{0.766170in}}{\pgfqpoint{1.689404in}{0.767377in}}%
\pgfpathquadraticcurveto{\pgfqpoint{1.684991in}{0.768584in}}{\pgfqpoint{1.680488in}{0.768584in}}%
\pgfpathquadraticcurveto{\pgfqpoint{1.670401in}{0.768584in}}{\pgfqpoint{1.664818in}{0.762189in}}%
\pgfpathquadraticcurveto{\pgfqpoint{1.659252in}{0.755813in}}{\pgfqpoint{1.659252in}{0.744267in}}%
\pgfpathquadraticcurveto{\pgfqpoint{1.659252in}{0.732721in}}{\pgfqpoint{1.664818in}{0.726327in}}%
\pgfpathquadraticcurveto{\pgfqpoint{1.670401in}{0.719951in}}{\pgfqpoint{1.680488in}{0.719951in}}%
\pgfpathquadraticcurveto{\pgfqpoint{1.684991in}{0.719951in}}{\pgfqpoint{1.689404in}{0.721158in}}%
\pgfpathquadraticcurveto{\pgfqpoint{1.693817in}{0.722364in}}{\pgfqpoint{1.698212in}{0.724796in}}%
\pgfpathlineto{\pgfqpoint{1.698212in}{0.715214in}}%
\pgfpathquadraticcurveto{\pgfqpoint{1.693871in}{0.713196in}}{\pgfqpoint{1.689224in}{0.712188in}}%
\pgfpathquadraticcurveto{\pgfqpoint{1.684595in}{0.711161in}}{\pgfqpoint{1.679353in}{0.711161in}}%
\pgfpathquadraticcurveto{\pgfqpoint{1.665106in}{0.711161in}}{\pgfqpoint{1.656712in}{0.720113in}}%
\pgfpathquadraticcurveto{\pgfqpoint{1.648337in}{0.729065in}}{\pgfqpoint{1.648337in}{0.744267in}}%
\pgfpathquadraticcurveto{\pgfqpoint{1.648337in}{0.759686in}}{\pgfqpoint{1.656802in}{0.768512in}}%
\pgfpathquadraticcurveto{\pgfqpoint{1.665286in}{0.777356in}}{\pgfqpoint{1.680038in}{0.777356in}}%
\pgfpathquadraticcurveto{\pgfqpoint{1.684811in}{0.777356in}}{\pgfqpoint{1.689368in}{0.776365in}}%
\pgfpathquadraticcurveto{\pgfqpoint{1.693925in}{0.775392in}}{\pgfqpoint{1.698212in}{0.773429in}}%
\pgfpathlineto{\pgfqpoint{1.698212in}{0.773429in}}%
\pgfpathclose%
\pgfpathmoveto{\pgfqpoint{1.740649in}{0.768584in}}%
\pgfpathquadraticcurveto{\pgfqpoint{1.732327in}{0.768584in}}{\pgfqpoint{1.727482in}{0.762081in}}%
\pgfpathquadraticcurveto{\pgfqpoint{1.722637in}{0.755579in}}{\pgfqpoint{1.722637in}{0.744267in}}%
\pgfpathquadraticcurveto{\pgfqpoint{1.722637in}{0.732956in}}{\pgfqpoint{1.727446in}{0.726453in}}%
\pgfpathquadraticcurveto{\pgfqpoint{1.732273in}{0.719951in}}{\pgfqpoint{1.740649in}{0.719951in}}%
\pgfpathquadraticcurveto{\pgfqpoint{1.748934in}{0.719951in}}{\pgfqpoint{1.753762in}{0.726471in}}%
\pgfpathquadraticcurveto{\pgfqpoint{1.758607in}{0.733010in}}{\pgfqpoint{1.758607in}{0.744267in}}%
\pgfpathquadraticcurveto{\pgfqpoint{1.758607in}{0.755471in}}{\pgfqpoint{1.753762in}{0.762027in}}%
\pgfpathquadraticcurveto{\pgfqpoint{1.748934in}{0.768584in}}{\pgfqpoint{1.740649in}{0.768584in}}%
\pgfpathlineto{\pgfqpoint{1.740649in}{0.768584in}}%
\pgfpathclose%
\pgfpathmoveto{\pgfqpoint{1.740649in}{0.777356in}}%
\pgfpathquadraticcurveto{\pgfqpoint{1.754158in}{0.777356in}}{\pgfqpoint{1.761867in}{0.768566in}}%
\pgfpathquadraticcurveto{\pgfqpoint{1.769594in}{0.759794in}}{\pgfqpoint{1.769594in}{0.744267in}}%
\pgfpathquadraticcurveto{\pgfqpoint{1.769594in}{0.728795in}}{\pgfqpoint{1.761867in}{0.719969in}}%
\pgfpathquadraticcurveto{\pgfqpoint{1.754158in}{0.711161in}}{\pgfqpoint{1.740649in}{0.711161in}}%
\pgfpathquadraticcurveto{\pgfqpoint{1.727086in}{0.711161in}}{\pgfqpoint{1.719394in}{0.719969in}}%
\pgfpathquadraticcurveto{\pgfqpoint{1.711721in}{0.728795in}}{\pgfqpoint{1.711721in}{0.744267in}}%
\pgfpathquadraticcurveto{\pgfqpoint{1.711721in}{0.759794in}}{\pgfqpoint{1.719394in}{0.768566in}}%
\pgfpathquadraticcurveto{\pgfqpoint{1.727086in}{0.777356in}}{\pgfqpoint{1.740649in}{0.777356in}}%
\pgfpathlineto{\pgfqpoint{1.740649in}{0.777356in}}%
\pgfpathclose%
\pgfpathmoveto{\pgfqpoint{1.835843in}{0.763738in}}%
\pgfpathquadraticcurveto{\pgfqpoint{1.839734in}{0.770727in}}{\pgfqpoint{1.845137in}{0.774041in}}%
\pgfpathquadraticcurveto{\pgfqpoint{1.850541in}{0.777356in}}{\pgfqpoint{1.857854in}{0.777356in}}%
\pgfpathquadraticcurveto{\pgfqpoint{1.867707in}{0.777356in}}{\pgfqpoint{1.873056in}{0.770457in}}%
\pgfpathquadraticcurveto{\pgfqpoint{1.878406in}{0.763576in}}{\pgfqpoint{1.878406in}{0.750860in}}%
\pgfpathlineto{\pgfqpoint{1.878406in}{0.712800in}}%
\pgfpathlineto{\pgfqpoint{1.867995in}{0.712800in}}%
\pgfpathlineto{\pgfqpoint{1.867995in}{0.750517in}}%
\pgfpathquadraticcurveto{\pgfqpoint{1.867995in}{0.759578in}}{\pgfqpoint{1.864771in}{0.763955in}}%
\pgfpathquadraticcurveto{\pgfqpoint{1.861564in}{0.768349in}}{\pgfqpoint{1.854990in}{0.768349in}}%
\pgfpathquadraticcurveto{\pgfqpoint{1.846939in}{0.768349in}}{\pgfqpoint{1.842255in}{0.763000in}}%
\pgfpathquadraticcurveto{\pgfqpoint{1.837590in}{0.757668in}}{\pgfqpoint{1.837590in}{0.748428in}}%
\pgfpathlineto{\pgfqpoint{1.837590in}{0.712800in}}%
\pgfpathlineto{\pgfqpoint{1.827179in}{0.712800in}}%
\pgfpathlineto{\pgfqpoint{1.827179in}{0.750517in}}%
\pgfpathquadraticcurveto{\pgfqpoint{1.827179in}{0.759632in}}{\pgfqpoint{1.823973in}{0.763991in}}%
\pgfpathquadraticcurveto{\pgfqpoint{1.820767in}{0.768349in}}{\pgfqpoint{1.814066in}{0.768349in}}%
\pgfpathquadraticcurveto{\pgfqpoint{1.806123in}{0.768349in}}{\pgfqpoint{1.801440in}{0.762982in}}%
\pgfpathquadraticcurveto{\pgfqpoint{1.796775in}{0.757614in}}{\pgfqpoint{1.796775in}{0.748428in}}%
\pgfpathlineto{\pgfqpoint{1.796775in}{0.712800in}}%
\pgfpathlineto{\pgfqpoint{1.786364in}{0.712800in}}%
\pgfpathlineto{\pgfqpoint{1.786364in}{0.775843in}}%
\pgfpathlineto{\pgfqpoint{1.796775in}{0.775843in}}%
\pgfpathlineto{\pgfqpoint{1.796775in}{0.766044in}}%
\pgfpathquadraticcurveto{\pgfqpoint{1.800323in}{0.771844in}}{\pgfqpoint{1.805276in}{0.774600in}}%
\pgfpathquadraticcurveto{\pgfqpoint{1.810230in}{0.777356in}}{\pgfqpoint{1.817038in}{0.777356in}}%
\pgfpathquadraticcurveto{\pgfqpoint{1.823919in}{0.777356in}}{\pgfqpoint{1.828728in}{0.773861in}}%
\pgfpathquadraticcurveto{\pgfqpoint{1.833538in}{0.770385in}}{\pgfqpoint{1.835843in}{0.763738in}}%
\pgfpathlineto{\pgfqpoint{1.835843in}{0.763738in}}%
\pgfpathclose%
\pgfpathmoveto{\pgfqpoint{1.952976in}{0.746915in}}%
\pgfpathlineto{\pgfqpoint{1.952976in}{0.741854in}}%
\pgfpathlineto{\pgfqpoint{1.905352in}{0.741854in}}%
\pgfpathquadraticcurveto{\pgfqpoint{1.906036in}{0.731154in}}{\pgfqpoint{1.911800in}{0.725553in}}%
\pgfpathquadraticcurveto{\pgfqpoint{1.917564in}{0.719951in}}{\pgfqpoint{1.927867in}{0.719951in}}%
\pgfpathquadraticcurveto{\pgfqpoint{1.933829in}{0.719951in}}{\pgfqpoint{1.939431in}{0.721410in}}%
\pgfpathquadraticcurveto{\pgfqpoint{1.945033in}{0.722869in}}{\pgfqpoint{1.950562in}{0.725805in}}%
\pgfpathlineto{\pgfqpoint{1.950562in}{0.716006in}}%
\pgfpathquadraticcurveto{\pgfqpoint{1.944979in}{0.713647in}}{\pgfqpoint{1.939125in}{0.712404in}}%
\pgfpathquadraticcurveto{\pgfqpoint{1.933271in}{0.711161in}}{\pgfqpoint{1.927255in}{0.711161in}}%
\pgfpathquadraticcurveto{\pgfqpoint{1.912161in}{0.711161in}}{\pgfqpoint{1.903353in}{0.719933in}}%
\pgfpathquadraticcurveto{\pgfqpoint{1.894545in}{0.728723in}}{\pgfqpoint{1.894545in}{0.743709in}}%
\pgfpathquadraticcurveto{\pgfqpoint{1.894545in}{0.759181in}}{\pgfqpoint{1.902902in}{0.768259in}}%
\pgfpathquadraticcurveto{\pgfqpoint{1.911260in}{0.777356in}}{\pgfqpoint{1.925454in}{0.777356in}}%
\pgfpathquadraticcurveto{\pgfqpoint{1.938170in}{0.777356in}}{\pgfqpoint{1.945573in}{0.769160in}}%
\pgfpathquadraticcurveto{\pgfqpoint{1.952976in}{0.760983in}}{\pgfqpoint{1.952976in}{0.746915in}}%
\pgfpathlineto{\pgfqpoint{1.952976in}{0.746915in}}%
\pgfpathclose%
\pgfpathmoveto{\pgfqpoint{1.942619in}{0.749959in}}%
\pgfpathquadraticcurveto{\pgfqpoint{1.942511in}{0.758443in}}{\pgfqpoint{1.937864in}{0.763504in}}%
\pgfpathquadraticcurveto{\pgfqpoint{1.933217in}{0.768584in}}{\pgfqpoint{1.925562in}{0.768584in}}%
\pgfpathquadraticcurveto{\pgfqpoint{1.916898in}{0.768584in}}{\pgfqpoint{1.911692in}{0.763684in}}%
\pgfpathquadraticcurveto{\pgfqpoint{1.906487in}{0.758785in}}{\pgfqpoint{1.905694in}{0.749887in}}%
\pgfpathlineto{\pgfqpoint{1.942619in}{0.749959in}}%
\pgfpathlineto{\pgfqpoint{1.942619in}{0.749959in}}%
\pgfpathclose%
\pgfpathmoveto{\pgfqpoint{1.964756in}{0.748986in}}%
\pgfpathlineto{\pgfqpoint{1.995089in}{0.748986in}}%
\pgfpathlineto{\pgfqpoint{1.995089in}{0.739764in}}%
\pgfpathlineto{\pgfqpoint{1.964756in}{0.739764in}}%
\pgfpathlineto{\pgfqpoint{1.964756in}{0.748986in}}%
\pgfpathlineto{\pgfqpoint{1.964756in}{0.748986in}}%
\pgfpathclose%
\pgfpathmoveto{\pgfqpoint{2.011570in}{0.800393in}}%
\pgfpathlineto{\pgfqpoint{2.021927in}{0.800393in}}%
\pgfpathlineto{\pgfqpoint{2.021927in}{0.712800in}}%
\pgfpathlineto{\pgfqpoint{2.011570in}{0.712800in}}%
\pgfpathlineto{\pgfqpoint{2.011570in}{0.800393in}}%
\pgfpathlineto{\pgfqpoint{2.011570in}{0.800393in}}%
\pgfpathclose%
\pgfpathmoveto{\pgfqpoint{2.043595in}{0.775843in}}%
\pgfpathlineto{\pgfqpoint{2.053952in}{0.775843in}}%
\pgfpathlineto{\pgfqpoint{2.053952in}{0.712800in}}%
\pgfpathlineto{\pgfqpoint{2.043595in}{0.712800in}}%
\pgfpathlineto{\pgfqpoint{2.043595in}{0.775843in}}%
\pgfpathlineto{\pgfqpoint{2.043595in}{0.775843in}}%
\pgfpathclose%
\pgfpathmoveto{\pgfqpoint{2.043595in}{0.800393in}}%
\pgfpathlineto{\pgfqpoint{2.053952in}{0.800393in}}%
\pgfpathlineto{\pgfqpoint{2.053952in}{0.787262in}}%
\pgfpathlineto{\pgfqpoint{2.043595in}{0.787262in}}%
\pgfpathlineto{\pgfqpoint{2.043595in}{0.800393in}}%
\pgfpathlineto{\pgfqpoint{2.043595in}{0.800393in}}%
\pgfpathclose%
\pgfpathmoveto{\pgfqpoint{2.128036in}{0.750860in}}%
\pgfpathlineto{\pgfqpoint{2.128036in}{0.712800in}}%
\pgfpathlineto{\pgfqpoint{2.117679in}{0.712800in}}%
\pgfpathlineto{\pgfqpoint{2.117679in}{0.750517in}}%
\pgfpathquadraticcurveto{\pgfqpoint{2.117679in}{0.759469in}}{\pgfqpoint{2.114185in}{0.763900in}}%
\pgfpathquadraticcurveto{\pgfqpoint{2.110691in}{0.768349in}}{\pgfqpoint{2.103720in}{0.768349in}}%
\pgfpathquadraticcurveto{\pgfqpoint{2.095326in}{0.768349in}}{\pgfqpoint{2.090481in}{0.763000in}}%
\pgfpathquadraticcurveto{\pgfqpoint{2.085636in}{0.757668in}}{\pgfqpoint{2.085636in}{0.748428in}}%
\pgfpathlineto{\pgfqpoint{2.085636in}{0.712800in}}%
\pgfpathlineto{\pgfqpoint{2.075225in}{0.712800in}}%
\pgfpathlineto{\pgfqpoint{2.075225in}{0.775843in}}%
\pgfpathlineto{\pgfqpoint{2.085636in}{0.775843in}}%
\pgfpathlineto{\pgfqpoint{2.085636in}{0.766044in}}%
\pgfpathquadraticcurveto{\pgfqpoint{2.089364in}{0.771736in}}{\pgfqpoint{2.094390in}{0.774546in}}%
\pgfpathquadraticcurveto{\pgfqpoint{2.099433in}{0.777356in}}{\pgfqpoint{2.106025in}{0.777356in}}%
\pgfpathquadraticcurveto{\pgfqpoint{2.116887in}{0.777356in}}{\pgfqpoint{2.122452in}{0.770637in}}%
\pgfpathquadraticcurveto{\pgfqpoint{2.128036in}{0.763918in}}{\pgfqpoint{2.128036in}{0.750860in}}%
\pgfpathlineto{\pgfqpoint{2.128036in}{0.750860in}}%
\pgfpathclose%
\pgfpathmoveto{\pgfqpoint{2.202607in}{0.746915in}}%
\pgfpathlineto{\pgfqpoint{2.202607in}{0.741854in}}%
\pgfpathlineto{\pgfqpoint{2.154982in}{0.741854in}}%
\pgfpathquadraticcurveto{\pgfqpoint{2.155667in}{0.731154in}}{\pgfqpoint{2.161431in}{0.725553in}}%
\pgfpathquadraticcurveto{\pgfqpoint{2.167195in}{0.719951in}}{\pgfqpoint{2.177498in}{0.719951in}}%
\pgfpathquadraticcurveto{\pgfqpoint{2.183460in}{0.719951in}}{\pgfqpoint{2.189061in}{0.721410in}}%
\pgfpathquadraticcurveto{\pgfqpoint{2.194663in}{0.722869in}}{\pgfqpoint{2.200193in}{0.725805in}}%
\pgfpathlineto{\pgfqpoint{2.200193in}{0.716006in}}%
\pgfpathquadraticcurveto{\pgfqpoint{2.194609in}{0.713647in}}{\pgfqpoint{2.188755in}{0.712404in}}%
\pgfpathquadraticcurveto{\pgfqpoint{2.182901in}{0.711161in}}{\pgfqpoint{2.176885in}{0.711161in}}%
\pgfpathquadraticcurveto{\pgfqpoint{2.161791in}{0.711161in}}{\pgfqpoint{2.152983in}{0.719933in}}%
\pgfpathquadraticcurveto{\pgfqpoint{2.144175in}{0.728723in}}{\pgfqpoint{2.144175in}{0.743709in}}%
\pgfpathquadraticcurveto{\pgfqpoint{2.144175in}{0.759181in}}{\pgfqpoint{2.152533in}{0.768259in}}%
\pgfpathquadraticcurveto{\pgfqpoint{2.160890in}{0.777356in}}{\pgfqpoint{2.175084in}{0.777356in}}%
\pgfpathquadraticcurveto{\pgfqpoint{2.187801in}{0.777356in}}{\pgfqpoint{2.195204in}{0.769160in}}%
\pgfpathquadraticcurveto{\pgfqpoint{2.202607in}{0.760983in}}{\pgfqpoint{2.202607in}{0.746915in}}%
\pgfpathlineto{\pgfqpoint{2.202607in}{0.746915in}}%
\pgfpathclose%
\pgfpathmoveto{\pgfqpoint{2.192250in}{0.749959in}}%
\pgfpathquadraticcurveto{\pgfqpoint{2.192141in}{0.758443in}}{\pgfqpoint{2.187494in}{0.763504in}}%
\pgfpathquadraticcurveto{\pgfqpoint{2.182847in}{0.768584in}}{\pgfqpoint{2.175192in}{0.768584in}}%
\pgfpathquadraticcurveto{\pgfqpoint{2.166528in}{0.768584in}}{\pgfqpoint{2.161323in}{0.763684in}}%
\pgfpathquadraticcurveto{\pgfqpoint{2.156117in}{0.758785in}}{\pgfqpoint{2.155325in}{0.749887in}}%
\pgfpathlineto{\pgfqpoint{2.192250in}{0.749959in}}%
\pgfpathlineto{\pgfqpoint{2.192250in}{0.749959in}}%
\pgfpathclose%
\pgfpathmoveto{\pgfqpoint{2.248267in}{0.744483in}}%
\pgfpathquadraticcurveto{\pgfqpoint{2.235713in}{0.744483in}}{\pgfqpoint{2.230868in}{0.741619in}}%
\pgfpathquadraticcurveto{\pgfqpoint{2.226022in}{0.738756in}}{\pgfqpoint{2.226022in}{0.731821in}}%
\pgfpathquadraticcurveto{\pgfqpoint{2.226022in}{0.726309in}}{\pgfqpoint{2.229661in}{0.723067in}}%
\pgfpathquadraticcurveto{\pgfqpoint{2.233299in}{0.719843in}}{\pgfqpoint{2.239531in}{0.719843in}}%
\pgfpathquadraticcurveto{\pgfqpoint{2.248159in}{0.719843in}}{\pgfqpoint{2.253365in}{0.725949in}}%
\pgfpathquadraticcurveto{\pgfqpoint{2.258570in}{0.732055in}}{\pgfqpoint{2.258570in}{0.742178in}}%
\pgfpathlineto{\pgfqpoint{2.258570in}{0.744483in}}%
\pgfpathlineto{\pgfqpoint{2.248267in}{0.744483in}}%
\pgfpathlineto{\pgfqpoint{2.248267in}{0.744483in}}%
\pgfpathclose%
\pgfpathmoveto{\pgfqpoint{2.268927in}{0.748770in}}%
\pgfpathlineto{\pgfqpoint{2.268927in}{0.712800in}}%
\pgfpathlineto{\pgfqpoint{2.258570in}{0.712800in}}%
\pgfpathlineto{\pgfqpoint{2.258570in}{0.722364in}}%
\pgfpathquadraticcurveto{\pgfqpoint{2.255022in}{0.716637in}}{\pgfqpoint{2.249726in}{0.713899in}}%
\pgfpathquadraticcurveto{\pgfqpoint{2.244431in}{0.711161in}}{\pgfqpoint{2.236776in}{0.711161in}}%
\pgfpathquadraticcurveto{\pgfqpoint{2.227103in}{0.711161in}}{\pgfqpoint{2.221375in}{0.716601in}}%
\pgfpathquadraticcurveto{\pgfqpoint{2.215665in}{0.722040in}}{\pgfqpoint{2.215665in}{0.731154in}}%
\pgfpathquadraticcurveto{\pgfqpoint{2.215665in}{0.741782in}}{\pgfqpoint{2.222780in}{0.747185in}}%
\pgfpathquadraticcurveto{\pgfqpoint{2.229913in}{0.752589in}}{\pgfqpoint{2.244034in}{0.752589in}}%
\pgfpathlineto{\pgfqpoint{2.258570in}{0.752589in}}%
\pgfpathlineto{\pgfqpoint{2.258570in}{0.753616in}}%
\pgfpathquadraticcurveto{\pgfqpoint{2.258570in}{0.760766in}}{\pgfqpoint{2.253869in}{0.764675in}}%
\pgfpathquadraticcurveto{\pgfqpoint{2.249168in}{0.768584in}}{\pgfqpoint{2.240666in}{0.768584in}}%
\pgfpathquadraticcurveto{\pgfqpoint{2.235263in}{0.768584in}}{\pgfqpoint{2.230129in}{0.767287in}}%
\pgfpathquadraticcurveto{\pgfqpoint{2.225014in}{0.765990in}}{\pgfqpoint{2.220294in}{0.763396in}}%
\pgfpathlineto{\pgfqpoint{2.220294in}{0.772979in}}%
\pgfpathquadraticcurveto{\pgfqpoint{2.225968in}{0.775176in}}{\pgfqpoint{2.231318in}{0.776257in}}%
\pgfpathquadraticcurveto{\pgfqpoint{2.236667in}{0.777356in}}{\pgfqpoint{2.241729in}{0.777356in}}%
\pgfpathquadraticcurveto{\pgfqpoint{2.255418in}{0.777356in}}{\pgfqpoint{2.262173in}{0.770259in}}%
\pgfpathquadraticcurveto{\pgfqpoint{2.268927in}{0.763180in}}{\pgfqpoint{2.268927in}{0.748770in}}%
\pgfpathlineto{\pgfqpoint{2.268927in}{0.748770in}}%
\pgfpathclose%
\pgfpathmoveto{\pgfqpoint{2.326782in}{0.766170in}}%
\pgfpathquadraticcurveto{\pgfqpoint{2.325035in}{0.767179in}}{\pgfqpoint{2.322982in}{0.767647in}}%
\pgfpathquadraticcurveto{\pgfqpoint{2.320928in}{0.768133in}}{\pgfqpoint{2.318461in}{0.768133in}}%
\pgfpathquadraticcurveto{\pgfqpoint{2.309671in}{0.768133in}}{\pgfqpoint{2.304970in}{0.762423in}}%
\pgfpathquadraticcurveto{\pgfqpoint{2.300268in}{0.756714in}}{\pgfqpoint{2.300268in}{0.746014in}}%
\pgfpathlineto{\pgfqpoint{2.300268in}{0.712800in}}%
\pgfpathlineto{\pgfqpoint{2.289857in}{0.712800in}}%
\pgfpathlineto{\pgfqpoint{2.289857in}{0.775843in}}%
\pgfpathlineto{\pgfqpoint{2.300268in}{0.775843in}}%
\pgfpathlineto{\pgfqpoint{2.300268in}{0.766044in}}%
\pgfpathquadraticcurveto{\pgfqpoint{2.303547in}{0.771790in}}{\pgfqpoint{2.308770in}{0.774564in}}%
\pgfpathquadraticcurveto{\pgfqpoint{2.314012in}{0.777356in}}{\pgfqpoint{2.321505in}{0.777356in}}%
\pgfpathquadraticcurveto{\pgfqpoint{2.322567in}{0.777356in}}{\pgfqpoint{2.323864in}{0.777211in}}%
\pgfpathquadraticcurveto{\pgfqpoint{2.325161in}{0.777085in}}{\pgfqpoint{2.326728in}{0.776797in}}%
\pgfpathlineto{\pgfqpoint{2.326782in}{0.766170in}}%
\pgfpathlineto{\pgfqpoint{2.326782in}{0.766170in}}%
\pgfpathclose%
\pgfusepath{fill}%
\end{pgfscope}%
\begin{pgfscope}%
\pgfsetbuttcap%
\pgfsetroundjoin%
\definecolor{currentfill}{rgb}{0.309804,0.372549,0.419608}%
\pgfsetfillcolor{currentfill}%
\pgfsetlinewidth{0.000000pt}%
\definecolor{currentstroke}{rgb}{0.309804,0.372549,0.419608}%
\pgfsetstrokecolor{currentstroke}%
\pgfsetdash{}{0pt}%
\pgfpathmoveto{\pgfqpoint{2.399731in}{0.775843in}}%
\pgfpathlineto{\pgfqpoint{2.410088in}{0.775843in}}%
\pgfpathlineto{\pgfqpoint{2.423039in}{0.726651in}}%
\pgfpathlineto{\pgfqpoint{2.435917in}{0.775843in}}%
\pgfpathlineto{\pgfqpoint{2.448130in}{0.775843in}}%
\pgfpathlineto{\pgfqpoint{2.461080in}{0.726651in}}%
\pgfpathlineto{\pgfqpoint{2.473977in}{0.775843in}}%
\pgfpathlineto{\pgfqpoint{2.484334in}{0.775843in}}%
\pgfpathlineto{\pgfqpoint{2.467835in}{0.712800in}}%
\pgfpathlineto{\pgfqpoint{2.455623in}{0.712800in}}%
\pgfpathlineto{\pgfqpoint{2.442060in}{0.764477in}}%
\pgfpathlineto{\pgfqpoint{2.428442in}{0.712800in}}%
\pgfpathlineto{\pgfqpoint{2.416212in}{0.712800in}}%
\pgfpathlineto{\pgfqpoint{2.399731in}{0.775843in}}%
\pgfpathlineto{\pgfqpoint{2.399731in}{0.775843in}}%
\pgfpathclose%
\pgfpathmoveto{\pgfqpoint{2.553951in}{0.746915in}}%
\pgfpathlineto{\pgfqpoint{2.553951in}{0.741854in}}%
\pgfpathlineto{\pgfqpoint{2.506327in}{0.741854in}}%
\pgfpathquadraticcurveto{\pgfqpoint{2.507011in}{0.731154in}}{\pgfqpoint{2.512775in}{0.725553in}}%
\pgfpathquadraticcurveto{\pgfqpoint{2.518539in}{0.719951in}}{\pgfqpoint{2.528842in}{0.719951in}}%
\pgfpathquadraticcurveto{\pgfqpoint{2.534804in}{0.719951in}}{\pgfqpoint{2.540406in}{0.721410in}}%
\pgfpathquadraticcurveto{\pgfqpoint{2.546008in}{0.722869in}}{\pgfqpoint{2.551537in}{0.725805in}}%
\pgfpathlineto{\pgfqpoint{2.551537in}{0.716006in}}%
\pgfpathquadraticcurveto{\pgfqpoint{2.545954in}{0.713647in}}{\pgfqpoint{2.540100in}{0.712404in}}%
\pgfpathquadraticcurveto{\pgfqpoint{2.534246in}{0.711161in}}{\pgfqpoint{2.528230in}{0.711161in}}%
\pgfpathquadraticcurveto{\pgfqpoint{2.513136in}{0.711161in}}{\pgfqpoint{2.504328in}{0.719933in}}%
\pgfpathquadraticcurveto{\pgfqpoint{2.495520in}{0.728723in}}{\pgfqpoint{2.495520in}{0.743709in}}%
\pgfpathquadraticcurveto{\pgfqpoint{2.495520in}{0.759181in}}{\pgfqpoint{2.503877in}{0.768259in}}%
\pgfpathquadraticcurveto{\pgfqpoint{2.512235in}{0.777356in}}{\pgfqpoint{2.526429in}{0.777356in}}%
\pgfpathquadraticcurveto{\pgfqpoint{2.539145in}{0.777356in}}{\pgfqpoint{2.546548in}{0.769160in}}%
\pgfpathquadraticcurveto{\pgfqpoint{2.553951in}{0.760983in}}{\pgfqpoint{2.553951in}{0.746915in}}%
\pgfpathlineto{\pgfqpoint{2.553951in}{0.746915in}}%
\pgfpathclose%
\pgfpathmoveto{\pgfqpoint{2.543594in}{0.749959in}}%
\pgfpathquadraticcurveto{\pgfqpoint{2.543486in}{0.758443in}}{\pgfqpoint{2.538839in}{0.763504in}}%
\pgfpathquadraticcurveto{\pgfqpoint{2.534192in}{0.768584in}}{\pgfqpoint{2.526537in}{0.768584in}}%
\pgfpathquadraticcurveto{\pgfqpoint{2.517873in}{0.768584in}}{\pgfqpoint{2.512667in}{0.763684in}}%
\pgfpathquadraticcurveto{\pgfqpoint{2.507462in}{0.758785in}}{\pgfqpoint{2.506669in}{0.749887in}}%
\pgfpathlineto{\pgfqpoint{2.543594in}{0.749959in}}%
\pgfpathlineto{\pgfqpoint{2.543594in}{0.749959in}}%
\pgfpathclose%
\pgfpathmoveto{\pgfqpoint{2.570955in}{0.775843in}}%
\pgfpathlineto{\pgfqpoint{2.581312in}{0.775843in}}%
\pgfpathlineto{\pgfqpoint{2.581312in}{0.712800in}}%
\pgfpathlineto{\pgfqpoint{2.570955in}{0.712800in}}%
\pgfpathlineto{\pgfqpoint{2.570955in}{0.775843in}}%
\pgfpathlineto{\pgfqpoint{2.570955in}{0.775843in}}%
\pgfpathclose%
\pgfpathmoveto{\pgfqpoint{2.570955in}{0.800393in}}%
\pgfpathlineto{\pgfqpoint{2.581312in}{0.800393in}}%
\pgfpathlineto{\pgfqpoint{2.581312in}{0.787262in}}%
\pgfpathlineto{\pgfqpoint{2.570955in}{0.787262in}}%
\pgfpathlineto{\pgfqpoint{2.570955in}{0.800393in}}%
\pgfpathlineto{\pgfqpoint{2.570955in}{0.800393in}}%
\pgfpathclose%
\pgfpathmoveto{\pgfqpoint{2.644462in}{0.745060in}}%
\pgfpathquadraticcurveto{\pgfqpoint{2.644462in}{0.756317in}}{\pgfqpoint{2.639815in}{0.762496in}}%
\pgfpathquadraticcurveto{\pgfqpoint{2.635186in}{0.768692in}}{\pgfqpoint{2.626792in}{0.768692in}}%
\pgfpathquadraticcurveto{\pgfqpoint{2.618471in}{0.768692in}}{\pgfqpoint{2.613823in}{0.762496in}}%
\pgfpathquadraticcurveto{\pgfqpoint{2.609176in}{0.756317in}}{\pgfqpoint{2.609176in}{0.745060in}}%
\pgfpathquadraticcurveto{\pgfqpoint{2.609176in}{0.733856in}}{\pgfqpoint{2.613823in}{0.727660in}}%
\pgfpathquadraticcurveto{\pgfqpoint{2.618471in}{0.721464in}}{\pgfqpoint{2.626792in}{0.721464in}}%
\pgfpathquadraticcurveto{\pgfqpoint{2.635186in}{0.721464in}}{\pgfqpoint{2.639815in}{0.727660in}}%
\pgfpathquadraticcurveto{\pgfqpoint{2.644462in}{0.733856in}}{\pgfqpoint{2.644462in}{0.745060in}}%
\pgfpathlineto{\pgfqpoint{2.644462in}{0.745060in}}%
\pgfpathclose%
\pgfpathmoveto{\pgfqpoint{2.654819in}{0.720617in}}%
\pgfpathquadraticcurveto{\pgfqpoint{2.654819in}{0.704532in}}{\pgfqpoint{2.647668in}{0.696679in}}%
\pgfpathquadraticcurveto{\pgfqpoint{2.640536in}{0.688826in}}{\pgfqpoint{2.625784in}{0.688826in}}%
\pgfpathquadraticcurveto{\pgfqpoint{2.620326in}{0.688826in}}{\pgfqpoint{2.615481in}{0.689636in}}%
\pgfpathquadraticcurveto{\pgfqpoint{2.610635in}{0.690447in}}{\pgfqpoint{2.606078in}{0.692140in}}%
\pgfpathlineto{\pgfqpoint{2.606078in}{0.702209in}}%
\pgfpathquadraticcurveto{\pgfqpoint{2.610635in}{0.699741in}}{\pgfqpoint{2.615084in}{0.698570in}}%
\pgfpathquadraticcurveto{\pgfqpoint{2.619533in}{0.697382in}}{\pgfqpoint{2.624144in}{0.697382in}}%
\pgfpathquadraticcurveto{\pgfqpoint{2.634339in}{0.697382in}}{\pgfqpoint{2.639401in}{0.702695in}}%
\pgfpathquadraticcurveto{\pgfqpoint{2.644462in}{0.708009in}}{\pgfqpoint{2.644462in}{0.718762in}}%
\pgfpathlineto{\pgfqpoint{2.644462in}{0.723895in}}%
\pgfpathquadraticcurveto{\pgfqpoint{2.641256in}{0.718312in}}{\pgfqpoint{2.636249in}{0.715556in}}%
\pgfpathquadraticcurveto{\pgfqpoint{2.631241in}{0.712800in}}{\pgfqpoint{2.624253in}{0.712800in}}%
\pgfpathquadraticcurveto{\pgfqpoint{2.612671in}{0.712800in}}{\pgfqpoint{2.605574in}{0.721626in}}%
\pgfpathquadraticcurveto{\pgfqpoint{2.598477in}{0.730470in}}{\pgfqpoint{2.598477in}{0.745060in}}%
\pgfpathquadraticcurveto{\pgfqpoint{2.598477in}{0.759686in}}{\pgfqpoint{2.605574in}{0.768512in}}%
\pgfpathquadraticcurveto{\pgfqpoint{2.612671in}{0.777356in}}{\pgfqpoint{2.624253in}{0.777356in}}%
\pgfpathquadraticcurveto{\pgfqpoint{2.631241in}{0.777356in}}{\pgfqpoint{2.636249in}{0.774600in}}%
\pgfpathquadraticcurveto{\pgfqpoint{2.641256in}{0.771844in}}{\pgfqpoint{2.644462in}{0.766278in}}%
\pgfpathlineto{\pgfqpoint{2.644462in}{0.775843in}}%
\pgfpathlineto{\pgfqpoint{2.654819in}{0.775843in}}%
\pgfpathlineto{\pgfqpoint{2.654819in}{0.720617in}}%
\pgfpathlineto{\pgfqpoint{2.654819in}{0.720617in}}%
\pgfpathclose%
\pgfpathmoveto{\pgfqpoint{2.728579in}{0.750860in}}%
\pgfpathlineto{\pgfqpoint{2.728579in}{0.712800in}}%
\pgfpathlineto{\pgfqpoint{2.718222in}{0.712800in}}%
\pgfpathlineto{\pgfqpoint{2.718222in}{0.750517in}}%
\pgfpathquadraticcurveto{\pgfqpoint{2.718222in}{0.759469in}}{\pgfqpoint{2.714728in}{0.763900in}}%
\pgfpathquadraticcurveto{\pgfqpoint{2.711233in}{0.768349in}}{\pgfqpoint{2.704263in}{0.768349in}}%
\pgfpathquadraticcurveto{\pgfqpoint{2.695869in}{0.768349in}}{\pgfqpoint{2.691024in}{0.763000in}}%
\pgfpathquadraticcurveto{\pgfqpoint{2.686178in}{0.757668in}}{\pgfqpoint{2.686178in}{0.748428in}}%
\pgfpathlineto{\pgfqpoint{2.686178in}{0.712800in}}%
\pgfpathlineto{\pgfqpoint{2.675767in}{0.712800in}}%
\pgfpathlineto{\pgfqpoint{2.675767in}{0.800393in}}%
\pgfpathlineto{\pgfqpoint{2.686178in}{0.800393in}}%
\pgfpathlineto{\pgfqpoint{2.686178in}{0.766044in}}%
\pgfpathquadraticcurveto{\pgfqpoint{2.689907in}{0.771736in}}{\pgfqpoint{2.694932in}{0.774546in}}%
\pgfpathquadraticcurveto{\pgfqpoint{2.699976in}{0.777356in}}{\pgfqpoint{2.706568in}{0.777356in}}%
\pgfpathquadraticcurveto{\pgfqpoint{2.717429in}{0.777356in}}{\pgfqpoint{2.722995in}{0.770637in}}%
\pgfpathquadraticcurveto{\pgfqpoint{2.728579in}{0.763918in}}{\pgfqpoint{2.728579in}{0.750860in}}%
\pgfpathlineto{\pgfqpoint{2.728579in}{0.750860in}}%
\pgfpathclose%
\pgfpathmoveto{\pgfqpoint{2.759470in}{0.793747in}}%
\pgfpathlineto{\pgfqpoint{2.759470in}{0.775843in}}%
\pgfpathlineto{\pgfqpoint{2.780796in}{0.775843in}}%
\pgfpathlineto{\pgfqpoint{2.780796in}{0.767791in}}%
\pgfpathlineto{\pgfqpoint{2.759470in}{0.767791in}}%
\pgfpathlineto{\pgfqpoint{2.759470in}{0.733568in}}%
\pgfpathquadraticcurveto{\pgfqpoint{2.759470in}{0.725859in}}{\pgfqpoint{2.761577in}{0.723661in}}%
\pgfpathquadraticcurveto{\pgfqpoint{2.763685in}{0.721464in}}{\pgfqpoint{2.770169in}{0.721464in}}%
\pgfpathlineto{\pgfqpoint{2.780796in}{0.721464in}}%
\pgfpathlineto{\pgfqpoint{2.780796in}{0.712800in}}%
\pgfpathlineto{\pgfqpoint{2.770169in}{0.712800in}}%
\pgfpathquadraticcurveto{\pgfqpoint{2.758173in}{0.712800in}}{\pgfqpoint{2.753616in}{0.717267in}}%
\pgfpathquadraticcurveto{\pgfqpoint{2.749059in}{0.721752in}}{\pgfqpoint{2.749059in}{0.733568in}}%
\pgfpathlineto{\pgfqpoint{2.749059in}{0.767791in}}%
\pgfpathlineto{\pgfqpoint{2.741458in}{0.767791in}}%
\pgfpathlineto{\pgfqpoint{2.741458in}{0.775843in}}%
\pgfpathlineto{\pgfqpoint{2.749059in}{0.775843in}}%
\pgfpathlineto{\pgfqpoint{2.749059in}{0.793747in}}%
\pgfpathlineto{\pgfqpoint{2.759470in}{0.793747in}}%
\pgfpathlineto{\pgfqpoint{2.759470in}{0.793747in}}%
\pgfpathclose%
\pgfpathmoveto{\pgfqpoint{2.834598in}{0.773987in}}%
\pgfpathlineto{\pgfqpoint{2.834598in}{0.764189in}}%
\pgfpathquadraticcurveto{\pgfqpoint{2.830221in}{0.766440in}}{\pgfqpoint{2.825484in}{0.767557in}}%
\pgfpathquadraticcurveto{\pgfqpoint{2.820765in}{0.768692in}}{\pgfqpoint{2.815686in}{0.768692in}}%
\pgfpathquadraticcurveto{\pgfqpoint{2.807976in}{0.768692in}}{\pgfqpoint{2.804122in}{0.766332in}}%
\pgfpathquadraticcurveto{\pgfqpoint{2.800267in}{0.763973in}}{\pgfqpoint{2.800267in}{0.759235in}}%
\pgfpathquadraticcurveto{\pgfqpoint{2.800267in}{0.755633in}}{\pgfqpoint{2.803023in}{0.753580in}}%
\pgfpathquadraticcurveto{\pgfqpoint{2.805779in}{0.751526in}}{\pgfqpoint{2.814119in}{0.749671in}}%
\pgfpathlineto{\pgfqpoint{2.817667in}{0.748878in}}%
\pgfpathquadraticcurveto{\pgfqpoint{2.828690in}{0.746519in}}{\pgfqpoint{2.833338in}{0.742214in}}%
\pgfpathquadraticcurveto{\pgfqpoint{2.837985in}{0.737909in}}{\pgfqpoint{2.837985in}{0.730200in}}%
\pgfpathquadraticcurveto{\pgfqpoint{2.837985in}{0.721410in}}{\pgfqpoint{2.831032in}{0.716276in}}%
\pgfpathquadraticcurveto{\pgfqpoint{2.824079in}{0.711161in}}{\pgfqpoint{2.811921in}{0.711161in}}%
\pgfpathquadraticcurveto{\pgfqpoint{2.806860in}{0.711161in}}{\pgfqpoint{2.801366in}{0.712152in}}%
\pgfpathquadraticcurveto{\pgfqpoint{2.795872in}{0.713142in}}{\pgfqpoint{2.789802in}{0.715106in}}%
\pgfpathlineto{\pgfqpoint{2.789802in}{0.725805in}}%
\pgfpathquadraticcurveto{\pgfqpoint{2.795548in}{0.722815in}}{\pgfqpoint{2.801114in}{0.721320in}}%
\pgfpathquadraticcurveto{\pgfqpoint{2.806680in}{0.719843in}}{\pgfqpoint{2.812155in}{0.719843in}}%
\pgfpathquadraticcurveto{\pgfqpoint{2.819468in}{0.719843in}}{\pgfqpoint{2.823395in}{0.722346in}}%
\pgfpathquadraticcurveto{\pgfqpoint{2.827340in}{0.724850in}}{\pgfqpoint{2.827340in}{0.729407in}}%
\pgfpathquadraticcurveto{\pgfqpoint{2.827340in}{0.733622in}}{\pgfqpoint{2.824494in}{0.735874in}}%
\pgfpathquadraticcurveto{\pgfqpoint{2.821666in}{0.738125in}}{\pgfqpoint{2.812029in}{0.740214in}}%
\pgfpathlineto{\pgfqpoint{2.808427in}{0.741061in}}%
\pgfpathquadraticcurveto{\pgfqpoint{2.798808in}{0.743078in}}{\pgfqpoint{2.794521in}{0.747275in}}%
\pgfpathquadraticcurveto{\pgfqpoint{2.790253in}{0.751472in}}{\pgfqpoint{2.790253in}{0.758785in}}%
\pgfpathquadraticcurveto{\pgfqpoint{2.790253in}{0.767683in}}{\pgfqpoint{2.796557in}{0.772510in}}%
\pgfpathquadraticcurveto{\pgfqpoint{2.802861in}{0.777356in}}{\pgfqpoint{2.814461in}{0.777356in}}%
\pgfpathquadraticcurveto{\pgfqpoint{2.820189in}{0.777356in}}{\pgfqpoint{2.825250in}{0.776509in}}%
\pgfpathquadraticcurveto{\pgfqpoint{2.830330in}{0.775680in}}{\pgfqpoint{2.834598in}{0.773987in}}%
\pgfpathlineto{\pgfqpoint{2.834598in}{0.773987in}}%
\pgfpathclose%
\pgfusepath{fill}%
\end{pgfscope}%
\begin{pgfscope}%
\pgfsetbuttcap%
\pgfsetroundjoin%
\definecolor{currentfill}{rgb}{0.309804,0.372549,0.419608}%
\pgfsetfillcolor{currentfill}%
\pgfsetlinewidth{0.000000pt}%
\definecolor{currentstroke}{rgb}{0.309804,0.372549,0.419608}%
\pgfsetstrokecolor{currentstroke}%
\pgfsetdash{}{0pt}%
\pgfpathmoveto{\pgfqpoint{2.962705in}{0.744267in}}%
\pgfpathquadraticcurveto{\pgfqpoint{2.962705in}{0.755687in}}{\pgfqpoint{2.958003in}{0.762189in}}%
\pgfpathquadraticcurveto{\pgfqpoint{2.953302in}{0.768692in}}{\pgfqpoint{2.945089in}{0.768692in}}%
\pgfpathquadraticcurveto{\pgfqpoint{2.936857in}{0.768692in}}{\pgfqpoint{2.932156in}{0.762189in}}%
\pgfpathquadraticcurveto{\pgfqpoint{2.927455in}{0.755687in}}{\pgfqpoint{2.927455in}{0.744267in}}%
\pgfpathquadraticcurveto{\pgfqpoint{2.927455in}{0.732848in}}{\pgfqpoint{2.932156in}{0.726345in}}%
\pgfpathquadraticcurveto{\pgfqpoint{2.936857in}{0.719843in}}{\pgfqpoint{2.945089in}{0.719843in}}%
\pgfpathquadraticcurveto{\pgfqpoint{2.953302in}{0.719843in}}{\pgfqpoint{2.958003in}{0.726345in}}%
\pgfpathquadraticcurveto{\pgfqpoint{2.962705in}{0.732848in}}{\pgfqpoint{2.962705in}{0.744267in}}%
\pgfpathlineto{\pgfqpoint{2.962705in}{0.744267in}}%
\pgfpathclose%
\pgfpathmoveto{\pgfqpoint{2.927455in}{0.766278in}}%
\pgfpathquadraticcurveto{\pgfqpoint{2.930733in}{0.771898in}}{\pgfqpoint{2.935704in}{0.774618in}}%
\pgfpathquadraticcurveto{\pgfqpoint{2.940694in}{0.777356in}}{\pgfqpoint{2.947610in}{0.777356in}}%
\pgfpathquadraticcurveto{\pgfqpoint{2.959102in}{0.777356in}}{\pgfqpoint{2.966271in}{0.768241in}}%
\pgfpathquadraticcurveto{\pgfqpoint{2.973458in}{0.759127in}}{\pgfqpoint{2.973458in}{0.744267in}}%
\pgfpathquadraticcurveto{\pgfqpoint{2.973458in}{0.729407in}}{\pgfqpoint{2.966271in}{0.720275in}}%
\pgfpathquadraticcurveto{\pgfqpoint{2.959102in}{0.711161in}}{\pgfqpoint{2.947610in}{0.711161in}}%
\pgfpathquadraticcurveto{\pgfqpoint{2.940694in}{0.711161in}}{\pgfqpoint{2.935704in}{0.713899in}}%
\pgfpathquadraticcurveto{\pgfqpoint{2.930733in}{0.716637in}}{\pgfqpoint{2.927455in}{0.722256in}}%
\pgfpathlineto{\pgfqpoint{2.927455in}{0.712800in}}%
\pgfpathlineto{\pgfqpoint{2.917044in}{0.712800in}}%
\pgfpathlineto{\pgfqpoint{2.917044in}{0.800393in}}%
\pgfpathlineto{\pgfqpoint{2.927455in}{0.800393in}}%
\pgfpathlineto{\pgfqpoint{2.927455in}{0.766278in}}%
\pgfpathlineto{\pgfqpoint{2.927455in}{0.766278in}}%
\pgfpathclose%
\pgfpathmoveto{\pgfqpoint{3.016849in}{0.706946in}}%
\pgfpathquadraticcurveto{\pgfqpoint{3.012472in}{0.695688in}}{\pgfqpoint{3.008293in}{0.692266in}}%
\pgfpathquadraticcurveto{\pgfqpoint{3.004132in}{0.688826in}}{\pgfqpoint{2.997162in}{0.688826in}}%
\pgfpathlineto{\pgfqpoint{2.988876in}{0.688826in}}%
\pgfpathlineto{\pgfqpoint{2.988876in}{0.697490in}}%
\pgfpathlineto{\pgfqpoint{2.994964in}{0.697490in}}%
\pgfpathquadraticcurveto{\pgfqpoint{2.999233in}{0.697490in}}{\pgfqpoint{3.001593in}{0.699525in}}%
\pgfpathquadraticcurveto{\pgfqpoint{3.003970in}{0.701542in}}{\pgfqpoint{3.006834in}{0.709089in}}%
\pgfpathlineto{\pgfqpoint{3.008690in}{0.713809in}}%
\pgfpathlineto{\pgfqpoint{2.983202in}{0.775843in}}%
\pgfpathlineto{\pgfqpoint{2.994172in}{0.775843in}}%
\pgfpathlineto{\pgfqpoint{3.013877in}{0.726543in}}%
\pgfpathlineto{\pgfqpoint{3.033582in}{0.775843in}}%
\pgfpathlineto{\pgfqpoint{3.044552in}{0.775843in}}%
\pgfpathlineto{\pgfqpoint{3.016849in}{0.706946in}}%
\pgfpathlineto{\pgfqpoint{3.016849in}{0.706946in}}%
\pgfpathclose%
\pgfusepath{fill}%
\end{pgfscope}%
\begin{pgfscope}%
\pgfsetbuttcap%
\pgfsetroundjoin%
\definecolor{currentfill}{rgb}{0.309804,0.372549,0.419608}%
\pgfsetfillcolor{currentfill}%
\pgfsetlinewidth{0.000000pt}%
\definecolor{currentstroke}{rgb}{0.309804,0.372549,0.419608}%
\pgfsetstrokecolor{currentstroke}%
\pgfsetdash{}{0pt}%
\pgfpathmoveto{\pgfqpoint{3.119382in}{0.793747in}}%
\pgfpathlineto{\pgfqpoint{3.119382in}{0.775843in}}%
\pgfpathlineto{\pgfqpoint{3.140708in}{0.775843in}}%
\pgfpathlineto{\pgfqpoint{3.140708in}{0.767791in}}%
\pgfpathlineto{\pgfqpoint{3.119382in}{0.767791in}}%
\pgfpathlineto{\pgfqpoint{3.119382in}{0.733568in}}%
\pgfpathquadraticcurveto{\pgfqpoint{3.119382in}{0.725859in}}{\pgfqpoint{3.121489in}{0.723661in}}%
\pgfpathquadraticcurveto{\pgfqpoint{3.123597in}{0.721464in}}{\pgfqpoint{3.130081in}{0.721464in}}%
\pgfpathlineto{\pgfqpoint{3.140708in}{0.721464in}}%
\pgfpathlineto{\pgfqpoint{3.140708in}{0.712800in}}%
\pgfpathlineto{\pgfqpoint{3.130081in}{0.712800in}}%
\pgfpathquadraticcurveto{\pgfqpoint{3.118085in}{0.712800in}}{\pgfqpoint{3.113528in}{0.717267in}}%
\pgfpathquadraticcurveto{\pgfqpoint{3.108971in}{0.721752in}}{\pgfqpoint{3.108971in}{0.733568in}}%
\pgfpathlineto{\pgfqpoint{3.108971in}{0.767791in}}%
\pgfpathlineto{\pgfqpoint{3.101370in}{0.767791in}}%
\pgfpathlineto{\pgfqpoint{3.101370in}{0.775843in}}%
\pgfpathlineto{\pgfqpoint{3.108971in}{0.775843in}}%
\pgfpathlineto{\pgfqpoint{3.108971in}{0.793747in}}%
\pgfpathlineto{\pgfqpoint{3.119382in}{0.793747in}}%
\pgfpathlineto{\pgfqpoint{3.119382in}{0.793747in}}%
\pgfpathclose%
\pgfpathmoveto{\pgfqpoint{3.206741in}{0.750860in}}%
\pgfpathlineto{\pgfqpoint{3.206741in}{0.712800in}}%
\pgfpathlineto{\pgfqpoint{3.196384in}{0.712800in}}%
\pgfpathlineto{\pgfqpoint{3.196384in}{0.750517in}}%
\pgfpathquadraticcurveto{\pgfqpoint{3.196384in}{0.759469in}}{\pgfqpoint{3.192889in}{0.763900in}}%
\pgfpathquadraticcurveto{\pgfqpoint{3.189395in}{0.768349in}}{\pgfqpoint{3.182424in}{0.768349in}}%
\pgfpathquadraticcurveto{\pgfqpoint{3.174031in}{0.768349in}}{\pgfqpoint{3.169185in}{0.763000in}}%
\pgfpathquadraticcurveto{\pgfqpoint{3.164340in}{0.757668in}}{\pgfqpoint{3.164340in}{0.748428in}}%
\pgfpathlineto{\pgfqpoint{3.164340in}{0.712800in}}%
\pgfpathlineto{\pgfqpoint{3.153929in}{0.712800in}}%
\pgfpathlineto{\pgfqpoint{3.153929in}{0.800393in}}%
\pgfpathlineto{\pgfqpoint{3.164340in}{0.800393in}}%
\pgfpathlineto{\pgfqpoint{3.164340in}{0.766044in}}%
\pgfpathquadraticcurveto{\pgfqpoint{3.168069in}{0.771736in}}{\pgfqpoint{3.173094in}{0.774546in}}%
\pgfpathquadraticcurveto{\pgfqpoint{3.178137in}{0.777356in}}{\pgfqpoint{3.184730in}{0.777356in}}%
\pgfpathquadraticcurveto{\pgfqpoint{3.195591in}{0.777356in}}{\pgfqpoint{3.201157in}{0.770637in}}%
\pgfpathquadraticcurveto{\pgfqpoint{3.206741in}{0.763918in}}{\pgfqpoint{3.206741in}{0.750860in}}%
\pgfpathlineto{\pgfqpoint{3.206741in}{0.750860in}}%
\pgfpathclose%
\pgfpathmoveto{\pgfqpoint{3.281311in}{0.746915in}}%
\pgfpathlineto{\pgfqpoint{3.281311in}{0.741854in}}%
\pgfpathlineto{\pgfqpoint{3.233687in}{0.741854in}}%
\pgfpathquadraticcurveto{\pgfqpoint{3.234371in}{0.731154in}}{\pgfqpoint{3.240135in}{0.725553in}}%
\pgfpathquadraticcurveto{\pgfqpoint{3.245899in}{0.719951in}}{\pgfqpoint{3.256202in}{0.719951in}}%
\pgfpathquadraticcurveto{\pgfqpoint{3.262164in}{0.719951in}}{\pgfqpoint{3.267766in}{0.721410in}}%
\pgfpathquadraticcurveto{\pgfqpoint{3.273368in}{0.722869in}}{\pgfqpoint{3.278897in}{0.725805in}}%
\pgfpathlineto{\pgfqpoint{3.278897in}{0.716006in}}%
\pgfpathquadraticcurveto{\pgfqpoint{3.273314in}{0.713647in}}{\pgfqpoint{3.267460in}{0.712404in}}%
\pgfpathquadraticcurveto{\pgfqpoint{3.261606in}{0.711161in}}{\pgfqpoint{3.255590in}{0.711161in}}%
\pgfpathquadraticcurveto{\pgfqpoint{3.240496in}{0.711161in}}{\pgfqpoint{3.231688in}{0.719933in}}%
\pgfpathquadraticcurveto{\pgfqpoint{3.222880in}{0.728723in}}{\pgfqpoint{3.222880in}{0.743709in}}%
\pgfpathquadraticcurveto{\pgfqpoint{3.222880in}{0.759181in}}{\pgfqpoint{3.231237in}{0.768259in}}%
\pgfpathquadraticcurveto{\pgfqpoint{3.239595in}{0.777356in}}{\pgfqpoint{3.253788in}{0.777356in}}%
\pgfpathquadraticcurveto{\pgfqpoint{3.266505in}{0.777356in}}{\pgfqpoint{3.273908in}{0.769160in}}%
\pgfpathquadraticcurveto{\pgfqpoint{3.281311in}{0.760983in}}{\pgfqpoint{3.281311in}{0.746915in}}%
\pgfpathlineto{\pgfqpoint{3.281311in}{0.746915in}}%
\pgfpathclose%
\pgfpathmoveto{\pgfqpoint{3.270954in}{0.749959in}}%
\pgfpathquadraticcurveto{\pgfqpoint{3.270846in}{0.758443in}}{\pgfqpoint{3.266199in}{0.763504in}}%
\pgfpathquadraticcurveto{\pgfqpoint{3.261552in}{0.768584in}}{\pgfqpoint{3.253897in}{0.768584in}}%
\pgfpathquadraticcurveto{\pgfqpoint{3.245233in}{0.768584in}}{\pgfqpoint{3.240027in}{0.763684in}}%
\pgfpathquadraticcurveto{\pgfqpoint{3.234822in}{0.758785in}}{\pgfqpoint{3.234029in}{0.749887in}}%
\pgfpathlineto{\pgfqpoint{3.270954in}{0.749959in}}%
\pgfpathlineto{\pgfqpoint{3.270954in}{0.749959in}}%
\pgfpathclose%
\pgfusepath{fill}%
\end{pgfscope}%
\begin{pgfscope}%
\pgfsetbuttcap%
\pgfsetroundjoin%
\definecolor{currentfill}{rgb}{0.309804,0.372549,0.419608}%
\pgfsetfillcolor{currentfill}%
\pgfsetlinewidth{0.000000pt}%
\definecolor{currentstroke}{rgb}{0.309804,0.372549,0.419608}%
\pgfsetstrokecolor{currentstroke}%
\pgfsetdash{}{0pt}%
\pgfpathmoveto{\pgfqpoint{3.404754in}{0.746915in}}%
\pgfpathlineto{\pgfqpoint{3.404754in}{0.741854in}}%
\pgfpathlineto{\pgfqpoint{3.357130in}{0.741854in}}%
\pgfpathquadraticcurveto{\pgfqpoint{3.357814in}{0.731154in}}{\pgfqpoint{3.363578in}{0.725553in}}%
\pgfpathquadraticcurveto{\pgfqpoint{3.369342in}{0.719951in}}{\pgfqpoint{3.379645in}{0.719951in}}%
\pgfpathquadraticcurveto{\pgfqpoint{3.385607in}{0.719951in}}{\pgfqpoint{3.391209in}{0.721410in}}%
\pgfpathquadraticcurveto{\pgfqpoint{3.396811in}{0.722869in}}{\pgfqpoint{3.402341in}{0.725805in}}%
\pgfpathlineto{\pgfqpoint{3.402341in}{0.716006in}}%
\pgfpathquadraticcurveto{\pgfqpoint{3.396757in}{0.713647in}}{\pgfqpoint{3.390903in}{0.712404in}}%
\pgfpathquadraticcurveto{\pgfqpoint{3.385049in}{0.711161in}}{\pgfqpoint{3.379033in}{0.711161in}}%
\pgfpathquadraticcurveto{\pgfqpoint{3.363939in}{0.711161in}}{\pgfqpoint{3.355131in}{0.719933in}}%
\pgfpathquadraticcurveto{\pgfqpoint{3.346323in}{0.728723in}}{\pgfqpoint{3.346323in}{0.743709in}}%
\pgfpathquadraticcurveto{\pgfqpoint{3.346323in}{0.759181in}}{\pgfqpoint{3.354680in}{0.768259in}}%
\pgfpathquadraticcurveto{\pgfqpoint{3.363038in}{0.777356in}}{\pgfqpoint{3.377232in}{0.777356in}}%
\pgfpathquadraticcurveto{\pgfqpoint{3.389948in}{0.777356in}}{\pgfqpoint{3.397351in}{0.769160in}}%
\pgfpathquadraticcurveto{\pgfqpoint{3.404754in}{0.760983in}}{\pgfqpoint{3.404754in}{0.746915in}}%
\pgfpathlineto{\pgfqpoint{3.404754in}{0.746915in}}%
\pgfpathclose%
\pgfpathmoveto{\pgfqpoint{3.394397in}{0.749959in}}%
\pgfpathquadraticcurveto{\pgfqpoint{3.394289in}{0.758443in}}{\pgfqpoint{3.389642in}{0.763504in}}%
\pgfpathquadraticcurveto{\pgfqpoint{3.384995in}{0.768584in}}{\pgfqpoint{3.377340in}{0.768584in}}%
\pgfpathquadraticcurveto{\pgfqpoint{3.368676in}{0.768584in}}{\pgfqpoint{3.363470in}{0.763684in}}%
\pgfpathquadraticcurveto{\pgfqpoint{3.358265in}{0.758785in}}{\pgfqpoint{3.357472in}{0.749887in}}%
\pgfpathlineto{\pgfqpoint{3.394397in}{0.749959in}}%
\pgfpathlineto{\pgfqpoint{3.394397in}{0.749959in}}%
\pgfpathclose%
\pgfpathmoveto{\pgfqpoint{3.472156in}{0.775843in}}%
\pgfpathlineto{\pgfqpoint{3.449352in}{0.745168in}}%
\pgfpathlineto{\pgfqpoint{3.473326in}{0.712800in}}%
\pgfpathlineto{\pgfqpoint{3.461114in}{0.712800in}}%
\pgfpathlineto{\pgfqpoint{3.442760in}{0.737567in}}%
\pgfpathlineto{\pgfqpoint{3.424423in}{0.712800in}}%
\pgfpathlineto{\pgfqpoint{3.412193in}{0.712800in}}%
\pgfpathlineto{\pgfqpoint{3.436690in}{0.745780in}}%
\pgfpathlineto{\pgfqpoint{3.414283in}{0.775843in}}%
\pgfpathlineto{\pgfqpoint{3.426495in}{0.775843in}}%
\pgfpathlineto{\pgfqpoint{3.443210in}{0.753381in}}%
\pgfpathlineto{\pgfqpoint{3.459925in}{0.775843in}}%
\pgfpathlineto{\pgfqpoint{3.472156in}{0.775843in}}%
\pgfpathlineto{\pgfqpoint{3.472156in}{0.775843in}}%
\pgfpathclose%
\pgfpathmoveto{\pgfqpoint{3.487970in}{0.775843in}}%
\pgfpathlineto{\pgfqpoint{3.498327in}{0.775843in}}%
\pgfpathlineto{\pgfqpoint{3.498327in}{0.712800in}}%
\pgfpathlineto{\pgfqpoint{3.487970in}{0.712800in}}%
\pgfpathlineto{\pgfqpoint{3.487970in}{0.775843in}}%
\pgfpathlineto{\pgfqpoint{3.487970in}{0.775843in}}%
\pgfpathclose%
\pgfpathmoveto{\pgfqpoint{3.487970in}{0.800393in}}%
\pgfpathlineto{\pgfqpoint{3.498327in}{0.800393in}}%
\pgfpathlineto{\pgfqpoint{3.498327in}{0.787262in}}%
\pgfpathlineto{\pgfqpoint{3.487970in}{0.787262in}}%
\pgfpathlineto{\pgfqpoint{3.487970in}{0.800393in}}%
\pgfpathlineto{\pgfqpoint{3.487970in}{0.800393in}}%
\pgfpathclose%
\pgfpathmoveto{\pgfqpoint{3.560181in}{0.773987in}}%
\pgfpathlineto{\pgfqpoint{3.560181in}{0.764189in}}%
\pgfpathquadraticcurveto{\pgfqpoint{3.555804in}{0.766440in}}{\pgfqpoint{3.551067in}{0.767557in}}%
\pgfpathquadraticcurveto{\pgfqpoint{3.546348in}{0.768692in}}{\pgfqpoint{3.541268in}{0.768692in}}%
\pgfpathquadraticcurveto{\pgfqpoint{3.533559in}{0.768692in}}{\pgfqpoint{3.529704in}{0.766332in}}%
\pgfpathquadraticcurveto{\pgfqpoint{3.525850in}{0.763973in}}{\pgfqpoint{3.525850in}{0.759235in}}%
\pgfpathquadraticcurveto{\pgfqpoint{3.525850in}{0.755633in}}{\pgfqpoint{3.528606in}{0.753580in}}%
\pgfpathquadraticcurveto{\pgfqpoint{3.531362in}{0.751526in}}{\pgfqpoint{3.539701in}{0.749671in}}%
\pgfpathlineto{\pgfqpoint{3.543250in}{0.748878in}}%
\pgfpathquadraticcurveto{\pgfqpoint{3.554273in}{0.746519in}}{\pgfqpoint{3.558920in}{0.742214in}}%
\pgfpathquadraticcurveto{\pgfqpoint{3.563567in}{0.737909in}}{\pgfqpoint{3.563567in}{0.730200in}}%
\pgfpathquadraticcurveto{\pgfqpoint{3.563567in}{0.721410in}}{\pgfqpoint{3.556615in}{0.716276in}}%
\pgfpathquadraticcurveto{\pgfqpoint{3.549662in}{0.711161in}}{\pgfqpoint{3.537504in}{0.711161in}}%
\pgfpathquadraticcurveto{\pgfqpoint{3.532442in}{0.711161in}}{\pgfqpoint{3.526949in}{0.712152in}}%
\pgfpathquadraticcurveto{\pgfqpoint{3.521455in}{0.713142in}}{\pgfqpoint{3.515385in}{0.715106in}}%
\pgfpathlineto{\pgfqpoint{3.515385in}{0.725805in}}%
\pgfpathquadraticcurveto{\pgfqpoint{3.521131in}{0.722815in}}{\pgfqpoint{3.526696in}{0.721320in}}%
\pgfpathquadraticcurveto{\pgfqpoint{3.532262in}{0.719843in}}{\pgfqpoint{3.537738in}{0.719843in}}%
\pgfpathquadraticcurveto{\pgfqpoint{3.545051in}{0.719843in}}{\pgfqpoint{3.548977in}{0.722346in}}%
\pgfpathquadraticcurveto{\pgfqpoint{3.552922in}{0.724850in}}{\pgfqpoint{3.552922in}{0.729407in}}%
\pgfpathquadraticcurveto{\pgfqpoint{3.552922in}{0.733622in}}{\pgfqpoint{3.550076in}{0.735874in}}%
\pgfpathquadraticcurveto{\pgfqpoint{3.547248in}{0.738125in}}{\pgfqpoint{3.537612in}{0.740214in}}%
\pgfpathlineto{\pgfqpoint{3.534009in}{0.741061in}}%
\pgfpathquadraticcurveto{\pgfqpoint{3.524391in}{0.743078in}}{\pgfqpoint{3.520104in}{0.747275in}}%
\pgfpathquadraticcurveto{\pgfqpoint{3.515835in}{0.751472in}}{\pgfqpoint{3.515835in}{0.758785in}}%
\pgfpathquadraticcurveto{\pgfqpoint{3.515835in}{0.767683in}}{\pgfqpoint{3.522139in}{0.772510in}}%
\pgfpathquadraticcurveto{\pgfqpoint{3.528444in}{0.777356in}}{\pgfqpoint{3.540043in}{0.777356in}}%
\pgfpathquadraticcurveto{\pgfqpoint{3.545771in}{0.777356in}}{\pgfqpoint{3.550833in}{0.776509in}}%
\pgfpathquadraticcurveto{\pgfqpoint{3.555912in}{0.775680in}}{\pgfqpoint{3.560181in}{0.773987in}}%
\pgfpathlineto{\pgfqpoint{3.560181in}{0.773987in}}%
\pgfpathclose%
\pgfpathmoveto{\pgfqpoint{3.590297in}{0.793747in}}%
\pgfpathlineto{\pgfqpoint{3.590297in}{0.775843in}}%
\pgfpathlineto{\pgfqpoint{3.611624in}{0.775843in}}%
\pgfpathlineto{\pgfqpoint{3.611624in}{0.767791in}}%
\pgfpathlineto{\pgfqpoint{3.590297in}{0.767791in}}%
\pgfpathlineto{\pgfqpoint{3.590297in}{0.733568in}}%
\pgfpathquadraticcurveto{\pgfqpoint{3.590297in}{0.725859in}}{\pgfqpoint{3.592405in}{0.723661in}}%
\pgfpathquadraticcurveto{\pgfqpoint{3.594512in}{0.721464in}}{\pgfqpoint{3.600997in}{0.721464in}}%
\pgfpathlineto{\pgfqpoint{3.611624in}{0.721464in}}%
\pgfpathlineto{\pgfqpoint{3.611624in}{0.712800in}}%
\pgfpathlineto{\pgfqpoint{3.600997in}{0.712800in}}%
\pgfpathquadraticcurveto{\pgfqpoint{3.589000in}{0.712800in}}{\pgfqpoint{3.584443in}{0.717267in}}%
\pgfpathquadraticcurveto{\pgfqpoint{3.579886in}{0.721752in}}{\pgfqpoint{3.579886in}{0.733568in}}%
\pgfpathlineto{\pgfqpoint{3.579886in}{0.767791in}}%
\pgfpathlineto{\pgfqpoint{3.572285in}{0.767791in}}%
\pgfpathlineto{\pgfqpoint{3.572285in}{0.775843in}}%
\pgfpathlineto{\pgfqpoint{3.579886in}{0.775843in}}%
\pgfpathlineto{\pgfqpoint{3.579886in}{0.793747in}}%
\pgfpathlineto{\pgfqpoint{3.590297in}{0.793747in}}%
\pgfpathlineto{\pgfqpoint{3.590297in}{0.793747in}}%
\pgfpathclose%
\pgfpathmoveto{\pgfqpoint{3.625241in}{0.775843in}}%
\pgfpathlineto{\pgfqpoint{3.635598in}{0.775843in}}%
\pgfpathlineto{\pgfqpoint{3.635598in}{0.712800in}}%
\pgfpathlineto{\pgfqpoint{3.625241in}{0.712800in}}%
\pgfpathlineto{\pgfqpoint{3.625241in}{0.775843in}}%
\pgfpathlineto{\pgfqpoint{3.625241in}{0.775843in}}%
\pgfpathclose%
\pgfpathmoveto{\pgfqpoint{3.625241in}{0.800393in}}%
\pgfpathlineto{\pgfqpoint{3.635598in}{0.800393in}}%
\pgfpathlineto{\pgfqpoint{3.635598in}{0.787262in}}%
\pgfpathlineto{\pgfqpoint{3.625241in}{0.787262in}}%
\pgfpathlineto{\pgfqpoint{3.625241in}{0.800393in}}%
\pgfpathlineto{\pgfqpoint{3.625241in}{0.800393in}}%
\pgfpathclose%
\pgfpathmoveto{\pgfqpoint{3.709682in}{0.750860in}}%
\pgfpathlineto{\pgfqpoint{3.709682in}{0.712800in}}%
\pgfpathlineto{\pgfqpoint{3.699325in}{0.712800in}}%
\pgfpathlineto{\pgfqpoint{3.699325in}{0.750517in}}%
\pgfpathquadraticcurveto{\pgfqpoint{3.699325in}{0.759469in}}{\pgfqpoint{3.695831in}{0.763900in}}%
\pgfpathquadraticcurveto{\pgfqpoint{3.692336in}{0.768349in}}{\pgfqpoint{3.685365in}{0.768349in}}%
\pgfpathquadraticcurveto{\pgfqpoint{3.676972in}{0.768349in}}{\pgfqpoint{3.672127in}{0.763000in}}%
\pgfpathquadraticcurveto{\pgfqpoint{3.667281in}{0.757668in}}{\pgfqpoint{3.667281in}{0.748428in}}%
\pgfpathlineto{\pgfqpoint{3.667281in}{0.712800in}}%
\pgfpathlineto{\pgfqpoint{3.656870in}{0.712800in}}%
\pgfpathlineto{\pgfqpoint{3.656870in}{0.775843in}}%
\pgfpathlineto{\pgfqpoint{3.667281in}{0.775843in}}%
\pgfpathlineto{\pgfqpoint{3.667281in}{0.766044in}}%
\pgfpathquadraticcurveto{\pgfqpoint{3.671010in}{0.771736in}}{\pgfqpoint{3.676035in}{0.774546in}}%
\pgfpathquadraticcurveto{\pgfqpoint{3.681079in}{0.777356in}}{\pgfqpoint{3.687671in}{0.777356in}}%
\pgfpathquadraticcurveto{\pgfqpoint{3.698532in}{0.777356in}}{\pgfqpoint{3.704098in}{0.770637in}}%
\pgfpathquadraticcurveto{\pgfqpoint{3.709682in}{0.763918in}}{\pgfqpoint{3.709682in}{0.750860in}}%
\pgfpathlineto{\pgfqpoint{3.709682in}{0.750860in}}%
\pgfpathclose%
\pgfpathmoveto{\pgfqpoint{3.771806in}{0.745060in}}%
\pgfpathquadraticcurveto{\pgfqpoint{3.771806in}{0.756317in}}{\pgfqpoint{3.767159in}{0.762496in}}%
\pgfpathquadraticcurveto{\pgfqpoint{3.762530in}{0.768692in}}{\pgfqpoint{3.754136in}{0.768692in}}%
\pgfpathquadraticcurveto{\pgfqpoint{3.745814in}{0.768692in}}{\pgfqpoint{3.741167in}{0.762496in}}%
\pgfpathquadraticcurveto{\pgfqpoint{3.736520in}{0.756317in}}{\pgfqpoint{3.736520in}{0.745060in}}%
\pgfpathquadraticcurveto{\pgfqpoint{3.736520in}{0.733856in}}{\pgfqpoint{3.741167in}{0.727660in}}%
\pgfpathquadraticcurveto{\pgfqpoint{3.745814in}{0.721464in}}{\pgfqpoint{3.754136in}{0.721464in}}%
\pgfpathquadraticcurveto{\pgfqpoint{3.762530in}{0.721464in}}{\pgfqpoint{3.767159in}{0.727660in}}%
\pgfpathquadraticcurveto{\pgfqpoint{3.771806in}{0.733856in}}{\pgfqpoint{3.771806in}{0.745060in}}%
\pgfpathlineto{\pgfqpoint{3.771806in}{0.745060in}}%
\pgfpathclose%
\pgfpathmoveto{\pgfqpoint{3.782163in}{0.720617in}}%
\pgfpathquadraticcurveto{\pgfqpoint{3.782163in}{0.704532in}}{\pgfqpoint{3.775012in}{0.696679in}}%
\pgfpathquadraticcurveto{\pgfqpoint{3.767879in}{0.688826in}}{\pgfqpoint{3.753127in}{0.688826in}}%
\pgfpathquadraticcurveto{\pgfqpoint{3.747670in}{0.688826in}}{\pgfqpoint{3.742824in}{0.689636in}}%
\pgfpathquadraticcurveto{\pgfqpoint{3.737979in}{0.690447in}}{\pgfqpoint{3.733422in}{0.692140in}}%
\pgfpathlineto{\pgfqpoint{3.733422in}{0.702209in}}%
\pgfpathquadraticcurveto{\pgfqpoint{3.737979in}{0.699741in}}{\pgfqpoint{3.742428in}{0.698570in}}%
\pgfpathquadraticcurveto{\pgfqpoint{3.746877in}{0.697382in}}{\pgfqpoint{3.751488in}{0.697382in}}%
\pgfpathquadraticcurveto{\pgfqpoint{3.761683in}{0.697382in}}{\pgfqpoint{3.766744in}{0.702695in}}%
\pgfpathquadraticcurveto{\pgfqpoint{3.771806in}{0.708009in}}{\pgfqpoint{3.771806in}{0.718762in}}%
\pgfpathlineto{\pgfqpoint{3.771806in}{0.723895in}}%
\pgfpathquadraticcurveto{\pgfqpoint{3.768600in}{0.718312in}}{\pgfqpoint{3.763592in}{0.715556in}}%
\pgfpathquadraticcurveto{\pgfqpoint{3.758585in}{0.712800in}}{\pgfqpoint{3.751596in}{0.712800in}}%
\pgfpathquadraticcurveto{\pgfqpoint{3.740014in}{0.712800in}}{\pgfqpoint{3.732918in}{0.721626in}}%
\pgfpathquadraticcurveto{\pgfqpoint{3.725821in}{0.730470in}}{\pgfqpoint{3.725821in}{0.745060in}}%
\pgfpathquadraticcurveto{\pgfqpoint{3.725821in}{0.759686in}}{\pgfqpoint{3.732918in}{0.768512in}}%
\pgfpathquadraticcurveto{\pgfqpoint{3.740014in}{0.777356in}}{\pgfqpoint{3.751596in}{0.777356in}}%
\pgfpathquadraticcurveto{\pgfqpoint{3.758585in}{0.777356in}}{\pgfqpoint{3.763592in}{0.774600in}}%
\pgfpathquadraticcurveto{\pgfqpoint{3.768600in}{0.771844in}}{\pgfqpoint{3.771806in}{0.766278in}}%
\pgfpathlineto{\pgfqpoint{3.771806in}{0.775843in}}%
\pgfpathlineto{\pgfqpoint{3.782163in}{0.775843in}}%
\pgfpathlineto{\pgfqpoint{3.782163in}{0.720617in}}%
\pgfpathlineto{\pgfqpoint{3.782163in}{0.720617in}}%
\pgfpathclose%
\pgfusepath{fill}%
\end{pgfscope}%
\begin{pgfscope}%
\pgfsetbuttcap%
\pgfsetroundjoin%
\definecolor{currentfill}{rgb}{0.309804,0.372549,0.419608}%
\pgfsetfillcolor{currentfill}%
\pgfsetlinewidth{0.000000pt}%
\definecolor{currentstroke}{rgb}{0.309804,0.372549,0.419608}%
\pgfsetstrokecolor{currentstroke}%
\pgfsetdash{}{0pt}%
\pgfpathmoveto{\pgfqpoint{3.905752in}{0.744267in}}%
\pgfpathquadraticcurveto{\pgfqpoint{3.905752in}{0.755687in}}{\pgfqpoint{3.901051in}{0.762189in}}%
\pgfpathquadraticcurveto{\pgfqpoint{3.896350in}{0.768692in}}{\pgfqpoint{3.888136in}{0.768692in}}%
\pgfpathquadraticcurveto{\pgfqpoint{3.879905in}{0.768692in}}{\pgfqpoint{3.875203in}{0.762189in}}%
\pgfpathquadraticcurveto{\pgfqpoint{3.870502in}{0.755687in}}{\pgfqpoint{3.870502in}{0.744267in}}%
\pgfpathquadraticcurveto{\pgfqpoint{3.870502in}{0.732848in}}{\pgfqpoint{3.875203in}{0.726345in}}%
\pgfpathquadraticcurveto{\pgfqpoint{3.879905in}{0.719843in}}{\pgfqpoint{3.888136in}{0.719843in}}%
\pgfpathquadraticcurveto{\pgfqpoint{3.896350in}{0.719843in}}{\pgfqpoint{3.901051in}{0.726345in}}%
\pgfpathquadraticcurveto{\pgfqpoint{3.905752in}{0.732848in}}{\pgfqpoint{3.905752in}{0.744267in}}%
\pgfpathlineto{\pgfqpoint{3.905752in}{0.744267in}}%
\pgfpathclose%
\pgfpathmoveto{\pgfqpoint{3.870502in}{0.766278in}}%
\pgfpathquadraticcurveto{\pgfqpoint{3.873780in}{0.771898in}}{\pgfqpoint{3.878752in}{0.774618in}}%
\pgfpathquadraticcurveto{\pgfqpoint{3.883741in}{0.777356in}}{\pgfqpoint{3.890658in}{0.777356in}}%
\pgfpathquadraticcurveto{\pgfqpoint{3.902150in}{0.777356in}}{\pgfqpoint{3.909318in}{0.768241in}}%
\pgfpathquadraticcurveto{\pgfqpoint{3.916505in}{0.759127in}}{\pgfqpoint{3.916505in}{0.744267in}}%
\pgfpathquadraticcurveto{\pgfqpoint{3.916505in}{0.729407in}}{\pgfqpoint{3.909318in}{0.720275in}}%
\pgfpathquadraticcurveto{\pgfqpoint{3.902150in}{0.711161in}}{\pgfqpoint{3.890658in}{0.711161in}}%
\pgfpathquadraticcurveto{\pgfqpoint{3.883741in}{0.711161in}}{\pgfqpoint{3.878752in}{0.713899in}}%
\pgfpathquadraticcurveto{\pgfqpoint{3.873780in}{0.716637in}}{\pgfqpoint{3.870502in}{0.722256in}}%
\pgfpathlineto{\pgfqpoint{3.870502in}{0.712800in}}%
\pgfpathlineto{\pgfqpoint{3.860091in}{0.712800in}}%
\pgfpathlineto{\pgfqpoint{3.860091in}{0.800393in}}%
\pgfpathlineto{\pgfqpoint{3.870502in}{0.800393in}}%
\pgfpathlineto{\pgfqpoint{3.870502in}{0.766278in}}%
\pgfpathlineto{\pgfqpoint{3.870502in}{0.766278in}}%
\pgfpathclose%
\pgfpathmoveto{\pgfqpoint{3.958095in}{0.768584in}}%
\pgfpathquadraticcurveto{\pgfqpoint{3.949774in}{0.768584in}}{\pgfqpoint{3.944928in}{0.762081in}}%
\pgfpathquadraticcurveto{\pgfqpoint{3.940083in}{0.755579in}}{\pgfqpoint{3.940083in}{0.744267in}}%
\pgfpathquadraticcurveto{\pgfqpoint{3.940083in}{0.732956in}}{\pgfqpoint{3.944892in}{0.726453in}}%
\pgfpathquadraticcurveto{\pgfqpoint{3.949720in}{0.719951in}}{\pgfqpoint{3.958095in}{0.719951in}}%
\pgfpathquadraticcurveto{\pgfqpoint{3.966381in}{0.719951in}}{\pgfqpoint{3.971208in}{0.726471in}}%
\pgfpathquadraticcurveto{\pgfqpoint{3.976053in}{0.733010in}}{\pgfqpoint{3.976053in}{0.744267in}}%
\pgfpathquadraticcurveto{\pgfqpoint{3.976053in}{0.755471in}}{\pgfqpoint{3.971208in}{0.762027in}}%
\pgfpathquadraticcurveto{\pgfqpoint{3.966381in}{0.768584in}}{\pgfqpoint{3.958095in}{0.768584in}}%
\pgfpathlineto{\pgfqpoint{3.958095in}{0.768584in}}%
\pgfpathclose%
\pgfpathmoveto{\pgfqpoint{3.958095in}{0.777356in}}%
\pgfpathquadraticcurveto{\pgfqpoint{3.971604in}{0.777356in}}{\pgfqpoint{3.979314in}{0.768566in}}%
\pgfpathquadraticcurveto{\pgfqpoint{3.987041in}{0.759794in}}{\pgfqpoint{3.987041in}{0.744267in}}%
\pgfpathquadraticcurveto{\pgfqpoint{3.987041in}{0.728795in}}{\pgfqpoint{3.979314in}{0.719969in}}%
\pgfpathquadraticcurveto{\pgfqpoint{3.971604in}{0.711161in}}{\pgfqpoint{3.958095in}{0.711161in}}%
\pgfpathquadraticcurveto{\pgfqpoint{3.944532in}{0.711161in}}{\pgfqpoint{3.936841in}{0.719969in}}%
\pgfpathquadraticcurveto{\pgfqpoint{3.929168in}{0.728795in}}{\pgfqpoint{3.929168in}{0.744267in}}%
\pgfpathquadraticcurveto{\pgfqpoint{3.929168in}{0.759794in}}{\pgfqpoint{3.936841in}{0.768566in}}%
\pgfpathquadraticcurveto{\pgfqpoint{3.944532in}{0.777356in}}{\pgfqpoint{3.958095in}{0.777356in}}%
\pgfpathlineto{\pgfqpoint{3.958095in}{0.777356in}}%
\pgfpathclose%
\pgfpathmoveto{\pgfqpoint{4.028631in}{0.768584in}}%
\pgfpathquadraticcurveto{\pgfqpoint{4.020309in}{0.768584in}}{\pgfqpoint{4.015464in}{0.762081in}}%
\pgfpathquadraticcurveto{\pgfqpoint{4.010619in}{0.755579in}}{\pgfqpoint{4.010619in}{0.744267in}}%
\pgfpathquadraticcurveto{\pgfqpoint{4.010619in}{0.732956in}}{\pgfqpoint{4.015428in}{0.726453in}}%
\pgfpathquadraticcurveto{\pgfqpoint{4.020255in}{0.719951in}}{\pgfqpoint{4.028631in}{0.719951in}}%
\pgfpathquadraticcurveto{\pgfqpoint{4.036916in}{0.719951in}}{\pgfqpoint{4.041744in}{0.726471in}}%
\pgfpathquadraticcurveto{\pgfqpoint{4.046589in}{0.733010in}}{\pgfqpoint{4.046589in}{0.744267in}}%
\pgfpathquadraticcurveto{\pgfqpoint{4.046589in}{0.755471in}}{\pgfqpoint{4.041744in}{0.762027in}}%
\pgfpathquadraticcurveto{\pgfqpoint{4.036916in}{0.768584in}}{\pgfqpoint{4.028631in}{0.768584in}}%
\pgfpathlineto{\pgfqpoint{4.028631in}{0.768584in}}%
\pgfpathclose%
\pgfpathmoveto{\pgfqpoint{4.028631in}{0.777356in}}%
\pgfpathquadraticcurveto{\pgfqpoint{4.042140in}{0.777356in}}{\pgfqpoint{4.049849in}{0.768566in}}%
\pgfpathquadraticcurveto{\pgfqpoint{4.057576in}{0.759794in}}{\pgfqpoint{4.057576in}{0.744267in}}%
\pgfpathquadraticcurveto{\pgfqpoint{4.057576in}{0.728795in}}{\pgfqpoint{4.049849in}{0.719969in}}%
\pgfpathquadraticcurveto{\pgfqpoint{4.042140in}{0.711161in}}{\pgfqpoint{4.028631in}{0.711161in}}%
\pgfpathquadraticcurveto{\pgfqpoint{4.015068in}{0.711161in}}{\pgfqpoint{4.007377in}{0.719969in}}%
\pgfpathquadraticcurveto{\pgfqpoint{3.999703in}{0.728795in}}{\pgfqpoint{3.999703in}{0.744267in}}%
\pgfpathquadraticcurveto{\pgfqpoint{3.999703in}{0.759794in}}{\pgfqpoint{4.007377in}{0.768566in}}%
\pgfpathquadraticcurveto{\pgfqpoint{4.015068in}{0.777356in}}{\pgfqpoint{4.028631in}{0.777356in}}%
\pgfpathlineto{\pgfqpoint{4.028631in}{0.777356in}}%
\pgfpathclose%
\pgfpathmoveto{\pgfqpoint{4.084991in}{0.793747in}}%
\pgfpathlineto{\pgfqpoint{4.084991in}{0.775843in}}%
\pgfpathlineto{\pgfqpoint{4.106317in}{0.775843in}}%
\pgfpathlineto{\pgfqpoint{4.106317in}{0.767791in}}%
\pgfpathlineto{\pgfqpoint{4.084991in}{0.767791in}}%
\pgfpathlineto{\pgfqpoint{4.084991in}{0.733568in}}%
\pgfpathquadraticcurveto{\pgfqpoint{4.084991in}{0.725859in}}{\pgfqpoint{4.087098in}{0.723661in}}%
\pgfpathquadraticcurveto{\pgfqpoint{4.089206in}{0.721464in}}{\pgfqpoint{4.095690in}{0.721464in}}%
\pgfpathlineto{\pgfqpoint{4.106317in}{0.721464in}}%
\pgfpathlineto{\pgfqpoint{4.106317in}{0.712800in}}%
\pgfpathlineto{\pgfqpoint{4.095690in}{0.712800in}}%
\pgfpathquadraticcurveto{\pgfqpoint{4.083694in}{0.712800in}}{\pgfqpoint{4.079137in}{0.717267in}}%
\pgfpathquadraticcurveto{\pgfqpoint{4.074580in}{0.721752in}}{\pgfqpoint{4.074580in}{0.733568in}}%
\pgfpathlineto{\pgfqpoint{4.074580in}{0.767791in}}%
\pgfpathlineto{\pgfqpoint{4.066979in}{0.767791in}}%
\pgfpathlineto{\pgfqpoint{4.066979in}{0.775843in}}%
\pgfpathlineto{\pgfqpoint{4.074580in}{0.775843in}}%
\pgfpathlineto{\pgfqpoint{4.074580in}{0.793747in}}%
\pgfpathlineto{\pgfqpoint{4.084991in}{0.793747in}}%
\pgfpathlineto{\pgfqpoint{4.084991in}{0.793747in}}%
\pgfpathclose%
\pgfpathmoveto{\pgfqpoint{4.160120in}{0.773987in}}%
\pgfpathlineto{\pgfqpoint{4.160120in}{0.764189in}}%
\pgfpathquadraticcurveto{\pgfqpoint{4.155743in}{0.766440in}}{\pgfqpoint{4.151005in}{0.767557in}}%
\pgfpathquadraticcurveto{\pgfqpoint{4.146286in}{0.768692in}}{\pgfqpoint{4.141207in}{0.768692in}}%
\pgfpathquadraticcurveto{\pgfqpoint{4.133498in}{0.768692in}}{\pgfqpoint{4.129643in}{0.766332in}}%
\pgfpathquadraticcurveto{\pgfqpoint{4.125788in}{0.763973in}}{\pgfqpoint{4.125788in}{0.759235in}}%
\pgfpathquadraticcurveto{\pgfqpoint{4.125788in}{0.755633in}}{\pgfqpoint{4.128544in}{0.753580in}}%
\pgfpathquadraticcurveto{\pgfqpoint{4.131300in}{0.751526in}}{\pgfqpoint{4.139640in}{0.749671in}}%
\pgfpathlineto{\pgfqpoint{4.143188in}{0.748878in}}%
\pgfpathquadraticcurveto{\pgfqpoint{4.154212in}{0.746519in}}{\pgfqpoint{4.158859in}{0.742214in}}%
\pgfpathquadraticcurveto{\pgfqpoint{4.163506in}{0.737909in}}{\pgfqpoint{4.163506in}{0.730200in}}%
\pgfpathquadraticcurveto{\pgfqpoint{4.163506in}{0.721410in}}{\pgfqpoint{4.156553in}{0.716276in}}%
\pgfpathquadraticcurveto{\pgfqpoint{4.149600in}{0.711161in}}{\pgfqpoint{4.137442in}{0.711161in}}%
\pgfpathquadraticcurveto{\pgfqpoint{4.132381in}{0.711161in}}{\pgfqpoint{4.126887in}{0.712152in}}%
\pgfpathquadraticcurveto{\pgfqpoint{4.121393in}{0.713142in}}{\pgfqpoint{4.115323in}{0.715106in}}%
\pgfpathlineto{\pgfqpoint{4.115323in}{0.725805in}}%
\pgfpathquadraticcurveto{\pgfqpoint{4.121069in}{0.722815in}}{\pgfqpoint{4.126635in}{0.721320in}}%
\pgfpathquadraticcurveto{\pgfqpoint{4.132201in}{0.719843in}}{\pgfqpoint{4.137676in}{0.719843in}}%
\pgfpathquadraticcurveto{\pgfqpoint{4.144989in}{0.719843in}}{\pgfqpoint{4.148916in}{0.722346in}}%
\pgfpathquadraticcurveto{\pgfqpoint{4.152861in}{0.724850in}}{\pgfqpoint{4.152861in}{0.729407in}}%
\pgfpathquadraticcurveto{\pgfqpoint{4.152861in}{0.733622in}}{\pgfqpoint{4.150015in}{0.735874in}}%
\pgfpathquadraticcurveto{\pgfqpoint{4.147187in}{0.738125in}}{\pgfqpoint{4.137550in}{0.740214in}}%
\pgfpathlineto{\pgfqpoint{4.133948in}{0.741061in}}%
\pgfpathquadraticcurveto{\pgfqpoint{4.124329in}{0.743078in}}{\pgfqpoint{4.120043in}{0.747275in}}%
\pgfpathquadraticcurveto{\pgfqpoint{4.115774in}{0.751472in}}{\pgfqpoint{4.115774in}{0.758785in}}%
\pgfpathquadraticcurveto{\pgfqpoint{4.115774in}{0.767683in}}{\pgfqpoint{4.122078in}{0.772510in}}%
\pgfpathquadraticcurveto{\pgfqpoint{4.128382in}{0.777356in}}{\pgfqpoint{4.139982in}{0.777356in}}%
\pgfpathquadraticcurveto{\pgfqpoint{4.145710in}{0.777356in}}{\pgfqpoint{4.150771in}{0.776509in}}%
\pgfpathquadraticcurveto{\pgfqpoint{4.155851in}{0.775680in}}{\pgfqpoint{4.160120in}{0.773987in}}%
\pgfpathlineto{\pgfqpoint{4.160120in}{0.773987in}}%
\pgfpathclose%
\pgfpathmoveto{\pgfqpoint{4.190236in}{0.793747in}}%
\pgfpathlineto{\pgfqpoint{4.190236in}{0.775843in}}%
\pgfpathlineto{\pgfqpoint{4.211562in}{0.775843in}}%
\pgfpathlineto{\pgfqpoint{4.211562in}{0.767791in}}%
\pgfpathlineto{\pgfqpoint{4.190236in}{0.767791in}}%
\pgfpathlineto{\pgfqpoint{4.190236in}{0.733568in}}%
\pgfpathquadraticcurveto{\pgfqpoint{4.190236in}{0.725859in}}{\pgfqpoint{4.192343in}{0.723661in}}%
\pgfpathquadraticcurveto{\pgfqpoint{4.194451in}{0.721464in}}{\pgfqpoint{4.200935in}{0.721464in}}%
\pgfpathlineto{\pgfqpoint{4.211562in}{0.721464in}}%
\pgfpathlineto{\pgfqpoint{4.211562in}{0.712800in}}%
\pgfpathlineto{\pgfqpoint{4.200935in}{0.712800in}}%
\pgfpathquadraticcurveto{\pgfqpoint{4.188939in}{0.712800in}}{\pgfqpoint{4.184382in}{0.717267in}}%
\pgfpathquadraticcurveto{\pgfqpoint{4.179825in}{0.721752in}}{\pgfqpoint{4.179825in}{0.733568in}}%
\pgfpathlineto{\pgfqpoint{4.179825in}{0.767791in}}%
\pgfpathlineto{\pgfqpoint{4.172224in}{0.767791in}}%
\pgfpathlineto{\pgfqpoint{4.172224in}{0.775843in}}%
\pgfpathlineto{\pgfqpoint{4.179825in}{0.775843in}}%
\pgfpathlineto{\pgfqpoint{4.179825in}{0.793747in}}%
\pgfpathlineto{\pgfqpoint{4.190236in}{0.793747in}}%
\pgfpathlineto{\pgfqpoint{4.190236in}{0.793747in}}%
\pgfpathclose%
\pgfpathmoveto{\pgfqpoint{4.261708in}{0.766170in}}%
\pgfpathquadraticcurveto{\pgfqpoint{4.259961in}{0.767179in}}{\pgfqpoint{4.257908in}{0.767647in}}%
\pgfpathquadraticcurveto{\pgfqpoint{4.255854in}{0.768133in}}{\pgfqpoint{4.253386in}{0.768133in}}%
\pgfpathquadraticcurveto{\pgfqpoint{4.244597in}{0.768133in}}{\pgfqpoint{4.239895in}{0.762423in}}%
\pgfpathquadraticcurveto{\pgfqpoint{4.235194in}{0.756714in}}{\pgfqpoint{4.235194in}{0.746014in}}%
\pgfpathlineto{\pgfqpoint{4.235194in}{0.712800in}}%
\pgfpathlineto{\pgfqpoint{4.224783in}{0.712800in}}%
\pgfpathlineto{\pgfqpoint{4.224783in}{0.775843in}}%
\pgfpathlineto{\pgfqpoint{4.235194in}{0.775843in}}%
\pgfpathlineto{\pgfqpoint{4.235194in}{0.766044in}}%
\pgfpathquadraticcurveto{\pgfqpoint{4.238472in}{0.771790in}}{\pgfqpoint{4.243696in}{0.774564in}}%
\pgfpathquadraticcurveto{\pgfqpoint{4.248937in}{0.777356in}}{\pgfqpoint{4.256431in}{0.777356in}}%
\pgfpathquadraticcurveto{\pgfqpoint{4.257493in}{0.777356in}}{\pgfqpoint{4.258790in}{0.777211in}}%
\pgfpathquadraticcurveto{\pgfqpoint{4.260087in}{0.777085in}}{\pgfqpoint{4.261654in}{0.776797in}}%
\pgfpathlineto{\pgfqpoint{4.261708in}{0.766170in}}%
\pgfpathlineto{\pgfqpoint{4.261708in}{0.766170in}}%
\pgfpathclose%
\pgfpathmoveto{\pgfqpoint{4.301227in}{0.744483in}}%
\pgfpathquadraticcurveto{\pgfqpoint{4.288672in}{0.744483in}}{\pgfqpoint{4.283827in}{0.741619in}}%
\pgfpathquadraticcurveto{\pgfqpoint{4.278982in}{0.738756in}}{\pgfqpoint{4.278982in}{0.731821in}}%
\pgfpathquadraticcurveto{\pgfqpoint{4.278982in}{0.726309in}}{\pgfqpoint{4.282620in}{0.723067in}}%
\pgfpathquadraticcurveto{\pgfqpoint{4.286259in}{0.719843in}}{\pgfqpoint{4.292491in}{0.719843in}}%
\pgfpathquadraticcurveto{\pgfqpoint{4.301119in}{0.719843in}}{\pgfqpoint{4.306324in}{0.725949in}}%
\pgfpathquadraticcurveto{\pgfqpoint{4.311530in}{0.732055in}}{\pgfqpoint{4.311530in}{0.742178in}}%
\pgfpathlineto{\pgfqpoint{4.311530in}{0.744483in}}%
\pgfpathlineto{\pgfqpoint{4.301227in}{0.744483in}}%
\pgfpathlineto{\pgfqpoint{4.301227in}{0.744483in}}%
\pgfpathclose%
\pgfpathmoveto{\pgfqpoint{4.321887in}{0.748770in}}%
\pgfpathlineto{\pgfqpoint{4.321887in}{0.712800in}}%
\pgfpathlineto{\pgfqpoint{4.311530in}{0.712800in}}%
\pgfpathlineto{\pgfqpoint{4.311530in}{0.722364in}}%
\pgfpathquadraticcurveto{\pgfqpoint{4.307981in}{0.716637in}}{\pgfqpoint{4.302686in}{0.713899in}}%
\pgfpathquadraticcurveto{\pgfqpoint{4.297390in}{0.711161in}}{\pgfqpoint{4.289735in}{0.711161in}}%
\pgfpathquadraticcurveto{\pgfqpoint{4.280062in}{0.711161in}}{\pgfqpoint{4.274335in}{0.716601in}}%
\pgfpathquadraticcurveto{\pgfqpoint{4.268625in}{0.722040in}}{\pgfqpoint{4.268625in}{0.731154in}}%
\pgfpathquadraticcurveto{\pgfqpoint{4.268625in}{0.741782in}}{\pgfqpoint{4.275740in}{0.747185in}}%
\pgfpathquadraticcurveto{\pgfqpoint{4.282872in}{0.752589in}}{\pgfqpoint{4.296994in}{0.752589in}}%
\pgfpathlineto{\pgfqpoint{4.311530in}{0.752589in}}%
\pgfpathlineto{\pgfqpoint{4.311530in}{0.753616in}}%
\pgfpathquadraticcurveto{\pgfqpoint{4.311530in}{0.760766in}}{\pgfqpoint{4.306829in}{0.764675in}}%
\pgfpathquadraticcurveto{\pgfqpoint{4.302127in}{0.768584in}}{\pgfqpoint{4.293626in}{0.768584in}}%
\pgfpathquadraticcurveto{\pgfqpoint{4.288222in}{0.768584in}}{\pgfqpoint{4.283089in}{0.767287in}}%
\pgfpathquadraticcurveto{\pgfqpoint{4.277973in}{0.765990in}}{\pgfqpoint{4.273254in}{0.763396in}}%
\pgfpathlineto{\pgfqpoint{4.273254in}{0.772979in}}%
\pgfpathquadraticcurveto{\pgfqpoint{4.278928in}{0.775176in}}{\pgfqpoint{4.284277in}{0.776257in}}%
\pgfpathquadraticcurveto{\pgfqpoint{4.289627in}{0.777356in}}{\pgfqpoint{4.294688in}{0.777356in}}%
\pgfpathquadraticcurveto{\pgfqpoint{4.308378in}{0.777356in}}{\pgfqpoint{4.315132in}{0.770259in}}%
\pgfpathquadraticcurveto{\pgfqpoint{4.321887in}{0.763180in}}{\pgfqpoint{4.321887in}{0.748770in}}%
\pgfpathlineto{\pgfqpoint{4.321887in}{0.748770in}}%
\pgfpathclose%
\pgfpathmoveto{\pgfqpoint{4.353228in}{0.722256in}}%
\pgfpathlineto{\pgfqpoint{4.353228in}{0.688826in}}%
\pgfpathlineto{\pgfqpoint{4.342817in}{0.688826in}}%
\pgfpathlineto{\pgfqpoint{4.342817in}{0.775843in}}%
\pgfpathlineto{\pgfqpoint{4.353228in}{0.775843in}}%
\pgfpathlineto{\pgfqpoint{4.353228in}{0.766278in}}%
\pgfpathquadraticcurveto{\pgfqpoint{4.356506in}{0.771898in}}{\pgfqpoint{4.361477in}{0.774618in}}%
\pgfpathquadraticcurveto{\pgfqpoint{4.366467in}{0.777356in}}{\pgfqpoint{4.373383in}{0.777356in}}%
\pgfpathquadraticcurveto{\pgfqpoint{4.384875in}{0.777356in}}{\pgfqpoint{4.392044in}{0.768241in}}%
\pgfpathquadraticcurveto{\pgfqpoint{4.399231in}{0.759127in}}{\pgfqpoint{4.399231in}{0.744267in}}%
\pgfpathquadraticcurveto{\pgfqpoint{4.399231in}{0.729407in}}{\pgfqpoint{4.392044in}{0.720275in}}%
\pgfpathquadraticcurveto{\pgfqpoint{4.384875in}{0.711161in}}{\pgfqpoint{4.373383in}{0.711161in}}%
\pgfpathquadraticcurveto{\pgfqpoint{4.366467in}{0.711161in}}{\pgfqpoint{4.361477in}{0.713899in}}%
\pgfpathquadraticcurveto{\pgfqpoint{4.356506in}{0.716637in}}{\pgfqpoint{4.353228in}{0.722256in}}%
\pgfpathlineto{\pgfqpoint{4.353228in}{0.722256in}}%
\pgfpathclose%
\pgfpathmoveto{\pgfqpoint{4.388478in}{0.744267in}}%
\pgfpathquadraticcurveto{\pgfqpoint{4.388478in}{0.755687in}}{\pgfqpoint{4.383776in}{0.762189in}}%
\pgfpathquadraticcurveto{\pgfqpoint{4.379075in}{0.768692in}}{\pgfqpoint{4.370862in}{0.768692in}}%
\pgfpathquadraticcurveto{\pgfqpoint{4.362630in}{0.768692in}}{\pgfqpoint{4.357929in}{0.762189in}}%
\pgfpathquadraticcurveto{\pgfqpoint{4.353228in}{0.755687in}}{\pgfqpoint{4.353228in}{0.744267in}}%
\pgfpathquadraticcurveto{\pgfqpoint{4.353228in}{0.732848in}}{\pgfqpoint{4.357929in}{0.726345in}}%
\pgfpathquadraticcurveto{\pgfqpoint{4.362630in}{0.719843in}}{\pgfqpoint{4.370862in}{0.719843in}}%
\pgfpathquadraticcurveto{\pgfqpoint{4.379075in}{0.719843in}}{\pgfqpoint{4.383776in}{0.726345in}}%
\pgfpathquadraticcurveto{\pgfqpoint{4.388478in}{0.732848in}}{\pgfqpoint{4.388478in}{0.744267in}}%
\pgfpathlineto{\pgfqpoint{4.388478in}{0.744267in}}%
\pgfpathclose%
\pgfusepath{fill}%
\end{pgfscope}%
\begin{pgfscope}%
\pgfsetbuttcap%
\pgfsetroundjoin%
\definecolor{currentfill}{rgb}{0.309804,0.372549,0.419608}%
\pgfsetfillcolor{currentfill}%
\pgfsetlinewidth{0.000000pt}%
\definecolor{currentstroke}{rgb}{0.309804,0.372549,0.419608}%
\pgfsetstrokecolor{currentstroke}%
\pgfsetdash{}{0pt}%
\pgfpathmoveto{\pgfqpoint{4.527553in}{0.773429in}}%
\pgfpathlineto{\pgfqpoint{4.527553in}{0.763738in}}%
\pgfpathquadraticcurveto{\pgfqpoint{4.523158in}{0.766170in}}{\pgfqpoint{4.518745in}{0.767377in}}%
\pgfpathquadraticcurveto{\pgfqpoint{4.514332in}{0.768584in}}{\pgfqpoint{4.509829in}{0.768584in}}%
\pgfpathquadraticcurveto{\pgfqpoint{4.499742in}{0.768584in}}{\pgfqpoint{4.494158in}{0.762189in}}%
\pgfpathquadraticcurveto{\pgfqpoint{4.488593in}{0.755813in}}{\pgfqpoint{4.488593in}{0.744267in}}%
\pgfpathquadraticcurveto{\pgfqpoint{4.488593in}{0.732721in}}{\pgfqpoint{4.494158in}{0.726327in}}%
\pgfpathquadraticcurveto{\pgfqpoint{4.499742in}{0.719951in}}{\pgfqpoint{4.509829in}{0.719951in}}%
\pgfpathquadraticcurveto{\pgfqpoint{4.514332in}{0.719951in}}{\pgfqpoint{4.518745in}{0.721158in}}%
\pgfpathquadraticcurveto{\pgfqpoint{4.523158in}{0.722364in}}{\pgfqpoint{4.527553in}{0.724796in}}%
\pgfpathlineto{\pgfqpoint{4.527553in}{0.715214in}}%
\pgfpathquadraticcurveto{\pgfqpoint{4.523212in}{0.713196in}}{\pgfqpoint{4.518565in}{0.712188in}}%
\pgfpathquadraticcurveto{\pgfqpoint{4.513936in}{0.711161in}}{\pgfqpoint{4.508694in}{0.711161in}}%
\pgfpathquadraticcurveto{\pgfqpoint{4.494447in}{0.711161in}}{\pgfqpoint{4.486053in}{0.720113in}}%
\pgfpathquadraticcurveto{\pgfqpoint{4.477677in}{0.729065in}}{\pgfqpoint{4.477677in}{0.744267in}}%
\pgfpathquadraticcurveto{\pgfqpoint{4.477677in}{0.759686in}}{\pgfqpoint{4.486143in}{0.768512in}}%
\pgfpathquadraticcurveto{\pgfqpoint{4.494627in}{0.777356in}}{\pgfqpoint{4.509379in}{0.777356in}}%
\pgfpathquadraticcurveto{\pgfqpoint{4.514152in}{0.777356in}}{\pgfqpoint{4.518709in}{0.776365in}}%
\pgfpathquadraticcurveto{\pgfqpoint{4.523266in}{0.775392in}}{\pgfqpoint{4.527553in}{0.773429in}}%
\pgfpathlineto{\pgfqpoint{4.527553in}{0.773429in}}%
\pgfpathclose%
\pgfpathmoveto{\pgfqpoint{4.545565in}{0.800393in}}%
\pgfpathlineto{\pgfqpoint{4.555922in}{0.800393in}}%
\pgfpathlineto{\pgfqpoint{4.555922in}{0.712800in}}%
\pgfpathlineto{\pgfqpoint{4.545565in}{0.712800in}}%
\pgfpathlineto{\pgfqpoint{4.545565in}{0.800393in}}%
\pgfpathlineto{\pgfqpoint{4.545565in}{0.800393in}}%
\pgfpathclose%
\pgfpathmoveto{\pgfqpoint{4.576528in}{0.737675in}}%
\pgfpathlineto{\pgfqpoint{4.576528in}{0.775843in}}%
\pgfpathlineto{\pgfqpoint{4.586885in}{0.775843in}}%
\pgfpathlineto{\pgfqpoint{4.586885in}{0.738071in}}%
\pgfpathquadraticcurveto{\pgfqpoint{4.586885in}{0.729119in}}{\pgfqpoint{4.590361in}{0.724634in}}%
\pgfpathquadraticcurveto{\pgfqpoint{4.593856in}{0.720167in}}{\pgfqpoint{4.600844in}{0.720167in}}%
\pgfpathquadraticcurveto{\pgfqpoint{4.609220in}{0.720167in}}{\pgfqpoint{4.614083in}{0.725517in}}%
\pgfpathquadraticcurveto{\pgfqpoint{4.618965in}{0.730866in}}{\pgfqpoint{4.618965in}{0.740106in}}%
\pgfpathlineto{\pgfqpoint{4.618965in}{0.775843in}}%
\pgfpathlineto{\pgfqpoint{4.629322in}{0.775843in}}%
\pgfpathlineto{\pgfqpoint{4.629322in}{0.712800in}}%
\pgfpathlineto{\pgfqpoint{4.618965in}{0.712800in}}%
\pgfpathlineto{\pgfqpoint{4.618965in}{0.722491in}}%
\pgfpathquadraticcurveto{\pgfqpoint{4.615200in}{0.716745in}}{\pgfqpoint{4.610211in}{0.713953in}}%
\pgfpathquadraticcurveto{\pgfqpoint{4.605239in}{0.711161in}}{\pgfqpoint{4.598647in}{0.711161in}}%
\pgfpathquadraticcurveto{\pgfqpoint{4.587786in}{0.711161in}}{\pgfqpoint{4.582148in}{0.717915in}}%
\pgfpathquadraticcurveto{\pgfqpoint{4.576528in}{0.724670in}}{\pgfqpoint{4.576528in}{0.737675in}}%
\pgfpathlineto{\pgfqpoint{4.576528in}{0.737675in}}%
\pgfpathclose%
\pgfpathmoveto{\pgfqpoint{4.602592in}{0.777356in}}%
\pgfpathlineto{\pgfqpoint{4.602592in}{0.777356in}}%
\pgfpathlineto{\pgfqpoint{4.602592in}{0.777356in}}%
\pgfpathclose%
\pgfpathmoveto{\pgfqpoint{4.690833in}{0.773987in}}%
\pgfpathlineto{\pgfqpoint{4.690833in}{0.764189in}}%
\pgfpathquadraticcurveto{\pgfqpoint{4.686456in}{0.766440in}}{\pgfqpoint{4.681719in}{0.767557in}}%
\pgfpathquadraticcurveto{\pgfqpoint{4.677000in}{0.768692in}}{\pgfqpoint{4.671920in}{0.768692in}}%
\pgfpathquadraticcurveto{\pgfqpoint{4.664211in}{0.768692in}}{\pgfqpoint{4.660357in}{0.766332in}}%
\pgfpathquadraticcurveto{\pgfqpoint{4.656502in}{0.763973in}}{\pgfqpoint{4.656502in}{0.759235in}}%
\pgfpathquadraticcurveto{\pgfqpoint{4.656502in}{0.755633in}}{\pgfqpoint{4.659258in}{0.753580in}}%
\pgfpathquadraticcurveto{\pgfqpoint{4.662014in}{0.751526in}}{\pgfqpoint{4.670353in}{0.749671in}}%
\pgfpathlineto{\pgfqpoint{4.673902in}{0.748878in}}%
\pgfpathquadraticcurveto{\pgfqpoint{4.684925in}{0.746519in}}{\pgfqpoint{4.689572in}{0.742214in}}%
\pgfpathquadraticcurveto{\pgfqpoint{4.694219in}{0.737909in}}{\pgfqpoint{4.694219in}{0.730200in}}%
\pgfpathquadraticcurveto{\pgfqpoint{4.694219in}{0.721410in}}{\pgfqpoint{4.687267in}{0.716276in}}%
\pgfpathquadraticcurveto{\pgfqpoint{4.680314in}{0.711161in}}{\pgfqpoint{4.668156in}{0.711161in}}%
\pgfpathquadraticcurveto{\pgfqpoint{4.663094in}{0.711161in}}{\pgfqpoint{4.657601in}{0.712152in}}%
\pgfpathquadraticcurveto{\pgfqpoint{4.652107in}{0.713142in}}{\pgfqpoint{4.646037in}{0.715106in}}%
\pgfpathlineto{\pgfqpoint{4.646037in}{0.725805in}}%
\pgfpathquadraticcurveto{\pgfqpoint{4.651783in}{0.722815in}}{\pgfqpoint{4.657348in}{0.721320in}}%
\pgfpathquadraticcurveto{\pgfqpoint{4.662914in}{0.719843in}}{\pgfqpoint{4.668390in}{0.719843in}}%
\pgfpathquadraticcurveto{\pgfqpoint{4.675703in}{0.719843in}}{\pgfqpoint{4.679630in}{0.722346in}}%
\pgfpathquadraticcurveto{\pgfqpoint{4.683574in}{0.724850in}}{\pgfqpoint{4.683574in}{0.729407in}}%
\pgfpathquadraticcurveto{\pgfqpoint{4.683574in}{0.733622in}}{\pgfqpoint{4.680728in}{0.735874in}}%
\pgfpathquadraticcurveto{\pgfqpoint{4.677900in}{0.738125in}}{\pgfqpoint{4.668264in}{0.740214in}}%
\pgfpathlineto{\pgfqpoint{4.664661in}{0.741061in}}%
\pgfpathquadraticcurveto{\pgfqpoint{4.655043in}{0.743078in}}{\pgfqpoint{4.650756in}{0.747275in}}%
\pgfpathquadraticcurveto{\pgfqpoint{4.646487in}{0.751472in}}{\pgfqpoint{4.646487in}{0.758785in}}%
\pgfpathquadraticcurveto{\pgfqpoint{4.646487in}{0.767683in}}{\pgfqpoint{4.652791in}{0.772510in}}%
\pgfpathquadraticcurveto{\pgfqpoint{4.659096in}{0.777356in}}{\pgfqpoint{4.670695in}{0.777356in}}%
\pgfpathquadraticcurveto{\pgfqpoint{4.676423in}{0.777356in}}{\pgfqpoint{4.681485in}{0.776509in}}%
\pgfpathquadraticcurveto{\pgfqpoint{4.686564in}{0.775680in}}{\pgfqpoint{4.690833in}{0.773987in}}%
\pgfpathlineto{\pgfqpoint{4.690833in}{0.773987in}}%
\pgfpathclose%
\pgfpathmoveto{\pgfqpoint{4.720949in}{0.793747in}}%
\pgfpathlineto{\pgfqpoint{4.720949in}{0.775843in}}%
\pgfpathlineto{\pgfqpoint{4.742276in}{0.775843in}}%
\pgfpathlineto{\pgfqpoint{4.742276in}{0.767791in}}%
\pgfpathlineto{\pgfqpoint{4.720949in}{0.767791in}}%
\pgfpathlineto{\pgfqpoint{4.720949in}{0.733568in}}%
\pgfpathquadraticcurveto{\pgfqpoint{4.720949in}{0.725859in}}{\pgfqpoint{4.723057in}{0.723661in}}%
\pgfpathquadraticcurveto{\pgfqpoint{4.725164in}{0.721464in}}{\pgfqpoint{4.731649in}{0.721464in}}%
\pgfpathlineto{\pgfqpoint{4.742276in}{0.721464in}}%
\pgfpathlineto{\pgfqpoint{4.742276in}{0.712800in}}%
\pgfpathlineto{\pgfqpoint{4.731649in}{0.712800in}}%
\pgfpathquadraticcurveto{\pgfqpoint{4.719653in}{0.712800in}}{\pgfqpoint{4.715095in}{0.717267in}}%
\pgfpathquadraticcurveto{\pgfqpoint{4.710538in}{0.721752in}}{\pgfqpoint{4.710538in}{0.733568in}}%
\pgfpathlineto{\pgfqpoint{4.710538in}{0.767791in}}%
\pgfpathlineto{\pgfqpoint{4.702937in}{0.767791in}}%
\pgfpathlineto{\pgfqpoint{4.702937in}{0.775843in}}%
\pgfpathlineto{\pgfqpoint{4.710538in}{0.775843in}}%
\pgfpathlineto{\pgfqpoint{4.710538in}{0.793747in}}%
\pgfpathlineto{\pgfqpoint{4.720949in}{0.793747in}}%
\pgfpathlineto{\pgfqpoint{4.720949in}{0.793747in}}%
\pgfpathclose%
\pgfpathmoveto{\pgfqpoint{4.809821in}{0.746915in}}%
\pgfpathlineto{\pgfqpoint{4.809821in}{0.741854in}}%
\pgfpathlineto{\pgfqpoint{4.762197in}{0.741854in}}%
\pgfpathquadraticcurveto{\pgfqpoint{4.762882in}{0.731154in}}{\pgfqpoint{4.768646in}{0.725553in}}%
\pgfpathquadraticcurveto{\pgfqpoint{4.774409in}{0.719951in}}{\pgfqpoint{4.784712in}{0.719951in}}%
\pgfpathquadraticcurveto{\pgfqpoint{4.790674in}{0.719951in}}{\pgfqpoint{4.796276in}{0.721410in}}%
\pgfpathquadraticcurveto{\pgfqpoint{4.801878in}{0.722869in}}{\pgfqpoint{4.807408in}{0.725805in}}%
\pgfpathlineto{\pgfqpoint{4.807408in}{0.716006in}}%
\pgfpathquadraticcurveto{\pgfqpoint{4.801824in}{0.713647in}}{\pgfqpoint{4.795970in}{0.712404in}}%
\pgfpathquadraticcurveto{\pgfqpoint{4.790116in}{0.711161in}}{\pgfqpoint{4.784100in}{0.711161in}}%
\pgfpathquadraticcurveto{\pgfqpoint{4.769006in}{0.711161in}}{\pgfqpoint{4.760198in}{0.719933in}}%
\pgfpathquadraticcurveto{\pgfqpoint{4.751390in}{0.728723in}}{\pgfqpoint{4.751390in}{0.743709in}}%
\pgfpathquadraticcurveto{\pgfqpoint{4.751390in}{0.759181in}}{\pgfqpoint{4.759748in}{0.768259in}}%
\pgfpathquadraticcurveto{\pgfqpoint{4.768105in}{0.777356in}}{\pgfqpoint{4.782299in}{0.777356in}}%
\pgfpathquadraticcurveto{\pgfqpoint{4.795015in}{0.777356in}}{\pgfqpoint{4.802418in}{0.769160in}}%
\pgfpathquadraticcurveto{\pgfqpoint{4.809821in}{0.760983in}}{\pgfqpoint{4.809821in}{0.746915in}}%
\pgfpathlineto{\pgfqpoint{4.809821in}{0.746915in}}%
\pgfpathclose%
\pgfpathmoveto{\pgfqpoint{4.799464in}{0.749959in}}%
\pgfpathquadraticcurveto{\pgfqpoint{4.799356in}{0.758443in}}{\pgfqpoint{4.794709in}{0.763504in}}%
\pgfpathquadraticcurveto{\pgfqpoint{4.790062in}{0.768584in}}{\pgfqpoint{4.782407in}{0.768584in}}%
\pgfpathquadraticcurveto{\pgfqpoint{4.773743in}{0.768584in}}{\pgfqpoint{4.768537in}{0.763684in}}%
\pgfpathquadraticcurveto{\pgfqpoint{4.763332in}{0.758785in}}{\pgfqpoint{4.762539in}{0.749887in}}%
\pgfpathlineto{\pgfqpoint{4.799464in}{0.749959in}}%
\pgfpathlineto{\pgfqpoint{4.799464in}{0.749959in}}%
\pgfpathclose%
\pgfpathmoveto{\pgfqpoint{4.863353in}{0.766170in}}%
\pgfpathquadraticcurveto{\pgfqpoint{4.861606in}{0.767179in}}{\pgfqpoint{4.859553in}{0.767647in}}%
\pgfpathquadraticcurveto{\pgfqpoint{4.857500in}{0.768133in}}{\pgfqpoint{4.855032in}{0.768133in}}%
\pgfpathquadraticcurveto{\pgfqpoint{4.846242in}{0.768133in}}{\pgfqpoint{4.841541in}{0.762423in}}%
\pgfpathquadraticcurveto{\pgfqpoint{4.836840in}{0.756714in}}{\pgfqpoint{4.836840in}{0.746014in}}%
\pgfpathlineto{\pgfqpoint{4.836840in}{0.712800in}}%
\pgfpathlineto{\pgfqpoint{4.826429in}{0.712800in}}%
\pgfpathlineto{\pgfqpoint{4.826429in}{0.775843in}}%
\pgfpathlineto{\pgfqpoint{4.836840in}{0.775843in}}%
\pgfpathlineto{\pgfqpoint{4.836840in}{0.766044in}}%
\pgfpathquadraticcurveto{\pgfqpoint{4.840118in}{0.771790in}}{\pgfqpoint{4.845341in}{0.774564in}}%
\pgfpathquadraticcurveto{\pgfqpoint{4.850583in}{0.777356in}}{\pgfqpoint{4.858076in}{0.777356in}}%
\pgfpathquadraticcurveto{\pgfqpoint{4.859139in}{0.777356in}}{\pgfqpoint{4.860436in}{0.777211in}}%
\pgfpathquadraticcurveto{\pgfqpoint{4.861732in}{0.777085in}}{\pgfqpoint{4.863299in}{0.776797in}}%
\pgfpathlineto{\pgfqpoint{4.863353in}{0.766170in}}%
\pgfpathlineto{\pgfqpoint{4.863353in}{0.766170in}}%
\pgfpathclose%
\pgfpathmoveto{\pgfqpoint{4.914400in}{0.773987in}}%
\pgfpathlineto{\pgfqpoint{4.914400in}{0.764189in}}%
\pgfpathquadraticcurveto{\pgfqpoint{4.910023in}{0.766440in}}{\pgfqpoint{4.905286in}{0.767557in}}%
\pgfpathquadraticcurveto{\pgfqpoint{4.900567in}{0.768692in}}{\pgfqpoint{4.895487in}{0.768692in}}%
\pgfpathquadraticcurveto{\pgfqpoint{4.887778in}{0.768692in}}{\pgfqpoint{4.883923in}{0.766332in}}%
\pgfpathquadraticcurveto{\pgfqpoint{4.880069in}{0.763973in}}{\pgfqpoint{4.880069in}{0.759235in}}%
\pgfpathquadraticcurveto{\pgfqpoint{4.880069in}{0.755633in}}{\pgfqpoint{4.882825in}{0.753580in}}%
\pgfpathquadraticcurveto{\pgfqpoint{4.885580in}{0.751526in}}{\pgfqpoint{4.893920in}{0.749671in}}%
\pgfpathlineto{\pgfqpoint{4.897468in}{0.748878in}}%
\pgfpathquadraticcurveto{\pgfqpoint{4.908492in}{0.746519in}}{\pgfqpoint{4.913139in}{0.742214in}}%
\pgfpathquadraticcurveto{\pgfqpoint{4.917786in}{0.737909in}}{\pgfqpoint{4.917786in}{0.730200in}}%
\pgfpathquadraticcurveto{\pgfqpoint{4.917786in}{0.721410in}}{\pgfqpoint{4.910834in}{0.716276in}}%
\pgfpathquadraticcurveto{\pgfqpoint{4.903881in}{0.711161in}}{\pgfqpoint{4.891723in}{0.711161in}}%
\pgfpathquadraticcurveto{\pgfqpoint{4.886661in}{0.711161in}}{\pgfqpoint{4.881167in}{0.712152in}}%
\pgfpathquadraticcurveto{\pgfqpoint{4.875674in}{0.713142in}}{\pgfqpoint{4.869604in}{0.715106in}}%
\pgfpathlineto{\pgfqpoint{4.869604in}{0.725805in}}%
\pgfpathquadraticcurveto{\pgfqpoint{4.875350in}{0.722815in}}{\pgfqpoint{4.880915in}{0.721320in}}%
\pgfpathquadraticcurveto{\pgfqpoint{4.886481in}{0.719843in}}{\pgfqpoint{4.891957in}{0.719843in}}%
\pgfpathquadraticcurveto{\pgfqpoint{4.899270in}{0.719843in}}{\pgfqpoint{4.903196in}{0.722346in}}%
\pgfpathquadraticcurveto{\pgfqpoint{4.907141in}{0.724850in}}{\pgfqpoint{4.907141in}{0.729407in}}%
\pgfpathquadraticcurveto{\pgfqpoint{4.907141in}{0.733622in}}{\pgfqpoint{4.904295in}{0.735874in}}%
\pgfpathquadraticcurveto{\pgfqpoint{4.901467in}{0.738125in}}{\pgfqpoint{4.891831in}{0.740214in}}%
\pgfpathlineto{\pgfqpoint{4.888228in}{0.741061in}}%
\pgfpathquadraticcurveto{\pgfqpoint{4.878610in}{0.743078in}}{\pgfqpoint{4.874323in}{0.747275in}}%
\pgfpathquadraticcurveto{\pgfqpoint{4.870054in}{0.751472in}}{\pgfqpoint{4.870054in}{0.758785in}}%
\pgfpathquadraticcurveto{\pgfqpoint{4.870054in}{0.767683in}}{\pgfqpoint{4.876358in}{0.772510in}}%
\pgfpathquadraticcurveto{\pgfqpoint{4.882662in}{0.777356in}}{\pgfqpoint{4.894262in}{0.777356in}}%
\pgfpathquadraticcurveto{\pgfqpoint{4.899990in}{0.777356in}}{\pgfqpoint{4.905052in}{0.776509in}}%
\pgfpathquadraticcurveto{\pgfqpoint{4.910131in}{0.775680in}}{\pgfqpoint{4.914400in}{0.773987in}}%
\pgfpathlineto{\pgfqpoint{4.914400in}{0.773987in}}%
\pgfpathclose%
\pgfpathmoveto{\pgfqpoint{4.936915in}{0.772402in}}%
\pgfpathlineto{\pgfqpoint{4.948785in}{0.772402in}}%
\pgfpathlineto{\pgfqpoint{4.948785in}{0.758119in}}%
\pgfpathlineto{\pgfqpoint{4.936915in}{0.758119in}}%
\pgfpathlineto{\pgfqpoint{4.936915in}{0.772402in}}%
\pgfpathlineto{\pgfqpoint{4.936915in}{0.772402in}}%
\pgfpathclose%
\pgfpathmoveto{\pgfqpoint{4.936915in}{0.727102in}}%
\pgfpathlineto{\pgfqpoint{4.948785in}{0.727102in}}%
\pgfpathlineto{\pgfqpoint{4.948785in}{0.717411in}}%
\pgfpathlineto{\pgfqpoint{4.939563in}{0.699399in}}%
\pgfpathlineto{\pgfqpoint{4.932304in}{0.699399in}}%
\pgfpathlineto{\pgfqpoint{4.936915in}{0.717411in}}%
\pgfpathlineto{\pgfqpoint{4.936915in}{0.727102in}}%
\pgfpathlineto{\pgfqpoint{4.936915in}{0.727102in}}%
\pgfpathclose%
\pgfusepath{fill}%
\end{pgfscope}%
\begin{pgfscope}%
\pgfsetbuttcap%
\pgfsetroundjoin%
\definecolor{currentfill}{rgb}{0.309804,0.372549,0.419608}%
\pgfsetfillcolor{currentfill}%
\pgfsetlinewidth{0.000000pt}%
\definecolor{currentstroke}{rgb}{0.309804,0.372549,0.419608}%
\pgfsetstrokecolor{currentstroke}%
\pgfsetdash{}{0pt}%
\pgfpathmoveto{\pgfqpoint{5.033873in}{0.775843in}}%
\pgfpathlineto{\pgfqpoint{5.044230in}{0.775843in}}%
\pgfpathlineto{\pgfqpoint{5.044230in}{0.712800in}}%
\pgfpathlineto{\pgfqpoint{5.033873in}{0.712800in}}%
\pgfpathlineto{\pgfqpoint{5.033873in}{0.775843in}}%
\pgfpathlineto{\pgfqpoint{5.033873in}{0.775843in}}%
\pgfpathclose%
\pgfpathmoveto{\pgfqpoint{5.033873in}{0.800393in}}%
\pgfpathlineto{\pgfqpoint{5.044230in}{0.800393in}}%
\pgfpathlineto{\pgfqpoint{5.044230in}{0.787262in}}%
\pgfpathlineto{\pgfqpoint{5.033873in}{0.787262in}}%
\pgfpathlineto{\pgfqpoint{5.033873in}{0.800393in}}%
\pgfpathlineto{\pgfqpoint{5.033873in}{0.800393in}}%
\pgfpathclose%
\pgfpathmoveto{\pgfqpoint{5.076148in}{0.793747in}}%
\pgfpathlineto{\pgfqpoint{5.076148in}{0.775843in}}%
\pgfpathlineto{\pgfqpoint{5.097474in}{0.775843in}}%
\pgfpathlineto{\pgfqpoint{5.097474in}{0.767791in}}%
\pgfpathlineto{\pgfqpoint{5.076148in}{0.767791in}}%
\pgfpathlineto{\pgfqpoint{5.076148in}{0.733568in}}%
\pgfpathquadraticcurveto{\pgfqpoint{5.076148in}{0.725859in}}{\pgfqpoint{5.078255in}{0.723661in}}%
\pgfpathquadraticcurveto{\pgfqpoint{5.080363in}{0.721464in}}{\pgfqpoint{5.086847in}{0.721464in}}%
\pgfpathlineto{\pgfqpoint{5.097474in}{0.721464in}}%
\pgfpathlineto{\pgfqpoint{5.097474in}{0.712800in}}%
\pgfpathlineto{\pgfqpoint{5.086847in}{0.712800in}}%
\pgfpathquadraticcurveto{\pgfqpoint{5.074851in}{0.712800in}}{\pgfqpoint{5.070294in}{0.717267in}}%
\pgfpathquadraticcurveto{\pgfqpoint{5.065737in}{0.721752in}}{\pgfqpoint{5.065737in}{0.733568in}}%
\pgfpathlineto{\pgfqpoint{5.065737in}{0.767791in}}%
\pgfpathlineto{\pgfqpoint{5.058136in}{0.767791in}}%
\pgfpathlineto{\pgfqpoint{5.058136in}{0.775843in}}%
\pgfpathlineto{\pgfqpoint{5.065737in}{0.775843in}}%
\pgfpathlineto{\pgfqpoint{5.065737in}{0.793747in}}%
\pgfpathlineto{\pgfqpoint{5.076148in}{0.793747in}}%
\pgfpathlineto{\pgfqpoint{5.076148in}{0.793747in}}%
\pgfpathclose%
\pgfusepath{fill}%
\end{pgfscope}%
\begin{pgfscope}%
\pgfsetbuttcap%
\pgfsetroundjoin%
\definecolor{currentfill}{rgb}{0.309804,0.372549,0.419608}%
\pgfsetfillcolor{currentfill}%
\pgfsetlinewidth{0.000000pt}%
\definecolor{currentstroke}{rgb}{0.309804,0.372549,0.419608}%
\pgfsetstrokecolor{currentstroke}%
\pgfsetdash{}{0pt}%
\pgfpathmoveto{\pgfqpoint{5.155566in}{0.775843in}}%
\pgfpathlineto{\pgfqpoint{5.165923in}{0.775843in}}%
\pgfpathlineto{\pgfqpoint{5.165923in}{0.712800in}}%
\pgfpathlineto{\pgfqpoint{5.155566in}{0.712800in}}%
\pgfpathlineto{\pgfqpoint{5.155566in}{0.775843in}}%
\pgfpathlineto{\pgfqpoint{5.155566in}{0.775843in}}%
\pgfpathclose%
\pgfpathmoveto{\pgfqpoint{5.155566in}{0.800393in}}%
\pgfpathlineto{\pgfqpoint{5.165923in}{0.800393in}}%
\pgfpathlineto{\pgfqpoint{5.165923in}{0.787262in}}%
\pgfpathlineto{\pgfqpoint{5.155566in}{0.787262in}}%
\pgfpathlineto{\pgfqpoint{5.155566in}{0.800393in}}%
\pgfpathlineto{\pgfqpoint{5.155566in}{0.800393in}}%
\pgfpathclose%
\pgfpathmoveto{\pgfqpoint{5.227777in}{0.773987in}}%
\pgfpathlineto{\pgfqpoint{5.227777in}{0.764189in}}%
\pgfpathquadraticcurveto{\pgfqpoint{5.223400in}{0.766440in}}{\pgfqpoint{5.218663in}{0.767557in}}%
\pgfpathquadraticcurveto{\pgfqpoint{5.213943in}{0.768692in}}{\pgfqpoint{5.208864in}{0.768692in}}%
\pgfpathquadraticcurveto{\pgfqpoint{5.201155in}{0.768692in}}{\pgfqpoint{5.197300in}{0.766332in}}%
\pgfpathquadraticcurveto{\pgfqpoint{5.193446in}{0.763973in}}{\pgfqpoint{5.193446in}{0.759235in}}%
\pgfpathquadraticcurveto{\pgfqpoint{5.193446in}{0.755633in}}{\pgfqpoint{5.196201in}{0.753580in}}%
\pgfpathquadraticcurveto{\pgfqpoint{5.198957in}{0.751526in}}{\pgfqpoint{5.207297in}{0.749671in}}%
\pgfpathlineto{\pgfqpoint{5.210845in}{0.748878in}}%
\pgfpathquadraticcurveto{\pgfqpoint{5.221869in}{0.746519in}}{\pgfqpoint{5.226516in}{0.742214in}}%
\pgfpathquadraticcurveto{\pgfqpoint{5.231163in}{0.737909in}}{\pgfqpoint{5.231163in}{0.730200in}}%
\pgfpathquadraticcurveto{\pgfqpoint{5.231163in}{0.721410in}}{\pgfqpoint{5.224210in}{0.716276in}}%
\pgfpathquadraticcurveto{\pgfqpoint{5.217258in}{0.711161in}}{\pgfqpoint{5.205099in}{0.711161in}}%
\pgfpathquadraticcurveto{\pgfqpoint{5.200038in}{0.711161in}}{\pgfqpoint{5.194544in}{0.712152in}}%
\pgfpathquadraticcurveto{\pgfqpoint{5.189051in}{0.713142in}}{\pgfqpoint{5.182981in}{0.715106in}}%
\pgfpathlineto{\pgfqpoint{5.182981in}{0.725805in}}%
\pgfpathquadraticcurveto{\pgfqpoint{5.188726in}{0.722815in}}{\pgfqpoint{5.194292in}{0.721320in}}%
\pgfpathquadraticcurveto{\pgfqpoint{5.199858in}{0.719843in}}{\pgfqpoint{5.205334in}{0.719843in}}%
\pgfpathquadraticcurveto{\pgfqpoint{5.212647in}{0.719843in}}{\pgfqpoint{5.216573in}{0.722346in}}%
\pgfpathquadraticcurveto{\pgfqpoint{5.220518in}{0.724850in}}{\pgfqpoint{5.220518in}{0.729407in}}%
\pgfpathquadraticcurveto{\pgfqpoint{5.220518in}{0.733622in}}{\pgfqpoint{5.217672in}{0.735874in}}%
\pgfpathquadraticcurveto{\pgfqpoint{5.214844in}{0.738125in}}{\pgfqpoint{5.205208in}{0.740214in}}%
\pgfpathlineto{\pgfqpoint{5.201605in}{0.741061in}}%
\pgfpathquadraticcurveto{\pgfqpoint{5.191987in}{0.743078in}}{\pgfqpoint{5.187700in}{0.747275in}}%
\pgfpathquadraticcurveto{\pgfqpoint{5.183431in}{0.751472in}}{\pgfqpoint{5.183431in}{0.758785in}}%
\pgfpathquadraticcurveto{\pgfqpoint{5.183431in}{0.767683in}}{\pgfqpoint{5.189735in}{0.772510in}}%
\pgfpathquadraticcurveto{\pgfqpoint{5.196039in}{0.777356in}}{\pgfqpoint{5.207639in}{0.777356in}}%
\pgfpathquadraticcurveto{\pgfqpoint{5.213367in}{0.777356in}}{\pgfqpoint{5.218428in}{0.776509in}}%
\pgfpathquadraticcurveto{\pgfqpoint{5.223508in}{0.775680in}}{\pgfqpoint{5.227777in}{0.773987in}}%
\pgfpathlineto{\pgfqpoint{5.227777in}{0.773987in}}%
\pgfpathclose%
\pgfusepath{fill}%
\end{pgfscope}%
\begin{pgfscope}%
\pgfsetbuttcap%
\pgfsetroundjoin%
\definecolor{currentfill}{rgb}{0.309804,0.372549,0.419608}%
\pgfsetfillcolor{currentfill}%
\pgfsetlinewidth{0.000000pt}%
\definecolor{currentstroke}{rgb}{0.309804,0.372549,0.419608}%
\pgfsetstrokecolor{currentstroke}%
\pgfsetdash{}{0pt}%
\pgfpathmoveto{\pgfqpoint{5.315916in}{0.744483in}}%
\pgfpathquadraticcurveto{\pgfqpoint{5.303362in}{0.744483in}}{\pgfqpoint{5.298517in}{0.741619in}}%
\pgfpathquadraticcurveto{\pgfqpoint{5.293671in}{0.738756in}}{\pgfqpoint{5.293671in}{0.731821in}}%
\pgfpathquadraticcurveto{\pgfqpoint{5.293671in}{0.726309in}}{\pgfqpoint{5.297310in}{0.723067in}}%
\pgfpathquadraticcurveto{\pgfqpoint{5.300948in}{0.719843in}}{\pgfqpoint{5.307180in}{0.719843in}}%
\pgfpathquadraticcurveto{\pgfqpoint{5.315808in}{0.719843in}}{\pgfqpoint{5.321014in}{0.725949in}}%
\pgfpathquadraticcurveto{\pgfqpoint{5.326219in}{0.732055in}}{\pgfqpoint{5.326219in}{0.742178in}}%
\pgfpathlineto{\pgfqpoint{5.326219in}{0.744483in}}%
\pgfpathlineto{\pgfqpoint{5.315916in}{0.744483in}}%
\pgfpathlineto{\pgfqpoint{5.315916in}{0.744483in}}%
\pgfpathclose%
\pgfpathmoveto{\pgfqpoint{5.336576in}{0.748770in}}%
\pgfpathlineto{\pgfqpoint{5.336576in}{0.712800in}}%
\pgfpathlineto{\pgfqpoint{5.326219in}{0.712800in}}%
\pgfpathlineto{\pgfqpoint{5.326219in}{0.722364in}}%
\pgfpathquadraticcurveto{\pgfqpoint{5.322671in}{0.716637in}}{\pgfqpoint{5.317375in}{0.713899in}}%
\pgfpathquadraticcurveto{\pgfqpoint{5.312080in}{0.711161in}}{\pgfqpoint{5.304425in}{0.711161in}}%
\pgfpathquadraticcurveto{\pgfqpoint{5.294752in}{0.711161in}}{\pgfqpoint{5.289024in}{0.716601in}}%
\pgfpathquadraticcurveto{\pgfqpoint{5.283314in}{0.722040in}}{\pgfqpoint{5.283314in}{0.731154in}}%
\pgfpathquadraticcurveto{\pgfqpoint{5.283314in}{0.741782in}}{\pgfqpoint{5.290429in}{0.747185in}}%
\pgfpathquadraticcurveto{\pgfqpoint{5.297562in}{0.752589in}}{\pgfqpoint{5.311683in}{0.752589in}}%
\pgfpathlineto{\pgfqpoint{5.326219in}{0.752589in}}%
\pgfpathlineto{\pgfqpoint{5.326219in}{0.753616in}}%
\pgfpathquadraticcurveto{\pgfqpoint{5.326219in}{0.760766in}}{\pgfqpoint{5.321518in}{0.764675in}}%
\pgfpathquadraticcurveto{\pgfqpoint{5.316817in}{0.768584in}}{\pgfqpoint{5.308315in}{0.768584in}}%
\pgfpathquadraticcurveto{\pgfqpoint{5.302912in}{0.768584in}}{\pgfqpoint{5.297778in}{0.767287in}}%
\pgfpathquadraticcurveto{\pgfqpoint{5.292663in}{0.765990in}}{\pgfqpoint{5.287943in}{0.763396in}}%
\pgfpathlineto{\pgfqpoint{5.287943in}{0.772979in}}%
\pgfpathquadraticcurveto{\pgfqpoint{5.293617in}{0.775176in}}{\pgfqpoint{5.298967in}{0.776257in}}%
\pgfpathquadraticcurveto{\pgfqpoint{5.304316in}{0.777356in}}{\pgfqpoint{5.309378in}{0.777356in}}%
\pgfpathquadraticcurveto{\pgfqpoint{5.323067in}{0.777356in}}{\pgfqpoint{5.329822in}{0.770259in}}%
\pgfpathquadraticcurveto{\pgfqpoint{5.336576in}{0.763180in}}{\pgfqpoint{5.336576in}{0.748770in}}%
\pgfpathlineto{\pgfqpoint{5.336576in}{0.748770in}}%
\pgfpathclose%
\pgfusepath{fill}%
\end{pgfscope}%
\begin{pgfscope}%
\pgfsetbuttcap%
\pgfsetroundjoin%
\definecolor{currentfill}{rgb}{0.309804,0.372549,0.419608}%
\pgfsetfillcolor{currentfill}%
\pgfsetlinewidth{0.000000pt}%
\definecolor{currentstroke}{rgb}{0.309804,0.372549,0.419608}%
\pgfsetstrokecolor{currentstroke}%
\pgfsetdash{}{0pt}%
\pgfpathmoveto{\pgfqpoint{5.440434in}{0.766278in}}%
\pgfpathlineto{\pgfqpoint{5.440434in}{0.800393in}}%
\pgfpathlineto{\pgfqpoint{5.450791in}{0.800393in}}%
\pgfpathlineto{\pgfqpoint{5.450791in}{0.712800in}}%
\pgfpathlineto{\pgfqpoint{5.440434in}{0.712800in}}%
\pgfpathlineto{\pgfqpoint{5.440434in}{0.722256in}}%
\pgfpathquadraticcurveto{\pgfqpoint{5.437174in}{0.716637in}}{\pgfqpoint{5.432184in}{0.713899in}}%
\pgfpathquadraticcurveto{\pgfqpoint{5.427213in}{0.711161in}}{\pgfqpoint{5.420224in}{0.711161in}}%
\pgfpathquadraticcurveto{\pgfqpoint{5.408804in}{0.711161in}}{\pgfqpoint{5.401618in}{0.720275in}}%
\pgfpathquadraticcurveto{\pgfqpoint{5.394449in}{0.729407in}}{\pgfqpoint{5.394449in}{0.744267in}}%
\pgfpathquadraticcurveto{\pgfqpoint{5.394449in}{0.759127in}}{\pgfqpoint{5.401618in}{0.768241in}}%
\pgfpathquadraticcurveto{\pgfqpoint{5.408804in}{0.777356in}}{\pgfqpoint{5.420224in}{0.777356in}}%
\pgfpathquadraticcurveto{\pgfqpoint{5.427213in}{0.777356in}}{\pgfqpoint{5.432184in}{0.774618in}}%
\pgfpathquadraticcurveto{\pgfqpoint{5.437174in}{0.771898in}}{\pgfqpoint{5.440434in}{0.766278in}}%
\pgfpathlineto{\pgfqpoint{5.440434in}{0.766278in}}%
\pgfpathclose%
\pgfpathmoveto{\pgfqpoint{5.405148in}{0.744267in}}%
\pgfpathquadraticcurveto{\pgfqpoint{5.405148in}{0.732848in}}{\pgfqpoint{5.409849in}{0.726345in}}%
\pgfpathquadraticcurveto{\pgfqpoint{5.414550in}{0.719843in}}{\pgfqpoint{5.422764in}{0.719843in}}%
\pgfpathquadraticcurveto{\pgfqpoint{5.430977in}{0.719843in}}{\pgfqpoint{5.435697in}{0.726345in}}%
\pgfpathquadraticcurveto{\pgfqpoint{5.440434in}{0.732848in}}{\pgfqpoint{5.440434in}{0.744267in}}%
\pgfpathquadraticcurveto{\pgfqpoint{5.440434in}{0.755687in}}{\pgfqpoint{5.435697in}{0.762189in}}%
\pgfpathquadraticcurveto{\pgfqpoint{5.430977in}{0.768692in}}{\pgfqpoint{5.422764in}{0.768692in}}%
\pgfpathquadraticcurveto{\pgfqpoint{5.414550in}{0.768692in}}{\pgfqpoint{5.409849in}{0.762189in}}%
\pgfpathquadraticcurveto{\pgfqpoint{5.405148in}{0.755687in}}{\pgfqpoint{5.405148in}{0.744267in}}%
\pgfpathlineto{\pgfqpoint{5.405148in}{0.744267in}}%
\pgfpathclose%
\pgfpathmoveto{\pgfqpoint{5.526064in}{0.746915in}}%
\pgfpathlineto{\pgfqpoint{5.526064in}{0.741854in}}%
\pgfpathlineto{\pgfqpoint{5.478439in}{0.741854in}}%
\pgfpathquadraticcurveto{\pgfqpoint{5.479124in}{0.731154in}}{\pgfqpoint{5.484888in}{0.725553in}}%
\pgfpathquadraticcurveto{\pgfqpoint{5.490652in}{0.719951in}}{\pgfqpoint{5.500955in}{0.719951in}}%
\pgfpathquadraticcurveto{\pgfqpoint{5.506917in}{0.719951in}}{\pgfqpoint{5.512518in}{0.721410in}}%
\pgfpathquadraticcurveto{\pgfqpoint{5.518120in}{0.722869in}}{\pgfqpoint{5.523650in}{0.725805in}}%
\pgfpathlineto{\pgfqpoint{5.523650in}{0.716006in}}%
\pgfpathquadraticcurveto{\pgfqpoint{5.518066in}{0.713647in}}{\pgfqpoint{5.512212in}{0.712404in}}%
\pgfpathquadraticcurveto{\pgfqpoint{5.506358in}{0.711161in}}{\pgfqpoint{5.500342in}{0.711161in}}%
\pgfpathquadraticcurveto{\pgfqpoint{5.485248in}{0.711161in}}{\pgfqpoint{5.476440in}{0.719933in}}%
\pgfpathquadraticcurveto{\pgfqpoint{5.467632in}{0.728723in}}{\pgfqpoint{5.467632in}{0.743709in}}%
\pgfpathquadraticcurveto{\pgfqpoint{5.467632in}{0.759181in}}{\pgfqpoint{5.475990in}{0.768259in}}%
\pgfpathquadraticcurveto{\pgfqpoint{5.484347in}{0.777356in}}{\pgfqpoint{5.498541in}{0.777356in}}%
\pgfpathquadraticcurveto{\pgfqpoint{5.511258in}{0.777356in}}{\pgfqpoint{5.518661in}{0.769160in}}%
\pgfpathquadraticcurveto{\pgfqpoint{5.526064in}{0.760983in}}{\pgfqpoint{5.526064in}{0.746915in}}%
\pgfpathlineto{\pgfqpoint{5.526064in}{0.746915in}}%
\pgfpathclose%
\pgfpathmoveto{\pgfqpoint{5.515707in}{0.749959in}}%
\pgfpathquadraticcurveto{\pgfqpoint{5.515599in}{0.758443in}}{\pgfqpoint{5.510951in}{0.763504in}}%
\pgfpathquadraticcurveto{\pgfqpoint{5.506304in}{0.768584in}}{\pgfqpoint{5.498649in}{0.768584in}}%
\pgfpathquadraticcurveto{\pgfqpoint{5.489985in}{0.768584in}}{\pgfqpoint{5.484780in}{0.763684in}}%
\pgfpathquadraticcurveto{\pgfqpoint{5.479574in}{0.758785in}}{\pgfqpoint{5.478782in}{0.749887in}}%
\pgfpathlineto{\pgfqpoint{5.515707in}{0.749959in}}%
\pgfpathlineto{\pgfqpoint{5.515707in}{0.749959in}}%
\pgfpathclose%
\pgfpathmoveto{\pgfqpoint{5.583252in}{0.773987in}}%
\pgfpathlineto{\pgfqpoint{5.583252in}{0.764189in}}%
\pgfpathquadraticcurveto{\pgfqpoint{5.578875in}{0.766440in}}{\pgfqpoint{5.574138in}{0.767557in}}%
\pgfpathquadraticcurveto{\pgfqpoint{5.569419in}{0.768692in}}{\pgfqpoint{5.564339in}{0.768692in}}%
\pgfpathquadraticcurveto{\pgfqpoint{5.556630in}{0.768692in}}{\pgfqpoint{5.552776in}{0.766332in}}%
\pgfpathquadraticcurveto{\pgfqpoint{5.548921in}{0.763973in}}{\pgfqpoint{5.548921in}{0.759235in}}%
\pgfpathquadraticcurveto{\pgfqpoint{5.548921in}{0.755633in}}{\pgfqpoint{5.551677in}{0.753580in}}%
\pgfpathquadraticcurveto{\pgfqpoint{5.554433in}{0.751526in}}{\pgfqpoint{5.562772in}{0.749671in}}%
\pgfpathlineto{\pgfqpoint{5.566321in}{0.748878in}}%
\pgfpathquadraticcurveto{\pgfqpoint{5.577344in}{0.746519in}}{\pgfqpoint{5.581991in}{0.742214in}}%
\pgfpathquadraticcurveto{\pgfqpoint{5.586638in}{0.737909in}}{\pgfqpoint{5.586638in}{0.730200in}}%
\pgfpathquadraticcurveto{\pgfqpoint{5.586638in}{0.721410in}}{\pgfqpoint{5.579686in}{0.716276in}}%
\pgfpathquadraticcurveto{\pgfqpoint{5.572733in}{0.711161in}}{\pgfqpoint{5.560575in}{0.711161in}}%
\pgfpathquadraticcurveto{\pgfqpoint{5.555513in}{0.711161in}}{\pgfqpoint{5.550020in}{0.712152in}}%
\pgfpathquadraticcurveto{\pgfqpoint{5.544526in}{0.713142in}}{\pgfqpoint{5.538456in}{0.715106in}}%
\pgfpathlineto{\pgfqpoint{5.538456in}{0.725805in}}%
\pgfpathquadraticcurveto{\pgfqpoint{5.544202in}{0.722815in}}{\pgfqpoint{5.549768in}{0.721320in}}%
\pgfpathquadraticcurveto{\pgfqpoint{5.555333in}{0.719843in}}{\pgfqpoint{5.560809in}{0.719843in}}%
\pgfpathquadraticcurveto{\pgfqpoint{5.568122in}{0.719843in}}{\pgfqpoint{5.572049in}{0.722346in}}%
\pgfpathquadraticcurveto{\pgfqpoint{5.575993in}{0.724850in}}{\pgfqpoint{5.575993in}{0.729407in}}%
\pgfpathquadraticcurveto{\pgfqpoint{5.575993in}{0.733622in}}{\pgfqpoint{5.573147in}{0.735874in}}%
\pgfpathquadraticcurveto{\pgfqpoint{5.570319in}{0.738125in}}{\pgfqpoint{5.560683in}{0.740214in}}%
\pgfpathlineto{\pgfqpoint{5.557081in}{0.741061in}}%
\pgfpathquadraticcurveto{\pgfqpoint{5.547462in}{0.743078in}}{\pgfqpoint{5.543175in}{0.747275in}}%
\pgfpathquadraticcurveto{\pgfqpoint{5.538906in}{0.751472in}}{\pgfqpoint{5.538906in}{0.758785in}}%
\pgfpathquadraticcurveto{\pgfqpoint{5.538906in}{0.767683in}}{\pgfqpoint{5.545211in}{0.772510in}}%
\pgfpathquadraticcurveto{\pgfqpoint{5.551515in}{0.777356in}}{\pgfqpoint{5.563115in}{0.777356in}}%
\pgfpathquadraticcurveto{\pgfqpoint{5.568842in}{0.777356in}}{\pgfqpoint{5.573904in}{0.776509in}}%
\pgfpathquadraticcurveto{\pgfqpoint{5.578983in}{0.775680in}}{\pgfqpoint{5.583252in}{0.773987in}}%
\pgfpathlineto{\pgfqpoint{5.583252in}{0.773987in}}%
\pgfpathclose%
\pgfpathmoveto{\pgfqpoint{5.603120in}{0.775843in}}%
\pgfpathlineto{\pgfqpoint{5.613477in}{0.775843in}}%
\pgfpathlineto{\pgfqpoint{5.613477in}{0.712800in}}%
\pgfpathlineto{\pgfqpoint{5.603120in}{0.712800in}}%
\pgfpathlineto{\pgfqpoint{5.603120in}{0.775843in}}%
\pgfpathlineto{\pgfqpoint{5.603120in}{0.775843in}}%
\pgfpathclose%
\pgfpathmoveto{\pgfqpoint{5.603120in}{0.800393in}}%
\pgfpathlineto{\pgfqpoint{5.613477in}{0.800393in}}%
\pgfpathlineto{\pgfqpoint{5.613477in}{0.787262in}}%
\pgfpathlineto{\pgfqpoint{5.603120in}{0.787262in}}%
\pgfpathlineto{\pgfqpoint{5.603120in}{0.800393in}}%
\pgfpathlineto{\pgfqpoint{5.603120in}{0.800393in}}%
\pgfpathclose%
\pgfpathmoveto{\pgfqpoint{5.676627in}{0.745060in}}%
\pgfpathquadraticcurveto{\pgfqpoint{5.676627in}{0.756317in}}{\pgfqpoint{5.671980in}{0.762496in}}%
\pgfpathquadraticcurveto{\pgfqpoint{5.667351in}{0.768692in}}{\pgfqpoint{5.658957in}{0.768692in}}%
\pgfpathquadraticcurveto{\pgfqpoint{5.650636in}{0.768692in}}{\pgfqpoint{5.645989in}{0.762496in}}%
\pgfpathquadraticcurveto{\pgfqpoint{5.641341in}{0.756317in}}{\pgfqpoint{5.641341in}{0.745060in}}%
\pgfpathquadraticcurveto{\pgfqpoint{5.641341in}{0.733856in}}{\pgfqpoint{5.645989in}{0.727660in}}%
\pgfpathquadraticcurveto{\pgfqpoint{5.650636in}{0.721464in}}{\pgfqpoint{5.658957in}{0.721464in}}%
\pgfpathquadraticcurveto{\pgfqpoint{5.667351in}{0.721464in}}{\pgfqpoint{5.671980in}{0.727660in}}%
\pgfpathquadraticcurveto{\pgfqpoint{5.676627in}{0.733856in}}{\pgfqpoint{5.676627in}{0.745060in}}%
\pgfpathlineto{\pgfqpoint{5.676627in}{0.745060in}}%
\pgfpathclose%
\pgfpathmoveto{\pgfqpoint{5.686984in}{0.720617in}}%
\pgfpathquadraticcurveto{\pgfqpoint{5.686984in}{0.704532in}}{\pgfqpoint{5.679833in}{0.696679in}}%
\pgfpathquadraticcurveto{\pgfqpoint{5.672701in}{0.688826in}}{\pgfqpoint{5.657949in}{0.688826in}}%
\pgfpathquadraticcurveto{\pgfqpoint{5.652491in}{0.688826in}}{\pgfqpoint{5.647646in}{0.689636in}}%
\pgfpathquadraticcurveto{\pgfqpoint{5.642800in}{0.690447in}}{\pgfqpoint{5.638243in}{0.692140in}}%
\pgfpathlineto{\pgfqpoint{5.638243in}{0.702209in}}%
\pgfpathquadraticcurveto{\pgfqpoint{5.642800in}{0.699741in}}{\pgfqpoint{5.647249in}{0.698570in}}%
\pgfpathquadraticcurveto{\pgfqpoint{5.651698in}{0.697382in}}{\pgfqpoint{5.656309in}{0.697382in}}%
\pgfpathquadraticcurveto{\pgfqpoint{5.666504in}{0.697382in}}{\pgfqpoint{5.671566in}{0.702695in}}%
\pgfpathquadraticcurveto{\pgfqpoint{5.676627in}{0.708009in}}{\pgfqpoint{5.676627in}{0.718762in}}%
\pgfpathlineto{\pgfqpoint{5.676627in}{0.723895in}}%
\pgfpathquadraticcurveto{\pgfqpoint{5.673421in}{0.718312in}}{\pgfqpoint{5.668414in}{0.715556in}}%
\pgfpathquadraticcurveto{\pgfqpoint{5.663406in}{0.712800in}}{\pgfqpoint{5.656418in}{0.712800in}}%
\pgfpathquadraticcurveto{\pgfqpoint{5.644836in}{0.712800in}}{\pgfqpoint{5.637739in}{0.721626in}}%
\pgfpathquadraticcurveto{\pgfqpoint{5.630642in}{0.730470in}}{\pgfqpoint{5.630642in}{0.745060in}}%
\pgfpathquadraticcurveto{\pgfqpoint{5.630642in}{0.759686in}}{\pgfqpoint{5.637739in}{0.768512in}}%
\pgfpathquadraticcurveto{\pgfqpoint{5.644836in}{0.777356in}}{\pgfqpoint{5.656418in}{0.777356in}}%
\pgfpathquadraticcurveto{\pgfqpoint{5.663406in}{0.777356in}}{\pgfqpoint{5.668414in}{0.774600in}}%
\pgfpathquadraticcurveto{\pgfqpoint{5.673421in}{0.771844in}}{\pgfqpoint{5.676627in}{0.766278in}}%
\pgfpathlineto{\pgfqpoint{5.676627in}{0.775843in}}%
\pgfpathlineto{\pgfqpoint{5.686984in}{0.775843in}}%
\pgfpathlineto{\pgfqpoint{5.686984in}{0.720617in}}%
\pgfpathlineto{\pgfqpoint{5.686984in}{0.720617in}}%
\pgfpathclose%
\pgfpathmoveto{\pgfqpoint{5.760744in}{0.750860in}}%
\pgfpathlineto{\pgfqpoint{5.760744in}{0.712800in}}%
\pgfpathlineto{\pgfqpoint{5.750387in}{0.712800in}}%
\pgfpathlineto{\pgfqpoint{5.750387in}{0.750517in}}%
\pgfpathquadraticcurveto{\pgfqpoint{5.750387in}{0.759469in}}{\pgfqpoint{5.746893in}{0.763900in}}%
\pgfpathquadraticcurveto{\pgfqpoint{5.743398in}{0.768349in}}{\pgfqpoint{5.736428in}{0.768349in}}%
\pgfpathquadraticcurveto{\pgfqpoint{5.728034in}{0.768349in}}{\pgfqpoint{5.723189in}{0.763000in}}%
\pgfpathquadraticcurveto{\pgfqpoint{5.718343in}{0.757668in}}{\pgfqpoint{5.718343in}{0.748428in}}%
\pgfpathlineto{\pgfqpoint{5.718343in}{0.712800in}}%
\pgfpathlineto{\pgfqpoint{5.707932in}{0.712800in}}%
\pgfpathlineto{\pgfqpoint{5.707932in}{0.775843in}}%
\pgfpathlineto{\pgfqpoint{5.718343in}{0.775843in}}%
\pgfpathlineto{\pgfqpoint{5.718343in}{0.766044in}}%
\pgfpathquadraticcurveto{\pgfqpoint{5.722072in}{0.771736in}}{\pgfqpoint{5.727097in}{0.774546in}}%
\pgfpathquadraticcurveto{\pgfqpoint{5.732141in}{0.777356in}}{\pgfqpoint{5.738733in}{0.777356in}}%
\pgfpathquadraticcurveto{\pgfqpoint{5.749594in}{0.777356in}}{\pgfqpoint{5.755160in}{0.770637in}}%
\pgfpathquadraticcurveto{\pgfqpoint{5.760744in}{0.763918in}}{\pgfqpoint{5.760744in}{0.750860in}}%
\pgfpathlineto{\pgfqpoint{5.760744in}{0.750860in}}%
\pgfpathclose%
\pgfusepath{fill}%
\end{pgfscope}%
\begin{pgfscope}%
\pgfsetbuttcap%
\pgfsetroundjoin%
\definecolor{currentfill}{rgb}{0.309804,0.372549,0.419608}%
\pgfsetfillcolor{currentfill}%
\pgfsetlinewidth{0.000000pt}%
\definecolor{currentstroke}{rgb}{0.309804,0.372549,0.419608}%
\pgfsetstrokecolor{currentstroke}%
\pgfsetdash{}{0pt}%
\pgfpathmoveto{\pgfqpoint{5.882127in}{0.766278in}}%
\pgfpathlineto{\pgfqpoint{5.882127in}{0.800393in}}%
\pgfpathlineto{\pgfqpoint{5.892484in}{0.800393in}}%
\pgfpathlineto{\pgfqpoint{5.892484in}{0.712800in}}%
\pgfpathlineto{\pgfqpoint{5.882127in}{0.712800in}}%
\pgfpathlineto{\pgfqpoint{5.882127in}{0.722256in}}%
\pgfpathquadraticcurveto{\pgfqpoint{5.878867in}{0.716637in}}{\pgfqpoint{5.873877in}{0.713899in}}%
\pgfpathquadraticcurveto{\pgfqpoint{5.868906in}{0.711161in}}{\pgfqpoint{5.861917in}{0.711161in}}%
\pgfpathquadraticcurveto{\pgfqpoint{5.850497in}{0.711161in}}{\pgfqpoint{5.843311in}{0.720275in}}%
\pgfpathquadraticcurveto{\pgfqpoint{5.836142in}{0.729407in}}{\pgfqpoint{5.836142in}{0.744267in}}%
\pgfpathquadraticcurveto{\pgfqpoint{5.836142in}{0.759127in}}{\pgfqpoint{5.843311in}{0.768241in}}%
\pgfpathquadraticcurveto{\pgfqpoint{5.850497in}{0.777356in}}{\pgfqpoint{5.861917in}{0.777356in}}%
\pgfpathquadraticcurveto{\pgfqpoint{5.868906in}{0.777356in}}{\pgfqpoint{5.873877in}{0.774618in}}%
\pgfpathquadraticcurveto{\pgfqpoint{5.878867in}{0.771898in}}{\pgfqpoint{5.882127in}{0.766278in}}%
\pgfpathlineto{\pgfqpoint{5.882127in}{0.766278in}}%
\pgfpathclose%
\pgfpathmoveto{\pgfqpoint{5.846841in}{0.744267in}}%
\pgfpathquadraticcurveto{\pgfqpoint{5.846841in}{0.732848in}}{\pgfqpoint{5.851542in}{0.726345in}}%
\pgfpathquadraticcurveto{\pgfqpoint{5.856243in}{0.719843in}}{\pgfqpoint{5.864457in}{0.719843in}}%
\pgfpathquadraticcurveto{\pgfqpoint{5.872670in}{0.719843in}}{\pgfqpoint{5.877390in}{0.726345in}}%
\pgfpathquadraticcurveto{\pgfqpoint{5.882127in}{0.732848in}}{\pgfqpoint{5.882127in}{0.744267in}}%
\pgfpathquadraticcurveto{\pgfqpoint{5.882127in}{0.755687in}}{\pgfqpoint{5.877390in}{0.762189in}}%
\pgfpathquadraticcurveto{\pgfqpoint{5.872670in}{0.768692in}}{\pgfqpoint{5.864457in}{0.768692in}}%
\pgfpathquadraticcurveto{\pgfqpoint{5.856243in}{0.768692in}}{\pgfqpoint{5.851542in}{0.762189in}}%
\pgfpathquadraticcurveto{\pgfqpoint{5.846841in}{0.755687in}}{\pgfqpoint{5.846841in}{0.744267in}}%
\pgfpathlineto{\pgfqpoint{5.846841in}{0.744267in}}%
\pgfpathclose%
\pgfpathmoveto{\pgfqpoint{5.913828in}{0.775843in}}%
\pgfpathlineto{\pgfqpoint{5.924185in}{0.775843in}}%
\pgfpathlineto{\pgfqpoint{5.924185in}{0.712800in}}%
\pgfpathlineto{\pgfqpoint{5.913828in}{0.712800in}}%
\pgfpathlineto{\pgfqpoint{5.913828in}{0.775843in}}%
\pgfpathlineto{\pgfqpoint{5.913828in}{0.775843in}}%
\pgfpathclose%
\pgfpathmoveto{\pgfqpoint{5.913828in}{0.800393in}}%
\pgfpathlineto{\pgfqpoint{5.924185in}{0.800393in}}%
\pgfpathlineto{\pgfqpoint{5.924185in}{0.787262in}}%
\pgfpathlineto{\pgfqpoint{5.913828in}{0.787262in}}%
\pgfpathlineto{\pgfqpoint{5.913828in}{0.800393in}}%
\pgfpathlineto{\pgfqpoint{5.913828in}{0.800393in}}%
\pgfpathclose%
\pgfpathmoveto{\pgfqpoint{5.974511in}{0.744483in}}%
\pgfpathquadraticcurveto{\pgfqpoint{5.961957in}{0.744483in}}{\pgfqpoint{5.957111in}{0.741619in}}%
\pgfpathquadraticcurveto{\pgfqpoint{5.952266in}{0.738756in}}{\pgfqpoint{5.952266in}{0.731821in}}%
\pgfpathquadraticcurveto{\pgfqpoint{5.952266in}{0.726309in}}{\pgfqpoint{5.955904in}{0.723067in}}%
\pgfpathquadraticcurveto{\pgfqpoint{5.959543in}{0.719843in}}{\pgfqpoint{5.965775in}{0.719843in}}%
\pgfpathquadraticcurveto{\pgfqpoint{5.974403in}{0.719843in}}{\pgfqpoint{5.979608in}{0.725949in}}%
\pgfpathquadraticcurveto{\pgfqpoint{5.984814in}{0.732055in}}{\pgfqpoint{5.984814in}{0.742178in}}%
\pgfpathlineto{\pgfqpoint{5.984814in}{0.744483in}}%
\pgfpathlineto{\pgfqpoint{5.974511in}{0.744483in}}%
\pgfpathlineto{\pgfqpoint{5.974511in}{0.744483in}}%
\pgfpathclose%
\pgfpathmoveto{\pgfqpoint{5.995171in}{0.748770in}}%
\pgfpathlineto{\pgfqpoint{5.995171in}{0.712800in}}%
\pgfpathlineto{\pgfqpoint{5.984814in}{0.712800in}}%
\pgfpathlineto{\pgfqpoint{5.984814in}{0.722364in}}%
\pgfpathquadraticcurveto{\pgfqpoint{5.981266in}{0.716637in}}{\pgfqpoint{5.975970in}{0.713899in}}%
\pgfpathquadraticcurveto{\pgfqpoint{5.970674in}{0.711161in}}{\pgfqpoint{5.963019in}{0.711161in}}%
\pgfpathquadraticcurveto{\pgfqpoint{5.953347in}{0.711161in}}{\pgfqpoint{5.947619in}{0.716601in}}%
\pgfpathquadraticcurveto{\pgfqpoint{5.941909in}{0.722040in}}{\pgfqpoint{5.941909in}{0.731154in}}%
\pgfpathquadraticcurveto{\pgfqpoint{5.941909in}{0.741782in}}{\pgfqpoint{5.949024in}{0.747185in}}%
\pgfpathquadraticcurveto{\pgfqpoint{5.956157in}{0.752589in}}{\pgfqpoint{5.970278in}{0.752589in}}%
\pgfpathlineto{\pgfqpoint{5.984814in}{0.752589in}}%
\pgfpathlineto{\pgfqpoint{5.984814in}{0.753616in}}%
\pgfpathquadraticcurveto{\pgfqpoint{5.984814in}{0.760766in}}{\pgfqpoint{5.980113in}{0.764675in}}%
\pgfpathquadraticcurveto{\pgfqpoint{5.975412in}{0.768584in}}{\pgfqpoint{5.966910in}{0.768584in}}%
\pgfpathquadraticcurveto{\pgfqpoint{5.961506in}{0.768584in}}{\pgfqpoint{5.956373in}{0.767287in}}%
\pgfpathquadraticcurveto{\pgfqpoint{5.951257in}{0.765990in}}{\pgfqpoint{5.946538in}{0.763396in}}%
\pgfpathlineto{\pgfqpoint{5.946538in}{0.772979in}}%
\pgfpathquadraticcurveto{\pgfqpoint{5.952212in}{0.775176in}}{\pgfqpoint{5.957562in}{0.776257in}}%
\pgfpathquadraticcurveto{\pgfqpoint{5.962911in}{0.777356in}}{\pgfqpoint{5.967973in}{0.777356in}}%
\pgfpathquadraticcurveto{\pgfqpoint{5.981662in}{0.777356in}}{\pgfqpoint{5.988416in}{0.770259in}}%
\pgfpathquadraticcurveto{\pgfqpoint{5.995171in}{0.763180in}}{\pgfqpoint{5.995171in}{0.748770in}}%
\pgfpathlineto{\pgfqpoint{5.995171in}{0.748770in}}%
\pgfpathclose%
\pgfpathmoveto{\pgfqpoint{6.057979in}{0.745060in}}%
\pgfpathquadraticcurveto{\pgfqpoint{6.057979in}{0.756317in}}{\pgfqpoint{6.053332in}{0.762496in}}%
\pgfpathquadraticcurveto{\pgfqpoint{6.048703in}{0.768692in}}{\pgfqpoint{6.040309in}{0.768692in}}%
\pgfpathquadraticcurveto{\pgfqpoint{6.031988in}{0.768692in}}{\pgfqpoint{6.027341in}{0.762496in}}%
\pgfpathquadraticcurveto{\pgfqpoint{6.022694in}{0.756317in}}{\pgfqpoint{6.022694in}{0.745060in}}%
\pgfpathquadraticcurveto{\pgfqpoint{6.022694in}{0.733856in}}{\pgfqpoint{6.027341in}{0.727660in}}%
\pgfpathquadraticcurveto{\pgfqpoint{6.031988in}{0.721464in}}{\pgfqpoint{6.040309in}{0.721464in}}%
\pgfpathquadraticcurveto{\pgfqpoint{6.048703in}{0.721464in}}{\pgfqpoint{6.053332in}{0.727660in}}%
\pgfpathquadraticcurveto{\pgfqpoint{6.057979in}{0.733856in}}{\pgfqpoint{6.057979in}{0.745060in}}%
\pgfpathlineto{\pgfqpoint{6.057979in}{0.745060in}}%
\pgfpathclose%
\pgfpathmoveto{\pgfqpoint{6.068336in}{0.720617in}}%
\pgfpathquadraticcurveto{\pgfqpoint{6.068336in}{0.704532in}}{\pgfqpoint{6.061186in}{0.696679in}}%
\pgfpathquadraticcurveto{\pgfqpoint{6.054053in}{0.688826in}}{\pgfqpoint{6.039301in}{0.688826in}}%
\pgfpathquadraticcurveto{\pgfqpoint{6.033843in}{0.688826in}}{\pgfqpoint{6.028998in}{0.689636in}}%
\pgfpathquadraticcurveto{\pgfqpoint{6.024153in}{0.690447in}}{\pgfqpoint{6.019595in}{0.692140in}}%
\pgfpathlineto{\pgfqpoint{6.019595in}{0.702209in}}%
\pgfpathquadraticcurveto{\pgfqpoint{6.024153in}{0.699741in}}{\pgfqpoint{6.028602in}{0.698570in}}%
\pgfpathquadraticcurveto{\pgfqpoint{6.033051in}{0.697382in}}{\pgfqpoint{6.037662in}{0.697382in}}%
\pgfpathquadraticcurveto{\pgfqpoint{6.047857in}{0.697382in}}{\pgfqpoint{6.052918in}{0.702695in}}%
\pgfpathquadraticcurveto{\pgfqpoint{6.057979in}{0.708009in}}{\pgfqpoint{6.057979in}{0.718762in}}%
\pgfpathlineto{\pgfqpoint{6.057979in}{0.723895in}}%
\pgfpathquadraticcurveto{\pgfqpoint{6.054773in}{0.718312in}}{\pgfqpoint{6.049766in}{0.715556in}}%
\pgfpathquadraticcurveto{\pgfqpoint{6.044758in}{0.712800in}}{\pgfqpoint{6.037770in}{0.712800in}}%
\pgfpathquadraticcurveto{\pgfqpoint{6.026188in}{0.712800in}}{\pgfqpoint{6.019091in}{0.721626in}}%
\pgfpathquadraticcurveto{\pgfqpoint{6.011994in}{0.730470in}}{\pgfqpoint{6.011994in}{0.745060in}}%
\pgfpathquadraticcurveto{\pgfqpoint{6.011994in}{0.759686in}}{\pgfqpoint{6.019091in}{0.768512in}}%
\pgfpathquadraticcurveto{\pgfqpoint{6.026188in}{0.777356in}}{\pgfqpoint{6.037770in}{0.777356in}}%
\pgfpathquadraticcurveto{\pgfqpoint{6.044758in}{0.777356in}}{\pgfqpoint{6.049766in}{0.774600in}}%
\pgfpathquadraticcurveto{\pgfqpoint{6.054773in}{0.771844in}}{\pgfqpoint{6.057979in}{0.766278in}}%
\pgfpathlineto{\pgfqpoint{6.057979in}{0.775843in}}%
\pgfpathlineto{\pgfqpoint{6.068336in}{0.775843in}}%
\pgfpathlineto{\pgfqpoint{6.068336in}{0.720617in}}%
\pgfpathlineto{\pgfqpoint{6.068336in}{0.720617in}}%
\pgfpathclose%
\pgfpathmoveto{\pgfqpoint{6.142096in}{0.750860in}}%
\pgfpathlineto{\pgfqpoint{6.142096in}{0.712800in}}%
\pgfpathlineto{\pgfqpoint{6.131739in}{0.712800in}}%
\pgfpathlineto{\pgfqpoint{6.131739in}{0.750517in}}%
\pgfpathquadraticcurveto{\pgfqpoint{6.131739in}{0.759469in}}{\pgfqpoint{6.128245in}{0.763900in}}%
\pgfpathquadraticcurveto{\pgfqpoint{6.124750in}{0.768349in}}{\pgfqpoint{6.117780in}{0.768349in}}%
\pgfpathquadraticcurveto{\pgfqpoint{6.109386in}{0.768349in}}{\pgfqpoint{6.104541in}{0.763000in}}%
\pgfpathquadraticcurveto{\pgfqpoint{6.099696in}{0.757668in}}{\pgfqpoint{6.099696in}{0.748428in}}%
\pgfpathlineto{\pgfqpoint{6.099696in}{0.712800in}}%
\pgfpathlineto{\pgfqpoint{6.089284in}{0.712800in}}%
\pgfpathlineto{\pgfqpoint{6.089284in}{0.775843in}}%
\pgfpathlineto{\pgfqpoint{6.099696in}{0.775843in}}%
\pgfpathlineto{\pgfqpoint{6.099696in}{0.766044in}}%
\pgfpathquadraticcurveto{\pgfqpoint{6.103424in}{0.771736in}}{\pgfqpoint{6.108449in}{0.774546in}}%
\pgfpathquadraticcurveto{\pgfqpoint{6.113493in}{0.777356in}}{\pgfqpoint{6.120085in}{0.777356in}}%
\pgfpathquadraticcurveto{\pgfqpoint{6.130947in}{0.777356in}}{\pgfqpoint{6.136512in}{0.770637in}}%
\pgfpathquadraticcurveto{\pgfqpoint{6.142096in}{0.763918in}}{\pgfqpoint{6.142096in}{0.750860in}}%
\pgfpathlineto{\pgfqpoint{6.142096in}{0.750860in}}%
\pgfpathclose%
\pgfpathmoveto{\pgfqpoint{6.187163in}{0.768584in}}%
\pgfpathquadraticcurveto{\pgfqpoint{6.178841in}{0.768584in}}{\pgfqpoint{6.173996in}{0.762081in}}%
\pgfpathquadraticcurveto{\pgfqpoint{6.169150in}{0.755579in}}{\pgfqpoint{6.169150in}{0.744267in}}%
\pgfpathquadraticcurveto{\pgfqpoint{6.169150in}{0.732956in}}{\pgfqpoint{6.173960in}{0.726453in}}%
\pgfpathquadraticcurveto{\pgfqpoint{6.178787in}{0.719951in}}{\pgfqpoint{6.187163in}{0.719951in}}%
\pgfpathquadraticcurveto{\pgfqpoint{6.195448in}{0.719951in}}{\pgfqpoint{6.200275in}{0.726471in}}%
\pgfpathquadraticcurveto{\pgfqpoint{6.205121in}{0.733010in}}{\pgfqpoint{6.205121in}{0.744267in}}%
\pgfpathquadraticcurveto{\pgfqpoint{6.205121in}{0.755471in}}{\pgfqpoint{6.200275in}{0.762027in}}%
\pgfpathquadraticcurveto{\pgfqpoint{6.195448in}{0.768584in}}{\pgfqpoint{6.187163in}{0.768584in}}%
\pgfpathlineto{\pgfqpoint{6.187163in}{0.768584in}}%
\pgfpathclose%
\pgfpathmoveto{\pgfqpoint{6.187163in}{0.777356in}}%
\pgfpathquadraticcurveto{\pgfqpoint{6.200672in}{0.777356in}}{\pgfqpoint{6.208381in}{0.768566in}}%
\pgfpathquadraticcurveto{\pgfqpoint{6.216108in}{0.759794in}}{\pgfqpoint{6.216108in}{0.744267in}}%
\pgfpathquadraticcurveto{\pgfqpoint{6.216108in}{0.728795in}}{\pgfqpoint{6.208381in}{0.719969in}}%
\pgfpathquadraticcurveto{\pgfqpoint{6.200672in}{0.711161in}}{\pgfqpoint{6.187163in}{0.711161in}}%
\pgfpathquadraticcurveto{\pgfqpoint{6.173599in}{0.711161in}}{\pgfqpoint{6.165908in}{0.719969in}}%
\pgfpathquadraticcurveto{\pgfqpoint{6.158235in}{0.728795in}}{\pgfqpoint{6.158235in}{0.744267in}}%
\pgfpathquadraticcurveto{\pgfqpoint{6.158235in}{0.759794in}}{\pgfqpoint{6.165908in}{0.768566in}}%
\pgfpathquadraticcurveto{\pgfqpoint{6.173599in}{0.777356in}}{\pgfqpoint{6.187163in}{0.777356in}}%
\pgfpathlineto{\pgfqpoint{6.187163in}{0.777356in}}%
\pgfpathclose%
\pgfpathmoveto{\pgfqpoint{6.273459in}{0.773987in}}%
\pgfpathlineto{\pgfqpoint{6.273459in}{0.764189in}}%
\pgfpathquadraticcurveto{\pgfqpoint{6.269082in}{0.766440in}}{\pgfqpoint{6.264345in}{0.767557in}}%
\pgfpathquadraticcurveto{\pgfqpoint{6.259625in}{0.768692in}}{\pgfqpoint{6.254546in}{0.768692in}}%
\pgfpathquadraticcurveto{\pgfqpoint{6.246837in}{0.768692in}}{\pgfqpoint{6.242982in}{0.766332in}}%
\pgfpathquadraticcurveto{\pgfqpoint{6.239128in}{0.763973in}}{\pgfqpoint{6.239128in}{0.759235in}}%
\pgfpathquadraticcurveto{\pgfqpoint{6.239128in}{0.755633in}}{\pgfqpoint{6.241883in}{0.753580in}}%
\pgfpathquadraticcurveto{\pgfqpoint{6.244639in}{0.751526in}}{\pgfqpoint{6.252979in}{0.749671in}}%
\pgfpathlineto{\pgfqpoint{6.256527in}{0.748878in}}%
\pgfpathquadraticcurveto{\pgfqpoint{6.267551in}{0.746519in}}{\pgfqpoint{6.272198in}{0.742214in}}%
\pgfpathquadraticcurveto{\pgfqpoint{6.276845in}{0.737909in}}{\pgfqpoint{6.276845in}{0.730200in}}%
\pgfpathquadraticcurveto{\pgfqpoint{6.276845in}{0.721410in}}{\pgfqpoint{6.269892in}{0.716276in}}%
\pgfpathquadraticcurveto{\pgfqpoint{6.262940in}{0.711161in}}{\pgfqpoint{6.250781in}{0.711161in}}%
\pgfpathquadraticcurveto{\pgfqpoint{6.245720in}{0.711161in}}{\pgfqpoint{6.240226in}{0.712152in}}%
\pgfpathquadraticcurveto{\pgfqpoint{6.234733in}{0.713142in}}{\pgfqpoint{6.228663in}{0.715106in}}%
\pgfpathlineto{\pgfqpoint{6.228663in}{0.725805in}}%
\pgfpathquadraticcurveto{\pgfqpoint{6.234408in}{0.722815in}}{\pgfqpoint{6.239974in}{0.721320in}}%
\pgfpathquadraticcurveto{\pgfqpoint{6.245540in}{0.719843in}}{\pgfqpoint{6.251016in}{0.719843in}}%
\pgfpathquadraticcurveto{\pgfqpoint{6.258329in}{0.719843in}}{\pgfqpoint{6.262255in}{0.722346in}}%
\pgfpathquadraticcurveto{\pgfqpoint{6.266200in}{0.724850in}}{\pgfqpoint{6.266200in}{0.729407in}}%
\pgfpathquadraticcurveto{\pgfqpoint{6.266200in}{0.733622in}}{\pgfqpoint{6.263354in}{0.735874in}}%
\pgfpathquadraticcurveto{\pgfqpoint{6.260526in}{0.738125in}}{\pgfqpoint{6.250890in}{0.740214in}}%
\pgfpathlineto{\pgfqpoint{6.247287in}{0.741061in}}%
\pgfpathquadraticcurveto{\pgfqpoint{6.237669in}{0.743078in}}{\pgfqpoint{6.233382in}{0.747275in}}%
\pgfpathquadraticcurveto{\pgfqpoint{6.229113in}{0.751472in}}{\pgfqpoint{6.229113in}{0.758785in}}%
\pgfpathquadraticcurveto{\pgfqpoint{6.229113in}{0.767683in}}{\pgfqpoint{6.235417in}{0.772510in}}%
\pgfpathquadraticcurveto{\pgfqpoint{6.241721in}{0.777356in}}{\pgfqpoint{6.253321in}{0.777356in}}%
\pgfpathquadraticcurveto{\pgfqpoint{6.259049in}{0.777356in}}{\pgfqpoint{6.264110in}{0.776509in}}%
\pgfpathquadraticcurveto{\pgfqpoint{6.269190in}{0.775680in}}{\pgfqpoint{6.273459in}{0.773987in}}%
\pgfpathlineto{\pgfqpoint{6.273459in}{0.773987in}}%
\pgfpathclose%
\pgfpathmoveto{\pgfqpoint{6.303575in}{0.793747in}}%
\pgfpathlineto{\pgfqpoint{6.303575in}{0.775843in}}%
\pgfpathlineto{\pgfqpoint{6.324901in}{0.775843in}}%
\pgfpathlineto{\pgfqpoint{6.324901in}{0.767791in}}%
\pgfpathlineto{\pgfqpoint{6.303575in}{0.767791in}}%
\pgfpathlineto{\pgfqpoint{6.303575in}{0.733568in}}%
\pgfpathquadraticcurveto{\pgfqpoint{6.303575in}{0.725859in}}{\pgfqpoint{6.305682in}{0.723661in}}%
\pgfpathquadraticcurveto{\pgfqpoint{6.307790in}{0.721464in}}{\pgfqpoint{6.314274in}{0.721464in}}%
\pgfpathlineto{\pgfqpoint{6.324901in}{0.721464in}}%
\pgfpathlineto{\pgfqpoint{6.324901in}{0.712800in}}%
\pgfpathlineto{\pgfqpoint{6.314274in}{0.712800in}}%
\pgfpathquadraticcurveto{\pgfqpoint{6.302278in}{0.712800in}}{\pgfqpoint{6.297721in}{0.717267in}}%
\pgfpathquadraticcurveto{\pgfqpoint{6.293164in}{0.721752in}}{\pgfqpoint{6.293164in}{0.733568in}}%
\pgfpathlineto{\pgfqpoint{6.293164in}{0.767791in}}%
\pgfpathlineto{\pgfqpoint{6.285563in}{0.767791in}}%
\pgfpathlineto{\pgfqpoint{6.285563in}{0.775843in}}%
\pgfpathlineto{\pgfqpoint{6.293164in}{0.775843in}}%
\pgfpathlineto{\pgfqpoint{6.293164in}{0.793747in}}%
\pgfpathlineto{\pgfqpoint{6.303575in}{0.793747in}}%
\pgfpathlineto{\pgfqpoint{6.303575in}{0.793747in}}%
\pgfpathclose%
\pgfpathmoveto{\pgfqpoint{6.338519in}{0.775843in}}%
\pgfpathlineto{\pgfqpoint{6.348876in}{0.775843in}}%
\pgfpathlineto{\pgfqpoint{6.348876in}{0.712800in}}%
\pgfpathlineto{\pgfqpoint{6.338519in}{0.712800in}}%
\pgfpathlineto{\pgfqpoint{6.338519in}{0.775843in}}%
\pgfpathlineto{\pgfqpoint{6.338519in}{0.775843in}}%
\pgfpathclose%
\pgfpathmoveto{\pgfqpoint{6.338519in}{0.800393in}}%
\pgfpathlineto{\pgfqpoint{6.348876in}{0.800393in}}%
\pgfpathlineto{\pgfqpoint{6.348876in}{0.787262in}}%
\pgfpathlineto{\pgfqpoint{6.338519in}{0.787262in}}%
\pgfpathlineto{\pgfqpoint{6.338519in}{0.800393in}}%
\pgfpathlineto{\pgfqpoint{6.338519in}{0.800393in}}%
\pgfpathclose%
\pgfpathmoveto{\pgfqpoint{6.415917in}{0.773429in}}%
\pgfpathlineto{\pgfqpoint{6.415917in}{0.763738in}}%
\pgfpathquadraticcurveto{\pgfqpoint{6.411522in}{0.766170in}}{\pgfqpoint{6.407109in}{0.767377in}}%
\pgfpathquadraticcurveto{\pgfqpoint{6.402696in}{0.768584in}}{\pgfqpoint{6.398193in}{0.768584in}}%
\pgfpathquadraticcurveto{\pgfqpoint{6.388106in}{0.768584in}}{\pgfqpoint{6.382522in}{0.762189in}}%
\pgfpathquadraticcurveto{\pgfqpoint{6.376957in}{0.755813in}}{\pgfqpoint{6.376957in}{0.744267in}}%
\pgfpathquadraticcurveto{\pgfqpoint{6.376957in}{0.732721in}}{\pgfqpoint{6.382522in}{0.726327in}}%
\pgfpathquadraticcurveto{\pgfqpoint{6.388106in}{0.719951in}}{\pgfqpoint{6.398193in}{0.719951in}}%
\pgfpathquadraticcurveto{\pgfqpoint{6.402696in}{0.719951in}}{\pgfqpoint{6.407109in}{0.721158in}}%
\pgfpathquadraticcurveto{\pgfqpoint{6.411522in}{0.722364in}}{\pgfqpoint{6.415917in}{0.724796in}}%
\pgfpathlineto{\pgfqpoint{6.415917in}{0.715214in}}%
\pgfpathquadraticcurveto{\pgfqpoint{6.411576in}{0.713196in}}{\pgfqpoint{6.406929in}{0.712188in}}%
\pgfpathquadraticcurveto{\pgfqpoint{6.402300in}{0.711161in}}{\pgfqpoint{6.397058in}{0.711161in}}%
\pgfpathquadraticcurveto{\pgfqpoint{6.382811in}{0.711161in}}{\pgfqpoint{6.374417in}{0.720113in}}%
\pgfpathquadraticcurveto{\pgfqpoint{6.366041in}{0.729065in}}{\pgfqpoint{6.366041in}{0.744267in}}%
\pgfpathquadraticcurveto{\pgfqpoint{6.366041in}{0.759686in}}{\pgfqpoint{6.374507in}{0.768512in}}%
\pgfpathquadraticcurveto{\pgfqpoint{6.382991in}{0.777356in}}{\pgfqpoint{6.397743in}{0.777356in}}%
\pgfpathquadraticcurveto{\pgfqpoint{6.402516in}{0.777356in}}{\pgfqpoint{6.407073in}{0.776365in}}%
\pgfpathquadraticcurveto{\pgfqpoint{6.411630in}{0.775392in}}{\pgfqpoint{6.415917in}{0.773429in}}%
\pgfpathlineto{\pgfqpoint{6.415917in}{0.773429in}}%
\pgfpathclose%
\pgfpathmoveto{\pgfqpoint{6.436577in}{0.727102in}}%
\pgfpathlineto{\pgfqpoint{6.448447in}{0.727102in}}%
\pgfpathlineto{\pgfqpoint{6.448447in}{0.717411in}}%
\pgfpathlineto{\pgfqpoint{6.439225in}{0.699399in}}%
\pgfpathlineto{\pgfqpoint{6.431966in}{0.699399in}}%
\pgfpathlineto{\pgfqpoint{6.436577in}{0.717411in}}%
\pgfpathlineto{\pgfqpoint{6.436577in}{0.727102in}}%
\pgfpathlineto{\pgfqpoint{6.436577in}{0.727102in}}%
\pgfpathclose%
\pgfusepath{fill}%
\end{pgfscope}%
\begin{pgfscope}%
\pgfsetbuttcap%
\pgfsetroundjoin%
\definecolor{currentfill}{rgb}{0.309804,0.372549,0.419608}%
\pgfsetfillcolor{currentfill}%
\pgfsetlinewidth{0.000000pt}%
\definecolor{currentstroke}{rgb}{0.309804,0.372549,0.419608}%
\pgfsetstrokecolor{currentstroke}%
\pgfsetdash{}{0pt}%
\pgfpathmoveto{\pgfqpoint{6.594753in}{0.750860in}}%
\pgfpathlineto{\pgfqpoint{6.594753in}{0.712800in}}%
\pgfpathlineto{\pgfqpoint{6.584396in}{0.712800in}}%
\pgfpathlineto{\pgfqpoint{6.584396in}{0.750517in}}%
\pgfpathquadraticcurveto{\pgfqpoint{6.584396in}{0.759469in}}{\pgfqpoint{6.580902in}{0.763900in}}%
\pgfpathquadraticcurveto{\pgfqpoint{6.577407in}{0.768349in}}{\pgfqpoint{6.570437in}{0.768349in}}%
\pgfpathquadraticcurveto{\pgfqpoint{6.562043in}{0.768349in}}{\pgfqpoint{6.557198in}{0.763000in}}%
\pgfpathquadraticcurveto{\pgfqpoint{6.552352in}{0.757668in}}{\pgfqpoint{6.552352in}{0.748428in}}%
\pgfpathlineto{\pgfqpoint{6.552352in}{0.712800in}}%
\pgfpathlineto{\pgfqpoint{6.541941in}{0.712800in}}%
\pgfpathlineto{\pgfqpoint{6.541941in}{0.775843in}}%
\pgfpathlineto{\pgfqpoint{6.552352in}{0.775843in}}%
\pgfpathlineto{\pgfqpoint{6.552352in}{0.766044in}}%
\pgfpathquadraticcurveto{\pgfqpoint{6.556081in}{0.771736in}}{\pgfqpoint{6.561106in}{0.774546in}}%
\pgfpathquadraticcurveto{\pgfqpoint{6.566150in}{0.777356in}}{\pgfqpoint{6.572742in}{0.777356in}}%
\pgfpathquadraticcurveto{\pgfqpoint{6.583603in}{0.777356in}}{\pgfqpoint{6.589169in}{0.770637in}}%
\pgfpathquadraticcurveto{\pgfqpoint{6.594753in}{0.763918in}}{\pgfqpoint{6.594753in}{0.750860in}}%
\pgfpathlineto{\pgfqpoint{6.594753in}{0.750860in}}%
\pgfpathclose%
\pgfpathmoveto{\pgfqpoint{6.639819in}{0.768584in}}%
\pgfpathquadraticcurveto{\pgfqpoint{6.631498in}{0.768584in}}{\pgfqpoint{6.626653in}{0.762081in}}%
\pgfpathquadraticcurveto{\pgfqpoint{6.621807in}{0.755579in}}{\pgfqpoint{6.621807in}{0.744267in}}%
\pgfpathquadraticcurveto{\pgfqpoint{6.621807in}{0.732956in}}{\pgfqpoint{6.626616in}{0.726453in}}%
\pgfpathquadraticcurveto{\pgfqpoint{6.631444in}{0.719951in}}{\pgfqpoint{6.639819in}{0.719951in}}%
\pgfpathquadraticcurveto{\pgfqpoint{6.648105in}{0.719951in}}{\pgfqpoint{6.652932in}{0.726471in}}%
\pgfpathquadraticcurveto{\pgfqpoint{6.657778in}{0.733010in}}{\pgfqpoint{6.657778in}{0.744267in}}%
\pgfpathquadraticcurveto{\pgfqpoint{6.657778in}{0.755471in}}{\pgfqpoint{6.652932in}{0.762027in}}%
\pgfpathquadraticcurveto{\pgfqpoint{6.648105in}{0.768584in}}{\pgfqpoint{6.639819in}{0.768584in}}%
\pgfpathlineto{\pgfqpoint{6.639819in}{0.768584in}}%
\pgfpathclose%
\pgfpathmoveto{\pgfqpoint{6.639819in}{0.777356in}}%
\pgfpathquadraticcurveto{\pgfqpoint{6.653329in}{0.777356in}}{\pgfqpoint{6.661038in}{0.768566in}}%
\pgfpathquadraticcurveto{\pgfqpoint{6.668765in}{0.759794in}}{\pgfqpoint{6.668765in}{0.744267in}}%
\pgfpathquadraticcurveto{\pgfqpoint{6.668765in}{0.728795in}}{\pgfqpoint{6.661038in}{0.719969in}}%
\pgfpathquadraticcurveto{\pgfqpoint{6.653329in}{0.711161in}}{\pgfqpoint{6.639819in}{0.711161in}}%
\pgfpathquadraticcurveto{\pgfqpoint{6.626256in}{0.711161in}}{\pgfqpoint{6.618565in}{0.719969in}}%
\pgfpathquadraticcurveto{\pgfqpoint{6.610892in}{0.728795in}}{\pgfqpoint{6.610892in}{0.744267in}}%
\pgfpathquadraticcurveto{\pgfqpoint{6.610892in}{0.759794in}}{\pgfqpoint{6.618565in}{0.768566in}}%
\pgfpathquadraticcurveto{\pgfqpoint{6.626256in}{0.777356in}}{\pgfqpoint{6.639819in}{0.777356in}}%
\pgfpathlineto{\pgfqpoint{6.639819in}{0.777356in}}%
\pgfpathclose%
\pgfpathmoveto{\pgfqpoint{6.696179in}{0.793747in}}%
\pgfpathlineto{\pgfqpoint{6.696179in}{0.775843in}}%
\pgfpathlineto{\pgfqpoint{6.717506in}{0.775843in}}%
\pgfpathlineto{\pgfqpoint{6.717506in}{0.767791in}}%
\pgfpathlineto{\pgfqpoint{6.696179in}{0.767791in}}%
\pgfpathlineto{\pgfqpoint{6.696179in}{0.733568in}}%
\pgfpathquadraticcurveto{\pgfqpoint{6.696179in}{0.725859in}}{\pgfqpoint{6.698287in}{0.723661in}}%
\pgfpathquadraticcurveto{\pgfqpoint{6.700394in}{0.721464in}}{\pgfqpoint{6.706879in}{0.721464in}}%
\pgfpathlineto{\pgfqpoint{6.717506in}{0.721464in}}%
\pgfpathlineto{\pgfqpoint{6.717506in}{0.712800in}}%
\pgfpathlineto{\pgfqpoint{6.706879in}{0.712800in}}%
\pgfpathquadraticcurveto{\pgfqpoint{6.694883in}{0.712800in}}{\pgfqpoint{6.690325in}{0.717267in}}%
\pgfpathquadraticcurveto{\pgfqpoint{6.685768in}{0.721752in}}{\pgfqpoint{6.685768in}{0.733568in}}%
\pgfpathlineto{\pgfqpoint{6.685768in}{0.767791in}}%
\pgfpathlineto{\pgfqpoint{6.678167in}{0.767791in}}%
\pgfpathlineto{\pgfqpoint{6.678167in}{0.775843in}}%
\pgfpathlineto{\pgfqpoint{6.685768in}{0.775843in}}%
\pgfpathlineto{\pgfqpoint{6.685768in}{0.793747in}}%
\pgfpathlineto{\pgfqpoint{6.696179in}{0.793747in}}%
\pgfpathlineto{\pgfqpoint{6.696179in}{0.793747in}}%
\pgfpathclose%
\pgfusepath{fill}%
\end{pgfscope}%
\begin{pgfscope}%
\pgfsetbuttcap%
\pgfsetroundjoin%
\definecolor{currentfill}{rgb}{0.309804,0.372549,0.419608}%
\pgfsetfillcolor{currentfill}%
\pgfsetlinewidth{0.000000pt}%
\definecolor{currentstroke}{rgb}{0.309804,0.372549,0.419608}%
\pgfsetstrokecolor{currentstroke}%
\pgfsetdash{}{0pt}%
\pgfpathmoveto{\pgfqpoint{6.812688in}{0.744483in}}%
\pgfpathquadraticcurveto{\pgfqpoint{6.800133in}{0.744483in}}{\pgfqpoint{6.795288in}{0.741619in}}%
\pgfpathquadraticcurveto{\pgfqpoint{6.790443in}{0.738756in}}{\pgfqpoint{6.790443in}{0.731821in}}%
\pgfpathquadraticcurveto{\pgfqpoint{6.790443in}{0.726309in}}{\pgfqpoint{6.794081in}{0.723067in}}%
\pgfpathquadraticcurveto{\pgfqpoint{6.797720in}{0.719843in}}{\pgfqpoint{6.803952in}{0.719843in}}%
\pgfpathquadraticcurveto{\pgfqpoint{6.812580in}{0.719843in}}{\pgfqpoint{6.817785in}{0.725949in}}%
\pgfpathquadraticcurveto{\pgfqpoint{6.822991in}{0.732055in}}{\pgfqpoint{6.822991in}{0.742178in}}%
\pgfpathlineto{\pgfqpoint{6.822991in}{0.744483in}}%
\pgfpathlineto{\pgfqpoint{6.812688in}{0.744483in}}%
\pgfpathlineto{\pgfqpoint{6.812688in}{0.744483in}}%
\pgfpathclose%
\pgfpathmoveto{\pgfqpoint{6.833348in}{0.748770in}}%
\pgfpathlineto{\pgfqpoint{6.833348in}{0.712800in}}%
\pgfpathlineto{\pgfqpoint{6.822991in}{0.712800in}}%
\pgfpathlineto{\pgfqpoint{6.822991in}{0.722364in}}%
\pgfpathquadraticcurveto{\pgfqpoint{6.819442in}{0.716637in}}{\pgfqpoint{6.814147in}{0.713899in}}%
\pgfpathquadraticcurveto{\pgfqpoint{6.808851in}{0.711161in}}{\pgfqpoint{6.801196in}{0.711161in}}%
\pgfpathquadraticcurveto{\pgfqpoint{6.791524in}{0.711161in}}{\pgfqpoint{6.785796in}{0.716601in}}%
\pgfpathquadraticcurveto{\pgfqpoint{6.780086in}{0.722040in}}{\pgfqpoint{6.780086in}{0.731154in}}%
\pgfpathquadraticcurveto{\pgfqpoint{6.780086in}{0.741782in}}{\pgfqpoint{6.787201in}{0.747185in}}%
\pgfpathquadraticcurveto{\pgfqpoint{6.794333in}{0.752589in}}{\pgfqpoint{6.808455in}{0.752589in}}%
\pgfpathlineto{\pgfqpoint{6.822991in}{0.752589in}}%
\pgfpathlineto{\pgfqpoint{6.822991in}{0.753616in}}%
\pgfpathquadraticcurveto{\pgfqpoint{6.822991in}{0.760766in}}{\pgfqpoint{6.818290in}{0.764675in}}%
\pgfpathquadraticcurveto{\pgfqpoint{6.813588in}{0.768584in}}{\pgfqpoint{6.805087in}{0.768584in}}%
\pgfpathquadraticcurveto{\pgfqpoint{6.799683in}{0.768584in}}{\pgfqpoint{6.794550in}{0.767287in}}%
\pgfpathquadraticcurveto{\pgfqpoint{6.789434in}{0.765990in}}{\pgfqpoint{6.784715in}{0.763396in}}%
\pgfpathlineto{\pgfqpoint{6.784715in}{0.772979in}}%
\pgfpathquadraticcurveto{\pgfqpoint{6.790389in}{0.775176in}}{\pgfqpoint{6.795738in}{0.776257in}}%
\pgfpathquadraticcurveto{\pgfqpoint{6.801088in}{0.777356in}}{\pgfqpoint{6.806149in}{0.777356in}}%
\pgfpathquadraticcurveto{\pgfqpoint{6.819839in}{0.777356in}}{\pgfqpoint{6.826593in}{0.770259in}}%
\pgfpathquadraticcurveto{\pgfqpoint{6.833348in}{0.763180in}}{\pgfqpoint{6.833348in}{0.748770in}}%
\pgfpathlineto{\pgfqpoint{6.833348in}{0.748770in}}%
\pgfpathclose%
\pgfusepath{fill}%
\end{pgfscope}%
\begin{pgfscope}%
\pgfsetbuttcap%
\pgfsetroundjoin%
\definecolor{currentfill}{rgb}{0.309804,0.372549,0.419608}%
\pgfsetfillcolor{currentfill}%
\pgfsetlinewidth{0.000000pt}%
\definecolor{currentstroke}{rgb}{0.309804,0.372549,0.419608}%
\pgfsetstrokecolor{currentstroke}%
\pgfsetdash{}{0pt}%
\pgfpathmoveto{\pgfqpoint{6.888302in}{0.775843in}}%
\pgfpathlineto{\pgfqpoint{6.899272in}{0.775843in}}%
\pgfpathlineto{\pgfqpoint{6.918977in}{0.722941in}}%
\pgfpathlineto{\pgfqpoint{6.938682in}{0.775843in}}%
\pgfpathlineto{\pgfqpoint{6.949652in}{0.775843in}}%
\pgfpathlineto{\pgfqpoint{6.926002in}{0.712800in}}%
\pgfpathlineto{\pgfqpoint{6.911934in}{0.712800in}}%
\pgfpathlineto{\pgfqpoint{6.888302in}{0.775843in}}%
\pgfpathlineto{\pgfqpoint{6.888302in}{0.775843in}}%
\pgfpathclose%
\pgfpathmoveto{\pgfqpoint{6.992611in}{0.744483in}}%
\pgfpathquadraticcurveto{\pgfqpoint{6.980056in}{0.744483in}}{\pgfqpoint{6.975211in}{0.741619in}}%
\pgfpathquadraticcurveto{\pgfqpoint{6.970366in}{0.738756in}}{\pgfqpoint{6.970366in}{0.731821in}}%
\pgfpathquadraticcurveto{\pgfqpoint{6.970366in}{0.726309in}}{\pgfqpoint{6.974004in}{0.723067in}}%
\pgfpathquadraticcurveto{\pgfqpoint{6.977643in}{0.719843in}}{\pgfqpoint{6.983875in}{0.719843in}}%
\pgfpathquadraticcurveto{\pgfqpoint{6.992503in}{0.719843in}}{\pgfqpoint{6.997708in}{0.725949in}}%
\pgfpathquadraticcurveto{\pgfqpoint{7.002914in}{0.732055in}}{\pgfqpoint{7.002914in}{0.742178in}}%
\pgfpathlineto{\pgfqpoint{7.002914in}{0.744483in}}%
\pgfpathlineto{\pgfqpoint{6.992611in}{0.744483in}}%
\pgfpathlineto{\pgfqpoint{6.992611in}{0.744483in}}%
\pgfpathclose%
\pgfpathmoveto{\pgfqpoint{7.013271in}{0.748770in}}%
\pgfpathlineto{\pgfqpoint{7.013271in}{0.712800in}}%
\pgfpathlineto{\pgfqpoint{7.002914in}{0.712800in}}%
\pgfpathlineto{\pgfqpoint{7.002914in}{0.722364in}}%
\pgfpathquadraticcurveto{\pgfqpoint{6.999365in}{0.716637in}}{\pgfqpoint{6.994070in}{0.713899in}}%
\pgfpathquadraticcurveto{\pgfqpoint{6.988774in}{0.711161in}}{\pgfqpoint{6.981119in}{0.711161in}}%
\pgfpathquadraticcurveto{\pgfqpoint{6.971446in}{0.711161in}}{\pgfqpoint{6.965719in}{0.716601in}}%
\pgfpathquadraticcurveto{\pgfqpoint{6.960009in}{0.722040in}}{\pgfqpoint{6.960009in}{0.731154in}}%
\pgfpathquadraticcurveto{\pgfqpoint{6.960009in}{0.741782in}}{\pgfqpoint{6.967124in}{0.747185in}}%
\pgfpathquadraticcurveto{\pgfqpoint{6.974256in}{0.752589in}}{\pgfqpoint{6.988378in}{0.752589in}}%
\pgfpathlineto{\pgfqpoint{7.002914in}{0.752589in}}%
\pgfpathlineto{\pgfqpoint{7.002914in}{0.753616in}}%
\pgfpathquadraticcurveto{\pgfqpoint{7.002914in}{0.760766in}}{\pgfqpoint{6.998213in}{0.764675in}}%
\pgfpathquadraticcurveto{\pgfqpoint{6.993511in}{0.768584in}}{\pgfqpoint{6.985010in}{0.768584in}}%
\pgfpathquadraticcurveto{\pgfqpoint{6.979606in}{0.768584in}}{\pgfqpoint{6.974473in}{0.767287in}}%
\pgfpathquadraticcurveto{\pgfqpoint{6.969357in}{0.765990in}}{\pgfqpoint{6.964638in}{0.763396in}}%
\pgfpathlineto{\pgfqpoint{6.964638in}{0.772979in}}%
\pgfpathquadraticcurveto{\pgfqpoint{6.970312in}{0.775176in}}{\pgfqpoint{6.975661in}{0.776257in}}%
\pgfpathquadraticcurveto{\pgfqpoint{6.981011in}{0.777356in}}{\pgfqpoint{6.986072in}{0.777356in}}%
\pgfpathquadraticcurveto{\pgfqpoint{6.999762in}{0.777356in}}{\pgfqpoint{7.006516in}{0.770259in}}%
\pgfpathquadraticcurveto{\pgfqpoint{7.013271in}{0.763180in}}{\pgfqpoint{7.013271in}{0.748770in}}%
\pgfpathlineto{\pgfqpoint{7.013271in}{0.748770in}}%
\pgfpathclose%
\pgfpathmoveto{\pgfqpoint{7.071126in}{0.766170in}}%
\pgfpathquadraticcurveto{\pgfqpoint{7.069379in}{0.767179in}}{\pgfqpoint{7.067325in}{0.767647in}}%
\pgfpathquadraticcurveto{\pgfqpoint{7.065272in}{0.768133in}}{\pgfqpoint{7.062804in}{0.768133in}}%
\pgfpathquadraticcurveto{\pgfqpoint{7.054014in}{0.768133in}}{\pgfqpoint{7.049313in}{0.762423in}}%
\pgfpathquadraticcurveto{\pgfqpoint{7.044612in}{0.756714in}}{\pgfqpoint{7.044612in}{0.746014in}}%
\pgfpathlineto{\pgfqpoint{7.044612in}{0.712800in}}%
\pgfpathlineto{\pgfqpoint{7.034201in}{0.712800in}}%
\pgfpathlineto{\pgfqpoint{7.034201in}{0.775843in}}%
\pgfpathlineto{\pgfqpoint{7.044612in}{0.775843in}}%
\pgfpathlineto{\pgfqpoint{7.044612in}{0.766044in}}%
\pgfpathquadraticcurveto{\pgfqpoint{7.047890in}{0.771790in}}{\pgfqpoint{7.053114in}{0.774564in}}%
\pgfpathquadraticcurveto{\pgfqpoint{7.058355in}{0.777356in}}{\pgfqpoint{7.065848in}{0.777356in}}%
\pgfpathquadraticcurveto{\pgfqpoint{7.066911in}{0.777356in}}{\pgfqpoint{7.068208in}{0.777211in}}%
\pgfpathquadraticcurveto{\pgfqpoint{7.069505in}{0.777085in}}{\pgfqpoint{7.071072in}{0.776797in}}%
\pgfpathlineto{\pgfqpoint{7.071126in}{0.766170in}}%
\pgfpathlineto{\pgfqpoint{7.071126in}{0.766170in}}%
\pgfpathclose%
\pgfpathmoveto{\pgfqpoint{7.081987in}{0.775843in}}%
\pgfpathlineto{\pgfqpoint{7.092344in}{0.775843in}}%
\pgfpathlineto{\pgfqpoint{7.092344in}{0.712800in}}%
\pgfpathlineto{\pgfqpoint{7.081987in}{0.712800in}}%
\pgfpathlineto{\pgfqpoint{7.081987in}{0.775843in}}%
\pgfpathlineto{\pgfqpoint{7.081987in}{0.775843in}}%
\pgfpathclose%
\pgfpathmoveto{\pgfqpoint{7.081987in}{0.800393in}}%
\pgfpathlineto{\pgfqpoint{7.092344in}{0.800393in}}%
\pgfpathlineto{\pgfqpoint{7.092344in}{0.787262in}}%
\pgfpathlineto{\pgfqpoint{7.081987in}{0.787262in}}%
\pgfpathlineto{\pgfqpoint{7.081987in}{0.800393in}}%
\pgfpathlineto{\pgfqpoint{7.081987in}{0.800393in}}%
\pgfpathclose%
\pgfpathmoveto{\pgfqpoint{7.142670in}{0.744483in}}%
\pgfpathquadraticcurveto{\pgfqpoint{7.130116in}{0.744483in}}{\pgfqpoint{7.125270in}{0.741619in}}%
\pgfpathquadraticcurveto{\pgfqpoint{7.120425in}{0.738756in}}{\pgfqpoint{7.120425in}{0.731821in}}%
\pgfpathquadraticcurveto{\pgfqpoint{7.120425in}{0.726309in}}{\pgfqpoint{7.124063in}{0.723067in}}%
\pgfpathquadraticcurveto{\pgfqpoint{7.127702in}{0.719843in}}{\pgfqpoint{7.133934in}{0.719843in}}%
\pgfpathquadraticcurveto{\pgfqpoint{7.142562in}{0.719843in}}{\pgfqpoint{7.147767in}{0.725949in}}%
\pgfpathquadraticcurveto{\pgfqpoint{7.152973in}{0.732055in}}{\pgfqpoint{7.152973in}{0.742178in}}%
\pgfpathlineto{\pgfqpoint{7.152973in}{0.744483in}}%
\pgfpathlineto{\pgfqpoint{7.142670in}{0.744483in}}%
\pgfpathlineto{\pgfqpoint{7.142670in}{0.744483in}}%
\pgfpathclose%
\pgfpathmoveto{\pgfqpoint{7.163330in}{0.748770in}}%
\pgfpathlineto{\pgfqpoint{7.163330in}{0.712800in}}%
\pgfpathlineto{\pgfqpoint{7.152973in}{0.712800in}}%
\pgfpathlineto{\pgfqpoint{7.152973in}{0.722364in}}%
\pgfpathquadraticcurveto{\pgfqpoint{7.149425in}{0.716637in}}{\pgfqpoint{7.144129in}{0.713899in}}%
\pgfpathquadraticcurveto{\pgfqpoint{7.138833in}{0.711161in}}{\pgfqpoint{7.131178in}{0.711161in}}%
\pgfpathquadraticcurveto{\pgfqpoint{7.121506in}{0.711161in}}{\pgfqpoint{7.115778in}{0.716601in}}%
\pgfpathquadraticcurveto{\pgfqpoint{7.110068in}{0.722040in}}{\pgfqpoint{7.110068in}{0.731154in}}%
\pgfpathquadraticcurveto{\pgfqpoint{7.110068in}{0.741782in}}{\pgfqpoint{7.117183in}{0.747185in}}%
\pgfpathquadraticcurveto{\pgfqpoint{7.124316in}{0.752589in}}{\pgfqpoint{7.138437in}{0.752589in}}%
\pgfpathlineto{\pgfqpoint{7.152973in}{0.752589in}}%
\pgfpathlineto{\pgfqpoint{7.152973in}{0.753616in}}%
\pgfpathquadraticcurveto{\pgfqpoint{7.152973in}{0.760766in}}{\pgfqpoint{7.148272in}{0.764675in}}%
\pgfpathquadraticcurveto{\pgfqpoint{7.143571in}{0.768584in}}{\pgfqpoint{7.135069in}{0.768584in}}%
\pgfpathquadraticcurveto{\pgfqpoint{7.129665in}{0.768584in}}{\pgfqpoint{7.124532in}{0.767287in}}%
\pgfpathquadraticcurveto{\pgfqpoint{7.119416in}{0.765990in}}{\pgfqpoint{7.114697in}{0.763396in}}%
\pgfpathlineto{\pgfqpoint{7.114697in}{0.772979in}}%
\pgfpathquadraticcurveto{\pgfqpoint{7.120371in}{0.775176in}}{\pgfqpoint{7.125721in}{0.776257in}}%
\pgfpathquadraticcurveto{\pgfqpoint{7.131070in}{0.777356in}}{\pgfqpoint{7.136132in}{0.777356in}}%
\pgfpathquadraticcurveto{\pgfqpoint{7.149821in}{0.777356in}}{\pgfqpoint{7.156575in}{0.770259in}}%
\pgfpathquadraticcurveto{\pgfqpoint{7.163330in}{0.763180in}}{\pgfqpoint{7.163330in}{0.748770in}}%
\pgfpathlineto{\pgfqpoint{7.163330in}{0.748770in}}%
\pgfpathclose%
\pgfpathmoveto{\pgfqpoint{7.237072in}{0.750860in}}%
\pgfpathlineto{\pgfqpoint{7.237072in}{0.712800in}}%
\pgfpathlineto{\pgfqpoint{7.226715in}{0.712800in}}%
\pgfpathlineto{\pgfqpoint{7.226715in}{0.750517in}}%
\pgfpathquadraticcurveto{\pgfqpoint{7.226715in}{0.759469in}}{\pgfqpoint{7.223220in}{0.763900in}}%
\pgfpathquadraticcurveto{\pgfqpoint{7.219726in}{0.768349in}}{\pgfqpoint{7.212755in}{0.768349in}}%
\pgfpathquadraticcurveto{\pgfqpoint{7.204362in}{0.768349in}}{\pgfqpoint{7.199516in}{0.763000in}}%
\pgfpathquadraticcurveto{\pgfqpoint{7.194671in}{0.757668in}}{\pgfqpoint{7.194671in}{0.748428in}}%
\pgfpathlineto{\pgfqpoint{7.194671in}{0.712800in}}%
\pgfpathlineto{\pgfqpoint{7.184260in}{0.712800in}}%
\pgfpathlineto{\pgfqpoint{7.184260in}{0.775843in}}%
\pgfpathlineto{\pgfqpoint{7.194671in}{0.775843in}}%
\pgfpathlineto{\pgfqpoint{7.194671in}{0.766044in}}%
\pgfpathquadraticcurveto{\pgfqpoint{7.198400in}{0.771736in}}{\pgfqpoint{7.203425in}{0.774546in}}%
\pgfpathquadraticcurveto{\pgfqpoint{7.208468in}{0.777356in}}{\pgfqpoint{7.215061in}{0.777356in}}%
\pgfpathquadraticcurveto{\pgfqpoint{7.225922in}{0.777356in}}{\pgfqpoint{7.231488in}{0.770637in}}%
\pgfpathquadraticcurveto{\pgfqpoint{7.237072in}{0.763918in}}{\pgfqpoint{7.237072in}{0.750860in}}%
\pgfpathlineto{\pgfqpoint{7.237072in}{0.750860in}}%
\pgfpathclose%
\pgfpathmoveto{\pgfqpoint{7.303086in}{0.773429in}}%
\pgfpathlineto{\pgfqpoint{7.303086in}{0.763738in}}%
\pgfpathquadraticcurveto{\pgfqpoint{7.298691in}{0.766170in}}{\pgfqpoint{7.294278in}{0.767377in}}%
\pgfpathquadraticcurveto{\pgfqpoint{7.289865in}{0.768584in}}{\pgfqpoint{7.285362in}{0.768584in}}%
\pgfpathquadraticcurveto{\pgfqpoint{7.275275in}{0.768584in}}{\pgfqpoint{7.269692in}{0.762189in}}%
\pgfpathquadraticcurveto{\pgfqpoint{7.264126in}{0.755813in}}{\pgfqpoint{7.264126in}{0.744267in}}%
\pgfpathquadraticcurveto{\pgfqpoint{7.264126in}{0.732721in}}{\pgfqpoint{7.269692in}{0.726327in}}%
\pgfpathquadraticcurveto{\pgfqpoint{7.275275in}{0.719951in}}{\pgfqpoint{7.285362in}{0.719951in}}%
\pgfpathquadraticcurveto{\pgfqpoint{7.289865in}{0.719951in}}{\pgfqpoint{7.294278in}{0.721158in}}%
\pgfpathquadraticcurveto{\pgfqpoint{7.298691in}{0.722364in}}{\pgfqpoint{7.303086in}{0.724796in}}%
\pgfpathlineto{\pgfqpoint{7.303086in}{0.715214in}}%
\pgfpathquadraticcurveto{\pgfqpoint{7.298745in}{0.713196in}}{\pgfqpoint{7.294098in}{0.712188in}}%
\pgfpathquadraticcurveto{\pgfqpoint{7.289469in}{0.711161in}}{\pgfqpoint{7.284228in}{0.711161in}}%
\pgfpathquadraticcurveto{\pgfqpoint{7.269980in}{0.711161in}}{\pgfqpoint{7.261586in}{0.720113in}}%
\pgfpathquadraticcurveto{\pgfqpoint{7.253211in}{0.729065in}}{\pgfqpoint{7.253211in}{0.744267in}}%
\pgfpathquadraticcurveto{\pgfqpoint{7.253211in}{0.759686in}}{\pgfqpoint{7.261676in}{0.768512in}}%
\pgfpathquadraticcurveto{\pgfqpoint{7.270160in}{0.777356in}}{\pgfqpoint{7.284912in}{0.777356in}}%
\pgfpathquadraticcurveto{\pgfqpoint{7.289685in}{0.777356in}}{\pgfqpoint{7.294242in}{0.776365in}}%
\pgfpathquadraticcurveto{\pgfqpoint{7.298799in}{0.775392in}}{\pgfqpoint{7.303086in}{0.773429in}}%
\pgfpathlineto{\pgfqpoint{7.303086in}{0.773429in}}%
\pgfpathclose%
\pgfpathmoveto{\pgfqpoint{7.375027in}{0.746915in}}%
\pgfpathlineto{\pgfqpoint{7.375027in}{0.741854in}}%
\pgfpathlineto{\pgfqpoint{7.327403in}{0.741854in}}%
\pgfpathquadraticcurveto{\pgfqpoint{7.328087in}{0.731154in}}{\pgfqpoint{7.333851in}{0.725553in}}%
\pgfpathquadraticcurveto{\pgfqpoint{7.339615in}{0.719951in}}{\pgfqpoint{7.349918in}{0.719951in}}%
\pgfpathquadraticcurveto{\pgfqpoint{7.355880in}{0.719951in}}{\pgfqpoint{7.361482in}{0.721410in}}%
\pgfpathquadraticcurveto{\pgfqpoint{7.367083in}{0.722869in}}{\pgfqpoint{7.372613in}{0.725805in}}%
\pgfpathlineto{\pgfqpoint{7.372613in}{0.716006in}}%
\pgfpathquadraticcurveto{\pgfqpoint{7.367029in}{0.713647in}}{\pgfqpoint{7.361175in}{0.712404in}}%
\pgfpathquadraticcurveto{\pgfqpoint{7.355321in}{0.711161in}}{\pgfqpoint{7.349305in}{0.711161in}}%
\pgfpathquadraticcurveto{\pgfqpoint{7.334211in}{0.711161in}}{\pgfqpoint{7.325403in}{0.719933in}}%
\pgfpathquadraticcurveto{\pgfqpoint{7.316595in}{0.728723in}}{\pgfqpoint{7.316595in}{0.743709in}}%
\pgfpathquadraticcurveto{\pgfqpoint{7.316595in}{0.759181in}}{\pgfqpoint{7.324953in}{0.768259in}}%
\pgfpathquadraticcurveto{\pgfqpoint{7.333311in}{0.777356in}}{\pgfqpoint{7.347504in}{0.777356in}}%
\pgfpathquadraticcurveto{\pgfqpoint{7.360221in}{0.777356in}}{\pgfqpoint{7.367624in}{0.769160in}}%
\pgfpathquadraticcurveto{\pgfqpoint{7.375027in}{0.760983in}}{\pgfqpoint{7.375027in}{0.746915in}}%
\pgfpathlineto{\pgfqpoint{7.375027in}{0.746915in}}%
\pgfpathclose%
\pgfpathmoveto{\pgfqpoint{7.364670in}{0.749959in}}%
\pgfpathquadraticcurveto{\pgfqpoint{7.364562in}{0.758443in}}{\pgfqpoint{7.359915in}{0.763504in}}%
\pgfpathquadraticcurveto{\pgfqpoint{7.355267in}{0.768584in}}{\pgfqpoint{7.347612in}{0.768584in}}%
\pgfpathquadraticcurveto{\pgfqpoint{7.338948in}{0.768584in}}{\pgfqpoint{7.333743in}{0.763684in}}%
\pgfpathquadraticcurveto{\pgfqpoint{7.328537in}{0.758785in}}{\pgfqpoint{7.327745in}{0.749887in}}%
\pgfpathlineto{\pgfqpoint{7.364670in}{0.749959in}}%
\pgfpathlineto{\pgfqpoint{7.364670in}{0.749959in}}%
\pgfpathclose%
\pgfusepath{fill}%
\end{pgfscope}%
\begin{pgfscope}%
\pgfsetbuttcap%
\pgfsetroundjoin%
\definecolor{currentfill}{rgb}{0.309804,0.372549,0.419608}%
\pgfsetfillcolor{currentfill}%
\pgfsetlinewidth{0.000000pt}%
\definecolor{currentstroke}{rgb}{0.309804,0.372549,0.419608}%
\pgfsetstrokecolor{currentstroke}%
\pgfsetdash{}{0pt}%
\pgfpathmoveto{\pgfqpoint{7.488898in}{0.745060in}}%
\pgfpathquadraticcurveto{\pgfqpoint{7.488898in}{0.756317in}}{\pgfqpoint{7.484251in}{0.762496in}}%
\pgfpathquadraticcurveto{\pgfqpoint{7.479622in}{0.768692in}}{\pgfqpoint{7.471228in}{0.768692in}}%
\pgfpathquadraticcurveto{\pgfqpoint{7.462907in}{0.768692in}}{\pgfqpoint{7.458260in}{0.762496in}}%
\pgfpathquadraticcurveto{\pgfqpoint{7.453612in}{0.756317in}}{\pgfqpoint{7.453612in}{0.745060in}}%
\pgfpathquadraticcurveto{\pgfqpoint{7.453612in}{0.733856in}}{\pgfqpoint{7.458260in}{0.727660in}}%
\pgfpathquadraticcurveto{\pgfqpoint{7.462907in}{0.721464in}}{\pgfqpoint{7.471228in}{0.721464in}}%
\pgfpathquadraticcurveto{\pgfqpoint{7.479622in}{0.721464in}}{\pgfqpoint{7.484251in}{0.727660in}}%
\pgfpathquadraticcurveto{\pgfqpoint{7.488898in}{0.733856in}}{\pgfqpoint{7.488898in}{0.745060in}}%
\pgfpathlineto{\pgfqpoint{7.488898in}{0.745060in}}%
\pgfpathclose%
\pgfpathmoveto{\pgfqpoint{7.499255in}{0.720617in}}%
\pgfpathquadraticcurveto{\pgfqpoint{7.499255in}{0.704532in}}{\pgfqpoint{7.492104in}{0.696679in}}%
\pgfpathquadraticcurveto{\pgfqpoint{7.484972in}{0.688826in}}{\pgfqpoint{7.470220in}{0.688826in}}%
\pgfpathquadraticcurveto{\pgfqpoint{7.464762in}{0.688826in}}{\pgfqpoint{7.459917in}{0.689636in}}%
\pgfpathquadraticcurveto{\pgfqpoint{7.455071in}{0.690447in}}{\pgfqpoint{7.450514in}{0.692140in}}%
\pgfpathlineto{\pgfqpoint{7.450514in}{0.702209in}}%
\pgfpathquadraticcurveto{\pgfqpoint{7.455071in}{0.699741in}}{\pgfqpoint{7.459520in}{0.698570in}}%
\pgfpathquadraticcurveto{\pgfqpoint{7.463969in}{0.697382in}}{\pgfqpoint{7.468581in}{0.697382in}}%
\pgfpathquadraticcurveto{\pgfqpoint{7.478775in}{0.697382in}}{\pgfqpoint{7.483837in}{0.702695in}}%
\pgfpathquadraticcurveto{\pgfqpoint{7.488898in}{0.708009in}}{\pgfqpoint{7.488898in}{0.718762in}}%
\pgfpathlineto{\pgfqpoint{7.488898in}{0.723895in}}%
\pgfpathquadraticcurveto{\pgfqpoint{7.485692in}{0.718312in}}{\pgfqpoint{7.480685in}{0.715556in}}%
\pgfpathquadraticcurveto{\pgfqpoint{7.475677in}{0.712800in}}{\pgfqpoint{7.468689in}{0.712800in}}%
\pgfpathquadraticcurveto{\pgfqpoint{7.457107in}{0.712800in}}{\pgfqpoint{7.450010in}{0.721626in}}%
\pgfpathquadraticcurveto{\pgfqpoint{7.442913in}{0.730470in}}{\pgfqpoint{7.442913in}{0.745060in}}%
\pgfpathquadraticcurveto{\pgfqpoint{7.442913in}{0.759686in}}{\pgfqpoint{7.450010in}{0.768512in}}%
\pgfpathquadraticcurveto{\pgfqpoint{7.457107in}{0.777356in}}{\pgfqpoint{7.468689in}{0.777356in}}%
\pgfpathquadraticcurveto{\pgfqpoint{7.475677in}{0.777356in}}{\pgfqpoint{7.480685in}{0.774600in}}%
\pgfpathquadraticcurveto{\pgfqpoint{7.485692in}{0.771844in}}{\pgfqpoint{7.488898in}{0.766278in}}%
\pgfpathlineto{\pgfqpoint{7.488898in}{0.775843in}}%
\pgfpathlineto{\pgfqpoint{7.499255in}{0.775843in}}%
\pgfpathlineto{\pgfqpoint{7.499255in}{0.720617in}}%
\pgfpathlineto{\pgfqpoint{7.499255in}{0.720617in}}%
\pgfpathclose%
\pgfpathmoveto{\pgfqpoint{7.519537in}{0.737675in}}%
\pgfpathlineto{\pgfqpoint{7.519537in}{0.775843in}}%
\pgfpathlineto{\pgfqpoint{7.529894in}{0.775843in}}%
\pgfpathlineto{\pgfqpoint{7.529894in}{0.738071in}}%
\pgfpathquadraticcurveto{\pgfqpoint{7.529894in}{0.729119in}}{\pgfqpoint{7.533370in}{0.724634in}}%
\pgfpathquadraticcurveto{\pgfqpoint{7.536865in}{0.720167in}}{\pgfqpoint{7.543853in}{0.720167in}}%
\pgfpathquadraticcurveto{\pgfqpoint{7.552229in}{0.720167in}}{\pgfqpoint{7.557092in}{0.725517in}}%
\pgfpathquadraticcurveto{\pgfqpoint{7.561974in}{0.730866in}}{\pgfqpoint{7.561974in}{0.740106in}}%
\pgfpathlineto{\pgfqpoint{7.561974in}{0.775843in}}%
\pgfpathlineto{\pgfqpoint{7.572331in}{0.775843in}}%
\pgfpathlineto{\pgfqpoint{7.572331in}{0.712800in}}%
\pgfpathlineto{\pgfqpoint{7.561974in}{0.712800in}}%
\pgfpathlineto{\pgfqpoint{7.561974in}{0.722491in}}%
\pgfpathquadraticcurveto{\pgfqpoint{7.558209in}{0.716745in}}{\pgfqpoint{7.553220in}{0.713953in}}%
\pgfpathquadraticcurveto{\pgfqpoint{7.548248in}{0.711161in}}{\pgfqpoint{7.541656in}{0.711161in}}%
\pgfpathquadraticcurveto{\pgfqpoint{7.530795in}{0.711161in}}{\pgfqpoint{7.525157in}{0.717915in}}%
\pgfpathquadraticcurveto{\pgfqpoint{7.519537in}{0.724670in}}{\pgfqpoint{7.519537in}{0.737675in}}%
\pgfpathlineto{\pgfqpoint{7.519537in}{0.737675in}}%
\pgfpathclose%
\pgfpathmoveto{\pgfqpoint{7.545601in}{0.777356in}}%
\pgfpathlineto{\pgfqpoint{7.545601in}{0.777356in}}%
\pgfpathlineto{\pgfqpoint{7.545601in}{0.777356in}}%
\pgfpathclose%
\pgfpathmoveto{\pgfqpoint{7.622314in}{0.744483in}}%
\pgfpathquadraticcurveto{\pgfqpoint{7.609760in}{0.744483in}}{\pgfqpoint{7.604915in}{0.741619in}}%
\pgfpathquadraticcurveto{\pgfqpoint{7.600069in}{0.738756in}}{\pgfqpoint{7.600069in}{0.731821in}}%
\pgfpathquadraticcurveto{\pgfqpoint{7.600069in}{0.726309in}}{\pgfqpoint{7.603708in}{0.723067in}}%
\pgfpathquadraticcurveto{\pgfqpoint{7.607346in}{0.719843in}}{\pgfqpoint{7.613578in}{0.719843in}}%
\pgfpathquadraticcurveto{\pgfqpoint{7.622206in}{0.719843in}}{\pgfqpoint{7.627412in}{0.725949in}}%
\pgfpathquadraticcurveto{\pgfqpoint{7.632617in}{0.732055in}}{\pgfqpoint{7.632617in}{0.742178in}}%
\pgfpathlineto{\pgfqpoint{7.632617in}{0.744483in}}%
\pgfpathlineto{\pgfqpoint{7.622314in}{0.744483in}}%
\pgfpathlineto{\pgfqpoint{7.622314in}{0.744483in}}%
\pgfpathclose%
\pgfpathmoveto{\pgfqpoint{7.642974in}{0.748770in}}%
\pgfpathlineto{\pgfqpoint{7.642974in}{0.712800in}}%
\pgfpathlineto{\pgfqpoint{7.632617in}{0.712800in}}%
\pgfpathlineto{\pgfqpoint{7.632617in}{0.722364in}}%
\pgfpathquadraticcurveto{\pgfqpoint{7.629069in}{0.716637in}}{\pgfqpoint{7.623773in}{0.713899in}}%
\pgfpathquadraticcurveto{\pgfqpoint{7.618478in}{0.711161in}}{\pgfqpoint{7.610823in}{0.711161in}}%
\pgfpathquadraticcurveto{\pgfqpoint{7.601150in}{0.711161in}}{\pgfqpoint{7.595422in}{0.716601in}}%
\pgfpathquadraticcurveto{\pgfqpoint{7.589712in}{0.722040in}}{\pgfqpoint{7.589712in}{0.731154in}}%
\pgfpathquadraticcurveto{\pgfqpoint{7.589712in}{0.741782in}}{\pgfqpoint{7.596827in}{0.747185in}}%
\pgfpathquadraticcurveto{\pgfqpoint{7.603960in}{0.752589in}}{\pgfqpoint{7.618081in}{0.752589in}}%
\pgfpathlineto{\pgfqpoint{7.632617in}{0.752589in}}%
\pgfpathlineto{\pgfqpoint{7.632617in}{0.753616in}}%
\pgfpathquadraticcurveto{\pgfqpoint{7.632617in}{0.760766in}}{\pgfqpoint{7.627916in}{0.764675in}}%
\pgfpathquadraticcurveto{\pgfqpoint{7.623215in}{0.768584in}}{\pgfqpoint{7.614713in}{0.768584in}}%
\pgfpathquadraticcurveto{\pgfqpoint{7.609310in}{0.768584in}}{\pgfqpoint{7.604176in}{0.767287in}}%
\pgfpathquadraticcurveto{\pgfqpoint{7.599061in}{0.765990in}}{\pgfqpoint{7.594341in}{0.763396in}}%
\pgfpathlineto{\pgfqpoint{7.594341in}{0.772979in}}%
\pgfpathquadraticcurveto{\pgfqpoint{7.600015in}{0.775176in}}{\pgfqpoint{7.605365in}{0.776257in}}%
\pgfpathquadraticcurveto{\pgfqpoint{7.610714in}{0.777356in}}{\pgfqpoint{7.615776in}{0.777356in}}%
\pgfpathquadraticcurveto{\pgfqpoint{7.629465in}{0.777356in}}{\pgfqpoint{7.636220in}{0.770259in}}%
\pgfpathquadraticcurveto{\pgfqpoint{7.642974in}{0.763180in}}{\pgfqpoint{7.642974in}{0.748770in}}%
\pgfpathlineto{\pgfqpoint{7.642974in}{0.748770in}}%
\pgfpathclose%
\pgfpathmoveto{\pgfqpoint{7.700829in}{0.766170in}}%
\pgfpathquadraticcurveto{\pgfqpoint{7.699082in}{0.767179in}}{\pgfqpoint{7.697029in}{0.767647in}}%
\pgfpathquadraticcurveto{\pgfqpoint{7.694975in}{0.768133in}}{\pgfqpoint{7.692508in}{0.768133in}}%
\pgfpathquadraticcurveto{\pgfqpoint{7.683718in}{0.768133in}}{\pgfqpoint{7.679017in}{0.762423in}}%
\pgfpathquadraticcurveto{\pgfqpoint{7.674315in}{0.756714in}}{\pgfqpoint{7.674315in}{0.746014in}}%
\pgfpathlineto{\pgfqpoint{7.674315in}{0.712800in}}%
\pgfpathlineto{\pgfqpoint{7.663904in}{0.712800in}}%
\pgfpathlineto{\pgfqpoint{7.663904in}{0.775843in}}%
\pgfpathlineto{\pgfqpoint{7.674315in}{0.775843in}}%
\pgfpathlineto{\pgfqpoint{7.674315in}{0.766044in}}%
\pgfpathquadraticcurveto{\pgfqpoint{7.677594in}{0.771790in}}{\pgfqpoint{7.682817in}{0.774564in}}%
\pgfpathquadraticcurveto{\pgfqpoint{7.688059in}{0.777356in}}{\pgfqpoint{7.695552in}{0.777356in}}%
\pgfpathquadraticcurveto{\pgfqpoint{7.696614in}{0.777356in}}{\pgfqpoint{7.697911in}{0.777211in}}%
\pgfpathquadraticcurveto{\pgfqpoint{7.699208in}{0.777085in}}{\pgfqpoint{7.700775in}{0.776797in}}%
\pgfpathlineto{\pgfqpoint{7.700829in}{0.766170in}}%
\pgfpathlineto{\pgfqpoint{7.700829in}{0.766170in}}%
\pgfpathclose%
\pgfpathmoveto{\pgfqpoint{7.740348in}{0.744483in}}%
\pgfpathquadraticcurveto{\pgfqpoint{7.727793in}{0.744483in}}{\pgfqpoint{7.722948in}{0.741619in}}%
\pgfpathquadraticcurveto{\pgfqpoint{7.718103in}{0.738756in}}{\pgfqpoint{7.718103in}{0.731821in}}%
\pgfpathquadraticcurveto{\pgfqpoint{7.718103in}{0.726309in}}{\pgfqpoint{7.721741in}{0.723067in}}%
\pgfpathquadraticcurveto{\pgfqpoint{7.725380in}{0.719843in}}{\pgfqpoint{7.731612in}{0.719843in}}%
\pgfpathquadraticcurveto{\pgfqpoint{7.740240in}{0.719843in}}{\pgfqpoint{7.745445in}{0.725949in}}%
\pgfpathquadraticcurveto{\pgfqpoint{7.750651in}{0.732055in}}{\pgfqpoint{7.750651in}{0.742178in}}%
\pgfpathlineto{\pgfqpoint{7.750651in}{0.744483in}}%
\pgfpathlineto{\pgfqpoint{7.740348in}{0.744483in}}%
\pgfpathlineto{\pgfqpoint{7.740348in}{0.744483in}}%
\pgfpathclose%
\pgfpathmoveto{\pgfqpoint{7.761008in}{0.748770in}}%
\pgfpathlineto{\pgfqpoint{7.761008in}{0.712800in}}%
\pgfpathlineto{\pgfqpoint{7.750651in}{0.712800in}}%
\pgfpathlineto{\pgfqpoint{7.750651in}{0.722364in}}%
\pgfpathquadraticcurveto{\pgfqpoint{7.747102in}{0.716637in}}{\pgfqpoint{7.741807in}{0.713899in}}%
\pgfpathquadraticcurveto{\pgfqpoint{7.736511in}{0.711161in}}{\pgfqpoint{7.728856in}{0.711161in}}%
\pgfpathquadraticcurveto{\pgfqpoint{7.719184in}{0.711161in}}{\pgfqpoint{7.713456in}{0.716601in}}%
\pgfpathquadraticcurveto{\pgfqpoint{7.707746in}{0.722040in}}{\pgfqpoint{7.707746in}{0.731154in}}%
\pgfpathquadraticcurveto{\pgfqpoint{7.707746in}{0.741782in}}{\pgfqpoint{7.714861in}{0.747185in}}%
\pgfpathquadraticcurveto{\pgfqpoint{7.721994in}{0.752589in}}{\pgfqpoint{7.736115in}{0.752589in}}%
\pgfpathlineto{\pgfqpoint{7.750651in}{0.752589in}}%
\pgfpathlineto{\pgfqpoint{7.750651in}{0.753616in}}%
\pgfpathquadraticcurveto{\pgfqpoint{7.750651in}{0.760766in}}{\pgfqpoint{7.745950in}{0.764675in}}%
\pgfpathquadraticcurveto{\pgfqpoint{7.741249in}{0.768584in}}{\pgfqpoint{7.732747in}{0.768584in}}%
\pgfpathquadraticcurveto{\pgfqpoint{7.727343in}{0.768584in}}{\pgfqpoint{7.722210in}{0.767287in}}%
\pgfpathquadraticcurveto{\pgfqpoint{7.717094in}{0.765990in}}{\pgfqpoint{7.712375in}{0.763396in}}%
\pgfpathlineto{\pgfqpoint{7.712375in}{0.772979in}}%
\pgfpathquadraticcurveto{\pgfqpoint{7.718049in}{0.775176in}}{\pgfqpoint{7.723398in}{0.776257in}}%
\pgfpathquadraticcurveto{\pgfqpoint{7.728748in}{0.777356in}}{\pgfqpoint{7.733810in}{0.777356in}}%
\pgfpathquadraticcurveto{\pgfqpoint{7.747499in}{0.777356in}}{\pgfqpoint{7.754253in}{0.770259in}}%
\pgfpathquadraticcurveto{\pgfqpoint{7.761008in}{0.763180in}}{\pgfqpoint{7.761008in}{0.748770in}}%
\pgfpathlineto{\pgfqpoint{7.761008in}{0.748770in}}%
\pgfpathclose%
\pgfpathmoveto{\pgfqpoint{7.834750in}{0.750860in}}%
\pgfpathlineto{\pgfqpoint{7.834750in}{0.712800in}}%
\pgfpathlineto{\pgfqpoint{7.824393in}{0.712800in}}%
\pgfpathlineto{\pgfqpoint{7.824393in}{0.750517in}}%
\pgfpathquadraticcurveto{\pgfqpoint{7.824393in}{0.759469in}}{\pgfqpoint{7.820898in}{0.763900in}}%
\pgfpathquadraticcurveto{\pgfqpoint{7.817404in}{0.768349in}}{\pgfqpoint{7.810433in}{0.768349in}}%
\pgfpathquadraticcurveto{\pgfqpoint{7.802040in}{0.768349in}}{\pgfqpoint{7.797194in}{0.763000in}}%
\pgfpathquadraticcurveto{\pgfqpoint{7.792349in}{0.757668in}}{\pgfqpoint{7.792349in}{0.748428in}}%
\pgfpathlineto{\pgfqpoint{7.792349in}{0.712800in}}%
\pgfpathlineto{\pgfqpoint{7.781938in}{0.712800in}}%
\pgfpathlineto{\pgfqpoint{7.781938in}{0.775843in}}%
\pgfpathlineto{\pgfqpoint{7.792349in}{0.775843in}}%
\pgfpathlineto{\pgfqpoint{7.792349in}{0.766044in}}%
\pgfpathquadraticcurveto{\pgfqpoint{7.796078in}{0.771736in}}{\pgfqpoint{7.801103in}{0.774546in}}%
\pgfpathquadraticcurveto{\pgfqpoint{7.806146in}{0.777356in}}{\pgfqpoint{7.812739in}{0.777356in}}%
\pgfpathquadraticcurveto{\pgfqpoint{7.823600in}{0.777356in}}{\pgfqpoint{7.829166in}{0.770637in}}%
\pgfpathquadraticcurveto{\pgfqpoint{7.834750in}{0.763918in}}{\pgfqpoint{7.834750in}{0.750860in}}%
\pgfpathlineto{\pgfqpoint{7.834750in}{0.750860in}}%
\pgfpathclose%
\pgfpathmoveto{\pgfqpoint{7.865640in}{0.793747in}}%
\pgfpathlineto{\pgfqpoint{7.865640in}{0.775843in}}%
\pgfpathlineto{\pgfqpoint{7.886967in}{0.775843in}}%
\pgfpathlineto{\pgfqpoint{7.886967in}{0.767791in}}%
\pgfpathlineto{\pgfqpoint{7.865640in}{0.767791in}}%
\pgfpathlineto{\pgfqpoint{7.865640in}{0.733568in}}%
\pgfpathquadraticcurveto{\pgfqpoint{7.865640in}{0.725859in}}{\pgfqpoint{7.867748in}{0.723661in}}%
\pgfpathquadraticcurveto{\pgfqpoint{7.869855in}{0.721464in}}{\pgfqpoint{7.876340in}{0.721464in}}%
\pgfpathlineto{\pgfqpoint{7.886967in}{0.721464in}}%
\pgfpathlineto{\pgfqpoint{7.886967in}{0.712800in}}%
\pgfpathlineto{\pgfqpoint{7.876340in}{0.712800in}}%
\pgfpathquadraticcurveto{\pgfqpoint{7.864344in}{0.712800in}}{\pgfqpoint{7.859787in}{0.717267in}}%
\pgfpathquadraticcurveto{\pgfqpoint{7.855229in}{0.721752in}}{\pgfqpoint{7.855229in}{0.733568in}}%
\pgfpathlineto{\pgfqpoint{7.855229in}{0.767791in}}%
\pgfpathlineto{\pgfqpoint{7.847628in}{0.767791in}}%
\pgfpathlineto{\pgfqpoint{7.847628in}{0.775843in}}%
\pgfpathlineto{\pgfqpoint{7.855229in}{0.775843in}}%
\pgfpathlineto{\pgfqpoint{7.855229in}{0.793747in}}%
\pgfpathlineto{\pgfqpoint{7.865640in}{0.793747in}}%
\pgfpathlineto{\pgfqpoint{7.865640in}{0.793747in}}%
\pgfpathclose%
\pgfpathmoveto{\pgfqpoint{7.954512in}{0.746915in}}%
\pgfpathlineto{\pgfqpoint{7.954512in}{0.741854in}}%
\pgfpathlineto{\pgfqpoint{7.906888in}{0.741854in}}%
\pgfpathquadraticcurveto{\pgfqpoint{7.907573in}{0.731154in}}{\pgfqpoint{7.913337in}{0.725553in}}%
\pgfpathquadraticcurveto{\pgfqpoint{7.919101in}{0.719951in}}{\pgfqpoint{7.929403in}{0.719951in}}%
\pgfpathquadraticcurveto{\pgfqpoint{7.935366in}{0.719951in}}{\pgfqpoint{7.940967in}{0.721410in}}%
\pgfpathquadraticcurveto{\pgfqpoint{7.946569in}{0.722869in}}{\pgfqpoint{7.952099in}{0.725805in}}%
\pgfpathlineto{\pgfqpoint{7.952099in}{0.716006in}}%
\pgfpathquadraticcurveto{\pgfqpoint{7.946515in}{0.713647in}}{\pgfqpoint{7.940661in}{0.712404in}}%
\pgfpathquadraticcurveto{\pgfqpoint{7.934807in}{0.711161in}}{\pgfqpoint{7.928791in}{0.711161in}}%
\pgfpathquadraticcurveto{\pgfqpoint{7.913697in}{0.711161in}}{\pgfqpoint{7.904889in}{0.719933in}}%
\pgfpathquadraticcurveto{\pgfqpoint{7.896081in}{0.728723in}}{\pgfqpoint{7.896081in}{0.743709in}}%
\pgfpathquadraticcurveto{\pgfqpoint{7.896081in}{0.759181in}}{\pgfqpoint{7.904439in}{0.768259in}}%
\pgfpathquadraticcurveto{\pgfqpoint{7.912796in}{0.777356in}}{\pgfqpoint{7.926990in}{0.777356in}}%
\pgfpathquadraticcurveto{\pgfqpoint{7.939706in}{0.777356in}}{\pgfqpoint{7.947109in}{0.769160in}}%
\pgfpathquadraticcurveto{\pgfqpoint{7.954512in}{0.760983in}}{\pgfqpoint{7.954512in}{0.746915in}}%
\pgfpathlineto{\pgfqpoint{7.954512in}{0.746915in}}%
\pgfpathclose%
\pgfpathmoveto{\pgfqpoint{7.944155in}{0.749959in}}%
\pgfpathquadraticcurveto{\pgfqpoint{7.944047in}{0.758443in}}{\pgfqpoint{7.939400in}{0.763504in}}%
\pgfpathquadraticcurveto{\pgfqpoint{7.934753in}{0.768584in}}{\pgfqpoint{7.927098in}{0.768584in}}%
\pgfpathquadraticcurveto{\pgfqpoint{7.918434in}{0.768584in}}{\pgfqpoint{7.913229in}{0.763684in}}%
\pgfpathquadraticcurveto{\pgfqpoint{7.908023in}{0.758785in}}{\pgfqpoint{7.907231in}{0.749887in}}%
\pgfpathlineto{\pgfqpoint{7.944155in}{0.749959in}}%
\pgfpathlineto{\pgfqpoint{7.944155in}{0.749959in}}%
\pgfpathclose%
\pgfpathmoveto{\pgfqpoint{8.025444in}{0.746915in}}%
\pgfpathlineto{\pgfqpoint{8.025444in}{0.741854in}}%
\pgfpathlineto{\pgfqpoint{7.977820in}{0.741854in}}%
\pgfpathquadraticcurveto{\pgfqpoint{7.978505in}{0.731154in}}{\pgfqpoint{7.984268in}{0.725553in}}%
\pgfpathquadraticcurveto{\pgfqpoint{7.990032in}{0.719951in}}{\pgfqpoint{8.000335in}{0.719951in}}%
\pgfpathquadraticcurveto{\pgfqpoint{8.006297in}{0.719951in}}{\pgfqpoint{8.011899in}{0.721410in}}%
\pgfpathquadraticcurveto{\pgfqpoint{8.017501in}{0.722869in}}{\pgfqpoint{8.023031in}{0.725805in}}%
\pgfpathlineto{\pgfqpoint{8.023031in}{0.716006in}}%
\pgfpathquadraticcurveto{\pgfqpoint{8.017447in}{0.713647in}}{\pgfqpoint{8.011593in}{0.712404in}}%
\pgfpathquadraticcurveto{\pgfqpoint{8.005739in}{0.711161in}}{\pgfqpoint{7.999723in}{0.711161in}}%
\pgfpathquadraticcurveto{\pgfqpoint{7.984629in}{0.711161in}}{\pgfqpoint{7.975821in}{0.719933in}}%
\pgfpathquadraticcurveto{\pgfqpoint{7.967013in}{0.728723in}}{\pgfqpoint{7.967013in}{0.743709in}}%
\pgfpathquadraticcurveto{\pgfqpoint{7.967013in}{0.759181in}}{\pgfqpoint{7.975370in}{0.768259in}}%
\pgfpathquadraticcurveto{\pgfqpoint{7.983728in}{0.777356in}}{\pgfqpoint{7.997922in}{0.777356in}}%
\pgfpathquadraticcurveto{\pgfqpoint{8.010638in}{0.777356in}}{\pgfqpoint{8.018041in}{0.769160in}}%
\pgfpathquadraticcurveto{\pgfqpoint{8.025444in}{0.760983in}}{\pgfqpoint{8.025444in}{0.746915in}}%
\pgfpathlineto{\pgfqpoint{8.025444in}{0.746915in}}%
\pgfpathclose%
\pgfpathmoveto{\pgfqpoint{8.015087in}{0.749959in}}%
\pgfpathquadraticcurveto{\pgfqpoint{8.014979in}{0.758443in}}{\pgfqpoint{8.010332in}{0.763504in}}%
\pgfpathquadraticcurveto{\pgfqpoint{8.005685in}{0.768584in}}{\pgfqpoint{7.998030in}{0.768584in}}%
\pgfpathquadraticcurveto{\pgfqpoint{7.989366in}{0.768584in}}{\pgfqpoint{7.984160in}{0.763684in}}%
\pgfpathquadraticcurveto{\pgfqpoint{7.978955in}{0.758785in}}{\pgfqpoint{7.978162in}{0.749887in}}%
\pgfpathlineto{\pgfqpoint{8.015087in}{0.749959in}}%
\pgfpathlineto{\pgfqpoint{8.015087in}{0.749959in}}%
\pgfpathclose%
\pgfpathmoveto{\pgfqpoint{8.043907in}{0.727102in}}%
\pgfpathlineto{\pgfqpoint{8.055795in}{0.727102in}}%
\pgfpathlineto{\pgfqpoint{8.055795in}{0.712800in}}%
\pgfpathlineto{\pgfqpoint{8.043907in}{0.712800in}}%
\pgfpathlineto{\pgfqpoint{8.043907in}{0.727102in}}%
\pgfpathlineto{\pgfqpoint{8.043907in}{0.727102in}}%
\pgfpathclose%
\pgfusepath{fill}%
\end{pgfscope}%
\begin{pgfscope}%
\pgfsetbuttcap%
\pgfsetroundjoin%
\definecolor{currentfill}{rgb}{0.309804,0.372549,0.419608}%
\pgfsetfillcolor{currentfill}%
\pgfsetlinewidth{0.000000pt}%
\definecolor{currentstroke}{rgb}{0.309804,0.372549,0.419608}%
\pgfsetstrokecolor{currentstroke}%
\pgfsetdash{}{0pt}%
\pgfpathmoveto{\pgfqpoint{8.137906in}{0.796845in}}%
\pgfpathlineto{\pgfqpoint{8.209000in}{0.796845in}}%
\pgfpathlineto{\pgfqpoint{8.209000in}{0.787262in}}%
\pgfpathlineto{\pgfqpoint{8.179171in}{0.787262in}}%
\pgfpathlineto{\pgfqpoint{8.179171in}{0.712800in}}%
\pgfpathlineto{\pgfqpoint{8.167752in}{0.712800in}}%
\pgfpathlineto{\pgfqpoint{8.167752in}{0.787262in}}%
\pgfpathlineto{\pgfqpoint{8.137906in}{0.787262in}}%
\pgfpathlineto{\pgfqpoint{8.137906in}{0.796845in}}%
\pgfpathlineto{\pgfqpoint{8.137906in}{0.796845in}}%
\pgfpathclose%
\pgfpathmoveto{\pgfqpoint{8.212494in}{0.796845in}}%
\pgfpathlineto{\pgfqpoint{8.223968in}{0.796845in}}%
\pgfpathlineto{\pgfqpoint{8.241638in}{0.725805in}}%
\pgfpathlineto{\pgfqpoint{8.259253in}{0.796845in}}%
\pgfpathlineto{\pgfqpoint{8.272042in}{0.796845in}}%
\pgfpathlineto{\pgfqpoint{8.289712in}{0.725805in}}%
\pgfpathlineto{\pgfqpoint{8.307328in}{0.796845in}}%
\pgfpathlineto{\pgfqpoint{8.318874in}{0.796845in}}%
\pgfpathlineto{\pgfqpoint{8.297763in}{0.712800in}}%
\pgfpathlineto{\pgfqpoint{8.283462in}{0.712800in}}%
\pgfpathlineto{\pgfqpoint{8.265738in}{0.785749in}}%
\pgfpathlineto{\pgfqpoint{8.247834in}{0.712800in}}%
\pgfpathlineto{\pgfqpoint{8.233532in}{0.712800in}}%
\pgfpathlineto{\pgfqpoint{8.212494in}{0.796845in}}%
\pgfpathlineto{\pgfqpoint{8.212494in}{0.796845in}}%
\pgfpathclose%
\pgfpathmoveto{\pgfqpoint{8.333950in}{0.796845in}}%
\pgfpathlineto{\pgfqpoint{8.382240in}{0.796845in}}%
\pgfpathlineto{\pgfqpoint{8.382240in}{0.787262in}}%
\pgfpathlineto{\pgfqpoint{8.345316in}{0.787262in}}%
\pgfpathlineto{\pgfqpoint{8.345316in}{0.762496in}}%
\pgfpathlineto{\pgfqpoint{8.378638in}{0.762496in}}%
\pgfpathlineto{\pgfqpoint{8.378638in}{0.752931in}}%
\pgfpathlineto{\pgfqpoint{8.345316in}{0.752931in}}%
\pgfpathlineto{\pgfqpoint{8.345316in}{0.712800in}}%
\pgfpathlineto{\pgfqpoint{8.333950in}{0.712800in}}%
\pgfpathlineto{\pgfqpoint{8.333950in}{0.796845in}}%
\pgfpathlineto{\pgfqpoint{8.333950in}{0.796845in}}%
\pgfpathclose%
\pgfpathmoveto{\pgfqpoint{8.400253in}{0.796845in}}%
\pgfpathlineto{\pgfqpoint{8.453388in}{0.796845in}}%
\pgfpathlineto{\pgfqpoint{8.453388in}{0.787262in}}%
\pgfpathlineto{\pgfqpoint{8.411618in}{0.787262in}}%
\pgfpathlineto{\pgfqpoint{8.411618in}{0.762387in}}%
\pgfpathlineto{\pgfqpoint{8.451641in}{0.762387in}}%
\pgfpathlineto{\pgfqpoint{8.451641in}{0.752823in}}%
\pgfpathlineto{\pgfqpoint{8.411618in}{0.752823in}}%
\pgfpathlineto{\pgfqpoint{8.411618in}{0.722364in}}%
\pgfpathlineto{\pgfqpoint{8.454397in}{0.722364in}}%
\pgfpathlineto{\pgfqpoint{8.454397in}{0.712800in}}%
\pgfpathlineto{\pgfqpoint{8.400253in}{0.712800in}}%
\pgfpathlineto{\pgfqpoint{8.400253in}{0.796845in}}%
\pgfpathlineto{\pgfqpoint{8.400253in}{0.796845in}}%
\pgfpathclose%
\pgfusepath{fill}%
\end{pgfscope}%
\begin{pgfscope}%
\pgfsetbuttcap%
\pgfsetroundjoin%
\definecolor{currentfill}{rgb}{0.309804,0.372549,0.419608}%
\pgfsetfillcolor{currentfill}%
\pgfsetlinewidth{0.000000pt}%
\definecolor{currentstroke}{rgb}{0.309804,0.372549,0.419608}%
\pgfsetstrokecolor{currentstroke}%
\pgfsetdash{}{0pt}%
\pgfpathmoveto{\pgfqpoint{8.567331in}{0.766170in}}%
\pgfpathquadraticcurveto{\pgfqpoint{8.565584in}{0.767179in}}{\pgfqpoint{8.563530in}{0.767647in}}%
\pgfpathquadraticcurveto{\pgfqpoint{8.561477in}{0.768133in}}{\pgfqpoint{8.559009in}{0.768133in}}%
\pgfpathquadraticcurveto{\pgfqpoint{8.550219in}{0.768133in}}{\pgfqpoint{8.545518in}{0.762423in}}%
\pgfpathquadraticcurveto{\pgfqpoint{8.540817in}{0.756714in}}{\pgfqpoint{8.540817in}{0.746014in}}%
\pgfpathlineto{\pgfqpoint{8.540817in}{0.712800in}}%
\pgfpathlineto{\pgfqpoint{8.530406in}{0.712800in}}%
\pgfpathlineto{\pgfqpoint{8.530406in}{0.775843in}}%
\pgfpathlineto{\pgfqpoint{8.540817in}{0.775843in}}%
\pgfpathlineto{\pgfqpoint{8.540817in}{0.766044in}}%
\pgfpathquadraticcurveto{\pgfqpoint{8.544095in}{0.771790in}}{\pgfqpoint{8.549319in}{0.774564in}}%
\pgfpathquadraticcurveto{\pgfqpoint{8.554560in}{0.777356in}}{\pgfqpoint{8.562053in}{0.777356in}}%
\pgfpathquadraticcurveto{\pgfqpoint{8.563116in}{0.777356in}}{\pgfqpoint{8.564413in}{0.777211in}}%
\pgfpathquadraticcurveto{\pgfqpoint{8.565710in}{0.777085in}}{\pgfqpoint{8.567277in}{0.776797in}}%
\pgfpathlineto{\pgfqpoint{8.567331in}{0.766170in}}%
\pgfpathlineto{\pgfqpoint{8.567331in}{0.766170in}}%
\pgfpathclose%
\pgfpathmoveto{\pgfqpoint{8.629581in}{0.746915in}}%
\pgfpathlineto{\pgfqpoint{8.629581in}{0.741854in}}%
\pgfpathlineto{\pgfqpoint{8.581957in}{0.741854in}}%
\pgfpathquadraticcurveto{\pgfqpoint{8.582641in}{0.731154in}}{\pgfqpoint{8.588405in}{0.725553in}}%
\pgfpathquadraticcurveto{\pgfqpoint{8.594169in}{0.719951in}}{\pgfqpoint{8.604472in}{0.719951in}}%
\pgfpathquadraticcurveto{\pgfqpoint{8.610434in}{0.719951in}}{\pgfqpoint{8.616036in}{0.721410in}}%
\pgfpathquadraticcurveto{\pgfqpoint{8.621637in}{0.722869in}}{\pgfqpoint{8.627167in}{0.725805in}}%
\pgfpathlineto{\pgfqpoint{8.627167in}{0.716006in}}%
\pgfpathquadraticcurveto{\pgfqpoint{8.621583in}{0.713647in}}{\pgfqpoint{8.615729in}{0.712404in}}%
\pgfpathquadraticcurveto{\pgfqpoint{8.609875in}{0.711161in}}{\pgfqpoint{8.603859in}{0.711161in}}%
\pgfpathquadraticcurveto{\pgfqpoint{8.588765in}{0.711161in}}{\pgfqpoint{8.579957in}{0.719933in}}%
\pgfpathquadraticcurveto{\pgfqpoint{8.571149in}{0.728723in}}{\pgfqpoint{8.571149in}{0.743709in}}%
\pgfpathquadraticcurveto{\pgfqpoint{8.571149in}{0.759181in}}{\pgfqpoint{8.579507in}{0.768259in}}%
\pgfpathquadraticcurveto{\pgfqpoint{8.587865in}{0.777356in}}{\pgfqpoint{8.602058in}{0.777356in}}%
\pgfpathquadraticcurveto{\pgfqpoint{8.614775in}{0.777356in}}{\pgfqpoint{8.622178in}{0.769160in}}%
\pgfpathquadraticcurveto{\pgfqpoint{8.629581in}{0.760983in}}{\pgfqpoint{8.629581in}{0.746915in}}%
\pgfpathlineto{\pgfqpoint{8.629581in}{0.746915in}}%
\pgfpathclose%
\pgfpathmoveto{\pgfqpoint{8.619224in}{0.749959in}}%
\pgfpathquadraticcurveto{\pgfqpoint{8.619116in}{0.758443in}}{\pgfqpoint{8.614469in}{0.763504in}}%
\pgfpathquadraticcurveto{\pgfqpoint{8.609821in}{0.768584in}}{\pgfqpoint{8.602166in}{0.768584in}}%
\pgfpathquadraticcurveto{\pgfqpoint{8.593502in}{0.768584in}}{\pgfqpoint{8.588297in}{0.763684in}}%
\pgfpathquadraticcurveto{\pgfqpoint{8.583091in}{0.758785in}}{\pgfqpoint{8.582299in}{0.749887in}}%
\pgfpathlineto{\pgfqpoint{8.619224in}{0.749959in}}%
\pgfpathlineto{\pgfqpoint{8.619224in}{0.749959in}}%
\pgfpathclose%
\pgfpathmoveto{\pgfqpoint{8.656833in}{0.793747in}}%
\pgfpathlineto{\pgfqpoint{8.656833in}{0.775843in}}%
\pgfpathlineto{\pgfqpoint{8.678159in}{0.775843in}}%
\pgfpathlineto{\pgfqpoint{8.678159in}{0.767791in}}%
\pgfpathlineto{\pgfqpoint{8.656833in}{0.767791in}}%
\pgfpathlineto{\pgfqpoint{8.656833in}{0.733568in}}%
\pgfpathquadraticcurveto{\pgfqpoint{8.656833in}{0.725859in}}{\pgfqpoint{8.658941in}{0.723661in}}%
\pgfpathquadraticcurveto{\pgfqpoint{8.661048in}{0.721464in}}{\pgfqpoint{8.667532in}{0.721464in}}%
\pgfpathlineto{\pgfqpoint{8.678159in}{0.721464in}}%
\pgfpathlineto{\pgfqpoint{8.678159in}{0.712800in}}%
\pgfpathlineto{\pgfqpoint{8.667532in}{0.712800in}}%
\pgfpathquadraticcurveto{\pgfqpoint{8.655536in}{0.712800in}}{\pgfqpoint{8.650979in}{0.717267in}}%
\pgfpathquadraticcurveto{\pgfqpoint{8.646422in}{0.721752in}}{\pgfqpoint{8.646422in}{0.733568in}}%
\pgfpathlineto{\pgfqpoint{8.646422in}{0.767791in}}%
\pgfpathlineto{\pgfqpoint{8.638821in}{0.767791in}}%
\pgfpathlineto{\pgfqpoint{8.638821in}{0.775843in}}%
\pgfpathlineto{\pgfqpoint{8.646422in}{0.775843in}}%
\pgfpathlineto{\pgfqpoint{8.646422in}{0.793747in}}%
\pgfpathlineto{\pgfqpoint{8.656833in}{0.793747in}}%
\pgfpathlineto{\pgfqpoint{8.656833in}{0.793747in}}%
\pgfpathclose%
\pgfpathmoveto{\pgfqpoint{8.720434in}{0.744483in}}%
\pgfpathquadraticcurveto{\pgfqpoint{8.707880in}{0.744483in}}{\pgfqpoint{8.703034in}{0.741619in}}%
\pgfpathquadraticcurveto{\pgfqpoint{8.698189in}{0.738756in}}{\pgfqpoint{8.698189in}{0.731821in}}%
\pgfpathquadraticcurveto{\pgfqpoint{8.698189in}{0.726309in}}{\pgfqpoint{8.701827in}{0.723067in}}%
\pgfpathquadraticcurveto{\pgfqpoint{8.705466in}{0.719843in}}{\pgfqpoint{8.711698in}{0.719843in}}%
\pgfpathquadraticcurveto{\pgfqpoint{8.720326in}{0.719843in}}{\pgfqpoint{8.725531in}{0.725949in}}%
\pgfpathquadraticcurveto{\pgfqpoint{8.730737in}{0.732055in}}{\pgfqpoint{8.730737in}{0.742178in}}%
\pgfpathlineto{\pgfqpoint{8.730737in}{0.744483in}}%
\pgfpathlineto{\pgfqpoint{8.720434in}{0.744483in}}%
\pgfpathlineto{\pgfqpoint{8.720434in}{0.744483in}}%
\pgfpathclose%
\pgfpathmoveto{\pgfqpoint{8.741094in}{0.748770in}}%
\pgfpathlineto{\pgfqpoint{8.741094in}{0.712800in}}%
\pgfpathlineto{\pgfqpoint{8.730737in}{0.712800in}}%
\pgfpathlineto{\pgfqpoint{8.730737in}{0.722364in}}%
\pgfpathquadraticcurveto{\pgfqpoint{8.727189in}{0.716637in}}{\pgfqpoint{8.721893in}{0.713899in}}%
\pgfpathquadraticcurveto{\pgfqpoint{8.716597in}{0.711161in}}{\pgfqpoint{8.708942in}{0.711161in}}%
\pgfpathquadraticcurveto{\pgfqpoint{8.699270in}{0.711161in}}{\pgfqpoint{8.693542in}{0.716601in}}%
\pgfpathquadraticcurveto{\pgfqpoint{8.687832in}{0.722040in}}{\pgfqpoint{8.687832in}{0.731154in}}%
\pgfpathquadraticcurveto{\pgfqpoint{8.687832in}{0.741782in}}{\pgfqpoint{8.694947in}{0.747185in}}%
\pgfpathquadraticcurveto{\pgfqpoint{8.702080in}{0.752589in}}{\pgfqpoint{8.716201in}{0.752589in}}%
\pgfpathlineto{\pgfqpoint{8.730737in}{0.752589in}}%
\pgfpathlineto{\pgfqpoint{8.730737in}{0.753616in}}%
\pgfpathquadraticcurveto{\pgfqpoint{8.730737in}{0.760766in}}{\pgfqpoint{8.726036in}{0.764675in}}%
\pgfpathquadraticcurveto{\pgfqpoint{8.721335in}{0.768584in}}{\pgfqpoint{8.712833in}{0.768584in}}%
\pgfpathquadraticcurveto{\pgfqpoint{8.707429in}{0.768584in}}{\pgfqpoint{8.702296in}{0.767287in}}%
\pgfpathquadraticcurveto{\pgfqpoint{8.697180in}{0.765990in}}{\pgfqpoint{8.692461in}{0.763396in}}%
\pgfpathlineto{\pgfqpoint{8.692461in}{0.772979in}}%
\pgfpathquadraticcurveto{\pgfqpoint{8.698135in}{0.775176in}}{\pgfqpoint{8.703485in}{0.776257in}}%
\pgfpathquadraticcurveto{\pgfqpoint{8.708834in}{0.777356in}}{\pgfqpoint{8.713896in}{0.777356in}}%
\pgfpathquadraticcurveto{\pgfqpoint{8.727585in}{0.777356in}}{\pgfqpoint{8.734339in}{0.770259in}}%
\pgfpathquadraticcurveto{\pgfqpoint{8.741094in}{0.763180in}}{\pgfqpoint{8.741094in}{0.748770in}}%
\pgfpathlineto{\pgfqpoint{8.741094in}{0.748770in}}%
\pgfpathclose%
\pgfpathmoveto{\pgfqpoint{8.762420in}{0.775843in}}%
\pgfpathlineto{\pgfqpoint{8.772777in}{0.775843in}}%
\pgfpathlineto{\pgfqpoint{8.772777in}{0.712800in}}%
\pgfpathlineto{\pgfqpoint{8.762420in}{0.712800in}}%
\pgfpathlineto{\pgfqpoint{8.762420in}{0.775843in}}%
\pgfpathlineto{\pgfqpoint{8.762420in}{0.775843in}}%
\pgfpathclose%
\pgfpathmoveto{\pgfqpoint{8.762420in}{0.800393in}}%
\pgfpathlineto{\pgfqpoint{8.772777in}{0.800393in}}%
\pgfpathlineto{\pgfqpoint{8.772777in}{0.787262in}}%
\pgfpathlineto{\pgfqpoint{8.762420in}{0.787262in}}%
\pgfpathlineto{\pgfqpoint{8.762420in}{0.800393in}}%
\pgfpathlineto{\pgfqpoint{8.762420in}{0.800393in}}%
\pgfpathclose%
\pgfpathmoveto{\pgfqpoint{8.846861in}{0.750860in}}%
\pgfpathlineto{\pgfqpoint{8.846861in}{0.712800in}}%
\pgfpathlineto{\pgfqpoint{8.836504in}{0.712800in}}%
\pgfpathlineto{\pgfqpoint{8.836504in}{0.750517in}}%
\pgfpathquadraticcurveto{\pgfqpoint{8.836504in}{0.759469in}}{\pgfqpoint{8.833010in}{0.763900in}}%
\pgfpathquadraticcurveto{\pgfqpoint{8.829516in}{0.768349in}}{\pgfqpoint{8.822545in}{0.768349in}}%
\pgfpathquadraticcurveto{\pgfqpoint{8.814151in}{0.768349in}}{\pgfqpoint{8.809306in}{0.763000in}}%
\pgfpathquadraticcurveto{\pgfqpoint{8.804461in}{0.757668in}}{\pgfqpoint{8.804461in}{0.748428in}}%
\pgfpathlineto{\pgfqpoint{8.804461in}{0.712800in}}%
\pgfpathlineto{\pgfqpoint{8.794050in}{0.712800in}}%
\pgfpathlineto{\pgfqpoint{8.794050in}{0.775843in}}%
\pgfpathlineto{\pgfqpoint{8.804461in}{0.775843in}}%
\pgfpathlineto{\pgfqpoint{8.804461in}{0.766044in}}%
\pgfpathquadraticcurveto{\pgfqpoint{8.808189in}{0.771736in}}{\pgfqpoint{8.813215in}{0.774546in}}%
\pgfpathquadraticcurveto{\pgfqpoint{8.818258in}{0.777356in}}{\pgfqpoint{8.824850in}{0.777356in}}%
\pgfpathquadraticcurveto{\pgfqpoint{8.835712in}{0.777356in}}{\pgfqpoint{8.841278in}{0.770637in}}%
\pgfpathquadraticcurveto{\pgfqpoint{8.846861in}{0.763918in}}{\pgfqpoint{8.846861in}{0.750860in}}%
\pgfpathlineto{\pgfqpoint{8.846861in}{0.750860in}}%
\pgfpathclose%
\pgfpathmoveto{\pgfqpoint{8.907688in}{0.773987in}}%
\pgfpathlineto{\pgfqpoint{8.907688in}{0.764189in}}%
\pgfpathquadraticcurveto{\pgfqpoint{8.903311in}{0.766440in}}{\pgfqpoint{8.898574in}{0.767557in}}%
\pgfpathquadraticcurveto{\pgfqpoint{8.893855in}{0.768692in}}{\pgfqpoint{8.888776in}{0.768692in}}%
\pgfpathquadraticcurveto{\pgfqpoint{8.881066in}{0.768692in}}{\pgfqpoint{8.877212in}{0.766332in}}%
\pgfpathquadraticcurveto{\pgfqpoint{8.873357in}{0.763973in}}{\pgfqpoint{8.873357in}{0.759235in}}%
\pgfpathquadraticcurveto{\pgfqpoint{8.873357in}{0.755633in}}{\pgfqpoint{8.876113in}{0.753580in}}%
\pgfpathquadraticcurveto{\pgfqpoint{8.878869in}{0.751526in}}{\pgfqpoint{8.887209in}{0.749671in}}%
\pgfpathlineto{\pgfqpoint{8.890757in}{0.748878in}}%
\pgfpathquadraticcurveto{\pgfqpoint{8.901780in}{0.746519in}}{\pgfqpoint{8.906427in}{0.742214in}}%
\pgfpathquadraticcurveto{\pgfqpoint{8.911075in}{0.737909in}}{\pgfqpoint{8.911075in}{0.730200in}}%
\pgfpathquadraticcurveto{\pgfqpoint{8.911075in}{0.721410in}}{\pgfqpoint{8.904122in}{0.716276in}}%
\pgfpathquadraticcurveto{\pgfqpoint{8.897169in}{0.711161in}}{\pgfqpoint{8.885011in}{0.711161in}}%
\pgfpathquadraticcurveto{\pgfqpoint{8.879950in}{0.711161in}}{\pgfqpoint{8.874456in}{0.712152in}}%
\pgfpathquadraticcurveto{\pgfqpoint{8.868962in}{0.713142in}}{\pgfqpoint{8.862892in}{0.715106in}}%
\pgfpathlineto{\pgfqpoint{8.862892in}{0.725805in}}%
\pgfpathquadraticcurveto{\pgfqpoint{8.868638in}{0.722815in}}{\pgfqpoint{8.874204in}{0.721320in}}%
\pgfpathquadraticcurveto{\pgfqpoint{8.879770in}{0.719843in}}{\pgfqpoint{8.885245in}{0.719843in}}%
\pgfpathquadraticcurveto{\pgfqpoint{8.892558in}{0.719843in}}{\pgfqpoint{8.896485in}{0.722346in}}%
\pgfpathquadraticcurveto{\pgfqpoint{8.900429in}{0.724850in}}{\pgfqpoint{8.900429in}{0.729407in}}%
\pgfpathquadraticcurveto{\pgfqpoint{8.900429in}{0.733622in}}{\pgfqpoint{8.897584in}{0.735874in}}%
\pgfpathquadraticcurveto{\pgfqpoint{8.894756in}{0.738125in}}{\pgfqpoint{8.885119in}{0.740214in}}%
\pgfpathlineto{\pgfqpoint{8.881517in}{0.741061in}}%
\pgfpathquadraticcurveto{\pgfqpoint{8.871898in}{0.743078in}}{\pgfqpoint{8.867611in}{0.747275in}}%
\pgfpathquadraticcurveto{\pgfqpoint{8.863342in}{0.751472in}}{\pgfqpoint{8.863342in}{0.758785in}}%
\pgfpathquadraticcurveto{\pgfqpoint{8.863342in}{0.767683in}}{\pgfqpoint{8.869647in}{0.772510in}}%
\pgfpathquadraticcurveto{\pgfqpoint{8.875951in}{0.777356in}}{\pgfqpoint{8.887551in}{0.777356in}}%
\pgfpathquadraticcurveto{\pgfqpoint{8.893279in}{0.777356in}}{\pgfqpoint{8.898340in}{0.776509in}}%
\pgfpathquadraticcurveto{\pgfqpoint{8.903419in}{0.775680in}}{\pgfqpoint{8.907688in}{0.773987in}}%
\pgfpathlineto{\pgfqpoint{8.907688in}{0.773987in}}%
\pgfpathclose%
\pgfusepath{fill}%
\end{pgfscope}%
\begin{pgfscope}%
\pgfsetbuttcap%
\pgfsetroundjoin%
\definecolor{currentfill}{rgb}{0.309804,0.372549,0.419608}%
\pgfsetfillcolor{currentfill}%
\pgfsetlinewidth{0.000000pt}%
\definecolor{currentstroke}{rgb}{0.309804,0.372549,0.419608}%
\pgfsetstrokecolor{currentstroke}%
\pgfsetdash{}{0pt}%
\pgfpathmoveto{\pgfqpoint{8.979995in}{0.775843in}}%
\pgfpathlineto{\pgfqpoint{8.990352in}{0.775843in}}%
\pgfpathlineto{\pgfqpoint{8.990352in}{0.712800in}}%
\pgfpathlineto{\pgfqpoint{8.979995in}{0.712800in}}%
\pgfpathlineto{\pgfqpoint{8.979995in}{0.775843in}}%
\pgfpathlineto{\pgfqpoint{8.979995in}{0.775843in}}%
\pgfpathclose%
\pgfpathmoveto{\pgfqpoint{8.979995in}{0.800393in}}%
\pgfpathlineto{\pgfqpoint{8.990352in}{0.800393in}}%
\pgfpathlineto{\pgfqpoint{8.990352in}{0.787262in}}%
\pgfpathlineto{\pgfqpoint{8.979995in}{0.787262in}}%
\pgfpathlineto{\pgfqpoint{8.979995in}{0.800393in}}%
\pgfpathlineto{\pgfqpoint{8.979995in}{0.800393in}}%
\pgfpathclose%
\pgfpathmoveto{\pgfqpoint{9.022269in}{0.793747in}}%
\pgfpathlineto{\pgfqpoint{9.022269in}{0.775843in}}%
\pgfpathlineto{\pgfqpoint{9.043596in}{0.775843in}}%
\pgfpathlineto{\pgfqpoint{9.043596in}{0.767791in}}%
\pgfpathlineto{\pgfqpoint{9.022269in}{0.767791in}}%
\pgfpathlineto{\pgfqpoint{9.022269in}{0.733568in}}%
\pgfpathquadraticcurveto{\pgfqpoint{9.022269in}{0.725859in}}{\pgfqpoint{9.024377in}{0.723661in}}%
\pgfpathquadraticcurveto{\pgfqpoint{9.026484in}{0.721464in}}{\pgfqpoint{9.032969in}{0.721464in}}%
\pgfpathlineto{\pgfqpoint{9.043596in}{0.721464in}}%
\pgfpathlineto{\pgfqpoint{9.043596in}{0.712800in}}%
\pgfpathlineto{\pgfqpoint{9.032969in}{0.712800in}}%
\pgfpathquadraticcurveto{\pgfqpoint{9.020973in}{0.712800in}}{\pgfqpoint{9.016416in}{0.717267in}}%
\pgfpathquadraticcurveto{\pgfqpoint{9.011858in}{0.721752in}}{\pgfqpoint{9.011858in}{0.733568in}}%
\pgfpathlineto{\pgfqpoint{9.011858in}{0.767791in}}%
\pgfpathlineto{\pgfqpoint{9.004257in}{0.767791in}}%
\pgfpathlineto{\pgfqpoint{9.004257in}{0.775843in}}%
\pgfpathlineto{\pgfqpoint{9.011858in}{0.775843in}}%
\pgfpathlineto{\pgfqpoint{9.011858in}{0.793747in}}%
\pgfpathlineto{\pgfqpoint{9.022269in}{0.793747in}}%
\pgfpathlineto{\pgfqpoint{9.022269in}{0.793747in}}%
\pgfpathclose%
\pgfpathmoveto{\pgfqpoint{9.097398in}{0.773987in}}%
\pgfpathlineto{\pgfqpoint{9.097398in}{0.764189in}}%
\pgfpathquadraticcurveto{\pgfqpoint{9.093021in}{0.766440in}}{\pgfqpoint{9.088284in}{0.767557in}}%
\pgfpathquadraticcurveto{\pgfqpoint{9.083565in}{0.768692in}}{\pgfqpoint{9.078485in}{0.768692in}}%
\pgfpathquadraticcurveto{\pgfqpoint{9.070776in}{0.768692in}}{\pgfqpoint{9.066922in}{0.766332in}}%
\pgfpathquadraticcurveto{\pgfqpoint{9.063067in}{0.763973in}}{\pgfqpoint{9.063067in}{0.759235in}}%
\pgfpathquadraticcurveto{\pgfqpoint{9.063067in}{0.755633in}}{\pgfqpoint{9.065823in}{0.753580in}}%
\pgfpathquadraticcurveto{\pgfqpoint{9.068579in}{0.751526in}}{\pgfqpoint{9.076918in}{0.749671in}}%
\pgfpathlineto{\pgfqpoint{9.080467in}{0.748878in}}%
\pgfpathquadraticcurveto{\pgfqpoint{9.091490in}{0.746519in}}{\pgfqpoint{9.096137in}{0.742214in}}%
\pgfpathquadraticcurveto{\pgfqpoint{9.100784in}{0.737909in}}{\pgfqpoint{9.100784in}{0.730200in}}%
\pgfpathquadraticcurveto{\pgfqpoint{9.100784in}{0.721410in}}{\pgfqpoint{9.093832in}{0.716276in}}%
\pgfpathquadraticcurveto{\pgfqpoint{9.086879in}{0.711161in}}{\pgfqpoint{9.074721in}{0.711161in}}%
\pgfpathquadraticcurveto{\pgfqpoint{9.069659in}{0.711161in}}{\pgfqpoint{9.064166in}{0.712152in}}%
\pgfpathquadraticcurveto{\pgfqpoint{9.058672in}{0.713142in}}{\pgfqpoint{9.052602in}{0.715106in}}%
\pgfpathlineto{\pgfqpoint{9.052602in}{0.725805in}}%
\pgfpathquadraticcurveto{\pgfqpoint{9.058348in}{0.722815in}}{\pgfqpoint{9.063914in}{0.721320in}}%
\pgfpathquadraticcurveto{\pgfqpoint{9.069479in}{0.719843in}}{\pgfqpoint{9.074955in}{0.719843in}}%
\pgfpathquadraticcurveto{\pgfqpoint{9.082268in}{0.719843in}}{\pgfqpoint{9.086195in}{0.722346in}}%
\pgfpathquadraticcurveto{\pgfqpoint{9.090139in}{0.724850in}}{\pgfqpoint{9.090139in}{0.729407in}}%
\pgfpathquadraticcurveto{\pgfqpoint{9.090139in}{0.733622in}}{\pgfqpoint{9.087293in}{0.735874in}}%
\pgfpathquadraticcurveto{\pgfqpoint{9.084465in}{0.738125in}}{\pgfqpoint{9.074829in}{0.740214in}}%
\pgfpathlineto{\pgfqpoint{9.071227in}{0.741061in}}%
\pgfpathquadraticcurveto{\pgfqpoint{9.061608in}{0.743078in}}{\pgfqpoint{9.057321in}{0.747275in}}%
\pgfpathquadraticcurveto{\pgfqpoint{9.053052in}{0.751472in}}{\pgfqpoint{9.053052in}{0.758785in}}%
\pgfpathquadraticcurveto{\pgfqpoint{9.053052in}{0.767683in}}{\pgfqpoint{9.059356in}{0.772510in}}%
\pgfpathquadraticcurveto{\pgfqpoint{9.065661in}{0.777356in}}{\pgfqpoint{9.077261in}{0.777356in}}%
\pgfpathquadraticcurveto{\pgfqpoint{9.082988in}{0.777356in}}{\pgfqpoint{9.088050in}{0.776509in}}%
\pgfpathquadraticcurveto{\pgfqpoint{9.093129in}{0.775680in}}{\pgfqpoint{9.097398in}{0.773987in}}%
\pgfpathlineto{\pgfqpoint{9.097398in}{0.773987in}}%
\pgfpathclose%
\pgfusepath{fill}%
\end{pgfscope}%
\begin{pgfscope}%
\pgfsetbuttcap%
\pgfsetroundjoin%
\definecolor{currentfill}{rgb}{0.309804,0.372549,0.419608}%
\pgfsetfillcolor{currentfill}%
\pgfsetlinewidth{0.000000pt}%
\definecolor{currentstroke}{rgb}{0.309804,0.372549,0.419608}%
\pgfsetstrokecolor{currentstroke}%
\pgfsetdash{}{0pt}%
\pgfpathmoveto{\pgfqpoint{9.215616in}{0.746915in}}%
\pgfpathlineto{\pgfqpoint{9.215616in}{0.741854in}}%
\pgfpathlineto{\pgfqpoint{9.167992in}{0.741854in}}%
\pgfpathquadraticcurveto{\pgfqpoint{9.168677in}{0.731154in}}{\pgfqpoint{9.174440in}{0.725553in}}%
\pgfpathquadraticcurveto{\pgfqpoint{9.180204in}{0.719951in}}{\pgfqpoint{9.190507in}{0.719951in}}%
\pgfpathquadraticcurveto{\pgfqpoint{9.196469in}{0.719951in}}{\pgfqpoint{9.202071in}{0.721410in}}%
\pgfpathquadraticcurveto{\pgfqpoint{9.207673in}{0.722869in}}{\pgfqpoint{9.213203in}{0.725805in}}%
\pgfpathlineto{\pgfqpoint{9.213203in}{0.716006in}}%
\pgfpathquadraticcurveto{\pgfqpoint{9.207619in}{0.713647in}}{\pgfqpoint{9.201765in}{0.712404in}}%
\pgfpathquadraticcurveto{\pgfqpoint{9.195911in}{0.711161in}}{\pgfqpoint{9.189895in}{0.711161in}}%
\pgfpathquadraticcurveto{\pgfqpoint{9.174801in}{0.711161in}}{\pgfqpoint{9.165993in}{0.719933in}}%
\pgfpathquadraticcurveto{\pgfqpoint{9.157185in}{0.728723in}}{\pgfqpoint{9.157185in}{0.743709in}}%
\pgfpathquadraticcurveto{\pgfqpoint{9.157185in}{0.759181in}}{\pgfqpoint{9.165542in}{0.768259in}}%
\pgfpathquadraticcurveto{\pgfqpoint{9.173900in}{0.777356in}}{\pgfqpoint{9.188094in}{0.777356in}}%
\pgfpathquadraticcurveto{\pgfqpoint{9.200810in}{0.777356in}}{\pgfqpoint{9.208213in}{0.769160in}}%
\pgfpathquadraticcurveto{\pgfqpoint{9.215616in}{0.760983in}}{\pgfqpoint{9.215616in}{0.746915in}}%
\pgfpathlineto{\pgfqpoint{9.215616in}{0.746915in}}%
\pgfpathclose%
\pgfpathmoveto{\pgfqpoint{9.205259in}{0.749959in}}%
\pgfpathquadraticcurveto{\pgfqpoint{9.205151in}{0.758443in}}{\pgfqpoint{9.200504in}{0.763504in}}%
\pgfpathquadraticcurveto{\pgfqpoint{9.195857in}{0.768584in}}{\pgfqpoint{9.188202in}{0.768584in}}%
\pgfpathquadraticcurveto{\pgfqpoint{9.179538in}{0.768584in}}{\pgfqpoint{9.174332in}{0.763684in}}%
\pgfpathquadraticcurveto{\pgfqpoint{9.169127in}{0.758785in}}{\pgfqpoint{9.168334in}{0.749887in}}%
\pgfpathlineto{\pgfqpoint{9.205259in}{0.749959in}}%
\pgfpathlineto{\pgfqpoint{9.205259in}{0.749959in}}%
\pgfpathclose%
\pgfpathmoveto{\pgfqpoint{9.232620in}{0.800393in}}%
\pgfpathlineto{\pgfqpoint{9.242977in}{0.800393in}}%
\pgfpathlineto{\pgfqpoint{9.242977in}{0.712800in}}%
\pgfpathlineto{\pgfqpoint{9.232620in}{0.712800in}}%
\pgfpathlineto{\pgfqpoint{9.232620in}{0.800393in}}%
\pgfpathlineto{\pgfqpoint{9.232620in}{0.800393in}}%
\pgfpathclose%
\pgfpathmoveto{\pgfqpoint{9.264645in}{0.775843in}}%
\pgfpathlineto{\pgfqpoint{9.275002in}{0.775843in}}%
\pgfpathlineto{\pgfqpoint{9.275002in}{0.712800in}}%
\pgfpathlineto{\pgfqpoint{9.264645in}{0.712800in}}%
\pgfpathlineto{\pgfqpoint{9.264645in}{0.775843in}}%
\pgfpathlineto{\pgfqpoint{9.264645in}{0.775843in}}%
\pgfpathclose%
\pgfpathmoveto{\pgfqpoint{9.264645in}{0.800393in}}%
\pgfpathlineto{\pgfqpoint{9.275002in}{0.800393in}}%
\pgfpathlineto{\pgfqpoint{9.275002in}{0.787262in}}%
\pgfpathlineto{\pgfqpoint{9.264645in}{0.787262in}}%
\pgfpathlineto{\pgfqpoint{9.264645in}{0.800393in}}%
\pgfpathlineto{\pgfqpoint{9.264645in}{0.800393in}}%
\pgfpathclose%
\pgfpathmoveto{\pgfqpoint{9.338153in}{0.745060in}}%
\pgfpathquadraticcurveto{\pgfqpoint{9.338153in}{0.756317in}}{\pgfqpoint{9.333506in}{0.762496in}}%
\pgfpathquadraticcurveto{\pgfqpoint{9.328877in}{0.768692in}}{\pgfqpoint{9.320483in}{0.768692in}}%
\pgfpathquadraticcurveto{\pgfqpoint{9.312161in}{0.768692in}}{\pgfqpoint{9.307514in}{0.762496in}}%
\pgfpathquadraticcurveto{\pgfqpoint{9.302867in}{0.756317in}}{\pgfqpoint{9.302867in}{0.745060in}}%
\pgfpathquadraticcurveto{\pgfqpoint{9.302867in}{0.733856in}}{\pgfqpoint{9.307514in}{0.727660in}}%
\pgfpathquadraticcurveto{\pgfqpoint{9.312161in}{0.721464in}}{\pgfqpoint{9.320483in}{0.721464in}}%
\pgfpathquadraticcurveto{\pgfqpoint{9.328877in}{0.721464in}}{\pgfqpoint{9.333506in}{0.727660in}}%
\pgfpathquadraticcurveto{\pgfqpoint{9.338153in}{0.733856in}}{\pgfqpoint{9.338153in}{0.745060in}}%
\pgfpathlineto{\pgfqpoint{9.338153in}{0.745060in}}%
\pgfpathclose%
\pgfpathmoveto{\pgfqpoint{9.348510in}{0.720617in}}%
\pgfpathquadraticcurveto{\pgfqpoint{9.348510in}{0.704532in}}{\pgfqpoint{9.341359in}{0.696679in}}%
\pgfpathquadraticcurveto{\pgfqpoint{9.334226in}{0.688826in}}{\pgfqpoint{9.319474in}{0.688826in}}%
\pgfpathquadraticcurveto{\pgfqpoint{9.314017in}{0.688826in}}{\pgfqpoint{9.309171in}{0.689636in}}%
\pgfpathquadraticcurveto{\pgfqpoint{9.304326in}{0.690447in}}{\pgfqpoint{9.299769in}{0.692140in}}%
\pgfpathlineto{\pgfqpoint{9.299769in}{0.702209in}}%
\pgfpathquadraticcurveto{\pgfqpoint{9.304326in}{0.699741in}}{\pgfqpoint{9.308775in}{0.698570in}}%
\pgfpathquadraticcurveto{\pgfqpoint{9.313224in}{0.697382in}}{\pgfqpoint{9.317835in}{0.697382in}}%
\pgfpathquadraticcurveto{\pgfqpoint{9.328030in}{0.697382in}}{\pgfqpoint{9.333091in}{0.702695in}}%
\pgfpathquadraticcurveto{\pgfqpoint{9.338153in}{0.708009in}}{\pgfqpoint{9.338153in}{0.718762in}}%
\pgfpathlineto{\pgfqpoint{9.338153in}{0.723895in}}%
\pgfpathquadraticcurveto{\pgfqpoint{9.334947in}{0.718312in}}{\pgfqpoint{9.329939in}{0.715556in}}%
\pgfpathquadraticcurveto{\pgfqpoint{9.324932in}{0.712800in}}{\pgfqpoint{9.317943in}{0.712800in}}%
\pgfpathquadraticcurveto{\pgfqpoint{9.306361in}{0.712800in}}{\pgfqpoint{9.299265in}{0.721626in}}%
\pgfpathquadraticcurveto{\pgfqpoint{9.292168in}{0.730470in}}{\pgfqpoint{9.292168in}{0.745060in}}%
\pgfpathquadraticcurveto{\pgfqpoint{9.292168in}{0.759686in}}{\pgfqpoint{9.299265in}{0.768512in}}%
\pgfpathquadraticcurveto{\pgfqpoint{9.306361in}{0.777356in}}{\pgfqpoint{9.317943in}{0.777356in}}%
\pgfpathquadraticcurveto{\pgfqpoint{9.324932in}{0.777356in}}{\pgfqpoint{9.329939in}{0.774600in}}%
\pgfpathquadraticcurveto{\pgfqpoint{9.334947in}{0.771844in}}{\pgfqpoint{9.338153in}{0.766278in}}%
\pgfpathlineto{\pgfqpoint{9.338153in}{0.775843in}}%
\pgfpathlineto{\pgfqpoint{9.348510in}{0.775843in}}%
\pgfpathlineto{\pgfqpoint{9.348510in}{0.720617in}}%
\pgfpathlineto{\pgfqpoint{9.348510in}{0.720617in}}%
\pgfpathclose%
\pgfpathmoveto{\pgfqpoint{9.369854in}{0.775843in}}%
\pgfpathlineto{\pgfqpoint{9.380211in}{0.775843in}}%
\pgfpathlineto{\pgfqpoint{9.380211in}{0.712800in}}%
\pgfpathlineto{\pgfqpoint{9.369854in}{0.712800in}}%
\pgfpathlineto{\pgfqpoint{9.369854in}{0.775843in}}%
\pgfpathlineto{\pgfqpoint{9.369854in}{0.775843in}}%
\pgfpathclose%
\pgfpathmoveto{\pgfqpoint{9.369854in}{0.800393in}}%
\pgfpathlineto{\pgfqpoint{9.380211in}{0.800393in}}%
\pgfpathlineto{\pgfqpoint{9.380211in}{0.787262in}}%
\pgfpathlineto{\pgfqpoint{9.369854in}{0.787262in}}%
\pgfpathlineto{\pgfqpoint{9.369854in}{0.800393in}}%
\pgfpathlineto{\pgfqpoint{9.369854in}{0.800393in}}%
\pgfpathclose%
\pgfpathmoveto{\pgfqpoint{9.447144in}{0.744267in}}%
\pgfpathquadraticcurveto{\pgfqpoint{9.447144in}{0.755687in}}{\pgfqpoint{9.442443in}{0.762189in}}%
\pgfpathquadraticcurveto{\pgfqpoint{9.437742in}{0.768692in}}{\pgfqpoint{9.429529in}{0.768692in}}%
\pgfpathquadraticcurveto{\pgfqpoint{9.421297in}{0.768692in}}{\pgfqpoint{9.416596in}{0.762189in}}%
\pgfpathquadraticcurveto{\pgfqpoint{9.411895in}{0.755687in}}{\pgfqpoint{9.411895in}{0.744267in}}%
\pgfpathquadraticcurveto{\pgfqpoint{9.411895in}{0.732848in}}{\pgfqpoint{9.416596in}{0.726345in}}%
\pgfpathquadraticcurveto{\pgfqpoint{9.421297in}{0.719843in}}{\pgfqpoint{9.429529in}{0.719843in}}%
\pgfpathquadraticcurveto{\pgfqpoint{9.437742in}{0.719843in}}{\pgfqpoint{9.442443in}{0.726345in}}%
\pgfpathquadraticcurveto{\pgfqpoint{9.447144in}{0.732848in}}{\pgfqpoint{9.447144in}{0.744267in}}%
\pgfpathlineto{\pgfqpoint{9.447144in}{0.744267in}}%
\pgfpathclose%
\pgfpathmoveto{\pgfqpoint{9.411895in}{0.766278in}}%
\pgfpathquadraticcurveto{\pgfqpoint{9.415173in}{0.771898in}}{\pgfqpoint{9.420144in}{0.774618in}}%
\pgfpathquadraticcurveto{\pgfqpoint{9.425134in}{0.777356in}}{\pgfqpoint{9.432050in}{0.777356in}}%
\pgfpathquadraticcurveto{\pgfqpoint{9.443542in}{0.777356in}}{\pgfqpoint{9.450711in}{0.768241in}}%
\pgfpathquadraticcurveto{\pgfqpoint{9.457898in}{0.759127in}}{\pgfqpoint{9.457898in}{0.744267in}}%
\pgfpathquadraticcurveto{\pgfqpoint{9.457898in}{0.729407in}}{\pgfqpoint{9.450711in}{0.720275in}}%
\pgfpathquadraticcurveto{\pgfqpoint{9.443542in}{0.711161in}}{\pgfqpoint{9.432050in}{0.711161in}}%
\pgfpathquadraticcurveto{\pgfqpoint{9.425134in}{0.711161in}}{\pgfqpoint{9.420144in}{0.713899in}}%
\pgfpathquadraticcurveto{\pgfqpoint{9.415173in}{0.716637in}}{\pgfqpoint{9.411895in}{0.722256in}}%
\pgfpathlineto{\pgfqpoint{9.411895in}{0.712800in}}%
\pgfpathlineto{\pgfqpoint{9.401484in}{0.712800in}}%
\pgfpathlineto{\pgfqpoint{9.401484in}{0.800393in}}%
\pgfpathlineto{\pgfqpoint{9.411895in}{0.800393in}}%
\pgfpathlineto{\pgfqpoint{9.411895in}{0.766278in}}%
\pgfpathlineto{\pgfqpoint{9.411895in}{0.766278in}}%
\pgfpathclose%
\pgfpathmoveto{\pgfqpoint{9.475063in}{0.800393in}}%
\pgfpathlineto{\pgfqpoint{9.485420in}{0.800393in}}%
\pgfpathlineto{\pgfqpoint{9.485420in}{0.712800in}}%
\pgfpathlineto{\pgfqpoint{9.475063in}{0.712800in}}%
\pgfpathlineto{\pgfqpoint{9.475063in}{0.800393in}}%
\pgfpathlineto{\pgfqpoint{9.475063in}{0.800393in}}%
\pgfpathclose%
\pgfpathmoveto{\pgfqpoint{9.561017in}{0.746915in}}%
\pgfpathlineto{\pgfqpoint{9.561017in}{0.741854in}}%
\pgfpathlineto{\pgfqpoint{9.513393in}{0.741854in}}%
\pgfpathquadraticcurveto{\pgfqpoint{9.514078in}{0.731154in}}{\pgfqpoint{9.519841in}{0.725553in}}%
\pgfpathquadraticcurveto{\pgfqpoint{9.525605in}{0.719951in}}{\pgfqpoint{9.535908in}{0.719951in}}%
\pgfpathquadraticcurveto{\pgfqpoint{9.541870in}{0.719951in}}{\pgfqpoint{9.547472in}{0.721410in}}%
\pgfpathquadraticcurveto{\pgfqpoint{9.553074in}{0.722869in}}{\pgfqpoint{9.558604in}{0.725805in}}%
\pgfpathlineto{\pgfqpoint{9.558604in}{0.716006in}}%
\pgfpathquadraticcurveto{\pgfqpoint{9.553020in}{0.713647in}}{\pgfqpoint{9.547166in}{0.712404in}}%
\pgfpathquadraticcurveto{\pgfqpoint{9.541312in}{0.711161in}}{\pgfqpoint{9.535296in}{0.711161in}}%
\pgfpathquadraticcurveto{\pgfqpoint{9.520202in}{0.711161in}}{\pgfqpoint{9.511394in}{0.719933in}}%
\pgfpathquadraticcurveto{\pgfqpoint{9.502586in}{0.728723in}}{\pgfqpoint{9.502586in}{0.743709in}}%
\pgfpathquadraticcurveto{\pgfqpoint{9.502586in}{0.759181in}}{\pgfqpoint{9.510943in}{0.768259in}}%
\pgfpathquadraticcurveto{\pgfqpoint{9.519301in}{0.777356in}}{\pgfqpoint{9.533495in}{0.777356in}}%
\pgfpathquadraticcurveto{\pgfqpoint{9.546211in}{0.777356in}}{\pgfqpoint{9.553614in}{0.769160in}}%
\pgfpathquadraticcurveto{\pgfqpoint{9.561017in}{0.760983in}}{\pgfqpoint{9.561017in}{0.746915in}}%
\pgfpathlineto{\pgfqpoint{9.561017in}{0.746915in}}%
\pgfpathclose%
\pgfpathmoveto{\pgfqpoint{9.550660in}{0.749959in}}%
\pgfpathquadraticcurveto{\pgfqpoint{9.550552in}{0.758443in}}{\pgfqpoint{9.545905in}{0.763504in}}%
\pgfpathquadraticcurveto{\pgfqpoint{9.541258in}{0.768584in}}{\pgfqpoint{9.533603in}{0.768584in}}%
\pgfpathquadraticcurveto{\pgfqpoint{9.524939in}{0.768584in}}{\pgfqpoint{9.519733in}{0.763684in}}%
\pgfpathquadraticcurveto{\pgfqpoint{9.514528in}{0.758785in}}{\pgfqpoint{9.513735in}{0.749887in}}%
\pgfpathlineto{\pgfqpoint{9.550660in}{0.749959in}}%
\pgfpathlineto{\pgfqpoint{9.550660in}{0.749959in}}%
\pgfpathclose%
\pgfusepath{fill}%
\end{pgfscope}%
\begin{pgfscope}%
\pgfsetbuttcap%
\pgfsetroundjoin%
\definecolor{currentfill}{rgb}{0.309804,0.372549,0.419608}%
\pgfsetfillcolor{currentfill}%
\pgfsetlinewidth{0.000000pt}%
\definecolor{currentstroke}{rgb}{0.309804,0.372549,0.419608}%
\pgfsetstrokecolor{currentstroke}%
\pgfsetdash{}{0pt}%
\pgfpathmoveto{\pgfqpoint{9.663566in}{0.773987in}}%
\pgfpathlineto{\pgfqpoint{9.663566in}{0.764189in}}%
\pgfpathquadraticcurveto{\pgfqpoint{9.659189in}{0.766440in}}{\pgfqpoint{9.654452in}{0.767557in}}%
\pgfpathquadraticcurveto{\pgfqpoint{9.649733in}{0.768692in}}{\pgfqpoint{9.644653in}{0.768692in}}%
\pgfpathquadraticcurveto{\pgfqpoint{9.636944in}{0.768692in}}{\pgfqpoint{9.633089in}{0.766332in}}%
\pgfpathquadraticcurveto{\pgfqpoint{9.629235in}{0.763973in}}{\pgfqpoint{9.629235in}{0.759235in}}%
\pgfpathquadraticcurveto{\pgfqpoint{9.629235in}{0.755633in}}{\pgfqpoint{9.631991in}{0.753580in}}%
\pgfpathquadraticcurveto{\pgfqpoint{9.634746in}{0.751526in}}{\pgfqpoint{9.643086in}{0.749671in}}%
\pgfpathlineto{\pgfqpoint{9.646634in}{0.748878in}}%
\pgfpathquadraticcurveto{\pgfqpoint{9.657658in}{0.746519in}}{\pgfqpoint{9.662305in}{0.742214in}}%
\pgfpathquadraticcurveto{\pgfqpoint{9.666952in}{0.737909in}}{\pgfqpoint{9.666952in}{0.730200in}}%
\pgfpathquadraticcurveto{\pgfqpoint{9.666952in}{0.721410in}}{\pgfqpoint{9.659999in}{0.716276in}}%
\pgfpathquadraticcurveto{\pgfqpoint{9.653047in}{0.711161in}}{\pgfqpoint{9.640889in}{0.711161in}}%
\pgfpathquadraticcurveto{\pgfqpoint{9.635827in}{0.711161in}}{\pgfqpoint{9.630333in}{0.712152in}}%
\pgfpathquadraticcurveto{\pgfqpoint{9.624840in}{0.713142in}}{\pgfqpoint{9.618770in}{0.715106in}}%
\pgfpathlineto{\pgfqpoint{9.618770in}{0.725805in}}%
\pgfpathquadraticcurveto{\pgfqpoint{9.624515in}{0.722815in}}{\pgfqpoint{9.630081in}{0.721320in}}%
\pgfpathquadraticcurveto{\pgfqpoint{9.635647in}{0.719843in}}{\pgfqpoint{9.641123in}{0.719843in}}%
\pgfpathquadraticcurveto{\pgfqpoint{9.648436in}{0.719843in}}{\pgfqpoint{9.652362in}{0.722346in}}%
\pgfpathquadraticcurveto{\pgfqpoint{9.656307in}{0.724850in}}{\pgfqpoint{9.656307in}{0.729407in}}%
\pgfpathquadraticcurveto{\pgfqpoint{9.656307in}{0.733622in}}{\pgfqpoint{9.653461in}{0.735874in}}%
\pgfpathquadraticcurveto{\pgfqpoint{9.650633in}{0.738125in}}{\pgfqpoint{9.640997in}{0.740214in}}%
\pgfpathlineto{\pgfqpoint{9.637394in}{0.741061in}}%
\pgfpathquadraticcurveto{\pgfqpoint{9.627776in}{0.743078in}}{\pgfqpoint{9.623489in}{0.747275in}}%
\pgfpathquadraticcurveto{\pgfqpoint{9.619220in}{0.751472in}}{\pgfqpoint{9.619220in}{0.758785in}}%
\pgfpathquadraticcurveto{\pgfqpoint{9.619220in}{0.767683in}}{\pgfqpoint{9.625524in}{0.772510in}}%
\pgfpathquadraticcurveto{\pgfqpoint{9.631828in}{0.777356in}}{\pgfqpoint{9.643428in}{0.777356in}}%
\pgfpathquadraticcurveto{\pgfqpoint{9.649156in}{0.777356in}}{\pgfqpoint{9.654218in}{0.776509in}}%
\pgfpathquadraticcurveto{\pgfqpoint{9.659297in}{0.775680in}}{\pgfqpoint{9.663566in}{0.773987in}}%
\pgfpathlineto{\pgfqpoint{9.663566in}{0.773987in}}%
\pgfpathclose%
\pgfpathmoveto{\pgfqpoint{9.712091in}{0.744483in}}%
\pgfpathquadraticcurveto{\pgfqpoint{9.699536in}{0.744483in}}{\pgfqpoint{9.694691in}{0.741619in}}%
\pgfpathquadraticcurveto{\pgfqpoint{9.689846in}{0.738756in}}{\pgfqpoint{9.689846in}{0.731821in}}%
\pgfpathquadraticcurveto{\pgfqpoint{9.689846in}{0.726309in}}{\pgfqpoint{9.693484in}{0.723067in}}%
\pgfpathquadraticcurveto{\pgfqpoint{9.697122in}{0.719843in}}{\pgfqpoint{9.703355in}{0.719843in}}%
\pgfpathquadraticcurveto{\pgfqpoint{9.711983in}{0.719843in}}{\pgfqpoint{9.717188in}{0.725949in}}%
\pgfpathquadraticcurveto{\pgfqpoint{9.722394in}{0.732055in}}{\pgfqpoint{9.722394in}{0.742178in}}%
\pgfpathlineto{\pgfqpoint{9.722394in}{0.744483in}}%
\pgfpathlineto{\pgfqpoint{9.712091in}{0.744483in}}%
\pgfpathlineto{\pgfqpoint{9.712091in}{0.744483in}}%
\pgfpathclose%
\pgfpathmoveto{\pgfqpoint{9.732751in}{0.748770in}}%
\pgfpathlineto{\pgfqpoint{9.732751in}{0.712800in}}%
\pgfpathlineto{\pgfqpoint{9.722394in}{0.712800in}}%
\pgfpathlineto{\pgfqpoint{9.722394in}{0.722364in}}%
\pgfpathquadraticcurveto{\pgfqpoint{9.718845in}{0.716637in}}{\pgfqpoint{9.713550in}{0.713899in}}%
\pgfpathquadraticcurveto{\pgfqpoint{9.708254in}{0.711161in}}{\pgfqpoint{9.700599in}{0.711161in}}%
\pgfpathquadraticcurveto{\pgfqpoint{9.690926in}{0.711161in}}{\pgfqpoint{9.685198in}{0.716601in}}%
\pgfpathquadraticcurveto{\pgfqpoint{9.679489in}{0.722040in}}{\pgfqpoint{9.679489in}{0.731154in}}%
\pgfpathquadraticcurveto{\pgfqpoint{9.679489in}{0.741782in}}{\pgfqpoint{9.686603in}{0.747185in}}%
\pgfpathquadraticcurveto{\pgfqpoint{9.693736in}{0.752589in}}{\pgfqpoint{9.707858in}{0.752589in}}%
\pgfpathlineto{\pgfqpoint{9.722394in}{0.752589in}}%
\pgfpathlineto{\pgfqpoint{9.722394in}{0.753616in}}%
\pgfpathquadraticcurveto{\pgfqpoint{9.722394in}{0.760766in}}{\pgfqpoint{9.717692in}{0.764675in}}%
\pgfpathquadraticcurveto{\pgfqpoint{9.712991in}{0.768584in}}{\pgfqpoint{9.704489in}{0.768584in}}%
\pgfpathquadraticcurveto{\pgfqpoint{9.699086in}{0.768584in}}{\pgfqpoint{9.693952in}{0.767287in}}%
\pgfpathquadraticcurveto{\pgfqpoint{9.688837in}{0.765990in}}{\pgfqpoint{9.684118in}{0.763396in}}%
\pgfpathlineto{\pgfqpoint{9.684118in}{0.772979in}}%
\pgfpathquadraticcurveto{\pgfqpoint{9.689792in}{0.775176in}}{\pgfqpoint{9.695141in}{0.776257in}}%
\pgfpathquadraticcurveto{\pgfqpoint{9.700491in}{0.777356in}}{\pgfqpoint{9.705552in}{0.777356in}}%
\pgfpathquadraticcurveto{\pgfqpoint{9.719241in}{0.777356in}}{\pgfqpoint{9.725996in}{0.770259in}}%
\pgfpathquadraticcurveto{\pgfqpoint{9.732751in}{0.763180in}}{\pgfqpoint{9.732751in}{0.748770in}}%
\pgfpathlineto{\pgfqpoint{9.732751in}{0.748770in}}%
\pgfpathclose%
\pgfpathmoveto{\pgfqpoint{9.803160in}{0.763738in}}%
\pgfpathquadraticcurveto{\pgfqpoint{9.807051in}{0.770727in}}{\pgfqpoint{9.812454in}{0.774041in}}%
\pgfpathquadraticcurveto{\pgfqpoint{9.817858in}{0.777356in}}{\pgfqpoint{9.825171in}{0.777356in}}%
\pgfpathquadraticcurveto{\pgfqpoint{9.835024in}{0.777356in}}{\pgfqpoint{9.840373in}{0.770457in}}%
\pgfpathquadraticcurveto{\pgfqpoint{9.845723in}{0.763576in}}{\pgfqpoint{9.845723in}{0.750860in}}%
\pgfpathlineto{\pgfqpoint{9.845723in}{0.712800in}}%
\pgfpathlineto{\pgfqpoint{9.835312in}{0.712800in}}%
\pgfpathlineto{\pgfqpoint{9.835312in}{0.750517in}}%
\pgfpathquadraticcurveto{\pgfqpoint{9.835312in}{0.759578in}}{\pgfqpoint{9.832088in}{0.763955in}}%
\pgfpathquadraticcurveto{\pgfqpoint{9.828881in}{0.768349in}}{\pgfqpoint{9.822307in}{0.768349in}}%
\pgfpathquadraticcurveto{\pgfqpoint{9.814256in}{0.768349in}}{\pgfqpoint{9.809572in}{0.763000in}}%
\pgfpathquadraticcurveto{\pgfqpoint{9.804907in}{0.757668in}}{\pgfqpoint{9.804907in}{0.748428in}}%
\pgfpathlineto{\pgfqpoint{9.804907in}{0.712800in}}%
\pgfpathlineto{\pgfqpoint{9.794496in}{0.712800in}}%
\pgfpathlineto{\pgfqpoint{9.794496in}{0.750517in}}%
\pgfpathquadraticcurveto{\pgfqpoint{9.794496in}{0.759632in}}{\pgfqpoint{9.791290in}{0.763991in}}%
\pgfpathquadraticcurveto{\pgfqpoint{9.788084in}{0.768349in}}{\pgfqpoint{9.781383in}{0.768349in}}%
\pgfpathquadraticcurveto{\pgfqpoint{9.773440in}{0.768349in}}{\pgfqpoint{9.768757in}{0.762982in}}%
\pgfpathquadraticcurveto{\pgfqpoint{9.764092in}{0.757614in}}{\pgfqpoint{9.764092in}{0.748428in}}%
\pgfpathlineto{\pgfqpoint{9.764092in}{0.712800in}}%
\pgfpathlineto{\pgfqpoint{9.753681in}{0.712800in}}%
\pgfpathlineto{\pgfqpoint{9.753681in}{0.775843in}}%
\pgfpathlineto{\pgfqpoint{9.764092in}{0.775843in}}%
\pgfpathlineto{\pgfqpoint{9.764092in}{0.766044in}}%
\pgfpathquadraticcurveto{\pgfqpoint{9.767640in}{0.771844in}}{\pgfqpoint{9.772593in}{0.774600in}}%
\pgfpathquadraticcurveto{\pgfqpoint{9.777547in}{0.777356in}}{\pgfqpoint{9.784355in}{0.777356in}}%
\pgfpathquadraticcurveto{\pgfqpoint{9.791236in}{0.777356in}}{\pgfqpoint{9.796045in}{0.773861in}}%
\pgfpathquadraticcurveto{\pgfqpoint{9.800854in}{0.770385in}}{\pgfqpoint{9.803160in}{0.763738in}}%
\pgfpathlineto{\pgfqpoint{9.803160in}{0.763738in}}%
\pgfpathclose%
\pgfpathmoveto{\pgfqpoint{9.876379in}{0.722256in}}%
\pgfpathlineto{\pgfqpoint{9.876379in}{0.688826in}}%
\pgfpathlineto{\pgfqpoint{9.865968in}{0.688826in}}%
\pgfpathlineto{\pgfqpoint{9.865968in}{0.775843in}}%
\pgfpathlineto{\pgfqpoint{9.876379in}{0.775843in}}%
\pgfpathlineto{\pgfqpoint{9.876379in}{0.766278in}}%
\pgfpathquadraticcurveto{\pgfqpoint{9.879658in}{0.771898in}}{\pgfqpoint{9.884629in}{0.774618in}}%
\pgfpathquadraticcurveto{\pgfqpoint{9.889618in}{0.777356in}}{\pgfqpoint{9.896535in}{0.777356in}}%
\pgfpathquadraticcurveto{\pgfqpoint{9.908027in}{0.777356in}}{\pgfqpoint{9.915196in}{0.768241in}}%
\pgfpathquadraticcurveto{\pgfqpoint{9.922382in}{0.759127in}}{\pgfqpoint{9.922382in}{0.744267in}}%
\pgfpathquadraticcurveto{\pgfqpoint{9.922382in}{0.729407in}}{\pgfqpoint{9.915196in}{0.720275in}}%
\pgfpathquadraticcurveto{\pgfqpoint{9.908027in}{0.711161in}}{\pgfqpoint{9.896535in}{0.711161in}}%
\pgfpathquadraticcurveto{\pgfqpoint{9.889618in}{0.711161in}}{\pgfqpoint{9.884629in}{0.713899in}}%
\pgfpathquadraticcurveto{\pgfqpoint{9.879658in}{0.716637in}}{\pgfqpoint{9.876379in}{0.722256in}}%
\pgfpathlineto{\pgfqpoint{9.876379in}{0.722256in}}%
\pgfpathclose%
\pgfpathmoveto{\pgfqpoint{9.911629in}{0.744267in}}%
\pgfpathquadraticcurveto{\pgfqpoint{9.911629in}{0.755687in}}{\pgfqpoint{9.906928in}{0.762189in}}%
\pgfpathquadraticcurveto{\pgfqpoint{9.902227in}{0.768692in}}{\pgfqpoint{9.894013in}{0.768692in}}%
\pgfpathquadraticcurveto{\pgfqpoint{9.885782in}{0.768692in}}{\pgfqpoint{9.881081in}{0.762189in}}%
\pgfpathquadraticcurveto{\pgfqpoint{9.876379in}{0.755687in}}{\pgfqpoint{9.876379in}{0.744267in}}%
\pgfpathquadraticcurveto{\pgfqpoint{9.876379in}{0.732848in}}{\pgfqpoint{9.881081in}{0.726345in}}%
\pgfpathquadraticcurveto{\pgfqpoint{9.885782in}{0.719843in}}{\pgfqpoint{9.894013in}{0.719843in}}%
\pgfpathquadraticcurveto{\pgfqpoint{9.902227in}{0.719843in}}{\pgfqpoint{9.906928in}{0.726345in}}%
\pgfpathquadraticcurveto{\pgfqpoint{9.911629in}{0.732848in}}{\pgfqpoint{9.911629in}{0.744267in}}%
\pgfpathlineto{\pgfqpoint{9.911629in}{0.744267in}}%
\pgfpathclose%
\pgfpathmoveto{\pgfqpoint{9.939548in}{0.800393in}}%
\pgfpathlineto{\pgfqpoint{9.949905in}{0.800393in}}%
\pgfpathlineto{\pgfqpoint{9.949905in}{0.712800in}}%
\pgfpathlineto{\pgfqpoint{9.939548in}{0.712800in}}%
\pgfpathlineto{\pgfqpoint{9.939548in}{0.800393in}}%
\pgfpathlineto{\pgfqpoint{9.939548in}{0.800393in}}%
\pgfpathclose%
\pgfpathmoveto{\pgfqpoint{10.025502in}{0.746915in}}%
\pgfpathlineto{\pgfqpoint{10.025502in}{0.741854in}}%
\pgfpathlineto{\pgfqpoint{9.977878in}{0.741854in}}%
\pgfpathquadraticcurveto{\pgfqpoint{9.978562in}{0.731154in}}{\pgfqpoint{9.984326in}{0.725553in}}%
\pgfpathquadraticcurveto{\pgfqpoint{9.990090in}{0.719951in}}{\pgfqpoint{10.000393in}{0.719951in}}%
\pgfpathquadraticcurveto{\pgfqpoint{10.006355in}{0.719951in}}{\pgfqpoint{10.011957in}{0.721410in}}%
\pgfpathquadraticcurveto{\pgfqpoint{10.017559in}{0.722869in}}{\pgfqpoint{10.023088in}{0.725805in}}%
\pgfpathlineto{\pgfqpoint{10.023088in}{0.716006in}}%
\pgfpathquadraticcurveto{\pgfqpoint{10.017505in}{0.713647in}}{\pgfqpoint{10.011651in}{0.712404in}}%
\pgfpathquadraticcurveto{\pgfqpoint{10.005797in}{0.711161in}}{\pgfqpoint{9.999781in}{0.711161in}}%
\pgfpathquadraticcurveto{\pgfqpoint{9.984686in}{0.711161in}}{\pgfqpoint{9.975879in}{0.719933in}}%
\pgfpathquadraticcurveto{\pgfqpoint{9.967071in}{0.728723in}}{\pgfqpoint{9.967071in}{0.743709in}}%
\pgfpathquadraticcurveto{\pgfqpoint{9.967071in}{0.759181in}}{\pgfqpoint{9.975428in}{0.768259in}}%
\pgfpathquadraticcurveto{\pgfqpoint{9.983786in}{0.777356in}}{\pgfqpoint{9.997979in}{0.777356in}}%
\pgfpathquadraticcurveto{\pgfqpoint{10.010696in}{0.777356in}}{\pgfqpoint{10.018099in}{0.769160in}}%
\pgfpathquadraticcurveto{\pgfqpoint{10.025502in}{0.760983in}}{\pgfqpoint{10.025502in}{0.746915in}}%
\pgfpathlineto{\pgfqpoint{10.025502in}{0.746915in}}%
\pgfpathclose%
\pgfpathmoveto{\pgfqpoint{10.015145in}{0.749959in}}%
\pgfpathquadraticcurveto{\pgfqpoint{10.015037in}{0.758443in}}{\pgfqpoint{10.010390in}{0.763504in}}%
\pgfpathquadraticcurveto{\pgfqpoint{10.005743in}{0.768584in}}{\pgfqpoint{9.998088in}{0.768584in}}%
\pgfpathquadraticcurveto{\pgfqpoint{9.989424in}{0.768584in}}{\pgfqpoint{9.984218in}{0.763684in}}%
\pgfpathquadraticcurveto{\pgfqpoint{9.979013in}{0.758785in}}{\pgfqpoint{9.978220in}{0.749887in}}%
\pgfpathlineto{\pgfqpoint{10.015145in}{0.749959in}}%
\pgfpathlineto{\pgfqpoint{10.015145in}{0.749959in}}%
\pgfpathclose%
\pgfpathmoveto{\pgfqpoint{10.045153in}{0.772402in}}%
\pgfpathlineto{\pgfqpoint{10.057023in}{0.772402in}}%
\pgfpathlineto{\pgfqpoint{10.057023in}{0.758119in}}%
\pgfpathlineto{\pgfqpoint{10.045153in}{0.758119in}}%
\pgfpathlineto{\pgfqpoint{10.045153in}{0.772402in}}%
\pgfpathlineto{\pgfqpoint{10.045153in}{0.772402in}}%
\pgfpathclose%
\pgfpathmoveto{\pgfqpoint{10.045153in}{0.727102in}}%
\pgfpathlineto{\pgfqpoint{10.057023in}{0.727102in}}%
\pgfpathlineto{\pgfqpoint{10.057023in}{0.717411in}}%
\pgfpathlineto{\pgfqpoint{10.047801in}{0.699399in}}%
\pgfpathlineto{\pgfqpoint{10.040542in}{0.699399in}}%
\pgfpathlineto{\pgfqpoint{10.045153in}{0.717411in}}%
\pgfpathlineto{\pgfqpoint{10.045153in}{0.727102in}}%
\pgfpathlineto{\pgfqpoint{10.045153in}{0.727102in}}%
\pgfpathclose%
\pgfusepath{fill}%
\end{pgfscope}%
\begin{pgfscope}%
\pgfsetbuttcap%
\pgfsetroundjoin%
\definecolor{currentfill}{rgb}{0.309804,0.372549,0.419608}%
\pgfsetfillcolor{currentfill}%
\pgfsetlinewidth{0.000000pt}%
\definecolor{currentstroke}{rgb}{0.309804,0.372549,0.419608}%
\pgfsetstrokecolor{currentstroke}%
\pgfsetdash{}{0pt}%
\pgfpathmoveto{\pgfqpoint{10.184157in}{0.763738in}}%
\pgfpathquadraticcurveto{\pgfqpoint{10.188047in}{0.770727in}}{\pgfqpoint{10.193451in}{0.774041in}}%
\pgfpathquadraticcurveto{\pgfqpoint{10.198855in}{0.777356in}}{\pgfqpoint{10.206168in}{0.777356in}}%
\pgfpathquadraticcurveto{\pgfqpoint{10.216020in}{0.777356in}}{\pgfqpoint{10.221370in}{0.770457in}}%
\pgfpathquadraticcurveto{\pgfqpoint{10.226719in}{0.763576in}}{\pgfqpoint{10.226719in}{0.750860in}}%
\pgfpathlineto{\pgfqpoint{10.226719in}{0.712800in}}%
\pgfpathlineto{\pgfqpoint{10.216308in}{0.712800in}}%
\pgfpathlineto{\pgfqpoint{10.216308in}{0.750517in}}%
\pgfpathquadraticcurveto{\pgfqpoint{10.216308in}{0.759578in}}{\pgfqpoint{10.213084in}{0.763955in}}%
\pgfpathquadraticcurveto{\pgfqpoint{10.209878in}{0.768349in}}{\pgfqpoint{10.203304in}{0.768349in}}%
\pgfpathquadraticcurveto{\pgfqpoint{10.195252in}{0.768349in}}{\pgfqpoint{10.190569in}{0.763000in}}%
\pgfpathquadraticcurveto{\pgfqpoint{10.185904in}{0.757668in}}{\pgfqpoint{10.185904in}{0.748428in}}%
\pgfpathlineto{\pgfqpoint{10.185904in}{0.712800in}}%
\pgfpathlineto{\pgfqpoint{10.175493in}{0.712800in}}%
\pgfpathlineto{\pgfqpoint{10.175493in}{0.750517in}}%
\pgfpathquadraticcurveto{\pgfqpoint{10.175493in}{0.759632in}}{\pgfqpoint{10.172287in}{0.763991in}}%
\pgfpathquadraticcurveto{\pgfqpoint{10.169081in}{0.768349in}}{\pgfqpoint{10.162380in}{0.768349in}}%
\pgfpathquadraticcurveto{\pgfqpoint{10.154437in}{0.768349in}}{\pgfqpoint{10.149754in}{0.762982in}}%
\pgfpathquadraticcurveto{\pgfqpoint{10.145088in}{0.757614in}}{\pgfqpoint{10.145088in}{0.748428in}}%
\pgfpathlineto{\pgfqpoint{10.145088in}{0.712800in}}%
\pgfpathlineto{\pgfqpoint{10.134677in}{0.712800in}}%
\pgfpathlineto{\pgfqpoint{10.134677in}{0.775843in}}%
\pgfpathlineto{\pgfqpoint{10.145088in}{0.775843in}}%
\pgfpathlineto{\pgfqpoint{10.145088in}{0.766044in}}%
\pgfpathquadraticcurveto{\pgfqpoint{10.148637in}{0.771844in}}{\pgfqpoint{10.153590in}{0.774600in}}%
\pgfpathquadraticcurveto{\pgfqpoint{10.158543in}{0.777356in}}{\pgfqpoint{10.165352in}{0.777356in}}%
\pgfpathquadraticcurveto{\pgfqpoint{10.172233in}{0.777356in}}{\pgfqpoint{10.177042in}{0.773861in}}%
\pgfpathquadraticcurveto{\pgfqpoint{10.181851in}{0.770385in}}{\pgfqpoint{10.184157in}{0.763738in}}%
\pgfpathlineto{\pgfqpoint{10.184157in}{0.763738in}}%
\pgfpathclose%
\pgfpathmoveto{\pgfqpoint{10.271786in}{0.768584in}}%
\pgfpathquadraticcurveto{\pgfqpoint{10.263464in}{0.768584in}}{\pgfqpoint{10.258619in}{0.762081in}}%
\pgfpathquadraticcurveto{\pgfqpoint{10.253774in}{0.755579in}}{\pgfqpoint{10.253774in}{0.744267in}}%
\pgfpathquadraticcurveto{\pgfqpoint{10.253774in}{0.732956in}}{\pgfqpoint{10.258583in}{0.726453in}}%
\pgfpathquadraticcurveto{\pgfqpoint{10.263410in}{0.719951in}}{\pgfqpoint{10.271786in}{0.719951in}}%
\pgfpathquadraticcurveto{\pgfqpoint{10.280071in}{0.719951in}}{\pgfqpoint{10.284899in}{0.726471in}}%
\pgfpathquadraticcurveto{\pgfqpoint{10.289744in}{0.733010in}}{\pgfqpoint{10.289744in}{0.744267in}}%
\pgfpathquadraticcurveto{\pgfqpoint{10.289744in}{0.755471in}}{\pgfqpoint{10.284899in}{0.762027in}}%
\pgfpathquadraticcurveto{\pgfqpoint{10.280071in}{0.768584in}}{\pgfqpoint{10.271786in}{0.768584in}}%
\pgfpathlineto{\pgfqpoint{10.271786in}{0.768584in}}%
\pgfpathclose%
\pgfpathmoveto{\pgfqpoint{10.271786in}{0.777356in}}%
\pgfpathquadraticcurveto{\pgfqpoint{10.285295in}{0.777356in}}{\pgfqpoint{10.293004in}{0.768566in}}%
\pgfpathquadraticcurveto{\pgfqpoint{10.300731in}{0.759794in}}{\pgfqpoint{10.300731in}{0.744267in}}%
\pgfpathquadraticcurveto{\pgfqpoint{10.300731in}{0.728795in}}{\pgfqpoint{10.293004in}{0.719969in}}%
\pgfpathquadraticcurveto{\pgfqpoint{10.285295in}{0.711161in}}{\pgfqpoint{10.271786in}{0.711161in}}%
\pgfpathquadraticcurveto{\pgfqpoint{10.258223in}{0.711161in}}{\pgfqpoint{10.250532in}{0.719969in}}%
\pgfpathquadraticcurveto{\pgfqpoint{10.242858in}{0.728795in}}{\pgfqpoint{10.242858in}{0.744267in}}%
\pgfpathquadraticcurveto{\pgfqpoint{10.242858in}{0.759794in}}{\pgfqpoint{10.250532in}{0.768566in}}%
\pgfpathquadraticcurveto{\pgfqpoint{10.258223in}{0.777356in}}{\pgfqpoint{10.271786in}{0.777356in}}%
\pgfpathlineto{\pgfqpoint{10.271786in}{0.777356in}}%
\pgfpathclose%
\pgfpathmoveto{\pgfqpoint{10.359379in}{0.766278in}}%
\pgfpathlineto{\pgfqpoint{10.359379in}{0.800393in}}%
\pgfpathlineto{\pgfqpoint{10.369736in}{0.800393in}}%
\pgfpathlineto{\pgfqpoint{10.369736in}{0.712800in}}%
\pgfpathlineto{\pgfqpoint{10.359379in}{0.712800in}}%
\pgfpathlineto{\pgfqpoint{10.359379in}{0.722256in}}%
\pgfpathquadraticcurveto{\pgfqpoint{10.356119in}{0.716637in}}{\pgfqpoint{10.351129in}{0.713899in}}%
\pgfpathquadraticcurveto{\pgfqpoint{10.346158in}{0.711161in}}{\pgfqpoint{10.339169in}{0.711161in}}%
\pgfpathquadraticcurveto{\pgfqpoint{10.327750in}{0.711161in}}{\pgfqpoint{10.320563in}{0.720275in}}%
\pgfpathquadraticcurveto{\pgfqpoint{10.313394in}{0.729407in}}{\pgfqpoint{10.313394in}{0.744267in}}%
\pgfpathquadraticcurveto{\pgfqpoint{10.313394in}{0.759127in}}{\pgfqpoint{10.320563in}{0.768241in}}%
\pgfpathquadraticcurveto{\pgfqpoint{10.327750in}{0.777356in}}{\pgfqpoint{10.339169in}{0.777356in}}%
\pgfpathquadraticcurveto{\pgfqpoint{10.346158in}{0.777356in}}{\pgfqpoint{10.351129in}{0.774618in}}%
\pgfpathquadraticcurveto{\pgfqpoint{10.356119in}{0.771898in}}{\pgfqpoint{10.359379in}{0.766278in}}%
\pgfpathlineto{\pgfqpoint{10.359379in}{0.766278in}}%
\pgfpathclose%
\pgfpathmoveto{\pgfqpoint{10.324093in}{0.744267in}}%
\pgfpathquadraticcurveto{\pgfqpoint{10.324093in}{0.732848in}}{\pgfqpoint{10.328794in}{0.726345in}}%
\pgfpathquadraticcurveto{\pgfqpoint{10.333495in}{0.719843in}}{\pgfqpoint{10.341709in}{0.719843in}}%
\pgfpathquadraticcurveto{\pgfqpoint{10.349923in}{0.719843in}}{\pgfqpoint{10.354642in}{0.726345in}}%
\pgfpathquadraticcurveto{\pgfqpoint{10.359379in}{0.732848in}}{\pgfqpoint{10.359379in}{0.744267in}}%
\pgfpathquadraticcurveto{\pgfqpoint{10.359379in}{0.755687in}}{\pgfqpoint{10.354642in}{0.762189in}}%
\pgfpathquadraticcurveto{\pgfqpoint{10.349923in}{0.768692in}}{\pgfqpoint{10.341709in}{0.768692in}}%
\pgfpathquadraticcurveto{\pgfqpoint{10.333495in}{0.768692in}}{\pgfqpoint{10.328794in}{0.762189in}}%
\pgfpathquadraticcurveto{\pgfqpoint{10.324093in}{0.755687in}}{\pgfqpoint{10.324093in}{0.744267in}}%
\pgfpathlineto{\pgfqpoint{10.324093in}{0.744267in}}%
\pgfpathclose%
\pgfpathmoveto{\pgfqpoint{10.445009in}{0.746915in}}%
\pgfpathlineto{\pgfqpoint{10.445009in}{0.741854in}}%
\pgfpathlineto{\pgfqpoint{10.397385in}{0.741854in}}%
\pgfpathquadraticcurveto{\pgfqpoint{10.398069in}{0.731154in}}{\pgfqpoint{10.403833in}{0.725553in}}%
\pgfpathquadraticcurveto{\pgfqpoint{10.409597in}{0.719951in}}{\pgfqpoint{10.419900in}{0.719951in}}%
\pgfpathquadraticcurveto{\pgfqpoint{10.425862in}{0.719951in}}{\pgfqpoint{10.431464in}{0.721410in}}%
\pgfpathquadraticcurveto{\pgfqpoint{10.437065in}{0.722869in}}{\pgfqpoint{10.442595in}{0.725805in}}%
\pgfpathlineto{\pgfqpoint{10.442595in}{0.716006in}}%
\pgfpathquadraticcurveto{\pgfqpoint{10.437011in}{0.713647in}}{\pgfqpoint{10.431157in}{0.712404in}}%
\pgfpathquadraticcurveto{\pgfqpoint{10.425303in}{0.711161in}}{\pgfqpoint{10.419287in}{0.711161in}}%
\pgfpathquadraticcurveto{\pgfqpoint{10.404193in}{0.711161in}}{\pgfqpoint{10.395385in}{0.719933in}}%
\pgfpathquadraticcurveto{\pgfqpoint{10.386577in}{0.728723in}}{\pgfqpoint{10.386577in}{0.743709in}}%
\pgfpathquadraticcurveto{\pgfqpoint{10.386577in}{0.759181in}}{\pgfqpoint{10.394935in}{0.768259in}}%
\pgfpathquadraticcurveto{\pgfqpoint{10.403293in}{0.777356in}}{\pgfqpoint{10.417486in}{0.777356in}}%
\pgfpathquadraticcurveto{\pgfqpoint{10.430203in}{0.777356in}}{\pgfqpoint{10.437606in}{0.769160in}}%
\pgfpathquadraticcurveto{\pgfqpoint{10.445009in}{0.760983in}}{\pgfqpoint{10.445009in}{0.746915in}}%
\pgfpathlineto{\pgfqpoint{10.445009in}{0.746915in}}%
\pgfpathclose%
\pgfpathmoveto{\pgfqpoint{10.434652in}{0.749959in}}%
\pgfpathquadraticcurveto{\pgfqpoint{10.434544in}{0.758443in}}{\pgfqpoint{10.429897in}{0.763504in}}%
\pgfpathquadraticcurveto{\pgfqpoint{10.425249in}{0.768584in}}{\pgfqpoint{10.417594in}{0.768584in}}%
\pgfpathquadraticcurveto{\pgfqpoint{10.408930in}{0.768584in}}{\pgfqpoint{10.403725in}{0.763684in}}%
\pgfpathquadraticcurveto{\pgfqpoint{10.398519in}{0.758785in}}{\pgfqpoint{10.397727in}{0.749887in}}%
\pgfpathlineto{\pgfqpoint{10.434652in}{0.749959in}}%
\pgfpathlineto{\pgfqpoint{10.434652in}{0.749959in}}%
\pgfpathclose%
\pgfpathmoveto{\pgfqpoint{10.498541in}{0.766170in}}%
\pgfpathquadraticcurveto{\pgfqpoint{10.496794in}{0.767179in}}{\pgfqpoint{10.494740in}{0.767647in}}%
\pgfpathquadraticcurveto{\pgfqpoint{10.492687in}{0.768133in}}{\pgfqpoint{10.490219in}{0.768133in}}%
\pgfpathquadraticcurveto{\pgfqpoint{10.481429in}{0.768133in}}{\pgfqpoint{10.476728in}{0.762423in}}%
\pgfpathquadraticcurveto{\pgfqpoint{10.472027in}{0.756714in}}{\pgfqpoint{10.472027in}{0.746014in}}%
\pgfpathlineto{\pgfqpoint{10.472027in}{0.712800in}}%
\pgfpathlineto{\pgfqpoint{10.461616in}{0.712800in}}%
\pgfpathlineto{\pgfqpoint{10.461616in}{0.775843in}}%
\pgfpathlineto{\pgfqpoint{10.472027in}{0.775843in}}%
\pgfpathlineto{\pgfqpoint{10.472027in}{0.766044in}}%
\pgfpathquadraticcurveto{\pgfqpoint{10.475305in}{0.771790in}}{\pgfqpoint{10.480529in}{0.774564in}}%
\pgfpathquadraticcurveto{\pgfqpoint{10.485770in}{0.777356in}}{\pgfqpoint{10.493263in}{0.777356in}}%
\pgfpathquadraticcurveto{\pgfqpoint{10.494326in}{0.777356in}}{\pgfqpoint{10.495623in}{0.777211in}}%
\pgfpathquadraticcurveto{\pgfqpoint{10.496920in}{0.777085in}}{\pgfqpoint{10.498487in}{0.776797in}}%
\pgfpathlineto{\pgfqpoint{10.498541in}{0.766170in}}%
\pgfpathlineto{\pgfqpoint{10.498541in}{0.766170in}}%
\pgfpathclose%
\pgfpathmoveto{\pgfqpoint{10.559800in}{0.750860in}}%
\pgfpathlineto{\pgfqpoint{10.559800in}{0.712800in}}%
\pgfpathlineto{\pgfqpoint{10.549443in}{0.712800in}}%
\pgfpathlineto{\pgfqpoint{10.549443in}{0.750517in}}%
\pgfpathquadraticcurveto{\pgfqpoint{10.549443in}{0.759469in}}{\pgfqpoint{10.545949in}{0.763900in}}%
\pgfpathquadraticcurveto{\pgfqpoint{10.542454in}{0.768349in}}{\pgfqpoint{10.535484in}{0.768349in}}%
\pgfpathquadraticcurveto{\pgfqpoint{10.527090in}{0.768349in}}{\pgfqpoint{10.522245in}{0.763000in}}%
\pgfpathquadraticcurveto{\pgfqpoint{10.517400in}{0.757668in}}{\pgfqpoint{10.517400in}{0.748428in}}%
\pgfpathlineto{\pgfqpoint{10.517400in}{0.712800in}}%
\pgfpathlineto{\pgfqpoint{10.506989in}{0.712800in}}%
\pgfpathlineto{\pgfqpoint{10.506989in}{0.775843in}}%
\pgfpathlineto{\pgfqpoint{10.517400in}{0.775843in}}%
\pgfpathlineto{\pgfqpoint{10.517400in}{0.766044in}}%
\pgfpathquadraticcurveto{\pgfqpoint{10.521128in}{0.771736in}}{\pgfqpoint{10.526153in}{0.774546in}}%
\pgfpathquadraticcurveto{\pgfqpoint{10.531197in}{0.777356in}}{\pgfqpoint{10.537789in}{0.777356in}}%
\pgfpathquadraticcurveto{\pgfqpoint{10.548651in}{0.777356in}}{\pgfqpoint{10.554216in}{0.770637in}}%
\pgfpathquadraticcurveto{\pgfqpoint{10.559800in}{0.763918in}}{\pgfqpoint{10.559800in}{0.750860in}}%
\pgfpathlineto{\pgfqpoint{10.559800in}{0.750860in}}%
\pgfpathclose%
\pgfusepath{fill}%
\end{pgfscope}%
\begin{pgfscope}%
\pgfsetbuttcap%
\pgfsetroundjoin%
\definecolor{currentfill}{rgb}{0.309804,0.372549,0.419608}%
\pgfsetfillcolor{currentfill}%
\pgfsetlinewidth{0.000000pt}%
\definecolor{currentstroke}{rgb}{0.309804,0.372549,0.419608}%
\pgfsetstrokecolor{currentstroke}%
\pgfsetdash{}{0pt}%
\pgfpathmoveto{\pgfqpoint{10.698664in}{0.746915in}}%
\pgfpathlineto{\pgfqpoint{10.698664in}{0.741854in}}%
\pgfpathlineto{\pgfqpoint{10.651040in}{0.741854in}}%
\pgfpathquadraticcurveto{\pgfqpoint{10.651724in}{0.731154in}}{\pgfqpoint{10.657488in}{0.725553in}}%
\pgfpathquadraticcurveto{\pgfqpoint{10.663252in}{0.719951in}}{\pgfqpoint{10.673555in}{0.719951in}}%
\pgfpathquadraticcurveto{\pgfqpoint{10.679517in}{0.719951in}}{\pgfqpoint{10.685118in}{0.721410in}}%
\pgfpathquadraticcurveto{\pgfqpoint{10.690720in}{0.722869in}}{\pgfqpoint{10.696250in}{0.725805in}}%
\pgfpathlineto{\pgfqpoint{10.696250in}{0.716006in}}%
\pgfpathquadraticcurveto{\pgfqpoint{10.690666in}{0.713647in}}{\pgfqpoint{10.684812in}{0.712404in}}%
\pgfpathquadraticcurveto{\pgfqpoint{10.678958in}{0.711161in}}{\pgfqpoint{10.672942in}{0.711161in}}%
\pgfpathquadraticcurveto{\pgfqpoint{10.657848in}{0.711161in}}{\pgfqpoint{10.649040in}{0.719933in}}%
\pgfpathquadraticcurveto{\pgfqpoint{10.640232in}{0.728723in}}{\pgfqpoint{10.640232in}{0.743709in}}%
\pgfpathquadraticcurveto{\pgfqpoint{10.640232in}{0.759181in}}{\pgfqpoint{10.648590in}{0.768259in}}%
\pgfpathquadraticcurveto{\pgfqpoint{10.656947in}{0.777356in}}{\pgfqpoint{10.671141in}{0.777356in}}%
\pgfpathquadraticcurveto{\pgfqpoint{10.683858in}{0.777356in}}{\pgfqpoint{10.691261in}{0.769160in}}%
\pgfpathquadraticcurveto{\pgfqpoint{10.698664in}{0.760983in}}{\pgfqpoint{10.698664in}{0.746915in}}%
\pgfpathlineto{\pgfqpoint{10.698664in}{0.746915in}}%
\pgfpathclose%
\pgfpathmoveto{\pgfqpoint{10.688307in}{0.749959in}}%
\pgfpathquadraticcurveto{\pgfqpoint{10.688199in}{0.758443in}}{\pgfqpoint{10.683551in}{0.763504in}}%
\pgfpathquadraticcurveto{\pgfqpoint{10.678904in}{0.768584in}}{\pgfqpoint{10.671249in}{0.768584in}}%
\pgfpathquadraticcurveto{\pgfqpoint{10.662585in}{0.768584in}}{\pgfqpoint{10.657380in}{0.763684in}}%
\pgfpathquadraticcurveto{\pgfqpoint{10.652174in}{0.758785in}}{\pgfqpoint{10.651382in}{0.749887in}}%
\pgfpathlineto{\pgfqpoint{10.688307in}{0.749959in}}%
\pgfpathlineto{\pgfqpoint{10.688307in}{0.749959in}}%
\pgfpathclose%
\pgfpathmoveto{\pgfqpoint{10.755852in}{0.773987in}}%
\pgfpathlineto{\pgfqpoint{10.755852in}{0.764189in}}%
\pgfpathquadraticcurveto{\pgfqpoint{10.751475in}{0.766440in}}{\pgfqpoint{10.746738in}{0.767557in}}%
\pgfpathquadraticcurveto{\pgfqpoint{10.742019in}{0.768692in}}{\pgfqpoint{10.736939in}{0.768692in}}%
\pgfpathquadraticcurveto{\pgfqpoint{10.729230in}{0.768692in}}{\pgfqpoint{10.725376in}{0.766332in}}%
\pgfpathquadraticcurveto{\pgfqpoint{10.721521in}{0.763973in}}{\pgfqpoint{10.721521in}{0.759235in}}%
\pgfpathquadraticcurveto{\pgfqpoint{10.721521in}{0.755633in}}{\pgfqpoint{10.724277in}{0.753580in}}%
\pgfpathquadraticcurveto{\pgfqpoint{10.727033in}{0.751526in}}{\pgfqpoint{10.735372in}{0.749671in}}%
\pgfpathlineto{\pgfqpoint{10.738921in}{0.748878in}}%
\pgfpathquadraticcurveto{\pgfqpoint{10.749944in}{0.746519in}}{\pgfqpoint{10.754591in}{0.742214in}}%
\pgfpathquadraticcurveto{\pgfqpoint{10.759239in}{0.737909in}}{\pgfqpoint{10.759239in}{0.730200in}}%
\pgfpathquadraticcurveto{\pgfqpoint{10.759239in}{0.721410in}}{\pgfqpoint{10.752286in}{0.716276in}}%
\pgfpathquadraticcurveto{\pgfqpoint{10.745333in}{0.711161in}}{\pgfqpoint{10.733175in}{0.711161in}}%
\pgfpathquadraticcurveto{\pgfqpoint{10.728114in}{0.711161in}}{\pgfqpoint{10.722620in}{0.712152in}}%
\pgfpathquadraticcurveto{\pgfqpoint{10.717126in}{0.713142in}}{\pgfqpoint{10.711056in}{0.715106in}}%
\pgfpathlineto{\pgfqpoint{10.711056in}{0.725805in}}%
\pgfpathquadraticcurveto{\pgfqpoint{10.716802in}{0.722815in}}{\pgfqpoint{10.722368in}{0.721320in}}%
\pgfpathquadraticcurveto{\pgfqpoint{10.727933in}{0.719843in}}{\pgfqpoint{10.733409in}{0.719843in}}%
\pgfpathquadraticcurveto{\pgfqpoint{10.740722in}{0.719843in}}{\pgfqpoint{10.744649in}{0.722346in}}%
\pgfpathquadraticcurveto{\pgfqpoint{10.748593in}{0.724850in}}{\pgfqpoint{10.748593in}{0.729407in}}%
\pgfpathquadraticcurveto{\pgfqpoint{10.748593in}{0.733622in}}{\pgfqpoint{10.745747in}{0.735874in}}%
\pgfpathquadraticcurveto{\pgfqpoint{10.742919in}{0.738125in}}{\pgfqpoint{10.733283in}{0.740214in}}%
\pgfpathlineto{\pgfqpoint{10.729681in}{0.741061in}}%
\pgfpathquadraticcurveto{\pgfqpoint{10.720062in}{0.743078in}}{\pgfqpoint{10.715775in}{0.747275in}}%
\pgfpathquadraticcurveto{\pgfqpoint{10.711506in}{0.751472in}}{\pgfqpoint{10.711506in}{0.758785in}}%
\pgfpathquadraticcurveto{\pgfqpoint{10.711506in}{0.767683in}}{\pgfqpoint{10.717811in}{0.772510in}}%
\pgfpathquadraticcurveto{\pgfqpoint{10.724115in}{0.777356in}}{\pgfqpoint{10.735715in}{0.777356in}}%
\pgfpathquadraticcurveto{\pgfqpoint{10.741442in}{0.777356in}}{\pgfqpoint{10.746504in}{0.776509in}}%
\pgfpathquadraticcurveto{\pgfqpoint{10.751583in}{0.775680in}}{\pgfqpoint{10.755852in}{0.773987in}}%
\pgfpathlineto{\pgfqpoint{10.755852in}{0.773987in}}%
\pgfpathclose%
\pgfpathmoveto{\pgfqpoint{10.785969in}{0.793747in}}%
\pgfpathlineto{\pgfqpoint{10.785969in}{0.775843in}}%
\pgfpathlineto{\pgfqpoint{10.807295in}{0.775843in}}%
\pgfpathlineto{\pgfqpoint{10.807295in}{0.767791in}}%
\pgfpathlineto{\pgfqpoint{10.785969in}{0.767791in}}%
\pgfpathlineto{\pgfqpoint{10.785969in}{0.733568in}}%
\pgfpathquadraticcurveto{\pgfqpoint{10.785969in}{0.725859in}}{\pgfqpoint{10.788076in}{0.723661in}}%
\pgfpathquadraticcurveto{\pgfqpoint{10.790183in}{0.721464in}}{\pgfqpoint{10.796668in}{0.721464in}}%
\pgfpathlineto{\pgfqpoint{10.807295in}{0.721464in}}%
\pgfpathlineto{\pgfqpoint{10.807295in}{0.712800in}}%
\pgfpathlineto{\pgfqpoint{10.796668in}{0.712800in}}%
\pgfpathquadraticcurveto{\pgfqpoint{10.784672in}{0.712800in}}{\pgfqpoint{10.780115in}{0.717267in}}%
\pgfpathquadraticcurveto{\pgfqpoint{10.775558in}{0.721752in}}{\pgfqpoint{10.775558in}{0.733568in}}%
\pgfpathlineto{\pgfqpoint{10.775558in}{0.767791in}}%
\pgfpathlineto{\pgfqpoint{10.767956in}{0.767791in}}%
\pgfpathlineto{\pgfqpoint{10.767956in}{0.775843in}}%
\pgfpathlineto{\pgfqpoint{10.775558in}{0.775843in}}%
\pgfpathlineto{\pgfqpoint{10.775558in}{0.793747in}}%
\pgfpathlineto{\pgfqpoint{10.785969in}{0.793747in}}%
\pgfpathlineto{\pgfqpoint{10.785969in}{0.793747in}}%
\pgfpathclose%
\pgfpathmoveto{\pgfqpoint{10.820912in}{0.775843in}}%
\pgfpathlineto{\pgfqpoint{10.831269in}{0.775843in}}%
\pgfpathlineto{\pgfqpoint{10.831269in}{0.712800in}}%
\pgfpathlineto{\pgfqpoint{10.820912in}{0.712800in}}%
\pgfpathlineto{\pgfqpoint{10.820912in}{0.775843in}}%
\pgfpathlineto{\pgfqpoint{10.820912in}{0.775843in}}%
\pgfpathclose%
\pgfpathmoveto{\pgfqpoint{10.820912in}{0.800393in}}%
\pgfpathlineto{\pgfqpoint{10.831269in}{0.800393in}}%
\pgfpathlineto{\pgfqpoint{10.831269in}{0.787262in}}%
\pgfpathlineto{\pgfqpoint{10.820912in}{0.787262in}}%
\pgfpathlineto{\pgfqpoint{10.820912in}{0.800393in}}%
\pgfpathlineto{\pgfqpoint{10.820912in}{0.800393in}}%
\pgfpathclose%
\pgfpathmoveto{\pgfqpoint{10.902021in}{0.763738in}}%
\pgfpathquadraticcurveto{\pgfqpoint{10.905911in}{0.770727in}}{\pgfqpoint{10.911315in}{0.774041in}}%
\pgfpathquadraticcurveto{\pgfqpoint{10.916719in}{0.777356in}}{\pgfqpoint{10.924032in}{0.777356in}}%
\pgfpathquadraticcurveto{\pgfqpoint{10.933884in}{0.777356in}}{\pgfqpoint{10.939234in}{0.770457in}}%
\pgfpathquadraticcurveto{\pgfqpoint{10.944584in}{0.763576in}}{\pgfqpoint{10.944584in}{0.750860in}}%
\pgfpathlineto{\pgfqpoint{10.944584in}{0.712800in}}%
\pgfpathlineto{\pgfqpoint{10.934173in}{0.712800in}}%
\pgfpathlineto{\pgfqpoint{10.934173in}{0.750517in}}%
\pgfpathquadraticcurveto{\pgfqpoint{10.934173in}{0.759578in}}{\pgfqpoint{10.930948in}{0.763955in}}%
\pgfpathquadraticcurveto{\pgfqpoint{10.927742in}{0.768349in}}{\pgfqpoint{10.921168in}{0.768349in}}%
\pgfpathquadraticcurveto{\pgfqpoint{10.913116in}{0.768349in}}{\pgfqpoint{10.908433in}{0.763000in}}%
\pgfpathquadraticcurveto{\pgfqpoint{10.903768in}{0.757668in}}{\pgfqpoint{10.903768in}{0.748428in}}%
\pgfpathlineto{\pgfqpoint{10.903768in}{0.712800in}}%
\pgfpathlineto{\pgfqpoint{10.893357in}{0.712800in}}%
\pgfpathlineto{\pgfqpoint{10.893357in}{0.750517in}}%
\pgfpathquadraticcurveto{\pgfqpoint{10.893357in}{0.759632in}}{\pgfqpoint{10.890151in}{0.763991in}}%
\pgfpathquadraticcurveto{\pgfqpoint{10.886945in}{0.768349in}}{\pgfqpoint{10.880244in}{0.768349in}}%
\pgfpathquadraticcurveto{\pgfqpoint{10.872301in}{0.768349in}}{\pgfqpoint{10.867618in}{0.762982in}}%
\pgfpathquadraticcurveto{\pgfqpoint{10.862952in}{0.757614in}}{\pgfqpoint{10.862952in}{0.748428in}}%
\pgfpathlineto{\pgfqpoint{10.862952in}{0.712800in}}%
\pgfpathlineto{\pgfqpoint{10.852541in}{0.712800in}}%
\pgfpathlineto{\pgfqpoint{10.852541in}{0.775843in}}%
\pgfpathlineto{\pgfqpoint{10.862952in}{0.775843in}}%
\pgfpathlineto{\pgfqpoint{10.862952in}{0.766044in}}%
\pgfpathquadraticcurveto{\pgfqpoint{10.866501in}{0.771844in}}{\pgfqpoint{10.871454in}{0.774600in}}%
\pgfpathquadraticcurveto{\pgfqpoint{10.876408in}{0.777356in}}{\pgfqpoint{10.883216in}{0.777356in}}%
\pgfpathquadraticcurveto{\pgfqpoint{10.890097in}{0.777356in}}{\pgfqpoint{10.894906in}{0.773861in}}%
\pgfpathquadraticcurveto{\pgfqpoint{10.899715in}{0.770385in}}{\pgfqpoint{10.902021in}{0.763738in}}%
\pgfpathlineto{\pgfqpoint{10.902021in}{0.763738in}}%
\pgfpathclose%
\pgfpathmoveto{\pgfqpoint{10.993883in}{0.744483in}}%
\pgfpathquadraticcurveto{\pgfqpoint{10.981328in}{0.744483in}}{\pgfqpoint{10.976483in}{0.741619in}}%
\pgfpathquadraticcurveto{\pgfqpoint{10.971638in}{0.738756in}}{\pgfqpoint{10.971638in}{0.731821in}}%
\pgfpathquadraticcurveto{\pgfqpoint{10.971638in}{0.726309in}}{\pgfqpoint{10.975276in}{0.723067in}}%
\pgfpathquadraticcurveto{\pgfqpoint{10.978915in}{0.719843in}}{\pgfqpoint{10.985147in}{0.719843in}}%
\pgfpathquadraticcurveto{\pgfqpoint{10.993775in}{0.719843in}}{\pgfqpoint{10.998980in}{0.725949in}}%
\pgfpathquadraticcurveto{\pgfqpoint{11.004186in}{0.732055in}}{\pgfqpoint{11.004186in}{0.742178in}}%
\pgfpathlineto{\pgfqpoint{11.004186in}{0.744483in}}%
\pgfpathlineto{\pgfqpoint{10.993883in}{0.744483in}}%
\pgfpathlineto{\pgfqpoint{10.993883in}{0.744483in}}%
\pgfpathclose%
\pgfpathmoveto{\pgfqpoint{11.014543in}{0.748770in}}%
\pgfpathlineto{\pgfqpoint{11.014543in}{0.712800in}}%
\pgfpathlineto{\pgfqpoint{11.004186in}{0.712800in}}%
\pgfpathlineto{\pgfqpoint{11.004186in}{0.722364in}}%
\pgfpathquadraticcurveto{\pgfqpoint{11.000637in}{0.716637in}}{\pgfqpoint{10.995342in}{0.713899in}}%
\pgfpathquadraticcurveto{\pgfqpoint{10.990046in}{0.711161in}}{\pgfqpoint{10.982391in}{0.711161in}}%
\pgfpathquadraticcurveto{\pgfqpoint{10.972719in}{0.711161in}}{\pgfqpoint{10.966991in}{0.716601in}}%
\pgfpathquadraticcurveto{\pgfqpoint{10.961281in}{0.722040in}}{\pgfqpoint{10.961281in}{0.731154in}}%
\pgfpathquadraticcurveto{\pgfqpoint{10.961281in}{0.741782in}}{\pgfqpoint{10.968396in}{0.747185in}}%
\pgfpathquadraticcurveto{\pgfqpoint{10.975528in}{0.752589in}}{\pgfqpoint{10.989650in}{0.752589in}}%
\pgfpathlineto{\pgfqpoint{11.004186in}{0.752589in}}%
\pgfpathlineto{\pgfqpoint{11.004186in}{0.753616in}}%
\pgfpathquadraticcurveto{\pgfqpoint{11.004186in}{0.760766in}}{\pgfqpoint{10.999485in}{0.764675in}}%
\pgfpathquadraticcurveto{\pgfqpoint{10.994783in}{0.768584in}}{\pgfqpoint{10.986282in}{0.768584in}}%
\pgfpathquadraticcurveto{\pgfqpoint{10.980878in}{0.768584in}}{\pgfqpoint{10.975745in}{0.767287in}}%
\pgfpathquadraticcurveto{\pgfqpoint{10.970629in}{0.765990in}}{\pgfqpoint{10.965910in}{0.763396in}}%
\pgfpathlineto{\pgfqpoint{10.965910in}{0.772979in}}%
\pgfpathquadraticcurveto{\pgfqpoint{10.971584in}{0.775176in}}{\pgfqpoint{10.976933in}{0.776257in}}%
\pgfpathquadraticcurveto{\pgfqpoint{10.982283in}{0.777356in}}{\pgfqpoint{10.987344in}{0.777356in}}%
\pgfpathquadraticcurveto{\pgfqpoint{11.001034in}{0.777356in}}{\pgfqpoint{11.007788in}{0.770259in}}%
\pgfpathquadraticcurveto{\pgfqpoint{11.014543in}{0.763180in}}{\pgfqpoint{11.014543in}{0.748770in}}%
\pgfpathlineto{\pgfqpoint{11.014543in}{0.748770in}}%
\pgfpathclose%
\pgfpathmoveto{\pgfqpoint{11.046118in}{0.793747in}}%
\pgfpathlineto{\pgfqpoint{11.046118in}{0.775843in}}%
\pgfpathlineto{\pgfqpoint{11.067444in}{0.775843in}}%
\pgfpathlineto{\pgfqpoint{11.067444in}{0.767791in}}%
\pgfpathlineto{\pgfqpoint{11.046118in}{0.767791in}}%
\pgfpathlineto{\pgfqpoint{11.046118in}{0.733568in}}%
\pgfpathquadraticcurveto{\pgfqpoint{11.046118in}{0.725859in}}{\pgfqpoint{11.048225in}{0.723661in}}%
\pgfpathquadraticcurveto{\pgfqpoint{11.050333in}{0.721464in}}{\pgfqpoint{11.056817in}{0.721464in}}%
\pgfpathlineto{\pgfqpoint{11.067444in}{0.721464in}}%
\pgfpathlineto{\pgfqpoint{11.067444in}{0.712800in}}%
\pgfpathlineto{\pgfqpoint{11.056817in}{0.712800in}}%
\pgfpathquadraticcurveto{\pgfqpoint{11.044821in}{0.712800in}}{\pgfqpoint{11.040264in}{0.717267in}}%
\pgfpathquadraticcurveto{\pgfqpoint{11.035707in}{0.721752in}}{\pgfqpoint{11.035707in}{0.733568in}}%
\pgfpathlineto{\pgfqpoint{11.035707in}{0.767791in}}%
\pgfpathlineto{\pgfqpoint{11.028106in}{0.767791in}}%
\pgfpathlineto{\pgfqpoint{11.028106in}{0.775843in}}%
\pgfpathlineto{\pgfqpoint{11.035707in}{0.775843in}}%
\pgfpathlineto{\pgfqpoint{11.035707in}{0.793747in}}%
\pgfpathlineto{\pgfqpoint{11.046118in}{0.793747in}}%
\pgfpathlineto{\pgfqpoint{11.046118in}{0.793747in}}%
\pgfpathclose%
\pgfpathmoveto{\pgfqpoint{11.105486in}{0.768584in}}%
\pgfpathquadraticcurveto{\pgfqpoint{11.097165in}{0.768584in}}{\pgfqpoint{11.092319in}{0.762081in}}%
\pgfpathquadraticcurveto{\pgfqpoint{11.087474in}{0.755579in}}{\pgfqpoint{11.087474in}{0.744267in}}%
\pgfpathquadraticcurveto{\pgfqpoint{11.087474in}{0.732956in}}{\pgfqpoint{11.092283in}{0.726453in}}%
\pgfpathquadraticcurveto{\pgfqpoint{11.097110in}{0.719951in}}{\pgfqpoint{11.105486in}{0.719951in}}%
\pgfpathquadraticcurveto{\pgfqpoint{11.113772in}{0.719951in}}{\pgfqpoint{11.118599in}{0.726471in}}%
\pgfpathquadraticcurveto{\pgfqpoint{11.123444in}{0.733010in}}{\pgfqpoint{11.123444in}{0.744267in}}%
\pgfpathquadraticcurveto{\pgfqpoint{11.123444in}{0.755471in}}{\pgfqpoint{11.118599in}{0.762027in}}%
\pgfpathquadraticcurveto{\pgfqpoint{11.113772in}{0.768584in}}{\pgfqpoint{11.105486in}{0.768584in}}%
\pgfpathlineto{\pgfqpoint{11.105486in}{0.768584in}}%
\pgfpathclose%
\pgfpathmoveto{\pgfqpoint{11.105486in}{0.777356in}}%
\pgfpathquadraticcurveto{\pgfqpoint{11.118995in}{0.777356in}}{\pgfqpoint{11.126704in}{0.768566in}}%
\pgfpathquadraticcurveto{\pgfqpoint{11.134432in}{0.759794in}}{\pgfqpoint{11.134432in}{0.744267in}}%
\pgfpathquadraticcurveto{\pgfqpoint{11.134432in}{0.728795in}}{\pgfqpoint{11.126704in}{0.719969in}}%
\pgfpathquadraticcurveto{\pgfqpoint{11.118995in}{0.711161in}}{\pgfqpoint{11.105486in}{0.711161in}}%
\pgfpathquadraticcurveto{\pgfqpoint{11.091923in}{0.711161in}}{\pgfqpoint{11.084232in}{0.719969in}}%
\pgfpathquadraticcurveto{\pgfqpoint{11.076559in}{0.728795in}}{\pgfqpoint{11.076559in}{0.744267in}}%
\pgfpathquadraticcurveto{\pgfqpoint{11.076559in}{0.759794in}}{\pgfqpoint{11.084232in}{0.768566in}}%
\pgfpathquadraticcurveto{\pgfqpoint{11.091923in}{0.777356in}}{\pgfqpoint{11.105486in}{0.777356in}}%
\pgfpathlineto{\pgfqpoint{11.105486in}{0.777356in}}%
\pgfpathclose%
\pgfpathmoveto{\pgfqpoint{11.188126in}{0.766170in}}%
\pgfpathquadraticcurveto{\pgfqpoint{11.186379in}{0.767179in}}{\pgfqpoint{11.184325in}{0.767647in}}%
\pgfpathquadraticcurveto{\pgfqpoint{11.182272in}{0.768133in}}{\pgfqpoint{11.179804in}{0.768133in}}%
\pgfpathquadraticcurveto{\pgfqpoint{11.171014in}{0.768133in}}{\pgfqpoint{11.166313in}{0.762423in}}%
\pgfpathquadraticcurveto{\pgfqpoint{11.161612in}{0.756714in}}{\pgfqpoint{11.161612in}{0.746014in}}%
\pgfpathlineto{\pgfqpoint{11.161612in}{0.712800in}}%
\pgfpathlineto{\pgfqpoint{11.151201in}{0.712800in}}%
\pgfpathlineto{\pgfqpoint{11.151201in}{0.775843in}}%
\pgfpathlineto{\pgfqpoint{11.161612in}{0.775843in}}%
\pgfpathlineto{\pgfqpoint{11.161612in}{0.766044in}}%
\pgfpathquadraticcurveto{\pgfqpoint{11.164890in}{0.771790in}}{\pgfqpoint{11.170114in}{0.774564in}}%
\pgfpathquadraticcurveto{\pgfqpoint{11.175355in}{0.777356in}}{\pgfqpoint{11.182848in}{0.777356in}}%
\pgfpathquadraticcurveto{\pgfqpoint{11.183911in}{0.777356in}}{\pgfqpoint{11.185208in}{0.777211in}}%
\pgfpathquadraticcurveto{\pgfqpoint{11.186505in}{0.777085in}}{\pgfqpoint{11.188072in}{0.776797in}}%
\pgfpathlineto{\pgfqpoint{11.188126in}{0.766170in}}%
\pgfpathlineto{\pgfqpoint{11.188126in}{0.766170in}}%
\pgfpathclose%
\pgfpathmoveto{\pgfqpoint{11.239172in}{0.773987in}}%
\pgfpathlineto{\pgfqpoint{11.239172in}{0.764189in}}%
\pgfpathquadraticcurveto{\pgfqpoint{11.234795in}{0.766440in}}{\pgfqpoint{11.230058in}{0.767557in}}%
\pgfpathquadraticcurveto{\pgfqpoint{11.225339in}{0.768692in}}{\pgfqpoint{11.220260in}{0.768692in}}%
\pgfpathquadraticcurveto{\pgfqpoint{11.212550in}{0.768692in}}{\pgfqpoint{11.208696in}{0.766332in}}%
\pgfpathquadraticcurveto{\pgfqpoint{11.204841in}{0.763973in}}{\pgfqpoint{11.204841in}{0.759235in}}%
\pgfpathquadraticcurveto{\pgfqpoint{11.204841in}{0.755633in}}{\pgfqpoint{11.207597in}{0.753580in}}%
\pgfpathquadraticcurveto{\pgfqpoint{11.210353in}{0.751526in}}{\pgfqpoint{11.218692in}{0.749671in}}%
\pgfpathlineto{\pgfqpoint{11.222241in}{0.748878in}}%
\pgfpathquadraticcurveto{\pgfqpoint{11.233264in}{0.746519in}}{\pgfqpoint{11.237911in}{0.742214in}}%
\pgfpathquadraticcurveto{\pgfqpoint{11.242559in}{0.737909in}}{\pgfqpoint{11.242559in}{0.730200in}}%
\pgfpathquadraticcurveto{\pgfqpoint{11.242559in}{0.721410in}}{\pgfqpoint{11.235606in}{0.716276in}}%
\pgfpathquadraticcurveto{\pgfqpoint{11.228653in}{0.711161in}}{\pgfqpoint{11.216495in}{0.711161in}}%
\pgfpathquadraticcurveto{\pgfqpoint{11.211434in}{0.711161in}}{\pgfqpoint{11.205940in}{0.712152in}}%
\pgfpathquadraticcurveto{\pgfqpoint{11.200446in}{0.713142in}}{\pgfqpoint{11.194376in}{0.715106in}}%
\pgfpathlineto{\pgfqpoint{11.194376in}{0.725805in}}%
\pgfpathquadraticcurveto{\pgfqpoint{11.200122in}{0.722815in}}{\pgfqpoint{11.205688in}{0.721320in}}%
\pgfpathquadraticcurveto{\pgfqpoint{11.211253in}{0.719843in}}{\pgfqpoint{11.216729in}{0.719843in}}%
\pgfpathquadraticcurveto{\pgfqpoint{11.224042in}{0.719843in}}{\pgfqpoint{11.227969in}{0.722346in}}%
\pgfpathquadraticcurveto{\pgfqpoint{11.231913in}{0.724850in}}{\pgfqpoint{11.231913in}{0.729407in}}%
\pgfpathquadraticcurveto{\pgfqpoint{11.231913in}{0.733622in}}{\pgfqpoint{11.229067in}{0.735874in}}%
\pgfpathquadraticcurveto{\pgfqpoint{11.226240in}{0.738125in}}{\pgfqpoint{11.216603in}{0.740214in}}%
\pgfpathlineto{\pgfqpoint{11.213001in}{0.741061in}}%
\pgfpathquadraticcurveto{\pgfqpoint{11.203382in}{0.743078in}}{\pgfqpoint{11.199095in}{0.747275in}}%
\pgfpathquadraticcurveto{\pgfqpoint{11.194826in}{0.751472in}}{\pgfqpoint{11.194826in}{0.758785in}}%
\pgfpathquadraticcurveto{\pgfqpoint{11.194826in}{0.767683in}}{\pgfqpoint{11.201131in}{0.772510in}}%
\pgfpathquadraticcurveto{\pgfqpoint{11.207435in}{0.777356in}}{\pgfqpoint{11.219035in}{0.777356in}}%
\pgfpathquadraticcurveto{\pgfqpoint{11.224763in}{0.777356in}}{\pgfqpoint{11.229824in}{0.776509in}}%
\pgfpathquadraticcurveto{\pgfqpoint{11.234903in}{0.775680in}}{\pgfqpoint{11.239172in}{0.773987in}}%
\pgfpathlineto{\pgfqpoint{11.239172in}{0.773987in}}%
\pgfpathclose%
\pgfusepath{fill}%
\end{pgfscope}%
\begin{pgfscope}%
\pgfsetbuttcap%
\pgfsetroundjoin%
\definecolor{currentfill}{rgb}{0.309804,0.372549,0.419608}%
\pgfsetfillcolor{currentfill}%
\pgfsetlinewidth{0.000000pt}%
\definecolor{currentstroke}{rgb}{0.309804,0.372549,0.419608}%
\pgfsetstrokecolor{currentstroke}%
\pgfsetdash{}{0pt}%
\pgfpathmoveto{\pgfqpoint{11.315365in}{0.737675in}}%
\pgfpathlineto{\pgfqpoint{11.315365in}{0.775843in}}%
\pgfpathlineto{\pgfqpoint{11.325722in}{0.775843in}}%
\pgfpathlineto{\pgfqpoint{11.325722in}{0.738071in}}%
\pgfpathquadraticcurveto{\pgfqpoint{11.325722in}{0.729119in}}{\pgfqpoint{11.329199in}{0.724634in}}%
\pgfpathquadraticcurveto{\pgfqpoint{11.332693in}{0.720167in}}{\pgfqpoint{11.339682in}{0.720167in}}%
\pgfpathquadraticcurveto{\pgfqpoint{11.348057in}{0.720167in}}{\pgfqpoint{11.352921in}{0.725517in}}%
\pgfpathquadraticcurveto{\pgfqpoint{11.357802in}{0.730866in}}{\pgfqpoint{11.357802in}{0.740106in}}%
\pgfpathlineto{\pgfqpoint{11.357802in}{0.775843in}}%
\pgfpathlineto{\pgfqpoint{11.368159in}{0.775843in}}%
\pgfpathlineto{\pgfqpoint{11.368159in}{0.712800in}}%
\pgfpathlineto{\pgfqpoint{11.357802in}{0.712800in}}%
\pgfpathlineto{\pgfqpoint{11.357802in}{0.722491in}}%
\pgfpathquadraticcurveto{\pgfqpoint{11.354038in}{0.716745in}}{\pgfqpoint{11.349048in}{0.713953in}}%
\pgfpathquadraticcurveto{\pgfqpoint{11.344077in}{0.711161in}}{\pgfqpoint{11.337484in}{0.711161in}}%
\pgfpathquadraticcurveto{\pgfqpoint{11.326623in}{0.711161in}}{\pgfqpoint{11.320985in}{0.717915in}}%
\pgfpathquadraticcurveto{\pgfqpoint{11.315365in}{0.724670in}}{\pgfqpoint{11.315365in}{0.737675in}}%
\pgfpathlineto{\pgfqpoint{11.315365in}{0.737675in}}%
\pgfpathclose%
\pgfpathmoveto{\pgfqpoint{11.341429in}{0.777356in}}%
\pgfpathlineto{\pgfqpoint{11.341429in}{0.777356in}}%
\pgfpathlineto{\pgfqpoint{11.341429in}{0.777356in}}%
\pgfpathclose%
\pgfpathmoveto{\pgfqpoint{11.429671in}{0.773987in}}%
\pgfpathlineto{\pgfqpoint{11.429671in}{0.764189in}}%
\pgfpathquadraticcurveto{\pgfqpoint{11.425294in}{0.766440in}}{\pgfqpoint{11.420556in}{0.767557in}}%
\pgfpathquadraticcurveto{\pgfqpoint{11.415837in}{0.768692in}}{\pgfqpoint{11.410758in}{0.768692in}}%
\pgfpathquadraticcurveto{\pgfqpoint{11.403049in}{0.768692in}}{\pgfqpoint{11.399194in}{0.766332in}}%
\pgfpathquadraticcurveto{\pgfqpoint{11.395339in}{0.763973in}}{\pgfqpoint{11.395339in}{0.759235in}}%
\pgfpathquadraticcurveto{\pgfqpoint{11.395339in}{0.755633in}}{\pgfqpoint{11.398095in}{0.753580in}}%
\pgfpathquadraticcurveto{\pgfqpoint{11.400851in}{0.751526in}}{\pgfqpoint{11.409191in}{0.749671in}}%
\pgfpathlineto{\pgfqpoint{11.412739in}{0.748878in}}%
\pgfpathquadraticcurveto{\pgfqpoint{11.423763in}{0.746519in}}{\pgfqpoint{11.428410in}{0.742214in}}%
\pgfpathquadraticcurveto{\pgfqpoint{11.433057in}{0.737909in}}{\pgfqpoint{11.433057in}{0.730200in}}%
\pgfpathquadraticcurveto{\pgfqpoint{11.433057in}{0.721410in}}{\pgfqpoint{11.426104in}{0.716276in}}%
\pgfpathquadraticcurveto{\pgfqpoint{11.419151in}{0.711161in}}{\pgfqpoint{11.406993in}{0.711161in}}%
\pgfpathquadraticcurveto{\pgfqpoint{11.401932in}{0.711161in}}{\pgfqpoint{11.396438in}{0.712152in}}%
\pgfpathquadraticcurveto{\pgfqpoint{11.390944in}{0.713142in}}{\pgfqpoint{11.384874in}{0.715106in}}%
\pgfpathlineto{\pgfqpoint{11.384874in}{0.725805in}}%
\pgfpathquadraticcurveto{\pgfqpoint{11.390620in}{0.722815in}}{\pgfqpoint{11.396186in}{0.721320in}}%
\pgfpathquadraticcurveto{\pgfqpoint{11.401752in}{0.719843in}}{\pgfqpoint{11.407227in}{0.719843in}}%
\pgfpathquadraticcurveto{\pgfqpoint{11.414540in}{0.719843in}}{\pgfqpoint{11.418467in}{0.722346in}}%
\pgfpathquadraticcurveto{\pgfqpoint{11.422412in}{0.724850in}}{\pgfqpoint{11.422412in}{0.729407in}}%
\pgfpathquadraticcurveto{\pgfqpoint{11.422412in}{0.733622in}}{\pgfqpoint{11.419566in}{0.735874in}}%
\pgfpathquadraticcurveto{\pgfqpoint{11.416738in}{0.738125in}}{\pgfqpoint{11.407101in}{0.740214in}}%
\pgfpathlineto{\pgfqpoint{11.403499in}{0.741061in}}%
\pgfpathquadraticcurveto{\pgfqpoint{11.393880in}{0.743078in}}{\pgfqpoint{11.389594in}{0.747275in}}%
\pgfpathquadraticcurveto{\pgfqpoint{11.385325in}{0.751472in}}{\pgfqpoint{11.385325in}{0.758785in}}%
\pgfpathquadraticcurveto{\pgfqpoint{11.385325in}{0.767683in}}{\pgfqpoint{11.391629in}{0.772510in}}%
\pgfpathquadraticcurveto{\pgfqpoint{11.397933in}{0.777356in}}{\pgfqpoint{11.409533in}{0.777356in}}%
\pgfpathquadraticcurveto{\pgfqpoint{11.415261in}{0.777356in}}{\pgfqpoint{11.420322in}{0.776509in}}%
\pgfpathquadraticcurveto{\pgfqpoint{11.425402in}{0.775680in}}{\pgfqpoint{11.429671in}{0.773987in}}%
\pgfpathlineto{\pgfqpoint{11.429671in}{0.773987in}}%
\pgfpathclose%
\pgfpathmoveto{\pgfqpoint{11.503466in}{0.746915in}}%
\pgfpathlineto{\pgfqpoint{11.503466in}{0.741854in}}%
\pgfpathlineto{\pgfqpoint{11.455842in}{0.741854in}}%
\pgfpathquadraticcurveto{\pgfqpoint{11.456527in}{0.731154in}}{\pgfqpoint{11.462291in}{0.725553in}}%
\pgfpathquadraticcurveto{\pgfqpoint{11.468054in}{0.719951in}}{\pgfqpoint{11.478357in}{0.719951in}}%
\pgfpathquadraticcurveto{\pgfqpoint{11.484319in}{0.719951in}}{\pgfqpoint{11.489921in}{0.721410in}}%
\pgfpathquadraticcurveto{\pgfqpoint{11.495523in}{0.722869in}}{\pgfqpoint{11.501053in}{0.725805in}}%
\pgfpathlineto{\pgfqpoint{11.501053in}{0.716006in}}%
\pgfpathquadraticcurveto{\pgfqpoint{11.495469in}{0.713647in}}{\pgfqpoint{11.489615in}{0.712404in}}%
\pgfpathquadraticcurveto{\pgfqpoint{11.483761in}{0.711161in}}{\pgfqpoint{11.477745in}{0.711161in}}%
\pgfpathquadraticcurveto{\pgfqpoint{11.462651in}{0.711161in}}{\pgfqpoint{11.453843in}{0.719933in}}%
\pgfpathquadraticcurveto{\pgfqpoint{11.445035in}{0.728723in}}{\pgfqpoint{11.445035in}{0.743709in}}%
\pgfpathquadraticcurveto{\pgfqpoint{11.445035in}{0.759181in}}{\pgfqpoint{11.453393in}{0.768259in}}%
\pgfpathquadraticcurveto{\pgfqpoint{11.461750in}{0.777356in}}{\pgfqpoint{11.475944in}{0.777356in}}%
\pgfpathquadraticcurveto{\pgfqpoint{11.488660in}{0.777356in}}{\pgfqpoint{11.496063in}{0.769160in}}%
\pgfpathquadraticcurveto{\pgfqpoint{11.503466in}{0.760983in}}{\pgfqpoint{11.503466in}{0.746915in}}%
\pgfpathlineto{\pgfqpoint{11.503466in}{0.746915in}}%
\pgfpathclose%
\pgfpathmoveto{\pgfqpoint{11.493109in}{0.749959in}}%
\pgfpathquadraticcurveto{\pgfqpoint{11.493001in}{0.758443in}}{\pgfqpoint{11.488354in}{0.763504in}}%
\pgfpathquadraticcurveto{\pgfqpoint{11.483707in}{0.768584in}}{\pgfqpoint{11.476052in}{0.768584in}}%
\pgfpathquadraticcurveto{\pgfqpoint{11.467388in}{0.768584in}}{\pgfqpoint{11.462182in}{0.763684in}}%
\pgfpathquadraticcurveto{\pgfqpoint{11.456977in}{0.758785in}}{\pgfqpoint{11.456184in}{0.749887in}}%
\pgfpathlineto{\pgfqpoint{11.493109in}{0.749959in}}%
\pgfpathlineto{\pgfqpoint{11.493109in}{0.749959in}}%
\pgfpathclose%
\pgfusepath{fill}%
\end{pgfscope}%
\begin{pgfscope}%
\pgfsetbuttcap%
\pgfsetroundjoin%
\definecolor{currentfill}{rgb}{0.309804,0.372549,0.419608}%
\pgfsetfillcolor{currentfill}%
\pgfsetlinewidth{0.000000pt}%
\definecolor{currentstroke}{rgb}{0.309804,0.372549,0.419608}%
\pgfsetstrokecolor{currentstroke}%
\pgfsetdash{}{0pt}%
\pgfpathmoveto{\pgfqpoint{11.578370in}{0.793747in}}%
\pgfpathlineto{\pgfqpoint{11.578370in}{0.775843in}}%
\pgfpathlineto{\pgfqpoint{11.599696in}{0.775843in}}%
\pgfpathlineto{\pgfqpoint{11.599696in}{0.767791in}}%
\pgfpathlineto{\pgfqpoint{11.578370in}{0.767791in}}%
\pgfpathlineto{\pgfqpoint{11.578370in}{0.733568in}}%
\pgfpathquadraticcurveto{\pgfqpoint{11.578370in}{0.725859in}}{\pgfqpoint{11.580477in}{0.723661in}}%
\pgfpathquadraticcurveto{\pgfqpoint{11.582585in}{0.721464in}}{\pgfqpoint{11.589069in}{0.721464in}}%
\pgfpathlineto{\pgfqpoint{11.599696in}{0.721464in}}%
\pgfpathlineto{\pgfqpoint{11.599696in}{0.712800in}}%
\pgfpathlineto{\pgfqpoint{11.589069in}{0.712800in}}%
\pgfpathquadraticcurveto{\pgfqpoint{11.577073in}{0.712800in}}{\pgfqpoint{11.572516in}{0.717267in}}%
\pgfpathquadraticcurveto{\pgfqpoint{11.567959in}{0.721752in}}{\pgfqpoint{11.567959in}{0.733568in}}%
\pgfpathlineto{\pgfqpoint{11.567959in}{0.767791in}}%
\pgfpathlineto{\pgfqpoint{11.560358in}{0.767791in}}%
\pgfpathlineto{\pgfqpoint{11.560358in}{0.775843in}}%
\pgfpathlineto{\pgfqpoint{11.567959in}{0.775843in}}%
\pgfpathlineto{\pgfqpoint{11.567959in}{0.793747in}}%
\pgfpathlineto{\pgfqpoint{11.578370in}{0.793747in}}%
\pgfpathlineto{\pgfqpoint{11.578370in}{0.793747in}}%
\pgfpathclose%
\pgfpathmoveto{\pgfqpoint{11.665729in}{0.750860in}}%
\pgfpathlineto{\pgfqpoint{11.665729in}{0.712800in}}%
\pgfpathlineto{\pgfqpoint{11.655372in}{0.712800in}}%
\pgfpathlineto{\pgfqpoint{11.655372in}{0.750517in}}%
\pgfpathquadraticcurveto{\pgfqpoint{11.655372in}{0.759469in}}{\pgfqpoint{11.651878in}{0.763900in}}%
\pgfpathquadraticcurveto{\pgfqpoint{11.648383in}{0.768349in}}{\pgfqpoint{11.641412in}{0.768349in}}%
\pgfpathquadraticcurveto{\pgfqpoint{11.633019in}{0.768349in}}{\pgfqpoint{11.628174in}{0.763000in}}%
\pgfpathquadraticcurveto{\pgfqpoint{11.623328in}{0.757668in}}{\pgfqpoint{11.623328in}{0.748428in}}%
\pgfpathlineto{\pgfqpoint{11.623328in}{0.712800in}}%
\pgfpathlineto{\pgfqpoint{11.612917in}{0.712800in}}%
\pgfpathlineto{\pgfqpoint{11.612917in}{0.800393in}}%
\pgfpathlineto{\pgfqpoint{11.623328in}{0.800393in}}%
\pgfpathlineto{\pgfqpoint{11.623328in}{0.766044in}}%
\pgfpathquadraticcurveto{\pgfqpoint{11.627057in}{0.771736in}}{\pgfqpoint{11.632082in}{0.774546in}}%
\pgfpathquadraticcurveto{\pgfqpoint{11.637126in}{0.777356in}}{\pgfqpoint{11.643718in}{0.777356in}}%
\pgfpathquadraticcurveto{\pgfqpoint{11.654579in}{0.777356in}}{\pgfqpoint{11.660145in}{0.770637in}}%
\pgfpathquadraticcurveto{\pgfqpoint{11.665729in}{0.763918in}}{\pgfqpoint{11.665729in}{0.750860in}}%
\pgfpathlineto{\pgfqpoint{11.665729in}{0.750860in}}%
\pgfpathclose%
\pgfpathmoveto{\pgfqpoint{11.740299in}{0.746915in}}%
\pgfpathlineto{\pgfqpoint{11.740299in}{0.741854in}}%
\pgfpathlineto{\pgfqpoint{11.692675in}{0.741854in}}%
\pgfpathquadraticcurveto{\pgfqpoint{11.693360in}{0.731154in}}{\pgfqpoint{11.699123in}{0.725553in}}%
\pgfpathquadraticcurveto{\pgfqpoint{11.704887in}{0.719951in}}{\pgfqpoint{11.715190in}{0.719951in}}%
\pgfpathquadraticcurveto{\pgfqpoint{11.721152in}{0.719951in}}{\pgfqpoint{11.726754in}{0.721410in}}%
\pgfpathquadraticcurveto{\pgfqpoint{11.732356in}{0.722869in}}{\pgfqpoint{11.737886in}{0.725805in}}%
\pgfpathlineto{\pgfqpoint{11.737886in}{0.716006in}}%
\pgfpathquadraticcurveto{\pgfqpoint{11.732302in}{0.713647in}}{\pgfqpoint{11.726448in}{0.712404in}}%
\pgfpathquadraticcurveto{\pgfqpoint{11.720594in}{0.711161in}}{\pgfqpoint{11.714578in}{0.711161in}}%
\pgfpathquadraticcurveto{\pgfqpoint{11.699484in}{0.711161in}}{\pgfqpoint{11.690676in}{0.719933in}}%
\pgfpathquadraticcurveto{\pgfqpoint{11.681868in}{0.728723in}}{\pgfqpoint{11.681868in}{0.743709in}}%
\pgfpathquadraticcurveto{\pgfqpoint{11.681868in}{0.759181in}}{\pgfqpoint{11.690225in}{0.768259in}}%
\pgfpathquadraticcurveto{\pgfqpoint{11.698583in}{0.777356in}}{\pgfqpoint{11.712777in}{0.777356in}}%
\pgfpathquadraticcurveto{\pgfqpoint{11.725493in}{0.777356in}}{\pgfqpoint{11.732896in}{0.769160in}}%
\pgfpathquadraticcurveto{\pgfqpoint{11.740299in}{0.760983in}}{\pgfqpoint{11.740299in}{0.746915in}}%
\pgfpathlineto{\pgfqpoint{11.740299in}{0.746915in}}%
\pgfpathclose%
\pgfpathmoveto{\pgfqpoint{11.729942in}{0.749959in}}%
\pgfpathquadraticcurveto{\pgfqpoint{11.729834in}{0.758443in}}{\pgfqpoint{11.725187in}{0.763504in}}%
\pgfpathquadraticcurveto{\pgfqpoint{11.720540in}{0.768584in}}{\pgfqpoint{11.712885in}{0.768584in}}%
\pgfpathquadraticcurveto{\pgfqpoint{11.704221in}{0.768584in}}{\pgfqpoint{11.699015in}{0.763684in}}%
\pgfpathquadraticcurveto{\pgfqpoint{11.693810in}{0.758785in}}{\pgfqpoint{11.693017in}{0.749887in}}%
\pgfpathlineto{\pgfqpoint{11.729942in}{0.749959in}}%
\pgfpathlineto{\pgfqpoint{11.729942in}{0.749959in}}%
\pgfpathclose%
\pgfpathmoveto{\pgfqpoint{11.757303in}{0.775843in}}%
\pgfpathlineto{\pgfqpoint{11.767660in}{0.775843in}}%
\pgfpathlineto{\pgfqpoint{11.767660in}{0.712800in}}%
\pgfpathlineto{\pgfqpoint{11.757303in}{0.712800in}}%
\pgfpathlineto{\pgfqpoint{11.757303in}{0.775843in}}%
\pgfpathlineto{\pgfqpoint{11.757303in}{0.775843in}}%
\pgfpathclose%
\pgfpathmoveto{\pgfqpoint{11.757303in}{0.800393in}}%
\pgfpathlineto{\pgfqpoint{11.767660in}{0.800393in}}%
\pgfpathlineto{\pgfqpoint{11.767660in}{0.787262in}}%
\pgfpathlineto{\pgfqpoint{11.757303in}{0.787262in}}%
\pgfpathlineto{\pgfqpoint{11.757303in}{0.800393in}}%
\pgfpathlineto{\pgfqpoint{11.757303in}{0.800393in}}%
\pgfpathclose%
\pgfpathmoveto{\pgfqpoint{11.825857in}{0.766170in}}%
\pgfpathquadraticcurveto{\pgfqpoint{11.824110in}{0.767179in}}{\pgfqpoint{11.822056in}{0.767647in}}%
\pgfpathquadraticcurveto{\pgfqpoint{11.820003in}{0.768133in}}{\pgfqpoint{11.817535in}{0.768133in}}%
\pgfpathquadraticcurveto{\pgfqpoint{11.808745in}{0.768133in}}{\pgfqpoint{11.804044in}{0.762423in}}%
\pgfpathquadraticcurveto{\pgfqpoint{11.799343in}{0.756714in}}{\pgfqpoint{11.799343in}{0.746014in}}%
\pgfpathlineto{\pgfqpoint{11.799343in}{0.712800in}}%
\pgfpathlineto{\pgfqpoint{11.788932in}{0.712800in}}%
\pgfpathlineto{\pgfqpoint{11.788932in}{0.775843in}}%
\pgfpathlineto{\pgfqpoint{11.799343in}{0.775843in}}%
\pgfpathlineto{\pgfqpoint{11.799343in}{0.766044in}}%
\pgfpathquadraticcurveto{\pgfqpoint{11.802621in}{0.771790in}}{\pgfqpoint{11.807845in}{0.774564in}}%
\pgfpathquadraticcurveto{\pgfqpoint{11.813086in}{0.777356in}}{\pgfqpoint{11.820579in}{0.777356in}}%
\pgfpathquadraticcurveto{\pgfqpoint{11.821642in}{0.777356in}}{\pgfqpoint{11.822939in}{0.777211in}}%
\pgfpathquadraticcurveto{\pgfqpoint{11.824236in}{0.777085in}}{\pgfqpoint{11.825803in}{0.776797in}}%
\pgfpathlineto{\pgfqpoint{11.825857in}{0.766170in}}%
\pgfpathlineto{\pgfqpoint{11.825857in}{0.766170in}}%
\pgfpathclose%
\pgfusepath{fill}%
\end{pgfscope}%
\begin{pgfscope}%
\pgfsetbuttcap%
\pgfsetroundjoin%
\definecolor{currentfill}{rgb}{0.309804,0.372549,0.419608}%
\pgfsetfillcolor{currentfill}%
\pgfsetlinewidth{0.000000pt}%
\definecolor{currentstroke}{rgb}{0.309804,0.372549,0.419608}%
\pgfsetstrokecolor{currentstroke}%
\pgfsetdash{}{0pt}%
\pgfpathmoveto{\pgfqpoint{11.914238in}{0.768584in}}%
\pgfpathquadraticcurveto{\pgfqpoint{11.905917in}{0.768584in}}{\pgfqpoint{11.901072in}{0.762081in}}%
\pgfpathquadraticcurveto{\pgfqpoint{11.896226in}{0.755579in}}{\pgfqpoint{11.896226in}{0.744267in}}%
\pgfpathquadraticcurveto{\pgfqpoint{11.896226in}{0.732956in}}{\pgfqpoint{11.901035in}{0.726453in}}%
\pgfpathquadraticcurveto{\pgfqpoint{11.905863in}{0.719951in}}{\pgfqpoint{11.914238in}{0.719951in}}%
\pgfpathquadraticcurveto{\pgfqpoint{11.922524in}{0.719951in}}{\pgfqpoint{11.927351in}{0.726471in}}%
\pgfpathquadraticcurveto{\pgfqpoint{11.932197in}{0.733010in}}{\pgfqpoint{11.932197in}{0.744267in}}%
\pgfpathquadraticcurveto{\pgfqpoint{11.932197in}{0.755471in}}{\pgfqpoint{11.927351in}{0.762027in}}%
\pgfpathquadraticcurveto{\pgfqpoint{11.922524in}{0.768584in}}{\pgfqpoint{11.914238in}{0.768584in}}%
\pgfpathlineto{\pgfqpoint{11.914238in}{0.768584in}}%
\pgfpathclose%
\pgfpathmoveto{\pgfqpoint{11.914238in}{0.777356in}}%
\pgfpathquadraticcurveto{\pgfqpoint{11.927748in}{0.777356in}}{\pgfqpoint{11.935457in}{0.768566in}}%
\pgfpathquadraticcurveto{\pgfqpoint{11.943184in}{0.759794in}}{\pgfqpoint{11.943184in}{0.744267in}}%
\pgfpathquadraticcurveto{\pgfqpoint{11.943184in}{0.728795in}}{\pgfqpoint{11.935457in}{0.719969in}}%
\pgfpathquadraticcurveto{\pgfqpoint{11.927748in}{0.711161in}}{\pgfqpoint{11.914238in}{0.711161in}}%
\pgfpathquadraticcurveto{\pgfqpoint{11.900675in}{0.711161in}}{\pgfqpoint{11.892984in}{0.719969in}}%
\pgfpathquadraticcurveto{\pgfqpoint{11.885311in}{0.728795in}}{\pgfqpoint{11.885311in}{0.744267in}}%
\pgfpathquadraticcurveto{\pgfqpoint{11.885311in}{0.759794in}}{\pgfqpoint{11.892984in}{0.768566in}}%
\pgfpathquadraticcurveto{\pgfqpoint{11.900675in}{0.777356in}}{\pgfqpoint{11.914238in}{0.777356in}}%
\pgfpathlineto{\pgfqpoint{11.914238in}{0.777356in}}%
\pgfpathclose%
\pgfpathmoveto{\pgfqpoint{11.954333in}{0.775843in}}%
\pgfpathlineto{\pgfqpoint{11.964690in}{0.775843in}}%
\pgfpathlineto{\pgfqpoint{11.977641in}{0.726651in}}%
\pgfpathlineto{\pgfqpoint{11.990520in}{0.775843in}}%
\pgfpathlineto{\pgfqpoint{12.002732in}{0.775843in}}%
\pgfpathlineto{\pgfqpoint{12.015683in}{0.726651in}}%
\pgfpathlineto{\pgfqpoint{12.028580in}{0.775843in}}%
\pgfpathlineto{\pgfqpoint{12.038937in}{0.775843in}}%
\pgfpathlineto{\pgfqpoint{12.022437in}{0.712800in}}%
\pgfpathlineto{\pgfqpoint{12.010225in}{0.712800in}}%
\pgfpathlineto{\pgfqpoint{11.996662in}{0.764477in}}%
\pgfpathlineto{\pgfqpoint{11.983045in}{0.712800in}}%
\pgfpathlineto{\pgfqpoint{11.970815in}{0.712800in}}%
\pgfpathlineto{\pgfqpoint{11.954333in}{0.775843in}}%
\pgfpathlineto{\pgfqpoint{11.954333in}{0.775843in}}%
\pgfpathclose%
\pgfpathmoveto{\pgfqpoint{12.107040in}{0.750860in}}%
\pgfpathlineto{\pgfqpoint{12.107040in}{0.712800in}}%
\pgfpathlineto{\pgfqpoint{12.096683in}{0.712800in}}%
\pgfpathlineto{\pgfqpoint{12.096683in}{0.750517in}}%
\pgfpathquadraticcurveto{\pgfqpoint{12.096683in}{0.759469in}}{\pgfqpoint{12.093189in}{0.763900in}}%
\pgfpathquadraticcurveto{\pgfqpoint{12.089695in}{0.768349in}}{\pgfqpoint{12.082724in}{0.768349in}}%
\pgfpathquadraticcurveto{\pgfqpoint{12.074330in}{0.768349in}}{\pgfqpoint{12.069485in}{0.763000in}}%
\pgfpathquadraticcurveto{\pgfqpoint{12.064640in}{0.757668in}}{\pgfqpoint{12.064640in}{0.748428in}}%
\pgfpathlineto{\pgfqpoint{12.064640in}{0.712800in}}%
\pgfpathlineto{\pgfqpoint{12.054229in}{0.712800in}}%
\pgfpathlineto{\pgfqpoint{12.054229in}{0.775843in}}%
\pgfpathlineto{\pgfqpoint{12.064640in}{0.775843in}}%
\pgfpathlineto{\pgfqpoint{12.064640in}{0.766044in}}%
\pgfpathquadraticcurveto{\pgfqpoint{12.068368in}{0.771736in}}{\pgfqpoint{12.073394in}{0.774546in}}%
\pgfpathquadraticcurveto{\pgfqpoint{12.078437in}{0.777356in}}{\pgfqpoint{12.085030in}{0.777356in}}%
\pgfpathquadraticcurveto{\pgfqpoint{12.095891in}{0.777356in}}{\pgfqpoint{12.101457in}{0.770637in}}%
\pgfpathquadraticcurveto{\pgfqpoint{12.107040in}{0.763918in}}{\pgfqpoint{12.107040in}{0.750860in}}%
\pgfpathlineto{\pgfqpoint{12.107040in}{0.750860in}}%
\pgfpathclose%
\pgfusepath{fill}%
\end{pgfscope}%
\begin{pgfscope}%
\pgfsetbuttcap%
\pgfsetroundjoin%
\definecolor{currentfill}{rgb}{0.309804,0.372549,0.419608}%
\pgfsetfillcolor{currentfill}%
\pgfsetlinewidth{0.000000pt}%
\definecolor{currentstroke}{rgb}{0.309804,0.372549,0.419608}%
\pgfsetstrokecolor{currentstroke}%
\pgfsetdash{}{0pt}%
\pgfpathmoveto{\pgfqpoint{12.221692in}{0.773987in}}%
\pgfpathlineto{\pgfqpoint{12.221692in}{0.764189in}}%
\pgfpathquadraticcurveto{\pgfqpoint{12.217315in}{0.766440in}}{\pgfqpoint{12.212578in}{0.767557in}}%
\pgfpathquadraticcurveto{\pgfqpoint{12.207859in}{0.768692in}}{\pgfqpoint{12.202779in}{0.768692in}}%
\pgfpathquadraticcurveto{\pgfqpoint{12.195070in}{0.768692in}}{\pgfqpoint{12.191215in}{0.766332in}}%
\pgfpathquadraticcurveto{\pgfqpoint{12.187361in}{0.763973in}}{\pgfqpoint{12.187361in}{0.759235in}}%
\pgfpathquadraticcurveto{\pgfqpoint{12.187361in}{0.755633in}}{\pgfqpoint{12.190117in}{0.753580in}}%
\pgfpathquadraticcurveto{\pgfqpoint{12.192872in}{0.751526in}}{\pgfqpoint{12.201212in}{0.749671in}}%
\pgfpathlineto{\pgfqpoint{12.204760in}{0.748878in}}%
\pgfpathquadraticcurveto{\pgfqpoint{12.215784in}{0.746519in}}{\pgfqpoint{12.220431in}{0.742214in}}%
\pgfpathquadraticcurveto{\pgfqpoint{12.225078in}{0.737909in}}{\pgfqpoint{12.225078in}{0.730200in}}%
\pgfpathquadraticcurveto{\pgfqpoint{12.225078in}{0.721410in}}{\pgfqpoint{12.218126in}{0.716276in}}%
\pgfpathquadraticcurveto{\pgfqpoint{12.211173in}{0.711161in}}{\pgfqpoint{12.199015in}{0.711161in}}%
\pgfpathquadraticcurveto{\pgfqpoint{12.193953in}{0.711161in}}{\pgfqpoint{12.188459in}{0.712152in}}%
\pgfpathquadraticcurveto{\pgfqpoint{12.182966in}{0.713142in}}{\pgfqpoint{12.176896in}{0.715106in}}%
\pgfpathlineto{\pgfqpoint{12.176896in}{0.725805in}}%
\pgfpathquadraticcurveto{\pgfqpoint{12.182642in}{0.722815in}}{\pgfqpoint{12.188207in}{0.721320in}}%
\pgfpathquadraticcurveto{\pgfqpoint{12.193773in}{0.719843in}}{\pgfqpoint{12.199249in}{0.719843in}}%
\pgfpathquadraticcurveto{\pgfqpoint{12.206562in}{0.719843in}}{\pgfqpoint{12.210488in}{0.722346in}}%
\pgfpathquadraticcurveto{\pgfqpoint{12.214433in}{0.724850in}}{\pgfqpoint{12.214433in}{0.729407in}}%
\pgfpathquadraticcurveto{\pgfqpoint{12.214433in}{0.733622in}}{\pgfqpoint{12.211587in}{0.735874in}}%
\pgfpathquadraticcurveto{\pgfqpoint{12.208759in}{0.738125in}}{\pgfqpoint{12.199123in}{0.740214in}}%
\pgfpathlineto{\pgfqpoint{12.195520in}{0.741061in}}%
\pgfpathquadraticcurveto{\pgfqpoint{12.185902in}{0.743078in}}{\pgfqpoint{12.181615in}{0.747275in}}%
\pgfpathquadraticcurveto{\pgfqpoint{12.177346in}{0.751472in}}{\pgfqpoint{12.177346in}{0.758785in}}%
\pgfpathquadraticcurveto{\pgfqpoint{12.177346in}{0.767683in}}{\pgfqpoint{12.183650in}{0.772510in}}%
\pgfpathquadraticcurveto{\pgfqpoint{12.189955in}{0.777356in}}{\pgfqpoint{12.201554in}{0.777356in}}%
\pgfpathquadraticcurveto{\pgfqpoint{12.207282in}{0.777356in}}{\pgfqpoint{12.212344in}{0.776509in}}%
\pgfpathquadraticcurveto{\pgfqpoint{12.217423in}{0.775680in}}{\pgfqpoint{12.221692in}{0.773987in}}%
\pgfpathlineto{\pgfqpoint{12.221692in}{0.773987in}}%
\pgfpathclose%
\pgfpathmoveto{\pgfqpoint{12.240497in}{0.737675in}}%
\pgfpathlineto{\pgfqpoint{12.240497in}{0.775843in}}%
\pgfpathlineto{\pgfqpoint{12.250854in}{0.775843in}}%
\pgfpathlineto{\pgfqpoint{12.250854in}{0.738071in}}%
\pgfpathquadraticcurveto{\pgfqpoint{12.250854in}{0.729119in}}{\pgfqpoint{12.254330in}{0.724634in}}%
\pgfpathquadraticcurveto{\pgfqpoint{12.257824in}{0.720167in}}{\pgfqpoint{12.264813in}{0.720167in}}%
\pgfpathquadraticcurveto{\pgfqpoint{12.273189in}{0.720167in}}{\pgfqpoint{12.278052in}{0.725517in}}%
\pgfpathquadraticcurveto{\pgfqpoint{12.282933in}{0.730866in}}{\pgfqpoint{12.282933in}{0.740106in}}%
\pgfpathlineto{\pgfqpoint{12.282933in}{0.775843in}}%
\pgfpathlineto{\pgfqpoint{12.293290in}{0.775843in}}%
\pgfpathlineto{\pgfqpoint{12.293290in}{0.712800in}}%
\pgfpathlineto{\pgfqpoint{12.282933in}{0.712800in}}%
\pgfpathlineto{\pgfqpoint{12.282933in}{0.722491in}}%
\pgfpathquadraticcurveto{\pgfqpoint{12.279169in}{0.716745in}}{\pgfqpoint{12.274179in}{0.713953in}}%
\pgfpathquadraticcurveto{\pgfqpoint{12.269208in}{0.711161in}}{\pgfqpoint{12.262616in}{0.711161in}}%
\pgfpathquadraticcurveto{\pgfqpoint{12.251754in}{0.711161in}}{\pgfqpoint{12.246116in}{0.717915in}}%
\pgfpathquadraticcurveto{\pgfqpoint{12.240497in}{0.724670in}}{\pgfqpoint{12.240497in}{0.737675in}}%
\pgfpathlineto{\pgfqpoint{12.240497in}{0.737675in}}%
\pgfpathclose%
\pgfpathmoveto{\pgfqpoint{12.266560in}{0.777356in}}%
\pgfpathlineto{\pgfqpoint{12.266560in}{0.777356in}}%
\pgfpathlineto{\pgfqpoint{12.266560in}{0.777356in}}%
\pgfpathclose%
\pgfpathmoveto{\pgfqpoint{12.324631in}{0.722256in}}%
\pgfpathlineto{\pgfqpoint{12.324631in}{0.688826in}}%
\pgfpathlineto{\pgfqpoint{12.314220in}{0.688826in}}%
\pgfpathlineto{\pgfqpoint{12.314220in}{0.775843in}}%
\pgfpathlineto{\pgfqpoint{12.324631in}{0.775843in}}%
\pgfpathlineto{\pgfqpoint{12.324631in}{0.766278in}}%
\pgfpathquadraticcurveto{\pgfqpoint{12.327910in}{0.771898in}}{\pgfqpoint{12.332881in}{0.774618in}}%
\pgfpathquadraticcurveto{\pgfqpoint{12.337870in}{0.777356in}}{\pgfqpoint{12.344787in}{0.777356in}}%
\pgfpathquadraticcurveto{\pgfqpoint{12.356279in}{0.777356in}}{\pgfqpoint{12.363448in}{0.768241in}}%
\pgfpathquadraticcurveto{\pgfqpoint{12.370634in}{0.759127in}}{\pgfqpoint{12.370634in}{0.744267in}}%
\pgfpathquadraticcurveto{\pgfqpoint{12.370634in}{0.729407in}}{\pgfqpoint{12.363448in}{0.720275in}}%
\pgfpathquadraticcurveto{\pgfqpoint{12.356279in}{0.711161in}}{\pgfqpoint{12.344787in}{0.711161in}}%
\pgfpathquadraticcurveto{\pgfqpoint{12.337870in}{0.711161in}}{\pgfqpoint{12.332881in}{0.713899in}}%
\pgfpathquadraticcurveto{\pgfqpoint{12.327910in}{0.716637in}}{\pgfqpoint{12.324631in}{0.722256in}}%
\pgfpathlineto{\pgfqpoint{12.324631in}{0.722256in}}%
\pgfpathclose%
\pgfpathmoveto{\pgfqpoint{12.359881in}{0.744267in}}%
\pgfpathquadraticcurveto{\pgfqpoint{12.359881in}{0.755687in}}{\pgfqpoint{12.355180in}{0.762189in}}%
\pgfpathquadraticcurveto{\pgfqpoint{12.350479in}{0.768692in}}{\pgfqpoint{12.342265in}{0.768692in}}%
\pgfpathquadraticcurveto{\pgfqpoint{12.334034in}{0.768692in}}{\pgfqpoint{12.329333in}{0.762189in}}%
\pgfpathquadraticcurveto{\pgfqpoint{12.324631in}{0.755687in}}{\pgfqpoint{12.324631in}{0.744267in}}%
\pgfpathquadraticcurveto{\pgfqpoint{12.324631in}{0.732848in}}{\pgfqpoint{12.329333in}{0.726345in}}%
\pgfpathquadraticcurveto{\pgfqpoint{12.334034in}{0.719843in}}{\pgfqpoint{12.342265in}{0.719843in}}%
\pgfpathquadraticcurveto{\pgfqpoint{12.350479in}{0.719843in}}{\pgfqpoint{12.355180in}{0.726345in}}%
\pgfpathquadraticcurveto{\pgfqpoint{12.359881in}{0.732848in}}{\pgfqpoint{12.359881in}{0.744267in}}%
\pgfpathlineto{\pgfqpoint{12.359881in}{0.744267in}}%
\pgfpathclose%
\pgfpathmoveto{\pgfqpoint{12.397815in}{0.722256in}}%
\pgfpathlineto{\pgfqpoint{12.397815in}{0.688826in}}%
\pgfpathlineto{\pgfqpoint{12.387404in}{0.688826in}}%
\pgfpathlineto{\pgfqpoint{12.387404in}{0.775843in}}%
\pgfpathlineto{\pgfqpoint{12.397815in}{0.775843in}}%
\pgfpathlineto{\pgfqpoint{12.397815in}{0.766278in}}%
\pgfpathquadraticcurveto{\pgfqpoint{12.401093in}{0.771898in}}{\pgfqpoint{12.406064in}{0.774618in}}%
\pgfpathquadraticcurveto{\pgfqpoint{12.411054in}{0.777356in}}{\pgfqpoint{12.417970in}{0.777356in}}%
\pgfpathquadraticcurveto{\pgfqpoint{12.429462in}{0.777356in}}{\pgfqpoint{12.436631in}{0.768241in}}%
\pgfpathquadraticcurveto{\pgfqpoint{12.443818in}{0.759127in}}{\pgfqpoint{12.443818in}{0.744267in}}%
\pgfpathquadraticcurveto{\pgfqpoint{12.443818in}{0.729407in}}{\pgfqpoint{12.436631in}{0.720275in}}%
\pgfpathquadraticcurveto{\pgfqpoint{12.429462in}{0.711161in}}{\pgfqpoint{12.417970in}{0.711161in}}%
\pgfpathquadraticcurveto{\pgfqpoint{12.411054in}{0.711161in}}{\pgfqpoint{12.406064in}{0.713899in}}%
\pgfpathquadraticcurveto{\pgfqpoint{12.401093in}{0.716637in}}{\pgfqpoint{12.397815in}{0.722256in}}%
\pgfpathlineto{\pgfqpoint{12.397815in}{0.722256in}}%
\pgfpathclose%
\pgfpathmoveto{\pgfqpoint{12.433065in}{0.744267in}}%
\pgfpathquadraticcurveto{\pgfqpoint{12.433065in}{0.755687in}}{\pgfqpoint{12.428363in}{0.762189in}}%
\pgfpathquadraticcurveto{\pgfqpoint{12.423662in}{0.768692in}}{\pgfqpoint{12.415449in}{0.768692in}}%
\pgfpathquadraticcurveto{\pgfqpoint{12.407217in}{0.768692in}}{\pgfqpoint{12.402516in}{0.762189in}}%
\pgfpathquadraticcurveto{\pgfqpoint{12.397815in}{0.755687in}}{\pgfqpoint{12.397815in}{0.744267in}}%
\pgfpathquadraticcurveto{\pgfqpoint{12.397815in}{0.732848in}}{\pgfqpoint{12.402516in}{0.726345in}}%
\pgfpathquadraticcurveto{\pgfqpoint{12.407217in}{0.719843in}}{\pgfqpoint{12.415449in}{0.719843in}}%
\pgfpathquadraticcurveto{\pgfqpoint{12.423662in}{0.719843in}}{\pgfqpoint{12.428363in}{0.726345in}}%
\pgfpathquadraticcurveto{\pgfqpoint{12.433065in}{0.732848in}}{\pgfqpoint{12.433065in}{0.744267in}}%
\pgfpathlineto{\pgfqpoint{12.433065in}{0.744267in}}%
\pgfpathclose%
\pgfpathmoveto{\pgfqpoint{12.485408in}{0.768584in}}%
\pgfpathquadraticcurveto{\pgfqpoint{12.477086in}{0.768584in}}{\pgfqpoint{12.472241in}{0.762081in}}%
\pgfpathquadraticcurveto{\pgfqpoint{12.467396in}{0.755579in}}{\pgfqpoint{12.467396in}{0.744267in}}%
\pgfpathquadraticcurveto{\pgfqpoint{12.467396in}{0.732956in}}{\pgfqpoint{12.472205in}{0.726453in}}%
\pgfpathquadraticcurveto{\pgfqpoint{12.477032in}{0.719951in}}{\pgfqpoint{12.485408in}{0.719951in}}%
\pgfpathquadraticcurveto{\pgfqpoint{12.493693in}{0.719951in}}{\pgfqpoint{12.498521in}{0.726471in}}%
\pgfpathquadraticcurveto{\pgfqpoint{12.503366in}{0.733010in}}{\pgfqpoint{12.503366in}{0.744267in}}%
\pgfpathquadraticcurveto{\pgfqpoint{12.503366in}{0.755471in}}{\pgfqpoint{12.498521in}{0.762027in}}%
\pgfpathquadraticcurveto{\pgfqpoint{12.493693in}{0.768584in}}{\pgfqpoint{12.485408in}{0.768584in}}%
\pgfpathlineto{\pgfqpoint{12.485408in}{0.768584in}}%
\pgfpathclose%
\pgfpathmoveto{\pgfqpoint{12.485408in}{0.777356in}}%
\pgfpathquadraticcurveto{\pgfqpoint{12.498917in}{0.777356in}}{\pgfqpoint{12.506626in}{0.768566in}}%
\pgfpathquadraticcurveto{\pgfqpoint{12.514353in}{0.759794in}}{\pgfqpoint{12.514353in}{0.744267in}}%
\pgfpathquadraticcurveto{\pgfqpoint{12.514353in}{0.728795in}}{\pgfqpoint{12.506626in}{0.719969in}}%
\pgfpathquadraticcurveto{\pgfqpoint{12.498917in}{0.711161in}}{\pgfqpoint{12.485408in}{0.711161in}}%
\pgfpathquadraticcurveto{\pgfqpoint{12.471845in}{0.711161in}}{\pgfqpoint{12.464154in}{0.719969in}}%
\pgfpathquadraticcurveto{\pgfqpoint{12.456480in}{0.728795in}}{\pgfqpoint{12.456480in}{0.744267in}}%
\pgfpathquadraticcurveto{\pgfqpoint{12.456480in}{0.759794in}}{\pgfqpoint{12.464154in}{0.768566in}}%
\pgfpathquadraticcurveto{\pgfqpoint{12.471845in}{0.777356in}}{\pgfqpoint{12.485408in}{0.777356in}}%
\pgfpathlineto{\pgfqpoint{12.485408in}{0.777356in}}%
\pgfpathclose%
\pgfpathmoveto{\pgfqpoint{12.568048in}{0.766170in}}%
\pgfpathquadraticcurveto{\pgfqpoint{12.566300in}{0.767179in}}{\pgfqpoint{12.564247in}{0.767647in}}%
\pgfpathquadraticcurveto{\pgfqpoint{12.562194in}{0.768133in}}{\pgfqpoint{12.559726in}{0.768133in}}%
\pgfpathquadraticcurveto{\pgfqpoint{12.550936in}{0.768133in}}{\pgfqpoint{12.546235in}{0.762423in}}%
\pgfpathquadraticcurveto{\pgfqpoint{12.541534in}{0.756714in}}{\pgfqpoint{12.541534in}{0.746014in}}%
\pgfpathlineto{\pgfqpoint{12.541534in}{0.712800in}}%
\pgfpathlineto{\pgfqpoint{12.531123in}{0.712800in}}%
\pgfpathlineto{\pgfqpoint{12.531123in}{0.775843in}}%
\pgfpathlineto{\pgfqpoint{12.541534in}{0.775843in}}%
\pgfpathlineto{\pgfqpoint{12.541534in}{0.766044in}}%
\pgfpathquadraticcurveto{\pgfqpoint{12.544812in}{0.771790in}}{\pgfqpoint{12.550035in}{0.774564in}}%
\pgfpathquadraticcurveto{\pgfqpoint{12.555277in}{0.777356in}}{\pgfqpoint{12.562770in}{0.777356in}}%
\pgfpathquadraticcurveto{\pgfqpoint{12.563833in}{0.777356in}}{\pgfqpoint{12.565130in}{0.777211in}}%
\pgfpathquadraticcurveto{\pgfqpoint{12.566427in}{0.777085in}}{\pgfqpoint{12.567994in}{0.776797in}}%
\pgfpathlineto{\pgfqpoint{12.568048in}{0.766170in}}%
\pgfpathlineto{\pgfqpoint{12.568048in}{0.766170in}}%
\pgfpathclose%
\pgfpathmoveto{\pgfqpoint{12.589158in}{0.793747in}}%
\pgfpathlineto{\pgfqpoint{12.589158in}{0.775843in}}%
\pgfpathlineto{\pgfqpoint{12.610484in}{0.775843in}}%
\pgfpathlineto{\pgfqpoint{12.610484in}{0.767791in}}%
\pgfpathlineto{\pgfqpoint{12.589158in}{0.767791in}}%
\pgfpathlineto{\pgfqpoint{12.589158in}{0.733568in}}%
\pgfpathquadraticcurveto{\pgfqpoint{12.589158in}{0.725859in}}{\pgfqpoint{12.591265in}{0.723661in}}%
\pgfpathquadraticcurveto{\pgfqpoint{12.593373in}{0.721464in}}{\pgfqpoint{12.599857in}{0.721464in}}%
\pgfpathlineto{\pgfqpoint{12.610484in}{0.721464in}}%
\pgfpathlineto{\pgfqpoint{12.610484in}{0.712800in}}%
\pgfpathlineto{\pgfqpoint{12.599857in}{0.712800in}}%
\pgfpathquadraticcurveto{\pgfqpoint{12.587861in}{0.712800in}}{\pgfqpoint{12.583304in}{0.717267in}}%
\pgfpathquadraticcurveto{\pgfqpoint{12.578747in}{0.721752in}}{\pgfqpoint{12.578747in}{0.733568in}}%
\pgfpathlineto{\pgfqpoint{12.578747in}{0.767791in}}%
\pgfpathlineto{\pgfqpoint{12.571146in}{0.767791in}}%
\pgfpathlineto{\pgfqpoint{12.571146in}{0.775843in}}%
\pgfpathlineto{\pgfqpoint{12.578747in}{0.775843in}}%
\pgfpathlineto{\pgfqpoint{12.578747in}{0.793747in}}%
\pgfpathlineto{\pgfqpoint{12.589158in}{0.793747in}}%
\pgfpathlineto{\pgfqpoint{12.589158in}{0.793747in}}%
\pgfpathclose%
\pgfpathmoveto{\pgfqpoint{12.678030in}{0.746915in}}%
\pgfpathlineto{\pgfqpoint{12.678030in}{0.741854in}}%
\pgfpathlineto{\pgfqpoint{12.630406in}{0.741854in}}%
\pgfpathquadraticcurveto{\pgfqpoint{12.631090in}{0.731154in}}{\pgfqpoint{12.636854in}{0.725553in}}%
\pgfpathquadraticcurveto{\pgfqpoint{12.642618in}{0.719951in}}{\pgfqpoint{12.652921in}{0.719951in}}%
\pgfpathquadraticcurveto{\pgfqpoint{12.658883in}{0.719951in}}{\pgfqpoint{12.664485in}{0.721410in}}%
\pgfpathquadraticcurveto{\pgfqpoint{12.670086in}{0.722869in}}{\pgfqpoint{12.675616in}{0.725805in}}%
\pgfpathlineto{\pgfqpoint{12.675616in}{0.716006in}}%
\pgfpathquadraticcurveto{\pgfqpoint{12.670032in}{0.713647in}}{\pgfqpoint{12.664178in}{0.712404in}}%
\pgfpathquadraticcurveto{\pgfqpoint{12.658325in}{0.711161in}}{\pgfqpoint{12.652308in}{0.711161in}}%
\pgfpathquadraticcurveto{\pgfqpoint{12.637214in}{0.711161in}}{\pgfqpoint{12.628406in}{0.719933in}}%
\pgfpathquadraticcurveto{\pgfqpoint{12.619598in}{0.728723in}}{\pgfqpoint{12.619598in}{0.743709in}}%
\pgfpathquadraticcurveto{\pgfqpoint{12.619598in}{0.759181in}}{\pgfqpoint{12.627956in}{0.768259in}}%
\pgfpathquadraticcurveto{\pgfqpoint{12.636314in}{0.777356in}}{\pgfqpoint{12.650507in}{0.777356in}}%
\pgfpathquadraticcurveto{\pgfqpoint{12.663224in}{0.777356in}}{\pgfqpoint{12.670627in}{0.769160in}}%
\pgfpathquadraticcurveto{\pgfqpoint{12.678030in}{0.760983in}}{\pgfqpoint{12.678030in}{0.746915in}}%
\pgfpathlineto{\pgfqpoint{12.678030in}{0.746915in}}%
\pgfpathclose%
\pgfpathmoveto{\pgfqpoint{12.667673in}{0.749959in}}%
\pgfpathquadraticcurveto{\pgfqpoint{12.667565in}{0.758443in}}{\pgfqpoint{12.662918in}{0.763504in}}%
\pgfpathquadraticcurveto{\pgfqpoint{12.658270in}{0.768584in}}{\pgfqpoint{12.650615in}{0.768584in}}%
\pgfpathquadraticcurveto{\pgfqpoint{12.641951in}{0.768584in}}{\pgfqpoint{12.636746in}{0.763684in}}%
\pgfpathquadraticcurveto{\pgfqpoint{12.631540in}{0.758785in}}{\pgfqpoint{12.630748in}{0.749887in}}%
\pgfpathlineto{\pgfqpoint{12.667673in}{0.749959in}}%
\pgfpathlineto{\pgfqpoint{12.667673in}{0.749959in}}%
\pgfpathclose%
\pgfpathmoveto{\pgfqpoint{12.736515in}{0.766278in}}%
\pgfpathlineto{\pgfqpoint{12.736515in}{0.800393in}}%
\pgfpathlineto{\pgfqpoint{12.746872in}{0.800393in}}%
\pgfpathlineto{\pgfqpoint{12.746872in}{0.712800in}}%
\pgfpathlineto{\pgfqpoint{12.736515in}{0.712800in}}%
\pgfpathlineto{\pgfqpoint{12.736515in}{0.722256in}}%
\pgfpathquadraticcurveto{\pgfqpoint{12.733255in}{0.716637in}}{\pgfqpoint{12.728266in}{0.713899in}}%
\pgfpathquadraticcurveto{\pgfqpoint{12.723294in}{0.711161in}}{\pgfqpoint{12.716306in}{0.711161in}}%
\pgfpathquadraticcurveto{\pgfqpoint{12.704886in}{0.711161in}}{\pgfqpoint{12.697699in}{0.720275in}}%
\pgfpathquadraticcurveto{\pgfqpoint{12.690530in}{0.729407in}}{\pgfqpoint{12.690530in}{0.744267in}}%
\pgfpathquadraticcurveto{\pgfqpoint{12.690530in}{0.759127in}}{\pgfqpoint{12.697699in}{0.768241in}}%
\pgfpathquadraticcurveto{\pgfqpoint{12.704886in}{0.777356in}}{\pgfqpoint{12.716306in}{0.777356in}}%
\pgfpathquadraticcurveto{\pgfqpoint{12.723294in}{0.777356in}}{\pgfqpoint{12.728266in}{0.774618in}}%
\pgfpathquadraticcurveto{\pgfqpoint{12.733255in}{0.771898in}}{\pgfqpoint{12.736515in}{0.766278in}}%
\pgfpathlineto{\pgfqpoint{12.736515in}{0.766278in}}%
\pgfpathclose%
\pgfpathmoveto{\pgfqpoint{12.701229in}{0.744267in}}%
\pgfpathquadraticcurveto{\pgfqpoint{12.701229in}{0.732848in}}{\pgfqpoint{12.705931in}{0.726345in}}%
\pgfpathquadraticcurveto{\pgfqpoint{12.710632in}{0.719843in}}{\pgfqpoint{12.718845in}{0.719843in}}%
\pgfpathquadraticcurveto{\pgfqpoint{12.727059in}{0.719843in}}{\pgfqpoint{12.731778in}{0.726345in}}%
\pgfpathquadraticcurveto{\pgfqpoint{12.736515in}{0.732848in}}{\pgfqpoint{12.736515in}{0.744267in}}%
\pgfpathquadraticcurveto{\pgfqpoint{12.736515in}{0.755687in}}{\pgfqpoint{12.731778in}{0.762189in}}%
\pgfpathquadraticcurveto{\pgfqpoint{12.727059in}{0.768692in}}{\pgfqpoint{12.718845in}{0.768692in}}%
\pgfpathquadraticcurveto{\pgfqpoint{12.710632in}{0.768692in}}{\pgfqpoint{12.705931in}{0.762189in}}%
\pgfpathquadraticcurveto{\pgfqpoint{12.701229in}{0.755687in}}{\pgfqpoint{12.701229in}{0.744267in}}%
\pgfpathlineto{\pgfqpoint{12.701229in}{0.744267in}}%
\pgfpathclose%
\pgfusepath{fill}%
\end{pgfscope}%
\begin{pgfscope}%
\pgfsetbuttcap%
\pgfsetroundjoin%
\definecolor{currentfill}{rgb}{0.309804,0.372549,0.419608}%
\pgfsetfillcolor{currentfill}%
\pgfsetlinewidth{0.000000pt}%
\definecolor{currentstroke}{rgb}{0.309804,0.372549,0.419608}%
\pgfsetstrokecolor{currentstroke}%
\pgfsetdash{}{0pt}%
\pgfpathmoveto{\pgfqpoint{12.853449in}{0.793747in}}%
\pgfpathlineto{\pgfqpoint{12.853449in}{0.775843in}}%
\pgfpathlineto{\pgfqpoint{12.874775in}{0.775843in}}%
\pgfpathlineto{\pgfqpoint{12.874775in}{0.767791in}}%
\pgfpathlineto{\pgfqpoint{12.853449in}{0.767791in}}%
\pgfpathlineto{\pgfqpoint{12.853449in}{0.733568in}}%
\pgfpathquadraticcurveto{\pgfqpoint{12.853449in}{0.725859in}}{\pgfqpoint{12.855556in}{0.723661in}}%
\pgfpathquadraticcurveto{\pgfqpoint{12.857663in}{0.721464in}}{\pgfqpoint{12.864148in}{0.721464in}}%
\pgfpathlineto{\pgfqpoint{12.874775in}{0.721464in}}%
\pgfpathlineto{\pgfqpoint{12.874775in}{0.712800in}}%
\pgfpathlineto{\pgfqpoint{12.864148in}{0.712800in}}%
\pgfpathquadraticcurveto{\pgfqpoint{12.852152in}{0.712800in}}{\pgfqpoint{12.847595in}{0.717267in}}%
\pgfpathquadraticcurveto{\pgfqpoint{12.843038in}{0.721752in}}{\pgfqpoint{12.843038in}{0.733568in}}%
\pgfpathlineto{\pgfqpoint{12.843038in}{0.767791in}}%
\pgfpathlineto{\pgfqpoint{12.835436in}{0.767791in}}%
\pgfpathlineto{\pgfqpoint{12.835436in}{0.775843in}}%
\pgfpathlineto{\pgfqpoint{12.843038in}{0.775843in}}%
\pgfpathlineto{\pgfqpoint{12.843038in}{0.793747in}}%
\pgfpathlineto{\pgfqpoint{12.853449in}{0.793747in}}%
\pgfpathlineto{\pgfqpoint{12.853449in}{0.793747in}}%
\pgfpathclose%
\pgfpathmoveto{\pgfqpoint{12.917050in}{0.744483in}}%
\pgfpathquadraticcurveto{\pgfqpoint{12.904495in}{0.744483in}}{\pgfqpoint{12.899650in}{0.741619in}}%
\pgfpathquadraticcurveto{\pgfqpoint{12.894805in}{0.738756in}}{\pgfqpoint{12.894805in}{0.731821in}}%
\pgfpathquadraticcurveto{\pgfqpoint{12.894805in}{0.726309in}}{\pgfqpoint{12.898443in}{0.723067in}}%
\pgfpathquadraticcurveto{\pgfqpoint{12.902081in}{0.719843in}}{\pgfqpoint{12.908314in}{0.719843in}}%
\pgfpathquadraticcurveto{\pgfqpoint{12.916941in}{0.719843in}}{\pgfqpoint{12.922147in}{0.725949in}}%
\pgfpathquadraticcurveto{\pgfqpoint{12.927352in}{0.732055in}}{\pgfqpoint{12.927352in}{0.742178in}}%
\pgfpathlineto{\pgfqpoint{12.927352in}{0.744483in}}%
\pgfpathlineto{\pgfqpoint{12.917050in}{0.744483in}}%
\pgfpathlineto{\pgfqpoint{12.917050in}{0.744483in}}%
\pgfpathclose%
\pgfpathmoveto{\pgfqpoint{12.937709in}{0.748770in}}%
\pgfpathlineto{\pgfqpoint{12.937709in}{0.712800in}}%
\pgfpathlineto{\pgfqpoint{12.927352in}{0.712800in}}%
\pgfpathlineto{\pgfqpoint{12.927352in}{0.722364in}}%
\pgfpathquadraticcurveto{\pgfqpoint{12.923804in}{0.716637in}}{\pgfqpoint{12.918509in}{0.713899in}}%
\pgfpathquadraticcurveto{\pgfqpoint{12.913213in}{0.711161in}}{\pgfqpoint{12.905558in}{0.711161in}}%
\pgfpathquadraticcurveto{\pgfqpoint{12.895885in}{0.711161in}}{\pgfqpoint{12.890157in}{0.716601in}}%
\pgfpathquadraticcurveto{\pgfqpoint{12.884448in}{0.722040in}}{\pgfqpoint{12.884448in}{0.731154in}}%
\pgfpathquadraticcurveto{\pgfqpoint{12.884448in}{0.741782in}}{\pgfqpoint{12.891562in}{0.747185in}}%
\pgfpathquadraticcurveto{\pgfqpoint{12.898695in}{0.752589in}}{\pgfqpoint{12.912817in}{0.752589in}}%
\pgfpathlineto{\pgfqpoint{12.927352in}{0.752589in}}%
\pgfpathlineto{\pgfqpoint{12.927352in}{0.753616in}}%
\pgfpathquadraticcurveto{\pgfqpoint{12.927352in}{0.760766in}}{\pgfqpoint{12.922651in}{0.764675in}}%
\pgfpathquadraticcurveto{\pgfqpoint{12.917950in}{0.768584in}}{\pgfqpoint{12.909448in}{0.768584in}}%
\pgfpathquadraticcurveto{\pgfqpoint{12.904045in}{0.768584in}}{\pgfqpoint{12.898911in}{0.767287in}}%
\pgfpathquadraticcurveto{\pgfqpoint{12.893796in}{0.765990in}}{\pgfqpoint{12.889077in}{0.763396in}}%
\pgfpathlineto{\pgfqpoint{12.889077in}{0.772979in}}%
\pgfpathquadraticcurveto{\pgfqpoint{12.894750in}{0.775176in}}{\pgfqpoint{12.900100in}{0.776257in}}%
\pgfpathquadraticcurveto{\pgfqpoint{12.905450in}{0.777356in}}{\pgfqpoint{12.910511in}{0.777356in}}%
\pgfpathquadraticcurveto{\pgfqpoint{12.924200in}{0.777356in}}{\pgfqpoint{12.930955in}{0.770259in}}%
\pgfpathquadraticcurveto{\pgfqpoint{12.937709in}{0.763180in}}{\pgfqpoint{12.937709in}{0.748770in}}%
\pgfpathlineto{\pgfqpoint{12.937709in}{0.748770in}}%
\pgfpathclose%
\pgfpathmoveto{\pgfqpoint{12.995564in}{0.766170in}}%
\pgfpathquadraticcurveto{\pgfqpoint{12.993817in}{0.767179in}}{\pgfqpoint{12.991764in}{0.767647in}}%
\pgfpathquadraticcurveto{\pgfqpoint{12.989711in}{0.768133in}}{\pgfqpoint{12.987243in}{0.768133in}}%
\pgfpathquadraticcurveto{\pgfqpoint{12.978453in}{0.768133in}}{\pgfqpoint{12.973752in}{0.762423in}}%
\pgfpathquadraticcurveto{\pgfqpoint{12.969051in}{0.756714in}}{\pgfqpoint{12.969051in}{0.746014in}}%
\pgfpathlineto{\pgfqpoint{12.969051in}{0.712800in}}%
\pgfpathlineto{\pgfqpoint{12.958640in}{0.712800in}}%
\pgfpathlineto{\pgfqpoint{12.958640in}{0.775843in}}%
\pgfpathlineto{\pgfqpoint{12.969051in}{0.775843in}}%
\pgfpathlineto{\pgfqpoint{12.969051in}{0.766044in}}%
\pgfpathquadraticcurveto{\pgfqpoint{12.972329in}{0.771790in}}{\pgfqpoint{12.977552in}{0.774564in}}%
\pgfpathquadraticcurveto{\pgfqpoint{12.982794in}{0.777356in}}{\pgfqpoint{12.990287in}{0.777356in}}%
\pgfpathquadraticcurveto{\pgfqpoint{12.991350in}{0.777356in}}{\pgfqpoint{12.992647in}{0.777211in}}%
\pgfpathquadraticcurveto{\pgfqpoint{12.993943in}{0.777085in}}{\pgfqpoint{12.995510in}{0.776797in}}%
\pgfpathlineto{\pgfqpoint{12.995564in}{0.766170in}}%
\pgfpathlineto{\pgfqpoint{12.995564in}{0.766170in}}%
\pgfpathclose%
\pgfpathmoveto{\pgfqpoint{13.045890in}{0.745060in}}%
\pgfpathquadraticcurveto{\pgfqpoint{13.045890in}{0.756317in}}{\pgfqpoint{13.041243in}{0.762496in}}%
\pgfpathquadraticcurveto{\pgfqpoint{13.036614in}{0.768692in}}{\pgfqpoint{13.028221in}{0.768692in}}%
\pgfpathquadraticcurveto{\pgfqpoint{13.019899in}{0.768692in}}{\pgfqpoint{13.015252in}{0.762496in}}%
\pgfpathquadraticcurveto{\pgfqpoint{13.010605in}{0.756317in}}{\pgfqpoint{13.010605in}{0.745060in}}%
\pgfpathquadraticcurveto{\pgfqpoint{13.010605in}{0.733856in}}{\pgfqpoint{13.015252in}{0.727660in}}%
\pgfpathquadraticcurveto{\pgfqpoint{13.019899in}{0.721464in}}{\pgfqpoint{13.028221in}{0.721464in}}%
\pgfpathquadraticcurveto{\pgfqpoint{13.036614in}{0.721464in}}{\pgfqpoint{13.041243in}{0.727660in}}%
\pgfpathquadraticcurveto{\pgfqpoint{13.045890in}{0.733856in}}{\pgfqpoint{13.045890in}{0.745060in}}%
\pgfpathlineto{\pgfqpoint{13.045890in}{0.745060in}}%
\pgfpathclose%
\pgfpathmoveto{\pgfqpoint{13.056247in}{0.720617in}}%
\pgfpathquadraticcurveto{\pgfqpoint{13.056247in}{0.704532in}}{\pgfqpoint{13.049097in}{0.696679in}}%
\pgfpathquadraticcurveto{\pgfqpoint{13.041964in}{0.688826in}}{\pgfqpoint{13.027212in}{0.688826in}}%
\pgfpathquadraticcurveto{\pgfqpoint{13.021754in}{0.688826in}}{\pgfqpoint{13.016909in}{0.689636in}}%
\pgfpathquadraticcurveto{\pgfqpoint{13.012064in}{0.690447in}}{\pgfqpoint{13.007507in}{0.692140in}}%
\pgfpathlineto{\pgfqpoint{13.007507in}{0.702209in}}%
\pgfpathquadraticcurveto{\pgfqpoint{13.012064in}{0.699741in}}{\pgfqpoint{13.016513in}{0.698570in}}%
\pgfpathquadraticcurveto{\pgfqpoint{13.020962in}{0.697382in}}{\pgfqpoint{13.025573in}{0.697382in}}%
\pgfpathquadraticcurveto{\pgfqpoint{13.035768in}{0.697382in}}{\pgfqpoint{13.040829in}{0.702695in}}%
\pgfpathquadraticcurveto{\pgfqpoint{13.045890in}{0.708009in}}{\pgfqpoint{13.045890in}{0.718762in}}%
\pgfpathlineto{\pgfqpoint{13.045890in}{0.723895in}}%
\pgfpathquadraticcurveto{\pgfqpoint{13.042684in}{0.718312in}}{\pgfqpoint{13.037677in}{0.715556in}}%
\pgfpathquadraticcurveto{\pgfqpoint{13.032670in}{0.712800in}}{\pgfqpoint{13.025681in}{0.712800in}}%
\pgfpathquadraticcurveto{\pgfqpoint{13.014099in}{0.712800in}}{\pgfqpoint{13.007002in}{0.721626in}}%
\pgfpathquadraticcurveto{\pgfqpoint{12.999905in}{0.730470in}}{\pgfqpoint{12.999905in}{0.745060in}}%
\pgfpathquadraticcurveto{\pgfqpoint{12.999905in}{0.759686in}}{\pgfqpoint{13.007002in}{0.768512in}}%
\pgfpathquadraticcurveto{\pgfqpoint{13.014099in}{0.777356in}}{\pgfqpoint{13.025681in}{0.777356in}}%
\pgfpathquadraticcurveto{\pgfqpoint{13.032670in}{0.777356in}}{\pgfqpoint{13.037677in}{0.774600in}}%
\pgfpathquadraticcurveto{\pgfqpoint{13.042684in}{0.771844in}}{\pgfqpoint{13.045890in}{0.766278in}}%
\pgfpathlineto{\pgfqpoint{13.045890in}{0.775843in}}%
\pgfpathlineto{\pgfqpoint{13.056247in}{0.775843in}}%
\pgfpathlineto{\pgfqpoint{13.056247in}{0.720617in}}%
\pgfpathlineto{\pgfqpoint{13.056247in}{0.720617in}}%
\pgfpathclose%
\pgfpathmoveto{\pgfqpoint{13.131520in}{0.746915in}}%
\pgfpathlineto{\pgfqpoint{13.131520in}{0.741854in}}%
\pgfpathlineto{\pgfqpoint{13.083896in}{0.741854in}}%
\pgfpathquadraticcurveto{\pgfqpoint{13.084581in}{0.731154in}}{\pgfqpoint{13.090344in}{0.725553in}}%
\pgfpathquadraticcurveto{\pgfqpoint{13.096108in}{0.719951in}}{\pgfqpoint{13.106411in}{0.719951in}}%
\pgfpathquadraticcurveto{\pgfqpoint{13.112373in}{0.719951in}}{\pgfqpoint{13.117975in}{0.721410in}}%
\pgfpathquadraticcurveto{\pgfqpoint{13.123577in}{0.722869in}}{\pgfqpoint{13.129107in}{0.725805in}}%
\pgfpathlineto{\pgfqpoint{13.129107in}{0.716006in}}%
\pgfpathquadraticcurveto{\pgfqpoint{13.123523in}{0.713647in}}{\pgfqpoint{13.117669in}{0.712404in}}%
\pgfpathquadraticcurveto{\pgfqpoint{13.111815in}{0.711161in}}{\pgfqpoint{13.105799in}{0.711161in}}%
\pgfpathquadraticcurveto{\pgfqpoint{13.090705in}{0.711161in}}{\pgfqpoint{13.081897in}{0.719933in}}%
\pgfpathquadraticcurveto{\pgfqpoint{13.073089in}{0.728723in}}{\pgfqpoint{13.073089in}{0.743709in}}%
\pgfpathquadraticcurveto{\pgfqpoint{13.073089in}{0.759181in}}{\pgfqpoint{13.081446in}{0.768259in}}%
\pgfpathquadraticcurveto{\pgfqpoint{13.089804in}{0.777356in}}{\pgfqpoint{13.103998in}{0.777356in}}%
\pgfpathquadraticcurveto{\pgfqpoint{13.116714in}{0.777356in}}{\pgfqpoint{13.124117in}{0.769160in}}%
\pgfpathquadraticcurveto{\pgfqpoint{13.131520in}{0.760983in}}{\pgfqpoint{13.131520in}{0.746915in}}%
\pgfpathlineto{\pgfqpoint{13.131520in}{0.746915in}}%
\pgfpathclose%
\pgfpathmoveto{\pgfqpoint{13.121163in}{0.749959in}}%
\pgfpathquadraticcurveto{\pgfqpoint{13.121055in}{0.758443in}}{\pgfqpoint{13.116408in}{0.763504in}}%
\pgfpathquadraticcurveto{\pgfqpoint{13.111761in}{0.768584in}}{\pgfqpoint{13.104106in}{0.768584in}}%
\pgfpathquadraticcurveto{\pgfqpoint{13.095442in}{0.768584in}}{\pgfqpoint{13.090236in}{0.763684in}}%
\pgfpathquadraticcurveto{\pgfqpoint{13.085031in}{0.758785in}}{\pgfqpoint{13.084238in}{0.749887in}}%
\pgfpathlineto{\pgfqpoint{13.121163in}{0.749959in}}%
\pgfpathlineto{\pgfqpoint{13.121163in}{0.749959in}}%
\pgfpathclose%
\pgfpathmoveto{\pgfqpoint{13.158773in}{0.793747in}}%
\pgfpathlineto{\pgfqpoint{13.158773in}{0.775843in}}%
\pgfpathlineto{\pgfqpoint{13.180099in}{0.775843in}}%
\pgfpathlineto{\pgfqpoint{13.180099in}{0.767791in}}%
\pgfpathlineto{\pgfqpoint{13.158773in}{0.767791in}}%
\pgfpathlineto{\pgfqpoint{13.158773in}{0.733568in}}%
\pgfpathquadraticcurveto{\pgfqpoint{13.158773in}{0.725859in}}{\pgfqpoint{13.160880in}{0.723661in}}%
\pgfpathquadraticcurveto{\pgfqpoint{13.162987in}{0.721464in}}{\pgfqpoint{13.169472in}{0.721464in}}%
\pgfpathlineto{\pgfqpoint{13.180099in}{0.721464in}}%
\pgfpathlineto{\pgfqpoint{13.180099in}{0.712800in}}%
\pgfpathlineto{\pgfqpoint{13.169472in}{0.712800in}}%
\pgfpathquadraticcurveto{\pgfqpoint{13.157476in}{0.712800in}}{\pgfqpoint{13.152919in}{0.717267in}}%
\pgfpathquadraticcurveto{\pgfqpoint{13.148362in}{0.721752in}}{\pgfqpoint{13.148362in}{0.733568in}}%
\pgfpathlineto{\pgfqpoint{13.148362in}{0.767791in}}%
\pgfpathlineto{\pgfqpoint{13.140760in}{0.767791in}}%
\pgfpathlineto{\pgfqpoint{13.140760in}{0.775843in}}%
\pgfpathlineto{\pgfqpoint{13.148362in}{0.775843in}}%
\pgfpathlineto{\pgfqpoint{13.148362in}{0.793747in}}%
\pgfpathlineto{\pgfqpoint{13.158773in}{0.793747in}}%
\pgfpathlineto{\pgfqpoint{13.158773in}{0.793747in}}%
\pgfpathclose%
\pgfpathmoveto{\pgfqpoint{13.233901in}{0.773987in}}%
\pgfpathlineto{\pgfqpoint{13.233901in}{0.764189in}}%
\pgfpathquadraticcurveto{\pgfqpoint{13.229524in}{0.766440in}}{\pgfqpoint{13.224787in}{0.767557in}}%
\pgfpathquadraticcurveto{\pgfqpoint{13.220068in}{0.768692in}}{\pgfqpoint{13.214989in}{0.768692in}}%
\pgfpathquadraticcurveto{\pgfqpoint{13.207279in}{0.768692in}}{\pgfqpoint{13.203425in}{0.766332in}}%
\pgfpathquadraticcurveto{\pgfqpoint{13.199570in}{0.763973in}}{\pgfqpoint{13.199570in}{0.759235in}}%
\pgfpathquadraticcurveto{\pgfqpoint{13.199570in}{0.755633in}}{\pgfqpoint{13.202326in}{0.753580in}}%
\pgfpathquadraticcurveto{\pgfqpoint{13.205082in}{0.751526in}}{\pgfqpoint{13.213421in}{0.749671in}}%
\pgfpathlineto{\pgfqpoint{13.216970in}{0.748878in}}%
\pgfpathquadraticcurveto{\pgfqpoint{13.227993in}{0.746519in}}{\pgfqpoint{13.232640in}{0.742214in}}%
\pgfpathquadraticcurveto{\pgfqpoint{13.237288in}{0.737909in}}{\pgfqpoint{13.237288in}{0.730200in}}%
\pgfpathquadraticcurveto{\pgfqpoint{13.237288in}{0.721410in}}{\pgfqpoint{13.230335in}{0.716276in}}%
\pgfpathquadraticcurveto{\pgfqpoint{13.223382in}{0.711161in}}{\pgfqpoint{13.211224in}{0.711161in}}%
\pgfpathquadraticcurveto{\pgfqpoint{13.206163in}{0.711161in}}{\pgfqpoint{13.200669in}{0.712152in}}%
\pgfpathquadraticcurveto{\pgfqpoint{13.195175in}{0.713142in}}{\pgfqpoint{13.189105in}{0.715106in}}%
\pgfpathlineto{\pgfqpoint{13.189105in}{0.725805in}}%
\pgfpathquadraticcurveto{\pgfqpoint{13.194851in}{0.722815in}}{\pgfqpoint{13.200417in}{0.721320in}}%
\pgfpathquadraticcurveto{\pgfqpoint{13.205982in}{0.719843in}}{\pgfqpoint{13.211458in}{0.719843in}}%
\pgfpathquadraticcurveto{\pgfqpoint{13.218771in}{0.719843in}}{\pgfqpoint{13.222698in}{0.722346in}}%
\pgfpathquadraticcurveto{\pgfqpoint{13.226642in}{0.724850in}}{\pgfqpoint{13.226642in}{0.729407in}}%
\pgfpathquadraticcurveto{\pgfqpoint{13.226642in}{0.733622in}}{\pgfqpoint{13.223796in}{0.735874in}}%
\pgfpathquadraticcurveto{\pgfqpoint{13.220969in}{0.738125in}}{\pgfqpoint{13.211332in}{0.740214in}}%
\pgfpathlineto{\pgfqpoint{13.207730in}{0.741061in}}%
\pgfpathquadraticcurveto{\pgfqpoint{13.198111in}{0.743078in}}{\pgfqpoint{13.193824in}{0.747275in}}%
\pgfpathquadraticcurveto{\pgfqpoint{13.189555in}{0.751472in}}{\pgfqpoint{13.189555in}{0.758785in}}%
\pgfpathquadraticcurveto{\pgfqpoint{13.189555in}{0.767683in}}{\pgfqpoint{13.195860in}{0.772510in}}%
\pgfpathquadraticcurveto{\pgfqpoint{13.202164in}{0.777356in}}{\pgfqpoint{13.213764in}{0.777356in}}%
\pgfpathquadraticcurveto{\pgfqpoint{13.219492in}{0.777356in}}{\pgfqpoint{13.224553in}{0.776509in}}%
\pgfpathquadraticcurveto{\pgfqpoint{13.229632in}{0.775680in}}{\pgfqpoint{13.233901in}{0.773987in}}%
\pgfpathlineto{\pgfqpoint{13.233901in}{0.773987in}}%
\pgfpathclose%
\pgfpathmoveto{\pgfqpoint{13.255228in}{0.727102in}}%
\pgfpathlineto{\pgfqpoint{13.267116in}{0.727102in}}%
\pgfpathlineto{\pgfqpoint{13.267116in}{0.712800in}}%
\pgfpathlineto{\pgfqpoint{13.255228in}{0.712800in}}%
\pgfpathlineto{\pgfqpoint{13.255228in}{0.727102in}}%
\pgfpathlineto{\pgfqpoint{13.255228in}{0.727102in}}%
\pgfpathclose%
\pgfusepath{fill}%
\end{pgfscope}%
\begin{pgfscope}%
\pgfsetbuttcap%
\pgfsetroundjoin%
\definecolor{currentfill}{rgb}{0.309804,0.372549,0.419608}%
\pgfsetfillcolor{currentfill}%
\pgfsetlinewidth{0.000000pt}%
\definecolor{currentstroke}{rgb}{0.309804,0.372549,0.419608}%
\pgfsetstrokecolor{currentstroke}%
\pgfsetdash{}{0pt}%
\pgfpathmoveto{\pgfqpoint{13.364677in}{0.752931in}}%
\pgfpathlineto{\pgfqpoint{13.364677in}{0.722148in}}%
\pgfpathlineto{\pgfqpoint{13.382924in}{0.722148in}}%
\pgfpathquadraticcurveto{\pgfqpoint{13.392092in}{0.722148in}}{\pgfqpoint{13.396505in}{0.725949in}}%
\pgfpathquadraticcurveto{\pgfqpoint{13.400936in}{0.729749in}}{\pgfqpoint{13.400936in}{0.737567in}}%
\pgfpathquadraticcurveto{\pgfqpoint{13.400936in}{0.745456in}}{\pgfqpoint{13.396505in}{0.749185in}}%
\pgfpathquadraticcurveto{\pgfqpoint{13.392092in}{0.752931in}}{\pgfqpoint{13.382924in}{0.752931in}}%
\pgfpathlineto{\pgfqpoint{13.364677in}{0.752931in}}%
\pgfpathlineto{\pgfqpoint{13.364677in}{0.752931in}}%
\pgfpathclose%
\pgfpathmoveto{\pgfqpoint{13.364677in}{0.787496in}}%
\pgfpathlineto{\pgfqpoint{13.364677in}{0.762171in}}%
\pgfpathlineto{\pgfqpoint{13.381519in}{0.762171in}}%
\pgfpathquadraticcurveto{\pgfqpoint{13.389840in}{0.762171in}}{\pgfqpoint{13.393911in}{0.765287in}}%
\pgfpathquadraticcurveto{\pgfqpoint{13.398000in}{0.768422in}}{\pgfqpoint{13.398000in}{0.774834in}}%
\pgfpathquadraticcurveto{\pgfqpoint{13.398000in}{0.781192in}}{\pgfqpoint{13.393911in}{0.784344in}}%
\pgfpathquadraticcurveto{\pgfqpoint{13.389840in}{0.787496in}}{\pgfqpoint{13.381519in}{0.787496in}}%
\pgfpathlineto{\pgfqpoint{13.364677in}{0.787496in}}%
\pgfpathlineto{\pgfqpoint{13.364677in}{0.787496in}}%
\pgfpathclose%
\pgfpathmoveto{\pgfqpoint{13.353312in}{0.796845in}}%
\pgfpathlineto{\pgfqpoint{13.382365in}{0.796845in}}%
\pgfpathquadraticcurveto{\pgfqpoint{13.395370in}{0.796845in}}{\pgfqpoint{13.402395in}{0.791441in}}%
\pgfpathquadraticcurveto{\pgfqpoint{13.409438in}{0.786037in}}{\pgfqpoint{13.409438in}{0.776077in}}%
\pgfpathquadraticcurveto{\pgfqpoint{13.409438in}{0.768349in}}{\pgfqpoint{13.405835in}{0.763792in}}%
\pgfpathquadraticcurveto{\pgfqpoint{13.402233in}{0.759235in}}{\pgfqpoint{13.395244in}{0.758119in}}%
\pgfpathquadraticcurveto{\pgfqpoint{13.403638in}{0.756317in}}{\pgfqpoint{13.408285in}{0.750589in}}%
\pgfpathquadraticcurveto{\pgfqpoint{13.412932in}{0.744880in}}{\pgfqpoint{13.412932in}{0.736324in}}%
\pgfpathquadraticcurveto{\pgfqpoint{13.412932in}{0.725066in}}{\pgfqpoint{13.405277in}{0.718924in}}%
\pgfpathquadraticcurveto{\pgfqpoint{13.397622in}{0.712800in}}{\pgfqpoint{13.383482in}{0.712800in}}%
\pgfpathlineto{\pgfqpoint{13.353312in}{0.712800in}}%
\pgfpathlineto{\pgfqpoint{13.353312in}{0.796845in}}%
\pgfpathlineto{\pgfqpoint{13.353312in}{0.796845in}}%
\pgfpathclose%
\pgfpathmoveto{\pgfqpoint{13.456377in}{0.768584in}}%
\pgfpathquadraticcurveto{\pgfqpoint{13.448056in}{0.768584in}}{\pgfqpoint{13.443210in}{0.762081in}}%
\pgfpathquadraticcurveto{\pgfqpoint{13.438365in}{0.755579in}}{\pgfqpoint{13.438365in}{0.744267in}}%
\pgfpathquadraticcurveto{\pgfqpoint{13.438365in}{0.732956in}}{\pgfqpoint{13.443174in}{0.726453in}}%
\pgfpathquadraticcurveto{\pgfqpoint{13.448002in}{0.719951in}}{\pgfqpoint{13.456377in}{0.719951in}}%
\pgfpathquadraticcurveto{\pgfqpoint{13.464663in}{0.719951in}}{\pgfqpoint{13.469490in}{0.726471in}}%
\pgfpathquadraticcurveto{\pgfqpoint{13.474335in}{0.733010in}}{\pgfqpoint{13.474335in}{0.744267in}}%
\pgfpathquadraticcurveto{\pgfqpoint{13.474335in}{0.755471in}}{\pgfqpoint{13.469490in}{0.762027in}}%
\pgfpathquadraticcurveto{\pgfqpoint{13.464663in}{0.768584in}}{\pgfqpoint{13.456377in}{0.768584in}}%
\pgfpathlineto{\pgfqpoint{13.456377in}{0.768584in}}%
\pgfpathclose%
\pgfpathmoveto{\pgfqpoint{13.456377in}{0.777356in}}%
\pgfpathquadraticcurveto{\pgfqpoint{13.469886in}{0.777356in}}{\pgfqpoint{13.477595in}{0.768566in}}%
\pgfpathquadraticcurveto{\pgfqpoint{13.485323in}{0.759794in}}{\pgfqpoint{13.485323in}{0.744267in}}%
\pgfpathquadraticcurveto{\pgfqpoint{13.485323in}{0.728795in}}{\pgfqpoint{13.477595in}{0.719969in}}%
\pgfpathquadraticcurveto{\pgfqpoint{13.469886in}{0.711161in}}{\pgfqpoint{13.456377in}{0.711161in}}%
\pgfpathquadraticcurveto{\pgfqpoint{13.442814in}{0.711161in}}{\pgfqpoint{13.435123in}{0.719969in}}%
\pgfpathquadraticcurveto{\pgfqpoint{13.427450in}{0.728795in}}{\pgfqpoint{13.427450in}{0.744267in}}%
\pgfpathquadraticcurveto{\pgfqpoint{13.427450in}{0.759794in}}{\pgfqpoint{13.435123in}{0.768566in}}%
\pgfpathquadraticcurveto{\pgfqpoint{13.442814in}{0.777356in}}{\pgfqpoint{13.456377in}{0.777356in}}%
\pgfpathlineto{\pgfqpoint{13.456377in}{0.777356in}}%
\pgfpathclose%
\pgfpathmoveto{\pgfqpoint{13.512737in}{0.793747in}}%
\pgfpathlineto{\pgfqpoint{13.512737in}{0.775843in}}%
\pgfpathlineto{\pgfqpoint{13.534064in}{0.775843in}}%
\pgfpathlineto{\pgfqpoint{13.534064in}{0.767791in}}%
\pgfpathlineto{\pgfqpoint{13.512737in}{0.767791in}}%
\pgfpathlineto{\pgfqpoint{13.512737in}{0.733568in}}%
\pgfpathquadraticcurveto{\pgfqpoint{13.512737in}{0.725859in}}{\pgfqpoint{13.514845in}{0.723661in}}%
\pgfpathquadraticcurveto{\pgfqpoint{13.516952in}{0.721464in}}{\pgfqpoint{13.523436in}{0.721464in}}%
\pgfpathlineto{\pgfqpoint{13.534064in}{0.721464in}}%
\pgfpathlineto{\pgfqpoint{13.534064in}{0.712800in}}%
\pgfpathlineto{\pgfqpoint{13.523436in}{0.712800in}}%
\pgfpathquadraticcurveto{\pgfqpoint{13.511440in}{0.712800in}}{\pgfqpoint{13.506883in}{0.717267in}}%
\pgfpathquadraticcurveto{\pgfqpoint{13.502326in}{0.721752in}}{\pgfqpoint{13.502326in}{0.733568in}}%
\pgfpathlineto{\pgfqpoint{13.502326in}{0.767791in}}%
\pgfpathlineto{\pgfqpoint{13.494725in}{0.767791in}}%
\pgfpathlineto{\pgfqpoint{13.494725in}{0.775843in}}%
\pgfpathlineto{\pgfqpoint{13.502326in}{0.775843in}}%
\pgfpathlineto{\pgfqpoint{13.502326in}{0.793747in}}%
\pgfpathlineto{\pgfqpoint{13.512737in}{0.793747in}}%
\pgfpathlineto{\pgfqpoint{13.512737in}{0.793747in}}%
\pgfpathclose%
\pgfpathmoveto{\pgfqpoint{13.600096in}{0.750860in}}%
\pgfpathlineto{\pgfqpoint{13.600096in}{0.712800in}}%
\pgfpathlineto{\pgfqpoint{13.589739in}{0.712800in}}%
\pgfpathlineto{\pgfqpoint{13.589739in}{0.750517in}}%
\pgfpathquadraticcurveto{\pgfqpoint{13.589739in}{0.759469in}}{\pgfqpoint{13.586245in}{0.763900in}}%
\pgfpathquadraticcurveto{\pgfqpoint{13.582750in}{0.768349in}}{\pgfqpoint{13.575780in}{0.768349in}}%
\pgfpathquadraticcurveto{\pgfqpoint{13.567386in}{0.768349in}}{\pgfqpoint{13.562541in}{0.763000in}}%
\pgfpathquadraticcurveto{\pgfqpoint{13.557696in}{0.757668in}}{\pgfqpoint{13.557696in}{0.748428in}}%
\pgfpathlineto{\pgfqpoint{13.557696in}{0.712800in}}%
\pgfpathlineto{\pgfqpoint{13.547285in}{0.712800in}}%
\pgfpathlineto{\pgfqpoint{13.547285in}{0.800393in}}%
\pgfpathlineto{\pgfqpoint{13.557696in}{0.800393in}}%
\pgfpathlineto{\pgfqpoint{13.557696in}{0.766044in}}%
\pgfpathquadraticcurveto{\pgfqpoint{13.561424in}{0.771736in}}{\pgfqpoint{13.566449in}{0.774546in}}%
\pgfpathquadraticcurveto{\pgfqpoint{13.571493in}{0.777356in}}{\pgfqpoint{13.578085in}{0.777356in}}%
\pgfpathquadraticcurveto{\pgfqpoint{13.588947in}{0.777356in}}{\pgfqpoint{13.594512in}{0.770637in}}%
\pgfpathquadraticcurveto{\pgfqpoint{13.600096in}{0.763918in}}{\pgfqpoint{13.600096in}{0.750860in}}%
\pgfpathlineto{\pgfqpoint{13.600096in}{0.750860in}}%
\pgfpathclose%
\pgfusepath{fill}%
\end{pgfscope}%
\begin{pgfscope}%
\pgfsetbuttcap%
\pgfsetroundjoin%
\definecolor{currentfill}{rgb}{0.309804,0.372549,0.419608}%
\pgfsetfillcolor{currentfill}%
\pgfsetlinewidth{0.000000pt}%
\definecolor{currentstroke}{rgb}{0.309804,0.372549,0.419608}%
\pgfsetstrokecolor{currentstroke}%
\pgfsetdash{}{0pt}%
\pgfpathmoveto{\pgfqpoint{0.963519in}{0.560883in}}%
\pgfpathquadraticcurveto{\pgfqpoint{0.950964in}{0.560883in}}{\pgfqpoint{0.946119in}{0.558019in}}%
\pgfpathquadraticcurveto{\pgfqpoint{0.941274in}{0.555156in}}{\pgfqpoint{0.941274in}{0.548221in}}%
\pgfpathquadraticcurveto{\pgfqpoint{0.941274in}{0.542709in}}{\pgfqpoint{0.944912in}{0.539467in}}%
\pgfpathquadraticcurveto{\pgfqpoint{0.948551in}{0.536243in}}{\pgfqpoint{0.954783in}{0.536243in}}%
\pgfpathquadraticcurveto{\pgfqpoint{0.963411in}{0.536243in}}{\pgfqpoint{0.968616in}{0.542349in}}%
\pgfpathquadraticcurveto{\pgfqpoint{0.973822in}{0.548455in}}{\pgfqpoint{0.973822in}{0.558578in}}%
\pgfpathlineto{\pgfqpoint{0.973822in}{0.560883in}}%
\pgfpathlineto{\pgfqpoint{0.963519in}{0.560883in}}%
\pgfpathlineto{\pgfqpoint{0.963519in}{0.560883in}}%
\pgfpathclose%
\pgfpathmoveto{\pgfqpoint{0.984179in}{0.565170in}}%
\pgfpathlineto{\pgfqpoint{0.984179in}{0.529200in}}%
\pgfpathlineto{\pgfqpoint{0.973822in}{0.529200in}}%
\pgfpathlineto{\pgfqpoint{0.973822in}{0.538764in}}%
\pgfpathquadraticcurveto{\pgfqpoint{0.970273in}{0.533037in}}{\pgfqpoint{0.964978in}{0.530299in}}%
\pgfpathquadraticcurveto{\pgfqpoint{0.959682in}{0.527561in}}{\pgfqpoint{0.952027in}{0.527561in}}%
\pgfpathquadraticcurveto{\pgfqpoint{0.942354in}{0.527561in}}{\pgfqpoint{0.936627in}{0.533001in}}%
\pgfpathquadraticcurveto{\pgfqpoint{0.930917in}{0.538440in}}{\pgfqpoint{0.930917in}{0.547554in}}%
\pgfpathquadraticcurveto{\pgfqpoint{0.930917in}{0.558182in}}{\pgfqpoint{0.938031in}{0.563585in}}%
\pgfpathquadraticcurveto{\pgfqpoint{0.945164in}{0.568989in}}{\pgfqpoint{0.959286in}{0.568989in}}%
\pgfpathlineto{\pgfqpoint{0.973822in}{0.568989in}}%
\pgfpathlineto{\pgfqpoint{0.973822in}{0.570016in}}%
\pgfpathquadraticcurveto{\pgfqpoint{0.973822in}{0.577166in}}{\pgfqpoint{0.969120in}{0.581075in}}%
\pgfpathquadraticcurveto{\pgfqpoint{0.964419in}{0.584984in}}{\pgfqpoint{0.955918in}{0.584984in}}%
\pgfpathquadraticcurveto{\pgfqpoint{0.950514in}{0.584984in}}{\pgfqpoint{0.945380in}{0.583687in}}%
\pgfpathquadraticcurveto{\pgfqpoint{0.940265in}{0.582390in}}{\pgfqpoint{0.935546in}{0.579796in}}%
\pgfpathlineto{\pgfqpoint{0.935546in}{0.589379in}}%
\pgfpathquadraticcurveto{\pgfqpoint{0.941220in}{0.591576in}}{\pgfqpoint{0.946569in}{0.592657in}}%
\pgfpathquadraticcurveto{\pgfqpoint{0.951919in}{0.593756in}}{\pgfqpoint{0.956980in}{0.593756in}}%
\pgfpathquadraticcurveto{\pgfqpoint{0.970669in}{0.593756in}}{\pgfqpoint{0.977424in}{0.586659in}}%
\pgfpathquadraticcurveto{\pgfqpoint{0.984179in}{0.579580in}}{\pgfqpoint{0.984179in}{0.565170in}}%
\pgfpathlineto{\pgfqpoint{0.984179in}{0.565170in}}%
\pgfpathclose%
\pgfpathmoveto{\pgfqpoint{1.057920in}{0.592243in}}%
\pgfpathlineto{\pgfqpoint{1.035117in}{0.561568in}}%
\pgfpathlineto{\pgfqpoint{1.059091in}{0.529200in}}%
\pgfpathlineto{\pgfqpoint{1.046879in}{0.529200in}}%
\pgfpathlineto{\pgfqpoint{1.028525in}{0.553967in}}%
\pgfpathlineto{\pgfqpoint{1.010188in}{0.529200in}}%
\pgfpathlineto{\pgfqpoint{0.997958in}{0.529200in}}%
\pgfpathlineto{\pgfqpoint{1.022454in}{0.562180in}}%
\pgfpathlineto{\pgfqpoint{1.000047in}{0.592243in}}%
\pgfpathlineto{\pgfqpoint{1.012260in}{0.592243in}}%
\pgfpathlineto{\pgfqpoint{1.028975in}{0.569781in}}%
\pgfpathlineto{\pgfqpoint{1.045690in}{0.592243in}}%
\pgfpathlineto{\pgfqpoint{1.057920in}{0.592243in}}%
\pgfpathlineto{\pgfqpoint{1.057920in}{0.592243in}}%
\pgfpathclose%
\pgfpathmoveto{\pgfqpoint{1.124115in}{0.563315in}}%
\pgfpathlineto{\pgfqpoint{1.124115in}{0.558254in}}%
\pgfpathlineto{\pgfqpoint{1.076491in}{0.558254in}}%
\pgfpathquadraticcurveto{\pgfqpoint{1.077175in}{0.547554in}}{\pgfqpoint{1.082939in}{0.541953in}}%
\pgfpathquadraticcurveto{\pgfqpoint{1.088703in}{0.536351in}}{\pgfqpoint{1.099006in}{0.536351in}}%
\pgfpathquadraticcurveto{\pgfqpoint{1.104968in}{0.536351in}}{\pgfqpoint{1.110570in}{0.537810in}}%
\pgfpathquadraticcurveto{\pgfqpoint{1.116172in}{0.539269in}}{\pgfqpoint{1.121701in}{0.542205in}}%
\pgfpathlineto{\pgfqpoint{1.121701in}{0.532406in}}%
\pgfpathquadraticcurveto{\pgfqpoint{1.116118in}{0.530047in}}{\pgfqpoint{1.110264in}{0.528804in}}%
\pgfpathquadraticcurveto{\pgfqpoint{1.104410in}{0.527561in}}{\pgfqpoint{1.098394in}{0.527561in}}%
\pgfpathquadraticcurveto{\pgfqpoint{1.083299in}{0.527561in}}{\pgfqpoint{1.074492in}{0.536333in}}%
\pgfpathquadraticcurveto{\pgfqpoint{1.065684in}{0.545123in}}{\pgfqpoint{1.065684in}{0.560109in}}%
\pgfpathquadraticcurveto{\pgfqpoint{1.065684in}{0.575581in}}{\pgfqpoint{1.074041in}{0.584659in}}%
\pgfpathquadraticcurveto{\pgfqpoint{1.082399in}{0.593756in}}{\pgfqpoint{1.096592in}{0.593756in}}%
\pgfpathquadraticcurveto{\pgfqpoint{1.109309in}{0.593756in}}{\pgfqpoint{1.116712in}{0.585560in}}%
\pgfpathquadraticcurveto{\pgfqpoint{1.124115in}{0.577383in}}{\pgfqpoint{1.124115in}{0.563315in}}%
\pgfpathlineto{\pgfqpoint{1.124115in}{0.563315in}}%
\pgfpathclose%
\pgfpathmoveto{\pgfqpoint{1.113758in}{0.566359in}}%
\pgfpathquadraticcurveto{\pgfqpoint{1.113650in}{0.574843in}}{\pgfqpoint{1.109003in}{0.579904in}}%
\pgfpathquadraticcurveto{\pgfqpoint{1.104356in}{0.584984in}}{\pgfqpoint{1.096701in}{0.584984in}}%
\pgfpathquadraticcurveto{\pgfqpoint{1.088037in}{0.584984in}}{\pgfqpoint{1.082831in}{0.580084in}}%
\pgfpathquadraticcurveto{\pgfqpoint{1.077626in}{0.575185in}}{\pgfqpoint{1.076833in}{0.566287in}}%
\pgfpathlineto{\pgfqpoint{1.113758in}{0.566359in}}%
\pgfpathlineto{\pgfqpoint{1.113758in}{0.566359in}}%
\pgfpathclose%
\pgfpathmoveto{\pgfqpoint{1.181304in}{0.590387in}}%
\pgfpathlineto{\pgfqpoint{1.181304in}{0.580589in}}%
\pgfpathquadraticcurveto{\pgfqpoint{1.176927in}{0.582840in}}{\pgfqpoint{1.172189in}{0.583957in}}%
\pgfpathquadraticcurveto{\pgfqpoint{1.167470in}{0.585092in}}{\pgfqpoint{1.162391in}{0.585092in}}%
\pgfpathquadraticcurveto{\pgfqpoint{1.154682in}{0.585092in}}{\pgfqpoint{1.150827in}{0.582732in}}%
\pgfpathquadraticcurveto{\pgfqpoint{1.146972in}{0.580373in}}{\pgfqpoint{1.146972in}{0.575635in}}%
\pgfpathquadraticcurveto{\pgfqpoint{1.146972in}{0.572033in}}{\pgfqpoint{1.149728in}{0.569980in}}%
\pgfpathquadraticcurveto{\pgfqpoint{1.152484in}{0.567926in}}{\pgfqpoint{1.160824in}{0.566071in}}%
\pgfpathlineto{\pgfqpoint{1.164372in}{0.565278in}}%
\pgfpathquadraticcurveto{\pgfqpoint{1.175396in}{0.562919in}}{\pgfqpoint{1.180043in}{0.558614in}}%
\pgfpathquadraticcurveto{\pgfqpoint{1.184690in}{0.554309in}}{\pgfqpoint{1.184690in}{0.546600in}}%
\pgfpathquadraticcurveto{\pgfqpoint{1.184690in}{0.537810in}}{\pgfqpoint{1.177737in}{0.532676in}}%
\pgfpathquadraticcurveto{\pgfqpoint{1.170785in}{0.527561in}}{\pgfqpoint{1.158626in}{0.527561in}}%
\pgfpathquadraticcurveto{\pgfqpoint{1.153565in}{0.527561in}}{\pgfqpoint{1.148071in}{0.528552in}}%
\pgfpathquadraticcurveto{\pgfqpoint{1.142577in}{0.529542in}}{\pgfqpoint{1.136507in}{0.531506in}}%
\pgfpathlineto{\pgfqpoint{1.136507in}{0.542205in}}%
\pgfpathquadraticcurveto{\pgfqpoint{1.142253in}{0.539215in}}{\pgfqpoint{1.147819in}{0.537720in}}%
\pgfpathquadraticcurveto{\pgfqpoint{1.153385in}{0.536243in}}{\pgfqpoint{1.158860in}{0.536243in}}%
\pgfpathquadraticcurveto{\pgfqpoint{1.166173in}{0.536243in}}{\pgfqpoint{1.170100in}{0.538746in}}%
\pgfpathquadraticcurveto{\pgfqpoint{1.174045in}{0.541250in}}{\pgfqpoint{1.174045in}{0.545807in}}%
\pgfpathquadraticcurveto{\pgfqpoint{1.174045in}{0.550022in}}{\pgfqpoint{1.171199in}{0.552274in}}%
\pgfpathquadraticcurveto{\pgfqpoint{1.168371in}{0.554525in}}{\pgfqpoint{1.158734in}{0.556614in}}%
\pgfpathlineto{\pgfqpoint{1.155132in}{0.557461in}}%
\pgfpathquadraticcurveto{\pgfqpoint{1.145513in}{0.559478in}}{\pgfqpoint{1.141227in}{0.563675in}}%
\pgfpathquadraticcurveto{\pgfqpoint{1.136958in}{0.567872in}}{\pgfqpoint{1.136958in}{0.575185in}}%
\pgfpathquadraticcurveto{\pgfqpoint{1.136958in}{0.584083in}}{\pgfqpoint{1.143262in}{0.588910in}}%
\pgfpathquadraticcurveto{\pgfqpoint{1.149566in}{0.593756in}}{\pgfqpoint{1.161166in}{0.593756in}}%
\pgfpathquadraticcurveto{\pgfqpoint{1.166894in}{0.593756in}}{\pgfqpoint{1.171955in}{0.592909in}}%
\pgfpathquadraticcurveto{\pgfqpoint{1.177035in}{0.592080in}}{\pgfqpoint{1.181304in}{0.590387in}}%
\pgfpathlineto{\pgfqpoint{1.181304in}{0.590387in}}%
\pgfpathclose%
\pgfusepath{fill}%
\end{pgfscope}%
\begin{pgfscope}%
\pgfsetbuttcap%
\pgfsetroundjoin%
\definecolor{currentfill}{rgb}{0.309804,0.372549,0.419608}%
\pgfsetfillcolor{currentfill}%
\pgfsetlinewidth{0.000000pt}%
\definecolor{currentstroke}{rgb}{0.309804,0.372549,0.419608}%
\pgfsetstrokecolor{currentstroke}%
\pgfsetdash{}{0pt}%
\pgfpathmoveto{\pgfqpoint{1.299252in}{0.582678in}}%
\pgfpathlineto{\pgfqpoint{1.299252in}{0.616793in}}%
\pgfpathlineto{\pgfqpoint{1.309609in}{0.616793in}}%
\pgfpathlineto{\pgfqpoint{1.309609in}{0.529200in}}%
\pgfpathlineto{\pgfqpoint{1.299252in}{0.529200in}}%
\pgfpathlineto{\pgfqpoint{1.299252in}{0.538656in}}%
\pgfpathquadraticcurveto{\pgfqpoint{1.295992in}{0.533037in}}{\pgfqpoint{1.291002in}{0.530299in}}%
\pgfpathquadraticcurveto{\pgfqpoint{1.286031in}{0.527561in}}{\pgfqpoint{1.279042in}{0.527561in}}%
\pgfpathquadraticcurveto{\pgfqpoint{1.267622in}{0.527561in}}{\pgfqpoint{1.260436in}{0.536675in}}%
\pgfpathquadraticcurveto{\pgfqpoint{1.253267in}{0.545807in}}{\pgfqpoint{1.253267in}{0.560667in}}%
\pgfpathquadraticcurveto{\pgfqpoint{1.253267in}{0.575527in}}{\pgfqpoint{1.260436in}{0.584641in}}%
\pgfpathquadraticcurveto{\pgfqpoint{1.267622in}{0.593756in}}{\pgfqpoint{1.279042in}{0.593756in}}%
\pgfpathquadraticcurveto{\pgfqpoint{1.286031in}{0.593756in}}{\pgfqpoint{1.291002in}{0.591018in}}%
\pgfpathquadraticcurveto{\pgfqpoint{1.295992in}{0.588298in}}{\pgfqpoint{1.299252in}{0.582678in}}%
\pgfpathlineto{\pgfqpoint{1.299252in}{0.582678in}}%
\pgfpathclose%
\pgfpathmoveto{\pgfqpoint{1.263966in}{0.560667in}}%
\pgfpathquadraticcurveto{\pgfqpoint{1.263966in}{0.549248in}}{\pgfqpoint{1.268667in}{0.542745in}}%
\pgfpathquadraticcurveto{\pgfqpoint{1.273368in}{0.536243in}}{\pgfqpoint{1.281582in}{0.536243in}}%
\pgfpathquadraticcurveto{\pgfqpoint{1.289795in}{0.536243in}}{\pgfqpoint{1.294515in}{0.542745in}}%
\pgfpathquadraticcurveto{\pgfqpoint{1.299252in}{0.549248in}}{\pgfqpoint{1.299252in}{0.560667in}}%
\pgfpathquadraticcurveto{\pgfqpoint{1.299252in}{0.572087in}}{\pgfqpoint{1.294515in}{0.578589in}}%
\pgfpathquadraticcurveto{\pgfqpoint{1.289795in}{0.585092in}}{\pgfqpoint{1.281582in}{0.585092in}}%
\pgfpathquadraticcurveto{\pgfqpoint{1.273368in}{0.585092in}}{\pgfqpoint{1.268667in}{0.578589in}}%
\pgfpathquadraticcurveto{\pgfqpoint{1.263966in}{0.572087in}}{\pgfqpoint{1.263966in}{0.560667in}}%
\pgfpathlineto{\pgfqpoint{1.263966in}{0.560667in}}%
\pgfpathclose%
\pgfpathmoveto{\pgfqpoint{1.330953in}{0.592243in}}%
\pgfpathlineto{\pgfqpoint{1.341310in}{0.592243in}}%
\pgfpathlineto{\pgfqpoint{1.341310in}{0.529200in}}%
\pgfpathlineto{\pgfqpoint{1.330953in}{0.529200in}}%
\pgfpathlineto{\pgfqpoint{1.330953in}{0.592243in}}%
\pgfpathlineto{\pgfqpoint{1.330953in}{0.592243in}}%
\pgfpathclose%
\pgfpathmoveto{\pgfqpoint{1.330953in}{0.616793in}}%
\pgfpathlineto{\pgfqpoint{1.341310in}{0.616793in}}%
\pgfpathlineto{\pgfqpoint{1.341310in}{0.603662in}}%
\pgfpathlineto{\pgfqpoint{1.330953in}{0.603662in}}%
\pgfpathlineto{\pgfqpoint{1.330953in}{0.616793in}}%
\pgfpathlineto{\pgfqpoint{1.330953in}{0.616793in}}%
\pgfpathclose%
\pgfpathmoveto{\pgfqpoint{1.355558in}{0.592243in}}%
\pgfpathlineto{\pgfqpoint{1.366527in}{0.592243in}}%
\pgfpathlineto{\pgfqpoint{1.386232in}{0.539341in}}%
\pgfpathlineto{\pgfqpoint{1.405938in}{0.592243in}}%
\pgfpathlineto{\pgfqpoint{1.416907in}{0.592243in}}%
\pgfpathlineto{\pgfqpoint{1.393257in}{0.529200in}}%
\pgfpathlineto{\pgfqpoint{1.379190in}{0.529200in}}%
\pgfpathlineto{\pgfqpoint{1.355558in}{0.592243in}}%
\pgfpathlineto{\pgfqpoint{1.355558in}{0.592243in}}%
\pgfpathclose%
\pgfpathmoveto{\pgfqpoint{1.431209in}{0.592243in}}%
\pgfpathlineto{\pgfqpoint{1.441566in}{0.592243in}}%
\pgfpathlineto{\pgfqpoint{1.441566in}{0.529200in}}%
\pgfpathlineto{\pgfqpoint{1.431209in}{0.529200in}}%
\pgfpathlineto{\pgfqpoint{1.431209in}{0.592243in}}%
\pgfpathlineto{\pgfqpoint{1.431209in}{0.592243in}}%
\pgfpathclose%
\pgfpathmoveto{\pgfqpoint{1.431209in}{0.616793in}}%
\pgfpathlineto{\pgfqpoint{1.441566in}{0.616793in}}%
\pgfpathlineto{\pgfqpoint{1.441566in}{0.603662in}}%
\pgfpathlineto{\pgfqpoint{1.431209in}{0.603662in}}%
\pgfpathlineto{\pgfqpoint{1.431209in}{0.616793in}}%
\pgfpathlineto{\pgfqpoint{1.431209in}{0.616793in}}%
\pgfpathclose%
\pgfpathmoveto{\pgfqpoint{1.504716in}{0.582678in}}%
\pgfpathlineto{\pgfqpoint{1.504716in}{0.616793in}}%
\pgfpathlineto{\pgfqpoint{1.515073in}{0.616793in}}%
\pgfpathlineto{\pgfqpoint{1.515073in}{0.529200in}}%
\pgfpathlineto{\pgfqpoint{1.504716in}{0.529200in}}%
\pgfpathlineto{\pgfqpoint{1.504716in}{0.538656in}}%
\pgfpathquadraticcurveto{\pgfqpoint{1.501456in}{0.533037in}}{\pgfqpoint{1.496467in}{0.530299in}}%
\pgfpathquadraticcurveto{\pgfqpoint{1.491496in}{0.527561in}}{\pgfqpoint{1.484507in}{0.527561in}}%
\pgfpathquadraticcurveto{\pgfqpoint{1.473087in}{0.527561in}}{\pgfqpoint{1.465900in}{0.536675in}}%
\pgfpathquadraticcurveto{\pgfqpoint{1.458731in}{0.545807in}}{\pgfqpoint{1.458731in}{0.560667in}}%
\pgfpathquadraticcurveto{\pgfqpoint{1.458731in}{0.575527in}}{\pgfqpoint{1.465900in}{0.584641in}}%
\pgfpathquadraticcurveto{\pgfqpoint{1.473087in}{0.593756in}}{\pgfqpoint{1.484507in}{0.593756in}}%
\pgfpathquadraticcurveto{\pgfqpoint{1.491496in}{0.593756in}}{\pgfqpoint{1.496467in}{0.591018in}}%
\pgfpathquadraticcurveto{\pgfqpoint{1.501456in}{0.588298in}}{\pgfqpoint{1.504716in}{0.582678in}}%
\pgfpathlineto{\pgfqpoint{1.504716in}{0.582678in}}%
\pgfpathclose%
\pgfpathmoveto{\pgfqpoint{1.469431in}{0.560667in}}%
\pgfpathquadraticcurveto{\pgfqpoint{1.469431in}{0.549248in}}{\pgfqpoint{1.474132in}{0.542745in}}%
\pgfpathquadraticcurveto{\pgfqpoint{1.478833in}{0.536243in}}{\pgfqpoint{1.487047in}{0.536243in}}%
\pgfpathquadraticcurveto{\pgfqpoint{1.495260in}{0.536243in}}{\pgfqpoint{1.499979in}{0.542745in}}%
\pgfpathquadraticcurveto{\pgfqpoint{1.504716in}{0.549248in}}{\pgfqpoint{1.504716in}{0.560667in}}%
\pgfpathquadraticcurveto{\pgfqpoint{1.504716in}{0.572087in}}{\pgfqpoint{1.499979in}{0.578589in}}%
\pgfpathquadraticcurveto{\pgfqpoint{1.495260in}{0.585092in}}{\pgfqpoint{1.487047in}{0.585092in}}%
\pgfpathquadraticcurveto{\pgfqpoint{1.478833in}{0.585092in}}{\pgfqpoint{1.474132in}{0.578589in}}%
\pgfpathquadraticcurveto{\pgfqpoint{1.469431in}{0.572087in}}{\pgfqpoint{1.469431in}{0.560667in}}%
\pgfpathlineto{\pgfqpoint{1.469431in}{0.560667in}}%
\pgfpathclose%
\pgfpathmoveto{\pgfqpoint{1.590346in}{0.563315in}}%
\pgfpathlineto{\pgfqpoint{1.590346in}{0.558254in}}%
\pgfpathlineto{\pgfqpoint{1.542722in}{0.558254in}}%
\pgfpathquadraticcurveto{\pgfqpoint{1.543407in}{0.547554in}}{\pgfqpoint{1.549170in}{0.541953in}}%
\pgfpathquadraticcurveto{\pgfqpoint{1.554934in}{0.536351in}}{\pgfqpoint{1.565237in}{0.536351in}}%
\pgfpathquadraticcurveto{\pgfqpoint{1.571199in}{0.536351in}}{\pgfqpoint{1.576801in}{0.537810in}}%
\pgfpathquadraticcurveto{\pgfqpoint{1.582403in}{0.539269in}}{\pgfqpoint{1.587933in}{0.542205in}}%
\pgfpathlineto{\pgfqpoint{1.587933in}{0.532406in}}%
\pgfpathquadraticcurveto{\pgfqpoint{1.582349in}{0.530047in}}{\pgfqpoint{1.576495in}{0.528804in}}%
\pgfpathquadraticcurveto{\pgfqpoint{1.570641in}{0.527561in}}{\pgfqpoint{1.564625in}{0.527561in}}%
\pgfpathquadraticcurveto{\pgfqpoint{1.549531in}{0.527561in}}{\pgfqpoint{1.540723in}{0.536333in}}%
\pgfpathquadraticcurveto{\pgfqpoint{1.531915in}{0.545123in}}{\pgfqpoint{1.531915in}{0.560109in}}%
\pgfpathquadraticcurveto{\pgfqpoint{1.531915in}{0.575581in}}{\pgfqpoint{1.540272in}{0.584659in}}%
\pgfpathquadraticcurveto{\pgfqpoint{1.548630in}{0.593756in}}{\pgfqpoint{1.562824in}{0.593756in}}%
\pgfpathquadraticcurveto{\pgfqpoint{1.575540in}{0.593756in}}{\pgfqpoint{1.582943in}{0.585560in}}%
\pgfpathquadraticcurveto{\pgfqpoint{1.590346in}{0.577383in}}{\pgfqpoint{1.590346in}{0.563315in}}%
\pgfpathlineto{\pgfqpoint{1.590346in}{0.563315in}}%
\pgfpathclose%
\pgfpathmoveto{\pgfqpoint{1.579989in}{0.566359in}}%
\pgfpathquadraticcurveto{\pgfqpoint{1.579881in}{0.574843in}}{\pgfqpoint{1.575234in}{0.579904in}}%
\pgfpathquadraticcurveto{\pgfqpoint{1.570587in}{0.584984in}}{\pgfqpoint{1.562932in}{0.584984in}}%
\pgfpathquadraticcurveto{\pgfqpoint{1.554268in}{0.584984in}}{\pgfqpoint{1.549062in}{0.580084in}}%
\pgfpathquadraticcurveto{\pgfqpoint{1.543857in}{0.575185in}}{\pgfqpoint{1.543064in}{0.566287in}}%
\pgfpathlineto{\pgfqpoint{1.579989in}{0.566359in}}%
\pgfpathlineto{\pgfqpoint{1.579989in}{0.566359in}}%
\pgfpathclose%
\pgfusepath{fill}%
\end{pgfscope}%
\begin{pgfscope}%
\pgfsetbuttcap%
\pgfsetroundjoin%
\definecolor{currentfill}{rgb}{0.309804,0.372549,0.419608}%
\pgfsetfillcolor{currentfill}%
\pgfsetlinewidth{0.000000pt}%
\definecolor{currentstroke}{rgb}{0.309804,0.372549,0.419608}%
\pgfsetstrokecolor{currentstroke}%
\pgfsetdash{}{0pt}%
\pgfpathmoveto{\pgfqpoint{1.719537in}{0.560667in}}%
\pgfpathquadraticcurveto{\pgfqpoint{1.719537in}{0.572087in}}{\pgfqpoint{1.714835in}{0.578589in}}%
\pgfpathquadraticcurveto{\pgfqpoint{1.710134in}{0.585092in}}{\pgfqpoint{1.701921in}{0.585092in}}%
\pgfpathquadraticcurveto{\pgfqpoint{1.693689in}{0.585092in}}{\pgfqpoint{1.688988in}{0.578589in}}%
\pgfpathquadraticcurveto{\pgfqpoint{1.684287in}{0.572087in}}{\pgfqpoint{1.684287in}{0.560667in}}%
\pgfpathquadraticcurveto{\pgfqpoint{1.684287in}{0.549248in}}{\pgfqpoint{1.688988in}{0.542745in}}%
\pgfpathquadraticcurveto{\pgfqpoint{1.693689in}{0.536243in}}{\pgfqpoint{1.701921in}{0.536243in}}%
\pgfpathquadraticcurveto{\pgfqpoint{1.710134in}{0.536243in}}{\pgfqpoint{1.714835in}{0.542745in}}%
\pgfpathquadraticcurveto{\pgfqpoint{1.719537in}{0.549248in}}{\pgfqpoint{1.719537in}{0.560667in}}%
\pgfpathlineto{\pgfqpoint{1.719537in}{0.560667in}}%
\pgfpathclose%
\pgfpathmoveto{\pgfqpoint{1.684287in}{0.582678in}}%
\pgfpathquadraticcurveto{\pgfqpoint{1.687565in}{0.588298in}}{\pgfqpoint{1.692536in}{0.591018in}}%
\pgfpathquadraticcurveto{\pgfqpoint{1.697526in}{0.593756in}}{\pgfqpoint{1.704442in}{0.593756in}}%
\pgfpathquadraticcurveto{\pgfqpoint{1.715934in}{0.593756in}}{\pgfqpoint{1.723103in}{0.584641in}}%
\pgfpathquadraticcurveto{\pgfqpoint{1.730290in}{0.575527in}}{\pgfqpoint{1.730290in}{0.560667in}}%
\pgfpathquadraticcurveto{\pgfqpoint{1.730290in}{0.545807in}}{\pgfqpoint{1.723103in}{0.536675in}}%
\pgfpathquadraticcurveto{\pgfqpoint{1.715934in}{0.527561in}}{\pgfqpoint{1.704442in}{0.527561in}}%
\pgfpathquadraticcurveto{\pgfqpoint{1.697526in}{0.527561in}}{\pgfqpoint{1.692536in}{0.530299in}}%
\pgfpathquadraticcurveto{\pgfqpoint{1.687565in}{0.533037in}}{\pgfqpoint{1.684287in}{0.538656in}}%
\pgfpathlineto{\pgfqpoint{1.684287in}{0.529200in}}%
\pgfpathlineto{\pgfqpoint{1.673876in}{0.529200in}}%
\pgfpathlineto{\pgfqpoint{1.673876in}{0.616793in}}%
\pgfpathlineto{\pgfqpoint{1.684287in}{0.616793in}}%
\pgfpathlineto{\pgfqpoint{1.684287in}{0.582678in}}%
\pgfpathlineto{\pgfqpoint{1.684287in}{0.582678in}}%
\pgfpathclose%
\pgfpathmoveto{\pgfqpoint{1.773681in}{0.523346in}}%
\pgfpathquadraticcurveto{\pgfqpoint{1.769304in}{0.512088in}}{\pgfqpoint{1.765125in}{0.508666in}}%
\pgfpathquadraticcurveto{\pgfqpoint{1.760965in}{0.505226in}}{\pgfqpoint{1.753994in}{0.505226in}}%
\pgfpathlineto{\pgfqpoint{1.745708in}{0.505226in}}%
\pgfpathlineto{\pgfqpoint{1.745708in}{0.513890in}}%
\pgfpathlineto{\pgfqpoint{1.751796in}{0.513890in}}%
\pgfpathquadraticcurveto{\pgfqpoint{1.756065in}{0.513890in}}{\pgfqpoint{1.758425in}{0.515925in}}%
\pgfpathquadraticcurveto{\pgfqpoint{1.760802in}{0.517942in}}{\pgfqpoint{1.763666in}{0.525489in}}%
\pgfpathlineto{\pgfqpoint{1.765522in}{0.530209in}}%
\pgfpathlineto{\pgfqpoint{1.740034in}{0.592243in}}%
\pgfpathlineto{\pgfqpoint{1.751004in}{0.592243in}}%
\pgfpathlineto{\pgfqpoint{1.770709in}{0.542943in}}%
\pgfpathlineto{\pgfqpoint{1.790414in}{0.592243in}}%
\pgfpathlineto{\pgfqpoint{1.801384in}{0.592243in}}%
\pgfpathlineto{\pgfqpoint{1.773681in}{0.523346in}}%
\pgfpathlineto{\pgfqpoint{1.773681in}{0.523346in}}%
\pgfpathclose%
\pgfusepath{fill}%
\end{pgfscope}%
\begin{pgfscope}%
\pgfsetbuttcap%
\pgfsetroundjoin%
\definecolor{currentfill}{rgb}{0.309804,0.372549,0.419608}%
\pgfsetfillcolor{currentfill}%
\pgfsetlinewidth{0.000000pt}%
\definecolor{currentstroke}{rgb}{0.309804,0.372549,0.419608}%
\pgfsetstrokecolor{currentstroke}%
\pgfsetdash{}{0pt}%
\pgfpathmoveto{\pgfqpoint{1.891023in}{0.610147in}}%
\pgfpathlineto{\pgfqpoint{1.891023in}{0.592243in}}%
\pgfpathlineto{\pgfqpoint{1.912350in}{0.592243in}}%
\pgfpathlineto{\pgfqpoint{1.912350in}{0.584191in}}%
\pgfpathlineto{\pgfqpoint{1.891023in}{0.584191in}}%
\pgfpathlineto{\pgfqpoint{1.891023in}{0.549968in}}%
\pgfpathquadraticcurveto{\pgfqpoint{1.891023in}{0.542259in}}{\pgfqpoint{1.893131in}{0.540061in}}%
\pgfpathquadraticcurveto{\pgfqpoint{1.895238in}{0.537864in}}{\pgfqpoint{1.901722in}{0.537864in}}%
\pgfpathlineto{\pgfqpoint{1.912350in}{0.537864in}}%
\pgfpathlineto{\pgfqpoint{1.912350in}{0.529200in}}%
\pgfpathlineto{\pgfqpoint{1.901722in}{0.529200in}}%
\pgfpathquadraticcurveto{\pgfqpoint{1.889726in}{0.529200in}}{\pgfqpoint{1.885169in}{0.533667in}}%
\pgfpathquadraticcurveto{\pgfqpoint{1.880612in}{0.538152in}}{\pgfqpoint{1.880612in}{0.549968in}}%
\pgfpathlineto{\pgfqpoint{1.880612in}{0.584191in}}%
\pgfpathlineto{\pgfqpoint{1.873011in}{0.584191in}}%
\pgfpathlineto{\pgfqpoint{1.873011in}{0.592243in}}%
\pgfpathlineto{\pgfqpoint{1.880612in}{0.592243in}}%
\pgfpathlineto{\pgfqpoint{1.880612in}{0.610147in}}%
\pgfpathlineto{\pgfqpoint{1.891023in}{0.610147in}}%
\pgfpathlineto{\pgfqpoint{1.891023in}{0.610147in}}%
\pgfpathclose%
\pgfpathmoveto{\pgfqpoint{1.978382in}{0.567260in}}%
\pgfpathlineto{\pgfqpoint{1.978382in}{0.529200in}}%
\pgfpathlineto{\pgfqpoint{1.968025in}{0.529200in}}%
\pgfpathlineto{\pgfqpoint{1.968025in}{0.566917in}}%
\pgfpathquadraticcurveto{\pgfqpoint{1.968025in}{0.575869in}}{\pgfqpoint{1.964531in}{0.580300in}}%
\pgfpathquadraticcurveto{\pgfqpoint{1.961036in}{0.584749in}}{\pgfqpoint{1.954066in}{0.584749in}}%
\pgfpathquadraticcurveto{\pgfqpoint{1.945672in}{0.584749in}}{\pgfqpoint{1.940827in}{0.579400in}}%
\pgfpathquadraticcurveto{\pgfqpoint{1.935982in}{0.574068in}}{\pgfqpoint{1.935982in}{0.564828in}}%
\pgfpathlineto{\pgfqpoint{1.935982in}{0.529200in}}%
\pgfpathlineto{\pgfqpoint{1.925570in}{0.529200in}}%
\pgfpathlineto{\pgfqpoint{1.925570in}{0.616793in}}%
\pgfpathlineto{\pgfqpoint{1.935982in}{0.616793in}}%
\pgfpathlineto{\pgfqpoint{1.935982in}{0.582444in}}%
\pgfpathquadraticcurveto{\pgfqpoint{1.939710in}{0.588136in}}{\pgfqpoint{1.944735in}{0.590946in}}%
\pgfpathquadraticcurveto{\pgfqpoint{1.949779in}{0.593756in}}{\pgfqpoint{1.956371in}{0.593756in}}%
\pgfpathquadraticcurveto{\pgfqpoint{1.967233in}{0.593756in}}{\pgfqpoint{1.972798in}{0.587037in}}%
\pgfpathquadraticcurveto{\pgfqpoint{1.978382in}{0.580318in}}{\pgfqpoint{1.978382in}{0.567260in}}%
\pgfpathlineto{\pgfqpoint{1.978382in}{0.567260in}}%
\pgfpathclose%
\pgfpathmoveto{\pgfqpoint{1.999024in}{0.592243in}}%
\pgfpathlineto{\pgfqpoint{2.009381in}{0.592243in}}%
\pgfpathlineto{\pgfqpoint{2.009381in}{0.529200in}}%
\pgfpathlineto{\pgfqpoint{1.999024in}{0.529200in}}%
\pgfpathlineto{\pgfqpoint{1.999024in}{0.592243in}}%
\pgfpathlineto{\pgfqpoint{1.999024in}{0.592243in}}%
\pgfpathclose%
\pgfpathmoveto{\pgfqpoint{1.999024in}{0.616793in}}%
\pgfpathlineto{\pgfqpoint{2.009381in}{0.616793in}}%
\pgfpathlineto{\pgfqpoint{2.009381in}{0.603662in}}%
\pgfpathlineto{\pgfqpoint{1.999024in}{0.603662in}}%
\pgfpathlineto{\pgfqpoint{1.999024in}{0.616793in}}%
\pgfpathlineto{\pgfqpoint{1.999024in}{0.616793in}}%
\pgfpathclose%
\pgfpathmoveto{\pgfqpoint{2.071235in}{0.590387in}}%
\pgfpathlineto{\pgfqpoint{2.071235in}{0.580589in}}%
\pgfpathquadraticcurveto{\pgfqpoint{2.066858in}{0.582840in}}{\pgfqpoint{2.062121in}{0.583957in}}%
\pgfpathquadraticcurveto{\pgfqpoint{2.057401in}{0.585092in}}{\pgfqpoint{2.052322in}{0.585092in}}%
\pgfpathquadraticcurveto{\pgfqpoint{2.044613in}{0.585092in}}{\pgfqpoint{2.040758in}{0.582732in}}%
\pgfpathquadraticcurveto{\pgfqpoint{2.036904in}{0.580373in}}{\pgfqpoint{2.036904in}{0.575635in}}%
\pgfpathquadraticcurveto{\pgfqpoint{2.036904in}{0.572033in}}{\pgfqpoint{2.039659in}{0.569980in}}%
\pgfpathquadraticcurveto{\pgfqpoint{2.042415in}{0.567926in}}{\pgfqpoint{2.050755in}{0.566071in}}%
\pgfpathlineto{\pgfqpoint{2.054303in}{0.565278in}}%
\pgfpathquadraticcurveto{\pgfqpoint{2.065327in}{0.562919in}}{\pgfqpoint{2.069974in}{0.558614in}}%
\pgfpathquadraticcurveto{\pgfqpoint{2.074621in}{0.554309in}}{\pgfqpoint{2.074621in}{0.546600in}}%
\pgfpathquadraticcurveto{\pgfqpoint{2.074621in}{0.537810in}}{\pgfqpoint{2.067668in}{0.532676in}}%
\pgfpathquadraticcurveto{\pgfqpoint{2.060716in}{0.527561in}}{\pgfqpoint{2.048557in}{0.527561in}}%
\pgfpathquadraticcurveto{\pgfqpoint{2.043496in}{0.527561in}}{\pgfqpoint{2.038002in}{0.528552in}}%
\pgfpathquadraticcurveto{\pgfqpoint{2.032509in}{0.529542in}}{\pgfqpoint{2.026439in}{0.531506in}}%
\pgfpathlineto{\pgfqpoint{2.026439in}{0.542205in}}%
\pgfpathquadraticcurveto{\pgfqpoint{2.032184in}{0.539215in}}{\pgfqpoint{2.037750in}{0.537720in}}%
\pgfpathquadraticcurveto{\pgfqpoint{2.043316in}{0.536243in}}{\pgfqpoint{2.048792in}{0.536243in}}%
\pgfpathquadraticcurveto{\pgfqpoint{2.056105in}{0.536243in}}{\pgfqpoint{2.060031in}{0.538746in}}%
\pgfpathquadraticcurveto{\pgfqpoint{2.063976in}{0.541250in}}{\pgfqpoint{2.063976in}{0.545807in}}%
\pgfpathquadraticcurveto{\pgfqpoint{2.063976in}{0.550022in}}{\pgfqpoint{2.061130in}{0.552274in}}%
\pgfpathquadraticcurveto{\pgfqpoint{2.058302in}{0.554525in}}{\pgfqpoint{2.048666in}{0.556614in}}%
\pgfpathlineto{\pgfqpoint{2.045063in}{0.557461in}}%
\pgfpathquadraticcurveto{\pgfqpoint{2.035445in}{0.559478in}}{\pgfqpoint{2.031158in}{0.563675in}}%
\pgfpathquadraticcurveto{\pgfqpoint{2.026889in}{0.567872in}}{\pgfqpoint{2.026889in}{0.575185in}}%
\pgfpathquadraticcurveto{\pgfqpoint{2.026889in}{0.584083in}}{\pgfqpoint{2.033193in}{0.588910in}}%
\pgfpathquadraticcurveto{\pgfqpoint{2.039497in}{0.593756in}}{\pgfqpoint{2.051097in}{0.593756in}}%
\pgfpathquadraticcurveto{\pgfqpoint{2.056825in}{0.593756in}}{\pgfqpoint{2.061886in}{0.592909in}}%
\pgfpathquadraticcurveto{\pgfqpoint{2.066966in}{0.592080in}}{\pgfqpoint{2.071235in}{0.590387in}}%
\pgfpathlineto{\pgfqpoint{2.071235in}{0.590387in}}%
\pgfpathclose%
\pgfusepath{fill}%
\end{pgfscope}%
\begin{pgfscope}%
\pgfsetbuttcap%
\pgfsetroundjoin%
\definecolor{currentfill}{rgb}{0.309804,0.372549,0.419608}%
\pgfsetfillcolor{currentfill}%
\pgfsetlinewidth{0.000000pt}%
\definecolor{currentstroke}{rgb}{0.309804,0.372549,0.419608}%
\pgfsetstrokecolor{currentstroke}%
\pgfsetdash{}{0pt}%
\pgfpathmoveto{\pgfqpoint{2.167291in}{0.538656in}}%
\pgfpathlineto{\pgfqpoint{2.167291in}{0.505226in}}%
\pgfpathlineto{\pgfqpoint{2.156880in}{0.505226in}}%
\pgfpathlineto{\pgfqpoint{2.156880in}{0.592243in}}%
\pgfpathlineto{\pgfqpoint{2.167291in}{0.592243in}}%
\pgfpathlineto{\pgfqpoint{2.167291in}{0.582678in}}%
\pgfpathquadraticcurveto{\pgfqpoint{2.170569in}{0.588298in}}{\pgfqpoint{2.175541in}{0.591018in}}%
\pgfpathquadraticcurveto{\pgfqpoint{2.180530in}{0.593756in}}{\pgfqpoint{2.187447in}{0.593756in}}%
\pgfpathquadraticcurveto{\pgfqpoint{2.198939in}{0.593756in}}{\pgfqpoint{2.206107in}{0.584641in}}%
\pgfpathquadraticcurveto{\pgfqpoint{2.213294in}{0.575527in}}{\pgfqpoint{2.213294in}{0.560667in}}%
\pgfpathquadraticcurveto{\pgfqpoint{2.213294in}{0.545807in}}{\pgfqpoint{2.206107in}{0.536675in}}%
\pgfpathquadraticcurveto{\pgfqpoint{2.198939in}{0.527561in}}{\pgfqpoint{2.187447in}{0.527561in}}%
\pgfpathquadraticcurveto{\pgfqpoint{2.180530in}{0.527561in}}{\pgfqpoint{2.175541in}{0.530299in}}%
\pgfpathquadraticcurveto{\pgfqpoint{2.170569in}{0.533037in}}{\pgfqpoint{2.167291in}{0.538656in}}%
\pgfpathlineto{\pgfqpoint{2.167291in}{0.538656in}}%
\pgfpathclose%
\pgfpathmoveto{\pgfqpoint{2.202541in}{0.560667in}}%
\pgfpathquadraticcurveto{\pgfqpoint{2.202541in}{0.572087in}}{\pgfqpoint{2.197840in}{0.578589in}}%
\pgfpathquadraticcurveto{\pgfqpoint{2.193139in}{0.585092in}}{\pgfqpoint{2.184925in}{0.585092in}}%
\pgfpathquadraticcurveto{\pgfqpoint{2.176694in}{0.585092in}}{\pgfqpoint{2.171992in}{0.578589in}}%
\pgfpathquadraticcurveto{\pgfqpoint{2.167291in}{0.572087in}}{\pgfqpoint{2.167291in}{0.560667in}}%
\pgfpathquadraticcurveto{\pgfqpoint{2.167291in}{0.549248in}}{\pgfqpoint{2.171992in}{0.542745in}}%
\pgfpathquadraticcurveto{\pgfqpoint{2.176694in}{0.536243in}}{\pgfqpoint{2.184925in}{0.536243in}}%
\pgfpathquadraticcurveto{\pgfqpoint{2.193139in}{0.536243in}}{\pgfqpoint{2.197840in}{0.542745in}}%
\pgfpathquadraticcurveto{\pgfqpoint{2.202541in}{0.549248in}}{\pgfqpoint{2.202541in}{0.560667in}}%
\pgfpathlineto{\pgfqpoint{2.202541in}{0.560667in}}%
\pgfpathclose%
\pgfpathmoveto{\pgfqpoint{2.259117in}{0.560883in}}%
\pgfpathquadraticcurveto{\pgfqpoint{2.246563in}{0.560883in}}{\pgfqpoint{2.241717in}{0.558019in}}%
\pgfpathquadraticcurveto{\pgfqpoint{2.236872in}{0.555156in}}{\pgfqpoint{2.236872in}{0.548221in}}%
\pgfpathquadraticcurveto{\pgfqpoint{2.236872in}{0.542709in}}{\pgfqpoint{2.240511in}{0.539467in}}%
\pgfpathquadraticcurveto{\pgfqpoint{2.244149in}{0.536243in}}{\pgfqpoint{2.250381in}{0.536243in}}%
\pgfpathquadraticcurveto{\pgfqpoint{2.259009in}{0.536243in}}{\pgfqpoint{2.264215in}{0.542349in}}%
\pgfpathquadraticcurveto{\pgfqpoint{2.269420in}{0.548455in}}{\pgfqpoint{2.269420in}{0.558578in}}%
\pgfpathlineto{\pgfqpoint{2.269420in}{0.560883in}}%
\pgfpathlineto{\pgfqpoint{2.259117in}{0.560883in}}%
\pgfpathlineto{\pgfqpoint{2.259117in}{0.560883in}}%
\pgfpathclose%
\pgfpathmoveto{\pgfqpoint{2.279777in}{0.565170in}}%
\pgfpathlineto{\pgfqpoint{2.279777in}{0.529200in}}%
\pgfpathlineto{\pgfqpoint{2.269420in}{0.529200in}}%
\pgfpathlineto{\pgfqpoint{2.269420in}{0.538764in}}%
\pgfpathquadraticcurveto{\pgfqpoint{2.265872in}{0.533037in}}{\pgfqpoint{2.260576in}{0.530299in}}%
\pgfpathquadraticcurveto{\pgfqpoint{2.255281in}{0.527561in}}{\pgfqpoint{2.247625in}{0.527561in}}%
\pgfpathquadraticcurveto{\pgfqpoint{2.237953in}{0.527561in}}{\pgfqpoint{2.232225in}{0.533001in}}%
\pgfpathquadraticcurveto{\pgfqpoint{2.226515in}{0.538440in}}{\pgfqpoint{2.226515in}{0.547554in}}%
\pgfpathquadraticcurveto{\pgfqpoint{2.226515in}{0.558182in}}{\pgfqpoint{2.233630in}{0.563585in}}%
\pgfpathquadraticcurveto{\pgfqpoint{2.240763in}{0.568989in}}{\pgfqpoint{2.254884in}{0.568989in}}%
\pgfpathlineto{\pgfqpoint{2.269420in}{0.568989in}}%
\pgfpathlineto{\pgfqpoint{2.269420in}{0.570016in}}%
\pgfpathquadraticcurveto{\pgfqpoint{2.269420in}{0.577166in}}{\pgfqpoint{2.264719in}{0.581075in}}%
\pgfpathquadraticcurveto{\pgfqpoint{2.260018in}{0.584984in}}{\pgfqpoint{2.251516in}{0.584984in}}%
\pgfpathquadraticcurveto{\pgfqpoint{2.246112in}{0.584984in}}{\pgfqpoint{2.240979in}{0.583687in}}%
\pgfpathquadraticcurveto{\pgfqpoint{2.235864in}{0.582390in}}{\pgfqpoint{2.231144in}{0.579796in}}%
\pgfpathlineto{\pgfqpoint{2.231144in}{0.589379in}}%
\pgfpathquadraticcurveto{\pgfqpoint{2.236818in}{0.591576in}}{\pgfqpoint{2.242168in}{0.592657in}}%
\pgfpathquadraticcurveto{\pgfqpoint{2.247517in}{0.593756in}}{\pgfqpoint{2.252579in}{0.593756in}}%
\pgfpathquadraticcurveto{\pgfqpoint{2.266268in}{0.593756in}}{\pgfqpoint{2.273023in}{0.586659in}}%
\pgfpathquadraticcurveto{\pgfqpoint{2.279777in}{0.579580in}}{\pgfqpoint{2.279777in}{0.565170in}}%
\pgfpathlineto{\pgfqpoint{2.279777in}{0.565170in}}%
\pgfpathclose%
\pgfpathmoveto{\pgfqpoint{2.301104in}{0.592243in}}%
\pgfpathlineto{\pgfqpoint{2.311461in}{0.592243in}}%
\pgfpathlineto{\pgfqpoint{2.311461in}{0.529200in}}%
\pgfpathlineto{\pgfqpoint{2.301104in}{0.529200in}}%
\pgfpathlineto{\pgfqpoint{2.301104in}{0.592243in}}%
\pgfpathlineto{\pgfqpoint{2.301104in}{0.592243in}}%
\pgfpathclose%
\pgfpathmoveto{\pgfqpoint{2.301104in}{0.616793in}}%
\pgfpathlineto{\pgfqpoint{2.311461in}{0.616793in}}%
\pgfpathlineto{\pgfqpoint{2.311461in}{0.603662in}}%
\pgfpathlineto{\pgfqpoint{2.301104in}{0.603662in}}%
\pgfpathlineto{\pgfqpoint{2.301104in}{0.616793in}}%
\pgfpathlineto{\pgfqpoint{2.301104in}{0.616793in}}%
\pgfpathclose%
\pgfpathmoveto{\pgfqpoint{2.369658in}{0.582570in}}%
\pgfpathquadraticcurveto{\pgfqpoint{2.367911in}{0.583579in}}{\pgfqpoint{2.365857in}{0.584047in}}%
\pgfpathquadraticcurveto{\pgfqpoint{2.363804in}{0.584533in}}{\pgfqpoint{2.361336in}{0.584533in}}%
\pgfpathquadraticcurveto{\pgfqpoint{2.352546in}{0.584533in}}{\pgfqpoint{2.347845in}{0.578823in}}%
\pgfpathquadraticcurveto{\pgfqpoint{2.343144in}{0.573114in}}{\pgfqpoint{2.343144in}{0.562414in}}%
\pgfpathlineto{\pgfqpoint{2.343144in}{0.529200in}}%
\pgfpathlineto{\pgfqpoint{2.332733in}{0.529200in}}%
\pgfpathlineto{\pgfqpoint{2.332733in}{0.592243in}}%
\pgfpathlineto{\pgfqpoint{2.343144in}{0.592243in}}%
\pgfpathlineto{\pgfqpoint{2.343144in}{0.582444in}}%
\pgfpathquadraticcurveto{\pgfqpoint{2.346422in}{0.588190in}}{\pgfqpoint{2.351646in}{0.590964in}}%
\pgfpathquadraticcurveto{\pgfqpoint{2.356887in}{0.593756in}}{\pgfqpoint{2.364380in}{0.593756in}}%
\pgfpathquadraticcurveto{\pgfqpoint{2.365443in}{0.593756in}}{\pgfqpoint{2.366740in}{0.593611in}}%
\pgfpathquadraticcurveto{\pgfqpoint{2.368037in}{0.593485in}}{\pgfqpoint{2.369604in}{0.593197in}}%
\pgfpathlineto{\pgfqpoint{2.369658in}{0.582570in}}%
\pgfpathlineto{\pgfqpoint{2.369658in}{0.582570in}}%
\pgfpathclose%
\pgfpathmoveto{\pgfqpoint{2.390318in}{0.613245in}}%
\pgfpathlineto{\pgfqpoint{2.390318in}{0.581994in}}%
\pgfpathlineto{\pgfqpoint{2.380753in}{0.581994in}}%
\pgfpathlineto{\pgfqpoint{2.380753in}{0.613245in}}%
\pgfpathlineto{\pgfqpoint{2.390318in}{0.613245in}}%
\pgfpathlineto{\pgfqpoint{2.390318in}{0.613245in}}%
\pgfpathclose%
\pgfpathmoveto{\pgfqpoint{2.452388in}{0.590387in}}%
\pgfpathlineto{\pgfqpoint{2.452388in}{0.580589in}}%
\pgfpathquadraticcurveto{\pgfqpoint{2.448011in}{0.582840in}}{\pgfqpoint{2.443273in}{0.583957in}}%
\pgfpathquadraticcurveto{\pgfqpoint{2.438554in}{0.585092in}}{\pgfqpoint{2.433475in}{0.585092in}}%
\pgfpathquadraticcurveto{\pgfqpoint{2.425766in}{0.585092in}}{\pgfqpoint{2.421911in}{0.582732in}}%
\pgfpathquadraticcurveto{\pgfqpoint{2.418056in}{0.580373in}}{\pgfqpoint{2.418056in}{0.575635in}}%
\pgfpathquadraticcurveto{\pgfqpoint{2.418056in}{0.572033in}}{\pgfqpoint{2.420812in}{0.569980in}}%
\pgfpathquadraticcurveto{\pgfqpoint{2.423568in}{0.567926in}}{\pgfqpoint{2.431908in}{0.566071in}}%
\pgfpathlineto{\pgfqpoint{2.435456in}{0.565278in}}%
\pgfpathquadraticcurveto{\pgfqpoint{2.446480in}{0.562919in}}{\pgfqpoint{2.451127in}{0.558614in}}%
\pgfpathquadraticcurveto{\pgfqpoint{2.455774in}{0.554309in}}{\pgfqpoint{2.455774in}{0.546600in}}%
\pgfpathquadraticcurveto{\pgfqpoint{2.455774in}{0.537810in}}{\pgfqpoint{2.448821in}{0.532676in}}%
\pgfpathquadraticcurveto{\pgfqpoint{2.441869in}{0.527561in}}{\pgfqpoint{2.429710in}{0.527561in}}%
\pgfpathquadraticcurveto{\pgfqpoint{2.424649in}{0.527561in}}{\pgfqpoint{2.419155in}{0.528552in}}%
\pgfpathquadraticcurveto{\pgfqpoint{2.413661in}{0.529542in}}{\pgfqpoint{2.407591in}{0.531506in}}%
\pgfpathlineto{\pgfqpoint{2.407591in}{0.542205in}}%
\pgfpathquadraticcurveto{\pgfqpoint{2.413337in}{0.539215in}}{\pgfqpoint{2.418903in}{0.537720in}}%
\pgfpathquadraticcurveto{\pgfqpoint{2.424469in}{0.536243in}}{\pgfqpoint{2.429944in}{0.536243in}}%
\pgfpathquadraticcurveto{\pgfqpoint{2.437257in}{0.536243in}}{\pgfqpoint{2.441184in}{0.538746in}}%
\pgfpathquadraticcurveto{\pgfqpoint{2.445129in}{0.541250in}}{\pgfqpoint{2.445129in}{0.545807in}}%
\pgfpathquadraticcurveto{\pgfqpoint{2.445129in}{0.550022in}}{\pgfqpoint{2.442283in}{0.552274in}}%
\pgfpathquadraticcurveto{\pgfqpoint{2.439455in}{0.554525in}}{\pgfqpoint{2.429818in}{0.556614in}}%
\pgfpathlineto{\pgfqpoint{2.426216in}{0.557461in}}%
\pgfpathquadraticcurveto{\pgfqpoint{2.416597in}{0.559478in}}{\pgfqpoint{2.412311in}{0.563675in}}%
\pgfpathquadraticcurveto{\pgfqpoint{2.408042in}{0.567872in}}{\pgfqpoint{2.408042in}{0.575185in}}%
\pgfpathquadraticcurveto{\pgfqpoint{2.408042in}{0.584083in}}{\pgfqpoint{2.414346in}{0.588910in}}%
\pgfpathquadraticcurveto{\pgfqpoint{2.420650in}{0.593756in}}{\pgfqpoint{2.432250in}{0.593756in}}%
\pgfpathquadraticcurveto{\pgfqpoint{2.437978in}{0.593756in}}{\pgfqpoint{2.443039in}{0.592909in}}%
\pgfpathquadraticcurveto{\pgfqpoint{2.448119in}{0.592080in}}{\pgfqpoint{2.452388in}{0.590387in}}%
\pgfpathlineto{\pgfqpoint{2.452388in}{0.590387in}}%
\pgfpathclose%
\pgfusepath{fill}%
\end{pgfscope}%
\begin{pgfscope}%
\pgfsetbuttcap%
\pgfsetroundjoin%
\definecolor{currentfill}{rgb}{0.309804,0.372549,0.419608}%
\pgfsetfillcolor{currentfill}%
\pgfsetlinewidth{0.000000pt}%
\definecolor{currentstroke}{rgb}{0.309804,0.372549,0.419608}%
\pgfsetstrokecolor{currentstroke}%
\pgfsetdash{}{0pt}%
\pgfpathmoveto{\pgfqpoint{2.573964in}{0.590387in}}%
\pgfpathlineto{\pgfqpoint{2.573964in}{0.580589in}}%
\pgfpathquadraticcurveto{\pgfqpoint{2.569587in}{0.582840in}}{\pgfqpoint{2.564850in}{0.583957in}}%
\pgfpathquadraticcurveto{\pgfqpoint{2.560131in}{0.585092in}}{\pgfqpoint{2.555051in}{0.585092in}}%
\pgfpathquadraticcurveto{\pgfqpoint{2.547342in}{0.585092in}}{\pgfqpoint{2.543487in}{0.582732in}}%
\pgfpathquadraticcurveto{\pgfqpoint{2.539633in}{0.580373in}}{\pgfqpoint{2.539633in}{0.575635in}}%
\pgfpathquadraticcurveto{\pgfqpoint{2.539633in}{0.572033in}}{\pgfqpoint{2.542389in}{0.569980in}}%
\pgfpathquadraticcurveto{\pgfqpoint{2.545144in}{0.567926in}}{\pgfqpoint{2.553484in}{0.566071in}}%
\pgfpathlineto{\pgfqpoint{2.557032in}{0.565278in}}%
\pgfpathquadraticcurveto{\pgfqpoint{2.568056in}{0.562919in}}{\pgfqpoint{2.572703in}{0.558614in}}%
\pgfpathquadraticcurveto{\pgfqpoint{2.577350in}{0.554309in}}{\pgfqpoint{2.577350in}{0.546600in}}%
\pgfpathquadraticcurveto{\pgfqpoint{2.577350in}{0.537810in}}{\pgfqpoint{2.570397in}{0.532676in}}%
\pgfpathquadraticcurveto{\pgfqpoint{2.563445in}{0.527561in}}{\pgfqpoint{2.551287in}{0.527561in}}%
\pgfpathquadraticcurveto{\pgfqpoint{2.546225in}{0.527561in}}{\pgfqpoint{2.540731in}{0.528552in}}%
\pgfpathquadraticcurveto{\pgfqpoint{2.535238in}{0.529542in}}{\pgfqpoint{2.529168in}{0.531506in}}%
\pgfpathlineto{\pgfqpoint{2.529168in}{0.542205in}}%
\pgfpathquadraticcurveto{\pgfqpoint{2.534914in}{0.539215in}}{\pgfqpoint{2.540479in}{0.537720in}}%
\pgfpathquadraticcurveto{\pgfqpoint{2.546045in}{0.536243in}}{\pgfqpoint{2.551521in}{0.536243in}}%
\pgfpathquadraticcurveto{\pgfqpoint{2.558834in}{0.536243in}}{\pgfqpoint{2.562760in}{0.538746in}}%
\pgfpathquadraticcurveto{\pgfqpoint{2.566705in}{0.541250in}}{\pgfqpoint{2.566705in}{0.545807in}}%
\pgfpathquadraticcurveto{\pgfqpoint{2.566705in}{0.550022in}}{\pgfqpoint{2.563859in}{0.552274in}}%
\pgfpathquadraticcurveto{\pgfqpoint{2.561031in}{0.554525in}}{\pgfqpoint{2.551395in}{0.556614in}}%
\pgfpathlineto{\pgfqpoint{2.547792in}{0.557461in}}%
\pgfpathquadraticcurveto{\pgfqpoint{2.538174in}{0.559478in}}{\pgfqpoint{2.533887in}{0.563675in}}%
\pgfpathquadraticcurveto{\pgfqpoint{2.529618in}{0.567872in}}{\pgfqpoint{2.529618in}{0.575185in}}%
\pgfpathquadraticcurveto{\pgfqpoint{2.529618in}{0.584083in}}{\pgfqpoint{2.535922in}{0.588910in}}%
\pgfpathquadraticcurveto{\pgfqpoint{2.542226in}{0.593756in}}{\pgfqpoint{2.553826in}{0.593756in}}%
\pgfpathquadraticcurveto{\pgfqpoint{2.559554in}{0.593756in}}{\pgfqpoint{2.564616in}{0.592909in}}%
\pgfpathquadraticcurveto{\pgfqpoint{2.569695in}{0.592080in}}{\pgfqpoint{2.573964in}{0.590387in}}%
\pgfpathlineto{\pgfqpoint{2.573964in}{0.590387in}}%
\pgfpathclose%
\pgfpathmoveto{\pgfqpoint{2.622489in}{0.560883in}}%
\pgfpathquadraticcurveto{\pgfqpoint{2.609934in}{0.560883in}}{\pgfqpoint{2.605089in}{0.558019in}}%
\pgfpathquadraticcurveto{\pgfqpoint{2.600244in}{0.555156in}}{\pgfqpoint{2.600244in}{0.548221in}}%
\pgfpathquadraticcurveto{\pgfqpoint{2.600244in}{0.542709in}}{\pgfqpoint{2.603882in}{0.539467in}}%
\pgfpathquadraticcurveto{\pgfqpoint{2.607520in}{0.536243in}}{\pgfqpoint{2.613753in}{0.536243in}}%
\pgfpathquadraticcurveto{\pgfqpoint{2.622381in}{0.536243in}}{\pgfqpoint{2.627586in}{0.542349in}}%
\pgfpathquadraticcurveto{\pgfqpoint{2.632792in}{0.548455in}}{\pgfqpoint{2.632792in}{0.558578in}}%
\pgfpathlineto{\pgfqpoint{2.632792in}{0.560883in}}%
\pgfpathlineto{\pgfqpoint{2.622489in}{0.560883in}}%
\pgfpathlineto{\pgfqpoint{2.622489in}{0.560883in}}%
\pgfpathclose%
\pgfpathmoveto{\pgfqpoint{2.643149in}{0.565170in}}%
\pgfpathlineto{\pgfqpoint{2.643149in}{0.529200in}}%
\pgfpathlineto{\pgfqpoint{2.632792in}{0.529200in}}%
\pgfpathlineto{\pgfqpoint{2.632792in}{0.538764in}}%
\pgfpathquadraticcurveto{\pgfqpoint{2.629243in}{0.533037in}}{\pgfqpoint{2.623948in}{0.530299in}}%
\pgfpathquadraticcurveto{\pgfqpoint{2.618652in}{0.527561in}}{\pgfqpoint{2.610997in}{0.527561in}}%
\pgfpathquadraticcurveto{\pgfqpoint{2.601324in}{0.527561in}}{\pgfqpoint{2.595596in}{0.533001in}}%
\pgfpathquadraticcurveto{\pgfqpoint{2.589887in}{0.538440in}}{\pgfqpoint{2.589887in}{0.547554in}}%
\pgfpathquadraticcurveto{\pgfqpoint{2.589887in}{0.558182in}}{\pgfqpoint{2.597001in}{0.563585in}}%
\pgfpathquadraticcurveto{\pgfqpoint{2.604134in}{0.568989in}}{\pgfqpoint{2.618256in}{0.568989in}}%
\pgfpathlineto{\pgfqpoint{2.632792in}{0.568989in}}%
\pgfpathlineto{\pgfqpoint{2.632792in}{0.570016in}}%
\pgfpathquadraticcurveto{\pgfqpoint{2.632792in}{0.577166in}}{\pgfqpoint{2.628090in}{0.581075in}}%
\pgfpathquadraticcurveto{\pgfqpoint{2.623389in}{0.584984in}}{\pgfqpoint{2.614887in}{0.584984in}}%
\pgfpathquadraticcurveto{\pgfqpoint{2.609484in}{0.584984in}}{\pgfqpoint{2.604350in}{0.583687in}}%
\pgfpathquadraticcurveto{\pgfqpoint{2.599235in}{0.582390in}}{\pgfqpoint{2.594516in}{0.579796in}}%
\pgfpathlineto{\pgfqpoint{2.594516in}{0.589379in}}%
\pgfpathquadraticcurveto{\pgfqpoint{2.600190in}{0.591576in}}{\pgfqpoint{2.605539in}{0.592657in}}%
\pgfpathquadraticcurveto{\pgfqpoint{2.610889in}{0.593756in}}{\pgfqpoint{2.615950in}{0.593756in}}%
\pgfpathquadraticcurveto{\pgfqpoint{2.629639in}{0.593756in}}{\pgfqpoint{2.636394in}{0.586659in}}%
\pgfpathquadraticcurveto{\pgfqpoint{2.643149in}{0.579580in}}{\pgfqpoint{2.643149in}{0.565170in}}%
\pgfpathlineto{\pgfqpoint{2.643149in}{0.565170in}}%
\pgfpathclose%
\pgfpathmoveto{\pgfqpoint{2.657054in}{0.592243in}}%
\pgfpathlineto{\pgfqpoint{2.668023in}{0.592243in}}%
\pgfpathlineto{\pgfqpoint{2.687729in}{0.539341in}}%
\pgfpathlineto{\pgfqpoint{2.707434in}{0.592243in}}%
\pgfpathlineto{\pgfqpoint{2.718403in}{0.592243in}}%
\pgfpathlineto{\pgfqpoint{2.694753in}{0.529200in}}%
\pgfpathlineto{\pgfqpoint{2.680686in}{0.529200in}}%
\pgfpathlineto{\pgfqpoint{2.657054in}{0.592243in}}%
\pgfpathlineto{\pgfqpoint{2.657054in}{0.592243in}}%
\pgfpathclose%
\pgfpathmoveto{\pgfqpoint{2.786633in}{0.563315in}}%
\pgfpathlineto{\pgfqpoint{2.786633in}{0.558254in}}%
\pgfpathlineto{\pgfqpoint{2.739009in}{0.558254in}}%
\pgfpathquadraticcurveto{\pgfqpoint{2.739694in}{0.547554in}}{\pgfqpoint{2.745458in}{0.541953in}}%
\pgfpathquadraticcurveto{\pgfqpoint{2.751221in}{0.536351in}}{\pgfqpoint{2.761524in}{0.536351in}}%
\pgfpathquadraticcurveto{\pgfqpoint{2.767486in}{0.536351in}}{\pgfqpoint{2.773088in}{0.537810in}}%
\pgfpathquadraticcurveto{\pgfqpoint{2.778690in}{0.539269in}}{\pgfqpoint{2.784220in}{0.542205in}}%
\pgfpathlineto{\pgfqpoint{2.784220in}{0.532406in}}%
\pgfpathquadraticcurveto{\pgfqpoint{2.778636in}{0.530047in}}{\pgfqpoint{2.772782in}{0.528804in}}%
\pgfpathquadraticcurveto{\pgfqpoint{2.766928in}{0.527561in}}{\pgfqpoint{2.760912in}{0.527561in}}%
\pgfpathquadraticcurveto{\pgfqpoint{2.745818in}{0.527561in}}{\pgfqpoint{2.737010in}{0.536333in}}%
\pgfpathquadraticcurveto{\pgfqpoint{2.728202in}{0.545123in}}{\pgfqpoint{2.728202in}{0.560109in}}%
\pgfpathquadraticcurveto{\pgfqpoint{2.728202in}{0.575581in}}{\pgfqpoint{2.736560in}{0.584659in}}%
\pgfpathquadraticcurveto{\pgfqpoint{2.744917in}{0.593756in}}{\pgfqpoint{2.759111in}{0.593756in}}%
\pgfpathquadraticcurveto{\pgfqpoint{2.771827in}{0.593756in}}{\pgfqpoint{2.779230in}{0.585560in}}%
\pgfpathquadraticcurveto{\pgfqpoint{2.786633in}{0.577383in}}{\pgfqpoint{2.786633in}{0.563315in}}%
\pgfpathlineto{\pgfqpoint{2.786633in}{0.563315in}}%
\pgfpathclose%
\pgfpathmoveto{\pgfqpoint{2.776276in}{0.566359in}}%
\pgfpathquadraticcurveto{\pgfqpoint{2.776168in}{0.574843in}}{\pgfqpoint{2.771521in}{0.579904in}}%
\pgfpathquadraticcurveto{\pgfqpoint{2.766874in}{0.584984in}}{\pgfqpoint{2.759219in}{0.584984in}}%
\pgfpathquadraticcurveto{\pgfqpoint{2.750555in}{0.584984in}}{\pgfqpoint{2.745349in}{0.580084in}}%
\pgfpathquadraticcurveto{\pgfqpoint{2.740144in}{0.575185in}}{\pgfqpoint{2.739351in}{0.566287in}}%
\pgfpathlineto{\pgfqpoint{2.776276in}{0.566359in}}%
\pgfpathlineto{\pgfqpoint{2.776276in}{0.566359in}}%
\pgfpathclose%
\pgfpathmoveto{\pgfqpoint{2.845119in}{0.582678in}}%
\pgfpathlineto{\pgfqpoint{2.845119in}{0.616793in}}%
\pgfpathlineto{\pgfqpoint{2.855476in}{0.616793in}}%
\pgfpathlineto{\pgfqpoint{2.855476in}{0.529200in}}%
\pgfpathlineto{\pgfqpoint{2.845119in}{0.529200in}}%
\pgfpathlineto{\pgfqpoint{2.845119in}{0.538656in}}%
\pgfpathquadraticcurveto{\pgfqpoint{2.841859in}{0.533037in}}{\pgfqpoint{2.836869in}{0.530299in}}%
\pgfpathquadraticcurveto{\pgfqpoint{2.831898in}{0.527561in}}{\pgfqpoint{2.824909in}{0.527561in}}%
\pgfpathquadraticcurveto{\pgfqpoint{2.813489in}{0.527561in}}{\pgfqpoint{2.806303in}{0.536675in}}%
\pgfpathquadraticcurveto{\pgfqpoint{2.799134in}{0.545807in}}{\pgfqpoint{2.799134in}{0.560667in}}%
\pgfpathquadraticcurveto{\pgfqpoint{2.799134in}{0.575527in}}{\pgfqpoint{2.806303in}{0.584641in}}%
\pgfpathquadraticcurveto{\pgfqpoint{2.813489in}{0.593756in}}{\pgfqpoint{2.824909in}{0.593756in}}%
\pgfpathquadraticcurveto{\pgfqpoint{2.831898in}{0.593756in}}{\pgfqpoint{2.836869in}{0.591018in}}%
\pgfpathquadraticcurveto{\pgfqpoint{2.841859in}{0.588298in}}{\pgfqpoint{2.845119in}{0.582678in}}%
\pgfpathlineto{\pgfqpoint{2.845119in}{0.582678in}}%
\pgfpathclose%
\pgfpathmoveto{\pgfqpoint{2.809833in}{0.560667in}}%
\pgfpathquadraticcurveto{\pgfqpoint{2.809833in}{0.549248in}}{\pgfqpoint{2.814534in}{0.542745in}}%
\pgfpathquadraticcurveto{\pgfqpoint{2.819235in}{0.536243in}}{\pgfqpoint{2.827449in}{0.536243in}}%
\pgfpathquadraticcurveto{\pgfqpoint{2.835662in}{0.536243in}}{\pgfqpoint{2.840382in}{0.542745in}}%
\pgfpathquadraticcurveto{\pgfqpoint{2.845119in}{0.549248in}}{\pgfqpoint{2.845119in}{0.560667in}}%
\pgfpathquadraticcurveto{\pgfqpoint{2.845119in}{0.572087in}}{\pgfqpoint{2.840382in}{0.578589in}}%
\pgfpathquadraticcurveto{\pgfqpoint{2.835662in}{0.585092in}}{\pgfqpoint{2.827449in}{0.585092in}}%
\pgfpathquadraticcurveto{\pgfqpoint{2.819235in}{0.585092in}}{\pgfqpoint{2.814534in}{0.578589in}}%
\pgfpathquadraticcurveto{\pgfqpoint{2.809833in}{0.572087in}}{\pgfqpoint{2.809833in}{0.560667in}}%
\pgfpathlineto{\pgfqpoint{2.809833in}{0.560667in}}%
\pgfpathclose%
\pgfusepath{fill}%
\end{pgfscope}%
\begin{pgfscope}%
\pgfsetbuttcap%
\pgfsetroundjoin%
\definecolor{currentfill}{rgb}{0.309804,0.372549,0.419608}%
\pgfsetfillcolor{currentfill}%
\pgfsetlinewidth{0.000000pt}%
\definecolor{currentstroke}{rgb}{0.309804,0.372549,0.419608}%
\pgfsetstrokecolor{currentstroke}%
\pgfsetdash{}{0pt}%
\pgfpathmoveto{\pgfqpoint{2.929077in}{0.613245in}}%
\pgfpathlineto{\pgfqpoint{3.000171in}{0.613245in}}%
\pgfpathlineto{\pgfqpoint{3.000171in}{0.603662in}}%
\pgfpathlineto{\pgfqpoint{2.970343in}{0.603662in}}%
\pgfpathlineto{\pgfqpoint{2.970343in}{0.529200in}}%
\pgfpathlineto{\pgfqpoint{2.958924in}{0.529200in}}%
\pgfpathlineto{\pgfqpoint{2.958924in}{0.603662in}}%
\pgfpathlineto{\pgfqpoint{2.929077in}{0.603662in}}%
\pgfpathlineto{\pgfqpoint{2.929077in}{0.613245in}}%
\pgfpathlineto{\pgfqpoint{2.929077in}{0.613245in}}%
\pgfpathclose%
\pgfpathmoveto{\pgfqpoint{3.003666in}{0.613245in}}%
\pgfpathlineto{\pgfqpoint{3.015139in}{0.613245in}}%
\pgfpathlineto{\pgfqpoint{3.032809in}{0.542205in}}%
\pgfpathlineto{\pgfqpoint{3.050425in}{0.613245in}}%
\pgfpathlineto{\pgfqpoint{3.063214in}{0.613245in}}%
\pgfpathlineto{\pgfqpoint{3.080884in}{0.542205in}}%
\pgfpathlineto{\pgfqpoint{3.098500in}{0.613245in}}%
\pgfpathlineto{\pgfqpoint{3.110046in}{0.613245in}}%
\pgfpathlineto{\pgfqpoint{3.088935in}{0.529200in}}%
\pgfpathlineto{\pgfqpoint{3.074634in}{0.529200in}}%
\pgfpathlineto{\pgfqpoint{3.056910in}{0.602149in}}%
\pgfpathlineto{\pgfqpoint{3.039006in}{0.529200in}}%
\pgfpathlineto{\pgfqpoint{3.024704in}{0.529200in}}%
\pgfpathlineto{\pgfqpoint{3.003666in}{0.613245in}}%
\pgfpathlineto{\pgfqpoint{3.003666in}{0.613245in}}%
\pgfpathclose%
\pgfpathmoveto{\pgfqpoint{3.125122in}{0.613245in}}%
\pgfpathlineto{\pgfqpoint{3.173412in}{0.613245in}}%
\pgfpathlineto{\pgfqpoint{3.173412in}{0.603662in}}%
\pgfpathlineto{\pgfqpoint{3.136487in}{0.603662in}}%
\pgfpathlineto{\pgfqpoint{3.136487in}{0.578896in}}%
\pgfpathlineto{\pgfqpoint{3.169810in}{0.578896in}}%
\pgfpathlineto{\pgfqpoint{3.169810in}{0.569331in}}%
\pgfpathlineto{\pgfqpoint{3.136487in}{0.569331in}}%
\pgfpathlineto{\pgfqpoint{3.136487in}{0.529200in}}%
\pgfpathlineto{\pgfqpoint{3.125122in}{0.529200in}}%
\pgfpathlineto{\pgfqpoint{3.125122in}{0.613245in}}%
\pgfpathlineto{\pgfqpoint{3.125122in}{0.613245in}}%
\pgfpathclose%
\pgfpathmoveto{\pgfqpoint{3.191424in}{0.613245in}}%
\pgfpathlineto{\pgfqpoint{3.244560in}{0.613245in}}%
\pgfpathlineto{\pgfqpoint{3.244560in}{0.603662in}}%
\pgfpathlineto{\pgfqpoint{3.202790in}{0.603662in}}%
\pgfpathlineto{\pgfqpoint{3.202790in}{0.578787in}}%
\pgfpathlineto{\pgfqpoint{3.242813in}{0.578787in}}%
\pgfpathlineto{\pgfqpoint{3.242813in}{0.569223in}}%
\pgfpathlineto{\pgfqpoint{3.202790in}{0.569223in}}%
\pgfpathlineto{\pgfqpoint{3.202790in}{0.538764in}}%
\pgfpathlineto{\pgfqpoint{3.245569in}{0.538764in}}%
\pgfpathlineto{\pgfqpoint{3.245569in}{0.529200in}}%
\pgfpathlineto{\pgfqpoint{3.191424in}{0.529200in}}%
\pgfpathlineto{\pgfqpoint{3.191424in}{0.613245in}}%
\pgfpathlineto{\pgfqpoint{3.191424in}{0.613245in}}%
\pgfpathclose%
\pgfusepath{fill}%
\end{pgfscope}%
\begin{pgfscope}%
\pgfsetbuttcap%
\pgfsetroundjoin%
\definecolor{currentfill}{rgb}{0.309804,0.372549,0.419608}%
\pgfsetfillcolor{currentfill}%
\pgfsetlinewidth{0.000000pt}%
\definecolor{currentstroke}{rgb}{0.309804,0.372549,0.419608}%
\pgfsetstrokecolor{currentstroke}%
\pgfsetdash{}{0pt}%
\pgfpathmoveto{\pgfqpoint{3.387613in}{0.610489in}}%
\pgfpathlineto{\pgfqpoint{3.387613in}{0.599393in}}%
\pgfpathquadraticcurveto{\pgfqpoint{3.381147in}{0.602491in}}{\pgfqpoint{3.375401in}{0.604004in}}%
\pgfpathquadraticcurveto{\pgfqpoint{3.369655in}{0.605536in}}{\pgfqpoint{3.364306in}{0.605536in}}%
\pgfpathquadraticcurveto{\pgfqpoint{3.355030in}{0.605536in}}{\pgfqpoint{3.349986in}{0.601933in}}%
\pgfpathquadraticcurveto{\pgfqpoint{3.344943in}{0.598331in}}{\pgfqpoint{3.344943in}{0.591684in}}%
\pgfpathquadraticcurveto{\pgfqpoint{3.344943in}{0.586100in}}{\pgfqpoint{3.348293in}{0.583254in}}%
\pgfpathquadraticcurveto{\pgfqpoint{3.351643in}{0.580427in}}{\pgfqpoint{3.360992in}{0.578679in}}%
\pgfpathlineto{\pgfqpoint{3.367854in}{0.577274in}}%
\pgfpathquadraticcurveto{\pgfqpoint{3.380571in}{0.574843in}}{\pgfqpoint{3.386623in}{0.568737in}}%
\pgfpathquadraticcurveto{\pgfqpoint{3.392675in}{0.562631in}}{\pgfqpoint{3.392675in}{0.552400in}}%
\pgfpathquadraticcurveto{\pgfqpoint{3.392675in}{0.540169in}}{\pgfqpoint{3.384479in}{0.533865in}}%
\pgfpathquadraticcurveto{\pgfqpoint{3.376302in}{0.527561in}}{\pgfqpoint{3.360487in}{0.527561in}}%
\pgfpathquadraticcurveto{\pgfqpoint{3.354525in}{0.527561in}}{\pgfqpoint{3.347789in}{0.528912in}}%
\pgfpathquadraticcurveto{\pgfqpoint{3.341070in}{0.530263in}}{\pgfqpoint{3.333865in}{0.532911in}}%
\pgfpathlineto{\pgfqpoint{3.333865in}{0.544618in}}%
\pgfpathquadraticcurveto{\pgfqpoint{3.340782in}{0.540746in}}{\pgfqpoint{3.347428in}{0.538764in}}%
\pgfpathquadraticcurveto{\pgfqpoint{3.354075in}{0.536801in}}{\pgfqpoint{3.360487in}{0.536801in}}%
\pgfpathquadraticcurveto{\pgfqpoint{3.370214in}{0.536801in}}{\pgfqpoint{3.375509in}{0.540620in}}%
\pgfpathquadraticcurveto{\pgfqpoint{3.380805in}{0.544456in}}{\pgfqpoint{3.380805in}{0.551553in}}%
\pgfpathquadraticcurveto{\pgfqpoint{3.380805in}{0.557731in}}{\pgfqpoint{3.377004in}{0.561226in}}%
\pgfpathquadraticcurveto{\pgfqpoint{3.373204in}{0.564720in}}{\pgfqpoint{3.364540in}{0.566467in}}%
\pgfpathlineto{\pgfqpoint{3.357605in}{0.567818in}}%
\pgfpathquadraticcurveto{\pgfqpoint{3.344889in}{0.570340in}}{\pgfqpoint{3.339197in}{0.575743in}}%
\pgfpathquadraticcurveto{\pgfqpoint{3.333523in}{0.581147in}}{\pgfqpoint{3.333523in}{0.590784in}}%
\pgfpathquadraticcurveto{\pgfqpoint{3.333523in}{0.601933in}}{\pgfqpoint{3.341376in}{0.608345in}}%
\pgfpathquadraticcurveto{\pgfqpoint{3.349230in}{0.614758in}}{\pgfqpoint{3.363009in}{0.614758in}}%
\pgfpathquadraticcurveto{\pgfqpoint{3.368935in}{0.614758in}}{\pgfqpoint{3.375059in}{0.613677in}}%
\pgfpathquadraticcurveto{\pgfqpoint{3.381201in}{0.612614in}}{\pgfqpoint{3.387613in}{0.610489in}}%
\pgfpathlineto{\pgfqpoint{3.387613in}{0.610489in}}%
\pgfpathclose%
\pgfpathmoveto{\pgfqpoint{3.410417in}{0.613245in}}%
\pgfpathlineto{\pgfqpoint{3.463553in}{0.613245in}}%
\pgfpathlineto{\pgfqpoint{3.463553in}{0.603662in}}%
\pgfpathlineto{\pgfqpoint{3.421783in}{0.603662in}}%
\pgfpathlineto{\pgfqpoint{3.421783in}{0.578787in}}%
\pgfpathlineto{\pgfqpoint{3.461806in}{0.578787in}}%
\pgfpathlineto{\pgfqpoint{3.461806in}{0.569223in}}%
\pgfpathlineto{\pgfqpoint{3.421783in}{0.569223in}}%
\pgfpathlineto{\pgfqpoint{3.421783in}{0.538764in}}%
\pgfpathlineto{\pgfqpoint{3.464561in}{0.538764in}}%
\pgfpathlineto{\pgfqpoint{3.464561in}{0.529200in}}%
\pgfpathlineto{\pgfqpoint{3.410417in}{0.529200in}}%
\pgfpathlineto{\pgfqpoint{3.410417in}{0.613245in}}%
\pgfpathlineto{\pgfqpoint{3.410417in}{0.613245in}}%
\pgfpathclose%
\pgfpathmoveto{\pgfqpoint{3.485456in}{0.588802in}}%
\pgfpathlineto{\pgfqpoint{3.497326in}{0.588802in}}%
\pgfpathlineto{\pgfqpoint{3.497326in}{0.574519in}}%
\pgfpathlineto{\pgfqpoint{3.485456in}{0.574519in}}%
\pgfpathlineto{\pgfqpoint{3.485456in}{0.588802in}}%
\pgfpathlineto{\pgfqpoint{3.485456in}{0.588802in}}%
\pgfpathclose%
\pgfpathmoveto{\pgfqpoint{3.485456in}{0.543502in}}%
\pgfpathlineto{\pgfqpoint{3.497326in}{0.543502in}}%
\pgfpathlineto{\pgfqpoint{3.497326in}{0.533811in}}%
\pgfpathlineto{\pgfqpoint{3.488103in}{0.515799in}}%
\pgfpathlineto{\pgfqpoint{3.480844in}{0.515799in}}%
\pgfpathlineto{\pgfqpoint{3.485456in}{0.533811in}}%
\pgfpathlineto{\pgfqpoint{3.485456in}{0.543502in}}%
\pgfpathlineto{\pgfqpoint{3.485456in}{0.543502in}}%
\pgfpathclose%
\pgfusepath{fill}%
\end{pgfscope}%
\begin{pgfscope}%
\pgfsetbuttcap%
\pgfsetroundjoin%
\definecolor{currentfill}{rgb}{0.309804,0.372549,0.419608}%
\pgfsetfillcolor{currentfill}%
\pgfsetlinewidth{0.000000pt}%
\definecolor{currentstroke}{rgb}{0.309804,0.372549,0.419608}%
\pgfsetstrokecolor{currentstroke}%
\pgfsetdash{}{0pt}%
\pgfpathmoveto{\pgfqpoint{3.603534in}{0.610147in}}%
\pgfpathlineto{\pgfqpoint{3.603534in}{0.592243in}}%
\pgfpathlineto{\pgfqpoint{3.624861in}{0.592243in}}%
\pgfpathlineto{\pgfqpoint{3.624861in}{0.584191in}}%
\pgfpathlineto{\pgfqpoint{3.603534in}{0.584191in}}%
\pgfpathlineto{\pgfqpoint{3.603534in}{0.549968in}}%
\pgfpathquadraticcurveto{\pgfqpoint{3.603534in}{0.542259in}}{\pgfqpoint{3.605642in}{0.540061in}}%
\pgfpathquadraticcurveto{\pgfqpoint{3.607749in}{0.537864in}}{\pgfqpoint{3.614234in}{0.537864in}}%
\pgfpathlineto{\pgfqpoint{3.624861in}{0.537864in}}%
\pgfpathlineto{\pgfqpoint{3.624861in}{0.529200in}}%
\pgfpathlineto{\pgfqpoint{3.614234in}{0.529200in}}%
\pgfpathquadraticcurveto{\pgfqpoint{3.602237in}{0.529200in}}{\pgfqpoint{3.597680in}{0.533667in}}%
\pgfpathquadraticcurveto{\pgfqpoint{3.593123in}{0.538152in}}{\pgfqpoint{3.593123in}{0.549968in}}%
\pgfpathlineto{\pgfqpoint{3.593123in}{0.584191in}}%
\pgfpathlineto{\pgfqpoint{3.585522in}{0.584191in}}%
\pgfpathlineto{\pgfqpoint{3.585522in}{0.592243in}}%
\pgfpathlineto{\pgfqpoint{3.593123in}{0.592243in}}%
\pgfpathlineto{\pgfqpoint{3.593123in}{0.610147in}}%
\pgfpathlineto{\pgfqpoint{3.603534in}{0.610147in}}%
\pgfpathlineto{\pgfqpoint{3.603534in}{0.610147in}}%
\pgfpathclose%
\pgfpathmoveto{\pgfqpoint{3.690893in}{0.567260in}}%
\pgfpathlineto{\pgfqpoint{3.690893in}{0.529200in}}%
\pgfpathlineto{\pgfqpoint{3.680536in}{0.529200in}}%
\pgfpathlineto{\pgfqpoint{3.680536in}{0.566917in}}%
\pgfpathquadraticcurveto{\pgfqpoint{3.680536in}{0.575869in}}{\pgfqpoint{3.677042in}{0.580300in}}%
\pgfpathquadraticcurveto{\pgfqpoint{3.673548in}{0.584749in}}{\pgfqpoint{3.666577in}{0.584749in}}%
\pgfpathquadraticcurveto{\pgfqpoint{3.658183in}{0.584749in}}{\pgfqpoint{3.653338in}{0.579400in}}%
\pgfpathquadraticcurveto{\pgfqpoint{3.648493in}{0.574068in}}{\pgfqpoint{3.648493in}{0.564828in}}%
\pgfpathlineto{\pgfqpoint{3.648493in}{0.529200in}}%
\pgfpathlineto{\pgfqpoint{3.638082in}{0.529200in}}%
\pgfpathlineto{\pgfqpoint{3.638082in}{0.616793in}}%
\pgfpathlineto{\pgfqpoint{3.648493in}{0.616793in}}%
\pgfpathlineto{\pgfqpoint{3.648493in}{0.582444in}}%
\pgfpathquadraticcurveto{\pgfqpoint{3.652221in}{0.588136in}}{\pgfqpoint{3.657247in}{0.590946in}}%
\pgfpathquadraticcurveto{\pgfqpoint{3.662290in}{0.593756in}}{\pgfqpoint{3.668882in}{0.593756in}}%
\pgfpathquadraticcurveto{\pgfqpoint{3.679744in}{0.593756in}}{\pgfqpoint{3.685310in}{0.587037in}}%
\pgfpathquadraticcurveto{\pgfqpoint{3.690893in}{0.580318in}}{\pgfqpoint{3.690893in}{0.567260in}}%
\pgfpathlineto{\pgfqpoint{3.690893in}{0.567260in}}%
\pgfpathclose%
\pgfpathmoveto{\pgfqpoint{3.765464in}{0.563315in}}%
\pgfpathlineto{\pgfqpoint{3.765464in}{0.558254in}}%
\pgfpathlineto{\pgfqpoint{3.717839in}{0.558254in}}%
\pgfpathquadraticcurveto{\pgfqpoint{3.718524in}{0.547554in}}{\pgfqpoint{3.724288in}{0.541953in}}%
\pgfpathquadraticcurveto{\pgfqpoint{3.730052in}{0.536351in}}{\pgfqpoint{3.740355in}{0.536351in}}%
\pgfpathquadraticcurveto{\pgfqpoint{3.746317in}{0.536351in}}{\pgfqpoint{3.751918in}{0.537810in}}%
\pgfpathquadraticcurveto{\pgfqpoint{3.757520in}{0.539269in}}{\pgfqpoint{3.763050in}{0.542205in}}%
\pgfpathlineto{\pgfqpoint{3.763050in}{0.532406in}}%
\pgfpathquadraticcurveto{\pgfqpoint{3.757466in}{0.530047in}}{\pgfqpoint{3.751612in}{0.528804in}}%
\pgfpathquadraticcurveto{\pgfqpoint{3.745758in}{0.527561in}}{\pgfqpoint{3.739742in}{0.527561in}}%
\pgfpathquadraticcurveto{\pgfqpoint{3.724648in}{0.527561in}}{\pgfqpoint{3.715840in}{0.536333in}}%
\pgfpathquadraticcurveto{\pgfqpoint{3.707032in}{0.545123in}}{\pgfqpoint{3.707032in}{0.560109in}}%
\pgfpathquadraticcurveto{\pgfqpoint{3.707032in}{0.575581in}}{\pgfqpoint{3.715390in}{0.584659in}}%
\pgfpathquadraticcurveto{\pgfqpoint{3.723747in}{0.593756in}}{\pgfqpoint{3.737941in}{0.593756in}}%
\pgfpathquadraticcurveto{\pgfqpoint{3.750658in}{0.593756in}}{\pgfqpoint{3.758061in}{0.585560in}}%
\pgfpathquadraticcurveto{\pgfqpoint{3.765464in}{0.577383in}}{\pgfqpoint{3.765464in}{0.563315in}}%
\pgfpathlineto{\pgfqpoint{3.765464in}{0.563315in}}%
\pgfpathclose%
\pgfpathmoveto{\pgfqpoint{3.755107in}{0.566359in}}%
\pgfpathquadraticcurveto{\pgfqpoint{3.754999in}{0.574843in}}{\pgfqpoint{3.750351in}{0.579904in}}%
\pgfpathquadraticcurveto{\pgfqpoint{3.745704in}{0.584984in}}{\pgfqpoint{3.738049in}{0.584984in}}%
\pgfpathquadraticcurveto{\pgfqpoint{3.729385in}{0.584984in}}{\pgfqpoint{3.724180in}{0.580084in}}%
\pgfpathquadraticcurveto{\pgfqpoint{3.718974in}{0.575185in}}{\pgfqpoint{3.718182in}{0.566287in}}%
\pgfpathlineto{\pgfqpoint{3.755107in}{0.566359in}}%
\pgfpathlineto{\pgfqpoint{3.755107in}{0.566359in}}%
\pgfpathclose%
\pgfusepath{fill}%
\end{pgfscope}%
\begin{pgfscope}%
\pgfsetbuttcap%
\pgfsetroundjoin%
\definecolor{currentfill}{rgb}{0.309804,0.372549,0.419608}%
\pgfsetfillcolor{currentfill}%
\pgfsetlinewidth{0.000000pt}%
\definecolor{currentstroke}{rgb}{0.309804,0.372549,0.419608}%
\pgfsetstrokecolor{currentstroke}%
\pgfsetdash{}{0pt}%
\pgfpathmoveto{\pgfqpoint{3.859802in}{0.538656in}}%
\pgfpathlineto{\pgfqpoint{3.859802in}{0.505226in}}%
\pgfpathlineto{\pgfqpoint{3.849391in}{0.505226in}}%
\pgfpathlineto{\pgfqpoint{3.849391in}{0.592243in}}%
\pgfpathlineto{\pgfqpoint{3.859802in}{0.592243in}}%
\pgfpathlineto{\pgfqpoint{3.859802in}{0.582678in}}%
\pgfpathquadraticcurveto{\pgfqpoint{3.863081in}{0.588298in}}{\pgfqpoint{3.868052in}{0.591018in}}%
\pgfpathquadraticcurveto{\pgfqpoint{3.873041in}{0.593756in}}{\pgfqpoint{3.879958in}{0.593756in}}%
\pgfpathquadraticcurveto{\pgfqpoint{3.891450in}{0.593756in}}{\pgfqpoint{3.898619in}{0.584641in}}%
\pgfpathquadraticcurveto{\pgfqpoint{3.905805in}{0.575527in}}{\pgfqpoint{3.905805in}{0.560667in}}%
\pgfpathquadraticcurveto{\pgfqpoint{3.905805in}{0.545807in}}{\pgfqpoint{3.898619in}{0.536675in}}%
\pgfpathquadraticcurveto{\pgfqpoint{3.891450in}{0.527561in}}{\pgfqpoint{3.879958in}{0.527561in}}%
\pgfpathquadraticcurveto{\pgfqpoint{3.873041in}{0.527561in}}{\pgfqpoint{3.868052in}{0.530299in}}%
\pgfpathquadraticcurveto{\pgfqpoint{3.863081in}{0.533037in}}{\pgfqpoint{3.859802in}{0.538656in}}%
\pgfpathlineto{\pgfqpoint{3.859802in}{0.538656in}}%
\pgfpathclose%
\pgfpathmoveto{\pgfqpoint{3.895052in}{0.560667in}}%
\pgfpathquadraticcurveto{\pgfqpoint{3.895052in}{0.572087in}}{\pgfqpoint{3.890351in}{0.578589in}}%
\pgfpathquadraticcurveto{\pgfqpoint{3.885650in}{0.585092in}}{\pgfqpoint{3.877436in}{0.585092in}}%
\pgfpathquadraticcurveto{\pgfqpoint{3.869205in}{0.585092in}}{\pgfqpoint{3.864504in}{0.578589in}}%
\pgfpathquadraticcurveto{\pgfqpoint{3.859802in}{0.572087in}}{\pgfqpoint{3.859802in}{0.560667in}}%
\pgfpathquadraticcurveto{\pgfqpoint{3.859802in}{0.549248in}}{\pgfqpoint{3.864504in}{0.542745in}}%
\pgfpathquadraticcurveto{\pgfqpoint{3.869205in}{0.536243in}}{\pgfqpoint{3.877436in}{0.536243in}}%
\pgfpathquadraticcurveto{\pgfqpoint{3.885650in}{0.536243in}}{\pgfqpoint{3.890351in}{0.542745in}}%
\pgfpathquadraticcurveto{\pgfqpoint{3.895052in}{0.549248in}}{\pgfqpoint{3.895052in}{0.560667in}}%
\pgfpathlineto{\pgfqpoint{3.895052in}{0.560667in}}%
\pgfpathclose%
\pgfpathmoveto{\pgfqpoint{3.922971in}{0.616793in}}%
\pgfpathlineto{\pgfqpoint{3.933328in}{0.616793in}}%
\pgfpathlineto{\pgfqpoint{3.933328in}{0.529200in}}%
\pgfpathlineto{\pgfqpoint{3.922971in}{0.529200in}}%
\pgfpathlineto{\pgfqpoint{3.922971in}{0.616793in}}%
\pgfpathlineto{\pgfqpoint{3.922971in}{0.616793in}}%
\pgfpathclose%
\pgfpathmoveto{\pgfqpoint{3.979421in}{0.584984in}}%
\pgfpathquadraticcurveto{\pgfqpoint{3.971100in}{0.584984in}}{\pgfqpoint{3.966254in}{0.578481in}}%
\pgfpathquadraticcurveto{\pgfqpoint{3.961409in}{0.571979in}}{\pgfqpoint{3.961409in}{0.560667in}}%
\pgfpathquadraticcurveto{\pgfqpoint{3.961409in}{0.549356in}}{\pgfqpoint{3.966218in}{0.542853in}}%
\pgfpathquadraticcurveto{\pgfqpoint{3.971045in}{0.536351in}}{\pgfqpoint{3.979421in}{0.536351in}}%
\pgfpathquadraticcurveto{\pgfqpoint{3.987707in}{0.536351in}}{\pgfqpoint{3.992534in}{0.542871in}}%
\pgfpathquadraticcurveto{\pgfqpoint{3.997379in}{0.549410in}}{\pgfqpoint{3.997379in}{0.560667in}}%
\pgfpathquadraticcurveto{\pgfqpoint{3.997379in}{0.571871in}}{\pgfqpoint{3.992534in}{0.578427in}}%
\pgfpathquadraticcurveto{\pgfqpoint{3.987707in}{0.584984in}}{\pgfqpoint{3.979421in}{0.584984in}}%
\pgfpathlineto{\pgfqpoint{3.979421in}{0.584984in}}%
\pgfpathclose%
\pgfpathmoveto{\pgfqpoint{3.979421in}{0.593756in}}%
\pgfpathquadraticcurveto{\pgfqpoint{3.992930in}{0.593756in}}{\pgfqpoint{4.000639in}{0.584966in}}%
\pgfpathquadraticcurveto{\pgfqpoint{4.008367in}{0.576194in}}{\pgfqpoint{4.008367in}{0.560667in}}%
\pgfpathquadraticcurveto{\pgfqpoint{4.008367in}{0.545195in}}{\pgfqpoint{4.000639in}{0.536369in}}%
\pgfpathquadraticcurveto{\pgfqpoint{3.992930in}{0.527561in}}{\pgfqpoint{3.979421in}{0.527561in}}%
\pgfpathquadraticcurveto{\pgfqpoint{3.965858in}{0.527561in}}{\pgfqpoint{3.958167in}{0.536369in}}%
\pgfpathquadraticcurveto{\pgfqpoint{3.950494in}{0.545195in}}{\pgfqpoint{3.950494in}{0.560667in}}%
\pgfpathquadraticcurveto{\pgfqpoint{3.950494in}{0.576194in}}{\pgfqpoint{3.958167in}{0.584966in}}%
\pgfpathquadraticcurveto{\pgfqpoint{3.965858in}{0.593756in}}{\pgfqpoint{3.979421in}{0.593756in}}%
\pgfpathlineto{\pgfqpoint{3.979421in}{0.593756in}}%
\pgfpathclose%
\pgfpathmoveto{\pgfqpoint{4.035781in}{0.610147in}}%
\pgfpathlineto{\pgfqpoint{4.035781in}{0.592243in}}%
\pgfpathlineto{\pgfqpoint{4.057108in}{0.592243in}}%
\pgfpathlineto{\pgfqpoint{4.057108in}{0.584191in}}%
\pgfpathlineto{\pgfqpoint{4.035781in}{0.584191in}}%
\pgfpathlineto{\pgfqpoint{4.035781in}{0.549968in}}%
\pgfpathquadraticcurveto{\pgfqpoint{4.035781in}{0.542259in}}{\pgfqpoint{4.037889in}{0.540061in}}%
\pgfpathquadraticcurveto{\pgfqpoint{4.039996in}{0.537864in}}{\pgfqpoint{4.046480in}{0.537864in}}%
\pgfpathlineto{\pgfqpoint{4.057108in}{0.537864in}}%
\pgfpathlineto{\pgfqpoint{4.057108in}{0.529200in}}%
\pgfpathlineto{\pgfqpoint{4.046480in}{0.529200in}}%
\pgfpathquadraticcurveto{\pgfqpoint{4.034484in}{0.529200in}}{\pgfqpoint{4.029927in}{0.533667in}}%
\pgfpathquadraticcurveto{\pgfqpoint{4.025370in}{0.538152in}}{\pgfqpoint{4.025370in}{0.549968in}}%
\pgfpathlineto{\pgfqpoint{4.025370in}{0.584191in}}%
\pgfpathlineto{\pgfqpoint{4.017769in}{0.584191in}}%
\pgfpathlineto{\pgfqpoint{4.017769in}{0.592243in}}%
\pgfpathlineto{\pgfqpoint{4.025370in}{0.592243in}}%
\pgfpathlineto{\pgfqpoint{4.025370in}{0.610147in}}%
\pgfpathlineto{\pgfqpoint{4.035781in}{0.610147in}}%
\pgfpathlineto{\pgfqpoint{4.035781in}{0.610147in}}%
\pgfpathclose%
\pgfpathmoveto{\pgfqpoint{4.080974in}{0.610147in}}%
\pgfpathlineto{\pgfqpoint{4.080974in}{0.592243in}}%
\pgfpathlineto{\pgfqpoint{4.102300in}{0.592243in}}%
\pgfpathlineto{\pgfqpoint{4.102300in}{0.584191in}}%
\pgfpathlineto{\pgfqpoint{4.080974in}{0.584191in}}%
\pgfpathlineto{\pgfqpoint{4.080974in}{0.549968in}}%
\pgfpathquadraticcurveto{\pgfqpoint{4.080974in}{0.542259in}}{\pgfqpoint{4.083081in}{0.540061in}}%
\pgfpathquadraticcurveto{\pgfqpoint{4.085188in}{0.537864in}}{\pgfqpoint{4.091673in}{0.537864in}}%
\pgfpathlineto{\pgfqpoint{4.102300in}{0.537864in}}%
\pgfpathlineto{\pgfqpoint{4.102300in}{0.529200in}}%
\pgfpathlineto{\pgfqpoint{4.091673in}{0.529200in}}%
\pgfpathquadraticcurveto{\pgfqpoint{4.079677in}{0.529200in}}{\pgfqpoint{4.075120in}{0.533667in}}%
\pgfpathquadraticcurveto{\pgfqpoint{4.070563in}{0.538152in}}{\pgfqpoint{4.070563in}{0.549968in}}%
\pgfpathlineto{\pgfqpoint{4.070563in}{0.584191in}}%
\pgfpathlineto{\pgfqpoint{4.062961in}{0.584191in}}%
\pgfpathlineto{\pgfqpoint{4.062961in}{0.592243in}}%
\pgfpathlineto{\pgfqpoint{4.070563in}{0.592243in}}%
\pgfpathlineto{\pgfqpoint{4.070563in}{0.610147in}}%
\pgfpathlineto{\pgfqpoint{4.080974in}{0.610147in}}%
\pgfpathlineto{\pgfqpoint{4.080974in}{0.610147in}}%
\pgfpathclose%
\pgfpathmoveto{\pgfqpoint{4.169846in}{0.563315in}}%
\pgfpathlineto{\pgfqpoint{4.169846in}{0.558254in}}%
\pgfpathlineto{\pgfqpoint{4.122221in}{0.558254in}}%
\pgfpathquadraticcurveto{\pgfqpoint{4.122906in}{0.547554in}}{\pgfqpoint{4.128670in}{0.541953in}}%
\pgfpathquadraticcurveto{\pgfqpoint{4.134434in}{0.536351in}}{\pgfqpoint{4.144737in}{0.536351in}}%
\pgfpathquadraticcurveto{\pgfqpoint{4.150699in}{0.536351in}}{\pgfqpoint{4.156300in}{0.537810in}}%
\pgfpathquadraticcurveto{\pgfqpoint{4.161902in}{0.539269in}}{\pgfqpoint{4.167432in}{0.542205in}}%
\pgfpathlineto{\pgfqpoint{4.167432in}{0.532406in}}%
\pgfpathquadraticcurveto{\pgfqpoint{4.161848in}{0.530047in}}{\pgfqpoint{4.155994in}{0.528804in}}%
\pgfpathquadraticcurveto{\pgfqpoint{4.150140in}{0.527561in}}{\pgfqpoint{4.144124in}{0.527561in}}%
\pgfpathquadraticcurveto{\pgfqpoint{4.129030in}{0.527561in}}{\pgfqpoint{4.120222in}{0.536333in}}%
\pgfpathquadraticcurveto{\pgfqpoint{4.111414in}{0.545123in}}{\pgfqpoint{4.111414in}{0.560109in}}%
\pgfpathquadraticcurveto{\pgfqpoint{4.111414in}{0.575581in}}{\pgfqpoint{4.119772in}{0.584659in}}%
\pgfpathquadraticcurveto{\pgfqpoint{4.128129in}{0.593756in}}{\pgfqpoint{4.142323in}{0.593756in}}%
\pgfpathquadraticcurveto{\pgfqpoint{4.155040in}{0.593756in}}{\pgfqpoint{4.162443in}{0.585560in}}%
\pgfpathquadraticcurveto{\pgfqpoint{4.169846in}{0.577383in}}{\pgfqpoint{4.169846in}{0.563315in}}%
\pgfpathlineto{\pgfqpoint{4.169846in}{0.563315in}}%
\pgfpathclose%
\pgfpathmoveto{\pgfqpoint{4.159489in}{0.566359in}}%
\pgfpathquadraticcurveto{\pgfqpoint{4.159381in}{0.574843in}}{\pgfqpoint{4.154733in}{0.579904in}}%
\pgfpathquadraticcurveto{\pgfqpoint{4.150086in}{0.584984in}}{\pgfqpoint{4.142431in}{0.584984in}}%
\pgfpathquadraticcurveto{\pgfqpoint{4.133767in}{0.584984in}}{\pgfqpoint{4.128562in}{0.580084in}}%
\pgfpathquadraticcurveto{\pgfqpoint{4.123356in}{0.575185in}}{\pgfqpoint{4.122564in}{0.566287in}}%
\pgfpathlineto{\pgfqpoint{4.159489in}{0.566359in}}%
\pgfpathlineto{\pgfqpoint{4.159489in}{0.566359in}}%
\pgfpathclose%
\pgfpathmoveto{\pgfqpoint{4.228331in}{0.582678in}}%
\pgfpathlineto{\pgfqpoint{4.228331in}{0.616793in}}%
\pgfpathlineto{\pgfqpoint{4.238688in}{0.616793in}}%
\pgfpathlineto{\pgfqpoint{4.238688in}{0.529200in}}%
\pgfpathlineto{\pgfqpoint{4.228331in}{0.529200in}}%
\pgfpathlineto{\pgfqpoint{4.228331in}{0.538656in}}%
\pgfpathquadraticcurveto{\pgfqpoint{4.225071in}{0.533037in}}{\pgfqpoint{4.220082in}{0.530299in}}%
\pgfpathquadraticcurveto{\pgfqpoint{4.215110in}{0.527561in}}{\pgfqpoint{4.208121in}{0.527561in}}%
\pgfpathquadraticcurveto{\pgfqpoint{4.196702in}{0.527561in}}{\pgfqpoint{4.189515in}{0.536675in}}%
\pgfpathquadraticcurveto{\pgfqpoint{4.182346in}{0.545807in}}{\pgfqpoint{4.182346in}{0.560667in}}%
\pgfpathquadraticcurveto{\pgfqpoint{4.182346in}{0.575527in}}{\pgfqpoint{4.189515in}{0.584641in}}%
\pgfpathquadraticcurveto{\pgfqpoint{4.196702in}{0.593756in}}{\pgfqpoint{4.208121in}{0.593756in}}%
\pgfpathquadraticcurveto{\pgfqpoint{4.215110in}{0.593756in}}{\pgfqpoint{4.220082in}{0.591018in}}%
\pgfpathquadraticcurveto{\pgfqpoint{4.225071in}{0.588298in}}{\pgfqpoint{4.228331in}{0.582678in}}%
\pgfpathlineto{\pgfqpoint{4.228331in}{0.582678in}}%
\pgfpathclose%
\pgfpathmoveto{\pgfqpoint{4.193045in}{0.560667in}}%
\pgfpathquadraticcurveto{\pgfqpoint{4.193045in}{0.549248in}}{\pgfqpoint{4.197746in}{0.542745in}}%
\pgfpathquadraticcurveto{\pgfqpoint{4.202448in}{0.536243in}}{\pgfqpoint{4.210661in}{0.536243in}}%
\pgfpathquadraticcurveto{\pgfqpoint{4.218875in}{0.536243in}}{\pgfqpoint{4.223594in}{0.542745in}}%
\pgfpathquadraticcurveto{\pgfqpoint{4.228331in}{0.549248in}}{\pgfqpoint{4.228331in}{0.560667in}}%
\pgfpathquadraticcurveto{\pgfqpoint{4.228331in}{0.572087in}}{\pgfqpoint{4.223594in}{0.578589in}}%
\pgfpathquadraticcurveto{\pgfqpoint{4.218875in}{0.585092in}}{\pgfqpoint{4.210661in}{0.585092in}}%
\pgfpathquadraticcurveto{\pgfqpoint{4.202448in}{0.585092in}}{\pgfqpoint{4.197746in}{0.578589in}}%
\pgfpathquadraticcurveto{\pgfqpoint{4.193045in}{0.572087in}}{\pgfqpoint{4.193045in}{0.560667in}}%
\pgfpathlineto{\pgfqpoint{4.193045in}{0.560667in}}%
\pgfpathclose%
\pgfusepath{fill}%
\end{pgfscope}%
\begin{pgfscope}%
\pgfsetbuttcap%
\pgfsetroundjoin%
\definecolor{currentfill}{rgb}{0.309804,0.372549,0.419608}%
\pgfsetfillcolor{currentfill}%
\pgfsetlinewidth{0.000000pt}%
\definecolor{currentstroke}{rgb}{0.309804,0.372549,0.419608}%
\pgfsetstrokecolor{currentstroke}%
\pgfsetdash{}{0pt}%
\pgfpathmoveto{\pgfqpoint{4.377772in}{0.561460in}}%
\pgfpathquadraticcurveto{\pgfqpoint{4.377772in}{0.572717in}}{\pgfqpoint{4.373125in}{0.578896in}}%
\pgfpathquadraticcurveto{\pgfqpoint{4.368496in}{0.585092in}}{\pgfqpoint{4.360102in}{0.585092in}}%
\pgfpathquadraticcurveto{\pgfqpoint{4.351780in}{0.585092in}}{\pgfqpoint{4.347133in}{0.578896in}}%
\pgfpathquadraticcurveto{\pgfqpoint{4.342486in}{0.572717in}}{\pgfqpoint{4.342486in}{0.561460in}}%
\pgfpathquadraticcurveto{\pgfqpoint{4.342486in}{0.550256in}}{\pgfqpoint{4.347133in}{0.544060in}}%
\pgfpathquadraticcurveto{\pgfqpoint{4.351780in}{0.537864in}}{\pgfqpoint{4.360102in}{0.537864in}}%
\pgfpathquadraticcurveto{\pgfqpoint{4.368496in}{0.537864in}}{\pgfqpoint{4.373125in}{0.544060in}}%
\pgfpathquadraticcurveto{\pgfqpoint{4.377772in}{0.550256in}}{\pgfqpoint{4.377772in}{0.561460in}}%
\pgfpathlineto{\pgfqpoint{4.377772in}{0.561460in}}%
\pgfpathclose%
\pgfpathmoveto{\pgfqpoint{4.388129in}{0.537017in}}%
\pgfpathquadraticcurveto{\pgfqpoint{4.388129in}{0.520932in}}{\pgfqpoint{4.380978in}{0.513079in}}%
\pgfpathquadraticcurveto{\pgfqpoint{4.373845in}{0.505226in}}{\pgfqpoint{4.359093in}{0.505226in}}%
\pgfpathquadraticcurveto{\pgfqpoint{4.353636in}{0.505226in}}{\pgfqpoint{4.348790in}{0.506036in}}%
\pgfpathquadraticcurveto{\pgfqpoint{4.343945in}{0.506847in}}{\pgfqpoint{4.339388in}{0.508540in}}%
\pgfpathlineto{\pgfqpoint{4.339388in}{0.518609in}}%
\pgfpathquadraticcurveto{\pgfqpoint{4.343945in}{0.516141in}}{\pgfqpoint{4.348394in}{0.514970in}}%
\pgfpathquadraticcurveto{\pgfqpoint{4.352843in}{0.513782in}}{\pgfqpoint{4.357454in}{0.513782in}}%
\pgfpathquadraticcurveto{\pgfqpoint{4.367649in}{0.513782in}}{\pgfqpoint{4.372710in}{0.519095in}}%
\pgfpathquadraticcurveto{\pgfqpoint{4.377772in}{0.524409in}}{\pgfqpoint{4.377772in}{0.535162in}}%
\pgfpathlineto{\pgfqpoint{4.377772in}{0.540295in}}%
\pgfpathquadraticcurveto{\pgfqpoint{4.374566in}{0.534712in}}{\pgfqpoint{4.369558in}{0.531956in}}%
\pgfpathquadraticcurveto{\pgfqpoint{4.364551in}{0.529200in}}{\pgfqpoint{4.357562in}{0.529200in}}%
\pgfpathquadraticcurveto{\pgfqpoint{4.345980in}{0.529200in}}{\pgfqpoint{4.338884in}{0.538026in}}%
\pgfpathquadraticcurveto{\pgfqpoint{4.331787in}{0.546870in}}{\pgfqpoint{4.331787in}{0.561460in}}%
\pgfpathquadraticcurveto{\pgfqpoint{4.331787in}{0.576086in}}{\pgfqpoint{4.338884in}{0.584912in}}%
\pgfpathquadraticcurveto{\pgfqpoint{4.345980in}{0.593756in}}{\pgfqpoint{4.357562in}{0.593756in}}%
\pgfpathquadraticcurveto{\pgfqpoint{4.364551in}{0.593756in}}{\pgfqpoint{4.369558in}{0.591000in}}%
\pgfpathquadraticcurveto{\pgfqpoint{4.374566in}{0.588244in}}{\pgfqpoint{4.377772in}{0.582678in}}%
\pgfpathlineto{\pgfqpoint{4.377772in}{0.592243in}}%
\pgfpathlineto{\pgfqpoint{4.388129in}{0.592243in}}%
\pgfpathlineto{\pgfqpoint{4.388129in}{0.537017in}}%
\pgfpathlineto{\pgfqpoint{4.388129in}{0.537017in}}%
\pgfpathclose%
\pgfpathmoveto{\pgfqpoint{4.438131in}{0.560883in}}%
\pgfpathquadraticcurveto{\pgfqpoint{4.425576in}{0.560883in}}{\pgfqpoint{4.420731in}{0.558019in}}%
\pgfpathquadraticcurveto{\pgfqpoint{4.415886in}{0.555156in}}{\pgfqpoint{4.415886in}{0.548221in}}%
\pgfpathquadraticcurveto{\pgfqpoint{4.415886in}{0.542709in}}{\pgfqpoint{4.419524in}{0.539467in}}%
\pgfpathquadraticcurveto{\pgfqpoint{4.423163in}{0.536243in}}{\pgfqpoint{4.429395in}{0.536243in}}%
\pgfpathquadraticcurveto{\pgfqpoint{4.438023in}{0.536243in}}{\pgfqpoint{4.443228in}{0.542349in}}%
\pgfpathquadraticcurveto{\pgfqpoint{4.448434in}{0.548455in}}{\pgfqpoint{4.448434in}{0.558578in}}%
\pgfpathlineto{\pgfqpoint{4.448434in}{0.560883in}}%
\pgfpathlineto{\pgfqpoint{4.438131in}{0.560883in}}%
\pgfpathlineto{\pgfqpoint{4.438131in}{0.560883in}}%
\pgfpathclose%
\pgfpathmoveto{\pgfqpoint{4.458791in}{0.565170in}}%
\pgfpathlineto{\pgfqpoint{4.458791in}{0.529200in}}%
\pgfpathlineto{\pgfqpoint{4.448434in}{0.529200in}}%
\pgfpathlineto{\pgfqpoint{4.448434in}{0.538764in}}%
\pgfpathquadraticcurveto{\pgfqpoint{4.444885in}{0.533037in}}{\pgfqpoint{4.439590in}{0.530299in}}%
\pgfpathquadraticcurveto{\pgfqpoint{4.434294in}{0.527561in}}{\pgfqpoint{4.426639in}{0.527561in}}%
\pgfpathquadraticcurveto{\pgfqpoint{4.416966in}{0.527561in}}{\pgfqpoint{4.411238in}{0.533001in}}%
\pgfpathquadraticcurveto{\pgfqpoint{4.405529in}{0.538440in}}{\pgfqpoint{4.405529in}{0.547554in}}%
\pgfpathquadraticcurveto{\pgfqpoint{4.405529in}{0.558182in}}{\pgfqpoint{4.412643in}{0.563585in}}%
\pgfpathquadraticcurveto{\pgfqpoint{4.419776in}{0.568989in}}{\pgfqpoint{4.433898in}{0.568989in}}%
\pgfpathlineto{\pgfqpoint{4.448434in}{0.568989in}}%
\pgfpathlineto{\pgfqpoint{4.448434in}{0.570016in}}%
\pgfpathquadraticcurveto{\pgfqpoint{4.448434in}{0.577166in}}{\pgfqpoint{4.443732in}{0.581075in}}%
\pgfpathquadraticcurveto{\pgfqpoint{4.439031in}{0.584984in}}{\pgfqpoint{4.430529in}{0.584984in}}%
\pgfpathquadraticcurveto{\pgfqpoint{4.425126in}{0.584984in}}{\pgfqpoint{4.419992in}{0.583687in}}%
\pgfpathquadraticcurveto{\pgfqpoint{4.414877in}{0.582390in}}{\pgfqpoint{4.410158in}{0.579796in}}%
\pgfpathlineto{\pgfqpoint{4.410158in}{0.589379in}}%
\pgfpathquadraticcurveto{\pgfqpoint{4.415832in}{0.591576in}}{\pgfqpoint{4.421181in}{0.592657in}}%
\pgfpathquadraticcurveto{\pgfqpoint{4.426531in}{0.593756in}}{\pgfqpoint{4.431592in}{0.593756in}}%
\pgfpathquadraticcurveto{\pgfqpoint{4.445281in}{0.593756in}}{\pgfqpoint{4.452036in}{0.586659in}}%
\pgfpathquadraticcurveto{\pgfqpoint{4.458791in}{0.579580in}}{\pgfqpoint{4.458791in}{0.565170in}}%
\pgfpathlineto{\pgfqpoint{4.458791in}{0.565170in}}%
\pgfpathclose%
\pgfpathmoveto{\pgfqpoint{4.490132in}{0.538656in}}%
\pgfpathlineto{\pgfqpoint{4.490132in}{0.505226in}}%
\pgfpathlineto{\pgfqpoint{4.479721in}{0.505226in}}%
\pgfpathlineto{\pgfqpoint{4.479721in}{0.592243in}}%
\pgfpathlineto{\pgfqpoint{4.490132in}{0.592243in}}%
\pgfpathlineto{\pgfqpoint{4.490132in}{0.582678in}}%
\pgfpathquadraticcurveto{\pgfqpoint{4.493410in}{0.588298in}}{\pgfqpoint{4.498381in}{0.591018in}}%
\pgfpathquadraticcurveto{\pgfqpoint{4.503371in}{0.593756in}}{\pgfqpoint{4.510287in}{0.593756in}}%
\pgfpathquadraticcurveto{\pgfqpoint{4.521779in}{0.593756in}}{\pgfqpoint{4.528948in}{0.584641in}}%
\pgfpathquadraticcurveto{\pgfqpoint{4.536135in}{0.575527in}}{\pgfqpoint{4.536135in}{0.560667in}}%
\pgfpathquadraticcurveto{\pgfqpoint{4.536135in}{0.545807in}}{\pgfqpoint{4.528948in}{0.536675in}}%
\pgfpathquadraticcurveto{\pgfqpoint{4.521779in}{0.527561in}}{\pgfqpoint{4.510287in}{0.527561in}}%
\pgfpathquadraticcurveto{\pgfqpoint{4.503371in}{0.527561in}}{\pgfqpoint{4.498381in}{0.530299in}}%
\pgfpathquadraticcurveto{\pgfqpoint{4.493410in}{0.533037in}}{\pgfqpoint{4.490132in}{0.538656in}}%
\pgfpathlineto{\pgfqpoint{4.490132in}{0.538656in}}%
\pgfpathclose%
\pgfpathmoveto{\pgfqpoint{4.525381in}{0.560667in}}%
\pgfpathquadraticcurveto{\pgfqpoint{4.525381in}{0.572087in}}{\pgfqpoint{4.520680in}{0.578589in}}%
\pgfpathquadraticcurveto{\pgfqpoint{4.515979in}{0.585092in}}{\pgfqpoint{4.507766in}{0.585092in}}%
\pgfpathquadraticcurveto{\pgfqpoint{4.499534in}{0.585092in}}{\pgfqpoint{4.494833in}{0.578589in}}%
\pgfpathquadraticcurveto{\pgfqpoint{4.490132in}{0.572087in}}{\pgfqpoint{4.490132in}{0.560667in}}%
\pgfpathquadraticcurveto{\pgfqpoint{4.490132in}{0.549248in}}{\pgfqpoint{4.494833in}{0.542745in}}%
\pgfpathquadraticcurveto{\pgfqpoint{4.499534in}{0.536243in}}{\pgfqpoint{4.507766in}{0.536243in}}%
\pgfpathquadraticcurveto{\pgfqpoint{4.515979in}{0.536243in}}{\pgfqpoint{4.520680in}{0.542745in}}%
\pgfpathquadraticcurveto{\pgfqpoint{4.525381in}{0.549248in}}{\pgfqpoint{4.525381in}{0.560667in}}%
\pgfpathlineto{\pgfqpoint{4.525381in}{0.560667in}}%
\pgfpathclose%
\pgfusepath{fill}%
\end{pgfscope}%
\begin{pgfscope}%
\pgfsetbuttcap%
\pgfsetroundjoin%
\definecolor{currentfill}{rgb}{0.309804,0.372549,0.419608}%
\pgfsetfillcolor{currentfill}%
\pgfsetlinewidth{0.000000pt}%
\definecolor{currentstroke}{rgb}{0.309804,0.372549,0.419608}%
\pgfsetstrokecolor{currentstroke}%
\pgfsetdash{}{0pt}%
\pgfpathmoveto{\pgfqpoint{4.622792in}{0.592243in}}%
\pgfpathlineto{\pgfqpoint{4.633149in}{0.592243in}}%
\pgfpathlineto{\pgfqpoint{4.633149in}{0.529200in}}%
\pgfpathlineto{\pgfqpoint{4.622792in}{0.529200in}}%
\pgfpathlineto{\pgfqpoint{4.622792in}{0.592243in}}%
\pgfpathlineto{\pgfqpoint{4.622792in}{0.592243in}}%
\pgfpathclose%
\pgfpathmoveto{\pgfqpoint{4.622792in}{0.616793in}}%
\pgfpathlineto{\pgfqpoint{4.633149in}{0.616793in}}%
\pgfpathlineto{\pgfqpoint{4.633149in}{0.603662in}}%
\pgfpathlineto{\pgfqpoint{4.622792in}{0.603662in}}%
\pgfpathlineto{\pgfqpoint{4.622792in}{0.616793in}}%
\pgfpathlineto{\pgfqpoint{4.622792in}{0.616793in}}%
\pgfpathclose%
\pgfpathmoveto{\pgfqpoint{4.695003in}{0.590387in}}%
\pgfpathlineto{\pgfqpoint{4.695003in}{0.580589in}}%
\pgfpathquadraticcurveto{\pgfqpoint{4.690626in}{0.582840in}}{\pgfqpoint{4.685889in}{0.583957in}}%
\pgfpathquadraticcurveto{\pgfqpoint{4.681170in}{0.585092in}}{\pgfqpoint{4.676090in}{0.585092in}}%
\pgfpathquadraticcurveto{\pgfqpoint{4.668381in}{0.585092in}}{\pgfqpoint{4.664526in}{0.582732in}}%
\pgfpathquadraticcurveto{\pgfqpoint{4.660672in}{0.580373in}}{\pgfqpoint{4.660672in}{0.575635in}}%
\pgfpathquadraticcurveto{\pgfqpoint{4.660672in}{0.572033in}}{\pgfqpoint{4.663428in}{0.569980in}}%
\pgfpathquadraticcurveto{\pgfqpoint{4.666183in}{0.567926in}}{\pgfqpoint{4.674523in}{0.566071in}}%
\pgfpathlineto{\pgfqpoint{4.678071in}{0.565278in}}%
\pgfpathquadraticcurveto{\pgfqpoint{4.689095in}{0.562919in}}{\pgfqpoint{4.693742in}{0.558614in}}%
\pgfpathquadraticcurveto{\pgfqpoint{4.698389in}{0.554309in}}{\pgfqpoint{4.698389in}{0.546600in}}%
\pgfpathquadraticcurveto{\pgfqpoint{4.698389in}{0.537810in}}{\pgfqpoint{4.691436in}{0.532676in}}%
\pgfpathquadraticcurveto{\pgfqpoint{4.684484in}{0.527561in}}{\pgfqpoint{4.672326in}{0.527561in}}%
\pgfpathquadraticcurveto{\pgfqpoint{4.667264in}{0.527561in}}{\pgfqpoint{4.661770in}{0.528552in}}%
\pgfpathquadraticcurveto{\pgfqpoint{4.656277in}{0.529542in}}{\pgfqpoint{4.650207in}{0.531506in}}%
\pgfpathlineto{\pgfqpoint{4.650207in}{0.542205in}}%
\pgfpathquadraticcurveto{\pgfqpoint{4.655953in}{0.539215in}}{\pgfqpoint{4.661518in}{0.537720in}}%
\pgfpathquadraticcurveto{\pgfqpoint{4.667084in}{0.536243in}}{\pgfqpoint{4.672560in}{0.536243in}}%
\pgfpathquadraticcurveto{\pgfqpoint{4.679873in}{0.536243in}}{\pgfqpoint{4.683799in}{0.538746in}}%
\pgfpathquadraticcurveto{\pgfqpoint{4.687744in}{0.541250in}}{\pgfqpoint{4.687744in}{0.545807in}}%
\pgfpathquadraticcurveto{\pgfqpoint{4.687744in}{0.550022in}}{\pgfqpoint{4.684898in}{0.552274in}}%
\pgfpathquadraticcurveto{\pgfqpoint{4.682070in}{0.554525in}}{\pgfqpoint{4.672434in}{0.556614in}}%
\pgfpathlineto{\pgfqpoint{4.668831in}{0.557461in}}%
\pgfpathquadraticcurveto{\pgfqpoint{4.659213in}{0.559478in}}{\pgfqpoint{4.654926in}{0.563675in}}%
\pgfpathquadraticcurveto{\pgfqpoint{4.650657in}{0.567872in}}{\pgfqpoint{4.650657in}{0.575185in}}%
\pgfpathquadraticcurveto{\pgfqpoint{4.650657in}{0.584083in}}{\pgfqpoint{4.656961in}{0.588910in}}%
\pgfpathquadraticcurveto{\pgfqpoint{4.663265in}{0.593756in}}{\pgfqpoint{4.674865in}{0.593756in}}%
\pgfpathquadraticcurveto{\pgfqpoint{4.680593in}{0.593756in}}{\pgfqpoint{4.685655in}{0.592909in}}%
\pgfpathquadraticcurveto{\pgfqpoint{4.690734in}{0.592080in}}{\pgfqpoint{4.695003in}{0.590387in}}%
\pgfpathlineto{\pgfqpoint{4.695003in}{0.590387in}}%
\pgfpathclose%
\pgfusepath{fill}%
\end{pgfscope}%
\begin{pgfscope}%
\pgfsetbuttcap%
\pgfsetroundjoin%
\definecolor{currentfill}{rgb}{0.309804,0.372549,0.419608}%
\pgfsetfillcolor{currentfill}%
\pgfsetlinewidth{0.000000pt}%
\definecolor{currentstroke}{rgb}{0.309804,0.372549,0.419608}%
\pgfsetstrokecolor{currentstroke}%
\pgfsetdash{}{0pt}%
\pgfpathmoveto{\pgfqpoint{4.821710in}{0.567260in}}%
\pgfpathlineto{\pgfqpoint{4.821710in}{0.529200in}}%
\pgfpathlineto{\pgfqpoint{4.811353in}{0.529200in}}%
\pgfpathlineto{\pgfqpoint{4.811353in}{0.566917in}}%
\pgfpathquadraticcurveto{\pgfqpoint{4.811353in}{0.575869in}}{\pgfqpoint{4.807858in}{0.580300in}}%
\pgfpathquadraticcurveto{\pgfqpoint{4.804364in}{0.584749in}}{\pgfqpoint{4.797393in}{0.584749in}}%
\pgfpathquadraticcurveto{\pgfqpoint{4.789000in}{0.584749in}}{\pgfqpoint{4.784154in}{0.579400in}}%
\pgfpathquadraticcurveto{\pgfqpoint{4.779309in}{0.574068in}}{\pgfqpoint{4.779309in}{0.564828in}}%
\pgfpathlineto{\pgfqpoint{4.779309in}{0.529200in}}%
\pgfpathlineto{\pgfqpoint{4.768898in}{0.529200in}}%
\pgfpathlineto{\pgfqpoint{4.768898in}{0.592243in}}%
\pgfpathlineto{\pgfqpoint{4.779309in}{0.592243in}}%
\pgfpathlineto{\pgfqpoint{4.779309in}{0.582444in}}%
\pgfpathquadraticcurveto{\pgfqpoint{4.783038in}{0.588136in}}{\pgfqpoint{4.788063in}{0.590946in}}%
\pgfpathquadraticcurveto{\pgfqpoint{4.793106in}{0.593756in}}{\pgfqpoint{4.799699in}{0.593756in}}%
\pgfpathquadraticcurveto{\pgfqpoint{4.810560in}{0.593756in}}{\pgfqpoint{4.816126in}{0.587037in}}%
\pgfpathquadraticcurveto{\pgfqpoint{4.821710in}{0.580318in}}{\pgfqpoint{4.821710in}{0.567260in}}%
\pgfpathlineto{\pgfqpoint{4.821710in}{0.567260in}}%
\pgfpathclose%
\pgfpathmoveto{\pgfqpoint{4.866776in}{0.584984in}}%
\pgfpathquadraticcurveto{\pgfqpoint{4.858455in}{0.584984in}}{\pgfqpoint{4.853609in}{0.578481in}}%
\pgfpathquadraticcurveto{\pgfqpoint{4.848764in}{0.571979in}}{\pgfqpoint{4.848764in}{0.560667in}}%
\pgfpathquadraticcurveto{\pgfqpoint{4.848764in}{0.549356in}}{\pgfqpoint{4.853573in}{0.542853in}}%
\pgfpathquadraticcurveto{\pgfqpoint{4.858400in}{0.536351in}}{\pgfqpoint{4.866776in}{0.536351in}}%
\pgfpathquadraticcurveto{\pgfqpoint{4.875062in}{0.536351in}}{\pgfqpoint{4.879889in}{0.542871in}}%
\pgfpathquadraticcurveto{\pgfqpoint{4.884734in}{0.549410in}}{\pgfqpoint{4.884734in}{0.560667in}}%
\pgfpathquadraticcurveto{\pgfqpoint{4.884734in}{0.571871in}}{\pgfqpoint{4.879889in}{0.578427in}}%
\pgfpathquadraticcurveto{\pgfqpoint{4.875062in}{0.584984in}}{\pgfqpoint{4.866776in}{0.584984in}}%
\pgfpathlineto{\pgfqpoint{4.866776in}{0.584984in}}%
\pgfpathclose%
\pgfpathmoveto{\pgfqpoint{4.866776in}{0.593756in}}%
\pgfpathquadraticcurveto{\pgfqpoint{4.880285in}{0.593756in}}{\pgfqpoint{4.887994in}{0.584966in}}%
\pgfpathquadraticcurveto{\pgfqpoint{4.895722in}{0.576194in}}{\pgfqpoint{4.895722in}{0.560667in}}%
\pgfpathquadraticcurveto{\pgfqpoint{4.895722in}{0.545195in}}{\pgfqpoint{4.887994in}{0.536369in}}%
\pgfpathquadraticcurveto{\pgfqpoint{4.880285in}{0.527561in}}{\pgfqpoint{4.866776in}{0.527561in}}%
\pgfpathquadraticcurveto{\pgfqpoint{4.853213in}{0.527561in}}{\pgfqpoint{4.845522in}{0.536369in}}%
\pgfpathquadraticcurveto{\pgfqpoint{4.837849in}{0.545195in}}{\pgfqpoint{4.837849in}{0.560667in}}%
\pgfpathquadraticcurveto{\pgfqpoint{4.837849in}{0.576194in}}{\pgfqpoint{4.845522in}{0.584966in}}%
\pgfpathquadraticcurveto{\pgfqpoint{4.853213in}{0.593756in}}{\pgfqpoint{4.866776in}{0.593756in}}%
\pgfpathlineto{\pgfqpoint{4.866776in}{0.593756in}}%
\pgfpathclose%
\pgfpathmoveto{\pgfqpoint{4.923136in}{0.610147in}}%
\pgfpathlineto{\pgfqpoint{4.923136in}{0.592243in}}%
\pgfpathlineto{\pgfqpoint{4.944463in}{0.592243in}}%
\pgfpathlineto{\pgfqpoint{4.944463in}{0.584191in}}%
\pgfpathlineto{\pgfqpoint{4.923136in}{0.584191in}}%
\pgfpathlineto{\pgfqpoint{4.923136in}{0.549968in}}%
\pgfpathquadraticcurveto{\pgfqpoint{4.923136in}{0.542259in}}{\pgfqpoint{4.925244in}{0.540061in}}%
\pgfpathquadraticcurveto{\pgfqpoint{4.927351in}{0.537864in}}{\pgfqpoint{4.933835in}{0.537864in}}%
\pgfpathlineto{\pgfqpoint{4.944463in}{0.537864in}}%
\pgfpathlineto{\pgfqpoint{4.944463in}{0.529200in}}%
\pgfpathlineto{\pgfqpoint{4.933835in}{0.529200in}}%
\pgfpathquadraticcurveto{\pgfqpoint{4.921839in}{0.529200in}}{\pgfqpoint{4.917282in}{0.533667in}}%
\pgfpathquadraticcurveto{\pgfqpoint{4.912725in}{0.538152in}}{\pgfqpoint{4.912725in}{0.549968in}}%
\pgfpathlineto{\pgfqpoint{4.912725in}{0.584191in}}%
\pgfpathlineto{\pgfqpoint{4.905124in}{0.584191in}}%
\pgfpathlineto{\pgfqpoint{4.905124in}{0.592243in}}%
\pgfpathlineto{\pgfqpoint{4.912725in}{0.592243in}}%
\pgfpathlineto{\pgfqpoint{4.912725in}{0.610147in}}%
\pgfpathlineto{\pgfqpoint{4.923136in}{0.610147in}}%
\pgfpathlineto{\pgfqpoint{4.923136in}{0.610147in}}%
\pgfpathclose%
\pgfusepath{fill}%
\end{pgfscope}%
\begin{pgfscope}%
\pgfsetbuttcap%
\pgfsetroundjoin%
\definecolor{currentfill}{rgb}{0.309804,0.372549,0.419608}%
\pgfsetfillcolor{currentfill}%
\pgfsetlinewidth{0.000000pt}%
\definecolor{currentstroke}{rgb}{0.309804,0.372549,0.419608}%
\pgfsetstrokecolor{currentstroke}%
\pgfsetdash{}{0pt}%
\pgfpathmoveto{\pgfqpoint{5.036046in}{0.610147in}}%
\pgfpathlineto{\pgfqpoint{5.036046in}{0.592243in}}%
\pgfpathlineto{\pgfqpoint{5.057372in}{0.592243in}}%
\pgfpathlineto{\pgfqpoint{5.057372in}{0.584191in}}%
\pgfpathlineto{\pgfqpoint{5.036046in}{0.584191in}}%
\pgfpathlineto{\pgfqpoint{5.036046in}{0.549968in}}%
\pgfpathquadraticcurveto{\pgfqpoint{5.036046in}{0.542259in}}{\pgfqpoint{5.038153in}{0.540061in}}%
\pgfpathquadraticcurveto{\pgfqpoint{5.040260in}{0.537864in}}{\pgfqpoint{5.046745in}{0.537864in}}%
\pgfpathlineto{\pgfqpoint{5.057372in}{0.537864in}}%
\pgfpathlineto{\pgfqpoint{5.057372in}{0.529200in}}%
\pgfpathlineto{\pgfqpoint{5.046745in}{0.529200in}}%
\pgfpathquadraticcurveto{\pgfqpoint{5.034749in}{0.529200in}}{\pgfqpoint{5.030192in}{0.533667in}}%
\pgfpathquadraticcurveto{\pgfqpoint{5.025634in}{0.538152in}}{\pgfqpoint{5.025634in}{0.549968in}}%
\pgfpathlineto{\pgfqpoint{5.025634in}{0.584191in}}%
\pgfpathlineto{\pgfqpoint{5.018033in}{0.584191in}}%
\pgfpathlineto{\pgfqpoint{5.018033in}{0.592243in}}%
\pgfpathlineto{\pgfqpoint{5.025634in}{0.592243in}}%
\pgfpathlineto{\pgfqpoint{5.025634in}{0.610147in}}%
\pgfpathlineto{\pgfqpoint{5.036046in}{0.610147in}}%
\pgfpathlineto{\pgfqpoint{5.036046in}{0.610147in}}%
\pgfpathclose%
\pgfpathmoveto{\pgfqpoint{5.123404in}{0.567260in}}%
\pgfpathlineto{\pgfqpoint{5.123404in}{0.529200in}}%
\pgfpathlineto{\pgfqpoint{5.113047in}{0.529200in}}%
\pgfpathlineto{\pgfqpoint{5.113047in}{0.566917in}}%
\pgfpathquadraticcurveto{\pgfqpoint{5.113047in}{0.575869in}}{\pgfqpoint{5.109553in}{0.580300in}}%
\pgfpathquadraticcurveto{\pgfqpoint{5.106059in}{0.584749in}}{\pgfqpoint{5.099088in}{0.584749in}}%
\pgfpathquadraticcurveto{\pgfqpoint{5.090694in}{0.584749in}}{\pgfqpoint{5.085849in}{0.579400in}}%
\pgfpathquadraticcurveto{\pgfqpoint{5.081004in}{0.574068in}}{\pgfqpoint{5.081004in}{0.564828in}}%
\pgfpathlineto{\pgfqpoint{5.081004in}{0.529200in}}%
\pgfpathlineto{\pgfqpoint{5.070593in}{0.529200in}}%
\pgfpathlineto{\pgfqpoint{5.070593in}{0.616793in}}%
\pgfpathlineto{\pgfqpoint{5.081004in}{0.616793in}}%
\pgfpathlineto{\pgfqpoint{5.081004in}{0.582444in}}%
\pgfpathquadraticcurveto{\pgfqpoint{5.084732in}{0.588136in}}{\pgfqpoint{5.089758in}{0.590946in}}%
\pgfpathquadraticcurveto{\pgfqpoint{5.094801in}{0.593756in}}{\pgfqpoint{5.101394in}{0.593756in}}%
\pgfpathquadraticcurveto{\pgfqpoint{5.112255in}{0.593756in}}{\pgfqpoint{5.117821in}{0.587037in}}%
\pgfpathquadraticcurveto{\pgfqpoint{5.123404in}{0.580318in}}{\pgfqpoint{5.123404in}{0.567260in}}%
\pgfpathlineto{\pgfqpoint{5.123404in}{0.567260in}}%
\pgfpathclose%
\pgfpathmoveto{\pgfqpoint{5.197975in}{0.563315in}}%
\pgfpathlineto{\pgfqpoint{5.197975in}{0.558254in}}%
\pgfpathlineto{\pgfqpoint{5.150351in}{0.558254in}}%
\pgfpathquadraticcurveto{\pgfqpoint{5.151035in}{0.547554in}}{\pgfqpoint{5.156799in}{0.541953in}}%
\pgfpathquadraticcurveto{\pgfqpoint{5.162563in}{0.536351in}}{\pgfqpoint{5.172866in}{0.536351in}}%
\pgfpathquadraticcurveto{\pgfqpoint{5.178828in}{0.536351in}}{\pgfqpoint{5.184430in}{0.537810in}}%
\pgfpathquadraticcurveto{\pgfqpoint{5.190031in}{0.539269in}}{\pgfqpoint{5.195561in}{0.542205in}}%
\pgfpathlineto{\pgfqpoint{5.195561in}{0.532406in}}%
\pgfpathquadraticcurveto{\pgfqpoint{5.189977in}{0.530047in}}{\pgfqpoint{5.184123in}{0.528804in}}%
\pgfpathquadraticcurveto{\pgfqpoint{5.178269in}{0.527561in}}{\pgfqpoint{5.172253in}{0.527561in}}%
\pgfpathquadraticcurveto{\pgfqpoint{5.157159in}{0.527561in}}{\pgfqpoint{5.148351in}{0.536333in}}%
\pgfpathquadraticcurveto{\pgfqpoint{5.139543in}{0.545123in}}{\pgfqpoint{5.139543in}{0.560109in}}%
\pgfpathquadraticcurveto{\pgfqpoint{5.139543in}{0.575581in}}{\pgfqpoint{5.147901in}{0.584659in}}%
\pgfpathquadraticcurveto{\pgfqpoint{5.156259in}{0.593756in}}{\pgfqpoint{5.170452in}{0.593756in}}%
\pgfpathquadraticcurveto{\pgfqpoint{5.183169in}{0.593756in}}{\pgfqpoint{5.190572in}{0.585560in}}%
\pgfpathquadraticcurveto{\pgfqpoint{5.197975in}{0.577383in}}{\pgfqpoint{5.197975in}{0.563315in}}%
\pgfpathlineto{\pgfqpoint{5.197975in}{0.563315in}}%
\pgfpathclose%
\pgfpathmoveto{\pgfqpoint{5.187618in}{0.566359in}}%
\pgfpathquadraticcurveto{\pgfqpoint{5.187510in}{0.574843in}}{\pgfqpoint{5.182863in}{0.579904in}}%
\pgfpathquadraticcurveto{\pgfqpoint{5.178215in}{0.584984in}}{\pgfqpoint{5.170560in}{0.584984in}}%
\pgfpathquadraticcurveto{\pgfqpoint{5.161896in}{0.584984in}}{\pgfqpoint{5.156691in}{0.580084in}}%
\pgfpathquadraticcurveto{\pgfqpoint{5.151485in}{0.575185in}}{\pgfqpoint{5.150693in}{0.566287in}}%
\pgfpathlineto{\pgfqpoint{5.187618in}{0.566359in}}%
\pgfpathlineto{\pgfqpoint{5.187618in}{0.566359in}}%
\pgfpathclose%
\pgfusepath{fill}%
\end{pgfscope}%
\begin{pgfscope}%
\pgfsetbuttcap%
\pgfsetroundjoin%
\definecolor{currentfill}{rgb}{0.309804,0.372549,0.419608}%
\pgfsetfillcolor{currentfill}%
\pgfsetlinewidth{0.000000pt}%
\definecolor{currentstroke}{rgb}{0.309804,0.372549,0.419608}%
\pgfsetstrokecolor{currentstroke}%
\pgfsetdash{}{0pt}%
\pgfpathmoveto{\pgfqpoint{5.292314in}{0.538656in}}%
\pgfpathlineto{\pgfqpoint{5.292314in}{0.505226in}}%
\pgfpathlineto{\pgfqpoint{5.281903in}{0.505226in}}%
\pgfpathlineto{\pgfqpoint{5.281903in}{0.592243in}}%
\pgfpathlineto{\pgfqpoint{5.292314in}{0.592243in}}%
\pgfpathlineto{\pgfqpoint{5.292314in}{0.582678in}}%
\pgfpathquadraticcurveto{\pgfqpoint{5.295592in}{0.588298in}}{\pgfqpoint{5.300563in}{0.591018in}}%
\pgfpathquadraticcurveto{\pgfqpoint{5.305553in}{0.593756in}}{\pgfqpoint{5.312469in}{0.593756in}}%
\pgfpathquadraticcurveto{\pgfqpoint{5.323961in}{0.593756in}}{\pgfqpoint{5.331130in}{0.584641in}}%
\pgfpathquadraticcurveto{\pgfqpoint{5.338317in}{0.575527in}}{\pgfqpoint{5.338317in}{0.560667in}}%
\pgfpathquadraticcurveto{\pgfqpoint{5.338317in}{0.545807in}}{\pgfqpoint{5.331130in}{0.536675in}}%
\pgfpathquadraticcurveto{\pgfqpoint{5.323961in}{0.527561in}}{\pgfqpoint{5.312469in}{0.527561in}}%
\pgfpathquadraticcurveto{\pgfqpoint{5.305553in}{0.527561in}}{\pgfqpoint{5.300563in}{0.530299in}}%
\pgfpathquadraticcurveto{\pgfqpoint{5.295592in}{0.533037in}}{\pgfqpoint{5.292314in}{0.538656in}}%
\pgfpathlineto{\pgfqpoint{5.292314in}{0.538656in}}%
\pgfpathclose%
\pgfpathmoveto{\pgfqpoint{5.327563in}{0.560667in}}%
\pgfpathquadraticcurveto{\pgfqpoint{5.327563in}{0.572087in}}{\pgfqpoint{5.322862in}{0.578589in}}%
\pgfpathquadraticcurveto{\pgfqpoint{5.318161in}{0.585092in}}{\pgfqpoint{5.309947in}{0.585092in}}%
\pgfpathquadraticcurveto{\pgfqpoint{5.301716in}{0.585092in}}{\pgfqpoint{5.297015in}{0.578589in}}%
\pgfpathquadraticcurveto{\pgfqpoint{5.292314in}{0.572087in}}{\pgfqpoint{5.292314in}{0.560667in}}%
\pgfpathquadraticcurveto{\pgfqpoint{5.292314in}{0.549248in}}{\pgfqpoint{5.297015in}{0.542745in}}%
\pgfpathquadraticcurveto{\pgfqpoint{5.301716in}{0.536243in}}{\pgfqpoint{5.309947in}{0.536243in}}%
\pgfpathquadraticcurveto{\pgfqpoint{5.318161in}{0.536243in}}{\pgfqpoint{5.322862in}{0.542745in}}%
\pgfpathquadraticcurveto{\pgfqpoint{5.327563in}{0.549248in}}{\pgfqpoint{5.327563in}{0.560667in}}%
\pgfpathlineto{\pgfqpoint{5.327563in}{0.560667in}}%
\pgfpathclose%
\pgfpathmoveto{\pgfqpoint{5.384140in}{0.560883in}}%
\pgfpathquadraticcurveto{\pgfqpoint{5.371585in}{0.560883in}}{\pgfqpoint{5.366740in}{0.558019in}}%
\pgfpathquadraticcurveto{\pgfqpoint{5.361895in}{0.555156in}}{\pgfqpoint{5.361895in}{0.548221in}}%
\pgfpathquadraticcurveto{\pgfqpoint{5.361895in}{0.542709in}}{\pgfqpoint{5.365533in}{0.539467in}}%
\pgfpathquadraticcurveto{\pgfqpoint{5.369171in}{0.536243in}}{\pgfqpoint{5.375404in}{0.536243in}}%
\pgfpathquadraticcurveto{\pgfqpoint{5.384031in}{0.536243in}}{\pgfqpoint{5.389237in}{0.542349in}}%
\pgfpathquadraticcurveto{\pgfqpoint{5.394442in}{0.548455in}}{\pgfqpoint{5.394442in}{0.558578in}}%
\pgfpathlineto{\pgfqpoint{5.394442in}{0.560883in}}%
\pgfpathlineto{\pgfqpoint{5.384140in}{0.560883in}}%
\pgfpathlineto{\pgfqpoint{5.384140in}{0.560883in}}%
\pgfpathclose%
\pgfpathmoveto{\pgfqpoint{5.404799in}{0.565170in}}%
\pgfpathlineto{\pgfqpoint{5.404799in}{0.529200in}}%
\pgfpathlineto{\pgfqpoint{5.394442in}{0.529200in}}%
\pgfpathlineto{\pgfqpoint{5.394442in}{0.538764in}}%
\pgfpathquadraticcurveto{\pgfqpoint{5.390894in}{0.533037in}}{\pgfqpoint{5.385599in}{0.530299in}}%
\pgfpathquadraticcurveto{\pgfqpoint{5.380303in}{0.527561in}}{\pgfqpoint{5.372648in}{0.527561in}}%
\pgfpathquadraticcurveto{\pgfqpoint{5.362975in}{0.527561in}}{\pgfqpoint{5.357247in}{0.533001in}}%
\pgfpathquadraticcurveto{\pgfqpoint{5.351538in}{0.538440in}}{\pgfqpoint{5.351538in}{0.547554in}}%
\pgfpathquadraticcurveto{\pgfqpoint{5.351538in}{0.558182in}}{\pgfqpoint{5.358652in}{0.563585in}}%
\pgfpathquadraticcurveto{\pgfqpoint{5.365785in}{0.568989in}}{\pgfqpoint{5.379907in}{0.568989in}}%
\pgfpathlineto{\pgfqpoint{5.394442in}{0.568989in}}%
\pgfpathlineto{\pgfqpoint{5.394442in}{0.570016in}}%
\pgfpathquadraticcurveto{\pgfqpoint{5.394442in}{0.577166in}}{\pgfqpoint{5.389741in}{0.581075in}}%
\pgfpathquadraticcurveto{\pgfqpoint{5.385040in}{0.584984in}}{\pgfqpoint{5.376538in}{0.584984in}}%
\pgfpathquadraticcurveto{\pgfqpoint{5.371135in}{0.584984in}}{\pgfqpoint{5.366001in}{0.583687in}}%
\pgfpathquadraticcurveto{\pgfqpoint{5.360886in}{0.582390in}}{\pgfqpoint{5.356167in}{0.579796in}}%
\pgfpathlineto{\pgfqpoint{5.356167in}{0.589379in}}%
\pgfpathquadraticcurveto{\pgfqpoint{5.361840in}{0.591576in}}{\pgfqpoint{5.367190in}{0.592657in}}%
\pgfpathquadraticcurveto{\pgfqpoint{5.372540in}{0.593756in}}{\pgfqpoint{5.377601in}{0.593756in}}%
\pgfpathquadraticcurveto{\pgfqpoint{5.391290in}{0.593756in}}{\pgfqpoint{5.398045in}{0.586659in}}%
\pgfpathquadraticcurveto{\pgfqpoint{5.404799in}{0.579580in}}{\pgfqpoint{5.404799in}{0.565170in}}%
\pgfpathlineto{\pgfqpoint{5.404799in}{0.565170in}}%
\pgfpathclose%
\pgfpathmoveto{\pgfqpoint{5.426126in}{0.592243in}}%
\pgfpathlineto{\pgfqpoint{5.436483in}{0.592243in}}%
\pgfpathlineto{\pgfqpoint{5.436483in}{0.529200in}}%
\pgfpathlineto{\pgfqpoint{5.426126in}{0.529200in}}%
\pgfpathlineto{\pgfqpoint{5.426126in}{0.592243in}}%
\pgfpathlineto{\pgfqpoint{5.426126in}{0.592243in}}%
\pgfpathclose%
\pgfpathmoveto{\pgfqpoint{5.426126in}{0.616793in}}%
\pgfpathlineto{\pgfqpoint{5.436483in}{0.616793in}}%
\pgfpathlineto{\pgfqpoint{5.436483in}{0.603662in}}%
\pgfpathlineto{\pgfqpoint{5.426126in}{0.603662in}}%
\pgfpathlineto{\pgfqpoint{5.426126in}{0.616793in}}%
\pgfpathlineto{\pgfqpoint{5.426126in}{0.616793in}}%
\pgfpathclose%
\pgfpathmoveto{\pgfqpoint{5.494680in}{0.582570in}}%
\pgfpathquadraticcurveto{\pgfqpoint{5.492933in}{0.583579in}}{\pgfqpoint{5.490880in}{0.584047in}}%
\pgfpathquadraticcurveto{\pgfqpoint{5.488826in}{0.584533in}}{\pgfqpoint{5.486359in}{0.584533in}}%
\pgfpathquadraticcurveto{\pgfqpoint{5.477569in}{0.584533in}}{\pgfqpoint{5.472867in}{0.578823in}}%
\pgfpathquadraticcurveto{\pgfqpoint{5.468166in}{0.573114in}}{\pgfqpoint{5.468166in}{0.562414in}}%
\pgfpathlineto{\pgfqpoint{5.468166in}{0.529200in}}%
\pgfpathlineto{\pgfqpoint{5.457755in}{0.529200in}}%
\pgfpathlineto{\pgfqpoint{5.457755in}{0.592243in}}%
\pgfpathlineto{\pgfqpoint{5.468166in}{0.592243in}}%
\pgfpathlineto{\pgfqpoint{5.468166in}{0.582444in}}%
\pgfpathquadraticcurveto{\pgfqpoint{5.471444in}{0.588190in}}{\pgfqpoint{5.476668in}{0.590964in}}%
\pgfpathquadraticcurveto{\pgfqpoint{5.481910in}{0.593756in}}{\pgfqpoint{5.489403in}{0.593756in}}%
\pgfpathquadraticcurveto{\pgfqpoint{5.490465in}{0.593756in}}{\pgfqpoint{5.491762in}{0.593611in}}%
\pgfpathquadraticcurveto{\pgfqpoint{5.493059in}{0.593485in}}{\pgfqpoint{5.494626in}{0.593197in}}%
\pgfpathlineto{\pgfqpoint{5.494680in}{0.582570in}}%
\pgfpathlineto{\pgfqpoint{5.494680in}{0.582570in}}%
\pgfpathclose%
\pgfpathmoveto{\pgfqpoint{5.556930in}{0.563315in}}%
\pgfpathlineto{\pgfqpoint{5.556930in}{0.558254in}}%
\pgfpathlineto{\pgfqpoint{5.509306in}{0.558254in}}%
\pgfpathquadraticcurveto{\pgfqpoint{5.509990in}{0.547554in}}{\pgfqpoint{5.515754in}{0.541953in}}%
\pgfpathquadraticcurveto{\pgfqpoint{5.521518in}{0.536351in}}{\pgfqpoint{5.531821in}{0.536351in}}%
\pgfpathquadraticcurveto{\pgfqpoint{5.537783in}{0.536351in}}{\pgfqpoint{5.543385in}{0.537810in}}%
\pgfpathquadraticcurveto{\pgfqpoint{5.548987in}{0.539269in}}{\pgfqpoint{5.554516in}{0.542205in}}%
\pgfpathlineto{\pgfqpoint{5.554516in}{0.532406in}}%
\pgfpathquadraticcurveto{\pgfqpoint{5.548933in}{0.530047in}}{\pgfqpoint{5.543079in}{0.528804in}}%
\pgfpathquadraticcurveto{\pgfqpoint{5.537225in}{0.527561in}}{\pgfqpoint{5.531209in}{0.527561in}}%
\pgfpathquadraticcurveto{\pgfqpoint{5.516115in}{0.527561in}}{\pgfqpoint{5.507307in}{0.536333in}}%
\pgfpathquadraticcurveto{\pgfqpoint{5.498499in}{0.545123in}}{\pgfqpoint{5.498499in}{0.560109in}}%
\pgfpathquadraticcurveto{\pgfqpoint{5.498499in}{0.575581in}}{\pgfqpoint{5.506856in}{0.584659in}}%
\pgfpathquadraticcurveto{\pgfqpoint{5.515214in}{0.593756in}}{\pgfqpoint{5.529408in}{0.593756in}}%
\pgfpathquadraticcurveto{\pgfqpoint{5.542124in}{0.593756in}}{\pgfqpoint{5.549527in}{0.585560in}}%
\pgfpathquadraticcurveto{\pgfqpoint{5.556930in}{0.577383in}}{\pgfqpoint{5.556930in}{0.563315in}}%
\pgfpathlineto{\pgfqpoint{5.556930in}{0.563315in}}%
\pgfpathclose%
\pgfpathmoveto{\pgfqpoint{5.546573in}{0.566359in}}%
\pgfpathquadraticcurveto{\pgfqpoint{5.546465in}{0.574843in}}{\pgfqpoint{5.541818in}{0.579904in}}%
\pgfpathquadraticcurveto{\pgfqpoint{5.537171in}{0.584984in}}{\pgfqpoint{5.529516in}{0.584984in}}%
\pgfpathquadraticcurveto{\pgfqpoint{5.520852in}{0.584984in}}{\pgfqpoint{5.515646in}{0.580084in}}%
\pgfpathquadraticcurveto{\pgfqpoint{5.510441in}{0.575185in}}{\pgfqpoint{5.509648in}{0.566287in}}%
\pgfpathlineto{\pgfqpoint{5.546573in}{0.566359in}}%
\pgfpathlineto{\pgfqpoint{5.546573in}{0.566359in}}%
\pgfpathclose%
\pgfpathmoveto{\pgfqpoint{5.615416in}{0.582678in}}%
\pgfpathlineto{\pgfqpoint{5.615416in}{0.616793in}}%
\pgfpathlineto{\pgfqpoint{5.625773in}{0.616793in}}%
\pgfpathlineto{\pgfqpoint{5.625773in}{0.529200in}}%
\pgfpathlineto{\pgfqpoint{5.615416in}{0.529200in}}%
\pgfpathlineto{\pgfqpoint{5.615416in}{0.538656in}}%
\pgfpathquadraticcurveto{\pgfqpoint{5.612155in}{0.533037in}}{\pgfqpoint{5.607166in}{0.530299in}}%
\pgfpathquadraticcurveto{\pgfqpoint{5.602195in}{0.527561in}}{\pgfqpoint{5.595206in}{0.527561in}}%
\pgfpathquadraticcurveto{\pgfqpoint{5.583786in}{0.527561in}}{\pgfqpoint{5.576599in}{0.536675in}}%
\pgfpathquadraticcurveto{\pgfqpoint{5.569431in}{0.545807in}}{\pgfqpoint{5.569431in}{0.560667in}}%
\pgfpathquadraticcurveto{\pgfqpoint{5.569431in}{0.575527in}}{\pgfqpoint{5.576599in}{0.584641in}}%
\pgfpathquadraticcurveto{\pgfqpoint{5.583786in}{0.593756in}}{\pgfqpoint{5.595206in}{0.593756in}}%
\pgfpathquadraticcurveto{\pgfqpoint{5.602195in}{0.593756in}}{\pgfqpoint{5.607166in}{0.591018in}}%
\pgfpathquadraticcurveto{\pgfqpoint{5.612155in}{0.588298in}}{\pgfqpoint{5.615416in}{0.582678in}}%
\pgfpathlineto{\pgfqpoint{5.615416in}{0.582678in}}%
\pgfpathclose%
\pgfpathmoveto{\pgfqpoint{5.580130in}{0.560667in}}%
\pgfpathquadraticcurveto{\pgfqpoint{5.580130in}{0.549248in}}{\pgfqpoint{5.584831in}{0.542745in}}%
\pgfpathquadraticcurveto{\pgfqpoint{5.589532in}{0.536243in}}{\pgfqpoint{5.597746in}{0.536243in}}%
\pgfpathquadraticcurveto{\pgfqpoint{5.605959in}{0.536243in}}{\pgfqpoint{5.610678in}{0.542745in}}%
\pgfpathquadraticcurveto{\pgfqpoint{5.615416in}{0.549248in}}{\pgfqpoint{5.615416in}{0.560667in}}%
\pgfpathquadraticcurveto{\pgfqpoint{5.615416in}{0.572087in}}{\pgfqpoint{5.610678in}{0.578589in}}%
\pgfpathquadraticcurveto{\pgfqpoint{5.605959in}{0.585092in}}{\pgfqpoint{5.597746in}{0.585092in}}%
\pgfpathquadraticcurveto{\pgfqpoint{5.589532in}{0.585092in}}{\pgfqpoint{5.584831in}{0.578589in}}%
\pgfpathquadraticcurveto{\pgfqpoint{5.580130in}{0.572087in}}{\pgfqpoint{5.580130in}{0.560667in}}%
\pgfpathlineto{\pgfqpoint{5.580130in}{0.560667in}}%
\pgfpathclose%
\pgfpathmoveto{\pgfqpoint{5.641893in}{0.565386in}}%
\pgfpathlineto{\pgfqpoint{5.672226in}{0.565386in}}%
\pgfpathlineto{\pgfqpoint{5.672226in}{0.556164in}}%
\pgfpathlineto{\pgfqpoint{5.641893in}{0.556164in}}%
\pgfpathlineto{\pgfqpoint{5.641893in}{0.565386in}}%
\pgfpathlineto{\pgfqpoint{5.641893in}{0.565386in}}%
\pgfpathclose%
\pgfpathmoveto{\pgfqpoint{5.730189in}{0.582678in}}%
\pgfpathlineto{\pgfqpoint{5.730189in}{0.616793in}}%
\pgfpathlineto{\pgfqpoint{5.740546in}{0.616793in}}%
\pgfpathlineto{\pgfqpoint{5.740546in}{0.529200in}}%
\pgfpathlineto{\pgfqpoint{5.730189in}{0.529200in}}%
\pgfpathlineto{\pgfqpoint{5.730189in}{0.538656in}}%
\pgfpathquadraticcurveto{\pgfqpoint{5.726929in}{0.533037in}}{\pgfqpoint{5.721939in}{0.530299in}}%
\pgfpathquadraticcurveto{\pgfqpoint{5.716968in}{0.527561in}}{\pgfqpoint{5.709979in}{0.527561in}}%
\pgfpathquadraticcurveto{\pgfqpoint{5.698560in}{0.527561in}}{\pgfqpoint{5.691373in}{0.536675in}}%
\pgfpathquadraticcurveto{\pgfqpoint{5.684204in}{0.545807in}}{\pgfqpoint{5.684204in}{0.560667in}}%
\pgfpathquadraticcurveto{\pgfqpoint{5.684204in}{0.575527in}}{\pgfqpoint{5.691373in}{0.584641in}}%
\pgfpathquadraticcurveto{\pgfqpoint{5.698560in}{0.593756in}}{\pgfqpoint{5.709979in}{0.593756in}}%
\pgfpathquadraticcurveto{\pgfqpoint{5.716968in}{0.593756in}}{\pgfqpoint{5.721939in}{0.591018in}}%
\pgfpathquadraticcurveto{\pgfqpoint{5.726929in}{0.588298in}}{\pgfqpoint{5.730189in}{0.582678in}}%
\pgfpathlineto{\pgfqpoint{5.730189in}{0.582678in}}%
\pgfpathclose%
\pgfpathmoveto{\pgfqpoint{5.694903in}{0.560667in}}%
\pgfpathquadraticcurveto{\pgfqpoint{5.694903in}{0.549248in}}{\pgfqpoint{5.699604in}{0.542745in}}%
\pgfpathquadraticcurveto{\pgfqpoint{5.704306in}{0.536243in}}{\pgfqpoint{5.712519in}{0.536243in}}%
\pgfpathquadraticcurveto{\pgfqpoint{5.720733in}{0.536243in}}{\pgfqpoint{5.725452in}{0.542745in}}%
\pgfpathquadraticcurveto{\pgfqpoint{5.730189in}{0.549248in}}{\pgfqpoint{5.730189in}{0.560667in}}%
\pgfpathquadraticcurveto{\pgfqpoint{5.730189in}{0.572087in}}{\pgfqpoint{5.725452in}{0.578589in}}%
\pgfpathquadraticcurveto{\pgfqpoint{5.720733in}{0.585092in}}{\pgfqpoint{5.712519in}{0.585092in}}%
\pgfpathquadraticcurveto{\pgfqpoint{5.704306in}{0.585092in}}{\pgfqpoint{5.699604in}{0.578589in}}%
\pgfpathquadraticcurveto{\pgfqpoint{5.694903in}{0.572087in}}{\pgfqpoint{5.694903in}{0.560667in}}%
\pgfpathlineto{\pgfqpoint{5.694903in}{0.560667in}}%
\pgfpathclose%
\pgfpathmoveto{\pgfqpoint{5.761890in}{0.592243in}}%
\pgfpathlineto{\pgfqpoint{5.772247in}{0.592243in}}%
\pgfpathlineto{\pgfqpoint{5.772247in}{0.529200in}}%
\pgfpathlineto{\pgfqpoint{5.761890in}{0.529200in}}%
\pgfpathlineto{\pgfqpoint{5.761890in}{0.592243in}}%
\pgfpathlineto{\pgfqpoint{5.761890in}{0.592243in}}%
\pgfpathclose%
\pgfpathmoveto{\pgfqpoint{5.761890in}{0.616793in}}%
\pgfpathlineto{\pgfqpoint{5.772247in}{0.616793in}}%
\pgfpathlineto{\pgfqpoint{5.772247in}{0.603662in}}%
\pgfpathlineto{\pgfqpoint{5.761890in}{0.603662in}}%
\pgfpathlineto{\pgfqpoint{5.761890in}{0.616793in}}%
\pgfpathlineto{\pgfqpoint{5.761890in}{0.616793in}}%
\pgfpathclose%
\pgfpathmoveto{\pgfqpoint{5.864668in}{0.616793in}}%
\pgfpathlineto{\pgfqpoint{5.864668in}{0.608165in}}%
\pgfpathlineto{\pgfqpoint{5.854761in}{0.608165in}}%
\pgfpathquadraticcurveto{\pgfqpoint{5.849195in}{0.608165in}}{\pgfqpoint{5.847016in}{0.605914in}}%
\pgfpathquadraticcurveto{\pgfqpoint{5.844854in}{0.603662in}}{\pgfqpoint{5.844854in}{0.597808in}}%
\pgfpathlineto{\pgfqpoint{5.844854in}{0.592243in}}%
\pgfpathlineto{\pgfqpoint{5.861912in}{0.592243in}}%
\pgfpathlineto{\pgfqpoint{5.861912in}{0.584191in}}%
\pgfpathlineto{\pgfqpoint{5.844854in}{0.584191in}}%
\pgfpathlineto{\pgfqpoint{5.844854in}{0.529200in}}%
\pgfpathlineto{\pgfqpoint{5.834443in}{0.529200in}}%
\pgfpathlineto{\pgfqpoint{5.834443in}{0.584191in}}%
\pgfpathlineto{\pgfqpoint{5.806020in}{0.584191in}}%
\pgfpathlineto{\pgfqpoint{5.806020in}{0.529200in}}%
\pgfpathlineto{\pgfqpoint{5.795609in}{0.529200in}}%
\pgfpathlineto{\pgfqpoint{5.795609in}{0.584191in}}%
\pgfpathlineto{\pgfqpoint{5.785702in}{0.584191in}}%
\pgfpathlineto{\pgfqpoint{5.785702in}{0.592243in}}%
\pgfpathlineto{\pgfqpoint{5.795609in}{0.592243in}}%
\pgfpathlineto{\pgfqpoint{5.795609in}{0.596638in}}%
\pgfpathquadraticcurveto{\pgfqpoint{5.795609in}{0.607157in}}{\pgfqpoint{5.800508in}{0.611966in}}%
\pgfpathquadraticcurveto{\pgfqpoint{5.805408in}{0.616793in}}{\pgfqpoint{5.816035in}{0.616793in}}%
\pgfpathlineto{\pgfqpoint{5.825834in}{0.616793in}}%
\pgfpathlineto{\pgfqpoint{5.825834in}{0.608165in}}%
\pgfpathlineto{\pgfqpoint{5.815927in}{0.608165in}}%
\pgfpathquadraticcurveto{\pgfqpoint{5.810361in}{0.608165in}}{\pgfqpoint{5.808164in}{0.605914in}}%
\pgfpathquadraticcurveto{\pgfqpoint{5.806020in}{0.603662in}}{\pgfqpoint{5.806020in}{0.597808in}}%
\pgfpathlineto{\pgfqpoint{5.806020in}{0.592243in}}%
\pgfpathlineto{\pgfqpoint{5.834443in}{0.592243in}}%
\pgfpathlineto{\pgfqpoint{5.834443in}{0.596638in}}%
\pgfpathquadraticcurveto{\pgfqpoint{5.834443in}{0.607157in}}{\pgfqpoint{5.839343in}{0.611966in}}%
\pgfpathquadraticcurveto{\pgfqpoint{5.844242in}{0.616793in}}{\pgfqpoint{5.854887in}{0.616793in}}%
\pgfpathlineto{\pgfqpoint{5.864668in}{0.616793in}}%
\pgfpathlineto{\pgfqpoint{5.864668in}{0.616793in}}%
\pgfpathclose%
\pgfpathmoveto{\pgfqpoint{5.927260in}{0.563315in}}%
\pgfpathlineto{\pgfqpoint{5.927260in}{0.558254in}}%
\pgfpathlineto{\pgfqpoint{5.879636in}{0.558254in}}%
\pgfpathquadraticcurveto{\pgfqpoint{5.880320in}{0.547554in}}{\pgfqpoint{5.886084in}{0.541953in}}%
\pgfpathquadraticcurveto{\pgfqpoint{5.891848in}{0.536351in}}{\pgfqpoint{5.902151in}{0.536351in}}%
\pgfpathquadraticcurveto{\pgfqpoint{5.908113in}{0.536351in}}{\pgfqpoint{5.913715in}{0.537810in}}%
\pgfpathquadraticcurveto{\pgfqpoint{5.919317in}{0.539269in}}{\pgfqpoint{5.924846in}{0.542205in}}%
\pgfpathlineto{\pgfqpoint{5.924846in}{0.532406in}}%
\pgfpathquadraticcurveto{\pgfqpoint{5.919263in}{0.530047in}}{\pgfqpoint{5.913409in}{0.528804in}}%
\pgfpathquadraticcurveto{\pgfqpoint{5.907555in}{0.527561in}}{\pgfqpoint{5.901539in}{0.527561in}}%
\pgfpathquadraticcurveto{\pgfqpoint{5.886444in}{0.527561in}}{\pgfqpoint{5.877637in}{0.536333in}}%
\pgfpathquadraticcurveto{\pgfqpoint{5.868829in}{0.545123in}}{\pgfqpoint{5.868829in}{0.560109in}}%
\pgfpathquadraticcurveto{\pgfqpoint{5.868829in}{0.575581in}}{\pgfqpoint{5.877186in}{0.584659in}}%
\pgfpathquadraticcurveto{\pgfqpoint{5.885544in}{0.593756in}}{\pgfqpoint{5.899737in}{0.593756in}}%
\pgfpathquadraticcurveto{\pgfqpoint{5.912454in}{0.593756in}}{\pgfqpoint{5.919857in}{0.585560in}}%
\pgfpathquadraticcurveto{\pgfqpoint{5.927260in}{0.577383in}}{\pgfqpoint{5.927260in}{0.563315in}}%
\pgfpathlineto{\pgfqpoint{5.927260in}{0.563315in}}%
\pgfpathclose%
\pgfpathmoveto{\pgfqpoint{5.916903in}{0.566359in}}%
\pgfpathquadraticcurveto{\pgfqpoint{5.916795in}{0.574843in}}{\pgfqpoint{5.912148in}{0.579904in}}%
\pgfpathquadraticcurveto{\pgfqpoint{5.907501in}{0.584984in}}{\pgfqpoint{5.899845in}{0.584984in}}%
\pgfpathquadraticcurveto{\pgfqpoint{5.891182in}{0.584984in}}{\pgfqpoint{5.885976in}{0.580084in}}%
\pgfpathquadraticcurveto{\pgfqpoint{5.880771in}{0.575185in}}{\pgfqpoint{5.879978in}{0.566287in}}%
\pgfpathlineto{\pgfqpoint{5.916903in}{0.566359in}}%
\pgfpathlineto{\pgfqpoint{5.916903in}{0.566359in}}%
\pgfpathclose%
\pgfpathmoveto{\pgfqpoint{5.980792in}{0.582570in}}%
\pgfpathquadraticcurveto{\pgfqpoint{5.979045in}{0.583579in}}{\pgfqpoint{5.976992in}{0.584047in}}%
\pgfpathquadraticcurveto{\pgfqpoint{5.974938in}{0.584533in}}{\pgfqpoint{5.972470in}{0.584533in}}%
\pgfpathquadraticcurveto{\pgfqpoint{5.963681in}{0.584533in}}{\pgfqpoint{5.958979in}{0.578823in}}%
\pgfpathquadraticcurveto{\pgfqpoint{5.954278in}{0.573114in}}{\pgfqpoint{5.954278in}{0.562414in}}%
\pgfpathlineto{\pgfqpoint{5.954278in}{0.529200in}}%
\pgfpathlineto{\pgfqpoint{5.943867in}{0.529200in}}%
\pgfpathlineto{\pgfqpoint{5.943867in}{0.592243in}}%
\pgfpathlineto{\pgfqpoint{5.954278in}{0.592243in}}%
\pgfpathlineto{\pgfqpoint{5.954278in}{0.582444in}}%
\pgfpathquadraticcurveto{\pgfqpoint{5.957556in}{0.588190in}}{\pgfqpoint{5.962780in}{0.590964in}}%
\pgfpathquadraticcurveto{\pgfqpoint{5.968021in}{0.593756in}}{\pgfqpoint{5.975515in}{0.593756in}}%
\pgfpathquadraticcurveto{\pgfqpoint{5.976577in}{0.593756in}}{\pgfqpoint{5.977874in}{0.593611in}}%
\pgfpathquadraticcurveto{\pgfqpoint{5.979171in}{0.593485in}}{\pgfqpoint{5.980738in}{0.593197in}}%
\pgfpathlineto{\pgfqpoint{5.980792in}{0.582570in}}%
\pgfpathlineto{\pgfqpoint{5.980792in}{0.582570in}}%
\pgfpathclose%
\pgfpathmoveto{\pgfqpoint{6.043042in}{0.563315in}}%
\pgfpathlineto{\pgfqpoint{6.043042in}{0.558254in}}%
\pgfpathlineto{\pgfqpoint{5.995418in}{0.558254in}}%
\pgfpathquadraticcurveto{\pgfqpoint{5.996102in}{0.547554in}}{\pgfqpoint{6.001866in}{0.541953in}}%
\pgfpathquadraticcurveto{\pgfqpoint{6.007630in}{0.536351in}}{\pgfqpoint{6.017933in}{0.536351in}}%
\pgfpathquadraticcurveto{\pgfqpoint{6.023895in}{0.536351in}}{\pgfqpoint{6.029497in}{0.537810in}}%
\pgfpathquadraticcurveto{\pgfqpoint{6.035099in}{0.539269in}}{\pgfqpoint{6.040628in}{0.542205in}}%
\pgfpathlineto{\pgfqpoint{6.040628in}{0.532406in}}%
\pgfpathquadraticcurveto{\pgfqpoint{6.035045in}{0.530047in}}{\pgfqpoint{6.029191in}{0.528804in}}%
\pgfpathquadraticcurveto{\pgfqpoint{6.023337in}{0.527561in}}{\pgfqpoint{6.017321in}{0.527561in}}%
\pgfpathquadraticcurveto{\pgfqpoint{6.002227in}{0.527561in}}{\pgfqpoint{5.993419in}{0.536333in}}%
\pgfpathquadraticcurveto{\pgfqpoint{5.984611in}{0.545123in}}{\pgfqpoint{5.984611in}{0.560109in}}%
\pgfpathquadraticcurveto{\pgfqpoint{5.984611in}{0.575581in}}{\pgfqpoint{5.992968in}{0.584659in}}%
\pgfpathquadraticcurveto{\pgfqpoint{6.001326in}{0.593756in}}{\pgfqpoint{6.015520in}{0.593756in}}%
\pgfpathquadraticcurveto{\pgfqpoint{6.028236in}{0.593756in}}{\pgfqpoint{6.035639in}{0.585560in}}%
\pgfpathquadraticcurveto{\pgfqpoint{6.043042in}{0.577383in}}{\pgfqpoint{6.043042in}{0.563315in}}%
\pgfpathlineto{\pgfqpoint{6.043042in}{0.563315in}}%
\pgfpathclose%
\pgfpathmoveto{\pgfqpoint{6.032685in}{0.566359in}}%
\pgfpathquadraticcurveto{\pgfqpoint{6.032577in}{0.574843in}}{\pgfqpoint{6.027930in}{0.579904in}}%
\pgfpathquadraticcurveto{\pgfqpoint{6.023283in}{0.584984in}}{\pgfqpoint{6.015628in}{0.584984in}}%
\pgfpathquadraticcurveto{\pgfqpoint{6.006964in}{0.584984in}}{\pgfqpoint{6.001758in}{0.580084in}}%
\pgfpathquadraticcurveto{\pgfqpoint{5.996553in}{0.575185in}}{\pgfqpoint{5.995760in}{0.566287in}}%
\pgfpathlineto{\pgfqpoint{6.032685in}{0.566359in}}%
\pgfpathlineto{\pgfqpoint{6.032685in}{0.566359in}}%
\pgfpathclose%
\pgfpathmoveto{\pgfqpoint{6.112461in}{0.567260in}}%
\pgfpathlineto{\pgfqpoint{6.112461in}{0.529200in}}%
\pgfpathlineto{\pgfqpoint{6.102104in}{0.529200in}}%
\pgfpathlineto{\pgfqpoint{6.102104in}{0.566917in}}%
\pgfpathquadraticcurveto{\pgfqpoint{6.102104in}{0.575869in}}{\pgfqpoint{6.098610in}{0.580300in}}%
\pgfpathquadraticcurveto{\pgfqpoint{6.095115in}{0.584749in}}{\pgfqpoint{6.088145in}{0.584749in}}%
\pgfpathquadraticcurveto{\pgfqpoint{6.079751in}{0.584749in}}{\pgfqpoint{6.074906in}{0.579400in}}%
\pgfpathquadraticcurveto{\pgfqpoint{6.070060in}{0.574068in}}{\pgfqpoint{6.070060in}{0.564828in}}%
\pgfpathlineto{\pgfqpoint{6.070060in}{0.529200in}}%
\pgfpathlineto{\pgfqpoint{6.059649in}{0.529200in}}%
\pgfpathlineto{\pgfqpoint{6.059649in}{0.592243in}}%
\pgfpathlineto{\pgfqpoint{6.070060in}{0.592243in}}%
\pgfpathlineto{\pgfqpoint{6.070060in}{0.582444in}}%
\pgfpathquadraticcurveto{\pgfqpoint{6.073789in}{0.588136in}}{\pgfqpoint{6.078814in}{0.590946in}}%
\pgfpathquadraticcurveto{\pgfqpoint{6.083858in}{0.593756in}}{\pgfqpoint{6.090450in}{0.593756in}}%
\pgfpathquadraticcurveto{\pgfqpoint{6.101311in}{0.593756in}}{\pgfqpoint{6.106877in}{0.587037in}}%
\pgfpathquadraticcurveto{\pgfqpoint{6.112461in}{0.580318in}}{\pgfqpoint{6.112461in}{0.567260in}}%
\pgfpathlineto{\pgfqpoint{6.112461in}{0.567260in}}%
\pgfpathclose%
\pgfpathmoveto{\pgfqpoint{6.178475in}{0.589829in}}%
\pgfpathlineto{\pgfqpoint{6.178475in}{0.580138in}}%
\pgfpathquadraticcurveto{\pgfqpoint{6.174081in}{0.582570in}}{\pgfqpoint{6.169668in}{0.583777in}}%
\pgfpathquadraticcurveto{\pgfqpoint{6.165255in}{0.584984in}}{\pgfqpoint{6.160752in}{0.584984in}}%
\pgfpathquadraticcurveto{\pgfqpoint{6.150665in}{0.584984in}}{\pgfqpoint{6.145081in}{0.578589in}}%
\pgfpathquadraticcurveto{\pgfqpoint{6.139515in}{0.572213in}}{\pgfqpoint{6.139515in}{0.560667in}}%
\pgfpathquadraticcurveto{\pgfqpoint{6.139515in}{0.549121in}}{\pgfqpoint{6.145081in}{0.542727in}}%
\pgfpathquadraticcurveto{\pgfqpoint{6.150665in}{0.536351in}}{\pgfqpoint{6.160752in}{0.536351in}}%
\pgfpathquadraticcurveto{\pgfqpoint{6.165255in}{0.536351in}}{\pgfqpoint{6.169668in}{0.537558in}}%
\pgfpathquadraticcurveto{\pgfqpoint{6.174081in}{0.538764in}}{\pgfqpoint{6.178475in}{0.541196in}}%
\pgfpathlineto{\pgfqpoint{6.178475in}{0.531614in}}%
\pgfpathquadraticcurveto{\pgfqpoint{6.174135in}{0.529596in}}{\pgfqpoint{6.169487in}{0.528588in}}%
\pgfpathquadraticcurveto{\pgfqpoint{6.164858in}{0.527561in}}{\pgfqpoint{6.159617in}{0.527561in}}%
\pgfpathquadraticcurveto{\pgfqpoint{6.145369in}{0.527561in}}{\pgfqpoint{6.136975in}{0.536513in}}%
\pgfpathquadraticcurveto{\pgfqpoint{6.128600in}{0.545465in}}{\pgfqpoint{6.128600in}{0.560667in}}%
\pgfpathquadraticcurveto{\pgfqpoint{6.128600in}{0.576086in}}{\pgfqpoint{6.137066in}{0.584912in}}%
\pgfpathquadraticcurveto{\pgfqpoint{6.145549in}{0.593756in}}{\pgfqpoint{6.160301in}{0.593756in}}%
\pgfpathquadraticcurveto{\pgfqpoint{6.165074in}{0.593756in}}{\pgfqpoint{6.169632in}{0.592765in}}%
\pgfpathquadraticcurveto{\pgfqpoint{6.174189in}{0.591792in}}{\pgfqpoint{6.178475in}{0.589829in}}%
\pgfpathlineto{\pgfqpoint{6.178475in}{0.589829in}}%
\pgfpathclose%
\pgfpathmoveto{\pgfqpoint{6.250416in}{0.563315in}}%
\pgfpathlineto{\pgfqpoint{6.250416in}{0.558254in}}%
\pgfpathlineto{\pgfqpoint{6.202792in}{0.558254in}}%
\pgfpathquadraticcurveto{\pgfqpoint{6.203476in}{0.547554in}}{\pgfqpoint{6.209240in}{0.541953in}}%
\pgfpathquadraticcurveto{\pgfqpoint{6.215004in}{0.536351in}}{\pgfqpoint{6.225307in}{0.536351in}}%
\pgfpathquadraticcurveto{\pgfqpoint{6.231269in}{0.536351in}}{\pgfqpoint{6.236871in}{0.537810in}}%
\pgfpathquadraticcurveto{\pgfqpoint{6.242473in}{0.539269in}}{\pgfqpoint{6.248002in}{0.542205in}}%
\pgfpathlineto{\pgfqpoint{6.248002in}{0.532406in}}%
\pgfpathquadraticcurveto{\pgfqpoint{6.242419in}{0.530047in}}{\pgfqpoint{6.236565in}{0.528804in}}%
\pgfpathquadraticcurveto{\pgfqpoint{6.230711in}{0.527561in}}{\pgfqpoint{6.224695in}{0.527561in}}%
\pgfpathquadraticcurveto{\pgfqpoint{6.209600in}{0.527561in}}{\pgfqpoint{6.200793in}{0.536333in}}%
\pgfpathquadraticcurveto{\pgfqpoint{6.191985in}{0.545123in}}{\pgfqpoint{6.191985in}{0.560109in}}%
\pgfpathquadraticcurveto{\pgfqpoint{6.191985in}{0.575581in}}{\pgfqpoint{6.200342in}{0.584659in}}%
\pgfpathquadraticcurveto{\pgfqpoint{6.208700in}{0.593756in}}{\pgfqpoint{6.222893in}{0.593756in}}%
\pgfpathquadraticcurveto{\pgfqpoint{6.235610in}{0.593756in}}{\pgfqpoint{6.243013in}{0.585560in}}%
\pgfpathquadraticcurveto{\pgfqpoint{6.250416in}{0.577383in}}{\pgfqpoint{6.250416in}{0.563315in}}%
\pgfpathlineto{\pgfqpoint{6.250416in}{0.563315in}}%
\pgfpathclose%
\pgfpathmoveto{\pgfqpoint{6.240059in}{0.566359in}}%
\pgfpathquadraticcurveto{\pgfqpoint{6.239951in}{0.574843in}}{\pgfqpoint{6.235304in}{0.579904in}}%
\pgfpathquadraticcurveto{\pgfqpoint{6.230657in}{0.584984in}}{\pgfqpoint{6.223002in}{0.584984in}}%
\pgfpathquadraticcurveto{\pgfqpoint{6.214338in}{0.584984in}}{\pgfqpoint{6.209132in}{0.580084in}}%
\pgfpathquadraticcurveto{\pgfqpoint{6.203927in}{0.575185in}}{\pgfqpoint{6.203134in}{0.566287in}}%
\pgfpathlineto{\pgfqpoint{6.240059in}{0.566359in}}%
\pgfpathlineto{\pgfqpoint{6.240059in}{0.566359in}}%
\pgfpathclose%
\pgfusepath{fill}%
\end{pgfscope}%
\begin{pgfscope}%
\pgfsetbuttcap%
\pgfsetroundjoin%
\definecolor{currentfill}{rgb}{0.309804,0.372549,0.419608}%
\pgfsetfillcolor{currentfill}%
\pgfsetlinewidth{0.000000pt}%
\definecolor{currentstroke}{rgb}{0.309804,0.372549,0.419608}%
\pgfsetstrokecolor{currentstroke}%
\pgfsetdash{}{0pt}%
\pgfpathmoveto{\pgfqpoint{6.364050in}{0.610147in}}%
\pgfpathlineto{\pgfqpoint{6.364050in}{0.592243in}}%
\pgfpathlineto{\pgfqpoint{6.385376in}{0.592243in}}%
\pgfpathlineto{\pgfqpoint{6.385376in}{0.584191in}}%
\pgfpathlineto{\pgfqpoint{6.364050in}{0.584191in}}%
\pgfpathlineto{\pgfqpoint{6.364050in}{0.549968in}}%
\pgfpathquadraticcurveto{\pgfqpoint{6.364050in}{0.542259in}}{\pgfqpoint{6.366157in}{0.540061in}}%
\pgfpathquadraticcurveto{\pgfqpoint{6.368265in}{0.537864in}}{\pgfqpoint{6.374749in}{0.537864in}}%
\pgfpathlineto{\pgfqpoint{6.385376in}{0.537864in}}%
\pgfpathlineto{\pgfqpoint{6.385376in}{0.529200in}}%
\pgfpathlineto{\pgfqpoint{6.374749in}{0.529200in}}%
\pgfpathquadraticcurveto{\pgfqpoint{6.362753in}{0.529200in}}{\pgfqpoint{6.358196in}{0.533667in}}%
\pgfpathquadraticcurveto{\pgfqpoint{6.353639in}{0.538152in}}{\pgfqpoint{6.353639in}{0.549968in}}%
\pgfpathlineto{\pgfqpoint{6.353639in}{0.584191in}}%
\pgfpathlineto{\pgfqpoint{6.346038in}{0.584191in}}%
\pgfpathlineto{\pgfqpoint{6.346038in}{0.592243in}}%
\pgfpathlineto{\pgfqpoint{6.353639in}{0.592243in}}%
\pgfpathlineto{\pgfqpoint{6.353639in}{0.610147in}}%
\pgfpathlineto{\pgfqpoint{6.364050in}{0.610147in}}%
\pgfpathlineto{\pgfqpoint{6.364050in}{0.610147in}}%
\pgfpathclose%
\pgfusepath{fill}%
\end{pgfscope}%
\begin{pgfscope}%
\pgfsetbuttcap%
\pgfsetroundjoin%
\definecolor{currentfill}{rgb}{0.309804,0.372549,0.419608}%
\pgfsetfillcolor{currentfill}%
\pgfsetlinewidth{0.000000pt}%
\definecolor{currentstroke}{rgb}{0.309804,0.372549,0.419608}%
\pgfsetstrokecolor{currentstroke}%
\pgfsetdash{}{0pt}%
\pgfpathmoveto{\pgfqpoint{6.500488in}{0.590387in}}%
\pgfpathlineto{\pgfqpoint{6.500488in}{0.580589in}}%
\pgfpathquadraticcurveto{\pgfqpoint{6.496111in}{0.582840in}}{\pgfqpoint{6.491374in}{0.583957in}}%
\pgfpathquadraticcurveto{\pgfqpoint{6.486655in}{0.585092in}}{\pgfqpoint{6.481576in}{0.585092in}}%
\pgfpathquadraticcurveto{\pgfqpoint{6.473866in}{0.585092in}}{\pgfqpoint{6.470012in}{0.582732in}}%
\pgfpathquadraticcurveto{\pgfqpoint{6.466157in}{0.580373in}}{\pgfqpoint{6.466157in}{0.575635in}}%
\pgfpathquadraticcurveto{\pgfqpoint{6.466157in}{0.572033in}}{\pgfqpoint{6.468913in}{0.569980in}}%
\pgfpathquadraticcurveto{\pgfqpoint{6.471669in}{0.567926in}}{\pgfqpoint{6.480009in}{0.566071in}}%
\pgfpathlineto{\pgfqpoint{6.483557in}{0.565278in}}%
\pgfpathquadraticcurveto{\pgfqpoint{6.494580in}{0.562919in}}{\pgfqpoint{6.499228in}{0.558614in}}%
\pgfpathquadraticcurveto{\pgfqpoint{6.503875in}{0.554309in}}{\pgfqpoint{6.503875in}{0.546600in}}%
\pgfpathquadraticcurveto{\pgfqpoint{6.503875in}{0.537810in}}{\pgfqpoint{6.496922in}{0.532676in}}%
\pgfpathquadraticcurveto{\pgfqpoint{6.489969in}{0.527561in}}{\pgfqpoint{6.477811in}{0.527561in}}%
\pgfpathquadraticcurveto{\pgfqpoint{6.472750in}{0.527561in}}{\pgfqpoint{6.467256in}{0.528552in}}%
\pgfpathquadraticcurveto{\pgfqpoint{6.461762in}{0.529542in}}{\pgfqpoint{6.455692in}{0.531506in}}%
\pgfpathlineto{\pgfqpoint{6.455692in}{0.542205in}}%
\pgfpathquadraticcurveto{\pgfqpoint{6.461438in}{0.539215in}}{\pgfqpoint{6.467004in}{0.537720in}}%
\pgfpathquadraticcurveto{\pgfqpoint{6.472570in}{0.536243in}}{\pgfqpoint{6.478045in}{0.536243in}}%
\pgfpathquadraticcurveto{\pgfqpoint{6.485358in}{0.536243in}}{\pgfqpoint{6.489285in}{0.538746in}}%
\pgfpathquadraticcurveto{\pgfqpoint{6.493230in}{0.541250in}}{\pgfqpoint{6.493230in}{0.545807in}}%
\pgfpathquadraticcurveto{\pgfqpoint{6.493230in}{0.550022in}}{\pgfqpoint{6.490384in}{0.552274in}}%
\pgfpathquadraticcurveto{\pgfqpoint{6.487556in}{0.554525in}}{\pgfqpoint{6.477919in}{0.556614in}}%
\pgfpathlineto{\pgfqpoint{6.474317in}{0.557461in}}%
\pgfpathquadraticcurveto{\pgfqpoint{6.464698in}{0.559478in}}{\pgfqpoint{6.460411in}{0.563675in}}%
\pgfpathquadraticcurveto{\pgfqpoint{6.456142in}{0.567872in}}{\pgfqpoint{6.456142in}{0.575185in}}%
\pgfpathquadraticcurveto{\pgfqpoint{6.456142in}{0.584083in}}{\pgfqpoint{6.462447in}{0.588910in}}%
\pgfpathquadraticcurveto{\pgfqpoint{6.468751in}{0.593756in}}{\pgfqpoint{6.480351in}{0.593756in}}%
\pgfpathquadraticcurveto{\pgfqpoint{6.486079in}{0.593756in}}{\pgfqpoint{6.491140in}{0.592909in}}%
\pgfpathquadraticcurveto{\pgfqpoint{6.496220in}{0.592080in}}{\pgfqpoint{6.500488in}{0.590387in}}%
\pgfpathlineto{\pgfqpoint{6.500488in}{0.590387in}}%
\pgfpathclose%
\pgfpathmoveto{\pgfqpoint{6.530605in}{0.610147in}}%
\pgfpathlineto{\pgfqpoint{6.530605in}{0.592243in}}%
\pgfpathlineto{\pgfqpoint{6.551931in}{0.592243in}}%
\pgfpathlineto{\pgfqpoint{6.551931in}{0.584191in}}%
\pgfpathlineto{\pgfqpoint{6.530605in}{0.584191in}}%
\pgfpathlineto{\pgfqpoint{6.530605in}{0.549968in}}%
\pgfpathquadraticcurveto{\pgfqpoint{6.530605in}{0.542259in}}{\pgfqpoint{6.532712in}{0.540061in}}%
\pgfpathquadraticcurveto{\pgfqpoint{6.534820in}{0.537864in}}{\pgfqpoint{6.541304in}{0.537864in}}%
\pgfpathlineto{\pgfqpoint{6.551931in}{0.537864in}}%
\pgfpathlineto{\pgfqpoint{6.551931in}{0.529200in}}%
\pgfpathlineto{\pgfqpoint{6.541304in}{0.529200in}}%
\pgfpathquadraticcurveto{\pgfqpoint{6.529308in}{0.529200in}}{\pgfqpoint{6.524751in}{0.533667in}}%
\pgfpathquadraticcurveto{\pgfqpoint{6.520194in}{0.538152in}}{\pgfqpoint{6.520194in}{0.549968in}}%
\pgfpathlineto{\pgfqpoint{6.520194in}{0.584191in}}%
\pgfpathlineto{\pgfqpoint{6.512593in}{0.584191in}}%
\pgfpathlineto{\pgfqpoint{6.512593in}{0.592243in}}%
\pgfpathlineto{\pgfqpoint{6.520194in}{0.592243in}}%
\pgfpathlineto{\pgfqpoint{6.520194in}{0.610147in}}%
\pgfpathlineto{\pgfqpoint{6.530605in}{0.610147in}}%
\pgfpathlineto{\pgfqpoint{6.530605in}{0.610147in}}%
\pgfpathclose%
\pgfpathmoveto{\pgfqpoint{6.594206in}{0.560883in}}%
\pgfpathquadraticcurveto{\pgfqpoint{6.581651in}{0.560883in}}{\pgfqpoint{6.576806in}{0.558019in}}%
\pgfpathquadraticcurveto{\pgfqpoint{6.571961in}{0.555156in}}{\pgfqpoint{6.571961in}{0.548221in}}%
\pgfpathquadraticcurveto{\pgfqpoint{6.571961in}{0.542709in}}{\pgfqpoint{6.575599in}{0.539467in}}%
\pgfpathquadraticcurveto{\pgfqpoint{6.579238in}{0.536243in}}{\pgfqpoint{6.585470in}{0.536243in}}%
\pgfpathquadraticcurveto{\pgfqpoint{6.594098in}{0.536243in}}{\pgfqpoint{6.599303in}{0.542349in}}%
\pgfpathquadraticcurveto{\pgfqpoint{6.604509in}{0.548455in}}{\pgfqpoint{6.604509in}{0.558578in}}%
\pgfpathlineto{\pgfqpoint{6.604509in}{0.560883in}}%
\pgfpathlineto{\pgfqpoint{6.594206in}{0.560883in}}%
\pgfpathlineto{\pgfqpoint{6.594206in}{0.560883in}}%
\pgfpathclose%
\pgfpathmoveto{\pgfqpoint{6.614866in}{0.565170in}}%
\pgfpathlineto{\pgfqpoint{6.614866in}{0.529200in}}%
\pgfpathlineto{\pgfqpoint{6.604509in}{0.529200in}}%
\pgfpathlineto{\pgfqpoint{6.604509in}{0.538764in}}%
\pgfpathquadraticcurveto{\pgfqpoint{6.600960in}{0.533037in}}{\pgfqpoint{6.595665in}{0.530299in}}%
\pgfpathquadraticcurveto{\pgfqpoint{6.590369in}{0.527561in}}{\pgfqpoint{6.582714in}{0.527561in}}%
\pgfpathquadraticcurveto{\pgfqpoint{6.573041in}{0.527561in}}{\pgfqpoint{6.567313in}{0.533001in}}%
\pgfpathquadraticcurveto{\pgfqpoint{6.561604in}{0.538440in}}{\pgfqpoint{6.561604in}{0.547554in}}%
\pgfpathquadraticcurveto{\pgfqpoint{6.561604in}{0.558182in}}{\pgfqpoint{6.568718in}{0.563585in}}%
\pgfpathquadraticcurveto{\pgfqpoint{6.575851in}{0.568989in}}{\pgfqpoint{6.589973in}{0.568989in}}%
\pgfpathlineto{\pgfqpoint{6.604509in}{0.568989in}}%
\pgfpathlineto{\pgfqpoint{6.604509in}{0.570016in}}%
\pgfpathquadraticcurveto{\pgfqpoint{6.604509in}{0.577166in}}{\pgfqpoint{6.599807in}{0.581075in}}%
\pgfpathquadraticcurveto{\pgfqpoint{6.595106in}{0.584984in}}{\pgfqpoint{6.586605in}{0.584984in}}%
\pgfpathquadraticcurveto{\pgfqpoint{6.581201in}{0.584984in}}{\pgfqpoint{6.576067in}{0.583687in}}%
\pgfpathquadraticcurveto{\pgfqpoint{6.570952in}{0.582390in}}{\pgfqpoint{6.566233in}{0.579796in}}%
\pgfpathlineto{\pgfqpoint{6.566233in}{0.589379in}}%
\pgfpathquadraticcurveto{\pgfqpoint{6.571907in}{0.591576in}}{\pgfqpoint{6.577256in}{0.592657in}}%
\pgfpathquadraticcurveto{\pgfqpoint{6.582606in}{0.593756in}}{\pgfqpoint{6.587667in}{0.593756in}}%
\pgfpathquadraticcurveto{\pgfqpoint{6.601356in}{0.593756in}}{\pgfqpoint{6.608111in}{0.586659in}}%
\pgfpathquadraticcurveto{\pgfqpoint{6.614866in}{0.579580in}}{\pgfqpoint{6.614866in}{0.565170in}}%
\pgfpathlineto{\pgfqpoint{6.614866in}{0.565170in}}%
\pgfpathclose%
\pgfpathmoveto{\pgfqpoint{6.646441in}{0.610147in}}%
\pgfpathlineto{\pgfqpoint{6.646441in}{0.592243in}}%
\pgfpathlineto{\pgfqpoint{6.667767in}{0.592243in}}%
\pgfpathlineto{\pgfqpoint{6.667767in}{0.584191in}}%
\pgfpathlineto{\pgfqpoint{6.646441in}{0.584191in}}%
\pgfpathlineto{\pgfqpoint{6.646441in}{0.549968in}}%
\pgfpathquadraticcurveto{\pgfqpoint{6.646441in}{0.542259in}}{\pgfqpoint{6.648548in}{0.540061in}}%
\pgfpathquadraticcurveto{\pgfqpoint{6.650656in}{0.537864in}}{\pgfqpoint{6.657140in}{0.537864in}}%
\pgfpathlineto{\pgfqpoint{6.667767in}{0.537864in}}%
\pgfpathlineto{\pgfqpoint{6.667767in}{0.529200in}}%
\pgfpathlineto{\pgfqpoint{6.657140in}{0.529200in}}%
\pgfpathquadraticcurveto{\pgfqpoint{6.645144in}{0.529200in}}{\pgfqpoint{6.640587in}{0.533667in}}%
\pgfpathquadraticcurveto{\pgfqpoint{6.636030in}{0.538152in}}{\pgfqpoint{6.636030in}{0.549968in}}%
\pgfpathlineto{\pgfqpoint{6.636030in}{0.584191in}}%
\pgfpathlineto{\pgfqpoint{6.628429in}{0.584191in}}%
\pgfpathlineto{\pgfqpoint{6.628429in}{0.592243in}}%
\pgfpathlineto{\pgfqpoint{6.636030in}{0.592243in}}%
\pgfpathlineto{\pgfqpoint{6.636030in}{0.610147in}}%
\pgfpathlineto{\pgfqpoint{6.646441in}{0.610147in}}%
\pgfpathlineto{\pgfqpoint{6.646441in}{0.610147in}}%
\pgfpathclose%
\pgfpathmoveto{\pgfqpoint{6.681384in}{0.592243in}}%
\pgfpathlineto{\pgfqpoint{6.691741in}{0.592243in}}%
\pgfpathlineto{\pgfqpoint{6.691741in}{0.529200in}}%
\pgfpathlineto{\pgfqpoint{6.681384in}{0.529200in}}%
\pgfpathlineto{\pgfqpoint{6.681384in}{0.592243in}}%
\pgfpathlineto{\pgfqpoint{6.681384in}{0.592243in}}%
\pgfpathclose%
\pgfpathmoveto{\pgfqpoint{6.681384in}{0.616793in}}%
\pgfpathlineto{\pgfqpoint{6.691741in}{0.616793in}}%
\pgfpathlineto{\pgfqpoint{6.691741in}{0.603662in}}%
\pgfpathlineto{\pgfqpoint{6.681384in}{0.603662in}}%
\pgfpathlineto{\pgfqpoint{6.681384in}{0.616793in}}%
\pgfpathlineto{\pgfqpoint{6.681384in}{0.616793in}}%
\pgfpathclose%
\pgfpathmoveto{\pgfqpoint{6.753595in}{0.590387in}}%
\pgfpathlineto{\pgfqpoint{6.753595in}{0.580589in}}%
\pgfpathquadraticcurveto{\pgfqpoint{6.749218in}{0.582840in}}{\pgfqpoint{6.744481in}{0.583957in}}%
\pgfpathquadraticcurveto{\pgfqpoint{6.739762in}{0.585092in}}{\pgfqpoint{6.734682in}{0.585092in}}%
\pgfpathquadraticcurveto{\pgfqpoint{6.726973in}{0.585092in}}{\pgfqpoint{6.723119in}{0.582732in}}%
\pgfpathquadraticcurveto{\pgfqpoint{6.719264in}{0.580373in}}{\pgfqpoint{6.719264in}{0.575635in}}%
\pgfpathquadraticcurveto{\pgfqpoint{6.719264in}{0.572033in}}{\pgfqpoint{6.722020in}{0.569980in}}%
\pgfpathquadraticcurveto{\pgfqpoint{6.724776in}{0.567926in}}{\pgfqpoint{6.733115in}{0.566071in}}%
\pgfpathlineto{\pgfqpoint{6.736664in}{0.565278in}}%
\pgfpathquadraticcurveto{\pgfqpoint{6.747687in}{0.562919in}}{\pgfqpoint{6.752334in}{0.558614in}}%
\pgfpathquadraticcurveto{\pgfqpoint{6.756981in}{0.554309in}}{\pgfqpoint{6.756981in}{0.546600in}}%
\pgfpathquadraticcurveto{\pgfqpoint{6.756981in}{0.537810in}}{\pgfqpoint{6.750029in}{0.532676in}}%
\pgfpathquadraticcurveto{\pgfqpoint{6.743076in}{0.527561in}}{\pgfqpoint{6.730918in}{0.527561in}}%
\pgfpathquadraticcurveto{\pgfqpoint{6.725856in}{0.527561in}}{\pgfqpoint{6.720363in}{0.528552in}}%
\pgfpathquadraticcurveto{\pgfqpoint{6.714869in}{0.529542in}}{\pgfqpoint{6.708799in}{0.531506in}}%
\pgfpathlineto{\pgfqpoint{6.708799in}{0.542205in}}%
\pgfpathquadraticcurveto{\pgfqpoint{6.714545in}{0.539215in}}{\pgfqpoint{6.720111in}{0.537720in}}%
\pgfpathquadraticcurveto{\pgfqpoint{6.725676in}{0.536243in}}{\pgfqpoint{6.731152in}{0.536243in}}%
\pgfpathquadraticcurveto{\pgfqpoint{6.738465in}{0.536243in}}{\pgfqpoint{6.742392in}{0.538746in}}%
\pgfpathquadraticcurveto{\pgfqpoint{6.746336in}{0.541250in}}{\pgfqpoint{6.746336in}{0.545807in}}%
\pgfpathquadraticcurveto{\pgfqpoint{6.746336in}{0.550022in}}{\pgfqpoint{6.743490in}{0.552274in}}%
\pgfpathquadraticcurveto{\pgfqpoint{6.740662in}{0.554525in}}{\pgfqpoint{6.731026in}{0.556614in}}%
\pgfpathlineto{\pgfqpoint{6.727424in}{0.557461in}}%
\pgfpathquadraticcurveto{\pgfqpoint{6.717805in}{0.559478in}}{\pgfqpoint{6.713518in}{0.563675in}}%
\pgfpathquadraticcurveto{\pgfqpoint{6.709249in}{0.567872in}}{\pgfqpoint{6.709249in}{0.575185in}}%
\pgfpathquadraticcurveto{\pgfqpoint{6.709249in}{0.584083in}}{\pgfqpoint{6.715554in}{0.588910in}}%
\pgfpathquadraticcurveto{\pgfqpoint{6.721858in}{0.593756in}}{\pgfqpoint{6.733458in}{0.593756in}}%
\pgfpathquadraticcurveto{\pgfqpoint{6.739185in}{0.593756in}}{\pgfqpoint{6.744247in}{0.592909in}}%
\pgfpathquadraticcurveto{\pgfqpoint{6.749326in}{0.592080in}}{\pgfqpoint{6.753595in}{0.590387in}}%
\pgfpathlineto{\pgfqpoint{6.753595in}{0.590387in}}%
\pgfpathclose%
\pgfpathmoveto{\pgfqpoint{6.783711in}{0.610147in}}%
\pgfpathlineto{\pgfqpoint{6.783711in}{0.592243in}}%
\pgfpathlineto{\pgfqpoint{6.805038in}{0.592243in}}%
\pgfpathlineto{\pgfqpoint{6.805038in}{0.584191in}}%
\pgfpathlineto{\pgfqpoint{6.783711in}{0.584191in}}%
\pgfpathlineto{\pgfqpoint{6.783711in}{0.549968in}}%
\pgfpathquadraticcurveto{\pgfqpoint{6.783711in}{0.542259in}}{\pgfqpoint{6.785819in}{0.540061in}}%
\pgfpathquadraticcurveto{\pgfqpoint{6.787926in}{0.537864in}}{\pgfqpoint{6.794411in}{0.537864in}}%
\pgfpathlineto{\pgfqpoint{6.805038in}{0.537864in}}%
\pgfpathlineto{\pgfqpoint{6.805038in}{0.529200in}}%
\pgfpathlineto{\pgfqpoint{6.794411in}{0.529200in}}%
\pgfpathquadraticcurveto{\pgfqpoint{6.782415in}{0.529200in}}{\pgfqpoint{6.777858in}{0.533667in}}%
\pgfpathquadraticcurveto{\pgfqpoint{6.773300in}{0.538152in}}{\pgfqpoint{6.773300in}{0.549968in}}%
\pgfpathlineto{\pgfqpoint{6.773300in}{0.584191in}}%
\pgfpathlineto{\pgfqpoint{6.765699in}{0.584191in}}%
\pgfpathlineto{\pgfqpoint{6.765699in}{0.592243in}}%
\pgfpathlineto{\pgfqpoint{6.773300in}{0.592243in}}%
\pgfpathlineto{\pgfqpoint{6.773300in}{0.610147in}}%
\pgfpathlineto{\pgfqpoint{6.783711in}{0.610147in}}%
\pgfpathlineto{\pgfqpoint{6.783711in}{0.610147in}}%
\pgfpathclose%
\pgfpathmoveto{\pgfqpoint{6.818655in}{0.592243in}}%
\pgfpathlineto{\pgfqpoint{6.829012in}{0.592243in}}%
\pgfpathlineto{\pgfqpoint{6.829012in}{0.529200in}}%
\pgfpathlineto{\pgfqpoint{6.818655in}{0.529200in}}%
\pgfpathlineto{\pgfqpoint{6.818655in}{0.592243in}}%
\pgfpathlineto{\pgfqpoint{6.818655in}{0.592243in}}%
\pgfpathclose%
\pgfpathmoveto{\pgfqpoint{6.818655in}{0.616793in}}%
\pgfpathlineto{\pgfqpoint{6.829012in}{0.616793in}}%
\pgfpathlineto{\pgfqpoint{6.829012in}{0.603662in}}%
\pgfpathlineto{\pgfqpoint{6.818655in}{0.603662in}}%
\pgfpathlineto{\pgfqpoint{6.818655in}{0.616793in}}%
\pgfpathlineto{\pgfqpoint{6.818655in}{0.616793in}}%
\pgfpathclose%
\pgfpathmoveto{\pgfqpoint{6.896053in}{0.589829in}}%
\pgfpathlineto{\pgfqpoint{6.896053in}{0.580138in}}%
\pgfpathquadraticcurveto{\pgfqpoint{6.891658in}{0.582570in}}{\pgfqpoint{6.887245in}{0.583777in}}%
\pgfpathquadraticcurveto{\pgfqpoint{6.882832in}{0.584984in}}{\pgfqpoint{6.878329in}{0.584984in}}%
\pgfpathquadraticcurveto{\pgfqpoint{6.868243in}{0.584984in}}{\pgfqpoint{6.862659in}{0.578589in}}%
\pgfpathquadraticcurveto{\pgfqpoint{6.857093in}{0.572213in}}{\pgfqpoint{6.857093in}{0.560667in}}%
\pgfpathquadraticcurveto{\pgfqpoint{6.857093in}{0.549121in}}{\pgfqpoint{6.862659in}{0.542727in}}%
\pgfpathquadraticcurveto{\pgfqpoint{6.868243in}{0.536351in}}{\pgfqpoint{6.878329in}{0.536351in}}%
\pgfpathquadraticcurveto{\pgfqpoint{6.882832in}{0.536351in}}{\pgfqpoint{6.887245in}{0.537558in}}%
\pgfpathquadraticcurveto{\pgfqpoint{6.891658in}{0.538764in}}{\pgfqpoint{6.896053in}{0.541196in}}%
\pgfpathlineto{\pgfqpoint{6.896053in}{0.531614in}}%
\pgfpathquadraticcurveto{\pgfqpoint{6.891712in}{0.529596in}}{\pgfqpoint{6.887065in}{0.528588in}}%
\pgfpathquadraticcurveto{\pgfqpoint{6.882436in}{0.527561in}}{\pgfqpoint{6.877195in}{0.527561in}}%
\pgfpathquadraticcurveto{\pgfqpoint{6.862947in}{0.527561in}}{\pgfqpoint{6.854553in}{0.536513in}}%
\pgfpathquadraticcurveto{\pgfqpoint{6.846178in}{0.545465in}}{\pgfqpoint{6.846178in}{0.560667in}}%
\pgfpathquadraticcurveto{\pgfqpoint{6.846178in}{0.576086in}}{\pgfqpoint{6.854643in}{0.584912in}}%
\pgfpathquadraticcurveto{\pgfqpoint{6.863127in}{0.593756in}}{\pgfqpoint{6.877879in}{0.593756in}}%
\pgfpathquadraticcurveto{\pgfqpoint{6.882652in}{0.593756in}}{\pgfqpoint{6.887209in}{0.592765in}}%
\pgfpathquadraticcurveto{\pgfqpoint{6.891766in}{0.591792in}}{\pgfqpoint{6.896053in}{0.589829in}}%
\pgfpathlineto{\pgfqpoint{6.896053in}{0.589829in}}%
\pgfpathclose%
\pgfpathmoveto{\pgfqpoint{6.915524in}{0.543502in}}%
\pgfpathlineto{\pgfqpoint{6.927412in}{0.543502in}}%
\pgfpathlineto{\pgfqpoint{6.927412in}{0.529200in}}%
\pgfpathlineto{\pgfqpoint{6.915524in}{0.529200in}}%
\pgfpathlineto{\pgfqpoint{6.915524in}{0.543502in}}%
\pgfpathlineto{\pgfqpoint{6.915524in}{0.543502in}}%
\pgfpathclose%
\pgfusepath{fill}%
\end{pgfscope}%
\begin{pgfscope}%
\pgfsetbuttcap%
\pgfsetroundjoin%
\definecolor{currentfill}{rgb}{0.309804,0.372549,0.419608}%
\pgfsetfillcolor{currentfill}%
\pgfsetlinewidth{0.000000pt}%
\definecolor{currentstroke}{rgb}{0.309804,0.372549,0.419608}%
\pgfsetstrokecolor{currentstroke}%
\pgfsetdash{}{0pt}%
\pgfpathmoveto{\pgfqpoint{7.090190in}{0.606760in}}%
\pgfpathlineto{\pgfqpoint{7.090190in}{0.594782in}}%
\pgfpathquadraticcurveto{\pgfqpoint{7.084444in}{0.600132in}}{\pgfqpoint{7.077942in}{0.602762in}}%
\pgfpathquadraticcurveto{\pgfqpoint{7.071440in}{0.605409in}}{\pgfqpoint{7.064127in}{0.605409in}}%
\pgfpathquadraticcurveto{\pgfqpoint{7.049717in}{0.605409in}}{\pgfqpoint{7.042062in}{0.596601in}}%
\pgfpathquadraticcurveto{\pgfqpoint{7.034407in}{0.587794in}}{\pgfqpoint{7.034407in}{0.571132in}}%
\pgfpathquadraticcurveto{\pgfqpoint{7.034407in}{0.554525in}}{\pgfqpoint{7.042062in}{0.545717in}}%
\pgfpathquadraticcurveto{\pgfqpoint{7.049717in}{0.536909in}}{\pgfqpoint{7.064127in}{0.536909in}}%
\pgfpathquadraticcurveto{\pgfqpoint{7.071440in}{0.536909in}}{\pgfqpoint{7.077942in}{0.539557in}}%
\pgfpathquadraticcurveto{\pgfqpoint{7.084444in}{0.542205in}}{\pgfqpoint{7.090190in}{0.547554in}}%
\pgfpathlineto{\pgfqpoint{7.090190in}{0.535666in}}%
\pgfpathquadraticcurveto{\pgfqpoint{7.084228in}{0.531614in}}{\pgfqpoint{7.077546in}{0.529578in}}%
\pgfpathquadraticcurveto{\pgfqpoint{7.070881in}{0.527561in}}{\pgfqpoint{7.063460in}{0.527561in}}%
\pgfpathquadraticcurveto{\pgfqpoint{7.044367in}{0.527561in}}{\pgfqpoint{7.033380in}{0.539233in}}%
\pgfpathquadraticcurveto{\pgfqpoint{7.022411in}{0.550923in}}{\pgfqpoint{7.022411in}{0.571132in}}%
\pgfpathquadraticcurveto{\pgfqpoint{7.022411in}{0.591396in}}{\pgfqpoint{7.033380in}{0.603068in}}%
\pgfpathquadraticcurveto{\pgfqpoint{7.044367in}{0.614758in}}{\pgfqpoint{7.063460in}{0.614758in}}%
\pgfpathquadraticcurveto{\pgfqpoint{7.070989in}{0.614758in}}{\pgfqpoint{7.077654in}{0.612758in}}%
\pgfpathquadraticcurveto{\pgfqpoint{7.084336in}{0.610759in}}{\pgfqpoint{7.090190in}{0.606760in}}%
\pgfpathlineto{\pgfqpoint{7.090190in}{0.606760in}}%
\pgfpathclose%
\pgfpathmoveto{\pgfqpoint{7.131726in}{0.584984in}}%
\pgfpathquadraticcurveto{\pgfqpoint{7.123405in}{0.584984in}}{\pgfqpoint{7.118559in}{0.578481in}}%
\pgfpathquadraticcurveto{\pgfqpoint{7.113714in}{0.571979in}}{\pgfqpoint{7.113714in}{0.560667in}}%
\pgfpathquadraticcurveto{\pgfqpoint{7.113714in}{0.549356in}}{\pgfqpoint{7.118523in}{0.542853in}}%
\pgfpathquadraticcurveto{\pgfqpoint{7.123351in}{0.536351in}}{\pgfqpoint{7.131726in}{0.536351in}}%
\pgfpathquadraticcurveto{\pgfqpoint{7.140012in}{0.536351in}}{\pgfqpoint{7.144839in}{0.542871in}}%
\pgfpathquadraticcurveto{\pgfqpoint{7.149684in}{0.549410in}}{\pgfqpoint{7.149684in}{0.560667in}}%
\pgfpathquadraticcurveto{\pgfqpoint{7.149684in}{0.571871in}}{\pgfqpoint{7.144839in}{0.578427in}}%
\pgfpathquadraticcurveto{\pgfqpoint{7.140012in}{0.584984in}}{\pgfqpoint{7.131726in}{0.584984in}}%
\pgfpathlineto{\pgfqpoint{7.131726in}{0.584984in}}%
\pgfpathclose%
\pgfpathmoveto{\pgfqpoint{7.131726in}{0.593756in}}%
\pgfpathquadraticcurveto{\pgfqpoint{7.145235in}{0.593756in}}{\pgfqpoint{7.152945in}{0.584966in}}%
\pgfpathquadraticcurveto{\pgfqpoint{7.160672in}{0.576194in}}{\pgfqpoint{7.160672in}{0.560667in}}%
\pgfpathquadraticcurveto{\pgfqpoint{7.160672in}{0.545195in}}{\pgfqpoint{7.152945in}{0.536369in}}%
\pgfpathquadraticcurveto{\pgfqpoint{7.145235in}{0.527561in}}{\pgfqpoint{7.131726in}{0.527561in}}%
\pgfpathquadraticcurveto{\pgfqpoint{7.118163in}{0.527561in}}{\pgfqpoint{7.110472in}{0.536369in}}%
\pgfpathquadraticcurveto{\pgfqpoint{7.102799in}{0.545195in}}{\pgfqpoint{7.102799in}{0.560667in}}%
\pgfpathquadraticcurveto{\pgfqpoint{7.102799in}{0.576194in}}{\pgfqpoint{7.110472in}{0.584966in}}%
\pgfpathquadraticcurveto{\pgfqpoint{7.118163in}{0.593756in}}{\pgfqpoint{7.131726in}{0.593756in}}%
\pgfpathlineto{\pgfqpoint{7.131726in}{0.593756in}}%
\pgfpathclose%
\pgfpathmoveto{\pgfqpoint{7.226921in}{0.580138in}}%
\pgfpathquadraticcurveto{\pgfqpoint{7.230811in}{0.587127in}}{\pgfqpoint{7.236215in}{0.590441in}}%
\pgfpathquadraticcurveto{\pgfqpoint{7.241618in}{0.593756in}}{\pgfqpoint{7.248931in}{0.593756in}}%
\pgfpathquadraticcurveto{\pgfqpoint{7.258784in}{0.593756in}}{\pgfqpoint{7.264134in}{0.586857in}}%
\pgfpathquadraticcurveto{\pgfqpoint{7.269483in}{0.579976in}}{\pgfqpoint{7.269483in}{0.567260in}}%
\pgfpathlineto{\pgfqpoint{7.269483in}{0.529200in}}%
\pgfpathlineto{\pgfqpoint{7.259072in}{0.529200in}}%
\pgfpathlineto{\pgfqpoint{7.259072in}{0.566917in}}%
\pgfpathquadraticcurveto{\pgfqpoint{7.259072in}{0.575978in}}{\pgfqpoint{7.255848in}{0.580355in}}%
\pgfpathquadraticcurveto{\pgfqpoint{7.252642in}{0.584749in}}{\pgfqpoint{7.246067in}{0.584749in}}%
\pgfpathquadraticcurveto{\pgfqpoint{7.238016in}{0.584749in}}{\pgfqpoint{7.233333in}{0.579400in}}%
\pgfpathquadraticcurveto{\pgfqpoint{7.228668in}{0.574068in}}{\pgfqpoint{7.228668in}{0.564828in}}%
\pgfpathlineto{\pgfqpoint{7.228668in}{0.529200in}}%
\pgfpathlineto{\pgfqpoint{7.218257in}{0.529200in}}%
\pgfpathlineto{\pgfqpoint{7.218257in}{0.566917in}}%
\pgfpathquadraticcurveto{\pgfqpoint{7.218257in}{0.576032in}}{\pgfqpoint{7.215051in}{0.580391in}}%
\pgfpathquadraticcurveto{\pgfqpoint{7.211844in}{0.584749in}}{\pgfqpoint{7.205144in}{0.584749in}}%
\pgfpathquadraticcurveto{\pgfqpoint{7.197200in}{0.584749in}}{\pgfqpoint{7.192517in}{0.579382in}}%
\pgfpathquadraticcurveto{\pgfqpoint{7.187852in}{0.574014in}}{\pgfqpoint{7.187852in}{0.564828in}}%
\pgfpathlineto{\pgfqpoint{7.187852in}{0.529200in}}%
\pgfpathlineto{\pgfqpoint{7.177441in}{0.529200in}}%
\pgfpathlineto{\pgfqpoint{7.177441in}{0.592243in}}%
\pgfpathlineto{\pgfqpoint{7.187852in}{0.592243in}}%
\pgfpathlineto{\pgfqpoint{7.187852in}{0.582444in}}%
\pgfpathquadraticcurveto{\pgfqpoint{7.191401in}{0.588244in}}{\pgfqpoint{7.196354in}{0.591000in}}%
\pgfpathquadraticcurveto{\pgfqpoint{7.201307in}{0.593756in}}{\pgfqpoint{7.208116in}{0.593756in}}%
\pgfpathquadraticcurveto{\pgfqpoint{7.214996in}{0.593756in}}{\pgfqpoint{7.219806in}{0.590261in}}%
\pgfpathquadraticcurveto{\pgfqpoint{7.224615in}{0.586785in}}{\pgfqpoint{7.226921in}{0.580138in}}%
\pgfpathlineto{\pgfqpoint{7.226921in}{0.580138in}}%
\pgfpathclose%
\pgfpathmoveto{\pgfqpoint{7.339208in}{0.580138in}}%
\pgfpathquadraticcurveto{\pgfqpoint{7.343099in}{0.587127in}}{\pgfqpoint{7.348503in}{0.590441in}}%
\pgfpathquadraticcurveto{\pgfqpoint{7.353906in}{0.593756in}}{\pgfqpoint{7.361219in}{0.593756in}}%
\pgfpathquadraticcurveto{\pgfqpoint{7.371072in}{0.593756in}}{\pgfqpoint{7.376421in}{0.586857in}}%
\pgfpathquadraticcurveto{\pgfqpoint{7.381771in}{0.579976in}}{\pgfqpoint{7.381771in}{0.567260in}}%
\pgfpathlineto{\pgfqpoint{7.381771in}{0.529200in}}%
\pgfpathlineto{\pgfqpoint{7.371360in}{0.529200in}}%
\pgfpathlineto{\pgfqpoint{7.371360in}{0.566917in}}%
\pgfpathquadraticcurveto{\pgfqpoint{7.371360in}{0.575978in}}{\pgfqpoint{7.368136in}{0.580355in}}%
\pgfpathquadraticcurveto{\pgfqpoint{7.364930in}{0.584749in}}{\pgfqpoint{7.358355in}{0.584749in}}%
\pgfpathquadraticcurveto{\pgfqpoint{7.350304in}{0.584749in}}{\pgfqpoint{7.345621in}{0.579400in}}%
\pgfpathquadraticcurveto{\pgfqpoint{7.340955in}{0.574068in}}{\pgfqpoint{7.340955in}{0.564828in}}%
\pgfpathlineto{\pgfqpoint{7.340955in}{0.529200in}}%
\pgfpathlineto{\pgfqpoint{7.330544in}{0.529200in}}%
\pgfpathlineto{\pgfqpoint{7.330544in}{0.566917in}}%
\pgfpathquadraticcurveto{\pgfqpoint{7.330544in}{0.576032in}}{\pgfqpoint{7.327338in}{0.580391in}}%
\pgfpathquadraticcurveto{\pgfqpoint{7.324132in}{0.584749in}}{\pgfqpoint{7.317432in}{0.584749in}}%
\pgfpathquadraticcurveto{\pgfqpoint{7.309488in}{0.584749in}}{\pgfqpoint{7.304805in}{0.579382in}}%
\pgfpathquadraticcurveto{\pgfqpoint{7.300140in}{0.574014in}}{\pgfqpoint{7.300140in}{0.564828in}}%
\pgfpathlineto{\pgfqpoint{7.300140in}{0.529200in}}%
\pgfpathlineto{\pgfqpoint{7.289729in}{0.529200in}}%
\pgfpathlineto{\pgfqpoint{7.289729in}{0.592243in}}%
\pgfpathlineto{\pgfqpoint{7.300140in}{0.592243in}}%
\pgfpathlineto{\pgfqpoint{7.300140in}{0.582444in}}%
\pgfpathquadraticcurveto{\pgfqpoint{7.303688in}{0.588244in}}{\pgfqpoint{7.308642in}{0.591000in}}%
\pgfpathquadraticcurveto{\pgfqpoint{7.313595in}{0.593756in}}{\pgfqpoint{7.320404in}{0.593756in}}%
\pgfpathquadraticcurveto{\pgfqpoint{7.327284in}{0.593756in}}{\pgfqpoint{7.332094in}{0.590261in}}%
\pgfpathquadraticcurveto{\pgfqpoint{7.336903in}{0.586785in}}{\pgfqpoint{7.339208in}{0.580138in}}%
\pgfpathlineto{\pgfqpoint{7.339208in}{0.580138in}}%
\pgfpathclose%
\pgfpathmoveto{\pgfqpoint{7.426837in}{0.584984in}}%
\pgfpathquadraticcurveto{\pgfqpoint{7.418516in}{0.584984in}}{\pgfqpoint{7.413671in}{0.578481in}}%
\pgfpathquadraticcurveto{\pgfqpoint{7.408825in}{0.571979in}}{\pgfqpoint{7.408825in}{0.560667in}}%
\pgfpathquadraticcurveto{\pgfqpoint{7.408825in}{0.549356in}}{\pgfqpoint{7.413635in}{0.542853in}}%
\pgfpathquadraticcurveto{\pgfqpoint{7.418462in}{0.536351in}}{\pgfqpoint{7.426837in}{0.536351in}}%
\pgfpathquadraticcurveto{\pgfqpoint{7.435123in}{0.536351in}}{\pgfqpoint{7.439950in}{0.542871in}}%
\pgfpathquadraticcurveto{\pgfqpoint{7.444796in}{0.549410in}}{\pgfqpoint{7.444796in}{0.560667in}}%
\pgfpathquadraticcurveto{\pgfqpoint{7.444796in}{0.571871in}}{\pgfqpoint{7.439950in}{0.578427in}}%
\pgfpathquadraticcurveto{\pgfqpoint{7.435123in}{0.584984in}}{\pgfqpoint{7.426837in}{0.584984in}}%
\pgfpathlineto{\pgfqpoint{7.426837in}{0.584984in}}%
\pgfpathclose%
\pgfpathmoveto{\pgfqpoint{7.426837in}{0.593756in}}%
\pgfpathquadraticcurveto{\pgfqpoint{7.440347in}{0.593756in}}{\pgfqpoint{7.448056in}{0.584966in}}%
\pgfpathquadraticcurveto{\pgfqpoint{7.455783in}{0.576194in}}{\pgfqpoint{7.455783in}{0.560667in}}%
\pgfpathquadraticcurveto{\pgfqpoint{7.455783in}{0.545195in}}{\pgfqpoint{7.448056in}{0.536369in}}%
\pgfpathquadraticcurveto{\pgfqpoint{7.440347in}{0.527561in}}{\pgfqpoint{7.426837in}{0.527561in}}%
\pgfpathquadraticcurveto{\pgfqpoint{7.413274in}{0.527561in}}{\pgfqpoint{7.405583in}{0.536369in}}%
\pgfpathquadraticcurveto{\pgfqpoint{7.397910in}{0.545195in}}{\pgfqpoint{7.397910in}{0.560667in}}%
\pgfpathquadraticcurveto{\pgfqpoint{7.397910in}{0.576194in}}{\pgfqpoint{7.405583in}{0.584966in}}%
\pgfpathquadraticcurveto{\pgfqpoint{7.413274in}{0.593756in}}{\pgfqpoint{7.426837in}{0.593756in}}%
\pgfpathlineto{\pgfqpoint{7.426837in}{0.593756in}}%
\pgfpathclose%
\pgfpathmoveto{\pgfqpoint{7.525364in}{0.567260in}}%
\pgfpathlineto{\pgfqpoint{7.525364in}{0.529200in}}%
\pgfpathlineto{\pgfqpoint{7.515007in}{0.529200in}}%
\pgfpathlineto{\pgfqpoint{7.515007in}{0.566917in}}%
\pgfpathquadraticcurveto{\pgfqpoint{7.515007in}{0.575869in}}{\pgfqpoint{7.511513in}{0.580300in}}%
\pgfpathquadraticcurveto{\pgfqpoint{7.508018in}{0.584749in}}{\pgfqpoint{7.501047in}{0.584749in}}%
\pgfpathquadraticcurveto{\pgfqpoint{7.492654in}{0.584749in}}{\pgfqpoint{7.487809in}{0.579400in}}%
\pgfpathquadraticcurveto{\pgfqpoint{7.482963in}{0.574068in}}{\pgfqpoint{7.482963in}{0.564828in}}%
\pgfpathlineto{\pgfqpoint{7.482963in}{0.529200in}}%
\pgfpathlineto{\pgfqpoint{7.472552in}{0.529200in}}%
\pgfpathlineto{\pgfqpoint{7.472552in}{0.592243in}}%
\pgfpathlineto{\pgfqpoint{7.482963in}{0.592243in}}%
\pgfpathlineto{\pgfqpoint{7.482963in}{0.582444in}}%
\pgfpathquadraticcurveto{\pgfqpoint{7.486692in}{0.588136in}}{\pgfqpoint{7.491717in}{0.590946in}}%
\pgfpathquadraticcurveto{\pgfqpoint{7.496761in}{0.593756in}}{\pgfqpoint{7.503353in}{0.593756in}}%
\pgfpathquadraticcurveto{\pgfqpoint{7.514214in}{0.593756in}}{\pgfqpoint{7.519780in}{0.587037in}}%
\pgfpathquadraticcurveto{\pgfqpoint{7.525364in}{0.580318in}}{\pgfqpoint{7.525364in}{0.567260in}}%
\pgfpathlineto{\pgfqpoint{7.525364in}{0.567260in}}%
\pgfpathclose%
\pgfusepath{fill}%
\end{pgfscope}%
\begin{pgfscope}%
\pgfsetbuttcap%
\pgfsetroundjoin%
\definecolor{currentfill}{rgb}{0.309804,0.372549,0.419608}%
\pgfsetfillcolor{currentfill}%
\pgfsetlinewidth{0.000000pt}%
\definecolor{currentstroke}{rgb}{0.309804,0.372549,0.419608}%
\pgfsetstrokecolor{currentstroke}%
\pgfsetdash{}{0pt}%
\pgfpathmoveto{\pgfqpoint{7.675723in}{0.567260in}}%
\pgfpathlineto{\pgfqpoint{7.675723in}{0.529200in}}%
\pgfpathlineto{\pgfqpoint{7.665366in}{0.529200in}}%
\pgfpathlineto{\pgfqpoint{7.665366in}{0.566917in}}%
\pgfpathquadraticcurveto{\pgfqpoint{7.665366in}{0.575869in}}{\pgfqpoint{7.661872in}{0.580300in}}%
\pgfpathquadraticcurveto{\pgfqpoint{7.658377in}{0.584749in}}{\pgfqpoint{7.651407in}{0.584749in}}%
\pgfpathquadraticcurveto{\pgfqpoint{7.643013in}{0.584749in}}{\pgfqpoint{7.638168in}{0.579400in}}%
\pgfpathquadraticcurveto{\pgfqpoint{7.633323in}{0.574068in}}{\pgfqpoint{7.633323in}{0.564828in}}%
\pgfpathlineto{\pgfqpoint{7.633323in}{0.529200in}}%
\pgfpathlineto{\pgfqpoint{7.622911in}{0.529200in}}%
\pgfpathlineto{\pgfqpoint{7.622911in}{0.616793in}}%
\pgfpathlineto{\pgfqpoint{7.633323in}{0.616793in}}%
\pgfpathlineto{\pgfqpoint{7.633323in}{0.582444in}}%
\pgfpathquadraticcurveto{\pgfqpoint{7.637051in}{0.588136in}}{\pgfqpoint{7.642076in}{0.590946in}}%
\pgfpathquadraticcurveto{\pgfqpoint{7.647120in}{0.593756in}}{\pgfqpoint{7.653712in}{0.593756in}}%
\pgfpathquadraticcurveto{\pgfqpoint{7.664574in}{0.593756in}}{\pgfqpoint{7.670139in}{0.587037in}}%
\pgfpathquadraticcurveto{\pgfqpoint{7.675723in}{0.580318in}}{\pgfqpoint{7.675723in}{0.567260in}}%
\pgfpathlineto{\pgfqpoint{7.675723in}{0.567260in}}%
\pgfpathclose%
\pgfpathmoveto{\pgfqpoint{7.720790in}{0.584984in}}%
\pgfpathquadraticcurveto{\pgfqpoint{7.712468in}{0.584984in}}{\pgfqpoint{7.707623in}{0.578481in}}%
\pgfpathquadraticcurveto{\pgfqpoint{7.702777in}{0.571979in}}{\pgfqpoint{7.702777in}{0.560667in}}%
\pgfpathquadraticcurveto{\pgfqpoint{7.702777in}{0.549356in}}{\pgfqpoint{7.707587in}{0.542853in}}%
\pgfpathquadraticcurveto{\pgfqpoint{7.712414in}{0.536351in}}{\pgfqpoint{7.720790in}{0.536351in}}%
\pgfpathquadraticcurveto{\pgfqpoint{7.729075in}{0.536351in}}{\pgfqpoint{7.733902in}{0.542871in}}%
\pgfpathquadraticcurveto{\pgfqpoint{7.738748in}{0.549410in}}{\pgfqpoint{7.738748in}{0.560667in}}%
\pgfpathquadraticcurveto{\pgfqpoint{7.738748in}{0.571871in}}{\pgfqpoint{7.733902in}{0.578427in}}%
\pgfpathquadraticcurveto{\pgfqpoint{7.729075in}{0.584984in}}{\pgfqpoint{7.720790in}{0.584984in}}%
\pgfpathlineto{\pgfqpoint{7.720790in}{0.584984in}}%
\pgfpathclose%
\pgfpathmoveto{\pgfqpoint{7.720790in}{0.593756in}}%
\pgfpathquadraticcurveto{\pgfqpoint{7.734299in}{0.593756in}}{\pgfqpoint{7.742008in}{0.584966in}}%
\pgfpathquadraticcurveto{\pgfqpoint{7.749735in}{0.576194in}}{\pgfqpoint{7.749735in}{0.560667in}}%
\pgfpathquadraticcurveto{\pgfqpoint{7.749735in}{0.545195in}}{\pgfqpoint{7.742008in}{0.536369in}}%
\pgfpathquadraticcurveto{\pgfqpoint{7.734299in}{0.527561in}}{\pgfqpoint{7.720790in}{0.527561in}}%
\pgfpathquadraticcurveto{\pgfqpoint{7.707226in}{0.527561in}}{\pgfqpoint{7.699535in}{0.536369in}}%
\pgfpathquadraticcurveto{\pgfqpoint{7.691862in}{0.545195in}}{\pgfqpoint{7.691862in}{0.560667in}}%
\pgfpathquadraticcurveto{\pgfqpoint{7.691862in}{0.576194in}}{\pgfqpoint{7.699535in}{0.584966in}}%
\pgfpathquadraticcurveto{\pgfqpoint{7.707226in}{0.593756in}}{\pgfqpoint{7.720790in}{0.593756in}}%
\pgfpathlineto{\pgfqpoint{7.720790in}{0.593756in}}%
\pgfpathclose%
\pgfpathmoveto{\pgfqpoint{7.803429in}{0.582570in}}%
\pgfpathquadraticcurveto{\pgfqpoint{7.801682in}{0.583579in}}{\pgfqpoint{7.799629in}{0.584047in}}%
\pgfpathquadraticcurveto{\pgfqpoint{7.797575in}{0.584533in}}{\pgfqpoint{7.795108in}{0.584533in}}%
\pgfpathquadraticcurveto{\pgfqpoint{7.786318in}{0.584533in}}{\pgfqpoint{7.781617in}{0.578823in}}%
\pgfpathquadraticcurveto{\pgfqpoint{7.776915in}{0.573114in}}{\pgfqpoint{7.776915in}{0.562414in}}%
\pgfpathlineto{\pgfqpoint{7.776915in}{0.529200in}}%
\pgfpathlineto{\pgfqpoint{7.766504in}{0.529200in}}%
\pgfpathlineto{\pgfqpoint{7.766504in}{0.592243in}}%
\pgfpathlineto{\pgfqpoint{7.776915in}{0.592243in}}%
\pgfpathlineto{\pgfqpoint{7.776915in}{0.582444in}}%
\pgfpathquadraticcurveto{\pgfqpoint{7.780194in}{0.588190in}}{\pgfqpoint{7.785417in}{0.590964in}}%
\pgfpathquadraticcurveto{\pgfqpoint{7.790659in}{0.593756in}}{\pgfqpoint{7.798152in}{0.593756in}}%
\pgfpathquadraticcurveto{\pgfqpoint{7.799214in}{0.593756in}}{\pgfqpoint{7.800511in}{0.593611in}}%
\pgfpathquadraticcurveto{\pgfqpoint{7.801808in}{0.593485in}}{\pgfqpoint{7.803375in}{0.593197in}}%
\pgfpathlineto{\pgfqpoint{7.803429in}{0.582570in}}%
\pgfpathlineto{\pgfqpoint{7.803429in}{0.582570in}}%
\pgfpathclose%
\pgfpathmoveto{\pgfqpoint{7.814291in}{0.592243in}}%
\pgfpathlineto{\pgfqpoint{7.824648in}{0.592243in}}%
\pgfpathlineto{\pgfqpoint{7.824648in}{0.529200in}}%
\pgfpathlineto{\pgfqpoint{7.814291in}{0.529200in}}%
\pgfpathlineto{\pgfqpoint{7.814291in}{0.592243in}}%
\pgfpathlineto{\pgfqpoint{7.814291in}{0.592243in}}%
\pgfpathclose%
\pgfpathmoveto{\pgfqpoint{7.814291in}{0.616793in}}%
\pgfpathlineto{\pgfqpoint{7.824648in}{0.616793in}}%
\pgfpathlineto{\pgfqpoint{7.824648in}{0.603662in}}%
\pgfpathlineto{\pgfqpoint{7.814291in}{0.603662in}}%
\pgfpathlineto{\pgfqpoint{7.814291in}{0.616793in}}%
\pgfpathlineto{\pgfqpoint{7.814291in}{0.616793in}}%
\pgfpathclose%
\pgfpathmoveto{\pgfqpoint{7.841813in}{0.592243in}}%
\pgfpathlineto{\pgfqpoint{7.891004in}{0.592243in}}%
\pgfpathlineto{\pgfqpoint{7.891004in}{0.582786in}}%
\pgfpathlineto{\pgfqpoint{7.852062in}{0.537468in}}%
\pgfpathlineto{\pgfqpoint{7.891004in}{0.537468in}}%
\pgfpathlineto{\pgfqpoint{7.891004in}{0.529200in}}%
\pgfpathlineto{\pgfqpoint{7.840408in}{0.529200in}}%
\pgfpathlineto{\pgfqpoint{7.840408in}{0.538656in}}%
\pgfpathlineto{\pgfqpoint{7.879369in}{0.583975in}}%
\pgfpathlineto{\pgfqpoint{7.841813in}{0.583975in}}%
\pgfpathlineto{\pgfqpoint{7.841813in}{0.592243in}}%
\pgfpathlineto{\pgfqpoint{7.841813in}{0.592243in}}%
\pgfpathclose%
\pgfpathmoveto{\pgfqpoint{7.931244in}{0.584984in}}%
\pgfpathquadraticcurveto{\pgfqpoint{7.922922in}{0.584984in}}{\pgfqpoint{7.918077in}{0.578481in}}%
\pgfpathquadraticcurveto{\pgfqpoint{7.913231in}{0.571979in}}{\pgfqpoint{7.913231in}{0.560667in}}%
\pgfpathquadraticcurveto{\pgfqpoint{7.913231in}{0.549356in}}{\pgfqpoint{7.918041in}{0.542853in}}%
\pgfpathquadraticcurveto{\pgfqpoint{7.922868in}{0.536351in}}{\pgfqpoint{7.931244in}{0.536351in}}%
\pgfpathquadraticcurveto{\pgfqpoint{7.939529in}{0.536351in}}{\pgfqpoint{7.944356in}{0.542871in}}%
\pgfpathquadraticcurveto{\pgfqpoint{7.949202in}{0.549410in}}{\pgfqpoint{7.949202in}{0.560667in}}%
\pgfpathquadraticcurveto{\pgfqpoint{7.949202in}{0.571871in}}{\pgfqpoint{7.944356in}{0.578427in}}%
\pgfpathquadraticcurveto{\pgfqpoint{7.939529in}{0.584984in}}{\pgfqpoint{7.931244in}{0.584984in}}%
\pgfpathlineto{\pgfqpoint{7.931244in}{0.584984in}}%
\pgfpathclose%
\pgfpathmoveto{\pgfqpoint{7.931244in}{0.593756in}}%
\pgfpathquadraticcurveto{\pgfqpoint{7.944753in}{0.593756in}}{\pgfqpoint{7.952462in}{0.584966in}}%
\pgfpathquadraticcurveto{\pgfqpoint{7.960189in}{0.576194in}}{\pgfqpoint{7.960189in}{0.560667in}}%
\pgfpathquadraticcurveto{\pgfqpoint{7.960189in}{0.545195in}}{\pgfqpoint{7.952462in}{0.536369in}}%
\pgfpathquadraticcurveto{\pgfqpoint{7.944753in}{0.527561in}}{\pgfqpoint{7.931244in}{0.527561in}}%
\pgfpathquadraticcurveto{\pgfqpoint{7.917680in}{0.527561in}}{\pgfqpoint{7.909989in}{0.536369in}}%
\pgfpathquadraticcurveto{\pgfqpoint{7.902316in}{0.545195in}}{\pgfqpoint{7.902316in}{0.560667in}}%
\pgfpathquadraticcurveto{\pgfqpoint{7.902316in}{0.576194in}}{\pgfqpoint{7.909989in}{0.584966in}}%
\pgfpathquadraticcurveto{\pgfqpoint{7.917680in}{0.593756in}}{\pgfqpoint{7.931244in}{0.593756in}}%
\pgfpathlineto{\pgfqpoint{7.931244in}{0.593756in}}%
\pgfpathclose%
\pgfpathmoveto{\pgfqpoint{8.029770in}{0.567260in}}%
\pgfpathlineto{\pgfqpoint{8.029770in}{0.529200in}}%
\pgfpathlineto{\pgfqpoint{8.019413in}{0.529200in}}%
\pgfpathlineto{\pgfqpoint{8.019413in}{0.566917in}}%
\pgfpathquadraticcurveto{\pgfqpoint{8.019413in}{0.575869in}}{\pgfqpoint{8.015919in}{0.580300in}}%
\pgfpathquadraticcurveto{\pgfqpoint{8.012424in}{0.584749in}}{\pgfqpoint{8.005454in}{0.584749in}}%
\pgfpathquadraticcurveto{\pgfqpoint{7.997060in}{0.584749in}}{\pgfqpoint{7.992215in}{0.579400in}}%
\pgfpathquadraticcurveto{\pgfqpoint{7.987369in}{0.574068in}}{\pgfqpoint{7.987369in}{0.564828in}}%
\pgfpathlineto{\pgfqpoint{7.987369in}{0.529200in}}%
\pgfpathlineto{\pgfqpoint{7.976958in}{0.529200in}}%
\pgfpathlineto{\pgfqpoint{7.976958in}{0.592243in}}%
\pgfpathlineto{\pgfqpoint{7.987369in}{0.592243in}}%
\pgfpathlineto{\pgfqpoint{7.987369in}{0.582444in}}%
\pgfpathquadraticcurveto{\pgfqpoint{7.991098in}{0.588136in}}{\pgfqpoint{7.996123in}{0.590946in}}%
\pgfpathquadraticcurveto{\pgfqpoint{8.001167in}{0.593756in}}{\pgfqpoint{8.007759in}{0.593756in}}%
\pgfpathquadraticcurveto{\pgfqpoint{8.018620in}{0.593756in}}{\pgfqpoint{8.024186in}{0.587037in}}%
\pgfpathquadraticcurveto{\pgfqpoint{8.029770in}{0.580318in}}{\pgfqpoint{8.029770in}{0.567260in}}%
\pgfpathlineto{\pgfqpoint{8.029770in}{0.567260in}}%
\pgfpathclose%
\pgfpathmoveto{\pgfqpoint{8.090597in}{0.590387in}}%
\pgfpathlineto{\pgfqpoint{8.090597in}{0.580589in}}%
\pgfpathquadraticcurveto{\pgfqpoint{8.086220in}{0.582840in}}{\pgfqpoint{8.081483in}{0.583957in}}%
\pgfpathquadraticcurveto{\pgfqpoint{8.076764in}{0.585092in}}{\pgfqpoint{8.071684in}{0.585092in}}%
\pgfpathquadraticcurveto{\pgfqpoint{8.063975in}{0.585092in}}{\pgfqpoint{8.060120in}{0.582732in}}%
\pgfpathquadraticcurveto{\pgfqpoint{8.056266in}{0.580373in}}{\pgfqpoint{8.056266in}{0.575635in}}%
\pgfpathquadraticcurveto{\pgfqpoint{8.056266in}{0.572033in}}{\pgfqpoint{8.059022in}{0.569980in}}%
\pgfpathquadraticcurveto{\pgfqpoint{8.061778in}{0.567926in}}{\pgfqpoint{8.070117in}{0.566071in}}%
\pgfpathlineto{\pgfqpoint{8.073666in}{0.565278in}}%
\pgfpathquadraticcurveto{\pgfqpoint{8.084689in}{0.562919in}}{\pgfqpoint{8.089336in}{0.558614in}}%
\pgfpathquadraticcurveto{\pgfqpoint{8.093983in}{0.554309in}}{\pgfqpoint{8.093983in}{0.546600in}}%
\pgfpathquadraticcurveto{\pgfqpoint{8.093983in}{0.537810in}}{\pgfqpoint{8.087031in}{0.532676in}}%
\pgfpathquadraticcurveto{\pgfqpoint{8.080078in}{0.527561in}}{\pgfqpoint{8.067920in}{0.527561in}}%
\pgfpathquadraticcurveto{\pgfqpoint{8.062858in}{0.527561in}}{\pgfqpoint{8.057365in}{0.528552in}}%
\pgfpathquadraticcurveto{\pgfqpoint{8.051871in}{0.529542in}}{\pgfqpoint{8.045801in}{0.531506in}}%
\pgfpathlineto{\pgfqpoint{8.045801in}{0.542205in}}%
\pgfpathquadraticcurveto{\pgfqpoint{8.051547in}{0.539215in}}{\pgfqpoint{8.057112in}{0.537720in}}%
\pgfpathquadraticcurveto{\pgfqpoint{8.062678in}{0.536243in}}{\pgfqpoint{8.068154in}{0.536243in}}%
\pgfpathquadraticcurveto{\pgfqpoint{8.075467in}{0.536243in}}{\pgfqpoint{8.079393in}{0.538746in}}%
\pgfpathquadraticcurveto{\pgfqpoint{8.083338in}{0.541250in}}{\pgfqpoint{8.083338in}{0.545807in}}%
\pgfpathquadraticcurveto{\pgfqpoint{8.083338in}{0.550022in}}{\pgfqpoint{8.080492in}{0.552274in}}%
\pgfpathquadraticcurveto{\pgfqpoint{8.077664in}{0.554525in}}{\pgfqpoint{8.068028in}{0.556614in}}%
\pgfpathlineto{\pgfqpoint{8.064425in}{0.557461in}}%
\pgfpathquadraticcurveto{\pgfqpoint{8.054807in}{0.559478in}}{\pgfqpoint{8.050520in}{0.563675in}}%
\pgfpathquadraticcurveto{\pgfqpoint{8.046251in}{0.567872in}}{\pgfqpoint{8.046251in}{0.575185in}}%
\pgfpathquadraticcurveto{\pgfqpoint{8.046251in}{0.584083in}}{\pgfqpoint{8.052555in}{0.588910in}}%
\pgfpathquadraticcurveto{\pgfqpoint{8.058860in}{0.593756in}}{\pgfqpoint{8.070459in}{0.593756in}}%
\pgfpathquadraticcurveto{\pgfqpoint{8.076187in}{0.593756in}}{\pgfqpoint{8.081249in}{0.592909in}}%
\pgfpathquadraticcurveto{\pgfqpoint{8.086328in}{0.592080in}}{\pgfqpoint{8.090597in}{0.590387in}}%
\pgfpathlineto{\pgfqpoint{8.090597in}{0.590387in}}%
\pgfpathclose%
\pgfusepath{fill}%
\end{pgfscope}%
\begin{pgfscope}%
\pgfsetbuttcap%
\pgfsetroundjoin%
\definecolor{currentfill}{rgb}{0.309804,0.372549,0.419608}%
\pgfsetfillcolor{currentfill}%
\pgfsetlinewidth{0.000000pt}%
\definecolor{currentstroke}{rgb}{0.309804,0.372549,0.419608}%
\pgfsetstrokecolor{currentstroke}%
\pgfsetdash{}{0pt}%
\pgfpathmoveto{\pgfqpoint{8.208467in}{0.560883in}}%
\pgfpathquadraticcurveto{\pgfqpoint{8.195913in}{0.560883in}}{\pgfqpoint{8.191068in}{0.558019in}}%
\pgfpathquadraticcurveto{\pgfqpoint{8.186222in}{0.555156in}}{\pgfqpoint{8.186222in}{0.548221in}}%
\pgfpathquadraticcurveto{\pgfqpoint{8.186222in}{0.542709in}}{\pgfqpoint{8.189861in}{0.539467in}}%
\pgfpathquadraticcurveto{\pgfqpoint{8.193499in}{0.536243in}}{\pgfqpoint{8.199731in}{0.536243in}}%
\pgfpathquadraticcurveto{\pgfqpoint{8.208359in}{0.536243in}}{\pgfqpoint{8.213565in}{0.542349in}}%
\pgfpathquadraticcurveto{\pgfqpoint{8.218770in}{0.548455in}}{\pgfqpoint{8.218770in}{0.558578in}}%
\pgfpathlineto{\pgfqpoint{8.218770in}{0.560883in}}%
\pgfpathlineto{\pgfqpoint{8.208467in}{0.560883in}}%
\pgfpathlineto{\pgfqpoint{8.208467in}{0.560883in}}%
\pgfpathclose%
\pgfpathmoveto{\pgfqpoint{8.229127in}{0.565170in}}%
\pgfpathlineto{\pgfqpoint{8.229127in}{0.529200in}}%
\pgfpathlineto{\pgfqpoint{8.218770in}{0.529200in}}%
\pgfpathlineto{\pgfqpoint{8.218770in}{0.538764in}}%
\pgfpathquadraticcurveto{\pgfqpoint{8.215222in}{0.533037in}}{\pgfqpoint{8.209926in}{0.530299in}}%
\pgfpathquadraticcurveto{\pgfqpoint{8.204631in}{0.527561in}}{\pgfqpoint{8.196976in}{0.527561in}}%
\pgfpathquadraticcurveto{\pgfqpoint{8.187303in}{0.527561in}}{\pgfqpoint{8.181575in}{0.533001in}}%
\pgfpathquadraticcurveto{\pgfqpoint{8.175865in}{0.538440in}}{\pgfqpoint{8.175865in}{0.547554in}}%
\pgfpathquadraticcurveto{\pgfqpoint{8.175865in}{0.558182in}}{\pgfqpoint{8.182980in}{0.563585in}}%
\pgfpathquadraticcurveto{\pgfqpoint{8.190113in}{0.568989in}}{\pgfqpoint{8.204234in}{0.568989in}}%
\pgfpathlineto{\pgfqpoint{8.218770in}{0.568989in}}%
\pgfpathlineto{\pgfqpoint{8.218770in}{0.570016in}}%
\pgfpathquadraticcurveto{\pgfqpoint{8.218770in}{0.577166in}}{\pgfqpoint{8.214069in}{0.581075in}}%
\pgfpathquadraticcurveto{\pgfqpoint{8.209368in}{0.584984in}}{\pgfqpoint{8.200866in}{0.584984in}}%
\pgfpathquadraticcurveto{\pgfqpoint{8.195463in}{0.584984in}}{\pgfqpoint{8.190329in}{0.583687in}}%
\pgfpathquadraticcurveto{\pgfqpoint{8.185214in}{0.582390in}}{\pgfqpoint{8.180494in}{0.579796in}}%
\pgfpathlineto{\pgfqpoint{8.180494in}{0.589379in}}%
\pgfpathquadraticcurveto{\pgfqpoint{8.186168in}{0.591576in}}{\pgfqpoint{8.191518in}{0.592657in}}%
\pgfpathquadraticcurveto{\pgfqpoint{8.196867in}{0.593756in}}{\pgfqpoint{8.201929in}{0.593756in}}%
\pgfpathquadraticcurveto{\pgfqpoint{8.215618in}{0.593756in}}{\pgfqpoint{8.222373in}{0.586659in}}%
\pgfpathquadraticcurveto{\pgfqpoint{8.229127in}{0.579580in}}{\pgfqpoint{8.229127in}{0.565170in}}%
\pgfpathlineto{\pgfqpoint{8.229127in}{0.565170in}}%
\pgfpathclose%
\pgfpathmoveto{\pgfqpoint{8.286982in}{0.582570in}}%
\pgfpathquadraticcurveto{\pgfqpoint{8.285235in}{0.583579in}}{\pgfqpoint{8.283182in}{0.584047in}}%
\pgfpathquadraticcurveto{\pgfqpoint{8.281128in}{0.584533in}}{\pgfqpoint{8.278661in}{0.584533in}}%
\pgfpathquadraticcurveto{\pgfqpoint{8.269871in}{0.584533in}}{\pgfqpoint{8.265170in}{0.578823in}}%
\pgfpathquadraticcurveto{\pgfqpoint{8.260468in}{0.573114in}}{\pgfqpoint{8.260468in}{0.562414in}}%
\pgfpathlineto{\pgfqpoint{8.260468in}{0.529200in}}%
\pgfpathlineto{\pgfqpoint{8.250057in}{0.529200in}}%
\pgfpathlineto{\pgfqpoint{8.250057in}{0.592243in}}%
\pgfpathlineto{\pgfqpoint{8.260468in}{0.592243in}}%
\pgfpathlineto{\pgfqpoint{8.260468in}{0.582444in}}%
\pgfpathquadraticcurveto{\pgfqpoint{8.263747in}{0.588190in}}{\pgfqpoint{8.268970in}{0.590964in}}%
\pgfpathquadraticcurveto{\pgfqpoint{8.274212in}{0.593756in}}{\pgfqpoint{8.281705in}{0.593756in}}%
\pgfpathquadraticcurveto{\pgfqpoint{8.282767in}{0.593756in}}{\pgfqpoint{8.284064in}{0.593611in}}%
\pgfpathquadraticcurveto{\pgfqpoint{8.285361in}{0.593485in}}{\pgfqpoint{8.286928in}{0.593197in}}%
\pgfpathlineto{\pgfqpoint{8.286982in}{0.582570in}}%
\pgfpathlineto{\pgfqpoint{8.286982in}{0.582570in}}%
\pgfpathclose%
\pgfpathmoveto{\pgfqpoint{8.349232in}{0.563315in}}%
\pgfpathlineto{\pgfqpoint{8.349232in}{0.558254in}}%
\pgfpathlineto{\pgfqpoint{8.301608in}{0.558254in}}%
\pgfpathquadraticcurveto{\pgfqpoint{8.302293in}{0.547554in}}{\pgfqpoint{8.308057in}{0.541953in}}%
\pgfpathquadraticcurveto{\pgfqpoint{8.313820in}{0.536351in}}{\pgfqpoint{8.324123in}{0.536351in}}%
\pgfpathquadraticcurveto{\pgfqpoint{8.330085in}{0.536351in}}{\pgfqpoint{8.335687in}{0.537810in}}%
\pgfpathquadraticcurveto{\pgfqpoint{8.341289in}{0.539269in}}{\pgfqpoint{8.346819in}{0.542205in}}%
\pgfpathlineto{\pgfqpoint{8.346819in}{0.532406in}}%
\pgfpathquadraticcurveto{\pgfqpoint{8.341235in}{0.530047in}}{\pgfqpoint{8.335381in}{0.528804in}}%
\pgfpathquadraticcurveto{\pgfqpoint{8.329527in}{0.527561in}}{\pgfqpoint{8.323511in}{0.527561in}}%
\pgfpathquadraticcurveto{\pgfqpoint{8.308417in}{0.527561in}}{\pgfqpoint{8.299609in}{0.536333in}}%
\pgfpathquadraticcurveto{\pgfqpoint{8.290801in}{0.545123in}}{\pgfqpoint{8.290801in}{0.560109in}}%
\pgfpathquadraticcurveto{\pgfqpoint{8.290801in}{0.575581in}}{\pgfqpoint{8.299159in}{0.584659in}}%
\pgfpathquadraticcurveto{\pgfqpoint{8.307516in}{0.593756in}}{\pgfqpoint{8.321710in}{0.593756in}}%
\pgfpathquadraticcurveto{\pgfqpoint{8.334426in}{0.593756in}}{\pgfqpoint{8.341829in}{0.585560in}}%
\pgfpathquadraticcurveto{\pgfqpoint{8.349232in}{0.577383in}}{\pgfqpoint{8.349232in}{0.563315in}}%
\pgfpathlineto{\pgfqpoint{8.349232in}{0.563315in}}%
\pgfpathclose%
\pgfpathmoveto{\pgfqpoint{8.338875in}{0.566359in}}%
\pgfpathquadraticcurveto{\pgfqpoint{8.338767in}{0.574843in}}{\pgfqpoint{8.334120in}{0.579904in}}%
\pgfpathquadraticcurveto{\pgfqpoint{8.329473in}{0.584984in}}{\pgfqpoint{8.321818in}{0.584984in}}%
\pgfpathquadraticcurveto{\pgfqpoint{8.313154in}{0.584984in}}{\pgfqpoint{8.307948in}{0.580084in}}%
\pgfpathquadraticcurveto{\pgfqpoint{8.302743in}{0.575185in}}{\pgfqpoint{8.301950in}{0.566287in}}%
\pgfpathlineto{\pgfqpoint{8.338875in}{0.566359in}}%
\pgfpathlineto{\pgfqpoint{8.338875in}{0.566359in}}%
\pgfpathclose%
\pgfusepath{fill}%
\end{pgfscope}%
\begin{pgfscope}%
\pgfsetbuttcap%
\pgfsetroundjoin%
\definecolor{currentfill}{rgb}{0.309804,0.372549,0.419608}%
\pgfsetfillcolor{currentfill}%
\pgfsetlinewidth{0.000000pt}%
\definecolor{currentstroke}{rgb}{0.309804,0.372549,0.419608}%
\pgfsetstrokecolor{currentstroke}%
\pgfsetdash{}{0pt}%
\pgfpathmoveto{\pgfqpoint{8.422750in}{0.554075in}}%
\pgfpathlineto{\pgfqpoint{8.422750in}{0.592243in}}%
\pgfpathlineto{\pgfqpoint{8.433106in}{0.592243in}}%
\pgfpathlineto{\pgfqpoint{8.433106in}{0.554471in}}%
\pgfpathquadraticcurveto{\pgfqpoint{8.433106in}{0.545519in}}{\pgfqpoint{8.436583in}{0.541034in}}%
\pgfpathquadraticcurveto{\pgfqpoint{8.440077in}{0.536567in}}{\pgfqpoint{8.447066in}{0.536567in}}%
\pgfpathquadraticcurveto{\pgfqpoint{8.455442in}{0.536567in}}{\pgfqpoint{8.460305in}{0.541917in}}%
\pgfpathquadraticcurveto{\pgfqpoint{8.465186in}{0.547266in}}{\pgfqpoint{8.465186in}{0.556506in}}%
\pgfpathlineto{\pgfqpoint{8.465186in}{0.592243in}}%
\pgfpathlineto{\pgfqpoint{8.475543in}{0.592243in}}%
\pgfpathlineto{\pgfqpoint{8.475543in}{0.529200in}}%
\pgfpathlineto{\pgfqpoint{8.465186in}{0.529200in}}%
\pgfpathlineto{\pgfqpoint{8.465186in}{0.538891in}}%
\pgfpathquadraticcurveto{\pgfqpoint{8.461422in}{0.533145in}}{\pgfqpoint{8.456432in}{0.530353in}}%
\pgfpathquadraticcurveto{\pgfqpoint{8.451461in}{0.527561in}}{\pgfqpoint{8.444868in}{0.527561in}}%
\pgfpathquadraticcurveto{\pgfqpoint{8.434007in}{0.527561in}}{\pgfqpoint{8.428369in}{0.534315in}}%
\pgfpathquadraticcurveto{\pgfqpoint{8.422750in}{0.541070in}}{\pgfqpoint{8.422750in}{0.554075in}}%
\pgfpathlineto{\pgfqpoint{8.422750in}{0.554075in}}%
\pgfpathclose%
\pgfpathmoveto{\pgfqpoint{8.448813in}{0.593756in}}%
\pgfpathlineto{\pgfqpoint{8.448813in}{0.593756in}}%
\pgfpathlineto{\pgfqpoint{8.448813in}{0.593756in}}%
\pgfpathclose%
\pgfpathmoveto{\pgfqpoint{8.549285in}{0.567260in}}%
\pgfpathlineto{\pgfqpoint{8.549285in}{0.529200in}}%
\pgfpathlineto{\pgfqpoint{8.538928in}{0.529200in}}%
\pgfpathlineto{\pgfqpoint{8.538928in}{0.566917in}}%
\pgfpathquadraticcurveto{\pgfqpoint{8.538928in}{0.575869in}}{\pgfqpoint{8.535434in}{0.580300in}}%
\pgfpathquadraticcurveto{\pgfqpoint{8.531939in}{0.584749in}}{\pgfqpoint{8.524968in}{0.584749in}}%
\pgfpathquadraticcurveto{\pgfqpoint{8.516575in}{0.584749in}}{\pgfqpoint{8.511730in}{0.579400in}}%
\pgfpathquadraticcurveto{\pgfqpoint{8.506884in}{0.574068in}}{\pgfqpoint{8.506884in}{0.564828in}}%
\pgfpathlineto{\pgfqpoint{8.506884in}{0.529200in}}%
\pgfpathlineto{\pgfqpoint{8.496473in}{0.529200in}}%
\pgfpathlineto{\pgfqpoint{8.496473in}{0.592243in}}%
\pgfpathlineto{\pgfqpoint{8.506884in}{0.592243in}}%
\pgfpathlineto{\pgfqpoint{8.506884in}{0.582444in}}%
\pgfpathquadraticcurveto{\pgfqpoint{8.510613in}{0.588136in}}{\pgfqpoint{8.515638in}{0.590946in}}%
\pgfpathquadraticcurveto{\pgfqpoint{8.520682in}{0.593756in}}{\pgfqpoint{8.527274in}{0.593756in}}%
\pgfpathquadraticcurveto{\pgfqpoint{8.538135in}{0.593756in}}{\pgfqpoint{8.543701in}{0.587037in}}%
\pgfpathquadraticcurveto{\pgfqpoint{8.549285in}{0.580318in}}{\pgfqpoint{8.549285in}{0.567260in}}%
\pgfpathlineto{\pgfqpoint{8.549285in}{0.567260in}}%
\pgfpathclose%
\pgfpathmoveto{\pgfqpoint{8.615299in}{0.589829in}}%
\pgfpathlineto{\pgfqpoint{8.615299in}{0.580138in}}%
\pgfpathquadraticcurveto{\pgfqpoint{8.610904in}{0.582570in}}{\pgfqpoint{8.606491in}{0.583777in}}%
\pgfpathquadraticcurveto{\pgfqpoint{8.602078in}{0.584984in}}{\pgfqpoint{8.597575in}{0.584984in}}%
\pgfpathquadraticcurveto{\pgfqpoint{8.587489in}{0.584984in}}{\pgfqpoint{8.581905in}{0.578589in}}%
\pgfpathquadraticcurveto{\pgfqpoint{8.576339in}{0.572213in}}{\pgfqpoint{8.576339in}{0.560667in}}%
\pgfpathquadraticcurveto{\pgfqpoint{8.576339in}{0.549121in}}{\pgfqpoint{8.581905in}{0.542727in}}%
\pgfpathquadraticcurveto{\pgfqpoint{8.587489in}{0.536351in}}{\pgfqpoint{8.597575in}{0.536351in}}%
\pgfpathquadraticcurveto{\pgfqpoint{8.602078in}{0.536351in}}{\pgfqpoint{8.606491in}{0.537558in}}%
\pgfpathquadraticcurveto{\pgfqpoint{8.610904in}{0.538764in}}{\pgfqpoint{8.615299in}{0.541196in}}%
\pgfpathlineto{\pgfqpoint{8.615299in}{0.531614in}}%
\pgfpathquadraticcurveto{\pgfqpoint{8.610958in}{0.529596in}}{\pgfqpoint{8.606311in}{0.528588in}}%
\pgfpathquadraticcurveto{\pgfqpoint{8.601682in}{0.527561in}}{\pgfqpoint{8.596441in}{0.527561in}}%
\pgfpathquadraticcurveto{\pgfqpoint{8.582193in}{0.527561in}}{\pgfqpoint{8.573799in}{0.536513in}}%
\pgfpathquadraticcurveto{\pgfqpoint{8.565424in}{0.545465in}}{\pgfqpoint{8.565424in}{0.560667in}}%
\pgfpathquadraticcurveto{\pgfqpoint{8.565424in}{0.576086in}}{\pgfqpoint{8.573889in}{0.584912in}}%
\pgfpathquadraticcurveto{\pgfqpoint{8.582373in}{0.593756in}}{\pgfqpoint{8.597125in}{0.593756in}}%
\pgfpathquadraticcurveto{\pgfqpoint{8.601898in}{0.593756in}}{\pgfqpoint{8.606455in}{0.592765in}}%
\pgfpathquadraticcurveto{\pgfqpoint{8.611013in}{0.591792in}}{\pgfqpoint{8.615299in}{0.589829in}}%
\pgfpathlineto{\pgfqpoint{8.615299in}{0.589829in}}%
\pgfpathclose%
\pgfpathmoveto{\pgfqpoint{8.685727in}{0.567260in}}%
\pgfpathlineto{\pgfqpoint{8.685727in}{0.529200in}}%
\pgfpathlineto{\pgfqpoint{8.675370in}{0.529200in}}%
\pgfpathlineto{\pgfqpoint{8.675370in}{0.566917in}}%
\pgfpathquadraticcurveto{\pgfqpoint{8.675370in}{0.575869in}}{\pgfqpoint{8.671876in}{0.580300in}}%
\pgfpathquadraticcurveto{\pgfqpoint{8.668381in}{0.584749in}}{\pgfqpoint{8.661411in}{0.584749in}}%
\pgfpathquadraticcurveto{\pgfqpoint{8.653017in}{0.584749in}}{\pgfqpoint{8.648172in}{0.579400in}}%
\pgfpathquadraticcurveto{\pgfqpoint{8.643326in}{0.574068in}}{\pgfqpoint{8.643326in}{0.564828in}}%
\pgfpathlineto{\pgfqpoint{8.643326in}{0.529200in}}%
\pgfpathlineto{\pgfqpoint{8.632915in}{0.529200in}}%
\pgfpathlineto{\pgfqpoint{8.632915in}{0.616793in}}%
\pgfpathlineto{\pgfqpoint{8.643326in}{0.616793in}}%
\pgfpathlineto{\pgfqpoint{8.643326in}{0.582444in}}%
\pgfpathquadraticcurveto{\pgfqpoint{8.647055in}{0.588136in}}{\pgfqpoint{8.652080in}{0.590946in}}%
\pgfpathquadraticcurveto{\pgfqpoint{8.657124in}{0.593756in}}{\pgfqpoint{8.663716in}{0.593756in}}%
\pgfpathquadraticcurveto{\pgfqpoint{8.674577in}{0.593756in}}{\pgfqpoint{8.680143in}{0.587037in}}%
\pgfpathquadraticcurveto{\pgfqpoint{8.685727in}{0.580318in}}{\pgfqpoint{8.685727in}{0.567260in}}%
\pgfpathlineto{\pgfqpoint{8.685727in}{0.567260in}}%
\pgfpathclose%
\pgfpathmoveto{\pgfqpoint{8.735026in}{0.560883in}}%
\pgfpathquadraticcurveto{\pgfqpoint{8.722472in}{0.560883in}}{\pgfqpoint{8.717626in}{0.558019in}}%
\pgfpathquadraticcurveto{\pgfqpoint{8.712781in}{0.555156in}}{\pgfqpoint{8.712781in}{0.548221in}}%
\pgfpathquadraticcurveto{\pgfqpoint{8.712781in}{0.542709in}}{\pgfqpoint{8.716420in}{0.539467in}}%
\pgfpathquadraticcurveto{\pgfqpoint{8.720058in}{0.536243in}}{\pgfqpoint{8.726290in}{0.536243in}}%
\pgfpathquadraticcurveto{\pgfqpoint{8.734918in}{0.536243in}}{\pgfqpoint{8.740124in}{0.542349in}}%
\pgfpathquadraticcurveto{\pgfqpoint{8.745329in}{0.548455in}}{\pgfqpoint{8.745329in}{0.558578in}}%
\pgfpathlineto{\pgfqpoint{8.745329in}{0.560883in}}%
\pgfpathlineto{\pgfqpoint{8.735026in}{0.560883in}}%
\pgfpathlineto{\pgfqpoint{8.735026in}{0.560883in}}%
\pgfpathclose%
\pgfpathmoveto{\pgfqpoint{8.755686in}{0.565170in}}%
\pgfpathlineto{\pgfqpoint{8.755686in}{0.529200in}}%
\pgfpathlineto{\pgfqpoint{8.745329in}{0.529200in}}%
\pgfpathlineto{\pgfqpoint{8.745329in}{0.538764in}}%
\pgfpathquadraticcurveto{\pgfqpoint{8.741781in}{0.533037in}}{\pgfqpoint{8.736485in}{0.530299in}}%
\pgfpathquadraticcurveto{\pgfqpoint{8.731190in}{0.527561in}}{\pgfqpoint{8.723534in}{0.527561in}}%
\pgfpathquadraticcurveto{\pgfqpoint{8.713862in}{0.527561in}}{\pgfqpoint{8.708134in}{0.533001in}}%
\pgfpathquadraticcurveto{\pgfqpoint{8.702424in}{0.538440in}}{\pgfqpoint{8.702424in}{0.547554in}}%
\pgfpathquadraticcurveto{\pgfqpoint{8.702424in}{0.558182in}}{\pgfqpoint{8.709539in}{0.563585in}}%
\pgfpathquadraticcurveto{\pgfqpoint{8.716672in}{0.568989in}}{\pgfqpoint{8.730793in}{0.568989in}}%
\pgfpathlineto{\pgfqpoint{8.745329in}{0.568989in}}%
\pgfpathlineto{\pgfqpoint{8.745329in}{0.570016in}}%
\pgfpathquadraticcurveto{\pgfqpoint{8.745329in}{0.577166in}}{\pgfqpoint{8.740628in}{0.581075in}}%
\pgfpathquadraticcurveto{\pgfqpoint{8.735927in}{0.584984in}}{\pgfqpoint{8.727425in}{0.584984in}}%
\pgfpathquadraticcurveto{\pgfqpoint{8.722021in}{0.584984in}}{\pgfqpoint{8.716888in}{0.583687in}}%
\pgfpathquadraticcurveto{\pgfqpoint{8.711773in}{0.582390in}}{\pgfqpoint{8.707053in}{0.579796in}}%
\pgfpathlineto{\pgfqpoint{8.707053in}{0.589379in}}%
\pgfpathquadraticcurveto{\pgfqpoint{8.712727in}{0.591576in}}{\pgfqpoint{8.718077in}{0.592657in}}%
\pgfpathquadraticcurveto{\pgfqpoint{8.723426in}{0.593756in}}{\pgfqpoint{8.728488in}{0.593756in}}%
\pgfpathquadraticcurveto{\pgfqpoint{8.742177in}{0.593756in}}{\pgfqpoint{8.748932in}{0.586659in}}%
\pgfpathquadraticcurveto{\pgfqpoint{8.755686in}{0.579580in}}{\pgfqpoint{8.755686in}{0.565170in}}%
\pgfpathlineto{\pgfqpoint{8.755686in}{0.565170in}}%
\pgfpathclose%
\pgfpathmoveto{\pgfqpoint{8.829428in}{0.567260in}}%
\pgfpathlineto{\pgfqpoint{8.829428in}{0.529200in}}%
\pgfpathlineto{\pgfqpoint{8.819071in}{0.529200in}}%
\pgfpathlineto{\pgfqpoint{8.819071in}{0.566917in}}%
\pgfpathquadraticcurveto{\pgfqpoint{8.819071in}{0.575869in}}{\pgfqpoint{8.815577in}{0.580300in}}%
\pgfpathquadraticcurveto{\pgfqpoint{8.812082in}{0.584749in}}{\pgfqpoint{8.805111in}{0.584749in}}%
\pgfpathquadraticcurveto{\pgfqpoint{8.796718in}{0.584749in}}{\pgfqpoint{8.791873in}{0.579400in}}%
\pgfpathquadraticcurveto{\pgfqpoint{8.787027in}{0.574068in}}{\pgfqpoint{8.787027in}{0.564828in}}%
\pgfpathlineto{\pgfqpoint{8.787027in}{0.529200in}}%
\pgfpathlineto{\pgfqpoint{8.776616in}{0.529200in}}%
\pgfpathlineto{\pgfqpoint{8.776616in}{0.592243in}}%
\pgfpathlineto{\pgfqpoint{8.787027in}{0.592243in}}%
\pgfpathlineto{\pgfqpoint{8.787027in}{0.582444in}}%
\pgfpathquadraticcurveto{\pgfqpoint{8.790756in}{0.588136in}}{\pgfqpoint{8.795781in}{0.590946in}}%
\pgfpathquadraticcurveto{\pgfqpoint{8.800825in}{0.593756in}}{\pgfqpoint{8.807417in}{0.593756in}}%
\pgfpathquadraticcurveto{\pgfqpoint{8.818278in}{0.593756in}}{\pgfqpoint{8.823844in}{0.587037in}}%
\pgfpathquadraticcurveto{\pgfqpoint{8.829428in}{0.580318in}}{\pgfqpoint{8.829428in}{0.567260in}}%
\pgfpathlineto{\pgfqpoint{8.829428in}{0.567260in}}%
\pgfpathclose%
\pgfpathmoveto{\pgfqpoint{8.891552in}{0.561460in}}%
\pgfpathquadraticcurveto{\pgfqpoint{8.891552in}{0.572717in}}{\pgfqpoint{8.886905in}{0.578896in}}%
\pgfpathquadraticcurveto{\pgfqpoint{8.882276in}{0.585092in}}{\pgfqpoint{8.873882in}{0.585092in}}%
\pgfpathquadraticcurveto{\pgfqpoint{8.865560in}{0.585092in}}{\pgfqpoint{8.860913in}{0.578896in}}%
\pgfpathquadraticcurveto{\pgfqpoint{8.856266in}{0.572717in}}{\pgfqpoint{8.856266in}{0.561460in}}%
\pgfpathquadraticcurveto{\pgfqpoint{8.856266in}{0.550256in}}{\pgfqpoint{8.860913in}{0.544060in}}%
\pgfpathquadraticcurveto{\pgfqpoint{8.865560in}{0.537864in}}{\pgfqpoint{8.873882in}{0.537864in}}%
\pgfpathquadraticcurveto{\pgfqpoint{8.882276in}{0.537864in}}{\pgfqpoint{8.886905in}{0.544060in}}%
\pgfpathquadraticcurveto{\pgfqpoint{8.891552in}{0.550256in}}{\pgfqpoint{8.891552in}{0.561460in}}%
\pgfpathlineto{\pgfqpoint{8.891552in}{0.561460in}}%
\pgfpathclose%
\pgfpathmoveto{\pgfqpoint{8.901909in}{0.537017in}}%
\pgfpathquadraticcurveto{\pgfqpoint{8.901909in}{0.520932in}}{\pgfqpoint{8.894758in}{0.513079in}}%
\pgfpathquadraticcurveto{\pgfqpoint{8.887625in}{0.505226in}}{\pgfqpoint{8.872873in}{0.505226in}}%
\pgfpathquadraticcurveto{\pgfqpoint{8.867416in}{0.505226in}}{\pgfqpoint{8.862570in}{0.506036in}}%
\pgfpathquadraticcurveto{\pgfqpoint{8.857725in}{0.506847in}}{\pgfqpoint{8.853168in}{0.508540in}}%
\pgfpathlineto{\pgfqpoint{8.853168in}{0.518609in}}%
\pgfpathquadraticcurveto{\pgfqpoint{8.857725in}{0.516141in}}{\pgfqpoint{8.862174in}{0.514970in}}%
\pgfpathquadraticcurveto{\pgfqpoint{8.866623in}{0.513782in}}{\pgfqpoint{8.871234in}{0.513782in}}%
\pgfpathquadraticcurveto{\pgfqpoint{8.881429in}{0.513782in}}{\pgfqpoint{8.886490in}{0.519095in}}%
\pgfpathquadraticcurveto{\pgfqpoint{8.891552in}{0.524409in}}{\pgfqpoint{8.891552in}{0.535162in}}%
\pgfpathlineto{\pgfqpoint{8.891552in}{0.540295in}}%
\pgfpathquadraticcurveto{\pgfqpoint{8.888346in}{0.534712in}}{\pgfqpoint{8.883338in}{0.531956in}}%
\pgfpathquadraticcurveto{\pgfqpoint{8.878331in}{0.529200in}}{\pgfqpoint{8.871342in}{0.529200in}}%
\pgfpathquadraticcurveto{\pgfqpoint{8.859760in}{0.529200in}}{\pgfqpoint{8.852664in}{0.538026in}}%
\pgfpathquadraticcurveto{\pgfqpoint{8.845567in}{0.546870in}}{\pgfqpoint{8.845567in}{0.561460in}}%
\pgfpathquadraticcurveto{\pgfqpoint{8.845567in}{0.576086in}}{\pgfqpoint{8.852664in}{0.584912in}}%
\pgfpathquadraticcurveto{\pgfqpoint{8.859760in}{0.593756in}}{\pgfqpoint{8.871342in}{0.593756in}}%
\pgfpathquadraticcurveto{\pgfqpoint{8.878331in}{0.593756in}}{\pgfqpoint{8.883338in}{0.591000in}}%
\pgfpathquadraticcurveto{\pgfqpoint{8.888346in}{0.588244in}}{\pgfqpoint{8.891552in}{0.582678in}}%
\pgfpathlineto{\pgfqpoint{8.891552in}{0.592243in}}%
\pgfpathlineto{\pgfqpoint{8.901909in}{0.592243in}}%
\pgfpathlineto{\pgfqpoint{8.901909in}{0.537017in}}%
\pgfpathlineto{\pgfqpoint{8.901909in}{0.537017in}}%
\pgfpathclose%
\pgfpathmoveto{\pgfqpoint{8.977182in}{0.563315in}}%
\pgfpathlineto{\pgfqpoint{8.977182in}{0.558254in}}%
\pgfpathlineto{\pgfqpoint{8.929557in}{0.558254in}}%
\pgfpathquadraticcurveto{\pgfqpoint{8.930242in}{0.547554in}}{\pgfqpoint{8.936006in}{0.541953in}}%
\pgfpathquadraticcurveto{\pgfqpoint{8.941770in}{0.536351in}}{\pgfqpoint{8.952073in}{0.536351in}}%
\pgfpathquadraticcurveto{\pgfqpoint{8.958035in}{0.536351in}}{\pgfqpoint{8.963636in}{0.537810in}}%
\pgfpathquadraticcurveto{\pgfqpoint{8.969238in}{0.539269in}}{\pgfqpoint{8.974768in}{0.542205in}}%
\pgfpathlineto{\pgfqpoint{8.974768in}{0.532406in}}%
\pgfpathquadraticcurveto{\pgfqpoint{8.969184in}{0.530047in}}{\pgfqpoint{8.963330in}{0.528804in}}%
\pgfpathquadraticcurveto{\pgfqpoint{8.957476in}{0.527561in}}{\pgfqpoint{8.951460in}{0.527561in}}%
\pgfpathquadraticcurveto{\pgfqpoint{8.936366in}{0.527561in}}{\pgfqpoint{8.927558in}{0.536333in}}%
\pgfpathquadraticcurveto{\pgfqpoint{8.918750in}{0.545123in}}{\pgfqpoint{8.918750in}{0.560109in}}%
\pgfpathquadraticcurveto{\pgfqpoint{8.918750in}{0.575581in}}{\pgfqpoint{8.927108in}{0.584659in}}%
\pgfpathquadraticcurveto{\pgfqpoint{8.935465in}{0.593756in}}{\pgfqpoint{8.949659in}{0.593756in}}%
\pgfpathquadraticcurveto{\pgfqpoint{8.962376in}{0.593756in}}{\pgfqpoint{8.969779in}{0.585560in}}%
\pgfpathquadraticcurveto{\pgfqpoint{8.977182in}{0.577383in}}{\pgfqpoint{8.977182in}{0.563315in}}%
\pgfpathlineto{\pgfqpoint{8.977182in}{0.563315in}}%
\pgfpathclose%
\pgfpathmoveto{\pgfqpoint{8.966825in}{0.566359in}}%
\pgfpathquadraticcurveto{\pgfqpoint{8.966717in}{0.574843in}}{\pgfqpoint{8.962069in}{0.579904in}}%
\pgfpathquadraticcurveto{\pgfqpoint{8.957422in}{0.584984in}}{\pgfqpoint{8.949767in}{0.584984in}}%
\pgfpathquadraticcurveto{\pgfqpoint{8.941103in}{0.584984in}}{\pgfqpoint{8.935898in}{0.580084in}}%
\pgfpathquadraticcurveto{\pgfqpoint{8.930692in}{0.575185in}}{\pgfqpoint{8.929900in}{0.566287in}}%
\pgfpathlineto{\pgfqpoint{8.966825in}{0.566359in}}%
\pgfpathlineto{\pgfqpoint{8.966825in}{0.566359in}}%
\pgfpathclose%
\pgfpathmoveto{\pgfqpoint{9.035667in}{0.582678in}}%
\pgfpathlineto{\pgfqpoint{9.035667in}{0.616793in}}%
\pgfpathlineto{\pgfqpoint{9.046024in}{0.616793in}}%
\pgfpathlineto{\pgfqpoint{9.046024in}{0.529200in}}%
\pgfpathlineto{\pgfqpoint{9.035667in}{0.529200in}}%
\pgfpathlineto{\pgfqpoint{9.035667in}{0.538656in}}%
\pgfpathquadraticcurveto{\pgfqpoint{9.032407in}{0.533037in}}{\pgfqpoint{9.027417in}{0.530299in}}%
\pgfpathquadraticcurveto{\pgfqpoint{9.022446in}{0.527561in}}{\pgfqpoint{9.015457in}{0.527561in}}%
\pgfpathquadraticcurveto{\pgfqpoint{9.004038in}{0.527561in}}{\pgfqpoint{8.996851in}{0.536675in}}%
\pgfpathquadraticcurveto{\pgfqpoint{8.989682in}{0.545807in}}{\pgfqpoint{8.989682in}{0.560667in}}%
\pgfpathquadraticcurveto{\pgfqpoint{8.989682in}{0.575527in}}{\pgfqpoint{8.996851in}{0.584641in}}%
\pgfpathquadraticcurveto{\pgfqpoint{9.004038in}{0.593756in}}{\pgfqpoint{9.015457in}{0.593756in}}%
\pgfpathquadraticcurveto{\pgfqpoint{9.022446in}{0.593756in}}{\pgfqpoint{9.027417in}{0.591018in}}%
\pgfpathquadraticcurveto{\pgfqpoint{9.032407in}{0.588298in}}{\pgfqpoint{9.035667in}{0.582678in}}%
\pgfpathlineto{\pgfqpoint{9.035667in}{0.582678in}}%
\pgfpathclose%
\pgfpathmoveto{\pgfqpoint{9.000381in}{0.560667in}}%
\pgfpathquadraticcurveto{\pgfqpoint{9.000381in}{0.549248in}}{\pgfqpoint{9.005082in}{0.542745in}}%
\pgfpathquadraticcurveto{\pgfqpoint{9.009784in}{0.536243in}}{\pgfqpoint{9.017997in}{0.536243in}}%
\pgfpathquadraticcurveto{\pgfqpoint{9.026211in}{0.536243in}}{\pgfqpoint{9.030930in}{0.542745in}}%
\pgfpathquadraticcurveto{\pgfqpoint{9.035667in}{0.549248in}}{\pgfqpoint{9.035667in}{0.560667in}}%
\pgfpathquadraticcurveto{\pgfqpoint{9.035667in}{0.572087in}}{\pgfqpoint{9.030930in}{0.578589in}}%
\pgfpathquadraticcurveto{\pgfqpoint{9.026211in}{0.585092in}}{\pgfqpoint{9.017997in}{0.585092in}}%
\pgfpathquadraticcurveto{\pgfqpoint{9.009784in}{0.585092in}}{\pgfqpoint{9.005082in}{0.578589in}}%
\pgfpathquadraticcurveto{\pgfqpoint{9.000381in}{0.572087in}}{\pgfqpoint{9.000381in}{0.560667in}}%
\pgfpathlineto{\pgfqpoint{9.000381in}{0.560667in}}%
\pgfpathclose%
\pgfpathmoveto{\pgfqpoint{9.068827in}{0.543502in}}%
\pgfpathlineto{\pgfqpoint{9.080715in}{0.543502in}}%
\pgfpathlineto{\pgfqpoint{9.080715in}{0.529200in}}%
\pgfpathlineto{\pgfqpoint{9.068827in}{0.529200in}}%
\pgfpathlineto{\pgfqpoint{9.068827in}{0.543502in}}%
\pgfpathlineto{\pgfqpoint{9.068827in}{0.543502in}}%
\pgfpathclose%
\pgfusepath{fill}%
\end{pgfscope}%
\begin{pgfscope}%
\pgfsetbuttcap%
\pgfsetroundjoin%
\definecolor{currentfill}{rgb}{0.309804,0.372549,0.419608}%
\pgfsetfillcolor{currentfill}%
\pgfsetlinewidth{0.000000pt}%
\definecolor{currentstroke}{rgb}{0.309804,0.372549,0.419608}%
\pgfsetstrokecolor{currentstroke}%
\pgfsetdash{}{0pt}%
\pgfpathmoveto{\pgfqpoint{9.210765in}{0.613245in}}%
\pgfpathlineto{\pgfqpoint{9.259055in}{0.613245in}}%
\pgfpathlineto{\pgfqpoint{9.259055in}{0.603662in}}%
\pgfpathlineto{\pgfqpoint{9.222130in}{0.603662in}}%
\pgfpathlineto{\pgfqpoint{9.222130in}{0.578896in}}%
\pgfpathlineto{\pgfqpoint{9.255453in}{0.578896in}}%
\pgfpathlineto{\pgfqpoint{9.255453in}{0.569331in}}%
\pgfpathlineto{\pgfqpoint{9.222130in}{0.569331in}}%
\pgfpathlineto{\pgfqpoint{9.222130in}{0.529200in}}%
\pgfpathlineto{\pgfqpoint{9.210765in}{0.529200in}}%
\pgfpathlineto{\pgfqpoint{9.210765in}{0.613245in}}%
\pgfpathlineto{\pgfqpoint{9.210765in}{0.613245in}}%
\pgfpathclose%
\pgfpathmoveto{\pgfqpoint{9.268224in}{0.592243in}}%
\pgfpathlineto{\pgfqpoint{9.278581in}{0.592243in}}%
\pgfpathlineto{\pgfqpoint{9.278581in}{0.529200in}}%
\pgfpathlineto{\pgfqpoint{9.268224in}{0.529200in}}%
\pgfpathlineto{\pgfqpoint{9.268224in}{0.592243in}}%
\pgfpathlineto{\pgfqpoint{9.268224in}{0.592243in}}%
\pgfpathclose%
\pgfpathmoveto{\pgfqpoint{9.268224in}{0.616793in}}%
\pgfpathlineto{\pgfqpoint{9.278581in}{0.616793in}}%
\pgfpathlineto{\pgfqpoint{9.278581in}{0.603662in}}%
\pgfpathlineto{\pgfqpoint{9.268224in}{0.603662in}}%
\pgfpathlineto{\pgfqpoint{9.268224in}{0.616793in}}%
\pgfpathlineto{\pgfqpoint{9.268224in}{0.616793in}}%
\pgfpathclose%
\pgfpathmoveto{\pgfqpoint{9.300249in}{0.616793in}}%
\pgfpathlineto{\pgfqpoint{9.310606in}{0.616793in}}%
\pgfpathlineto{\pgfqpoint{9.310606in}{0.529200in}}%
\pgfpathlineto{\pgfqpoint{9.300249in}{0.529200in}}%
\pgfpathlineto{\pgfqpoint{9.300249in}{0.616793in}}%
\pgfpathlineto{\pgfqpoint{9.300249in}{0.616793in}}%
\pgfpathclose%
\pgfpathmoveto{\pgfqpoint{9.332275in}{0.616793in}}%
\pgfpathlineto{\pgfqpoint{9.342632in}{0.616793in}}%
\pgfpathlineto{\pgfqpoint{9.342632in}{0.529200in}}%
\pgfpathlineto{\pgfqpoint{9.332275in}{0.529200in}}%
\pgfpathlineto{\pgfqpoint{9.332275in}{0.616793in}}%
\pgfpathlineto{\pgfqpoint{9.332275in}{0.616793in}}%
\pgfpathclose%
\pgfusepath{fill}%
\end{pgfscope}%
\begin{pgfscope}%
\pgfsetbuttcap%
\pgfsetroundjoin%
\definecolor{currentfill}{rgb}{0.309804,0.372549,0.419608}%
\pgfsetfillcolor{currentfill}%
\pgfsetlinewidth{0.000000pt}%
\definecolor{currentstroke}{rgb}{0.309804,0.372549,0.419608}%
\pgfsetstrokecolor{currentstroke}%
\pgfsetdash{}{0pt}%
\pgfpathmoveto{\pgfqpoint{9.417316in}{0.554075in}}%
\pgfpathlineto{\pgfqpoint{9.417316in}{0.592243in}}%
\pgfpathlineto{\pgfqpoint{9.427673in}{0.592243in}}%
\pgfpathlineto{\pgfqpoint{9.427673in}{0.554471in}}%
\pgfpathquadraticcurveto{\pgfqpoint{9.427673in}{0.545519in}}{\pgfqpoint{9.431150in}{0.541034in}}%
\pgfpathquadraticcurveto{\pgfqpoint{9.434644in}{0.536567in}}{\pgfqpoint{9.441633in}{0.536567in}}%
\pgfpathquadraticcurveto{\pgfqpoint{9.450009in}{0.536567in}}{\pgfqpoint{9.454872in}{0.541917in}}%
\pgfpathquadraticcurveto{\pgfqpoint{9.459753in}{0.547266in}}{\pgfqpoint{9.459753in}{0.556506in}}%
\pgfpathlineto{\pgfqpoint{9.459753in}{0.592243in}}%
\pgfpathlineto{\pgfqpoint{9.470110in}{0.592243in}}%
\pgfpathlineto{\pgfqpoint{9.470110in}{0.529200in}}%
\pgfpathlineto{\pgfqpoint{9.459753in}{0.529200in}}%
\pgfpathlineto{\pgfqpoint{9.459753in}{0.538891in}}%
\pgfpathquadraticcurveto{\pgfqpoint{9.455989in}{0.533145in}}{\pgfqpoint{9.450999in}{0.530353in}}%
\pgfpathquadraticcurveto{\pgfqpoint{9.446028in}{0.527561in}}{\pgfqpoint{9.439435in}{0.527561in}}%
\pgfpathquadraticcurveto{\pgfqpoint{9.428574in}{0.527561in}}{\pgfqpoint{9.422936in}{0.534315in}}%
\pgfpathquadraticcurveto{\pgfqpoint{9.417316in}{0.541070in}}{\pgfqpoint{9.417316in}{0.554075in}}%
\pgfpathlineto{\pgfqpoint{9.417316in}{0.554075in}}%
\pgfpathclose%
\pgfpathmoveto{\pgfqpoint{9.443380in}{0.593756in}}%
\pgfpathlineto{\pgfqpoint{9.443380in}{0.593756in}}%
\pgfpathlineto{\pgfqpoint{9.443380in}{0.593756in}}%
\pgfpathclose%
\pgfpathmoveto{\pgfqpoint{9.531622in}{0.590387in}}%
\pgfpathlineto{\pgfqpoint{9.531622in}{0.580589in}}%
\pgfpathquadraticcurveto{\pgfqpoint{9.527245in}{0.582840in}}{\pgfqpoint{9.522507in}{0.583957in}}%
\pgfpathquadraticcurveto{\pgfqpoint{9.517788in}{0.585092in}}{\pgfqpoint{9.512709in}{0.585092in}}%
\pgfpathquadraticcurveto{\pgfqpoint{9.505000in}{0.585092in}}{\pgfqpoint{9.501145in}{0.582732in}}%
\pgfpathquadraticcurveto{\pgfqpoint{9.497290in}{0.580373in}}{\pgfqpoint{9.497290in}{0.575635in}}%
\pgfpathquadraticcurveto{\pgfqpoint{9.497290in}{0.572033in}}{\pgfqpoint{9.500046in}{0.569980in}}%
\pgfpathquadraticcurveto{\pgfqpoint{9.502802in}{0.567926in}}{\pgfqpoint{9.511142in}{0.566071in}}%
\pgfpathlineto{\pgfqpoint{9.514690in}{0.565278in}}%
\pgfpathquadraticcurveto{\pgfqpoint{9.525714in}{0.562919in}}{\pgfqpoint{9.530361in}{0.558614in}}%
\pgfpathquadraticcurveto{\pgfqpoint{9.535008in}{0.554309in}}{\pgfqpoint{9.535008in}{0.546600in}}%
\pgfpathquadraticcurveto{\pgfqpoint{9.535008in}{0.537810in}}{\pgfqpoint{9.528055in}{0.532676in}}%
\pgfpathquadraticcurveto{\pgfqpoint{9.521102in}{0.527561in}}{\pgfqpoint{9.508944in}{0.527561in}}%
\pgfpathquadraticcurveto{\pgfqpoint{9.503883in}{0.527561in}}{\pgfqpoint{9.498389in}{0.528552in}}%
\pgfpathquadraticcurveto{\pgfqpoint{9.492895in}{0.529542in}}{\pgfqpoint{9.486825in}{0.531506in}}%
\pgfpathlineto{\pgfqpoint{9.486825in}{0.542205in}}%
\pgfpathquadraticcurveto{\pgfqpoint{9.492571in}{0.539215in}}{\pgfqpoint{9.498137in}{0.537720in}}%
\pgfpathquadraticcurveto{\pgfqpoint{9.503703in}{0.536243in}}{\pgfqpoint{9.509178in}{0.536243in}}%
\pgfpathquadraticcurveto{\pgfqpoint{9.516491in}{0.536243in}}{\pgfqpoint{9.520418in}{0.538746in}}%
\pgfpathquadraticcurveto{\pgfqpoint{9.524363in}{0.541250in}}{\pgfqpoint{9.524363in}{0.545807in}}%
\pgfpathquadraticcurveto{\pgfqpoint{9.524363in}{0.550022in}}{\pgfqpoint{9.521517in}{0.552274in}}%
\pgfpathquadraticcurveto{\pgfqpoint{9.518689in}{0.554525in}}{\pgfqpoint{9.509052in}{0.556614in}}%
\pgfpathlineto{\pgfqpoint{9.505450in}{0.557461in}}%
\pgfpathquadraticcurveto{\pgfqpoint{9.495831in}{0.559478in}}{\pgfqpoint{9.491545in}{0.563675in}}%
\pgfpathquadraticcurveto{\pgfqpoint{9.487276in}{0.567872in}}{\pgfqpoint{9.487276in}{0.575185in}}%
\pgfpathquadraticcurveto{\pgfqpoint{9.487276in}{0.584083in}}{\pgfqpoint{9.493580in}{0.588910in}}%
\pgfpathquadraticcurveto{\pgfqpoint{9.499884in}{0.593756in}}{\pgfqpoint{9.511484in}{0.593756in}}%
\pgfpathquadraticcurveto{\pgfqpoint{9.517212in}{0.593756in}}{\pgfqpoint{9.522273in}{0.592909in}}%
\pgfpathquadraticcurveto{\pgfqpoint{9.527353in}{0.592080in}}{\pgfqpoint{9.531622in}{0.590387in}}%
\pgfpathlineto{\pgfqpoint{9.531622in}{0.590387in}}%
\pgfpathclose%
\pgfpathmoveto{\pgfqpoint{9.605417in}{0.563315in}}%
\pgfpathlineto{\pgfqpoint{9.605417in}{0.558254in}}%
\pgfpathlineto{\pgfqpoint{9.557793in}{0.558254in}}%
\pgfpathquadraticcurveto{\pgfqpoint{9.558478in}{0.547554in}}{\pgfqpoint{9.564242in}{0.541953in}}%
\pgfpathquadraticcurveto{\pgfqpoint{9.570005in}{0.536351in}}{\pgfqpoint{9.580308in}{0.536351in}}%
\pgfpathquadraticcurveto{\pgfqpoint{9.586270in}{0.536351in}}{\pgfqpoint{9.591872in}{0.537810in}}%
\pgfpathquadraticcurveto{\pgfqpoint{9.597474in}{0.539269in}}{\pgfqpoint{9.603004in}{0.542205in}}%
\pgfpathlineto{\pgfqpoint{9.603004in}{0.532406in}}%
\pgfpathquadraticcurveto{\pgfqpoint{9.597420in}{0.530047in}}{\pgfqpoint{9.591566in}{0.528804in}}%
\pgfpathquadraticcurveto{\pgfqpoint{9.585712in}{0.527561in}}{\pgfqpoint{9.579696in}{0.527561in}}%
\pgfpathquadraticcurveto{\pgfqpoint{9.564602in}{0.527561in}}{\pgfqpoint{9.555794in}{0.536333in}}%
\pgfpathquadraticcurveto{\pgfqpoint{9.546986in}{0.545123in}}{\pgfqpoint{9.546986in}{0.560109in}}%
\pgfpathquadraticcurveto{\pgfqpoint{9.546986in}{0.575581in}}{\pgfqpoint{9.555344in}{0.584659in}}%
\pgfpathquadraticcurveto{\pgfqpoint{9.563701in}{0.593756in}}{\pgfqpoint{9.577895in}{0.593756in}}%
\pgfpathquadraticcurveto{\pgfqpoint{9.590611in}{0.593756in}}{\pgfqpoint{9.598014in}{0.585560in}}%
\pgfpathquadraticcurveto{\pgfqpoint{9.605417in}{0.577383in}}{\pgfqpoint{9.605417in}{0.563315in}}%
\pgfpathlineto{\pgfqpoint{9.605417in}{0.563315in}}%
\pgfpathclose%
\pgfpathmoveto{\pgfqpoint{9.595060in}{0.566359in}}%
\pgfpathquadraticcurveto{\pgfqpoint{9.594952in}{0.574843in}}{\pgfqpoint{9.590305in}{0.579904in}}%
\pgfpathquadraticcurveto{\pgfqpoint{9.585658in}{0.584984in}}{\pgfqpoint{9.578003in}{0.584984in}}%
\pgfpathquadraticcurveto{\pgfqpoint{9.569339in}{0.584984in}}{\pgfqpoint{9.564134in}{0.580084in}}%
\pgfpathquadraticcurveto{\pgfqpoint{9.558928in}{0.575185in}}{\pgfqpoint{9.558135in}{0.566287in}}%
\pgfpathlineto{\pgfqpoint{9.595060in}{0.566359in}}%
\pgfpathlineto{\pgfqpoint{9.595060in}{0.566359in}}%
\pgfpathclose%
\pgfpathmoveto{\pgfqpoint{9.662606in}{0.590387in}}%
\pgfpathlineto{\pgfqpoint{9.662606in}{0.580589in}}%
\pgfpathquadraticcurveto{\pgfqpoint{9.658229in}{0.582840in}}{\pgfqpoint{9.653492in}{0.583957in}}%
\pgfpathquadraticcurveto{\pgfqpoint{9.648773in}{0.585092in}}{\pgfqpoint{9.643693in}{0.585092in}}%
\pgfpathquadraticcurveto{\pgfqpoint{9.635984in}{0.585092in}}{\pgfqpoint{9.632129in}{0.582732in}}%
\pgfpathquadraticcurveto{\pgfqpoint{9.628275in}{0.580373in}}{\pgfqpoint{9.628275in}{0.575635in}}%
\pgfpathquadraticcurveto{\pgfqpoint{9.628275in}{0.572033in}}{\pgfqpoint{9.631031in}{0.569980in}}%
\pgfpathquadraticcurveto{\pgfqpoint{9.633787in}{0.567926in}}{\pgfqpoint{9.642126in}{0.566071in}}%
\pgfpathlineto{\pgfqpoint{9.645675in}{0.565278in}}%
\pgfpathquadraticcurveto{\pgfqpoint{9.656698in}{0.562919in}}{\pgfqpoint{9.661345in}{0.558614in}}%
\pgfpathquadraticcurveto{\pgfqpoint{9.665992in}{0.554309in}}{\pgfqpoint{9.665992in}{0.546600in}}%
\pgfpathquadraticcurveto{\pgfqpoint{9.665992in}{0.537810in}}{\pgfqpoint{9.659040in}{0.532676in}}%
\pgfpathquadraticcurveto{\pgfqpoint{9.652087in}{0.527561in}}{\pgfqpoint{9.639929in}{0.527561in}}%
\pgfpathquadraticcurveto{\pgfqpoint{9.634867in}{0.527561in}}{\pgfqpoint{9.629374in}{0.528552in}}%
\pgfpathquadraticcurveto{\pgfqpoint{9.623880in}{0.529542in}}{\pgfqpoint{9.617810in}{0.531506in}}%
\pgfpathlineto{\pgfqpoint{9.617810in}{0.542205in}}%
\pgfpathquadraticcurveto{\pgfqpoint{9.623556in}{0.539215in}}{\pgfqpoint{9.629121in}{0.537720in}}%
\pgfpathquadraticcurveto{\pgfqpoint{9.634687in}{0.536243in}}{\pgfqpoint{9.640163in}{0.536243in}}%
\pgfpathquadraticcurveto{\pgfqpoint{9.647476in}{0.536243in}}{\pgfqpoint{9.651402in}{0.538746in}}%
\pgfpathquadraticcurveto{\pgfqpoint{9.655347in}{0.541250in}}{\pgfqpoint{9.655347in}{0.545807in}}%
\pgfpathquadraticcurveto{\pgfqpoint{9.655347in}{0.550022in}}{\pgfqpoint{9.652501in}{0.552274in}}%
\pgfpathquadraticcurveto{\pgfqpoint{9.649673in}{0.554525in}}{\pgfqpoint{9.640037in}{0.556614in}}%
\pgfpathlineto{\pgfqpoint{9.636434in}{0.557461in}}%
\pgfpathquadraticcurveto{\pgfqpoint{9.626816in}{0.559478in}}{\pgfqpoint{9.622529in}{0.563675in}}%
\pgfpathquadraticcurveto{\pgfqpoint{9.618260in}{0.567872in}}{\pgfqpoint{9.618260in}{0.575185in}}%
\pgfpathquadraticcurveto{\pgfqpoint{9.618260in}{0.584083in}}{\pgfqpoint{9.624564in}{0.588910in}}%
\pgfpathquadraticcurveto{\pgfqpoint{9.630869in}{0.593756in}}{\pgfqpoint{9.642468in}{0.593756in}}%
\pgfpathquadraticcurveto{\pgfqpoint{9.648196in}{0.593756in}}{\pgfqpoint{9.653258in}{0.592909in}}%
\pgfpathquadraticcurveto{\pgfqpoint{9.658337in}{0.592080in}}{\pgfqpoint{9.662606in}{0.590387in}}%
\pgfpathlineto{\pgfqpoint{9.662606in}{0.590387in}}%
\pgfpathclose%
\pgfusepath{fill}%
\end{pgfscope}%
\begin{pgfscope}%
\pgfsetbuttcap%
\pgfsetroundjoin%
\definecolor{currentfill}{rgb}{0.309804,0.372549,0.419608}%
\pgfsetfillcolor{currentfill}%
\pgfsetlinewidth{0.000000pt}%
\definecolor{currentstroke}{rgb}{0.309804,0.372549,0.419608}%
\pgfsetstrokecolor{currentstroke}%
\pgfsetdash{}{0pt}%
\pgfpathmoveto{\pgfqpoint{9.754896in}{0.538656in}}%
\pgfpathlineto{\pgfqpoint{9.754896in}{0.505226in}}%
\pgfpathlineto{\pgfqpoint{9.744485in}{0.505226in}}%
\pgfpathlineto{\pgfqpoint{9.744485in}{0.592243in}}%
\pgfpathlineto{\pgfqpoint{9.754896in}{0.592243in}}%
\pgfpathlineto{\pgfqpoint{9.754896in}{0.582678in}}%
\pgfpathquadraticcurveto{\pgfqpoint{9.758174in}{0.588298in}}{\pgfqpoint{9.763146in}{0.591018in}}%
\pgfpathquadraticcurveto{\pgfqpoint{9.768135in}{0.593756in}}{\pgfqpoint{9.775052in}{0.593756in}}%
\pgfpathquadraticcurveto{\pgfqpoint{9.786544in}{0.593756in}}{\pgfqpoint{9.793712in}{0.584641in}}%
\pgfpathquadraticcurveto{\pgfqpoint{9.800899in}{0.575527in}}{\pgfqpoint{9.800899in}{0.560667in}}%
\pgfpathquadraticcurveto{\pgfqpoint{9.800899in}{0.545807in}}{\pgfqpoint{9.793712in}{0.536675in}}%
\pgfpathquadraticcurveto{\pgfqpoint{9.786544in}{0.527561in}}{\pgfqpoint{9.775052in}{0.527561in}}%
\pgfpathquadraticcurveto{\pgfqpoint{9.768135in}{0.527561in}}{\pgfqpoint{9.763146in}{0.530299in}}%
\pgfpathquadraticcurveto{\pgfqpoint{9.758174in}{0.533037in}}{\pgfqpoint{9.754896in}{0.538656in}}%
\pgfpathlineto{\pgfqpoint{9.754896in}{0.538656in}}%
\pgfpathclose%
\pgfpathmoveto{\pgfqpoint{9.790146in}{0.560667in}}%
\pgfpathquadraticcurveto{\pgfqpoint{9.790146in}{0.572087in}}{\pgfqpoint{9.785445in}{0.578589in}}%
\pgfpathquadraticcurveto{\pgfqpoint{9.780744in}{0.585092in}}{\pgfqpoint{9.772530in}{0.585092in}}%
\pgfpathquadraticcurveto{\pgfqpoint{9.764299in}{0.585092in}}{\pgfqpoint{9.759597in}{0.578589in}}%
\pgfpathquadraticcurveto{\pgfqpoint{9.754896in}{0.572087in}}{\pgfqpoint{9.754896in}{0.560667in}}%
\pgfpathquadraticcurveto{\pgfqpoint{9.754896in}{0.549248in}}{\pgfqpoint{9.759597in}{0.542745in}}%
\pgfpathquadraticcurveto{\pgfqpoint{9.764299in}{0.536243in}}{\pgfqpoint{9.772530in}{0.536243in}}%
\pgfpathquadraticcurveto{\pgfqpoint{9.780744in}{0.536243in}}{\pgfqpoint{9.785445in}{0.542745in}}%
\pgfpathquadraticcurveto{\pgfqpoint{9.790146in}{0.549248in}}{\pgfqpoint{9.790146in}{0.560667in}}%
\pgfpathlineto{\pgfqpoint{9.790146in}{0.560667in}}%
\pgfpathclose%
\pgfpathmoveto{\pgfqpoint{9.846722in}{0.560883in}}%
\pgfpathquadraticcurveto{\pgfqpoint{9.834168in}{0.560883in}}{\pgfqpoint{9.829322in}{0.558019in}}%
\pgfpathquadraticcurveto{\pgfqpoint{9.824477in}{0.555156in}}{\pgfqpoint{9.824477in}{0.548221in}}%
\pgfpathquadraticcurveto{\pgfqpoint{9.824477in}{0.542709in}}{\pgfqpoint{9.828116in}{0.539467in}}%
\pgfpathquadraticcurveto{\pgfqpoint{9.831754in}{0.536243in}}{\pgfqpoint{9.837986in}{0.536243in}}%
\pgfpathquadraticcurveto{\pgfqpoint{9.846614in}{0.536243in}}{\pgfqpoint{9.851820in}{0.542349in}}%
\pgfpathquadraticcurveto{\pgfqpoint{9.857025in}{0.548455in}}{\pgfqpoint{9.857025in}{0.558578in}}%
\pgfpathlineto{\pgfqpoint{9.857025in}{0.560883in}}%
\pgfpathlineto{\pgfqpoint{9.846722in}{0.560883in}}%
\pgfpathlineto{\pgfqpoint{9.846722in}{0.560883in}}%
\pgfpathclose%
\pgfpathmoveto{\pgfqpoint{9.867382in}{0.565170in}}%
\pgfpathlineto{\pgfqpoint{9.867382in}{0.529200in}}%
\pgfpathlineto{\pgfqpoint{9.857025in}{0.529200in}}%
\pgfpathlineto{\pgfqpoint{9.857025in}{0.538764in}}%
\pgfpathquadraticcurveto{\pgfqpoint{9.853477in}{0.533037in}}{\pgfqpoint{9.848181in}{0.530299in}}%
\pgfpathquadraticcurveto{\pgfqpoint{9.842886in}{0.527561in}}{\pgfqpoint{9.835230in}{0.527561in}}%
\pgfpathquadraticcurveto{\pgfqpoint{9.825558in}{0.527561in}}{\pgfqpoint{9.819830in}{0.533001in}}%
\pgfpathquadraticcurveto{\pgfqpoint{9.814120in}{0.538440in}}{\pgfqpoint{9.814120in}{0.547554in}}%
\pgfpathquadraticcurveto{\pgfqpoint{9.814120in}{0.558182in}}{\pgfqpoint{9.821235in}{0.563585in}}%
\pgfpathquadraticcurveto{\pgfqpoint{9.828368in}{0.568989in}}{\pgfqpoint{9.842489in}{0.568989in}}%
\pgfpathlineto{\pgfqpoint{9.857025in}{0.568989in}}%
\pgfpathlineto{\pgfqpoint{9.857025in}{0.570016in}}%
\pgfpathquadraticcurveto{\pgfqpoint{9.857025in}{0.577166in}}{\pgfqpoint{9.852324in}{0.581075in}}%
\pgfpathquadraticcurveto{\pgfqpoint{9.847623in}{0.584984in}}{\pgfqpoint{9.839121in}{0.584984in}}%
\pgfpathquadraticcurveto{\pgfqpoint{9.833717in}{0.584984in}}{\pgfqpoint{9.828584in}{0.583687in}}%
\pgfpathquadraticcurveto{\pgfqpoint{9.823468in}{0.582390in}}{\pgfqpoint{9.818749in}{0.579796in}}%
\pgfpathlineto{\pgfqpoint{9.818749in}{0.589379in}}%
\pgfpathquadraticcurveto{\pgfqpoint{9.824423in}{0.591576in}}{\pgfqpoint{9.829773in}{0.592657in}}%
\pgfpathquadraticcurveto{\pgfqpoint{9.835122in}{0.593756in}}{\pgfqpoint{9.840184in}{0.593756in}}%
\pgfpathquadraticcurveto{\pgfqpoint{9.853873in}{0.593756in}}{\pgfqpoint{9.860628in}{0.586659in}}%
\pgfpathquadraticcurveto{\pgfqpoint{9.867382in}{0.579580in}}{\pgfqpoint{9.867382in}{0.565170in}}%
\pgfpathlineto{\pgfqpoint{9.867382in}{0.565170in}}%
\pgfpathclose%
\pgfpathmoveto{\pgfqpoint{9.888708in}{0.592243in}}%
\pgfpathlineto{\pgfqpoint{9.899065in}{0.592243in}}%
\pgfpathlineto{\pgfqpoint{9.899065in}{0.529200in}}%
\pgfpathlineto{\pgfqpoint{9.888708in}{0.529200in}}%
\pgfpathlineto{\pgfqpoint{9.888708in}{0.592243in}}%
\pgfpathlineto{\pgfqpoint{9.888708in}{0.592243in}}%
\pgfpathclose%
\pgfpathmoveto{\pgfqpoint{9.888708in}{0.616793in}}%
\pgfpathlineto{\pgfqpoint{9.899065in}{0.616793in}}%
\pgfpathlineto{\pgfqpoint{9.899065in}{0.603662in}}%
\pgfpathlineto{\pgfqpoint{9.888708in}{0.603662in}}%
\pgfpathlineto{\pgfqpoint{9.888708in}{0.616793in}}%
\pgfpathlineto{\pgfqpoint{9.888708in}{0.616793in}}%
\pgfpathclose%
\pgfpathmoveto{\pgfqpoint{9.957263in}{0.582570in}}%
\pgfpathquadraticcurveto{\pgfqpoint{9.955516in}{0.583579in}}{\pgfqpoint{9.953462in}{0.584047in}}%
\pgfpathquadraticcurveto{\pgfqpoint{9.951409in}{0.584533in}}{\pgfqpoint{9.948941in}{0.584533in}}%
\pgfpathquadraticcurveto{\pgfqpoint{9.940151in}{0.584533in}}{\pgfqpoint{9.935450in}{0.578823in}}%
\pgfpathquadraticcurveto{\pgfqpoint{9.930749in}{0.573114in}}{\pgfqpoint{9.930749in}{0.562414in}}%
\pgfpathlineto{\pgfqpoint{9.930749in}{0.529200in}}%
\pgfpathlineto{\pgfqpoint{9.920338in}{0.529200in}}%
\pgfpathlineto{\pgfqpoint{9.920338in}{0.592243in}}%
\pgfpathlineto{\pgfqpoint{9.930749in}{0.592243in}}%
\pgfpathlineto{\pgfqpoint{9.930749in}{0.582444in}}%
\pgfpathquadraticcurveto{\pgfqpoint{9.934027in}{0.588190in}}{\pgfqpoint{9.939251in}{0.590964in}}%
\pgfpathquadraticcurveto{\pgfqpoint{9.944492in}{0.593756in}}{\pgfqpoint{9.951985in}{0.593756in}}%
\pgfpathquadraticcurveto{\pgfqpoint{9.953048in}{0.593756in}}{\pgfqpoint{9.954345in}{0.593611in}}%
\pgfpathquadraticcurveto{\pgfqpoint{9.955642in}{0.593485in}}{\pgfqpoint{9.957209in}{0.593197in}}%
\pgfpathlineto{\pgfqpoint{9.957263in}{0.582570in}}%
\pgfpathlineto{\pgfqpoint{9.957263in}{0.582570in}}%
\pgfpathclose%
\pgfpathmoveto{\pgfqpoint{10.019513in}{0.563315in}}%
\pgfpathlineto{\pgfqpoint{10.019513in}{0.558254in}}%
\pgfpathlineto{\pgfqpoint{9.971889in}{0.558254in}}%
\pgfpathquadraticcurveto{\pgfqpoint{9.972573in}{0.547554in}}{\pgfqpoint{9.978337in}{0.541953in}}%
\pgfpathquadraticcurveto{\pgfqpoint{9.984101in}{0.536351in}}{\pgfqpoint{9.994404in}{0.536351in}}%
\pgfpathquadraticcurveto{\pgfqpoint{10.000366in}{0.536351in}}{\pgfqpoint{10.005968in}{0.537810in}}%
\pgfpathquadraticcurveto{\pgfqpoint{10.011569in}{0.539269in}}{\pgfqpoint{10.017099in}{0.542205in}}%
\pgfpathlineto{\pgfqpoint{10.017099in}{0.532406in}}%
\pgfpathquadraticcurveto{\pgfqpoint{10.011515in}{0.530047in}}{\pgfqpoint{10.005661in}{0.528804in}}%
\pgfpathquadraticcurveto{\pgfqpoint{9.999807in}{0.527561in}}{\pgfqpoint{9.993791in}{0.527561in}}%
\pgfpathquadraticcurveto{\pgfqpoint{9.978697in}{0.527561in}}{\pgfqpoint{9.969889in}{0.536333in}}%
\pgfpathquadraticcurveto{\pgfqpoint{9.961081in}{0.545123in}}{\pgfqpoint{9.961081in}{0.560109in}}%
\pgfpathquadraticcurveto{\pgfqpoint{9.961081in}{0.575581in}}{\pgfqpoint{9.969439in}{0.584659in}}%
\pgfpathquadraticcurveto{\pgfqpoint{9.977797in}{0.593756in}}{\pgfqpoint{9.991990in}{0.593756in}}%
\pgfpathquadraticcurveto{\pgfqpoint{10.004707in}{0.593756in}}{\pgfqpoint{10.012110in}{0.585560in}}%
\pgfpathquadraticcurveto{\pgfqpoint{10.019513in}{0.577383in}}{\pgfqpoint{10.019513in}{0.563315in}}%
\pgfpathlineto{\pgfqpoint{10.019513in}{0.563315in}}%
\pgfpathclose%
\pgfpathmoveto{\pgfqpoint{10.009156in}{0.566359in}}%
\pgfpathquadraticcurveto{\pgfqpoint{10.009048in}{0.574843in}}{\pgfqpoint{10.004401in}{0.579904in}}%
\pgfpathquadraticcurveto{\pgfqpoint{9.999753in}{0.584984in}}{\pgfqpoint{9.992098in}{0.584984in}}%
\pgfpathquadraticcurveto{\pgfqpoint{9.983434in}{0.584984in}}{\pgfqpoint{9.978229in}{0.580084in}}%
\pgfpathquadraticcurveto{\pgfqpoint{9.973023in}{0.575185in}}{\pgfqpoint{9.972231in}{0.566287in}}%
\pgfpathlineto{\pgfqpoint{10.009156in}{0.566359in}}%
\pgfpathlineto{\pgfqpoint{10.009156in}{0.566359in}}%
\pgfpathclose%
\pgfpathmoveto{\pgfqpoint{10.077998in}{0.582678in}}%
\pgfpathlineto{\pgfqpoint{10.077998in}{0.616793in}}%
\pgfpathlineto{\pgfqpoint{10.088355in}{0.616793in}}%
\pgfpathlineto{\pgfqpoint{10.088355in}{0.529200in}}%
\pgfpathlineto{\pgfqpoint{10.077998in}{0.529200in}}%
\pgfpathlineto{\pgfqpoint{10.077998in}{0.538656in}}%
\pgfpathquadraticcurveto{\pgfqpoint{10.074738in}{0.533037in}}{\pgfqpoint{10.069749in}{0.530299in}}%
\pgfpathquadraticcurveto{\pgfqpoint{10.064777in}{0.527561in}}{\pgfqpoint{10.057789in}{0.527561in}}%
\pgfpathquadraticcurveto{\pgfqpoint{10.046369in}{0.527561in}}{\pgfqpoint{10.039182in}{0.536675in}}%
\pgfpathquadraticcurveto{\pgfqpoint{10.032013in}{0.545807in}}{\pgfqpoint{10.032013in}{0.560667in}}%
\pgfpathquadraticcurveto{\pgfqpoint{10.032013in}{0.575527in}}{\pgfqpoint{10.039182in}{0.584641in}}%
\pgfpathquadraticcurveto{\pgfqpoint{10.046369in}{0.593756in}}{\pgfqpoint{10.057789in}{0.593756in}}%
\pgfpathquadraticcurveto{\pgfqpoint{10.064777in}{0.593756in}}{\pgfqpoint{10.069749in}{0.591018in}}%
\pgfpathquadraticcurveto{\pgfqpoint{10.074738in}{0.588298in}}{\pgfqpoint{10.077998in}{0.582678in}}%
\pgfpathlineto{\pgfqpoint{10.077998in}{0.582678in}}%
\pgfpathclose%
\pgfpathmoveto{\pgfqpoint{10.042712in}{0.560667in}}%
\pgfpathquadraticcurveto{\pgfqpoint{10.042712in}{0.549248in}}{\pgfqpoint{10.047414in}{0.542745in}}%
\pgfpathquadraticcurveto{\pgfqpoint{10.052115in}{0.536243in}}{\pgfqpoint{10.060328in}{0.536243in}}%
\pgfpathquadraticcurveto{\pgfqpoint{10.068542in}{0.536243in}}{\pgfqpoint{10.073261in}{0.542745in}}%
\pgfpathquadraticcurveto{\pgfqpoint{10.077998in}{0.549248in}}{\pgfqpoint{10.077998in}{0.560667in}}%
\pgfpathquadraticcurveto{\pgfqpoint{10.077998in}{0.572087in}}{\pgfqpoint{10.073261in}{0.578589in}}%
\pgfpathquadraticcurveto{\pgfqpoint{10.068542in}{0.585092in}}{\pgfqpoint{10.060328in}{0.585092in}}%
\pgfpathquadraticcurveto{\pgfqpoint{10.052115in}{0.585092in}}{\pgfqpoint{10.047414in}{0.578589in}}%
\pgfpathquadraticcurveto{\pgfqpoint{10.042712in}{0.572087in}}{\pgfqpoint{10.042712in}{0.560667in}}%
\pgfpathlineto{\pgfqpoint{10.042712in}{0.560667in}}%
\pgfpathclose%
\pgfpathmoveto{\pgfqpoint{10.104476in}{0.565386in}}%
\pgfpathlineto{\pgfqpoint{10.134809in}{0.565386in}}%
\pgfpathlineto{\pgfqpoint{10.134809in}{0.556164in}}%
\pgfpathlineto{\pgfqpoint{10.104476in}{0.556164in}}%
\pgfpathlineto{\pgfqpoint{10.104476in}{0.565386in}}%
\pgfpathlineto{\pgfqpoint{10.104476in}{0.565386in}}%
\pgfpathclose%
\pgfpathmoveto{\pgfqpoint{10.196554in}{0.560667in}}%
\pgfpathquadraticcurveto{\pgfqpoint{10.196554in}{0.572087in}}{\pgfqpoint{10.191853in}{0.578589in}}%
\pgfpathquadraticcurveto{\pgfqpoint{10.187152in}{0.585092in}}{\pgfqpoint{10.178938in}{0.585092in}}%
\pgfpathquadraticcurveto{\pgfqpoint{10.170707in}{0.585092in}}{\pgfqpoint{10.166006in}{0.578589in}}%
\pgfpathquadraticcurveto{\pgfqpoint{10.161304in}{0.572087in}}{\pgfqpoint{10.161304in}{0.560667in}}%
\pgfpathquadraticcurveto{\pgfqpoint{10.161304in}{0.549248in}}{\pgfqpoint{10.166006in}{0.542745in}}%
\pgfpathquadraticcurveto{\pgfqpoint{10.170707in}{0.536243in}}{\pgfqpoint{10.178938in}{0.536243in}}%
\pgfpathquadraticcurveto{\pgfqpoint{10.187152in}{0.536243in}}{\pgfqpoint{10.191853in}{0.542745in}}%
\pgfpathquadraticcurveto{\pgfqpoint{10.196554in}{0.549248in}}{\pgfqpoint{10.196554in}{0.560667in}}%
\pgfpathlineto{\pgfqpoint{10.196554in}{0.560667in}}%
\pgfpathclose%
\pgfpathmoveto{\pgfqpoint{10.161304in}{0.582678in}}%
\pgfpathquadraticcurveto{\pgfqpoint{10.164583in}{0.588298in}}{\pgfqpoint{10.169554in}{0.591018in}}%
\pgfpathquadraticcurveto{\pgfqpoint{10.174543in}{0.593756in}}{\pgfqpoint{10.181460in}{0.593756in}}%
\pgfpathquadraticcurveto{\pgfqpoint{10.192952in}{0.593756in}}{\pgfqpoint{10.200121in}{0.584641in}}%
\pgfpathquadraticcurveto{\pgfqpoint{10.207307in}{0.575527in}}{\pgfqpoint{10.207307in}{0.560667in}}%
\pgfpathquadraticcurveto{\pgfqpoint{10.207307in}{0.545807in}}{\pgfqpoint{10.200121in}{0.536675in}}%
\pgfpathquadraticcurveto{\pgfqpoint{10.192952in}{0.527561in}}{\pgfqpoint{10.181460in}{0.527561in}}%
\pgfpathquadraticcurveto{\pgfqpoint{10.174543in}{0.527561in}}{\pgfqpoint{10.169554in}{0.530299in}}%
\pgfpathquadraticcurveto{\pgfqpoint{10.164583in}{0.533037in}}{\pgfqpoint{10.161304in}{0.538656in}}%
\pgfpathlineto{\pgfqpoint{10.161304in}{0.529200in}}%
\pgfpathlineto{\pgfqpoint{10.150893in}{0.529200in}}%
\pgfpathlineto{\pgfqpoint{10.150893in}{0.616793in}}%
\pgfpathlineto{\pgfqpoint{10.161304in}{0.616793in}}%
\pgfpathlineto{\pgfqpoint{10.161304in}{0.582678in}}%
\pgfpathlineto{\pgfqpoint{10.161304in}{0.582678in}}%
\pgfpathclose%
\pgfpathmoveto{\pgfqpoint{10.248897in}{0.584984in}}%
\pgfpathquadraticcurveto{\pgfqpoint{10.240576in}{0.584984in}}{\pgfqpoint{10.235731in}{0.578481in}}%
\pgfpathquadraticcurveto{\pgfqpoint{10.230885in}{0.571979in}}{\pgfqpoint{10.230885in}{0.560667in}}%
\pgfpathquadraticcurveto{\pgfqpoint{10.230885in}{0.549356in}}{\pgfqpoint{10.235695in}{0.542853in}}%
\pgfpathquadraticcurveto{\pgfqpoint{10.240522in}{0.536351in}}{\pgfqpoint{10.248897in}{0.536351in}}%
\pgfpathquadraticcurveto{\pgfqpoint{10.257183in}{0.536351in}}{\pgfqpoint{10.262010in}{0.542871in}}%
\pgfpathquadraticcurveto{\pgfqpoint{10.266856in}{0.549410in}}{\pgfqpoint{10.266856in}{0.560667in}}%
\pgfpathquadraticcurveto{\pgfqpoint{10.266856in}{0.571871in}}{\pgfqpoint{10.262010in}{0.578427in}}%
\pgfpathquadraticcurveto{\pgfqpoint{10.257183in}{0.584984in}}{\pgfqpoint{10.248897in}{0.584984in}}%
\pgfpathlineto{\pgfqpoint{10.248897in}{0.584984in}}%
\pgfpathclose%
\pgfpathmoveto{\pgfqpoint{10.248897in}{0.593756in}}%
\pgfpathquadraticcurveto{\pgfqpoint{10.262407in}{0.593756in}}{\pgfqpoint{10.270116in}{0.584966in}}%
\pgfpathquadraticcurveto{\pgfqpoint{10.277843in}{0.576194in}}{\pgfqpoint{10.277843in}{0.560667in}}%
\pgfpathquadraticcurveto{\pgfqpoint{10.277843in}{0.545195in}}{\pgfqpoint{10.270116in}{0.536369in}}%
\pgfpathquadraticcurveto{\pgfqpoint{10.262407in}{0.527561in}}{\pgfqpoint{10.248897in}{0.527561in}}%
\pgfpathquadraticcurveto{\pgfqpoint{10.235334in}{0.527561in}}{\pgfqpoint{10.227643in}{0.536369in}}%
\pgfpathquadraticcurveto{\pgfqpoint{10.219970in}{0.545195in}}{\pgfqpoint{10.219970in}{0.560667in}}%
\pgfpathquadraticcurveto{\pgfqpoint{10.219970in}{0.576194in}}{\pgfqpoint{10.227643in}{0.584966in}}%
\pgfpathquadraticcurveto{\pgfqpoint{10.235334in}{0.593756in}}{\pgfqpoint{10.248897in}{0.593756in}}%
\pgfpathlineto{\pgfqpoint{10.248897in}{0.593756in}}%
\pgfpathclose%
\pgfpathmoveto{\pgfqpoint{10.319433in}{0.584984in}}%
\pgfpathquadraticcurveto{\pgfqpoint{10.311111in}{0.584984in}}{\pgfqpoint{10.306266in}{0.578481in}}%
\pgfpathquadraticcurveto{\pgfqpoint{10.301421in}{0.571979in}}{\pgfqpoint{10.301421in}{0.560667in}}%
\pgfpathquadraticcurveto{\pgfqpoint{10.301421in}{0.549356in}}{\pgfqpoint{10.306230in}{0.542853in}}%
\pgfpathquadraticcurveto{\pgfqpoint{10.311057in}{0.536351in}}{\pgfqpoint{10.319433in}{0.536351in}}%
\pgfpathquadraticcurveto{\pgfqpoint{10.327719in}{0.536351in}}{\pgfqpoint{10.332546in}{0.542871in}}%
\pgfpathquadraticcurveto{\pgfqpoint{10.337391in}{0.549410in}}{\pgfqpoint{10.337391in}{0.560667in}}%
\pgfpathquadraticcurveto{\pgfqpoint{10.337391in}{0.571871in}}{\pgfqpoint{10.332546in}{0.578427in}}%
\pgfpathquadraticcurveto{\pgfqpoint{10.327719in}{0.584984in}}{\pgfqpoint{10.319433in}{0.584984in}}%
\pgfpathlineto{\pgfqpoint{10.319433in}{0.584984in}}%
\pgfpathclose%
\pgfpathmoveto{\pgfqpoint{10.319433in}{0.593756in}}%
\pgfpathquadraticcurveto{\pgfqpoint{10.332942in}{0.593756in}}{\pgfqpoint{10.340651in}{0.584966in}}%
\pgfpathquadraticcurveto{\pgfqpoint{10.348379in}{0.576194in}}{\pgfqpoint{10.348379in}{0.560667in}}%
\pgfpathquadraticcurveto{\pgfqpoint{10.348379in}{0.545195in}}{\pgfqpoint{10.340651in}{0.536369in}}%
\pgfpathquadraticcurveto{\pgfqpoint{10.332942in}{0.527561in}}{\pgfqpoint{10.319433in}{0.527561in}}%
\pgfpathquadraticcurveto{\pgfqpoint{10.305870in}{0.527561in}}{\pgfqpoint{10.298179in}{0.536369in}}%
\pgfpathquadraticcurveto{\pgfqpoint{10.290506in}{0.545195in}}{\pgfqpoint{10.290506in}{0.560667in}}%
\pgfpathquadraticcurveto{\pgfqpoint{10.290506in}{0.576194in}}{\pgfqpoint{10.298179in}{0.584966in}}%
\pgfpathquadraticcurveto{\pgfqpoint{10.305870in}{0.593756in}}{\pgfqpoint{10.319433in}{0.593756in}}%
\pgfpathlineto{\pgfqpoint{10.319433in}{0.593756in}}%
\pgfpathclose%
\pgfpathmoveto{\pgfqpoint{10.375793in}{0.610147in}}%
\pgfpathlineto{\pgfqpoint{10.375793in}{0.592243in}}%
\pgfpathlineto{\pgfqpoint{10.397119in}{0.592243in}}%
\pgfpathlineto{\pgfqpoint{10.397119in}{0.584191in}}%
\pgfpathlineto{\pgfqpoint{10.375793in}{0.584191in}}%
\pgfpathlineto{\pgfqpoint{10.375793in}{0.549968in}}%
\pgfpathquadraticcurveto{\pgfqpoint{10.375793in}{0.542259in}}{\pgfqpoint{10.377901in}{0.540061in}}%
\pgfpathquadraticcurveto{\pgfqpoint{10.380008in}{0.537864in}}{\pgfqpoint{10.386492in}{0.537864in}}%
\pgfpathlineto{\pgfqpoint{10.397119in}{0.537864in}}%
\pgfpathlineto{\pgfqpoint{10.397119in}{0.529200in}}%
\pgfpathlineto{\pgfqpoint{10.386492in}{0.529200in}}%
\pgfpathquadraticcurveto{\pgfqpoint{10.374496in}{0.529200in}}{\pgfqpoint{10.369939in}{0.533667in}}%
\pgfpathquadraticcurveto{\pgfqpoint{10.365382in}{0.538152in}}{\pgfqpoint{10.365382in}{0.549968in}}%
\pgfpathlineto{\pgfqpoint{10.365382in}{0.584191in}}%
\pgfpathlineto{\pgfqpoint{10.357781in}{0.584191in}}%
\pgfpathlineto{\pgfqpoint{10.357781in}{0.592243in}}%
\pgfpathlineto{\pgfqpoint{10.365382in}{0.592243in}}%
\pgfpathlineto{\pgfqpoint{10.365382in}{0.610147in}}%
\pgfpathlineto{\pgfqpoint{10.375793in}{0.610147in}}%
\pgfpathlineto{\pgfqpoint{10.375793in}{0.610147in}}%
\pgfpathclose%
\pgfpathmoveto{\pgfqpoint{10.450922in}{0.590387in}}%
\pgfpathlineto{\pgfqpoint{10.450922in}{0.580589in}}%
\pgfpathquadraticcurveto{\pgfqpoint{10.446545in}{0.582840in}}{\pgfqpoint{10.441808in}{0.583957in}}%
\pgfpathquadraticcurveto{\pgfqpoint{10.437088in}{0.585092in}}{\pgfqpoint{10.432009in}{0.585092in}}%
\pgfpathquadraticcurveto{\pgfqpoint{10.424300in}{0.585092in}}{\pgfqpoint{10.420445in}{0.582732in}}%
\pgfpathquadraticcurveto{\pgfqpoint{10.416591in}{0.580373in}}{\pgfqpoint{10.416591in}{0.575635in}}%
\pgfpathquadraticcurveto{\pgfqpoint{10.416591in}{0.572033in}}{\pgfqpoint{10.419346in}{0.569980in}}%
\pgfpathquadraticcurveto{\pgfqpoint{10.422102in}{0.567926in}}{\pgfqpoint{10.430442in}{0.566071in}}%
\pgfpathlineto{\pgfqpoint{10.433990in}{0.565278in}}%
\pgfpathquadraticcurveto{\pgfqpoint{10.445014in}{0.562919in}}{\pgfqpoint{10.449661in}{0.558614in}}%
\pgfpathquadraticcurveto{\pgfqpoint{10.454308in}{0.554309in}}{\pgfqpoint{10.454308in}{0.546600in}}%
\pgfpathquadraticcurveto{\pgfqpoint{10.454308in}{0.537810in}}{\pgfqpoint{10.447355in}{0.532676in}}%
\pgfpathquadraticcurveto{\pgfqpoint{10.440403in}{0.527561in}}{\pgfqpoint{10.428244in}{0.527561in}}%
\pgfpathquadraticcurveto{\pgfqpoint{10.423183in}{0.527561in}}{\pgfqpoint{10.417689in}{0.528552in}}%
\pgfpathquadraticcurveto{\pgfqpoint{10.412196in}{0.529542in}}{\pgfqpoint{10.406126in}{0.531506in}}%
\pgfpathlineto{\pgfqpoint{10.406126in}{0.542205in}}%
\pgfpathquadraticcurveto{\pgfqpoint{10.411871in}{0.539215in}}{\pgfqpoint{10.417437in}{0.537720in}}%
\pgfpathquadraticcurveto{\pgfqpoint{10.423003in}{0.536243in}}{\pgfqpoint{10.428479in}{0.536243in}}%
\pgfpathquadraticcurveto{\pgfqpoint{10.435792in}{0.536243in}}{\pgfqpoint{10.439718in}{0.538746in}}%
\pgfpathquadraticcurveto{\pgfqpoint{10.443663in}{0.541250in}}{\pgfqpoint{10.443663in}{0.545807in}}%
\pgfpathquadraticcurveto{\pgfqpoint{10.443663in}{0.550022in}}{\pgfqpoint{10.440817in}{0.552274in}}%
\pgfpathquadraticcurveto{\pgfqpoint{10.437989in}{0.554525in}}{\pgfqpoint{10.428353in}{0.556614in}}%
\pgfpathlineto{\pgfqpoint{10.424750in}{0.557461in}}%
\pgfpathquadraticcurveto{\pgfqpoint{10.415132in}{0.559478in}}{\pgfqpoint{10.410845in}{0.563675in}}%
\pgfpathquadraticcurveto{\pgfqpoint{10.406576in}{0.567872in}}{\pgfqpoint{10.406576in}{0.575185in}}%
\pgfpathquadraticcurveto{\pgfqpoint{10.406576in}{0.584083in}}{\pgfqpoint{10.412880in}{0.588910in}}%
\pgfpathquadraticcurveto{\pgfqpoint{10.419184in}{0.593756in}}{\pgfqpoint{10.430784in}{0.593756in}}%
\pgfpathquadraticcurveto{\pgfqpoint{10.436512in}{0.593756in}}{\pgfqpoint{10.441573in}{0.592909in}}%
\pgfpathquadraticcurveto{\pgfqpoint{10.446653in}{0.592080in}}{\pgfqpoint{10.450922in}{0.590387in}}%
\pgfpathlineto{\pgfqpoint{10.450922in}{0.590387in}}%
\pgfpathclose%
\pgfpathmoveto{\pgfqpoint{10.481038in}{0.610147in}}%
\pgfpathlineto{\pgfqpoint{10.481038in}{0.592243in}}%
\pgfpathlineto{\pgfqpoint{10.502364in}{0.592243in}}%
\pgfpathlineto{\pgfqpoint{10.502364in}{0.584191in}}%
\pgfpathlineto{\pgfqpoint{10.481038in}{0.584191in}}%
\pgfpathlineto{\pgfqpoint{10.481038in}{0.549968in}}%
\pgfpathquadraticcurveto{\pgfqpoint{10.481038in}{0.542259in}}{\pgfqpoint{10.483146in}{0.540061in}}%
\pgfpathquadraticcurveto{\pgfqpoint{10.485253in}{0.537864in}}{\pgfqpoint{10.491737in}{0.537864in}}%
\pgfpathlineto{\pgfqpoint{10.502364in}{0.537864in}}%
\pgfpathlineto{\pgfqpoint{10.502364in}{0.529200in}}%
\pgfpathlineto{\pgfqpoint{10.491737in}{0.529200in}}%
\pgfpathquadraticcurveto{\pgfqpoint{10.479741in}{0.529200in}}{\pgfqpoint{10.475184in}{0.533667in}}%
\pgfpathquadraticcurveto{\pgfqpoint{10.470627in}{0.538152in}}{\pgfqpoint{10.470627in}{0.549968in}}%
\pgfpathlineto{\pgfqpoint{10.470627in}{0.584191in}}%
\pgfpathlineto{\pgfqpoint{10.463026in}{0.584191in}}%
\pgfpathlineto{\pgfqpoint{10.463026in}{0.592243in}}%
\pgfpathlineto{\pgfqpoint{10.470627in}{0.592243in}}%
\pgfpathlineto{\pgfqpoint{10.470627in}{0.610147in}}%
\pgfpathlineto{\pgfqpoint{10.481038in}{0.610147in}}%
\pgfpathlineto{\pgfqpoint{10.481038in}{0.610147in}}%
\pgfpathclose%
\pgfpathmoveto{\pgfqpoint{10.552510in}{0.582570in}}%
\pgfpathquadraticcurveto{\pgfqpoint{10.550763in}{0.583579in}}{\pgfqpoint{10.548710in}{0.584047in}}%
\pgfpathquadraticcurveto{\pgfqpoint{10.546656in}{0.584533in}}{\pgfqpoint{10.544189in}{0.584533in}}%
\pgfpathquadraticcurveto{\pgfqpoint{10.535399in}{0.584533in}}{\pgfqpoint{10.530698in}{0.578823in}}%
\pgfpathquadraticcurveto{\pgfqpoint{10.525996in}{0.573114in}}{\pgfqpoint{10.525996in}{0.562414in}}%
\pgfpathlineto{\pgfqpoint{10.525996in}{0.529200in}}%
\pgfpathlineto{\pgfqpoint{10.515585in}{0.529200in}}%
\pgfpathlineto{\pgfqpoint{10.515585in}{0.592243in}}%
\pgfpathlineto{\pgfqpoint{10.525996in}{0.592243in}}%
\pgfpathlineto{\pgfqpoint{10.525996in}{0.582444in}}%
\pgfpathquadraticcurveto{\pgfqpoint{10.529275in}{0.588190in}}{\pgfqpoint{10.534498in}{0.590964in}}%
\pgfpathquadraticcurveto{\pgfqpoint{10.539740in}{0.593756in}}{\pgfqpoint{10.547233in}{0.593756in}}%
\pgfpathquadraticcurveto{\pgfqpoint{10.548295in}{0.593756in}}{\pgfqpoint{10.549592in}{0.593611in}}%
\pgfpathquadraticcurveto{\pgfqpoint{10.550889in}{0.593485in}}{\pgfqpoint{10.552456in}{0.593197in}}%
\pgfpathlineto{\pgfqpoint{10.552510in}{0.582570in}}%
\pgfpathlineto{\pgfqpoint{10.552510in}{0.582570in}}%
\pgfpathclose%
\pgfpathmoveto{\pgfqpoint{10.592029in}{0.560883in}}%
\pgfpathquadraticcurveto{\pgfqpoint{10.579475in}{0.560883in}}{\pgfqpoint{10.574629in}{0.558019in}}%
\pgfpathquadraticcurveto{\pgfqpoint{10.569784in}{0.555156in}}{\pgfqpoint{10.569784in}{0.548221in}}%
\pgfpathquadraticcurveto{\pgfqpoint{10.569784in}{0.542709in}}{\pgfqpoint{10.573422in}{0.539467in}}%
\pgfpathquadraticcurveto{\pgfqpoint{10.577061in}{0.536243in}}{\pgfqpoint{10.583293in}{0.536243in}}%
\pgfpathquadraticcurveto{\pgfqpoint{10.591921in}{0.536243in}}{\pgfqpoint{10.597126in}{0.542349in}}%
\pgfpathquadraticcurveto{\pgfqpoint{10.602332in}{0.548455in}}{\pgfqpoint{10.602332in}{0.558578in}}%
\pgfpathlineto{\pgfqpoint{10.602332in}{0.560883in}}%
\pgfpathlineto{\pgfqpoint{10.592029in}{0.560883in}}%
\pgfpathlineto{\pgfqpoint{10.592029in}{0.560883in}}%
\pgfpathclose%
\pgfpathmoveto{\pgfqpoint{10.612689in}{0.565170in}}%
\pgfpathlineto{\pgfqpoint{10.612689in}{0.529200in}}%
\pgfpathlineto{\pgfqpoint{10.602332in}{0.529200in}}%
\pgfpathlineto{\pgfqpoint{10.602332in}{0.538764in}}%
\pgfpathquadraticcurveto{\pgfqpoint{10.598784in}{0.533037in}}{\pgfqpoint{10.593488in}{0.530299in}}%
\pgfpathquadraticcurveto{\pgfqpoint{10.588192in}{0.527561in}}{\pgfqpoint{10.580537in}{0.527561in}}%
\pgfpathquadraticcurveto{\pgfqpoint{10.570865in}{0.527561in}}{\pgfqpoint{10.565137in}{0.533001in}}%
\pgfpathquadraticcurveto{\pgfqpoint{10.559427in}{0.538440in}}{\pgfqpoint{10.559427in}{0.547554in}}%
\pgfpathquadraticcurveto{\pgfqpoint{10.559427in}{0.558182in}}{\pgfqpoint{10.566542in}{0.563585in}}%
\pgfpathquadraticcurveto{\pgfqpoint{10.573675in}{0.568989in}}{\pgfqpoint{10.587796in}{0.568989in}}%
\pgfpathlineto{\pgfqpoint{10.602332in}{0.568989in}}%
\pgfpathlineto{\pgfqpoint{10.602332in}{0.570016in}}%
\pgfpathquadraticcurveto{\pgfqpoint{10.602332in}{0.577166in}}{\pgfqpoint{10.597631in}{0.581075in}}%
\pgfpathquadraticcurveto{\pgfqpoint{10.592930in}{0.584984in}}{\pgfqpoint{10.584428in}{0.584984in}}%
\pgfpathquadraticcurveto{\pgfqpoint{10.579024in}{0.584984in}}{\pgfqpoint{10.573891in}{0.583687in}}%
\pgfpathquadraticcurveto{\pgfqpoint{10.568775in}{0.582390in}}{\pgfqpoint{10.564056in}{0.579796in}}%
\pgfpathlineto{\pgfqpoint{10.564056in}{0.589379in}}%
\pgfpathquadraticcurveto{\pgfqpoint{10.569730in}{0.591576in}}{\pgfqpoint{10.575080in}{0.592657in}}%
\pgfpathquadraticcurveto{\pgfqpoint{10.580429in}{0.593756in}}{\pgfqpoint{10.585491in}{0.593756in}}%
\pgfpathquadraticcurveto{\pgfqpoint{10.599180in}{0.593756in}}{\pgfqpoint{10.605934in}{0.586659in}}%
\pgfpathquadraticcurveto{\pgfqpoint{10.612689in}{0.579580in}}{\pgfqpoint{10.612689in}{0.565170in}}%
\pgfpathlineto{\pgfqpoint{10.612689in}{0.565170in}}%
\pgfpathclose%
\pgfpathmoveto{\pgfqpoint{10.644030in}{0.538656in}}%
\pgfpathlineto{\pgfqpoint{10.644030in}{0.505226in}}%
\pgfpathlineto{\pgfqpoint{10.633619in}{0.505226in}}%
\pgfpathlineto{\pgfqpoint{10.633619in}{0.592243in}}%
\pgfpathlineto{\pgfqpoint{10.644030in}{0.592243in}}%
\pgfpathlineto{\pgfqpoint{10.644030in}{0.582678in}}%
\pgfpathquadraticcurveto{\pgfqpoint{10.647308in}{0.588298in}}{\pgfqpoint{10.652280in}{0.591018in}}%
\pgfpathquadraticcurveto{\pgfqpoint{10.657269in}{0.593756in}}{\pgfqpoint{10.664186in}{0.593756in}}%
\pgfpathquadraticcurveto{\pgfqpoint{10.675677in}{0.593756in}}{\pgfqpoint{10.682846in}{0.584641in}}%
\pgfpathquadraticcurveto{\pgfqpoint{10.690033in}{0.575527in}}{\pgfqpoint{10.690033in}{0.560667in}}%
\pgfpathquadraticcurveto{\pgfqpoint{10.690033in}{0.545807in}}{\pgfqpoint{10.682846in}{0.536675in}}%
\pgfpathquadraticcurveto{\pgfqpoint{10.675677in}{0.527561in}}{\pgfqpoint{10.664186in}{0.527561in}}%
\pgfpathquadraticcurveto{\pgfqpoint{10.657269in}{0.527561in}}{\pgfqpoint{10.652280in}{0.530299in}}%
\pgfpathquadraticcurveto{\pgfqpoint{10.647308in}{0.533037in}}{\pgfqpoint{10.644030in}{0.538656in}}%
\pgfpathlineto{\pgfqpoint{10.644030in}{0.538656in}}%
\pgfpathclose%
\pgfpathmoveto{\pgfqpoint{10.679280in}{0.560667in}}%
\pgfpathquadraticcurveto{\pgfqpoint{10.679280in}{0.572087in}}{\pgfqpoint{10.674579in}{0.578589in}}%
\pgfpathquadraticcurveto{\pgfqpoint{10.669878in}{0.585092in}}{\pgfqpoint{10.661664in}{0.585092in}}%
\pgfpathquadraticcurveto{\pgfqpoint{10.653432in}{0.585092in}}{\pgfqpoint{10.648731in}{0.578589in}}%
\pgfpathquadraticcurveto{\pgfqpoint{10.644030in}{0.572087in}}{\pgfqpoint{10.644030in}{0.560667in}}%
\pgfpathquadraticcurveto{\pgfqpoint{10.644030in}{0.549248in}}{\pgfqpoint{10.648731in}{0.542745in}}%
\pgfpathquadraticcurveto{\pgfqpoint{10.653432in}{0.536243in}}{\pgfqpoint{10.661664in}{0.536243in}}%
\pgfpathquadraticcurveto{\pgfqpoint{10.669878in}{0.536243in}}{\pgfqpoint{10.674579in}{0.542745in}}%
\pgfpathquadraticcurveto{\pgfqpoint{10.679280in}{0.549248in}}{\pgfqpoint{10.679280in}{0.560667in}}%
\pgfpathlineto{\pgfqpoint{10.679280in}{0.560667in}}%
\pgfpathclose%
\pgfusepath{fill}%
\end{pgfscope}%
\begin{pgfscope}%
\pgfsetbuttcap%
\pgfsetroundjoin%
\definecolor{currentfill}{rgb}{0.309804,0.372549,0.419608}%
\pgfsetfillcolor{currentfill}%
\pgfsetlinewidth{0.000000pt}%
\definecolor{currentstroke}{rgb}{0.309804,0.372549,0.419608}%
\pgfsetstrokecolor{currentstroke}%
\pgfsetdash{}{0pt}%
\pgfpathmoveto{\pgfqpoint{10.851299in}{0.567260in}}%
\pgfpathlineto{\pgfqpoint{10.851299in}{0.529200in}}%
\pgfpathlineto{\pgfqpoint{10.840942in}{0.529200in}}%
\pgfpathlineto{\pgfqpoint{10.840942in}{0.566917in}}%
\pgfpathquadraticcurveto{\pgfqpoint{10.840942in}{0.575869in}}{\pgfqpoint{10.837448in}{0.580300in}}%
\pgfpathquadraticcurveto{\pgfqpoint{10.833953in}{0.584749in}}{\pgfqpoint{10.826983in}{0.584749in}}%
\pgfpathquadraticcurveto{\pgfqpoint{10.818589in}{0.584749in}}{\pgfqpoint{10.813744in}{0.579400in}}%
\pgfpathquadraticcurveto{\pgfqpoint{10.808898in}{0.574068in}}{\pgfqpoint{10.808898in}{0.564828in}}%
\pgfpathlineto{\pgfqpoint{10.808898in}{0.529200in}}%
\pgfpathlineto{\pgfqpoint{10.798487in}{0.529200in}}%
\pgfpathlineto{\pgfqpoint{10.798487in}{0.592243in}}%
\pgfpathlineto{\pgfqpoint{10.808898in}{0.592243in}}%
\pgfpathlineto{\pgfqpoint{10.808898in}{0.582444in}}%
\pgfpathquadraticcurveto{\pgfqpoint{10.812627in}{0.588136in}}{\pgfqpoint{10.817652in}{0.590946in}}%
\pgfpathquadraticcurveto{\pgfqpoint{10.822696in}{0.593756in}}{\pgfqpoint{10.829288in}{0.593756in}}%
\pgfpathquadraticcurveto{\pgfqpoint{10.840149in}{0.593756in}}{\pgfqpoint{10.845715in}{0.587037in}}%
\pgfpathquadraticcurveto{\pgfqpoint{10.851299in}{0.580318in}}{\pgfqpoint{10.851299in}{0.567260in}}%
\pgfpathlineto{\pgfqpoint{10.851299in}{0.567260in}}%
\pgfpathclose%
\pgfpathmoveto{\pgfqpoint{10.896365in}{0.584984in}}%
\pgfpathquadraticcurveto{\pgfqpoint{10.888044in}{0.584984in}}{\pgfqpoint{10.883199in}{0.578481in}}%
\pgfpathquadraticcurveto{\pgfqpoint{10.878353in}{0.571979in}}{\pgfqpoint{10.878353in}{0.560667in}}%
\pgfpathquadraticcurveto{\pgfqpoint{10.878353in}{0.549356in}}{\pgfqpoint{10.883163in}{0.542853in}}%
\pgfpathquadraticcurveto{\pgfqpoint{10.887990in}{0.536351in}}{\pgfqpoint{10.896365in}{0.536351in}}%
\pgfpathquadraticcurveto{\pgfqpoint{10.904651in}{0.536351in}}{\pgfqpoint{10.909478in}{0.542871in}}%
\pgfpathquadraticcurveto{\pgfqpoint{10.914324in}{0.549410in}}{\pgfqpoint{10.914324in}{0.560667in}}%
\pgfpathquadraticcurveto{\pgfqpoint{10.914324in}{0.571871in}}{\pgfqpoint{10.909478in}{0.578427in}}%
\pgfpathquadraticcurveto{\pgfqpoint{10.904651in}{0.584984in}}{\pgfqpoint{10.896365in}{0.584984in}}%
\pgfpathlineto{\pgfqpoint{10.896365in}{0.584984in}}%
\pgfpathclose%
\pgfpathmoveto{\pgfqpoint{10.896365in}{0.593756in}}%
\pgfpathquadraticcurveto{\pgfqpoint{10.909875in}{0.593756in}}{\pgfqpoint{10.917584in}{0.584966in}}%
\pgfpathquadraticcurveto{\pgfqpoint{10.925311in}{0.576194in}}{\pgfqpoint{10.925311in}{0.560667in}}%
\pgfpathquadraticcurveto{\pgfqpoint{10.925311in}{0.545195in}}{\pgfqpoint{10.917584in}{0.536369in}}%
\pgfpathquadraticcurveto{\pgfqpoint{10.909875in}{0.527561in}}{\pgfqpoint{10.896365in}{0.527561in}}%
\pgfpathquadraticcurveto{\pgfqpoint{10.882802in}{0.527561in}}{\pgfqpoint{10.875111in}{0.536369in}}%
\pgfpathquadraticcurveto{\pgfqpoint{10.867438in}{0.545195in}}{\pgfqpoint{10.867438in}{0.560667in}}%
\pgfpathquadraticcurveto{\pgfqpoint{10.867438in}{0.576194in}}{\pgfqpoint{10.875111in}{0.584966in}}%
\pgfpathquadraticcurveto{\pgfqpoint{10.882802in}{0.593756in}}{\pgfqpoint{10.896365in}{0.593756in}}%
\pgfpathlineto{\pgfqpoint{10.896365in}{0.593756in}}%
\pgfpathclose%
\pgfpathmoveto{\pgfqpoint{10.979005in}{0.582570in}}%
\pgfpathquadraticcurveto{\pgfqpoint{10.977258in}{0.583579in}}{\pgfqpoint{10.975205in}{0.584047in}}%
\pgfpathquadraticcurveto{\pgfqpoint{10.973151in}{0.584533in}}{\pgfqpoint{10.970684in}{0.584533in}}%
\pgfpathquadraticcurveto{\pgfqpoint{10.961894in}{0.584533in}}{\pgfqpoint{10.957192in}{0.578823in}}%
\pgfpathquadraticcurveto{\pgfqpoint{10.952491in}{0.573114in}}{\pgfqpoint{10.952491in}{0.562414in}}%
\pgfpathlineto{\pgfqpoint{10.952491in}{0.529200in}}%
\pgfpathlineto{\pgfqpoint{10.942080in}{0.529200in}}%
\pgfpathlineto{\pgfqpoint{10.942080in}{0.592243in}}%
\pgfpathlineto{\pgfqpoint{10.952491in}{0.592243in}}%
\pgfpathlineto{\pgfqpoint{10.952491in}{0.582444in}}%
\pgfpathquadraticcurveto{\pgfqpoint{10.955770in}{0.588190in}}{\pgfqpoint{10.960993in}{0.590964in}}%
\pgfpathquadraticcurveto{\pgfqpoint{10.966235in}{0.593756in}}{\pgfqpoint{10.973728in}{0.593756in}}%
\pgfpathquadraticcurveto{\pgfqpoint{10.974790in}{0.593756in}}{\pgfqpoint{10.976087in}{0.593611in}}%
\pgfpathquadraticcurveto{\pgfqpoint{10.977384in}{0.593485in}}{\pgfqpoint{10.978951in}{0.593197in}}%
\pgfpathlineto{\pgfqpoint{10.979005in}{0.582570in}}%
\pgfpathlineto{\pgfqpoint{10.979005in}{0.582570in}}%
\pgfpathclose%
\pgfpathmoveto{\pgfqpoint{11.036932in}{0.580138in}}%
\pgfpathquadraticcurveto{\pgfqpoint{11.040823in}{0.587127in}}{\pgfqpoint{11.046227in}{0.590441in}}%
\pgfpathquadraticcurveto{\pgfqpoint{11.051630in}{0.593756in}}{\pgfqpoint{11.058943in}{0.593756in}}%
\pgfpathquadraticcurveto{\pgfqpoint{11.068796in}{0.593756in}}{\pgfqpoint{11.074145in}{0.586857in}}%
\pgfpathquadraticcurveto{\pgfqpoint{11.079495in}{0.579976in}}{\pgfqpoint{11.079495in}{0.567260in}}%
\pgfpathlineto{\pgfqpoint{11.079495in}{0.529200in}}%
\pgfpathlineto{\pgfqpoint{11.069084in}{0.529200in}}%
\pgfpathlineto{\pgfqpoint{11.069084in}{0.566917in}}%
\pgfpathquadraticcurveto{\pgfqpoint{11.069084in}{0.575978in}}{\pgfqpoint{11.065860in}{0.580355in}}%
\pgfpathquadraticcurveto{\pgfqpoint{11.062654in}{0.584749in}}{\pgfqpoint{11.056079in}{0.584749in}}%
\pgfpathquadraticcurveto{\pgfqpoint{11.048028in}{0.584749in}}{\pgfqpoint{11.043345in}{0.579400in}}%
\pgfpathquadraticcurveto{\pgfqpoint{11.038679in}{0.574068in}}{\pgfqpoint{11.038679in}{0.564828in}}%
\pgfpathlineto{\pgfqpoint{11.038679in}{0.529200in}}%
\pgfpathlineto{\pgfqpoint{11.028268in}{0.529200in}}%
\pgfpathlineto{\pgfqpoint{11.028268in}{0.566917in}}%
\pgfpathquadraticcurveto{\pgfqpoint{11.028268in}{0.576032in}}{\pgfqpoint{11.025062in}{0.580391in}}%
\pgfpathquadraticcurveto{\pgfqpoint{11.021856in}{0.584749in}}{\pgfqpoint{11.015156in}{0.584749in}}%
\pgfpathquadraticcurveto{\pgfqpoint{11.007212in}{0.584749in}}{\pgfqpoint{11.002529in}{0.579382in}}%
\pgfpathquadraticcurveto{\pgfqpoint{10.997864in}{0.574014in}}{\pgfqpoint{10.997864in}{0.564828in}}%
\pgfpathlineto{\pgfqpoint{10.997864in}{0.529200in}}%
\pgfpathlineto{\pgfqpoint{10.987453in}{0.529200in}}%
\pgfpathlineto{\pgfqpoint{10.987453in}{0.592243in}}%
\pgfpathlineto{\pgfqpoint{10.997864in}{0.592243in}}%
\pgfpathlineto{\pgfqpoint{10.997864in}{0.582444in}}%
\pgfpathquadraticcurveto{\pgfqpoint{11.001412in}{0.588244in}}{\pgfqpoint{11.006366in}{0.591000in}}%
\pgfpathquadraticcurveto{\pgfqpoint{11.011319in}{0.593756in}}{\pgfqpoint{11.018128in}{0.593756in}}%
\pgfpathquadraticcurveto{\pgfqpoint{11.025008in}{0.593756in}}{\pgfqpoint{11.029817in}{0.590261in}}%
\pgfpathquadraticcurveto{\pgfqpoint{11.034627in}{0.586785in}}{\pgfqpoint{11.036932in}{0.580138in}}%
\pgfpathlineto{\pgfqpoint{11.036932in}{0.580138in}}%
\pgfpathclose%
\pgfpathmoveto{\pgfqpoint{11.128794in}{0.560883in}}%
\pgfpathquadraticcurveto{\pgfqpoint{11.116240in}{0.560883in}}{\pgfqpoint{11.111395in}{0.558019in}}%
\pgfpathquadraticcurveto{\pgfqpoint{11.106549in}{0.555156in}}{\pgfqpoint{11.106549in}{0.548221in}}%
\pgfpathquadraticcurveto{\pgfqpoint{11.106549in}{0.542709in}}{\pgfqpoint{11.110188in}{0.539467in}}%
\pgfpathquadraticcurveto{\pgfqpoint{11.113826in}{0.536243in}}{\pgfqpoint{11.120058in}{0.536243in}}%
\pgfpathquadraticcurveto{\pgfqpoint{11.128686in}{0.536243in}}{\pgfqpoint{11.133892in}{0.542349in}}%
\pgfpathquadraticcurveto{\pgfqpoint{11.139097in}{0.548455in}}{\pgfqpoint{11.139097in}{0.558578in}}%
\pgfpathlineto{\pgfqpoint{11.139097in}{0.560883in}}%
\pgfpathlineto{\pgfqpoint{11.128794in}{0.560883in}}%
\pgfpathlineto{\pgfqpoint{11.128794in}{0.560883in}}%
\pgfpathclose%
\pgfpathmoveto{\pgfqpoint{11.149454in}{0.565170in}}%
\pgfpathlineto{\pgfqpoint{11.149454in}{0.529200in}}%
\pgfpathlineto{\pgfqpoint{11.139097in}{0.529200in}}%
\pgfpathlineto{\pgfqpoint{11.139097in}{0.538764in}}%
\pgfpathquadraticcurveto{\pgfqpoint{11.135549in}{0.533037in}}{\pgfqpoint{11.130253in}{0.530299in}}%
\pgfpathquadraticcurveto{\pgfqpoint{11.124958in}{0.527561in}}{\pgfqpoint{11.117302in}{0.527561in}}%
\pgfpathquadraticcurveto{\pgfqpoint{11.107630in}{0.527561in}}{\pgfqpoint{11.101902in}{0.533001in}}%
\pgfpathquadraticcurveto{\pgfqpoint{11.096192in}{0.538440in}}{\pgfqpoint{11.096192in}{0.547554in}}%
\pgfpathquadraticcurveto{\pgfqpoint{11.096192in}{0.558182in}}{\pgfqpoint{11.103307in}{0.563585in}}%
\pgfpathquadraticcurveto{\pgfqpoint{11.110440in}{0.568989in}}{\pgfqpoint{11.124561in}{0.568989in}}%
\pgfpathlineto{\pgfqpoint{11.139097in}{0.568989in}}%
\pgfpathlineto{\pgfqpoint{11.139097in}{0.570016in}}%
\pgfpathquadraticcurveto{\pgfqpoint{11.139097in}{0.577166in}}{\pgfqpoint{11.134396in}{0.581075in}}%
\pgfpathquadraticcurveto{\pgfqpoint{11.129695in}{0.584984in}}{\pgfqpoint{11.121193in}{0.584984in}}%
\pgfpathquadraticcurveto{\pgfqpoint{11.115789in}{0.584984in}}{\pgfqpoint{11.110656in}{0.583687in}}%
\pgfpathquadraticcurveto{\pgfqpoint{11.105541in}{0.582390in}}{\pgfqpoint{11.100821in}{0.579796in}}%
\pgfpathlineto{\pgfqpoint{11.100821in}{0.589379in}}%
\pgfpathquadraticcurveto{\pgfqpoint{11.106495in}{0.591576in}}{\pgfqpoint{11.111845in}{0.592657in}}%
\pgfpathquadraticcurveto{\pgfqpoint{11.117194in}{0.593756in}}{\pgfqpoint{11.122256in}{0.593756in}}%
\pgfpathquadraticcurveto{\pgfqpoint{11.135945in}{0.593756in}}{\pgfqpoint{11.142700in}{0.586659in}}%
\pgfpathquadraticcurveto{\pgfqpoint{11.149454in}{0.579580in}}{\pgfqpoint{11.149454in}{0.565170in}}%
\pgfpathlineto{\pgfqpoint{11.149454in}{0.565170in}}%
\pgfpathclose%
\pgfpathmoveto{\pgfqpoint{11.170781in}{0.616793in}}%
\pgfpathlineto{\pgfqpoint{11.181138in}{0.616793in}}%
\pgfpathlineto{\pgfqpoint{11.181138in}{0.529200in}}%
\pgfpathlineto{\pgfqpoint{11.170781in}{0.529200in}}%
\pgfpathlineto{\pgfqpoint{11.170781in}{0.616793in}}%
\pgfpathlineto{\pgfqpoint{11.170781in}{0.616793in}}%
\pgfpathclose%
\pgfusepath{fill}%
\end{pgfscope}%
\begin{pgfscope}%
\pgfsetbuttcap%
\pgfsetroundjoin%
\definecolor{currentfill}{rgb}{0.309804,0.372549,0.419608}%
\pgfsetfillcolor{currentfill}%
\pgfsetlinewidth{0.000000pt}%
\definecolor{currentstroke}{rgb}{0.309804,0.372549,0.419608}%
\pgfsetstrokecolor{currentstroke}%
\pgfsetdash{}{0pt}%
\pgfpathmoveto{\pgfqpoint{11.283604in}{0.610147in}}%
\pgfpathlineto{\pgfqpoint{11.283604in}{0.592243in}}%
\pgfpathlineto{\pgfqpoint{11.304930in}{0.592243in}}%
\pgfpathlineto{\pgfqpoint{11.304930in}{0.584191in}}%
\pgfpathlineto{\pgfqpoint{11.283604in}{0.584191in}}%
\pgfpathlineto{\pgfqpoint{11.283604in}{0.549968in}}%
\pgfpathquadraticcurveto{\pgfqpoint{11.283604in}{0.542259in}}{\pgfqpoint{11.285711in}{0.540061in}}%
\pgfpathquadraticcurveto{\pgfqpoint{11.287818in}{0.537864in}}{\pgfqpoint{11.294303in}{0.537864in}}%
\pgfpathlineto{\pgfqpoint{11.304930in}{0.537864in}}%
\pgfpathlineto{\pgfqpoint{11.304930in}{0.529200in}}%
\pgfpathlineto{\pgfqpoint{11.294303in}{0.529200in}}%
\pgfpathquadraticcurveto{\pgfqpoint{11.282307in}{0.529200in}}{\pgfqpoint{11.277750in}{0.533667in}}%
\pgfpathquadraticcurveto{\pgfqpoint{11.273193in}{0.538152in}}{\pgfqpoint{11.273193in}{0.549968in}}%
\pgfpathlineto{\pgfqpoint{11.273193in}{0.584191in}}%
\pgfpathlineto{\pgfqpoint{11.265591in}{0.584191in}}%
\pgfpathlineto{\pgfqpoint{11.265591in}{0.592243in}}%
\pgfpathlineto{\pgfqpoint{11.273193in}{0.592243in}}%
\pgfpathlineto{\pgfqpoint{11.273193in}{0.610147in}}%
\pgfpathlineto{\pgfqpoint{11.283604in}{0.610147in}}%
\pgfpathlineto{\pgfqpoint{11.283604in}{0.610147in}}%
\pgfpathclose%
\pgfpathmoveto{\pgfqpoint{11.372476in}{0.563315in}}%
\pgfpathlineto{\pgfqpoint{11.372476in}{0.558254in}}%
\pgfpathlineto{\pgfqpoint{11.324851in}{0.558254in}}%
\pgfpathquadraticcurveto{\pgfqpoint{11.325536in}{0.547554in}}{\pgfqpoint{11.331300in}{0.541953in}}%
\pgfpathquadraticcurveto{\pgfqpoint{11.337064in}{0.536351in}}{\pgfqpoint{11.347367in}{0.536351in}}%
\pgfpathquadraticcurveto{\pgfqpoint{11.353329in}{0.536351in}}{\pgfqpoint{11.358930in}{0.537810in}}%
\pgfpathquadraticcurveto{\pgfqpoint{11.364532in}{0.539269in}}{\pgfqpoint{11.370062in}{0.542205in}}%
\pgfpathlineto{\pgfqpoint{11.370062in}{0.532406in}}%
\pgfpathquadraticcurveto{\pgfqpoint{11.364478in}{0.530047in}}{\pgfqpoint{11.358624in}{0.528804in}}%
\pgfpathquadraticcurveto{\pgfqpoint{11.352770in}{0.527561in}}{\pgfqpoint{11.346754in}{0.527561in}}%
\pgfpathquadraticcurveto{\pgfqpoint{11.331660in}{0.527561in}}{\pgfqpoint{11.322852in}{0.536333in}}%
\pgfpathquadraticcurveto{\pgfqpoint{11.314044in}{0.545123in}}{\pgfqpoint{11.314044in}{0.560109in}}%
\pgfpathquadraticcurveto{\pgfqpoint{11.314044in}{0.575581in}}{\pgfqpoint{11.322402in}{0.584659in}}%
\pgfpathquadraticcurveto{\pgfqpoint{11.330759in}{0.593756in}}{\pgfqpoint{11.344953in}{0.593756in}}%
\pgfpathquadraticcurveto{\pgfqpoint{11.357670in}{0.593756in}}{\pgfqpoint{11.365073in}{0.585560in}}%
\pgfpathquadraticcurveto{\pgfqpoint{11.372476in}{0.577383in}}{\pgfqpoint{11.372476in}{0.563315in}}%
\pgfpathlineto{\pgfqpoint{11.372476in}{0.563315in}}%
\pgfpathclose%
\pgfpathmoveto{\pgfqpoint{11.362119in}{0.566359in}}%
\pgfpathquadraticcurveto{\pgfqpoint{11.362010in}{0.574843in}}{\pgfqpoint{11.357363in}{0.579904in}}%
\pgfpathquadraticcurveto{\pgfqpoint{11.352716in}{0.584984in}}{\pgfqpoint{11.345061in}{0.584984in}}%
\pgfpathquadraticcurveto{\pgfqpoint{11.336397in}{0.584984in}}{\pgfqpoint{11.331192in}{0.580084in}}%
\pgfpathquadraticcurveto{\pgfqpoint{11.325986in}{0.575185in}}{\pgfqpoint{11.325194in}{0.566287in}}%
\pgfpathlineto{\pgfqpoint{11.362119in}{0.566359in}}%
\pgfpathlineto{\pgfqpoint{11.362119in}{0.566359in}}%
\pgfpathclose%
\pgfpathmoveto{\pgfqpoint{11.429664in}{0.590387in}}%
\pgfpathlineto{\pgfqpoint{11.429664in}{0.580589in}}%
\pgfpathquadraticcurveto{\pgfqpoint{11.425287in}{0.582840in}}{\pgfqpoint{11.420550in}{0.583957in}}%
\pgfpathquadraticcurveto{\pgfqpoint{11.415831in}{0.585092in}}{\pgfqpoint{11.410751in}{0.585092in}}%
\pgfpathquadraticcurveto{\pgfqpoint{11.403042in}{0.585092in}}{\pgfqpoint{11.399188in}{0.582732in}}%
\pgfpathquadraticcurveto{\pgfqpoint{11.395333in}{0.580373in}}{\pgfqpoint{11.395333in}{0.575635in}}%
\pgfpathquadraticcurveto{\pgfqpoint{11.395333in}{0.572033in}}{\pgfqpoint{11.398089in}{0.569980in}}%
\pgfpathquadraticcurveto{\pgfqpoint{11.400845in}{0.567926in}}{\pgfqpoint{11.409184in}{0.566071in}}%
\pgfpathlineto{\pgfqpoint{11.412733in}{0.565278in}}%
\pgfpathquadraticcurveto{\pgfqpoint{11.423756in}{0.562919in}}{\pgfqpoint{11.428403in}{0.558614in}}%
\pgfpathquadraticcurveto{\pgfqpoint{11.433050in}{0.554309in}}{\pgfqpoint{11.433050in}{0.546600in}}%
\pgfpathquadraticcurveto{\pgfqpoint{11.433050in}{0.537810in}}{\pgfqpoint{11.426098in}{0.532676in}}%
\pgfpathquadraticcurveto{\pgfqpoint{11.419145in}{0.527561in}}{\pgfqpoint{11.406987in}{0.527561in}}%
\pgfpathquadraticcurveto{\pgfqpoint{11.401925in}{0.527561in}}{\pgfqpoint{11.396432in}{0.528552in}}%
\pgfpathquadraticcurveto{\pgfqpoint{11.390938in}{0.529542in}}{\pgfqpoint{11.384868in}{0.531506in}}%
\pgfpathlineto{\pgfqpoint{11.384868in}{0.542205in}}%
\pgfpathquadraticcurveto{\pgfqpoint{11.390614in}{0.539215in}}{\pgfqpoint{11.396180in}{0.537720in}}%
\pgfpathquadraticcurveto{\pgfqpoint{11.401745in}{0.536243in}}{\pgfqpoint{11.407221in}{0.536243in}}%
\pgfpathquadraticcurveto{\pgfqpoint{11.414534in}{0.536243in}}{\pgfqpoint{11.418461in}{0.538746in}}%
\pgfpathquadraticcurveto{\pgfqpoint{11.422405in}{0.541250in}}{\pgfqpoint{11.422405in}{0.545807in}}%
\pgfpathquadraticcurveto{\pgfqpoint{11.422405in}{0.550022in}}{\pgfqpoint{11.419559in}{0.552274in}}%
\pgfpathquadraticcurveto{\pgfqpoint{11.416731in}{0.554525in}}{\pgfqpoint{11.407095in}{0.556614in}}%
\pgfpathlineto{\pgfqpoint{11.403492in}{0.557461in}}%
\pgfpathquadraticcurveto{\pgfqpoint{11.393874in}{0.559478in}}{\pgfqpoint{11.389587in}{0.563675in}}%
\pgfpathquadraticcurveto{\pgfqpoint{11.385318in}{0.567872in}}{\pgfqpoint{11.385318in}{0.575185in}}%
\pgfpathquadraticcurveto{\pgfqpoint{11.385318in}{0.584083in}}{\pgfqpoint{11.391622in}{0.588910in}}%
\pgfpathquadraticcurveto{\pgfqpoint{11.397927in}{0.593756in}}{\pgfqpoint{11.409527in}{0.593756in}}%
\pgfpathquadraticcurveto{\pgfqpoint{11.415254in}{0.593756in}}{\pgfqpoint{11.420316in}{0.592909in}}%
\pgfpathquadraticcurveto{\pgfqpoint{11.425395in}{0.592080in}}{\pgfqpoint{11.429664in}{0.590387in}}%
\pgfpathlineto{\pgfqpoint{11.429664in}{0.590387in}}%
\pgfpathclose%
\pgfpathmoveto{\pgfqpoint{11.459780in}{0.610147in}}%
\pgfpathlineto{\pgfqpoint{11.459780in}{0.592243in}}%
\pgfpathlineto{\pgfqpoint{11.481107in}{0.592243in}}%
\pgfpathlineto{\pgfqpoint{11.481107in}{0.584191in}}%
\pgfpathlineto{\pgfqpoint{11.459780in}{0.584191in}}%
\pgfpathlineto{\pgfqpoint{11.459780in}{0.549968in}}%
\pgfpathquadraticcurveto{\pgfqpoint{11.459780in}{0.542259in}}{\pgfqpoint{11.461888in}{0.540061in}}%
\pgfpathquadraticcurveto{\pgfqpoint{11.463995in}{0.537864in}}{\pgfqpoint{11.470480in}{0.537864in}}%
\pgfpathlineto{\pgfqpoint{11.481107in}{0.537864in}}%
\pgfpathlineto{\pgfqpoint{11.481107in}{0.529200in}}%
\pgfpathlineto{\pgfqpoint{11.470480in}{0.529200in}}%
\pgfpathquadraticcurveto{\pgfqpoint{11.458484in}{0.529200in}}{\pgfqpoint{11.453926in}{0.533667in}}%
\pgfpathquadraticcurveto{\pgfqpoint{11.449369in}{0.538152in}}{\pgfqpoint{11.449369in}{0.549968in}}%
\pgfpathlineto{\pgfqpoint{11.449369in}{0.584191in}}%
\pgfpathlineto{\pgfqpoint{11.441768in}{0.584191in}}%
\pgfpathlineto{\pgfqpoint{11.441768in}{0.592243in}}%
\pgfpathlineto{\pgfqpoint{11.449369in}{0.592243in}}%
\pgfpathlineto{\pgfqpoint{11.449369in}{0.610147in}}%
\pgfpathlineto{\pgfqpoint{11.459780in}{0.610147in}}%
\pgfpathlineto{\pgfqpoint{11.459780in}{0.610147in}}%
\pgfpathclose%
\pgfpathmoveto{\pgfqpoint{11.534909in}{0.590387in}}%
\pgfpathlineto{\pgfqpoint{11.534909in}{0.580589in}}%
\pgfpathquadraticcurveto{\pgfqpoint{11.530532in}{0.582840in}}{\pgfqpoint{11.525795in}{0.583957in}}%
\pgfpathquadraticcurveto{\pgfqpoint{11.521076in}{0.585092in}}{\pgfqpoint{11.515996in}{0.585092in}}%
\pgfpathquadraticcurveto{\pgfqpoint{11.508287in}{0.585092in}}{\pgfqpoint{11.504433in}{0.582732in}}%
\pgfpathquadraticcurveto{\pgfqpoint{11.500578in}{0.580373in}}{\pgfqpoint{11.500578in}{0.575635in}}%
\pgfpathquadraticcurveto{\pgfqpoint{11.500578in}{0.572033in}}{\pgfqpoint{11.503334in}{0.569980in}}%
\pgfpathquadraticcurveto{\pgfqpoint{11.506090in}{0.567926in}}{\pgfqpoint{11.514429in}{0.566071in}}%
\pgfpathlineto{\pgfqpoint{11.517978in}{0.565278in}}%
\pgfpathquadraticcurveto{\pgfqpoint{11.529001in}{0.562919in}}{\pgfqpoint{11.533648in}{0.558614in}}%
\pgfpathquadraticcurveto{\pgfqpoint{11.538295in}{0.554309in}}{\pgfqpoint{11.538295in}{0.546600in}}%
\pgfpathquadraticcurveto{\pgfqpoint{11.538295in}{0.537810in}}{\pgfqpoint{11.531343in}{0.532676in}}%
\pgfpathquadraticcurveto{\pgfqpoint{11.524390in}{0.527561in}}{\pgfqpoint{11.512232in}{0.527561in}}%
\pgfpathquadraticcurveto{\pgfqpoint{11.507170in}{0.527561in}}{\pgfqpoint{11.501677in}{0.528552in}}%
\pgfpathquadraticcurveto{\pgfqpoint{11.496183in}{0.529542in}}{\pgfqpoint{11.490113in}{0.531506in}}%
\pgfpathlineto{\pgfqpoint{11.490113in}{0.542205in}}%
\pgfpathquadraticcurveto{\pgfqpoint{11.495859in}{0.539215in}}{\pgfqpoint{11.501425in}{0.537720in}}%
\pgfpathquadraticcurveto{\pgfqpoint{11.506990in}{0.536243in}}{\pgfqpoint{11.512466in}{0.536243in}}%
\pgfpathquadraticcurveto{\pgfqpoint{11.519779in}{0.536243in}}{\pgfqpoint{11.523706in}{0.538746in}}%
\pgfpathquadraticcurveto{\pgfqpoint{11.527650in}{0.541250in}}{\pgfqpoint{11.527650in}{0.545807in}}%
\pgfpathquadraticcurveto{\pgfqpoint{11.527650in}{0.550022in}}{\pgfqpoint{11.524804in}{0.552274in}}%
\pgfpathquadraticcurveto{\pgfqpoint{11.521976in}{0.554525in}}{\pgfqpoint{11.512340in}{0.556614in}}%
\pgfpathlineto{\pgfqpoint{11.508737in}{0.557461in}}%
\pgfpathquadraticcurveto{\pgfqpoint{11.499119in}{0.559478in}}{\pgfqpoint{11.494832in}{0.563675in}}%
\pgfpathquadraticcurveto{\pgfqpoint{11.490563in}{0.567872in}}{\pgfqpoint{11.490563in}{0.575185in}}%
\pgfpathquadraticcurveto{\pgfqpoint{11.490563in}{0.584083in}}{\pgfqpoint{11.496867in}{0.588910in}}%
\pgfpathquadraticcurveto{\pgfqpoint{11.503172in}{0.593756in}}{\pgfqpoint{11.514772in}{0.593756in}}%
\pgfpathquadraticcurveto{\pgfqpoint{11.520499in}{0.593756in}}{\pgfqpoint{11.525561in}{0.592909in}}%
\pgfpathquadraticcurveto{\pgfqpoint{11.530640in}{0.592080in}}{\pgfqpoint{11.534909in}{0.590387in}}%
\pgfpathlineto{\pgfqpoint{11.534909in}{0.590387in}}%
\pgfpathclose%
\pgfpathmoveto{\pgfqpoint{11.557424in}{0.543502in}}%
\pgfpathlineto{\pgfqpoint{11.569294in}{0.543502in}}%
\pgfpathlineto{\pgfqpoint{11.569294in}{0.533811in}}%
\pgfpathlineto{\pgfqpoint{11.560072in}{0.515799in}}%
\pgfpathlineto{\pgfqpoint{11.552813in}{0.515799in}}%
\pgfpathlineto{\pgfqpoint{11.557424in}{0.533811in}}%
\pgfpathlineto{\pgfqpoint{11.557424in}{0.543502in}}%
\pgfpathlineto{\pgfqpoint{11.557424in}{0.543502in}}%
\pgfpathclose%
\pgfusepath{fill}%
\end{pgfscope}%
\begin{pgfscope}%
\pgfsetbuttcap%
\pgfsetroundjoin%
\definecolor{currentfill}{rgb}{0.309804,0.372549,0.419608}%
\pgfsetfillcolor{currentfill}%
\pgfsetlinewidth{0.000000pt}%
\definecolor{currentstroke}{rgb}{0.309804,0.372549,0.419608}%
\pgfsetstrokecolor{currentstroke}%
\pgfsetdash{}{0pt}%
\pgfpathmoveto{\pgfqpoint{11.653841in}{0.592243in}}%
\pgfpathlineto{\pgfqpoint{11.664198in}{0.592243in}}%
\pgfpathlineto{\pgfqpoint{11.677149in}{0.543051in}}%
\pgfpathlineto{\pgfqpoint{11.690027in}{0.592243in}}%
\pgfpathlineto{\pgfqpoint{11.702239in}{0.592243in}}%
\pgfpathlineto{\pgfqpoint{11.715190in}{0.543051in}}%
\pgfpathlineto{\pgfqpoint{11.728087in}{0.592243in}}%
\pgfpathlineto{\pgfqpoint{11.738444in}{0.592243in}}%
\pgfpathlineto{\pgfqpoint{11.721945in}{0.529200in}}%
\pgfpathlineto{\pgfqpoint{11.709733in}{0.529200in}}%
\pgfpathlineto{\pgfqpoint{11.696169in}{0.580877in}}%
\pgfpathlineto{\pgfqpoint{11.682552in}{0.529200in}}%
\pgfpathlineto{\pgfqpoint{11.670322in}{0.529200in}}%
\pgfpathlineto{\pgfqpoint{11.653841in}{0.592243in}}%
\pgfpathlineto{\pgfqpoint{11.653841in}{0.592243in}}%
\pgfpathclose%
\pgfpathmoveto{\pgfqpoint{11.754132in}{0.592243in}}%
\pgfpathlineto{\pgfqpoint{11.764489in}{0.592243in}}%
\pgfpathlineto{\pgfqpoint{11.764489in}{0.529200in}}%
\pgfpathlineto{\pgfqpoint{11.754132in}{0.529200in}}%
\pgfpathlineto{\pgfqpoint{11.754132in}{0.592243in}}%
\pgfpathlineto{\pgfqpoint{11.754132in}{0.592243in}}%
\pgfpathclose%
\pgfpathmoveto{\pgfqpoint{11.754132in}{0.616793in}}%
\pgfpathlineto{\pgfqpoint{11.764489in}{0.616793in}}%
\pgfpathlineto{\pgfqpoint{11.764489in}{0.603662in}}%
\pgfpathlineto{\pgfqpoint{11.754132in}{0.603662in}}%
\pgfpathlineto{\pgfqpoint{11.754132in}{0.616793in}}%
\pgfpathlineto{\pgfqpoint{11.754132in}{0.616793in}}%
\pgfpathclose%
\pgfpathmoveto{\pgfqpoint{11.796407in}{0.610147in}}%
\pgfpathlineto{\pgfqpoint{11.796407in}{0.592243in}}%
\pgfpathlineto{\pgfqpoint{11.817733in}{0.592243in}}%
\pgfpathlineto{\pgfqpoint{11.817733in}{0.584191in}}%
\pgfpathlineto{\pgfqpoint{11.796407in}{0.584191in}}%
\pgfpathlineto{\pgfqpoint{11.796407in}{0.549968in}}%
\pgfpathquadraticcurveto{\pgfqpoint{11.796407in}{0.542259in}}{\pgfqpoint{11.798514in}{0.540061in}}%
\pgfpathquadraticcurveto{\pgfqpoint{11.800622in}{0.537864in}}{\pgfqpoint{11.807106in}{0.537864in}}%
\pgfpathlineto{\pgfqpoint{11.817733in}{0.537864in}}%
\pgfpathlineto{\pgfqpoint{11.817733in}{0.529200in}}%
\pgfpathlineto{\pgfqpoint{11.807106in}{0.529200in}}%
\pgfpathquadraticcurveto{\pgfqpoint{11.795110in}{0.529200in}}{\pgfqpoint{11.790553in}{0.533667in}}%
\pgfpathquadraticcurveto{\pgfqpoint{11.785996in}{0.538152in}}{\pgfqpoint{11.785996in}{0.549968in}}%
\pgfpathlineto{\pgfqpoint{11.785996in}{0.584191in}}%
\pgfpathlineto{\pgfqpoint{11.778395in}{0.584191in}}%
\pgfpathlineto{\pgfqpoint{11.778395in}{0.592243in}}%
\pgfpathlineto{\pgfqpoint{11.785996in}{0.592243in}}%
\pgfpathlineto{\pgfqpoint{11.785996in}{0.610147in}}%
\pgfpathlineto{\pgfqpoint{11.796407in}{0.610147in}}%
\pgfpathlineto{\pgfqpoint{11.796407in}{0.610147in}}%
\pgfpathclose%
\pgfpathmoveto{\pgfqpoint{11.883766in}{0.567260in}}%
\pgfpathlineto{\pgfqpoint{11.883766in}{0.529200in}}%
\pgfpathlineto{\pgfqpoint{11.873409in}{0.529200in}}%
\pgfpathlineto{\pgfqpoint{11.873409in}{0.566917in}}%
\pgfpathquadraticcurveto{\pgfqpoint{11.873409in}{0.575869in}}{\pgfqpoint{11.869915in}{0.580300in}}%
\pgfpathquadraticcurveto{\pgfqpoint{11.866420in}{0.584749in}}{\pgfqpoint{11.859450in}{0.584749in}}%
\pgfpathquadraticcurveto{\pgfqpoint{11.851056in}{0.584749in}}{\pgfqpoint{11.846211in}{0.579400in}}%
\pgfpathquadraticcurveto{\pgfqpoint{11.841365in}{0.574068in}}{\pgfqpoint{11.841365in}{0.564828in}}%
\pgfpathlineto{\pgfqpoint{11.841365in}{0.529200in}}%
\pgfpathlineto{\pgfqpoint{11.830954in}{0.529200in}}%
\pgfpathlineto{\pgfqpoint{11.830954in}{0.616793in}}%
\pgfpathlineto{\pgfqpoint{11.841365in}{0.616793in}}%
\pgfpathlineto{\pgfqpoint{11.841365in}{0.582444in}}%
\pgfpathquadraticcurveto{\pgfqpoint{11.845094in}{0.588136in}}{\pgfqpoint{11.850119in}{0.590946in}}%
\pgfpathquadraticcurveto{\pgfqpoint{11.855163in}{0.593756in}}{\pgfqpoint{11.861755in}{0.593756in}}%
\pgfpathquadraticcurveto{\pgfqpoint{11.872616in}{0.593756in}}{\pgfqpoint{11.878182in}{0.587037in}}%
\pgfpathquadraticcurveto{\pgfqpoint{11.883766in}{0.580318in}}{\pgfqpoint{11.883766in}{0.567260in}}%
\pgfpathlineto{\pgfqpoint{11.883766in}{0.567260in}}%
\pgfpathclose%
\pgfpathmoveto{\pgfqpoint{11.928832in}{0.584984in}}%
\pgfpathquadraticcurveto{\pgfqpoint{11.920511in}{0.584984in}}{\pgfqpoint{11.915665in}{0.578481in}}%
\pgfpathquadraticcurveto{\pgfqpoint{11.910820in}{0.571979in}}{\pgfqpoint{11.910820in}{0.560667in}}%
\pgfpathquadraticcurveto{\pgfqpoint{11.910820in}{0.549356in}}{\pgfqpoint{11.915629in}{0.542853in}}%
\pgfpathquadraticcurveto{\pgfqpoint{11.920457in}{0.536351in}}{\pgfqpoint{11.928832in}{0.536351in}}%
\pgfpathquadraticcurveto{\pgfqpoint{11.937118in}{0.536351in}}{\pgfqpoint{11.941945in}{0.542871in}}%
\pgfpathquadraticcurveto{\pgfqpoint{11.946790in}{0.549410in}}{\pgfqpoint{11.946790in}{0.560667in}}%
\pgfpathquadraticcurveto{\pgfqpoint{11.946790in}{0.571871in}}{\pgfqpoint{11.941945in}{0.578427in}}%
\pgfpathquadraticcurveto{\pgfqpoint{11.937118in}{0.584984in}}{\pgfqpoint{11.928832in}{0.584984in}}%
\pgfpathlineto{\pgfqpoint{11.928832in}{0.584984in}}%
\pgfpathclose%
\pgfpathmoveto{\pgfqpoint{11.928832in}{0.593756in}}%
\pgfpathquadraticcurveto{\pgfqpoint{11.942341in}{0.593756in}}{\pgfqpoint{11.950051in}{0.584966in}}%
\pgfpathquadraticcurveto{\pgfqpoint{11.957778in}{0.576194in}}{\pgfqpoint{11.957778in}{0.560667in}}%
\pgfpathquadraticcurveto{\pgfqpoint{11.957778in}{0.545195in}}{\pgfqpoint{11.950051in}{0.536369in}}%
\pgfpathquadraticcurveto{\pgfqpoint{11.942341in}{0.527561in}}{\pgfqpoint{11.928832in}{0.527561in}}%
\pgfpathquadraticcurveto{\pgfqpoint{11.915269in}{0.527561in}}{\pgfqpoint{11.907578in}{0.536369in}}%
\pgfpathquadraticcurveto{\pgfqpoint{11.899905in}{0.545195in}}{\pgfqpoint{11.899905in}{0.560667in}}%
\pgfpathquadraticcurveto{\pgfqpoint{11.899905in}{0.576194in}}{\pgfqpoint{11.907578in}{0.584966in}}%
\pgfpathquadraticcurveto{\pgfqpoint{11.915269in}{0.593756in}}{\pgfqpoint{11.928832in}{0.593756in}}%
\pgfpathlineto{\pgfqpoint{11.928832in}{0.593756in}}%
\pgfpathclose%
\pgfpathmoveto{\pgfqpoint{11.973881in}{0.554075in}}%
\pgfpathlineto{\pgfqpoint{11.973881in}{0.592243in}}%
\pgfpathlineto{\pgfqpoint{11.984238in}{0.592243in}}%
\pgfpathlineto{\pgfqpoint{11.984238in}{0.554471in}}%
\pgfpathquadraticcurveto{\pgfqpoint{11.984238in}{0.545519in}}{\pgfqpoint{11.987714in}{0.541034in}}%
\pgfpathquadraticcurveto{\pgfqpoint{11.991208in}{0.536567in}}{\pgfqpoint{11.998197in}{0.536567in}}%
\pgfpathquadraticcurveto{\pgfqpoint{12.006573in}{0.536567in}}{\pgfqpoint{12.011436in}{0.541917in}}%
\pgfpathquadraticcurveto{\pgfqpoint{12.016317in}{0.547266in}}{\pgfqpoint{12.016317in}{0.556506in}}%
\pgfpathlineto{\pgfqpoint{12.016317in}{0.592243in}}%
\pgfpathlineto{\pgfqpoint{12.026674in}{0.592243in}}%
\pgfpathlineto{\pgfqpoint{12.026674in}{0.529200in}}%
\pgfpathlineto{\pgfqpoint{12.016317in}{0.529200in}}%
\pgfpathlineto{\pgfqpoint{12.016317in}{0.538891in}}%
\pgfpathquadraticcurveto{\pgfqpoint{12.012553in}{0.533145in}}{\pgfqpoint{12.007563in}{0.530353in}}%
\pgfpathquadraticcurveto{\pgfqpoint{12.002592in}{0.527561in}}{\pgfqpoint{11.996000in}{0.527561in}}%
\pgfpathquadraticcurveto{\pgfqpoint{11.985138in}{0.527561in}}{\pgfqpoint{11.979501in}{0.534315in}}%
\pgfpathquadraticcurveto{\pgfqpoint{11.973881in}{0.541070in}}{\pgfqpoint{11.973881in}{0.554075in}}%
\pgfpathlineto{\pgfqpoint{11.973881in}{0.554075in}}%
\pgfpathclose%
\pgfpathmoveto{\pgfqpoint{11.999944in}{0.593756in}}%
\pgfpathlineto{\pgfqpoint{11.999944in}{0.593756in}}%
\pgfpathlineto{\pgfqpoint{11.999944in}{0.593756in}}%
\pgfpathclose%
\pgfpathmoveto{\pgfqpoint{12.058250in}{0.610147in}}%
\pgfpathlineto{\pgfqpoint{12.058250in}{0.592243in}}%
\pgfpathlineto{\pgfqpoint{12.079576in}{0.592243in}}%
\pgfpathlineto{\pgfqpoint{12.079576in}{0.584191in}}%
\pgfpathlineto{\pgfqpoint{12.058250in}{0.584191in}}%
\pgfpathlineto{\pgfqpoint{12.058250in}{0.549968in}}%
\pgfpathquadraticcurveto{\pgfqpoint{12.058250in}{0.542259in}}{\pgfqpoint{12.060357in}{0.540061in}}%
\pgfpathquadraticcurveto{\pgfqpoint{12.062465in}{0.537864in}}{\pgfqpoint{12.068949in}{0.537864in}}%
\pgfpathlineto{\pgfqpoint{12.079576in}{0.537864in}}%
\pgfpathlineto{\pgfqpoint{12.079576in}{0.529200in}}%
\pgfpathlineto{\pgfqpoint{12.068949in}{0.529200in}}%
\pgfpathquadraticcurveto{\pgfqpoint{12.056953in}{0.529200in}}{\pgfqpoint{12.052396in}{0.533667in}}%
\pgfpathquadraticcurveto{\pgfqpoint{12.047839in}{0.538152in}}{\pgfqpoint{12.047839in}{0.549968in}}%
\pgfpathlineto{\pgfqpoint{12.047839in}{0.584191in}}%
\pgfpathlineto{\pgfqpoint{12.040238in}{0.584191in}}%
\pgfpathlineto{\pgfqpoint{12.040238in}{0.592243in}}%
\pgfpathlineto{\pgfqpoint{12.047839in}{0.592243in}}%
\pgfpathlineto{\pgfqpoint{12.047839in}{0.610147in}}%
\pgfpathlineto{\pgfqpoint{12.058250in}{0.610147in}}%
\pgfpathlineto{\pgfqpoint{12.058250in}{0.610147in}}%
\pgfpathclose%
\pgfusepath{fill}%
\end{pgfscope}%
\begin{pgfscope}%
\pgfsetbuttcap%
\pgfsetroundjoin%
\definecolor{currentfill}{rgb}{0.309804,0.372549,0.419608}%
\pgfsetfillcolor{currentfill}%
\pgfsetlinewidth{0.000000pt}%
\definecolor{currentstroke}{rgb}{0.309804,0.372549,0.419608}%
\pgfsetstrokecolor{currentstroke}%
\pgfsetdash{}{0pt}%
\pgfpathmoveto{\pgfqpoint{12.225442in}{0.580138in}}%
\pgfpathquadraticcurveto{\pgfqpoint{12.229333in}{0.587127in}}{\pgfqpoint{12.234736in}{0.590441in}}%
\pgfpathquadraticcurveto{\pgfqpoint{12.240140in}{0.593756in}}{\pgfqpoint{12.247453in}{0.593756in}}%
\pgfpathquadraticcurveto{\pgfqpoint{12.257306in}{0.593756in}}{\pgfqpoint{12.262655in}{0.586857in}}%
\pgfpathquadraticcurveto{\pgfqpoint{12.268005in}{0.579976in}}{\pgfqpoint{12.268005in}{0.567260in}}%
\pgfpathlineto{\pgfqpoint{12.268005in}{0.529200in}}%
\pgfpathlineto{\pgfqpoint{12.257594in}{0.529200in}}%
\pgfpathlineto{\pgfqpoint{12.257594in}{0.566917in}}%
\pgfpathquadraticcurveto{\pgfqpoint{12.257594in}{0.575978in}}{\pgfqpoint{12.254370in}{0.580355in}}%
\pgfpathquadraticcurveto{\pgfqpoint{12.251164in}{0.584749in}}{\pgfqpoint{12.244589in}{0.584749in}}%
\pgfpathquadraticcurveto{\pgfqpoint{12.236538in}{0.584749in}}{\pgfqpoint{12.231855in}{0.579400in}}%
\pgfpathquadraticcurveto{\pgfqpoint{12.227189in}{0.574068in}}{\pgfqpoint{12.227189in}{0.564828in}}%
\pgfpathlineto{\pgfqpoint{12.227189in}{0.529200in}}%
\pgfpathlineto{\pgfqpoint{12.216778in}{0.529200in}}%
\pgfpathlineto{\pgfqpoint{12.216778in}{0.566917in}}%
\pgfpathquadraticcurveto{\pgfqpoint{12.216778in}{0.576032in}}{\pgfqpoint{12.213572in}{0.580391in}}%
\pgfpathquadraticcurveto{\pgfqpoint{12.210366in}{0.584749in}}{\pgfqpoint{12.203666in}{0.584749in}}%
\pgfpathquadraticcurveto{\pgfqpoint{12.195722in}{0.584749in}}{\pgfqpoint{12.191039in}{0.579382in}}%
\pgfpathquadraticcurveto{\pgfqpoint{12.186374in}{0.574014in}}{\pgfqpoint{12.186374in}{0.564828in}}%
\pgfpathlineto{\pgfqpoint{12.186374in}{0.529200in}}%
\pgfpathlineto{\pgfqpoint{12.175963in}{0.529200in}}%
\pgfpathlineto{\pgfqpoint{12.175963in}{0.592243in}}%
\pgfpathlineto{\pgfqpoint{12.186374in}{0.592243in}}%
\pgfpathlineto{\pgfqpoint{12.186374in}{0.582444in}}%
\pgfpathquadraticcurveto{\pgfqpoint{12.189922in}{0.588244in}}{\pgfqpoint{12.194876in}{0.591000in}}%
\pgfpathquadraticcurveto{\pgfqpoint{12.199829in}{0.593756in}}{\pgfqpoint{12.206638in}{0.593756in}}%
\pgfpathquadraticcurveto{\pgfqpoint{12.213518in}{0.593756in}}{\pgfqpoint{12.218327in}{0.590261in}}%
\pgfpathquadraticcurveto{\pgfqpoint{12.223137in}{0.586785in}}{\pgfqpoint{12.225442in}{0.580138in}}%
\pgfpathlineto{\pgfqpoint{12.225442in}{0.580138in}}%
\pgfpathclose%
\pgfpathmoveto{\pgfqpoint{12.287584in}{0.554075in}}%
\pgfpathlineto{\pgfqpoint{12.287584in}{0.592243in}}%
\pgfpathlineto{\pgfqpoint{12.297941in}{0.592243in}}%
\pgfpathlineto{\pgfqpoint{12.297941in}{0.554471in}}%
\pgfpathquadraticcurveto{\pgfqpoint{12.297941in}{0.545519in}}{\pgfqpoint{12.301417in}{0.541034in}}%
\pgfpathquadraticcurveto{\pgfqpoint{12.304912in}{0.536567in}}{\pgfqpoint{12.311901in}{0.536567in}}%
\pgfpathquadraticcurveto{\pgfqpoint{12.320276in}{0.536567in}}{\pgfqpoint{12.325139in}{0.541917in}}%
\pgfpathquadraticcurveto{\pgfqpoint{12.330021in}{0.547266in}}{\pgfqpoint{12.330021in}{0.556506in}}%
\pgfpathlineto{\pgfqpoint{12.330021in}{0.592243in}}%
\pgfpathlineto{\pgfqpoint{12.340378in}{0.592243in}}%
\pgfpathlineto{\pgfqpoint{12.340378in}{0.529200in}}%
\pgfpathlineto{\pgfqpoint{12.330021in}{0.529200in}}%
\pgfpathlineto{\pgfqpoint{12.330021in}{0.538891in}}%
\pgfpathquadraticcurveto{\pgfqpoint{12.326256in}{0.533145in}}{\pgfqpoint{12.321267in}{0.530353in}}%
\pgfpathquadraticcurveto{\pgfqpoint{12.316296in}{0.527561in}}{\pgfqpoint{12.309703in}{0.527561in}}%
\pgfpathquadraticcurveto{\pgfqpoint{12.298842in}{0.527561in}}{\pgfqpoint{12.293204in}{0.534315in}}%
\pgfpathquadraticcurveto{\pgfqpoint{12.287584in}{0.541070in}}{\pgfqpoint{12.287584in}{0.554075in}}%
\pgfpathlineto{\pgfqpoint{12.287584in}{0.554075in}}%
\pgfpathclose%
\pgfpathmoveto{\pgfqpoint{12.313648in}{0.593756in}}%
\pgfpathlineto{\pgfqpoint{12.313648in}{0.593756in}}%
\pgfpathlineto{\pgfqpoint{12.313648in}{0.593756in}}%
\pgfpathclose%
\pgfpathmoveto{\pgfqpoint{12.361704in}{0.616793in}}%
\pgfpathlineto{\pgfqpoint{12.372061in}{0.616793in}}%
\pgfpathlineto{\pgfqpoint{12.372061in}{0.529200in}}%
\pgfpathlineto{\pgfqpoint{12.361704in}{0.529200in}}%
\pgfpathlineto{\pgfqpoint{12.361704in}{0.616793in}}%
\pgfpathlineto{\pgfqpoint{12.361704in}{0.616793in}}%
\pgfpathclose%
\pgfpathmoveto{\pgfqpoint{12.403979in}{0.610147in}}%
\pgfpathlineto{\pgfqpoint{12.403979in}{0.592243in}}%
\pgfpathlineto{\pgfqpoint{12.425305in}{0.592243in}}%
\pgfpathlineto{\pgfqpoint{12.425305in}{0.584191in}}%
\pgfpathlineto{\pgfqpoint{12.403979in}{0.584191in}}%
\pgfpathlineto{\pgfqpoint{12.403979in}{0.549968in}}%
\pgfpathquadraticcurveto{\pgfqpoint{12.403979in}{0.542259in}}{\pgfqpoint{12.406086in}{0.540061in}}%
\pgfpathquadraticcurveto{\pgfqpoint{12.408194in}{0.537864in}}{\pgfqpoint{12.414678in}{0.537864in}}%
\pgfpathlineto{\pgfqpoint{12.425305in}{0.537864in}}%
\pgfpathlineto{\pgfqpoint{12.425305in}{0.529200in}}%
\pgfpathlineto{\pgfqpoint{12.414678in}{0.529200in}}%
\pgfpathquadraticcurveto{\pgfqpoint{12.402682in}{0.529200in}}{\pgfqpoint{12.398125in}{0.533667in}}%
\pgfpathquadraticcurveto{\pgfqpoint{12.393568in}{0.538152in}}{\pgfqpoint{12.393568in}{0.549968in}}%
\pgfpathlineto{\pgfqpoint{12.393568in}{0.584191in}}%
\pgfpathlineto{\pgfqpoint{12.385967in}{0.584191in}}%
\pgfpathlineto{\pgfqpoint{12.385967in}{0.592243in}}%
\pgfpathlineto{\pgfqpoint{12.393568in}{0.592243in}}%
\pgfpathlineto{\pgfqpoint{12.393568in}{0.610147in}}%
\pgfpathlineto{\pgfqpoint{12.403979in}{0.610147in}}%
\pgfpathlineto{\pgfqpoint{12.403979in}{0.610147in}}%
\pgfpathclose%
\pgfpathmoveto{\pgfqpoint{12.438922in}{0.592243in}}%
\pgfpathlineto{\pgfqpoint{12.449279in}{0.592243in}}%
\pgfpathlineto{\pgfqpoint{12.449279in}{0.529200in}}%
\pgfpathlineto{\pgfqpoint{12.438922in}{0.529200in}}%
\pgfpathlineto{\pgfqpoint{12.438922in}{0.592243in}}%
\pgfpathlineto{\pgfqpoint{12.438922in}{0.592243in}}%
\pgfpathclose%
\pgfpathmoveto{\pgfqpoint{12.438922in}{0.616793in}}%
\pgfpathlineto{\pgfqpoint{12.449279in}{0.616793in}}%
\pgfpathlineto{\pgfqpoint{12.449279in}{0.603662in}}%
\pgfpathlineto{\pgfqpoint{12.438922in}{0.603662in}}%
\pgfpathlineto{\pgfqpoint{12.438922in}{0.616793in}}%
\pgfpathlineto{\pgfqpoint{12.438922in}{0.616793in}}%
\pgfpathclose%
\pgfpathmoveto{\pgfqpoint{12.480963in}{0.538656in}}%
\pgfpathlineto{\pgfqpoint{12.480963in}{0.505226in}}%
\pgfpathlineto{\pgfqpoint{12.470552in}{0.505226in}}%
\pgfpathlineto{\pgfqpoint{12.470552in}{0.592243in}}%
\pgfpathlineto{\pgfqpoint{12.480963in}{0.592243in}}%
\pgfpathlineto{\pgfqpoint{12.480963in}{0.582678in}}%
\pgfpathquadraticcurveto{\pgfqpoint{12.484241in}{0.588298in}}{\pgfqpoint{12.489212in}{0.591018in}}%
\pgfpathquadraticcurveto{\pgfqpoint{12.494202in}{0.593756in}}{\pgfqpoint{12.501118in}{0.593756in}}%
\pgfpathquadraticcurveto{\pgfqpoint{12.512610in}{0.593756in}}{\pgfqpoint{12.519779in}{0.584641in}}%
\pgfpathquadraticcurveto{\pgfqpoint{12.526966in}{0.575527in}}{\pgfqpoint{12.526966in}{0.560667in}}%
\pgfpathquadraticcurveto{\pgfqpoint{12.526966in}{0.545807in}}{\pgfqpoint{12.519779in}{0.536675in}}%
\pgfpathquadraticcurveto{\pgfqpoint{12.512610in}{0.527561in}}{\pgfqpoint{12.501118in}{0.527561in}}%
\pgfpathquadraticcurveto{\pgfqpoint{12.494202in}{0.527561in}}{\pgfqpoint{12.489212in}{0.530299in}}%
\pgfpathquadraticcurveto{\pgfqpoint{12.484241in}{0.533037in}}{\pgfqpoint{12.480963in}{0.538656in}}%
\pgfpathlineto{\pgfqpoint{12.480963in}{0.538656in}}%
\pgfpathclose%
\pgfpathmoveto{\pgfqpoint{12.516212in}{0.560667in}}%
\pgfpathquadraticcurveto{\pgfqpoint{12.516212in}{0.572087in}}{\pgfqpoint{12.511511in}{0.578589in}}%
\pgfpathquadraticcurveto{\pgfqpoint{12.506810in}{0.585092in}}{\pgfqpoint{12.498597in}{0.585092in}}%
\pgfpathquadraticcurveto{\pgfqpoint{12.490365in}{0.585092in}}{\pgfqpoint{12.485664in}{0.578589in}}%
\pgfpathquadraticcurveto{\pgfqpoint{12.480963in}{0.572087in}}{\pgfqpoint{12.480963in}{0.560667in}}%
\pgfpathquadraticcurveto{\pgfqpoint{12.480963in}{0.549248in}}{\pgfqpoint{12.485664in}{0.542745in}}%
\pgfpathquadraticcurveto{\pgfqpoint{12.490365in}{0.536243in}}{\pgfqpoint{12.498597in}{0.536243in}}%
\pgfpathquadraticcurveto{\pgfqpoint{12.506810in}{0.536243in}}{\pgfqpoint{12.511511in}{0.542745in}}%
\pgfpathquadraticcurveto{\pgfqpoint{12.516212in}{0.549248in}}{\pgfqpoint{12.516212in}{0.560667in}}%
\pgfpathlineto{\pgfqpoint{12.516212in}{0.560667in}}%
\pgfpathclose%
\pgfpathmoveto{\pgfqpoint{12.544131in}{0.616793in}}%
\pgfpathlineto{\pgfqpoint{12.554488in}{0.616793in}}%
\pgfpathlineto{\pgfqpoint{12.554488in}{0.529200in}}%
\pgfpathlineto{\pgfqpoint{12.544131in}{0.529200in}}%
\pgfpathlineto{\pgfqpoint{12.544131in}{0.616793in}}%
\pgfpathlineto{\pgfqpoint{12.544131in}{0.616793in}}%
\pgfpathclose%
\pgfpathmoveto{\pgfqpoint{12.576157in}{0.592243in}}%
\pgfpathlineto{\pgfqpoint{12.586514in}{0.592243in}}%
\pgfpathlineto{\pgfqpoint{12.586514in}{0.529200in}}%
\pgfpathlineto{\pgfqpoint{12.576157in}{0.529200in}}%
\pgfpathlineto{\pgfqpoint{12.576157in}{0.592243in}}%
\pgfpathlineto{\pgfqpoint{12.576157in}{0.592243in}}%
\pgfpathclose%
\pgfpathmoveto{\pgfqpoint{12.576157in}{0.616793in}}%
\pgfpathlineto{\pgfqpoint{12.586514in}{0.616793in}}%
\pgfpathlineto{\pgfqpoint{12.586514in}{0.603662in}}%
\pgfpathlineto{\pgfqpoint{12.576157in}{0.603662in}}%
\pgfpathlineto{\pgfqpoint{12.576157in}{0.616793in}}%
\pgfpathlineto{\pgfqpoint{12.576157in}{0.616793in}}%
\pgfpathclose%
\pgfpathmoveto{\pgfqpoint{12.653555in}{0.589829in}}%
\pgfpathlineto{\pgfqpoint{12.653555in}{0.580138in}}%
\pgfpathquadraticcurveto{\pgfqpoint{12.649160in}{0.582570in}}{\pgfqpoint{12.644747in}{0.583777in}}%
\pgfpathquadraticcurveto{\pgfqpoint{12.640334in}{0.584984in}}{\pgfqpoint{12.635831in}{0.584984in}}%
\pgfpathquadraticcurveto{\pgfqpoint{12.625744in}{0.584984in}}{\pgfqpoint{12.620161in}{0.578589in}}%
\pgfpathquadraticcurveto{\pgfqpoint{12.614595in}{0.572213in}}{\pgfqpoint{12.614595in}{0.560667in}}%
\pgfpathquadraticcurveto{\pgfqpoint{12.614595in}{0.549121in}}{\pgfqpoint{12.620161in}{0.542727in}}%
\pgfpathquadraticcurveto{\pgfqpoint{12.625744in}{0.536351in}}{\pgfqpoint{12.635831in}{0.536351in}}%
\pgfpathquadraticcurveto{\pgfqpoint{12.640334in}{0.536351in}}{\pgfqpoint{12.644747in}{0.537558in}}%
\pgfpathquadraticcurveto{\pgfqpoint{12.649160in}{0.538764in}}{\pgfqpoint{12.653555in}{0.541196in}}%
\pgfpathlineto{\pgfqpoint{12.653555in}{0.531614in}}%
\pgfpathquadraticcurveto{\pgfqpoint{12.649214in}{0.529596in}}{\pgfqpoint{12.644567in}{0.528588in}}%
\pgfpathquadraticcurveto{\pgfqpoint{12.639938in}{0.527561in}}{\pgfqpoint{12.634696in}{0.527561in}}%
\pgfpathquadraticcurveto{\pgfqpoint{12.620449in}{0.527561in}}{\pgfqpoint{12.612055in}{0.536513in}}%
\pgfpathquadraticcurveto{\pgfqpoint{12.603679in}{0.545465in}}{\pgfqpoint{12.603679in}{0.560667in}}%
\pgfpathquadraticcurveto{\pgfqpoint{12.603679in}{0.576086in}}{\pgfqpoint{12.612145in}{0.584912in}}%
\pgfpathquadraticcurveto{\pgfqpoint{12.620629in}{0.593756in}}{\pgfqpoint{12.635381in}{0.593756in}}%
\pgfpathquadraticcurveto{\pgfqpoint{12.640154in}{0.593756in}}{\pgfqpoint{12.644711in}{0.592765in}}%
\pgfpathquadraticcurveto{\pgfqpoint{12.649268in}{0.591792in}}{\pgfqpoint{12.653555in}{0.589829in}}%
\pgfpathlineto{\pgfqpoint{12.653555in}{0.589829in}}%
\pgfpathclose%
\pgfpathmoveto{\pgfqpoint{12.671567in}{0.592243in}}%
\pgfpathlineto{\pgfqpoint{12.681924in}{0.592243in}}%
\pgfpathlineto{\pgfqpoint{12.681924in}{0.529200in}}%
\pgfpathlineto{\pgfqpoint{12.671567in}{0.529200in}}%
\pgfpathlineto{\pgfqpoint{12.671567in}{0.592243in}}%
\pgfpathlineto{\pgfqpoint{12.671567in}{0.592243in}}%
\pgfpathclose%
\pgfpathmoveto{\pgfqpoint{12.671567in}{0.616793in}}%
\pgfpathlineto{\pgfqpoint{12.681924in}{0.616793in}}%
\pgfpathlineto{\pgfqpoint{12.681924in}{0.603662in}}%
\pgfpathlineto{\pgfqpoint{12.671567in}{0.603662in}}%
\pgfpathlineto{\pgfqpoint{12.671567in}{0.616793in}}%
\pgfpathlineto{\pgfqpoint{12.671567in}{0.616793in}}%
\pgfpathclose%
\pgfpathmoveto{\pgfqpoint{12.713842in}{0.610147in}}%
\pgfpathlineto{\pgfqpoint{12.713842in}{0.592243in}}%
\pgfpathlineto{\pgfqpoint{12.735168in}{0.592243in}}%
\pgfpathlineto{\pgfqpoint{12.735168in}{0.584191in}}%
\pgfpathlineto{\pgfqpoint{12.713842in}{0.584191in}}%
\pgfpathlineto{\pgfqpoint{12.713842in}{0.549968in}}%
\pgfpathquadraticcurveto{\pgfqpoint{12.713842in}{0.542259in}}{\pgfqpoint{12.715949in}{0.540061in}}%
\pgfpathquadraticcurveto{\pgfqpoint{12.718057in}{0.537864in}}{\pgfqpoint{12.724541in}{0.537864in}}%
\pgfpathlineto{\pgfqpoint{12.735168in}{0.537864in}}%
\pgfpathlineto{\pgfqpoint{12.735168in}{0.529200in}}%
\pgfpathlineto{\pgfqpoint{12.724541in}{0.529200in}}%
\pgfpathquadraticcurveto{\pgfqpoint{12.712545in}{0.529200in}}{\pgfqpoint{12.707988in}{0.533667in}}%
\pgfpathquadraticcurveto{\pgfqpoint{12.703431in}{0.538152in}}{\pgfqpoint{12.703431in}{0.549968in}}%
\pgfpathlineto{\pgfqpoint{12.703431in}{0.584191in}}%
\pgfpathlineto{\pgfqpoint{12.695830in}{0.584191in}}%
\pgfpathlineto{\pgfqpoint{12.695830in}{0.592243in}}%
\pgfpathlineto{\pgfqpoint{12.703431in}{0.592243in}}%
\pgfpathlineto{\pgfqpoint{12.703431in}{0.610147in}}%
\pgfpathlineto{\pgfqpoint{12.713842in}{0.610147in}}%
\pgfpathlineto{\pgfqpoint{12.713842in}{0.610147in}}%
\pgfpathclose%
\pgfpathmoveto{\pgfqpoint{12.775011in}{0.523346in}}%
\pgfpathquadraticcurveto{\pgfqpoint{12.770634in}{0.512088in}}{\pgfqpoint{12.766455in}{0.508666in}}%
\pgfpathquadraticcurveto{\pgfqpoint{12.762294in}{0.505226in}}{\pgfqpoint{12.755324in}{0.505226in}}%
\pgfpathlineto{\pgfqpoint{12.747038in}{0.505226in}}%
\pgfpathlineto{\pgfqpoint{12.747038in}{0.513890in}}%
\pgfpathlineto{\pgfqpoint{12.753126in}{0.513890in}}%
\pgfpathquadraticcurveto{\pgfqpoint{12.757395in}{0.513890in}}{\pgfqpoint{12.759755in}{0.515925in}}%
\pgfpathquadraticcurveto{\pgfqpoint{12.762132in}{0.517942in}}{\pgfqpoint{12.764996in}{0.525489in}}%
\pgfpathlineto{\pgfqpoint{12.766852in}{0.530209in}}%
\pgfpathlineto{\pgfqpoint{12.741364in}{0.592243in}}%
\pgfpathlineto{\pgfqpoint{12.752334in}{0.592243in}}%
\pgfpathlineto{\pgfqpoint{12.772039in}{0.542943in}}%
\pgfpathlineto{\pgfqpoint{12.791744in}{0.592243in}}%
\pgfpathlineto{\pgfqpoint{12.802714in}{0.592243in}}%
\pgfpathlineto{\pgfqpoint{12.775011in}{0.523346in}}%
\pgfpathlineto{\pgfqpoint{12.775011in}{0.523346in}}%
\pgfpathclose%
\pgfusepath{fill}%
\end{pgfscope}%
\begin{pgfscope}%
\pgfsetbuttcap%
\pgfsetroundjoin%
\definecolor{currentfill}{rgb}{0.309804,0.372549,0.419608}%
\pgfsetfillcolor{currentfill}%
\pgfsetlinewidth{0.000000pt}%
\definecolor{currentstroke}{rgb}{0.309804,0.372549,0.419608}%
\pgfsetstrokecolor{currentstroke}%
\pgfsetdash{}{0pt}%
\pgfpathmoveto{\pgfqpoint{12.931519in}{0.560883in}}%
\pgfpathquadraticcurveto{\pgfqpoint{12.918964in}{0.560883in}}{\pgfqpoint{12.914119in}{0.558019in}}%
\pgfpathquadraticcurveto{\pgfqpoint{12.909274in}{0.555156in}}{\pgfqpoint{12.909274in}{0.548221in}}%
\pgfpathquadraticcurveto{\pgfqpoint{12.909274in}{0.542709in}}{\pgfqpoint{12.912912in}{0.539467in}}%
\pgfpathquadraticcurveto{\pgfqpoint{12.916551in}{0.536243in}}{\pgfqpoint{12.922783in}{0.536243in}}%
\pgfpathquadraticcurveto{\pgfqpoint{12.931411in}{0.536243in}}{\pgfqpoint{12.936616in}{0.542349in}}%
\pgfpathquadraticcurveto{\pgfqpoint{12.941822in}{0.548455in}}{\pgfqpoint{12.941822in}{0.558578in}}%
\pgfpathlineto{\pgfqpoint{12.941822in}{0.560883in}}%
\pgfpathlineto{\pgfqpoint{12.931519in}{0.560883in}}%
\pgfpathlineto{\pgfqpoint{12.931519in}{0.560883in}}%
\pgfpathclose%
\pgfpathmoveto{\pgfqpoint{12.952179in}{0.565170in}}%
\pgfpathlineto{\pgfqpoint{12.952179in}{0.529200in}}%
\pgfpathlineto{\pgfqpoint{12.941822in}{0.529200in}}%
\pgfpathlineto{\pgfqpoint{12.941822in}{0.538764in}}%
\pgfpathquadraticcurveto{\pgfqpoint{12.938273in}{0.533037in}}{\pgfqpoint{12.932978in}{0.530299in}}%
\pgfpathquadraticcurveto{\pgfqpoint{12.927682in}{0.527561in}}{\pgfqpoint{12.920027in}{0.527561in}}%
\pgfpathquadraticcurveto{\pgfqpoint{12.910354in}{0.527561in}}{\pgfqpoint{12.904627in}{0.533001in}}%
\pgfpathquadraticcurveto{\pgfqpoint{12.898917in}{0.538440in}}{\pgfqpoint{12.898917in}{0.547554in}}%
\pgfpathquadraticcurveto{\pgfqpoint{12.898917in}{0.558182in}}{\pgfqpoint{12.906031in}{0.563585in}}%
\pgfpathquadraticcurveto{\pgfqpoint{12.913164in}{0.568989in}}{\pgfqpoint{12.927286in}{0.568989in}}%
\pgfpathlineto{\pgfqpoint{12.941822in}{0.568989in}}%
\pgfpathlineto{\pgfqpoint{12.941822in}{0.570016in}}%
\pgfpathquadraticcurveto{\pgfqpoint{12.941822in}{0.577166in}}{\pgfqpoint{12.937120in}{0.581075in}}%
\pgfpathquadraticcurveto{\pgfqpoint{12.932419in}{0.584984in}}{\pgfqpoint{12.923918in}{0.584984in}}%
\pgfpathquadraticcurveto{\pgfqpoint{12.918514in}{0.584984in}}{\pgfqpoint{12.913380in}{0.583687in}}%
\pgfpathquadraticcurveto{\pgfqpoint{12.908265in}{0.582390in}}{\pgfqpoint{12.903546in}{0.579796in}}%
\pgfpathlineto{\pgfqpoint{12.903546in}{0.589379in}}%
\pgfpathquadraticcurveto{\pgfqpoint{12.909220in}{0.591576in}}{\pgfqpoint{12.914569in}{0.592657in}}%
\pgfpathquadraticcurveto{\pgfqpoint{12.919919in}{0.593756in}}{\pgfqpoint{12.924980in}{0.593756in}}%
\pgfpathquadraticcurveto{\pgfqpoint{12.938669in}{0.593756in}}{\pgfqpoint{12.945424in}{0.586659in}}%
\pgfpathquadraticcurveto{\pgfqpoint{12.952179in}{0.579580in}}{\pgfqpoint{12.952179in}{0.565170in}}%
\pgfpathlineto{\pgfqpoint{12.952179in}{0.565170in}}%
\pgfpathclose%
\pgfpathmoveto{\pgfqpoint{13.014987in}{0.582678in}}%
\pgfpathlineto{\pgfqpoint{13.014987in}{0.616793in}}%
\pgfpathlineto{\pgfqpoint{13.025344in}{0.616793in}}%
\pgfpathlineto{\pgfqpoint{13.025344in}{0.529200in}}%
\pgfpathlineto{\pgfqpoint{13.014987in}{0.529200in}}%
\pgfpathlineto{\pgfqpoint{13.014987in}{0.538656in}}%
\pgfpathquadraticcurveto{\pgfqpoint{13.011727in}{0.533037in}}{\pgfqpoint{13.006737in}{0.530299in}}%
\pgfpathquadraticcurveto{\pgfqpoint{13.001766in}{0.527561in}}{\pgfqpoint{12.994777in}{0.527561in}}%
\pgfpathquadraticcurveto{\pgfqpoint{12.983358in}{0.527561in}}{\pgfqpoint{12.976171in}{0.536675in}}%
\pgfpathquadraticcurveto{\pgfqpoint{12.969002in}{0.545807in}}{\pgfqpoint{12.969002in}{0.560667in}}%
\pgfpathquadraticcurveto{\pgfqpoint{12.969002in}{0.575527in}}{\pgfqpoint{12.976171in}{0.584641in}}%
\pgfpathquadraticcurveto{\pgfqpoint{12.983358in}{0.593756in}}{\pgfqpoint{12.994777in}{0.593756in}}%
\pgfpathquadraticcurveto{\pgfqpoint{13.001766in}{0.593756in}}{\pgfqpoint{13.006737in}{0.591018in}}%
\pgfpathquadraticcurveto{\pgfqpoint{13.011727in}{0.588298in}}{\pgfqpoint{13.014987in}{0.582678in}}%
\pgfpathlineto{\pgfqpoint{13.014987in}{0.582678in}}%
\pgfpathclose%
\pgfpathmoveto{\pgfqpoint{12.979701in}{0.560667in}}%
\pgfpathquadraticcurveto{\pgfqpoint{12.979701in}{0.549248in}}{\pgfqpoint{12.984402in}{0.542745in}}%
\pgfpathquadraticcurveto{\pgfqpoint{12.989104in}{0.536243in}}{\pgfqpoint{12.997317in}{0.536243in}}%
\pgfpathquadraticcurveto{\pgfqpoint{13.005531in}{0.536243in}}{\pgfqpoint{13.010250in}{0.542745in}}%
\pgfpathquadraticcurveto{\pgfqpoint{13.014987in}{0.549248in}}{\pgfqpoint{13.014987in}{0.560667in}}%
\pgfpathquadraticcurveto{\pgfqpoint{13.014987in}{0.572087in}}{\pgfqpoint{13.010250in}{0.578589in}}%
\pgfpathquadraticcurveto{\pgfqpoint{13.005531in}{0.585092in}}{\pgfqpoint{12.997317in}{0.585092in}}%
\pgfpathquadraticcurveto{\pgfqpoint{12.989104in}{0.585092in}}{\pgfqpoint{12.984402in}{0.578589in}}%
\pgfpathquadraticcurveto{\pgfqpoint{12.979701in}{0.572087in}}{\pgfqpoint{12.979701in}{0.560667in}}%
\pgfpathlineto{\pgfqpoint{12.979701in}{0.560667in}}%
\pgfpathclose%
\pgfpathmoveto{\pgfqpoint{13.046688in}{0.592243in}}%
\pgfpathlineto{\pgfqpoint{13.057045in}{0.592243in}}%
\pgfpathlineto{\pgfqpoint{13.057045in}{0.528065in}}%
\pgfpathquadraticcurveto{\pgfqpoint{13.057045in}{0.516033in}}{\pgfqpoint{13.052452in}{0.510629in}}%
\pgfpathquadraticcurveto{\pgfqpoint{13.047877in}{0.505226in}}{\pgfqpoint{13.037682in}{0.505226in}}%
\pgfpathlineto{\pgfqpoint{13.033738in}{0.505226in}}%
\pgfpathlineto{\pgfqpoint{13.033738in}{0.513998in}}%
\pgfpathlineto{\pgfqpoint{13.036512in}{0.513998in}}%
\pgfpathquadraticcurveto{\pgfqpoint{13.042419in}{0.513998in}}{\pgfqpoint{13.044545in}{0.516736in}}%
\pgfpathquadraticcurveto{\pgfqpoint{13.046688in}{0.519455in}}{\pgfqpoint{13.046688in}{0.528065in}}%
\pgfpathlineto{\pgfqpoint{13.046688in}{0.592243in}}%
\pgfpathlineto{\pgfqpoint{13.046688in}{0.592243in}}%
\pgfpathclose%
\pgfpathmoveto{\pgfqpoint{13.046688in}{0.616793in}}%
\pgfpathlineto{\pgfqpoint{13.057045in}{0.616793in}}%
\pgfpathlineto{\pgfqpoint{13.057045in}{0.603662in}}%
\pgfpathlineto{\pgfqpoint{13.046688in}{0.603662in}}%
\pgfpathlineto{\pgfqpoint{13.046688in}{0.616793in}}%
\pgfpathlineto{\pgfqpoint{13.046688in}{0.616793in}}%
\pgfpathclose%
\pgfpathmoveto{\pgfqpoint{13.077651in}{0.554075in}}%
\pgfpathlineto{\pgfqpoint{13.077651in}{0.592243in}}%
\pgfpathlineto{\pgfqpoint{13.088008in}{0.592243in}}%
\pgfpathlineto{\pgfqpoint{13.088008in}{0.554471in}}%
\pgfpathquadraticcurveto{\pgfqpoint{13.088008in}{0.545519in}}{\pgfqpoint{13.091485in}{0.541034in}}%
\pgfpathquadraticcurveto{\pgfqpoint{13.094979in}{0.536567in}}{\pgfqpoint{13.101968in}{0.536567in}}%
\pgfpathquadraticcurveto{\pgfqpoint{13.110343in}{0.536567in}}{\pgfqpoint{13.115207in}{0.541917in}}%
\pgfpathquadraticcurveto{\pgfqpoint{13.120088in}{0.547266in}}{\pgfqpoint{13.120088in}{0.556506in}}%
\pgfpathlineto{\pgfqpoint{13.120088in}{0.592243in}}%
\pgfpathlineto{\pgfqpoint{13.130445in}{0.592243in}}%
\pgfpathlineto{\pgfqpoint{13.130445in}{0.529200in}}%
\pgfpathlineto{\pgfqpoint{13.120088in}{0.529200in}}%
\pgfpathlineto{\pgfqpoint{13.120088in}{0.538891in}}%
\pgfpathquadraticcurveto{\pgfqpoint{13.116323in}{0.533145in}}{\pgfqpoint{13.111334in}{0.530353in}}%
\pgfpathquadraticcurveto{\pgfqpoint{13.106363in}{0.527561in}}{\pgfqpoint{13.099770in}{0.527561in}}%
\pgfpathquadraticcurveto{\pgfqpoint{13.088909in}{0.527561in}}{\pgfqpoint{13.083271in}{0.534315in}}%
\pgfpathquadraticcurveto{\pgfqpoint{13.077651in}{0.541070in}}{\pgfqpoint{13.077651in}{0.554075in}}%
\pgfpathlineto{\pgfqpoint{13.077651in}{0.554075in}}%
\pgfpathclose%
\pgfpathmoveto{\pgfqpoint{13.103715in}{0.593756in}}%
\pgfpathlineto{\pgfqpoint{13.103715in}{0.593756in}}%
\pgfpathlineto{\pgfqpoint{13.103715in}{0.593756in}}%
\pgfpathclose%
\pgfpathmoveto{\pgfqpoint{13.191956in}{0.590387in}}%
\pgfpathlineto{\pgfqpoint{13.191956in}{0.580589in}}%
\pgfpathquadraticcurveto{\pgfqpoint{13.187579in}{0.582840in}}{\pgfqpoint{13.182842in}{0.583957in}}%
\pgfpathquadraticcurveto{\pgfqpoint{13.178123in}{0.585092in}}{\pgfqpoint{13.173044in}{0.585092in}}%
\pgfpathquadraticcurveto{\pgfqpoint{13.165334in}{0.585092in}}{\pgfqpoint{13.161480in}{0.582732in}}%
\pgfpathquadraticcurveto{\pgfqpoint{13.157625in}{0.580373in}}{\pgfqpoint{13.157625in}{0.575635in}}%
\pgfpathquadraticcurveto{\pgfqpoint{13.157625in}{0.572033in}}{\pgfqpoint{13.160381in}{0.569980in}}%
\pgfpathquadraticcurveto{\pgfqpoint{13.163137in}{0.567926in}}{\pgfqpoint{13.171477in}{0.566071in}}%
\pgfpathlineto{\pgfqpoint{13.175025in}{0.565278in}}%
\pgfpathquadraticcurveto{\pgfqpoint{13.186048in}{0.562919in}}{\pgfqpoint{13.190696in}{0.558614in}}%
\pgfpathquadraticcurveto{\pgfqpoint{13.195343in}{0.554309in}}{\pgfqpoint{13.195343in}{0.546600in}}%
\pgfpathquadraticcurveto{\pgfqpoint{13.195343in}{0.537810in}}{\pgfqpoint{13.188390in}{0.532676in}}%
\pgfpathquadraticcurveto{\pgfqpoint{13.181437in}{0.527561in}}{\pgfqpoint{13.169279in}{0.527561in}}%
\pgfpathquadraticcurveto{\pgfqpoint{13.164218in}{0.527561in}}{\pgfqpoint{13.158724in}{0.528552in}}%
\pgfpathquadraticcurveto{\pgfqpoint{13.153230in}{0.529542in}}{\pgfqpoint{13.147160in}{0.531506in}}%
\pgfpathlineto{\pgfqpoint{13.147160in}{0.542205in}}%
\pgfpathquadraticcurveto{\pgfqpoint{13.152906in}{0.539215in}}{\pgfqpoint{13.158472in}{0.537720in}}%
\pgfpathquadraticcurveto{\pgfqpoint{13.164038in}{0.536243in}}{\pgfqpoint{13.169513in}{0.536243in}}%
\pgfpathquadraticcurveto{\pgfqpoint{13.176826in}{0.536243in}}{\pgfqpoint{13.180753in}{0.538746in}}%
\pgfpathquadraticcurveto{\pgfqpoint{13.184697in}{0.541250in}}{\pgfqpoint{13.184697in}{0.545807in}}%
\pgfpathquadraticcurveto{\pgfqpoint{13.184697in}{0.550022in}}{\pgfqpoint{13.181852in}{0.552274in}}%
\pgfpathquadraticcurveto{\pgfqpoint{13.179024in}{0.554525in}}{\pgfqpoint{13.169387in}{0.556614in}}%
\pgfpathlineto{\pgfqpoint{13.165785in}{0.557461in}}%
\pgfpathquadraticcurveto{\pgfqpoint{13.156166in}{0.559478in}}{\pgfqpoint{13.151879in}{0.563675in}}%
\pgfpathquadraticcurveto{\pgfqpoint{13.147610in}{0.567872in}}{\pgfqpoint{13.147610in}{0.575185in}}%
\pgfpathquadraticcurveto{\pgfqpoint{13.147610in}{0.584083in}}{\pgfqpoint{13.153915in}{0.588910in}}%
\pgfpathquadraticcurveto{\pgfqpoint{13.160219in}{0.593756in}}{\pgfqpoint{13.171819in}{0.593756in}}%
\pgfpathquadraticcurveto{\pgfqpoint{13.177547in}{0.593756in}}{\pgfqpoint{13.182608in}{0.592909in}}%
\pgfpathquadraticcurveto{\pgfqpoint{13.187687in}{0.592080in}}{\pgfqpoint{13.191956in}{0.590387in}}%
\pgfpathlineto{\pgfqpoint{13.191956in}{0.590387in}}%
\pgfpathclose%
\pgfpathmoveto{\pgfqpoint{13.222073in}{0.610147in}}%
\pgfpathlineto{\pgfqpoint{13.222073in}{0.592243in}}%
\pgfpathlineto{\pgfqpoint{13.243399in}{0.592243in}}%
\pgfpathlineto{\pgfqpoint{13.243399in}{0.584191in}}%
\pgfpathlineto{\pgfqpoint{13.222073in}{0.584191in}}%
\pgfpathlineto{\pgfqpoint{13.222073in}{0.549968in}}%
\pgfpathquadraticcurveto{\pgfqpoint{13.222073in}{0.542259in}}{\pgfqpoint{13.224180in}{0.540061in}}%
\pgfpathquadraticcurveto{\pgfqpoint{13.226288in}{0.537864in}}{\pgfqpoint{13.232772in}{0.537864in}}%
\pgfpathlineto{\pgfqpoint{13.243399in}{0.537864in}}%
\pgfpathlineto{\pgfqpoint{13.243399in}{0.529200in}}%
\pgfpathlineto{\pgfqpoint{13.232772in}{0.529200in}}%
\pgfpathquadraticcurveto{\pgfqpoint{13.220776in}{0.529200in}}{\pgfqpoint{13.216219in}{0.533667in}}%
\pgfpathquadraticcurveto{\pgfqpoint{13.211662in}{0.538152in}}{\pgfqpoint{13.211662in}{0.549968in}}%
\pgfpathlineto{\pgfqpoint{13.211662in}{0.584191in}}%
\pgfpathlineto{\pgfqpoint{13.204061in}{0.584191in}}%
\pgfpathlineto{\pgfqpoint{13.204061in}{0.592243in}}%
\pgfpathlineto{\pgfqpoint{13.211662in}{0.592243in}}%
\pgfpathlineto{\pgfqpoint{13.211662in}{0.610147in}}%
\pgfpathlineto{\pgfqpoint{13.222073in}{0.610147in}}%
\pgfpathlineto{\pgfqpoint{13.222073in}{0.610147in}}%
\pgfpathclose%
\pgfpathmoveto{\pgfqpoint{13.306099in}{0.580138in}}%
\pgfpathquadraticcurveto{\pgfqpoint{13.309990in}{0.587127in}}{\pgfqpoint{13.315394in}{0.590441in}}%
\pgfpathquadraticcurveto{\pgfqpoint{13.320797in}{0.593756in}}{\pgfqpoint{13.328110in}{0.593756in}}%
\pgfpathquadraticcurveto{\pgfqpoint{13.337963in}{0.593756in}}{\pgfqpoint{13.343312in}{0.586857in}}%
\pgfpathquadraticcurveto{\pgfqpoint{13.348662in}{0.579976in}}{\pgfqpoint{13.348662in}{0.567260in}}%
\pgfpathlineto{\pgfqpoint{13.348662in}{0.529200in}}%
\pgfpathlineto{\pgfqpoint{13.338251in}{0.529200in}}%
\pgfpathlineto{\pgfqpoint{13.338251in}{0.566917in}}%
\pgfpathquadraticcurveto{\pgfqpoint{13.338251in}{0.575978in}}{\pgfqpoint{13.335027in}{0.580355in}}%
\pgfpathquadraticcurveto{\pgfqpoint{13.331821in}{0.584749in}}{\pgfqpoint{13.325246in}{0.584749in}}%
\pgfpathquadraticcurveto{\pgfqpoint{13.317195in}{0.584749in}}{\pgfqpoint{13.312512in}{0.579400in}}%
\pgfpathquadraticcurveto{\pgfqpoint{13.307847in}{0.574068in}}{\pgfqpoint{13.307847in}{0.564828in}}%
\pgfpathlineto{\pgfqpoint{13.307847in}{0.529200in}}%
\pgfpathlineto{\pgfqpoint{13.297436in}{0.529200in}}%
\pgfpathlineto{\pgfqpoint{13.297436in}{0.566917in}}%
\pgfpathquadraticcurveto{\pgfqpoint{13.297436in}{0.576032in}}{\pgfqpoint{13.294229in}{0.580391in}}%
\pgfpathquadraticcurveto{\pgfqpoint{13.291023in}{0.584749in}}{\pgfqpoint{13.284323in}{0.584749in}}%
\pgfpathquadraticcurveto{\pgfqpoint{13.276379in}{0.584749in}}{\pgfqpoint{13.271696in}{0.579382in}}%
\pgfpathquadraticcurveto{\pgfqpoint{13.267031in}{0.574014in}}{\pgfqpoint{13.267031in}{0.564828in}}%
\pgfpathlineto{\pgfqpoint{13.267031in}{0.529200in}}%
\pgfpathlineto{\pgfqpoint{13.256620in}{0.529200in}}%
\pgfpathlineto{\pgfqpoint{13.256620in}{0.592243in}}%
\pgfpathlineto{\pgfqpoint{13.267031in}{0.592243in}}%
\pgfpathlineto{\pgfqpoint{13.267031in}{0.582444in}}%
\pgfpathquadraticcurveto{\pgfqpoint{13.270579in}{0.588244in}}{\pgfqpoint{13.275533in}{0.591000in}}%
\pgfpathquadraticcurveto{\pgfqpoint{13.280486in}{0.593756in}}{\pgfqpoint{13.287295in}{0.593756in}}%
\pgfpathquadraticcurveto{\pgfqpoint{13.294175in}{0.593756in}}{\pgfqpoint{13.298985in}{0.590261in}}%
\pgfpathquadraticcurveto{\pgfqpoint{13.303794in}{0.586785in}}{\pgfqpoint{13.306099in}{0.580138in}}%
\pgfpathlineto{\pgfqpoint{13.306099in}{0.580138in}}%
\pgfpathclose%
\pgfpathmoveto{\pgfqpoint{13.423232in}{0.563315in}}%
\pgfpathlineto{\pgfqpoint{13.423232in}{0.558254in}}%
\pgfpathlineto{\pgfqpoint{13.375608in}{0.558254in}}%
\pgfpathquadraticcurveto{\pgfqpoint{13.376293in}{0.547554in}}{\pgfqpoint{13.382057in}{0.541953in}}%
\pgfpathquadraticcurveto{\pgfqpoint{13.387821in}{0.536351in}}{\pgfqpoint{13.398123in}{0.536351in}}%
\pgfpathquadraticcurveto{\pgfqpoint{13.404086in}{0.536351in}}{\pgfqpoint{13.409687in}{0.537810in}}%
\pgfpathquadraticcurveto{\pgfqpoint{13.415289in}{0.539269in}}{\pgfqpoint{13.420819in}{0.542205in}}%
\pgfpathlineto{\pgfqpoint{13.420819in}{0.532406in}}%
\pgfpathquadraticcurveto{\pgfqpoint{13.415235in}{0.530047in}}{\pgfqpoint{13.409381in}{0.528804in}}%
\pgfpathquadraticcurveto{\pgfqpoint{13.403527in}{0.527561in}}{\pgfqpoint{13.397511in}{0.527561in}}%
\pgfpathquadraticcurveto{\pgfqpoint{13.382417in}{0.527561in}}{\pgfqpoint{13.373609in}{0.536333in}}%
\pgfpathquadraticcurveto{\pgfqpoint{13.364801in}{0.545123in}}{\pgfqpoint{13.364801in}{0.560109in}}%
\pgfpathquadraticcurveto{\pgfqpoint{13.364801in}{0.575581in}}{\pgfqpoint{13.373159in}{0.584659in}}%
\pgfpathquadraticcurveto{\pgfqpoint{13.381516in}{0.593756in}}{\pgfqpoint{13.395710in}{0.593756in}}%
\pgfpathquadraticcurveto{\pgfqpoint{13.408426in}{0.593756in}}{\pgfqpoint{13.415829in}{0.585560in}}%
\pgfpathquadraticcurveto{\pgfqpoint{13.423232in}{0.577383in}}{\pgfqpoint{13.423232in}{0.563315in}}%
\pgfpathlineto{\pgfqpoint{13.423232in}{0.563315in}}%
\pgfpathclose%
\pgfpathmoveto{\pgfqpoint{13.412875in}{0.566359in}}%
\pgfpathquadraticcurveto{\pgfqpoint{13.412767in}{0.574843in}}{\pgfqpoint{13.408120in}{0.579904in}}%
\pgfpathquadraticcurveto{\pgfqpoint{13.403473in}{0.584984in}}{\pgfqpoint{13.395818in}{0.584984in}}%
\pgfpathquadraticcurveto{\pgfqpoint{13.387154in}{0.584984in}}{\pgfqpoint{13.381949in}{0.580084in}}%
\pgfpathquadraticcurveto{\pgfqpoint{13.376743in}{0.575185in}}{\pgfqpoint{13.375951in}{0.566287in}}%
\pgfpathlineto{\pgfqpoint{13.412875in}{0.566359in}}%
\pgfpathlineto{\pgfqpoint{13.412875in}{0.566359in}}%
\pgfpathclose%
\pgfpathmoveto{\pgfqpoint{13.492651in}{0.567260in}}%
\pgfpathlineto{\pgfqpoint{13.492651in}{0.529200in}}%
\pgfpathlineto{\pgfqpoint{13.482294in}{0.529200in}}%
\pgfpathlineto{\pgfqpoint{13.482294in}{0.566917in}}%
\pgfpathquadraticcurveto{\pgfqpoint{13.482294in}{0.575869in}}{\pgfqpoint{13.478800in}{0.580300in}}%
\pgfpathquadraticcurveto{\pgfqpoint{13.475306in}{0.584749in}}{\pgfqpoint{13.468335in}{0.584749in}}%
\pgfpathquadraticcurveto{\pgfqpoint{13.459941in}{0.584749in}}{\pgfqpoint{13.455096in}{0.579400in}}%
\pgfpathquadraticcurveto{\pgfqpoint{13.450251in}{0.574068in}}{\pgfqpoint{13.450251in}{0.564828in}}%
\pgfpathlineto{\pgfqpoint{13.450251in}{0.529200in}}%
\pgfpathlineto{\pgfqpoint{13.439840in}{0.529200in}}%
\pgfpathlineto{\pgfqpoint{13.439840in}{0.592243in}}%
\pgfpathlineto{\pgfqpoint{13.450251in}{0.592243in}}%
\pgfpathlineto{\pgfqpoint{13.450251in}{0.582444in}}%
\pgfpathquadraticcurveto{\pgfqpoint{13.453979in}{0.588136in}}{\pgfqpoint{13.459005in}{0.590946in}}%
\pgfpathquadraticcurveto{\pgfqpoint{13.464048in}{0.593756in}}{\pgfqpoint{13.470640in}{0.593756in}}%
\pgfpathquadraticcurveto{\pgfqpoint{13.481502in}{0.593756in}}{\pgfqpoint{13.487067in}{0.587037in}}%
\pgfpathquadraticcurveto{\pgfqpoint{13.492651in}{0.580318in}}{\pgfqpoint{13.492651in}{0.567260in}}%
\pgfpathlineto{\pgfqpoint{13.492651in}{0.567260in}}%
\pgfpathclose%
\pgfpathmoveto{\pgfqpoint{13.523542in}{0.610147in}}%
\pgfpathlineto{\pgfqpoint{13.523542in}{0.592243in}}%
\pgfpathlineto{\pgfqpoint{13.544868in}{0.592243in}}%
\pgfpathlineto{\pgfqpoint{13.544868in}{0.584191in}}%
\pgfpathlineto{\pgfqpoint{13.523542in}{0.584191in}}%
\pgfpathlineto{\pgfqpoint{13.523542in}{0.549968in}}%
\pgfpathquadraticcurveto{\pgfqpoint{13.523542in}{0.542259in}}{\pgfqpoint{13.525650in}{0.540061in}}%
\pgfpathquadraticcurveto{\pgfqpoint{13.527757in}{0.537864in}}{\pgfqpoint{13.534241in}{0.537864in}}%
\pgfpathlineto{\pgfqpoint{13.544868in}{0.537864in}}%
\pgfpathlineto{\pgfqpoint{13.544868in}{0.529200in}}%
\pgfpathlineto{\pgfqpoint{13.534241in}{0.529200in}}%
\pgfpathquadraticcurveto{\pgfqpoint{13.522245in}{0.529200in}}{\pgfqpoint{13.517688in}{0.533667in}}%
\pgfpathquadraticcurveto{\pgfqpoint{13.513131in}{0.538152in}}{\pgfqpoint{13.513131in}{0.549968in}}%
\pgfpathlineto{\pgfqpoint{13.513131in}{0.584191in}}%
\pgfpathlineto{\pgfqpoint{13.505530in}{0.584191in}}%
\pgfpathlineto{\pgfqpoint{13.505530in}{0.592243in}}%
\pgfpathlineto{\pgfqpoint{13.513131in}{0.592243in}}%
\pgfpathlineto{\pgfqpoint{13.513131in}{0.610147in}}%
\pgfpathlineto{\pgfqpoint{13.523542in}{0.610147in}}%
\pgfpathlineto{\pgfqpoint{13.523542in}{0.610147in}}%
\pgfpathclose%
\pgfpathmoveto{\pgfqpoint{13.559945in}{0.543502in}}%
\pgfpathlineto{\pgfqpoint{13.571833in}{0.543502in}}%
\pgfpathlineto{\pgfqpoint{13.571833in}{0.529200in}}%
\pgfpathlineto{\pgfqpoint{13.559945in}{0.529200in}}%
\pgfpathlineto{\pgfqpoint{13.559945in}{0.543502in}}%
\pgfpathlineto{\pgfqpoint{13.559945in}{0.543502in}}%
\pgfpathclose%
\pgfusepath{fill}%
\end{pgfscope}%
\begin{pgfscope}%
\pgfsetbuttcap%
\pgfsetroundjoin%
\definecolor{currentfill}{rgb}{0.309804,0.372549,0.419608}%
\pgfsetfillcolor{currentfill}%
\pgfsetlinewidth{0.000000pt}%
\definecolor{currentstroke}{rgb}{0.309804,0.372549,0.419608}%
\pgfsetstrokecolor{currentstroke}%
\pgfsetdash{}{0pt}%
\pgfpathmoveto{\pgfqpoint{0.946677in}{0.420296in}}%
\pgfpathlineto{\pgfqpoint{0.946677in}{0.354948in}}%
\pgfpathlineto{\pgfqpoint{0.960421in}{0.354948in}}%
\pgfpathquadraticcurveto{\pgfqpoint{0.977820in}{0.354948in}}{\pgfqpoint{0.985890in}{0.362820in}}%
\pgfpathquadraticcurveto{\pgfqpoint{0.993959in}{0.370709in}}{\pgfqpoint{0.993959in}{0.387712in}}%
\pgfpathquadraticcurveto{\pgfqpoint{0.993959in}{0.404590in}}{\pgfqpoint{0.985890in}{0.412443in}}%
\pgfpathquadraticcurveto{\pgfqpoint{0.977820in}{0.420296in}}{\pgfqpoint{0.960421in}{0.420296in}}%
\pgfpathlineto{\pgfqpoint{0.946677in}{0.420296in}}%
\pgfpathlineto{\pgfqpoint{0.946677in}{0.420296in}}%
\pgfpathclose%
\pgfpathmoveto{\pgfqpoint{0.935312in}{0.429645in}}%
\pgfpathlineto{\pgfqpoint{0.958673in}{0.429645in}}%
\pgfpathquadraticcurveto{\pgfqpoint{0.983098in}{0.429645in}}{\pgfqpoint{0.994518in}{0.419486in}}%
\pgfpathquadraticcurveto{\pgfqpoint{1.005955in}{0.409327in}}{\pgfqpoint{1.005955in}{0.387712in}}%
\pgfpathquadraticcurveto{\pgfqpoint{1.005955in}{0.365972in}}{\pgfqpoint{0.994464in}{0.355777in}}%
\pgfpathquadraticcurveto{\pgfqpoint{0.982990in}{0.345600in}}{\pgfqpoint{0.958673in}{0.345600in}}%
\pgfpathlineto{\pgfqpoint{0.935312in}{0.345600in}}%
\pgfpathlineto{\pgfqpoint{0.935312in}{0.429645in}}%
\pgfpathlineto{\pgfqpoint{0.935312in}{0.429645in}}%
\pgfpathclose%
\pgfpathmoveto{\pgfqpoint{1.052283in}{0.377283in}}%
\pgfpathquadraticcurveto{\pgfqpoint{1.039728in}{0.377283in}}{\pgfqpoint{1.034883in}{0.374419in}}%
\pgfpathquadraticcurveto{\pgfqpoint{1.030038in}{0.371556in}}{\pgfqpoint{1.030038in}{0.364621in}}%
\pgfpathquadraticcurveto{\pgfqpoint{1.030038in}{0.359109in}}{\pgfqpoint{1.033676in}{0.355867in}}%
\pgfpathquadraticcurveto{\pgfqpoint{1.037314in}{0.352643in}}{\pgfqpoint{1.043547in}{0.352643in}}%
\pgfpathquadraticcurveto{\pgfqpoint{1.052174in}{0.352643in}}{\pgfqpoint{1.057380in}{0.358749in}}%
\pgfpathquadraticcurveto{\pgfqpoint{1.062586in}{0.364855in}}{\pgfqpoint{1.062586in}{0.374978in}}%
\pgfpathlineto{\pgfqpoint{1.062586in}{0.377283in}}%
\pgfpathlineto{\pgfqpoint{1.052283in}{0.377283in}}%
\pgfpathlineto{\pgfqpoint{1.052283in}{0.377283in}}%
\pgfpathclose%
\pgfpathmoveto{\pgfqpoint{1.072942in}{0.381570in}}%
\pgfpathlineto{\pgfqpoint{1.072942in}{0.345600in}}%
\pgfpathlineto{\pgfqpoint{1.062586in}{0.345600in}}%
\pgfpathlineto{\pgfqpoint{1.062586in}{0.355164in}}%
\pgfpathquadraticcurveto{\pgfqpoint{1.059037in}{0.349437in}}{\pgfqpoint{1.053742in}{0.346699in}}%
\pgfpathquadraticcurveto{\pgfqpoint{1.048446in}{0.343961in}}{\pgfqpoint{1.040791in}{0.343961in}}%
\pgfpathquadraticcurveto{\pgfqpoint{1.031118in}{0.343961in}}{\pgfqpoint{1.025390in}{0.349401in}}%
\pgfpathquadraticcurveto{\pgfqpoint{1.019681in}{0.354840in}}{\pgfqpoint{1.019681in}{0.363954in}}%
\pgfpathquadraticcurveto{\pgfqpoint{1.019681in}{0.374582in}}{\pgfqpoint{1.026795in}{0.379985in}}%
\pgfpathquadraticcurveto{\pgfqpoint{1.033928in}{0.385389in}}{\pgfqpoint{1.048050in}{0.385389in}}%
\pgfpathlineto{\pgfqpoint{1.062586in}{0.385389in}}%
\pgfpathlineto{\pgfqpoint{1.062586in}{0.386416in}}%
\pgfpathquadraticcurveto{\pgfqpoint{1.062586in}{0.393566in}}{\pgfqpoint{1.057884in}{0.397475in}}%
\pgfpathquadraticcurveto{\pgfqpoint{1.053183in}{0.401384in}}{\pgfqpoint{1.044681in}{0.401384in}}%
\pgfpathquadraticcurveto{\pgfqpoint{1.039278in}{0.401384in}}{\pgfqpoint{1.034144in}{0.400087in}}%
\pgfpathquadraticcurveto{\pgfqpoint{1.029029in}{0.398790in}}{\pgfqpoint{1.024310in}{0.396196in}}%
\pgfpathlineto{\pgfqpoint{1.024310in}{0.405779in}}%
\pgfpathquadraticcurveto{\pgfqpoint{1.029984in}{0.407976in}}{\pgfqpoint{1.035333in}{0.409057in}}%
\pgfpathquadraticcurveto{\pgfqpoint{1.040683in}{0.410156in}}{\pgfqpoint{1.045744in}{0.410156in}}%
\pgfpathquadraticcurveto{\pgfqpoint{1.059433in}{0.410156in}}{\pgfqpoint{1.066188in}{0.403059in}}%
\pgfpathquadraticcurveto{\pgfqpoint{1.072942in}{0.395980in}}{\pgfqpoint{1.072942in}{0.381570in}}%
\pgfpathlineto{\pgfqpoint{1.072942in}{0.381570in}}%
\pgfpathclose%
\pgfpathmoveto{\pgfqpoint{1.134454in}{0.406787in}}%
\pgfpathlineto{\pgfqpoint{1.134454in}{0.396989in}}%
\pgfpathquadraticcurveto{\pgfqpoint{1.130077in}{0.399240in}}{\pgfqpoint{1.125340in}{0.400357in}}%
\pgfpathquadraticcurveto{\pgfqpoint{1.120621in}{0.401492in}}{\pgfqpoint{1.115541in}{0.401492in}}%
\pgfpathquadraticcurveto{\pgfqpoint{1.107832in}{0.401492in}}{\pgfqpoint{1.103977in}{0.399132in}}%
\pgfpathquadraticcurveto{\pgfqpoint{1.100123in}{0.396773in}}{\pgfqpoint{1.100123in}{0.392035in}}%
\pgfpathquadraticcurveto{\pgfqpoint{1.100123in}{0.388433in}}{\pgfqpoint{1.102879in}{0.386380in}}%
\pgfpathquadraticcurveto{\pgfqpoint{1.105635in}{0.384326in}}{\pgfqpoint{1.113974in}{0.382471in}}%
\pgfpathlineto{\pgfqpoint{1.117523in}{0.381678in}}%
\pgfpathquadraticcurveto{\pgfqpoint{1.128546in}{0.379319in}}{\pgfqpoint{1.133193in}{0.375014in}}%
\pgfpathquadraticcurveto{\pgfqpoint{1.137840in}{0.370709in}}{\pgfqpoint{1.137840in}{0.363000in}}%
\pgfpathquadraticcurveto{\pgfqpoint{1.137840in}{0.354210in}}{\pgfqpoint{1.130888in}{0.349076in}}%
\pgfpathquadraticcurveto{\pgfqpoint{1.123935in}{0.343961in}}{\pgfqpoint{1.111777in}{0.343961in}}%
\pgfpathquadraticcurveto{\pgfqpoint{1.106715in}{0.343961in}}{\pgfqpoint{1.101222in}{0.344952in}}%
\pgfpathquadraticcurveto{\pgfqpoint{1.095728in}{0.345942in}}{\pgfqpoint{1.089658in}{0.347906in}}%
\pgfpathlineto{\pgfqpoint{1.089658in}{0.358605in}}%
\pgfpathquadraticcurveto{\pgfqpoint{1.095404in}{0.355615in}}{\pgfqpoint{1.100969in}{0.354120in}}%
\pgfpathquadraticcurveto{\pgfqpoint{1.106535in}{0.352643in}}{\pgfqpoint{1.112011in}{0.352643in}}%
\pgfpathquadraticcurveto{\pgfqpoint{1.119324in}{0.352643in}}{\pgfqpoint{1.123250in}{0.355146in}}%
\pgfpathquadraticcurveto{\pgfqpoint{1.127195in}{0.357650in}}{\pgfqpoint{1.127195in}{0.362207in}}%
\pgfpathquadraticcurveto{\pgfqpoint{1.127195in}{0.366422in}}{\pgfqpoint{1.124349in}{0.368674in}}%
\pgfpathquadraticcurveto{\pgfqpoint{1.121521in}{0.370925in}}{\pgfqpoint{1.111885in}{0.373014in}}%
\pgfpathlineto{\pgfqpoint{1.108282in}{0.373861in}}%
\pgfpathquadraticcurveto{\pgfqpoint{1.098664in}{0.375878in}}{\pgfqpoint{1.094377in}{0.380075in}}%
\pgfpathquadraticcurveto{\pgfqpoint{1.090108in}{0.384272in}}{\pgfqpoint{1.090108in}{0.391585in}}%
\pgfpathquadraticcurveto{\pgfqpoint{1.090108in}{0.400483in}}{\pgfqpoint{1.096412in}{0.405310in}}%
\pgfpathquadraticcurveto{\pgfqpoint{1.102717in}{0.410156in}}{\pgfqpoint{1.114316in}{0.410156in}}%
\pgfpathquadraticcurveto{\pgfqpoint{1.120044in}{0.410156in}}{\pgfqpoint{1.125106in}{0.409309in}}%
\pgfpathquadraticcurveto{\pgfqpoint{1.130185in}{0.408480in}}{\pgfqpoint{1.134454in}{0.406787in}}%
\pgfpathlineto{\pgfqpoint{1.134454in}{0.406787in}}%
\pgfpathclose%
\pgfpathmoveto{\pgfqpoint{1.206737in}{0.383660in}}%
\pgfpathlineto{\pgfqpoint{1.206737in}{0.345600in}}%
\pgfpathlineto{\pgfqpoint{1.196380in}{0.345600in}}%
\pgfpathlineto{\pgfqpoint{1.196380in}{0.383317in}}%
\pgfpathquadraticcurveto{\pgfqpoint{1.196380in}{0.392269in}}{\pgfqpoint{1.192885in}{0.396700in}}%
\pgfpathquadraticcurveto{\pgfqpoint{1.189391in}{0.401149in}}{\pgfqpoint{1.182420in}{0.401149in}}%
\pgfpathquadraticcurveto{\pgfqpoint{1.174027in}{0.401149in}}{\pgfqpoint{1.169181in}{0.395800in}}%
\pgfpathquadraticcurveto{\pgfqpoint{1.164336in}{0.390468in}}{\pgfqpoint{1.164336in}{0.381228in}}%
\pgfpathlineto{\pgfqpoint{1.164336in}{0.345600in}}%
\pgfpathlineto{\pgfqpoint{1.153925in}{0.345600in}}%
\pgfpathlineto{\pgfqpoint{1.153925in}{0.433193in}}%
\pgfpathlineto{\pgfqpoint{1.164336in}{0.433193in}}%
\pgfpathlineto{\pgfqpoint{1.164336in}{0.398844in}}%
\pgfpathquadraticcurveto{\pgfqpoint{1.168065in}{0.404536in}}{\pgfqpoint{1.173090in}{0.407346in}}%
\pgfpathquadraticcurveto{\pgfqpoint{1.178133in}{0.410156in}}{\pgfqpoint{1.184726in}{0.410156in}}%
\pgfpathquadraticcurveto{\pgfqpoint{1.195587in}{0.410156in}}{\pgfqpoint{1.201153in}{0.403437in}}%
\pgfpathquadraticcurveto{\pgfqpoint{1.206737in}{0.396718in}}{\pgfqpoint{1.206737in}{0.383660in}}%
\pgfpathlineto{\pgfqpoint{1.206737in}{0.383660in}}%
\pgfpathclose%
\pgfpathmoveto{\pgfqpoint{1.281307in}{0.379715in}}%
\pgfpathlineto{\pgfqpoint{1.281307in}{0.374654in}}%
\pgfpathlineto{\pgfqpoint{1.233683in}{0.374654in}}%
\pgfpathquadraticcurveto{\pgfqpoint{1.234367in}{0.363954in}}{\pgfqpoint{1.240131in}{0.358353in}}%
\pgfpathquadraticcurveto{\pgfqpoint{1.245895in}{0.352751in}}{\pgfqpoint{1.256198in}{0.352751in}}%
\pgfpathquadraticcurveto{\pgfqpoint{1.262160in}{0.352751in}}{\pgfqpoint{1.267762in}{0.354210in}}%
\pgfpathquadraticcurveto{\pgfqpoint{1.273364in}{0.355669in}}{\pgfqpoint{1.278893in}{0.358605in}}%
\pgfpathlineto{\pgfqpoint{1.278893in}{0.348806in}}%
\pgfpathquadraticcurveto{\pgfqpoint{1.273310in}{0.346447in}}{\pgfqpoint{1.267456in}{0.345204in}}%
\pgfpathquadraticcurveto{\pgfqpoint{1.261602in}{0.343961in}}{\pgfqpoint{1.255586in}{0.343961in}}%
\pgfpathquadraticcurveto{\pgfqpoint{1.240492in}{0.343961in}}{\pgfqpoint{1.231684in}{0.352733in}}%
\pgfpathquadraticcurveto{\pgfqpoint{1.222876in}{0.361523in}}{\pgfqpoint{1.222876in}{0.376509in}}%
\pgfpathquadraticcurveto{\pgfqpoint{1.222876in}{0.391981in}}{\pgfqpoint{1.231233in}{0.401059in}}%
\pgfpathquadraticcurveto{\pgfqpoint{1.239591in}{0.410156in}}{\pgfqpoint{1.253785in}{0.410156in}}%
\pgfpathquadraticcurveto{\pgfqpoint{1.266501in}{0.410156in}}{\pgfqpoint{1.273904in}{0.401960in}}%
\pgfpathquadraticcurveto{\pgfqpoint{1.281307in}{0.393783in}}{\pgfqpoint{1.281307in}{0.379715in}}%
\pgfpathlineto{\pgfqpoint{1.281307in}{0.379715in}}%
\pgfpathclose%
\pgfpathmoveto{\pgfqpoint{1.270950in}{0.382759in}}%
\pgfpathquadraticcurveto{\pgfqpoint{1.270842in}{0.391243in}}{\pgfqpoint{1.266195in}{0.396304in}}%
\pgfpathquadraticcurveto{\pgfqpoint{1.261548in}{0.401384in}}{\pgfqpoint{1.253893in}{0.401384in}}%
\pgfpathquadraticcurveto{\pgfqpoint{1.245229in}{0.401384in}}{\pgfqpoint{1.240023in}{0.396484in}}%
\pgfpathquadraticcurveto{\pgfqpoint{1.234818in}{0.391585in}}{\pgfqpoint{1.234025in}{0.382687in}}%
\pgfpathlineto{\pgfqpoint{1.270950in}{0.382759in}}%
\pgfpathlineto{\pgfqpoint{1.270950in}{0.382759in}}%
\pgfpathclose%
\pgfpathmoveto{\pgfqpoint{1.339793in}{0.399078in}}%
\pgfpathlineto{\pgfqpoint{1.339793in}{0.433193in}}%
\pgfpathlineto{\pgfqpoint{1.350150in}{0.433193in}}%
\pgfpathlineto{\pgfqpoint{1.350150in}{0.345600in}}%
\pgfpathlineto{\pgfqpoint{1.339793in}{0.345600in}}%
\pgfpathlineto{\pgfqpoint{1.339793in}{0.355056in}}%
\pgfpathquadraticcurveto{\pgfqpoint{1.336532in}{0.349437in}}{\pgfqpoint{1.331543in}{0.346699in}}%
\pgfpathquadraticcurveto{\pgfqpoint{1.326572in}{0.343961in}}{\pgfqpoint{1.319583in}{0.343961in}}%
\pgfpathquadraticcurveto{\pgfqpoint{1.308163in}{0.343961in}}{\pgfqpoint{1.300976in}{0.353075in}}%
\pgfpathquadraticcurveto{\pgfqpoint{1.293808in}{0.362207in}}{\pgfqpoint{1.293808in}{0.377067in}}%
\pgfpathquadraticcurveto{\pgfqpoint{1.293808in}{0.391927in}}{\pgfqpoint{1.300976in}{0.401041in}}%
\pgfpathquadraticcurveto{\pgfqpoint{1.308163in}{0.410156in}}{\pgfqpoint{1.319583in}{0.410156in}}%
\pgfpathquadraticcurveto{\pgfqpoint{1.326572in}{0.410156in}}{\pgfqpoint{1.331543in}{0.407418in}}%
\pgfpathquadraticcurveto{\pgfqpoint{1.336532in}{0.404698in}}{\pgfqpoint{1.339793in}{0.399078in}}%
\pgfpathlineto{\pgfqpoint{1.339793in}{0.399078in}}%
\pgfpathclose%
\pgfpathmoveto{\pgfqpoint{1.304507in}{0.377067in}}%
\pgfpathquadraticcurveto{\pgfqpoint{1.304507in}{0.365648in}}{\pgfqpoint{1.309208in}{0.359145in}}%
\pgfpathquadraticcurveto{\pgfqpoint{1.313909in}{0.352643in}}{\pgfqpoint{1.322123in}{0.352643in}}%
\pgfpathquadraticcurveto{\pgfqpoint{1.330336in}{0.352643in}}{\pgfqpoint{1.335055in}{0.359145in}}%
\pgfpathquadraticcurveto{\pgfqpoint{1.339793in}{0.365648in}}{\pgfqpoint{1.339793in}{0.377067in}}%
\pgfpathquadraticcurveto{\pgfqpoint{1.339793in}{0.388487in}}{\pgfqpoint{1.335055in}{0.394989in}}%
\pgfpathquadraticcurveto{\pgfqpoint{1.330336in}{0.401492in}}{\pgfqpoint{1.322123in}{0.401492in}}%
\pgfpathquadraticcurveto{\pgfqpoint{1.313909in}{0.401492in}}{\pgfqpoint{1.309208in}{0.394989in}}%
\pgfpathquadraticcurveto{\pgfqpoint{1.304507in}{0.388487in}}{\pgfqpoint{1.304507in}{0.377067in}}%
\pgfpathlineto{\pgfqpoint{1.304507in}{0.377067in}}%
\pgfpathclose%
\pgfpathmoveto{\pgfqpoint{1.408131in}{0.433193in}}%
\pgfpathlineto{\pgfqpoint{1.418488in}{0.433193in}}%
\pgfpathlineto{\pgfqpoint{1.418488in}{0.345600in}}%
\pgfpathlineto{\pgfqpoint{1.408131in}{0.345600in}}%
\pgfpathlineto{\pgfqpoint{1.408131in}{0.433193in}}%
\pgfpathlineto{\pgfqpoint{1.408131in}{0.433193in}}%
\pgfpathclose%
\pgfpathmoveto{\pgfqpoint{1.440156in}{0.408643in}}%
\pgfpathlineto{\pgfqpoint{1.450513in}{0.408643in}}%
\pgfpathlineto{\pgfqpoint{1.450513in}{0.345600in}}%
\pgfpathlineto{\pgfqpoint{1.440156in}{0.345600in}}%
\pgfpathlineto{\pgfqpoint{1.440156in}{0.408643in}}%
\pgfpathlineto{\pgfqpoint{1.440156in}{0.408643in}}%
\pgfpathclose%
\pgfpathmoveto{\pgfqpoint{1.440156in}{0.433193in}}%
\pgfpathlineto{\pgfqpoint{1.450513in}{0.433193in}}%
\pgfpathlineto{\pgfqpoint{1.450513in}{0.420062in}}%
\pgfpathlineto{\pgfqpoint{1.440156in}{0.420062in}}%
\pgfpathlineto{\pgfqpoint{1.440156in}{0.433193in}}%
\pgfpathlineto{\pgfqpoint{1.440156in}{0.433193in}}%
\pgfpathclose%
\pgfpathmoveto{\pgfqpoint{1.524597in}{0.383660in}}%
\pgfpathlineto{\pgfqpoint{1.524597in}{0.345600in}}%
\pgfpathlineto{\pgfqpoint{1.514240in}{0.345600in}}%
\pgfpathlineto{\pgfqpoint{1.514240in}{0.383317in}}%
\pgfpathquadraticcurveto{\pgfqpoint{1.514240in}{0.392269in}}{\pgfqpoint{1.510746in}{0.396700in}}%
\pgfpathquadraticcurveto{\pgfqpoint{1.507252in}{0.401149in}}{\pgfqpoint{1.500281in}{0.401149in}}%
\pgfpathquadraticcurveto{\pgfqpoint{1.491887in}{0.401149in}}{\pgfqpoint{1.487042in}{0.395800in}}%
\pgfpathquadraticcurveto{\pgfqpoint{1.482197in}{0.390468in}}{\pgfqpoint{1.482197in}{0.381228in}}%
\pgfpathlineto{\pgfqpoint{1.482197in}{0.345600in}}%
\pgfpathlineto{\pgfqpoint{1.471786in}{0.345600in}}%
\pgfpathlineto{\pgfqpoint{1.471786in}{0.408643in}}%
\pgfpathlineto{\pgfqpoint{1.482197in}{0.408643in}}%
\pgfpathlineto{\pgfqpoint{1.482197in}{0.398844in}}%
\pgfpathquadraticcurveto{\pgfqpoint{1.485925in}{0.404536in}}{\pgfqpoint{1.490951in}{0.407346in}}%
\pgfpathquadraticcurveto{\pgfqpoint{1.495994in}{0.410156in}}{\pgfqpoint{1.502586in}{0.410156in}}%
\pgfpathquadraticcurveto{\pgfqpoint{1.513448in}{0.410156in}}{\pgfqpoint{1.519013in}{0.403437in}}%
\pgfpathquadraticcurveto{\pgfqpoint{1.524597in}{0.396718in}}{\pgfqpoint{1.524597in}{0.383660in}}%
\pgfpathlineto{\pgfqpoint{1.524597in}{0.383660in}}%
\pgfpathclose%
\pgfpathmoveto{\pgfqpoint{1.599168in}{0.379715in}}%
\pgfpathlineto{\pgfqpoint{1.599168in}{0.374654in}}%
\pgfpathlineto{\pgfqpoint{1.551543in}{0.374654in}}%
\pgfpathquadraticcurveto{\pgfqpoint{1.552228in}{0.363954in}}{\pgfqpoint{1.557992in}{0.358353in}}%
\pgfpathquadraticcurveto{\pgfqpoint{1.563756in}{0.352751in}}{\pgfqpoint{1.574059in}{0.352751in}}%
\pgfpathquadraticcurveto{\pgfqpoint{1.580021in}{0.352751in}}{\pgfqpoint{1.585622in}{0.354210in}}%
\pgfpathquadraticcurveto{\pgfqpoint{1.591224in}{0.355669in}}{\pgfqpoint{1.596754in}{0.358605in}}%
\pgfpathlineto{\pgfqpoint{1.596754in}{0.348806in}}%
\pgfpathquadraticcurveto{\pgfqpoint{1.591170in}{0.346447in}}{\pgfqpoint{1.585316in}{0.345204in}}%
\pgfpathquadraticcurveto{\pgfqpoint{1.579462in}{0.343961in}}{\pgfqpoint{1.573446in}{0.343961in}}%
\pgfpathquadraticcurveto{\pgfqpoint{1.558352in}{0.343961in}}{\pgfqpoint{1.549544in}{0.352733in}}%
\pgfpathquadraticcurveto{\pgfqpoint{1.540736in}{0.361523in}}{\pgfqpoint{1.540736in}{0.376509in}}%
\pgfpathquadraticcurveto{\pgfqpoint{1.540736in}{0.391981in}}{\pgfqpoint{1.549094in}{0.401059in}}%
\pgfpathquadraticcurveto{\pgfqpoint{1.557451in}{0.410156in}}{\pgfqpoint{1.571645in}{0.410156in}}%
\pgfpathquadraticcurveto{\pgfqpoint{1.584362in}{0.410156in}}{\pgfqpoint{1.591765in}{0.401960in}}%
\pgfpathquadraticcurveto{\pgfqpoint{1.599168in}{0.393783in}}{\pgfqpoint{1.599168in}{0.379715in}}%
\pgfpathlineto{\pgfqpoint{1.599168in}{0.379715in}}%
\pgfpathclose%
\pgfpathmoveto{\pgfqpoint{1.588811in}{0.382759in}}%
\pgfpathquadraticcurveto{\pgfqpoint{1.588702in}{0.391243in}}{\pgfqpoint{1.584055in}{0.396304in}}%
\pgfpathquadraticcurveto{\pgfqpoint{1.579408in}{0.401384in}}{\pgfqpoint{1.571753in}{0.401384in}}%
\pgfpathquadraticcurveto{\pgfqpoint{1.563089in}{0.401384in}}{\pgfqpoint{1.557884in}{0.396484in}}%
\pgfpathquadraticcurveto{\pgfqpoint{1.552678in}{0.391585in}}{\pgfqpoint{1.551886in}{0.382687in}}%
\pgfpathlineto{\pgfqpoint{1.588811in}{0.382759in}}%
\pgfpathlineto{\pgfqpoint{1.588811in}{0.382759in}}%
\pgfpathclose%
\pgfpathmoveto{\pgfqpoint{1.618819in}{0.359902in}}%
\pgfpathlineto{\pgfqpoint{1.630689in}{0.359902in}}%
\pgfpathlineto{\pgfqpoint{1.630689in}{0.345600in}}%
\pgfpathlineto{\pgfqpoint{1.618819in}{0.345600in}}%
\pgfpathlineto{\pgfqpoint{1.618819in}{0.359902in}}%
\pgfpathlineto{\pgfqpoint{1.618819in}{0.359902in}}%
\pgfpathclose%
\pgfpathmoveto{\pgfqpoint{1.618819in}{0.405202in}}%
\pgfpathlineto{\pgfqpoint{1.630689in}{0.405202in}}%
\pgfpathlineto{\pgfqpoint{1.630689in}{0.390919in}}%
\pgfpathlineto{\pgfqpoint{1.618819in}{0.390919in}}%
\pgfpathlineto{\pgfqpoint{1.618819in}{0.405202in}}%
\pgfpathlineto{\pgfqpoint{1.618819in}{0.405202in}}%
\pgfpathclose%
\pgfpathmoveto{\pgfqpoint{1.745570in}{0.379715in}}%
\pgfpathlineto{\pgfqpoint{1.745570in}{0.374654in}}%
\pgfpathlineto{\pgfqpoint{1.697946in}{0.374654in}}%
\pgfpathquadraticcurveto{\pgfqpoint{1.698631in}{0.363954in}}{\pgfqpoint{1.704395in}{0.358353in}}%
\pgfpathquadraticcurveto{\pgfqpoint{1.710158in}{0.352751in}}{\pgfqpoint{1.720461in}{0.352751in}}%
\pgfpathquadraticcurveto{\pgfqpoint{1.726423in}{0.352751in}}{\pgfqpoint{1.732025in}{0.354210in}}%
\pgfpathquadraticcurveto{\pgfqpoint{1.737627in}{0.355669in}}{\pgfqpoint{1.743157in}{0.358605in}}%
\pgfpathlineto{\pgfqpoint{1.743157in}{0.348806in}}%
\pgfpathquadraticcurveto{\pgfqpoint{1.737573in}{0.346447in}}{\pgfqpoint{1.731719in}{0.345204in}}%
\pgfpathquadraticcurveto{\pgfqpoint{1.725865in}{0.343961in}}{\pgfqpoint{1.719849in}{0.343961in}}%
\pgfpathquadraticcurveto{\pgfqpoint{1.704755in}{0.343961in}}{\pgfqpoint{1.695947in}{0.352733in}}%
\pgfpathquadraticcurveto{\pgfqpoint{1.687139in}{0.361523in}}{\pgfqpoint{1.687139in}{0.376509in}}%
\pgfpathquadraticcurveto{\pgfqpoint{1.687139in}{0.391981in}}{\pgfqpoint{1.695497in}{0.401059in}}%
\pgfpathquadraticcurveto{\pgfqpoint{1.703854in}{0.410156in}}{\pgfqpoint{1.718048in}{0.410156in}}%
\pgfpathquadraticcurveto{\pgfqpoint{1.730764in}{0.410156in}}{\pgfqpoint{1.738167in}{0.401960in}}%
\pgfpathquadraticcurveto{\pgfqpoint{1.745570in}{0.393783in}}{\pgfqpoint{1.745570in}{0.379715in}}%
\pgfpathlineto{\pgfqpoint{1.745570in}{0.379715in}}%
\pgfpathclose%
\pgfpathmoveto{\pgfqpoint{1.735213in}{0.382759in}}%
\pgfpathquadraticcurveto{\pgfqpoint{1.735105in}{0.391243in}}{\pgfqpoint{1.730458in}{0.396304in}}%
\pgfpathquadraticcurveto{\pgfqpoint{1.725811in}{0.401384in}}{\pgfqpoint{1.718156in}{0.401384in}}%
\pgfpathquadraticcurveto{\pgfqpoint{1.709492in}{0.401384in}}{\pgfqpoint{1.704286in}{0.396484in}}%
\pgfpathquadraticcurveto{\pgfqpoint{1.699081in}{0.391585in}}{\pgfqpoint{1.698288in}{0.382687in}}%
\pgfpathlineto{\pgfqpoint{1.735213in}{0.382759in}}%
\pgfpathlineto{\pgfqpoint{1.735213in}{0.382759in}}%
\pgfpathclose%
\pgfpathmoveto{\pgfqpoint{1.768770in}{0.377067in}}%
\pgfpathquadraticcurveto{\pgfqpoint{1.768770in}{0.365648in}}{\pgfqpoint{1.773471in}{0.359145in}}%
\pgfpathquadraticcurveto{\pgfqpoint{1.778172in}{0.352643in}}{\pgfqpoint{1.786386in}{0.352643in}}%
\pgfpathquadraticcurveto{\pgfqpoint{1.794599in}{0.352643in}}{\pgfqpoint{1.799319in}{0.359145in}}%
\pgfpathquadraticcurveto{\pgfqpoint{1.804056in}{0.365648in}}{\pgfqpoint{1.804056in}{0.377067in}}%
\pgfpathquadraticcurveto{\pgfqpoint{1.804056in}{0.388487in}}{\pgfqpoint{1.799319in}{0.394989in}}%
\pgfpathquadraticcurveto{\pgfqpoint{1.794599in}{0.401492in}}{\pgfqpoint{1.786386in}{0.401492in}}%
\pgfpathquadraticcurveto{\pgfqpoint{1.778172in}{0.401492in}}{\pgfqpoint{1.773471in}{0.394989in}}%
\pgfpathquadraticcurveto{\pgfqpoint{1.768770in}{0.388487in}}{\pgfqpoint{1.768770in}{0.377067in}}%
\pgfpathlineto{\pgfqpoint{1.768770in}{0.377067in}}%
\pgfpathclose%
\pgfpathmoveto{\pgfqpoint{1.804056in}{0.355056in}}%
\pgfpathquadraticcurveto{\pgfqpoint{1.800796in}{0.349437in}}{\pgfqpoint{1.795806in}{0.346699in}}%
\pgfpathquadraticcurveto{\pgfqpoint{1.790835in}{0.343961in}}{\pgfqpoint{1.783846in}{0.343961in}}%
\pgfpathquadraticcurveto{\pgfqpoint{1.772426in}{0.343961in}}{\pgfqpoint{1.765240in}{0.353075in}}%
\pgfpathquadraticcurveto{\pgfqpoint{1.758071in}{0.362207in}}{\pgfqpoint{1.758071in}{0.377067in}}%
\pgfpathquadraticcurveto{\pgfqpoint{1.758071in}{0.391927in}}{\pgfqpoint{1.765240in}{0.401041in}}%
\pgfpathquadraticcurveto{\pgfqpoint{1.772426in}{0.410156in}}{\pgfqpoint{1.783846in}{0.410156in}}%
\pgfpathquadraticcurveto{\pgfqpoint{1.790835in}{0.410156in}}{\pgfqpoint{1.795806in}{0.407418in}}%
\pgfpathquadraticcurveto{\pgfqpoint{1.800796in}{0.404698in}}{\pgfqpoint{1.804056in}{0.399078in}}%
\pgfpathlineto{\pgfqpoint{1.804056in}{0.408643in}}%
\pgfpathlineto{\pgfqpoint{1.814413in}{0.408643in}}%
\pgfpathlineto{\pgfqpoint{1.814413in}{0.321626in}}%
\pgfpathlineto{\pgfqpoint{1.804056in}{0.321626in}}%
\pgfpathlineto{\pgfqpoint{1.804056in}{0.355056in}}%
\pgfpathlineto{\pgfqpoint{1.804056in}{0.355056in}}%
\pgfpathclose%
\pgfpathmoveto{\pgfqpoint{1.834694in}{0.370475in}}%
\pgfpathlineto{\pgfqpoint{1.834694in}{0.408643in}}%
\pgfpathlineto{\pgfqpoint{1.845051in}{0.408643in}}%
\pgfpathlineto{\pgfqpoint{1.845051in}{0.370871in}}%
\pgfpathquadraticcurveto{\pgfqpoint{1.845051in}{0.361919in}}{\pgfqpoint{1.848528in}{0.357434in}}%
\pgfpathquadraticcurveto{\pgfqpoint{1.852022in}{0.352967in}}{\pgfqpoint{1.859011in}{0.352967in}}%
\pgfpathquadraticcurveto{\pgfqpoint{1.867387in}{0.352967in}}{\pgfqpoint{1.872250in}{0.358317in}}%
\pgfpathquadraticcurveto{\pgfqpoint{1.877131in}{0.363666in}}{\pgfqpoint{1.877131in}{0.372906in}}%
\pgfpathlineto{\pgfqpoint{1.877131in}{0.408643in}}%
\pgfpathlineto{\pgfqpoint{1.887488in}{0.408643in}}%
\pgfpathlineto{\pgfqpoint{1.887488in}{0.345600in}}%
\pgfpathlineto{\pgfqpoint{1.877131in}{0.345600in}}%
\pgfpathlineto{\pgfqpoint{1.877131in}{0.355291in}}%
\pgfpathquadraticcurveto{\pgfqpoint{1.873367in}{0.349545in}}{\pgfqpoint{1.868377in}{0.346753in}}%
\pgfpathquadraticcurveto{\pgfqpoint{1.863406in}{0.343961in}}{\pgfqpoint{1.856813in}{0.343961in}}%
\pgfpathquadraticcurveto{\pgfqpoint{1.845952in}{0.343961in}}{\pgfqpoint{1.840314in}{0.350715in}}%
\pgfpathquadraticcurveto{\pgfqpoint{1.834694in}{0.357470in}}{\pgfqpoint{1.834694in}{0.370475in}}%
\pgfpathlineto{\pgfqpoint{1.834694in}{0.370475in}}%
\pgfpathclose%
\pgfpathmoveto{\pgfqpoint{1.860758in}{0.410156in}}%
\pgfpathlineto{\pgfqpoint{1.860758in}{0.410156in}}%
\pgfpathlineto{\pgfqpoint{1.860758in}{0.410156in}}%
\pgfpathclose%
\pgfpathmoveto{\pgfqpoint{1.937472in}{0.377283in}}%
\pgfpathquadraticcurveto{\pgfqpoint{1.924917in}{0.377283in}}{\pgfqpoint{1.920072in}{0.374419in}}%
\pgfpathquadraticcurveto{\pgfqpoint{1.915227in}{0.371556in}}{\pgfqpoint{1.915227in}{0.364621in}}%
\pgfpathquadraticcurveto{\pgfqpoint{1.915227in}{0.359109in}}{\pgfqpoint{1.918865in}{0.355867in}}%
\pgfpathquadraticcurveto{\pgfqpoint{1.922504in}{0.352643in}}{\pgfqpoint{1.928736in}{0.352643in}}%
\pgfpathquadraticcurveto{\pgfqpoint{1.937364in}{0.352643in}}{\pgfqpoint{1.942569in}{0.358749in}}%
\pgfpathquadraticcurveto{\pgfqpoint{1.947775in}{0.364855in}}{\pgfqpoint{1.947775in}{0.374978in}}%
\pgfpathlineto{\pgfqpoint{1.947775in}{0.377283in}}%
\pgfpathlineto{\pgfqpoint{1.937472in}{0.377283in}}%
\pgfpathlineto{\pgfqpoint{1.937472in}{0.377283in}}%
\pgfpathclose%
\pgfpathmoveto{\pgfqpoint{1.958132in}{0.381570in}}%
\pgfpathlineto{\pgfqpoint{1.958132in}{0.345600in}}%
\pgfpathlineto{\pgfqpoint{1.947775in}{0.345600in}}%
\pgfpathlineto{\pgfqpoint{1.947775in}{0.355164in}}%
\pgfpathquadraticcurveto{\pgfqpoint{1.944226in}{0.349437in}}{\pgfqpoint{1.938931in}{0.346699in}}%
\pgfpathquadraticcurveto{\pgfqpoint{1.933635in}{0.343961in}}{\pgfqpoint{1.925980in}{0.343961in}}%
\pgfpathquadraticcurveto{\pgfqpoint{1.916308in}{0.343961in}}{\pgfqpoint{1.910580in}{0.349401in}}%
\pgfpathquadraticcurveto{\pgfqpoint{1.904870in}{0.354840in}}{\pgfqpoint{1.904870in}{0.363954in}}%
\pgfpathquadraticcurveto{\pgfqpoint{1.904870in}{0.374582in}}{\pgfqpoint{1.911985in}{0.379985in}}%
\pgfpathquadraticcurveto{\pgfqpoint{1.919117in}{0.385389in}}{\pgfqpoint{1.933239in}{0.385389in}}%
\pgfpathlineto{\pgfqpoint{1.947775in}{0.385389in}}%
\pgfpathlineto{\pgfqpoint{1.947775in}{0.386416in}}%
\pgfpathquadraticcurveto{\pgfqpoint{1.947775in}{0.393566in}}{\pgfqpoint{1.943074in}{0.397475in}}%
\pgfpathquadraticcurveto{\pgfqpoint{1.938372in}{0.401384in}}{\pgfqpoint{1.929871in}{0.401384in}}%
\pgfpathquadraticcurveto{\pgfqpoint{1.924467in}{0.401384in}}{\pgfqpoint{1.919334in}{0.400087in}}%
\pgfpathquadraticcurveto{\pgfqpoint{1.914218in}{0.398790in}}{\pgfqpoint{1.909499in}{0.396196in}}%
\pgfpathlineto{\pgfqpoint{1.909499in}{0.405779in}}%
\pgfpathquadraticcurveto{\pgfqpoint{1.915173in}{0.407976in}}{\pgfqpoint{1.920522in}{0.409057in}}%
\pgfpathquadraticcurveto{\pgfqpoint{1.925872in}{0.410156in}}{\pgfqpoint{1.930933in}{0.410156in}}%
\pgfpathquadraticcurveto{\pgfqpoint{1.944623in}{0.410156in}}{\pgfqpoint{1.951377in}{0.403059in}}%
\pgfpathquadraticcurveto{\pgfqpoint{1.958132in}{0.395980in}}{\pgfqpoint{1.958132in}{0.381570in}}%
\pgfpathlineto{\pgfqpoint{1.958132in}{0.381570in}}%
\pgfpathclose%
\pgfpathmoveto{\pgfqpoint{1.979458in}{0.433193in}}%
\pgfpathlineto{\pgfqpoint{1.989815in}{0.433193in}}%
\pgfpathlineto{\pgfqpoint{1.989815in}{0.345600in}}%
\pgfpathlineto{\pgfqpoint{1.979458in}{0.345600in}}%
\pgfpathlineto{\pgfqpoint{1.979458in}{0.433193in}}%
\pgfpathlineto{\pgfqpoint{1.979458in}{0.433193in}}%
\pgfpathclose%
\pgfpathmoveto{\pgfqpoint{2.011484in}{0.408643in}}%
\pgfpathlineto{\pgfqpoint{2.021841in}{0.408643in}}%
\pgfpathlineto{\pgfqpoint{2.021841in}{0.345600in}}%
\pgfpathlineto{\pgfqpoint{2.011484in}{0.345600in}}%
\pgfpathlineto{\pgfqpoint{2.011484in}{0.408643in}}%
\pgfpathlineto{\pgfqpoint{2.011484in}{0.408643in}}%
\pgfpathclose%
\pgfpathmoveto{\pgfqpoint{2.011484in}{0.433193in}}%
\pgfpathlineto{\pgfqpoint{2.021841in}{0.433193in}}%
\pgfpathlineto{\pgfqpoint{2.021841in}{0.420062in}}%
\pgfpathlineto{\pgfqpoint{2.011484in}{0.420062in}}%
\pgfpathlineto{\pgfqpoint{2.011484in}{0.433193in}}%
\pgfpathlineto{\pgfqpoint{2.011484in}{0.433193in}}%
\pgfpathclose%
\pgfpathmoveto{\pgfqpoint{2.053758in}{0.426547in}}%
\pgfpathlineto{\pgfqpoint{2.053758in}{0.408643in}}%
\pgfpathlineto{\pgfqpoint{2.075085in}{0.408643in}}%
\pgfpathlineto{\pgfqpoint{2.075085in}{0.400591in}}%
\pgfpathlineto{\pgfqpoint{2.053758in}{0.400591in}}%
\pgfpathlineto{\pgfqpoint{2.053758in}{0.366368in}}%
\pgfpathquadraticcurveto{\pgfqpoint{2.053758in}{0.358659in}}{\pgfqpoint{2.055866in}{0.356461in}}%
\pgfpathquadraticcurveto{\pgfqpoint{2.057973in}{0.354264in}}{\pgfqpoint{2.064457in}{0.354264in}}%
\pgfpathlineto{\pgfqpoint{2.075085in}{0.354264in}}%
\pgfpathlineto{\pgfqpoint{2.075085in}{0.345600in}}%
\pgfpathlineto{\pgfqpoint{2.064457in}{0.345600in}}%
\pgfpathquadraticcurveto{\pgfqpoint{2.052461in}{0.345600in}}{\pgfqpoint{2.047904in}{0.350067in}}%
\pgfpathquadraticcurveto{\pgfqpoint{2.043347in}{0.354552in}}{\pgfqpoint{2.043347in}{0.366368in}}%
\pgfpathlineto{\pgfqpoint{2.043347in}{0.400591in}}%
\pgfpathlineto{\pgfqpoint{2.035746in}{0.400591in}}%
\pgfpathlineto{\pgfqpoint{2.035746in}{0.408643in}}%
\pgfpathlineto{\pgfqpoint{2.043347in}{0.408643in}}%
\pgfpathlineto{\pgfqpoint{2.043347in}{0.426547in}}%
\pgfpathlineto{\pgfqpoint{2.053758in}{0.426547in}}%
\pgfpathlineto{\pgfqpoint{2.053758in}{0.426547in}}%
\pgfpathclose%
\pgfpathmoveto{\pgfqpoint{2.114928in}{0.339746in}}%
\pgfpathquadraticcurveto{\pgfqpoint{2.110551in}{0.328488in}}{\pgfqpoint{2.106372in}{0.325066in}}%
\pgfpathquadraticcurveto{\pgfqpoint{2.102211in}{0.321626in}}{\pgfqpoint{2.095240in}{0.321626in}}%
\pgfpathlineto{\pgfqpoint{2.086955in}{0.321626in}}%
\pgfpathlineto{\pgfqpoint{2.086955in}{0.330290in}}%
\pgfpathlineto{\pgfqpoint{2.093043in}{0.330290in}}%
\pgfpathquadraticcurveto{\pgfqpoint{2.097312in}{0.330290in}}{\pgfqpoint{2.099671in}{0.332325in}}%
\pgfpathquadraticcurveto{\pgfqpoint{2.102049in}{0.334342in}}{\pgfqpoint{2.104913in}{0.341889in}}%
\pgfpathlineto{\pgfqpoint{2.106768in}{0.346609in}}%
\pgfpathlineto{\pgfqpoint{2.081281in}{0.408643in}}%
\pgfpathlineto{\pgfqpoint{2.092250in}{0.408643in}}%
\pgfpathlineto{\pgfqpoint{2.111956in}{0.359343in}}%
\pgfpathlineto{\pgfqpoint{2.131661in}{0.408643in}}%
\pgfpathlineto{\pgfqpoint{2.142630in}{0.408643in}}%
\pgfpathlineto{\pgfqpoint{2.114928in}{0.339746in}}%
\pgfpathlineto{\pgfqpoint{2.114928in}{0.339746in}}%
\pgfpathclose%
\pgfpathmoveto{\pgfqpoint{2.141964in}{0.359902in}}%
\pgfpathlineto{\pgfqpoint{2.153852in}{0.359902in}}%
\pgfpathlineto{\pgfqpoint{2.153852in}{0.345600in}}%
\pgfpathlineto{\pgfqpoint{2.141964in}{0.345600in}}%
\pgfpathlineto{\pgfqpoint{2.141964in}{0.359902in}}%
\pgfpathlineto{\pgfqpoint{2.141964in}{0.359902in}}%
\pgfpathclose%
\pgfpathmoveto{\pgfqpoint{2.264609in}{0.426889in}}%
\pgfpathlineto{\pgfqpoint{2.264609in}{0.415793in}}%
\pgfpathquadraticcurveto{\pgfqpoint{2.258142in}{0.418891in}}{\pgfqpoint{2.252396in}{0.420404in}}%
\pgfpathquadraticcurveto{\pgfqpoint{2.246650in}{0.421936in}}{\pgfqpoint{2.241301in}{0.421936in}}%
\pgfpathquadraticcurveto{\pgfqpoint{2.232025in}{0.421936in}}{\pgfqpoint{2.226981in}{0.418333in}}%
\pgfpathquadraticcurveto{\pgfqpoint{2.221938in}{0.414731in}}{\pgfqpoint{2.221938in}{0.408084in}}%
\pgfpathquadraticcurveto{\pgfqpoint{2.221938in}{0.402500in}}{\pgfqpoint{2.225288in}{0.399654in}}%
\pgfpathquadraticcurveto{\pgfqpoint{2.228638in}{0.396827in}}{\pgfqpoint{2.237987in}{0.395079in}}%
\pgfpathlineto{\pgfqpoint{2.244849in}{0.393674in}}%
\pgfpathquadraticcurveto{\pgfqpoint{2.257566in}{0.391243in}}{\pgfqpoint{2.263618in}{0.385137in}}%
\pgfpathquadraticcurveto{\pgfqpoint{2.269670in}{0.379031in}}{\pgfqpoint{2.269670in}{0.368800in}}%
\pgfpathquadraticcurveto{\pgfqpoint{2.269670in}{0.356569in}}{\pgfqpoint{2.261474in}{0.350265in}}%
\pgfpathquadraticcurveto{\pgfqpoint{2.253297in}{0.343961in}}{\pgfqpoint{2.237482in}{0.343961in}}%
\pgfpathquadraticcurveto{\pgfqpoint{2.231520in}{0.343961in}}{\pgfqpoint{2.224784in}{0.345312in}}%
\pgfpathquadraticcurveto{\pgfqpoint{2.218065in}{0.346663in}}{\pgfqpoint{2.210860in}{0.349311in}}%
\pgfpathlineto{\pgfqpoint{2.210860in}{0.361018in}}%
\pgfpathquadraticcurveto{\pgfqpoint{2.217777in}{0.357146in}}{\pgfqpoint{2.224423in}{0.355164in}}%
\pgfpathquadraticcurveto{\pgfqpoint{2.231070in}{0.353201in}}{\pgfqpoint{2.237482in}{0.353201in}}%
\pgfpathquadraticcurveto{\pgfqpoint{2.247209in}{0.353201in}}{\pgfqpoint{2.252504in}{0.357020in}}%
\pgfpathquadraticcurveto{\pgfqpoint{2.257800in}{0.360856in}}{\pgfqpoint{2.257800in}{0.367953in}}%
\pgfpathquadraticcurveto{\pgfqpoint{2.257800in}{0.374131in}}{\pgfqpoint{2.253999in}{0.377626in}}%
\pgfpathquadraticcurveto{\pgfqpoint{2.250199in}{0.381120in}}{\pgfqpoint{2.241535in}{0.382867in}}%
\pgfpathlineto{\pgfqpoint{2.234600in}{0.384218in}}%
\pgfpathquadraticcurveto{\pgfqpoint{2.221884in}{0.386740in}}{\pgfqpoint{2.216192in}{0.392143in}}%
\pgfpathquadraticcurveto{\pgfqpoint{2.210518in}{0.397547in}}{\pgfqpoint{2.210518in}{0.407184in}}%
\pgfpathquadraticcurveto{\pgfqpoint{2.210518in}{0.418333in}}{\pgfqpoint{2.218371in}{0.424745in}}%
\pgfpathquadraticcurveto{\pgfqpoint{2.226225in}{0.431158in}}{\pgfqpoint{2.240004in}{0.431158in}}%
\pgfpathquadraticcurveto{\pgfqpoint{2.245930in}{0.431158in}}{\pgfqpoint{2.252054in}{0.430077in}}%
\pgfpathquadraticcurveto{\pgfqpoint{2.258196in}{0.429014in}}{\pgfqpoint{2.264609in}{0.426889in}}%
\pgfpathlineto{\pgfqpoint{2.264609in}{0.426889in}}%
\pgfpathclose%
\pgfpathmoveto{\pgfqpoint{2.297211in}{0.426547in}}%
\pgfpathlineto{\pgfqpoint{2.297211in}{0.408643in}}%
\pgfpathlineto{\pgfqpoint{2.318537in}{0.408643in}}%
\pgfpathlineto{\pgfqpoint{2.318537in}{0.400591in}}%
\pgfpathlineto{\pgfqpoint{2.297211in}{0.400591in}}%
\pgfpathlineto{\pgfqpoint{2.297211in}{0.366368in}}%
\pgfpathquadraticcurveto{\pgfqpoint{2.297211in}{0.358659in}}{\pgfqpoint{2.299318in}{0.356461in}}%
\pgfpathquadraticcurveto{\pgfqpoint{2.301425in}{0.354264in}}{\pgfqpoint{2.307910in}{0.354264in}}%
\pgfpathlineto{\pgfqpoint{2.318537in}{0.354264in}}%
\pgfpathlineto{\pgfqpoint{2.318537in}{0.345600in}}%
\pgfpathlineto{\pgfqpoint{2.307910in}{0.345600in}}%
\pgfpathquadraticcurveto{\pgfqpoint{2.295914in}{0.345600in}}{\pgfqpoint{2.291357in}{0.350067in}}%
\pgfpathquadraticcurveto{\pgfqpoint{2.286799in}{0.354552in}}{\pgfqpoint{2.286799in}{0.366368in}}%
\pgfpathlineto{\pgfqpoint{2.286799in}{0.400591in}}%
\pgfpathlineto{\pgfqpoint{2.279198in}{0.400591in}}%
\pgfpathlineto{\pgfqpoint{2.279198in}{0.408643in}}%
\pgfpathlineto{\pgfqpoint{2.286799in}{0.408643in}}%
\pgfpathlineto{\pgfqpoint{2.286799in}{0.426547in}}%
\pgfpathlineto{\pgfqpoint{2.297211in}{0.426547in}}%
\pgfpathlineto{\pgfqpoint{2.297211in}{0.426547in}}%
\pgfpathclose%
\pgfpathmoveto{\pgfqpoint{2.331091in}{0.370475in}}%
\pgfpathlineto{\pgfqpoint{2.331091in}{0.408643in}}%
\pgfpathlineto{\pgfqpoint{2.341448in}{0.408643in}}%
\pgfpathlineto{\pgfqpoint{2.341448in}{0.370871in}}%
\pgfpathquadraticcurveto{\pgfqpoint{2.341448in}{0.361919in}}{\pgfqpoint{2.344925in}{0.357434in}}%
\pgfpathquadraticcurveto{\pgfqpoint{2.348419in}{0.352967in}}{\pgfqpoint{2.355408in}{0.352967in}}%
\pgfpathquadraticcurveto{\pgfqpoint{2.363783in}{0.352967in}}{\pgfqpoint{2.368647in}{0.358317in}}%
\pgfpathquadraticcurveto{\pgfqpoint{2.373528in}{0.363666in}}{\pgfqpoint{2.373528in}{0.372906in}}%
\pgfpathlineto{\pgfqpoint{2.373528in}{0.408643in}}%
\pgfpathlineto{\pgfqpoint{2.383885in}{0.408643in}}%
\pgfpathlineto{\pgfqpoint{2.383885in}{0.345600in}}%
\pgfpathlineto{\pgfqpoint{2.373528in}{0.345600in}}%
\pgfpathlineto{\pgfqpoint{2.373528in}{0.355291in}}%
\pgfpathquadraticcurveto{\pgfqpoint{2.369763in}{0.349545in}}{\pgfqpoint{2.364774in}{0.346753in}}%
\pgfpathquadraticcurveto{\pgfqpoint{2.359803in}{0.343961in}}{\pgfqpoint{2.353210in}{0.343961in}}%
\pgfpathquadraticcurveto{\pgfqpoint{2.342349in}{0.343961in}}{\pgfqpoint{2.336711in}{0.350715in}}%
\pgfpathquadraticcurveto{\pgfqpoint{2.331091in}{0.357470in}}{\pgfqpoint{2.331091in}{0.370475in}}%
\pgfpathlineto{\pgfqpoint{2.331091in}{0.370475in}}%
\pgfpathclose%
\pgfpathmoveto{\pgfqpoint{2.357155in}{0.410156in}}%
\pgfpathlineto{\pgfqpoint{2.357155in}{0.410156in}}%
\pgfpathlineto{\pgfqpoint{2.357155in}{0.410156in}}%
\pgfpathclose%
\pgfpathmoveto{\pgfqpoint{2.446693in}{0.399078in}}%
\pgfpathlineto{\pgfqpoint{2.446693in}{0.433193in}}%
\pgfpathlineto{\pgfqpoint{2.457050in}{0.433193in}}%
\pgfpathlineto{\pgfqpoint{2.457050in}{0.345600in}}%
\pgfpathlineto{\pgfqpoint{2.446693in}{0.345600in}}%
\pgfpathlineto{\pgfqpoint{2.446693in}{0.355056in}}%
\pgfpathquadraticcurveto{\pgfqpoint{2.443433in}{0.349437in}}{\pgfqpoint{2.438444in}{0.346699in}}%
\pgfpathquadraticcurveto{\pgfqpoint{2.433472in}{0.343961in}}{\pgfqpoint{2.426484in}{0.343961in}}%
\pgfpathquadraticcurveto{\pgfqpoint{2.415064in}{0.343961in}}{\pgfqpoint{2.407877in}{0.353075in}}%
\pgfpathquadraticcurveto{\pgfqpoint{2.400708in}{0.362207in}}{\pgfqpoint{2.400708in}{0.377067in}}%
\pgfpathquadraticcurveto{\pgfqpoint{2.400708in}{0.391927in}}{\pgfqpoint{2.407877in}{0.401041in}}%
\pgfpathquadraticcurveto{\pgfqpoint{2.415064in}{0.410156in}}{\pgfqpoint{2.426484in}{0.410156in}}%
\pgfpathquadraticcurveto{\pgfqpoint{2.433472in}{0.410156in}}{\pgfqpoint{2.438444in}{0.407418in}}%
\pgfpathquadraticcurveto{\pgfqpoint{2.443433in}{0.404698in}}{\pgfqpoint{2.446693in}{0.399078in}}%
\pgfpathlineto{\pgfqpoint{2.446693in}{0.399078in}}%
\pgfpathclose%
\pgfpathmoveto{\pgfqpoint{2.411408in}{0.377067in}}%
\pgfpathquadraticcurveto{\pgfqpoint{2.411408in}{0.365648in}}{\pgfqpoint{2.416109in}{0.359145in}}%
\pgfpathquadraticcurveto{\pgfqpoint{2.420810in}{0.352643in}}{\pgfqpoint{2.429023in}{0.352643in}}%
\pgfpathquadraticcurveto{\pgfqpoint{2.437237in}{0.352643in}}{\pgfqpoint{2.441956in}{0.359145in}}%
\pgfpathquadraticcurveto{\pgfqpoint{2.446693in}{0.365648in}}{\pgfqpoint{2.446693in}{0.377067in}}%
\pgfpathquadraticcurveto{\pgfqpoint{2.446693in}{0.388487in}}{\pgfqpoint{2.441956in}{0.394989in}}%
\pgfpathquadraticcurveto{\pgfqpoint{2.437237in}{0.401492in}}{\pgfqpoint{2.429023in}{0.401492in}}%
\pgfpathquadraticcurveto{\pgfqpoint{2.420810in}{0.401492in}}{\pgfqpoint{2.416109in}{0.394989in}}%
\pgfpathquadraticcurveto{\pgfqpoint{2.411408in}{0.388487in}}{\pgfqpoint{2.411408in}{0.377067in}}%
\pgfpathlineto{\pgfqpoint{2.411408in}{0.377067in}}%
\pgfpathclose%
\pgfpathmoveto{\pgfqpoint{2.504620in}{0.339746in}}%
\pgfpathquadraticcurveto{\pgfqpoint{2.500243in}{0.328488in}}{\pgfqpoint{2.496065in}{0.325066in}}%
\pgfpathquadraticcurveto{\pgfqpoint{2.491904in}{0.321626in}}{\pgfqpoint{2.484933in}{0.321626in}}%
\pgfpathlineto{\pgfqpoint{2.476648in}{0.321626in}}%
\pgfpathlineto{\pgfqpoint{2.476648in}{0.330290in}}%
\pgfpathlineto{\pgfqpoint{2.482736in}{0.330290in}}%
\pgfpathquadraticcurveto{\pgfqpoint{2.487005in}{0.330290in}}{\pgfqpoint{2.489364in}{0.332325in}}%
\pgfpathquadraticcurveto{\pgfqpoint{2.491742in}{0.334342in}}{\pgfqpoint{2.494606in}{0.341889in}}%
\pgfpathlineto{\pgfqpoint{2.496461in}{0.346609in}}%
\pgfpathlineto{\pgfqpoint{2.470974in}{0.408643in}}%
\pgfpathlineto{\pgfqpoint{2.481943in}{0.408643in}}%
\pgfpathlineto{\pgfqpoint{2.501648in}{0.359343in}}%
\pgfpathlineto{\pgfqpoint{2.521354in}{0.408643in}}%
\pgfpathlineto{\pgfqpoint{2.532323in}{0.408643in}}%
\pgfpathlineto{\pgfqpoint{2.504620in}{0.339746in}}%
\pgfpathlineto{\pgfqpoint{2.504620in}{0.339746in}}%
\pgfpathclose%
\pgfpathmoveto{\pgfqpoint{2.609037in}{0.385515in}}%
\pgfpathquadraticcurveto{\pgfqpoint{2.600931in}{0.385515in}}{\pgfqpoint{2.596284in}{0.381174in}}%
\pgfpathquadraticcurveto{\pgfqpoint{2.591655in}{0.376833in}}{\pgfqpoint{2.591655in}{0.369250in}}%
\pgfpathquadraticcurveto{\pgfqpoint{2.591655in}{0.361649in}}{\pgfqpoint{2.596284in}{0.357308in}}%
\pgfpathquadraticcurveto{\pgfqpoint{2.600931in}{0.352967in}}{\pgfqpoint{2.609037in}{0.352967in}}%
\pgfpathquadraticcurveto{\pgfqpoint{2.617142in}{0.352967in}}{\pgfqpoint{2.621808in}{0.357326in}}%
\pgfpathquadraticcurveto{\pgfqpoint{2.626491in}{0.361703in}}{\pgfqpoint{2.626491in}{0.369250in}}%
\pgfpathquadraticcurveto{\pgfqpoint{2.626491in}{0.376833in}}{\pgfqpoint{2.621844in}{0.381174in}}%
\pgfpathquadraticcurveto{\pgfqpoint{2.617214in}{0.385515in}}{\pgfqpoint{2.609037in}{0.385515in}}%
\pgfpathlineto{\pgfqpoint{2.609037in}{0.385515in}}%
\pgfpathclose%
\pgfpathmoveto{\pgfqpoint{2.597671in}{0.390342in}}%
\pgfpathquadraticcurveto{\pgfqpoint{2.590358in}{0.392143in}}{\pgfqpoint{2.586270in}{0.397151in}}%
\pgfpathquadraticcurveto{\pgfqpoint{2.582199in}{0.402176in}}{\pgfqpoint{2.582199in}{0.409381in}}%
\pgfpathquadraticcurveto{\pgfqpoint{2.582199in}{0.419450in}}{\pgfqpoint{2.589368in}{0.425304in}}%
\pgfpathquadraticcurveto{\pgfqpoint{2.596554in}{0.431158in}}{\pgfqpoint{2.609037in}{0.431158in}}%
\pgfpathquadraticcurveto{\pgfqpoint{2.621591in}{0.431158in}}{\pgfqpoint{2.628742in}{0.425304in}}%
\pgfpathquadraticcurveto{\pgfqpoint{2.635893in}{0.419450in}}{\pgfqpoint{2.635893in}{0.409381in}}%
\pgfpathquadraticcurveto{\pgfqpoint{2.635893in}{0.402176in}}{\pgfqpoint{2.631804in}{0.397151in}}%
\pgfpathquadraticcurveto{\pgfqpoint{2.627734in}{0.392143in}}{\pgfqpoint{2.620475in}{0.390342in}}%
\pgfpathquadraticcurveto{\pgfqpoint{2.628688in}{0.388433in}}{\pgfqpoint{2.633263in}{0.382849in}}%
\pgfpathquadraticcurveto{\pgfqpoint{2.637856in}{0.377283in}}{\pgfqpoint{2.637856in}{0.369250in}}%
\pgfpathquadraticcurveto{\pgfqpoint{2.637856in}{0.357020in}}{\pgfqpoint{2.630399in}{0.350481in}}%
\pgfpathquadraticcurveto{\pgfqpoint{2.622942in}{0.343961in}}{\pgfqpoint{2.609037in}{0.343961in}}%
\pgfpathquadraticcurveto{\pgfqpoint{2.595150in}{0.343961in}}{\pgfqpoint{2.587674in}{0.350481in}}%
\pgfpathquadraticcurveto{\pgfqpoint{2.580217in}{0.357020in}}{\pgfqpoint{2.580217in}{0.369250in}}%
\pgfpathquadraticcurveto{\pgfqpoint{2.580217in}{0.377283in}}{\pgfqpoint{2.584829in}{0.382849in}}%
\pgfpathquadraticcurveto{\pgfqpoint{2.589458in}{0.388433in}}{\pgfqpoint{2.597671in}{0.390342in}}%
\pgfpathlineto{\pgfqpoint{2.597671in}{0.390342in}}%
\pgfpathclose%
\pgfpathmoveto{\pgfqpoint{2.593510in}{0.408300in}}%
\pgfpathquadraticcurveto{\pgfqpoint{2.593510in}{0.401780in}}{\pgfqpoint{2.597581in}{0.398123in}}%
\pgfpathquadraticcurveto{\pgfqpoint{2.601670in}{0.394467in}}{\pgfqpoint{2.609037in}{0.394467in}}%
\pgfpathquadraticcurveto{\pgfqpoint{2.616368in}{0.394467in}}{\pgfqpoint{2.620493in}{0.398123in}}%
\pgfpathquadraticcurveto{\pgfqpoint{2.624635in}{0.401780in}}{\pgfqpoint{2.624635in}{0.408300in}}%
\pgfpathquadraticcurveto{\pgfqpoint{2.624635in}{0.414839in}}{\pgfqpoint{2.620493in}{0.418495in}}%
\pgfpathquadraticcurveto{\pgfqpoint{2.616368in}{0.422152in}}{\pgfqpoint{2.609037in}{0.422152in}}%
\pgfpathquadraticcurveto{\pgfqpoint{2.601670in}{0.422152in}}{\pgfqpoint{2.597581in}{0.418495in}}%
\pgfpathquadraticcurveto{\pgfqpoint{2.593510in}{0.414839in}}{\pgfqpoint{2.593510in}{0.408300in}}%
\pgfpathlineto{\pgfqpoint{2.593510in}{0.408300in}}%
\pgfpathclose%
\pgfpathmoveto{\pgfqpoint{2.692181in}{0.370475in}}%
\pgfpathlineto{\pgfqpoint{2.692181in}{0.408643in}}%
\pgfpathlineto{\pgfqpoint{2.702538in}{0.408643in}}%
\pgfpathlineto{\pgfqpoint{2.702538in}{0.370871in}}%
\pgfpathquadraticcurveto{\pgfqpoint{2.702538in}{0.361919in}}{\pgfqpoint{2.706014in}{0.357434in}}%
\pgfpathquadraticcurveto{\pgfqpoint{2.709509in}{0.352967in}}{\pgfqpoint{2.716497in}{0.352967in}}%
\pgfpathquadraticcurveto{\pgfqpoint{2.724873in}{0.352967in}}{\pgfqpoint{2.729736in}{0.358317in}}%
\pgfpathquadraticcurveto{\pgfqpoint{2.734618in}{0.363666in}}{\pgfqpoint{2.734618in}{0.372906in}}%
\pgfpathlineto{\pgfqpoint{2.734618in}{0.408643in}}%
\pgfpathlineto{\pgfqpoint{2.744975in}{0.408643in}}%
\pgfpathlineto{\pgfqpoint{2.744975in}{0.345600in}}%
\pgfpathlineto{\pgfqpoint{2.734618in}{0.345600in}}%
\pgfpathlineto{\pgfqpoint{2.734618in}{0.355291in}}%
\pgfpathquadraticcurveto{\pgfqpoint{2.730853in}{0.349545in}}{\pgfqpoint{2.725864in}{0.346753in}}%
\pgfpathquadraticcurveto{\pgfqpoint{2.720892in}{0.343961in}}{\pgfqpoint{2.714300in}{0.343961in}}%
\pgfpathquadraticcurveto{\pgfqpoint{2.703439in}{0.343961in}}{\pgfqpoint{2.697801in}{0.350715in}}%
\pgfpathquadraticcurveto{\pgfqpoint{2.692181in}{0.357470in}}{\pgfqpoint{2.692181in}{0.370475in}}%
\pgfpathlineto{\pgfqpoint{2.692181in}{0.370475in}}%
\pgfpathclose%
\pgfpathmoveto{\pgfqpoint{2.718245in}{0.410156in}}%
\pgfpathlineto{\pgfqpoint{2.718245in}{0.410156in}}%
\pgfpathlineto{\pgfqpoint{2.718245in}{0.410156in}}%
\pgfpathclose%
\pgfpathmoveto{\pgfqpoint{2.806486in}{0.406787in}}%
\pgfpathlineto{\pgfqpoint{2.806486in}{0.396989in}}%
\pgfpathquadraticcurveto{\pgfqpoint{2.802109in}{0.399240in}}{\pgfqpoint{2.797372in}{0.400357in}}%
\pgfpathquadraticcurveto{\pgfqpoint{2.792653in}{0.401492in}}{\pgfqpoint{2.787573in}{0.401492in}}%
\pgfpathquadraticcurveto{\pgfqpoint{2.779864in}{0.401492in}}{\pgfqpoint{2.776010in}{0.399132in}}%
\pgfpathquadraticcurveto{\pgfqpoint{2.772155in}{0.396773in}}{\pgfqpoint{2.772155in}{0.392035in}}%
\pgfpathquadraticcurveto{\pgfqpoint{2.772155in}{0.388433in}}{\pgfqpoint{2.774911in}{0.386380in}}%
\pgfpathquadraticcurveto{\pgfqpoint{2.777667in}{0.384326in}}{\pgfqpoint{2.786006in}{0.382471in}}%
\pgfpathlineto{\pgfqpoint{2.789555in}{0.381678in}}%
\pgfpathquadraticcurveto{\pgfqpoint{2.800578in}{0.379319in}}{\pgfqpoint{2.805225in}{0.375014in}}%
\pgfpathquadraticcurveto{\pgfqpoint{2.809872in}{0.370709in}}{\pgfqpoint{2.809872in}{0.363000in}}%
\pgfpathquadraticcurveto{\pgfqpoint{2.809872in}{0.354210in}}{\pgfqpoint{2.802920in}{0.349076in}}%
\pgfpathquadraticcurveto{\pgfqpoint{2.795967in}{0.343961in}}{\pgfqpoint{2.783809in}{0.343961in}}%
\pgfpathquadraticcurveto{\pgfqpoint{2.778747in}{0.343961in}}{\pgfqpoint{2.773254in}{0.344952in}}%
\pgfpathquadraticcurveto{\pgfqpoint{2.767760in}{0.345942in}}{\pgfqpoint{2.761690in}{0.347906in}}%
\pgfpathlineto{\pgfqpoint{2.761690in}{0.358605in}}%
\pgfpathquadraticcurveto{\pgfqpoint{2.767436in}{0.355615in}}{\pgfqpoint{2.773002in}{0.354120in}}%
\pgfpathquadraticcurveto{\pgfqpoint{2.778567in}{0.352643in}}{\pgfqpoint{2.784043in}{0.352643in}}%
\pgfpathquadraticcurveto{\pgfqpoint{2.791356in}{0.352643in}}{\pgfqpoint{2.795283in}{0.355146in}}%
\pgfpathquadraticcurveto{\pgfqpoint{2.799227in}{0.357650in}}{\pgfqpoint{2.799227in}{0.362207in}}%
\pgfpathquadraticcurveto{\pgfqpoint{2.799227in}{0.366422in}}{\pgfqpoint{2.796381in}{0.368674in}}%
\pgfpathquadraticcurveto{\pgfqpoint{2.793553in}{0.370925in}}{\pgfqpoint{2.783917in}{0.373014in}}%
\pgfpathlineto{\pgfqpoint{2.780314in}{0.373861in}}%
\pgfpathquadraticcurveto{\pgfqpoint{2.770696in}{0.375878in}}{\pgfqpoint{2.766409in}{0.380075in}}%
\pgfpathquadraticcurveto{\pgfqpoint{2.762140in}{0.384272in}}{\pgfqpoint{2.762140in}{0.391585in}}%
\pgfpathquadraticcurveto{\pgfqpoint{2.762140in}{0.400483in}}{\pgfqpoint{2.768444in}{0.405310in}}%
\pgfpathquadraticcurveto{\pgfqpoint{2.774749in}{0.410156in}}{\pgfqpoint{2.786349in}{0.410156in}}%
\pgfpathquadraticcurveto{\pgfqpoint{2.792076in}{0.410156in}}{\pgfqpoint{2.797138in}{0.409309in}}%
\pgfpathquadraticcurveto{\pgfqpoint{2.802217in}{0.408480in}}{\pgfqpoint{2.806486in}{0.406787in}}%
\pgfpathlineto{\pgfqpoint{2.806486in}{0.406787in}}%
\pgfpathclose%
\pgfpathmoveto{\pgfqpoint{2.880282in}{0.379715in}}%
\pgfpathlineto{\pgfqpoint{2.880282in}{0.374654in}}%
\pgfpathlineto{\pgfqpoint{2.832658in}{0.374654in}}%
\pgfpathquadraticcurveto{\pgfqpoint{2.833342in}{0.363954in}}{\pgfqpoint{2.839106in}{0.358353in}}%
\pgfpathquadraticcurveto{\pgfqpoint{2.844870in}{0.352751in}}{\pgfqpoint{2.855173in}{0.352751in}}%
\pgfpathquadraticcurveto{\pgfqpoint{2.861135in}{0.352751in}}{\pgfqpoint{2.866737in}{0.354210in}}%
\pgfpathquadraticcurveto{\pgfqpoint{2.872339in}{0.355669in}}{\pgfqpoint{2.877868in}{0.358605in}}%
\pgfpathlineto{\pgfqpoint{2.877868in}{0.348806in}}%
\pgfpathquadraticcurveto{\pgfqpoint{2.872285in}{0.346447in}}{\pgfqpoint{2.866431in}{0.345204in}}%
\pgfpathquadraticcurveto{\pgfqpoint{2.860577in}{0.343961in}}{\pgfqpoint{2.854561in}{0.343961in}}%
\pgfpathquadraticcurveto{\pgfqpoint{2.839466in}{0.343961in}}{\pgfqpoint{2.830658in}{0.352733in}}%
\pgfpathquadraticcurveto{\pgfqpoint{2.821850in}{0.361523in}}{\pgfqpoint{2.821850in}{0.376509in}}%
\pgfpathquadraticcurveto{\pgfqpoint{2.821850in}{0.391981in}}{\pgfqpoint{2.830208in}{0.401059in}}%
\pgfpathquadraticcurveto{\pgfqpoint{2.838566in}{0.410156in}}{\pgfqpoint{2.852759in}{0.410156in}}%
\pgfpathquadraticcurveto{\pgfqpoint{2.865476in}{0.410156in}}{\pgfqpoint{2.872879in}{0.401960in}}%
\pgfpathquadraticcurveto{\pgfqpoint{2.880282in}{0.393783in}}{\pgfqpoint{2.880282in}{0.379715in}}%
\pgfpathlineto{\pgfqpoint{2.880282in}{0.379715in}}%
\pgfpathclose%
\pgfpathmoveto{\pgfqpoint{2.869925in}{0.382759in}}%
\pgfpathquadraticcurveto{\pgfqpoint{2.869817in}{0.391243in}}{\pgfqpoint{2.865170in}{0.396304in}}%
\pgfpathquadraticcurveto{\pgfqpoint{2.860523in}{0.401384in}}{\pgfqpoint{2.852867in}{0.401384in}}%
\pgfpathquadraticcurveto{\pgfqpoint{2.844204in}{0.401384in}}{\pgfqpoint{2.838998in}{0.396484in}}%
\pgfpathquadraticcurveto{\pgfqpoint{2.833793in}{0.391585in}}{\pgfqpoint{2.833000in}{0.382687in}}%
\pgfpathlineto{\pgfqpoint{2.869925in}{0.382759in}}%
\pgfpathlineto{\pgfqpoint{2.869925in}{0.382759in}}%
\pgfpathclose%
\pgfpathmoveto{\pgfqpoint{2.937470in}{0.406787in}}%
\pgfpathlineto{\pgfqpoint{2.937470in}{0.396989in}}%
\pgfpathquadraticcurveto{\pgfqpoint{2.933094in}{0.399240in}}{\pgfqpoint{2.928356in}{0.400357in}}%
\pgfpathquadraticcurveto{\pgfqpoint{2.923637in}{0.401492in}}{\pgfqpoint{2.918558in}{0.401492in}}%
\pgfpathquadraticcurveto{\pgfqpoint{2.910849in}{0.401492in}}{\pgfqpoint{2.906994in}{0.399132in}}%
\pgfpathquadraticcurveto{\pgfqpoint{2.903139in}{0.396773in}}{\pgfqpoint{2.903139in}{0.392035in}}%
\pgfpathquadraticcurveto{\pgfqpoint{2.903139in}{0.388433in}}{\pgfqpoint{2.905895in}{0.386380in}}%
\pgfpathquadraticcurveto{\pgfqpoint{2.908651in}{0.384326in}}{\pgfqpoint{2.916991in}{0.382471in}}%
\pgfpathlineto{\pgfqpoint{2.920539in}{0.381678in}}%
\pgfpathquadraticcurveto{\pgfqpoint{2.931563in}{0.379319in}}{\pgfqpoint{2.936210in}{0.375014in}}%
\pgfpathquadraticcurveto{\pgfqpoint{2.940857in}{0.370709in}}{\pgfqpoint{2.940857in}{0.363000in}}%
\pgfpathquadraticcurveto{\pgfqpoint{2.940857in}{0.354210in}}{\pgfqpoint{2.933904in}{0.349076in}}%
\pgfpathquadraticcurveto{\pgfqpoint{2.926951in}{0.343961in}}{\pgfqpoint{2.914793in}{0.343961in}}%
\pgfpathquadraticcurveto{\pgfqpoint{2.909732in}{0.343961in}}{\pgfqpoint{2.904238in}{0.344952in}}%
\pgfpathquadraticcurveto{\pgfqpoint{2.898744in}{0.345942in}}{\pgfqpoint{2.892674in}{0.347906in}}%
\pgfpathlineto{\pgfqpoint{2.892674in}{0.358605in}}%
\pgfpathquadraticcurveto{\pgfqpoint{2.898420in}{0.355615in}}{\pgfqpoint{2.903986in}{0.354120in}}%
\pgfpathquadraticcurveto{\pgfqpoint{2.909552in}{0.352643in}}{\pgfqpoint{2.915027in}{0.352643in}}%
\pgfpathquadraticcurveto{\pgfqpoint{2.922340in}{0.352643in}}{\pgfqpoint{2.926267in}{0.355146in}}%
\pgfpathquadraticcurveto{\pgfqpoint{2.930212in}{0.357650in}}{\pgfqpoint{2.930212in}{0.362207in}}%
\pgfpathquadraticcurveto{\pgfqpoint{2.930212in}{0.366422in}}{\pgfqpoint{2.927366in}{0.368674in}}%
\pgfpathquadraticcurveto{\pgfqpoint{2.924538in}{0.370925in}}{\pgfqpoint{2.914901in}{0.373014in}}%
\pgfpathlineto{\pgfqpoint{2.911299in}{0.373861in}}%
\pgfpathquadraticcurveto{\pgfqpoint{2.901680in}{0.375878in}}{\pgfqpoint{2.897393in}{0.380075in}}%
\pgfpathquadraticcurveto{\pgfqpoint{2.893125in}{0.384272in}}{\pgfqpoint{2.893125in}{0.391585in}}%
\pgfpathquadraticcurveto{\pgfqpoint{2.893125in}{0.400483in}}{\pgfqpoint{2.899429in}{0.405310in}}%
\pgfpathquadraticcurveto{\pgfqpoint{2.905733in}{0.410156in}}{\pgfqpoint{2.917333in}{0.410156in}}%
\pgfpathquadraticcurveto{\pgfqpoint{2.923061in}{0.410156in}}{\pgfqpoint{2.928122in}{0.409309in}}%
\pgfpathquadraticcurveto{\pgfqpoint{2.933202in}{0.408480in}}{\pgfqpoint{2.937470in}{0.406787in}}%
\pgfpathlineto{\pgfqpoint{2.937470in}{0.406787in}}%
\pgfpathclose%
\pgfpathmoveto{\pgfqpoint{2.992912in}{0.370475in}}%
\pgfpathlineto{\pgfqpoint{2.992912in}{0.408643in}}%
\pgfpathlineto{\pgfqpoint{3.003269in}{0.408643in}}%
\pgfpathlineto{\pgfqpoint{3.003269in}{0.370871in}}%
\pgfpathquadraticcurveto{\pgfqpoint{3.003269in}{0.361919in}}{\pgfqpoint{3.006745in}{0.357434in}}%
\pgfpathquadraticcurveto{\pgfqpoint{3.010240in}{0.352967in}}{\pgfqpoint{3.017228in}{0.352967in}}%
\pgfpathquadraticcurveto{\pgfqpoint{3.025604in}{0.352967in}}{\pgfqpoint{3.030467in}{0.358317in}}%
\pgfpathquadraticcurveto{\pgfqpoint{3.035349in}{0.363666in}}{\pgfqpoint{3.035349in}{0.372906in}}%
\pgfpathlineto{\pgfqpoint{3.035349in}{0.408643in}}%
\pgfpathlineto{\pgfqpoint{3.045706in}{0.408643in}}%
\pgfpathlineto{\pgfqpoint{3.045706in}{0.345600in}}%
\pgfpathlineto{\pgfqpoint{3.035349in}{0.345600in}}%
\pgfpathlineto{\pgfqpoint{3.035349in}{0.355291in}}%
\pgfpathquadraticcurveto{\pgfqpoint{3.031584in}{0.349545in}}{\pgfqpoint{3.026595in}{0.346753in}}%
\pgfpathquadraticcurveto{\pgfqpoint{3.021623in}{0.343961in}}{\pgfqpoint{3.015031in}{0.343961in}}%
\pgfpathquadraticcurveto{\pgfqpoint{3.004169in}{0.343961in}}{\pgfqpoint{2.998532in}{0.350715in}}%
\pgfpathquadraticcurveto{\pgfqpoint{2.992912in}{0.357470in}}{\pgfqpoint{2.992912in}{0.370475in}}%
\pgfpathlineto{\pgfqpoint{2.992912in}{0.370475in}}%
\pgfpathclose%
\pgfpathmoveto{\pgfqpoint{3.018975in}{0.410156in}}%
\pgfpathlineto{\pgfqpoint{3.018975in}{0.410156in}}%
\pgfpathlineto{\pgfqpoint{3.018975in}{0.410156in}}%
\pgfpathclose%
\pgfpathmoveto{\pgfqpoint{3.119447in}{0.383660in}}%
\pgfpathlineto{\pgfqpoint{3.119447in}{0.345600in}}%
\pgfpathlineto{\pgfqpoint{3.109090in}{0.345600in}}%
\pgfpathlineto{\pgfqpoint{3.109090in}{0.383317in}}%
\pgfpathquadraticcurveto{\pgfqpoint{3.109090in}{0.392269in}}{\pgfqpoint{3.105596in}{0.396700in}}%
\pgfpathquadraticcurveto{\pgfqpoint{3.102102in}{0.401149in}}{\pgfqpoint{3.095131in}{0.401149in}}%
\pgfpathquadraticcurveto{\pgfqpoint{3.086737in}{0.401149in}}{\pgfqpoint{3.081892in}{0.395800in}}%
\pgfpathquadraticcurveto{\pgfqpoint{3.077047in}{0.390468in}}{\pgfqpoint{3.077047in}{0.381228in}}%
\pgfpathlineto{\pgfqpoint{3.077047in}{0.345600in}}%
\pgfpathlineto{\pgfqpoint{3.066636in}{0.345600in}}%
\pgfpathlineto{\pgfqpoint{3.066636in}{0.408643in}}%
\pgfpathlineto{\pgfqpoint{3.077047in}{0.408643in}}%
\pgfpathlineto{\pgfqpoint{3.077047in}{0.398844in}}%
\pgfpathquadraticcurveto{\pgfqpoint{3.080775in}{0.404536in}}{\pgfqpoint{3.085801in}{0.407346in}}%
\pgfpathquadraticcurveto{\pgfqpoint{3.090844in}{0.410156in}}{\pgfqpoint{3.097436in}{0.410156in}}%
\pgfpathquadraticcurveto{\pgfqpoint{3.108298in}{0.410156in}}{\pgfqpoint{3.113863in}{0.403437in}}%
\pgfpathquadraticcurveto{\pgfqpoint{3.119447in}{0.396718in}}{\pgfqpoint{3.119447in}{0.383660in}}%
\pgfpathlineto{\pgfqpoint{3.119447in}{0.383660in}}%
\pgfpathclose%
\pgfpathmoveto{\pgfqpoint{3.176618in}{0.398970in}}%
\pgfpathquadraticcurveto{\pgfqpoint{3.174871in}{0.399979in}}{\pgfqpoint{3.172817in}{0.400447in}}%
\pgfpathquadraticcurveto{\pgfqpoint{3.170764in}{0.400933in}}{\pgfqpoint{3.168296in}{0.400933in}}%
\pgfpathquadraticcurveto{\pgfqpoint{3.159506in}{0.400933in}}{\pgfqpoint{3.154805in}{0.395223in}}%
\pgfpathquadraticcurveto{\pgfqpoint{3.150104in}{0.389514in}}{\pgfqpoint{3.150104in}{0.378814in}}%
\pgfpathlineto{\pgfqpoint{3.150104in}{0.345600in}}%
\pgfpathlineto{\pgfqpoint{3.139693in}{0.345600in}}%
\pgfpathlineto{\pgfqpoint{3.139693in}{0.408643in}}%
\pgfpathlineto{\pgfqpoint{3.150104in}{0.408643in}}%
\pgfpathlineto{\pgfqpoint{3.150104in}{0.398844in}}%
\pgfpathquadraticcurveto{\pgfqpoint{3.153382in}{0.404590in}}{\pgfqpoint{3.158606in}{0.407364in}}%
\pgfpathquadraticcurveto{\pgfqpoint{3.163847in}{0.410156in}}{\pgfqpoint{3.171340in}{0.410156in}}%
\pgfpathquadraticcurveto{\pgfqpoint{3.172403in}{0.410156in}}{\pgfqpoint{3.173700in}{0.410011in}}%
\pgfpathquadraticcurveto{\pgfqpoint{3.174997in}{0.409885in}}{\pgfqpoint{3.176564in}{0.409597in}}%
\pgfpathlineto{\pgfqpoint{3.176618in}{0.398970in}}%
\pgfpathlineto{\pgfqpoint{3.176618in}{0.398970in}}%
\pgfpathclose%
\pgfpathmoveto{\pgfqpoint{3.238868in}{0.379715in}}%
\pgfpathlineto{\pgfqpoint{3.238868in}{0.374654in}}%
\pgfpathlineto{\pgfqpoint{3.191244in}{0.374654in}}%
\pgfpathquadraticcurveto{\pgfqpoint{3.191928in}{0.363954in}}{\pgfqpoint{3.197692in}{0.358353in}}%
\pgfpathquadraticcurveto{\pgfqpoint{3.203456in}{0.352751in}}{\pgfqpoint{3.213759in}{0.352751in}}%
\pgfpathquadraticcurveto{\pgfqpoint{3.219721in}{0.352751in}}{\pgfqpoint{3.225323in}{0.354210in}}%
\pgfpathquadraticcurveto{\pgfqpoint{3.230924in}{0.355669in}}{\pgfqpoint{3.236454in}{0.358605in}}%
\pgfpathlineto{\pgfqpoint{3.236454in}{0.348806in}}%
\pgfpathquadraticcurveto{\pgfqpoint{3.230870in}{0.346447in}}{\pgfqpoint{3.225016in}{0.345204in}}%
\pgfpathquadraticcurveto{\pgfqpoint{3.219163in}{0.343961in}}{\pgfqpoint{3.213146in}{0.343961in}}%
\pgfpathquadraticcurveto{\pgfqpoint{3.198052in}{0.343961in}}{\pgfqpoint{3.189244in}{0.352733in}}%
\pgfpathquadraticcurveto{\pgfqpoint{3.180436in}{0.361523in}}{\pgfqpoint{3.180436in}{0.376509in}}%
\pgfpathquadraticcurveto{\pgfqpoint{3.180436in}{0.391981in}}{\pgfqpoint{3.188794in}{0.401059in}}%
\pgfpathquadraticcurveto{\pgfqpoint{3.197152in}{0.410156in}}{\pgfqpoint{3.211345in}{0.410156in}}%
\pgfpathquadraticcurveto{\pgfqpoint{3.224062in}{0.410156in}}{\pgfqpoint{3.231465in}{0.401960in}}%
\pgfpathquadraticcurveto{\pgfqpoint{3.238868in}{0.393783in}}{\pgfqpoint{3.238868in}{0.379715in}}%
\pgfpathlineto{\pgfqpoint{3.238868in}{0.379715in}}%
\pgfpathclose%
\pgfpathmoveto{\pgfqpoint{3.228511in}{0.382759in}}%
\pgfpathquadraticcurveto{\pgfqpoint{3.228403in}{0.391243in}}{\pgfqpoint{3.223756in}{0.396304in}}%
\pgfpathquadraticcurveto{\pgfqpoint{3.219109in}{0.401384in}}{\pgfqpoint{3.211453in}{0.401384in}}%
\pgfpathquadraticcurveto{\pgfqpoint{3.202789in}{0.401384in}}{\pgfqpoint{3.197584in}{0.396484in}}%
\pgfpathquadraticcurveto{\pgfqpoint{3.192378in}{0.391585in}}{\pgfqpoint{3.191586in}{0.382687in}}%
\pgfpathlineto{\pgfqpoint{3.228511in}{0.382759in}}%
\pgfpathlineto{\pgfqpoint{3.228511in}{0.382759in}}%
\pgfpathclose%
\pgfpathmoveto{\pgfqpoint{3.296056in}{0.406787in}}%
\pgfpathlineto{\pgfqpoint{3.296056in}{0.396989in}}%
\pgfpathquadraticcurveto{\pgfqpoint{3.291679in}{0.399240in}}{\pgfqpoint{3.286942in}{0.400357in}}%
\pgfpathquadraticcurveto{\pgfqpoint{3.282223in}{0.401492in}}{\pgfqpoint{3.277144in}{0.401492in}}%
\pgfpathquadraticcurveto{\pgfqpoint{3.269434in}{0.401492in}}{\pgfqpoint{3.265580in}{0.399132in}}%
\pgfpathquadraticcurveto{\pgfqpoint{3.261725in}{0.396773in}}{\pgfqpoint{3.261725in}{0.392035in}}%
\pgfpathquadraticcurveto{\pgfqpoint{3.261725in}{0.388433in}}{\pgfqpoint{3.264481in}{0.386380in}}%
\pgfpathquadraticcurveto{\pgfqpoint{3.267237in}{0.384326in}}{\pgfqpoint{3.275577in}{0.382471in}}%
\pgfpathlineto{\pgfqpoint{3.279125in}{0.381678in}}%
\pgfpathquadraticcurveto{\pgfqpoint{3.290148in}{0.379319in}}{\pgfqpoint{3.294796in}{0.375014in}}%
\pgfpathquadraticcurveto{\pgfqpoint{3.299443in}{0.370709in}}{\pgfqpoint{3.299443in}{0.363000in}}%
\pgfpathquadraticcurveto{\pgfqpoint{3.299443in}{0.354210in}}{\pgfqpoint{3.292490in}{0.349076in}}%
\pgfpathquadraticcurveto{\pgfqpoint{3.285537in}{0.343961in}}{\pgfqpoint{3.273379in}{0.343961in}}%
\pgfpathquadraticcurveto{\pgfqpoint{3.268318in}{0.343961in}}{\pgfqpoint{3.262824in}{0.344952in}}%
\pgfpathquadraticcurveto{\pgfqpoint{3.257330in}{0.345942in}}{\pgfqpoint{3.251260in}{0.347906in}}%
\pgfpathlineto{\pgfqpoint{3.251260in}{0.358605in}}%
\pgfpathquadraticcurveto{\pgfqpoint{3.257006in}{0.355615in}}{\pgfqpoint{3.262572in}{0.354120in}}%
\pgfpathquadraticcurveto{\pgfqpoint{3.268138in}{0.352643in}}{\pgfqpoint{3.273613in}{0.352643in}}%
\pgfpathquadraticcurveto{\pgfqpoint{3.280926in}{0.352643in}}{\pgfqpoint{3.284853in}{0.355146in}}%
\pgfpathquadraticcurveto{\pgfqpoint{3.288798in}{0.357650in}}{\pgfqpoint{3.288798in}{0.362207in}}%
\pgfpathquadraticcurveto{\pgfqpoint{3.288798in}{0.366422in}}{\pgfqpoint{3.285952in}{0.368674in}}%
\pgfpathquadraticcurveto{\pgfqpoint{3.283124in}{0.370925in}}{\pgfqpoint{3.273487in}{0.373014in}}%
\pgfpathlineto{\pgfqpoint{3.269885in}{0.373861in}}%
\pgfpathquadraticcurveto{\pgfqpoint{3.260266in}{0.375878in}}{\pgfqpoint{3.255979in}{0.380075in}}%
\pgfpathquadraticcurveto{\pgfqpoint{3.251711in}{0.384272in}}{\pgfqpoint{3.251711in}{0.391585in}}%
\pgfpathquadraticcurveto{\pgfqpoint{3.251711in}{0.400483in}}{\pgfqpoint{3.258015in}{0.405310in}}%
\pgfpathquadraticcurveto{\pgfqpoint{3.264319in}{0.410156in}}{\pgfqpoint{3.275919in}{0.410156in}}%
\pgfpathquadraticcurveto{\pgfqpoint{3.281647in}{0.410156in}}{\pgfqpoint{3.286708in}{0.409309in}}%
\pgfpathquadraticcurveto{\pgfqpoint{3.291788in}{0.408480in}}{\pgfqpoint{3.296056in}{0.406787in}}%
\pgfpathlineto{\pgfqpoint{3.296056in}{0.406787in}}%
\pgfpathclose%
\pgfpathmoveto{\pgfqpoint{3.326173in}{0.426547in}}%
\pgfpathlineto{\pgfqpoint{3.326173in}{0.408643in}}%
\pgfpathlineto{\pgfqpoint{3.347499in}{0.408643in}}%
\pgfpathlineto{\pgfqpoint{3.347499in}{0.400591in}}%
\pgfpathlineto{\pgfqpoint{3.326173in}{0.400591in}}%
\pgfpathlineto{\pgfqpoint{3.326173in}{0.366368in}}%
\pgfpathquadraticcurveto{\pgfqpoint{3.326173in}{0.358659in}}{\pgfqpoint{3.328280in}{0.356461in}}%
\pgfpathquadraticcurveto{\pgfqpoint{3.330388in}{0.354264in}}{\pgfqpoint{3.336872in}{0.354264in}}%
\pgfpathlineto{\pgfqpoint{3.347499in}{0.354264in}}%
\pgfpathlineto{\pgfqpoint{3.347499in}{0.345600in}}%
\pgfpathlineto{\pgfqpoint{3.336872in}{0.345600in}}%
\pgfpathquadraticcurveto{\pgfqpoint{3.324876in}{0.345600in}}{\pgfqpoint{3.320319in}{0.350067in}}%
\pgfpathquadraticcurveto{\pgfqpoint{3.315762in}{0.354552in}}{\pgfqpoint{3.315762in}{0.366368in}}%
\pgfpathlineto{\pgfqpoint{3.315762in}{0.400591in}}%
\pgfpathlineto{\pgfqpoint{3.308161in}{0.400591in}}%
\pgfpathlineto{\pgfqpoint{3.308161in}{0.408643in}}%
\pgfpathlineto{\pgfqpoint{3.315762in}{0.408643in}}%
\pgfpathlineto{\pgfqpoint{3.315762in}{0.426547in}}%
\pgfpathlineto{\pgfqpoint{3.326173in}{0.426547in}}%
\pgfpathlineto{\pgfqpoint{3.326173in}{0.426547in}}%
\pgfpathclose%
\pgfpathmoveto{\pgfqpoint{3.397645in}{0.398970in}}%
\pgfpathquadraticcurveto{\pgfqpoint{3.395898in}{0.399979in}}{\pgfqpoint{3.393844in}{0.400447in}}%
\pgfpathquadraticcurveto{\pgfqpoint{3.391791in}{0.400933in}}{\pgfqpoint{3.389323in}{0.400933in}}%
\pgfpathquadraticcurveto{\pgfqpoint{3.380533in}{0.400933in}}{\pgfqpoint{3.375832in}{0.395223in}}%
\pgfpathquadraticcurveto{\pgfqpoint{3.371131in}{0.389514in}}{\pgfqpoint{3.371131in}{0.378814in}}%
\pgfpathlineto{\pgfqpoint{3.371131in}{0.345600in}}%
\pgfpathlineto{\pgfqpoint{3.360720in}{0.345600in}}%
\pgfpathlineto{\pgfqpoint{3.360720in}{0.408643in}}%
\pgfpathlineto{\pgfqpoint{3.371131in}{0.408643in}}%
\pgfpathlineto{\pgfqpoint{3.371131in}{0.398844in}}%
\pgfpathquadraticcurveto{\pgfqpoint{3.374409in}{0.404590in}}{\pgfqpoint{3.379633in}{0.407364in}}%
\pgfpathquadraticcurveto{\pgfqpoint{3.384874in}{0.410156in}}{\pgfqpoint{3.392367in}{0.410156in}}%
\pgfpathquadraticcurveto{\pgfqpoint{3.393430in}{0.410156in}}{\pgfqpoint{3.394727in}{0.410011in}}%
\pgfpathquadraticcurveto{\pgfqpoint{3.396024in}{0.409885in}}{\pgfqpoint{3.397591in}{0.409597in}}%
\pgfpathlineto{\pgfqpoint{3.397645in}{0.398970in}}%
\pgfpathlineto{\pgfqpoint{3.397645in}{0.398970in}}%
\pgfpathclose%
\pgfpathmoveto{\pgfqpoint{3.408506in}{0.408643in}}%
\pgfpathlineto{\pgfqpoint{3.418863in}{0.408643in}}%
\pgfpathlineto{\pgfqpoint{3.418863in}{0.345600in}}%
\pgfpathlineto{\pgfqpoint{3.408506in}{0.345600in}}%
\pgfpathlineto{\pgfqpoint{3.408506in}{0.408643in}}%
\pgfpathlineto{\pgfqpoint{3.408506in}{0.408643in}}%
\pgfpathclose%
\pgfpathmoveto{\pgfqpoint{3.408506in}{0.433193in}}%
\pgfpathlineto{\pgfqpoint{3.418863in}{0.433193in}}%
\pgfpathlineto{\pgfqpoint{3.418863in}{0.420062in}}%
\pgfpathlineto{\pgfqpoint{3.408506in}{0.420062in}}%
\pgfpathlineto{\pgfqpoint{3.408506in}{0.433193in}}%
\pgfpathlineto{\pgfqpoint{3.408506in}{0.433193in}}%
\pgfpathclose%
\pgfpathmoveto{\pgfqpoint{3.485905in}{0.406229in}}%
\pgfpathlineto{\pgfqpoint{3.485905in}{0.396538in}}%
\pgfpathquadraticcurveto{\pgfqpoint{3.481510in}{0.398970in}}{\pgfqpoint{3.477097in}{0.400177in}}%
\pgfpathquadraticcurveto{\pgfqpoint{3.472684in}{0.401384in}}{\pgfqpoint{3.468181in}{0.401384in}}%
\pgfpathquadraticcurveto{\pgfqpoint{3.458094in}{0.401384in}}{\pgfqpoint{3.452510in}{0.394989in}}%
\pgfpathquadraticcurveto{\pgfqpoint{3.446944in}{0.388613in}}{\pgfqpoint{3.446944in}{0.377067in}}%
\pgfpathquadraticcurveto{\pgfqpoint{3.446944in}{0.365521in}}{\pgfqpoint{3.452510in}{0.359127in}}%
\pgfpathquadraticcurveto{\pgfqpoint{3.458094in}{0.352751in}}{\pgfqpoint{3.468181in}{0.352751in}}%
\pgfpathquadraticcurveto{\pgfqpoint{3.472684in}{0.352751in}}{\pgfqpoint{3.477097in}{0.353958in}}%
\pgfpathquadraticcurveto{\pgfqpoint{3.481510in}{0.355164in}}{\pgfqpoint{3.485905in}{0.357596in}}%
\pgfpathlineto{\pgfqpoint{3.485905in}{0.348014in}}%
\pgfpathquadraticcurveto{\pgfqpoint{3.481564in}{0.345996in}}{\pgfqpoint{3.476916in}{0.344988in}}%
\pgfpathquadraticcurveto{\pgfqpoint{3.472287in}{0.343961in}}{\pgfqpoint{3.467046in}{0.343961in}}%
\pgfpathquadraticcurveto{\pgfqpoint{3.452798in}{0.343961in}}{\pgfqpoint{3.444405in}{0.352913in}}%
\pgfpathquadraticcurveto{\pgfqpoint{3.436029in}{0.361865in}}{\pgfqpoint{3.436029in}{0.377067in}}%
\pgfpathquadraticcurveto{\pgfqpoint{3.436029in}{0.392486in}}{\pgfqpoint{3.444495in}{0.401312in}}%
\pgfpathquadraticcurveto{\pgfqpoint{3.452978in}{0.410156in}}{\pgfqpoint{3.467730in}{0.410156in}}%
\pgfpathquadraticcurveto{\pgfqpoint{3.472503in}{0.410156in}}{\pgfqpoint{3.477061in}{0.409165in}}%
\pgfpathquadraticcurveto{\pgfqpoint{3.481618in}{0.408192in}}{\pgfqpoint{3.485905in}{0.406229in}}%
\pgfpathlineto{\pgfqpoint{3.485905in}{0.406229in}}%
\pgfpathclose%
\pgfpathmoveto{\pgfqpoint{3.514166in}{0.426547in}}%
\pgfpathlineto{\pgfqpoint{3.514166in}{0.408643in}}%
\pgfpathlineto{\pgfqpoint{3.535492in}{0.408643in}}%
\pgfpathlineto{\pgfqpoint{3.535492in}{0.400591in}}%
\pgfpathlineto{\pgfqpoint{3.514166in}{0.400591in}}%
\pgfpathlineto{\pgfqpoint{3.514166in}{0.366368in}}%
\pgfpathquadraticcurveto{\pgfqpoint{3.514166in}{0.358659in}}{\pgfqpoint{3.516273in}{0.356461in}}%
\pgfpathquadraticcurveto{\pgfqpoint{3.518380in}{0.354264in}}{\pgfqpoint{3.524865in}{0.354264in}}%
\pgfpathlineto{\pgfqpoint{3.535492in}{0.354264in}}%
\pgfpathlineto{\pgfqpoint{3.535492in}{0.345600in}}%
\pgfpathlineto{\pgfqpoint{3.524865in}{0.345600in}}%
\pgfpathquadraticcurveto{\pgfqpoint{3.512869in}{0.345600in}}{\pgfqpoint{3.508312in}{0.350067in}}%
\pgfpathquadraticcurveto{\pgfqpoint{3.503755in}{0.354552in}}{\pgfqpoint{3.503755in}{0.366368in}}%
\pgfpathlineto{\pgfqpoint{3.503755in}{0.400591in}}%
\pgfpathlineto{\pgfqpoint{3.496153in}{0.400591in}}%
\pgfpathlineto{\pgfqpoint{3.496153in}{0.408643in}}%
\pgfpathlineto{\pgfqpoint{3.503755in}{0.408643in}}%
\pgfpathlineto{\pgfqpoint{3.503755in}{0.426547in}}%
\pgfpathlineto{\pgfqpoint{3.514166in}{0.426547in}}%
\pgfpathlineto{\pgfqpoint{3.514166in}{0.426547in}}%
\pgfpathclose%
\pgfpathmoveto{\pgfqpoint{3.603038in}{0.379715in}}%
\pgfpathlineto{\pgfqpoint{3.603038in}{0.374654in}}%
\pgfpathlineto{\pgfqpoint{3.555413in}{0.374654in}}%
\pgfpathquadraticcurveto{\pgfqpoint{3.556098in}{0.363954in}}{\pgfqpoint{3.561862in}{0.358353in}}%
\pgfpathquadraticcurveto{\pgfqpoint{3.567626in}{0.352751in}}{\pgfqpoint{3.577929in}{0.352751in}}%
\pgfpathquadraticcurveto{\pgfqpoint{3.583891in}{0.352751in}}{\pgfqpoint{3.589492in}{0.354210in}}%
\pgfpathquadraticcurveto{\pgfqpoint{3.595094in}{0.355669in}}{\pgfqpoint{3.600624in}{0.358605in}}%
\pgfpathlineto{\pgfqpoint{3.600624in}{0.348806in}}%
\pgfpathquadraticcurveto{\pgfqpoint{3.595040in}{0.346447in}}{\pgfqpoint{3.589186in}{0.345204in}}%
\pgfpathquadraticcurveto{\pgfqpoint{3.583332in}{0.343961in}}{\pgfqpoint{3.577316in}{0.343961in}}%
\pgfpathquadraticcurveto{\pgfqpoint{3.562222in}{0.343961in}}{\pgfqpoint{3.553414in}{0.352733in}}%
\pgfpathquadraticcurveto{\pgfqpoint{3.544606in}{0.361523in}}{\pgfqpoint{3.544606in}{0.376509in}}%
\pgfpathquadraticcurveto{\pgfqpoint{3.544606in}{0.391981in}}{\pgfqpoint{3.552964in}{0.401059in}}%
\pgfpathquadraticcurveto{\pgfqpoint{3.561321in}{0.410156in}}{\pgfqpoint{3.575515in}{0.410156in}}%
\pgfpathquadraticcurveto{\pgfqpoint{3.588232in}{0.410156in}}{\pgfqpoint{3.595635in}{0.401960in}}%
\pgfpathquadraticcurveto{\pgfqpoint{3.603038in}{0.393783in}}{\pgfqpoint{3.603038in}{0.379715in}}%
\pgfpathlineto{\pgfqpoint{3.603038in}{0.379715in}}%
\pgfpathclose%
\pgfpathmoveto{\pgfqpoint{3.592681in}{0.382759in}}%
\pgfpathquadraticcurveto{\pgfqpoint{3.592572in}{0.391243in}}{\pgfqpoint{3.587925in}{0.396304in}}%
\pgfpathquadraticcurveto{\pgfqpoint{3.583278in}{0.401384in}}{\pgfqpoint{3.575623in}{0.401384in}}%
\pgfpathquadraticcurveto{\pgfqpoint{3.566959in}{0.401384in}}{\pgfqpoint{3.561754in}{0.396484in}}%
\pgfpathquadraticcurveto{\pgfqpoint{3.556548in}{0.391585in}}{\pgfqpoint{3.555756in}{0.382687in}}%
\pgfpathlineto{\pgfqpoint{3.592681in}{0.382759in}}%
\pgfpathlineto{\pgfqpoint{3.592681in}{0.382759in}}%
\pgfpathclose%
\pgfpathmoveto{\pgfqpoint{3.661523in}{0.399078in}}%
\pgfpathlineto{\pgfqpoint{3.661523in}{0.433193in}}%
\pgfpathlineto{\pgfqpoint{3.671880in}{0.433193in}}%
\pgfpathlineto{\pgfqpoint{3.671880in}{0.345600in}}%
\pgfpathlineto{\pgfqpoint{3.661523in}{0.345600in}}%
\pgfpathlineto{\pgfqpoint{3.661523in}{0.355056in}}%
\pgfpathquadraticcurveto{\pgfqpoint{3.658263in}{0.349437in}}{\pgfqpoint{3.653273in}{0.346699in}}%
\pgfpathquadraticcurveto{\pgfqpoint{3.648302in}{0.343961in}}{\pgfqpoint{3.641313in}{0.343961in}}%
\pgfpathquadraticcurveto{\pgfqpoint{3.629894in}{0.343961in}}{\pgfqpoint{3.622707in}{0.353075in}}%
\pgfpathquadraticcurveto{\pgfqpoint{3.615538in}{0.362207in}}{\pgfqpoint{3.615538in}{0.377067in}}%
\pgfpathquadraticcurveto{\pgfqpoint{3.615538in}{0.391927in}}{\pgfqpoint{3.622707in}{0.401041in}}%
\pgfpathquadraticcurveto{\pgfqpoint{3.629894in}{0.410156in}}{\pgfqpoint{3.641313in}{0.410156in}}%
\pgfpathquadraticcurveto{\pgfqpoint{3.648302in}{0.410156in}}{\pgfqpoint{3.653273in}{0.407418in}}%
\pgfpathquadraticcurveto{\pgfqpoint{3.658263in}{0.404698in}}{\pgfqpoint{3.661523in}{0.399078in}}%
\pgfpathlineto{\pgfqpoint{3.661523in}{0.399078in}}%
\pgfpathclose%
\pgfpathmoveto{\pgfqpoint{3.626237in}{0.377067in}}%
\pgfpathquadraticcurveto{\pgfqpoint{3.626237in}{0.365648in}}{\pgfqpoint{3.630938in}{0.359145in}}%
\pgfpathquadraticcurveto{\pgfqpoint{3.635640in}{0.352643in}}{\pgfqpoint{3.643853in}{0.352643in}}%
\pgfpathquadraticcurveto{\pgfqpoint{3.652067in}{0.352643in}}{\pgfqpoint{3.656786in}{0.359145in}}%
\pgfpathquadraticcurveto{\pgfqpoint{3.661523in}{0.365648in}}{\pgfqpoint{3.661523in}{0.377067in}}%
\pgfpathquadraticcurveto{\pgfqpoint{3.661523in}{0.388487in}}{\pgfqpoint{3.656786in}{0.394989in}}%
\pgfpathquadraticcurveto{\pgfqpoint{3.652067in}{0.401492in}}{\pgfqpoint{3.643853in}{0.401492in}}%
\pgfpathquadraticcurveto{\pgfqpoint{3.635640in}{0.401492in}}{\pgfqpoint{3.630938in}{0.394989in}}%
\pgfpathquadraticcurveto{\pgfqpoint{3.626237in}{0.388487in}}{\pgfqpoint{3.626237in}{0.377067in}}%
\pgfpathlineto{\pgfqpoint{3.626237in}{0.377067in}}%
\pgfpathclose%
\pgfpathmoveto{\pgfqpoint{3.718658in}{0.429645in}}%
\pgfpathlineto{\pgfqpoint{3.789752in}{0.429645in}}%
\pgfpathlineto{\pgfqpoint{3.789752in}{0.420062in}}%
\pgfpathlineto{\pgfqpoint{3.759923in}{0.420062in}}%
\pgfpathlineto{\pgfqpoint{3.759923in}{0.345600in}}%
\pgfpathlineto{\pgfqpoint{3.748504in}{0.345600in}}%
\pgfpathlineto{\pgfqpoint{3.748504in}{0.420062in}}%
\pgfpathlineto{\pgfqpoint{3.718658in}{0.420062in}}%
\pgfpathlineto{\pgfqpoint{3.718658in}{0.429645in}}%
\pgfpathlineto{\pgfqpoint{3.718658in}{0.429645in}}%
\pgfpathclose%
\pgfpathmoveto{\pgfqpoint{3.793246in}{0.429645in}}%
\pgfpathlineto{\pgfqpoint{3.804720in}{0.429645in}}%
\pgfpathlineto{\pgfqpoint{3.822390in}{0.358605in}}%
\pgfpathlineto{\pgfqpoint{3.840005in}{0.429645in}}%
\pgfpathlineto{\pgfqpoint{3.852794in}{0.429645in}}%
\pgfpathlineto{\pgfqpoint{3.870464in}{0.358605in}}%
\pgfpathlineto{\pgfqpoint{3.888080in}{0.429645in}}%
\pgfpathlineto{\pgfqpoint{3.899626in}{0.429645in}}%
\pgfpathlineto{\pgfqpoint{3.878515in}{0.345600in}}%
\pgfpathlineto{\pgfqpoint{3.864214in}{0.345600in}}%
\pgfpathlineto{\pgfqpoint{3.846490in}{0.418549in}}%
\pgfpathlineto{\pgfqpoint{3.828586in}{0.345600in}}%
\pgfpathlineto{\pgfqpoint{3.814284in}{0.345600in}}%
\pgfpathlineto{\pgfqpoint{3.793246in}{0.429645in}}%
\pgfpathlineto{\pgfqpoint{3.793246in}{0.429645in}}%
\pgfpathclose%
\pgfpathmoveto{\pgfqpoint{3.914702in}{0.429645in}}%
\pgfpathlineto{\pgfqpoint{3.962992in}{0.429645in}}%
\pgfpathlineto{\pgfqpoint{3.962992in}{0.420062in}}%
\pgfpathlineto{\pgfqpoint{3.926067in}{0.420062in}}%
\pgfpathlineto{\pgfqpoint{3.926067in}{0.395296in}}%
\pgfpathlineto{\pgfqpoint{3.959390in}{0.395296in}}%
\pgfpathlineto{\pgfqpoint{3.959390in}{0.385731in}}%
\pgfpathlineto{\pgfqpoint{3.926067in}{0.385731in}}%
\pgfpathlineto{\pgfqpoint{3.926067in}{0.345600in}}%
\pgfpathlineto{\pgfqpoint{3.914702in}{0.345600in}}%
\pgfpathlineto{\pgfqpoint{3.914702in}{0.429645in}}%
\pgfpathlineto{\pgfqpoint{3.914702in}{0.429645in}}%
\pgfpathclose%
\pgfpathmoveto{\pgfqpoint{3.981005in}{0.429645in}}%
\pgfpathlineto{\pgfqpoint{4.034140in}{0.429645in}}%
\pgfpathlineto{\pgfqpoint{4.034140in}{0.420062in}}%
\pgfpathlineto{\pgfqpoint{3.992370in}{0.420062in}}%
\pgfpathlineto{\pgfqpoint{3.992370in}{0.395187in}}%
\pgfpathlineto{\pgfqpoint{4.032393in}{0.395187in}}%
\pgfpathlineto{\pgfqpoint{4.032393in}{0.385623in}}%
\pgfpathlineto{\pgfqpoint{3.992370in}{0.385623in}}%
\pgfpathlineto{\pgfqpoint{3.992370in}{0.355164in}}%
\pgfpathlineto{\pgfqpoint{4.035149in}{0.355164in}}%
\pgfpathlineto{\pgfqpoint{4.035149in}{0.345600in}}%
\pgfpathlineto{\pgfqpoint{3.981005in}{0.345600in}}%
\pgfpathlineto{\pgfqpoint{3.981005in}{0.429645in}}%
\pgfpathlineto{\pgfqpoint{3.981005in}{0.429645in}}%
\pgfpathclose%
\pgfpathmoveto{\pgfqpoint{4.056043in}{0.405202in}}%
\pgfpathlineto{\pgfqpoint{4.067913in}{0.405202in}}%
\pgfpathlineto{\pgfqpoint{4.067913in}{0.390919in}}%
\pgfpathlineto{\pgfqpoint{4.056043in}{0.390919in}}%
\pgfpathlineto{\pgfqpoint{4.056043in}{0.405202in}}%
\pgfpathlineto{\pgfqpoint{4.056043in}{0.405202in}}%
\pgfpathclose%
\pgfpathmoveto{\pgfqpoint{4.056043in}{0.359902in}}%
\pgfpathlineto{\pgfqpoint{4.067913in}{0.359902in}}%
\pgfpathlineto{\pgfqpoint{4.067913in}{0.350211in}}%
\pgfpathlineto{\pgfqpoint{4.058691in}{0.332199in}}%
\pgfpathlineto{\pgfqpoint{4.051432in}{0.332199in}}%
\pgfpathlineto{\pgfqpoint{4.056043in}{0.350211in}}%
\pgfpathlineto{\pgfqpoint{4.056043in}{0.359902in}}%
\pgfpathlineto{\pgfqpoint{4.056043in}{0.359902in}}%
\pgfpathclose%
\pgfpathmoveto{\pgfqpoint{4.128866in}{0.408643in}}%
\pgfpathlineto{\pgfqpoint{4.139223in}{0.408643in}}%
\pgfpathlineto{\pgfqpoint{4.139223in}{0.345600in}}%
\pgfpathlineto{\pgfqpoint{4.128866in}{0.345600in}}%
\pgfpathlineto{\pgfqpoint{4.128866in}{0.408643in}}%
\pgfpathlineto{\pgfqpoint{4.128866in}{0.408643in}}%
\pgfpathclose%
\pgfpathmoveto{\pgfqpoint{4.128866in}{0.433193in}}%
\pgfpathlineto{\pgfqpoint{4.139223in}{0.433193in}}%
\pgfpathlineto{\pgfqpoint{4.139223in}{0.420062in}}%
\pgfpathlineto{\pgfqpoint{4.128866in}{0.420062in}}%
\pgfpathlineto{\pgfqpoint{4.128866in}{0.433193in}}%
\pgfpathlineto{\pgfqpoint{4.128866in}{0.433193in}}%
\pgfpathclose%
\pgfpathmoveto{\pgfqpoint{4.171141in}{0.426547in}}%
\pgfpathlineto{\pgfqpoint{4.171141in}{0.408643in}}%
\pgfpathlineto{\pgfqpoint{4.192467in}{0.408643in}}%
\pgfpathlineto{\pgfqpoint{4.192467in}{0.400591in}}%
\pgfpathlineto{\pgfqpoint{4.171141in}{0.400591in}}%
\pgfpathlineto{\pgfqpoint{4.171141in}{0.366368in}}%
\pgfpathquadraticcurveto{\pgfqpoint{4.171141in}{0.358659in}}{\pgfqpoint{4.173248in}{0.356461in}}%
\pgfpathquadraticcurveto{\pgfqpoint{4.175356in}{0.354264in}}{\pgfqpoint{4.181840in}{0.354264in}}%
\pgfpathlineto{\pgfqpoint{4.192467in}{0.354264in}}%
\pgfpathlineto{\pgfqpoint{4.192467in}{0.345600in}}%
\pgfpathlineto{\pgfqpoint{4.181840in}{0.345600in}}%
\pgfpathquadraticcurveto{\pgfqpoint{4.169844in}{0.345600in}}{\pgfqpoint{4.165287in}{0.350067in}}%
\pgfpathquadraticcurveto{\pgfqpoint{4.160730in}{0.354552in}}{\pgfqpoint{4.160730in}{0.366368in}}%
\pgfpathlineto{\pgfqpoint{4.160730in}{0.400591in}}%
\pgfpathlineto{\pgfqpoint{4.153129in}{0.400591in}}%
\pgfpathlineto{\pgfqpoint{4.153129in}{0.408643in}}%
\pgfpathlineto{\pgfqpoint{4.160730in}{0.408643in}}%
\pgfpathlineto{\pgfqpoint{4.160730in}{0.426547in}}%
\pgfpathlineto{\pgfqpoint{4.171141in}{0.426547in}}%
\pgfpathlineto{\pgfqpoint{4.171141in}{0.426547in}}%
\pgfpathclose%
\pgfpathmoveto{\pgfqpoint{4.246270in}{0.406787in}}%
\pgfpathlineto{\pgfqpoint{4.246270in}{0.396989in}}%
\pgfpathquadraticcurveto{\pgfqpoint{4.241893in}{0.399240in}}{\pgfqpoint{4.237155in}{0.400357in}}%
\pgfpathquadraticcurveto{\pgfqpoint{4.232436in}{0.401492in}}{\pgfqpoint{4.227357in}{0.401492in}}%
\pgfpathquadraticcurveto{\pgfqpoint{4.219648in}{0.401492in}}{\pgfqpoint{4.215793in}{0.399132in}}%
\pgfpathquadraticcurveto{\pgfqpoint{4.211938in}{0.396773in}}{\pgfqpoint{4.211938in}{0.392035in}}%
\pgfpathquadraticcurveto{\pgfqpoint{4.211938in}{0.388433in}}{\pgfqpoint{4.214694in}{0.386380in}}%
\pgfpathquadraticcurveto{\pgfqpoint{4.217450in}{0.384326in}}{\pgfqpoint{4.225790in}{0.382471in}}%
\pgfpathlineto{\pgfqpoint{4.229338in}{0.381678in}}%
\pgfpathquadraticcurveto{\pgfqpoint{4.240362in}{0.379319in}}{\pgfqpoint{4.245009in}{0.375014in}}%
\pgfpathquadraticcurveto{\pgfqpoint{4.249656in}{0.370709in}}{\pgfqpoint{4.249656in}{0.363000in}}%
\pgfpathquadraticcurveto{\pgfqpoint{4.249656in}{0.354210in}}{\pgfqpoint{4.242703in}{0.349076in}}%
\pgfpathquadraticcurveto{\pgfqpoint{4.235750in}{0.343961in}}{\pgfqpoint{4.223592in}{0.343961in}}%
\pgfpathquadraticcurveto{\pgfqpoint{4.218531in}{0.343961in}}{\pgfqpoint{4.213037in}{0.344952in}}%
\pgfpathquadraticcurveto{\pgfqpoint{4.207543in}{0.345942in}}{\pgfqpoint{4.201473in}{0.347906in}}%
\pgfpathlineto{\pgfqpoint{4.201473in}{0.358605in}}%
\pgfpathquadraticcurveto{\pgfqpoint{4.207219in}{0.355615in}}{\pgfqpoint{4.212785in}{0.354120in}}%
\pgfpathquadraticcurveto{\pgfqpoint{4.218351in}{0.352643in}}{\pgfqpoint{4.223826in}{0.352643in}}%
\pgfpathquadraticcurveto{\pgfqpoint{4.231139in}{0.352643in}}{\pgfqpoint{4.235066in}{0.355146in}}%
\pgfpathquadraticcurveto{\pgfqpoint{4.239011in}{0.357650in}}{\pgfqpoint{4.239011in}{0.362207in}}%
\pgfpathquadraticcurveto{\pgfqpoint{4.239011in}{0.366422in}}{\pgfqpoint{4.236165in}{0.368674in}}%
\pgfpathquadraticcurveto{\pgfqpoint{4.233337in}{0.370925in}}{\pgfqpoint{4.223700in}{0.373014in}}%
\pgfpathlineto{\pgfqpoint{4.220098in}{0.373861in}}%
\pgfpathquadraticcurveto{\pgfqpoint{4.210479in}{0.375878in}}{\pgfqpoint{4.206192in}{0.380075in}}%
\pgfpathquadraticcurveto{\pgfqpoint{4.201924in}{0.384272in}}{\pgfqpoint{4.201924in}{0.391585in}}%
\pgfpathquadraticcurveto{\pgfqpoint{4.201924in}{0.400483in}}{\pgfqpoint{4.208228in}{0.405310in}}%
\pgfpathquadraticcurveto{\pgfqpoint{4.214532in}{0.410156in}}{\pgfqpoint{4.226132in}{0.410156in}}%
\pgfpathquadraticcurveto{\pgfqpoint{4.231860in}{0.410156in}}{\pgfqpoint{4.236921in}{0.409309in}}%
\pgfpathquadraticcurveto{\pgfqpoint{4.242001in}{0.408480in}}{\pgfqpoint{4.246270in}{0.406787in}}%
\pgfpathlineto{\pgfqpoint{4.246270in}{0.406787in}}%
\pgfpathclose%
\pgfpathmoveto{\pgfqpoint{4.339302in}{0.398970in}}%
\pgfpathquadraticcurveto{\pgfqpoint{4.337555in}{0.399979in}}{\pgfqpoint{4.335502in}{0.400447in}}%
\pgfpathquadraticcurveto{\pgfqpoint{4.333448in}{0.400933in}}{\pgfqpoint{4.330981in}{0.400933in}}%
\pgfpathquadraticcurveto{\pgfqpoint{4.322191in}{0.400933in}}{\pgfqpoint{4.317490in}{0.395223in}}%
\pgfpathquadraticcurveto{\pgfqpoint{4.312788in}{0.389514in}}{\pgfqpoint{4.312788in}{0.378814in}}%
\pgfpathlineto{\pgfqpoint{4.312788in}{0.345600in}}%
\pgfpathlineto{\pgfqpoint{4.302377in}{0.345600in}}%
\pgfpathlineto{\pgfqpoint{4.302377in}{0.408643in}}%
\pgfpathlineto{\pgfqpoint{4.312788in}{0.408643in}}%
\pgfpathlineto{\pgfqpoint{4.312788in}{0.398844in}}%
\pgfpathquadraticcurveto{\pgfqpoint{4.316067in}{0.404590in}}{\pgfqpoint{4.321290in}{0.407364in}}%
\pgfpathquadraticcurveto{\pgfqpoint{4.326532in}{0.410156in}}{\pgfqpoint{4.334025in}{0.410156in}}%
\pgfpathquadraticcurveto{\pgfqpoint{4.335087in}{0.410156in}}{\pgfqpoint{4.336384in}{0.410011in}}%
\pgfpathquadraticcurveto{\pgfqpoint{4.337681in}{0.409885in}}{\pgfqpoint{4.339248in}{0.409597in}}%
\pgfpathlineto{\pgfqpoint{4.339302in}{0.398970in}}%
\pgfpathlineto{\pgfqpoint{4.339302in}{0.398970in}}%
\pgfpathclose%
\pgfpathmoveto{\pgfqpoint{4.401552in}{0.379715in}}%
\pgfpathlineto{\pgfqpoint{4.401552in}{0.374654in}}%
\pgfpathlineto{\pgfqpoint{4.353928in}{0.374654in}}%
\pgfpathquadraticcurveto{\pgfqpoint{4.354613in}{0.363954in}}{\pgfqpoint{4.360377in}{0.358353in}}%
\pgfpathquadraticcurveto{\pgfqpoint{4.366140in}{0.352751in}}{\pgfqpoint{4.376443in}{0.352751in}}%
\pgfpathquadraticcurveto{\pgfqpoint{4.382405in}{0.352751in}}{\pgfqpoint{4.388007in}{0.354210in}}%
\pgfpathquadraticcurveto{\pgfqpoint{4.393609in}{0.355669in}}{\pgfqpoint{4.399139in}{0.358605in}}%
\pgfpathlineto{\pgfqpoint{4.399139in}{0.348806in}}%
\pgfpathquadraticcurveto{\pgfqpoint{4.393555in}{0.346447in}}{\pgfqpoint{4.387701in}{0.345204in}}%
\pgfpathquadraticcurveto{\pgfqpoint{4.381847in}{0.343961in}}{\pgfqpoint{4.375831in}{0.343961in}}%
\pgfpathquadraticcurveto{\pgfqpoint{4.360737in}{0.343961in}}{\pgfqpoint{4.351929in}{0.352733in}}%
\pgfpathquadraticcurveto{\pgfqpoint{4.343121in}{0.361523in}}{\pgfqpoint{4.343121in}{0.376509in}}%
\pgfpathquadraticcurveto{\pgfqpoint{4.343121in}{0.391981in}}{\pgfqpoint{4.351479in}{0.401059in}}%
\pgfpathquadraticcurveto{\pgfqpoint{4.359836in}{0.410156in}}{\pgfqpoint{4.374030in}{0.410156in}}%
\pgfpathquadraticcurveto{\pgfqpoint{4.386746in}{0.410156in}}{\pgfqpoint{4.394149in}{0.401960in}}%
\pgfpathquadraticcurveto{\pgfqpoint{4.401552in}{0.393783in}}{\pgfqpoint{4.401552in}{0.379715in}}%
\pgfpathlineto{\pgfqpoint{4.401552in}{0.379715in}}%
\pgfpathclose%
\pgfpathmoveto{\pgfqpoint{4.391195in}{0.382759in}}%
\pgfpathquadraticcurveto{\pgfqpoint{4.391087in}{0.391243in}}{\pgfqpoint{4.386440in}{0.396304in}}%
\pgfpathquadraticcurveto{\pgfqpoint{4.381793in}{0.401384in}}{\pgfqpoint{4.374138in}{0.401384in}}%
\pgfpathquadraticcurveto{\pgfqpoint{4.365474in}{0.401384in}}{\pgfqpoint{4.360268in}{0.396484in}}%
\pgfpathquadraticcurveto{\pgfqpoint{4.355063in}{0.391585in}}{\pgfqpoint{4.354270in}{0.382687in}}%
\pgfpathlineto{\pgfqpoint{4.391195in}{0.382759in}}%
\pgfpathlineto{\pgfqpoint{4.391195in}{0.382759in}}%
\pgfpathclose%
\pgfpathmoveto{\pgfqpoint{4.458741in}{0.406787in}}%
\pgfpathlineto{\pgfqpoint{4.458741in}{0.396989in}}%
\pgfpathquadraticcurveto{\pgfqpoint{4.454364in}{0.399240in}}{\pgfqpoint{4.449627in}{0.400357in}}%
\pgfpathquadraticcurveto{\pgfqpoint{4.444908in}{0.401492in}}{\pgfqpoint{4.439828in}{0.401492in}}%
\pgfpathquadraticcurveto{\pgfqpoint{4.432119in}{0.401492in}}{\pgfqpoint{4.428264in}{0.399132in}}%
\pgfpathquadraticcurveto{\pgfqpoint{4.424410in}{0.396773in}}{\pgfqpoint{4.424410in}{0.392035in}}%
\pgfpathquadraticcurveto{\pgfqpoint{4.424410in}{0.388433in}}{\pgfqpoint{4.427166in}{0.386380in}}%
\pgfpathquadraticcurveto{\pgfqpoint{4.429921in}{0.384326in}}{\pgfqpoint{4.438261in}{0.382471in}}%
\pgfpathlineto{\pgfqpoint{4.441809in}{0.381678in}}%
\pgfpathquadraticcurveto{\pgfqpoint{4.452833in}{0.379319in}}{\pgfqpoint{4.457480in}{0.375014in}}%
\pgfpathquadraticcurveto{\pgfqpoint{4.462127in}{0.370709in}}{\pgfqpoint{4.462127in}{0.363000in}}%
\pgfpathquadraticcurveto{\pgfqpoint{4.462127in}{0.354210in}}{\pgfqpoint{4.455174in}{0.349076in}}%
\pgfpathquadraticcurveto{\pgfqpoint{4.448222in}{0.343961in}}{\pgfqpoint{4.436064in}{0.343961in}}%
\pgfpathquadraticcurveto{\pgfqpoint{4.431002in}{0.343961in}}{\pgfqpoint{4.425508in}{0.344952in}}%
\pgfpathquadraticcurveto{\pgfqpoint{4.420015in}{0.345942in}}{\pgfqpoint{4.413945in}{0.347906in}}%
\pgfpathlineto{\pgfqpoint{4.413945in}{0.358605in}}%
\pgfpathquadraticcurveto{\pgfqpoint{4.419691in}{0.355615in}}{\pgfqpoint{4.425256in}{0.354120in}}%
\pgfpathquadraticcurveto{\pgfqpoint{4.430822in}{0.352643in}}{\pgfqpoint{4.436298in}{0.352643in}}%
\pgfpathquadraticcurveto{\pgfqpoint{4.443611in}{0.352643in}}{\pgfqpoint{4.447537in}{0.355146in}}%
\pgfpathquadraticcurveto{\pgfqpoint{4.451482in}{0.357650in}}{\pgfqpoint{4.451482in}{0.362207in}}%
\pgfpathquadraticcurveto{\pgfqpoint{4.451482in}{0.366422in}}{\pgfqpoint{4.448636in}{0.368674in}}%
\pgfpathquadraticcurveto{\pgfqpoint{4.445808in}{0.370925in}}{\pgfqpoint{4.436172in}{0.373014in}}%
\pgfpathlineto{\pgfqpoint{4.432569in}{0.373861in}}%
\pgfpathquadraticcurveto{\pgfqpoint{4.422951in}{0.375878in}}{\pgfqpoint{4.418664in}{0.380075in}}%
\pgfpathquadraticcurveto{\pgfqpoint{4.414395in}{0.384272in}}{\pgfqpoint{4.414395in}{0.391585in}}%
\pgfpathquadraticcurveto{\pgfqpoint{4.414395in}{0.400483in}}{\pgfqpoint{4.420699in}{0.405310in}}%
\pgfpathquadraticcurveto{\pgfqpoint{4.427003in}{0.410156in}}{\pgfqpoint{4.438603in}{0.410156in}}%
\pgfpathquadraticcurveto{\pgfqpoint{4.444331in}{0.410156in}}{\pgfqpoint{4.449393in}{0.409309in}}%
\pgfpathquadraticcurveto{\pgfqpoint{4.454472in}{0.408480in}}{\pgfqpoint{4.458741in}{0.406787in}}%
\pgfpathlineto{\pgfqpoint{4.458741in}{0.406787in}}%
\pgfpathclose%
\pgfpathmoveto{\pgfqpoint{4.488857in}{0.426547in}}%
\pgfpathlineto{\pgfqpoint{4.488857in}{0.408643in}}%
\pgfpathlineto{\pgfqpoint{4.510184in}{0.408643in}}%
\pgfpathlineto{\pgfqpoint{4.510184in}{0.400591in}}%
\pgfpathlineto{\pgfqpoint{4.488857in}{0.400591in}}%
\pgfpathlineto{\pgfqpoint{4.488857in}{0.366368in}}%
\pgfpathquadraticcurveto{\pgfqpoint{4.488857in}{0.358659in}}{\pgfqpoint{4.490965in}{0.356461in}}%
\pgfpathquadraticcurveto{\pgfqpoint{4.493072in}{0.354264in}}{\pgfqpoint{4.499556in}{0.354264in}}%
\pgfpathlineto{\pgfqpoint{4.510184in}{0.354264in}}%
\pgfpathlineto{\pgfqpoint{4.510184in}{0.345600in}}%
\pgfpathlineto{\pgfqpoint{4.499556in}{0.345600in}}%
\pgfpathquadraticcurveto{\pgfqpoint{4.487560in}{0.345600in}}{\pgfqpoint{4.483003in}{0.350067in}}%
\pgfpathquadraticcurveto{\pgfqpoint{4.478446in}{0.354552in}}{\pgfqpoint{4.478446in}{0.366368in}}%
\pgfpathlineto{\pgfqpoint{4.478446in}{0.400591in}}%
\pgfpathlineto{\pgfqpoint{4.470845in}{0.400591in}}%
\pgfpathlineto{\pgfqpoint{4.470845in}{0.408643in}}%
\pgfpathlineto{\pgfqpoint{4.478446in}{0.408643in}}%
\pgfpathlineto{\pgfqpoint{4.478446in}{0.426547in}}%
\pgfpathlineto{\pgfqpoint{4.488857in}{0.426547in}}%
\pgfpathlineto{\pgfqpoint{4.488857in}{0.426547in}}%
\pgfpathclose%
\pgfpathmoveto{\pgfqpoint{4.560329in}{0.398970in}}%
\pgfpathquadraticcurveto{\pgfqpoint{4.558582in}{0.399979in}}{\pgfqpoint{4.556529in}{0.400447in}}%
\pgfpathquadraticcurveto{\pgfqpoint{4.554475in}{0.400933in}}{\pgfqpoint{4.552008in}{0.400933in}}%
\pgfpathquadraticcurveto{\pgfqpoint{4.543218in}{0.400933in}}{\pgfqpoint{4.538517in}{0.395223in}}%
\pgfpathquadraticcurveto{\pgfqpoint{4.533816in}{0.389514in}}{\pgfqpoint{4.533816in}{0.378814in}}%
\pgfpathlineto{\pgfqpoint{4.533816in}{0.345600in}}%
\pgfpathlineto{\pgfqpoint{4.523405in}{0.345600in}}%
\pgfpathlineto{\pgfqpoint{4.523405in}{0.408643in}}%
\pgfpathlineto{\pgfqpoint{4.533816in}{0.408643in}}%
\pgfpathlineto{\pgfqpoint{4.533816in}{0.398844in}}%
\pgfpathquadraticcurveto{\pgfqpoint{4.537094in}{0.404590in}}{\pgfqpoint{4.542317in}{0.407364in}}%
\pgfpathquadraticcurveto{\pgfqpoint{4.547559in}{0.410156in}}{\pgfqpoint{4.555052in}{0.410156in}}%
\pgfpathquadraticcurveto{\pgfqpoint{4.556115in}{0.410156in}}{\pgfqpoint{4.557411in}{0.410011in}}%
\pgfpathquadraticcurveto{\pgfqpoint{4.558708in}{0.409885in}}{\pgfqpoint{4.560275in}{0.409597in}}%
\pgfpathlineto{\pgfqpoint{4.560329in}{0.398970in}}%
\pgfpathlineto{\pgfqpoint{4.560329in}{0.398970in}}%
\pgfpathclose%
\pgfpathmoveto{\pgfqpoint{4.571191in}{0.408643in}}%
\pgfpathlineto{\pgfqpoint{4.581548in}{0.408643in}}%
\pgfpathlineto{\pgfqpoint{4.581548in}{0.345600in}}%
\pgfpathlineto{\pgfqpoint{4.571191in}{0.345600in}}%
\pgfpathlineto{\pgfqpoint{4.571191in}{0.408643in}}%
\pgfpathlineto{\pgfqpoint{4.571191in}{0.408643in}}%
\pgfpathclose%
\pgfpathmoveto{\pgfqpoint{4.571191in}{0.433193in}}%
\pgfpathlineto{\pgfqpoint{4.581548in}{0.433193in}}%
\pgfpathlineto{\pgfqpoint{4.581548in}{0.420062in}}%
\pgfpathlineto{\pgfqpoint{4.571191in}{0.420062in}}%
\pgfpathlineto{\pgfqpoint{4.571191in}{0.433193in}}%
\pgfpathlineto{\pgfqpoint{4.571191in}{0.433193in}}%
\pgfpathclose%
\pgfpathmoveto{\pgfqpoint{4.648589in}{0.406229in}}%
\pgfpathlineto{\pgfqpoint{4.648589in}{0.396538in}}%
\pgfpathquadraticcurveto{\pgfqpoint{4.644194in}{0.398970in}}{\pgfqpoint{4.639781in}{0.400177in}}%
\pgfpathquadraticcurveto{\pgfqpoint{4.635368in}{0.401384in}}{\pgfqpoint{4.630865in}{0.401384in}}%
\pgfpathquadraticcurveto{\pgfqpoint{4.620778in}{0.401384in}}{\pgfqpoint{4.615194in}{0.394989in}}%
\pgfpathquadraticcurveto{\pgfqpoint{4.609629in}{0.388613in}}{\pgfqpoint{4.609629in}{0.377067in}}%
\pgfpathquadraticcurveto{\pgfqpoint{4.609629in}{0.365521in}}{\pgfqpoint{4.615194in}{0.359127in}}%
\pgfpathquadraticcurveto{\pgfqpoint{4.620778in}{0.352751in}}{\pgfqpoint{4.630865in}{0.352751in}}%
\pgfpathquadraticcurveto{\pgfqpoint{4.635368in}{0.352751in}}{\pgfqpoint{4.639781in}{0.353958in}}%
\pgfpathquadraticcurveto{\pgfqpoint{4.644194in}{0.355164in}}{\pgfqpoint{4.648589in}{0.357596in}}%
\pgfpathlineto{\pgfqpoint{4.648589in}{0.348014in}}%
\pgfpathquadraticcurveto{\pgfqpoint{4.644248in}{0.345996in}}{\pgfqpoint{4.639601in}{0.344988in}}%
\pgfpathquadraticcurveto{\pgfqpoint{4.634972in}{0.343961in}}{\pgfqpoint{4.629730in}{0.343961in}}%
\pgfpathquadraticcurveto{\pgfqpoint{4.615483in}{0.343961in}}{\pgfqpoint{4.607089in}{0.352913in}}%
\pgfpathquadraticcurveto{\pgfqpoint{4.598713in}{0.361865in}}{\pgfqpoint{4.598713in}{0.377067in}}%
\pgfpathquadraticcurveto{\pgfqpoint{4.598713in}{0.392486in}}{\pgfqpoint{4.607179in}{0.401312in}}%
\pgfpathquadraticcurveto{\pgfqpoint{4.615663in}{0.410156in}}{\pgfqpoint{4.630415in}{0.410156in}}%
\pgfpathquadraticcurveto{\pgfqpoint{4.635188in}{0.410156in}}{\pgfqpoint{4.639745in}{0.409165in}}%
\pgfpathquadraticcurveto{\pgfqpoint{4.644302in}{0.408192in}}{\pgfqpoint{4.648589in}{0.406229in}}%
\pgfpathlineto{\pgfqpoint{4.648589in}{0.406229in}}%
\pgfpathclose%
\pgfpathmoveto{\pgfqpoint{4.676850in}{0.426547in}}%
\pgfpathlineto{\pgfqpoint{4.676850in}{0.408643in}}%
\pgfpathlineto{\pgfqpoint{4.698176in}{0.408643in}}%
\pgfpathlineto{\pgfqpoint{4.698176in}{0.400591in}}%
\pgfpathlineto{\pgfqpoint{4.676850in}{0.400591in}}%
\pgfpathlineto{\pgfqpoint{4.676850in}{0.366368in}}%
\pgfpathquadraticcurveto{\pgfqpoint{4.676850in}{0.358659in}}{\pgfqpoint{4.678957in}{0.356461in}}%
\pgfpathquadraticcurveto{\pgfqpoint{4.681065in}{0.354264in}}{\pgfqpoint{4.687549in}{0.354264in}}%
\pgfpathlineto{\pgfqpoint{4.698176in}{0.354264in}}%
\pgfpathlineto{\pgfqpoint{4.698176in}{0.345600in}}%
\pgfpathlineto{\pgfqpoint{4.687549in}{0.345600in}}%
\pgfpathquadraticcurveto{\pgfqpoint{4.675553in}{0.345600in}}{\pgfqpoint{4.670996in}{0.350067in}}%
\pgfpathquadraticcurveto{\pgfqpoint{4.666439in}{0.354552in}}{\pgfqpoint{4.666439in}{0.366368in}}%
\pgfpathlineto{\pgfqpoint{4.666439in}{0.400591in}}%
\pgfpathlineto{\pgfqpoint{4.658838in}{0.400591in}}%
\pgfpathlineto{\pgfqpoint{4.658838in}{0.408643in}}%
\pgfpathlineto{\pgfqpoint{4.666439in}{0.408643in}}%
\pgfpathlineto{\pgfqpoint{4.666439in}{0.426547in}}%
\pgfpathlineto{\pgfqpoint{4.676850in}{0.426547in}}%
\pgfpathlineto{\pgfqpoint{4.676850in}{0.426547in}}%
\pgfpathclose%
\pgfpathmoveto{\pgfqpoint{4.765722in}{0.379715in}}%
\pgfpathlineto{\pgfqpoint{4.765722in}{0.374654in}}%
\pgfpathlineto{\pgfqpoint{4.718098in}{0.374654in}}%
\pgfpathquadraticcurveto{\pgfqpoint{4.718782in}{0.363954in}}{\pgfqpoint{4.724546in}{0.358353in}}%
\pgfpathquadraticcurveto{\pgfqpoint{4.730310in}{0.352751in}}{\pgfqpoint{4.740613in}{0.352751in}}%
\pgfpathquadraticcurveto{\pgfqpoint{4.746575in}{0.352751in}}{\pgfqpoint{4.752177in}{0.354210in}}%
\pgfpathquadraticcurveto{\pgfqpoint{4.757779in}{0.355669in}}{\pgfqpoint{4.763308in}{0.358605in}}%
\pgfpathlineto{\pgfqpoint{4.763308in}{0.348806in}}%
\pgfpathquadraticcurveto{\pgfqpoint{4.757725in}{0.346447in}}{\pgfqpoint{4.751871in}{0.345204in}}%
\pgfpathquadraticcurveto{\pgfqpoint{4.746017in}{0.343961in}}{\pgfqpoint{4.740001in}{0.343961in}}%
\pgfpathquadraticcurveto{\pgfqpoint{4.724906in}{0.343961in}}{\pgfqpoint{4.716099in}{0.352733in}}%
\pgfpathquadraticcurveto{\pgfqpoint{4.707291in}{0.361523in}}{\pgfqpoint{4.707291in}{0.376509in}}%
\pgfpathquadraticcurveto{\pgfqpoint{4.707291in}{0.391981in}}{\pgfqpoint{4.715648in}{0.401059in}}%
\pgfpathquadraticcurveto{\pgfqpoint{4.724006in}{0.410156in}}{\pgfqpoint{4.738199in}{0.410156in}}%
\pgfpathquadraticcurveto{\pgfqpoint{4.750916in}{0.410156in}}{\pgfqpoint{4.758319in}{0.401960in}}%
\pgfpathquadraticcurveto{\pgfqpoint{4.765722in}{0.393783in}}{\pgfqpoint{4.765722in}{0.379715in}}%
\pgfpathlineto{\pgfqpoint{4.765722in}{0.379715in}}%
\pgfpathclose%
\pgfpathmoveto{\pgfqpoint{4.755365in}{0.382759in}}%
\pgfpathquadraticcurveto{\pgfqpoint{4.755257in}{0.391243in}}{\pgfqpoint{4.750610in}{0.396304in}}%
\pgfpathquadraticcurveto{\pgfqpoint{4.745963in}{0.401384in}}{\pgfqpoint{4.738308in}{0.401384in}}%
\pgfpathquadraticcurveto{\pgfqpoint{4.729644in}{0.401384in}}{\pgfqpoint{4.724438in}{0.396484in}}%
\pgfpathquadraticcurveto{\pgfqpoint{4.719233in}{0.391585in}}{\pgfqpoint{4.718440in}{0.382687in}}%
\pgfpathlineto{\pgfqpoint{4.755365in}{0.382759in}}%
\pgfpathlineto{\pgfqpoint{4.755365in}{0.382759in}}%
\pgfpathclose%
\pgfpathmoveto{\pgfqpoint{4.824207in}{0.399078in}}%
\pgfpathlineto{\pgfqpoint{4.824207in}{0.433193in}}%
\pgfpathlineto{\pgfqpoint{4.834564in}{0.433193in}}%
\pgfpathlineto{\pgfqpoint{4.834564in}{0.345600in}}%
\pgfpathlineto{\pgfqpoint{4.824207in}{0.345600in}}%
\pgfpathlineto{\pgfqpoint{4.824207in}{0.355056in}}%
\pgfpathquadraticcurveto{\pgfqpoint{4.820947in}{0.349437in}}{\pgfqpoint{4.815958in}{0.346699in}}%
\pgfpathquadraticcurveto{\pgfqpoint{4.810987in}{0.343961in}}{\pgfqpoint{4.803998in}{0.343961in}}%
\pgfpathquadraticcurveto{\pgfqpoint{4.792578in}{0.343961in}}{\pgfqpoint{4.785391in}{0.353075in}}%
\pgfpathquadraticcurveto{\pgfqpoint{4.778222in}{0.362207in}}{\pgfqpoint{4.778222in}{0.377067in}}%
\pgfpathquadraticcurveto{\pgfqpoint{4.778222in}{0.391927in}}{\pgfqpoint{4.785391in}{0.401041in}}%
\pgfpathquadraticcurveto{\pgfqpoint{4.792578in}{0.410156in}}{\pgfqpoint{4.803998in}{0.410156in}}%
\pgfpathquadraticcurveto{\pgfqpoint{4.810987in}{0.410156in}}{\pgfqpoint{4.815958in}{0.407418in}}%
\pgfpathquadraticcurveto{\pgfqpoint{4.820947in}{0.404698in}}{\pgfqpoint{4.824207in}{0.399078in}}%
\pgfpathlineto{\pgfqpoint{4.824207in}{0.399078in}}%
\pgfpathclose%
\pgfpathmoveto{\pgfqpoint{4.788922in}{0.377067in}}%
\pgfpathquadraticcurveto{\pgfqpoint{4.788922in}{0.365648in}}{\pgfqpoint{4.793623in}{0.359145in}}%
\pgfpathquadraticcurveto{\pgfqpoint{4.798324in}{0.352643in}}{\pgfqpoint{4.806538in}{0.352643in}}%
\pgfpathquadraticcurveto{\pgfqpoint{4.814751in}{0.352643in}}{\pgfqpoint{4.819470in}{0.359145in}}%
\pgfpathquadraticcurveto{\pgfqpoint{4.824207in}{0.365648in}}{\pgfqpoint{4.824207in}{0.377067in}}%
\pgfpathquadraticcurveto{\pgfqpoint{4.824207in}{0.388487in}}{\pgfqpoint{4.819470in}{0.394989in}}%
\pgfpathquadraticcurveto{\pgfqpoint{4.814751in}{0.401492in}}{\pgfqpoint{4.806538in}{0.401492in}}%
\pgfpathquadraticcurveto{\pgfqpoint{4.798324in}{0.401492in}}{\pgfqpoint{4.793623in}{0.394989in}}%
\pgfpathquadraticcurveto{\pgfqpoint{4.788922in}{0.388487in}}{\pgfqpoint{4.788922in}{0.377067in}}%
\pgfpathlineto{\pgfqpoint{4.788922in}{0.377067in}}%
\pgfpathclose%
\pgfpathmoveto{\pgfqpoint{4.902560in}{0.355056in}}%
\pgfpathlineto{\pgfqpoint{4.902560in}{0.321626in}}%
\pgfpathlineto{\pgfqpoint{4.892149in}{0.321626in}}%
\pgfpathlineto{\pgfqpoint{4.892149in}{0.408643in}}%
\pgfpathlineto{\pgfqpoint{4.902560in}{0.408643in}}%
\pgfpathlineto{\pgfqpoint{4.902560in}{0.399078in}}%
\pgfpathquadraticcurveto{\pgfqpoint{4.905839in}{0.404698in}}{\pgfqpoint{4.910810in}{0.407418in}}%
\pgfpathquadraticcurveto{\pgfqpoint{4.915799in}{0.410156in}}{\pgfqpoint{4.922716in}{0.410156in}}%
\pgfpathquadraticcurveto{\pgfqpoint{4.934208in}{0.410156in}}{\pgfqpoint{4.941377in}{0.401041in}}%
\pgfpathquadraticcurveto{\pgfqpoint{4.948563in}{0.391927in}}{\pgfqpoint{4.948563in}{0.377067in}}%
\pgfpathquadraticcurveto{\pgfqpoint{4.948563in}{0.362207in}}{\pgfqpoint{4.941377in}{0.353075in}}%
\pgfpathquadraticcurveto{\pgfqpoint{4.934208in}{0.343961in}}{\pgfqpoint{4.922716in}{0.343961in}}%
\pgfpathquadraticcurveto{\pgfqpoint{4.915799in}{0.343961in}}{\pgfqpoint{4.910810in}{0.346699in}}%
\pgfpathquadraticcurveto{\pgfqpoint{4.905839in}{0.349437in}}{\pgfqpoint{4.902560in}{0.355056in}}%
\pgfpathlineto{\pgfqpoint{4.902560in}{0.355056in}}%
\pgfpathclose%
\pgfpathmoveto{\pgfqpoint{4.937810in}{0.377067in}}%
\pgfpathquadraticcurveto{\pgfqpoint{4.937810in}{0.388487in}}{\pgfqpoint{4.933109in}{0.394989in}}%
\pgfpathquadraticcurveto{\pgfqpoint{4.928408in}{0.401492in}}{\pgfqpoint{4.920194in}{0.401492in}}%
\pgfpathquadraticcurveto{\pgfqpoint{4.911963in}{0.401492in}}{\pgfqpoint{4.907262in}{0.394989in}}%
\pgfpathquadraticcurveto{\pgfqpoint{4.902560in}{0.388487in}}{\pgfqpoint{4.902560in}{0.377067in}}%
\pgfpathquadraticcurveto{\pgfqpoint{4.902560in}{0.365648in}}{\pgfqpoint{4.907262in}{0.359145in}}%
\pgfpathquadraticcurveto{\pgfqpoint{4.911963in}{0.352643in}}{\pgfqpoint{4.920194in}{0.352643in}}%
\pgfpathquadraticcurveto{\pgfqpoint{4.928408in}{0.352643in}}{\pgfqpoint{4.933109in}{0.359145in}}%
\pgfpathquadraticcurveto{\pgfqpoint{4.937810in}{0.365648in}}{\pgfqpoint{4.937810in}{0.377067in}}%
\pgfpathlineto{\pgfqpoint{4.937810in}{0.377067in}}%
\pgfpathclose%
\pgfpathmoveto{\pgfqpoint{4.964666in}{0.370475in}}%
\pgfpathlineto{\pgfqpoint{4.964666in}{0.408643in}}%
\pgfpathlineto{\pgfqpoint{4.975023in}{0.408643in}}%
\pgfpathlineto{\pgfqpoint{4.975023in}{0.370871in}}%
\pgfpathquadraticcurveto{\pgfqpoint{4.975023in}{0.361919in}}{\pgfqpoint{4.978500in}{0.357434in}}%
\pgfpathquadraticcurveto{\pgfqpoint{4.981994in}{0.352967in}}{\pgfqpoint{4.988983in}{0.352967in}}%
\pgfpathquadraticcurveto{\pgfqpoint{4.997358in}{0.352967in}}{\pgfqpoint{5.002222in}{0.358317in}}%
\pgfpathquadraticcurveto{\pgfqpoint{5.007103in}{0.363666in}}{\pgfqpoint{5.007103in}{0.372906in}}%
\pgfpathlineto{\pgfqpoint{5.007103in}{0.408643in}}%
\pgfpathlineto{\pgfqpoint{5.017460in}{0.408643in}}%
\pgfpathlineto{\pgfqpoint{5.017460in}{0.345600in}}%
\pgfpathlineto{\pgfqpoint{5.007103in}{0.345600in}}%
\pgfpathlineto{\pgfqpoint{5.007103in}{0.355291in}}%
\pgfpathquadraticcurveto{\pgfqpoint{5.003338in}{0.349545in}}{\pgfqpoint{4.998349in}{0.346753in}}%
\pgfpathquadraticcurveto{\pgfqpoint{4.993378in}{0.343961in}}{\pgfqpoint{4.986785in}{0.343961in}}%
\pgfpathquadraticcurveto{\pgfqpoint{4.975924in}{0.343961in}}{\pgfqpoint{4.970286in}{0.350715in}}%
\pgfpathquadraticcurveto{\pgfqpoint{4.964666in}{0.357470in}}{\pgfqpoint{4.964666in}{0.370475in}}%
\pgfpathlineto{\pgfqpoint{4.964666in}{0.370475in}}%
\pgfpathclose%
\pgfpathmoveto{\pgfqpoint{4.990730in}{0.410156in}}%
\pgfpathlineto{\pgfqpoint{4.990730in}{0.410156in}}%
\pgfpathlineto{\pgfqpoint{4.990730in}{0.410156in}}%
\pgfpathclose%
\pgfpathmoveto{\pgfqpoint{5.084051in}{0.377067in}}%
\pgfpathquadraticcurveto{\pgfqpoint{5.084051in}{0.388487in}}{\pgfqpoint{5.079350in}{0.394989in}}%
\pgfpathquadraticcurveto{\pgfqpoint{5.074648in}{0.401492in}}{\pgfqpoint{5.066435in}{0.401492in}}%
\pgfpathquadraticcurveto{\pgfqpoint{5.058203in}{0.401492in}}{\pgfqpoint{5.053502in}{0.394989in}}%
\pgfpathquadraticcurveto{\pgfqpoint{5.048801in}{0.388487in}}{\pgfqpoint{5.048801in}{0.377067in}}%
\pgfpathquadraticcurveto{\pgfqpoint{5.048801in}{0.365648in}}{\pgfqpoint{5.053502in}{0.359145in}}%
\pgfpathquadraticcurveto{\pgfqpoint{5.058203in}{0.352643in}}{\pgfqpoint{5.066435in}{0.352643in}}%
\pgfpathquadraticcurveto{\pgfqpoint{5.074648in}{0.352643in}}{\pgfqpoint{5.079350in}{0.359145in}}%
\pgfpathquadraticcurveto{\pgfqpoint{5.084051in}{0.365648in}}{\pgfqpoint{5.084051in}{0.377067in}}%
\pgfpathlineto{\pgfqpoint{5.084051in}{0.377067in}}%
\pgfpathclose%
\pgfpathmoveto{\pgfqpoint{5.048801in}{0.399078in}}%
\pgfpathquadraticcurveto{\pgfqpoint{5.052079in}{0.404698in}}{\pgfqpoint{5.057051in}{0.407418in}}%
\pgfpathquadraticcurveto{\pgfqpoint{5.062040in}{0.410156in}}{\pgfqpoint{5.068957in}{0.410156in}}%
\pgfpathquadraticcurveto{\pgfqpoint{5.080448in}{0.410156in}}{\pgfqpoint{5.087617in}{0.401041in}}%
\pgfpathquadraticcurveto{\pgfqpoint{5.094804in}{0.391927in}}{\pgfqpoint{5.094804in}{0.377067in}}%
\pgfpathquadraticcurveto{\pgfqpoint{5.094804in}{0.362207in}}{\pgfqpoint{5.087617in}{0.353075in}}%
\pgfpathquadraticcurveto{\pgfqpoint{5.080448in}{0.343961in}}{\pgfqpoint{5.068957in}{0.343961in}}%
\pgfpathquadraticcurveto{\pgfqpoint{5.062040in}{0.343961in}}{\pgfqpoint{5.057051in}{0.346699in}}%
\pgfpathquadraticcurveto{\pgfqpoint{5.052079in}{0.349437in}}{\pgfqpoint{5.048801in}{0.355056in}}%
\pgfpathlineto{\pgfqpoint{5.048801in}{0.345600in}}%
\pgfpathlineto{\pgfqpoint{5.038390in}{0.345600in}}%
\pgfpathlineto{\pgfqpoint{5.038390in}{0.433193in}}%
\pgfpathlineto{\pgfqpoint{5.048801in}{0.433193in}}%
\pgfpathlineto{\pgfqpoint{5.048801in}{0.399078in}}%
\pgfpathlineto{\pgfqpoint{5.048801in}{0.399078in}}%
\pgfpathclose%
\pgfpathmoveto{\pgfqpoint{5.111970in}{0.433193in}}%
\pgfpathlineto{\pgfqpoint{5.122327in}{0.433193in}}%
\pgfpathlineto{\pgfqpoint{5.122327in}{0.345600in}}%
\pgfpathlineto{\pgfqpoint{5.111970in}{0.345600in}}%
\pgfpathlineto{\pgfqpoint{5.111970in}{0.433193in}}%
\pgfpathlineto{\pgfqpoint{5.111970in}{0.433193in}}%
\pgfpathclose%
\pgfpathmoveto{\pgfqpoint{5.143995in}{0.408643in}}%
\pgfpathlineto{\pgfqpoint{5.154352in}{0.408643in}}%
\pgfpathlineto{\pgfqpoint{5.154352in}{0.345600in}}%
\pgfpathlineto{\pgfqpoint{5.143995in}{0.345600in}}%
\pgfpathlineto{\pgfqpoint{5.143995in}{0.408643in}}%
\pgfpathlineto{\pgfqpoint{5.143995in}{0.408643in}}%
\pgfpathclose%
\pgfpathmoveto{\pgfqpoint{5.143995in}{0.433193in}}%
\pgfpathlineto{\pgfqpoint{5.154352in}{0.433193in}}%
\pgfpathlineto{\pgfqpoint{5.154352in}{0.420062in}}%
\pgfpathlineto{\pgfqpoint{5.143995in}{0.420062in}}%
\pgfpathlineto{\pgfqpoint{5.143995in}{0.433193in}}%
\pgfpathlineto{\pgfqpoint{5.143995in}{0.433193in}}%
\pgfpathclose%
\pgfpathmoveto{\pgfqpoint{5.216206in}{0.406787in}}%
\pgfpathlineto{\pgfqpoint{5.216206in}{0.396989in}}%
\pgfpathquadraticcurveto{\pgfqpoint{5.211829in}{0.399240in}}{\pgfqpoint{5.207092in}{0.400357in}}%
\pgfpathquadraticcurveto{\pgfqpoint{5.202373in}{0.401492in}}{\pgfqpoint{5.197293in}{0.401492in}}%
\pgfpathquadraticcurveto{\pgfqpoint{5.189584in}{0.401492in}}{\pgfqpoint{5.185729in}{0.399132in}}%
\pgfpathquadraticcurveto{\pgfqpoint{5.181875in}{0.396773in}}{\pgfqpoint{5.181875in}{0.392035in}}%
\pgfpathquadraticcurveto{\pgfqpoint{5.181875in}{0.388433in}}{\pgfqpoint{5.184631in}{0.386380in}}%
\pgfpathquadraticcurveto{\pgfqpoint{5.187387in}{0.384326in}}{\pgfqpoint{5.195726in}{0.382471in}}%
\pgfpathlineto{\pgfqpoint{5.199275in}{0.381678in}}%
\pgfpathquadraticcurveto{\pgfqpoint{5.210298in}{0.379319in}}{\pgfqpoint{5.214945in}{0.375014in}}%
\pgfpathquadraticcurveto{\pgfqpoint{5.219592in}{0.370709in}}{\pgfqpoint{5.219592in}{0.363000in}}%
\pgfpathquadraticcurveto{\pgfqpoint{5.219592in}{0.354210in}}{\pgfqpoint{5.212640in}{0.349076in}}%
\pgfpathquadraticcurveto{\pgfqpoint{5.205687in}{0.343961in}}{\pgfqpoint{5.193529in}{0.343961in}}%
\pgfpathquadraticcurveto{\pgfqpoint{5.188467in}{0.343961in}}{\pgfqpoint{5.182974in}{0.344952in}}%
\pgfpathquadraticcurveto{\pgfqpoint{5.177480in}{0.345942in}}{\pgfqpoint{5.171410in}{0.347906in}}%
\pgfpathlineto{\pgfqpoint{5.171410in}{0.358605in}}%
\pgfpathquadraticcurveto{\pgfqpoint{5.177156in}{0.355615in}}{\pgfqpoint{5.182721in}{0.354120in}}%
\pgfpathquadraticcurveto{\pgfqpoint{5.188287in}{0.352643in}}{\pgfqpoint{5.193763in}{0.352643in}}%
\pgfpathquadraticcurveto{\pgfqpoint{5.201076in}{0.352643in}}{\pgfqpoint{5.205002in}{0.355146in}}%
\pgfpathquadraticcurveto{\pgfqpoint{5.208947in}{0.357650in}}{\pgfqpoint{5.208947in}{0.362207in}}%
\pgfpathquadraticcurveto{\pgfqpoint{5.208947in}{0.366422in}}{\pgfqpoint{5.206101in}{0.368674in}}%
\pgfpathquadraticcurveto{\pgfqpoint{5.203273in}{0.370925in}}{\pgfqpoint{5.193637in}{0.373014in}}%
\pgfpathlineto{\pgfqpoint{5.190034in}{0.373861in}}%
\pgfpathquadraticcurveto{\pgfqpoint{5.180416in}{0.375878in}}{\pgfqpoint{5.176129in}{0.380075in}}%
\pgfpathquadraticcurveto{\pgfqpoint{5.171860in}{0.384272in}}{\pgfqpoint{5.171860in}{0.391585in}}%
\pgfpathquadraticcurveto{\pgfqpoint{5.171860in}{0.400483in}}{\pgfqpoint{5.178164in}{0.405310in}}%
\pgfpathquadraticcurveto{\pgfqpoint{5.184469in}{0.410156in}}{\pgfqpoint{5.196068in}{0.410156in}}%
\pgfpathquadraticcurveto{\pgfqpoint{5.201796in}{0.410156in}}{\pgfqpoint{5.206858in}{0.409309in}}%
\pgfpathquadraticcurveto{\pgfqpoint{5.211937in}{0.408480in}}{\pgfqpoint{5.216206in}{0.406787in}}%
\pgfpathlineto{\pgfqpoint{5.216206in}{0.406787in}}%
\pgfpathclose%
\pgfpathmoveto{\pgfqpoint{5.288489in}{0.383660in}}%
\pgfpathlineto{\pgfqpoint{5.288489in}{0.345600in}}%
\pgfpathlineto{\pgfqpoint{5.278132in}{0.345600in}}%
\pgfpathlineto{\pgfqpoint{5.278132in}{0.383317in}}%
\pgfpathquadraticcurveto{\pgfqpoint{5.278132in}{0.392269in}}{\pgfqpoint{5.274637in}{0.396700in}}%
\pgfpathquadraticcurveto{\pgfqpoint{5.271143in}{0.401149in}}{\pgfqpoint{5.264172in}{0.401149in}}%
\pgfpathquadraticcurveto{\pgfqpoint{5.255779in}{0.401149in}}{\pgfqpoint{5.250933in}{0.395800in}}%
\pgfpathquadraticcurveto{\pgfqpoint{5.246088in}{0.390468in}}{\pgfqpoint{5.246088in}{0.381228in}}%
\pgfpathlineto{\pgfqpoint{5.246088in}{0.345600in}}%
\pgfpathlineto{\pgfqpoint{5.235677in}{0.345600in}}%
\pgfpathlineto{\pgfqpoint{5.235677in}{0.433193in}}%
\pgfpathlineto{\pgfqpoint{5.246088in}{0.433193in}}%
\pgfpathlineto{\pgfqpoint{5.246088in}{0.398844in}}%
\pgfpathquadraticcurveto{\pgfqpoint{5.249817in}{0.404536in}}{\pgfqpoint{5.254842in}{0.407346in}}%
\pgfpathquadraticcurveto{\pgfqpoint{5.259885in}{0.410156in}}{\pgfqpoint{5.266478in}{0.410156in}}%
\pgfpathquadraticcurveto{\pgfqpoint{5.277339in}{0.410156in}}{\pgfqpoint{5.282905in}{0.403437in}}%
\pgfpathquadraticcurveto{\pgfqpoint{5.288489in}{0.396718in}}{\pgfqpoint{5.288489in}{0.383660in}}%
\pgfpathlineto{\pgfqpoint{5.288489in}{0.383660in}}%
\pgfpathclose%
\pgfpathmoveto{\pgfqpoint{5.363059in}{0.379715in}}%
\pgfpathlineto{\pgfqpoint{5.363059in}{0.374654in}}%
\pgfpathlineto{\pgfqpoint{5.315435in}{0.374654in}}%
\pgfpathquadraticcurveto{\pgfqpoint{5.316119in}{0.363954in}}{\pgfqpoint{5.321883in}{0.358353in}}%
\pgfpathquadraticcurveto{\pgfqpoint{5.327647in}{0.352751in}}{\pgfqpoint{5.337950in}{0.352751in}}%
\pgfpathquadraticcurveto{\pgfqpoint{5.343912in}{0.352751in}}{\pgfqpoint{5.349514in}{0.354210in}}%
\pgfpathquadraticcurveto{\pgfqpoint{5.355116in}{0.355669in}}{\pgfqpoint{5.360645in}{0.358605in}}%
\pgfpathlineto{\pgfqpoint{5.360645in}{0.348806in}}%
\pgfpathquadraticcurveto{\pgfqpoint{5.355062in}{0.346447in}}{\pgfqpoint{5.349208in}{0.345204in}}%
\pgfpathquadraticcurveto{\pgfqpoint{5.343354in}{0.343961in}}{\pgfqpoint{5.337338in}{0.343961in}}%
\pgfpathquadraticcurveto{\pgfqpoint{5.322243in}{0.343961in}}{\pgfqpoint{5.313436in}{0.352733in}}%
\pgfpathquadraticcurveto{\pgfqpoint{5.304628in}{0.361523in}}{\pgfqpoint{5.304628in}{0.376509in}}%
\pgfpathquadraticcurveto{\pgfqpoint{5.304628in}{0.391981in}}{\pgfqpoint{5.312985in}{0.401059in}}%
\pgfpathquadraticcurveto{\pgfqpoint{5.321343in}{0.410156in}}{\pgfqpoint{5.335536in}{0.410156in}}%
\pgfpathquadraticcurveto{\pgfqpoint{5.348253in}{0.410156in}}{\pgfqpoint{5.355656in}{0.401960in}}%
\pgfpathquadraticcurveto{\pgfqpoint{5.363059in}{0.393783in}}{\pgfqpoint{5.363059in}{0.379715in}}%
\pgfpathlineto{\pgfqpoint{5.363059in}{0.379715in}}%
\pgfpathclose%
\pgfpathmoveto{\pgfqpoint{5.352702in}{0.382759in}}%
\pgfpathquadraticcurveto{\pgfqpoint{5.352594in}{0.391243in}}{\pgfqpoint{5.347947in}{0.396304in}}%
\pgfpathquadraticcurveto{\pgfqpoint{5.343300in}{0.401384in}}{\pgfqpoint{5.335645in}{0.401384in}}%
\pgfpathquadraticcurveto{\pgfqpoint{5.326981in}{0.401384in}}{\pgfqpoint{5.321775in}{0.396484in}}%
\pgfpathquadraticcurveto{\pgfqpoint{5.316570in}{0.391585in}}{\pgfqpoint{5.315777in}{0.382687in}}%
\pgfpathlineto{\pgfqpoint{5.352702in}{0.382759in}}%
\pgfpathlineto{\pgfqpoint{5.352702in}{0.382759in}}%
\pgfpathclose%
\pgfpathmoveto{\pgfqpoint{5.421544in}{0.399078in}}%
\pgfpathlineto{\pgfqpoint{5.421544in}{0.433193in}}%
\pgfpathlineto{\pgfqpoint{5.431901in}{0.433193in}}%
\pgfpathlineto{\pgfqpoint{5.431901in}{0.345600in}}%
\pgfpathlineto{\pgfqpoint{5.421544in}{0.345600in}}%
\pgfpathlineto{\pgfqpoint{5.421544in}{0.355056in}}%
\pgfpathquadraticcurveto{\pgfqpoint{5.418284in}{0.349437in}}{\pgfqpoint{5.413295in}{0.346699in}}%
\pgfpathquadraticcurveto{\pgfqpoint{5.408324in}{0.343961in}}{\pgfqpoint{5.401335in}{0.343961in}}%
\pgfpathquadraticcurveto{\pgfqpoint{5.389915in}{0.343961in}}{\pgfqpoint{5.382728in}{0.353075in}}%
\pgfpathquadraticcurveto{\pgfqpoint{5.375559in}{0.362207in}}{\pgfqpoint{5.375559in}{0.377067in}}%
\pgfpathquadraticcurveto{\pgfqpoint{5.375559in}{0.391927in}}{\pgfqpoint{5.382728in}{0.401041in}}%
\pgfpathquadraticcurveto{\pgfqpoint{5.389915in}{0.410156in}}{\pgfqpoint{5.401335in}{0.410156in}}%
\pgfpathquadraticcurveto{\pgfqpoint{5.408324in}{0.410156in}}{\pgfqpoint{5.413295in}{0.407418in}}%
\pgfpathquadraticcurveto{\pgfqpoint{5.418284in}{0.404698in}}{\pgfqpoint{5.421544in}{0.399078in}}%
\pgfpathlineto{\pgfqpoint{5.421544in}{0.399078in}}%
\pgfpathclose%
\pgfpathmoveto{\pgfqpoint{5.386259in}{0.377067in}}%
\pgfpathquadraticcurveto{\pgfqpoint{5.386259in}{0.365648in}}{\pgfqpoint{5.390960in}{0.359145in}}%
\pgfpathquadraticcurveto{\pgfqpoint{5.395661in}{0.352643in}}{\pgfqpoint{5.403875in}{0.352643in}}%
\pgfpathquadraticcurveto{\pgfqpoint{5.412088in}{0.352643in}}{\pgfqpoint{5.416807in}{0.359145in}}%
\pgfpathquadraticcurveto{\pgfqpoint{5.421544in}{0.365648in}}{\pgfqpoint{5.421544in}{0.377067in}}%
\pgfpathquadraticcurveto{\pgfqpoint{5.421544in}{0.388487in}}{\pgfqpoint{5.416807in}{0.394989in}}%
\pgfpathquadraticcurveto{\pgfqpoint{5.412088in}{0.401492in}}{\pgfqpoint{5.403875in}{0.401492in}}%
\pgfpathquadraticcurveto{\pgfqpoint{5.395661in}{0.401492in}}{\pgfqpoint{5.390960in}{0.394989in}}%
\pgfpathquadraticcurveto{\pgfqpoint{5.386259in}{0.388487in}}{\pgfqpoint{5.386259in}{0.377067in}}%
\pgfpathlineto{\pgfqpoint{5.386259in}{0.377067in}}%
\pgfpathclose%
\pgfpathmoveto{\pgfqpoint{5.540821in}{0.408643in}}%
\pgfpathlineto{\pgfqpoint{5.540821in}{0.345600in}}%
\pgfpathlineto{\pgfqpoint{5.530410in}{0.345600in}}%
\pgfpathlineto{\pgfqpoint{5.530410in}{0.400591in}}%
\pgfpathlineto{\pgfqpoint{5.501987in}{0.400591in}}%
\pgfpathlineto{\pgfqpoint{5.501987in}{0.345600in}}%
\pgfpathlineto{\pgfqpoint{5.491576in}{0.345600in}}%
\pgfpathlineto{\pgfqpoint{5.491576in}{0.400591in}}%
\pgfpathlineto{\pgfqpoint{5.481669in}{0.400591in}}%
\pgfpathlineto{\pgfqpoint{5.481669in}{0.408643in}}%
\pgfpathlineto{\pgfqpoint{5.491576in}{0.408643in}}%
\pgfpathlineto{\pgfqpoint{5.491576in}{0.413038in}}%
\pgfpathquadraticcurveto{\pgfqpoint{5.491576in}{0.423340in}}{\pgfqpoint{5.496439in}{0.428258in}}%
\pgfpathquadraticcurveto{\pgfqpoint{5.501320in}{0.433193in}}{\pgfqpoint{5.511389in}{0.433193in}}%
\pgfpathlineto{\pgfqpoint{5.521800in}{0.433193in}}%
\pgfpathlineto{\pgfqpoint{5.521800in}{0.424565in}}%
\pgfpathlineto{\pgfqpoint{5.511893in}{0.424565in}}%
\pgfpathquadraticcurveto{\pgfqpoint{5.506328in}{0.424565in}}{\pgfqpoint{5.504148in}{0.422314in}}%
\pgfpathquadraticcurveto{\pgfqpoint{5.501987in}{0.420062in}}{\pgfqpoint{5.501987in}{0.414208in}}%
\pgfpathlineto{\pgfqpoint{5.501987in}{0.408643in}}%
\pgfpathlineto{\pgfqpoint{5.540821in}{0.408643in}}%
\pgfpathlineto{\pgfqpoint{5.540821in}{0.408643in}}%
\pgfpathclose%
\pgfpathmoveto{\pgfqpoint{5.530410in}{0.433067in}}%
\pgfpathlineto{\pgfqpoint{5.540821in}{0.433067in}}%
\pgfpathlineto{\pgfqpoint{5.540821in}{0.419954in}}%
\pgfpathlineto{\pgfqpoint{5.530410in}{0.419954in}}%
\pgfpathlineto{\pgfqpoint{5.530410in}{0.433067in}}%
\pgfpathlineto{\pgfqpoint{5.530410in}{0.433067in}}%
\pgfpathclose%
\pgfpathmoveto{\pgfqpoint{5.572738in}{0.426547in}}%
\pgfpathlineto{\pgfqpoint{5.572738in}{0.408643in}}%
\pgfpathlineto{\pgfqpoint{5.594065in}{0.408643in}}%
\pgfpathlineto{\pgfqpoint{5.594065in}{0.400591in}}%
\pgfpathlineto{\pgfqpoint{5.572738in}{0.400591in}}%
\pgfpathlineto{\pgfqpoint{5.572738in}{0.366368in}}%
\pgfpathquadraticcurveto{\pgfqpoint{5.572738in}{0.358659in}}{\pgfqpoint{5.574846in}{0.356461in}}%
\pgfpathquadraticcurveto{\pgfqpoint{5.576953in}{0.354264in}}{\pgfqpoint{5.583438in}{0.354264in}}%
\pgfpathlineto{\pgfqpoint{5.594065in}{0.354264in}}%
\pgfpathlineto{\pgfqpoint{5.594065in}{0.345600in}}%
\pgfpathlineto{\pgfqpoint{5.583438in}{0.345600in}}%
\pgfpathquadraticcurveto{\pgfqpoint{5.571442in}{0.345600in}}{\pgfqpoint{5.566885in}{0.350067in}}%
\pgfpathquadraticcurveto{\pgfqpoint{5.562327in}{0.354552in}}{\pgfqpoint{5.562327in}{0.366368in}}%
\pgfpathlineto{\pgfqpoint{5.562327in}{0.400591in}}%
\pgfpathlineto{\pgfqpoint{5.554726in}{0.400591in}}%
\pgfpathlineto{\pgfqpoint{5.554726in}{0.408643in}}%
\pgfpathlineto{\pgfqpoint{5.562327in}{0.408643in}}%
\pgfpathlineto{\pgfqpoint{5.562327in}{0.426547in}}%
\pgfpathlineto{\pgfqpoint{5.572738in}{0.426547in}}%
\pgfpathlineto{\pgfqpoint{5.572738in}{0.426547in}}%
\pgfpathclose%
\pgfpathmoveto{\pgfqpoint{5.644319in}{0.408643in}}%
\pgfpathlineto{\pgfqpoint{5.654676in}{0.408643in}}%
\pgfpathlineto{\pgfqpoint{5.654676in}{0.345600in}}%
\pgfpathlineto{\pgfqpoint{5.644319in}{0.345600in}}%
\pgfpathlineto{\pgfqpoint{5.644319in}{0.408643in}}%
\pgfpathlineto{\pgfqpoint{5.644319in}{0.408643in}}%
\pgfpathclose%
\pgfpathmoveto{\pgfqpoint{5.644319in}{0.433193in}}%
\pgfpathlineto{\pgfqpoint{5.654676in}{0.433193in}}%
\pgfpathlineto{\pgfqpoint{5.654676in}{0.420062in}}%
\pgfpathlineto{\pgfqpoint{5.644319in}{0.420062in}}%
\pgfpathlineto{\pgfqpoint{5.644319in}{0.433193in}}%
\pgfpathlineto{\pgfqpoint{5.644319in}{0.433193in}}%
\pgfpathclose%
\pgfpathmoveto{\pgfqpoint{5.716530in}{0.406787in}}%
\pgfpathlineto{\pgfqpoint{5.716530in}{0.396989in}}%
\pgfpathquadraticcurveto{\pgfqpoint{5.712153in}{0.399240in}}{\pgfqpoint{5.707415in}{0.400357in}}%
\pgfpathquadraticcurveto{\pgfqpoint{5.702696in}{0.401492in}}{\pgfqpoint{5.697617in}{0.401492in}}%
\pgfpathquadraticcurveto{\pgfqpoint{5.689908in}{0.401492in}}{\pgfqpoint{5.686053in}{0.399132in}}%
\pgfpathquadraticcurveto{\pgfqpoint{5.682198in}{0.396773in}}{\pgfqpoint{5.682198in}{0.392035in}}%
\pgfpathquadraticcurveto{\pgfqpoint{5.682198in}{0.388433in}}{\pgfqpoint{5.684954in}{0.386380in}}%
\pgfpathquadraticcurveto{\pgfqpoint{5.687710in}{0.384326in}}{\pgfqpoint{5.696050in}{0.382471in}}%
\pgfpathlineto{\pgfqpoint{5.699598in}{0.381678in}}%
\pgfpathquadraticcurveto{\pgfqpoint{5.710622in}{0.379319in}}{\pgfqpoint{5.715269in}{0.375014in}}%
\pgfpathquadraticcurveto{\pgfqpoint{5.719916in}{0.370709in}}{\pgfqpoint{5.719916in}{0.363000in}}%
\pgfpathquadraticcurveto{\pgfqpoint{5.719916in}{0.354210in}}{\pgfqpoint{5.712963in}{0.349076in}}%
\pgfpathquadraticcurveto{\pgfqpoint{5.706010in}{0.343961in}}{\pgfqpoint{5.693852in}{0.343961in}}%
\pgfpathquadraticcurveto{\pgfqpoint{5.688791in}{0.343961in}}{\pgfqpoint{5.683297in}{0.344952in}}%
\pgfpathquadraticcurveto{\pgfqpoint{5.677803in}{0.345942in}}{\pgfqpoint{5.671733in}{0.347906in}}%
\pgfpathlineto{\pgfqpoint{5.671733in}{0.358605in}}%
\pgfpathquadraticcurveto{\pgfqpoint{5.677479in}{0.355615in}}{\pgfqpoint{5.683045in}{0.354120in}}%
\pgfpathquadraticcurveto{\pgfqpoint{5.688611in}{0.352643in}}{\pgfqpoint{5.694086in}{0.352643in}}%
\pgfpathquadraticcurveto{\pgfqpoint{5.701399in}{0.352643in}}{\pgfqpoint{5.705326in}{0.355146in}}%
\pgfpathquadraticcurveto{\pgfqpoint{5.709271in}{0.357650in}}{\pgfqpoint{5.709271in}{0.362207in}}%
\pgfpathquadraticcurveto{\pgfqpoint{5.709271in}{0.366422in}}{\pgfqpoint{5.706425in}{0.368674in}}%
\pgfpathquadraticcurveto{\pgfqpoint{5.703597in}{0.370925in}}{\pgfqpoint{5.693960in}{0.373014in}}%
\pgfpathlineto{\pgfqpoint{5.690358in}{0.373861in}}%
\pgfpathquadraticcurveto{\pgfqpoint{5.680739in}{0.375878in}}{\pgfqpoint{5.676452in}{0.380075in}}%
\pgfpathquadraticcurveto{\pgfqpoint{5.672184in}{0.384272in}}{\pgfqpoint{5.672184in}{0.391585in}}%
\pgfpathquadraticcurveto{\pgfqpoint{5.672184in}{0.400483in}}{\pgfqpoint{5.678488in}{0.405310in}}%
\pgfpathquadraticcurveto{\pgfqpoint{5.684792in}{0.410156in}}{\pgfqpoint{5.696392in}{0.410156in}}%
\pgfpathquadraticcurveto{\pgfqpoint{5.702120in}{0.410156in}}{\pgfqpoint{5.707181in}{0.409309in}}%
\pgfpathquadraticcurveto{\pgfqpoint{5.712261in}{0.408480in}}{\pgfqpoint{5.716530in}{0.406787in}}%
\pgfpathlineto{\pgfqpoint{5.716530in}{0.406787in}}%
\pgfpathclose%
\pgfpathmoveto{\pgfqpoint{5.809562in}{0.398970in}}%
\pgfpathquadraticcurveto{\pgfqpoint{5.807815in}{0.399979in}}{\pgfqpoint{5.805762in}{0.400447in}}%
\pgfpathquadraticcurveto{\pgfqpoint{5.803708in}{0.400933in}}{\pgfqpoint{5.801241in}{0.400933in}}%
\pgfpathquadraticcurveto{\pgfqpoint{5.792451in}{0.400933in}}{\pgfqpoint{5.787750in}{0.395223in}}%
\pgfpathquadraticcurveto{\pgfqpoint{5.783048in}{0.389514in}}{\pgfqpoint{5.783048in}{0.378814in}}%
\pgfpathlineto{\pgfqpoint{5.783048in}{0.345600in}}%
\pgfpathlineto{\pgfqpoint{5.772637in}{0.345600in}}%
\pgfpathlineto{\pgfqpoint{5.772637in}{0.408643in}}%
\pgfpathlineto{\pgfqpoint{5.783048in}{0.408643in}}%
\pgfpathlineto{\pgfqpoint{5.783048in}{0.398844in}}%
\pgfpathquadraticcurveto{\pgfqpoint{5.786327in}{0.404590in}}{\pgfqpoint{5.791550in}{0.407364in}}%
\pgfpathquadraticcurveto{\pgfqpoint{5.796792in}{0.410156in}}{\pgfqpoint{5.804285in}{0.410156in}}%
\pgfpathquadraticcurveto{\pgfqpoint{5.805347in}{0.410156in}}{\pgfqpoint{5.806644in}{0.410011in}}%
\pgfpathquadraticcurveto{\pgfqpoint{5.807941in}{0.409885in}}{\pgfqpoint{5.809508in}{0.409597in}}%
\pgfpathlineto{\pgfqpoint{5.809562in}{0.398970in}}%
\pgfpathlineto{\pgfqpoint{5.809562in}{0.398970in}}%
\pgfpathclose%
\pgfpathmoveto{\pgfqpoint{5.871812in}{0.379715in}}%
\pgfpathlineto{\pgfqpoint{5.871812in}{0.374654in}}%
\pgfpathlineto{\pgfqpoint{5.824188in}{0.374654in}}%
\pgfpathquadraticcurveto{\pgfqpoint{5.824873in}{0.363954in}}{\pgfqpoint{5.830637in}{0.358353in}}%
\pgfpathquadraticcurveto{\pgfqpoint{5.836400in}{0.352751in}}{\pgfqpoint{5.846703in}{0.352751in}}%
\pgfpathquadraticcurveto{\pgfqpoint{5.852665in}{0.352751in}}{\pgfqpoint{5.858267in}{0.354210in}}%
\pgfpathquadraticcurveto{\pgfqpoint{5.863869in}{0.355669in}}{\pgfqpoint{5.869399in}{0.358605in}}%
\pgfpathlineto{\pgfqpoint{5.869399in}{0.348806in}}%
\pgfpathquadraticcurveto{\pgfqpoint{5.863815in}{0.346447in}}{\pgfqpoint{5.857961in}{0.345204in}}%
\pgfpathquadraticcurveto{\pgfqpoint{5.852107in}{0.343961in}}{\pgfqpoint{5.846091in}{0.343961in}}%
\pgfpathquadraticcurveto{\pgfqpoint{5.830997in}{0.343961in}}{\pgfqpoint{5.822189in}{0.352733in}}%
\pgfpathquadraticcurveto{\pgfqpoint{5.813381in}{0.361523in}}{\pgfqpoint{5.813381in}{0.376509in}}%
\pgfpathquadraticcurveto{\pgfqpoint{5.813381in}{0.391981in}}{\pgfqpoint{5.821738in}{0.401059in}}%
\pgfpathquadraticcurveto{\pgfqpoint{5.830096in}{0.410156in}}{\pgfqpoint{5.844290in}{0.410156in}}%
\pgfpathquadraticcurveto{\pgfqpoint{5.857006in}{0.410156in}}{\pgfqpoint{5.864409in}{0.401960in}}%
\pgfpathquadraticcurveto{\pgfqpoint{5.871812in}{0.393783in}}{\pgfqpoint{5.871812in}{0.379715in}}%
\pgfpathlineto{\pgfqpoint{5.871812in}{0.379715in}}%
\pgfpathclose%
\pgfpathmoveto{\pgfqpoint{5.861455in}{0.382759in}}%
\pgfpathquadraticcurveto{\pgfqpoint{5.861347in}{0.391243in}}{\pgfqpoint{5.856700in}{0.396304in}}%
\pgfpathquadraticcurveto{\pgfqpoint{5.852053in}{0.401384in}}{\pgfqpoint{5.844398in}{0.401384in}}%
\pgfpathquadraticcurveto{\pgfqpoint{5.835734in}{0.401384in}}{\pgfqpoint{5.830528in}{0.396484in}}%
\pgfpathquadraticcurveto{\pgfqpoint{5.825323in}{0.391585in}}{\pgfqpoint{5.824530in}{0.382687in}}%
\pgfpathlineto{\pgfqpoint{5.861455in}{0.382759in}}%
\pgfpathlineto{\pgfqpoint{5.861455in}{0.382759in}}%
\pgfpathclose%
\pgfpathmoveto{\pgfqpoint{5.899065in}{0.426547in}}%
\pgfpathlineto{\pgfqpoint{5.899065in}{0.408643in}}%
\pgfpathlineto{\pgfqpoint{5.920391in}{0.408643in}}%
\pgfpathlineto{\pgfqpoint{5.920391in}{0.400591in}}%
\pgfpathlineto{\pgfqpoint{5.899065in}{0.400591in}}%
\pgfpathlineto{\pgfqpoint{5.899065in}{0.366368in}}%
\pgfpathquadraticcurveto{\pgfqpoint{5.899065in}{0.358659in}}{\pgfqpoint{5.901172in}{0.356461in}}%
\pgfpathquadraticcurveto{\pgfqpoint{5.903280in}{0.354264in}}{\pgfqpoint{5.909764in}{0.354264in}}%
\pgfpathlineto{\pgfqpoint{5.920391in}{0.354264in}}%
\pgfpathlineto{\pgfqpoint{5.920391in}{0.345600in}}%
\pgfpathlineto{\pgfqpoint{5.909764in}{0.345600in}}%
\pgfpathquadraticcurveto{\pgfqpoint{5.897768in}{0.345600in}}{\pgfqpoint{5.893211in}{0.350067in}}%
\pgfpathquadraticcurveto{\pgfqpoint{5.888654in}{0.354552in}}{\pgfqpoint{5.888654in}{0.366368in}}%
\pgfpathlineto{\pgfqpoint{5.888654in}{0.400591in}}%
\pgfpathlineto{\pgfqpoint{5.881053in}{0.400591in}}%
\pgfpathlineto{\pgfqpoint{5.881053in}{0.408643in}}%
\pgfpathlineto{\pgfqpoint{5.888654in}{0.408643in}}%
\pgfpathlineto{\pgfqpoint{5.888654in}{0.426547in}}%
\pgfpathlineto{\pgfqpoint{5.899065in}{0.426547in}}%
\pgfpathlineto{\pgfqpoint{5.899065in}{0.426547in}}%
\pgfpathclose%
\pgfpathmoveto{\pgfqpoint{5.962666in}{0.377283in}}%
\pgfpathquadraticcurveto{\pgfqpoint{5.950111in}{0.377283in}}{\pgfqpoint{5.945266in}{0.374419in}}%
\pgfpathquadraticcurveto{\pgfqpoint{5.940421in}{0.371556in}}{\pgfqpoint{5.940421in}{0.364621in}}%
\pgfpathquadraticcurveto{\pgfqpoint{5.940421in}{0.359109in}}{\pgfqpoint{5.944059in}{0.355867in}}%
\pgfpathquadraticcurveto{\pgfqpoint{5.947697in}{0.352643in}}{\pgfqpoint{5.953930in}{0.352643in}}%
\pgfpathquadraticcurveto{\pgfqpoint{5.962558in}{0.352643in}}{\pgfqpoint{5.967763in}{0.358749in}}%
\pgfpathquadraticcurveto{\pgfqpoint{5.972969in}{0.364855in}}{\pgfqpoint{5.972969in}{0.374978in}}%
\pgfpathlineto{\pgfqpoint{5.972969in}{0.377283in}}%
\pgfpathlineto{\pgfqpoint{5.962666in}{0.377283in}}%
\pgfpathlineto{\pgfqpoint{5.962666in}{0.377283in}}%
\pgfpathclose%
\pgfpathmoveto{\pgfqpoint{5.983326in}{0.381570in}}%
\pgfpathlineto{\pgfqpoint{5.983326in}{0.345600in}}%
\pgfpathlineto{\pgfqpoint{5.972969in}{0.345600in}}%
\pgfpathlineto{\pgfqpoint{5.972969in}{0.355164in}}%
\pgfpathquadraticcurveto{\pgfqpoint{5.969420in}{0.349437in}}{\pgfqpoint{5.964125in}{0.346699in}}%
\pgfpathquadraticcurveto{\pgfqpoint{5.958829in}{0.343961in}}{\pgfqpoint{5.951174in}{0.343961in}}%
\pgfpathquadraticcurveto{\pgfqpoint{5.941501in}{0.343961in}}{\pgfqpoint{5.935773in}{0.349401in}}%
\pgfpathquadraticcurveto{\pgfqpoint{5.930064in}{0.354840in}}{\pgfqpoint{5.930064in}{0.363954in}}%
\pgfpathquadraticcurveto{\pgfqpoint{5.930064in}{0.374582in}}{\pgfqpoint{5.937178in}{0.379985in}}%
\pgfpathquadraticcurveto{\pgfqpoint{5.944311in}{0.385389in}}{\pgfqpoint{5.958433in}{0.385389in}}%
\pgfpathlineto{\pgfqpoint{5.972969in}{0.385389in}}%
\pgfpathlineto{\pgfqpoint{5.972969in}{0.386416in}}%
\pgfpathquadraticcurveto{\pgfqpoint{5.972969in}{0.393566in}}{\pgfqpoint{5.968267in}{0.397475in}}%
\pgfpathquadraticcurveto{\pgfqpoint{5.963566in}{0.401384in}}{\pgfqpoint{5.955064in}{0.401384in}}%
\pgfpathquadraticcurveto{\pgfqpoint{5.949661in}{0.401384in}}{\pgfqpoint{5.944527in}{0.400087in}}%
\pgfpathquadraticcurveto{\pgfqpoint{5.939412in}{0.398790in}}{\pgfqpoint{5.934693in}{0.396196in}}%
\pgfpathlineto{\pgfqpoint{5.934693in}{0.405779in}}%
\pgfpathquadraticcurveto{\pgfqpoint{5.940367in}{0.407976in}}{\pgfqpoint{5.945716in}{0.409057in}}%
\pgfpathquadraticcurveto{\pgfqpoint{5.951066in}{0.410156in}}{\pgfqpoint{5.956127in}{0.410156in}}%
\pgfpathquadraticcurveto{\pgfqpoint{5.969816in}{0.410156in}}{\pgfqpoint{5.976571in}{0.403059in}}%
\pgfpathquadraticcurveto{\pgfqpoint{5.983326in}{0.395980in}}{\pgfqpoint{5.983326in}{0.381570in}}%
\pgfpathlineto{\pgfqpoint{5.983326in}{0.381570in}}%
\pgfpathclose%
\pgfpathmoveto{\pgfqpoint{6.004652in}{0.408643in}}%
\pgfpathlineto{\pgfqpoint{6.015009in}{0.408643in}}%
\pgfpathlineto{\pgfqpoint{6.015009in}{0.345600in}}%
\pgfpathlineto{\pgfqpoint{6.004652in}{0.345600in}}%
\pgfpathlineto{\pgfqpoint{6.004652in}{0.408643in}}%
\pgfpathlineto{\pgfqpoint{6.004652in}{0.408643in}}%
\pgfpathclose%
\pgfpathmoveto{\pgfqpoint{6.004652in}{0.433193in}}%
\pgfpathlineto{\pgfqpoint{6.015009in}{0.433193in}}%
\pgfpathlineto{\pgfqpoint{6.015009in}{0.420062in}}%
\pgfpathlineto{\pgfqpoint{6.004652in}{0.420062in}}%
\pgfpathlineto{\pgfqpoint{6.004652in}{0.433193in}}%
\pgfpathlineto{\pgfqpoint{6.004652in}{0.433193in}}%
\pgfpathclose%
\pgfpathmoveto{\pgfqpoint{6.089093in}{0.383660in}}%
\pgfpathlineto{\pgfqpoint{6.089093in}{0.345600in}}%
\pgfpathlineto{\pgfqpoint{6.078736in}{0.345600in}}%
\pgfpathlineto{\pgfqpoint{6.078736in}{0.383317in}}%
\pgfpathquadraticcurveto{\pgfqpoint{6.078736in}{0.392269in}}{\pgfqpoint{6.075242in}{0.396700in}}%
\pgfpathquadraticcurveto{\pgfqpoint{6.071747in}{0.401149in}}{\pgfqpoint{6.064776in}{0.401149in}}%
\pgfpathquadraticcurveto{\pgfqpoint{6.056383in}{0.401149in}}{\pgfqpoint{6.051538in}{0.395800in}}%
\pgfpathquadraticcurveto{\pgfqpoint{6.046692in}{0.390468in}}{\pgfqpoint{6.046692in}{0.381228in}}%
\pgfpathlineto{\pgfqpoint{6.046692in}{0.345600in}}%
\pgfpathlineto{\pgfqpoint{6.036281in}{0.345600in}}%
\pgfpathlineto{\pgfqpoint{6.036281in}{0.408643in}}%
\pgfpathlineto{\pgfqpoint{6.046692in}{0.408643in}}%
\pgfpathlineto{\pgfqpoint{6.046692in}{0.398844in}}%
\pgfpathquadraticcurveto{\pgfqpoint{6.050421in}{0.404536in}}{\pgfqpoint{6.055446in}{0.407346in}}%
\pgfpathquadraticcurveto{\pgfqpoint{6.060490in}{0.410156in}}{\pgfqpoint{6.067082in}{0.410156in}}%
\pgfpathquadraticcurveto{\pgfqpoint{6.077943in}{0.410156in}}{\pgfqpoint{6.083509in}{0.403437in}}%
\pgfpathquadraticcurveto{\pgfqpoint{6.089093in}{0.396718in}}{\pgfqpoint{6.089093in}{0.383660in}}%
\pgfpathlineto{\pgfqpoint{6.089093in}{0.383660in}}%
\pgfpathclose%
\pgfpathmoveto{\pgfqpoint{6.163663in}{0.379715in}}%
\pgfpathlineto{\pgfqpoint{6.163663in}{0.374654in}}%
\pgfpathlineto{\pgfqpoint{6.116039in}{0.374654in}}%
\pgfpathquadraticcurveto{\pgfqpoint{6.116724in}{0.363954in}}{\pgfqpoint{6.122487in}{0.358353in}}%
\pgfpathquadraticcurveto{\pgfqpoint{6.128251in}{0.352751in}}{\pgfqpoint{6.138554in}{0.352751in}}%
\pgfpathquadraticcurveto{\pgfqpoint{6.144516in}{0.352751in}}{\pgfqpoint{6.150118in}{0.354210in}}%
\pgfpathquadraticcurveto{\pgfqpoint{6.155720in}{0.355669in}}{\pgfqpoint{6.161250in}{0.358605in}}%
\pgfpathlineto{\pgfqpoint{6.161250in}{0.348806in}}%
\pgfpathquadraticcurveto{\pgfqpoint{6.155666in}{0.346447in}}{\pgfqpoint{6.149812in}{0.345204in}}%
\pgfpathquadraticcurveto{\pgfqpoint{6.143958in}{0.343961in}}{\pgfqpoint{6.137942in}{0.343961in}}%
\pgfpathquadraticcurveto{\pgfqpoint{6.122848in}{0.343961in}}{\pgfqpoint{6.114040in}{0.352733in}}%
\pgfpathquadraticcurveto{\pgfqpoint{6.105232in}{0.361523in}}{\pgfqpoint{6.105232in}{0.376509in}}%
\pgfpathquadraticcurveto{\pgfqpoint{6.105232in}{0.391981in}}{\pgfqpoint{6.113589in}{0.401059in}}%
\pgfpathquadraticcurveto{\pgfqpoint{6.121947in}{0.410156in}}{\pgfqpoint{6.136141in}{0.410156in}}%
\pgfpathquadraticcurveto{\pgfqpoint{6.148857in}{0.410156in}}{\pgfqpoint{6.156260in}{0.401960in}}%
\pgfpathquadraticcurveto{\pgfqpoint{6.163663in}{0.393783in}}{\pgfqpoint{6.163663in}{0.379715in}}%
\pgfpathlineto{\pgfqpoint{6.163663in}{0.379715in}}%
\pgfpathclose%
\pgfpathmoveto{\pgfqpoint{6.153306in}{0.382759in}}%
\pgfpathquadraticcurveto{\pgfqpoint{6.153198in}{0.391243in}}{\pgfqpoint{6.148551in}{0.396304in}}%
\pgfpathquadraticcurveto{\pgfqpoint{6.143904in}{0.401384in}}{\pgfqpoint{6.136249in}{0.401384in}}%
\pgfpathquadraticcurveto{\pgfqpoint{6.127585in}{0.401384in}}{\pgfqpoint{6.122379in}{0.396484in}}%
\pgfpathquadraticcurveto{\pgfqpoint{6.117174in}{0.391585in}}{\pgfqpoint{6.116381in}{0.382687in}}%
\pgfpathlineto{\pgfqpoint{6.153306in}{0.382759in}}%
\pgfpathlineto{\pgfqpoint{6.153306in}{0.382759in}}%
\pgfpathclose%
\pgfpathmoveto{\pgfqpoint{6.222149in}{0.399078in}}%
\pgfpathlineto{\pgfqpoint{6.222149in}{0.433193in}}%
\pgfpathlineto{\pgfqpoint{6.232506in}{0.433193in}}%
\pgfpathlineto{\pgfqpoint{6.232506in}{0.345600in}}%
\pgfpathlineto{\pgfqpoint{6.222149in}{0.345600in}}%
\pgfpathlineto{\pgfqpoint{6.222149in}{0.355056in}}%
\pgfpathquadraticcurveto{\pgfqpoint{6.218888in}{0.349437in}}{\pgfqpoint{6.213899in}{0.346699in}}%
\pgfpathquadraticcurveto{\pgfqpoint{6.208928in}{0.343961in}}{\pgfqpoint{6.201939in}{0.343961in}}%
\pgfpathquadraticcurveto{\pgfqpoint{6.190519in}{0.343961in}}{\pgfqpoint{6.183332in}{0.353075in}}%
\pgfpathquadraticcurveto{\pgfqpoint{6.176164in}{0.362207in}}{\pgfqpoint{6.176164in}{0.377067in}}%
\pgfpathquadraticcurveto{\pgfqpoint{6.176164in}{0.391927in}}{\pgfqpoint{6.183332in}{0.401041in}}%
\pgfpathquadraticcurveto{\pgfqpoint{6.190519in}{0.410156in}}{\pgfqpoint{6.201939in}{0.410156in}}%
\pgfpathquadraticcurveto{\pgfqpoint{6.208928in}{0.410156in}}{\pgfqpoint{6.213899in}{0.407418in}}%
\pgfpathquadraticcurveto{\pgfqpoint{6.218888in}{0.404698in}}{\pgfqpoint{6.222149in}{0.399078in}}%
\pgfpathlineto{\pgfqpoint{6.222149in}{0.399078in}}%
\pgfpathclose%
\pgfpathmoveto{\pgfqpoint{6.186863in}{0.377067in}}%
\pgfpathquadraticcurveto{\pgfqpoint{6.186863in}{0.365648in}}{\pgfqpoint{6.191564in}{0.359145in}}%
\pgfpathquadraticcurveto{\pgfqpoint{6.196265in}{0.352643in}}{\pgfqpoint{6.204479in}{0.352643in}}%
\pgfpathquadraticcurveto{\pgfqpoint{6.212692in}{0.352643in}}{\pgfqpoint{6.217411in}{0.359145in}}%
\pgfpathquadraticcurveto{\pgfqpoint{6.222149in}{0.365648in}}{\pgfqpoint{6.222149in}{0.377067in}}%
\pgfpathquadraticcurveto{\pgfqpoint{6.222149in}{0.388487in}}{\pgfqpoint{6.217411in}{0.394989in}}%
\pgfpathquadraticcurveto{\pgfqpoint{6.212692in}{0.401492in}}{\pgfqpoint{6.204479in}{0.401492in}}%
\pgfpathquadraticcurveto{\pgfqpoint{6.196265in}{0.401492in}}{\pgfqpoint{6.191564in}{0.394989in}}%
\pgfpathquadraticcurveto{\pgfqpoint{6.186863in}{0.388487in}}{\pgfqpoint{6.186863in}{0.377067in}}%
\pgfpathlineto{\pgfqpoint{6.186863in}{0.377067in}}%
\pgfpathclose%
\pgfpathmoveto{\pgfqpoint{6.330672in}{0.406787in}}%
\pgfpathlineto{\pgfqpoint{6.330672in}{0.396989in}}%
\pgfpathquadraticcurveto{\pgfqpoint{6.326295in}{0.399240in}}{\pgfqpoint{6.321558in}{0.400357in}}%
\pgfpathquadraticcurveto{\pgfqpoint{6.316839in}{0.401492in}}{\pgfqpoint{6.311759in}{0.401492in}}%
\pgfpathquadraticcurveto{\pgfqpoint{6.304050in}{0.401492in}}{\pgfqpoint{6.300195in}{0.399132in}}%
\pgfpathquadraticcurveto{\pgfqpoint{6.296341in}{0.396773in}}{\pgfqpoint{6.296341in}{0.392035in}}%
\pgfpathquadraticcurveto{\pgfqpoint{6.296341in}{0.388433in}}{\pgfqpoint{6.299097in}{0.386380in}}%
\pgfpathquadraticcurveto{\pgfqpoint{6.301852in}{0.384326in}}{\pgfqpoint{6.310192in}{0.382471in}}%
\pgfpathlineto{\pgfqpoint{6.313740in}{0.381678in}}%
\pgfpathquadraticcurveto{\pgfqpoint{6.324764in}{0.379319in}}{\pgfqpoint{6.329411in}{0.375014in}}%
\pgfpathquadraticcurveto{\pgfqpoint{6.334058in}{0.370709in}}{\pgfqpoint{6.334058in}{0.363000in}}%
\pgfpathquadraticcurveto{\pgfqpoint{6.334058in}{0.354210in}}{\pgfqpoint{6.327105in}{0.349076in}}%
\pgfpathquadraticcurveto{\pgfqpoint{6.320153in}{0.343961in}}{\pgfqpoint{6.307995in}{0.343961in}}%
\pgfpathquadraticcurveto{\pgfqpoint{6.302933in}{0.343961in}}{\pgfqpoint{6.297439in}{0.344952in}}%
\pgfpathquadraticcurveto{\pgfqpoint{6.291946in}{0.345942in}}{\pgfqpoint{6.285876in}{0.347906in}}%
\pgfpathlineto{\pgfqpoint{6.285876in}{0.358605in}}%
\pgfpathquadraticcurveto{\pgfqpoint{6.291622in}{0.355615in}}{\pgfqpoint{6.297187in}{0.354120in}}%
\pgfpathquadraticcurveto{\pgfqpoint{6.302753in}{0.352643in}}{\pgfqpoint{6.308229in}{0.352643in}}%
\pgfpathquadraticcurveto{\pgfqpoint{6.315542in}{0.352643in}}{\pgfqpoint{6.319468in}{0.355146in}}%
\pgfpathquadraticcurveto{\pgfqpoint{6.323413in}{0.357650in}}{\pgfqpoint{6.323413in}{0.362207in}}%
\pgfpathquadraticcurveto{\pgfqpoint{6.323413in}{0.366422in}}{\pgfqpoint{6.320567in}{0.368674in}}%
\pgfpathquadraticcurveto{\pgfqpoint{6.317739in}{0.370925in}}{\pgfqpoint{6.308103in}{0.373014in}}%
\pgfpathlineto{\pgfqpoint{6.304500in}{0.373861in}}%
\pgfpathquadraticcurveto{\pgfqpoint{6.294882in}{0.375878in}}{\pgfqpoint{6.290595in}{0.380075in}}%
\pgfpathquadraticcurveto{\pgfqpoint{6.286326in}{0.384272in}}{\pgfqpoint{6.286326in}{0.391585in}}%
\pgfpathquadraticcurveto{\pgfqpoint{6.286326in}{0.400483in}}{\pgfqpoint{6.292630in}{0.405310in}}%
\pgfpathquadraticcurveto{\pgfqpoint{6.298934in}{0.410156in}}{\pgfqpoint{6.310534in}{0.410156in}}%
\pgfpathquadraticcurveto{\pgfqpoint{6.316262in}{0.410156in}}{\pgfqpoint{6.321324in}{0.409309in}}%
\pgfpathquadraticcurveto{\pgfqpoint{6.326403in}{0.408480in}}{\pgfqpoint{6.330672in}{0.406787in}}%
\pgfpathlineto{\pgfqpoint{6.330672in}{0.406787in}}%
\pgfpathclose%
\pgfpathmoveto{\pgfqpoint{6.404468in}{0.379715in}}%
\pgfpathlineto{\pgfqpoint{6.404468in}{0.374654in}}%
\pgfpathlineto{\pgfqpoint{6.356844in}{0.374654in}}%
\pgfpathquadraticcurveto{\pgfqpoint{6.357528in}{0.363954in}}{\pgfqpoint{6.363292in}{0.358353in}}%
\pgfpathquadraticcurveto{\pgfqpoint{6.369056in}{0.352751in}}{\pgfqpoint{6.379359in}{0.352751in}}%
\pgfpathquadraticcurveto{\pgfqpoint{6.385321in}{0.352751in}}{\pgfqpoint{6.390923in}{0.354210in}}%
\pgfpathquadraticcurveto{\pgfqpoint{6.396524in}{0.355669in}}{\pgfqpoint{6.402054in}{0.358605in}}%
\pgfpathlineto{\pgfqpoint{6.402054in}{0.348806in}}%
\pgfpathquadraticcurveto{\pgfqpoint{6.396470in}{0.346447in}}{\pgfqpoint{6.390616in}{0.345204in}}%
\pgfpathquadraticcurveto{\pgfqpoint{6.384762in}{0.343961in}}{\pgfqpoint{6.378746in}{0.343961in}}%
\pgfpathquadraticcurveto{\pgfqpoint{6.363652in}{0.343961in}}{\pgfqpoint{6.354844in}{0.352733in}}%
\pgfpathquadraticcurveto{\pgfqpoint{6.346036in}{0.361523in}}{\pgfqpoint{6.346036in}{0.376509in}}%
\pgfpathquadraticcurveto{\pgfqpoint{6.346036in}{0.391981in}}{\pgfqpoint{6.354394in}{0.401059in}}%
\pgfpathquadraticcurveto{\pgfqpoint{6.362752in}{0.410156in}}{\pgfqpoint{6.376945in}{0.410156in}}%
\pgfpathquadraticcurveto{\pgfqpoint{6.389662in}{0.410156in}}{\pgfqpoint{6.397065in}{0.401960in}}%
\pgfpathquadraticcurveto{\pgfqpoint{6.404468in}{0.393783in}}{\pgfqpoint{6.404468in}{0.379715in}}%
\pgfpathlineto{\pgfqpoint{6.404468in}{0.379715in}}%
\pgfpathclose%
\pgfpathmoveto{\pgfqpoint{6.394111in}{0.382759in}}%
\pgfpathquadraticcurveto{\pgfqpoint{6.394003in}{0.391243in}}{\pgfqpoint{6.389355in}{0.396304in}}%
\pgfpathquadraticcurveto{\pgfqpoint{6.384708in}{0.401384in}}{\pgfqpoint{6.377053in}{0.401384in}}%
\pgfpathquadraticcurveto{\pgfqpoint{6.368389in}{0.401384in}}{\pgfqpoint{6.363184in}{0.396484in}}%
\pgfpathquadraticcurveto{\pgfqpoint{6.357978in}{0.391585in}}{\pgfqpoint{6.357186in}{0.382687in}}%
\pgfpathlineto{\pgfqpoint{6.394111in}{0.382759in}}%
\pgfpathlineto{\pgfqpoint{6.394111in}{0.382759in}}%
\pgfpathclose%
\pgfpathmoveto{\pgfqpoint{6.431486in}{0.355056in}}%
\pgfpathlineto{\pgfqpoint{6.431486in}{0.321626in}}%
\pgfpathlineto{\pgfqpoint{6.421075in}{0.321626in}}%
\pgfpathlineto{\pgfqpoint{6.421075in}{0.408643in}}%
\pgfpathlineto{\pgfqpoint{6.431486in}{0.408643in}}%
\pgfpathlineto{\pgfqpoint{6.431486in}{0.399078in}}%
\pgfpathquadraticcurveto{\pgfqpoint{6.434764in}{0.404698in}}{\pgfqpoint{6.439735in}{0.407418in}}%
\pgfpathquadraticcurveto{\pgfqpoint{6.444725in}{0.410156in}}{\pgfqpoint{6.451641in}{0.410156in}}%
\pgfpathquadraticcurveto{\pgfqpoint{6.463133in}{0.410156in}}{\pgfqpoint{6.470302in}{0.401041in}}%
\pgfpathquadraticcurveto{\pgfqpoint{6.477489in}{0.391927in}}{\pgfqpoint{6.477489in}{0.377067in}}%
\pgfpathquadraticcurveto{\pgfqpoint{6.477489in}{0.362207in}}{\pgfqpoint{6.470302in}{0.353075in}}%
\pgfpathquadraticcurveto{\pgfqpoint{6.463133in}{0.343961in}}{\pgfqpoint{6.451641in}{0.343961in}}%
\pgfpathquadraticcurveto{\pgfqpoint{6.444725in}{0.343961in}}{\pgfqpoint{6.439735in}{0.346699in}}%
\pgfpathquadraticcurveto{\pgfqpoint{6.434764in}{0.349437in}}{\pgfqpoint{6.431486in}{0.355056in}}%
\pgfpathlineto{\pgfqpoint{6.431486in}{0.355056in}}%
\pgfpathclose%
\pgfpathmoveto{\pgfqpoint{6.466736in}{0.377067in}}%
\pgfpathquadraticcurveto{\pgfqpoint{6.466736in}{0.388487in}}{\pgfqpoint{6.462035in}{0.394989in}}%
\pgfpathquadraticcurveto{\pgfqpoint{6.457333in}{0.401492in}}{\pgfqpoint{6.449120in}{0.401492in}}%
\pgfpathquadraticcurveto{\pgfqpoint{6.440888in}{0.401492in}}{\pgfqpoint{6.436187in}{0.394989in}}%
\pgfpathquadraticcurveto{\pgfqpoint{6.431486in}{0.388487in}}{\pgfqpoint{6.431486in}{0.377067in}}%
\pgfpathquadraticcurveto{\pgfqpoint{6.431486in}{0.365648in}}{\pgfqpoint{6.436187in}{0.359145in}}%
\pgfpathquadraticcurveto{\pgfqpoint{6.440888in}{0.352643in}}{\pgfqpoint{6.449120in}{0.352643in}}%
\pgfpathquadraticcurveto{\pgfqpoint{6.457333in}{0.352643in}}{\pgfqpoint{6.462035in}{0.359145in}}%
\pgfpathquadraticcurveto{\pgfqpoint{6.466736in}{0.365648in}}{\pgfqpoint{6.466736in}{0.377067in}}%
\pgfpathlineto{\pgfqpoint{6.466736in}{0.377067in}}%
\pgfpathclose%
\pgfpathmoveto{\pgfqpoint{6.523312in}{0.377283in}}%
\pgfpathquadraticcurveto{\pgfqpoint{6.510757in}{0.377283in}}{\pgfqpoint{6.505912in}{0.374419in}}%
\pgfpathquadraticcurveto{\pgfqpoint{6.501067in}{0.371556in}}{\pgfqpoint{6.501067in}{0.364621in}}%
\pgfpathquadraticcurveto{\pgfqpoint{6.501067in}{0.359109in}}{\pgfqpoint{6.504705in}{0.355867in}}%
\pgfpathquadraticcurveto{\pgfqpoint{6.508344in}{0.352643in}}{\pgfqpoint{6.514576in}{0.352643in}}%
\pgfpathquadraticcurveto{\pgfqpoint{6.523204in}{0.352643in}}{\pgfqpoint{6.528409in}{0.358749in}}%
\pgfpathquadraticcurveto{\pgfqpoint{6.533615in}{0.364855in}}{\pgfqpoint{6.533615in}{0.374978in}}%
\pgfpathlineto{\pgfqpoint{6.533615in}{0.377283in}}%
\pgfpathlineto{\pgfqpoint{6.523312in}{0.377283in}}%
\pgfpathlineto{\pgfqpoint{6.523312in}{0.377283in}}%
\pgfpathclose%
\pgfpathmoveto{\pgfqpoint{6.543972in}{0.381570in}}%
\pgfpathlineto{\pgfqpoint{6.543972in}{0.345600in}}%
\pgfpathlineto{\pgfqpoint{6.533615in}{0.345600in}}%
\pgfpathlineto{\pgfqpoint{6.533615in}{0.355164in}}%
\pgfpathquadraticcurveto{\pgfqpoint{6.530066in}{0.349437in}}{\pgfqpoint{6.524771in}{0.346699in}}%
\pgfpathquadraticcurveto{\pgfqpoint{6.519475in}{0.343961in}}{\pgfqpoint{6.511820in}{0.343961in}}%
\pgfpathquadraticcurveto{\pgfqpoint{6.502148in}{0.343961in}}{\pgfqpoint{6.496420in}{0.349401in}}%
\pgfpathquadraticcurveto{\pgfqpoint{6.490710in}{0.354840in}}{\pgfqpoint{6.490710in}{0.363954in}}%
\pgfpathquadraticcurveto{\pgfqpoint{6.490710in}{0.374582in}}{\pgfqpoint{6.497825in}{0.379985in}}%
\pgfpathquadraticcurveto{\pgfqpoint{6.504957in}{0.385389in}}{\pgfqpoint{6.519079in}{0.385389in}}%
\pgfpathlineto{\pgfqpoint{6.533615in}{0.385389in}}%
\pgfpathlineto{\pgfqpoint{6.533615in}{0.386416in}}%
\pgfpathquadraticcurveto{\pgfqpoint{6.533615in}{0.393566in}}{\pgfqpoint{6.528914in}{0.397475in}}%
\pgfpathquadraticcurveto{\pgfqpoint{6.524212in}{0.401384in}}{\pgfqpoint{6.515711in}{0.401384in}}%
\pgfpathquadraticcurveto{\pgfqpoint{6.510307in}{0.401384in}}{\pgfqpoint{6.505174in}{0.400087in}}%
\pgfpathquadraticcurveto{\pgfqpoint{6.500058in}{0.398790in}}{\pgfqpoint{6.495339in}{0.396196in}}%
\pgfpathlineto{\pgfqpoint{6.495339in}{0.405779in}}%
\pgfpathquadraticcurveto{\pgfqpoint{6.501013in}{0.407976in}}{\pgfqpoint{6.506362in}{0.409057in}}%
\pgfpathquadraticcurveto{\pgfqpoint{6.511712in}{0.410156in}}{\pgfqpoint{6.516773in}{0.410156in}}%
\pgfpathquadraticcurveto{\pgfqpoint{6.530463in}{0.410156in}}{\pgfqpoint{6.537217in}{0.403059in}}%
\pgfpathquadraticcurveto{\pgfqpoint{6.543972in}{0.395980in}}{\pgfqpoint{6.543972in}{0.381570in}}%
\pgfpathlineto{\pgfqpoint{6.543972in}{0.381570in}}%
\pgfpathclose%
\pgfpathmoveto{\pgfqpoint{6.601827in}{0.398970in}}%
\pgfpathquadraticcurveto{\pgfqpoint{6.600080in}{0.399979in}}{\pgfqpoint{6.598026in}{0.400447in}}%
\pgfpathquadraticcurveto{\pgfqpoint{6.595973in}{0.400933in}}{\pgfqpoint{6.593505in}{0.400933in}}%
\pgfpathquadraticcurveto{\pgfqpoint{6.584715in}{0.400933in}}{\pgfqpoint{6.580014in}{0.395223in}}%
\pgfpathquadraticcurveto{\pgfqpoint{6.575313in}{0.389514in}}{\pgfqpoint{6.575313in}{0.378814in}}%
\pgfpathlineto{\pgfqpoint{6.575313in}{0.345600in}}%
\pgfpathlineto{\pgfqpoint{6.564902in}{0.345600in}}%
\pgfpathlineto{\pgfqpoint{6.564902in}{0.408643in}}%
\pgfpathlineto{\pgfqpoint{6.575313in}{0.408643in}}%
\pgfpathlineto{\pgfqpoint{6.575313in}{0.398844in}}%
\pgfpathquadraticcurveto{\pgfqpoint{6.578591in}{0.404590in}}{\pgfqpoint{6.583815in}{0.407364in}}%
\pgfpathquadraticcurveto{\pgfqpoint{6.589056in}{0.410156in}}{\pgfqpoint{6.596549in}{0.410156in}}%
\pgfpathquadraticcurveto{\pgfqpoint{6.597612in}{0.410156in}}{\pgfqpoint{6.598909in}{0.410011in}}%
\pgfpathquadraticcurveto{\pgfqpoint{6.600206in}{0.409885in}}{\pgfqpoint{6.601773in}{0.409597in}}%
\pgfpathlineto{\pgfqpoint{6.601827in}{0.398970in}}%
\pgfpathlineto{\pgfqpoint{6.601827in}{0.398970in}}%
\pgfpathclose%
\pgfpathmoveto{\pgfqpoint{6.641345in}{0.377283in}}%
\pgfpathquadraticcurveto{\pgfqpoint{6.628791in}{0.377283in}}{\pgfqpoint{6.623946in}{0.374419in}}%
\pgfpathquadraticcurveto{\pgfqpoint{6.619100in}{0.371556in}}{\pgfqpoint{6.619100in}{0.364621in}}%
\pgfpathquadraticcurveto{\pgfqpoint{6.619100in}{0.359109in}}{\pgfqpoint{6.622739in}{0.355867in}}%
\pgfpathquadraticcurveto{\pgfqpoint{6.626377in}{0.352643in}}{\pgfqpoint{6.632610in}{0.352643in}}%
\pgfpathquadraticcurveto{\pgfqpoint{6.641237in}{0.352643in}}{\pgfqpoint{6.646443in}{0.358749in}}%
\pgfpathquadraticcurveto{\pgfqpoint{6.651648in}{0.364855in}}{\pgfqpoint{6.651648in}{0.374978in}}%
\pgfpathlineto{\pgfqpoint{6.651648in}{0.377283in}}%
\pgfpathlineto{\pgfqpoint{6.641345in}{0.377283in}}%
\pgfpathlineto{\pgfqpoint{6.641345in}{0.377283in}}%
\pgfpathclose%
\pgfpathmoveto{\pgfqpoint{6.662005in}{0.381570in}}%
\pgfpathlineto{\pgfqpoint{6.662005in}{0.345600in}}%
\pgfpathlineto{\pgfqpoint{6.651648in}{0.345600in}}%
\pgfpathlineto{\pgfqpoint{6.651648in}{0.355164in}}%
\pgfpathquadraticcurveto{\pgfqpoint{6.648100in}{0.349437in}}{\pgfqpoint{6.642804in}{0.346699in}}%
\pgfpathquadraticcurveto{\pgfqpoint{6.637509in}{0.343961in}}{\pgfqpoint{6.629854in}{0.343961in}}%
\pgfpathquadraticcurveto{\pgfqpoint{6.620181in}{0.343961in}}{\pgfqpoint{6.614453in}{0.349401in}}%
\pgfpathquadraticcurveto{\pgfqpoint{6.608743in}{0.354840in}}{\pgfqpoint{6.608743in}{0.363954in}}%
\pgfpathquadraticcurveto{\pgfqpoint{6.608743in}{0.374582in}}{\pgfqpoint{6.615858in}{0.379985in}}%
\pgfpathquadraticcurveto{\pgfqpoint{6.622991in}{0.385389in}}{\pgfqpoint{6.637113in}{0.385389in}}%
\pgfpathlineto{\pgfqpoint{6.651648in}{0.385389in}}%
\pgfpathlineto{\pgfqpoint{6.651648in}{0.386416in}}%
\pgfpathquadraticcurveto{\pgfqpoint{6.651648in}{0.393566in}}{\pgfqpoint{6.646947in}{0.397475in}}%
\pgfpathquadraticcurveto{\pgfqpoint{6.642246in}{0.401384in}}{\pgfqpoint{6.633744in}{0.401384in}}%
\pgfpathquadraticcurveto{\pgfqpoint{6.628341in}{0.401384in}}{\pgfqpoint{6.623207in}{0.400087in}}%
\pgfpathquadraticcurveto{\pgfqpoint{6.618092in}{0.398790in}}{\pgfqpoint{6.613373in}{0.396196in}}%
\pgfpathlineto{\pgfqpoint{6.613373in}{0.405779in}}%
\pgfpathquadraticcurveto{\pgfqpoint{6.619046in}{0.407976in}}{\pgfqpoint{6.624396in}{0.409057in}}%
\pgfpathquadraticcurveto{\pgfqpoint{6.629746in}{0.410156in}}{\pgfqpoint{6.634807in}{0.410156in}}%
\pgfpathquadraticcurveto{\pgfqpoint{6.648496in}{0.410156in}}{\pgfqpoint{6.655251in}{0.403059in}}%
\pgfpathquadraticcurveto{\pgfqpoint{6.662005in}{0.395980in}}{\pgfqpoint{6.662005in}{0.381570in}}%
\pgfpathlineto{\pgfqpoint{6.662005in}{0.381570in}}%
\pgfpathclose%
\pgfpathmoveto{\pgfqpoint{6.693581in}{0.426547in}}%
\pgfpathlineto{\pgfqpoint{6.693581in}{0.408643in}}%
\pgfpathlineto{\pgfqpoint{6.714907in}{0.408643in}}%
\pgfpathlineto{\pgfqpoint{6.714907in}{0.400591in}}%
\pgfpathlineto{\pgfqpoint{6.693581in}{0.400591in}}%
\pgfpathlineto{\pgfqpoint{6.693581in}{0.366368in}}%
\pgfpathquadraticcurveto{\pgfqpoint{6.693581in}{0.358659in}}{\pgfqpoint{6.695688in}{0.356461in}}%
\pgfpathquadraticcurveto{\pgfqpoint{6.697796in}{0.354264in}}{\pgfqpoint{6.704280in}{0.354264in}}%
\pgfpathlineto{\pgfqpoint{6.714907in}{0.354264in}}%
\pgfpathlineto{\pgfqpoint{6.714907in}{0.345600in}}%
\pgfpathlineto{\pgfqpoint{6.704280in}{0.345600in}}%
\pgfpathquadraticcurveto{\pgfqpoint{6.692284in}{0.345600in}}{\pgfqpoint{6.687727in}{0.350067in}}%
\pgfpathquadraticcurveto{\pgfqpoint{6.683170in}{0.354552in}}{\pgfqpoint{6.683170in}{0.366368in}}%
\pgfpathlineto{\pgfqpoint{6.683170in}{0.400591in}}%
\pgfpathlineto{\pgfqpoint{6.675569in}{0.400591in}}%
\pgfpathlineto{\pgfqpoint{6.675569in}{0.408643in}}%
\pgfpathlineto{\pgfqpoint{6.683170in}{0.408643in}}%
\pgfpathlineto{\pgfqpoint{6.683170in}{0.426547in}}%
\pgfpathlineto{\pgfqpoint{6.693581in}{0.426547in}}%
\pgfpathlineto{\pgfqpoint{6.693581in}{0.426547in}}%
\pgfpathclose%
\pgfpathmoveto{\pgfqpoint{6.782453in}{0.379715in}}%
\pgfpathlineto{\pgfqpoint{6.782453in}{0.374654in}}%
\pgfpathlineto{\pgfqpoint{6.734829in}{0.374654in}}%
\pgfpathquadraticcurveto{\pgfqpoint{6.735513in}{0.363954in}}{\pgfqpoint{6.741277in}{0.358353in}}%
\pgfpathquadraticcurveto{\pgfqpoint{6.747041in}{0.352751in}}{\pgfqpoint{6.757344in}{0.352751in}}%
\pgfpathquadraticcurveto{\pgfqpoint{6.763306in}{0.352751in}}{\pgfqpoint{6.768908in}{0.354210in}}%
\pgfpathquadraticcurveto{\pgfqpoint{6.774509in}{0.355669in}}{\pgfqpoint{6.780039in}{0.358605in}}%
\pgfpathlineto{\pgfqpoint{6.780039in}{0.348806in}}%
\pgfpathquadraticcurveto{\pgfqpoint{6.774455in}{0.346447in}}{\pgfqpoint{6.768601in}{0.345204in}}%
\pgfpathquadraticcurveto{\pgfqpoint{6.762747in}{0.343961in}}{\pgfqpoint{6.756731in}{0.343961in}}%
\pgfpathquadraticcurveto{\pgfqpoint{6.741637in}{0.343961in}}{\pgfqpoint{6.732829in}{0.352733in}}%
\pgfpathquadraticcurveto{\pgfqpoint{6.724021in}{0.361523in}}{\pgfqpoint{6.724021in}{0.376509in}}%
\pgfpathquadraticcurveto{\pgfqpoint{6.724021in}{0.391981in}}{\pgfqpoint{6.732379in}{0.401059in}}%
\pgfpathquadraticcurveto{\pgfqpoint{6.740737in}{0.410156in}}{\pgfqpoint{6.754930in}{0.410156in}}%
\pgfpathquadraticcurveto{\pgfqpoint{6.767647in}{0.410156in}}{\pgfqpoint{6.775050in}{0.401960in}}%
\pgfpathquadraticcurveto{\pgfqpoint{6.782453in}{0.393783in}}{\pgfqpoint{6.782453in}{0.379715in}}%
\pgfpathlineto{\pgfqpoint{6.782453in}{0.379715in}}%
\pgfpathclose%
\pgfpathmoveto{\pgfqpoint{6.772096in}{0.382759in}}%
\pgfpathquadraticcurveto{\pgfqpoint{6.771988in}{0.391243in}}{\pgfqpoint{6.767340in}{0.396304in}}%
\pgfpathquadraticcurveto{\pgfqpoint{6.762693in}{0.401384in}}{\pgfqpoint{6.755038in}{0.401384in}}%
\pgfpathquadraticcurveto{\pgfqpoint{6.746374in}{0.401384in}}{\pgfqpoint{6.741169in}{0.396484in}}%
\pgfpathquadraticcurveto{\pgfqpoint{6.735963in}{0.391585in}}{\pgfqpoint{6.735171in}{0.382687in}}%
\pgfpathlineto{\pgfqpoint{6.772096in}{0.382759in}}%
\pgfpathlineto{\pgfqpoint{6.772096in}{0.382759in}}%
\pgfpathclose%
\pgfpathmoveto{\pgfqpoint{6.799456in}{0.433193in}}%
\pgfpathlineto{\pgfqpoint{6.809813in}{0.433193in}}%
\pgfpathlineto{\pgfqpoint{6.809813in}{0.345600in}}%
\pgfpathlineto{\pgfqpoint{6.799456in}{0.345600in}}%
\pgfpathlineto{\pgfqpoint{6.799456in}{0.433193in}}%
\pgfpathlineto{\pgfqpoint{6.799456in}{0.433193in}}%
\pgfpathclose%
\pgfpathmoveto{\pgfqpoint{6.857707in}{0.339746in}}%
\pgfpathquadraticcurveto{\pgfqpoint{6.853331in}{0.328488in}}{\pgfqpoint{6.849152in}{0.325066in}}%
\pgfpathquadraticcurveto{\pgfqpoint{6.844991in}{0.321626in}}{\pgfqpoint{6.838020in}{0.321626in}}%
\pgfpathlineto{\pgfqpoint{6.829735in}{0.321626in}}%
\pgfpathlineto{\pgfqpoint{6.829735in}{0.330290in}}%
\pgfpathlineto{\pgfqpoint{6.835823in}{0.330290in}}%
\pgfpathquadraticcurveto{\pgfqpoint{6.840092in}{0.330290in}}{\pgfqpoint{6.842451in}{0.332325in}}%
\pgfpathquadraticcurveto{\pgfqpoint{6.844829in}{0.334342in}}{\pgfqpoint{6.847693in}{0.341889in}}%
\pgfpathlineto{\pgfqpoint{6.849548in}{0.346609in}}%
\pgfpathlineto{\pgfqpoint{6.824061in}{0.408643in}}%
\pgfpathlineto{\pgfqpoint{6.835030in}{0.408643in}}%
\pgfpathlineto{\pgfqpoint{6.854735in}{0.359343in}}%
\pgfpathlineto{\pgfqpoint{6.874441in}{0.408643in}}%
\pgfpathlineto{\pgfqpoint{6.885410in}{0.408643in}}%
\pgfpathlineto{\pgfqpoint{6.857707in}{0.339746in}}%
\pgfpathlineto{\pgfqpoint{6.857707in}{0.339746in}}%
\pgfpathclose%
\pgfpathmoveto{\pgfqpoint{6.884744in}{0.359902in}}%
\pgfpathlineto{\pgfqpoint{6.896632in}{0.359902in}}%
\pgfpathlineto{\pgfqpoint{6.896632in}{0.345600in}}%
\pgfpathlineto{\pgfqpoint{6.884744in}{0.345600in}}%
\pgfpathlineto{\pgfqpoint{6.884744in}{0.359902in}}%
\pgfpathlineto{\pgfqpoint{6.884744in}{0.359902in}}%
\pgfpathclose%
\pgfpathmoveto{\pgfqpoint{7.007388in}{0.426889in}}%
\pgfpathlineto{\pgfqpoint{7.007388in}{0.415793in}}%
\pgfpathquadraticcurveto{\pgfqpoint{7.000922in}{0.418891in}}{\pgfqpoint{6.995176in}{0.420404in}}%
\pgfpathquadraticcurveto{\pgfqpoint{6.989430in}{0.421936in}}{\pgfqpoint{6.984081in}{0.421936in}}%
\pgfpathquadraticcurveto{\pgfqpoint{6.974804in}{0.421936in}}{\pgfqpoint{6.969761in}{0.418333in}}%
\pgfpathquadraticcurveto{\pgfqpoint{6.964718in}{0.414731in}}{\pgfqpoint{6.964718in}{0.408084in}}%
\pgfpathquadraticcurveto{\pgfqpoint{6.964718in}{0.402500in}}{\pgfqpoint{6.968068in}{0.399654in}}%
\pgfpathquadraticcurveto{\pgfqpoint{6.971418in}{0.396827in}}{\pgfqpoint{6.980766in}{0.395079in}}%
\pgfpathlineto{\pgfqpoint{6.987629in}{0.393674in}}%
\pgfpathquadraticcurveto{\pgfqpoint{7.000346in}{0.391243in}}{\pgfqpoint{7.006398in}{0.385137in}}%
\pgfpathquadraticcurveto{\pgfqpoint{7.012450in}{0.379031in}}{\pgfqpoint{7.012450in}{0.368800in}}%
\pgfpathquadraticcurveto{\pgfqpoint{7.012450in}{0.356569in}}{\pgfqpoint{7.004254in}{0.350265in}}%
\pgfpathquadraticcurveto{\pgfqpoint{6.996077in}{0.343961in}}{\pgfqpoint{6.980262in}{0.343961in}}%
\pgfpathquadraticcurveto{\pgfqpoint{6.974300in}{0.343961in}}{\pgfqpoint{6.967564in}{0.345312in}}%
\pgfpathquadraticcurveto{\pgfqpoint{6.960845in}{0.346663in}}{\pgfqpoint{6.953640in}{0.349311in}}%
\pgfpathlineto{\pgfqpoint{6.953640in}{0.361018in}}%
\pgfpathquadraticcurveto{\pgfqpoint{6.960557in}{0.357146in}}{\pgfqpoint{6.967203in}{0.355164in}}%
\pgfpathquadraticcurveto{\pgfqpoint{6.973850in}{0.353201in}}{\pgfqpoint{6.980262in}{0.353201in}}%
\pgfpathquadraticcurveto{\pgfqpoint{6.989989in}{0.353201in}}{\pgfqpoint{6.995284in}{0.357020in}}%
\pgfpathquadraticcurveto{\pgfqpoint{7.000580in}{0.360856in}}{\pgfqpoint{7.000580in}{0.367953in}}%
\pgfpathquadraticcurveto{\pgfqpoint{7.000580in}{0.374131in}}{\pgfqpoint{6.996779in}{0.377626in}}%
\pgfpathquadraticcurveto{\pgfqpoint{6.992979in}{0.381120in}}{\pgfqpoint{6.984315in}{0.382867in}}%
\pgfpathlineto{\pgfqpoint{6.977380in}{0.384218in}}%
\pgfpathquadraticcurveto{\pgfqpoint{6.964664in}{0.386740in}}{\pgfqpoint{6.958972in}{0.392143in}}%
\pgfpathquadraticcurveto{\pgfqpoint{6.953298in}{0.397547in}}{\pgfqpoint{6.953298in}{0.407184in}}%
\pgfpathquadraticcurveto{\pgfqpoint{6.953298in}{0.418333in}}{\pgfqpoint{6.961151in}{0.424745in}}%
\pgfpathquadraticcurveto{\pgfqpoint{6.969005in}{0.431158in}}{\pgfqpoint{6.982784in}{0.431158in}}%
\pgfpathquadraticcurveto{\pgfqpoint{6.988710in}{0.431158in}}{\pgfqpoint{6.994834in}{0.430077in}}%
\pgfpathquadraticcurveto{\pgfqpoint{7.000976in}{0.429014in}}{\pgfqpoint{7.007388in}{0.426889in}}%
\pgfpathlineto{\pgfqpoint{7.007388in}{0.426889in}}%
\pgfpathclose%
\pgfpathmoveto{\pgfqpoint{7.039990in}{0.426547in}}%
\pgfpathlineto{\pgfqpoint{7.039990in}{0.408643in}}%
\pgfpathlineto{\pgfqpoint{7.061317in}{0.408643in}}%
\pgfpathlineto{\pgfqpoint{7.061317in}{0.400591in}}%
\pgfpathlineto{\pgfqpoint{7.039990in}{0.400591in}}%
\pgfpathlineto{\pgfqpoint{7.039990in}{0.366368in}}%
\pgfpathquadraticcurveto{\pgfqpoint{7.039990in}{0.358659in}}{\pgfqpoint{7.042098in}{0.356461in}}%
\pgfpathquadraticcurveto{\pgfqpoint{7.044205in}{0.354264in}}{\pgfqpoint{7.050690in}{0.354264in}}%
\pgfpathlineto{\pgfqpoint{7.061317in}{0.354264in}}%
\pgfpathlineto{\pgfqpoint{7.061317in}{0.345600in}}%
\pgfpathlineto{\pgfqpoint{7.050690in}{0.345600in}}%
\pgfpathquadraticcurveto{\pgfqpoint{7.038694in}{0.345600in}}{\pgfqpoint{7.034137in}{0.350067in}}%
\pgfpathquadraticcurveto{\pgfqpoint{7.029579in}{0.354552in}}{\pgfqpoint{7.029579in}{0.366368in}}%
\pgfpathlineto{\pgfqpoint{7.029579in}{0.400591in}}%
\pgfpathlineto{\pgfqpoint{7.021978in}{0.400591in}}%
\pgfpathlineto{\pgfqpoint{7.021978in}{0.408643in}}%
\pgfpathlineto{\pgfqpoint{7.029579in}{0.408643in}}%
\pgfpathlineto{\pgfqpoint{7.029579in}{0.426547in}}%
\pgfpathlineto{\pgfqpoint{7.039990in}{0.426547in}}%
\pgfpathlineto{\pgfqpoint{7.039990in}{0.426547in}}%
\pgfpathclose%
\pgfpathmoveto{\pgfqpoint{7.073871in}{0.370475in}}%
\pgfpathlineto{\pgfqpoint{7.073871in}{0.408643in}}%
\pgfpathlineto{\pgfqpoint{7.084228in}{0.408643in}}%
\pgfpathlineto{\pgfqpoint{7.084228in}{0.370871in}}%
\pgfpathquadraticcurveto{\pgfqpoint{7.084228in}{0.361919in}}{\pgfqpoint{7.087705in}{0.357434in}}%
\pgfpathquadraticcurveto{\pgfqpoint{7.091199in}{0.352967in}}{\pgfqpoint{7.098188in}{0.352967in}}%
\pgfpathquadraticcurveto{\pgfqpoint{7.106563in}{0.352967in}}{\pgfqpoint{7.111427in}{0.358317in}}%
\pgfpathquadraticcurveto{\pgfqpoint{7.116308in}{0.363666in}}{\pgfqpoint{7.116308in}{0.372906in}}%
\pgfpathlineto{\pgfqpoint{7.116308in}{0.408643in}}%
\pgfpathlineto{\pgfqpoint{7.126665in}{0.408643in}}%
\pgfpathlineto{\pgfqpoint{7.126665in}{0.345600in}}%
\pgfpathlineto{\pgfqpoint{7.116308in}{0.345600in}}%
\pgfpathlineto{\pgfqpoint{7.116308in}{0.355291in}}%
\pgfpathquadraticcurveto{\pgfqpoint{7.112543in}{0.349545in}}{\pgfqpoint{7.107554in}{0.346753in}}%
\pgfpathquadraticcurveto{\pgfqpoint{7.102583in}{0.343961in}}{\pgfqpoint{7.095990in}{0.343961in}}%
\pgfpathquadraticcurveto{\pgfqpoint{7.085129in}{0.343961in}}{\pgfqpoint{7.079491in}{0.350715in}}%
\pgfpathquadraticcurveto{\pgfqpoint{7.073871in}{0.357470in}}{\pgfqpoint{7.073871in}{0.370475in}}%
\pgfpathlineto{\pgfqpoint{7.073871in}{0.370475in}}%
\pgfpathclose%
\pgfpathmoveto{\pgfqpoint{7.099935in}{0.410156in}}%
\pgfpathlineto{\pgfqpoint{7.099935in}{0.410156in}}%
\pgfpathlineto{\pgfqpoint{7.099935in}{0.410156in}}%
\pgfpathclose%
\pgfpathmoveto{\pgfqpoint{7.189473in}{0.399078in}}%
\pgfpathlineto{\pgfqpoint{7.189473in}{0.433193in}}%
\pgfpathlineto{\pgfqpoint{7.199830in}{0.433193in}}%
\pgfpathlineto{\pgfqpoint{7.199830in}{0.345600in}}%
\pgfpathlineto{\pgfqpoint{7.189473in}{0.345600in}}%
\pgfpathlineto{\pgfqpoint{7.189473in}{0.355056in}}%
\pgfpathquadraticcurveto{\pgfqpoint{7.186213in}{0.349437in}}{\pgfqpoint{7.181224in}{0.346699in}}%
\pgfpathquadraticcurveto{\pgfqpoint{7.176252in}{0.343961in}}{\pgfqpoint{7.169264in}{0.343961in}}%
\pgfpathquadraticcurveto{\pgfqpoint{7.157844in}{0.343961in}}{\pgfqpoint{7.150657in}{0.353075in}}%
\pgfpathquadraticcurveto{\pgfqpoint{7.143488in}{0.362207in}}{\pgfqpoint{7.143488in}{0.377067in}}%
\pgfpathquadraticcurveto{\pgfqpoint{7.143488in}{0.391927in}}{\pgfqpoint{7.150657in}{0.401041in}}%
\pgfpathquadraticcurveto{\pgfqpoint{7.157844in}{0.410156in}}{\pgfqpoint{7.169264in}{0.410156in}}%
\pgfpathquadraticcurveto{\pgfqpoint{7.176252in}{0.410156in}}{\pgfqpoint{7.181224in}{0.407418in}}%
\pgfpathquadraticcurveto{\pgfqpoint{7.186213in}{0.404698in}}{\pgfqpoint{7.189473in}{0.399078in}}%
\pgfpathlineto{\pgfqpoint{7.189473in}{0.399078in}}%
\pgfpathclose%
\pgfpathmoveto{\pgfqpoint{7.154188in}{0.377067in}}%
\pgfpathquadraticcurveto{\pgfqpoint{7.154188in}{0.365648in}}{\pgfqpoint{7.158889in}{0.359145in}}%
\pgfpathquadraticcurveto{\pgfqpoint{7.163590in}{0.352643in}}{\pgfqpoint{7.171803in}{0.352643in}}%
\pgfpathquadraticcurveto{\pgfqpoint{7.180017in}{0.352643in}}{\pgfqpoint{7.184736in}{0.359145in}}%
\pgfpathquadraticcurveto{\pgfqpoint{7.189473in}{0.365648in}}{\pgfqpoint{7.189473in}{0.377067in}}%
\pgfpathquadraticcurveto{\pgfqpoint{7.189473in}{0.388487in}}{\pgfqpoint{7.184736in}{0.394989in}}%
\pgfpathquadraticcurveto{\pgfqpoint{7.180017in}{0.401492in}}{\pgfqpoint{7.171803in}{0.401492in}}%
\pgfpathquadraticcurveto{\pgfqpoint{7.163590in}{0.401492in}}{\pgfqpoint{7.158889in}{0.394989in}}%
\pgfpathquadraticcurveto{\pgfqpoint{7.154188in}{0.388487in}}{\pgfqpoint{7.154188in}{0.377067in}}%
\pgfpathlineto{\pgfqpoint{7.154188in}{0.377067in}}%
\pgfpathclose%
\pgfpathmoveto{\pgfqpoint{7.247400in}{0.339746in}}%
\pgfpathquadraticcurveto{\pgfqpoint{7.243023in}{0.328488in}}{\pgfqpoint{7.238845in}{0.325066in}}%
\pgfpathquadraticcurveto{\pgfqpoint{7.234684in}{0.321626in}}{\pgfqpoint{7.227713in}{0.321626in}}%
\pgfpathlineto{\pgfqpoint{7.219428in}{0.321626in}}%
\pgfpathlineto{\pgfqpoint{7.219428in}{0.330290in}}%
\pgfpathlineto{\pgfqpoint{7.225516in}{0.330290in}}%
\pgfpathquadraticcurveto{\pgfqpoint{7.229785in}{0.330290in}}{\pgfqpoint{7.232144in}{0.332325in}}%
\pgfpathquadraticcurveto{\pgfqpoint{7.234522in}{0.334342in}}{\pgfqpoint{7.237386in}{0.341889in}}%
\pgfpathlineto{\pgfqpoint{7.239241in}{0.346609in}}%
\pgfpathlineto{\pgfqpoint{7.213754in}{0.408643in}}%
\pgfpathlineto{\pgfqpoint{7.224723in}{0.408643in}}%
\pgfpathlineto{\pgfqpoint{7.244428in}{0.359343in}}%
\pgfpathlineto{\pgfqpoint{7.264134in}{0.408643in}}%
\pgfpathlineto{\pgfqpoint{7.275103in}{0.408643in}}%
\pgfpathlineto{\pgfqpoint{7.247400in}{0.339746in}}%
\pgfpathlineto{\pgfqpoint{7.247400in}{0.339746in}}%
\pgfpathclose%
\pgfpathmoveto{\pgfqpoint{7.358752in}{0.419738in}}%
\pgfpathlineto{\pgfqpoint{7.330040in}{0.374870in}}%
\pgfpathlineto{\pgfqpoint{7.358752in}{0.374870in}}%
\pgfpathlineto{\pgfqpoint{7.358752in}{0.419738in}}%
\pgfpathlineto{\pgfqpoint{7.358752in}{0.419738in}}%
\pgfpathclose%
\pgfpathmoveto{\pgfqpoint{7.355762in}{0.429645in}}%
\pgfpathlineto{\pgfqpoint{7.370063in}{0.429645in}}%
\pgfpathlineto{\pgfqpoint{7.370063in}{0.374870in}}%
\pgfpathlineto{\pgfqpoint{7.382059in}{0.374870in}}%
\pgfpathlineto{\pgfqpoint{7.382059in}{0.365413in}}%
\pgfpathlineto{\pgfqpoint{7.370063in}{0.365413in}}%
\pgfpathlineto{\pgfqpoint{7.370063in}{0.345600in}}%
\pgfpathlineto{\pgfqpoint{7.358752in}{0.345600in}}%
\pgfpathlineto{\pgfqpoint{7.358752in}{0.365413in}}%
\pgfpathlineto{\pgfqpoint{7.320818in}{0.365413in}}%
\pgfpathlineto{\pgfqpoint{7.320818in}{0.376383in}}%
\pgfpathlineto{\pgfqpoint{7.355762in}{0.429645in}}%
\pgfpathlineto{\pgfqpoint{7.355762in}{0.429645in}}%
\pgfpathclose%
\pgfpathmoveto{\pgfqpoint{7.481396in}{0.406229in}}%
\pgfpathlineto{\pgfqpoint{7.481396in}{0.396538in}}%
\pgfpathquadraticcurveto{\pgfqpoint{7.477001in}{0.398970in}}{\pgfqpoint{7.472588in}{0.400177in}}%
\pgfpathquadraticcurveto{\pgfqpoint{7.468175in}{0.401384in}}{\pgfqpoint{7.463672in}{0.401384in}}%
\pgfpathquadraticcurveto{\pgfqpoint{7.453586in}{0.401384in}}{\pgfqpoint{7.448002in}{0.394989in}}%
\pgfpathquadraticcurveto{\pgfqpoint{7.442436in}{0.388613in}}{\pgfqpoint{7.442436in}{0.377067in}}%
\pgfpathquadraticcurveto{\pgfqpoint{7.442436in}{0.365521in}}{\pgfqpoint{7.448002in}{0.359127in}}%
\pgfpathquadraticcurveto{\pgfqpoint{7.453586in}{0.352751in}}{\pgfqpoint{7.463672in}{0.352751in}}%
\pgfpathquadraticcurveto{\pgfqpoint{7.468175in}{0.352751in}}{\pgfqpoint{7.472588in}{0.353958in}}%
\pgfpathquadraticcurveto{\pgfqpoint{7.477001in}{0.355164in}}{\pgfqpoint{7.481396in}{0.357596in}}%
\pgfpathlineto{\pgfqpoint{7.481396in}{0.348014in}}%
\pgfpathquadraticcurveto{\pgfqpoint{7.477055in}{0.345996in}}{\pgfqpoint{7.472408in}{0.344988in}}%
\pgfpathquadraticcurveto{\pgfqpoint{7.467779in}{0.343961in}}{\pgfqpoint{7.462538in}{0.343961in}}%
\pgfpathquadraticcurveto{\pgfqpoint{7.448290in}{0.343961in}}{\pgfqpoint{7.439896in}{0.352913in}}%
\pgfpathquadraticcurveto{\pgfqpoint{7.431521in}{0.361865in}}{\pgfqpoint{7.431521in}{0.377067in}}%
\pgfpathquadraticcurveto{\pgfqpoint{7.431521in}{0.392486in}}{\pgfqpoint{7.439986in}{0.401312in}}%
\pgfpathquadraticcurveto{\pgfqpoint{7.448470in}{0.410156in}}{\pgfqpoint{7.463222in}{0.410156in}}%
\pgfpathquadraticcurveto{\pgfqpoint{7.467995in}{0.410156in}}{\pgfqpoint{7.472552in}{0.409165in}}%
\pgfpathquadraticcurveto{\pgfqpoint{7.477109in}{0.408192in}}{\pgfqpoint{7.481396in}{0.406229in}}%
\pgfpathlineto{\pgfqpoint{7.481396in}{0.406229in}}%
\pgfpathclose%
\pgfpathmoveto{\pgfqpoint{7.499408in}{0.433193in}}%
\pgfpathlineto{\pgfqpoint{7.509765in}{0.433193in}}%
\pgfpathlineto{\pgfqpoint{7.509765in}{0.345600in}}%
\pgfpathlineto{\pgfqpoint{7.499408in}{0.345600in}}%
\pgfpathlineto{\pgfqpoint{7.499408in}{0.433193in}}%
\pgfpathlineto{\pgfqpoint{7.499408in}{0.433193in}}%
\pgfpathclose%
\pgfpathmoveto{\pgfqpoint{7.530371in}{0.370475in}}%
\pgfpathlineto{\pgfqpoint{7.530371in}{0.408643in}}%
\pgfpathlineto{\pgfqpoint{7.540728in}{0.408643in}}%
\pgfpathlineto{\pgfqpoint{7.540728in}{0.370871in}}%
\pgfpathquadraticcurveto{\pgfqpoint{7.540728in}{0.361919in}}{\pgfqpoint{7.544205in}{0.357434in}}%
\pgfpathquadraticcurveto{\pgfqpoint{7.547699in}{0.352967in}}{\pgfqpoint{7.554688in}{0.352967in}}%
\pgfpathquadraticcurveto{\pgfqpoint{7.563063in}{0.352967in}}{\pgfqpoint{7.567927in}{0.358317in}}%
\pgfpathquadraticcurveto{\pgfqpoint{7.572808in}{0.363666in}}{\pgfqpoint{7.572808in}{0.372906in}}%
\pgfpathlineto{\pgfqpoint{7.572808in}{0.408643in}}%
\pgfpathlineto{\pgfqpoint{7.583165in}{0.408643in}}%
\pgfpathlineto{\pgfqpoint{7.583165in}{0.345600in}}%
\pgfpathlineto{\pgfqpoint{7.572808in}{0.345600in}}%
\pgfpathlineto{\pgfqpoint{7.572808in}{0.355291in}}%
\pgfpathquadraticcurveto{\pgfqpoint{7.569043in}{0.349545in}}{\pgfqpoint{7.564054in}{0.346753in}}%
\pgfpathquadraticcurveto{\pgfqpoint{7.559083in}{0.343961in}}{\pgfqpoint{7.552490in}{0.343961in}}%
\pgfpathquadraticcurveto{\pgfqpoint{7.541629in}{0.343961in}}{\pgfqpoint{7.535991in}{0.350715in}}%
\pgfpathquadraticcurveto{\pgfqpoint{7.530371in}{0.357470in}}{\pgfqpoint{7.530371in}{0.370475in}}%
\pgfpathlineto{\pgfqpoint{7.530371in}{0.370475in}}%
\pgfpathclose%
\pgfpathmoveto{\pgfqpoint{7.556435in}{0.410156in}}%
\pgfpathlineto{\pgfqpoint{7.556435in}{0.410156in}}%
\pgfpathlineto{\pgfqpoint{7.556435in}{0.410156in}}%
\pgfpathclose%
\pgfpathmoveto{\pgfqpoint{7.644676in}{0.406787in}}%
\pgfpathlineto{\pgfqpoint{7.644676in}{0.396989in}}%
\pgfpathquadraticcurveto{\pgfqpoint{7.640299in}{0.399240in}}{\pgfqpoint{7.635562in}{0.400357in}}%
\pgfpathquadraticcurveto{\pgfqpoint{7.630843in}{0.401492in}}{\pgfqpoint{7.625764in}{0.401492in}}%
\pgfpathquadraticcurveto{\pgfqpoint{7.618054in}{0.401492in}}{\pgfqpoint{7.614200in}{0.399132in}}%
\pgfpathquadraticcurveto{\pgfqpoint{7.610345in}{0.396773in}}{\pgfqpoint{7.610345in}{0.392035in}}%
\pgfpathquadraticcurveto{\pgfqpoint{7.610345in}{0.388433in}}{\pgfqpoint{7.613101in}{0.386380in}}%
\pgfpathquadraticcurveto{\pgfqpoint{7.615857in}{0.384326in}}{\pgfqpoint{7.624197in}{0.382471in}}%
\pgfpathlineto{\pgfqpoint{7.627745in}{0.381678in}}%
\pgfpathquadraticcurveto{\pgfqpoint{7.638768in}{0.379319in}}{\pgfqpoint{7.643416in}{0.375014in}}%
\pgfpathquadraticcurveto{\pgfqpoint{7.648063in}{0.370709in}}{\pgfqpoint{7.648063in}{0.363000in}}%
\pgfpathquadraticcurveto{\pgfqpoint{7.648063in}{0.354210in}}{\pgfqpoint{7.641110in}{0.349076in}}%
\pgfpathquadraticcurveto{\pgfqpoint{7.634157in}{0.343961in}}{\pgfqpoint{7.621999in}{0.343961in}}%
\pgfpathquadraticcurveto{\pgfqpoint{7.616938in}{0.343961in}}{\pgfqpoint{7.611444in}{0.344952in}}%
\pgfpathquadraticcurveto{\pgfqpoint{7.605950in}{0.345942in}}{\pgfqpoint{7.599880in}{0.347906in}}%
\pgfpathlineto{\pgfqpoint{7.599880in}{0.358605in}}%
\pgfpathquadraticcurveto{\pgfqpoint{7.605626in}{0.355615in}}{\pgfqpoint{7.611192in}{0.354120in}}%
\pgfpathquadraticcurveto{\pgfqpoint{7.616758in}{0.352643in}}{\pgfqpoint{7.622233in}{0.352643in}}%
\pgfpathquadraticcurveto{\pgfqpoint{7.629546in}{0.352643in}}{\pgfqpoint{7.633473in}{0.355146in}}%
\pgfpathquadraticcurveto{\pgfqpoint{7.637418in}{0.357650in}}{\pgfqpoint{7.637418in}{0.362207in}}%
\pgfpathquadraticcurveto{\pgfqpoint{7.637418in}{0.366422in}}{\pgfqpoint{7.634572in}{0.368674in}}%
\pgfpathquadraticcurveto{\pgfqpoint{7.631744in}{0.370925in}}{\pgfqpoint{7.622107in}{0.373014in}}%
\pgfpathlineto{\pgfqpoint{7.618505in}{0.373861in}}%
\pgfpathquadraticcurveto{\pgfqpoint{7.608886in}{0.375878in}}{\pgfqpoint{7.604599in}{0.380075in}}%
\pgfpathquadraticcurveto{\pgfqpoint{7.600331in}{0.384272in}}{\pgfqpoint{7.600331in}{0.391585in}}%
\pgfpathquadraticcurveto{\pgfqpoint{7.600331in}{0.400483in}}{\pgfqpoint{7.606635in}{0.405310in}}%
\pgfpathquadraticcurveto{\pgfqpoint{7.612939in}{0.410156in}}{\pgfqpoint{7.624539in}{0.410156in}}%
\pgfpathquadraticcurveto{\pgfqpoint{7.630267in}{0.410156in}}{\pgfqpoint{7.635328in}{0.409309in}}%
\pgfpathquadraticcurveto{\pgfqpoint{7.640408in}{0.408480in}}{\pgfqpoint{7.644676in}{0.406787in}}%
\pgfpathlineto{\pgfqpoint{7.644676in}{0.406787in}}%
\pgfpathclose%
\pgfpathmoveto{\pgfqpoint{7.674793in}{0.426547in}}%
\pgfpathlineto{\pgfqpoint{7.674793in}{0.408643in}}%
\pgfpathlineto{\pgfqpoint{7.696119in}{0.408643in}}%
\pgfpathlineto{\pgfqpoint{7.696119in}{0.400591in}}%
\pgfpathlineto{\pgfqpoint{7.674793in}{0.400591in}}%
\pgfpathlineto{\pgfqpoint{7.674793in}{0.366368in}}%
\pgfpathquadraticcurveto{\pgfqpoint{7.674793in}{0.358659in}}{\pgfqpoint{7.676900in}{0.356461in}}%
\pgfpathquadraticcurveto{\pgfqpoint{7.679008in}{0.354264in}}{\pgfqpoint{7.685492in}{0.354264in}}%
\pgfpathlineto{\pgfqpoint{7.696119in}{0.354264in}}%
\pgfpathlineto{\pgfqpoint{7.696119in}{0.345600in}}%
\pgfpathlineto{\pgfqpoint{7.685492in}{0.345600in}}%
\pgfpathquadraticcurveto{\pgfqpoint{7.673496in}{0.345600in}}{\pgfqpoint{7.668939in}{0.350067in}}%
\pgfpathquadraticcurveto{\pgfqpoint{7.664382in}{0.354552in}}{\pgfqpoint{7.664382in}{0.366368in}}%
\pgfpathlineto{\pgfqpoint{7.664382in}{0.400591in}}%
\pgfpathlineto{\pgfqpoint{7.656781in}{0.400591in}}%
\pgfpathlineto{\pgfqpoint{7.656781in}{0.408643in}}%
\pgfpathlineto{\pgfqpoint{7.664382in}{0.408643in}}%
\pgfpathlineto{\pgfqpoint{7.664382in}{0.426547in}}%
\pgfpathlineto{\pgfqpoint{7.674793in}{0.426547in}}%
\pgfpathlineto{\pgfqpoint{7.674793in}{0.426547in}}%
\pgfpathclose%
\pgfpathmoveto{\pgfqpoint{7.763665in}{0.379715in}}%
\pgfpathlineto{\pgfqpoint{7.763665in}{0.374654in}}%
\pgfpathlineto{\pgfqpoint{7.716041in}{0.374654in}}%
\pgfpathquadraticcurveto{\pgfqpoint{7.716725in}{0.363954in}}{\pgfqpoint{7.722489in}{0.358353in}}%
\pgfpathquadraticcurveto{\pgfqpoint{7.728253in}{0.352751in}}{\pgfqpoint{7.738556in}{0.352751in}}%
\pgfpathquadraticcurveto{\pgfqpoint{7.744518in}{0.352751in}}{\pgfqpoint{7.750120in}{0.354210in}}%
\pgfpathquadraticcurveto{\pgfqpoint{7.755721in}{0.355669in}}{\pgfqpoint{7.761251in}{0.358605in}}%
\pgfpathlineto{\pgfqpoint{7.761251in}{0.348806in}}%
\pgfpathquadraticcurveto{\pgfqpoint{7.755667in}{0.346447in}}{\pgfqpoint{7.749813in}{0.345204in}}%
\pgfpathquadraticcurveto{\pgfqpoint{7.743959in}{0.343961in}}{\pgfqpoint{7.737943in}{0.343961in}}%
\pgfpathquadraticcurveto{\pgfqpoint{7.722849in}{0.343961in}}{\pgfqpoint{7.714041in}{0.352733in}}%
\pgfpathquadraticcurveto{\pgfqpoint{7.705233in}{0.361523in}}{\pgfqpoint{7.705233in}{0.376509in}}%
\pgfpathquadraticcurveto{\pgfqpoint{7.705233in}{0.391981in}}{\pgfqpoint{7.713591in}{0.401059in}}%
\pgfpathquadraticcurveto{\pgfqpoint{7.721949in}{0.410156in}}{\pgfqpoint{7.736142in}{0.410156in}}%
\pgfpathquadraticcurveto{\pgfqpoint{7.748859in}{0.410156in}}{\pgfqpoint{7.756262in}{0.401960in}}%
\pgfpathquadraticcurveto{\pgfqpoint{7.763665in}{0.393783in}}{\pgfqpoint{7.763665in}{0.379715in}}%
\pgfpathlineto{\pgfqpoint{7.763665in}{0.379715in}}%
\pgfpathclose%
\pgfpathmoveto{\pgfqpoint{7.753308in}{0.382759in}}%
\pgfpathquadraticcurveto{\pgfqpoint{7.753200in}{0.391243in}}{\pgfqpoint{7.748553in}{0.396304in}}%
\pgfpathquadraticcurveto{\pgfqpoint{7.743905in}{0.401384in}}{\pgfqpoint{7.736250in}{0.401384in}}%
\pgfpathquadraticcurveto{\pgfqpoint{7.727586in}{0.401384in}}{\pgfqpoint{7.722381in}{0.396484in}}%
\pgfpathquadraticcurveto{\pgfqpoint{7.717175in}{0.391585in}}{\pgfqpoint{7.716383in}{0.382687in}}%
\pgfpathlineto{\pgfqpoint{7.753308in}{0.382759in}}%
\pgfpathlineto{\pgfqpoint{7.753308in}{0.382759in}}%
\pgfpathclose%
\pgfpathmoveto{\pgfqpoint{7.817197in}{0.398970in}}%
\pgfpathquadraticcurveto{\pgfqpoint{7.815450in}{0.399979in}}{\pgfqpoint{7.813396in}{0.400447in}}%
\pgfpathquadraticcurveto{\pgfqpoint{7.811343in}{0.400933in}}{\pgfqpoint{7.808875in}{0.400933in}}%
\pgfpathquadraticcurveto{\pgfqpoint{7.800085in}{0.400933in}}{\pgfqpoint{7.795384in}{0.395223in}}%
\pgfpathquadraticcurveto{\pgfqpoint{7.790683in}{0.389514in}}{\pgfqpoint{7.790683in}{0.378814in}}%
\pgfpathlineto{\pgfqpoint{7.790683in}{0.345600in}}%
\pgfpathlineto{\pgfqpoint{7.780272in}{0.345600in}}%
\pgfpathlineto{\pgfqpoint{7.780272in}{0.408643in}}%
\pgfpathlineto{\pgfqpoint{7.790683in}{0.408643in}}%
\pgfpathlineto{\pgfqpoint{7.790683in}{0.398844in}}%
\pgfpathquadraticcurveto{\pgfqpoint{7.793961in}{0.404590in}}{\pgfqpoint{7.799185in}{0.407364in}}%
\pgfpathquadraticcurveto{\pgfqpoint{7.804426in}{0.410156in}}{\pgfqpoint{7.811919in}{0.410156in}}%
\pgfpathquadraticcurveto{\pgfqpoint{7.812982in}{0.410156in}}{\pgfqpoint{7.814279in}{0.410011in}}%
\pgfpathquadraticcurveto{\pgfqpoint{7.815576in}{0.409885in}}{\pgfqpoint{7.817143in}{0.409597in}}%
\pgfpathlineto{\pgfqpoint{7.817197in}{0.398970in}}%
\pgfpathlineto{\pgfqpoint{7.817197in}{0.398970in}}%
\pgfpathclose%
\pgfpathmoveto{\pgfqpoint{7.868243in}{0.406787in}}%
\pgfpathlineto{\pgfqpoint{7.868243in}{0.396989in}}%
\pgfpathquadraticcurveto{\pgfqpoint{7.863866in}{0.399240in}}{\pgfqpoint{7.859129in}{0.400357in}}%
\pgfpathquadraticcurveto{\pgfqpoint{7.854410in}{0.401492in}}{\pgfqpoint{7.849331in}{0.401492in}}%
\pgfpathquadraticcurveto{\pgfqpoint{7.841621in}{0.401492in}}{\pgfqpoint{7.837767in}{0.399132in}}%
\pgfpathquadraticcurveto{\pgfqpoint{7.833912in}{0.396773in}}{\pgfqpoint{7.833912in}{0.392035in}}%
\pgfpathquadraticcurveto{\pgfqpoint{7.833912in}{0.388433in}}{\pgfqpoint{7.836668in}{0.386380in}}%
\pgfpathquadraticcurveto{\pgfqpoint{7.839424in}{0.384326in}}{\pgfqpoint{7.847763in}{0.382471in}}%
\pgfpathlineto{\pgfqpoint{7.851312in}{0.381678in}}%
\pgfpathquadraticcurveto{\pgfqpoint{7.862335in}{0.379319in}}{\pgfqpoint{7.866982in}{0.375014in}}%
\pgfpathquadraticcurveto{\pgfqpoint{7.871630in}{0.370709in}}{\pgfqpoint{7.871630in}{0.363000in}}%
\pgfpathquadraticcurveto{\pgfqpoint{7.871630in}{0.354210in}}{\pgfqpoint{7.864677in}{0.349076in}}%
\pgfpathquadraticcurveto{\pgfqpoint{7.857724in}{0.343961in}}{\pgfqpoint{7.845566in}{0.343961in}}%
\pgfpathquadraticcurveto{\pgfqpoint{7.840505in}{0.343961in}}{\pgfqpoint{7.835011in}{0.344952in}}%
\pgfpathquadraticcurveto{\pgfqpoint{7.829517in}{0.345942in}}{\pgfqpoint{7.823447in}{0.347906in}}%
\pgfpathlineto{\pgfqpoint{7.823447in}{0.358605in}}%
\pgfpathquadraticcurveto{\pgfqpoint{7.829193in}{0.355615in}}{\pgfqpoint{7.834759in}{0.354120in}}%
\pgfpathquadraticcurveto{\pgfqpoint{7.840324in}{0.352643in}}{\pgfqpoint{7.845800in}{0.352643in}}%
\pgfpathquadraticcurveto{\pgfqpoint{7.853113in}{0.352643in}}{\pgfqpoint{7.857040in}{0.355146in}}%
\pgfpathquadraticcurveto{\pgfqpoint{7.860984in}{0.357650in}}{\pgfqpoint{7.860984in}{0.362207in}}%
\pgfpathquadraticcurveto{\pgfqpoint{7.860984in}{0.366422in}}{\pgfqpoint{7.858138in}{0.368674in}}%
\pgfpathquadraticcurveto{\pgfqpoint{7.855311in}{0.370925in}}{\pgfqpoint{7.845674in}{0.373014in}}%
\pgfpathlineto{\pgfqpoint{7.842072in}{0.373861in}}%
\pgfpathquadraticcurveto{\pgfqpoint{7.832453in}{0.375878in}}{\pgfqpoint{7.828166in}{0.380075in}}%
\pgfpathquadraticcurveto{\pgfqpoint{7.823897in}{0.384272in}}{\pgfqpoint{7.823897in}{0.391585in}}%
\pgfpathquadraticcurveto{\pgfqpoint{7.823897in}{0.400483in}}{\pgfqpoint{7.830202in}{0.405310in}}%
\pgfpathquadraticcurveto{\pgfqpoint{7.836506in}{0.410156in}}{\pgfqpoint{7.848106in}{0.410156in}}%
\pgfpathquadraticcurveto{\pgfqpoint{7.853834in}{0.410156in}}{\pgfqpoint{7.858895in}{0.409309in}}%
\pgfpathquadraticcurveto{\pgfqpoint{7.863974in}{0.408480in}}{\pgfqpoint{7.868243in}{0.406787in}}%
\pgfpathlineto{\pgfqpoint{7.868243in}{0.406787in}}%
\pgfpathclose%
\pgfpathmoveto{\pgfqpoint{7.928188in}{0.355164in}}%
\pgfpathlineto{\pgfqpoint{7.946758in}{0.355164in}}%
\pgfpathlineto{\pgfqpoint{7.946758in}{0.419288in}}%
\pgfpathlineto{\pgfqpoint{7.926549in}{0.415235in}}%
\pgfpathlineto{\pgfqpoint{7.926549in}{0.425592in}}%
\pgfpathlineto{\pgfqpoint{7.946650in}{0.429645in}}%
\pgfpathlineto{\pgfqpoint{7.958016in}{0.429645in}}%
\pgfpathlineto{\pgfqpoint{7.958016in}{0.355164in}}%
\pgfpathlineto{\pgfqpoint{7.976586in}{0.355164in}}%
\pgfpathlineto{\pgfqpoint{7.976586in}{0.345600in}}%
\pgfpathlineto{\pgfqpoint{7.928188in}{0.345600in}}%
\pgfpathlineto{\pgfqpoint{7.928188in}{0.355164in}}%
\pgfpathlineto{\pgfqpoint{7.928188in}{0.355164in}}%
\pgfpathclose%
\pgfpathmoveto{\pgfqpoint{8.030803in}{0.419738in}}%
\pgfpathlineto{\pgfqpoint{8.002092in}{0.374870in}}%
\pgfpathlineto{\pgfqpoint{8.030803in}{0.374870in}}%
\pgfpathlineto{\pgfqpoint{8.030803in}{0.419738in}}%
\pgfpathlineto{\pgfqpoint{8.030803in}{0.419738in}}%
\pgfpathclose%
\pgfpathmoveto{\pgfqpoint{8.027813in}{0.429645in}}%
\pgfpathlineto{\pgfqpoint{8.042115in}{0.429645in}}%
\pgfpathlineto{\pgfqpoint{8.042115in}{0.374870in}}%
\pgfpathlineto{\pgfqpoint{8.054111in}{0.374870in}}%
\pgfpathlineto{\pgfqpoint{8.054111in}{0.365413in}}%
\pgfpathlineto{\pgfqpoint{8.042115in}{0.365413in}}%
\pgfpathlineto{\pgfqpoint{8.042115in}{0.345600in}}%
\pgfpathlineto{\pgfqpoint{8.030803in}{0.345600in}}%
\pgfpathlineto{\pgfqpoint{8.030803in}{0.365413in}}%
\pgfpathlineto{\pgfqpoint{7.992869in}{0.365413in}}%
\pgfpathlineto{\pgfqpoint{7.992869in}{0.376383in}}%
\pgfpathlineto{\pgfqpoint{8.027813in}{0.429645in}}%
\pgfpathlineto{\pgfqpoint{8.027813in}{0.429645in}}%
\pgfpathclose%
\pgfpathmoveto{\pgfqpoint{8.107355in}{0.390919in}}%
\pgfpathquadraticcurveto{\pgfqpoint{8.115514in}{0.389171in}}{\pgfqpoint{8.120089in}{0.383642in}}%
\pgfpathquadraticcurveto{\pgfqpoint{8.124682in}{0.378130in}}{\pgfqpoint{8.124682in}{0.370024in}}%
\pgfpathquadraticcurveto{\pgfqpoint{8.124682in}{0.357596in}}{\pgfqpoint{8.116127in}{0.350769in}}%
\pgfpathquadraticcurveto{\pgfqpoint{8.107571in}{0.343961in}}{\pgfqpoint{8.091810in}{0.343961in}}%
\pgfpathquadraticcurveto{\pgfqpoint{8.086533in}{0.343961in}}{\pgfqpoint{8.080931in}{0.345006in}}%
\pgfpathquadraticcurveto{\pgfqpoint{8.075329in}{0.346050in}}{\pgfqpoint{8.069367in}{0.348140in}}%
\pgfpathlineto{\pgfqpoint{8.069367in}{0.359109in}}%
\pgfpathquadraticcurveto{\pgfqpoint{8.074086in}{0.356353in}}{\pgfqpoint{8.079706in}{0.354948in}}%
\pgfpathquadraticcurveto{\pgfqpoint{8.085344in}{0.353543in}}{\pgfqpoint{8.091486in}{0.353543in}}%
\pgfpathquadraticcurveto{\pgfqpoint{8.102167in}{0.353543in}}{\pgfqpoint{8.107769in}{0.357758in}}%
\pgfpathquadraticcurveto{\pgfqpoint{8.113371in}{0.361973in}}{\pgfqpoint{8.113371in}{0.370024in}}%
\pgfpathquadraticcurveto{\pgfqpoint{8.113371in}{0.377463in}}{\pgfqpoint{8.108165in}{0.381642in}}%
\pgfpathquadraticcurveto{\pgfqpoint{8.102960in}{0.385839in}}{\pgfqpoint{8.093683in}{0.385839in}}%
\pgfpathlineto{\pgfqpoint{8.083885in}{0.385839in}}%
\pgfpathlineto{\pgfqpoint{8.083885in}{0.395187in}}%
\pgfpathlineto{\pgfqpoint{8.094134in}{0.395187in}}%
\pgfpathquadraticcurveto{\pgfqpoint{8.102509in}{0.395187in}}{\pgfqpoint{8.106958in}{0.398538in}}%
\pgfpathquadraticcurveto{\pgfqpoint{8.111407in}{0.401888in}}{\pgfqpoint{8.111407in}{0.408192in}}%
\pgfpathquadraticcurveto{\pgfqpoint{8.111407in}{0.414659in}}{\pgfqpoint{8.106814in}{0.418117in}}%
\pgfpathquadraticcurveto{\pgfqpoint{8.102239in}{0.421593in}}{\pgfqpoint{8.093683in}{0.421593in}}%
\pgfpathquadraticcurveto{\pgfqpoint{8.089000in}{0.421593in}}{\pgfqpoint{8.083651in}{0.420567in}}%
\pgfpathquadraticcurveto{\pgfqpoint{8.078301in}{0.419558in}}{\pgfqpoint{8.071889in}{0.417432in}}%
\pgfpathlineto{\pgfqpoint{8.071889in}{0.427555in}}%
\pgfpathquadraticcurveto{\pgfqpoint{8.078373in}{0.429357in}}{\pgfqpoint{8.084029in}{0.430257in}}%
\pgfpathquadraticcurveto{\pgfqpoint{8.089685in}{0.431158in}}{\pgfqpoint{8.094692in}{0.431158in}}%
\pgfpathquadraticcurveto{\pgfqpoint{8.107643in}{0.431158in}}{\pgfqpoint{8.115172in}{0.425268in}}%
\pgfpathquadraticcurveto{\pgfqpoint{8.122719in}{0.419396in}}{\pgfqpoint{8.122719in}{0.409381in}}%
\pgfpathquadraticcurveto{\pgfqpoint{8.122719in}{0.402392in}}{\pgfqpoint{8.118720in}{0.397583in}}%
\pgfpathquadraticcurveto{\pgfqpoint{8.114722in}{0.392774in}}{\pgfqpoint{8.107355in}{0.390919in}}%
\pgfpathlineto{\pgfqpoint{8.107355in}{0.390919in}}%
\pgfpathclose%
\pgfpathmoveto{\pgfqpoint{8.226685in}{0.377067in}}%
\pgfpathquadraticcurveto{\pgfqpoint{8.226685in}{0.388487in}}{\pgfqpoint{8.221984in}{0.394989in}}%
\pgfpathquadraticcurveto{\pgfqpoint{8.217283in}{0.401492in}}{\pgfqpoint{8.209069in}{0.401492in}}%
\pgfpathquadraticcurveto{\pgfqpoint{8.200838in}{0.401492in}}{\pgfqpoint{8.196137in}{0.394989in}}%
\pgfpathquadraticcurveto{\pgfqpoint{8.191435in}{0.388487in}}{\pgfqpoint{8.191435in}{0.377067in}}%
\pgfpathquadraticcurveto{\pgfqpoint{8.191435in}{0.365648in}}{\pgfqpoint{8.196137in}{0.359145in}}%
\pgfpathquadraticcurveto{\pgfqpoint{8.200838in}{0.352643in}}{\pgfqpoint{8.209069in}{0.352643in}}%
\pgfpathquadraticcurveto{\pgfqpoint{8.217283in}{0.352643in}}{\pgfqpoint{8.221984in}{0.359145in}}%
\pgfpathquadraticcurveto{\pgfqpoint{8.226685in}{0.365648in}}{\pgfqpoint{8.226685in}{0.377067in}}%
\pgfpathlineto{\pgfqpoint{8.226685in}{0.377067in}}%
\pgfpathclose%
\pgfpathmoveto{\pgfqpoint{8.191435in}{0.399078in}}%
\pgfpathquadraticcurveto{\pgfqpoint{8.194714in}{0.404698in}}{\pgfqpoint{8.199685in}{0.407418in}}%
\pgfpathquadraticcurveto{\pgfqpoint{8.204674in}{0.410156in}}{\pgfqpoint{8.211591in}{0.410156in}}%
\pgfpathquadraticcurveto{\pgfqpoint{8.223083in}{0.410156in}}{\pgfqpoint{8.230252in}{0.401041in}}%
\pgfpathquadraticcurveto{\pgfqpoint{8.237438in}{0.391927in}}{\pgfqpoint{8.237438in}{0.377067in}}%
\pgfpathquadraticcurveto{\pgfqpoint{8.237438in}{0.362207in}}{\pgfqpoint{8.230252in}{0.353075in}}%
\pgfpathquadraticcurveto{\pgfqpoint{8.223083in}{0.343961in}}{\pgfqpoint{8.211591in}{0.343961in}}%
\pgfpathquadraticcurveto{\pgfqpoint{8.204674in}{0.343961in}}{\pgfqpoint{8.199685in}{0.346699in}}%
\pgfpathquadraticcurveto{\pgfqpoint{8.194714in}{0.349437in}}{\pgfqpoint{8.191435in}{0.355056in}}%
\pgfpathlineto{\pgfqpoint{8.191435in}{0.345600in}}%
\pgfpathlineto{\pgfqpoint{8.181024in}{0.345600in}}%
\pgfpathlineto{\pgfqpoint{8.181024in}{0.433193in}}%
\pgfpathlineto{\pgfqpoint{8.191435in}{0.433193in}}%
\pgfpathlineto{\pgfqpoint{8.191435in}{0.399078in}}%
\pgfpathlineto{\pgfqpoint{8.191435in}{0.399078in}}%
\pgfpathclose%
\pgfpathmoveto{\pgfqpoint{8.253541in}{0.370475in}}%
\pgfpathlineto{\pgfqpoint{8.253541in}{0.408643in}}%
\pgfpathlineto{\pgfqpoint{8.263898in}{0.408643in}}%
\pgfpathlineto{\pgfqpoint{8.263898in}{0.370871in}}%
\pgfpathquadraticcurveto{\pgfqpoint{8.263898in}{0.361919in}}{\pgfqpoint{8.267375in}{0.357434in}}%
\pgfpathquadraticcurveto{\pgfqpoint{8.270869in}{0.352967in}}{\pgfqpoint{8.277858in}{0.352967in}}%
\pgfpathquadraticcurveto{\pgfqpoint{8.286233in}{0.352967in}}{\pgfqpoint{8.291097in}{0.358317in}}%
\pgfpathquadraticcurveto{\pgfqpoint{8.295978in}{0.363666in}}{\pgfqpoint{8.295978in}{0.372906in}}%
\pgfpathlineto{\pgfqpoint{8.295978in}{0.408643in}}%
\pgfpathlineto{\pgfqpoint{8.306335in}{0.408643in}}%
\pgfpathlineto{\pgfqpoint{8.306335in}{0.345600in}}%
\pgfpathlineto{\pgfqpoint{8.295978in}{0.345600in}}%
\pgfpathlineto{\pgfqpoint{8.295978in}{0.355291in}}%
\pgfpathquadraticcurveto{\pgfqpoint{8.292213in}{0.349545in}}{\pgfqpoint{8.287224in}{0.346753in}}%
\pgfpathquadraticcurveto{\pgfqpoint{8.282253in}{0.343961in}}{\pgfqpoint{8.275660in}{0.343961in}}%
\pgfpathquadraticcurveto{\pgfqpoint{8.264799in}{0.343961in}}{\pgfqpoint{8.259161in}{0.350715in}}%
\pgfpathquadraticcurveto{\pgfqpoint{8.253541in}{0.357470in}}{\pgfqpoint{8.253541in}{0.370475in}}%
\pgfpathlineto{\pgfqpoint{8.253541in}{0.370475in}}%
\pgfpathclose%
\pgfpathmoveto{\pgfqpoint{8.279605in}{0.410156in}}%
\pgfpathlineto{\pgfqpoint{8.279605in}{0.410156in}}%
\pgfpathlineto{\pgfqpoint{8.279605in}{0.410156in}}%
\pgfpathclose%
\pgfpathmoveto{\pgfqpoint{8.327661in}{0.408643in}}%
\pgfpathlineto{\pgfqpoint{8.338018in}{0.408643in}}%
\pgfpathlineto{\pgfqpoint{8.338018in}{0.345600in}}%
\pgfpathlineto{\pgfqpoint{8.327661in}{0.345600in}}%
\pgfpathlineto{\pgfqpoint{8.327661in}{0.408643in}}%
\pgfpathlineto{\pgfqpoint{8.327661in}{0.408643in}}%
\pgfpathclose%
\pgfpathmoveto{\pgfqpoint{8.327661in}{0.433193in}}%
\pgfpathlineto{\pgfqpoint{8.338018in}{0.433193in}}%
\pgfpathlineto{\pgfqpoint{8.338018in}{0.420062in}}%
\pgfpathlineto{\pgfqpoint{8.327661in}{0.420062in}}%
\pgfpathlineto{\pgfqpoint{8.327661in}{0.433193in}}%
\pgfpathlineto{\pgfqpoint{8.327661in}{0.433193in}}%
\pgfpathclose%
\pgfpathmoveto{\pgfqpoint{8.359687in}{0.433193in}}%
\pgfpathlineto{\pgfqpoint{8.370044in}{0.433193in}}%
\pgfpathlineto{\pgfqpoint{8.370044in}{0.345600in}}%
\pgfpathlineto{\pgfqpoint{8.359687in}{0.345600in}}%
\pgfpathlineto{\pgfqpoint{8.359687in}{0.433193in}}%
\pgfpathlineto{\pgfqpoint{8.359687in}{0.433193in}}%
\pgfpathclose%
\pgfpathmoveto{\pgfqpoint{8.433194in}{0.399078in}}%
\pgfpathlineto{\pgfqpoint{8.433194in}{0.433193in}}%
\pgfpathlineto{\pgfqpoint{8.443551in}{0.433193in}}%
\pgfpathlineto{\pgfqpoint{8.443551in}{0.345600in}}%
\pgfpathlineto{\pgfqpoint{8.433194in}{0.345600in}}%
\pgfpathlineto{\pgfqpoint{8.433194in}{0.355056in}}%
\pgfpathquadraticcurveto{\pgfqpoint{8.429934in}{0.349437in}}{\pgfqpoint{8.424945in}{0.346699in}}%
\pgfpathquadraticcurveto{\pgfqpoint{8.419974in}{0.343961in}}{\pgfqpoint{8.412985in}{0.343961in}}%
\pgfpathquadraticcurveto{\pgfqpoint{8.401565in}{0.343961in}}{\pgfqpoint{8.394378in}{0.353075in}}%
\pgfpathquadraticcurveto{\pgfqpoint{8.387209in}{0.362207in}}{\pgfqpoint{8.387209in}{0.377067in}}%
\pgfpathquadraticcurveto{\pgfqpoint{8.387209in}{0.391927in}}{\pgfqpoint{8.394378in}{0.401041in}}%
\pgfpathquadraticcurveto{\pgfqpoint{8.401565in}{0.410156in}}{\pgfqpoint{8.412985in}{0.410156in}}%
\pgfpathquadraticcurveto{\pgfqpoint{8.419974in}{0.410156in}}{\pgfqpoint{8.424945in}{0.407418in}}%
\pgfpathquadraticcurveto{\pgfqpoint{8.429934in}{0.404698in}}{\pgfqpoint{8.433194in}{0.399078in}}%
\pgfpathlineto{\pgfqpoint{8.433194in}{0.399078in}}%
\pgfpathclose%
\pgfpathmoveto{\pgfqpoint{8.397909in}{0.377067in}}%
\pgfpathquadraticcurveto{\pgfqpoint{8.397909in}{0.365648in}}{\pgfqpoint{8.402610in}{0.359145in}}%
\pgfpathquadraticcurveto{\pgfqpoint{8.407311in}{0.352643in}}{\pgfqpoint{8.415525in}{0.352643in}}%
\pgfpathquadraticcurveto{\pgfqpoint{8.423738in}{0.352643in}}{\pgfqpoint{8.428457in}{0.359145in}}%
\pgfpathquadraticcurveto{\pgfqpoint{8.433194in}{0.365648in}}{\pgfqpoint{8.433194in}{0.377067in}}%
\pgfpathquadraticcurveto{\pgfqpoint{8.433194in}{0.388487in}}{\pgfqpoint{8.428457in}{0.394989in}}%
\pgfpathquadraticcurveto{\pgfqpoint{8.423738in}{0.401492in}}{\pgfqpoint{8.415525in}{0.401492in}}%
\pgfpathquadraticcurveto{\pgfqpoint{8.407311in}{0.401492in}}{\pgfqpoint{8.402610in}{0.394989in}}%
\pgfpathquadraticcurveto{\pgfqpoint{8.397909in}{0.388487in}}{\pgfqpoint{8.397909in}{0.377067in}}%
\pgfpathlineto{\pgfqpoint{8.397909in}{0.377067in}}%
\pgfpathclose%
\pgfpathmoveto{\pgfqpoint{8.518824in}{0.379715in}}%
\pgfpathlineto{\pgfqpoint{8.518824in}{0.374654in}}%
\pgfpathlineto{\pgfqpoint{8.471200in}{0.374654in}}%
\pgfpathquadraticcurveto{\pgfqpoint{8.471885in}{0.363954in}}{\pgfqpoint{8.477648in}{0.358353in}}%
\pgfpathquadraticcurveto{\pgfqpoint{8.483412in}{0.352751in}}{\pgfqpoint{8.493715in}{0.352751in}}%
\pgfpathquadraticcurveto{\pgfqpoint{8.499677in}{0.352751in}}{\pgfqpoint{8.505279in}{0.354210in}}%
\pgfpathquadraticcurveto{\pgfqpoint{8.510881in}{0.355669in}}{\pgfqpoint{8.516411in}{0.358605in}}%
\pgfpathlineto{\pgfqpoint{8.516411in}{0.348806in}}%
\pgfpathquadraticcurveto{\pgfqpoint{8.510827in}{0.346447in}}{\pgfqpoint{8.504973in}{0.345204in}}%
\pgfpathquadraticcurveto{\pgfqpoint{8.499119in}{0.343961in}}{\pgfqpoint{8.493103in}{0.343961in}}%
\pgfpathquadraticcurveto{\pgfqpoint{8.478009in}{0.343961in}}{\pgfqpoint{8.469201in}{0.352733in}}%
\pgfpathquadraticcurveto{\pgfqpoint{8.460393in}{0.361523in}}{\pgfqpoint{8.460393in}{0.376509in}}%
\pgfpathquadraticcurveto{\pgfqpoint{8.460393in}{0.391981in}}{\pgfqpoint{8.468750in}{0.401059in}}%
\pgfpathquadraticcurveto{\pgfqpoint{8.477108in}{0.410156in}}{\pgfqpoint{8.491302in}{0.410156in}}%
\pgfpathquadraticcurveto{\pgfqpoint{8.504018in}{0.410156in}}{\pgfqpoint{8.511421in}{0.401960in}}%
\pgfpathquadraticcurveto{\pgfqpoint{8.518824in}{0.393783in}}{\pgfqpoint{8.518824in}{0.379715in}}%
\pgfpathlineto{\pgfqpoint{8.518824in}{0.379715in}}%
\pgfpathclose%
\pgfpathmoveto{\pgfqpoint{8.508467in}{0.382759in}}%
\pgfpathquadraticcurveto{\pgfqpoint{8.508359in}{0.391243in}}{\pgfqpoint{8.503712in}{0.396304in}}%
\pgfpathquadraticcurveto{\pgfqpoint{8.499065in}{0.401384in}}{\pgfqpoint{8.491410in}{0.401384in}}%
\pgfpathquadraticcurveto{\pgfqpoint{8.482746in}{0.401384in}}{\pgfqpoint{8.477540in}{0.396484in}}%
\pgfpathquadraticcurveto{\pgfqpoint{8.472335in}{0.391585in}}{\pgfqpoint{8.471542in}{0.382687in}}%
\pgfpathlineto{\pgfqpoint{8.508467in}{0.382759in}}%
\pgfpathlineto{\pgfqpoint{8.508467in}{0.382759in}}%
\pgfpathclose%
\pgfpathmoveto{\pgfqpoint{8.572356in}{0.398970in}}%
\pgfpathquadraticcurveto{\pgfqpoint{8.570609in}{0.399979in}}{\pgfqpoint{8.568556in}{0.400447in}}%
\pgfpathquadraticcurveto{\pgfqpoint{8.566502in}{0.400933in}}{\pgfqpoint{8.564035in}{0.400933in}}%
\pgfpathquadraticcurveto{\pgfqpoint{8.555245in}{0.400933in}}{\pgfqpoint{8.550544in}{0.395223in}}%
\pgfpathquadraticcurveto{\pgfqpoint{8.545842in}{0.389514in}}{\pgfqpoint{8.545842in}{0.378814in}}%
\pgfpathlineto{\pgfqpoint{8.545842in}{0.345600in}}%
\pgfpathlineto{\pgfqpoint{8.535431in}{0.345600in}}%
\pgfpathlineto{\pgfqpoint{8.535431in}{0.408643in}}%
\pgfpathlineto{\pgfqpoint{8.545842in}{0.408643in}}%
\pgfpathlineto{\pgfqpoint{8.545842in}{0.398844in}}%
\pgfpathquadraticcurveto{\pgfqpoint{8.549121in}{0.404590in}}{\pgfqpoint{8.554344in}{0.407364in}}%
\pgfpathquadraticcurveto{\pgfqpoint{8.559586in}{0.410156in}}{\pgfqpoint{8.567079in}{0.410156in}}%
\pgfpathquadraticcurveto{\pgfqpoint{8.568141in}{0.410156in}}{\pgfqpoint{8.569438in}{0.410011in}}%
\pgfpathquadraticcurveto{\pgfqpoint{8.570735in}{0.409885in}}{\pgfqpoint{8.572302in}{0.409597in}}%
\pgfpathlineto{\pgfqpoint{8.572356in}{0.398970in}}%
\pgfpathlineto{\pgfqpoint{8.572356in}{0.398970in}}%
\pgfpathclose%
\pgfpathmoveto{\pgfqpoint{8.623403in}{0.406787in}}%
\pgfpathlineto{\pgfqpoint{8.623403in}{0.396989in}}%
\pgfpathquadraticcurveto{\pgfqpoint{8.619026in}{0.399240in}}{\pgfqpoint{8.614289in}{0.400357in}}%
\pgfpathquadraticcurveto{\pgfqpoint{8.609569in}{0.401492in}}{\pgfqpoint{8.604490in}{0.401492in}}%
\pgfpathquadraticcurveto{\pgfqpoint{8.596781in}{0.401492in}}{\pgfqpoint{8.592926in}{0.399132in}}%
\pgfpathquadraticcurveto{\pgfqpoint{8.589072in}{0.396773in}}{\pgfqpoint{8.589072in}{0.392035in}}%
\pgfpathquadraticcurveto{\pgfqpoint{8.589072in}{0.388433in}}{\pgfqpoint{8.591827in}{0.386380in}}%
\pgfpathquadraticcurveto{\pgfqpoint{8.594583in}{0.384326in}}{\pgfqpoint{8.602923in}{0.382471in}}%
\pgfpathlineto{\pgfqpoint{8.606471in}{0.381678in}}%
\pgfpathquadraticcurveto{\pgfqpoint{8.617495in}{0.379319in}}{\pgfqpoint{8.622142in}{0.375014in}}%
\pgfpathquadraticcurveto{\pgfqpoint{8.626789in}{0.370709in}}{\pgfqpoint{8.626789in}{0.363000in}}%
\pgfpathquadraticcurveto{\pgfqpoint{8.626789in}{0.354210in}}{\pgfqpoint{8.619836in}{0.349076in}}%
\pgfpathquadraticcurveto{\pgfqpoint{8.612884in}{0.343961in}}{\pgfqpoint{8.600725in}{0.343961in}}%
\pgfpathquadraticcurveto{\pgfqpoint{8.595664in}{0.343961in}}{\pgfqpoint{8.590170in}{0.344952in}}%
\pgfpathquadraticcurveto{\pgfqpoint{8.584677in}{0.345942in}}{\pgfqpoint{8.578607in}{0.347906in}}%
\pgfpathlineto{\pgfqpoint{8.578607in}{0.358605in}}%
\pgfpathquadraticcurveto{\pgfqpoint{8.584352in}{0.355615in}}{\pgfqpoint{8.589918in}{0.354120in}}%
\pgfpathquadraticcurveto{\pgfqpoint{8.595484in}{0.352643in}}{\pgfqpoint{8.600960in}{0.352643in}}%
\pgfpathquadraticcurveto{\pgfqpoint{8.608273in}{0.352643in}}{\pgfqpoint{8.612199in}{0.355146in}}%
\pgfpathquadraticcurveto{\pgfqpoint{8.616144in}{0.357650in}}{\pgfqpoint{8.616144in}{0.362207in}}%
\pgfpathquadraticcurveto{\pgfqpoint{8.616144in}{0.366422in}}{\pgfqpoint{8.613298in}{0.368674in}}%
\pgfpathquadraticcurveto{\pgfqpoint{8.610470in}{0.370925in}}{\pgfqpoint{8.600834in}{0.373014in}}%
\pgfpathlineto{\pgfqpoint{8.597231in}{0.373861in}}%
\pgfpathquadraticcurveto{\pgfqpoint{8.587613in}{0.375878in}}{\pgfqpoint{8.583326in}{0.380075in}}%
\pgfpathquadraticcurveto{\pgfqpoint{8.579057in}{0.384272in}}{\pgfqpoint{8.579057in}{0.391585in}}%
\pgfpathquadraticcurveto{\pgfqpoint{8.579057in}{0.400483in}}{\pgfqpoint{8.585361in}{0.405310in}}%
\pgfpathquadraticcurveto{\pgfqpoint{8.591665in}{0.410156in}}{\pgfqpoint{8.603265in}{0.410156in}}%
\pgfpathquadraticcurveto{\pgfqpoint{8.608993in}{0.410156in}}{\pgfqpoint{8.614054in}{0.409309in}}%
\pgfpathquadraticcurveto{\pgfqpoint{8.619134in}{0.408480in}}{\pgfqpoint{8.623403in}{0.406787in}}%
\pgfpathlineto{\pgfqpoint{8.623403in}{0.406787in}}%
\pgfpathclose%
\pgfpathmoveto{\pgfqpoint{8.721389in}{0.399078in}}%
\pgfpathlineto{\pgfqpoint{8.721389in}{0.433193in}}%
\pgfpathlineto{\pgfqpoint{8.731746in}{0.433193in}}%
\pgfpathlineto{\pgfqpoint{8.731746in}{0.345600in}}%
\pgfpathlineto{\pgfqpoint{8.721389in}{0.345600in}}%
\pgfpathlineto{\pgfqpoint{8.721389in}{0.355056in}}%
\pgfpathquadraticcurveto{\pgfqpoint{8.718129in}{0.349437in}}{\pgfqpoint{8.713139in}{0.346699in}}%
\pgfpathquadraticcurveto{\pgfqpoint{8.708168in}{0.343961in}}{\pgfqpoint{8.701179in}{0.343961in}}%
\pgfpathquadraticcurveto{\pgfqpoint{8.689760in}{0.343961in}}{\pgfqpoint{8.682573in}{0.353075in}}%
\pgfpathquadraticcurveto{\pgfqpoint{8.675404in}{0.362207in}}{\pgfqpoint{8.675404in}{0.377067in}}%
\pgfpathquadraticcurveto{\pgfqpoint{8.675404in}{0.391927in}}{\pgfqpoint{8.682573in}{0.401041in}}%
\pgfpathquadraticcurveto{\pgfqpoint{8.689760in}{0.410156in}}{\pgfqpoint{8.701179in}{0.410156in}}%
\pgfpathquadraticcurveto{\pgfqpoint{8.708168in}{0.410156in}}{\pgfqpoint{8.713139in}{0.407418in}}%
\pgfpathquadraticcurveto{\pgfqpoint{8.718129in}{0.404698in}}{\pgfqpoint{8.721389in}{0.399078in}}%
\pgfpathlineto{\pgfqpoint{8.721389in}{0.399078in}}%
\pgfpathclose%
\pgfpathmoveto{\pgfqpoint{8.686103in}{0.377067in}}%
\pgfpathquadraticcurveto{\pgfqpoint{8.686103in}{0.365648in}}{\pgfqpoint{8.690804in}{0.359145in}}%
\pgfpathquadraticcurveto{\pgfqpoint{8.695505in}{0.352643in}}{\pgfqpoint{8.703719in}{0.352643in}}%
\pgfpathquadraticcurveto{\pgfqpoint{8.711933in}{0.352643in}}{\pgfqpoint{8.716652in}{0.359145in}}%
\pgfpathquadraticcurveto{\pgfqpoint{8.721389in}{0.365648in}}{\pgfqpoint{8.721389in}{0.377067in}}%
\pgfpathquadraticcurveto{\pgfqpoint{8.721389in}{0.388487in}}{\pgfqpoint{8.716652in}{0.394989in}}%
\pgfpathquadraticcurveto{\pgfqpoint{8.711933in}{0.401492in}}{\pgfqpoint{8.703719in}{0.401492in}}%
\pgfpathquadraticcurveto{\pgfqpoint{8.695505in}{0.401492in}}{\pgfqpoint{8.690804in}{0.394989in}}%
\pgfpathquadraticcurveto{\pgfqpoint{8.686103in}{0.388487in}}{\pgfqpoint{8.686103in}{0.377067in}}%
\pgfpathlineto{\pgfqpoint{8.686103in}{0.377067in}}%
\pgfpathclose%
\pgfpathmoveto{\pgfqpoint{8.807019in}{0.379715in}}%
\pgfpathlineto{\pgfqpoint{8.807019in}{0.374654in}}%
\pgfpathlineto{\pgfqpoint{8.759395in}{0.374654in}}%
\pgfpathquadraticcurveto{\pgfqpoint{8.760079in}{0.363954in}}{\pgfqpoint{8.765843in}{0.358353in}}%
\pgfpathquadraticcurveto{\pgfqpoint{8.771607in}{0.352751in}}{\pgfqpoint{8.781910in}{0.352751in}}%
\pgfpathquadraticcurveto{\pgfqpoint{8.787872in}{0.352751in}}{\pgfqpoint{8.793474in}{0.354210in}}%
\pgfpathquadraticcurveto{\pgfqpoint{8.799075in}{0.355669in}}{\pgfqpoint{8.804605in}{0.358605in}}%
\pgfpathlineto{\pgfqpoint{8.804605in}{0.348806in}}%
\pgfpathquadraticcurveto{\pgfqpoint{8.799021in}{0.346447in}}{\pgfqpoint{8.793167in}{0.345204in}}%
\pgfpathquadraticcurveto{\pgfqpoint{8.787313in}{0.343961in}}{\pgfqpoint{8.781297in}{0.343961in}}%
\pgfpathquadraticcurveto{\pgfqpoint{8.766203in}{0.343961in}}{\pgfqpoint{8.757395in}{0.352733in}}%
\pgfpathquadraticcurveto{\pgfqpoint{8.748587in}{0.361523in}}{\pgfqpoint{8.748587in}{0.376509in}}%
\pgfpathquadraticcurveto{\pgfqpoint{8.748587in}{0.391981in}}{\pgfqpoint{8.756945in}{0.401059in}}%
\pgfpathquadraticcurveto{\pgfqpoint{8.765303in}{0.410156in}}{\pgfqpoint{8.779496in}{0.410156in}}%
\pgfpathquadraticcurveto{\pgfqpoint{8.792213in}{0.410156in}}{\pgfqpoint{8.799616in}{0.401960in}}%
\pgfpathquadraticcurveto{\pgfqpoint{8.807019in}{0.393783in}}{\pgfqpoint{8.807019in}{0.379715in}}%
\pgfpathlineto{\pgfqpoint{8.807019in}{0.379715in}}%
\pgfpathclose%
\pgfpathmoveto{\pgfqpoint{8.796662in}{0.382759in}}%
\pgfpathquadraticcurveto{\pgfqpoint{8.796554in}{0.391243in}}{\pgfqpoint{8.791906in}{0.396304in}}%
\pgfpathquadraticcurveto{\pgfqpoint{8.787259in}{0.401384in}}{\pgfqpoint{8.779604in}{0.401384in}}%
\pgfpathquadraticcurveto{\pgfqpoint{8.770940in}{0.401384in}}{\pgfqpoint{8.765735in}{0.396484in}}%
\pgfpathquadraticcurveto{\pgfqpoint{8.760529in}{0.391585in}}{\pgfqpoint{8.759737in}{0.382687in}}%
\pgfpathlineto{\pgfqpoint{8.796662in}{0.382759in}}%
\pgfpathlineto{\pgfqpoint{8.796662in}{0.382759in}}%
\pgfpathclose%
\pgfpathmoveto{\pgfqpoint{8.864207in}{0.406787in}}%
\pgfpathlineto{\pgfqpoint{8.864207in}{0.396989in}}%
\pgfpathquadraticcurveto{\pgfqpoint{8.859830in}{0.399240in}}{\pgfqpoint{8.855093in}{0.400357in}}%
\pgfpathquadraticcurveto{\pgfqpoint{8.850374in}{0.401492in}}{\pgfqpoint{8.845294in}{0.401492in}}%
\pgfpathquadraticcurveto{\pgfqpoint{8.837585in}{0.401492in}}{\pgfqpoint{8.833731in}{0.399132in}}%
\pgfpathquadraticcurveto{\pgfqpoint{8.829876in}{0.396773in}}{\pgfqpoint{8.829876in}{0.392035in}}%
\pgfpathquadraticcurveto{\pgfqpoint{8.829876in}{0.388433in}}{\pgfqpoint{8.832632in}{0.386380in}}%
\pgfpathquadraticcurveto{\pgfqpoint{8.835388in}{0.384326in}}{\pgfqpoint{8.843727in}{0.382471in}}%
\pgfpathlineto{\pgfqpoint{8.847276in}{0.381678in}}%
\pgfpathquadraticcurveto{\pgfqpoint{8.858299in}{0.379319in}}{\pgfqpoint{8.862946in}{0.375014in}}%
\pgfpathquadraticcurveto{\pgfqpoint{8.867594in}{0.370709in}}{\pgfqpoint{8.867594in}{0.363000in}}%
\pgfpathquadraticcurveto{\pgfqpoint{8.867594in}{0.354210in}}{\pgfqpoint{8.860641in}{0.349076in}}%
\pgfpathquadraticcurveto{\pgfqpoint{8.853688in}{0.343961in}}{\pgfqpoint{8.841530in}{0.343961in}}%
\pgfpathquadraticcurveto{\pgfqpoint{8.836469in}{0.343961in}}{\pgfqpoint{8.830975in}{0.344952in}}%
\pgfpathquadraticcurveto{\pgfqpoint{8.825481in}{0.345942in}}{\pgfqpoint{8.819411in}{0.347906in}}%
\pgfpathlineto{\pgfqpoint{8.819411in}{0.358605in}}%
\pgfpathquadraticcurveto{\pgfqpoint{8.825157in}{0.355615in}}{\pgfqpoint{8.830723in}{0.354120in}}%
\pgfpathquadraticcurveto{\pgfqpoint{8.836288in}{0.352643in}}{\pgfqpoint{8.841764in}{0.352643in}}%
\pgfpathquadraticcurveto{\pgfqpoint{8.849077in}{0.352643in}}{\pgfqpoint{8.853004in}{0.355146in}}%
\pgfpathquadraticcurveto{\pgfqpoint{8.856948in}{0.357650in}}{\pgfqpoint{8.856948in}{0.362207in}}%
\pgfpathquadraticcurveto{\pgfqpoint{8.856948in}{0.366422in}}{\pgfqpoint{8.854102in}{0.368674in}}%
\pgfpathquadraticcurveto{\pgfqpoint{8.851275in}{0.370925in}}{\pgfqpoint{8.841638in}{0.373014in}}%
\pgfpathlineto{\pgfqpoint{8.838036in}{0.373861in}}%
\pgfpathquadraticcurveto{\pgfqpoint{8.828417in}{0.375878in}}{\pgfqpoint{8.824130in}{0.380075in}}%
\pgfpathquadraticcurveto{\pgfqpoint{8.819861in}{0.384272in}}{\pgfqpoint{8.819861in}{0.391585in}}%
\pgfpathquadraticcurveto{\pgfqpoint{8.819861in}{0.400483in}}{\pgfqpoint{8.826166in}{0.405310in}}%
\pgfpathquadraticcurveto{\pgfqpoint{8.832470in}{0.410156in}}{\pgfqpoint{8.844070in}{0.410156in}}%
\pgfpathquadraticcurveto{\pgfqpoint{8.849798in}{0.410156in}}{\pgfqpoint{8.854859in}{0.409309in}}%
\pgfpathquadraticcurveto{\pgfqpoint{8.859938in}{0.408480in}}{\pgfqpoint{8.864207in}{0.406787in}}%
\pgfpathlineto{\pgfqpoint{8.864207in}{0.406787in}}%
\pgfpathclose%
\pgfpathmoveto{\pgfqpoint{8.894089in}{0.355056in}}%
\pgfpathlineto{\pgfqpoint{8.894089in}{0.321626in}}%
\pgfpathlineto{\pgfqpoint{8.883678in}{0.321626in}}%
\pgfpathlineto{\pgfqpoint{8.883678in}{0.408643in}}%
\pgfpathlineto{\pgfqpoint{8.894089in}{0.408643in}}%
\pgfpathlineto{\pgfqpoint{8.894089in}{0.399078in}}%
\pgfpathquadraticcurveto{\pgfqpoint{8.897368in}{0.404698in}}{\pgfqpoint{8.902339in}{0.407418in}}%
\pgfpathquadraticcurveto{\pgfqpoint{8.907328in}{0.410156in}}{\pgfqpoint{8.914245in}{0.410156in}}%
\pgfpathquadraticcurveto{\pgfqpoint{8.925737in}{0.410156in}}{\pgfqpoint{8.932906in}{0.401041in}}%
\pgfpathquadraticcurveto{\pgfqpoint{8.940092in}{0.391927in}}{\pgfqpoint{8.940092in}{0.377067in}}%
\pgfpathquadraticcurveto{\pgfqpoint{8.940092in}{0.362207in}}{\pgfqpoint{8.932906in}{0.353075in}}%
\pgfpathquadraticcurveto{\pgfqpoint{8.925737in}{0.343961in}}{\pgfqpoint{8.914245in}{0.343961in}}%
\pgfpathquadraticcurveto{\pgfqpoint{8.907328in}{0.343961in}}{\pgfqpoint{8.902339in}{0.346699in}}%
\pgfpathquadraticcurveto{\pgfqpoint{8.897368in}{0.349437in}}{\pgfqpoint{8.894089in}{0.355056in}}%
\pgfpathlineto{\pgfqpoint{8.894089in}{0.355056in}}%
\pgfpathclose%
\pgfpathmoveto{\pgfqpoint{8.929339in}{0.377067in}}%
\pgfpathquadraticcurveto{\pgfqpoint{8.929339in}{0.388487in}}{\pgfqpoint{8.924638in}{0.394989in}}%
\pgfpathquadraticcurveto{\pgfqpoint{8.919937in}{0.401492in}}{\pgfqpoint{8.911723in}{0.401492in}}%
\pgfpathquadraticcurveto{\pgfqpoint{8.903492in}{0.401492in}}{\pgfqpoint{8.898791in}{0.394989in}}%
\pgfpathquadraticcurveto{\pgfqpoint{8.894089in}{0.388487in}}{\pgfqpoint{8.894089in}{0.377067in}}%
\pgfpathquadraticcurveto{\pgfqpoint{8.894089in}{0.365648in}}{\pgfqpoint{8.898791in}{0.359145in}}%
\pgfpathquadraticcurveto{\pgfqpoint{8.903492in}{0.352643in}}{\pgfqpoint{8.911723in}{0.352643in}}%
\pgfpathquadraticcurveto{\pgfqpoint{8.919937in}{0.352643in}}{\pgfqpoint{8.924638in}{0.359145in}}%
\pgfpathquadraticcurveto{\pgfqpoint{8.929339in}{0.365648in}}{\pgfqpoint{8.929339in}{0.377067in}}%
\pgfpathlineto{\pgfqpoint{8.929339in}{0.377067in}}%
\pgfpathclose%
\pgfpathmoveto{\pgfqpoint{8.957258in}{0.408643in}}%
\pgfpathlineto{\pgfqpoint{8.967615in}{0.408643in}}%
\pgfpathlineto{\pgfqpoint{8.967615in}{0.345600in}}%
\pgfpathlineto{\pgfqpoint{8.957258in}{0.345600in}}%
\pgfpathlineto{\pgfqpoint{8.957258in}{0.408643in}}%
\pgfpathlineto{\pgfqpoint{8.957258in}{0.408643in}}%
\pgfpathclose%
\pgfpathmoveto{\pgfqpoint{8.957258in}{0.433193in}}%
\pgfpathlineto{\pgfqpoint{8.967615in}{0.433193in}}%
\pgfpathlineto{\pgfqpoint{8.967615in}{0.420062in}}%
\pgfpathlineto{\pgfqpoint{8.957258in}{0.420062in}}%
\pgfpathlineto{\pgfqpoint{8.957258in}{0.433193in}}%
\pgfpathlineto{\pgfqpoint{8.957258in}{0.433193in}}%
\pgfpathclose%
\pgfpathmoveto{\pgfqpoint{8.999533in}{0.426547in}}%
\pgfpathlineto{\pgfqpoint{8.999533in}{0.408643in}}%
\pgfpathlineto{\pgfqpoint{9.020859in}{0.408643in}}%
\pgfpathlineto{\pgfqpoint{9.020859in}{0.400591in}}%
\pgfpathlineto{\pgfqpoint{8.999533in}{0.400591in}}%
\pgfpathlineto{\pgfqpoint{8.999533in}{0.366368in}}%
\pgfpathquadraticcurveto{\pgfqpoint{8.999533in}{0.358659in}}{\pgfqpoint{9.001640in}{0.356461in}}%
\pgfpathquadraticcurveto{\pgfqpoint{9.003747in}{0.354264in}}{\pgfqpoint{9.010232in}{0.354264in}}%
\pgfpathlineto{\pgfqpoint{9.020859in}{0.354264in}}%
\pgfpathlineto{\pgfqpoint{9.020859in}{0.345600in}}%
\pgfpathlineto{\pgfqpoint{9.010232in}{0.345600in}}%
\pgfpathquadraticcurveto{\pgfqpoint{8.998236in}{0.345600in}}{\pgfqpoint{8.993679in}{0.350067in}}%
\pgfpathquadraticcurveto{\pgfqpoint{8.989122in}{0.354552in}}{\pgfqpoint{8.989122in}{0.366368in}}%
\pgfpathlineto{\pgfqpoint{8.989122in}{0.400591in}}%
\pgfpathlineto{\pgfqpoint{8.981520in}{0.400591in}}%
\pgfpathlineto{\pgfqpoint{8.981520in}{0.408643in}}%
\pgfpathlineto{\pgfqpoint{8.989122in}{0.408643in}}%
\pgfpathlineto{\pgfqpoint{8.989122in}{0.426547in}}%
\pgfpathlineto{\pgfqpoint{8.999533in}{0.426547in}}%
\pgfpathlineto{\pgfqpoint{8.999533in}{0.426547in}}%
\pgfpathclose%
\pgfpathmoveto{\pgfqpoint{9.088405in}{0.379715in}}%
\pgfpathlineto{\pgfqpoint{9.088405in}{0.374654in}}%
\pgfpathlineto{\pgfqpoint{9.040780in}{0.374654in}}%
\pgfpathquadraticcurveto{\pgfqpoint{9.041465in}{0.363954in}}{\pgfqpoint{9.047229in}{0.358353in}}%
\pgfpathquadraticcurveto{\pgfqpoint{9.052993in}{0.352751in}}{\pgfqpoint{9.063296in}{0.352751in}}%
\pgfpathquadraticcurveto{\pgfqpoint{9.069258in}{0.352751in}}{\pgfqpoint{9.074859in}{0.354210in}}%
\pgfpathquadraticcurveto{\pgfqpoint{9.080461in}{0.355669in}}{\pgfqpoint{9.085991in}{0.358605in}}%
\pgfpathlineto{\pgfqpoint{9.085991in}{0.348806in}}%
\pgfpathquadraticcurveto{\pgfqpoint{9.080407in}{0.346447in}}{\pgfqpoint{9.074553in}{0.345204in}}%
\pgfpathquadraticcurveto{\pgfqpoint{9.068699in}{0.343961in}}{\pgfqpoint{9.062683in}{0.343961in}}%
\pgfpathquadraticcurveto{\pgfqpoint{9.047589in}{0.343961in}}{\pgfqpoint{9.038781in}{0.352733in}}%
\pgfpathquadraticcurveto{\pgfqpoint{9.029973in}{0.361523in}}{\pgfqpoint{9.029973in}{0.376509in}}%
\pgfpathquadraticcurveto{\pgfqpoint{9.029973in}{0.391981in}}{\pgfqpoint{9.038331in}{0.401059in}}%
\pgfpathquadraticcurveto{\pgfqpoint{9.046688in}{0.410156in}}{\pgfqpoint{9.060882in}{0.410156in}}%
\pgfpathquadraticcurveto{\pgfqpoint{9.073599in}{0.410156in}}{\pgfqpoint{9.081002in}{0.401960in}}%
\pgfpathquadraticcurveto{\pgfqpoint{9.088405in}{0.393783in}}{\pgfqpoint{9.088405in}{0.379715in}}%
\pgfpathlineto{\pgfqpoint{9.088405in}{0.379715in}}%
\pgfpathclose%
\pgfpathmoveto{\pgfqpoint{9.078048in}{0.382759in}}%
\pgfpathquadraticcurveto{\pgfqpoint{9.077939in}{0.391243in}}{\pgfqpoint{9.073292in}{0.396304in}}%
\pgfpathquadraticcurveto{\pgfqpoint{9.068645in}{0.401384in}}{\pgfqpoint{9.060990in}{0.401384in}}%
\pgfpathquadraticcurveto{\pgfqpoint{9.052326in}{0.401384in}}{\pgfqpoint{9.047121in}{0.396484in}}%
\pgfpathquadraticcurveto{\pgfqpoint{9.041915in}{0.391585in}}{\pgfqpoint{9.041123in}{0.382687in}}%
\pgfpathlineto{\pgfqpoint{9.078048in}{0.382759in}}%
\pgfpathlineto{\pgfqpoint{9.078048in}{0.382759in}}%
\pgfpathclose%
\pgfpathmoveto{\pgfqpoint{9.166469in}{0.401384in}}%
\pgfpathquadraticcurveto{\pgfqpoint{9.158148in}{0.401384in}}{\pgfqpoint{9.153302in}{0.394881in}}%
\pgfpathquadraticcurveto{\pgfqpoint{9.148457in}{0.388379in}}{\pgfqpoint{9.148457in}{0.377067in}}%
\pgfpathquadraticcurveto{\pgfqpoint{9.148457in}{0.365756in}}{\pgfqpoint{9.153266in}{0.359253in}}%
\pgfpathquadraticcurveto{\pgfqpoint{9.158094in}{0.352751in}}{\pgfqpoint{9.166469in}{0.352751in}}%
\pgfpathquadraticcurveto{\pgfqpoint{9.174755in}{0.352751in}}{\pgfqpoint{9.179582in}{0.359271in}}%
\pgfpathquadraticcurveto{\pgfqpoint{9.184427in}{0.365810in}}{\pgfqpoint{9.184427in}{0.377067in}}%
\pgfpathquadraticcurveto{\pgfqpoint{9.184427in}{0.388271in}}{\pgfqpoint{9.179582in}{0.394827in}}%
\pgfpathquadraticcurveto{\pgfqpoint{9.174755in}{0.401384in}}{\pgfqpoint{9.166469in}{0.401384in}}%
\pgfpathlineto{\pgfqpoint{9.166469in}{0.401384in}}%
\pgfpathclose%
\pgfpathmoveto{\pgfqpoint{9.166469in}{0.410156in}}%
\pgfpathquadraticcurveto{\pgfqpoint{9.179978in}{0.410156in}}{\pgfqpoint{9.187688in}{0.401366in}}%
\pgfpathquadraticcurveto{\pgfqpoint{9.195415in}{0.392594in}}{\pgfqpoint{9.195415in}{0.377067in}}%
\pgfpathquadraticcurveto{\pgfqpoint{9.195415in}{0.361595in}}{\pgfqpoint{9.187688in}{0.352769in}}%
\pgfpathquadraticcurveto{\pgfqpoint{9.179978in}{0.343961in}}{\pgfqpoint{9.166469in}{0.343961in}}%
\pgfpathquadraticcurveto{\pgfqpoint{9.152906in}{0.343961in}}{\pgfqpoint{9.145215in}{0.352769in}}%
\pgfpathquadraticcurveto{\pgfqpoint{9.137542in}{0.361595in}}{\pgfqpoint{9.137542in}{0.377067in}}%
\pgfpathquadraticcurveto{\pgfqpoint{9.137542in}{0.392594in}}{\pgfqpoint{9.145215in}{0.401366in}}%
\pgfpathquadraticcurveto{\pgfqpoint{9.152906in}{0.410156in}}{\pgfqpoint{9.166469in}{0.410156in}}%
\pgfpathlineto{\pgfqpoint{9.166469in}{0.410156in}}%
\pgfpathclose%
\pgfpathmoveto{\pgfqpoint{9.264996in}{0.383660in}}%
\pgfpathlineto{\pgfqpoint{9.264996in}{0.345600in}}%
\pgfpathlineto{\pgfqpoint{9.254639in}{0.345600in}}%
\pgfpathlineto{\pgfqpoint{9.254639in}{0.383317in}}%
\pgfpathquadraticcurveto{\pgfqpoint{9.254639in}{0.392269in}}{\pgfqpoint{9.251144in}{0.396700in}}%
\pgfpathquadraticcurveto{\pgfqpoint{9.247650in}{0.401149in}}{\pgfqpoint{9.240679in}{0.401149in}}%
\pgfpathquadraticcurveto{\pgfqpoint{9.232286in}{0.401149in}}{\pgfqpoint{9.227440in}{0.395800in}}%
\pgfpathquadraticcurveto{\pgfqpoint{9.222595in}{0.390468in}}{\pgfqpoint{9.222595in}{0.381228in}}%
\pgfpathlineto{\pgfqpoint{9.222595in}{0.345600in}}%
\pgfpathlineto{\pgfqpoint{9.212184in}{0.345600in}}%
\pgfpathlineto{\pgfqpoint{9.212184in}{0.408643in}}%
\pgfpathlineto{\pgfqpoint{9.222595in}{0.408643in}}%
\pgfpathlineto{\pgfqpoint{9.222595in}{0.398844in}}%
\pgfpathquadraticcurveto{\pgfqpoint{9.226324in}{0.404536in}}{\pgfqpoint{9.231349in}{0.407346in}}%
\pgfpathquadraticcurveto{\pgfqpoint{9.236392in}{0.410156in}}{\pgfqpoint{9.242985in}{0.410156in}}%
\pgfpathquadraticcurveto{\pgfqpoint{9.253846in}{0.410156in}}{\pgfqpoint{9.259412in}{0.403437in}}%
\pgfpathquadraticcurveto{\pgfqpoint{9.264996in}{0.396718in}}{\pgfqpoint{9.264996in}{0.383660in}}%
\pgfpathlineto{\pgfqpoint{9.264996in}{0.383660in}}%
\pgfpathclose%
\pgfpathmoveto{\pgfqpoint{9.285638in}{0.433193in}}%
\pgfpathlineto{\pgfqpoint{9.295995in}{0.433193in}}%
\pgfpathlineto{\pgfqpoint{9.295995in}{0.345600in}}%
\pgfpathlineto{\pgfqpoint{9.285638in}{0.345600in}}%
\pgfpathlineto{\pgfqpoint{9.285638in}{0.433193in}}%
\pgfpathlineto{\pgfqpoint{9.285638in}{0.433193in}}%
\pgfpathclose%
\pgfpathmoveto{\pgfqpoint{9.343889in}{0.339746in}}%
\pgfpathquadraticcurveto{\pgfqpoint{9.339512in}{0.328488in}}{\pgfqpoint{9.335333in}{0.325066in}}%
\pgfpathquadraticcurveto{\pgfqpoint{9.331172in}{0.321626in}}{\pgfqpoint{9.324202in}{0.321626in}}%
\pgfpathlineto{\pgfqpoint{9.315916in}{0.321626in}}%
\pgfpathlineto{\pgfqpoint{9.315916in}{0.330290in}}%
\pgfpathlineto{\pgfqpoint{9.322004in}{0.330290in}}%
\pgfpathquadraticcurveto{\pgfqpoint{9.326273in}{0.330290in}}{\pgfqpoint{9.328633in}{0.332325in}}%
\pgfpathquadraticcurveto{\pgfqpoint{9.331010in}{0.334342in}}{\pgfqpoint{9.333874in}{0.341889in}}%
\pgfpathlineto{\pgfqpoint{9.335729in}{0.346609in}}%
\pgfpathlineto{\pgfqpoint{9.310242in}{0.408643in}}%
\pgfpathlineto{\pgfqpoint{9.321212in}{0.408643in}}%
\pgfpathlineto{\pgfqpoint{9.340917in}{0.359343in}}%
\pgfpathlineto{\pgfqpoint{9.360622in}{0.408643in}}%
\pgfpathlineto{\pgfqpoint{9.371592in}{0.408643in}}%
\pgfpathlineto{\pgfqpoint{9.343889in}{0.339746in}}%
\pgfpathlineto{\pgfqpoint{9.343889in}{0.339746in}}%
\pgfpathclose%
\pgfpathmoveto{\pgfqpoint{9.432779in}{0.426547in}}%
\pgfpathlineto{\pgfqpoint{9.432779in}{0.408643in}}%
\pgfpathlineto{\pgfqpoint{9.454105in}{0.408643in}}%
\pgfpathlineto{\pgfqpoint{9.454105in}{0.400591in}}%
\pgfpathlineto{\pgfqpoint{9.432779in}{0.400591in}}%
\pgfpathlineto{\pgfqpoint{9.432779in}{0.366368in}}%
\pgfpathquadraticcurveto{\pgfqpoint{9.432779in}{0.358659in}}{\pgfqpoint{9.434886in}{0.356461in}}%
\pgfpathquadraticcurveto{\pgfqpoint{9.436994in}{0.354264in}}{\pgfqpoint{9.443478in}{0.354264in}}%
\pgfpathlineto{\pgfqpoint{9.454105in}{0.354264in}}%
\pgfpathlineto{\pgfqpoint{9.454105in}{0.345600in}}%
\pgfpathlineto{\pgfqpoint{9.443478in}{0.345600in}}%
\pgfpathquadraticcurveto{\pgfqpoint{9.431482in}{0.345600in}}{\pgfqpoint{9.426925in}{0.350067in}}%
\pgfpathquadraticcurveto{\pgfqpoint{9.422368in}{0.354552in}}{\pgfqpoint{9.422368in}{0.366368in}}%
\pgfpathlineto{\pgfqpoint{9.422368in}{0.400591in}}%
\pgfpathlineto{\pgfqpoint{9.414767in}{0.400591in}}%
\pgfpathlineto{\pgfqpoint{9.414767in}{0.408643in}}%
\pgfpathlineto{\pgfqpoint{9.422368in}{0.408643in}}%
\pgfpathlineto{\pgfqpoint{9.422368in}{0.426547in}}%
\pgfpathlineto{\pgfqpoint{9.432779in}{0.426547in}}%
\pgfpathlineto{\pgfqpoint{9.432779in}{0.426547in}}%
\pgfpathclose%
\pgfpathmoveto{\pgfqpoint{9.520138in}{0.383660in}}%
\pgfpathlineto{\pgfqpoint{9.520138in}{0.345600in}}%
\pgfpathlineto{\pgfqpoint{9.509781in}{0.345600in}}%
\pgfpathlineto{\pgfqpoint{9.509781in}{0.383317in}}%
\pgfpathquadraticcurveto{\pgfqpoint{9.509781in}{0.392269in}}{\pgfqpoint{9.506286in}{0.396700in}}%
\pgfpathquadraticcurveto{\pgfqpoint{9.502792in}{0.401149in}}{\pgfqpoint{9.495821in}{0.401149in}}%
\pgfpathquadraticcurveto{\pgfqpoint{9.487428in}{0.401149in}}{\pgfqpoint{9.482582in}{0.395800in}}%
\pgfpathquadraticcurveto{\pgfqpoint{9.477737in}{0.390468in}}{\pgfqpoint{9.477737in}{0.381228in}}%
\pgfpathlineto{\pgfqpoint{9.477737in}{0.345600in}}%
\pgfpathlineto{\pgfqpoint{9.467326in}{0.345600in}}%
\pgfpathlineto{\pgfqpoint{9.467326in}{0.433193in}}%
\pgfpathlineto{\pgfqpoint{9.477737in}{0.433193in}}%
\pgfpathlineto{\pgfqpoint{9.477737in}{0.398844in}}%
\pgfpathquadraticcurveto{\pgfqpoint{9.481466in}{0.404536in}}{\pgfqpoint{9.486491in}{0.407346in}}%
\pgfpathquadraticcurveto{\pgfqpoint{9.491535in}{0.410156in}}{\pgfqpoint{9.498127in}{0.410156in}}%
\pgfpathquadraticcurveto{\pgfqpoint{9.508988in}{0.410156in}}{\pgfqpoint{9.514554in}{0.403437in}}%
\pgfpathquadraticcurveto{\pgfqpoint{9.520138in}{0.396718in}}{\pgfqpoint{9.520138in}{0.383660in}}%
\pgfpathlineto{\pgfqpoint{9.520138in}{0.383660in}}%
\pgfpathclose%
\pgfpathmoveto{\pgfqpoint{9.577308in}{0.398970in}}%
\pgfpathquadraticcurveto{\pgfqpoint{9.575561in}{0.399979in}}{\pgfqpoint{9.573508in}{0.400447in}}%
\pgfpathquadraticcurveto{\pgfqpoint{9.571454in}{0.400933in}}{\pgfqpoint{9.568987in}{0.400933in}}%
\pgfpathquadraticcurveto{\pgfqpoint{9.560197in}{0.400933in}}{\pgfqpoint{9.555496in}{0.395223in}}%
\pgfpathquadraticcurveto{\pgfqpoint{9.550794in}{0.389514in}}{\pgfqpoint{9.550794in}{0.378814in}}%
\pgfpathlineto{\pgfqpoint{9.550794in}{0.345600in}}%
\pgfpathlineto{\pgfqpoint{9.540383in}{0.345600in}}%
\pgfpathlineto{\pgfqpoint{9.540383in}{0.408643in}}%
\pgfpathlineto{\pgfqpoint{9.550794in}{0.408643in}}%
\pgfpathlineto{\pgfqpoint{9.550794in}{0.398844in}}%
\pgfpathquadraticcurveto{\pgfqpoint{9.554073in}{0.404590in}}{\pgfqpoint{9.559296in}{0.407364in}}%
\pgfpathquadraticcurveto{\pgfqpoint{9.564538in}{0.410156in}}{\pgfqpoint{9.572031in}{0.410156in}}%
\pgfpathquadraticcurveto{\pgfqpoint{9.573094in}{0.410156in}}{\pgfqpoint{9.574390in}{0.410011in}}%
\pgfpathquadraticcurveto{\pgfqpoint{9.575687in}{0.409885in}}{\pgfqpoint{9.577254in}{0.409597in}}%
\pgfpathlineto{\pgfqpoint{9.577308in}{0.398970in}}%
\pgfpathlineto{\pgfqpoint{9.577308in}{0.398970in}}%
\pgfpathclose%
\pgfpathmoveto{\pgfqpoint{9.639558in}{0.379715in}}%
\pgfpathlineto{\pgfqpoint{9.639558in}{0.374654in}}%
\pgfpathlineto{\pgfqpoint{9.591934in}{0.374654in}}%
\pgfpathquadraticcurveto{\pgfqpoint{9.592619in}{0.363954in}}{\pgfqpoint{9.598383in}{0.358353in}}%
\pgfpathquadraticcurveto{\pgfqpoint{9.604146in}{0.352751in}}{\pgfqpoint{9.614449in}{0.352751in}}%
\pgfpathquadraticcurveto{\pgfqpoint{9.620411in}{0.352751in}}{\pgfqpoint{9.626013in}{0.354210in}}%
\pgfpathquadraticcurveto{\pgfqpoint{9.631615in}{0.355669in}}{\pgfqpoint{9.637145in}{0.358605in}}%
\pgfpathlineto{\pgfqpoint{9.637145in}{0.348806in}}%
\pgfpathquadraticcurveto{\pgfqpoint{9.631561in}{0.346447in}}{\pgfqpoint{9.625707in}{0.345204in}}%
\pgfpathquadraticcurveto{\pgfqpoint{9.619853in}{0.343961in}}{\pgfqpoint{9.613837in}{0.343961in}}%
\pgfpathquadraticcurveto{\pgfqpoint{9.598743in}{0.343961in}}{\pgfqpoint{9.589935in}{0.352733in}}%
\pgfpathquadraticcurveto{\pgfqpoint{9.581127in}{0.361523in}}{\pgfqpoint{9.581127in}{0.376509in}}%
\pgfpathquadraticcurveto{\pgfqpoint{9.581127in}{0.391981in}}{\pgfqpoint{9.589485in}{0.401059in}}%
\pgfpathquadraticcurveto{\pgfqpoint{9.597842in}{0.410156in}}{\pgfqpoint{9.612036in}{0.410156in}}%
\pgfpathquadraticcurveto{\pgfqpoint{9.624752in}{0.410156in}}{\pgfqpoint{9.632155in}{0.401960in}}%
\pgfpathquadraticcurveto{\pgfqpoint{9.639558in}{0.393783in}}{\pgfqpoint{9.639558in}{0.379715in}}%
\pgfpathlineto{\pgfqpoint{9.639558in}{0.379715in}}%
\pgfpathclose%
\pgfpathmoveto{\pgfqpoint{9.629201in}{0.382759in}}%
\pgfpathquadraticcurveto{\pgfqpoint{9.629093in}{0.391243in}}{\pgfqpoint{9.624446in}{0.396304in}}%
\pgfpathquadraticcurveto{\pgfqpoint{9.619799in}{0.401384in}}{\pgfqpoint{9.612144in}{0.401384in}}%
\pgfpathquadraticcurveto{\pgfqpoint{9.603480in}{0.401384in}}{\pgfqpoint{9.598275in}{0.396484in}}%
\pgfpathquadraticcurveto{\pgfqpoint{9.593069in}{0.391585in}}{\pgfqpoint{9.592276in}{0.382687in}}%
\pgfpathlineto{\pgfqpoint{9.629201in}{0.382759in}}%
\pgfpathlineto{\pgfqpoint{9.629201in}{0.382759in}}%
\pgfpathclose%
\pgfpathmoveto{\pgfqpoint{9.710490in}{0.379715in}}%
\pgfpathlineto{\pgfqpoint{9.710490in}{0.374654in}}%
\pgfpathlineto{\pgfqpoint{9.662866in}{0.374654in}}%
\pgfpathquadraticcurveto{\pgfqpoint{9.663551in}{0.363954in}}{\pgfqpoint{9.669314in}{0.358353in}}%
\pgfpathquadraticcurveto{\pgfqpoint{9.675078in}{0.352751in}}{\pgfqpoint{9.685381in}{0.352751in}}%
\pgfpathquadraticcurveto{\pgfqpoint{9.691343in}{0.352751in}}{\pgfqpoint{9.696945in}{0.354210in}}%
\pgfpathquadraticcurveto{\pgfqpoint{9.702547in}{0.355669in}}{\pgfqpoint{9.708077in}{0.358605in}}%
\pgfpathlineto{\pgfqpoint{9.708077in}{0.348806in}}%
\pgfpathquadraticcurveto{\pgfqpoint{9.702493in}{0.346447in}}{\pgfqpoint{9.696639in}{0.345204in}}%
\pgfpathquadraticcurveto{\pgfqpoint{9.690785in}{0.343961in}}{\pgfqpoint{9.684769in}{0.343961in}}%
\pgfpathquadraticcurveto{\pgfqpoint{9.669675in}{0.343961in}}{\pgfqpoint{9.660867in}{0.352733in}}%
\pgfpathquadraticcurveto{\pgfqpoint{9.652059in}{0.361523in}}{\pgfqpoint{9.652059in}{0.376509in}}%
\pgfpathquadraticcurveto{\pgfqpoint{9.652059in}{0.391981in}}{\pgfqpoint{9.660416in}{0.401059in}}%
\pgfpathquadraticcurveto{\pgfqpoint{9.668774in}{0.410156in}}{\pgfqpoint{9.682968in}{0.410156in}}%
\pgfpathquadraticcurveto{\pgfqpoint{9.695684in}{0.410156in}}{\pgfqpoint{9.703087in}{0.401960in}}%
\pgfpathquadraticcurveto{\pgfqpoint{9.710490in}{0.393783in}}{\pgfqpoint{9.710490in}{0.379715in}}%
\pgfpathlineto{\pgfqpoint{9.710490in}{0.379715in}}%
\pgfpathclose%
\pgfpathmoveto{\pgfqpoint{9.700133in}{0.382759in}}%
\pgfpathquadraticcurveto{\pgfqpoint{9.700025in}{0.391243in}}{\pgfqpoint{9.695378in}{0.396304in}}%
\pgfpathquadraticcurveto{\pgfqpoint{9.690731in}{0.401384in}}{\pgfqpoint{9.683076in}{0.401384in}}%
\pgfpathquadraticcurveto{\pgfqpoint{9.674412in}{0.401384in}}{\pgfqpoint{9.669206in}{0.396484in}}%
\pgfpathquadraticcurveto{\pgfqpoint{9.664001in}{0.391585in}}{\pgfqpoint{9.663208in}{0.382687in}}%
\pgfpathlineto{\pgfqpoint{9.700133in}{0.382759in}}%
\pgfpathlineto{\pgfqpoint{9.700133in}{0.382759in}}%
\pgfpathclose%
\pgfpathmoveto{\pgfqpoint{9.774379in}{0.426547in}}%
\pgfpathlineto{\pgfqpoint{9.774379in}{0.408643in}}%
\pgfpathlineto{\pgfqpoint{9.795706in}{0.408643in}}%
\pgfpathlineto{\pgfqpoint{9.795706in}{0.400591in}}%
\pgfpathlineto{\pgfqpoint{9.774379in}{0.400591in}}%
\pgfpathlineto{\pgfqpoint{9.774379in}{0.366368in}}%
\pgfpathquadraticcurveto{\pgfqpoint{9.774379in}{0.358659in}}{\pgfqpoint{9.776487in}{0.356461in}}%
\pgfpathquadraticcurveto{\pgfqpoint{9.778594in}{0.354264in}}{\pgfqpoint{9.785079in}{0.354264in}}%
\pgfpathlineto{\pgfqpoint{9.795706in}{0.354264in}}%
\pgfpathlineto{\pgfqpoint{9.795706in}{0.345600in}}%
\pgfpathlineto{\pgfqpoint{9.785079in}{0.345600in}}%
\pgfpathquadraticcurveto{\pgfqpoint{9.773082in}{0.345600in}}{\pgfqpoint{9.768525in}{0.350067in}}%
\pgfpathquadraticcurveto{\pgfqpoint{9.763968in}{0.354552in}}{\pgfqpoint{9.763968in}{0.366368in}}%
\pgfpathlineto{\pgfqpoint{9.763968in}{0.400591in}}%
\pgfpathlineto{\pgfqpoint{9.756367in}{0.400591in}}%
\pgfpathlineto{\pgfqpoint{9.756367in}{0.408643in}}%
\pgfpathlineto{\pgfqpoint{9.763968in}{0.408643in}}%
\pgfpathlineto{\pgfqpoint{9.763968in}{0.426547in}}%
\pgfpathlineto{\pgfqpoint{9.774379in}{0.426547in}}%
\pgfpathlineto{\pgfqpoint{9.774379in}{0.426547in}}%
\pgfpathclose%
\pgfpathmoveto{\pgfqpoint{9.845852in}{0.398970in}}%
\pgfpathquadraticcurveto{\pgfqpoint{9.844104in}{0.399979in}}{\pgfqpoint{9.842051in}{0.400447in}}%
\pgfpathquadraticcurveto{\pgfqpoint{9.839998in}{0.400933in}}{\pgfqpoint{9.837530in}{0.400933in}}%
\pgfpathquadraticcurveto{\pgfqpoint{9.828740in}{0.400933in}}{\pgfqpoint{9.824039in}{0.395223in}}%
\pgfpathquadraticcurveto{\pgfqpoint{9.819338in}{0.389514in}}{\pgfqpoint{9.819338in}{0.378814in}}%
\pgfpathlineto{\pgfqpoint{9.819338in}{0.345600in}}%
\pgfpathlineto{\pgfqpoint{9.808927in}{0.345600in}}%
\pgfpathlineto{\pgfqpoint{9.808927in}{0.408643in}}%
\pgfpathlineto{\pgfqpoint{9.819338in}{0.408643in}}%
\pgfpathlineto{\pgfqpoint{9.819338in}{0.398844in}}%
\pgfpathquadraticcurveto{\pgfqpoint{9.822616in}{0.404590in}}{\pgfqpoint{9.827839in}{0.407364in}}%
\pgfpathquadraticcurveto{\pgfqpoint{9.833081in}{0.410156in}}{\pgfqpoint{9.840574in}{0.410156in}}%
\pgfpathquadraticcurveto{\pgfqpoint{9.841637in}{0.410156in}}{\pgfqpoint{9.842934in}{0.410011in}}%
\pgfpathquadraticcurveto{\pgfqpoint{9.844230in}{0.409885in}}{\pgfqpoint{9.845798in}{0.409597in}}%
\pgfpathlineto{\pgfqpoint{9.845852in}{0.398970in}}%
\pgfpathlineto{\pgfqpoint{9.845852in}{0.398970in}}%
\pgfpathclose%
\pgfpathmoveto{\pgfqpoint{9.908102in}{0.379715in}}%
\pgfpathlineto{\pgfqpoint{9.908102in}{0.374654in}}%
\pgfpathlineto{\pgfqpoint{9.860477in}{0.374654in}}%
\pgfpathquadraticcurveto{\pgfqpoint{9.861162in}{0.363954in}}{\pgfqpoint{9.866926in}{0.358353in}}%
\pgfpathquadraticcurveto{\pgfqpoint{9.872690in}{0.352751in}}{\pgfqpoint{9.882993in}{0.352751in}}%
\pgfpathquadraticcurveto{\pgfqpoint{9.888955in}{0.352751in}}{\pgfqpoint{9.894556in}{0.354210in}}%
\pgfpathquadraticcurveto{\pgfqpoint{9.900158in}{0.355669in}}{\pgfqpoint{9.905688in}{0.358605in}}%
\pgfpathlineto{\pgfqpoint{9.905688in}{0.348806in}}%
\pgfpathquadraticcurveto{\pgfqpoint{9.900104in}{0.346447in}}{\pgfqpoint{9.894250in}{0.345204in}}%
\pgfpathquadraticcurveto{\pgfqpoint{9.888396in}{0.343961in}}{\pgfqpoint{9.882380in}{0.343961in}}%
\pgfpathquadraticcurveto{\pgfqpoint{9.867286in}{0.343961in}}{\pgfqpoint{9.858478in}{0.352733in}}%
\pgfpathquadraticcurveto{\pgfqpoint{9.849670in}{0.361523in}}{\pgfqpoint{9.849670in}{0.376509in}}%
\pgfpathquadraticcurveto{\pgfqpoint{9.849670in}{0.391981in}}{\pgfqpoint{9.858028in}{0.401059in}}%
\pgfpathquadraticcurveto{\pgfqpoint{9.866385in}{0.410156in}}{\pgfqpoint{9.880579in}{0.410156in}}%
\pgfpathquadraticcurveto{\pgfqpoint{9.893296in}{0.410156in}}{\pgfqpoint{9.900699in}{0.401960in}}%
\pgfpathquadraticcurveto{\pgfqpoint{9.908102in}{0.393783in}}{\pgfqpoint{9.908102in}{0.379715in}}%
\pgfpathlineto{\pgfqpoint{9.908102in}{0.379715in}}%
\pgfpathclose%
\pgfpathmoveto{\pgfqpoint{9.897745in}{0.382759in}}%
\pgfpathquadraticcurveto{\pgfqpoint{9.897637in}{0.391243in}}{\pgfqpoint{9.892989in}{0.396304in}}%
\pgfpathquadraticcurveto{\pgfqpoint{9.888342in}{0.401384in}}{\pgfqpoint{9.880687in}{0.401384in}}%
\pgfpathquadraticcurveto{\pgfqpoint{9.872023in}{0.401384in}}{\pgfqpoint{9.866818in}{0.396484in}}%
\pgfpathquadraticcurveto{\pgfqpoint{9.861612in}{0.391585in}}{\pgfqpoint{9.860820in}{0.382687in}}%
\pgfpathlineto{\pgfqpoint{9.897745in}{0.382759in}}%
\pgfpathlineto{\pgfqpoint{9.897745in}{0.382759in}}%
\pgfpathclose%
\pgfpathmoveto{\pgfqpoint{9.953762in}{0.377283in}}%
\pgfpathquadraticcurveto{\pgfqpoint{9.941208in}{0.377283in}}{\pgfqpoint{9.936363in}{0.374419in}}%
\pgfpathquadraticcurveto{\pgfqpoint{9.931517in}{0.371556in}}{\pgfqpoint{9.931517in}{0.364621in}}%
\pgfpathquadraticcurveto{\pgfqpoint{9.931517in}{0.359109in}}{\pgfqpoint{9.935156in}{0.355867in}}%
\pgfpathquadraticcurveto{\pgfqpoint{9.938794in}{0.352643in}}{\pgfqpoint{9.945026in}{0.352643in}}%
\pgfpathquadraticcurveto{\pgfqpoint{9.953654in}{0.352643in}}{\pgfqpoint{9.958860in}{0.358749in}}%
\pgfpathquadraticcurveto{\pgfqpoint{9.964065in}{0.364855in}}{\pgfqpoint{9.964065in}{0.374978in}}%
\pgfpathlineto{\pgfqpoint{9.964065in}{0.377283in}}%
\pgfpathlineto{\pgfqpoint{9.953762in}{0.377283in}}%
\pgfpathlineto{\pgfqpoint{9.953762in}{0.377283in}}%
\pgfpathclose%
\pgfpathmoveto{\pgfqpoint{9.974422in}{0.381570in}}%
\pgfpathlineto{\pgfqpoint{9.974422in}{0.345600in}}%
\pgfpathlineto{\pgfqpoint{9.964065in}{0.345600in}}%
\pgfpathlineto{\pgfqpoint{9.964065in}{0.355164in}}%
\pgfpathquadraticcurveto{\pgfqpoint{9.960517in}{0.349437in}}{\pgfqpoint{9.955221in}{0.346699in}}%
\pgfpathquadraticcurveto{\pgfqpoint{9.949926in}{0.343961in}}{\pgfqpoint{9.942271in}{0.343961in}}%
\pgfpathquadraticcurveto{\pgfqpoint{9.932598in}{0.343961in}}{\pgfqpoint{9.926870in}{0.349401in}}%
\pgfpathquadraticcurveto{\pgfqpoint{9.921160in}{0.354840in}}{\pgfqpoint{9.921160in}{0.363954in}}%
\pgfpathquadraticcurveto{\pgfqpoint{9.921160in}{0.374582in}}{\pgfqpoint{9.928275in}{0.379985in}}%
\pgfpathquadraticcurveto{\pgfqpoint{9.935408in}{0.385389in}}{\pgfqpoint{9.949530in}{0.385389in}}%
\pgfpathlineto{\pgfqpoint{9.964065in}{0.385389in}}%
\pgfpathlineto{\pgfqpoint{9.964065in}{0.386416in}}%
\pgfpathquadraticcurveto{\pgfqpoint{9.964065in}{0.393566in}}{\pgfqpoint{9.959364in}{0.397475in}}%
\pgfpathquadraticcurveto{\pgfqpoint{9.954663in}{0.401384in}}{\pgfqpoint{9.946161in}{0.401384in}}%
\pgfpathquadraticcurveto{\pgfqpoint{9.940758in}{0.401384in}}{\pgfqpoint{9.935624in}{0.400087in}}%
\pgfpathquadraticcurveto{\pgfqpoint{9.930509in}{0.398790in}}{\pgfqpoint{9.925789in}{0.396196in}}%
\pgfpathlineto{\pgfqpoint{9.925789in}{0.405779in}}%
\pgfpathquadraticcurveto{\pgfqpoint{9.931463in}{0.407976in}}{\pgfqpoint{9.936813in}{0.409057in}}%
\pgfpathquadraticcurveto{\pgfqpoint{9.942163in}{0.410156in}}{\pgfqpoint{9.947224in}{0.410156in}}%
\pgfpathquadraticcurveto{\pgfqpoint{9.960913in}{0.410156in}}{\pgfqpoint{9.967668in}{0.403059in}}%
\pgfpathquadraticcurveto{\pgfqpoint{9.974422in}{0.395980in}}{\pgfqpoint{9.974422in}{0.381570in}}%
\pgfpathlineto{\pgfqpoint{9.974422in}{0.381570in}}%
\pgfpathclose%
\pgfpathmoveto{\pgfqpoint{10.005998in}{0.426547in}}%
\pgfpathlineto{\pgfqpoint{10.005998in}{0.408643in}}%
\pgfpathlineto{\pgfqpoint{10.027324in}{0.408643in}}%
\pgfpathlineto{\pgfqpoint{10.027324in}{0.400591in}}%
\pgfpathlineto{\pgfqpoint{10.005998in}{0.400591in}}%
\pgfpathlineto{\pgfqpoint{10.005998in}{0.366368in}}%
\pgfpathquadraticcurveto{\pgfqpoint{10.005998in}{0.358659in}}{\pgfqpoint{10.008105in}{0.356461in}}%
\pgfpathquadraticcurveto{\pgfqpoint{10.010212in}{0.354264in}}{\pgfqpoint{10.016697in}{0.354264in}}%
\pgfpathlineto{\pgfqpoint{10.027324in}{0.354264in}}%
\pgfpathlineto{\pgfqpoint{10.027324in}{0.345600in}}%
\pgfpathlineto{\pgfqpoint{10.016697in}{0.345600in}}%
\pgfpathquadraticcurveto{\pgfqpoint{10.004701in}{0.345600in}}{\pgfqpoint{10.000144in}{0.350067in}}%
\pgfpathquadraticcurveto{\pgfqpoint{9.995587in}{0.354552in}}{\pgfqpoint{9.995587in}{0.366368in}}%
\pgfpathlineto{\pgfqpoint{9.995587in}{0.400591in}}%
\pgfpathlineto{\pgfqpoint{9.987985in}{0.400591in}}%
\pgfpathlineto{\pgfqpoint{9.987985in}{0.408643in}}%
\pgfpathlineto{\pgfqpoint{9.995587in}{0.408643in}}%
\pgfpathlineto{\pgfqpoint{9.995587in}{0.426547in}}%
\pgfpathlineto{\pgfqpoint{10.005998in}{0.426547in}}%
\pgfpathlineto{\pgfqpoint{10.005998in}{0.426547in}}%
\pgfpathclose%
\pgfpathmoveto{\pgfqpoint{10.090024in}{0.396538in}}%
\pgfpathquadraticcurveto{\pgfqpoint{10.093915in}{0.403527in}}{\pgfqpoint{10.099319in}{0.406841in}}%
\pgfpathquadraticcurveto{\pgfqpoint{10.104722in}{0.410156in}}{\pgfqpoint{10.112035in}{0.410156in}}%
\pgfpathquadraticcurveto{\pgfqpoint{10.121888in}{0.410156in}}{\pgfqpoint{10.127237in}{0.403257in}}%
\pgfpathquadraticcurveto{\pgfqpoint{10.132587in}{0.396376in}}{\pgfqpoint{10.132587in}{0.383660in}}%
\pgfpathlineto{\pgfqpoint{10.132587in}{0.345600in}}%
\pgfpathlineto{\pgfqpoint{10.122176in}{0.345600in}}%
\pgfpathlineto{\pgfqpoint{10.122176in}{0.383317in}}%
\pgfpathquadraticcurveto{\pgfqpoint{10.122176in}{0.392378in}}{\pgfqpoint{10.118952in}{0.396755in}}%
\pgfpathquadraticcurveto{\pgfqpoint{10.115746in}{0.401149in}}{\pgfqpoint{10.109171in}{0.401149in}}%
\pgfpathquadraticcurveto{\pgfqpoint{10.101120in}{0.401149in}}{\pgfqpoint{10.096437in}{0.395800in}}%
\pgfpathquadraticcurveto{\pgfqpoint{10.091771in}{0.390468in}}{\pgfqpoint{10.091771in}{0.381228in}}%
\pgfpathlineto{\pgfqpoint{10.091771in}{0.345600in}}%
\pgfpathlineto{\pgfqpoint{10.081360in}{0.345600in}}%
\pgfpathlineto{\pgfqpoint{10.081360in}{0.383317in}}%
\pgfpathquadraticcurveto{\pgfqpoint{10.081360in}{0.392432in}}{\pgfqpoint{10.078154in}{0.396791in}}%
\pgfpathquadraticcurveto{\pgfqpoint{10.074948in}{0.401149in}}{\pgfqpoint{10.068248in}{0.401149in}}%
\pgfpathquadraticcurveto{\pgfqpoint{10.060304in}{0.401149in}}{\pgfqpoint{10.055621in}{0.395782in}}%
\pgfpathquadraticcurveto{\pgfqpoint{10.050956in}{0.390414in}}{\pgfqpoint{10.050956in}{0.381228in}}%
\pgfpathlineto{\pgfqpoint{10.050956in}{0.345600in}}%
\pgfpathlineto{\pgfqpoint{10.040545in}{0.345600in}}%
\pgfpathlineto{\pgfqpoint{10.040545in}{0.408643in}}%
\pgfpathlineto{\pgfqpoint{10.050956in}{0.408643in}}%
\pgfpathlineto{\pgfqpoint{10.050956in}{0.398844in}}%
\pgfpathquadraticcurveto{\pgfqpoint{10.054504in}{0.404644in}}{\pgfqpoint{10.059458in}{0.407400in}}%
\pgfpathquadraticcurveto{\pgfqpoint{10.064411in}{0.410156in}}{\pgfqpoint{10.071220in}{0.410156in}}%
\pgfpathquadraticcurveto{\pgfqpoint{10.078100in}{0.410156in}}{\pgfqpoint{10.082910in}{0.406661in}}%
\pgfpathquadraticcurveto{\pgfqpoint{10.087719in}{0.403185in}}{\pgfqpoint{10.090024in}{0.396538in}}%
\pgfpathlineto{\pgfqpoint{10.090024in}{0.396538in}}%
\pgfpathclose%
\pgfpathmoveto{\pgfqpoint{10.207157in}{0.379715in}}%
\pgfpathlineto{\pgfqpoint{10.207157in}{0.374654in}}%
\pgfpathlineto{\pgfqpoint{10.159533in}{0.374654in}}%
\pgfpathquadraticcurveto{\pgfqpoint{10.160218in}{0.363954in}}{\pgfqpoint{10.165982in}{0.358353in}}%
\pgfpathquadraticcurveto{\pgfqpoint{10.171745in}{0.352751in}}{\pgfqpoint{10.182048in}{0.352751in}}%
\pgfpathquadraticcurveto{\pgfqpoint{10.188010in}{0.352751in}}{\pgfqpoint{10.193612in}{0.354210in}}%
\pgfpathquadraticcurveto{\pgfqpoint{10.199214in}{0.355669in}}{\pgfqpoint{10.204744in}{0.358605in}}%
\pgfpathlineto{\pgfqpoint{10.204744in}{0.348806in}}%
\pgfpathquadraticcurveto{\pgfqpoint{10.199160in}{0.346447in}}{\pgfqpoint{10.193306in}{0.345204in}}%
\pgfpathquadraticcurveto{\pgfqpoint{10.187452in}{0.343961in}}{\pgfqpoint{10.181436in}{0.343961in}}%
\pgfpathquadraticcurveto{\pgfqpoint{10.166342in}{0.343961in}}{\pgfqpoint{10.157534in}{0.352733in}}%
\pgfpathquadraticcurveto{\pgfqpoint{10.148726in}{0.361523in}}{\pgfqpoint{10.148726in}{0.376509in}}%
\pgfpathquadraticcurveto{\pgfqpoint{10.148726in}{0.391981in}}{\pgfqpoint{10.157084in}{0.401059in}}%
\pgfpathquadraticcurveto{\pgfqpoint{10.165441in}{0.410156in}}{\pgfqpoint{10.179635in}{0.410156in}}%
\pgfpathquadraticcurveto{\pgfqpoint{10.192351in}{0.410156in}}{\pgfqpoint{10.199754in}{0.401960in}}%
\pgfpathquadraticcurveto{\pgfqpoint{10.207157in}{0.393783in}}{\pgfqpoint{10.207157in}{0.379715in}}%
\pgfpathlineto{\pgfqpoint{10.207157in}{0.379715in}}%
\pgfpathclose%
\pgfpathmoveto{\pgfqpoint{10.196800in}{0.382759in}}%
\pgfpathquadraticcurveto{\pgfqpoint{10.196692in}{0.391243in}}{\pgfqpoint{10.192045in}{0.396304in}}%
\pgfpathquadraticcurveto{\pgfqpoint{10.187398in}{0.401384in}}{\pgfqpoint{10.179743in}{0.401384in}}%
\pgfpathquadraticcurveto{\pgfqpoint{10.171079in}{0.401384in}}{\pgfqpoint{10.165873in}{0.396484in}}%
\pgfpathquadraticcurveto{\pgfqpoint{10.160668in}{0.391585in}}{\pgfqpoint{10.159875in}{0.382687in}}%
\pgfpathlineto{\pgfqpoint{10.196800in}{0.382759in}}%
\pgfpathlineto{\pgfqpoint{10.196800in}{0.382759in}}%
\pgfpathclose%
\pgfpathmoveto{\pgfqpoint{10.276576in}{0.383660in}}%
\pgfpathlineto{\pgfqpoint{10.276576in}{0.345600in}}%
\pgfpathlineto{\pgfqpoint{10.266219in}{0.345600in}}%
\pgfpathlineto{\pgfqpoint{10.266219in}{0.383317in}}%
\pgfpathquadraticcurveto{\pgfqpoint{10.266219in}{0.392269in}}{\pgfqpoint{10.262725in}{0.396700in}}%
\pgfpathquadraticcurveto{\pgfqpoint{10.259230in}{0.401149in}}{\pgfqpoint{10.252260in}{0.401149in}}%
\pgfpathquadraticcurveto{\pgfqpoint{10.243866in}{0.401149in}}{\pgfqpoint{10.239021in}{0.395800in}}%
\pgfpathquadraticcurveto{\pgfqpoint{10.234176in}{0.390468in}}{\pgfqpoint{10.234176in}{0.381228in}}%
\pgfpathlineto{\pgfqpoint{10.234176in}{0.345600in}}%
\pgfpathlineto{\pgfqpoint{10.223765in}{0.345600in}}%
\pgfpathlineto{\pgfqpoint{10.223765in}{0.408643in}}%
\pgfpathlineto{\pgfqpoint{10.234176in}{0.408643in}}%
\pgfpathlineto{\pgfqpoint{10.234176in}{0.398844in}}%
\pgfpathquadraticcurveto{\pgfqpoint{10.237904in}{0.404536in}}{\pgfqpoint{10.242929in}{0.407346in}}%
\pgfpathquadraticcurveto{\pgfqpoint{10.247973in}{0.410156in}}{\pgfqpoint{10.254565in}{0.410156in}}%
\pgfpathquadraticcurveto{\pgfqpoint{10.265427in}{0.410156in}}{\pgfqpoint{10.270992in}{0.403437in}}%
\pgfpathquadraticcurveto{\pgfqpoint{10.276576in}{0.396718in}}{\pgfqpoint{10.276576in}{0.383660in}}%
\pgfpathlineto{\pgfqpoint{10.276576in}{0.383660in}}%
\pgfpathclose%
\pgfpathmoveto{\pgfqpoint{10.307467in}{0.426547in}}%
\pgfpathlineto{\pgfqpoint{10.307467in}{0.408643in}}%
\pgfpathlineto{\pgfqpoint{10.328793in}{0.408643in}}%
\pgfpathlineto{\pgfqpoint{10.328793in}{0.400591in}}%
\pgfpathlineto{\pgfqpoint{10.307467in}{0.400591in}}%
\pgfpathlineto{\pgfqpoint{10.307467in}{0.366368in}}%
\pgfpathquadraticcurveto{\pgfqpoint{10.307467in}{0.358659in}}{\pgfqpoint{10.309574in}{0.356461in}}%
\pgfpathquadraticcurveto{\pgfqpoint{10.311682in}{0.354264in}}{\pgfqpoint{10.318166in}{0.354264in}}%
\pgfpathlineto{\pgfqpoint{10.328793in}{0.354264in}}%
\pgfpathlineto{\pgfqpoint{10.328793in}{0.345600in}}%
\pgfpathlineto{\pgfqpoint{10.318166in}{0.345600in}}%
\pgfpathquadraticcurveto{\pgfqpoint{10.306170in}{0.345600in}}{\pgfqpoint{10.301613in}{0.350067in}}%
\pgfpathquadraticcurveto{\pgfqpoint{10.297056in}{0.354552in}}{\pgfqpoint{10.297056in}{0.366368in}}%
\pgfpathlineto{\pgfqpoint{10.297056in}{0.400591in}}%
\pgfpathlineto{\pgfqpoint{10.289455in}{0.400591in}}%
\pgfpathlineto{\pgfqpoint{10.289455in}{0.408643in}}%
\pgfpathlineto{\pgfqpoint{10.297056in}{0.408643in}}%
\pgfpathlineto{\pgfqpoint{10.297056in}{0.426547in}}%
\pgfpathlineto{\pgfqpoint{10.307467in}{0.426547in}}%
\pgfpathlineto{\pgfqpoint{10.307467in}{0.426547in}}%
\pgfpathclose%
\pgfpathmoveto{\pgfqpoint{10.337187in}{0.381786in}}%
\pgfpathlineto{\pgfqpoint{10.367520in}{0.381786in}}%
\pgfpathlineto{\pgfqpoint{10.367520in}{0.372564in}}%
\pgfpathlineto{\pgfqpoint{10.337187in}{0.372564in}}%
\pgfpathlineto{\pgfqpoint{10.337187in}{0.381786in}}%
\pgfpathlineto{\pgfqpoint{10.337187in}{0.381786in}}%
\pgfpathclose%
\pgfpathmoveto{\pgfqpoint{10.420529in}{0.398970in}}%
\pgfpathquadraticcurveto{\pgfqpoint{10.418782in}{0.399979in}}{\pgfqpoint{10.416729in}{0.400447in}}%
\pgfpathquadraticcurveto{\pgfqpoint{10.414675in}{0.400933in}}{\pgfqpoint{10.412208in}{0.400933in}}%
\pgfpathquadraticcurveto{\pgfqpoint{10.403418in}{0.400933in}}{\pgfqpoint{10.398717in}{0.395223in}}%
\pgfpathquadraticcurveto{\pgfqpoint{10.394015in}{0.389514in}}{\pgfqpoint{10.394015in}{0.378814in}}%
\pgfpathlineto{\pgfqpoint{10.394015in}{0.345600in}}%
\pgfpathlineto{\pgfqpoint{10.383604in}{0.345600in}}%
\pgfpathlineto{\pgfqpoint{10.383604in}{0.408643in}}%
\pgfpathlineto{\pgfqpoint{10.394015in}{0.408643in}}%
\pgfpathlineto{\pgfqpoint{10.394015in}{0.398844in}}%
\pgfpathquadraticcurveto{\pgfqpoint{10.397294in}{0.404590in}}{\pgfqpoint{10.402517in}{0.407364in}}%
\pgfpathquadraticcurveto{\pgfqpoint{10.407759in}{0.410156in}}{\pgfqpoint{10.415252in}{0.410156in}}%
\pgfpathquadraticcurveto{\pgfqpoint{10.416314in}{0.410156in}}{\pgfqpoint{10.417611in}{0.410011in}}%
\pgfpathquadraticcurveto{\pgfqpoint{10.418908in}{0.409885in}}{\pgfqpoint{10.420475in}{0.409597in}}%
\pgfpathlineto{\pgfqpoint{10.420529in}{0.398970in}}%
\pgfpathlineto{\pgfqpoint{10.420529in}{0.398970in}}%
\pgfpathclose%
\pgfpathmoveto{\pgfqpoint{10.453275in}{0.401384in}}%
\pgfpathquadraticcurveto{\pgfqpoint{10.444954in}{0.401384in}}{\pgfqpoint{10.440109in}{0.394881in}}%
\pgfpathquadraticcurveto{\pgfqpoint{10.435263in}{0.388379in}}{\pgfqpoint{10.435263in}{0.377067in}}%
\pgfpathquadraticcurveto{\pgfqpoint{10.435263in}{0.365756in}}{\pgfqpoint{10.440072in}{0.359253in}}%
\pgfpathquadraticcurveto{\pgfqpoint{10.444900in}{0.352751in}}{\pgfqpoint{10.453275in}{0.352751in}}%
\pgfpathquadraticcurveto{\pgfqpoint{10.461561in}{0.352751in}}{\pgfqpoint{10.466388in}{0.359271in}}%
\pgfpathquadraticcurveto{\pgfqpoint{10.471234in}{0.365810in}}{\pgfqpoint{10.471234in}{0.377067in}}%
\pgfpathquadraticcurveto{\pgfqpoint{10.471234in}{0.388271in}}{\pgfqpoint{10.466388in}{0.394827in}}%
\pgfpathquadraticcurveto{\pgfqpoint{10.461561in}{0.401384in}}{\pgfqpoint{10.453275in}{0.401384in}}%
\pgfpathlineto{\pgfqpoint{10.453275in}{0.401384in}}%
\pgfpathclose%
\pgfpathmoveto{\pgfqpoint{10.453275in}{0.410156in}}%
\pgfpathquadraticcurveto{\pgfqpoint{10.466785in}{0.410156in}}{\pgfqpoint{10.474494in}{0.401366in}}%
\pgfpathquadraticcurveto{\pgfqpoint{10.482221in}{0.392594in}}{\pgfqpoint{10.482221in}{0.377067in}}%
\pgfpathquadraticcurveto{\pgfqpoint{10.482221in}{0.361595in}}{\pgfqpoint{10.474494in}{0.352769in}}%
\pgfpathquadraticcurveto{\pgfqpoint{10.466785in}{0.343961in}}{\pgfqpoint{10.453275in}{0.343961in}}%
\pgfpathquadraticcurveto{\pgfqpoint{10.439712in}{0.343961in}}{\pgfqpoint{10.432021in}{0.352769in}}%
\pgfpathquadraticcurveto{\pgfqpoint{10.424348in}{0.361595in}}{\pgfqpoint{10.424348in}{0.377067in}}%
\pgfpathquadraticcurveto{\pgfqpoint{10.424348in}{0.392594in}}{\pgfqpoint{10.432021in}{0.401366in}}%
\pgfpathquadraticcurveto{\pgfqpoint{10.439712in}{0.410156in}}{\pgfqpoint{10.453275in}{0.410156in}}%
\pgfpathlineto{\pgfqpoint{10.453275in}{0.410156in}}%
\pgfpathclose%
\pgfpathmoveto{\pgfqpoint{10.499387in}{0.433193in}}%
\pgfpathlineto{\pgfqpoint{10.509743in}{0.433193in}}%
\pgfpathlineto{\pgfqpoint{10.509743in}{0.345600in}}%
\pgfpathlineto{\pgfqpoint{10.499387in}{0.345600in}}%
\pgfpathlineto{\pgfqpoint{10.499387in}{0.433193in}}%
\pgfpathlineto{\pgfqpoint{10.499387in}{0.433193in}}%
\pgfpathclose%
\pgfpathmoveto{\pgfqpoint{10.531412in}{0.433193in}}%
\pgfpathlineto{\pgfqpoint{10.541769in}{0.433193in}}%
\pgfpathlineto{\pgfqpoint{10.541769in}{0.345600in}}%
\pgfpathlineto{\pgfqpoint{10.531412in}{0.345600in}}%
\pgfpathlineto{\pgfqpoint{10.531412in}{0.433193in}}%
\pgfpathlineto{\pgfqpoint{10.531412in}{0.433193in}}%
\pgfpathclose%
\pgfpathmoveto{\pgfqpoint{10.587862in}{0.401384in}}%
\pgfpathquadraticcurveto{\pgfqpoint{10.579541in}{0.401384in}}{\pgfqpoint{10.574695in}{0.394881in}}%
\pgfpathquadraticcurveto{\pgfqpoint{10.569850in}{0.388379in}}{\pgfqpoint{10.569850in}{0.377067in}}%
\pgfpathquadraticcurveto{\pgfqpoint{10.569850in}{0.365756in}}{\pgfqpoint{10.574659in}{0.359253in}}%
\pgfpathquadraticcurveto{\pgfqpoint{10.579487in}{0.352751in}}{\pgfqpoint{10.587862in}{0.352751in}}%
\pgfpathquadraticcurveto{\pgfqpoint{10.596148in}{0.352751in}}{\pgfqpoint{10.600975in}{0.359271in}}%
\pgfpathquadraticcurveto{\pgfqpoint{10.605820in}{0.365810in}}{\pgfqpoint{10.605820in}{0.377067in}}%
\pgfpathquadraticcurveto{\pgfqpoint{10.605820in}{0.388271in}}{\pgfqpoint{10.600975in}{0.394827in}}%
\pgfpathquadraticcurveto{\pgfqpoint{10.596148in}{0.401384in}}{\pgfqpoint{10.587862in}{0.401384in}}%
\pgfpathlineto{\pgfqpoint{10.587862in}{0.401384in}}%
\pgfpathclose%
\pgfpathmoveto{\pgfqpoint{10.587862in}{0.410156in}}%
\pgfpathquadraticcurveto{\pgfqpoint{10.601371in}{0.410156in}}{\pgfqpoint{10.609081in}{0.401366in}}%
\pgfpathquadraticcurveto{\pgfqpoint{10.616808in}{0.392594in}}{\pgfqpoint{10.616808in}{0.377067in}}%
\pgfpathquadraticcurveto{\pgfqpoint{10.616808in}{0.361595in}}{\pgfqpoint{10.609081in}{0.352769in}}%
\pgfpathquadraticcurveto{\pgfqpoint{10.601371in}{0.343961in}}{\pgfqpoint{10.587862in}{0.343961in}}%
\pgfpathquadraticcurveto{\pgfqpoint{10.574299in}{0.343961in}}{\pgfqpoint{10.566608in}{0.352769in}}%
\pgfpathquadraticcurveto{\pgfqpoint{10.558935in}{0.361595in}}{\pgfqpoint{10.558935in}{0.377067in}}%
\pgfpathquadraticcurveto{\pgfqpoint{10.558935in}{0.392594in}}{\pgfqpoint{10.566608in}{0.401366in}}%
\pgfpathquadraticcurveto{\pgfqpoint{10.574299in}{0.410156in}}{\pgfqpoint{10.587862in}{0.410156in}}%
\pgfpathlineto{\pgfqpoint{10.587862in}{0.410156in}}%
\pgfpathclose%
\pgfpathmoveto{\pgfqpoint{10.632911in}{0.370475in}}%
\pgfpathlineto{\pgfqpoint{10.632911in}{0.408643in}}%
\pgfpathlineto{\pgfqpoint{10.643268in}{0.408643in}}%
\pgfpathlineto{\pgfqpoint{10.643268in}{0.370871in}}%
\pgfpathquadraticcurveto{\pgfqpoint{10.643268in}{0.361919in}}{\pgfqpoint{10.646744in}{0.357434in}}%
\pgfpathquadraticcurveto{\pgfqpoint{10.650238in}{0.352967in}}{\pgfqpoint{10.657227in}{0.352967in}}%
\pgfpathquadraticcurveto{\pgfqpoint{10.665603in}{0.352967in}}{\pgfqpoint{10.670466in}{0.358317in}}%
\pgfpathquadraticcurveto{\pgfqpoint{10.675347in}{0.363666in}}{\pgfqpoint{10.675347in}{0.372906in}}%
\pgfpathlineto{\pgfqpoint{10.675347in}{0.408643in}}%
\pgfpathlineto{\pgfqpoint{10.685704in}{0.408643in}}%
\pgfpathlineto{\pgfqpoint{10.685704in}{0.345600in}}%
\pgfpathlineto{\pgfqpoint{10.675347in}{0.345600in}}%
\pgfpathlineto{\pgfqpoint{10.675347in}{0.355291in}}%
\pgfpathquadraticcurveto{\pgfqpoint{10.671583in}{0.349545in}}{\pgfqpoint{10.666593in}{0.346753in}}%
\pgfpathquadraticcurveto{\pgfqpoint{10.661622in}{0.343961in}}{\pgfqpoint{10.655030in}{0.343961in}}%
\pgfpathquadraticcurveto{\pgfqpoint{10.644168in}{0.343961in}}{\pgfqpoint{10.638530in}{0.350715in}}%
\pgfpathquadraticcurveto{\pgfqpoint{10.632911in}{0.357470in}}{\pgfqpoint{10.632911in}{0.370475in}}%
\pgfpathlineto{\pgfqpoint{10.632911in}{0.370475in}}%
\pgfpathclose%
\pgfpathmoveto{\pgfqpoint{10.658974in}{0.410156in}}%
\pgfpathlineto{\pgfqpoint{10.658974in}{0.410156in}}%
\pgfpathlineto{\pgfqpoint{10.658974in}{0.410156in}}%
\pgfpathclose%
\pgfpathmoveto{\pgfqpoint{10.717280in}{0.426547in}}%
\pgfpathlineto{\pgfqpoint{10.717280in}{0.408643in}}%
\pgfpathlineto{\pgfqpoint{10.738606in}{0.408643in}}%
\pgfpathlineto{\pgfqpoint{10.738606in}{0.400591in}}%
\pgfpathlineto{\pgfqpoint{10.717280in}{0.400591in}}%
\pgfpathlineto{\pgfqpoint{10.717280in}{0.366368in}}%
\pgfpathquadraticcurveto{\pgfqpoint{10.717280in}{0.358659in}}{\pgfqpoint{10.719387in}{0.356461in}}%
\pgfpathquadraticcurveto{\pgfqpoint{10.721494in}{0.354264in}}{\pgfqpoint{10.727979in}{0.354264in}}%
\pgfpathlineto{\pgfqpoint{10.738606in}{0.354264in}}%
\pgfpathlineto{\pgfqpoint{10.738606in}{0.345600in}}%
\pgfpathlineto{\pgfqpoint{10.727979in}{0.345600in}}%
\pgfpathquadraticcurveto{\pgfqpoint{10.715983in}{0.345600in}}{\pgfqpoint{10.711426in}{0.350067in}}%
\pgfpathquadraticcurveto{\pgfqpoint{10.706868in}{0.354552in}}{\pgfqpoint{10.706868in}{0.366368in}}%
\pgfpathlineto{\pgfqpoint{10.706868in}{0.400591in}}%
\pgfpathlineto{\pgfqpoint{10.699267in}{0.400591in}}%
\pgfpathlineto{\pgfqpoint{10.699267in}{0.408643in}}%
\pgfpathlineto{\pgfqpoint{10.706868in}{0.408643in}}%
\pgfpathlineto{\pgfqpoint{10.706868in}{0.426547in}}%
\pgfpathlineto{\pgfqpoint{10.717280in}{0.426547in}}%
\pgfpathlineto{\pgfqpoint{10.717280in}{0.426547in}}%
\pgfpathclose%
\pgfpathmoveto{\pgfqpoint{10.825388in}{0.398970in}}%
\pgfpathquadraticcurveto{\pgfqpoint{10.823641in}{0.399979in}}{\pgfqpoint{10.821588in}{0.400447in}}%
\pgfpathquadraticcurveto{\pgfqpoint{10.819535in}{0.400933in}}{\pgfqpoint{10.817067in}{0.400933in}}%
\pgfpathquadraticcurveto{\pgfqpoint{10.808277in}{0.400933in}}{\pgfqpoint{10.803576in}{0.395223in}}%
\pgfpathquadraticcurveto{\pgfqpoint{10.798875in}{0.389514in}}{\pgfqpoint{10.798875in}{0.378814in}}%
\pgfpathlineto{\pgfqpoint{10.798875in}{0.345600in}}%
\pgfpathlineto{\pgfqpoint{10.788464in}{0.345600in}}%
\pgfpathlineto{\pgfqpoint{10.788464in}{0.408643in}}%
\pgfpathlineto{\pgfqpoint{10.798875in}{0.408643in}}%
\pgfpathlineto{\pgfqpoint{10.798875in}{0.398844in}}%
\pgfpathquadraticcurveto{\pgfqpoint{10.802153in}{0.404590in}}{\pgfqpoint{10.807376in}{0.407364in}}%
\pgfpathquadraticcurveto{\pgfqpoint{10.812618in}{0.410156in}}{\pgfqpoint{10.820111in}{0.410156in}}%
\pgfpathquadraticcurveto{\pgfqpoint{10.821174in}{0.410156in}}{\pgfqpoint{10.822470in}{0.410011in}}%
\pgfpathquadraticcurveto{\pgfqpoint{10.823767in}{0.409885in}}{\pgfqpoint{10.825334in}{0.409597in}}%
\pgfpathlineto{\pgfqpoint{10.825388in}{0.398970in}}%
\pgfpathlineto{\pgfqpoint{10.825388in}{0.398970in}}%
\pgfpathclose%
\pgfpathmoveto{\pgfqpoint{10.887638in}{0.379715in}}%
\pgfpathlineto{\pgfqpoint{10.887638in}{0.374654in}}%
\pgfpathlineto{\pgfqpoint{10.840014in}{0.374654in}}%
\pgfpathquadraticcurveto{\pgfqpoint{10.840699in}{0.363954in}}{\pgfqpoint{10.846463in}{0.358353in}}%
\pgfpathquadraticcurveto{\pgfqpoint{10.852227in}{0.352751in}}{\pgfqpoint{10.862530in}{0.352751in}}%
\pgfpathquadraticcurveto{\pgfqpoint{10.868492in}{0.352751in}}{\pgfqpoint{10.874093in}{0.354210in}}%
\pgfpathquadraticcurveto{\pgfqpoint{10.879695in}{0.355669in}}{\pgfqpoint{10.885225in}{0.358605in}}%
\pgfpathlineto{\pgfqpoint{10.885225in}{0.348806in}}%
\pgfpathquadraticcurveto{\pgfqpoint{10.879641in}{0.346447in}}{\pgfqpoint{10.873787in}{0.345204in}}%
\pgfpathquadraticcurveto{\pgfqpoint{10.867933in}{0.343961in}}{\pgfqpoint{10.861917in}{0.343961in}}%
\pgfpathquadraticcurveto{\pgfqpoint{10.846823in}{0.343961in}}{\pgfqpoint{10.838015in}{0.352733in}}%
\pgfpathquadraticcurveto{\pgfqpoint{10.829207in}{0.361523in}}{\pgfqpoint{10.829207in}{0.376509in}}%
\pgfpathquadraticcurveto{\pgfqpoint{10.829207in}{0.391981in}}{\pgfqpoint{10.837565in}{0.401059in}}%
\pgfpathquadraticcurveto{\pgfqpoint{10.845922in}{0.410156in}}{\pgfqpoint{10.860116in}{0.410156in}}%
\pgfpathquadraticcurveto{\pgfqpoint{10.872832in}{0.410156in}}{\pgfqpoint{10.880235in}{0.401960in}}%
\pgfpathquadraticcurveto{\pgfqpoint{10.887638in}{0.393783in}}{\pgfqpoint{10.887638in}{0.379715in}}%
\pgfpathlineto{\pgfqpoint{10.887638in}{0.379715in}}%
\pgfpathclose%
\pgfpathmoveto{\pgfqpoint{10.877281in}{0.382759in}}%
\pgfpathquadraticcurveto{\pgfqpoint{10.877173in}{0.391243in}}{\pgfqpoint{10.872526in}{0.396304in}}%
\pgfpathquadraticcurveto{\pgfqpoint{10.867879in}{0.401384in}}{\pgfqpoint{10.860224in}{0.401384in}}%
\pgfpathquadraticcurveto{\pgfqpoint{10.851560in}{0.401384in}}{\pgfqpoint{10.846355in}{0.396484in}}%
\pgfpathquadraticcurveto{\pgfqpoint{10.841149in}{0.391585in}}{\pgfqpoint{10.840357in}{0.382687in}}%
\pgfpathlineto{\pgfqpoint{10.877281in}{0.382759in}}%
\pgfpathlineto{\pgfqpoint{10.877281in}{0.382759in}}%
\pgfpathclose%
\pgfpathmoveto{\pgfqpoint{10.946124in}{0.377860in}}%
\pgfpathquadraticcurveto{\pgfqpoint{10.946124in}{0.389117in}}{\pgfqpoint{10.941477in}{0.395296in}}%
\pgfpathquadraticcurveto{\pgfqpoint{10.936848in}{0.401492in}}{\pgfqpoint{10.928454in}{0.401492in}}%
\pgfpathquadraticcurveto{\pgfqpoint{10.920132in}{0.401492in}}{\pgfqpoint{10.915485in}{0.395296in}}%
\pgfpathquadraticcurveto{\pgfqpoint{10.910838in}{0.389117in}}{\pgfqpoint{10.910838in}{0.377860in}}%
\pgfpathquadraticcurveto{\pgfqpoint{10.910838in}{0.366656in}}{\pgfqpoint{10.915485in}{0.360460in}}%
\pgfpathquadraticcurveto{\pgfqpoint{10.920132in}{0.354264in}}{\pgfqpoint{10.928454in}{0.354264in}}%
\pgfpathquadraticcurveto{\pgfqpoint{10.936848in}{0.354264in}}{\pgfqpoint{10.941477in}{0.360460in}}%
\pgfpathquadraticcurveto{\pgfqpoint{10.946124in}{0.366656in}}{\pgfqpoint{10.946124in}{0.377860in}}%
\pgfpathlineto{\pgfqpoint{10.946124in}{0.377860in}}%
\pgfpathclose%
\pgfpathmoveto{\pgfqpoint{10.956481in}{0.353417in}}%
\pgfpathquadraticcurveto{\pgfqpoint{10.956481in}{0.337332in}}{\pgfqpoint{10.949330in}{0.329479in}}%
\pgfpathquadraticcurveto{\pgfqpoint{10.942197in}{0.321626in}}{\pgfqpoint{10.927445in}{0.321626in}}%
\pgfpathquadraticcurveto{\pgfqpoint{10.921988in}{0.321626in}}{\pgfqpoint{10.917142in}{0.322436in}}%
\pgfpathquadraticcurveto{\pgfqpoint{10.912297in}{0.323247in}}{\pgfqpoint{10.907740in}{0.324940in}}%
\pgfpathlineto{\pgfqpoint{10.907740in}{0.335009in}}%
\pgfpathquadraticcurveto{\pgfqpoint{10.912297in}{0.332541in}}{\pgfqpoint{10.916746in}{0.331370in}}%
\pgfpathquadraticcurveto{\pgfqpoint{10.921195in}{0.330182in}}{\pgfqpoint{10.925806in}{0.330182in}}%
\pgfpathquadraticcurveto{\pgfqpoint{10.936001in}{0.330182in}}{\pgfqpoint{10.941063in}{0.335495in}}%
\pgfpathquadraticcurveto{\pgfqpoint{10.946124in}{0.340809in}}{\pgfqpoint{10.946124in}{0.351562in}}%
\pgfpathlineto{\pgfqpoint{10.946124in}{0.356695in}}%
\pgfpathquadraticcurveto{\pgfqpoint{10.942918in}{0.351112in}}{\pgfqpoint{10.937910in}{0.348356in}}%
\pgfpathquadraticcurveto{\pgfqpoint{10.932903in}{0.345600in}}{\pgfqpoint{10.925914in}{0.345600in}}%
\pgfpathquadraticcurveto{\pgfqpoint{10.914332in}{0.345600in}}{\pgfqpoint{10.907236in}{0.354426in}}%
\pgfpathquadraticcurveto{\pgfqpoint{10.900139in}{0.363270in}}{\pgfqpoint{10.900139in}{0.377860in}}%
\pgfpathquadraticcurveto{\pgfqpoint{10.900139in}{0.392486in}}{\pgfqpoint{10.907236in}{0.401312in}}%
\pgfpathquadraticcurveto{\pgfqpoint{10.914332in}{0.410156in}}{\pgfqpoint{10.925914in}{0.410156in}}%
\pgfpathquadraticcurveto{\pgfqpoint{10.932903in}{0.410156in}}{\pgfqpoint{10.937910in}{0.407400in}}%
\pgfpathquadraticcurveto{\pgfqpoint{10.942918in}{0.404644in}}{\pgfqpoint{10.946124in}{0.399078in}}%
\pgfpathlineto{\pgfqpoint{10.946124in}{0.408643in}}%
\pgfpathlineto{\pgfqpoint{10.956481in}{0.408643in}}%
\pgfpathlineto{\pgfqpoint{10.956481in}{0.353417in}}%
\pgfpathlineto{\pgfqpoint{10.956481in}{0.353417in}}%
\pgfpathclose%
\pgfpathmoveto{\pgfqpoint{10.977825in}{0.408643in}}%
\pgfpathlineto{\pgfqpoint{10.988182in}{0.408643in}}%
\pgfpathlineto{\pgfqpoint{10.988182in}{0.345600in}}%
\pgfpathlineto{\pgfqpoint{10.977825in}{0.345600in}}%
\pgfpathlineto{\pgfqpoint{10.977825in}{0.408643in}}%
\pgfpathlineto{\pgfqpoint{10.977825in}{0.408643in}}%
\pgfpathclose%
\pgfpathmoveto{\pgfqpoint{10.977825in}{0.433193in}}%
\pgfpathlineto{\pgfqpoint{10.988182in}{0.433193in}}%
\pgfpathlineto{\pgfqpoint{10.988182in}{0.420062in}}%
\pgfpathlineto{\pgfqpoint{10.977825in}{0.420062in}}%
\pgfpathlineto{\pgfqpoint{10.977825in}{0.433193in}}%
\pgfpathlineto{\pgfqpoint{10.977825in}{0.433193in}}%
\pgfpathclose%
\pgfpathmoveto{\pgfqpoint{11.034275in}{0.401384in}}%
\pgfpathquadraticcurveto{\pgfqpoint{11.025954in}{0.401384in}}{\pgfqpoint{11.021109in}{0.394881in}}%
\pgfpathquadraticcurveto{\pgfqpoint{11.016263in}{0.388379in}}{\pgfqpoint{11.016263in}{0.377067in}}%
\pgfpathquadraticcurveto{\pgfqpoint{11.016263in}{0.365756in}}{\pgfqpoint{11.021072in}{0.359253in}}%
\pgfpathquadraticcurveto{\pgfqpoint{11.025900in}{0.352751in}}{\pgfqpoint{11.034275in}{0.352751in}}%
\pgfpathquadraticcurveto{\pgfqpoint{11.042561in}{0.352751in}}{\pgfqpoint{11.047388in}{0.359271in}}%
\pgfpathquadraticcurveto{\pgfqpoint{11.052234in}{0.365810in}}{\pgfqpoint{11.052234in}{0.377067in}}%
\pgfpathquadraticcurveto{\pgfqpoint{11.052234in}{0.388271in}}{\pgfqpoint{11.047388in}{0.394827in}}%
\pgfpathquadraticcurveto{\pgfqpoint{11.042561in}{0.401384in}}{\pgfqpoint{11.034275in}{0.401384in}}%
\pgfpathlineto{\pgfqpoint{11.034275in}{0.401384in}}%
\pgfpathclose%
\pgfpathmoveto{\pgfqpoint{11.034275in}{0.410156in}}%
\pgfpathquadraticcurveto{\pgfqpoint{11.047785in}{0.410156in}}{\pgfqpoint{11.055494in}{0.401366in}}%
\pgfpathquadraticcurveto{\pgfqpoint{11.063221in}{0.392594in}}{\pgfqpoint{11.063221in}{0.377067in}}%
\pgfpathquadraticcurveto{\pgfqpoint{11.063221in}{0.361595in}}{\pgfqpoint{11.055494in}{0.352769in}}%
\pgfpathquadraticcurveto{\pgfqpoint{11.047785in}{0.343961in}}{\pgfqpoint{11.034275in}{0.343961in}}%
\pgfpathquadraticcurveto{\pgfqpoint{11.020712in}{0.343961in}}{\pgfqpoint{11.013021in}{0.352769in}}%
\pgfpathquadraticcurveto{\pgfqpoint{11.005348in}{0.361595in}}{\pgfqpoint{11.005348in}{0.377067in}}%
\pgfpathquadraticcurveto{\pgfqpoint{11.005348in}{0.392594in}}{\pgfqpoint{11.013021in}{0.401366in}}%
\pgfpathquadraticcurveto{\pgfqpoint{11.020712in}{0.410156in}}{\pgfqpoint{11.034275in}{0.410156in}}%
\pgfpathlineto{\pgfqpoint{11.034275in}{0.410156in}}%
\pgfpathclose%
\pgfpathmoveto{\pgfqpoint{11.132802in}{0.383660in}}%
\pgfpathlineto{\pgfqpoint{11.132802in}{0.345600in}}%
\pgfpathlineto{\pgfqpoint{11.122445in}{0.345600in}}%
\pgfpathlineto{\pgfqpoint{11.122445in}{0.383317in}}%
\pgfpathquadraticcurveto{\pgfqpoint{11.122445in}{0.392269in}}{\pgfqpoint{11.118951in}{0.396700in}}%
\pgfpathquadraticcurveto{\pgfqpoint{11.115456in}{0.401149in}}{\pgfqpoint{11.108485in}{0.401149in}}%
\pgfpathquadraticcurveto{\pgfqpoint{11.100092in}{0.401149in}}{\pgfqpoint{11.095247in}{0.395800in}}%
\pgfpathquadraticcurveto{\pgfqpoint{11.090401in}{0.390468in}}{\pgfqpoint{11.090401in}{0.381228in}}%
\pgfpathlineto{\pgfqpoint{11.090401in}{0.345600in}}%
\pgfpathlineto{\pgfqpoint{11.079990in}{0.345600in}}%
\pgfpathlineto{\pgfqpoint{11.079990in}{0.408643in}}%
\pgfpathlineto{\pgfqpoint{11.090401in}{0.408643in}}%
\pgfpathlineto{\pgfqpoint{11.090401in}{0.398844in}}%
\pgfpathquadraticcurveto{\pgfqpoint{11.094130in}{0.404536in}}{\pgfqpoint{11.099155in}{0.407346in}}%
\pgfpathquadraticcurveto{\pgfqpoint{11.104199in}{0.410156in}}{\pgfqpoint{11.110791in}{0.410156in}}%
\pgfpathquadraticcurveto{\pgfqpoint{11.121652in}{0.410156in}}{\pgfqpoint{11.127218in}{0.403437in}}%
\pgfpathquadraticcurveto{\pgfqpoint{11.132802in}{0.396718in}}{\pgfqpoint{11.132802in}{0.383660in}}%
\pgfpathlineto{\pgfqpoint{11.132802in}{0.383660in}}%
\pgfpathclose%
\pgfpathmoveto{\pgfqpoint{11.193629in}{0.406787in}}%
\pgfpathlineto{\pgfqpoint{11.193629in}{0.396989in}}%
\pgfpathquadraticcurveto{\pgfqpoint{11.189252in}{0.399240in}}{\pgfqpoint{11.184515in}{0.400357in}}%
\pgfpathquadraticcurveto{\pgfqpoint{11.179796in}{0.401492in}}{\pgfqpoint{11.174716in}{0.401492in}}%
\pgfpathquadraticcurveto{\pgfqpoint{11.167007in}{0.401492in}}{\pgfqpoint{11.163152in}{0.399132in}}%
\pgfpathquadraticcurveto{\pgfqpoint{11.159298in}{0.396773in}}{\pgfqpoint{11.159298in}{0.392035in}}%
\pgfpathquadraticcurveto{\pgfqpoint{11.159298in}{0.388433in}}{\pgfqpoint{11.162054in}{0.386380in}}%
\pgfpathquadraticcurveto{\pgfqpoint{11.164809in}{0.384326in}}{\pgfqpoint{11.173149in}{0.382471in}}%
\pgfpathlineto{\pgfqpoint{11.176697in}{0.381678in}}%
\pgfpathquadraticcurveto{\pgfqpoint{11.187721in}{0.379319in}}{\pgfqpoint{11.192368in}{0.375014in}}%
\pgfpathquadraticcurveto{\pgfqpoint{11.197015in}{0.370709in}}{\pgfqpoint{11.197015in}{0.363000in}}%
\pgfpathquadraticcurveto{\pgfqpoint{11.197015in}{0.354210in}}{\pgfqpoint{11.190063in}{0.349076in}}%
\pgfpathquadraticcurveto{\pgfqpoint{11.183110in}{0.343961in}}{\pgfqpoint{11.170952in}{0.343961in}}%
\pgfpathquadraticcurveto{\pgfqpoint{11.165890in}{0.343961in}}{\pgfqpoint{11.160396in}{0.344952in}}%
\pgfpathquadraticcurveto{\pgfqpoint{11.154903in}{0.345942in}}{\pgfqpoint{11.148833in}{0.347906in}}%
\pgfpathlineto{\pgfqpoint{11.148833in}{0.358605in}}%
\pgfpathquadraticcurveto{\pgfqpoint{11.154579in}{0.355615in}}{\pgfqpoint{11.160144in}{0.354120in}}%
\pgfpathquadraticcurveto{\pgfqpoint{11.165710in}{0.352643in}}{\pgfqpoint{11.171186in}{0.352643in}}%
\pgfpathquadraticcurveto{\pgfqpoint{11.178499in}{0.352643in}}{\pgfqpoint{11.182425in}{0.355146in}}%
\pgfpathquadraticcurveto{\pgfqpoint{11.186370in}{0.357650in}}{\pgfqpoint{11.186370in}{0.362207in}}%
\pgfpathquadraticcurveto{\pgfqpoint{11.186370in}{0.366422in}}{\pgfqpoint{11.183524in}{0.368674in}}%
\pgfpathquadraticcurveto{\pgfqpoint{11.180696in}{0.370925in}}{\pgfqpoint{11.171060in}{0.373014in}}%
\pgfpathlineto{\pgfqpoint{11.167457in}{0.373861in}}%
\pgfpathquadraticcurveto{\pgfqpoint{11.157839in}{0.375878in}}{\pgfqpoint{11.153552in}{0.380075in}}%
\pgfpathquadraticcurveto{\pgfqpoint{11.149283in}{0.384272in}}{\pgfqpoint{11.149283in}{0.391585in}}%
\pgfpathquadraticcurveto{\pgfqpoint{11.149283in}{0.400483in}}{\pgfqpoint{11.155587in}{0.405310in}}%
\pgfpathquadraticcurveto{\pgfqpoint{11.161891in}{0.410156in}}{\pgfqpoint{11.173491in}{0.410156in}}%
\pgfpathquadraticcurveto{\pgfqpoint{11.179219in}{0.410156in}}{\pgfqpoint{11.184281in}{0.409309in}}%
\pgfpathquadraticcurveto{\pgfqpoint{11.189360in}{0.408480in}}{\pgfqpoint{11.193629in}{0.406787in}}%
\pgfpathlineto{\pgfqpoint{11.193629in}{0.406787in}}%
\pgfpathclose%
\pgfpathmoveto{\pgfqpoint{11.214955in}{0.359902in}}%
\pgfpathlineto{\pgfqpoint{11.226843in}{0.359902in}}%
\pgfpathlineto{\pgfqpoint{11.226843in}{0.345600in}}%
\pgfpathlineto{\pgfqpoint{11.214955in}{0.345600in}}%
\pgfpathlineto{\pgfqpoint{11.214955in}{0.359902in}}%
\pgfpathlineto{\pgfqpoint{11.214955in}{0.359902in}}%
\pgfpathclose%
\pgfusepath{fill}%
\end{pgfscope}%
\end{pgfpicture}%
\makeatother%
\endgroup%

\end{lrbox}
\noindent\resizebox{!}{6.85in}{\usebox{\writeupfigure}}
\end{figure}
\end{landscape}
\clearpage
\begin{landscape}
\begin{figure}[p]
\centering
\caption{TWFE refitted on estimator-specific aligned samples}
\label{fig:aligned}
\begin{lrbox}{\writeupfigure}
%% Creator: Matplotlib, PGF backend
%%
%% To include the figure in your LaTeX document, write
%%   \input{<filename>.pgf}
%%
%% Make sure the required packages are loaded in your preamble
%%   \usepackage{pgf}
%%
%% Also ensure that all the required font packages are loaded; for instance,
%% the lmodern package is sometimes necessary when using math font.
%%   \usepackage{lmodern}
%%
%% Figures using additional raster images can only be included by \input if
%% they are in the same directory as the main LaTeX file. For loading figures
%% from other directories you can use the `import` package
%%   \usepackage{import}
%%
%% and then include the figures with
%%   \import{<path to file>}{<filename>.pgf}
%%
%% Matplotlib used the following preamble
%%   \def\mathdefault#1{#1}
%%   \everymath=\expandafter{\the\everymath\displaystyle}
%%   \IfFileExists{scrextend.sty}{
%%     \usepackage[fontsize=9.000000pt]{scrextend}
%%   }{
%%     \renewcommand{\normalsize}{\fontsize{9.000000}{10.800000}\selectfont}
%%     \normalsize
%%   }
%%   
%%   \makeatletter\@ifpackageloaded{underscore}{}{\usepackage[strings]{underscore}}\makeatother
%%
\begingroup%
\makeatletter%
\begin{pgfpicture}%
\pgfpathrectangle{\pgfpointorigin}{\pgfqpoint{14.000000in}{9.600000in}}%
\pgfusepath{use as bounding box, clip}%
\begin{pgfscope}%
\pgfsetbuttcap%
\pgfsetmiterjoin%
\definecolor{currentfill}{rgb}{1.000000,1.000000,1.000000}%
\pgfsetfillcolor{currentfill}%
\pgfsetlinewidth{0.000000pt}%
\definecolor{currentstroke}{rgb}{1.000000,1.000000,1.000000}%
\pgfsetstrokecolor{currentstroke}%
\pgfsetdash{}{0pt}%
\pgfpathmoveto{\pgfqpoint{0.000000in}{0.000000in}}%
\pgfpathlineto{\pgfqpoint{14.000000in}{0.000000in}}%
\pgfpathlineto{\pgfqpoint{14.000000in}{9.600000in}}%
\pgfpathlineto{\pgfqpoint{0.000000in}{9.600000in}}%
\pgfpathlineto{\pgfqpoint{0.000000in}{0.000000in}}%
\pgfpathclose%
\pgfusepath{fill}%
\end{pgfscope}%
\begin{pgfscope}%
\pgfsetbuttcap%
\pgfsetmiterjoin%
\definecolor{currentfill}{rgb}{1.000000,1.000000,1.000000}%
\pgfsetfillcolor{currentfill}%
\pgfsetlinewidth{0.000000pt}%
\definecolor{currentstroke}{rgb}{0.000000,0.000000,0.000000}%
\pgfsetstrokecolor{currentstroke}%
\pgfsetstrokeopacity{0.000000}%
\pgfsetdash{}{0pt}%
\pgfpathmoveto{\pgfqpoint{0.977909in}{2.928000in}}%
\pgfpathlineto{\pgfqpoint{6.641909in}{2.928000in}}%
\pgfpathlineto{\pgfqpoint{6.641909in}{8.592000in}}%
\pgfpathlineto{\pgfqpoint{0.977909in}{8.592000in}}%
\pgfpathlineto{\pgfqpoint{0.977909in}{2.928000in}}%
\pgfpathclose%
\pgfusepath{fill}%
\end{pgfscope}%
\begin{pgfscope}%
\pgfpathrectangle{\pgfqpoint{0.977909in}{2.928000in}}{\pgfqpoint{5.664000in}{5.664000in}}%
\pgfusepath{clip}%
\pgfsetrectcap%
\pgfsetroundjoin%
\pgfsetlinewidth{0.602250pt}%
\definecolor{currentstroke}{rgb}{0.898039,0.913725,0.929412}%
\pgfsetstrokecolor{currentstroke}%
\pgfsetdash{}{0pt}%
\pgfpathmoveto{\pgfqpoint{0.977909in}{2.928000in}}%
\pgfpathlineto{\pgfqpoint{0.977909in}{8.592000in}}%
\pgfusepath{stroke}%
\end{pgfscope}%
\begin{pgfscope}%
\pgfsetbuttcap%
\pgfsetroundjoin%
\definecolor{currentfill}{rgb}{0.090196,0.168627,0.227451}%
\pgfsetfillcolor{currentfill}%
\pgfsetlinewidth{0.803000pt}%
\definecolor{currentstroke}{rgb}{0.090196,0.168627,0.227451}%
\pgfsetstrokecolor{currentstroke}%
\pgfsetdash{}{0pt}%
\pgfsys@defobject{currentmarker}{\pgfqpoint{0.000000in}{-0.048611in}}{\pgfqpoint{0.000000in}{0.000000in}}{%
\pgfpathmoveto{\pgfqpoint{0.000000in}{0.000000in}}%
\pgfpathlineto{\pgfqpoint{0.000000in}{-0.048611in}}%
\pgfusepath{stroke,fill}%
}%
\begin{pgfscope}%
\pgfsys@transformshift{0.977909in}{2.928000in}%
\pgfsys@useobject{currentmarker}{}%
\end{pgfscope}%
\end{pgfscope}%
\begin{pgfscope}%
\pgfsetbuttcap%
\pgfsetroundjoin%
\definecolor{currentfill}{rgb}{0.090196,0.168627,0.227451}%
\pgfsetfillcolor{currentfill}%
\pgfsetlinewidth{0.000000pt}%
\definecolor{currentstroke}{rgb}{0.090196,0.168627,0.227451}%
\pgfsetstrokecolor{currentstroke}%
\pgfsetdash{}{0pt}%
\pgfpathmoveto{\pgfqpoint{0.836151in}{2.780182in}}%
\pgfpathlineto{\pgfqpoint{0.914393in}{2.780182in}}%
\pgfpathlineto{\pgfqpoint{0.914393in}{2.769811in}}%
\pgfpathlineto{\pgfqpoint{0.836151in}{2.769811in}}%
\pgfpathlineto{\pgfqpoint{0.836151in}{2.780182in}}%
\pgfpathlineto{\pgfqpoint{0.836151in}{2.780182in}}%
\pgfpathclose%
\pgfpathmoveto{\pgfqpoint{0.941151in}{2.826940in}}%
\pgfpathlineto{\pgfqpoint{0.989550in}{2.826940in}}%
\pgfpathlineto{\pgfqpoint{0.989550in}{2.816549in}}%
\pgfpathlineto{\pgfqpoint{0.952440in}{2.816549in}}%
\pgfpathlineto{\pgfqpoint{0.952440in}{2.794225in}}%
\pgfpathquadraticcurveto{\pgfqpoint{0.955116in}{2.795143in}}{\pgfqpoint{0.957792in}{2.795592in}}%
\pgfpathquadraticcurveto{\pgfqpoint{0.960487in}{2.796041in}}{\pgfqpoint{0.963183in}{2.796041in}}%
\pgfpathquadraticcurveto{\pgfqpoint{0.978436in}{2.796041in}}{\pgfqpoint{0.987343in}{2.787682in}}%
\pgfpathquadraticcurveto{\pgfqpoint{0.996268in}{2.779323in}}{\pgfqpoint{0.996268in}{2.765045in}}%
\pgfpathquadraticcurveto{\pgfqpoint{0.996268in}{2.750338in}}{\pgfqpoint{0.987108in}{2.742174in}}%
\pgfpathquadraticcurveto{\pgfqpoint{0.977948in}{2.734030in}}{\pgfqpoint{0.961288in}{2.734030in}}%
\pgfpathquadraticcurveto{\pgfqpoint{0.955546in}{2.734030in}}{\pgfqpoint{0.949589in}{2.735006in}}%
\pgfpathquadraticcurveto{\pgfqpoint{0.943651in}{2.735983in}}{\pgfqpoint{0.937304in}{2.737936in}}%
\pgfpathlineto{\pgfqpoint{0.937304in}{2.750338in}}%
\pgfpathquadraticcurveto{\pgfqpoint{0.942792in}{2.747350in}}{\pgfqpoint{0.948651in}{2.745885in}}%
\pgfpathquadraticcurveto{\pgfqpoint{0.954511in}{2.744420in}}{\pgfqpoint{0.961034in}{2.744420in}}%
\pgfpathquadraticcurveto{\pgfqpoint{0.971600in}{2.744420in}}{\pgfqpoint{0.977753in}{2.749967in}}%
\pgfpathquadraticcurveto{\pgfqpoint{0.983925in}{2.755514in}}{\pgfqpoint{0.983925in}{2.765045in}}%
\pgfpathquadraticcurveto{\pgfqpoint{0.983925in}{2.774557in}}{\pgfqpoint{0.977753in}{2.780104in}}%
\pgfpathquadraticcurveto{\pgfqpoint{0.971600in}{2.785670in}}{\pgfqpoint{0.961034in}{2.785670in}}%
\pgfpathquadraticcurveto{\pgfqpoint{0.956093in}{2.785670in}}{\pgfqpoint{0.951171in}{2.784577in}}%
\pgfpathquadraticcurveto{\pgfqpoint{0.946268in}{2.783483in}}{\pgfqpoint{0.941151in}{2.781159in}}%
\pgfpathlineto{\pgfqpoint{0.941151in}{2.826940in}}%
\pgfpathlineto{\pgfqpoint{0.941151in}{2.826940in}}%
\pgfpathclose%
\pgfpathmoveto{\pgfqpoint{1.020546in}{2.751315in}}%
\pgfpathlineto{\pgfqpoint{1.033436in}{2.751315in}}%
\pgfpathlineto{\pgfqpoint{1.033436in}{2.735807in}}%
\pgfpathlineto{\pgfqpoint{1.020546in}{2.735807in}}%
\pgfpathlineto{\pgfqpoint{1.020546in}{2.751315in}}%
\pgfpathlineto{\pgfqpoint{1.020546in}{2.751315in}}%
\pgfpathclose%
\pgfpathmoveto{\pgfqpoint{1.086640in}{2.818815in}}%
\pgfpathquadraticcurveto{\pgfqpoint{1.077128in}{2.818815in}}{\pgfqpoint{1.072323in}{2.809440in}}%
\pgfpathquadraticcurveto{\pgfqpoint{1.067538in}{2.800084in}}{\pgfqpoint{1.067538in}{2.781276in}}%
\pgfpathquadraticcurveto{\pgfqpoint{1.067538in}{2.762545in}}{\pgfqpoint{1.072323in}{2.753170in}}%
\pgfpathquadraticcurveto{\pgfqpoint{1.077128in}{2.743795in}}{\pgfqpoint{1.086640in}{2.743795in}}%
\pgfpathquadraticcurveto{\pgfqpoint{1.096229in}{2.743795in}}{\pgfqpoint{1.101015in}{2.753170in}}%
\pgfpathquadraticcurveto{\pgfqpoint{1.105819in}{2.762545in}}{\pgfqpoint{1.105819in}{2.781276in}}%
\pgfpathquadraticcurveto{\pgfqpoint{1.105819in}{2.800084in}}{\pgfqpoint{1.101015in}{2.809440in}}%
\pgfpathquadraticcurveto{\pgfqpoint{1.096229in}{2.818815in}}{\pgfqpoint{1.086640in}{2.818815in}}%
\pgfpathlineto{\pgfqpoint{1.086640in}{2.818815in}}%
\pgfpathclose%
\pgfpathmoveto{\pgfqpoint{1.086640in}{2.828581in}}%
\pgfpathquadraticcurveto{\pgfqpoint{1.101972in}{2.828581in}}{\pgfqpoint{1.110058in}{2.816452in}}%
\pgfpathquadraticcurveto{\pgfqpoint{1.118143in}{2.804342in}}{\pgfqpoint{1.118143in}{2.781276in}}%
\pgfpathquadraticcurveto{\pgfqpoint{1.118143in}{2.758268in}}{\pgfqpoint{1.110058in}{2.746139in}}%
\pgfpathquadraticcurveto{\pgfqpoint{1.101972in}{2.734030in}}{\pgfqpoint{1.086640in}{2.734030in}}%
\pgfpathquadraticcurveto{\pgfqpoint{1.071327in}{2.734030in}}{\pgfqpoint{1.063241in}{2.746139in}}%
\pgfpathquadraticcurveto{\pgfqpoint{1.055155in}{2.758268in}}{\pgfqpoint{1.055155in}{2.781276in}}%
\pgfpathquadraticcurveto{\pgfqpoint{1.055155in}{2.804342in}}{\pgfqpoint{1.063241in}{2.816452in}}%
\pgfpathquadraticcurveto{\pgfqpoint{1.071327in}{2.828581in}}{\pgfqpoint{1.086640in}{2.828581in}}%
\pgfpathlineto{\pgfqpoint{1.086640in}{2.828581in}}%
\pgfpathclose%
\pgfusepath{fill}%
\end{pgfscope}%
\begin{pgfscope}%
\pgfpathrectangle{\pgfqpoint{0.977909in}{2.928000in}}{\pgfqpoint{5.664000in}{5.664000in}}%
\pgfusepath{clip}%
\pgfsetrectcap%
\pgfsetroundjoin%
\pgfsetlinewidth{0.602250pt}%
\definecolor{currentstroke}{rgb}{0.898039,0.913725,0.929412}%
\pgfsetstrokecolor{currentstroke}%
\pgfsetdash{}{0pt}%
\pgfpathmoveto{\pgfqpoint{2.393909in}{2.928000in}}%
\pgfpathlineto{\pgfqpoint{2.393909in}{8.592000in}}%
\pgfusepath{stroke}%
\end{pgfscope}%
\begin{pgfscope}%
\pgfsetbuttcap%
\pgfsetroundjoin%
\definecolor{currentfill}{rgb}{0.090196,0.168627,0.227451}%
\pgfsetfillcolor{currentfill}%
\pgfsetlinewidth{0.803000pt}%
\definecolor{currentstroke}{rgb}{0.090196,0.168627,0.227451}%
\pgfsetstrokecolor{currentstroke}%
\pgfsetdash{}{0pt}%
\pgfsys@defobject{currentmarker}{\pgfqpoint{0.000000in}{-0.048611in}}{\pgfqpoint{0.000000in}{0.000000in}}{%
\pgfpathmoveto{\pgfqpoint{0.000000in}{0.000000in}}%
\pgfpathlineto{\pgfqpoint{0.000000in}{-0.048611in}}%
\pgfusepath{stroke,fill}%
}%
\begin{pgfscope}%
\pgfsys@transformshift{2.393909in}{2.928000in}%
\pgfsys@useobject{currentmarker}{}%
\end{pgfscope}%
\end{pgfscope}%
\begin{pgfscope}%
\pgfsetbuttcap%
\pgfsetroundjoin%
\definecolor{currentfill}{rgb}{0.090196,0.168627,0.227451}%
\pgfsetfillcolor{currentfill}%
\pgfsetlinewidth{0.000000pt}%
\definecolor{currentstroke}{rgb}{0.090196,0.168627,0.227451}%
\pgfsetstrokecolor{currentstroke}%
\pgfsetdash{}{0pt}%
\pgfpathmoveto{\pgfqpoint{2.252151in}{2.780182in}}%
\pgfpathlineto{\pgfqpoint{2.330393in}{2.780182in}}%
\pgfpathlineto{\pgfqpoint{2.330393in}{2.769811in}}%
\pgfpathlineto{\pgfqpoint{2.252151in}{2.769811in}}%
\pgfpathlineto{\pgfqpoint{2.252151in}{2.780182in}}%
\pgfpathlineto{\pgfqpoint{2.252151in}{2.780182in}}%
\pgfpathclose%
\pgfpathmoveto{\pgfqpoint{2.367640in}{2.746178in}}%
\pgfpathlineto{\pgfqpoint{2.410667in}{2.746178in}}%
\pgfpathlineto{\pgfqpoint{2.410667in}{2.735807in}}%
\pgfpathlineto{\pgfqpoint{2.352815in}{2.735807in}}%
\pgfpathlineto{\pgfqpoint{2.352815in}{2.746178in}}%
\pgfpathquadraticcurveto{\pgfqpoint{2.359827in}{2.753444in}}{\pgfqpoint{2.371936in}{2.765670in}}%
\pgfpathquadraticcurveto{\pgfqpoint{2.384065in}{2.777916in}}{\pgfqpoint{2.387171in}{2.781471in}}%
\pgfpathquadraticcurveto{\pgfqpoint{2.393089in}{2.788112in}}{\pgfqpoint{2.395433in}{2.792721in}}%
\pgfpathquadraticcurveto{\pgfqpoint{2.397796in}{2.797331in}}{\pgfqpoint{2.397796in}{2.801784in}}%
\pgfpathquadraticcurveto{\pgfqpoint{2.397796in}{2.809049in}}{\pgfqpoint{2.392698in}{2.813620in}}%
\pgfpathquadraticcurveto{\pgfqpoint{2.387600in}{2.818209in}}{\pgfqpoint{2.379417in}{2.818209in}}%
\pgfpathquadraticcurveto{\pgfqpoint{2.373616in}{2.818209in}}{\pgfqpoint{2.367171in}{2.816198in}}%
\pgfpathquadraticcurveto{\pgfqpoint{2.360745in}{2.814186in}}{\pgfqpoint{2.353421in}{2.810084in}}%
\pgfpathlineto{\pgfqpoint{2.353421in}{2.822545in}}%
\pgfpathquadraticcurveto{\pgfqpoint{2.360862in}{2.825534in}}{\pgfqpoint{2.367327in}{2.827057in}}%
\pgfpathquadraticcurveto{\pgfqpoint{2.373811in}{2.828581in}}{\pgfqpoint{2.379183in}{2.828581in}}%
\pgfpathquadraticcurveto{\pgfqpoint{2.393343in}{2.828581in}}{\pgfqpoint{2.401761in}{2.821491in}}%
\pgfpathquadraticcurveto{\pgfqpoint{2.410179in}{2.814420in}}{\pgfqpoint{2.410179in}{2.802584in}}%
\pgfpathquadraticcurveto{\pgfqpoint{2.410179in}{2.796959in}}{\pgfqpoint{2.408069in}{2.791920in}}%
\pgfpathquadraticcurveto{\pgfqpoint{2.405979in}{2.786901in}}{\pgfqpoint{2.400413in}{2.780065in}}%
\pgfpathquadraticcurveto{\pgfqpoint{2.398890in}{2.778288in}}{\pgfqpoint{2.390706in}{2.769831in}}%
\pgfpathquadraticcurveto{\pgfqpoint{2.382542in}{2.761373in}}{\pgfqpoint{2.367640in}{2.746178in}}%
\pgfpathlineto{\pgfqpoint{2.367640in}{2.746178in}}%
\pgfpathclose%
\pgfpathmoveto{\pgfqpoint{2.436546in}{2.751315in}}%
\pgfpathlineto{\pgfqpoint{2.449436in}{2.751315in}}%
\pgfpathlineto{\pgfqpoint{2.449436in}{2.735807in}}%
\pgfpathlineto{\pgfqpoint{2.436546in}{2.735807in}}%
\pgfpathlineto{\pgfqpoint{2.436546in}{2.751315in}}%
\pgfpathlineto{\pgfqpoint{2.436546in}{2.751315in}}%
\pgfpathclose%
\pgfpathmoveto{\pgfqpoint{2.476409in}{2.826940in}}%
\pgfpathlineto{\pgfqpoint{2.524808in}{2.826940in}}%
\pgfpathlineto{\pgfqpoint{2.524808in}{2.816549in}}%
\pgfpathlineto{\pgfqpoint{2.487698in}{2.816549in}}%
\pgfpathlineto{\pgfqpoint{2.487698in}{2.794225in}}%
\pgfpathquadraticcurveto{\pgfqpoint{2.490374in}{2.795143in}}{\pgfqpoint{2.493050in}{2.795592in}}%
\pgfpathquadraticcurveto{\pgfqpoint{2.495745in}{2.796041in}}{\pgfqpoint{2.498440in}{2.796041in}}%
\pgfpathquadraticcurveto{\pgfqpoint{2.513694in}{2.796041in}}{\pgfqpoint{2.522600in}{2.787682in}}%
\pgfpathquadraticcurveto{\pgfqpoint{2.531526in}{2.779323in}}{\pgfqpoint{2.531526in}{2.765045in}}%
\pgfpathquadraticcurveto{\pgfqpoint{2.531526in}{2.750338in}}{\pgfqpoint{2.522366in}{2.742174in}}%
\pgfpathquadraticcurveto{\pgfqpoint{2.513206in}{2.734030in}}{\pgfqpoint{2.496546in}{2.734030in}}%
\pgfpathquadraticcurveto{\pgfqpoint{2.490804in}{2.734030in}}{\pgfqpoint{2.484847in}{2.735006in}}%
\pgfpathquadraticcurveto{\pgfqpoint{2.478909in}{2.735983in}}{\pgfqpoint{2.472561in}{2.737936in}}%
\pgfpathlineto{\pgfqpoint{2.472561in}{2.750338in}}%
\pgfpathquadraticcurveto{\pgfqpoint{2.478050in}{2.747350in}}{\pgfqpoint{2.483909in}{2.745885in}}%
\pgfpathquadraticcurveto{\pgfqpoint{2.489768in}{2.744420in}}{\pgfqpoint{2.496292in}{2.744420in}}%
\pgfpathquadraticcurveto{\pgfqpoint{2.506858in}{2.744420in}}{\pgfqpoint{2.513011in}{2.749967in}}%
\pgfpathquadraticcurveto{\pgfqpoint{2.519183in}{2.755514in}}{\pgfqpoint{2.519183in}{2.765045in}}%
\pgfpathquadraticcurveto{\pgfqpoint{2.519183in}{2.774557in}}{\pgfqpoint{2.513011in}{2.780104in}}%
\pgfpathquadraticcurveto{\pgfqpoint{2.506858in}{2.785670in}}{\pgfqpoint{2.496292in}{2.785670in}}%
\pgfpathquadraticcurveto{\pgfqpoint{2.491350in}{2.785670in}}{\pgfqpoint{2.486429in}{2.784577in}}%
\pgfpathquadraticcurveto{\pgfqpoint{2.481526in}{2.783483in}}{\pgfqpoint{2.476409in}{2.781159in}}%
\pgfpathlineto{\pgfqpoint{2.476409in}{2.826940in}}%
\pgfpathlineto{\pgfqpoint{2.476409in}{2.826940in}}%
\pgfpathclose%
\pgfusepath{fill}%
\end{pgfscope}%
\begin{pgfscope}%
\pgfpathrectangle{\pgfqpoint{0.977909in}{2.928000in}}{\pgfqpoint{5.664000in}{5.664000in}}%
\pgfusepath{clip}%
\pgfsetrectcap%
\pgfsetroundjoin%
\pgfsetlinewidth{0.602250pt}%
\definecolor{currentstroke}{rgb}{0.898039,0.913725,0.929412}%
\pgfsetstrokecolor{currentstroke}%
\pgfsetdash{}{0pt}%
\pgfpathmoveto{\pgfqpoint{3.809909in}{2.928000in}}%
\pgfpathlineto{\pgfqpoint{3.809909in}{8.592000in}}%
\pgfusepath{stroke}%
\end{pgfscope}%
\begin{pgfscope}%
\pgfsetbuttcap%
\pgfsetroundjoin%
\definecolor{currentfill}{rgb}{0.090196,0.168627,0.227451}%
\pgfsetfillcolor{currentfill}%
\pgfsetlinewidth{0.803000pt}%
\definecolor{currentstroke}{rgb}{0.090196,0.168627,0.227451}%
\pgfsetstrokecolor{currentstroke}%
\pgfsetdash{}{0pt}%
\pgfsys@defobject{currentmarker}{\pgfqpoint{0.000000in}{-0.048611in}}{\pgfqpoint{0.000000in}{0.000000in}}{%
\pgfpathmoveto{\pgfqpoint{0.000000in}{0.000000in}}%
\pgfpathlineto{\pgfqpoint{0.000000in}{-0.048611in}}%
\pgfusepath{stroke,fill}%
}%
\begin{pgfscope}%
\pgfsys@transformshift{3.809909in}{2.928000in}%
\pgfsys@useobject{currentmarker}{}%
\end{pgfscope}%
\end{pgfscope}%
\begin{pgfscope}%
\pgfsetbuttcap%
\pgfsetroundjoin%
\definecolor{currentfill}{rgb}{0.090196,0.168627,0.227451}%
\pgfsetfillcolor{currentfill}%
\pgfsetlinewidth{0.000000pt}%
\definecolor{currentstroke}{rgb}{0.090196,0.168627,0.227451}%
\pgfsetstrokecolor{currentstroke}%
\pgfsetdash{}{0pt}%
\pgfpathmoveto{\pgfqpoint{3.749636in}{2.818815in}}%
\pgfpathquadraticcurveto{\pgfqpoint{3.740124in}{2.818815in}}{\pgfqpoint{3.735319in}{2.809440in}}%
\pgfpathquadraticcurveto{\pgfqpoint{3.730534in}{2.800084in}}{\pgfqpoint{3.730534in}{2.781276in}}%
\pgfpathquadraticcurveto{\pgfqpoint{3.730534in}{2.762545in}}{\pgfqpoint{3.735319in}{2.753170in}}%
\pgfpathquadraticcurveto{\pgfqpoint{3.740124in}{2.743795in}}{\pgfqpoint{3.749636in}{2.743795in}}%
\pgfpathquadraticcurveto{\pgfqpoint{3.759225in}{2.743795in}}{\pgfqpoint{3.764011in}{2.753170in}}%
\pgfpathquadraticcurveto{\pgfqpoint{3.768815in}{2.762545in}}{\pgfqpoint{3.768815in}{2.781276in}}%
\pgfpathquadraticcurveto{\pgfqpoint{3.768815in}{2.800084in}}{\pgfqpoint{3.764011in}{2.809440in}}%
\pgfpathquadraticcurveto{\pgfqpoint{3.759225in}{2.818815in}}{\pgfqpoint{3.749636in}{2.818815in}}%
\pgfpathlineto{\pgfqpoint{3.749636in}{2.818815in}}%
\pgfpathclose%
\pgfpathmoveto{\pgfqpoint{3.749636in}{2.828581in}}%
\pgfpathquadraticcurveto{\pgfqpoint{3.764968in}{2.828581in}}{\pgfqpoint{3.773054in}{2.816452in}}%
\pgfpathquadraticcurveto{\pgfqpoint{3.781140in}{2.804342in}}{\pgfqpoint{3.781140in}{2.781276in}}%
\pgfpathquadraticcurveto{\pgfqpoint{3.781140in}{2.758268in}}{\pgfqpoint{3.773054in}{2.746139in}}%
\pgfpathquadraticcurveto{\pgfqpoint{3.764968in}{2.734030in}}{\pgfqpoint{3.749636in}{2.734030in}}%
\pgfpathquadraticcurveto{\pgfqpoint{3.734323in}{2.734030in}}{\pgfqpoint{3.726237in}{2.746139in}}%
\pgfpathquadraticcurveto{\pgfqpoint{3.718151in}{2.758268in}}{\pgfqpoint{3.718151in}{2.781276in}}%
\pgfpathquadraticcurveto{\pgfqpoint{3.718151in}{2.804342in}}{\pgfqpoint{3.726237in}{2.816452in}}%
\pgfpathquadraticcurveto{\pgfqpoint{3.734323in}{2.828581in}}{\pgfqpoint{3.749636in}{2.828581in}}%
\pgfpathlineto{\pgfqpoint{3.749636in}{2.828581in}}%
\pgfpathclose%
\pgfpathmoveto{\pgfqpoint{3.802800in}{2.751315in}}%
\pgfpathlineto{\pgfqpoint{3.815690in}{2.751315in}}%
\pgfpathlineto{\pgfqpoint{3.815690in}{2.735807in}}%
\pgfpathlineto{\pgfqpoint{3.802800in}{2.735807in}}%
\pgfpathlineto{\pgfqpoint{3.802800in}{2.751315in}}%
\pgfpathlineto{\pgfqpoint{3.802800in}{2.751315in}}%
\pgfpathclose%
\pgfpathmoveto{\pgfqpoint{3.868893in}{2.818815in}}%
\pgfpathquadraticcurveto{\pgfqpoint{3.859382in}{2.818815in}}{\pgfqpoint{3.854577in}{2.809440in}}%
\pgfpathquadraticcurveto{\pgfqpoint{3.849792in}{2.800084in}}{\pgfqpoint{3.849792in}{2.781276in}}%
\pgfpathquadraticcurveto{\pgfqpoint{3.849792in}{2.762545in}}{\pgfqpoint{3.854577in}{2.753170in}}%
\pgfpathquadraticcurveto{\pgfqpoint{3.859382in}{2.743795in}}{\pgfqpoint{3.868893in}{2.743795in}}%
\pgfpathquadraticcurveto{\pgfqpoint{3.878483in}{2.743795in}}{\pgfqpoint{3.883268in}{2.753170in}}%
\pgfpathquadraticcurveto{\pgfqpoint{3.888073in}{2.762545in}}{\pgfqpoint{3.888073in}{2.781276in}}%
\pgfpathquadraticcurveto{\pgfqpoint{3.888073in}{2.800084in}}{\pgfqpoint{3.883268in}{2.809440in}}%
\pgfpathquadraticcurveto{\pgfqpoint{3.878483in}{2.818815in}}{\pgfqpoint{3.868893in}{2.818815in}}%
\pgfpathlineto{\pgfqpoint{3.868893in}{2.818815in}}%
\pgfpathclose%
\pgfpathmoveto{\pgfqpoint{3.868893in}{2.828581in}}%
\pgfpathquadraticcurveto{\pgfqpoint{3.884225in}{2.828581in}}{\pgfqpoint{3.892311in}{2.816452in}}%
\pgfpathquadraticcurveto{\pgfqpoint{3.900397in}{2.804342in}}{\pgfqpoint{3.900397in}{2.781276in}}%
\pgfpathquadraticcurveto{\pgfqpoint{3.900397in}{2.758268in}}{\pgfqpoint{3.892311in}{2.746139in}}%
\pgfpathquadraticcurveto{\pgfqpoint{3.884225in}{2.734030in}}{\pgfqpoint{3.868893in}{2.734030in}}%
\pgfpathquadraticcurveto{\pgfqpoint{3.853581in}{2.734030in}}{\pgfqpoint{3.845495in}{2.746139in}}%
\pgfpathquadraticcurveto{\pgfqpoint{3.837409in}{2.758268in}}{\pgfqpoint{3.837409in}{2.781276in}}%
\pgfpathquadraticcurveto{\pgfqpoint{3.837409in}{2.804342in}}{\pgfqpoint{3.845495in}{2.816452in}}%
\pgfpathquadraticcurveto{\pgfqpoint{3.853581in}{2.828581in}}{\pgfqpoint{3.868893in}{2.828581in}}%
\pgfpathlineto{\pgfqpoint{3.868893in}{2.828581in}}%
\pgfpathclose%
\pgfusepath{fill}%
\end{pgfscope}%
\begin{pgfscope}%
\pgfpathrectangle{\pgfqpoint{0.977909in}{2.928000in}}{\pgfqpoint{5.664000in}{5.664000in}}%
\pgfusepath{clip}%
\pgfsetrectcap%
\pgfsetroundjoin%
\pgfsetlinewidth{0.602250pt}%
\definecolor{currentstroke}{rgb}{0.898039,0.913725,0.929412}%
\pgfsetstrokecolor{currentstroke}%
\pgfsetdash{}{0pt}%
\pgfpathmoveto{\pgfqpoint{5.225909in}{2.928000in}}%
\pgfpathlineto{\pgfqpoint{5.225909in}{8.592000in}}%
\pgfusepath{stroke}%
\end{pgfscope}%
\begin{pgfscope}%
\pgfsetbuttcap%
\pgfsetroundjoin%
\definecolor{currentfill}{rgb}{0.090196,0.168627,0.227451}%
\pgfsetfillcolor{currentfill}%
\pgfsetlinewidth{0.803000pt}%
\definecolor{currentstroke}{rgb}{0.090196,0.168627,0.227451}%
\pgfsetstrokecolor{currentstroke}%
\pgfsetdash{}{0pt}%
\pgfsys@defobject{currentmarker}{\pgfqpoint{0.000000in}{-0.048611in}}{\pgfqpoint{0.000000in}{0.000000in}}{%
\pgfpathmoveto{\pgfqpoint{0.000000in}{0.000000in}}%
\pgfpathlineto{\pgfqpoint{0.000000in}{-0.048611in}}%
\pgfusepath{stroke,fill}%
}%
\begin{pgfscope}%
\pgfsys@transformshift{5.225909in}{2.928000in}%
\pgfsys@useobject{currentmarker}{}%
\end{pgfscope}%
\end{pgfscope}%
\begin{pgfscope}%
\pgfsetbuttcap%
\pgfsetroundjoin%
\definecolor{currentfill}{rgb}{0.090196,0.168627,0.227451}%
\pgfsetfillcolor{currentfill}%
\pgfsetlinewidth{0.000000pt}%
\definecolor{currentstroke}{rgb}{0.090196,0.168627,0.227451}%
\pgfsetstrokecolor{currentstroke}%
\pgfsetdash{}{0pt}%
\pgfpathmoveto{\pgfqpoint{5.149893in}{2.746178in}}%
\pgfpathlineto{\pgfqpoint{5.192921in}{2.746178in}}%
\pgfpathlineto{\pgfqpoint{5.192921in}{2.735807in}}%
\pgfpathlineto{\pgfqpoint{5.135069in}{2.735807in}}%
\pgfpathlineto{\pgfqpoint{5.135069in}{2.746178in}}%
\pgfpathquadraticcurveto{\pgfqpoint{5.142081in}{2.753444in}}{\pgfqpoint{5.154190in}{2.765670in}}%
\pgfpathquadraticcurveto{\pgfqpoint{5.166319in}{2.777916in}}{\pgfqpoint{5.169425in}{2.781471in}}%
\pgfpathquadraticcurveto{\pgfqpoint{5.175343in}{2.788112in}}{\pgfqpoint{5.177686in}{2.792721in}}%
\pgfpathquadraticcurveto{\pgfqpoint{5.180050in}{2.797331in}}{\pgfqpoint{5.180050in}{2.801784in}}%
\pgfpathquadraticcurveto{\pgfqpoint{5.180050in}{2.809049in}}{\pgfqpoint{5.174952in}{2.813620in}}%
\pgfpathquadraticcurveto{\pgfqpoint{5.169854in}{2.818209in}}{\pgfqpoint{5.161671in}{2.818209in}}%
\pgfpathquadraticcurveto{\pgfqpoint{5.155870in}{2.818209in}}{\pgfqpoint{5.149425in}{2.816198in}}%
\pgfpathquadraticcurveto{\pgfqpoint{5.142999in}{2.814186in}}{\pgfqpoint{5.135675in}{2.810084in}}%
\pgfpathlineto{\pgfqpoint{5.135675in}{2.822545in}}%
\pgfpathquadraticcurveto{\pgfqpoint{5.143116in}{2.825534in}}{\pgfqpoint{5.149581in}{2.827057in}}%
\pgfpathquadraticcurveto{\pgfqpoint{5.156065in}{2.828581in}}{\pgfqpoint{5.161436in}{2.828581in}}%
\pgfpathquadraticcurveto{\pgfqpoint{5.175597in}{2.828581in}}{\pgfqpoint{5.184015in}{2.821491in}}%
\pgfpathquadraticcurveto{\pgfqpoint{5.192433in}{2.814420in}}{\pgfqpoint{5.192433in}{2.802584in}}%
\pgfpathquadraticcurveto{\pgfqpoint{5.192433in}{2.796959in}}{\pgfqpoint{5.190323in}{2.791920in}}%
\pgfpathquadraticcurveto{\pgfqpoint{5.188233in}{2.786901in}}{\pgfqpoint{5.182667in}{2.780065in}}%
\pgfpathquadraticcurveto{\pgfqpoint{5.181143in}{2.778288in}}{\pgfqpoint{5.172960in}{2.769831in}}%
\pgfpathquadraticcurveto{\pgfqpoint{5.164796in}{2.761373in}}{\pgfqpoint{5.149893in}{2.746178in}}%
\pgfpathlineto{\pgfqpoint{5.149893in}{2.746178in}}%
\pgfpathclose%
\pgfpathmoveto{\pgfqpoint{5.218800in}{2.751315in}}%
\pgfpathlineto{\pgfqpoint{5.231690in}{2.751315in}}%
\pgfpathlineto{\pgfqpoint{5.231690in}{2.735807in}}%
\pgfpathlineto{\pgfqpoint{5.218800in}{2.735807in}}%
\pgfpathlineto{\pgfqpoint{5.218800in}{2.751315in}}%
\pgfpathlineto{\pgfqpoint{5.218800in}{2.751315in}}%
\pgfpathclose%
\pgfpathmoveto{\pgfqpoint{5.258663in}{2.826940in}}%
\pgfpathlineto{\pgfqpoint{5.307061in}{2.826940in}}%
\pgfpathlineto{\pgfqpoint{5.307061in}{2.816549in}}%
\pgfpathlineto{\pgfqpoint{5.269952in}{2.816549in}}%
\pgfpathlineto{\pgfqpoint{5.269952in}{2.794225in}}%
\pgfpathquadraticcurveto{\pgfqpoint{5.272628in}{2.795143in}}{\pgfqpoint{5.275304in}{2.795592in}}%
\pgfpathquadraticcurveto{\pgfqpoint{5.277999in}{2.796041in}}{\pgfqpoint{5.280694in}{2.796041in}}%
\pgfpathquadraticcurveto{\pgfqpoint{5.295948in}{2.796041in}}{\pgfqpoint{5.304854in}{2.787682in}}%
\pgfpathquadraticcurveto{\pgfqpoint{5.313780in}{2.779323in}}{\pgfqpoint{5.313780in}{2.765045in}}%
\pgfpathquadraticcurveto{\pgfqpoint{5.313780in}{2.750338in}}{\pgfqpoint{5.304620in}{2.742174in}}%
\pgfpathquadraticcurveto{\pgfqpoint{5.295460in}{2.734030in}}{\pgfqpoint{5.278800in}{2.734030in}}%
\pgfpathquadraticcurveto{\pgfqpoint{5.273058in}{2.734030in}}{\pgfqpoint{5.267100in}{2.735006in}}%
\pgfpathquadraticcurveto{\pgfqpoint{5.261163in}{2.735983in}}{\pgfqpoint{5.254815in}{2.737936in}}%
\pgfpathlineto{\pgfqpoint{5.254815in}{2.750338in}}%
\pgfpathquadraticcurveto{\pgfqpoint{5.260304in}{2.747350in}}{\pgfqpoint{5.266163in}{2.745885in}}%
\pgfpathquadraticcurveto{\pgfqpoint{5.272022in}{2.744420in}}{\pgfqpoint{5.278546in}{2.744420in}}%
\pgfpathquadraticcurveto{\pgfqpoint{5.289112in}{2.744420in}}{\pgfqpoint{5.295265in}{2.749967in}}%
\pgfpathquadraticcurveto{\pgfqpoint{5.301436in}{2.755514in}}{\pgfqpoint{5.301436in}{2.765045in}}%
\pgfpathquadraticcurveto{\pgfqpoint{5.301436in}{2.774557in}}{\pgfqpoint{5.295265in}{2.780104in}}%
\pgfpathquadraticcurveto{\pgfqpoint{5.289112in}{2.785670in}}{\pgfqpoint{5.278546in}{2.785670in}}%
\pgfpathquadraticcurveto{\pgfqpoint{5.273604in}{2.785670in}}{\pgfqpoint{5.268683in}{2.784577in}}%
\pgfpathquadraticcurveto{\pgfqpoint{5.263780in}{2.783483in}}{\pgfqpoint{5.258663in}{2.781159in}}%
\pgfpathlineto{\pgfqpoint{5.258663in}{2.826940in}}%
\pgfpathlineto{\pgfqpoint{5.258663in}{2.826940in}}%
\pgfpathclose%
\pgfusepath{fill}%
\end{pgfscope}%
\begin{pgfscope}%
\pgfpathrectangle{\pgfqpoint{0.977909in}{2.928000in}}{\pgfqpoint{5.664000in}{5.664000in}}%
\pgfusepath{clip}%
\pgfsetrectcap%
\pgfsetroundjoin%
\pgfsetlinewidth{0.602250pt}%
\definecolor{currentstroke}{rgb}{0.898039,0.913725,0.929412}%
\pgfsetstrokecolor{currentstroke}%
\pgfsetdash{}{0pt}%
\pgfpathmoveto{\pgfqpoint{6.641909in}{2.928000in}}%
\pgfpathlineto{\pgfqpoint{6.641909in}{8.592000in}}%
\pgfusepath{stroke}%
\end{pgfscope}%
\begin{pgfscope}%
\pgfsetbuttcap%
\pgfsetroundjoin%
\definecolor{currentfill}{rgb}{0.090196,0.168627,0.227451}%
\pgfsetfillcolor{currentfill}%
\pgfsetlinewidth{0.803000pt}%
\definecolor{currentstroke}{rgb}{0.090196,0.168627,0.227451}%
\pgfsetstrokecolor{currentstroke}%
\pgfsetdash{}{0pt}%
\pgfsys@defobject{currentmarker}{\pgfqpoint{0.000000in}{-0.048611in}}{\pgfqpoint{0.000000in}{0.000000in}}{%
\pgfpathmoveto{\pgfqpoint{0.000000in}{0.000000in}}%
\pgfpathlineto{\pgfqpoint{0.000000in}{-0.048611in}}%
\pgfusepath{stroke,fill}%
}%
\begin{pgfscope}%
\pgfsys@transformshift{6.641909in}{2.928000in}%
\pgfsys@useobject{currentmarker}{}%
\end{pgfscope}%
\end{pgfscope}%
\begin{pgfscope}%
\pgfsetbuttcap%
\pgfsetroundjoin%
\definecolor{currentfill}{rgb}{0.090196,0.168627,0.227451}%
\pgfsetfillcolor{currentfill}%
\pgfsetlinewidth{0.000000pt}%
\definecolor{currentstroke}{rgb}{0.090196,0.168627,0.227451}%
\pgfsetstrokecolor{currentstroke}%
\pgfsetdash{}{0pt}%
\pgfpathmoveto{\pgfqpoint{6.555405in}{2.826940in}}%
\pgfpathlineto{\pgfqpoint{6.603804in}{2.826940in}}%
\pgfpathlineto{\pgfqpoint{6.603804in}{2.816549in}}%
\pgfpathlineto{\pgfqpoint{6.566694in}{2.816549in}}%
\pgfpathlineto{\pgfqpoint{6.566694in}{2.794225in}}%
\pgfpathquadraticcurveto{\pgfqpoint{6.569370in}{2.795143in}}{\pgfqpoint{6.572046in}{2.795592in}}%
\pgfpathquadraticcurveto{\pgfqpoint{6.574741in}{2.796041in}}{\pgfqpoint{6.577436in}{2.796041in}}%
\pgfpathquadraticcurveto{\pgfqpoint{6.592690in}{2.796041in}}{\pgfqpoint{6.601597in}{2.787682in}}%
\pgfpathquadraticcurveto{\pgfqpoint{6.610522in}{2.779323in}}{\pgfqpoint{6.610522in}{2.765045in}}%
\pgfpathquadraticcurveto{\pgfqpoint{6.610522in}{2.750338in}}{\pgfqpoint{6.601362in}{2.742174in}}%
\pgfpathquadraticcurveto{\pgfqpoint{6.592202in}{2.734030in}}{\pgfqpoint{6.575542in}{2.734030in}}%
\pgfpathquadraticcurveto{\pgfqpoint{6.569800in}{2.734030in}}{\pgfqpoint{6.563843in}{2.735006in}}%
\pgfpathquadraticcurveto{\pgfqpoint{6.557905in}{2.735983in}}{\pgfqpoint{6.551558in}{2.737936in}}%
\pgfpathlineto{\pgfqpoint{6.551558in}{2.750338in}}%
\pgfpathquadraticcurveto{\pgfqpoint{6.557046in}{2.747350in}}{\pgfqpoint{6.562905in}{2.745885in}}%
\pgfpathquadraticcurveto{\pgfqpoint{6.568765in}{2.744420in}}{\pgfqpoint{6.575288in}{2.744420in}}%
\pgfpathquadraticcurveto{\pgfqpoint{6.585854in}{2.744420in}}{\pgfqpoint{6.592007in}{2.749967in}}%
\pgfpathquadraticcurveto{\pgfqpoint{6.598179in}{2.755514in}}{\pgfqpoint{6.598179in}{2.765045in}}%
\pgfpathquadraticcurveto{\pgfqpoint{6.598179in}{2.774557in}}{\pgfqpoint{6.592007in}{2.780104in}}%
\pgfpathquadraticcurveto{\pgfqpoint{6.585854in}{2.785670in}}{\pgfqpoint{6.575288in}{2.785670in}}%
\pgfpathquadraticcurveto{\pgfqpoint{6.570347in}{2.785670in}}{\pgfqpoint{6.565425in}{2.784577in}}%
\pgfpathquadraticcurveto{\pgfqpoint{6.560522in}{2.783483in}}{\pgfqpoint{6.555405in}{2.781159in}}%
\pgfpathlineto{\pgfqpoint{6.555405in}{2.826940in}}%
\pgfpathlineto{\pgfqpoint{6.555405in}{2.826940in}}%
\pgfpathclose%
\pgfpathmoveto{\pgfqpoint{6.634800in}{2.751315in}}%
\pgfpathlineto{\pgfqpoint{6.647690in}{2.751315in}}%
\pgfpathlineto{\pgfqpoint{6.647690in}{2.735807in}}%
\pgfpathlineto{\pgfqpoint{6.634800in}{2.735807in}}%
\pgfpathlineto{\pgfqpoint{6.634800in}{2.751315in}}%
\pgfpathlineto{\pgfqpoint{6.634800in}{2.751315in}}%
\pgfpathclose%
\pgfpathmoveto{\pgfqpoint{6.700893in}{2.818815in}}%
\pgfpathquadraticcurveto{\pgfqpoint{6.691382in}{2.818815in}}{\pgfqpoint{6.686577in}{2.809440in}}%
\pgfpathquadraticcurveto{\pgfqpoint{6.681792in}{2.800084in}}{\pgfqpoint{6.681792in}{2.781276in}}%
\pgfpathquadraticcurveto{\pgfqpoint{6.681792in}{2.762545in}}{\pgfqpoint{6.686577in}{2.753170in}}%
\pgfpathquadraticcurveto{\pgfqpoint{6.691382in}{2.743795in}}{\pgfqpoint{6.700893in}{2.743795in}}%
\pgfpathquadraticcurveto{\pgfqpoint{6.710483in}{2.743795in}}{\pgfqpoint{6.715268in}{2.753170in}}%
\pgfpathquadraticcurveto{\pgfqpoint{6.720073in}{2.762545in}}{\pgfqpoint{6.720073in}{2.781276in}}%
\pgfpathquadraticcurveto{\pgfqpoint{6.720073in}{2.800084in}}{\pgfqpoint{6.715268in}{2.809440in}}%
\pgfpathquadraticcurveto{\pgfqpoint{6.710483in}{2.818815in}}{\pgfqpoint{6.700893in}{2.818815in}}%
\pgfpathlineto{\pgfqpoint{6.700893in}{2.818815in}}%
\pgfpathclose%
\pgfpathmoveto{\pgfqpoint{6.700893in}{2.828581in}}%
\pgfpathquadraticcurveto{\pgfqpoint{6.716225in}{2.828581in}}{\pgfqpoint{6.724311in}{2.816452in}}%
\pgfpathquadraticcurveto{\pgfqpoint{6.732397in}{2.804342in}}{\pgfqpoint{6.732397in}{2.781276in}}%
\pgfpathquadraticcurveto{\pgfqpoint{6.732397in}{2.758268in}}{\pgfqpoint{6.724311in}{2.746139in}}%
\pgfpathquadraticcurveto{\pgfqpoint{6.716225in}{2.734030in}}{\pgfqpoint{6.700893in}{2.734030in}}%
\pgfpathquadraticcurveto{\pgfqpoint{6.685581in}{2.734030in}}{\pgfqpoint{6.677495in}{2.746139in}}%
\pgfpathquadraticcurveto{\pgfqpoint{6.669409in}{2.758268in}}{\pgfqpoint{6.669409in}{2.781276in}}%
\pgfpathquadraticcurveto{\pgfqpoint{6.669409in}{2.804342in}}{\pgfqpoint{6.677495in}{2.816452in}}%
\pgfpathquadraticcurveto{\pgfqpoint{6.685581in}{2.828581in}}{\pgfqpoint{6.700893in}{2.828581in}}%
\pgfpathlineto{\pgfqpoint{6.700893in}{2.828581in}}%
\pgfpathclose%
\pgfusepath{fill}%
\end{pgfscope}%
\begin{pgfscope}%
\pgfsetbuttcap%
\pgfsetroundjoin%
\definecolor{currentfill}{rgb}{0.090196,0.168627,0.227451}%
\pgfsetfillcolor{currentfill}%
\pgfsetlinewidth{0.000000pt}%
\definecolor{currentstroke}{rgb}{0.090196,0.168627,0.227451}%
\pgfsetstrokecolor{currentstroke}%
\pgfsetdash{}{0pt}%
\pgfpathmoveto{\pgfqpoint{2.256669in}{2.572435in}}%
\pgfpathlineto{\pgfqpoint{2.256669in}{2.557730in}}%
\pgfpathquadraticcurveto{\pgfqpoint{2.248099in}{2.561836in}}{\pgfqpoint{2.240484in}{2.563841in}}%
\pgfpathquadraticcurveto{\pgfqpoint{2.232869in}{2.565870in}}{\pgfqpoint{2.225779in}{2.565870in}}%
\pgfpathquadraticcurveto{\pgfqpoint{2.213485in}{2.565870in}}{\pgfqpoint{2.206801in}{2.561096in}}%
\pgfpathquadraticcurveto{\pgfqpoint{2.200117in}{2.556321in}}{\pgfqpoint{2.200117in}{2.547513in}}%
\pgfpathquadraticcurveto{\pgfqpoint{2.200117in}{2.540113in}}{\pgfqpoint{2.204558in}{2.536341in}}%
\pgfpathquadraticcurveto{\pgfqpoint{2.208998in}{2.532593in}}{\pgfqpoint{2.221387in}{2.530278in}}%
\pgfpathlineto{\pgfqpoint{2.230482in}{2.528416in}}%
\pgfpathquadraticcurveto{\pgfqpoint{2.247335in}{2.525193in}}{\pgfqpoint{2.255356in}{2.517100in}}%
\pgfpathquadraticcurveto{\pgfqpoint{2.263377in}{2.509008in}}{\pgfqpoint{2.263377in}{2.495449in}}%
\pgfpathquadraticcurveto{\pgfqpoint{2.263377in}{2.479240in}}{\pgfqpoint{2.252515in}{2.470885in}}%
\pgfpathquadraticcurveto{\pgfqpoint{2.241678in}{2.462530in}}{\pgfqpoint{2.220719in}{2.462530in}}%
\pgfpathquadraticcurveto{\pgfqpoint{2.212817in}{2.462530in}}{\pgfqpoint{2.203889in}{2.464321in}}%
\pgfpathquadraticcurveto{\pgfqpoint{2.194985in}{2.466111in}}{\pgfqpoint{2.185436in}{2.469620in}}%
\pgfpathlineto{\pgfqpoint{2.185436in}{2.485137in}}%
\pgfpathquadraticcurveto{\pgfqpoint{2.194603in}{2.480004in}}{\pgfqpoint{2.203412in}{2.477378in}}%
\pgfpathquadraticcurveto{\pgfqpoint{2.212220in}{2.474776in}}{\pgfqpoint{2.220719in}{2.474776in}}%
\pgfpathquadraticcurveto{\pgfqpoint{2.233609in}{2.474776in}}{\pgfqpoint{2.240627in}{2.479837in}}%
\pgfpathquadraticcurveto{\pgfqpoint{2.247646in}{2.484922in}}{\pgfqpoint{2.247646in}{2.494327in}}%
\pgfpathquadraticcurveto{\pgfqpoint{2.247646in}{2.502515in}}{\pgfqpoint{2.242609in}{2.507146in}}%
\pgfpathquadraticcurveto{\pgfqpoint{2.237572in}{2.511777in}}{\pgfqpoint{2.226090in}{2.514093in}}%
\pgfpathlineto{\pgfqpoint{2.216899in}{2.515883in}}%
\pgfpathquadraticcurveto{\pgfqpoint{2.200046in}{2.519225in}}{\pgfqpoint{2.192502in}{2.526387in}}%
\pgfpathquadraticcurveto{\pgfqpoint{2.184983in}{2.533548in}}{\pgfqpoint{2.184983in}{2.546319in}}%
\pgfpathquadraticcurveto{\pgfqpoint{2.184983in}{2.561096in}}{\pgfqpoint{2.195391in}{2.569594in}}%
\pgfpathquadraticcurveto{\pgfqpoint{2.205799in}{2.578092in}}{\pgfqpoint{2.224061in}{2.578092in}}%
\pgfpathquadraticcurveto{\pgfqpoint{2.231914in}{2.578092in}}{\pgfqpoint{2.240031in}{2.576660in}}%
\pgfpathquadraticcurveto{\pgfqpoint{2.248171in}{2.575252in}}{\pgfqpoint{2.256669in}{2.572435in}}%
\pgfpathlineto{\pgfqpoint{2.256669in}{2.572435in}}%
\pgfpathclose%
\pgfpathmoveto{\pgfqpoint{2.324273in}{2.506692in}}%
\pgfpathquadraticcurveto{\pgfqpoint{2.307635in}{2.506692in}}{\pgfqpoint{2.301213in}{2.502897in}}%
\pgfpathquadraticcurveto{\pgfqpoint{2.294792in}{2.499101in}}{\pgfqpoint{2.294792in}{2.489911in}}%
\pgfpathquadraticcurveto{\pgfqpoint{2.294792in}{2.482606in}}{\pgfqpoint{2.299614in}{2.478309in}}%
\pgfpathquadraticcurveto{\pgfqpoint{2.304436in}{2.474036in}}{\pgfqpoint{2.312696in}{2.474036in}}%
\pgfpathquadraticcurveto{\pgfqpoint{2.324130in}{2.474036in}}{\pgfqpoint{2.331029in}{2.482129in}}%
\pgfpathquadraticcurveto{\pgfqpoint{2.337928in}{2.490221in}}{\pgfqpoint{2.337928in}{2.503637in}}%
\pgfpathlineto{\pgfqpoint{2.337928in}{2.506692in}}%
\pgfpathlineto{\pgfqpoint{2.324273in}{2.506692in}}%
\pgfpathlineto{\pgfqpoint{2.324273in}{2.506692in}}%
\pgfpathclose%
\pgfpathmoveto{\pgfqpoint{2.351654in}{2.512374in}}%
\pgfpathlineto{\pgfqpoint{2.351654in}{2.464702in}}%
\pgfpathlineto{\pgfqpoint{2.337928in}{2.464702in}}%
\pgfpathlineto{\pgfqpoint{2.337928in}{2.477378in}}%
\pgfpathquadraticcurveto{\pgfqpoint{2.333225in}{2.469787in}}{\pgfqpoint{2.326207in}{2.466159in}}%
\pgfpathquadraticcurveto{\pgfqpoint{2.319189in}{2.462530in}}{\pgfqpoint{2.309043in}{2.462530in}}%
\pgfpathquadraticcurveto{\pgfqpoint{2.296224in}{2.462530in}}{\pgfqpoint{2.288633in}{2.469739in}}%
\pgfpathquadraticcurveto{\pgfqpoint{2.281066in}{2.476949in}}{\pgfqpoint{2.281066in}{2.489028in}}%
\pgfpathquadraticcurveto{\pgfqpoint{2.281066in}{2.503112in}}{\pgfqpoint{2.290495in}{2.510273in}}%
\pgfpathquadraticcurveto{\pgfqpoint{2.299948in}{2.517435in}}{\pgfqpoint{2.318663in}{2.517435in}}%
\pgfpathlineto{\pgfqpoint{2.337928in}{2.517435in}}%
\pgfpathlineto{\pgfqpoint{2.337928in}{2.518795in}}%
\pgfpathquadraticcurveto{\pgfqpoint{2.337928in}{2.528272in}}{\pgfqpoint{2.331697in}{2.533452in}}%
\pgfpathquadraticcurveto{\pgfqpoint{2.325467in}{2.538633in}}{\pgfqpoint{2.314199in}{2.538633in}}%
\pgfpathquadraticcurveto{\pgfqpoint{2.307038in}{2.538633in}}{\pgfqpoint{2.300235in}{2.536914in}}%
\pgfpathquadraticcurveto{\pgfqpoint{2.293455in}{2.535195in}}{\pgfqpoint{2.287201in}{2.531758in}}%
\pgfpathlineto{\pgfqpoint{2.287201in}{2.544457in}}%
\pgfpathquadraticcurveto{\pgfqpoint{2.294720in}{2.547370in}}{\pgfqpoint{2.301810in}{2.548802in}}%
\pgfpathquadraticcurveto{\pgfqpoint{2.308900in}{2.550258in}}{\pgfqpoint{2.315608in}{2.550258in}}%
\pgfpathquadraticcurveto{\pgfqpoint{2.333750in}{2.550258in}}{\pgfqpoint{2.342702in}{2.540853in}}%
\pgfpathquadraticcurveto{\pgfqpoint{2.351654in}{2.531471in}}{\pgfqpoint{2.351654in}{2.512374in}}%
\pgfpathlineto{\pgfqpoint{2.351654in}{2.512374in}}%
\pgfpathclose%
\pgfpathmoveto{\pgfqpoint{2.444968in}{2.532211in}}%
\pgfpathquadraticcurveto{\pgfqpoint{2.450124in}{2.541473in}}{\pgfqpoint{2.457285in}{2.545866in}}%
\pgfpathquadraticcurveto{\pgfqpoint{2.464447in}{2.550258in}}{\pgfqpoint{2.474139in}{2.550258in}}%
\pgfpathquadraticcurveto{\pgfqpoint{2.487196in}{2.550258in}}{\pgfqpoint{2.494286in}{2.541115in}}%
\pgfpathquadraticcurveto{\pgfqpoint{2.501376in}{2.531996in}}{\pgfqpoint{2.501376in}{2.515143in}}%
\pgfpathlineto{\pgfqpoint{2.501376in}{2.464702in}}%
\pgfpathlineto{\pgfqpoint{2.487578in}{2.464702in}}%
\pgfpathlineto{\pgfqpoint{2.487578in}{2.514689in}}%
\pgfpathquadraticcurveto{\pgfqpoint{2.487578in}{2.526697in}}{\pgfqpoint{2.483305in}{2.532498in}}%
\pgfpathquadraticcurveto{\pgfqpoint{2.479056in}{2.538322in}}{\pgfqpoint{2.470343in}{2.538322in}}%
\pgfpathquadraticcurveto{\pgfqpoint{2.459673in}{2.538322in}}{\pgfqpoint{2.453466in}{2.531232in}}%
\pgfpathquadraticcurveto{\pgfqpoint{2.447283in}{2.524166in}}{\pgfqpoint{2.447283in}{2.511920in}}%
\pgfpathlineto{\pgfqpoint{2.447283in}{2.464702in}}%
\pgfpathlineto{\pgfqpoint{2.433485in}{2.464702in}}%
\pgfpathlineto{\pgfqpoint{2.433485in}{2.514689in}}%
\pgfpathquadraticcurveto{\pgfqpoint{2.433485in}{2.526768in}}{\pgfqpoint{2.429236in}{2.532545in}}%
\pgfpathquadraticcurveto{\pgfqpoint{2.424987in}{2.538322in}}{\pgfqpoint{2.416107in}{2.538322in}}%
\pgfpathquadraticcurveto{\pgfqpoint{2.405580in}{2.538322in}}{\pgfqpoint{2.399373in}{2.531209in}}%
\pgfpathquadraticcurveto{\pgfqpoint{2.393190in}{2.524095in}}{\pgfqpoint{2.393190in}{2.511920in}}%
\pgfpathlineto{\pgfqpoint{2.393190in}{2.464702in}}%
\pgfpathlineto{\pgfqpoint{2.379393in}{2.464702in}}%
\pgfpathlineto{\pgfqpoint{2.379393in}{2.548253in}}%
\pgfpathlineto{\pgfqpoint{2.393190in}{2.548253in}}%
\pgfpathlineto{\pgfqpoint{2.393190in}{2.535267in}}%
\pgfpathquadraticcurveto{\pgfqpoint{2.397893in}{2.542953in}}{\pgfqpoint{2.404458in}{2.546606in}}%
\pgfpathquadraticcurveto{\pgfqpoint{2.411022in}{2.550258in}}{\pgfqpoint{2.420046in}{2.550258in}}%
\pgfpathquadraticcurveto{\pgfqpoint{2.429165in}{2.550258in}}{\pgfqpoint{2.435538in}{2.545627in}}%
\pgfpathquadraticcurveto{\pgfqpoint{2.441912in}{2.541020in}}{\pgfqpoint{2.444968in}{2.532211in}}%
\pgfpathlineto{\pgfqpoint{2.444968in}{2.532211in}}%
\pgfpathclose%
\pgfpathmoveto{\pgfqpoint{2.542005in}{2.477235in}}%
\pgfpathlineto{\pgfqpoint{2.542005in}{2.432929in}}%
\pgfpathlineto{\pgfqpoint{2.528208in}{2.432929in}}%
\pgfpathlineto{\pgfqpoint{2.528208in}{2.548253in}}%
\pgfpathlineto{\pgfqpoint{2.542005in}{2.548253in}}%
\pgfpathlineto{\pgfqpoint{2.542005in}{2.535577in}}%
\pgfpathquadraticcurveto{\pgfqpoint{2.546350in}{2.543025in}}{\pgfqpoint{2.552939in}{2.546630in}}%
\pgfpathquadraticcurveto{\pgfqpoint{2.559551in}{2.550258in}}{\pgfqpoint{2.568718in}{2.550258in}}%
\pgfpathquadraticcurveto{\pgfqpoint{2.583948in}{2.550258in}}{\pgfqpoint{2.593449in}{2.538179in}}%
\pgfpathquadraticcurveto{\pgfqpoint{2.602973in}{2.526100in}}{\pgfqpoint{2.602973in}{2.506406in}}%
\pgfpathquadraticcurveto{\pgfqpoint{2.602973in}{2.486712in}}{\pgfqpoint{2.593449in}{2.474609in}}%
\pgfpathquadraticcurveto{\pgfqpoint{2.583948in}{2.462530in}}{\pgfqpoint{2.568718in}{2.462530in}}%
\pgfpathquadraticcurveto{\pgfqpoint{2.559551in}{2.462530in}}{\pgfqpoint{2.552939in}{2.466159in}}%
\pgfpathquadraticcurveto{\pgfqpoint{2.546350in}{2.469787in}}{\pgfqpoint{2.542005in}{2.477235in}}%
\pgfpathlineto{\pgfqpoint{2.542005in}{2.477235in}}%
\pgfpathclose%
\pgfpathmoveto{\pgfqpoint{2.588722in}{2.506406in}}%
\pgfpathquadraticcurveto{\pgfqpoint{2.588722in}{2.521541in}}{\pgfqpoint{2.582492in}{2.530158in}}%
\pgfpathquadraticcurveto{\pgfqpoint{2.576261in}{2.538776in}}{\pgfqpoint{2.565376in}{2.538776in}}%
\pgfpathquadraticcurveto{\pgfqpoint{2.554466in}{2.538776in}}{\pgfqpoint{2.548236in}{2.530158in}}%
\pgfpathquadraticcurveto{\pgfqpoint{2.542005in}{2.521541in}}{\pgfqpoint{2.542005in}{2.506406in}}%
\pgfpathquadraticcurveto{\pgfqpoint{2.542005in}{2.491271in}}{\pgfqpoint{2.548236in}{2.482654in}}%
\pgfpathquadraticcurveto{\pgfqpoint{2.554466in}{2.474036in}}{\pgfqpoint{2.565376in}{2.474036in}}%
\pgfpathquadraticcurveto{\pgfqpoint{2.576261in}{2.474036in}}{\pgfqpoint{2.582492in}{2.482654in}}%
\pgfpathquadraticcurveto{\pgfqpoint{2.588722in}{2.491271in}}{\pgfqpoint{2.588722in}{2.506406in}}%
\pgfpathlineto{\pgfqpoint{2.588722in}{2.506406in}}%
\pgfpathclose%
\pgfpathmoveto{\pgfqpoint{2.625723in}{2.580790in}}%
\pgfpathlineto{\pgfqpoint{2.639449in}{2.580790in}}%
\pgfpathlineto{\pgfqpoint{2.639449in}{2.464702in}}%
\pgfpathlineto{\pgfqpoint{2.625723in}{2.464702in}}%
\pgfpathlineto{\pgfqpoint{2.625723in}{2.580790in}}%
\pgfpathlineto{\pgfqpoint{2.625723in}{2.580790in}}%
\pgfpathclose%
\pgfpathmoveto{\pgfqpoint{2.739638in}{2.509915in}}%
\pgfpathlineto{\pgfqpoint{2.739638in}{2.503207in}}%
\pgfpathlineto{\pgfqpoint{2.676522in}{2.503207in}}%
\pgfpathquadraticcurveto{\pgfqpoint{2.677429in}{2.489028in}}{\pgfqpoint{2.685068in}{2.481604in}}%
\pgfpathquadraticcurveto{\pgfqpoint{2.692706in}{2.474179in}}{\pgfqpoint{2.706361in}{2.474179in}}%
\pgfpathquadraticcurveto{\pgfqpoint{2.714262in}{2.474179in}}{\pgfqpoint{2.721686in}{2.476113in}}%
\pgfpathquadraticcurveto{\pgfqpoint{2.729110in}{2.478047in}}{\pgfqpoint{2.736439in}{2.481938in}}%
\pgfpathlineto{\pgfqpoint{2.736439in}{2.468952in}}%
\pgfpathquadraticcurveto{\pgfqpoint{2.729039in}{2.465824in}}{\pgfqpoint{2.721281in}{2.464177in}}%
\pgfpathquadraticcurveto{\pgfqpoint{2.713522in}{2.462530in}}{\pgfqpoint{2.705549in}{2.462530in}}%
\pgfpathquadraticcurveto{\pgfqpoint{2.685545in}{2.462530in}}{\pgfqpoint{2.673872in}{2.474156in}}%
\pgfpathquadraticcurveto{\pgfqpoint{2.662199in}{2.485805in}}{\pgfqpoint{2.662199in}{2.505666in}}%
\pgfpathquadraticcurveto{\pgfqpoint{2.662199in}{2.526172in}}{\pgfqpoint{2.673275in}{2.538203in}}%
\pgfpathquadraticcurveto{\pgfqpoint{2.684351in}{2.550258in}}{\pgfqpoint{2.703162in}{2.550258in}}%
\pgfpathquadraticcurveto{\pgfqpoint{2.720015in}{2.550258in}}{\pgfqpoint{2.729827in}{2.539396in}}%
\pgfpathquadraticcurveto{\pgfqpoint{2.739638in}{2.528559in}}{\pgfqpoint{2.739638in}{2.509915in}}%
\pgfpathlineto{\pgfqpoint{2.739638in}{2.509915in}}%
\pgfpathclose%
\pgfpathmoveto{\pgfqpoint{2.725912in}{2.513949in}}%
\pgfpathquadraticcurveto{\pgfqpoint{2.725768in}{2.525193in}}{\pgfqpoint{2.719610in}{2.531901in}}%
\pgfpathquadraticcurveto{\pgfqpoint{2.713451in}{2.538633in}}{\pgfqpoint{2.703305in}{2.538633in}}%
\pgfpathquadraticcurveto{\pgfqpoint{2.691823in}{2.538633in}}{\pgfqpoint{2.684924in}{2.532140in}}%
\pgfpathquadraticcurveto{\pgfqpoint{2.678025in}{2.525646in}}{\pgfqpoint{2.676975in}{2.513854in}}%
\pgfpathlineto{\pgfqpoint{2.725912in}{2.513949in}}%
\pgfpathlineto{\pgfqpoint{2.725912in}{2.513949in}}%
\pgfpathclose%
\pgfpathmoveto{\pgfqpoint{2.755250in}{2.512660in}}%
\pgfpathlineto{\pgfqpoint{2.795449in}{2.512660in}}%
\pgfpathlineto{\pgfqpoint{2.795449in}{2.500438in}}%
\pgfpathlineto{\pgfqpoint{2.755250in}{2.500438in}}%
\pgfpathlineto{\pgfqpoint{2.755250in}{2.512660in}}%
\pgfpathlineto{\pgfqpoint{2.755250in}{2.512660in}}%
\pgfpathclose%
\pgfpathmoveto{\pgfqpoint{2.855272in}{2.506692in}}%
\pgfpathquadraticcurveto{\pgfqpoint{2.838633in}{2.506692in}}{\pgfqpoint{2.832212in}{2.502897in}}%
\pgfpathquadraticcurveto{\pgfqpoint{2.825790in}{2.499101in}}{\pgfqpoint{2.825790in}{2.489911in}}%
\pgfpathquadraticcurveto{\pgfqpoint{2.825790in}{2.482606in}}{\pgfqpoint{2.830612in}{2.478309in}}%
\pgfpathquadraticcurveto{\pgfqpoint{2.835434in}{2.474036in}}{\pgfqpoint{2.843694in}{2.474036in}}%
\pgfpathquadraticcurveto{\pgfqpoint{2.855128in}{2.474036in}}{\pgfqpoint{2.862027in}{2.482129in}}%
\pgfpathquadraticcurveto{\pgfqpoint{2.868926in}{2.490221in}}{\pgfqpoint{2.868926in}{2.503637in}}%
\pgfpathlineto{\pgfqpoint{2.868926in}{2.506692in}}%
\pgfpathlineto{\pgfqpoint{2.855272in}{2.506692in}}%
\pgfpathlineto{\pgfqpoint{2.855272in}{2.506692in}}%
\pgfpathclose%
\pgfpathmoveto{\pgfqpoint{2.882652in}{2.512374in}}%
\pgfpathlineto{\pgfqpoint{2.882652in}{2.464702in}}%
\pgfpathlineto{\pgfqpoint{2.868926in}{2.464702in}}%
\pgfpathlineto{\pgfqpoint{2.868926in}{2.477378in}}%
\pgfpathquadraticcurveto{\pgfqpoint{2.864223in}{2.469787in}}{\pgfqpoint{2.857205in}{2.466159in}}%
\pgfpathquadraticcurveto{\pgfqpoint{2.850187in}{2.462530in}}{\pgfqpoint{2.840041in}{2.462530in}}%
\pgfpathquadraticcurveto{\pgfqpoint{2.827222in}{2.462530in}}{\pgfqpoint{2.819631in}{2.469739in}}%
\pgfpathquadraticcurveto{\pgfqpoint{2.812064in}{2.476949in}}{\pgfqpoint{2.812064in}{2.489028in}}%
\pgfpathquadraticcurveto{\pgfqpoint{2.812064in}{2.503112in}}{\pgfqpoint{2.821493in}{2.510273in}}%
\pgfpathquadraticcurveto{\pgfqpoint{2.830946in}{2.517435in}}{\pgfqpoint{2.849662in}{2.517435in}}%
\pgfpathlineto{\pgfqpoint{2.868926in}{2.517435in}}%
\pgfpathlineto{\pgfqpoint{2.868926in}{2.518795in}}%
\pgfpathquadraticcurveto{\pgfqpoint{2.868926in}{2.528272in}}{\pgfqpoint{2.862696in}{2.533452in}}%
\pgfpathquadraticcurveto{\pgfqpoint{2.856465in}{2.538633in}}{\pgfqpoint{2.845198in}{2.538633in}}%
\pgfpathquadraticcurveto{\pgfqpoint{2.838036in}{2.538633in}}{\pgfqpoint{2.831233in}{2.536914in}}%
\pgfpathquadraticcurveto{\pgfqpoint{2.824453in}{2.535195in}}{\pgfqpoint{2.818199in}{2.531758in}}%
\pgfpathlineto{\pgfqpoint{2.818199in}{2.544457in}}%
\pgfpathquadraticcurveto{\pgfqpoint{2.825719in}{2.547370in}}{\pgfqpoint{2.832808in}{2.548802in}}%
\pgfpathquadraticcurveto{\pgfqpoint{2.839898in}{2.550258in}}{\pgfqpoint{2.846606in}{2.550258in}}%
\pgfpathquadraticcurveto{\pgfqpoint{2.864749in}{2.550258in}}{\pgfqpoint{2.873700in}{2.540853in}}%
\pgfpathquadraticcurveto{\pgfqpoint{2.882652in}{2.531471in}}{\pgfqpoint{2.882652in}{2.512374in}}%
\pgfpathlineto{\pgfqpoint{2.882652in}{2.512374in}}%
\pgfpathclose%
\pgfpathmoveto{\pgfqpoint{2.910916in}{2.580790in}}%
\pgfpathlineto{\pgfqpoint{2.924642in}{2.580790in}}%
\pgfpathlineto{\pgfqpoint{2.924642in}{2.464702in}}%
\pgfpathlineto{\pgfqpoint{2.910916in}{2.464702in}}%
\pgfpathlineto{\pgfqpoint{2.910916in}{2.580790in}}%
\pgfpathlineto{\pgfqpoint{2.910916in}{2.580790in}}%
\pgfpathclose%
\pgfpathmoveto{\pgfqpoint{2.953360in}{2.548253in}}%
\pgfpathlineto{\pgfqpoint{2.967086in}{2.548253in}}%
\pgfpathlineto{\pgfqpoint{2.967086in}{2.464702in}}%
\pgfpathlineto{\pgfqpoint{2.953360in}{2.464702in}}%
\pgfpathlineto{\pgfqpoint{2.953360in}{2.548253in}}%
\pgfpathlineto{\pgfqpoint{2.953360in}{2.548253in}}%
\pgfpathclose%
\pgfpathmoveto{\pgfqpoint{2.953360in}{2.580790in}}%
\pgfpathlineto{\pgfqpoint{2.967086in}{2.580790in}}%
\pgfpathlineto{\pgfqpoint{2.967086in}{2.563387in}}%
\pgfpathlineto{\pgfqpoint{2.953360in}{2.563387in}}%
\pgfpathlineto{\pgfqpoint{2.953360in}{2.580790in}}%
\pgfpathlineto{\pgfqpoint{2.953360in}{2.580790in}}%
\pgfpathclose%
\pgfpathmoveto{\pgfqpoint{3.050779in}{2.507456in}}%
\pgfpathquadraticcurveto{\pgfqpoint{3.050779in}{2.522376in}}{\pgfqpoint{3.044620in}{2.530564in}}%
\pgfpathquadraticcurveto{\pgfqpoint{3.038485in}{2.538776in}}{\pgfqpoint{3.027361in}{2.538776in}}%
\pgfpathquadraticcurveto{\pgfqpoint{3.016333in}{2.538776in}}{\pgfqpoint{3.010174in}{2.530564in}}%
\pgfpathquadraticcurveto{\pgfqpoint{3.004015in}{2.522376in}}{\pgfqpoint{3.004015in}{2.507456in}}%
\pgfpathquadraticcurveto{\pgfqpoint{3.004015in}{2.492608in}}{\pgfqpoint{3.010174in}{2.484396in}}%
\pgfpathquadraticcurveto{\pgfqpoint{3.016333in}{2.476185in}}{\pgfqpoint{3.027361in}{2.476185in}}%
\pgfpathquadraticcurveto{\pgfqpoint{3.038485in}{2.476185in}}{\pgfqpoint{3.044620in}{2.484396in}}%
\pgfpathquadraticcurveto{\pgfqpoint{3.050779in}{2.492608in}}{\pgfqpoint{3.050779in}{2.507456in}}%
\pgfpathlineto{\pgfqpoint{3.050779in}{2.507456in}}%
\pgfpathclose%
\pgfpathmoveto{\pgfqpoint{3.064505in}{2.475063in}}%
\pgfpathquadraticcurveto{\pgfqpoint{3.064505in}{2.453745in}}{\pgfqpoint{3.055028in}{2.443337in}}%
\pgfpathquadraticcurveto{\pgfqpoint{3.045575in}{2.432929in}}{\pgfqpoint{3.026025in}{2.432929in}}%
\pgfpathquadraticcurveto{\pgfqpoint{3.018791in}{2.432929in}}{\pgfqpoint{3.012370in}{2.434004in}}%
\pgfpathquadraticcurveto{\pgfqpoint{3.005949in}{2.435078in}}{\pgfqpoint{2.999909in}{2.437322in}}%
\pgfpathlineto{\pgfqpoint{2.999909in}{2.450666in}}%
\pgfpathquadraticcurveto{\pgfqpoint{3.005949in}{2.447396in}}{\pgfqpoint{3.011845in}{2.445844in}}%
\pgfpathquadraticcurveto{\pgfqpoint{3.017741in}{2.444268in}}{\pgfqpoint{3.023852in}{2.444268in}}%
\pgfpathquadraticcurveto{\pgfqpoint{3.037364in}{2.444268in}}{\pgfqpoint{3.044071in}{2.451311in}}%
\pgfpathquadraticcurveto{\pgfqpoint{3.050779in}{2.458353in}}{\pgfqpoint{3.050779in}{2.472604in}}%
\pgfpathlineto{\pgfqpoint{3.050779in}{2.479407in}}%
\pgfpathquadraticcurveto{\pgfqpoint{3.046530in}{2.472007in}}{\pgfqpoint{3.039894in}{2.468355in}}%
\pgfpathquadraticcurveto{\pgfqpoint{3.033258in}{2.464702in}}{\pgfqpoint{3.023995in}{2.464702in}}%
\pgfpathquadraticcurveto{\pgfqpoint{3.008646in}{2.464702in}}{\pgfqpoint{2.999241in}{2.476400in}}%
\pgfpathquadraticcurveto{\pgfqpoint{2.989835in}{2.488120in}}{\pgfqpoint{2.989835in}{2.507456in}}%
\pgfpathquadraticcurveto{\pgfqpoint{2.989835in}{2.526840in}}{\pgfqpoint{2.999241in}{2.538537in}}%
\pgfpathquadraticcurveto{\pgfqpoint{3.008646in}{2.550258in}}{\pgfqpoint{3.023995in}{2.550258in}}%
\pgfpathquadraticcurveto{\pgfqpoint{3.033258in}{2.550258in}}{\pgfqpoint{3.039894in}{2.546606in}}%
\pgfpathquadraticcurveto{\pgfqpoint{3.046530in}{2.542953in}}{\pgfqpoint{3.050779in}{2.535577in}}%
\pgfpathlineto{\pgfqpoint{3.050779in}{2.548253in}}%
\pgfpathlineto{\pgfqpoint{3.064505in}{2.548253in}}%
\pgfpathlineto{\pgfqpoint{3.064505in}{2.475063in}}%
\pgfpathlineto{\pgfqpoint{3.064505in}{2.475063in}}%
\pgfpathclose%
\pgfpathmoveto{\pgfqpoint{3.162259in}{2.515143in}}%
\pgfpathlineto{\pgfqpoint{3.162259in}{2.464702in}}%
\pgfpathlineto{\pgfqpoint{3.148533in}{2.464702in}}%
\pgfpathlineto{\pgfqpoint{3.148533in}{2.514689in}}%
\pgfpathquadraticcurveto{\pgfqpoint{3.148533in}{2.526554in}}{\pgfqpoint{3.143902in}{2.532426in}}%
\pgfpathquadraticcurveto{\pgfqpoint{3.139271in}{2.538322in}}{\pgfqpoint{3.130033in}{2.538322in}}%
\pgfpathquadraticcurveto{\pgfqpoint{3.118909in}{2.538322in}}{\pgfqpoint{3.112487in}{2.531232in}}%
\pgfpathquadraticcurveto{\pgfqpoint{3.106066in}{2.524166in}}{\pgfqpoint{3.106066in}{2.511920in}}%
\pgfpathlineto{\pgfqpoint{3.106066in}{2.464702in}}%
\pgfpathlineto{\pgfqpoint{3.092268in}{2.464702in}}%
\pgfpathlineto{\pgfqpoint{3.092268in}{2.548253in}}%
\pgfpathlineto{\pgfqpoint{3.106066in}{2.548253in}}%
\pgfpathlineto{\pgfqpoint{3.106066in}{2.535267in}}%
\pgfpathquadraticcurveto{\pgfqpoint{3.111007in}{2.542810in}}{\pgfqpoint{3.117667in}{2.546534in}}%
\pgfpathquadraticcurveto{\pgfqpoint{3.124351in}{2.550258in}}{\pgfqpoint{3.133088in}{2.550258in}}%
\pgfpathquadraticcurveto{\pgfqpoint{3.147483in}{2.550258in}}{\pgfqpoint{3.154859in}{2.541354in}}%
\pgfpathquadraticcurveto{\pgfqpoint{3.162259in}{2.532450in}}{\pgfqpoint{3.162259in}{2.515143in}}%
\pgfpathlineto{\pgfqpoint{3.162259in}{2.515143in}}%
\pgfpathclose%
\pgfpathmoveto{\pgfqpoint{3.261087in}{2.509915in}}%
\pgfpathlineto{\pgfqpoint{3.261087in}{2.503207in}}%
\pgfpathlineto{\pgfqpoint{3.197971in}{2.503207in}}%
\pgfpathquadraticcurveto{\pgfqpoint{3.198878in}{2.489028in}}{\pgfqpoint{3.206517in}{2.481604in}}%
\pgfpathquadraticcurveto{\pgfqpoint{3.214156in}{2.474179in}}{\pgfqpoint{3.227811in}{2.474179in}}%
\pgfpathquadraticcurveto{\pgfqpoint{3.235712in}{2.474179in}}{\pgfqpoint{3.243136in}{2.476113in}}%
\pgfpathquadraticcurveto{\pgfqpoint{3.250560in}{2.478047in}}{\pgfqpoint{3.257889in}{2.481938in}}%
\pgfpathlineto{\pgfqpoint{3.257889in}{2.468952in}}%
\pgfpathquadraticcurveto{\pgfqpoint{3.250489in}{2.465824in}}{\pgfqpoint{3.242730in}{2.464177in}}%
\pgfpathquadraticcurveto{\pgfqpoint{3.234972in}{2.462530in}}{\pgfqpoint{3.226999in}{2.462530in}}%
\pgfpathquadraticcurveto{\pgfqpoint{3.206995in}{2.462530in}}{\pgfqpoint{3.195321in}{2.474156in}}%
\pgfpathquadraticcurveto{\pgfqpoint{3.183648in}{2.485805in}}{\pgfqpoint{3.183648in}{2.505666in}}%
\pgfpathquadraticcurveto{\pgfqpoint{3.183648in}{2.526172in}}{\pgfqpoint{3.194725in}{2.538203in}}%
\pgfpathquadraticcurveto{\pgfqpoint{3.205801in}{2.550258in}}{\pgfqpoint{3.224612in}{2.550258in}}%
\pgfpathquadraticcurveto{\pgfqpoint{3.241465in}{2.550258in}}{\pgfqpoint{3.251276in}{2.539396in}}%
\pgfpathquadraticcurveto{\pgfqpoint{3.261087in}{2.528559in}}{\pgfqpoint{3.261087in}{2.509915in}}%
\pgfpathlineto{\pgfqpoint{3.261087in}{2.509915in}}%
\pgfpathclose%
\pgfpathmoveto{\pgfqpoint{3.247361in}{2.513949in}}%
\pgfpathquadraticcurveto{\pgfqpoint{3.247218in}{2.525193in}}{\pgfqpoint{3.241059in}{2.531901in}}%
\pgfpathquadraticcurveto{\pgfqpoint{3.234900in}{2.538633in}}{\pgfqpoint{3.224755in}{2.538633in}}%
\pgfpathquadraticcurveto{\pgfqpoint{3.213273in}{2.538633in}}{\pgfqpoint{3.206374in}{2.532140in}}%
\pgfpathquadraticcurveto{\pgfqpoint{3.199475in}{2.525646in}}{\pgfqpoint{3.198425in}{2.513854in}}%
\pgfpathlineto{\pgfqpoint{3.247361in}{2.513949in}}%
\pgfpathlineto{\pgfqpoint{3.247361in}{2.513949in}}%
\pgfpathclose%
\pgfpathmoveto{\pgfqpoint{3.338598in}{2.535577in}}%
\pgfpathlineto{\pgfqpoint{3.338598in}{2.580790in}}%
\pgfpathlineto{\pgfqpoint{3.352324in}{2.580790in}}%
\pgfpathlineto{\pgfqpoint{3.352324in}{2.464702in}}%
\pgfpathlineto{\pgfqpoint{3.338598in}{2.464702in}}%
\pgfpathlineto{\pgfqpoint{3.338598in}{2.477235in}}%
\pgfpathquadraticcurveto{\pgfqpoint{3.334278in}{2.469787in}}{\pgfqpoint{3.327665in}{2.466159in}}%
\pgfpathquadraticcurveto{\pgfqpoint{3.321077in}{2.462530in}}{\pgfqpoint{3.311814in}{2.462530in}}%
\pgfpathquadraticcurveto{\pgfqpoint{3.296680in}{2.462530in}}{\pgfqpoint{3.287155in}{2.474609in}}%
\pgfpathquadraticcurveto{\pgfqpoint{3.277654in}{2.486712in}}{\pgfqpoint{3.277654in}{2.506406in}}%
\pgfpathquadraticcurveto{\pgfqpoint{3.277654in}{2.526100in}}{\pgfqpoint{3.287155in}{2.538179in}}%
\pgfpathquadraticcurveto{\pgfqpoint{3.296680in}{2.550258in}}{\pgfqpoint{3.311814in}{2.550258in}}%
\pgfpathquadraticcurveto{\pgfqpoint{3.321077in}{2.550258in}}{\pgfqpoint{3.327665in}{2.546630in}}%
\pgfpathquadraticcurveto{\pgfqpoint{3.334278in}{2.543025in}}{\pgfqpoint{3.338598in}{2.535577in}}%
\pgfpathlineto{\pgfqpoint{3.338598in}{2.535577in}}%
\pgfpathclose%
\pgfpathmoveto{\pgfqpoint{3.291834in}{2.506406in}}%
\pgfpathquadraticcurveto{\pgfqpoint{3.291834in}{2.491271in}}{\pgfqpoint{3.298064in}{2.482654in}}%
\pgfpathquadraticcurveto{\pgfqpoint{3.304295in}{2.474036in}}{\pgfqpoint{3.315180in}{2.474036in}}%
\pgfpathquadraticcurveto{\pgfqpoint{3.326066in}{2.474036in}}{\pgfqpoint{3.332320in}{2.482654in}}%
\pgfpathquadraticcurveto{\pgfqpoint{3.338598in}{2.491271in}}{\pgfqpoint{3.338598in}{2.506406in}}%
\pgfpathquadraticcurveto{\pgfqpoint{3.338598in}{2.521541in}}{\pgfqpoint{3.332320in}{2.530158in}}%
\pgfpathquadraticcurveto{\pgfqpoint{3.326066in}{2.538776in}}{\pgfqpoint{3.315180in}{2.538776in}}%
\pgfpathquadraticcurveto{\pgfqpoint{3.304295in}{2.538776in}}{\pgfqpoint{3.298064in}{2.530158in}}%
\pgfpathquadraticcurveto{\pgfqpoint{3.291834in}{2.521541in}}{\pgfqpoint{3.291834in}{2.506406in}}%
\pgfpathlineto{\pgfqpoint{3.291834in}{2.506406in}}%
\pgfpathclose%
\pgfpathmoveto{\pgfqpoint{3.414319in}{2.576087in}}%
\pgfpathlineto{\pgfqpoint{3.508540in}{2.576087in}}%
\pgfpathlineto{\pgfqpoint{3.508540in}{2.563387in}}%
\pgfpathlineto{\pgfqpoint{3.469008in}{2.563387in}}%
\pgfpathlineto{\pgfqpoint{3.469008in}{2.464702in}}%
\pgfpathlineto{\pgfqpoint{3.453874in}{2.464702in}}%
\pgfpathlineto{\pgfqpoint{3.453874in}{2.563387in}}%
\pgfpathlineto{\pgfqpoint{3.414319in}{2.563387in}}%
\pgfpathlineto{\pgfqpoint{3.414319in}{2.576087in}}%
\pgfpathlineto{\pgfqpoint{3.414319in}{2.576087in}}%
\pgfpathclose%
\pgfpathmoveto{\pgfqpoint{3.513171in}{2.576087in}}%
\pgfpathlineto{\pgfqpoint{3.528377in}{2.576087in}}%
\pgfpathlineto{\pgfqpoint{3.551795in}{2.481938in}}%
\pgfpathlineto{\pgfqpoint{3.575141in}{2.576087in}}%
\pgfpathlineto{\pgfqpoint{3.592090in}{2.576087in}}%
\pgfpathlineto{\pgfqpoint{3.615508in}{2.481938in}}%
\pgfpathlineto{\pgfqpoint{3.638854in}{2.576087in}}%
\pgfpathlineto{\pgfqpoint{3.654156in}{2.576087in}}%
\pgfpathlineto{\pgfqpoint{3.626179in}{2.464702in}}%
\pgfpathlineto{\pgfqpoint{3.607225in}{2.464702in}}%
\pgfpathlineto{\pgfqpoint{3.583735in}{2.561382in}}%
\pgfpathlineto{\pgfqpoint{3.560007in}{2.464702in}}%
\pgfpathlineto{\pgfqpoint{3.541053in}{2.464702in}}%
\pgfpathlineto{\pgfqpoint{3.513171in}{2.576087in}}%
\pgfpathlineto{\pgfqpoint{3.513171in}{2.576087in}}%
\pgfpathclose%
\pgfpathmoveto{\pgfqpoint{3.674137in}{2.576087in}}%
\pgfpathlineto{\pgfqpoint{3.738136in}{2.576087in}}%
\pgfpathlineto{\pgfqpoint{3.738136in}{2.563387in}}%
\pgfpathlineto{\pgfqpoint{3.689199in}{2.563387in}}%
\pgfpathlineto{\pgfqpoint{3.689199in}{2.530564in}}%
\pgfpathlineto{\pgfqpoint{3.733362in}{2.530564in}}%
\pgfpathlineto{\pgfqpoint{3.733362in}{2.517888in}}%
\pgfpathlineto{\pgfqpoint{3.689199in}{2.517888in}}%
\pgfpathlineto{\pgfqpoint{3.689199in}{2.464702in}}%
\pgfpathlineto{\pgfqpoint{3.674137in}{2.464702in}}%
\pgfpathlineto{\pgfqpoint{3.674137in}{2.576087in}}%
\pgfpathlineto{\pgfqpoint{3.674137in}{2.576087in}}%
\pgfpathclose%
\pgfpathmoveto{\pgfqpoint{3.762008in}{2.576087in}}%
\pgfpathlineto{\pgfqpoint{3.832429in}{2.576087in}}%
\pgfpathlineto{\pgfqpoint{3.832429in}{2.563387in}}%
\pgfpathlineto{\pgfqpoint{3.777071in}{2.563387in}}%
\pgfpathlineto{\pgfqpoint{3.777071in}{2.530421in}}%
\pgfpathlineto{\pgfqpoint{3.830113in}{2.530421in}}%
\pgfpathlineto{\pgfqpoint{3.830113in}{2.517745in}}%
\pgfpathlineto{\pgfqpoint{3.777071in}{2.517745in}}%
\pgfpathlineto{\pgfqpoint{3.777071in}{2.477378in}}%
\pgfpathlineto{\pgfqpoint{3.833765in}{2.477378in}}%
\pgfpathlineto{\pgfqpoint{3.833765in}{2.464702in}}%
\pgfpathlineto{\pgfqpoint{3.762008in}{2.464702in}}%
\pgfpathlineto{\pgfqpoint{3.762008in}{2.576087in}}%
\pgfpathlineto{\pgfqpoint{3.762008in}{2.576087in}}%
\pgfpathclose%
\pgfpathmoveto{\pgfqpoint{3.977973in}{2.509915in}}%
\pgfpathlineto{\pgfqpoint{3.977973in}{2.503207in}}%
\pgfpathlineto{\pgfqpoint{3.914857in}{2.503207in}}%
\pgfpathquadraticcurveto{\pgfqpoint{3.915764in}{2.489028in}}{\pgfqpoint{3.923403in}{2.481604in}}%
\pgfpathquadraticcurveto{\pgfqpoint{3.931042in}{2.474179in}}{\pgfqpoint{3.944696in}{2.474179in}}%
\pgfpathquadraticcurveto{\pgfqpoint{3.952598in}{2.474179in}}{\pgfqpoint{3.960022in}{2.476113in}}%
\pgfpathquadraticcurveto{\pgfqpoint{3.967446in}{2.478047in}}{\pgfqpoint{3.974775in}{2.481938in}}%
\pgfpathlineto{\pgfqpoint{3.974775in}{2.468952in}}%
\pgfpathquadraticcurveto{\pgfqpoint{3.967374in}{2.465824in}}{\pgfqpoint{3.959616in}{2.464177in}}%
\pgfpathquadraticcurveto{\pgfqpoint{3.951858in}{2.462530in}}{\pgfqpoint{3.943885in}{2.462530in}}%
\pgfpathquadraticcurveto{\pgfqpoint{3.923880in}{2.462530in}}{\pgfqpoint{3.912207in}{2.474156in}}%
\pgfpathquadraticcurveto{\pgfqpoint{3.900534in}{2.485805in}}{\pgfqpoint{3.900534in}{2.505666in}}%
\pgfpathquadraticcurveto{\pgfqpoint{3.900534in}{2.526172in}}{\pgfqpoint{3.911610in}{2.538203in}}%
\pgfpathquadraticcurveto{\pgfqpoint{3.922687in}{2.550258in}}{\pgfqpoint{3.941498in}{2.550258in}}%
\pgfpathquadraticcurveto{\pgfqpoint{3.958351in}{2.550258in}}{\pgfqpoint{3.968162in}{2.539396in}}%
\pgfpathquadraticcurveto{\pgfqpoint{3.977973in}{2.528559in}}{\pgfqpoint{3.977973in}{2.509915in}}%
\pgfpathlineto{\pgfqpoint{3.977973in}{2.509915in}}%
\pgfpathclose%
\pgfpathmoveto{\pgfqpoint{3.964247in}{2.513949in}}%
\pgfpathquadraticcurveto{\pgfqpoint{3.964104in}{2.525193in}}{\pgfqpoint{3.957945in}{2.531901in}}%
\pgfpathquadraticcurveto{\pgfqpoint{3.951786in}{2.538633in}}{\pgfqpoint{3.941641in}{2.538633in}}%
\pgfpathquadraticcurveto{\pgfqpoint{3.930159in}{2.538633in}}{\pgfqpoint{3.923260in}{2.532140in}}%
\pgfpathquadraticcurveto{\pgfqpoint{3.916361in}{2.525646in}}{\pgfqpoint{3.915311in}{2.513854in}}%
\pgfpathlineto{\pgfqpoint{3.964247in}{2.513949in}}%
\pgfpathlineto{\pgfqpoint{3.964247in}{2.513949in}}%
\pgfpathclose%
\pgfpathmoveto{\pgfqpoint{4.053765in}{2.545794in}}%
\pgfpathlineto{\pgfqpoint{4.053765in}{2.532808in}}%
\pgfpathquadraticcurveto{\pgfqpoint{4.047965in}{2.535792in}}{\pgfqpoint{4.041686in}{2.537272in}}%
\pgfpathquadraticcurveto{\pgfqpoint{4.035432in}{2.538776in}}{\pgfqpoint{4.028700in}{2.538776in}}%
\pgfpathquadraticcurveto{\pgfqpoint{4.018483in}{2.538776in}}{\pgfqpoint{4.013375in}{2.535649in}}%
\pgfpathquadraticcurveto{\pgfqpoint{4.008266in}{2.532521in}}{\pgfqpoint{4.008266in}{2.526243in}}%
\pgfpathquadraticcurveto{\pgfqpoint{4.008266in}{2.521469in}}{\pgfqpoint{4.011919in}{2.518748in}}%
\pgfpathquadraticcurveto{\pgfqpoint{4.015571in}{2.516026in}}{\pgfqpoint{4.026624in}{2.513567in}}%
\pgfpathlineto{\pgfqpoint{4.031326in}{2.512517in}}%
\pgfpathquadraticcurveto{\pgfqpoint{4.045936in}{2.509390in}}{\pgfqpoint{4.052094in}{2.503685in}}%
\pgfpathquadraticcurveto{\pgfqpoint{4.058253in}{2.497979in}}{\pgfqpoint{4.058253in}{2.487762in}}%
\pgfpathquadraticcurveto{\pgfqpoint{4.058253in}{2.476113in}}{\pgfqpoint{4.049039in}{2.469310in}}%
\pgfpathquadraticcurveto{\pgfqpoint{4.039824in}{2.462530in}}{\pgfqpoint{4.023711in}{2.462530in}}%
\pgfpathquadraticcurveto{\pgfqpoint{4.017003in}{2.462530in}}{\pgfqpoint{4.009722in}{2.463843in}}%
\pgfpathquadraticcurveto{\pgfqpoint{4.002442in}{2.465156in}}{\pgfqpoint{3.994397in}{2.467758in}}%
\pgfpathlineto{\pgfqpoint{3.994397in}{2.481938in}}%
\pgfpathquadraticcurveto{\pgfqpoint{4.002012in}{2.477975in}}{\pgfqpoint{4.009388in}{2.475994in}}%
\pgfpathquadraticcurveto{\pgfqpoint{4.016765in}{2.474036in}}{\pgfqpoint{4.024022in}{2.474036in}}%
\pgfpathquadraticcurveto{\pgfqpoint{4.033713in}{2.474036in}}{\pgfqpoint{4.038917in}{2.477354in}}%
\pgfpathquadraticcurveto{\pgfqpoint{4.044145in}{2.480673in}}{\pgfqpoint{4.044145in}{2.486712in}}%
\pgfpathquadraticcurveto{\pgfqpoint{4.044145in}{2.492298in}}{\pgfqpoint{4.040374in}{2.495282in}}%
\pgfpathquadraticcurveto{\pgfqpoint{4.036626in}{2.498266in}}{\pgfqpoint{4.023854in}{2.501035in}}%
\pgfpathlineto{\pgfqpoint{4.019080in}{2.502157in}}%
\pgfpathquadraticcurveto{\pgfqpoint{4.006333in}{2.504831in}}{\pgfqpoint{4.000651in}{2.510393in}}%
\pgfpathquadraticcurveto{\pgfqpoint{3.994994in}{2.515955in}}{\pgfqpoint{3.994994in}{2.525646in}}%
\pgfpathquadraticcurveto{\pgfqpoint{3.994994in}{2.537439in}}{\pgfqpoint{4.003349in}{2.543837in}}%
\pgfpathquadraticcurveto{\pgfqpoint{4.011704in}{2.550258in}}{\pgfqpoint{4.027077in}{2.550258in}}%
\pgfpathquadraticcurveto{\pgfqpoint{4.034668in}{2.550258in}}{\pgfqpoint{4.041376in}{2.549136in}}%
\pgfpathquadraticcurveto{\pgfqpoint{4.048108in}{2.548038in}}{\pgfqpoint{4.053765in}{2.545794in}}%
\pgfpathlineto{\pgfqpoint{4.053765in}{2.545794in}}%
\pgfpathclose%
\pgfpathmoveto{\pgfqpoint{4.093679in}{2.571981in}}%
\pgfpathlineto{\pgfqpoint{4.093679in}{2.548253in}}%
\pgfpathlineto{\pgfqpoint{4.121943in}{2.548253in}}%
\pgfpathlineto{\pgfqpoint{4.121943in}{2.537582in}}%
\pgfpathlineto{\pgfqpoint{4.093679in}{2.537582in}}%
\pgfpathlineto{\pgfqpoint{4.093679in}{2.492226in}}%
\pgfpathquadraticcurveto{\pgfqpoint{4.093679in}{2.482009in}}{\pgfqpoint{4.096472in}{2.479097in}}%
\pgfpathquadraticcurveto{\pgfqpoint{4.099265in}{2.476185in}}{\pgfqpoint{4.107858in}{2.476185in}}%
\pgfpathlineto{\pgfqpoint{4.121943in}{2.476185in}}%
\pgfpathlineto{\pgfqpoint{4.121943in}{2.464702in}}%
\pgfpathlineto{\pgfqpoint{4.107858in}{2.464702in}}%
\pgfpathquadraticcurveto{\pgfqpoint{4.091960in}{2.464702in}}{\pgfqpoint{4.085920in}{2.470623in}}%
\pgfpathquadraticcurveto{\pgfqpoint{4.079881in}{2.476567in}}{\pgfqpoint{4.079881in}{2.492226in}}%
\pgfpathlineto{\pgfqpoint{4.079881in}{2.537582in}}%
\pgfpathlineto{\pgfqpoint{4.069807in}{2.537582in}}%
\pgfpathlineto{\pgfqpoint{4.069807in}{2.548253in}}%
\pgfpathlineto{\pgfqpoint{4.079881in}{2.548253in}}%
\pgfpathlineto{\pgfqpoint{4.079881in}{2.571981in}}%
\pgfpathlineto{\pgfqpoint{4.093679in}{2.571981in}}%
\pgfpathlineto{\pgfqpoint{4.093679in}{2.571981in}}%
\pgfpathclose%
\pgfpathmoveto{\pgfqpoint{4.139989in}{2.548253in}}%
\pgfpathlineto{\pgfqpoint{4.153716in}{2.548253in}}%
\pgfpathlineto{\pgfqpoint{4.153716in}{2.464702in}}%
\pgfpathlineto{\pgfqpoint{4.139989in}{2.464702in}}%
\pgfpathlineto{\pgfqpoint{4.139989in}{2.548253in}}%
\pgfpathlineto{\pgfqpoint{4.139989in}{2.548253in}}%
\pgfpathclose%
\pgfpathmoveto{\pgfqpoint{4.139989in}{2.580790in}}%
\pgfpathlineto{\pgfqpoint{4.153716in}{2.580790in}}%
\pgfpathlineto{\pgfqpoint{4.153716in}{2.563387in}}%
\pgfpathlineto{\pgfqpoint{4.139989in}{2.563387in}}%
\pgfpathlineto{\pgfqpoint{4.139989in}{2.580790in}}%
\pgfpathlineto{\pgfqpoint{4.139989in}{2.580790in}}%
\pgfpathclose%
\pgfpathmoveto{\pgfqpoint{4.247483in}{2.532211in}}%
\pgfpathquadraticcurveto{\pgfqpoint{4.252639in}{2.541473in}}{\pgfqpoint{4.259801in}{2.545866in}}%
\pgfpathquadraticcurveto{\pgfqpoint{4.266962in}{2.550258in}}{\pgfqpoint{4.276654in}{2.550258in}}%
\pgfpathquadraticcurveto{\pgfqpoint{4.289712in}{2.550258in}}{\pgfqpoint{4.296801in}{2.541115in}}%
\pgfpathquadraticcurveto{\pgfqpoint{4.303891in}{2.531996in}}{\pgfqpoint{4.303891in}{2.515143in}}%
\pgfpathlineto{\pgfqpoint{4.303891in}{2.464702in}}%
\pgfpathlineto{\pgfqpoint{4.290094in}{2.464702in}}%
\pgfpathlineto{\pgfqpoint{4.290094in}{2.514689in}}%
\pgfpathquadraticcurveto{\pgfqpoint{4.290094in}{2.526697in}}{\pgfqpoint{4.285821in}{2.532498in}}%
\pgfpathquadraticcurveto{\pgfqpoint{4.281571in}{2.538322in}}{\pgfqpoint{4.272858in}{2.538322in}}%
\pgfpathquadraticcurveto{\pgfqpoint{4.262188in}{2.538322in}}{\pgfqpoint{4.255981in}{2.531232in}}%
\pgfpathquadraticcurveto{\pgfqpoint{4.249798in}{2.524166in}}{\pgfqpoint{4.249798in}{2.511920in}}%
\pgfpathlineto{\pgfqpoint{4.249798in}{2.464702in}}%
\pgfpathlineto{\pgfqpoint{4.236001in}{2.464702in}}%
\pgfpathlineto{\pgfqpoint{4.236001in}{2.514689in}}%
\pgfpathquadraticcurveto{\pgfqpoint{4.236001in}{2.526768in}}{\pgfqpoint{4.231752in}{2.532545in}}%
\pgfpathquadraticcurveto{\pgfqpoint{4.227502in}{2.538322in}}{\pgfqpoint{4.218622in}{2.538322in}}%
\pgfpathquadraticcurveto{\pgfqpoint{4.208095in}{2.538322in}}{\pgfqpoint{4.201888in}{2.531209in}}%
\pgfpathquadraticcurveto{\pgfqpoint{4.195706in}{2.524095in}}{\pgfqpoint{4.195706in}{2.511920in}}%
\pgfpathlineto{\pgfqpoint{4.195706in}{2.464702in}}%
\pgfpathlineto{\pgfqpoint{4.181908in}{2.464702in}}%
\pgfpathlineto{\pgfqpoint{4.181908in}{2.548253in}}%
\pgfpathlineto{\pgfqpoint{4.195706in}{2.548253in}}%
\pgfpathlineto{\pgfqpoint{4.195706in}{2.535267in}}%
\pgfpathquadraticcurveto{\pgfqpoint{4.200408in}{2.542953in}}{\pgfqpoint{4.206973in}{2.546606in}}%
\pgfpathquadraticcurveto{\pgfqpoint{4.213538in}{2.550258in}}{\pgfqpoint{4.222561in}{2.550258in}}%
\pgfpathquadraticcurveto{\pgfqpoint{4.231680in}{2.550258in}}{\pgfqpoint{4.238054in}{2.545627in}}%
\pgfpathquadraticcurveto{\pgfqpoint{4.244427in}{2.541020in}}{\pgfqpoint{4.247483in}{2.532211in}}%
\pgfpathlineto{\pgfqpoint{4.247483in}{2.532211in}}%
\pgfpathclose%
\pgfpathmoveto{\pgfqpoint{4.369228in}{2.506692in}}%
\pgfpathquadraticcurveto{\pgfqpoint{4.352589in}{2.506692in}}{\pgfqpoint{4.346168in}{2.502897in}}%
\pgfpathquadraticcurveto{\pgfqpoint{4.339746in}{2.499101in}}{\pgfqpoint{4.339746in}{2.489911in}}%
\pgfpathquadraticcurveto{\pgfqpoint{4.339746in}{2.482606in}}{\pgfqpoint{4.344568in}{2.478309in}}%
\pgfpathquadraticcurveto{\pgfqpoint{4.349390in}{2.474036in}}{\pgfqpoint{4.357650in}{2.474036in}}%
\pgfpathquadraticcurveto{\pgfqpoint{4.369084in}{2.474036in}}{\pgfqpoint{4.375983in}{2.482129in}}%
\pgfpathquadraticcurveto{\pgfqpoint{4.382882in}{2.490221in}}{\pgfqpoint{4.382882in}{2.503637in}}%
\pgfpathlineto{\pgfqpoint{4.382882in}{2.506692in}}%
\pgfpathlineto{\pgfqpoint{4.369228in}{2.506692in}}%
\pgfpathlineto{\pgfqpoint{4.369228in}{2.506692in}}%
\pgfpathclose%
\pgfpathmoveto{\pgfqpoint{4.396608in}{2.512374in}}%
\pgfpathlineto{\pgfqpoint{4.396608in}{2.464702in}}%
\pgfpathlineto{\pgfqpoint{4.382882in}{2.464702in}}%
\pgfpathlineto{\pgfqpoint{4.382882in}{2.477378in}}%
\pgfpathquadraticcurveto{\pgfqpoint{4.378179in}{2.469787in}}{\pgfqpoint{4.371161in}{2.466159in}}%
\pgfpathquadraticcurveto{\pgfqpoint{4.364143in}{2.462530in}}{\pgfqpoint{4.353998in}{2.462530in}}%
\pgfpathquadraticcurveto{\pgfqpoint{4.341179in}{2.462530in}}{\pgfqpoint{4.333587in}{2.469739in}}%
\pgfpathquadraticcurveto{\pgfqpoint{4.326020in}{2.476949in}}{\pgfqpoint{4.326020in}{2.489028in}}%
\pgfpathquadraticcurveto{\pgfqpoint{4.326020in}{2.503112in}}{\pgfqpoint{4.335449in}{2.510273in}}%
\pgfpathquadraticcurveto{\pgfqpoint{4.344903in}{2.517435in}}{\pgfqpoint{4.363618in}{2.517435in}}%
\pgfpathlineto{\pgfqpoint{4.382882in}{2.517435in}}%
\pgfpathlineto{\pgfqpoint{4.382882in}{2.518795in}}%
\pgfpathquadraticcurveto{\pgfqpoint{4.382882in}{2.528272in}}{\pgfqpoint{4.376652in}{2.533452in}}%
\pgfpathquadraticcurveto{\pgfqpoint{4.370421in}{2.538633in}}{\pgfqpoint{4.359154in}{2.538633in}}%
\pgfpathquadraticcurveto{\pgfqpoint{4.351992in}{2.538633in}}{\pgfqpoint{4.345189in}{2.536914in}}%
\pgfpathquadraticcurveto{\pgfqpoint{4.338410in}{2.535195in}}{\pgfqpoint{4.332155in}{2.531758in}}%
\pgfpathlineto{\pgfqpoint{4.332155in}{2.544457in}}%
\pgfpathquadraticcurveto{\pgfqpoint{4.339675in}{2.547370in}}{\pgfqpoint{4.346765in}{2.548802in}}%
\pgfpathquadraticcurveto{\pgfqpoint{4.353854in}{2.550258in}}{\pgfqpoint{4.360562in}{2.550258in}}%
\pgfpathquadraticcurveto{\pgfqpoint{4.378705in}{2.550258in}}{\pgfqpoint{4.387656in}{2.540853in}}%
\pgfpathquadraticcurveto{\pgfqpoint{4.396608in}{2.531471in}}{\pgfqpoint{4.396608in}{2.512374in}}%
\pgfpathlineto{\pgfqpoint{4.396608in}{2.512374in}}%
\pgfpathclose%
\pgfpathmoveto{\pgfqpoint{4.438455in}{2.571981in}}%
\pgfpathlineto{\pgfqpoint{4.438455in}{2.548253in}}%
\pgfpathlineto{\pgfqpoint{4.466719in}{2.548253in}}%
\pgfpathlineto{\pgfqpoint{4.466719in}{2.537582in}}%
\pgfpathlineto{\pgfqpoint{4.438455in}{2.537582in}}%
\pgfpathlineto{\pgfqpoint{4.438455in}{2.492226in}}%
\pgfpathquadraticcurveto{\pgfqpoint{4.438455in}{2.482009in}}{\pgfqpoint{4.441248in}{2.479097in}}%
\pgfpathquadraticcurveto{\pgfqpoint{4.444041in}{2.476185in}}{\pgfqpoint{4.452635in}{2.476185in}}%
\pgfpathlineto{\pgfqpoint{4.466719in}{2.476185in}}%
\pgfpathlineto{\pgfqpoint{4.466719in}{2.464702in}}%
\pgfpathlineto{\pgfqpoint{4.452635in}{2.464702in}}%
\pgfpathquadraticcurveto{\pgfqpoint{4.436736in}{2.464702in}}{\pgfqpoint{4.430697in}{2.470623in}}%
\pgfpathquadraticcurveto{\pgfqpoint{4.424657in}{2.476567in}}{\pgfqpoint{4.424657in}{2.492226in}}%
\pgfpathlineto{\pgfqpoint{4.424657in}{2.537582in}}%
\pgfpathlineto{\pgfqpoint{4.414584in}{2.537582in}}%
\pgfpathlineto{\pgfqpoint{4.414584in}{2.548253in}}%
\pgfpathlineto{\pgfqpoint{4.424657in}{2.548253in}}%
\pgfpathlineto{\pgfqpoint{4.424657in}{2.571981in}}%
\pgfpathlineto{\pgfqpoint{4.438455in}{2.571981in}}%
\pgfpathlineto{\pgfqpoint{4.438455in}{2.571981in}}%
\pgfpathclose%
\pgfpathmoveto{\pgfqpoint{4.556237in}{2.509915in}}%
\pgfpathlineto{\pgfqpoint{4.556237in}{2.503207in}}%
\pgfpathlineto{\pgfqpoint{4.493121in}{2.503207in}}%
\pgfpathquadraticcurveto{\pgfqpoint{4.494028in}{2.489028in}}{\pgfqpoint{4.501667in}{2.481604in}}%
\pgfpathquadraticcurveto{\pgfqpoint{4.509306in}{2.474179in}}{\pgfqpoint{4.522960in}{2.474179in}}%
\pgfpathquadraticcurveto{\pgfqpoint{4.530862in}{2.474179in}}{\pgfqpoint{4.538286in}{2.476113in}}%
\pgfpathquadraticcurveto{\pgfqpoint{4.545710in}{2.478047in}}{\pgfqpoint{4.553038in}{2.481938in}}%
\pgfpathlineto{\pgfqpoint{4.553038in}{2.468952in}}%
\pgfpathquadraticcurveto{\pgfqpoint{4.545638in}{2.465824in}}{\pgfqpoint{4.537880in}{2.464177in}}%
\pgfpathquadraticcurveto{\pgfqpoint{4.530122in}{2.462530in}}{\pgfqpoint{4.522149in}{2.462530in}}%
\pgfpathquadraticcurveto{\pgfqpoint{4.502144in}{2.462530in}}{\pgfqpoint{4.490471in}{2.474156in}}%
\pgfpathquadraticcurveto{\pgfqpoint{4.478798in}{2.485805in}}{\pgfqpoint{4.478798in}{2.505666in}}%
\pgfpathquadraticcurveto{\pgfqpoint{4.478798in}{2.526172in}}{\pgfqpoint{4.489874in}{2.538203in}}%
\pgfpathquadraticcurveto{\pgfqpoint{4.500951in}{2.550258in}}{\pgfqpoint{4.519762in}{2.550258in}}%
\pgfpathquadraticcurveto{\pgfqpoint{4.536615in}{2.550258in}}{\pgfqpoint{4.546426in}{2.539396in}}%
\pgfpathquadraticcurveto{\pgfqpoint{4.556237in}{2.528559in}}{\pgfqpoint{4.556237in}{2.509915in}}%
\pgfpathlineto{\pgfqpoint{4.556237in}{2.509915in}}%
\pgfpathclose%
\pgfpathmoveto{\pgfqpoint{4.542511in}{2.513949in}}%
\pgfpathquadraticcurveto{\pgfqpoint{4.542368in}{2.525193in}}{\pgfqpoint{4.536209in}{2.531901in}}%
\pgfpathquadraticcurveto{\pgfqpoint{4.530050in}{2.538633in}}{\pgfqpoint{4.519905in}{2.538633in}}%
\pgfpathquadraticcurveto{\pgfqpoint{4.508423in}{2.538633in}}{\pgfqpoint{4.501524in}{2.532140in}}%
\pgfpathquadraticcurveto{\pgfqpoint{4.494625in}{2.525646in}}{\pgfqpoint{4.493574in}{2.513854in}}%
\pgfpathlineto{\pgfqpoint{4.542511in}{2.513949in}}%
\pgfpathlineto{\pgfqpoint{4.542511in}{2.513949in}}%
\pgfpathclose%
\pgfpathmoveto{\pgfqpoint{4.660293in}{2.580623in}}%
\pgfpathquadraticcurveto{\pgfqpoint{4.650315in}{2.563483in}}{\pgfqpoint{4.645445in}{2.546677in}}%
\pgfpathquadraticcurveto{\pgfqpoint{4.640599in}{2.529896in}}{\pgfqpoint{4.640599in}{2.512660in}}%
\pgfpathquadraticcurveto{\pgfqpoint{4.640599in}{2.495449in}}{\pgfqpoint{4.645493in}{2.478548in}}%
\pgfpathquadraticcurveto{\pgfqpoint{4.650387in}{2.461647in}}{\pgfqpoint{4.660293in}{2.444555in}}%
\pgfpathlineto{\pgfqpoint{4.648357in}{2.444555in}}%
\pgfpathquadraticcurveto{\pgfqpoint{4.637186in}{2.462100in}}{\pgfqpoint{4.631624in}{2.479025in}}%
\pgfpathquadraticcurveto{\pgfqpoint{4.626061in}{2.495950in}}{\pgfqpoint{4.626061in}{2.512660in}}%
\pgfpathquadraticcurveto{\pgfqpoint{4.626061in}{2.529299in}}{\pgfqpoint{4.631576in}{2.546152in}}%
\pgfpathquadraticcurveto{\pgfqpoint{4.637114in}{2.563029in}}{\pgfqpoint{4.648357in}{2.580623in}}%
\pgfpathlineto{\pgfqpoint{4.660293in}{2.580623in}}%
\pgfpathlineto{\pgfqpoint{4.660293in}{2.580623in}}%
\pgfpathclose%
\pgfpathmoveto{\pgfqpoint{4.672086in}{2.576087in}}%
\pgfpathlineto{\pgfqpoint{4.766307in}{2.576087in}}%
\pgfpathlineto{\pgfqpoint{4.766307in}{2.563387in}}%
\pgfpathlineto{\pgfqpoint{4.726775in}{2.563387in}}%
\pgfpathlineto{\pgfqpoint{4.726775in}{2.464702in}}%
\pgfpathlineto{\pgfqpoint{4.711641in}{2.464702in}}%
\pgfpathlineto{\pgfqpoint{4.711641in}{2.563387in}}%
\pgfpathlineto{\pgfqpoint{4.672086in}{2.563387in}}%
\pgfpathlineto{\pgfqpoint{4.672086in}{2.576087in}}%
\pgfpathlineto{\pgfqpoint{4.672086in}{2.576087in}}%
\pgfpathclose%
\pgfpathmoveto{\pgfqpoint{4.770938in}{2.576087in}}%
\pgfpathlineto{\pgfqpoint{4.786144in}{2.576087in}}%
\pgfpathlineto{\pgfqpoint{4.809562in}{2.481938in}}%
\pgfpathlineto{\pgfqpoint{4.832908in}{2.576087in}}%
\pgfpathlineto{\pgfqpoint{4.849857in}{2.576087in}}%
\pgfpathlineto{\pgfqpoint{4.873275in}{2.481938in}}%
\pgfpathlineto{\pgfqpoint{4.896621in}{2.576087in}}%
\pgfpathlineto{\pgfqpoint{4.911923in}{2.576087in}}%
\pgfpathlineto{\pgfqpoint{4.883946in}{2.464702in}}%
\pgfpathlineto{\pgfqpoint{4.864992in}{2.464702in}}%
\pgfpathlineto{\pgfqpoint{4.841502in}{2.561382in}}%
\pgfpathlineto{\pgfqpoint{4.817774in}{2.464702in}}%
\pgfpathlineto{\pgfqpoint{4.798820in}{2.464702in}}%
\pgfpathlineto{\pgfqpoint{4.770938in}{2.576087in}}%
\pgfpathlineto{\pgfqpoint{4.770938in}{2.576087in}}%
\pgfpathclose%
\pgfpathmoveto{\pgfqpoint{4.931903in}{2.576087in}}%
\pgfpathlineto{\pgfqpoint{4.995903in}{2.576087in}}%
\pgfpathlineto{\pgfqpoint{4.995903in}{2.563387in}}%
\pgfpathlineto{\pgfqpoint{4.946966in}{2.563387in}}%
\pgfpathlineto{\pgfqpoint{4.946966in}{2.530564in}}%
\pgfpathlineto{\pgfqpoint{4.991129in}{2.530564in}}%
\pgfpathlineto{\pgfqpoint{4.991129in}{2.517888in}}%
\pgfpathlineto{\pgfqpoint{4.946966in}{2.517888in}}%
\pgfpathlineto{\pgfqpoint{4.946966in}{2.464702in}}%
\pgfpathlineto{\pgfqpoint{4.931903in}{2.464702in}}%
\pgfpathlineto{\pgfqpoint{4.931903in}{2.576087in}}%
\pgfpathlineto{\pgfqpoint{4.931903in}{2.576087in}}%
\pgfpathclose%
\pgfpathmoveto{\pgfqpoint{5.019775in}{2.576087in}}%
\pgfpathlineto{\pgfqpoint{5.090196in}{2.576087in}}%
\pgfpathlineto{\pgfqpoint{5.090196in}{2.563387in}}%
\pgfpathlineto{\pgfqpoint{5.034837in}{2.563387in}}%
\pgfpathlineto{\pgfqpoint{5.034837in}{2.530421in}}%
\pgfpathlineto{\pgfqpoint{5.087880in}{2.530421in}}%
\pgfpathlineto{\pgfqpoint{5.087880in}{2.517745in}}%
\pgfpathlineto{\pgfqpoint{5.034837in}{2.517745in}}%
\pgfpathlineto{\pgfqpoint{5.034837in}{2.477378in}}%
\pgfpathlineto{\pgfqpoint{5.091532in}{2.477378in}}%
\pgfpathlineto{\pgfqpoint{5.091532in}{2.464702in}}%
\pgfpathlineto{\pgfqpoint{5.019775in}{2.464702in}}%
\pgfpathlineto{\pgfqpoint{5.019775in}{2.576087in}}%
\pgfpathlineto{\pgfqpoint{5.019775in}{2.576087in}}%
\pgfpathclose%
\pgfpathmoveto{\pgfqpoint{5.231634in}{2.572435in}}%
\pgfpathlineto{\pgfqpoint{5.231634in}{2.557730in}}%
\pgfpathquadraticcurveto{\pgfqpoint{5.223064in}{2.561836in}}{\pgfqpoint{5.215449in}{2.563841in}}%
\pgfpathquadraticcurveto{\pgfqpoint{5.207834in}{2.565870in}}{\pgfqpoint{5.200745in}{2.565870in}}%
\pgfpathquadraticcurveto{\pgfqpoint{5.188451in}{2.565870in}}{\pgfqpoint{5.181767in}{2.561096in}}%
\pgfpathquadraticcurveto{\pgfqpoint{5.175083in}{2.556321in}}{\pgfqpoint{5.175083in}{2.547513in}}%
\pgfpathquadraticcurveto{\pgfqpoint{5.175083in}{2.540113in}}{\pgfqpoint{5.179523in}{2.536341in}}%
\pgfpathquadraticcurveto{\pgfqpoint{5.183963in}{2.532593in}}{\pgfqpoint{5.196352in}{2.530278in}}%
\pgfpathlineto{\pgfqpoint{5.205447in}{2.528416in}}%
\pgfpathquadraticcurveto{\pgfqpoint{5.222301in}{2.525193in}}{\pgfqpoint{5.230321in}{2.517100in}}%
\pgfpathquadraticcurveto{\pgfqpoint{5.238342in}{2.509008in}}{\pgfqpoint{5.238342in}{2.495449in}}%
\pgfpathquadraticcurveto{\pgfqpoint{5.238342in}{2.479240in}}{\pgfqpoint{5.227481in}{2.470885in}}%
\pgfpathquadraticcurveto{\pgfqpoint{5.216643in}{2.462530in}}{\pgfqpoint{5.195684in}{2.462530in}}%
\pgfpathquadraticcurveto{\pgfqpoint{5.187782in}{2.462530in}}{\pgfqpoint{5.178854in}{2.464321in}}%
\pgfpathquadraticcurveto{\pgfqpoint{5.169950in}{2.466111in}}{\pgfqpoint{5.160402in}{2.469620in}}%
\pgfpathlineto{\pgfqpoint{5.160402in}{2.485137in}}%
\pgfpathquadraticcurveto{\pgfqpoint{5.169568in}{2.480004in}}{\pgfqpoint{5.178377in}{2.477378in}}%
\pgfpathquadraticcurveto{\pgfqpoint{5.187186in}{2.474776in}}{\pgfqpoint{5.195684in}{2.474776in}}%
\pgfpathquadraticcurveto{\pgfqpoint{5.208574in}{2.474776in}}{\pgfqpoint{5.215593in}{2.479837in}}%
\pgfpathquadraticcurveto{\pgfqpoint{5.222611in}{2.484922in}}{\pgfqpoint{5.222611in}{2.494327in}}%
\pgfpathquadraticcurveto{\pgfqpoint{5.222611in}{2.502515in}}{\pgfqpoint{5.217574in}{2.507146in}}%
\pgfpathquadraticcurveto{\pgfqpoint{5.212537in}{2.511777in}}{\pgfqpoint{5.201055in}{2.514093in}}%
\pgfpathlineto{\pgfqpoint{5.191864in}{2.515883in}}%
\pgfpathquadraticcurveto{\pgfqpoint{5.175011in}{2.519225in}}{\pgfqpoint{5.167468in}{2.526387in}}%
\pgfpathquadraticcurveto{\pgfqpoint{5.159948in}{2.533548in}}{\pgfqpoint{5.159948in}{2.546319in}}%
\pgfpathquadraticcurveto{\pgfqpoint{5.159948in}{2.561096in}}{\pgfqpoint{5.170356in}{2.569594in}}%
\pgfpathquadraticcurveto{\pgfqpoint{5.180764in}{2.578092in}}{\pgfqpoint{5.199026in}{2.578092in}}%
\pgfpathquadraticcurveto{\pgfqpoint{5.206880in}{2.578092in}}{\pgfqpoint{5.214996in}{2.576660in}}%
\pgfpathquadraticcurveto{\pgfqpoint{5.223136in}{2.575252in}}{\pgfqpoint{5.231634in}{2.572435in}}%
\pgfpathlineto{\pgfqpoint{5.231634in}{2.572435in}}%
\pgfpathclose%
\pgfpathmoveto{\pgfqpoint{5.261856in}{2.576087in}}%
\pgfpathlineto{\pgfqpoint{5.332277in}{2.576087in}}%
\pgfpathlineto{\pgfqpoint{5.332277in}{2.563387in}}%
\pgfpathlineto{\pgfqpoint{5.276919in}{2.563387in}}%
\pgfpathlineto{\pgfqpoint{5.276919in}{2.530421in}}%
\pgfpathlineto{\pgfqpoint{5.329961in}{2.530421in}}%
\pgfpathlineto{\pgfqpoint{5.329961in}{2.517745in}}%
\pgfpathlineto{\pgfqpoint{5.276919in}{2.517745in}}%
\pgfpathlineto{\pgfqpoint{5.276919in}{2.477378in}}%
\pgfpathlineto{\pgfqpoint{5.333614in}{2.477378in}}%
\pgfpathlineto{\pgfqpoint{5.333614in}{2.464702in}}%
\pgfpathlineto{\pgfqpoint{5.261856in}{2.464702in}}%
\pgfpathlineto{\pgfqpoint{5.261856in}{2.576087in}}%
\pgfpathlineto{\pgfqpoint{5.261856in}{2.576087in}}%
\pgfpathclose%
\pgfpathmoveto{\pgfqpoint{5.411053in}{2.545794in}}%
\pgfpathlineto{\pgfqpoint{5.411053in}{2.532808in}}%
\pgfpathquadraticcurveto{\pgfqpoint{5.405252in}{2.535792in}}{\pgfqpoint{5.398974in}{2.537272in}}%
\pgfpathquadraticcurveto{\pgfqpoint{5.392719in}{2.538776in}}{\pgfqpoint{5.385988in}{2.538776in}}%
\pgfpathquadraticcurveto{\pgfqpoint{5.375771in}{2.538776in}}{\pgfqpoint{5.370662in}{2.535649in}}%
\pgfpathquadraticcurveto{\pgfqpoint{5.365554in}{2.532521in}}{\pgfqpoint{5.365554in}{2.526243in}}%
\pgfpathquadraticcurveto{\pgfqpoint{5.365554in}{2.521469in}}{\pgfqpoint{5.369206in}{2.518748in}}%
\pgfpathquadraticcurveto{\pgfqpoint{5.372858in}{2.516026in}}{\pgfqpoint{5.383911in}{2.513567in}}%
\pgfpathlineto{\pgfqpoint{5.388614in}{2.512517in}}%
\pgfpathquadraticcurveto{\pgfqpoint{5.403223in}{2.509390in}}{\pgfqpoint{5.409382in}{2.503685in}}%
\pgfpathquadraticcurveto{\pgfqpoint{5.415541in}{2.497979in}}{\pgfqpoint{5.415541in}{2.487762in}}%
\pgfpathquadraticcurveto{\pgfqpoint{5.415541in}{2.476113in}}{\pgfqpoint{5.406326in}{2.469310in}}%
\pgfpathquadraticcurveto{\pgfqpoint{5.397112in}{2.462530in}}{\pgfqpoint{5.380999in}{2.462530in}}%
\pgfpathquadraticcurveto{\pgfqpoint{5.374291in}{2.462530in}}{\pgfqpoint{5.367010in}{2.463843in}}%
\pgfpathquadraticcurveto{\pgfqpoint{5.359729in}{2.465156in}}{\pgfqpoint{5.351684in}{2.467758in}}%
\pgfpathlineto{\pgfqpoint{5.351684in}{2.481938in}}%
\pgfpathquadraticcurveto{\pgfqpoint{5.359299in}{2.477975in}}{\pgfqpoint{5.366676in}{2.475994in}}%
\pgfpathquadraticcurveto{\pgfqpoint{5.374052in}{2.474036in}}{\pgfqpoint{5.381309in}{2.474036in}}%
\pgfpathquadraticcurveto{\pgfqpoint{5.391001in}{2.474036in}}{\pgfqpoint{5.396205in}{2.477354in}}%
\pgfpathquadraticcurveto{\pgfqpoint{5.401433in}{2.480673in}}{\pgfqpoint{5.401433in}{2.486712in}}%
\pgfpathquadraticcurveto{\pgfqpoint{5.401433in}{2.492298in}}{\pgfqpoint{5.397661in}{2.495282in}}%
\pgfpathquadraticcurveto{\pgfqpoint{5.393913in}{2.498266in}}{\pgfqpoint{5.381142in}{2.501035in}}%
\pgfpathlineto{\pgfqpoint{5.376367in}{2.502157in}}%
\pgfpathquadraticcurveto{\pgfqpoint{5.363620in}{2.504831in}}{\pgfqpoint{5.357939in}{2.510393in}}%
\pgfpathquadraticcurveto{\pgfqpoint{5.352281in}{2.515955in}}{\pgfqpoint{5.352281in}{2.525646in}}%
\pgfpathquadraticcurveto{\pgfqpoint{5.352281in}{2.537439in}}{\pgfqpoint{5.360636in}{2.543837in}}%
\pgfpathquadraticcurveto{\pgfqpoint{5.368991in}{2.550258in}}{\pgfqpoint{5.384364in}{2.550258in}}%
\pgfpathquadraticcurveto{\pgfqpoint{5.391956in}{2.550258in}}{\pgfqpoint{5.398663in}{2.549136in}}%
\pgfpathquadraticcurveto{\pgfqpoint{5.405395in}{2.548038in}}{\pgfqpoint{5.411053in}{2.545794in}}%
\pgfpathlineto{\pgfqpoint{5.411053in}{2.545794in}}%
\pgfpathclose%
\pgfpathmoveto{\pgfqpoint{5.435235in}{2.580623in}}%
\pgfpathlineto{\pgfqpoint{5.447170in}{2.580623in}}%
\pgfpathquadraticcurveto{\pgfqpoint{5.458342in}{2.563029in}}{\pgfqpoint{5.463904in}{2.546152in}}%
\pgfpathquadraticcurveto{\pgfqpoint{5.469466in}{2.529299in}}{\pgfqpoint{5.469466in}{2.512660in}}%
\pgfpathquadraticcurveto{\pgfqpoint{5.469466in}{2.495950in}}{\pgfqpoint{5.463904in}{2.479025in}}%
\pgfpathquadraticcurveto{\pgfqpoint{5.458342in}{2.462100in}}{\pgfqpoint{5.447170in}{2.444555in}}%
\pgfpathlineto{\pgfqpoint{5.435235in}{2.444555in}}%
\pgfpathquadraticcurveto{\pgfqpoint{5.445141in}{2.461647in}}{\pgfqpoint{5.450035in}{2.478548in}}%
\pgfpathquadraticcurveto{\pgfqpoint{5.454929in}{2.495449in}}{\pgfqpoint{5.454929in}{2.512660in}}%
\pgfpathquadraticcurveto{\pgfqpoint{5.454929in}{2.529896in}}{\pgfqpoint{5.450035in}{2.546677in}}%
\pgfpathquadraticcurveto{\pgfqpoint{5.445141in}{2.563483in}}{\pgfqpoint{5.435235in}{2.580623in}}%
\pgfpathlineto{\pgfqpoint{5.435235in}{2.580623in}}%
\pgfpathclose%
\pgfusepath{fill}%
\end{pgfscope}%
\begin{pgfscope}%
\pgfpathrectangle{\pgfqpoint{0.977909in}{2.928000in}}{\pgfqpoint{5.664000in}{5.664000in}}%
\pgfusepath{clip}%
\pgfsetrectcap%
\pgfsetroundjoin%
\pgfsetlinewidth{0.602250pt}%
\definecolor{currentstroke}{rgb}{0.898039,0.913725,0.929412}%
\pgfsetstrokecolor{currentstroke}%
\pgfsetdash{}{0pt}%
\pgfpathmoveto{\pgfqpoint{0.977909in}{2.928000in}}%
\pgfpathlineto{\pgfqpoint{6.641909in}{2.928000in}}%
\pgfusepath{stroke}%
\end{pgfscope}%
\begin{pgfscope}%
\pgfsetbuttcap%
\pgfsetroundjoin%
\definecolor{currentfill}{rgb}{0.090196,0.168627,0.227451}%
\pgfsetfillcolor{currentfill}%
\pgfsetlinewidth{0.803000pt}%
\definecolor{currentstroke}{rgb}{0.090196,0.168627,0.227451}%
\pgfsetstrokecolor{currentstroke}%
\pgfsetdash{}{0pt}%
\pgfsys@defobject{currentmarker}{\pgfqpoint{-0.048611in}{0.000000in}}{\pgfqpoint{-0.000000in}{0.000000in}}{%
\pgfpathmoveto{\pgfqpoint{-0.000000in}{0.000000in}}%
\pgfpathlineto{\pgfqpoint{-0.048611in}{0.000000in}}%
\pgfusepath{stroke,fill}%
}%
\begin{pgfscope}%
\pgfsys@transformshift{0.977909in}{2.928000in}%
\pgfsys@useobject{currentmarker}{}%
\end{pgfscope}%
\end{pgfscope}%
\begin{pgfscope}%
\pgfsetbuttcap%
\pgfsetroundjoin%
\definecolor{currentfill}{rgb}{0.090196,0.168627,0.227451}%
\pgfsetfillcolor{currentfill}%
\pgfsetlinewidth{0.000000pt}%
\definecolor{currentstroke}{rgb}{0.090196,0.168627,0.227451}%
\pgfsetstrokecolor{currentstroke}%
\pgfsetdash{}{0pt}%
\pgfpathmoveto{\pgfqpoint{0.583929in}{2.924890in}}%
\pgfpathlineto{\pgfqpoint{0.662171in}{2.924890in}}%
\pgfpathlineto{\pgfqpoint{0.662171in}{2.914519in}}%
\pgfpathlineto{\pgfqpoint{0.583929in}{2.914519in}}%
\pgfpathlineto{\pgfqpoint{0.583929in}{2.924890in}}%
\pgfpathlineto{\pgfqpoint{0.583929in}{2.924890in}}%
\pgfpathclose%
\pgfpathmoveto{\pgfqpoint{0.688929in}{2.971647in}}%
\pgfpathlineto{\pgfqpoint{0.737327in}{2.971647in}}%
\pgfpathlineto{\pgfqpoint{0.737327in}{2.961257in}}%
\pgfpathlineto{\pgfqpoint{0.700218in}{2.961257in}}%
\pgfpathlineto{\pgfqpoint{0.700218in}{2.938933in}}%
\pgfpathquadraticcurveto{\pgfqpoint{0.702894in}{2.939851in}}{\pgfqpoint{0.705570in}{2.940300in}}%
\pgfpathquadraticcurveto{\pgfqpoint{0.708265in}{2.940749in}}{\pgfqpoint{0.710960in}{2.940749in}}%
\pgfpathquadraticcurveto{\pgfqpoint{0.726214in}{2.940749in}}{\pgfqpoint{0.735120in}{2.932390in}}%
\pgfpathquadraticcurveto{\pgfqpoint{0.744046in}{2.924030in}}{\pgfqpoint{0.744046in}{2.909753in}}%
\pgfpathquadraticcurveto{\pgfqpoint{0.744046in}{2.895046in}}{\pgfqpoint{0.734886in}{2.886882in}}%
\pgfpathquadraticcurveto{\pgfqpoint{0.725726in}{2.878737in}}{\pgfqpoint{0.709066in}{2.878737in}}%
\pgfpathquadraticcurveto{\pgfqpoint{0.703324in}{2.878737in}}{\pgfqpoint{0.697367in}{2.879714in}}%
\pgfpathquadraticcurveto{\pgfqpoint{0.691429in}{2.880690in}}{\pgfqpoint{0.685081in}{2.882644in}}%
\pgfpathlineto{\pgfqpoint{0.685081in}{2.895046in}}%
\pgfpathquadraticcurveto{\pgfqpoint{0.690570in}{2.892058in}}{\pgfqpoint{0.696429in}{2.890593in}}%
\pgfpathquadraticcurveto{\pgfqpoint{0.702288in}{2.889128in}}{\pgfqpoint{0.708812in}{2.889128in}}%
\pgfpathquadraticcurveto{\pgfqpoint{0.719378in}{2.889128in}}{\pgfqpoint{0.725531in}{2.894675in}}%
\pgfpathquadraticcurveto{\pgfqpoint{0.731702in}{2.900222in}}{\pgfqpoint{0.731702in}{2.909753in}}%
\pgfpathquadraticcurveto{\pgfqpoint{0.731702in}{2.919265in}}{\pgfqpoint{0.725531in}{2.924812in}}%
\pgfpathquadraticcurveto{\pgfqpoint{0.719378in}{2.930378in}}{\pgfqpoint{0.708812in}{2.930378in}}%
\pgfpathquadraticcurveto{\pgfqpoint{0.703870in}{2.930378in}}{\pgfqpoint{0.698949in}{2.929284in}}%
\pgfpathquadraticcurveto{\pgfqpoint{0.694046in}{2.928190in}}{\pgfqpoint{0.688929in}{2.925866in}}%
\pgfpathlineto{\pgfqpoint{0.688929in}{2.971647in}}%
\pgfpathlineto{\pgfqpoint{0.688929in}{2.971647in}}%
\pgfpathclose%
\pgfpathmoveto{\pgfqpoint{0.768324in}{2.896022in}}%
\pgfpathlineto{\pgfqpoint{0.781214in}{2.896022in}}%
\pgfpathlineto{\pgfqpoint{0.781214in}{2.880515in}}%
\pgfpathlineto{\pgfqpoint{0.768324in}{2.880515in}}%
\pgfpathlineto{\pgfqpoint{0.768324in}{2.896022in}}%
\pgfpathlineto{\pgfqpoint{0.768324in}{2.896022in}}%
\pgfpathclose%
\pgfpathmoveto{\pgfqpoint{0.834417in}{2.963522in}}%
\pgfpathquadraticcurveto{\pgfqpoint{0.824906in}{2.963522in}}{\pgfqpoint{0.820101in}{2.954147in}}%
\pgfpathquadraticcurveto{\pgfqpoint{0.815316in}{2.944792in}}{\pgfqpoint{0.815316in}{2.925983in}}%
\pgfpathquadraticcurveto{\pgfqpoint{0.815316in}{2.907253in}}{\pgfqpoint{0.820101in}{2.897878in}}%
\pgfpathquadraticcurveto{\pgfqpoint{0.824906in}{2.888503in}}{\pgfqpoint{0.834417in}{2.888503in}}%
\pgfpathquadraticcurveto{\pgfqpoint{0.844007in}{2.888503in}}{\pgfqpoint{0.848792in}{2.897878in}}%
\pgfpathquadraticcurveto{\pgfqpoint{0.853597in}{2.907253in}}{\pgfqpoint{0.853597in}{2.925983in}}%
\pgfpathquadraticcurveto{\pgfqpoint{0.853597in}{2.944792in}}{\pgfqpoint{0.848792in}{2.954147in}}%
\pgfpathquadraticcurveto{\pgfqpoint{0.844007in}{2.963522in}}{\pgfqpoint{0.834417in}{2.963522in}}%
\pgfpathlineto{\pgfqpoint{0.834417in}{2.963522in}}%
\pgfpathclose%
\pgfpathmoveto{\pgfqpoint{0.834417in}{2.973288in}}%
\pgfpathquadraticcurveto{\pgfqpoint{0.849749in}{2.973288in}}{\pgfqpoint{0.857835in}{2.961159in}}%
\pgfpathquadraticcurveto{\pgfqpoint{0.865921in}{2.949050in}}{\pgfqpoint{0.865921in}{2.925983in}}%
\pgfpathquadraticcurveto{\pgfqpoint{0.865921in}{2.902976in}}{\pgfqpoint{0.857835in}{2.890847in}}%
\pgfpathquadraticcurveto{\pgfqpoint{0.849749in}{2.878737in}}{\pgfqpoint{0.834417in}{2.878737in}}%
\pgfpathquadraticcurveto{\pgfqpoint{0.819105in}{2.878737in}}{\pgfqpoint{0.811019in}{2.890847in}}%
\pgfpathquadraticcurveto{\pgfqpoint{0.802933in}{2.902976in}}{\pgfqpoint{0.802933in}{2.925983in}}%
\pgfpathquadraticcurveto{\pgfqpoint{0.802933in}{2.949050in}}{\pgfqpoint{0.811019in}{2.961159in}}%
\pgfpathquadraticcurveto{\pgfqpoint{0.819105in}{2.973288in}}{\pgfqpoint{0.834417in}{2.973288in}}%
\pgfpathlineto{\pgfqpoint{0.834417in}{2.973288in}}%
\pgfpathclose%
\pgfusepath{fill}%
\end{pgfscope}%
\begin{pgfscope}%
\pgfpathrectangle{\pgfqpoint{0.977909in}{2.928000in}}{\pgfqpoint{5.664000in}{5.664000in}}%
\pgfusepath{clip}%
\pgfsetrectcap%
\pgfsetroundjoin%
\pgfsetlinewidth{0.602250pt}%
\definecolor{currentstroke}{rgb}{0.898039,0.913725,0.929412}%
\pgfsetstrokecolor{currentstroke}%
\pgfsetdash{}{0pt}%
\pgfpathmoveto{\pgfqpoint{0.977909in}{4.344000in}}%
\pgfpathlineto{\pgfqpoint{6.641909in}{4.344000in}}%
\pgfusepath{stroke}%
\end{pgfscope}%
\begin{pgfscope}%
\pgfsetbuttcap%
\pgfsetroundjoin%
\definecolor{currentfill}{rgb}{0.090196,0.168627,0.227451}%
\pgfsetfillcolor{currentfill}%
\pgfsetlinewidth{0.803000pt}%
\definecolor{currentstroke}{rgb}{0.090196,0.168627,0.227451}%
\pgfsetstrokecolor{currentstroke}%
\pgfsetdash{}{0pt}%
\pgfsys@defobject{currentmarker}{\pgfqpoint{-0.048611in}{0.000000in}}{\pgfqpoint{-0.000000in}{0.000000in}}{%
\pgfpathmoveto{\pgfqpoint{-0.000000in}{0.000000in}}%
\pgfpathlineto{\pgfqpoint{-0.048611in}{0.000000in}}%
\pgfusepath{stroke,fill}%
}%
\begin{pgfscope}%
\pgfsys@transformshift{0.977909in}{4.344000in}%
\pgfsys@useobject{currentmarker}{}%
\end{pgfscope}%
\end{pgfscope}%
\begin{pgfscope}%
\pgfsetbuttcap%
\pgfsetroundjoin%
\definecolor{currentfill}{rgb}{0.090196,0.168627,0.227451}%
\pgfsetfillcolor{currentfill}%
\pgfsetlinewidth{0.000000pt}%
\definecolor{currentstroke}{rgb}{0.090196,0.168627,0.227451}%
\pgfsetstrokecolor{currentstroke}%
\pgfsetdash{}{0pt}%
\pgfpathmoveto{\pgfqpoint{0.583929in}{4.340890in}}%
\pgfpathlineto{\pgfqpoint{0.662171in}{4.340890in}}%
\pgfpathlineto{\pgfqpoint{0.662171in}{4.330519in}}%
\pgfpathlineto{\pgfqpoint{0.583929in}{4.330519in}}%
\pgfpathlineto{\pgfqpoint{0.583929in}{4.340890in}}%
\pgfpathlineto{\pgfqpoint{0.583929in}{4.340890in}}%
\pgfpathclose%
\pgfpathmoveto{\pgfqpoint{0.699417in}{4.306886in}}%
\pgfpathlineto{\pgfqpoint{0.742445in}{4.306886in}}%
\pgfpathlineto{\pgfqpoint{0.742445in}{4.296515in}}%
\pgfpathlineto{\pgfqpoint{0.684593in}{4.296515in}}%
\pgfpathlineto{\pgfqpoint{0.684593in}{4.306886in}}%
\pgfpathquadraticcurveto{\pgfqpoint{0.691605in}{4.314151in}}{\pgfqpoint{0.703714in}{4.326378in}}%
\pgfpathquadraticcurveto{\pgfqpoint{0.715843in}{4.338624in}}{\pgfqpoint{0.718949in}{4.342179in}}%
\pgfpathquadraticcurveto{\pgfqpoint{0.724867in}{4.348819in}}{\pgfqpoint{0.727210in}{4.353429in}}%
\pgfpathquadraticcurveto{\pgfqpoint{0.729574in}{4.358038in}}{\pgfqpoint{0.729574in}{4.362491in}}%
\pgfpathquadraticcurveto{\pgfqpoint{0.729574in}{4.369757in}}{\pgfqpoint{0.724476in}{4.374327in}}%
\pgfpathquadraticcurveto{\pgfqpoint{0.719378in}{4.378917in}}{\pgfqpoint{0.711195in}{4.378917in}}%
\pgfpathquadraticcurveto{\pgfqpoint{0.705394in}{4.378917in}}{\pgfqpoint{0.698949in}{4.376905in}}%
\pgfpathquadraticcurveto{\pgfqpoint{0.692523in}{4.374894in}}{\pgfqpoint{0.685199in}{4.370792in}}%
\pgfpathlineto{\pgfqpoint{0.685199in}{4.383253in}}%
\pgfpathquadraticcurveto{\pgfqpoint{0.692640in}{4.386241in}}{\pgfqpoint{0.699105in}{4.387765in}}%
\pgfpathquadraticcurveto{\pgfqpoint{0.705589in}{4.389288in}}{\pgfqpoint{0.710960in}{4.389288in}}%
\pgfpathquadraticcurveto{\pgfqpoint{0.725120in}{4.389288in}}{\pgfqpoint{0.733538in}{4.382198in}}%
\pgfpathquadraticcurveto{\pgfqpoint{0.741956in}{4.375128in}}{\pgfqpoint{0.741956in}{4.363292in}}%
\pgfpathquadraticcurveto{\pgfqpoint{0.741956in}{4.357667in}}{\pgfqpoint{0.739847in}{4.352628in}}%
\pgfpathquadraticcurveto{\pgfqpoint{0.737757in}{4.347608in}}{\pgfqpoint{0.732191in}{4.340772in}}%
\pgfpathquadraticcurveto{\pgfqpoint{0.730667in}{4.338995in}}{\pgfqpoint{0.722484in}{4.330538in}}%
\pgfpathquadraticcurveto{\pgfqpoint{0.714320in}{4.322081in}}{\pgfqpoint{0.699417in}{4.306886in}}%
\pgfpathlineto{\pgfqpoint{0.699417in}{4.306886in}}%
\pgfpathclose%
\pgfpathmoveto{\pgfqpoint{0.768324in}{4.312022in}}%
\pgfpathlineto{\pgfqpoint{0.781214in}{4.312022in}}%
\pgfpathlineto{\pgfqpoint{0.781214in}{4.296515in}}%
\pgfpathlineto{\pgfqpoint{0.768324in}{4.296515in}}%
\pgfpathlineto{\pgfqpoint{0.768324in}{4.312022in}}%
\pgfpathlineto{\pgfqpoint{0.768324in}{4.312022in}}%
\pgfpathclose%
\pgfpathmoveto{\pgfqpoint{0.808187in}{4.387647in}}%
\pgfpathlineto{\pgfqpoint{0.856585in}{4.387647in}}%
\pgfpathlineto{\pgfqpoint{0.856585in}{4.377257in}}%
\pgfpathlineto{\pgfqpoint{0.819476in}{4.377257in}}%
\pgfpathlineto{\pgfqpoint{0.819476in}{4.354933in}}%
\pgfpathquadraticcurveto{\pgfqpoint{0.822152in}{4.355851in}}{\pgfqpoint{0.824827in}{4.356300in}}%
\pgfpathquadraticcurveto{\pgfqpoint{0.827523in}{4.356749in}}{\pgfqpoint{0.830218in}{4.356749in}}%
\pgfpathquadraticcurveto{\pgfqpoint{0.845472in}{4.356749in}}{\pgfqpoint{0.854378in}{4.348390in}}%
\pgfpathquadraticcurveto{\pgfqpoint{0.863304in}{4.340030in}}{\pgfqpoint{0.863304in}{4.325753in}}%
\pgfpathquadraticcurveto{\pgfqpoint{0.863304in}{4.311046in}}{\pgfqpoint{0.854144in}{4.302882in}}%
\pgfpathquadraticcurveto{\pgfqpoint{0.844984in}{4.294737in}}{\pgfqpoint{0.828324in}{4.294737in}}%
\pgfpathquadraticcurveto{\pgfqpoint{0.822581in}{4.294737in}}{\pgfqpoint{0.816624in}{4.295714in}}%
\pgfpathquadraticcurveto{\pgfqpoint{0.810687in}{4.296690in}}{\pgfqpoint{0.804339in}{4.298644in}}%
\pgfpathlineto{\pgfqpoint{0.804339in}{4.311046in}}%
\pgfpathquadraticcurveto{\pgfqpoint{0.809827in}{4.308058in}}{\pgfqpoint{0.815687in}{4.306593in}}%
\pgfpathquadraticcurveto{\pgfqpoint{0.821546in}{4.305128in}}{\pgfqpoint{0.828070in}{4.305128in}}%
\pgfpathquadraticcurveto{\pgfqpoint{0.838636in}{4.305128in}}{\pgfqpoint{0.844788in}{4.310675in}}%
\pgfpathquadraticcurveto{\pgfqpoint{0.850960in}{4.316222in}}{\pgfqpoint{0.850960in}{4.325753in}}%
\pgfpathquadraticcurveto{\pgfqpoint{0.850960in}{4.335265in}}{\pgfqpoint{0.844788in}{4.340812in}}%
\pgfpathquadraticcurveto{\pgfqpoint{0.838636in}{4.346378in}}{\pgfqpoint{0.828070in}{4.346378in}}%
\pgfpathquadraticcurveto{\pgfqpoint{0.823128in}{4.346378in}}{\pgfqpoint{0.818206in}{4.345284in}}%
\pgfpathquadraticcurveto{\pgfqpoint{0.813304in}{4.344190in}}{\pgfqpoint{0.808187in}{4.341866in}}%
\pgfpathlineto{\pgfqpoint{0.808187in}{4.387647in}}%
\pgfpathlineto{\pgfqpoint{0.808187in}{4.387647in}}%
\pgfpathclose%
\pgfusepath{fill}%
\end{pgfscope}%
\begin{pgfscope}%
\pgfpathrectangle{\pgfqpoint{0.977909in}{2.928000in}}{\pgfqpoint{5.664000in}{5.664000in}}%
\pgfusepath{clip}%
\pgfsetrectcap%
\pgfsetroundjoin%
\pgfsetlinewidth{0.602250pt}%
\definecolor{currentstroke}{rgb}{0.898039,0.913725,0.929412}%
\pgfsetstrokecolor{currentstroke}%
\pgfsetdash{}{0pt}%
\pgfpathmoveto{\pgfqpoint{0.977909in}{5.760000in}}%
\pgfpathlineto{\pgfqpoint{6.641909in}{5.760000in}}%
\pgfusepath{stroke}%
\end{pgfscope}%
\begin{pgfscope}%
\pgfsetbuttcap%
\pgfsetroundjoin%
\definecolor{currentfill}{rgb}{0.090196,0.168627,0.227451}%
\pgfsetfillcolor{currentfill}%
\pgfsetlinewidth{0.803000pt}%
\definecolor{currentstroke}{rgb}{0.090196,0.168627,0.227451}%
\pgfsetstrokecolor{currentstroke}%
\pgfsetdash{}{0pt}%
\pgfsys@defobject{currentmarker}{\pgfqpoint{-0.048611in}{0.000000in}}{\pgfqpoint{-0.000000in}{0.000000in}}{%
\pgfpathmoveto{\pgfqpoint{-0.000000in}{0.000000in}}%
\pgfpathlineto{\pgfqpoint{-0.048611in}{0.000000in}}%
\pgfusepath{stroke,fill}%
}%
\begin{pgfscope}%
\pgfsys@transformshift{0.977909in}{5.760000in}%
\pgfsys@useobject{currentmarker}{}%
\end{pgfscope}%
\end{pgfscope}%
\begin{pgfscope}%
\pgfsetbuttcap%
\pgfsetroundjoin%
\definecolor{currentfill}{rgb}{0.090196,0.168627,0.227451}%
\pgfsetfillcolor{currentfill}%
\pgfsetlinewidth{0.000000pt}%
\definecolor{currentstroke}{rgb}{0.090196,0.168627,0.227451}%
\pgfsetstrokecolor{currentstroke}%
\pgfsetdash{}{0pt}%
\pgfpathmoveto{\pgfqpoint{0.720413in}{5.795522in}}%
\pgfpathquadraticcurveto{\pgfqpoint{0.710902in}{5.795522in}}{\pgfqpoint{0.706097in}{5.786147in}}%
\pgfpathquadraticcurveto{\pgfqpoint{0.701312in}{5.776792in}}{\pgfqpoint{0.701312in}{5.757983in}}%
\pgfpathquadraticcurveto{\pgfqpoint{0.701312in}{5.739253in}}{\pgfqpoint{0.706097in}{5.729878in}}%
\pgfpathquadraticcurveto{\pgfqpoint{0.710902in}{5.720503in}}{\pgfqpoint{0.720413in}{5.720503in}}%
\pgfpathquadraticcurveto{\pgfqpoint{0.730003in}{5.720503in}}{\pgfqpoint{0.734788in}{5.729878in}}%
\pgfpathquadraticcurveto{\pgfqpoint{0.739593in}{5.739253in}}{\pgfqpoint{0.739593in}{5.757983in}}%
\pgfpathquadraticcurveto{\pgfqpoint{0.739593in}{5.776792in}}{\pgfqpoint{0.734788in}{5.786147in}}%
\pgfpathquadraticcurveto{\pgfqpoint{0.730003in}{5.795522in}}{\pgfqpoint{0.720413in}{5.795522in}}%
\pgfpathlineto{\pgfqpoint{0.720413in}{5.795522in}}%
\pgfpathclose%
\pgfpathmoveto{\pgfqpoint{0.720413in}{5.805288in}}%
\pgfpathquadraticcurveto{\pgfqpoint{0.735745in}{5.805288in}}{\pgfqpoint{0.743831in}{5.793159in}}%
\pgfpathquadraticcurveto{\pgfqpoint{0.751917in}{5.781050in}}{\pgfqpoint{0.751917in}{5.757983in}}%
\pgfpathquadraticcurveto{\pgfqpoint{0.751917in}{5.734976in}}{\pgfqpoint{0.743831in}{5.722847in}}%
\pgfpathquadraticcurveto{\pgfqpoint{0.735745in}{5.710737in}}{\pgfqpoint{0.720413in}{5.710737in}}%
\pgfpathquadraticcurveto{\pgfqpoint{0.705101in}{5.710737in}}{\pgfqpoint{0.697015in}{5.722847in}}%
\pgfpathquadraticcurveto{\pgfqpoint{0.688929in}{5.734976in}}{\pgfqpoint{0.688929in}{5.757983in}}%
\pgfpathquadraticcurveto{\pgfqpoint{0.688929in}{5.781050in}}{\pgfqpoint{0.697015in}{5.793159in}}%
\pgfpathquadraticcurveto{\pgfqpoint{0.705101in}{5.805288in}}{\pgfqpoint{0.720413in}{5.805288in}}%
\pgfpathlineto{\pgfqpoint{0.720413in}{5.805288in}}%
\pgfpathclose%
\pgfpathmoveto{\pgfqpoint{0.773577in}{5.728022in}}%
\pgfpathlineto{\pgfqpoint{0.786468in}{5.728022in}}%
\pgfpathlineto{\pgfqpoint{0.786468in}{5.712515in}}%
\pgfpathlineto{\pgfqpoint{0.773577in}{5.712515in}}%
\pgfpathlineto{\pgfqpoint{0.773577in}{5.728022in}}%
\pgfpathlineto{\pgfqpoint{0.773577in}{5.728022in}}%
\pgfpathclose%
\pgfpathmoveto{\pgfqpoint{0.839671in}{5.795522in}}%
\pgfpathquadraticcurveto{\pgfqpoint{0.830160in}{5.795522in}}{\pgfqpoint{0.825355in}{5.786147in}}%
\pgfpathquadraticcurveto{\pgfqpoint{0.820570in}{5.776792in}}{\pgfqpoint{0.820570in}{5.757983in}}%
\pgfpathquadraticcurveto{\pgfqpoint{0.820570in}{5.739253in}}{\pgfqpoint{0.825355in}{5.729878in}}%
\pgfpathquadraticcurveto{\pgfqpoint{0.830160in}{5.720503in}}{\pgfqpoint{0.839671in}{5.720503in}}%
\pgfpathquadraticcurveto{\pgfqpoint{0.849261in}{5.720503in}}{\pgfqpoint{0.854046in}{5.729878in}}%
\pgfpathquadraticcurveto{\pgfqpoint{0.858851in}{5.739253in}}{\pgfqpoint{0.858851in}{5.757983in}}%
\pgfpathquadraticcurveto{\pgfqpoint{0.858851in}{5.776792in}}{\pgfqpoint{0.854046in}{5.786147in}}%
\pgfpathquadraticcurveto{\pgfqpoint{0.849261in}{5.795522in}}{\pgfqpoint{0.839671in}{5.795522in}}%
\pgfpathlineto{\pgfqpoint{0.839671in}{5.795522in}}%
\pgfpathclose%
\pgfpathmoveto{\pgfqpoint{0.839671in}{5.805288in}}%
\pgfpathquadraticcurveto{\pgfqpoint{0.855003in}{5.805288in}}{\pgfqpoint{0.863089in}{5.793159in}}%
\pgfpathquadraticcurveto{\pgfqpoint{0.871175in}{5.781050in}}{\pgfqpoint{0.871175in}{5.757983in}}%
\pgfpathquadraticcurveto{\pgfqpoint{0.871175in}{5.734976in}}{\pgfqpoint{0.863089in}{5.722847in}}%
\pgfpathquadraticcurveto{\pgfqpoint{0.855003in}{5.710737in}}{\pgfqpoint{0.839671in}{5.710737in}}%
\pgfpathquadraticcurveto{\pgfqpoint{0.824359in}{5.710737in}}{\pgfqpoint{0.816273in}{5.722847in}}%
\pgfpathquadraticcurveto{\pgfqpoint{0.808187in}{5.734976in}}{\pgfqpoint{0.808187in}{5.757983in}}%
\pgfpathquadraticcurveto{\pgfqpoint{0.808187in}{5.781050in}}{\pgfqpoint{0.816273in}{5.793159in}}%
\pgfpathquadraticcurveto{\pgfqpoint{0.824359in}{5.805288in}}{\pgfqpoint{0.839671in}{5.805288in}}%
\pgfpathlineto{\pgfqpoint{0.839671in}{5.805288in}}%
\pgfpathclose%
\pgfusepath{fill}%
\end{pgfscope}%
\begin{pgfscope}%
\pgfpathrectangle{\pgfqpoint{0.977909in}{2.928000in}}{\pgfqpoint{5.664000in}{5.664000in}}%
\pgfusepath{clip}%
\pgfsetrectcap%
\pgfsetroundjoin%
\pgfsetlinewidth{0.602250pt}%
\definecolor{currentstroke}{rgb}{0.898039,0.913725,0.929412}%
\pgfsetstrokecolor{currentstroke}%
\pgfsetdash{}{0pt}%
\pgfpathmoveto{\pgfqpoint{0.977909in}{7.176000in}}%
\pgfpathlineto{\pgfqpoint{6.641909in}{7.176000in}}%
\pgfusepath{stroke}%
\end{pgfscope}%
\begin{pgfscope}%
\pgfsetbuttcap%
\pgfsetroundjoin%
\definecolor{currentfill}{rgb}{0.090196,0.168627,0.227451}%
\pgfsetfillcolor{currentfill}%
\pgfsetlinewidth{0.803000pt}%
\definecolor{currentstroke}{rgb}{0.090196,0.168627,0.227451}%
\pgfsetstrokecolor{currentstroke}%
\pgfsetdash{}{0pt}%
\pgfsys@defobject{currentmarker}{\pgfqpoint{-0.048611in}{0.000000in}}{\pgfqpoint{-0.000000in}{0.000000in}}{%
\pgfpathmoveto{\pgfqpoint{-0.000000in}{0.000000in}}%
\pgfpathlineto{\pgfqpoint{-0.048611in}{0.000000in}}%
\pgfusepath{stroke,fill}%
}%
\begin{pgfscope}%
\pgfsys@transformshift{0.977909in}{7.176000in}%
\pgfsys@useobject{currentmarker}{}%
\end{pgfscope}%
\end{pgfscope}%
\begin{pgfscope}%
\pgfsetbuttcap%
\pgfsetroundjoin%
\definecolor{currentfill}{rgb}{0.090196,0.168627,0.227451}%
\pgfsetfillcolor{currentfill}%
\pgfsetlinewidth{0.000000pt}%
\definecolor{currentstroke}{rgb}{0.090196,0.168627,0.227451}%
\pgfsetstrokecolor{currentstroke}%
\pgfsetdash{}{0pt}%
\pgfpathmoveto{\pgfqpoint{0.704671in}{7.138886in}}%
\pgfpathlineto{\pgfqpoint{0.747699in}{7.138886in}}%
\pgfpathlineto{\pgfqpoint{0.747699in}{7.128515in}}%
\pgfpathlineto{\pgfqpoint{0.689847in}{7.128515in}}%
\pgfpathlineto{\pgfqpoint{0.689847in}{7.138886in}}%
\pgfpathquadraticcurveto{\pgfqpoint{0.696859in}{7.146151in}}{\pgfqpoint{0.708968in}{7.158378in}}%
\pgfpathquadraticcurveto{\pgfqpoint{0.721097in}{7.170624in}}{\pgfqpoint{0.724202in}{7.174179in}}%
\pgfpathquadraticcurveto{\pgfqpoint{0.730120in}{7.180819in}}{\pgfqpoint{0.732464in}{7.185429in}}%
\pgfpathquadraticcurveto{\pgfqpoint{0.734827in}{7.190038in}}{\pgfqpoint{0.734827in}{7.194491in}}%
\pgfpathquadraticcurveto{\pgfqpoint{0.734827in}{7.201757in}}{\pgfqpoint{0.729730in}{7.206327in}}%
\pgfpathquadraticcurveto{\pgfqpoint{0.724632in}{7.210917in}}{\pgfqpoint{0.716449in}{7.210917in}}%
\pgfpathquadraticcurveto{\pgfqpoint{0.710648in}{7.210917in}}{\pgfqpoint{0.704202in}{7.208905in}}%
\pgfpathquadraticcurveto{\pgfqpoint{0.697777in}{7.206894in}}{\pgfqpoint{0.690452in}{7.202792in}}%
\pgfpathlineto{\pgfqpoint{0.690452in}{7.215253in}}%
\pgfpathquadraticcurveto{\pgfqpoint{0.697894in}{7.218241in}}{\pgfqpoint{0.704359in}{7.219765in}}%
\pgfpathquadraticcurveto{\pgfqpoint{0.710843in}{7.221288in}}{\pgfqpoint{0.716214in}{7.221288in}}%
\pgfpathquadraticcurveto{\pgfqpoint{0.730374in}{7.221288in}}{\pgfqpoint{0.738792in}{7.214198in}}%
\pgfpathquadraticcurveto{\pgfqpoint{0.747210in}{7.207128in}}{\pgfqpoint{0.747210in}{7.195292in}}%
\pgfpathquadraticcurveto{\pgfqpoint{0.747210in}{7.189667in}}{\pgfqpoint{0.745101in}{7.184628in}}%
\pgfpathquadraticcurveto{\pgfqpoint{0.743011in}{7.179608in}}{\pgfqpoint{0.737445in}{7.172772in}}%
\pgfpathquadraticcurveto{\pgfqpoint{0.735921in}{7.170995in}}{\pgfqpoint{0.727738in}{7.162538in}}%
\pgfpathquadraticcurveto{\pgfqpoint{0.719574in}{7.154081in}}{\pgfqpoint{0.704671in}{7.138886in}}%
\pgfpathlineto{\pgfqpoint{0.704671in}{7.138886in}}%
\pgfpathclose%
\pgfpathmoveto{\pgfqpoint{0.773577in}{7.144022in}}%
\pgfpathlineto{\pgfqpoint{0.786468in}{7.144022in}}%
\pgfpathlineto{\pgfqpoint{0.786468in}{7.128515in}}%
\pgfpathlineto{\pgfqpoint{0.773577in}{7.128515in}}%
\pgfpathlineto{\pgfqpoint{0.773577in}{7.144022in}}%
\pgfpathlineto{\pgfqpoint{0.773577in}{7.144022in}}%
\pgfpathclose%
\pgfpathmoveto{\pgfqpoint{0.813441in}{7.219647in}}%
\pgfpathlineto{\pgfqpoint{0.861839in}{7.219647in}}%
\pgfpathlineto{\pgfqpoint{0.861839in}{7.209257in}}%
\pgfpathlineto{\pgfqpoint{0.824730in}{7.209257in}}%
\pgfpathlineto{\pgfqpoint{0.824730in}{7.186933in}}%
\pgfpathquadraticcurveto{\pgfqpoint{0.827406in}{7.187851in}}{\pgfqpoint{0.830081in}{7.188300in}}%
\pgfpathquadraticcurveto{\pgfqpoint{0.832777in}{7.188749in}}{\pgfqpoint{0.835472in}{7.188749in}}%
\pgfpathquadraticcurveto{\pgfqpoint{0.850726in}{7.188749in}}{\pgfqpoint{0.859632in}{7.180390in}}%
\pgfpathquadraticcurveto{\pgfqpoint{0.868558in}{7.172030in}}{\pgfqpoint{0.868558in}{7.157753in}}%
\pgfpathquadraticcurveto{\pgfqpoint{0.868558in}{7.143046in}}{\pgfqpoint{0.859398in}{7.134882in}}%
\pgfpathquadraticcurveto{\pgfqpoint{0.850238in}{7.126737in}}{\pgfqpoint{0.833577in}{7.126737in}}%
\pgfpathquadraticcurveto{\pgfqpoint{0.827835in}{7.126737in}}{\pgfqpoint{0.821878in}{7.127714in}}%
\pgfpathquadraticcurveto{\pgfqpoint{0.815941in}{7.128690in}}{\pgfqpoint{0.809593in}{7.130644in}}%
\pgfpathlineto{\pgfqpoint{0.809593in}{7.143046in}}%
\pgfpathquadraticcurveto{\pgfqpoint{0.815081in}{7.140058in}}{\pgfqpoint{0.820941in}{7.138593in}}%
\pgfpathquadraticcurveto{\pgfqpoint{0.826800in}{7.137128in}}{\pgfqpoint{0.833324in}{7.137128in}}%
\pgfpathquadraticcurveto{\pgfqpoint{0.843890in}{7.137128in}}{\pgfqpoint{0.850042in}{7.142675in}}%
\pgfpathquadraticcurveto{\pgfqpoint{0.856214in}{7.148222in}}{\pgfqpoint{0.856214in}{7.157753in}}%
\pgfpathquadraticcurveto{\pgfqpoint{0.856214in}{7.167265in}}{\pgfqpoint{0.850042in}{7.172812in}}%
\pgfpathquadraticcurveto{\pgfqpoint{0.843890in}{7.178378in}}{\pgfqpoint{0.833324in}{7.178378in}}%
\pgfpathquadraticcurveto{\pgfqpoint{0.828382in}{7.178378in}}{\pgfqpoint{0.823460in}{7.177284in}}%
\pgfpathquadraticcurveto{\pgfqpoint{0.818558in}{7.176190in}}{\pgfqpoint{0.813441in}{7.173866in}}%
\pgfpathlineto{\pgfqpoint{0.813441in}{7.219647in}}%
\pgfpathlineto{\pgfqpoint{0.813441in}{7.219647in}}%
\pgfpathclose%
\pgfusepath{fill}%
\end{pgfscope}%
\begin{pgfscope}%
\pgfpathrectangle{\pgfqpoint{0.977909in}{2.928000in}}{\pgfqpoint{5.664000in}{5.664000in}}%
\pgfusepath{clip}%
\pgfsetrectcap%
\pgfsetroundjoin%
\pgfsetlinewidth{0.602250pt}%
\definecolor{currentstroke}{rgb}{0.898039,0.913725,0.929412}%
\pgfsetstrokecolor{currentstroke}%
\pgfsetdash{}{0pt}%
\pgfpathmoveto{\pgfqpoint{0.977909in}{8.592000in}}%
\pgfpathlineto{\pgfqpoint{6.641909in}{8.592000in}}%
\pgfusepath{stroke}%
\end{pgfscope}%
\begin{pgfscope}%
\pgfsetbuttcap%
\pgfsetroundjoin%
\definecolor{currentfill}{rgb}{0.090196,0.168627,0.227451}%
\pgfsetfillcolor{currentfill}%
\pgfsetlinewidth{0.803000pt}%
\definecolor{currentstroke}{rgb}{0.090196,0.168627,0.227451}%
\pgfsetstrokecolor{currentstroke}%
\pgfsetdash{}{0pt}%
\pgfsys@defobject{currentmarker}{\pgfqpoint{-0.048611in}{0.000000in}}{\pgfqpoint{-0.000000in}{0.000000in}}{%
\pgfpathmoveto{\pgfqpoint{-0.000000in}{0.000000in}}%
\pgfpathlineto{\pgfqpoint{-0.048611in}{0.000000in}}%
\pgfusepath{stroke,fill}%
}%
\begin{pgfscope}%
\pgfsys@transformshift{0.977909in}{8.592000in}%
\pgfsys@useobject{currentmarker}{}%
\end{pgfscope}%
\end{pgfscope}%
\begin{pgfscope}%
\pgfsetbuttcap%
\pgfsetroundjoin%
\definecolor{currentfill}{rgb}{0.090196,0.168627,0.227451}%
\pgfsetfillcolor{currentfill}%
\pgfsetlinewidth{0.000000pt}%
\definecolor{currentstroke}{rgb}{0.090196,0.168627,0.227451}%
\pgfsetstrokecolor{currentstroke}%
\pgfsetdash{}{0pt}%
\pgfpathmoveto{\pgfqpoint{0.694183in}{8.635647in}}%
\pgfpathlineto{\pgfqpoint{0.742581in}{8.635647in}}%
\pgfpathlineto{\pgfqpoint{0.742581in}{8.625257in}}%
\pgfpathlineto{\pgfqpoint{0.705472in}{8.625257in}}%
\pgfpathlineto{\pgfqpoint{0.705472in}{8.602933in}}%
\pgfpathquadraticcurveto{\pgfqpoint{0.708148in}{8.603851in}}{\pgfqpoint{0.710824in}{8.604300in}}%
\pgfpathquadraticcurveto{\pgfqpoint{0.713519in}{8.604749in}}{\pgfqpoint{0.716214in}{8.604749in}}%
\pgfpathquadraticcurveto{\pgfqpoint{0.731468in}{8.604749in}}{\pgfqpoint{0.740374in}{8.596390in}}%
\pgfpathquadraticcurveto{\pgfqpoint{0.749300in}{8.588030in}}{\pgfqpoint{0.749300in}{8.573753in}}%
\pgfpathquadraticcurveto{\pgfqpoint{0.749300in}{8.559046in}}{\pgfqpoint{0.740140in}{8.550882in}}%
\pgfpathquadraticcurveto{\pgfqpoint{0.730980in}{8.542737in}}{\pgfqpoint{0.714320in}{8.542737in}}%
\pgfpathquadraticcurveto{\pgfqpoint{0.708577in}{8.542737in}}{\pgfqpoint{0.702620in}{8.543714in}}%
\pgfpathquadraticcurveto{\pgfqpoint{0.696683in}{8.544690in}}{\pgfqpoint{0.690335in}{8.546644in}}%
\pgfpathlineto{\pgfqpoint{0.690335in}{8.559046in}}%
\pgfpathquadraticcurveto{\pgfqpoint{0.695824in}{8.556058in}}{\pgfqpoint{0.701683in}{8.554593in}}%
\pgfpathquadraticcurveto{\pgfqpoint{0.707542in}{8.553128in}}{\pgfqpoint{0.714066in}{8.553128in}}%
\pgfpathquadraticcurveto{\pgfqpoint{0.724632in}{8.553128in}}{\pgfqpoint{0.730785in}{8.558675in}}%
\pgfpathquadraticcurveto{\pgfqpoint{0.736956in}{8.564222in}}{\pgfqpoint{0.736956in}{8.573753in}}%
\pgfpathquadraticcurveto{\pgfqpoint{0.736956in}{8.583265in}}{\pgfqpoint{0.730785in}{8.588812in}}%
\pgfpathquadraticcurveto{\pgfqpoint{0.724632in}{8.594378in}}{\pgfqpoint{0.714066in}{8.594378in}}%
\pgfpathquadraticcurveto{\pgfqpoint{0.709124in}{8.594378in}}{\pgfqpoint{0.704202in}{8.593284in}}%
\pgfpathquadraticcurveto{\pgfqpoint{0.699300in}{8.592190in}}{\pgfqpoint{0.694183in}{8.589866in}}%
\pgfpathlineto{\pgfqpoint{0.694183in}{8.635647in}}%
\pgfpathlineto{\pgfqpoint{0.694183in}{8.635647in}}%
\pgfpathclose%
\pgfpathmoveto{\pgfqpoint{0.773577in}{8.560022in}}%
\pgfpathlineto{\pgfqpoint{0.786468in}{8.560022in}}%
\pgfpathlineto{\pgfqpoint{0.786468in}{8.544515in}}%
\pgfpathlineto{\pgfqpoint{0.773577in}{8.544515in}}%
\pgfpathlineto{\pgfqpoint{0.773577in}{8.560022in}}%
\pgfpathlineto{\pgfqpoint{0.773577in}{8.560022in}}%
\pgfpathclose%
\pgfpathmoveto{\pgfqpoint{0.839671in}{8.627522in}}%
\pgfpathquadraticcurveto{\pgfqpoint{0.830160in}{8.627522in}}{\pgfqpoint{0.825355in}{8.618147in}}%
\pgfpathquadraticcurveto{\pgfqpoint{0.820570in}{8.608792in}}{\pgfqpoint{0.820570in}{8.589983in}}%
\pgfpathquadraticcurveto{\pgfqpoint{0.820570in}{8.571253in}}{\pgfqpoint{0.825355in}{8.561878in}}%
\pgfpathquadraticcurveto{\pgfqpoint{0.830160in}{8.552503in}}{\pgfqpoint{0.839671in}{8.552503in}}%
\pgfpathquadraticcurveto{\pgfqpoint{0.849261in}{8.552503in}}{\pgfqpoint{0.854046in}{8.561878in}}%
\pgfpathquadraticcurveto{\pgfqpoint{0.858851in}{8.571253in}}{\pgfqpoint{0.858851in}{8.589983in}}%
\pgfpathquadraticcurveto{\pgfqpoint{0.858851in}{8.608792in}}{\pgfqpoint{0.854046in}{8.618147in}}%
\pgfpathquadraticcurveto{\pgfqpoint{0.849261in}{8.627522in}}{\pgfqpoint{0.839671in}{8.627522in}}%
\pgfpathlineto{\pgfqpoint{0.839671in}{8.627522in}}%
\pgfpathclose%
\pgfpathmoveto{\pgfqpoint{0.839671in}{8.637288in}}%
\pgfpathquadraticcurveto{\pgfqpoint{0.855003in}{8.637288in}}{\pgfqpoint{0.863089in}{8.625159in}}%
\pgfpathquadraticcurveto{\pgfqpoint{0.871175in}{8.613050in}}{\pgfqpoint{0.871175in}{8.589983in}}%
\pgfpathquadraticcurveto{\pgfqpoint{0.871175in}{8.566976in}}{\pgfqpoint{0.863089in}{8.554847in}}%
\pgfpathquadraticcurveto{\pgfqpoint{0.855003in}{8.542737in}}{\pgfqpoint{0.839671in}{8.542737in}}%
\pgfpathquadraticcurveto{\pgfqpoint{0.824359in}{8.542737in}}{\pgfqpoint{0.816273in}{8.554847in}}%
\pgfpathquadraticcurveto{\pgfqpoint{0.808187in}{8.566976in}}{\pgfqpoint{0.808187in}{8.589983in}}%
\pgfpathquadraticcurveto{\pgfqpoint{0.808187in}{8.613050in}}{\pgfqpoint{0.816273in}{8.625159in}}%
\pgfpathquadraticcurveto{\pgfqpoint{0.824359in}{8.637288in}}{\pgfqpoint{0.839671in}{8.637288in}}%
\pgfpathlineto{\pgfqpoint{0.839671in}{8.637288in}}%
\pgfpathclose%
\pgfusepath{fill}%
\end{pgfscope}%
\begin{pgfscope}%
\pgfsetbuttcap%
\pgfsetroundjoin%
\definecolor{currentfill}{rgb}{0.090196,0.168627,0.227451}%
\pgfsetfillcolor{currentfill}%
\pgfsetlinewidth{0.000000pt}%
\definecolor{currentstroke}{rgb}{0.090196,0.168627,0.227451}%
\pgfsetstrokecolor{currentstroke}%
\pgfsetdash{}{0pt}%
\pgfpathmoveto{\pgfqpoint{0.320082in}{4.896992in}}%
\pgfpathlineto{\pgfqpoint{0.335957in}{4.896992in}}%
\pgfpathquadraticcurveto{\pgfqpoint{0.328867in}{4.889377in}}{\pgfqpoint{0.325382in}{4.880760in}}%
\pgfpathquadraticcurveto{\pgfqpoint{0.321873in}{4.872142in}}{\pgfqpoint{0.321873in}{4.862450in}}%
\pgfpathquadraticcurveto{\pgfqpoint{0.321873in}{4.843353in}}{\pgfqpoint{0.333546in}{4.833207in}}%
\pgfpathquadraticcurveto{\pgfqpoint{0.345219in}{4.823062in}}{\pgfqpoint{0.367300in}{4.823062in}}%
\pgfpathquadraticcurveto{\pgfqpoint{0.389310in}{4.823062in}}{\pgfqpoint{0.400983in}{4.833207in}}%
\pgfpathquadraticcurveto{\pgfqpoint{0.412656in}{4.843353in}}{\pgfqpoint{0.412656in}{4.862450in}}%
\pgfpathquadraticcurveto{\pgfqpoint{0.412656in}{4.872142in}}{\pgfqpoint{0.409147in}{4.880760in}}%
\pgfpathquadraticcurveto{\pgfqpoint{0.405638in}{4.889377in}}{\pgfqpoint{0.398548in}{4.896992in}}%
\pgfpathlineto{\pgfqpoint{0.414303in}{4.896992in}}%
\pgfpathquadraticcurveto{\pgfqpoint{0.419674in}{4.889091in}}{\pgfqpoint{0.422372in}{4.880234in}}%
\pgfpathquadraticcurveto{\pgfqpoint{0.425046in}{4.871402in}}{\pgfqpoint{0.425046in}{4.861567in}}%
\pgfpathquadraticcurveto{\pgfqpoint{0.425046in}{4.836263in}}{\pgfqpoint{0.409577in}{4.821701in}}%
\pgfpathquadraticcurveto{\pgfqpoint{0.394084in}{4.807164in}}{\pgfqpoint{0.367300in}{4.807164in}}%
\pgfpathquadraticcurveto{\pgfqpoint{0.340445in}{4.807164in}}{\pgfqpoint{0.324976in}{4.821701in}}%
\pgfpathquadraticcurveto{\pgfqpoint{0.309484in}{4.836263in}}{\pgfqpoint{0.309484in}{4.861567in}}%
\pgfpathquadraticcurveto{\pgfqpoint{0.309484in}{4.871545in}}{\pgfqpoint{0.312133in}{4.880378in}}%
\pgfpathquadraticcurveto{\pgfqpoint{0.314783in}{4.889234in}}{\pgfqpoint{0.320082in}{4.896992in}}%
\pgfpathlineto{\pgfqpoint{0.320082in}{4.896992in}}%
\pgfpathclose%
\pgfpathmoveto{\pgfqpoint{0.315141in}{4.987036in}}%
\pgfpathlineto{\pgfqpoint{0.329846in}{4.987036in}}%
\pgfpathquadraticcurveto{\pgfqpoint{0.325740in}{4.978466in}}{\pgfqpoint{0.323735in}{4.970851in}}%
\pgfpathquadraticcurveto{\pgfqpoint{0.321706in}{4.963236in}}{\pgfqpoint{0.321706in}{4.956146in}}%
\pgfpathquadraticcurveto{\pgfqpoint{0.321706in}{4.943852in}}{\pgfqpoint{0.326480in}{4.937168in}}%
\pgfpathquadraticcurveto{\pgfqpoint{0.331254in}{4.930484in}}{\pgfqpoint{0.340063in}{4.930484in}}%
\pgfpathquadraticcurveto{\pgfqpoint{0.347463in}{4.930484in}}{\pgfqpoint{0.351235in}{4.934924in}}%
\pgfpathquadraticcurveto{\pgfqpoint{0.354983in}{4.939364in}}{\pgfqpoint{0.357298in}{4.951753in}}%
\pgfpathlineto{\pgfqpoint{0.359160in}{4.960849in}}%
\pgfpathquadraticcurveto{\pgfqpoint{0.362383in}{4.977702in}}{\pgfqpoint{0.370475in}{4.985723in}}%
\pgfpathquadraticcurveto{\pgfqpoint{0.378568in}{4.993743in}}{\pgfqpoint{0.392127in}{4.993743in}}%
\pgfpathquadraticcurveto{\pgfqpoint{0.408336in}{4.993743in}}{\pgfqpoint{0.416691in}{4.982882in}}%
\pgfpathquadraticcurveto{\pgfqpoint{0.425046in}{4.972044in}}{\pgfqpoint{0.425046in}{4.951085in}}%
\pgfpathquadraticcurveto{\pgfqpoint{0.425046in}{4.943184in}}{\pgfqpoint{0.423255in}{4.934256in}}%
\pgfpathquadraticcurveto{\pgfqpoint{0.421465in}{4.925352in}}{\pgfqpoint{0.417956in}{4.915803in}}%
\pgfpathlineto{\pgfqpoint{0.402439in}{4.915803in}}%
\pgfpathquadraticcurveto{\pgfqpoint{0.407572in}{4.924970in}}{\pgfqpoint{0.410198in}{4.933778in}}%
\pgfpathquadraticcurveto{\pgfqpoint{0.412799in}{4.942587in}}{\pgfqpoint{0.412799in}{4.951085in}}%
\pgfpathquadraticcurveto{\pgfqpoint{0.412799in}{4.963976in}}{\pgfqpoint{0.407739in}{4.970994in}}%
\pgfpathquadraticcurveto{\pgfqpoint{0.402654in}{4.978012in}}{\pgfqpoint{0.393249in}{4.978012in}}%
\pgfpathquadraticcurveto{\pgfqpoint{0.385061in}{4.978012in}}{\pgfqpoint{0.380430in}{4.972975in}}%
\pgfpathquadraticcurveto{\pgfqpoint{0.375799in}{4.967938in}}{\pgfqpoint{0.373483in}{4.956456in}}%
\pgfpathlineto{\pgfqpoint{0.371693in}{4.947266in}}%
\pgfpathquadraticcurveto{\pgfqpoint{0.368351in}{4.930412in}}{\pgfqpoint{0.361189in}{4.922869in}}%
\pgfpathquadraticcurveto{\pgfqpoint{0.354028in}{4.915349in}}{\pgfqpoint{0.341257in}{4.915349in}}%
\pgfpathquadraticcurveto{\pgfqpoint{0.326480in}{4.915349in}}{\pgfqpoint{0.317982in}{4.925757in}}%
\pgfpathquadraticcurveto{\pgfqpoint{0.309484in}{4.936165in}}{\pgfqpoint{0.309484in}{4.954427in}}%
\pgfpathquadraticcurveto{\pgfqpoint{0.309484in}{4.962281in}}{\pgfqpoint{0.310916in}{4.970397in}}%
\pgfpathquadraticcurveto{\pgfqpoint{0.312324in}{4.978537in}}{\pgfqpoint{0.315141in}{4.987036in}}%
\pgfpathlineto{\pgfqpoint{0.315141in}{4.987036in}}%
\pgfpathclose%
\pgfpathmoveto{\pgfqpoint{0.326337in}{5.057337in}}%
\pgfpathlineto{\pgfqpoint{0.381767in}{5.036879in}}%
\pgfpathlineto{\pgfqpoint{0.381767in}{5.077843in}}%
\pgfpathlineto{\pgfqpoint{0.326337in}{5.057337in}}%
\pgfpathlineto{\pgfqpoint{0.326337in}{5.057337in}}%
\pgfpathclose%
\pgfpathmoveto{\pgfqpoint{0.311489in}{5.048815in}}%
\pgfpathlineto{\pgfqpoint{0.311489in}{5.065907in}}%
\pgfpathlineto{\pgfqpoint{0.422873in}{5.108351in}}%
\pgfpathlineto{\pgfqpoint{0.422873in}{5.092691in}}%
\pgfpathlineto{\pgfqpoint{0.394299in}{5.082546in}}%
\pgfpathlineto{\pgfqpoint{0.394299in}{5.032344in}}%
\pgfpathlineto{\pgfqpoint{0.422873in}{5.022198in}}%
\pgfpathlineto{\pgfqpoint{0.422873in}{5.006300in}}%
\pgfpathlineto{\pgfqpoint{0.311489in}{5.048815in}}%
\pgfpathlineto{\pgfqpoint{0.311489in}{5.048815in}}%
\pgfpathclose%
\pgfpathmoveto{\pgfqpoint{0.377661in}{5.244036in}}%
\pgfpathlineto{\pgfqpoint{0.384369in}{5.244036in}}%
\pgfpathlineto{\pgfqpoint{0.384369in}{5.180920in}}%
\pgfpathquadraticcurveto{\pgfqpoint{0.398548in}{5.181827in}}{\pgfqpoint{0.405972in}{5.189466in}}%
\pgfpathquadraticcurveto{\pgfqpoint{0.413396in}{5.197105in}}{\pgfqpoint{0.413396in}{5.210760in}}%
\pgfpathquadraticcurveto{\pgfqpoint{0.413396in}{5.218661in}}{\pgfqpoint{0.411463in}{5.226085in}}%
\pgfpathquadraticcurveto{\pgfqpoint{0.409529in}{5.233509in}}{\pgfqpoint{0.405638in}{5.240838in}}%
\pgfpathlineto{\pgfqpoint{0.418624in}{5.240838in}}%
\pgfpathquadraticcurveto{\pgfqpoint{0.421751in}{5.233437in}}{\pgfqpoint{0.423398in}{5.225679in}}%
\pgfpathquadraticcurveto{\pgfqpoint{0.425046in}{5.217921in}}{\pgfqpoint{0.425046in}{5.209948in}}%
\pgfpathquadraticcurveto{\pgfqpoint{0.425046in}{5.189944in}}{\pgfqpoint{0.413420in}{5.178270in}}%
\pgfpathquadraticcurveto{\pgfqpoint{0.401771in}{5.166597in}}{\pgfqpoint{0.381910in}{5.166597in}}%
\pgfpathquadraticcurveto{\pgfqpoint{0.361404in}{5.166597in}}{\pgfqpoint{0.349373in}{5.177674in}}%
\pgfpathquadraticcurveto{\pgfqpoint{0.337318in}{5.188750in}}{\pgfqpoint{0.337318in}{5.207561in}}%
\pgfpathquadraticcurveto{\pgfqpoint{0.337318in}{5.224414in}}{\pgfqpoint{0.348179in}{5.234225in}}%
\pgfpathquadraticcurveto{\pgfqpoint{0.359017in}{5.244036in}}{\pgfqpoint{0.377661in}{5.244036in}}%
\pgfpathlineto{\pgfqpoint{0.377661in}{5.244036in}}%
\pgfpathclose%
\pgfpathmoveto{\pgfqpoint{0.373626in}{5.230310in}}%
\pgfpathquadraticcurveto{\pgfqpoint{0.362383in}{5.230167in}}{\pgfqpoint{0.355675in}{5.224008in}}%
\pgfpathquadraticcurveto{\pgfqpoint{0.348943in}{5.217849in}}{\pgfqpoint{0.348943in}{5.207704in}}%
\pgfpathquadraticcurveto{\pgfqpoint{0.348943in}{5.196222in}}{\pgfqpoint{0.355436in}{5.189323in}}%
\pgfpathquadraticcurveto{\pgfqpoint{0.361929in}{5.182424in}}{\pgfqpoint{0.373722in}{5.181374in}}%
\pgfpathlineto{\pgfqpoint{0.373626in}{5.230310in}}%
\pgfpathlineto{\pgfqpoint{0.373626in}{5.230310in}}%
\pgfpathclose%
\pgfpathmoveto{\pgfqpoint{0.341782in}{5.319829in}}%
\pgfpathlineto{\pgfqpoint{0.354768in}{5.319829in}}%
\pgfpathquadraticcurveto{\pgfqpoint{0.351784in}{5.314028in}}{\pgfqpoint{0.350304in}{5.307750in}}%
\pgfpathquadraticcurveto{\pgfqpoint{0.348800in}{5.301495in}}{\pgfqpoint{0.348800in}{5.294763in}}%
\pgfpathquadraticcurveto{\pgfqpoint{0.348800in}{5.284546in}}{\pgfqpoint{0.351927in}{5.279438in}}%
\pgfpathquadraticcurveto{\pgfqpoint{0.355054in}{5.274329in}}{\pgfqpoint{0.361332in}{5.274329in}}%
\pgfpathquadraticcurveto{\pgfqpoint{0.366107in}{5.274329in}}{\pgfqpoint{0.368828in}{5.277982in}}%
\pgfpathquadraticcurveto{\pgfqpoint{0.371549in}{5.281634in}}{\pgfqpoint{0.374008in}{5.292687in}}%
\pgfpathlineto{\pgfqpoint{0.375059in}{5.297389in}}%
\pgfpathquadraticcurveto{\pgfqpoint{0.378186in}{5.311999in}}{\pgfqpoint{0.383891in}{5.318158in}}%
\pgfpathquadraticcurveto{\pgfqpoint{0.389596in}{5.324316in}}{\pgfqpoint{0.399813in}{5.324316in}}%
\pgfpathquadraticcurveto{\pgfqpoint{0.411463in}{5.324316in}}{\pgfqpoint{0.418266in}{5.315102in}}%
\pgfpathquadraticcurveto{\pgfqpoint{0.425046in}{5.305888in}}{\pgfqpoint{0.425046in}{5.289774in}}%
\pgfpathquadraticcurveto{\pgfqpoint{0.425046in}{5.283066in}}{\pgfqpoint{0.423733in}{5.275786in}}%
\pgfpathquadraticcurveto{\pgfqpoint{0.422420in}{5.268505in}}{\pgfqpoint{0.419818in}{5.260460in}}%
\pgfpathlineto{\pgfqpoint{0.405638in}{5.260460in}}%
\pgfpathquadraticcurveto{\pgfqpoint{0.409601in}{5.268075in}}{\pgfqpoint{0.411582in}{5.275451in}}%
\pgfpathquadraticcurveto{\pgfqpoint{0.413540in}{5.282828in}}{\pgfqpoint{0.413540in}{5.290085in}}%
\pgfpathquadraticcurveto{\pgfqpoint{0.413540in}{5.299776in}}{\pgfqpoint{0.410221in}{5.304980in}}%
\pgfpathquadraticcurveto{\pgfqpoint{0.406903in}{5.310208in}}{\pgfqpoint{0.400864in}{5.310208in}}%
\pgfpathquadraticcurveto{\pgfqpoint{0.395278in}{5.310208in}}{\pgfqpoint{0.392294in}{5.306437in}}%
\pgfpathquadraticcurveto{\pgfqpoint{0.389310in}{5.302689in}}{\pgfqpoint{0.386541in}{5.289918in}}%
\pgfpathlineto{\pgfqpoint{0.385419in}{5.285143in}}%
\pgfpathquadraticcurveto{\pgfqpoint{0.382745in}{5.272396in}}{\pgfqpoint{0.377183in}{5.266714in}}%
\pgfpathquadraticcurveto{\pgfqpoint{0.371621in}{5.261057in}}{\pgfqpoint{0.361929in}{5.261057in}}%
\pgfpathquadraticcurveto{\pgfqpoint{0.350137in}{5.261057in}}{\pgfqpoint{0.343739in}{5.269412in}}%
\pgfpathquadraticcurveto{\pgfqpoint{0.337318in}{5.277767in}}{\pgfqpoint{0.337318in}{5.293140in}}%
\pgfpathquadraticcurveto{\pgfqpoint{0.337318in}{5.300731in}}{\pgfqpoint{0.338440in}{5.307439in}}%
\pgfpathquadraticcurveto{\pgfqpoint{0.339538in}{5.314171in}}{\pgfqpoint{0.341782in}{5.319829in}}%
\pgfpathlineto{\pgfqpoint{0.341782in}{5.319829in}}%
\pgfpathclose%
\pgfpathmoveto{\pgfqpoint{0.315595in}{5.359742in}}%
\pgfpathlineto{\pgfqpoint{0.339323in}{5.359742in}}%
\pgfpathlineto{\pgfqpoint{0.339323in}{5.388006in}}%
\pgfpathlineto{\pgfqpoint{0.349994in}{5.388006in}}%
\pgfpathlineto{\pgfqpoint{0.349994in}{5.359742in}}%
\pgfpathlineto{\pgfqpoint{0.395349in}{5.359742in}}%
\pgfpathquadraticcurveto{\pgfqpoint{0.405566in}{5.359742in}}{\pgfqpoint{0.408479in}{5.362535in}}%
\pgfpathquadraticcurveto{\pgfqpoint{0.411391in}{5.365328in}}{\pgfqpoint{0.411391in}{5.373921in}}%
\pgfpathlineto{\pgfqpoint{0.411391in}{5.388006in}}%
\pgfpathlineto{\pgfqpoint{0.422873in}{5.388006in}}%
\pgfpathlineto{\pgfqpoint{0.422873in}{5.373921in}}%
\pgfpathquadraticcurveto{\pgfqpoint{0.422873in}{5.358023in}}{\pgfqpoint{0.416953in}{5.351984in}}%
\pgfpathquadraticcurveto{\pgfqpoint{0.411009in}{5.345944in}}{\pgfqpoint{0.395349in}{5.345944in}}%
\pgfpathlineto{\pgfqpoint{0.349994in}{5.345944in}}%
\pgfpathlineto{\pgfqpoint{0.349994in}{5.335870in}}%
\pgfpathlineto{\pgfqpoint{0.339323in}{5.335870in}}%
\pgfpathlineto{\pgfqpoint{0.339323in}{5.345944in}}%
\pgfpathlineto{\pgfqpoint{0.315595in}{5.345944in}}%
\pgfpathlineto{\pgfqpoint{0.315595in}{5.359742in}}%
\pgfpathlineto{\pgfqpoint{0.315595in}{5.359742in}}%
\pgfpathclose%
\pgfpathmoveto{\pgfqpoint{0.339323in}{5.406053in}}%
\pgfpathlineto{\pgfqpoint{0.339323in}{5.419779in}}%
\pgfpathlineto{\pgfqpoint{0.422873in}{5.419779in}}%
\pgfpathlineto{\pgfqpoint{0.422873in}{5.406053in}}%
\pgfpathlineto{\pgfqpoint{0.339323in}{5.406053in}}%
\pgfpathlineto{\pgfqpoint{0.339323in}{5.406053in}}%
\pgfpathclose%
\pgfpathmoveto{\pgfqpoint{0.306786in}{5.406053in}}%
\pgfpathlineto{\pgfqpoint{0.306786in}{5.419779in}}%
\pgfpathlineto{\pgfqpoint{0.324188in}{5.419779in}}%
\pgfpathlineto{\pgfqpoint{0.324188in}{5.406053in}}%
\pgfpathlineto{\pgfqpoint{0.306786in}{5.406053in}}%
\pgfpathlineto{\pgfqpoint{0.306786in}{5.406053in}}%
\pgfpathclose%
\pgfpathmoveto{\pgfqpoint{0.355365in}{5.513546in}}%
\pgfpathquadraticcurveto{\pgfqpoint{0.346102in}{5.518702in}}{\pgfqpoint{0.341710in}{5.525864in}}%
\pgfpathquadraticcurveto{\pgfqpoint{0.337318in}{5.533025in}}{\pgfqpoint{0.337318in}{5.542717in}}%
\pgfpathquadraticcurveto{\pgfqpoint{0.337318in}{5.555775in}}{\pgfqpoint{0.346461in}{5.562865in}}%
\pgfpathquadraticcurveto{\pgfqpoint{0.355579in}{5.569954in}}{\pgfqpoint{0.372433in}{5.569954in}}%
\pgfpathlineto{\pgfqpoint{0.422873in}{5.569954in}}%
\pgfpathlineto{\pgfqpoint{0.422873in}{5.556157in}}%
\pgfpathlineto{\pgfqpoint{0.372886in}{5.556157in}}%
\pgfpathquadraticcurveto{\pgfqpoint{0.360879in}{5.556157in}}{\pgfqpoint{0.355078in}{5.551884in}}%
\pgfpathquadraticcurveto{\pgfqpoint{0.349253in}{5.547635in}}{\pgfqpoint{0.349253in}{5.538921in}}%
\pgfpathquadraticcurveto{\pgfqpoint{0.349253in}{5.528251in}}{\pgfqpoint{0.356343in}{5.522044in}}%
\pgfpathquadraticcurveto{\pgfqpoint{0.363409in}{5.515862in}}{\pgfqpoint{0.375655in}{5.515862in}}%
\pgfpathlineto{\pgfqpoint{0.422873in}{5.515862in}}%
\pgfpathlineto{\pgfqpoint{0.422873in}{5.502064in}}%
\pgfpathlineto{\pgfqpoint{0.372886in}{5.502064in}}%
\pgfpathquadraticcurveto{\pgfqpoint{0.360807in}{5.502064in}}{\pgfqpoint{0.355030in}{5.497815in}}%
\pgfpathquadraticcurveto{\pgfqpoint{0.349253in}{5.493566in}}{\pgfqpoint{0.349253in}{5.484685in}}%
\pgfpathquadraticcurveto{\pgfqpoint{0.349253in}{5.474158in}}{\pgfqpoint{0.356367in}{5.467951in}}%
\pgfpathquadraticcurveto{\pgfqpoint{0.363481in}{5.461769in}}{\pgfqpoint{0.375655in}{5.461769in}}%
\pgfpathlineto{\pgfqpoint{0.422873in}{5.461769in}}%
\pgfpathlineto{\pgfqpoint{0.422873in}{5.447971in}}%
\pgfpathlineto{\pgfqpoint{0.339323in}{5.447971in}}%
\pgfpathlineto{\pgfqpoint{0.339323in}{5.461769in}}%
\pgfpathlineto{\pgfqpoint{0.352309in}{5.461769in}}%
\pgfpathquadraticcurveto{\pgfqpoint{0.344622in}{5.466471in}}{\pgfqpoint{0.340970in}{5.473036in}}%
\pgfpathquadraticcurveto{\pgfqpoint{0.337318in}{5.479601in}}{\pgfqpoint{0.337318in}{5.488624in}}%
\pgfpathquadraticcurveto{\pgfqpoint{0.337318in}{5.497743in}}{\pgfqpoint{0.341949in}{5.504117in}}%
\pgfpathquadraticcurveto{\pgfqpoint{0.346556in}{5.510490in}}{\pgfqpoint{0.355365in}{5.513546in}}%
\pgfpathlineto{\pgfqpoint{0.355365in}{5.513546in}}%
\pgfpathclose%
\pgfpathmoveto{\pgfqpoint{0.380883in}{5.635291in}}%
\pgfpathquadraticcurveto{\pgfqpoint{0.380883in}{5.618652in}}{\pgfqpoint{0.384679in}{5.612231in}}%
\pgfpathquadraticcurveto{\pgfqpoint{0.388474in}{5.605809in}}{\pgfqpoint{0.397665in}{5.605809in}}%
\pgfpathquadraticcurveto{\pgfqpoint{0.404970in}{5.605809in}}{\pgfqpoint{0.409267in}{5.610632in}}%
\pgfpathquadraticcurveto{\pgfqpoint{0.413540in}{5.615454in}}{\pgfqpoint{0.413540in}{5.623713in}}%
\pgfpathquadraticcurveto{\pgfqpoint{0.413540in}{5.635148in}}{\pgfqpoint{0.405447in}{5.642046in}}%
\pgfpathquadraticcurveto{\pgfqpoint{0.397355in}{5.648945in}}{\pgfqpoint{0.383939in}{5.648945in}}%
\pgfpathlineto{\pgfqpoint{0.380883in}{5.648945in}}%
\pgfpathlineto{\pgfqpoint{0.380883in}{5.635291in}}%
\pgfpathlineto{\pgfqpoint{0.380883in}{5.635291in}}%
\pgfpathclose%
\pgfpathmoveto{\pgfqpoint{0.375202in}{5.662671in}}%
\pgfpathlineto{\pgfqpoint{0.422873in}{5.662671in}}%
\pgfpathlineto{\pgfqpoint{0.422873in}{5.648945in}}%
\pgfpathlineto{\pgfqpoint{0.410198in}{5.648945in}}%
\pgfpathquadraticcurveto{\pgfqpoint{0.417789in}{5.644243in}}{\pgfqpoint{0.421417in}{5.637224in}}%
\pgfpathquadraticcurveto{\pgfqpoint{0.425046in}{5.630206in}}{\pgfqpoint{0.425046in}{5.620061in}}%
\pgfpathquadraticcurveto{\pgfqpoint{0.425046in}{5.607242in}}{\pgfqpoint{0.417836in}{5.599651in}}%
\pgfpathquadraticcurveto{\pgfqpoint{0.410627in}{5.592083in}}{\pgfqpoint{0.398548in}{5.592083in}}%
\pgfpathquadraticcurveto{\pgfqpoint{0.384464in}{5.592083in}}{\pgfqpoint{0.377303in}{5.601513in}}%
\pgfpathquadraticcurveto{\pgfqpoint{0.370141in}{5.610966in}}{\pgfqpoint{0.370141in}{5.629681in}}%
\pgfpathlineto{\pgfqpoint{0.370141in}{5.648945in}}%
\pgfpathlineto{\pgfqpoint{0.368780in}{5.648945in}}%
\pgfpathquadraticcurveto{\pgfqpoint{0.359303in}{5.648945in}}{\pgfqpoint{0.354123in}{5.642715in}}%
\pgfpathquadraticcurveto{\pgfqpoint{0.348943in}{5.636484in}}{\pgfqpoint{0.348943in}{5.625217in}}%
\pgfpathquadraticcurveto{\pgfqpoint{0.348943in}{5.618056in}}{\pgfqpoint{0.350662in}{5.611252in}}%
\pgfpathquadraticcurveto{\pgfqpoint{0.352381in}{5.604473in}}{\pgfqpoint{0.355818in}{5.598218in}}%
\pgfpathlineto{\pgfqpoint{0.343119in}{5.598218in}}%
\pgfpathquadraticcurveto{\pgfqpoint{0.340206in}{5.605738in}}{\pgfqpoint{0.338774in}{5.612828in}}%
\pgfpathquadraticcurveto{\pgfqpoint{0.337318in}{5.619918in}}{\pgfqpoint{0.337318in}{5.626625in}}%
\pgfpathquadraticcurveto{\pgfqpoint{0.337318in}{5.644768in}}{\pgfqpoint{0.346723in}{5.653720in}}%
\pgfpathquadraticcurveto{\pgfqpoint{0.356105in}{5.662671in}}{\pgfqpoint{0.375202in}{5.662671in}}%
\pgfpathlineto{\pgfqpoint{0.375202in}{5.662671in}}%
\pgfpathclose%
\pgfpathmoveto{\pgfqpoint{0.315595in}{5.704518in}}%
\pgfpathlineto{\pgfqpoint{0.339323in}{5.704518in}}%
\pgfpathlineto{\pgfqpoint{0.339323in}{5.732782in}}%
\pgfpathlineto{\pgfqpoint{0.349994in}{5.732782in}}%
\pgfpathlineto{\pgfqpoint{0.349994in}{5.704518in}}%
\pgfpathlineto{\pgfqpoint{0.395349in}{5.704518in}}%
\pgfpathquadraticcurveto{\pgfqpoint{0.405566in}{5.704518in}}{\pgfqpoint{0.408479in}{5.707311in}}%
\pgfpathquadraticcurveto{\pgfqpoint{0.411391in}{5.710104in}}{\pgfqpoint{0.411391in}{5.718698in}}%
\pgfpathlineto{\pgfqpoint{0.411391in}{5.732782in}}%
\pgfpathlineto{\pgfqpoint{0.422873in}{5.732782in}}%
\pgfpathlineto{\pgfqpoint{0.422873in}{5.718698in}}%
\pgfpathquadraticcurveto{\pgfqpoint{0.422873in}{5.702799in}}{\pgfqpoint{0.416953in}{5.696760in}}%
\pgfpathquadraticcurveto{\pgfqpoint{0.411009in}{5.690720in}}{\pgfqpoint{0.395349in}{5.690720in}}%
\pgfpathlineto{\pgfqpoint{0.349994in}{5.690720in}}%
\pgfpathlineto{\pgfqpoint{0.349994in}{5.680647in}}%
\pgfpathlineto{\pgfqpoint{0.339323in}{5.680647in}}%
\pgfpathlineto{\pgfqpoint{0.339323in}{5.690720in}}%
\pgfpathlineto{\pgfqpoint{0.315595in}{5.690720in}}%
\pgfpathlineto{\pgfqpoint{0.315595in}{5.704518in}}%
\pgfpathlineto{\pgfqpoint{0.315595in}{5.704518in}}%
\pgfpathclose%
\pgfpathmoveto{\pgfqpoint{0.377661in}{5.822300in}}%
\pgfpathlineto{\pgfqpoint{0.384369in}{5.822300in}}%
\pgfpathlineto{\pgfqpoint{0.384369in}{5.759184in}}%
\pgfpathquadraticcurveto{\pgfqpoint{0.398548in}{5.760091in}}{\pgfqpoint{0.405972in}{5.767730in}}%
\pgfpathquadraticcurveto{\pgfqpoint{0.413396in}{5.775369in}}{\pgfqpoint{0.413396in}{5.789023in}}%
\pgfpathquadraticcurveto{\pgfqpoint{0.413396in}{5.796925in}}{\pgfqpoint{0.411463in}{5.804349in}}%
\pgfpathquadraticcurveto{\pgfqpoint{0.409529in}{5.811773in}}{\pgfqpoint{0.405638in}{5.819102in}}%
\pgfpathlineto{\pgfqpoint{0.418624in}{5.819102in}}%
\pgfpathquadraticcurveto{\pgfqpoint{0.421751in}{5.811701in}}{\pgfqpoint{0.423398in}{5.803943in}}%
\pgfpathquadraticcurveto{\pgfqpoint{0.425046in}{5.796185in}}{\pgfqpoint{0.425046in}{5.788212in}}%
\pgfpathquadraticcurveto{\pgfqpoint{0.425046in}{5.768207in}}{\pgfqpoint{0.413420in}{5.756534in}}%
\pgfpathquadraticcurveto{\pgfqpoint{0.401771in}{5.744861in}}{\pgfqpoint{0.381910in}{5.744861in}}%
\pgfpathquadraticcurveto{\pgfqpoint{0.361404in}{5.744861in}}{\pgfqpoint{0.349373in}{5.755937in}}%
\pgfpathquadraticcurveto{\pgfqpoint{0.337318in}{5.767014in}}{\pgfqpoint{0.337318in}{5.785825in}}%
\pgfpathquadraticcurveto{\pgfqpoint{0.337318in}{5.802678in}}{\pgfqpoint{0.348179in}{5.812489in}}%
\pgfpathquadraticcurveto{\pgfqpoint{0.359017in}{5.822300in}}{\pgfqpoint{0.377661in}{5.822300in}}%
\pgfpathlineto{\pgfqpoint{0.377661in}{5.822300in}}%
\pgfpathclose%
\pgfpathmoveto{\pgfqpoint{0.373626in}{5.808574in}}%
\pgfpathquadraticcurveto{\pgfqpoint{0.362383in}{5.808431in}}{\pgfqpoint{0.355675in}{5.802272in}}%
\pgfpathquadraticcurveto{\pgfqpoint{0.348943in}{5.796113in}}{\pgfqpoint{0.348943in}{5.785968in}}%
\pgfpathquadraticcurveto{\pgfqpoint{0.348943in}{5.774486in}}{\pgfqpoint{0.355436in}{5.767587in}}%
\pgfpathquadraticcurveto{\pgfqpoint{0.361929in}{5.760688in}}{\pgfqpoint{0.373722in}{5.759638in}}%
\pgfpathlineto{\pgfqpoint{0.373626in}{5.808574in}}%
\pgfpathlineto{\pgfqpoint{0.373626in}{5.808574in}}%
\pgfpathclose%
\pgfpathmoveto{\pgfqpoint{0.306953in}{5.926356in}}%
\pgfpathquadraticcurveto{\pgfqpoint{0.324093in}{5.916378in}}{\pgfqpoint{0.340898in}{5.911508in}}%
\pgfpathquadraticcurveto{\pgfqpoint{0.357680in}{5.906662in}}{\pgfqpoint{0.374915in}{5.906662in}}%
\pgfpathquadraticcurveto{\pgfqpoint{0.392127in}{5.906662in}}{\pgfqpoint{0.409028in}{5.911556in}}%
\pgfpathquadraticcurveto{\pgfqpoint{0.425929in}{5.916450in}}{\pgfqpoint{0.443021in}{5.926356in}}%
\pgfpathlineto{\pgfqpoint{0.443021in}{5.914421in}}%
\pgfpathquadraticcurveto{\pgfqpoint{0.425475in}{5.903249in}}{\pgfqpoint{0.408550in}{5.897687in}}%
\pgfpathquadraticcurveto{\pgfqpoint{0.391625in}{5.892125in}}{\pgfqpoint{0.374915in}{5.892125in}}%
\pgfpathquadraticcurveto{\pgfqpoint{0.358277in}{5.892125in}}{\pgfqpoint{0.341424in}{5.897639in}}%
\pgfpathquadraticcurveto{\pgfqpoint{0.324546in}{5.903177in}}{\pgfqpoint{0.306953in}{5.914421in}}%
\pgfpathlineto{\pgfqpoint{0.306953in}{5.926356in}}%
\pgfpathlineto{\pgfqpoint{0.306953in}{5.926356in}}%
\pgfpathclose%
\pgfpathmoveto{\pgfqpoint{0.311489in}{5.938149in}}%
\pgfpathlineto{\pgfqpoint{0.311489in}{6.032370in}}%
\pgfpathlineto{\pgfqpoint{0.324188in}{6.032370in}}%
\pgfpathlineto{\pgfqpoint{0.324188in}{5.992839in}}%
\pgfpathlineto{\pgfqpoint{0.422873in}{5.992839in}}%
\pgfpathlineto{\pgfqpoint{0.422873in}{5.977704in}}%
\pgfpathlineto{\pgfqpoint{0.324188in}{5.977704in}}%
\pgfpathlineto{\pgfqpoint{0.324188in}{5.938149in}}%
\pgfpathlineto{\pgfqpoint{0.311489in}{5.938149in}}%
\pgfpathlineto{\pgfqpoint{0.311489in}{5.938149in}}%
\pgfpathclose%
\pgfpathmoveto{\pgfqpoint{0.311489in}{6.037001in}}%
\pgfpathlineto{\pgfqpoint{0.311489in}{6.052207in}}%
\pgfpathlineto{\pgfqpoint{0.405638in}{6.075625in}}%
\pgfpathlineto{\pgfqpoint{0.311489in}{6.098971in}}%
\pgfpathlineto{\pgfqpoint{0.311489in}{6.115920in}}%
\pgfpathlineto{\pgfqpoint{0.405638in}{6.139338in}}%
\pgfpathlineto{\pgfqpoint{0.311489in}{6.162684in}}%
\pgfpathlineto{\pgfqpoint{0.311489in}{6.177986in}}%
\pgfpathlineto{\pgfqpoint{0.422873in}{6.150009in}}%
\pgfpathlineto{\pgfqpoint{0.422873in}{6.131055in}}%
\pgfpathlineto{\pgfqpoint{0.326194in}{6.107565in}}%
\pgfpathlineto{\pgfqpoint{0.422873in}{6.083837in}}%
\pgfpathlineto{\pgfqpoint{0.422873in}{6.064883in}}%
\pgfpathlineto{\pgfqpoint{0.311489in}{6.037001in}}%
\pgfpathlineto{\pgfqpoint{0.311489in}{6.037001in}}%
\pgfpathclose%
\pgfpathmoveto{\pgfqpoint{0.311489in}{6.197967in}}%
\pgfpathlineto{\pgfqpoint{0.311489in}{6.261966in}}%
\pgfpathlineto{\pgfqpoint{0.324188in}{6.261966in}}%
\pgfpathlineto{\pgfqpoint{0.324188in}{6.213030in}}%
\pgfpathlineto{\pgfqpoint{0.357012in}{6.213030in}}%
\pgfpathlineto{\pgfqpoint{0.357012in}{6.257192in}}%
\pgfpathlineto{\pgfqpoint{0.369688in}{6.257192in}}%
\pgfpathlineto{\pgfqpoint{0.369688in}{6.213030in}}%
\pgfpathlineto{\pgfqpoint{0.422873in}{6.213030in}}%
\pgfpathlineto{\pgfqpoint{0.422873in}{6.197967in}}%
\pgfpathlineto{\pgfqpoint{0.311489in}{6.197967in}}%
\pgfpathlineto{\pgfqpoint{0.311489in}{6.197967in}}%
\pgfpathclose%
\pgfpathmoveto{\pgfqpoint{0.311489in}{6.285838in}}%
\pgfpathlineto{\pgfqpoint{0.311489in}{6.356259in}}%
\pgfpathlineto{\pgfqpoint{0.324188in}{6.356259in}}%
\pgfpathlineto{\pgfqpoint{0.324188in}{6.300901in}}%
\pgfpathlineto{\pgfqpoint{0.357155in}{6.300901in}}%
\pgfpathlineto{\pgfqpoint{0.357155in}{6.353943in}}%
\pgfpathlineto{\pgfqpoint{0.369831in}{6.353943in}}%
\pgfpathlineto{\pgfqpoint{0.369831in}{6.300901in}}%
\pgfpathlineto{\pgfqpoint{0.410198in}{6.300901in}}%
\pgfpathlineto{\pgfqpoint{0.410198in}{6.357595in}}%
\pgfpathlineto{\pgfqpoint{0.422873in}{6.357595in}}%
\pgfpathlineto{\pgfqpoint{0.422873in}{6.285838in}}%
\pgfpathlineto{\pgfqpoint{0.311489in}{6.285838in}}%
\pgfpathlineto{\pgfqpoint{0.311489in}{6.285838in}}%
\pgfpathclose%
\pgfpathmoveto{\pgfqpoint{0.315141in}{6.497697in}}%
\pgfpathlineto{\pgfqpoint{0.329846in}{6.497697in}}%
\pgfpathquadraticcurveto{\pgfqpoint{0.325740in}{6.489128in}}{\pgfqpoint{0.323735in}{6.481513in}}%
\pgfpathquadraticcurveto{\pgfqpoint{0.321706in}{6.473898in}}{\pgfqpoint{0.321706in}{6.466808in}}%
\pgfpathquadraticcurveto{\pgfqpoint{0.321706in}{6.454514in}}{\pgfqpoint{0.326480in}{6.447830in}}%
\pgfpathquadraticcurveto{\pgfqpoint{0.331254in}{6.441146in}}{\pgfqpoint{0.340063in}{6.441146in}}%
\pgfpathquadraticcurveto{\pgfqpoint{0.347463in}{6.441146in}}{\pgfqpoint{0.351235in}{6.445586in}}%
\pgfpathquadraticcurveto{\pgfqpoint{0.354983in}{6.450026in}}{\pgfqpoint{0.357298in}{6.462415in}}%
\pgfpathlineto{\pgfqpoint{0.359160in}{6.471510in}}%
\pgfpathquadraticcurveto{\pgfqpoint{0.362383in}{6.488364in}}{\pgfqpoint{0.370475in}{6.496385in}}%
\pgfpathquadraticcurveto{\pgfqpoint{0.378568in}{6.504405in}}{\pgfqpoint{0.392127in}{6.504405in}}%
\pgfpathquadraticcurveto{\pgfqpoint{0.408336in}{6.504405in}}{\pgfqpoint{0.416691in}{6.493544in}}%
\pgfpathquadraticcurveto{\pgfqpoint{0.425046in}{6.482706in}}{\pgfqpoint{0.425046in}{6.461747in}}%
\pgfpathquadraticcurveto{\pgfqpoint{0.425046in}{6.453845in}}{\pgfqpoint{0.423255in}{6.444918in}}%
\pgfpathquadraticcurveto{\pgfqpoint{0.421465in}{6.436013in}}{\pgfqpoint{0.417956in}{6.426465in}}%
\pgfpathlineto{\pgfqpoint{0.402439in}{6.426465in}}%
\pgfpathquadraticcurveto{\pgfqpoint{0.407572in}{6.435632in}}{\pgfqpoint{0.410198in}{6.444440in}}%
\pgfpathquadraticcurveto{\pgfqpoint{0.412799in}{6.453249in}}{\pgfqpoint{0.412799in}{6.461747in}}%
\pgfpathquadraticcurveto{\pgfqpoint{0.412799in}{6.474638in}}{\pgfqpoint{0.407739in}{6.481656in}}%
\pgfpathquadraticcurveto{\pgfqpoint{0.402654in}{6.488674in}}{\pgfqpoint{0.393249in}{6.488674in}}%
\pgfpathquadraticcurveto{\pgfqpoint{0.385061in}{6.488674in}}{\pgfqpoint{0.380430in}{6.483637in}}%
\pgfpathquadraticcurveto{\pgfqpoint{0.375799in}{6.478600in}}{\pgfqpoint{0.373483in}{6.467118in}}%
\pgfpathlineto{\pgfqpoint{0.371693in}{6.457928in}}%
\pgfpathquadraticcurveto{\pgfqpoint{0.368351in}{6.441074in}}{\pgfqpoint{0.361189in}{6.433531in}}%
\pgfpathquadraticcurveto{\pgfqpoint{0.354028in}{6.426011in}}{\pgfqpoint{0.341257in}{6.426011in}}%
\pgfpathquadraticcurveto{\pgfqpoint{0.326480in}{6.426011in}}{\pgfqpoint{0.317982in}{6.436419in}}%
\pgfpathquadraticcurveto{\pgfqpoint{0.309484in}{6.446827in}}{\pgfqpoint{0.309484in}{6.465089in}}%
\pgfpathquadraticcurveto{\pgfqpoint{0.309484in}{6.472943in}}{\pgfqpoint{0.310916in}{6.481059in}}%
\pgfpathquadraticcurveto{\pgfqpoint{0.312324in}{6.489199in}}{\pgfqpoint{0.315141in}{6.497697in}}%
\pgfpathlineto{\pgfqpoint{0.315141in}{6.497697in}}%
\pgfpathclose%
\pgfpathmoveto{\pgfqpoint{0.311489in}{6.527919in}}%
\pgfpathlineto{\pgfqpoint{0.311489in}{6.598340in}}%
\pgfpathlineto{\pgfqpoint{0.324188in}{6.598340in}}%
\pgfpathlineto{\pgfqpoint{0.324188in}{6.542982in}}%
\pgfpathlineto{\pgfqpoint{0.357155in}{6.542982in}}%
\pgfpathlineto{\pgfqpoint{0.357155in}{6.596024in}}%
\pgfpathlineto{\pgfqpoint{0.369831in}{6.596024in}}%
\pgfpathlineto{\pgfqpoint{0.369831in}{6.542982in}}%
\pgfpathlineto{\pgfqpoint{0.410198in}{6.542982in}}%
\pgfpathlineto{\pgfqpoint{0.410198in}{6.599677in}}%
\pgfpathlineto{\pgfqpoint{0.422873in}{6.599677in}}%
\pgfpathlineto{\pgfqpoint{0.422873in}{6.527919in}}%
\pgfpathlineto{\pgfqpoint{0.311489in}{6.527919in}}%
\pgfpathlineto{\pgfqpoint{0.311489in}{6.527919in}}%
\pgfpathclose%
\pgfpathmoveto{\pgfqpoint{0.341782in}{6.677116in}}%
\pgfpathlineto{\pgfqpoint{0.354768in}{6.677116in}}%
\pgfpathquadraticcurveto{\pgfqpoint{0.351784in}{6.671315in}}{\pgfqpoint{0.350304in}{6.665037in}}%
\pgfpathquadraticcurveto{\pgfqpoint{0.348800in}{6.658783in}}{\pgfqpoint{0.348800in}{6.652051in}}%
\pgfpathquadraticcurveto{\pgfqpoint{0.348800in}{6.641834in}}{\pgfqpoint{0.351927in}{6.636725in}}%
\pgfpathquadraticcurveto{\pgfqpoint{0.355054in}{6.631617in}}{\pgfqpoint{0.361332in}{6.631617in}}%
\pgfpathquadraticcurveto{\pgfqpoint{0.366107in}{6.631617in}}{\pgfqpoint{0.368828in}{6.635269in}}%
\pgfpathquadraticcurveto{\pgfqpoint{0.371549in}{6.638921in}}{\pgfqpoint{0.374008in}{6.649974in}}%
\pgfpathlineto{\pgfqpoint{0.375059in}{6.654677in}}%
\pgfpathquadraticcurveto{\pgfqpoint{0.378186in}{6.669286in}}{\pgfqpoint{0.383891in}{6.675445in}}%
\pgfpathquadraticcurveto{\pgfqpoint{0.389596in}{6.681604in}}{\pgfqpoint{0.399813in}{6.681604in}}%
\pgfpathquadraticcurveto{\pgfqpoint{0.411463in}{6.681604in}}{\pgfqpoint{0.418266in}{6.672389in}}%
\pgfpathquadraticcurveto{\pgfqpoint{0.425046in}{6.663175in}}{\pgfqpoint{0.425046in}{6.647062in}}%
\pgfpathquadraticcurveto{\pgfqpoint{0.425046in}{6.640354in}}{\pgfqpoint{0.423733in}{6.633073in}}%
\pgfpathquadraticcurveto{\pgfqpoint{0.422420in}{6.625792in}}{\pgfqpoint{0.419818in}{6.617747in}}%
\pgfpathlineto{\pgfqpoint{0.405638in}{6.617747in}}%
\pgfpathquadraticcurveto{\pgfqpoint{0.409601in}{6.625362in}}{\pgfqpoint{0.411582in}{6.632739in}}%
\pgfpathquadraticcurveto{\pgfqpoint{0.413540in}{6.640115in}}{\pgfqpoint{0.413540in}{6.647372in}}%
\pgfpathquadraticcurveto{\pgfqpoint{0.413540in}{6.657064in}}{\pgfqpoint{0.410221in}{6.662268in}}%
\pgfpathquadraticcurveto{\pgfqpoint{0.406903in}{6.667496in}}{\pgfqpoint{0.400864in}{6.667496in}}%
\pgfpathquadraticcurveto{\pgfqpoint{0.395278in}{6.667496in}}{\pgfqpoint{0.392294in}{6.663724in}}%
\pgfpathquadraticcurveto{\pgfqpoint{0.389310in}{6.659976in}}{\pgfqpoint{0.386541in}{6.647205in}}%
\pgfpathlineto{\pgfqpoint{0.385419in}{6.642431in}}%
\pgfpathquadraticcurveto{\pgfqpoint{0.382745in}{6.629683in}}{\pgfqpoint{0.377183in}{6.624002in}}%
\pgfpathquadraticcurveto{\pgfqpoint{0.371621in}{6.618344in}}{\pgfqpoint{0.361929in}{6.618344in}}%
\pgfpathquadraticcurveto{\pgfqpoint{0.350137in}{6.618344in}}{\pgfqpoint{0.343739in}{6.626699in}}%
\pgfpathquadraticcurveto{\pgfqpoint{0.337318in}{6.635054in}}{\pgfqpoint{0.337318in}{6.650428in}}%
\pgfpathquadraticcurveto{\pgfqpoint{0.337318in}{6.658019in}}{\pgfqpoint{0.338440in}{6.664727in}}%
\pgfpathquadraticcurveto{\pgfqpoint{0.339538in}{6.671458in}}{\pgfqpoint{0.341782in}{6.677116in}}%
\pgfpathlineto{\pgfqpoint{0.341782in}{6.677116in}}%
\pgfpathclose%
\pgfpathmoveto{\pgfqpoint{0.306953in}{6.701298in}}%
\pgfpathlineto{\pgfqpoint{0.306953in}{6.713234in}}%
\pgfpathquadraticcurveto{\pgfqpoint{0.324546in}{6.724405in}}{\pgfqpoint{0.341424in}{6.729967in}}%
\pgfpathquadraticcurveto{\pgfqpoint{0.358277in}{6.735530in}}{\pgfqpoint{0.374915in}{6.735530in}}%
\pgfpathquadraticcurveto{\pgfqpoint{0.391625in}{6.735530in}}{\pgfqpoint{0.408550in}{6.729967in}}%
\pgfpathquadraticcurveto{\pgfqpoint{0.425475in}{6.724405in}}{\pgfqpoint{0.443021in}{6.713234in}}%
\pgfpathlineto{\pgfqpoint{0.443021in}{6.701298in}}%
\pgfpathquadraticcurveto{\pgfqpoint{0.425929in}{6.711204in}}{\pgfqpoint{0.409028in}{6.716098in}}%
\pgfpathquadraticcurveto{\pgfqpoint{0.392127in}{6.720992in}}{\pgfqpoint{0.374915in}{6.720992in}}%
\pgfpathquadraticcurveto{\pgfqpoint{0.357680in}{6.720992in}}{\pgfqpoint{0.340898in}{6.716098in}}%
\pgfpathquadraticcurveto{\pgfqpoint{0.324093in}{6.711204in}}{\pgfqpoint{0.306953in}{6.701298in}}%
\pgfpathlineto{\pgfqpoint{0.306953in}{6.701298in}}%
\pgfpathclose%
\pgfusepath{fill}%
\end{pgfscope}%
\begin{pgfscope}%
\pgfpathrectangle{\pgfqpoint{0.977909in}{2.928000in}}{\pgfqpoint{5.664000in}{5.664000in}}%
\pgfusepath{clip}%
\pgfsetrectcap%
\pgfsetroundjoin%
\pgfsetlinewidth{0.853187pt}%
\definecolor{currentstroke}{rgb}{0.745098,0.776471,0.803922}%
\pgfsetstrokecolor{currentstroke}%
\pgfsetdash{}{0pt}%
\pgfpathmoveto{\pgfqpoint{0.977909in}{5.760000in}}%
\pgfpathlineto{\pgfqpoint{6.641909in}{5.760000in}}%
\pgfusepath{stroke}%
\end{pgfscope}%
\begin{pgfscope}%
\pgfpathrectangle{\pgfqpoint{0.977909in}{2.928000in}}{\pgfqpoint{5.664000in}{5.664000in}}%
\pgfusepath{clip}%
\pgfsetrectcap%
\pgfsetroundjoin%
\pgfsetlinewidth{0.853187pt}%
\definecolor{currentstroke}{rgb}{0.745098,0.776471,0.803922}%
\pgfsetstrokecolor{currentstroke}%
\pgfsetdash{}{0pt}%
\pgfpathmoveto{\pgfqpoint{3.809909in}{2.928000in}}%
\pgfpathlineto{\pgfqpoint{3.809909in}{8.592000in}}%
\pgfusepath{stroke}%
\end{pgfscope}%
\begin{pgfscope}%
\pgfpathrectangle{\pgfqpoint{0.977909in}{2.928000in}}{\pgfqpoint{5.664000in}{5.664000in}}%
\pgfusepath{clip}%
\pgfsetbuttcap%
\pgfsetroundjoin%
\pgfsetlinewidth{1.104125pt}%
\definecolor{currentstroke}{rgb}{0.203922,0.243137,0.278431}%
\pgfsetstrokecolor{currentstroke}%
\pgfsetdash{{5.500000pt}{4.400000pt}}{0.000000pt}%
\pgfpathmoveto{\pgfqpoint{0.977909in}{2.928000in}}%
\pgfpathlineto{\pgfqpoint{6.641909in}{8.592000in}}%
\pgfusepath{stroke}%
\end{pgfscope}%
\begin{pgfscope}%
\pgfsetrectcap%
\pgfsetmiterjoin%
\pgfsetlinewidth{0.803000pt}%
\definecolor{currentstroke}{rgb}{0.627451,0.670588,0.705882}%
\pgfsetstrokecolor{currentstroke}%
\pgfsetdash{}{0pt}%
\pgfpathmoveto{\pgfqpoint{0.977909in}{2.928000in}}%
\pgfpathlineto{\pgfqpoint{0.977909in}{8.592000in}}%
\pgfusepath{stroke}%
\end{pgfscope}%
\begin{pgfscope}%
\pgfsetrectcap%
\pgfsetmiterjoin%
\pgfsetlinewidth{0.803000pt}%
\definecolor{currentstroke}{rgb}{0.627451,0.670588,0.705882}%
\pgfsetstrokecolor{currentstroke}%
\pgfsetdash{}{0pt}%
\pgfpathmoveto{\pgfqpoint{0.977909in}{2.928000in}}%
\pgfpathlineto{\pgfqpoint{6.641909in}{2.928000in}}%
\pgfusepath{stroke}%
\end{pgfscope}%
\begin{pgfscope}%
\pgfpathrectangle{\pgfqpoint{0.977909in}{2.928000in}}{\pgfqpoint{5.664000in}{5.664000in}}%
\pgfusepath{clip}%
\pgfsetbuttcap%
\pgfsetroundjoin%
\pgfsetlinewidth{1.003750pt}%
\definecolor{currentstroke}{rgb}{0.541176,0.380392,0.658824}%
\pgfsetstrokecolor{currentstroke}%
\pgfsetstrokeopacity{0.900000}%
\pgfsetdash{}{0pt}%
\pgfpathmoveto{\pgfqpoint{5.127260in}{7.303561in}}%
\pgfpathcurveto{\pgfqpoint{5.137514in}{7.303561in}}{\pgfqpoint{5.147350in}{7.307635in}}{\pgfqpoint{5.154601in}{7.314886in}}%
\pgfpathcurveto{\pgfqpoint{5.161851in}{7.322136in}}{\pgfqpoint{5.165925in}{7.331972in}}{\pgfqpoint{5.165925in}{7.342226in}}%
\pgfpathcurveto{\pgfqpoint{5.165925in}{7.352480in}}{\pgfqpoint{5.161851in}{7.362316in}}{\pgfqpoint{5.154601in}{7.369566in}}%
\pgfpathcurveto{\pgfqpoint{5.147350in}{7.376817in}}{\pgfqpoint{5.137514in}{7.380891in}}{\pgfqpoint{5.127260in}{7.380891in}}%
\pgfpathcurveto{\pgfqpoint{5.117006in}{7.380891in}}{\pgfqpoint{5.107171in}{7.376817in}}{\pgfqpoint{5.099920in}{7.369566in}}%
\pgfpathcurveto{\pgfqpoint{5.092669in}{7.362316in}}{\pgfqpoint{5.088595in}{7.352480in}}{\pgfqpoint{5.088595in}{7.342226in}}%
\pgfpathcurveto{\pgfqpoint{5.088595in}{7.331972in}}{\pgfqpoint{5.092669in}{7.322136in}}{\pgfqpoint{5.099920in}{7.314886in}}%
\pgfpathcurveto{\pgfqpoint{5.107171in}{7.307635in}}{\pgfqpoint{5.117006in}{7.303561in}}{\pgfqpoint{5.127260in}{7.303561in}}%
\pgfpathlineto{\pgfqpoint{5.127260in}{7.303561in}}%
\pgfpathclose%
\pgfusepath{stroke}%
\end{pgfscope}%
\begin{pgfscope}%
\pgfpathrectangle{\pgfqpoint{0.977909in}{2.928000in}}{\pgfqpoint{5.664000in}{5.664000in}}%
\pgfusepath{clip}%
\pgfsetbuttcap%
\pgfsetroundjoin%
\pgfsetlinewidth{1.003750pt}%
\definecolor{currentstroke}{rgb}{0.541176,0.380392,0.658824}%
\pgfsetstrokecolor{currentstroke}%
\pgfsetstrokeopacity{0.900000}%
\pgfsetdash{}{0pt}%
\pgfpathmoveto{\pgfqpoint{5.009922in}{7.114794in}}%
\pgfpathcurveto{\pgfqpoint{5.020176in}{7.114794in}}{\pgfqpoint{5.030011in}{7.118868in}}{\pgfqpoint{5.037262in}{7.126119in}}%
\pgfpathcurveto{\pgfqpoint{5.044513in}{7.133370in}}{\pgfqpoint{5.048587in}{7.143205in}}{\pgfqpoint{5.048587in}{7.153459in}}%
\pgfpathcurveto{\pgfqpoint{5.048587in}{7.163713in}}{\pgfqpoint{5.044513in}{7.173549in}}{\pgfqpoint{5.037262in}{7.180800in}}%
\pgfpathcurveto{\pgfqpoint{5.030011in}{7.188050in}}{\pgfqpoint{5.020176in}{7.192124in}}{\pgfqpoint{5.009922in}{7.192124in}}%
\pgfpathcurveto{\pgfqpoint{4.999668in}{7.192124in}}{\pgfqpoint{4.989832in}{7.188050in}}{\pgfqpoint{4.982581in}{7.180800in}}%
\pgfpathcurveto{\pgfqpoint{4.975331in}{7.173549in}}{\pgfqpoint{4.971257in}{7.163713in}}{\pgfqpoint{4.971257in}{7.153459in}}%
\pgfpathcurveto{\pgfqpoint{4.971257in}{7.143205in}}{\pgfqpoint{4.975331in}{7.133370in}}{\pgfqpoint{4.982581in}{7.126119in}}%
\pgfpathcurveto{\pgfqpoint{4.989832in}{7.118868in}}{\pgfqpoint{4.999668in}{7.114794in}}{\pgfqpoint{5.009922in}{7.114794in}}%
\pgfpathlineto{\pgfqpoint{5.009922in}{7.114794in}}%
\pgfpathclose%
\pgfusepath{stroke}%
\end{pgfscope}%
\begin{pgfscope}%
\pgfpathrectangle{\pgfqpoint{0.977909in}{2.928000in}}{\pgfqpoint{5.664000in}{5.664000in}}%
\pgfusepath{clip}%
\pgfsetbuttcap%
\pgfsetroundjoin%
\pgfsetlinewidth{1.003750pt}%
\definecolor{currentstroke}{rgb}{0.541176,0.380392,0.658824}%
\pgfsetstrokecolor{currentstroke}%
\pgfsetstrokeopacity{0.900000}%
\pgfsetdash{}{0pt}%
\pgfpathmoveto{\pgfqpoint{6.010907in}{7.413420in}}%
\pgfpathcurveto{\pgfqpoint{6.021161in}{7.413420in}}{\pgfqpoint{6.030997in}{7.417494in}}{\pgfqpoint{6.038248in}{7.424745in}}%
\pgfpathcurveto{\pgfqpoint{6.045498in}{7.431996in}}{\pgfqpoint{6.049572in}{7.441831in}}{\pgfqpoint{6.049572in}{7.452085in}}%
\pgfpathcurveto{\pgfqpoint{6.049572in}{7.462339in}}{\pgfqpoint{6.045498in}{7.472175in}}{\pgfqpoint{6.038248in}{7.479426in}}%
\pgfpathcurveto{\pgfqpoint{6.030997in}{7.486676in}}{\pgfqpoint{6.021161in}{7.490750in}}{\pgfqpoint{6.010907in}{7.490750in}}%
\pgfpathcurveto{\pgfqpoint{6.000653in}{7.490750in}}{\pgfqpoint{5.990818in}{7.486676in}}{\pgfqpoint{5.983567in}{7.479426in}}%
\pgfpathcurveto{\pgfqpoint{5.976316in}{7.472175in}}{\pgfqpoint{5.972242in}{7.462339in}}{\pgfqpoint{5.972242in}{7.452085in}}%
\pgfpathcurveto{\pgfqpoint{5.972242in}{7.441831in}}{\pgfqpoint{5.976316in}{7.431996in}}{\pgfqpoint{5.983567in}{7.424745in}}%
\pgfpathcurveto{\pgfqpoint{5.990818in}{7.417494in}}{\pgfqpoint{6.000653in}{7.413420in}}{\pgfqpoint{6.010907in}{7.413420in}}%
\pgfpathlineto{\pgfqpoint{6.010907in}{7.413420in}}%
\pgfpathclose%
\pgfusepath{stroke}%
\end{pgfscope}%
\begin{pgfscope}%
\pgfpathrectangle{\pgfqpoint{0.977909in}{2.928000in}}{\pgfqpoint{5.664000in}{5.664000in}}%
\pgfusepath{clip}%
\pgfsetbuttcap%
\pgfsetroundjoin%
\pgfsetlinewidth{1.003750pt}%
\definecolor{currentstroke}{rgb}{0.541176,0.380392,0.658824}%
\pgfsetstrokecolor{currentstroke}%
\pgfsetstrokeopacity{0.900000}%
\pgfsetdash{}{0pt}%
\pgfpathmoveto{\pgfqpoint{4.941413in}{6.373192in}}%
\pgfpathcurveto{\pgfqpoint{4.951667in}{6.373192in}}{\pgfqpoint{4.961503in}{6.377266in}}{\pgfqpoint{4.968753in}{6.384517in}}%
\pgfpathcurveto{\pgfqpoint{4.976004in}{6.391768in}}{\pgfqpoint{4.980078in}{6.401603in}}{\pgfqpoint{4.980078in}{6.411857in}}%
\pgfpathcurveto{\pgfqpoint{4.980078in}{6.422112in}}{\pgfqpoint{4.976004in}{6.431947in}}{\pgfqpoint{4.968753in}{6.439198in}}%
\pgfpathcurveto{\pgfqpoint{4.961503in}{6.446449in}}{\pgfqpoint{4.951667in}{6.450523in}}{\pgfqpoint{4.941413in}{6.450523in}}%
\pgfpathcurveto{\pgfqpoint{4.931159in}{6.450523in}}{\pgfqpoint{4.921324in}{6.446449in}}{\pgfqpoint{4.914073in}{6.439198in}}%
\pgfpathcurveto{\pgfqpoint{4.906822in}{6.431947in}}{\pgfqpoint{4.902748in}{6.422112in}}{\pgfqpoint{4.902748in}{6.411857in}}%
\pgfpathcurveto{\pgfqpoint{4.902748in}{6.401603in}}{\pgfqpoint{4.906822in}{6.391768in}}{\pgfqpoint{4.914073in}{6.384517in}}%
\pgfpathcurveto{\pgfqpoint{4.921324in}{6.377266in}}{\pgfqpoint{4.931159in}{6.373192in}}{\pgfqpoint{4.941413in}{6.373192in}}%
\pgfpathlineto{\pgfqpoint{4.941413in}{6.373192in}}%
\pgfpathclose%
\pgfusepath{stroke}%
\end{pgfscope}%
\begin{pgfscope}%
\pgfpathrectangle{\pgfqpoint{0.977909in}{2.928000in}}{\pgfqpoint{5.664000in}{5.664000in}}%
\pgfusepath{clip}%
\pgfsetbuttcap%
\pgfsetroundjoin%
\pgfsetlinewidth{1.003750pt}%
\definecolor{currentstroke}{rgb}{0.541176,0.380392,0.658824}%
\pgfsetstrokecolor{currentstroke}%
\pgfsetstrokeopacity{0.900000}%
\pgfsetdash{}{0pt}%
\pgfpathmoveto{\pgfqpoint{5.682610in}{7.540059in}}%
\pgfpathcurveto{\pgfqpoint{5.692864in}{7.540059in}}{\pgfqpoint{5.702700in}{7.544133in}}{\pgfqpoint{5.709950in}{7.551384in}}%
\pgfpathcurveto{\pgfqpoint{5.717201in}{7.558635in}}{\pgfqpoint{5.721275in}{7.568470in}}{\pgfqpoint{5.721275in}{7.578724in}}%
\pgfpathcurveto{\pgfqpoint{5.721275in}{7.588978in}}{\pgfqpoint{5.717201in}{7.598814in}}{\pgfqpoint{5.709950in}{7.606064in}}%
\pgfpathcurveto{\pgfqpoint{5.702700in}{7.613315in}}{\pgfqpoint{5.692864in}{7.617389in}}{\pgfqpoint{5.682610in}{7.617389in}}%
\pgfpathcurveto{\pgfqpoint{5.672356in}{7.617389in}}{\pgfqpoint{5.662521in}{7.613315in}}{\pgfqpoint{5.655270in}{7.606064in}}%
\pgfpathcurveto{\pgfqpoint{5.648019in}{7.598814in}}{\pgfqpoint{5.643945in}{7.588978in}}{\pgfqpoint{5.643945in}{7.578724in}}%
\pgfpathcurveto{\pgfqpoint{5.643945in}{7.568470in}}{\pgfqpoint{5.648019in}{7.558635in}}{\pgfqpoint{5.655270in}{7.551384in}}%
\pgfpathcurveto{\pgfqpoint{5.662521in}{7.544133in}}{\pgfqpoint{5.672356in}{7.540059in}}{\pgfqpoint{5.682610in}{7.540059in}}%
\pgfpathlineto{\pgfqpoint{5.682610in}{7.540059in}}%
\pgfpathclose%
\pgfusepath{stroke}%
\end{pgfscope}%
\begin{pgfscope}%
\pgfpathrectangle{\pgfqpoint{0.977909in}{2.928000in}}{\pgfqpoint{5.664000in}{5.664000in}}%
\pgfusepath{clip}%
\pgfsetbuttcap%
\pgfsetroundjoin%
\pgfsetlinewidth{1.003750pt}%
\definecolor{currentstroke}{rgb}{0.541176,0.380392,0.658824}%
\pgfsetstrokecolor{currentstroke}%
\pgfsetstrokeopacity{0.900000}%
\pgfsetdash{}{0pt}%
\pgfpathmoveto{\pgfqpoint{4.734480in}{6.266923in}}%
\pgfpathcurveto{\pgfqpoint{4.744734in}{6.266923in}}{\pgfqpoint{4.754569in}{6.270997in}}{\pgfqpoint{4.761820in}{6.278247in}}%
\pgfpathcurveto{\pgfqpoint{4.769071in}{6.285498in}}{\pgfqpoint{4.773145in}{6.295334in}}{\pgfqpoint{4.773145in}{6.305588in}}%
\pgfpathcurveto{\pgfqpoint{4.773145in}{6.315842in}}{\pgfqpoint{4.769071in}{6.325677in}}{\pgfqpoint{4.761820in}{6.332928in}}%
\pgfpathcurveto{\pgfqpoint{4.754569in}{6.340179in}}{\pgfqpoint{4.744734in}{6.344253in}}{\pgfqpoint{4.734480in}{6.344253in}}%
\pgfpathcurveto{\pgfqpoint{4.724226in}{6.344253in}}{\pgfqpoint{4.714390in}{6.340179in}}{\pgfqpoint{4.707139in}{6.332928in}}%
\pgfpathcurveto{\pgfqpoint{4.699889in}{6.325677in}}{\pgfqpoint{4.695815in}{6.315842in}}{\pgfqpoint{4.695815in}{6.305588in}}%
\pgfpathcurveto{\pgfqpoint{4.695815in}{6.295334in}}{\pgfqpoint{4.699889in}{6.285498in}}{\pgfqpoint{4.707139in}{6.278247in}}%
\pgfpathcurveto{\pgfqpoint{4.714390in}{6.270997in}}{\pgfqpoint{4.724226in}{6.266923in}}{\pgfqpoint{4.734480in}{6.266923in}}%
\pgfpathlineto{\pgfqpoint{4.734480in}{6.266923in}}%
\pgfpathclose%
\pgfusepath{stroke}%
\end{pgfscope}%
\begin{pgfscope}%
\pgfpathrectangle{\pgfqpoint{0.977909in}{2.928000in}}{\pgfqpoint{5.664000in}{5.664000in}}%
\pgfusepath{clip}%
\pgfsetbuttcap%
\pgfsetroundjoin%
\pgfsetlinewidth{1.003750pt}%
\definecolor{currentstroke}{rgb}{0.541176,0.380392,0.658824}%
\pgfsetstrokecolor{currentstroke}%
\pgfsetstrokeopacity{0.900000}%
\pgfsetdash{}{0pt}%
\pgfpathmoveto{\pgfqpoint{5.589652in}{7.231294in}}%
\pgfpathcurveto{\pgfqpoint{5.599906in}{7.231294in}}{\pgfqpoint{5.609742in}{7.235368in}}{\pgfqpoint{5.616993in}{7.242618in}}%
\pgfpathcurveto{\pgfqpoint{5.624243in}{7.249869in}}{\pgfqpoint{5.628317in}{7.259705in}}{\pgfqpoint{5.628317in}{7.269959in}}%
\pgfpathcurveto{\pgfqpoint{5.628317in}{7.280213in}}{\pgfqpoint{5.624243in}{7.290048in}}{\pgfqpoint{5.616993in}{7.297299in}}%
\pgfpathcurveto{\pgfqpoint{5.609742in}{7.304550in}}{\pgfqpoint{5.599906in}{7.308624in}}{\pgfqpoint{5.589652in}{7.308624in}}%
\pgfpathcurveto{\pgfqpoint{5.579398in}{7.308624in}}{\pgfqpoint{5.569563in}{7.304550in}}{\pgfqpoint{5.562312in}{7.297299in}}%
\pgfpathcurveto{\pgfqpoint{5.555061in}{7.290048in}}{\pgfqpoint{5.550987in}{7.280213in}}{\pgfqpoint{5.550987in}{7.269959in}}%
\pgfpathcurveto{\pgfqpoint{5.550987in}{7.259705in}}{\pgfqpoint{5.555061in}{7.249869in}}{\pgfqpoint{5.562312in}{7.242618in}}%
\pgfpathcurveto{\pgfqpoint{5.569563in}{7.235368in}}{\pgfqpoint{5.579398in}{7.231294in}}{\pgfqpoint{5.589652in}{7.231294in}}%
\pgfpathlineto{\pgfqpoint{5.589652in}{7.231294in}}%
\pgfpathclose%
\pgfusepath{stroke}%
\end{pgfscope}%
\begin{pgfscope}%
\pgfpathrectangle{\pgfqpoint{0.977909in}{2.928000in}}{\pgfqpoint{5.664000in}{5.664000in}}%
\pgfusepath{clip}%
\pgfsetbuttcap%
\pgfsetroundjoin%
\pgfsetlinewidth{1.003750pt}%
\definecolor{currentstroke}{rgb}{0.541176,0.380392,0.658824}%
\pgfsetstrokecolor{currentstroke}%
\pgfsetstrokeopacity{0.900000}%
\pgfsetdash{}{0pt}%
\pgfpathmoveto{\pgfqpoint{3.930945in}{6.081865in}}%
\pgfpathcurveto{\pgfqpoint{3.941199in}{6.081865in}}{\pgfqpoint{3.951034in}{6.085939in}}{\pgfqpoint{3.958285in}{6.093190in}}%
\pgfpathcurveto{\pgfqpoint{3.965536in}{6.100441in}}{\pgfqpoint{3.969610in}{6.110276in}}{\pgfqpoint{3.969610in}{6.120530in}}%
\pgfpathcurveto{\pgfqpoint{3.969610in}{6.130784in}}{\pgfqpoint{3.965536in}{6.140620in}}{\pgfqpoint{3.958285in}{6.147871in}}%
\pgfpathcurveto{\pgfqpoint{3.951034in}{6.155121in}}{\pgfqpoint{3.941199in}{6.159195in}}{\pgfqpoint{3.930945in}{6.159195in}}%
\pgfpathcurveto{\pgfqpoint{3.920691in}{6.159195in}}{\pgfqpoint{3.910855in}{6.155121in}}{\pgfqpoint{3.903604in}{6.147871in}}%
\pgfpathcurveto{\pgfqpoint{3.896354in}{6.140620in}}{\pgfqpoint{3.892280in}{6.130784in}}{\pgfqpoint{3.892280in}{6.120530in}}%
\pgfpathcurveto{\pgfqpoint{3.892280in}{6.110276in}}{\pgfqpoint{3.896354in}{6.100441in}}{\pgfqpoint{3.903604in}{6.093190in}}%
\pgfpathcurveto{\pgfqpoint{3.910855in}{6.085939in}}{\pgfqpoint{3.920691in}{6.081865in}}{\pgfqpoint{3.930945in}{6.081865in}}%
\pgfpathlineto{\pgfqpoint{3.930945in}{6.081865in}}%
\pgfpathclose%
\pgfusepath{stroke}%
\end{pgfscope}%
\begin{pgfscope}%
\pgfpathrectangle{\pgfqpoint{0.977909in}{2.928000in}}{\pgfqpoint{5.664000in}{5.664000in}}%
\pgfusepath{clip}%
\pgfsetbuttcap%
\pgfsetroundjoin%
\pgfsetlinewidth{1.003750pt}%
\definecolor{currentstroke}{rgb}{0.541176,0.380392,0.658824}%
\pgfsetstrokecolor{currentstroke}%
\pgfsetstrokeopacity{0.900000}%
\pgfsetdash{}{0pt}%
\pgfpathmoveto{\pgfqpoint{4.038363in}{5.416770in}}%
\pgfpathcurveto{\pgfqpoint{4.048617in}{5.416770in}}{\pgfqpoint{4.058452in}{5.420844in}}{\pgfqpoint{4.065703in}{5.428095in}}%
\pgfpathcurveto{\pgfqpoint{4.072954in}{5.435346in}}{\pgfqpoint{4.077028in}{5.445181in}}{\pgfqpoint{4.077028in}{5.455436in}}%
\pgfpathcurveto{\pgfqpoint{4.077028in}{5.465690in}}{\pgfqpoint{4.072954in}{5.475525in}}{\pgfqpoint{4.065703in}{5.482776in}}%
\pgfpathcurveto{\pgfqpoint{4.058452in}{5.490027in}}{\pgfqpoint{4.048617in}{5.494101in}}{\pgfqpoint{4.038363in}{5.494101in}}%
\pgfpathcurveto{\pgfqpoint{4.028109in}{5.494101in}}{\pgfqpoint{4.018273in}{5.490027in}}{\pgfqpoint{4.011023in}{5.482776in}}%
\pgfpathcurveto{\pgfqpoint{4.003772in}{5.475525in}}{\pgfqpoint{3.999698in}{5.465690in}}{\pgfqpoint{3.999698in}{5.455436in}}%
\pgfpathcurveto{\pgfqpoint{3.999698in}{5.445181in}}{\pgfqpoint{4.003772in}{5.435346in}}{\pgfqpoint{4.011023in}{5.428095in}}%
\pgfpathcurveto{\pgfqpoint{4.018273in}{5.420844in}}{\pgfqpoint{4.028109in}{5.416770in}}{\pgfqpoint{4.038363in}{5.416770in}}%
\pgfpathlineto{\pgfqpoint{4.038363in}{5.416770in}}%
\pgfpathclose%
\pgfusepath{stroke}%
\end{pgfscope}%
\begin{pgfscope}%
\pgfpathrectangle{\pgfqpoint{0.977909in}{2.928000in}}{\pgfqpoint{5.664000in}{5.664000in}}%
\pgfusepath{clip}%
\pgfsetbuttcap%
\pgfsetroundjoin%
\pgfsetlinewidth{1.003750pt}%
\definecolor{currentstroke}{rgb}{0.541176,0.380392,0.658824}%
\pgfsetstrokecolor{currentstroke}%
\pgfsetstrokeopacity{0.900000}%
\pgfsetdash{}{0pt}%
\pgfpathmoveto{\pgfqpoint{4.811275in}{6.660347in}}%
\pgfpathcurveto{\pgfqpoint{4.821529in}{6.660347in}}{\pgfqpoint{4.831365in}{6.664421in}}{\pgfqpoint{4.838616in}{6.671672in}}%
\pgfpathcurveto{\pgfqpoint{4.845866in}{6.678922in}}{\pgfqpoint{4.849940in}{6.688758in}}{\pgfqpoint{4.849940in}{6.699012in}}%
\pgfpathcurveto{\pgfqpoint{4.849940in}{6.709266in}}{\pgfqpoint{4.845866in}{6.719101in}}{\pgfqpoint{4.838616in}{6.726352in}}%
\pgfpathcurveto{\pgfqpoint{4.831365in}{6.733603in}}{\pgfqpoint{4.821529in}{6.737677in}}{\pgfqpoint{4.811275in}{6.737677in}}%
\pgfpathcurveto{\pgfqpoint{4.801021in}{6.737677in}}{\pgfqpoint{4.791186in}{6.733603in}}{\pgfqpoint{4.783935in}{6.726352in}}%
\pgfpathcurveto{\pgfqpoint{4.776684in}{6.719101in}}{\pgfqpoint{4.772610in}{6.709266in}}{\pgfqpoint{4.772610in}{6.699012in}}%
\pgfpathcurveto{\pgfqpoint{4.772610in}{6.688758in}}{\pgfqpoint{4.776684in}{6.678922in}}{\pgfqpoint{4.783935in}{6.671672in}}%
\pgfpathcurveto{\pgfqpoint{4.791186in}{6.664421in}}{\pgfqpoint{4.801021in}{6.660347in}}{\pgfqpoint{4.811275in}{6.660347in}}%
\pgfpathlineto{\pgfqpoint{4.811275in}{6.660347in}}%
\pgfpathclose%
\pgfusepath{stroke}%
\end{pgfscope}%
\begin{pgfscope}%
\pgfpathrectangle{\pgfqpoint{0.977909in}{2.928000in}}{\pgfqpoint{5.664000in}{5.664000in}}%
\pgfusepath{clip}%
\pgfsetbuttcap%
\pgfsetroundjoin%
\pgfsetlinewidth{1.003750pt}%
\definecolor{currentstroke}{rgb}{0.541176,0.380392,0.658824}%
\pgfsetstrokecolor{currentstroke}%
\pgfsetstrokeopacity{0.900000}%
\pgfsetdash{}{0pt}%
\pgfpathmoveto{\pgfqpoint{4.291224in}{5.640433in}}%
\pgfpathcurveto{\pgfqpoint{4.301478in}{5.640433in}}{\pgfqpoint{4.311314in}{5.644507in}}{\pgfqpoint{4.318565in}{5.651757in}}%
\pgfpathcurveto{\pgfqpoint{4.325815in}{5.659008in}}{\pgfqpoint{4.329889in}{5.668844in}}{\pgfqpoint{4.329889in}{5.679098in}}%
\pgfpathcurveto{\pgfqpoint{4.329889in}{5.689352in}}{\pgfqpoint{4.325815in}{5.699187in}}{\pgfqpoint{4.318565in}{5.706438in}}%
\pgfpathcurveto{\pgfqpoint{4.311314in}{5.713689in}}{\pgfqpoint{4.301478in}{5.717763in}}{\pgfqpoint{4.291224in}{5.717763in}}%
\pgfpathcurveto{\pgfqpoint{4.280970in}{5.717763in}}{\pgfqpoint{4.271135in}{5.713689in}}{\pgfqpoint{4.263884in}{5.706438in}}%
\pgfpathcurveto{\pgfqpoint{4.256633in}{5.699187in}}{\pgfqpoint{4.252559in}{5.689352in}}{\pgfqpoint{4.252559in}{5.679098in}}%
\pgfpathcurveto{\pgfqpoint{4.252559in}{5.668844in}}{\pgfqpoint{4.256633in}{5.659008in}}{\pgfqpoint{4.263884in}{5.651757in}}%
\pgfpathcurveto{\pgfqpoint{4.271135in}{5.644507in}}{\pgfqpoint{4.280970in}{5.640433in}}{\pgfqpoint{4.291224in}{5.640433in}}%
\pgfpathlineto{\pgfqpoint{4.291224in}{5.640433in}}%
\pgfpathclose%
\pgfusepath{stroke}%
\end{pgfscope}%
\begin{pgfscope}%
\pgfpathrectangle{\pgfqpoint{0.977909in}{2.928000in}}{\pgfqpoint{5.664000in}{5.664000in}}%
\pgfusepath{clip}%
\pgfsetbuttcap%
\pgfsetroundjoin%
\pgfsetlinewidth{1.003750pt}%
\definecolor{currentstroke}{rgb}{0.541176,0.380392,0.658824}%
\pgfsetstrokecolor{currentstroke}%
\pgfsetstrokeopacity{0.900000}%
\pgfsetdash{}{0pt}%
\pgfpathmoveto{\pgfqpoint{4.841451in}{6.797467in}}%
\pgfpathcurveto{\pgfqpoint{4.851705in}{6.797467in}}{\pgfqpoint{4.861540in}{6.801541in}}{\pgfqpoint{4.868791in}{6.808792in}}%
\pgfpathcurveto{\pgfqpoint{4.876042in}{6.816042in}}{\pgfqpoint{4.880116in}{6.825878in}}{\pgfqpoint{4.880116in}{6.836132in}}%
\pgfpathcurveto{\pgfqpoint{4.880116in}{6.846386in}}{\pgfqpoint{4.876042in}{6.856221in}}{\pgfqpoint{4.868791in}{6.863472in}}%
\pgfpathcurveto{\pgfqpoint{4.861540in}{6.870723in}}{\pgfqpoint{4.851705in}{6.874797in}}{\pgfqpoint{4.841451in}{6.874797in}}%
\pgfpathcurveto{\pgfqpoint{4.831197in}{6.874797in}}{\pgfqpoint{4.821361in}{6.870723in}}{\pgfqpoint{4.814111in}{6.863472in}}%
\pgfpathcurveto{\pgfqpoint{4.806860in}{6.856221in}}{\pgfqpoint{4.802786in}{6.846386in}}{\pgfqpoint{4.802786in}{6.836132in}}%
\pgfpathcurveto{\pgfqpoint{4.802786in}{6.825878in}}{\pgfqpoint{4.806860in}{6.816042in}}{\pgfqpoint{4.814111in}{6.808792in}}%
\pgfpathcurveto{\pgfqpoint{4.821361in}{6.801541in}}{\pgfqpoint{4.831197in}{6.797467in}}{\pgfqpoint{4.841451in}{6.797467in}}%
\pgfpathlineto{\pgfqpoint{4.841451in}{6.797467in}}%
\pgfpathclose%
\pgfusepath{stroke}%
\end{pgfscope}%
\begin{pgfscope}%
\pgfpathrectangle{\pgfqpoint{0.977909in}{2.928000in}}{\pgfqpoint{5.664000in}{5.664000in}}%
\pgfusepath{clip}%
\pgfsetbuttcap%
\pgfsetroundjoin%
\pgfsetlinewidth{1.003750pt}%
\definecolor{currentstroke}{rgb}{0.541176,0.380392,0.658824}%
\pgfsetstrokecolor{currentstroke}%
\pgfsetstrokeopacity{0.900000}%
\pgfsetdash{}{0pt}%
\pgfpathmoveto{\pgfqpoint{4.733350in}{6.465960in}}%
\pgfpathcurveto{\pgfqpoint{4.743604in}{6.465960in}}{\pgfqpoint{4.753439in}{6.470034in}}{\pgfqpoint{4.760690in}{6.477284in}}%
\pgfpathcurveto{\pgfqpoint{4.767941in}{6.484535in}}{\pgfqpoint{4.772015in}{6.494371in}}{\pgfqpoint{4.772015in}{6.504625in}}%
\pgfpathcurveto{\pgfqpoint{4.772015in}{6.514879in}}{\pgfqpoint{4.767941in}{6.524714in}}{\pgfqpoint{4.760690in}{6.531965in}}%
\pgfpathcurveto{\pgfqpoint{4.753439in}{6.539216in}}{\pgfqpoint{4.743604in}{6.543290in}}{\pgfqpoint{4.733350in}{6.543290in}}%
\pgfpathcurveto{\pgfqpoint{4.723096in}{6.543290in}}{\pgfqpoint{4.713260in}{6.539216in}}{\pgfqpoint{4.706009in}{6.531965in}}%
\pgfpathcurveto{\pgfqpoint{4.698759in}{6.524714in}}{\pgfqpoint{4.694685in}{6.514879in}}{\pgfqpoint{4.694685in}{6.504625in}}%
\pgfpathcurveto{\pgfqpoint{4.694685in}{6.494371in}}{\pgfqpoint{4.698759in}{6.484535in}}{\pgfqpoint{4.706009in}{6.477284in}}%
\pgfpathcurveto{\pgfqpoint{4.713260in}{6.470034in}}{\pgfqpoint{4.723096in}{6.465960in}}{\pgfqpoint{4.733350in}{6.465960in}}%
\pgfpathlineto{\pgfqpoint{4.733350in}{6.465960in}}%
\pgfpathclose%
\pgfusepath{stroke}%
\end{pgfscope}%
\begin{pgfscope}%
\pgfpathrectangle{\pgfqpoint{0.977909in}{2.928000in}}{\pgfqpoint{5.664000in}{5.664000in}}%
\pgfusepath{clip}%
\pgfsetbuttcap%
\pgfsetroundjoin%
\pgfsetlinewidth{1.003750pt}%
\definecolor{currentstroke}{rgb}{0.541176,0.380392,0.658824}%
\pgfsetstrokecolor{currentstroke}%
\pgfsetstrokeopacity{0.900000}%
\pgfsetdash{}{0pt}%
\pgfpathmoveto{\pgfqpoint{4.223308in}{6.277735in}}%
\pgfpathcurveto{\pgfqpoint{4.233562in}{6.277735in}}{\pgfqpoint{4.243397in}{6.281809in}}{\pgfqpoint{4.250648in}{6.289060in}}%
\pgfpathcurveto{\pgfqpoint{4.257899in}{6.296311in}}{\pgfqpoint{4.261973in}{6.306146in}}{\pgfqpoint{4.261973in}{6.316400in}}%
\pgfpathcurveto{\pgfqpoint{4.261973in}{6.326654in}}{\pgfqpoint{4.257899in}{6.336490in}}{\pgfqpoint{4.250648in}{6.343741in}}%
\pgfpathcurveto{\pgfqpoint{4.243397in}{6.350991in}}{\pgfqpoint{4.233562in}{6.355065in}}{\pgfqpoint{4.223308in}{6.355065in}}%
\pgfpathcurveto{\pgfqpoint{4.213054in}{6.355065in}}{\pgfqpoint{4.203218in}{6.350991in}}{\pgfqpoint{4.195968in}{6.343741in}}%
\pgfpathcurveto{\pgfqpoint{4.188717in}{6.336490in}}{\pgfqpoint{4.184643in}{6.326654in}}{\pgfqpoint{4.184643in}{6.316400in}}%
\pgfpathcurveto{\pgfqpoint{4.184643in}{6.306146in}}{\pgfqpoint{4.188717in}{6.296311in}}{\pgfqpoint{4.195968in}{6.289060in}}%
\pgfpathcurveto{\pgfqpoint{4.203218in}{6.281809in}}{\pgfqpoint{4.213054in}{6.277735in}}{\pgfqpoint{4.223308in}{6.277735in}}%
\pgfpathlineto{\pgfqpoint{4.223308in}{6.277735in}}%
\pgfpathclose%
\pgfusepath{stroke}%
\end{pgfscope}%
\begin{pgfscope}%
\pgfpathrectangle{\pgfqpoint{0.977909in}{2.928000in}}{\pgfqpoint{5.664000in}{5.664000in}}%
\pgfusepath{clip}%
\pgfsetbuttcap%
\pgfsetroundjoin%
\pgfsetlinewidth{1.003750pt}%
\definecolor{currentstroke}{rgb}{0.541176,0.380392,0.658824}%
\pgfsetstrokecolor{currentstroke}%
\pgfsetstrokeopacity{0.900000}%
\pgfsetdash{}{0pt}%
\pgfpathmoveto{\pgfqpoint{4.037968in}{5.167673in}}%
\pgfpathcurveto{\pgfqpoint{4.048222in}{5.167673in}}{\pgfqpoint{4.058058in}{5.171747in}}{\pgfqpoint{4.065308in}{5.178998in}}%
\pgfpathcurveto{\pgfqpoint{4.072559in}{5.186249in}}{\pgfqpoint{4.076633in}{5.196084in}}{\pgfqpoint{4.076633in}{5.206338in}}%
\pgfpathcurveto{\pgfqpoint{4.076633in}{5.216592in}}{\pgfqpoint{4.072559in}{5.226428in}}{\pgfqpoint{4.065308in}{5.233679in}}%
\pgfpathcurveto{\pgfqpoint{4.058058in}{5.240929in}}{\pgfqpoint{4.048222in}{5.245003in}}{\pgfqpoint{4.037968in}{5.245003in}}%
\pgfpathcurveto{\pgfqpoint{4.027714in}{5.245003in}}{\pgfqpoint{4.017878in}{5.240929in}}{\pgfqpoint{4.010628in}{5.233679in}}%
\pgfpathcurveto{\pgfqpoint{4.003377in}{5.226428in}}{\pgfqpoint{3.999303in}{5.216592in}}{\pgfqpoint{3.999303in}{5.206338in}}%
\pgfpathcurveto{\pgfqpoint{3.999303in}{5.196084in}}{\pgfqpoint{4.003377in}{5.186249in}}{\pgfqpoint{4.010628in}{5.178998in}}%
\pgfpathcurveto{\pgfqpoint{4.017878in}{5.171747in}}{\pgfqpoint{4.027714in}{5.167673in}}{\pgfqpoint{4.037968in}{5.167673in}}%
\pgfpathlineto{\pgfqpoint{4.037968in}{5.167673in}}%
\pgfpathclose%
\pgfusepath{stroke}%
\end{pgfscope}%
\begin{pgfscope}%
\pgfpathrectangle{\pgfqpoint{0.977909in}{2.928000in}}{\pgfqpoint{5.664000in}{5.664000in}}%
\pgfusepath{clip}%
\pgfsetbuttcap%
\pgfsetroundjoin%
\pgfsetlinewidth{1.003750pt}%
\definecolor{currentstroke}{rgb}{0.541176,0.380392,0.658824}%
\pgfsetstrokecolor{currentstroke}%
\pgfsetstrokeopacity{0.900000}%
\pgfsetdash{}{0pt}%
\pgfpathmoveto{\pgfqpoint{4.211180in}{6.422771in}}%
\pgfpathcurveto{\pgfqpoint{4.221434in}{6.422771in}}{\pgfqpoint{4.231269in}{6.426845in}}{\pgfqpoint{4.238520in}{6.434096in}}%
\pgfpathcurveto{\pgfqpoint{4.245771in}{6.441346in}}{\pgfqpoint{4.249845in}{6.451182in}}{\pgfqpoint{4.249845in}{6.461436in}}%
\pgfpathcurveto{\pgfqpoint{4.249845in}{6.471690in}}{\pgfqpoint{4.245771in}{6.481526in}}{\pgfqpoint{4.238520in}{6.488776in}}%
\pgfpathcurveto{\pgfqpoint{4.231269in}{6.496027in}}{\pgfqpoint{4.221434in}{6.500101in}}{\pgfqpoint{4.211180in}{6.500101in}}%
\pgfpathcurveto{\pgfqpoint{4.200926in}{6.500101in}}{\pgfqpoint{4.191090in}{6.496027in}}{\pgfqpoint{4.183839in}{6.488776in}}%
\pgfpathcurveto{\pgfqpoint{4.176589in}{6.481526in}}{\pgfqpoint{4.172515in}{6.471690in}}{\pgfqpoint{4.172515in}{6.461436in}}%
\pgfpathcurveto{\pgfqpoint{4.172515in}{6.451182in}}{\pgfqpoint{4.176589in}{6.441346in}}{\pgfqpoint{4.183839in}{6.434096in}}%
\pgfpathcurveto{\pgfqpoint{4.191090in}{6.426845in}}{\pgfqpoint{4.200926in}{6.422771in}}{\pgfqpoint{4.211180in}{6.422771in}}%
\pgfpathlineto{\pgfqpoint{4.211180in}{6.422771in}}%
\pgfpathclose%
\pgfusepath{stroke}%
\end{pgfscope}%
\begin{pgfscope}%
\pgfpathrectangle{\pgfqpoint{0.977909in}{2.928000in}}{\pgfqpoint{5.664000in}{5.664000in}}%
\pgfusepath{clip}%
\pgfsetbuttcap%
\pgfsetroundjoin%
\pgfsetlinewidth{1.003750pt}%
\definecolor{currentstroke}{rgb}{0.541176,0.380392,0.658824}%
\pgfsetstrokecolor{currentstroke}%
\pgfsetstrokeopacity{0.900000}%
\pgfsetdash{}{0pt}%
\pgfpathmoveto{\pgfqpoint{4.602626in}{6.293329in}}%
\pgfpathcurveto{\pgfqpoint{4.612881in}{6.293329in}}{\pgfqpoint{4.622716in}{6.297403in}}{\pgfqpoint{4.629967in}{6.304654in}}%
\pgfpathcurveto{\pgfqpoint{4.637217in}{6.311905in}}{\pgfqpoint{4.641291in}{6.321740in}}{\pgfqpoint{4.641291in}{6.331994in}}%
\pgfpathcurveto{\pgfqpoint{4.641291in}{6.342248in}}{\pgfqpoint{4.637217in}{6.352084in}}{\pgfqpoint{4.629967in}{6.359334in}}%
\pgfpathcurveto{\pgfqpoint{4.622716in}{6.366585in}}{\pgfqpoint{4.612881in}{6.370659in}}{\pgfqpoint{4.602626in}{6.370659in}}%
\pgfpathcurveto{\pgfqpoint{4.592372in}{6.370659in}}{\pgfqpoint{4.582537in}{6.366585in}}{\pgfqpoint{4.575286in}{6.359334in}}%
\pgfpathcurveto{\pgfqpoint{4.568035in}{6.352084in}}{\pgfqpoint{4.563961in}{6.342248in}}{\pgfqpoint{4.563961in}{6.331994in}}%
\pgfpathcurveto{\pgfqpoint{4.563961in}{6.321740in}}{\pgfqpoint{4.568035in}{6.311905in}}{\pgfqpoint{4.575286in}{6.304654in}}%
\pgfpathcurveto{\pgfqpoint{4.582537in}{6.297403in}}{\pgfqpoint{4.592372in}{6.293329in}}{\pgfqpoint{4.602626in}{6.293329in}}%
\pgfpathlineto{\pgfqpoint{4.602626in}{6.293329in}}%
\pgfpathclose%
\pgfusepath{stroke}%
\end{pgfscope}%
\begin{pgfscope}%
\pgfpathrectangle{\pgfqpoint{0.977909in}{2.928000in}}{\pgfqpoint{5.664000in}{5.664000in}}%
\pgfusepath{clip}%
\pgfsetbuttcap%
\pgfsetroundjoin%
\pgfsetlinewidth{1.003750pt}%
\definecolor{currentstroke}{rgb}{0.541176,0.380392,0.658824}%
\pgfsetstrokecolor{currentstroke}%
\pgfsetstrokeopacity{0.900000}%
\pgfsetdash{}{0pt}%
\pgfpathmoveto{\pgfqpoint{3.853758in}{5.568276in}}%
\pgfpathcurveto{\pgfqpoint{3.864012in}{5.568276in}}{\pgfqpoint{3.873847in}{5.572350in}}{\pgfqpoint{3.881098in}{5.579601in}}%
\pgfpathcurveto{\pgfqpoint{3.888349in}{5.586851in}}{\pgfqpoint{3.892423in}{5.596687in}}{\pgfqpoint{3.892423in}{5.606941in}}%
\pgfpathcurveto{\pgfqpoint{3.892423in}{5.617195in}}{\pgfqpoint{3.888349in}{5.627031in}}{\pgfqpoint{3.881098in}{5.634281in}}%
\pgfpathcurveto{\pgfqpoint{3.873847in}{5.641532in}}{\pgfqpoint{3.864012in}{5.645606in}}{\pgfqpoint{3.853758in}{5.645606in}}%
\pgfpathcurveto{\pgfqpoint{3.843504in}{5.645606in}}{\pgfqpoint{3.833668in}{5.641532in}}{\pgfqpoint{3.826417in}{5.634281in}}%
\pgfpathcurveto{\pgfqpoint{3.819167in}{5.627031in}}{\pgfqpoint{3.815093in}{5.617195in}}{\pgfqpoint{3.815093in}{5.606941in}}%
\pgfpathcurveto{\pgfqpoint{3.815093in}{5.596687in}}{\pgfqpoint{3.819167in}{5.586851in}}{\pgfqpoint{3.826417in}{5.579601in}}%
\pgfpathcurveto{\pgfqpoint{3.833668in}{5.572350in}}{\pgfqpoint{3.843504in}{5.568276in}}{\pgfqpoint{3.853758in}{5.568276in}}%
\pgfpathlineto{\pgfqpoint{3.853758in}{5.568276in}}%
\pgfpathclose%
\pgfusepath{stroke}%
\end{pgfscope}%
\begin{pgfscope}%
\pgfpathrectangle{\pgfqpoint{0.977909in}{2.928000in}}{\pgfqpoint{5.664000in}{5.664000in}}%
\pgfusepath{clip}%
\pgfsetbuttcap%
\pgfsetroundjoin%
\pgfsetlinewidth{1.003750pt}%
\definecolor{currentstroke}{rgb}{0.541176,0.380392,0.658824}%
\pgfsetstrokecolor{currentstroke}%
\pgfsetstrokeopacity{0.900000}%
\pgfsetdash{}{0pt}%
\pgfpathmoveto{\pgfqpoint{4.764293in}{6.439341in}}%
\pgfpathcurveto{\pgfqpoint{4.774547in}{6.439341in}}{\pgfqpoint{4.784383in}{6.443414in}}{\pgfqpoint{4.791633in}{6.450665in}}%
\pgfpathcurveto{\pgfqpoint{4.798884in}{6.457916in}}{\pgfqpoint{4.802958in}{6.467751in}}{\pgfqpoint{4.802958in}{6.478006in}}%
\pgfpathcurveto{\pgfqpoint{4.802958in}{6.488260in}}{\pgfqpoint{4.798884in}{6.498095in}}{\pgfqpoint{4.791633in}{6.505346in}}%
\pgfpathcurveto{\pgfqpoint{4.784383in}{6.512597in}}{\pgfqpoint{4.774547in}{6.516671in}}{\pgfqpoint{4.764293in}{6.516671in}}%
\pgfpathcurveto{\pgfqpoint{4.754039in}{6.516671in}}{\pgfqpoint{4.744204in}{6.512597in}}{\pgfqpoint{4.736953in}{6.505346in}}%
\pgfpathcurveto{\pgfqpoint{4.729702in}{6.498095in}}{\pgfqpoint{4.725628in}{6.488260in}}{\pgfqpoint{4.725628in}{6.478006in}}%
\pgfpathcurveto{\pgfqpoint{4.725628in}{6.467751in}}{\pgfqpoint{4.729702in}{6.457916in}}{\pgfqpoint{4.736953in}{6.450665in}}%
\pgfpathcurveto{\pgfqpoint{4.744204in}{6.443414in}}{\pgfqpoint{4.754039in}{6.439341in}}{\pgfqpoint{4.764293in}{6.439341in}}%
\pgfpathlineto{\pgfqpoint{4.764293in}{6.439341in}}%
\pgfpathclose%
\pgfusepath{stroke}%
\end{pgfscope}%
\begin{pgfscope}%
\pgfpathrectangle{\pgfqpoint{0.977909in}{2.928000in}}{\pgfqpoint{5.664000in}{5.664000in}}%
\pgfusepath{clip}%
\pgfsetbuttcap%
\pgfsetroundjoin%
\pgfsetlinewidth{1.003750pt}%
\definecolor{currentstroke}{rgb}{0.541176,0.380392,0.658824}%
\pgfsetstrokecolor{currentstroke}%
\pgfsetstrokeopacity{0.900000}%
\pgfsetdash{}{0pt}%
\pgfpathmoveto{\pgfqpoint{4.593846in}{6.402268in}}%
\pgfpathcurveto{\pgfqpoint{4.604100in}{6.402268in}}{\pgfqpoint{4.613935in}{6.406342in}}{\pgfqpoint{4.621186in}{6.413593in}}%
\pgfpathcurveto{\pgfqpoint{4.628437in}{6.420844in}}{\pgfqpoint{4.632511in}{6.430679in}}{\pgfqpoint{4.632511in}{6.440933in}}%
\pgfpathcurveto{\pgfqpoint{4.632511in}{6.451187in}}{\pgfqpoint{4.628437in}{6.461023in}}{\pgfqpoint{4.621186in}{6.468274in}}%
\pgfpathcurveto{\pgfqpoint{4.613935in}{6.475524in}}{\pgfqpoint{4.604100in}{6.479598in}}{\pgfqpoint{4.593846in}{6.479598in}}%
\pgfpathcurveto{\pgfqpoint{4.583591in}{6.479598in}}{\pgfqpoint{4.573756in}{6.475524in}}{\pgfqpoint{4.566505in}{6.468274in}}%
\pgfpathcurveto{\pgfqpoint{4.559254in}{6.461023in}}{\pgfqpoint{4.555180in}{6.451187in}}{\pgfqpoint{4.555180in}{6.440933in}}%
\pgfpathcurveto{\pgfqpoint{4.555180in}{6.430679in}}{\pgfqpoint{4.559254in}{6.420844in}}{\pgfqpoint{4.566505in}{6.413593in}}%
\pgfpathcurveto{\pgfqpoint{4.573756in}{6.406342in}}{\pgfqpoint{4.583591in}{6.402268in}}{\pgfqpoint{4.593846in}{6.402268in}}%
\pgfpathlineto{\pgfqpoint{4.593846in}{6.402268in}}%
\pgfpathclose%
\pgfusepath{stroke}%
\end{pgfscope}%
\begin{pgfscope}%
\pgfpathrectangle{\pgfqpoint{0.977909in}{2.928000in}}{\pgfqpoint{5.664000in}{5.664000in}}%
\pgfusepath{clip}%
\pgfsetbuttcap%
\pgfsetroundjoin%
\pgfsetlinewidth{1.003750pt}%
\definecolor{currentstroke}{rgb}{0.541176,0.380392,0.658824}%
\pgfsetstrokecolor{currentstroke}%
\pgfsetstrokeopacity{0.900000}%
\pgfsetdash{}{0pt}%
\pgfpathmoveto{\pgfqpoint{4.577808in}{5.821761in}}%
\pgfpathcurveto{\pgfqpoint{4.588062in}{5.821761in}}{\pgfqpoint{4.597898in}{5.825835in}}{\pgfqpoint{4.605149in}{5.833086in}}%
\pgfpathcurveto{\pgfqpoint{4.612399in}{5.840337in}}{\pgfqpoint{4.616473in}{5.850172in}}{\pgfqpoint{4.616473in}{5.860426in}}%
\pgfpathcurveto{\pgfqpoint{4.616473in}{5.870680in}}{\pgfqpoint{4.612399in}{5.880516in}}{\pgfqpoint{4.605149in}{5.887766in}}%
\pgfpathcurveto{\pgfqpoint{4.597898in}{5.895017in}}{\pgfqpoint{4.588062in}{5.899091in}}{\pgfqpoint{4.577808in}{5.899091in}}%
\pgfpathcurveto{\pgfqpoint{4.567554in}{5.899091in}}{\pgfqpoint{4.557719in}{5.895017in}}{\pgfqpoint{4.550468in}{5.887766in}}%
\pgfpathcurveto{\pgfqpoint{4.543217in}{5.880516in}}{\pgfqpoint{4.539143in}{5.870680in}}{\pgfqpoint{4.539143in}{5.860426in}}%
\pgfpathcurveto{\pgfqpoint{4.539143in}{5.850172in}}{\pgfqpoint{4.543217in}{5.840337in}}{\pgfqpoint{4.550468in}{5.833086in}}%
\pgfpathcurveto{\pgfqpoint{4.557719in}{5.825835in}}{\pgfqpoint{4.567554in}{5.821761in}}{\pgfqpoint{4.577808in}{5.821761in}}%
\pgfpathlineto{\pgfqpoint{4.577808in}{5.821761in}}%
\pgfpathclose%
\pgfusepath{stroke}%
\end{pgfscope}%
\begin{pgfscope}%
\pgfpathrectangle{\pgfqpoint{0.977909in}{2.928000in}}{\pgfqpoint{5.664000in}{5.664000in}}%
\pgfusepath{clip}%
\pgfsetbuttcap%
\pgfsetroundjoin%
\pgfsetlinewidth{1.003750pt}%
\definecolor{currentstroke}{rgb}{0.541176,0.380392,0.658824}%
\pgfsetstrokecolor{currentstroke}%
\pgfsetstrokeopacity{0.900000}%
\pgfsetdash{}{0pt}%
\pgfpathmoveto{\pgfqpoint{4.790636in}{6.143245in}}%
\pgfpathcurveto{\pgfqpoint{4.800890in}{6.143245in}}{\pgfqpoint{4.810725in}{6.147319in}}{\pgfqpoint{4.817976in}{6.154569in}}%
\pgfpathcurveto{\pgfqpoint{4.825227in}{6.161820in}}{\pgfqpoint{4.829301in}{6.171656in}}{\pgfqpoint{4.829301in}{6.181910in}}%
\pgfpathcurveto{\pgfqpoint{4.829301in}{6.192164in}}{\pgfqpoint{4.825227in}{6.201999in}}{\pgfqpoint{4.817976in}{6.209250in}}%
\pgfpathcurveto{\pgfqpoint{4.810725in}{6.216501in}}{\pgfqpoint{4.800890in}{6.220575in}}{\pgfqpoint{4.790636in}{6.220575in}}%
\pgfpathcurveto{\pgfqpoint{4.780382in}{6.220575in}}{\pgfqpoint{4.770546in}{6.216501in}}{\pgfqpoint{4.763295in}{6.209250in}}%
\pgfpathcurveto{\pgfqpoint{4.756045in}{6.201999in}}{\pgfqpoint{4.751971in}{6.192164in}}{\pgfqpoint{4.751971in}{6.181910in}}%
\pgfpathcurveto{\pgfqpoint{4.751971in}{6.171656in}}{\pgfqpoint{4.756045in}{6.161820in}}{\pgfqpoint{4.763295in}{6.154569in}}%
\pgfpathcurveto{\pgfqpoint{4.770546in}{6.147319in}}{\pgfqpoint{4.780382in}{6.143245in}}{\pgfqpoint{4.790636in}{6.143245in}}%
\pgfpathlineto{\pgfqpoint{4.790636in}{6.143245in}}%
\pgfpathclose%
\pgfusepath{stroke}%
\end{pgfscope}%
\begin{pgfscope}%
\pgfpathrectangle{\pgfqpoint{0.977909in}{2.928000in}}{\pgfqpoint{5.664000in}{5.664000in}}%
\pgfusepath{clip}%
\pgfsetbuttcap%
\pgfsetroundjoin%
\pgfsetlinewidth{1.003750pt}%
\definecolor{currentstroke}{rgb}{0.541176,0.380392,0.658824}%
\pgfsetstrokecolor{currentstroke}%
\pgfsetstrokeopacity{0.900000}%
\pgfsetdash{}{0pt}%
\pgfpathmoveto{\pgfqpoint{6.088101in}{8.056421in}}%
\pgfpathcurveto{\pgfqpoint{6.098355in}{8.056421in}}{\pgfqpoint{6.108190in}{8.060495in}}{\pgfqpoint{6.115441in}{8.067745in}}%
\pgfpathcurveto{\pgfqpoint{6.122692in}{8.074996in}}{\pgfqpoint{6.126766in}{8.084832in}}{\pgfqpoint{6.126766in}{8.095086in}}%
\pgfpathcurveto{\pgfqpoint{6.126766in}{8.105340in}}{\pgfqpoint{6.122692in}{8.115175in}}{\pgfqpoint{6.115441in}{8.122426in}}%
\pgfpathcurveto{\pgfqpoint{6.108190in}{8.129677in}}{\pgfqpoint{6.098355in}{8.133751in}}{\pgfqpoint{6.088101in}{8.133751in}}%
\pgfpathcurveto{\pgfqpoint{6.077847in}{8.133751in}}{\pgfqpoint{6.068011in}{8.129677in}}{\pgfqpoint{6.060761in}{8.122426in}}%
\pgfpathcurveto{\pgfqpoint{6.053510in}{8.115175in}}{\pgfqpoint{6.049436in}{8.105340in}}{\pgfqpoint{6.049436in}{8.095086in}}%
\pgfpathcurveto{\pgfqpoint{6.049436in}{8.084832in}}{\pgfqpoint{6.053510in}{8.074996in}}{\pgfqpoint{6.060761in}{8.067745in}}%
\pgfpathcurveto{\pgfqpoint{6.068011in}{8.060495in}}{\pgfqpoint{6.077847in}{8.056421in}}{\pgfqpoint{6.088101in}{8.056421in}}%
\pgfpathlineto{\pgfqpoint{6.088101in}{8.056421in}}%
\pgfpathclose%
\pgfusepath{stroke}%
\end{pgfscope}%
\begin{pgfscope}%
\pgfpathrectangle{\pgfqpoint{0.977909in}{2.928000in}}{\pgfqpoint{5.664000in}{5.664000in}}%
\pgfusepath{clip}%
\pgfsetbuttcap%
\pgfsetroundjoin%
\pgfsetlinewidth{1.003750pt}%
\definecolor{currentstroke}{rgb}{0.541176,0.380392,0.658824}%
\pgfsetstrokecolor{currentstroke}%
\pgfsetstrokeopacity{0.900000}%
\pgfsetdash{}{0pt}%
\pgfpathmoveto{\pgfqpoint{4.034079in}{5.886271in}}%
\pgfpathcurveto{\pgfqpoint{4.044333in}{5.886271in}}{\pgfqpoint{4.054168in}{5.890345in}}{\pgfqpoint{4.061419in}{5.897596in}}%
\pgfpathcurveto{\pgfqpoint{4.068670in}{5.904847in}}{\pgfqpoint{4.072744in}{5.914682in}}{\pgfqpoint{4.072744in}{5.924936in}}%
\pgfpathcurveto{\pgfqpoint{4.072744in}{5.935190in}}{\pgfqpoint{4.068670in}{5.945026in}}{\pgfqpoint{4.061419in}{5.952277in}}%
\pgfpathcurveto{\pgfqpoint{4.054168in}{5.959527in}}{\pgfqpoint{4.044333in}{5.963601in}}{\pgfqpoint{4.034079in}{5.963601in}}%
\pgfpathcurveto{\pgfqpoint{4.023825in}{5.963601in}}{\pgfqpoint{4.013989in}{5.959527in}}{\pgfqpoint{4.006738in}{5.952277in}}%
\pgfpathcurveto{\pgfqpoint{3.999488in}{5.945026in}}{\pgfqpoint{3.995414in}{5.935190in}}{\pgfqpoint{3.995414in}{5.924936in}}%
\pgfpathcurveto{\pgfqpoint{3.995414in}{5.914682in}}{\pgfqpoint{3.999488in}{5.904847in}}{\pgfqpoint{4.006738in}{5.897596in}}%
\pgfpathcurveto{\pgfqpoint{4.013989in}{5.890345in}}{\pgfqpoint{4.023825in}{5.886271in}}{\pgfqpoint{4.034079in}{5.886271in}}%
\pgfpathlineto{\pgfqpoint{4.034079in}{5.886271in}}%
\pgfpathclose%
\pgfusepath{stroke}%
\end{pgfscope}%
\begin{pgfscope}%
\pgfpathrectangle{\pgfqpoint{0.977909in}{2.928000in}}{\pgfqpoint{5.664000in}{5.664000in}}%
\pgfusepath{clip}%
\pgfsetbuttcap%
\pgfsetroundjoin%
\pgfsetlinewidth{1.003750pt}%
\definecolor{currentstroke}{rgb}{0.541176,0.380392,0.658824}%
\pgfsetstrokecolor{currentstroke}%
\pgfsetstrokeopacity{0.900000}%
\pgfsetdash{}{0pt}%
\pgfpathmoveto{\pgfqpoint{4.439546in}{6.724045in}}%
\pgfpathcurveto{\pgfqpoint{4.449800in}{6.724045in}}{\pgfqpoint{4.459635in}{6.728119in}}{\pgfqpoint{4.466886in}{6.735369in}}%
\pgfpathcurveto{\pgfqpoint{4.474137in}{6.742620in}}{\pgfqpoint{4.478211in}{6.752456in}}{\pgfqpoint{4.478211in}{6.762710in}}%
\pgfpathcurveto{\pgfqpoint{4.478211in}{6.772964in}}{\pgfqpoint{4.474137in}{6.782799in}}{\pgfqpoint{4.466886in}{6.790050in}}%
\pgfpathcurveto{\pgfqpoint{4.459635in}{6.797301in}}{\pgfqpoint{4.449800in}{6.801375in}}{\pgfqpoint{4.439546in}{6.801375in}}%
\pgfpathcurveto{\pgfqpoint{4.429292in}{6.801375in}}{\pgfqpoint{4.419456in}{6.797301in}}{\pgfqpoint{4.412205in}{6.790050in}}%
\pgfpathcurveto{\pgfqpoint{4.404955in}{6.782799in}}{\pgfqpoint{4.400881in}{6.772964in}}{\pgfqpoint{4.400881in}{6.762710in}}%
\pgfpathcurveto{\pgfqpoint{4.400881in}{6.752456in}}{\pgfqpoint{4.404955in}{6.742620in}}{\pgfqpoint{4.412205in}{6.735369in}}%
\pgfpathcurveto{\pgfqpoint{4.419456in}{6.728119in}}{\pgfqpoint{4.429292in}{6.724045in}}{\pgfqpoint{4.439546in}{6.724045in}}%
\pgfpathlineto{\pgfqpoint{4.439546in}{6.724045in}}%
\pgfpathclose%
\pgfusepath{stroke}%
\end{pgfscope}%
\begin{pgfscope}%
\pgfpathrectangle{\pgfqpoint{0.977909in}{2.928000in}}{\pgfqpoint{5.664000in}{5.664000in}}%
\pgfusepath{clip}%
\pgfsetbuttcap%
\pgfsetroundjoin%
\pgfsetlinewidth{1.003750pt}%
\definecolor{currentstroke}{rgb}{0.541176,0.380392,0.658824}%
\pgfsetstrokecolor{currentstroke}%
\pgfsetstrokeopacity{0.900000}%
\pgfsetdash{}{0pt}%
\pgfpathmoveto{\pgfqpoint{4.558086in}{6.358854in}}%
\pgfpathcurveto{\pgfqpoint{4.568340in}{6.358854in}}{\pgfqpoint{4.578176in}{6.362928in}}{\pgfqpoint{4.585427in}{6.370179in}}%
\pgfpathcurveto{\pgfqpoint{4.592677in}{6.377430in}}{\pgfqpoint{4.596751in}{6.387265in}}{\pgfqpoint{4.596751in}{6.397520in}}%
\pgfpathcurveto{\pgfqpoint{4.596751in}{6.407774in}}{\pgfqpoint{4.592677in}{6.417609in}}{\pgfqpoint{4.585427in}{6.424860in}}%
\pgfpathcurveto{\pgfqpoint{4.578176in}{6.432111in}}{\pgfqpoint{4.568340in}{6.436185in}}{\pgfqpoint{4.558086in}{6.436185in}}%
\pgfpathcurveto{\pgfqpoint{4.547832in}{6.436185in}}{\pgfqpoint{4.537997in}{6.432111in}}{\pgfqpoint{4.530746in}{6.424860in}}%
\pgfpathcurveto{\pgfqpoint{4.523495in}{6.417609in}}{\pgfqpoint{4.519421in}{6.407774in}}{\pgfqpoint{4.519421in}{6.397520in}}%
\pgfpathcurveto{\pgfqpoint{4.519421in}{6.387265in}}{\pgfqpoint{4.523495in}{6.377430in}}{\pgfqpoint{4.530746in}{6.370179in}}%
\pgfpathcurveto{\pgfqpoint{4.537997in}{6.362928in}}{\pgfqpoint{4.547832in}{6.358854in}}{\pgfqpoint{4.558086in}{6.358854in}}%
\pgfpathlineto{\pgfqpoint{4.558086in}{6.358854in}}%
\pgfpathclose%
\pgfusepath{stroke}%
\end{pgfscope}%
\begin{pgfscope}%
\pgfpathrectangle{\pgfqpoint{0.977909in}{2.928000in}}{\pgfqpoint{5.664000in}{5.664000in}}%
\pgfusepath{clip}%
\pgfsetbuttcap%
\pgfsetroundjoin%
\pgfsetlinewidth{1.003750pt}%
\definecolor{currentstroke}{rgb}{0.541176,0.380392,0.658824}%
\pgfsetstrokecolor{currentstroke}%
\pgfsetstrokeopacity{0.900000}%
\pgfsetdash{}{0pt}%
\pgfpathmoveto{\pgfqpoint{5.152360in}{7.016782in}}%
\pgfpathcurveto{\pgfqpoint{5.162614in}{7.016782in}}{\pgfqpoint{5.172449in}{7.020856in}}{\pgfqpoint{5.179700in}{7.028107in}}%
\pgfpathcurveto{\pgfqpoint{5.186951in}{7.035357in}}{\pgfqpoint{5.191025in}{7.045193in}}{\pgfqpoint{5.191025in}{7.055447in}}%
\pgfpathcurveto{\pgfqpoint{5.191025in}{7.065701in}}{\pgfqpoint{5.186951in}{7.075536in}}{\pgfqpoint{5.179700in}{7.082787in}}%
\pgfpathcurveto{\pgfqpoint{5.172449in}{7.090038in}}{\pgfqpoint{5.162614in}{7.094112in}}{\pgfqpoint{5.152360in}{7.094112in}}%
\pgfpathcurveto{\pgfqpoint{5.142106in}{7.094112in}}{\pgfqpoint{5.132270in}{7.090038in}}{\pgfqpoint{5.125020in}{7.082787in}}%
\pgfpathcurveto{\pgfqpoint{5.117769in}{7.075536in}}{\pgfqpoint{5.113695in}{7.065701in}}{\pgfqpoint{5.113695in}{7.055447in}}%
\pgfpathcurveto{\pgfqpoint{5.113695in}{7.045193in}}{\pgfqpoint{5.117769in}{7.035357in}}{\pgfqpoint{5.125020in}{7.028107in}}%
\pgfpathcurveto{\pgfqpoint{5.132270in}{7.020856in}}{\pgfqpoint{5.142106in}{7.016782in}}{\pgfqpoint{5.152360in}{7.016782in}}%
\pgfpathlineto{\pgfqpoint{5.152360in}{7.016782in}}%
\pgfpathclose%
\pgfusepath{stroke}%
\end{pgfscope}%
\begin{pgfscope}%
\pgfpathrectangle{\pgfqpoint{0.977909in}{2.928000in}}{\pgfqpoint{5.664000in}{5.664000in}}%
\pgfusepath{clip}%
\pgfsetbuttcap%
\pgfsetroundjoin%
\pgfsetlinewidth{1.003750pt}%
\definecolor{currentstroke}{rgb}{0.541176,0.380392,0.658824}%
\pgfsetstrokecolor{currentstroke}%
\pgfsetstrokeopacity{0.900000}%
\pgfsetdash{}{0pt}%
\pgfpathmoveto{\pgfqpoint{5.343651in}{7.342307in}}%
\pgfpathcurveto{\pgfqpoint{5.353905in}{7.342307in}}{\pgfqpoint{5.363741in}{7.346381in}}{\pgfqpoint{5.370992in}{7.353631in}}%
\pgfpathcurveto{\pgfqpoint{5.378242in}{7.360882in}}{\pgfqpoint{5.382316in}{7.370718in}}{\pgfqpoint{5.382316in}{7.380972in}}%
\pgfpathcurveto{\pgfqpoint{5.382316in}{7.391226in}}{\pgfqpoint{5.378242in}{7.401061in}}{\pgfqpoint{5.370992in}{7.408312in}}%
\pgfpathcurveto{\pgfqpoint{5.363741in}{7.415563in}}{\pgfqpoint{5.353905in}{7.419637in}}{\pgfqpoint{5.343651in}{7.419637in}}%
\pgfpathcurveto{\pgfqpoint{5.333397in}{7.419637in}}{\pgfqpoint{5.323562in}{7.415563in}}{\pgfqpoint{5.316311in}{7.408312in}}%
\pgfpathcurveto{\pgfqpoint{5.309060in}{7.401061in}}{\pgfqpoint{5.304986in}{7.391226in}}{\pgfqpoint{5.304986in}{7.380972in}}%
\pgfpathcurveto{\pgfqpoint{5.304986in}{7.370718in}}{\pgfqpoint{5.309060in}{7.360882in}}{\pgfqpoint{5.316311in}{7.353631in}}%
\pgfpathcurveto{\pgfqpoint{5.323562in}{7.346381in}}{\pgfqpoint{5.333397in}{7.342307in}}{\pgfqpoint{5.343651in}{7.342307in}}%
\pgfpathlineto{\pgfqpoint{5.343651in}{7.342307in}}%
\pgfpathclose%
\pgfusepath{stroke}%
\end{pgfscope}%
\begin{pgfscope}%
\pgfpathrectangle{\pgfqpoint{0.977909in}{2.928000in}}{\pgfqpoint{5.664000in}{5.664000in}}%
\pgfusepath{clip}%
\pgfsetbuttcap%
\pgfsetroundjoin%
\pgfsetlinewidth{1.003750pt}%
\definecolor{currentstroke}{rgb}{0.541176,0.380392,0.658824}%
\pgfsetstrokecolor{currentstroke}%
\pgfsetstrokeopacity{0.900000}%
\pgfsetdash{}{0pt}%
\pgfpathmoveto{\pgfqpoint{4.694815in}{7.204501in}}%
\pgfpathcurveto{\pgfqpoint{4.705069in}{7.204501in}}{\pgfqpoint{4.714905in}{7.208575in}}{\pgfqpoint{4.722155in}{7.215826in}}%
\pgfpathcurveto{\pgfqpoint{4.729406in}{7.223077in}}{\pgfqpoint{4.733480in}{7.232912in}}{\pgfqpoint{4.733480in}{7.243166in}}%
\pgfpathcurveto{\pgfqpoint{4.733480in}{7.253421in}}{\pgfqpoint{4.729406in}{7.263256in}}{\pgfqpoint{4.722155in}{7.270507in}}%
\pgfpathcurveto{\pgfqpoint{4.714905in}{7.277757in}}{\pgfqpoint{4.705069in}{7.281831in}}{\pgfqpoint{4.694815in}{7.281831in}}%
\pgfpathcurveto{\pgfqpoint{4.684561in}{7.281831in}}{\pgfqpoint{4.674725in}{7.277757in}}{\pgfqpoint{4.667475in}{7.270507in}}%
\pgfpathcurveto{\pgfqpoint{4.660224in}{7.263256in}}{\pgfqpoint{4.656150in}{7.253421in}}{\pgfqpoint{4.656150in}{7.243166in}}%
\pgfpathcurveto{\pgfqpoint{4.656150in}{7.232912in}}{\pgfqpoint{4.660224in}{7.223077in}}{\pgfqpoint{4.667475in}{7.215826in}}%
\pgfpathcurveto{\pgfqpoint{4.674725in}{7.208575in}}{\pgfqpoint{4.684561in}{7.204501in}}{\pgfqpoint{4.694815in}{7.204501in}}%
\pgfpathlineto{\pgfqpoint{4.694815in}{7.204501in}}%
\pgfpathclose%
\pgfusepath{stroke}%
\end{pgfscope}%
\begin{pgfscope}%
\pgfpathrectangle{\pgfqpoint{0.977909in}{2.928000in}}{\pgfqpoint{5.664000in}{5.664000in}}%
\pgfusepath{clip}%
\pgfsetbuttcap%
\pgfsetroundjoin%
\pgfsetlinewidth{1.003750pt}%
\definecolor{currentstroke}{rgb}{0.541176,0.380392,0.658824}%
\pgfsetstrokecolor{currentstroke}%
\pgfsetstrokeopacity{0.900000}%
\pgfsetdash{}{0pt}%
\pgfpathmoveto{\pgfqpoint{4.601604in}{7.088276in}}%
\pgfpathcurveto{\pgfqpoint{4.611858in}{7.088276in}}{\pgfqpoint{4.621694in}{7.092350in}}{\pgfqpoint{4.628945in}{7.099601in}}%
\pgfpathcurveto{\pgfqpoint{4.636195in}{7.106851in}}{\pgfqpoint{4.640269in}{7.116687in}}{\pgfqpoint{4.640269in}{7.126941in}}%
\pgfpathcurveto{\pgfqpoint{4.640269in}{7.137195in}}{\pgfqpoint{4.636195in}{7.147030in}}{\pgfqpoint{4.628945in}{7.154281in}}%
\pgfpathcurveto{\pgfqpoint{4.621694in}{7.161532in}}{\pgfqpoint{4.611858in}{7.165606in}}{\pgfqpoint{4.601604in}{7.165606in}}%
\pgfpathcurveto{\pgfqpoint{4.591350in}{7.165606in}}{\pgfqpoint{4.581515in}{7.161532in}}{\pgfqpoint{4.574264in}{7.154281in}}%
\pgfpathcurveto{\pgfqpoint{4.567013in}{7.147030in}}{\pgfqpoint{4.562939in}{7.137195in}}{\pgfqpoint{4.562939in}{7.126941in}}%
\pgfpathcurveto{\pgfqpoint{4.562939in}{7.116687in}}{\pgfqpoint{4.567013in}{7.106851in}}{\pgfqpoint{4.574264in}{7.099601in}}%
\pgfpathcurveto{\pgfqpoint{4.581515in}{7.092350in}}{\pgfqpoint{4.591350in}{7.088276in}}{\pgfqpoint{4.601604in}{7.088276in}}%
\pgfpathlineto{\pgfqpoint{4.601604in}{7.088276in}}%
\pgfpathclose%
\pgfusepath{stroke}%
\end{pgfscope}%
\begin{pgfscope}%
\pgfpathrectangle{\pgfqpoint{0.977909in}{2.928000in}}{\pgfqpoint{5.664000in}{5.664000in}}%
\pgfusepath{clip}%
\pgfsetbuttcap%
\pgfsetroundjoin%
\pgfsetlinewidth{1.003750pt}%
\definecolor{currentstroke}{rgb}{0.541176,0.380392,0.658824}%
\pgfsetstrokecolor{currentstroke}%
\pgfsetstrokeopacity{0.900000}%
\pgfsetdash{}{0pt}%
\pgfpathmoveto{\pgfqpoint{3.791508in}{5.801385in}}%
\pgfpathcurveto{\pgfqpoint{3.801762in}{5.801385in}}{\pgfqpoint{3.811598in}{5.805459in}}{\pgfqpoint{3.818848in}{5.812710in}}%
\pgfpathcurveto{\pgfqpoint{3.826099in}{5.819960in}}{\pgfqpoint{3.830173in}{5.829796in}}{\pgfqpoint{3.830173in}{5.840050in}}%
\pgfpathcurveto{\pgfqpoint{3.830173in}{5.850304in}}{\pgfqpoint{3.826099in}{5.860140in}}{\pgfqpoint{3.818848in}{5.867390in}}%
\pgfpathcurveto{\pgfqpoint{3.811598in}{5.874641in}}{\pgfqpoint{3.801762in}{5.878715in}}{\pgfqpoint{3.791508in}{5.878715in}}%
\pgfpathcurveto{\pgfqpoint{3.781254in}{5.878715in}}{\pgfqpoint{3.771418in}{5.874641in}}{\pgfqpoint{3.764168in}{5.867390in}}%
\pgfpathcurveto{\pgfqpoint{3.756917in}{5.860140in}}{\pgfqpoint{3.752843in}{5.850304in}}{\pgfqpoint{3.752843in}{5.840050in}}%
\pgfpathcurveto{\pgfqpoint{3.752843in}{5.829796in}}{\pgfqpoint{3.756917in}{5.819960in}}{\pgfqpoint{3.764168in}{5.812710in}}%
\pgfpathcurveto{\pgfqpoint{3.771418in}{5.805459in}}{\pgfqpoint{3.781254in}{5.801385in}}{\pgfqpoint{3.791508in}{5.801385in}}%
\pgfpathlineto{\pgfqpoint{3.791508in}{5.801385in}}%
\pgfpathclose%
\pgfusepath{stroke}%
\end{pgfscope}%
\begin{pgfscope}%
\pgfpathrectangle{\pgfqpoint{0.977909in}{2.928000in}}{\pgfqpoint{5.664000in}{5.664000in}}%
\pgfusepath{clip}%
\pgfsetbuttcap%
\pgfsetroundjoin%
\pgfsetlinewidth{1.003750pt}%
\definecolor{currentstroke}{rgb}{0.541176,0.380392,0.658824}%
\pgfsetstrokecolor{currentstroke}%
\pgfsetstrokeopacity{0.900000}%
\pgfsetdash{}{0pt}%
\pgfpathmoveto{\pgfqpoint{4.605556in}{6.297182in}}%
\pgfpathcurveto{\pgfqpoint{4.615810in}{6.297182in}}{\pgfqpoint{4.625646in}{6.301256in}}{\pgfqpoint{4.632897in}{6.308507in}}%
\pgfpathcurveto{\pgfqpoint{4.640147in}{6.315758in}}{\pgfqpoint{4.644221in}{6.325593in}}{\pgfqpoint{4.644221in}{6.335848in}}%
\pgfpathcurveto{\pgfqpoint{4.644221in}{6.346102in}}{\pgfqpoint{4.640147in}{6.355937in}}{\pgfqpoint{4.632897in}{6.363188in}}%
\pgfpathcurveto{\pgfqpoint{4.625646in}{6.370439in}}{\pgfqpoint{4.615810in}{6.374513in}}{\pgfqpoint{4.605556in}{6.374513in}}%
\pgfpathcurveto{\pgfqpoint{4.595302in}{6.374513in}}{\pgfqpoint{4.585467in}{6.370439in}}{\pgfqpoint{4.578216in}{6.363188in}}%
\pgfpathcurveto{\pgfqpoint{4.570965in}{6.355937in}}{\pgfqpoint{4.566891in}{6.346102in}}{\pgfqpoint{4.566891in}{6.335848in}}%
\pgfpathcurveto{\pgfqpoint{4.566891in}{6.325593in}}{\pgfqpoint{4.570965in}{6.315758in}}{\pgfqpoint{4.578216in}{6.308507in}}%
\pgfpathcurveto{\pgfqpoint{4.585467in}{6.301256in}}{\pgfqpoint{4.595302in}{6.297182in}}{\pgfqpoint{4.605556in}{6.297182in}}%
\pgfpathlineto{\pgfqpoint{4.605556in}{6.297182in}}%
\pgfpathclose%
\pgfusepath{stroke}%
\end{pgfscope}%
\begin{pgfscope}%
\pgfpathrectangle{\pgfqpoint{0.977909in}{2.928000in}}{\pgfqpoint{5.664000in}{5.664000in}}%
\pgfusepath{clip}%
\pgfsetbuttcap%
\pgfsetroundjoin%
\pgfsetlinewidth{1.003750pt}%
\definecolor{currentstroke}{rgb}{0.541176,0.380392,0.658824}%
\pgfsetstrokecolor{currentstroke}%
\pgfsetstrokeopacity{0.900000}%
\pgfsetdash{}{0pt}%
\pgfpathmoveto{\pgfqpoint{3.549283in}{6.020549in}}%
\pgfpathcurveto{\pgfqpoint{3.559537in}{6.020549in}}{\pgfqpoint{3.569372in}{6.024623in}}{\pgfqpoint{3.576623in}{6.031874in}}%
\pgfpathcurveto{\pgfqpoint{3.583874in}{6.039124in}}{\pgfqpoint{3.587948in}{6.048960in}}{\pgfqpoint{3.587948in}{6.059214in}}%
\pgfpathcurveto{\pgfqpoint{3.587948in}{6.069468in}}{\pgfqpoint{3.583874in}{6.079303in}}{\pgfqpoint{3.576623in}{6.086554in}}%
\pgfpathcurveto{\pgfqpoint{3.569372in}{6.093805in}}{\pgfqpoint{3.559537in}{6.097879in}}{\pgfqpoint{3.549283in}{6.097879in}}%
\pgfpathcurveto{\pgfqpoint{3.539029in}{6.097879in}}{\pgfqpoint{3.529193in}{6.093805in}}{\pgfqpoint{3.521942in}{6.086554in}}%
\pgfpathcurveto{\pgfqpoint{3.514692in}{6.079303in}}{\pgfqpoint{3.510618in}{6.069468in}}{\pgfqpoint{3.510618in}{6.059214in}}%
\pgfpathcurveto{\pgfqpoint{3.510618in}{6.048960in}}{\pgfqpoint{3.514692in}{6.039124in}}{\pgfqpoint{3.521942in}{6.031874in}}%
\pgfpathcurveto{\pgfqpoint{3.529193in}{6.024623in}}{\pgfqpoint{3.539029in}{6.020549in}}{\pgfqpoint{3.549283in}{6.020549in}}%
\pgfpathlineto{\pgfqpoint{3.549283in}{6.020549in}}%
\pgfpathclose%
\pgfusepath{stroke}%
\end{pgfscope}%
\begin{pgfscope}%
\pgfpathrectangle{\pgfqpoint{0.977909in}{2.928000in}}{\pgfqpoint{5.664000in}{5.664000in}}%
\pgfusepath{clip}%
\pgfsetbuttcap%
\pgfsetroundjoin%
\pgfsetlinewidth{1.003750pt}%
\definecolor{currentstroke}{rgb}{0.541176,0.380392,0.658824}%
\pgfsetstrokecolor{currentstroke}%
\pgfsetstrokeopacity{0.900000}%
\pgfsetdash{}{0pt}%
\pgfpathmoveto{\pgfqpoint{4.690315in}{6.939896in}}%
\pgfpathcurveto{\pgfqpoint{4.700569in}{6.939896in}}{\pgfqpoint{4.710404in}{6.943970in}}{\pgfqpoint{4.717655in}{6.951221in}}%
\pgfpathcurveto{\pgfqpoint{4.724906in}{6.958472in}}{\pgfqpoint{4.728980in}{6.968307in}}{\pgfqpoint{4.728980in}{6.978561in}}%
\pgfpathcurveto{\pgfqpoint{4.728980in}{6.988816in}}{\pgfqpoint{4.724906in}{6.998651in}}{\pgfqpoint{4.717655in}{7.005902in}}%
\pgfpathcurveto{\pgfqpoint{4.710404in}{7.013152in}}{\pgfqpoint{4.700569in}{7.017226in}}{\pgfqpoint{4.690315in}{7.017226in}}%
\pgfpathcurveto{\pgfqpoint{4.680060in}{7.017226in}}{\pgfqpoint{4.670225in}{7.013152in}}{\pgfqpoint{4.662974in}{7.005902in}}%
\pgfpathcurveto{\pgfqpoint{4.655724in}{6.998651in}}{\pgfqpoint{4.651650in}{6.988816in}}{\pgfqpoint{4.651650in}{6.978561in}}%
\pgfpathcurveto{\pgfqpoint{4.651650in}{6.968307in}}{\pgfqpoint{4.655724in}{6.958472in}}{\pgfqpoint{4.662974in}{6.951221in}}%
\pgfpathcurveto{\pgfqpoint{4.670225in}{6.943970in}}{\pgfqpoint{4.680060in}{6.939896in}}{\pgfqpoint{4.690315in}{6.939896in}}%
\pgfpathlineto{\pgfqpoint{4.690315in}{6.939896in}}%
\pgfpathclose%
\pgfusepath{stroke}%
\end{pgfscope}%
\begin{pgfscope}%
\pgfpathrectangle{\pgfqpoint{0.977909in}{2.928000in}}{\pgfqpoint{5.664000in}{5.664000in}}%
\pgfusepath{clip}%
\pgfsetbuttcap%
\pgfsetroundjoin%
\pgfsetlinewidth{1.003750pt}%
\definecolor{currentstroke}{rgb}{0.541176,0.380392,0.658824}%
\pgfsetstrokecolor{currentstroke}%
\pgfsetstrokeopacity{0.900000}%
\pgfsetdash{}{0pt}%
\pgfpathmoveto{\pgfqpoint{4.385248in}{6.616304in}}%
\pgfpathcurveto{\pgfqpoint{4.395502in}{6.616304in}}{\pgfqpoint{4.405337in}{6.620378in}}{\pgfqpoint{4.412588in}{6.627629in}}%
\pgfpathcurveto{\pgfqpoint{4.419839in}{6.634880in}}{\pgfqpoint{4.423913in}{6.644715in}}{\pgfqpoint{4.423913in}{6.654969in}}%
\pgfpathcurveto{\pgfqpoint{4.423913in}{6.665223in}}{\pgfqpoint{4.419839in}{6.675059in}}{\pgfqpoint{4.412588in}{6.682310in}}%
\pgfpathcurveto{\pgfqpoint{4.405337in}{6.689560in}}{\pgfqpoint{4.395502in}{6.693634in}}{\pgfqpoint{4.385248in}{6.693634in}}%
\pgfpathcurveto{\pgfqpoint{4.374993in}{6.693634in}}{\pgfqpoint{4.365158in}{6.689560in}}{\pgfqpoint{4.357907in}{6.682310in}}%
\pgfpathcurveto{\pgfqpoint{4.350656in}{6.675059in}}{\pgfqpoint{4.346582in}{6.665223in}}{\pgfqpoint{4.346582in}{6.654969in}}%
\pgfpathcurveto{\pgfqpoint{4.346582in}{6.644715in}}{\pgfqpoint{4.350656in}{6.634880in}}{\pgfqpoint{4.357907in}{6.627629in}}%
\pgfpathcurveto{\pgfqpoint{4.365158in}{6.620378in}}{\pgfqpoint{4.374993in}{6.616304in}}{\pgfqpoint{4.385248in}{6.616304in}}%
\pgfpathlineto{\pgfqpoint{4.385248in}{6.616304in}}%
\pgfpathclose%
\pgfusepath{stroke}%
\end{pgfscope}%
\begin{pgfscope}%
\pgfpathrectangle{\pgfqpoint{0.977909in}{2.928000in}}{\pgfqpoint{5.664000in}{5.664000in}}%
\pgfusepath{clip}%
\pgfsetbuttcap%
\pgfsetroundjoin%
\pgfsetlinewidth{1.003750pt}%
\definecolor{currentstroke}{rgb}{0.541176,0.380392,0.658824}%
\pgfsetstrokecolor{currentstroke}%
\pgfsetstrokeopacity{0.900000}%
\pgfsetdash{}{0pt}%
\pgfpathmoveto{\pgfqpoint{4.320617in}{5.562815in}}%
\pgfpathcurveto{\pgfqpoint{4.330871in}{5.562815in}}{\pgfqpoint{4.340706in}{5.566889in}}{\pgfqpoint{4.347957in}{5.574140in}}%
\pgfpathcurveto{\pgfqpoint{4.355208in}{5.581391in}}{\pgfqpoint{4.359282in}{5.591226in}}{\pgfqpoint{4.359282in}{5.601480in}}%
\pgfpathcurveto{\pgfqpoint{4.359282in}{5.611734in}}{\pgfqpoint{4.355208in}{5.621570in}}{\pgfqpoint{4.347957in}{5.628820in}}%
\pgfpathcurveto{\pgfqpoint{4.340706in}{5.636071in}}{\pgfqpoint{4.330871in}{5.640145in}}{\pgfqpoint{4.320617in}{5.640145in}}%
\pgfpathcurveto{\pgfqpoint{4.310363in}{5.640145in}}{\pgfqpoint{4.300527in}{5.636071in}}{\pgfqpoint{4.293277in}{5.628820in}}%
\pgfpathcurveto{\pgfqpoint{4.286026in}{5.621570in}}{\pgfqpoint{4.281952in}{5.611734in}}{\pgfqpoint{4.281952in}{5.601480in}}%
\pgfpathcurveto{\pgfqpoint{4.281952in}{5.591226in}}{\pgfqpoint{4.286026in}{5.581391in}}{\pgfqpoint{4.293277in}{5.574140in}}%
\pgfpathcurveto{\pgfqpoint{4.300527in}{5.566889in}}{\pgfqpoint{4.310363in}{5.562815in}}{\pgfqpoint{4.320617in}{5.562815in}}%
\pgfpathlineto{\pgfqpoint{4.320617in}{5.562815in}}%
\pgfpathclose%
\pgfusepath{stroke}%
\end{pgfscope}%
\begin{pgfscope}%
\pgfpathrectangle{\pgfqpoint{0.977909in}{2.928000in}}{\pgfqpoint{5.664000in}{5.664000in}}%
\pgfusepath{clip}%
\pgfsetbuttcap%
\pgfsetroundjoin%
\pgfsetlinewidth{1.003750pt}%
\definecolor{currentstroke}{rgb}{0.541176,0.380392,0.658824}%
\pgfsetstrokecolor{currentstroke}%
\pgfsetstrokeopacity{0.900000}%
\pgfsetdash{}{0pt}%
\pgfpathmoveto{\pgfqpoint{3.061608in}{4.418055in}}%
\pgfpathcurveto{\pgfqpoint{3.071862in}{4.418055in}}{\pgfqpoint{3.081698in}{4.422129in}}{\pgfqpoint{3.088948in}{4.429380in}}%
\pgfpathcurveto{\pgfqpoint{3.096199in}{4.436631in}}{\pgfqpoint{3.100273in}{4.446466in}}{\pgfqpoint{3.100273in}{4.456720in}}%
\pgfpathcurveto{\pgfqpoint{3.100273in}{4.466974in}}{\pgfqpoint{3.096199in}{4.476810in}}{\pgfqpoint{3.088948in}{4.484060in}}%
\pgfpathcurveto{\pgfqpoint{3.081698in}{4.491311in}}{\pgfqpoint{3.071862in}{4.495385in}}{\pgfqpoint{3.061608in}{4.495385in}}%
\pgfpathcurveto{\pgfqpoint{3.051354in}{4.495385in}}{\pgfqpoint{3.041519in}{4.491311in}}{\pgfqpoint{3.034268in}{4.484060in}}%
\pgfpathcurveto{\pgfqpoint{3.027017in}{4.476810in}}{\pgfqpoint{3.022943in}{4.466974in}}{\pgfqpoint{3.022943in}{4.456720in}}%
\pgfpathcurveto{\pgfqpoint{3.022943in}{4.446466in}}{\pgfqpoint{3.027017in}{4.436631in}}{\pgfqpoint{3.034268in}{4.429380in}}%
\pgfpathcurveto{\pgfqpoint{3.041519in}{4.422129in}}{\pgfqpoint{3.051354in}{4.418055in}}{\pgfqpoint{3.061608in}{4.418055in}}%
\pgfpathlineto{\pgfqpoint{3.061608in}{4.418055in}}%
\pgfpathclose%
\pgfusepath{stroke}%
\end{pgfscope}%
\begin{pgfscope}%
\pgfpathrectangle{\pgfqpoint{0.977909in}{2.928000in}}{\pgfqpoint{5.664000in}{5.664000in}}%
\pgfusepath{clip}%
\pgfsetbuttcap%
\pgfsetroundjoin%
\pgfsetlinewidth{1.003750pt}%
\definecolor{currentstroke}{rgb}{0.541176,0.380392,0.658824}%
\pgfsetstrokecolor{currentstroke}%
\pgfsetstrokeopacity{0.900000}%
\pgfsetdash{}{0pt}%
\pgfpathmoveto{\pgfqpoint{3.734344in}{5.489765in}}%
\pgfpathcurveto{\pgfqpoint{3.744598in}{5.489765in}}{\pgfqpoint{3.754434in}{5.493839in}}{\pgfqpoint{3.761684in}{5.501090in}}%
\pgfpathcurveto{\pgfqpoint{3.768935in}{5.508341in}}{\pgfqpoint{3.773009in}{5.518176in}}{\pgfqpoint{3.773009in}{5.528430in}}%
\pgfpathcurveto{\pgfqpoint{3.773009in}{5.538684in}}{\pgfqpoint{3.768935in}{5.548520in}}{\pgfqpoint{3.761684in}{5.555771in}}%
\pgfpathcurveto{\pgfqpoint{3.754434in}{5.563021in}}{\pgfqpoint{3.744598in}{5.567095in}}{\pgfqpoint{3.734344in}{5.567095in}}%
\pgfpathcurveto{\pgfqpoint{3.724090in}{5.567095in}}{\pgfqpoint{3.714254in}{5.563021in}}{\pgfqpoint{3.707004in}{5.555771in}}%
\pgfpathcurveto{\pgfqpoint{3.699753in}{5.548520in}}{\pgfqpoint{3.695679in}{5.538684in}}{\pgfqpoint{3.695679in}{5.528430in}}%
\pgfpathcurveto{\pgfqpoint{3.695679in}{5.518176in}}{\pgfqpoint{3.699753in}{5.508341in}}{\pgfqpoint{3.707004in}{5.501090in}}%
\pgfpathcurveto{\pgfqpoint{3.714254in}{5.493839in}}{\pgfqpoint{3.724090in}{5.489765in}}{\pgfqpoint{3.734344in}{5.489765in}}%
\pgfpathlineto{\pgfqpoint{3.734344in}{5.489765in}}%
\pgfpathclose%
\pgfusepath{stroke}%
\end{pgfscope}%
\begin{pgfscope}%
\pgfpathrectangle{\pgfqpoint{0.977909in}{2.928000in}}{\pgfqpoint{5.664000in}{5.664000in}}%
\pgfusepath{clip}%
\pgfsetbuttcap%
\pgfsetroundjoin%
\pgfsetlinewidth{1.003750pt}%
\definecolor{currentstroke}{rgb}{0.541176,0.380392,0.658824}%
\pgfsetstrokecolor{currentstroke}%
\pgfsetstrokeopacity{0.900000}%
\pgfsetdash{}{0pt}%
\pgfpathmoveto{\pgfqpoint{3.726761in}{5.011499in}}%
\pgfpathcurveto{\pgfqpoint{3.737015in}{5.011499in}}{\pgfqpoint{3.746850in}{5.015573in}}{\pgfqpoint{3.754101in}{5.022823in}}%
\pgfpathcurveto{\pgfqpoint{3.761352in}{5.030074in}}{\pgfqpoint{3.765426in}{5.039910in}}{\pgfqpoint{3.765426in}{5.050164in}}%
\pgfpathcurveto{\pgfqpoint{3.765426in}{5.060418in}}{\pgfqpoint{3.761352in}{5.070253in}}{\pgfqpoint{3.754101in}{5.077504in}}%
\pgfpathcurveto{\pgfqpoint{3.746850in}{5.084755in}}{\pgfqpoint{3.737015in}{5.088829in}}{\pgfqpoint{3.726761in}{5.088829in}}%
\pgfpathcurveto{\pgfqpoint{3.716507in}{5.088829in}}{\pgfqpoint{3.706671in}{5.084755in}}{\pgfqpoint{3.699421in}{5.077504in}}%
\pgfpathcurveto{\pgfqpoint{3.692170in}{5.070253in}}{\pgfqpoint{3.688096in}{5.060418in}}{\pgfqpoint{3.688096in}{5.050164in}}%
\pgfpathcurveto{\pgfqpoint{3.688096in}{5.039910in}}{\pgfqpoint{3.692170in}{5.030074in}}{\pgfqpoint{3.699421in}{5.022823in}}%
\pgfpathcurveto{\pgfqpoint{3.706671in}{5.015573in}}{\pgfqpoint{3.716507in}{5.011499in}}{\pgfqpoint{3.726761in}{5.011499in}}%
\pgfpathlineto{\pgfqpoint{3.726761in}{5.011499in}}%
\pgfpathclose%
\pgfusepath{stroke}%
\end{pgfscope}%
\begin{pgfscope}%
\pgfpathrectangle{\pgfqpoint{0.977909in}{2.928000in}}{\pgfqpoint{5.664000in}{5.664000in}}%
\pgfusepath{clip}%
\pgfsetbuttcap%
\pgfsetroundjoin%
\pgfsetlinewidth{1.003750pt}%
\definecolor{currentstroke}{rgb}{0.541176,0.380392,0.658824}%
\pgfsetstrokecolor{currentstroke}%
\pgfsetstrokeopacity{0.900000}%
\pgfsetdash{}{0pt}%
\pgfpathmoveto{\pgfqpoint{3.333858in}{4.691835in}}%
\pgfpathcurveto{\pgfqpoint{3.344112in}{4.691835in}}{\pgfqpoint{3.353947in}{4.695909in}}{\pgfqpoint{3.361198in}{4.703160in}}%
\pgfpathcurveto{\pgfqpoint{3.368449in}{4.710411in}}{\pgfqpoint{3.372523in}{4.720246in}}{\pgfqpoint{3.372523in}{4.730500in}}%
\pgfpathcurveto{\pgfqpoint{3.372523in}{4.740755in}}{\pgfqpoint{3.368449in}{4.750590in}}{\pgfqpoint{3.361198in}{4.757841in}}%
\pgfpathcurveto{\pgfqpoint{3.353947in}{4.765091in}}{\pgfqpoint{3.344112in}{4.769165in}}{\pgfqpoint{3.333858in}{4.769165in}}%
\pgfpathcurveto{\pgfqpoint{3.323604in}{4.769165in}}{\pgfqpoint{3.313768in}{4.765091in}}{\pgfqpoint{3.306518in}{4.757841in}}%
\pgfpathcurveto{\pgfqpoint{3.299267in}{4.750590in}}{\pgfqpoint{3.295193in}{4.740755in}}{\pgfqpoint{3.295193in}{4.730500in}}%
\pgfpathcurveto{\pgfqpoint{3.295193in}{4.720246in}}{\pgfqpoint{3.299267in}{4.710411in}}{\pgfqpoint{3.306518in}{4.703160in}}%
\pgfpathcurveto{\pgfqpoint{3.313768in}{4.695909in}}{\pgfqpoint{3.323604in}{4.691835in}}{\pgfqpoint{3.333858in}{4.691835in}}%
\pgfpathlineto{\pgfqpoint{3.333858in}{4.691835in}}%
\pgfpathclose%
\pgfusepath{stroke}%
\end{pgfscope}%
\begin{pgfscope}%
\pgfpathrectangle{\pgfqpoint{0.977909in}{2.928000in}}{\pgfqpoint{5.664000in}{5.664000in}}%
\pgfusepath{clip}%
\pgfsetbuttcap%
\pgfsetroundjoin%
\pgfsetlinewidth{1.003750pt}%
\definecolor{currentstroke}{rgb}{0.541176,0.380392,0.658824}%
\pgfsetstrokecolor{currentstroke}%
\pgfsetstrokeopacity{0.900000}%
\pgfsetdash{}{0pt}%
\pgfpathmoveto{\pgfqpoint{3.452489in}{4.780858in}}%
\pgfpathcurveto{\pgfqpoint{3.462743in}{4.780858in}}{\pgfqpoint{3.472578in}{4.784932in}}{\pgfqpoint{3.479829in}{4.792182in}}%
\pgfpathcurveto{\pgfqpoint{3.487080in}{4.799433in}}{\pgfqpoint{3.491154in}{4.809269in}}{\pgfqpoint{3.491154in}{4.819523in}}%
\pgfpathcurveto{\pgfqpoint{3.491154in}{4.829777in}}{\pgfqpoint{3.487080in}{4.839612in}}{\pgfqpoint{3.479829in}{4.846863in}}%
\pgfpathcurveto{\pgfqpoint{3.472578in}{4.854114in}}{\pgfqpoint{3.462743in}{4.858188in}}{\pgfqpoint{3.452489in}{4.858188in}}%
\pgfpathcurveto{\pgfqpoint{3.442235in}{4.858188in}}{\pgfqpoint{3.432399in}{4.854114in}}{\pgfqpoint{3.425149in}{4.846863in}}%
\pgfpathcurveto{\pgfqpoint{3.417898in}{4.839612in}}{\pgfqpoint{3.413824in}{4.829777in}}{\pgfqpoint{3.413824in}{4.819523in}}%
\pgfpathcurveto{\pgfqpoint{3.413824in}{4.809269in}}{\pgfqpoint{3.417898in}{4.799433in}}{\pgfqpoint{3.425149in}{4.792182in}}%
\pgfpathcurveto{\pgfqpoint{3.432399in}{4.784932in}}{\pgfqpoint{3.442235in}{4.780858in}}{\pgfqpoint{3.452489in}{4.780858in}}%
\pgfpathlineto{\pgfqpoint{3.452489in}{4.780858in}}%
\pgfpathclose%
\pgfusepath{stroke}%
\end{pgfscope}%
\begin{pgfscope}%
\pgfpathrectangle{\pgfqpoint{0.977909in}{2.928000in}}{\pgfqpoint{5.664000in}{5.664000in}}%
\pgfusepath{clip}%
\pgfsetbuttcap%
\pgfsetroundjoin%
\pgfsetlinewidth{1.003750pt}%
\definecolor{currentstroke}{rgb}{0.541176,0.380392,0.658824}%
\pgfsetstrokecolor{currentstroke}%
\pgfsetstrokeopacity{0.900000}%
\pgfsetdash{}{0pt}%
\pgfpathmoveto{\pgfqpoint{3.708651in}{5.646340in}}%
\pgfpathcurveto{\pgfqpoint{3.718905in}{5.646340in}}{\pgfqpoint{3.728741in}{5.650414in}}{\pgfqpoint{3.735991in}{5.657665in}}%
\pgfpathcurveto{\pgfqpoint{3.743242in}{5.664916in}}{\pgfqpoint{3.747316in}{5.674751in}}{\pgfqpoint{3.747316in}{5.685005in}}%
\pgfpathcurveto{\pgfqpoint{3.747316in}{5.695259in}}{\pgfqpoint{3.743242in}{5.705095in}}{\pgfqpoint{3.735991in}{5.712346in}}%
\pgfpathcurveto{\pgfqpoint{3.728741in}{5.719596in}}{\pgfqpoint{3.718905in}{5.723670in}}{\pgfqpoint{3.708651in}{5.723670in}}%
\pgfpathcurveto{\pgfqpoint{3.698397in}{5.723670in}}{\pgfqpoint{3.688562in}{5.719596in}}{\pgfqpoint{3.681311in}{5.712346in}}%
\pgfpathcurveto{\pgfqpoint{3.674060in}{5.705095in}}{\pgfqpoint{3.669986in}{5.695259in}}{\pgfqpoint{3.669986in}{5.685005in}}%
\pgfpathcurveto{\pgfqpoint{3.669986in}{5.674751in}}{\pgfqpoint{3.674060in}{5.664916in}}{\pgfqpoint{3.681311in}{5.657665in}}%
\pgfpathcurveto{\pgfqpoint{3.688562in}{5.650414in}}{\pgfqpoint{3.698397in}{5.646340in}}{\pgfqpoint{3.708651in}{5.646340in}}%
\pgfpathlineto{\pgfqpoint{3.708651in}{5.646340in}}%
\pgfpathclose%
\pgfusepath{stroke}%
\end{pgfscope}%
\begin{pgfscope}%
\pgfpathrectangle{\pgfqpoint{0.977909in}{2.928000in}}{\pgfqpoint{5.664000in}{5.664000in}}%
\pgfusepath{clip}%
\pgfsetbuttcap%
\pgfsetroundjoin%
\pgfsetlinewidth{1.003750pt}%
\definecolor{currentstroke}{rgb}{0.541176,0.380392,0.658824}%
\pgfsetstrokecolor{currentstroke}%
\pgfsetstrokeopacity{0.900000}%
\pgfsetdash{}{0pt}%
\pgfpathmoveto{\pgfqpoint{4.423459in}{5.662073in}}%
\pgfpathcurveto{\pgfqpoint{4.433713in}{5.662073in}}{\pgfqpoint{4.443549in}{5.666147in}}{\pgfqpoint{4.450799in}{5.673398in}}%
\pgfpathcurveto{\pgfqpoint{4.458050in}{5.680649in}}{\pgfqpoint{4.462124in}{5.690484in}}{\pgfqpoint{4.462124in}{5.700738in}}%
\pgfpathcurveto{\pgfqpoint{4.462124in}{5.710992in}}{\pgfqpoint{4.458050in}{5.720828in}}{\pgfqpoint{4.450799in}{5.728078in}}%
\pgfpathcurveto{\pgfqpoint{4.443549in}{5.735329in}}{\pgfqpoint{4.433713in}{5.739403in}}{\pgfqpoint{4.423459in}{5.739403in}}%
\pgfpathcurveto{\pgfqpoint{4.413205in}{5.739403in}}{\pgfqpoint{4.403369in}{5.735329in}}{\pgfqpoint{4.396119in}{5.728078in}}%
\pgfpathcurveto{\pgfqpoint{4.388868in}{5.720828in}}{\pgfqpoint{4.384794in}{5.710992in}}{\pgfqpoint{4.384794in}{5.700738in}}%
\pgfpathcurveto{\pgfqpoint{4.384794in}{5.690484in}}{\pgfqpoint{4.388868in}{5.680649in}}{\pgfqpoint{4.396119in}{5.673398in}}%
\pgfpathcurveto{\pgfqpoint{4.403369in}{5.666147in}}{\pgfqpoint{4.413205in}{5.662073in}}{\pgfqpoint{4.423459in}{5.662073in}}%
\pgfpathlineto{\pgfqpoint{4.423459in}{5.662073in}}%
\pgfpathclose%
\pgfusepath{stroke}%
\end{pgfscope}%
\begin{pgfscope}%
\pgfpathrectangle{\pgfqpoint{0.977909in}{2.928000in}}{\pgfqpoint{5.664000in}{5.664000in}}%
\pgfusepath{clip}%
\pgfsetbuttcap%
\pgfsetroundjoin%
\pgfsetlinewidth{1.003750pt}%
\definecolor{currentstroke}{rgb}{0.541176,0.380392,0.658824}%
\pgfsetstrokecolor{currentstroke}%
\pgfsetstrokeopacity{0.900000}%
\pgfsetdash{}{0pt}%
\pgfpathmoveto{\pgfqpoint{5.027273in}{6.763126in}}%
\pgfpathcurveto{\pgfqpoint{5.037527in}{6.763126in}}{\pgfqpoint{5.047363in}{6.767200in}}{\pgfqpoint{5.054614in}{6.774450in}}%
\pgfpathcurveto{\pgfqpoint{5.061864in}{6.781701in}}{\pgfqpoint{5.065938in}{6.791536in}}{\pgfqpoint{5.065938in}{6.801791in}}%
\pgfpathcurveto{\pgfqpoint{5.065938in}{6.812045in}}{\pgfqpoint{5.061864in}{6.821880in}}{\pgfqpoint{5.054614in}{6.829131in}}%
\pgfpathcurveto{\pgfqpoint{5.047363in}{6.836382in}}{\pgfqpoint{5.037527in}{6.840456in}}{\pgfqpoint{5.027273in}{6.840456in}}%
\pgfpathcurveto{\pgfqpoint{5.017019in}{6.840456in}}{\pgfqpoint{5.007184in}{6.836382in}}{\pgfqpoint{4.999933in}{6.829131in}}%
\pgfpathcurveto{\pgfqpoint{4.992682in}{6.821880in}}{\pgfqpoint{4.988608in}{6.812045in}}{\pgfqpoint{4.988608in}{6.801791in}}%
\pgfpathcurveto{\pgfqpoint{4.988608in}{6.791536in}}{\pgfqpoint{4.992682in}{6.781701in}}{\pgfqpoint{4.999933in}{6.774450in}}%
\pgfpathcurveto{\pgfqpoint{5.007184in}{6.767200in}}{\pgfqpoint{5.017019in}{6.763126in}}{\pgfqpoint{5.027273in}{6.763126in}}%
\pgfpathlineto{\pgfqpoint{5.027273in}{6.763126in}}%
\pgfpathclose%
\pgfusepath{stroke}%
\end{pgfscope}%
\begin{pgfscope}%
\pgfpathrectangle{\pgfqpoint{0.977909in}{2.928000in}}{\pgfqpoint{5.664000in}{5.664000in}}%
\pgfusepath{clip}%
\pgfsetbuttcap%
\pgfsetroundjoin%
\pgfsetlinewidth{1.003750pt}%
\definecolor{currentstroke}{rgb}{0.541176,0.380392,0.658824}%
\pgfsetstrokecolor{currentstroke}%
\pgfsetstrokeopacity{0.900000}%
\pgfsetdash{}{0pt}%
\pgfpathmoveto{\pgfqpoint{4.448443in}{5.762369in}}%
\pgfpathcurveto{\pgfqpoint{4.458697in}{5.762369in}}{\pgfqpoint{4.468532in}{5.766443in}}{\pgfqpoint{4.475783in}{5.773694in}}%
\pgfpathcurveto{\pgfqpoint{4.483034in}{5.780944in}}{\pgfqpoint{4.487108in}{5.790780in}}{\pgfqpoint{4.487108in}{5.801034in}}%
\pgfpathcurveto{\pgfqpoint{4.487108in}{5.811288in}}{\pgfqpoint{4.483034in}{5.821124in}}{\pgfqpoint{4.475783in}{5.828374in}}%
\pgfpathcurveto{\pgfqpoint{4.468532in}{5.835625in}}{\pgfqpoint{4.458697in}{5.839699in}}{\pgfqpoint{4.448443in}{5.839699in}}%
\pgfpathcurveto{\pgfqpoint{4.438189in}{5.839699in}}{\pgfqpoint{4.428353in}{5.835625in}}{\pgfqpoint{4.421102in}{5.828374in}}%
\pgfpathcurveto{\pgfqpoint{4.413852in}{5.821124in}}{\pgfqpoint{4.409778in}{5.811288in}}{\pgfqpoint{4.409778in}{5.801034in}}%
\pgfpathcurveto{\pgfqpoint{4.409778in}{5.790780in}}{\pgfqpoint{4.413852in}{5.780944in}}{\pgfqpoint{4.421102in}{5.773694in}}%
\pgfpathcurveto{\pgfqpoint{4.428353in}{5.766443in}}{\pgfqpoint{4.438189in}{5.762369in}}{\pgfqpoint{4.448443in}{5.762369in}}%
\pgfpathlineto{\pgfqpoint{4.448443in}{5.762369in}}%
\pgfpathclose%
\pgfusepath{stroke}%
\end{pgfscope}%
\begin{pgfscope}%
\pgfpathrectangle{\pgfqpoint{0.977909in}{2.928000in}}{\pgfqpoint{5.664000in}{5.664000in}}%
\pgfusepath{clip}%
\pgfsetbuttcap%
\pgfsetroundjoin%
\pgfsetlinewidth{1.003750pt}%
\definecolor{currentstroke}{rgb}{0.541176,0.380392,0.658824}%
\pgfsetstrokecolor{currentstroke}%
\pgfsetstrokeopacity{0.900000}%
\pgfsetdash{}{0pt}%
\pgfpathmoveto{\pgfqpoint{4.411407in}{6.380197in}}%
\pgfpathcurveto{\pgfqpoint{4.421661in}{6.380197in}}{\pgfqpoint{4.431496in}{6.384271in}}{\pgfqpoint{4.438747in}{6.391521in}}%
\pgfpathcurveto{\pgfqpoint{4.445998in}{6.398772in}}{\pgfqpoint{4.450072in}{6.408608in}}{\pgfqpoint{4.450072in}{6.418862in}}%
\pgfpathcurveto{\pgfqpoint{4.450072in}{6.429116in}}{\pgfqpoint{4.445998in}{6.438951in}}{\pgfqpoint{4.438747in}{6.446202in}}%
\pgfpathcurveto{\pgfqpoint{4.431496in}{6.453453in}}{\pgfqpoint{4.421661in}{6.457527in}}{\pgfqpoint{4.411407in}{6.457527in}}%
\pgfpathcurveto{\pgfqpoint{4.401153in}{6.457527in}}{\pgfqpoint{4.391317in}{6.453453in}}{\pgfqpoint{4.384067in}{6.446202in}}%
\pgfpathcurveto{\pgfqpoint{4.376816in}{6.438951in}}{\pgfqpoint{4.372742in}{6.429116in}}{\pgfqpoint{4.372742in}{6.418862in}}%
\pgfpathcurveto{\pgfqpoint{4.372742in}{6.408608in}}{\pgfqpoint{4.376816in}{6.398772in}}{\pgfqpoint{4.384067in}{6.391521in}}%
\pgfpathcurveto{\pgfqpoint{4.391317in}{6.384271in}}{\pgfqpoint{4.401153in}{6.380197in}}{\pgfqpoint{4.411407in}{6.380197in}}%
\pgfpathlineto{\pgfqpoint{4.411407in}{6.380197in}}%
\pgfpathclose%
\pgfusepath{stroke}%
\end{pgfscope}%
\begin{pgfscope}%
\pgfpathrectangle{\pgfqpoint{0.977909in}{2.928000in}}{\pgfqpoint{5.664000in}{5.664000in}}%
\pgfusepath{clip}%
\pgfsetbuttcap%
\pgfsetroundjoin%
\pgfsetlinewidth{1.003750pt}%
\definecolor{currentstroke}{rgb}{0.541176,0.380392,0.658824}%
\pgfsetstrokecolor{currentstroke}%
\pgfsetstrokeopacity{0.900000}%
\pgfsetdash{}{0pt}%
\pgfpathmoveto{\pgfqpoint{4.622010in}{5.830172in}}%
\pgfpathcurveto{\pgfqpoint{4.632265in}{5.830172in}}{\pgfqpoint{4.642100in}{5.834246in}}{\pgfqpoint{4.649351in}{5.841496in}}%
\pgfpathcurveto{\pgfqpoint{4.656602in}{5.848747in}}{\pgfqpoint{4.660675in}{5.858582in}}{\pgfqpoint{4.660675in}{5.868837in}}%
\pgfpathcurveto{\pgfqpoint{4.660675in}{5.879091in}}{\pgfqpoint{4.656602in}{5.888926in}}{\pgfqpoint{4.649351in}{5.896177in}}%
\pgfpathcurveto{\pgfqpoint{4.642100in}{5.903428in}}{\pgfqpoint{4.632265in}{5.907502in}}{\pgfqpoint{4.622010in}{5.907502in}}%
\pgfpathcurveto{\pgfqpoint{4.611756in}{5.907502in}}{\pgfqpoint{4.601921in}{5.903428in}}{\pgfqpoint{4.594670in}{5.896177in}}%
\pgfpathcurveto{\pgfqpoint{4.587419in}{5.888926in}}{\pgfqpoint{4.583345in}{5.879091in}}{\pgfqpoint{4.583345in}{5.868837in}}%
\pgfpathcurveto{\pgfqpoint{4.583345in}{5.858582in}}{\pgfqpoint{4.587419in}{5.848747in}}{\pgfqpoint{4.594670in}{5.841496in}}%
\pgfpathcurveto{\pgfqpoint{4.601921in}{5.834246in}}{\pgfqpoint{4.611756in}{5.830172in}}{\pgfqpoint{4.622010in}{5.830172in}}%
\pgfpathlineto{\pgfqpoint{4.622010in}{5.830172in}}%
\pgfpathclose%
\pgfusepath{stroke}%
\end{pgfscope}%
\begin{pgfscope}%
\pgfpathrectangle{\pgfqpoint{0.977909in}{2.928000in}}{\pgfqpoint{5.664000in}{5.664000in}}%
\pgfusepath{clip}%
\pgfsetbuttcap%
\pgfsetroundjoin%
\pgfsetlinewidth{1.003750pt}%
\definecolor{currentstroke}{rgb}{0.541176,0.380392,0.658824}%
\pgfsetstrokecolor{currentstroke}%
\pgfsetstrokeopacity{0.900000}%
\pgfsetdash{}{0pt}%
\pgfpathmoveto{\pgfqpoint{4.728887in}{6.238009in}}%
\pgfpathcurveto{\pgfqpoint{4.739141in}{6.238009in}}{\pgfqpoint{4.748977in}{6.242083in}}{\pgfqpoint{4.756227in}{6.249334in}}%
\pgfpathcurveto{\pgfqpoint{4.763478in}{6.256584in}}{\pgfqpoint{4.767552in}{6.266420in}}{\pgfqpoint{4.767552in}{6.276674in}}%
\pgfpathcurveto{\pgfqpoint{4.767552in}{6.286928in}}{\pgfqpoint{4.763478in}{6.296763in}}{\pgfqpoint{4.756227in}{6.304014in}}%
\pgfpathcurveto{\pgfqpoint{4.748977in}{6.311265in}}{\pgfqpoint{4.739141in}{6.315339in}}{\pgfqpoint{4.728887in}{6.315339in}}%
\pgfpathcurveto{\pgfqpoint{4.718633in}{6.315339in}}{\pgfqpoint{4.708797in}{6.311265in}}{\pgfqpoint{4.701547in}{6.304014in}}%
\pgfpathcurveto{\pgfqpoint{4.694296in}{6.296763in}}{\pgfqpoint{4.690222in}{6.286928in}}{\pgfqpoint{4.690222in}{6.276674in}}%
\pgfpathcurveto{\pgfqpoint{4.690222in}{6.266420in}}{\pgfqpoint{4.694296in}{6.256584in}}{\pgfqpoint{4.701547in}{6.249334in}}%
\pgfpathcurveto{\pgfqpoint{4.708797in}{6.242083in}}{\pgfqpoint{4.718633in}{6.238009in}}{\pgfqpoint{4.728887in}{6.238009in}}%
\pgfpathlineto{\pgfqpoint{4.728887in}{6.238009in}}%
\pgfpathclose%
\pgfusepath{stroke}%
\end{pgfscope}%
\begin{pgfscope}%
\pgfpathrectangle{\pgfqpoint{0.977909in}{2.928000in}}{\pgfqpoint{5.664000in}{5.664000in}}%
\pgfusepath{clip}%
\pgfsetbuttcap%
\pgfsetroundjoin%
\pgfsetlinewidth{1.003750pt}%
\definecolor{currentstroke}{rgb}{0.541176,0.380392,0.658824}%
\pgfsetstrokecolor{currentstroke}%
\pgfsetstrokeopacity{0.900000}%
\pgfsetdash{}{0pt}%
\pgfpathmoveto{\pgfqpoint{3.948540in}{4.943214in}}%
\pgfpathcurveto{\pgfqpoint{3.958794in}{4.943214in}}{\pgfqpoint{3.968630in}{4.947288in}}{\pgfqpoint{3.975880in}{4.954539in}}%
\pgfpathcurveto{\pgfqpoint{3.983131in}{4.961789in}}{\pgfqpoint{3.987205in}{4.971625in}}{\pgfqpoint{3.987205in}{4.981879in}}%
\pgfpathcurveto{\pgfqpoint{3.987205in}{4.992133in}}{\pgfqpoint{3.983131in}{5.001969in}}{\pgfqpoint{3.975880in}{5.009219in}}%
\pgfpathcurveto{\pgfqpoint{3.968630in}{5.016470in}}{\pgfqpoint{3.958794in}{5.020544in}}{\pgfqpoint{3.948540in}{5.020544in}}%
\pgfpathcurveto{\pgfqpoint{3.938286in}{5.020544in}}{\pgfqpoint{3.928451in}{5.016470in}}{\pgfqpoint{3.921200in}{5.009219in}}%
\pgfpathcurveto{\pgfqpoint{3.913949in}{5.001969in}}{\pgfqpoint{3.909875in}{4.992133in}}{\pgfqpoint{3.909875in}{4.981879in}}%
\pgfpathcurveto{\pgfqpoint{3.909875in}{4.971625in}}{\pgfqpoint{3.913949in}{4.961789in}}{\pgfqpoint{3.921200in}{4.954539in}}%
\pgfpathcurveto{\pgfqpoint{3.928451in}{4.947288in}}{\pgfqpoint{3.938286in}{4.943214in}}{\pgfqpoint{3.948540in}{4.943214in}}%
\pgfpathlineto{\pgfqpoint{3.948540in}{4.943214in}}%
\pgfpathclose%
\pgfusepath{stroke}%
\end{pgfscope}%
\begin{pgfscope}%
\pgfpathrectangle{\pgfqpoint{0.977909in}{2.928000in}}{\pgfqpoint{5.664000in}{5.664000in}}%
\pgfusepath{clip}%
\pgfsetbuttcap%
\pgfsetroundjoin%
\pgfsetlinewidth{1.003750pt}%
\definecolor{currentstroke}{rgb}{0.541176,0.380392,0.658824}%
\pgfsetstrokecolor{currentstroke}%
\pgfsetstrokeopacity{0.900000}%
\pgfsetdash{}{0pt}%
\pgfpathmoveto{\pgfqpoint{5.876832in}{7.470310in}}%
\pgfpathcurveto{\pgfqpoint{5.887086in}{7.470310in}}{\pgfqpoint{5.896922in}{7.474384in}}{\pgfqpoint{5.904173in}{7.481634in}}%
\pgfpathcurveto{\pgfqpoint{5.911423in}{7.488885in}}{\pgfqpoint{5.915497in}{7.498721in}}{\pgfqpoint{5.915497in}{7.508975in}}%
\pgfpathcurveto{\pgfqpoint{5.915497in}{7.519229in}}{\pgfqpoint{5.911423in}{7.529064in}}{\pgfqpoint{5.904173in}{7.536315in}}%
\pgfpathcurveto{\pgfqpoint{5.896922in}{7.543566in}}{\pgfqpoint{5.887086in}{7.547640in}}{\pgfqpoint{5.876832in}{7.547640in}}%
\pgfpathcurveto{\pgfqpoint{5.866578in}{7.547640in}}{\pgfqpoint{5.856743in}{7.543566in}}{\pgfqpoint{5.849492in}{7.536315in}}%
\pgfpathcurveto{\pgfqpoint{5.842241in}{7.529064in}}{\pgfqpoint{5.838167in}{7.519229in}}{\pgfqpoint{5.838167in}{7.508975in}}%
\pgfpathcurveto{\pgfqpoint{5.838167in}{7.498721in}}{\pgfqpoint{5.842241in}{7.488885in}}{\pgfqpoint{5.849492in}{7.481634in}}%
\pgfpathcurveto{\pgfqpoint{5.856743in}{7.474384in}}{\pgfqpoint{5.866578in}{7.470310in}}{\pgfqpoint{5.876832in}{7.470310in}}%
\pgfpathlineto{\pgfqpoint{5.876832in}{7.470310in}}%
\pgfpathclose%
\pgfusepath{stroke}%
\end{pgfscope}%
\begin{pgfscope}%
\pgfpathrectangle{\pgfqpoint{0.977909in}{2.928000in}}{\pgfqpoint{5.664000in}{5.664000in}}%
\pgfusepath{clip}%
\pgfsetbuttcap%
\pgfsetroundjoin%
\pgfsetlinewidth{1.003750pt}%
\definecolor{currentstroke}{rgb}{0.541176,0.380392,0.658824}%
\pgfsetstrokecolor{currentstroke}%
\pgfsetstrokeopacity{0.900000}%
\pgfsetdash{}{0pt}%
\pgfpathmoveto{\pgfqpoint{5.440429in}{7.234399in}}%
\pgfpathcurveto{\pgfqpoint{5.450683in}{7.234399in}}{\pgfqpoint{5.460519in}{7.238473in}}{\pgfqpoint{5.467769in}{7.245724in}}%
\pgfpathcurveto{\pgfqpoint{5.475020in}{7.252974in}}{\pgfqpoint{5.479094in}{7.262810in}}{\pgfqpoint{5.479094in}{7.273064in}}%
\pgfpathcurveto{\pgfqpoint{5.479094in}{7.283318in}}{\pgfqpoint{5.475020in}{7.293154in}}{\pgfqpoint{5.467769in}{7.300404in}}%
\pgfpathcurveto{\pgfqpoint{5.460519in}{7.307655in}}{\pgfqpoint{5.450683in}{7.311729in}}{\pgfqpoint{5.440429in}{7.311729in}}%
\pgfpathcurveto{\pgfqpoint{5.430175in}{7.311729in}}{\pgfqpoint{5.420339in}{7.307655in}}{\pgfqpoint{5.413089in}{7.300404in}}%
\pgfpathcurveto{\pgfqpoint{5.405838in}{7.293154in}}{\pgfqpoint{5.401764in}{7.283318in}}{\pgfqpoint{5.401764in}{7.273064in}}%
\pgfpathcurveto{\pgfqpoint{5.401764in}{7.262810in}}{\pgfqpoint{5.405838in}{7.252974in}}{\pgfqpoint{5.413089in}{7.245724in}}%
\pgfpathcurveto{\pgfqpoint{5.420339in}{7.238473in}}{\pgfqpoint{5.430175in}{7.234399in}}{\pgfqpoint{5.440429in}{7.234399in}}%
\pgfpathlineto{\pgfqpoint{5.440429in}{7.234399in}}%
\pgfpathclose%
\pgfusepath{stroke}%
\end{pgfscope}%
\begin{pgfscope}%
\pgfpathrectangle{\pgfqpoint{0.977909in}{2.928000in}}{\pgfqpoint{5.664000in}{5.664000in}}%
\pgfusepath{clip}%
\pgfsetbuttcap%
\pgfsetroundjoin%
\pgfsetlinewidth{1.003750pt}%
\definecolor{currentstroke}{rgb}{0.541176,0.380392,0.658824}%
\pgfsetstrokecolor{currentstroke}%
\pgfsetstrokeopacity{0.900000}%
\pgfsetdash{}{0pt}%
\pgfpathmoveto{\pgfqpoint{5.218824in}{6.663041in}}%
\pgfpathcurveto{\pgfqpoint{5.229078in}{6.663041in}}{\pgfqpoint{5.238914in}{6.667115in}}{\pgfqpoint{5.246164in}{6.674366in}}%
\pgfpathcurveto{\pgfqpoint{5.253415in}{6.681617in}}{\pgfqpoint{5.257489in}{6.691452in}}{\pgfqpoint{5.257489in}{6.701706in}}%
\pgfpathcurveto{\pgfqpoint{5.257489in}{6.711961in}}{\pgfqpoint{5.253415in}{6.721796in}}{\pgfqpoint{5.246164in}{6.729047in}}%
\pgfpathcurveto{\pgfqpoint{5.238914in}{6.736298in}}{\pgfqpoint{5.229078in}{6.740372in}}{\pgfqpoint{5.218824in}{6.740372in}}%
\pgfpathcurveto{\pgfqpoint{5.208570in}{6.740372in}}{\pgfqpoint{5.198735in}{6.736298in}}{\pgfqpoint{5.191484in}{6.729047in}}%
\pgfpathcurveto{\pgfqpoint{5.184233in}{6.721796in}}{\pgfqpoint{5.180159in}{6.711961in}}{\pgfqpoint{5.180159in}{6.701706in}}%
\pgfpathcurveto{\pgfqpoint{5.180159in}{6.691452in}}{\pgfqpoint{5.184233in}{6.681617in}}{\pgfqpoint{5.191484in}{6.674366in}}%
\pgfpathcurveto{\pgfqpoint{5.198735in}{6.667115in}}{\pgfqpoint{5.208570in}{6.663041in}}{\pgfqpoint{5.218824in}{6.663041in}}%
\pgfpathlineto{\pgfqpoint{5.218824in}{6.663041in}}%
\pgfpathclose%
\pgfusepath{stroke}%
\end{pgfscope}%
\begin{pgfscope}%
\pgfpathrectangle{\pgfqpoint{0.977909in}{2.928000in}}{\pgfqpoint{5.664000in}{5.664000in}}%
\pgfusepath{clip}%
\pgfsetbuttcap%
\pgfsetroundjoin%
\pgfsetlinewidth{1.003750pt}%
\definecolor{currentstroke}{rgb}{0.541176,0.380392,0.658824}%
\pgfsetstrokecolor{currentstroke}%
\pgfsetstrokeopacity{0.900000}%
\pgfsetdash{}{0pt}%
\pgfpathmoveto{\pgfqpoint{4.852290in}{6.768363in}}%
\pgfpathcurveto{\pgfqpoint{4.862544in}{6.768363in}}{\pgfqpoint{4.872379in}{6.772437in}}{\pgfqpoint{4.879630in}{6.779688in}}%
\pgfpathcurveto{\pgfqpoint{4.886881in}{6.786939in}}{\pgfqpoint{4.890955in}{6.796774in}}{\pgfqpoint{4.890955in}{6.807028in}}%
\pgfpathcurveto{\pgfqpoint{4.890955in}{6.817283in}}{\pgfqpoint{4.886881in}{6.827118in}}{\pgfqpoint{4.879630in}{6.834369in}}%
\pgfpathcurveto{\pgfqpoint{4.872379in}{6.841620in}}{\pgfqpoint{4.862544in}{6.845694in}}{\pgfqpoint{4.852290in}{6.845694in}}%
\pgfpathcurveto{\pgfqpoint{4.842036in}{6.845694in}}{\pgfqpoint{4.832200in}{6.841620in}}{\pgfqpoint{4.824950in}{6.834369in}}%
\pgfpathcurveto{\pgfqpoint{4.817699in}{6.827118in}}{\pgfqpoint{4.813625in}{6.817283in}}{\pgfqpoint{4.813625in}{6.807028in}}%
\pgfpathcurveto{\pgfqpoint{4.813625in}{6.796774in}}{\pgfqpoint{4.817699in}{6.786939in}}{\pgfqpoint{4.824950in}{6.779688in}}%
\pgfpathcurveto{\pgfqpoint{4.832200in}{6.772437in}}{\pgfqpoint{4.842036in}{6.768363in}}{\pgfqpoint{4.852290in}{6.768363in}}%
\pgfpathlineto{\pgfqpoint{4.852290in}{6.768363in}}%
\pgfpathclose%
\pgfusepath{stroke}%
\end{pgfscope}%
\begin{pgfscope}%
\pgfpathrectangle{\pgfqpoint{0.977909in}{2.928000in}}{\pgfqpoint{5.664000in}{5.664000in}}%
\pgfusepath{clip}%
\pgfsetbuttcap%
\pgfsetroundjoin%
\pgfsetlinewidth{1.003750pt}%
\definecolor{currentstroke}{rgb}{0.541176,0.380392,0.658824}%
\pgfsetstrokecolor{currentstroke}%
\pgfsetstrokeopacity{0.900000}%
\pgfsetdash{}{0pt}%
\pgfpathmoveto{\pgfqpoint{5.187430in}{7.328925in}}%
\pgfpathcurveto{\pgfqpoint{5.197684in}{7.328925in}}{\pgfqpoint{5.207519in}{7.332999in}}{\pgfqpoint{5.214770in}{7.340249in}}%
\pgfpathcurveto{\pgfqpoint{5.222021in}{7.347500in}}{\pgfqpoint{5.226095in}{7.357336in}}{\pgfqpoint{5.226095in}{7.367590in}}%
\pgfpathcurveto{\pgfqpoint{5.226095in}{7.377844in}}{\pgfqpoint{5.222021in}{7.387679in}}{\pgfqpoint{5.214770in}{7.394930in}}%
\pgfpathcurveto{\pgfqpoint{5.207519in}{7.402181in}}{\pgfqpoint{5.197684in}{7.406255in}}{\pgfqpoint{5.187430in}{7.406255in}}%
\pgfpathcurveto{\pgfqpoint{5.177176in}{7.406255in}}{\pgfqpoint{5.167340in}{7.402181in}}{\pgfqpoint{5.160090in}{7.394930in}}%
\pgfpathcurveto{\pgfqpoint{5.152839in}{7.387679in}}{\pgfqpoint{5.148765in}{7.377844in}}{\pgfqpoint{5.148765in}{7.367590in}}%
\pgfpathcurveto{\pgfqpoint{5.148765in}{7.357336in}}{\pgfqpoint{5.152839in}{7.347500in}}{\pgfqpoint{5.160090in}{7.340249in}}%
\pgfpathcurveto{\pgfqpoint{5.167340in}{7.332999in}}{\pgfqpoint{5.177176in}{7.328925in}}{\pgfqpoint{5.187430in}{7.328925in}}%
\pgfpathlineto{\pgfqpoint{5.187430in}{7.328925in}}%
\pgfpathclose%
\pgfusepath{stroke}%
\end{pgfscope}%
\begin{pgfscope}%
\pgfpathrectangle{\pgfqpoint{0.977909in}{2.928000in}}{\pgfqpoint{5.664000in}{5.664000in}}%
\pgfusepath{clip}%
\pgfsetbuttcap%
\pgfsetroundjoin%
\pgfsetlinewidth{1.003750pt}%
\definecolor{currentstroke}{rgb}{0.541176,0.380392,0.658824}%
\pgfsetstrokecolor{currentstroke}%
\pgfsetstrokeopacity{0.900000}%
\pgfsetdash{}{0pt}%
\pgfpathmoveto{\pgfqpoint{3.898851in}{6.280871in}}%
\pgfpathcurveto{\pgfqpoint{3.909105in}{6.280871in}}{\pgfqpoint{3.918940in}{6.284945in}}{\pgfqpoint{3.926191in}{6.292196in}}%
\pgfpathcurveto{\pgfqpoint{3.933442in}{6.299447in}}{\pgfqpoint{3.937516in}{6.309282in}}{\pgfqpoint{3.937516in}{6.319536in}}%
\pgfpathcurveto{\pgfqpoint{3.937516in}{6.329791in}}{\pgfqpoint{3.933442in}{6.339626in}}{\pgfqpoint{3.926191in}{6.346877in}}%
\pgfpathcurveto{\pgfqpoint{3.918940in}{6.354127in}}{\pgfqpoint{3.909105in}{6.358201in}}{\pgfqpoint{3.898851in}{6.358201in}}%
\pgfpathcurveto{\pgfqpoint{3.888597in}{6.358201in}}{\pgfqpoint{3.878761in}{6.354127in}}{\pgfqpoint{3.871510in}{6.346877in}}%
\pgfpathcurveto{\pgfqpoint{3.864260in}{6.339626in}}{\pgfqpoint{3.860186in}{6.329791in}}{\pgfqpoint{3.860186in}{6.319536in}}%
\pgfpathcurveto{\pgfqpoint{3.860186in}{6.309282in}}{\pgfqpoint{3.864260in}{6.299447in}}{\pgfqpoint{3.871510in}{6.292196in}}%
\pgfpathcurveto{\pgfqpoint{3.878761in}{6.284945in}}{\pgfqpoint{3.888597in}{6.280871in}}{\pgfqpoint{3.898851in}{6.280871in}}%
\pgfpathlineto{\pgfqpoint{3.898851in}{6.280871in}}%
\pgfpathclose%
\pgfusepath{stroke}%
\end{pgfscope}%
\begin{pgfscope}%
\pgfpathrectangle{\pgfqpoint{0.977909in}{2.928000in}}{\pgfqpoint{5.664000in}{5.664000in}}%
\pgfusepath{clip}%
\pgfsetbuttcap%
\pgfsetroundjoin%
\pgfsetlinewidth{1.003750pt}%
\definecolor{currentstroke}{rgb}{0.541176,0.380392,0.658824}%
\pgfsetstrokecolor{currentstroke}%
\pgfsetstrokeopacity{0.900000}%
\pgfsetdash{}{0pt}%
\pgfpathmoveto{\pgfqpoint{4.446540in}{6.778426in}}%
\pgfpathcurveto{\pgfqpoint{4.456794in}{6.778426in}}{\pgfqpoint{4.466629in}{6.782500in}}{\pgfqpoint{4.473880in}{6.789751in}}%
\pgfpathcurveto{\pgfqpoint{4.481131in}{6.797002in}}{\pgfqpoint{4.485205in}{6.806837in}}{\pgfqpoint{4.485205in}{6.817091in}}%
\pgfpathcurveto{\pgfqpoint{4.485205in}{6.827346in}}{\pgfqpoint{4.481131in}{6.837181in}}{\pgfqpoint{4.473880in}{6.844432in}}%
\pgfpathcurveto{\pgfqpoint{4.466629in}{6.851683in}}{\pgfqpoint{4.456794in}{6.855757in}}{\pgfqpoint{4.446540in}{6.855757in}}%
\pgfpathcurveto{\pgfqpoint{4.436286in}{6.855757in}}{\pgfqpoint{4.426450in}{6.851683in}}{\pgfqpoint{4.419199in}{6.844432in}}%
\pgfpathcurveto{\pgfqpoint{4.411949in}{6.837181in}}{\pgfqpoint{4.407875in}{6.827346in}}{\pgfqpoint{4.407875in}{6.817091in}}%
\pgfpathcurveto{\pgfqpoint{4.407875in}{6.806837in}}{\pgfqpoint{4.411949in}{6.797002in}}{\pgfqpoint{4.419199in}{6.789751in}}%
\pgfpathcurveto{\pgfqpoint{4.426450in}{6.782500in}}{\pgfqpoint{4.436286in}{6.778426in}}{\pgfqpoint{4.446540in}{6.778426in}}%
\pgfpathlineto{\pgfqpoint{4.446540in}{6.778426in}}%
\pgfpathclose%
\pgfusepath{stroke}%
\end{pgfscope}%
\begin{pgfscope}%
\pgfpathrectangle{\pgfqpoint{0.977909in}{2.928000in}}{\pgfqpoint{5.664000in}{5.664000in}}%
\pgfusepath{clip}%
\pgfsetbuttcap%
\pgfsetroundjoin%
\pgfsetlinewidth{1.003750pt}%
\definecolor{currentstroke}{rgb}{0.541176,0.380392,0.658824}%
\pgfsetstrokecolor{currentstroke}%
\pgfsetstrokeopacity{0.900000}%
\pgfsetdash{}{0pt}%
\pgfpathmoveto{\pgfqpoint{4.389463in}{5.539220in}}%
\pgfpathcurveto{\pgfqpoint{4.399717in}{5.539220in}}{\pgfqpoint{4.409553in}{5.543294in}}{\pgfqpoint{4.416803in}{5.550545in}}%
\pgfpathcurveto{\pgfqpoint{4.424054in}{5.557796in}}{\pgfqpoint{4.428128in}{5.567631in}}{\pgfqpoint{4.428128in}{5.577885in}}%
\pgfpathcurveto{\pgfqpoint{4.428128in}{5.588140in}}{\pgfqpoint{4.424054in}{5.597975in}}{\pgfqpoint{4.416803in}{5.605226in}}%
\pgfpathcurveto{\pgfqpoint{4.409553in}{5.612476in}}{\pgfqpoint{4.399717in}{5.616550in}}{\pgfqpoint{4.389463in}{5.616550in}}%
\pgfpathcurveto{\pgfqpoint{4.379209in}{5.616550in}}{\pgfqpoint{4.369373in}{5.612476in}}{\pgfqpoint{4.362123in}{5.605226in}}%
\pgfpathcurveto{\pgfqpoint{4.354872in}{5.597975in}}{\pgfqpoint{4.350798in}{5.588140in}}{\pgfqpoint{4.350798in}{5.577885in}}%
\pgfpathcurveto{\pgfqpoint{4.350798in}{5.567631in}}{\pgfqpoint{4.354872in}{5.557796in}}{\pgfqpoint{4.362123in}{5.550545in}}%
\pgfpathcurveto{\pgfqpoint{4.369373in}{5.543294in}}{\pgfqpoint{4.379209in}{5.539220in}}{\pgfqpoint{4.389463in}{5.539220in}}%
\pgfpathlineto{\pgfqpoint{4.389463in}{5.539220in}}%
\pgfpathclose%
\pgfusepath{stroke}%
\end{pgfscope}%
\begin{pgfscope}%
\pgfpathrectangle{\pgfqpoint{0.977909in}{2.928000in}}{\pgfqpoint{5.664000in}{5.664000in}}%
\pgfusepath{clip}%
\pgfsetbuttcap%
\pgfsetroundjoin%
\pgfsetlinewidth{1.003750pt}%
\definecolor{currentstroke}{rgb}{0.541176,0.380392,0.658824}%
\pgfsetstrokecolor{currentstroke}%
\pgfsetstrokeopacity{0.900000}%
\pgfsetdash{}{0pt}%
\pgfpathmoveto{\pgfqpoint{3.311040in}{5.110057in}}%
\pgfpathcurveto{\pgfqpoint{3.321294in}{5.110057in}}{\pgfqpoint{3.331129in}{5.114131in}}{\pgfqpoint{3.338380in}{5.121382in}}%
\pgfpathcurveto{\pgfqpoint{3.345631in}{5.128633in}}{\pgfqpoint{3.349705in}{5.138468in}}{\pgfqpoint{3.349705in}{5.148722in}}%
\pgfpathcurveto{\pgfqpoint{3.349705in}{5.158976in}}{\pgfqpoint{3.345631in}{5.168812in}}{\pgfqpoint{3.338380in}{5.176063in}}%
\pgfpathcurveto{\pgfqpoint{3.331129in}{5.183313in}}{\pgfqpoint{3.321294in}{5.187387in}}{\pgfqpoint{3.311040in}{5.187387in}}%
\pgfpathcurveto{\pgfqpoint{3.300786in}{5.187387in}}{\pgfqpoint{3.290950in}{5.183313in}}{\pgfqpoint{3.283700in}{5.176063in}}%
\pgfpathcurveto{\pgfqpoint{3.276449in}{5.168812in}}{\pgfqpoint{3.272375in}{5.158976in}}{\pgfqpoint{3.272375in}{5.148722in}}%
\pgfpathcurveto{\pgfqpoint{3.272375in}{5.138468in}}{\pgfqpoint{3.276449in}{5.128633in}}{\pgfqpoint{3.283700in}{5.121382in}}%
\pgfpathcurveto{\pgfqpoint{3.290950in}{5.114131in}}{\pgfqpoint{3.300786in}{5.110057in}}{\pgfqpoint{3.311040in}{5.110057in}}%
\pgfpathlineto{\pgfqpoint{3.311040in}{5.110057in}}%
\pgfpathclose%
\pgfusepath{stroke}%
\end{pgfscope}%
\begin{pgfscope}%
\pgfpathrectangle{\pgfqpoint{0.977909in}{2.928000in}}{\pgfqpoint{5.664000in}{5.664000in}}%
\pgfusepath{clip}%
\pgfsetbuttcap%
\pgfsetroundjoin%
\pgfsetlinewidth{1.003750pt}%
\definecolor{currentstroke}{rgb}{0.541176,0.380392,0.658824}%
\pgfsetstrokecolor{currentstroke}%
\pgfsetstrokeopacity{0.900000}%
\pgfsetdash{}{0pt}%
\pgfpathmoveto{\pgfqpoint{4.052117in}{5.491438in}}%
\pgfpathcurveto{\pgfqpoint{4.062371in}{5.491438in}}{\pgfqpoint{4.072207in}{5.495512in}}{\pgfqpoint{4.079458in}{5.502762in}}%
\pgfpathcurveto{\pgfqpoint{4.086708in}{5.510013in}}{\pgfqpoint{4.090782in}{5.519848in}}{\pgfqpoint{4.090782in}{5.530103in}}%
\pgfpathcurveto{\pgfqpoint{4.090782in}{5.540357in}}{\pgfqpoint{4.086708in}{5.550192in}}{\pgfqpoint{4.079458in}{5.557443in}}%
\pgfpathcurveto{\pgfqpoint{4.072207in}{5.564694in}}{\pgfqpoint{4.062371in}{5.568768in}}{\pgfqpoint{4.052117in}{5.568768in}}%
\pgfpathcurveto{\pgfqpoint{4.041863in}{5.568768in}}{\pgfqpoint{4.032028in}{5.564694in}}{\pgfqpoint{4.024777in}{5.557443in}}%
\pgfpathcurveto{\pgfqpoint{4.017526in}{5.550192in}}{\pgfqpoint{4.013452in}{5.540357in}}{\pgfqpoint{4.013452in}{5.530103in}}%
\pgfpathcurveto{\pgfqpoint{4.013452in}{5.519848in}}{\pgfqpoint{4.017526in}{5.510013in}}{\pgfqpoint{4.024777in}{5.502762in}}%
\pgfpathcurveto{\pgfqpoint{4.032028in}{5.495512in}}{\pgfqpoint{4.041863in}{5.491438in}}{\pgfqpoint{4.052117in}{5.491438in}}%
\pgfpathlineto{\pgfqpoint{4.052117in}{5.491438in}}%
\pgfpathclose%
\pgfusepath{stroke}%
\end{pgfscope}%
\begin{pgfscope}%
\pgfpathrectangle{\pgfqpoint{0.977909in}{2.928000in}}{\pgfqpoint{5.664000in}{5.664000in}}%
\pgfusepath{clip}%
\pgfsetbuttcap%
\pgfsetroundjoin%
\pgfsetlinewidth{1.003750pt}%
\definecolor{currentstroke}{rgb}{0.541176,0.380392,0.658824}%
\pgfsetstrokecolor{currentstroke}%
\pgfsetstrokeopacity{0.900000}%
\pgfsetdash{}{0pt}%
\pgfpathmoveto{\pgfqpoint{4.587894in}{6.097784in}}%
\pgfpathcurveto{\pgfqpoint{4.598148in}{6.097784in}}{\pgfqpoint{4.607984in}{6.101858in}}{\pgfqpoint{4.615235in}{6.109109in}}%
\pgfpathcurveto{\pgfqpoint{4.622485in}{6.116359in}}{\pgfqpoint{4.626559in}{6.126195in}}{\pgfqpoint{4.626559in}{6.136449in}}%
\pgfpathcurveto{\pgfqpoint{4.626559in}{6.146703in}}{\pgfqpoint{4.622485in}{6.156539in}}{\pgfqpoint{4.615235in}{6.163789in}}%
\pgfpathcurveto{\pgfqpoint{4.607984in}{6.171040in}}{\pgfqpoint{4.598148in}{6.175114in}}{\pgfqpoint{4.587894in}{6.175114in}}%
\pgfpathcurveto{\pgfqpoint{4.577640in}{6.175114in}}{\pgfqpoint{4.567805in}{6.171040in}}{\pgfqpoint{4.560554in}{6.163789in}}%
\pgfpathcurveto{\pgfqpoint{4.553303in}{6.156539in}}{\pgfqpoint{4.549229in}{6.146703in}}{\pgfqpoint{4.549229in}{6.136449in}}%
\pgfpathcurveto{\pgfqpoint{4.549229in}{6.126195in}}{\pgfqpoint{4.553303in}{6.116359in}}{\pgfqpoint{4.560554in}{6.109109in}}%
\pgfpathcurveto{\pgfqpoint{4.567805in}{6.101858in}}{\pgfqpoint{4.577640in}{6.097784in}}{\pgfqpoint{4.587894in}{6.097784in}}%
\pgfpathlineto{\pgfqpoint{4.587894in}{6.097784in}}%
\pgfpathclose%
\pgfusepath{stroke}%
\end{pgfscope}%
\begin{pgfscope}%
\pgfpathrectangle{\pgfqpoint{0.977909in}{2.928000in}}{\pgfqpoint{5.664000in}{5.664000in}}%
\pgfusepath{clip}%
\pgfsetbuttcap%
\pgfsetroundjoin%
\pgfsetlinewidth{1.003750pt}%
\definecolor{currentstroke}{rgb}{0.541176,0.380392,0.658824}%
\pgfsetstrokecolor{currentstroke}%
\pgfsetstrokeopacity{0.900000}%
\pgfsetdash{}{0pt}%
\pgfpathmoveto{\pgfqpoint{3.959609in}{5.265231in}}%
\pgfpathcurveto{\pgfqpoint{3.969863in}{5.265231in}}{\pgfqpoint{3.979699in}{5.269305in}}{\pgfqpoint{3.986950in}{5.276556in}}%
\pgfpathcurveto{\pgfqpoint{3.994200in}{5.283806in}}{\pgfqpoint{3.998274in}{5.293642in}}{\pgfqpoint{3.998274in}{5.303896in}}%
\pgfpathcurveto{\pgfqpoint{3.998274in}{5.314150in}}{\pgfqpoint{3.994200in}{5.323985in}}{\pgfqpoint{3.986950in}{5.331236in}}%
\pgfpathcurveto{\pgfqpoint{3.979699in}{5.338487in}}{\pgfqpoint{3.969863in}{5.342561in}}{\pgfqpoint{3.959609in}{5.342561in}}%
\pgfpathcurveto{\pgfqpoint{3.949355in}{5.342561in}}{\pgfqpoint{3.939520in}{5.338487in}}{\pgfqpoint{3.932269in}{5.331236in}}%
\pgfpathcurveto{\pgfqpoint{3.925018in}{5.323985in}}{\pgfqpoint{3.920944in}{5.314150in}}{\pgfqpoint{3.920944in}{5.303896in}}%
\pgfpathcurveto{\pgfqpoint{3.920944in}{5.293642in}}{\pgfqpoint{3.925018in}{5.283806in}}{\pgfqpoint{3.932269in}{5.276556in}}%
\pgfpathcurveto{\pgfqpoint{3.939520in}{5.269305in}}{\pgfqpoint{3.949355in}{5.265231in}}{\pgfqpoint{3.959609in}{5.265231in}}%
\pgfpathlineto{\pgfqpoint{3.959609in}{5.265231in}}%
\pgfpathclose%
\pgfusepath{stroke}%
\end{pgfscope}%
\begin{pgfscope}%
\pgfpathrectangle{\pgfqpoint{0.977909in}{2.928000in}}{\pgfqpoint{5.664000in}{5.664000in}}%
\pgfusepath{clip}%
\pgfsetbuttcap%
\pgfsetroundjoin%
\pgfsetlinewidth{1.003750pt}%
\definecolor{currentstroke}{rgb}{0.541176,0.380392,0.658824}%
\pgfsetstrokecolor{currentstroke}%
\pgfsetstrokeopacity{0.900000}%
\pgfsetdash{}{0pt}%
\pgfpathmoveto{\pgfqpoint{5.305204in}{6.708904in}}%
\pgfpathcurveto{\pgfqpoint{5.315458in}{6.708904in}}{\pgfqpoint{5.325293in}{6.712978in}}{\pgfqpoint{5.332544in}{6.720229in}}%
\pgfpathcurveto{\pgfqpoint{5.339795in}{6.727479in}}{\pgfqpoint{5.343869in}{6.737315in}}{\pgfqpoint{5.343869in}{6.747569in}}%
\pgfpathcurveto{\pgfqpoint{5.343869in}{6.757823in}}{\pgfqpoint{5.339795in}{6.767659in}}{\pgfqpoint{5.332544in}{6.774909in}}%
\pgfpathcurveto{\pgfqpoint{5.325293in}{6.782160in}}{\pgfqpoint{5.315458in}{6.786234in}}{\pgfqpoint{5.305204in}{6.786234in}}%
\pgfpathcurveto{\pgfqpoint{5.294950in}{6.786234in}}{\pgfqpoint{5.285114in}{6.782160in}}{\pgfqpoint{5.277863in}{6.774909in}}%
\pgfpathcurveto{\pgfqpoint{5.270613in}{6.767659in}}{\pgfqpoint{5.266539in}{6.757823in}}{\pgfqpoint{5.266539in}{6.747569in}}%
\pgfpathcurveto{\pgfqpoint{5.266539in}{6.737315in}}{\pgfqpoint{5.270613in}{6.727479in}}{\pgfqpoint{5.277863in}{6.720229in}}%
\pgfpathcurveto{\pgfqpoint{5.285114in}{6.712978in}}{\pgfqpoint{5.294950in}{6.708904in}}{\pgfqpoint{5.305204in}{6.708904in}}%
\pgfpathlineto{\pgfqpoint{5.305204in}{6.708904in}}%
\pgfpathclose%
\pgfusepath{stroke}%
\end{pgfscope}%
\begin{pgfscope}%
\pgfpathrectangle{\pgfqpoint{0.977909in}{2.928000in}}{\pgfqpoint{5.664000in}{5.664000in}}%
\pgfusepath{clip}%
\pgfsetbuttcap%
\pgfsetroundjoin%
\pgfsetlinewidth{1.003750pt}%
\definecolor{currentstroke}{rgb}{0.541176,0.380392,0.658824}%
\pgfsetstrokecolor{currentstroke}%
\pgfsetstrokeopacity{0.900000}%
\pgfsetdash{}{0pt}%
\pgfpathmoveto{\pgfqpoint{4.293573in}{5.766423in}}%
\pgfpathcurveto{\pgfqpoint{4.303827in}{5.766423in}}{\pgfqpoint{4.313663in}{5.770497in}}{\pgfqpoint{4.320914in}{5.777748in}}%
\pgfpathcurveto{\pgfqpoint{4.328164in}{5.784998in}}{\pgfqpoint{4.332238in}{5.794834in}}{\pgfqpoint{4.332238in}{5.805088in}}%
\pgfpathcurveto{\pgfqpoint{4.332238in}{5.815342in}}{\pgfqpoint{4.328164in}{5.825178in}}{\pgfqpoint{4.320914in}{5.832428in}}%
\pgfpathcurveto{\pgfqpoint{4.313663in}{5.839679in}}{\pgfqpoint{4.303827in}{5.843753in}}{\pgfqpoint{4.293573in}{5.843753in}}%
\pgfpathcurveto{\pgfqpoint{4.283319in}{5.843753in}}{\pgfqpoint{4.273484in}{5.839679in}}{\pgfqpoint{4.266233in}{5.832428in}}%
\pgfpathcurveto{\pgfqpoint{4.258982in}{5.825178in}}{\pgfqpoint{4.254908in}{5.815342in}}{\pgfqpoint{4.254908in}{5.805088in}}%
\pgfpathcurveto{\pgfqpoint{4.254908in}{5.794834in}}{\pgfqpoint{4.258982in}{5.784998in}}{\pgfqpoint{4.266233in}{5.777748in}}%
\pgfpathcurveto{\pgfqpoint{4.273484in}{5.770497in}}{\pgfqpoint{4.283319in}{5.766423in}}{\pgfqpoint{4.293573in}{5.766423in}}%
\pgfpathlineto{\pgfqpoint{4.293573in}{5.766423in}}%
\pgfpathclose%
\pgfusepath{stroke}%
\end{pgfscope}%
\begin{pgfscope}%
\pgfpathrectangle{\pgfqpoint{0.977909in}{2.928000in}}{\pgfqpoint{5.664000in}{5.664000in}}%
\pgfusepath{clip}%
\pgfsetbuttcap%
\pgfsetroundjoin%
\pgfsetlinewidth{1.003750pt}%
\definecolor{currentstroke}{rgb}{0.541176,0.380392,0.658824}%
\pgfsetstrokecolor{currentstroke}%
\pgfsetstrokeopacity{0.900000}%
\pgfsetdash{}{0pt}%
\pgfpathmoveto{\pgfqpoint{4.852364in}{5.831458in}}%
\pgfpathcurveto{\pgfqpoint{4.862618in}{5.831458in}}{\pgfqpoint{4.872454in}{5.835532in}}{\pgfqpoint{4.879704in}{5.842783in}}%
\pgfpathcurveto{\pgfqpoint{4.886955in}{5.850033in}}{\pgfqpoint{4.891029in}{5.859869in}}{\pgfqpoint{4.891029in}{5.870123in}}%
\pgfpathcurveto{\pgfqpoint{4.891029in}{5.880377in}}{\pgfqpoint{4.886955in}{5.890213in}}{\pgfqpoint{4.879704in}{5.897463in}}%
\pgfpathcurveto{\pgfqpoint{4.872454in}{5.904714in}}{\pgfqpoint{4.862618in}{5.908788in}}{\pgfqpoint{4.852364in}{5.908788in}}%
\pgfpathcurveto{\pgfqpoint{4.842110in}{5.908788in}}{\pgfqpoint{4.832275in}{5.904714in}}{\pgfqpoint{4.825024in}{5.897463in}}%
\pgfpathcurveto{\pgfqpoint{4.817773in}{5.890213in}}{\pgfqpoint{4.813699in}{5.880377in}}{\pgfqpoint{4.813699in}{5.870123in}}%
\pgfpathcurveto{\pgfqpoint{4.813699in}{5.859869in}}{\pgfqpoint{4.817773in}{5.850033in}}{\pgfqpoint{4.825024in}{5.842783in}}%
\pgfpathcurveto{\pgfqpoint{4.832275in}{5.835532in}}{\pgfqpoint{4.842110in}{5.831458in}}{\pgfqpoint{4.852364in}{5.831458in}}%
\pgfpathlineto{\pgfqpoint{4.852364in}{5.831458in}}%
\pgfpathclose%
\pgfusepath{stroke}%
\end{pgfscope}%
\begin{pgfscope}%
\pgfpathrectangle{\pgfqpoint{0.977909in}{2.928000in}}{\pgfqpoint{5.664000in}{5.664000in}}%
\pgfusepath{clip}%
\pgfsetbuttcap%
\pgfsetroundjoin%
\pgfsetlinewidth{1.003750pt}%
\definecolor{currentstroke}{rgb}{0.541176,0.380392,0.658824}%
\pgfsetstrokecolor{currentstroke}%
\pgfsetstrokeopacity{0.900000}%
\pgfsetdash{}{0pt}%
\pgfpathmoveto{\pgfqpoint{4.938666in}{6.885654in}}%
\pgfpathcurveto{\pgfqpoint{4.948920in}{6.885654in}}{\pgfqpoint{4.958755in}{6.889728in}}{\pgfqpoint{4.966006in}{6.896978in}}%
\pgfpathcurveto{\pgfqpoint{4.973257in}{6.904229in}}{\pgfqpoint{4.977331in}{6.914065in}}{\pgfqpoint{4.977331in}{6.924319in}}%
\pgfpathcurveto{\pgfqpoint{4.977331in}{6.934573in}}{\pgfqpoint{4.973257in}{6.944408in}}{\pgfqpoint{4.966006in}{6.951659in}}%
\pgfpathcurveto{\pgfqpoint{4.958755in}{6.958910in}}{\pgfqpoint{4.948920in}{6.962984in}}{\pgfqpoint{4.938666in}{6.962984in}}%
\pgfpathcurveto{\pgfqpoint{4.928412in}{6.962984in}}{\pgfqpoint{4.918576in}{6.958910in}}{\pgfqpoint{4.911325in}{6.951659in}}%
\pgfpathcurveto{\pgfqpoint{4.904075in}{6.944408in}}{\pgfqpoint{4.900001in}{6.934573in}}{\pgfqpoint{4.900001in}{6.924319in}}%
\pgfpathcurveto{\pgfqpoint{4.900001in}{6.914065in}}{\pgfqpoint{4.904075in}{6.904229in}}{\pgfqpoint{4.911325in}{6.896978in}}%
\pgfpathcurveto{\pgfqpoint{4.918576in}{6.889728in}}{\pgfqpoint{4.928412in}{6.885654in}}{\pgfqpoint{4.938666in}{6.885654in}}%
\pgfpathlineto{\pgfqpoint{4.938666in}{6.885654in}}%
\pgfpathclose%
\pgfusepath{stroke}%
\end{pgfscope}%
\begin{pgfscope}%
\pgfpathrectangle{\pgfqpoint{0.977909in}{2.928000in}}{\pgfqpoint{5.664000in}{5.664000in}}%
\pgfusepath{clip}%
\pgfsetbuttcap%
\pgfsetroundjoin%
\pgfsetlinewidth{1.003750pt}%
\definecolor{currentstroke}{rgb}{0.541176,0.380392,0.658824}%
\pgfsetstrokecolor{currentstroke}%
\pgfsetstrokeopacity{0.900000}%
\pgfsetdash{}{0pt}%
\pgfpathmoveto{\pgfqpoint{3.683854in}{5.946016in}}%
\pgfpathcurveto{\pgfqpoint{3.694108in}{5.946016in}}{\pgfqpoint{3.703943in}{5.950090in}}{\pgfqpoint{3.711194in}{5.957341in}}%
\pgfpathcurveto{\pgfqpoint{3.718445in}{5.964592in}}{\pgfqpoint{3.722519in}{5.974427in}}{\pgfqpoint{3.722519in}{5.984681in}}%
\pgfpathcurveto{\pgfqpoint{3.722519in}{5.994935in}}{\pgfqpoint{3.718445in}{6.004771in}}{\pgfqpoint{3.711194in}{6.012021in}}%
\pgfpathcurveto{\pgfqpoint{3.703943in}{6.019272in}}{\pgfqpoint{3.694108in}{6.023346in}}{\pgfqpoint{3.683854in}{6.023346in}}%
\pgfpathcurveto{\pgfqpoint{3.673600in}{6.023346in}}{\pgfqpoint{3.663764in}{6.019272in}}{\pgfqpoint{3.656513in}{6.012021in}}%
\pgfpathcurveto{\pgfqpoint{3.649263in}{6.004771in}}{\pgfqpoint{3.645189in}{5.994935in}}{\pgfqpoint{3.645189in}{5.984681in}}%
\pgfpathcurveto{\pgfqpoint{3.645189in}{5.974427in}}{\pgfqpoint{3.649263in}{5.964592in}}{\pgfqpoint{3.656513in}{5.957341in}}%
\pgfpathcurveto{\pgfqpoint{3.663764in}{5.950090in}}{\pgfqpoint{3.673600in}{5.946016in}}{\pgfqpoint{3.683854in}{5.946016in}}%
\pgfpathlineto{\pgfqpoint{3.683854in}{5.946016in}}%
\pgfpathclose%
\pgfusepath{stroke}%
\end{pgfscope}%
\begin{pgfscope}%
\pgfpathrectangle{\pgfqpoint{0.977909in}{2.928000in}}{\pgfqpoint{5.664000in}{5.664000in}}%
\pgfusepath{clip}%
\pgfsetbuttcap%
\pgfsetroundjoin%
\pgfsetlinewidth{1.003750pt}%
\definecolor{currentstroke}{rgb}{0.541176,0.380392,0.658824}%
\pgfsetstrokecolor{currentstroke}%
\pgfsetstrokeopacity{0.900000}%
\pgfsetdash{}{0pt}%
\pgfpathmoveto{\pgfqpoint{5.607901in}{6.541198in}}%
\pgfpathcurveto{\pgfqpoint{5.618156in}{6.541198in}}{\pgfqpoint{5.627991in}{6.545272in}}{\pgfqpoint{5.635242in}{6.552522in}}%
\pgfpathcurveto{\pgfqpoint{5.642492in}{6.559773in}}{\pgfqpoint{5.646566in}{6.569609in}}{\pgfqpoint{5.646566in}{6.579863in}}%
\pgfpathcurveto{\pgfqpoint{5.646566in}{6.590117in}}{\pgfqpoint{5.642492in}{6.599952in}}{\pgfqpoint{5.635242in}{6.607203in}}%
\pgfpathcurveto{\pgfqpoint{5.627991in}{6.614454in}}{\pgfqpoint{5.618156in}{6.618528in}}{\pgfqpoint{5.607901in}{6.618528in}}%
\pgfpathcurveto{\pgfqpoint{5.597647in}{6.618528in}}{\pgfqpoint{5.587812in}{6.614454in}}{\pgfqpoint{5.580561in}{6.607203in}}%
\pgfpathcurveto{\pgfqpoint{5.573310in}{6.599952in}}{\pgfqpoint{5.569236in}{6.590117in}}{\pgfqpoint{5.569236in}{6.579863in}}%
\pgfpathcurveto{\pgfqpoint{5.569236in}{6.569609in}}{\pgfqpoint{5.573310in}{6.559773in}}{\pgfqpoint{5.580561in}{6.552522in}}%
\pgfpathcurveto{\pgfqpoint{5.587812in}{6.545272in}}{\pgfqpoint{5.597647in}{6.541198in}}{\pgfqpoint{5.607901in}{6.541198in}}%
\pgfpathlineto{\pgfqpoint{5.607901in}{6.541198in}}%
\pgfpathclose%
\pgfusepath{stroke}%
\end{pgfscope}%
\begin{pgfscope}%
\pgfpathrectangle{\pgfqpoint{0.977909in}{2.928000in}}{\pgfqpoint{5.664000in}{5.664000in}}%
\pgfusepath{clip}%
\pgfsetbuttcap%
\pgfsetroundjoin%
\pgfsetlinewidth{1.003750pt}%
\definecolor{currentstroke}{rgb}{0.541176,0.380392,0.658824}%
\pgfsetstrokecolor{currentstroke}%
\pgfsetstrokeopacity{0.900000}%
\pgfsetdash{}{0pt}%
\pgfpathmoveto{\pgfqpoint{4.037287in}{6.099199in}}%
\pgfpathcurveto{\pgfqpoint{4.047541in}{6.099199in}}{\pgfqpoint{4.057377in}{6.103273in}}{\pgfqpoint{4.064627in}{6.110524in}}%
\pgfpathcurveto{\pgfqpoint{4.071878in}{6.117775in}}{\pgfqpoint{4.075952in}{6.127610in}}{\pgfqpoint{4.075952in}{6.137864in}}%
\pgfpathcurveto{\pgfqpoint{4.075952in}{6.148118in}}{\pgfqpoint{4.071878in}{6.157954in}}{\pgfqpoint{4.064627in}{6.165205in}}%
\pgfpathcurveto{\pgfqpoint{4.057377in}{6.172455in}}{\pgfqpoint{4.047541in}{6.176529in}}{\pgfqpoint{4.037287in}{6.176529in}}%
\pgfpathcurveto{\pgfqpoint{4.027033in}{6.176529in}}{\pgfqpoint{4.017198in}{6.172455in}}{\pgfqpoint{4.009947in}{6.165205in}}%
\pgfpathcurveto{\pgfqpoint{4.002696in}{6.157954in}}{\pgfqpoint{3.998622in}{6.148118in}}{\pgfqpoint{3.998622in}{6.137864in}}%
\pgfpathcurveto{\pgfqpoint{3.998622in}{6.127610in}}{\pgfqpoint{4.002696in}{6.117775in}}{\pgfqpoint{4.009947in}{6.110524in}}%
\pgfpathcurveto{\pgfqpoint{4.017198in}{6.103273in}}{\pgfqpoint{4.027033in}{6.099199in}}{\pgfqpoint{4.037287in}{6.099199in}}%
\pgfpathlineto{\pgfqpoint{4.037287in}{6.099199in}}%
\pgfpathclose%
\pgfusepath{stroke}%
\end{pgfscope}%
\begin{pgfscope}%
\pgfpathrectangle{\pgfqpoint{0.977909in}{2.928000in}}{\pgfqpoint{5.664000in}{5.664000in}}%
\pgfusepath{clip}%
\pgfsetbuttcap%
\pgfsetroundjoin%
\pgfsetlinewidth{1.003750pt}%
\definecolor{currentstroke}{rgb}{0.541176,0.380392,0.658824}%
\pgfsetstrokecolor{currentstroke}%
\pgfsetstrokeopacity{0.900000}%
\pgfsetdash{}{0pt}%
\pgfpathmoveto{\pgfqpoint{3.230992in}{5.135233in}}%
\pgfpathcurveto{\pgfqpoint{3.241246in}{5.135233in}}{\pgfqpoint{3.251082in}{5.139307in}}{\pgfqpoint{3.258332in}{5.146558in}}%
\pgfpathcurveto{\pgfqpoint{3.265583in}{5.153809in}}{\pgfqpoint{3.269657in}{5.163644in}}{\pgfqpoint{3.269657in}{5.173898in}}%
\pgfpathcurveto{\pgfqpoint{3.269657in}{5.184152in}}{\pgfqpoint{3.265583in}{5.193988in}}{\pgfqpoint{3.258332in}{5.201238in}}%
\pgfpathcurveto{\pgfqpoint{3.251082in}{5.208489in}}{\pgfqpoint{3.241246in}{5.212563in}}{\pgfqpoint{3.230992in}{5.212563in}}%
\pgfpathcurveto{\pgfqpoint{3.220738in}{5.212563in}}{\pgfqpoint{3.210902in}{5.208489in}}{\pgfqpoint{3.203652in}{5.201238in}}%
\pgfpathcurveto{\pgfqpoint{3.196401in}{5.193988in}}{\pgfqpoint{3.192327in}{5.184152in}}{\pgfqpoint{3.192327in}{5.173898in}}%
\pgfpathcurveto{\pgfqpoint{3.192327in}{5.163644in}}{\pgfqpoint{3.196401in}{5.153809in}}{\pgfqpoint{3.203652in}{5.146558in}}%
\pgfpathcurveto{\pgfqpoint{3.210902in}{5.139307in}}{\pgfqpoint{3.220738in}{5.135233in}}{\pgfqpoint{3.230992in}{5.135233in}}%
\pgfpathlineto{\pgfqpoint{3.230992in}{5.135233in}}%
\pgfpathclose%
\pgfusepath{stroke}%
\end{pgfscope}%
\begin{pgfscope}%
\pgfpathrectangle{\pgfqpoint{0.977909in}{2.928000in}}{\pgfqpoint{5.664000in}{5.664000in}}%
\pgfusepath{clip}%
\pgfsetbuttcap%
\pgfsetroundjoin%
\pgfsetlinewidth{1.003750pt}%
\definecolor{currentstroke}{rgb}{0.541176,0.380392,0.658824}%
\pgfsetstrokecolor{currentstroke}%
\pgfsetstrokeopacity{0.900000}%
\pgfsetdash{}{0pt}%
\pgfpathmoveto{\pgfqpoint{4.405957in}{5.599941in}}%
\pgfpathcurveto{\pgfqpoint{4.416211in}{5.599941in}}{\pgfqpoint{4.426047in}{5.604015in}}{\pgfqpoint{4.433297in}{5.611266in}}%
\pgfpathcurveto{\pgfqpoint{4.440548in}{5.618516in}}{\pgfqpoint{4.444622in}{5.628352in}}{\pgfqpoint{4.444622in}{5.638606in}}%
\pgfpathcurveto{\pgfqpoint{4.444622in}{5.648860in}}{\pgfqpoint{4.440548in}{5.658696in}}{\pgfqpoint{4.433297in}{5.665946in}}%
\pgfpathcurveto{\pgfqpoint{4.426047in}{5.673197in}}{\pgfqpoint{4.416211in}{5.677271in}}{\pgfqpoint{4.405957in}{5.677271in}}%
\pgfpathcurveto{\pgfqpoint{4.395703in}{5.677271in}}{\pgfqpoint{4.385868in}{5.673197in}}{\pgfqpoint{4.378617in}{5.665946in}}%
\pgfpathcurveto{\pgfqpoint{4.371366in}{5.658696in}}{\pgfqpoint{4.367292in}{5.648860in}}{\pgfqpoint{4.367292in}{5.638606in}}%
\pgfpathcurveto{\pgfqpoint{4.367292in}{5.628352in}}{\pgfqpoint{4.371366in}{5.618516in}}{\pgfqpoint{4.378617in}{5.611266in}}%
\pgfpathcurveto{\pgfqpoint{4.385868in}{5.604015in}}{\pgfqpoint{4.395703in}{5.599941in}}{\pgfqpoint{4.405957in}{5.599941in}}%
\pgfpathlineto{\pgfqpoint{4.405957in}{5.599941in}}%
\pgfpathclose%
\pgfusepath{stroke}%
\end{pgfscope}%
\begin{pgfscope}%
\pgfpathrectangle{\pgfqpoint{0.977909in}{2.928000in}}{\pgfqpoint{5.664000in}{5.664000in}}%
\pgfusepath{clip}%
\pgfsetbuttcap%
\pgfsetroundjoin%
\pgfsetlinewidth{1.003750pt}%
\definecolor{currentstroke}{rgb}{0.541176,0.380392,0.658824}%
\pgfsetstrokecolor{currentstroke}%
\pgfsetstrokeopacity{0.900000}%
\pgfsetdash{}{0pt}%
\pgfpathmoveto{\pgfqpoint{4.355415in}{6.358045in}}%
\pgfpathcurveto{\pgfqpoint{4.365669in}{6.358045in}}{\pgfqpoint{4.375505in}{6.362119in}}{\pgfqpoint{4.382756in}{6.369370in}}%
\pgfpathcurveto{\pgfqpoint{4.390006in}{6.376621in}}{\pgfqpoint{4.394080in}{6.386456in}}{\pgfqpoint{4.394080in}{6.396710in}}%
\pgfpathcurveto{\pgfqpoint{4.394080in}{6.406964in}}{\pgfqpoint{4.390006in}{6.416800in}}{\pgfqpoint{4.382756in}{6.424051in}}%
\pgfpathcurveto{\pgfqpoint{4.375505in}{6.431301in}}{\pgfqpoint{4.365669in}{6.435375in}}{\pgfqpoint{4.355415in}{6.435375in}}%
\pgfpathcurveto{\pgfqpoint{4.345161in}{6.435375in}}{\pgfqpoint{4.335326in}{6.431301in}}{\pgfqpoint{4.328075in}{6.424051in}}%
\pgfpathcurveto{\pgfqpoint{4.320824in}{6.416800in}}{\pgfqpoint{4.316750in}{6.406964in}}{\pgfqpoint{4.316750in}{6.396710in}}%
\pgfpathcurveto{\pgfqpoint{4.316750in}{6.386456in}}{\pgfqpoint{4.320824in}{6.376621in}}{\pgfqpoint{4.328075in}{6.369370in}}%
\pgfpathcurveto{\pgfqpoint{4.335326in}{6.362119in}}{\pgfqpoint{4.345161in}{6.358045in}}{\pgfqpoint{4.355415in}{6.358045in}}%
\pgfpathlineto{\pgfqpoint{4.355415in}{6.358045in}}%
\pgfpathclose%
\pgfusepath{stroke}%
\end{pgfscope}%
\begin{pgfscope}%
\pgfpathrectangle{\pgfqpoint{0.977909in}{2.928000in}}{\pgfqpoint{5.664000in}{5.664000in}}%
\pgfusepath{clip}%
\pgfsetbuttcap%
\pgfsetroundjoin%
\pgfsetlinewidth{1.003750pt}%
\definecolor{currentstroke}{rgb}{0.541176,0.380392,0.658824}%
\pgfsetstrokecolor{currentstroke}%
\pgfsetstrokeopacity{0.900000}%
\pgfsetdash{}{0pt}%
\pgfpathmoveto{\pgfqpoint{4.408546in}{5.619322in}}%
\pgfpathcurveto{\pgfqpoint{4.418800in}{5.619322in}}{\pgfqpoint{4.428636in}{5.623396in}}{\pgfqpoint{4.435886in}{5.630647in}}%
\pgfpathcurveto{\pgfqpoint{4.443137in}{5.637897in}}{\pgfqpoint{4.447211in}{5.647733in}}{\pgfqpoint{4.447211in}{5.657987in}}%
\pgfpathcurveto{\pgfqpoint{4.447211in}{5.668241in}}{\pgfqpoint{4.443137in}{5.678076in}}{\pgfqpoint{4.435886in}{5.685327in}}%
\pgfpathcurveto{\pgfqpoint{4.428636in}{5.692578in}}{\pgfqpoint{4.418800in}{5.696652in}}{\pgfqpoint{4.408546in}{5.696652in}}%
\pgfpathcurveto{\pgfqpoint{4.398292in}{5.696652in}}{\pgfqpoint{4.388457in}{5.692578in}}{\pgfqpoint{4.381206in}{5.685327in}}%
\pgfpathcurveto{\pgfqpoint{4.373955in}{5.678076in}}{\pgfqpoint{4.369881in}{5.668241in}}{\pgfqpoint{4.369881in}{5.657987in}}%
\pgfpathcurveto{\pgfqpoint{4.369881in}{5.647733in}}{\pgfqpoint{4.373955in}{5.637897in}}{\pgfqpoint{4.381206in}{5.630647in}}%
\pgfpathcurveto{\pgfqpoint{4.388457in}{5.623396in}}{\pgfqpoint{4.398292in}{5.619322in}}{\pgfqpoint{4.408546in}{5.619322in}}%
\pgfpathlineto{\pgfqpoint{4.408546in}{5.619322in}}%
\pgfpathclose%
\pgfusepath{stroke}%
\end{pgfscope}%
\begin{pgfscope}%
\pgfpathrectangle{\pgfqpoint{0.977909in}{2.928000in}}{\pgfqpoint{5.664000in}{5.664000in}}%
\pgfusepath{clip}%
\pgfsetbuttcap%
\pgfsetroundjoin%
\pgfsetlinewidth{1.003750pt}%
\definecolor{currentstroke}{rgb}{0.541176,0.380392,0.658824}%
\pgfsetstrokecolor{currentstroke}%
\pgfsetstrokeopacity{0.900000}%
\pgfsetdash{}{0pt}%
\pgfpathmoveto{\pgfqpoint{4.387972in}{5.687708in}}%
\pgfpathcurveto{\pgfqpoint{4.398226in}{5.687708in}}{\pgfqpoint{4.408061in}{5.691782in}}{\pgfqpoint{4.415312in}{5.699032in}}%
\pgfpathcurveto{\pgfqpoint{4.422563in}{5.706283in}}{\pgfqpoint{4.426637in}{5.716118in}}{\pgfqpoint{4.426637in}{5.726373in}}%
\pgfpathcurveto{\pgfqpoint{4.426637in}{5.736627in}}{\pgfqpoint{4.422563in}{5.746462in}}{\pgfqpoint{4.415312in}{5.753713in}}%
\pgfpathcurveto{\pgfqpoint{4.408061in}{5.760964in}}{\pgfqpoint{4.398226in}{5.765038in}}{\pgfqpoint{4.387972in}{5.765038in}}%
\pgfpathcurveto{\pgfqpoint{4.377718in}{5.765038in}}{\pgfqpoint{4.367882in}{5.760964in}}{\pgfqpoint{4.360631in}{5.753713in}}%
\pgfpathcurveto{\pgfqpoint{4.353381in}{5.746462in}}{\pgfqpoint{4.349307in}{5.736627in}}{\pgfqpoint{4.349307in}{5.726373in}}%
\pgfpathcurveto{\pgfqpoint{4.349307in}{5.716118in}}{\pgfqpoint{4.353381in}{5.706283in}}{\pgfqpoint{4.360631in}{5.699032in}}%
\pgfpathcurveto{\pgfqpoint{4.367882in}{5.691782in}}{\pgfqpoint{4.377718in}{5.687708in}}{\pgfqpoint{4.387972in}{5.687708in}}%
\pgfpathlineto{\pgfqpoint{4.387972in}{5.687708in}}%
\pgfpathclose%
\pgfusepath{stroke}%
\end{pgfscope}%
\begin{pgfscope}%
\pgfpathrectangle{\pgfqpoint{0.977909in}{2.928000in}}{\pgfqpoint{5.664000in}{5.664000in}}%
\pgfusepath{clip}%
\pgfsetbuttcap%
\pgfsetroundjoin%
\pgfsetlinewidth{1.003750pt}%
\definecolor{currentstroke}{rgb}{0.541176,0.380392,0.658824}%
\pgfsetstrokecolor{currentstroke}%
\pgfsetstrokeopacity{0.900000}%
\pgfsetdash{}{0pt}%
\pgfpathmoveto{\pgfqpoint{4.665367in}{6.867129in}}%
\pgfpathcurveto{\pgfqpoint{4.675621in}{6.867129in}}{\pgfqpoint{4.685456in}{6.871203in}}{\pgfqpoint{4.692707in}{6.878453in}}%
\pgfpathcurveto{\pgfqpoint{4.699958in}{6.885704in}}{\pgfqpoint{4.704032in}{6.895540in}}{\pgfqpoint{4.704032in}{6.905794in}}%
\pgfpathcurveto{\pgfqpoint{4.704032in}{6.916048in}}{\pgfqpoint{4.699958in}{6.925883in}}{\pgfqpoint{4.692707in}{6.933134in}}%
\pgfpathcurveto{\pgfqpoint{4.685456in}{6.940385in}}{\pgfqpoint{4.675621in}{6.944459in}}{\pgfqpoint{4.665367in}{6.944459in}}%
\pgfpathcurveto{\pgfqpoint{4.655113in}{6.944459in}}{\pgfqpoint{4.645277in}{6.940385in}}{\pgfqpoint{4.638027in}{6.933134in}}%
\pgfpathcurveto{\pgfqpoint{4.630776in}{6.925883in}}{\pgfqpoint{4.626702in}{6.916048in}}{\pgfqpoint{4.626702in}{6.905794in}}%
\pgfpathcurveto{\pgfqpoint{4.626702in}{6.895540in}}{\pgfqpoint{4.630776in}{6.885704in}}{\pgfqpoint{4.638027in}{6.878453in}}%
\pgfpathcurveto{\pgfqpoint{4.645277in}{6.871203in}}{\pgfqpoint{4.655113in}{6.867129in}}{\pgfqpoint{4.665367in}{6.867129in}}%
\pgfpathlineto{\pgfqpoint{4.665367in}{6.867129in}}%
\pgfpathclose%
\pgfusepath{stroke}%
\end{pgfscope}%
\begin{pgfscope}%
\pgfpathrectangle{\pgfqpoint{0.977909in}{2.928000in}}{\pgfqpoint{5.664000in}{5.664000in}}%
\pgfusepath{clip}%
\pgfsetbuttcap%
\pgfsetroundjoin%
\pgfsetlinewidth{1.003750pt}%
\definecolor{currentstroke}{rgb}{0.541176,0.380392,0.658824}%
\pgfsetstrokecolor{currentstroke}%
\pgfsetstrokeopacity{0.900000}%
\pgfsetdash{}{0pt}%
\pgfpathmoveto{\pgfqpoint{4.394838in}{6.373884in}}%
\pgfpathcurveto{\pgfqpoint{4.405092in}{6.373884in}}{\pgfqpoint{4.414927in}{6.377958in}}{\pgfqpoint{4.422178in}{6.385208in}}%
\pgfpathcurveto{\pgfqpoint{4.429429in}{6.392459in}}{\pgfqpoint{4.433503in}{6.402295in}}{\pgfqpoint{4.433503in}{6.412549in}}%
\pgfpathcurveto{\pgfqpoint{4.433503in}{6.422803in}}{\pgfqpoint{4.429429in}{6.432638in}}{\pgfqpoint{4.422178in}{6.439889in}}%
\pgfpathcurveto{\pgfqpoint{4.414927in}{6.447140in}}{\pgfqpoint{4.405092in}{6.451214in}}{\pgfqpoint{4.394838in}{6.451214in}}%
\pgfpathcurveto{\pgfqpoint{4.384584in}{6.451214in}}{\pgfqpoint{4.374748in}{6.447140in}}{\pgfqpoint{4.367498in}{6.439889in}}%
\pgfpathcurveto{\pgfqpoint{4.360247in}{6.432638in}}{\pgfqpoint{4.356173in}{6.422803in}}{\pgfqpoint{4.356173in}{6.412549in}}%
\pgfpathcurveto{\pgfqpoint{4.356173in}{6.402295in}}{\pgfqpoint{4.360247in}{6.392459in}}{\pgfqpoint{4.367498in}{6.385208in}}%
\pgfpathcurveto{\pgfqpoint{4.374748in}{6.377958in}}{\pgfqpoint{4.384584in}{6.373884in}}{\pgfqpoint{4.394838in}{6.373884in}}%
\pgfpathlineto{\pgfqpoint{4.394838in}{6.373884in}}%
\pgfpathclose%
\pgfusepath{stroke}%
\end{pgfscope}%
\begin{pgfscope}%
\pgfpathrectangle{\pgfqpoint{0.977909in}{2.928000in}}{\pgfqpoint{5.664000in}{5.664000in}}%
\pgfusepath{clip}%
\pgfsetbuttcap%
\pgfsetroundjoin%
\pgfsetlinewidth{1.003750pt}%
\definecolor{currentstroke}{rgb}{0.541176,0.380392,0.658824}%
\pgfsetstrokecolor{currentstroke}%
\pgfsetstrokeopacity{0.900000}%
\pgfsetdash{}{0pt}%
\pgfpathmoveto{\pgfqpoint{4.839284in}{7.001671in}}%
\pgfpathcurveto{\pgfqpoint{4.849538in}{7.001671in}}{\pgfqpoint{4.859373in}{7.005745in}}{\pgfqpoint{4.866624in}{7.012996in}}%
\pgfpathcurveto{\pgfqpoint{4.873875in}{7.020246in}}{\pgfqpoint{4.877949in}{7.030082in}}{\pgfqpoint{4.877949in}{7.040336in}}%
\pgfpathcurveto{\pgfqpoint{4.877949in}{7.050590in}}{\pgfqpoint{4.873875in}{7.060426in}}{\pgfqpoint{4.866624in}{7.067676in}}%
\pgfpathcurveto{\pgfqpoint{4.859373in}{7.074927in}}{\pgfqpoint{4.849538in}{7.079001in}}{\pgfqpoint{4.839284in}{7.079001in}}%
\pgfpathcurveto{\pgfqpoint{4.829030in}{7.079001in}}{\pgfqpoint{4.819194in}{7.074927in}}{\pgfqpoint{4.811943in}{7.067676in}}%
\pgfpathcurveto{\pgfqpoint{4.804693in}{7.060426in}}{\pgfqpoint{4.800619in}{7.050590in}}{\pgfqpoint{4.800619in}{7.040336in}}%
\pgfpathcurveto{\pgfqpoint{4.800619in}{7.030082in}}{\pgfqpoint{4.804693in}{7.020246in}}{\pgfqpoint{4.811943in}{7.012996in}}%
\pgfpathcurveto{\pgfqpoint{4.819194in}{7.005745in}}{\pgfqpoint{4.829030in}{7.001671in}}{\pgfqpoint{4.839284in}{7.001671in}}%
\pgfpathlineto{\pgfqpoint{4.839284in}{7.001671in}}%
\pgfpathclose%
\pgfusepath{stroke}%
\end{pgfscope}%
\begin{pgfscope}%
\pgfpathrectangle{\pgfqpoint{0.977909in}{2.928000in}}{\pgfqpoint{5.664000in}{5.664000in}}%
\pgfusepath{clip}%
\pgfsetbuttcap%
\pgfsetroundjoin%
\pgfsetlinewidth{1.003750pt}%
\definecolor{currentstroke}{rgb}{0.541176,0.380392,0.658824}%
\pgfsetstrokecolor{currentstroke}%
\pgfsetstrokeopacity{0.900000}%
\pgfsetdash{}{0pt}%
\pgfpathmoveto{\pgfqpoint{3.404876in}{4.758469in}}%
\pgfpathcurveto{\pgfqpoint{3.415130in}{4.758469in}}{\pgfqpoint{3.424966in}{4.762543in}}{\pgfqpoint{3.432216in}{4.769794in}}%
\pgfpathcurveto{\pgfqpoint{3.439467in}{4.777045in}}{\pgfqpoint{3.443541in}{4.786880in}}{\pgfqpoint{3.443541in}{4.797134in}}%
\pgfpathcurveto{\pgfqpoint{3.443541in}{4.807389in}}{\pgfqpoint{3.439467in}{4.817224in}}{\pgfqpoint{3.432216in}{4.824475in}}%
\pgfpathcurveto{\pgfqpoint{3.424966in}{4.831725in}}{\pgfqpoint{3.415130in}{4.835799in}}{\pgfqpoint{3.404876in}{4.835799in}}%
\pgfpathcurveto{\pgfqpoint{3.394622in}{4.835799in}}{\pgfqpoint{3.384787in}{4.831725in}}{\pgfqpoint{3.377536in}{4.824475in}}%
\pgfpathcurveto{\pgfqpoint{3.370285in}{4.817224in}}{\pgfqpoint{3.366211in}{4.807389in}}{\pgfqpoint{3.366211in}{4.797134in}}%
\pgfpathcurveto{\pgfqpoint{3.366211in}{4.786880in}}{\pgfqpoint{3.370285in}{4.777045in}}{\pgfqpoint{3.377536in}{4.769794in}}%
\pgfpathcurveto{\pgfqpoint{3.384787in}{4.762543in}}{\pgfqpoint{3.394622in}{4.758469in}}{\pgfqpoint{3.404876in}{4.758469in}}%
\pgfpathlineto{\pgfqpoint{3.404876in}{4.758469in}}%
\pgfpathclose%
\pgfusepath{stroke}%
\end{pgfscope}%
\begin{pgfscope}%
\pgfpathrectangle{\pgfqpoint{0.977909in}{2.928000in}}{\pgfqpoint{5.664000in}{5.664000in}}%
\pgfusepath{clip}%
\pgfsetbuttcap%
\pgfsetroundjoin%
\pgfsetlinewidth{1.003750pt}%
\definecolor{currentstroke}{rgb}{0.541176,0.380392,0.658824}%
\pgfsetstrokecolor{currentstroke}%
\pgfsetstrokeopacity{0.900000}%
\pgfsetdash{}{0pt}%
\pgfpathmoveto{\pgfqpoint{4.775949in}{6.293959in}}%
\pgfpathcurveto{\pgfqpoint{4.786203in}{6.293959in}}{\pgfqpoint{4.796039in}{6.298033in}}{\pgfqpoint{4.803290in}{6.305284in}}%
\pgfpathcurveto{\pgfqpoint{4.810540in}{6.312535in}}{\pgfqpoint{4.814614in}{6.322370in}}{\pgfqpoint{4.814614in}{6.332624in}}%
\pgfpathcurveto{\pgfqpoint{4.814614in}{6.342878in}}{\pgfqpoint{4.810540in}{6.352714in}}{\pgfqpoint{4.803290in}{6.359965in}}%
\pgfpathcurveto{\pgfqpoint{4.796039in}{6.367215in}}{\pgfqpoint{4.786203in}{6.371289in}}{\pgfqpoint{4.775949in}{6.371289in}}%
\pgfpathcurveto{\pgfqpoint{4.765695in}{6.371289in}}{\pgfqpoint{4.755860in}{6.367215in}}{\pgfqpoint{4.748609in}{6.359965in}}%
\pgfpathcurveto{\pgfqpoint{4.741358in}{6.352714in}}{\pgfqpoint{4.737284in}{6.342878in}}{\pgfqpoint{4.737284in}{6.332624in}}%
\pgfpathcurveto{\pgfqpoint{4.737284in}{6.322370in}}{\pgfqpoint{4.741358in}{6.312535in}}{\pgfqpoint{4.748609in}{6.305284in}}%
\pgfpathcurveto{\pgfqpoint{4.755860in}{6.298033in}}{\pgfqpoint{4.765695in}{6.293959in}}{\pgfqpoint{4.775949in}{6.293959in}}%
\pgfpathlineto{\pgfqpoint{4.775949in}{6.293959in}}%
\pgfpathclose%
\pgfusepath{stroke}%
\end{pgfscope}%
\begin{pgfscope}%
\pgfsetbuttcap%
\pgfsetroundjoin%
\pgfsetlinewidth{1.003750pt}%
\definecolor{currentstroke}{rgb}{0.541176,0.380392,0.658824}%
\pgfsetstrokecolor{currentstroke}%
\pgfsetstrokeopacity{0.900000}%
\pgfsetdash{}{0pt}%
\pgfpathmoveto{\pgfqpoint{6.641909in}{8.532262in}}%
\pgfpathlineto{\pgfqpoint{6.701647in}{8.592000in}}%
\pgfpathlineto{\pgfqpoint{6.641909in}{8.651738in}}%
\pgfpathlineto{\pgfqpoint{6.582171in}{8.592000in}}%
\pgfpathlineto{\pgfqpoint{6.641909in}{8.532262in}}%
\pgfpathclose%
\pgfusepath{stroke}%
\end{pgfscope}%
\begin{pgfscope}%
\pgfpathrectangle{\pgfqpoint{0.977909in}{2.928000in}}{\pgfqpoint{5.664000in}{5.664000in}}%
\pgfusepath{clip}%
\pgfsetbuttcap%
\pgfsetroundjoin%
\pgfsetlinewidth{1.003750pt}%
\definecolor{currentstroke}{rgb}{0.321569,0.321569,0.321569}%
\pgfsetstrokecolor{currentstroke}%
\pgfsetstrokeopacity{0.900000}%
\pgfsetdash{}{0pt}%
\pgfpathmoveto{\pgfqpoint{2.798558in}{4.758553in}}%
\pgfpathcurveto{\pgfqpoint{2.808812in}{4.758553in}}{\pgfqpoint{2.818647in}{4.762627in}}{\pgfqpoint{2.825898in}{4.769878in}}%
\pgfpathcurveto{\pgfqpoint{2.833149in}{4.777129in}}{\pgfqpoint{2.837223in}{4.786964in}}{\pgfqpoint{2.837223in}{4.797218in}}%
\pgfpathcurveto{\pgfqpoint{2.837223in}{4.807472in}}{\pgfqpoint{2.833149in}{4.817308in}}{\pgfqpoint{2.825898in}{4.824559in}}%
\pgfpathcurveto{\pgfqpoint{2.818647in}{4.831809in}}{\pgfqpoint{2.808812in}{4.835883in}}{\pgfqpoint{2.798558in}{4.835883in}}%
\pgfpathcurveto{\pgfqpoint{2.788304in}{4.835883in}}{\pgfqpoint{2.778468in}{4.831809in}}{\pgfqpoint{2.771217in}{4.824559in}}%
\pgfpathcurveto{\pgfqpoint{2.763967in}{4.817308in}}{\pgfqpoint{2.759893in}{4.807472in}}{\pgfqpoint{2.759893in}{4.797218in}}%
\pgfpathcurveto{\pgfqpoint{2.759893in}{4.786964in}}{\pgfqpoint{2.763967in}{4.777129in}}{\pgfqpoint{2.771217in}{4.769878in}}%
\pgfpathcurveto{\pgfqpoint{2.778468in}{4.762627in}}{\pgfqpoint{2.788304in}{4.758553in}}{\pgfqpoint{2.798558in}{4.758553in}}%
\pgfpathlineto{\pgfqpoint{2.798558in}{4.758553in}}%
\pgfpathclose%
\pgfusepath{stroke}%
\end{pgfscope}%
\begin{pgfscope}%
\pgfpathrectangle{\pgfqpoint{0.977909in}{2.928000in}}{\pgfqpoint{5.664000in}{5.664000in}}%
\pgfusepath{clip}%
\pgfsetbuttcap%
\pgfsetroundjoin%
\pgfsetlinewidth{1.003750pt}%
\definecolor{currentstroke}{rgb}{0.321569,0.321569,0.321569}%
\pgfsetstrokecolor{currentstroke}%
\pgfsetstrokeopacity{0.900000}%
\pgfsetdash{}{0pt}%
\pgfpathmoveto{\pgfqpoint{3.111602in}{4.979750in}}%
\pgfpathcurveto{\pgfqpoint{3.121856in}{4.979750in}}{\pgfqpoint{3.131692in}{4.983824in}}{\pgfqpoint{3.138943in}{4.991074in}}%
\pgfpathcurveto{\pgfqpoint{3.146193in}{4.998325in}}{\pgfqpoint{3.150267in}{5.008161in}}{\pgfqpoint{3.150267in}{5.018415in}}%
\pgfpathcurveto{\pgfqpoint{3.150267in}{5.028669in}}{\pgfqpoint{3.146193in}{5.038504in}}{\pgfqpoint{3.138943in}{5.045755in}}%
\pgfpathcurveto{\pgfqpoint{3.131692in}{5.053006in}}{\pgfqpoint{3.121856in}{5.057080in}}{\pgfqpoint{3.111602in}{5.057080in}}%
\pgfpathcurveto{\pgfqpoint{3.101348in}{5.057080in}}{\pgfqpoint{3.091513in}{5.053006in}}{\pgfqpoint{3.084262in}{5.045755in}}%
\pgfpathcurveto{\pgfqpoint{3.077011in}{5.038504in}}{\pgfqpoint{3.072937in}{5.028669in}}{\pgfqpoint{3.072937in}{5.018415in}}%
\pgfpathcurveto{\pgfqpoint{3.072937in}{5.008161in}}{\pgfqpoint{3.077011in}{4.998325in}}{\pgfqpoint{3.084262in}{4.991074in}}%
\pgfpathcurveto{\pgfqpoint{3.091513in}{4.983824in}}{\pgfqpoint{3.101348in}{4.979750in}}{\pgfqpoint{3.111602in}{4.979750in}}%
\pgfpathlineto{\pgfqpoint{3.111602in}{4.979750in}}%
\pgfpathclose%
\pgfusepath{stroke}%
\end{pgfscope}%
\begin{pgfscope}%
\pgfpathrectangle{\pgfqpoint{0.977909in}{2.928000in}}{\pgfqpoint{5.664000in}{5.664000in}}%
\pgfusepath{clip}%
\pgfsetbuttcap%
\pgfsetroundjoin%
\pgfsetlinewidth{1.003750pt}%
\definecolor{currentstroke}{rgb}{0.321569,0.321569,0.321569}%
\pgfsetstrokecolor{currentstroke}%
\pgfsetstrokeopacity{0.900000}%
\pgfsetdash{}{0pt}%
\pgfpathmoveto{\pgfqpoint{3.845101in}{5.652485in}}%
\pgfpathcurveto{\pgfqpoint{3.855355in}{5.652485in}}{\pgfqpoint{3.865191in}{5.656559in}}{\pgfqpoint{3.872441in}{5.663810in}}%
\pgfpathcurveto{\pgfqpoint{3.879692in}{5.671061in}}{\pgfqpoint{3.883766in}{5.680896in}}{\pgfqpoint{3.883766in}{5.691150in}}%
\pgfpathcurveto{\pgfqpoint{3.883766in}{5.701404in}}{\pgfqpoint{3.879692in}{5.711240in}}{\pgfqpoint{3.872441in}{5.718490in}}%
\pgfpathcurveto{\pgfqpoint{3.865191in}{5.725741in}}{\pgfqpoint{3.855355in}{5.729815in}}{\pgfqpoint{3.845101in}{5.729815in}}%
\pgfpathcurveto{\pgfqpoint{3.834847in}{5.729815in}}{\pgfqpoint{3.825012in}{5.725741in}}{\pgfqpoint{3.817761in}{5.718490in}}%
\pgfpathcurveto{\pgfqpoint{3.810510in}{5.711240in}}{\pgfqpoint{3.806436in}{5.701404in}}{\pgfqpoint{3.806436in}{5.691150in}}%
\pgfpathcurveto{\pgfqpoint{3.806436in}{5.680896in}}{\pgfqpoint{3.810510in}{5.671061in}}{\pgfqpoint{3.817761in}{5.663810in}}%
\pgfpathcurveto{\pgfqpoint{3.825012in}{5.656559in}}{\pgfqpoint{3.834847in}{5.652485in}}{\pgfqpoint{3.845101in}{5.652485in}}%
\pgfpathlineto{\pgfqpoint{3.845101in}{5.652485in}}%
\pgfpathclose%
\pgfusepath{stroke}%
\end{pgfscope}%
\begin{pgfscope}%
\pgfpathrectangle{\pgfqpoint{0.977909in}{2.928000in}}{\pgfqpoint{5.664000in}{5.664000in}}%
\pgfusepath{clip}%
\pgfsetbuttcap%
\pgfsetroundjoin%
\pgfsetlinewidth{1.003750pt}%
\definecolor{currentstroke}{rgb}{0.321569,0.321569,0.321569}%
\pgfsetstrokecolor{currentstroke}%
\pgfsetstrokeopacity{0.900000}%
\pgfsetdash{}{0pt}%
\pgfpathmoveto{\pgfqpoint{3.165675in}{5.019816in}}%
\pgfpathcurveto{\pgfqpoint{3.175929in}{5.019816in}}{\pgfqpoint{3.185764in}{5.023890in}}{\pgfqpoint{3.193015in}{5.031141in}}%
\pgfpathcurveto{\pgfqpoint{3.200266in}{5.038392in}}{\pgfqpoint{3.204340in}{5.048227in}}{\pgfqpoint{3.204340in}{5.058481in}}%
\pgfpathcurveto{\pgfqpoint{3.204340in}{5.068735in}}{\pgfqpoint{3.200266in}{5.078571in}}{\pgfqpoint{3.193015in}{5.085822in}}%
\pgfpathcurveto{\pgfqpoint{3.185764in}{5.093072in}}{\pgfqpoint{3.175929in}{5.097146in}}{\pgfqpoint{3.165675in}{5.097146in}}%
\pgfpathcurveto{\pgfqpoint{3.155420in}{5.097146in}}{\pgfqpoint{3.145585in}{5.093072in}}{\pgfqpoint{3.138334in}{5.085822in}}%
\pgfpathcurveto{\pgfqpoint{3.131083in}{5.078571in}}{\pgfqpoint{3.127010in}{5.068735in}}{\pgfqpoint{3.127010in}{5.058481in}}%
\pgfpathcurveto{\pgfqpoint{3.127010in}{5.048227in}}{\pgfqpoint{3.131083in}{5.038392in}}{\pgfqpoint{3.138334in}{5.031141in}}%
\pgfpathcurveto{\pgfqpoint{3.145585in}{5.023890in}}{\pgfqpoint{3.155420in}{5.019816in}}{\pgfqpoint{3.165675in}{5.019816in}}%
\pgfpathlineto{\pgfqpoint{3.165675in}{5.019816in}}%
\pgfpathclose%
\pgfusepath{stroke}%
\end{pgfscope}%
\begin{pgfscope}%
\pgfpathrectangle{\pgfqpoint{0.977909in}{2.928000in}}{\pgfqpoint{5.664000in}{5.664000in}}%
\pgfusepath{clip}%
\pgfsetbuttcap%
\pgfsetroundjoin%
\pgfsetlinewidth{1.003750pt}%
\definecolor{currentstroke}{rgb}{0.321569,0.321569,0.321569}%
\pgfsetstrokecolor{currentstroke}%
\pgfsetstrokeopacity{0.900000}%
\pgfsetdash{}{0pt}%
\pgfpathmoveto{\pgfqpoint{3.594171in}{5.288571in}}%
\pgfpathcurveto{\pgfqpoint{3.604425in}{5.288571in}}{\pgfqpoint{3.614261in}{5.292645in}}{\pgfqpoint{3.621512in}{5.299896in}}%
\pgfpathcurveto{\pgfqpoint{3.628762in}{5.307147in}}{\pgfqpoint{3.632836in}{5.316982in}}{\pgfqpoint{3.632836in}{5.327236in}}%
\pgfpathcurveto{\pgfqpoint{3.632836in}{5.337491in}}{\pgfqpoint{3.628762in}{5.347326in}}{\pgfqpoint{3.621512in}{5.354577in}}%
\pgfpathcurveto{\pgfqpoint{3.614261in}{5.361827in}}{\pgfqpoint{3.604425in}{5.365901in}}{\pgfqpoint{3.594171in}{5.365901in}}%
\pgfpathcurveto{\pgfqpoint{3.583917in}{5.365901in}}{\pgfqpoint{3.574082in}{5.361827in}}{\pgfqpoint{3.566831in}{5.354577in}}%
\pgfpathcurveto{\pgfqpoint{3.559580in}{5.347326in}}{\pgfqpoint{3.555506in}{5.337491in}}{\pgfqpoint{3.555506in}{5.327236in}}%
\pgfpathcurveto{\pgfqpoint{3.555506in}{5.316982in}}{\pgfqpoint{3.559580in}{5.307147in}}{\pgfqpoint{3.566831in}{5.299896in}}%
\pgfpathcurveto{\pgfqpoint{3.574082in}{5.292645in}}{\pgfqpoint{3.583917in}{5.288571in}}{\pgfqpoint{3.594171in}{5.288571in}}%
\pgfpathlineto{\pgfqpoint{3.594171in}{5.288571in}}%
\pgfpathclose%
\pgfusepath{stroke}%
\end{pgfscope}%
\begin{pgfscope}%
\pgfpathrectangle{\pgfqpoint{0.977909in}{2.928000in}}{\pgfqpoint{5.664000in}{5.664000in}}%
\pgfusepath{clip}%
\pgfsetbuttcap%
\pgfsetroundjoin%
\pgfsetlinewidth{1.003750pt}%
\definecolor{currentstroke}{rgb}{0.321569,0.321569,0.321569}%
\pgfsetstrokecolor{currentstroke}%
\pgfsetstrokeopacity{0.900000}%
\pgfsetdash{}{0pt}%
\pgfpathmoveto{\pgfqpoint{2.752270in}{4.824274in}}%
\pgfpathcurveto{\pgfqpoint{2.762524in}{4.824274in}}{\pgfqpoint{2.772360in}{4.828348in}}{\pgfqpoint{2.779610in}{4.835599in}}%
\pgfpathcurveto{\pgfqpoint{2.786861in}{4.842850in}}{\pgfqpoint{2.790935in}{4.852685in}}{\pgfqpoint{2.790935in}{4.862939in}}%
\pgfpathcurveto{\pgfqpoint{2.790935in}{4.873193in}}{\pgfqpoint{2.786861in}{4.883029in}}{\pgfqpoint{2.779610in}{4.890279in}}%
\pgfpathcurveto{\pgfqpoint{2.772360in}{4.897530in}}{\pgfqpoint{2.762524in}{4.901604in}}{\pgfqpoint{2.752270in}{4.901604in}}%
\pgfpathcurveto{\pgfqpoint{2.742016in}{4.901604in}}{\pgfqpoint{2.732180in}{4.897530in}}{\pgfqpoint{2.724930in}{4.890279in}}%
\pgfpathcurveto{\pgfqpoint{2.717679in}{4.883029in}}{\pgfqpoint{2.713605in}{4.873193in}}{\pgfqpoint{2.713605in}{4.862939in}}%
\pgfpathcurveto{\pgfqpoint{2.713605in}{4.852685in}}{\pgfqpoint{2.717679in}{4.842850in}}{\pgfqpoint{2.724930in}{4.835599in}}%
\pgfpathcurveto{\pgfqpoint{2.732180in}{4.828348in}}{\pgfqpoint{2.742016in}{4.824274in}}{\pgfqpoint{2.752270in}{4.824274in}}%
\pgfpathlineto{\pgfqpoint{2.752270in}{4.824274in}}%
\pgfpathclose%
\pgfusepath{stroke}%
\end{pgfscope}%
\begin{pgfscope}%
\pgfpathrectangle{\pgfqpoint{0.977909in}{2.928000in}}{\pgfqpoint{5.664000in}{5.664000in}}%
\pgfusepath{clip}%
\pgfsetbuttcap%
\pgfsetroundjoin%
\pgfsetlinewidth{1.003750pt}%
\definecolor{currentstroke}{rgb}{0.321569,0.321569,0.321569}%
\pgfsetstrokecolor{currentstroke}%
\pgfsetstrokeopacity{0.900000}%
\pgfsetdash{}{0pt}%
\pgfpathmoveto{\pgfqpoint{3.346145in}{5.147620in}}%
\pgfpathcurveto{\pgfqpoint{3.356400in}{5.147620in}}{\pgfqpoint{3.366235in}{5.151694in}}{\pgfqpoint{3.373486in}{5.158944in}}%
\pgfpathcurveto{\pgfqpoint{3.380737in}{5.166195in}}{\pgfqpoint{3.384811in}{5.176031in}}{\pgfqpoint{3.384811in}{5.186285in}}%
\pgfpathcurveto{\pgfqpoint{3.384811in}{5.196539in}}{\pgfqpoint{3.380737in}{5.206374in}}{\pgfqpoint{3.373486in}{5.213625in}}%
\pgfpathcurveto{\pgfqpoint{3.366235in}{5.220876in}}{\pgfqpoint{3.356400in}{5.224950in}}{\pgfqpoint{3.346145in}{5.224950in}}%
\pgfpathcurveto{\pgfqpoint{3.335891in}{5.224950in}}{\pgfqpoint{3.326056in}{5.220876in}}{\pgfqpoint{3.318805in}{5.213625in}}%
\pgfpathcurveto{\pgfqpoint{3.311554in}{5.206374in}}{\pgfqpoint{3.307480in}{5.196539in}}{\pgfqpoint{3.307480in}{5.186285in}}%
\pgfpathcurveto{\pgfqpoint{3.307480in}{5.176031in}}{\pgfqpoint{3.311554in}{5.166195in}}{\pgfqpoint{3.318805in}{5.158944in}}%
\pgfpathcurveto{\pgfqpoint{3.326056in}{5.151694in}}{\pgfqpoint{3.335891in}{5.147620in}}{\pgfqpoint{3.346145in}{5.147620in}}%
\pgfpathlineto{\pgfqpoint{3.346145in}{5.147620in}}%
\pgfpathclose%
\pgfusepath{stroke}%
\end{pgfscope}%
\begin{pgfscope}%
\pgfpathrectangle{\pgfqpoint{0.977909in}{2.928000in}}{\pgfqpoint{5.664000in}{5.664000in}}%
\pgfusepath{clip}%
\pgfsetbuttcap%
\pgfsetroundjoin%
\pgfsetlinewidth{1.003750pt}%
\definecolor{currentstroke}{rgb}{0.321569,0.321569,0.321569}%
\pgfsetstrokecolor{currentstroke}%
\pgfsetstrokeopacity{0.900000}%
\pgfsetdash{}{0pt}%
\pgfpathmoveto{\pgfqpoint{3.886015in}{5.583507in}}%
\pgfpathcurveto{\pgfqpoint{3.896269in}{5.583507in}}{\pgfqpoint{3.906105in}{5.587581in}}{\pgfqpoint{3.913356in}{5.594832in}}%
\pgfpathcurveto{\pgfqpoint{3.920606in}{5.602082in}}{\pgfqpoint{3.924680in}{5.611918in}}{\pgfqpoint{3.924680in}{5.622172in}}%
\pgfpathcurveto{\pgfqpoint{3.924680in}{5.632426in}}{\pgfqpoint{3.920606in}{5.642261in}}{\pgfqpoint{3.913356in}{5.649512in}}%
\pgfpathcurveto{\pgfqpoint{3.906105in}{5.656763in}}{\pgfqpoint{3.896269in}{5.660837in}}{\pgfqpoint{3.886015in}{5.660837in}}%
\pgfpathcurveto{\pgfqpoint{3.875761in}{5.660837in}}{\pgfqpoint{3.865926in}{5.656763in}}{\pgfqpoint{3.858675in}{5.649512in}}%
\pgfpathcurveto{\pgfqpoint{3.851424in}{5.642261in}}{\pgfqpoint{3.847350in}{5.632426in}}{\pgfqpoint{3.847350in}{5.622172in}}%
\pgfpathcurveto{\pgfqpoint{3.847350in}{5.611918in}}{\pgfqpoint{3.851424in}{5.602082in}}{\pgfqpoint{3.858675in}{5.594832in}}%
\pgfpathcurveto{\pgfqpoint{3.865926in}{5.587581in}}{\pgfqpoint{3.875761in}{5.583507in}}{\pgfqpoint{3.886015in}{5.583507in}}%
\pgfpathlineto{\pgfqpoint{3.886015in}{5.583507in}}%
\pgfpathclose%
\pgfusepath{stroke}%
\end{pgfscope}%
\begin{pgfscope}%
\pgfpathrectangle{\pgfqpoint{0.977909in}{2.928000in}}{\pgfqpoint{5.664000in}{5.664000in}}%
\pgfusepath{clip}%
\pgfsetbuttcap%
\pgfsetroundjoin%
\pgfsetlinewidth{1.003750pt}%
\definecolor{currentstroke}{rgb}{0.321569,0.321569,0.321569}%
\pgfsetstrokecolor{currentstroke}%
\pgfsetstrokeopacity{0.900000}%
\pgfsetdash{}{0pt}%
\pgfpathmoveto{\pgfqpoint{5.018110in}{7.459529in}}%
\pgfpathcurveto{\pgfqpoint{5.028364in}{7.459529in}}{\pgfqpoint{5.038200in}{7.463603in}}{\pgfqpoint{5.045450in}{7.470854in}}%
\pgfpathcurveto{\pgfqpoint{5.052701in}{7.478104in}}{\pgfqpoint{5.056775in}{7.487940in}}{\pgfqpoint{5.056775in}{7.498194in}}%
\pgfpathcurveto{\pgfqpoint{5.056775in}{7.508448in}}{\pgfqpoint{5.052701in}{7.518284in}}{\pgfqpoint{5.045450in}{7.525534in}}%
\pgfpathcurveto{\pgfqpoint{5.038200in}{7.532785in}}{\pgfqpoint{5.028364in}{7.536859in}}{\pgfqpoint{5.018110in}{7.536859in}}%
\pgfpathcurveto{\pgfqpoint{5.007856in}{7.536859in}}{\pgfqpoint{4.998021in}{7.532785in}}{\pgfqpoint{4.990770in}{7.525534in}}%
\pgfpathcurveto{\pgfqpoint{4.983519in}{7.518284in}}{\pgfqpoint{4.979445in}{7.508448in}}{\pgfqpoint{4.979445in}{7.498194in}}%
\pgfpathcurveto{\pgfqpoint{4.979445in}{7.487940in}}{\pgfqpoint{4.983519in}{7.478104in}}{\pgfqpoint{4.990770in}{7.470854in}}%
\pgfpathcurveto{\pgfqpoint{4.998021in}{7.463603in}}{\pgfqpoint{5.007856in}{7.459529in}}{\pgfqpoint{5.018110in}{7.459529in}}%
\pgfpathlineto{\pgfqpoint{5.018110in}{7.459529in}}%
\pgfpathclose%
\pgfusepath{stroke}%
\end{pgfscope}%
\begin{pgfscope}%
\pgfpathrectangle{\pgfqpoint{0.977909in}{2.928000in}}{\pgfqpoint{5.664000in}{5.664000in}}%
\pgfusepath{clip}%
\pgfsetbuttcap%
\pgfsetroundjoin%
\pgfsetlinewidth{1.003750pt}%
\definecolor{currentstroke}{rgb}{0.321569,0.321569,0.321569}%
\pgfsetstrokecolor{currentstroke}%
\pgfsetstrokeopacity{0.900000}%
\pgfsetdash{}{0pt}%
\pgfpathmoveto{\pgfqpoint{3.172390in}{5.102918in}}%
\pgfpathcurveto{\pgfqpoint{3.182644in}{5.102918in}}{\pgfqpoint{3.192480in}{5.106992in}}{\pgfqpoint{3.199730in}{5.114243in}}%
\pgfpathcurveto{\pgfqpoint{3.206981in}{5.121494in}}{\pgfqpoint{3.211055in}{5.131329in}}{\pgfqpoint{3.211055in}{5.141583in}}%
\pgfpathcurveto{\pgfqpoint{3.211055in}{5.151837in}}{\pgfqpoint{3.206981in}{5.161673in}}{\pgfqpoint{3.199730in}{5.168924in}}%
\pgfpathcurveto{\pgfqpoint{3.192480in}{5.176174in}}{\pgfqpoint{3.182644in}{5.180248in}}{\pgfqpoint{3.172390in}{5.180248in}}%
\pgfpathcurveto{\pgfqpoint{3.162136in}{5.180248in}}{\pgfqpoint{3.152300in}{5.176174in}}{\pgfqpoint{3.145050in}{5.168924in}}%
\pgfpathcurveto{\pgfqpoint{3.137799in}{5.161673in}}{\pgfqpoint{3.133725in}{5.151837in}}{\pgfqpoint{3.133725in}{5.141583in}}%
\pgfpathcurveto{\pgfqpoint{3.133725in}{5.131329in}}{\pgfqpoint{3.137799in}{5.121494in}}{\pgfqpoint{3.145050in}{5.114243in}}%
\pgfpathcurveto{\pgfqpoint{3.152300in}{5.106992in}}{\pgfqpoint{3.162136in}{5.102918in}}{\pgfqpoint{3.172390in}{5.102918in}}%
\pgfpathlineto{\pgfqpoint{3.172390in}{5.102918in}}%
\pgfpathclose%
\pgfusepath{stroke}%
\end{pgfscope}%
\begin{pgfscope}%
\pgfpathrectangle{\pgfqpoint{0.977909in}{2.928000in}}{\pgfqpoint{5.664000in}{5.664000in}}%
\pgfusepath{clip}%
\pgfsetbuttcap%
\pgfsetroundjoin%
\pgfsetlinewidth{1.003750pt}%
\definecolor{currentstroke}{rgb}{0.321569,0.321569,0.321569}%
\pgfsetstrokecolor{currentstroke}%
\pgfsetstrokeopacity{0.900000}%
\pgfsetdash{}{0pt}%
\pgfpathmoveto{\pgfqpoint{3.340865in}{5.060541in}}%
\pgfpathcurveto{\pgfqpoint{3.351119in}{5.060541in}}{\pgfqpoint{3.360955in}{5.064615in}}{\pgfqpoint{3.368206in}{5.071866in}}%
\pgfpathcurveto{\pgfqpoint{3.375456in}{5.079116in}}{\pgfqpoint{3.379530in}{5.088952in}}{\pgfqpoint{3.379530in}{5.099206in}}%
\pgfpathcurveto{\pgfqpoint{3.379530in}{5.109460in}}{\pgfqpoint{3.375456in}{5.119295in}}{\pgfqpoint{3.368206in}{5.126546in}}%
\pgfpathcurveto{\pgfqpoint{3.360955in}{5.133797in}}{\pgfqpoint{3.351119in}{5.137871in}}{\pgfqpoint{3.340865in}{5.137871in}}%
\pgfpathcurveto{\pgfqpoint{3.330611in}{5.137871in}}{\pgfqpoint{3.320776in}{5.133797in}}{\pgfqpoint{3.313525in}{5.126546in}}%
\pgfpathcurveto{\pgfqpoint{3.306274in}{5.119295in}}{\pgfqpoint{3.302200in}{5.109460in}}{\pgfqpoint{3.302200in}{5.099206in}}%
\pgfpathcurveto{\pgfqpoint{3.302200in}{5.088952in}}{\pgfqpoint{3.306274in}{5.079116in}}{\pgfqpoint{3.313525in}{5.071866in}}%
\pgfpathcurveto{\pgfqpoint{3.320776in}{5.064615in}}{\pgfqpoint{3.330611in}{5.060541in}}{\pgfqpoint{3.340865in}{5.060541in}}%
\pgfpathlineto{\pgfqpoint{3.340865in}{5.060541in}}%
\pgfpathclose%
\pgfusepath{stroke}%
\end{pgfscope}%
\begin{pgfscope}%
\pgfpathrectangle{\pgfqpoint{0.977909in}{2.928000in}}{\pgfqpoint{5.664000in}{5.664000in}}%
\pgfusepath{clip}%
\pgfsetbuttcap%
\pgfsetroundjoin%
\pgfsetlinewidth{1.003750pt}%
\definecolor{currentstroke}{rgb}{0.321569,0.321569,0.321569}%
\pgfsetstrokecolor{currentstroke}%
\pgfsetstrokeopacity{0.900000}%
\pgfsetdash{}{0pt}%
\pgfpathmoveto{\pgfqpoint{4.017407in}{5.570974in}}%
\pgfpathcurveto{\pgfqpoint{4.027662in}{5.570974in}}{\pgfqpoint{4.037497in}{5.575048in}}{\pgfqpoint{4.044748in}{5.582298in}}%
\pgfpathcurveto{\pgfqpoint{4.051998in}{5.589549in}}{\pgfqpoint{4.056072in}{5.599384in}}{\pgfqpoint{4.056072in}{5.609639in}}%
\pgfpathcurveto{\pgfqpoint{4.056072in}{5.619893in}}{\pgfqpoint{4.051998in}{5.629728in}}{\pgfqpoint{4.044748in}{5.636979in}}%
\pgfpathcurveto{\pgfqpoint{4.037497in}{5.644230in}}{\pgfqpoint{4.027662in}{5.648304in}}{\pgfqpoint{4.017407in}{5.648304in}}%
\pgfpathcurveto{\pgfqpoint{4.007153in}{5.648304in}}{\pgfqpoint{3.997318in}{5.644230in}}{\pgfqpoint{3.990067in}{5.636979in}}%
\pgfpathcurveto{\pgfqpoint{3.982816in}{5.629728in}}{\pgfqpoint{3.978742in}{5.619893in}}{\pgfqpoint{3.978742in}{5.609639in}}%
\pgfpathcurveto{\pgfqpoint{3.978742in}{5.599384in}}{\pgfqpoint{3.982816in}{5.589549in}}{\pgfqpoint{3.990067in}{5.582298in}}%
\pgfpathcurveto{\pgfqpoint{3.997318in}{5.575048in}}{\pgfqpoint{4.007153in}{5.570974in}}{\pgfqpoint{4.017407in}{5.570974in}}%
\pgfpathlineto{\pgfqpoint{4.017407in}{5.570974in}}%
\pgfpathclose%
\pgfusepath{stroke}%
\end{pgfscope}%
\begin{pgfscope}%
\pgfpathrectangle{\pgfqpoint{0.977909in}{2.928000in}}{\pgfqpoint{5.664000in}{5.664000in}}%
\pgfusepath{clip}%
\pgfsetbuttcap%
\pgfsetroundjoin%
\pgfsetlinewidth{1.003750pt}%
\definecolor{currentstroke}{rgb}{0.321569,0.321569,0.321569}%
\pgfsetstrokecolor{currentstroke}%
\pgfsetstrokeopacity{0.900000}%
\pgfsetdash{}{0pt}%
\pgfpathmoveto{\pgfqpoint{2.820454in}{4.999695in}}%
\pgfpathcurveto{\pgfqpoint{2.830708in}{4.999695in}}{\pgfqpoint{2.840543in}{5.003769in}}{\pgfqpoint{2.847794in}{5.011020in}}%
\pgfpathcurveto{\pgfqpoint{2.855045in}{5.018271in}}{\pgfqpoint{2.859119in}{5.028106in}}{\pgfqpoint{2.859119in}{5.038360in}}%
\pgfpathcurveto{\pgfqpoint{2.859119in}{5.048614in}}{\pgfqpoint{2.855045in}{5.058450in}}{\pgfqpoint{2.847794in}{5.065700in}}%
\pgfpathcurveto{\pgfqpoint{2.840543in}{5.072951in}}{\pgfqpoint{2.830708in}{5.077025in}}{\pgfqpoint{2.820454in}{5.077025in}}%
\pgfpathcurveto{\pgfqpoint{2.810200in}{5.077025in}}{\pgfqpoint{2.800364in}{5.072951in}}{\pgfqpoint{2.793113in}{5.065700in}}%
\pgfpathcurveto{\pgfqpoint{2.785863in}{5.058450in}}{\pgfqpoint{2.781789in}{5.048614in}}{\pgfqpoint{2.781789in}{5.038360in}}%
\pgfpathcurveto{\pgfqpoint{2.781789in}{5.028106in}}{\pgfqpoint{2.785863in}{5.018271in}}{\pgfqpoint{2.793113in}{5.011020in}}%
\pgfpathcurveto{\pgfqpoint{2.800364in}{5.003769in}}{\pgfqpoint{2.810200in}{4.999695in}}{\pgfqpoint{2.820454in}{4.999695in}}%
\pgfpathlineto{\pgfqpoint{2.820454in}{4.999695in}}%
\pgfpathclose%
\pgfusepath{stroke}%
\end{pgfscope}%
\begin{pgfscope}%
\pgfpathrectangle{\pgfqpoint{0.977909in}{2.928000in}}{\pgfqpoint{5.664000in}{5.664000in}}%
\pgfusepath{clip}%
\pgfsetbuttcap%
\pgfsetroundjoin%
\pgfsetlinewidth{1.003750pt}%
\definecolor{currentstroke}{rgb}{0.321569,0.321569,0.321569}%
\pgfsetstrokecolor{currentstroke}%
\pgfsetstrokeopacity{0.900000}%
\pgfsetdash{}{0pt}%
\pgfpathmoveto{\pgfqpoint{3.219125in}{5.318706in}}%
\pgfpathcurveto{\pgfqpoint{3.229379in}{5.318706in}}{\pgfqpoint{3.239215in}{5.322780in}}{\pgfqpoint{3.246465in}{5.330030in}}%
\pgfpathcurveto{\pgfqpoint{3.253716in}{5.337281in}}{\pgfqpoint{3.257790in}{5.347117in}}{\pgfqpoint{3.257790in}{5.357371in}}%
\pgfpathcurveto{\pgfqpoint{3.257790in}{5.367625in}}{\pgfqpoint{3.253716in}{5.377460in}}{\pgfqpoint{3.246465in}{5.384711in}}%
\pgfpathcurveto{\pgfqpoint{3.239215in}{5.391962in}}{\pgfqpoint{3.229379in}{5.396036in}}{\pgfqpoint{3.219125in}{5.396036in}}%
\pgfpathcurveto{\pgfqpoint{3.208871in}{5.396036in}}{\pgfqpoint{3.199036in}{5.391962in}}{\pgfqpoint{3.191785in}{5.384711in}}%
\pgfpathcurveto{\pgfqpoint{3.184534in}{5.377460in}}{\pgfqpoint{3.180460in}{5.367625in}}{\pgfqpoint{3.180460in}{5.357371in}}%
\pgfpathcurveto{\pgfqpoint{3.180460in}{5.347117in}}{\pgfqpoint{3.184534in}{5.337281in}}{\pgfqpoint{3.191785in}{5.330030in}}%
\pgfpathcurveto{\pgfqpoint{3.199036in}{5.322780in}}{\pgfqpoint{3.208871in}{5.318706in}}{\pgfqpoint{3.219125in}{5.318706in}}%
\pgfpathlineto{\pgfqpoint{3.219125in}{5.318706in}}%
\pgfpathclose%
\pgfusepath{stroke}%
\end{pgfscope}%
\begin{pgfscope}%
\pgfpathrectangle{\pgfqpoint{0.977909in}{2.928000in}}{\pgfqpoint{5.664000in}{5.664000in}}%
\pgfusepath{clip}%
\pgfsetbuttcap%
\pgfsetroundjoin%
\pgfsetlinewidth{1.003750pt}%
\definecolor{currentstroke}{rgb}{0.321569,0.321569,0.321569}%
\pgfsetstrokecolor{currentstroke}%
\pgfsetstrokeopacity{0.900000}%
\pgfsetdash{}{0pt}%
\pgfpathmoveto{\pgfqpoint{2.568245in}{4.853763in}}%
\pgfpathcurveto{\pgfqpoint{2.578500in}{4.853763in}}{\pgfqpoint{2.588335in}{4.857837in}}{\pgfqpoint{2.595586in}{4.865088in}}%
\pgfpathcurveto{\pgfqpoint{2.602837in}{4.872339in}}{\pgfqpoint{2.606911in}{4.882174in}}{\pgfqpoint{2.606911in}{4.892428in}}%
\pgfpathcurveto{\pgfqpoint{2.606911in}{4.902682in}}{\pgfqpoint{2.602837in}{4.912518in}}{\pgfqpoint{2.595586in}{4.919769in}}%
\pgfpathcurveto{\pgfqpoint{2.588335in}{4.927019in}}{\pgfqpoint{2.578500in}{4.931093in}}{\pgfqpoint{2.568245in}{4.931093in}}%
\pgfpathcurveto{\pgfqpoint{2.557991in}{4.931093in}}{\pgfqpoint{2.548156in}{4.927019in}}{\pgfqpoint{2.540905in}{4.919769in}}%
\pgfpathcurveto{\pgfqpoint{2.533654in}{4.912518in}}{\pgfqpoint{2.529580in}{4.902682in}}{\pgfqpoint{2.529580in}{4.892428in}}%
\pgfpathcurveto{\pgfqpoint{2.529580in}{4.882174in}}{\pgfqpoint{2.533654in}{4.872339in}}{\pgfqpoint{2.540905in}{4.865088in}}%
\pgfpathcurveto{\pgfqpoint{2.548156in}{4.857837in}}{\pgfqpoint{2.557991in}{4.853763in}}{\pgfqpoint{2.568245in}{4.853763in}}%
\pgfpathlineto{\pgfqpoint{2.568245in}{4.853763in}}%
\pgfpathclose%
\pgfusepath{stroke}%
\end{pgfscope}%
\begin{pgfscope}%
\pgfpathrectangle{\pgfqpoint{0.977909in}{2.928000in}}{\pgfqpoint{5.664000in}{5.664000in}}%
\pgfusepath{clip}%
\pgfsetbuttcap%
\pgfsetroundjoin%
\pgfsetlinewidth{1.003750pt}%
\definecolor{currentstroke}{rgb}{0.321569,0.321569,0.321569}%
\pgfsetstrokecolor{currentstroke}%
\pgfsetstrokeopacity{0.900000}%
\pgfsetdash{}{0pt}%
\pgfpathmoveto{\pgfqpoint{2.961030in}{5.093804in}}%
\pgfpathcurveto{\pgfqpoint{2.971284in}{5.093804in}}{\pgfqpoint{2.981120in}{5.097878in}}{\pgfqpoint{2.988371in}{5.105129in}}%
\pgfpathcurveto{\pgfqpoint{2.995621in}{5.112380in}}{\pgfqpoint{2.999695in}{5.122215in}}{\pgfqpoint{2.999695in}{5.132469in}}%
\pgfpathcurveto{\pgfqpoint{2.999695in}{5.142723in}}{\pgfqpoint{2.995621in}{5.152559in}}{\pgfqpoint{2.988371in}{5.159810in}}%
\pgfpathcurveto{\pgfqpoint{2.981120in}{5.167060in}}{\pgfqpoint{2.971284in}{5.171134in}}{\pgfqpoint{2.961030in}{5.171134in}}%
\pgfpathcurveto{\pgfqpoint{2.950776in}{5.171134in}}{\pgfqpoint{2.940941in}{5.167060in}}{\pgfqpoint{2.933690in}{5.159810in}}%
\pgfpathcurveto{\pgfqpoint{2.926439in}{5.152559in}}{\pgfqpoint{2.922365in}{5.142723in}}{\pgfqpoint{2.922365in}{5.132469in}}%
\pgfpathcurveto{\pgfqpoint{2.922365in}{5.122215in}}{\pgfqpoint{2.926439in}{5.112380in}}{\pgfqpoint{2.933690in}{5.105129in}}%
\pgfpathcurveto{\pgfqpoint{2.940941in}{5.097878in}}{\pgfqpoint{2.950776in}{5.093804in}}{\pgfqpoint{2.961030in}{5.093804in}}%
\pgfpathlineto{\pgfqpoint{2.961030in}{5.093804in}}%
\pgfpathclose%
\pgfusepath{stroke}%
\end{pgfscope}%
\begin{pgfscope}%
\pgfpathrectangle{\pgfqpoint{0.977909in}{2.928000in}}{\pgfqpoint{5.664000in}{5.664000in}}%
\pgfusepath{clip}%
\pgfsetbuttcap%
\pgfsetroundjoin%
\pgfsetlinewidth{1.003750pt}%
\definecolor{currentstroke}{rgb}{0.321569,0.321569,0.321569}%
\pgfsetstrokecolor{currentstroke}%
\pgfsetstrokeopacity{0.900000}%
\pgfsetdash{}{0pt}%
\pgfpathmoveto{\pgfqpoint{3.666799in}{5.516748in}}%
\pgfpathcurveto{\pgfqpoint{3.677053in}{5.516748in}}{\pgfqpoint{3.686888in}{5.520822in}}{\pgfqpoint{3.694139in}{5.528073in}}%
\pgfpathcurveto{\pgfqpoint{3.701390in}{5.535323in}}{\pgfqpoint{3.705464in}{5.545159in}}{\pgfqpoint{3.705464in}{5.555413in}}%
\pgfpathcurveto{\pgfqpoint{3.705464in}{5.565667in}}{\pgfqpoint{3.701390in}{5.575503in}}{\pgfqpoint{3.694139in}{5.582753in}}%
\pgfpathcurveto{\pgfqpoint{3.686888in}{5.590004in}}{\pgfqpoint{3.677053in}{5.594078in}}{\pgfqpoint{3.666799in}{5.594078in}}%
\pgfpathcurveto{\pgfqpoint{3.656544in}{5.594078in}}{\pgfqpoint{3.646709in}{5.590004in}}{\pgfqpoint{3.639458in}{5.582753in}}%
\pgfpathcurveto{\pgfqpoint{3.632208in}{5.575503in}}{\pgfqpoint{3.628134in}{5.565667in}}{\pgfqpoint{3.628134in}{5.555413in}}%
\pgfpathcurveto{\pgfqpoint{3.628134in}{5.545159in}}{\pgfqpoint{3.632208in}{5.535323in}}{\pgfqpoint{3.639458in}{5.528073in}}%
\pgfpathcurveto{\pgfqpoint{3.646709in}{5.520822in}}{\pgfqpoint{3.656544in}{5.516748in}}{\pgfqpoint{3.666799in}{5.516748in}}%
\pgfpathlineto{\pgfqpoint{3.666799in}{5.516748in}}%
\pgfpathclose%
\pgfusepath{stroke}%
\end{pgfscope}%
\begin{pgfscope}%
\pgfpathrectangle{\pgfqpoint{0.977909in}{2.928000in}}{\pgfqpoint{5.664000in}{5.664000in}}%
\pgfusepath{clip}%
\pgfsetbuttcap%
\pgfsetroundjoin%
\pgfsetlinewidth{1.003750pt}%
\definecolor{currentstroke}{rgb}{0.321569,0.321569,0.321569}%
\pgfsetstrokecolor{currentstroke}%
\pgfsetstrokeopacity{0.900000}%
\pgfsetdash{}{0pt}%
\pgfpathmoveto{\pgfqpoint{2.791540in}{4.440985in}}%
\pgfpathcurveto{\pgfqpoint{2.801794in}{4.440985in}}{\pgfqpoint{2.811630in}{4.445059in}}{\pgfqpoint{2.818880in}{4.452309in}}%
\pgfpathcurveto{\pgfqpoint{2.826131in}{4.459560in}}{\pgfqpoint{2.830205in}{4.469396in}}{\pgfqpoint{2.830205in}{4.479650in}}%
\pgfpathcurveto{\pgfqpoint{2.830205in}{4.489904in}}{\pgfqpoint{2.826131in}{4.499739in}}{\pgfqpoint{2.818880in}{4.506990in}}%
\pgfpathcurveto{\pgfqpoint{2.811630in}{4.514241in}}{\pgfqpoint{2.801794in}{4.518315in}}{\pgfqpoint{2.791540in}{4.518315in}}%
\pgfpathcurveto{\pgfqpoint{2.781286in}{4.518315in}}{\pgfqpoint{2.771451in}{4.514241in}}{\pgfqpoint{2.764200in}{4.506990in}}%
\pgfpathcurveto{\pgfqpoint{2.756949in}{4.499739in}}{\pgfqpoint{2.752875in}{4.489904in}}{\pgfqpoint{2.752875in}{4.479650in}}%
\pgfpathcurveto{\pgfqpoint{2.752875in}{4.469396in}}{\pgfqpoint{2.756949in}{4.459560in}}{\pgfqpoint{2.764200in}{4.452309in}}%
\pgfpathcurveto{\pgfqpoint{2.771451in}{4.445059in}}{\pgfqpoint{2.781286in}{4.440985in}}{\pgfqpoint{2.791540in}{4.440985in}}%
\pgfpathlineto{\pgfqpoint{2.791540in}{4.440985in}}%
\pgfpathclose%
\pgfusepath{stroke}%
\end{pgfscope}%
\begin{pgfscope}%
\pgfpathrectangle{\pgfqpoint{0.977909in}{2.928000in}}{\pgfqpoint{5.664000in}{5.664000in}}%
\pgfusepath{clip}%
\pgfsetbuttcap%
\pgfsetroundjoin%
\pgfsetlinewidth{1.003750pt}%
\definecolor{currentstroke}{rgb}{0.321569,0.321569,0.321569}%
\pgfsetstrokecolor{currentstroke}%
\pgfsetstrokeopacity{0.900000}%
\pgfsetdash{}{0pt}%
\pgfpathmoveto{\pgfqpoint{3.004852in}{4.625181in}}%
\pgfpathcurveto{\pgfqpoint{3.015106in}{4.625181in}}{\pgfqpoint{3.024941in}{4.629255in}}{\pgfqpoint{3.032192in}{4.636505in}}%
\pgfpathcurveto{\pgfqpoint{3.039443in}{4.643756in}}{\pgfqpoint{3.043517in}{4.653592in}}{\pgfqpoint{3.043517in}{4.663846in}}%
\pgfpathcurveto{\pgfqpoint{3.043517in}{4.674100in}}{\pgfqpoint{3.039443in}{4.683935in}}{\pgfqpoint{3.032192in}{4.691186in}}%
\pgfpathcurveto{\pgfqpoint{3.024941in}{4.698437in}}{\pgfqpoint{3.015106in}{4.702511in}}{\pgfqpoint{3.004852in}{4.702511in}}%
\pgfpathcurveto{\pgfqpoint{2.994598in}{4.702511in}}{\pgfqpoint{2.984762in}{4.698437in}}{\pgfqpoint{2.977511in}{4.691186in}}%
\pgfpathcurveto{\pgfqpoint{2.970261in}{4.683935in}}{\pgfqpoint{2.966187in}{4.674100in}}{\pgfqpoint{2.966187in}{4.663846in}}%
\pgfpathcurveto{\pgfqpoint{2.966187in}{4.653592in}}{\pgfqpoint{2.970261in}{4.643756in}}{\pgfqpoint{2.977511in}{4.636505in}}%
\pgfpathcurveto{\pgfqpoint{2.984762in}{4.629255in}}{\pgfqpoint{2.994598in}{4.625181in}}{\pgfqpoint{3.004852in}{4.625181in}}%
\pgfpathlineto{\pgfqpoint{3.004852in}{4.625181in}}%
\pgfpathclose%
\pgfusepath{stroke}%
\end{pgfscope}%
\begin{pgfscope}%
\pgfpathrectangle{\pgfqpoint{0.977909in}{2.928000in}}{\pgfqpoint{5.664000in}{5.664000in}}%
\pgfusepath{clip}%
\pgfsetbuttcap%
\pgfsetroundjoin%
\pgfsetlinewidth{1.003750pt}%
\definecolor{currentstroke}{rgb}{0.321569,0.321569,0.321569}%
\pgfsetstrokecolor{currentstroke}%
\pgfsetstrokeopacity{0.900000}%
\pgfsetdash{}{0pt}%
\pgfpathmoveto{\pgfqpoint{1.158935in}{3.663157in}}%
\pgfpathcurveto{\pgfqpoint{1.169189in}{3.663157in}}{\pgfqpoint{1.179025in}{3.667231in}}{\pgfqpoint{1.186276in}{3.674482in}}%
\pgfpathcurveto{\pgfqpoint{1.193526in}{3.681733in}}{\pgfqpoint{1.197600in}{3.691568in}}{\pgfqpoint{1.197600in}{3.701822in}}%
\pgfpathcurveto{\pgfqpoint{1.197600in}{3.712077in}}{\pgfqpoint{1.193526in}{3.721912in}}{\pgfqpoint{1.186276in}{3.729163in}}%
\pgfpathcurveto{\pgfqpoint{1.179025in}{3.736414in}}{\pgfqpoint{1.169189in}{3.740488in}}{\pgfqpoint{1.158935in}{3.740488in}}%
\pgfpathcurveto{\pgfqpoint{1.148681in}{3.740488in}}{\pgfqpoint{1.138846in}{3.736414in}}{\pgfqpoint{1.131595in}{3.729163in}}%
\pgfpathcurveto{\pgfqpoint{1.124344in}{3.721912in}}{\pgfqpoint{1.120270in}{3.712077in}}{\pgfqpoint{1.120270in}{3.701822in}}%
\pgfpathcurveto{\pgfqpoint{1.120270in}{3.691568in}}{\pgfqpoint{1.124344in}{3.681733in}}{\pgfqpoint{1.131595in}{3.674482in}}%
\pgfpathcurveto{\pgfqpoint{1.138846in}{3.667231in}}{\pgfqpoint{1.148681in}{3.663157in}}{\pgfqpoint{1.158935in}{3.663157in}}%
\pgfpathlineto{\pgfqpoint{1.158935in}{3.663157in}}%
\pgfpathclose%
\pgfusepath{stroke}%
\end{pgfscope}%
\begin{pgfscope}%
\pgfpathrectangle{\pgfqpoint{0.977909in}{2.928000in}}{\pgfqpoint{5.664000in}{5.664000in}}%
\pgfusepath{clip}%
\pgfsetbuttcap%
\pgfsetroundjoin%
\pgfsetlinewidth{1.003750pt}%
\definecolor{currentstroke}{rgb}{0.321569,0.321569,0.321569}%
\pgfsetstrokecolor{currentstroke}%
\pgfsetstrokeopacity{0.900000}%
\pgfsetdash{}{0pt}%
\pgfpathmoveto{\pgfqpoint{3.999695in}{5.490495in}}%
\pgfpathcurveto{\pgfqpoint{4.009949in}{5.490495in}}{\pgfqpoint{4.019785in}{5.494569in}}{\pgfqpoint{4.027035in}{5.501819in}}%
\pgfpathcurveto{\pgfqpoint{4.034286in}{5.509070in}}{\pgfqpoint{4.038360in}{5.518906in}}{\pgfqpoint{4.038360in}{5.529160in}}%
\pgfpathcurveto{\pgfqpoint{4.038360in}{5.539414in}}{\pgfqpoint{4.034286in}{5.549249in}}{\pgfqpoint{4.027035in}{5.556500in}}%
\pgfpathcurveto{\pgfqpoint{4.019785in}{5.563751in}}{\pgfqpoint{4.009949in}{5.567825in}}{\pgfqpoint{3.999695in}{5.567825in}}%
\pgfpathcurveto{\pgfqpoint{3.989441in}{5.567825in}}{\pgfqpoint{3.979605in}{5.563751in}}{\pgfqpoint{3.972355in}{5.556500in}}%
\pgfpathcurveto{\pgfqpoint{3.965104in}{5.549249in}}{\pgfqpoint{3.961030in}{5.539414in}}{\pgfqpoint{3.961030in}{5.529160in}}%
\pgfpathcurveto{\pgfqpoint{3.961030in}{5.518906in}}{\pgfqpoint{3.965104in}{5.509070in}}{\pgfqpoint{3.972355in}{5.501819in}}%
\pgfpathcurveto{\pgfqpoint{3.979605in}{5.494569in}}{\pgfqpoint{3.989441in}{5.490495in}}{\pgfqpoint{3.999695in}{5.490495in}}%
\pgfpathlineto{\pgfqpoint{3.999695in}{5.490495in}}%
\pgfpathclose%
\pgfusepath{stroke}%
\end{pgfscope}%
\begin{pgfscope}%
\pgfpathrectangle{\pgfqpoint{0.977909in}{2.928000in}}{\pgfqpoint{5.664000in}{5.664000in}}%
\pgfusepath{clip}%
\pgfsetbuttcap%
\pgfsetroundjoin%
\pgfsetlinewidth{1.003750pt}%
\definecolor{currentstroke}{rgb}{0.321569,0.321569,0.321569}%
\pgfsetstrokecolor{currentstroke}%
\pgfsetstrokeopacity{0.900000}%
\pgfsetdash{}{0pt}%
\pgfpathmoveto{\pgfqpoint{2.385500in}{4.492592in}}%
\pgfpathcurveto{\pgfqpoint{2.395754in}{4.492592in}}{\pgfqpoint{2.405590in}{4.496666in}}{\pgfqpoint{2.412841in}{4.503916in}}%
\pgfpathcurveto{\pgfqpoint{2.420091in}{4.511167in}}{\pgfqpoint{2.424165in}{4.521002in}}{\pgfqpoint{2.424165in}{4.531257in}}%
\pgfpathcurveto{\pgfqpoint{2.424165in}{4.541511in}}{\pgfqpoint{2.420091in}{4.551346in}}{\pgfqpoint{2.412841in}{4.558597in}}%
\pgfpathcurveto{\pgfqpoint{2.405590in}{4.565848in}}{\pgfqpoint{2.395754in}{4.569922in}}{\pgfqpoint{2.385500in}{4.569922in}}%
\pgfpathcurveto{\pgfqpoint{2.375246in}{4.569922in}}{\pgfqpoint{2.365411in}{4.565848in}}{\pgfqpoint{2.358160in}{4.558597in}}%
\pgfpathcurveto{\pgfqpoint{2.350909in}{4.551346in}}{\pgfqpoint{2.346835in}{4.541511in}}{\pgfqpoint{2.346835in}{4.531257in}}%
\pgfpathcurveto{\pgfqpoint{2.346835in}{4.521002in}}{\pgfqpoint{2.350909in}{4.511167in}}{\pgfqpoint{2.358160in}{4.503916in}}%
\pgfpathcurveto{\pgfqpoint{2.365411in}{4.496666in}}{\pgfqpoint{2.375246in}{4.492592in}}{\pgfqpoint{2.385500in}{4.492592in}}%
\pgfpathlineto{\pgfqpoint{2.385500in}{4.492592in}}%
\pgfpathclose%
\pgfusepath{stroke}%
\end{pgfscope}%
\begin{pgfscope}%
\pgfpathrectangle{\pgfqpoint{0.977909in}{2.928000in}}{\pgfqpoint{5.664000in}{5.664000in}}%
\pgfusepath{clip}%
\pgfsetbuttcap%
\pgfsetroundjoin%
\pgfsetlinewidth{1.003750pt}%
\definecolor{currentstroke}{rgb}{0.321569,0.321569,0.321569}%
\pgfsetstrokecolor{currentstroke}%
\pgfsetstrokeopacity{0.900000}%
\pgfsetdash{}{0pt}%
\pgfpathmoveto{\pgfqpoint{2.690812in}{4.986640in}}%
\pgfpathcurveto{\pgfqpoint{2.701066in}{4.986640in}}{\pgfqpoint{2.710901in}{4.990714in}}{\pgfqpoint{2.718152in}{4.997965in}}%
\pgfpathcurveto{\pgfqpoint{2.725403in}{5.005215in}}{\pgfqpoint{2.729477in}{5.015051in}}{\pgfqpoint{2.729477in}{5.025305in}}%
\pgfpathcurveto{\pgfqpoint{2.729477in}{5.035559in}}{\pgfqpoint{2.725403in}{5.045395in}}{\pgfqpoint{2.718152in}{5.052645in}}%
\pgfpathcurveto{\pgfqpoint{2.710901in}{5.059896in}}{\pgfqpoint{2.701066in}{5.063970in}}{\pgfqpoint{2.690812in}{5.063970in}}%
\pgfpathcurveto{\pgfqpoint{2.680558in}{5.063970in}}{\pgfqpoint{2.670722in}{5.059896in}}{\pgfqpoint{2.663471in}{5.052645in}}%
\pgfpathcurveto{\pgfqpoint{2.656221in}{5.045395in}}{\pgfqpoint{2.652147in}{5.035559in}}{\pgfqpoint{2.652147in}{5.025305in}}%
\pgfpathcurveto{\pgfqpoint{2.652147in}{5.015051in}}{\pgfqpoint{2.656221in}{5.005215in}}{\pgfqpoint{2.663471in}{4.997965in}}%
\pgfpathcurveto{\pgfqpoint{2.670722in}{4.990714in}}{\pgfqpoint{2.680558in}{4.986640in}}{\pgfqpoint{2.690812in}{4.986640in}}%
\pgfpathlineto{\pgfqpoint{2.690812in}{4.986640in}}%
\pgfpathclose%
\pgfusepath{stroke}%
\end{pgfscope}%
\begin{pgfscope}%
\pgfpathrectangle{\pgfqpoint{0.977909in}{2.928000in}}{\pgfqpoint{5.664000in}{5.664000in}}%
\pgfusepath{clip}%
\pgfsetbuttcap%
\pgfsetroundjoin%
\pgfsetlinewidth{1.003750pt}%
\definecolor{currentstroke}{rgb}{0.321569,0.321569,0.321569}%
\pgfsetstrokecolor{currentstroke}%
\pgfsetstrokeopacity{0.900000}%
\pgfsetdash{}{0pt}%
\pgfpathmoveto{\pgfqpoint{2.387127in}{4.698630in}}%
\pgfpathcurveto{\pgfqpoint{2.397381in}{4.698630in}}{\pgfqpoint{2.407217in}{4.702704in}}{\pgfqpoint{2.414467in}{4.709954in}}%
\pgfpathcurveto{\pgfqpoint{2.421718in}{4.717205in}}{\pgfqpoint{2.425792in}{4.727040in}}{\pgfqpoint{2.425792in}{4.737295in}}%
\pgfpathcurveto{\pgfqpoint{2.425792in}{4.747549in}}{\pgfqpoint{2.421718in}{4.757384in}}{\pgfqpoint{2.414467in}{4.764635in}}%
\pgfpathcurveto{\pgfqpoint{2.407217in}{4.771886in}}{\pgfqpoint{2.397381in}{4.775960in}}{\pgfqpoint{2.387127in}{4.775960in}}%
\pgfpathcurveto{\pgfqpoint{2.376873in}{4.775960in}}{\pgfqpoint{2.367037in}{4.771886in}}{\pgfqpoint{2.359787in}{4.764635in}}%
\pgfpathcurveto{\pgfqpoint{2.352536in}{4.757384in}}{\pgfqpoint{2.348462in}{4.747549in}}{\pgfqpoint{2.348462in}{4.737295in}}%
\pgfpathcurveto{\pgfqpoint{2.348462in}{4.727040in}}{\pgfqpoint{2.352536in}{4.717205in}}{\pgfqpoint{2.359787in}{4.709954in}}%
\pgfpathcurveto{\pgfqpoint{2.367037in}{4.702704in}}{\pgfqpoint{2.376873in}{4.698630in}}{\pgfqpoint{2.387127in}{4.698630in}}%
\pgfpathlineto{\pgfqpoint{2.387127in}{4.698630in}}%
\pgfpathclose%
\pgfusepath{stroke}%
\end{pgfscope}%
\begin{pgfscope}%
\pgfpathrectangle{\pgfqpoint{0.977909in}{2.928000in}}{\pgfqpoint{5.664000in}{5.664000in}}%
\pgfusepath{clip}%
\pgfsetbuttcap%
\pgfsetroundjoin%
\pgfsetlinewidth{1.003750pt}%
\definecolor{currentstroke}{rgb}{0.321569,0.321569,0.321569}%
\pgfsetstrokecolor{currentstroke}%
\pgfsetstrokeopacity{0.900000}%
\pgfsetdash{}{0pt}%
\pgfpathmoveto{\pgfqpoint{3.190054in}{5.290002in}}%
\pgfpathcurveto{\pgfqpoint{3.200308in}{5.290002in}}{\pgfqpoint{3.210143in}{5.294076in}}{\pgfqpoint{3.217394in}{5.301327in}}%
\pgfpathcurveto{\pgfqpoint{3.224645in}{5.308578in}}{\pgfqpoint{3.228719in}{5.318413in}}{\pgfqpoint{3.228719in}{5.328667in}}%
\pgfpathcurveto{\pgfqpoint{3.228719in}{5.338922in}}{\pgfqpoint{3.224645in}{5.348757in}}{\pgfqpoint{3.217394in}{5.356008in}}%
\pgfpathcurveto{\pgfqpoint{3.210143in}{5.363258in}}{\pgfqpoint{3.200308in}{5.367332in}}{\pgfqpoint{3.190054in}{5.367332in}}%
\pgfpathcurveto{\pgfqpoint{3.179800in}{5.367332in}}{\pgfqpoint{3.169964in}{5.363258in}}{\pgfqpoint{3.162714in}{5.356008in}}%
\pgfpathcurveto{\pgfqpoint{3.155463in}{5.348757in}}{\pgfqpoint{3.151389in}{5.338922in}}{\pgfqpoint{3.151389in}{5.328667in}}%
\pgfpathcurveto{\pgfqpoint{3.151389in}{5.318413in}}{\pgfqpoint{3.155463in}{5.308578in}}{\pgfqpoint{3.162714in}{5.301327in}}%
\pgfpathcurveto{\pgfqpoint{3.169964in}{5.294076in}}{\pgfqpoint{3.179800in}{5.290002in}}{\pgfqpoint{3.190054in}{5.290002in}}%
\pgfpathlineto{\pgfqpoint{3.190054in}{5.290002in}}%
\pgfpathclose%
\pgfusepath{stroke}%
\end{pgfscope}%
\begin{pgfscope}%
\pgfpathrectangle{\pgfqpoint{0.977909in}{2.928000in}}{\pgfqpoint{5.664000in}{5.664000in}}%
\pgfusepath{clip}%
\pgfsetbuttcap%
\pgfsetroundjoin%
\pgfsetlinewidth{1.003750pt}%
\definecolor{currentstroke}{rgb}{0.321569,0.321569,0.321569}%
\pgfsetstrokecolor{currentstroke}%
\pgfsetstrokeopacity{0.900000}%
\pgfsetdash{}{0pt}%
\pgfpathmoveto{\pgfqpoint{4.955025in}{6.747333in}}%
\pgfpathcurveto{\pgfqpoint{4.965279in}{6.747333in}}{\pgfqpoint{4.975114in}{6.751407in}}{\pgfqpoint{4.982365in}{6.758658in}}%
\pgfpathcurveto{\pgfqpoint{4.989616in}{6.765909in}}{\pgfqpoint{4.993690in}{6.775744in}}{\pgfqpoint{4.993690in}{6.785998in}}%
\pgfpathcurveto{\pgfqpoint{4.993690in}{6.796252in}}{\pgfqpoint{4.989616in}{6.806088in}}{\pgfqpoint{4.982365in}{6.813339in}}%
\pgfpathcurveto{\pgfqpoint{4.975114in}{6.820589in}}{\pgfqpoint{4.965279in}{6.824663in}}{\pgfqpoint{4.955025in}{6.824663in}}%
\pgfpathcurveto{\pgfqpoint{4.944771in}{6.824663in}}{\pgfqpoint{4.934935in}{6.820589in}}{\pgfqpoint{4.927685in}{6.813339in}}%
\pgfpathcurveto{\pgfqpoint{4.920434in}{6.806088in}}{\pgfqpoint{4.916360in}{6.796252in}}{\pgfqpoint{4.916360in}{6.785998in}}%
\pgfpathcurveto{\pgfqpoint{4.916360in}{6.775744in}}{\pgfqpoint{4.920434in}{6.765909in}}{\pgfqpoint{4.927685in}{6.758658in}}%
\pgfpathcurveto{\pgfqpoint{4.934935in}{6.751407in}}{\pgfqpoint{4.944771in}{6.747333in}}{\pgfqpoint{4.955025in}{6.747333in}}%
\pgfpathlineto{\pgfqpoint{4.955025in}{6.747333in}}%
\pgfpathclose%
\pgfusepath{stroke}%
\end{pgfscope}%
\begin{pgfscope}%
\pgfpathrectangle{\pgfqpoint{0.977909in}{2.928000in}}{\pgfqpoint{5.664000in}{5.664000in}}%
\pgfusepath{clip}%
\pgfsetbuttcap%
\pgfsetroundjoin%
\pgfsetlinewidth{1.003750pt}%
\definecolor{currentstroke}{rgb}{0.321569,0.321569,0.321569}%
\pgfsetstrokecolor{currentstroke}%
\pgfsetstrokeopacity{0.900000}%
\pgfsetdash{}{0pt}%
\pgfpathmoveto{\pgfqpoint{1.949611in}{4.076503in}}%
\pgfpathcurveto{\pgfqpoint{1.959865in}{4.076503in}}{\pgfqpoint{1.969700in}{4.080577in}}{\pgfqpoint{1.976951in}{4.087828in}}%
\pgfpathcurveto{\pgfqpoint{1.984202in}{4.095078in}}{\pgfqpoint{1.988276in}{4.104914in}}{\pgfqpoint{1.988276in}{4.115168in}}%
\pgfpathcurveto{\pgfqpoint{1.988276in}{4.125422in}}{\pgfqpoint{1.984202in}{4.135258in}}{\pgfqpoint{1.976951in}{4.142508in}}%
\pgfpathcurveto{\pgfqpoint{1.969700in}{4.149759in}}{\pgfqpoint{1.959865in}{4.153833in}}{\pgfqpoint{1.949611in}{4.153833in}}%
\pgfpathcurveto{\pgfqpoint{1.939357in}{4.153833in}}{\pgfqpoint{1.929521in}{4.149759in}}{\pgfqpoint{1.922271in}{4.142508in}}%
\pgfpathcurveto{\pgfqpoint{1.915020in}{4.135258in}}{\pgfqpoint{1.910946in}{4.125422in}}{\pgfqpoint{1.910946in}{4.115168in}}%
\pgfpathcurveto{\pgfqpoint{1.910946in}{4.104914in}}{\pgfqpoint{1.915020in}{4.095078in}}{\pgfqpoint{1.922271in}{4.087828in}}%
\pgfpathcurveto{\pgfqpoint{1.929521in}{4.080577in}}{\pgfqpoint{1.939357in}{4.076503in}}{\pgfqpoint{1.949611in}{4.076503in}}%
\pgfpathlineto{\pgfqpoint{1.949611in}{4.076503in}}%
\pgfpathclose%
\pgfusepath{stroke}%
\end{pgfscope}%
\begin{pgfscope}%
\pgfpathrectangle{\pgfqpoint{0.977909in}{2.928000in}}{\pgfqpoint{5.664000in}{5.664000in}}%
\pgfusepath{clip}%
\pgfsetbuttcap%
\pgfsetroundjoin%
\pgfsetlinewidth{1.003750pt}%
\definecolor{currentstroke}{rgb}{0.321569,0.321569,0.321569}%
\pgfsetstrokecolor{currentstroke}%
\pgfsetstrokeopacity{0.900000}%
\pgfsetdash{}{0pt}%
\pgfpathmoveto{\pgfqpoint{1.111178in}{3.193859in}}%
\pgfpathcurveto{\pgfqpoint{1.121432in}{3.193859in}}{\pgfqpoint{1.131267in}{3.197933in}}{\pgfqpoint{1.138518in}{3.205183in}}%
\pgfpathcurveto{\pgfqpoint{1.145769in}{3.212434in}}{\pgfqpoint{1.149843in}{3.222269in}}{\pgfqpoint{1.149843in}{3.232524in}}%
\pgfpathcurveto{\pgfqpoint{1.149843in}{3.242778in}}{\pgfqpoint{1.145769in}{3.252613in}}{\pgfqpoint{1.138518in}{3.259864in}}%
\pgfpathcurveto{\pgfqpoint{1.131267in}{3.267115in}}{\pgfqpoint{1.121432in}{3.271189in}}{\pgfqpoint{1.111178in}{3.271189in}}%
\pgfpathcurveto{\pgfqpoint{1.100923in}{3.271189in}}{\pgfqpoint{1.091088in}{3.267115in}}{\pgfqpoint{1.083837in}{3.259864in}}%
\pgfpathcurveto{\pgfqpoint{1.076587in}{3.252613in}}{\pgfqpoint{1.072513in}{3.242778in}}{\pgfqpoint{1.072513in}{3.232524in}}%
\pgfpathcurveto{\pgfqpoint{1.072513in}{3.222269in}}{\pgfqpoint{1.076587in}{3.212434in}}{\pgfqpoint{1.083837in}{3.205183in}}%
\pgfpathcurveto{\pgfqpoint{1.091088in}{3.197933in}}{\pgfqpoint{1.100923in}{3.193859in}}{\pgfqpoint{1.111178in}{3.193859in}}%
\pgfpathlineto{\pgfqpoint{1.111178in}{3.193859in}}%
\pgfpathclose%
\pgfusepath{stroke}%
\end{pgfscope}%
\begin{pgfscope}%
\pgfpathrectangle{\pgfqpoint{0.977909in}{2.928000in}}{\pgfqpoint{5.664000in}{5.664000in}}%
\pgfusepath{clip}%
\pgfsetbuttcap%
\pgfsetroundjoin%
\pgfsetlinewidth{1.003750pt}%
\definecolor{currentstroke}{rgb}{0.321569,0.321569,0.321569}%
\pgfsetstrokecolor{currentstroke}%
\pgfsetstrokeopacity{0.900000}%
\pgfsetdash{}{0pt}%
\pgfpathmoveto{\pgfqpoint{4.198226in}{6.899800in}}%
\pgfpathcurveto{\pgfqpoint{4.208480in}{6.899800in}}{\pgfqpoint{4.218316in}{6.903874in}}{\pgfqpoint{4.225566in}{6.911125in}}%
\pgfpathcurveto{\pgfqpoint{4.232817in}{6.918375in}}{\pgfqpoint{4.236891in}{6.928211in}}{\pgfqpoint{4.236891in}{6.938465in}}%
\pgfpathcurveto{\pgfqpoint{4.236891in}{6.948719in}}{\pgfqpoint{4.232817in}{6.958554in}}{\pgfqpoint{4.225566in}{6.965805in}}%
\pgfpathcurveto{\pgfqpoint{4.218316in}{6.973056in}}{\pgfqpoint{4.208480in}{6.977130in}}{\pgfqpoint{4.198226in}{6.977130in}}%
\pgfpathcurveto{\pgfqpoint{4.187972in}{6.977130in}}{\pgfqpoint{4.178136in}{6.973056in}}{\pgfqpoint{4.170886in}{6.965805in}}%
\pgfpathcurveto{\pgfqpoint{4.163635in}{6.958554in}}{\pgfqpoint{4.159561in}{6.948719in}}{\pgfqpoint{4.159561in}{6.938465in}}%
\pgfpathcurveto{\pgfqpoint{4.159561in}{6.928211in}}{\pgfqpoint{4.163635in}{6.918375in}}{\pgfqpoint{4.170886in}{6.911125in}}%
\pgfpathcurveto{\pgfqpoint{4.178136in}{6.903874in}}{\pgfqpoint{4.187972in}{6.899800in}}{\pgfqpoint{4.198226in}{6.899800in}}%
\pgfpathlineto{\pgfqpoint{4.198226in}{6.899800in}}%
\pgfpathclose%
\pgfusepath{stroke}%
\end{pgfscope}%
\begin{pgfscope}%
\pgfpathrectangle{\pgfqpoint{0.977909in}{2.928000in}}{\pgfqpoint{5.664000in}{5.664000in}}%
\pgfusepath{clip}%
\pgfsetbuttcap%
\pgfsetroundjoin%
\pgfsetlinewidth{1.003750pt}%
\definecolor{currentstroke}{rgb}{0.321569,0.321569,0.321569}%
\pgfsetstrokecolor{currentstroke}%
\pgfsetstrokeopacity{0.900000}%
\pgfsetdash{}{0pt}%
\pgfpathmoveto{\pgfqpoint{3.616469in}{4.873611in}}%
\pgfpathcurveto{\pgfqpoint{3.626723in}{4.873611in}}{\pgfqpoint{3.636558in}{4.877685in}}{\pgfqpoint{3.643809in}{4.884936in}}%
\pgfpathcurveto{\pgfqpoint{3.651060in}{4.892187in}}{\pgfqpoint{3.655134in}{4.902022in}}{\pgfqpoint{3.655134in}{4.912276in}}%
\pgfpathcurveto{\pgfqpoint{3.655134in}{4.922530in}}{\pgfqpoint{3.651060in}{4.932366in}}{\pgfqpoint{3.643809in}{4.939616in}}%
\pgfpathcurveto{\pgfqpoint{3.636558in}{4.946867in}}{\pgfqpoint{3.626723in}{4.950941in}}{\pgfqpoint{3.616469in}{4.950941in}}%
\pgfpathcurveto{\pgfqpoint{3.606214in}{4.950941in}}{\pgfqpoint{3.596379in}{4.946867in}}{\pgfqpoint{3.589128in}{4.939616in}}%
\pgfpathcurveto{\pgfqpoint{3.581878in}{4.932366in}}{\pgfqpoint{3.577804in}{4.922530in}}{\pgfqpoint{3.577804in}{4.912276in}}%
\pgfpathcurveto{\pgfqpoint{3.577804in}{4.902022in}}{\pgfqpoint{3.581878in}{4.892187in}}{\pgfqpoint{3.589128in}{4.884936in}}%
\pgfpathcurveto{\pgfqpoint{3.596379in}{4.877685in}}{\pgfqpoint{3.606214in}{4.873611in}}{\pgfqpoint{3.616469in}{4.873611in}}%
\pgfpathlineto{\pgfqpoint{3.616469in}{4.873611in}}%
\pgfpathclose%
\pgfusepath{stroke}%
\end{pgfscope}%
\begin{pgfscope}%
\pgfsetbuttcap%
\pgfsetroundjoin%
\pgfsetlinewidth{1.003750pt}%
\definecolor{currentstroke}{rgb}{0.321569,0.321569,0.321569}%
\pgfsetstrokecolor{currentstroke}%
\pgfsetstrokeopacity{0.900000}%
\pgfsetdash{}{0pt}%
\pgfpathmoveto{\pgfqpoint{0.977909in}{2.868262in}}%
\pgfpathlineto{\pgfqpoint{1.037647in}{2.928000in}}%
\pgfpathlineto{\pgfqpoint{0.977909in}{2.987738in}}%
\pgfpathlineto{\pgfqpoint{0.918171in}{2.928000in}}%
\pgfpathlineto{\pgfqpoint{0.977909in}{2.868262in}}%
\pgfpathclose%
\pgfusepath{stroke}%
\end{pgfscope}%
\begin{pgfscope}%
\pgfsetbuttcap%
\pgfsetroundjoin%
\pgfsetlinewidth{1.003750pt}%
\definecolor{currentstroke}{rgb}{0.321569,0.321569,0.321569}%
\pgfsetstrokecolor{currentstroke}%
\pgfsetstrokeopacity{0.900000}%
\pgfsetdash{}{0pt}%
\pgfpathmoveto{\pgfqpoint{0.977909in}{2.868262in}}%
\pgfpathlineto{\pgfqpoint{1.037647in}{2.928000in}}%
\pgfpathlineto{\pgfqpoint{0.977909in}{2.987738in}}%
\pgfpathlineto{\pgfqpoint{0.918171in}{2.928000in}}%
\pgfpathlineto{\pgfqpoint{0.977909in}{2.868262in}}%
\pgfpathclose%
\pgfusepath{stroke}%
\end{pgfscope}%
\begin{pgfscope}%
\pgfsetbuttcap%
\pgfsetroundjoin%
\pgfsetlinewidth{1.003750pt}%
\definecolor{currentstroke}{rgb}{0.321569,0.321569,0.321569}%
\pgfsetstrokecolor{currentstroke}%
\pgfsetstrokeopacity{0.900000}%
\pgfsetdash{}{0pt}%
\pgfpathmoveto{\pgfqpoint{0.977909in}{2.868262in}}%
\pgfpathlineto{\pgfqpoint{1.037647in}{2.928000in}}%
\pgfpathlineto{\pgfqpoint{0.977909in}{2.987738in}}%
\pgfpathlineto{\pgfqpoint{0.918171in}{2.928000in}}%
\pgfpathlineto{\pgfqpoint{0.977909in}{2.868262in}}%
\pgfpathclose%
\pgfusepath{stroke}%
\end{pgfscope}%
\begin{pgfscope}%
\pgfsetbuttcap%
\pgfsetroundjoin%
\pgfsetlinewidth{1.003750pt}%
\definecolor{currentstroke}{rgb}{0.321569,0.321569,0.321569}%
\pgfsetstrokecolor{currentstroke}%
\pgfsetstrokeopacity{0.900000}%
\pgfsetdash{}{0pt}%
\pgfpathmoveto{\pgfqpoint{0.977909in}{3.377848in}}%
\pgfpathlineto{\pgfqpoint{1.037647in}{3.437586in}}%
\pgfpathlineto{\pgfqpoint{0.977909in}{3.497324in}}%
\pgfpathlineto{\pgfqpoint{0.918171in}{3.437586in}}%
\pgfpathlineto{\pgfqpoint{0.977909in}{3.377848in}}%
\pgfpathclose%
\pgfusepath{stroke}%
\end{pgfscope}%
\begin{pgfscope}%
\pgfsetbuttcap%
\pgfsetroundjoin%
\pgfsetlinewidth{1.003750pt}%
\definecolor{currentstroke}{rgb}{0.321569,0.321569,0.321569}%
\pgfsetstrokecolor{currentstroke}%
\pgfsetstrokeopacity{0.900000}%
\pgfsetdash{}{0pt}%
\pgfpathmoveto{\pgfqpoint{0.977909in}{3.223976in}}%
\pgfpathlineto{\pgfqpoint{1.037647in}{3.283715in}}%
\pgfpathlineto{\pgfqpoint{0.977909in}{3.343453in}}%
\pgfpathlineto{\pgfqpoint{0.918171in}{3.283715in}}%
\pgfpathlineto{\pgfqpoint{0.977909in}{3.223976in}}%
\pgfpathclose%
\pgfusepath{stroke}%
\end{pgfscope}%
\begin{pgfscope}%
\pgfpathrectangle{\pgfqpoint{0.977909in}{2.928000in}}{\pgfqpoint{5.664000in}{5.664000in}}%
\pgfusepath{clip}%
\pgfsetbuttcap%
\pgfsetroundjoin%
\pgfsetlinewidth{1.003750pt}%
\definecolor{currentstroke}{rgb}{0.768627,0.603922,0.000000}%
\pgfsetstrokecolor{currentstroke}%
\pgfsetstrokeopacity{0.900000}%
\pgfsetdash{}{0pt}%
\pgfpathmoveto{\pgfqpoint{1.412553in}{7.616874in}}%
\pgfpathcurveto{\pgfqpoint{1.422807in}{7.616874in}}{\pgfqpoint{1.432643in}{7.620948in}}{\pgfqpoint{1.439893in}{7.628199in}}%
\pgfpathcurveto{\pgfqpoint{1.447144in}{7.635449in}}{\pgfqpoint{1.451218in}{7.645285in}}{\pgfqpoint{1.451218in}{7.655539in}}%
\pgfpathcurveto{\pgfqpoint{1.451218in}{7.665793in}}{\pgfqpoint{1.447144in}{7.675629in}}{\pgfqpoint{1.439893in}{7.682879in}}%
\pgfpathcurveto{\pgfqpoint{1.432643in}{7.690130in}}{\pgfqpoint{1.422807in}{7.694204in}}{\pgfqpoint{1.412553in}{7.694204in}}%
\pgfpathcurveto{\pgfqpoint{1.402299in}{7.694204in}}{\pgfqpoint{1.392464in}{7.690130in}}{\pgfqpoint{1.385213in}{7.682879in}}%
\pgfpathcurveto{\pgfqpoint{1.377962in}{7.675629in}}{\pgfqpoint{1.373888in}{7.665793in}}{\pgfqpoint{1.373888in}{7.655539in}}%
\pgfpathcurveto{\pgfqpoint{1.373888in}{7.645285in}}{\pgfqpoint{1.377962in}{7.635449in}}{\pgfqpoint{1.385213in}{7.628199in}}%
\pgfpathcurveto{\pgfqpoint{1.392464in}{7.620948in}}{\pgfqpoint{1.402299in}{7.616874in}}{\pgfqpoint{1.412553in}{7.616874in}}%
\pgfpathlineto{\pgfqpoint{1.412553in}{7.616874in}}%
\pgfpathclose%
\pgfusepath{stroke}%
\end{pgfscope}%
\begin{pgfscope}%
\pgfpathrectangle{\pgfqpoint{0.977909in}{2.928000in}}{\pgfqpoint{5.664000in}{5.664000in}}%
\pgfusepath{clip}%
\pgfsetbuttcap%
\pgfsetroundjoin%
\pgfsetlinewidth{1.003750pt}%
\definecolor{currentstroke}{rgb}{0.768627,0.603922,0.000000}%
\pgfsetstrokecolor{currentstroke}%
\pgfsetstrokeopacity{0.900000}%
\pgfsetdash{}{0pt}%
\pgfpathmoveto{\pgfqpoint{4.144153in}{6.453877in}}%
\pgfpathcurveto{\pgfqpoint{4.154407in}{6.453877in}}{\pgfqpoint{4.164243in}{6.457951in}}{\pgfqpoint{4.171494in}{6.465202in}}%
\pgfpathcurveto{\pgfqpoint{4.178744in}{6.472453in}}{\pgfqpoint{4.182818in}{6.482288in}}{\pgfqpoint{4.182818in}{6.492542in}}%
\pgfpathcurveto{\pgfqpoint{4.182818in}{6.502796in}}{\pgfqpoint{4.178744in}{6.512632in}}{\pgfqpoint{4.171494in}{6.519883in}}%
\pgfpathcurveto{\pgfqpoint{4.164243in}{6.527133in}}{\pgfqpoint{4.154407in}{6.531207in}}{\pgfqpoint{4.144153in}{6.531207in}}%
\pgfpathcurveto{\pgfqpoint{4.133899in}{6.531207in}}{\pgfqpoint{4.124064in}{6.527133in}}{\pgfqpoint{4.116813in}{6.519883in}}%
\pgfpathcurveto{\pgfqpoint{4.109562in}{6.512632in}}{\pgfqpoint{4.105488in}{6.502796in}}{\pgfqpoint{4.105488in}{6.492542in}}%
\pgfpathcurveto{\pgfqpoint{4.105488in}{6.482288in}}{\pgfqpoint{4.109562in}{6.472453in}}{\pgfqpoint{4.116813in}{6.465202in}}%
\pgfpathcurveto{\pgfqpoint{4.124064in}{6.457951in}}{\pgfqpoint{4.133899in}{6.453877in}}{\pgfqpoint{4.144153in}{6.453877in}}%
\pgfpathlineto{\pgfqpoint{4.144153in}{6.453877in}}%
\pgfpathclose%
\pgfusepath{stroke}%
\end{pgfscope}%
\begin{pgfscope}%
\pgfpathrectangle{\pgfqpoint{0.977909in}{2.928000in}}{\pgfqpoint{5.664000in}{5.664000in}}%
\pgfusepath{clip}%
\pgfsetbuttcap%
\pgfsetroundjoin%
\pgfsetlinewidth{1.003750pt}%
\definecolor{currentstroke}{rgb}{0.768627,0.603922,0.000000}%
\pgfsetstrokecolor{currentstroke}%
\pgfsetstrokeopacity{0.900000}%
\pgfsetdash{}{0pt}%
\pgfpathmoveto{\pgfqpoint{5.038011in}{6.711265in}}%
\pgfpathcurveto{\pgfqpoint{5.048266in}{6.711265in}}{\pgfqpoint{5.058101in}{6.715339in}}{\pgfqpoint{5.065352in}{6.722590in}}%
\pgfpathcurveto{\pgfqpoint{5.072602in}{6.729841in}}{\pgfqpoint{5.076676in}{6.739676in}}{\pgfqpoint{5.076676in}{6.749930in}}%
\pgfpathcurveto{\pgfqpoint{5.076676in}{6.760184in}}{\pgfqpoint{5.072602in}{6.770020in}}{\pgfqpoint{5.065352in}{6.777270in}}%
\pgfpathcurveto{\pgfqpoint{5.058101in}{6.784521in}}{\pgfqpoint{5.048266in}{6.788595in}}{\pgfqpoint{5.038011in}{6.788595in}}%
\pgfpathcurveto{\pgfqpoint{5.027757in}{6.788595in}}{\pgfqpoint{5.017922in}{6.784521in}}{\pgfqpoint{5.010671in}{6.777270in}}%
\pgfpathcurveto{\pgfqpoint{5.003420in}{6.770020in}}{\pgfqpoint{4.999346in}{6.760184in}}{\pgfqpoint{4.999346in}{6.749930in}}%
\pgfpathcurveto{\pgfqpoint{4.999346in}{6.739676in}}{\pgfqpoint{5.003420in}{6.729841in}}{\pgfqpoint{5.010671in}{6.722590in}}%
\pgfpathcurveto{\pgfqpoint{5.017922in}{6.715339in}}{\pgfqpoint{5.027757in}{6.711265in}}{\pgfqpoint{5.038011in}{6.711265in}}%
\pgfpathlineto{\pgfqpoint{5.038011in}{6.711265in}}%
\pgfpathclose%
\pgfusepath{stroke}%
\end{pgfscope}%
\begin{pgfscope}%
\pgfpathrectangle{\pgfqpoint{0.977909in}{2.928000in}}{\pgfqpoint{5.664000in}{5.664000in}}%
\pgfusepath{clip}%
\pgfsetbuttcap%
\pgfsetroundjoin%
\pgfsetlinewidth{1.003750pt}%
\definecolor{currentstroke}{rgb}{0.768627,0.603922,0.000000}%
\pgfsetstrokecolor{currentstroke}%
\pgfsetstrokeopacity{0.900000}%
\pgfsetdash{}{0pt}%
\pgfpathmoveto{\pgfqpoint{5.049550in}{6.791913in}}%
\pgfpathcurveto{\pgfqpoint{5.059804in}{6.791913in}}{\pgfqpoint{5.069639in}{6.795987in}}{\pgfqpoint{5.076890in}{6.803238in}}%
\pgfpathcurveto{\pgfqpoint{5.084141in}{6.810489in}}{\pgfqpoint{5.088215in}{6.820324in}}{\pgfqpoint{5.088215in}{6.830578in}}%
\pgfpathcurveto{\pgfqpoint{5.088215in}{6.840832in}}{\pgfqpoint{5.084141in}{6.850668in}}{\pgfqpoint{5.076890in}{6.857918in}}%
\pgfpathcurveto{\pgfqpoint{5.069639in}{6.865169in}}{\pgfqpoint{5.059804in}{6.869243in}}{\pgfqpoint{5.049550in}{6.869243in}}%
\pgfpathcurveto{\pgfqpoint{5.039296in}{6.869243in}}{\pgfqpoint{5.029460in}{6.865169in}}{\pgfqpoint{5.022210in}{6.857918in}}%
\pgfpathcurveto{\pgfqpoint{5.014959in}{6.850668in}}{\pgfqpoint{5.010885in}{6.840832in}}{\pgfqpoint{5.010885in}{6.830578in}}%
\pgfpathcurveto{\pgfqpoint{5.010885in}{6.820324in}}{\pgfqpoint{5.014959in}{6.810489in}}{\pgfqpoint{5.022210in}{6.803238in}}%
\pgfpathcurveto{\pgfqpoint{5.029460in}{6.795987in}}{\pgfqpoint{5.039296in}{6.791913in}}{\pgfqpoint{5.049550in}{6.791913in}}%
\pgfpathlineto{\pgfqpoint{5.049550in}{6.791913in}}%
\pgfpathclose%
\pgfusepath{stroke}%
\end{pgfscope}%
\begin{pgfscope}%
\pgfpathrectangle{\pgfqpoint{0.977909in}{2.928000in}}{\pgfqpoint{5.664000in}{5.664000in}}%
\pgfusepath{clip}%
\pgfsetbuttcap%
\pgfsetroundjoin%
\pgfsetlinewidth{1.003750pt}%
\definecolor{currentstroke}{rgb}{0.768627,0.603922,0.000000}%
\pgfsetstrokecolor{currentstroke}%
\pgfsetstrokeopacity{0.900000}%
\pgfsetdash{}{0pt}%
\pgfpathmoveto{\pgfqpoint{4.971123in}{6.732874in}}%
\pgfpathcurveto{\pgfqpoint{4.981377in}{6.732874in}}{\pgfqpoint{4.991212in}{6.736948in}}{\pgfqpoint{4.998463in}{6.744199in}}%
\pgfpathcurveto{\pgfqpoint{5.005714in}{6.751450in}}{\pgfqpoint{5.009788in}{6.761285in}}{\pgfqpoint{5.009788in}{6.771539in}}%
\pgfpathcurveto{\pgfqpoint{5.009788in}{6.781793in}}{\pgfqpoint{5.005714in}{6.791629in}}{\pgfqpoint{4.998463in}{6.798880in}}%
\pgfpathcurveto{\pgfqpoint{4.991212in}{6.806130in}}{\pgfqpoint{4.981377in}{6.810204in}}{\pgfqpoint{4.971123in}{6.810204in}}%
\pgfpathcurveto{\pgfqpoint{4.960868in}{6.810204in}}{\pgfqpoint{4.951033in}{6.806130in}}{\pgfqpoint{4.943782in}{6.798880in}}%
\pgfpathcurveto{\pgfqpoint{4.936531in}{6.791629in}}{\pgfqpoint{4.932457in}{6.781793in}}{\pgfqpoint{4.932457in}{6.771539in}}%
\pgfpathcurveto{\pgfqpoint{4.932457in}{6.761285in}}{\pgfqpoint{4.936531in}{6.751450in}}{\pgfqpoint{4.943782in}{6.744199in}}%
\pgfpathcurveto{\pgfqpoint{4.951033in}{6.736948in}}{\pgfqpoint{4.960868in}{6.732874in}}{\pgfqpoint{4.971123in}{6.732874in}}%
\pgfpathlineto{\pgfqpoint{4.971123in}{6.732874in}}%
\pgfpathclose%
\pgfusepath{stroke}%
\end{pgfscope}%
\begin{pgfscope}%
\pgfpathrectangle{\pgfqpoint{0.977909in}{2.928000in}}{\pgfqpoint{5.664000in}{5.664000in}}%
\pgfusepath{clip}%
\pgfsetbuttcap%
\pgfsetroundjoin%
\pgfsetlinewidth{1.003750pt}%
\definecolor{currentstroke}{rgb}{0.768627,0.603922,0.000000}%
\pgfsetstrokecolor{currentstroke}%
\pgfsetstrokeopacity{0.900000}%
\pgfsetdash{}{0pt}%
\pgfpathmoveto{\pgfqpoint{4.948712in}{6.665444in}}%
\pgfpathcurveto{\pgfqpoint{4.958966in}{6.665444in}}{\pgfqpoint{4.968802in}{6.669518in}}{\pgfqpoint{4.976053in}{6.676769in}}%
\pgfpathcurveto{\pgfqpoint{4.983303in}{6.684019in}}{\pgfqpoint{4.987377in}{6.693855in}}{\pgfqpoint{4.987377in}{6.704109in}}%
\pgfpathcurveto{\pgfqpoint{4.987377in}{6.714363in}}{\pgfqpoint{4.983303in}{6.724199in}}{\pgfqpoint{4.976053in}{6.731449in}}%
\pgfpathcurveto{\pgfqpoint{4.968802in}{6.738700in}}{\pgfqpoint{4.958966in}{6.742774in}}{\pgfqpoint{4.948712in}{6.742774in}}%
\pgfpathcurveto{\pgfqpoint{4.938458in}{6.742774in}}{\pgfqpoint{4.928623in}{6.738700in}}{\pgfqpoint{4.921372in}{6.731449in}}%
\pgfpathcurveto{\pgfqpoint{4.914121in}{6.724199in}}{\pgfqpoint{4.910047in}{6.714363in}}{\pgfqpoint{4.910047in}{6.704109in}}%
\pgfpathcurveto{\pgfqpoint{4.910047in}{6.693855in}}{\pgfqpoint{4.914121in}{6.684019in}}{\pgfqpoint{4.921372in}{6.676769in}}%
\pgfpathcurveto{\pgfqpoint{4.928623in}{6.669518in}}{\pgfqpoint{4.938458in}{6.665444in}}{\pgfqpoint{4.948712in}{6.665444in}}%
\pgfpathlineto{\pgfqpoint{4.948712in}{6.665444in}}%
\pgfpathclose%
\pgfusepath{stroke}%
\end{pgfscope}%
\begin{pgfscope}%
\pgfpathrectangle{\pgfqpoint{0.977909in}{2.928000in}}{\pgfqpoint{5.664000in}{5.664000in}}%
\pgfusepath{clip}%
\pgfsetbuttcap%
\pgfsetroundjoin%
\pgfsetlinewidth{1.003750pt}%
\definecolor{currentstroke}{rgb}{0.200000,0.133333,0.533333}%
\pgfsetstrokecolor{currentstroke}%
\pgfsetstrokeopacity{0.900000}%
\pgfsetdash{}{0pt}%
\pgfpathmoveto{\pgfqpoint{1.801670in}{3.427583in}}%
\pgfpathcurveto{\pgfqpoint{1.811924in}{3.427583in}}{\pgfqpoint{1.821759in}{3.431657in}}{\pgfqpoint{1.829010in}{3.438907in}}%
\pgfpathcurveto{\pgfqpoint{1.836261in}{3.446158in}}{\pgfqpoint{1.840335in}{3.455993in}}{\pgfqpoint{1.840335in}{3.466248in}}%
\pgfpathcurveto{\pgfqpoint{1.840335in}{3.476502in}}{\pgfqpoint{1.836261in}{3.486337in}}{\pgfqpoint{1.829010in}{3.493588in}}%
\pgfpathcurveto{\pgfqpoint{1.821759in}{3.500839in}}{\pgfqpoint{1.811924in}{3.504913in}}{\pgfqpoint{1.801670in}{3.504913in}}%
\pgfpathcurveto{\pgfqpoint{1.791416in}{3.504913in}}{\pgfqpoint{1.781580in}{3.500839in}}{\pgfqpoint{1.774330in}{3.493588in}}%
\pgfpathcurveto{\pgfqpoint{1.767079in}{3.486337in}}{\pgfqpoint{1.763005in}{3.476502in}}{\pgfqpoint{1.763005in}{3.466248in}}%
\pgfpathcurveto{\pgfqpoint{1.763005in}{3.455993in}}{\pgfqpoint{1.767079in}{3.446158in}}{\pgfqpoint{1.774330in}{3.438907in}}%
\pgfpathcurveto{\pgfqpoint{1.781580in}{3.431657in}}{\pgfqpoint{1.791416in}{3.427583in}}{\pgfqpoint{1.801670in}{3.427583in}}%
\pgfpathlineto{\pgfqpoint{1.801670in}{3.427583in}}%
\pgfpathclose%
\pgfusepath{stroke}%
\end{pgfscope}%
\begin{pgfscope}%
\pgfpathrectangle{\pgfqpoint{0.977909in}{2.928000in}}{\pgfqpoint{5.664000in}{5.664000in}}%
\pgfusepath{clip}%
\pgfsetbuttcap%
\pgfsetroundjoin%
\pgfsetlinewidth{1.003750pt}%
\definecolor{currentstroke}{rgb}{0.200000,0.133333,0.533333}%
\pgfsetstrokecolor{currentstroke}%
\pgfsetstrokeopacity{0.900000}%
\pgfsetdash{}{0pt}%
\pgfpathmoveto{\pgfqpoint{2.086038in}{3.551767in}}%
\pgfpathcurveto{\pgfqpoint{2.096292in}{3.551767in}}{\pgfqpoint{2.106128in}{3.555841in}}{\pgfqpoint{2.113379in}{3.563091in}}%
\pgfpathcurveto{\pgfqpoint{2.120629in}{3.570342in}}{\pgfqpoint{2.124703in}{3.580177in}}{\pgfqpoint{2.124703in}{3.590432in}}%
\pgfpathcurveto{\pgfqpoint{2.124703in}{3.600686in}}{\pgfqpoint{2.120629in}{3.610521in}}{\pgfqpoint{2.113379in}{3.617772in}}%
\pgfpathcurveto{\pgfqpoint{2.106128in}{3.625023in}}{\pgfqpoint{2.096292in}{3.629097in}}{\pgfqpoint{2.086038in}{3.629097in}}%
\pgfpathcurveto{\pgfqpoint{2.075784in}{3.629097in}}{\pgfqpoint{2.065949in}{3.625023in}}{\pgfqpoint{2.058698in}{3.617772in}}%
\pgfpathcurveto{\pgfqpoint{2.051447in}{3.610521in}}{\pgfqpoint{2.047373in}{3.600686in}}{\pgfqpoint{2.047373in}{3.590432in}}%
\pgfpathcurveto{\pgfqpoint{2.047373in}{3.580177in}}{\pgfqpoint{2.051447in}{3.570342in}}{\pgfqpoint{2.058698in}{3.563091in}}%
\pgfpathcurveto{\pgfqpoint{2.065949in}{3.555841in}}{\pgfqpoint{2.075784in}{3.551767in}}{\pgfqpoint{2.086038in}{3.551767in}}%
\pgfpathlineto{\pgfqpoint{2.086038in}{3.551767in}}%
\pgfpathclose%
\pgfusepath{stroke}%
\end{pgfscope}%
\begin{pgfscope}%
\pgfpathrectangle{\pgfqpoint{0.977909in}{2.928000in}}{\pgfqpoint{5.664000in}{5.664000in}}%
\pgfusepath{clip}%
\pgfsetbuttcap%
\pgfsetroundjoin%
\pgfsetlinewidth{1.003750pt}%
\definecolor{currentstroke}{rgb}{0.200000,0.133333,0.533333}%
\pgfsetstrokecolor{currentstroke}%
\pgfsetstrokeopacity{0.900000}%
\pgfsetdash{}{0pt}%
\pgfpathmoveto{\pgfqpoint{2.693764in}{3.586556in}}%
\pgfpathcurveto{\pgfqpoint{2.704018in}{3.586556in}}{\pgfqpoint{2.713853in}{3.590630in}}{\pgfqpoint{2.721104in}{3.597881in}}%
\pgfpathcurveto{\pgfqpoint{2.728355in}{3.605132in}}{\pgfqpoint{2.732429in}{3.614967in}}{\pgfqpoint{2.732429in}{3.625221in}}%
\pgfpathcurveto{\pgfqpoint{2.732429in}{3.635475in}}{\pgfqpoint{2.728355in}{3.645311in}}{\pgfqpoint{2.721104in}{3.652562in}}%
\pgfpathcurveto{\pgfqpoint{2.713853in}{3.659812in}}{\pgfqpoint{2.704018in}{3.663886in}}{\pgfqpoint{2.693764in}{3.663886in}}%
\pgfpathcurveto{\pgfqpoint{2.683509in}{3.663886in}}{\pgfqpoint{2.673674in}{3.659812in}}{\pgfqpoint{2.666423in}{3.652562in}}%
\pgfpathcurveto{\pgfqpoint{2.659172in}{3.645311in}}{\pgfqpoint{2.655099in}{3.635475in}}{\pgfqpoint{2.655099in}{3.625221in}}%
\pgfpathcurveto{\pgfqpoint{2.655099in}{3.614967in}}{\pgfqpoint{2.659172in}{3.605132in}}{\pgfqpoint{2.666423in}{3.597881in}}%
\pgfpathcurveto{\pgfqpoint{2.673674in}{3.590630in}}{\pgfqpoint{2.683509in}{3.586556in}}{\pgfqpoint{2.693764in}{3.586556in}}%
\pgfpathlineto{\pgfqpoint{2.693764in}{3.586556in}}%
\pgfpathclose%
\pgfusepath{stroke}%
\end{pgfscope}%
\begin{pgfscope}%
\pgfpathrectangle{\pgfqpoint{0.977909in}{2.928000in}}{\pgfqpoint{5.664000in}{5.664000in}}%
\pgfusepath{clip}%
\pgfsetbuttcap%
\pgfsetroundjoin%
\pgfsetlinewidth{1.003750pt}%
\definecolor{currentstroke}{rgb}{0.200000,0.133333,0.533333}%
\pgfsetstrokecolor{currentstroke}%
\pgfsetstrokeopacity{0.900000}%
\pgfsetdash{}{0pt}%
\pgfpathmoveto{\pgfqpoint{3.571020in}{4.220784in}}%
\pgfpathcurveto{\pgfqpoint{3.581274in}{4.220784in}}{\pgfqpoint{3.591110in}{4.224858in}}{\pgfqpoint{3.598361in}{4.232109in}}%
\pgfpathcurveto{\pgfqpoint{3.605611in}{4.239359in}}{\pgfqpoint{3.609685in}{4.249195in}}{\pgfqpoint{3.609685in}{4.259449in}}%
\pgfpathcurveto{\pgfqpoint{3.609685in}{4.269703in}}{\pgfqpoint{3.605611in}{4.279539in}}{\pgfqpoint{3.598361in}{4.286789in}}%
\pgfpathcurveto{\pgfqpoint{3.591110in}{4.294040in}}{\pgfqpoint{3.581274in}{4.298114in}}{\pgfqpoint{3.571020in}{4.298114in}}%
\pgfpathcurveto{\pgfqpoint{3.560766in}{4.298114in}}{\pgfqpoint{3.550931in}{4.294040in}}{\pgfqpoint{3.543680in}{4.286789in}}%
\pgfpathcurveto{\pgfqpoint{3.536429in}{4.279539in}}{\pgfqpoint{3.532355in}{4.269703in}}{\pgfqpoint{3.532355in}{4.259449in}}%
\pgfpathcurveto{\pgfqpoint{3.532355in}{4.249195in}}{\pgfqpoint{3.536429in}{4.239359in}}{\pgfqpoint{3.543680in}{4.232109in}}%
\pgfpathcurveto{\pgfqpoint{3.550931in}{4.224858in}}{\pgfqpoint{3.560766in}{4.220784in}}{\pgfqpoint{3.571020in}{4.220784in}}%
\pgfpathlineto{\pgfqpoint{3.571020in}{4.220784in}}%
\pgfpathclose%
\pgfusepath{stroke}%
\end{pgfscope}%
\begin{pgfscope}%
\pgfpathrectangle{\pgfqpoint{0.977909in}{2.928000in}}{\pgfqpoint{5.664000in}{5.664000in}}%
\pgfusepath{clip}%
\pgfsetbuttcap%
\pgfsetroundjoin%
\pgfsetlinewidth{1.003750pt}%
\definecolor{currentstroke}{rgb}{0.200000,0.133333,0.533333}%
\pgfsetstrokecolor{currentstroke}%
\pgfsetstrokeopacity{0.900000}%
\pgfsetdash{}{0pt}%
\pgfpathmoveto{\pgfqpoint{2.517118in}{3.494914in}}%
\pgfpathcurveto{\pgfqpoint{2.527372in}{3.494914in}}{\pgfqpoint{2.537208in}{3.498988in}}{\pgfqpoint{2.544459in}{3.506238in}}%
\pgfpathcurveto{\pgfqpoint{2.551709in}{3.513489in}}{\pgfqpoint{2.555783in}{3.523324in}}{\pgfqpoint{2.555783in}{3.533579in}}%
\pgfpathcurveto{\pgfqpoint{2.555783in}{3.543833in}}{\pgfqpoint{2.551709in}{3.553668in}}{\pgfqpoint{2.544459in}{3.560919in}}%
\pgfpathcurveto{\pgfqpoint{2.537208in}{3.568170in}}{\pgfqpoint{2.527372in}{3.572244in}}{\pgfqpoint{2.517118in}{3.572244in}}%
\pgfpathcurveto{\pgfqpoint{2.506864in}{3.572244in}}{\pgfqpoint{2.497029in}{3.568170in}}{\pgfqpoint{2.489778in}{3.560919in}}%
\pgfpathcurveto{\pgfqpoint{2.482527in}{3.553668in}}{\pgfqpoint{2.478453in}{3.543833in}}{\pgfqpoint{2.478453in}{3.533579in}}%
\pgfpathcurveto{\pgfqpoint{2.478453in}{3.523324in}}{\pgfqpoint{2.482527in}{3.513489in}}{\pgfqpoint{2.489778in}{3.506238in}}%
\pgfpathcurveto{\pgfqpoint{2.497029in}{3.498988in}}{\pgfqpoint{2.506864in}{3.494914in}}{\pgfqpoint{2.517118in}{3.494914in}}%
\pgfpathlineto{\pgfqpoint{2.517118in}{3.494914in}}%
\pgfpathclose%
\pgfusepath{stroke}%
\end{pgfscope}%
\begin{pgfscope}%
\pgfpathrectangle{\pgfqpoint{0.977909in}{2.928000in}}{\pgfqpoint{5.664000in}{5.664000in}}%
\pgfusepath{clip}%
\pgfsetbuttcap%
\pgfsetroundjoin%
\pgfsetlinewidth{1.003750pt}%
\definecolor{currentstroke}{rgb}{0.200000,0.133333,0.533333}%
\pgfsetstrokecolor{currentstroke}%
\pgfsetstrokeopacity{0.900000}%
\pgfsetdash{}{0pt}%
\pgfpathmoveto{\pgfqpoint{3.405967in}{4.237148in}}%
\pgfpathcurveto{\pgfqpoint{3.416221in}{4.237148in}}{\pgfqpoint{3.426056in}{4.241222in}}{\pgfqpoint{3.433307in}{4.248472in}}%
\pgfpathcurveto{\pgfqpoint{3.440558in}{4.255723in}}{\pgfqpoint{3.444632in}{4.265559in}}{\pgfqpoint{3.444632in}{4.275813in}}%
\pgfpathcurveto{\pgfqpoint{3.444632in}{4.286067in}}{\pgfqpoint{3.440558in}{4.295902in}}{\pgfqpoint{3.433307in}{4.303153in}}%
\pgfpathcurveto{\pgfqpoint{3.426056in}{4.310404in}}{\pgfqpoint{3.416221in}{4.314478in}}{\pgfqpoint{3.405967in}{4.314478in}}%
\pgfpathcurveto{\pgfqpoint{3.395713in}{4.314478in}}{\pgfqpoint{3.385877in}{4.310404in}}{\pgfqpoint{3.378626in}{4.303153in}}%
\pgfpathcurveto{\pgfqpoint{3.371376in}{4.295902in}}{\pgfqpoint{3.367302in}{4.286067in}}{\pgfqpoint{3.367302in}{4.275813in}}%
\pgfpathcurveto{\pgfqpoint{3.367302in}{4.265559in}}{\pgfqpoint{3.371376in}{4.255723in}}{\pgfqpoint{3.378626in}{4.248472in}}%
\pgfpathcurveto{\pgfqpoint{3.385877in}{4.241222in}}{\pgfqpoint{3.395713in}{4.237148in}}{\pgfqpoint{3.405967in}{4.237148in}}%
\pgfpathlineto{\pgfqpoint{3.405967in}{4.237148in}}%
\pgfpathclose%
\pgfusepath{stroke}%
\end{pgfscope}%
\begin{pgfscope}%
\pgfpathrectangle{\pgfqpoint{0.977909in}{2.928000in}}{\pgfqpoint{5.664000in}{5.664000in}}%
\pgfusepath{clip}%
\pgfsetbuttcap%
\pgfsetroundjoin%
\pgfsetlinewidth{1.003750pt}%
\definecolor{currentstroke}{rgb}{0.176471,0.713725,0.643137}%
\pgfsetstrokecolor{currentstroke}%
\pgfsetstrokeopacity{0.900000}%
\pgfsetdash{}{0pt}%
\pgfpathmoveto{\pgfqpoint{5.401022in}{7.552858in}}%
\pgfpathcurveto{\pgfqpoint{5.411276in}{7.552858in}}{\pgfqpoint{5.421111in}{7.556932in}}{\pgfqpoint{5.428362in}{7.564183in}}%
\pgfpathcurveto{\pgfqpoint{5.435613in}{7.571434in}}{\pgfqpoint{5.439687in}{7.581269in}}{\pgfqpoint{5.439687in}{7.591523in}}%
\pgfpathcurveto{\pgfqpoint{5.439687in}{7.601777in}}{\pgfqpoint{5.435613in}{7.611613in}}{\pgfqpoint{5.428362in}{7.618864in}}%
\pgfpathcurveto{\pgfqpoint{5.421111in}{7.626114in}}{\pgfqpoint{5.411276in}{7.630188in}}{\pgfqpoint{5.401022in}{7.630188in}}%
\pgfpathcurveto{\pgfqpoint{5.390768in}{7.630188in}}{\pgfqpoint{5.380932in}{7.626114in}}{\pgfqpoint{5.373681in}{7.618864in}}%
\pgfpathcurveto{\pgfqpoint{5.366431in}{7.611613in}}{\pgfqpoint{5.362357in}{7.601777in}}{\pgfqpoint{5.362357in}{7.591523in}}%
\pgfpathcurveto{\pgfqpoint{5.362357in}{7.581269in}}{\pgfqpoint{5.366431in}{7.571434in}}{\pgfqpoint{5.373681in}{7.564183in}}%
\pgfpathcurveto{\pgfqpoint{5.380932in}{7.556932in}}{\pgfqpoint{5.390768in}{7.552858in}}{\pgfqpoint{5.401022in}{7.552858in}}%
\pgfpathlineto{\pgfqpoint{5.401022in}{7.552858in}}%
\pgfpathclose%
\pgfusepath{stroke}%
\end{pgfscope}%
\begin{pgfscope}%
\pgfpathrectangle{\pgfqpoint{0.977909in}{2.928000in}}{\pgfqpoint{5.664000in}{5.664000in}}%
\pgfusepath{clip}%
\pgfsetbuttcap%
\pgfsetroundjoin%
\pgfsetlinewidth{1.003750pt}%
\definecolor{currentstroke}{rgb}{0.176471,0.713725,0.643137}%
\pgfsetstrokecolor{currentstroke}%
\pgfsetstrokeopacity{0.900000}%
\pgfsetdash{}{0pt}%
\pgfpathmoveto{\pgfqpoint{5.376212in}{7.787616in}}%
\pgfpathcurveto{\pgfqpoint{5.386466in}{7.787616in}}{\pgfqpoint{5.396302in}{7.791690in}}{\pgfqpoint{5.403552in}{7.798941in}}%
\pgfpathcurveto{\pgfqpoint{5.410803in}{7.806192in}}{\pgfqpoint{5.414877in}{7.816027in}}{\pgfqpoint{5.414877in}{7.826281in}}%
\pgfpathcurveto{\pgfqpoint{5.414877in}{7.836535in}}{\pgfqpoint{5.410803in}{7.846371in}}{\pgfqpoint{5.403552in}{7.853621in}}%
\pgfpathcurveto{\pgfqpoint{5.396302in}{7.860872in}}{\pgfqpoint{5.386466in}{7.864946in}}{\pgfqpoint{5.376212in}{7.864946in}}%
\pgfpathcurveto{\pgfqpoint{5.365958in}{7.864946in}}{\pgfqpoint{5.356123in}{7.860872in}}{\pgfqpoint{5.348872in}{7.853621in}}%
\pgfpathcurveto{\pgfqpoint{5.341621in}{7.846371in}}{\pgfqpoint{5.337547in}{7.836535in}}{\pgfqpoint{5.337547in}{7.826281in}}%
\pgfpathcurveto{\pgfqpoint{5.337547in}{7.816027in}}{\pgfqpoint{5.341621in}{7.806192in}}{\pgfqpoint{5.348872in}{7.798941in}}%
\pgfpathcurveto{\pgfqpoint{5.356123in}{7.791690in}}{\pgfqpoint{5.365958in}{7.787616in}}{\pgfqpoint{5.376212in}{7.787616in}}%
\pgfpathlineto{\pgfqpoint{5.376212in}{7.787616in}}%
\pgfpathclose%
\pgfusepath{stroke}%
\end{pgfscope}%
\begin{pgfscope}%
\pgfpathrectangle{\pgfqpoint{0.977909in}{2.928000in}}{\pgfqpoint{5.664000in}{5.664000in}}%
\pgfusepath{clip}%
\pgfsetbuttcap%
\pgfsetroundjoin%
\pgfsetlinewidth{1.003750pt}%
\definecolor{currentstroke}{rgb}{0.176471,0.713725,0.643137}%
\pgfsetstrokecolor{currentstroke}%
\pgfsetstrokeopacity{0.900000}%
\pgfsetdash{}{0pt}%
\pgfpathmoveto{\pgfqpoint{5.744517in}{7.985280in}}%
\pgfpathcurveto{\pgfqpoint{5.754771in}{7.985280in}}{\pgfqpoint{5.764607in}{7.989354in}}{\pgfqpoint{5.771858in}{7.996604in}}%
\pgfpathcurveto{\pgfqpoint{5.779108in}{8.003855in}}{\pgfqpoint{5.783182in}{8.013691in}}{\pgfqpoint{5.783182in}{8.023945in}}%
\pgfpathcurveto{\pgfqpoint{5.783182in}{8.034199in}}{\pgfqpoint{5.779108in}{8.044034in}}{\pgfqpoint{5.771858in}{8.051285in}}%
\pgfpathcurveto{\pgfqpoint{5.764607in}{8.058536in}}{\pgfqpoint{5.754771in}{8.062610in}}{\pgfqpoint{5.744517in}{8.062610in}}%
\pgfpathcurveto{\pgfqpoint{5.734263in}{8.062610in}}{\pgfqpoint{5.724428in}{8.058536in}}{\pgfqpoint{5.717177in}{8.051285in}}%
\pgfpathcurveto{\pgfqpoint{5.709926in}{8.044034in}}{\pgfqpoint{5.705852in}{8.034199in}}{\pgfqpoint{5.705852in}{8.023945in}}%
\pgfpathcurveto{\pgfqpoint{5.705852in}{8.013691in}}{\pgfqpoint{5.709926in}{8.003855in}}{\pgfqpoint{5.717177in}{7.996604in}}%
\pgfpathcurveto{\pgfqpoint{5.724428in}{7.989354in}}{\pgfqpoint{5.734263in}{7.985280in}}{\pgfqpoint{5.744517in}{7.985280in}}%
\pgfpathlineto{\pgfqpoint{5.744517in}{7.985280in}}%
\pgfpathclose%
\pgfusepath{stroke}%
\end{pgfscope}%
\begin{pgfscope}%
\pgfpathrectangle{\pgfqpoint{0.977909in}{2.928000in}}{\pgfqpoint{5.664000in}{5.664000in}}%
\pgfusepath{clip}%
\pgfsetbuttcap%
\pgfsetroundjoin%
\pgfsetlinewidth{1.003750pt}%
\definecolor{currentstroke}{rgb}{0.176471,0.713725,0.643137}%
\pgfsetstrokecolor{currentstroke}%
\pgfsetstrokeopacity{0.900000}%
\pgfsetdash{}{0pt}%
\pgfpathmoveto{\pgfqpoint{4.034825in}{5.191452in}}%
\pgfpathcurveto{\pgfqpoint{4.045079in}{5.191452in}}{\pgfqpoint{4.054914in}{5.195526in}}{\pgfqpoint{4.062165in}{5.202777in}}%
\pgfpathcurveto{\pgfqpoint{4.069416in}{5.210027in}}{\pgfqpoint{4.073490in}{5.219863in}}{\pgfqpoint{4.073490in}{5.230117in}}%
\pgfpathcurveto{\pgfqpoint{4.073490in}{5.240371in}}{\pgfqpoint{4.069416in}{5.250207in}}{\pgfqpoint{4.062165in}{5.257457in}}%
\pgfpathcurveto{\pgfqpoint{4.054914in}{5.264708in}}{\pgfqpoint{4.045079in}{5.268782in}}{\pgfqpoint{4.034825in}{5.268782in}}%
\pgfpathcurveto{\pgfqpoint{4.024571in}{5.268782in}}{\pgfqpoint{4.014735in}{5.264708in}}{\pgfqpoint{4.007484in}{5.257457in}}%
\pgfpathcurveto{\pgfqpoint{4.000234in}{5.250207in}}{\pgfqpoint{3.996160in}{5.240371in}}{\pgfqpoint{3.996160in}{5.230117in}}%
\pgfpathcurveto{\pgfqpoint{3.996160in}{5.219863in}}{\pgfqpoint{4.000234in}{5.210027in}}{\pgfqpoint{4.007484in}{5.202777in}}%
\pgfpathcurveto{\pgfqpoint{4.014735in}{5.195526in}}{\pgfqpoint{4.024571in}{5.191452in}}{\pgfqpoint{4.034825in}{5.191452in}}%
\pgfpathlineto{\pgfqpoint{4.034825in}{5.191452in}}%
\pgfpathclose%
\pgfusepath{stroke}%
\end{pgfscope}%
\begin{pgfscope}%
\pgfsetbuttcap%
\pgfsetroundjoin%
\pgfsetlinewidth{1.003750pt}%
\definecolor{currentstroke}{rgb}{0.176471,0.713725,0.643137}%
\pgfsetstrokecolor{currentstroke}%
\pgfsetstrokeopacity{0.900000}%
\pgfsetdash{}{0pt}%
\pgfpathmoveto{\pgfqpoint{5.804792in}{8.532262in}}%
\pgfpathlineto{\pgfqpoint{5.864530in}{8.592000in}}%
\pgfpathlineto{\pgfqpoint{5.804792in}{8.651738in}}%
\pgfpathlineto{\pgfqpoint{5.745054in}{8.592000in}}%
\pgfpathlineto{\pgfqpoint{5.804792in}{8.532262in}}%
\pgfpathclose%
\pgfusepath{stroke}%
\end{pgfscope}%
\begin{pgfscope}%
\pgfpathrectangle{\pgfqpoint{0.977909in}{2.928000in}}{\pgfqpoint{5.664000in}{5.664000in}}%
\pgfusepath{clip}%
\pgfsetbuttcap%
\pgfsetroundjoin%
\pgfsetlinewidth{1.003750pt}%
\definecolor{currentstroke}{rgb}{0.000000,0.447059,0.698039}%
\pgfsetstrokecolor{currentstroke}%
\pgfsetstrokeopacity{0.900000}%
\pgfsetdash{}{0pt}%
\pgfpathmoveto{\pgfqpoint{3.718566in}{6.064948in}}%
\pgfpathcurveto{\pgfqpoint{3.728820in}{6.064948in}}{\pgfqpoint{3.738655in}{6.069022in}}{\pgfqpoint{3.745906in}{6.076272in}}%
\pgfpathcurveto{\pgfqpoint{3.753157in}{6.083523in}}{\pgfqpoint{3.757231in}{6.093359in}}{\pgfqpoint{3.757231in}{6.103613in}}%
\pgfpathcurveto{\pgfqpoint{3.757231in}{6.113867in}}{\pgfqpoint{3.753157in}{6.123702in}}{\pgfqpoint{3.745906in}{6.130953in}}%
\pgfpathcurveto{\pgfqpoint{3.738655in}{6.138204in}}{\pgfqpoint{3.728820in}{6.142278in}}{\pgfqpoint{3.718566in}{6.142278in}}%
\pgfpathcurveto{\pgfqpoint{3.708312in}{6.142278in}}{\pgfqpoint{3.698476in}{6.138204in}}{\pgfqpoint{3.691226in}{6.130953in}}%
\pgfpathcurveto{\pgfqpoint{3.683975in}{6.123702in}}{\pgfqpoint{3.679901in}{6.113867in}}{\pgfqpoint{3.679901in}{6.103613in}}%
\pgfpathcurveto{\pgfqpoint{3.679901in}{6.093359in}}{\pgfqpoint{3.683975in}{6.083523in}}{\pgfqpoint{3.691226in}{6.076272in}}%
\pgfpathcurveto{\pgfqpoint{3.698476in}{6.069022in}}{\pgfqpoint{3.708312in}{6.064948in}}{\pgfqpoint{3.718566in}{6.064948in}}%
\pgfpathlineto{\pgfqpoint{3.718566in}{6.064948in}}%
\pgfpathclose%
\pgfusepath{stroke}%
\end{pgfscope}%
\begin{pgfscope}%
\pgfpathrectangle{\pgfqpoint{0.977909in}{2.928000in}}{\pgfqpoint{5.664000in}{5.664000in}}%
\pgfusepath{clip}%
\pgfsetbuttcap%
\pgfsetroundjoin%
\pgfsetlinewidth{1.003750pt}%
\definecolor{currentstroke}{rgb}{0.901961,0.403922,0.619608}%
\pgfsetstrokecolor{currentstroke}%
\pgfsetstrokeopacity{0.900000}%
\pgfsetdash{}{0pt}%
\pgfpathmoveto{\pgfqpoint{3.354265in}{4.152510in}}%
\pgfpathcurveto{\pgfqpoint{3.364519in}{4.152510in}}{\pgfqpoint{3.374354in}{4.156584in}}{\pgfqpoint{3.381605in}{4.163835in}}%
\pgfpathcurveto{\pgfqpoint{3.388856in}{4.171086in}}{\pgfqpoint{3.392930in}{4.180921in}}{\pgfqpoint{3.392930in}{4.191175in}}%
\pgfpathcurveto{\pgfqpoint{3.392930in}{4.201429in}}{\pgfqpoint{3.388856in}{4.211265in}}{\pgfqpoint{3.381605in}{4.218516in}}%
\pgfpathcurveto{\pgfqpoint{3.374354in}{4.225766in}}{\pgfqpoint{3.364519in}{4.229840in}}{\pgfqpoint{3.354265in}{4.229840in}}%
\pgfpathcurveto{\pgfqpoint{3.344011in}{4.229840in}}{\pgfqpoint{3.334175in}{4.225766in}}{\pgfqpoint{3.326924in}{4.218516in}}%
\pgfpathcurveto{\pgfqpoint{3.319674in}{4.211265in}}{\pgfqpoint{3.315600in}{4.201429in}}{\pgfqpoint{3.315600in}{4.191175in}}%
\pgfpathcurveto{\pgfqpoint{3.315600in}{4.180921in}}{\pgfqpoint{3.319674in}{4.171086in}}{\pgfqpoint{3.326924in}{4.163835in}}%
\pgfpathcurveto{\pgfqpoint{3.334175in}{4.156584in}}{\pgfqpoint{3.344011in}{4.152510in}}{\pgfqpoint{3.354265in}{4.152510in}}%
\pgfpathlineto{\pgfqpoint{3.354265in}{4.152510in}}%
\pgfpathclose%
\pgfusepath{stroke}%
\end{pgfscope}%
\begin{pgfscope}%
\pgfpathrectangle{\pgfqpoint{0.977909in}{2.928000in}}{\pgfqpoint{5.664000in}{5.664000in}}%
\pgfusepath{clip}%
\pgfsetbuttcap%
\pgfsetroundjoin%
\pgfsetlinewidth{1.003750pt}%
\definecolor{currentstroke}{rgb}{0.901961,0.403922,0.619608}%
\pgfsetstrokecolor{currentstroke}%
\pgfsetstrokeopacity{0.900000}%
\pgfsetdash{}{0pt}%
\pgfpathmoveto{\pgfqpoint{4.447912in}{7.202983in}}%
\pgfpathcurveto{\pgfqpoint{4.458166in}{7.202983in}}{\pgfqpoint{4.468002in}{7.207057in}}{\pgfqpoint{4.475252in}{7.214308in}}%
\pgfpathcurveto{\pgfqpoint{4.482503in}{7.221559in}}{\pgfqpoint{4.486577in}{7.231394in}}{\pgfqpoint{4.486577in}{7.241648in}}%
\pgfpathcurveto{\pgfqpoint{4.486577in}{7.251902in}}{\pgfqpoint{4.482503in}{7.261738in}}{\pgfqpoint{4.475252in}{7.268988in}}%
\pgfpathcurveto{\pgfqpoint{4.468002in}{7.276239in}}{\pgfqpoint{4.458166in}{7.280313in}}{\pgfqpoint{4.447912in}{7.280313in}}%
\pgfpathcurveto{\pgfqpoint{4.437658in}{7.280313in}}{\pgfqpoint{4.427822in}{7.276239in}}{\pgfqpoint{4.420572in}{7.268988in}}%
\pgfpathcurveto{\pgfqpoint{4.413321in}{7.261738in}}{\pgfqpoint{4.409247in}{7.251902in}}{\pgfqpoint{4.409247in}{7.241648in}}%
\pgfpathcurveto{\pgfqpoint{4.409247in}{7.231394in}}{\pgfqpoint{4.413321in}{7.221559in}}{\pgfqpoint{4.420572in}{7.214308in}}%
\pgfpathcurveto{\pgfqpoint{4.427822in}{7.207057in}}{\pgfqpoint{4.437658in}{7.202983in}}{\pgfqpoint{4.447912in}{7.202983in}}%
\pgfpathlineto{\pgfqpoint{4.447912in}{7.202983in}}%
\pgfpathclose%
\pgfusepath{stroke}%
\end{pgfscope}%
\begin{pgfscope}%
\pgfpathrectangle{\pgfqpoint{0.977909in}{2.928000in}}{\pgfqpoint{5.664000in}{5.664000in}}%
\pgfusepath{clip}%
\pgfsetbuttcap%
\pgfsetroundjoin%
\pgfsetlinewidth{1.003750pt}%
\definecolor{currentstroke}{rgb}{0.901961,0.403922,0.619608}%
\pgfsetstrokecolor{currentstroke}%
\pgfsetstrokeopacity{0.900000}%
\pgfsetdash{}{0pt}%
\pgfpathmoveto{\pgfqpoint{4.431879in}{6.684557in}}%
\pgfpathcurveto{\pgfqpoint{4.442133in}{6.684557in}}{\pgfqpoint{4.451969in}{6.688631in}}{\pgfqpoint{4.459220in}{6.695882in}}%
\pgfpathcurveto{\pgfqpoint{4.466470in}{6.703133in}}{\pgfqpoint{4.470544in}{6.712968in}}{\pgfqpoint{4.470544in}{6.723222in}}%
\pgfpathcurveto{\pgfqpoint{4.470544in}{6.733477in}}{\pgfqpoint{4.466470in}{6.743312in}}{\pgfqpoint{4.459220in}{6.750563in}}%
\pgfpathcurveto{\pgfqpoint{4.451969in}{6.757814in}}{\pgfqpoint{4.442133in}{6.761888in}}{\pgfqpoint{4.431879in}{6.761888in}}%
\pgfpathcurveto{\pgfqpoint{4.421625in}{6.761888in}}{\pgfqpoint{4.411790in}{6.757814in}}{\pgfqpoint{4.404539in}{6.750563in}}%
\pgfpathcurveto{\pgfqpoint{4.397288in}{6.743312in}}{\pgfqpoint{4.393214in}{6.733477in}}{\pgfqpoint{4.393214in}{6.723222in}}%
\pgfpathcurveto{\pgfqpoint{4.393214in}{6.712968in}}{\pgfqpoint{4.397288in}{6.703133in}}{\pgfqpoint{4.404539in}{6.695882in}}%
\pgfpathcurveto{\pgfqpoint{4.411790in}{6.688631in}}{\pgfqpoint{4.421625in}{6.684557in}}{\pgfqpoint{4.431879in}{6.684557in}}%
\pgfpathlineto{\pgfqpoint{4.431879in}{6.684557in}}%
\pgfpathclose%
\pgfusepath{stroke}%
\end{pgfscope}%
\begin{pgfscope}%
\pgfpathrectangle{\pgfqpoint{0.977909in}{2.928000in}}{\pgfqpoint{5.664000in}{5.664000in}}%
\pgfusepath{clip}%
\pgfsetbuttcap%
\pgfsetroundjoin%
\pgfsetlinewidth{1.003750pt}%
\definecolor{currentstroke}{rgb}{0.901961,0.403922,0.619608}%
\pgfsetstrokecolor{currentstroke}%
\pgfsetstrokeopacity{0.900000}%
\pgfsetdash{}{0pt}%
\pgfpathmoveto{\pgfqpoint{4.098156in}{5.261997in}}%
\pgfpathcurveto{\pgfqpoint{4.108410in}{5.261997in}}{\pgfqpoint{4.118246in}{5.266071in}}{\pgfqpoint{4.125496in}{5.273322in}}%
\pgfpathcurveto{\pgfqpoint{4.132747in}{5.280573in}}{\pgfqpoint{4.136821in}{5.290408in}}{\pgfqpoint{4.136821in}{5.300662in}}%
\pgfpathcurveto{\pgfqpoint{4.136821in}{5.310916in}}{\pgfqpoint{4.132747in}{5.320752in}}{\pgfqpoint{4.125496in}{5.328002in}}%
\pgfpathcurveto{\pgfqpoint{4.118246in}{5.335253in}}{\pgfqpoint{4.108410in}{5.339327in}}{\pgfqpoint{4.098156in}{5.339327in}}%
\pgfpathcurveto{\pgfqpoint{4.087902in}{5.339327in}}{\pgfqpoint{4.078066in}{5.335253in}}{\pgfqpoint{4.070816in}{5.328002in}}%
\pgfpathcurveto{\pgfqpoint{4.063565in}{5.320752in}}{\pgfqpoint{4.059491in}{5.310916in}}{\pgfqpoint{4.059491in}{5.300662in}}%
\pgfpathcurveto{\pgfqpoint{4.059491in}{5.290408in}}{\pgfqpoint{4.063565in}{5.280573in}}{\pgfqpoint{4.070816in}{5.273322in}}%
\pgfpathcurveto{\pgfqpoint{4.078066in}{5.266071in}}{\pgfqpoint{4.087902in}{5.261997in}}{\pgfqpoint{4.098156in}{5.261997in}}%
\pgfpathlineto{\pgfqpoint{4.098156in}{5.261997in}}%
\pgfpathclose%
\pgfusepath{stroke}%
\end{pgfscope}%
\begin{pgfscope}%
\pgfpathrectangle{\pgfqpoint{0.977909in}{2.928000in}}{\pgfqpoint{5.664000in}{5.664000in}}%
\pgfusepath{clip}%
\pgfsetbuttcap%
\pgfsetroundjoin%
\pgfsetlinewidth{1.003750pt}%
\definecolor{currentstroke}{rgb}{0.941176,0.584314,0.337255}%
\pgfsetstrokecolor{currentstroke}%
\pgfsetstrokeopacity{0.900000}%
\pgfsetdash{}{0pt}%
\pgfpathmoveto{\pgfqpoint{3.517164in}{7.859497in}}%
\pgfpathcurveto{\pgfqpoint{3.527418in}{7.859497in}}{\pgfqpoint{3.537254in}{7.863571in}}{\pgfqpoint{3.544504in}{7.870822in}}%
\pgfpathcurveto{\pgfqpoint{3.551755in}{7.878072in}}{\pgfqpoint{3.555829in}{7.887908in}}{\pgfqpoint{3.555829in}{7.898162in}}%
\pgfpathcurveto{\pgfqpoint{3.555829in}{7.908416in}}{\pgfqpoint{3.551755in}{7.918251in}}{\pgfqpoint{3.544504in}{7.925502in}}%
\pgfpathcurveto{\pgfqpoint{3.537254in}{7.932753in}}{\pgfqpoint{3.527418in}{7.936827in}}{\pgfqpoint{3.517164in}{7.936827in}}%
\pgfpathcurveto{\pgfqpoint{3.506910in}{7.936827in}}{\pgfqpoint{3.497075in}{7.932753in}}{\pgfqpoint{3.489824in}{7.925502in}}%
\pgfpathcurveto{\pgfqpoint{3.482573in}{7.918251in}}{\pgfqpoint{3.478499in}{7.908416in}}{\pgfqpoint{3.478499in}{7.898162in}}%
\pgfpathcurveto{\pgfqpoint{3.478499in}{7.887908in}}{\pgfqpoint{3.482573in}{7.878072in}}{\pgfqpoint{3.489824in}{7.870822in}}%
\pgfpathcurveto{\pgfqpoint{3.497075in}{7.863571in}}{\pgfqpoint{3.506910in}{7.859497in}}{\pgfqpoint{3.517164in}{7.859497in}}%
\pgfpathlineto{\pgfqpoint{3.517164in}{7.859497in}}%
\pgfpathclose%
\pgfusepath{stroke}%
\end{pgfscope}%
\begin{pgfscope}%
\pgfpathrectangle{\pgfqpoint{0.977909in}{2.928000in}}{\pgfqpoint{5.664000in}{5.664000in}}%
\pgfusepath{clip}%
\pgfsetbuttcap%
\pgfsetroundjoin%
\pgfsetlinewidth{1.003750pt}%
\definecolor{currentstroke}{rgb}{0.941176,0.584314,0.337255}%
\pgfsetstrokecolor{currentstroke}%
\pgfsetstrokeopacity{0.900000}%
\pgfsetdash{}{0pt}%
\pgfpathmoveto{\pgfqpoint{3.961945in}{7.014558in}}%
\pgfpathcurveto{\pgfqpoint{3.972199in}{7.014558in}}{\pgfqpoint{3.982035in}{7.018632in}}{\pgfqpoint{3.989286in}{7.025883in}}%
\pgfpathcurveto{\pgfqpoint{3.996536in}{7.033134in}}{\pgfqpoint{4.000610in}{7.042969in}}{\pgfqpoint{4.000610in}{7.053223in}}%
\pgfpathcurveto{\pgfqpoint{4.000610in}{7.063478in}}{\pgfqpoint{3.996536in}{7.073313in}}{\pgfqpoint{3.989286in}{7.080564in}}%
\pgfpathcurveto{\pgfqpoint{3.982035in}{7.087814in}}{\pgfqpoint{3.972199in}{7.091888in}}{\pgfqpoint{3.961945in}{7.091888in}}%
\pgfpathcurveto{\pgfqpoint{3.951691in}{7.091888in}}{\pgfqpoint{3.941856in}{7.087814in}}{\pgfqpoint{3.934605in}{7.080564in}}%
\pgfpathcurveto{\pgfqpoint{3.927354in}{7.073313in}}{\pgfqpoint{3.923280in}{7.063478in}}{\pgfqpoint{3.923280in}{7.053223in}}%
\pgfpathcurveto{\pgfqpoint{3.923280in}{7.042969in}}{\pgfqpoint{3.927354in}{7.033134in}}{\pgfqpoint{3.934605in}{7.025883in}}%
\pgfpathcurveto{\pgfqpoint{3.941856in}{7.018632in}}{\pgfqpoint{3.951691in}{7.014558in}}{\pgfqpoint{3.961945in}{7.014558in}}%
\pgfpathlineto{\pgfqpoint{3.961945in}{7.014558in}}%
\pgfpathclose%
\pgfusepath{stroke}%
\end{pgfscope}%
\begin{pgfscope}%
\pgfpathrectangle{\pgfqpoint{0.977909in}{2.928000in}}{\pgfqpoint{5.664000in}{5.664000in}}%
\pgfusepath{clip}%
\pgfsetbuttcap%
\pgfsetroundjoin%
\pgfsetlinewidth{1.003750pt}%
\definecolor{currentstroke}{rgb}{0.835294,0.368627,0.000000}%
\pgfsetstrokecolor{currentstroke}%
\pgfsetstrokeopacity{0.900000}%
\pgfsetdash{}{0pt}%
\pgfpathmoveto{\pgfqpoint{5.658202in}{6.534007in}}%
\pgfpathcurveto{\pgfqpoint{5.668456in}{6.534007in}}{\pgfqpoint{5.678291in}{6.538081in}}{\pgfqpoint{5.685542in}{6.545332in}}%
\pgfpathcurveto{\pgfqpoint{5.692793in}{6.552582in}}{\pgfqpoint{5.696867in}{6.562418in}}{\pgfqpoint{5.696867in}{6.572672in}}%
\pgfpathcurveto{\pgfqpoint{5.696867in}{6.582926in}}{\pgfqpoint{5.692793in}{6.592762in}}{\pgfqpoint{5.685542in}{6.600012in}}%
\pgfpathcurveto{\pgfqpoint{5.678291in}{6.607263in}}{\pgfqpoint{5.668456in}{6.611337in}}{\pgfqpoint{5.658202in}{6.611337in}}%
\pgfpathcurveto{\pgfqpoint{5.647948in}{6.611337in}}{\pgfqpoint{5.638112in}{6.607263in}}{\pgfqpoint{5.630861in}{6.600012in}}%
\pgfpathcurveto{\pgfqpoint{5.623611in}{6.592762in}}{\pgfqpoint{5.619537in}{6.582926in}}{\pgfqpoint{5.619537in}{6.572672in}}%
\pgfpathcurveto{\pgfqpoint{5.619537in}{6.562418in}}{\pgfqpoint{5.623611in}{6.552582in}}{\pgfqpoint{5.630861in}{6.545332in}}%
\pgfpathcurveto{\pgfqpoint{5.638112in}{6.538081in}}{\pgfqpoint{5.647948in}{6.534007in}}{\pgfqpoint{5.658202in}{6.534007in}}%
\pgfpathlineto{\pgfqpoint{5.658202in}{6.534007in}}%
\pgfpathclose%
\pgfusepath{stroke}%
\end{pgfscope}%
\begin{pgfscope}%
\pgfpathrectangle{\pgfqpoint{0.977909in}{2.928000in}}{\pgfqpoint{5.664000in}{5.664000in}}%
\pgfusepath{clip}%
\pgfsetbuttcap%
\pgfsetroundjoin%
\pgfsetlinewidth{1.003750pt}%
\definecolor{currentstroke}{rgb}{0.835294,0.368627,0.000000}%
\pgfsetstrokecolor{currentstroke}%
\pgfsetstrokeopacity{0.900000}%
\pgfsetdash{}{0pt}%
\pgfpathmoveto{\pgfqpoint{5.085996in}{6.296399in}}%
\pgfpathcurveto{\pgfqpoint{5.096250in}{6.296399in}}{\pgfqpoint{5.106085in}{6.300473in}}{\pgfqpoint{5.113336in}{6.307724in}}%
\pgfpathcurveto{\pgfqpoint{5.120587in}{6.314974in}}{\pgfqpoint{5.124661in}{6.324810in}}{\pgfqpoint{5.124661in}{6.335064in}}%
\pgfpathcurveto{\pgfqpoint{5.124661in}{6.345318in}}{\pgfqpoint{5.120587in}{6.355153in}}{\pgfqpoint{5.113336in}{6.362404in}}%
\pgfpathcurveto{\pgfqpoint{5.106085in}{6.369655in}}{\pgfqpoint{5.096250in}{6.373729in}}{\pgfqpoint{5.085996in}{6.373729in}}%
\pgfpathcurveto{\pgfqpoint{5.075742in}{6.373729in}}{\pgfqpoint{5.065906in}{6.369655in}}{\pgfqpoint{5.058655in}{6.362404in}}%
\pgfpathcurveto{\pgfqpoint{5.051405in}{6.355153in}}{\pgfqpoint{5.047331in}{6.345318in}}{\pgfqpoint{5.047331in}{6.335064in}}%
\pgfpathcurveto{\pgfqpoint{5.047331in}{6.324810in}}{\pgfqpoint{5.051405in}{6.314974in}}{\pgfqpoint{5.058655in}{6.307724in}}%
\pgfpathcurveto{\pgfqpoint{5.065906in}{6.300473in}}{\pgfqpoint{5.075742in}{6.296399in}}{\pgfqpoint{5.085996in}{6.296399in}}%
\pgfpathlineto{\pgfqpoint{5.085996in}{6.296399in}}%
\pgfpathclose%
\pgfusepath{stroke}%
\end{pgfscope}%
\begin{pgfscope}%
\pgfpathrectangle{\pgfqpoint{0.977909in}{2.928000in}}{\pgfqpoint{5.664000in}{5.664000in}}%
\pgfusepath{clip}%
\pgfsetbuttcap%
\pgfsetroundjoin%
\pgfsetlinewidth{1.003750pt}%
\definecolor{currentstroke}{rgb}{0.549020,0.160784,0.505882}%
\pgfsetstrokecolor{currentstroke}%
\pgfsetstrokeopacity{0.900000}%
\pgfsetdash{}{0pt}%
\pgfpathmoveto{\pgfqpoint{4.526598in}{6.752782in}}%
\pgfpathcurveto{\pgfqpoint{4.536852in}{6.752782in}}{\pgfqpoint{4.546687in}{6.756856in}}{\pgfqpoint{4.553938in}{6.764106in}}%
\pgfpathcurveto{\pgfqpoint{4.561189in}{6.771357in}}{\pgfqpoint{4.565263in}{6.781193in}}{\pgfqpoint{4.565263in}{6.791447in}}%
\pgfpathcurveto{\pgfqpoint{4.565263in}{6.801701in}}{\pgfqpoint{4.561189in}{6.811536in}}{\pgfqpoint{4.553938in}{6.818787in}}%
\pgfpathcurveto{\pgfqpoint{4.546687in}{6.826038in}}{\pgfqpoint{4.536852in}{6.830112in}}{\pgfqpoint{4.526598in}{6.830112in}}%
\pgfpathcurveto{\pgfqpoint{4.516344in}{6.830112in}}{\pgfqpoint{4.506508in}{6.826038in}}{\pgfqpoint{4.499258in}{6.818787in}}%
\pgfpathcurveto{\pgfqpoint{4.492007in}{6.811536in}}{\pgfqpoint{4.487933in}{6.801701in}}{\pgfqpoint{4.487933in}{6.791447in}}%
\pgfpathcurveto{\pgfqpoint{4.487933in}{6.781193in}}{\pgfqpoint{4.492007in}{6.771357in}}{\pgfqpoint{4.499258in}{6.764106in}}%
\pgfpathcurveto{\pgfqpoint{4.506508in}{6.756856in}}{\pgfqpoint{4.516344in}{6.752782in}}{\pgfqpoint{4.526598in}{6.752782in}}%
\pgfpathlineto{\pgfqpoint{4.526598in}{6.752782in}}%
\pgfpathclose%
\pgfusepath{stroke}%
\end{pgfscope}%
\begin{pgfscope}%
\pgfpathrectangle{\pgfqpoint{0.977909in}{2.928000in}}{\pgfqpoint{5.664000in}{5.664000in}}%
\pgfusepath{clip}%
\pgfsetbuttcap%
\pgfsetroundjoin%
\pgfsetlinewidth{1.003750pt}%
\definecolor{currentstroke}{rgb}{0.149020,0.274510,0.325490}%
\pgfsetstrokecolor{currentstroke}%
\pgfsetstrokeopacity{0.900000}%
\pgfsetdash{}{0pt}%
\pgfpathmoveto{\pgfqpoint{4.605728in}{6.492294in}}%
\pgfpathcurveto{\pgfqpoint{4.615982in}{6.492294in}}{\pgfqpoint{4.625817in}{6.496368in}}{\pgfqpoint{4.633068in}{6.503619in}}%
\pgfpathcurveto{\pgfqpoint{4.640319in}{6.510870in}}{\pgfqpoint{4.644393in}{6.520705in}}{\pgfqpoint{4.644393in}{6.530959in}}%
\pgfpathcurveto{\pgfqpoint{4.644393in}{6.541213in}}{\pgfqpoint{4.640319in}{6.551049in}}{\pgfqpoint{4.633068in}{6.558300in}}%
\pgfpathcurveto{\pgfqpoint{4.625817in}{6.565550in}}{\pgfqpoint{4.615982in}{6.569624in}}{\pgfqpoint{4.605728in}{6.569624in}}%
\pgfpathcurveto{\pgfqpoint{4.595474in}{6.569624in}}{\pgfqpoint{4.585638in}{6.565550in}}{\pgfqpoint{4.578387in}{6.558300in}}%
\pgfpathcurveto{\pgfqpoint{4.571137in}{6.551049in}}{\pgfqpoint{4.567063in}{6.541213in}}{\pgfqpoint{4.567063in}{6.530959in}}%
\pgfpathcurveto{\pgfqpoint{4.567063in}{6.520705in}}{\pgfqpoint{4.571137in}{6.510870in}}{\pgfqpoint{4.578387in}{6.503619in}}%
\pgfpathcurveto{\pgfqpoint{4.585638in}{6.496368in}}{\pgfqpoint{4.595474in}{6.492294in}}{\pgfqpoint{4.605728in}{6.492294in}}%
\pgfpathlineto{\pgfqpoint{4.605728in}{6.492294in}}%
\pgfpathclose%
\pgfusepath{stroke}%
\end{pgfscope}%
\begin{pgfscope}%
\pgfpathrectangle{\pgfqpoint{0.977909in}{2.928000in}}{\pgfqpoint{5.664000in}{5.664000in}}%
\pgfusepath{clip}%
\pgfsetbuttcap%
\pgfsetroundjoin%
\pgfsetlinewidth{1.003750pt}%
\definecolor{currentstroke}{rgb}{0.149020,0.274510,0.325490}%
\pgfsetstrokecolor{currentstroke}%
\pgfsetstrokeopacity{0.900000}%
\pgfsetdash{}{0pt}%
\pgfpathmoveto{\pgfqpoint{4.913455in}{6.697165in}}%
\pgfpathcurveto{\pgfqpoint{4.923709in}{6.697165in}}{\pgfqpoint{4.933545in}{6.701239in}}{\pgfqpoint{4.940796in}{6.708490in}}%
\pgfpathcurveto{\pgfqpoint{4.948046in}{6.715740in}}{\pgfqpoint{4.952120in}{6.725576in}}{\pgfqpoint{4.952120in}{6.735830in}}%
\pgfpathcurveto{\pgfqpoint{4.952120in}{6.746084in}}{\pgfqpoint{4.948046in}{6.755919in}}{\pgfqpoint{4.940796in}{6.763170in}}%
\pgfpathcurveto{\pgfqpoint{4.933545in}{6.770421in}}{\pgfqpoint{4.923709in}{6.774495in}}{\pgfqpoint{4.913455in}{6.774495in}}%
\pgfpathcurveto{\pgfqpoint{4.903201in}{6.774495in}}{\pgfqpoint{4.893366in}{6.770421in}}{\pgfqpoint{4.886115in}{6.763170in}}%
\pgfpathcurveto{\pgfqpoint{4.878864in}{6.755919in}}{\pgfqpoint{4.874790in}{6.746084in}}{\pgfqpoint{4.874790in}{6.735830in}}%
\pgfpathcurveto{\pgfqpoint{4.874790in}{6.725576in}}{\pgfqpoint{4.878864in}{6.715740in}}{\pgfqpoint{4.886115in}{6.708490in}}%
\pgfpathcurveto{\pgfqpoint{4.893366in}{6.701239in}}{\pgfqpoint{4.903201in}{6.697165in}}{\pgfqpoint{4.913455in}{6.697165in}}%
\pgfpathlineto{\pgfqpoint{4.913455in}{6.697165in}}%
\pgfpathclose%
\pgfusepath{stroke}%
\end{pgfscope}%
\begin{pgfscope}%
\pgfpathrectangle{\pgfqpoint{0.977909in}{2.928000in}}{\pgfqpoint{5.664000in}{5.664000in}}%
\pgfusepath{clip}%
\pgfsetbuttcap%
\pgfsetroundjoin%
\pgfsetlinewidth{1.003750pt}%
\definecolor{currentstroke}{rgb}{0.149020,0.274510,0.325490}%
\pgfsetstrokecolor{currentstroke}%
\pgfsetstrokeopacity{0.900000}%
\pgfsetdash{}{0pt}%
\pgfpathmoveto{\pgfqpoint{5.123027in}{6.927470in}}%
\pgfpathcurveto{\pgfqpoint{5.133281in}{6.927470in}}{\pgfqpoint{5.143117in}{6.931544in}}{\pgfqpoint{5.150367in}{6.938795in}}%
\pgfpathcurveto{\pgfqpoint{5.157618in}{6.946045in}}{\pgfqpoint{5.161692in}{6.955881in}}{\pgfqpoint{5.161692in}{6.966135in}}%
\pgfpathcurveto{\pgfqpoint{5.161692in}{6.976389in}}{\pgfqpoint{5.157618in}{6.986224in}}{\pgfqpoint{5.150367in}{6.993475in}}%
\pgfpathcurveto{\pgfqpoint{5.143117in}{7.000726in}}{\pgfqpoint{5.133281in}{7.004800in}}{\pgfqpoint{5.123027in}{7.004800in}}%
\pgfpathcurveto{\pgfqpoint{5.112773in}{7.004800in}}{\pgfqpoint{5.102937in}{7.000726in}}{\pgfqpoint{5.095687in}{6.993475in}}%
\pgfpathcurveto{\pgfqpoint{5.088436in}{6.986224in}}{\pgfqpoint{5.084362in}{6.976389in}}{\pgfqpoint{5.084362in}{6.966135in}}%
\pgfpathcurveto{\pgfqpoint{5.084362in}{6.955881in}}{\pgfqpoint{5.088436in}{6.946045in}}{\pgfqpoint{5.095687in}{6.938795in}}%
\pgfpathcurveto{\pgfqpoint{5.102937in}{6.931544in}}{\pgfqpoint{5.112773in}{6.927470in}}{\pgfqpoint{5.123027in}{6.927470in}}%
\pgfpathlineto{\pgfqpoint{5.123027in}{6.927470in}}%
\pgfpathclose%
\pgfusepath{stroke}%
\end{pgfscope}%
\begin{pgfscope}%
\pgfpathrectangle{\pgfqpoint{0.977909in}{2.928000in}}{\pgfqpoint{5.664000in}{5.664000in}}%
\pgfusepath{clip}%
\pgfsetbuttcap%
\pgfsetroundjoin%
\pgfsetlinewidth{1.003750pt}%
\definecolor{currentstroke}{rgb}{0.000000,0.619608,0.450980}%
\pgfsetstrokecolor{currentstroke}%
\pgfsetstrokeopacity{0.900000}%
\pgfsetdash{}{0pt}%
\pgfpathmoveto{\pgfqpoint{5.247216in}{6.754705in}}%
\pgfpathcurveto{\pgfqpoint{5.257470in}{6.754705in}}{\pgfqpoint{5.267306in}{6.758779in}}{\pgfqpoint{5.274557in}{6.766030in}}%
\pgfpathcurveto{\pgfqpoint{5.281807in}{6.773281in}}{\pgfqpoint{5.285881in}{6.783116in}}{\pgfqpoint{5.285881in}{6.793370in}}%
\pgfpathcurveto{\pgfqpoint{5.285881in}{6.803624in}}{\pgfqpoint{5.281807in}{6.813460in}}{\pgfqpoint{5.274557in}{6.820711in}}%
\pgfpathcurveto{\pgfqpoint{5.267306in}{6.827961in}}{\pgfqpoint{5.257470in}{6.832035in}}{\pgfqpoint{5.247216in}{6.832035in}}%
\pgfpathcurveto{\pgfqpoint{5.236962in}{6.832035in}}{\pgfqpoint{5.227127in}{6.827961in}}{\pgfqpoint{5.219876in}{6.820711in}}%
\pgfpathcurveto{\pgfqpoint{5.212625in}{6.813460in}}{\pgfqpoint{5.208551in}{6.803624in}}{\pgfqpoint{5.208551in}{6.793370in}}%
\pgfpathcurveto{\pgfqpoint{5.208551in}{6.783116in}}{\pgfqpoint{5.212625in}{6.773281in}}{\pgfqpoint{5.219876in}{6.766030in}}%
\pgfpathcurveto{\pgfqpoint{5.227127in}{6.758779in}}{\pgfqpoint{5.236962in}{6.754705in}}{\pgfqpoint{5.247216in}{6.754705in}}%
\pgfpathlineto{\pgfqpoint{5.247216in}{6.754705in}}%
\pgfpathclose%
\pgfusepath{stroke}%
\end{pgfscope}%
\begin{pgfscope}%
\pgfpathrectangle{\pgfqpoint{0.977909in}{2.928000in}}{\pgfqpoint{5.664000in}{5.664000in}}%
\pgfusepath{clip}%
\pgfsetbuttcap%
\pgfsetroundjoin%
\pgfsetlinewidth{1.003750pt}%
\definecolor{currentstroke}{rgb}{0.000000,0.619608,0.450980}%
\pgfsetstrokecolor{currentstroke}%
\pgfsetstrokeopacity{0.900000}%
\pgfsetdash{}{0pt}%
\pgfpathmoveto{\pgfqpoint{5.080940in}{7.129300in}}%
\pgfpathcurveto{\pgfqpoint{5.091194in}{7.129300in}}{\pgfqpoint{5.101030in}{7.133374in}}{\pgfqpoint{5.108280in}{7.140625in}}%
\pgfpathcurveto{\pgfqpoint{5.115531in}{7.147876in}}{\pgfqpoint{5.119605in}{7.157711in}}{\pgfqpoint{5.119605in}{7.167965in}}%
\pgfpathcurveto{\pgfqpoint{5.119605in}{7.178220in}}{\pgfqpoint{5.115531in}{7.188055in}}{\pgfqpoint{5.108280in}{7.195306in}}%
\pgfpathcurveto{\pgfqpoint{5.101030in}{7.202557in}}{\pgfqpoint{5.091194in}{7.206631in}}{\pgfqpoint{5.080940in}{7.206631in}}%
\pgfpathcurveto{\pgfqpoint{5.070686in}{7.206631in}}{\pgfqpoint{5.060850in}{7.202557in}}{\pgfqpoint{5.053600in}{7.195306in}}%
\pgfpathcurveto{\pgfqpoint{5.046349in}{7.188055in}}{\pgfqpoint{5.042275in}{7.178220in}}{\pgfqpoint{5.042275in}{7.167965in}}%
\pgfpathcurveto{\pgfqpoint{5.042275in}{7.157711in}}{\pgfqpoint{5.046349in}{7.147876in}}{\pgfqpoint{5.053600in}{7.140625in}}%
\pgfpathcurveto{\pgfqpoint{5.060850in}{7.133374in}}{\pgfqpoint{5.070686in}{7.129300in}}{\pgfqpoint{5.080940in}{7.129300in}}%
\pgfpathlineto{\pgfqpoint{5.080940in}{7.129300in}}%
\pgfpathclose%
\pgfusepath{stroke}%
\end{pgfscope}%
\begin{pgfscope}%
\pgfpathrectangle{\pgfqpoint{0.977909in}{2.928000in}}{\pgfqpoint{5.664000in}{5.664000in}}%
\pgfusepath{clip}%
\pgfsetbuttcap%
\pgfsetroundjoin%
\pgfsetlinewidth{1.003750pt}%
\definecolor{currentstroke}{rgb}{0.478431,0.243137,0.615686}%
\pgfsetstrokecolor{currentstroke}%
\pgfsetstrokeopacity{0.900000}%
\pgfsetdash{}{0pt}%
\pgfpathmoveto{\pgfqpoint{4.513394in}{5.172227in}}%
\pgfpathcurveto{\pgfqpoint{4.523648in}{5.172227in}}{\pgfqpoint{4.533483in}{5.176301in}}{\pgfqpoint{4.540734in}{5.183552in}}%
\pgfpathcurveto{\pgfqpoint{4.547985in}{5.190803in}}{\pgfqpoint{4.552059in}{5.200638in}}{\pgfqpoint{4.552059in}{5.210892in}}%
\pgfpathcurveto{\pgfqpoint{4.552059in}{5.221146in}}{\pgfqpoint{4.547985in}{5.230982in}}{\pgfqpoint{4.540734in}{5.238233in}}%
\pgfpathcurveto{\pgfqpoint{4.533483in}{5.245483in}}{\pgfqpoint{4.523648in}{5.249557in}}{\pgfqpoint{4.513394in}{5.249557in}}%
\pgfpathcurveto{\pgfqpoint{4.503140in}{5.249557in}}{\pgfqpoint{4.493304in}{5.245483in}}{\pgfqpoint{4.486053in}{5.238233in}}%
\pgfpathcurveto{\pgfqpoint{4.478803in}{5.230982in}}{\pgfqpoint{4.474729in}{5.221146in}}{\pgfqpoint{4.474729in}{5.210892in}}%
\pgfpathcurveto{\pgfqpoint{4.474729in}{5.200638in}}{\pgfqpoint{4.478803in}{5.190803in}}{\pgfqpoint{4.486053in}{5.183552in}}%
\pgfpathcurveto{\pgfqpoint{4.493304in}{5.176301in}}{\pgfqpoint{4.503140in}{5.172227in}}{\pgfqpoint{4.513394in}{5.172227in}}%
\pgfpathlineto{\pgfqpoint{4.513394in}{5.172227in}}%
\pgfpathclose%
\pgfusepath{stroke}%
\end{pgfscope}%
\begin{pgfscope}%
\pgfpathrectangle{\pgfqpoint{0.977909in}{2.928000in}}{\pgfqpoint{5.664000in}{5.664000in}}%
\pgfusepath{clip}%
\pgfsetbuttcap%
\pgfsetroundjoin%
\pgfsetlinewidth{1.003750pt}%
\definecolor{currentstroke}{rgb}{0.478431,0.243137,0.615686}%
\pgfsetstrokecolor{currentstroke}%
\pgfsetstrokeopacity{0.900000}%
\pgfsetdash{}{0pt}%
\pgfpathmoveto{\pgfqpoint{5.175638in}{6.113717in}}%
\pgfpathcurveto{\pgfqpoint{5.185892in}{6.113717in}}{\pgfqpoint{5.195728in}{6.117791in}}{\pgfqpoint{5.202979in}{6.125042in}}%
\pgfpathcurveto{\pgfqpoint{5.210229in}{6.132293in}}{\pgfqpoint{5.214303in}{6.142128in}}{\pgfqpoint{5.214303in}{6.152382in}}%
\pgfpathcurveto{\pgfqpoint{5.214303in}{6.162636in}}{\pgfqpoint{5.210229in}{6.172472in}}{\pgfqpoint{5.202979in}{6.179723in}}%
\pgfpathcurveto{\pgfqpoint{5.195728in}{6.186973in}}{\pgfqpoint{5.185892in}{6.191047in}}{\pgfqpoint{5.175638in}{6.191047in}}%
\pgfpathcurveto{\pgfqpoint{5.165384in}{6.191047in}}{\pgfqpoint{5.155549in}{6.186973in}}{\pgfqpoint{5.148298in}{6.179723in}}%
\pgfpathcurveto{\pgfqpoint{5.141047in}{6.172472in}}{\pgfqpoint{5.136973in}{6.162636in}}{\pgfqpoint{5.136973in}{6.152382in}}%
\pgfpathcurveto{\pgfqpoint{5.136973in}{6.142128in}}{\pgfqpoint{5.141047in}{6.132293in}}{\pgfqpoint{5.148298in}{6.125042in}}%
\pgfpathcurveto{\pgfqpoint{5.155549in}{6.117791in}}{\pgfqpoint{5.165384in}{6.113717in}}{\pgfqpoint{5.175638in}{6.113717in}}%
\pgfpathlineto{\pgfqpoint{5.175638in}{6.113717in}}%
\pgfpathclose%
\pgfusepath{stroke}%
\end{pgfscope}%
\begin{pgfscope}%
\pgfpathrectangle{\pgfqpoint{0.977909in}{2.928000in}}{\pgfqpoint{5.664000in}{5.664000in}}%
\pgfusepath{clip}%
\pgfsetbuttcap%
\pgfsetroundjoin%
\pgfsetlinewidth{1.003750pt}%
\definecolor{currentstroke}{rgb}{0.549020,0.337255,0.294118}%
\pgfsetstrokecolor{currentstroke}%
\pgfsetstrokeopacity{0.900000}%
\pgfsetdash{}{0pt}%
\pgfpathmoveto{\pgfqpoint{3.989524in}{4.514553in}}%
\pgfpathcurveto{\pgfqpoint{3.999778in}{4.514553in}}{\pgfqpoint{4.009614in}{4.518627in}}{\pgfqpoint{4.016864in}{4.525878in}}%
\pgfpathcurveto{\pgfqpoint{4.024115in}{4.533129in}}{\pgfqpoint{4.028189in}{4.542964in}}{\pgfqpoint{4.028189in}{4.553218in}}%
\pgfpathcurveto{\pgfqpoint{4.028189in}{4.563472in}}{\pgfqpoint{4.024115in}{4.573308in}}{\pgfqpoint{4.016864in}{4.580559in}}%
\pgfpathcurveto{\pgfqpoint{4.009614in}{4.587809in}}{\pgfqpoint{3.999778in}{4.591883in}}{\pgfqpoint{3.989524in}{4.591883in}}%
\pgfpathcurveto{\pgfqpoint{3.979270in}{4.591883in}}{\pgfqpoint{3.969434in}{4.587809in}}{\pgfqpoint{3.962184in}{4.580559in}}%
\pgfpathcurveto{\pgfqpoint{3.954933in}{4.573308in}}{\pgfqpoint{3.950859in}{4.563472in}}{\pgfqpoint{3.950859in}{4.553218in}}%
\pgfpathcurveto{\pgfqpoint{3.950859in}{4.542964in}}{\pgfqpoint{3.954933in}{4.533129in}}{\pgfqpoint{3.962184in}{4.525878in}}%
\pgfpathcurveto{\pgfqpoint{3.969434in}{4.518627in}}{\pgfqpoint{3.979270in}{4.514553in}}{\pgfqpoint{3.989524in}{4.514553in}}%
\pgfpathlineto{\pgfqpoint{3.989524in}{4.514553in}}%
\pgfpathclose%
\pgfusepath{stroke}%
\end{pgfscope}%
\begin{pgfscope}%
\pgfsetbuttcap%
\pgfsetroundjoin%
\definecolor{currentfill}{rgb}{0.352941,0.407843,0.450980}%
\pgfsetfillcolor{currentfill}%
\pgfsetlinewidth{0.000000pt}%
\definecolor{currentstroke}{rgb}{0.352941,0.407843,0.450980}%
\pgfsetstrokecolor{currentstroke}%
\pgfsetdash{}{0pt}%
\pgfpathmoveto{\pgfqpoint{1.003226in}{8.744547in}}%
\pgfpathlineto{\pgfqpoint{1.048644in}{8.744547in}}%
\pgfpathlineto{\pgfqpoint{1.048644in}{8.733600in}}%
\pgfpathlineto{\pgfqpoint{0.987578in}{8.733600in}}%
\pgfpathlineto{\pgfqpoint{0.987578in}{8.744547in}}%
\pgfpathquadraticcurveto{\pgfqpoint{0.994979in}{8.752217in}}{\pgfqpoint{1.007762in}{8.765122in}}%
\pgfpathquadraticcurveto{\pgfqpoint{1.020564in}{8.778049in}}{\pgfqpoint{1.023842in}{8.781801in}}%
\pgfpathquadraticcurveto{\pgfqpoint{1.030089in}{8.788811in}}{\pgfqpoint{1.032563in}{8.793676in}}%
\pgfpathquadraticcurveto{\pgfqpoint{1.035058in}{8.798541in}}{\pgfqpoint{1.035058in}{8.803242in}}%
\pgfpathquadraticcurveto{\pgfqpoint{1.035058in}{8.810911in}}{\pgfqpoint{1.029677in}{8.815735in}}%
\pgfpathquadraticcurveto{\pgfqpoint{1.024296in}{8.820580in}}{\pgfqpoint{1.015658in}{8.820580in}}%
\pgfpathquadraticcurveto{\pgfqpoint{1.009535in}{8.820580in}}{\pgfqpoint{1.002731in}{8.818457in}}%
\pgfpathquadraticcurveto{\pgfqpoint{0.995948in}{8.816333in}}{\pgfqpoint{0.988217in}{8.812004in}}%
\pgfpathlineto{\pgfqpoint{0.988217in}{8.825157in}}%
\pgfpathquadraticcurveto{\pgfqpoint{0.996072in}{8.828311in}}{\pgfqpoint{1.002896in}{8.829919in}}%
\pgfpathquadraticcurveto{\pgfqpoint{1.009741in}{8.831528in}}{\pgfqpoint{1.015410in}{8.831528in}}%
\pgfpathquadraticcurveto{\pgfqpoint{1.030357in}{8.831528in}}{\pgfqpoint{1.039243in}{8.824044in}}%
\pgfpathquadraticcurveto{\pgfqpoint{1.048128in}{8.816581in}}{\pgfqpoint{1.048128in}{8.804087in}}%
\pgfpathquadraticcurveto{\pgfqpoint{1.048128in}{8.798150in}}{\pgfqpoint{1.045902in}{8.792831in}}%
\pgfpathquadraticcurveto{\pgfqpoint{1.043696in}{8.787532in}}{\pgfqpoint{1.037820in}{8.780317in}}%
\pgfpathquadraticcurveto{\pgfqpoint{1.036212in}{8.778440in}}{\pgfqpoint{1.027574in}{8.769514in}}%
\pgfpathquadraticcurveto{\pgfqpoint{1.018956in}{8.760587in}}{\pgfqpoint{1.003226in}{8.744547in}}%
\pgfpathlineto{\pgfqpoint{1.003226in}{8.744547in}}%
\pgfpathclose%
\pgfpathmoveto{\pgfqpoint{1.072682in}{8.829796in}}%
\pgfpathlineto{\pgfqpoint{1.134531in}{8.829796in}}%
\pgfpathlineto{\pgfqpoint{1.134531in}{8.824250in}}%
\pgfpathlineto{\pgfqpoint{1.099607in}{8.733600in}}%
\pgfpathlineto{\pgfqpoint{1.086021in}{8.733600in}}%
\pgfpathlineto{\pgfqpoint{1.118883in}{8.818828in}}%
\pgfpathlineto{\pgfqpoint{1.072682in}{8.818828in}}%
\pgfpathlineto{\pgfqpoint{1.072682in}{8.829796in}}%
\pgfpathlineto{\pgfqpoint{1.072682in}{8.829796in}}%
\pgfpathclose%
\pgfpathmoveto{\pgfqpoint{1.199349in}{8.785471in}}%
\pgfpathquadraticcurveto{\pgfqpoint{1.208688in}{8.783471in}}{\pgfqpoint{1.213925in}{8.777142in}}%
\pgfpathquadraticcurveto{\pgfqpoint{1.219182in}{8.770833in}}{\pgfqpoint{1.219182in}{8.761556in}}%
\pgfpathquadraticcurveto{\pgfqpoint{1.219182in}{8.747330in}}{\pgfqpoint{1.209389in}{8.739517in}}%
\pgfpathquadraticcurveto{\pgfqpoint{1.199596in}{8.731724in}}{\pgfqpoint{1.181557in}{8.731724in}}%
\pgfpathquadraticcurveto{\pgfqpoint{1.175517in}{8.731724in}}{\pgfqpoint{1.169105in}{8.732920in}}%
\pgfpathquadraticcurveto{\pgfqpoint{1.162693in}{8.734115in}}{\pgfqpoint{1.155869in}{8.736507in}}%
\pgfpathlineto{\pgfqpoint{1.155869in}{8.749062in}}%
\pgfpathquadraticcurveto{\pgfqpoint{1.161271in}{8.745908in}}{\pgfqpoint{1.167703in}{8.744300in}}%
\pgfpathquadraticcurveto{\pgfqpoint{1.174156in}{8.742692in}}{\pgfqpoint{1.181186in}{8.742692in}}%
\pgfpathquadraticcurveto{\pgfqpoint{1.193411in}{8.742692in}}{\pgfqpoint{1.199823in}{8.747516in}}%
\pgfpathquadraticcurveto{\pgfqpoint{1.206235in}{8.752340in}}{\pgfqpoint{1.206235in}{8.761556in}}%
\pgfpathquadraticcurveto{\pgfqpoint{1.206235in}{8.770070in}}{\pgfqpoint{1.200277in}{8.774853in}}%
\pgfpathquadraticcurveto{\pgfqpoint{1.194319in}{8.779657in}}{\pgfqpoint{1.183701in}{8.779657in}}%
\pgfpathlineto{\pgfqpoint{1.172486in}{8.779657in}}%
\pgfpathlineto{\pgfqpoint{1.172486in}{8.790357in}}%
\pgfpathlineto{\pgfqpoint{1.184217in}{8.790357in}}%
\pgfpathquadraticcurveto{\pgfqpoint{1.193803in}{8.790357in}}{\pgfqpoint{1.198895in}{8.794191in}}%
\pgfpathquadraticcurveto{\pgfqpoint{1.203988in}{8.798026in}}{\pgfqpoint{1.203988in}{8.805242in}}%
\pgfpathquadraticcurveto{\pgfqpoint{1.203988in}{8.812643in}}{\pgfqpoint{1.198730in}{8.816601in}}%
\pgfpathquadraticcurveto{\pgfqpoint{1.193494in}{8.820580in}}{\pgfqpoint{1.183701in}{8.820580in}}%
\pgfpathquadraticcurveto{\pgfqpoint{1.178341in}{8.820580in}}{\pgfqpoint{1.172218in}{8.819405in}}%
\pgfpathquadraticcurveto{\pgfqpoint{1.166095in}{8.818251in}}{\pgfqpoint{1.158755in}{8.815818in}}%
\pgfpathlineto{\pgfqpoint{1.158755in}{8.827404in}}%
\pgfpathquadraticcurveto{\pgfqpoint{1.166177in}{8.829466in}}{\pgfqpoint{1.172651in}{8.830497in}}%
\pgfpathquadraticcurveto{\pgfqpoint{1.179124in}{8.831528in}}{\pgfqpoint{1.184856in}{8.831528in}}%
\pgfpathquadraticcurveto{\pgfqpoint{1.199679in}{8.831528in}}{\pgfqpoint{1.208296in}{8.824786in}}%
\pgfpathquadraticcurveto{\pgfqpoint{1.216935in}{8.818065in}}{\pgfqpoint{1.216935in}{8.806602in}}%
\pgfpathquadraticcurveto{\pgfqpoint{1.216935in}{8.798603in}}{\pgfqpoint{1.212358in}{8.793099in}}%
\pgfpathquadraticcurveto{\pgfqpoint{1.207781in}{8.787594in}}{\pgfqpoint{1.199349in}{8.785471in}}%
\pgfpathlineto{\pgfqpoint{1.199349in}{8.785471in}}%
\pgfpathclose%
\pgfpathmoveto{\pgfqpoint{1.325933in}{8.794686in}}%
\pgfpathquadraticcurveto{\pgfqpoint{1.323933in}{8.795841in}}{\pgfqpoint{1.321583in}{8.796377in}}%
\pgfpathquadraticcurveto{\pgfqpoint{1.319233in}{8.796933in}}{\pgfqpoint{1.316408in}{8.796933in}}%
\pgfpathquadraticcurveto{\pgfqpoint{1.306348in}{8.796933in}}{\pgfqpoint{1.300967in}{8.790398in}}%
\pgfpathquadraticcurveto{\pgfqpoint{1.295586in}{8.783863in}}{\pgfqpoint{1.295586in}{8.771616in}}%
\pgfpathlineto{\pgfqpoint{1.295586in}{8.733600in}}%
\pgfpathlineto{\pgfqpoint{1.283670in}{8.733600in}}%
\pgfpathlineto{\pgfqpoint{1.283670in}{8.805757in}}%
\pgfpathlineto{\pgfqpoint{1.295586in}{8.805757in}}%
\pgfpathlineto{\pgfqpoint{1.295586in}{8.794542in}}%
\pgfpathquadraticcurveto{\pgfqpoint{1.299338in}{8.801118in}}{\pgfqpoint{1.305317in}{8.804293in}}%
\pgfpathquadraticcurveto{\pgfqpoint{1.311316in}{8.807489in}}{\pgfqpoint{1.319893in}{8.807489in}}%
\pgfpathquadraticcurveto{\pgfqpoint{1.321109in}{8.807489in}}{\pgfqpoint{1.322593in}{8.807324in}}%
\pgfpathquadraticcurveto{\pgfqpoint{1.324078in}{8.807180in}}{\pgfqpoint{1.325871in}{8.806850in}}%
\pgfpathlineto{\pgfqpoint{1.325933in}{8.794686in}}%
\pgfpathlineto{\pgfqpoint{1.325933in}{8.794686in}}%
\pgfpathclose%
\pgfpathmoveto{\pgfqpoint{1.397183in}{8.772647in}}%
\pgfpathlineto{\pgfqpoint{1.397183in}{8.766854in}}%
\pgfpathlineto{\pgfqpoint{1.342674in}{8.766854in}}%
\pgfpathquadraticcurveto{\pgfqpoint{1.343457in}{8.754608in}}{\pgfqpoint{1.350054in}{8.748196in}}%
\pgfpathquadraticcurveto{\pgfqpoint{1.356651in}{8.741785in}}{\pgfqpoint{1.368444in}{8.741785in}}%
\pgfpathquadraticcurveto{\pgfqpoint{1.375268in}{8.741785in}}{\pgfqpoint{1.381680in}{8.743455in}}%
\pgfpathquadraticcurveto{\pgfqpoint{1.388091in}{8.745125in}}{\pgfqpoint{1.394421in}{8.748485in}}%
\pgfpathlineto{\pgfqpoint{1.394421in}{8.737270in}}%
\pgfpathquadraticcurveto{\pgfqpoint{1.388030in}{8.734569in}}{\pgfqpoint{1.381329in}{8.733146in}}%
\pgfpathquadraticcurveto{\pgfqpoint{1.374629in}{8.731724in}}{\pgfqpoint{1.367743in}{8.731724in}}%
\pgfpathquadraticcurveto{\pgfqpoint{1.350467in}{8.731724in}}{\pgfqpoint{1.340385in}{8.741764in}}%
\pgfpathquadraticcurveto{\pgfqpoint{1.330304in}{8.751825in}}{\pgfqpoint{1.330304in}{8.768978in}}%
\pgfpathquadraticcurveto{\pgfqpoint{1.330304in}{8.786687in}}{\pgfqpoint{1.339870in}{8.797078in}}%
\pgfpathquadraticcurveto{\pgfqpoint{1.349436in}{8.807489in}}{\pgfqpoint{1.365681in}{8.807489in}}%
\pgfpathquadraticcurveto{\pgfqpoint{1.380237in}{8.807489in}}{\pgfqpoint{1.388710in}{8.798108in}}%
\pgfpathquadraticcurveto{\pgfqpoint{1.397183in}{8.788749in}}{\pgfqpoint{1.397183in}{8.772647in}}%
\pgfpathlineto{\pgfqpoint{1.397183in}{8.772647in}}%
\pgfpathclose%
\pgfpathmoveto{\pgfqpoint{1.385329in}{8.776131in}}%
\pgfpathquadraticcurveto{\pgfqpoint{1.385205in}{8.785842in}}{\pgfqpoint{1.379886in}{8.791635in}}%
\pgfpathquadraticcurveto{\pgfqpoint{1.374567in}{8.797449in}}{\pgfqpoint{1.365805in}{8.797449in}}%
\pgfpathquadraticcurveto{\pgfqpoint{1.355889in}{8.797449in}}{\pgfqpoint{1.349931in}{8.791841in}}%
\pgfpathquadraticcurveto{\pgfqpoint{1.343972in}{8.786233in}}{\pgfqpoint{1.343065in}{8.776049in}}%
\pgfpathlineto{\pgfqpoint{1.385329in}{8.776131in}}%
\pgfpathlineto{\pgfqpoint{1.385329in}{8.776131in}}%
\pgfpathclose%
\pgfpathmoveto{\pgfqpoint{1.428376in}{8.826250in}}%
\pgfpathlineto{\pgfqpoint{1.428376in}{8.805757in}}%
\pgfpathlineto{\pgfqpoint{1.452785in}{8.805757in}}%
\pgfpathlineto{\pgfqpoint{1.452785in}{8.796542in}}%
\pgfpathlineto{\pgfqpoint{1.428376in}{8.796542in}}%
\pgfpathlineto{\pgfqpoint{1.428376in}{8.757371in}}%
\pgfpathquadraticcurveto{\pgfqpoint{1.428376in}{8.748547in}}{\pgfqpoint{1.430788in}{8.746032in}}%
\pgfpathquadraticcurveto{\pgfqpoint{1.433200in}{8.743516in}}{\pgfqpoint{1.440622in}{8.743516in}}%
\pgfpathlineto{\pgfqpoint{1.452785in}{8.743516in}}%
\pgfpathlineto{\pgfqpoint{1.452785in}{8.733600in}}%
\pgfpathlineto{\pgfqpoint{1.440622in}{8.733600in}}%
\pgfpathquadraticcurveto{\pgfqpoint{1.426891in}{8.733600in}}{\pgfqpoint{1.421675in}{8.738713in}}%
\pgfpathquadraticcurveto{\pgfqpoint{1.416459in}{8.743846in}}{\pgfqpoint{1.416459in}{8.757371in}}%
\pgfpathlineto{\pgfqpoint{1.416459in}{8.796542in}}%
\pgfpathlineto{\pgfqpoint{1.407759in}{8.796542in}}%
\pgfpathlineto{\pgfqpoint{1.407759in}{8.805757in}}%
\pgfpathlineto{\pgfqpoint{1.416459in}{8.805757in}}%
\pgfpathlineto{\pgfqpoint{1.416459in}{8.826250in}}%
\pgfpathlineto{\pgfqpoint{1.428376in}{8.826250in}}%
\pgfpathlineto{\pgfqpoint{1.428376in}{8.826250in}}%
\pgfpathclose%
\pgfpathmoveto{\pgfqpoint{1.501172in}{8.769864in}}%
\pgfpathquadraticcurveto{\pgfqpoint{1.486802in}{8.769864in}}{\pgfqpoint{1.481257in}{8.766586in}}%
\pgfpathquadraticcurveto{\pgfqpoint{1.475711in}{8.763308in}}{\pgfqpoint{1.475711in}{8.755371in}}%
\pgfpathquadraticcurveto{\pgfqpoint{1.475711in}{8.749062in}}{\pgfqpoint{1.479875in}{8.745351in}}%
\pgfpathquadraticcurveto{\pgfqpoint{1.484040in}{8.741661in}}{\pgfqpoint{1.491173in}{8.741661in}}%
\pgfpathquadraticcurveto{\pgfqpoint{1.501048in}{8.741661in}}{\pgfqpoint{1.507006in}{8.748650in}}%
\pgfpathquadraticcurveto{\pgfqpoint{1.512964in}{8.755639in}}{\pgfqpoint{1.512964in}{8.767225in}}%
\pgfpathlineto{\pgfqpoint{1.512964in}{8.769864in}}%
\pgfpathlineto{\pgfqpoint{1.501172in}{8.769864in}}%
\pgfpathlineto{\pgfqpoint{1.501172in}{8.769864in}}%
\pgfpathclose%
\pgfpathmoveto{\pgfqpoint{1.524819in}{8.774771in}}%
\pgfpathlineto{\pgfqpoint{1.524819in}{8.733600in}}%
\pgfpathlineto{\pgfqpoint{1.512964in}{8.733600in}}%
\pgfpathlineto{\pgfqpoint{1.512964in}{8.744547in}}%
\pgfpathquadraticcurveto{\pgfqpoint{1.508903in}{8.737991in}}{\pgfqpoint{1.502842in}{8.734858in}}%
\pgfpathquadraticcurveto{\pgfqpoint{1.496781in}{8.731724in}}{\pgfqpoint{1.488019in}{8.731724in}}%
\pgfpathquadraticcurveto{\pgfqpoint{1.476948in}{8.731724in}}{\pgfqpoint{1.470392in}{8.737950in}}%
\pgfpathquadraticcurveto{\pgfqpoint{1.463856in}{8.744176in}}{\pgfqpoint{1.463856in}{8.754608in}}%
\pgfpathquadraticcurveto{\pgfqpoint{1.463856in}{8.766772in}}{\pgfqpoint{1.472000in}{8.772957in}}%
\pgfpathquadraticcurveto{\pgfqpoint{1.480164in}{8.779141in}}{\pgfqpoint{1.496327in}{8.779141in}}%
\pgfpathlineto{\pgfqpoint{1.512964in}{8.779141in}}%
\pgfpathlineto{\pgfqpoint{1.512964in}{8.780317in}}%
\pgfpathquadraticcurveto{\pgfqpoint{1.512964in}{8.788501in}}{\pgfqpoint{1.507584in}{8.792975in}}%
\pgfpathquadraticcurveto{\pgfqpoint{1.502203in}{8.797449in}}{\pgfqpoint{1.492472in}{8.797449in}}%
\pgfpathquadraticcurveto{\pgfqpoint{1.486287in}{8.797449in}}{\pgfqpoint{1.480411in}{8.795964in}}%
\pgfpathquadraticcurveto{\pgfqpoint{1.474556in}{8.794480in}}{\pgfqpoint{1.469155in}{8.791511in}}%
\pgfpathlineto{\pgfqpoint{1.469155in}{8.802479in}}%
\pgfpathquadraticcurveto{\pgfqpoint{1.475649in}{8.804994in}}{\pgfqpoint{1.481772in}{8.806231in}}%
\pgfpathquadraticcurveto{\pgfqpoint{1.487895in}{8.807489in}}{\pgfqpoint{1.493688in}{8.807489in}}%
\pgfpathquadraticcurveto{\pgfqpoint{1.509357in}{8.807489in}}{\pgfqpoint{1.517088in}{8.799366in}}%
\pgfpathquadraticcurveto{\pgfqpoint{1.524819in}{8.791264in}}{\pgfqpoint{1.524819in}{8.774771in}}%
\pgfpathlineto{\pgfqpoint{1.524819in}{8.774771in}}%
\pgfpathclose%
\pgfpathmoveto{\pgfqpoint{1.549229in}{8.805757in}}%
\pgfpathlineto{\pgfqpoint{1.561083in}{8.805757in}}%
\pgfpathlineto{\pgfqpoint{1.561083in}{8.733600in}}%
\pgfpathlineto{\pgfqpoint{1.549229in}{8.733600in}}%
\pgfpathlineto{\pgfqpoint{1.549229in}{8.805757in}}%
\pgfpathlineto{\pgfqpoint{1.549229in}{8.805757in}}%
\pgfpathclose%
\pgfpathmoveto{\pgfqpoint{1.549229in}{8.833857in}}%
\pgfpathlineto{\pgfqpoint{1.561083in}{8.833857in}}%
\pgfpathlineto{\pgfqpoint{1.561083in}{8.818828in}}%
\pgfpathlineto{\pgfqpoint{1.549229in}{8.818828in}}%
\pgfpathlineto{\pgfqpoint{1.549229in}{8.833857in}}%
\pgfpathlineto{\pgfqpoint{1.549229in}{8.833857in}}%
\pgfpathclose%
\pgfpathmoveto{\pgfqpoint{1.645878in}{8.777162in}}%
\pgfpathlineto{\pgfqpoint{1.645878in}{8.733600in}}%
\pgfpathlineto{\pgfqpoint{1.634023in}{8.733600in}}%
\pgfpathlineto{\pgfqpoint{1.634023in}{8.776771in}}%
\pgfpathquadraticcurveto{\pgfqpoint{1.634023in}{8.787017in}}{\pgfqpoint{1.630024in}{8.792088in}}%
\pgfpathquadraticcurveto{\pgfqpoint{1.626024in}{8.797181in}}{\pgfqpoint{1.618046in}{8.797181in}}%
\pgfpathquadraticcurveto{\pgfqpoint{1.608439in}{8.797181in}}{\pgfqpoint{1.602893in}{8.791058in}}%
\pgfpathquadraticcurveto{\pgfqpoint{1.597347in}{8.784955in}}{\pgfqpoint{1.597347in}{8.774379in}}%
\pgfpathlineto{\pgfqpoint{1.597347in}{8.733600in}}%
\pgfpathlineto{\pgfqpoint{1.585431in}{8.733600in}}%
\pgfpathlineto{\pgfqpoint{1.585431in}{8.805757in}}%
\pgfpathlineto{\pgfqpoint{1.597347in}{8.805757in}}%
\pgfpathlineto{\pgfqpoint{1.597347in}{8.794542in}}%
\pgfpathquadraticcurveto{\pgfqpoint{1.601615in}{8.801057in}}{\pgfqpoint{1.607367in}{8.804273in}}%
\pgfpathquadraticcurveto{\pgfqpoint{1.613139in}{8.807489in}}{\pgfqpoint{1.620685in}{8.807489in}}%
\pgfpathquadraticcurveto{\pgfqpoint{1.633116in}{8.807489in}}{\pgfqpoint{1.639487in}{8.799799in}}%
\pgfpathquadraticcurveto{\pgfqpoint{1.645878in}{8.792109in}}{\pgfqpoint{1.645878in}{8.777162in}}%
\pgfpathlineto{\pgfqpoint{1.645878in}{8.777162in}}%
\pgfpathclose%
\pgfpathmoveto{\pgfqpoint{1.731229in}{8.772647in}}%
\pgfpathlineto{\pgfqpoint{1.731229in}{8.766854in}}%
\pgfpathlineto{\pgfqpoint{1.676720in}{8.766854in}}%
\pgfpathquadraticcurveto{\pgfqpoint{1.677503in}{8.754608in}}{\pgfqpoint{1.684100in}{8.748196in}}%
\pgfpathquadraticcurveto{\pgfqpoint{1.690698in}{8.741785in}}{\pgfqpoint{1.702490in}{8.741785in}}%
\pgfpathquadraticcurveto{\pgfqpoint{1.709314in}{8.741785in}}{\pgfqpoint{1.715726in}{8.743455in}}%
\pgfpathquadraticcurveto{\pgfqpoint{1.722138in}{8.745125in}}{\pgfqpoint{1.728467in}{8.748485in}}%
\pgfpathlineto{\pgfqpoint{1.728467in}{8.737270in}}%
\pgfpathquadraticcurveto{\pgfqpoint{1.722076in}{8.734569in}}{\pgfqpoint{1.715375in}{8.733146in}}%
\pgfpathquadraticcurveto{\pgfqpoint{1.708675in}{8.731724in}}{\pgfqpoint{1.701789in}{8.731724in}}%
\pgfpathquadraticcurveto{\pgfqpoint{1.684513in}{8.731724in}}{\pgfqpoint{1.674431in}{8.741764in}}%
\pgfpathquadraticcurveto{\pgfqpoint{1.664350in}{8.751825in}}{\pgfqpoint{1.664350in}{8.768978in}}%
\pgfpathquadraticcurveto{\pgfqpoint{1.664350in}{8.786687in}}{\pgfqpoint{1.673916in}{8.797078in}}%
\pgfpathquadraticcurveto{\pgfqpoint{1.683482in}{8.807489in}}{\pgfqpoint{1.699728in}{8.807489in}}%
\pgfpathquadraticcurveto{\pgfqpoint{1.714283in}{8.807489in}}{\pgfqpoint{1.722756in}{8.798108in}}%
\pgfpathquadraticcurveto{\pgfqpoint{1.731229in}{8.788749in}}{\pgfqpoint{1.731229in}{8.772647in}}%
\pgfpathlineto{\pgfqpoint{1.731229in}{8.772647in}}%
\pgfpathclose%
\pgfpathmoveto{\pgfqpoint{1.719375in}{8.776131in}}%
\pgfpathquadraticcurveto{\pgfqpoint{1.719251in}{8.785842in}}{\pgfqpoint{1.713932in}{8.791635in}}%
\pgfpathquadraticcurveto{\pgfqpoint{1.708613in}{8.797449in}}{\pgfqpoint{1.699851in}{8.797449in}}%
\pgfpathquadraticcurveto{\pgfqpoint{1.689935in}{8.797449in}}{\pgfqpoint{1.683977in}{8.791841in}}%
\pgfpathquadraticcurveto{\pgfqpoint{1.678019in}{8.786233in}}{\pgfqpoint{1.677112in}{8.776049in}}%
\pgfpathlineto{\pgfqpoint{1.719375in}{8.776131in}}%
\pgfpathlineto{\pgfqpoint{1.719375in}{8.776131in}}%
\pgfpathclose%
\pgfpathmoveto{\pgfqpoint{1.798171in}{8.794810in}}%
\pgfpathlineto{\pgfqpoint{1.798171in}{8.833857in}}%
\pgfpathlineto{\pgfqpoint{1.810025in}{8.833857in}}%
\pgfpathlineto{\pgfqpoint{1.810025in}{8.733600in}}%
\pgfpathlineto{\pgfqpoint{1.798171in}{8.733600in}}%
\pgfpathlineto{\pgfqpoint{1.798171in}{8.744424in}}%
\pgfpathquadraticcurveto{\pgfqpoint{1.794439in}{8.737991in}}{\pgfqpoint{1.788728in}{8.734858in}}%
\pgfpathquadraticcurveto{\pgfqpoint{1.783038in}{8.731724in}}{\pgfqpoint{1.775039in}{8.731724in}}%
\pgfpathquadraticcurveto{\pgfqpoint{1.761968in}{8.731724in}}{\pgfqpoint{1.753742in}{8.742156in}}%
\pgfpathquadraticcurveto{\pgfqpoint{1.745537in}{8.752608in}}{\pgfqpoint{1.745537in}{8.769617in}}%
\pgfpathquadraticcurveto{\pgfqpoint{1.745537in}{8.786625in}}{\pgfqpoint{1.753742in}{8.797057in}}%
\pgfpathquadraticcurveto{\pgfqpoint{1.761968in}{8.807489in}}{\pgfqpoint{1.775039in}{8.807489in}}%
\pgfpathquadraticcurveto{\pgfqpoint{1.783038in}{8.807489in}}{\pgfqpoint{1.788728in}{8.804355in}}%
\pgfpathquadraticcurveto{\pgfqpoint{1.794439in}{8.801242in}}{\pgfqpoint{1.798171in}{8.794810in}}%
\pgfpathlineto{\pgfqpoint{1.798171in}{8.794810in}}%
\pgfpathclose%
\pgfpathmoveto{\pgfqpoint{1.757783in}{8.769617in}}%
\pgfpathquadraticcurveto{\pgfqpoint{1.757783in}{8.756546in}}{\pgfqpoint{1.763164in}{8.749103in}}%
\pgfpathquadraticcurveto{\pgfqpoint{1.768545in}{8.741661in}}{\pgfqpoint{1.777946in}{8.741661in}}%
\pgfpathquadraticcurveto{\pgfqpoint{1.787347in}{8.741661in}}{\pgfqpoint{1.792749in}{8.749103in}}%
\pgfpathquadraticcurveto{\pgfqpoint{1.798171in}{8.756546in}}{\pgfqpoint{1.798171in}{8.769617in}}%
\pgfpathquadraticcurveto{\pgfqpoint{1.798171in}{8.782687in}}{\pgfqpoint{1.792749in}{8.790130in}}%
\pgfpathquadraticcurveto{\pgfqpoint{1.787347in}{8.797572in}}{\pgfqpoint{1.777946in}{8.797572in}}%
\pgfpathquadraticcurveto{\pgfqpoint{1.768545in}{8.797572in}}{\pgfqpoint{1.763164in}{8.790130in}}%
\pgfpathquadraticcurveto{\pgfqpoint{1.757783in}{8.782687in}}{\pgfqpoint{1.757783in}{8.769617in}}%
\pgfpathlineto{\pgfqpoint{1.757783in}{8.769617in}}%
\pgfpathclose%
\pgfpathmoveto{\pgfqpoint{1.897459in}{8.829796in}}%
\pgfpathlineto{\pgfqpoint{1.908406in}{8.829796in}}%
\pgfpathlineto{\pgfqpoint{1.874905in}{8.721354in}}%
\pgfpathlineto{\pgfqpoint{1.863957in}{8.721354in}}%
\pgfpathlineto{\pgfqpoint{1.897459in}{8.829796in}}%
\pgfpathlineto{\pgfqpoint{1.897459in}{8.829796in}}%
\pgfpathclose%
\pgfpathmoveto{\pgfqpoint{2.003880in}{8.785471in}}%
\pgfpathquadraticcurveto{\pgfqpoint{2.013219in}{8.783471in}}{\pgfqpoint{2.018456in}{8.777142in}}%
\pgfpathquadraticcurveto{\pgfqpoint{2.023713in}{8.770833in}}{\pgfqpoint{2.023713in}{8.761556in}}%
\pgfpathquadraticcurveto{\pgfqpoint{2.023713in}{8.747330in}}{\pgfqpoint{2.013920in}{8.739517in}}%
\pgfpathquadraticcurveto{\pgfqpoint{2.004128in}{8.731724in}}{\pgfqpoint{1.986088in}{8.731724in}}%
\pgfpathquadraticcurveto{\pgfqpoint{1.980048in}{8.731724in}}{\pgfqpoint{1.973636in}{8.732920in}}%
\pgfpathquadraticcurveto{\pgfqpoint{1.967224in}{8.734115in}}{\pgfqpoint{1.960400in}{8.736507in}}%
\pgfpathlineto{\pgfqpoint{1.960400in}{8.749062in}}%
\pgfpathquadraticcurveto{\pgfqpoint{1.965802in}{8.745908in}}{\pgfqpoint{1.972234in}{8.744300in}}%
\pgfpathquadraticcurveto{\pgfqpoint{1.978687in}{8.742692in}}{\pgfqpoint{1.985717in}{8.742692in}}%
\pgfpathquadraticcurveto{\pgfqpoint{1.997943in}{8.742692in}}{\pgfqpoint{2.004354in}{8.747516in}}%
\pgfpathquadraticcurveto{\pgfqpoint{2.010766in}{8.752340in}}{\pgfqpoint{2.010766in}{8.761556in}}%
\pgfpathquadraticcurveto{\pgfqpoint{2.010766in}{8.770070in}}{\pgfqpoint{2.004808in}{8.774853in}}%
\pgfpathquadraticcurveto{\pgfqpoint{1.998850in}{8.779657in}}{\pgfqpoint{1.988232in}{8.779657in}}%
\pgfpathlineto{\pgfqpoint{1.977017in}{8.779657in}}%
\pgfpathlineto{\pgfqpoint{1.977017in}{8.790357in}}%
\pgfpathlineto{\pgfqpoint{1.988748in}{8.790357in}}%
\pgfpathquadraticcurveto{\pgfqpoint{1.998334in}{8.790357in}}{\pgfqpoint{2.003427in}{8.794191in}}%
\pgfpathquadraticcurveto{\pgfqpoint{2.008519in}{8.798026in}}{\pgfqpoint{2.008519in}{8.805242in}}%
\pgfpathquadraticcurveto{\pgfqpoint{2.008519in}{8.812643in}}{\pgfqpoint{2.003262in}{8.816601in}}%
\pgfpathquadraticcurveto{\pgfqpoint{1.998025in}{8.820580in}}{\pgfqpoint{1.988232in}{8.820580in}}%
\pgfpathquadraticcurveto{\pgfqpoint{1.982872in}{8.820580in}}{\pgfqpoint{1.976749in}{8.819405in}}%
\pgfpathquadraticcurveto{\pgfqpoint{1.970626in}{8.818251in}}{\pgfqpoint{1.963287in}{8.815818in}}%
\pgfpathlineto{\pgfqpoint{1.963287in}{8.827404in}}%
\pgfpathquadraticcurveto{\pgfqpoint{1.970709in}{8.829466in}}{\pgfqpoint{1.977182in}{8.830497in}}%
\pgfpathquadraticcurveto{\pgfqpoint{1.983656in}{8.831528in}}{\pgfqpoint{1.989387in}{8.831528in}}%
\pgfpathquadraticcurveto{\pgfqpoint{2.004210in}{8.831528in}}{\pgfqpoint{2.012828in}{8.824786in}}%
\pgfpathquadraticcurveto{\pgfqpoint{2.021466in}{8.818065in}}{\pgfqpoint{2.021466in}{8.806602in}}%
\pgfpathquadraticcurveto{\pgfqpoint{2.021466in}{8.798603in}}{\pgfqpoint{2.016889in}{8.793099in}}%
\pgfpathquadraticcurveto{\pgfqpoint{2.012312in}{8.787594in}}{\pgfqpoint{2.003880in}{8.785471in}}%
\pgfpathlineto{\pgfqpoint{2.003880in}{8.785471in}}%
\pgfpathclose%
\pgfpathmoveto{\pgfqpoint{2.077852in}{8.786873in}}%
\pgfpathquadraticcurveto{\pgfqpoint{2.069090in}{8.786873in}}{\pgfqpoint{2.063956in}{8.780873in}}%
\pgfpathquadraticcurveto{\pgfqpoint{2.058843in}{8.774894in}}{\pgfqpoint{2.058843in}{8.764463in}}%
\pgfpathquadraticcurveto{\pgfqpoint{2.058843in}{8.754093in}}{\pgfqpoint{2.063956in}{8.748052in}}%
\pgfpathquadraticcurveto{\pgfqpoint{2.069090in}{8.742032in}}{\pgfqpoint{2.077852in}{8.742032in}}%
\pgfpathquadraticcurveto{\pgfqpoint{2.086614in}{8.742032in}}{\pgfqpoint{2.091726in}{8.748052in}}%
\pgfpathquadraticcurveto{\pgfqpoint{2.096839in}{8.754093in}}{\pgfqpoint{2.096839in}{8.764463in}}%
\pgfpathquadraticcurveto{\pgfqpoint{2.096839in}{8.774894in}}{\pgfqpoint{2.091726in}{8.780873in}}%
\pgfpathquadraticcurveto{\pgfqpoint{2.086614in}{8.786873in}}{\pgfqpoint{2.077852in}{8.786873in}}%
\pgfpathlineto{\pgfqpoint{2.077852in}{8.786873in}}%
\pgfpathclose%
\pgfpathmoveto{\pgfqpoint{2.103684in}{8.827672in}}%
\pgfpathlineto{\pgfqpoint{2.103684in}{8.815818in}}%
\pgfpathquadraticcurveto{\pgfqpoint{2.098777in}{8.818127in}}{\pgfqpoint{2.093788in}{8.819343in}}%
\pgfpathquadraticcurveto{\pgfqpoint{2.088799in}{8.820580in}}{\pgfqpoint{2.083892in}{8.820580in}}%
\pgfpathquadraticcurveto{\pgfqpoint{2.071007in}{8.820580in}}{\pgfqpoint{2.064204in}{8.811880in}}%
\pgfpathquadraticcurveto{\pgfqpoint{2.057421in}{8.803180in}}{\pgfqpoint{2.056452in}{8.785594in}}%
\pgfpathquadraticcurveto{\pgfqpoint{2.060245in}{8.791202in}}{\pgfqpoint{2.065977in}{8.794191in}}%
\pgfpathquadraticcurveto{\pgfqpoint{2.071729in}{8.797181in}}{\pgfqpoint{2.078614in}{8.797181in}}%
\pgfpathquadraticcurveto{\pgfqpoint{2.093108in}{8.797181in}}{\pgfqpoint{2.101519in}{8.788378in}}%
\pgfpathquadraticcurveto{\pgfqpoint{2.109931in}{8.779595in}}{\pgfqpoint{2.109931in}{8.764463in}}%
\pgfpathquadraticcurveto{\pgfqpoint{2.109931in}{8.749639in}}{\pgfqpoint{2.101169in}{8.740671in}}%
\pgfpathquadraticcurveto{\pgfqpoint{2.092407in}{8.731724in}}{\pgfqpoint{2.077852in}{8.731724in}}%
\pgfpathquadraticcurveto{\pgfqpoint{2.061152in}{8.731724in}}{\pgfqpoint{2.052329in}{8.744506in}}%
\pgfpathquadraticcurveto{\pgfqpoint{2.043505in}{8.757309in}}{\pgfqpoint{2.043505in}{8.781595in}}%
\pgfpathquadraticcurveto{\pgfqpoint{2.043505in}{8.804396in}}{\pgfqpoint{2.054328in}{8.817962in}}%
\pgfpathquadraticcurveto{\pgfqpoint{2.065152in}{8.831528in}}{\pgfqpoint{2.083377in}{8.831528in}}%
\pgfpathquadraticcurveto{\pgfqpoint{2.088283in}{8.831528in}}{\pgfqpoint{2.093273in}{8.830559in}}%
\pgfpathquadraticcurveto{\pgfqpoint{2.098262in}{8.829590in}}{\pgfqpoint{2.103684in}{8.827672in}}%
\pgfpathlineto{\pgfqpoint{2.103684in}{8.827672in}}%
\pgfpathclose%
\pgfpathmoveto{\pgfqpoint{2.132732in}{8.735600in}}%
\pgfpathlineto{\pgfqpoint{2.132732in}{8.747454in}}%
\pgfpathquadraticcurveto{\pgfqpoint{2.137639in}{8.745125in}}{\pgfqpoint{2.142649in}{8.743908in}}%
\pgfpathquadraticcurveto{\pgfqpoint{2.147679in}{8.742692in}}{\pgfqpoint{2.152524in}{8.742692in}}%
\pgfpathquadraticcurveto{\pgfqpoint{2.165409in}{8.742692in}}{\pgfqpoint{2.172192in}{8.751351in}}%
\pgfpathquadraticcurveto{\pgfqpoint{2.178995in}{8.760010in}}{\pgfqpoint{2.179964in}{8.777678in}}%
\pgfpathquadraticcurveto{\pgfqpoint{2.176233in}{8.772132in}}{\pgfqpoint{2.170481in}{8.769163in}}%
\pgfpathquadraticcurveto{\pgfqpoint{2.164749in}{8.766194in}}{\pgfqpoint{2.157802in}{8.766194in}}%
\pgfpathquadraticcurveto{\pgfqpoint{2.143370in}{8.766194in}}{\pgfqpoint{2.134959in}{8.774915in}}%
\pgfpathquadraticcurveto{\pgfqpoint{2.126547in}{8.783656in}}{\pgfqpoint{2.126547in}{8.798809in}}%
\pgfpathquadraticcurveto{\pgfqpoint{2.126547in}{8.813612in}}{\pgfqpoint{2.135309in}{8.822559in}}%
\pgfpathquadraticcurveto{\pgfqpoint{2.144071in}{8.831528in}}{\pgfqpoint{2.158626in}{8.831528in}}%
\pgfpathquadraticcurveto{\pgfqpoint{2.175326in}{8.831528in}}{\pgfqpoint{2.184108in}{8.818725in}}%
\pgfpathquadraticcurveto{\pgfqpoint{2.192911in}{8.805943in}}{\pgfqpoint{2.192911in}{8.781595in}}%
\pgfpathquadraticcurveto{\pgfqpoint{2.192911in}{8.758855in}}{\pgfqpoint{2.182108in}{8.745289in}}%
\pgfpathquadraticcurveto{\pgfqpoint{2.171326in}{8.731724in}}{\pgfqpoint{2.153101in}{8.731724in}}%
\pgfpathquadraticcurveto{\pgfqpoint{2.148194in}{8.731724in}}{\pgfqpoint{2.143164in}{8.732693in}}%
\pgfpathquadraticcurveto{\pgfqpoint{2.138154in}{8.733662in}}{\pgfqpoint{2.132732in}{8.735600in}}%
\pgfpathlineto{\pgfqpoint{2.132732in}{8.735600in}}%
\pgfpathclose%
\pgfpathmoveto{\pgfqpoint{2.158626in}{8.776379in}}%
\pgfpathquadraticcurveto{\pgfqpoint{2.167388in}{8.776379in}}{\pgfqpoint{2.172501in}{8.782358in}}%
\pgfpathquadraticcurveto{\pgfqpoint{2.177635in}{8.788357in}}{\pgfqpoint{2.177635in}{8.798809in}}%
\pgfpathquadraticcurveto{\pgfqpoint{2.177635in}{8.809179in}}{\pgfqpoint{2.172501in}{8.815199in}}%
\pgfpathquadraticcurveto{\pgfqpoint{2.167388in}{8.821219in}}{\pgfqpoint{2.158626in}{8.821219in}}%
\pgfpathquadraticcurveto{\pgfqpoint{2.149864in}{8.821219in}}{\pgfqpoint{2.144752in}{8.815199in}}%
\pgfpathquadraticcurveto{\pgfqpoint{2.139639in}{8.809179in}}{\pgfqpoint{2.139639in}{8.798809in}}%
\pgfpathquadraticcurveto{\pgfqpoint{2.139639in}{8.788357in}}{\pgfqpoint{2.144752in}{8.782358in}}%
\pgfpathquadraticcurveto{\pgfqpoint{2.149864in}{8.776379in}}{\pgfqpoint{2.158626in}{8.776379in}}%
\pgfpathlineto{\pgfqpoint{2.158626in}{8.776379in}}%
\pgfpathclose%
\pgfpathmoveto{\pgfqpoint{2.308486in}{8.802995in}}%
\pgfpathlineto{\pgfqpoint{2.308486in}{8.791903in}}%
\pgfpathquadraticcurveto{\pgfqpoint{2.303456in}{8.794686in}}{\pgfqpoint{2.298405in}{8.796067in}}%
\pgfpathquadraticcurveto{\pgfqpoint{2.293354in}{8.797449in}}{\pgfqpoint{2.288200in}{8.797449in}}%
\pgfpathquadraticcurveto{\pgfqpoint{2.276655in}{8.797449in}}{\pgfqpoint{2.270264in}{8.790130in}}%
\pgfpathquadraticcurveto{\pgfqpoint{2.263893in}{8.782832in}}{\pgfqpoint{2.263893in}{8.769617in}}%
\pgfpathquadraticcurveto{\pgfqpoint{2.263893in}{8.756402in}}{\pgfqpoint{2.270264in}{8.749083in}}%
\pgfpathquadraticcurveto{\pgfqpoint{2.276655in}{8.741785in}}{\pgfqpoint{2.288200in}{8.741785in}}%
\pgfpathquadraticcurveto{\pgfqpoint{2.293354in}{8.741785in}}{\pgfqpoint{2.298405in}{8.743166in}}%
\pgfpathquadraticcurveto{\pgfqpoint{2.303456in}{8.744547in}}{\pgfqpoint{2.308486in}{8.747330in}}%
\pgfpathlineto{\pgfqpoint{2.308486in}{8.736363in}}%
\pgfpathquadraticcurveto{\pgfqpoint{2.303518in}{8.734054in}}{\pgfqpoint{2.298199in}{8.732899in}}%
\pgfpathquadraticcurveto{\pgfqpoint{2.292900in}{8.731724in}}{\pgfqpoint{2.286901in}{8.731724in}}%
\pgfpathquadraticcurveto{\pgfqpoint{2.270594in}{8.731724in}}{\pgfqpoint{2.260986in}{8.741970in}}%
\pgfpathquadraticcurveto{\pgfqpoint{2.251400in}{8.752217in}}{\pgfqpoint{2.251400in}{8.769617in}}%
\pgfpathquadraticcurveto{\pgfqpoint{2.251400in}{8.787264in}}{\pgfqpoint{2.261089in}{8.797366in}}%
\pgfpathquadraticcurveto{\pgfqpoint{2.270800in}{8.807489in}}{\pgfqpoint{2.287684in}{8.807489in}}%
\pgfpathquadraticcurveto{\pgfqpoint{2.293148in}{8.807489in}}{\pgfqpoint{2.298364in}{8.806355in}}%
\pgfpathquadraticcurveto{\pgfqpoint{2.303580in}{8.805242in}}{\pgfqpoint{2.308486in}{8.802995in}}%
\pgfpathlineto{\pgfqpoint{2.308486in}{8.802995in}}%
\pgfpathclose%
\pgfpathmoveto{\pgfqpoint{2.357058in}{8.797449in}}%
\pgfpathquadraticcurveto{\pgfqpoint{2.347534in}{8.797449in}}{\pgfqpoint{2.341988in}{8.790006in}}%
\pgfpathquadraticcurveto{\pgfqpoint{2.336442in}{8.782564in}}{\pgfqpoint{2.336442in}{8.769617in}}%
\pgfpathquadraticcurveto{\pgfqpoint{2.336442in}{8.756670in}}{\pgfqpoint{2.341947in}{8.749227in}}%
\pgfpathquadraticcurveto{\pgfqpoint{2.347472in}{8.741785in}}{\pgfqpoint{2.357058in}{8.741785in}}%
\pgfpathquadraticcurveto{\pgfqpoint{2.366542in}{8.741785in}}{\pgfqpoint{2.372067in}{8.749248in}}%
\pgfpathquadraticcurveto{\pgfqpoint{2.377613in}{8.756732in}}{\pgfqpoint{2.377613in}{8.769617in}}%
\pgfpathquadraticcurveto{\pgfqpoint{2.377613in}{8.782440in}}{\pgfqpoint{2.372067in}{8.789944in}}%
\pgfpathquadraticcurveto{\pgfqpoint{2.366542in}{8.797449in}}{\pgfqpoint{2.357058in}{8.797449in}}%
\pgfpathlineto{\pgfqpoint{2.357058in}{8.797449in}}%
\pgfpathclose%
\pgfpathmoveto{\pgfqpoint{2.357058in}{8.807489in}}%
\pgfpathquadraticcurveto{\pgfqpoint{2.372521in}{8.807489in}}{\pgfqpoint{2.381344in}{8.797428in}}%
\pgfpathquadraticcurveto{\pgfqpoint{2.390189in}{8.787388in}}{\pgfqpoint{2.390189in}{8.769617in}}%
\pgfpathquadraticcurveto{\pgfqpoint{2.390189in}{8.751907in}}{\pgfqpoint{2.381344in}{8.741805in}}%
\pgfpathquadraticcurveto{\pgfqpoint{2.372521in}{8.731724in}}{\pgfqpoint{2.357058in}{8.731724in}}%
\pgfpathquadraticcurveto{\pgfqpoint{2.341534in}{8.731724in}}{\pgfqpoint{2.332731in}{8.741805in}}%
\pgfpathquadraticcurveto{\pgfqpoint{2.323949in}{8.751907in}}{\pgfqpoint{2.323949in}{8.769617in}}%
\pgfpathquadraticcurveto{\pgfqpoint{2.323949in}{8.787388in}}{\pgfqpoint{2.332731in}{8.797428in}}%
\pgfpathquadraticcurveto{\pgfqpoint{2.341534in}{8.807489in}}{\pgfqpoint{2.357058in}{8.807489in}}%
\pgfpathlineto{\pgfqpoint{2.357058in}{8.807489in}}%
\pgfpathclose%
\pgfpathmoveto{\pgfqpoint{2.466016in}{8.791903in}}%
\pgfpathquadraticcurveto{\pgfqpoint{2.470469in}{8.799902in}}{\pgfqpoint{2.476654in}{8.803695in}}%
\pgfpathquadraticcurveto{\pgfqpoint{2.482839in}{8.807489in}}{\pgfqpoint{2.491209in}{8.807489in}}%
\pgfpathquadraticcurveto{\pgfqpoint{2.502486in}{8.807489in}}{\pgfqpoint{2.508609in}{8.799593in}}%
\pgfpathquadraticcurveto{\pgfqpoint{2.514732in}{8.791717in}}{\pgfqpoint{2.514732in}{8.777162in}}%
\pgfpathlineto{\pgfqpoint{2.514732in}{8.733600in}}%
\pgfpathlineto{\pgfqpoint{2.502816in}{8.733600in}}%
\pgfpathlineto{\pgfqpoint{2.502816in}{8.776771in}}%
\pgfpathquadraticcurveto{\pgfqpoint{2.502816in}{8.787141in}}{\pgfqpoint{2.499125in}{8.792150in}}%
\pgfpathquadraticcurveto{\pgfqpoint{2.495456in}{8.797181in}}{\pgfqpoint{2.487931in}{8.797181in}}%
\pgfpathquadraticcurveto{\pgfqpoint{2.478715in}{8.797181in}}{\pgfqpoint{2.473355in}{8.791058in}}%
\pgfpathquadraticcurveto{\pgfqpoint{2.468015in}{8.784955in}}{\pgfqpoint{2.468015in}{8.774379in}}%
\pgfpathlineto{\pgfqpoint{2.468015in}{8.733600in}}%
\pgfpathlineto{\pgfqpoint{2.456099in}{8.733600in}}%
\pgfpathlineto{\pgfqpoint{2.456099in}{8.776771in}}%
\pgfpathquadraticcurveto{\pgfqpoint{2.456099in}{8.787202in}}{\pgfqpoint{2.452429in}{8.792192in}}%
\pgfpathquadraticcurveto{\pgfqpoint{2.448760in}{8.797181in}}{\pgfqpoint{2.441091in}{8.797181in}}%
\pgfpathquadraticcurveto{\pgfqpoint{2.431999in}{8.797181in}}{\pgfqpoint{2.426638in}{8.791037in}}%
\pgfpathquadraticcurveto{\pgfqpoint{2.421299in}{8.784893in}}{\pgfqpoint{2.421299in}{8.774379in}}%
\pgfpathlineto{\pgfqpoint{2.421299in}{8.733600in}}%
\pgfpathlineto{\pgfqpoint{2.409383in}{8.733600in}}%
\pgfpathlineto{\pgfqpoint{2.409383in}{8.805757in}}%
\pgfpathlineto{\pgfqpoint{2.421299in}{8.805757in}}%
\pgfpathlineto{\pgfqpoint{2.421299in}{8.794542in}}%
\pgfpathquadraticcurveto{\pgfqpoint{2.425360in}{8.801180in}}{\pgfqpoint{2.431030in}{8.804335in}}%
\pgfpathquadraticcurveto{\pgfqpoint{2.436699in}{8.807489in}}{\pgfqpoint{2.444492in}{8.807489in}}%
\pgfpathquadraticcurveto{\pgfqpoint{2.452368in}{8.807489in}}{\pgfqpoint{2.457872in}{8.803489in}}%
\pgfpathquadraticcurveto{\pgfqpoint{2.463377in}{8.799510in}}{\pgfqpoint{2.466016in}{8.791903in}}%
\pgfpathlineto{\pgfqpoint{2.466016in}{8.791903in}}%
\pgfpathclose%
\pgfpathmoveto{\pgfqpoint{2.549821in}{8.744424in}}%
\pgfpathlineto{\pgfqpoint{2.549821in}{8.706160in}}%
\pgfpathlineto{\pgfqpoint{2.537905in}{8.706160in}}%
\pgfpathlineto{\pgfqpoint{2.537905in}{8.805757in}}%
\pgfpathlineto{\pgfqpoint{2.549821in}{8.805757in}}%
\pgfpathlineto{\pgfqpoint{2.549821in}{8.794810in}}%
\pgfpathquadraticcurveto{\pgfqpoint{2.553573in}{8.801242in}}{\pgfqpoint{2.559263in}{8.804355in}}%
\pgfpathquadraticcurveto{\pgfqpoint{2.564974in}{8.807489in}}{\pgfqpoint{2.572891in}{8.807489in}}%
\pgfpathquadraticcurveto{\pgfqpoint{2.586044in}{8.807489in}}{\pgfqpoint{2.594249in}{8.797057in}}%
\pgfpathquadraticcurveto{\pgfqpoint{2.602475in}{8.786625in}}{\pgfqpoint{2.602475in}{8.769617in}}%
\pgfpathquadraticcurveto{\pgfqpoint{2.602475in}{8.752608in}}{\pgfqpoint{2.594249in}{8.742156in}}%
\pgfpathquadraticcurveto{\pgfqpoint{2.586044in}{8.731724in}}{\pgfqpoint{2.572891in}{8.731724in}}%
\pgfpathquadraticcurveto{\pgfqpoint{2.564974in}{8.731724in}}{\pgfqpoint{2.559263in}{8.734858in}}%
\pgfpathquadraticcurveto{\pgfqpoint{2.553573in}{8.737991in}}{\pgfqpoint{2.549821in}{8.744424in}}%
\pgfpathlineto{\pgfqpoint{2.549821in}{8.744424in}}%
\pgfpathclose%
\pgfpathmoveto{\pgfqpoint{2.590167in}{8.769617in}}%
\pgfpathquadraticcurveto{\pgfqpoint{2.590167in}{8.782687in}}{\pgfqpoint{2.584786in}{8.790130in}}%
\pgfpathquadraticcurveto{\pgfqpoint{2.579405in}{8.797572in}}{\pgfqpoint{2.570004in}{8.797572in}}%
\pgfpathquadraticcurveto{\pgfqpoint{2.560583in}{8.797572in}}{\pgfqpoint{2.555202in}{8.790130in}}%
\pgfpathquadraticcurveto{\pgfqpoint{2.549821in}{8.782687in}}{\pgfqpoint{2.549821in}{8.769617in}}%
\pgfpathquadraticcurveto{\pgfqpoint{2.549821in}{8.756546in}}{\pgfqpoint{2.555202in}{8.749103in}}%
\pgfpathquadraticcurveto{\pgfqpoint{2.560583in}{8.741661in}}{\pgfqpoint{2.570004in}{8.741661in}}%
\pgfpathquadraticcurveto{\pgfqpoint{2.579405in}{8.741661in}}{\pgfqpoint{2.584786in}{8.749103in}}%
\pgfpathquadraticcurveto{\pgfqpoint{2.590167in}{8.756546in}}{\pgfqpoint{2.590167in}{8.769617in}}%
\pgfpathlineto{\pgfqpoint{2.590167in}{8.769617in}}%
\pgfpathclose%
\pgfpathmoveto{\pgfqpoint{2.654923in}{8.769864in}}%
\pgfpathquadraticcurveto{\pgfqpoint{2.640553in}{8.769864in}}{\pgfqpoint{2.635008in}{8.766586in}}%
\pgfpathquadraticcurveto{\pgfqpoint{2.629462in}{8.763308in}}{\pgfqpoint{2.629462in}{8.755371in}}%
\pgfpathquadraticcurveto{\pgfqpoint{2.629462in}{8.749062in}}{\pgfqpoint{2.633626in}{8.745351in}}%
\pgfpathquadraticcurveto{\pgfqpoint{2.637791in}{8.741661in}}{\pgfqpoint{2.644924in}{8.741661in}}%
\pgfpathquadraticcurveto{\pgfqpoint{2.654799in}{8.741661in}}{\pgfqpoint{2.660757in}{8.748650in}}%
\pgfpathquadraticcurveto{\pgfqpoint{2.666716in}{8.755639in}}{\pgfqpoint{2.666716in}{8.767225in}}%
\pgfpathlineto{\pgfqpoint{2.666716in}{8.769864in}}%
\pgfpathlineto{\pgfqpoint{2.654923in}{8.769864in}}%
\pgfpathlineto{\pgfqpoint{2.654923in}{8.769864in}}%
\pgfpathclose%
\pgfpathmoveto{\pgfqpoint{2.678570in}{8.774771in}}%
\pgfpathlineto{\pgfqpoint{2.678570in}{8.733600in}}%
\pgfpathlineto{\pgfqpoint{2.666716in}{8.733600in}}%
\pgfpathlineto{\pgfqpoint{2.666716in}{8.744547in}}%
\pgfpathquadraticcurveto{\pgfqpoint{2.662654in}{8.737991in}}{\pgfqpoint{2.656593in}{8.734858in}}%
\pgfpathquadraticcurveto{\pgfqpoint{2.650532in}{8.731724in}}{\pgfqpoint{2.641770in}{8.731724in}}%
\pgfpathquadraticcurveto{\pgfqpoint{2.630699in}{8.731724in}}{\pgfqpoint{2.624143in}{8.737950in}}%
\pgfpathquadraticcurveto{\pgfqpoint{2.617607in}{8.744176in}}{\pgfqpoint{2.617607in}{8.754608in}}%
\pgfpathquadraticcurveto{\pgfqpoint{2.617607in}{8.766772in}}{\pgfqpoint{2.625751in}{8.772957in}}%
\pgfpathquadraticcurveto{\pgfqpoint{2.633915in}{8.779141in}}{\pgfqpoint{2.650078in}{8.779141in}}%
\pgfpathlineto{\pgfqpoint{2.666716in}{8.779141in}}%
\pgfpathlineto{\pgfqpoint{2.666716in}{8.780317in}}%
\pgfpathquadraticcurveto{\pgfqpoint{2.666716in}{8.788501in}}{\pgfqpoint{2.661335in}{8.792975in}}%
\pgfpathquadraticcurveto{\pgfqpoint{2.655954in}{8.797449in}}{\pgfqpoint{2.646223in}{8.797449in}}%
\pgfpathquadraticcurveto{\pgfqpoint{2.640038in}{8.797449in}}{\pgfqpoint{2.634162in}{8.795964in}}%
\pgfpathquadraticcurveto{\pgfqpoint{2.628307in}{8.794480in}}{\pgfqpoint{2.622906in}{8.791511in}}%
\pgfpathlineto{\pgfqpoint{2.622906in}{8.802479in}}%
\pgfpathquadraticcurveto{\pgfqpoint{2.629400in}{8.804994in}}{\pgfqpoint{2.635523in}{8.806231in}}%
\pgfpathquadraticcurveto{\pgfqpoint{2.641646in}{8.807489in}}{\pgfqpoint{2.647439in}{8.807489in}}%
\pgfpathquadraticcurveto{\pgfqpoint{2.663108in}{8.807489in}}{\pgfqpoint{2.670839in}{8.799366in}}%
\pgfpathquadraticcurveto{\pgfqpoint{2.678570in}{8.791264in}}{\pgfqpoint{2.678570in}{8.774771in}}%
\pgfpathlineto{\pgfqpoint{2.678570in}{8.774771in}}%
\pgfpathclose%
\pgfpathmoveto{\pgfqpoint{2.744790in}{8.794686in}}%
\pgfpathquadraticcurveto{\pgfqpoint{2.742790in}{8.795841in}}{\pgfqpoint{2.740439in}{8.796377in}}%
\pgfpathquadraticcurveto{\pgfqpoint{2.738089in}{8.796933in}}{\pgfqpoint{2.735265in}{8.796933in}}%
\pgfpathquadraticcurveto{\pgfqpoint{2.725204in}{8.796933in}}{\pgfqpoint{2.719823in}{8.790398in}}%
\pgfpathquadraticcurveto{\pgfqpoint{2.714442in}{8.783863in}}{\pgfqpoint{2.714442in}{8.771616in}}%
\pgfpathlineto{\pgfqpoint{2.714442in}{8.733600in}}%
\pgfpathlineto{\pgfqpoint{2.702526in}{8.733600in}}%
\pgfpathlineto{\pgfqpoint{2.702526in}{8.805757in}}%
\pgfpathlineto{\pgfqpoint{2.714442in}{8.805757in}}%
\pgfpathlineto{\pgfqpoint{2.714442in}{8.794542in}}%
\pgfpathquadraticcurveto{\pgfqpoint{2.718194in}{8.801118in}}{\pgfqpoint{2.724173in}{8.804293in}}%
\pgfpathquadraticcurveto{\pgfqpoint{2.730173in}{8.807489in}}{\pgfqpoint{2.738749in}{8.807489in}}%
\pgfpathquadraticcurveto{\pgfqpoint{2.739965in}{8.807489in}}{\pgfqpoint{2.741450in}{8.807324in}}%
\pgfpathquadraticcurveto{\pgfqpoint{2.742934in}{8.807180in}}{\pgfqpoint{2.744728in}{8.806850in}}%
\pgfpathlineto{\pgfqpoint{2.744790in}{8.794686in}}%
\pgfpathlineto{\pgfqpoint{2.744790in}{8.794686in}}%
\pgfpathclose%
\pgfpathmoveto{\pgfqpoint{2.757221in}{8.805757in}}%
\pgfpathlineto{\pgfqpoint{2.769076in}{8.805757in}}%
\pgfpathlineto{\pgfqpoint{2.769076in}{8.733600in}}%
\pgfpathlineto{\pgfqpoint{2.757221in}{8.733600in}}%
\pgfpathlineto{\pgfqpoint{2.757221in}{8.805757in}}%
\pgfpathlineto{\pgfqpoint{2.757221in}{8.805757in}}%
\pgfpathclose%
\pgfpathmoveto{\pgfqpoint{2.757221in}{8.833857in}}%
\pgfpathlineto{\pgfqpoint{2.769076in}{8.833857in}}%
\pgfpathlineto{\pgfqpoint{2.769076in}{8.818828in}}%
\pgfpathlineto{\pgfqpoint{2.757221in}{8.818828in}}%
\pgfpathlineto{\pgfqpoint{2.757221in}{8.833857in}}%
\pgfpathlineto{\pgfqpoint{2.757221in}{8.833857in}}%
\pgfpathclose%
\pgfpathmoveto{\pgfqpoint{2.839872in}{8.803634in}}%
\pgfpathlineto{\pgfqpoint{2.839872in}{8.792418in}}%
\pgfpathquadraticcurveto{\pgfqpoint{2.834862in}{8.794995in}}{\pgfqpoint{2.829440in}{8.796274in}}%
\pgfpathquadraticcurveto{\pgfqpoint{2.824039in}{8.797572in}}{\pgfqpoint{2.818225in}{8.797572in}}%
\pgfpathquadraticcurveto{\pgfqpoint{2.809401in}{8.797572in}}{\pgfqpoint{2.804989in}{8.794872in}}%
\pgfpathquadraticcurveto{\pgfqpoint{2.800577in}{8.792171in}}{\pgfqpoint{2.800577in}{8.786749in}}%
\pgfpathquadraticcurveto{\pgfqpoint{2.800577in}{8.782626in}}{\pgfqpoint{2.803732in}{8.780275in}}%
\pgfpathquadraticcurveto{\pgfqpoint{2.806886in}{8.777925in}}{\pgfqpoint{2.816431in}{8.775802in}}%
\pgfpathlineto{\pgfqpoint{2.820493in}{8.774894in}}%
\pgfpathquadraticcurveto{\pgfqpoint{2.833110in}{8.772194in}}{\pgfqpoint{2.838429in}{8.767266in}}%
\pgfpathquadraticcurveto{\pgfqpoint{2.843748in}{8.762339in}}{\pgfqpoint{2.843748in}{8.753515in}}%
\pgfpathquadraticcurveto{\pgfqpoint{2.843748in}{8.743455in}}{\pgfqpoint{2.835790in}{8.737579in}}%
\pgfpathquadraticcurveto{\pgfqpoint{2.827832in}{8.731724in}}{\pgfqpoint{2.813916in}{8.731724in}}%
\pgfpathquadraticcurveto{\pgfqpoint{2.808123in}{8.731724in}}{\pgfqpoint{2.801835in}{8.732858in}}%
\pgfpathquadraticcurveto{\pgfqpoint{2.795547in}{8.733992in}}{\pgfqpoint{2.788599in}{8.736239in}}%
\pgfpathlineto{\pgfqpoint{2.788599in}{8.748485in}}%
\pgfpathquadraticcurveto{\pgfqpoint{2.795176in}{8.745063in}}{\pgfqpoint{2.801546in}{8.743352in}}%
\pgfpathquadraticcurveto{\pgfqpoint{2.807917in}{8.741661in}}{\pgfqpoint{2.814184in}{8.741661in}}%
\pgfpathquadraticcurveto{\pgfqpoint{2.822554in}{8.741661in}}{\pgfqpoint{2.827049in}{8.744527in}}%
\pgfpathquadraticcurveto{\pgfqpoint{2.831564in}{8.747392in}}{\pgfqpoint{2.831564in}{8.752608in}}%
\pgfpathquadraticcurveto{\pgfqpoint{2.831564in}{8.757432in}}{\pgfqpoint{2.828306in}{8.760010in}}%
\pgfpathquadraticcurveto{\pgfqpoint{2.825069in}{8.762587in}}{\pgfqpoint{2.814040in}{8.764978in}}%
\pgfpathlineto{\pgfqpoint{2.809916in}{8.765947in}}%
\pgfpathquadraticcurveto{\pgfqpoint{2.798907in}{8.768256in}}{\pgfqpoint{2.794001in}{8.773060in}}%
\pgfpathquadraticcurveto{\pgfqpoint{2.789115in}{8.777863in}}{\pgfqpoint{2.789115in}{8.786233in}}%
\pgfpathquadraticcurveto{\pgfqpoint{2.789115in}{8.796418in}}{\pgfqpoint{2.796330in}{8.801943in}}%
\pgfpathquadraticcurveto{\pgfqpoint{2.803546in}{8.807489in}}{\pgfqpoint{2.816823in}{8.807489in}}%
\pgfpathquadraticcurveto{\pgfqpoint{2.823379in}{8.807489in}}{\pgfqpoint{2.829172in}{8.806520in}}%
\pgfpathquadraticcurveto{\pgfqpoint{2.834986in}{8.805572in}}{\pgfqpoint{2.839872in}{8.803634in}}%
\pgfpathlineto{\pgfqpoint{2.839872in}{8.803634in}}%
\pgfpathclose%
\pgfpathmoveto{\pgfqpoint{2.890568in}{8.797449in}}%
\pgfpathquadraticcurveto{\pgfqpoint{2.881043in}{8.797449in}}{\pgfqpoint{2.875497in}{8.790006in}}%
\pgfpathquadraticcurveto{\pgfqpoint{2.869951in}{8.782564in}}{\pgfqpoint{2.869951in}{8.769617in}}%
\pgfpathquadraticcurveto{\pgfqpoint{2.869951in}{8.756670in}}{\pgfqpoint{2.875456in}{8.749227in}}%
\pgfpathquadraticcurveto{\pgfqpoint{2.880981in}{8.741785in}}{\pgfqpoint{2.890568in}{8.741785in}}%
\pgfpathquadraticcurveto{\pgfqpoint{2.900051in}{8.741785in}}{\pgfqpoint{2.905576in}{8.749248in}}%
\pgfpathquadraticcurveto{\pgfqpoint{2.911122in}{8.756732in}}{\pgfqpoint{2.911122in}{8.769617in}}%
\pgfpathquadraticcurveto{\pgfqpoint{2.911122in}{8.782440in}}{\pgfqpoint{2.905576in}{8.789944in}}%
\pgfpathquadraticcurveto{\pgfqpoint{2.900051in}{8.797449in}}{\pgfqpoint{2.890568in}{8.797449in}}%
\pgfpathlineto{\pgfqpoint{2.890568in}{8.797449in}}%
\pgfpathclose%
\pgfpathmoveto{\pgfqpoint{2.890568in}{8.807489in}}%
\pgfpathquadraticcurveto{\pgfqpoint{2.906030in}{8.807489in}}{\pgfqpoint{2.914854in}{8.797428in}}%
\pgfpathquadraticcurveto{\pgfqpoint{2.923698in}{8.787388in}}{\pgfqpoint{2.923698in}{8.769617in}}%
\pgfpathquadraticcurveto{\pgfqpoint{2.923698in}{8.751907in}}{\pgfqpoint{2.914854in}{8.741805in}}%
\pgfpathquadraticcurveto{\pgfqpoint{2.906030in}{8.731724in}}{\pgfqpoint{2.890568in}{8.731724in}}%
\pgfpathquadraticcurveto{\pgfqpoint{2.875043in}{8.731724in}}{\pgfqpoint{2.866240in}{8.741805in}}%
\pgfpathquadraticcurveto{\pgfqpoint{2.857458in}{8.751907in}}{\pgfqpoint{2.857458in}{8.769617in}}%
\pgfpathquadraticcurveto{\pgfqpoint{2.857458in}{8.787388in}}{\pgfqpoint{2.866240in}{8.797428in}}%
\pgfpathquadraticcurveto{\pgfqpoint{2.875043in}{8.807489in}}{\pgfqpoint{2.890568in}{8.807489in}}%
\pgfpathlineto{\pgfqpoint{2.890568in}{8.807489in}}%
\pgfpathclose%
\pgfpathmoveto{\pgfqpoint{3.003339in}{8.777162in}}%
\pgfpathlineto{\pgfqpoint{3.003339in}{8.733600in}}%
\pgfpathlineto{\pgfqpoint{2.991484in}{8.733600in}}%
\pgfpathlineto{\pgfqpoint{2.991484in}{8.776771in}}%
\pgfpathquadraticcurveto{\pgfqpoint{2.991484in}{8.787017in}}{\pgfqpoint{2.987485in}{8.792088in}}%
\pgfpathquadraticcurveto{\pgfqpoint{2.983485in}{8.797181in}}{\pgfqpoint{2.975507in}{8.797181in}}%
\pgfpathquadraticcurveto{\pgfqpoint{2.965900in}{8.797181in}}{\pgfqpoint{2.960354in}{8.791058in}}%
\pgfpathquadraticcurveto{\pgfqpoint{2.954808in}{8.784955in}}{\pgfqpoint{2.954808in}{8.774379in}}%
\pgfpathlineto{\pgfqpoint{2.954808in}{8.733600in}}%
\pgfpathlineto{\pgfqpoint{2.942892in}{8.733600in}}%
\pgfpathlineto{\pgfqpoint{2.942892in}{8.805757in}}%
\pgfpathlineto{\pgfqpoint{2.954808in}{8.805757in}}%
\pgfpathlineto{\pgfqpoint{2.954808in}{8.794542in}}%
\pgfpathquadraticcurveto{\pgfqpoint{2.959076in}{8.801057in}}{\pgfqpoint{2.964827in}{8.804273in}}%
\pgfpathquadraticcurveto{\pgfqpoint{2.970600in}{8.807489in}}{\pgfqpoint{2.978146in}{8.807489in}}%
\pgfpathquadraticcurveto{\pgfqpoint{2.990577in}{8.807489in}}{\pgfqpoint{2.996948in}{8.799799in}}%
\pgfpathquadraticcurveto{\pgfqpoint{3.003339in}{8.792109in}}{\pgfqpoint{3.003339in}{8.777162in}}%
\pgfpathlineto{\pgfqpoint{3.003339in}{8.777162in}}%
\pgfpathclose%
\pgfpathmoveto{\pgfqpoint{3.072960in}{8.803634in}}%
\pgfpathlineto{\pgfqpoint{3.072960in}{8.792418in}}%
\pgfpathquadraticcurveto{\pgfqpoint{3.067950in}{8.794995in}}{\pgfqpoint{3.062528in}{8.796274in}}%
\pgfpathquadraticcurveto{\pgfqpoint{3.057127in}{8.797572in}}{\pgfqpoint{3.051313in}{8.797572in}}%
\pgfpathquadraticcurveto{\pgfqpoint{3.042489in}{8.797572in}}{\pgfqpoint{3.038077in}{8.794872in}}%
\pgfpathquadraticcurveto{\pgfqpoint{3.033665in}{8.792171in}}{\pgfqpoint{3.033665in}{8.786749in}}%
\pgfpathquadraticcurveto{\pgfqpoint{3.033665in}{8.782626in}}{\pgfqpoint{3.036820in}{8.780275in}}%
\pgfpathquadraticcurveto{\pgfqpoint{3.039974in}{8.777925in}}{\pgfqpoint{3.049519in}{8.775802in}}%
\pgfpathlineto{\pgfqpoint{3.053581in}{8.774894in}}%
\pgfpathquadraticcurveto{\pgfqpoint{3.066198in}{8.772194in}}{\pgfqpoint{3.071517in}{8.767266in}}%
\pgfpathquadraticcurveto{\pgfqpoint{3.076836in}{8.762339in}}{\pgfqpoint{3.076836in}{8.753515in}}%
\pgfpathquadraticcurveto{\pgfqpoint{3.076836in}{8.743455in}}{\pgfqpoint{3.068878in}{8.737579in}}%
\pgfpathquadraticcurveto{\pgfqpoint{3.060920in}{8.731724in}}{\pgfqpoint{3.047004in}{8.731724in}}%
\pgfpathquadraticcurveto{\pgfqpoint{3.041211in}{8.731724in}}{\pgfqpoint{3.034923in}{8.732858in}}%
\pgfpathquadraticcurveto{\pgfqpoint{3.028635in}{8.733992in}}{\pgfqpoint{3.021687in}{8.736239in}}%
\pgfpathlineto{\pgfqpoint{3.021687in}{8.748485in}}%
\pgfpathquadraticcurveto{\pgfqpoint{3.028264in}{8.745063in}}{\pgfqpoint{3.034634in}{8.743352in}}%
\pgfpathquadraticcurveto{\pgfqpoint{3.041005in}{8.741661in}}{\pgfqpoint{3.047272in}{8.741661in}}%
\pgfpathquadraticcurveto{\pgfqpoint{3.055642in}{8.741661in}}{\pgfqpoint{3.060137in}{8.744527in}}%
\pgfpathquadraticcurveto{\pgfqpoint{3.064652in}{8.747392in}}{\pgfqpoint{3.064652in}{8.752608in}}%
\pgfpathquadraticcurveto{\pgfqpoint{3.064652in}{8.757432in}}{\pgfqpoint{3.061394in}{8.760010in}}%
\pgfpathquadraticcurveto{\pgfqpoint{3.058158in}{8.762587in}}{\pgfqpoint{3.047128in}{8.764978in}}%
\pgfpathlineto{\pgfqpoint{3.043005in}{8.765947in}}%
\pgfpathquadraticcurveto{\pgfqpoint{3.031995in}{8.768256in}}{\pgfqpoint{3.027089in}{8.773060in}}%
\pgfpathquadraticcurveto{\pgfqpoint{3.022203in}{8.777863in}}{\pgfqpoint{3.022203in}{8.786233in}}%
\pgfpathquadraticcurveto{\pgfqpoint{3.022203in}{8.796418in}}{\pgfqpoint{3.029418in}{8.801943in}}%
\pgfpathquadraticcurveto{\pgfqpoint{3.036634in}{8.807489in}}{\pgfqpoint{3.049911in}{8.807489in}}%
\pgfpathquadraticcurveto{\pgfqpoint{3.056467in}{8.807489in}}{\pgfqpoint{3.062260in}{8.806520in}}%
\pgfpathquadraticcurveto{\pgfqpoint{3.068074in}{8.805572in}}{\pgfqpoint{3.072960in}{8.803634in}}%
\pgfpathlineto{\pgfqpoint{3.072960in}{8.803634in}}%
\pgfpathclose%
\pgfpathmoveto{\pgfqpoint{3.098730in}{8.801819in}}%
\pgfpathlineto{\pgfqpoint{3.112317in}{8.801819in}}%
\pgfpathlineto{\pgfqpoint{3.112317in}{8.785471in}}%
\pgfpathlineto{\pgfqpoint{3.098730in}{8.785471in}}%
\pgfpathlineto{\pgfqpoint{3.098730in}{8.801819in}}%
\pgfpathlineto{\pgfqpoint{3.098730in}{8.801819in}}%
\pgfpathclose%
\pgfpathmoveto{\pgfqpoint{3.098730in}{8.749969in}}%
\pgfpathlineto{\pgfqpoint{3.112317in}{8.749969in}}%
\pgfpathlineto{\pgfqpoint{3.112317in}{8.738878in}}%
\pgfpathlineto{\pgfqpoint{3.101761in}{8.718261in}}%
\pgfpathlineto{\pgfqpoint{3.093453in}{8.718261in}}%
\pgfpathlineto{\pgfqpoint{3.098730in}{8.738878in}}%
\pgfpathlineto{\pgfqpoint{3.098730in}{8.749969in}}%
\pgfpathlineto{\pgfqpoint{3.098730in}{8.749969in}}%
\pgfpathclose%
\pgfpathmoveto{\pgfqpoint{3.223191in}{8.785471in}}%
\pgfpathquadraticcurveto{\pgfqpoint{3.232530in}{8.783471in}}{\pgfqpoint{3.237767in}{8.777142in}}%
\pgfpathquadraticcurveto{\pgfqpoint{3.243024in}{8.770833in}}{\pgfqpoint{3.243024in}{8.761556in}}%
\pgfpathquadraticcurveto{\pgfqpoint{3.243024in}{8.747330in}}{\pgfqpoint{3.233231in}{8.739517in}}%
\pgfpathquadraticcurveto{\pgfqpoint{3.223439in}{8.731724in}}{\pgfqpoint{3.205399in}{8.731724in}}%
\pgfpathquadraticcurveto{\pgfqpoint{3.199359in}{8.731724in}}{\pgfqpoint{3.192947in}{8.732920in}}%
\pgfpathquadraticcurveto{\pgfqpoint{3.186535in}{8.734115in}}{\pgfqpoint{3.179711in}{8.736507in}}%
\pgfpathlineto{\pgfqpoint{3.179711in}{8.749062in}}%
\pgfpathquadraticcurveto{\pgfqpoint{3.185113in}{8.745908in}}{\pgfqpoint{3.191545in}{8.744300in}}%
\pgfpathquadraticcurveto{\pgfqpoint{3.197998in}{8.742692in}}{\pgfqpoint{3.205028in}{8.742692in}}%
\pgfpathquadraticcurveto{\pgfqpoint{3.217254in}{8.742692in}}{\pgfqpoint{3.223665in}{8.747516in}}%
\pgfpathquadraticcurveto{\pgfqpoint{3.230077in}{8.752340in}}{\pgfqpoint{3.230077in}{8.761556in}}%
\pgfpathquadraticcurveto{\pgfqpoint{3.230077in}{8.770070in}}{\pgfqpoint{3.224119in}{8.774853in}}%
\pgfpathquadraticcurveto{\pgfqpoint{3.218161in}{8.779657in}}{\pgfqpoint{3.207543in}{8.779657in}}%
\pgfpathlineto{\pgfqpoint{3.196328in}{8.779657in}}%
\pgfpathlineto{\pgfqpoint{3.196328in}{8.790357in}}%
\pgfpathlineto{\pgfqpoint{3.208059in}{8.790357in}}%
\pgfpathquadraticcurveto{\pgfqpoint{3.217645in}{8.790357in}}{\pgfqpoint{3.222738in}{8.794191in}}%
\pgfpathquadraticcurveto{\pgfqpoint{3.227830in}{8.798026in}}{\pgfqpoint{3.227830in}{8.805242in}}%
\pgfpathquadraticcurveto{\pgfqpoint{3.227830in}{8.812643in}}{\pgfqpoint{3.222573in}{8.816601in}}%
\pgfpathquadraticcurveto{\pgfqpoint{3.217336in}{8.820580in}}{\pgfqpoint{3.207543in}{8.820580in}}%
\pgfpathquadraticcurveto{\pgfqpoint{3.202183in}{8.820580in}}{\pgfqpoint{3.196060in}{8.819405in}}%
\pgfpathquadraticcurveto{\pgfqpoint{3.189937in}{8.818251in}}{\pgfqpoint{3.182598in}{8.815818in}}%
\pgfpathlineto{\pgfqpoint{3.182598in}{8.827404in}}%
\pgfpathquadraticcurveto{\pgfqpoint{3.190020in}{8.829466in}}{\pgfqpoint{3.196493in}{8.830497in}}%
\pgfpathquadraticcurveto{\pgfqpoint{3.202967in}{8.831528in}}{\pgfqpoint{3.208698in}{8.831528in}}%
\pgfpathquadraticcurveto{\pgfqpoint{3.223521in}{8.831528in}}{\pgfqpoint{3.232139in}{8.824786in}}%
\pgfpathquadraticcurveto{\pgfqpoint{3.240777in}{8.818065in}}{\pgfqpoint{3.240777in}{8.806602in}}%
\pgfpathquadraticcurveto{\pgfqpoint{3.240777in}{8.798603in}}{\pgfqpoint{3.236200in}{8.793099in}}%
\pgfpathquadraticcurveto{\pgfqpoint{3.231623in}{8.787594in}}{\pgfqpoint{3.223191in}{8.785471in}}%
\pgfpathlineto{\pgfqpoint{3.223191in}{8.785471in}}%
\pgfpathclose%
\pgfpathmoveto{\pgfqpoint{3.269970in}{8.744547in}}%
\pgfpathlineto{\pgfqpoint{3.291225in}{8.744547in}}%
\pgfpathlineto{\pgfqpoint{3.291225in}{8.817941in}}%
\pgfpathlineto{\pgfqpoint{3.268094in}{8.813303in}}%
\pgfpathlineto{\pgfqpoint{3.268094in}{8.825157in}}%
\pgfpathlineto{\pgfqpoint{3.291101in}{8.829796in}}%
\pgfpathlineto{\pgfqpoint{3.304110in}{8.829796in}}%
\pgfpathlineto{\pgfqpoint{3.304110in}{8.744547in}}%
\pgfpathlineto{\pgfqpoint{3.325366in}{8.744547in}}%
\pgfpathlineto{\pgfqpoint{3.325366in}{8.733600in}}%
\pgfpathlineto{\pgfqpoint{3.269970in}{8.733600in}}%
\pgfpathlineto{\pgfqpoint{3.269970in}{8.744547in}}%
\pgfpathlineto{\pgfqpoint{3.269970in}{8.744547in}}%
\pgfpathclose%
\pgfpathmoveto{\pgfqpoint{3.443848in}{8.802995in}}%
\pgfpathlineto{\pgfqpoint{3.443848in}{8.791903in}}%
\pgfpathquadraticcurveto{\pgfqpoint{3.438817in}{8.794686in}}{\pgfqpoint{3.433766in}{8.796067in}}%
\pgfpathquadraticcurveto{\pgfqpoint{3.428715in}{8.797449in}}{\pgfqpoint{3.423561in}{8.797449in}}%
\pgfpathquadraticcurveto{\pgfqpoint{3.412016in}{8.797449in}}{\pgfqpoint{3.405625in}{8.790130in}}%
\pgfpathquadraticcurveto{\pgfqpoint{3.399255in}{8.782832in}}{\pgfqpoint{3.399255in}{8.769617in}}%
\pgfpathquadraticcurveto{\pgfqpoint{3.399255in}{8.756402in}}{\pgfqpoint{3.405625in}{8.749083in}}%
\pgfpathquadraticcurveto{\pgfqpoint{3.412016in}{8.741785in}}{\pgfqpoint{3.423561in}{8.741785in}}%
\pgfpathquadraticcurveto{\pgfqpoint{3.428715in}{8.741785in}}{\pgfqpoint{3.433766in}{8.743166in}}%
\pgfpathquadraticcurveto{\pgfqpoint{3.438817in}{8.744547in}}{\pgfqpoint{3.443848in}{8.747330in}}%
\pgfpathlineto{\pgfqpoint{3.443848in}{8.736363in}}%
\pgfpathquadraticcurveto{\pgfqpoint{3.438879in}{8.734054in}}{\pgfqpoint{3.433560in}{8.732899in}}%
\pgfpathquadraticcurveto{\pgfqpoint{3.428262in}{8.731724in}}{\pgfqpoint{3.422262in}{8.731724in}}%
\pgfpathquadraticcurveto{\pgfqpoint{3.405955in}{8.731724in}}{\pgfqpoint{3.396348in}{8.741970in}}%
\pgfpathquadraticcurveto{\pgfqpoint{3.386761in}{8.752217in}}{\pgfqpoint{3.386761in}{8.769617in}}%
\pgfpathquadraticcurveto{\pgfqpoint{3.386761in}{8.787264in}}{\pgfqpoint{3.396451in}{8.797366in}}%
\pgfpathquadraticcurveto{\pgfqpoint{3.406161in}{8.807489in}}{\pgfqpoint{3.423046in}{8.807489in}}%
\pgfpathquadraticcurveto{\pgfqpoint{3.428509in}{8.807489in}}{\pgfqpoint{3.433725in}{8.806355in}}%
\pgfpathquadraticcurveto{\pgfqpoint{3.438941in}{8.805242in}}{\pgfqpoint{3.443848in}{8.802995in}}%
\pgfpathlineto{\pgfqpoint{3.443848in}{8.802995in}}%
\pgfpathclose%
\pgfpathmoveto{\pgfqpoint{3.497265in}{8.769864in}}%
\pgfpathquadraticcurveto{\pgfqpoint{3.482895in}{8.769864in}}{\pgfqpoint{3.477349in}{8.766586in}}%
\pgfpathquadraticcurveto{\pgfqpoint{3.471803in}{8.763308in}}{\pgfqpoint{3.471803in}{8.755371in}}%
\pgfpathquadraticcurveto{\pgfqpoint{3.471803in}{8.749062in}}{\pgfqpoint{3.475968in}{8.745351in}}%
\pgfpathquadraticcurveto{\pgfqpoint{3.480132in}{8.741661in}}{\pgfqpoint{3.487266in}{8.741661in}}%
\pgfpathquadraticcurveto{\pgfqpoint{3.497141in}{8.741661in}}{\pgfqpoint{3.503099in}{8.748650in}}%
\pgfpathquadraticcurveto{\pgfqpoint{3.509057in}{8.755639in}}{\pgfqpoint{3.509057in}{8.767225in}}%
\pgfpathlineto{\pgfqpoint{3.509057in}{8.769864in}}%
\pgfpathlineto{\pgfqpoint{3.497265in}{8.769864in}}%
\pgfpathlineto{\pgfqpoint{3.497265in}{8.769864in}}%
\pgfpathclose%
\pgfpathmoveto{\pgfqpoint{3.520911in}{8.774771in}}%
\pgfpathlineto{\pgfqpoint{3.520911in}{8.733600in}}%
\pgfpathlineto{\pgfqpoint{3.509057in}{8.733600in}}%
\pgfpathlineto{\pgfqpoint{3.509057in}{8.744547in}}%
\pgfpathquadraticcurveto{\pgfqpoint{3.504996in}{8.737991in}}{\pgfqpoint{3.498934in}{8.734858in}}%
\pgfpathquadraticcurveto{\pgfqpoint{3.492873in}{8.731724in}}{\pgfqpoint{3.484111in}{8.731724in}}%
\pgfpathquadraticcurveto{\pgfqpoint{3.473040in}{8.731724in}}{\pgfqpoint{3.466484in}{8.737950in}}%
\pgfpathquadraticcurveto{\pgfqpoint{3.459949in}{8.744176in}}{\pgfqpoint{3.459949in}{8.754608in}}%
\pgfpathquadraticcurveto{\pgfqpoint{3.459949in}{8.766772in}}{\pgfqpoint{3.468092in}{8.772957in}}%
\pgfpathquadraticcurveto{\pgfqpoint{3.476257in}{8.779141in}}{\pgfqpoint{3.492420in}{8.779141in}}%
\pgfpathlineto{\pgfqpoint{3.509057in}{8.779141in}}%
\pgfpathlineto{\pgfqpoint{3.509057in}{8.780317in}}%
\pgfpathquadraticcurveto{\pgfqpoint{3.509057in}{8.788501in}}{\pgfqpoint{3.503676in}{8.792975in}}%
\pgfpathquadraticcurveto{\pgfqpoint{3.498295in}{8.797449in}}{\pgfqpoint{3.488564in}{8.797449in}}%
\pgfpathquadraticcurveto{\pgfqpoint{3.482380in}{8.797449in}}{\pgfqpoint{3.476504in}{8.795964in}}%
\pgfpathquadraticcurveto{\pgfqpoint{3.470649in}{8.794480in}}{\pgfqpoint{3.465247in}{8.791511in}}%
\pgfpathlineto{\pgfqpoint{3.465247in}{8.802479in}}%
\pgfpathquadraticcurveto{\pgfqpoint{3.471742in}{8.804994in}}{\pgfqpoint{3.477865in}{8.806231in}}%
\pgfpathquadraticcurveto{\pgfqpoint{3.483988in}{8.807489in}}{\pgfqpoint{3.489781in}{8.807489in}}%
\pgfpathquadraticcurveto{\pgfqpoint{3.505449in}{8.807489in}}{\pgfqpoint{3.513180in}{8.799366in}}%
\pgfpathquadraticcurveto{\pgfqpoint{3.520911in}{8.791264in}}{\pgfqpoint{3.520911in}{8.774771in}}%
\pgfpathlineto{\pgfqpoint{3.520911in}{8.774771in}}%
\pgfpathclose%
\pgfpathmoveto{\pgfqpoint{3.556784in}{8.744424in}}%
\pgfpathlineto{\pgfqpoint{3.556784in}{8.706160in}}%
\pgfpathlineto{\pgfqpoint{3.544868in}{8.706160in}}%
\pgfpathlineto{\pgfqpoint{3.544868in}{8.805757in}}%
\pgfpathlineto{\pgfqpoint{3.556784in}{8.805757in}}%
\pgfpathlineto{\pgfqpoint{3.556784in}{8.794810in}}%
\pgfpathquadraticcurveto{\pgfqpoint{3.560536in}{8.801242in}}{\pgfqpoint{3.566226in}{8.804355in}}%
\pgfpathquadraticcurveto{\pgfqpoint{3.571937in}{8.807489in}}{\pgfqpoint{3.579854in}{8.807489in}}%
\pgfpathquadraticcurveto{\pgfqpoint{3.593007in}{8.807489in}}{\pgfqpoint{3.601212in}{8.797057in}}%
\pgfpathquadraticcurveto{\pgfqpoint{3.609438in}{8.786625in}}{\pgfqpoint{3.609438in}{8.769617in}}%
\pgfpathquadraticcurveto{\pgfqpoint{3.609438in}{8.752608in}}{\pgfqpoint{3.601212in}{8.742156in}}%
\pgfpathquadraticcurveto{\pgfqpoint{3.593007in}{8.731724in}}{\pgfqpoint{3.579854in}{8.731724in}}%
\pgfpathquadraticcurveto{\pgfqpoint{3.571937in}{8.731724in}}{\pgfqpoint{3.566226in}{8.734858in}}%
\pgfpathquadraticcurveto{\pgfqpoint{3.560536in}{8.737991in}}{\pgfqpoint{3.556784in}{8.744424in}}%
\pgfpathlineto{\pgfqpoint{3.556784in}{8.744424in}}%
\pgfpathclose%
\pgfpathmoveto{\pgfqpoint{3.597130in}{8.769617in}}%
\pgfpathquadraticcurveto{\pgfqpoint{3.597130in}{8.782687in}}{\pgfqpoint{3.591749in}{8.790130in}}%
\pgfpathquadraticcurveto{\pgfqpoint{3.586368in}{8.797572in}}{\pgfqpoint{3.576967in}{8.797572in}}%
\pgfpathquadraticcurveto{\pgfqpoint{3.567546in}{8.797572in}}{\pgfqpoint{3.562165in}{8.790130in}}%
\pgfpathquadraticcurveto{\pgfqpoint{3.556784in}{8.782687in}}{\pgfqpoint{3.556784in}{8.769617in}}%
\pgfpathquadraticcurveto{\pgfqpoint{3.556784in}{8.756546in}}{\pgfqpoint{3.562165in}{8.749103in}}%
\pgfpathquadraticcurveto{\pgfqpoint{3.567546in}{8.741661in}}{\pgfqpoint{3.576967in}{8.741661in}}%
\pgfpathquadraticcurveto{\pgfqpoint{3.586368in}{8.741661in}}{\pgfqpoint{3.591749in}{8.749103in}}%
\pgfpathquadraticcurveto{\pgfqpoint{3.597130in}{8.756546in}}{\pgfqpoint{3.597130in}{8.769617in}}%
\pgfpathlineto{\pgfqpoint{3.597130in}{8.769617in}}%
\pgfpathclose%
\pgfpathmoveto{\pgfqpoint{3.640548in}{8.744424in}}%
\pgfpathlineto{\pgfqpoint{3.640548in}{8.706160in}}%
\pgfpathlineto{\pgfqpoint{3.628632in}{8.706160in}}%
\pgfpathlineto{\pgfqpoint{3.628632in}{8.805757in}}%
\pgfpathlineto{\pgfqpoint{3.640548in}{8.805757in}}%
\pgfpathlineto{\pgfqpoint{3.640548in}{8.794810in}}%
\pgfpathquadraticcurveto{\pgfqpoint{3.644300in}{8.801242in}}{\pgfqpoint{3.649990in}{8.804355in}}%
\pgfpathquadraticcurveto{\pgfqpoint{3.655701in}{8.807489in}}{\pgfqpoint{3.663618in}{8.807489in}}%
\pgfpathquadraticcurveto{\pgfqpoint{3.676771in}{8.807489in}}{\pgfqpoint{3.684976in}{8.797057in}}%
\pgfpathquadraticcurveto{\pgfqpoint{3.693202in}{8.786625in}}{\pgfqpoint{3.693202in}{8.769617in}}%
\pgfpathquadraticcurveto{\pgfqpoint{3.693202in}{8.752608in}}{\pgfqpoint{3.684976in}{8.742156in}}%
\pgfpathquadraticcurveto{\pgfqpoint{3.676771in}{8.731724in}}{\pgfqpoint{3.663618in}{8.731724in}}%
\pgfpathquadraticcurveto{\pgfqpoint{3.655701in}{8.731724in}}{\pgfqpoint{3.649990in}{8.734858in}}%
\pgfpathquadraticcurveto{\pgfqpoint{3.644300in}{8.737991in}}{\pgfqpoint{3.640548in}{8.744424in}}%
\pgfpathlineto{\pgfqpoint{3.640548in}{8.744424in}}%
\pgfpathclose%
\pgfpathmoveto{\pgfqpoint{3.680894in}{8.769617in}}%
\pgfpathquadraticcurveto{\pgfqpoint{3.680894in}{8.782687in}}{\pgfqpoint{3.675513in}{8.790130in}}%
\pgfpathquadraticcurveto{\pgfqpoint{3.670132in}{8.797572in}}{\pgfqpoint{3.660731in}{8.797572in}}%
\pgfpathquadraticcurveto{\pgfqpoint{3.651310in}{8.797572in}}{\pgfqpoint{3.645929in}{8.790130in}}%
\pgfpathquadraticcurveto{\pgfqpoint{3.640548in}{8.782687in}}{\pgfqpoint{3.640548in}{8.769617in}}%
\pgfpathquadraticcurveto{\pgfqpoint{3.640548in}{8.756546in}}{\pgfqpoint{3.645929in}{8.749103in}}%
\pgfpathquadraticcurveto{\pgfqpoint{3.651310in}{8.741661in}}{\pgfqpoint{3.660731in}{8.741661in}}%
\pgfpathquadraticcurveto{\pgfqpoint{3.670132in}{8.741661in}}{\pgfqpoint{3.675513in}{8.749103in}}%
\pgfpathquadraticcurveto{\pgfqpoint{3.680894in}{8.756546in}}{\pgfqpoint{3.680894in}{8.769617in}}%
\pgfpathlineto{\pgfqpoint{3.680894in}{8.769617in}}%
\pgfpathclose%
\pgfpathmoveto{\pgfqpoint{3.774575in}{8.772647in}}%
\pgfpathlineto{\pgfqpoint{3.774575in}{8.766854in}}%
\pgfpathlineto{\pgfqpoint{3.720065in}{8.766854in}}%
\pgfpathquadraticcurveto{\pgfqpoint{3.720849in}{8.754608in}}{\pgfqpoint{3.727446in}{8.748196in}}%
\pgfpathquadraticcurveto{\pgfqpoint{3.734043in}{8.741785in}}{\pgfqpoint{3.745836in}{8.741785in}}%
\pgfpathquadraticcurveto{\pgfqpoint{3.752660in}{8.741785in}}{\pgfqpoint{3.759071in}{8.743455in}}%
\pgfpathquadraticcurveto{\pgfqpoint{3.765483in}{8.745125in}}{\pgfqpoint{3.771812in}{8.748485in}}%
\pgfpathlineto{\pgfqpoint{3.771812in}{8.737270in}}%
\pgfpathquadraticcurveto{\pgfqpoint{3.765421in}{8.734569in}}{\pgfqpoint{3.758721in}{8.733146in}}%
\pgfpathquadraticcurveto{\pgfqpoint{3.752020in}{8.731724in}}{\pgfqpoint{3.745135in}{8.731724in}}%
\pgfpathquadraticcurveto{\pgfqpoint{3.727858in}{8.731724in}}{\pgfqpoint{3.717777in}{8.741764in}}%
\pgfpathquadraticcurveto{\pgfqpoint{3.707695in}{8.751825in}}{\pgfqpoint{3.707695in}{8.768978in}}%
\pgfpathquadraticcurveto{\pgfqpoint{3.707695in}{8.786687in}}{\pgfqpoint{3.717261in}{8.797078in}}%
\pgfpathquadraticcurveto{\pgfqpoint{3.726827in}{8.807489in}}{\pgfqpoint{3.743073in}{8.807489in}}%
\pgfpathquadraticcurveto{\pgfqpoint{3.757628in}{8.807489in}}{\pgfqpoint{3.766101in}{8.798108in}}%
\pgfpathquadraticcurveto{\pgfqpoint{3.774575in}{8.788749in}}{\pgfqpoint{3.774575in}{8.772647in}}%
\pgfpathlineto{\pgfqpoint{3.774575in}{8.772647in}}%
\pgfpathclose%
\pgfpathmoveto{\pgfqpoint{3.762720in}{8.776131in}}%
\pgfpathquadraticcurveto{\pgfqpoint{3.762597in}{8.785842in}}{\pgfqpoint{3.757278in}{8.791635in}}%
\pgfpathquadraticcurveto{\pgfqpoint{3.751959in}{8.797449in}}{\pgfqpoint{3.743197in}{8.797449in}}%
\pgfpathquadraticcurveto{\pgfqpoint{3.733280in}{8.797449in}}{\pgfqpoint{3.727322in}{8.791841in}}%
\pgfpathquadraticcurveto{\pgfqpoint{3.721364in}{8.786233in}}{\pgfqpoint{3.720457in}{8.776049in}}%
\pgfpathlineto{\pgfqpoint{3.762720in}{8.776131in}}%
\pgfpathlineto{\pgfqpoint{3.762720in}{8.776131in}}%
\pgfpathclose%
\pgfpathmoveto{\pgfqpoint{3.841516in}{8.794810in}}%
\pgfpathlineto{\pgfqpoint{3.841516in}{8.833857in}}%
\pgfpathlineto{\pgfqpoint{3.853370in}{8.833857in}}%
\pgfpathlineto{\pgfqpoint{3.853370in}{8.733600in}}%
\pgfpathlineto{\pgfqpoint{3.841516in}{8.733600in}}%
\pgfpathlineto{\pgfqpoint{3.841516in}{8.744424in}}%
\pgfpathquadraticcurveto{\pgfqpoint{3.837784in}{8.737991in}}{\pgfqpoint{3.832074in}{8.734858in}}%
\pgfpathquadraticcurveto{\pgfqpoint{3.826383in}{8.731724in}}{\pgfqpoint{3.818384in}{8.731724in}}%
\pgfpathquadraticcurveto{\pgfqpoint{3.805314in}{8.731724in}}{\pgfqpoint{3.797088in}{8.742156in}}%
\pgfpathquadraticcurveto{\pgfqpoint{3.788882in}{8.752608in}}{\pgfqpoint{3.788882in}{8.769617in}}%
\pgfpathquadraticcurveto{\pgfqpoint{3.788882in}{8.786625in}}{\pgfqpoint{3.797088in}{8.797057in}}%
\pgfpathquadraticcurveto{\pgfqpoint{3.805314in}{8.807489in}}{\pgfqpoint{3.818384in}{8.807489in}}%
\pgfpathquadraticcurveto{\pgfqpoint{3.826383in}{8.807489in}}{\pgfqpoint{3.832074in}{8.804355in}}%
\pgfpathquadraticcurveto{\pgfqpoint{3.837784in}{8.801242in}}{\pgfqpoint{3.841516in}{8.794810in}}%
\pgfpathlineto{\pgfqpoint{3.841516in}{8.794810in}}%
\pgfpathclose%
\pgfpathmoveto{\pgfqpoint{3.801128in}{8.769617in}}%
\pgfpathquadraticcurveto{\pgfqpoint{3.801128in}{8.756546in}}{\pgfqpoint{3.806509in}{8.749103in}}%
\pgfpathquadraticcurveto{\pgfqpoint{3.811890in}{8.741661in}}{\pgfqpoint{3.821291in}{8.741661in}}%
\pgfpathquadraticcurveto{\pgfqpoint{3.830692in}{8.741661in}}{\pgfqpoint{3.836094in}{8.749103in}}%
\pgfpathquadraticcurveto{\pgfqpoint{3.841516in}{8.756546in}}{\pgfqpoint{3.841516in}{8.769617in}}%
\pgfpathquadraticcurveto{\pgfqpoint{3.841516in}{8.782687in}}{\pgfqpoint{3.836094in}{8.790130in}}%
\pgfpathquadraticcurveto{\pgfqpoint{3.830692in}{8.797572in}}{\pgfqpoint{3.821291in}{8.797572in}}%
\pgfpathquadraticcurveto{\pgfqpoint{3.811890in}{8.797572in}}{\pgfqpoint{3.806509in}{8.790130in}}%
\pgfpathquadraticcurveto{\pgfqpoint{3.801128in}{8.782687in}}{\pgfqpoint{3.801128in}{8.769617in}}%
\pgfpathlineto{\pgfqpoint{3.801128in}{8.769617in}}%
\pgfpathclose%
\pgfpathmoveto{\pgfqpoint{3.948205in}{8.833713in}}%
\pgfpathquadraticcurveto{\pgfqpoint{3.939588in}{8.818910in}}{\pgfqpoint{3.935382in}{8.804396in}}%
\pgfpathquadraticcurveto{\pgfqpoint{3.931197in}{8.789903in}}{\pgfqpoint{3.931197in}{8.775018in}}%
\pgfpathquadraticcurveto{\pgfqpoint{3.931197in}{8.760154in}}{\pgfqpoint{3.935423in}{8.745557in}}%
\pgfpathquadraticcurveto{\pgfqpoint{3.939650in}{8.730961in}}{\pgfqpoint{3.948205in}{8.716200in}}%
\pgfpathlineto{\pgfqpoint{3.937897in}{8.716200in}}%
\pgfpathquadraticcurveto{\pgfqpoint{3.928249in}{8.731353in}}{\pgfqpoint{3.923445in}{8.745970in}}%
\pgfpathquadraticcurveto{\pgfqpoint{3.918642in}{8.760587in}}{\pgfqpoint{3.918642in}{8.775018in}}%
\pgfpathquadraticcurveto{\pgfqpoint{3.918642in}{8.789388in}}{\pgfqpoint{3.923404in}{8.803943in}}%
\pgfpathquadraticcurveto{\pgfqpoint{3.928187in}{8.818519in}}{\pgfqpoint{3.937897in}{8.833713in}}%
\pgfpathlineto{\pgfqpoint{3.948205in}{8.833713in}}%
\pgfpathlineto{\pgfqpoint{3.948205in}{8.833713in}}%
\pgfpathclose%
\pgfpathmoveto{\pgfqpoint{4.018693in}{8.794810in}}%
\pgfpathlineto{\pgfqpoint{4.018693in}{8.833857in}}%
\pgfpathlineto{\pgfqpoint{4.030547in}{8.833857in}}%
\pgfpathlineto{\pgfqpoint{4.030547in}{8.733600in}}%
\pgfpathlineto{\pgfqpoint{4.018693in}{8.733600in}}%
\pgfpathlineto{\pgfqpoint{4.018693in}{8.744424in}}%
\pgfpathquadraticcurveto{\pgfqpoint{4.014961in}{8.737991in}}{\pgfqpoint{4.009250in}{8.734858in}}%
\pgfpathquadraticcurveto{\pgfqpoint{4.003560in}{8.731724in}}{\pgfqpoint{3.995561in}{8.731724in}}%
\pgfpathquadraticcurveto{\pgfqpoint{3.982490in}{8.731724in}}{\pgfqpoint{3.974264in}{8.742156in}}%
\pgfpathquadraticcurveto{\pgfqpoint{3.966059in}{8.752608in}}{\pgfqpoint{3.966059in}{8.769617in}}%
\pgfpathquadraticcurveto{\pgfqpoint{3.966059in}{8.786625in}}{\pgfqpoint{3.974264in}{8.797057in}}%
\pgfpathquadraticcurveto{\pgfqpoint{3.982490in}{8.807489in}}{\pgfqpoint{3.995561in}{8.807489in}}%
\pgfpathquadraticcurveto{\pgfqpoint{4.003560in}{8.807489in}}{\pgfqpoint{4.009250in}{8.804355in}}%
\pgfpathquadraticcurveto{\pgfqpoint{4.014961in}{8.801242in}}{\pgfqpoint{4.018693in}{8.794810in}}%
\pgfpathlineto{\pgfqpoint{4.018693in}{8.794810in}}%
\pgfpathclose%
\pgfpathmoveto{\pgfqpoint{3.978305in}{8.769617in}}%
\pgfpathquadraticcurveto{\pgfqpoint{3.978305in}{8.756546in}}{\pgfqpoint{3.983686in}{8.749103in}}%
\pgfpathquadraticcurveto{\pgfqpoint{3.989067in}{8.741661in}}{\pgfqpoint{3.998468in}{8.741661in}}%
\pgfpathquadraticcurveto{\pgfqpoint{4.007869in}{8.741661in}}{\pgfqpoint{4.013270in}{8.749103in}}%
\pgfpathquadraticcurveto{\pgfqpoint{4.018693in}{8.756546in}}{\pgfqpoint{4.018693in}{8.769617in}}%
\pgfpathquadraticcurveto{\pgfqpoint{4.018693in}{8.782687in}}{\pgfqpoint{4.013270in}{8.790130in}}%
\pgfpathquadraticcurveto{\pgfqpoint{4.007869in}{8.797572in}}{\pgfqpoint{3.998468in}{8.797572in}}%
\pgfpathquadraticcurveto{\pgfqpoint{3.989067in}{8.797572in}}{\pgfqpoint{3.983686in}{8.790130in}}%
\pgfpathquadraticcurveto{\pgfqpoint{3.978305in}{8.782687in}}{\pgfqpoint{3.978305in}{8.769617in}}%
\pgfpathlineto{\pgfqpoint{3.978305in}{8.769617in}}%
\pgfpathclose%
\pgfpathmoveto{\pgfqpoint{4.054977in}{8.805757in}}%
\pgfpathlineto{\pgfqpoint{4.066832in}{8.805757in}}%
\pgfpathlineto{\pgfqpoint{4.066832in}{8.733600in}}%
\pgfpathlineto{\pgfqpoint{4.054977in}{8.733600in}}%
\pgfpathlineto{\pgfqpoint{4.054977in}{8.805757in}}%
\pgfpathlineto{\pgfqpoint{4.054977in}{8.805757in}}%
\pgfpathclose%
\pgfpathmoveto{\pgfqpoint{4.054977in}{8.833857in}}%
\pgfpathlineto{\pgfqpoint{4.066832in}{8.833857in}}%
\pgfpathlineto{\pgfqpoint{4.066832in}{8.818828in}}%
\pgfpathlineto{\pgfqpoint{4.054977in}{8.818828in}}%
\pgfpathlineto{\pgfqpoint{4.054977in}{8.833857in}}%
\pgfpathlineto{\pgfqpoint{4.054977in}{8.833857in}}%
\pgfpathclose%
\pgfpathmoveto{\pgfqpoint{4.124434in}{8.769864in}}%
\pgfpathquadraticcurveto{\pgfqpoint{4.110064in}{8.769864in}}{\pgfqpoint{4.104518in}{8.766586in}}%
\pgfpathquadraticcurveto{\pgfqpoint{4.098972in}{8.763308in}}{\pgfqpoint{4.098972in}{8.755371in}}%
\pgfpathquadraticcurveto{\pgfqpoint{4.098972in}{8.749062in}}{\pgfqpoint{4.103137in}{8.745351in}}%
\pgfpathquadraticcurveto{\pgfqpoint{4.107301in}{8.741661in}}{\pgfqpoint{4.114435in}{8.741661in}}%
\pgfpathquadraticcurveto{\pgfqpoint{4.124310in}{8.741661in}}{\pgfqpoint{4.130268in}{8.748650in}}%
\pgfpathquadraticcurveto{\pgfqpoint{4.136226in}{8.755639in}}{\pgfqpoint{4.136226in}{8.767225in}}%
\pgfpathlineto{\pgfqpoint{4.136226in}{8.769864in}}%
\pgfpathlineto{\pgfqpoint{4.124434in}{8.769864in}}%
\pgfpathlineto{\pgfqpoint{4.124434in}{8.769864in}}%
\pgfpathclose%
\pgfpathmoveto{\pgfqpoint{4.148081in}{8.774771in}}%
\pgfpathlineto{\pgfqpoint{4.148081in}{8.733600in}}%
\pgfpathlineto{\pgfqpoint{4.136226in}{8.733600in}}%
\pgfpathlineto{\pgfqpoint{4.136226in}{8.744547in}}%
\pgfpathquadraticcurveto{\pgfqpoint{4.132165in}{8.737991in}}{\pgfqpoint{4.126104in}{8.734858in}}%
\pgfpathquadraticcurveto{\pgfqpoint{4.120042in}{8.731724in}}{\pgfqpoint{4.111280in}{8.731724in}}%
\pgfpathquadraticcurveto{\pgfqpoint{4.100209in}{8.731724in}}{\pgfqpoint{4.093653in}{8.737950in}}%
\pgfpathquadraticcurveto{\pgfqpoint{4.087118in}{8.744176in}}{\pgfqpoint{4.087118in}{8.754608in}}%
\pgfpathquadraticcurveto{\pgfqpoint{4.087118in}{8.766772in}}{\pgfqpoint{4.095262in}{8.772957in}}%
\pgfpathquadraticcurveto{\pgfqpoint{4.103426in}{8.779141in}}{\pgfqpoint{4.119589in}{8.779141in}}%
\pgfpathlineto{\pgfqpoint{4.136226in}{8.779141in}}%
\pgfpathlineto{\pgfqpoint{4.136226in}{8.780317in}}%
\pgfpathquadraticcurveto{\pgfqpoint{4.136226in}{8.788501in}}{\pgfqpoint{4.130845in}{8.792975in}}%
\pgfpathquadraticcurveto{\pgfqpoint{4.125464in}{8.797449in}}{\pgfqpoint{4.115734in}{8.797449in}}%
\pgfpathquadraticcurveto{\pgfqpoint{4.109549in}{8.797449in}}{\pgfqpoint{4.103673in}{8.795964in}}%
\pgfpathquadraticcurveto{\pgfqpoint{4.097818in}{8.794480in}}{\pgfqpoint{4.092416in}{8.791511in}}%
\pgfpathlineto{\pgfqpoint{4.092416in}{8.802479in}}%
\pgfpathquadraticcurveto{\pgfqpoint{4.098911in}{8.804994in}}{\pgfqpoint{4.105034in}{8.806231in}}%
\pgfpathquadraticcurveto{\pgfqpoint{4.111157in}{8.807489in}}{\pgfqpoint{4.116950in}{8.807489in}}%
\pgfpathquadraticcurveto{\pgfqpoint{4.132618in}{8.807489in}}{\pgfqpoint{4.140349in}{8.799366in}}%
\pgfpathquadraticcurveto{\pgfqpoint{4.148081in}{8.791264in}}{\pgfqpoint{4.148081in}{8.774771in}}%
\pgfpathlineto{\pgfqpoint{4.148081in}{8.774771in}}%
\pgfpathclose%
\pgfpathmoveto{\pgfqpoint{4.228670in}{8.791903in}}%
\pgfpathquadraticcurveto{\pgfqpoint{4.233123in}{8.799902in}}{\pgfqpoint{4.239308in}{8.803695in}}%
\pgfpathquadraticcurveto{\pgfqpoint{4.245493in}{8.807489in}}{\pgfqpoint{4.253863in}{8.807489in}}%
\pgfpathquadraticcurveto{\pgfqpoint{4.265140in}{8.807489in}}{\pgfqpoint{4.271263in}{8.799593in}}%
\pgfpathquadraticcurveto{\pgfqpoint{4.277386in}{8.791717in}}{\pgfqpoint{4.277386in}{8.777162in}}%
\pgfpathlineto{\pgfqpoint{4.277386in}{8.733600in}}%
\pgfpathlineto{\pgfqpoint{4.265470in}{8.733600in}}%
\pgfpathlineto{\pgfqpoint{4.265470in}{8.776771in}}%
\pgfpathquadraticcurveto{\pgfqpoint{4.265470in}{8.787141in}}{\pgfqpoint{4.261780in}{8.792150in}}%
\pgfpathquadraticcurveto{\pgfqpoint{4.258110in}{8.797181in}}{\pgfqpoint{4.250585in}{8.797181in}}%
\pgfpathquadraticcurveto{\pgfqpoint{4.241369in}{8.797181in}}{\pgfqpoint{4.236009in}{8.791058in}}%
\pgfpathquadraticcurveto{\pgfqpoint{4.230670in}{8.784955in}}{\pgfqpoint{4.230670in}{8.774379in}}%
\pgfpathlineto{\pgfqpoint{4.230670in}{8.733600in}}%
\pgfpathlineto{\pgfqpoint{4.218753in}{8.733600in}}%
\pgfpathlineto{\pgfqpoint{4.218753in}{8.776771in}}%
\pgfpathquadraticcurveto{\pgfqpoint{4.218753in}{8.787202in}}{\pgfqpoint{4.215084in}{8.792192in}}%
\pgfpathquadraticcurveto{\pgfqpoint{4.211414in}{8.797181in}}{\pgfqpoint{4.203745in}{8.797181in}}%
\pgfpathquadraticcurveto{\pgfqpoint{4.194653in}{8.797181in}}{\pgfqpoint{4.189293in}{8.791037in}}%
\pgfpathquadraticcurveto{\pgfqpoint{4.183953in}{8.784893in}}{\pgfqpoint{4.183953in}{8.774379in}}%
\pgfpathlineto{\pgfqpoint{4.183953in}{8.733600in}}%
\pgfpathlineto{\pgfqpoint{4.172037in}{8.733600in}}%
\pgfpathlineto{\pgfqpoint{4.172037in}{8.805757in}}%
\pgfpathlineto{\pgfqpoint{4.183953in}{8.805757in}}%
\pgfpathlineto{\pgfqpoint{4.183953in}{8.794542in}}%
\pgfpathquadraticcurveto{\pgfqpoint{4.188014in}{8.801180in}}{\pgfqpoint{4.193684in}{8.804335in}}%
\pgfpathquadraticcurveto{\pgfqpoint{4.199353in}{8.807489in}}{\pgfqpoint{4.207146in}{8.807489in}}%
\pgfpathquadraticcurveto{\pgfqpoint{4.215022in}{8.807489in}}{\pgfqpoint{4.220526in}{8.803489in}}%
\pgfpathquadraticcurveto{\pgfqpoint{4.226031in}{8.799510in}}{\pgfqpoint{4.228670in}{8.791903in}}%
\pgfpathlineto{\pgfqpoint{4.228670in}{8.791903in}}%
\pgfpathclose%
\pgfpathmoveto{\pgfqpoint{4.328968in}{8.797449in}}%
\pgfpathquadraticcurveto{\pgfqpoint{4.319443in}{8.797449in}}{\pgfqpoint{4.313898in}{8.790006in}}%
\pgfpathquadraticcurveto{\pgfqpoint{4.308352in}{8.782564in}}{\pgfqpoint{4.308352in}{8.769617in}}%
\pgfpathquadraticcurveto{\pgfqpoint{4.308352in}{8.756670in}}{\pgfqpoint{4.313856in}{8.749227in}}%
\pgfpathquadraticcurveto{\pgfqpoint{4.319382in}{8.741785in}}{\pgfqpoint{4.328968in}{8.741785in}}%
\pgfpathquadraticcurveto{\pgfqpoint{4.338452in}{8.741785in}}{\pgfqpoint{4.343977in}{8.749248in}}%
\pgfpathquadraticcurveto{\pgfqpoint{4.349523in}{8.756732in}}{\pgfqpoint{4.349523in}{8.769617in}}%
\pgfpathquadraticcurveto{\pgfqpoint{4.349523in}{8.782440in}}{\pgfqpoint{4.343977in}{8.789944in}}%
\pgfpathquadraticcurveto{\pgfqpoint{4.338452in}{8.797449in}}{\pgfqpoint{4.328968in}{8.797449in}}%
\pgfpathlineto{\pgfqpoint{4.328968in}{8.797449in}}%
\pgfpathclose%
\pgfpathmoveto{\pgfqpoint{4.328968in}{8.807489in}}%
\pgfpathquadraticcurveto{\pgfqpoint{4.344430in}{8.807489in}}{\pgfqpoint{4.353254in}{8.797428in}}%
\pgfpathquadraticcurveto{\pgfqpoint{4.362099in}{8.787388in}}{\pgfqpoint{4.362099in}{8.769617in}}%
\pgfpathquadraticcurveto{\pgfqpoint{4.362099in}{8.751907in}}{\pgfqpoint{4.353254in}{8.741805in}}%
\pgfpathquadraticcurveto{\pgfqpoint{4.344430in}{8.731724in}}{\pgfqpoint{4.328968in}{8.731724in}}%
\pgfpathquadraticcurveto{\pgfqpoint{4.313444in}{8.731724in}}{\pgfqpoint{4.304641in}{8.741805in}}%
\pgfpathquadraticcurveto{\pgfqpoint{4.295858in}{8.751907in}}{\pgfqpoint{4.295858in}{8.769617in}}%
\pgfpathquadraticcurveto{\pgfqpoint{4.295858in}{8.787388in}}{\pgfqpoint{4.304641in}{8.797428in}}%
\pgfpathquadraticcurveto{\pgfqpoint{4.313444in}{8.807489in}}{\pgfqpoint{4.328968in}{8.807489in}}%
\pgfpathlineto{\pgfqpoint{4.328968in}{8.807489in}}%
\pgfpathclose%
\pgfpathmoveto{\pgfqpoint{4.441739in}{8.777162in}}%
\pgfpathlineto{\pgfqpoint{4.441739in}{8.733600in}}%
\pgfpathlineto{\pgfqpoint{4.429885in}{8.733600in}}%
\pgfpathlineto{\pgfqpoint{4.429885in}{8.776771in}}%
\pgfpathquadraticcurveto{\pgfqpoint{4.429885in}{8.787017in}}{\pgfqpoint{4.425885in}{8.792088in}}%
\pgfpathquadraticcurveto{\pgfqpoint{4.421886in}{8.797181in}}{\pgfqpoint{4.413907in}{8.797181in}}%
\pgfpathquadraticcurveto{\pgfqpoint{4.404300in}{8.797181in}}{\pgfqpoint{4.398754in}{8.791058in}}%
\pgfpathquadraticcurveto{\pgfqpoint{4.393209in}{8.784955in}}{\pgfqpoint{4.393209in}{8.774379in}}%
\pgfpathlineto{\pgfqpoint{4.393209in}{8.733600in}}%
\pgfpathlineto{\pgfqpoint{4.381292in}{8.733600in}}%
\pgfpathlineto{\pgfqpoint{4.381292in}{8.805757in}}%
\pgfpathlineto{\pgfqpoint{4.393209in}{8.805757in}}%
\pgfpathlineto{\pgfqpoint{4.393209in}{8.794542in}}%
\pgfpathquadraticcurveto{\pgfqpoint{4.397476in}{8.801057in}}{\pgfqpoint{4.403228in}{8.804273in}}%
\pgfpathquadraticcurveto{\pgfqpoint{4.409001in}{8.807489in}}{\pgfqpoint{4.416546in}{8.807489in}}%
\pgfpathquadraticcurveto{\pgfqpoint{4.428978in}{8.807489in}}{\pgfqpoint{4.435348in}{8.799799in}}%
\pgfpathquadraticcurveto{\pgfqpoint{4.441739in}{8.792109in}}{\pgfqpoint{4.441739in}{8.777162in}}%
\pgfpathlineto{\pgfqpoint{4.441739in}{8.777162in}}%
\pgfpathclose%
\pgfpathmoveto{\pgfqpoint{4.512845in}{8.794810in}}%
\pgfpathlineto{\pgfqpoint{4.512845in}{8.833857in}}%
\pgfpathlineto{\pgfqpoint{4.524699in}{8.833857in}}%
\pgfpathlineto{\pgfqpoint{4.524699in}{8.733600in}}%
\pgfpathlineto{\pgfqpoint{4.512845in}{8.733600in}}%
\pgfpathlineto{\pgfqpoint{4.512845in}{8.744424in}}%
\pgfpathquadraticcurveto{\pgfqpoint{4.509114in}{8.737991in}}{\pgfqpoint{4.503403in}{8.734858in}}%
\pgfpathquadraticcurveto{\pgfqpoint{4.497713in}{8.731724in}}{\pgfqpoint{4.489714in}{8.731724in}}%
\pgfpathquadraticcurveto{\pgfqpoint{4.476643in}{8.731724in}}{\pgfqpoint{4.468417in}{8.742156in}}%
\pgfpathquadraticcurveto{\pgfqpoint{4.460212in}{8.752608in}}{\pgfqpoint{4.460212in}{8.769617in}}%
\pgfpathquadraticcurveto{\pgfqpoint{4.460212in}{8.786625in}}{\pgfqpoint{4.468417in}{8.797057in}}%
\pgfpathquadraticcurveto{\pgfqpoint{4.476643in}{8.807489in}}{\pgfqpoint{4.489714in}{8.807489in}}%
\pgfpathquadraticcurveto{\pgfqpoint{4.497713in}{8.807489in}}{\pgfqpoint{4.503403in}{8.804355in}}%
\pgfpathquadraticcurveto{\pgfqpoint{4.509114in}{8.801242in}}{\pgfqpoint{4.512845in}{8.794810in}}%
\pgfpathlineto{\pgfqpoint{4.512845in}{8.794810in}}%
\pgfpathclose%
\pgfpathmoveto{\pgfqpoint{4.472458in}{8.769617in}}%
\pgfpathquadraticcurveto{\pgfqpoint{4.472458in}{8.756546in}}{\pgfqpoint{4.477839in}{8.749103in}}%
\pgfpathquadraticcurveto{\pgfqpoint{4.483219in}{8.741661in}}{\pgfqpoint{4.492620in}{8.741661in}}%
\pgfpathquadraticcurveto{\pgfqpoint{4.502022in}{8.741661in}}{\pgfqpoint{4.507423in}{8.749103in}}%
\pgfpathquadraticcurveto{\pgfqpoint{4.512845in}{8.756546in}}{\pgfqpoint{4.512845in}{8.769617in}}%
\pgfpathquadraticcurveto{\pgfqpoint{4.512845in}{8.782687in}}{\pgfqpoint{4.507423in}{8.790130in}}%
\pgfpathquadraticcurveto{\pgfqpoint{4.502022in}{8.797572in}}{\pgfqpoint{4.492620in}{8.797572in}}%
\pgfpathquadraticcurveto{\pgfqpoint{4.483219in}{8.797572in}}{\pgfqpoint{4.477839in}{8.790130in}}%
\pgfpathquadraticcurveto{\pgfqpoint{4.472458in}{8.782687in}}{\pgfqpoint{4.472458in}{8.769617in}}%
\pgfpathlineto{\pgfqpoint{4.472458in}{8.769617in}}%
\pgfpathclose%
\pgfpathmoveto{\pgfqpoint{4.595125in}{8.803634in}}%
\pgfpathlineto{\pgfqpoint{4.595125in}{8.792418in}}%
\pgfpathquadraticcurveto{\pgfqpoint{4.590115in}{8.794995in}}{\pgfqpoint{4.584693in}{8.796274in}}%
\pgfpathquadraticcurveto{\pgfqpoint{4.579291in}{8.797572in}}{\pgfqpoint{4.573478in}{8.797572in}}%
\pgfpathquadraticcurveto{\pgfqpoint{4.564654in}{8.797572in}}{\pgfqpoint{4.560242in}{8.794872in}}%
\pgfpathquadraticcurveto{\pgfqpoint{4.555830in}{8.792171in}}{\pgfqpoint{4.555830in}{8.786749in}}%
\pgfpathquadraticcurveto{\pgfqpoint{4.555830in}{8.782626in}}{\pgfqpoint{4.558984in}{8.780275in}}%
\pgfpathquadraticcurveto{\pgfqpoint{4.562139in}{8.777925in}}{\pgfqpoint{4.571684in}{8.775802in}}%
\pgfpathlineto{\pgfqpoint{4.575745in}{8.774894in}}%
\pgfpathquadraticcurveto{\pgfqpoint{4.588363in}{8.772194in}}{\pgfqpoint{4.593682in}{8.767266in}}%
\pgfpathquadraticcurveto{\pgfqpoint{4.599001in}{8.762339in}}{\pgfqpoint{4.599001in}{8.753515in}}%
\pgfpathquadraticcurveto{\pgfqpoint{4.599001in}{8.743455in}}{\pgfqpoint{4.591043in}{8.737579in}}%
\pgfpathquadraticcurveto{\pgfqpoint{4.583085in}{8.731724in}}{\pgfqpoint{4.569169in}{8.731724in}}%
\pgfpathquadraticcurveto{\pgfqpoint{4.563376in}{8.731724in}}{\pgfqpoint{4.557088in}{8.732858in}}%
\pgfpathquadraticcurveto{\pgfqpoint{4.550800in}{8.733992in}}{\pgfqpoint{4.543852in}{8.736239in}}%
\pgfpathlineto{\pgfqpoint{4.543852in}{8.748485in}}%
\pgfpathquadraticcurveto{\pgfqpoint{4.550429in}{8.745063in}}{\pgfqpoint{4.556799in}{8.743352in}}%
\pgfpathquadraticcurveto{\pgfqpoint{4.563170in}{8.741661in}}{\pgfqpoint{4.569437in}{8.741661in}}%
\pgfpathquadraticcurveto{\pgfqpoint{4.577807in}{8.741661in}}{\pgfqpoint{4.582301in}{8.744527in}}%
\pgfpathquadraticcurveto{\pgfqpoint{4.586816in}{8.747392in}}{\pgfqpoint{4.586816in}{8.752608in}}%
\pgfpathquadraticcurveto{\pgfqpoint{4.586816in}{8.757432in}}{\pgfqpoint{4.583559in}{8.760010in}}%
\pgfpathquadraticcurveto{\pgfqpoint{4.580322in}{8.762587in}}{\pgfqpoint{4.569293in}{8.764978in}}%
\pgfpathlineto{\pgfqpoint{4.565169in}{8.765947in}}%
\pgfpathquadraticcurveto{\pgfqpoint{4.554160in}{8.768256in}}{\pgfqpoint{4.549253in}{8.773060in}}%
\pgfpathquadraticcurveto{\pgfqpoint{4.544367in}{8.777863in}}{\pgfqpoint{4.544367in}{8.786233in}}%
\pgfpathquadraticcurveto{\pgfqpoint{4.544367in}{8.796418in}}{\pgfqpoint{4.551583in}{8.801943in}}%
\pgfpathquadraticcurveto{\pgfqpoint{4.558799in}{8.807489in}}{\pgfqpoint{4.572076in}{8.807489in}}%
\pgfpathquadraticcurveto{\pgfqpoint{4.578632in}{8.807489in}}{\pgfqpoint{4.584425in}{8.806520in}}%
\pgfpathquadraticcurveto{\pgfqpoint{4.590239in}{8.805572in}}{\pgfqpoint{4.595125in}{8.803634in}}%
\pgfpathlineto{\pgfqpoint{4.595125in}{8.803634in}}%
\pgfpathclose%
\pgfpathmoveto{\pgfqpoint{4.616009in}{8.833713in}}%
\pgfpathlineto{\pgfqpoint{4.626317in}{8.833713in}}%
\pgfpathquadraticcurveto{\pgfqpoint{4.635966in}{8.818519in}}{\pgfqpoint{4.640769in}{8.803943in}}%
\pgfpathquadraticcurveto{\pgfqpoint{4.645573in}{8.789388in}}{\pgfqpoint{4.645573in}{8.775018in}}%
\pgfpathquadraticcurveto{\pgfqpoint{4.645573in}{8.760587in}}{\pgfqpoint{4.640769in}{8.745970in}}%
\pgfpathquadraticcurveto{\pgfqpoint{4.635966in}{8.731353in}}{\pgfqpoint{4.626317in}{8.716200in}}%
\pgfpathlineto{\pgfqpoint{4.616009in}{8.716200in}}%
\pgfpathquadraticcurveto{\pgfqpoint{4.624565in}{8.730961in}}{\pgfqpoint{4.628791in}{8.745557in}}%
\pgfpathquadraticcurveto{\pgfqpoint{4.633018in}{8.760154in}}{\pgfqpoint{4.633018in}{8.775018in}}%
\pgfpathquadraticcurveto{\pgfqpoint{4.633018in}{8.789903in}}{\pgfqpoint{4.628791in}{8.804396in}}%
\pgfpathquadraticcurveto{\pgfqpoint{4.624565in}{8.818910in}}{\pgfqpoint{4.616009in}{8.833713in}}%
\pgfpathlineto{\pgfqpoint{4.616009in}{8.833713in}}%
\pgfpathclose%
\pgfusepath{fill}%
\end{pgfscope}%
\begin{pgfscope}%
\pgfsetbuttcap%
\pgfsetroundjoin%
\definecolor{currentfill}{rgb}{0.090196,0.168627,0.227451}%
\pgfsetfillcolor{currentfill}%
\pgfsetlinewidth{0.000000pt}%
\definecolor{currentstroke}{rgb}{0.090196,0.168627,0.227451}%
\pgfsetstrokecolor{currentstroke}%
\pgfsetdash{}{0pt}%
\pgfpathmoveto{\pgfqpoint{1.108187in}{8.988667in}}%
\pgfpathquadraticcurveto{\pgfqpoint{1.098100in}{8.983441in}}{\pgfqpoint{1.087193in}{8.980798in}}%
\pgfpathquadraticcurveto{\pgfqpoint{1.076286in}{8.978124in}}{\pgfqpoint{1.064406in}{8.978124in}}%
\pgfpathquadraticcurveto{\pgfqpoint{1.028981in}{8.978124in}}{\pgfqpoint{1.008291in}{8.997933in}}%
\pgfpathquadraticcurveto{\pgfqpoint{0.987601in}{9.017742in}}{\pgfqpoint{0.987601in}{9.051618in}}%
\pgfpathquadraticcurveto{\pgfqpoint{0.987601in}{9.085615in}}{\pgfqpoint{1.008291in}{9.105394in}}%
\pgfpathquadraticcurveto{\pgfqpoint{1.028981in}{9.125203in}}{\pgfqpoint{1.064406in}{9.125203in}}%
\pgfpathquadraticcurveto{\pgfqpoint{1.076286in}{9.125203in}}{\pgfqpoint{1.087193in}{9.122530in}}%
\pgfpathquadraticcurveto{\pgfqpoint{1.098100in}{9.119886in}}{\pgfqpoint{1.108187in}{9.114661in}}%
\pgfpathlineto{\pgfqpoint{1.108187in}{9.085342in}}%
\pgfpathquadraticcurveto{\pgfqpoint{1.098009in}{9.092269in}}{\pgfqpoint{1.088135in}{9.095490in}}%
\pgfpathquadraticcurveto{\pgfqpoint{1.078261in}{9.098710in}}{\pgfqpoint{1.067354in}{9.098710in}}%
\pgfpathquadraticcurveto{\pgfqpoint{1.047788in}{9.098710in}}{\pgfqpoint{1.036577in}{9.086162in}}%
\pgfpathquadraticcurveto{\pgfqpoint{1.025396in}{9.073645in}}{\pgfqpoint{1.025396in}{9.051618in}}%
\pgfpathquadraticcurveto{\pgfqpoint{1.025396in}{9.029682in}}{\pgfqpoint{1.036577in}{9.017135in}}%
\pgfpathquadraticcurveto{\pgfqpoint{1.047788in}{9.004617in}}{\pgfqpoint{1.067354in}{9.004617in}}%
\pgfpathquadraticcurveto{\pgfqpoint{1.078261in}{9.004617in}}{\pgfqpoint{1.088135in}{9.007838in}}%
\pgfpathquadraticcurveto{\pgfqpoint{1.098009in}{9.011089in}}{\pgfqpoint{1.108187in}{9.018016in}}%
\pgfpathlineto{\pgfqpoint{1.108187in}{8.988667in}}%
\pgfpathlineto{\pgfqpoint{1.108187in}{8.988667in}}%
\pgfpathclose%
\pgfpathmoveto{\pgfqpoint{1.240713in}{9.118185in}}%
\pgfpathlineto{\pgfqpoint{1.240713in}{9.088168in}}%
\pgfpathquadraticcurveto{\pgfqpoint{1.229046in}{9.093393in}}{\pgfqpoint{1.217926in}{9.096036in}}%
\pgfpathquadraticcurveto{\pgfqpoint{1.206837in}{9.098710in}}{\pgfqpoint{1.196963in}{9.098710in}}%
\pgfpathquadraticcurveto{\pgfqpoint{1.183868in}{9.098710in}}{\pgfqpoint{1.177579in}{9.095095in}}%
\pgfpathquadraticcurveto{\pgfqpoint{1.171321in}{9.091510in}}{\pgfqpoint{1.171321in}{9.083914in}}%
\pgfpathquadraticcurveto{\pgfqpoint{1.171321in}{9.078202in}}{\pgfqpoint{1.175544in}{9.075012in}}%
\pgfpathquadraticcurveto{\pgfqpoint{1.179767in}{9.071852in}}{\pgfqpoint{1.190887in}{9.069574in}}%
\pgfpathlineto{\pgfqpoint{1.206442in}{9.066444in}}%
\pgfpathquadraticcurveto{\pgfqpoint{1.230079in}{9.061674in}}{\pgfqpoint{1.240045in}{9.051983in}}%
\pgfpathquadraticcurveto{\pgfqpoint{1.250040in}{9.042321in}}{\pgfqpoint{1.250040in}{9.024457in}}%
\pgfpathquadraticcurveto{\pgfqpoint{1.250040in}{9.001032in}}{\pgfqpoint{1.236125in}{8.989578in}}%
\pgfpathquadraticcurveto{\pgfqpoint{1.222210in}{8.978124in}}{\pgfqpoint{1.193621in}{8.978124in}}%
\pgfpathquadraticcurveto{\pgfqpoint{1.180162in}{8.978124in}}{\pgfqpoint{1.166581in}{8.980707in}}%
\pgfpathquadraticcurveto{\pgfqpoint{1.153000in}{8.983259in}}{\pgfqpoint{1.139420in}{8.988302in}}%
\pgfpathlineto{\pgfqpoint{1.139420in}{9.019140in}}%
\pgfpathquadraticcurveto{\pgfqpoint{1.153000in}{9.011939in}}{\pgfqpoint{1.165670in}{9.008263in}}%
\pgfpathquadraticcurveto{\pgfqpoint{1.178339in}{9.004617in}}{\pgfqpoint{1.190127in}{9.004617in}}%
\pgfpathquadraticcurveto{\pgfqpoint{1.202097in}{9.004617in}}{\pgfqpoint{1.208447in}{9.008597in}}%
\pgfpathquadraticcurveto{\pgfqpoint{1.214797in}{9.012608in}}{\pgfqpoint{1.214797in}{9.020021in}}%
\pgfpathquadraticcurveto{\pgfqpoint{1.214797in}{9.026644in}}{\pgfqpoint{1.210483in}{9.030260in}}%
\pgfpathquadraticcurveto{\pgfqpoint{1.206169in}{9.033875in}}{\pgfqpoint{1.193256in}{9.036731in}}%
\pgfpathlineto{\pgfqpoint{1.179098in}{9.039860in}}%
\pgfpathquadraticcurveto{\pgfqpoint{1.157831in}{9.044418in}}{\pgfqpoint{1.147987in}{9.054383in}}%
\pgfpathquadraticcurveto{\pgfqpoint{1.138174in}{9.064348in}}{\pgfqpoint{1.138174in}{9.081240in}}%
\pgfpathquadraticcurveto{\pgfqpoint{1.138174in}{9.102417in}}{\pgfqpoint{1.151846in}{9.113810in}}%
\pgfpathquadraticcurveto{\pgfqpoint{1.165518in}{9.125203in}}{\pgfqpoint{1.191160in}{9.125203in}}%
\pgfpathquadraticcurveto{\pgfqpoint{1.202857in}{9.125203in}}{\pgfqpoint{1.215192in}{9.123441in}}%
\pgfpathquadraticcurveto{\pgfqpoint{1.227527in}{9.121679in}}{\pgfqpoint{1.240713in}{9.118185in}}%
\pgfpathlineto{\pgfqpoint{1.240713in}{9.118185in}}%
\pgfpathclose%
\pgfpathmoveto{\pgfqpoint{1.368135in}{9.006714in}}%
\pgfpathlineto{\pgfqpoint{1.310986in}{9.006714in}}%
\pgfpathlineto{\pgfqpoint{1.301963in}{8.980889in}}%
\pgfpathlineto{\pgfqpoint{1.265201in}{8.980889in}}%
\pgfpathlineto{\pgfqpoint{1.317701in}{9.122651in}}%
\pgfpathlineto{\pgfqpoint{1.361299in}{9.122651in}}%
\pgfpathlineto{\pgfqpoint{1.413799in}{8.980889in}}%
\pgfpathlineto{\pgfqpoint{1.377067in}{8.980889in}}%
\pgfpathlineto{\pgfqpoint{1.368135in}{9.006714in}}%
\pgfpathlineto{\pgfqpoint{1.368135in}{9.006714in}}%
\pgfpathclose%
\pgfpathmoveto{\pgfqpoint{1.320101in}{9.033024in}}%
\pgfpathlineto{\pgfqpoint{1.358929in}{9.033024in}}%
\pgfpathlineto{\pgfqpoint{1.339545in}{9.089413in}}%
\pgfpathlineto{\pgfqpoint{1.320101in}{9.033024in}}%
\pgfpathlineto{\pgfqpoint{1.320101in}{9.033024in}}%
\pgfpathclose%
\pgfpathmoveto{\pgfqpoint{1.485379in}{9.087226in}}%
\pgfpathlineto{\pgfqpoint{1.519376in}{9.087226in}}%
\pgfpathlineto{\pgfqpoint{1.545869in}{9.013732in}}%
\pgfpathlineto{\pgfqpoint{1.572241in}{9.087226in}}%
\pgfpathlineto{\pgfqpoint{1.606329in}{9.087226in}}%
\pgfpathlineto{\pgfqpoint{1.564463in}{8.980889in}}%
\pgfpathlineto{\pgfqpoint{1.527154in}{8.980889in}}%
\pgfpathlineto{\pgfqpoint{1.485379in}{9.087226in}}%
\pgfpathlineto{\pgfqpoint{1.485379in}{9.087226in}}%
\pgfpathclose%
\pgfpathmoveto{\pgfqpoint{1.708595in}{9.083914in}}%
\pgfpathlineto{\pgfqpoint{1.708595in}{9.058089in}}%
\pgfpathquadraticcurveto{\pgfqpoint{1.697688in}{9.062647in}}{\pgfqpoint{1.687510in}{9.064925in}}%
\pgfpathquadraticcurveto{\pgfqpoint{1.677362in}{9.067204in}}{\pgfqpoint{1.668339in}{9.067204in}}%
\pgfpathquadraticcurveto{\pgfqpoint{1.658647in}{9.067204in}}{\pgfqpoint{1.653938in}{9.064773in}}%
\pgfpathquadraticcurveto{\pgfqpoint{1.649259in}{9.062343in}}{\pgfqpoint{1.649259in}{9.057330in}}%
\pgfpathquadraticcurveto{\pgfqpoint{1.649259in}{9.053228in}}{\pgfqpoint{1.652814in}{9.051041in}}%
\pgfpathquadraticcurveto{\pgfqpoint{1.656368in}{9.048884in}}{\pgfqpoint{1.665574in}{9.047820in}}%
\pgfpathlineto{\pgfqpoint{1.671559in}{9.046970in}}%
\pgfpathquadraticcurveto{\pgfqpoint{1.697688in}{9.043658in}}{\pgfqpoint{1.706681in}{9.036062in}}%
\pgfpathquadraticcurveto{\pgfqpoint{1.715704in}{9.028467in}}{\pgfqpoint{1.715704in}{9.012213in}}%
\pgfpathquadraticcurveto{\pgfqpoint{1.715704in}{8.995229in}}{\pgfqpoint{1.703156in}{8.986661in}}%
\pgfpathquadraticcurveto{\pgfqpoint{1.690639in}{8.978124in}}{\pgfqpoint{1.665787in}{8.978124in}}%
\pgfpathquadraticcurveto{\pgfqpoint{1.655244in}{8.978124in}}{\pgfqpoint{1.643972in}{8.979795in}}%
\pgfpathquadraticcurveto{\pgfqpoint{1.632731in}{8.981466in}}{\pgfqpoint{1.620852in}{8.984778in}}%
\pgfpathlineto{\pgfqpoint{1.620852in}{9.010602in}}%
\pgfpathquadraticcurveto{\pgfqpoint{1.631030in}{9.005681in}}{\pgfqpoint{1.641694in}{9.003189in}}%
\pgfpathquadraticcurveto{\pgfqpoint{1.652388in}{9.000728in}}{\pgfqpoint{1.663387in}{9.000728in}}%
\pgfpathquadraticcurveto{\pgfqpoint{1.673382in}{9.000728in}}{\pgfqpoint{1.678395in}{9.003463in}}%
\pgfpathquadraticcurveto{\pgfqpoint{1.683439in}{9.006227in}}{\pgfqpoint{1.683439in}{9.011666in}}%
\pgfpathquadraticcurveto{\pgfqpoint{1.683439in}{9.016223in}}{\pgfqpoint{1.679975in}{9.018441in}}%
\pgfpathquadraticcurveto{\pgfqpoint{1.676512in}{9.020659in}}{\pgfqpoint{1.666151in}{9.021905in}}%
\pgfpathlineto{\pgfqpoint{1.660166in}{9.022664in}}%
\pgfpathquadraticcurveto{\pgfqpoint{1.637471in}{9.025520in}}{\pgfqpoint{1.628356in}{9.033207in}}%
\pgfpathquadraticcurveto{\pgfqpoint{1.619242in}{9.040893in}}{\pgfqpoint{1.619242in}{9.056570in}}%
\pgfpathquadraticcurveto{\pgfqpoint{1.619242in}{9.073463in}}{\pgfqpoint{1.630817in}{9.081605in}}%
\pgfpathquadraticcurveto{\pgfqpoint{1.642423in}{9.089778in}}{\pgfqpoint{1.666334in}{9.089778in}}%
\pgfpathquadraticcurveto{\pgfqpoint{1.675752in}{9.089778in}}{\pgfqpoint{1.686082in}{9.088350in}}%
\pgfpathquadraticcurveto{\pgfqpoint{1.696442in}{9.086952in}}{\pgfqpoint{1.708595in}{9.083914in}}%
\pgfpathlineto{\pgfqpoint{1.708595in}{9.083914in}}%
\pgfpathclose%
\pgfpathmoveto{\pgfqpoint{1.793543in}{9.122651in}}%
\pgfpathlineto{\pgfqpoint{1.924185in}{9.122651in}}%
\pgfpathlineto{\pgfqpoint{1.924185in}{9.095003in}}%
\pgfpathlineto{\pgfqpoint{1.877184in}{9.095003in}}%
\pgfpathlineto{\pgfqpoint{1.877184in}{8.980889in}}%
\pgfpathlineto{\pgfqpoint{1.840635in}{8.980889in}}%
\pgfpathlineto{\pgfqpoint{1.840635in}{9.095003in}}%
\pgfpathlineto{\pgfqpoint{1.793543in}{9.095003in}}%
\pgfpathlineto{\pgfqpoint{1.793543in}{9.122651in}}%
\pgfpathlineto{\pgfqpoint{1.793543in}{9.122651in}}%
\pgfpathclose%
\pgfpathmoveto{\pgfqpoint{1.931051in}{9.122651in}}%
\pgfpathlineto{\pgfqpoint{1.966082in}{9.122651in}}%
\pgfpathlineto{\pgfqpoint{1.990570in}{9.019626in}}%
\pgfpathlineto{\pgfqpoint{2.014875in}{9.122651in}}%
\pgfpathlineto{\pgfqpoint{2.050088in}{9.122651in}}%
\pgfpathlineto{\pgfqpoint{2.074393in}{9.019626in}}%
\pgfpathlineto{\pgfqpoint{2.098912in}{9.122651in}}%
\pgfpathlineto{\pgfqpoint{2.133638in}{9.122651in}}%
\pgfpathlineto{\pgfqpoint{2.100218in}{8.980889in}}%
\pgfpathlineto{\pgfqpoint{2.058078in}{8.980889in}}%
\pgfpathlineto{\pgfqpoint{2.032345in}{9.088654in}}%
\pgfpathlineto{\pgfqpoint{2.006915in}{8.980889in}}%
\pgfpathlineto{\pgfqpoint{1.964745in}{8.980889in}}%
\pgfpathlineto{\pgfqpoint{1.931051in}{9.122651in}}%
\pgfpathlineto{\pgfqpoint{1.931051in}{9.122651in}}%
\pgfpathclose%
\pgfpathmoveto{\pgfqpoint{2.157579in}{9.122651in}}%
\pgfpathlineto{\pgfqpoint{2.256199in}{9.122651in}}%
\pgfpathlineto{\pgfqpoint{2.256199in}{9.095003in}}%
\pgfpathlineto{\pgfqpoint{2.194129in}{9.095003in}}%
\pgfpathlineto{\pgfqpoint{2.194129in}{9.068632in}}%
\pgfpathlineto{\pgfqpoint{2.252523in}{9.068632in}}%
\pgfpathlineto{\pgfqpoint{2.252523in}{9.040984in}}%
\pgfpathlineto{\pgfqpoint{2.194129in}{9.040984in}}%
\pgfpathlineto{\pgfqpoint{2.194129in}{8.980889in}}%
\pgfpathlineto{\pgfqpoint{2.157579in}{8.980889in}}%
\pgfpathlineto{\pgfqpoint{2.157579in}{9.122651in}}%
\pgfpathlineto{\pgfqpoint{2.157579in}{9.122651in}}%
\pgfpathclose%
\pgfpathmoveto{\pgfqpoint{2.290409in}{9.122651in}}%
\pgfpathlineto{\pgfqpoint{2.389029in}{9.122651in}}%
\pgfpathlineto{\pgfqpoint{2.389029in}{9.095003in}}%
\pgfpathlineto{\pgfqpoint{2.326959in}{9.095003in}}%
\pgfpathlineto{\pgfqpoint{2.326959in}{9.068632in}}%
\pgfpathlineto{\pgfqpoint{2.385353in}{9.068632in}}%
\pgfpathlineto{\pgfqpoint{2.385353in}{9.040984in}}%
\pgfpathlineto{\pgfqpoint{2.326959in}{9.040984in}}%
\pgfpathlineto{\pgfqpoint{2.326959in}{9.008506in}}%
\pgfpathlineto{\pgfqpoint{2.391125in}{9.008506in}}%
\pgfpathlineto{\pgfqpoint{2.391125in}{8.980889in}}%
\pgfpathlineto{\pgfqpoint{2.290409in}{8.980889in}}%
\pgfpathlineto{\pgfqpoint{2.290409in}{9.122651in}}%
\pgfpathlineto{\pgfqpoint{2.290409in}{9.122651in}}%
\pgfpathclose%
\pgfusepath{fill}%
\end{pgfscope}%
\begin{pgfscope}%
\pgfsetbuttcap%
\pgfsetroundjoin%
\definecolor{currentfill}{rgb}{0.321569,0.321569,0.321569}%
\pgfsetfillcolor{currentfill}%
\pgfsetfillopacity{0.900000}%
\pgfsetlinewidth{0.803000pt}%
\definecolor{currentstroke}{rgb}{0.321569,0.321569,0.321569}%
\pgfsetstrokecolor{currentstroke}%
\pgfsetstrokeopacity{0.900000}%
\pgfsetdash{}{0pt}%
\pgfsys@defobject{currentmarker}{\pgfqpoint{-0.059738in}{-0.059738in}}{\pgfqpoint{0.059738in}{0.059738in}}{%
\pgfpathmoveto{\pgfqpoint{0.000000in}{-0.059738in}}%
\pgfpathlineto{\pgfqpoint{0.059738in}{0.000000in}}%
\pgfpathlineto{\pgfqpoint{0.000000in}{0.059738in}}%
\pgfpathlineto{\pgfqpoint{-0.059738in}{0.000000in}}%
\pgfpathlineto{\pgfqpoint{0.000000in}{-0.059738in}}%
\pgfpathclose%
\pgfusepath{stroke,fill}%
}%
\begin{pgfscope}%
\pgfsys@transformshift{0.977909in}{3.486827in}%
\pgfsys@useobject{currentmarker}{}%
\end{pgfscope}%
\begin{pgfscope}%
\pgfsys@transformshift{0.977909in}{3.703646in}%
\pgfsys@useobject{currentmarker}{}%
\end{pgfscope}%
\end{pgfscope}%
\begin{pgfscope}%
\pgfsetbuttcap%
\pgfsetroundjoin%
\definecolor{currentfill}{rgb}{0.768627,0.603922,0.000000}%
\pgfsetfillcolor{currentfill}%
\pgfsetfillopacity{0.900000}%
\pgfsetlinewidth{0.803000pt}%
\definecolor{currentstroke}{rgb}{0.768627,0.603922,0.000000}%
\pgfsetstrokecolor{currentstroke}%
\pgfsetstrokeopacity{0.900000}%
\pgfsetdash{}{0pt}%
\pgfsys@defobject{currentmarker}{\pgfqpoint{-0.059738in}{-0.059738in}}{\pgfqpoint{0.059738in}{0.059738in}}{%
\pgfpathmoveto{\pgfqpoint{0.000000in}{-0.059738in}}%
\pgfpathlineto{\pgfqpoint{0.059738in}{0.000000in}}%
\pgfpathlineto{\pgfqpoint{0.000000in}{0.059738in}}%
\pgfpathlineto{\pgfqpoint{-0.059738in}{0.000000in}}%
\pgfpathlineto{\pgfqpoint{0.000000in}{-0.059738in}}%
\pgfpathclose%
\pgfusepath{stroke,fill}%
}%
\begin{pgfscope}%
\pgfsys@transformshift{0.977909in}{4.202747in}%
\pgfsys@useobject{currentmarker}{}%
\end{pgfscope}%
\begin{pgfscope}%
\pgfsys@transformshift{0.977909in}{5.341813in}%
\pgfsys@useobject{currentmarker}{}%
\end{pgfscope}%
\end{pgfscope}%
\begin{pgfscope}%
\pgfpathrectangle{\pgfqpoint{0.977909in}{2.928000in}}{\pgfqpoint{5.664000in}{5.664000in}}%
\pgfusepath{clip}%
\pgfsetbuttcap%
\pgfsetroundjoin%
\definecolor{currentfill}{rgb}{0.176471,0.713725,0.643137}%
\pgfsetfillcolor{currentfill}%
\pgfsetfillopacity{0.900000}%
\pgfsetlinewidth{0.803000pt}%
\definecolor{currentstroke}{rgb}{0.176471,0.713725,0.643137}%
\pgfsetstrokecolor{currentstroke}%
\pgfsetstrokeopacity{0.900000}%
\pgfsetdash{}{0pt}%
\pgfsys@defobject{currentmarker}{\pgfqpoint{-0.038665in}{-0.038665in}}{\pgfqpoint{0.038665in}{0.038665in}}{%
\pgfpathmoveto{\pgfqpoint{0.000000in}{-0.038665in}}%
\pgfpathcurveto{\pgfqpoint{0.010254in}{-0.038665in}}{\pgfqpoint{0.020090in}{-0.034591in}}{\pgfqpoint{0.027340in}{-0.027340in}}%
\pgfpathcurveto{\pgfqpoint{0.034591in}{-0.020090in}}{\pgfqpoint{0.038665in}{-0.010254in}}{\pgfqpoint{0.038665in}{0.000000in}}%
\pgfpathcurveto{\pgfqpoint{0.038665in}{0.010254in}}{\pgfqpoint{0.034591in}{0.020090in}}{\pgfqpoint{0.027340in}{0.027340in}}%
\pgfpathcurveto{\pgfqpoint{0.020090in}{0.034591in}}{\pgfqpoint{0.010254in}{0.038665in}}{\pgfqpoint{0.000000in}{0.038665in}}%
\pgfpathcurveto{\pgfqpoint{-0.010254in}{0.038665in}}{\pgfqpoint{-0.020090in}{0.034591in}}{\pgfqpoint{-0.027340in}{0.027340in}}%
\pgfpathcurveto{\pgfqpoint{-0.034591in}{0.020090in}}{\pgfqpoint{-0.038665in}{0.010254in}}{\pgfqpoint{-0.038665in}{0.000000in}}%
\pgfpathcurveto{\pgfqpoint{-0.038665in}{-0.010254in}}{\pgfqpoint{-0.034591in}{-0.020090in}}{\pgfqpoint{-0.027340in}{-0.027340in}}%
\pgfpathcurveto{\pgfqpoint{-0.020090in}{-0.034591in}}{\pgfqpoint{-0.010254in}{-0.038665in}}{\pgfqpoint{0.000000in}{-0.038665in}}%
\pgfpathlineto{\pgfqpoint{0.000000in}{-0.038665in}}%
\pgfpathclose%
\pgfusepath{stroke,fill}%
}%
\begin{pgfscope}%
\pgfsys@transformshift{4.589601in}{8.376059in}%
\pgfsys@useobject{currentmarker}{}%
\end{pgfscope}%
\end{pgfscope}%
\begin{pgfscope}%
\pgfsetbuttcap%
\pgfsetroundjoin%
\definecolor{currentfill}{rgb}{0.000000,0.447059,0.698039}%
\pgfsetfillcolor{currentfill}%
\pgfsetfillopacity{0.900000}%
\pgfsetlinewidth{0.803000pt}%
\definecolor{currentstroke}{rgb}{0.000000,0.447059,0.698039}%
\pgfsetstrokecolor{currentstroke}%
\pgfsetstrokeopacity{0.900000}%
\pgfsetdash{}{0pt}%
\pgfsys@defobject{currentmarker}{\pgfqpoint{-0.059738in}{-0.059738in}}{\pgfqpoint{0.059738in}{0.059738in}}{%
\pgfpathmoveto{\pgfqpoint{0.000000in}{-0.059738in}}%
\pgfpathlineto{\pgfqpoint{0.059738in}{0.000000in}}%
\pgfpathlineto{\pgfqpoint{0.000000in}{0.059738in}}%
\pgfpathlineto{\pgfqpoint{-0.059738in}{0.000000in}}%
\pgfpathlineto{\pgfqpoint{0.000000in}{-0.059738in}}%
\pgfpathclose%
\pgfusepath{stroke,fill}%
}%
\begin{pgfscope}%
\pgfsys@transformshift{3.828962in}{8.592000in}%
\pgfsys@useobject{currentmarker}{}%
\end{pgfscope}%
\begin{pgfscope}%
\pgfsys@transformshift{2.263567in}{2.928000in}%
\pgfsys@useobject{currentmarker}{}%
\end{pgfscope}%
\begin{pgfscope}%
\pgfsys@transformshift{2.285323in}{2.928000in}%
\pgfsys@useobject{currentmarker}{}%
\end{pgfscope}%
\begin{pgfscope}%
\pgfsys@transformshift{2.743448in}{2.928000in}%
\pgfsys@useobject{currentmarker}{}%
\end{pgfscope}%
\begin{pgfscope}%
\pgfsys@transformshift{3.743835in}{2.928000in}%
\pgfsys@useobject{currentmarker}{}%
\end{pgfscope}%
\begin{pgfscope}%
\pgfsys@transformshift{2.806739in}{2.928000in}%
\pgfsys@useobject{currentmarker}{}%
\end{pgfscope}%
\end{pgfscope}%
\begin{pgfscope}%
\pgfsetbuttcap%
\pgfsetroundjoin%
\definecolor{currentfill}{rgb}{0.941176,0.584314,0.337255}%
\pgfsetfillcolor{currentfill}%
\pgfsetfillopacity{0.900000}%
\pgfsetlinewidth{0.803000pt}%
\definecolor{currentstroke}{rgb}{0.941176,0.584314,0.337255}%
\pgfsetstrokecolor{currentstroke}%
\pgfsetstrokeopacity{0.900000}%
\pgfsetdash{}{0pt}%
\pgfsys@defobject{currentmarker}{\pgfqpoint{-0.059738in}{-0.059738in}}{\pgfqpoint{0.059738in}{0.059738in}}{%
\pgfpathmoveto{\pgfqpoint{0.000000in}{-0.059738in}}%
\pgfpathlineto{\pgfqpoint{0.059738in}{0.000000in}}%
\pgfpathlineto{\pgfqpoint{0.000000in}{0.059738in}}%
\pgfpathlineto{\pgfqpoint{-0.059738in}{0.000000in}}%
\pgfpathlineto{\pgfqpoint{0.000000in}{-0.059738in}}%
\pgfpathclose%
\pgfusepath{stroke,fill}%
}%
\begin{pgfscope}%
\pgfsys@transformshift{5.115066in}{8.592000in}%
\pgfsys@useobject{currentmarker}{}%
\end{pgfscope}%
\begin{pgfscope}%
\pgfsys@transformshift{4.212351in}{8.592000in}%
\pgfsys@useobject{currentmarker}{}%
\end{pgfscope}%
\end{pgfscope}%
\begin{pgfscope}%
\pgfsetbuttcap%
\pgfsetroundjoin%
\definecolor{currentfill}{rgb}{0.549020,0.160784,0.505882}%
\pgfsetfillcolor{currentfill}%
\pgfsetfillopacity{0.900000}%
\pgfsetlinewidth{0.803000pt}%
\definecolor{currentstroke}{rgb}{0.549020,0.160784,0.505882}%
\pgfsetstrokecolor{currentstroke}%
\pgfsetstrokeopacity{0.900000}%
\pgfsetdash{}{0pt}%
\pgfsys@defobject{currentmarker}{\pgfqpoint{-0.059738in}{-0.059738in}}{\pgfqpoint{0.059738in}{0.059738in}}{%
\pgfpathmoveto{\pgfqpoint{0.000000in}{-0.059738in}}%
\pgfpathlineto{\pgfqpoint{0.059738in}{0.000000in}}%
\pgfpathlineto{\pgfqpoint{0.000000in}{0.059738in}}%
\pgfpathlineto{\pgfqpoint{-0.059738in}{0.000000in}}%
\pgfpathlineto{\pgfqpoint{0.000000in}{-0.059738in}}%
\pgfpathclose%
\pgfusepath{stroke,fill}%
}%
\begin{pgfscope}%
\pgfsys@transformshift{6.641909in}{8.494328in}%
\pgfsys@useobject{currentmarker}{}%
\end{pgfscope}%
\begin{pgfscope}%
\pgfsys@transformshift{6.641909in}{8.592000in}%
\pgfsys@useobject{currentmarker}{}%
\end{pgfscope}%
\end{pgfscope}%
\begin{pgfscope}%
\pgfpathrectangle{\pgfqpoint{0.977909in}{2.928000in}}{\pgfqpoint{5.664000in}{5.664000in}}%
\pgfusepath{clip}%
\pgfsetbuttcap%
\pgfsetroundjoin%
\definecolor{currentfill}{rgb}{0.000000,0.560784,0.611765}%
\pgfsetfillcolor{currentfill}%
\pgfsetfillopacity{0.900000}%
\pgfsetlinewidth{0.803000pt}%
\definecolor{currentstroke}{rgb}{0.000000,0.560784,0.611765}%
\pgfsetstrokecolor{currentstroke}%
\pgfsetstrokeopacity{0.900000}%
\pgfsetdash{}{0pt}%
\pgfsys@defobject{currentmarker}{\pgfqpoint{-0.038665in}{-0.038665in}}{\pgfqpoint{0.038665in}{0.038665in}}{%
\pgfpathmoveto{\pgfqpoint{0.000000in}{-0.038665in}}%
\pgfpathcurveto{\pgfqpoint{0.010254in}{-0.038665in}}{\pgfqpoint{0.020090in}{-0.034591in}}{\pgfqpoint{0.027340in}{-0.027340in}}%
\pgfpathcurveto{\pgfqpoint{0.034591in}{-0.020090in}}{\pgfqpoint{0.038665in}{-0.010254in}}{\pgfqpoint{0.038665in}{0.000000in}}%
\pgfpathcurveto{\pgfqpoint{0.038665in}{0.010254in}}{\pgfqpoint{0.034591in}{0.020090in}}{\pgfqpoint{0.027340in}{0.027340in}}%
\pgfpathcurveto{\pgfqpoint{0.020090in}{0.034591in}}{\pgfqpoint{0.010254in}{0.038665in}}{\pgfqpoint{0.000000in}{0.038665in}}%
\pgfpathcurveto{\pgfqpoint{-0.010254in}{0.038665in}}{\pgfqpoint{-0.020090in}{0.034591in}}{\pgfqpoint{-0.027340in}{0.027340in}}%
\pgfpathcurveto{\pgfqpoint{-0.034591in}{0.020090in}}{\pgfqpoint{-0.038665in}{0.010254in}}{\pgfqpoint{-0.038665in}{0.000000in}}%
\pgfpathcurveto{\pgfqpoint{-0.038665in}{-0.010254in}}{\pgfqpoint{-0.034591in}{-0.020090in}}{\pgfqpoint{-0.027340in}{-0.027340in}}%
\pgfpathcurveto{\pgfqpoint{-0.020090in}{-0.034591in}}{\pgfqpoint{-0.010254in}{-0.038665in}}{\pgfqpoint{0.000000in}{-0.038665in}}%
\pgfpathlineto{\pgfqpoint{0.000000in}{-0.038665in}}%
\pgfpathclose%
\pgfusepath{stroke,fill}%
}%
\begin{pgfscope}%
\pgfsys@transformshift{1.790598in}{5.216460in}%
\pgfsys@useobject{currentmarker}{}%
\end{pgfscope}%
\end{pgfscope}%
\begin{pgfscope}%
\pgfsetbuttcap%
\pgfsetroundjoin%
\definecolor{currentfill}{rgb}{0.000000,0.560784,0.611765}%
\pgfsetfillcolor{currentfill}%
\pgfsetfillopacity{0.900000}%
\pgfsetlinewidth{0.803000pt}%
\definecolor{currentstroke}{rgb}{0.000000,0.560784,0.611765}%
\pgfsetstrokecolor{currentstroke}%
\pgfsetstrokeopacity{0.900000}%
\pgfsetdash{}{0pt}%
\pgfsys@defobject{currentmarker}{\pgfqpoint{-0.059738in}{-0.059738in}}{\pgfqpoint{0.059738in}{0.059738in}}{%
\pgfpathmoveto{\pgfqpoint{0.000000in}{-0.059738in}}%
\pgfpathlineto{\pgfqpoint{0.059738in}{0.000000in}}%
\pgfpathlineto{\pgfqpoint{0.000000in}{0.059738in}}%
\pgfpathlineto{\pgfqpoint{-0.059738in}{0.000000in}}%
\pgfpathlineto{\pgfqpoint{0.000000in}{-0.059738in}}%
\pgfpathclose%
\pgfusepath{stroke,fill}%
}%
\begin{pgfscope}%
\pgfsys@transformshift{0.977909in}{3.282932in}%
\pgfsys@useobject{currentmarker}{}%
\end{pgfscope}%
\end{pgfscope}%
\begin{pgfscope}%
\pgfsetbuttcap%
\pgfsetroundjoin%
\pgfsetlinewidth{0.903375pt}%
\definecolor{currentstroke}{rgb}{0.321569,0.321569,0.321569}%
\pgfsetstrokecolor{currentstroke}%
\pgfsetdash{}{0pt}%
\pgfpathmoveto{\pgfqpoint{0.977909in}{2.911471in}}%
\pgfpathlineto{\pgfqpoint{1.050743in}{2.984305in}}%
\pgfpathlineto{\pgfqpoint{0.977909in}{3.057139in}}%
\pgfpathlineto{\pgfqpoint{0.905075in}{2.984305in}}%
\pgfpathlineto{\pgfqpoint{0.977909in}{2.911471in}}%
\pgfpathclose%
\pgfusepath{stroke}%
\end{pgfscope}%
\begin{pgfscope}%
\pgfsetbuttcap%
\pgfsetroundjoin%
\pgfsetlinewidth{0.903375pt}%
\definecolor{currentstroke}{rgb}{0.321569,0.321569,0.321569}%
\pgfsetstrokecolor{currentstroke}%
\pgfsetdash{}{0pt}%
\pgfpathmoveto{\pgfqpoint{0.977909in}{3.649807in}}%
\pgfpathlineto{\pgfqpoint{1.050743in}{3.722641in}}%
\pgfpathlineto{\pgfqpoint{0.977909in}{3.795474in}}%
\pgfpathlineto{\pgfqpoint{0.905075in}{3.722641in}}%
\pgfpathlineto{\pgfqpoint{0.977909in}{3.649807in}}%
\pgfpathclose%
\pgfusepath{stroke}%
\end{pgfscope}%
\begin{pgfscope}%
\pgfsetbuttcap%
\pgfsetroundjoin%
\pgfsetlinewidth{0.903375pt}%
\definecolor{currentstroke}{rgb}{0.321569,0.321569,0.321569}%
\pgfsetstrokecolor{currentstroke}%
\pgfsetdash{}{0pt}%
\pgfpathmoveto{\pgfqpoint{0.977909in}{3.013932in}}%
\pgfpathlineto{\pgfqpoint{1.050743in}{3.086766in}}%
\pgfpathlineto{\pgfqpoint{0.977909in}{3.159600in}}%
\pgfpathlineto{\pgfqpoint{0.905075in}{3.086766in}}%
\pgfpathlineto{\pgfqpoint{0.977909in}{3.013932in}}%
\pgfpathclose%
\pgfusepath{stroke}%
\end{pgfscope}%
\begin{pgfscope}%
\pgfsetbuttcap%
\pgfsetroundjoin%
\pgfsetlinewidth{0.903375pt}%
\definecolor{currentstroke}{rgb}{0.321569,0.321569,0.321569}%
\pgfsetstrokecolor{currentstroke}%
\pgfsetdash{}{0pt}%
\pgfpathmoveto{\pgfqpoint{6.641909in}{6.439903in}}%
\pgfpathlineto{\pgfqpoint{6.714743in}{6.512736in}}%
\pgfpathlineto{\pgfqpoint{6.641909in}{6.585570in}}%
\pgfpathlineto{\pgfqpoint{6.569075in}{6.512736in}}%
\pgfpathlineto{\pgfqpoint{6.641909in}{6.439903in}}%
\pgfpathclose%
\pgfusepath{stroke}%
\end{pgfscope}%
\begin{pgfscope}%
\pgfpathrectangle{\pgfqpoint{0.977909in}{2.928000in}}{\pgfqpoint{5.664000in}{5.664000in}}%
\pgfusepath{clip}%
\pgfsetbuttcap%
\pgfsetroundjoin%
\pgfsetlinewidth{0.803000pt}%
\definecolor{currentstroke}{rgb}{0.768627,0.603922,0.000000}%
\pgfsetstrokecolor{currentstroke}%
\pgfsetdash{}{0pt}%
\pgfpathmoveto{\pgfqpoint{5.287964in}{5.968075in}}%
\pgfpathlineto{\pgfqpoint{5.381133in}{5.968075in}}%
\pgfpathlineto{\pgfqpoint{5.381133in}{6.061244in}}%
\pgfpathlineto{\pgfqpoint{5.287964in}{6.061244in}}%
\pgfpathlineto{\pgfqpoint{5.287964in}{5.968075in}}%
\pgfpathclose%
\pgfusepath{stroke}%
\end{pgfscope}%
\begin{pgfscope}%
\pgfpathrectangle{\pgfqpoint{0.977909in}{2.928000in}}{\pgfqpoint{5.664000in}{5.664000in}}%
\pgfusepath{clip}%
\pgfsetbuttcap%
\pgfsetroundjoin%
\pgfsetlinewidth{0.803000pt}%
\definecolor{currentstroke}{rgb}{0.768627,0.603922,0.000000}%
\pgfsetstrokecolor{currentstroke}%
\pgfsetdash{}{0pt}%
\pgfpathmoveto{\pgfqpoint{5.386376in}{5.883058in}}%
\pgfpathlineto{\pgfqpoint{5.479545in}{5.883058in}}%
\pgfpathlineto{\pgfqpoint{5.479545in}{5.976228in}}%
\pgfpathlineto{\pgfqpoint{5.386376in}{5.976228in}}%
\pgfpathlineto{\pgfqpoint{5.386376in}{5.883058in}}%
\pgfpathclose%
\pgfusepath{stroke}%
\end{pgfscope}%
\begin{pgfscope}%
\pgfpathrectangle{\pgfqpoint{0.977909in}{2.928000in}}{\pgfqpoint{5.664000in}{5.664000in}}%
\pgfusepath{clip}%
\pgfsetbuttcap%
\pgfsetroundjoin%
\pgfsetlinewidth{0.803000pt}%
\definecolor{currentstroke}{rgb}{0.768627,0.603922,0.000000}%
\pgfsetstrokecolor{currentstroke}%
\pgfsetdash{}{0pt}%
\pgfpathmoveto{\pgfqpoint{3.577261in}{6.216615in}}%
\pgfpathlineto{\pgfqpoint{3.670430in}{6.216615in}}%
\pgfpathlineto{\pgfqpoint{3.670430in}{6.309784in}}%
\pgfpathlineto{\pgfqpoint{3.577261in}{6.309784in}}%
\pgfpathlineto{\pgfqpoint{3.577261in}{6.216615in}}%
\pgfpathclose%
\pgfusepath{stroke}%
\end{pgfscope}%
\begin{pgfscope}%
\pgfpathrectangle{\pgfqpoint{0.977909in}{2.928000in}}{\pgfqpoint{5.664000in}{5.664000in}}%
\pgfusepath{clip}%
\pgfsetbuttcap%
\pgfsetroundjoin%
\pgfsetlinewidth{0.803000pt}%
\definecolor{currentstroke}{rgb}{0.768627,0.603922,0.000000}%
\pgfsetstrokecolor{currentstroke}%
\pgfsetdash{}{0pt}%
\pgfpathmoveto{\pgfqpoint{4.092515in}{5.607736in}}%
\pgfpathlineto{\pgfqpoint{4.185684in}{5.607736in}}%
\pgfpathlineto{\pgfqpoint{4.185684in}{5.700905in}}%
\pgfpathlineto{\pgfqpoint{4.092515in}{5.700905in}}%
\pgfpathlineto{\pgfqpoint{4.092515in}{5.607736in}}%
\pgfpathclose%
\pgfusepath{stroke}%
\end{pgfscope}%
\begin{pgfscope}%
\pgfpathrectangle{\pgfqpoint{0.977909in}{2.928000in}}{\pgfqpoint{5.664000in}{5.664000in}}%
\pgfusepath{clip}%
\pgfsetbuttcap%
\pgfsetroundjoin%
\pgfsetlinewidth{0.803000pt}%
\definecolor{currentstroke}{rgb}{0.768627,0.603922,0.000000}%
\pgfsetstrokecolor{currentstroke}%
\pgfsetdash{}{0pt}%
\pgfpathmoveto{\pgfqpoint{4.000478in}{5.510280in}}%
\pgfpathlineto{\pgfqpoint{4.093647in}{5.510280in}}%
\pgfpathlineto{\pgfqpoint{4.093647in}{5.603449in}}%
\pgfpathlineto{\pgfqpoint{4.000478in}{5.603449in}}%
\pgfpathlineto{\pgfqpoint{4.000478in}{5.510280in}}%
\pgfpathclose%
\pgfusepath{stroke}%
\end{pgfscope}%
\begin{pgfscope}%
\pgfpathrectangle{\pgfqpoint{0.977909in}{2.928000in}}{\pgfqpoint{5.664000in}{5.664000in}}%
\pgfusepath{clip}%
\pgfsetbuttcap%
\pgfsetroundjoin%
\pgfsetlinewidth{0.803000pt}%
\definecolor{currentstroke}{rgb}{0.768627,0.603922,0.000000}%
\pgfsetstrokecolor{currentstroke}%
\pgfsetdash{}{0pt}%
\pgfpathmoveto{\pgfqpoint{3.613003in}{5.402820in}}%
\pgfpathlineto{\pgfqpoint{3.706173in}{5.402820in}}%
\pgfpathlineto{\pgfqpoint{3.706173in}{5.495990in}}%
\pgfpathlineto{\pgfqpoint{3.613003in}{5.495990in}}%
\pgfpathlineto{\pgfqpoint{3.613003in}{5.402820in}}%
\pgfpathclose%
\pgfusepath{stroke}%
\end{pgfscope}%
\begin{pgfscope}%
\pgfsetbuttcap%
\pgfsetroundjoin%
\pgfsetlinewidth{0.903375pt}%
\definecolor{currentstroke}{rgb}{0.337255,0.705882,0.913725}%
\pgfsetstrokecolor{currentstroke}%
\pgfsetdash{}{0pt}%
\pgfpathmoveto{\pgfqpoint{6.641909in}{8.519166in}}%
\pgfpathlineto{\pgfqpoint{6.714743in}{8.592000in}}%
\pgfpathlineto{\pgfqpoint{6.641909in}{8.664834in}}%
\pgfpathlineto{\pgfqpoint{6.569075in}{8.592000in}}%
\pgfpathlineto{\pgfqpoint{6.641909in}{8.519166in}}%
\pgfpathclose%
\pgfusepath{stroke}%
\end{pgfscope}%
\begin{pgfscope}%
\pgfpathrectangle{\pgfqpoint{0.977909in}{2.928000in}}{\pgfqpoint{5.664000in}{5.664000in}}%
\pgfusepath{clip}%
\pgfsetbuttcap%
\pgfsetroundjoin%
\pgfsetlinewidth{0.803000pt}%
\definecolor{currentstroke}{rgb}{0.901961,0.196078,0.156863}%
\pgfsetstrokecolor{currentstroke}%
\pgfsetdash{}{0pt}%
\pgfpathmoveto{\pgfqpoint{4.756458in}{7.763238in}}%
\pgfpathlineto{\pgfqpoint{4.849627in}{7.763238in}}%
\pgfpathlineto{\pgfqpoint{4.849627in}{7.856407in}}%
\pgfpathlineto{\pgfqpoint{4.756458in}{7.856407in}}%
\pgfpathlineto{\pgfqpoint{4.756458in}{7.763238in}}%
\pgfpathclose%
\pgfusepath{stroke}%
\end{pgfscope}%
\begin{pgfscope}%
\pgfpathrectangle{\pgfqpoint{0.977909in}{2.928000in}}{\pgfqpoint{5.664000in}{5.664000in}}%
\pgfusepath{clip}%
\pgfsetbuttcap%
\pgfsetroundjoin%
\pgfsetlinewidth{0.803000pt}%
\definecolor{currentstroke}{rgb}{0.901961,0.196078,0.156863}%
\pgfsetstrokecolor{currentstroke}%
\pgfsetdash{}{0pt}%
\pgfpathmoveto{\pgfqpoint{3.805026in}{5.993820in}}%
\pgfpathlineto{\pgfqpoint{3.898195in}{5.993820in}}%
\pgfpathlineto{\pgfqpoint{3.898195in}{6.086990in}}%
\pgfpathlineto{\pgfqpoint{3.805026in}{6.086990in}}%
\pgfpathlineto{\pgfqpoint{3.805026in}{5.993820in}}%
\pgfpathclose%
\pgfusepath{stroke}%
\end{pgfscope}%
\begin{pgfscope}%
\pgfsetbuttcap%
\pgfsetroundjoin%
\pgfsetlinewidth{0.803000pt}%
\definecolor{currentstroke}{rgb}{0.901961,0.196078,0.156863}%
\pgfsetstrokecolor{currentstroke}%
\pgfsetdash{}{0pt}%
\pgfpathmoveto{\pgfqpoint{5.875013in}{8.498314in}}%
\pgfpathlineto{\pgfqpoint{5.968699in}{8.592000in}}%
\pgfpathlineto{\pgfqpoint{5.875013in}{8.685686in}}%
\pgfpathlineto{\pgfqpoint{5.781327in}{8.592000in}}%
\pgfpathlineto{\pgfqpoint{5.875013in}{8.498314in}}%
\pgfpathclose%
\pgfusepath{stroke}%
\end{pgfscope}%
\begin{pgfscope}%
\pgfsetbuttcap%
\pgfsetroundjoin%
\pgfsetlinewidth{0.803000pt}%
\definecolor{currentstroke}{rgb}{0.901961,0.196078,0.156863}%
\pgfsetstrokecolor{currentstroke}%
\pgfsetdash{}{0pt}%
\pgfpathmoveto{\pgfqpoint{6.641909in}{8.195180in}}%
\pgfpathlineto{\pgfqpoint{6.735595in}{8.288866in}}%
\pgfpathlineto{\pgfqpoint{6.641909in}{8.382552in}}%
\pgfpathlineto{\pgfqpoint{6.548223in}{8.288866in}}%
\pgfpathlineto{\pgfqpoint{6.641909in}{8.195180in}}%
\pgfpathclose%
\pgfusepath{stroke}%
\end{pgfscope}%
\begin{pgfscope}%
\pgfsetbuttcap%
\pgfsetroundjoin%
\pgfsetlinewidth{0.803000pt}%
\definecolor{currentstroke}{rgb}{0.901961,0.196078,0.156863}%
\pgfsetstrokecolor{currentstroke}%
\pgfsetdash{}{0pt}%
\pgfpathmoveto{\pgfqpoint{6.595324in}{8.242281in}}%
\pgfpathlineto{\pgfqpoint{6.688494in}{8.242281in}}%
\pgfpathlineto{\pgfqpoint{6.688494in}{8.335451in}}%
\pgfpathlineto{\pgfqpoint{6.595324in}{8.335451in}}%
\pgfpathlineto{\pgfqpoint{6.595324in}{8.242281in}}%
\pgfpathclose%
\pgfusepath{stroke}%
\end{pgfscope}%
\begin{pgfscope}%
\pgfsetbuttcap%
\pgfsetroundjoin%
\pgfsetlinewidth{0.903375pt}%
\definecolor{currentstroke}{rgb}{0.901961,0.196078,0.156863}%
\pgfsetstrokecolor{currentstroke}%
\pgfsetdash{}{0pt}%
\pgfpathmoveto{\pgfqpoint{6.248891in}{8.519166in}}%
\pgfpathlineto{\pgfqpoint{6.321725in}{8.592000in}}%
\pgfpathlineto{\pgfqpoint{6.248891in}{8.664834in}}%
\pgfpathlineto{\pgfqpoint{6.176057in}{8.592000in}}%
\pgfpathlineto{\pgfqpoint{6.248891in}{8.519166in}}%
\pgfpathclose%
\pgfusepath{stroke}%
\end{pgfscope}%
\begin{pgfscope}%
\pgfsetbuttcap%
\pgfsetroundjoin%
\pgfsetlinewidth{0.803000pt}%
\definecolor{currentstroke}{rgb}{0.498039,0.560784,0.678431}%
\pgfsetstrokecolor{currentstroke}%
\pgfsetdash{}{0pt}%
\pgfpathmoveto{\pgfqpoint{0.977909in}{2.834314in}}%
\pgfpathlineto{\pgfqpoint{1.071595in}{2.928000in}}%
\pgfpathlineto{\pgfqpoint{0.977909in}{3.021686in}}%
\pgfpathlineto{\pgfqpoint{0.884223in}{2.928000in}}%
\pgfpathlineto{\pgfqpoint{0.977909in}{2.834314in}}%
\pgfpathclose%
\pgfusepath{stroke}%
\end{pgfscope}%
\begin{pgfscope}%
\pgfpathrectangle{\pgfqpoint{0.977909in}{2.928000in}}{\pgfqpoint{5.664000in}{5.664000in}}%
\pgfusepath{clip}%
\pgfsetbuttcap%
\pgfsetroundjoin%
\definecolor{currentfill}{rgb}{0.541176,0.380392,0.658824}%
\pgfsetfillcolor{currentfill}%
\pgfsetlinewidth{1.254687pt}%
\definecolor{currentstroke}{rgb}{0.541176,0.380392,0.658824}%
\pgfsetstrokecolor{currentstroke}%
\pgfsetdash{}{0pt}%
\pgfsys@defobject{currentmarker}{\pgfqpoint{-0.043368in}{-0.043368in}}{\pgfqpoint{0.043368in}{0.043368in}}{%
\pgfpathmoveto{\pgfqpoint{-0.043368in}{-0.043368in}}%
\pgfpathlineto{\pgfqpoint{0.043368in}{0.043368in}}%
\pgfpathmoveto{\pgfqpoint{-0.043368in}{0.043368in}}%
\pgfpathlineto{\pgfqpoint{0.043368in}{-0.043368in}}%
\pgfusepath{stroke,fill}%
}%
\begin{pgfscope}%
\pgfsys@transformshift{4.438991in}{6.942726in}%
\pgfsys@useobject{currentmarker}{}%
\end{pgfscope}%
\begin{pgfscope}%
\pgfsys@transformshift{3.511169in}{5.142631in}%
\pgfsys@useobject{currentmarker}{}%
\end{pgfscope}%
\begin{pgfscope}%
\pgfsys@transformshift{4.031743in}{5.249998in}%
\pgfsys@useobject{currentmarker}{}%
\end{pgfscope}%
\end{pgfscope}%
\begin{pgfscope}%
\pgfpathrectangle{\pgfqpoint{0.977909in}{2.928000in}}{\pgfqpoint{5.664000in}{5.664000in}}%
\pgfusepath{clip}%
\pgfsetbuttcap%
\pgfsetroundjoin%
\definecolor{currentfill}{rgb}{0.321569,0.321569,0.321569}%
\pgfsetfillcolor{currentfill}%
\pgfsetlinewidth{1.254687pt}%
\definecolor{currentstroke}{rgb}{0.321569,0.321569,0.321569}%
\pgfsetstrokecolor{currentstroke}%
\pgfsetdash{}{0pt}%
\pgfsys@defobject{currentmarker}{\pgfqpoint{-0.043368in}{-0.043368in}}{\pgfqpoint{0.043368in}{0.043368in}}{%
\pgfpathmoveto{\pgfqpoint{-0.043368in}{-0.043368in}}%
\pgfpathlineto{\pgfqpoint{0.043368in}{0.043368in}}%
\pgfpathmoveto{\pgfqpoint{-0.043368in}{0.043368in}}%
\pgfpathlineto{\pgfqpoint{0.043368in}{-0.043368in}}%
\pgfusepath{stroke,fill}%
}%
\begin{pgfscope}%
\pgfsys@transformshift{2.491897in}{4.546643in}%
\pgfsys@useobject{currentmarker}{}%
\end{pgfscope}%
\begin{pgfscope}%
\pgfsys@transformshift{2.365396in}{4.511315in}%
\pgfsys@useobject{currentmarker}{}%
\end{pgfscope}%
\begin{pgfscope}%
\pgfsys@transformshift{3.268953in}{6.082665in}%
\pgfsys@useobject{currentmarker}{}%
\end{pgfscope}%
\begin{pgfscope}%
\pgfsys@transformshift{3.144551in}{5.532329in}%
\pgfsys@useobject{currentmarker}{}%
\end{pgfscope}%
\begin{pgfscope}%
\pgfsys@transformshift{5.411713in}{7.105527in}%
\pgfsys@useobject{currentmarker}{}%
\end{pgfscope}%
\begin{pgfscope}%
\pgfsys@transformshift{4.989252in}{6.705418in}%
\pgfsys@useobject{currentmarker}{}%
\end{pgfscope}%
\begin{pgfscope}%
\pgfsys@transformshift{5.042329in}{6.763732in}%
\pgfsys@useobject{currentmarker}{}%
\end{pgfscope}%
\begin{pgfscope}%
\pgfsys@transformshift{5.319875in}{7.078714in}%
\pgfsys@useobject{currentmarker}{}%
\end{pgfscope}%
\begin{pgfscope}%
\pgfsys@transformshift{4.997231in}{6.669286in}%
\pgfsys@useobject{currentmarker}{}%
\end{pgfscope}%
\begin{pgfscope}%
\pgfsys@transformshift{5.507568in}{6.842235in}%
\pgfsys@useobject{currentmarker}{}%
\end{pgfscope}%
\begin{pgfscope}%
\pgfsys@transformshift{5.206402in}{6.821625in}%
\pgfsys@useobject{currentmarker}{}%
\end{pgfscope}%
\begin{pgfscope}%
\pgfsys@transformshift{5.450477in}{6.857856in}%
\pgfsys@useobject{currentmarker}{}%
\end{pgfscope}%
\begin{pgfscope}%
\pgfsys@transformshift{6.442271in}{8.502535in}%
\pgfsys@useobject{currentmarker}{}%
\end{pgfscope}%
\begin{pgfscope}%
\pgfsys@transformshift{5.845639in}{7.830814in}%
\pgfsys@useobject{currentmarker}{}%
\end{pgfscope}%
\begin{pgfscope}%
\pgfsys@transformshift{2.521073in}{4.962680in}%
\pgfsys@useobject{currentmarker}{}%
\end{pgfscope}%
\begin{pgfscope}%
\pgfsys@transformshift{3.292495in}{5.199830in}%
\pgfsys@useobject{currentmarker}{}%
\end{pgfscope}%
\begin{pgfscope}%
\pgfsys@transformshift{2.456002in}{4.820261in}%
\pgfsys@useobject{currentmarker}{}%
\end{pgfscope}%
\begin{pgfscope}%
\pgfsys@transformshift{2.886227in}{5.120928in}%
\pgfsys@useobject{currentmarker}{}%
\end{pgfscope}%
\begin{pgfscope}%
\pgfsys@transformshift{3.899586in}{6.036239in}%
\pgfsys@useobject{currentmarker}{}%
\end{pgfscope}%
\begin{pgfscope}%
\pgfsys@transformshift{3.643746in}{5.904673in}%
\pgfsys@useobject{currentmarker}{}%
\end{pgfscope}%
\begin{pgfscope}%
\pgfsys@transformshift{2.372660in}{3.705328in}%
\pgfsys@useobject{currentmarker}{}%
\end{pgfscope}%
\begin{pgfscope}%
\pgfsys@transformshift{3.392716in}{5.288067in}%
\pgfsys@useobject{currentmarker}{}%
\end{pgfscope}%
\begin{pgfscope}%
\pgfsys@transformshift{2.993675in}{5.274084in}%
\pgfsys@useobject{currentmarker}{}%
\end{pgfscope}%
\begin{pgfscope}%
\pgfsys@transformshift{3.336193in}{5.693737in}%
\pgfsys@useobject{currentmarker}{}%
\end{pgfscope}%
\begin{pgfscope}%
\pgfsys@transformshift{2.791022in}{5.778231in}%
\pgfsys@useobject{currentmarker}{}%
\end{pgfscope}%
\begin{pgfscope}%
\pgfsys@transformshift{4.723748in}{7.315767in}%
\pgfsys@useobject{currentmarker}{}%
\end{pgfscope}%
\begin{pgfscope}%
\pgfsys@transformshift{4.888340in}{7.217833in}%
\pgfsys@useobject{currentmarker}{}%
\end{pgfscope}%
\begin{pgfscope}%
\pgfsys@transformshift{2.220301in}{4.005870in}%
\pgfsys@useobject{currentmarker}{}%
\end{pgfscope}%
\begin{pgfscope}%
\pgfsys@transformshift{4.688696in}{6.967367in}%
\pgfsys@useobject{currentmarker}{}%
\end{pgfscope}%
\begin{pgfscope}%
\pgfsys@transformshift{4.700792in}{6.976322in}%
\pgfsys@useobject{currentmarker}{}%
\end{pgfscope}%
\begin{pgfscope}%
\pgfsys@transformshift{4.900389in}{7.095557in}%
\pgfsys@useobject{currentmarker}{}%
\end{pgfscope}%
\begin{pgfscope}%
\pgfsys@transformshift{5.128247in}{7.189578in}%
\pgfsys@useobject{currentmarker}{}%
\end{pgfscope}%
\begin{pgfscope}%
\pgfsys@transformshift{2.194820in}{4.490861in}%
\pgfsys@useobject{currentmarker}{}%
\end{pgfscope}%
\begin{pgfscope}%
\pgfsys@transformshift{2.067306in}{3.953384in}%
\pgfsys@useobject{currentmarker}{}%
\end{pgfscope}%
\begin{pgfscope}%
\pgfsys@transformshift{4.481209in}{6.152264in}%
\pgfsys@useobject{currentmarker}{}%
\end{pgfscope}%
\begin{pgfscope}%
\pgfsys@transformshift{4.725971in}{5.874931in}%
\pgfsys@useobject{currentmarker}{}%
\end{pgfscope}%
\begin{pgfscope}%
\pgfsys@transformshift{6.623462in}{7.210034in}%
\pgfsys@useobject{currentmarker}{}%
\end{pgfscope}%
\end{pgfscope}%
\begin{pgfscope}%
\pgfsetbuttcap%
\pgfsetroundjoin%
\definecolor{currentfill}{rgb}{0.321569,0.321569,0.321569}%
\pgfsetfillcolor{currentfill}%
\pgfsetlinewidth{1.254687pt}%
\definecolor{currentstroke}{rgb}{0.321569,0.321569,0.321569}%
\pgfsetstrokecolor{currentstroke}%
\pgfsetdash{}{0pt}%
\pgfsys@defobject{currentmarker}{\pgfqpoint{-0.043368in}{-0.043368in}}{\pgfqpoint{0.043368in}{0.043368in}}{%
\pgfpathmoveto{\pgfqpoint{-0.043368in}{-0.043368in}}%
\pgfpathlineto{\pgfqpoint{0.043368in}{0.043368in}}%
\pgfpathmoveto{\pgfqpoint{-0.043368in}{0.043368in}}%
\pgfpathlineto{\pgfqpoint{0.043368in}{-0.043368in}}%
\pgfusepath{stroke,fill}%
}%
\begin{pgfscope}%
\pgfsys@transformshift{0.977909in}{2.984305in}%
\pgfsys@useobject{currentmarker}{}%
\end{pgfscope}%
\begin{pgfscope}%
\pgfsys@transformshift{0.977909in}{3.722641in}%
\pgfsys@useobject{currentmarker}{}%
\end{pgfscope}%
\begin{pgfscope}%
\pgfsys@transformshift{0.977909in}{3.086766in}%
\pgfsys@useobject{currentmarker}{}%
\end{pgfscope}%
\begin{pgfscope}%
\pgfsys@transformshift{6.641909in}{6.512736in}%
\pgfsys@useobject{currentmarker}{}%
\end{pgfscope}%
\end{pgfscope}%
\begin{pgfscope}%
\pgfpathrectangle{\pgfqpoint{0.977909in}{2.928000in}}{\pgfqpoint{5.664000in}{5.664000in}}%
\pgfusepath{clip}%
\pgfsetbuttcap%
\pgfsetroundjoin%
\definecolor{currentfill}{rgb}{0.768627,0.603922,0.000000}%
\pgfsetfillcolor{currentfill}%
\pgfsetlinewidth{1.254687pt}%
\definecolor{currentstroke}{rgb}{0.768627,0.603922,0.000000}%
\pgfsetstrokecolor{currentstroke}%
\pgfsetdash{}{0pt}%
\pgfsys@defobject{currentmarker}{\pgfqpoint{-0.043368in}{-0.043368in}}{\pgfqpoint{0.043368in}{0.043368in}}{%
\pgfpathmoveto{\pgfqpoint{-0.043368in}{-0.043368in}}%
\pgfpathlineto{\pgfqpoint{0.043368in}{0.043368in}}%
\pgfpathmoveto{\pgfqpoint{-0.043368in}{0.043368in}}%
\pgfpathlineto{\pgfqpoint{0.043368in}{-0.043368in}}%
\pgfusepath{stroke,fill}%
}%
\begin{pgfscope}%
\pgfsys@transformshift{5.334548in}{6.014660in}%
\pgfsys@useobject{currentmarker}{}%
\end{pgfscope}%
\begin{pgfscope}%
\pgfsys@transformshift{5.432961in}{5.929643in}%
\pgfsys@useobject{currentmarker}{}%
\end{pgfscope}%
\begin{pgfscope}%
\pgfsys@transformshift{3.623846in}{6.263199in}%
\pgfsys@useobject{currentmarker}{}%
\end{pgfscope}%
\begin{pgfscope}%
\pgfsys@transformshift{4.139099in}{5.654321in}%
\pgfsys@useobject{currentmarker}{}%
\end{pgfscope}%
\end{pgfscope}%
\begin{pgfscope}%
\pgfpathrectangle{\pgfqpoint{0.977909in}{2.928000in}}{\pgfqpoint{5.664000in}{5.664000in}}%
\pgfusepath{clip}%
\pgfsetbuttcap%
\pgfsetroundjoin%
\definecolor{currentfill}{rgb}{0.768627,0.603922,0.000000}%
\pgfsetfillcolor{currentfill}%
\pgfsetlinewidth{1.254687pt}%
\definecolor{currentstroke}{rgb}{0.768627,0.603922,0.000000}%
\pgfsetstrokecolor{currentstroke}%
\pgfsetdash{}{0pt}%
\pgfsys@defobject{currentmarker}{\pgfqpoint{-0.043368in}{-0.043368in}}{\pgfqpoint{0.043368in}{0.043368in}}{%
\pgfpathmoveto{\pgfqpoint{-0.043368in}{-0.043368in}}%
\pgfpathlineto{\pgfqpoint{0.043368in}{0.043368in}}%
\pgfpathmoveto{\pgfqpoint{-0.043368in}{0.043368in}}%
\pgfpathlineto{\pgfqpoint{0.043368in}{-0.043368in}}%
\pgfusepath{stroke,fill}%
}%
\begin{pgfscope}%
\pgfsys@transformshift{4.047063in}{5.556865in}%
\pgfsys@useobject{currentmarker}{}%
\end{pgfscope}%
\begin{pgfscope}%
\pgfsys@transformshift{3.659588in}{5.449405in}%
\pgfsys@useobject{currentmarker}{}%
\end{pgfscope}%
\end{pgfscope}%
\begin{pgfscope}%
\pgfpathrectangle{\pgfqpoint{0.977909in}{2.928000in}}{\pgfqpoint{5.664000in}{5.664000in}}%
\pgfusepath{clip}%
\pgfsetbuttcap%
\pgfsetroundjoin%
\definecolor{currentfill}{rgb}{0.768627,0.603922,0.000000}%
\pgfsetfillcolor{currentfill}%
\pgfsetlinewidth{1.254687pt}%
\definecolor{currentstroke}{rgb}{0.768627,0.603922,0.000000}%
\pgfsetstrokecolor{currentstroke}%
\pgfsetdash{}{0pt}%
\pgfsys@defobject{currentmarker}{\pgfqpoint{-0.043368in}{-0.043368in}}{\pgfqpoint{0.043368in}{0.043368in}}{%
\pgfpathmoveto{\pgfqpoint{-0.043368in}{-0.043368in}}%
\pgfpathlineto{\pgfqpoint{0.043368in}{0.043368in}}%
\pgfpathmoveto{\pgfqpoint{-0.043368in}{0.043368in}}%
\pgfpathlineto{\pgfqpoint{0.043368in}{-0.043368in}}%
\pgfusepath{stroke,fill}%
}%
\begin{pgfscope}%
\pgfsys@transformshift{4.557334in}{6.497306in}%
\pgfsys@useobject{currentmarker}{}%
\end{pgfscope}%
\begin{pgfscope}%
\pgfsys@transformshift{4.230764in}{6.098239in}%
\pgfsys@useobject{currentmarker}{}%
\end{pgfscope}%
\end{pgfscope}%
\begin{pgfscope}%
\pgfpathrectangle{\pgfqpoint{0.977909in}{2.928000in}}{\pgfqpoint{5.664000in}{5.664000in}}%
\pgfusepath{clip}%
\pgfsetbuttcap%
\pgfsetroundjoin%
\definecolor{currentfill}{rgb}{0.058824,0.541176,0.372549}%
\pgfsetfillcolor{currentfill}%
\pgfsetlinewidth{1.254687pt}%
\definecolor{currentstroke}{rgb}{0.058824,0.541176,0.372549}%
\pgfsetstrokecolor{currentstroke}%
\pgfsetdash{}{0pt}%
\pgfsys@defobject{currentmarker}{\pgfqpoint{-0.043368in}{-0.043368in}}{\pgfqpoint{0.043368in}{0.043368in}}{%
\pgfpathmoveto{\pgfqpoint{-0.043368in}{-0.043368in}}%
\pgfpathlineto{\pgfqpoint{0.043368in}{0.043368in}}%
\pgfpathmoveto{\pgfqpoint{-0.043368in}{0.043368in}}%
\pgfpathlineto{\pgfqpoint{0.043368in}{-0.043368in}}%
\pgfusepath{stroke,fill}%
}%
\begin{pgfscope}%
\pgfsys@transformshift{4.192989in}{6.668654in}%
\pgfsys@useobject{currentmarker}{}%
\end{pgfscope}%
\begin{pgfscope}%
\pgfsys@transformshift{3.880342in}{5.751343in}%
\pgfsys@useobject{currentmarker}{}%
\end{pgfscope}%
\begin{pgfscope}%
\pgfsys@transformshift{3.986397in}{5.180637in}%
\pgfsys@useobject{currentmarker}{}%
\end{pgfscope}%
\begin{pgfscope}%
\pgfsys@transformshift{4.466248in}{7.584130in}%
\pgfsys@useobject{currentmarker}{}%
\end{pgfscope}%
\begin{pgfscope}%
\pgfsys@transformshift{4.775456in}{7.151391in}%
\pgfsys@useobject{currentmarker}{}%
\end{pgfscope}%
\begin{pgfscope}%
\pgfsys@transformshift{3.919142in}{5.561933in}%
\pgfsys@useobject{currentmarker}{}%
\end{pgfscope}%
\begin{pgfscope}%
\pgfsys@transformshift{3.677338in}{4.977537in}%
\pgfsys@useobject{currentmarker}{}%
\end{pgfscope}%
\begin{pgfscope}%
\pgfsys@transformshift{3.898086in}{4.998285in}%
\pgfsys@useobject{currentmarker}{}%
\end{pgfscope}%
\end{pgfscope}%
\begin{pgfscope}%
\pgfsetbuttcap%
\pgfsetroundjoin%
\definecolor{currentfill}{rgb}{0.337255,0.705882,0.913725}%
\pgfsetfillcolor{currentfill}%
\pgfsetlinewidth{1.254687pt}%
\definecolor{currentstroke}{rgb}{0.337255,0.705882,0.913725}%
\pgfsetstrokecolor{currentstroke}%
\pgfsetdash{}{0pt}%
\pgfsys@defobject{currentmarker}{\pgfqpoint{-0.043368in}{-0.043368in}}{\pgfqpoint{0.043368in}{0.043368in}}{%
\pgfpathmoveto{\pgfqpoint{-0.043368in}{-0.043368in}}%
\pgfpathlineto{\pgfqpoint{0.043368in}{0.043368in}}%
\pgfpathmoveto{\pgfqpoint{-0.043368in}{0.043368in}}%
\pgfpathlineto{\pgfqpoint{0.043368in}{-0.043368in}}%
\pgfusepath{stroke,fill}%
}%
\begin{pgfscope}%
\pgfsys@transformshift{6.641909in}{8.592000in}%
\pgfsys@useobject{currentmarker}{}%
\end{pgfscope}%
\end{pgfscope}%
\begin{pgfscope}%
\pgfpathrectangle{\pgfqpoint{0.977909in}{2.928000in}}{\pgfqpoint{5.664000in}{5.664000in}}%
\pgfusepath{clip}%
\pgfsetbuttcap%
\pgfsetroundjoin%
\definecolor{currentfill}{rgb}{0.901961,0.196078,0.156863}%
\pgfsetfillcolor{currentfill}%
\pgfsetlinewidth{1.254687pt}%
\definecolor{currentstroke}{rgb}{0.901961,0.196078,0.156863}%
\pgfsetstrokecolor{currentstroke}%
\pgfsetdash{}{0pt}%
\pgfsys@defobject{currentmarker}{\pgfqpoint{-0.043368in}{-0.043368in}}{\pgfqpoint{0.043368in}{0.043368in}}{%
\pgfpathmoveto{\pgfqpoint{-0.043368in}{-0.043368in}}%
\pgfpathlineto{\pgfqpoint{0.043368in}{0.043368in}}%
\pgfpathmoveto{\pgfqpoint{-0.043368in}{0.043368in}}%
\pgfpathlineto{\pgfqpoint{0.043368in}{-0.043368in}}%
\pgfusepath{stroke,fill}%
}%
\begin{pgfscope}%
\pgfsys@transformshift{4.803042in}{7.809822in}%
\pgfsys@useobject{currentmarker}{}%
\end{pgfscope}%
\begin{pgfscope}%
\pgfsys@transformshift{3.851610in}{6.040405in}%
\pgfsys@useobject{currentmarker}{}%
\end{pgfscope}%
\end{pgfscope}%
\begin{pgfscope}%
\pgfsetbuttcap%
\pgfsetroundjoin%
\definecolor{currentfill}{rgb}{0.901961,0.196078,0.156863}%
\pgfsetfillcolor{currentfill}%
\pgfsetlinewidth{1.254687pt}%
\definecolor{currentstroke}{rgb}{0.901961,0.196078,0.156863}%
\pgfsetstrokecolor{currentstroke}%
\pgfsetdash{}{0pt}%
\pgfsys@defobject{currentmarker}{\pgfqpoint{-0.043368in}{-0.043368in}}{\pgfqpoint{0.043368in}{0.043368in}}{%
\pgfpathmoveto{\pgfqpoint{-0.043368in}{-0.043368in}}%
\pgfpathlineto{\pgfqpoint{0.043368in}{0.043368in}}%
\pgfpathmoveto{\pgfqpoint{-0.043368in}{0.043368in}}%
\pgfpathlineto{\pgfqpoint{0.043368in}{-0.043368in}}%
\pgfusepath{stroke,fill}%
}%
\begin{pgfscope}%
\pgfsys@transformshift{6.641909in}{8.288866in}%
\pgfsys@useobject{currentmarker}{}%
\end{pgfscope}%
\end{pgfscope}%
\begin{pgfscope}%
\pgfsetbuttcap%
\pgfsetroundjoin%
\definecolor{currentfill}{rgb}{0.901961,0.196078,0.156863}%
\pgfsetfillcolor{currentfill}%
\pgfsetlinewidth{1.254687pt}%
\definecolor{currentstroke}{rgb}{0.901961,0.196078,0.156863}%
\pgfsetstrokecolor{currentstroke}%
\pgfsetdash{}{0pt}%
\pgfsys@defobject{currentmarker}{\pgfqpoint{-0.043368in}{-0.043368in}}{\pgfqpoint{0.043368in}{0.043368in}}{%
\pgfpathmoveto{\pgfqpoint{-0.043368in}{-0.043368in}}%
\pgfpathlineto{\pgfqpoint{0.043368in}{0.043368in}}%
\pgfpathmoveto{\pgfqpoint{-0.043368in}{0.043368in}}%
\pgfpathlineto{\pgfqpoint{0.043368in}{-0.043368in}}%
\pgfusepath{stroke,fill}%
}%
\begin{pgfscope}%
\pgfsys@transformshift{6.248891in}{8.592000in}%
\pgfsys@useobject{currentmarker}{}%
\end{pgfscope}%
\end{pgfscope}%
\begin{pgfscope}%
\pgfpathrectangle{\pgfqpoint{0.977909in}{2.928000in}}{\pgfqpoint{5.664000in}{5.664000in}}%
\pgfusepath{clip}%
\pgfsetbuttcap%
\pgfsetroundjoin%
\definecolor{currentfill}{rgb}{0.835294,0.368627,0.000000}%
\pgfsetfillcolor{currentfill}%
\pgfsetlinewidth{1.254687pt}%
\definecolor{currentstroke}{rgb}{0.835294,0.368627,0.000000}%
\pgfsetstrokecolor{currentstroke}%
\pgfsetdash{}{0pt}%
\pgfsys@defobject{currentmarker}{\pgfqpoint{-0.043368in}{-0.043368in}}{\pgfqpoint{0.043368in}{0.043368in}}{%
\pgfpathmoveto{\pgfqpoint{-0.043368in}{-0.043368in}}%
\pgfpathlineto{\pgfqpoint{0.043368in}{0.043368in}}%
\pgfpathmoveto{\pgfqpoint{-0.043368in}{0.043368in}}%
\pgfpathlineto{\pgfqpoint{0.043368in}{-0.043368in}}%
\pgfusepath{stroke,fill}%
}%
\begin{pgfscope}%
\pgfsys@transformshift{4.839167in}{5.622188in}%
\pgfsys@useobject{currentmarker}{}%
\end{pgfscope}%
\end{pgfscope}%
\begin{pgfscope}%
\pgfpathrectangle{\pgfqpoint{0.977909in}{2.928000in}}{\pgfqpoint{5.664000in}{5.664000in}}%
\pgfusepath{clip}%
\pgfsetbuttcap%
\pgfsetroundjoin%
\definecolor{currentfill}{rgb}{0.541176,0.380392,0.658824}%
\pgfsetfillcolor{currentfill}%
\pgfsetfillopacity{0.900000}%
\pgfsetlinewidth{0.803000pt}%
\definecolor{currentstroke}{rgb}{0.541176,0.380392,0.658824}%
\pgfsetstrokecolor{currentstroke}%
\pgfsetstrokeopacity{0.900000}%
\pgfsetdash{}{0pt}%
\pgfsys@defobject{currentmarker}{\pgfqpoint{-0.042241in}{-0.042241in}}{\pgfqpoint{0.042241in}{0.042241in}}{%
\pgfpathmoveto{\pgfqpoint{-0.042241in}{-0.042241in}}%
\pgfpathlineto{\pgfqpoint{0.042241in}{-0.042241in}}%
\pgfpathlineto{\pgfqpoint{0.042241in}{0.042241in}}%
\pgfpathlineto{\pgfqpoint{-0.042241in}{0.042241in}}%
\pgfpathlineto{\pgfqpoint{-0.042241in}{-0.042241in}}%
\pgfpathclose%
\pgfusepath{stroke,fill}%
}%
\begin{pgfscope}%
\pgfsys@transformshift{4.985794in}{5.789507in}%
\pgfsys@useobject{currentmarker}{}%
\end{pgfscope}%
\end{pgfscope}%
\begin{pgfscope}%
\pgfpathrectangle{\pgfqpoint{0.977909in}{2.928000in}}{\pgfqpoint{5.664000in}{5.664000in}}%
\pgfusepath{clip}%
\pgfsetbuttcap%
\pgfsetroundjoin%
\pgfsetlinewidth{1.003750pt}%
\definecolor{currentstroke}{rgb}{0.541176,0.380392,0.658824}%
\pgfsetstrokecolor{currentstroke}%
\pgfsetstrokeopacity{0.900000}%
\pgfsetdash{}{0pt}%
\pgfpathmoveto{\pgfqpoint{5.510810in}{6.979136in}}%
\pgfpathlineto{\pgfqpoint{5.595293in}{6.979136in}}%
\pgfpathlineto{\pgfqpoint{5.595293in}{7.063619in}}%
\pgfpathlineto{\pgfqpoint{5.510810in}{7.063619in}}%
\pgfpathlineto{\pgfqpoint{5.510810in}{6.979136in}}%
\pgfpathclose%
\pgfusepath{stroke}%
\end{pgfscope}%
\begin{pgfscope}%
\pgfpathrectangle{\pgfqpoint{0.977909in}{2.928000in}}{\pgfqpoint{5.664000in}{5.664000in}}%
\pgfusepath{clip}%
\pgfsetbuttcap%
\pgfsetroundjoin%
\pgfsetlinewidth{1.003750pt}%
\definecolor{currentstroke}{rgb}{0.541176,0.380392,0.658824}%
\pgfsetstrokecolor{currentstroke}%
\pgfsetstrokeopacity{0.900000}%
\pgfsetdash{}{0pt}%
\pgfpathmoveto{\pgfqpoint{4.780765in}{6.324979in}}%
\pgfpathlineto{\pgfqpoint{4.865248in}{6.324979in}}%
\pgfpathlineto{\pgfqpoint{4.865248in}{6.409462in}}%
\pgfpathlineto{\pgfqpoint{4.780765in}{6.409462in}}%
\pgfpathlineto{\pgfqpoint{4.780765in}{6.324979in}}%
\pgfpathclose%
\pgfusepath{stroke}%
\end{pgfscope}%
\begin{pgfscope}%
\pgfpathrectangle{\pgfqpoint{0.977909in}{2.928000in}}{\pgfqpoint{5.664000in}{5.664000in}}%
\pgfusepath{clip}%
\pgfsetbuttcap%
\pgfsetroundjoin%
\pgfsetlinewidth{1.003750pt}%
\definecolor{currentstroke}{rgb}{0.541176,0.380392,0.658824}%
\pgfsetstrokecolor{currentstroke}%
\pgfsetstrokeopacity{0.900000}%
\pgfsetdash{}{0pt}%
\pgfpathmoveto{\pgfqpoint{4.549835in}{6.092251in}}%
\pgfpathlineto{\pgfqpoint{4.634318in}{6.092251in}}%
\pgfpathlineto{\pgfqpoint{4.634318in}{6.176734in}}%
\pgfpathlineto{\pgfqpoint{4.549835in}{6.176734in}}%
\pgfpathlineto{\pgfqpoint{4.549835in}{6.092251in}}%
\pgfpathclose%
\pgfusepath{stroke}%
\end{pgfscope}%
\begin{pgfscope}%
\pgfpathrectangle{\pgfqpoint{0.977909in}{2.928000in}}{\pgfqpoint{5.664000in}{5.664000in}}%
\pgfusepath{clip}%
\pgfsetbuttcap%
\pgfsetroundjoin%
\pgfsetlinewidth{1.003750pt}%
\definecolor{currentstroke}{rgb}{0.541176,0.380392,0.658824}%
\pgfsetstrokecolor{currentstroke}%
\pgfsetstrokeopacity{0.900000}%
\pgfsetdash{}{0pt}%
\pgfpathmoveto{\pgfqpoint{4.812923in}{7.075340in}}%
\pgfpathlineto{\pgfqpoint{4.897406in}{7.075340in}}%
\pgfpathlineto{\pgfqpoint{4.897406in}{7.159823in}}%
\pgfpathlineto{\pgfqpoint{4.812923in}{7.159823in}}%
\pgfpathlineto{\pgfqpoint{4.812923in}{7.075340in}}%
\pgfpathclose%
\pgfusepath{stroke}%
\end{pgfscope}%
\begin{pgfscope}%
\pgfpathrectangle{\pgfqpoint{0.977909in}{2.928000in}}{\pgfqpoint{5.664000in}{5.664000in}}%
\pgfusepath{clip}%
\pgfsetbuttcap%
\pgfsetroundjoin%
\pgfsetlinewidth{1.003750pt}%
\definecolor{currentstroke}{rgb}{0.541176,0.380392,0.658824}%
\pgfsetstrokecolor{currentstroke}%
\pgfsetstrokeopacity{0.900000}%
\pgfsetdash{}{0pt}%
\pgfpathmoveto{\pgfqpoint{4.372014in}{6.826848in}}%
\pgfpathlineto{\pgfqpoint{4.456496in}{6.826848in}}%
\pgfpathlineto{\pgfqpoint{4.456496in}{6.911331in}}%
\pgfpathlineto{\pgfqpoint{4.372014in}{6.911331in}}%
\pgfpathlineto{\pgfqpoint{4.372014in}{6.826848in}}%
\pgfpathclose%
\pgfusepath{stroke}%
\end{pgfscope}%
\begin{pgfscope}%
\pgfpathrectangle{\pgfqpoint{0.977909in}{2.928000in}}{\pgfqpoint{5.664000in}{5.664000in}}%
\pgfusepath{clip}%
\pgfsetbuttcap%
\pgfsetroundjoin%
\pgfsetlinewidth{1.003750pt}%
\definecolor{currentstroke}{rgb}{0.541176,0.380392,0.658824}%
\pgfsetstrokecolor{currentstroke}%
\pgfsetstrokeopacity{0.900000}%
\pgfsetdash{}{0pt}%
\pgfpathmoveto{\pgfqpoint{4.426589in}{6.886374in}}%
\pgfpathlineto{\pgfqpoint{4.511071in}{6.886374in}}%
\pgfpathlineto{\pgfqpoint{4.511071in}{6.970857in}}%
\pgfpathlineto{\pgfqpoint{4.426589in}{6.970857in}}%
\pgfpathlineto{\pgfqpoint{4.426589in}{6.886374in}}%
\pgfpathclose%
\pgfusepath{stroke}%
\end{pgfscope}%
\begin{pgfscope}%
\pgfpathrectangle{\pgfqpoint{0.977909in}{2.928000in}}{\pgfqpoint{5.664000in}{5.664000in}}%
\pgfusepath{clip}%
\pgfsetbuttcap%
\pgfsetroundjoin%
\pgfsetlinewidth{1.003750pt}%
\definecolor{currentstroke}{rgb}{0.541176,0.380392,0.658824}%
\pgfsetstrokecolor{currentstroke}%
\pgfsetstrokeopacity{0.900000}%
\pgfsetdash{}{0pt}%
\pgfpathmoveto{\pgfqpoint{4.919224in}{6.412552in}}%
\pgfpathlineto{\pgfqpoint{5.003707in}{6.412552in}}%
\pgfpathlineto{\pgfqpoint{5.003707in}{6.497035in}}%
\pgfpathlineto{\pgfqpoint{4.919224in}{6.497035in}}%
\pgfpathlineto{\pgfqpoint{4.919224in}{6.412552in}}%
\pgfpathclose%
\pgfusepath{stroke}%
\end{pgfscope}%
\begin{pgfscope}%
\pgfpathrectangle{\pgfqpoint{0.977909in}{2.928000in}}{\pgfqpoint{5.664000in}{5.664000in}}%
\pgfusepath{clip}%
\pgfsetbuttcap%
\pgfsetroundjoin%
\pgfsetlinewidth{1.003750pt}%
\definecolor{currentstroke}{rgb}{0.541176,0.380392,0.658824}%
\pgfsetstrokecolor{currentstroke}%
\pgfsetstrokeopacity{0.900000}%
\pgfsetdash{}{0pt}%
\pgfpathmoveto{\pgfqpoint{4.068964in}{5.604697in}}%
\pgfpathlineto{\pgfqpoint{4.153447in}{5.604697in}}%
\pgfpathlineto{\pgfqpoint{4.153447in}{5.689180in}}%
\pgfpathlineto{\pgfqpoint{4.068964in}{5.689180in}}%
\pgfpathlineto{\pgfqpoint{4.068964in}{5.604697in}}%
\pgfpathclose%
\pgfusepath{stroke}%
\end{pgfscope}%
\begin{pgfscope}%
\pgfpathrectangle{\pgfqpoint{0.977909in}{2.928000in}}{\pgfqpoint{5.664000in}{5.664000in}}%
\pgfusepath{clip}%
\pgfsetbuttcap%
\pgfsetroundjoin%
\pgfsetlinewidth{1.003750pt}%
\definecolor{currentstroke}{rgb}{0.541176,0.380392,0.658824}%
\pgfsetstrokecolor{currentstroke}%
\pgfsetstrokeopacity{0.900000}%
\pgfsetdash{}{0pt}%
\pgfpathmoveto{\pgfqpoint{3.362202in}{4.963458in}}%
\pgfpathlineto{\pgfqpoint{3.446684in}{4.963458in}}%
\pgfpathlineto{\pgfqpoint{3.446684in}{5.047941in}}%
\pgfpathlineto{\pgfqpoint{3.362202in}{5.047941in}}%
\pgfpathlineto{\pgfqpoint{3.362202in}{4.963458in}}%
\pgfpathclose%
\pgfusepath{stroke}%
\end{pgfscope}%
\begin{pgfscope}%
\pgfpathrectangle{\pgfqpoint{0.977909in}{2.928000in}}{\pgfqpoint{5.664000in}{5.664000in}}%
\pgfusepath{clip}%
\pgfsetbuttcap%
\pgfsetroundjoin%
\pgfsetlinewidth{1.003750pt}%
\definecolor{currentstroke}{rgb}{0.541176,0.380392,0.658824}%
\pgfsetstrokecolor{currentstroke}%
\pgfsetstrokeopacity{0.900000}%
\pgfsetdash{}{0pt}%
\pgfpathmoveto{\pgfqpoint{4.525413in}{6.153606in}}%
\pgfpathlineto{\pgfqpoint{4.609896in}{6.153606in}}%
\pgfpathlineto{\pgfqpoint{4.609896in}{6.238089in}}%
\pgfpathlineto{\pgfqpoint{4.525413in}{6.238089in}}%
\pgfpathlineto{\pgfqpoint{4.525413in}{6.153606in}}%
\pgfpathclose%
\pgfusepath{stroke}%
\end{pgfscope}%
\begin{pgfscope}%
\pgfpathrectangle{\pgfqpoint{0.977909in}{2.928000in}}{\pgfqpoint{5.664000in}{5.664000in}}%
\pgfusepath{clip}%
\pgfsetbuttcap%
\pgfsetroundjoin%
\pgfsetlinewidth{1.003750pt}%
\definecolor{currentstroke}{rgb}{0.541176,0.380392,0.658824}%
\pgfsetstrokecolor{currentstroke}%
\pgfsetstrokeopacity{0.900000}%
\pgfsetdash{}{0pt}%
\pgfpathmoveto{\pgfqpoint{3.336530in}{5.086554in}}%
\pgfpathlineto{\pgfqpoint{3.421012in}{5.086554in}}%
\pgfpathlineto{\pgfqpoint{3.421012in}{5.171037in}}%
\pgfpathlineto{\pgfqpoint{3.336530in}{5.171037in}}%
\pgfpathlineto{\pgfqpoint{3.336530in}{5.086554in}}%
\pgfpathclose%
\pgfusepath{stroke}%
\end{pgfscope}%
\begin{pgfscope}%
\pgfpathrectangle{\pgfqpoint{0.977909in}{2.928000in}}{\pgfqpoint{5.664000in}{5.664000in}}%
\pgfusepath{clip}%
\pgfsetbuttcap%
\pgfsetroundjoin%
\pgfsetlinewidth{1.003750pt}%
\definecolor{currentstroke}{rgb}{0.541176,0.380392,0.658824}%
\pgfsetstrokecolor{currentstroke}%
\pgfsetstrokeopacity{0.900000}%
\pgfsetdash{}{0pt}%
\pgfpathmoveto{\pgfqpoint{2.876142in}{4.829447in}}%
\pgfpathlineto{\pgfqpoint{2.960625in}{4.829447in}}%
\pgfpathlineto{\pgfqpoint{2.960625in}{4.913930in}}%
\pgfpathlineto{\pgfqpoint{2.876142in}{4.913930in}}%
\pgfpathlineto{\pgfqpoint{2.876142in}{4.829447in}}%
\pgfpathclose%
\pgfusepath{stroke}%
\end{pgfscope}%
\begin{pgfscope}%
\pgfpathrectangle{\pgfqpoint{0.977909in}{2.928000in}}{\pgfqpoint{5.664000in}{5.664000in}}%
\pgfusepath{clip}%
\pgfsetbuttcap%
\pgfsetroundjoin%
\pgfsetlinewidth{1.003750pt}%
\definecolor{currentstroke}{rgb}{0.541176,0.380392,0.658824}%
\pgfsetstrokecolor{currentstroke}%
\pgfsetstrokeopacity{0.900000}%
\pgfsetdash{}{0pt}%
\pgfpathmoveto{\pgfqpoint{3.284443in}{4.957330in}}%
\pgfpathlineto{\pgfqpoint{3.368926in}{4.957330in}}%
\pgfpathlineto{\pgfqpoint{3.368926in}{5.041813in}}%
\pgfpathlineto{\pgfqpoint{3.284443in}{5.041813in}}%
\pgfpathlineto{\pgfqpoint{3.284443in}{4.957330in}}%
\pgfpathclose%
\pgfusepath{stroke}%
\end{pgfscope}%
\begin{pgfscope}%
\pgfpathrectangle{\pgfqpoint{0.977909in}{2.928000in}}{\pgfqpoint{5.664000in}{5.664000in}}%
\pgfusepath{clip}%
\pgfsetbuttcap%
\pgfsetroundjoin%
\pgfsetlinewidth{1.003750pt}%
\definecolor{currentstroke}{rgb}{0.541176,0.380392,0.658824}%
\pgfsetstrokecolor{currentstroke}%
\pgfsetstrokeopacity{0.900000}%
\pgfsetdash{}{0pt}%
\pgfpathmoveto{\pgfqpoint{2.836008in}{4.709914in}}%
\pgfpathlineto{\pgfqpoint{2.920490in}{4.709914in}}%
\pgfpathlineto{\pgfqpoint{2.920490in}{4.794397in}}%
\pgfpathlineto{\pgfqpoint{2.836008in}{4.794397in}}%
\pgfpathlineto{\pgfqpoint{2.836008in}{4.709914in}}%
\pgfpathclose%
\pgfusepath{stroke}%
\end{pgfscope}%
\begin{pgfscope}%
\pgfpathrectangle{\pgfqpoint{0.977909in}{2.928000in}}{\pgfqpoint{5.664000in}{5.664000in}}%
\pgfusepath{clip}%
\pgfsetbuttcap%
\pgfsetroundjoin%
\pgfsetlinewidth{1.003750pt}%
\definecolor{currentstroke}{rgb}{0.541176,0.380392,0.658824}%
\pgfsetstrokecolor{currentstroke}%
\pgfsetstrokeopacity{0.900000}%
\pgfsetdash{}{0pt}%
\pgfpathmoveto{\pgfqpoint{4.375470in}{6.301232in}}%
\pgfpathlineto{\pgfqpoint{4.459953in}{6.301232in}}%
\pgfpathlineto{\pgfqpoint{4.459953in}{6.385715in}}%
\pgfpathlineto{\pgfqpoint{4.375470in}{6.385715in}}%
\pgfpathlineto{\pgfqpoint{4.375470in}{6.301232in}}%
\pgfpathclose%
\pgfusepath{stroke}%
\end{pgfscope}%
\begin{pgfscope}%
\pgfpathrectangle{\pgfqpoint{0.977909in}{2.928000in}}{\pgfqpoint{5.664000in}{5.664000in}}%
\pgfusepath{clip}%
\pgfsetbuttcap%
\pgfsetroundjoin%
\pgfsetlinewidth{1.003750pt}%
\definecolor{currentstroke}{rgb}{0.541176,0.380392,0.658824}%
\pgfsetstrokecolor{currentstroke}%
\pgfsetstrokeopacity{0.900000}%
\pgfsetdash{}{0pt}%
\pgfpathmoveto{\pgfqpoint{3.876498in}{5.876046in}}%
\pgfpathlineto{\pgfqpoint{3.960981in}{5.876046in}}%
\pgfpathlineto{\pgfqpoint{3.960981in}{5.960529in}}%
\pgfpathlineto{\pgfqpoint{3.876498in}{5.960529in}}%
\pgfpathlineto{\pgfqpoint{3.876498in}{5.876046in}}%
\pgfpathclose%
\pgfusepath{stroke}%
\end{pgfscope}%
\begin{pgfscope}%
\pgfpathrectangle{\pgfqpoint{0.977909in}{2.928000in}}{\pgfqpoint{5.664000in}{5.664000in}}%
\pgfusepath{clip}%
\pgfsetbuttcap%
\pgfsetroundjoin%
\pgfsetlinewidth{1.003750pt}%
\definecolor{currentstroke}{rgb}{0.541176,0.380392,0.658824}%
\pgfsetstrokecolor{currentstroke}%
\pgfsetstrokeopacity{0.900000}%
\pgfsetdash{}{0pt}%
\pgfpathmoveto{\pgfqpoint{4.905921in}{6.730014in}}%
\pgfpathlineto{\pgfqpoint{4.990404in}{6.730014in}}%
\pgfpathlineto{\pgfqpoint{4.990404in}{6.814497in}}%
\pgfpathlineto{\pgfqpoint{4.905921in}{6.814497in}}%
\pgfpathlineto{\pgfqpoint{4.905921in}{6.730014in}}%
\pgfpathclose%
\pgfusepath{stroke}%
\end{pgfscope}%
\begin{pgfscope}%
\pgfpathrectangle{\pgfqpoint{0.977909in}{2.928000in}}{\pgfqpoint{5.664000in}{5.664000in}}%
\pgfusepath{clip}%
\pgfsetbuttcap%
\pgfsetroundjoin%
\pgfsetlinewidth{1.003750pt}%
\definecolor{currentstroke}{rgb}{0.541176,0.380392,0.658824}%
\pgfsetstrokecolor{currentstroke}%
\pgfsetstrokeopacity{0.900000}%
\pgfsetdash{}{0pt}%
\pgfpathmoveto{\pgfqpoint{3.215293in}{5.000692in}}%
\pgfpathlineto{\pgfqpoint{3.299776in}{5.000692in}}%
\pgfpathlineto{\pgfqpoint{3.299776in}{5.085175in}}%
\pgfpathlineto{\pgfqpoint{3.215293in}{5.085175in}}%
\pgfpathlineto{\pgfqpoint{3.215293in}{5.000692in}}%
\pgfpathclose%
\pgfusepath{stroke}%
\end{pgfscope}%
\begin{pgfscope}%
\pgfpathrectangle{\pgfqpoint{0.977909in}{2.928000in}}{\pgfqpoint{5.664000in}{5.664000in}}%
\pgfusepath{clip}%
\pgfsetbuttcap%
\pgfsetroundjoin%
\pgfsetlinewidth{1.003750pt}%
\definecolor{currentstroke}{rgb}{0.541176,0.380392,0.658824}%
\pgfsetstrokecolor{currentstroke}%
\pgfsetstrokeopacity{0.900000}%
\pgfsetdash{}{0pt}%
\pgfpathmoveto{\pgfqpoint{2.412577in}{4.691543in}}%
\pgfpathlineto{\pgfqpoint{2.497060in}{4.691543in}}%
\pgfpathlineto{\pgfqpoint{2.497060in}{4.776025in}}%
\pgfpathlineto{\pgfqpoint{2.412577in}{4.776025in}}%
\pgfpathlineto{\pgfqpoint{2.412577in}{4.691543in}}%
\pgfpathclose%
\pgfusepath{stroke}%
\end{pgfscope}%
\begin{pgfscope}%
\pgfpathrectangle{\pgfqpoint{0.977909in}{2.928000in}}{\pgfqpoint{5.664000in}{5.664000in}}%
\pgfusepath{clip}%
\pgfsetbuttcap%
\pgfsetroundjoin%
\pgfsetlinewidth{1.003750pt}%
\definecolor{currentstroke}{rgb}{0.541176,0.380392,0.658824}%
\pgfsetstrokecolor{currentstroke}%
\pgfsetstrokeopacity{0.900000}%
\pgfsetdash{}{0pt}%
\pgfpathmoveto{\pgfqpoint{3.683680in}{5.423582in}}%
\pgfpathlineto{\pgfqpoint{3.768163in}{5.423582in}}%
\pgfpathlineto{\pgfqpoint{3.768163in}{5.508065in}}%
\pgfpathlineto{\pgfqpoint{3.683680in}{5.508065in}}%
\pgfpathlineto{\pgfqpoint{3.683680in}{5.423582in}}%
\pgfpathclose%
\pgfusepath{stroke}%
\end{pgfscope}%
\begin{pgfscope}%
\pgfpathrectangle{\pgfqpoint{0.977909in}{2.928000in}}{\pgfqpoint{5.664000in}{5.664000in}}%
\pgfusepath{clip}%
\pgfsetbuttcap%
\pgfsetroundjoin%
\pgfsetlinewidth{1.003750pt}%
\definecolor{currentstroke}{rgb}{0.541176,0.380392,0.658824}%
\pgfsetstrokecolor{currentstroke}%
\pgfsetstrokeopacity{0.900000}%
\pgfsetdash{}{0pt}%
\pgfpathmoveto{\pgfqpoint{2.725034in}{4.597229in}}%
\pgfpathlineto{\pgfqpoint{2.809517in}{4.597229in}}%
\pgfpathlineto{\pgfqpoint{2.809517in}{4.681712in}}%
\pgfpathlineto{\pgfqpoint{2.725034in}{4.681712in}}%
\pgfpathlineto{\pgfqpoint{2.725034in}{4.597229in}}%
\pgfpathclose%
\pgfusepath{stroke}%
\end{pgfscope}%
\begin{pgfscope}%
\pgfpathrectangle{\pgfqpoint{0.977909in}{2.928000in}}{\pgfqpoint{5.664000in}{5.664000in}}%
\pgfusepath{clip}%
\pgfsetbuttcap%
\pgfsetroundjoin%
\pgfsetlinewidth{1.003750pt}%
\definecolor{currentstroke}{rgb}{0.541176,0.380392,0.658824}%
\pgfsetstrokecolor{currentstroke}%
\pgfsetstrokeopacity{0.900000}%
\pgfsetdash{}{0pt}%
\pgfpathmoveto{\pgfqpoint{3.762370in}{5.858481in}}%
\pgfpathlineto{\pgfqpoint{3.846853in}{5.858481in}}%
\pgfpathlineto{\pgfqpoint{3.846853in}{5.942964in}}%
\pgfpathlineto{\pgfqpoint{3.762370in}{5.942964in}}%
\pgfpathlineto{\pgfqpoint{3.762370in}{5.858481in}}%
\pgfpathclose%
\pgfusepath{stroke}%
\end{pgfscope}%
\begin{pgfscope}%
\pgfpathrectangle{\pgfqpoint{0.977909in}{2.928000in}}{\pgfqpoint{5.664000in}{5.664000in}}%
\pgfusepath{clip}%
\pgfsetbuttcap%
\pgfsetroundjoin%
\pgfsetlinewidth{1.003750pt}%
\definecolor{currentstroke}{rgb}{0.541176,0.380392,0.658824}%
\pgfsetstrokecolor{currentstroke}%
\pgfsetstrokeopacity{0.900000}%
\pgfsetdash{}{0pt}%
\pgfpathmoveto{\pgfqpoint{3.830371in}{5.378958in}}%
\pgfpathlineto{\pgfqpoint{3.914854in}{5.378958in}}%
\pgfpathlineto{\pgfqpoint{3.914854in}{5.463441in}}%
\pgfpathlineto{\pgfqpoint{3.830371in}{5.463441in}}%
\pgfpathlineto{\pgfqpoint{3.830371in}{5.378958in}}%
\pgfpathclose%
\pgfusepath{stroke}%
\end{pgfscope}%
\begin{pgfscope}%
\pgfpathrectangle{\pgfqpoint{0.977909in}{2.928000in}}{\pgfqpoint{5.664000in}{5.664000in}}%
\pgfusepath{clip}%
\pgfsetbuttcap%
\pgfsetroundjoin%
\definecolor{currentfill}{rgb}{0.768627,0.603922,0.000000}%
\pgfsetfillcolor{currentfill}%
\pgfsetfillopacity{0.900000}%
\pgfsetlinewidth{0.803000pt}%
\definecolor{currentstroke}{rgb}{0.768627,0.603922,0.000000}%
\pgfsetstrokecolor{currentstroke}%
\pgfsetstrokeopacity{0.900000}%
\pgfsetdash{}{0pt}%
\pgfsys@defobject{currentmarker}{\pgfqpoint{-0.042241in}{-0.042241in}}{\pgfqpoint{0.042241in}{0.042241in}}{%
\pgfpathmoveto{\pgfqpoint{-0.042241in}{-0.042241in}}%
\pgfpathlineto{\pgfqpoint{0.042241in}{-0.042241in}}%
\pgfpathlineto{\pgfqpoint{0.042241in}{0.042241in}}%
\pgfpathlineto{\pgfqpoint{-0.042241in}{0.042241in}}%
\pgfpathlineto{\pgfqpoint{-0.042241in}{-0.042241in}}%
\pgfpathclose%
\pgfusepath{stroke,fill}%
}%
\begin{pgfscope}%
\pgfsys@transformshift{2.474273in}{5.336620in}%
\pgfsys@useobject{currentmarker}{}%
\end{pgfscope}%
\end{pgfscope}%
\begin{pgfscope}%
\pgfpathrectangle{\pgfqpoint{0.977909in}{2.928000in}}{\pgfqpoint{5.664000in}{5.664000in}}%
\pgfusepath{clip}%
\pgfsetbuttcap%
\pgfsetroundjoin%
\pgfsetlinewidth{1.003750pt}%
\definecolor{currentstroke}{rgb}{0.768627,0.603922,0.000000}%
\pgfsetstrokecolor{currentstroke}%
\pgfsetstrokeopacity{0.900000}%
\pgfsetdash{}{0pt}%
\pgfpathmoveto{\pgfqpoint{3.502931in}{5.852022in}}%
\pgfpathlineto{\pgfqpoint{3.587414in}{5.852022in}}%
\pgfpathlineto{\pgfqpoint{3.587414in}{5.936504in}}%
\pgfpathlineto{\pgfqpoint{3.502931in}{5.936504in}}%
\pgfpathlineto{\pgfqpoint{3.502931in}{5.852022in}}%
\pgfpathclose%
\pgfusepath{stroke}%
\end{pgfscope}%
\begin{pgfscope}%
\pgfpathrectangle{\pgfqpoint{0.977909in}{2.928000in}}{\pgfqpoint{5.664000in}{5.664000in}}%
\pgfusepath{clip}%
\pgfsetbuttcap%
\pgfsetroundjoin%
\pgfsetlinewidth{1.003750pt}%
\definecolor{currentstroke}{rgb}{0.200000,0.133333,0.533333}%
\pgfsetstrokecolor{currentstroke}%
\pgfsetstrokeopacity{0.900000}%
\pgfsetdash{}{0pt}%
\pgfpathmoveto{\pgfqpoint{2.408245in}{4.231161in}}%
\pgfpathlineto{\pgfqpoint{2.492728in}{4.231161in}}%
\pgfpathlineto{\pgfqpoint{2.492728in}{4.315644in}}%
\pgfpathlineto{\pgfqpoint{2.408245in}{4.315644in}}%
\pgfpathlineto{\pgfqpoint{2.408245in}{4.231161in}}%
\pgfpathclose%
\pgfusepath{stroke}%
\end{pgfscope}%
\begin{pgfscope}%
\pgfpathrectangle{\pgfqpoint{0.977909in}{2.928000in}}{\pgfqpoint{5.664000in}{5.664000in}}%
\pgfusepath{clip}%
\pgfsetbuttcap%
\pgfsetroundjoin%
\pgfsetlinewidth{1.003750pt}%
\definecolor{currentstroke}{rgb}{0.200000,0.133333,0.533333}%
\pgfsetstrokecolor{currentstroke}%
\pgfsetstrokeopacity{0.900000}%
\pgfsetdash{}{0pt}%
\pgfpathmoveto{\pgfqpoint{2.546047in}{4.467109in}}%
\pgfpathlineto{\pgfqpoint{2.630529in}{4.467109in}}%
\pgfpathlineto{\pgfqpoint{2.630529in}{4.551591in}}%
\pgfpathlineto{\pgfqpoint{2.546047in}{4.551591in}}%
\pgfpathlineto{\pgfqpoint{2.546047in}{4.467109in}}%
\pgfpathclose%
\pgfusepath{stroke}%
\end{pgfscope}%
\begin{pgfscope}%
\pgfpathrectangle{\pgfqpoint{0.977909in}{2.928000in}}{\pgfqpoint{5.664000in}{5.664000in}}%
\pgfusepath{clip}%
\pgfsetbuttcap%
\pgfsetroundjoin%
\pgfsetlinewidth{1.003750pt}%
\definecolor{currentstroke}{rgb}{0.200000,0.133333,0.533333}%
\pgfsetstrokecolor{currentstroke}%
\pgfsetstrokeopacity{0.900000}%
\pgfsetdash{}{0pt}%
\pgfpathmoveto{\pgfqpoint{3.111152in}{5.037182in}}%
\pgfpathlineto{\pgfqpoint{3.195634in}{5.037182in}}%
\pgfpathlineto{\pgfqpoint{3.195634in}{5.121665in}}%
\pgfpathlineto{\pgfqpoint{3.111152in}{5.121665in}}%
\pgfpathlineto{\pgfqpoint{3.111152in}{5.037182in}}%
\pgfpathclose%
\pgfusepath{stroke}%
\end{pgfscope}%
\begin{pgfscope}%
\pgfpathrectangle{\pgfqpoint{0.977909in}{2.928000in}}{\pgfqpoint{5.664000in}{5.664000in}}%
\pgfusepath{clip}%
\pgfsetbuttcap%
\pgfsetroundjoin%
\pgfsetlinewidth{1.003750pt}%
\definecolor{currentstroke}{rgb}{0.200000,0.133333,0.533333}%
\pgfsetstrokecolor{currentstroke}%
\pgfsetstrokeopacity{0.900000}%
\pgfsetdash{}{0pt}%
\pgfpathmoveto{\pgfqpoint{4.422149in}{6.239688in}}%
\pgfpathlineto{\pgfqpoint{4.506632in}{6.239688in}}%
\pgfpathlineto{\pgfqpoint{4.506632in}{6.324171in}}%
\pgfpathlineto{\pgfqpoint{4.422149in}{6.324171in}}%
\pgfpathlineto{\pgfqpoint{4.422149in}{6.239688in}}%
\pgfpathclose%
\pgfusepath{stroke}%
\end{pgfscope}%
\begin{pgfscope}%
\pgfpathrectangle{\pgfqpoint{0.977909in}{2.928000in}}{\pgfqpoint{5.664000in}{5.664000in}}%
\pgfusepath{clip}%
\pgfsetbuttcap%
\pgfsetroundjoin%
\pgfsetlinewidth{1.003750pt}%
\definecolor{currentstroke}{rgb}{0.200000,0.133333,0.533333}%
\pgfsetstrokecolor{currentstroke}%
\pgfsetstrokeopacity{0.900000}%
\pgfsetdash{}{0pt}%
\pgfpathmoveto{\pgfqpoint{4.092751in}{6.091788in}}%
\pgfpathlineto{\pgfqpoint{4.177234in}{6.091788in}}%
\pgfpathlineto{\pgfqpoint{4.177234in}{6.176271in}}%
\pgfpathlineto{\pgfqpoint{4.092751in}{6.176271in}}%
\pgfpathlineto{\pgfqpoint{4.092751in}{6.091788in}}%
\pgfpathclose%
\pgfusepath{stroke}%
\end{pgfscope}%
\begin{pgfscope}%
\pgfpathrectangle{\pgfqpoint{0.977909in}{2.928000in}}{\pgfqpoint{5.664000in}{5.664000in}}%
\pgfusepath{clip}%
\pgfsetbuttcap%
\pgfsetroundjoin%
\pgfsetlinewidth{1.003750pt}%
\definecolor{currentstroke}{rgb}{0.200000,0.133333,0.533333}%
\pgfsetstrokecolor{currentstroke}%
\pgfsetstrokeopacity{0.900000}%
\pgfsetdash{}{0pt}%
\pgfpathmoveto{\pgfqpoint{3.696664in}{5.727907in}}%
\pgfpathlineto{\pgfqpoint{3.781147in}{5.727907in}}%
\pgfpathlineto{\pgfqpoint{3.781147in}{5.812389in}}%
\pgfpathlineto{\pgfqpoint{3.696664in}{5.812389in}}%
\pgfpathlineto{\pgfqpoint{3.696664in}{5.727907in}}%
\pgfpathclose%
\pgfusepath{stroke}%
\end{pgfscope}%
\begin{pgfscope}%
\pgfpathrectangle{\pgfqpoint{0.977909in}{2.928000in}}{\pgfqpoint{5.664000in}{5.664000in}}%
\pgfusepath{clip}%
\pgfsetbuttcap%
\pgfsetroundjoin%
\definecolor{currentfill}{rgb}{0.176471,0.713725,0.643137}%
\pgfsetfillcolor{currentfill}%
\pgfsetfillopacity{0.900000}%
\pgfsetlinewidth{0.803000pt}%
\definecolor{currentstroke}{rgb}{0.176471,0.713725,0.643137}%
\pgfsetstrokecolor{currentstroke}%
\pgfsetstrokeopacity{0.900000}%
\pgfsetdash{}{0pt}%
\pgfsys@defobject{currentmarker}{\pgfqpoint{-0.042241in}{-0.042241in}}{\pgfqpoint{0.042241in}{0.042241in}}{%
\pgfpathmoveto{\pgfqpoint{-0.042241in}{-0.042241in}}%
\pgfpathlineto{\pgfqpoint{0.042241in}{-0.042241in}}%
\pgfpathlineto{\pgfqpoint{0.042241in}{0.042241in}}%
\pgfpathlineto{\pgfqpoint{-0.042241in}{0.042241in}}%
\pgfpathlineto{\pgfqpoint{-0.042241in}{-0.042241in}}%
\pgfpathclose%
\pgfusepath{stroke,fill}%
}%
\begin{pgfscope}%
\pgfsys@transformshift{2.578341in}{5.644152in}%
\pgfsys@useobject{currentmarker}{}%
\end{pgfscope}%
\end{pgfscope}%
\begin{pgfscope}%
\pgfpathrectangle{\pgfqpoint{0.977909in}{2.928000in}}{\pgfqpoint{5.664000in}{5.664000in}}%
\pgfusepath{clip}%
\pgfsetbuttcap%
\pgfsetroundjoin%
\pgfsetlinewidth{1.003750pt}%
\definecolor{currentstroke}{rgb}{0.176471,0.713725,0.643137}%
\pgfsetstrokecolor{currentstroke}%
\pgfsetstrokeopacity{0.900000}%
\pgfsetdash{}{0pt}%
\pgfpathmoveto{\pgfqpoint{5.291014in}{6.888872in}}%
\pgfpathlineto{\pgfqpoint{5.375497in}{6.888872in}}%
\pgfpathlineto{\pgfqpoint{5.375497in}{6.973355in}}%
\pgfpathlineto{\pgfqpoint{5.291014in}{6.973355in}}%
\pgfpathlineto{\pgfqpoint{5.291014in}{6.888872in}}%
\pgfpathclose%
\pgfusepath{stroke}%
\end{pgfscope}%
\begin{pgfscope}%
\pgfpathrectangle{\pgfqpoint{0.977909in}{2.928000in}}{\pgfqpoint{5.664000in}{5.664000in}}%
\pgfusepath{clip}%
\pgfsetbuttcap%
\pgfsetroundjoin%
\pgfsetlinewidth{1.003750pt}%
\definecolor{currentstroke}{rgb}{0.176471,0.713725,0.643137}%
\pgfsetstrokecolor{currentstroke}%
\pgfsetstrokeopacity{0.900000}%
\pgfsetdash{}{0pt}%
\pgfpathmoveto{\pgfqpoint{4.688051in}{7.198833in}}%
\pgfpathlineto{\pgfqpoint{4.772534in}{7.198833in}}%
\pgfpathlineto{\pgfqpoint{4.772534in}{7.283316in}}%
\pgfpathlineto{\pgfqpoint{4.688051in}{7.283316in}}%
\pgfpathlineto{\pgfqpoint{4.688051in}{7.198833in}}%
\pgfpathclose%
\pgfusepath{stroke}%
\end{pgfscope}%
\begin{pgfscope}%
\pgfpathrectangle{\pgfqpoint{0.977909in}{2.928000in}}{\pgfqpoint{5.664000in}{5.664000in}}%
\pgfusepath{clip}%
\pgfsetbuttcap%
\pgfsetroundjoin%
\pgfsetlinewidth{1.003750pt}%
\definecolor{currentstroke}{rgb}{0.176471,0.713725,0.643137}%
\pgfsetstrokecolor{currentstroke}%
\pgfsetstrokeopacity{0.900000}%
\pgfsetdash{}{0pt}%
\pgfpathmoveto{\pgfqpoint{4.954821in}{6.711419in}}%
\pgfpathlineto{\pgfqpoint{5.039304in}{6.711419in}}%
\pgfpathlineto{\pgfqpoint{5.039304in}{6.795901in}}%
\pgfpathlineto{\pgfqpoint{4.954821in}{6.795901in}}%
\pgfpathlineto{\pgfqpoint{4.954821in}{6.711419in}}%
\pgfpathclose%
\pgfusepath{stroke}%
\end{pgfscope}%
\begin{pgfscope}%
\pgfpathrectangle{\pgfqpoint{0.977909in}{2.928000in}}{\pgfqpoint{5.664000in}{5.664000in}}%
\pgfusepath{clip}%
\pgfsetbuttcap%
\pgfsetroundjoin%
\pgfsetlinewidth{1.003750pt}%
\definecolor{currentstroke}{rgb}{0.176471,0.713725,0.643137}%
\pgfsetstrokecolor{currentstroke}%
\pgfsetstrokeopacity{0.900000}%
\pgfsetdash{}{0pt}%
\pgfpathmoveto{\pgfqpoint{5.030667in}{5.056427in}}%
\pgfpathlineto{\pgfqpoint{5.115149in}{5.056427in}}%
\pgfpathlineto{\pgfqpoint{5.115149in}{5.140910in}}%
\pgfpathlineto{\pgfqpoint{5.030667in}{5.140910in}}%
\pgfpathlineto{\pgfqpoint{5.030667in}{5.056427in}}%
\pgfpathclose%
\pgfusepath{stroke}%
\end{pgfscope}%
\begin{pgfscope}%
\pgfsetbuttcap%
\pgfsetroundjoin%
\definecolor{currentfill}{rgb}{0.901961,0.196078,0.156863}%
\pgfsetfillcolor{currentfill}%
\pgfsetfillopacity{0.900000}%
\pgfsetlinewidth{0.803000pt}%
\definecolor{currentstroke}{rgb}{0.901961,0.196078,0.156863}%
\pgfsetstrokecolor{currentstroke}%
\pgfsetstrokeopacity{0.900000}%
\pgfsetdash{}{0pt}%
\pgfsys@defobject{currentmarker}{\pgfqpoint{-0.042241in}{-0.042241in}}{\pgfqpoint{0.042241in}{0.042241in}}{%
\pgfpathmoveto{\pgfqpoint{-0.042241in}{-0.042241in}}%
\pgfpathlineto{\pgfqpoint{0.042241in}{-0.042241in}}%
\pgfpathlineto{\pgfqpoint{0.042241in}{0.042241in}}%
\pgfpathlineto{\pgfqpoint{-0.042241in}{0.042241in}}%
\pgfpathlineto{\pgfqpoint{-0.042241in}{-0.042241in}}%
\pgfpathclose%
\pgfusepath{stroke,fill}%
}%
\begin{pgfscope}%
\pgfsys@transformshift{5.875013in}{8.592000in}%
\pgfsys@useobject{currentmarker}{}%
\end{pgfscope}%
\end{pgfscope}%
\begin{pgfscope}%
\pgfpathrectangle{\pgfqpoint{0.977909in}{2.928000in}}{\pgfqpoint{5.664000in}{5.664000in}}%
\pgfusepath{clip}%
\pgfsetbuttcap%
\pgfsetroundjoin%
\pgfsetlinewidth{1.003750pt}%
\definecolor{currentstroke}{rgb}{0.498039,0.560784,0.678431}%
\pgfsetstrokecolor{currentstroke}%
\pgfsetstrokeopacity{0.900000}%
\pgfsetdash{}{0pt}%
\pgfpathmoveto{\pgfqpoint{4.359374in}{6.063292in}}%
\pgfpathlineto{\pgfqpoint{4.443856in}{6.063292in}}%
\pgfpathlineto{\pgfqpoint{4.443856in}{6.147775in}}%
\pgfpathlineto{\pgfqpoint{4.359374in}{6.147775in}}%
\pgfpathlineto{\pgfqpoint{4.359374in}{6.063292in}}%
\pgfpathclose%
\pgfusepath{stroke}%
\end{pgfscope}%
\begin{pgfscope}%
\pgfsetbuttcap%
\pgfsetroundjoin%
\pgfsetlinewidth{1.003750pt}%
\definecolor{currentstroke}{rgb}{0.498039,0.560784,0.678431}%
\pgfsetstrokecolor{currentstroke}%
\pgfsetstrokeopacity{0.900000}%
\pgfsetdash{}{0pt}%
\pgfpathmoveto{\pgfqpoint{0.935668in}{2.885759in}}%
\pgfpathlineto{\pgfqpoint{1.020150in}{2.885759in}}%
\pgfpathlineto{\pgfqpoint{1.020150in}{2.970241in}}%
\pgfpathlineto{\pgfqpoint{0.935668in}{2.970241in}}%
\pgfpathlineto{\pgfqpoint{0.935668in}{2.885759in}}%
\pgfpathclose%
\pgfusepath{stroke}%
\end{pgfscope}%
\begin{pgfscope}%
\pgfpathrectangle{\pgfqpoint{0.977909in}{2.928000in}}{\pgfqpoint{5.664000in}{5.664000in}}%
\pgfusepath{clip}%
\pgfsetbuttcap%
\pgfsetroundjoin%
\pgfsetlinewidth{1.003750pt}%
\definecolor{currentstroke}{rgb}{0.709804,0.121569,0.270588}%
\pgfsetstrokecolor{currentstroke}%
\pgfsetstrokeopacity{0.900000}%
\pgfsetdash{}{0pt}%
\pgfpathmoveto{\pgfqpoint{6.320728in}{8.449256in}}%
\pgfpathlineto{\pgfqpoint{6.405211in}{8.449256in}}%
\pgfpathlineto{\pgfqpoint{6.405211in}{8.533739in}}%
\pgfpathlineto{\pgfqpoint{6.320728in}{8.533739in}}%
\pgfpathlineto{\pgfqpoint{6.320728in}{8.449256in}}%
\pgfpathclose%
\pgfusepath{stroke}%
\end{pgfscope}%
\begin{pgfscope}%
\pgfsetbuttcap%
\pgfsetmiterjoin%
\definecolor{currentfill}{rgb}{1.000000,1.000000,1.000000}%
\pgfsetfillcolor{currentfill}%
\pgfsetfillopacity{0.970000}%
\pgfsetlinewidth{0.702625pt}%
\definecolor{currentstroke}{rgb}{0.796078,0.827451,0.850980}%
\pgfsetstrokecolor{currentstroke}%
\pgfsetstrokeopacity{0.970000}%
\pgfsetdash{}{0pt}%
\pgfpathmoveto{\pgfqpoint{1.176149in}{6.665893in}}%
\pgfpathlineto{\pgfqpoint{3.203024in}{6.665893in}}%
\pgfpathquadraticcurveto{\pgfqpoint{3.280941in}{6.665893in}}{\pgfqpoint{3.280941in}{6.743810in}}%
\pgfpathlineto{\pgfqpoint{3.280941in}{7.402560in}}%
\pgfpathquadraticcurveto{\pgfqpoint{3.280941in}{7.480477in}}{\pgfqpoint{3.203024in}{7.480477in}}%
\pgfpathlineto{\pgfqpoint{1.176149in}{7.480477in}}%
\pgfpathquadraticcurveto{\pgfqpoint{1.098232in}{7.480477in}}{\pgfqpoint{1.098232in}{7.402560in}}%
\pgfpathlineto{\pgfqpoint{1.098232in}{6.743810in}}%
\pgfpathquadraticcurveto{\pgfqpoint{1.098232in}{6.665893in}}{\pgfqpoint{1.176149in}{6.665893in}}%
\pgfpathlineto{\pgfqpoint{1.176149in}{6.665893in}}%
\pgfpathclose%
\pgfusepath{stroke,fill}%
\end{pgfscope}%
\begin{pgfscope}%
\pgfsetbuttcap%
\pgfsetroundjoin%
\definecolor{currentfill}{rgb}{0.090196,0.168627,0.227451}%
\pgfsetfillcolor{currentfill}%
\pgfsetlinewidth{0.000000pt}%
\definecolor{currentstroke}{rgb}{0.090196,0.168627,0.227451}%
\pgfsetstrokecolor{currentstroke}%
\pgfsetdash{}{0pt}%
\pgfpathmoveto{\pgfqpoint{1.251963in}{7.352665in}}%
\pgfpathlineto{\pgfqpoint{1.251963in}{7.339030in}}%
\pgfpathquadraticcurveto{\pgfqpoint{1.244016in}{7.342837in}}{\pgfqpoint{1.236955in}{7.344697in}}%
\pgfpathquadraticcurveto{\pgfqpoint{1.229894in}{7.346578in}}{\pgfqpoint{1.223320in}{7.346578in}}%
\pgfpathquadraticcurveto{\pgfqpoint{1.211920in}{7.346578in}}{\pgfqpoint{1.205722in}{7.342151in}}%
\pgfpathquadraticcurveto{\pgfqpoint{1.199524in}{7.337724in}}{\pgfqpoint{1.199524in}{7.329556in}}%
\pgfpathquadraticcurveto{\pgfqpoint{1.199524in}{7.322694in}}{\pgfqpoint{1.203641in}{7.319197in}}%
\pgfpathquadraticcurveto{\pgfqpoint{1.207758in}{7.315721in}}{\pgfqpoint{1.219247in}{7.313574in}}%
\pgfpathlineto{\pgfqpoint{1.227680in}{7.311848in}}%
\pgfpathquadraticcurveto{\pgfqpoint{1.243308in}{7.308859in}}{\pgfqpoint{1.250745in}{7.301356in}}%
\pgfpathquadraticcurveto{\pgfqpoint{1.258183in}{7.293852in}}{\pgfqpoint{1.258183in}{7.281279in}}%
\pgfpathquadraticcurveto{\pgfqpoint{1.258183in}{7.266249in}}{\pgfqpoint{1.248111in}{7.258501in}}%
\pgfpathquadraticcurveto{\pgfqpoint{1.238062in}{7.250754in}}{\pgfqpoint{1.218627in}{7.250754in}}%
\pgfpathquadraticcurveto{\pgfqpoint{1.211300in}{7.250754in}}{\pgfqpoint{1.203021in}{7.252414in}}%
\pgfpathquadraticcurveto{\pgfqpoint{1.194765in}{7.254074in}}{\pgfqpoint{1.185911in}{7.257328in}}%
\pgfpathlineto{\pgfqpoint{1.185911in}{7.271716in}}%
\pgfpathquadraticcurveto{\pgfqpoint{1.194411in}{7.266957in}}{\pgfqpoint{1.202579in}{7.264522in}}%
\pgfpathquadraticcurveto{\pgfqpoint{1.210747in}{7.262109in}}{\pgfqpoint{1.218627in}{7.262109in}}%
\pgfpathquadraticcurveto{\pgfqpoint{1.230580in}{7.262109in}}{\pgfqpoint{1.237088in}{7.266802in}}%
\pgfpathquadraticcurveto{\pgfqpoint{1.243596in}{7.271517in}}{\pgfqpoint{1.243596in}{7.280238in}}%
\pgfpathquadraticcurveto{\pgfqpoint{1.243596in}{7.287831in}}{\pgfqpoint{1.238925in}{7.292125in}}%
\pgfpathquadraticcurveto{\pgfqpoint{1.234255in}{7.296419in}}{\pgfqpoint{1.223607in}{7.298567in}}%
\pgfpathlineto{\pgfqpoint{1.215085in}{7.300227in}}%
\pgfpathquadraticcurveto{\pgfqpoint{1.199458in}{7.303326in}}{\pgfqpoint{1.192463in}{7.309966in}}%
\pgfpathquadraticcurveto{\pgfqpoint{1.185490in}{7.316607in}}{\pgfqpoint{1.185490in}{7.328449in}}%
\pgfpathquadraticcurveto{\pgfqpoint{1.185490in}{7.342151in}}{\pgfqpoint{1.195141in}{7.350031in}}%
\pgfpathquadraticcurveto{\pgfqpoint{1.204792in}{7.357912in}}{\pgfqpoint{1.221726in}{7.357912in}}%
\pgfpathquadraticcurveto{\pgfqpoint{1.229008in}{7.357912in}}{\pgfqpoint{1.236535in}{7.356583in}}%
\pgfpathquadraticcurveto{\pgfqpoint{1.244083in}{7.355277in}}{\pgfqpoint{1.251963in}{7.352665in}}%
\pgfpathlineto{\pgfqpoint{1.251963in}{7.352665in}}%
\pgfpathclose%
\pgfpathmoveto{\pgfqpoint{1.279433in}{7.330242in}}%
\pgfpathlineto{\pgfqpoint{1.292161in}{7.330242in}}%
\pgfpathlineto{\pgfqpoint{1.292161in}{7.252768in}}%
\pgfpathlineto{\pgfqpoint{1.279433in}{7.252768in}}%
\pgfpathlineto{\pgfqpoint{1.279433in}{7.330242in}}%
\pgfpathlineto{\pgfqpoint{1.279433in}{7.330242in}}%
\pgfpathclose%
\pgfpathmoveto{\pgfqpoint{1.279433in}{7.360413in}}%
\pgfpathlineto{\pgfqpoint{1.292161in}{7.360413in}}%
\pgfpathlineto{\pgfqpoint{1.292161in}{7.344276in}}%
\pgfpathlineto{\pgfqpoint{1.279433in}{7.344276in}}%
\pgfpathlineto{\pgfqpoint{1.279433in}{7.360413in}}%
\pgfpathlineto{\pgfqpoint{1.279433in}{7.360413in}}%
\pgfpathclose%
\pgfpathmoveto{\pgfqpoint{1.369768in}{7.292413in}}%
\pgfpathquadraticcurveto{\pgfqpoint{1.369768in}{7.306248in}}{\pgfqpoint{1.364057in}{7.313840in}}%
\pgfpathquadraticcurveto{\pgfqpoint{1.358368in}{7.321455in}}{\pgfqpoint{1.348053in}{7.321455in}}%
\pgfpathquadraticcurveto{\pgfqpoint{1.337826in}{7.321455in}}{\pgfqpoint{1.332115in}{7.313840in}}%
\pgfpathquadraticcurveto{\pgfqpoint{1.326404in}{7.306248in}}{\pgfqpoint{1.326404in}{7.292413in}}%
\pgfpathquadraticcurveto{\pgfqpoint{1.326404in}{7.278645in}}{\pgfqpoint{1.332115in}{7.271030in}}%
\pgfpathquadraticcurveto{\pgfqpoint{1.337826in}{7.263415in}}{\pgfqpoint{1.348053in}{7.263415in}}%
\pgfpathquadraticcurveto{\pgfqpoint{1.358368in}{7.263415in}}{\pgfqpoint{1.364057in}{7.271030in}}%
\pgfpathquadraticcurveto{\pgfqpoint{1.369768in}{7.278645in}}{\pgfqpoint{1.369768in}{7.292413in}}%
\pgfpathlineto{\pgfqpoint{1.369768in}{7.292413in}}%
\pgfpathclose%
\pgfpathmoveto{\pgfqpoint{1.382495in}{7.262375in}}%
\pgfpathquadraticcurveto{\pgfqpoint{1.382495in}{7.242608in}}{\pgfqpoint{1.373708in}{7.232957in}}%
\pgfpathquadraticcurveto{\pgfqpoint{1.364942in}{7.223306in}}{\pgfqpoint{1.346813in}{7.223306in}}%
\pgfpathquadraticcurveto{\pgfqpoint{1.340106in}{7.223306in}}{\pgfqpoint{1.334152in}{7.224302in}}%
\pgfpathquadraticcurveto{\pgfqpoint{1.328197in}{7.225298in}}{\pgfqpoint{1.322597in}{7.227379in}}%
\pgfpathlineto{\pgfqpoint{1.322597in}{7.239753in}}%
\pgfpathquadraticcurveto{\pgfqpoint{1.328197in}{7.236720in}}{\pgfqpoint{1.333665in}{7.235281in}}%
\pgfpathquadraticcurveto{\pgfqpoint{1.339132in}{7.233820in}}{\pgfqpoint{1.344799in}{7.233820in}}%
\pgfpathquadraticcurveto{\pgfqpoint{1.357327in}{7.233820in}}{\pgfqpoint{1.363548in}{7.240350in}}%
\pgfpathquadraticcurveto{\pgfqpoint{1.369768in}{7.246880in}}{\pgfqpoint{1.369768in}{7.260095in}}%
\pgfpathlineto{\pgfqpoint{1.369768in}{7.266404in}}%
\pgfpathquadraticcurveto{\pgfqpoint{1.365827in}{7.259542in}}{\pgfqpoint{1.359674in}{7.256155in}}%
\pgfpathquadraticcurveto{\pgfqpoint{1.353520in}{7.252768in}}{\pgfqpoint{1.344932in}{7.252768in}}%
\pgfpathquadraticcurveto{\pgfqpoint{1.330699in}{7.252768in}}{\pgfqpoint{1.321977in}{7.263615in}}%
\pgfpathquadraticcurveto{\pgfqpoint{1.313256in}{7.274483in}}{\pgfqpoint{1.313256in}{7.292413in}}%
\pgfpathquadraticcurveto{\pgfqpoint{1.313256in}{7.310387in}}{\pgfqpoint{1.321977in}{7.321233in}}%
\pgfpathquadraticcurveto{\pgfqpoint{1.330699in}{7.332102in}}{\pgfqpoint{1.344932in}{7.332102in}}%
\pgfpathquadraticcurveto{\pgfqpoint{1.353520in}{7.332102in}}{\pgfqpoint{1.359674in}{7.328715in}}%
\pgfpathquadraticcurveto{\pgfqpoint{1.365827in}{7.325328in}}{\pgfqpoint{1.369768in}{7.318488in}}%
\pgfpathlineto{\pgfqpoint{1.369768in}{7.330242in}}%
\pgfpathlineto{\pgfqpoint{1.382495in}{7.330242in}}%
\pgfpathlineto{\pgfqpoint{1.382495in}{7.262375in}}%
\pgfpathlineto{\pgfqpoint{1.382495in}{7.262375in}}%
\pgfpathclose%
\pgfpathmoveto{\pgfqpoint{1.473140in}{7.299540in}}%
\pgfpathlineto{\pgfqpoint{1.473140in}{7.252768in}}%
\pgfpathlineto{\pgfqpoint{1.460412in}{7.252768in}}%
\pgfpathlineto{\pgfqpoint{1.460412in}{7.299120in}}%
\pgfpathquadraticcurveto{\pgfqpoint{1.460412in}{7.310121in}}{\pgfqpoint{1.456118in}{7.315567in}}%
\pgfpathquadraticcurveto{\pgfqpoint{1.451824in}{7.321034in}}{\pgfqpoint{1.443257in}{7.321034in}}%
\pgfpathquadraticcurveto{\pgfqpoint{1.432942in}{7.321034in}}{\pgfqpoint{1.426988in}{7.314460in}}%
\pgfpathquadraticcurveto{\pgfqpoint{1.421033in}{7.307908in}}{\pgfqpoint{1.421033in}{7.296552in}}%
\pgfpathlineto{\pgfqpoint{1.421033in}{7.252768in}}%
\pgfpathlineto{\pgfqpoint{1.408239in}{7.252768in}}%
\pgfpathlineto{\pgfqpoint{1.408239in}{7.330242in}}%
\pgfpathlineto{\pgfqpoint{1.421033in}{7.330242in}}%
\pgfpathlineto{\pgfqpoint{1.421033in}{7.318201in}}%
\pgfpathquadraticcurveto{\pgfqpoint{1.425615in}{7.325195in}}{\pgfqpoint{1.431791in}{7.328649in}}%
\pgfpathquadraticcurveto{\pgfqpoint{1.437989in}{7.332102in}}{\pgfqpoint{1.446090in}{7.332102in}}%
\pgfpathquadraticcurveto{\pgfqpoint{1.459438in}{7.332102in}}{\pgfqpoint{1.466278in}{7.323845in}}%
\pgfpathquadraticcurveto{\pgfqpoint{1.473140in}{7.315589in}}{\pgfqpoint{1.473140in}{7.299540in}}%
\pgfpathlineto{\pgfqpoint{1.473140in}{7.299540in}}%
\pgfpathclose%
\pgfpathmoveto{\pgfqpoint{1.492088in}{7.297238in}}%
\pgfpathlineto{\pgfqpoint{1.529364in}{7.297238in}}%
\pgfpathlineto{\pgfqpoint{1.529364in}{7.285905in}}%
\pgfpathlineto{\pgfqpoint{1.492088in}{7.285905in}}%
\pgfpathlineto{\pgfqpoint{1.492088in}{7.297238in}}%
\pgfpathlineto{\pgfqpoint{1.492088in}{7.297238in}}%
\pgfpathclose%
\pgfpathmoveto{\pgfqpoint{1.594508in}{7.318356in}}%
\pgfpathquadraticcurveto{\pgfqpoint{1.592361in}{7.319595in}}{\pgfqpoint{1.589838in}{7.320171in}}%
\pgfpathquadraticcurveto{\pgfqpoint{1.587314in}{7.320768in}}{\pgfqpoint{1.584282in}{7.320768in}}%
\pgfpathquadraticcurveto{\pgfqpoint{1.573480in}{7.320768in}}{\pgfqpoint{1.567702in}{7.313751in}}%
\pgfpathquadraticcurveto{\pgfqpoint{1.561925in}{7.306734in}}{\pgfqpoint{1.561925in}{7.293586in}}%
\pgfpathlineto{\pgfqpoint{1.561925in}{7.252768in}}%
\pgfpathlineto{\pgfqpoint{1.549131in}{7.252768in}}%
\pgfpathlineto{\pgfqpoint{1.549131in}{7.330242in}}%
\pgfpathlineto{\pgfqpoint{1.561925in}{7.330242in}}%
\pgfpathlineto{\pgfqpoint{1.561925in}{7.318201in}}%
\pgfpathquadraticcurveto{\pgfqpoint{1.565954in}{7.325262in}}{\pgfqpoint{1.572373in}{7.328671in}}%
\pgfpathquadraticcurveto{\pgfqpoint{1.578814in}{7.332102in}}{\pgfqpoint{1.588023in}{7.332102in}}%
\pgfpathquadraticcurveto{\pgfqpoint{1.589329in}{7.332102in}}{\pgfqpoint{1.590923in}{7.331925in}}%
\pgfpathquadraticcurveto{\pgfqpoint{1.592516in}{7.331770in}}{\pgfqpoint{1.594442in}{7.331415in}}%
\pgfpathlineto{\pgfqpoint{1.594508in}{7.318356in}}%
\pgfpathlineto{\pgfqpoint{1.594508in}{7.318356in}}%
\pgfpathclose%
\pgfpathmoveto{\pgfqpoint{1.671008in}{7.294693in}}%
\pgfpathlineto{\pgfqpoint{1.671008in}{7.288473in}}%
\pgfpathlineto{\pgfqpoint{1.612482in}{7.288473in}}%
\pgfpathquadraticcurveto{\pgfqpoint{1.613324in}{7.275324in}}{\pgfqpoint{1.620407in}{7.268440in}}%
\pgfpathquadraticcurveto{\pgfqpoint{1.627490in}{7.261556in}}{\pgfqpoint{1.640152in}{7.261556in}}%
\pgfpathquadraticcurveto{\pgfqpoint{1.647479in}{7.261556in}}{\pgfqpoint{1.654363in}{7.263349in}}%
\pgfpathquadraticcurveto{\pgfqpoint{1.661247in}{7.265142in}}{\pgfqpoint{1.668042in}{7.268750in}}%
\pgfpathlineto{\pgfqpoint{1.668042in}{7.256708in}}%
\pgfpathquadraticcurveto{\pgfqpoint{1.661180in}{7.253809in}}{\pgfqpoint{1.653986in}{7.252281in}}%
\pgfpathquadraticcurveto{\pgfqpoint{1.646792in}{7.250754in}}{\pgfqpoint{1.639399in}{7.250754in}}%
\pgfpathquadraticcurveto{\pgfqpoint{1.620850in}{7.250754in}}{\pgfqpoint{1.610025in}{7.261534in}}%
\pgfpathquadraticcurveto{\pgfqpoint{1.599201in}{7.272336in}}{\pgfqpoint{1.599201in}{7.290753in}}%
\pgfpathquadraticcurveto{\pgfqpoint{1.599201in}{7.309767in}}{\pgfqpoint{1.609472in}{7.320923in}}%
\pgfpathquadraticcurveto{\pgfqpoint{1.619743in}{7.332102in}}{\pgfqpoint{1.637186in}{7.332102in}}%
\pgfpathquadraticcurveto{\pgfqpoint{1.652813in}{7.332102in}}{\pgfqpoint{1.661911in}{7.322030in}}%
\pgfpathquadraticcurveto{\pgfqpoint{1.671008in}{7.311981in}}{\pgfqpoint{1.671008in}{7.294693in}}%
\pgfpathlineto{\pgfqpoint{1.671008in}{7.294693in}}%
\pgfpathclose%
\pgfpathmoveto{\pgfqpoint{1.658281in}{7.298434in}}%
\pgfpathquadraticcurveto{\pgfqpoint{1.658148in}{7.308859in}}{\pgfqpoint{1.652437in}{7.315080in}}%
\pgfpathquadraticcurveto{\pgfqpoint{1.646726in}{7.321322in}}{\pgfqpoint{1.637318in}{7.321322in}}%
\pgfpathquadraticcurveto{\pgfqpoint{1.626671in}{7.321322in}}{\pgfqpoint{1.620274in}{7.315301in}}%
\pgfpathquadraticcurveto{\pgfqpoint{1.613877in}{7.309280in}}{\pgfqpoint{1.612903in}{7.298345in}}%
\pgfpathlineto{\pgfqpoint{1.658281in}{7.298434in}}%
\pgfpathlineto{\pgfqpoint{1.658281in}{7.298434in}}%
\pgfpathclose%
\pgfpathmoveto{\pgfqpoint{1.682785in}{7.330242in}}%
\pgfpathlineto{\pgfqpoint{1.696265in}{7.330242in}}%
\pgfpathlineto{\pgfqpoint{1.720481in}{7.265231in}}%
\pgfpathlineto{\pgfqpoint{1.744697in}{7.330242in}}%
\pgfpathlineto{\pgfqpoint{1.758178in}{7.330242in}}%
\pgfpathlineto{\pgfqpoint{1.729114in}{7.252768in}}%
\pgfpathlineto{\pgfqpoint{1.711826in}{7.252768in}}%
\pgfpathlineto{\pgfqpoint{1.682785in}{7.330242in}}%
\pgfpathlineto{\pgfqpoint{1.682785in}{7.330242in}}%
\pgfpathclose%
\pgfpathmoveto{\pgfqpoint{1.842027in}{7.294693in}}%
\pgfpathlineto{\pgfqpoint{1.842027in}{7.288473in}}%
\pgfpathlineto{\pgfqpoint{1.783501in}{7.288473in}}%
\pgfpathquadraticcurveto{\pgfqpoint{1.784342in}{7.275324in}}{\pgfqpoint{1.791425in}{7.268440in}}%
\pgfpathquadraticcurveto{\pgfqpoint{1.798508in}{7.261556in}}{\pgfqpoint{1.811170in}{7.261556in}}%
\pgfpathquadraticcurveto{\pgfqpoint{1.818497in}{7.261556in}}{\pgfqpoint{1.825381in}{7.263349in}}%
\pgfpathquadraticcurveto{\pgfqpoint{1.832265in}{7.265142in}}{\pgfqpoint{1.839061in}{7.268750in}}%
\pgfpathlineto{\pgfqpoint{1.839061in}{7.256708in}}%
\pgfpathquadraticcurveto{\pgfqpoint{1.832199in}{7.253809in}}{\pgfqpoint{1.825005in}{7.252281in}}%
\pgfpathquadraticcurveto{\pgfqpoint{1.817811in}{7.250754in}}{\pgfqpoint{1.810417in}{7.250754in}}%
\pgfpathquadraticcurveto{\pgfqpoint{1.791868in}{7.250754in}}{\pgfqpoint{1.781044in}{7.261534in}}%
\pgfpathquadraticcurveto{\pgfqpoint{1.770219in}{7.272336in}}{\pgfqpoint{1.770219in}{7.290753in}}%
\pgfpathquadraticcurveto{\pgfqpoint{1.770219in}{7.309767in}}{\pgfqpoint{1.780490in}{7.320923in}}%
\pgfpathquadraticcurveto{\pgfqpoint{1.790761in}{7.332102in}}{\pgfqpoint{1.808204in}{7.332102in}}%
\pgfpathquadraticcurveto{\pgfqpoint{1.823831in}{7.332102in}}{\pgfqpoint{1.832929in}{7.322030in}}%
\pgfpathquadraticcurveto{\pgfqpoint{1.842027in}{7.311981in}}{\pgfqpoint{1.842027in}{7.294693in}}%
\pgfpathlineto{\pgfqpoint{1.842027in}{7.294693in}}%
\pgfpathclose%
\pgfpathmoveto{\pgfqpoint{1.829299in}{7.298434in}}%
\pgfpathquadraticcurveto{\pgfqpoint{1.829166in}{7.308859in}}{\pgfqpoint{1.823455in}{7.315080in}}%
\pgfpathquadraticcurveto{\pgfqpoint{1.817744in}{7.321322in}}{\pgfqpoint{1.808337in}{7.321322in}}%
\pgfpathquadraticcurveto{\pgfqpoint{1.797689in}{7.321322in}}{\pgfqpoint{1.791292in}{7.315301in}}%
\pgfpathquadraticcurveto{\pgfqpoint{1.784895in}{7.309280in}}{\pgfqpoint{1.783921in}{7.298345in}}%
\pgfpathlineto{\pgfqpoint{1.829299in}{7.298434in}}%
\pgfpathlineto{\pgfqpoint{1.829299in}{7.298434in}}%
\pgfpathclose%
\pgfpathmoveto{\pgfqpoint{1.907813in}{7.318356in}}%
\pgfpathquadraticcurveto{\pgfqpoint{1.905666in}{7.319595in}}{\pgfqpoint{1.903143in}{7.320171in}}%
\pgfpathquadraticcurveto{\pgfqpoint{1.900619in}{7.320768in}}{\pgfqpoint{1.897587in}{7.320768in}}%
\pgfpathquadraticcurveto{\pgfqpoint{1.886785in}{7.320768in}}{\pgfqpoint{1.881007in}{7.313751in}}%
\pgfpathquadraticcurveto{\pgfqpoint{1.875230in}{7.306734in}}{\pgfqpoint{1.875230in}{7.293586in}}%
\pgfpathlineto{\pgfqpoint{1.875230in}{7.252768in}}%
\pgfpathlineto{\pgfqpoint{1.862436in}{7.252768in}}%
\pgfpathlineto{\pgfqpoint{1.862436in}{7.330242in}}%
\pgfpathlineto{\pgfqpoint{1.875230in}{7.330242in}}%
\pgfpathlineto{\pgfqpoint{1.875230in}{7.318201in}}%
\pgfpathquadraticcurveto{\pgfqpoint{1.879258in}{7.325262in}}{\pgfqpoint{1.885678in}{7.328671in}}%
\pgfpathquadraticcurveto{\pgfqpoint{1.892119in}{7.332102in}}{\pgfqpoint{1.901327in}{7.332102in}}%
\pgfpathquadraticcurveto{\pgfqpoint{1.902633in}{7.332102in}}{\pgfqpoint{1.904227in}{7.331925in}}%
\pgfpathquadraticcurveto{\pgfqpoint{1.905821in}{7.331770in}}{\pgfqpoint{1.907747in}{7.331415in}}%
\pgfpathlineto{\pgfqpoint{1.907813in}{7.318356in}}%
\pgfpathlineto{\pgfqpoint{1.907813in}{7.318356in}}%
\pgfpathclose%
\pgfpathmoveto{\pgfqpoint{1.970545in}{7.327962in}}%
\pgfpathlineto{\pgfqpoint{1.970545in}{7.315921in}}%
\pgfpathquadraticcurveto{\pgfqpoint{1.965166in}{7.318688in}}{\pgfqpoint{1.959344in}{7.320060in}}%
\pgfpathquadraticcurveto{\pgfqpoint{1.953545in}{7.321455in}}{\pgfqpoint{1.947303in}{7.321455in}}%
\pgfpathquadraticcurveto{\pgfqpoint{1.937829in}{7.321455in}}{\pgfqpoint{1.933092in}{7.318555in}}%
\pgfpathquadraticcurveto{\pgfqpoint{1.928355in}{7.315655in}}{\pgfqpoint{1.928355in}{7.309833in}}%
\pgfpathquadraticcurveto{\pgfqpoint{1.928355in}{7.305406in}}{\pgfqpoint{1.931742in}{7.302883in}}%
\pgfpathquadraticcurveto{\pgfqpoint{1.935128in}{7.300359in}}{\pgfqpoint{1.945377in}{7.298080in}}%
\pgfpathlineto{\pgfqpoint{1.949738in}{7.297106in}}%
\pgfpathquadraticcurveto{\pgfqpoint{1.963285in}{7.294206in}}{\pgfqpoint{1.968995in}{7.288915in}}%
\pgfpathquadraticcurveto{\pgfqpoint{1.974706in}{7.283625in}}{\pgfqpoint{1.974706in}{7.274151in}}%
\pgfpathquadraticcurveto{\pgfqpoint{1.974706in}{7.263349in}}{\pgfqpoint{1.966162in}{7.257040in}}%
\pgfpathquadraticcurveto{\pgfqpoint{1.957618in}{7.250754in}}{\pgfqpoint{1.942676in}{7.250754in}}%
\pgfpathquadraticcurveto{\pgfqpoint{1.936456in}{7.250754in}}{\pgfqpoint{1.929705in}{7.251971in}}%
\pgfpathquadraticcurveto{\pgfqpoint{1.922954in}{7.253189in}}{\pgfqpoint{1.915494in}{7.255602in}}%
\pgfpathlineto{\pgfqpoint{1.915494in}{7.268750in}}%
\pgfpathquadraticcurveto{\pgfqpoint{1.922555in}{7.265076in}}{\pgfqpoint{1.929395in}{7.263238in}}%
\pgfpathquadraticcurveto{\pgfqpoint{1.936235in}{7.261423in}}{\pgfqpoint{1.942964in}{7.261423in}}%
\pgfpathquadraticcurveto{\pgfqpoint{1.951951in}{7.261423in}}{\pgfqpoint{1.956777in}{7.264500in}}%
\pgfpathquadraticcurveto{\pgfqpoint{1.961624in}{7.267577in}}{\pgfqpoint{1.961624in}{7.273177in}}%
\pgfpathquadraticcurveto{\pgfqpoint{1.961624in}{7.278357in}}{\pgfqpoint{1.958127in}{7.281124in}}%
\pgfpathquadraticcurveto{\pgfqpoint{1.954652in}{7.283891in}}{\pgfqpoint{1.942809in}{7.286458in}}%
\pgfpathlineto{\pgfqpoint{1.938382in}{7.287499in}}%
\pgfpathquadraticcurveto{\pgfqpoint{1.926562in}{7.289978in}}{\pgfqpoint{1.921294in}{7.295136in}}%
\pgfpathquadraticcurveto{\pgfqpoint{1.916048in}{7.300293in}}{\pgfqpoint{1.916048in}{7.309280in}}%
\pgfpathquadraticcurveto{\pgfqpoint{1.916048in}{7.320215in}}{\pgfqpoint{1.923795in}{7.326147in}}%
\pgfpathquadraticcurveto{\pgfqpoint{1.931542in}{7.332102in}}{\pgfqpoint{1.945798in}{7.332102in}}%
\pgfpathquadraticcurveto{\pgfqpoint{1.952837in}{7.332102in}}{\pgfqpoint{1.959057in}{7.331061in}}%
\pgfpathquadraticcurveto{\pgfqpoint{1.965299in}{7.330043in}}{\pgfqpoint{1.970545in}{7.327962in}}%
\pgfpathlineto{\pgfqpoint{1.970545in}{7.327962in}}%
\pgfpathclose%
\pgfpathmoveto{\pgfqpoint{2.030178in}{7.291705in}}%
\pgfpathquadraticcurveto{\pgfqpoint{2.014749in}{7.291705in}}{\pgfqpoint{2.008795in}{7.288185in}}%
\pgfpathquadraticcurveto{\pgfqpoint{2.002840in}{7.284665in}}{\pgfqpoint{2.002840in}{7.276143in}}%
\pgfpathquadraticcurveto{\pgfqpoint{2.002840in}{7.269370in}}{\pgfqpoint{2.007312in}{7.265386in}}%
\pgfpathquadraticcurveto{\pgfqpoint{2.011783in}{7.261423in}}{\pgfqpoint{2.019442in}{7.261423in}}%
\pgfpathquadraticcurveto{\pgfqpoint{2.030045in}{7.261423in}}{\pgfqpoint{2.036442in}{7.268927in}}%
\pgfpathquadraticcurveto{\pgfqpoint{2.042839in}{7.276431in}}{\pgfqpoint{2.042839in}{7.288871in}}%
\pgfpathlineto{\pgfqpoint{2.042839in}{7.291705in}}%
\pgfpathlineto{\pgfqpoint{2.030178in}{7.291705in}}%
\pgfpathlineto{\pgfqpoint{2.030178in}{7.291705in}}%
\pgfpathclose%
\pgfpathmoveto{\pgfqpoint{2.055567in}{7.296973in}}%
\pgfpathlineto{\pgfqpoint{2.055567in}{7.252768in}}%
\pgfpathlineto{\pgfqpoint{2.042839in}{7.252768in}}%
\pgfpathlineto{\pgfqpoint{2.042839in}{7.264522in}}%
\pgfpathquadraticcurveto{\pgfqpoint{2.038479in}{7.257483in}}{\pgfqpoint{2.031971in}{7.254119in}}%
\pgfpathquadraticcurveto{\pgfqpoint{2.025463in}{7.250754in}}{\pgfqpoint{2.016055in}{7.250754in}}%
\pgfpathquadraticcurveto{\pgfqpoint{2.004169in}{7.250754in}}{\pgfqpoint{1.997130in}{7.257439in}}%
\pgfpathquadraticcurveto{\pgfqpoint{1.990113in}{7.264124in}}{\pgfqpoint{1.990113in}{7.275324in}}%
\pgfpathquadraticcurveto{\pgfqpoint{1.990113in}{7.288384in}}{\pgfqpoint{1.998856in}{7.295025in}}%
\pgfpathquadraticcurveto{\pgfqpoint{2.007622in}{7.301665in}}{\pgfqpoint{2.024976in}{7.301665in}}%
\pgfpathlineto{\pgfqpoint{2.042839in}{7.301665in}}%
\pgfpathlineto{\pgfqpoint{2.042839in}{7.302927in}}%
\pgfpathquadraticcurveto{\pgfqpoint{2.042839in}{7.311715in}}{\pgfqpoint{2.037062in}{7.316518in}}%
\pgfpathquadraticcurveto{\pgfqpoint{2.031285in}{7.321322in}}{\pgfqpoint{2.020837in}{7.321322in}}%
\pgfpathquadraticcurveto{\pgfqpoint{2.014196in}{7.321322in}}{\pgfqpoint{2.007887in}{7.319728in}}%
\pgfpathquadraticcurveto{\pgfqpoint{2.001601in}{7.318134in}}{\pgfqpoint{1.995801in}{7.314947in}}%
\pgfpathlineto{\pgfqpoint{1.995801in}{7.326723in}}%
\pgfpathquadraticcurveto{\pgfqpoint{2.002774in}{7.329423in}}{\pgfqpoint{2.009348in}{7.330751in}}%
\pgfpathquadraticcurveto{\pgfqpoint{2.015923in}{7.332102in}}{\pgfqpoint{2.022143in}{7.332102in}}%
\pgfpathquadraticcurveto{\pgfqpoint{2.038965in}{7.332102in}}{\pgfqpoint{2.047266in}{7.323380in}}%
\pgfpathquadraticcurveto{\pgfqpoint{2.055567in}{7.314681in}}{\pgfqpoint{2.055567in}{7.296973in}}%
\pgfpathlineto{\pgfqpoint{2.055567in}{7.296973in}}%
\pgfpathclose%
\pgfpathmoveto{\pgfqpoint{2.081775in}{7.360413in}}%
\pgfpathlineto{\pgfqpoint{2.094503in}{7.360413in}}%
\pgfpathlineto{\pgfqpoint{2.094503in}{7.252768in}}%
\pgfpathlineto{\pgfqpoint{2.081775in}{7.252768in}}%
\pgfpathlineto{\pgfqpoint{2.081775in}{7.360413in}}%
\pgfpathlineto{\pgfqpoint{2.081775in}{7.360413in}}%
\pgfpathclose%
\pgfpathmoveto{\pgfqpoint{2.170516in}{7.327962in}}%
\pgfpathlineto{\pgfqpoint{2.170516in}{7.315921in}}%
\pgfpathquadraticcurveto{\pgfqpoint{2.165137in}{7.318688in}}{\pgfqpoint{2.159316in}{7.320060in}}%
\pgfpathquadraticcurveto{\pgfqpoint{2.153516in}{7.321455in}}{\pgfqpoint{2.147274in}{7.321455in}}%
\pgfpathquadraticcurveto{\pgfqpoint{2.137800in}{7.321455in}}{\pgfqpoint{2.133063in}{7.318555in}}%
\pgfpathquadraticcurveto{\pgfqpoint{2.128326in}{7.315655in}}{\pgfqpoint{2.128326in}{7.309833in}}%
\pgfpathquadraticcurveto{\pgfqpoint{2.128326in}{7.305406in}}{\pgfqpoint{2.131713in}{7.302883in}}%
\pgfpathquadraticcurveto{\pgfqpoint{2.135100in}{7.300359in}}{\pgfqpoint{2.145348in}{7.298080in}}%
\pgfpathlineto{\pgfqpoint{2.149709in}{7.297106in}}%
\pgfpathquadraticcurveto{\pgfqpoint{2.163256in}{7.294206in}}{\pgfqpoint{2.168967in}{7.288915in}}%
\pgfpathquadraticcurveto{\pgfqpoint{2.174678in}{7.283625in}}{\pgfqpoint{2.174678in}{7.274151in}}%
\pgfpathquadraticcurveto{\pgfqpoint{2.174678in}{7.263349in}}{\pgfqpoint{2.166133in}{7.257040in}}%
\pgfpathquadraticcurveto{\pgfqpoint{2.157589in}{7.250754in}}{\pgfqpoint{2.142648in}{7.250754in}}%
\pgfpathquadraticcurveto{\pgfqpoint{2.136428in}{7.250754in}}{\pgfqpoint{2.129676in}{7.251971in}}%
\pgfpathquadraticcurveto{\pgfqpoint{2.122925in}{7.253189in}}{\pgfqpoint{2.115465in}{7.255602in}}%
\pgfpathlineto{\pgfqpoint{2.115465in}{7.268750in}}%
\pgfpathquadraticcurveto{\pgfqpoint{2.122527in}{7.265076in}}{\pgfqpoint{2.129367in}{7.263238in}}%
\pgfpathquadraticcurveto{\pgfqpoint{2.136206in}{7.261423in}}{\pgfqpoint{2.142936in}{7.261423in}}%
\pgfpathquadraticcurveto{\pgfqpoint{2.151923in}{7.261423in}}{\pgfqpoint{2.156748in}{7.264500in}}%
\pgfpathquadraticcurveto{\pgfqpoint{2.161596in}{7.267577in}}{\pgfqpoint{2.161596in}{7.273177in}}%
\pgfpathquadraticcurveto{\pgfqpoint{2.161596in}{7.278357in}}{\pgfqpoint{2.158098in}{7.281124in}}%
\pgfpathquadraticcurveto{\pgfqpoint{2.154623in}{7.283891in}}{\pgfqpoint{2.142781in}{7.286458in}}%
\pgfpathlineto{\pgfqpoint{2.138354in}{7.287499in}}%
\pgfpathquadraticcurveto{\pgfqpoint{2.126533in}{7.289978in}}{\pgfqpoint{2.121265in}{7.295136in}}%
\pgfpathquadraticcurveto{\pgfqpoint{2.116019in}{7.300293in}}{\pgfqpoint{2.116019in}{7.309280in}}%
\pgfpathquadraticcurveto{\pgfqpoint{2.116019in}{7.320215in}}{\pgfqpoint{2.123766in}{7.326147in}}%
\pgfpathquadraticcurveto{\pgfqpoint{2.131514in}{7.332102in}}{\pgfqpoint{2.145769in}{7.332102in}}%
\pgfpathquadraticcurveto{\pgfqpoint{2.152808in}{7.332102in}}{\pgfqpoint{2.159028in}{7.331061in}}%
\pgfpathquadraticcurveto{\pgfqpoint{2.165270in}{7.330043in}}{\pgfqpoint{2.170516in}{7.327962in}}%
\pgfpathlineto{\pgfqpoint{2.170516in}{7.327962in}}%
\pgfpathclose%
\pgfpathmoveto{\pgfqpoint{2.198186in}{7.270344in}}%
\pgfpathlineto{\pgfqpoint{2.212773in}{7.270344in}}%
\pgfpathlineto{\pgfqpoint{2.212773in}{7.252768in}}%
\pgfpathlineto{\pgfqpoint{2.198186in}{7.252768in}}%
\pgfpathlineto{\pgfqpoint{2.198186in}{7.270344in}}%
\pgfpathlineto{\pgfqpoint{2.198186in}{7.270344in}}%
\pgfpathclose%
\pgfpathmoveto{\pgfqpoint{2.198186in}{7.326014in}}%
\pgfpathlineto{\pgfqpoint{2.212773in}{7.326014in}}%
\pgfpathlineto{\pgfqpoint{2.212773in}{7.308461in}}%
\pgfpathlineto{\pgfqpoint{2.198186in}{7.308461in}}%
\pgfpathlineto{\pgfqpoint{2.198186in}{7.326014in}}%
\pgfpathlineto{\pgfqpoint{2.198186in}{7.326014in}}%
\pgfpathclose%
\pgfpathmoveto{\pgfqpoint{2.291907in}{7.264522in}}%
\pgfpathlineto{\pgfqpoint{2.314729in}{7.264522in}}%
\pgfpathlineto{\pgfqpoint{2.314729in}{7.343324in}}%
\pgfpathlineto{\pgfqpoint{2.289893in}{7.338344in}}%
\pgfpathlineto{\pgfqpoint{2.289893in}{7.351072in}}%
\pgfpathlineto{\pgfqpoint{2.314596in}{7.356052in}}%
\pgfpathlineto{\pgfqpoint{2.328563in}{7.356052in}}%
\pgfpathlineto{\pgfqpoint{2.328563in}{7.264522in}}%
\pgfpathlineto{\pgfqpoint{2.351385in}{7.264522in}}%
\pgfpathlineto{\pgfqpoint{2.351385in}{7.252768in}}%
\pgfpathlineto{\pgfqpoint{2.291907in}{7.252768in}}%
\pgfpathlineto{\pgfqpoint{2.291907in}{7.264522in}}%
\pgfpathlineto{\pgfqpoint{2.291907in}{7.264522in}}%
\pgfpathclose%
\pgfpathmoveto{\pgfqpoint{2.411239in}{7.309966in}}%
\pgfpathquadraticcurveto{\pgfqpoint{2.401831in}{7.309966in}}{\pgfqpoint{2.396320in}{7.303525in}}%
\pgfpathquadraticcurveto{\pgfqpoint{2.390830in}{7.297106in}}{\pgfqpoint{2.390830in}{7.285905in}}%
\pgfpathquadraticcurveto{\pgfqpoint{2.390830in}{7.274771in}}{\pgfqpoint{2.396320in}{7.268285in}}%
\pgfpathquadraticcurveto{\pgfqpoint{2.401831in}{7.261822in}}{\pgfqpoint{2.411239in}{7.261822in}}%
\pgfpathquadraticcurveto{\pgfqpoint{2.420646in}{7.261822in}}{\pgfqpoint{2.426136in}{7.268285in}}%
\pgfpathquadraticcurveto{\pgfqpoint{2.431626in}{7.274771in}}{\pgfqpoint{2.431626in}{7.285905in}}%
\pgfpathquadraticcurveto{\pgfqpoint{2.431626in}{7.297106in}}{\pgfqpoint{2.426136in}{7.303525in}}%
\pgfpathquadraticcurveto{\pgfqpoint{2.420646in}{7.309966in}}{\pgfqpoint{2.411239in}{7.309966in}}%
\pgfpathlineto{\pgfqpoint{2.411239in}{7.309966in}}%
\pgfpathclose%
\pgfpathmoveto{\pgfqpoint{2.438975in}{7.353772in}}%
\pgfpathlineto{\pgfqpoint{2.438975in}{7.341044in}}%
\pgfpathquadraticcurveto{\pgfqpoint{2.433706in}{7.343524in}}{\pgfqpoint{2.428350in}{7.344830in}}%
\pgfpathquadraticcurveto{\pgfqpoint{2.422993in}{7.346158in}}{\pgfqpoint{2.417725in}{7.346158in}}%
\pgfpathquadraticcurveto{\pgfqpoint{2.403890in}{7.346158in}}{\pgfqpoint{2.396585in}{7.336817in}}%
\pgfpathquadraticcurveto{\pgfqpoint{2.389303in}{7.327475in}}{\pgfqpoint{2.388262in}{7.308594in}}%
\pgfpathquadraticcurveto{\pgfqpoint{2.392335in}{7.314615in}}{\pgfqpoint{2.398489in}{7.317824in}}%
\pgfpathquadraticcurveto{\pgfqpoint{2.404665in}{7.321034in}}{\pgfqpoint{2.412058in}{7.321034in}}%
\pgfpathquadraticcurveto{\pgfqpoint{2.427619in}{7.321034in}}{\pgfqpoint{2.436650in}{7.311582in}}%
\pgfpathquadraticcurveto{\pgfqpoint{2.445682in}{7.302152in}}{\pgfqpoint{2.445682in}{7.285905in}}%
\pgfpathquadraticcurveto{\pgfqpoint{2.445682in}{7.269990in}}{\pgfqpoint{2.436274in}{7.260361in}}%
\pgfpathquadraticcurveto{\pgfqpoint{2.426867in}{7.250754in}}{\pgfqpoint{2.411239in}{7.250754in}}%
\pgfpathquadraticcurveto{\pgfqpoint{2.393309in}{7.250754in}}{\pgfqpoint{2.383835in}{7.264478in}}%
\pgfpathquadraticcurveto{\pgfqpoint{2.374361in}{7.278224in}}{\pgfqpoint{2.374361in}{7.304300in}}%
\pgfpathquadraticcurveto{\pgfqpoint{2.374361in}{7.328781in}}{\pgfqpoint{2.385982in}{7.343346in}}%
\pgfpathquadraticcurveto{\pgfqpoint{2.397604in}{7.357912in}}{\pgfqpoint{2.417171in}{7.357912in}}%
\pgfpathquadraticcurveto{\pgfqpoint{2.422439in}{7.357912in}}{\pgfqpoint{2.427796in}{7.356871in}}%
\pgfpathquadraticcurveto{\pgfqpoint{2.433153in}{7.355831in}}{\pgfqpoint{2.438975in}{7.353772in}}%
\pgfpathlineto{\pgfqpoint{2.438975in}{7.353772in}}%
\pgfpathclose%
\pgfpathmoveto{\pgfqpoint{2.469743in}{7.270344in}}%
\pgfpathlineto{\pgfqpoint{2.484352in}{7.270344in}}%
\pgfpathlineto{\pgfqpoint{2.484352in}{7.252768in}}%
\pgfpathlineto{\pgfqpoint{2.469743in}{7.252768in}}%
\pgfpathlineto{\pgfqpoint{2.469743in}{7.270344in}}%
\pgfpathlineto{\pgfqpoint{2.469743in}{7.270344in}}%
\pgfpathclose%
\pgfpathmoveto{\pgfqpoint{2.514921in}{7.356052in}}%
\pgfpathlineto{\pgfqpoint{2.569773in}{7.356052in}}%
\pgfpathlineto{\pgfqpoint{2.569773in}{7.344276in}}%
\pgfpathlineto{\pgfqpoint{2.527715in}{7.344276in}}%
\pgfpathlineto{\pgfqpoint{2.527715in}{7.318975in}}%
\pgfpathquadraticcurveto{\pgfqpoint{2.530748in}{7.320016in}}{\pgfqpoint{2.533781in}{7.320525in}}%
\pgfpathquadraticcurveto{\pgfqpoint{2.536835in}{7.321034in}}{\pgfqpoint{2.539890in}{7.321034in}}%
\pgfpathquadraticcurveto{\pgfqpoint{2.557178in}{7.321034in}}{\pgfqpoint{2.567271in}{7.311560in}}%
\pgfpathquadraticcurveto{\pgfqpoint{2.577387in}{7.302086in}}{\pgfqpoint{2.577387in}{7.285905in}}%
\pgfpathquadraticcurveto{\pgfqpoint{2.577387in}{7.269237in}}{\pgfqpoint{2.567006in}{7.259984in}}%
\pgfpathquadraticcurveto{\pgfqpoint{2.556624in}{7.250754in}}{\pgfqpoint{2.537743in}{7.250754in}}%
\pgfpathquadraticcurveto{\pgfqpoint{2.531235in}{7.250754in}}{\pgfqpoint{2.524484in}{7.251861in}}%
\pgfpathquadraticcurveto{\pgfqpoint{2.517755in}{7.252968in}}{\pgfqpoint{2.510561in}{7.255181in}}%
\pgfpathlineto{\pgfqpoint{2.510561in}{7.269237in}}%
\pgfpathquadraticcurveto{\pgfqpoint{2.516781in}{7.265850in}}{\pgfqpoint{2.523421in}{7.264190in}}%
\pgfpathquadraticcurveto{\pgfqpoint{2.530062in}{7.262530in}}{\pgfqpoint{2.537455in}{7.262530in}}%
\pgfpathquadraticcurveto{\pgfqpoint{2.549430in}{7.262530in}}{\pgfqpoint{2.556403in}{7.268817in}}%
\pgfpathquadraticcurveto{\pgfqpoint{2.563398in}{7.275103in}}{\pgfqpoint{2.563398in}{7.285905in}}%
\pgfpathquadraticcurveto{\pgfqpoint{2.563398in}{7.296685in}}{\pgfqpoint{2.556403in}{7.302971in}}%
\pgfpathquadraticcurveto{\pgfqpoint{2.549430in}{7.309280in}}{\pgfqpoint{2.537455in}{7.309280in}}%
\pgfpathquadraticcurveto{\pgfqpoint{2.531855in}{7.309280in}}{\pgfqpoint{2.526277in}{7.308040in}}%
\pgfpathquadraticcurveto{\pgfqpoint{2.520721in}{7.306801in}}{\pgfqpoint{2.514921in}{7.304167in}}%
\pgfpathlineto{\pgfqpoint{2.514921in}{7.356052in}}%
\pgfpathlineto{\pgfqpoint{2.514921in}{7.356052in}}%
\pgfpathclose%
\pgfpathmoveto{\pgfqpoint{2.692757in}{7.298212in}}%
\pgfpathquadraticcurveto{\pgfqpoint{2.686736in}{7.298212in}}{\pgfqpoint{2.683305in}{7.293099in}}%
\pgfpathquadraticcurveto{\pgfqpoint{2.679896in}{7.287986in}}{\pgfqpoint{2.679896in}{7.278844in}}%
\pgfpathquadraticcurveto{\pgfqpoint{2.679896in}{7.269857in}}{\pgfqpoint{2.683305in}{7.264699in}}%
\pgfpathquadraticcurveto{\pgfqpoint{2.686736in}{7.259542in}}{\pgfqpoint{2.692757in}{7.259542in}}%
\pgfpathquadraticcurveto{\pgfqpoint{2.698645in}{7.259542in}}{\pgfqpoint{2.702054in}{7.264699in}}%
\pgfpathquadraticcurveto{\pgfqpoint{2.705485in}{7.269857in}}{\pgfqpoint{2.705485in}{7.278844in}}%
\pgfpathquadraticcurveto{\pgfqpoint{2.705485in}{7.287919in}}{\pgfqpoint{2.702054in}{7.293055in}}%
\pgfpathquadraticcurveto{\pgfqpoint{2.698645in}{7.298212in}}{\pgfqpoint{2.692757in}{7.298212in}}%
\pgfpathlineto{\pgfqpoint{2.692757in}{7.298212in}}%
\pgfpathclose%
\pgfpathmoveto{\pgfqpoint{2.692757in}{7.307000in}}%
\pgfpathquadraticcurveto{\pgfqpoint{2.703692in}{7.307000in}}{\pgfqpoint{2.710111in}{7.299386in}}%
\pgfpathquadraticcurveto{\pgfqpoint{2.716553in}{7.291793in}}{\pgfqpoint{2.716553in}{7.278844in}}%
\pgfpathquadraticcurveto{\pgfqpoint{2.716553in}{7.265917in}}{\pgfqpoint{2.710089in}{7.258324in}}%
\pgfpathquadraticcurveto{\pgfqpoint{2.703626in}{7.250754in}}{\pgfqpoint{2.692757in}{7.250754in}}%
\pgfpathquadraticcurveto{\pgfqpoint{2.681689in}{7.250754in}}{\pgfqpoint{2.675248in}{7.258324in}}%
\pgfpathquadraticcurveto{\pgfqpoint{2.668829in}{7.265917in}}{\pgfqpoint{2.668829in}{7.278844in}}%
\pgfpathquadraticcurveto{\pgfqpoint{2.668829in}{7.291859in}}{\pgfqpoint{2.675292in}{7.299430in}}%
\pgfpathquadraticcurveto{\pgfqpoint{2.681756in}{7.307000in}}{\pgfqpoint{2.692757in}{7.307000in}}%
\pgfpathlineto{\pgfqpoint{2.692757in}{7.307000in}}%
\pgfpathclose%
\pgfpathmoveto{\pgfqpoint{2.621370in}{7.349124in}}%
\pgfpathquadraticcurveto{\pgfqpoint{2.615416in}{7.349124in}}{\pgfqpoint{2.611985in}{7.343966in}}%
\pgfpathquadraticcurveto{\pgfqpoint{2.608576in}{7.338831in}}{\pgfqpoint{2.608576in}{7.329822in}}%
\pgfpathquadraticcurveto{\pgfqpoint{2.608576in}{7.320702in}}{\pgfqpoint{2.611963in}{7.315567in}}%
\pgfpathquadraticcurveto{\pgfqpoint{2.615350in}{7.310453in}}{\pgfqpoint{2.621370in}{7.310453in}}%
\pgfpathquadraticcurveto{\pgfqpoint{2.627391in}{7.310453in}}{\pgfqpoint{2.630800in}{7.315567in}}%
\pgfpathquadraticcurveto{\pgfqpoint{2.634231in}{7.320702in}}{\pgfqpoint{2.634231in}{7.329822in}}%
\pgfpathquadraticcurveto{\pgfqpoint{2.634231in}{7.338742in}}{\pgfqpoint{2.630778in}{7.343922in}}%
\pgfpathquadraticcurveto{\pgfqpoint{2.627325in}{7.349124in}}{\pgfqpoint{2.621370in}{7.349124in}}%
\pgfpathlineto{\pgfqpoint{2.621370in}{7.349124in}}%
\pgfpathclose%
\pgfpathmoveto{\pgfqpoint{2.683837in}{7.357912in}}%
\pgfpathlineto{\pgfqpoint{2.694904in}{7.357912in}}%
\pgfpathlineto{\pgfqpoint{2.630291in}{7.250754in}}%
\pgfpathlineto{\pgfqpoint{2.619223in}{7.250754in}}%
\pgfpathlineto{\pgfqpoint{2.683837in}{7.357912in}}%
\pgfpathlineto{\pgfqpoint{2.683837in}{7.357912in}}%
\pgfpathclose%
\pgfpathmoveto{\pgfqpoint{2.621370in}{7.357912in}}%
\pgfpathquadraticcurveto{\pgfqpoint{2.632305in}{7.357912in}}{\pgfqpoint{2.638791in}{7.350341in}}%
\pgfpathquadraticcurveto{\pgfqpoint{2.645299in}{7.342771in}}{\pgfqpoint{2.645299in}{7.329822in}}%
\pgfpathquadraticcurveto{\pgfqpoint{2.645299in}{7.316762in}}{\pgfqpoint{2.638835in}{7.309214in}}%
\pgfpathquadraticcurveto{\pgfqpoint{2.632372in}{7.301665in}}{\pgfqpoint{2.621370in}{7.301665in}}%
\pgfpathquadraticcurveto{\pgfqpoint{2.610369in}{7.301665in}}{\pgfqpoint{2.603972in}{7.309236in}}%
\pgfpathquadraticcurveto{\pgfqpoint{2.597575in}{7.316828in}}{\pgfqpoint{2.597575in}{7.329822in}}%
\pgfpathquadraticcurveto{\pgfqpoint{2.597575in}{7.342705in}}{\pgfqpoint{2.603994in}{7.350297in}}%
\pgfpathquadraticcurveto{\pgfqpoint{2.610436in}{7.357912in}}{\pgfqpoint{2.621370in}{7.357912in}}%
\pgfpathlineto{\pgfqpoint{2.621370in}{7.357912in}}%
\pgfpathclose%
\pgfusepath{fill}%
\end{pgfscope}%
\begin{pgfscope}%
\pgfsetbuttcap%
\pgfsetroundjoin%
\definecolor{currentfill}{rgb}{0.090196,0.168627,0.227451}%
\pgfsetfillcolor{currentfill}%
\pgfsetlinewidth{0.000000pt}%
\definecolor{currentstroke}{rgb}{0.090196,0.168627,0.227451}%
\pgfsetstrokecolor{currentstroke}%
\pgfsetdash{}{0pt}%
\pgfpathmoveto{\pgfqpoint{1.251963in}{7.133082in}}%
\pgfpathlineto{\pgfqpoint{1.251963in}{7.119447in}}%
\pgfpathquadraticcurveto{\pgfqpoint{1.244016in}{7.123254in}}{\pgfqpoint{1.236955in}{7.125113in}}%
\pgfpathquadraticcurveto{\pgfqpoint{1.229894in}{7.126995in}}{\pgfqpoint{1.223320in}{7.126995in}}%
\pgfpathquadraticcurveto{\pgfqpoint{1.211920in}{7.126995in}}{\pgfqpoint{1.205722in}{7.122568in}}%
\pgfpathquadraticcurveto{\pgfqpoint{1.199524in}{7.118141in}}{\pgfqpoint{1.199524in}{7.109973in}}%
\pgfpathquadraticcurveto{\pgfqpoint{1.199524in}{7.103111in}}{\pgfqpoint{1.203641in}{7.099613in}}%
\pgfpathquadraticcurveto{\pgfqpoint{1.207758in}{7.096138in}}{\pgfqpoint{1.219247in}{7.093991in}}%
\pgfpathlineto{\pgfqpoint{1.227680in}{7.092264in}}%
\pgfpathquadraticcurveto{\pgfqpoint{1.243308in}{7.089276in}}{\pgfqpoint{1.250745in}{7.081772in}}%
\pgfpathquadraticcurveto{\pgfqpoint{1.258183in}{7.074268in}}{\pgfqpoint{1.258183in}{7.061695in}}%
\pgfpathquadraticcurveto{\pgfqpoint{1.258183in}{7.046665in}}{\pgfqpoint{1.248111in}{7.038918in}}%
\pgfpathquadraticcurveto{\pgfqpoint{1.238062in}{7.031171in}}{\pgfqpoint{1.218627in}{7.031171in}}%
\pgfpathquadraticcurveto{\pgfqpoint{1.211300in}{7.031171in}}{\pgfqpoint{1.203021in}{7.032831in}}%
\pgfpathquadraticcurveto{\pgfqpoint{1.194765in}{7.034491in}}{\pgfqpoint{1.185911in}{7.037745in}}%
\pgfpathlineto{\pgfqpoint{1.185911in}{7.052133in}}%
\pgfpathquadraticcurveto{\pgfqpoint{1.194411in}{7.047374in}}{\pgfqpoint{1.202579in}{7.044939in}}%
\pgfpathquadraticcurveto{\pgfqpoint{1.210747in}{7.042526in}}{\pgfqpoint{1.218627in}{7.042526in}}%
\pgfpathquadraticcurveto{\pgfqpoint{1.230580in}{7.042526in}}{\pgfqpoint{1.237088in}{7.047219in}}%
\pgfpathquadraticcurveto{\pgfqpoint{1.243596in}{7.051934in}}{\pgfqpoint{1.243596in}{7.060655in}}%
\pgfpathquadraticcurveto{\pgfqpoint{1.243596in}{7.068248in}}{\pgfqpoint{1.238925in}{7.072542in}}%
\pgfpathquadraticcurveto{\pgfqpoint{1.234255in}{7.076836in}}{\pgfqpoint{1.223607in}{7.078983in}}%
\pgfpathlineto{\pgfqpoint{1.215085in}{7.080643in}}%
\pgfpathquadraticcurveto{\pgfqpoint{1.199458in}{7.083742in}}{\pgfqpoint{1.192463in}{7.090383in}}%
\pgfpathquadraticcurveto{\pgfqpoint{1.185490in}{7.097024in}}{\pgfqpoint{1.185490in}{7.108866in}}%
\pgfpathquadraticcurveto{\pgfqpoint{1.185490in}{7.122568in}}{\pgfqpoint{1.195141in}{7.130448in}}%
\pgfpathquadraticcurveto{\pgfqpoint{1.204792in}{7.138328in}}{\pgfqpoint{1.221726in}{7.138328in}}%
\pgfpathquadraticcurveto{\pgfqpoint{1.229008in}{7.138328in}}{\pgfqpoint{1.236535in}{7.137000in}}%
\pgfpathquadraticcurveto{\pgfqpoint{1.244083in}{7.135694in}}{\pgfqpoint{1.251963in}{7.133082in}}%
\pgfpathlineto{\pgfqpoint{1.251963in}{7.133082in}}%
\pgfpathclose%
\pgfpathmoveto{\pgfqpoint{1.279433in}{7.110659in}}%
\pgfpathlineto{\pgfqpoint{1.292161in}{7.110659in}}%
\pgfpathlineto{\pgfqpoint{1.292161in}{7.033185in}}%
\pgfpathlineto{\pgfqpoint{1.279433in}{7.033185in}}%
\pgfpathlineto{\pgfqpoint{1.279433in}{7.110659in}}%
\pgfpathlineto{\pgfqpoint{1.279433in}{7.110659in}}%
\pgfpathclose%
\pgfpathmoveto{\pgfqpoint{1.279433in}{7.140830in}}%
\pgfpathlineto{\pgfqpoint{1.292161in}{7.140830in}}%
\pgfpathlineto{\pgfqpoint{1.292161in}{7.124693in}}%
\pgfpathlineto{\pgfqpoint{1.279433in}{7.124693in}}%
\pgfpathlineto{\pgfqpoint{1.279433in}{7.140830in}}%
\pgfpathlineto{\pgfqpoint{1.279433in}{7.140830in}}%
\pgfpathclose%
\pgfpathmoveto{\pgfqpoint{1.369768in}{7.072830in}}%
\pgfpathquadraticcurveto{\pgfqpoint{1.369768in}{7.086664in}}{\pgfqpoint{1.364057in}{7.094257in}}%
\pgfpathquadraticcurveto{\pgfqpoint{1.358368in}{7.101871in}}{\pgfqpoint{1.348053in}{7.101871in}}%
\pgfpathquadraticcurveto{\pgfqpoint{1.337826in}{7.101871in}}{\pgfqpoint{1.332115in}{7.094257in}}%
\pgfpathquadraticcurveto{\pgfqpoint{1.326404in}{7.086664in}}{\pgfqpoint{1.326404in}{7.072830in}}%
\pgfpathquadraticcurveto{\pgfqpoint{1.326404in}{7.059061in}}{\pgfqpoint{1.332115in}{7.051447in}}%
\pgfpathquadraticcurveto{\pgfqpoint{1.337826in}{7.043832in}}{\pgfqpoint{1.348053in}{7.043832in}}%
\pgfpathquadraticcurveto{\pgfqpoint{1.358368in}{7.043832in}}{\pgfqpoint{1.364057in}{7.051447in}}%
\pgfpathquadraticcurveto{\pgfqpoint{1.369768in}{7.059061in}}{\pgfqpoint{1.369768in}{7.072830in}}%
\pgfpathlineto{\pgfqpoint{1.369768in}{7.072830in}}%
\pgfpathclose%
\pgfpathmoveto{\pgfqpoint{1.382495in}{7.042792in}}%
\pgfpathquadraticcurveto{\pgfqpoint{1.382495in}{7.023025in}}{\pgfqpoint{1.373708in}{7.013374in}}%
\pgfpathquadraticcurveto{\pgfqpoint{1.364942in}{7.003723in}}{\pgfqpoint{1.346813in}{7.003723in}}%
\pgfpathquadraticcurveto{\pgfqpoint{1.340106in}{7.003723in}}{\pgfqpoint{1.334152in}{7.004719in}}%
\pgfpathquadraticcurveto{\pgfqpoint{1.328197in}{7.005715in}}{\pgfqpoint{1.322597in}{7.007796in}}%
\pgfpathlineto{\pgfqpoint{1.322597in}{7.020169in}}%
\pgfpathquadraticcurveto{\pgfqpoint{1.328197in}{7.017137in}}{\pgfqpoint{1.333665in}{7.015698in}}%
\pgfpathquadraticcurveto{\pgfqpoint{1.339132in}{7.014237in}}{\pgfqpoint{1.344799in}{7.014237in}}%
\pgfpathquadraticcurveto{\pgfqpoint{1.357327in}{7.014237in}}{\pgfqpoint{1.363548in}{7.020767in}}%
\pgfpathquadraticcurveto{\pgfqpoint{1.369768in}{7.027297in}}{\pgfqpoint{1.369768in}{7.040512in}}%
\pgfpathlineto{\pgfqpoint{1.369768in}{7.046820in}}%
\pgfpathquadraticcurveto{\pgfqpoint{1.365827in}{7.039958in}}{\pgfqpoint{1.359674in}{7.036572in}}%
\pgfpathquadraticcurveto{\pgfqpoint{1.353520in}{7.033185in}}{\pgfqpoint{1.344932in}{7.033185in}}%
\pgfpathquadraticcurveto{\pgfqpoint{1.330699in}{7.033185in}}{\pgfqpoint{1.321977in}{7.044031in}}%
\pgfpathquadraticcurveto{\pgfqpoint{1.313256in}{7.054900in}}{\pgfqpoint{1.313256in}{7.072830in}}%
\pgfpathquadraticcurveto{\pgfqpoint{1.313256in}{7.090803in}}{\pgfqpoint{1.321977in}{7.101650in}}%
\pgfpathquadraticcurveto{\pgfqpoint{1.330699in}{7.112518in}}{\pgfqpoint{1.344932in}{7.112518in}}%
\pgfpathquadraticcurveto{\pgfqpoint{1.353520in}{7.112518in}}{\pgfqpoint{1.359674in}{7.109132in}}%
\pgfpathquadraticcurveto{\pgfqpoint{1.365827in}{7.105745in}}{\pgfqpoint{1.369768in}{7.098905in}}%
\pgfpathlineto{\pgfqpoint{1.369768in}{7.110659in}}%
\pgfpathlineto{\pgfqpoint{1.382495in}{7.110659in}}%
\pgfpathlineto{\pgfqpoint{1.382495in}{7.042792in}}%
\pgfpathlineto{\pgfqpoint{1.382495in}{7.042792in}}%
\pgfpathclose%
\pgfpathmoveto{\pgfqpoint{1.473140in}{7.079957in}}%
\pgfpathlineto{\pgfqpoint{1.473140in}{7.033185in}}%
\pgfpathlineto{\pgfqpoint{1.460412in}{7.033185in}}%
\pgfpathlineto{\pgfqpoint{1.460412in}{7.079537in}}%
\pgfpathquadraticcurveto{\pgfqpoint{1.460412in}{7.090538in}}{\pgfqpoint{1.456118in}{7.095983in}}%
\pgfpathquadraticcurveto{\pgfqpoint{1.451824in}{7.101451in}}{\pgfqpoint{1.443257in}{7.101451in}}%
\pgfpathquadraticcurveto{\pgfqpoint{1.432942in}{7.101451in}}{\pgfqpoint{1.426988in}{7.094876in}}%
\pgfpathquadraticcurveto{\pgfqpoint{1.421033in}{7.088324in}}{\pgfqpoint{1.421033in}{7.076969in}}%
\pgfpathlineto{\pgfqpoint{1.421033in}{7.033185in}}%
\pgfpathlineto{\pgfqpoint{1.408239in}{7.033185in}}%
\pgfpathlineto{\pgfqpoint{1.408239in}{7.110659in}}%
\pgfpathlineto{\pgfqpoint{1.421033in}{7.110659in}}%
\pgfpathlineto{\pgfqpoint{1.421033in}{7.098617in}}%
\pgfpathquadraticcurveto{\pgfqpoint{1.425615in}{7.105612in}}{\pgfqpoint{1.431791in}{7.109065in}}%
\pgfpathquadraticcurveto{\pgfqpoint{1.437989in}{7.112518in}}{\pgfqpoint{1.446090in}{7.112518in}}%
\pgfpathquadraticcurveto{\pgfqpoint{1.459438in}{7.112518in}}{\pgfqpoint{1.466278in}{7.104262in}}%
\pgfpathquadraticcurveto{\pgfqpoint{1.473140in}{7.096005in}}{\pgfqpoint{1.473140in}{7.079957in}}%
\pgfpathlineto{\pgfqpoint{1.473140in}{7.079957in}}%
\pgfpathclose%
\pgfpathmoveto{\pgfqpoint{1.498507in}{7.110659in}}%
\pgfpathlineto{\pgfqpoint{1.511235in}{7.110659in}}%
\pgfpathlineto{\pgfqpoint{1.511235in}{7.033185in}}%
\pgfpathlineto{\pgfqpoint{1.498507in}{7.033185in}}%
\pgfpathlineto{\pgfqpoint{1.498507in}{7.110659in}}%
\pgfpathlineto{\pgfqpoint{1.498507in}{7.110659in}}%
\pgfpathclose%
\pgfpathmoveto{\pgfqpoint{1.498507in}{7.140830in}}%
\pgfpathlineto{\pgfqpoint{1.511235in}{7.140830in}}%
\pgfpathlineto{\pgfqpoint{1.511235in}{7.124693in}}%
\pgfpathlineto{\pgfqpoint{1.498507in}{7.124693in}}%
\pgfpathlineto{\pgfqpoint{1.498507in}{7.140830in}}%
\pgfpathlineto{\pgfqpoint{1.498507in}{7.140830in}}%
\pgfpathclose%
\pgfpathmoveto{\pgfqpoint{1.600463in}{7.110659in}}%
\pgfpathlineto{\pgfqpoint{1.600463in}{7.033185in}}%
\pgfpathlineto{\pgfqpoint{1.587669in}{7.033185in}}%
\pgfpathlineto{\pgfqpoint{1.587669in}{7.100764in}}%
\pgfpathlineto{\pgfqpoint{1.552739in}{7.100764in}}%
\pgfpathlineto{\pgfqpoint{1.552739in}{7.033185in}}%
\pgfpathlineto{\pgfqpoint{1.539945in}{7.033185in}}%
\pgfpathlineto{\pgfqpoint{1.539945in}{7.100764in}}%
\pgfpathlineto{\pgfqpoint{1.527770in}{7.100764in}}%
\pgfpathlineto{\pgfqpoint{1.527770in}{7.110659in}}%
\pgfpathlineto{\pgfqpoint{1.539945in}{7.110659in}}%
\pgfpathlineto{\pgfqpoint{1.539945in}{7.116060in}}%
\pgfpathquadraticcurveto{\pgfqpoint{1.539945in}{7.128721in}}{\pgfqpoint{1.545921in}{7.134764in}}%
\pgfpathquadraticcurveto{\pgfqpoint{1.551920in}{7.140830in}}{\pgfqpoint{1.564294in}{7.140830in}}%
\pgfpathlineto{\pgfqpoint{1.577088in}{7.140830in}}%
\pgfpathlineto{\pgfqpoint{1.577088in}{7.130227in}}%
\pgfpathlineto{\pgfqpoint{1.564913in}{7.130227in}}%
\pgfpathquadraticcurveto{\pgfqpoint{1.558074in}{7.130227in}}{\pgfqpoint{1.555395in}{7.127460in}}%
\pgfpathquadraticcurveto{\pgfqpoint{1.552739in}{7.124693in}}{\pgfqpoint{1.552739in}{7.117499in}}%
\pgfpathlineto{\pgfqpoint{1.552739in}{7.110659in}}%
\pgfpathlineto{\pgfqpoint{1.600463in}{7.110659in}}%
\pgfpathlineto{\pgfqpoint{1.600463in}{7.110659in}}%
\pgfpathclose%
\pgfpathmoveto{\pgfqpoint{1.587669in}{7.140675in}}%
\pgfpathlineto{\pgfqpoint{1.600463in}{7.140675in}}%
\pgfpathlineto{\pgfqpoint{1.600463in}{7.124560in}}%
\pgfpathlineto{\pgfqpoint{1.587669in}{7.124560in}}%
\pgfpathlineto{\pgfqpoint{1.587669in}{7.140675in}}%
\pgfpathlineto{\pgfqpoint{1.587669in}{7.140675in}}%
\pgfpathclose%
\pgfpathmoveto{\pgfqpoint{1.682851in}{7.107693in}}%
\pgfpathlineto{\pgfqpoint{1.682851in}{7.095784in}}%
\pgfpathquadraticcurveto{\pgfqpoint{1.677450in}{7.098772in}}{\pgfqpoint{1.672027in}{7.100255in}}%
\pgfpathquadraticcurveto{\pgfqpoint{1.666604in}{7.101738in}}{\pgfqpoint{1.661070in}{7.101738in}}%
\pgfpathquadraticcurveto{\pgfqpoint{1.648674in}{7.101738in}}{\pgfqpoint{1.641812in}{7.093880in}}%
\pgfpathquadraticcurveto{\pgfqpoint{1.634972in}{7.086044in}}{\pgfqpoint{1.634972in}{7.071856in}}%
\pgfpathquadraticcurveto{\pgfqpoint{1.634972in}{7.057667in}}{\pgfqpoint{1.641812in}{7.049809in}}%
\pgfpathquadraticcurveto{\pgfqpoint{1.648674in}{7.041973in}}{\pgfqpoint{1.661070in}{7.041973in}}%
\pgfpathquadraticcurveto{\pgfqpoint{1.666604in}{7.041973in}}{\pgfqpoint{1.672027in}{7.043456in}}%
\pgfpathquadraticcurveto{\pgfqpoint{1.677450in}{7.044939in}}{\pgfqpoint{1.682851in}{7.047927in}}%
\pgfpathlineto{\pgfqpoint{1.682851in}{7.036151in}}%
\pgfpathquadraticcurveto{\pgfqpoint{1.677516in}{7.033672in}}{\pgfqpoint{1.671805in}{7.032432in}}%
\pgfpathquadraticcurveto{\pgfqpoint{1.666117in}{7.031171in}}{\pgfqpoint{1.659675in}{7.031171in}}%
\pgfpathquadraticcurveto{\pgfqpoint{1.642166in}{7.031171in}}{\pgfqpoint{1.631851in}{7.042172in}}%
\pgfpathquadraticcurveto{\pgfqpoint{1.621558in}{7.053173in}}{\pgfqpoint{1.621558in}{7.071856in}}%
\pgfpathquadraticcurveto{\pgfqpoint{1.621558in}{7.090803in}}{\pgfqpoint{1.631962in}{7.101650in}}%
\pgfpathquadraticcurveto{\pgfqpoint{1.642387in}{7.112518in}}{\pgfqpoint{1.660516in}{7.112518in}}%
\pgfpathquadraticcurveto{\pgfqpoint{1.666382in}{7.112518in}}{\pgfqpoint{1.671982in}{7.111301in}}%
\pgfpathquadraticcurveto{\pgfqpoint{1.677583in}{7.110106in}}{\pgfqpoint{1.682851in}{7.107693in}}%
\pgfpathlineto{\pgfqpoint{1.682851in}{7.107693in}}%
\pgfpathclose%
\pgfpathmoveto{\pgfqpoint{1.740204in}{7.072121in}}%
\pgfpathquadraticcurveto{\pgfqpoint{1.724775in}{7.072121in}}{\pgfqpoint{1.718821in}{7.068602in}}%
\pgfpathquadraticcurveto{\pgfqpoint{1.712867in}{7.065082in}}{\pgfqpoint{1.712867in}{7.056560in}}%
\pgfpathquadraticcurveto{\pgfqpoint{1.712867in}{7.049787in}}{\pgfqpoint{1.717338in}{7.045802in}}%
\pgfpathquadraticcurveto{\pgfqpoint{1.721809in}{7.041840in}}{\pgfqpoint{1.729468in}{7.041840in}}%
\pgfpathquadraticcurveto{\pgfqpoint{1.740071in}{7.041840in}}{\pgfqpoint{1.746468in}{7.049344in}}%
\pgfpathquadraticcurveto{\pgfqpoint{1.752865in}{7.056848in}}{\pgfqpoint{1.752865in}{7.069288in}}%
\pgfpathlineto{\pgfqpoint{1.752865in}{7.072121in}}%
\pgfpathlineto{\pgfqpoint{1.740204in}{7.072121in}}%
\pgfpathlineto{\pgfqpoint{1.740204in}{7.072121in}}%
\pgfpathclose%
\pgfpathmoveto{\pgfqpoint{1.765593in}{7.077389in}}%
\pgfpathlineto{\pgfqpoint{1.765593in}{7.033185in}}%
\pgfpathlineto{\pgfqpoint{1.752865in}{7.033185in}}%
\pgfpathlineto{\pgfqpoint{1.752865in}{7.044939in}}%
\pgfpathquadraticcurveto{\pgfqpoint{1.748505in}{7.037900in}}{\pgfqpoint{1.741997in}{7.034535in}}%
\pgfpathquadraticcurveto{\pgfqpoint{1.735489in}{7.031171in}}{\pgfqpoint{1.726081in}{7.031171in}}%
\pgfpathquadraticcurveto{\pgfqpoint{1.714195in}{7.031171in}}{\pgfqpoint{1.707156in}{7.037856in}}%
\pgfpathquadraticcurveto{\pgfqpoint{1.700139in}{7.044540in}}{\pgfqpoint{1.700139in}{7.055741in}}%
\pgfpathquadraticcurveto{\pgfqpoint{1.700139in}{7.068801in}}{\pgfqpoint{1.708882in}{7.075442in}}%
\pgfpathquadraticcurveto{\pgfqpoint{1.717648in}{7.082082in}}{\pgfqpoint{1.735002in}{7.082082in}}%
\pgfpathlineto{\pgfqpoint{1.752865in}{7.082082in}}%
\pgfpathlineto{\pgfqpoint{1.752865in}{7.083344in}}%
\pgfpathquadraticcurveto{\pgfqpoint{1.752865in}{7.092132in}}{\pgfqpoint{1.747088in}{7.096935in}}%
\pgfpathquadraticcurveto{\pgfqpoint{1.741311in}{7.101738in}}{\pgfqpoint{1.730863in}{7.101738in}}%
\pgfpathquadraticcurveto{\pgfqpoint{1.724222in}{7.101738in}}{\pgfqpoint{1.717913in}{7.100145in}}%
\pgfpathquadraticcurveto{\pgfqpoint{1.711627in}{7.098551in}}{\pgfqpoint{1.705827in}{7.095363in}}%
\pgfpathlineto{\pgfqpoint{1.705827in}{7.107139in}}%
\pgfpathquadraticcurveto{\pgfqpoint{1.712800in}{7.109840in}}{\pgfqpoint{1.719374in}{7.111168in}}%
\pgfpathquadraticcurveto{\pgfqpoint{1.725949in}{7.112518in}}{\pgfqpoint{1.732169in}{7.112518in}}%
\pgfpathquadraticcurveto{\pgfqpoint{1.748992in}{7.112518in}}{\pgfqpoint{1.757292in}{7.103797in}}%
\pgfpathquadraticcurveto{\pgfqpoint{1.765593in}{7.095098in}}{\pgfqpoint{1.765593in}{7.077389in}}%
\pgfpathlineto{\pgfqpoint{1.765593in}{7.077389in}}%
\pgfpathclose%
\pgfpathmoveto{\pgfqpoint{1.856215in}{7.079957in}}%
\pgfpathlineto{\pgfqpoint{1.856215in}{7.033185in}}%
\pgfpathlineto{\pgfqpoint{1.843488in}{7.033185in}}%
\pgfpathlineto{\pgfqpoint{1.843488in}{7.079537in}}%
\pgfpathquadraticcurveto{\pgfqpoint{1.843488in}{7.090538in}}{\pgfqpoint{1.839193in}{7.095983in}}%
\pgfpathquadraticcurveto{\pgfqpoint{1.834899in}{7.101451in}}{\pgfqpoint{1.826333in}{7.101451in}}%
\pgfpathquadraticcurveto{\pgfqpoint{1.816018in}{7.101451in}}{\pgfqpoint{1.810063in}{7.094876in}}%
\pgfpathquadraticcurveto{\pgfqpoint{1.804109in}{7.088324in}}{\pgfqpoint{1.804109in}{7.076969in}}%
\pgfpathlineto{\pgfqpoint{1.804109in}{7.033185in}}%
\pgfpathlineto{\pgfqpoint{1.791314in}{7.033185in}}%
\pgfpathlineto{\pgfqpoint{1.791314in}{7.110659in}}%
\pgfpathlineto{\pgfqpoint{1.804109in}{7.110659in}}%
\pgfpathlineto{\pgfqpoint{1.804109in}{7.098617in}}%
\pgfpathquadraticcurveto{\pgfqpoint{1.808691in}{7.105612in}}{\pgfqpoint{1.814867in}{7.109065in}}%
\pgfpathquadraticcurveto{\pgfqpoint{1.821064in}{7.112518in}}{\pgfqpoint{1.829166in}{7.112518in}}%
\pgfpathquadraticcurveto{\pgfqpoint{1.842514in}{7.112518in}}{\pgfqpoint{1.849354in}{7.104262in}}%
\pgfpathquadraticcurveto{\pgfqpoint{1.856215in}{7.096005in}}{\pgfqpoint{1.856215in}{7.079957in}}%
\pgfpathlineto{\pgfqpoint{1.856215in}{7.079957in}}%
\pgfpathclose%
\pgfpathmoveto{\pgfqpoint{1.894178in}{7.132662in}}%
\pgfpathlineto{\pgfqpoint{1.894178in}{7.110659in}}%
\pgfpathlineto{\pgfqpoint{1.920386in}{7.110659in}}%
\pgfpathlineto{\pgfqpoint{1.920386in}{7.100764in}}%
\pgfpathlineto{\pgfqpoint{1.894178in}{7.100764in}}%
\pgfpathlineto{\pgfqpoint{1.894178in}{7.058707in}}%
\pgfpathquadraticcurveto{\pgfqpoint{1.894178in}{7.049233in}}{\pgfqpoint{1.896768in}{7.046533in}}%
\pgfpathquadraticcurveto{\pgfqpoint{1.899357in}{7.043832in}}{\pgfqpoint{1.907326in}{7.043832in}}%
\pgfpathlineto{\pgfqpoint{1.920386in}{7.043832in}}%
\pgfpathlineto{\pgfqpoint{1.920386in}{7.033185in}}%
\pgfpathlineto{\pgfqpoint{1.907326in}{7.033185in}}%
\pgfpathquadraticcurveto{\pgfqpoint{1.892584in}{7.033185in}}{\pgfqpoint{1.886984in}{7.038675in}}%
\pgfpathquadraticcurveto{\pgfqpoint{1.881383in}{7.044186in}}{\pgfqpoint{1.881383in}{7.058707in}}%
\pgfpathlineto{\pgfqpoint{1.881383in}{7.100764in}}%
\pgfpathlineto{\pgfqpoint{1.872042in}{7.100764in}}%
\pgfpathlineto{\pgfqpoint{1.872042in}{7.110659in}}%
\pgfpathlineto{\pgfqpoint{1.881383in}{7.110659in}}%
\pgfpathlineto{\pgfqpoint{1.881383in}{7.132662in}}%
\pgfpathlineto{\pgfqpoint{1.894178in}{7.132662in}}%
\pgfpathlineto{\pgfqpoint{1.894178in}{7.132662in}}%
\pgfpathclose%
\pgfpathmoveto{\pgfqpoint{2.033122in}{7.098905in}}%
\pgfpathlineto{\pgfqpoint{2.033122in}{7.140830in}}%
\pgfpathlineto{\pgfqpoint{2.045850in}{7.140830in}}%
\pgfpathlineto{\pgfqpoint{2.045850in}{7.033185in}}%
\pgfpathlineto{\pgfqpoint{2.033122in}{7.033185in}}%
\pgfpathlineto{\pgfqpoint{2.033122in}{7.044806in}}%
\pgfpathquadraticcurveto{\pgfqpoint{2.029115in}{7.037900in}}{\pgfqpoint{2.022984in}{7.034535in}}%
\pgfpathquadraticcurveto{\pgfqpoint{2.016874in}{7.031171in}}{\pgfqpoint{2.008286in}{7.031171in}}%
\pgfpathquadraticcurveto{\pgfqpoint{1.994252in}{7.031171in}}{\pgfqpoint{1.985420in}{7.042371in}}%
\pgfpathquadraticcurveto{\pgfqpoint{1.976610in}{7.053594in}}{\pgfqpoint{1.976610in}{7.071856in}}%
\pgfpathquadraticcurveto{\pgfqpoint{1.976610in}{7.090117in}}{\pgfqpoint{1.985420in}{7.101318in}}%
\pgfpathquadraticcurveto{\pgfqpoint{1.994252in}{7.112518in}}{\pgfqpoint{2.008286in}{7.112518in}}%
\pgfpathquadraticcurveto{\pgfqpoint{2.016874in}{7.112518in}}{\pgfqpoint{2.022984in}{7.109154in}}%
\pgfpathquadraticcurveto{\pgfqpoint{2.029115in}{7.105811in}}{\pgfqpoint{2.033122in}{7.098905in}}%
\pgfpathlineto{\pgfqpoint{2.033122in}{7.098905in}}%
\pgfpathclose%
\pgfpathmoveto{\pgfqpoint{1.989758in}{7.071856in}}%
\pgfpathquadraticcurveto{\pgfqpoint{1.989758in}{7.057822in}}{\pgfqpoint{1.995536in}{7.049831in}}%
\pgfpathquadraticcurveto{\pgfqpoint{2.001313in}{7.041840in}}{\pgfqpoint{2.011407in}{7.041840in}}%
\pgfpathquadraticcurveto{\pgfqpoint{2.021501in}{7.041840in}}{\pgfqpoint{2.027300in}{7.049831in}}%
\pgfpathquadraticcurveto{\pgfqpoint{2.033122in}{7.057822in}}{\pgfqpoint{2.033122in}{7.071856in}}%
\pgfpathquadraticcurveto{\pgfqpoint{2.033122in}{7.085889in}}{\pgfqpoint{2.027300in}{7.093880in}}%
\pgfpathquadraticcurveto{\pgfqpoint{2.021501in}{7.101871in}}{\pgfqpoint{2.011407in}{7.101871in}}%
\pgfpathquadraticcurveto{\pgfqpoint{2.001313in}{7.101871in}}{\pgfqpoint{1.995536in}{7.093880in}}%
\pgfpathquadraticcurveto{\pgfqpoint{1.989758in}{7.085889in}}{\pgfqpoint{1.989758in}{7.071856in}}%
\pgfpathlineto{\pgfqpoint{1.989758in}{7.071856in}}%
\pgfpathclose%
\pgfpathmoveto{\pgfqpoint{2.072080in}{7.110659in}}%
\pgfpathlineto{\pgfqpoint{2.084808in}{7.110659in}}%
\pgfpathlineto{\pgfqpoint{2.084808in}{7.033185in}}%
\pgfpathlineto{\pgfqpoint{2.072080in}{7.033185in}}%
\pgfpathlineto{\pgfqpoint{2.072080in}{7.110659in}}%
\pgfpathlineto{\pgfqpoint{2.072080in}{7.110659in}}%
\pgfpathclose%
\pgfpathmoveto{\pgfqpoint{2.072080in}{7.140830in}}%
\pgfpathlineto{\pgfqpoint{2.084808in}{7.140830in}}%
\pgfpathlineto{\pgfqpoint{2.084808in}{7.124693in}}%
\pgfpathlineto{\pgfqpoint{2.072080in}{7.124693in}}%
\pgfpathlineto{\pgfqpoint{2.072080in}{7.140830in}}%
\pgfpathlineto{\pgfqpoint{2.072080in}{7.140830in}}%
\pgfpathclose%
\pgfpathmoveto{\pgfqpoint{2.198385in}{7.140830in}}%
\pgfpathlineto{\pgfqpoint{2.198385in}{7.130227in}}%
\pgfpathlineto{\pgfqpoint{2.186210in}{7.130227in}}%
\pgfpathquadraticcurveto{\pgfqpoint{2.179370in}{7.130227in}}{\pgfqpoint{2.176692in}{7.127460in}}%
\pgfpathquadraticcurveto{\pgfqpoint{2.174036in}{7.124693in}}{\pgfqpoint{2.174036in}{7.117499in}}%
\pgfpathlineto{\pgfqpoint{2.174036in}{7.110659in}}%
\pgfpathlineto{\pgfqpoint{2.194998in}{7.110659in}}%
\pgfpathlineto{\pgfqpoint{2.194998in}{7.100764in}}%
\pgfpathlineto{\pgfqpoint{2.174036in}{7.100764in}}%
\pgfpathlineto{\pgfqpoint{2.174036in}{7.033185in}}%
\pgfpathlineto{\pgfqpoint{2.161242in}{7.033185in}}%
\pgfpathlineto{\pgfqpoint{2.161242in}{7.100764in}}%
\pgfpathlineto{\pgfqpoint{2.126312in}{7.100764in}}%
\pgfpathlineto{\pgfqpoint{2.126312in}{7.033185in}}%
\pgfpathlineto{\pgfqpoint{2.113518in}{7.033185in}}%
\pgfpathlineto{\pgfqpoint{2.113518in}{7.100764in}}%
\pgfpathlineto{\pgfqpoint{2.101343in}{7.100764in}}%
\pgfpathlineto{\pgfqpoint{2.101343in}{7.110659in}}%
\pgfpathlineto{\pgfqpoint{2.113518in}{7.110659in}}%
\pgfpathlineto{\pgfqpoint{2.113518in}{7.116060in}}%
\pgfpathquadraticcurveto{\pgfqpoint{2.113518in}{7.128987in}}{\pgfqpoint{2.119538in}{7.134897in}}%
\pgfpathquadraticcurveto{\pgfqpoint{2.125559in}{7.140830in}}{\pgfqpoint{2.138619in}{7.140830in}}%
\pgfpathlineto{\pgfqpoint{2.150661in}{7.140830in}}%
\pgfpathlineto{\pgfqpoint{2.150661in}{7.130227in}}%
\pgfpathlineto{\pgfqpoint{2.138486in}{7.130227in}}%
\pgfpathquadraticcurveto{\pgfqpoint{2.131646in}{7.130227in}}{\pgfqpoint{2.128946in}{7.127460in}}%
\pgfpathquadraticcurveto{\pgfqpoint{2.126312in}{7.124693in}}{\pgfqpoint{2.126312in}{7.117499in}}%
\pgfpathlineto{\pgfqpoint{2.126312in}{7.110659in}}%
\pgfpathlineto{\pgfqpoint{2.161242in}{7.110659in}}%
\pgfpathlineto{\pgfqpoint{2.161242in}{7.116060in}}%
\pgfpathquadraticcurveto{\pgfqpoint{2.161242in}{7.128987in}}{\pgfqpoint{2.167262in}{7.134897in}}%
\pgfpathquadraticcurveto{\pgfqpoint{2.173283in}{7.140830in}}{\pgfqpoint{2.186365in}{7.140830in}}%
\pgfpathlineto{\pgfqpoint{2.198385in}{7.140830in}}%
\pgfpathlineto{\pgfqpoint{2.198385in}{7.140830in}}%
\pgfpathclose%
\pgfpathmoveto{\pgfqpoint{2.275305in}{7.075109in}}%
\pgfpathlineto{\pgfqpoint{2.275305in}{7.068889in}}%
\pgfpathlineto{\pgfqpoint{2.216779in}{7.068889in}}%
\pgfpathquadraticcurveto{\pgfqpoint{2.217620in}{7.055741in}}{\pgfqpoint{2.224704in}{7.048857in}}%
\pgfpathquadraticcurveto{\pgfqpoint{2.231787in}{7.041973in}}{\pgfqpoint{2.244449in}{7.041973in}}%
\pgfpathquadraticcurveto{\pgfqpoint{2.251775in}{7.041973in}}{\pgfqpoint{2.258660in}{7.043766in}}%
\pgfpathquadraticcurveto{\pgfqpoint{2.265544in}{7.045559in}}{\pgfqpoint{2.272339in}{7.049167in}}%
\pgfpathlineto{\pgfqpoint{2.272339in}{7.037125in}}%
\pgfpathquadraticcurveto{\pgfqpoint{2.265477in}{7.034225in}}{\pgfqpoint{2.258283in}{7.032698in}}%
\pgfpathquadraticcurveto{\pgfqpoint{2.251089in}{7.031171in}}{\pgfqpoint{2.243696in}{7.031171in}}%
\pgfpathquadraticcurveto{\pgfqpoint{2.225146in}{7.031171in}}{\pgfqpoint{2.214322in}{7.041951in}}%
\pgfpathquadraticcurveto{\pgfqpoint{2.203498in}{7.052753in}}{\pgfqpoint{2.203498in}{7.071169in}}%
\pgfpathquadraticcurveto{\pgfqpoint{2.203498in}{7.090184in}}{\pgfqpoint{2.213769in}{7.101340in}}%
\pgfpathquadraticcurveto{\pgfqpoint{2.224040in}{7.112518in}}{\pgfqpoint{2.241482in}{7.112518in}}%
\pgfpathquadraticcurveto{\pgfqpoint{2.257110in}{7.112518in}}{\pgfqpoint{2.266208in}{7.102447in}}%
\pgfpathquadraticcurveto{\pgfqpoint{2.275305in}{7.092397in}}{\pgfqpoint{2.275305in}{7.075109in}}%
\pgfpathlineto{\pgfqpoint{2.275305in}{7.075109in}}%
\pgfpathclose%
\pgfpathmoveto{\pgfqpoint{2.262577in}{7.078850in}}%
\pgfpathquadraticcurveto{\pgfqpoint{2.262445in}{7.089276in}}{\pgfqpoint{2.256734in}{7.095496in}}%
\pgfpathquadraticcurveto{\pgfqpoint{2.251023in}{7.101738in}}{\pgfqpoint{2.241615in}{7.101738in}}%
\pgfpathquadraticcurveto{\pgfqpoint{2.230968in}{7.101738in}}{\pgfqpoint{2.224571in}{7.095718in}}%
\pgfpathquadraticcurveto{\pgfqpoint{2.218174in}{7.089697in}}{\pgfqpoint{2.217200in}{7.078762in}}%
\pgfpathlineto{\pgfqpoint{2.262577in}{7.078850in}}%
\pgfpathlineto{\pgfqpoint{2.262577in}{7.078850in}}%
\pgfpathclose%
\pgfpathmoveto{\pgfqpoint{2.341092in}{7.098772in}}%
\pgfpathquadraticcurveto{\pgfqpoint{2.338945in}{7.100012in}}{\pgfqpoint{2.336421in}{7.100587in}}%
\pgfpathquadraticcurveto{\pgfqpoint{2.333898in}{7.101185in}}{\pgfqpoint{2.330865in}{7.101185in}}%
\pgfpathquadraticcurveto{\pgfqpoint{2.320063in}{7.101185in}}{\pgfqpoint{2.314286in}{7.094168in}}%
\pgfpathquadraticcurveto{\pgfqpoint{2.308508in}{7.087151in}}{\pgfqpoint{2.308508in}{7.074003in}}%
\pgfpathlineto{\pgfqpoint{2.308508in}{7.033185in}}%
\pgfpathlineto{\pgfqpoint{2.295714in}{7.033185in}}%
\pgfpathlineto{\pgfqpoint{2.295714in}{7.110659in}}%
\pgfpathlineto{\pgfqpoint{2.308508in}{7.110659in}}%
\pgfpathlineto{\pgfqpoint{2.308508in}{7.098617in}}%
\pgfpathquadraticcurveto{\pgfqpoint{2.312537in}{7.105678in}}{\pgfqpoint{2.318956in}{7.109087in}}%
\pgfpathquadraticcurveto{\pgfqpoint{2.325398in}{7.112518in}}{\pgfqpoint{2.334606in}{7.112518in}}%
\pgfpathquadraticcurveto{\pgfqpoint{2.335912in}{7.112518in}}{\pgfqpoint{2.337506in}{7.112341in}}%
\pgfpathquadraticcurveto{\pgfqpoint{2.339100in}{7.112186in}}{\pgfqpoint{2.341025in}{7.111832in}}%
\pgfpathlineto{\pgfqpoint{2.341092in}{7.098772in}}%
\pgfpathlineto{\pgfqpoint{2.341092in}{7.098772in}}%
\pgfpathclose%
\pgfpathmoveto{\pgfqpoint{2.417592in}{7.075109in}}%
\pgfpathlineto{\pgfqpoint{2.417592in}{7.068889in}}%
\pgfpathlineto{\pgfqpoint{2.359066in}{7.068889in}}%
\pgfpathquadraticcurveto{\pgfqpoint{2.359907in}{7.055741in}}{\pgfqpoint{2.366990in}{7.048857in}}%
\pgfpathquadraticcurveto{\pgfqpoint{2.374074in}{7.041973in}}{\pgfqpoint{2.386735in}{7.041973in}}%
\pgfpathquadraticcurveto{\pgfqpoint{2.394062in}{7.041973in}}{\pgfqpoint{2.400946in}{7.043766in}}%
\pgfpathquadraticcurveto{\pgfqpoint{2.407830in}{7.045559in}}{\pgfqpoint{2.414626in}{7.049167in}}%
\pgfpathlineto{\pgfqpoint{2.414626in}{7.037125in}}%
\pgfpathquadraticcurveto{\pgfqpoint{2.407764in}{7.034225in}}{\pgfqpoint{2.400570in}{7.032698in}}%
\pgfpathquadraticcurveto{\pgfqpoint{2.393376in}{7.031171in}}{\pgfqpoint{2.385982in}{7.031171in}}%
\pgfpathquadraticcurveto{\pgfqpoint{2.367433in}{7.031171in}}{\pgfqpoint{2.356609in}{7.041951in}}%
\pgfpathquadraticcurveto{\pgfqpoint{2.345785in}{7.052753in}}{\pgfqpoint{2.345785in}{7.071169in}}%
\pgfpathquadraticcurveto{\pgfqpoint{2.345785in}{7.090184in}}{\pgfqpoint{2.356055in}{7.101340in}}%
\pgfpathquadraticcurveto{\pgfqpoint{2.366326in}{7.112518in}}{\pgfqpoint{2.383769in}{7.112518in}}%
\pgfpathquadraticcurveto{\pgfqpoint{2.399396in}{7.112518in}}{\pgfqpoint{2.408494in}{7.102447in}}%
\pgfpathquadraticcurveto{\pgfqpoint{2.417592in}{7.092397in}}{\pgfqpoint{2.417592in}{7.075109in}}%
\pgfpathlineto{\pgfqpoint{2.417592in}{7.075109in}}%
\pgfpathclose%
\pgfpathmoveto{\pgfqpoint{2.404864in}{7.078850in}}%
\pgfpathquadraticcurveto{\pgfqpoint{2.404731in}{7.089276in}}{\pgfqpoint{2.399020in}{7.095496in}}%
\pgfpathquadraticcurveto{\pgfqpoint{2.393309in}{7.101738in}}{\pgfqpoint{2.383902in}{7.101738in}}%
\pgfpathquadraticcurveto{\pgfqpoint{2.373255in}{7.101738in}}{\pgfqpoint{2.366857in}{7.095718in}}%
\pgfpathquadraticcurveto{\pgfqpoint{2.360460in}{7.089697in}}{\pgfqpoint{2.359486in}{7.078762in}}%
\pgfpathlineto{\pgfqpoint{2.404864in}{7.078850in}}%
\pgfpathlineto{\pgfqpoint{2.404864in}{7.078850in}}%
\pgfpathclose%
\pgfpathmoveto{\pgfqpoint{2.502902in}{7.079957in}}%
\pgfpathlineto{\pgfqpoint{2.502902in}{7.033185in}}%
\pgfpathlineto{\pgfqpoint{2.490174in}{7.033185in}}%
\pgfpathlineto{\pgfqpoint{2.490174in}{7.079537in}}%
\pgfpathquadraticcurveto{\pgfqpoint{2.490174in}{7.090538in}}{\pgfqpoint{2.485880in}{7.095983in}}%
\pgfpathquadraticcurveto{\pgfqpoint{2.481585in}{7.101451in}}{\pgfqpoint{2.473019in}{7.101451in}}%
\pgfpathquadraticcurveto{\pgfqpoint{2.462704in}{7.101451in}}{\pgfqpoint{2.456749in}{7.094876in}}%
\pgfpathquadraticcurveto{\pgfqpoint{2.450795in}{7.088324in}}{\pgfqpoint{2.450795in}{7.076969in}}%
\pgfpathlineto{\pgfqpoint{2.450795in}{7.033185in}}%
\pgfpathlineto{\pgfqpoint{2.438001in}{7.033185in}}%
\pgfpathlineto{\pgfqpoint{2.438001in}{7.110659in}}%
\pgfpathlineto{\pgfqpoint{2.450795in}{7.110659in}}%
\pgfpathlineto{\pgfqpoint{2.450795in}{7.098617in}}%
\pgfpathquadraticcurveto{\pgfqpoint{2.455377in}{7.105612in}}{\pgfqpoint{2.461553in}{7.109065in}}%
\pgfpathquadraticcurveto{\pgfqpoint{2.467751in}{7.112518in}}{\pgfqpoint{2.475852in}{7.112518in}}%
\pgfpathquadraticcurveto{\pgfqpoint{2.489200in}{7.112518in}}{\pgfqpoint{2.496040in}{7.104262in}}%
\pgfpathquadraticcurveto{\pgfqpoint{2.502902in}{7.096005in}}{\pgfqpoint{2.502902in}{7.079957in}}%
\pgfpathlineto{\pgfqpoint{2.502902in}{7.079957in}}%
\pgfpathclose%
\pgfpathmoveto{\pgfqpoint{2.584028in}{7.107693in}}%
\pgfpathlineto{\pgfqpoint{2.584028in}{7.095784in}}%
\pgfpathquadraticcurveto{\pgfqpoint{2.578627in}{7.098772in}}{\pgfqpoint{2.573204in}{7.100255in}}%
\pgfpathquadraticcurveto{\pgfqpoint{2.567781in}{7.101738in}}{\pgfqpoint{2.562247in}{7.101738in}}%
\pgfpathquadraticcurveto{\pgfqpoint{2.549851in}{7.101738in}}{\pgfqpoint{2.542989in}{7.093880in}}%
\pgfpathquadraticcurveto{\pgfqpoint{2.536149in}{7.086044in}}{\pgfqpoint{2.536149in}{7.071856in}}%
\pgfpathquadraticcurveto{\pgfqpoint{2.536149in}{7.057667in}}{\pgfqpoint{2.542989in}{7.049809in}}%
\pgfpathquadraticcurveto{\pgfqpoint{2.549851in}{7.041973in}}{\pgfqpoint{2.562247in}{7.041973in}}%
\pgfpathquadraticcurveto{\pgfqpoint{2.567781in}{7.041973in}}{\pgfqpoint{2.573204in}{7.043456in}}%
\pgfpathquadraticcurveto{\pgfqpoint{2.578627in}{7.044939in}}{\pgfqpoint{2.584028in}{7.047927in}}%
\pgfpathlineto{\pgfqpoint{2.584028in}{7.036151in}}%
\pgfpathquadraticcurveto{\pgfqpoint{2.578693in}{7.033672in}}{\pgfqpoint{2.572982in}{7.032432in}}%
\pgfpathquadraticcurveto{\pgfqpoint{2.567294in}{7.031171in}}{\pgfqpoint{2.560852in}{7.031171in}}%
\pgfpathquadraticcurveto{\pgfqpoint{2.543343in}{7.031171in}}{\pgfqpoint{2.533028in}{7.042172in}}%
\pgfpathquadraticcurveto{\pgfqpoint{2.522735in}{7.053173in}}{\pgfqpoint{2.522735in}{7.071856in}}%
\pgfpathquadraticcurveto{\pgfqpoint{2.522735in}{7.090803in}}{\pgfqpoint{2.533139in}{7.101650in}}%
\pgfpathquadraticcurveto{\pgfqpoint{2.543564in}{7.112518in}}{\pgfqpoint{2.561693in}{7.112518in}}%
\pgfpathquadraticcurveto{\pgfqpoint{2.567559in}{7.112518in}}{\pgfqpoint{2.573160in}{7.111301in}}%
\pgfpathquadraticcurveto{\pgfqpoint{2.578760in}{7.110106in}}{\pgfqpoint{2.584028in}{7.107693in}}%
\pgfpathlineto{\pgfqpoint{2.584028in}{7.107693in}}%
\pgfpathclose%
\pgfpathmoveto{\pgfqpoint{2.672437in}{7.075109in}}%
\pgfpathlineto{\pgfqpoint{2.672437in}{7.068889in}}%
\pgfpathlineto{\pgfqpoint{2.613911in}{7.068889in}}%
\pgfpathquadraticcurveto{\pgfqpoint{2.614752in}{7.055741in}}{\pgfqpoint{2.621835in}{7.048857in}}%
\pgfpathquadraticcurveto{\pgfqpoint{2.628919in}{7.041973in}}{\pgfqpoint{2.641580in}{7.041973in}}%
\pgfpathquadraticcurveto{\pgfqpoint{2.648907in}{7.041973in}}{\pgfqpoint{2.655791in}{7.043766in}}%
\pgfpathquadraticcurveto{\pgfqpoint{2.662675in}{7.045559in}}{\pgfqpoint{2.669471in}{7.049167in}}%
\pgfpathlineto{\pgfqpoint{2.669471in}{7.037125in}}%
\pgfpathquadraticcurveto{\pgfqpoint{2.662609in}{7.034225in}}{\pgfqpoint{2.655415in}{7.032698in}}%
\pgfpathquadraticcurveto{\pgfqpoint{2.648221in}{7.031171in}}{\pgfqpoint{2.640827in}{7.031171in}}%
\pgfpathquadraticcurveto{\pgfqpoint{2.622278in}{7.031171in}}{\pgfqpoint{2.611454in}{7.041951in}}%
\pgfpathquadraticcurveto{\pgfqpoint{2.600630in}{7.052753in}}{\pgfqpoint{2.600630in}{7.071169in}}%
\pgfpathquadraticcurveto{\pgfqpoint{2.600630in}{7.090184in}}{\pgfqpoint{2.610900in}{7.101340in}}%
\pgfpathquadraticcurveto{\pgfqpoint{2.621171in}{7.112518in}}{\pgfqpoint{2.638614in}{7.112518in}}%
\pgfpathquadraticcurveto{\pgfqpoint{2.654242in}{7.112518in}}{\pgfqpoint{2.663339in}{7.102447in}}%
\pgfpathquadraticcurveto{\pgfqpoint{2.672437in}{7.092397in}}{\pgfqpoint{2.672437in}{7.075109in}}%
\pgfpathlineto{\pgfqpoint{2.672437in}{7.075109in}}%
\pgfpathclose%
\pgfpathmoveto{\pgfqpoint{2.659709in}{7.078850in}}%
\pgfpathquadraticcurveto{\pgfqpoint{2.659576in}{7.089276in}}{\pgfqpoint{2.653865in}{7.095496in}}%
\pgfpathquadraticcurveto{\pgfqpoint{2.648154in}{7.101738in}}{\pgfqpoint{2.638747in}{7.101738in}}%
\pgfpathquadraticcurveto{\pgfqpoint{2.628100in}{7.101738in}}{\pgfqpoint{2.621702in}{7.095718in}}%
\pgfpathquadraticcurveto{\pgfqpoint{2.615305in}{7.089697in}}{\pgfqpoint{2.614331in}{7.078762in}}%
\pgfpathlineto{\pgfqpoint{2.659709in}{7.078850in}}%
\pgfpathlineto{\pgfqpoint{2.659709in}{7.078850in}}%
\pgfpathclose%
\pgfpathmoveto{\pgfqpoint{2.742717in}{7.108379in}}%
\pgfpathlineto{\pgfqpoint{2.742717in}{7.096337in}}%
\pgfpathquadraticcurveto{\pgfqpoint{2.737338in}{7.099104in}}{\pgfqpoint{2.731516in}{7.100477in}}%
\pgfpathquadraticcurveto{\pgfqpoint{2.725717in}{7.101871in}}{\pgfqpoint{2.719475in}{7.101871in}}%
\pgfpathquadraticcurveto{\pgfqpoint{2.710001in}{7.101871in}}{\pgfqpoint{2.705264in}{7.098971in}}%
\pgfpathquadraticcurveto{\pgfqpoint{2.700527in}{7.096072in}}{\pgfqpoint{2.700527in}{7.090250in}}%
\pgfpathquadraticcurveto{\pgfqpoint{2.700527in}{7.085823in}}{\pgfqpoint{2.703913in}{7.083300in}}%
\pgfpathquadraticcurveto{\pgfqpoint{2.707300in}{7.080776in}}{\pgfqpoint{2.717549in}{7.078496in}}%
\pgfpathlineto{\pgfqpoint{2.721910in}{7.077522in}}%
\pgfpathquadraticcurveto{\pgfqpoint{2.735456in}{7.074623in}}{\pgfqpoint{2.741167in}{7.069332in}}%
\pgfpathquadraticcurveto{\pgfqpoint{2.746878in}{7.064042in}}{\pgfqpoint{2.746878in}{7.054568in}}%
\pgfpathquadraticcurveto{\pgfqpoint{2.746878in}{7.043766in}}{\pgfqpoint{2.738334in}{7.037457in}}%
\pgfpathquadraticcurveto{\pgfqpoint{2.729790in}{7.031171in}}{\pgfqpoint{2.714848in}{7.031171in}}%
\pgfpathquadraticcurveto{\pgfqpoint{2.708628in}{7.031171in}}{\pgfqpoint{2.701877in}{7.032388in}}%
\pgfpathquadraticcurveto{\pgfqpoint{2.695126in}{7.033606in}}{\pgfqpoint{2.687666in}{7.036018in}}%
\pgfpathlineto{\pgfqpoint{2.687666in}{7.049167in}}%
\pgfpathquadraticcurveto{\pgfqpoint{2.694727in}{7.045492in}}{\pgfqpoint{2.701567in}{7.043655in}}%
\pgfpathquadraticcurveto{\pgfqpoint{2.708407in}{7.041840in}}{\pgfqpoint{2.715136in}{7.041840in}}%
\pgfpathquadraticcurveto{\pgfqpoint{2.724123in}{7.041840in}}{\pgfqpoint{2.728949in}{7.044917in}}%
\pgfpathquadraticcurveto{\pgfqpoint{2.733796in}{7.047994in}}{\pgfqpoint{2.733796in}{7.053594in}}%
\pgfpathquadraticcurveto{\pgfqpoint{2.733796in}{7.058774in}}{\pgfqpoint{2.730299in}{7.061540in}}%
\pgfpathquadraticcurveto{\pgfqpoint{2.726824in}{7.064307in}}{\pgfqpoint{2.714981in}{7.066875in}}%
\pgfpathlineto{\pgfqpoint{2.710554in}{7.067915in}}%
\pgfpathquadraticcurveto{\pgfqpoint{2.698734in}{7.070395in}}{\pgfqpoint{2.693465in}{7.075552in}}%
\pgfpathquadraticcurveto{\pgfqpoint{2.688219in}{7.080710in}}{\pgfqpoint{2.688219in}{7.089697in}}%
\pgfpathquadraticcurveto{\pgfqpoint{2.688219in}{7.100632in}}{\pgfqpoint{2.695967in}{7.106564in}}%
\pgfpathquadraticcurveto{\pgfqpoint{2.703714in}{7.112518in}}{\pgfqpoint{2.717969in}{7.112518in}}%
\pgfpathquadraticcurveto{\pgfqpoint{2.725008in}{7.112518in}}{\pgfqpoint{2.731229in}{7.111478in}}%
\pgfpathquadraticcurveto{\pgfqpoint{2.737471in}{7.110460in}}{\pgfqpoint{2.742717in}{7.108379in}}%
\pgfpathlineto{\pgfqpoint{2.742717in}{7.108379in}}%
\pgfpathclose%
\pgfpathmoveto{\pgfqpoint{2.770386in}{7.050761in}}%
\pgfpathlineto{\pgfqpoint{2.784973in}{7.050761in}}%
\pgfpathlineto{\pgfqpoint{2.784973in}{7.033185in}}%
\pgfpathlineto{\pgfqpoint{2.770386in}{7.033185in}}%
\pgfpathlineto{\pgfqpoint{2.770386in}{7.050761in}}%
\pgfpathlineto{\pgfqpoint{2.770386in}{7.050761in}}%
\pgfpathclose%
\pgfpathmoveto{\pgfqpoint{2.770386in}{7.106431in}}%
\pgfpathlineto{\pgfqpoint{2.784973in}{7.106431in}}%
\pgfpathlineto{\pgfqpoint{2.784973in}{7.088878in}}%
\pgfpathlineto{\pgfqpoint{2.770386in}{7.088878in}}%
\pgfpathlineto{\pgfqpoint{2.770386in}{7.106431in}}%
\pgfpathlineto{\pgfqpoint{2.770386in}{7.106431in}}%
\pgfpathclose%
\pgfpathmoveto{\pgfqpoint{2.858153in}{7.136469in}}%
\pgfpathlineto{\pgfqpoint{2.924559in}{7.136469in}}%
\pgfpathlineto{\pgfqpoint{2.924559in}{7.130514in}}%
\pgfpathlineto{\pgfqpoint{2.887062in}{7.033185in}}%
\pgfpathlineto{\pgfqpoint{2.872475in}{7.033185in}}%
\pgfpathlineto{\pgfqpoint{2.907758in}{7.124693in}}%
\pgfpathlineto{\pgfqpoint{2.858153in}{7.124693in}}%
\pgfpathlineto{\pgfqpoint{2.858153in}{7.136469in}}%
\pgfpathlineto{\pgfqpoint{2.858153in}{7.136469in}}%
\pgfpathclose%
\pgfpathmoveto{\pgfqpoint{2.951808in}{7.050761in}}%
\pgfpathlineto{\pgfqpoint{2.966417in}{7.050761in}}%
\pgfpathlineto{\pgfqpoint{2.966417in}{7.033185in}}%
\pgfpathlineto{\pgfqpoint{2.951808in}{7.033185in}}%
\pgfpathlineto{\pgfqpoint{2.951808in}{7.050761in}}%
\pgfpathlineto{\pgfqpoint{2.951808in}{7.050761in}}%
\pgfpathclose%
\pgfpathmoveto{\pgfqpoint{2.993312in}{7.136469in}}%
\pgfpathlineto{\pgfqpoint{3.059718in}{7.136469in}}%
\pgfpathlineto{\pgfqpoint{3.059718in}{7.130514in}}%
\pgfpathlineto{\pgfqpoint{3.022221in}{7.033185in}}%
\pgfpathlineto{\pgfqpoint{3.007633in}{7.033185in}}%
\pgfpathlineto{\pgfqpoint{3.042917in}{7.124693in}}%
\pgfpathlineto{\pgfqpoint{2.993312in}{7.124693in}}%
\pgfpathlineto{\pgfqpoint{2.993312in}{7.136469in}}%
\pgfpathlineto{\pgfqpoint{2.993312in}{7.136469in}}%
\pgfpathclose%
\pgfpathmoveto{\pgfqpoint{3.174822in}{7.078629in}}%
\pgfpathquadraticcurveto{\pgfqpoint{3.168801in}{7.078629in}}{\pgfqpoint{3.165370in}{7.073516in}}%
\pgfpathquadraticcurveto{\pgfqpoint{3.161962in}{7.068402in}}{\pgfqpoint{3.161962in}{7.059261in}}%
\pgfpathquadraticcurveto{\pgfqpoint{3.161962in}{7.050274in}}{\pgfqpoint{3.165370in}{7.045116in}}%
\pgfpathquadraticcurveto{\pgfqpoint{3.168801in}{7.039958in}}{\pgfqpoint{3.174822in}{7.039958in}}%
\pgfpathquadraticcurveto{\pgfqpoint{3.180710in}{7.039958in}}{\pgfqpoint{3.184119in}{7.045116in}}%
\pgfpathquadraticcurveto{\pgfqpoint{3.187550in}{7.050274in}}{\pgfqpoint{3.187550in}{7.059261in}}%
\pgfpathquadraticcurveto{\pgfqpoint{3.187550in}{7.068336in}}{\pgfqpoint{3.184119in}{7.073471in}}%
\pgfpathquadraticcurveto{\pgfqpoint{3.180710in}{7.078629in}}{\pgfqpoint{3.174822in}{7.078629in}}%
\pgfpathlineto{\pgfqpoint{3.174822in}{7.078629in}}%
\pgfpathclose%
\pgfpathmoveto{\pgfqpoint{3.174822in}{7.087417in}}%
\pgfpathquadraticcurveto{\pgfqpoint{3.185757in}{7.087417in}}{\pgfqpoint{3.192176in}{7.079802in}}%
\pgfpathquadraticcurveto{\pgfqpoint{3.198618in}{7.072210in}}{\pgfqpoint{3.198618in}{7.059261in}}%
\pgfpathquadraticcurveto{\pgfqpoint{3.198618in}{7.046333in}}{\pgfqpoint{3.192154in}{7.038741in}}%
\pgfpathquadraticcurveto{\pgfqpoint{3.185691in}{7.031171in}}{\pgfqpoint{3.174822in}{7.031171in}}%
\pgfpathquadraticcurveto{\pgfqpoint{3.163755in}{7.031171in}}{\pgfqpoint{3.157313in}{7.038741in}}%
\pgfpathquadraticcurveto{\pgfqpoint{3.150894in}{7.046333in}}{\pgfqpoint{3.150894in}{7.059261in}}%
\pgfpathquadraticcurveto{\pgfqpoint{3.150894in}{7.072276in}}{\pgfqpoint{3.157357in}{7.079846in}}%
\pgfpathquadraticcurveto{\pgfqpoint{3.163821in}{7.087417in}}{\pgfqpoint{3.174822in}{7.087417in}}%
\pgfpathlineto{\pgfqpoint{3.174822in}{7.087417in}}%
\pgfpathclose%
\pgfpathmoveto{\pgfqpoint{3.103436in}{7.129540in}}%
\pgfpathquadraticcurveto{\pgfqpoint{3.097481in}{7.129540in}}{\pgfqpoint{3.094050in}{7.124383in}}%
\pgfpathquadraticcurveto{\pgfqpoint{3.090641in}{7.119248in}}{\pgfqpoint{3.090641in}{7.110238in}}%
\pgfpathquadraticcurveto{\pgfqpoint{3.090641in}{7.101119in}}{\pgfqpoint{3.094028in}{7.095983in}}%
\pgfpathquadraticcurveto{\pgfqpoint{3.097415in}{7.090870in}}{\pgfqpoint{3.103436in}{7.090870in}}%
\pgfpathquadraticcurveto{\pgfqpoint{3.109456in}{7.090870in}}{\pgfqpoint{3.112865in}{7.095983in}}%
\pgfpathquadraticcurveto{\pgfqpoint{3.116296in}{7.101119in}}{\pgfqpoint{3.116296in}{7.110238in}}%
\pgfpathquadraticcurveto{\pgfqpoint{3.116296in}{7.119159in}}{\pgfqpoint{3.112843in}{7.124339in}}%
\pgfpathquadraticcurveto{\pgfqpoint{3.109390in}{7.129540in}}{\pgfqpoint{3.103436in}{7.129540in}}%
\pgfpathlineto{\pgfqpoint{3.103436in}{7.129540in}}%
\pgfpathclose%
\pgfpathmoveto{\pgfqpoint{3.165902in}{7.138328in}}%
\pgfpathlineto{\pgfqpoint{3.176969in}{7.138328in}}%
\pgfpathlineto{\pgfqpoint{3.112356in}{7.031171in}}%
\pgfpathlineto{\pgfqpoint{3.101288in}{7.031171in}}%
\pgfpathlineto{\pgfqpoint{3.165902in}{7.138328in}}%
\pgfpathlineto{\pgfqpoint{3.165902in}{7.138328in}}%
\pgfpathclose%
\pgfpathmoveto{\pgfqpoint{3.103436in}{7.138328in}}%
\pgfpathquadraticcurveto{\pgfqpoint{3.114370in}{7.138328in}}{\pgfqpoint{3.120856in}{7.130758in}}%
\pgfpathquadraticcurveto{\pgfqpoint{3.127364in}{7.123188in}}{\pgfqpoint{3.127364in}{7.110238in}}%
\pgfpathquadraticcurveto{\pgfqpoint{3.127364in}{7.097178in}}{\pgfqpoint{3.120900in}{7.089630in}}%
\pgfpathquadraticcurveto{\pgfqpoint{3.114437in}{7.082082in}}{\pgfqpoint{3.103436in}{7.082082in}}%
\pgfpathquadraticcurveto{\pgfqpoint{3.092434in}{7.082082in}}{\pgfqpoint{3.086037in}{7.089652in}}%
\pgfpathquadraticcurveto{\pgfqpoint{3.079640in}{7.097245in}}{\pgfqpoint{3.079640in}{7.110238in}}%
\pgfpathquadraticcurveto{\pgfqpoint{3.079640in}{7.123121in}}{\pgfqpoint{3.086059in}{7.130714in}}%
\pgfpathquadraticcurveto{\pgfqpoint{3.092501in}{7.138328in}}{\pgfqpoint{3.103436in}{7.138328in}}%
\pgfpathlineto{\pgfqpoint{3.103436in}{7.138328in}}%
\pgfpathclose%
\pgfusepath{fill}%
\end{pgfscope}%
\begin{pgfscope}%
\pgfsetbuttcap%
\pgfsetroundjoin%
\definecolor{currentfill}{rgb}{0.090196,0.168627,0.227451}%
\pgfsetfillcolor{currentfill}%
\pgfsetlinewidth{0.000000pt}%
\definecolor{currentstroke}{rgb}{0.090196,0.168627,0.227451}%
\pgfsetstrokecolor{currentstroke}%
\pgfsetdash{}{0pt}%
\pgfpathmoveto{\pgfqpoint{1.251963in}{6.913499in}}%
\pgfpathlineto{\pgfqpoint{1.251963in}{6.899863in}}%
\pgfpathquadraticcurveto{\pgfqpoint{1.244016in}{6.903671in}}{\pgfqpoint{1.236955in}{6.905530in}}%
\pgfpathquadraticcurveto{\pgfqpoint{1.229894in}{6.907412in}}{\pgfqpoint{1.223320in}{6.907412in}}%
\pgfpathquadraticcurveto{\pgfqpoint{1.211920in}{6.907412in}}{\pgfqpoint{1.205722in}{6.902984in}}%
\pgfpathquadraticcurveto{\pgfqpoint{1.199524in}{6.898557in}}{\pgfqpoint{1.199524in}{6.890389in}}%
\pgfpathquadraticcurveto{\pgfqpoint{1.199524in}{6.883527in}}{\pgfqpoint{1.203641in}{6.880030in}}%
\pgfpathquadraticcurveto{\pgfqpoint{1.207758in}{6.876555in}}{\pgfqpoint{1.219247in}{6.874408in}}%
\pgfpathlineto{\pgfqpoint{1.227680in}{6.872681in}}%
\pgfpathquadraticcurveto{\pgfqpoint{1.243308in}{6.869693in}}{\pgfqpoint{1.250745in}{6.862189in}}%
\pgfpathquadraticcurveto{\pgfqpoint{1.258183in}{6.854685in}}{\pgfqpoint{1.258183in}{6.842112in}}%
\pgfpathquadraticcurveto{\pgfqpoint{1.258183in}{6.827082in}}{\pgfqpoint{1.248111in}{6.819335in}}%
\pgfpathquadraticcurveto{\pgfqpoint{1.238062in}{6.811587in}}{\pgfqpoint{1.218627in}{6.811587in}}%
\pgfpathquadraticcurveto{\pgfqpoint{1.211300in}{6.811587in}}{\pgfqpoint{1.203021in}{6.813248in}}%
\pgfpathquadraticcurveto{\pgfqpoint{1.194765in}{6.814908in}}{\pgfqpoint{1.185911in}{6.818162in}}%
\pgfpathlineto{\pgfqpoint{1.185911in}{6.832550in}}%
\pgfpathquadraticcurveto{\pgfqpoint{1.194411in}{6.827790in}}{\pgfqpoint{1.202579in}{6.825356in}}%
\pgfpathquadraticcurveto{\pgfqpoint{1.210747in}{6.822943in}}{\pgfqpoint{1.218627in}{6.822943in}}%
\pgfpathquadraticcurveto{\pgfqpoint{1.230580in}{6.822943in}}{\pgfqpoint{1.237088in}{6.827636in}}%
\pgfpathquadraticcurveto{\pgfqpoint{1.243596in}{6.832350in}}{\pgfqpoint{1.243596in}{6.841072in}}%
\pgfpathquadraticcurveto{\pgfqpoint{1.243596in}{6.848664in}}{\pgfqpoint{1.238925in}{6.852958in}}%
\pgfpathquadraticcurveto{\pgfqpoint{1.234255in}{6.857253in}}{\pgfqpoint{1.223607in}{6.859400in}}%
\pgfpathlineto{\pgfqpoint{1.215085in}{6.861060in}}%
\pgfpathquadraticcurveto{\pgfqpoint{1.199458in}{6.864159in}}{\pgfqpoint{1.192463in}{6.870800in}}%
\pgfpathquadraticcurveto{\pgfqpoint{1.185490in}{6.877440in}}{\pgfqpoint{1.185490in}{6.889283in}}%
\pgfpathquadraticcurveto{\pgfqpoint{1.185490in}{6.902984in}}{\pgfqpoint{1.195141in}{6.910865in}}%
\pgfpathquadraticcurveto{\pgfqpoint{1.204792in}{6.918745in}}{\pgfqpoint{1.221726in}{6.918745in}}%
\pgfpathquadraticcurveto{\pgfqpoint{1.229008in}{6.918745in}}{\pgfqpoint{1.236535in}{6.917417in}}%
\pgfpathquadraticcurveto{\pgfqpoint{1.244083in}{6.916111in}}{\pgfqpoint{1.251963in}{6.913499in}}%
\pgfpathlineto{\pgfqpoint{1.251963in}{6.913499in}}%
\pgfpathclose%
\pgfpathmoveto{\pgfqpoint{1.279433in}{6.891076in}}%
\pgfpathlineto{\pgfqpoint{1.292161in}{6.891076in}}%
\pgfpathlineto{\pgfqpoint{1.292161in}{6.813602in}}%
\pgfpathlineto{\pgfqpoint{1.279433in}{6.813602in}}%
\pgfpathlineto{\pgfqpoint{1.279433in}{6.891076in}}%
\pgfpathlineto{\pgfqpoint{1.279433in}{6.891076in}}%
\pgfpathclose%
\pgfpathmoveto{\pgfqpoint{1.279433in}{6.921246in}}%
\pgfpathlineto{\pgfqpoint{1.292161in}{6.921246in}}%
\pgfpathlineto{\pgfqpoint{1.292161in}{6.905109in}}%
\pgfpathlineto{\pgfqpoint{1.279433in}{6.905109in}}%
\pgfpathlineto{\pgfqpoint{1.279433in}{6.921246in}}%
\pgfpathlineto{\pgfqpoint{1.279433in}{6.921246in}}%
\pgfpathclose%
\pgfpathmoveto{\pgfqpoint{1.369768in}{6.853246in}}%
\pgfpathquadraticcurveto{\pgfqpoint{1.369768in}{6.867081in}}{\pgfqpoint{1.364057in}{6.874673in}}%
\pgfpathquadraticcurveto{\pgfqpoint{1.358368in}{6.882288in}}{\pgfqpoint{1.348053in}{6.882288in}}%
\pgfpathquadraticcurveto{\pgfqpoint{1.337826in}{6.882288in}}{\pgfqpoint{1.332115in}{6.874673in}}%
\pgfpathquadraticcurveto{\pgfqpoint{1.326404in}{6.867081in}}{\pgfqpoint{1.326404in}{6.853246in}}%
\pgfpathquadraticcurveto{\pgfqpoint{1.326404in}{6.839478in}}{\pgfqpoint{1.332115in}{6.831863in}}%
\pgfpathquadraticcurveto{\pgfqpoint{1.337826in}{6.824249in}}{\pgfqpoint{1.348053in}{6.824249in}}%
\pgfpathquadraticcurveto{\pgfqpoint{1.358368in}{6.824249in}}{\pgfqpoint{1.364057in}{6.831863in}}%
\pgfpathquadraticcurveto{\pgfqpoint{1.369768in}{6.839478in}}{\pgfqpoint{1.369768in}{6.853246in}}%
\pgfpathlineto{\pgfqpoint{1.369768in}{6.853246in}}%
\pgfpathclose%
\pgfpathmoveto{\pgfqpoint{1.382495in}{6.823208in}}%
\pgfpathquadraticcurveto{\pgfqpoint{1.382495in}{6.803442in}}{\pgfqpoint{1.373708in}{6.793790in}}%
\pgfpathquadraticcurveto{\pgfqpoint{1.364942in}{6.784139in}}{\pgfqpoint{1.346813in}{6.784139in}}%
\pgfpathquadraticcurveto{\pgfqpoint{1.340106in}{6.784139in}}{\pgfqpoint{1.334152in}{6.785136in}}%
\pgfpathquadraticcurveto{\pgfqpoint{1.328197in}{6.786132in}}{\pgfqpoint{1.322597in}{6.788212in}}%
\pgfpathlineto{\pgfqpoint{1.322597in}{6.800586in}}%
\pgfpathquadraticcurveto{\pgfqpoint{1.328197in}{6.797553in}}{\pgfqpoint{1.333665in}{6.796115in}}%
\pgfpathquadraticcurveto{\pgfqpoint{1.339132in}{6.794654in}}{\pgfqpoint{1.344799in}{6.794654in}}%
\pgfpathquadraticcurveto{\pgfqpoint{1.357327in}{6.794654in}}{\pgfqpoint{1.363548in}{6.801184in}}%
\pgfpathquadraticcurveto{\pgfqpoint{1.369768in}{6.807714in}}{\pgfqpoint{1.369768in}{6.820928in}}%
\pgfpathlineto{\pgfqpoint{1.369768in}{6.827237in}}%
\pgfpathquadraticcurveto{\pgfqpoint{1.365827in}{6.820375in}}{\pgfqpoint{1.359674in}{6.816988in}}%
\pgfpathquadraticcurveto{\pgfqpoint{1.353520in}{6.813602in}}{\pgfqpoint{1.344932in}{6.813602in}}%
\pgfpathquadraticcurveto{\pgfqpoint{1.330699in}{6.813602in}}{\pgfqpoint{1.321977in}{6.824448in}}%
\pgfpathquadraticcurveto{\pgfqpoint{1.313256in}{6.835317in}}{\pgfqpoint{1.313256in}{6.853246in}}%
\pgfpathquadraticcurveto{\pgfqpoint{1.313256in}{6.871220in}}{\pgfqpoint{1.321977in}{6.882067in}}%
\pgfpathquadraticcurveto{\pgfqpoint{1.330699in}{6.892935in}}{\pgfqpoint{1.344932in}{6.892935in}}%
\pgfpathquadraticcurveto{\pgfqpoint{1.353520in}{6.892935in}}{\pgfqpoint{1.359674in}{6.889548in}}%
\pgfpathquadraticcurveto{\pgfqpoint{1.365827in}{6.886162in}}{\pgfqpoint{1.369768in}{6.879322in}}%
\pgfpathlineto{\pgfqpoint{1.369768in}{6.891076in}}%
\pgfpathlineto{\pgfqpoint{1.382495in}{6.891076in}}%
\pgfpathlineto{\pgfqpoint{1.382495in}{6.823208in}}%
\pgfpathlineto{\pgfqpoint{1.382495in}{6.823208in}}%
\pgfpathclose%
\pgfpathmoveto{\pgfqpoint{1.473140in}{6.860374in}}%
\pgfpathlineto{\pgfqpoint{1.473140in}{6.813602in}}%
\pgfpathlineto{\pgfqpoint{1.460412in}{6.813602in}}%
\pgfpathlineto{\pgfqpoint{1.460412in}{6.859953in}}%
\pgfpathquadraticcurveto{\pgfqpoint{1.460412in}{6.870955in}}{\pgfqpoint{1.456118in}{6.876400in}}%
\pgfpathquadraticcurveto{\pgfqpoint{1.451824in}{6.881867in}}{\pgfqpoint{1.443257in}{6.881867in}}%
\pgfpathquadraticcurveto{\pgfqpoint{1.432942in}{6.881867in}}{\pgfqpoint{1.426988in}{6.875293in}}%
\pgfpathquadraticcurveto{\pgfqpoint{1.421033in}{6.868741in}}{\pgfqpoint{1.421033in}{6.857386in}}%
\pgfpathlineto{\pgfqpoint{1.421033in}{6.813602in}}%
\pgfpathlineto{\pgfqpoint{1.408239in}{6.813602in}}%
\pgfpathlineto{\pgfqpoint{1.408239in}{6.891076in}}%
\pgfpathlineto{\pgfqpoint{1.421033in}{6.891076in}}%
\pgfpathlineto{\pgfqpoint{1.421033in}{6.879034in}}%
\pgfpathquadraticcurveto{\pgfqpoint{1.425615in}{6.886029in}}{\pgfqpoint{1.431791in}{6.889482in}}%
\pgfpathquadraticcurveto{\pgfqpoint{1.437989in}{6.892935in}}{\pgfqpoint{1.446090in}{6.892935in}}%
\pgfpathquadraticcurveto{\pgfqpoint{1.459438in}{6.892935in}}{\pgfqpoint{1.466278in}{6.884678in}}%
\pgfpathquadraticcurveto{\pgfqpoint{1.473140in}{6.876422in}}{\pgfqpoint{1.473140in}{6.860374in}}%
\pgfpathlineto{\pgfqpoint{1.473140in}{6.860374in}}%
\pgfpathclose%
\pgfpathmoveto{\pgfqpoint{1.498507in}{6.891076in}}%
\pgfpathlineto{\pgfqpoint{1.511235in}{6.891076in}}%
\pgfpathlineto{\pgfqpoint{1.511235in}{6.813602in}}%
\pgfpathlineto{\pgfqpoint{1.498507in}{6.813602in}}%
\pgfpathlineto{\pgfqpoint{1.498507in}{6.891076in}}%
\pgfpathlineto{\pgfqpoint{1.498507in}{6.891076in}}%
\pgfpathclose%
\pgfpathmoveto{\pgfqpoint{1.498507in}{6.921246in}}%
\pgfpathlineto{\pgfqpoint{1.511235in}{6.921246in}}%
\pgfpathlineto{\pgfqpoint{1.511235in}{6.905109in}}%
\pgfpathlineto{\pgfqpoint{1.498507in}{6.905109in}}%
\pgfpathlineto{\pgfqpoint{1.498507in}{6.921246in}}%
\pgfpathlineto{\pgfqpoint{1.498507in}{6.921246in}}%
\pgfpathclose%
\pgfpathmoveto{\pgfqpoint{1.600463in}{6.891076in}}%
\pgfpathlineto{\pgfqpoint{1.600463in}{6.813602in}}%
\pgfpathlineto{\pgfqpoint{1.587669in}{6.813602in}}%
\pgfpathlineto{\pgfqpoint{1.587669in}{6.881181in}}%
\pgfpathlineto{\pgfqpoint{1.552739in}{6.881181in}}%
\pgfpathlineto{\pgfqpoint{1.552739in}{6.813602in}}%
\pgfpathlineto{\pgfqpoint{1.539945in}{6.813602in}}%
\pgfpathlineto{\pgfqpoint{1.539945in}{6.881181in}}%
\pgfpathlineto{\pgfqpoint{1.527770in}{6.881181in}}%
\pgfpathlineto{\pgfqpoint{1.527770in}{6.891076in}}%
\pgfpathlineto{\pgfqpoint{1.539945in}{6.891076in}}%
\pgfpathlineto{\pgfqpoint{1.539945in}{6.896477in}}%
\pgfpathquadraticcurveto{\pgfqpoint{1.539945in}{6.909138in}}{\pgfqpoint{1.545921in}{6.915181in}}%
\pgfpathquadraticcurveto{\pgfqpoint{1.551920in}{6.921246in}}{\pgfqpoint{1.564294in}{6.921246in}}%
\pgfpathlineto{\pgfqpoint{1.577088in}{6.921246in}}%
\pgfpathlineto{\pgfqpoint{1.577088in}{6.910643in}}%
\pgfpathlineto{\pgfqpoint{1.564913in}{6.910643in}}%
\pgfpathquadraticcurveto{\pgfqpoint{1.558074in}{6.910643in}}{\pgfqpoint{1.555395in}{6.907876in}}%
\pgfpathquadraticcurveto{\pgfqpoint{1.552739in}{6.905109in}}{\pgfqpoint{1.552739in}{6.897915in}}%
\pgfpathlineto{\pgfqpoint{1.552739in}{6.891076in}}%
\pgfpathlineto{\pgfqpoint{1.600463in}{6.891076in}}%
\pgfpathlineto{\pgfqpoint{1.600463in}{6.891076in}}%
\pgfpathclose%
\pgfpathmoveto{\pgfqpoint{1.587669in}{6.921091in}}%
\pgfpathlineto{\pgfqpoint{1.600463in}{6.921091in}}%
\pgfpathlineto{\pgfqpoint{1.600463in}{6.904977in}}%
\pgfpathlineto{\pgfqpoint{1.587669in}{6.904977in}}%
\pgfpathlineto{\pgfqpoint{1.587669in}{6.921091in}}%
\pgfpathlineto{\pgfqpoint{1.587669in}{6.921091in}}%
\pgfpathclose%
\pgfpathmoveto{\pgfqpoint{1.682851in}{6.888109in}}%
\pgfpathlineto{\pgfqpoint{1.682851in}{6.876201in}}%
\pgfpathquadraticcurveto{\pgfqpoint{1.677450in}{6.879189in}}{\pgfqpoint{1.672027in}{6.880672in}}%
\pgfpathquadraticcurveto{\pgfqpoint{1.666604in}{6.882155in}}{\pgfqpoint{1.661070in}{6.882155in}}%
\pgfpathquadraticcurveto{\pgfqpoint{1.648674in}{6.882155in}}{\pgfqpoint{1.641812in}{6.874297in}}%
\pgfpathquadraticcurveto{\pgfqpoint{1.634972in}{6.866461in}}{\pgfqpoint{1.634972in}{6.852272in}}%
\pgfpathquadraticcurveto{\pgfqpoint{1.634972in}{6.838083in}}{\pgfqpoint{1.641812in}{6.830225in}}%
\pgfpathquadraticcurveto{\pgfqpoint{1.648674in}{6.822389in}}{\pgfqpoint{1.661070in}{6.822389in}}%
\pgfpathquadraticcurveto{\pgfqpoint{1.666604in}{6.822389in}}{\pgfqpoint{1.672027in}{6.823873in}}%
\pgfpathquadraticcurveto{\pgfqpoint{1.677450in}{6.825356in}}{\pgfqpoint{1.682851in}{6.828344in}}%
\pgfpathlineto{\pgfqpoint{1.682851in}{6.816568in}}%
\pgfpathquadraticcurveto{\pgfqpoint{1.677516in}{6.814089in}}{\pgfqpoint{1.671805in}{6.812849in}}%
\pgfpathquadraticcurveto{\pgfqpoint{1.666117in}{6.811587in}}{\pgfqpoint{1.659675in}{6.811587in}}%
\pgfpathquadraticcurveto{\pgfqpoint{1.642166in}{6.811587in}}{\pgfqpoint{1.631851in}{6.822589in}}%
\pgfpathquadraticcurveto{\pgfqpoint{1.621558in}{6.833590in}}{\pgfqpoint{1.621558in}{6.852272in}}%
\pgfpathquadraticcurveto{\pgfqpoint{1.621558in}{6.871220in}}{\pgfqpoint{1.631962in}{6.882067in}}%
\pgfpathquadraticcurveto{\pgfqpoint{1.642387in}{6.892935in}}{\pgfqpoint{1.660516in}{6.892935in}}%
\pgfpathquadraticcurveto{\pgfqpoint{1.666382in}{6.892935in}}{\pgfqpoint{1.671982in}{6.891718in}}%
\pgfpathquadraticcurveto{\pgfqpoint{1.677583in}{6.890522in}}{\pgfqpoint{1.682851in}{6.888109in}}%
\pgfpathlineto{\pgfqpoint{1.682851in}{6.888109in}}%
\pgfpathclose%
\pgfpathmoveto{\pgfqpoint{1.740204in}{6.852538in}}%
\pgfpathquadraticcurveto{\pgfqpoint{1.724775in}{6.852538in}}{\pgfqpoint{1.718821in}{6.849018in}}%
\pgfpathquadraticcurveto{\pgfqpoint{1.712867in}{6.845499in}}{\pgfqpoint{1.712867in}{6.836977in}}%
\pgfpathquadraticcurveto{\pgfqpoint{1.712867in}{6.830203in}}{\pgfqpoint{1.717338in}{6.826219in}}%
\pgfpathquadraticcurveto{\pgfqpoint{1.721809in}{6.822257in}}{\pgfqpoint{1.729468in}{6.822257in}}%
\pgfpathquadraticcurveto{\pgfqpoint{1.740071in}{6.822257in}}{\pgfqpoint{1.746468in}{6.829761in}}%
\pgfpathquadraticcurveto{\pgfqpoint{1.752865in}{6.837264in}}{\pgfqpoint{1.752865in}{6.849705in}}%
\pgfpathlineto{\pgfqpoint{1.752865in}{6.852538in}}%
\pgfpathlineto{\pgfqpoint{1.740204in}{6.852538in}}%
\pgfpathlineto{\pgfqpoint{1.740204in}{6.852538in}}%
\pgfpathclose%
\pgfpathmoveto{\pgfqpoint{1.765593in}{6.857806in}}%
\pgfpathlineto{\pgfqpoint{1.765593in}{6.813602in}}%
\pgfpathlineto{\pgfqpoint{1.752865in}{6.813602in}}%
\pgfpathlineto{\pgfqpoint{1.752865in}{6.825356in}}%
\pgfpathquadraticcurveto{\pgfqpoint{1.748505in}{6.818317in}}{\pgfqpoint{1.741997in}{6.814952in}}%
\pgfpathquadraticcurveto{\pgfqpoint{1.735489in}{6.811587in}}{\pgfqpoint{1.726081in}{6.811587in}}%
\pgfpathquadraticcurveto{\pgfqpoint{1.714195in}{6.811587in}}{\pgfqpoint{1.707156in}{6.818272in}}%
\pgfpathquadraticcurveto{\pgfqpoint{1.700139in}{6.824957in}}{\pgfqpoint{1.700139in}{6.836158in}}%
\pgfpathquadraticcurveto{\pgfqpoint{1.700139in}{6.849218in}}{\pgfqpoint{1.708882in}{6.855858in}}%
\pgfpathquadraticcurveto{\pgfqpoint{1.717648in}{6.862499in}}{\pgfqpoint{1.735002in}{6.862499in}}%
\pgfpathlineto{\pgfqpoint{1.752865in}{6.862499in}}%
\pgfpathlineto{\pgfqpoint{1.752865in}{6.863761in}}%
\pgfpathquadraticcurveto{\pgfqpoint{1.752865in}{6.872548in}}{\pgfqpoint{1.747088in}{6.877352in}}%
\pgfpathquadraticcurveto{\pgfqpoint{1.741311in}{6.882155in}}{\pgfqpoint{1.730863in}{6.882155in}}%
\pgfpathquadraticcurveto{\pgfqpoint{1.724222in}{6.882155in}}{\pgfqpoint{1.717913in}{6.880561in}}%
\pgfpathquadraticcurveto{\pgfqpoint{1.711627in}{6.878968in}}{\pgfqpoint{1.705827in}{6.875780in}}%
\pgfpathlineto{\pgfqpoint{1.705827in}{6.887556in}}%
\pgfpathquadraticcurveto{\pgfqpoint{1.712800in}{6.890257in}}{\pgfqpoint{1.719374in}{6.891585in}}%
\pgfpathquadraticcurveto{\pgfqpoint{1.725949in}{6.892935in}}{\pgfqpoint{1.732169in}{6.892935in}}%
\pgfpathquadraticcurveto{\pgfqpoint{1.748992in}{6.892935in}}{\pgfqpoint{1.757292in}{6.884214in}}%
\pgfpathquadraticcurveto{\pgfqpoint{1.765593in}{6.875514in}}{\pgfqpoint{1.765593in}{6.857806in}}%
\pgfpathlineto{\pgfqpoint{1.765593in}{6.857806in}}%
\pgfpathclose%
\pgfpathmoveto{\pgfqpoint{1.856215in}{6.860374in}}%
\pgfpathlineto{\pgfqpoint{1.856215in}{6.813602in}}%
\pgfpathlineto{\pgfqpoint{1.843488in}{6.813602in}}%
\pgfpathlineto{\pgfqpoint{1.843488in}{6.859953in}}%
\pgfpathquadraticcurveto{\pgfqpoint{1.843488in}{6.870955in}}{\pgfqpoint{1.839193in}{6.876400in}}%
\pgfpathquadraticcurveto{\pgfqpoint{1.834899in}{6.881867in}}{\pgfqpoint{1.826333in}{6.881867in}}%
\pgfpathquadraticcurveto{\pgfqpoint{1.816018in}{6.881867in}}{\pgfqpoint{1.810063in}{6.875293in}}%
\pgfpathquadraticcurveto{\pgfqpoint{1.804109in}{6.868741in}}{\pgfqpoint{1.804109in}{6.857386in}}%
\pgfpathlineto{\pgfqpoint{1.804109in}{6.813602in}}%
\pgfpathlineto{\pgfqpoint{1.791314in}{6.813602in}}%
\pgfpathlineto{\pgfqpoint{1.791314in}{6.891076in}}%
\pgfpathlineto{\pgfqpoint{1.804109in}{6.891076in}}%
\pgfpathlineto{\pgfqpoint{1.804109in}{6.879034in}}%
\pgfpathquadraticcurveto{\pgfqpoint{1.808691in}{6.886029in}}{\pgfqpoint{1.814867in}{6.889482in}}%
\pgfpathquadraticcurveto{\pgfqpoint{1.821064in}{6.892935in}}{\pgfqpoint{1.829166in}{6.892935in}}%
\pgfpathquadraticcurveto{\pgfqpoint{1.842514in}{6.892935in}}{\pgfqpoint{1.849354in}{6.884678in}}%
\pgfpathquadraticcurveto{\pgfqpoint{1.856215in}{6.876422in}}{\pgfqpoint{1.856215in}{6.860374in}}%
\pgfpathlineto{\pgfqpoint{1.856215in}{6.860374in}}%
\pgfpathclose%
\pgfpathmoveto{\pgfqpoint{1.894178in}{6.913078in}}%
\pgfpathlineto{\pgfqpoint{1.894178in}{6.891076in}}%
\pgfpathlineto{\pgfqpoint{1.920386in}{6.891076in}}%
\pgfpathlineto{\pgfqpoint{1.920386in}{6.881181in}}%
\pgfpathlineto{\pgfqpoint{1.894178in}{6.881181in}}%
\pgfpathlineto{\pgfqpoint{1.894178in}{6.839124in}}%
\pgfpathquadraticcurveto{\pgfqpoint{1.894178in}{6.829650in}}{\pgfqpoint{1.896768in}{6.826949in}}%
\pgfpathquadraticcurveto{\pgfqpoint{1.899357in}{6.824249in}}{\pgfqpoint{1.907326in}{6.824249in}}%
\pgfpathlineto{\pgfqpoint{1.920386in}{6.824249in}}%
\pgfpathlineto{\pgfqpoint{1.920386in}{6.813602in}}%
\pgfpathlineto{\pgfqpoint{1.907326in}{6.813602in}}%
\pgfpathquadraticcurveto{\pgfqpoint{1.892584in}{6.813602in}}{\pgfqpoint{1.886984in}{6.819091in}}%
\pgfpathquadraticcurveto{\pgfqpoint{1.881383in}{6.824603in}}{\pgfqpoint{1.881383in}{6.839124in}}%
\pgfpathlineto{\pgfqpoint{1.881383in}{6.881181in}}%
\pgfpathlineto{\pgfqpoint{1.872042in}{6.881181in}}%
\pgfpathlineto{\pgfqpoint{1.872042in}{6.891076in}}%
\pgfpathlineto{\pgfqpoint{1.881383in}{6.891076in}}%
\pgfpathlineto{\pgfqpoint{1.881383in}{6.913078in}}%
\pgfpathlineto{\pgfqpoint{1.894178in}{6.913078in}}%
\pgfpathlineto{\pgfqpoint{1.894178in}{6.913078in}}%
\pgfpathclose%
\pgfpathmoveto{\pgfqpoint{2.033963in}{6.902431in}}%
\pgfpathlineto{\pgfqpoint{2.033963in}{6.863893in}}%
\pgfpathlineto{\pgfqpoint{2.072479in}{6.863893in}}%
\pgfpathlineto{\pgfqpoint{2.072479in}{6.852139in}}%
\pgfpathlineto{\pgfqpoint{2.033963in}{6.852139in}}%
\pgfpathlineto{\pgfqpoint{2.033963in}{6.813602in}}%
\pgfpathlineto{\pgfqpoint{2.022342in}{6.813602in}}%
\pgfpathlineto{\pgfqpoint{2.022342in}{6.852139in}}%
\pgfpathlineto{\pgfqpoint{1.983804in}{6.852139in}}%
\pgfpathlineto{\pgfqpoint{1.983804in}{6.863893in}}%
\pgfpathlineto{\pgfqpoint{2.022342in}{6.863893in}}%
\pgfpathlineto{\pgfqpoint{2.022342in}{6.902431in}}%
\pgfpathlineto{\pgfqpoint{2.033963in}{6.902431in}}%
\pgfpathlineto{\pgfqpoint{2.033963in}{6.902431in}}%
\pgfpathclose%
\pgfpathmoveto{\pgfqpoint{2.195419in}{6.862034in}}%
\pgfpathquadraticcurveto{\pgfqpoint{2.199912in}{6.860507in}}{\pgfqpoint{2.204162in}{6.855526in}}%
\pgfpathquadraticcurveto{\pgfqpoint{2.208412in}{6.850546in}}{\pgfqpoint{2.212706in}{6.841824in}}%
\pgfpathlineto{\pgfqpoint{2.226895in}{6.813602in}}%
\pgfpathlineto{\pgfqpoint{2.211865in}{6.813602in}}%
\pgfpathlineto{\pgfqpoint{2.198673in}{6.840098in}}%
\pgfpathquadraticcurveto{\pgfqpoint{2.193537in}{6.850479in}}{\pgfqpoint{2.188734in}{6.853866in}}%
\pgfpathquadraticcurveto{\pgfqpoint{2.183930in}{6.857253in}}{\pgfqpoint{2.175630in}{6.857253in}}%
\pgfpathlineto{\pgfqpoint{2.160400in}{6.857253in}}%
\pgfpathlineto{\pgfqpoint{2.160400in}{6.813602in}}%
\pgfpathlineto{\pgfqpoint{2.146433in}{6.813602in}}%
\pgfpathlineto{\pgfqpoint{2.146433in}{6.916886in}}%
\pgfpathlineto{\pgfqpoint{2.177976in}{6.916886in}}%
\pgfpathquadraticcurveto{\pgfqpoint{2.195684in}{6.916886in}}{\pgfqpoint{2.204406in}{6.909470in}}%
\pgfpathquadraticcurveto{\pgfqpoint{2.213127in}{6.902077in}}{\pgfqpoint{2.213127in}{6.887136in}}%
\pgfpathquadraticcurveto{\pgfqpoint{2.213127in}{6.877374in}}{\pgfqpoint{2.208589in}{6.870932in}}%
\pgfpathquadraticcurveto{\pgfqpoint{2.204051in}{6.864513in}}{\pgfqpoint{2.195419in}{6.862034in}}%
\pgfpathlineto{\pgfqpoint{2.195419in}{6.862034in}}%
\pgfpathclose%
\pgfpathmoveto{\pgfqpoint{2.160400in}{6.905397in}}%
\pgfpathlineto{\pgfqpoint{2.160400in}{6.868741in}}%
\pgfpathlineto{\pgfqpoint{2.177976in}{6.868741in}}%
\pgfpathquadraticcurveto{\pgfqpoint{2.188070in}{6.868741in}}{\pgfqpoint{2.193227in}{6.873412in}}%
\pgfpathquadraticcurveto{\pgfqpoint{2.198385in}{6.878082in}}{\pgfqpoint{2.198385in}{6.887136in}}%
\pgfpathquadraticcurveto{\pgfqpoint{2.198385in}{6.896189in}}{\pgfqpoint{2.193227in}{6.900793in}}%
\pgfpathquadraticcurveto{\pgfqpoint{2.188070in}{6.905397in}}{\pgfqpoint{2.177976in}{6.905397in}}%
\pgfpathlineto{\pgfqpoint{2.160400in}{6.905397in}}%
\pgfpathlineto{\pgfqpoint{2.160400in}{6.905397in}}%
\pgfpathclose%
\pgfpathmoveto{\pgfqpoint{2.304236in}{6.855526in}}%
\pgfpathlineto{\pgfqpoint{2.304236in}{6.849306in}}%
\pgfpathlineto{\pgfqpoint{2.245710in}{6.849306in}}%
\pgfpathquadraticcurveto{\pgfqpoint{2.246551in}{6.836158in}}{\pgfqpoint{2.253635in}{6.829274in}}%
\pgfpathquadraticcurveto{\pgfqpoint{2.260718in}{6.822389in}}{\pgfqpoint{2.273380in}{6.822389in}}%
\pgfpathquadraticcurveto{\pgfqpoint{2.280706in}{6.822389in}}{\pgfqpoint{2.287590in}{6.824182in}}%
\pgfpathquadraticcurveto{\pgfqpoint{2.294475in}{6.825975in}}{\pgfqpoint{2.301270in}{6.829583in}}%
\pgfpathlineto{\pgfqpoint{2.301270in}{6.817542in}}%
\pgfpathquadraticcurveto{\pgfqpoint{2.294408in}{6.814642in}}{\pgfqpoint{2.287214in}{6.813115in}}%
\pgfpathquadraticcurveto{\pgfqpoint{2.280020in}{6.811587in}}{\pgfqpoint{2.272627in}{6.811587in}}%
\pgfpathquadraticcurveto{\pgfqpoint{2.254077in}{6.811587in}}{\pgfqpoint{2.243253in}{6.822367in}}%
\pgfpathquadraticcurveto{\pgfqpoint{2.232429in}{6.833169in}}{\pgfqpoint{2.232429in}{6.851586in}}%
\pgfpathquadraticcurveto{\pgfqpoint{2.232429in}{6.870600in}}{\pgfqpoint{2.242700in}{6.881757in}}%
\pgfpathquadraticcurveto{\pgfqpoint{2.252971in}{6.892935in}}{\pgfqpoint{2.270413in}{6.892935in}}%
\pgfpathquadraticcurveto{\pgfqpoint{2.286041in}{6.892935in}}{\pgfqpoint{2.295139in}{6.882863in}}%
\pgfpathquadraticcurveto{\pgfqpoint{2.304236in}{6.872814in}}{\pgfqpoint{2.304236in}{6.855526in}}%
\pgfpathlineto{\pgfqpoint{2.304236in}{6.855526in}}%
\pgfpathclose%
\pgfpathmoveto{\pgfqpoint{2.291508in}{6.859267in}}%
\pgfpathquadraticcurveto{\pgfqpoint{2.291376in}{6.869693in}}{\pgfqpoint{2.285665in}{6.875913in}}%
\pgfpathquadraticcurveto{\pgfqpoint{2.279954in}{6.882155in}}{\pgfqpoint{2.270546in}{6.882155in}}%
\pgfpathquadraticcurveto{\pgfqpoint{2.259899in}{6.882155in}}{\pgfqpoint{2.253502in}{6.876134in}}%
\pgfpathquadraticcurveto{\pgfqpoint{2.247105in}{6.870113in}}{\pgfqpoint{2.246131in}{6.859178in}}%
\pgfpathlineto{\pgfqpoint{2.291508in}{6.859267in}}%
\pgfpathlineto{\pgfqpoint{2.291508in}{6.859267in}}%
\pgfpathclose%
\pgfpathmoveto{\pgfqpoint{2.316012in}{6.891076in}}%
\pgfpathlineto{\pgfqpoint{2.329493in}{6.891076in}}%
\pgfpathlineto{\pgfqpoint{2.353709in}{6.826064in}}%
\pgfpathlineto{\pgfqpoint{2.377925in}{6.891076in}}%
\pgfpathlineto{\pgfqpoint{2.391406in}{6.891076in}}%
\pgfpathlineto{\pgfqpoint{2.362342in}{6.813602in}}%
\pgfpathlineto{\pgfqpoint{2.345054in}{6.813602in}}%
\pgfpathlineto{\pgfqpoint{2.316012in}{6.891076in}}%
\pgfpathlineto{\pgfqpoint{2.316012in}{6.891076in}}%
\pgfpathclose%
\pgfpathmoveto{\pgfqpoint{2.475255in}{6.855526in}}%
\pgfpathlineto{\pgfqpoint{2.475255in}{6.849306in}}%
\pgfpathlineto{\pgfqpoint{2.416729in}{6.849306in}}%
\pgfpathquadraticcurveto{\pgfqpoint{2.417570in}{6.836158in}}{\pgfqpoint{2.424653in}{6.829274in}}%
\pgfpathquadraticcurveto{\pgfqpoint{2.431736in}{6.822389in}}{\pgfqpoint{2.444398in}{6.822389in}}%
\pgfpathquadraticcurveto{\pgfqpoint{2.451725in}{6.822389in}}{\pgfqpoint{2.458609in}{6.824182in}}%
\pgfpathquadraticcurveto{\pgfqpoint{2.465493in}{6.825975in}}{\pgfqpoint{2.472288in}{6.829583in}}%
\pgfpathlineto{\pgfqpoint{2.472288in}{6.817542in}}%
\pgfpathquadraticcurveto{\pgfqpoint{2.465426in}{6.814642in}}{\pgfqpoint{2.458232in}{6.813115in}}%
\pgfpathquadraticcurveto{\pgfqpoint{2.451038in}{6.811587in}}{\pgfqpoint{2.443645in}{6.811587in}}%
\pgfpathquadraticcurveto{\pgfqpoint{2.425096in}{6.811587in}}{\pgfqpoint{2.414271in}{6.822367in}}%
\pgfpathquadraticcurveto{\pgfqpoint{2.403447in}{6.833169in}}{\pgfqpoint{2.403447in}{6.851586in}}%
\pgfpathquadraticcurveto{\pgfqpoint{2.403447in}{6.870600in}}{\pgfqpoint{2.413718in}{6.881757in}}%
\pgfpathquadraticcurveto{\pgfqpoint{2.423989in}{6.892935in}}{\pgfqpoint{2.441432in}{6.892935in}}%
\pgfpathquadraticcurveto{\pgfqpoint{2.457059in}{6.892935in}}{\pgfqpoint{2.466157in}{6.882863in}}%
\pgfpathquadraticcurveto{\pgfqpoint{2.475255in}{6.872814in}}{\pgfqpoint{2.475255in}{6.855526in}}%
\pgfpathlineto{\pgfqpoint{2.475255in}{6.855526in}}%
\pgfpathclose%
\pgfpathmoveto{\pgfqpoint{2.462527in}{6.859267in}}%
\pgfpathquadraticcurveto{\pgfqpoint{2.462394in}{6.869693in}}{\pgfqpoint{2.456683in}{6.875913in}}%
\pgfpathquadraticcurveto{\pgfqpoint{2.450972in}{6.882155in}}{\pgfqpoint{2.441564in}{6.882155in}}%
\pgfpathquadraticcurveto{\pgfqpoint{2.430917in}{6.882155in}}{\pgfqpoint{2.424520in}{6.876134in}}%
\pgfpathquadraticcurveto{\pgfqpoint{2.418123in}{6.870113in}}{\pgfqpoint{2.417149in}{6.859178in}}%
\pgfpathlineto{\pgfqpoint{2.462527in}{6.859267in}}%
\pgfpathlineto{\pgfqpoint{2.462527in}{6.859267in}}%
\pgfpathclose%
\pgfpathmoveto{\pgfqpoint{2.541041in}{6.879189in}}%
\pgfpathquadraticcurveto{\pgfqpoint{2.538894in}{6.880428in}}{\pgfqpoint{2.536370in}{6.881004in}}%
\pgfpathquadraticcurveto{\pgfqpoint{2.533847in}{6.881602in}}{\pgfqpoint{2.530814in}{6.881602in}}%
\pgfpathquadraticcurveto{\pgfqpoint{2.520012in}{6.881602in}}{\pgfqpoint{2.514235in}{6.874585in}}%
\pgfpathquadraticcurveto{\pgfqpoint{2.508458in}{6.867568in}}{\pgfqpoint{2.508458in}{6.854419in}}%
\pgfpathlineto{\pgfqpoint{2.508458in}{6.813602in}}%
\pgfpathlineto{\pgfqpoint{2.495663in}{6.813602in}}%
\pgfpathlineto{\pgfqpoint{2.495663in}{6.891076in}}%
\pgfpathlineto{\pgfqpoint{2.508458in}{6.891076in}}%
\pgfpathlineto{\pgfqpoint{2.508458in}{6.879034in}}%
\pgfpathquadraticcurveto{\pgfqpoint{2.512486in}{6.886095in}}{\pgfqpoint{2.518906in}{6.889504in}}%
\pgfpathquadraticcurveto{\pgfqpoint{2.525347in}{6.892935in}}{\pgfqpoint{2.534555in}{6.892935in}}%
\pgfpathquadraticcurveto{\pgfqpoint{2.535861in}{6.892935in}}{\pgfqpoint{2.537455in}{6.892758in}}%
\pgfpathquadraticcurveto{\pgfqpoint{2.539049in}{6.892603in}}{\pgfqpoint{2.540975in}{6.892249in}}%
\pgfpathlineto{\pgfqpoint{2.541041in}{6.879189in}}%
\pgfpathlineto{\pgfqpoint{2.541041in}{6.879189in}}%
\pgfpathclose%
\pgfpathmoveto{\pgfqpoint{2.603773in}{6.888796in}}%
\pgfpathlineto{\pgfqpoint{2.603773in}{6.876754in}}%
\pgfpathquadraticcurveto{\pgfqpoint{2.598394in}{6.879521in}}{\pgfqpoint{2.592572in}{6.880893in}}%
\pgfpathquadraticcurveto{\pgfqpoint{2.586773in}{6.882288in}}{\pgfqpoint{2.580531in}{6.882288in}}%
\pgfpathquadraticcurveto{\pgfqpoint{2.571057in}{6.882288in}}{\pgfqpoint{2.566320in}{6.879388in}}%
\pgfpathquadraticcurveto{\pgfqpoint{2.561583in}{6.876488in}}{\pgfqpoint{2.561583in}{6.870667in}}%
\pgfpathquadraticcurveto{\pgfqpoint{2.561583in}{6.866240in}}{\pgfqpoint{2.564969in}{6.863716in}}%
\pgfpathquadraticcurveto{\pgfqpoint{2.568356in}{6.861193in}}{\pgfqpoint{2.578605in}{6.858913in}}%
\pgfpathlineto{\pgfqpoint{2.582965in}{6.857939in}}%
\pgfpathquadraticcurveto{\pgfqpoint{2.596512in}{6.855039in}}{\pgfqpoint{2.602223in}{6.849749in}}%
\pgfpathquadraticcurveto{\pgfqpoint{2.607934in}{6.844458in}}{\pgfqpoint{2.607934in}{6.834984in}}%
\pgfpathquadraticcurveto{\pgfqpoint{2.607934in}{6.824182in}}{\pgfqpoint{2.599390in}{6.817874in}}%
\pgfpathquadraticcurveto{\pgfqpoint{2.590846in}{6.811587in}}{\pgfqpoint{2.575904in}{6.811587in}}%
\pgfpathquadraticcurveto{\pgfqpoint{2.569684in}{6.811587in}}{\pgfqpoint{2.562933in}{6.812805in}}%
\pgfpathquadraticcurveto{\pgfqpoint{2.556182in}{6.814022in}}{\pgfqpoint{2.548722in}{6.816435in}}%
\pgfpathlineto{\pgfqpoint{2.548722in}{6.829583in}}%
\pgfpathquadraticcurveto{\pgfqpoint{2.555783in}{6.825909in}}{\pgfqpoint{2.562623in}{6.824072in}}%
\pgfpathquadraticcurveto{\pgfqpoint{2.569463in}{6.822257in}}{\pgfqpoint{2.576192in}{6.822257in}}%
\pgfpathquadraticcurveto{\pgfqpoint{2.585179in}{6.822257in}}{\pgfqpoint{2.590005in}{6.825333in}}%
\pgfpathquadraticcurveto{\pgfqpoint{2.594852in}{6.828410in}}{\pgfqpoint{2.594852in}{6.834011in}}%
\pgfpathquadraticcurveto{\pgfqpoint{2.594852in}{6.839190in}}{\pgfqpoint{2.591355in}{6.841957in}}%
\pgfpathquadraticcurveto{\pgfqpoint{2.587880in}{6.844724in}}{\pgfqpoint{2.576037in}{6.847292in}}%
\pgfpathlineto{\pgfqpoint{2.571610in}{6.848332in}}%
\pgfpathquadraticcurveto{\pgfqpoint{2.559790in}{6.850811in}}{\pgfqpoint{2.554521in}{6.855969in}}%
\pgfpathquadraticcurveto{\pgfqpoint{2.549275in}{6.861126in}}{\pgfqpoint{2.549275in}{6.870113in}}%
\pgfpathquadraticcurveto{\pgfqpoint{2.549275in}{6.881048in}}{\pgfqpoint{2.557023in}{6.886981in}}%
\pgfpathquadraticcurveto{\pgfqpoint{2.564770in}{6.892935in}}{\pgfqpoint{2.579025in}{6.892935in}}%
\pgfpathquadraticcurveto{\pgfqpoint{2.586064in}{6.892935in}}{\pgfqpoint{2.592285in}{6.891895in}}%
\pgfpathquadraticcurveto{\pgfqpoint{2.598527in}{6.890876in}}{\pgfqpoint{2.603773in}{6.888796in}}%
\pgfpathlineto{\pgfqpoint{2.603773in}{6.888796in}}%
\pgfpathclose%
\pgfpathmoveto{\pgfqpoint{2.663406in}{6.852538in}}%
\pgfpathquadraticcurveto{\pgfqpoint{2.647977in}{6.852538in}}{\pgfqpoint{2.642023in}{6.849018in}}%
\pgfpathquadraticcurveto{\pgfqpoint{2.636068in}{6.845499in}}{\pgfqpoint{2.636068in}{6.836977in}}%
\pgfpathquadraticcurveto{\pgfqpoint{2.636068in}{6.830203in}}{\pgfqpoint{2.640540in}{6.826219in}}%
\pgfpathquadraticcurveto{\pgfqpoint{2.645011in}{6.822257in}}{\pgfqpoint{2.652670in}{6.822257in}}%
\pgfpathquadraticcurveto{\pgfqpoint{2.663273in}{6.822257in}}{\pgfqpoint{2.669670in}{6.829761in}}%
\pgfpathquadraticcurveto{\pgfqpoint{2.676067in}{6.837264in}}{\pgfqpoint{2.676067in}{6.849705in}}%
\pgfpathlineto{\pgfqpoint{2.676067in}{6.852538in}}%
\pgfpathlineto{\pgfqpoint{2.663406in}{6.852538in}}%
\pgfpathlineto{\pgfqpoint{2.663406in}{6.852538in}}%
\pgfpathclose%
\pgfpathmoveto{\pgfqpoint{2.688795in}{6.857806in}}%
\pgfpathlineto{\pgfqpoint{2.688795in}{6.813602in}}%
\pgfpathlineto{\pgfqpoint{2.676067in}{6.813602in}}%
\pgfpathlineto{\pgfqpoint{2.676067in}{6.825356in}}%
\pgfpathquadraticcurveto{\pgfqpoint{2.671706in}{6.818317in}}{\pgfqpoint{2.665199in}{6.814952in}}%
\pgfpathquadraticcurveto{\pgfqpoint{2.658691in}{6.811587in}}{\pgfqpoint{2.649283in}{6.811587in}}%
\pgfpathquadraticcurveto{\pgfqpoint{2.637396in}{6.811587in}}{\pgfqpoint{2.630357in}{6.818272in}}%
\pgfpathquadraticcurveto{\pgfqpoint{2.623340in}{6.824957in}}{\pgfqpoint{2.623340in}{6.836158in}}%
\pgfpathquadraticcurveto{\pgfqpoint{2.623340in}{6.849218in}}{\pgfqpoint{2.632084in}{6.855858in}}%
\pgfpathquadraticcurveto{\pgfqpoint{2.640850in}{6.862499in}}{\pgfqpoint{2.658204in}{6.862499in}}%
\pgfpathlineto{\pgfqpoint{2.676067in}{6.862499in}}%
\pgfpathlineto{\pgfqpoint{2.676067in}{6.863761in}}%
\pgfpathquadraticcurveto{\pgfqpoint{2.676067in}{6.872548in}}{\pgfqpoint{2.670290in}{6.877352in}}%
\pgfpathquadraticcurveto{\pgfqpoint{2.664512in}{6.882155in}}{\pgfqpoint{2.654064in}{6.882155in}}%
\pgfpathquadraticcurveto{\pgfqpoint{2.647424in}{6.882155in}}{\pgfqpoint{2.641115in}{6.880561in}}%
\pgfpathquadraticcurveto{\pgfqpoint{2.634829in}{6.878968in}}{\pgfqpoint{2.629029in}{6.875780in}}%
\pgfpathlineto{\pgfqpoint{2.629029in}{6.887556in}}%
\pgfpathquadraticcurveto{\pgfqpoint{2.636002in}{6.890257in}}{\pgfqpoint{2.642576in}{6.891585in}}%
\pgfpathquadraticcurveto{\pgfqpoint{2.649150in}{6.892935in}}{\pgfqpoint{2.655370in}{6.892935in}}%
\pgfpathquadraticcurveto{\pgfqpoint{2.672193in}{6.892935in}}{\pgfqpoint{2.680494in}{6.884214in}}%
\pgfpathquadraticcurveto{\pgfqpoint{2.688795in}{6.875514in}}{\pgfqpoint{2.688795in}{6.857806in}}%
\pgfpathlineto{\pgfqpoint{2.688795in}{6.857806in}}%
\pgfpathclose%
\pgfpathmoveto{\pgfqpoint{2.715003in}{6.921246in}}%
\pgfpathlineto{\pgfqpoint{2.727731in}{6.921246in}}%
\pgfpathlineto{\pgfqpoint{2.727731in}{6.813602in}}%
\pgfpathlineto{\pgfqpoint{2.715003in}{6.813602in}}%
\pgfpathlineto{\pgfqpoint{2.715003in}{6.921246in}}%
\pgfpathlineto{\pgfqpoint{2.715003in}{6.921246in}}%
\pgfpathclose%
\pgfpathmoveto{\pgfqpoint{2.757614in}{6.831177in}}%
\pgfpathlineto{\pgfqpoint{2.772201in}{6.831177in}}%
\pgfpathlineto{\pgfqpoint{2.772201in}{6.813602in}}%
\pgfpathlineto{\pgfqpoint{2.757614in}{6.813602in}}%
\pgfpathlineto{\pgfqpoint{2.757614in}{6.831177in}}%
\pgfpathlineto{\pgfqpoint{2.757614in}{6.831177in}}%
\pgfpathclose%
\pgfpathmoveto{\pgfqpoint{2.757614in}{6.886848in}}%
\pgfpathlineto{\pgfqpoint{2.772201in}{6.886848in}}%
\pgfpathlineto{\pgfqpoint{2.772201in}{6.869294in}}%
\pgfpathlineto{\pgfqpoint{2.757614in}{6.869294in}}%
\pgfpathlineto{\pgfqpoint{2.757614in}{6.886848in}}%
\pgfpathlineto{\pgfqpoint{2.757614in}{6.886848in}}%
\pgfpathclose%
\pgfpathmoveto{\pgfqpoint{2.878783in}{6.907677in}}%
\pgfpathquadraticcurveto{\pgfqpoint{2.868003in}{6.907677in}}{\pgfqpoint{2.862558in}{6.897052in}}%
\pgfpathquadraticcurveto{\pgfqpoint{2.857135in}{6.886449in}}{\pgfqpoint{2.857135in}{6.865133in}}%
\pgfpathquadraticcurveto{\pgfqpoint{2.857135in}{6.843905in}}{\pgfqpoint{2.862558in}{6.833280in}}%
\pgfpathquadraticcurveto{\pgfqpoint{2.868003in}{6.822655in}}{\pgfqpoint{2.878783in}{6.822655in}}%
\pgfpathquadraticcurveto{\pgfqpoint{2.889652in}{6.822655in}}{\pgfqpoint{2.895075in}{6.833280in}}%
\pgfpathquadraticcurveto{\pgfqpoint{2.900520in}{6.843905in}}{\pgfqpoint{2.900520in}{6.865133in}}%
\pgfpathquadraticcurveto{\pgfqpoint{2.900520in}{6.886449in}}{\pgfqpoint{2.895075in}{6.897052in}}%
\pgfpathquadraticcurveto{\pgfqpoint{2.889652in}{6.907677in}}{\pgfqpoint{2.878783in}{6.907677in}}%
\pgfpathlineto{\pgfqpoint{2.878783in}{6.907677in}}%
\pgfpathclose%
\pgfpathmoveto{\pgfqpoint{2.878783in}{6.918745in}}%
\pgfpathquadraticcurveto{\pgfqpoint{2.896160in}{6.918745in}}{\pgfqpoint{2.905324in}{6.904999in}}%
\pgfpathquadraticcurveto{\pgfqpoint{2.914488in}{6.891275in}}{\pgfqpoint{2.914488in}{6.865133in}}%
\pgfpathquadraticcurveto{\pgfqpoint{2.914488in}{6.839057in}}{\pgfqpoint{2.905324in}{6.825311in}}%
\pgfpathquadraticcurveto{\pgfqpoint{2.896160in}{6.811587in}}{\pgfqpoint{2.878783in}{6.811587in}}%
\pgfpathquadraticcurveto{\pgfqpoint{2.861429in}{6.811587in}}{\pgfqpoint{2.852265in}{6.825311in}}%
\pgfpathquadraticcurveto{\pgfqpoint{2.843101in}{6.839057in}}{\pgfqpoint{2.843101in}{6.865133in}}%
\pgfpathquadraticcurveto{\pgfqpoint{2.843101in}{6.891275in}}{\pgfqpoint{2.852265in}{6.904999in}}%
\pgfpathquadraticcurveto{\pgfqpoint{2.861429in}{6.918745in}}{\pgfqpoint{2.878783in}{6.918745in}}%
\pgfpathlineto{\pgfqpoint{2.878783in}{6.918745in}}%
\pgfpathclose%
\pgfpathmoveto{\pgfqpoint{2.939036in}{6.831177in}}%
\pgfpathlineto{\pgfqpoint{2.953645in}{6.831177in}}%
\pgfpathlineto{\pgfqpoint{2.953645in}{6.813602in}}%
\pgfpathlineto{\pgfqpoint{2.939036in}{6.813602in}}%
\pgfpathlineto{\pgfqpoint{2.939036in}{6.831177in}}%
\pgfpathlineto{\pgfqpoint{2.939036in}{6.831177in}}%
\pgfpathclose%
\pgfpathmoveto{\pgfqpoint{3.013942in}{6.907677in}}%
\pgfpathquadraticcurveto{\pgfqpoint{3.003162in}{6.907677in}}{\pgfqpoint{2.997717in}{6.897052in}}%
\pgfpathquadraticcurveto{\pgfqpoint{2.992294in}{6.886449in}}{\pgfqpoint{2.992294in}{6.865133in}}%
\pgfpathquadraticcurveto{\pgfqpoint{2.992294in}{6.843905in}}{\pgfqpoint{2.997717in}{6.833280in}}%
\pgfpathquadraticcurveto{\pgfqpoint{3.003162in}{6.822655in}}{\pgfqpoint{3.013942in}{6.822655in}}%
\pgfpathquadraticcurveto{\pgfqpoint{3.024811in}{6.822655in}}{\pgfqpoint{3.030234in}{6.833280in}}%
\pgfpathquadraticcurveto{\pgfqpoint{3.035679in}{6.843905in}}{\pgfqpoint{3.035679in}{6.865133in}}%
\pgfpathquadraticcurveto{\pgfqpoint{3.035679in}{6.886449in}}{\pgfqpoint{3.030234in}{6.897052in}}%
\pgfpathquadraticcurveto{\pgfqpoint{3.024811in}{6.907677in}}{\pgfqpoint{3.013942in}{6.907677in}}%
\pgfpathlineto{\pgfqpoint{3.013942in}{6.907677in}}%
\pgfpathclose%
\pgfpathmoveto{\pgfqpoint{3.013942in}{6.918745in}}%
\pgfpathquadraticcurveto{\pgfqpoint{3.031318in}{6.918745in}}{\pgfqpoint{3.040482in}{6.904999in}}%
\pgfpathquadraticcurveto{\pgfqpoint{3.049646in}{6.891275in}}{\pgfqpoint{3.049646in}{6.865133in}}%
\pgfpathquadraticcurveto{\pgfqpoint{3.049646in}{6.839057in}}{\pgfqpoint{3.040482in}{6.825311in}}%
\pgfpathquadraticcurveto{\pgfqpoint{3.031318in}{6.811587in}}{\pgfqpoint{3.013942in}{6.811587in}}%
\pgfpathquadraticcurveto{\pgfqpoint{2.996588in}{6.811587in}}{\pgfqpoint{2.987424in}{6.825311in}}%
\pgfpathquadraticcurveto{\pgfqpoint{2.978260in}{6.839057in}}{\pgfqpoint{2.978260in}{6.865133in}}%
\pgfpathquadraticcurveto{\pgfqpoint{2.978260in}{6.891275in}}{\pgfqpoint{2.987424in}{6.904999in}}%
\pgfpathquadraticcurveto{\pgfqpoint{2.996588in}{6.918745in}}{\pgfqpoint{3.013942in}{6.918745in}}%
\pgfpathlineto{\pgfqpoint{3.013942in}{6.918745in}}%
\pgfpathclose%
\pgfpathmoveto{\pgfqpoint{3.162050in}{6.859046in}}%
\pgfpathquadraticcurveto{\pgfqpoint{3.156029in}{6.859046in}}{\pgfqpoint{3.152598in}{6.853932in}}%
\pgfpathquadraticcurveto{\pgfqpoint{3.149189in}{6.848819in}}{\pgfqpoint{3.149189in}{6.839677in}}%
\pgfpathquadraticcurveto{\pgfqpoint{3.149189in}{6.830690in}}{\pgfqpoint{3.152598in}{6.825533in}}%
\pgfpathquadraticcurveto{\pgfqpoint{3.156029in}{6.820375in}}{\pgfqpoint{3.162050in}{6.820375in}}%
\pgfpathquadraticcurveto{\pgfqpoint{3.167938in}{6.820375in}}{\pgfqpoint{3.171347in}{6.825533in}}%
\pgfpathquadraticcurveto{\pgfqpoint{3.174778in}{6.830690in}}{\pgfqpoint{3.174778in}{6.839677in}}%
\pgfpathquadraticcurveto{\pgfqpoint{3.174778in}{6.848753in}}{\pgfqpoint{3.171347in}{6.853888in}}%
\pgfpathquadraticcurveto{\pgfqpoint{3.167938in}{6.859046in}}{\pgfqpoint{3.162050in}{6.859046in}}%
\pgfpathlineto{\pgfqpoint{3.162050in}{6.859046in}}%
\pgfpathclose%
\pgfpathmoveto{\pgfqpoint{3.162050in}{6.867833in}}%
\pgfpathquadraticcurveto{\pgfqpoint{3.172985in}{6.867833in}}{\pgfqpoint{3.179404in}{6.860219in}}%
\pgfpathquadraticcurveto{\pgfqpoint{3.185846in}{6.852626in}}{\pgfqpoint{3.185846in}{6.839677in}}%
\pgfpathquadraticcurveto{\pgfqpoint{3.185846in}{6.826750in}}{\pgfqpoint{3.179382in}{6.819158in}}%
\pgfpathquadraticcurveto{\pgfqpoint{3.172919in}{6.811587in}}{\pgfqpoint{3.162050in}{6.811587in}}%
\pgfpathquadraticcurveto{\pgfqpoint{3.150982in}{6.811587in}}{\pgfqpoint{3.144541in}{6.819158in}}%
\pgfpathquadraticcurveto{\pgfqpoint{3.138122in}{6.826750in}}{\pgfqpoint{3.138122in}{6.839677in}}%
\pgfpathquadraticcurveto{\pgfqpoint{3.138122in}{6.852693in}}{\pgfqpoint{3.144585in}{6.860263in}}%
\pgfpathquadraticcurveto{\pgfqpoint{3.151049in}{6.867833in}}{\pgfqpoint{3.162050in}{6.867833in}}%
\pgfpathlineto{\pgfqpoint{3.162050in}{6.867833in}}%
\pgfpathclose%
\pgfpathmoveto{\pgfqpoint{3.090663in}{6.909957in}}%
\pgfpathquadraticcurveto{\pgfqpoint{3.084709in}{6.909957in}}{\pgfqpoint{3.081278in}{6.904800in}}%
\pgfpathquadraticcurveto{\pgfqpoint{3.077869in}{6.899664in}}{\pgfqpoint{3.077869in}{6.890655in}}%
\pgfpathquadraticcurveto{\pgfqpoint{3.077869in}{6.881535in}}{\pgfqpoint{3.081256in}{6.876400in}}%
\pgfpathquadraticcurveto{\pgfqpoint{3.084643in}{6.871287in}}{\pgfqpoint{3.090663in}{6.871287in}}%
\pgfpathquadraticcurveto{\pgfqpoint{3.096684in}{6.871287in}}{\pgfqpoint{3.100093in}{6.876400in}}%
\pgfpathquadraticcurveto{\pgfqpoint{3.103524in}{6.881535in}}{\pgfqpoint{3.103524in}{6.890655in}}%
\pgfpathquadraticcurveto{\pgfqpoint{3.103524in}{6.899576in}}{\pgfqpoint{3.100071in}{6.904755in}}%
\pgfpathquadraticcurveto{\pgfqpoint{3.096618in}{6.909957in}}{\pgfqpoint{3.090663in}{6.909957in}}%
\pgfpathlineto{\pgfqpoint{3.090663in}{6.909957in}}%
\pgfpathclose%
\pgfpathmoveto{\pgfqpoint{3.153130in}{6.918745in}}%
\pgfpathlineto{\pgfqpoint{3.164197in}{6.918745in}}%
\pgfpathlineto{\pgfqpoint{3.099584in}{6.811587in}}%
\pgfpathlineto{\pgfqpoint{3.088516in}{6.811587in}}%
\pgfpathlineto{\pgfqpoint{3.153130in}{6.918745in}}%
\pgfpathlineto{\pgfqpoint{3.153130in}{6.918745in}}%
\pgfpathclose%
\pgfpathmoveto{\pgfqpoint{3.090663in}{6.918745in}}%
\pgfpathquadraticcurveto{\pgfqpoint{3.101598in}{6.918745in}}{\pgfqpoint{3.108084in}{6.911175in}}%
\pgfpathquadraticcurveto{\pgfqpoint{3.114592in}{6.903604in}}{\pgfqpoint{3.114592in}{6.890655in}}%
\pgfpathquadraticcurveto{\pgfqpoint{3.114592in}{6.877595in}}{\pgfqpoint{3.108128in}{6.870047in}}%
\pgfpathquadraticcurveto{\pgfqpoint{3.101665in}{6.862499in}}{\pgfqpoint{3.090663in}{6.862499in}}%
\pgfpathquadraticcurveto{\pgfqpoint{3.079662in}{6.862499in}}{\pgfqpoint{3.073265in}{6.870069in}}%
\pgfpathquadraticcurveto{\pgfqpoint{3.066868in}{6.877662in}}{\pgfqpoint{3.066868in}{6.890655in}}%
\pgfpathquadraticcurveto{\pgfqpoint{3.066868in}{6.903538in}}{\pgfqpoint{3.073287in}{6.911130in}}%
\pgfpathquadraticcurveto{\pgfqpoint{3.079729in}{6.918745in}}{\pgfqpoint{3.090663in}{6.918745in}}%
\pgfpathlineto{\pgfqpoint{3.090663in}{6.918745in}}%
\pgfpathclose%
\pgfusepath{fill}%
\end{pgfscope}%
\begin{pgfscope}%
\pgfsetbuttcap%
\pgfsetmiterjoin%
\definecolor{currentfill}{rgb}{1.000000,1.000000,1.000000}%
\pgfsetfillcolor{currentfill}%
\pgfsetlinewidth{0.000000pt}%
\definecolor{currentstroke}{rgb}{0.000000,0.000000,0.000000}%
\pgfsetstrokecolor{currentstroke}%
\pgfsetstrokeopacity{0.000000}%
\pgfsetdash{}{0pt}%
\pgfpathmoveto{\pgfqpoint{7.904091in}{2.928000in}}%
\pgfpathlineto{\pgfqpoint{13.568091in}{2.928000in}}%
\pgfpathlineto{\pgfqpoint{13.568091in}{8.592000in}}%
\pgfpathlineto{\pgfqpoint{7.904091in}{8.592000in}}%
\pgfpathlineto{\pgfqpoint{7.904091in}{2.928000in}}%
\pgfpathclose%
\pgfusepath{fill}%
\end{pgfscope}%
\begin{pgfscope}%
\pgfpathrectangle{\pgfqpoint{7.904091in}{2.928000in}}{\pgfqpoint{5.664000in}{5.664000in}}%
\pgfusepath{clip}%
\pgfsetrectcap%
\pgfsetroundjoin%
\pgfsetlinewidth{0.602250pt}%
\definecolor{currentstroke}{rgb}{0.898039,0.913725,0.929412}%
\pgfsetstrokecolor{currentstroke}%
\pgfsetdash{}{0pt}%
\pgfpathmoveto{\pgfqpoint{7.904091in}{2.928000in}}%
\pgfpathlineto{\pgfqpoint{7.904091in}{8.592000in}}%
\pgfusepath{stroke}%
\end{pgfscope}%
\begin{pgfscope}%
\pgfsetbuttcap%
\pgfsetroundjoin%
\definecolor{currentfill}{rgb}{0.090196,0.168627,0.227451}%
\pgfsetfillcolor{currentfill}%
\pgfsetlinewidth{0.803000pt}%
\definecolor{currentstroke}{rgb}{0.090196,0.168627,0.227451}%
\pgfsetstrokecolor{currentstroke}%
\pgfsetdash{}{0pt}%
\pgfsys@defobject{currentmarker}{\pgfqpoint{0.000000in}{-0.048611in}}{\pgfqpoint{0.000000in}{0.000000in}}{%
\pgfpathmoveto{\pgfqpoint{0.000000in}{0.000000in}}%
\pgfpathlineto{\pgfqpoint{0.000000in}{-0.048611in}}%
\pgfusepath{stroke,fill}%
}%
\begin{pgfscope}%
\pgfsys@transformshift{7.904091in}{2.928000in}%
\pgfsys@useobject{currentmarker}{}%
\end{pgfscope}%
\end{pgfscope}%
\begin{pgfscope}%
\pgfsetbuttcap%
\pgfsetroundjoin%
\definecolor{currentfill}{rgb}{0.090196,0.168627,0.227451}%
\pgfsetfillcolor{currentfill}%
\pgfsetlinewidth{0.000000pt}%
\definecolor{currentstroke}{rgb}{0.090196,0.168627,0.227451}%
\pgfsetstrokecolor{currentstroke}%
\pgfsetdash{}{0pt}%
\pgfpathmoveto{\pgfqpoint{7.762333in}{2.780182in}}%
\pgfpathlineto{\pgfqpoint{7.840575in}{2.780182in}}%
\pgfpathlineto{\pgfqpoint{7.840575in}{2.769811in}}%
\pgfpathlineto{\pgfqpoint{7.762333in}{2.769811in}}%
\pgfpathlineto{\pgfqpoint{7.762333in}{2.780182in}}%
\pgfpathlineto{\pgfqpoint{7.762333in}{2.780182in}}%
\pgfpathclose%
\pgfpathmoveto{\pgfqpoint{7.867333in}{2.826940in}}%
\pgfpathlineto{\pgfqpoint{7.915732in}{2.826940in}}%
\pgfpathlineto{\pgfqpoint{7.915732in}{2.816549in}}%
\pgfpathlineto{\pgfqpoint{7.878622in}{2.816549in}}%
\pgfpathlineto{\pgfqpoint{7.878622in}{2.794225in}}%
\pgfpathquadraticcurveto{\pgfqpoint{7.881298in}{2.795143in}}{\pgfqpoint{7.883974in}{2.795592in}}%
\pgfpathquadraticcurveto{\pgfqpoint{7.886669in}{2.796041in}}{\pgfqpoint{7.889364in}{2.796041in}}%
\pgfpathquadraticcurveto{\pgfqpoint{7.904618in}{2.796041in}}{\pgfqpoint{7.913525in}{2.787682in}}%
\pgfpathquadraticcurveto{\pgfqpoint{7.922450in}{2.779323in}}{\pgfqpoint{7.922450in}{2.765045in}}%
\pgfpathquadraticcurveto{\pgfqpoint{7.922450in}{2.750338in}}{\pgfqpoint{7.913290in}{2.742174in}}%
\pgfpathquadraticcurveto{\pgfqpoint{7.904130in}{2.734030in}}{\pgfqpoint{7.887470in}{2.734030in}}%
\pgfpathquadraticcurveto{\pgfqpoint{7.881728in}{2.734030in}}{\pgfqpoint{7.875771in}{2.735006in}}%
\pgfpathquadraticcurveto{\pgfqpoint{7.869833in}{2.735983in}}{\pgfqpoint{7.863485in}{2.737936in}}%
\pgfpathlineto{\pgfqpoint{7.863485in}{2.750338in}}%
\pgfpathquadraticcurveto{\pgfqpoint{7.868974in}{2.747350in}}{\pgfqpoint{7.874833in}{2.745885in}}%
\pgfpathquadraticcurveto{\pgfqpoint{7.880692in}{2.744420in}}{\pgfqpoint{7.887216in}{2.744420in}}%
\pgfpathquadraticcurveto{\pgfqpoint{7.897782in}{2.744420in}}{\pgfqpoint{7.903935in}{2.749967in}}%
\pgfpathquadraticcurveto{\pgfqpoint{7.910107in}{2.755514in}}{\pgfqpoint{7.910107in}{2.765045in}}%
\pgfpathquadraticcurveto{\pgfqpoint{7.910107in}{2.774557in}}{\pgfqpoint{7.903935in}{2.780104in}}%
\pgfpathquadraticcurveto{\pgfqpoint{7.897782in}{2.785670in}}{\pgfqpoint{7.887216in}{2.785670in}}%
\pgfpathquadraticcurveto{\pgfqpoint{7.882275in}{2.785670in}}{\pgfqpoint{7.877353in}{2.784577in}}%
\pgfpathquadraticcurveto{\pgfqpoint{7.872450in}{2.783483in}}{\pgfqpoint{7.867333in}{2.781159in}}%
\pgfpathlineto{\pgfqpoint{7.867333in}{2.826940in}}%
\pgfpathlineto{\pgfqpoint{7.867333in}{2.826940in}}%
\pgfpathclose%
\pgfpathmoveto{\pgfqpoint{7.946728in}{2.751315in}}%
\pgfpathlineto{\pgfqpoint{7.959618in}{2.751315in}}%
\pgfpathlineto{\pgfqpoint{7.959618in}{2.735807in}}%
\pgfpathlineto{\pgfqpoint{7.946728in}{2.735807in}}%
\pgfpathlineto{\pgfqpoint{7.946728in}{2.751315in}}%
\pgfpathlineto{\pgfqpoint{7.946728in}{2.751315in}}%
\pgfpathclose%
\pgfpathmoveto{\pgfqpoint{8.012821in}{2.818815in}}%
\pgfpathquadraticcurveto{\pgfqpoint{8.003310in}{2.818815in}}{\pgfqpoint{7.998505in}{2.809440in}}%
\pgfpathquadraticcurveto{\pgfqpoint{7.993720in}{2.800084in}}{\pgfqpoint{7.993720in}{2.781276in}}%
\pgfpathquadraticcurveto{\pgfqpoint{7.993720in}{2.762545in}}{\pgfqpoint{7.998505in}{2.753170in}}%
\pgfpathquadraticcurveto{\pgfqpoint{8.003310in}{2.743795in}}{\pgfqpoint{8.012821in}{2.743795in}}%
\pgfpathquadraticcurveto{\pgfqpoint{8.022411in}{2.743795in}}{\pgfqpoint{8.027196in}{2.753170in}}%
\pgfpathquadraticcurveto{\pgfqpoint{8.032001in}{2.762545in}}{\pgfqpoint{8.032001in}{2.781276in}}%
\pgfpathquadraticcurveto{\pgfqpoint{8.032001in}{2.800084in}}{\pgfqpoint{8.027196in}{2.809440in}}%
\pgfpathquadraticcurveto{\pgfqpoint{8.022411in}{2.818815in}}{\pgfqpoint{8.012821in}{2.818815in}}%
\pgfpathlineto{\pgfqpoint{8.012821in}{2.818815in}}%
\pgfpathclose%
\pgfpathmoveto{\pgfqpoint{8.012821in}{2.828581in}}%
\pgfpathquadraticcurveto{\pgfqpoint{8.028153in}{2.828581in}}{\pgfqpoint{8.036239in}{2.816452in}}%
\pgfpathquadraticcurveto{\pgfqpoint{8.044325in}{2.804342in}}{\pgfqpoint{8.044325in}{2.781276in}}%
\pgfpathquadraticcurveto{\pgfqpoint{8.044325in}{2.758268in}}{\pgfqpoint{8.036239in}{2.746139in}}%
\pgfpathquadraticcurveto{\pgfqpoint{8.028153in}{2.734030in}}{\pgfqpoint{8.012821in}{2.734030in}}%
\pgfpathquadraticcurveto{\pgfqpoint{7.997509in}{2.734030in}}{\pgfqpoint{7.989423in}{2.746139in}}%
\pgfpathquadraticcurveto{\pgfqpoint{7.981337in}{2.758268in}}{\pgfqpoint{7.981337in}{2.781276in}}%
\pgfpathquadraticcurveto{\pgfqpoint{7.981337in}{2.804342in}}{\pgfqpoint{7.989423in}{2.816452in}}%
\pgfpathquadraticcurveto{\pgfqpoint{7.997509in}{2.828581in}}{\pgfqpoint{8.012821in}{2.828581in}}%
\pgfpathlineto{\pgfqpoint{8.012821in}{2.828581in}}%
\pgfpathclose%
\pgfusepath{fill}%
\end{pgfscope}%
\begin{pgfscope}%
\pgfpathrectangle{\pgfqpoint{7.904091in}{2.928000in}}{\pgfqpoint{5.664000in}{5.664000in}}%
\pgfusepath{clip}%
\pgfsetrectcap%
\pgfsetroundjoin%
\pgfsetlinewidth{0.602250pt}%
\definecolor{currentstroke}{rgb}{0.898039,0.913725,0.929412}%
\pgfsetstrokecolor{currentstroke}%
\pgfsetdash{}{0pt}%
\pgfpathmoveto{\pgfqpoint{9.320091in}{2.928000in}}%
\pgfpathlineto{\pgfqpoint{9.320091in}{8.592000in}}%
\pgfusepath{stroke}%
\end{pgfscope}%
\begin{pgfscope}%
\pgfsetbuttcap%
\pgfsetroundjoin%
\definecolor{currentfill}{rgb}{0.090196,0.168627,0.227451}%
\pgfsetfillcolor{currentfill}%
\pgfsetlinewidth{0.803000pt}%
\definecolor{currentstroke}{rgb}{0.090196,0.168627,0.227451}%
\pgfsetstrokecolor{currentstroke}%
\pgfsetdash{}{0pt}%
\pgfsys@defobject{currentmarker}{\pgfqpoint{0.000000in}{-0.048611in}}{\pgfqpoint{0.000000in}{0.000000in}}{%
\pgfpathmoveto{\pgfqpoint{0.000000in}{0.000000in}}%
\pgfpathlineto{\pgfqpoint{0.000000in}{-0.048611in}}%
\pgfusepath{stroke,fill}%
}%
\begin{pgfscope}%
\pgfsys@transformshift{9.320091in}{2.928000in}%
\pgfsys@useobject{currentmarker}{}%
\end{pgfscope}%
\end{pgfscope}%
\begin{pgfscope}%
\pgfsetbuttcap%
\pgfsetroundjoin%
\definecolor{currentfill}{rgb}{0.090196,0.168627,0.227451}%
\pgfsetfillcolor{currentfill}%
\pgfsetlinewidth{0.000000pt}%
\definecolor{currentstroke}{rgb}{0.090196,0.168627,0.227451}%
\pgfsetstrokecolor{currentstroke}%
\pgfsetdash{}{0pt}%
\pgfpathmoveto{\pgfqpoint{9.178333in}{2.780182in}}%
\pgfpathlineto{\pgfqpoint{9.256575in}{2.780182in}}%
\pgfpathlineto{\pgfqpoint{9.256575in}{2.769811in}}%
\pgfpathlineto{\pgfqpoint{9.178333in}{2.769811in}}%
\pgfpathlineto{\pgfqpoint{9.178333in}{2.780182in}}%
\pgfpathlineto{\pgfqpoint{9.178333in}{2.780182in}}%
\pgfpathclose%
\pgfpathmoveto{\pgfqpoint{9.293821in}{2.746178in}}%
\pgfpathlineto{\pgfqpoint{9.336849in}{2.746178in}}%
\pgfpathlineto{\pgfqpoint{9.336849in}{2.735807in}}%
\pgfpathlineto{\pgfqpoint{9.278997in}{2.735807in}}%
\pgfpathlineto{\pgfqpoint{9.278997in}{2.746178in}}%
\pgfpathquadraticcurveto{\pgfqpoint{9.286009in}{2.753444in}}{\pgfqpoint{9.298118in}{2.765670in}}%
\pgfpathquadraticcurveto{\pgfqpoint{9.310247in}{2.777916in}}{\pgfqpoint{9.313353in}{2.781471in}}%
\pgfpathquadraticcurveto{\pgfqpoint{9.319271in}{2.788112in}}{\pgfqpoint{9.321614in}{2.792721in}}%
\pgfpathquadraticcurveto{\pgfqpoint{9.323978in}{2.797331in}}{\pgfqpoint{9.323978in}{2.801784in}}%
\pgfpathquadraticcurveto{\pgfqpoint{9.323978in}{2.809049in}}{\pgfqpoint{9.318880in}{2.813620in}}%
\pgfpathquadraticcurveto{\pgfqpoint{9.313782in}{2.818209in}}{\pgfqpoint{9.305599in}{2.818209in}}%
\pgfpathquadraticcurveto{\pgfqpoint{9.299798in}{2.818209in}}{\pgfqpoint{9.293353in}{2.816198in}}%
\pgfpathquadraticcurveto{\pgfqpoint{9.286927in}{2.814186in}}{\pgfqpoint{9.279603in}{2.810084in}}%
\pgfpathlineto{\pgfqpoint{9.279603in}{2.822545in}}%
\pgfpathquadraticcurveto{\pgfqpoint{9.287044in}{2.825534in}}{\pgfqpoint{9.293509in}{2.827057in}}%
\pgfpathquadraticcurveto{\pgfqpoint{9.299993in}{2.828581in}}{\pgfqpoint{9.305364in}{2.828581in}}%
\pgfpathquadraticcurveto{\pgfqpoint{9.319525in}{2.828581in}}{\pgfqpoint{9.327942in}{2.821491in}}%
\pgfpathquadraticcurveto{\pgfqpoint{9.336360in}{2.814420in}}{\pgfqpoint{9.336360in}{2.802584in}}%
\pgfpathquadraticcurveto{\pgfqpoint{9.336360in}{2.796959in}}{\pgfqpoint{9.334251in}{2.791920in}}%
\pgfpathquadraticcurveto{\pgfqpoint{9.332161in}{2.786901in}}{\pgfqpoint{9.326595in}{2.780065in}}%
\pgfpathquadraticcurveto{\pgfqpoint{9.325071in}{2.778288in}}{\pgfqpoint{9.316888in}{2.769831in}}%
\pgfpathquadraticcurveto{\pgfqpoint{9.308724in}{2.761373in}}{\pgfqpoint{9.293821in}{2.746178in}}%
\pgfpathlineto{\pgfqpoint{9.293821in}{2.746178in}}%
\pgfpathclose%
\pgfpathmoveto{\pgfqpoint{9.362728in}{2.751315in}}%
\pgfpathlineto{\pgfqpoint{9.375618in}{2.751315in}}%
\pgfpathlineto{\pgfqpoint{9.375618in}{2.735807in}}%
\pgfpathlineto{\pgfqpoint{9.362728in}{2.735807in}}%
\pgfpathlineto{\pgfqpoint{9.362728in}{2.751315in}}%
\pgfpathlineto{\pgfqpoint{9.362728in}{2.751315in}}%
\pgfpathclose%
\pgfpathmoveto{\pgfqpoint{9.402591in}{2.826940in}}%
\pgfpathlineto{\pgfqpoint{9.450989in}{2.826940in}}%
\pgfpathlineto{\pgfqpoint{9.450989in}{2.816549in}}%
\pgfpathlineto{\pgfqpoint{9.413880in}{2.816549in}}%
\pgfpathlineto{\pgfqpoint{9.413880in}{2.794225in}}%
\pgfpathquadraticcurveto{\pgfqpoint{9.416556in}{2.795143in}}{\pgfqpoint{9.419232in}{2.795592in}}%
\pgfpathquadraticcurveto{\pgfqpoint{9.421927in}{2.796041in}}{\pgfqpoint{9.424622in}{2.796041in}}%
\pgfpathquadraticcurveto{\pgfqpoint{9.439876in}{2.796041in}}{\pgfqpoint{9.448782in}{2.787682in}}%
\pgfpathquadraticcurveto{\pgfqpoint{9.457708in}{2.779323in}}{\pgfqpoint{9.457708in}{2.765045in}}%
\pgfpathquadraticcurveto{\pgfqpoint{9.457708in}{2.750338in}}{\pgfqpoint{9.448548in}{2.742174in}}%
\pgfpathquadraticcurveto{\pgfqpoint{9.439388in}{2.734030in}}{\pgfqpoint{9.422728in}{2.734030in}}%
\pgfpathquadraticcurveto{\pgfqpoint{9.416985in}{2.734030in}}{\pgfqpoint{9.411028in}{2.735006in}}%
\pgfpathquadraticcurveto{\pgfqpoint{9.405091in}{2.735983in}}{\pgfqpoint{9.398743in}{2.737936in}}%
\pgfpathlineto{\pgfqpoint{9.398743in}{2.750338in}}%
\pgfpathquadraticcurveto{\pgfqpoint{9.404232in}{2.747350in}}{\pgfqpoint{9.410091in}{2.745885in}}%
\pgfpathquadraticcurveto{\pgfqpoint{9.415950in}{2.744420in}}{\pgfqpoint{9.422474in}{2.744420in}}%
\pgfpathquadraticcurveto{\pgfqpoint{9.433040in}{2.744420in}}{\pgfqpoint{9.439192in}{2.749967in}}%
\pgfpathquadraticcurveto{\pgfqpoint{9.445364in}{2.755514in}}{\pgfqpoint{9.445364in}{2.765045in}}%
\pgfpathquadraticcurveto{\pgfqpoint{9.445364in}{2.774557in}}{\pgfqpoint{9.439192in}{2.780104in}}%
\pgfpathquadraticcurveto{\pgfqpoint{9.433040in}{2.785670in}}{\pgfqpoint{9.422474in}{2.785670in}}%
\pgfpathquadraticcurveto{\pgfqpoint{9.417532in}{2.785670in}}{\pgfqpoint{9.412610in}{2.784577in}}%
\pgfpathquadraticcurveto{\pgfqpoint{9.407708in}{2.783483in}}{\pgfqpoint{9.402591in}{2.781159in}}%
\pgfpathlineto{\pgfqpoint{9.402591in}{2.826940in}}%
\pgfpathlineto{\pgfqpoint{9.402591in}{2.826940in}}%
\pgfpathclose%
\pgfusepath{fill}%
\end{pgfscope}%
\begin{pgfscope}%
\pgfpathrectangle{\pgfqpoint{7.904091in}{2.928000in}}{\pgfqpoint{5.664000in}{5.664000in}}%
\pgfusepath{clip}%
\pgfsetrectcap%
\pgfsetroundjoin%
\pgfsetlinewidth{0.602250pt}%
\definecolor{currentstroke}{rgb}{0.898039,0.913725,0.929412}%
\pgfsetstrokecolor{currentstroke}%
\pgfsetdash{}{0pt}%
\pgfpathmoveto{\pgfqpoint{10.736091in}{2.928000in}}%
\pgfpathlineto{\pgfqpoint{10.736091in}{8.592000in}}%
\pgfusepath{stroke}%
\end{pgfscope}%
\begin{pgfscope}%
\pgfsetbuttcap%
\pgfsetroundjoin%
\definecolor{currentfill}{rgb}{0.090196,0.168627,0.227451}%
\pgfsetfillcolor{currentfill}%
\pgfsetlinewidth{0.803000pt}%
\definecolor{currentstroke}{rgb}{0.090196,0.168627,0.227451}%
\pgfsetstrokecolor{currentstroke}%
\pgfsetdash{}{0pt}%
\pgfsys@defobject{currentmarker}{\pgfqpoint{0.000000in}{-0.048611in}}{\pgfqpoint{0.000000in}{0.000000in}}{%
\pgfpathmoveto{\pgfqpoint{0.000000in}{0.000000in}}%
\pgfpathlineto{\pgfqpoint{0.000000in}{-0.048611in}}%
\pgfusepath{stroke,fill}%
}%
\begin{pgfscope}%
\pgfsys@transformshift{10.736091in}{2.928000in}%
\pgfsys@useobject{currentmarker}{}%
\end{pgfscope}%
\end{pgfscope}%
\begin{pgfscope}%
\pgfsetbuttcap%
\pgfsetroundjoin%
\definecolor{currentfill}{rgb}{0.090196,0.168627,0.227451}%
\pgfsetfillcolor{currentfill}%
\pgfsetlinewidth{0.000000pt}%
\definecolor{currentstroke}{rgb}{0.090196,0.168627,0.227451}%
\pgfsetstrokecolor{currentstroke}%
\pgfsetdash{}{0pt}%
\pgfpathmoveto{\pgfqpoint{10.675817in}{2.818815in}}%
\pgfpathquadraticcurveto{\pgfqpoint{10.666306in}{2.818815in}}{\pgfqpoint{10.661501in}{2.809440in}}%
\pgfpathquadraticcurveto{\pgfqpoint{10.656716in}{2.800084in}}{\pgfqpoint{10.656716in}{2.781276in}}%
\pgfpathquadraticcurveto{\pgfqpoint{10.656716in}{2.762545in}}{\pgfqpoint{10.661501in}{2.753170in}}%
\pgfpathquadraticcurveto{\pgfqpoint{10.666306in}{2.743795in}}{\pgfqpoint{10.675817in}{2.743795in}}%
\pgfpathquadraticcurveto{\pgfqpoint{10.685407in}{2.743795in}}{\pgfqpoint{10.690192in}{2.753170in}}%
\pgfpathquadraticcurveto{\pgfqpoint{10.694997in}{2.762545in}}{\pgfqpoint{10.694997in}{2.781276in}}%
\pgfpathquadraticcurveto{\pgfqpoint{10.694997in}{2.800084in}}{\pgfqpoint{10.690192in}{2.809440in}}%
\pgfpathquadraticcurveto{\pgfqpoint{10.685407in}{2.818815in}}{\pgfqpoint{10.675817in}{2.818815in}}%
\pgfpathlineto{\pgfqpoint{10.675817in}{2.818815in}}%
\pgfpathclose%
\pgfpathmoveto{\pgfqpoint{10.675817in}{2.828581in}}%
\pgfpathquadraticcurveto{\pgfqpoint{10.691150in}{2.828581in}}{\pgfqpoint{10.699235in}{2.816452in}}%
\pgfpathquadraticcurveto{\pgfqpoint{10.707321in}{2.804342in}}{\pgfqpoint{10.707321in}{2.781276in}}%
\pgfpathquadraticcurveto{\pgfqpoint{10.707321in}{2.758268in}}{\pgfqpoint{10.699235in}{2.746139in}}%
\pgfpathquadraticcurveto{\pgfqpoint{10.691150in}{2.734030in}}{\pgfqpoint{10.675817in}{2.734030in}}%
\pgfpathquadraticcurveto{\pgfqpoint{10.660505in}{2.734030in}}{\pgfqpoint{10.652419in}{2.746139in}}%
\pgfpathquadraticcurveto{\pgfqpoint{10.644333in}{2.758268in}}{\pgfqpoint{10.644333in}{2.781276in}}%
\pgfpathquadraticcurveto{\pgfqpoint{10.644333in}{2.804342in}}{\pgfqpoint{10.652419in}{2.816452in}}%
\pgfpathquadraticcurveto{\pgfqpoint{10.660505in}{2.828581in}}{\pgfqpoint{10.675817in}{2.828581in}}%
\pgfpathlineto{\pgfqpoint{10.675817in}{2.828581in}}%
\pgfpathclose%
\pgfpathmoveto{\pgfqpoint{10.728982in}{2.751315in}}%
\pgfpathlineto{\pgfqpoint{10.741872in}{2.751315in}}%
\pgfpathlineto{\pgfqpoint{10.741872in}{2.735807in}}%
\pgfpathlineto{\pgfqpoint{10.728982in}{2.735807in}}%
\pgfpathlineto{\pgfqpoint{10.728982in}{2.751315in}}%
\pgfpathlineto{\pgfqpoint{10.728982in}{2.751315in}}%
\pgfpathclose%
\pgfpathmoveto{\pgfqpoint{10.795075in}{2.818815in}}%
\pgfpathquadraticcurveto{\pgfqpoint{10.785564in}{2.818815in}}{\pgfqpoint{10.780759in}{2.809440in}}%
\pgfpathquadraticcurveto{\pgfqpoint{10.775974in}{2.800084in}}{\pgfqpoint{10.775974in}{2.781276in}}%
\pgfpathquadraticcurveto{\pgfqpoint{10.775974in}{2.762545in}}{\pgfqpoint{10.780759in}{2.753170in}}%
\pgfpathquadraticcurveto{\pgfqpoint{10.785564in}{2.743795in}}{\pgfqpoint{10.795075in}{2.743795in}}%
\pgfpathquadraticcurveto{\pgfqpoint{10.804665in}{2.743795in}}{\pgfqpoint{10.809450in}{2.753170in}}%
\pgfpathquadraticcurveto{\pgfqpoint{10.814255in}{2.762545in}}{\pgfqpoint{10.814255in}{2.781276in}}%
\pgfpathquadraticcurveto{\pgfqpoint{10.814255in}{2.800084in}}{\pgfqpoint{10.809450in}{2.809440in}}%
\pgfpathquadraticcurveto{\pgfqpoint{10.804665in}{2.818815in}}{\pgfqpoint{10.795075in}{2.818815in}}%
\pgfpathlineto{\pgfqpoint{10.795075in}{2.818815in}}%
\pgfpathclose%
\pgfpathmoveto{\pgfqpoint{10.795075in}{2.828581in}}%
\pgfpathquadraticcurveto{\pgfqpoint{10.810407in}{2.828581in}}{\pgfqpoint{10.818493in}{2.816452in}}%
\pgfpathquadraticcurveto{\pgfqpoint{10.826579in}{2.804342in}}{\pgfqpoint{10.826579in}{2.781276in}}%
\pgfpathquadraticcurveto{\pgfqpoint{10.826579in}{2.758268in}}{\pgfqpoint{10.818493in}{2.746139in}}%
\pgfpathquadraticcurveto{\pgfqpoint{10.810407in}{2.734030in}}{\pgfqpoint{10.795075in}{2.734030in}}%
\pgfpathquadraticcurveto{\pgfqpoint{10.779763in}{2.734030in}}{\pgfqpoint{10.771677in}{2.746139in}}%
\pgfpathquadraticcurveto{\pgfqpoint{10.763591in}{2.758268in}}{\pgfqpoint{10.763591in}{2.781276in}}%
\pgfpathquadraticcurveto{\pgfqpoint{10.763591in}{2.804342in}}{\pgfqpoint{10.771677in}{2.816452in}}%
\pgfpathquadraticcurveto{\pgfqpoint{10.779763in}{2.828581in}}{\pgfqpoint{10.795075in}{2.828581in}}%
\pgfpathlineto{\pgfqpoint{10.795075in}{2.828581in}}%
\pgfpathclose%
\pgfusepath{fill}%
\end{pgfscope}%
\begin{pgfscope}%
\pgfpathrectangle{\pgfqpoint{7.904091in}{2.928000in}}{\pgfqpoint{5.664000in}{5.664000in}}%
\pgfusepath{clip}%
\pgfsetrectcap%
\pgfsetroundjoin%
\pgfsetlinewidth{0.602250pt}%
\definecolor{currentstroke}{rgb}{0.898039,0.913725,0.929412}%
\pgfsetstrokecolor{currentstroke}%
\pgfsetdash{}{0pt}%
\pgfpathmoveto{\pgfqpoint{12.152091in}{2.928000in}}%
\pgfpathlineto{\pgfqpoint{12.152091in}{8.592000in}}%
\pgfusepath{stroke}%
\end{pgfscope}%
\begin{pgfscope}%
\pgfsetbuttcap%
\pgfsetroundjoin%
\definecolor{currentfill}{rgb}{0.090196,0.168627,0.227451}%
\pgfsetfillcolor{currentfill}%
\pgfsetlinewidth{0.803000pt}%
\definecolor{currentstroke}{rgb}{0.090196,0.168627,0.227451}%
\pgfsetstrokecolor{currentstroke}%
\pgfsetdash{}{0pt}%
\pgfsys@defobject{currentmarker}{\pgfqpoint{0.000000in}{-0.048611in}}{\pgfqpoint{0.000000in}{0.000000in}}{%
\pgfpathmoveto{\pgfqpoint{0.000000in}{0.000000in}}%
\pgfpathlineto{\pgfqpoint{0.000000in}{-0.048611in}}%
\pgfusepath{stroke,fill}%
}%
\begin{pgfscope}%
\pgfsys@transformshift{12.152091in}{2.928000in}%
\pgfsys@useobject{currentmarker}{}%
\end{pgfscope}%
\end{pgfscope}%
\begin{pgfscope}%
\pgfsetbuttcap%
\pgfsetroundjoin%
\definecolor{currentfill}{rgb}{0.090196,0.168627,0.227451}%
\pgfsetfillcolor{currentfill}%
\pgfsetlinewidth{0.000000pt}%
\definecolor{currentstroke}{rgb}{0.090196,0.168627,0.227451}%
\pgfsetstrokecolor{currentstroke}%
\pgfsetdash{}{0pt}%
\pgfpathmoveto{\pgfqpoint{12.076075in}{2.746178in}}%
\pgfpathlineto{\pgfqpoint{12.119103in}{2.746178in}}%
\pgfpathlineto{\pgfqpoint{12.119103in}{2.735807in}}%
\pgfpathlineto{\pgfqpoint{12.061251in}{2.735807in}}%
\pgfpathlineto{\pgfqpoint{12.061251in}{2.746178in}}%
\pgfpathquadraticcurveto{\pgfqpoint{12.068263in}{2.753444in}}{\pgfqpoint{12.080372in}{2.765670in}}%
\pgfpathquadraticcurveto{\pgfqpoint{12.092501in}{2.777916in}}{\pgfqpoint{12.095607in}{2.781471in}}%
\pgfpathquadraticcurveto{\pgfqpoint{12.101525in}{2.788112in}}{\pgfqpoint{12.103868in}{2.792721in}}%
\pgfpathquadraticcurveto{\pgfqpoint{12.106232in}{2.797331in}}{\pgfqpoint{12.106232in}{2.801784in}}%
\pgfpathquadraticcurveto{\pgfqpoint{12.106232in}{2.809049in}}{\pgfqpoint{12.101134in}{2.813620in}}%
\pgfpathquadraticcurveto{\pgfqpoint{12.096036in}{2.818209in}}{\pgfqpoint{12.087853in}{2.818209in}}%
\pgfpathquadraticcurveto{\pgfqpoint{12.082052in}{2.818209in}}{\pgfqpoint{12.075607in}{2.816198in}}%
\pgfpathquadraticcurveto{\pgfqpoint{12.069181in}{2.814186in}}{\pgfqpoint{12.061857in}{2.810084in}}%
\pgfpathlineto{\pgfqpoint{12.061857in}{2.822545in}}%
\pgfpathquadraticcurveto{\pgfqpoint{12.069298in}{2.825534in}}{\pgfqpoint{12.075763in}{2.827057in}}%
\pgfpathquadraticcurveto{\pgfqpoint{12.082247in}{2.828581in}}{\pgfqpoint{12.087618in}{2.828581in}}%
\pgfpathquadraticcurveto{\pgfqpoint{12.101778in}{2.828581in}}{\pgfqpoint{12.110196in}{2.821491in}}%
\pgfpathquadraticcurveto{\pgfqpoint{12.118614in}{2.814420in}}{\pgfqpoint{12.118614in}{2.802584in}}%
\pgfpathquadraticcurveto{\pgfqpoint{12.118614in}{2.796959in}}{\pgfqpoint{12.116505in}{2.791920in}}%
\pgfpathquadraticcurveto{\pgfqpoint{12.114415in}{2.786901in}}{\pgfqpoint{12.108849in}{2.780065in}}%
\pgfpathquadraticcurveto{\pgfqpoint{12.107325in}{2.778288in}}{\pgfqpoint{12.099142in}{2.769831in}}%
\pgfpathquadraticcurveto{\pgfqpoint{12.090978in}{2.761373in}}{\pgfqpoint{12.076075in}{2.746178in}}%
\pgfpathlineto{\pgfqpoint{12.076075in}{2.746178in}}%
\pgfpathclose%
\pgfpathmoveto{\pgfqpoint{12.144982in}{2.751315in}}%
\pgfpathlineto{\pgfqpoint{12.157872in}{2.751315in}}%
\pgfpathlineto{\pgfqpoint{12.157872in}{2.735807in}}%
\pgfpathlineto{\pgfqpoint{12.144982in}{2.735807in}}%
\pgfpathlineto{\pgfqpoint{12.144982in}{2.751315in}}%
\pgfpathlineto{\pgfqpoint{12.144982in}{2.751315in}}%
\pgfpathclose%
\pgfpathmoveto{\pgfqpoint{12.184845in}{2.826940in}}%
\pgfpathlineto{\pgfqpoint{12.233243in}{2.826940in}}%
\pgfpathlineto{\pgfqpoint{12.233243in}{2.816549in}}%
\pgfpathlineto{\pgfqpoint{12.196134in}{2.816549in}}%
\pgfpathlineto{\pgfqpoint{12.196134in}{2.794225in}}%
\pgfpathquadraticcurveto{\pgfqpoint{12.198810in}{2.795143in}}{\pgfqpoint{12.201485in}{2.795592in}}%
\pgfpathquadraticcurveto{\pgfqpoint{12.204181in}{2.796041in}}{\pgfqpoint{12.206876in}{2.796041in}}%
\pgfpathquadraticcurveto{\pgfqpoint{12.222130in}{2.796041in}}{\pgfqpoint{12.231036in}{2.787682in}}%
\pgfpathquadraticcurveto{\pgfqpoint{12.239962in}{2.779323in}}{\pgfqpoint{12.239962in}{2.765045in}}%
\pgfpathquadraticcurveto{\pgfqpoint{12.239962in}{2.750338in}}{\pgfqpoint{12.230802in}{2.742174in}}%
\pgfpathquadraticcurveto{\pgfqpoint{12.221642in}{2.734030in}}{\pgfqpoint{12.204982in}{2.734030in}}%
\pgfpathquadraticcurveto{\pgfqpoint{12.199239in}{2.734030in}}{\pgfqpoint{12.193282in}{2.735006in}}%
\pgfpathquadraticcurveto{\pgfqpoint{12.187345in}{2.735983in}}{\pgfqpoint{12.180997in}{2.737936in}}%
\pgfpathlineto{\pgfqpoint{12.180997in}{2.750338in}}%
\pgfpathquadraticcurveto{\pgfqpoint{12.186485in}{2.747350in}}{\pgfqpoint{12.192345in}{2.745885in}}%
\pgfpathquadraticcurveto{\pgfqpoint{12.198204in}{2.744420in}}{\pgfqpoint{12.204728in}{2.744420in}}%
\pgfpathquadraticcurveto{\pgfqpoint{12.215294in}{2.744420in}}{\pgfqpoint{12.221446in}{2.749967in}}%
\pgfpathquadraticcurveto{\pgfqpoint{12.227618in}{2.755514in}}{\pgfqpoint{12.227618in}{2.765045in}}%
\pgfpathquadraticcurveto{\pgfqpoint{12.227618in}{2.774557in}}{\pgfqpoint{12.221446in}{2.780104in}}%
\pgfpathquadraticcurveto{\pgfqpoint{12.215294in}{2.785670in}}{\pgfqpoint{12.204728in}{2.785670in}}%
\pgfpathquadraticcurveto{\pgfqpoint{12.199786in}{2.785670in}}{\pgfqpoint{12.194864in}{2.784577in}}%
\pgfpathquadraticcurveto{\pgfqpoint{12.189962in}{2.783483in}}{\pgfqpoint{12.184845in}{2.781159in}}%
\pgfpathlineto{\pgfqpoint{12.184845in}{2.826940in}}%
\pgfpathlineto{\pgfqpoint{12.184845in}{2.826940in}}%
\pgfpathclose%
\pgfusepath{fill}%
\end{pgfscope}%
\begin{pgfscope}%
\pgfpathrectangle{\pgfqpoint{7.904091in}{2.928000in}}{\pgfqpoint{5.664000in}{5.664000in}}%
\pgfusepath{clip}%
\pgfsetrectcap%
\pgfsetroundjoin%
\pgfsetlinewidth{0.602250pt}%
\definecolor{currentstroke}{rgb}{0.898039,0.913725,0.929412}%
\pgfsetstrokecolor{currentstroke}%
\pgfsetdash{}{0pt}%
\pgfpathmoveto{\pgfqpoint{13.568091in}{2.928000in}}%
\pgfpathlineto{\pgfqpoint{13.568091in}{8.592000in}}%
\pgfusepath{stroke}%
\end{pgfscope}%
\begin{pgfscope}%
\pgfsetbuttcap%
\pgfsetroundjoin%
\definecolor{currentfill}{rgb}{0.090196,0.168627,0.227451}%
\pgfsetfillcolor{currentfill}%
\pgfsetlinewidth{0.803000pt}%
\definecolor{currentstroke}{rgb}{0.090196,0.168627,0.227451}%
\pgfsetstrokecolor{currentstroke}%
\pgfsetdash{}{0pt}%
\pgfsys@defobject{currentmarker}{\pgfqpoint{0.000000in}{-0.048611in}}{\pgfqpoint{0.000000in}{0.000000in}}{%
\pgfpathmoveto{\pgfqpoint{0.000000in}{0.000000in}}%
\pgfpathlineto{\pgfqpoint{0.000000in}{-0.048611in}}%
\pgfusepath{stroke,fill}%
}%
\begin{pgfscope}%
\pgfsys@transformshift{13.568091in}{2.928000in}%
\pgfsys@useobject{currentmarker}{}%
\end{pgfscope}%
\end{pgfscope}%
\begin{pgfscope}%
\pgfsetbuttcap%
\pgfsetroundjoin%
\definecolor{currentfill}{rgb}{0.090196,0.168627,0.227451}%
\pgfsetfillcolor{currentfill}%
\pgfsetlinewidth{0.000000pt}%
\definecolor{currentstroke}{rgb}{0.090196,0.168627,0.227451}%
\pgfsetstrokecolor{currentstroke}%
\pgfsetdash{}{0pt}%
\pgfpathmoveto{\pgfqpoint{13.481587in}{2.826940in}}%
\pgfpathlineto{\pgfqpoint{13.529985in}{2.826940in}}%
\pgfpathlineto{\pgfqpoint{13.529985in}{2.816549in}}%
\pgfpathlineto{\pgfqpoint{13.492876in}{2.816549in}}%
\pgfpathlineto{\pgfqpoint{13.492876in}{2.794225in}}%
\pgfpathquadraticcurveto{\pgfqpoint{13.495552in}{2.795143in}}{\pgfqpoint{13.498228in}{2.795592in}}%
\pgfpathquadraticcurveto{\pgfqpoint{13.500923in}{2.796041in}}{\pgfqpoint{13.503618in}{2.796041in}}%
\pgfpathquadraticcurveto{\pgfqpoint{13.518872in}{2.796041in}}{\pgfqpoint{13.527778in}{2.787682in}}%
\pgfpathquadraticcurveto{\pgfqpoint{13.536704in}{2.779323in}}{\pgfqpoint{13.536704in}{2.765045in}}%
\pgfpathquadraticcurveto{\pgfqpoint{13.536704in}{2.750338in}}{\pgfqpoint{13.527544in}{2.742174in}}%
\pgfpathquadraticcurveto{\pgfqpoint{13.518384in}{2.734030in}}{\pgfqpoint{13.501724in}{2.734030in}}%
\pgfpathquadraticcurveto{\pgfqpoint{13.495982in}{2.734030in}}{\pgfqpoint{13.490025in}{2.735006in}}%
\pgfpathquadraticcurveto{\pgfqpoint{13.484087in}{2.735983in}}{\pgfqpoint{13.477739in}{2.737936in}}%
\pgfpathlineto{\pgfqpoint{13.477739in}{2.750338in}}%
\pgfpathquadraticcurveto{\pgfqpoint{13.483228in}{2.747350in}}{\pgfqpoint{13.489087in}{2.745885in}}%
\pgfpathquadraticcurveto{\pgfqpoint{13.494946in}{2.744420in}}{\pgfqpoint{13.501470in}{2.744420in}}%
\pgfpathquadraticcurveto{\pgfqpoint{13.512036in}{2.744420in}}{\pgfqpoint{13.518189in}{2.749967in}}%
\pgfpathquadraticcurveto{\pgfqpoint{13.524360in}{2.755514in}}{\pgfqpoint{13.524360in}{2.765045in}}%
\pgfpathquadraticcurveto{\pgfqpoint{13.524360in}{2.774557in}}{\pgfqpoint{13.518189in}{2.780104in}}%
\pgfpathquadraticcurveto{\pgfqpoint{13.512036in}{2.785670in}}{\pgfqpoint{13.501470in}{2.785670in}}%
\pgfpathquadraticcurveto{\pgfqpoint{13.496528in}{2.785670in}}{\pgfqpoint{13.491607in}{2.784577in}}%
\pgfpathquadraticcurveto{\pgfqpoint{13.486704in}{2.783483in}}{\pgfqpoint{13.481587in}{2.781159in}}%
\pgfpathlineto{\pgfqpoint{13.481587in}{2.826940in}}%
\pgfpathlineto{\pgfqpoint{13.481587in}{2.826940in}}%
\pgfpathclose%
\pgfpathmoveto{\pgfqpoint{13.560982in}{2.751315in}}%
\pgfpathlineto{\pgfqpoint{13.573872in}{2.751315in}}%
\pgfpathlineto{\pgfqpoint{13.573872in}{2.735807in}}%
\pgfpathlineto{\pgfqpoint{13.560982in}{2.735807in}}%
\pgfpathlineto{\pgfqpoint{13.560982in}{2.751315in}}%
\pgfpathlineto{\pgfqpoint{13.560982in}{2.751315in}}%
\pgfpathclose%
\pgfpathmoveto{\pgfqpoint{13.627075in}{2.818815in}}%
\pgfpathquadraticcurveto{\pgfqpoint{13.617564in}{2.818815in}}{\pgfqpoint{13.612759in}{2.809440in}}%
\pgfpathquadraticcurveto{\pgfqpoint{13.607974in}{2.800084in}}{\pgfqpoint{13.607974in}{2.781276in}}%
\pgfpathquadraticcurveto{\pgfqpoint{13.607974in}{2.762545in}}{\pgfqpoint{13.612759in}{2.753170in}}%
\pgfpathquadraticcurveto{\pgfqpoint{13.617564in}{2.743795in}}{\pgfqpoint{13.627075in}{2.743795in}}%
\pgfpathquadraticcurveto{\pgfqpoint{13.636665in}{2.743795in}}{\pgfqpoint{13.641450in}{2.753170in}}%
\pgfpathquadraticcurveto{\pgfqpoint{13.646255in}{2.762545in}}{\pgfqpoint{13.646255in}{2.781276in}}%
\pgfpathquadraticcurveto{\pgfqpoint{13.646255in}{2.800084in}}{\pgfqpoint{13.641450in}{2.809440in}}%
\pgfpathquadraticcurveto{\pgfqpoint{13.636665in}{2.818815in}}{\pgfqpoint{13.627075in}{2.818815in}}%
\pgfpathlineto{\pgfqpoint{13.627075in}{2.818815in}}%
\pgfpathclose%
\pgfpathmoveto{\pgfqpoint{13.627075in}{2.828581in}}%
\pgfpathquadraticcurveto{\pgfqpoint{13.642407in}{2.828581in}}{\pgfqpoint{13.650493in}{2.816452in}}%
\pgfpathquadraticcurveto{\pgfqpoint{13.658579in}{2.804342in}}{\pgfqpoint{13.658579in}{2.781276in}}%
\pgfpathquadraticcurveto{\pgfqpoint{13.658579in}{2.758268in}}{\pgfqpoint{13.650493in}{2.746139in}}%
\pgfpathquadraticcurveto{\pgfqpoint{13.642407in}{2.734030in}}{\pgfqpoint{13.627075in}{2.734030in}}%
\pgfpathquadraticcurveto{\pgfqpoint{13.611763in}{2.734030in}}{\pgfqpoint{13.603677in}{2.746139in}}%
\pgfpathquadraticcurveto{\pgfqpoint{13.595591in}{2.758268in}}{\pgfqpoint{13.595591in}{2.781276in}}%
\pgfpathquadraticcurveto{\pgfqpoint{13.595591in}{2.804342in}}{\pgfqpoint{13.603677in}{2.816452in}}%
\pgfpathquadraticcurveto{\pgfqpoint{13.611763in}{2.828581in}}{\pgfqpoint{13.627075in}{2.828581in}}%
\pgfpathlineto{\pgfqpoint{13.627075in}{2.828581in}}%
\pgfpathclose%
\pgfusepath{fill}%
\end{pgfscope}%
\begin{pgfscope}%
\pgfsetbuttcap%
\pgfsetroundjoin%
\definecolor{currentfill}{rgb}{0.090196,0.168627,0.227451}%
\pgfsetfillcolor{currentfill}%
\pgfsetlinewidth{0.000000pt}%
\definecolor{currentstroke}{rgb}{0.090196,0.168627,0.227451}%
\pgfsetstrokecolor{currentstroke}%
\pgfsetdash{}{0pt}%
\pgfpathmoveto{\pgfqpoint{9.182851in}{2.572435in}}%
\pgfpathlineto{\pgfqpoint{9.182851in}{2.557730in}}%
\pgfpathquadraticcurveto{\pgfqpoint{9.174281in}{2.561836in}}{\pgfqpoint{9.166666in}{2.563841in}}%
\pgfpathquadraticcurveto{\pgfqpoint{9.159051in}{2.565870in}}{\pgfqpoint{9.151961in}{2.565870in}}%
\pgfpathquadraticcurveto{\pgfqpoint{9.139667in}{2.565870in}}{\pgfqpoint{9.132983in}{2.561096in}}%
\pgfpathquadraticcurveto{\pgfqpoint{9.126299in}{2.556321in}}{\pgfqpoint{9.126299in}{2.547513in}}%
\pgfpathquadraticcurveto{\pgfqpoint{9.126299in}{2.540113in}}{\pgfqpoint{9.130739in}{2.536341in}}%
\pgfpathquadraticcurveto{\pgfqpoint{9.135179in}{2.532593in}}{\pgfqpoint{9.147569in}{2.530278in}}%
\pgfpathlineto{\pgfqpoint{9.156664in}{2.528416in}}%
\pgfpathquadraticcurveto{\pgfqpoint{9.173517in}{2.525193in}}{\pgfqpoint{9.181538in}{2.517100in}}%
\pgfpathquadraticcurveto{\pgfqpoint{9.189559in}{2.509008in}}{\pgfqpoint{9.189559in}{2.495449in}}%
\pgfpathquadraticcurveto{\pgfqpoint{9.189559in}{2.479240in}}{\pgfqpoint{9.178697in}{2.470885in}}%
\pgfpathquadraticcurveto{\pgfqpoint{9.167860in}{2.462530in}}{\pgfqpoint{9.146900in}{2.462530in}}%
\pgfpathquadraticcurveto{\pgfqpoint{9.138999in}{2.462530in}}{\pgfqpoint{9.130071in}{2.464321in}}%
\pgfpathquadraticcurveto{\pgfqpoint{9.121167in}{2.466111in}}{\pgfqpoint{9.111618in}{2.469620in}}%
\pgfpathlineto{\pgfqpoint{9.111618in}{2.485137in}}%
\pgfpathquadraticcurveto{\pgfqpoint{9.120785in}{2.480004in}}{\pgfqpoint{9.129594in}{2.477378in}}%
\pgfpathquadraticcurveto{\pgfqpoint{9.138402in}{2.474776in}}{\pgfqpoint{9.146900in}{2.474776in}}%
\pgfpathquadraticcurveto{\pgfqpoint{9.159791in}{2.474776in}}{\pgfqpoint{9.166809in}{2.479837in}}%
\pgfpathquadraticcurveto{\pgfqpoint{9.173827in}{2.484922in}}{\pgfqpoint{9.173827in}{2.494327in}}%
\pgfpathquadraticcurveto{\pgfqpoint{9.173827in}{2.502515in}}{\pgfqpoint{9.168791in}{2.507146in}}%
\pgfpathquadraticcurveto{\pgfqpoint{9.163754in}{2.511777in}}{\pgfqpoint{9.152271in}{2.514093in}}%
\pgfpathlineto{\pgfqpoint{9.143081in}{2.515883in}}%
\pgfpathquadraticcurveto{\pgfqpoint{9.126228in}{2.519225in}}{\pgfqpoint{9.118684in}{2.526387in}}%
\pgfpathquadraticcurveto{\pgfqpoint{9.111165in}{2.533548in}}{\pgfqpoint{9.111165in}{2.546319in}}%
\pgfpathquadraticcurveto{\pgfqpoint{9.111165in}{2.561096in}}{\pgfqpoint{9.121573in}{2.569594in}}%
\pgfpathquadraticcurveto{\pgfqpoint{9.131981in}{2.578092in}}{\pgfqpoint{9.150242in}{2.578092in}}%
\pgfpathquadraticcurveto{\pgfqpoint{9.158096in}{2.578092in}}{\pgfqpoint{9.166212in}{2.576660in}}%
\pgfpathquadraticcurveto{\pgfqpoint{9.174353in}{2.575252in}}{\pgfqpoint{9.182851in}{2.572435in}}%
\pgfpathlineto{\pgfqpoint{9.182851in}{2.572435in}}%
\pgfpathclose%
\pgfpathmoveto{\pgfqpoint{9.250455in}{2.506692in}}%
\pgfpathquadraticcurveto{\pgfqpoint{9.233817in}{2.506692in}}{\pgfqpoint{9.227395in}{2.502897in}}%
\pgfpathquadraticcurveto{\pgfqpoint{9.220974in}{2.499101in}}{\pgfqpoint{9.220974in}{2.489911in}}%
\pgfpathquadraticcurveto{\pgfqpoint{9.220974in}{2.482606in}}{\pgfqpoint{9.225796in}{2.478309in}}%
\pgfpathquadraticcurveto{\pgfqpoint{9.230618in}{2.474036in}}{\pgfqpoint{9.238877in}{2.474036in}}%
\pgfpathquadraticcurveto{\pgfqpoint{9.250312in}{2.474036in}}{\pgfqpoint{9.257211in}{2.482129in}}%
\pgfpathquadraticcurveto{\pgfqpoint{9.264110in}{2.490221in}}{\pgfqpoint{9.264110in}{2.503637in}}%
\pgfpathlineto{\pgfqpoint{9.264110in}{2.506692in}}%
\pgfpathlineto{\pgfqpoint{9.250455in}{2.506692in}}%
\pgfpathlineto{\pgfqpoint{9.250455in}{2.506692in}}%
\pgfpathclose%
\pgfpathmoveto{\pgfqpoint{9.277836in}{2.512374in}}%
\pgfpathlineto{\pgfqpoint{9.277836in}{2.464702in}}%
\pgfpathlineto{\pgfqpoint{9.264110in}{2.464702in}}%
\pgfpathlineto{\pgfqpoint{9.264110in}{2.477378in}}%
\pgfpathquadraticcurveto{\pgfqpoint{9.259407in}{2.469787in}}{\pgfqpoint{9.252389in}{2.466159in}}%
\pgfpathquadraticcurveto{\pgfqpoint{9.245370in}{2.462530in}}{\pgfqpoint{9.235225in}{2.462530in}}%
\pgfpathquadraticcurveto{\pgfqpoint{9.222406in}{2.462530in}}{\pgfqpoint{9.214815in}{2.469739in}}%
\pgfpathquadraticcurveto{\pgfqpoint{9.207248in}{2.476949in}}{\pgfqpoint{9.207248in}{2.489028in}}%
\pgfpathquadraticcurveto{\pgfqpoint{9.207248in}{2.503112in}}{\pgfqpoint{9.216677in}{2.510273in}}%
\pgfpathquadraticcurveto{\pgfqpoint{9.226130in}{2.517435in}}{\pgfqpoint{9.244845in}{2.517435in}}%
\pgfpathlineto{\pgfqpoint{9.264110in}{2.517435in}}%
\pgfpathlineto{\pgfqpoint{9.264110in}{2.518795in}}%
\pgfpathquadraticcurveto{\pgfqpoint{9.264110in}{2.528272in}}{\pgfqpoint{9.257879in}{2.533452in}}%
\pgfpathquadraticcurveto{\pgfqpoint{9.251649in}{2.538633in}}{\pgfqpoint{9.240381in}{2.538633in}}%
\pgfpathquadraticcurveto{\pgfqpoint{9.233220in}{2.538633in}}{\pgfqpoint{9.226416in}{2.536914in}}%
\pgfpathquadraticcurveto{\pgfqpoint{9.219637in}{2.535195in}}{\pgfqpoint{9.213383in}{2.531758in}}%
\pgfpathlineto{\pgfqpoint{9.213383in}{2.544457in}}%
\pgfpathquadraticcurveto{\pgfqpoint{9.220902in}{2.547370in}}{\pgfqpoint{9.227992in}{2.548802in}}%
\pgfpathquadraticcurveto{\pgfqpoint{9.235082in}{2.550258in}}{\pgfqpoint{9.241790in}{2.550258in}}%
\pgfpathquadraticcurveto{\pgfqpoint{9.259932in}{2.550258in}}{\pgfqpoint{9.268884in}{2.540853in}}%
\pgfpathquadraticcurveto{\pgfqpoint{9.277836in}{2.531471in}}{\pgfqpoint{9.277836in}{2.512374in}}%
\pgfpathlineto{\pgfqpoint{9.277836in}{2.512374in}}%
\pgfpathclose%
\pgfpathmoveto{\pgfqpoint{9.371150in}{2.532211in}}%
\pgfpathquadraticcurveto{\pgfqpoint{9.376306in}{2.541473in}}{\pgfqpoint{9.383467in}{2.545866in}}%
\pgfpathquadraticcurveto{\pgfqpoint{9.390629in}{2.550258in}}{\pgfqpoint{9.400321in}{2.550258in}}%
\pgfpathquadraticcurveto{\pgfqpoint{9.413378in}{2.550258in}}{\pgfqpoint{9.420468in}{2.541115in}}%
\pgfpathquadraticcurveto{\pgfqpoint{9.427558in}{2.531996in}}{\pgfqpoint{9.427558in}{2.515143in}}%
\pgfpathlineto{\pgfqpoint{9.427558in}{2.464702in}}%
\pgfpathlineto{\pgfqpoint{9.413760in}{2.464702in}}%
\pgfpathlineto{\pgfqpoint{9.413760in}{2.514689in}}%
\pgfpathquadraticcurveto{\pgfqpoint{9.413760in}{2.526697in}}{\pgfqpoint{9.409487in}{2.532498in}}%
\pgfpathquadraticcurveto{\pgfqpoint{9.405238in}{2.538322in}}{\pgfqpoint{9.396525in}{2.538322in}}%
\pgfpathquadraticcurveto{\pgfqpoint{9.385854in}{2.538322in}}{\pgfqpoint{9.379648in}{2.531232in}}%
\pgfpathquadraticcurveto{\pgfqpoint{9.373465in}{2.524166in}}{\pgfqpoint{9.373465in}{2.511920in}}%
\pgfpathlineto{\pgfqpoint{9.373465in}{2.464702in}}%
\pgfpathlineto{\pgfqpoint{9.359667in}{2.464702in}}%
\pgfpathlineto{\pgfqpoint{9.359667in}{2.514689in}}%
\pgfpathquadraticcurveto{\pgfqpoint{9.359667in}{2.526768in}}{\pgfqpoint{9.355418in}{2.532545in}}%
\pgfpathquadraticcurveto{\pgfqpoint{9.351169in}{2.538322in}}{\pgfqpoint{9.342289in}{2.538322in}}%
\pgfpathquadraticcurveto{\pgfqpoint{9.331761in}{2.538322in}}{\pgfqpoint{9.325555in}{2.531209in}}%
\pgfpathquadraticcurveto{\pgfqpoint{9.319372in}{2.524095in}}{\pgfqpoint{9.319372in}{2.511920in}}%
\pgfpathlineto{\pgfqpoint{9.319372in}{2.464702in}}%
\pgfpathlineto{\pgfqpoint{9.305574in}{2.464702in}}%
\pgfpathlineto{\pgfqpoint{9.305574in}{2.548253in}}%
\pgfpathlineto{\pgfqpoint{9.319372in}{2.548253in}}%
\pgfpathlineto{\pgfqpoint{9.319372in}{2.535267in}}%
\pgfpathquadraticcurveto{\pgfqpoint{9.324075in}{2.542953in}}{\pgfqpoint{9.330640in}{2.546606in}}%
\pgfpathquadraticcurveto{\pgfqpoint{9.337204in}{2.550258in}}{\pgfqpoint{9.346228in}{2.550258in}}%
\pgfpathquadraticcurveto{\pgfqpoint{9.355347in}{2.550258in}}{\pgfqpoint{9.361720in}{2.545627in}}%
\pgfpathquadraticcurveto{\pgfqpoint{9.368094in}{2.541020in}}{\pgfqpoint{9.371150in}{2.532211in}}%
\pgfpathlineto{\pgfqpoint{9.371150in}{2.532211in}}%
\pgfpathclose%
\pgfpathmoveto{\pgfqpoint{9.468187in}{2.477235in}}%
\pgfpathlineto{\pgfqpoint{9.468187in}{2.432929in}}%
\pgfpathlineto{\pgfqpoint{9.454390in}{2.432929in}}%
\pgfpathlineto{\pgfqpoint{9.454390in}{2.548253in}}%
\pgfpathlineto{\pgfqpoint{9.468187in}{2.548253in}}%
\pgfpathlineto{\pgfqpoint{9.468187in}{2.535577in}}%
\pgfpathquadraticcurveto{\pgfqpoint{9.472532in}{2.543025in}}{\pgfqpoint{9.479120in}{2.546630in}}%
\pgfpathquadraticcurveto{\pgfqpoint{9.485733in}{2.550258in}}{\pgfqpoint{9.494900in}{2.550258in}}%
\pgfpathquadraticcurveto{\pgfqpoint{9.510130in}{2.550258in}}{\pgfqpoint{9.519630in}{2.538179in}}%
\pgfpathquadraticcurveto{\pgfqpoint{9.529155in}{2.526100in}}{\pgfqpoint{9.529155in}{2.506406in}}%
\pgfpathquadraticcurveto{\pgfqpoint{9.529155in}{2.486712in}}{\pgfqpoint{9.519630in}{2.474609in}}%
\pgfpathquadraticcurveto{\pgfqpoint{9.510130in}{2.462530in}}{\pgfqpoint{9.494900in}{2.462530in}}%
\pgfpathquadraticcurveto{\pgfqpoint{9.485733in}{2.462530in}}{\pgfqpoint{9.479120in}{2.466159in}}%
\pgfpathquadraticcurveto{\pgfqpoint{9.472532in}{2.469787in}}{\pgfqpoint{9.468187in}{2.477235in}}%
\pgfpathlineto{\pgfqpoint{9.468187in}{2.477235in}}%
\pgfpathclose%
\pgfpathmoveto{\pgfqpoint{9.514904in}{2.506406in}}%
\pgfpathquadraticcurveto{\pgfqpoint{9.514904in}{2.521541in}}{\pgfqpoint{9.508673in}{2.530158in}}%
\pgfpathquadraticcurveto{\pgfqpoint{9.502443in}{2.538776in}}{\pgfqpoint{9.491557in}{2.538776in}}%
\pgfpathquadraticcurveto{\pgfqpoint{9.480648in}{2.538776in}}{\pgfqpoint{9.474418in}{2.530158in}}%
\pgfpathquadraticcurveto{\pgfqpoint{9.468187in}{2.521541in}}{\pgfqpoint{9.468187in}{2.506406in}}%
\pgfpathquadraticcurveto{\pgfqpoint{9.468187in}{2.491271in}}{\pgfqpoint{9.474418in}{2.482654in}}%
\pgfpathquadraticcurveto{\pgfqpoint{9.480648in}{2.474036in}}{\pgfqpoint{9.491557in}{2.474036in}}%
\pgfpathquadraticcurveto{\pgfqpoint{9.502443in}{2.474036in}}{\pgfqpoint{9.508673in}{2.482654in}}%
\pgfpathquadraticcurveto{\pgfqpoint{9.514904in}{2.491271in}}{\pgfqpoint{9.514904in}{2.506406in}}%
\pgfpathlineto{\pgfqpoint{9.514904in}{2.506406in}}%
\pgfpathclose%
\pgfpathmoveto{\pgfqpoint{9.551905in}{2.580790in}}%
\pgfpathlineto{\pgfqpoint{9.565631in}{2.580790in}}%
\pgfpathlineto{\pgfqpoint{9.565631in}{2.464702in}}%
\pgfpathlineto{\pgfqpoint{9.551905in}{2.464702in}}%
\pgfpathlineto{\pgfqpoint{9.551905in}{2.580790in}}%
\pgfpathlineto{\pgfqpoint{9.551905in}{2.580790in}}%
\pgfpathclose%
\pgfpathmoveto{\pgfqpoint{9.665820in}{2.509915in}}%
\pgfpathlineto{\pgfqpoint{9.665820in}{2.503207in}}%
\pgfpathlineto{\pgfqpoint{9.602703in}{2.503207in}}%
\pgfpathquadraticcurveto{\pgfqpoint{9.603610in}{2.489028in}}{\pgfqpoint{9.611249in}{2.481604in}}%
\pgfpathquadraticcurveto{\pgfqpoint{9.618888in}{2.474179in}}{\pgfqpoint{9.632543in}{2.474179in}}%
\pgfpathquadraticcurveto{\pgfqpoint{9.640444in}{2.474179in}}{\pgfqpoint{9.647868in}{2.476113in}}%
\pgfpathquadraticcurveto{\pgfqpoint{9.655292in}{2.478047in}}{\pgfqpoint{9.662621in}{2.481938in}}%
\pgfpathlineto{\pgfqpoint{9.662621in}{2.468952in}}%
\pgfpathquadraticcurveto{\pgfqpoint{9.655221in}{2.465824in}}{\pgfqpoint{9.647462in}{2.464177in}}%
\pgfpathquadraticcurveto{\pgfqpoint{9.639704in}{2.462530in}}{\pgfqpoint{9.631731in}{2.462530in}}%
\pgfpathquadraticcurveto{\pgfqpoint{9.611727in}{2.462530in}}{\pgfqpoint{9.600054in}{2.474156in}}%
\pgfpathquadraticcurveto{\pgfqpoint{9.588380in}{2.485805in}}{\pgfqpoint{9.588380in}{2.505666in}}%
\pgfpathquadraticcurveto{\pgfqpoint{9.588380in}{2.526172in}}{\pgfqpoint{9.599457in}{2.538203in}}%
\pgfpathquadraticcurveto{\pgfqpoint{9.610533in}{2.550258in}}{\pgfqpoint{9.629344in}{2.550258in}}%
\pgfpathquadraticcurveto{\pgfqpoint{9.646197in}{2.550258in}}{\pgfqpoint{9.656008in}{2.539396in}}%
\pgfpathquadraticcurveto{\pgfqpoint{9.665820in}{2.528559in}}{\pgfqpoint{9.665820in}{2.509915in}}%
\pgfpathlineto{\pgfqpoint{9.665820in}{2.509915in}}%
\pgfpathclose%
\pgfpathmoveto{\pgfqpoint{9.652094in}{2.513949in}}%
\pgfpathquadraticcurveto{\pgfqpoint{9.651950in}{2.525193in}}{\pgfqpoint{9.645791in}{2.531901in}}%
\pgfpathquadraticcurveto{\pgfqpoint{9.639633in}{2.538633in}}{\pgfqpoint{9.629487in}{2.538633in}}%
\pgfpathquadraticcurveto{\pgfqpoint{9.618005in}{2.538633in}}{\pgfqpoint{9.611106in}{2.532140in}}%
\pgfpathquadraticcurveto{\pgfqpoint{9.604207in}{2.525646in}}{\pgfqpoint{9.603157in}{2.513854in}}%
\pgfpathlineto{\pgfqpoint{9.652094in}{2.513949in}}%
\pgfpathlineto{\pgfqpoint{9.652094in}{2.513949in}}%
\pgfpathclose%
\pgfpathmoveto{\pgfqpoint{9.681432in}{2.512660in}}%
\pgfpathlineto{\pgfqpoint{9.721631in}{2.512660in}}%
\pgfpathlineto{\pgfqpoint{9.721631in}{2.500438in}}%
\pgfpathlineto{\pgfqpoint{9.681432in}{2.500438in}}%
\pgfpathlineto{\pgfqpoint{9.681432in}{2.512660in}}%
\pgfpathlineto{\pgfqpoint{9.681432in}{2.512660in}}%
\pgfpathclose%
\pgfpathmoveto{\pgfqpoint{9.781453in}{2.506692in}}%
\pgfpathquadraticcurveto{\pgfqpoint{9.764815in}{2.506692in}}{\pgfqpoint{9.758393in}{2.502897in}}%
\pgfpathquadraticcurveto{\pgfqpoint{9.751972in}{2.499101in}}{\pgfqpoint{9.751972in}{2.489911in}}%
\pgfpathquadraticcurveto{\pgfqpoint{9.751972in}{2.482606in}}{\pgfqpoint{9.756794in}{2.478309in}}%
\pgfpathquadraticcurveto{\pgfqpoint{9.761616in}{2.474036in}}{\pgfqpoint{9.769876in}{2.474036in}}%
\pgfpathquadraticcurveto{\pgfqpoint{9.781310in}{2.474036in}}{\pgfqpoint{9.788209in}{2.482129in}}%
\pgfpathquadraticcurveto{\pgfqpoint{9.795108in}{2.490221in}}{\pgfqpoint{9.795108in}{2.503637in}}%
\pgfpathlineto{\pgfqpoint{9.795108in}{2.506692in}}%
\pgfpathlineto{\pgfqpoint{9.781453in}{2.506692in}}%
\pgfpathlineto{\pgfqpoint{9.781453in}{2.506692in}}%
\pgfpathclose%
\pgfpathmoveto{\pgfqpoint{9.808834in}{2.512374in}}%
\pgfpathlineto{\pgfqpoint{9.808834in}{2.464702in}}%
\pgfpathlineto{\pgfqpoint{9.795108in}{2.464702in}}%
\pgfpathlineto{\pgfqpoint{9.795108in}{2.477378in}}%
\pgfpathquadraticcurveto{\pgfqpoint{9.790405in}{2.469787in}}{\pgfqpoint{9.783387in}{2.466159in}}%
\pgfpathquadraticcurveto{\pgfqpoint{9.776369in}{2.462530in}}{\pgfqpoint{9.766223in}{2.462530in}}%
\pgfpathquadraticcurveto{\pgfqpoint{9.753404in}{2.462530in}}{\pgfqpoint{9.745813in}{2.469739in}}%
\pgfpathquadraticcurveto{\pgfqpoint{9.738246in}{2.476949in}}{\pgfqpoint{9.738246in}{2.489028in}}%
\pgfpathquadraticcurveto{\pgfqpoint{9.738246in}{2.503112in}}{\pgfqpoint{9.747675in}{2.510273in}}%
\pgfpathquadraticcurveto{\pgfqpoint{9.757128in}{2.517435in}}{\pgfqpoint{9.775844in}{2.517435in}}%
\pgfpathlineto{\pgfqpoint{9.795108in}{2.517435in}}%
\pgfpathlineto{\pgfqpoint{9.795108in}{2.518795in}}%
\pgfpathquadraticcurveto{\pgfqpoint{9.795108in}{2.528272in}}{\pgfqpoint{9.788877in}{2.533452in}}%
\pgfpathquadraticcurveto{\pgfqpoint{9.782647in}{2.538633in}}{\pgfqpoint{9.771380in}{2.538633in}}%
\pgfpathquadraticcurveto{\pgfqpoint{9.764218in}{2.538633in}}{\pgfqpoint{9.757415in}{2.536914in}}%
\pgfpathquadraticcurveto{\pgfqpoint{9.750635in}{2.535195in}}{\pgfqpoint{9.744381in}{2.531758in}}%
\pgfpathlineto{\pgfqpoint{9.744381in}{2.544457in}}%
\pgfpathquadraticcurveto{\pgfqpoint{9.751900in}{2.547370in}}{\pgfqpoint{9.758990in}{2.548802in}}%
\pgfpathquadraticcurveto{\pgfqpoint{9.766080in}{2.550258in}}{\pgfqpoint{9.772788in}{2.550258in}}%
\pgfpathquadraticcurveto{\pgfqpoint{9.790930in}{2.550258in}}{\pgfqpoint{9.799882in}{2.540853in}}%
\pgfpathquadraticcurveto{\pgfqpoint{9.808834in}{2.531471in}}{\pgfqpoint{9.808834in}{2.512374in}}%
\pgfpathlineto{\pgfqpoint{9.808834in}{2.512374in}}%
\pgfpathclose%
\pgfpathmoveto{\pgfqpoint{9.837098in}{2.580790in}}%
\pgfpathlineto{\pgfqpoint{9.850824in}{2.580790in}}%
\pgfpathlineto{\pgfqpoint{9.850824in}{2.464702in}}%
\pgfpathlineto{\pgfqpoint{9.837098in}{2.464702in}}%
\pgfpathlineto{\pgfqpoint{9.837098in}{2.580790in}}%
\pgfpathlineto{\pgfqpoint{9.837098in}{2.580790in}}%
\pgfpathclose%
\pgfpathmoveto{\pgfqpoint{9.879541in}{2.548253in}}%
\pgfpathlineto{\pgfqpoint{9.893268in}{2.548253in}}%
\pgfpathlineto{\pgfqpoint{9.893268in}{2.464702in}}%
\pgfpathlineto{\pgfqpoint{9.879541in}{2.464702in}}%
\pgfpathlineto{\pgfqpoint{9.879541in}{2.548253in}}%
\pgfpathlineto{\pgfqpoint{9.879541in}{2.548253in}}%
\pgfpathclose%
\pgfpathmoveto{\pgfqpoint{9.879541in}{2.580790in}}%
\pgfpathlineto{\pgfqpoint{9.893268in}{2.580790in}}%
\pgfpathlineto{\pgfqpoint{9.893268in}{2.563387in}}%
\pgfpathlineto{\pgfqpoint{9.879541in}{2.563387in}}%
\pgfpathlineto{\pgfqpoint{9.879541in}{2.580790in}}%
\pgfpathlineto{\pgfqpoint{9.879541in}{2.580790in}}%
\pgfpathclose%
\pgfpathmoveto{\pgfqpoint{9.976961in}{2.507456in}}%
\pgfpathquadraticcurveto{\pgfqpoint{9.976961in}{2.522376in}}{\pgfqpoint{9.970802in}{2.530564in}}%
\pgfpathquadraticcurveto{\pgfqpoint{9.964667in}{2.538776in}}{\pgfqpoint{9.953543in}{2.538776in}}%
\pgfpathquadraticcurveto{\pgfqpoint{9.942515in}{2.538776in}}{\pgfqpoint{9.936356in}{2.530564in}}%
\pgfpathquadraticcurveto{\pgfqpoint{9.930197in}{2.522376in}}{\pgfqpoint{9.930197in}{2.507456in}}%
\pgfpathquadraticcurveto{\pgfqpoint{9.930197in}{2.492608in}}{\pgfqpoint{9.936356in}{2.484396in}}%
\pgfpathquadraticcurveto{\pgfqpoint{9.942515in}{2.476185in}}{\pgfqpoint{9.953543in}{2.476185in}}%
\pgfpathquadraticcurveto{\pgfqpoint{9.964667in}{2.476185in}}{\pgfqpoint{9.970802in}{2.484396in}}%
\pgfpathquadraticcurveto{\pgfqpoint{9.976961in}{2.492608in}}{\pgfqpoint{9.976961in}{2.507456in}}%
\pgfpathlineto{\pgfqpoint{9.976961in}{2.507456in}}%
\pgfpathclose%
\pgfpathmoveto{\pgfqpoint{9.990687in}{2.475063in}}%
\pgfpathquadraticcurveto{\pgfqpoint{9.990687in}{2.453745in}}{\pgfqpoint{9.981210in}{2.443337in}}%
\pgfpathquadraticcurveto{\pgfqpoint{9.971757in}{2.432929in}}{\pgfqpoint{9.952206in}{2.432929in}}%
\pgfpathquadraticcurveto{\pgfqpoint{9.944973in}{2.432929in}}{\pgfqpoint{9.938552in}{2.434004in}}%
\pgfpathquadraticcurveto{\pgfqpoint{9.932130in}{2.435078in}}{\pgfqpoint{9.926091in}{2.437322in}}%
\pgfpathlineto{\pgfqpoint{9.926091in}{2.450666in}}%
\pgfpathquadraticcurveto{\pgfqpoint{9.932130in}{2.447396in}}{\pgfqpoint{9.938027in}{2.445844in}}%
\pgfpathquadraticcurveto{\pgfqpoint{9.943923in}{2.444268in}}{\pgfqpoint{9.950034in}{2.444268in}}%
\pgfpathquadraticcurveto{\pgfqpoint{9.963545in}{2.444268in}}{\pgfqpoint{9.970253in}{2.451311in}}%
\pgfpathquadraticcurveto{\pgfqpoint{9.976961in}{2.458353in}}{\pgfqpoint{9.976961in}{2.472604in}}%
\pgfpathlineto{\pgfqpoint{9.976961in}{2.479407in}}%
\pgfpathquadraticcurveto{\pgfqpoint{9.972712in}{2.472007in}}{\pgfqpoint{9.966076in}{2.468355in}}%
\pgfpathquadraticcurveto{\pgfqpoint{9.959439in}{2.464702in}}{\pgfqpoint{9.950177in}{2.464702in}}%
\pgfpathquadraticcurveto{\pgfqpoint{9.934828in}{2.464702in}}{\pgfqpoint{9.925423in}{2.476400in}}%
\pgfpathquadraticcurveto{\pgfqpoint{9.916017in}{2.488120in}}{\pgfqpoint{9.916017in}{2.507456in}}%
\pgfpathquadraticcurveto{\pgfqpoint{9.916017in}{2.526840in}}{\pgfqpoint{9.925423in}{2.538537in}}%
\pgfpathquadraticcurveto{\pgfqpoint{9.934828in}{2.550258in}}{\pgfqpoint{9.950177in}{2.550258in}}%
\pgfpathquadraticcurveto{\pgfqpoint{9.959439in}{2.550258in}}{\pgfqpoint{9.966076in}{2.546606in}}%
\pgfpathquadraticcurveto{\pgfqpoint{9.972712in}{2.542953in}}{\pgfqpoint{9.976961in}{2.535577in}}%
\pgfpathlineto{\pgfqpoint{9.976961in}{2.548253in}}%
\pgfpathlineto{\pgfqpoint{9.990687in}{2.548253in}}%
\pgfpathlineto{\pgfqpoint{9.990687in}{2.475063in}}%
\pgfpathlineto{\pgfqpoint{9.990687in}{2.475063in}}%
\pgfpathclose%
\pgfpathmoveto{\pgfqpoint{10.088441in}{2.515143in}}%
\pgfpathlineto{\pgfqpoint{10.088441in}{2.464702in}}%
\pgfpathlineto{\pgfqpoint{10.074715in}{2.464702in}}%
\pgfpathlineto{\pgfqpoint{10.074715in}{2.514689in}}%
\pgfpathquadraticcurveto{\pgfqpoint{10.074715in}{2.526554in}}{\pgfqpoint{10.070084in}{2.532426in}}%
\pgfpathquadraticcurveto{\pgfqpoint{10.065453in}{2.538322in}}{\pgfqpoint{10.056215in}{2.538322in}}%
\pgfpathquadraticcurveto{\pgfqpoint{10.045090in}{2.538322in}}{\pgfqpoint{10.038669in}{2.531232in}}%
\pgfpathquadraticcurveto{\pgfqpoint{10.032248in}{2.524166in}}{\pgfqpoint{10.032248in}{2.511920in}}%
\pgfpathlineto{\pgfqpoint{10.032248in}{2.464702in}}%
\pgfpathlineto{\pgfqpoint{10.018450in}{2.464702in}}%
\pgfpathlineto{\pgfqpoint{10.018450in}{2.548253in}}%
\pgfpathlineto{\pgfqpoint{10.032248in}{2.548253in}}%
\pgfpathlineto{\pgfqpoint{10.032248in}{2.535267in}}%
\pgfpathquadraticcurveto{\pgfqpoint{10.037189in}{2.542810in}}{\pgfqpoint{10.043849in}{2.546534in}}%
\pgfpathquadraticcurveto{\pgfqpoint{10.050533in}{2.550258in}}{\pgfqpoint{10.059270in}{2.550258in}}%
\pgfpathquadraticcurveto{\pgfqpoint{10.073665in}{2.550258in}}{\pgfqpoint{10.081041in}{2.541354in}}%
\pgfpathquadraticcurveto{\pgfqpoint{10.088441in}{2.532450in}}{\pgfqpoint{10.088441in}{2.515143in}}%
\pgfpathlineto{\pgfqpoint{10.088441in}{2.515143in}}%
\pgfpathclose%
\pgfpathmoveto{\pgfqpoint{10.187269in}{2.509915in}}%
\pgfpathlineto{\pgfqpoint{10.187269in}{2.503207in}}%
\pgfpathlineto{\pgfqpoint{10.124153in}{2.503207in}}%
\pgfpathquadraticcurveto{\pgfqpoint{10.125060in}{2.489028in}}{\pgfqpoint{10.132699in}{2.481604in}}%
\pgfpathquadraticcurveto{\pgfqpoint{10.140338in}{2.474179in}}{\pgfqpoint{10.153992in}{2.474179in}}%
\pgfpathquadraticcurveto{\pgfqpoint{10.161894in}{2.474179in}}{\pgfqpoint{10.169318in}{2.476113in}}%
\pgfpathquadraticcurveto{\pgfqpoint{10.176742in}{2.478047in}}{\pgfqpoint{10.184071in}{2.481938in}}%
\pgfpathlineto{\pgfqpoint{10.184071in}{2.468952in}}%
\pgfpathquadraticcurveto{\pgfqpoint{10.176670in}{2.465824in}}{\pgfqpoint{10.168912in}{2.464177in}}%
\pgfpathquadraticcurveto{\pgfqpoint{10.161154in}{2.462530in}}{\pgfqpoint{10.153181in}{2.462530in}}%
\pgfpathquadraticcurveto{\pgfqpoint{10.133176in}{2.462530in}}{\pgfqpoint{10.121503in}{2.474156in}}%
\pgfpathquadraticcurveto{\pgfqpoint{10.109830in}{2.485805in}}{\pgfqpoint{10.109830in}{2.505666in}}%
\pgfpathquadraticcurveto{\pgfqpoint{10.109830in}{2.526172in}}{\pgfqpoint{10.120906in}{2.538203in}}%
\pgfpathquadraticcurveto{\pgfqpoint{10.131983in}{2.550258in}}{\pgfqpoint{10.150794in}{2.550258in}}%
\pgfpathquadraticcurveto{\pgfqpoint{10.167647in}{2.550258in}}{\pgfqpoint{10.177458in}{2.539396in}}%
\pgfpathquadraticcurveto{\pgfqpoint{10.187269in}{2.528559in}}{\pgfqpoint{10.187269in}{2.509915in}}%
\pgfpathlineto{\pgfqpoint{10.187269in}{2.509915in}}%
\pgfpathclose%
\pgfpathmoveto{\pgfqpoint{10.173543in}{2.513949in}}%
\pgfpathquadraticcurveto{\pgfqpoint{10.173400in}{2.525193in}}{\pgfqpoint{10.167241in}{2.531901in}}%
\pgfpathquadraticcurveto{\pgfqpoint{10.161082in}{2.538633in}}{\pgfqpoint{10.150937in}{2.538633in}}%
\pgfpathquadraticcurveto{\pgfqpoint{10.139455in}{2.538633in}}{\pgfqpoint{10.132556in}{2.532140in}}%
\pgfpathquadraticcurveto{\pgfqpoint{10.125657in}{2.525646in}}{\pgfqpoint{10.124607in}{2.513854in}}%
\pgfpathlineto{\pgfqpoint{10.173543in}{2.513949in}}%
\pgfpathlineto{\pgfqpoint{10.173543in}{2.513949in}}%
\pgfpathclose%
\pgfpathmoveto{\pgfqpoint{10.264780in}{2.535577in}}%
\pgfpathlineto{\pgfqpoint{10.264780in}{2.580790in}}%
\pgfpathlineto{\pgfqpoint{10.278506in}{2.580790in}}%
\pgfpathlineto{\pgfqpoint{10.278506in}{2.464702in}}%
\pgfpathlineto{\pgfqpoint{10.264780in}{2.464702in}}%
\pgfpathlineto{\pgfqpoint{10.264780in}{2.477235in}}%
\pgfpathquadraticcurveto{\pgfqpoint{10.260459in}{2.469787in}}{\pgfqpoint{10.253847in}{2.466159in}}%
\pgfpathquadraticcurveto{\pgfqpoint{10.247258in}{2.462530in}}{\pgfqpoint{10.237996in}{2.462530in}}%
\pgfpathquadraticcurveto{\pgfqpoint{10.222862in}{2.462530in}}{\pgfqpoint{10.213337in}{2.474609in}}%
\pgfpathquadraticcurveto{\pgfqpoint{10.203836in}{2.486712in}}{\pgfqpoint{10.203836in}{2.506406in}}%
\pgfpathquadraticcurveto{\pgfqpoint{10.203836in}{2.526100in}}{\pgfqpoint{10.213337in}{2.538179in}}%
\pgfpathquadraticcurveto{\pgfqpoint{10.222862in}{2.550258in}}{\pgfqpoint{10.237996in}{2.550258in}}%
\pgfpathquadraticcurveto{\pgfqpoint{10.247258in}{2.550258in}}{\pgfqpoint{10.253847in}{2.546630in}}%
\pgfpathquadraticcurveto{\pgfqpoint{10.260459in}{2.543025in}}{\pgfqpoint{10.264780in}{2.535577in}}%
\pgfpathlineto{\pgfqpoint{10.264780in}{2.535577in}}%
\pgfpathclose%
\pgfpathmoveto{\pgfqpoint{10.218016in}{2.506406in}}%
\pgfpathquadraticcurveto{\pgfqpoint{10.218016in}{2.491271in}}{\pgfqpoint{10.224246in}{2.482654in}}%
\pgfpathquadraticcurveto{\pgfqpoint{10.230477in}{2.474036in}}{\pgfqpoint{10.241362in}{2.474036in}}%
\pgfpathquadraticcurveto{\pgfqpoint{10.252248in}{2.474036in}}{\pgfqpoint{10.258502in}{2.482654in}}%
\pgfpathquadraticcurveto{\pgfqpoint{10.264780in}{2.491271in}}{\pgfqpoint{10.264780in}{2.506406in}}%
\pgfpathquadraticcurveto{\pgfqpoint{10.264780in}{2.521541in}}{\pgfqpoint{10.258502in}{2.530158in}}%
\pgfpathquadraticcurveto{\pgfqpoint{10.252248in}{2.538776in}}{\pgfqpoint{10.241362in}{2.538776in}}%
\pgfpathquadraticcurveto{\pgfqpoint{10.230477in}{2.538776in}}{\pgfqpoint{10.224246in}{2.530158in}}%
\pgfpathquadraticcurveto{\pgfqpoint{10.218016in}{2.521541in}}{\pgfqpoint{10.218016in}{2.506406in}}%
\pgfpathlineto{\pgfqpoint{10.218016in}{2.506406in}}%
\pgfpathclose%
\pgfpathmoveto{\pgfqpoint{10.340501in}{2.576087in}}%
\pgfpathlineto{\pgfqpoint{10.434722in}{2.576087in}}%
\pgfpathlineto{\pgfqpoint{10.434722in}{2.563387in}}%
\pgfpathlineto{\pgfqpoint{10.395190in}{2.563387in}}%
\pgfpathlineto{\pgfqpoint{10.395190in}{2.464702in}}%
\pgfpathlineto{\pgfqpoint{10.380056in}{2.464702in}}%
\pgfpathlineto{\pgfqpoint{10.380056in}{2.563387in}}%
\pgfpathlineto{\pgfqpoint{10.340501in}{2.563387in}}%
\pgfpathlineto{\pgfqpoint{10.340501in}{2.576087in}}%
\pgfpathlineto{\pgfqpoint{10.340501in}{2.576087in}}%
\pgfpathclose%
\pgfpathmoveto{\pgfqpoint{10.439353in}{2.576087in}}%
\pgfpathlineto{\pgfqpoint{10.454559in}{2.576087in}}%
\pgfpathlineto{\pgfqpoint{10.477977in}{2.481938in}}%
\pgfpathlineto{\pgfqpoint{10.501323in}{2.576087in}}%
\pgfpathlineto{\pgfqpoint{10.518272in}{2.576087in}}%
\pgfpathlineto{\pgfqpoint{10.541690in}{2.481938in}}%
\pgfpathlineto{\pgfqpoint{10.565036in}{2.576087in}}%
\pgfpathlineto{\pgfqpoint{10.580338in}{2.576087in}}%
\pgfpathlineto{\pgfqpoint{10.552360in}{2.464702in}}%
\pgfpathlineto{\pgfqpoint{10.533406in}{2.464702in}}%
\pgfpathlineto{\pgfqpoint{10.509917in}{2.561382in}}%
\pgfpathlineto{\pgfqpoint{10.486189in}{2.464702in}}%
\pgfpathlineto{\pgfqpoint{10.467235in}{2.464702in}}%
\pgfpathlineto{\pgfqpoint{10.439353in}{2.576087in}}%
\pgfpathlineto{\pgfqpoint{10.439353in}{2.576087in}}%
\pgfpathclose%
\pgfpathmoveto{\pgfqpoint{10.600318in}{2.576087in}}%
\pgfpathlineto{\pgfqpoint{10.664318in}{2.576087in}}%
\pgfpathlineto{\pgfqpoint{10.664318in}{2.563387in}}%
\pgfpathlineto{\pgfqpoint{10.615381in}{2.563387in}}%
\pgfpathlineto{\pgfqpoint{10.615381in}{2.530564in}}%
\pgfpathlineto{\pgfqpoint{10.659544in}{2.530564in}}%
\pgfpathlineto{\pgfqpoint{10.659544in}{2.517888in}}%
\pgfpathlineto{\pgfqpoint{10.615381in}{2.517888in}}%
\pgfpathlineto{\pgfqpoint{10.615381in}{2.464702in}}%
\pgfpathlineto{\pgfqpoint{10.600318in}{2.464702in}}%
\pgfpathlineto{\pgfqpoint{10.600318in}{2.576087in}}%
\pgfpathlineto{\pgfqpoint{10.600318in}{2.576087in}}%
\pgfpathclose%
\pgfpathmoveto{\pgfqpoint{10.688189in}{2.576087in}}%
\pgfpathlineto{\pgfqpoint{10.758610in}{2.576087in}}%
\pgfpathlineto{\pgfqpoint{10.758610in}{2.563387in}}%
\pgfpathlineto{\pgfqpoint{10.703252in}{2.563387in}}%
\pgfpathlineto{\pgfqpoint{10.703252in}{2.530421in}}%
\pgfpathlineto{\pgfqpoint{10.756295in}{2.530421in}}%
\pgfpathlineto{\pgfqpoint{10.756295in}{2.517745in}}%
\pgfpathlineto{\pgfqpoint{10.703252in}{2.517745in}}%
\pgfpathlineto{\pgfqpoint{10.703252in}{2.477378in}}%
\pgfpathlineto{\pgfqpoint{10.759947in}{2.477378in}}%
\pgfpathlineto{\pgfqpoint{10.759947in}{2.464702in}}%
\pgfpathlineto{\pgfqpoint{10.688189in}{2.464702in}}%
\pgfpathlineto{\pgfqpoint{10.688189in}{2.576087in}}%
\pgfpathlineto{\pgfqpoint{10.688189in}{2.576087in}}%
\pgfpathclose%
\pgfpathmoveto{\pgfqpoint{10.904155in}{2.509915in}}%
\pgfpathlineto{\pgfqpoint{10.904155in}{2.503207in}}%
\pgfpathlineto{\pgfqpoint{10.841039in}{2.503207in}}%
\pgfpathquadraticcurveto{\pgfqpoint{10.841946in}{2.489028in}}{\pgfqpoint{10.849585in}{2.481604in}}%
\pgfpathquadraticcurveto{\pgfqpoint{10.857224in}{2.474179in}}{\pgfqpoint{10.870878in}{2.474179in}}%
\pgfpathquadraticcurveto{\pgfqpoint{10.878780in}{2.474179in}}{\pgfqpoint{10.886204in}{2.476113in}}%
\pgfpathquadraticcurveto{\pgfqpoint{10.893628in}{2.478047in}}{\pgfqpoint{10.900956in}{2.481938in}}%
\pgfpathlineto{\pgfqpoint{10.900956in}{2.468952in}}%
\pgfpathquadraticcurveto{\pgfqpoint{10.893556in}{2.465824in}}{\pgfqpoint{10.885798in}{2.464177in}}%
\pgfpathquadraticcurveto{\pgfqpoint{10.878040in}{2.462530in}}{\pgfqpoint{10.870067in}{2.462530in}}%
\pgfpathquadraticcurveto{\pgfqpoint{10.850062in}{2.462530in}}{\pgfqpoint{10.838389in}{2.474156in}}%
\pgfpathquadraticcurveto{\pgfqpoint{10.826716in}{2.485805in}}{\pgfqpoint{10.826716in}{2.505666in}}%
\pgfpathquadraticcurveto{\pgfqpoint{10.826716in}{2.526172in}}{\pgfqpoint{10.837792in}{2.538203in}}%
\pgfpathquadraticcurveto{\pgfqpoint{10.848869in}{2.550258in}}{\pgfqpoint{10.867679in}{2.550258in}}%
\pgfpathquadraticcurveto{\pgfqpoint{10.884533in}{2.550258in}}{\pgfqpoint{10.894344in}{2.539396in}}%
\pgfpathquadraticcurveto{\pgfqpoint{10.904155in}{2.528559in}}{\pgfqpoint{10.904155in}{2.509915in}}%
\pgfpathlineto{\pgfqpoint{10.904155in}{2.509915in}}%
\pgfpathclose%
\pgfpathmoveto{\pgfqpoint{10.890429in}{2.513949in}}%
\pgfpathquadraticcurveto{\pgfqpoint{10.890286in}{2.525193in}}{\pgfqpoint{10.884127in}{2.531901in}}%
\pgfpathquadraticcurveto{\pgfqpoint{10.877968in}{2.538633in}}{\pgfqpoint{10.867823in}{2.538633in}}%
\pgfpathquadraticcurveto{\pgfqpoint{10.856340in}{2.538633in}}{\pgfqpoint{10.849442in}{2.532140in}}%
\pgfpathquadraticcurveto{\pgfqpoint{10.842543in}{2.525646in}}{\pgfqpoint{10.841492in}{2.513854in}}%
\pgfpathlineto{\pgfqpoint{10.890429in}{2.513949in}}%
\pgfpathlineto{\pgfqpoint{10.890429in}{2.513949in}}%
\pgfpathclose%
\pgfpathmoveto{\pgfqpoint{10.979947in}{2.545794in}}%
\pgfpathlineto{\pgfqpoint{10.979947in}{2.532808in}}%
\pgfpathquadraticcurveto{\pgfqpoint{10.974146in}{2.535792in}}{\pgfqpoint{10.967868in}{2.537272in}}%
\pgfpathquadraticcurveto{\pgfqpoint{10.961614in}{2.538776in}}{\pgfqpoint{10.954882in}{2.538776in}}%
\pgfpathquadraticcurveto{\pgfqpoint{10.944665in}{2.538776in}}{\pgfqpoint{10.939557in}{2.535649in}}%
\pgfpathquadraticcurveto{\pgfqpoint{10.934448in}{2.532521in}}{\pgfqpoint{10.934448in}{2.526243in}}%
\pgfpathquadraticcurveto{\pgfqpoint{10.934448in}{2.521469in}}{\pgfqpoint{10.938100in}{2.518748in}}%
\pgfpathquadraticcurveto{\pgfqpoint{10.941753in}{2.516026in}}{\pgfqpoint{10.952805in}{2.513567in}}%
\pgfpathlineto{\pgfqpoint{10.957508in}{2.512517in}}%
\pgfpathquadraticcurveto{\pgfqpoint{10.972117in}{2.509390in}}{\pgfqpoint{10.978276in}{2.503685in}}%
\pgfpathquadraticcurveto{\pgfqpoint{10.984435in}{2.497979in}}{\pgfqpoint{10.984435in}{2.487762in}}%
\pgfpathquadraticcurveto{\pgfqpoint{10.984435in}{2.476113in}}{\pgfqpoint{10.975221in}{2.469310in}}%
\pgfpathquadraticcurveto{\pgfqpoint{10.966006in}{2.462530in}}{\pgfqpoint{10.949893in}{2.462530in}}%
\pgfpathquadraticcurveto{\pgfqpoint{10.943185in}{2.462530in}}{\pgfqpoint{10.935904in}{2.463843in}}%
\pgfpathquadraticcurveto{\pgfqpoint{10.928623in}{2.465156in}}{\pgfqpoint{10.920579in}{2.467758in}}%
\pgfpathlineto{\pgfqpoint{10.920579in}{2.481938in}}%
\pgfpathquadraticcurveto{\pgfqpoint{10.928194in}{2.477975in}}{\pgfqpoint{10.935570in}{2.475994in}}%
\pgfpathquadraticcurveto{\pgfqpoint{10.942946in}{2.474036in}}{\pgfqpoint{10.950203in}{2.474036in}}%
\pgfpathquadraticcurveto{\pgfqpoint{10.959895in}{2.474036in}}{\pgfqpoint{10.965099in}{2.477354in}}%
\pgfpathquadraticcurveto{\pgfqpoint{10.970327in}{2.480673in}}{\pgfqpoint{10.970327in}{2.486712in}}%
\pgfpathquadraticcurveto{\pgfqpoint{10.970327in}{2.492298in}}{\pgfqpoint{10.966555in}{2.495282in}}%
\pgfpathquadraticcurveto{\pgfqpoint{10.962807in}{2.498266in}}{\pgfqpoint{10.950036in}{2.501035in}}%
\pgfpathlineto{\pgfqpoint{10.945262in}{2.502157in}}%
\pgfpathquadraticcurveto{\pgfqpoint{10.932515in}{2.504831in}}{\pgfqpoint{10.926833in}{2.510393in}}%
\pgfpathquadraticcurveto{\pgfqpoint{10.921176in}{2.515955in}}{\pgfqpoint{10.921176in}{2.525646in}}%
\pgfpathquadraticcurveto{\pgfqpoint{10.921176in}{2.537439in}}{\pgfqpoint{10.929531in}{2.543837in}}%
\pgfpathquadraticcurveto{\pgfqpoint{10.937886in}{2.550258in}}{\pgfqpoint{10.953259in}{2.550258in}}%
\pgfpathquadraticcurveto{\pgfqpoint{10.960850in}{2.550258in}}{\pgfqpoint{10.967558in}{2.549136in}}%
\pgfpathquadraticcurveto{\pgfqpoint{10.974290in}{2.548038in}}{\pgfqpoint{10.979947in}{2.545794in}}%
\pgfpathlineto{\pgfqpoint{10.979947in}{2.545794in}}%
\pgfpathclose%
\pgfpathmoveto{\pgfqpoint{11.019860in}{2.571981in}}%
\pgfpathlineto{\pgfqpoint{11.019860in}{2.548253in}}%
\pgfpathlineto{\pgfqpoint{11.048124in}{2.548253in}}%
\pgfpathlineto{\pgfqpoint{11.048124in}{2.537582in}}%
\pgfpathlineto{\pgfqpoint{11.019860in}{2.537582in}}%
\pgfpathlineto{\pgfqpoint{11.019860in}{2.492226in}}%
\pgfpathquadraticcurveto{\pgfqpoint{11.019860in}{2.482009in}}{\pgfqpoint{11.022653in}{2.479097in}}%
\pgfpathquadraticcurveto{\pgfqpoint{11.025446in}{2.476185in}}{\pgfqpoint{11.034040in}{2.476185in}}%
\pgfpathlineto{\pgfqpoint{11.048124in}{2.476185in}}%
\pgfpathlineto{\pgfqpoint{11.048124in}{2.464702in}}%
\pgfpathlineto{\pgfqpoint{11.034040in}{2.464702in}}%
\pgfpathquadraticcurveto{\pgfqpoint{11.018142in}{2.464702in}}{\pgfqpoint{11.012102in}{2.470623in}}%
\pgfpathquadraticcurveto{\pgfqpoint{11.006063in}{2.476567in}}{\pgfqpoint{11.006063in}{2.492226in}}%
\pgfpathlineto{\pgfqpoint{11.006063in}{2.537582in}}%
\pgfpathlineto{\pgfqpoint{10.995989in}{2.537582in}}%
\pgfpathlineto{\pgfqpoint{10.995989in}{2.548253in}}%
\pgfpathlineto{\pgfqpoint{11.006063in}{2.548253in}}%
\pgfpathlineto{\pgfqpoint{11.006063in}{2.571981in}}%
\pgfpathlineto{\pgfqpoint{11.019860in}{2.571981in}}%
\pgfpathlineto{\pgfqpoint{11.019860in}{2.571981in}}%
\pgfpathclose%
\pgfpathmoveto{\pgfqpoint{11.066171in}{2.548253in}}%
\pgfpathlineto{\pgfqpoint{11.079897in}{2.548253in}}%
\pgfpathlineto{\pgfqpoint{11.079897in}{2.464702in}}%
\pgfpathlineto{\pgfqpoint{11.066171in}{2.464702in}}%
\pgfpathlineto{\pgfqpoint{11.066171in}{2.548253in}}%
\pgfpathlineto{\pgfqpoint{11.066171in}{2.548253in}}%
\pgfpathclose%
\pgfpathmoveto{\pgfqpoint{11.066171in}{2.580790in}}%
\pgfpathlineto{\pgfqpoint{11.079897in}{2.580790in}}%
\pgfpathlineto{\pgfqpoint{11.079897in}{2.563387in}}%
\pgfpathlineto{\pgfqpoint{11.066171in}{2.563387in}}%
\pgfpathlineto{\pgfqpoint{11.066171in}{2.580790in}}%
\pgfpathlineto{\pgfqpoint{11.066171in}{2.580790in}}%
\pgfpathclose%
\pgfpathmoveto{\pgfqpoint{11.173665in}{2.532211in}}%
\pgfpathquadraticcurveto{\pgfqpoint{11.178821in}{2.541473in}}{\pgfqpoint{11.185982in}{2.545866in}}%
\pgfpathquadraticcurveto{\pgfqpoint{11.193144in}{2.550258in}}{\pgfqpoint{11.202836in}{2.550258in}}%
\pgfpathquadraticcurveto{\pgfqpoint{11.215893in}{2.550258in}}{\pgfqpoint{11.222983in}{2.541115in}}%
\pgfpathquadraticcurveto{\pgfqpoint{11.230073in}{2.531996in}}{\pgfqpoint{11.230073in}{2.515143in}}%
\pgfpathlineto{\pgfqpoint{11.230073in}{2.464702in}}%
\pgfpathlineto{\pgfqpoint{11.216275in}{2.464702in}}%
\pgfpathlineto{\pgfqpoint{11.216275in}{2.514689in}}%
\pgfpathquadraticcurveto{\pgfqpoint{11.216275in}{2.526697in}}{\pgfqpoint{11.212002in}{2.532498in}}%
\pgfpathquadraticcurveto{\pgfqpoint{11.207753in}{2.538322in}}{\pgfqpoint{11.199040in}{2.538322in}}%
\pgfpathquadraticcurveto{\pgfqpoint{11.188370in}{2.538322in}}{\pgfqpoint{11.182163in}{2.531232in}}%
\pgfpathquadraticcurveto{\pgfqpoint{11.175980in}{2.524166in}}{\pgfqpoint{11.175980in}{2.511920in}}%
\pgfpathlineto{\pgfqpoint{11.175980in}{2.464702in}}%
\pgfpathlineto{\pgfqpoint{11.162182in}{2.464702in}}%
\pgfpathlineto{\pgfqpoint{11.162182in}{2.514689in}}%
\pgfpathquadraticcurveto{\pgfqpoint{11.162182in}{2.526768in}}{\pgfqpoint{11.157933in}{2.532545in}}%
\pgfpathquadraticcurveto{\pgfqpoint{11.153684in}{2.538322in}}{\pgfqpoint{11.144804in}{2.538322in}}%
\pgfpathquadraticcurveto{\pgfqpoint{11.134277in}{2.538322in}}{\pgfqpoint{11.128070in}{2.531209in}}%
\pgfpathquadraticcurveto{\pgfqpoint{11.121887in}{2.524095in}}{\pgfqpoint{11.121887in}{2.511920in}}%
\pgfpathlineto{\pgfqpoint{11.121887in}{2.464702in}}%
\pgfpathlineto{\pgfqpoint{11.108090in}{2.464702in}}%
\pgfpathlineto{\pgfqpoint{11.108090in}{2.548253in}}%
\pgfpathlineto{\pgfqpoint{11.121887in}{2.548253in}}%
\pgfpathlineto{\pgfqpoint{11.121887in}{2.535267in}}%
\pgfpathquadraticcurveto{\pgfqpoint{11.126590in}{2.542953in}}{\pgfqpoint{11.133155in}{2.546606in}}%
\pgfpathquadraticcurveto{\pgfqpoint{11.139719in}{2.550258in}}{\pgfqpoint{11.148743in}{2.550258in}}%
\pgfpathquadraticcurveto{\pgfqpoint{11.157862in}{2.550258in}}{\pgfqpoint{11.164235in}{2.545627in}}%
\pgfpathquadraticcurveto{\pgfqpoint{11.170609in}{2.541020in}}{\pgfqpoint{11.173665in}{2.532211in}}%
\pgfpathlineto{\pgfqpoint{11.173665in}{2.532211in}}%
\pgfpathclose%
\pgfpathmoveto{\pgfqpoint{11.295409in}{2.506692in}}%
\pgfpathquadraticcurveto{\pgfqpoint{11.278771in}{2.506692in}}{\pgfqpoint{11.272350in}{2.502897in}}%
\pgfpathquadraticcurveto{\pgfqpoint{11.265928in}{2.499101in}}{\pgfqpoint{11.265928in}{2.489911in}}%
\pgfpathquadraticcurveto{\pgfqpoint{11.265928in}{2.482606in}}{\pgfqpoint{11.270750in}{2.478309in}}%
\pgfpathquadraticcurveto{\pgfqpoint{11.275572in}{2.474036in}}{\pgfqpoint{11.283832in}{2.474036in}}%
\pgfpathquadraticcurveto{\pgfqpoint{11.295266in}{2.474036in}}{\pgfqpoint{11.302165in}{2.482129in}}%
\pgfpathquadraticcurveto{\pgfqpoint{11.309064in}{2.490221in}}{\pgfqpoint{11.309064in}{2.503637in}}%
\pgfpathlineto{\pgfqpoint{11.309064in}{2.506692in}}%
\pgfpathlineto{\pgfqpoint{11.295409in}{2.506692in}}%
\pgfpathlineto{\pgfqpoint{11.295409in}{2.506692in}}%
\pgfpathclose%
\pgfpathmoveto{\pgfqpoint{11.322790in}{2.512374in}}%
\pgfpathlineto{\pgfqpoint{11.322790in}{2.464702in}}%
\pgfpathlineto{\pgfqpoint{11.309064in}{2.464702in}}%
\pgfpathlineto{\pgfqpoint{11.309064in}{2.477378in}}%
\pgfpathquadraticcurveto{\pgfqpoint{11.304361in}{2.469787in}}{\pgfqpoint{11.297343in}{2.466159in}}%
\pgfpathquadraticcurveto{\pgfqpoint{11.290325in}{2.462530in}}{\pgfqpoint{11.280179in}{2.462530in}}%
\pgfpathquadraticcurveto{\pgfqpoint{11.267360in}{2.462530in}}{\pgfqpoint{11.259769in}{2.469739in}}%
\pgfpathquadraticcurveto{\pgfqpoint{11.252202in}{2.476949in}}{\pgfqpoint{11.252202in}{2.489028in}}%
\pgfpathquadraticcurveto{\pgfqpoint{11.252202in}{2.503112in}}{\pgfqpoint{11.261631in}{2.510273in}}%
\pgfpathquadraticcurveto{\pgfqpoint{11.271084in}{2.517435in}}{\pgfqpoint{11.289800in}{2.517435in}}%
\pgfpathlineto{\pgfqpoint{11.309064in}{2.517435in}}%
\pgfpathlineto{\pgfqpoint{11.309064in}{2.518795in}}%
\pgfpathquadraticcurveto{\pgfqpoint{11.309064in}{2.528272in}}{\pgfqpoint{11.302834in}{2.533452in}}%
\pgfpathquadraticcurveto{\pgfqpoint{11.296603in}{2.538633in}}{\pgfqpoint{11.285336in}{2.538633in}}%
\pgfpathquadraticcurveto{\pgfqpoint{11.278174in}{2.538633in}}{\pgfqpoint{11.271371in}{2.536914in}}%
\pgfpathquadraticcurveto{\pgfqpoint{11.264591in}{2.535195in}}{\pgfqpoint{11.258337in}{2.531758in}}%
\pgfpathlineto{\pgfqpoint{11.258337in}{2.544457in}}%
\pgfpathquadraticcurveto{\pgfqpoint{11.265857in}{2.547370in}}{\pgfqpoint{11.272946in}{2.548802in}}%
\pgfpathquadraticcurveto{\pgfqpoint{11.280036in}{2.550258in}}{\pgfqpoint{11.286744in}{2.550258in}}%
\pgfpathquadraticcurveto{\pgfqpoint{11.304886in}{2.550258in}}{\pgfqpoint{11.313838in}{2.540853in}}%
\pgfpathquadraticcurveto{\pgfqpoint{11.322790in}{2.531471in}}{\pgfqpoint{11.322790in}{2.512374in}}%
\pgfpathlineto{\pgfqpoint{11.322790in}{2.512374in}}%
\pgfpathclose%
\pgfpathmoveto{\pgfqpoint{11.364637in}{2.571981in}}%
\pgfpathlineto{\pgfqpoint{11.364637in}{2.548253in}}%
\pgfpathlineto{\pgfqpoint{11.392901in}{2.548253in}}%
\pgfpathlineto{\pgfqpoint{11.392901in}{2.537582in}}%
\pgfpathlineto{\pgfqpoint{11.364637in}{2.537582in}}%
\pgfpathlineto{\pgfqpoint{11.364637in}{2.492226in}}%
\pgfpathquadraticcurveto{\pgfqpoint{11.364637in}{2.482009in}}{\pgfqpoint{11.367430in}{2.479097in}}%
\pgfpathquadraticcurveto{\pgfqpoint{11.370223in}{2.476185in}}{\pgfqpoint{11.378817in}{2.476185in}}%
\pgfpathlineto{\pgfqpoint{11.392901in}{2.476185in}}%
\pgfpathlineto{\pgfqpoint{11.392901in}{2.464702in}}%
\pgfpathlineto{\pgfqpoint{11.378817in}{2.464702in}}%
\pgfpathquadraticcurveto{\pgfqpoint{11.362918in}{2.464702in}}{\pgfqpoint{11.356879in}{2.470623in}}%
\pgfpathquadraticcurveto{\pgfqpoint{11.350839in}{2.476567in}}{\pgfqpoint{11.350839in}{2.492226in}}%
\pgfpathlineto{\pgfqpoint{11.350839in}{2.537582in}}%
\pgfpathlineto{\pgfqpoint{11.340765in}{2.537582in}}%
\pgfpathlineto{\pgfqpoint{11.340765in}{2.548253in}}%
\pgfpathlineto{\pgfqpoint{11.350839in}{2.548253in}}%
\pgfpathlineto{\pgfqpoint{11.350839in}{2.571981in}}%
\pgfpathlineto{\pgfqpoint{11.364637in}{2.571981in}}%
\pgfpathlineto{\pgfqpoint{11.364637in}{2.571981in}}%
\pgfpathclose%
\pgfpathmoveto{\pgfqpoint{11.482419in}{2.509915in}}%
\pgfpathlineto{\pgfqpoint{11.482419in}{2.503207in}}%
\pgfpathlineto{\pgfqpoint{11.419303in}{2.503207in}}%
\pgfpathquadraticcurveto{\pgfqpoint{11.420210in}{2.489028in}}{\pgfqpoint{11.427849in}{2.481604in}}%
\pgfpathquadraticcurveto{\pgfqpoint{11.435488in}{2.474179in}}{\pgfqpoint{11.449142in}{2.474179in}}%
\pgfpathquadraticcurveto{\pgfqpoint{11.457044in}{2.474179in}}{\pgfqpoint{11.464468in}{2.476113in}}%
\pgfpathquadraticcurveto{\pgfqpoint{11.471892in}{2.478047in}}{\pgfqpoint{11.479220in}{2.481938in}}%
\pgfpathlineto{\pgfqpoint{11.479220in}{2.468952in}}%
\pgfpathquadraticcurveto{\pgfqpoint{11.471820in}{2.465824in}}{\pgfqpoint{11.464062in}{2.464177in}}%
\pgfpathquadraticcurveto{\pgfqpoint{11.456304in}{2.462530in}}{\pgfqpoint{11.448330in}{2.462530in}}%
\pgfpathquadraticcurveto{\pgfqpoint{11.428326in}{2.462530in}}{\pgfqpoint{11.416653in}{2.474156in}}%
\pgfpathquadraticcurveto{\pgfqpoint{11.404980in}{2.485805in}}{\pgfqpoint{11.404980in}{2.505666in}}%
\pgfpathquadraticcurveto{\pgfqpoint{11.404980in}{2.526172in}}{\pgfqpoint{11.416056in}{2.538203in}}%
\pgfpathquadraticcurveto{\pgfqpoint{11.427133in}{2.550258in}}{\pgfqpoint{11.445943in}{2.550258in}}%
\pgfpathquadraticcurveto{\pgfqpoint{11.462797in}{2.550258in}}{\pgfqpoint{11.472608in}{2.539396in}}%
\pgfpathquadraticcurveto{\pgfqpoint{11.482419in}{2.528559in}}{\pgfqpoint{11.482419in}{2.509915in}}%
\pgfpathlineto{\pgfqpoint{11.482419in}{2.509915in}}%
\pgfpathclose%
\pgfpathmoveto{\pgfqpoint{11.468693in}{2.513949in}}%
\pgfpathquadraticcurveto{\pgfqpoint{11.468550in}{2.525193in}}{\pgfqpoint{11.462391in}{2.531901in}}%
\pgfpathquadraticcurveto{\pgfqpoint{11.456232in}{2.538633in}}{\pgfqpoint{11.446087in}{2.538633in}}%
\pgfpathquadraticcurveto{\pgfqpoint{11.434604in}{2.538633in}}{\pgfqpoint{11.427705in}{2.532140in}}%
\pgfpathquadraticcurveto{\pgfqpoint{11.420807in}{2.525646in}}{\pgfqpoint{11.419756in}{2.513854in}}%
\pgfpathlineto{\pgfqpoint{11.468693in}{2.513949in}}%
\pgfpathlineto{\pgfqpoint{11.468693in}{2.513949in}}%
\pgfpathclose%
\pgfpathmoveto{\pgfqpoint{11.586475in}{2.580623in}}%
\pgfpathquadraticcurveto{\pgfqpoint{11.576497in}{2.563483in}}{\pgfqpoint{11.571627in}{2.546677in}}%
\pgfpathquadraticcurveto{\pgfqpoint{11.566781in}{2.529896in}}{\pgfqpoint{11.566781in}{2.512660in}}%
\pgfpathquadraticcurveto{\pgfqpoint{11.566781in}{2.495449in}}{\pgfqpoint{11.571675in}{2.478548in}}%
\pgfpathquadraticcurveto{\pgfqpoint{11.576568in}{2.461647in}}{\pgfqpoint{11.586475in}{2.444555in}}%
\pgfpathlineto{\pgfqpoint{11.574539in}{2.444555in}}%
\pgfpathquadraticcurveto{\pgfqpoint{11.563367in}{2.462100in}}{\pgfqpoint{11.557805in}{2.479025in}}%
\pgfpathquadraticcurveto{\pgfqpoint{11.552243in}{2.495950in}}{\pgfqpoint{11.552243in}{2.512660in}}%
\pgfpathquadraticcurveto{\pgfqpoint{11.552243in}{2.529299in}}{\pgfqpoint{11.557758in}{2.546152in}}%
\pgfpathquadraticcurveto{\pgfqpoint{11.563296in}{2.563029in}}{\pgfqpoint{11.574539in}{2.580623in}}%
\pgfpathlineto{\pgfqpoint{11.586475in}{2.580623in}}%
\pgfpathlineto{\pgfqpoint{11.586475in}{2.580623in}}%
\pgfpathclose%
\pgfpathmoveto{\pgfqpoint{11.598268in}{2.576087in}}%
\pgfpathlineto{\pgfqpoint{11.692488in}{2.576087in}}%
\pgfpathlineto{\pgfqpoint{11.692488in}{2.563387in}}%
\pgfpathlineto{\pgfqpoint{11.652957in}{2.563387in}}%
\pgfpathlineto{\pgfqpoint{11.652957in}{2.464702in}}%
\pgfpathlineto{\pgfqpoint{11.637823in}{2.464702in}}%
\pgfpathlineto{\pgfqpoint{11.637823in}{2.563387in}}%
\pgfpathlineto{\pgfqpoint{11.598268in}{2.563387in}}%
\pgfpathlineto{\pgfqpoint{11.598268in}{2.576087in}}%
\pgfpathlineto{\pgfqpoint{11.598268in}{2.576087in}}%
\pgfpathclose%
\pgfpathmoveto{\pgfqpoint{11.697120in}{2.576087in}}%
\pgfpathlineto{\pgfqpoint{11.712326in}{2.576087in}}%
\pgfpathlineto{\pgfqpoint{11.735744in}{2.481938in}}%
\pgfpathlineto{\pgfqpoint{11.759090in}{2.576087in}}%
\pgfpathlineto{\pgfqpoint{11.776039in}{2.576087in}}%
\pgfpathlineto{\pgfqpoint{11.799457in}{2.481938in}}%
\pgfpathlineto{\pgfqpoint{11.822803in}{2.576087in}}%
\pgfpathlineto{\pgfqpoint{11.838105in}{2.576087in}}%
\pgfpathlineto{\pgfqpoint{11.810127in}{2.464702in}}%
\pgfpathlineto{\pgfqpoint{11.791173in}{2.464702in}}%
\pgfpathlineto{\pgfqpoint{11.767684in}{2.561382in}}%
\pgfpathlineto{\pgfqpoint{11.743955in}{2.464702in}}%
\pgfpathlineto{\pgfqpoint{11.725001in}{2.464702in}}%
\pgfpathlineto{\pgfqpoint{11.697120in}{2.576087in}}%
\pgfpathlineto{\pgfqpoint{11.697120in}{2.576087in}}%
\pgfpathclose%
\pgfpathmoveto{\pgfqpoint{11.858085in}{2.576087in}}%
\pgfpathlineto{\pgfqpoint{11.922085in}{2.576087in}}%
\pgfpathlineto{\pgfqpoint{11.922085in}{2.563387in}}%
\pgfpathlineto{\pgfqpoint{11.873148in}{2.563387in}}%
\pgfpathlineto{\pgfqpoint{11.873148in}{2.530564in}}%
\pgfpathlineto{\pgfqpoint{11.917311in}{2.530564in}}%
\pgfpathlineto{\pgfqpoint{11.917311in}{2.517888in}}%
\pgfpathlineto{\pgfqpoint{11.873148in}{2.517888in}}%
\pgfpathlineto{\pgfqpoint{11.873148in}{2.464702in}}%
\pgfpathlineto{\pgfqpoint{11.858085in}{2.464702in}}%
\pgfpathlineto{\pgfqpoint{11.858085in}{2.576087in}}%
\pgfpathlineto{\pgfqpoint{11.858085in}{2.576087in}}%
\pgfpathclose%
\pgfpathmoveto{\pgfqpoint{11.945956in}{2.576087in}}%
\pgfpathlineto{\pgfqpoint{12.016377in}{2.576087in}}%
\pgfpathlineto{\pgfqpoint{12.016377in}{2.563387in}}%
\pgfpathlineto{\pgfqpoint{11.961019in}{2.563387in}}%
\pgfpathlineto{\pgfqpoint{11.961019in}{2.530421in}}%
\pgfpathlineto{\pgfqpoint{12.014062in}{2.530421in}}%
\pgfpathlineto{\pgfqpoint{12.014062in}{2.517745in}}%
\pgfpathlineto{\pgfqpoint{11.961019in}{2.517745in}}%
\pgfpathlineto{\pgfqpoint{11.961019in}{2.477378in}}%
\pgfpathlineto{\pgfqpoint{12.017714in}{2.477378in}}%
\pgfpathlineto{\pgfqpoint{12.017714in}{2.464702in}}%
\pgfpathlineto{\pgfqpoint{11.945956in}{2.464702in}}%
\pgfpathlineto{\pgfqpoint{11.945956in}{2.576087in}}%
\pgfpathlineto{\pgfqpoint{11.945956in}{2.576087in}}%
\pgfpathclose%
\pgfpathmoveto{\pgfqpoint{12.157816in}{2.572435in}}%
\pgfpathlineto{\pgfqpoint{12.157816in}{2.557730in}}%
\pgfpathquadraticcurveto{\pgfqpoint{12.149246in}{2.561836in}}{\pgfqpoint{12.141631in}{2.563841in}}%
\pgfpathquadraticcurveto{\pgfqpoint{12.134016in}{2.565870in}}{\pgfqpoint{12.126926in}{2.565870in}}%
\pgfpathquadraticcurveto{\pgfqpoint{12.114633in}{2.565870in}}{\pgfqpoint{12.107949in}{2.561096in}}%
\pgfpathquadraticcurveto{\pgfqpoint{12.101265in}{2.556321in}}{\pgfqpoint{12.101265in}{2.547513in}}%
\pgfpathquadraticcurveto{\pgfqpoint{12.101265in}{2.540113in}}{\pgfqpoint{12.105705in}{2.536341in}}%
\pgfpathquadraticcurveto{\pgfqpoint{12.110145in}{2.532593in}}{\pgfqpoint{12.122534in}{2.530278in}}%
\pgfpathlineto{\pgfqpoint{12.131629in}{2.528416in}}%
\pgfpathquadraticcurveto{\pgfqpoint{12.148482in}{2.525193in}}{\pgfqpoint{12.156503in}{2.517100in}}%
\pgfpathquadraticcurveto{\pgfqpoint{12.164524in}{2.509008in}}{\pgfqpoint{12.164524in}{2.495449in}}%
\pgfpathquadraticcurveto{\pgfqpoint{12.164524in}{2.479240in}}{\pgfqpoint{12.153663in}{2.470885in}}%
\pgfpathquadraticcurveto{\pgfqpoint{12.142825in}{2.462530in}}{\pgfqpoint{12.121866in}{2.462530in}}%
\pgfpathquadraticcurveto{\pgfqpoint{12.113964in}{2.462530in}}{\pgfqpoint{12.105036in}{2.464321in}}%
\pgfpathquadraticcurveto{\pgfqpoint{12.096132in}{2.466111in}}{\pgfqpoint{12.086584in}{2.469620in}}%
\pgfpathlineto{\pgfqpoint{12.086584in}{2.485137in}}%
\pgfpathquadraticcurveto{\pgfqpoint{12.095750in}{2.480004in}}{\pgfqpoint{12.104559in}{2.477378in}}%
\pgfpathquadraticcurveto{\pgfqpoint{12.113367in}{2.474776in}}{\pgfqpoint{12.121866in}{2.474776in}}%
\pgfpathquadraticcurveto{\pgfqpoint{12.134756in}{2.474776in}}{\pgfqpoint{12.141775in}{2.479837in}}%
\pgfpathquadraticcurveto{\pgfqpoint{12.148793in}{2.484922in}}{\pgfqpoint{12.148793in}{2.494327in}}%
\pgfpathquadraticcurveto{\pgfqpoint{12.148793in}{2.502515in}}{\pgfqpoint{12.143756in}{2.507146in}}%
\pgfpathquadraticcurveto{\pgfqpoint{12.138719in}{2.511777in}}{\pgfqpoint{12.127237in}{2.514093in}}%
\pgfpathlineto{\pgfqpoint{12.118046in}{2.515883in}}%
\pgfpathquadraticcurveto{\pgfqpoint{12.101193in}{2.519225in}}{\pgfqpoint{12.093650in}{2.526387in}}%
\pgfpathquadraticcurveto{\pgfqpoint{12.086130in}{2.533548in}}{\pgfqpoint{12.086130in}{2.546319in}}%
\pgfpathquadraticcurveto{\pgfqpoint{12.086130in}{2.561096in}}{\pgfqpoint{12.096538in}{2.569594in}}%
\pgfpathquadraticcurveto{\pgfqpoint{12.106946in}{2.578092in}}{\pgfqpoint{12.125208in}{2.578092in}}%
\pgfpathquadraticcurveto{\pgfqpoint{12.133061in}{2.578092in}}{\pgfqpoint{12.141178in}{2.576660in}}%
\pgfpathquadraticcurveto{\pgfqpoint{12.149318in}{2.575252in}}{\pgfqpoint{12.157816in}{2.572435in}}%
\pgfpathlineto{\pgfqpoint{12.157816in}{2.572435in}}%
\pgfpathclose%
\pgfpathmoveto{\pgfqpoint{12.188038in}{2.576087in}}%
\pgfpathlineto{\pgfqpoint{12.258459in}{2.576087in}}%
\pgfpathlineto{\pgfqpoint{12.258459in}{2.563387in}}%
\pgfpathlineto{\pgfqpoint{12.203100in}{2.563387in}}%
\pgfpathlineto{\pgfqpoint{12.203100in}{2.530421in}}%
\pgfpathlineto{\pgfqpoint{12.256143in}{2.530421in}}%
\pgfpathlineto{\pgfqpoint{12.256143in}{2.517745in}}%
\pgfpathlineto{\pgfqpoint{12.203100in}{2.517745in}}%
\pgfpathlineto{\pgfqpoint{12.203100in}{2.477378in}}%
\pgfpathlineto{\pgfqpoint{12.259795in}{2.477378in}}%
\pgfpathlineto{\pgfqpoint{12.259795in}{2.464702in}}%
\pgfpathlineto{\pgfqpoint{12.188038in}{2.464702in}}%
\pgfpathlineto{\pgfqpoint{12.188038in}{2.576087in}}%
\pgfpathlineto{\pgfqpoint{12.188038in}{2.576087in}}%
\pgfpathclose%
\pgfpathmoveto{\pgfqpoint{12.337235in}{2.545794in}}%
\pgfpathlineto{\pgfqpoint{12.337235in}{2.532808in}}%
\pgfpathquadraticcurveto{\pgfqpoint{12.331434in}{2.535792in}}{\pgfqpoint{12.325156in}{2.537272in}}%
\pgfpathquadraticcurveto{\pgfqpoint{12.318901in}{2.538776in}}{\pgfqpoint{12.312169in}{2.538776in}}%
\pgfpathquadraticcurveto{\pgfqpoint{12.301952in}{2.538776in}}{\pgfqpoint{12.296844in}{2.535649in}}%
\pgfpathquadraticcurveto{\pgfqpoint{12.291735in}{2.532521in}}{\pgfqpoint{12.291735in}{2.526243in}}%
\pgfpathquadraticcurveto{\pgfqpoint{12.291735in}{2.521469in}}{\pgfqpoint{12.295388in}{2.518748in}}%
\pgfpathquadraticcurveto{\pgfqpoint{12.299040in}{2.516026in}}{\pgfqpoint{12.310093in}{2.513567in}}%
\pgfpathlineto{\pgfqpoint{12.314795in}{2.512517in}}%
\pgfpathquadraticcurveto{\pgfqpoint{12.329405in}{2.509390in}}{\pgfqpoint{12.335564in}{2.503685in}}%
\pgfpathquadraticcurveto{\pgfqpoint{12.341722in}{2.497979in}}{\pgfqpoint{12.341722in}{2.487762in}}%
\pgfpathquadraticcurveto{\pgfqpoint{12.341722in}{2.476113in}}{\pgfqpoint{12.332508in}{2.469310in}}%
\pgfpathquadraticcurveto{\pgfqpoint{12.323294in}{2.462530in}}{\pgfqpoint{12.307180in}{2.462530in}}%
\pgfpathquadraticcurveto{\pgfqpoint{12.300472in}{2.462530in}}{\pgfqpoint{12.293192in}{2.463843in}}%
\pgfpathquadraticcurveto{\pgfqpoint{12.285911in}{2.465156in}}{\pgfqpoint{12.277866in}{2.467758in}}%
\pgfpathlineto{\pgfqpoint{12.277866in}{2.481938in}}%
\pgfpathquadraticcurveto{\pgfqpoint{12.285481in}{2.477975in}}{\pgfqpoint{12.292857in}{2.475994in}}%
\pgfpathquadraticcurveto{\pgfqpoint{12.300234in}{2.474036in}}{\pgfqpoint{12.307491in}{2.474036in}}%
\pgfpathquadraticcurveto{\pgfqpoint{12.317182in}{2.474036in}}{\pgfqpoint{12.322386in}{2.477354in}}%
\pgfpathquadraticcurveto{\pgfqpoint{12.327614in}{2.480673in}}{\pgfqpoint{12.327614in}{2.486712in}}%
\pgfpathquadraticcurveto{\pgfqpoint{12.327614in}{2.492298in}}{\pgfqpoint{12.323843in}{2.495282in}}%
\pgfpathquadraticcurveto{\pgfqpoint{12.320095in}{2.498266in}}{\pgfqpoint{12.307324in}{2.501035in}}%
\pgfpathlineto{\pgfqpoint{12.302549in}{2.502157in}}%
\pgfpathquadraticcurveto{\pgfqpoint{12.289802in}{2.504831in}}{\pgfqpoint{12.284120in}{2.510393in}}%
\pgfpathquadraticcurveto{\pgfqpoint{12.278463in}{2.515955in}}{\pgfqpoint{12.278463in}{2.525646in}}%
\pgfpathquadraticcurveto{\pgfqpoint{12.278463in}{2.537439in}}{\pgfqpoint{12.286818in}{2.543837in}}%
\pgfpathquadraticcurveto{\pgfqpoint{12.295173in}{2.550258in}}{\pgfqpoint{12.310546in}{2.550258in}}%
\pgfpathquadraticcurveto{\pgfqpoint{12.318137in}{2.550258in}}{\pgfqpoint{12.324845in}{2.549136in}}%
\pgfpathquadraticcurveto{\pgfqpoint{12.331577in}{2.548038in}}{\pgfqpoint{12.337235in}{2.545794in}}%
\pgfpathlineto{\pgfqpoint{12.337235in}{2.545794in}}%
\pgfpathclose%
\pgfpathmoveto{\pgfqpoint{12.361416in}{2.580623in}}%
\pgfpathlineto{\pgfqpoint{12.373352in}{2.580623in}}%
\pgfpathquadraticcurveto{\pgfqpoint{12.384524in}{2.563029in}}{\pgfqpoint{12.390086in}{2.546152in}}%
\pgfpathquadraticcurveto{\pgfqpoint{12.395648in}{2.529299in}}{\pgfqpoint{12.395648in}{2.512660in}}%
\pgfpathquadraticcurveto{\pgfqpoint{12.395648in}{2.495950in}}{\pgfqpoint{12.390086in}{2.479025in}}%
\pgfpathquadraticcurveto{\pgfqpoint{12.384524in}{2.462100in}}{\pgfqpoint{12.373352in}{2.444555in}}%
\pgfpathlineto{\pgfqpoint{12.361416in}{2.444555in}}%
\pgfpathquadraticcurveto{\pgfqpoint{12.371323in}{2.461647in}}{\pgfqpoint{12.376217in}{2.478548in}}%
\pgfpathquadraticcurveto{\pgfqpoint{12.381110in}{2.495449in}}{\pgfqpoint{12.381110in}{2.512660in}}%
\pgfpathquadraticcurveto{\pgfqpoint{12.381110in}{2.529896in}}{\pgfqpoint{12.376217in}{2.546677in}}%
\pgfpathquadraticcurveto{\pgfqpoint{12.371323in}{2.563483in}}{\pgfqpoint{12.361416in}{2.580623in}}%
\pgfpathlineto{\pgfqpoint{12.361416in}{2.580623in}}%
\pgfpathclose%
\pgfusepath{fill}%
\end{pgfscope}%
\begin{pgfscope}%
\pgfpathrectangle{\pgfqpoint{7.904091in}{2.928000in}}{\pgfqpoint{5.664000in}{5.664000in}}%
\pgfusepath{clip}%
\pgfsetrectcap%
\pgfsetroundjoin%
\pgfsetlinewidth{0.602250pt}%
\definecolor{currentstroke}{rgb}{0.898039,0.913725,0.929412}%
\pgfsetstrokecolor{currentstroke}%
\pgfsetdash{}{0pt}%
\pgfpathmoveto{\pgfqpoint{7.904091in}{2.928000in}}%
\pgfpathlineto{\pgfqpoint{13.568091in}{2.928000in}}%
\pgfusepath{stroke}%
\end{pgfscope}%
\begin{pgfscope}%
\pgfsetbuttcap%
\pgfsetroundjoin%
\definecolor{currentfill}{rgb}{0.090196,0.168627,0.227451}%
\pgfsetfillcolor{currentfill}%
\pgfsetlinewidth{0.803000pt}%
\definecolor{currentstroke}{rgb}{0.090196,0.168627,0.227451}%
\pgfsetstrokecolor{currentstroke}%
\pgfsetdash{}{0pt}%
\pgfsys@defobject{currentmarker}{\pgfqpoint{-0.048611in}{0.000000in}}{\pgfqpoint{-0.000000in}{0.000000in}}{%
\pgfpathmoveto{\pgfqpoint{-0.000000in}{0.000000in}}%
\pgfpathlineto{\pgfqpoint{-0.048611in}{0.000000in}}%
\pgfusepath{stroke,fill}%
}%
\begin{pgfscope}%
\pgfsys@transformshift{7.904091in}{2.928000in}%
\pgfsys@useobject{currentmarker}{}%
\end{pgfscope}%
\end{pgfscope}%
\begin{pgfscope}%
\pgfsetbuttcap%
\pgfsetroundjoin%
\definecolor{currentfill}{rgb}{0.090196,0.168627,0.227451}%
\pgfsetfillcolor{currentfill}%
\pgfsetlinewidth{0.000000pt}%
\definecolor{currentstroke}{rgb}{0.090196,0.168627,0.227451}%
\pgfsetstrokecolor{currentstroke}%
\pgfsetdash{}{0pt}%
\pgfpathmoveto{\pgfqpoint{7.510111in}{2.924890in}}%
\pgfpathlineto{\pgfqpoint{7.588353in}{2.924890in}}%
\pgfpathlineto{\pgfqpoint{7.588353in}{2.914519in}}%
\pgfpathlineto{\pgfqpoint{7.510111in}{2.914519in}}%
\pgfpathlineto{\pgfqpoint{7.510111in}{2.924890in}}%
\pgfpathlineto{\pgfqpoint{7.510111in}{2.924890in}}%
\pgfpathclose%
\pgfpathmoveto{\pgfqpoint{7.615111in}{2.971647in}}%
\pgfpathlineto{\pgfqpoint{7.663509in}{2.971647in}}%
\pgfpathlineto{\pgfqpoint{7.663509in}{2.961257in}}%
\pgfpathlineto{\pgfqpoint{7.626400in}{2.961257in}}%
\pgfpathlineto{\pgfqpoint{7.626400in}{2.938933in}}%
\pgfpathquadraticcurveto{\pgfqpoint{7.629076in}{2.939851in}}{\pgfqpoint{7.631751in}{2.940300in}}%
\pgfpathquadraticcurveto{\pgfqpoint{7.634447in}{2.940749in}}{\pgfqpoint{7.637142in}{2.940749in}}%
\pgfpathquadraticcurveto{\pgfqpoint{7.652396in}{2.940749in}}{\pgfqpoint{7.661302in}{2.932390in}}%
\pgfpathquadraticcurveto{\pgfqpoint{7.670228in}{2.924030in}}{\pgfqpoint{7.670228in}{2.909753in}}%
\pgfpathquadraticcurveto{\pgfqpoint{7.670228in}{2.895046in}}{\pgfqpoint{7.661068in}{2.886882in}}%
\pgfpathquadraticcurveto{\pgfqpoint{7.651908in}{2.878737in}}{\pgfqpoint{7.635248in}{2.878737in}}%
\pgfpathquadraticcurveto{\pgfqpoint{7.629505in}{2.878737in}}{\pgfqpoint{7.623548in}{2.879714in}}%
\pgfpathquadraticcurveto{\pgfqpoint{7.617611in}{2.880690in}}{\pgfqpoint{7.611263in}{2.882644in}}%
\pgfpathlineto{\pgfqpoint{7.611263in}{2.895046in}}%
\pgfpathquadraticcurveto{\pgfqpoint{7.616751in}{2.892058in}}{\pgfqpoint{7.622611in}{2.890593in}}%
\pgfpathquadraticcurveto{\pgfqpoint{7.628470in}{2.889128in}}{\pgfqpoint{7.634994in}{2.889128in}}%
\pgfpathquadraticcurveto{\pgfqpoint{7.645560in}{2.889128in}}{\pgfqpoint{7.651712in}{2.894675in}}%
\pgfpathquadraticcurveto{\pgfqpoint{7.657884in}{2.900222in}}{\pgfqpoint{7.657884in}{2.909753in}}%
\pgfpathquadraticcurveto{\pgfqpoint{7.657884in}{2.919265in}}{\pgfqpoint{7.651712in}{2.924812in}}%
\pgfpathquadraticcurveto{\pgfqpoint{7.645560in}{2.930378in}}{\pgfqpoint{7.634994in}{2.930378in}}%
\pgfpathquadraticcurveto{\pgfqpoint{7.630052in}{2.930378in}}{\pgfqpoint{7.625130in}{2.929284in}}%
\pgfpathquadraticcurveto{\pgfqpoint{7.620228in}{2.928190in}}{\pgfqpoint{7.615111in}{2.925866in}}%
\pgfpathlineto{\pgfqpoint{7.615111in}{2.971647in}}%
\pgfpathlineto{\pgfqpoint{7.615111in}{2.971647in}}%
\pgfpathclose%
\pgfpathmoveto{\pgfqpoint{7.694505in}{2.896022in}}%
\pgfpathlineto{\pgfqpoint{7.707396in}{2.896022in}}%
\pgfpathlineto{\pgfqpoint{7.707396in}{2.880515in}}%
\pgfpathlineto{\pgfqpoint{7.694505in}{2.880515in}}%
\pgfpathlineto{\pgfqpoint{7.694505in}{2.896022in}}%
\pgfpathlineto{\pgfqpoint{7.694505in}{2.896022in}}%
\pgfpathclose%
\pgfpathmoveto{\pgfqpoint{7.760599in}{2.963522in}}%
\pgfpathquadraticcurveto{\pgfqpoint{7.751087in}{2.963522in}}{\pgfqpoint{7.746283in}{2.954147in}}%
\pgfpathquadraticcurveto{\pgfqpoint{7.741498in}{2.944792in}}{\pgfqpoint{7.741498in}{2.925983in}}%
\pgfpathquadraticcurveto{\pgfqpoint{7.741498in}{2.907253in}}{\pgfqpoint{7.746283in}{2.897878in}}%
\pgfpathquadraticcurveto{\pgfqpoint{7.751087in}{2.888503in}}{\pgfqpoint{7.760599in}{2.888503in}}%
\pgfpathquadraticcurveto{\pgfqpoint{7.770189in}{2.888503in}}{\pgfqpoint{7.774974in}{2.897878in}}%
\pgfpathquadraticcurveto{\pgfqpoint{7.779779in}{2.907253in}}{\pgfqpoint{7.779779in}{2.925983in}}%
\pgfpathquadraticcurveto{\pgfqpoint{7.779779in}{2.944792in}}{\pgfqpoint{7.774974in}{2.954147in}}%
\pgfpathquadraticcurveto{\pgfqpoint{7.770189in}{2.963522in}}{\pgfqpoint{7.760599in}{2.963522in}}%
\pgfpathlineto{\pgfqpoint{7.760599in}{2.963522in}}%
\pgfpathclose%
\pgfpathmoveto{\pgfqpoint{7.760599in}{2.973288in}}%
\pgfpathquadraticcurveto{\pgfqpoint{7.775931in}{2.973288in}}{\pgfqpoint{7.784017in}{2.961159in}}%
\pgfpathquadraticcurveto{\pgfqpoint{7.792103in}{2.949050in}}{\pgfqpoint{7.792103in}{2.925983in}}%
\pgfpathquadraticcurveto{\pgfqpoint{7.792103in}{2.902976in}}{\pgfqpoint{7.784017in}{2.890847in}}%
\pgfpathquadraticcurveto{\pgfqpoint{7.775931in}{2.878737in}}{\pgfqpoint{7.760599in}{2.878737in}}%
\pgfpathquadraticcurveto{\pgfqpoint{7.745287in}{2.878737in}}{\pgfqpoint{7.737201in}{2.890847in}}%
\pgfpathquadraticcurveto{\pgfqpoint{7.729115in}{2.902976in}}{\pgfqpoint{7.729115in}{2.925983in}}%
\pgfpathquadraticcurveto{\pgfqpoint{7.729115in}{2.949050in}}{\pgfqpoint{7.737201in}{2.961159in}}%
\pgfpathquadraticcurveto{\pgfqpoint{7.745287in}{2.973288in}}{\pgfqpoint{7.760599in}{2.973288in}}%
\pgfpathlineto{\pgfqpoint{7.760599in}{2.973288in}}%
\pgfpathclose%
\pgfusepath{fill}%
\end{pgfscope}%
\begin{pgfscope}%
\pgfpathrectangle{\pgfqpoint{7.904091in}{2.928000in}}{\pgfqpoint{5.664000in}{5.664000in}}%
\pgfusepath{clip}%
\pgfsetrectcap%
\pgfsetroundjoin%
\pgfsetlinewidth{0.602250pt}%
\definecolor{currentstroke}{rgb}{0.898039,0.913725,0.929412}%
\pgfsetstrokecolor{currentstroke}%
\pgfsetdash{}{0pt}%
\pgfpathmoveto{\pgfqpoint{7.904091in}{4.344000in}}%
\pgfpathlineto{\pgfqpoint{13.568091in}{4.344000in}}%
\pgfusepath{stroke}%
\end{pgfscope}%
\begin{pgfscope}%
\pgfsetbuttcap%
\pgfsetroundjoin%
\definecolor{currentfill}{rgb}{0.090196,0.168627,0.227451}%
\pgfsetfillcolor{currentfill}%
\pgfsetlinewidth{0.803000pt}%
\definecolor{currentstroke}{rgb}{0.090196,0.168627,0.227451}%
\pgfsetstrokecolor{currentstroke}%
\pgfsetdash{}{0pt}%
\pgfsys@defobject{currentmarker}{\pgfqpoint{-0.048611in}{0.000000in}}{\pgfqpoint{-0.000000in}{0.000000in}}{%
\pgfpathmoveto{\pgfqpoint{-0.000000in}{0.000000in}}%
\pgfpathlineto{\pgfqpoint{-0.048611in}{0.000000in}}%
\pgfusepath{stroke,fill}%
}%
\begin{pgfscope}%
\pgfsys@transformshift{7.904091in}{4.344000in}%
\pgfsys@useobject{currentmarker}{}%
\end{pgfscope}%
\end{pgfscope}%
\begin{pgfscope}%
\pgfsetbuttcap%
\pgfsetroundjoin%
\definecolor{currentfill}{rgb}{0.090196,0.168627,0.227451}%
\pgfsetfillcolor{currentfill}%
\pgfsetlinewidth{0.000000pt}%
\definecolor{currentstroke}{rgb}{0.090196,0.168627,0.227451}%
\pgfsetstrokecolor{currentstroke}%
\pgfsetdash{}{0pt}%
\pgfpathmoveto{\pgfqpoint{7.510111in}{4.340890in}}%
\pgfpathlineto{\pgfqpoint{7.588353in}{4.340890in}}%
\pgfpathlineto{\pgfqpoint{7.588353in}{4.330519in}}%
\pgfpathlineto{\pgfqpoint{7.510111in}{4.330519in}}%
\pgfpathlineto{\pgfqpoint{7.510111in}{4.340890in}}%
\pgfpathlineto{\pgfqpoint{7.510111in}{4.340890in}}%
\pgfpathclose%
\pgfpathmoveto{\pgfqpoint{7.625599in}{4.306886in}}%
\pgfpathlineto{\pgfqpoint{7.668626in}{4.306886in}}%
\pgfpathlineto{\pgfqpoint{7.668626in}{4.296515in}}%
\pgfpathlineto{\pgfqpoint{7.610775in}{4.296515in}}%
\pgfpathlineto{\pgfqpoint{7.610775in}{4.306886in}}%
\pgfpathquadraticcurveto{\pgfqpoint{7.617787in}{4.314151in}}{\pgfqpoint{7.629896in}{4.326378in}}%
\pgfpathquadraticcurveto{\pgfqpoint{7.642025in}{4.338624in}}{\pgfqpoint{7.645130in}{4.342179in}}%
\pgfpathquadraticcurveto{\pgfqpoint{7.651048in}{4.348819in}}{\pgfqpoint{7.653392in}{4.353429in}}%
\pgfpathquadraticcurveto{\pgfqpoint{7.655755in}{4.358038in}}{\pgfqpoint{7.655755in}{4.362491in}}%
\pgfpathquadraticcurveto{\pgfqpoint{7.655755in}{4.369757in}}{\pgfqpoint{7.650658in}{4.374327in}}%
\pgfpathquadraticcurveto{\pgfqpoint{7.645560in}{4.378917in}}{\pgfqpoint{7.637376in}{4.378917in}}%
\pgfpathquadraticcurveto{\pgfqpoint{7.631576in}{4.378917in}}{\pgfqpoint{7.625130in}{4.376905in}}%
\pgfpathquadraticcurveto{\pgfqpoint{7.618705in}{4.374894in}}{\pgfqpoint{7.611380in}{4.370792in}}%
\pgfpathlineto{\pgfqpoint{7.611380in}{4.383253in}}%
\pgfpathquadraticcurveto{\pgfqpoint{7.618822in}{4.386241in}}{\pgfqpoint{7.625287in}{4.387765in}}%
\pgfpathquadraticcurveto{\pgfqpoint{7.631771in}{4.389288in}}{\pgfqpoint{7.637142in}{4.389288in}}%
\pgfpathquadraticcurveto{\pgfqpoint{7.651302in}{4.389288in}}{\pgfqpoint{7.659720in}{4.382198in}}%
\pgfpathquadraticcurveto{\pgfqpoint{7.668138in}{4.375128in}}{\pgfqpoint{7.668138in}{4.363292in}}%
\pgfpathquadraticcurveto{\pgfqpoint{7.668138in}{4.357667in}}{\pgfqpoint{7.666029in}{4.352628in}}%
\pgfpathquadraticcurveto{\pgfqpoint{7.663939in}{4.347608in}}{\pgfqpoint{7.658373in}{4.340772in}}%
\pgfpathquadraticcurveto{\pgfqpoint{7.656849in}{4.338995in}}{\pgfqpoint{7.648666in}{4.330538in}}%
\pgfpathquadraticcurveto{\pgfqpoint{7.640501in}{4.322081in}}{\pgfqpoint{7.625599in}{4.306886in}}%
\pgfpathlineto{\pgfqpoint{7.625599in}{4.306886in}}%
\pgfpathclose%
\pgfpathmoveto{\pgfqpoint{7.694505in}{4.312022in}}%
\pgfpathlineto{\pgfqpoint{7.707396in}{4.312022in}}%
\pgfpathlineto{\pgfqpoint{7.707396in}{4.296515in}}%
\pgfpathlineto{\pgfqpoint{7.694505in}{4.296515in}}%
\pgfpathlineto{\pgfqpoint{7.694505in}{4.312022in}}%
\pgfpathlineto{\pgfqpoint{7.694505in}{4.312022in}}%
\pgfpathclose%
\pgfpathmoveto{\pgfqpoint{7.734369in}{4.387647in}}%
\pgfpathlineto{\pgfqpoint{7.782767in}{4.387647in}}%
\pgfpathlineto{\pgfqpoint{7.782767in}{4.377257in}}%
\pgfpathlineto{\pgfqpoint{7.745658in}{4.377257in}}%
\pgfpathlineto{\pgfqpoint{7.745658in}{4.354933in}}%
\pgfpathquadraticcurveto{\pgfqpoint{7.748334in}{4.355851in}}{\pgfqpoint{7.751009in}{4.356300in}}%
\pgfpathquadraticcurveto{\pgfqpoint{7.753705in}{4.356749in}}{\pgfqpoint{7.756400in}{4.356749in}}%
\pgfpathquadraticcurveto{\pgfqpoint{7.771654in}{4.356749in}}{\pgfqpoint{7.780560in}{4.348390in}}%
\pgfpathquadraticcurveto{\pgfqpoint{7.789486in}{4.340030in}}{\pgfqpoint{7.789486in}{4.325753in}}%
\pgfpathquadraticcurveto{\pgfqpoint{7.789486in}{4.311046in}}{\pgfqpoint{7.780326in}{4.302882in}}%
\pgfpathquadraticcurveto{\pgfqpoint{7.771166in}{4.294737in}}{\pgfqpoint{7.754505in}{4.294737in}}%
\pgfpathquadraticcurveto{\pgfqpoint{7.748763in}{4.294737in}}{\pgfqpoint{7.742806in}{4.295714in}}%
\pgfpathquadraticcurveto{\pgfqpoint{7.736869in}{4.296690in}}{\pgfqpoint{7.730521in}{4.298644in}}%
\pgfpathlineto{\pgfqpoint{7.730521in}{4.311046in}}%
\pgfpathquadraticcurveto{\pgfqpoint{7.736009in}{4.308058in}}{\pgfqpoint{7.741869in}{4.306593in}}%
\pgfpathquadraticcurveto{\pgfqpoint{7.747728in}{4.305128in}}{\pgfqpoint{7.754251in}{4.305128in}}%
\pgfpathquadraticcurveto{\pgfqpoint{7.764818in}{4.305128in}}{\pgfqpoint{7.770970in}{4.310675in}}%
\pgfpathquadraticcurveto{\pgfqpoint{7.777142in}{4.316222in}}{\pgfqpoint{7.777142in}{4.325753in}}%
\pgfpathquadraticcurveto{\pgfqpoint{7.777142in}{4.335265in}}{\pgfqpoint{7.770970in}{4.340812in}}%
\pgfpathquadraticcurveto{\pgfqpoint{7.764818in}{4.346378in}}{\pgfqpoint{7.754251in}{4.346378in}}%
\pgfpathquadraticcurveto{\pgfqpoint{7.749310in}{4.346378in}}{\pgfqpoint{7.744388in}{4.345284in}}%
\pgfpathquadraticcurveto{\pgfqpoint{7.739486in}{4.344190in}}{\pgfqpoint{7.734369in}{4.341866in}}%
\pgfpathlineto{\pgfqpoint{7.734369in}{4.387647in}}%
\pgfpathlineto{\pgfqpoint{7.734369in}{4.387647in}}%
\pgfpathclose%
\pgfusepath{fill}%
\end{pgfscope}%
\begin{pgfscope}%
\pgfpathrectangle{\pgfqpoint{7.904091in}{2.928000in}}{\pgfqpoint{5.664000in}{5.664000in}}%
\pgfusepath{clip}%
\pgfsetrectcap%
\pgfsetroundjoin%
\pgfsetlinewidth{0.602250pt}%
\definecolor{currentstroke}{rgb}{0.898039,0.913725,0.929412}%
\pgfsetstrokecolor{currentstroke}%
\pgfsetdash{}{0pt}%
\pgfpathmoveto{\pgfqpoint{7.904091in}{5.760000in}}%
\pgfpathlineto{\pgfqpoint{13.568091in}{5.760000in}}%
\pgfusepath{stroke}%
\end{pgfscope}%
\begin{pgfscope}%
\pgfsetbuttcap%
\pgfsetroundjoin%
\definecolor{currentfill}{rgb}{0.090196,0.168627,0.227451}%
\pgfsetfillcolor{currentfill}%
\pgfsetlinewidth{0.803000pt}%
\definecolor{currentstroke}{rgb}{0.090196,0.168627,0.227451}%
\pgfsetstrokecolor{currentstroke}%
\pgfsetdash{}{0pt}%
\pgfsys@defobject{currentmarker}{\pgfqpoint{-0.048611in}{0.000000in}}{\pgfqpoint{-0.000000in}{0.000000in}}{%
\pgfpathmoveto{\pgfqpoint{-0.000000in}{0.000000in}}%
\pgfpathlineto{\pgfqpoint{-0.048611in}{0.000000in}}%
\pgfusepath{stroke,fill}%
}%
\begin{pgfscope}%
\pgfsys@transformshift{7.904091in}{5.760000in}%
\pgfsys@useobject{currentmarker}{}%
\end{pgfscope}%
\end{pgfscope}%
\begin{pgfscope}%
\pgfsetbuttcap%
\pgfsetroundjoin%
\definecolor{currentfill}{rgb}{0.090196,0.168627,0.227451}%
\pgfsetfillcolor{currentfill}%
\pgfsetlinewidth{0.000000pt}%
\definecolor{currentstroke}{rgb}{0.090196,0.168627,0.227451}%
\pgfsetstrokecolor{currentstroke}%
\pgfsetdash{}{0pt}%
\pgfpathmoveto{\pgfqpoint{7.646595in}{5.795522in}}%
\pgfpathquadraticcurveto{\pgfqpoint{7.637084in}{5.795522in}}{\pgfqpoint{7.632279in}{5.786147in}}%
\pgfpathquadraticcurveto{\pgfqpoint{7.627494in}{5.776792in}}{\pgfqpoint{7.627494in}{5.757983in}}%
\pgfpathquadraticcurveto{\pgfqpoint{7.627494in}{5.739253in}}{\pgfqpoint{7.632279in}{5.729878in}}%
\pgfpathquadraticcurveto{\pgfqpoint{7.637084in}{5.720503in}}{\pgfqpoint{7.646595in}{5.720503in}}%
\pgfpathquadraticcurveto{\pgfqpoint{7.656185in}{5.720503in}}{\pgfqpoint{7.660970in}{5.729878in}}%
\pgfpathquadraticcurveto{\pgfqpoint{7.665775in}{5.739253in}}{\pgfqpoint{7.665775in}{5.757983in}}%
\pgfpathquadraticcurveto{\pgfqpoint{7.665775in}{5.776792in}}{\pgfqpoint{7.660970in}{5.786147in}}%
\pgfpathquadraticcurveto{\pgfqpoint{7.656185in}{5.795522in}}{\pgfqpoint{7.646595in}{5.795522in}}%
\pgfpathlineto{\pgfqpoint{7.646595in}{5.795522in}}%
\pgfpathclose%
\pgfpathmoveto{\pgfqpoint{7.646595in}{5.805288in}}%
\pgfpathquadraticcurveto{\pgfqpoint{7.661927in}{5.805288in}}{\pgfqpoint{7.670013in}{5.793159in}}%
\pgfpathquadraticcurveto{\pgfqpoint{7.678099in}{5.781050in}}{\pgfqpoint{7.678099in}{5.757983in}}%
\pgfpathquadraticcurveto{\pgfqpoint{7.678099in}{5.734976in}}{\pgfqpoint{7.670013in}{5.722847in}}%
\pgfpathquadraticcurveto{\pgfqpoint{7.661927in}{5.710737in}}{\pgfqpoint{7.646595in}{5.710737in}}%
\pgfpathquadraticcurveto{\pgfqpoint{7.631283in}{5.710737in}}{\pgfqpoint{7.623197in}{5.722847in}}%
\pgfpathquadraticcurveto{\pgfqpoint{7.615111in}{5.734976in}}{\pgfqpoint{7.615111in}{5.757983in}}%
\pgfpathquadraticcurveto{\pgfqpoint{7.615111in}{5.781050in}}{\pgfqpoint{7.623197in}{5.793159in}}%
\pgfpathquadraticcurveto{\pgfqpoint{7.631283in}{5.805288in}}{\pgfqpoint{7.646595in}{5.805288in}}%
\pgfpathlineto{\pgfqpoint{7.646595in}{5.805288in}}%
\pgfpathclose%
\pgfpathmoveto{\pgfqpoint{7.699759in}{5.728022in}}%
\pgfpathlineto{\pgfqpoint{7.712650in}{5.728022in}}%
\pgfpathlineto{\pgfqpoint{7.712650in}{5.712515in}}%
\pgfpathlineto{\pgfqpoint{7.699759in}{5.712515in}}%
\pgfpathlineto{\pgfqpoint{7.699759in}{5.728022in}}%
\pgfpathlineto{\pgfqpoint{7.699759in}{5.728022in}}%
\pgfpathclose%
\pgfpathmoveto{\pgfqpoint{7.765853in}{5.795522in}}%
\pgfpathquadraticcurveto{\pgfqpoint{7.756341in}{5.795522in}}{\pgfqpoint{7.751537in}{5.786147in}}%
\pgfpathquadraticcurveto{\pgfqpoint{7.746751in}{5.776792in}}{\pgfqpoint{7.746751in}{5.757983in}}%
\pgfpathquadraticcurveto{\pgfqpoint{7.746751in}{5.739253in}}{\pgfqpoint{7.751537in}{5.729878in}}%
\pgfpathquadraticcurveto{\pgfqpoint{7.756341in}{5.720503in}}{\pgfqpoint{7.765853in}{5.720503in}}%
\pgfpathquadraticcurveto{\pgfqpoint{7.775443in}{5.720503in}}{\pgfqpoint{7.780228in}{5.729878in}}%
\pgfpathquadraticcurveto{\pgfqpoint{7.785033in}{5.739253in}}{\pgfqpoint{7.785033in}{5.757983in}}%
\pgfpathquadraticcurveto{\pgfqpoint{7.785033in}{5.776792in}}{\pgfqpoint{7.780228in}{5.786147in}}%
\pgfpathquadraticcurveto{\pgfqpoint{7.775443in}{5.795522in}}{\pgfqpoint{7.765853in}{5.795522in}}%
\pgfpathlineto{\pgfqpoint{7.765853in}{5.795522in}}%
\pgfpathclose%
\pgfpathmoveto{\pgfqpoint{7.765853in}{5.805288in}}%
\pgfpathquadraticcurveto{\pgfqpoint{7.781185in}{5.805288in}}{\pgfqpoint{7.789271in}{5.793159in}}%
\pgfpathquadraticcurveto{\pgfqpoint{7.797357in}{5.781050in}}{\pgfqpoint{7.797357in}{5.757983in}}%
\pgfpathquadraticcurveto{\pgfqpoint{7.797357in}{5.734976in}}{\pgfqpoint{7.789271in}{5.722847in}}%
\pgfpathquadraticcurveto{\pgfqpoint{7.781185in}{5.710737in}}{\pgfqpoint{7.765853in}{5.710737in}}%
\pgfpathquadraticcurveto{\pgfqpoint{7.750541in}{5.710737in}}{\pgfqpoint{7.742455in}{5.722847in}}%
\pgfpathquadraticcurveto{\pgfqpoint{7.734369in}{5.734976in}}{\pgfqpoint{7.734369in}{5.757983in}}%
\pgfpathquadraticcurveto{\pgfqpoint{7.734369in}{5.781050in}}{\pgfqpoint{7.742455in}{5.793159in}}%
\pgfpathquadraticcurveto{\pgfqpoint{7.750541in}{5.805288in}}{\pgfqpoint{7.765853in}{5.805288in}}%
\pgfpathlineto{\pgfqpoint{7.765853in}{5.805288in}}%
\pgfpathclose%
\pgfusepath{fill}%
\end{pgfscope}%
\begin{pgfscope}%
\pgfpathrectangle{\pgfqpoint{7.904091in}{2.928000in}}{\pgfqpoint{5.664000in}{5.664000in}}%
\pgfusepath{clip}%
\pgfsetrectcap%
\pgfsetroundjoin%
\pgfsetlinewidth{0.602250pt}%
\definecolor{currentstroke}{rgb}{0.898039,0.913725,0.929412}%
\pgfsetstrokecolor{currentstroke}%
\pgfsetdash{}{0pt}%
\pgfpathmoveto{\pgfqpoint{7.904091in}{7.176000in}}%
\pgfpathlineto{\pgfqpoint{13.568091in}{7.176000in}}%
\pgfusepath{stroke}%
\end{pgfscope}%
\begin{pgfscope}%
\pgfsetbuttcap%
\pgfsetroundjoin%
\definecolor{currentfill}{rgb}{0.090196,0.168627,0.227451}%
\pgfsetfillcolor{currentfill}%
\pgfsetlinewidth{0.803000pt}%
\definecolor{currentstroke}{rgb}{0.090196,0.168627,0.227451}%
\pgfsetstrokecolor{currentstroke}%
\pgfsetdash{}{0pt}%
\pgfsys@defobject{currentmarker}{\pgfqpoint{-0.048611in}{0.000000in}}{\pgfqpoint{-0.000000in}{0.000000in}}{%
\pgfpathmoveto{\pgfqpoint{-0.000000in}{0.000000in}}%
\pgfpathlineto{\pgfqpoint{-0.048611in}{0.000000in}}%
\pgfusepath{stroke,fill}%
}%
\begin{pgfscope}%
\pgfsys@transformshift{7.904091in}{7.176000in}%
\pgfsys@useobject{currentmarker}{}%
\end{pgfscope}%
\end{pgfscope}%
\begin{pgfscope}%
\pgfsetbuttcap%
\pgfsetroundjoin%
\definecolor{currentfill}{rgb}{0.090196,0.168627,0.227451}%
\pgfsetfillcolor{currentfill}%
\pgfsetlinewidth{0.000000pt}%
\definecolor{currentstroke}{rgb}{0.090196,0.168627,0.227451}%
\pgfsetstrokecolor{currentstroke}%
\pgfsetdash{}{0pt}%
\pgfpathmoveto{\pgfqpoint{7.630853in}{7.138886in}}%
\pgfpathlineto{\pgfqpoint{7.673880in}{7.138886in}}%
\pgfpathlineto{\pgfqpoint{7.673880in}{7.128515in}}%
\pgfpathlineto{\pgfqpoint{7.616029in}{7.128515in}}%
\pgfpathlineto{\pgfqpoint{7.616029in}{7.138886in}}%
\pgfpathquadraticcurveto{\pgfqpoint{7.623041in}{7.146151in}}{\pgfqpoint{7.635150in}{7.158378in}}%
\pgfpathquadraticcurveto{\pgfqpoint{7.647279in}{7.170624in}}{\pgfqpoint{7.650384in}{7.174179in}}%
\pgfpathquadraticcurveto{\pgfqpoint{7.656302in}{7.180819in}}{\pgfqpoint{7.658646in}{7.185429in}}%
\pgfpathquadraticcurveto{\pgfqpoint{7.661009in}{7.190038in}}{\pgfqpoint{7.661009in}{7.194491in}}%
\pgfpathquadraticcurveto{\pgfqpoint{7.661009in}{7.201757in}}{\pgfqpoint{7.655912in}{7.206327in}}%
\pgfpathquadraticcurveto{\pgfqpoint{7.650814in}{7.210917in}}{\pgfqpoint{7.642630in}{7.210917in}}%
\pgfpathquadraticcurveto{\pgfqpoint{7.636830in}{7.210917in}}{\pgfqpoint{7.630384in}{7.208905in}}%
\pgfpathquadraticcurveto{\pgfqpoint{7.623959in}{7.206894in}}{\pgfqpoint{7.616634in}{7.202792in}}%
\pgfpathlineto{\pgfqpoint{7.616634in}{7.215253in}}%
\pgfpathquadraticcurveto{\pgfqpoint{7.624076in}{7.218241in}}{\pgfqpoint{7.630541in}{7.219765in}}%
\pgfpathquadraticcurveto{\pgfqpoint{7.637025in}{7.221288in}}{\pgfqpoint{7.642396in}{7.221288in}}%
\pgfpathquadraticcurveto{\pgfqpoint{7.656556in}{7.221288in}}{\pgfqpoint{7.664974in}{7.214198in}}%
\pgfpathquadraticcurveto{\pgfqpoint{7.673392in}{7.207128in}}{\pgfqpoint{7.673392in}{7.195292in}}%
\pgfpathquadraticcurveto{\pgfqpoint{7.673392in}{7.189667in}}{\pgfqpoint{7.671283in}{7.184628in}}%
\pgfpathquadraticcurveto{\pgfqpoint{7.669193in}{7.179608in}}{\pgfqpoint{7.663626in}{7.172772in}}%
\pgfpathquadraticcurveto{\pgfqpoint{7.662103in}{7.170995in}}{\pgfqpoint{7.653919in}{7.162538in}}%
\pgfpathquadraticcurveto{\pgfqpoint{7.645755in}{7.154081in}}{\pgfqpoint{7.630853in}{7.138886in}}%
\pgfpathlineto{\pgfqpoint{7.630853in}{7.138886in}}%
\pgfpathclose%
\pgfpathmoveto{\pgfqpoint{7.699759in}{7.144022in}}%
\pgfpathlineto{\pgfqpoint{7.712650in}{7.144022in}}%
\pgfpathlineto{\pgfqpoint{7.712650in}{7.128515in}}%
\pgfpathlineto{\pgfqpoint{7.699759in}{7.128515in}}%
\pgfpathlineto{\pgfqpoint{7.699759in}{7.144022in}}%
\pgfpathlineto{\pgfqpoint{7.699759in}{7.144022in}}%
\pgfpathclose%
\pgfpathmoveto{\pgfqpoint{7.739623in}{7.219647in}}%
\pgfpathlineto{\pgfqpoint{7.788021in}{7.219647in}}%
\pgfpathlineto{\pgfqpoint{7.788021in}{7.209257in}}%
\pgfpathlineto{\pgfqpoint{7.750912in}{7.209257in}}%
\pgfpathlineto{\pgfqpoint{7.750912in}{7.186933in}}%
\pgfpathquadraticcurveto{\pgfqpoint{7.753587in}{7.187851in}}{\pgfqpoint{7.756263in}{7.188300in}}%
\pgfpathquadraticcurveto{\pgfqpoint{7.758959in}{7.188749in}}{\pgfqpoint{7.761654in}{7.188749in}}%
\pgfpathquadraticcurveto{\pgfqpoint{7.776908in}{7.188749in}}{\pgfqpoint{7.785814in}{7.180390in}}%
\pgfpathquadraticcurveto{\pgfqpoint{7.794740in}{7.172030in}}{\pgfqpoint{7.794740in}{7.157753in}}%
\pgfpathquadraticcurveto{\pgfqpoint{7.794740in}{7.143046in}}{\pgfqpoint{7.785580in}{7.134882in}}%
\pgfpathquadraticcurveto{\pgfqpoint{7.776419in}{7.126737in}}{\pgfqpoint{7.759759in}{7.126737in}}%
\pgfpathquadraticcurveto{\pgfqpoint{7.754017in}{7.126737in}}{\pgfqpoint{7.748060in}{7.127714in}}%
\pgfpathquadraticcurveto{\pgfqpoint{7.742123in}{7.128690in}}{\pgfqpoint{7.735775in}{7.130644in}}%
\pgfpathlineto{\pgfqpoint{7.735775in}{7.143046in}}%
\pgfpathquadraticcurveto{\pgfqpoint{7.741263in}{7.140058in}}{\pgfqpoint{7.747123in}{7.138593in}}%
\pgfpathquadraticcurveto{\pgfqpoint{7.752982in}{7.137128in}}{\pgfqpoint{7.759505in}{7.137128in}}%
\pgfpathquadraticcurveto{\pgfqpoint{7.770072in}{7.137128in}}{\pgfqpoint{7.776224in}{7.142675in}}%
\pgfpathquadraticcurveto{\pgfqpoint{7.782396in}{7.148222in}}{\pgfqpoint{7.782396in}{7.157753in}}%
\pgfpathquadraticcurveto{\pgfqpoint{7.782396in}{7.167265in}}{\pgfqpoint{7.776224in}{7.172812in}}%
\pgfpathquadraticcurveto{\pgfqpoint{7.770072in}{7.178378in}}{\pgfqpoint{7.759505in}{7.178378in}}%
\pgfpathquadraticcurveto{\pgfqpoint{7.754564in}{7.178378in}}{\pgfqpoint{7.749642in}{7.177284in}}%
\pgfpathquadraticcurveto{\pgfqpoint{7.744740in}{7.176190in}}{\pgfqpoint{7.739623in}{7.173866in}}%
\pgfpathlineto{\pgfqpoint{7.739623in}{7.219647in}}%
\pgfpathlineto{\pgfqpoint{7.739623in}{7.219647in}}%
\pgfpathclose%
\pgfusepath{fill}%
\end{pgfscope}%
\begin{pgfscope}%
\pgfpathrectangle{\pgfqpoint{7.904091in}{2.928000in}}{\pgfqpoint{5.664000in}{5.664000in}}%
\pgfusepath{clip}%
\pgfsetrectcap%
\pgfsetroundjoin%
\pgfsetlinewidth{0.602250pt}%
\definecolor{currentstroke}{rgb}{0.898039,0.913725,0.929412}%
\pgfsetstrokecolor{currentstroke}%
\pgfsetdash{}{0pt}%
\pgfpathmoveto{\pgfqpoint{7.904091in}{8.592000in}}%
\pgfpathlineto{\pgfqpoint{13.568091in}{8.592000in}}%
\pgfusepath{stroke}%
\end{pgfscope}%
\begin{pgfscope}%
\pgfsetbuttcap%
\pgfsetroundjoin%
\definecolor{currentfill}{rgb}{0.090196,0.168627,0.227451}%
\pgfsetfillcolor{currentfill}%
\pgfsetlinewidth{0.803000pt}%
\definecolor{currentstroke}{rgb}{0.090196,0.168627,0.227451}%
\pgfsetstrokecolor{currentstroke}%
\pgfsetdash{}{0pt}%
\pgfsys@defobject{currentmarker}{\pgfqpoint{-0.048611in}{0.000000in}}{\pgfqpoint{-0.000000in}{0.000000in}}{%
\pgfpathmoveto{\pgfqpoint{-0.000000in}{0.000000in}}%
\pgfpathlineto{\pgfqpoint{-0.048611in}{0.000000in}}%
\pgfusepath{stroke,fill}%
}%
\begin{pgfscope}%
\pgfsys@transformshift{7.904091in}{8.592000in}%
\pgfsys@useobject{currentmarker}{}%
\end{pgfscope}%
\end{pgfscope}%
\begin{pgfscope}%
\pgfsetbuttcap%
\pgfsetroundjoin%
\definecolor{currentfill}{rgb}{0.090196,0.168627,0.227451}%
\pgfsetfillcolor{currentfill}%
\pgfsetlinewidth{0.000000pt}%
\definecolor{currentstroke}{rgb}{0.090196,0.168627,0.227451}%
\pgfsetstrokecolor{currentstroke}%
\pgfsetdash{}{0pt}%
\pgfpathmoveto{\pgfqpoint{7.620365in}{8.635647in}}%
\pgfpathlineto{\pgfqpoint{7.668763in}{8.635647in}}%
\pgfpathlineto{\pgfqpoint{7.668763in}{8.625257in}}%
\pgfpathlineto{\pgfqpoint{7.631654in}{8.625257in}}%
\pgfpathlineto{\pgfqpoint{7.631654in}{8.602933in}}%
\pgfpathquadraticcurveto{\pgfqpoint{7.634330in}{8.603851in}}{\pgfqpoint{7.637005in}{8.604300in}}%
\pgfpathquadraticcurveto{\pgfqpoint{7.639701in}{8.604749in}}{\pgfqpoint{7.642396in}{8.604749in}}%
\pgfpathquadraticcurveto{\pgfqpoint{7.657650in}{8.604749in}}{\pgfqpoint{7.666556in}{8.596390in}}%
\pgfpathquadraticcurveto{\pgfqpoint{7.675482in}{8.588030in}}{\pgfqpoint{7.675482in}{8.573753in}}%
\pgfpathquadraticcurveto{\pgfqpoint{7.675482in}{8.559046in}}{\pgfqpoint{7.666322in}{8.550882in}}%
\pgfpathquadraticcurveto{\pgfqpoint{7.657162in}{8.542737in}}{\pgfqpoint{7.640501in}{8.542737in}}%
\pgfpathquadraticcurveto{\pgfqpoint{7.634759in}{8.542737in}}{\pgfqpoint{7.628802in}{8.543714in}}%
\pgfpathquadraticcurveto{\pgfqpoint{7.622865in}{8.544690in}}{\pgfqpoint{7.616517in}{8.546644in}}%
\pgfpathlineto{\pgfqpoint{7.616517in}{8.559046in}}%
\pgfpathquadraticcurveto{\pgfqpoint{7.622005in}{8.556058in}}{\pgfqpoint{7.627865in}{8.554593in}}%
\pgfpathquadraticcurveto{\pgfqpoint{7.633724in}{8.553128in}}{\pgfqpoint{7.640248in}{8.553128in}}%
\pgfpathquadraticcurveto{\pgfqpoint{7.650814in}{8.553128in}}{\pgfqpoint{7.656966in}{8.558675in}}%
\pgfpathquadraticcurveto{\pgfqpoint{7.663138in}{8.564222in}}{\pgfqpoint{7.663138in}{8.573753in}}%
\pgfpathquadraticcurveto{\pgfqpoint{7.663138in}{8.583265in}}{\pgfqpoint{7.656966in}{8.588812in}}%
\pgfpathquadraticcurveto{\pgfqpoint{7.650814in}{8.594378in}}{\pgfqpoint{7.640248in}{8.594378in}}%
\pgfpathquadraticcurveto{\pgfqpoint{7.635306in}{8.594378in}}{\pgfqpoint{7.630384in}{8.593284in}}%
\pgfpathquadraticcurveto{\pgfqpoint{7.625482in}{8.592190in}}{\pgfqpoint{7.620365in}{8.589866in}}%
\pgfpathlineto{\pgfqpoint{7.620365in}{8.635647in}}%
\pgfpathlineto{\pgfqpoint{7.620365in}{8.635647in}}%
\pgfpathclose%
\pgfpathmoveto{\pgfqpoint{7.699759in}{8.560022in}}%
\pgfpathlineto{\pgfqpoint{7.712650in}{8.560022in}}%
\pgfpathlineto{\pgfqpoint{7.712650in}{8.544515in}}%
\pgfpathlineto{\pgfqpoint{7.699759in}{8.544515in}}%
\pgfpathlineto{\pgfqpoint{7.699759in}{8.560022in}}%
\pgfpathlineto{\pgfqpoint{7.699759in}{8.560022in}}%
\pgfpathclose%
\pgfpathmoveto{\pgfqpoint{7.765853in}{8.627522in}}%
\pgfpathquadraticcurveto{\pgfqpoint{7.756341in}{8.627522in}}{\pgfqpoint{7.751537in}{8.618147in}}%
\pgfpathquadraticcurveto{\pgfqpoint{7.746751in}{8.608792in}}{\pgfqpoint{7.746751in}{8.589983in}}%
\pgfpathquadraticcurveto{\pgfqpoint{7.746751in}{8.571253in}}{\pgfqpoint{7.751537in}{8.561878in}}%
\pgfpathquadraticcurveto{\pgfqpoint{7.756341in}{8.552503in}}{\pgfqpoint{7.765853in}{8.552503in}}%
\pgfpathquadraticcurveto{\pgfqpoint{7.775443in}{8.552503in}}{\pgfqpoint{7.780228in}{8.561878in}}%
\pgfpathquadraticcurveto{\pgfqpoint{7.785033in}{8.571253in}}{\pgfqpoint{7.785033in}{8.589983in}}%
\pgfpathquadraticcurveto{\pgfqpoint{7.785033in}{8.608792in}}{\pgfqpoint{7.780228in}{8.618147in}}%
\pgfpathquadraticcurveto{\pgfqpoint{7.775443in}{8.627522in}}{\pgfqpoint{7.765853in}{8.627522in}}%
\pgfpathlineto{\pgfqpoint{7.765853in}{8.627522in}}%
\pgfpathclose%
\pgfpathmoveto{\pgfqpoint{7.765853in}{8.637288in}}%
\pgfpathquadraticcurveto{\pgfqpoint{7.781185in}{8.637288in}}{\pgfqpoint{7.789271in}{8.625159in}}%
\pgfpathquadraticcurveto{\pgfqpoint{7.797357in}{8.613050in}}{\pgfqpoint{7.797357in}{8.589983in}}%
\pgfpathquadraticcurveto{\pgfqpoint{7.797357in}{8.566976in}}{\pgfqpoint{7.789271in}{8.554847in}}%
\pgfpathquadraticcurveto{\pgfqpoint{7.781185in}{8.542737in}}{\pgfqpoint{7.765853in}{8.542737in}}%
\pgfpathquadraticcurveto{\pgfqpoint{7.750541in}{8.542737in}}{\pgfqpoint{7.742455in}{8.554847in}}%
\pgfpathquadraticcurveto{\pgfqpoint{7.734369in}{8.566976in}}{\pgfqpoint{7.734369in}{8.589983in}}%
\pgfpathquadraticcurveto{\pgfqpoint{7.734369in}{8.613050in}}{\pgfqpoint{7.742455in}{8.625159in}}%
\pgfpathquadraticcurveto{\pgfqpoint{7.750541in}{8.637288in}}{\pgfqpoint{7.765853in}{8.637288in}}%
\pgfpathlineto{\pgfqpoint{7.765853in}{8.637288in}}%
\pgfpathclose%
\pgfusepath{fill}%
\end{pgfscope}%
\begin{pgfscope}%
\pgfsetbuttcap%
\pgfsetroundjoin%
\definecolor{currentfill}{rgb}{0.090196,0.168627,0.227451}%
\pgfsetfillcolor{currentfill}%
\pgfsetlinewidth{0.000000pt}%
\definecolor{currentstroke}{rgb}{0.090196,0.168627,0.227451}%
\pgfsetstrokecolor{currentstroke}%
\pgfsetdash{}{0pt}%
\pgfpathmoveto{\pgfqpoint{7.295869in}{4.860054in}}%
\pgfpathlineto{\pgfqpoint{7.336666in}{4.860054in}}%
\pgfpathlineto{\pgfqpoint{7.336666in}{4.884236in}}%
\pgfpathquadraticcurveto{\pgfqpoint{7.336666in}{4.896387in}}{\pgfqpoint{7.331629in}{4.902235in}}%
\pgfpathquadraticcurveto{\pgfqpoint{7.326592in}{4.908108in}}{\pgfqpoint{7.316232in}{4.908108in}}%
\pgfpathquadraticcurveto{\pgfqpoint{7.305776in}{4.908108in}}{\pgfqpoint{7.300835in}{4.902235in}}%
\pgfpathquadraticcurveto{\pgfqpoint{7.295869in}{4.896387in}}{\pgfqpoint{7.295869in}{4.884236in}}%
\pgfpathlineto{\pgfqpoint{7.295869in}{4.860054in}}%
\pgfpathlineto{\pgfqpoint{7.295869in}{4.860054in}}%
\pgfpathclose%
\pgfpathmoveto{\pgfqpoint{7.250060in}{4.860054in}}%
\pgfpathlineto{\pgfqpoint{7.283623in}{4.860054in}}%
\pgfpathlineto{\pgfqpoint{7.283623in}{4.882374in}}%
\pgfpathquadraticcurveto{\pgfqpoint{7.283623in}{4.893403in}}{\pgfqpoint{7.279493in}{4.898798in}}%
\pgfpathquadraticcurveto{\pgfqpoint{7.275340in}{4.904217in}}{\pgfqpoint{7.266842in}{4.904217in}}%
\pgfpathquadraticcurveto{\pgfqpoint{7.258415in}{4.904217in}}{\pgfqpoint{7.254237in}{4.898798in}}%
\pgfpathquadraticcurveto{\pgfqpoint{7.250060in}{4.893403in}}{\pgfqpoint{7.250060in}{4.882374in}}%
\pgfpathlineto{\pgfqpoint{7.250060in}{4.860054in}}%
\pgfpathlineto{\pgfqpoint{7.250060in}{4.860054in}}%
\pgfpathclose%
\pgfpathmoveto{\pgfqpoint{7.237671in}{4.844991in}}%
\pgfpathlineto{\pgfqpoint{7.237671in}{4.883496in}}%
\pgfpathquadraticcurveto{\pgfqpoint{7.237671in}{4.900731in}}{\pgfqpoint{7.244832in}{4.910041in}}%
\pgfpathquadraticcurveto{\pgfqpoint{7.251993in}{4.919375in}}{\pgfqpoint{7.265194in}{4.919375in}}%
\pgfpathquadraticcurveto{\pgfqpoint{7.275435in}{4.919375in}}{\pgfqpoint{7.281475in}{4.914601in}}%
\pgfpathquadraticcurveto{\pgfqpoint{7.287514in}{4.909826in}}{\pgfqpoint{7.288994in}{4.900564in}}%
\pgfpathquadraticcurveto{\pgfqpoint{7.291381in}{4.911688in}}{\pgfqpoint{7.298973in}{4.917847in}}%
\pgfpathquadraticcurveto{\pgfqpoint{7.306540in}{4.924006in}}{\pgfqpoint{7.317879in}{4.924006in}}%
\pgfpathquadraticcurveto{\pgfqpoint{7.332799in}{4.924006in}}{\pgfqpoint{7.340939in}{4.913861in}}%
\pgfpathquadraticcurveto{\pgfqpoint{7.349055in}{4.903715in}}{\pgfqpoint{7.349055in}{4.884976in}}%
\pgfpathlineto{\pgfqpoint{7.349055in}{4.844991in}}%
\pgfpathlineto{\pgfqpoint{7.237671in}{4.844991in}}%
\pgfpathlineto{\pgfqpoint{7.237671in}{4.844991in}}%
\pgfpathclose%
\pgfpathmoveto{\pgfqpoint{7.237671in}{4.949811in}}%
\pgfpathlineto{\pgfqpoint{7.237671in}{4.964874in}}%
\pgfpathlineto{\pgfqpoint{7.341297in}{4.964874in}}%
\pgfpathquadraticcurveto{\pgfqpoint{7.361444in}{4.964874in}}{\pgfqpoint{7.370539in}{4.957235in}}%
\pgfpathquadraticcurveto{\pgfqpoint{7.379635in}{4.949596in}}{\pgfqpoint{7.379635in}{4.932648in}}%
\pgfpathlineto{\pgfqpoint{7.379635in}{4.926918in}}%
\pgfpathlineto{\pgfqpoint{7.366959in}{4.926918in}}%
\pgfpathlineto{\pgfqpoint{7.366959in}{4.931621in}}%
\pgfpathquadraticcurveto{\pgfqpoint{7.366959in}{4.941599in}}{\pgfqpoint{7.361349in}{4.945705in}}%
\pgfpathquadraticcurveto{\pgfqpoint{7.355763in}{4.949811in}}{\pgfqpoint{7.341297in}{4.949811in}}%
\pgfpathlineto{\pgfqpoint{7.237671in}{4.949811in}}%
\pgfpathlineto{\pgfqpoint{7.237671in}{4.949811in}}%
\pgfpathclose%
\pgfpathmoveto{\pgfqpoint{7.241323in}{5.061649in}}%
\pgfpathlineto{\pgfqpoint{7.256028in}{5.061649in}}%
\pgfpathquadraticcurveto{\pgfqpoint{7.251922in}{5.053079in}}{\pgfqpoint{7.249917in}{5.045464in}}%
\pgfpathquadraticcurveto{\pgfqpoint{7.247888in}{5.037849in}}{\pgfqpoint{7.247888in}{5.030760in}}%
\pgfpathquadraticcurveto{\pgfqpoint{7.247888in}{5.018466in}}{\pgfqpoint{7.252662in}{5.011782in}}%
\pgfpathquadraticcurveto{\pgfqpoint{7.257436in}{5.005098in}}{\pgfqpoint{7.266245in}{5.005098in}}%
\pgfpathquadraticcurveto{\pgfqpoint{7.273645in}{5.005098in}}{\pgfqpoint{7.277417in}{5.009538in}}%
\pgfpathquadraticcurveto{\pgfqpoint{7.281164in}{5.013978in}}{\pgfqpoint{7.283480in}{5.026367in}}%
\pgfpathlineto{\pgfqpoint{7.285342in}{5.035462in}}%
\pgfpathquadraticcurveto{\pgfqpoint{7.288565in}{5.052316in}}{\pgfqpoint{7.296657in}{5.060336in}}%
\pgfpathquadraticcurveto{\pgfqpoint{7.304750in}{5.068357in}}{\pgfqpoint{7.318309in}{5.068357in}}%
\pgfpathquadraticcurveto{\pgfqpoint{7.334517in}{5.068357in}}{\pgfqpoint{7.342872in}{5.057496in}}%
\pgfpathquadraticcurveto{\pgfqpoint{7.351227in}{5.046658in}}{\pgfqpoint{7.351227in}{5.025699in}}%
\pgfpathquadraticcurveto{\pgfqpoint{7.351227in}{5.017797in}}{\pgfqpoint{7.349437in}{5.008869in}}%
\pgfpathquadraticcurveto{\pgfqpoint{7.347647in}{4.999965in}}{\pgfqpoint{7.344138in}{4.990417in}}%
\pgfpathlineto{\pgfqpoint{7.328621in}{4.990417in}}%
\pgfpathquadraticcurveto{\pgfqpoint{7.333753in}{4.999583in}}{\pgfqpoint{7.336379in}{5.008392in}}%
\pgfpathquadraticcurveto{\pgfqpoint{7.338981in}{5.017201in}}{\pgfqpoint{7.338981in}{5.025699in}}%
\pgfpathquadraticcurveto{\pgfqpoint{7.338981in}{5.038589in}}{\pgfqpoint{7.333921in}{5.045608in}}%
\pgfpathquadraticcurveto{\pgfqpoint{7.328836in}{5.052626in}}{\pgfqpoint{7.319431in}{5.052626in}}%
\pgfpathquadraticcurveto{\pgfqpoint{7.311243in}{5.052626in}}{\pgfqpoint{7.306612in}{5.047589in}}%
\pgfpathquadraticcurveto{\pgfqpoint{7.301980in}{5.042552in}}{\pgfqpoint{7.299665in}{5.031070in}}%
\pgfpathlineto{\pgfqpoint{7.297875in}{5.021879in}}%
\pgfpathquadraticcurveto{\pgfqpoint{7.294533in}{5.005026in}}{\pgfqpoint{7.287371in}{4.997483in}}%
\pgfpathquadraticcurveto{\pgfqpoint{7.280210in}{4.989963in}}{\pgfqpoint{7.267438in}{4.989963in}}%
\pgfpathquadraticcurveto{\pgfqpoint{7.252662in}{4.989963in}}{\pgfqpoint{7.244164in}{5.000371in}}%
\pgfpathquadraticcurveto{\pgfqpoint{7.235665in}{5.010779in}}{\pgfqpoint{7.235665in}{5.029041in}}%
\pgfpathquadraticcurveto{\pgfqpoint{7.235665in}{5.036895in}}{\pgfqpoint{7.237098in}{5.045011in}}%
\pgfpathquadraticcurveto{\pgfqpoint{7.238506in}{5.053151in}}{\pgfqpoint{7.241323in}{5.061649in}}%
\pgfpathlineto{\pgfqpoint{7.241323in}{5.061649in}}%
\pgfpathclose%
\pgfpathmoveto{\pgfqpoint{7.303842in}{5.211300in}}%
\pgfpathlineto{\pgfqpoint{7.310550in}{5.211300in}}%
\pgfpathlineto{\pgfqpoint{7.310550in}{5.148184in}}%
\pgfpathquadraticcurveto{\pgfqpoint{7.324730in}{5.149091in}}{\pgfqpoint{7.332154in}{5.156730in}}%
\pgfpathquadraticcurveto{\pgfqpoint{7.339578in}{5.164368in}}{\pgfqpoint{7.339578in}{5.178023in}}%
\pgfpathquadraticcurveto{\pgfqpoint{7.339578in}{5.185924in}}{\pgfqpoint{7.337645in}{5.193349in}}%
\pgfpathquadraticcurveto{\pgfqpoint{7.335711in}{5.200773in}}{\pgfqpoint{7.331820in}{5.208101in}}%
\pgfpathlineto{\pgfqpoint{7.344806in}{5.208101in}}%
\pgfpathquadraticcurveto{\pgfqpoint{7.347933in}{5.200701in}}{\pgfqpoint{7.349580in}{5.192943in}}%
\pgfpathquadraticcurveto{\pgfqpoint{7.351227in}{5.185184in}}{\pgfqpoint{7.351227in}{5.177211in}}%
\pgfpathquadraticcurveto{\pgfqpoint{7.351227in}{5.157207in}}{\pgfqpoint{7.339602in}{5.145534in}}%
\pgfpathquadraticcurveto{\pgfqpoint{7.327953in}{5.133861in}}{\pgfqpoint{7.308092in}{5.133861in}}%
\pgfpathquadraticcurveto{\pgfqpoint{7.287586in}{5.133861in}}{\pgfqpoint{7.275555in}{5.144937in}}%
\pgfpathquadraticcurveto{\pgfqpoint{7.263500in}{5.156013in}}{\pgfqpoint{7.263500in}{5.174824in}}%
\pgfpathquadraticcurveto{\pgfqpoint{7.263500in}{5.191678in}}{\pgfqpoint{7.274361in}{5.201489in}}%
\pgfpathquadraticcurveto{\pgfqpoint{7.285199in}{5.211300in}}{\pgfqpoint{7.303842in}{5.211300in}}%
\pgfpathlineto{\pgfqpoint{7.303842in}{5.211300in}}%
\pgfpathclose%
\pgfpathmoveto{\pgfqpoint{7.299808in}{5.197574in}}%
\pgfpathquadraticcurveto{\pgfqpoint{7.288565in}{5.197431in}}{\pgfqpoint{7.281857in}{5.191272in}}%
\pgfpathquadraticcurveto{\pgfqpoint{7.275125in}{5.185113in}}{\pgfqpoint{7.275125in}{5.174967in}}%
\pgfpathquadraticcurveto{\pgfqpoint{7.275125in}{5.163485in}}{\pgfqpoint{7.281618in}{5.156586in}}%
\pgfpathquadraticcurveto{\pgfqpoint{7.288111in}{5.149687in}}{\pgfqpoint{7.299904in}{5.148637in}}%
\pgfpathlineto{\pgfqpoint{7.299808in}{5.197574in}}%
\pgfpathlineto{\pgfqpoint{7.299808in}{5.197574in}}%
\pgfpathclose%
\pgfpathmoveto{\pgfqpoint{7.267964in}{5.287092in}}%
\pgfpathlineto{\pgfqpoint{7.280950in}{5.287092in}}%
\pgfpathquadraticcurveto{\pgfqpoint{7.277966in}{5.281291in}}{\pgfqpoint{7.276486in}{5.275013in}}%
\pgfpathquadraticcurveto{\pgfqpoint{7.274982in}{5.268759in}}{\pgfqpoint{7.274982in}{5.262027in}}%
\pgfpathquadraticcurveto{\pgfqpoint{7.274982in}{5.251810in}}{\pgfqpoint{7.278109in}{5.246701in}}%
\pgfpathquadraticcurveto{\pgfqpoint{7.281236in}{5.241593in}}{\pgfqpoint{7.287514in}{5.241593in}}%
\pgfpathquadraticcurveto{\pgfqpoint{7.292289in}{5.241593in}}{\pgfqpoint{7.295010in}{5.245245in}}%
\pgfpathquadraticcurveto{\pgfqpoint{7.297731in}{5.248898in}}{\pgfqpoint{7.300190in}{5.259950in}}%
\pgfpathlineto{\pgfqpoint{7.301240in}{5.264653in}}%
\pgfpathquadraticcurveto{\pgfqpoint{7.304368in}{5.279262in}}{\pgfqpoint{7.310073in}{5.285421in}}%
\pgfpathquadraticcurveto{\pgfqpoint{7.315778in}{5.291580in}}{\pgfqpoint{7.325995in}{5.291580in}}%
\pgfpathquadraticcurveto{\pgfqpoint{7.337645in}{5.291580in}}{\pgfqpoint{7.344448in}{5.282365in}}%
\pgfpathquadraticcurveto{\pgfqpoint{7.351227in}{5.273151in}}{\pgfqpoint{7.351227in}{5.257038in}}%
\pgfpathquadraticcurveto{\pgfqpoint{7.351227in}{5.250330in}}{\pgfqpoint{7.349914in}{5.243049in}}%
\pgfpathquadraticcurveto{\pgfqpoint{7.348602in}{5.235768in}}{\pgfqpoint{7.346000in}{5.227724in}}%
\pgfpathlineto{\pgfqpoint{7.331820in}{5.227724in}}%
\pgfpathquadraticcurveto{\pgfqpoint{7.335783in}{5.235339in}}{\pgfqpoint{7.337764in}{5.242715in}}%
\pgfpathquadraticcurveto{\pgfqpoint{7.339721in}{5.250091in}}{\pgfqpoint{7.339721in}{5.257348in}}%
\pgfpathquadraticcurveto{\pgfqpoint{7.339721in}{5.267040in}}{\pgfqpoint{7.336403in}{5.272244in}}%
\pgfpathquadraticcurveto{\pgfqpoint{7.333085in}{5.277472in}}{\pgfqpoint{7.327046in}{5.277472in}}%
\pgfpathquadraticcurveto{\pgfqpoint{7.321460in}{5.277472in}}{\pgfqpoint{7.318476in}{5.273700in}}%
\pgfpathquadraticcurveto{\pgfqpoint{7.315492in}{5.269952in}}{\pgfqpoint{7.312723in}{5.257181in}}%
\pgfpathlineto{\pgfqpoint{7.311601in}{5.252407in}}%
\pgfpathquadraticcurveto{\pgfqpoint{7.308927in}{5.239659in}}{\pgfqpoint{7.303365in}{5.233978in}}%
\pgfpathquadraticcurveto{\pgfqpoint{7.297803in}{5.228320in}}{\pgfqpoint{7.288111in}{5.228320in}}%
\pgfpathquadraticcurveto{\pgfqpoint{7.276319in}{5.228320in}}{\pgfqpoint{7.269921in}{5.236675in}}%
\pgfpathquadraticcurveto{\pgfqpoint{7.263500in}{5.245030in}}{\pgfqpoint{7.263500in}{5.260404in}}%
\pgfpathquadraticcurveto{\pgfqpoint{7.263500in}{5.267995in}}{\pgfqpoint{7.264622in}{5.274703in}}%
\pgfpathquadraticcurveto{\pgfqpoint{7.265720in}{5.281434in}}{\pgfqpoint{7.267964in}{5.287092in}}%
\pgfpathlineto{\pgfqpoint{7.267964in}{5.287092in}}%
\pgfpathclose%
\pgfpathmoveto{\pgfqpoint{7.241776in}{5.327005in}}%
\pgfpathlineto{\pgfqpoint{7.265505in}{5.327005in}}%
\pgfpathlineto{\pgfqpoint{7.265505in}{5.355269in}}%
\pgfpathlineto{\pgfqpoint{7.276175in}{5.355269in}}%
\pgfpathlineto{\pgfqpoint{7.276175in}{5.327005in}}%
\pgfpathlineto{\pgfqpoint{7.321531in}{5.327005in}}%
\pgfpathquadraticcurveto{\pgfqpoint{7.331748in}{5.327005in}}{\pgfqpoint{7.334661in}{5.329798in}}%
\pgfpathquadraticcurveto{\pgfqpoint{7.337573in}{5.332591in}}{\pgfqpoint{7.337573in}{5.341185in}}%
\pgfpathlineto{\pgfqpoint{7.337573in}{5.355269in}}%
\pgfpathlineto{\pgfqpoint{7.349055in}{5.355269in}}%
\pgfpathlineto{\pgfqpoint{7.349055in}{5.341185in}}%
\pgfpathquadraticcurveto{\pgfqpoint{7.349055in}{5.325286in}}{\pgfqpoint{7.343135in}{5.319247in}}%
\pgfpathquadraticcurveto{\pgfqpoint{7.337191in}{5.313207in}}{\pgfqpoint{7.321531in}{5.313207in}}%
\pgfpathlineto{\pgfqpoint{7.276175in}{5.313207in}}%
\pgfpathlineto{\pgfqpoint{7.276175in}{5.303134in}}%
\pgfpathlineto{\pgfqpoint{7.265505in}{5.303134in}}%
\pgfpathlineto{\pgfqpoint{7.265505in}{5.313207in}}%
\pgfpathlineto{\pgfqpoint{7.241776in}{5.313207in}}%
\pgfpathlineto{\pgfqpoint{7.241776in}{5.327005in}}%
\pgfpathlineto{\pgfqpoint{7.241776in}{5.327005in}}%
\pgfpathclose%
\pgfpathmoveto{\pgfqpoint{7.265505in}{5.373316in}}%
\pgfpathlineto{\pgfqpoint{7.265505in}{5.387042in}}%
\pgfpathlineto{\pgfqpoint{7.349055in}{5.387042in}}%
\pgfpathlineto{\pgfqpoint{7.349055in}{5.373316in}}%
\pgfpathlineto{\pgfqpoint{7.265505in}{5.373316in}}%
\pgfpathlineto{\pgfqpoint{7.265505in}{5.373316in}}%
\pgfpathclose%
\pgfpathmoveto{\pgfqpoint{7.232968in}{5.373316in}}%
\pgfpathlineto{\pgfqpoint{7.232968in}{5.387042in}}%
\pgfpathlineto{\pgfqpoint{7.250370in}{5.387042in}}%
\pgfpathlineto{\pgfqpoint{7.250370in}{5.373316in}}%
\pgfpathlineto{\pgfqpoint{7.232968in}{5.373316in}}%
\pgfpathlineto{\pgfqpoint{7.232968in}{5.373316in}}%
\pgfpathclose%
\pgfpathmoveto{\pgfqpoint{7.281546in}{5.480809in}}%
\pgfpathquadraticcurveto{\pgfqpoint{7.272284in}{5.485966in}}{\pgfqpoint{7.267892in}{5.493127in}}%
\pgfpathquadraticcurveto{\pgfqpoint{7.263500in}{5.500289in}}{\pgfqpoint{7.263500in}{5.509980in}}%
\pgfpathquadraticcurveto{\pgfqpoint{7.263500in}{5.523038in}}{\pgfqpoint{7.272642in}{5.530128in}}%
\pgfpathquadraticcurveto{\pgfqpoint{7.281761in}{5.537218in}}{\pgfqpoint{7.298615in}{5.537218in}}%
\pgfpathlineto{\pgfqpoint{7.349055in}{5.537218in}}%
\pgfpathlineto{\pgfqpoint{7.349055in}{5.523420in}}%
\pgfpathlineto{\pgfqpoint{7.299068in}{5.523420in}}%
\pgfpathquadraticcurveto{\pgfqpoint{7.287061in}{5.523420in}}{\pgfqpoint{7.281260in}{5.519147in}}%
\pgfpathquadraticcurveto{\pgfqpoint{7.275435in}{5.514898in}}{\pgfqpoint{7.275435in}{5.506185in}}%
\pgfpathquadraticcurveto{\pgfqpoint{7.275435in}{5.495514in}}{\pgfqpoint{7.282525in}{5.489308in}}%
\pgfpathquadraticcurveto{\pgfqpoint{7.289591in}{5.483125in}}{\pgfqpoint{7.301837in}{5.483125in}}%
\pgfpathlineto{\pgfqpoint{7.349055in}{5.483125in}}%
\pgfpathlineto{\pgfqpoint{7.349055in}{5.469327in}}%
\pgfpathlineto{\pgfqpoint{7.299068in}{5.469327in}}%
\pgfpathquadraticcurveto{\pgfqpoint{7.286989in}{5.469327in}}{\pgfqpoint{7.281212in}{5.465078in}}%
\pgfpathquadraticcurveto{\pgfqpoint{7.275435in}{5.460829in}}{\pgfqpoint{7.275435in}{5.451949in}}%
\pgfpathquadraticcurveto{\pgfqpoint{7.275435in}{5.441421in}}{\pgfqpoint{7.282549in}{5.435215in}}%
\pgfpathquadraticcurveto{\pgfqpoint{7.289663in}{5.429032in}}{\pgfqpoint{7.301837in}{5.429032in}}%
\pgfpathlineto{\pgfqpoint{7.349055in}{5.429032in}}%
\pgfpathlineto{\pgfqpoint{7.349055in}{5.415234in}}%
\pgfpathlineto{\pgfqpoint{7.265505in}{5.415234in}}%
\pgfpathlineto{\pgfqpoint{7.265505in}{5.429032in}}%
\pgfpathlineto{\pgfqpoint{7.278491in}{5.429032in}}%
\pgfpathquadraticcurveto{\pgfqpoint{7.270804in}{5.433735in}}{\pgfqpoint{7.267152in}{5.440299in}}%
\pgfpathquadraticcurveto{\pgfqpoint{7.263500in}{5.446864in}}{\pgfqpoint{7.263500in}{5.455888in}}%
\pgfpathquadraticcurveto{\pgfqpoint{7.263500in}{5.465007in}}{\pgfqpoint{7.268131in}{5.471380in}}%
\pgfpathquadraticcurveto{\pgfqpoint{7.272738in}{5.477754in}}{\pgfqpoint{7.281546in}{5.480809in}}%
\pgfpathlineto{\pgfqpoint{7.281546in}{5.480809in}}%
\pgfpathclose%
\pgfpathmoveto{\pgfqpoint{7.307065in}{5.602554in}}%
\pgfpathquadraticcurveto{\pgfqpoint{7.307065in}{5.585916in}}{\pgfqpoint{7.310861in}{5.579494in}}%
\pgfpathquadraticcurveto{\pgfqpoint{7.314656in}{5.573073in}}{\pgfqpoint{7.323847in}{5.573073in}}%
\pgfpathquadraticcurveto{\pgfqpoint{7.331151in}{5.573073in}}{\pgfqpoint{7.335448in}{5.577895in}}%
\pgfpathquadraticcurveto{\pgfqpoint{7.339721in}{5.582717in}}{\pgfqpoint{7.339721in}{5.590977in}}%
\pgfpathquadraticcurveto{\pgfqpoint{7.339721in}{5.602411in}}{\pgfqpoint{7.331629in}{5.609310in}}%
\pgfpathquadraticcurveto{\pgfqpoint{7.323536in}{5.616209in}}{\pgfqpoint{7.310121in}{5.616209in}}%
\pgfpathlineto{\pgfqpoint{7.307065in}{5.616209in}}%
\pgfpathlineto{\pgfqpoint{7.307065in}{5.602554in}}%
\pgfpathlineto{\pgfqpoint{7.307065in}{5.602554in}}%
\pgfpathclose%
\pgfpathmoveto{\pgfqpoint{7.301384in}{5.629935in}}%
\pgfpathlineto{\pgfqpoint{7.349055in}{5.629935in}}%
\pgfpathlineto{\pgfqpoint{7.349055in}{5.616209in}}%
\pgfpathlineto{\pgfqpoint{7.336379in}{5.616209in}}%
\pgfpathquadraticcurveto{\pgfqpoint{7.343970in}{5.611506in}}{\pgfqpoint{7.347599in}{5.604488in}}%
\pgfpathquadraticcurveto{\pgfqpoint{7.351227in}{5.597470in}}{\pgfqpoint{7.351227in}{5.587324in}}%
\pgfpathquadraticcurveto{\pgfqpoint{7.351227in}{5.574505in}}{\pgfqpoint{7.344018in}{5.566914in}}%
\pgfpathquadraticcurveto{\pgfqpoint{7.336809in}{5.559347in}}{\pgfqpoint{7.324730in}{5.559347in}}%
\pgfpathquadraticcurveto{\pgfqpoint{7.310646in}{5.559347in}}{\pgfqpoint{7.303484in}{5.568776in}}%
\pgfpathquadraticcurveto{\pgfqpoint{7.296323in}{5.578229in}}{\pgfqpoint{7.296323in}{5.596944in}}%
\pgfpathlineto{\pgfqpoint{7.296323in}{5.616209in}}%
\pgfpathlineto{\pgfqpoint{7.294962in}{5.616209in}}%
\pgfpathquadraticcurveto{\pgfqpoint{7.285485in}{5.616209in}}{\pgfqpoint{7.280305in}{5.609978in}}%
\pgfpathquadraticcurveto{\pgfqpoint{7.275125in}{5.603748in}}{\pgfqpoint{7.275125in}{5.592480in}}%
\pgfpathquadraticcurveto{\pgfqpoint{7.275125in}{5.585319in}}{\pgfqpoint{7.276844in}{5.578516in}}%
\pgfpathquadraticcurveto{\pgfqpoint{7.278562in}{5.571736in}}{\pgfqpoint{7.282000in}{5.565482in}}%
\pgfpathlineto{\pgfqpoint{7.269300in}{5.565482in}}%
\pgfpathquadraticcurveto{\pgfqpoint{7.266388in}{5.573001in}}{\pgfqpoint{7.264956in}{5.580091in}}%
\pgfpathquadraticcurveto{\pgfqpoint{7.263500in}{5.587181in}}{\pgfqpoint{7.263500in}{5.593889in}}%
\pgfpathquadraticcurveto{\pgfqpoint{7.263500in}{5.612031in}}{\pgfqpoint{7.272905in}{5.620983in}}%
\pgfpathquadraticcurveto{\pgfqpoint{7.282286in}{5.629935in}}{\pgfqpoint{7.301384in}{5.629935in}}%
\pgfpathlineto{\pgfqpoint{7.301384in}{5.629935in}}%
\pgfpathclose%
\pgfpathmoveto{\pgfqpoint{7.241776in}{5.671782in}}%
\pgfpathlineto{\pgfqpoint{7.265505in}{5.671782in}}%
\pgfpathlineto{\pgfqpoint{7.265505in}{5.700046in}}%
\pgfpathlineto{\pgfqpoint{7.276175in}{5.700046in}}%
\pgfpathlineto{\pgfqpoint{7.276175in}{5.671782in}}%
\pgfpathlineto{\pgfqpoint{7.321531in}{5.671782in}}%
\pgfpathquadraticcurveto{\pgfqpoint{7.331748in}{5.671782in}}{\pgfqpoint{7.334661in}{5.674575in}}%
\pgfpathquadraticcurveto{\pgfqpoint{7.337573in}{5.677368in}}{\pgfqpoint{7.337573in}{5.685961in}}%
\pgfpathlineto{\pgfqpoint{7.337573in}{5.700046in}}%
\pgfpathlineto{\pgfqpoint{7.349055in}{5.700046in}}%
\pgfpathlineto{\pgfqpoint{7.349055in}{5.685961in}}%
\pgfpathquadraticcurveto{\pgfqpoint{7.349055in}{5.670063in}}{\pgfqpoint{7.343135in}{5.664023in}}%
\pgfpathquadraticcurveto{\pgfqpoint{7.337191in}{5.657984in}}{\pgfqpoint{7.321531in}{5.657984in}}%
\pgfpathlineto{\pgfqpoint{7.276175in}{5.657984in}}%
\pgfpathlineto{\pgfqpoint{7.276175in}{5.647910in}}%
\pgfpathlineto{\pgfqpoint{7.265505in}{5.647910in}}%
\pgfpathlineto{\pgfqpoint{7.265505in}{5.657984in}}%
\pgfpathlineto{\pgfqpoint{7.241776in}{5.657984in}}%
\pgfpathlineto{\pgfqpoint{7.241776in}{5.671782in}}%
\pgfpathlineto{\pgfqpoint{7.241776in}{5.671782in}}%
\pgfpathclose%
\pgfpathmoveto{\pgfqpoint{7.303842in}{5.789564in}}%
\pgfpathlineto{\pgfqpoint{7.310550in}{5.789564in}}%
\pgfpathlineto{\pgfqpoint{7.310550in}{5.726447in}}%
\pgfpathquadraticcurveto{\pgfqpoint{7.324730in}{5.727355in}}{\pgfqpoint{7.332154in}{5.734993in}}%
\pgfpathquadraticcurveto{\pgfqpoint{7.339578in}{5.742632in}}{\pgfqpoint{7.339578in}{5.756287in}}%
\pgfpathquadraticcurveto{\pgfqpoint{7.339578in}{5.764188in}}{\pgfqpoint{7.337645in}{5.771612in}}%
\pgfpathquadraticcurveto{\pgfqpoint{7.335711in}{5.779036in}}{\pgfqpoint{7.331820in}{5.786365in}}%
\pgfpathlineto{\pgfqpoint{7.344806in}{5.786365in}}%
\pgfpathquadraticcurveto{\pgfqpoint{7.347933in}{5.778965in}}{\pgfqpoint{7.349580in}{5.771207in}}%
\pgfpathquadraticcurveto{\pgfqpoint{7.351227in}{5.763448in}}{\pgfqpoint{7.351227in}{5.755475in}}%
\pgfpathquadraticcurveto{\pgfqpoint{7.351227in}{5.735471in}}{\pgfqpoint{7.339602in}{5.723798in}}%
\pgfpathquadraticcurveto{\pgfqpoint{7.327953in}{5.712125in}}{\pgfqpoint{7.308092in}{5.712125in}}%
\pgfpathquadraticcurveto{\pgfqpoint{7.287586in}{5.712125in}}{\pgfqpoint{7.275555in}{5.723201in}}%
\pgfpathquadraticcurveto{\pgfqpoint{7.263500in}{5.734277in}}{\pgfqpoint{7.263500in}{5.753088in}}%
\pgfpathquadraticcurveto{\pgfqpoint{7.263500in}{5.769941in}}{\pgfqpoint{7.274361in}{5.779753in}}%
\pgfpathquadraticcurveto{\pgfqpoint{7.285199in}{5.789564in}}{\pgfqpoint{7.303842in}{5.789564in}}%
\pgfpathlineto{\pgfqpoint{7.303842in}{5.789564in}}%
\pgfpathclose%
\pgfpathmoveto{\pgfqpoint{7.299808in}{5.775838in}}%
\pgfpathquadraticcurveto{\pgfqpoint{7.288565in}{5.775694in}}{\pgfqpoint{7.281857in}{5.769536in}}%
\pgfpathquadraticcurveto{\pgfqpoint{7.275125in}{5.763377in}}{\pgfqpoint{7.275125in}{5.753231in}}%
\pgfpathquadraticcurveto{\pgfqpoint{7.275125in}{5.741749in}}{\pgfqpoint{7.281618in}{5.734850in}}%
\pgfpathquadraticcurveto{\pgfqpoint{7.288111in}{5.727951in}}{\pgfqpoint{7.299904in}{5.726901in}}%
\pgfpathlineto{\pgfqpoint{7.299808in}{5.775838in}}%
\pgfpathlineto{\pgfqpoint{7.299808in}{5.775838in}}%
\pgfpathclose%
\pgfpathmoveto{\pgfqpoint{7.233135in}{5.893620in}}%
\pgfpathquadraticcurveto{\pgfqpoint{7.250275in}{5.883641in}}{\pgfqpoint{7.267080in}{5.878772in}}%
\pgfpathquadraticcurveto{\pgfqpoint{7.283862in}{5.873926in}}{\pgfqpoint{7.301097in}{5.873926in}}%
\pgfpathquadraticcurveto{\pgfqpoint{7.318309in}{5.873926in}}{\pgfqpoint{7.335210in}{5.878819in}}%
\pgfpathquadraticcurveto{\pgfqpoint{7.352111in}{5.883713in}}{\pgfqpoint{7.369203in}{5.893620in}}%
\pgfpathlineto{\pgfqpoint{7.369203in}{5.881684in}}%
\pgfpathquadraticcurveto{\pgfqpoint{7.351657in}{5.870512in}}{\pgfqpoint{7.334732in}{5.864950in}}%
\pgfpathquadraticcurveto{\pgfqpoint{7.317807in}{5.859388in}}{\pgfqpoint{7.301097in}{5.859388in}}%
\pgfpathquadraticcurveto{\pgfqpoint{7.284459in}{5.859388in}}{\pgfqpoint{7.267605in}{5.864902in}}%
\pgfpathquadraticcurveto{\pgfqpoint{7.250728in}{5.870441in}}{\pgfqpoint{7.233135in}{5.881684in}}%
\pgfpathlineto{\pgfqpoint{7.233135in}{5.893620in}}%
\pgfpathlineto{\pgfqpoint{7.233135in}{5.893620in}}%
\pgfpathclose%
\pgfpathmoveto{\pgfqpoint{7.237671in}{5.905412in}}%
\pgfpathlineto{\pgfqpoint{7.237671in}{5.999633in}}%
\pgfpathlineto{\pgfqpoint{7.250370in}{5.999633in}}%
\pgfpathlineto{\pgfqpoint{7.250370in}{5.960102in}}%
\pgfpathlineto{\pgfqpoint{7.349055in}{5.960102in}}%
\pgfpathlineto{\pgfqpoint{7.349055in}{5.944967in}}%
\pgfpathlineto{\pgfqpoint{7.250370in}{5.944967in}}%
\pgfpathlineto{\pgfqpoint{7.250370in}{5.905412in}}%
\pgfpathlineto{\pgfqpoint{7.237671in}{5.905412in}}%
\pgfpathlineto{\pgfqpoint{7.237671in}{5.905412in}}%
\pgfpathclose%
\pgfpathmoveto{\pgfqpoint{7.237671in}{6.004264in}}%
\pgfpathlineto{\pgfqpoint{7.237671in}{6.019470in}}%
\pgfpathlineto{\pgfqpoint{7.331820in}{6.042888in}}%
\pgfpathlineto{\pgfqpoint{7.237671in}{6.066235in}}%
\pgfpathlineto{\pgfqpoint{7.237671in}{6.083184in}}%
\pgfpathlineto{\pgfqpoint{7.331820in}{6.106602in}}%
\pgfpathlineto{\pgfqpoint{7.237671in}{6.129948in}}%
\pgfpathlineto{\pgfqpoint{7.237671in}{6.145250in}}%
\pgfpathlineto{\pgfqpoint{7.349055in}{6.117272in}}%
\pgfpathlineto{\pgfqpoint{7.349055in}{6.098318in}}%
\pgfpathlineto{\pgfqpoint{7.252375in}{6.074829in}}%
\pgfpathlineto{\pgfqpoint{7.349055in}{6.051100in}}%
\pgfpathlineto{\pgfqpoint{7.349055in}{6.032146in}}%
\pgfpathlineto{\pgfqpoint{7.237671in}{6.004264in}}%
\pgfpathlineto{\pgfqpoint{7.237671in}{6.004264in}}%
\pgfpathclose%
\pgfpathmoveto{\pgfqpoint{7.237671in}{6.165230in}}%
\pgfpathlineto{\pgfqpoint{7.237671in}{6.229230in}}%
\pgfpathlineto{\pgfqpoint{7.250370in}{6.229230in}}%
\pgfpathlineto{\pgfqpoint{7.250370in}{6.180293in}}%
\pgfpathlineto{\pgfqpoint{7.283194in}{6.180293in}}%
\pgfpathlineto{\pgfqpoint{7.283194in}{6.224455in}}%
\pgfpathlineto{\pgfqpoint{7.295869in}{6.224455in}}%
\pgfpathlineto{\pgfqpoint{7.295869in}{6.180293in}}%
\pgfpathlineto{\pgfqpoint{7.349055in}{6.180293in}}%
\pgfpathlineto{\pgfqpoint{7.349055in}{6.165230in}}%
\pgfpathlineto{\pgfqpoint{7.237671in}{6.165230in}}%
\pgfpathlineto{\pgfqpoint{7.237671in}{6.165230in}}%
\pgfpathclose%
\pgfpathmoveto{\pgfqpoint{7.237671in}{6.253101in}}%
\pgfpathlineto{\pgfqpoint{7.237671in}{6.323522in}}%
\pgfpathlineto{\pgfqpoint{7.250370in}{6.323522in}}%
\pgfpathlineto{\pgfqpoint{7.250370in}{6.268164in}}%
\pgfpathlineto{\pgfqpoint{7.283337in}{6.268164in}}%
\pgfpathlineto{\pgfqpoint{7.283337in}{6.321207in}}%
\pgfpathlineto{\pgfqpoint{7.296013in}{6.321207in}}%
\pgfpathlineto{\pgfqpoint{7.296013in}{6.268164in}}%
\pgfpathlineto{\pgfqpoint{7.336379in}{6.268164in}}%
\pgfpathlineto{\pgfqpoint{7.336379in}{6.324859in}}%
\pgfpathlineto{\pgfqpoint{7.349055in}{6.324859in}}%
\pgfpathlineto{\pgfqpoint{7.349055in}{6.253101in}}%
\pgfpathlineto{\pgfqpoint{7.237671in}{6.253101in}}%
\pgfpathlineto{\pgfqpoint{7.237671in}{6.253101in}}%
\pgfpathclose%
\pgfpathmoveto{\pgfqpoint{7.241323in}{6.464961in}}%
\pgfpathlineto{\pgfqpoint{7.256028in}{6.464961in}}%
\pgfpathquadraticcurveto{\pgfqpoint{7.251922in}{6.456391in}}{\pgfqpoint{7.249917in}{6.448776in}}%
\pgfpathquadraticcurveto{\pgfqpoint{7.247888in}{6.441161in}}{\pgfqpoint{7.247888in}{6.434071in}}%
\pgfpathquadraticcurveto{\pgfqpoint{7.247888in}{6.421777in}}{\pgfqpoint{7.252662in}{6.415093in}}%
\pgfpathquadraticcurveto{\pgfqpoint{7.257436in}{6.408409in}}{\pgfqpoint{7.266245in}{6.408409in}}%
\pgfpathquadraticcurveto{\pgfqpoint{7.273645in}{6.408409in}}{\pgfqpoint{7.277417in}{6.412849in}}%
\pgfpathquadraticcurveto{\pgfqpoint{7.281164in}{6.417289in}}{\pgfqpoint{7.283480in}{6.429679in}}%
\pgfpathlineto{\pgfqpoint{7.285342in}{6.438774in}}%
\pgfpathquadraticcurveto{\pgfqpoint{7.288565in}{6.455627in}}{\pgfqpoint{7.296657in}{6.463648in}}%
\pgfpathquadraticcurveto{\pgfqpoint{7.304750in}{6.471669in}}{\pgfqpoint{7.318309in}{6.471669in}}%
\pgfpathquadraticcurveto{\pgfqpoint{7.334517in}{6.471669in}}{\pgfqpoint{7.342872in}{6.460807in}}%
\pgfpathquadraticcurveto{\pgfqpoint{7.351227in}{6.449970in}}{\pgfqpoint{7.351227in}{6.429010in}}%
\pgfpathquadraticcurveto{\pgfqpoint{7.351227in}{6.421109in}}{\pgfqpoint{7.349437in}{6.412181in}}%
\pgfpathquadraticcurveto{\pgfqpoint{7.347647in}{6.403277in}}{\pgfqpoint{7.344138in}{6.393728in}}%
\pgfpathlineto{\pgfqpoint{7.328621in}{6.393728in}}%
\pgfpathquadraticcurveto{\pgfqpoint{7.333753in}{6.402895in}}{\pgfqpoint{7.336379in}{6.411704in}}%
\pgfpathquadraticcurveto{\pgfqpoint{7.338981in}{6.420512in}}{\pgfqpoint{7.338981in}{6.429010in}}%
\pgfpathquadraticcurveto{\pgfqpoint{7.338981in}{6.441901in}}{\pgfqpoint{7.333921in}{6.448919in}}%
\pgfpathquadraticcurveto{\pgfqpoint{7.328836in}{6.455938in}}{\pgfqpoint{7.319431in}{6.455938in}}%
\pgfpathquadraticcurveto{\pgfqpoint{7.311243in}{6.455938in}}{\pgfqpoint{7.306612in}{6.450901in}}%
\pgfpathquadraticcurveto{\pgfqpoint{7.301980in}{6.445864in}}{\pgfqpoint{7.299665in}{6.434382in}}%
\pgfpathlineto{\pgfqpoint{7.297875in}{6.425191in}}%
\pgfpathquadraticcurveto{\pgfqpoint{7.294533in}{6.408338in}}{\pgfqpoint{7.287371in}{6.400794in}}%
\pgfpathquadraticcurveto{\pgfqpoint{7.280210in}{6.393275in}}{\pgfqpoint{7.267438in}{6.393275in}}%
\pgfpathquadraticcurveto{\pgfqpoint{7.252662in}{6.393275in}}{\pgfqpoint{7.244164in}{6.403683in}}%
\pgfpathquadraticcurveto{\pgfqpoint{7.235665in}{6.414091in}}{\pgfqpoint{7.235665in}{6.432352in}}%
\pgfpathquadraticcurveto{\pgfqpoint{7.235665in}{6.440206in}}{\pgfqpoint{7.237098in}{6.448322in}}%
\pgfpathquadraticcurveto{\pgfqpoint{7.238506in}{6.456463in}}{\pgfqpoint{7.241323in}{6.464961in}}%
\pgfpathlineto{\pgfqpoint{7.241323in}{6.464961in}}%
\pgfpathclose%
\pgfpathmoveto{\pgfqpoint{7.237671in}{6.495182in}}%
\pgfpathlineto{\pgfqpoint{7.237671in}{6.565603in}}%
\pgfpathlineto{\pgfqpoint{7.250370in}{6.565603in}}%
\pgfpathlineto{\pgfqpoint{7.250370in}{6.510245in}}%
\pgfpathlineto{\pgfqpoint{7.283337in}{6.510245in}}%
\pgfpathlineto{\pgfqpoint{7.283337in}{6.563288in}}%
\pgfpathlineto{\pgfqpoint{7.296013in}{6.563288in}}%
\pgfpathlineto{\pgfqpoint{7.296013in}{6.510245in}}%
\pgfpathlineto{\pgfqpoint{7.336379in}{6.510245in}}%
\pgfpathlineto{\pgfqpoint{7.336379in}{6.566940in}}%
\pgfpathlineto{\pgfqpoint{7.349055in}{6.566940in}}%
\pgfpathlineto{\pgfqpoint{7.349055in}{6.495182in}}%
\pgfpathlineto{\pgfqpoint{7.237671in}{6.495182in}}%
\pgfpathlineto{\pgfqpoint{7.237671in}{6.495182in}}%
\pgfpathclose%
\pgfpathmoveto{\pgfqpoint{7.267964in}{6.644379in}}%
\pgfpathlineto{\pgfqpoint{7.280950in}{6.644379in}}%
\pgfpathquadraticcurveto{\pgfqpoint{7.277966in}{6.638579in}}{\pgfqpoint{7.276486in}{6.632300in}}%
\pgfpathquadraticcurveto{\pgfqpoint{7.274982in}{6.626046in}}{\pgfqpoint{7.274982in}{6.619314in}}%
\pgfpathquadraticcurveto{\pgfqpoint{7.274982in}{6.609097in}}{\pgfqpoint{7.278109in}{6.603989in}}%
\pgfpathquadraticcurveto{\pgfqpoint{7.281236in}{6.598880in}}{\pgfqpoint{7.287514in}{6.598880in}}%
\pgfpathquadraticcurveto{\pgfqpoint{7.292289in}{6.598880in}}{\pgfqpoint{7.295010in}{6.602533in}}%
\pgfpathquadraticcurveto{\pgfqpoint{7.297731in}{6.606185in}}{\pgfqpoint{7.300190in}{6.617237in}}%
\pgfpathlineto{\pgfqpoint{7.301240in}{6.621940in}}%
\pgfpathquadraticcurveto{\pgfqpoint{7.304368in}{6.636549in}}{\pgfqpoint{7.310073in}{6.642708in}}%
\pgfpathquadraticcurveto{\pgfqpoint{7.315778in}{6.648867in}}{\pgfqpoint{7.325995in}{6.648867in}}%
\pgfpathquadraticcurveto{\pgfqpoint{7.337645in}{6.648867in}}{\pgfqpoint{7.344448in}{6.639653in}}%
\pgfpathquadraticcurveto{\pgfqpoint{7.351227in}{6.630438in}}{\pgfqpoint{7.351227in}{6.614325in}}%
\pgfpathquadraticcurveto{\pgfqpoint{7.351227in}{6.607617in}}{\pgfqpoint{7.349914in}{6.600336in}}%
\pgfpathquadraticcurveto{\pgfqpoint{7.348602in}{6.593056in}}{\pgfqpoint{7.346000in}{6.585011in}}%
\pgfpathlineto{\pgfqpoint{7.331820in}{6.585011in}}%
\pgfpathquadraticcurveto{\pgfqpoint{7.335783in}{6.592626in}}{\pgfqpoint{7.337764in}{6.600002in}}%
\pgfpathquadraticcurveto{\pgfqpoint{7.339721in}{6.607378in}}{\pgfqpoint{7.339721in}{6.614635in}}%
\pgfpathquadraticcurveto{\pgfqpoint{7.339721in}{6.624327in}}{\pgfqpoint{7.336403in}{6.629531in}}%
\pgfpathquadraticcurveto{\pgfqpoint{7.333085in}{6.634759in}}{\pgfqpoint{7.327046in}{6.634759in}}%
\pgfpathquadraticcurveto{\pgfqpoint{7.321460in}{6.634759in}}{\pgfqpoint{7.318476in}{6.630987in}}%
\pgfpathquadraticcurveto{\pgfqpoint{7.315492in}{6.627240in}}{\pgfqpoint{7.312723in}{6.614468in}}%
\pgfpathlineto{\pgfqpoint{7.311601in}{6.609694in}}%
\pgfpathquadraticcurveto{\pgfqpoint{7.308927in}{6.596947in}}{\pgfqpoint{7.303365in}{6.591265in}}%
\pgfpathquadraticcurveto{\pgfqpoint{7.297803in}{6.585608in}}{\pgfqpoint{7.288111in}{6.585608in}}%
\pgfpathquadraticcurveto{\pgfqpoint{7.276319in}{6.585608in}}{\pgfqpoint{7.269921in}{6.593963in}}%
\pgfpathquadraticcurveto{\pgfqpoint{7.263500in}{6.602318in}}{\pgfqpoint{7.263500in}{6.617691in}}%
\pgfpathquadraticcurveto{\pgfqpoint{7.263500in}{6.625282in}}{\pgfqpoint{7.264622in}{6.631990in}}%
\pgfpathquadraticcurveto{\pgfqpoint{7.265720in}{6.638722in}}{\pgfqpoint{7.267964in}{6.644379in}}%
\pgfpathlineto{\pgfqpoint{7.267964in}{6.644379in}}%
\pgfpathclose%
\pgfpathmoveto{\pgfqpoint{7.233135in}{6.668561in}}%
\pgfpathlineto{\pgfqpoint{7.233135in}{6.680497in}}%
\pgfpathquadraticcurveto{\pgfqpoint{7.250728in}{6.691669in}}{\pgfqpoint{7.267605in}{6.697231in}}%
\pgfpathquadraticcurveto{\pgfqpoint{7.284459in}{6.702793in}}{\pgfqpoint{7.301097in}{6.702793in}}%
\pgfpathquadraticcurveto{\pgfqpoint{7.317807in}{6.702793in}}{\pgfqpoint{7.334732in}{6.697231in}}%
\pgfpathquadraticcurveto{\pgfqpoint{7.351657in}{6.691669in}}{\pgfqpoint{7.369203in}{6.680497in}}%
\pgfpathlineto{\pgfqpoint{7.369203in}{6.668561in}}%
\pgfpathquadraticcurveto{\pgfqpoint{7.352111in}{6.678468in}}{\pgfqpoint{7.335210in}{6.683362in}}%
\pgfpathquadraticcurveto{\pgfqpoint{7.318309in}{6.688255in}}{\pgfqpoint{7.301097in}{6.688255in}}%
\pgfpathquadraticcurveto{\pgfqpoint{7.283862in}{6.688255in}}{\pgfqpoint{7.267080in}{6.683362in}}%
\pgfpathquadraticcurveto{\pgfqpoint{7.250275in}{6.678468in}}{\pgfqpoint{7.233135in}{6.668561in}}%
\pgfpathlineto{\pgfqpoint{7.233135in}{6.668561in}}%
\pgfpathclose%
\pgfusepath{fill}%
\end{pgfscope}%
\begin{pgfscope}%
\pgfpathrectangle{\pgfqpoint{7.904091in}{2.928000in}}{\pgfqpoint{5.664000in}{5.664000in}}%
\pgfusepath{clip}%
\pgfsetrectcap%
\pgfsetroundjoin%
\pgfsetlinewidth{0.853187pt}%
\definecolor{currentstroke}{rgb}{0.745098,0.776471,0.803922}%
\pgfsetstrokecolor{currentstroke}%
\pgfsetdash{}{0pt}%
\pgfpathmoveto{\pgfqpoint{7.904091in}{5.760000in}}%
\pgfpathlineto{\pgfqpoint{13.568091in}{5.760000in}}%
\pgfusepath{stroke}%
\end{pgfscope}%
\begin{pgfscope}%
\pgfpathrectangle{\pgfqpoint{7.904091in}{2.928000in}}{\pgfqpoint{5.664000in}{5.664000in}}%
\pgfusepath{clip}%
\pgfsetrectcap%
\pgfsetroundjoin%
\pgfsetlinewidth{0.853187pt}%
\definecolor{currentstroke}{rgb}{0.745098,0.776471,0.803922}%
\pgfsetstrokecolor{currentstroke}%
\pgfsetdash{}{0pt}%
\pgfpathmoveto{\pgfqpoint{10.736091in}{2.928000in}}%
\pgfpathlineto{\pgfqpoint{10.736091in}{8.592000in}}%
\pgfusepath{stroke}%
\end{pgfscope}%
\begin{pgfscope}%
\pgfpathrectangle{\pgfqpoint{7.904091in}{2.928000in}}{\pgfqpoint{5.664000in}{5.664000in}}%
\pgfusepath{clip}%
\pgfsetbuttcap%
\pgfsetroundjoin%
\pgfsetlinewidth{1.104125pt}%
\definecolor{currentstroke}{rgb}{0.203922,0.243137,0.278431}%
\pgfsetstrokecolor{currentstroke}%
\pgfsetdash{{5.500000pt}{4.400000pt}}{0.000000pt}%
\pgfpathmoveto{\pgfqpoint{7.904091in}{2.928000in}}%
\pgfpathlineto{\pgfqpoint{13.568091in}{8.592000in}}%
\pgfusepath{stroke}%
\end{pgfscope}%
\begin{pgfscope}%
\pgfsetrectcap%
\pgfsetmiterjoin%
\pgfsetlinewidth{0.803000pt}%
\definecolor{currentstroke}{rgb}{0.627451,0.670588,0.705882}%
\pgfsetstrokecolor{currentstroke}%
\pgfsetdash{}{0pt}%
\pgfpathmoveto{\pgfqpoint{7.904091in}{2.928000in}}%
\pgfpathlineto{\pgfqpoint{7.904091in}{8.592000in}}%
\pgfusepath{stroke}%
\end{pgfscope}%
\begin{pgfscope}%
\pgfsetrectcap%
\pgfsetmiterjoin%
\pgfsetlinewidth{0.803000pt}%
\definecolor{currentstroke}{rgb}{0.627451,0.670588,0.705882}%
\pgfsetstrokecolor{currentstroke}%
\pgfsetdash{}{0pt}%
\pgfpathmoveto{\pgfqpoint{7.904091in}{2.928000in}}%
\pgfpathlineto{\pgfqpoint{13.568091in}{2.928000in}}%
\pgfusepath{stroke}%
\end{pgfscope}%
\begin{pgfscope}%
\pgfpathrectangle{\pgfqpoint{7.904091in}{2.928000in}}{\pgfqpoint{5.664000in}{5.664000in}}%
\pgfusepath{clip}%
\pgfsetbuttcap%
\pgfsetroundjoin%
\pgfsetlinewidth{1.003750pt}%
\definecolor{currentstroke}{rgb}{0.541176,0.380392,0.658824}%
\pgfsetstrokecolor{currentstroke}%
\pgfsetstrokeopacity{0.900000}%
\pgfsetdash{}{0pt}%
\pgfpathmoveto{\pgfqpoint{12.053442in}{7.491156in}}%
\pgfpathcurveto{\pgfqpoint{12.063696in}{7.491156in}}{\pgfqpoint{12.073532in}{7.495230in}}{\pgfqpoint{12.080782in}{7.502481in}}%
\pgfpathcurveto{\pgfqpoint{12.088033in}{7.509732in}}{\pgfqpoint{12.092107in}{7.519567in}}{\pgfqpoint{12.092107in}{7.529821in}}%
\pgfpathcurveto{\pgfqpoint{12.092107in}{7.540075in}}{\pgfqpoint{12.088033in}{7.549911in}}{\pgfqpoint{12.080782in}{7.557161in}}%
\pgfpathcurveto{\pgfqpoint{12.073532in}{7.564412in}}{\pgfqpoint{12.063696in}{7.568486in}}{\pgfqpoint{12.053442in}{7.568486in}}%
\pgfpathcurveto{\pgfqpoint{12.043188in}{7.568486in}}{\pgfqpoint{12.033353in}{7.564412in}}{\pgfqpoint{12.026102in}{7.557161in}}%
\pgfpathcurveto{\pgfqpoint{12.018851in}{7.549911in}}{\pgfqpoint{12.014777in}{7.540075in}}{\pgfqpoint{12.014777in}{7.529821in}}%
\pgfpathcurveto{\pgfqpoint{12.014777in}{7.519567in}}{\pgfqpoint{12.018851in}{7.509732in}}{\pgfqpoint{12.026102in}{7.502481in}}%
\pgfpathcurveto{\pgfqpoint{12.033353in}{7.495230in}}{\pgfqpoint{12.043188in}{7.491156in}}{\pgfqpoint{12.053442in}{7.491156in}}%
\pgfpathlineto{\pgfqpoint{12.053442in}{7.491156in}}%
\pgfpathclose%
\pgfusepath{stroke}%
\end{pgfscope}%
\begin{pgfscope}%
\pgfpathrectangle{\pgfqpoint{7.904091in}{2.928000in}}{\pgfqpoint{5.664000in}{5.664000in}}%
\pgfusepath{clip}%
\pgfsetbuttcap%
\pgfsetroundjoin%
\pgfsetlinewidth{1.003750pt}%
\definecolor{currentstroke}{rgb}{0.541176,0.380392,0.658824}%
\pgfsetstrokecolor{currentstroke}%
\pgfsetstrokeopacity{0.900000}%
\pgfsetdash{}{0pt}%
\pgfpathmoveto{\pgfqpoint{11.936104in}{7.288882in}}%
\pgfpathcurveto{\pgfqpoint{11.946358in}{7.288882in}}{\pgfqpoint{11.956193in}{7.292956in}}{\pgfqpoint{11.963444in}{7.300206in}}%
\pgfpathcurveto{\pgfqpoint{11.970695in}{7.307457in}}{\pgfqpoint{11.974769in}{7.317292in}}{\pgfqpoint{11.974769in}{7.327547in}}%
\pgfpathcurveto{\pgfqpoint{11.974769in}{7.337801in}}{\pgfqpoint{11.970695in}{7.347636in}}{\pgfqpoint{11.963444in}{7.354887in}}%
\pgfpathcurveto{\pgfqpoint{11.956193in}{7.362138in}}{\pgfqpoint{11.946358in}{7.366212in}}{\pgfqpoint{11.936104in}{7.366212in}}%
\pgfpathcurveto{\pgfqpoint{11.925849in}{7.366212in}}{\pgfqpoint{11.916014in}{7.362138in}}{\pgfqpoint{11.908763in}{7.354887in}}%
\pgfpathcurveto{\pgfqpoint{11.901512in}{7.347636in}}{\pgfqpoint{11.897438in}{7.337801in}}{\pgfqpoint{11.897438in}{7.327547in}}%
\pgfpathcurveto{\pgfqpoint{11.897438in}{7.317292in}}{\pgfqpoint{11.901512in}{7.307457in}}{\pgfqpoint{11.908763in}{7.300206in}}%
\pgfpathcurveto{\pgfqpoint{11.916014in}{7.292956in}}{\pgfqpoint{11.925849in}{7.288882in}}{\pgfqpoint{11.936104in}{7.288882in}}%
\pgfpathlineto{\pgfqpoint{11.936104in}{7.288882in}}%
\pgfpathclose%
\pgfusepath{stroke}%
\end{pgfscope}%
\begin{pgfscope}%
\pgfpathrectangle{\pgfqpoint{7.904091in}{2.928000in}}{\pgfqpoint{5.664000in}{5.664000in}}%
\pgfusepath{clip}%
\pgfsetbuttcap%
\pgfsetroundjoin%
\pgfsetlinewidth{1.003750pt}%
\definecolor{currentstroke}{rgb}{0.541176,0.380392,0.658824}%
\pgfsetstrokecolor{currentstroke}%
\pgfsetstrokeopacity{0.900000}%
\pgfsetdash{}{0pt}%
\pgfpathmoveto{\pgfqpoint{12.937089in}{7.608684in}}%
\pgfpathcurveto{\pgfqpoint{12.947343in}{7.608684in}}{\pgfqpoint{12.957179in}{7.612758in}}{\pgfqpoint{12.964430in}{7.620008in}}%
\pgfpathcurveto{\pgfqpoint{12.971680in}{7.627259in}}{\pgfqpoint{12.975754in}{7.637095in}}{\pgfqpoint{12.975754in}{7.647349in}}%
\pgfpathcurveto{\pgfqpoint{12.975754in}{7.657603in}}{\pgfqpoint{12.971680in}{7.667438in}}{\pgfqpoint{12.964430in}{7.674689in}}%
\pgfpathcurveto{\pgfqpoint{12.957179in}{7.681940in}}{\pgfqpoint{12.947343in}{7.686014in}}{\pgfqpoint{12.937089in}{7.686014in}}%
\pgfpathcurveto{\pgfqpoint{12.926835in}{7.686014in}}{\pgfqpoint{12.917000in}{7.681940in}}{\pgfqpoint{12.909749in}{7.674689in}}%
\pgfpathcurveto{\pgfqpoint{12.902498in}{7.667438in}}{\pgfqpoint{12.898424in}{7.657603in}}{\pgfqpoint{12.898424in}{7.647349in}}%
\pgfpathcurveto{\pgfqpoint{12.898424in}{7.637095in}}{\pgfqpoint{12.902498in}{7.627259in}}{\pgfqpoint{12.909749in}{7.620008in}}%
\pgfpathcurveto{\pgfqpoint{12.917000in}{7.612758in}}{\pgfqpoint{12.926835in}{7.608684in}}{\pgfqpoint{12.937089in}{7.608684in}}%
\pgfpathlineto{\pgfqpoint{12.937089in}{7.608684in}}%
\pgfpathclose%
\pgfusepath{stroke}%
\end{pgfscope}%
\begin{pgfscope}%
\pgfpathrectangle{\pgfqpoint{7.904091in}{2.928000in}}{\pgfqpoint{5.664000in}{5.664000in}}%
\pgfusepath{clip}%
\pgfsetbuttcap%
\pgfsetroundjoin%
\pgfsetlinewidth{1.003750pt}%
\definecolor{currentstroke}{rgb}{0.541176,0.380392,0.658824}%
\pgfsetstrokecolor{currentstroke}%
\pgfsetstrokeopacity{0.900000}%
\pgfsetdash{}{0pt}%
\pgfpathmoveto{\pgfqpoint{11.867595in}{6.528435in}}%
\pgfpathcurveto{\pgfqpoint{11.877849in}{6.528435in}}{\pgfqpoint{11.887685in}{6.532509in}}{\pgfqpoint{11.894935in}{6.539760in}}%
\pgfpathcurveto{\pgfqpoint{11.902186in}{6.547011in}}{\pgfqpoint{11.906260in}{6.556846in}}{\pgfqpoint{11.906260in}{6.567100in}}%
\pgfpathcurveto{\pgfqpoint{11.906260in}{6.577355in}}{\pgfqpoint{11.902186in}{6.587190in}}{\pgfqpoint{11.894935in}{6.594441in}}%
\pgfpathcurveto{\pgfqpoint{11.887685in}{6.601691in}}{\pgfqpoint{11.877849in}{6.605765in}}{\pgfqpoint{11.867595in}{6.605765in}}%
\pgfpathcurveto{\pgfqpoint{11.857341in}{6.605765in}}{\pgfqpoint{11.847505in}{6.601691in}}{\pgfqpoint{11.840255in}{6.594441in}}%
\pgfpathcurveto{\pgfqpoint{11.833004in}{6.587190in}}{\pgfqpoint{11.828930in}{6.577355in}}{\pgfqpoint{11.828930in}{6.567100in}}%
\pgfpathcurveto{\pgfqpoint{11.828930in}{6.556846in}}{\pgfqpoint{11.833004in}{6.547011in}}{\pgfqpoint{11.840255in}{6.539760in}}%
\pgfpathcurveto{\pgfqpoint{11.847505in}{6.532509in}}{\pgfqpoint{11.857341in}{6.528435in}}{\pgfqpoint{11.867595in}{6.528435in}}%
\pgfpathlineto{\pgfqpoint{11.867595in}{6.528435in}}%
\pgfpathclose%
\pgfusepath{stroke}%
\end{pgfscope}%
\begin{pgfscope}%
\pgfpathrectangle{\pgfqpoint{7.904091in}{2.928000in}}{\pgfqpoint{5.664000in}{5.664000in}}%
\pgfusepath{clip}%
\pgfsetbuttcap%
\pgfsetroundjoin%
\pgfsetlinewidth{1.003750pt}%
\definecolor{currentstroke}{rgb}{0.541176,0.380392,0.658824}%
\pgfsetstrokecolor{currentstroke}%
\pgfsetstrokeopacity{0.900000}%
\pgfsetdash{}{0pt}%
\pgfpathmoveto{\pgfqpoint{12.608792in}{7.778714in}}%
\pgfpathcurveto{\pgfqpoint{12.619046in}{7.778714in}}{\pgfqpoint{12.628881in}{7.782788in}}{\pgfqpoint{12.636132in}{7.790038in}}%
\pgfpathcurveto{\pgfqpoint{12.643383in}{7.797289in}}{\pgfqpoint{12.647457in}{7.807125in}}{\pgfqpoint{12.647457in}{7.817379in}}%
\pgfpathcurveto{\pgfqpoint{12.647457in}{7.827633in}}{\pgfqpoint{12.643383in}{7.837468in}}{\pgfqpoint{12.636132in}{7.844719in}}%
\pgfpathcurveto{\pgfqpoint{12.628881in}{7.851970in}}{\pgfqpoint{12.619046in}{7.856044in}}{\pgfqpoint{12.608792in}{7.856044in}}%
\pgfpathcurveto{\pgfqpoint{12.598538in}{7.856044in}}{\pgfqpoint{12.588702in}{7.851970in}}{\pgfqpoint{12.581452in}{7.844719in}}%
\pgfpathcurveto{\pgfqpoint{12.574201in}{7.837468in}}{\pgfqpoint{12.570127in}{7.827633in}}{\pgfqpoint{12.570127in}{7.817379in}}%
\pgfpathcurveto{\pgfqpoint{12.570127in}{7.807125in}}{\pgfqpoint{12.574201in}{7.797289in}}{\pgfqpoint{12.581452in}{7.790038in}}%
\pgfpathcurveto{\pgfqpoint{12.588702in}{7.782788in}}{\pgfqpoint{12.598538in}{7.778714in}}{\pgfqpoint{12.608792in}{7.778714in}}%
\pgfpathlineto{\pgfqpoint{12.608792in}{7.778714in}}%
\pgfpathclose%
\pgfusepath{stroke}%
\end{pgfscope}%
\begin{pgfscope}%
\pgfpathrectangle{\pgfqpoint{7.904091in}{2.928000in}}{\pgfqpoint{5.664000in}{5.664000in}}%
\pgfusepath{clip}%
\pgfsetbuttcap%
\pgfsetroundjoin%
\pgfsetlinewidth{1.003750pt}%
\definecolor{currentstroke}{rgb}{0.541176,0.380392,0.658824}%
\pgfsetstrokecolor{currentstroke}%
\pgfsetstrokeopacity{0.900000}%
\pgfsetdash{}{0pt}%
\pgfpathmoveto{\pgfqpoint{11.660662in}{6.350161in}}%
\pgfpathcurveto{\pgfqpoint{11.670916in}{6.350161in}}{\pgfqpoint{11.680751in}{6.354235in}}{\pgfqpoint{11.688002in}{6.361486in}}%
\pgfpathcurveto{\pgfqpoint{11.695253in}{6.368737in}}{\pgfqpoint{11.699327in}{6.378572in}}{\pgfqpoint{11.699327in}{6.388826in}}%
\pgfpathcurveto{\pgfqpoint{11.699327in}{6.399080in}}{\pgfqpoint{11.695253in}{6.408916in}}{\pgfqpoint{11.688002in}{6.416166in}}%
\pgfpathcurveto{\pgfqpoint{11.680751in}{6.423417in}}{\pgfqpoint{11.670916in}{6.427491in}}{\pgfqpoint{11.660662in}{6.427491in}}%
\pgfpathcurveto{\pgfqpoint{11.650407in}{6.427491in}}{\pgfqpoint{11.640572in}{6.423417in}}{\pgfqpoint{11.633321in}{6.416166in}}%
\pgfpathcurveto{\pgfqpoint{11.626071in}{6.408916in}}{\pgfqpoint{11.621997in}{6.399080in}}{\pgfqpoint{11.621997in}{6.388826in}}%
\pgfpathcurveto{\pgfqpoint{11.621997in}{6.378572in}}{\pgfqpoint{11.626071in}{6.368737in}}{\pgfqpoint{11.633321in}{6.361486in}}%
\pgfpathcurveto{\pgfqpoint{11.640572in}{6.354235in}}{\pgfqpoint{11.650407in}{6.350161in}}{\pgfqpoint{11.660662in}{6.350161in}}%
\pgfpathlineto{\pgfqpoint{11.660662in}{6.350161in}}%
\pgfpathclose%
\pgfusepath{stroke}%
\end{pgfscope}%
\begin{pgfscope}%
\pgfpathrectangle{\pgfqpoint{7.904091in}{2.928000in}}{\pgfqpoint{5.664000in}{5.664000in}}%
\pgfusepath{clip}%
\pgfsetbuttcap%
\pgfsetroundjoin%
\pgfsetlinewidth{1.003750pt}%
\definecolor{currentstroke}{rgb}{0.541176,0.380392,0.658824}%
\pgfsetstrokecolor{currentstroke}%
\pgfsetstrokeopacity{0.900000}%
\pgfsetdash{}{0pt}%
\pgfpathmoveto{\pgfqpoint{12.515834in}{7.441492in}}%
\pgfpathcurveto{\pgfqpoint{12.526088in}{7.441492in}}{\pgfqpoint{12.535924in}{7.445566in}}{\pgfqpoint{12.543174in}{7.452816in}}%
\pgfpathcurveto{\pgfqpoint{12.550425in}{7.460067in}}{\pgfqpoint{12.554499in}{7.469903in}}{\pgfqpoint{12.554499in}{7.480157in}}%
\pgfpathcurveto{\pgfqpoint{12.554499in}{7.490411in}}{\pgfqpoint{12.550425in}{7.500246in}}{\pgfqpoint{12.543174in}{7.507497in}}%
\pgfpathcurveto{\pgfqpoint{12.535924in}{7.514748in}}{\pgfqpoint{12.526088in}{7.518822in}}{\pgfqpoint{12.515834in}{7.518822in}}%
\pgfpathcurveto{\pgfqpoint{12.505580in}{7.518822in}}{\pgfqpoint{12.495744in}{7.514748in}}{\pgfqpoint{12.488494in}{7.507497in}}%
\pgfpathcurveto{\pgfqpoint{12.481243in}{7.500246in}}{\pgfqpoint{12.477169in}{7.490411in}}{\pgfqpoint{12.477169in}{7.480157in}}%
\pgfpathcurveto{\pgfqpoint{12.477169in}{7.469903in}}{\pgfqpoint{12.481243in}{7.460067in}}{\pgfqpoint{12.488494in}{7.452816in}}%
\pgfpathcurveto{\pgfqpoint{12.495744in}{7.445566in}}{\pgfqpoint{12.505580in}{7.441492in}}{\pgfqpoint{12.515834in}{7.441492in}}%
\pgfpathlineto{\pgfqpoint{12.515834in}{7.441492in}}%
\pgfpathclose%
\pgfusepath{stroke}%
\end{pgfscope}%
\begin{pgfscope}%
\pgfpathrectangle{\pgfqpoint{7.904091in}{2.928000in}}{\pgfqpoint{5.664000in}{5.664000in}}%
\pgfusepath{clip}%
\pgfsetbuttcap%
\pgfsetroundjoin%
\pgfsetlinewidth{1.003750pt}%
\definecolor{currentstroke}{rgb}{0.541176,0.380392,0.658824}%
\pgfsetstrokecolor{currentstroke}%
\pgfsetstrokeopacity{0.900000}%
\pgfsetdash{}{0pt}%
\pgfpathmoveto{\pgfqpoint{10.857126in}{6.120525in}}%
\pgfpathcurveto{\pgfqpoint{10.867381in}{6.120525in}}{\pgfqpoint{10.877216in}{6.124599in}}{\pgfqpoint{10.884467in}{6.131849in}}%
\pgfpathcurveto{\pgfqpoint{10.891717in}{6.139100in}}{\pgfqpoint{10.895791in}{6.148936in}}{\pgfqpoint{10.895791in}{6.159190in}}%
\pgfpathcurveto{\pgfqpoint{10.895791in}{6.169444in}}{\pgfqpoint{10.891717in}{6.179279in}}{\pgfqpoint{10.884467in}{6.186530in}}%
\pgfpathcurveto{\pgfqpoint{10.877216in}{6.193781in}}{\pgfqpoint{10.867381in}{6.197855in}}{\pgfqpoint{10.857126in}{6.197855in}}%
\pgfpathcurveto{\pgfqpoint{10.846872in}{6.197855in}}{\pgfqpoint{10.837037in}{6.193781in}}{\pgfqpoint{10.829786in}{6.186530in}}%
\pgfpathcurveto{\pgfqpoint{10.822535in}{6.179279in}}{\pgfqpoint{10.818461in}{6.169444in}}{\pgfqpoint{10.818461in}{6.159190in}}%
\pgfpathcurveto{\pgfqpoint{10.818461in}{6.148936in}}{\pgfqpoint{10.822535in}{6.139100in}}{\pgfqpoint{10.829786in}{6.131849in}}%
\pgfpathcurveto{\pgfqpoint{10.837037in}{6.124599in}}{\pgfqpoint{10.846872in}{6.120525in}}{\pgfqpoint{10.857126in}{6.120525in}}%
\pgfpathlineto{\pgfqpoint{10.857126in}{6.120525in}}%
\pgfpathclose%
\pgfusepath{stroke}%
\end{pgfscope}%
\begin{pgfscope}%
\pgfpathrectangle{\pgfqpoint{7.904091in}{2.928000in}}{\pgfqpoint{5.664000in}{5.664000in}}%
\pgfusepath{clip}%
\pgfsetbuttcap%
\pgfsetroundjoin%
\pgfsetlinewidth{1.003750pt}%
\definecolor{currentstroke}{rgb}{0.541176,0.380392,0.658824}%
\pgfsetstrokecolor{currentstroke}%
\pgfsetstrokeopacity{0.900000}%
\pgfsetdash{}{0pt}%
\pgfpathmoveto{\pgfqpoint{10.964545in}{5.419147in}}%
\pgfpathcurveto{\pgfqpoint{10.974799in}{5.419147in}}{\pgfqpoint{10.984634in}{5.423221in}}{\pgfqpoint{10.991885in}{5.430472in}}%
\pgfpathcurveto{\pgfqpoint{10.999136in}{5.437723in}}{\pgfqpoint{11.003210in}{5.447558in}}{\pgfqpoint{11.003210in}{5.457812in}}%
\pgfpathcurveto{\pgfqpoint{11.003210in}{5.468066in}}{\pgfqpoint{10.999136in}{5.477902in}}{\pgfqpoint{10.991885in}{5.485153in}}%
\pgfpathcurveto{\pgfqpoint{10.984634in}{5.492403in}}{\pgfqpoint{10.974799in}{5.496477in}}{\pgfqpoint{10.964545in}{5.496477in}}%
\pgfpathcurveto{\pgfqpoint{10.954291in}{5.496477in}}{\pgfqpoint{10.944455in}{5.492403in}}{\pgfqpoint{10.937204in}{5.485153in}}%
\pgfpathcurveto{\pgfqpoint{10.929954in}{5.477902in}}{\pgfqpoint{10.925880in}{5.468066in}}{\pgfqpoint{10.925880in}{5.457812in}}%
\pgfpathcurveto{\pgfqpoint{10.925880in}{5.447558in}}{\pgfqpoint{10.929954in}{5.437723in}}{\pgfqpoint{10.937204in}{5.430472in}}%
\pgfpathcurveto{\pgfqpoint{10.944455in}{5.423221in}}{\pgfqpoint{10.954291in}{5.419147in}}{\pgfqpoint{10.964545in}{5.419147in}}%
\pgfpathlineto{\pgfqpoint{10.964545in}{5.419147in}}%
\pgfpathclose%
\pgfusepath{stroke}%
\end{pgfscope}%
\begin{pgfscope}%
\pgfpathrectangle{\pgfqpoint{7.904091in}{2.928000in}}{\pgfqpoint{5.664000in}{5.664000in}}%
\pgfusepath{clip}%
\pgfsetbuttcap%
\pgfsetroundjoin%
\pgfsetlinewidth{1.003750pt}%
\definecolor{currentstroke}{rgb}{0.541176,0.380392,0.658824}%
\pgfsetstrokecolor{currentstroke}%
\pgfsetstrokeopacity{0.900000}%
\pgfsetdash{}{0pt}%
\pgfpathmoveto{\pgfqpoint{11.737457in}{6.557479in}}%
\pgfpathcurveto{\pgfqpoint{11.747711in}{6.557479in}}{\pgfqpoint{11.757547in}{6.561553in}}{\pgfqpoint{11.764797in}{6.568803in}}%
\pgfpathcurveto{\pgfqpoint{11.772048in}{6.576054in}}{\pgfqpoint{11.776122in}{6.585890in}}{\pgfqpoint{11.776122in}{6.596144in}}%
\pgfpathcurveto{\pgfqpoint{11.776122in}{6.606398in}}{\pgfqpoint{11.772048in}{6.616233in}}{\pgfqpoint{11.764797in}{6.623484in}}%
\pgfpathcurveto{\pgfqpoint{11.757547in}{6.630735in}}{\pgfqpoint{11.747711in}{6.634809in}}{\pgfqpoint{11.737457in}{6.634809in}}%
\pgfpathcurveto{\pgfqpoint{11.727203in}{6.634809in}}{\pgfqpoint{11.717367in}{6.630735in}}{\pgfqpoint{11.710117in}{6.623484in}}%
\pgfpathcurveto{\pgfqpoint{11.702866in}{6.616233in}}{\pgfqpoint{11.698792in}{6.606398in}}{\pgfqpoint{11.698792in}{6.596144in}}%
\pgfpathcurveto{\pgfqpoint{11.698792in}{6.585890in}}{\pgfqpoint{11.702866in}{6.576054in}}{\pgfqpoint{11.710117in}{6.568803in}}%
\pgfpathcurveto{\pgfqpoint{11.717367in}{6.561553in}}{\pgfqpoint{11.727203in}{6.557479in}}{\pgfqpoint{11.737457in}{6.557479in}}%
\pgfpathlineto{\pgfqpoint{11.737457in}{6.557479in}}%
\pgfpathclose%
\pgfusepath{stroke}%
\end{pgfscope}%
\begin{pgfscope}%
\pgfpathrectangle{\pgfqpoint{7.904091in}{2.928000in}}{\pgfqpoint{5.664000in}{5.664000in}}%
\pgfusepath{clip}%
\pgfsetbuttcap%
\pgfsetroundjoin%
\pgfsetlinewidth{1.003750pt}%
\definecolor{currentstroke}{rgb}{0.541176,0.380392,0.658824}%
\pgfsetstrokecolor{currentstroke}%
\pgfsetstrokeopacity{0.900000}%
\pgfsetdash{}{0pt}%
\pgfpathmoveto{\pgfqpoint{11.217406in}{5.874031in}}%
\pgfpathcurveto{\pgfqpoint{11.227660in}{5.874031in}}{\pgfqpoint{11.237496in}{5.878105in}}{\pgfqpoint{11.244746in}{5.885356in}}%
\pgfpathcurveto{\pgfqpoint{11.251997in}{5.892606in}}{\pgfqpoint{11.256071in}{5.902442in}}{\pgfqpoint{11.256071in}{5.912696in}}%
\pgfpathcurveto{\pgfqpoint{11.256071in}{5.922950in}}{\pgfqpoint{11.251997in}{5.932786in}}{\pgfqpoint{11.244746in}{5.940036in}}%
\pgfpathcurveto{\pgfqpoint{11.237496in}{5.947287in}}{\pgfqpoint{11.227660in}{5.951361in}}{\pgfqpoint{11.217406in}{5.951361in}}%
\pgfpathcurveto{\pgfqpoint{11.207152in}{5.951361in}}{\pgfqpoint{11.197317in}{5.947287in}}{\pgfqpoint{11.190066in}{5.940036in}}%
\pgfpathcurveto{\pgfqpoint{11.182815in}{5.932786in}}{\pgfqpoint{11.178741in}{5.922950in}}{\pgfqpoint{11.178741in}{5.912696in}}%
\pgfpathcurveto{\pgfqpoint{11.178741in}{5.902442in}}{\pgfqpoint{11.182815in}{5.892606in}}{\pgfqpoint{11.190066in}{5.885356in}}%
\pgfpathcurveto{\pgfqpoint{11.197317in}{5.878105in}}{\pgfqpoint{11.207152in}{5.874031in}}{\pgfqpoint{11.217406in}{5.874031in}}%
\pgfpathlineto{\pgfqpoint{11.217406in}{5.874031in}}%
\pgfpathclose%
\pgfusepath{stroke}%
\end{pgfscope}%
\begin{pgfscope}%
\pgfpathrectangle{\pgfqpoint{7.904091in}{2.928000in}}{\pgfqpoint{5.664000in}{5.664000in}}%
\pgfusepath{clip}%
\pgfsetbuttcap%
\pgfsetroundjoin%
\pgfsetlinewidth{1.003750pt}%
\definecolor{currentstroke}{rgb}{0.541176,0.380392,0.658824}%
\pgfsetstrokecolor{currentstroke}%
\pgfsetstrokeopacity{0.900000}%
\pgfsetdash{}{0pt}%
\pgfpathmoveto{\pgfqpoint{11.767633in}{6.971409in}}%
\pgfpathcurveto{\pgfqpoint{11.777887in}{6.971409in}}{\pgfqpoint{11.787722in}{6.975483in}}{\pgfqpoint{11.794973in}{6.982733in}}%
\pgfpathcurveto{\pgfqpoint{11.802224in}{6.989984in}}{\pgfqpoint{11.806298in}{6.999820in}}{\pgfqpoint{11.806298in}{7.010074in}}%
\pgfpathcurveto{\pgfqpoint{11.806298in}{7.020328in}}{\pgfqpoint{11.802224in}{7.030163in}}{\pgfqpoint{11.794973in}{7.037414in}}%
\pgfpathcurveto{\pgfqpoint{11.787722in}{7.044665in}}{\pgfqpoint{11.777887in}{7.048739in}}{\pgfqpoint{11.767633in}{7.048739in}}%
\pgfpathcurveto{\pgfqpoint{11.757379in}{7.048739in}}{\pgfqpoint{11.747543in}{7.044665in}}{\pgfqpoint{11.740292in}{7.037414in}}%
\pgfpathcurveto{\pgfqpoint{11.733042in}{7.030163in}}{\pgfqpoint{11.728968in}{7.020328in}}{\pgfqpoint{11.728968in}{7.010074in}}%
\pgfpathcurveto{\pgfqpoint{11.728968in}{6.999820in}}{\pgfqpoint{11.733042in}{6.989984in}}{\pgfqpoint{11.740292in}{6.982733in}}%
\pgfpathcurveto{\pgfqpoint{11.747543in}{6.975483in}}{\pgfqpoint{11.757379in}{6.971409in}}{\pgfqpoint{11.767633in}{6.971409in}}%
\pgfpathlineto{\pgfqpoint{11.767633in}{6.971409in}}%
\pgfpathclose%
\pgfusepath{stroke}%
\end{pgfscope}%
\begin{pgfscope}%
\pgfpathrectangle{\pgfqpoint{7.904091in}{2.928000in}}{\pgfqpoint{5.664000in}{5.664000in}}%
\pgfusepath{clip}%
\pgfsetbuttcap%
\pgfsetroundjoin%
\pgfsetlinewidth{1.003750pt}%
\definecolor{currentstroke}{rgb}{0.541176,0.380392,0.658824}%
\pgfsetstrokecolor{currentstroke}%
\pgfsetstrokeopacity{0.900000}%
\pgfsetdash{}{0pt}%
\pgfpathmoveto{\pgfqpoint{11.659531in}{6.643178in}}%
\pgfpathcurveto{\pgfqpoint{11.669786in}{6.643178in}}{\pgfqpoint{11.679621in}{6.647252in}}{\pgfqpoint{11.686872in}{6.654503in}}%
\pgfpathcurveto{\pgfqpoint{11.694122in}{6.661754in}}{\pgfqpoint{11.698196in}{6.671589in}}{\pgfqpoint{11.698196in}{6.681843in}}%
\pgfpathcurveto{\pgfqpoint{11.698196in}{6.692098in}}{\pgfqpoint{11.694122in}{6.701933in}}{\pgfqpoint{11.686872in}{6.709184in}}%
\pgfpathcurveto{\pgfqpoint{11.679621in}{6.716434in}}{\pgfqpoint{11.669786in}{6.720508in}}{\pgfqpoint{11.659531in}{6.720508in}}%
\pgfpathcurveto{\pgfqpoint{11.649277in}{6.720508in}}{\pgfqpoint{11.639442in}{6.716434in}}{\pgfqpoint{11.632191in}{6.709184in}}%
\pgfpathcurveto{\pgfqpoint{11.624940in}{6.701933in}}{\pgfqpoint{11.620866in}{6.692098in}}{\pgfqpoint{11.620866in}{6.681843in}}%
\pgfpathcurveto{\pgfqpoint{11.620866in}{6.671589in}}{\pgfqpoint{11.624940in}{6.661754in}}{\pgfqpoint{11.632191in}{6.654503in}}%
\pgfpathcurveto{\pgfqpoint{11.639442in}{6.647252in}}{\pgfqpoint{11.649277in}{6.643178in}}{\pgfqpoint{11.659531in}{6.643178in}}%
\pgfpathlineto{\pgfqpoint{11.659531in}{6.643178in}}%
\pgfpathclose%
\pgfusepath{stroke}%
\end{pgfscope}%
\begin{pgfscope}%
\pgfpathrectangle{\pgfqpoint{7.904091in}{2.928000in}}{\pgfqpoint{5.664000in}{5.664000in}}%
\pgfusepath{clip}%
\pgfsetbuttcap%
\pgfsetroundjoin%
\pgfsetlinewidth{1.003750pt}%
\definecolor{currentstroke}{rgb}{0.541176,0.380392,0.658824}%
\pgfsetstrokecolor{currentstroke}%
\pgfsetstrokeopacity{0.900000}%
\pgfsetdash{}{0pt}%
\pgfpathmoveto{\pgfqpoint{11.149490in}{6.366506in}}%
\pgfpathcurveto{\pgfqpoint{11.159744in}{6.366506in}}{\pgfqpoint{11.169579in}{6.370580in}}{\pgfqpoint{11.176830in}{6.377831in}}%
\pgfpathcurveto{\pgfqpoint{11.184081in}{6.385081in}}{\pgfqpoint{11.188155in}{6.394917in}}{\pgfqpoint{11.188155in}{6.405171in}}%
\pgfpathcurveto{\pgfqpoint{11.188155in}{6.415425in}}{\pgfqpoint{11.184081in}{6.425261in}}{\pgfqpoint{11.176830in}{6.432511in}}%
\pgfpathcurveto{\pgfqpoint{11.169579in}{6.439762in}}{\pgfqpoint{11.159744in}{6.443836in}}{\pgfqpoint{11.149490in}{6.443836in}}%
\pgfpathcurveto{\pgfqpoint{11.139236in}{6.443836in}}{\pgfqpoint{11.129400in}{6.439762in}}{\pgfqpoint{11.122149in}{6.432511in}}%
\pgfpathcurveto{\pgfqpoint{11.114899in}{6.425261in}}{\pgfqpoint{11.110825in}{6.415425in}}{\pgfqpoint{11.110825in}{6.405171in}}%
\pgfpathcurveto{\pgfqpoint{11.110825in}{6.394917in}}{\pgfqpoint{11.114899in}{6.385081in}}{\pgfqpoint{11.122149in}{6.377831in}}%
\pgfpathcurveto{\pgfqpoint{11.129400in}{6.370580in}}{\pgfqpoint{11.139236in}{6.366506in}}{\pgfqpoint{11.149490in}{6.366506in}}%
\pgfpathlineto{\pgfqpoint{11.149490in}{6.366506in}}%
\pgfpathclose%
\pgfusepath{stroke}%
\end{pgfscope}%
\begin{pgfscope}%
\pgfpathrectangle{\pgfqpoint{7.904091in}{2.928000in}}{\pgfqpoint{5.664000in}{5.664000in}}%
\pgfusepath{clip}%
\pgfsetbuttcap%
\pgfsetroundjoin%
\pgfsetlinewidth{1.003750pt}%
\definecolor{currentstroke}{rgb}{0.541176,0.380392,0.658824}%
\pgfsetstrokecolor{currentstroke}%
\pgfsetstrokeopacity{0.900000}%
\pgfsetdash{}{0pt}%
\pgfpathmoveto{\pgfqpoint{10.964150in}{5.389422in}}%
\pgfpathcurveto{\pgfqpoint{10.974404in}{5.389422in}}{\pgfqpoint{10.984239in}{5.393496in}}{\pgfqpoint{10.991490in}{5.400747in}}%
\pgfpathcurveto{\pgfqpoint{10.998741in}{5.407997in}}{\pgfqpoint{11.002815in}{5.417833in}}{\pgfqpoint{11.002815in}{5.428087in}}%
\pgfpathcurveto{\pgfqpoint{11.002815in}{5.438341in}}{\pgfqpoint{10.998741in}{5.448176in}}{\pgfqpoint{10.991490in}{5.455427in}}%
\pgfpathcurveto{\pgfqpoint{10.984239in}{5.462678in}}{\pgfqpoint{10.974404in}{5.466752in}}{\pgfqpoint{10.964150in}{5.466752in}}%
\pgfpathcurveto{\pgfqpoint{10.953896in}{5.466752in}}{\pgfqpoint{10.944060in}{5.462678in}}{\pgfqpoint{10.936810in}{5.455427in}}%
\pgfpathcurveto{\pgfqpoint{10.929559in}{5.448176in}}{\pgfqpoint{10.925485in}{5.438341in}}{\pgfqpoint{10.925485in}{5.428087in}}%
\pgfpathcurveto{\pgfqpoint{10.925485in}{5.417833in}}{\pgfqpoint{10.929559in}{5.407997in}}{\pgfqpoint{10.936810in}{5.400747in}}%
\pgfpathcurveto{\pgfqpoint{10.944060in}{5.393496in}}{\pgfqpoint{10.953896in}{5.389422in}}{\pgfqpoint{10.964150in}{5.389422in}}%
\pgfpathlineto{\pgfqpoint{10.964150in}{5.389422in}}%
\pgfpathclose%
\pgfusepath{stroke}%
\end{pgfscope}%
\begin{pgfscope}%
\pgfpathrectangle{\pgfqpoint{7.904091in}{2.928000in}}{\pgfqpoint{5.664000in}{5.664000in}}%
\pgfusepath{clip}%
\pgfsetbuttcap%
\pgfsetroundjoin%
\pgfsetlinewidth{1.003750pt}%
\definecolor{currentstroke}{rgb}{0.541176,0.380392,0.658824}%
\pgfsetstrokecolor{currentstroke}%
\pgfsetstrokeopacity{0.900000}%
\pgfsetdash{}{0pt}%
\pgfpathmoveto{\pgfqpoint{11.137361in}{6.434831in}}%
\pgfpathcurveto{\pgfqpoint{11.147616in}{6.434831in}}{\pgfqpoint{11.157451in}{6.438905in}}{\pgfqpoint{11.164702in}{6.446156in}}%
\pgfpathcurveto{\pgfqpoint{11.171953in}{6.453407in}}{\pgfqpoint{11.176027in}{6.463242in}}{\pgfqpoint{11.176027in}{6.473496in}}%
\pgfpathcurveto{\pgfqpoint{11.176027in}{6.483751in}}{\pgfqpoint{11.171953in}{6.493586in}}{\pgfqpoint{11.164702in}{6.500837in}}%
\pgfpathcurveto{\pgfqpoint{11.157451in}{6.508088in}}{\pgfqpoint{11.147616in}{6.512161in}}{\pgfqpoint{11.137361in}{6.512161in}}%
\pgfpathcurveto{\pgfqpoint{11.127107in}{6.512161in}}{\pgfqpoint{11.117272in}{6.508088in}}{\pgfqpoint{11.110021in}{6.500837in}}%
\pgfpathcurveto{\pgfqpoint{11.102770in}{6.493586in}}{\pgfqpoint{11.098696in}{6.483751in}}{\pgfqpoint{11.098696in}{6.473496in}}%
\pgfpathcurveto{\pgfqpoint{11.098696in}{6.463242in}}{\pgfqpoint{11.102770in}{6.453407in}}{\pgfqpoint{11.110021in}{6.446156in}}%
\pgfpathcurveto{\pgfqpoint{11.117272in}{6.438905in}}{\pgfqpoint{11.127107in}{6.434831in}}{\pgfqpoint{11.137361in}{6.434831in}}%
\pgfpathlineto{\pgfqpoint{11.137361in}{6.434831in}}%
\pgfpathclose%
\pgfusepath{stroke}%
\end{pgfscope}%
\begin{pgfscope}%
\pgfpathrectangle{\pgfqpoint{7.904091in}{2.928000in}}{\pgfqpoint{5.664000in}{5.664000in}}%
\pgfusepath{clip}%
\pgfsetbuttcap%
\pgfsetroundjoin%
\pgfsetlinewidth{1.003750pt}%
\definecolor{currentstroke}{rgb}{0.541176,0.380392,0.658824}%
\pgfsetstrokecolor{currentstroke}%
\pgfsetstrokeopacity{0.900000}%
\pgfsetdash{}{0pt}%
\pgfpathmoveto{\pgfqpoint{11.528808in}{6.229591in}}%
\pgfpathcurveto{\pgfqpoint{11.539062in}{6.229591in}}{\pgfqpoint{11.548898in}{6.233665in}}{\pgfqpoint{11.556149in}{6.240916in}}%
\pgfpathcurveto{\pgfqpoint{11.563399in}{6.248167in}}{\pgfqpoint{11.567473in}{6.258002in}}{\pgfqpoint{11.567473in}{6.268257in}}%
\pgfpathcurveto{\pgfqpoint{11.567473in}{6.278511in}}{\pgfqpoint{11.563399in}{6.288346in}}{\pgfqpoint{11.556149in}{6.295597in}}%
\pgfpathcurveto{\pgfqpoint{11.548898in}{6.302848in}}{\pgfqpoint{11.539062in}{6.306922in}}{\pgfqpoint{11.528808in}{6.306922in}}%
\pgfpathcurveto{\pgfqpoint{11.518554in}{6.306922in}}{\pgfqpoint{11.508719in}{6.302848in}}{\pgfqpoint{11.501468in}{6.295597in}}%
\pgfpathcurveto{\pgfqpoint{11.494217in}{6.288346in}}{\pgfqpoint{11.490143in}{6.278511in}}{\pgfqpoint{11.490143in}{6.268257in}}%
\pgfpathcurveto{\pgfqpoint{11.490143in}{6.258002in}}{\pgfqpoint{11.494217in}{6.248167in}}{\pgfqpoint{11.501468in}{6.240916in}}%
\pgfpathcurveto{\pgfqpoint{11.508719in}{6.233665in}}{\pgfqpoint{11.518554in}{6.229591in}}{\pgfqpoint{11.528808in}{6.229591in}}%
\pgfpathlineto{\pgfqpoint{11.528808in}{6.229591in}}%
\pgfpathclose%
\pgfusepath{stroke}%
\end{pgfscope}%
\begin{pgfscope}%
\pgfpathrectangle{\pgfqpoint{7.904091in}{2.928000in}}{\pgfqpoint{5.664000in}{5.664000in}}%
\pgfusepath{clip}%
\pgfsetbuttcap%
\pgfsetroundjoin%
\pgfsetlinewidth{1.003750pt}%
\definecolor{currentstroke}{rgb}{0.541176,0.380392,0.658824}%
\pgfsetstrokecolor{currentstroke}%
\pgfsetstrokeopacity{0.900000}%
\pgfsetdash{}{0pt}%
\pgfpathmoveto{\pgfqpoint{10.779939in}{5.248223in}}%
\pgfpathcurveto{\pgfqpoint{10.790193in}{5.248223in}}{\pgfqpoint{10.800029in}{5.252297in}}{\pgfqpoint{10.807280in}{5.259548in}}%
\pgfpathcurveto{\pgfqpoint{10.814530in}{5.266798in}}{\pgfqpoint{10.818604in}{5.276634in}}{\pgfqpoint{10.818604in}{5.286888in}}%
\pgfpathcurveto{\pgfqpoint{10.818604in}{5.297142in}}{\pgfqpoint{10.814530in}{5.306977in}}{\pgfqpoint{10.807280in}{5.314228in}}%
\pgfpathcurveto{\pgfqpoint{10.800029in}{5.321479in}}{\pgfqpoint{10.790193in}{5.325553in}}{\pgfqpoint{10.779939in}{5.325553in}}%
\pgfpathcurveto{\pgfqpoint{10.769685in}{5.325553in}}{\pgfqpoint{10.759850in}{5.321479in}}{\pgfqpoint{10.752599in}{5.314228in}}%
\pgfpathcurveto{\pgfqpoint{10.745348in}{5.306977in}}{\pgfqpoint{10.741274in}{5.297142in}}{\pgfqpoint{10.741274in}{5.286888in}}%
\pgfpathcurveto{\pgfqpoint{10.741274in}{5.276634in}}{\pgfqpoint{10.745348in}{5.266798in}}{\pgfqpoint{10.752599in}{5.259548in}}%
\pgfpathcurveto{\pgfqpoint{10.759850in}{5.252297in}}{\pgfqpoint{10.769685in}{5.248223in}}{\pgfqpoint{10.779939in}{5.248223in}}%
\pgfpathlineto{\pgfqpoint{10.779939in}{5.248223in}}%
\pgfpathclose%
\pgfusepath{stroke}%
\end{pgfscope}%
\begin{pgfscope}%
\pgfpathrectangle{\pgfqpoint{7.904091in}{2.928000in}}{\pgfqpoint{5.664000in}{5.664000in}}%
\pgfusepath{clip}%
\pgfsetbuttcap%
\pgfsetroundjoin%
\pgfsetlinewidth{1.003750pt}%
\definecolor{currentstroke}{rgb}{0.541176,0.380392,0.658824}%
\pgfsetstrokecolor{currentstroke}%
\pgfsetstrokeopacity{0.900000}%
\pgfsetdash{}{0pt}%
\pgfpathmoveto{\pgfqpoint{11.690475in}{6.402584in}}%
\pgfpathcurveto{\pgfqpoint{11.700729in}{6.402584in}}{\pgfqpoint{11.710565in}{6.406658in}}{\pgfqpoint{11.717815in}{6.413909in}}%
\pgfpathcurveto{\pgfqpoint{11.725066in}{6.421160in}}{\pgfqpoint{11.729140in}{6.430995in}}{\pgfqpoint{11.729140in}{6.441250in}}%
\pgfpathcurveto{\pgfqpoint{11.729140in}{6.451504in}}{\pgfqpoint{11.725066in}{6.461339in}}{\pgfqpoint{11.717815in}{6.468590in}}%
\pgfpathcurveto{\pgfqpoint{11.710565in}{6.475841in}}{\pgfqpoint{11.700729in}{6.479915in}}{\pgfqpoint{11.690475in}{6.479915in}}%
\pgfpathcurveto{\pgfqpoint{11.680221in}{6.479915in}}{\pgfqpoint{11.670385in}{6.475841in}}{\pgfqpoint{11.663135in}{6.468590in}}%
\pgfpathcurveto{\pgfqpoint{11.655884in}{6.461339in}}{\pgfqpoint{11.651810in}{6.451504in}}{\pgfqpoint{11.651810in}{6.441250in}}%
\pgfpathcurveto{\pgfqpoint{11.651810in}{6.430995in}}{\pgfqpoint{11.655884in}{6.421160in}}{\pgfqpoint{11.663135in}{6.413909in}}%
\pgfpathcurveto{\pgfqpoint{11.670385in}{6.406658in}}{\pgfqpoint{11.680221in}{6.402584in}}{\pgfqpoint{11.690475in}{6.402584in}}%
\pgfpathlineto{\pgfqpoint{11.690475in}{6.402584in}}%
\pgfpathclose%
\pgfusepath{stroke}%
\end{pgfscope}%
\begin{pgfscope}%
\pgfpathrectangle{\pgfqpoint{7.904091in}{2.928000in}}{\pgfqpoint{5.664000in}{5.664000in}}%
\pgfusepath{clip}%
\pgfsetbuttcap%
\pgfsetroundjoin%
\pgfsetlinewidth{1.003750pt}%
\definecolor{currentstroke}{rgb}{0.541176,0.380392,0.658824}%
\pgfsetstrokecolor{currentstroke}%
\pgfsetstrokeopacity{0.900000}%
\pgfsetdash{}{0pt}%
\pgfpathmoveto{\pgfqpoint{11.520027in}{6.449942in}}%
\pgfpathcurveto{\pgfqpoint{11.530281in}{6.449942in}}{\pgfqpoint{11.540117in}{6.454016in}}{\pgfqpoint{11.547368in}{6.461267in}}%
\pgfpathcurveto{\pgfqpoint{11.554618in}{6.468518in}}{\pgfqpoint{11.558692in}{6.478353in}}{\pgfqpoint{11.558692in}{6.488607in}}%
\pgfpathcurveto{\pgfqpoint{11.558692in}{6.498861in}}{\pgfqpoint{11.554618in}{6.508697in}}{\pgfqpoint{11.547368in}{6.515948in}}%
\pgfpathcurveto{\pgfqpoint{11.540117in}{6.523198in}}{\pgfqpoint{11.530281in}{6.527272in}}{\pgfqpoint{11.520027in}{6.527272in}}%
\pgfpathcurveto{\pgfqpoint{11.509773in}{6.527272in}}{\pgfqpoint{11.499938in}{6.523198in}}{\pgfqpoint{11.492687in}{6.515948in}}%
\pgfpathcurveto{\pgfqpoint{11.485436in}{6.508697in}}{\pgfqpoint{11.481362in}{6.498861in}}{\pgfqpoint{11.481362in}{6.488607in}}%
\pgfpathcurveto{\pgfqpoint{11.481362in}{6.478353in}}{\pgfqpoint{11.485436in}{6.468518in}}{\pgfqpoint{11.492687in}{6.461267in}}%
\pgfpathcurveto{\pgfqpoint{11.499938in}{6.454016in}}{\pgfqpoint{11.509773in}{6.449942in}}{\pgfqpoint{11.520027in}{6.449942in}}%
\pgfpathlineto{\pgfqpoint{11.520027in}{6.449942in}}%
\pgfpathclose%
\pgfusepath{stroke}%
\end{pgfscope}%
\begin{pgfscope}%
\pgfpathrectangle{\pgfqpoint{7.904091in}{2.928000in}}{\pgfqpoint{5.664000in}{5.664000in}}%
\pgfusepath{clip}%
\pgfsetbuttcap%
\pgfsetroundjoin%
\pgfsetlinewidth{1.003750pt}%
\definecolor{currentstroke}{rgb}{0.541176,0.380392,0.658824}%
\pgfsetstrokecolor{currentstroke}%
\pgfsetstrokeopacity{0.900000}%
\pgfsetdash{}{0pt}%
\pgfpathmoveto{\pgfqpoint{11.503990in}{6.030404in}}%
\pgfpathcurveto{\pgfqpoint{11.514244in}{6.030404in}}{\pgfqpoint{11.524080in}{6.034478in}}{\pgfqpoint{11.531330in}{6.041729in}}%
\pgfpathcurveto{\pgfqpoint{11.538581in}{6.048980in}}{\pgfqpoint{11.542655in}{6.058815in}}{\pgfqpoint{11.542655in}{6.069069in}}%
\pgfpathcurveto{\pgfqpoint{11.542655in}{6.079324in}}{\pgfqpoint{11.538581in}{6.089159in}}{\pgfqpoint{11.531330in}{6.096410in}}%
\pgfpathcurveto{\pgfqpoint{11.524080in}{6.103661in}}{\pgfqpoint{11.514244in}{6.107734in}}{\pgfqpoint{11.503990in}{6.107734in}}%
\pgfpathcurveto{\pgfqpoint{11.493736in}{6.107734in}}{\pgfqpoint{11.483901in}{6.103661in}}{\pgfqpoint{11.476650in}{6.096410in}}%
\pgfpathcurveto{\pgfqpoint{11.469399in}{6.089159in}}{\pgfqpoint{11.465325in}{6.079324in}}{\pgfqpoint{11.465325in}{6.069069in}}%
\pgfpathcurveto{\pgfqpoint{11.465325in}{6.058815in}}{\pgfqpoint{11.469399in}{6.048980in}}{\pgfqpoint{11.476650in}{6.041729in}}%
\pgfpathcurveto{\pgfqpoint{11.483901in}{6.034478in}}{\pgfqpoint{11.493736in}{6.030404in}}{\pgfqpoint{11.503990in}{6.030404in}}%
\pgfpathlineto{\pgfqpoint{11.503990in}{6.030404in}}%
\pgfpathclose%
\pgfusepath{stroke}%
\end{pgfscope}%
\begin{pgfscope}%
\pgfpathrectangle{\pgfqpoint{7.904091in}{2.928000in}}{\pgfqpoint{5.664000in}{5.664000in}}%
\pgfusepath{clip}%
\pgfsetbuttcap%
\pgfsetroundjoin%
\pgfsetlinewidth{1.003750pt}%
\definecolor{currentstroke}{rgb}{0.541176,0.380392,0.658824}%
\pgfsetstrokecolor{currentstroke}%
\pgfsetstrokeopacity{0.900000}%
\pgfsetdash{}{0pt}%
\pgfpathmoveto{\pgfqpoint{11.716818in}{6.296396in}}%
\pgfpathcurveto{\pgfqpoint{11.727072in}{6.296396in}}{\pgfqpoint{11.736907in}{6.300470in}}{\pgfqpoint{11.744158in}{6.307720in}}%
\pgfpathcurveto{\pgfqpoint{11.751409in}{6.314971in}}{\pgfqpoint{11.755483in}{6.324807in}}{\pgfqpoint{11.755483in}{6.335061in}}%
\pgfpathcurveto{\pgfqpoint{11.755483in}{6.345315in}}{\pgfqpoint{11.751409in}{6.355150in}}{\pgfqpoint{11.744158in}{6.362401in}}%
\pgfpathcurveto{\pgfqpoint{11.736907in}{6.369652in}}{\pgfqpoint{11.727072in}{6.373726in}}{\pgfqpoint{11.716818in}{6.373726in}}%
\pgfpathcurveto{\pgfqpoint{11.706564in}{6.373726in}}{\pgfqpoint{11.696728in}{6.369652in}}{\pgfqpoint{11.689477in}{6.362401in}}%
\pgfpathcurveto{\pgfqpoint{11.682227in}{6.355150in}}{\pgfqpoint{11.678153in}{6.345315in}}{\pgfqpoint{11.678153in}{6.335061in}}%
\pgfpathcurveto{\pgfqpoint{11.678153in}{6.324807in}}{\pgfqpoint{11.682227in}{6.314971in}}{\pgfqpoint{11.689477in}{6.307720in}}%
\pgfpathcurveto{\pgfqpoint{11.696728in}{6.300470in}}{\pgfqpoint{11.706564in}{6.296396in}}{\pgfqpoint{11.716818in}{6.296396in}}%
\pgfpathlineto{\pgfqpoint{11.716818in}{6.296396in}}%
\pgfpathclose%
\pgfusepath{stroke}%
\end{pgfscope}%
\begin{pgfscope}%
\pgfpathrectangle{\pgfqpoint{7.904091in}{2.928000in}}{\pgfqpoint{5.664000in}{5.664000in}}%
\pgfusepath{clip}%
\pgfsetbuttcap%
\pgfsetroundjoin%
\pgfsetlinewidth{1.003750pt}%
\definecolor{currentstroke}{rgb}{0.541176,0.380392,0.658824}%
\pgfsetstrokecolor{currentstroke}%
\pgfsetstrokeopacity{0.900000}%
\pgfsetdash{}{0pt}%
\pgfpathmoveto{\pgfqpoint{13.014283in}{8.382764in}}%
\pgfpathcurveto{\pgfqpoint{13.024537in}{8.382764in}}{\pgfqpoint{13.034372in}{8.386838in}}{\pgfqpoint{13.041623in}{8.394089in}}%
\pgfpathcurveto{\pgfqpoint{13.048874in}{8.401340in}}{\pgfqpoint{13.052948in}{8.411175in}}{\pgfqpoint{13.052948in}{8.421429in}}%
\pgfpathcurveto{\pgfqpoint{13.052948in}{8.431684in}}{\pgfqpoint{13.048874in}{8.441519in}}{\pgfqpoint{13.041623in}{8.448770in}}%
\pgfpathcurveto{\pgfqpoint{13.034372in}{8.456020in}}{\pgfqpoint{13.024537in}{8.460094in}}{\pgfqpoint{13.014283in}{8.460094in}}%
\pgfpathcurveto{\pgfqpoint{13.004029in}{8.460094in}}{\pgfqpoint{12.994193in}{8.456020in}}{\pgfqpoint{12.986942in}{8.448770in}}%
\pgfpathcurveto{\pgfqpoint{12.979692in}{8.441519in}}{\pgfqpoint{12.975618in}{8.431684in}}{\pgfqpoint{12.975618in}{8.421429in}}%
\pgfpathcurveto{\pgfqpoint{12.975618in}{8.411175in}}{\pgfqpoint{12.979692in}{8.401340in}}{\pgfqpoint{12.986942in}{8.394089in}}%
\pgfpathcurveto{\pgfqpoint{12.994193in}{8.386838in}}{\pgfqpoint{13.004029in}{8.382764in}}{\pgfqpoint{13.014283in}{8.382764in}}%
\pgfpathlineto{\pgfqpoint{13.014283in}{8.382764in}}%
\pgfpathclose%
\pgfusepath{stroke}%
\end{pgfscope}%
\begin{pgfscope}%
\pgfpathrectangle{\pgfqpoint{7.904091in}{2.928000in}}{\pgfqpoint{5.664000in}{5.664000in}}%
\pgfusepath{clip}%
\pgfsetbuttcap%
\pgfsetroundjoin%
\pgfsetlinewidth{1.003750pt}%
\definecolor{currentstroke}{rgb}{0.541176,0.380392,0.658824}%
\pgfsetstrokecolor{currentstroke}%
\pgfsetstrokeopacity{0.900000}%
\pgfsetdash{}{0pt}%
\pgfpathmoveto{\pgfqpoint{10.960260in}{5.873860in}}%
\pgfpathcurveto{\pgfqpoint{10.970515in}{5.873860in}}{\pgfqpoint{10.980350in}{5.877934in}}{\pgfqpoint{10.987601in}{5.885185in}}%
\pgfpathcurveto{\pgfqpoint{10.994851in}{5.892436in}}{\pgfqpoint{10.998925in}{5.902271in}}{\pgfqpoint{10.998925in}{5.912525in}}%
\pgfpathcurveto{\pgfqpoint{10.998925in}{5.922779in}}{\pgfqpoint{10.994851in}{5.932615in}}{\pgfqpoint{10.987601in}{5.939865in}}%
\pgfpathcurveto{\pgfqpoint{10.980350in}{5.947116in}}{\pgfqpoint{10.970515in}{5.951190in}}{\pgfqpoint{10.960260in}{5.951190in}}%
\pgfpathcurveto{\pgfqpoint{10.950006in}{5.951190in}}{\pgfqpoint{10.940171in}{5.947116in}}{\pgfqpoint{10.932920in}{5.939865in}}%
\pgfpathcurveto{\pgfqpoint{10.925669in}{5.932615in}}{\pgfqpoint{10.921595in}{5.922779in}}{\pgfqpoint{10.921595in}{5.912525in}}%
\pgfpathcurveto{\pgfqpoint{10.921595in}{5.902271in}}{\pgfqpoint{10.925669in}{5.892436in}}{\pgfqpoint{10.932920in}{5.885185in}}%
\pgfpathcurveto{\pgfqpoint{10.940171in}{5.877934in}}{\pgfqpoint{10.950006in}{5.873860in}}{\pgfqpoint{10.960260in}{5.873860in}}%
\pgfpathlineto{\pgfqpoint{10.960260in}{5.873860in}}%
\pgfpathclose%
\pgfusepath{stroke}%
\end{pgfscope}%
\begin{pgfscope}%
\pgfpathrectangle{\pgfqpoint{7.904091in}{2.928000in}}{\pgfqpoint{5.664000in}{5.664000in}}%
\pgfusepath{clip}%
\pgfsetbuttcap%
\pgfsetroundjoin%
\pgfsetlinewidth{1.003750pt}%
\definecolor{currentstroke}{rgb}{0.541176,0.380392,0.658824}%
\pgfsetstrokecolor{currentstroke}%
\pgfsetstrokeopacity{0.900000}%
\pgfsetdash{}{0pt}%
\pgfpathmoveto{\pgfqpoint{11.365728in}{7.005316in}}%
\pgfpathcurveto{\pgfqpoint{11.375982in}{7.005316in}}{\pgfqpoint{11.385817in}{7.009390in}}{\pgfqpoint{11.393068in}{7.016641in}}%
\pgfpathcurveto{\pgfqpoint{11.400319in}{7.023892in}}{\pgfqpoint{11.404393in}{7.033727in}}{\pgfqpoint{11.404393in}{7.043981in}}%
\pgfpathcurveto{\pgfqpoint{11.404393in}{7.054235in}}{\pgfqpoint{11.400319in}{7.064071in}}{\pgfqpoint{11.393068in}{7.071322in}}%
\pgfpathcurveto{\pgfqpoint{11.385817in}{7.078572in}}{\pgfqpoint{11.375982in}{7.082646in}}{\pgfqpoint{11.365728in}{7.082646in}}%
\pgfpathcurveto{\pgfqpoint{11.355474in}{7.082646in}}{\pgfqpoint{11.345638in}{7.078572in}}{\pgfqpoint{11.338387in}{7.071322in}}%
\pgfpathcurveto{\pgfqpoint{11.331137in}{7.064071in}}{\pgfqpoint{11.327063in}{7.054235in}}{\pgfqpoint{11.327063in}{7.043981in}}%
\pgfpathcurveto{\pgfqpoint{11.327063in}{7.033727in}}{\pgfqpoint{11.331137in}{7.023892in}}{\pgfqpoint{11.338387in}{7.016641in}}%
\pgfpathcurveto{\pgfqpoint{11.345638in}{7.009390in}}{\pgfqpoint{11.355474in}{7.005316in}}{\pgfqpoint{11.365728in}{7.005316in}}%
\pgfpathlineto{\pgfqpoint{11.365728in}{7.005316in}}%
\pgfpathclose%
\pgfusepath{stroke}%
\end{pgfscope}%
\begin{pgfscope}%
\pgfpathrectangle{\pgfqpoint{7.904091in}{2.928000in}}{\pgfqpoint{5.664000in}{5.664000in}}%
\pgfusepath{clip}%
\pgfsetbuttcap%
\pgfsetroundjoin%
\pgfsetlinewidth{1.003750pt}%
\definecolor{currentstroke}{rgb}{0.541176,0.380392,0.658824}%
\pgfsetstrokecolor{currentstroke}%
\pgfsetstrokeopacity{0.900000}%
\pgfsetdash{}{0pt}%
\pgfpathmoveto{\pgfqpoint{11.484268in}{6.424236in}}%
\pgfpathcurveto{\pgfqpoint{11.494522in}{6.424236in}}{\pgfqpoint{11.504358in}{6.428310in}}{\pgfqpoint{11.511608in}{6.435561in}}%
\pgfpathcurveto{\pgfqpoint{11.518859in}{6.442811in}}{\pgfqpoint{11.522933in}{6.452647in}}{\pgfqpoint{11.522933in}{6.462901in}}%
\pgfpathcurveto{\pgfqpoint{11.522933in}{6.473155in}}{\pgfqpoint{11.518859in}{6.482991in}}{\pgfqpoint{11.511608in}{6.490241in}}%
\pgfpathcurveto{\pgfqpoint{11.504358in}{6.497492in}}{\pgfqpoint{11.494522in}{6.501566in}}{\pgfqpoint{11.484268in}{6.501566in}}%
\pgfpathcurveto{\pgfqpoint{11.474014in}{6.501566in}}{\pgfqpoint{11.464179in}{6.497492in}}{\pgfqpoint{11.456928in}{6.490241in}}%
\pgfpathcurveto{\pgfqpoint{11.449677in}{6.482991in}}{\pgfqpoint{11.445603in}{6.473155in}}{\pgfqpoint{11.445603in}{6.462901in}}%
\pgfpathcurveto{\pgfqpoint{11.445603in}{6.452647in}}{\pgfqpoint{11.449677in}{6.442811in}}{\pgfqpoint{11.456928in}{6.435561in}}%
\pgfpathcurveto{\pgfqpoint{11.464179in}{6.428310in}}{\pgfqpoint{11.474014in}{6.424236in}}{\pgfqpoint{11.484268in}{6.424236in}}%
\pgfpathlineto{\pgfqpoint{11.484268in}{6.424236in}}%
\pgfpathclose%
\pgfusepath{stroke}%
\end{pgfscope}%
\begin{pgfscope}%
\pgfpathrectangle{\pgfqpoint{7.904091in}{2.928000in}}{\pgfqpoint{5.664000in}{5.664000in}}%
\pgfusepath{clip}%
\pgfsetbuttcap%
\pgfsetroundjoin%
\pgfsetlinewidth{1.003750pt}%
\definecolor{currentstroke}{rgb}{0.541176,0.380392,0.658824}%
\pgfsetstrokecolor{currentstroke}%
\pgfsetstrokeopacity{0.900000}%
\pgfsetdash{}{0pt}%
\pgfpathmoveto{\pgfqpoint{12.078542in}{7.046945in}}%
\pgfpathcurveto{\pgfqpoint{12.088796in}{7.046945in}}{\pgfqpoint{12.098631in}{7.051019in}}{\pgfqpoint{12.105882in}{7.058270in}}%
\pgfpathcurveto{\pgfqpoint{12.113133in}{7.065520in}}{\pgfqpoint{12.117207in}{7.075356in}}{\pgfqpoint{12.117207in}{7.085610in}}%
\pgfpathcurveto{\pgfqpoint{12.117207in}{7.095864in}}{\pgfqpoint{12.113133in}{7.105700in}}{\pgfqpoint{12.105882in}{7.112950in}}%
\pgfpathcurveto{\pgfqpoint{12.098631in}{7.120201in}}{\pgfqpoint{12.088796in}{7.124275in}}{\pgfqpoint{12.078542in}{7.124275in}}%
\pgfpathcurveto{\pgfqpoint{12.068288in}{7.124275in}}{\pgfqpoint{12.058452in}{7.120201in}}{\pgfqpoint{12.051201in}{7.112950in}}%
\pgfpathcurveto{\pgfqpoint{12.043951in}{7.105700in}}{\pgfqpoint{12.039877in}{7.095864in}}{\pgfqpoint{12.039877in}{7.085610in}}%
\pgfpathcurveto{\pgfqpoint{12.039877in}{7.075356in}}{\pgfqpoint{12.043951in}{7.065520in}}{\pgfqpoint{12.051201in}{7.058270in}}%
\pgfpathcurveto{\pgfqpoint{12.058452in}{7.051019in}}{\pgfqpoint{12.068288in}{7.046945in}}{\pgfqpoint{12.078542in}{7.046945in}}%
\pgfpathlineto{\pgfqpoint{12.078542in}{7.046945in}}%
\pgfpathclose%
\pgfusepath{stroke}%
\end{pgfscope}%
\begin{pgfscope}%
\pgfpathrectangle{\pgfqpoint{7.904091in}{2.928000in}}{\pgfqpoint{5.664000in}{5.664000in}}%
\pgfusepath{clip}%
\pgfsetbuttcap%
\pgfsetroundjoin%
\pgfsetlinewidth{1.003750pt}%
\definecolor{currentstroke}{rgb}{0.541176,0.380392,0.658824}%
\pgfsetstrokecolor{currentstroke}%
\pgfsetstrokeopacity{0.900000}%
\pgfsetdash{}{0pt}%
\pgfpathmoveto{\pgfqpoint{12.269833in}{7.474511in}}%
\pgfpathcurveto{\pgfqpoint{12.280087in}{7.474511in}}{\pgfqpoint{12.289923in}{7.478585in}}{\pgfqpoint{12.297173in}{7.485836in}}%
\pgfpathcurveto{\pgfqpoint{12.304424in}{7.493087in}}{\pgfqpoint{12.308498in}{7.502922in}}{\pgfqpoint{12.308498in}{7.513176in}}%
\pgfpathcurveto{\pgfqpoint{12.308498in}{7.523430in}}{\pgfqpoint{12.304424in}{7.533266in}}{\pgfqpoint{12.297173in}{7.540517in}}%
\pgfpathcurveto{\pgfqpoint{12.289923in}{7.547767in}}{\pgfqpoint{12.280087in}{7.551841in}}{\pgfqpoint{12.269833in}{7.551841in}}%
\pgfpathcurveto{\pgfqpoint{12.259579in}{7.551841in}}{\pgfqpoint{12.249744in}{7.547767in}}{\pgfqpoint{12.242493in}{7.540517in}}%
\pgfpathcurveto{\pgfqpoint{12.235242in}{7.533266in}}{\pgfqpoint{12.231168in}{7.523430in}}{\pgfqpoint{12.231168in}{7.513176in}}%
\pgfpathcurveto{\pgfqpoint{12.231168in}{7.502922in}}{\pgfqpoint{12.235242in}{7.493087in}}{\pgfqpoint{12.242493in}{7.485836in}}%
\pgfpathcurveto{\pgfqpoint{12.249744in}{7.478585in}}{\pgfqpoint{12.259579in}{7.474511in}}{\pgfqpoint{12.269833in}{7.474511in}}%
\pgfpathlineto{\pgfqpoint{12.269833in}{7.474511in}}%
\pgfpathclose%
\pgfusepath{stroke}%
\end{pgfscope}%
\begin{pgfscope}%
\pgfpathrectangle{\pgfqpoint{7.904091in}{2.928000in}}{\pgfqpoint{5.664000in}{5.664000in}}%
\pgfusepath{clip}%
\pgfsetbuttcap%
\pgfsetroundjoin%
\pgfsetlinewidth{1.003750pt}%
\definecolor{currentstroke}{rgb}{0.541176,0.380392,0.658824}%
\pgfsetstrokecolor{currentstroke}%
\pgfsetstrokeopacity{0.900000}%
\pgfsetdash{}{0pt}%
\pgfpathmoveto{\pgfqpoint{11.620997in}{6.857049in}}%
\pgfpathcurveto{\pgfqpoint{11.631251in}{6.857049in}}{\pgfqpoint{11.641086in}{6.861123in}}{\pgfqpoint{11.648337in}{6.868374in}}%
\pgfpathcurveto{\pgfqpoint{11.655588in}{6.875625in}}{\pgfqpoint{11.659662in}{6.885460in}}{\pgfqpoint{11.659662in}{6.895714in}}%
\pgfpathcurveto{\pgfqpoint{11.659662in}{6.905968in}}{\pgfqpoint{11.655588in}{6.915804in}}{\pgfqpoint{11.648337in}{6.923054in}}%
\pgfpathcurveto{\pgfqpoint{11.641086in}{6.930305in}}{\pgfqpoint{11.631251in}{6.934379in}}{\pgfqpoint{11.620997in}{6.934379in}}%
\pgfpathcurveto{\pgfqpoint{11.610743in}{6.934379in}}{\pgfqpoint{11.600907in}{6.930305in}}{\pgfqpoint{11.593657in}{6.923054in}}%
\pgfpathcurveto{\pgfqpoint{11.586406in}{6.915804in}}{\pgfqpoint{11.582332in}{6.905968in}}{\pgfqpoint{11.582332in}{6.895714in}}%
\pgfpathcurveto{\pgfqpoint{11.582332in}{6.885460in}}{\pgfqpoint{11.586406in}{6.875625in}}{\pgfqpoint{11.593657in}{6.868374in}}%
\pgfpathcurveto{\pgfqpoint{11.600907in}{6.861123in}}{\pgfqpoint{11.610743in}{6.857049in}}{\pgfqpoint{11.620997in}{6.857049in}}%
\pgfpathlineto{\pgfqpoint{11.620997in}{6.857049in}}%
\pgfpathclose%
\pgfusepath{stroke}%
\end{pgfscope}%
\begin{pgfscope}%
\pgfpathrectangle{\pgfqpoint{7.904091in}{2.928000in}}{\pgfqpoint{5.664000in}{5.664000in}}%
\pgfusepath{clip}%
\pgfsetbuttcap%
\pgfsetroundjoin%
\pgfsetlinewidth{1.003750pt}%
\definecolor{currentstroke}{rgb}{0.541176,0.380392,0.658824}%
\pgfsetstrokecolor{currentstroke}%
\pgfsetstrokeopacity{0.900000}%
\pgfsetdash{}{0pt}%
\pgfpathmoveto{\pgfqpoint{11.527786in}{6.938215in}}%
\pgfpathcurveto{\pgfqpoint{11.538040in}{6.938215in}}{\pgfqpoint{11.547876in}{6.942289in}}{\pgfqpoint{11.555127in}{6.949540in}}%
\pgfpathcurveto{\pgfqpoint{11.562377in}{6.956791in}}{\pgfqpoint{11.566451in}{6.966626in}}{\pgfqpoint{11.566451in}{6.976880in}}%
\pgfpathcurveto{\pgfqpoint{11.566451in}{6.987134in}}{\pgfqpoint{11.562377in}{6.996970in}}{\pgfqpoint{11.555127in}{7.004220in}}%
\pgfpathcurveto{\pgfqpoint{11.547876in}{7.011471in}}{\pgfqpoint{11.538040in}{7.015545in}}{\pgfqpoint{11.527786in}{7.015545in}}%
\pgfpathcurveto{\pgfqpoint{11.517532in}{7.015545in}}{\pgfqpoint{11.507697in}{7.011471in}}{\pgfqpoint{11.500446in}{7.004220in}}%
\pgfpathcurveto{\pgfqpoint{11.493195in}{6.996970in}}{\pgfqpoint{11.489121in}{6.987134in}}{\pgfqpoint{11.489121in}{6.976880in}}%
\pgfpathcurveto{\pgfqpoint{11.489121in}{6.966626in}}{\pgfqpoint{11.493195in}{6.956791in}}{\pgfqpoint{11.500446in}{6.949540in}}%
\pgfpathcurveto{\pgfqpoint{11.507697in}{6.942289in}}{\pgfqpoint{11.517532in}{6.938215in}}{\pgfqpoint{11.527786in}{6.938215in}}%
\pgfpathlineto{\pgfqpoint{11.527786in}{6.938215in}}%
\pgfpathclose%
\pgfusepath{stroke}%
\end{pgfscope}%
\begin{pgfscope}%
\pgfpathrectangle{\pgfqpoint{7.904091in}{2.928000in}}{\pgfqpoint{5.664000in}{5.664000in}}%
\pgfusepath{clip}%
\pgfsetbuttcap%
\pgfsetroundjoin%
\pgfsetlinewidth{1.003750pt}%
\definecolor{currentstroke}{rgb}{0.541176,0.380392,0.658824}%
\pgfsetstrokecolor{currentstroke}%
\pgfsetstrokeopacity{0.900000}%
\pgfsetdash{}{0pt}%
\pgfpathmoveto{\pgfqpoint{10.717690in}{5.375895in}}%
\pgfpathcurveto{\pgfqpoint{10.727944in}{5.375895in}}{\pgfqpoint{10.737779in}{5.379969in}}{\pgfqpoint{10.745030in}{5.387219in}}%
\pgfpathcurveto{\pgfqpoint{10.752281in}{5.394470in}}{\pgfqpoint{10.756355in}{5.404306in}}{\pgfqpoint{10.756355in}{5.414560in}}%
\pgfpathcurveto{\pgfqpoint{10.756355in}{5.424814in}}{\pgfqpoint{10.752281in}{5.434649in}}{\pgfqpoint{10.745030in}{5.441900in}}%
\pgfpathcurveto{\pgfqpoint{10.737779in}{5.449151in}}{\pgfqpoint{10.727944in}{5.453225in}}{\pgfqpoint{10.717690in}{5.453225in}}%
\pgfpathcurveto{\pgfqpoint{10.707436in}{5.453225in}}{\pgfqpoint{10.697600in}{5.449151in}}{\pgfqpoint{10.690350in}{5.441900in}}%
\pgfpathcurveto{\pgfqpoint{10.683099in}{5.434649in}}{\pgfqpoint{10.679025in}{5.424814in}}{\pgfqpoint{10.679025in}{5.414560in}}%
\pgfpathcurveto{\pgfqpoint{10.679025in}{5.404306in}}{\pgfqpoint{10.683099in}{5.394470in}}{\pgfqpoint{10.690350in}{5.387219in}}%
\pgfpathcurveto{\pgfqpoint{10.697600in}{5.379969in}}{\pgfqpoint{10.707436in}{5.375895in}}{\pgfqpoint{10.717690in}{5.375895in}}%
\pgfpathlineto{\pgfqpoint{10.717690in}{5.375895in}}%
\pgfpathclose%
\pgfusepath{stroke}%
\end{pgfscope}%
\begin{pgfscope}%
\pgfpathrectangle{\pgfqpoint{7.904091in}{2.928000in}}{\pgfqpoint{5.664000in}{5.664000in}}%
\pgfusepath{clip}%
\pgfsetbuttcap%
\pgfsetroundjoin%
\pgfsetlinewidth{1.003750pt}%
\definecolor{currentstroke}{rgb}{0.541176,0.380392,0.658824}%
\pgfsetstrokecolor{currentstroke}%
\pgfsetstrokeopacity{0.900000}%
\pgfsetdash{}{0pt}%
\pgfpathmoveto{\pgfqpoint{11.531738in}{6.419357in}}%
\pgfpathcurveto{\pgfqpoint{11.541992in}{6.419357in}}{\pgfqpoint{11.551828in}{6.423431in}}{\pgfqpoint{11.559078in}{6.430682in}}%
\pgfpathcurveto{\pgfqpoint{11.566329in}{6.437932in}}{\pgfqpoint{11.570403in}{6.447768in}}{\pgfqpoint{11.570403in}{6.458022in}}%
\pgfpathcurveto{\pgfqpoint{11.570403in}{6.468276in}}{\pgfqpoint{11.566329in}{6.478111in}}{\pgfqpoint{11.559078in}{6.485362in}}%
\pgfpathcurveto{\pgfqpoint{11.551828in}{6.492613in}}{\pgfqpoint{11.541992in}{6.496687in}}{\pgfqpoint{11.531738in}{6.496687in}}%
\pgfpathcurveto{\pgfqpoint{11.521484in}{6.496687in}}{\pgfqpoint{11.511649in}{6.492613in}}{\pgfqpoint{11.504398in}{6.485362in}}%
\pgfpathcurveto{\pgfqpoint{11.497147in}{6.478111in}}{\pgfqpoint{11.493073in}{6.468276in}}{\pgfqpoint{11.493073in}{6.458022in}}%
\pgfpathcurveto{\pgfqpoint{11.493073in}{6.447768in}}{\pgfqpoint{11.497147in}{6.437932in}}{\pgfqpoint{11.504398in}{6.430682in}}%
\pgfpathcurveto{\pgfqpoint{11.511649in}{6.423431in}}{\pgfqpoint{11.521484in}{6.419357in}}{\pgfqpoint{11.531738in}{6.419357in}}%
\pgfpathlineto{\pgfqpoint{11.531738in}{6.419357in}}%
\pgfpathclose%
\pgfusepath{stroke}%
\end{pgfscope}%
\begin{pgfscope}%
\pgfpathrectangle{\pgfqpoint{7.904091in}{2.928000in}}{\pgfqpoint{5.664000in}{5.664000in}}%
\pgfusepath{clip}%
\pgfsetbuttcap%
\pgfsetroundjoin%
\pgfsetlinewidth{1.003750pt}%
\definecolor{currentstroke}{rgb}{0.541176,0.380392,0.658824}%
\pgfsetstrokecolor{currentstroke}%
\pgfsetstrokeopacity{0.900000}%
\pgfsetdash{}{0pt}%
\pgfpathmoveto{\pgfqpoint{10.475465in}{6.206588in}}%
\pgfpathcurveto{\pgfqpoint{10.485719in}{6.206588in}}{\pgfqpoint{10.495554in}{6.210662in}}{\pgfqpoint{10.502805in}{6.217913in}}%
\pgfpathcurveto{\pgfqpoint{10.510056in}{6.225163in}}{\pgfqpoint{10.514130in}{6.234999in}}{\pgfqpoint{10.514130in}{6.245253in}}%
\pgfpathcurveto{\pgfqpoint{10.514130in}{6.255507in}}{\pgfqpoint{10.510056in}{6.265343in}}{\pgfqpoint{10.502805in}{6.272593in}}%
\pgfpathcurveto{\pgfqpoint{10.495554in}{6.279844in}}{\pgfqpoint{10.485719in}{6.283918in}}{\pgfqpoint{10.475465in}{6.283918in}}%
\pgfpathcurveto{\pgfqpoint{10.465210in}{6.283918in}}{\pgfqpoint{10.455375in}{6.279844in}}{\pgfqpoint{10.448124in}{6.272593in}}%
\pgfpathcurveto{\pgfqpoint{10.440873in}{6.265343in}}{\pgfqpoint{10.436799in}{6.255507in}}{\pgfqpoint{10.436799in}{6.245253in}}%
\pgfpathcurveto{\pgfqpoint{10.436799in}{6.234999in}}{\pgfqpoint{10.440873in}{6.225163in}}{\pgfqpoint{10.448124in}{6.217913in}}%
\pgfpathcurveto{\pgfqpoint{10.455375in}{6.210662in}}{\pgfqpoint{10.465210in}{6.206588in}}{\pgfqpoint{10.475465in}{6.206588in}}%
\pgfpathlineto{\pgfqpoint{10.475465in}{6.206588in}}%
\pgfpathclose%
\pgfusepath{stroke}%
\end{pgfscope}%
\begin{pgfscope}%
\pgfpathrectangle{\pgfqpoint{7.904091in}{2.928000in}}{\pgfqpoint{5.664000in}{5.664000in}}%
\pgfusepath{clip}%
\pgfsetbuttcap%
\pgfsetroundjoin%
\pgfsetlinewidth{1.003750pt}%
\definecolor{currentstroke}{rgb}{0.541176,0.380392,0.658824}%
\pgfsetstrokecolor{currentstroke}%
\pgfsetstrokeopacity{0.900000}%
\pgfsetdash{}{0pt}%
\pgfpathmoveto{\pgfqpoint{11.616496in}{6.986493in}}%
\pgfpathcurveto{\pgfqpoint{11.626750in}{6.986493in}}{\pgfqpoint{11.636586in}{6.990567in}}{\pgfqpoint{11.643837in}{6.997818in}}%
\pgfpathcurveto{\pgfqpoint{11.651087in}{7.005068in}}{\pgfqpoint{11.655161in}{7.014904in}}{\pgfqpoint{11.655161in}{7.025158in}}%
\pgfpathcurveto{\pgfqpoint{11.655161in}{7.035412in}}{\pgfqpoint{11.651087in}{7.045248in}}{\pgfqpoint{11.643837in}{7.052498in}}%
\pgfpathcurveto{\pgfqpoint{11.636586in}{7.059749in}}{\pgfqpoint{11.626750in}{7.063823in}}{\pgfqpoint{11.616496in}{7.063823in}}%
\pgfpathcurveto{\pgfqpoint{11.606242in}{7.063823in}}{\pgfqpoint{11.596407in}{7.059749in}}{\pgfqpoint{11.589156in}{7.052498in}}%
\pgfpathcurveto{\pgfqpoint{11.581905in}{7.045248in}}{\pgfqpoint{11.577831in}{7.035412in}}{\pgfqpoint{11.577831in}{7.025158in}}%
\pgfpathcurveto{\pgfqpoint{11.577831in}{7.014904in}}{\pgfqpoint{11.581905in}{7.005068in}}{\pgfqpoint{11.589156in}{6.997818in}}%
\pgfpathcurveto{\pgfqpoint{11.596407in}{6.990567in}}{\pgfqpoint{11.606242in}{6.986493in}}{\pgfqpoint{11.616496in}{6.986493in}}%
\pgfpathlineto{\pgfqpoint{11.616496in}{6.986493in}}%
\pgfpathclose%
\pgfusepath{stroke}%
\end{pgfscope}%
\begin{pgfscope}%
\pgfpathrectangle{\pgfqpoint{7.904091in}{2.928000in}}{\pgfqpoint{5.664000in}{5.664000in}}%
\pgfusepath{clip}%
\pgfsetbuttcap%
\pgfsetroundjoin%
\pgfsetlinewidth{1.003750pt}%
\definecolor{currentstroke}{rgb}{0.541176,0.380392,0.658824}%
\pgfsetstrokecolor{currentstroke}%
\pgfsetstrokeopacity{0.900000}%
\pgfsetdash{}{0pt}%
\pgfpathmoveto{\pgfqpoint{11.311429in}{6.662568in}}%
\pgfpathcurveto{\pgfqpoint{11.321683in}{6.662568in}}{\pgfqpoint{11.331519in}{6.666642in}}{\pgfqpoint{11.338770in}{6.673893in}}%
\pgfpathcurveto{\pgfqpoint{11.346020in}{6.681143in}}{\pgfqpoint{11.350094in}{6.690979in}}{\pgfqpoint{11.350094in}{6.701233in}}%
\pgfpathcurveto{\pgfqpoint{11.350094in}{6.711487in}}{\pgfqpoint{11.346020in}{6.721322in}}{\pgfqpoint{11.338770in}{6.728573in}}%
\pgfpathcurveto{\pgfqpoint{11.331519in}{6.735824in}}{\pgfqpoint{11.321683in}{6.739898in}}{\pgfqpoint{11.311429in}{6.739898in}}%
\pgfpathcurveto{\pgfqpoint{11.301175in}{6.739898in}}{\pgfqpoint{11.291340in}{6.735824in}}{\pgfqpoint{11.284089in}{6.728573in}}%
\pgfpathcurveto{\pgfqpoint{11.276838in}{6.721322in}}{\pgfqpoint{11.272764in}{6.711487in}}{\pgfqpoint{11.272764in}{6.701233in}}%
\pgfpathcurveto{\pgfqpoint{11.272764in}{6.690979in}}{\pgfqpoint{11.276838in}{6.681143in}}{\pgfqpoint{11.284089in}{6.673893in}}%
\pgfpathcurveto{\pgfqpoint{11.291340in}{6.666642in}}{\pgfqpoint{11.301175in}{6.662568in}}{\pgfqpoint{11.311429in}{6.662568in}}%
\pgfpathlineto{\pgfqpoint{11.311429in}{6.662568in}}%
\pgfpathclose%
\pgfusepath{stroke}%
\end{pgfscope}%
\begin{pgfscope}%
\pgfpathrectangle{\pgfqpoint{7.904091in}{2.928000in}}{\pgfqpoint{5.664000in}{5.664000in}}%
\pgfusepath{clip}%
\pgfsetbuttcap%
\pgfsetroundjoin%
\pgfsetlinewidth{1.003750pt}%
\definecolor{currentstroke}{rgb}{0.541176,0.380392,0.658824}%
\pgfsetstrokecolor{currentstroke}%
\pgfsetstrokeopacity{0.900000}%
\pgfsetdash{}{0pt}%
\pgfpathmoveto{\pgfqpoint{11.246799in}{5.896718in}}%
\pgfpathcurveto{\pgfqpoint{11.257053in}{5.896718in}}{\pgfqpoint{11.266888in}{5.900792in}}{\pgfqpoint{11.274139in}{5.908042in}}%
\pgfpathcurveto{\pgfqpoint{11.281390in}{5.915293in}}{\pgfqpoint{11.285464in}{5.925129in}}{\pgfqpoint{11.285464in}{5.935383in}}%
\pgfpathcurveto{\pgfqpoint{11.285464in}{5.945637in}}{\pgfqpoint{11.281390in}{5.955472in}}{\pgfqpoint{11.274139in}{5.962723in}}%
\pgfpathcurveto{\pgfqpoint{11.266888in}{5.969974in}}{\pgfqpoint{11.257053in}{5.974048in}}{\pgfqpoint{11.246799in}{5.974048in}}%
\pgfpathcurveto{\pgfqpoint{11.236545in}{5.974048in}}{\pgfqpoint{11.226709in}{5.969974in}}{\pgfqpoint{11.219458in}{5.962723in}}%
\pgfpathcurveto{\pgfqpoint{11.212208in}{5.955472in}}{\pgfqpoint{11.208134in}{5.945637in}}{\pgfqpoint{11.208134in}{5.935383in}}%
\pgfpathcurveto{\pgfqpoint{11.208134in}{5.925129in}}{\pgfqpoint{11.212208in}{5.915293in}}{\pgfqpoint{11.219458in}{5.908042in}}%
\pgfpathcurveto{\pgfqpoint{11.226709in}{5.900792in}}{\pgfqpoint{11.236545in}{5.896718in}}{\pgfqpoint{11.246799in}{5.896718in}}%
\pgfpathlineto{\pgfqpoint{11.246799in}{5.896718in}}%
\pgfpathclose%
\pgfusepath{stroke}%
\end{pgfscope}%
\begin{pgfscope}%
\pgfpathrectangle{\pgfqpoint{7.904091in}{2.928000in}}{\pgfqpoint{5.664000in}{5.664000in}}%
\pgfusepath{clip}%
\pgfsetbuttcap%
\pgfsetroundjoin%
\pgfsetlinewidth{1.003750pt}%
\definecolor{currentstroke}{rgb}{0.541176,0.380392,0.658824}%
\pgfsetstrokecolor{currentstroke}%
\pgfsetstrokeopacity{0.900000}%
\pgfsetdash{}{0pt}%
\pgfpathmoveto{\pgfqpoint{9.987790in}{4.131346in}}%
\pgfpathcurveto{\pgfqpoint{9.998044in}{4.131346in}}{\pgfqpoint{10.007880in}{4.135420in}}{\pgfqpoint{10.015130in}{4.142671in}}%
\pgfpathcurveto{\pgfqpoint{10.022381in}{4.149922in}}{\pgfqpoint{10.026455in}{4.159757in}}{\pgfqpoint{10.026455in}{4.170011in}}%
\pgfpathcurveto{\pgfqpoint{10.026455in}{4.180265in}}{\pgfqpoint{10.022381in}{4.190101in}}{\pgfqpoint{10.015130in}{4.197352in}}%
\pgfpathcurveto{\pgfqpoint{10.007880in}{4.204602in}}{\pgfqpoint{9.998044in}{4.208676in}}{\pgfqpoint{9.987790in}{4.208676in}}%
\pgfpathcurveto{\pgfqpoint{9.977536in}{4.208676in}}{\pgfqpoint{9.967700in}{4.204602in}}{\pgfqpoint{9.960450in}{4.197352in}}%
\pgfpathcurveto{\pgfqpoint{9.953199in}{4.190101in}}{\pgfqpoint{9.949125in}{4.180265in}}{\pgfqpoint{9.949125in}{4.170011in}}%
\pgfpathcurveto{\pgfqpoint{9.949125in}{4.159757in}}{\pgfqpoint{9.953199in}{4.149922in}}{\pgfqpoint{9.960450in}{4.142671in}}%
\pgfpathcurveto{\pgfqpoint{9.967700in}{4.135420in}}{\pgfqpoint{9.977536in}{4.131346in}}{\pgfqpoint{9.987790in}{4.131346in}}%
\pgfpathlineto{\pgfqpoint{9.987790in}{4.131346in}}%
\pgfpathclose%
\pgfusepath{stroke}%
\end{pgfscope}%
\begin{pgfscope}%
\pgfpathrectangle{\pgfqpoint{7.904091in}{2.928000in}}{\pgfqpoint{5.664000in}{5.664000in}}%
\pgfusepath{clip}%
\pgfsetbuttcap%
\pgfsetroundjoin%
\pgfsetlinewidth{1.003750pt}%
\definecolor{currentstroke}{rgb}{0.541176,0.380392,0.658824}%
\pgfsetstrokecolor{currentstroke}%
\pgfsetstrokeopacity{0.900000}%
\pgfsetdash{}{0pt}%
\pgfpathmoveto{\pgfqpoint{10.660526in}{5.581221in}}%
\pgfpathcurveto{\pgfqpoint{10.670780in}{5.581221in}}{\pgfqpoint{10.680615in}{5.585295in}}{\pgfqpoint{10.687866in}{5.592546in}}%
\pgfpathcurveto{\pgfqpoint{10.695117in}{5.599797in}}{\pgfqpoint{10.699191in}{5.609632in}}{\pgfqpoint{10.699191in}{5.619886in}}%
\pgfpathcurveto{\pgfqpoint{10.699191in}{5.630140in}}{\pgfqpoint{10.695117in}{5.639976in}}{\pgfqpoint{10.687866in}{5.647226in}}%
\pgfpathcurveto{\pgfqpoint{10.680615in}{5.654477in}}{\pgfqpoint{10.670780in}{5.658551in}}{\pgfqpoint{10.660526in}{5.658551in}}%
\pgfpathcurveto{\pgfqpoint{10.650272in}{5.658551in}}{\pgfqpoint{10.640436in}{5.654477in}}{\pgfqpoint{10.633186in}{5.647226in}}%
\pgfpathcurveto{\pgfqpoint{10.625935in}{5.639976in}}{\pgfqpoint{10.621861in}{5.630140in}}{\pgfqpoint{10.621861in}{5.619886in}}%
\pgfpathcurveto{\pgfqpoint{10.621861in}{5.609632in}}{\pgfqpoint{10.625935in}{5.599797in}}{\pgfqpoint{10.633186in}{5.592546in}}%
\pgfpathcurveto{\pgfqpoint{10.640436in}{5.585295in}}{\pgfqpoint{10.650272in}{5.581221in}}{\pgfqpoint{10.660526in}{5.581221in}}%
\pgfpathlineto{\pgfqpoint{10.660526in}{5.581221in}}%
\pgfpathclose%
\pgfusepath{stroke}%
\end{pgfscope}%
\begin{pgfscope}%
\pgfpathrectangle{\pgfqpoint{7.904091in}{2.928000in}}{\pgfqpoint{5.664000in}{5.664000in}}%
\pgfusepath{clip}%
\pgfsetbuttcap%
\pgfsetroundjoin%
\pgfsetlinewidth{1.003750pt}%
\definecolor{currentstroke}{rgb}{0.541176,0.380392,0.658824}%
\pgfsetstrokecolor{currentstroke}%
\pgfsetstrokeopacity{0.900000}%
\pgfsetdash{}{0pt}%
\pgfpathmoveto{\pgfqpoint{10.652943in}{4.983778in}}%
\pgfpathcurveto{\pgfqpoint{10.663197in}{4.983778in}}{\pgfqpoint{10.673032in}{4.987852in}}{\pgfqpoint{10.680283in}{4.995103in}}%
\pgfpathcurveto{\pgfqpoint{10.687534in}{5.002353in}}{\pgfqpoint{10.691608in}{5.012189in}}{\pgfqpoint{10.691608in}{5.022443in}}%
\pgfpathcurveto{\pgfqpoint{10.691608in}{5.032697in}}{\pgfqpoint{10.687534in}{5.042532in}}{\pgfqpoint{10.680283in}{5.049783in}}%
\pgfpathcurveto{\pgfqpoint{10.673032in}{5.057034in}}{\pgfqpoint{10.663197in}{5.061108in}}{\pgfqpoint{10.652943in}{5.061108in}}%
\pgfpathcurveto{\pgfqpoint{10.642689in}{5.061108in}}{\pgfqpoint{10.632853in}{5.057034in}}{\pgfqpoint{10.625602in}{5.049783in}}%
\pgfpathcurveto{\pgfqpoint{10.618352in}{5.042532in}}{\pgfqpoint{10.614278in}{5.032697in}}{\pgfqpoint{10.614278in}{5.022443in}}%
\pgfpathcurveto{\pgfqpoint{10.614278in}{5.012189in}}{\pgfqpoint{10.618352in}{5.002353in}}{\pgfqpoint{10.625602in}{4.995103in}}%
\pgfpathcurveto{\pgfqpoint{10.632853in}{4.987852in}}{\pgfqpoint{10.642689in}{4.983778in}}{\pgfqpoint{10.652943in}{4.983778in}}%
\pgfpathlineto{\pgfqpoint{10.652943in}{4.983778in}}%
\pgfpathclose%
\pgfusepath{stroke}%
\end{pgfscope}%
\begin{pgfscope}%
\pgfpathrectangle{\pgfqpoint{7.904091in}{2.928000in}}{\pgfqpoint{5.664000in}{5.664000in}}%
\pgfusepath{clip}%
\pgfsetbuttcap%
\pgfsetroundjoin%
\pgfsetlinewidth{1.003750pt}%
\definecolor{currentstroke}{rgb}{0.541176,0.380392,0.658824}%
\pgfsetstrokecolor{currentstroke}%
\pgfsetstrokeopacity{0.900000}%
\pgfsetdash{}{0pt}%
\pgfpathmoveto{\pgfqpoint{10.260040in}{4.638356in}}%
\pgfpathcurveto{\pgfqpoint{10.270294in}{4.638356in}}{\pgfqpoint{10.280129in}{4.642430in}}{\pgfqpoint{10.287380in}{4.649681in}}%
\pgfpathcurveto{\pgfqpoint{10.294631in}{4.656931in}}{\pgfqpoint{10.298705in}{4.666767in}}{\pgfqpoint{10.298705in}{4.677021in}}%
\pgfpathcurveto{\pgfqpoint{10.298705in}{4.687275in}}{\pgfqpoint{10.294631in}{4.697110in}}{\pgfqpoint{10.287380in}{4.704361in}}%
\pgfpathcurveto{\pgfqpoint{10.280129in}{4.711612in}}{\pgfqpoint{10.270294in}{4.715686in}}{\pgfqpoint{10.260040in}{4.715686in}}%
\pgfpathcurveto{\pgfqpoint{10.249786in}{4.715686in}}{\pgfqpoint{10.239950in}{4.711612in}}{\pgfqpoint{10.232699in}{4.704361in}}%
\pgfpathcurveto{\pgfqpoint{10.225449in}{4.697110in}}{\pgfqpoint{10.221375in}{4.687275in}}{\pgfqpoint{10.221375in}{4.677021in}}%
\pgfpathcurveto{\pgfqpoint{10.221375in}{4.666767in}}{\pgfqpoint{10.225449in}{4.656931in}}{\pgfqpoint{10.232699in}{4.649681in}}%
\pgfpathcurveto{\pgfqpoint{10.239950in}{4.642430in}}{\pgfqpoint{10.249786in}{4.638356in}}{\pgfqpoint{10.260040in}{4.638356in}}%
\pgfpathlineto{\pgfqpoint{10.260040in}{4.638356in}}%
\pgfpathclose%
\pgfusepath{stroke}%
\end{pgfscope}%
\begin{pgfscope}%
\pgfpathrectangle{\pgfqpoint{7.904091in}{2.928000in}}{\pgfqpoint{5.664000in}{5.664000in}}%
\pgfusepath{clip}%
\pgfsetbuttcap%
\pgfsetroundjoin%
\pgfsetlinewidth{1.003750pt}%
\definecolor{currentstroke}{rgb}{0.541176,0.380392,0.658824}%
\pgfsetstrokecolor{currentstroke}%
\pgfsetstrokeopacity{0.900000}%
\pgfsetdash{}{0pt}%
\pgfpathmoveto{\pgfqpoint{10.378671in}{4.994038in}}%
\pgfpathcurveto{\pgfqpoint{10.388925in}{4.994038in}}{\pgfqpoint{10.398760in}{4.998112in}}{\pgfqpoint{10.406011in}{5.005362in}}%
\pgfpathcurveto{\pgfqpoint{10.413262in}{5.012613in}}{\pgfqpoint{10.417336in}{5.022449in}}{\pgfqpoint{10.417336in}{5.032703in}}%
\pgfpathcurveto{\pgfqpoint{10.417336in}{5.042957in}}{\pgfqpoint{10.413262in}{5.052792in}}{\pgfqpoint{10.406011in}{5.060043in}}%
\pgfpathcurveto{\pgfqpoint{10.398760in}{5.067294in}}{\pgfqpoint{10.388925in}{5.071368in}}{\pgfqpoint{10.378671in}{5.071368in}}%
\pgfpathcurveto{\pgfqpoint{10.368417in}{5.071368in}}{\pgfqpoint{10.358581in}{5.067294in}}{\pgfqpoint{10.351330in}{5.060043in}}%
\pgfpathcurveto{\pgfqpoint{10.344080in}{5.052792in}}{\pgfqpoint{10.340006in}{5.042957in}}{\pgfqpoint{10.340006in}{5.032703in}}%
\pgfpathcurveto{\pgfqpoint{10.340006in}{5.022449in}}{\pgfqpoint{10.344080in}{5.012613in}}{\pgfqpoint{10.351330in}{5.005362in}}%
\pgfpathcurveto{\pgfqpoint{10.358581in}{4.998112in}}{\pgfqpoint{10.368417in}{4.994038in}}{\pgfqpoint{10.378671in}{4.994038in}}%
\pgfpathlineto{\pgfqpoint{10.378671in}{4.994038in}}%
\pgfpathclose%
\pgfusepath{stroke}%
\end{pgfscope}%
\begin{pgfscope}%
\pgfpathrectangle{\pgfqpoint{7.904091in}{2.928000in}}{\pgfqpoint{5.664000in}{5.664000in}}%
\pgfusepath{clip}%
\pgfsetbuttcap%
\pgfsetroundjoin%
\pgfsetlinewidth{1.003750pt}%
\definecolor{currentstroke}{rgb}{0.541176,0.380392,0.658824}%
\pgfsetstrokecolor{currentstroke}%
\pgfsetstrokeopacity{0.900000}%
\pgfsetdash{}{0pt}%
\pgfpathmoveto{\pgfqpoint{10.634833in}{5.642632in}}%
\pgfpathcurveto{\pgfqpoint{10.645087in}{5.642632in}}{\pgfqpoint{10.654922in}{5.646706in}}{\pgfqpoint{10.662173in}{5.653957in}}%
\pgfpathcurveto{\pgfqpoint{10.669424in}{5.661208in}}{\pgfqpoint{10.673498in}{5.671043in}}{\pgfqpoint{10.673498in}{5.681297in}}%
\pgfpathcurveto{\pgfqpoint{10.673498in}{5.691552in}}{\pgfqpoint{10.669424in}{5.701387in}}{\pgfqpoint{10.662173in}{5.708638in}}%
\pgfpathcurveto{\pgfqpoint{10.654922in}{5.715888in}}{\pgfqpoint{10.645087in}{5.719962in}}{\pgfqpoint{10.634833in}{5.719962in}}%
\pgfpathcurveto{\pgfqpoint{10.624579in}{5.719962in}}{\pgfqpoint{10.614743in}{5.715888in}}{\pgfqpoint{10.607493in}{5.708638in}}%
\pgfpathcurveto{\pgfqpoint{10.600242in}{5.701387in}}{\pgfqpoint{10.596168in}{5.691552in}}{\pgfqpoint{10.596168in}{5.681297in}}%
\pgfpathcurveto{\pgfqpoint{10.596168in}{5.671043in}}{\pgfqpoint{10.600242in}{5.661208in}}{\pgfqpoint{10.607493in}{5.653957in}}%
\pgfpathcurveto{\pgfqpoint{10.614743in}{5.646706in}}{\pgfqpoint{10.624579in}{5.642632in}}{\pgfqpoint{10.634833in}{5.642632in}}%
\pgfpathlineto{\pgfqpoint{10.634833in}{5.642632in}}%
\pgfpathclose%
\pgfusepath{stroke}%
\end{pgfscope}%
\begin{pgfscope}%
\pgfpathrectangle{\pgfqpoint{7.904091in}{2.928000in}}{\pgfqpoint{5.664000in}{5.664000in}}%
\pgfusepath{clip}%
\pgfsetbuttcap%
\pgfsetroundjoin%
\pgfsetlinewidth{1.003750pt}%
\definecolor{currentstroke}{rgb}{0.541176,0.380392,0.658824}%
\pgfsetstrokecolor{currentstroke}%
\pgfsetstrokeopacity{0.900000}%
\pgfsetdash{}{0pt}%
\pgfpathmoveto{\pgfqpoint{11.349641in}{6.064223in}}%
\pgfpathcurveto{\pgfqpoint{11.359895in}{6.064223in}}{\pgfqpoint{11.369730in}{6.068297in}}{\pgfqpoint{11.376981in}{6.075548in}}%
\pgfpathcurveto{\pgfqpoint{11.384232in}{6.082799in}}{\pgfqpoint{11.388306in}{6.092634in}}{\pgfqpoint{11.388306in}{6.102888in}}%
\pgfpathcurveto{\pgfqpoint{11.388306in}{6.113142in}}{\pgfqpoint{11.384232in}{6.122978in}}{\pgfqpoint{11.376981in}{6.130229in}}%
\pgfpathcurveto{\pgfqpoint{11.369730in}{6.137479in}}{\pgfqpoint{11.359895in}{6.141553in}}{\pgfqpoint{11.349641in}{6.141553in}}%
\pgfpathcurveto{\pgfqpoint{11.339387in}{6.141553in}}{\pgfqpoint{11.329551in}{6.137479in}}{\pgfqpoint{11.322301in}{6.130229in}}%
\pgfpathcurveto{\pgfqpoint{11.315050in}{6.122978in}}{\pgfqpoint{11.310976in}{6.113142in}}{\pgfqpoint{11.310976in}{6.102888in}}%
\pgfpathcurveto{\pgfqpoint{11.310976in}{6.092634in}}{\pgfqpoint{11.315050in}{6.082799in}}{\pgfqpoint{11.322301in}{6.075548in}}%
\pgfpathcurveto{\pgfqpoint{11.329551in}{6.068297in}}{\pgfqpoint{11.339387in}{6.064223in}}{\pgfqpoint{11.349641in}{6.064223in}}%
\pgfpathlineto{\pgfqpoint{11.349641in}{6.064223in}}%
\pgfpathclose%
\pgfusepath{stroke}%
\end{pgfscope}%
\begin{pgfscope}%
\pgfpathrectangle{\pgfqpoint{7.904091in}{2.928000in}}{\pgfqpoint{5.664000in}{5.664000in}}%
\pgfusepath{clip}%
\pgfsetbuttcap%
\pgfsetroundjoin%
\pgfsetlinewidth{1.003750pt}%
\definecolor{currentstroke}{rgb}{0.541176,0.380392,0.658824}%
\pgfsetstrokecolor{currentstroke}%
\pgfsetstrokeopacity{0.900000}%
\pgfsetdash{}{0pt}%
\pgfpathmoveto{\pgfqpoint{11.953455in}{6.733978in}}%
\pgfpathcurveto{\pgfqpoint{11.963709in}{6.733978in}}{\pgfqpoint{11.973545in}{6.738052in}}{\pgfqpoint{11.980795in}{6.745302in}}%
\pgfpathcurveto{\pgfqpoint{11.988046in}{6.752553in}}{\pgfqpoint{11.992120in}{6.762389in}}{\pgfqpoint{11.992120in}{6.772643in}}%
\pgfpathcurveto{\pgfqpoint{11.992120in}{6.782897in}}{\pgfqpoint{11.988046in}{6.792732in}}{\pgfqpoint{11.980795in}{6.799983in}}%
\pgfpathcurveto{\pgfqpoint{11.973545in}{6.807234in}}{\pgfqpoint{11.963709in}{6.811308in}}{\pgfqpoint{11.953455in}{6.811308in}}%
\pgfpathcurveto{\pgfqpoint{11.943201in}{6.811308in}}{\pgfqpoint{11.933365in}{6.807234in}}{\pgfqpoint{11.926115in}{6.799983in}}%
\pgfpathcurveto{\pgfqpoint{11.918864in}{6.792732in}}{\pgfqpoint{11.914790in}{6.782897in}}{\pgfqpoint{11.914790in}{6.772643in}}%
\pgfpathcurveto{\pgfqpoint{11.914790in}{6.762389in}}{\pgfqpoint{11.918864in}{6.752553in}}{\pgfqpoint{11.926115in}{6.745302in}}%
\pgfpathcurveto{\pgfqpoint{11.933365in}{6.738052in}}{\pgfqpoint{11.943201in}{6.733978in}}{\pgfqpoint{11.953455in}{6.733978in}}%
\pgfpathlineto{\pgfqpoint{11.953455in}{6.733978in}}%
\pgfpathclose%
\pgfusepath{stroke}%
\end{pgfscope}%
\begin{pgfscope}%
\pgfpathrectangle{\pgfqpoint{7.904091in}{2.928000in}}{\pgfqpoint{5.664000in}{5.664000in}}%
\pgfusepath{clip}%
\pgfsetbuttcap%
\pgfsetroundjoin%
\pgfsetlinewidth{1.003750pt}%
\definecolor{currentstroke}{rgb}{0.541176,0.380392,0.658824}%
\pgfsetstrokecolor{currentstroke}%
\pgfsetstrokeopacity{0.900000}%
\pgfsetdash{}{0pt}%
\pgfpathmoveto{\pgfqpoint{11.374624in}{5.967054in}}%
\pgfpathcurveto{\pgfqpoint{11.384879in}{5.967054in}}{\pgfqpoint{11.394714in}{5.971128in}}{\pgfqpoint{11.401965in}{5.978379in}}%
\pgfpathcurveto{\pgfqpoint{11.409215in}{5.985630in}}{\pgfqpoint{11.413289in}{5.995465in}}{\pgfqpoint{11.413289in}{6.005719in}}%
\pgfpathcurveto{\pgfqpoint{11.413289in}{6.015973in}}{\pgfqpoint{11.409215in}{6.025809in}}{\pgfqpoint{11.401965in}{6.033059in}}%
\pgfpathcurveto{\pgfqpoint{11.394714in}{6.040310in}}{\pgfqpoint{11.384879in}{6.044384in}}{\pgfqpoint{11.374624in}{6.044384in}}%
\pgfpathcurveto{\pgfqpoint{11.364370in}{6.044384in}}{\pgfqpoint{11.354535in}{6.040310in}}{\pgfqpoint{11.347284in}{6.033059in}}%
\pgfpathcurveto{\pgfqpoint{11.340033in}{6.025809in}}{\pgfqpoint{11.335959in}{6.015973in}}{\pgfqpoint{11.335959in}{6.005719in}}%
\pgfpathcurveto{\pgfqpoint{11.335959in}{5.995465in}}{\pgfqpoint{11.340033in}{5.985630in}}{\pgfqpoint{11.347284in}{5.978379in}}%
\pgfpathcurveto{\pgfqpoint{11.354535in}{5.971128in}}{\pgfqpoint{11.364370in}{5.967054in}}{\pgfqpoint{11.374624in}{5.967054in}}%
\pgfpathlineto{\pgfqpoint{11.374624in}{5.967054in}}%
\pgfpathclose%
\pgfusepath{stroke}%
\end{pgfscope}%
\begin{pgfscope}%
\pgfpathrectangle{\pgfqpoint{7.904091in}{2.928000in}}{\pgfqpoint{5.664000in}{5.664000in}}%
\pgfusepath{clip}%
\pgfsetbuttcap%
\pgfsetroundjoin%
\pgfsetlinewidth{1.003750pt}%
\definecolor{currentstroke}{rgb}{0.541176,0.380392,0.658824}%
\pgfsetstrokecolor{currentstroke}%
\pgfsetstrokeopacity{0.900000}%
\pgfsetdash{}{0pt}%
\pgfpathmoveto{\pgfqpoint{11.337589in}{6.359740in}}%
\pgfpathcurveto{\pgfqpoint{11.347843in}{6.359740in}}{\pgfqpoint{11.357678in}{6.363814in}}{\pgfqpoint{11.364929in}{6.371065in}}%
\pgfpathcurveto{\pgfqpoint{11.372180in}{6.378315in}}{\pgfqpoint{11.376254in}{6.388151in}}{\pgfqpoint{11.376254in}{6.398405in}}%
\pgfpathcurveto{\pgfqpoint{11.376254in}{6.408659in}}{\pgfqpoint{11.372180in}{6.418495in}}{\pgfqpoint{11.364929in}{6.425745in}}%
\pgfpathcurveto{\pgfqpoint{11.357678in}{6.432996in}}{\pgfqpoint{11.347843in}{6.437070in}}{\pgfqpoint{11.337589in}{6.437070in}}%
\pgfpathcurveto{\pgfqpoint{11.327335in}{6.437070in}}{\pgfqpoint{11.317499in}{6.432996in}}{\pgfqpoint{11.310248in}{6.425745in}}%
\pgfpathcurveto{\pgfqpoint{11.302998in}{6.418495in}}{\pgfqpoint{11.298924in}{6.408659in}}{\pgfqpoint{11.298924in}{6.398405in}}%
\pgfpathcurveto{\pgfqpoint{11.298924in}{6.388151in}}{\pgfqpoint{11.302998in}{6.378315in}}{\pgfqpoint{11.310248in}{6.371065in}}%
\pgfpathcurveto{\pgfqpoint{11.317499in}{6.363814in}}{\pgfqpoint{11.327335in}{6.359740in}}{\pgfqpoint{11.337589in}{6.359740in}}%
\pgfpathlineto{\pgfqpoint{11.337589in}{6.359740in}}%
\pgfpathclose%
\pgfusepath{stroke}%
\end{pgfscope}%
\begin{pgfscope}%
\pgfpathrectangle{\pgfqpoint{7.904091in}{2.928000in}}{\pgfqpoint{5.664000in}{5.664000in}}%
\pgfusepath{clip}%
\pgfsetbuttcap%
\pgfsetroundjoin%
\pgfsetlinewidth{1.003750pt}%
\definecolor{currentstroke}{rgb}{0.541176,0.380392,0.658824}%
\pgfsetstrokecolor{currentstroke}%
\pgfsetstrokeopacity{0.900000}%
\pgfsetdash{}{0pt}%
\pgfpathmoveto{\pgfqpoint{11.548192in}{6.182534in}}%
\pgfpathcurveto{\pgfqpoint{11.558446in}{6.182534in}}{\pgfqpoint{11.568282in}{6.186608in}}{\pgfqpoint{11.575533in}{6.193858in}}%
\pgfpathcurveto{\pgfqpoint{11.582783in}{6.201109in}}{\pgfqpoint{11.586857in}{6.210945in}}{\pgfqpoint{11.586857in}{6.221199in}}%
\pgfpathcurveto{\pgfqpoint{11.586857in}{6.231453in}}{\pgfqpoint{11.582783in}{6.241288in}}{\pgfqpoint{11.575533in}{6.248539in}}%
\pgfpathcurveto{\pgfqpoint{11.568282in}{6.255790in}}{\pgfqpoint{11.558446in}{6.259864in}}{\pgfqpoint{11.548192in}{6.259864in}}%
\pgfpathcurveto{\pgfqpoint{11.537938in}{6.259864in}}{\pgfqpoint{11.528103in}{6.255790in}}{\pgfqpoint{11.520852in}{6.248539in}}%
\pgfpathcurveto{\pgfqpoint{11.513601in}{6.241288in}}{\pgfqpoint{11.509527in}{6.231453in}}{\pgfqpoint{11.509527in}{6.221199in}}%
\pgfpathcurveto{\pgfqpoint{11.509527in}{6.210945in}}{\pgfqpoint{11.513601in}{6.201109in}}{\pgfqpoint{11.520852in}{6.193858in}}%
\pgfpathcurveto{\pgfqpoint{11.528103in}{6.186608in}}{\pgfqpoint{11.537938in}{6.182534in}}{\pgfqpoint{11.548192in}{6.182534in}}%
\pgfpathlineto{\pgfqpoint{11.548192in}{6.182534in}}%
\pgfpathclose%
\pgfusepath{stroke}%
\end{pgfscope}%
\begin{pgfscope}%
\pgfpathrectangle{\pgfqpoint{7.904091in}{2.928000in}}{\pgfqpoint{5.664000in}{5.664000in}}%
\pgfusepath{clip}%
\pgfsetbuttcap%
\pgfsetroundjoin%
\pgfsetlinewidth{1.003750pt}%
\definecolor{currentstroke}{rgb}{0.541176,0.380392,0.658824}%
\pgfsetstrokecolor{currentstroke}%
\pgfsetstrokeopacity{0.900000}%
\pgfsetdash{}{0pt}%
\pgfpathmoveto{\pgfqpoint{11.655069in}{6.445216in}}%
\pgfpathcurveto{\pgfqpoint{11.665323in}{6.445216in}}{\pgfqpoint{11.675158in}{6.449290in}}{\pgfqpoint{11.682409in}{6.456541in}}%
\pgfpathcurveto{\pgfqpoint{11.689660in}{6.463792in}}{\pgfqpoint{11.693734in}{6.473627in}}{\pgfqpoint{11.693734in}{6.483881in}}%
\pgfpathcurveto{\pgfqpoint{11.693734in}{6.494135in}}{\pgfqpoint{11.689660in}{6.503971in}}{\pgfqpoint{11.682409in}{6.511222in}}%
\pgfpathcurveto{\pgfqpoint{11.675158in}{6.518472in}}{\pgfqpoint{11.665323in}{6.522546in}}{\pgfqpoint{11.655069in}{6.522546in}}%
\pgfpathcurveto{\pgfqpoint{11.644815in}{6.522546in}}{\pgfqpoint{11.634979in}{6.518472in}}{\pgfqpoint{11.627728in}{6.511222in}}%
\pgfpathcurveto{\pgfqpoint{11.620478in}{6.503971in}}{\pgfqpoint{11.616404in}{6.494135in}}{\pgfqpoint{11.616404in}{6.483881in}}%
\pgfpathcurveto{\pgfqpoint{11.616404in}{6.473627in}}{\pgfqpoint{11.620478in}{6.463792in}}{\pgfqpoint{11.627728in}{6.456541in}}%
\pgfpathcurveto{\pgfqpoint{11.634979in}{6.449290in}}{\pgfqpoint{11.644815in}{6.445216in}}{\pgfqpoint{11.655069in}{6.445216in}}%
\pgfpathlineto{\pgfqpoint{11.655069in}{6.445216in}}%
\pgfpathclose%
\pgfusepath{stroke}%
\end{pgfscope}%
\begin{pgfscope}%
\pgfpathrectangle{\pgfqpoint{7.904091in}{2.928000in}}{\pgfqpoint{5.664000in}{5.664000in}}%
\pgfusepath{clip}%
\pgfsetbuttcap%
\pgfsetroundjoin%
\pgfsetlinewidth{1.003750pt}%
\definecolor{currentstroke}{rgb}{0.541176,0.380392,0.658824}%
\pgfsetstrokecolor{currentstroke}%
\pgfsetstrokeopacity{0.900000}%
\pgfsetdash{}{0pt}%
\pgfpathmoveto{\pgfqpoint{12.803014in}{7.586508in}}%
\pgfpathcurveto{\pgfqpoint{12.813268in}{7.586508in}}{\pgfqpoint{12.823104in}{7.590582in}}{\pgfqpoint{12.830354in}{7.597833in}}%
\pgfpathcurveto{\pgfqpoint{12.837605in}{7.605083in}}{\pgfqpoint{12.841679in}{7.614919in}}{\pgfqpoint{12.841679in}{7.625173in}}%
\pgfpathcurveto{\pgfqpoint{12.841679in}{7.635427in}}{\pgfqpoint{12.837605in}{7.645263in}}{\pgfqpoint{12.830354in}{7.652513in}}%
\pgfpathcurveto{\pgfqpoint{12.823104in}{7.659764in}}{\pgfqpoint{12.813268in}{7.663838in}}{\pgfqpoint{12.803014in}{7.663838in}}%
\pgfpathcurveto{\pgfqpoint{12.792760in}{7.663838in}}{\pgfqpoint{12.782925in}{7.659764in}}{\pgfqpoint{12.775674in}{7.652513in}}%
\pgfpathcurveto{\pgfqpoint{12.768423in}{7.645263in}}{\pgfqpoint{12.764349in}{7.635427in}}{\pgfqpoint{12.764349in}{7.625173in}}%
\pgfpathcurveto{\pgfqpoint{12.764349in}{7.614919in}}{\pgfqpoint{12.768423in}{7.605083in}}{\pgfqpoint{12.775674in}{7.597833in}}%
\pgfpathcurveto{\pgfqpoint{12.782925in}{7.590582in}}{\pgfqpoint{12.792760in}{7.586508in}}{\pgfqpoint{12.803014in}{7.586508in}}%
\pgfpathlineto{\pgfqpoint{12.803014in}{7.586508in}}%
\pgfpathclose%
\pgfusepath{stroke}%
\end{pgfscope}%
\begin{pgfscope}%
\pgfpathrectangle{\pgfqpoint{7.904091in}{2.928000in}}{\pgfqpoint{5.664000in}{5.664000in}}%
\pgfusepath{clip}%
\pgfsetbuttcap%
\pgfsetroundjoin%
\pgfsetlinewidth{1.003750pt}%
\definecolor{currentstroke}{rgb}{0.541176,0.380392,0.658824}%
\pgfsetstrokecolor{currentstroke}%
\pgfsetstrokeopacity{0.900000}%
\pgfsetdash{}{0pt}%
\pgfpathmoveto{\pgfqpoint{12.366611in}{7.305598in}}%
\pgfpathcurveto{\pgfqpoint{12.376865in}{7.305598in}}{\pgfqpoint{12.386700in}{7.309672in}}{\pgfqpoint{12.393951in}{7.316923in}}%
\pgfpathcurveto{\pgfqpoint{12.401202in}{7.324173in}}{\pgfqpoint{12.405276in}{7.334009in}}{\pgfqpoint{12.405276in}{7.344263in}}%
\pgfpathcurveto{\pgfqpoint{12.405276in}{7.354517in}}{\pgfqpoint{12.401202in}{7.364352in}}{\pgfqpoint{12.393951in}{7.371603in}}%
\pgfpathcurveto{\pgfqpoint{12.386700in}{7.378854in}}{\pgfqpoint{12.376865in}{7.382928in}}{\pgfqpoint{12.366611in}{7.382928in}}%
\pgfpathcurveto{\pgfqpoint{12.356357in}{7.382928in}}{\pgfqpoint{12.346521in}{7.378854in}}{\pgfqpoint{12.339270in}{7.371603in}}%
\pgfpathcurveto{\pgfqpoint{12.332020in}{7.364352in}}{\pgfqpoint{12.327946in}{7.354517in}}{\pgfqpoint{12.327946in}{7.344263in}}%
\pgfpathcurveto{\pgfqpoint{12.327946in}{7.334009in}}{\pgfqpoint{12.332020in}{7.324173in}}{\pgfqpoint{12.339270in}{7.316923in}}%
\pgfpathcurveto{\pgfqpoint{12.346521in}{7.309672in}}{\pgfqpoint{12.356357in}{7.305598in}}{\pgfqpoint{12.366611in}{7.305598in}}%
\pgfpathlineto{\pgfqpoint{12.366611in}{7.305598in}}%
\pgfpathclose%
\pgfusepath{stroke}%
\end{pgfscope}%
\begin{pgfscope}%
\pgfpathrectangle{\pgfqpoint{7.904091in}{2.928000in}}{\pgfqpoint{5.664000in}{5.664000in}}%
\pgfusepath{clip}%
\pgfsetbuttcap%
\pgfsetroundjoin%
\pgfsetlinewidth{1.003750pt}%
\definecolor{currentstroke}{rgb}{0.541176,0.380392,0.658824}%
\pgfsetstrokecolor{currentstroke}%
\pgfsetstrokeopacity{0.900000}%
\pgfsetdash{}{0pt}%
\pgfpathmoveto{\pgfqpoint{12.145006in}{6.787124in}}%
\pgfpathcurveto{\pgfqpoint{12.155260in}{6.787124in}}{\pgfqpoint{12.165096in}{6.791198in}}{\pgfqpoint{12.172346in}{6.798448in}}%
\pgfpathcurveto{\pgfqpoint{12.179597in}{6.805699in}}{\pgfqpoint{12.183671in}{6.815535in}}{\pgfqpoint{12.183671in}{6.825789in}}%
\pgfpathcurveto{\pgfqpoint{12.183671in}{6.836043in}}{\pgfqpoint{12.179597in}{6.845878in}}{\pgfqpoint{12.172346in}{6.853129in}}%
\pgfpathcurveto{\pgfqpoint{12.165096in}{6.860380in}}{\pgfqpoint{12.155260in}{6.864454in}}{\pgfqpoint{12.145006in}{6.864454in}}%
\pgfpathcurveto{\pgfqpoint{12.134752in}{6.864454in}}{\pgfqpoint{12.124916in}{6.860380in}}{\pgfqpoint{12.117666in}{6.853129in}}%
\pgfpathcurveto{\pgfqpoint{12.110415in}{6.845878in}}{\pgfqpoint{12.106341in}{6.836043in}}{\pgfqpoint{12.106341in}{6.825789in}}%
\pgfpathcurveto{\pgfqpoint{12.106341in}{6.815535in}}{\pgfqpoint{12.110415in}{6.805699in}}{\pgfqpoint{12.117666in}{6.798448in}}%
\pgfpathcurveto{\pgfqpoint{12.124916in}{6.791198in}}{\pgfqpoint{12.134752in}{6.787124in}}{\pgfqpoint{12.145006in}{6.787124in}}%
\pgfpathlineto{\pgfqpoint{12.145006in}{6.787124in}}%
\pgfpathclose%
\pgfusepath{stroke}%
\end{pgfscope}%
\begin{pgfscope}%
\pgfpathrectangle{\pgfqpoint{7.904091in}{2.928000in}}{\pgfqpoint{5.664000in}{5.664000in}}%
\pgfusepath{clip}%
\pgfsetbuttcap%
\pgfsetroundjoin%
\pgfsetlinewidth{1.003750pt}%
\definecolor{currentstroke}{rgb}{0.541176,0.380392,0.658824}%
\pgfsetstrokecolor{currentstroke}%
\pgfsetstrokeopacity{0.900000}%
\pgfsetdash{}{0pt}%
\pgfpathmoveto{\pgfqpoint{11.778472in}{6.612164in}}%
\pgfpathcurveto{\pgfqpoint{11.788726in}{6.612164in}}{\pgfqpoint{11.798561in}{6.616238in}}{\pgfqpoint{11.805812in}{6.623489in}}%
\pgfpathcurveto{\pgfqpoint{11.813063in}{6.630740in}}{\pgfqpoint{11.817137in}{6.640575in}}{\pgfqpoint{11.817137in}{6.650829in}}%
\pgfpathcurveto{\pgfqpoint{11.817137in}{6.661083in}}{\pgfqpoint{11.813063in}{6.670919in}}{\pgfqpoint{11.805812in}{6.678169in}}%
\pgfpathcurveto{\pgfqpoint{11.798561in}{6.685420in}}{\pgfqpoint{11.788726in}{6.689494in}}{\pgfqpoint{11.778472in}{6.689494in}}%
\pgfpathcurveto{\pgfqpoint{11.768218in}{6.689494in}}{\pgfqpoint{11.758382in}{6.685420in}}{\pgfqpoint{11.751131in}{6.678169in}}%
\pgfpathcurveto{\pgfqpoint{11.743881in}{6.670919in}}{\pgfqpoint{11.739807in}{6.661083in}}{\pgfqpoint{11.739807in}{6.650829in}}%
\pgfpathcurveto{\pgfqpoint{11.739807in}{6.640575in}}{\pgfqpoint{11.743881in}{6.630740in}}{\pgfqpoint{11.751131in}{6.623489in}}%
\pgfpathcurveto{\pgfqpoint{11.758382in}{6.616238in}}{\pgfqpoint{11.768218in}{6.612164in}}{\pgfqpoint{11.778472in}{6.612164in}}%
\pgfpathlineto{\pgfqpoint{11.778472in}{6.612164in}}%
\pgfpathclose%
\pgfusepath{stroke}%
\end{pgfscope}%
\begin{pgfscope}%
\pgfpathrectangle{\pgfqpoint{7.904091in}{2.928000in}}{\pgfqpoint{5.664000in}{5.664000in}}%
\pgfusepath{clip}%
\pgfsetbuttcap%
\pgfsetroundjoin%
\pgfsetlinewidth{1.003750pt}%
\definecolor{currentstroke}{rgb}{0.541176,0.380392,0.658824}%
\pgfsetstrokecolor{currentstroke}%
\pgfsetstrokeopacity{0.900000}%
\pgfsetdash{}{0pt}%
\pgfpathmoveto{\pgfqpoint{12.113612in}{7.256771in}}%
\pgfpathcurveto{\pgfqpoint{12.123866in}{7.256771in}}{\pgfqpoint{12.133701in}{7.260845in}}{\pgfqpoint{12.140952in}{7.268095in}}%
\pgfpathcurveto{\pgfqpoint{12.148203in}{7.275346in}}{\pgfqpoint{12.152277in}{7.285182in}}{\pgfqpoint{12.152277in}{7.295436in}}%
\pgfpathcurveto{\pgfqpoint{12.152277in}{7.305690in}}{\pgfqpoint{12.148203in}{7.315525in}}{\pgfqpoint{12.140952in}{7.322776in}}%
\pgfpathcurveto{\pgfqpoint{12.133701in}{7.330027in}}{\pgfqpoint{12.123866in}{7.334101in}}{\pgfqpoint{12.113612in}{7.334101in}}%
\pgfpathcurveto{\pgfqpoint{12.103358in}{7.334101in}}{\pgfqpoint{12.093522in}{7.330027in}}{\pgfqpoint{12.086271in}{7.322776in}}%
\pgfpathcurveto{\pgfqpoint{12.079021in}{7.315525in}}{\pgfqpoint{12.074947in}{7.305690in}}{\pgfqpoint{12.074947in}{7.295436in}}%
\pgfpathcurveto{\pgfqpoint{12.074947in}{7.285182in}}{\pgfqpoint{12.079021in}{7.275346in}}{\pgfqpoint{12.086271in}{7.268095in}}%
\pgfpathcurveto{\pgfqpoint{12.093522in}{7.260845in}}{\pgfqpoint{12.103358in}{7.256771in}}{\pgfqpoint{12.113612in}{7.256771in}}%
\pgfpathlineto{\pgfqpoint{12.113612in}{7.256771in}}%
\pgfpathclose%
\pgfusepath{stroke}%
\end{pgfscope}%
\begin{pgfscope}%
\pgfpathrectangle{\pgfqpoint{7.904091in}{2.928000in}}{\pgfqpoint{5.664000in}{5.664000in}}%
\pgfusepath{clip}%
\pgfsetbuttcap%
\pgfsetroundjoin%
\pgfsetlinewidth{1.003750pt}%
\definecolor{currentstroke}{rgb}{0.541176,0.380392,0.658824}%
\pgfsetstrokecolor{currentstroke}%
\pgfsetstrokeopacity{0.900000}%
\pgfsetdash{}{0pt}%
\pgfpathmoveto{\pgfqpoint{10.825033in}{5.940521in}}%
\pgfpathcurveto{\pgfqpoint{10.835287in}{5.940521in}}{\pgfqpoint{10.845122in}{5.944595in}}{\pgfqpoint{10.852373in}{5.951845in}}%
\pgfpathcurveto{\pgfqpoint{10.859624in}{5.959096in}}{\pgfqpoint{10.863698in}{5.968932in}}{\pgfqpoint{10.863698in}{5.979186in}}%
\pgfpathcurveto{\pgfqpoint{10.863698in}{5.989440in}}{\pgfqpoint{10.859624in}{5.999275in}}{\pgfqpoint{10.852373in}{6.006526in}}%
\pgfpathcurveto{\pgfqpoint{10.845122in}{6.013777in}}{\pgfqpoint{10.835287in}{6.017851in}}{\pgfqpoint{10.825033in}{6.017851in}}%
\pgfpathcurveto{\pgfqpoint{10.814778in}{6.017851in}}{\pgfqpoint{10.804943in}{6.013777in}}{\pgfqpoint{10.797692in}{6.006526in}}%
\pgfpathcurveto{\pgfqpoint{10.790441in}{5.999275in}}{\pgfqpoint{10.786367in}{5.989440in}}{\pgfqpoint{10.786367in}{5.979186in}}%
\pgfpathcurveto{\pgfqpoint{10.786367in}{5.968932in}}{\pgfqpoint{10.790441in}{5.959096in}}{\pgfqpoint{10.797692in}{5.951845in}}%
\pgfpathcurveto{\pgfqpoint{10.804943in}{5.944595in}}{\pgfqpoint{10.814778in}{5.940521in}}{\pgfqpoint{10.825033in}{5.940521in}}%
\pgfpathlineto{\pgfqpoint{10.825033in}{5.940521in}}%
\pgfpathclose%
\pgfusepath{stroke}%
\end{pgfscope}%
\begin{pgfscope}%
\pgfpathrectangle{\pgfqpoint{7.904091in}{2.928000in}}{\pgfqpoint{5.664000in}{5.664000in}}%
\pgfusepath{clip}%
\pgfsetbuttcap%
\pgfsetroundjoin%
\pgfsetlinewidth{1.003750pt}%
\definecolor{currentstroke}{rgb}{0.541176,0.380392,0.658824}%
\pgfsetstrokecolor{currentstroke}%
\pgfsetstrokeopacity{0.900000}%
\pgfsetdash{}{0pt}%
\pgfpathmoveto{\pgfqpoint{11.372722in}{6.650151in}}%
\pgfpathcurveto{\pgfqpoint{11.382976in}{6.650151in}}{\pgfqpoint{11.392811in}{6.654225in}}{\pgfqpoint{11.400062in}{6.661476in}}%
\pgfpathcurveto{\pgfqpoint{11.407313in}{6.668727in}}{\pgfqpoint{11.411387in}{6.678562in}}{\pgfqpoint{11.411387in}{6.688816in}}%
\pgfpathcurveto{\pgfqpoint{11.411387in}{6.699070in}}{\pgfqpoint{11.407313in}{6.708906in}}{\pgfqpoint{11.400062in}{6.716156in}}%
\pgfpathcurveto{\pgfqpoint{11.392811in}{6.723407in}}{\pgfqpoint{11.382976in}{6.727481in}}{\pgfqpoint{11.372722in}{6.727481in}}%
\pgfpathcurveto{\pgfqpoint{11.362467in}{6.727481in}}{\pgfqpoint{11.352632in}{6.723407in}}{\pgfqpoint{11.345381in}{6.716156in}}%
\pgfpathcurveto{\pgfqpoint{11.338130in}{6.708906in}}{\pgfqpoint{11.334056in}{6.699070in}}{\pgfqpoint{11.334056in}{6.688816in}}%
\pgfpathcurveto{\pgfqpoint{11.334056in}{6.678562in}}{\pgfqpoint{11.338130in}{6.668727in}}{\pgfqpoint{11.345381in}{6.661476in}}%
\pgfpathcurveto{\pgfqpoint{11.352632in}{6.654225in}}{\pgfqpoint{11.362467in}{6.650151in}}{\pgfqpoint{11.372722in}{6.650151in}}%
\pgfpathlineto{\pgfqpoint{11.372722in}{6.650151in}}%
\pgfpathclose%
\pgfusepath{stroke}%
\end{pgfscope}%
\begin{pgfscope}%
\pgfpathrectangle{\pgfqpoint{7.904091in}{2.928000in}}{\pgfqpoint{5.664000in}{5.664000in}}%
\pgfusepath{clip}%
\pgfsetbuttcap%
\pgfsetroundjoin%
\pgfsetlinewidth{1.003750pt}%
\definecolor{currentstroke}{rgb}{0.541176,0.380392,0.658824}%
\pgfsetstrokecolor{currentstroke}%
\pgfsetstrokeopacity{0.900000}%
\pgfsetdash{}{0pt}%
\pgfpathmoveto{\pgfqpoint{11.315645in}{5.825798in}}%
\pgfpathcurveto{\pgfqpoint{11.325899in}{5.825798in}}{\pgfqpoint{11.335734in}{5.829872in}}{\pgfqpoint{11.342985in}{5.837123in}}%
\pgfpathcurveto{\pgfqpoint{11.350236in}{5.844374in}}{\pgfqpoint{11.354310in}{5.854209in}}{\pgfqpoint{11.354310in}{5.864463in}}%
\pgfpathcurveto{\pgfqpoint{11.354310in}{5.874717in}}{\pgfqpoint{11.350236in}{5.884553in}}{\pgfqpoint{11.342985in}{5.891804in}}%
\pgfpathcurveto{\pgfqpoint{11.335734in}{5.899054in}}{\pgfqpoint{11.325899in}{5.903128in}}{\pgfqpoint{11.315645in}{5.903128in}}%
\pgfpathcurveto{\pgfqpoint{11.305391in}{5.903128in}}{\pgfqpoint{11.295555in}{5.899054in}}{\pgfqpoint{11.288304in}{5.891804in}}%
\pgfpathcurveto{\pgfqpoint{11.281054in}{5.884553in}}{\pgfqpoint{11.276980in}{5.874717in}}{\pgfqpoint{11.276980in}{5.864463in}}%
\pgfpathcurveto{\pgfqpoint{11.276980in}{5.854209in}}{\pgfqpoint{11.281054in}{5.844374in}}{\pgfqpoint{11.288304in}{5.837123in}}%
\pgfpathcurveto{\pgfqpoint{11.295555in}{5.829872in}}{\pgfqpoint{11.305391in}{5.825798in}}{\pgfqpoint{11.315645in}{5.825798in}}%
\pgfpathlineto{\pgfqpoint{11.315645in}{5.825798in}}%
\pgfpathclose%
\pgfusepath{stroke}%
\end{pgfscope}%
\begin{pgfscope}%
\pgfpathrectangle{\pgfqpoint{7.904091in}{2.928000in}}{\pgfqpoint{5.664000in}{5.664000in}}%
\pgfusepath{clip}%
\pgfsetbuttcap%
\pgfsetroundjoin%
\pgfsetlinewidth{1.003750pt}%
\definecolor{currentstroke}{rgb}{0.541176,0.380392,0.658824}%
\pgfsetstrokecolor{currentstroke}%
\pgfsetstrokeopacity{0.900000}%
\pgfsetdash{}{0pt}%
\pgfpathmoveto{\pgfqpoint{10.237222in}{5.087656in}}%
\pgfpathcurveto{\pgfqpoint{10.247476in}{5.087656in}}{\pgfqpoint{10.257311in}{5.091730in}}{\pgfqpoint{10.264562in}{5.098981in}}%
\pgfpathcurveto{\pgfqpoint{10.271813in}{5.106231in}}{\pgfqpoint{10.275887in}{5.116067in}}{\pgfqpoint{10.275887in}{5.126321in}}%
\pgfpathcurveto{\pgfqpoint{10.275887in}{5.136575in}}{\pgfqpoint{10.271813in}{5.146411in}}{\pgfqpoint{10.264562in}{5.153661in}}%
\pgfpathcurveto{\pgfqpoint{10.257311in}{5.160912in}}{\pgfqpoint{10.247476in}{5.164986in}}{\pgfqpoint{10.237222in}{5.164986in}}%
\pgfpathcurveto{\pgfqpoint{10.226968in}{5.164986in}}{\pgfqpoint{10.217132in}{5.160912in}}{\pgfqpoint{10.209881in}{5.153661in}}%
\pgfpathcurveto{\pgfqpoint{10.202631in}{5.146411in}}{\pgfqpoint{10.198557in}{5.136575in}}{\pgfqpoint{10.198557in}{5.126321in}}%
\pgfpathcurveto{\pgfqpoint{10.198557in}{5.116067in}}{\pgfqpoint{10.202631in}{5.106231in}}{\pgfqpoint{10.209881in}{5.098981in}}%
\pgfpathcurveto{\pgfqpoint{10.217132in}{5.091730in}}{\pgfqpoint{10.226968in}{5.087656in}}{\pgfqpoint{10.237222in}{5.087656in}}%
\pgfpathlineto{\pgfqpoint{10.237222in}{5.087656in}}%
\pgfpathclose%
\pgfusepath{stroke}%
\end{pgfscope}%
\begin{pgfscope}%
\pgfpathrectangle{\pgfqpoint{7.904091in}{2.928000in}}{\pgfqpoint{5.664000in}{5.664000in}}%
\pgfusepath{clip}%
\pgfsetbuttcap%
\pgfsetroundjoin%
\pgfsetlinewidth{1.003750pt}%
\definecolor{currentstroke}{rgb}{0.541176,0.380392,0.658824}%
\pgfsetstrokecolor{currentstroke}%
\pgfsetstrokeopacity{0.900000}%
\pgfsetdash{}{0pt}%
\pgfpathmoveto{\pgfqpoint{10.978299in}{5.405731in}}%
\pgfpathcurveto{\pgfqpoint{10.988553in}{5.405731in}}{\pgfqpoint{10.998389in}{5.409805in}}{\pgfqpoint{11.005639in}{5.417056in}}%
\pgfpathcurveto{\pgfqpoint{11.012890in}{5.424306in}}{\pgfqpoint{11.016964in}{5.434142in}}{\pgfqpoint{11.016964in}{5.444396in}}%
\pgfpathcurveto{\pgfqpoint{11.016964in}{5.454650in}}{\pgfqpoint{11.012890in}{5.464485in}}{\pgfqpoint{11.005639in}{5.471736in}}%
\pgfpathcurveto{\pgfqpoint{10.998389in}{5.478987in}}{\pgfqpoint{10.988553in}{5.483061in}}{\pgfqpoint{10.978299in}{5.483061in}}%
\pgfpathcurveto{\pgfqpoint{10.968045in}{5.483061in}}{\pgfqpoint{10.958209in}{5.478987in}}{\pgfqpoint{10.950959in}{5.471736in}}%
\pgfpathcurveto{\pgfqpoint{10.943708in}{5.464485in}}{\pgfqpoint{10.939634in}{5.454650in}}{\pgfqpoint{10.939634in}{5.444396in}}%
\pgfpathcurveto{\pgfqpoint{10.939634in}{5.434142in}}{\pgfqpoint{10.943708in}{5.424306in}}{\pgfqpoint{10.950959in}{5.417056in}}%
\pgfpathcurveto{\pgfqpoint{10.958209in}{5.409805in}}{\pgfqpoint{10.968045in}{5.405731in}}{\pgfqpoint{10.978299in}{5.405731in}}%
\pgfpathlineto{\pgfqpoint{10.978299in}{5.405731in}}%
\pgfpathclose%
\pgfusepath{stroke}%
\end{pgfscope}%
\begin{pgfscope}%
\pgfpathrectangle{\pgfqpoint{7.904091in}{2.928000in}}{\pgfqpoint{5.664000in}{5.664000in}}%
\pgfusepath{clip}%
\pgfsetbuttcap%
\pgfsetroundjoin%
\pgfsetlinewidth{1.003750pt}%
\definecolor{currentstroke}{rgb}{0.541176,0.380392,0.658824}%
\pgfsetstrokecolor{currentstroke}%
\pgfsetstrokeopacity{0.900000}%
\pgfsetdash{}{0pt}%
\pgfpathmoveto{\pgfqpoint{11.514076in}{6.358441in}}%
\pgfpathcurveto{\pgfqpoint{11.524330in}{6.358441in}}{\pgfqpoint{11.534166in}{6.362515in}}{\pgfqpoint{11.541417in}{6.369766in}}%
\pgfpathcurveto{\pgfqpoint{11.548667in}{6.377017in}}{\pgfqpoint{11.552741in}{6.386852in}}{\pgfqpoint{11.552741in}{6.397106in}}%
\pgfpathcurveto{\pgfqpoint{11.552741in}{6.407360in}}{\pgfqpoint{11.548667in}{6.417196in}}{\pgfqpoint{11.541417in}{6.424447in}}%
\pgfpathcurveto{\pgfqpoint{11.534166in}{6.431697in}}{\pgfqpoint{11.524330in}{6.435771in}}{\pgfqpoint{11.514076in}{6.435771in}}%
\pgfpathcurveto{\pgfqpoint{11.503822in}{6.435771in}}{\pgfqpoint{11.493987in}{6.431697in}}{\pgfqpoint{11.486736in}{6.424447in}}%
\pgfpathcurveto{\pgfqpoint{11.479485in}{6.417196in}}{\pgfqpoint{11.475411in}{6.407360in}}{\pgfqpoint{11.475411in}{6.397106in}}%
\pgfpathcurveto{\pgfqpoint{11.475411in}{6.386852in}}{\pgfqpoint{11.479485in}{6.377017in}}{\pgfqpoint{11.486736in}{6.369766in}}%
\pgfpathcurveto{\pgfqpoint{11.493987in}{6.362515in}}{\pgfqpoint{11.503822in}{6.358441in}}{\pgfqpoint{11.514076in}{6.358441in}}%
\pgfpathlineto{\pgfqpoint{11.514076in}{6.358441in}}%
\pgfpathclose%
\pgfusepath{stroke}%
\end{pgfscope}%
\begin{pgfscope}%
\pgfpathrectangle{\pgfqpoint{7.904091in}{2.928000in}}{\pgfqpoint{5.664000in}{5.664000in}}%
\pgfusepath{clip}%
\pgfsetbuttcap%
\pgfsetroundjoin%
\pgfsetlinewidth{1.003750pt}%
\definecolor{currentstroke}{rgb}{0.541176,0.380392,0.658824}%
\pgfsetstrokecolor{currentstroke}%
\pgfsetstrokeopacity{0.900000}%
\pgfsetdash{}{0pt}%
\pgfpathmoveto{\pgfqpoint{10.885791in}{4.989107in}}%
\pgfpathcurveto{\pgfqpoint{10.896045in}{4.989107in}}{\pgfqpoint{10.905881in}{4.993181in}}{\pgfqpoint{10.913131in}{5.000431in}}%
\pgfpathcurveto{\pgfqpoint{10.920382in}{5.007682in}}{\pgfqpoint{10.924456in}{5.017518in}}{\pgfqpoint{10.924456in}{5.027772in}}%
\pgfpathcurveto{\pgfqpoint{10.924456in}{5.038026in}}{\pgfqpoint{10.920382in}{5.047861in}}{\pgfqpoint{10.913131in}{5.055112in}}%
\pgfpathcurveto{\pgfqpoint{10.905881in}{5.062363in}}{\pgfqpoint{10.896045in}{5.066437in}}{\pgfqpoint{10.885791in}{5.066437in}}%
\pgfpathcurveto{\pgfqpoint{10.875537in}{5.066437in}}{\pgfqpoint{10.865701in}{5.062363in}}{\pgfqpoint{10.858451in}{5.055112in}}%
\pgfpathcurveto{\pgfqpoint{10.851200in}{5.047861in}}{\pgfqpoint{10.847126in}{5.038026in}}{\pgfqpoint{10.847126in}{5.027772in}}%
\pgfpathcurveto{\pgfqpoint{10.847126in}{5.017518in}}{\pgfqpoint{10.851200in}{5.007682in}}{\pgfqpoint{10.858451in}{5.000431in}}%
\pgfpathcurveto{\pgfqpoint{10.865701in}{4.993181in}}{\pgfqpoint{10.875537in}{4.989107in}}{\pgfqpoint{10.885791in}{4.989107in}}%
\pgfpathlineto{\pgfqpoint{10.885791in}{4.989107in}}%
\pgfpathclose%
\pgfusepath{stroke}%
\end{pgfscope}%
\begin{pgfscope}%
\pgfpathrectangle{\pgfqpoint{7.904091in}{2.928000in}}{\pgfqpoint{5.664000in}{5.664000in}}%
\pgfusepath{clip}%
\pgfsetbuttcap%
\pgfsetroundjoin%
\pgfsetlinewidth{1.003750pt}%
\definecolor{currentstroke}{rgb}{0.541176,0.380392,0.658824}%
\pgfsetstrokecolor{currentstroke}%
\pgfsetstrokeopacity{0.900000}%
\pgfsetdash{}{0pt}%
\pgfpathmoveto{\pgfqpoint{12.231386in}{7.309786in}}%
\pgfpathcurveto{\pgfqpoint{12.241640in}{7.309786in}}{\pgfqpoint{12.251475in}{7.313860in}}{\pgfqpoint{12.258726in}{7.321111in}}%
\pgfpathcurveto{\pgfqpoint{12.265977in}{7.328362in}}{\pgfqpoint{12.270051in}{7.338197in}}{\pgfqpoint{12.270051in}{7.348451in}}%
\pgfpathcurveto{\pgfqpoint{12.270051in}{7.358705in}}{\pgfqpoint{12.265977in}{7.368541in}}{\pgfqpoint{12.258726in}{7.375792in}}%
\pgfpathcurveto{\pgfqpoint{12.251475in}{7.383042in}}{\pgfqpoint{12.241640in}{7.387116in}}{\pgfqpoint{12.231386in}{7.387116in}}%
\pgfpathcurveto{\pgfqpoint{12.221131in}{7.387116in}}{\pgfqpoint{12.211296in}{7.383042in}}{\pgfqpoint{12.204045in}{7.375792in}}%
\pgfpathcurveto{\pgfqpoint{12.196794in}{7.368541in}}{\pgfqpoint{12.192720in}{7.358705in}}{\pgfqpoint{12.192720in}{7.348451in}}%
\pgfpathcurveto{\pgfqpoint{12.192720in}{7.338197in}}{\pgfqpoint{12.196794in}{7.328362in}}{\pgfqpoint{12.204045in}{7.321111in}}%
\pgfpathcurveto{\pgfqpoint{12.211296in}{7.313860in}}{\pgfqpoint{12.221131in}{7.309786in}}{\pgfqpoint{12.231386in}{7.309786in}}%
\pgfpathlineto{\pgfqpoint{12.231386in}{7.309786in}}%
\pgfpathclose%
\pgfusepath{stroke}%
\end{pgfscope}%
\begin{pgfscope}%
\pgfpathrectangle{\pgfqpoint{7.904091in}{2.928000in}}{\pgfqpoint{5.664000in}{5.664000in}}%
\pgfusepath{clip}%
\pgfsetbuttcap%
\pgfsetroundjoin%
\pgfsetlinewidth{1.003750pt}%
\definecolor{currentstroke}{rgb}{0.541176,0.380392,0.658824}%
\pgfsetstrokecolor{currentstroke}%
\pgfsetstrokeopacity{0.900000}%
\pgfsetdash{}{0pt}%
\pgfpathmoveto{\pgfqpoint{11.219755in}{5.980800in}}%
\pgfpathcurveto{\pgfqpoint{11.230009in}{5.980800in}}{\pgfqpoint{11.239845in}{5.984874in}}{\pgfqpoint{11.247095in}{5.992124in}}%
\pgfpathcurveto{\pgfqpoint{11.254346in}{5.999375in}}{\pgfqpoint{11.258420in}{6.009211in}}{\pgfqpoint{11.258420in}{6.019465in}}%
\pgfpathcurveto{\pgfqpoint{11.258420in}{6.029719in}}{\pgfqpoint{11.254346in}{6.039554in}}{\pgfqpoint{11.247095in}{6.046805in}}%
\pgfpathcurveto{\pgfqpoint{11.239845in}{6.054056in}}{\pgfqpoint{11.230009in}{6.058130in}}{\pgfqpoint{11.219755in}{6.058130in}}%
\pgfpathcurveto{\pgfqpoint{11.209501in}{6.058130in}}{\pgfqpoint{11.199665in}{6.054056in}}{\pgfqpoint{11.192415in}{6.046805in}}%
\pgfpathcurveto{\pgfqpoint{11.185164in}{6.039554in}}{\pgfqpoint{11.181090in}{6.029719in}}{\pgfqpoint{11.181090in}{6.019465in}}%
\pgfpathcurveto{\pgfqpoint{11.181090in}{6.009211in}}{\pgfqpoint{11.185164in}{5.999375in}}{\pgfqpoint{11.192415in}{5.992124in}}%
\pgfpathcurveto{\pgfqpoint{11.199665in}{5.984874in}}{\pgfqpoint{11.209501in}{5.980800in}}{\pgfqpoint{11.219755in}{5.980800in}}%
\pgfpathlineto{\pgfqpoint{11.219755in}{5.980800in}}%
\pgfpathclose%
\pgfusepath{stroke}%
\end{pgfscope}%
\begin{pgfscope}%
\pgfpathrectangle{\pgfqpoint{7.904091in}{2.928000in}}{\pgfqpoint{5.664000in}{5.664000in}}%
\pgfusepath{clip}%
\pgfsetbuttcap%
\pgfsetroundjoin%
\pgfsetlinewidth{1.003750pt}%
\definecolor{currentstroke}{rgb}{0.541176,0.380392,0.658824}%
\pgfsetstrokecolor{currentstroke}%
\pgfsetstrokeopacity{0.900000}%
\pgfsetdash{}{0pt}%
\pgfpathmoveto{\pgfqpoint{11.778546in}{6.486212in}}%
\pgfpathcurveto{\pgfqpoint{11.788800in}{6.486212in}}{\pgfqpoint{11.798636in}{6.490286in}}{\pgfqpoint{11.805886in}{6.497537in}}%
\pgfpathcurveto{\pgfqpoint{11.813137in}{6.504787in}}{\pgfqpoint{11.817211in}{6.514623in}}{\pgfqpoint{11.817211in}{6.524877in}}%
\pgfpathcurveto{\pgfqpoint{11.817211in}{6.535131in}}{\pgfqpoint{11.813137in}{6.544967in}}{\pgfqpoint{11.805886in}{6.552217in}}%
\pgfpathcurveto{\pgfqpoint{11.798636in}{6.559468in}}{\pgfqpoint{11.788800in}{6.563542in}}{\pgfqpoint{11.778546in}{6.563542in}}%
\pgfpathcurveto{\pgfqpoint{11.768292in}{6.563542in}}{\pgfqpoint{11.758456in}{6.559468in}}{\pgfqpoint{11.751206in}{6.552217in}}%
\pgfpathcurveto{\pgfqpoint{11.743955in}{6.544967in}}{\pgfqpoint{11.739881in}{6.535131in}}{\pgfqpoint{11.739881in}{6.524877in}}%
\pgfpathcurveto{\pgfqpoint{11.739881in}{6.514623in}}{\pgfqpoint{11.743955in}{6.504787in}}{\pgfqpoint{11.751206in}{6.497537in}}%
\pgfpathcurveto{\pgfqpoint{11.758456in}{6.490286in}}{\pgfqpoint{11.768292in}{6.486212in}}{\pgfqpoint{11.778546in}{6.486212in}}%
\pgfpathlineto{\pgfqpoint{11.778546in}{6.486212in}}%
\pgfpathclose%
\pgfusepath{stroke}%
\end{pgfscope}%
\begin{pgfscope}%
\pgfpathrectangle{\pgfqpoint{7.904091in}{2.928000in}}{\pgfqpoint{5.664000in}{5.664000in}}%
\pgfusepath{clip}%
\pgfsetbuttcap%
\pgfsetroundjoin%
\pgfsetlinewidth{1.003750pt}%
\definecolor{currentstroke}{rgb}{0.541176,0.380392,0.658824}%
\pgfsetstrokecolor{currentstroke}%
\pgfsetstrokeopacity{0.900000}%
\pgfsetdash{}{0pt}%
\pgfpathmoveto{\pgfqpoint{11.864848in}{7.237257in}}%
\pgfpathcurveto{\pgfqpoint{11.875102in}{7.237257in}}{\pgfqpoint{11.884937in}{7.241331in}}{\pgfqpoint{11.892188in}{7.248582in}}%
\pgfpathcurveto{\pgfqpoint{11.899439in}{7.255833in}}{\pgfqpoint{11.903513in}{7.265668in}}{\pgfqpoint{11.903513in}{7.275922in}}%
\pgfpathcurveto{\pgfqpoint{11.903513in}{7.286176in}}{\pgfqpoint{11.899439in}{7.296012in}}{\pgfqpoint{11.892188in}{7.303262in}}%
\pgfpathcurveto{\pgfqpoint{11.884937in}{7.310513in}}{\pgfqpoint{11.875102in}{7.314587in}}{\pgfqpoint{11.864848in}{7.314587in}}%
\pgfpathcurveto{\pgfqpoint{11.854594in}{7.314587in}}{\pgfqpoint{11.844758in}{7.310513in}}{\pgfqpoint{11.837507in}{7.303262in}}%
\pgfpathcurveto{\pgfqpoint{11.830257in}{7.296012in}}{\pgfqpoint{11.826183in}{7.286176in}}{\pgfqpoint{11.826183in}{7.275922in}}%
\pgfpathcurveto{\pgfqpoint{11.826183in}{7.265668in}}{\pgfqpoint{11.830257in}{7.255833in}}{\pgfqpoint{11.837507in}{7.248582in}}%
\pgfpathcurveto{\pgfqpoint{11.844758in}{7.241331in}}{\pgfqpoint{11.854594in}{7.237257in}}{\pgfqpoint{11.864848in}{7.237257in}}%
\pgfpathlineto{\pgfqpoint{11.864848in}{7.237257in}}%
\pgfpathclose%
\pgfusepath{stroke}%
\end{pgfscope}%
\begin{pgfscope}%
\pgfpathrectangle{\pgfqpoint{7.904091in}{2.928000in}}{\pgfqpoint{5.664000in}{5.664000in}}%
\pgfusepath{clip}%
\pgfsetbuttcap%
\pgfsetroundjoin%
\pgfsetlinewidth{1.003750pt}%
\definecolor{currentstroke}{rgb}{0.541176,0.380392,0.658824}%
\pgfsetstrokecolor{currentstroke}%
\pgfsetstrokeopacity{0.900000}%
\pgfsetdash{}{0pt}%
\pgfpathmoveto{\pgfqpoint{10.610035in}{5.621086in}}%
\pgfpathcurveto{\pgfqpoint{10.620290in}{5.621086in}}{\pgfqpoint{10.630125in}{5.625160in}}{\pgfqpoint{10.637376in}{5.632411in}}%
\pgfpathcurveto{\pgfqpoint{10.644627in}{5.639662in}}{\pgfqpoint{10.648701in}{5.649497in}}{\pgfqpoint{10.648701in}{5.659751in}}%
\pgfpathcurveto{\pgfqpoint{10.648701in}{5.670005in}}{\pgfqpoint{10.644627in}{5.679841in}}{\pgfqpoint{10.637376in}{5.687092in}}%
\pgfpathcurveto{\pgfqpoint{10.630125in}{5.694342in}}{\pgfqpoint{10.620290in}{5.698416in}}{\pgfqpoint{10.610035in}{5.698416in}}%
\pgfpathcurveto{\pgfqpoint{10.599781in}{5.698416in}}{\pgfqpoint{10.589946in}{5.694342in}}{\pgfqpoint{10.582695in}{5.687092in}}%
\pgfpathcurveto{\pgfqpoint{10.575444in}{5.679841in}}{\pgfqpoint{10.571370in}{5.670005in}}{\pgfqpoint{10.571370in}{5.659751in}}%
\pgfpathcurveto{\pgfqpoint{10.571370in}{5.649497in}}{\pgfqpoint{10.575444in}{5.639662in}}{\pgfqpoint{10.582695in}{5.632411in}}%
\pgfpathcurveto{\pgfqpoint{10.589946in}{5.625160in}}{\pgfqpoint{10.599781in}{5.621086in}}{\pgfqpoint{10.610035in}{5.621086in}}%
\pgfpathlineto{\pgfqpoint{10.610035in}{5.621086in}}%
\pgfpathclose%
\pgfusepath{stroke}%
\end{pgfscope}%
\begin{pgfscope}%
\pgfpathrectangle{\pgfqpoint{7.904091in}{2.928000in}}{\pgfqpoint{5.664000in}{5.664000in}}%
\pgfusepath{clip}%
\pgfsetbuttcap%
\pgfsetroundjoin%
\pgfsetlinewidth{1.003750pt}%
\definecolor{currentstroke}{rgb}{0.541176,0.380392,0.658824}%
\pgfsetstrokecolor{currentstroke}%
\pgfsetstrokeopacity{0.900000}%
\pgfsetdash{}{0pt}%
\pgfpathmoveto{\pgfqpoint{12.534083in}{7.116165in}}%
\pgfpathcurveto{\pgfqpoint{12.544337in}{7.116165in}}{\pgfqpoint{12.554173in}{7.120239in}}{\pgfqpoint{12.561424in}{7.127490in}}%
\pgfpathcurveto{\pgfqpoint{12.568674in}{7.134740in}}{\pgfqpoint{12.572748in}{7.144576in}}{\pgfqpoint{12.572748in}{7.154830in}}%
\pgfpathcurveto{\pgfqpoint{12.572748in}{7.165084in}}{\pgfqpoint{12.568674in}{7.174919in}}{\pgfqpoint{12.561424in}{7.182170in}}%
\pgfpathcurveto{\pgfqpoint{12.554173in}{7.189421in}}{\pgfqpoint{12.544337in}{7.193495in}}{\pgfqpoint{12.534083in}{7.193495in}}%
\pgfpathcurveto{\pgfqpoint{12.523829in}{7.193495in}}{\pgfqpoint{12.513994in}{7.189421in}}{\pgfqpoint{12.506743in}{7.182170in}}%
\pgfpathcurveto{\pgfqpoint{12.499492in}{7.174919in}}{\pgfqpoint{12.495418in}{7.165084in}}{\pgfqpoint{12.495418in}{7.154830in}}%
\pgfpathcurveto{\pgfqpoint{12.495418in}{7.144576in}}{\pgfqpoint{12.499492in}{7.134740in}}{\pgfqpoint{12.506743in}{7.127490in}}%
\pgfpathcurveto{\pgfqpoint{12.513994in}{7.120239in}}{\pgfqpoint{12.523829in}{7.116165in}}{\pgfqpoint{12.534083in}{7.116165in}}%
\pgfpathlineto{\pgfqpoint{12.534083in}{7.116165in}}%
\pgfpathclose%
\pgfusepath{stroke}%
\end{pgfscope}%
\begin{pgfscope}%
\pgfpathrectangle{\pgfqpoint{7.904091in}{2.928000in}}{\pgfqpoint{5.664000in}{5.664000in}}%
\pgfusepath{clip}%
\pgfsetbuttcap%
\pgfsetroundjoin%
\pgfsetlinewidth{1.003750pt}%
\definecolor{currentstroke}{rgb}{0.541176,0.380392,0.658824}%
\pgfsetstrokecolor{currentstroke}%
\pgfsetstrokeopacity{0.900000}%
\pgfsetdash{}{0pt}%
\pgfpathmoveto{\pgfqpoint{10.963469in}{6.454312in}}%
\pgfpathcurveto{\pgfqpoint{10.973723in}{6.454312in}}{\pgfqpoint{10.983559in}{6.458386in}}{\pgfqpoint{10.990809in}{6.465636in}}%
\pgfpathcurveto{\pgfqpoint{10.998060in}{6.472887in}}{\pgfqpoint{11.002134in}{6.482723in}}{\pgfqpoint{11.002134in}{6.492977in}}%
\pgfpathcurveto{\pgfqpoint{11.002134in}{6.503231in}}{\pgfqpoint{10.998060in}{6.513066in}}{\pgfqpoint{10.990809in}{6.520317in}}%
\pgfpathcurveto{\pgfqpoint{10.983559in}{6.527568in}}{\pgfqpoint{10.973723in}{6.531642in}}{\pgfqpoint{10.963469in}{6.531642in}}%
\pgfpathcurveto{\pgfqpoint{10.953215in}{6.531642in}}{\pgfqpoint{10.943379in}{6.527568in}}{\pgfqpoint{10.936129in}{6.520317in}}%
\pgfpathcurveto{\pgfqpoint{10.928878in}{6.513066in}}{\pgfqpoint{10.924804in}{6.503231in}}{\pgfqpoint{10.924804in}{6.492977in}}%
\pgfpathcurveto{\pgfqpoint{10.924804in}{6.482723in}}{\pgfqpoint{10.928878in}{6.472887in}}{\pgfqpoint{10.936129in}{6.465636in}}%
\pgfpathcurveto{\pgfqpoint{10.943379in}{6.458386in}}{\pgfqpoint{10.953215in}{6.454312in}}{\pgfqpoint{10.963469in}{6.454312in}}%
\pgfpathlineto{\pgfqpoint{10.963469in}{6.454312in}}%
\pgfpathclose%
\pgfusepath{stroke}%
\end{pgfscope}%
\begin{pgfscope}%
\pgfpathrectangle{\pgfqpoint{7.904091in}{2.928000in}}{\pgfqpoint{5.664000in}{5.664000in}}%
\pgfusepath{clip}%
\pgfsetbuttcap%
\pgfsetroundjoin%
\pgfsetlinewidth{1.003750pt}%
\definecolor{currentstroke}{rgb}{0.541176,0.380392,0.658824}%
\pgfsetstrokecolor{currentstroke}%
\pgfsetstrokeopacity{0.900000}%
\pgfsetdash{}{0pt}%
\pgfpathmoveto{\pgfqpoint{10.157174in}{5.237235in}}%
\pgfpathcurveto{\pgfqpoint{10.167428in}{5.237235in}}{\pgfqpoint{10.177263in}{5.241309in}}{\pgfqpoint{10.184514in}{5.248560in}}%
\pgfpathcurveto{\pgfqpoint{10.191765in}{5.255810in}}{\pgfqpoint{10.195839in}{5.265646in}}{\pgfqpoint{10.195839in}{5.275900in}}%
\pgfpathcurveto{\pgfqpoint{10.195839in}{5.286154in}}{\pgfqpoint{10.191765in}{5.295989in}}{\pgfqpoint{10.184514in}{5.303240in}}%
\pgfpathcurveto{\pgfqpoint{10.177263in}{5.310491in}}{\pgfqpoint{10.167428in}{5.314565in}}{\pgfqpoint{10.157174in}{5.314565in}}%
\pgfpathcurveto{\pgfqpoint{10.146920in}{5.314565in}}{\pgfqpoint{10.137084in}{5.310491in}}{\pgfqpoint{10.129834in}{5.303240in}}%
\pgfpathcurveto{\pgfqpoint{10.122583in}{5.295989in}}{\pgfqpoint{10.118509in}{5.286154in}}{\pgfqpoint{10.118509in}{5.275900in}}%
\pgfpathcurveto{\pgfqpoint{10.118509in}{5.265646in}}{\pgfqpoint{10.122583in}{5.255810in}}{\pgfqpoint{10.129834in}{5.248560in}}%
\pgfpathcurveto{\pgfqpoint{10.137084in}{5.241309in}}{\pgfqpoint{10.146920in}{5.237235in}}{\pgfqpoint{10.157174in}{5.237235in}}%
\pgfpathlineto{\pgfqpoint{10.157174in}{5.237235in}}%
\pgfpathclose%
\pgfusepath{stroke}%
\end{pgfscope}%
\begin{pgfscope}%
\pgfpathrectangle{\pgfqpoint{7.904091in}{2.928000in}}{\pgfqpoint{5.664000in}{5.664000in}}%
\pgfusepath{clip}%
\pgfsetbuttcap%
\pgfsetroundjoin%
\pgfsetlinewidth{1.003750pt}%
\definecolor{currentstroke}{rgb}{0.541176,0.380392,0.658824}%
\pgfsetstrokecolor{currentstroke}%
\pgfsetstrokeopacity{0.900000}%
\pgfsetdash{}{0pt}%
\pgfpathmoveto{\pgfqpoint{11.332139in}{5.993923in}}%
\pgfpathcurveto{\pgfqpoint{11.342393in}{5.993923in}}{\pgfqpoint{11.352229in}{5.997997in}}{\pgfqpoint{11.359479in}{6.005247in}}%
\pgfpathcurveto{\pgfqpoint{11.366730in}{6.012498in}}{\pgfqpoint{11.370804in}{6.022334in}}{\pgfqpoint{11.370804in}{6.032588in}}%
\pgfpathcurveto{\pgfqpoint{11.370804in}{6.042842in}}{\pgfqpoint{11.366730in}{6.052677in}}{\pgfqpoint{11.359479in}{6.059928in}}%
\pgfpathcurveto{\pgfqpoint{11.352229in}{6.067179in}}{\pgfqpoint{11.342393in}{6.071253in}}{\pgfqpoint{11.332139in}{6.071253in}}%
\pgfpathcurveto{\pgfqpoint{11.321885in}{6.071253in}}{\pgfqpoint{11.312049in}{6.067179in}}{\pgfqpoint{11.304799in}{6.059928in}}%
\pgfpathcurveto{\pgfqpoint{11.297548in}{6.052677in}}{\pgfqpoint{11.293474in}{6.042842in}}{\pgfqpoint{11.293474in}{6.032588in}}%
\pgfpathcurveto{\pgfqpoint{11.293474in}{6.022334in}}{\pgfqpoint{11.297548in}{6.012498in}}{\pgfqpoint{11.304799in}{6.005247in}}%
\pgfpathcurveto{\pgfqpoint{11.312049in}{5.997997in}}{\pgfqpoint{11.321885in}{5.993923in}}{\pgfqpoint{11.332139in}{5.993923in}}%
\pgfpathlineto{\pgfqpoint{11.332139in}{5.993923in}}%
\pgfpathclose%
\pgfusepath{stroke}%
\end{pgfscope}%
\begin{pgfscope}%
\pgfpathrectangle{\pgfqpoint{7.904091in}{2.928000in}}{\pgfqpoint{5.664000in}{5.664000in}}%
\pgfusepath{clip}%
\pgfsetbuttcap%
\pgfsetroundjoin%
\pgfsetlinewidth{1.003750pt}%
\definecolor{currentstroke}{rgb}{0.541176,0.380392,0.658824}%
\pgfsetstrokecolor{currentstroke}%
\pgfsetstrokeopacity{0.900000}%
\pgfsetdash{}{0pt}%
\pgfpathmoveto{\pgfqpoint{11.281597in}{6.648760in}}%
\pgfpathcurveto{\pgfqpoint{11.291851in}{6.648760in}}{\pgfqpoint{11.301687in}{6.652834in}}{\pgfqpoint{11.308937in}{6.660085in}}%
\pgfpathcurveto{\pgfqpoint{11.316188in}{6.667335in}}{\pgfqpoint{11.320262in}{6.677171in}}{\pgfqpoint{11.320262in}{6.687425in}}%
\pgfpathcurveto{\pgfqpoint{11.320262in}{6.697679in}}{\pgfqpoint{11.316188in}{6.707514in}}{\pgfqpoint{11.308937in}{6.714765in}}%
\pgfpathcurveto{\pgfqpoint{11.301687in}{6.722016in}}{\pgfqpoint{11.291851in}{6.726090in}}{\pgfqpoint{11.281597in}{6.726090in}}%
\pgfpathcurveto{\pgfqpoint{11.271343in}{6.726090in}}{\pgfqpoint{11.261507in}{6.722016in}}{\pgfqpoint{11.254257in}{6.714765in}}%
\pgfpathcurveto{\pgfqpoint{11.247006in}{6.707514in}}{\pgfqpoint{11.242932in}{6.697679in}}{\pgfqpoint{11.242932in}{6.687425in}}%
\pgfpathcurveto{\pgfqpoint{11.242932in}{6.677171in}}{\pgfqpoint{11.247006in}{6.667335in}}{\pgfqpoint{11.254257in}{6.660085in}}%
\pgfpathcurveto{\pgfqpoint{11.261507in}{6.652834in}}{\pgfqpoint{11.271343in}{6.648760in}}{\pgfqpoint{11.281597in}{6.648760in}}%
\pgfpathlineto{\pgfqpoint{11.281597in}{6.648760in}}%
\pgfpathclose%
\pgfusepath{stroke}%
\end{pgfscope}%
\begin{pgfscope}%
\pgfpathrectangle{\pgfqpoint{7.904091in}{2.928000in}}{\pgfqpoint{5.664000in}{5.664000in}}%
\pgfusepath{clip}%
\pgfsetbuttcap%
\pgfsetroundjoin%
\pgfsetlinewidth{1.003750pt}%
\definecolor{currentstroke}{rgb}{0.541176,0.380392,0.658824}%
\pgfsetstrokecolor{currentstroke}%
\pgfsetstrokeopacity{0.900000}%
\pgfsetdash{}{0pt}%
\pgfpathmoveto{\pgfqpoint{11.334728in}{5.849021in}}%
\pgfpathcurveto{\pgfqpoint{11.344982in}{5.849021in}}{\pgfqpoint{11.354817in}{5.853095in}}{\pgfqpoint{11.362068in}{5.860346in}}%
\pgfpathcurveto{\pgfqpoint{11.369319in}{5.867597in}}{\pgfqpoint{11.373393in}{5.877432in}}{\pgfqpoint{11.373393in}{5.887686in}}%
\pgfpathcurveto{\pgfqpoint{11.373393in}{5.897941in}}{\pgfqpoint{11.369319in}{5.907776in}}{\pgfqpoint{11.362068in}{5.915027in}}%
\pgfpathcurveto{\pgfqpoint{11.354817in}{5.922277in}}{\pgfqpoint{11.344982in}{5.926351in}}{\pgfqpoint{11.334728in}{5.926351in}}%
\pgfpathcurveto{\pgfqpoint{11.324474in}{5.926351in}}{\pgfqpoint{11.314638in}{5.922277in}}{\pgfqpoint{11.307388in}{5.915027in}}%
\pgfpathcurveto{\pgfqpoint{11.300137in}{5.907776in}}{\pgfqpoint{11.296063in}{5.897941in}}{\pgfqpoint{11.296063in}{5.887686in}}%
\pgfpathcurveto{\pgfqpoint{11.296063in}{5.877432in}}{\pgfqpoint{11.300137in}{5.867597in}}{\pgfqpoint{11.307388in}{5.860346in}}%
\pgfpathcurveto{\pgfqpoint{11.314638in}{5.853095in}}{\pgfqpoint{11.324474in}{5.849021in}}{\pgfqpoint{11.334728in}{5.849021in}}%
\pgfpathlineto{\pgfqpoint{11.334728in}{5.849021in}}%
\pgfpathclose%
\pgfusepath{stroke}%
\end{pgfscope}%
\begin{pgfscope}%
\pgfpathrectangle{\pgfqpoint{7.904091in}{2.928000in}}{\pgfqpoint{5.664000in}{5.664000in}}%
\pgfusepath{clip}%
\pgfsetbuttcap%
\pgfsetroundjoin%
\pgfsetlinewidth{1.003750pt}%
\definecolor{currentstroke}{rgb}{0.541176,0.380392,0.658824}%
\pgfsetstrokecolor{currentstroke}%
\pgfsetstrokeopacity{0.900000}%
\pgfsetdash{}{0pt}%
\pgfpathmoveto{\pgfqpoint{11.314154in}{5.978579in}}%
\pgfpathcurveto{\pgfqpoint{11.324408in}{5.978579in}}{\pgfqpoint{11.334243in}{5.982653in}}{\pgfqpoint{11.341494in}{5.989904in}}%
\pgfpathcurveto{\pgfqpoint{11.348745in}{5.997155in}}{\pgfqpoint{11.352819in}{6.006990in}}{\pgfqpoint{11.352819in}{6.017244in}}%
\pgfpathcurveto{\pgfqpoint{11.352819in}{6.027498in}}{\pgfqpoint{11.348745in}{6.037334in}}{\pgfqpoint{11.341494in}{6.044585in}}%
\pgfpathcurveto{\pgfqpoint{11.334243in}{6.051835in}}{\pgfqpoint{11.324408in}{6.055909in}}{\pgfqpoint{11.314154in}{6.055909in}}%
\pgfpathcurveto{\pgfqpoint{11.303899in}{6.055909in}}{\pgfqpoint{11.294064in}{6.051835in}}{\pgfqpoint{11.286813in}{6.044585in}}%
\pgfpathcurveto{\pgfqpoint{11.279562in}{6.037334in}}{\pgfqpoint{11.275488in}{6.027498in}}{\pgfqpoint{11.275488in}{6.017244in}}%
\pgfpathcurveto{\pgfqpoint{11.275488in}{6.006990in}}{\pgfqpoint{11.279562in}{5.997155in}}{\pgfqpoint{11.286813in}{5.989904in}}%
\pgfpathcurveto{\pgfqpoint{11.294064in}{5.982653in}}{\pgfqpoint{11.303899in}{5.978579in}}{\pgfqpoint{11.314154in}{5.978579in}}%
\pgfpathlineto{\pgfqpoint{11.314154in}{5.978579in}}%
\pgfpathclose%
\pgfusepath{stroke}%
\end{pgfscope}%
\begin{pgfscope}%
\pgfpathrectangle{\pgfqpoint{7.904091in}{2.928000in}}{\pgfqpoint{5.664000in}{5.664000in}}%
\pgfusepath{clip}%
\pgfsetbuttcap%
\pgfsetroundjoin%
\pgfsetlinewidth{1.003750pt}%
\definecolor{currentstroke}{rgb}{0.541176,0.380392,0.658824}%
\pgfsetstrokecolor{currentstroke}%
\pgfsetstrokeopacity{0.900000}%
\pgfsetdash{}{0pt}%
\pgfpathmoveto{\pgfqpoint{11.591549in}{6.978810in}}%
\pgfpathcurveto{\pgfqpoint{11.601803in}{6.978810in}}{\pgfqpoint{11.611638in}{6.982884in}}{\pgfqpoint{11.618889in}{6.990135in}}%
\pgfpathcurveto{\pgfqpoint{11.626140in}{6.997386in}}{\pgfqpoint{11.630214in}{7.007221in}}{\pgfqpoint{11.630214in}{7.017475in}}%
\pgfpathcurveto{\pgfqpoint{11.630214in}{7.027729in}}{\pgfqpoint{11.626140in}{7.037565in}}{\pgfqpoint{11.618889in}{7.044815in}}%
\pgfpathcurveto{\pgfqpoint{11.611638in}{7.052066in}}{\pgfqpoint{11.601803in}{7.056140in}}{\pgfqpoint{11.591549in}{7.056140in}}%
\pgfpathcurveto{\pgfqpoint{11.581295in}{7.056140in}}{\pgfqpoint{11.571459in}{7.052066in}}{\pgfqpoint{11.564208in}{7.044815in}}%
\pgfpathcurveto{\pgfqpoint{11.556958in}{7.037565in}}{\pgfqpoint{11.552884in}{7.027729in}}{\pgfqpoint{11.552884in}{7.017475in}}%
\pgfpathcurveto{\pgfqpoint{11.552884in}{7.007221in}}{\pgfqpoint{11.556958in}{6.997386in}}{\pgfqpoint{11.564208in}{6.990135in}}%
\pgfpathcurveto{\pgfqpoint{11.571459in}{6.982884in}}{\pgfqpoint{11.581295in}{6.978810in}}{\pgfqpoint{11.591549in}{6.978810in}}%
\pgfpathlineto{\pgfqpoint{11.591549in}{6.978810in}}%
\pgfpathclose%
\pgfusepath{stroke}%
\end{pgfscope}%
\begin{pgfscope}%
\pgfpathrectangle{\pgfqpoint{7.904091in}{2.928000in}}{\pgfqpoint{5.664000in}{5.664000in}}%
\pgfusepath{clip}%
\pgfsetbuttcap%
\pgfsetroundjoin%
\pgfsetlinewidth{1.003750pt}%
\definecolor{currentstroke}{rgb}{0.541176,0.380392,0.658824}%
\pgfsetstrokecolor{currentstroke}%
\pgfsetstrokeopacity{0.900000}%
\pgfsetdash{}{0pt}%
\pgfpathmoveto{\pgfqpoint{11.321020in}{6.779867in}}%
\pgfpathcurveto{\pgfqpoint{11.331274in}{6.779867in}}{\pgfqpoint{11.341109in}{6.783941in}}{\pgfqpoint{11.348360in}{6.791192in}}%
\pgfpathcurveto{\pgfqpoint{11.355611in}{6.798443in}}{\pgfqpoint{11.359685in}{6.808278in}}{\pgfqpoint{11.359685in}{6.818532in}}%
\pgfpathcurveto{\pgfqpoint{11.359685in}{6.828786in}}{\pgfqpoint{11.355611in}{6.838622in}}{\pgfqpoint{11.348360in}{6.845873in}}%
\pgfpathcurveto{\pgfqpoint{11.341109in}{6.853123in}}{\pgfqpoint{11.331274in}{6.857197in}}{\pgfqpoint{11.321020in}{6.857197in}}%
\pgfpathcurveto{\pgfqpoint{11.310766in}{6.857197in}}{\pgfqpoint{11.300930in}{6.853123in}}{\pgfqpoint{11.293679in}{6.845873in}}%
\pgfpathcurveto{\pgfqpoint{11.286429in}{6.838622in}}{\pgfqpoint{11.282355in}{6.828786in}}{\pgfqpoint{11.282355in}{6.818532in}}%
\pgfpathcurveto{\pgfqpoint{11.282355in}{6.808278in}}{\pgfqpoint{11.286429in}{6.798443in}}{\pgfqpoint{11.293679in}{6.791192in}}%
\pgfpathcurveto{\pgfqpoint{11.300930in}{6.783941in}}{\pgfqpoint{11.310766in}{6.779867in}}{\pgfqpoint{11.321020in}{6.779867in}}%
\pgfpathlineto{\pgfqpoint{11.321020in}{6.779867in}}%
\pgfpathclose%
\pgfusepath{stroke}%
\end{pgfscope}%
\begin{pgfscope}%
\pgfpathrectangle{\pgfqpoint{7.904091in}{2.928000in}}{\pgfqpoint{5.664000in}{5.664000in}}%
\pgfusepath{clip}%
\pgfsetbuttcap%
\pgfsetroundjoin%
\pgfsetlinewidth{1.003750pt}%
\definecolor{currentstroke}{rgb}{0.541176,0.380392,0.658824}%
\pgfsetstrokecolor{currentstroke}%
\pgfsetstrokeopacity{0.900000}%
\pgfsetdash{}{0pt}%
\pgfpathmoveto{\pgfqpoint{11.765466in}{7.383840in}}%
\pgfpathcurveto{\pgfqpoint{11.775720in}{7.383840in}}{\pgfqpoint{11.785555in}{7.387914in}}{\pgfqpoint{11.792806in}{7.395165in}}%
\pgfpathcurveto{\pgfqpoint{11.800057in}{7.402415in}}{\pgfqpoint{11.804131in}{7.412251in}}{\pgfqpoint{11.804131in}{7.422505in}}%
\pgfpathcurveto{\pgfqpoint{11.804131in}{7.432759in}}{\pgfqpoint{11.800057in}{7.442595in}}{\pgfqpoint{11.792806in}{7.449845in}}%
\pgfpathcurveto{\pgfqpoint{11.785555in}{7.457096in}}{\pgfqpoint{11.775720in}{7.461170in}}{\pgfqpoint{11.765466in}{7.461170in}}%
\pgfpathcurveto{\pgfqpoint{11.755212in}{7.461170in}}{\pgfqpoint{11.745376in}{7.457096in}}{\pgfqpoint{11.738125in}{7.449845in}}%
\pgfpathcurveto{\pgfqpoint{11.730875in}{7.442595in}}{\pgfqpoint{11.726801in}{7.432759in}}{\pgfqpoint{11.726801in}{7.422505in}}%
\pgfpathcurveto{\pgfqpoint{11.726801in}{7.412251in}}{\pgfqpoint{11.730875in}{7.402415in}}{\pgfqpoint{11.738125in}{7.395165in}}%
\pgfpathcurveto{\pgfqpoint{11.745376in}{7.387914in}}{\pgfqpoint{11.755212in}{7.383840in}}{\pgfqpoint{11.765466in}{7.383840in}}%
\pgfpathlineto{\pgfqpoint{11.765466in}{7.383840in}}%
\pgfpathclose%
\pgfusepath{stroke}%
\end{pgfscope}%
\begin{pgfscope}%
\pgfpathrectangle{\pgfqpoint{7.904091in}{2.928000in}}{\pgfqpoint{5.664000in}{5.664000in}}%
\pgfusepath{clip}%
\pgfsetbuttcap%
\pgfsetroundjoin%
\pgfsetlinewidth{1.003750pt}%
\definecolor{currentstroke}{rgb}{0.541176,0.380392,0.658824}%
\pgfsetstrokecolor{currentstroke}%
\pgfsetstrokeopacity{0.900000}%
\pgfsetdash{}{0pt}%
\pgfpathmoveto{\pgfqpoint{10.331058in}{5.002465in}}%
\pgfpathcurveto{\pgfqpoint{10.341312in}{5.002465in}}{\pgfqpoint{10.351148in}{5.006539in}}{\pgfqpoint{10.358398in}{5.013790in}}%
\pgfpathcurveto{\pgfqpoint{10.365649in}{5.021040in}}{\pgfqpoint{10.369723in}{5.030876in}}{\pgfqpoint{10.369723in}{5.041130in}}%
\pgfpathcurveto{\pgfqpoint{10.369723in}{5.051384in}}{\pgfqpoint{10.365649in}{5.061219in}}{\pgfqpoint{10.358398in}{5.068470in}}%
\pgfpathcurveto{\pgfqpoint{10.351148in}{5.075721in}}{\pgfqpoint{10.341312in}{5.079795in}}{\pgfqpoint{10.331058in}{5.079795in}}%
\pgfpathcurveto{\pgfqpoint{10.320804in}{5.079795in}}{\pgfqpoint{10.310968in}{5.075721in}}{\pgfqpoint{10.303718in}{5.068470in}}%
\pgfpathcurveto{\pgfqpoint{10.296467in}{5.061219in}}{\pgfqpoint{10.292393in}{5.051384in}}{\pgfqpoint{10.292393in}{5.041130in}}%
\pgfpathcurveto{\pgfqpoint{10.292393in}{5.030876in}}{\pgfqpoint{10.296467in}{5.021040in}}{\pgfqpoint{10.303718in}{5.013790in}}%
\pgfpathcurveto{\pgfqpoint{10.310968in}{5.006539in}}{\pgfqpoint{10.320804in}{5.002465in}}{\pgfqpoint{10.331058in}{5.002465in}}%
\pgfpathlineto{\pgfqpoint{10.331058in}{5.002465in}}%
\pgfpathclose%
\pgfusepath{stroke}%
\end{pgfscope}%
\begin{pgfscope}%
\pgfpathrectangle{\pgfqpoint{7.904091in}{2.928000in}}{\pgfqpoint{5.664000in}{5.664000in}}%
\pgfusepath{clip}%
\pgfsetbuttcap%
\pgfsetroundjoin%
\pgfsetlinewidth{1.003750pt}%
\definecolor{currentstroke}{rgb}{0.541176,0.380392,0.658824}%
\pgfsetstrokecolor{currentstroke}%
\pgfsetstrokeopacity{0.900000}%
\pgfsetdash{}{0pt}%
\pgfpathmoveto{\pgfqpoint{11.702131in}{6.516364in}}%
\pgfpathcurveto{\pgfqpoint{11.712385in}{6.516364in}}{\pgfqpoint{11.722221in}{6.520438in}}{\pgfqpoint{11.729472in}{6.527688in}}%
\pgfpathcurveto{\pgfqpoint{11.736722in}{6.534939in}}{\pgfqpoint{11.740796in}{6.544775in}}{\pgfqpoint{11.740796in}{6.555029in}}%
\pgfpathcurveto{\pgfqpoint{11.740796in}{6.565283in}}{\pgfqpoint{11.736722in}{6.575118in}}{\pgfqpoint{11.729472in}{6.582369in}}%
\pgfpathcurveto{\pgfqpoint{11.722221in}{6.589620in}}{\pgfqpoint{11.712385in}{6.593694in}}{\pgfqpoint{11.702131in}{6.593694in}}%
\pgfpathcurveto{\pgfqpoint{11.691877in}{6.593694in}}{\pgfqpoint{11.682042in}{6.589620in}}{\pgfqpoint{11.674791in}{6.582369in}}%
\pgfpathcurveto{\pgfqpoint{11.667540in}{6.575118in}}{\pgfqpoint{11.663466in}{6.565283in}}{\pgfqpoint{11.663466in}{6.555029in}}%
\pgfpathcurveto{\pgfqpoint{11.663466in}{6.544775in}}{\pgfqpoint{11.667540in}{6.534939in}}{\pgfqpoint{11.674791in}{6.527688in}}%
\pgfpathcurveto{\pgfqpoint{11.682042in}{6.520438in}}{\pgfqpoint{11.691877in}{6.516364in}}{\pgfqpoint{11.702131in}{6.516364in}}%
\pgfpathlineto{\pgfqpoint{11.702131in}{6.516364in}}%
\pgfpathclose%
\pgfusepath{stroke}%
\end{pgfscope}%
\begin{pgfscope}%
\pgfsetbuttcap%
\pgfsetroundjoin%
\pgfsetlinewidth{1.003750pt}%
\definecolor{currentstroke}{rgb}{0.541176,0.380392,0.658824}%
\pgfsetstrokecolor{currentstroke}%
\pgfsetstrokeopacity{0.900000}%
\pgfsetdash{}{0pt}%
\pgfpathmoveto{\pgfqpoint{13.568091in}{8.532262in}}%
\pgfpathlineto{\pgfqpoint{13.627829in}{8.592000in}}%
\pgfpathlineto{\pgfqpoint{13.568091in}{8.651738in}}%
\pgfpathlineto{\pgfqpoint{13.508353in}{8.592000in}}%
\pgfpathlineto{\pgfqpoint{13.568091in}{8.532262in}}%
\pgfpathclose%
\pgfusepath{stroke}%
\end{pgfscope}%
\begin{pgfscope}%
\pgfpathrectangle{\pgfqpoint{7.904091in}{2.928000in}}{\pgfqpoint{5.664000in}{5.664000in}}%
\pgfusepath{clip}%
\pgfsetbuttcap%
\pgfsetroundjoin%
\pgfsetlinewidth{1.003750pt}%
\definecolor{currentstroke}{rgb}{0.321569,0.321569,0.321569}%
\pgfsetstrokecolor{currentstroke}%
\pgfsetstrokeopacity{0.900000}%
\pgfsetdash{}{0pt}%
\pgfpathmoveto{\pgfqpoint{9.813675in}{4.636051in}}%
\pgfpathcurveto{\pgfqpoint{9.823929in}{4.636051in}}{\pgfqpoint{9.833765in}{4.640125in}}{\pgfqpoint{9.841016in}{4.647376in}}%
\pgfpathcurveto{\pgfqpoint{9.848266in}{4.654626in}}{\pgfqpoint{9.852340in}{4.664462in}}{\pgfqpoint{9.852340in}{4.674716in}}%
\pgfpathcurveto{\pgfqpoint{9.852340in}{4.684970in}}{\pgfqpoint{9.848266in}{4.694805in}}{\pgfqpoint{9.841016in}{4.702056in}}%
\pgfpathcurveto{\pgfqpoint{9.833765in}{4.709307in}}{\pgfqpoint{9.823929in}{4.713381in}}{\pgfqpoint{9.813675in}{4.713381in}}%
\pgfpathcurveto{\pgfqpoint{9.803421in}{4.713381in}}{\pgfqpoint{9.793586in}{4.709307in}}{\pgfqpoint{9.786335in}{4.702056in}}%
\pgfpathcurveto{\pgfqpoint{9.779084in}{4.694805in}}{\pgfqpoint{9.775010in}{4.684970in}}{\pgfqpoint{9.775010in}{4.674716in}}%
\pgfpathcurveto{\pgfqpoint{9.775010in}{4.664462in}}{\pgfqpoint{9.779084in}{4.654626in}}{\pgfqpoint{9.786335in}{4.647376in}}%
\pgfpathcurveto{\pgfqpoint{9.793586in}{4.640125in}}{\pgfqpoint{9.803421in}{4.636051in}}{\pgfqpoint{9.813675in}{4.636051in}}%
\pgfpathlineto{\pgfqpoint{9.813675in}{4.636051in}}%
\pgfpathclose%
\pgfusepath{stroke}%
\end{pgfscope}%
\begin{pgfscope}%
\pgfpathrectangle{\pgfqpoint{7.904091in}{2.928000in}}{\pgfqpoint{5.664000in}{5.664000in}}%
\pgfusepath{clip}%
\pgfsetbuttcap%
\pgfsetroundjoin%
\pgfsetlinewidth{1.003750pt}%
\definecolor{currentstroke}{rgb}{0.321569,0.321569,0.321569}%
\pgfsetstrokecolor{currentstroke}%
\pgfsetstrokeopacity{0.900000}%
\pgfsetdash{}{0pt}%
\pgfpathmoveto{\pgfqpoint{10.010366in}{4.923966in}}%
\pgfpathcurveto{\pgfqpoint{10.020620in}{4.923966in}}{\pgfqpoint{10.030455in}{4.928040in}}{\pgfqpoint{10.037706in}{4.935291in}}%
\pgfpathcurveto{\pgfqpoint{10.044957in}{4.942542in}}{\pgfqpoint{10.049031in}{4.952377in}}{\pgfqpoint{10.049031in}{4.962631in}}%
\pgfpathcurveto{\pgfqpoint{10.049031in}{4.972886in}}{\pgfqpoint{10.044957in}{4.982721in}}{\pgfqpoint{10.037706in}{4.989972in}}%
\pgfpathcurveto{\pgfqpoint{10.030455in}{4.997223in}}{\pgfqpoint{10.020620in}{5.001296in}}{\pgfqpoint{10.010366in}{5.001296in}}%
\pgfpathcurveto{\pgfqpoint{10.000112in}{5.001296in}}{\pgfqpoint{9.990276in}{4.997223in}}{\pgfqpoint{9.983025in}{4.989972in}}%
\pgfpathcurveto{\pgfqpoint{9.975775in}{4.982721in}}{\pgfqpoint{9.971701in}{4.972886in}}{\pgfqpoint{9.971701in}{4.962631in}}%
\pgfpathcurveto{\pgfqpoint{9.971701in}{4.952377in}}{\pgfqpoint{9.975775in}{4.942542in}}{\pgfqpoint{9.983025in}{4.935291in}}%
\pgfpathcurveto{\pgfqpoint{9.990276in}{4.928040in}}{\pgfqpoint{10.000112in}{4.923966in}}{\pgfqpoint{10.010366in}{4.923966in}}%
\pgfpathlineto{\pgfqpoint{10.010366in}{4.923966in}}%
\pgfpathclose%
\pgfusepath{stroke}%
\end{pgfscope}%
\begin{pgfscope}%
\pgfpathrectangle{\pgfqpoint{7.904091in}{2.928000in}}{\pgfqpoint{5.664000in}{5.664000in}}%
\pgfusepath{clip}%
\pgfsetbuttcap%
\pgfsetroundjoin%
\pgfsetlinewidth{1.003750pt}%
\definecolor{currentstroke}{rgb}{0.321569,0.321569,0.321569}%
\pgfsetstrokecolor{currentstroke}%
\pgfsetstrokeopacity{0.900000}%
\pgfsetdash{}{0pt}%
\pgfpathmoveto{\pgfqpoint{9.234809in}{3.957465in}}%
\pgfpathcurveto{\pgfqpoint{9.245063in}{3.957465in}}{\pgfqpoint{9.254899in}{3.961539in}}{\pgfqpoint{9.262149in}{3.968790in}}%
\pgfpathcurveto{\pgfqpoint{9.269400in}{3.976040in}}{\pgfqpoint{9.273474in}{3.985876in}}{\pgfqpoint{9.273474in}{3.996130in}}%
\pgfpathcurveto{\pgfqpoint{9.273474in}{4.006384in}}{\pgfqpoint{9.269400in}{4.016219in}}{\pgfqpoint{9.262149in}{4.023470in}}%
\pgfpathcurveto{\pgfqpoint{9.254899in}{4.030721in}}{\pgfqpoint{9.245063in}{4.034795in}}{\pgfqpoint{9.234809in}{4.034795in}}%
\pgfpathcurveto{\pgfqpoint{9.224555in}{4.034795in}}{\pgfqpoint{9.214719in}{4.030721in}}{\pgfqpoint{9.207469in}{4.023470in}}%
\pgfpathcurveto{\pgfqpoint{9.200218in}{4.016219in}}{\pgfqpoint{9.196144in}{4.006384in}}{\pgfqpoint{9.196144in}{3.996130in}}%
\pgfpathcurveto{\pgfqpoint{9.196144in}{3.985876in}}{\pgfqpoint{9.200218in}{3.976040in}}{\pgfqpoint{9.207469in}{3.968790in}}%
\pgfpathcurveto{\pgfqpoint{9.214719in}{3.961539in}}{\pgfqpoint{9.224555in}{3.957465in}}{\pgfqpoint{9.234809in}{3.957465in}}%
\pgfpathlineto{\pgfqpoint{9.234809in}{3.957465in}}%
\pgfpathclose%
\pgfusepath{stroke}%
\end{pgfscope}%
\begin{pgfscope}%
\pgfpathrectangle{\pgfqpoint{7.904091in}{2.928000in}}{\pgfqpoint{5.664000in}{5.664000in}}%
\pgfusepath{clip}%
\pgfsetbuttcap%
\pgfsetroundjoin%
\pgfsetlinewidth{1.003750pt}%
\definecolor{currentstroke}{rgb}{0.321569,0.321569,0.321569}%
\pgfsetstrokecolor{currentstroke}%
\pgfsetstrokeopacity{0.900000}%
\pgfsetdash{}{0pt}%
\pgfpathmoveto{\pgfqpoint{9.370673in}{4.402648in}}%
\pgfpathcurveto{\pgfqpoint{9.380927in}{4.402648in}}{\pgfqpoint{9.390762in}{4.406722in}}{\pgfqpoint{9.398013in}{4.413973in}}%
\pgfpathcurveto{\pgfqpoint{9.405264in}{4.421223in}}{\pgfqpoint{9.409338in}{4.431059in}}{\pgfqpoint{9.409338in}{4.441313in}}%
\pgfpathcurveto{\pgfqpoint{9.409338in}{4.451567in}}{\pgfqpoint{9.405264in}{4.461402in}}{\pgfqpoint{9.398013in}{4.468653in}}%
\pgfpathcurveto{\pgfqpoint{9.390762in}{4.475904in}}{\pgfqpoint{9.380927in}{4.479978in}}{\pgfqpoint{9.370673in}{4.479978in}}%
\pgfpathcurveto{\pgfqpoint{9.360419in}{4.479978in}}{\pgfqpoint{9.350583in}{4.475904in}}{\pgfqpoint{9.343332in}{4.468653in}}%
\pgfpathcurveto{\pgfqpoint{9.336082in}{4.461402in}}{\pgfqpoint{9.332008in}{4.451567in}}{\pgfqpoint{9.332008in}{4.441313in}}%
\pgfpathcurveto{\pgfqpoint{9.332008in}{4.431059in}}{\pgfqpoint{9.336082in}{4.421223in}}{\pgfqpoint{9.343332in}{4.413973in}}%
\pgfpathcurveto{\pgfqpoint{9.350583in}{4.406722in}}{\pgfqpoint{9.360419in}{4.402648in}}{\pgfqpoint{9.370673in}{4.402648in}}%
\pgfpathlineto{\pgfqpoint{9.370673in}{4.402648in}}%
\pgfpathclose%
\pgfusepath{stroke}%
\end{pgfscope}%
\begin{pgfscope}%
\pgfpathrectangle{\pgfqpoint{7.904091in}{2.928000in}}{\pgfqpoint{5.664000in}{5.664000in}}%
\pgfusepath{clip}%
\pgfsetbuttcap%
\pgfsetroundjoin%
\pgfsetlinewidth{1.003750pt}%
\definecolor{currentstroke}{rgb}{0.321569,0.321569,0.321569}%
\pgfsetstrokecolor{currentstroke}%
\pgfsetstrokeopacity{0.900000}%
\pgfsetdash{}{0pt}%
\pgfpathmoveto{\pgfqpoint{9.436933in}{4.444323in}}%
\pgfpathcurveto{\pgfqpoint{9.447187in}{4.444323in}}{\pgfqpoint{9.457023in}{4.448397in}}{\pgfqpoint{9.464274in}{4.455648in}}%
\pgfpathcurveto{\pgfqpoint{9.471524in}{4.462899in}}{\pgfqpoint{9.475598in}{4.472734in}}{\pgfqpoint{9.475598in}{4.482988in}}%
\pgfpathcurveto{\pgfqpoint{9.475598in}{4.493242in}}{\pgfqpoint{9.471524in}{4.503078in}}{\pgfqpoint{9.464274in}{4.510328in}}%
\pgfpathcurveto{\pgfqpoint{9.457023in}{4.517579in}}{\pgfqpoint{9.447187in}{4.521653in}}{\pgfqpoint{9.436933in}{4.521653in}}%
\pgfpathcurveto{\pgfqpoint{9.426679in}{4.521653in}}{\pgfqpoint{9.416844in}{4.517579in}}{\pgfqpoint{9.409593in}{4.510328in}}%
\pgfpathcurveto{\pgfqpoint{9.402342in}{4.503078in}}{\pgfqpoint{9.398268in}{4.493242in}}{\pgfqpoint{9.398268in}{4.482988in}}%
\pgfpathcurveto{\pgfqpoint{9.398268in}{4.472734in}}{\pgfqpoint{9.402342in}{4.462899in}}{\pgfqpoint{9.409593in}{4.455648in}}%
\pgfpathcurveto{\pgfqpoint{9.416844in}{4.448397in}}{\pgfqpoint{9.426679in}{4.444323in}}{\pgfqpoint{9.436933in}{4.444323in}}%
\pgfpathlineto{\pgfqpoint{9.436933in}{4.444323in}}%
\pgfpathclose%
\pgfusepath{stroke}%
\end{pgfscope}%
\begin{pgfscope}%
\pgfpathrectangle{\pgfqpoint{7.904091in}{2.928000in}}{\pgfqpoint{5.664000in}{5.664000in}}%
\pgfusepath{clip}%
\pgfsetbuttcap%
\pgfsetroundjoin%
\pgfsetlinewidth{1.003750pt}%
\definecolor{currentstroke}{rgb}{0.321569,0.321569,0.321569}%
\pgfsetstrokecolor{currentstroke}%
\pgfsetstrokeopacity{0.900000}%
\pgfsetdash{}{0pt}%
\pgfpathmoveto{\pgfqpoint{8.927501in}{4.028134in}}%
\pgfpathcurveto{\pgfqpoint{8.937755in}{4.028134in}}{\pgfqpoint{8.947590in}{4.032208in}}{\pgfqpoint{8.954841in}{4.039458in}}%
\pgfpathcurveto{\pgfqpoint{8.962092in}{4.046709in}}{\pgfqpoint{8.966166in}{4.056545in}}{\pgfqpoint{8.966166in}{4.066799in}}%
\pgfpathcurveto{\pgfqpoint{8.966166in}{4.077053in}}{\pgfqpoint{8.962092in}{4.086888in}}{\pgfqpoint{8.954841in}{4.094139in}}%
\pgfpathcurveto{\pgfqpoint{8.947590in}{4.101390in}}{\pgfqpoint{8.937755in}{4.105464in}}{\pgfqpoint{8.927501in}{4.105464in}}%
\pgfpathcurveto{\pgfqpoint{8.917247in}{4.105464in}}{\pgfqpoint{8.907411in}{4.101390in}}{\pgfqpoint{8.900160in}{4.094139in}}%
\pgfpathcurveto{\pgfqpoint{8.892910in}{4.086888in}}{\pgfqpoint{8.888836in}{4.077053in}}{\pgfqpoint{8.888836in}{4.066799in}}%
\pgfpathcurveto{\pgfqpoint{8.888836in}{4.056545in}}{\pgfqpoint{8.892910in}{4.046709in}}{\pgfqpoint{8.900160in}{4.039458in}}%
\pgfpathcurveto{\pgfqpoint{8.907411in}{4.032208in}}{\pgfqpoint{8.917247in}{4.028134in}}{\pgfqpoint{8.927501in}{4.028134in}}%
\pgfpathlineto{\pgfqpoint{8.927501in}{4.028134in}}%
\pgfpathclose%
\pgfusepath{stroke}%
\end{pgfscope}%
\begin{pgfscope}%
\pgfpathrectangle{\pgfqpoint{7.904091in}{2.928000in}}{\pgfqpoint{5.664000in}{5.664000in}}%
\pgfusepath{clip}%
\pgfsetbuttcap%
\pgfsetroundjoin%
\pgfsetlinewidth{1.003750pt}%
\definecolor{currentstroke}{rgb}{0.321569,0.321569,0.321569}%
\pgfsetstrokecolor{currentstroke}%
\pgfsetstrokeopacity{0.900000}%
\pgfsetdash{}{0pt}%
\pgfpathmoveto{\pgfqpoint{9.123445in}{4.169725in}}%
\pgfpathcurveto{\pgfqpoint{9.133699in}{4.169725in}}{\pgfqpoint{9.143535in}{4.173799in}}{\pgfqpoint{9.150785in}{4.181050in}}%
\pgfpathcurveto{\pgfqpoint{9.158036in}{4.188301in}}{\pgfqpoint{9.162110in}{4.198136in}}{\pgfqpoint{9.162110in}{4.208390in}}%
\pgfpathcurveto{\pgfqpoint{9.162110in}{4.218645in}}{\pgfqpoint{9.158036in}{4.228480in}}{\pgfqpoint{9.150785in}{4.235731in}}%
\pgfpathcurveto{\pgfqpoint{9.143535in}{4.242981in}}{\pgfqpoint{9.133699in}{4.247055in}}{\pgfqpoint{9.123445in}{4.247055in}}%
\pgfpathcurveto{\pgfqpoint{9.113191in}{4.247055in}}{\pgfqpoint{9.103356in}{4.242981in}}{\pgfqpoint{9.096105in}{4.235731in}}%
\pgfpathcurveto{\pgfqpoint{9.088854in}{4.228480in}}{\pgfqpoint{9.084780in}{4.218645in}}{\pgfqpoint{9.084780in}{4.208390in}}%
\pgfpathcurveto{\pgfqpoint{9.084780in}{4.198136in}}{\pgfqpoint{9.088854in}{4.188301in}}{\pgfqpoint{9.096105in}{4.181050in}}%
\pgfpathcurveto{\pgfqpoint{9.103356in}{4.173799in}}{\pgfqpoint{9.113191in}{4.169725in}}{\pgfqpoint{9.123445in}{4.169725in}}%
\pgfpathlineto{\pgfqpoint{9.123445in}{4.169725in}}%
\pgfpathclose%
\pgfusepath{stroke}%
\end{pgfscope}%
\begin{pgfscope}%
\pgfpathrectangle{\pgfqpoint{7.904091in}{2.928000in}}{\pgfqpoint{5.664000in}{5.664000in}}%
\pgfusepath{clip}%
\pgfsetbuttcap%
\pgfsetroundjoin%
\pgfsetlinewidth{1.003750pt}%
\definecolor{currentstroke}{rgb}{0.321569,0.321569,0.321569}%
\pgfsetstrokecolor{currentstroke}%
\pgfsetstrokeopacity{0.900000}%
\pgfsetdash{}{0pt}%
\pgfpathmoveto{\pgfqpoint{9.612128in}{4.431866in}}%
\pgfpathcurveto{\pgfqpoint{9.622382in}{4.431866in}}{\pgfqpoint{9.632217in}{4.435940in}}{\pgfqpoint{9.639468in}{4.443191in}}%
\pgfpathcurveto{\pgfqpoint{9.646719in}{4.450442in}}{\pgfqpoint{9.650793in}{4.460277in}}{\pgfqpoint{9.650793in}{4.470531in}}%
\pgfpathcurveto{\pgfqpoint{9.650793in}{4.480785in}}{\pgfqpoint{9.646719in}{4.490621in}}{\pgfqpoint{9.639468in}{4.497872in}}%
\pgfpathcurveto{\pgfqpoint{9.632217in}{4.505122in}}{\pgfqpoint{9.622382in}{4.509196in}}{\pgfqpoint{9.612128in}{4.509196in}}%
\pgfpathcurveto{\pgfqpoint{9.601874in}{4.509196in}}{\pgfqpoint{9.592038in}{4.505122in}}{\pgfqpoint{9.584787in}{4.497872in}}%
\pgfpathcurveto{\pgfqpoint{9.577537in}{4.490621in}}{\pgfqpoint{9.573463in}{4.480785in}}{\pgfqpoint{9.573463in}{4.470531in}}%
\pgfpathcurveto{\pgfqpoint{9.573463in}{4.460277in}}{\pgfqpoint{9.577537in}{4.450442in}}{\pgfqpoint{9.584787in}{4.443191in}}%
\pgfpathcurveto{\pgfqpoint{9.592038in}{4.435940in}}{\pgfqpoint{9.601874in}{4.431866in}}{\pgfqpoint{9.612128in}{4.431866in}}%
\pgfpathlineto{\pgfqpoint{9.612128in}{4.431866in}}%
\pgfpathclose%
\pgfusepath{stroke}%
\end{pgfscope}%
\begin{pgfscope}%
\pgfpathrectangle{\pgfqpoint{7.904091in}{2.928000in}}{\pgfqpoint{5.664000in}{5.664000in}}%
\pgfusepath{clip}%
\pgfsetbuttcap%
\pgfsetroundjoin%
\pgfsetlinewidth{1.003750pt}%
\definecolor{currentstroke}{rgb}{0.321569,0.321569,0.321569}%
\pgfsetstrokecolor{currentstroke}%
\pgfsetstrokeopacity{0.900000}%
\pgfsetdash{}{0pt}%
\pgfpathmoveto{\pgfqpoint{10.081384in}{4.865887in}}%
\pgfpathcurveto{\pgfqpoint{10.091638in}{4.865887in}}{\pgfqpoint{10.101474in}{4.869961in}}{\pgfqpoint{10.108724in}{4.877211in}}%
\pgfpathcurveto{\pgfqpoint{10.115975in}{4.884462in}}{\pgfqpoint{10.120049in}{4.894298in}}{\pgfqpoint{10.120049in}{4.904552in}}%
\pgfpathcurveto{\pgfqpoint{10.120049in}{4.914806in}}{\pgfqpoint{10.115975in}{4.924641in}}{\pgfqpoint{10.108724in}{4.931892in}}%
\pgfpathcurveto{\pgfqpoint{10.101474in}{4.939143in}}{\pgfqpoint{10.091638in}{4.943217in}}{\pgfqpoint{10.081384in}{4.943217in}}%
\pgfpathcurveto{\pgfqpoint{10.071130in}{4.943217in}}{\pgfqpoint{10.061295in}{4.939143in}}{\pgfqpoint{10.054044in}{4.931892in}}%
\pgfpathcurveto{\pgfqpoint{10.046793in}{4.924641in}}{\pgfqpoint{10.042719in}{4.914806in}}{\pgfqpoint{10.042719in}{4.904552in}}%
\pgfpathcurveto{\pgfqpoint{10.042719in}{4.894298in}}{\pgfqpoint{10.046793in}{4.884462in}}{\pgfqpoint{10.054044in}{4.877211in}}%
\pgfpathcurveto{\pgfqpoint{10.061295in}{4.869961in}}{\pgfqpoint{10.071130in}{4.865887in}}{\pgfqpoint{10.081384in}{4.865887in}}%
\pgfpathlineto{\pgfqpoint{10.081384in}{4.865887in}}%
\pgfpathclose%
\pgfusepath{stroke}%
\end{pgfscope}%
\begin{pgfscope}%
\pgfpathrectangle{\pgfqpoint{7.904091in}{2.928000in}}{\pgfqpoint{5.664000in}{5.664000in}}%
\pgfusepath{clip}%
\pgfsetbuttcap%
\pgfsetroundjoin%
\pgfsetlinewidth{1.003750pt}%
\definecolor{currentstroke}{rgb}{0.321569,0.321569,0.321569}%
\pgfsetstrokecolor{currentstroke}%
\pgfsetstrokeopacity{0.900000}%
\pgfsetdash{}{0pt}%
\pgfpathmoveto{\pgfqpoint{12.440013in}{7.856118in}}%
\pgfpathcurveto{\pgfqpoint{12.450267in}{7.856118in}}{\pgfqpoint{12.460103in}{7.860192in}}{\pgfqpoint{12.467354in}{7.867443in}}%
\pgfpathcurveto{\pgfqpoint{12.474604in}{7.874693in}}{\pgfqpoint{12.478678in}{7.884529in}}{\pgfqpoint{12.478678in}{7.894783in}}%
\pgfpathcurveto{\pgfqpoint{12.478678in}{7.905037in}}{\pgfqpoint{12.474604in}{7.914873in}}{\pgfqpoint{12.467354in}{7.922123in}}%
\pgfpathcurveto{\pgfqpoint{12.460103in}{7.929374in}}{\pgfqpoint{12.450267in}{7.933448in}}{\pgfqpoint{12.440013in}{7.933448in}}%
\pgfpathcurveto{\pgfqpoint{12.429759in}{7.933448in}}{\pgfqpoint{12.419924in}{7.929374in}}{\pgfqpoint{12.412673in}{7.922123in}}%
\pgfpathcurveto{\pgfqpoint{12.405422in}{7.914873in}}{\pgfqpoint{12.401348in}{7.905037in}}{\pgfqpoint{12.401348in}{7.894783in}}%
\pgfpathcurveto{\pgfqpoint{12.401348in}{7.884529in}}{\pgfqpoint{12.405422in}{7.874693in}}{\pgfqpoint{12.412673in}{7.867443in}}%
\pgfpathcurveto{\pgfqpoint{12.419924in}{7.860192in}}{\pgfqpoint{12.429759in}{7.856118in}}{\pgfqpoint{12.440013in}{7.856118in}}%
\pgfpathlineto{\pgfqpoint{12.440013in}{7.856118in}}%
\pgfpathclose%
\pgfusepath{stroke}%
\end{pgfscope}%
\begin{pgfscope}%
\pgfpathrectangle{\pgfqpoint{7.904091in}{2.928000in}}{\pgfqpoint{5.664000in}{5.664000in}}%
\pgfusepath{clip}%
\pgfsetbuttcap%
\pgfsetroundjoin%
\pgfsetlinewidth{1.003750pt}%
\definecolor{currentstroke}{rgb}{0.321569,0.321569,0.321569}%
\pgfsetstrokecolor{currentstroke}%
\pgfsetstrokeopacity{0.900000}%
\pgfsetdash{}{0pt}%
\pgfpathmoveto{\pgfqpoint{12.208833in}{7.648010in}}%
\pgfpathcurveto{\pgfqpoint{12.219087in}{7.648010in}}{\pgfqpoint{12.228922in}{7.652084in}}{\pgfqpoint{12.236173in}{7.659335in}}%
\pgfpathcurveto{\pgfqpoint{12.243424in}{7.666586in}}{\pgfqpoint{12.247498in}{7.676421in}}{\pgfqpoint{12.247498in}{7.686675in}}%
\pgfpathcurveto{\pgfqpoint{12.247498in}{7.696929in}}{\pgfqpoint{12.243424in}{7.706765in}}{\pgfqpoint{12.236173in}{7.714015in}}%
\pgfpathcurveto{\pgfqpoint{12.228922in}{7.721266in}}{\pgfqpoint{12.219087in}{7.725340in}}{\pgfqpoint{12.208833in}{7.725340in}}%
\pgfpathcurveto{\pgfqpoint{12.198579in}{7.725340in}}{\pgfqpoint{12.188743in}{7.721266in}}{\pgfqpoint{12.181493in}{7.714015in}}%
\pgfpathcurveto{\pgfqpoint{12.174242in}{7.706765in}}{\pgfqpoint{12.170168in}{7.696929in}}{\pgfqpoint{12.170168in}{7.686675in}}%
\pgfpathcurveto{\pgfqpoint{12.170168in}{7.676421in}}{\pgfqpoint{12.174242in}{7.666586in}}{\pgfqpoint{12.181493in}{7.659335in}}%
\pgfpathcurveto{\pgfqpoint{12.188743in}{7.652084in}}{\pgfqpoint{12.198579in}{7.648010in}}{\pgfqpoint{12.208833in}{7.648010in}}%
\pgfpathlineto{\pgfqpoint{12.208833in}{7.648010in}}%
\pgfpathclose%
\pgfusepath{stroke}%
\end{pgfscope}%
\begin{pgfscope}%
\pgfpathrectangle{\pgfqpoint{7.904091in}{2.928000in}}{\pgfqpoint{5.664000in}{5.664000in}}%
\pgfusepath{clip}%
\pgfsetbuttcap%
\pgfsetroundjoin%
\pgfsetlinewidth{1.003750pt}%
\definecolor{currentstroke}{rgb}{0.321569,0.321569,0.321569}%
\pgfsetstrokecolor{currentstroke}%
\pgfsetstrokeopacity{0.900000}%
\pgfsetdash{}{0pt}%
\pgfpathmoveto{\pgfqpoint{12.401716in}{7.832953in}}%
\pgfpathcurveto{\pgfqpoint{12.411970in}{7.832953in}}{\pgfqpoint{12.421806in}{7.837027in}}{\pgfqpoint{12.429057in}{7.844278in}}%
\pgfpathcurveto{\pgfqpoint{12.436307in}{7.851529in}}{\pgfqpoint{12.440381in}{7.861364in}}{\pgfqpoint{12.440381in}{7.871618in}}%
\pgfpathcurveto{\pgfqpoint{12.440381in}{7.881873in}}{\pgfqpoint{12.436307in}{7.891708in}}{\pgfqpoint{12.429057in}{7.898959in}}%
\pgfpathcurveto{\pgfqpoint{12.421806in}{7.906210in}}{\pgfqpoint{12.411970in}{7.910283in}}{\pgfqpoint{12.401716in}{7.910283in}}%
\pgfpathcurveto{\pgfqpoint{12.391462in}{7.910283in}}{\pgfqpoint{12.381627in}{7.906210in}}{\pgfqpoint{12.374376in}{7.898959in}}%
\pgfpathcurveto{\pgfqpoint{12.367125in}{7.891708in}}{\pgfqpoint{12.363051in}{7.881873in}}{\pgfqpoint{12.363051in}{7.871618in}}%
\pgfpathcurveto{\pgfqpoint{12.363051in}{7.861364in}}{\pgfqpoint{12.367125in}{7.851529in}}{\pgfqpoint{12.374376in}{7.844278in}}%
\pgfpathcurveto{\pgfqpoint{12.381627in}{7.837027in}}{\pgfqpoint{12.391462in}{7.832953in}}{\pgfqpoint{12.401716in}{7.832953in}}%
\pgfpathlineto{\pgfqpoint{12.401716in}{7.832953in}}%
\pgfpathclose%
\pgfusepath{stroke}%
\end{pgfscope}%
\begin{pgfscope}%
\pgfpathrectangle{\pgfqpoint{7.904091in}{2.928000in}}{\pgfqpoint{5.664000in}{5.664000in}}%
\pgfusepath{clip}%
\pgfsetbuttcap%
\pgfsetroundjoin%
\pgfsetlinewidth{1.003750pt}%
\definecolor{currentstroke}{rgb}{0.321569,0.321569,0.321569}%
\pgfsetstrokecolor{currentstroke}%
\pgfsetstrokeopacity{0.900000}%
\pgfsetdash{}{0pt}%
\pgfpathmoveto{\pgfqpoint{13.014440in}{8.213497in}}%
\pgfpathcurveto{\pgfqpoint{13.024694in}{8.213497in}}{\pgfqpoint{13.034529in}{8.217571in}}{\pgfqpoint{13.041780in}{8.224822in}}%
\pgfpathcurveto{\pgfqpoint{13.049031in}{8.232072in}}{\pgfqpoint{13.053105in}{8.241908in}}{\pgfqpoint{13.053105in}{8.252162in}}%
\pgfpathcurveto{\pgfqpoint{13.053105in}{8.262416in}}{\pgfqpoint{13.049031in}{8.272251in}}{\pgfqpoint{13.041780in}{8.279502in}}%
\pgfpathcurveto{\pgfqpoint{13.034529in}{8.286753in}}{\pgfqpoint{13.024694in}{8.290827in}}{\pgfqpoint{13.014440in}{8.290827in}}%
\pgfpathcurveto{\pgfqpoint{13.004186in}{8.290827in}}{\pgfqpoint{12.994350in}{8.286753in}}{\pgfqpoint{12.987099in}{8.279502in}}%
\pgfpathcurveto{\pgfqpoint{12.979849in}{8.272251in}}{\pgfqpoint{12.975775in}{8.262416in}}{\pgfqpoint{12.975775in}{8.252162in}}%
\pgfpathcurveto{\pgfqpoint{12.975775in}{8.241908in}}{\pgfqpoint{12.979849in}{8.232072in}}{\pgfqpoint{12.987099in}{8.224822in}}%
\pgfpathcurveto{\pgfqpoint{12.994350in}{8.217571in}}{\pgfqpoint{13.004186in}{8.213497in}}{\pgfqpoint{13.014440in}{8.213497in}}%
\pgfpathlineto{\pgfqpoint{13.014440in}{8.213497in}}%
\pgfpathclose%
\pgfusepath{stroke}%
\end{pgfscope}%
\begin{pgfscope}%
\pgfpathrectangle{\pgfqpoint{7.904091in}{2.928000in}}{\pgfqpoint{5.664000in}{5.664000in}}%
\pgfusepath{clip}%
\pgfsetbuttcap%
\pgfsetroundjoin%
\pgfsetlinewidth{1.003750pt}%
\definecolor{currentstroke}{rgb}{0.321569,0.321569,0.321569}%
\pgfsetstrokecolor{currentstroke}%
\pgfsetstrokeopacity{0.900000}%
\pgfsetdash{}{0pt}%
\pgfpathmoveto{\pgfqpoint{12.588935in}{7.871090in}}%
\pgfpathcurveto{\pgfqpoint{12.599189in}{7.871090in}}{\pgfqpoint{12.609024in}{7.875164in}}{\pgfqpoint{12.616275in}{7.882415in}}%
\pgfpathcurveto{\pgfqpoint{12.623526in}{7.889666in}}{\pgfqpoint{12.627600in}{7.899501in}}{\pgfqpoint{12.627600in}{7.909755in}}%
\pgfpathcurveto{\pgfqpoint{12.627600in}{7.920009in}}{\pgfqpoint{12.623526in}{7.929845in}}{\pgfqpoint{12.616275in}{7.937095in}}%
\pgfpathcurveto{\pgfqpoint{12.609024in}{7.944346in}}{\pgfqpoint{12.599189in}{7.948420in}}{\pgfqpoint{12.588935in}{7.948420in}}%
\pgfpathcurveto{\pgfqpoint{12.578681in}{7.948420in}}{\pgfqpoint{12.568845in}{7.944346in}}{\pgfqpoint{12.561595in}{7.937095in}}%
\pgfpathcurveto{\pgfqpoint{12.554344in}{7.929845in}}{\pgfqpoint{12.550270in}{7.920009in}}{\pgfqpoint{12.550270in}{7.909755in}}%
\pgfpathcurveto{\pgfqpoint{12.550270in}{7.899501in}}{\pgfqpoint{12.554344in}{7.889666in}}{\pgfqpoint{12.561595in}{7.882415in}}%
\pgfpathcurveto{\pgfqpoint{12.568845in}{7.875164in}}{\pgfqpoint{12.578681in}{7.871090in}}{\pgfqpoint{12.588935in}{7.871090in}}%
\pgfpathlineto{\pgfqpoint{12.588935in}{7.871090in}}%
\pgfpathclose%
\pgfusepath{stroke}%
\end{pgfscope}%
\begin{pgfscope}%
\pgfpathrectangle{\pgfqpoint{7.904091in}{2.928000in}}{\pgfqpoint{5.664000in}{5.664000in}}%
\pgfusepath{clip}%
\pgfsetbuttcap%
\pgfsetroundjoin%
\pgfsetlinewidth{1.003750pt}%
\definecolor{currentstroke}{rgb}{0.321569,0.321569,0.321569}%
\pgfsetstrokecolor{currentstroke}%
\pgfsetstrokeopacity{0.900000}%
\pgfsetdash{}{0pt}%
\pgfpathmoveto{\pgfqpoint{10.049441in}{4.936371in}}%
\pgfpathcurveto{\pgfqpoint{10.059695in}{4.936371in}}{\pgfqpoint{10.069531in}{4.940445in}}{\pgfqpoint{10.076781in}{4.947696in}}%
\pgfpathcurveto{\pgfqpoint{10.084032in}{4.954947in}}{\pgfqpoint{10.088106in}{4.964782in}}{\pgfqpoint{10.088106in}{4.975036in}}%
\pgfpathcurveto{\pgfqpoint{10.088106in}{4.985290in}}{\pgfqpoint{10.084032in}{4.995126in}}{\pgfqpoint{10.076781in}{5.002376in}}%
\pgfpathcurveto{\pgfqpoint{10.069531in}{5.009627in}}{\pgfqpoint{10.059695in}{5.013701in}}{\pgfqpoint{10.049441in}{5.013701in}}%
\pgfpathcurveto{\pgfqpoint{10.039187in}{5.013701in}}{\pgfqpoint{10.029352in}{5.009627in}}{\pgfqpoint{10.022101in}{5.002376in}}%
\pgfpathcurveto{\pgfqpoint{10.014850in}{4.995126in}}{\pgfqpoint{10.010776in}{4.985290in}}{\pgfqpoint{10.010776in}{4.975036in}}%
\pgfpathcurveto{\pgfqpoint{10.010776in}{4.964782in}}{\pgfqpoint{10.014850in}{4.954947in}}{\pgfqpoint{10.022101in}{4.947696in}}%
\pgfpathcurveto{\pgfqpoint{10.029352in}{4.940445in}}{\pgfqpoint{10.039187in}{4.936371in}}{\pgfqpoint{10.049441in}{4.936371in}}%
\pgfpathlineto{\pgfqpoint{10.049441in}{4.936371in}}%
\pgfpathclose%
\pgfusepath{stroke}%
\end{pgfscope}%
\begin{pgfscope}%
\pgfpathrectangle{\pgfqpoint{7.904091in}{2.928000in}}{\pgfqpoint{5.664000in}{5.664000in}}%
\pgfusepath{clip}%
\pgfsetbuttcap%
\pgfsetroundjoin%
\pgfsetlinewidth{1.003750pt}%
\definecolor{currentstroke}{rgb}{0.321569,0.321569,0.321569}%
\pgfsetstrokecolor{currentstroke}%
\pgfsetstrokeopacity{0.900000}%
\pgfsetdash{}{0pt}%
\pgfpathmoveto{\pgfqpoint{10.091640in}{5.044579in}}%
\pgfpathcurveto{\pgfqpoint{10.101894in}{5.044579in}}{\pgfqpoint{10.111730in}{5.048653in}}{\pgfqpoint{10.118981in}{5.055904in}}%
\pgfpathcurveto{\pgfqpoint{10.126231in}{5.063155in}}{\pgfqpoint{10.130305in}{5.072990in}}{\pgfqpoint{10.130305in}{5.083244in}}%
\pgfpathcurveto{\pgfqpoint{10.130305in}{5.093499in}}{\pgfqpoint{10.126231in}{5.103334in}}{\pgfqpoint{10.118981in}{5.110585in}}%
\pgfpathcurveto{\pgfqpoint{10.111730in}{5.117836in}}{\pgfqpoint{10.101894in}{5.121910in}}{\pgfqpoint{10.091640in}{5.121910in}}%
\pgfpathcurveto{\pgfqpoint{10.081386in}{5.121910in}}{\pgfqpoint{10.071551in}{5.117836in}}{\pgfqpoint{10.064300in}{5.110585in}}%
\pgfpathcurveto{\pgfqpoint{10.057049in}{5.103334in}}{\pgfqpoint{10.052975in}{5.093499in}}{\pgfqpoint{10.052975in}{5.083244in}}%
\pgfpathcurveto{\pgfqpoint{10.052975in}{5.072990in}}{\pgfqpoint{10.057049in}{5.063155in}}{\pgfqpoint{10.064300in}{5.055904in}}%
\pgfpathcurveto{\pgfqpoint{10.071551in}{5.048653in}}{\pgfqpoint{10.081386in}{5.044579in}}{\pgfqpoint{10.091640in}{5.044579in}}%
\pgfpathlineto{\pgfqpoint{10.091640in}{5.044579in}}%
\pgfpathclose%
\pgfusepath{stroke}%
\end{pgfscope}%
\begin{pgfscope}%
\pgfpathrectangle{\pgfqpoint{7.904091in}{2.928000in}}{\pgfqpoint{5.664000in}{5.664000in}}%
\pgfusepath{clip}%
\pgfsetbuttcap%
\pgfsetroundjoin%
\pgfsetlinewidth{1.003750pt}%
\definecolor{currentstroke}{rgb}{0.321569,0.321569,0.321569}%
\pgfsetstrokecolor{currentstroke}%
\pgfsetstrokeopacity{0.900000}%
\pgfsetdash{}{0pt}%
\pgfpathmoveto{\pgfqpoint{9.261261in}{4.297888in}}%
\pgfpathcurveto{\pgfqpoint{9.271515in}{4.297888in}}{\pgfqpoint{9.281350in}{4.301962in}}{\pgfqpoint{9.288601in}{4.309213in}}%
\pgfpathcurveto{\pgfqpoint{9.295852in}{4.316464in}}{\pgfqpoint{9.299926in}{4.326299in}}{\pgfqpoint{9.299926in}{4.336553in}}%
\pgfpathcurveto{\pgfqpoint{9.299926in}{4.346807in}}{\pgfqpoint{9.295852in}{4.356643in}}{\pgfqpoint{9.288601in}{4.363894in}}%
\pgfpathcurveto{\pgfqpoint{9.281350in}{4.371144in}}{\pgfqpoint{9.271515in}{4.375218in}}{\pgfqpoint{9.261261in}{4.375218in}}%
\pgfpathcurveto{\pgfqpoint{9.251006in}{4.375218in}}{\pgfqpoint{9.241171in}{4.371144in}}{\pgfqpoint{9.233920in}{4.363894in}}%
\pgfpathcurveto{\pgfqpoint{9.226669in}{4.356643in}}{\pgfqpoint{9.222595in}{4.346807in}}{\pgfqpoint{9.222595in}{4.336553in}}%
\pgfpathcurveto{\pgfqpoint{9.222595in}{4.326299in}}{\pgfqpoint{9.226669in}{4.316464in}}{\pgfqpoint{9.233920in}{4.309213in}}%
\pgfpathcurveto{\pgfqpoint{9.241171in}{4.301962in}}{\pgfqpoint{9.251006in}{4.297888in}}{\pgfqpoint{9.261261in}{4.297888in}}%
\pgfpathlineto{\pgfqpoint{9.261261in}{4.297888in}}%
\pgfpathclose%
\pgfusepath{stroke}%
\end{pgfscope}%
\begin{pgfscope}%
\pgfpathrectangle{\pgfqpoint{7.904091in}{2.928000in}}{\pgfqpoint{5.664000in}{5.664000in}}%
\pgfusepath{clip}%
\pgfsetbuttcap%
\pgfsetroundjoin%
\pgfsetlinewidth{1.003750pt}%
\definecolor{currentstroke}{rgb}{0.321569,0.321569,0.321569}%
\pgfsetstrokecolor{currentstroke}%
\pgfsetstrokeopacity{0.900000}%
\pgfsetdash{}{0pt}%
\pgfpathmoveto{\pgfqpoint{9.545860in}{4.556174in}}%
\pgfpathcurveto{\pgfqpoint{9.556114in}{4.556174in}}{\pgfqpoint{9.565949in}{4.560248in}}{\pgfqpoint{9.573200in}{4.567499in}}%
\pgfpathcurveto{\pgfqpoint{9.580451in}{4.574750in}}{\pgfqpoint{9.584525in}{4.584585in}}{\pgfqpoint{9.584525in}{4.594839in}}%
\pgfpathcurveto{\pgfqpoint{9.584525in}{4.605093in}}{\pgfqpoint{9.580451in}{4.614929in}}{\pgfqpoint{9.573200in}{4.622180in}}%
\pgfpathcurveto{\pgfqpoint{9.565949in}{4.629430in}}{\pgfqpoint{9.556114in}{4.633504in}}{\pgfqpoint{9.545860in}{4.633504in}}%
\pgfpathcurveto{\pgfqpoint{9.535606in}{4.633504in}}{\pgfqpoint{9.525770in}{4.629430in}}{\pgfqpoint{9.518519in}{4.622180in}}%
\pgfpathcurveto{\pgfqpoint{9.511269in}{4.614929in}}{\pgfqpoint{9.507195in}{4.605093in}}{\pgfqpoint{9.507195in}{4.594839in}}%
\pgfpathcurveto{\pgfqpoint{9.507195in}{4.584585in}}{\pgfqpoint{9.511269in}{4.574750in}}{\pgfqpoint{9.518519in}{4.567499in}}%
\pgfpathcurveto{\pgfqpoint{9.525770in}{4.560248in}}{\pgfqpoint{9.535606in}{4.556174in}}{\pgfqpoint{9.545860in}{4.556174in}}%
\pgfpathlineto{\pgfqpoint{9.545860in}{4.556174in}}%
\pgfpathclose%
\pgfusepath{stroke}%
\end{pgfscope}%
\begin{pgfscope}%
\pgfpathrectangle{\pgfqpoint{7.904091in}{2.928000in}}{\pgfqpoint{5.664000in}{5.664000in}}%
\pgfusepath{clip}%
\pgfsetbuttcap%
\pgfsetroundjoin%
\pgfsetlinewidth{1.003750pt}%
\definecolor{currentstroke}{rgb}{0.321569,0.321569,0.321569}%
\pgfsetstrokecolor{currentstroke}%
\pgfsetstrokeopacity{0.900000}%
\pgfsetdash{}{0pt}%
\pgfpathmoveto{\pgfqpoint{9.030166in}{4.101798in}}%
\pgfpathcurveto{\pgfqpoint{9.040420in}{4.101798in}}{\pgfqpoint{9.050256in}{4.105872in}}{\pgfqpoint{9.057506in}{4.113122in}}%
\pgfpathcurveto{\pgfqpoint{9.064757in}{4.120373in}}{\pgfqpoint{9.068831in}{4.130209in}}{\pgfqpoint{9.068831in}{4.140463in}}%
\pgfpathcurveto{\pgfqpoint{9.068831in}{4.150717in}}{\pgfqpoint{9.064757in}{4.160552in}}{\pgfqpoint{9.057506in}{4.167803in}}%
\pgfpathcurveto{\pgfqpoint{9.050256in}{4.175054in}}{\pgfqpoint{9.040420in}{4.179128in}}{\pgfqpoint{9.030166in}{4.179128in}}%
\pgfpathcurveto{\pgfqpoint{9.019912in}{4.179128in}}{\pgfqpoint{9.010076in}{4.175054in}}{\pgfqpoint{9.002826in}{4.167803in}}%
\pgfpathcurveto{\pgfqpoint{8.995575in}{4.160552in}}{\pgfqpoint{8.991501in}{4.150717in}}{\pgfqpoint{8.991501in}{4.140463in}}%
\pgfpathcurveto{\pgfqpoint{8.991501in}{4.130209in}}{\pgfqpoint{8.995575in}{4.120373in}}{\pgfqpoint{9.002826in}{4.113122in}}%
\pgfpathcurveto{\pgfqpoint{9.010076in}{4.105872in}}{\pgfqpoint{9.019912in}{4.101798in}}{\pgfqpoint{9.030166in}{4.101798in}}%
\pgfpathlineto{\pgfqpoint{9.030166in}{4.101798in}}%
\pgfpathclose%
\pgfusepath{stroke}%
\end{pgfscope}%
\begin{pgfscope}%
\pgfpathrectangle{\pgfqpoint{7.904091in}{2.928000in}}{\pgfqpoint{5.664000in}{5.664000in}}%
\pgfusepath{clip}%
\pgfsetbuttcap%
\pgfsetroundjoin%
\pgfsetlinewidth{1.003750pt}%
\definecolor{currentstroke}{rgb}{0.321569,0.321569,0.321569}%
\pgfsetstrokecolor{currentstroke}%
\pgfsetstrokeopacity{0.900000}%
\pgfsetdash{}{0pt}%
\pgfpathmoveto{\pgfqpoint{9.323819in}{4.344975in}}%
\pgfpathcurveto{\pgfqpoint{9.334073in}{4.344975in}}{\pgfqpoint{9.343908in}{4.349049in}}{\pgfqpoint{9.351159in}{4.356299in}}%
\pgfpathcurveto{\pgfqpoint{9.358410in}{4.363550in}}{\pgfqpoint{9.362484in}{4.373386in}}{\pgfqpoint{9.362484in}{4.383640in}}%
\pgfpathcurveto{\pgfqpoint{9.362484in}{4.393894in}}{\pgfqpoint{9.358410in}{4.403729in}}{\pgfqpoint{9.351159in}{4.410980in}}%
\pgfpathcurveto{\pgfqpoint{9.343908in}{4.418231in}}{\pgfqpoint{9.334073in}{4.422305in}}{\pgfqpoint{9.323819in}{4.422305in}}%
\pgfpathcurveto{\pgfqpoint{9.313565in}{4.422305in}}{\pgfqpoint{9.303729in}{4.418231in}}{\pgfqpoint{9.296479in}{4.410980in}}%
\pgfpathcurveto{\pgfqpoint{9.289228in}{4.403729in}}{\pgfqpoint{9.285154in}{4.393894in}}{\pgfqpoint{9.285154in}{4.383640in}}%
\pgfpathcurveto{\pgfqpoint{9.285154in}{4.373386in}}{\pgfqpoint{9.289228in}{4.363550in}}{\pgfqpoint{9.296479in}{4.356299in}}%
\pgfpathcurveto{\pgfqpoint{9.303729in}{4.349049in}}{\pgfqpoint{9.313565in}{4.344975in}}{\pgfqpoint{9.323819in}{4.344975in}}%
\pgfpathlineto{\pgfqpoint{9.323819in}{4.344975in}}%
\pgfpathclose%
\pgfusepath{stroke}%
\end{pgfscope}%
\begin{pgfscope}%
\pgfpathrectangle{\pgfqpoint{7.904091in}{2.928000in}}{\pgfqpoint{5.664000in}{5.664000in}}%
\pgfusepath{clip}%
\pgfsetbuttcap%
\pgfsetroundjoin%
\pgfsetlinewidth{1.003750pt}%
\definecolor{currentstroke}{rgb}{0.321569,0.321569,0.321569}%
\pgfsetstrokecolor{currentstroke}%
\pgfsetstrokeopacity{0.900000}%
\pgfsetdash{}{0pt}%
\pgfpathmoveto{\pgfqpoint{9.832617in}{4.566010in}}%
\pgfpathcurveto{\pgfqpoint{9.842871in}{4.566010in}}{\pgfqpoint{9.852707in}{4.570084in}}{\pgfqpoint{9.859957in}{4.577334in}}%
\pgfpathcurveto{\pgfqpoint{9.867208in}{4.584585in}}{\pgfqpoint{9.871282in}{4.594421in}}{\pgfqpoint{9.871282in}{4.604675in}}%
\pgfpathcurveto{\pgfqpoint{9.871282in}{4.614929in}}{\pgfqpoint{9.867208in}{4.624764in}}{\pgfqpoint{9.859957in}{4.632015in}}%
\pgfpathcurveto{\pgfqpoint{9.852707in}{4.639266in}}{\pgfqpoint{9.842871in}{4.643340in}}{\pgfqpoint{9.832617in}{4.643340in}}%
\pgfpathcurveto{\pgfqpoint{9.822363in}{4.643340in}}{\pgfqpoint{9.812527in}{4.639266in}}{\pgfqpoint{9.805277in}{4.632015in}}%
\pgfpathcurveto{\pgfqpoint{9.798026in}{4.624764in}}{\pgfqpoint{9.793952in}{4.614929in}}{\pgfqpoint{9.793952in}{4.604675in}}%
\pgfpathcurveto{\pgfqpoint{9.793952in}{4.594421in}}{\pgfqpoint{9.798026in}{4.584585in}}{\pgfqpoint{9.805277in}{4.577334in}}%
\pgfpathcurveto{\pgfqpoint{9.812527in}{4.570084in}}{\pgfqpoint{9.822363in}{4.566010in}}{\pgfqpoint{9.832617in}{4.566010in}}%
\pgfpathlineto{\pgfqpoint{9.832617in}{4.566010in}}%
\pgfpathclose%
\pgfusepath{stroke}%
\end{pgfscope}%
\begin{pgfscope}%
\pgfpathrectangle{\pgfqpoint{7.904091in}{2.928000in}}{\pgfqpoint{5.664000in}{5.664000in}}%
\pgfusepath{clip}%
\pgfsetbuttcap%
\pgfsetroundjoin%
\pgfsetlinewidth{1.003750pt}%
\definecolor{currentstroke}{rgb}{0.321569,0.321569,0.321569}%
\pgfsetstrokecolor{currentstroke}%
\pgfsetstrokeopacity{0.900000}%
\pgfsetdash{}{0pt}%
\pgfpathmoveto{\pgfqpoint{9.997721in}{4.747063in}}%
\pgfpathcurveto{\pgfqpoint{10.007975in}{4.747063in}}{\pgfqpoint{10.017811in}{4.751137in}}{\pgfqpoint{10.025062in}{4.758387in}}%
\pgfpathcurveto{\pgfqpoint{10.032312in}{4.765638in}}{\pgfqpoint{10.036386in}{4.775474in}}{\pgfqpoint{10.036386in}{4.785728in}}%
\pgfpathcurveto{\pgfqpoint{10.036386in}{4.795982in}}{\pgfqpoint{10.032312in}{4.805817in}}{\pgfqpoint{10.025062in}{4.813068in}}%
\pgfpathcurveto{\pgfqpoint{10.017811in}{4.820319in}}{\pgfqpoint{10.007975in}{4.824393in}}{\pgfqpoint{9.997721in}{4.824393in}}%
\pgfpathcurveto{\pgfqpoint{9.987467in}{4.824393in}}{\pgfqpoint{9.977632in}{4.820319in}}{\pgfqpoint{9.970381in}{4.813068in}}%
\pgfpathcurveto{\pgfqpoint{9.963130in}{4.805817in}}{\pgfqpoint{9.959056in}{4.795982in}}{\pgfqpoint{9.959056in}{4.785728in}}%
\pgfpathcurveto{\pgfqpoint{9.959056in}{4.775474in}}{\pgfqpoint{9.963130in}{4.765638in}}{\pgfqpoint{9.970381in}{4.758387in}}%
\pgfpathcurveto{\pgfqpoint{9.977632in}{4.751137in}}{\pgfqpoint{9.987467in}{4.747063in}}{\pgfqpoint{9.997721in}{4.747063in}}%
\pgfpathlineto{\pgfqpoint{9.997721in}{4.747063in}}%
\pgfpathclose%
\pgfusepath{stroke}%
\end{pgfscope}%
\begin{pgfscope}%
\pgfpathrectangle{\pgfqpoint{7.904091in}{2.928000in}}{\pgfqpoint{5.664000in}{5.664000in}}%
\pgfusepath{clip}%
\pgfsetbuttcap%
\pgfsetroundjoin%
\pgfsetlinewidth{1.003750pt}%
\definecolor{currentstroke}{rgb}{0.321569,0.321569,0.321569}%
\pgfsetstrokecolor{currentstroke}%
\pgfsetstrokeopacity{0.900000}%
\pgfsetdash{}{0pt}%
\pgfpathmoveto{\pgfqpoint{9.796041in}{4.580844in}}%
\pgfpathcurveto{\pgfqpoint{9.806295in}{4.580844in}}{\pgfqpoint{9.816130in}{4.584918in}}{\pgfqpoint{9.823381in}{4.592168in}}%
\pgfpathcurveto{\pgfqpoint{9.830632in}{4.599419in}}{\pgfqpoint{9.834706in}{4.609255in}}{\pgfqpoint{9.834706in}{4.619509in}}%
\pgfpathcurveto{\pgfqpoint{9.834706in}{4.629763in}}{\pgfqpoint{9.830632in}{4.639598in}}{\pgfqpoint{9.823381in}{4.646849in}}%
\pgfpathcurveto{\pgfqpoint{9.816130in}{4.654100in}}{\pgfqpoint{9.806295in}{4.658174in}}{\pgfqpoint{9.796041in}{4.658174in}}%
\pgfpathcurveto{\pgfqpoint{9.785787in}{4.658174in}}{\pgfqpoint{9.775951in}{4.654100in}}{\pgfqpoint{9.768700in}{4.646849in}}%
\pgfpathcurveto{\pgfqpoint{9.761450in}{4.639598in}}{\pgfqpoint{9.757376in}{4.629763in}}{\pgfqpoint{9.757376in}{4.619509in}}%
\pgfpathcurveto{\pgfqpoint{9.757376in}{4.609255in}}{\pgfqpoint{9.761450in}{4.599419in}}{\pgfqpoint{9.768700in}{4.592168in}}%
\pgfpathcurveto{\pgfqpoint{9.775951in}{4.584918in}}{\pgfqpoint{9.785787in}{4.580844in}}{\pgfqpoint{9.796041in}{4.580844in}}%
\pgfpathlineto{\pgfqpoint{9.796041in}{4.580844in}}%
\pgfpathclose%
\pgfusepath{stroke}%
\end{pgfscope}%
\begin{pgfscope}%
\pgfpathrectangle{\pgfqpoint{7.904091in}{2.928000in}}{\pgfqpoint{5.664000in}{5.664000in}}%
\pgfusepath{clip}%
\pgfsetbuttcap%
\pgfsetroundjoin%
\pgfsetlinewidth{1.003750pt}%
\definecolor{currentstroke}{rgb}{0.321569,0.321569,0.321569}%
\pgfsetstrokecolor{currentstroke}%
\pgfsetstrokeopacity{0.900000}%
\pgfsetdash{}{0pt}%
\pgfpathmoveto{\pgfqpoint{10.019274in}{4.856820in}}%
\pgfpathcurveto{\pgfqpoint{10.029528in}{4.856820in}}{\pgfqpoint{10.039363in}{4.860894in}}{\pgfqpoint{10.046614in}{4.868145in}}%
\pgfpathcurveto{\pgfqpoint{10.053865in}{4.875396in}}{\pgfqpoint{10.057939in}{4.885231in}}{\pgfqpoint{10.057939in}{4.895485in}}%
\pgfpathcurveto{\pgfqpoint{10.057939in}{4.905739in}}{\pgfqpoint{10.053865in}{4.915575in}}{\pgfqpoint{10.046614in}{4.922826in}}%
\pgfpathcurveto{\pgfqpoint{10.039363in}{4.930076in}}{\pgfqpoint{10.029528in}{4.934150in}}{\pgfqpoint{10.019274in}{4.934150in}}%
\pgfpathcurveto{\pgfqpoint{10.009019in}{4.934150in}}{\pgfqpoint{9.999184in}{4.930076in}}{\pgfqpoint{9.991933in}{4.922826in}}%
\pgfpathcurveto{\pgfqpoint{9.984682in}{4.915575in}}{\pgfqpoint{9.980609in}{4.905739in}}{\pgfqpoint{9.980609in}{4.895485in}}%
\pgfpathcurveto{\pgfqpoint{9.980609in}{4.885231in}}{\pgfqpoint{9.984682in}{4.875396in}}{\pgfqpoint{9.991933in}{4.868145in}}%
\pgfpathcurveto{\pgfqpoint{9.999184in}{4.860894in}}{\pgfqpoint{10.009019in}{4.856820in}}{\pgfqpoint{10.019274in}{4.856820in}}%
\pgfpathlineto{\pgfqpoint{10.019274in}{4.856820in}}%
\pgfpathclose%
\pgfusepath{stroke}%
\end{pgfscope}%
\begin{pgfscope}%
\pgfpathrectangle{\pgfqpoint{7.904091in}{2.928000in}}{\pgfqpoint{5.664000in}{5.664000in}}%
\pgfusepath{clip}%
\pgfsetbuttcap%
\pgfsetroundjoin%
\pgfsetlinewidth{1.003750pt}%
\definecolor{currentstroke}{rgb}{0.321569,0.321569,0.321569}%
\pgfsetstrokecolor{currentstroke}%
\pgfsetstrokeopacity{0.900000}%
\pgfsetdash{}{0pt}%
\pgfpathmoveto{\pgfqpoint{9.261633in}{4.578995in}}%
\pgfpathcurveto{\pgfqpoint{9.271887in}{4.578995in}}{\pgfqpoint{9.281722in}{4.583069in}}{\pgfqpoint{9.288973in}{4.590320in}}%
\pgfpathcurveto{\pgfqpoint{9.296224in}{4.597570in}}{\pgfqpoint{9.300298in}{4.607406in}}{\pgfqpoint{9.300298in}{4.617660in}}%
\pgfpathcurveto{\pgfqpoint{9.300298in}{4.627914in}}{\pgfqpoint{9.296224in}{4.637750in}}{\pgfqpoint{9.288973in}{4.645000in}}%
\pgfpathcurveto{\pgfqpoint{9.281722in}{4.652251in}}{\pgfqpoint{9.271887in}{4.656325in}}{\pgfqpoint{9.261633in}{4.656325in}}%
\pgfpathcurveto{\pgfqpoint{9.251378in}{4.656325in}}{\pgfqpoint{9.241543in}{4.652251in}}{\pgfqpoint{9.234292in}{4.645000in}}%
\pgfpathcurveto{\pgfqpoint{9.227041in}{4.637750in}}{\pgfqpoint{9.222968in}{4.627914in}}{\pgfqpoint{9.222968in}{4.617660in}}%
\pgfpathcurveto{\pgfqpoint{9.222968in}{4.607406in}}{\pgfqpoint{9.227041in}{4.597570in}}{\pgfqpoint{9.234292in}{4.590320in}}%
\pgfpathcurveto{\pgfqpoint{9.241543in}{4.583069in}}{\pgfqpoint{9.251378in}{4.578995in}}{\pgfqpoint{9.261633in}{4.578995in}}%
\pgfpathlineto{\pgfqpoint{9.261633in}{4.578995in}}%
\pgfpathclose%
\pgfusepath{stroke}%
\end{pgfscope}%
\begin{pgfscope}%
\pgfpathrectangle{\pgfqpoint{7.904091in}{2.928000in}}{\pgfqpoint{5.664000in}{5.664000in}}%
\pgfusepath{clip}%
\pgfsetbuttcap%
\pgfsetroundjoin%
\pgfsetlinewidth{1.003750pt}%
\definecolor{currentstroke}{rgb}{0.321569,0.321569,0.321569}%
\pgfsetstrokecolor{currentstroke}%
\pgfsetstrokeopacity{0.900000}%
\pgfsetdash{}{0pt}%
\pgfpathmoveto{\pgfqpoint{9.389137in}{4.612445in}}%
\pgfpathcurveto{\pgfqpoint{9.399391in}{4.612445in}}{\pgfqpoint{9.409227in}{4.616519in}}{\pgfqpoint{9.416478in}{4.623769in}}%
\pgfpathcurveto{\pgfqpoint{9.423728in}{4.631020in}}{\pgfqpoint{9.427802in}{4.640855in}}{\pgfqpoint{9.427802in}{4.651110in}}%
\pgfpathcurveto{\pgfqpoint{9.427802in}{4.661364in}}{\pgfqpoint{9.423728in}{4.671199in}}{\pgfqpoint{9.416478in}{4.678450in}}%
\pgfpathcurveto{\pgfqpoint{9.409227in}{4.685701in}}{\pgfqpoint{9.399391in}{4.689775in}}{\pgfqpoint{9.389137in}{4.689775in}}%
\pgfpathcurveto{\pgfqpoint{9.378883in}{4.689775in}}{\pgfqpoint{9.369048in}{4.685701in}}{\pgfqpoint{9.361797in}{4.678450in}}%
\pgfpathcurveto{\pgfqpoint{9.354546in}{4.671199in}}{\pgfqpoint{9.350472in}{4.661364in}}{\pgfqpoint{9.350472in}{4.651110in}}%
\pgfpathcurveto{\pgfqpoint{9.350472in}{4.640855in}}{\pgfqpoint{9.354546in}{4.631020in}}{\pgfqpoint{9.361797in}{4.623769in}}%
\pgfpathcurveto{\pgfqpoint{9.369048in}{4.616519in}}{\pgfqpoint{9.378883in}{4.612445in}}{\pgfqpoint{9.389137in}{4.612445in}}%
\pgfpathlineto{\pgfqpoint{9.389137in}{4.612445in}}%
\pgfpathclose%
\pgfusepath{stroke}%
\end{pgfscope}%
\begin{pgfscope}%
\pgfpathrectangle{\pgfqpoint{7.904091in}{2.928000in}}{\pgfqpoint{5.664000in}{5.664000in}}%
\pgfusepath{clip}%
\pgfsetbuttcap%
\pgfsetroundjoin%
\pgfsetlinewidth{1.003750pt}%
\definecolor{currentstroke}{rgb}{0.321569,0.321569,0.321569}%
\pgfsetstrokecolor{currentstroke}%
\pgfsetstrokeopacity{0.900000}%
\pgfsetdash{}{0pt}%
\pgfpathmoveto{\pgfqpoint{9.082725in}{4.263301in}}%
\pgfpathcurveto{\pgfqpoint{9.092979in}{4.263301in}}{\pgfqpoint{9.102815in}{4.267375in}}{\pgfqpoint{9.110065in}{4.274626in}}%
\pgfpathcurveto{\pgfqpoint{9.117316in}{4.281876in}}{\pgfqpoint{9.121390in}{4.291712in}}{\pgfqpoint{9.121390in}{4.301966in}}%
\pgfpathcurveto{\pgfqpoint{9.121390in}{4.312220in}}{\pgfqpoint{9.117316in}{4.322055in}}{\pgfqpoint{9.110065in}{4.329306in}}%
\pgfpathcurveto{\pgfqpoint{9.102815in}{4.336557in}}{\pgfqpoint{9.092979in}{4.340631in}}{\pgfqpoint{9.082725in}{4.340631in}}%
\pgfpathcurveto{\pgfqpoint{9.072471in}{4.340631in}}{\pgfqpoint{9.062636in}{4.336557in}}{\pgfqpoint{9.055385in}{4.329306in}}%
\pgfpathcurveto{\pgfqpoint{9.048134in}{4.322055in}}{\pgfqpoint{9.044060in}{4.312220in}}{\pgfqpoint{9.044060in}{4.301966in}}%
\pgfpathcurveto{\pgfqpoint{9.044060in}{4.291712in}}{\pgfqpoint{9.048134in}{4.281876in}}{\pgfqpoint{9.055385in}{4.274626in}}%
\pgfpathcurveto{\pgfqpoint{9.062636in}{4.267375in}}{\pgfqpoint{9.072471in}{4.263301in}}{\pgfqpoint{9.082725in}{4.263301in}}%
\pgfpathlineto{\pgfqpoint{9.082725in}{4.263301in}}%
\pgfpathclose%
\pgfusepath{stroke}%
\end{pgfscope}%
\begin{pgfscope}%
\pgfpathrectangle{\pgfqpoint{7.904091in}{2.928000in}}{\pgfqpoint{5.664000in}{5.664000in}}%
\pgfusepath{clip}%
\pgfsetbuttcap%
\pgfsetroundjoin%
\pgfsetlinewidth{1.003750pt}%
\definecolor{currentstroke}{rgb}{0.321569,0.321569,0.321569}%
\pgfsetstrokecolor{currentstroke}%
\pgfsetstrokeopacity{0.900000}%
\pgfsetdash{}{0pt}%
\pgfpathmoveto{\pgfqpoint{9.524233in}{4.471290in}}%
\pgfpathcurveto{\pgfqpoint{9.534487in}{4.471290in}}{\pgfqpoint{9.544322in}{4.475364in}}{\pgfqpoint{9.551573in}{4.482614in}}%
\pgfpathcurveto{\pgfqpoint{9.558824in}{4.489865in}}{\pgfqpoint{9.562898in}{4.499700in}}{\pgfqpoint{9.562898in}{4.509955in}}%
\pgfpathcurveto{\pgfqpoint{9.562898in}{4.520209in}}{\pgfqpoint{9.558824in}{4.530044in}}{\pgfqpoint{9.551573in}{4.537295in}}%
\pgfpathcurveto{\pgfqpoint{9.544322in}{4.544546in}}{\pgfqpoint{9.534487in}{4.548620in}}{\pgfqpoint{9.524233in}{4.548620in}}%
\pgfpathcurveto{\pgfqpoint{9.513979in}{4.548620in}}{\pgfqpoint{9.504143in}{4.544546in}}{\pgfqpoint{9.496893in}{4.537295in}}%
\pgfpathcurveto{\pgfqpoint{9.489642in}{4.530044in}}{\pgfqpoint{9.485568in}{4.520209in}}{\pgfqpoint{9.485568in}{4.509955in}}%
\pgfpathcurveto{\pgfqpoint{9.485568in}{4.499700in}}{\pgfqpoint{9.489642in}{4.489865in}}{\pgfqpoint{9.496893in}{4.482614in}}%
\pgfpathcurveto{\pgfqpoint{9.504143in}{4.475364in}}{\pgfqpoint{9.513979in}{4.471290in}}{\pgfqpoint{9.524233in}{4.471290in}}%
\pgfpathlineto{\pgfqpoint{9.524233in}{4.471290in}}%
\pgfpathclose%
\pgfusepath{stroke}%
\end{pgfscope}%
\begin{pgfscope}%
\pgfpathrectangle{\pgfqpoint{7.904091in}{2.928000in}}{\pgfqpoint{5.664000in}{5.664000in}}%
\pgfusepath{clip}%
\pgfsetbuttcap%
\pgfsetroundjoin%
\pgfsetlinewidth{1.003750pt}%
\definecolor{currentstroke}{rgb}{0.321569,0.321569,0.321569}%
\pgfsetstrokecolor{currentstroke}%
\pgfsetstrokeopacity{0.900000}%
\pgfsetdash{}{0pt}%
\pgfpathmoveto{\pgfqpoint{9.966819in}{4.957491in}}%
\pgfpathcurveto{\pgfqpoint{9.977073in}{4.957491in}}{\pgfqpoint{9.986908in}{4.961565in}}{\pgfqpoint{9.994159in}{4.968816in}}%
\pgfpathcurveto{\pgfqpoint{10.001410in}{4.976067in}}{\pgfqpoint{10.005484in}{4.985902in}}{\pgfqpoint{10.005484in}{4.996156in}}%
\pgfpathcurveto{\pgfqpoint{10.005484in}{5.006410in}}{\pgfqpoint{10.001410in}{5.016246in}}{\pgfqpoint{9.994159in}{5.023497in}}%
\pgfpathcurveto{\pgfqpoint{9.986908in}{5.030747in}}{\pgfqpoint{9.977073in}{5.034821in}}{\pgfqpoint{9.966819in}{5.034821in}}%
\pgfpathcurveto{\pgfqpoint{9.956564in}{5.034821in}}{\pgfqpoint{9.946729in}{5.030747in}}{\pgfqpoint{9.939478in}{5.023497in}}%
\pgfpathcurveto{\pgfqpoint{9.932227in}{5.016246in}}{\pgfqpoint{9.928153in}{5.006410in}}{\pgfqpoint{9.928153in}{4.996156in}}%
\pgfpathcurveto{\pgfqpoint{9.928153in}{4.985902in}}{\pgfqpoint{9.932227in}{4.976067in}}{\pgfqpoint{9.939478in}{4.968816in}}%
\pgfpathcurveto{\pgfqpoint{9.946729in}{4.961565in}}{\pgfqpoint{9.956564in}{4.957491in}}{\pgfqpoint{9.966819in}{4.957491in}}%
\pgfpathlineto{\pgfqpoint{9.966819in}{4.957491in}}%
\pgfpathclose%
\pgfusepath{stroke}%
\end{pgfscope}%
\begin{pgfscope}%
\pgfpathrectangle{\pgfqpoint{7.904091in}{2.928000in}}{\pgfqpoint{5.664000in}{5.664000in}}%
\pgfusepath{clip}%
\pgfsetbuttcap%
\pgfsetroundjoin%
\pgfsetlinewidth{1.003750pt}%
\definecolor{currentstroke}{rgb}{0.321569,0.321569,0.321569}%
\pgfsetstrokecolor{currentstroke}%
\pgfsetstrokeopacity{0.900000}%
\pgfsetdash{}{0pt}%
\pgfpathmoveto{\pgfqpoint{8.827171in}{4.194714in}}%
\pgfpathcurveto{\pgfqpoint{8.837425in}{4.194714in}}{\pgfqpoint{8.847260in}{4.198788in}}{\pgfqpoint{8.854511in}{4.206039in}}%
\pgfpathcurveto{\pgfqpoint{8.861762in}{4.213290in}}{\pgfqpoint{8.865836in}{4.223125in}}{\pgfqpoint{8.865836in}{4.233379in}}%
\pgfpathcurveto{\pgfqpoint{8.865836in}{4.243633in}}{\pgfqpoint{8.861762in}{4.253469in}}{\pgfqpoint{8.854511in}{4.260719in}}%
\pgfpathcurveto{\pgfqpoint{8.847260in}{4.267970in}}{\pgfqpoint{8.837425in}{4.272044in}}{\pgfqpoint{8.827171in}{4.272044in}}%
\pgfpathcurveto{\pgfqpoint{8.816917in}{4.272044in}}{\pgfqpoint{8.807081in}{4.267970in}}{\pgfqpoint{8.799830in}{4.260719in}}%
\pgfpathcurveto{\pgfqpoint{8.792580in}{4.253469in}}{\pgfqpoint{8.788506in}{4.243633in}}{\pgfqpoint{8.788506in}{4.233379in}}%
\pgfpathcurveto{\pgfqpoint{8.788506in}{4.223125in}}{\pgfqpoint{8.792580in}{4.213290in}}{\pgfqpoint{8.799830in}{4.206039in}}%
\pgfpathcurveto{\pgfqpoint{8.807081in}{4.198788in}}{\pgfqpoint{8.816917in}{4.194714in}}{\pgfqpoint{8.827171in}{4.194714in}}%
\pgfpathlineto{\pgfqpoint{8.827171in}{4.194714in}}%
\pgfpathclose%
\pgfusepath{stroke}%
\end{pgfscope}%
\begin{pgfscope}%
\pgfpathrectangle{\pgfqpoint{7.904091in}{2.928000in}}{\pgfqpoint{5.664000in}{5.664000in}}%
\pgfusepath{clip}%
\pgfsetbuttcap%
\pgfsetroundjoin%
\pgfsetlinewidth{1.003750pt}%
\definecolor{currentstroke}{rgb}{0.321569,0.321569,0.321569}%
\pgfsetstrokecolor{currentstroke}%
\pgfsetstrokeopacity{0.900000}%
\pgfsetdash{}{0pt}%
\pgfpathmoveto{\pgfqpoint{11.482403in}{6.381648in}}%
\pgfpathcurveto{\pgfqpoint{11.492657in}{6.381648in}}{\pgfqpoint{11.502492in}{6.385722in}}{\pgfqpoint{11.509743in}{6.392973in}}%
\pgfpathcurveto{\pgfqpoint{11.516994in}{6.400223in}}{\pgfqpoint{11.521068in}{6.410059in}}{\pgfqpoint{11.521068in}{6.420313in}}%
\pgfpathcurveto{\pgfqpoint{11.521068in}{6.430567in}}{\pgfqpoint{11.516994in}{6.440403in}}{\pgfqpoint{11.509743in}{6.447653in}}%
\pgfpathcurveto{\pgfqpoint{11.502492in}{6.454904in}}{\pgfqpoint{11.492657in}{6.458978in}}{\pgfqpoint{11.482403in}{6.458978in}}%
\pgfpathcurveto{\pgfqpoint{11.472149in}{6.458978in}}{\pgfqpoint{11.462313in}{6.454904in}}{\pgfqpoint{11.455063in}{6.447653in}}%
\pgfpathcurveto{\pgfqpoint{11.447812in}{6.440403in}}{\pgfqpoint{11.443738in}{6.430567in}}{\pgfqpoint{11.443738in}{6.420313in}}%
\pgfpathcurveto{\pgfqpoint{11.443738in}{6.410059in}}{\pgfqpoint{11.447812in}{6.400223in}}{\pgfqpoint{11.455063in}{6.392973in}}%
\pgfpathcurveto{\pgfqpoint{11.462313in}{6.385722in}}{\pgfqpoint{11.472149in}{6.381648in}}{\pgfqpoint{11.482403in}{6.381648in}}%
\pgfpathlineto{\pgfqpoint{11.482403in}{6.381648in}}%
\pgfpathclose%
\pgfusepath{stroke}%
\end{pgfscope}%
\begin{pgfscope}%
\pgfpathrectangle{\pgfqpoint{7.904091in}{2.928000in}}{\pgfqpoint{5.664000in}{5.664000in}}%
\pgfusepath{clip}%
\pgfsetbuttcap%
\pgfsetroundjoin%
\pgfsetlinewidth{1.003750pt}%
\definecolor{currentstroke}{rgb}{0.321569,0.321569,0.321569}%
\pgfsetstrokecolor{currentstroke}%
\pgfsetstrokeopacity{0.900000}%
\pgfsetdash{}{0pt}%
\pgfpathmoveto{\pgfqpoint{11.625314in}{6.491428in}}%
\pgfpathcurveto{\pgfqpoint{11.635568in}{6.491428in}}{\pgfqpoint{11.645404in}{6.495502in}}{\pgfqpoint{11.652655in}{6.502753in}}%
\pgfpathcurveto{\pgfqpoint{11.659905in}{6.510004in}}{\pgfqpoint{11.663979in}{6.519839in}}{\pgfqpoint{11.663979in}{6.530093in}}%
\pgfpathcurveto{\pgfqpoint{11.663979in}{6.540347in}}{\pgfqpoint{11.659905in}{6.550183in}}{\pgfqpoint{11.652655in}{6.557434in}}%
\pgfpathcurveto{\pgfqpoint{11.645404in}{6.564684in}}{\pgfqpoint{11.635568in}{6.568758in}}{\pgfqpoint{11.625314in}{6.568758in}}%
\pgfpathcurveto{\pgfqpoint{11.615060in}{6.568758in}}{\pgfqpoint{11.605225in}{6.564684in}}{\pgfqpoint{11.597974in}{6.557434in}}%
\pgfpathcurveto{\pgfqpoint{11.590723in}{6.550183in}}{\pgfqpoint{11.586649in}{6.540347in}}{\pgfqpoint{11.586649in}{6.530093in}}%
\pgfpathcurveto{\pgfqpoint{11.586649in}{6.519839in}}{\pgfqpoint{11.590723in}{6.510004in}}{\pgfqpoint{11.597974in}{6.502753in}}%
\pgfpathcurveto{\pgfqpoint{11.605225in}{6.495502in}}{\pgfqpoint{11.615060in}{6.491428in}}{\pgfqpoint{11.625314in}{6.491428in}}%
\pgfpathlineto{\pgfqpoint{11.625314in}{6.491428in}}%
\pgfpathclose%
\pgfusepath{stroke}%
\end{pgfscope}%
\begin{pgfscope}%
\pgfpathrectangle{\pgfqpoint{7.904091in}{2.928000in}}{\pgfqpoint{5.664000in}{5.664000in}}%
\pgfusepath{clip}%
\pgfsetbuttcap%
\pgfsetroundjoin%
\pgfsetlinewidth{1.003750pt}%
\definecolor{currentstroke}{rgb}{0.321569,0.321569,0.321569}%
\pgfsetstrokecolor{currentstroke}%
\pgfsetstrokeopacity{0.900000}%
\pgfsetdash{}{0pt}%
\pgfpathmoveto{\pgfqpoint{11.740454in}{6.591594in}}%
\pgfpathcurveto{\pgfqpoint{11.750709in}{6.591594in}}{\pgfqpoint{11.760544in}{6.595668in}}{\pgfqpoint{11.767795in}{6.602919in}}%
\pgfpathcurveto{\pgfqpoint{11.775046in}{6.610169in}}{\pgfqpoint{11.779119in}{6.620005in}}{\pgfqpoint{11.779119in}{6.630259in}}%
\pgfpathcurveto{\pgfqpoint{11.779119in}{6.640513in}}{\pgfqpoint{11.775046in}{6.650348in}}{\pgfqpoint{11.767795in}{6.657599in}}%
\pgfpathcurveto{\pgfqpoint{11.760544in}{6.664850in}}{\pgfqpoint{11.750709in}{6.668924in}}{\pgfqpoint{11.740454in}{6.668924in}}%
\pgfpathcurveto{\pgfqpoint{11.730200in}{6.668924in}}{\pgfqpoint{11.720365in}{6.664850in}}{\pgfqpoint{11.713114in}{6.657599in}}%
\pgfpathcurveto{\pgfqpoint{11.705863in}{6.650348in}}{\pgfqpoint{11.701789in}{6.640513in}}{\pgfqpoint{11.701789in}{6.630259in}}%
\pgfpathcurveto{\pgfqpoint{11.701789in}{6.620005in}}{\pgfqpoint{11.705863in}{6.610169in}}{\pgfqpoint{11.713114in}{6.602919in}}%
\pgfpathcurveto{\pgfqpoint{11.720365in}{6.595668in}}{\pgfqpoint{11.730200in}{6.591594in}}{\pgfqpoint{11.740454in}{6.591594in}}%
\pgfpathlineto{\pgfqpoint{11.740454in}{6.591594in}}%
\pgfpathclose%
\pgfusepath{stroke}%
\end{pgfscope}%
\begin{pgfscope}%
\pgfpathrectangle{\pgfqpoint{7.904091in}{2.928000in}}{\pgfqpoint{5.664000in}{5.664000in}}%
\pgfusepath{clip}%
\pgfsetbuttcap%
\pgfsetroundjoin%
\pgfsetlinewidth{1.003750pt}%
\definecolor{currentstroke}{rgb}{0.321569,0.321569,0.321569}%
\pgfsetstrokecolor{currentstroke}%
\pgfsetstrokeopacity{0.900000}%
\pgfsetdash{}{0pt}%
\pgfpathmoveto{\pgfqpoint{11.956847in}{6.859985in}}%
\pgfpathcurveto{\pgfqpoint{11.967101in}{6.859985in}}{\pgfqpoint{11.976937in}{6.864059in}}{\pgfqpoint{11.984188in}{6.871309in}}%
\pgfpathcurveto{\pgfqpoint{11.991438in}{6.878560in}}{\pgfqpoint{11.995512in}{6.888396in}}{\pgfqpoint{11.995512in}{6.898650in}}%
\pgfpathcurveto{\pgfqpoint{11.995512in}{6.908904in}}{\pgfqpoint{11.991438in}{6.918739in}}{\pgfqpoint{11.984188in}{6.925990in}}%
\pgfpathcurveto{\pgfqpoint{11.976937in}{6.933241in}}{\pgfqpoint{11.967101in}{6.937315in}}{\pgfqpoint{11.956847in}{6.937315in}}%
\pgfpathcurveto{\pgfqpoint{11.946593in}{6.937315in}}{\pgfqpoint{11.936758in}{6.933241in}}{\pgfqpoint{11.929507in}{6.925990in}}%
\pgfpathcurveto{\pgfqpoint{11.922256in}{6.918739in}}{\pgfqpoint{11.918182in}{6.908904in}}{\pgfqpoint{11.918182in}{6.898650in}}%
\pgfpathcurveto{\pgfqpoint{11.918182in}{6.888396in}}{\pgfqpoint{11.922256in}{6.878560in}}{\pgfqpoint{11.929507in}{6.871309in}}%
\pgfpathcurveto{\pgfqpoint{11.936758in}{6.864059in}}{\pgfqpoint{11.946593in}{6.859985in}}{\pgfqpoint{11.956847in}{6.859985in}}%
\pgfpathlineto{\pgfqpoint{11.956847in}{6.859985in}}%
\pgfpathclose%
\pgfusepath{stroke}%
\end{pgfscope}%
\begin{pgfscope}%
\pgfsetbuttcap%
\pgfsetroundjoin%
\pgfsetlinewidth{1.003750pt}%
\definecolor{currentstroke}{rgb}{0.321569,0.321569,0.321569}%
\pgfsetstrokecolor{currentstroke}%
\pgfsetstrokeopacity{0.900000}%
\pgfsetdash{}{0pt}%
\pgfpathmoveto{\pgfqpoint{7.904091in}{2.868262in}}%
\pgfpathlineto{\pgfqpoint{7.963829in}{2.928000in}}%
\pgfpathlineto{\pgfqpoint{7.904091in}{2.987738in}}%
\pgfpathlineto{\pgfqpoint{7.844353in}{2.928000in}}%
\pgfpathlineto{\pgfqpoint{7.904091in}{2.868262in}}%
\pgfpathclose%
\pgfusepath{stroke}%
\end{pgfscope}%
\begin{pgfscope}%
\pgfsetbuttcap%
\pgfsetroundjoin%
\pgfsetlinewidth{1.003750pt}%
\definecolor{currentstroke}{rgb}{0.321569,0.321569,0.321569}%
\pgfsetstrokecolor{currentstroke}%
\pgfsetstrokeopacity{0.900000}%
\pgfsetdash{}{0pt}%
\pgfpathmoveto{\pgfqpoint{7.904091in}{2.868262in}}%
\pgfpathlineto{\pgfqpoint{7.963829in}{2.928000in}}%
\pgfpathlineto{\pgfqpoint{7.904091in}{2.987738in}}%
\pgfpathlineto{\pgfqpoint{7.844353in}{2.928000in}}%
\pgfpathlineto{\pgfqpoint{7.904091in}{2.868262in}}%
\pgfpathclose%
\pgfusepath{stroke}%
\end{pgfscope}%
\begin{pgfscope}%
\pgfsetbuttcap%
\pgfsetroundjoin%
\pgfsetlinewidth{1.003750pt}%
\definecolor{currentstroke}{rgb}{0.321569,0.321569,0.321569}%
\pgfsetstrokecolor{currentstroke}%
\pgfsetstrokeopacity{0.900000}%
\pgfsetdash{}{0pt}%
\pgfpathmoveto{\pgfqpoint{7.904091in}{2.868262in}}%
\pgfpathlineto{\pgfqpoint{7.963829in}{2.928000in}}%
\pgfpathlineto{\pgfqpoint{7.904091in}{2.987738in}}%
\pgfpathlineto{\pgfqpoint{7.844353in}{2.928000in}}%
\pgfpathlineto{\pgfqpoint{7.904091in}{2.868262in}}%
\pgfpathclose%
\pgfusepath{stroke}%
\end{pgfscope}%
\begin{pgfscope}%
\pgfsetbuttcap%
\pgfsetroundjoin%
\pgfsetlinewidth{1.003750pt}%
\definecolor{currentstroke}{rgb}{0.321569,0.321569,0.321569}%
\pgfsetstrokecolor{currentstroke}%
\pgfsetstrokeopacity{0.900000}%
\pgfsetdash{}{0pt}%
\pgfpathmoveto{\pgfqpoint{7.904091in}{2.868262in}}%
\pgfpathlineto{\pgfqpoint{7.963829in}{2.928000in}}%
\pgfpathlineto{\pgfqpoint{7.904091in}{2.987738in}}%
\pgfpathlineto{\pgfqpoint{7.844353in}{2.928000in}}%
\pgfpathlineto{\pgfqpoint{7.904091in}{2.868262in}}%
\pgfpathclose%
\pgfusepath{stroke}%
\end{pgfscope}%
\begin{pgfscope}%
\pgfsetbuttcap%
\pgfsetroundjoin%
\pgfsetlinewidth{1.003750pt}%
\definecolor{currentstroke}{rgb}{0.321569,0.321569,0.321569}%
\pgfsetstrokecolor{currentstroke}%
\pgfsetstrokeopacity{0.900000}%
\pgfsetdash{}{0pt}%
\pgfpathmoveto{\pgfqpoint{7.904091in}{2.868262in}}%
\pgfpathlineto{\pgfqpoint{7.963829in}{2.928000in}}%
\pgfpathlineto{\pgfqpoint{7.904091in}{2.987738in}}%
\pgfpathlineto{\pgfqpoint{7.844353in}{2.928000in}}%
\pgfpathlineto{\pgfqpoint{7.904091in}{2.868262in}}%
\pgfpathclose%
\pgfusepath{stroke}%
\end{pgfscope}%
\begin{pgfscope}%
\pgfsetbuttcap%
\pgfsetroundjoin%
\pgfsetlinewidth{1.003750pt}%
\definecolor{currentstroke}{rgb}{0.321569,0.321569,0.321569}%
\pgfsetstrokecolor{currentstroke}%
\pgfsetstrokeopacity{0.900000}%
\pgfsetdash{}{0pt}%
\pgfpathmoveto{\pgfqpoint{7.904091in}{2.868262in}}%
\pgfpathlineto{\pgfqpoint{7.963829in}{2.928000in}}%
\pgfpathlineto{\pgfqpoint{7.904091in}{2.987738in}}%
\pgfpathlineto{\pgfqpoint{7.844353in}{2.928000in}}%
\pgfpathlineto{\pgfqpoint{7.904091in}{2.868262in}}%
\pgfpathclose%
\pgfusepath{stroke}%
\end{pgfscope}%
\begin{pgfscope}%
\pgfpathrectangle{\pgfqpoint{7.904091in}{2.928000in}}{\pgfqpoint{5.664000in}{5.664000in}}%
\pgfusepath{clip}%
\pgfsetbuttcap%
\pgfsetroundjoin%
\pgfsetlinewidth{1.003750pt}%
\definecolor{currentstroke}{rgb}{0.176471,0.713725,0.643137}%
\pgfsetstrokecolor{currentstroke}%
\pgfsetstrokeopacity{0.900000}%
\pgfsetdash{}{0pt}%
\pgfpathmoveto{\pgfqpoint{12.934901in}{8.285454in}}%
\pgfpathcurveto{\pgfqpoint{12.945155in}{8.285454in}}{\pgfqpoint{12.954990in}{8.289528in}}{\pgfqpoint{12.962241in}{8.296779in}}%
\pgfpathcurveto{\pgfqpoint{12.969492in}{8.304029in}}{\pgfqpoint{12.973566in}{8.313865in}}{\pgfqpoint{12.973566in}{8.324119in}}%
\pgfpathcurveto{\pgfqpoint{12.973566in}{8.334373in}}{\pgfqpoint{12.969492in}{8.344209in}}{\pgfqpoint{12.962241in}{8.351459in}}%
\pgfpathcurveto{\pgfqpoint{12.954990in}{8.358710in}}{\pgfqpoint{12.945155in}{8.362784in}}{\pgfqpoint{12.934901in}{8.362784in}}%
\pgfpathcurveto{\pgfqpoint{12.924647in}{8.362784in}}{\pgfqpoint{12.914811in}{8.358710in}}{\pgfqpoint{12.907561in}{8.351459in}}%
\pgfpathcurveto{\pgfqpoint{12.900310in}{8.344209in}}{\pgfqpoint{12.896236in}{8.334373in}}{\pgfqpoint{12.896236in}{8.324119in}}%
\pgfpathcurveto{\pgfqpoint{12.896236in}{8.313865in}}{\pgfqpoint{12.900310in}{8.304029in}}{\pgfqpoint{12.907561in}{8.296779in}}%
\pgfpathcurveto{\pgfqpoint{12.914811in}{8.289528in}}{\pgfqpoint{12.924647in}{8.285454in}}{\pgfqpoint{12.934901in}{8.285454in}}%
\pgfpathlineto{\pgfqpoint{12.934901in}{8.285454in}}%
\pgfpathclose%
\pgfusepath{stroke}%
\end{pgfscope}%
\begin{pgfscope}%
\pgfpathrectangle{\pgfqpoint{7.904091in}{2.928000in}}{\pgfqpoint{5.664000in}{5.664000in}}%
\pgfusepath{clip}%
\pgfsetbuttcap%
\pgfsetroundjoin%
\pgfsetlinewidth{1.003750pt}%
\definecolor{currentstroke}{rgb}{0.176471,0.713725,0.643137}%
\pgfsetstrokecolor{currentstroke}%
\pgfsetstrokeopacity{0.900000}%
\pgfsetdash{}{0pt}%
\pgfpathmoveto{\pgfqpoint{12.909857in}{8.244449in}}%
\pgfpathcurveto{\pgfqpoint{12.920111in}{8.244449in}}{\pgfqpoint{12.929947in}{8.248523in}}{\pgfqpoint{12.937198in}{8.255774in}}%
\pgfpathcurveto{\pgfqpoint{12.944448in}{8.263025in}}{\pgfqpoint{12.948522in}{8.272860in}}{\pgfqpoint{12.948522in}{8.283114in}}%
\pgfpathcurveto{\pgfqpoint{12.948522in}{8.293368in}}{\pgfqpoint{12.944448in}{8.303204in}}{\pgfqpoint{12.937198in}{8.310455in}}%
\pgfpathcurveto{\pgfqpoint{12.929947in}{8.317705in}}{\pgfqpoint{12.920111in}{8.321779in}}{\pgfqpoint{12.909857in}{8.321779in}}%
\pgfpathcurveto{\pgfqpoint{12.899603in}{8.321779in}}{\pgfqpoint{12.889768in}{8.317705in}}{\pgfqpoint{12.882517in}{8.310455in}}%
\pgfpathcurveto{\pgfqpoint{12.875266in}{8.303204in}}{\pgfqpoint{12.871192in}{8.293368in}}{\pgfqpoint{12.871192in}{8.283114in}}%
\pgfpathcurveto{\pgfqpoint{12.871192in}{8.272860in}}{\pgfqpoint{12.875266in}{8.263025in}}{\pgfqpoint{12.882517in}{8.255774in}}%
\pgfpathcurveto{\pgfqpoint{12.889768in}{8.248523in}}{\pgfqpoint{12.899603in}{8.244449in}}{\pgfqpoint{12.909857in}{8.244449in}}%
\pgfpathlineto{\pgfqpoint{12.909857in}{8.244449in}}%
\pgfpathclose%
\pgfusepath{stroke}%
\end{pgfscope}%
\begin{pgfscope}%
\pgfpathrectangle{\pgfqpoint{7.904091in}{2.928000in}}{\pgfqpoint{5.664000in}{5.664000in}}%
\pgfusepath{clip}%
\pgfsetbuttcap%
\pgfsetroundjoin%
\pgfsetlinewidth{1.003750pt}%
\definecolor{currentstroke}{rgb}{0.176471,0.713725,0.643137}%
\pgfsetstrokecolor{currentstroke}%
\pgfsetstrokeopacity{0.900000}%
\pgfsetdash{}{0pt}%
\pgfpathmoveto{\pgfqpoint{12.970615in}{8.239079in}}%
\pgfpathcurveto{\pgfqpoint{12.980869in}{8.239079in}}{\pgfqpoint{12.990704in}{8.243153in}}{\pgfqpoint{12.997955in}{8.250404in}}%
\pgfpathcurveto{\pgfqpoint{13.005206in}{8.257655in}}{\pgfqpoint{13.009280in}{8.267490in}}{\pgfqpoint{13.009280in}{8.277744in}}%
\pgfpathcurveto{\pgfqpoint{13.009280in}{8.287998in}}{\pgfqpoint{13.005206in}{8.297834in}}{\pgfqpoint{12.997955in}{8.305085in}}%
\pgfpathcurveto{\pgfqpoint{12.990704in}{8.312335in}}{\pgfqpoint{12.980869in}{8.316409in}}{\pgfqpoint{12.970615in}{8.316409in}}%
\pgfpathcurveto{\pgfqpoint{12.960361in}{8.316409in}}{\pgfqpoint{12.950525in}{8.312335in}}{\pgfqpoint{12.943274in}{8.305085in}}%
\pgfpathcurveto{\pgfqpoint{12.936024in}{8.297834in}}{\pgfqpoint{12.931950in}{8.287998in}}{\pgfqpoint{12.931950in}{8.277744in}}%
\pgfpathcurveto{\pgfqpoint{12.931950in}{8.267490in}}{\pgfqpoint{12.936024in}{8.257655in}}{\pgfqpoint{12.943274in}{8.250404in}}%
\pgfpathcurveto{\pgfqpoint{12.950525in}{8.243153in}}{\pgfqpoint{12.960361in}{8.239079in}}{\pgfqpoint{12.970615in}{8.239079in}}%
\pgfpathlineto{\pgfqpoint{12.970615in}{8.239079in}}%
\pgfpathclose%
\pgfusepath{stroke}%
\end{pgfscope}%
\begin{pgfscope}%
\pgfpathrectangle{\pgfqpoint{7.904091in}{2.928000in}}{\pgfqpoint{5.664000in}{5.664000in}}%
\pgfusepath{clip}%
\pgfsetbuttcap%
\pgfsetroundjoin%
\pgfsetlinewidth{1.003750pt}%
\definecolor{currentstroke}{rgb}{0.176471,0.713725,0.643137}%
\pgfsetstrokecolor{currentstroke}%
\pgfsetstrokeopacity{0.900000}%
\pgfsetdash{}{0pt}%
\pgfpathmoveto{\pgfqpoint{13.221283in}{8.183660in}}%
\pgfpathcurveto{\pgfqpoint{13.231537in}{8.183660in}}{\pgfqpoint{13.241372in}{8.187734in}}{\pgfqpoint{13.248623in}{8.194984in}}%
\pgfpathcurveto{\pgfqpoint{13.255874in}{8.202235in}}{\pgfqpoint{13.259948in}{8.212071in}}{\pgfqpoint{13.259948in}{8.222325in}}%
\pgfpathcurveto{\pgfqpoint{13.259948in}{8.232579in}}{\pgfqpoint{13.255874in}{8.242414in}}{\pgfqpoint{13.248623in}{8.249665in}}%
\pgfpathcurveto{\pgfqpoint{13.241372in}{8.256916in}}{\pgfqpoint{13.231537in}{8.260990in}}{\pgfqpoint{13.221283in}{8.260990in}}%
\pgfpathcurveto{\pgfqpoint{13.211029in}{8.260990in}}{\pgfqpoint{13.201193in}{8.256916in}}{\pgfqpoint{13.193942in}{8.249665in}}%
\pgfpathcurveto{\pgfqpoint{13.186692in}{8.242414in}}{\pgfqpoint{13.182618in}{8.232579in}}{\pgfqpoint{13.182618in}{8.222325in}}%
\pgfpathcurveto{\pgfqpoint{13.182618in}{8.212071in}}{\pgfqpoint{13.186692in}{8.202235in}}{\pgfqpoint{13.193942in}{8.194984in}}%
\pgfpathcurveto{\pgfqpoint{13.201193in}{8.187734in}}{\pgfqpoint{13.211029in}{8.183660in}}{\pgfqpoint{13.221283in}{8.183660in}}%
\pgfpathlineto{\pgfqpoint{13.221283in}{8.183660in}}%
\pgfpathclose%
\pgfusepath{stroke}%
\end{pgfscope}%
\begin{pgfscope}%
\pgfpathrectangle{\pgfqpoint{7.904091in}{2.928000in}}{\pgfqpoint{5.664000in}{5.664000in}}%
\pgfusepath{clip}%
\pgfsetbuttcap%
\pgfsetroundjoin%
\pgfsetlinewidth{1.003750pt}%
\definecolor{currentstroke}{rgb}{0.176471,0.713725,0.643137}%
\pgfsetstrokecolor{currentstroke}%
\pgfsetstrokeopacity{0.900000}%
\pgfsetdash{}{0pt}%
\pgfpathmoveto{\pgfqpoint{12.685630in}{8.002372in}}%
\pgfpathcurveto{\pgfqpoint{12.695885in}{8.002372in}}{\pgfqpoint{12.705720in}{8.006446in}}{\pgfqpoint{12.712971in}{8.013697in}}%
\pgfpathcurveto{\pgfqpoint{12.720222in}{8.020947in}}{\pgfqpoint{12.724296in}{8.030783in}}{\pgfqpoint{12.724296in}{8.041037in}}%
\pgfpathcurveto{\pgfqpoint{12.724296in}{8.051291in}}{\pgfqpoint{12.720222in}{8.061126in}}{\pgfqpoint{12.712971in}{8.068377in}}%
\pgfpathcurveto{\pgfqpoint{12.705720in}{8.075628in}}{\pgfqpoint{12.695885in}{8.079702in}}{\pgfqpoint{12.685630in}{8.079702in}}%
\pgfpathcurveto{\pgfqpoint{12.675376in}{8.079702in}}{\pgfqpoint{12.665541in}{8.075628in}}{\pgfqpoint{12.658290in}{8.068377in}}%
\pgfpathcurveto{\pgfqpoint{12.651039in}{8.061126in}}{\pgfqpoint{12.646965in}{8.051291in}}{\pgfqpoint{12.646965in}{8.041037in}}%
\pgfpathcurveto{\pgfqpoint{12.646965in}{8.030783in}}{\pgfqpoint{12.651039in}{8.020947in}}{\pgfqpoint{12.658290in}{8.013697in}}%
\pgfpathcurveto{\pgfqpoint{12.665541in}{8.006446in}}{\pgfqpoint{12.675376in}{8.002372in}}{\pgfqpoint{12.685630in}{8.002372in}}%
\pgfpathlineto{\pgfqpoint{12.685630in}{8.002372in}}%
\pgfpathclose%
\pgfusepath{stroke}%
\end{pgfscope}%
\begin{pgfscope}%
\pgfpathrectangle{\pgfqpoint{7.904091in}{2.928000in}}{\pgfqpoint{5.664000in}{5.664000in}}%
\pgfusepath{clip}%
\pgfsetbuttcap%
\pgfsetroundjoin%
\pgfsetlinewidth{1.003750pt}%
\definecolor{currentstroke}{rgb}{0.000000,0.447059,0.698039}%
\pgfsetstrokecolor{currentstroke}%
\pgfsetstrokeopacity{0.900000}%
\pgfsetdash{}{0pt}%
\pgfpathmoveto{\pgfqpoint{10.644748in}{5.706285in}}%
\pgfpathcurveto{\pgfqpoint{10.655002in}{5.706285in}}{\pgfqpoint{10.664837in}{5.710359in}}{\pgfqpoint{10.672088in}{5.717610in}}%
\pgfpathcurveto{\pgfqpoint{10.679339in}{5.724861in}}{\pgfqpoint{10.683413in}{5.734696in}}{\pgfqpoint{10.683413in}{5.744950in}}%
\pgfpathcurveto{\pgfqpoint{10.683413in}{5.755205in}}{\pgfqpoint{10.679339in}{5.765040in}}{\pgfqpoint{10.672088in}{5.772291in}}%
\pgfpathcurveto{\pgfqpoint{10.664837in}{5.779541in}}{\pgfqpoint{10.655002in}{5.783615in}}{\pgfqpoint{10.644748in}{5.783615in}}%
\pgfpathcurveto{\pgfqpoint{10.634494in}{5.783615in}}{\pgfqpoint{10.624658in}{5.779541in}}{\pgfqpoint{10.617407in}{5.772291in}}%
\pgfpathcurveto{\pgfqpoint{10.610157in}{5.765040in}}{\pgfqpoint{10.606083in}{5.755205in}}{\pgfqpoint{10.606083in}{5.744950in}}%
\pgfpathcurveto{\pgfqpoint{10.606083in}{5.734696in}}{\pgfqpoint{10.610157in}{5.724861in}}{\pgfqpoint{10.617407in}{5.717610in}}%
\pgfpathcurveto{\pgfqpoint{10.624658in}{5.710359in}}{\pgfqpoint{10.634494in}{5.706285in}}{\pgfqpoint{10.644748in}{5.706285in}}%
\pgfpathlineto{\pgfqpoint{10.644748in}{5.706285in}}%
\pgfpathclose%
\pgfusepath{stroke}%
\end{pgfscope}%
\begin{pgfscope}%
\pgfpathrectangle{\pgfqpoint{7.904091in}{2.928000in}}{\pgfqpoint{5.664000in}{5.664000in}}%
\pgfusepath{clip}%
\pgfsetbuttcap%
\pgfsetroundjoin%
\pgfsetlinewidth{1.003750pt}%
\definecolor{currentstroke}{rgb}{0.000000,0.447059,0.698039}%
\pgfsetstrokecolor{currentstroke}%
\pgfsetstrokeopacity{0.900000}%
\pgfsetdash{}{0pt}%
\pgfpathmoveto{\pgfqpoint{10.670017in}{4.968920in}}%
\pgfpathcurveto{\pgfqpoint{10.680271in}{4.968920in}}{\pgfqpoint{10.690106in}{4.972994in}}{\pgfqpoint{10.697357in}{4.980245in}}%
\pgfpathcurveto{\pgfqpoint{10.704608in}{4.987496in}}{\pgfqpoint{10.708682in}{4.997331in}}{\pgfqpoint{10.708682in}{5.007585in}}%
\pgfpathcurveto{\pgfqpoint{10.708682in}{5.017839in}}{\pgfqpoint{10.704608in}{5.027675in}}{\pgfqpoint{10.697357in}{5.034926in}}%
\pgfpathcurveto{\pgfqpoint{10.690106in}{5.042176in}}{\pgfqpoint{10.680271in}{5.046250in}}{\pgfqpoint{10.670017in}{5.046250in}}%
\pgfpathcurveto{\pgfqpoint{10.659762in}{5.046250in}}{\pgfqpoint{10.649927in}{5.042176in}}{\pgfqpoint{10.642676in}{5.034926in}}%
\pgfpathcurveto{\pgfqpoint{10.635425in}{5.027675in}}{\pgfqpoint{10.631351in}{5.017839in}}{\pgfqpoint{10.631351in}{5.007585in}}%
\pgfpathcurveto{\pgfqpoint{10.631351in}{4.997331in}}{\pgfqpoint{10.635425in}{4.987496in}}{\pgfqpoint{10.642676in}{4.980245in}}%
\pgfpathcurveto{\pgfqpoint{10.649927in}{4.972994in}}{\pgfqpoint{10.659762in}{4.968920in}}{\pgfqpoint{10.670017in}{4.968920in}}%
\pgfpathlineto{\pgfqpoint{10.670017in}{4.968920in}}%
\pgfpathclose%
\pgfusepath{stroke}%
\end{pgfscope}%
\begin{pgfscope}%
\pgfpathrectangle{\pgfqpoint{7.904091in}{2.928000in}}{\pgfqpoint{5.664000in}{5.664000in}}%
\pgfusepath{clip}%
\pgfsetbuttcap%
\pgfsetroundjoin%
\pgfsetlinewidth{1.003750pt}%
\definecolor{currentstroke}{rgb}{0.337255,0.705882,0.913725}%
\pgfsetstrokecolor{currentstroke}%
\pgfsetstrokeopacity{0.900000}%
\pgfsetdash{}{0pt}%
\pgfpathmoveto{\pgfqpoint{8.306979in}{3.476842in}}%
\pgfpathcurveto{\pgfqpoint{8.317234in}{3.476842in}}{\pgfqpoint{8.327069in}{3.480916in}}{\pgfqpoint{8.334320in}{3.488167in}}%
\pgfpathcurveto{\pgfqpoint{8.341571in}{3.495418in}}{\pgfqpoint{8.345645in}{3.505253in}}{\pgfqpoint{8.345645in}{3.515507in}}%
\pgfpathcurveto{\pgfqpoint{8.345645in}{3.525761in}}{\pgfqpoint{8.341571in}{3.535597in}}{\pgfqpoint{8.334320in}{3.542848in}}%
\pgfpathcurveto{\pgfqpoint{8.327069in}{3.550098in}}{\pgfqpoint{8.317234in}{3.554172in}}{\pgfqpoint{8.306979in}{3.554172in}}%
\pgfpathcurveto{\pgfqpoint{8.296725in}{3.554172in}}{\pgfqpoint{8.286890in}{3.550098in}}{\pgfqpoint{8.279639in}{3.542848in}}%
\pgfpathcurveto{\pgfqpoint{8.272388in}{3.535597in}}{\pgfqpoint{8.268314in}{3.525761in}}{\pgfqpoint{8.268314in}{3.515507in}}%
\pgfpathcurveto{\pgfqpoint{8.268314in}{3.505253in}}{\pgfqpoint{8.272388in}{3.495418in}}{\pgfqpoint{8.279639in}{3.488167in}}%
\pgfpathcurveto{\pgfqpoint{8.286890in}{3.480916in}}{\pgfqpoint{8.296725in}{3.476842in}}{\pgfqpoint{8.306979in}{3.476842in}}%
\pgfpathlineto{\pgfqpoint{8.306979in}{3.476842in}}%
\pgfpathclose%
\pgfusepath{stroke}%
\end{pgfscope}%
\begin{pgfscope}%
\pgfpathrectangle{\pgfqpoint{7.904091in}{2.928000in}}{\pgfqpoint{5.664000in}{5.664000in}}%
\pgfusepath{clip}%
\pgfsetbuttcap%
\pgfsetroundjoin%
\pgfsetlinewidth{1.003750pt}%
\definecolor{currentstroke}{rgb}{0.337255,0.705882,0.913725}%
\pgfsetstrokecolor{currentstroke}%
\pgfsetstrokeopacity{0.900000}%
\pgfsetdash{}{0pt}%
\pgfpathmoveto{\pgfqpoint{9.362400in}{4.089292in}}%
\pgfpathcurveto{\pgfqpoint{9.372654in}{4.089292in}}{\pgfqpoint{9.382489in}{4.093366in}}{\pgfqpoint{9.389740in}{4.100617in}}%
\pgfpathcurveto{\pgfqpoint{9.396991in}{4.107867in}}{\pgfqpoint{9.401065in}{4.117703in}}{\pgfqpoint{9.401065in}{4.127957in}}%
\pgfpathcurveto{\pgfqpoint{9.401065in}{4.138211in}}{\pgfqpoint{9.396991in}{4.148047in}}{\pgfqpoint{9.389740in}{4.155297in}}%
\pgfpathcurveto{\pgfqpoint{9.382489in}{4.162548in}}{\pgfqpoint{9.372654in}{4.166622in}}{\pgfqpoint{9.362400in}{4.166622in}}%
\pgfpathcurveto{\pgfqpoint{9.352146in}{4.166622in}}{\pgfqpoint{9.342310in}{4.162548in}}{\pgfqpoint{9.335059in}{4.155297in}}%
\pgfpathcurveto{\pgfqpoint{9.327809in}{4.148047in}}{\pgfqpoint{9.323735in}{4.138211in}}{\pgfqpoint{9.323735in}{4.127957in}}%
\pgfpathcurveto{\pgfqpoint{9.323735in}{4.117703in}}{\pgfqpoint{9.327809in}{4.107867in}}{\pgfqpoint{9.335059in}{4.100617in}}%
\pgfpathcurveto{\pgfqpoint{9.342310in}{4.093366in}}{\pgfqpoint{9.352146in}{4.089292in}}{\pgfqpoint{9.362400in}{4.089292in}}%
\pgfpathlineto{\pgfqpoint{9.362400in}{4.089292in}}%
\pgfpathclose%
\pgfusepath{stroke}%
\end{pgfscope}%
\begin{pgfscope}%
\pgfsetbuttcap%
\pgfsetroundjoin%
\pgfsetlinewidth{1.003750pt}%
\definecolor{currentstroke}{rgb}{0.337255,0.705882,0.913725}%
\pgfsetstrokecolor{currentstroke}%
\pgfsetstrokeopacity{0.900000}%
\pgfsetdash{}{0pt}%
\pgfpathmoveto{\pgfqpoint{13.568091in}{8.532262in}}%
\pgfpathlineto{\pgfqpoint{13.627829in}{8.592000in}}%
\pgfpathlineto{\pgfqpoint{13.568091in}{8.651738in}}%
\pgfpathlineto{\pgfqpoint{13.508353in}{8.592000in}}%
\pgfpathlineto{\pgfqpoint{13.568091in}{8.532262in}}%
\pgfpathclose%
\pgfusepath{stroke}%
\end{pgfscope}%
\begin{pgfscope}%
\pgfsetbuttcap%
\pgfsetroundjoin%
\pgfsetlinewidth{1.003750pt}%
\definecolor{currentstroke}{rgb}{0.337255,0.705882,0.913725}%
\pgfsetstrokecolor{currentstroke}%
\pgfsetstrokeopacity{0.900000}%
\pgfsetdash{}{0pt}%
\pgfpathmoveto{\pgfqpoint{7.904091in}{2.868262in}}%
\pgfpathlineto{\pgfqpoint{7.963829in}{2.928000in}}%
\pgfpathlineto{\pgfqpoint{7.904091in}{2.987738in}}%
\pgfpathlineto{\pgfqpoint{7.844353in}{2.928000in}}%
\pgfpathlineto{\pgfqpoint{7.904091in}{2.868262in}}%
\pgfpathclose%
\pgfusepath{stroke}%
\end{pgfscope}%
\begin{pgfscope}%
\pgfsetbuttcap%
\pgfsetroundjoin%
\pgfsetlinewidth{1.003750pt}%
\definecolor{currentstroke}{rgb}{0.337255,0.705882,0.913725}%
\pgfsetstrokecolor{currentstroke}%
\pgfsetstrokeopacity{0.900000}%
\pgfsetdash{}{0pt}%
\pgfpathmoveto{\pgfqpoint{7.904091in}{2.868262in}}%
\pgfpathlineto{\pgfqpoint{7.963829in}{2.928000in}}%
\pgfpathlineto{\pgfqpoint{7.904091in}{2.987738in}}%
\pgfpathlineto{\pgfqpoint{7.844353in}{2.928000in}}%
\pgfpathlineto{\pgfqpoint{7.904091in}{2.868262in}}%
\pgfpathclose%
\pgfusepath{stroke}%
\end{pgfscope}%
\begin{pgfscope}%
\pgfpathrectangle{\pgfqpoint{7.904091in}{2.928000in}}{\pgfqpoint{5.664000in}{5.664000in}}%
\pgfusepath{clip}%
\pgfsetbuttcap%
\pgfsetroundjoin%
\pgfsetlinewidth{1.003750pt}%
\definecolor{currentstroke}{rgb}{0.662745,0.368627,0.000000}%
\pgfsetstrokecolor{currentstroke}%
\pgfsetstrokeopacity{0.900000}%
\pgfsetdash{}{0pt}%
\pgfpathmoveto{\pgfqpoint{12.453693in}{7.620286in}}%
\pgfpathcurveto{\pgfqpoint{12.463947in}{7.620286in}}{\pgfqpoint{12.473782in}{7.624360in}}{\pgfqpoint{12.481033in}{7.631610in}}%
\pgfpathcurveto{\pgfqpoint{12.488284in}{7.638861in}}{\pgfqpoint{12.492358in}{7.648697in}}{\pgfqpoint{12.492358in}{7.658951in}}%
\pgfpathcurveto{\pgfqpoint{12.492358in}{7.669205in}}{\pgfqpoint{12.488284in}{7.679040in}}{\pgfqpoint{12.481033in}{7.686291in}}%
\pgfpathcurveto{\pgfqpoint{12.473782in}{7.693542in}}{\pgfqpoint{12.463947in}{7.697616in}}{\pgfqpoint{12.453693in}{7.697616in}}%
\pgfpathcurveto{\pgfqpoint{12.443438in}{7.697616in}}{\pgfqpoint{12.433603in}{7.693542in}}{\pgfqpoint{12.426352in}{7.686291in}}%
\pgfpathcurveto{\pgfqpoint{12.419102in}{7.679040in}}{\pgfqpoint{12.415028in}{7.669205in}}{\pgfqpoint{12.415028in}{7.658951in}}%
\pgfpathcurveto{\pgfqpoint{12.415028in}{7.648697in}}{\pgfqpoint{12.419102in}{7.638861in}}{\pgfqpoint{12.426352in}{7.631610in}}%
\pgfpathcurveto{\pgfqpoint{12.433603in}{7.624360in}}{\pgfqpoint{12.443438in}{7.620286in}}{\pgfqpoint{12.453693in}{7.620286in}}%
\pgfpathlineto{\pgfqpoint{12.453693in}{7.620286in}}%
\pgfpathclose%
\pgfusepath{stroke}%
\end{pgfscope}%
\begin{pgfscope}%
\pgfpathrectangle{\pgfqpoint{7.904091in}{2.928000in}}{\pgfqpoint{5.664000in}{5.664000in}}%
\pgfusepath{clip}%
\pgfsetbuttcap%
\pgfsetroundjoin%
\pgfsetlinewidth{1.003750pt}%
\definecolor{currentstroke}{rgb}{0.662745,0.368627,0.000000}%
\pgfsetstrokecolor{currentstroke}%
\pgfsetstrokeopacity{0.900000}%
\pgfsetdash{}{0pt}%
\pgfpathmoveto{\pgfqpoint{12.914593in}{7.899708in}}%
\pgfpathcurveto{\pgfqpoint{12.924847in}{7.899708in}}{\pgfqpoint{12.934682in}{7.903782in}}{\pgfqpoint{12.941933in}{7.911033in}}%
\pgfpathcurveto{\pgfqpoint{12.949184in}{7.918283in}}{\pgfqpoint{12.953258in}{7.928119in}}{\pgfqpoint{12.953258in}{7.938373in}}%
\pgfpathcurveto{\pgfqpoint{12.953258in}{7.948627in}}{\pgfqpoint{12.949184in}{7.958462in}}{\pgfqpoint{12.941933in}{7.965713in}}%
\pgfpathcurveto{\pgfqpoint{12.934682in}{7.972964in}}{\pgfqpoint{12.924847in}{7.977038in}}{\pgfqpoint{12.914593in}{7.977038in}}%
\pgfpathcurveto{\pgfqpoint{12.904339in}{7.977038in}}{\pgfqpoint{12.894503in}{7.972964in}}{\pgfqpoint{12.887253in}{7.965713in}}%
\pgfpathcurveto{\pgfqpoint{12.880002in}{7.958462in}}{\pgfqpoint{12.875928in}{7.948627in}}{\pgfqpoint{12.875928in}{7.938373in}}%
\pgfpathcurveto{\pgfqpoint{12.875928in}{7.928119in}}{\pgfqpoint{12.880002in}{7.918283in}}{\pgfqpoint{12.887253in}{7.911033in}}%
\pgfpathcurveto{\pgfqpoint{12.894503in}{7.903782in}}{\pgfqpoint{12.904339in}{7.899708in}}{\pgfqpoint{12.914593in}{7.899708in}}%
\pgfpathlineto{\pgfqpoint{12.914593in}{7.899708in}}%
\pgfpathclose%
\pgfusepath{stroke}%
\end{pgfscope}%
\begin{pgfscope}%
\pgfpathrectangle{\pgfqpoint{7.904091in}{2.928000in}}{\pgfqpoint{5.664000in}{5.664000in}}%
\pgfusepath{clip}%
\pgfsetbuttcap%
\pgfsetroundjoin%
\pgfsetlinewidth{1.003750pt}%
\definecolor{currentstroke}{rgb}{0.662745,0.368627,0.000000}%
\pgfsetstrokecolor{currentstroke}%
\pgfsetstrokeopacity{0.900000}%
\pgfsetdash{}{0pt}%
\pgfpathmoveto{\pgfqpoint{12.495888in}{7.535294in}}%
\pgfpathcurveto{\pgfqpoint{12.506142in}{7.535294in}}{\pgfqpoint{12.515977in}{7.539368in}}{\pgfqpoint{12.523228in}{7.546618in}}%
\pgfpathcurveto{\pgfqpoint{12.530479in}{7.553869in}}{\pgfqpoint{12.534553in}{7.563704in}}{\pgfqpoint{12.534553in}{7.573959in}}%
\pgfpathcurveto{\pgfqpoint{12.534553in}{7.584213in}}{\pgfqpoint{12.530479in}{7.594048in}}{\pgfqpoint{12.523228in}{7.601299in}}%
\pgfpathcurveto{\pgfqpoint{12.515977in}{7.608550in}}{\pgfqpoint{12.506142in}{7.612624in}}{\pgfqpoint{12.495888in}{7.612624in}}%
\pgfpathcurveto{\pgfqpoint{12.485633in}{7.612624in}}{\pgfqpoint{12.475798in}{7.608550in}}{\pgfqpoint{12.468547in}{7.601299in}}%
\pgfpathcurveto{\pgfqpoint{12.461297in}{7.594048in}}{\pgfqpoint{12.457223in}{7.584213in}}{\pgfqpoint{12.457223in}{7.573959in}}%
\pgfpathcurveto{\pgfqpoint{12.457223in}{7.563704in}}{\pgfqpoint{12.461297in}{7.553869in}}{\pgfqpoint{12.468547in}{7.546618in}}%
\pgfpathcurveto{\pgfqpoint{12.475798in}{7.539368in}}{\pgfqpoint{12.485633in}{7.535294in}}{\pgfqpoint{12.495888in}{7.535294in}}%
\pgfpathlineto{\pgfqpoint{12.495888in}{7.535294in}}%
\pgfpathclose%
\pgfusepath{stroke}%
\end{pgfscope}%
\begin{pgfscope}%
\pgfsetbuttcap%
\pgfsetroundjoin%
\pgfsetlinewidth{1.003750pt}%
\definecolor{currentstroke}{rgb}{0.662745,0.368627,0.000000}%
\pgfsetstrokecolor{currentstroke}%
\pgfsetstrokeopacity{0.900000}%
\pgfsetdash{}{0pt}%
\pgfpathmoveto{\pgfqpoint{13.568091in}{8.532262in}}%
\pgfpathlineto{\pgfqpoint{13.627829in}{8.592000in}}%
\pgfpathlineto{\pgfqpoint{13.568091in}{8.651738in}}%
\pgfpathlineto{\pgfqpoint{13.508353in}{8.592000in}}%
\pgfpathlineto{\pgfqpoint{13.568091in}{8.532262in}}%
\pgfpathclose%
\pgfusepath{stroke}%
\end{pgfscope}%
\begin{pgfscope}%
\pgfpathrectangle{\pgfqpoint{7.904091in}{2.928000in}}{\pgfqpoint{5.664000in}{5.664000in}}%
\pgfusepath{clip}%
\pgfsetbuttcap%
\pgfsetroundjoin%
\pgfsetlinewidth{1.003750pt}%
\definecolor{currentstroke}{rgb}{0.901961,0.403922,0.619608}%
\pgfsetstrokecolor{currentstroke}%
\pgfsetstrokeopacity{0.900000}%
\pgfsetdash{}{0pt}%
\pgfpathmoveto{\pgfqpoint{10.201622in}{5.029299in}}%
\pgfpathcurveto{\pgfqpoint{10.211876in}{5.029299in}}{\pgfqpoint{10.221712in}{5.033373in}}{\pgfqpoint{10.228962in}{5.040624in}}%
\pgfpathcurveto{\pgfqpoint{10.236213in}{5.047875in}}{\pgfqpoint{10.240287in}{5.057710in}}{\pgfqpoint{10.240287in}{5.067964in}}%
\pgfpathcurveto{\pgfqpoint{10.240287in}{5.078218in}}{\pgfqpoint{10.236213in}{5.088054in}}{\pgfqpoint{10.228962in}{5.095304in}}%
\pgfpathcurveto{\pgfqpoint{10.221712in}{5.102555in}}{\pgfqpoint{10.211876in}{5.106629in}}{\pgfqpoint{10.201622in}{5.106629in}}%
\pgfpathcurveto{\pgfqpoint{10.191368in}{5.106629in}}{\pgfqpoint{10.181532in}{5.102555in}}{\pgfqpoint{10.174282in}{5.095304in}}%
\pgfpathcurveto{\pgfqpoint{10.167031in}{5.088054in}}{\pgfqpoint{10.162957in}{5.078218in}}{\pgfqpoint{10.162957in}{5.067964in}}%
\pgfpathcurveto{\pgfqpoint{10.162957in}{5.057710in}}{\pgfqpoint{10.167031in}{5.047875in}}{\pgfqpoint{10.174282in}{5.040624in}}%
\pgfpathcurveto{\pgfqpoint{10.181532in}{5.033373in}}{\pgfqpoint{10.191368in}{5.029299in}}{\pgfqpoint{10.201622in}{5.029299in}}%
\pgfpathlineto{\pgfqpoint{10.201622in}{5.029299in}}%
\pgfpathclose%
\pgfusepath{stroke}%
\end{pgfscope}%
\begin{pgfscope}%
\pgfpathrectangle{\pgfqpoint{7.904091in}{2.928000in}}{\pgfqpoint{5.664000in}{5.664000in}}%
\pgfusepath{clip}%
\pgfsetbuttcap%
\pgfsetroundjoin%
\pgfsetlinewidth{1.003750pt}%
\definecolor{currentstroke}{rgb}{0.901961,0.403922,0.619608}%
\pgfsetstrokecolor{currentstroke}%
\pgfsetstrokeopacity{0.900000}%
\pgfsetdash{}{0pt}%
\pgfpathmoveto{\pgfqpoint{11.609820in}{6.875668in}}%
\pgfpathcurveto{\pgfqpoint{11.620074in}{6.875668in}}{\pgfqpoint{11.629910in}{6.879742in}}{\pgfqpoint{11.637161in}{6.886993in}}%
\pgfpathcurveto{\pgfqpoint{11.644411in}{6.894243in}}{\pgfqpoint{11.648485in}{6.904079in}}{\pgfqpoint{11.648485in}{6.914333in}}%
\pgfpathcurveto{\pgfqpoint{11.648485in}{6.924587in}}{\pgfqpoint{11.644411in}{6.934422in}}{\pgfqpoint{11.637161in}{6.941673in}}%
\pgfpathcurveto{\pgfqpoint{11.629910in}{6.948924in}}{\pgfqpoint{11.620074in}{6.952998in}}{\pgfqpoint{11.609820in}{6.952998in}}%
\pgfpathcurveto{\pgfqpoint{11.599566in}{6.952998in}}{\pgfqpoint{11.589731in}{6.948924in}}{\pgfqpoint{11.582480in}{6.941673in}}%
\pgfpathcurveto{\pgfqpoint{11.575229in}{6.934422in}}{\pgfqpoint{11.571155in}{6.924587in}}{\pgfqpoint{11.571155in}{6.914333in}}%
\pgfpathcurveto{\pgfqpoint{11.571155in}{6.904079in}}{\pgfqpoint{11.575229in}{6.894243in}}{\pgfqpoint{11.582480in}{6.886993in}}%
\pgfpathcurveto{\pgfqpoint{11.589731in}{6.879742in}}{\pgfqpoint{11.599566in}{6.875668in}}{\pgfqpoint{11.609820in}{6.875668in}}%
\pgfpathlineto{\pgfqpoint{11.609820in}{6.875668in}}%
\pgfpathclose%
\pgfusepath{stroke}%
\end{pgfscope}%
\begin{pgfscope}%
\pgfpathrectangle{\pgfqpoint{7.904091in}{2.928000in}}{\pgfqpoint{5.664000in}{5.664000in}}%
\pgfusepath{clip}%
\pgfsetbuttcap%
\pgfsetroundjoin%
\pgfsetlinewidth{1.003750pt}%
\definecolor{currentstroke}{rgb}{0.901961,0.403922,0.619608}%
\pgfsetstrokecolor{currentstroke}%
\pgfsetstrokeopacity{0.900000}%
\pgfsetdash{}{0pt}%
\pgfpathmoveto{\pgfqpoint{11.318311in}{6.609233in}}%
\pgfpathcurveto{\pgfqpoint{11.328565in}{6.609233in}}{\pgfqpoint{11.338400in}{6.613307in}}{\pgfqpoint{11.345651in}{6.620558in}}%
\pgfpathcurveto{\pgfqpoint{11.352902in}{6.627809in}}{\pgfqpoint{11.356976in}{6.637644in}}{\pgfqpoint{11.356976in}{6.647898in}}%
\pgfpathcurveto{\pgfqpoint{11.356976in}{6.658152in}}{\pgfqpoint{11.352902in}{6.667988in}}{\pgfqpoint{11.345651in}{6.675239in}}%
\pgfpathcurveto{\pgfqpoint{11.338400in}{6.682489in}}{\pgfqpoint{11.328565in}{6.686563in}}{\pgfqpoint{11.318311in}{6.686563in}}%
\pgfpathcurveto{\pgfqpoint{11.308057in}{6.686563in}}{\pgfqpoint{11.298221in}{6.682489in}}{\pgfqpoint{11.290970in}{6.675239in}}%
\pgfpathcurveto{\pgfqpoint{11.283720in}{6.667988in}}{\pgfqpoint{11.279646in}{6.658152in}}{\pgfqpoint{11.279646in}{6.647898in}}%
\pgfpathcurveto{\pgfqpoint{11.279646in}{6.637644in}}{\pgfqpoint{11.283720in}{6.627809in}}{\pgfqpoint{11.290970in}{6.620558in}}%
\pgfpathcurveto{\pgfqpoint{11.298221in}{6.613307in}}{\pgfqpoint{11.308057in}{6.609233in}}{\pgfqpoint{11.318311in}{6.609233in}}%
\pgfpathlineto{\pgfqpoint{11.318311in}{6.609233in}}%
\pgfpathclose%
\pgfusepath{stroke}%
\end{pgfscope}%
\begin{pgfscope}%
\pgfpathrectangle{\pgfqpoint{7.904091in}{2.928000in}}{\pgfqpoint{5.664000in}{5.664000in}}%
\pgfusepath{clip}%
\pgfsetbuttcap%
\pgfsetroundjoin%
\pgfsetlinewidth{1.003750pt}%
\definecolor{currentstroke}{rgb}{0.901961,0.403922,0.619608}%
\pgfsetstrokecolor{currentstroke}%
\pgfsetstrokeopacity{0.900000}%
\pgfsetdash{}{0pt}%
\pgfpathmoveto{\pgfqpoint{11.421887in}{5.519686in}}%
\pgfpathcurveto{\pgfqpoint{11.432141in}{5.519686in}}{\pgfqpoint{11.441977in}{5.523760in}}{\pgfqpoint{11.449227in}{5.531011in}}%
\pgfpathcurveto{\pgfqpoint{11.456478in}{5.538262in}}{\pgfqpoint{11.460552in}{5.548097in}}{\pgfqpoint{11.460552in}{5.558352in}}%
\pgfpathcurveto{\pgfqpoint{11.460552in}{5.568606in}}{\pgfqpoint{11.456478in}{5.578441in}}{\pgfqpoint{11.449227in}{5.585692in}}%
\pgfpathcurveto{\pgfqpoint{11.441977in}{5.592943in}}{\pgfqpoint{11.432141in}{5.597017in}}{\pgfqpoint{11.421887in}{5.597017in}}%
\pgfpathcurveto{\pgfqpoint{11.411633in}{5.597017in}}{\pgfqpoint{11.401798in}{5.592943in}}{\pgfqpoint{11.394547in}{5.585692in}}%
\pgfpathcurveto{\pgfqpoint{11.387296in}{5.578441in}}{\pgfqpoint{11.383222in}{5.568606in}}{\pgfqpoint{11.383222in}{5.558352in}}%
\pgfpathcurveto{\pgfqpoint{11.383222in}{5.548097in}}{\pgfqpoint{11.387296in}{5.538262in}}{\pgfqpoint{11.394547in}{5.531011in}}%
\pgfpathcurveto{\pgfqpoint{11.401798in}{5.523760in}}{\pgfqpoint{11.411633in}{5.519686in}}{\pgfqpoint{11.421887in}{5.519686in}}%
\pgfpathlineto{\pgfqpoint{11.421887in}{5.519686in}}%
\pgfpathclose%
\pgfusepath{stroke}%
\end{pgfscope}%
\begin{pgfscope}%
\pgfpathrectangle{\pgfqpoint{7.904091in}{2.928000in}}{\pgfqpoint{5.664000in}{5.664000in}}%
\pgfusepath{clip}%
\pgfsetbuttcap%
\pgfsetroundjoin%
\pgfsetlinewidth{1.003750pt}%
\definecolor{currentstroke}{rgb}{0.941176,0.584314,0.337255}%
\pgfsetstrokecolor{currentstroke}%
\pgfsetstrokeopacity{0.900000}%
\pgfsetdash{}{0pt}%
\pgfpathmoveto{\pgfqpoint{13.276102in}{8.332602in}}%
\pgfpathcurveto{\pgfqpoint{13.286356in}{8.332602in}}{\pgfqpoint{13.296192in}{8.336676in}}{\pgfqpoint{13.303443in}{8.343927in}}%
\pgfpathcurveto{\pgfqpoint{13.310693in}{8.351177in}}{\pgfqpoint{13.314767in}{8.361013in}}{\pgfqpoint{13.314767in}{8.371267in}}%
\pgfpathcurveto{\pgfqpoint{13.314767in}{8.381521in}}{\pgfqpoint{13.310693in}{8.391357in}}{\pgfqpoint{13.303443in}{8.398607in}}%
\pgfpathcurveto{\pgfqpoint{13.296192in}{8.405858in}}{\pgfqpoint{13.286356in}{8.409932in}}{\pgfqpoint{13.276102in}{8.409932in}}%
\pgfpathcurveto{\pgfqpoint{13.265848in}{8.409932in}}{\pgfqpoint{13.256013in}{8.405858in}}{\pgfqpoint{13.248762in}{8.398607in}}%
\pgfpathcurveto{\pgfqpoint{13.241511in}{8.391357in}}{\pgfqpoint{13.237437in}{8.381521in}}{\pgfqpoint{13.237437in}{8.371267in}}%
\pgfpathcurveto{\pgfqpoint{13.237437in}{8.361013in}}{\pgfqpoint{13.241511in}{8.351177in}}{\pgfqpoint{13.248762in}{8.343927in}}%
\pgfpathcurveto{\pgfqpoint{13.256013in}{8.336676in}}{\pgfqpoint{13.265848in}{8.332602in}}{\pgfqpoint{13.276102in}{8.332602in}}%
\pgfpathlineto{\pgfqpoint{13.276102in}{8.332602in}}%
\pgfpathclose%
\pgfusepath{stroke}%
\end{pgfscope}%
\begin{pgfscope}%
\pgfpathrectangle{\pgfqpoint{7.904091in}{2.928000in}}{\pgfqpoint{5.664000in}{5.664000in}}%
\pgfusepath{clip}%
\pgfsetbuttcap%
\pgfsetroundjoin%
\pgfsetlinewidth{1.003750pt}%
\definecolor{currentstroke}{rgb}{0.835294,0.368627,0.000000}%
\pgfsetstrokecolor{currentstroke}%
\pgfsetstrokeopacity{0.900000}%
\pgfsetdash{}{0pt}%
\pgfpathmoveto{\pgfqpoint{12.285154in}{7.174966in}}%
\pgfpathcurveto{\pgfqpoint{12.295408in}{7.174966in}}{\pgfqpoint{12.305244in}{7.179040in}}{\pgfqpoint{12.312494in}{7.186291in}}%
\pgfpathcurveto{\pgfqpoint{12.319745in}{7.193542in}}{\pgfqpoint{12.323819in}{7.203377in}}{\pgfqpoint{12.323819in}{7.213631in}}%
\pgfpathcurveto{\pgfqpoint{12.323819in}{7.223885in}}{\pgfqpoint{12.319745in}{7.233721in}}{\pgfqpoint{12.312494in}{7.240972in}}%
\pgfpathcurveto{\pgfqpoint{12.305244in}{7.248222in}}{\pgfqpoint{12.295408in}{7.252296in}}{\pgfqpoint{12.285154in}{7.252296in}}%
\pgfpathcurveto{\pgfqpoint{12.274900in}{7.252296in}}{\pgfqpoint{12.265065in}{7.248222in}}{\pgfqpoint{12.257814in}{7.240972in}}%
\pgfpathcurveto{\pgfqpoint{12.250563in}{7.233721in}}{\pgfqpoint{12.246489in}{7.223885in}}{\pgfqpoint{12.246489in}{7.213631in}}%
\pgfpathcurveto{\pgfqpoint{12.246489in}{7.203377in}}{\pgfqpoint{12.250563in}{7.193542in}}{\pgfqpoint{12.257814in}{7.186291in}}%
\pgfpathcurveto{\pgfqpoint{12.265065in}{7.179040in}}{\pgfqpoint{12.274900in}{7.174966in}}{\pgfqpoint{12.285154in}{7.174966in}}%
\pgfpathlineto{\pgfqpoint{12.285154in}{7.174966in}}%
\pgfpathclose%
\pgfusepath{stroke}%
\end{pgfscope}%
\begin{pgfscope}%
\pgfpathrectangle{\pgfqpoint{7.904091in}{2.928000in}}{\pgfqpoint{5.664000in}{5.664000in}}%
\pgfusepath{clip}%
\pgfsetbuttcap%
\pgfsetroundjoin%
\pgfsetlinewidth{1.003750pt}%
\definecolor{currentstroke}{rgb}{0.835294,0.368627,0.000000}%
\pgfsetstrokecolor{currentstroke}%
\pgfsetstrokeopacity{0.900000}%
\pgfsetdash{}{0pt}%
\pgfpathmoveto{\pgfqpoint{11.964226in}{7.038747in}}%
\pgfpathcurveto{\pgfqpoint{11.974480in}{7.038747in}}{\pgfqpoint{11.984315in}{7.042821in}}{\pgfqpoint{11.991566in}{7.050072in}}%
\pgfpathcurveto{\pgfqpoint{11.998817in}{7.057322in}}{\pgfqpoint{12.002891in}{7.067158in}}{\pgfqpoint{12.002891in}{7.077412in}}%
\pgfpathcurveto{\pgfqpoint{12.002891in}{7.087666in}}{\pgfqpoint{11.998817in}{7.097502in}}{\pgfqpoint{11.991566in}{7.104752in}}%
\pgfpathcurveto{\pgfqpoint{11.984315in}{7.112003in}}{\pgfqpoint{11.974480in}{7.116077in}}{\pgfqpoint{11.964226in}{7.116077in}}%
\pgfpathcurveto{\pgfqpoint{11.953972in}{7.116077in}}{\pgfqpoint{11.944136in}{7.112003in}}{\pgfqpoint{11.936886in}{7.104752in}}%
\pgfpathcurveto{\pgfqpoint{11.929635in}{7.097502in}}{\pgfqpoint{11.925561in}{7.087666in}}{\pgfqpoint{11.925561in}{7.077412in}}%
\pgfpathcurveto{\pgfqpoint{11.925561in}{7.067158in}}{\pgfqpoint{11.929635in}{7.057322in}}{\pgfqpoint{11.936886in}{7.050072in}}%
\pgfpathcurveto{\pgfqpoint{11.944136in}{7.042821in}}{\pgfqpoint{11.953972in}{7.038747in}}{\pgfqpoint{11.964226in}{7.038747in}}%
\pgfpathlineto{\pgfqpoint{11.964226in}{7.038747in}}%
\pgfpathclose%
\pgfusepath{stroke}%
\end{pgfscope}%
\begin{pgfscope}%
\pgfpathrectangle{\pgfqpoint{7.904091in}{2.928000in}}{\pgfqpoint{5.664000in}{5.664000in}}%
\pgfusepath{clip}%
\pgfsetbuttcap%
\pgfsetroundjoin%
\pgfsetlinewidth{1.003750pt}%
\definecolor{currentstroke}{rgb}{0.149020,0.274510,0.325490}%
\pgfsetstrokecolor{currentstroke}%
\pgfsetstrokeopacity{0.900000}%
\pgfsetdash{}{0pt}%
\pgfpathmoveto{\pgfqpoint{11.531909in}{6.570419in}}%
\pgfpathcurveto{\pgfqpoint{11.542164in}{6.570419in}}{\pgfqpoint{11.551999in}{6.574493in}}{\pgfqpoint{11.559250in}{6.581743in}}%
\pgfpathcurveto{\pgfqpoint{11.566501in}{6.588994in}}{\pgfqpoint{11.570575in}{6.598830in}}{\pgfqpoint{11.570575in}{6.609084in}}%
\pgfpathcurveto{\pgfqpoint{11.570575in}{6.619338in}}{\pgfqpoint{11.566501in}{6.629173in}}{\pgfqpoint{11.559250in}{6.636424in}}%
\pgfpathcurveto{\pgfqpoint{11.551999in}{6.643675in}}{\pgfqpoint{11.542164in}{6.647749in}}{\pgfqpoint{11.531909in}{6.647749in}}%
\pgfpathcurveto{\pgfqpoint{11.521655in}{6.647749in}}{\pgfqpoint{11.511820in}{6.643675in}}{\pgfqpoint{11.504569in}{6.636424in}}%
\pgfpathcurveto{\pgfqpoint{11.497318in}{6.629173in}}{\pgfqpoint{11.493244in}{6.619338in}}{\pgfqpoint{11.493244in}{6.609084in}}%
\pgfpathcurveto{\pgfqpoint{11.493244in}{6.598830in}}{\pgfqpoint{11.497318in}{6.588994in}}{\pgfqpoint{11.504569in}{6.581743in}}%
\pgfpathcurveto{\pgfqpoint{11.511820in}{6.574493in}}{\pgfqpoint{11.521655in}{6.570419in}}{\pgfqpoint{11.531909in}{6.570419in}}%
\pgfpathlineto{\pgfqpoint{11.531909in}{6.570419in}}%
\pgfpathclose%
\pgfusepath{stroke}%
\end{pgfscope}%
\begin{pgfscope}%
\pgfpathrectangle{\pgfqpoint{7.904091in}{2.928000in}}{\pgfqpoint{5.664000in}{5.664000in}}%
\pgfusepath{clip}%
\pgfsetbuttcap%
\pgfsetroundjoin%
\pgfsetlinewidth{1.003750pt}%
\definecolor{currentstroke}{rgb}{0.149020,0.274510,0.325490}%
\pgfsetstrokecolor{currentstroke}%
\pgfsetstrokeopacity{0.900000}%
\pgfsetdash{}{0pt}%
\pgfpathmoveto{\pgfqpoint{11.839637in}{6.874190in}}%
\pgfpathcurveto{\pgfqpoint{11.849891in}{6.874190in}}{\pgfqpoint{11.859727in}{6.878264in}}{\pgfqpoint{11.866977in}{6.885515in}}%
\pgfpathcurveto{\pgfqpoint{11.874228in}{6.892765in}}{\pgfqpoint{11.878302in}{6.902601in}}{\pgfqpoint{11.878302in}{6.912855in}}%
\pgfpathcurveto{\pgfqpoint{11.878302in}{6.923109in}}{\pgfqpoint{11.874228in}{6.932944in}}{\pgfqpoint{11.866977in}{6.940195in}}%
\pgfpathcurveto{\pgfqpoint{11.859727in}{6.947446in}}{\pgfqpoint{11.849891in}{6.951520in}}{\pgfqpoint{11.839637in}{6.951520in}}%
\pgfpathcurveto{\pgfqpoint{11.829383in}{6.951520in}}{\pgfqpoint{11.819547in}{6.947446in}}{\pgfqpoint{11.812297in}{6.940195in}}%
\pgfpathcurveto{\pgfqpoint{11.805046in}{6.932944in}}{\pgfqpoint{11.800972in}{6.923109in}}{\pgfqpoint{11.800972in}{6.912855in}}%
\pgfpathcurveto{\pgfqpoint{11.800972in}{6.902601in}}{\pgfqpoint{11.805046in}{6.892765in}}{\pgfqpoint{11.812297in}{6.885515in}}%
\pgfpathcurveto{\pgfqpoint{11.819547in}{6.878264in}}{\pgfqpoint{11.829383in}{6.874190in}}{\pgfqpoint{11.839637in}{6.874190in}}%
\pgfpathlineto{\pgfqpoint{11.839637in}{6.874190in}}%
\pgfpathclose%
\pgfusepath{stroke}%
\end{pgfscope}%
\begin{pgfscope}%
\pgfpathrectangle{\pgfqpoint{7.904091in}{2.928000in}}{\pgfqpoint{5.664000in}{5.664000in}}%
\pgfusepath{clip}%
\pgfsetbuttcap%
\pgfsetroundjoin%
\pgfsetlinewidth{1.003750pt}%
\definecolor{currentstroke}{rgb}{0.149020,0.274510,0.325490}%
\pgfsetstrokecolor{currentstroke}%
\pgfsetstrokeopacity{0.900000}%
\pgfsetdash{}{0pt}%
\pgfpathmoveto{\pgfqpoint{12.049209in}{7.132584in}}%
\pgfpathcurveto{\pgfqpoint{12.059463in}{7.132584in}}{\pgfqpoint{12.069298in}{7.136658in}}{\pgfqpoint{12.076549in}{7.143909in}}%
\pgfpathcurveto{\pgfqpoint{12.083800in}{7.151160in}}{\pgfqpoint{12.087874in}{7.160995in}}{\pgfqpoint{12.087874in}{7.171249in}}%
\pgfpathcurveto{\pgfqpoint{12.087874in}{7.181503in}}{\pgfqpoint{12.083800in}{7.191339in}}{\pgfqpoint{12.076549in}{7.198589in}}%
\pgfpathcurveto{\pgfqpoint{12.069298in}{7.205840in}}{\pgfqpoint{12.059463in}{7.209914in}}{\pgfqpoint{12.049209in}{7.209914in}}%
\pgfpathcurveto{\pgfqpoint{12.038955in}{7.209914in}}{\pgfqpoint{12.029119in}{7.205840in}}{\pgfqpoint{12.021868in}{7.198589in}}%
\pgfpathcurveto{\pgfqpoint{12.014618in}{7.191339in}}{\pgfqpoint{12.010544in}{7.181503in}}{\pgfqpoint{12.010544in}{7.171249in}}%
\pgfpathcurveto{\pgfqpoint{12.010544in}{7.160995in}}{\pgfqpoint{12.014618in}{7.151160in}}{\pgfqpoint{12.021868in}{7.143909in}}%
\pgfpathcurveto{\pgfqpoint{12.029119in}{7.136658in}}{\pgfqpoint{12.038955in}{7.132584in}}{\pgfqpoint{12.049209in}{7.132584in}}%
\pgfpathlineto{\pgfqpoint{12.049209in}{7.132584in}}%
\pgfpathclose%
\pgfusepath{stroke}%
\end{pgfscope}%
\begin{pgfscope}%
\pgfpathrectangle{\pgfqpoint{7.904091in}{2.928000in}}{\pgfqpoint{5.664000in}{5.664000in}}%
\pgfusepath{clip}%
\pgfsetbuttcap%
\pgfsetroundjoin%
\pgfsetlinewidth{1.003750pt}%
\definecolor{currentstroke}{rgb}{0.000000,0.619608,0.450980}%
\pgfsetstrokecolor{currentstroke}%
\pgfsetstrokeopacity{0.900000}%
\pgfsetdash{}{0pt}%
\pgfpathmoveto{\pgfqpoint{12.351290in}{7.298092in}}%
\pgfpathcurveto{\pgfqpoint{12.361544in}{7.298092in}}{\pgfqpoint{12.371379in}{7.302166in}}{\pgfqpoint{12.378630in}{7.309417in}}%
\pgfpathcurveto{\pgfqpoint{12.385881in}{7.316667in}}{\pgfqpoint{12.389955in}{7.326503in}}{\pgfqpoint{12.389955in}{7.336757in}}%
\pgfpathcurveto{\pgfqpoint{12.389955in}{7.347011in}}{\pgfqpoint{12.385881in}{7.356846in}}{\pgfqpoint{12.378630in}{7.364097in}}%
\pgfpathcurveto{\pgfqpoint{12.371379in}{7.371348in}}{\pgfqpoint{12.361544in}{7.375422in}}{\pgfqpoint{12.351290in}{7.375422in}}%
\pgfpathcurveto{\pgfqpoint{12.341036in}{7.375422in}}{\pgfqpoint{12.331200in}{7.371348in}}{\pgfqpoint{12.323949in}{7.364097in}}%
\pgfpathcurveto{\pgfqpoint{12.316699in}{7.356846in}}{\pgfqpoint{12.312625in}{7.347011in}}{\pgfqpoint{12.312625in}{7.336757in}}%
\pgfpathcurveto{\pgfqpoint{12.312625in}{7.326503in}}{\pgfqpoint{12.316699in}{7.316667in}}{\pgfqpoint{12.323949in}{7.309417in}}%
\pgfpathcurveto{\pgfqpoint{12.331200in}{7.302166in}}{\pgfqpoint{12.341036in}{7.298092in}}{\pgfqpoint{12.351290in}{7.298092in}}%
\pgfpathlineto{\pgfqpoint{12.351290in}{7.298092in}}%
\pgfpathclose%
\pgfusepath{stroke}%
\end{pgfscope}%
\begin{pgfscope}%
\pgfpathrectangle{\pgfqpoint{7.904091in}{2.928000in}}{\pgfqpoint{5.664000in}{5.664000in}}%
\pgfusepath{clip}%
\pgfsetbuttcap%
\pgfsetroundjoin%
\pgfsetlinewidth{1.003750pt}%
\definecolor{currentstroke}{rgb}{0.000000,0.619608,0.450980}%
\pgfsetstrokecolor{currentstroke}%
\pgfsetstrokeopacity{0.900000}%
\pgfsetdash{}{0pt}%
\pgfpathmoveto{\pgfqpoint{12.657091in}{7.788950in}}%
\pgfpathcurveto{\pgfqpoint{12.667345in}{7.788950in}}{\pgfqpoint{12.677180in}{7.793024in}}{\pgfqpoint{12.684431in}{7.800275in}}%
\pgfpathcurveto{\pgfqpoint{12.691682in}{7.807526in}}{\pgfqpoint{12.695756in}{7.817361in}}{\pgfqpoint{12.695756in}{7.827615in}}%
\pgfpathcurveto{\pgfqpoint{12.695756in}{7.837869in}}{\pgfqpoint{12.691682in}{7.847705in}}{\pgfqpoint{12.684431in}{7.854955in}}%
\pgfpathcurveto{\pgfqpoint{12.677180in}{7.862206in}}{\pgfqpoint{12.667345in}{7.866280in}}{\pgfqpoint{12.657091in}{7.866280in}}%
\pgfpathcurveto{\pgfqpoint{12.646837in}{7.866280in}}{\pgfqpoint{12.637001in}{7.862206in}}{\pgfqpoint{12.629751in}{7.854955in}}%
\pgfpathcurveto{\pgfqpoint{12.622500in}{7.847705in}}{\pgfqpoint{12.618426in}{7.837869in}}{\pgfqpoint{12.618426in}{7.827615in}}%
\pgfpathcurveto{\pgfqpoint{12.618426in}{7.817361in}}{\pgfqpoint{12.622500in}{7.807526in}}{\pgfqpoint{12.629751in}{7.800275in}}%
\pgfpathcurveto{\pgfqpoint{12.637001in}{7.793024in}}{\pgfqpoint{12.646837in}{7.788950in}}{\pgfqpoint{12.657091in}{7.788950in}}%
\pgfpathlineto{\pgfqpoint{12.657091in}{7.788950in}}%
\pgfpathclose%
\pgfusepath{stroke}%
\end{pgfscope}%
\begin{pgfscope}%
\pgfpathrectangle{\pgfqpoint{7.904091in}{2.928000in}}{\pgfqpoint{5.664000in}{5.664000in}}%
\pgfusepath{clip}%
\pgfsetbuttcap%
\pgfsetroundjoin%
\pgfsetlinewidth{1.003750pt}%
\definecolor{currentstroke}{rgb}{0.478431,0.243137,0.615686}%
\pgfsetstrokecolor{currentstroke}%
\pgfsetstrokeopacity{0.900000}%
\pgfsetdash{}{0pt}%
\pgfpathmoveto{\pgfqpoint{11.710071in}{6.788965in}}%
\pgfpathcurveto{\pgfqpoint{11.720325in}{6.788965in}}{\pgfqpoint{11.730160in}{6.793039in}}{\pgfqpoint{11.737411in}{6.800290in}}%
\pgfpathcurveto{\pgfqpoint{11.744662in}{6.807541in}}{\pgfqpoint{11.748736in}{6.817376in}}{\pgfqpoint{11.748736in}{6.827631in}}%
\pgfpathcurveto{\pgfqpoint{11.748736in}{6.837885in}}{\pgfqpoint{11.744662in}{6.847720in}}{\pgfqpoint{11.737411in}{6.854971in}}%
\pgfpathcurveto{\pgfqpoint{11.730160in}{6.862222in}}{\pgfqpoint{11.720325in}{6.866296in}}{\pgfqpoint{11.710071in}{6.866296in}}%
\pgfpathcurveto{\pgfqpoint{11.699817in}{6.866296in}}{\pgfqpoint{11.689981in}{6.862222in}}{\pgfqpoint{11.682730in}{6.854971in}}%
\pgfpathcurveto{\pgfqpoint{11.675480in}{6.847720in}}{\pgfqpoint{11.671406in}{6.837885in}}{\pgfqpoint{11.671406in}{6.827631in}}%
\pgfpathcurveto{\pgfqpoint{11.671406in}{6.817376in}}{\pgfqpoint{11.675480in}{6.807541in}}{\pgfqpoint{11.682730in}{6.800290in}}%
\pgfpathcurveto{\pgfqpoint{11.689981in}{6.793039in}}{\pgfqpoint{11.699817in}{6.788965in}}{\pgfqpoint{11.710071in}{6.788965in}}%
\pgfpathlineto{\pgfqpoint{11.710071in}{6.788965in}}%
\pgfpathclose%
\pgfusepath{stroke}%
\end{pgfscope}%
\begin{pgfscope}%
\pgfpathrectangle{\pgfqpoint{7.904091in}{2.928000in}}{\pgfqpoint{5.664000in}{5.664000in}}%
\pgfusepath{clip}%
\pgfsetbuttcap%
\pgfsetroundjoin%
\pgfsetlinewidth{1.003750pt}%
\definecolor{currentstroke}{rgb}{0.478431,0.243137,0.615686}%
\pgfsetstrokecolor{currentstroke}%
\pgfsetstrokeopacity{0.900000}%
\pgfsetdash{}{0pt}%
\pgfpathmoveto{\pgfqpoint{12.177147in}{7.291164in}}%
\pgfpathcurveto{\pgfqpoint{12.187401in}{7.291164in}}{\pgfqpoint{12.197237in}{7.295238in}}{\pgfqpoint{12.204488in}{7.302489in}}%
\pgfpathcurveto{\pgfqpoint{12.211738in}{7.309740in}}{\pgfqpoint{12.215812in}{7.319575in}}{\pgfqpoint{12.215812in}{7.329829in}}%
\pgfpathcurveto{\pgfqpoint{12.215812in}{7.340083in}}{\pgfqpoint{12.211738in}{7.349919in}}{\pgfqpoint{12.204488in}{7.357170in}}%
\pgfpathcurveto{\pgfqpoint{12.197237in}{7.364420in}}{\pgfqpoint{12.187401in}{7.368494in}}{\pgfqpoint{12.177147in}{7.368494in}}%
\pgfpathcurveto{\pgfqpoint{12.166893in}{7.368494in}}{\pgfqpoint{12.157058in}{7.364420in}}{\pgfqpoint{12.149807in}{7.357170in}}%
\pgfpathcurveto{\pgfqpoint{12.142556in}{7.349919in}}{\pgfqpoint{12.138482in}{7.340083in}}{\pgfqpoint{12.138482in}{7.329829in}}%
\pgfpathcurveto{\pgfqpoint{12.138482in}{7.319575in}}{\pgfqpoint{12.142556in}{7.309740in}}{\pgfqpoint{12.149807in}{7.302489in}}%
\pgfpathcurveto{\pgfqpoint{12.157058in}{7.295238in}}{\pgfqpoint{12.166893in}{7.291164in}}{\pgfqpoint{12.177147in}{7.291164in}}%
\pgfpathlineto{\pgfqpoint{12.177147in}{7.291164in}}%
\pgfpathclose%
\pgfusepath{stroke}%
\end{pgfscope}%
\begin{pgfscope}%
\pgfpathrectangle{\pgfqpoint{7.904091in}{2.928000in}}{\pgfqpoint{5.664000in}{5.664000in}}%
\pgfusepath{clip}%
\pgfsetbuttcap%
\pgfsetroundjoin%
\pgfsetlinewidth{1.003750pt}%
\definecolor{currentstroke}{rgb}{0.549020,0.337255,0.294118}%
\pgfsetstrokecolor{currentstroke}%
\pgfsetstrokeopacity{0.900000}%
\pgfsetdash{}{0pt}%
\pgfpathmoveto{\pgfqpoint{11.359022in}{5.647762in}}%
\pgfpathcurveto{\pgfqpoint{11.369276in}{5.647762in}}{\pgfqpoint{11.379112in}{5.651836in}}{\pgfqpoint{11.386363in}{5.659086in}}%
\pgfpathcurveto{\pgfqpoint{11.393613in}{5.666337in}}{\pgfqpoint{11.397687in}{5.676172in}}{\pgfqpoint{11.397687in}{5.686427in}}%
\pgfpathcurveto{\pgfqpoint{11.397687in}{5.696681in}}{\pgfqpoint{11.393613in}{5.706516in}}{\pgfqpoint{11.386363in}{5.713767in}}%
\pgfpathcurveto{\pgfqpoint{11.379112in}{5.721018in}}{\pgfqpoint{11.369276in}{5.725092in}}{\pgfqpoint{11.359022in}{5.725092in}}%
\pgfpathcurveto{\pgfqpoint{11.348768in}{5.725092in}}{\pgfqpoint{11.338933in}{5.721018in}}{\pgfqpoint{11.331682in}{5.713767in}}%
\pgfpathcurveto{\pgfqpoint{11.324431in}{5.706516in}}{\pgfqpoint{11.320357in}{5.696681in}}{\pgfqpoint{11.320357in}{5.686427in}}%
\pgfpathcurveto{\pgfqpoint{11.320357in}{5.676172in}}{\pgfqpoint{11.324431in}{5.666337in}}{\pgfqpoint{11.331682in}{5.659086in}}%
\pgfpathcurveto{\pgfqpoint{11.338933in}{5.651836in}}{\pgfqpoint{11.348768in}{5.647762in}}{\pgfqpoint{11.359022in}{5.647762in}}%
\pgfpathlineto{\pgfqpoint{11.359022in}{5.647762in}}%
\pgfpathclose%
\pgfusepath{stroke}%
\end{pgfscope}%
\begin{pgfscope}%
\pgfsetbuttcap%
\pgfsetroundjoin%
\definecolor{currentfill}{rgb}{0.352941,0.407843,0.450980}%
\pgfsetfillcolor{currentfill}%
\pgfsetlinewidth{0.000000pt}%
\definecolor{currentstroke}{rgb}{0.352941,0.407843,0.450980}%
\pgfsetstrokecolor{currentstroke}%
\pgfsetdash{}{0pt}%
\pgfpathmoveto{\pgfqpoint{7.929408in}{8.744547in}}%
\pgfpathlineto{\pgfqpoint{7.974826in}{8.744547in}}%
\pgfpathlineto{\pgfqpoint{7.974826in}{8.733600in}}%
\pgfpathlineto{\pgfqpoint{7.913760in}{8.733600in}}%
\pgfpathlineto{\pgfqpoint{7.913760in}{8.744547in}}%
\pgfpathquadraticcurveto{\pgfqpoint{7.921161in}{8.752217in}}{\pgfqpoint{7.933943in}{8.765122in}}%
\pgfpathquadraticcurveto{\pgfqpoint{7.946746in}{8.778049in}}{\pgfqpoint{7.950024in}{8.781801in}}%
\pgfpathquadraticcurveto{\pgfqpoint{7.956271in}{8.788811in}}{\pgfqpoint{7.958745in}{8.793676in}}%
\pgfpathquadraticcurveto{\pgfqpoint{7.961239in}{8.798541in}}{\pgfqpoint{7.961239in}{8.803242in}}%
\pgfpathquadraticcurveto{\pgfqpoint{7.961239in}{8.810911in}}{\pgfqpoint{7.955858in}{8.815735in}}%
\pgfpathquadraticcurveto{\pgfqpoint{7.950478in}{8.820580in}}{\pgfqpoint{7.941839in}{8.820580in}}%
\pgfpathquadraticcurveto{\pgfqpoint{7.935716in}{8.820580in}}{\pgfqpoint{7.928913in}{8.818457in}}%
\pgfpathquadraticcurveto{\pgfqpoint{7.922130in}{8.816333in}}{\pgfqpoint{7.914399in}{8.812004in}}%
\pgfpathlineto{\pgfqpoint{7.914399in}{8.825157in}}%
\pgfpathquadraticcurveto{\pgfqpoint{7.922254in}{8.828311in}}{\pgfqpoint{7.929078in}{8.829919in}}%
\pgfpathquadraticcurveto{\pgfqpoint{7.935923in}{8.831528in}}{\pgfqpoint{7.941592in}{8.831528in}}%
\pgfpathquadraticcurveto{\pgfqpoint{7.956539in}{8.831528in}}{\pgfqpoint{7.965424in}{8.824044in}}%
\pgfpathquadraticcurveto{\pgfqpoint{7.974310in}{8.816581in}}{\pgfqpoint{7.974310in}{8.804087in}}%
\pgfpathquadraticcurveto{\pgfqpoint{7.974310in}{8.798150in}}{\pgfqpoint{7.972084in}{8.792831in}}%
\pgfpathquadraticcurveto{\pgfqpoint{7.969878in}{8.787532in}}{\pgfqpoint{7.964002in}{8.780317in}}%
\pgfpathquadraticcurveto{\pgfqpoint{7.962394in}{8.778440in}}{\pgfqpoint{7.953756in}{8.769514in}}%
\pgfpathquadraticcurveto{\pgfqpoint{7.945138in}{8.760587in}}{\pgfqpoint{7.929408in}{8.744547in}}%
\pgfpathlineto{\pgfqpoint{7.929408in}{8.744547in}}%
\pgfpathclose%
\pgfpathmoveto{\pgfqpoint{8.002534in}{8.735600in}}%
\pgfpathlineto{\pgfqpoint{8.002534in}{8.747454in}}%
\pgfpathquadraticcurveto{\pgfqpoint{8.007441in}{8.745125in}}{\pgfqpoint{8.012450in}{8.743908in}}%
\pgfpathquadraticcurveto{\pgfqpoint{8.017481in}{8.742692in}}{\pgfqpoint{8.022326in}{8.742692in}}%
\pgfpathquadraticcurveto{\pgfqpoint{8.035211in}{8.742692in}}{\pgfqpoint{8.041993in}{8.751351in}}%
\pgfpathquadraticcurveto{\pgfqpoint{8.048797in}{8.760010in}}{\pgfqpoint{8.049766in}{8.777678in}}%
\pgfpathquadraticcurveto{\pgfqpoint{8.046034in}{8.772132in}}{\pgfqpoint{8.040282in}{8.769163in}}%
\pgfpathquadraticcurveto{\pgfqpoint{8.034551in}{8.766194in}}{\pgfqpoint{8.027603in}{8.766194in}}%
\pgfpathquadraticcurveto{\pgfqpoint{8.013172in}{8.766194in}}{\pgfqpoint{8.004760in}{8.774915in}}%
\pgfpathquadraticcurveto{\pgfqpoint{7.996349in}{8.783656in}}{\pgfqpoint{7.996349in}{8.798809in}}%
\pgfpathquadraticcurveto{\pgfqpoint{7.996349in}{8.813612in}}{\pgfqpoint{8.005111in}{8.822559in}}%
\pgfpathquadraticcurveto{\pgfqpoint{8.013873in}{8.831528in}}{\pgfqpoint{8.028428in}{8.831528in}}%
\pgfpathquadraticcurveto{\pgfqpoint{8.045127in}{8.831528in}}{\pgfqpoint{8.053910in}{8.818725in}}%
\pgfpathquadraticcurveto{\pgfqpoint{8.062713in}{8.805943in}}{\pgfqpoint{8.062713in}{8.781595in}}%
\pgfpathquadraticcurveto{\pgfqpoint{8.062713in}{8.758855in}}{\pgfqpoint{8.051910in}{8.745289in}}%
\pgfpathquadraticcurveto{\pgfqpoint{8.041128in}{8.731724in}}{\pgfqpoint{8.022903in}{8.731724in}}%
\pgfpathquadraticcurveto{\pgfqpoint{8.017996in}{8.731724in}}{\pgfqpoint{8.012966in}{8.732693in}}%
\pgfpathquadraticcurveto{\pgfqpoint{8.007956in}{8.733662in}}{\pgfqpoint{8.002534in}{8.735600in}}%
\pgfpathlineto{\pgfqpoint{8.002534in}{8.735600in}}%
\pgfpathclose%
\pgfpathmoveto{\pgfqpoint{8.028428in}{8.776379in}}%
\pgfpathquadraticcurveto{\pgfqpoint{8.037190in}{8.776379in}}{\pgfqpoint{8.042303in}{8.782358in}}%
\pgfpathquadraticcurveto{\pgfqpoint{8.047436in}{8.788357in}}{\pgfqpoint{8.047436in}{8.798809in}}%
\pgfpathquadraticcurveto{\pgfqpoint{8.047436in}{8.809179in}}{\pgfqpoint{8.042303in}{8.815199in}}%
\pgfpathquadraticcurveto{\pgfqpoint{8.037190in}{8.821219in}}{\pgfqpoint{8.028428in}{8.821219in}}%
\pgfpathquadraticcurveto{\pgfqpoint{8.019666in}{8.821219in}}{\pgfqpoint{8.014553in}{8.815199in}}%
\pgfpathquadraticcurveto{\pgfqpoint{8.009440in}{8.809179in}}{\pgfqpoint{8.009440in}{8.798809in}}%
\pgfpathquadraticcurveto{\pgfqpoint{8.009440in}{8.788357in}}{\pgfqpoint{8.014553in}{8.782358in}}%
\pgfpathquadraticcurveto{\pgfqpoint{8.019666in}{8.776379in}}{\pgfqpoint{8.028428in}{8.776379in}}%
\pgfpathlineto{\pgfqpoint{8.028428in}{8.776379in}}%
\pgfpathclose%
\pgfpathmoveto{\pgfqpoint{8.082814in}{8.829796in}}%
\pgfpathlineto{\pgfqpoint{8.144663in}{8.829796in}}%
\pgfpathlineto{\pgfqpoint{8.144663in}{8.824250in}}%
\pgfpathlineto{\pgfqpoint{8.109739in}{8.733600in}}%
\pgfpathlineto{\pgfqpoint{8.096153in}{8.733600in}}%
\pgfpathlineto{\pgfqpoint{8.129015in}{8.818828in}}%
\pgfpathlineto{\pgfqpoint{8.082814in}{8.818828in}}%
\pgfpathlineto{\pgfqpoint{8.082814in}{8.829796in}}%
\pgfpathlineto{\pgfqpoint{8.082814in}{8.829796in}}%
\pgfpathclose%
\pgfpathmoveto{\pgfqpoint{8.252115in}{8.794686in}}%
\pgfpathquadraticcurveto{\pgfqpoint{8.250115in}{8.795841in}}{\pgfqpoint{8.247765in}{8.796377in}}%
\pgfpathquadraticcurveto{\pgfqpoint{8.245415in}{8.796933in}}{\pgfqpoint{8.242590in}{8.796933in}}%
\pgfpathquadraticcurveto{\pgfqpoint{8.232529in}{8.796933in}}{\pgfqpoint{8.227149in}{8.790398in}}%
\pgfpathquadraticcurveto{\pgfqpoint{8.221768in}{8.783863in}}{\pgfqpoint{8.221768in}{8.771616in}}%
\pgfpathlineto{\pgfqpoint{8.221768in}{8.733600in}}%
\pgfpathlineto{\pgfqpoint{8.209852in}{8.733600in}}%
\pgfpathlineto{\pgfqpoint{8.209852in}{8.805757in}}%
\pgfpathlineto{\pgfqpoint{8.221768in}{8.805757in}}%
\pgfpathlineto{\pgfqpoint{8.221768in}{8.794542in}}%
\pgfpathquadraticcurveto{\pgfqpoint{8.225520in}{8.801118in}}{\pgfqpoint{8.231499in}{8.804293in}}%
\pgfpathquadraticcurveto{\pgfqpoint{8.237498in}{8.807489in}}{\pgfqpoint{8.246074in}{8.807489in}}%
\pgfpathquadraticcurveto{\pgfqpoint{8.247291in}{8.807489in}}{\pgfqpoint{8.248775in}{8.807324in}}%
\pgfpathquadraticcurveto{\pgfqpoint{8.250260in}{8.807180in}}{\pgfqpoint{8.252053in}{8.806850in}}%
\pgfpathlineto{\pgfqpoint{8.252115in}{8.794686in}}%
\pgfpathlineto{\pgfqpoint{8.252115in}{8.794686in}}%
\pgfpathclose%
\pgfpathmoveto{\pgfqpoint{8.323365in}{8.772647in}}%
\pgfpathlineto{\pgfqpoint{8.323365in}{8.766854in}}%
\pgfpathlineto{\pgfqpoint{8.268855in}{8.766854in}}%
\pgfpathquadraticcurveto{\pgfqpoint{8.269639in}{8.754608in}}{\pgfqpoint{8.276236in}{8.748196in}}%
\pgfpathquadraticcurveto{\pgfqpoint{8.282833in}{8.741785in}}{\pgfqpoint{8.294626in}{8.741785in}}%
\pgfpathquadraticcurveto{\pgfqpoint{8.301450in}{8.741785in}}{\pgfqpoint{8.307862in}{8.743455in}}%
\pgfpathquadraticcurveto{\pgfqpoint{8.314273in}{8.745125in}}{\pgfqpoint{8.320602in}{8.748485in}}%
\pgfpathlineto{\pgfqpoint{8.320602in}{8.737270in}}%
\pgfpathquadraticcurveto{\pgfqpoint{8.314211in}{8.734569in}}{\pgfqpoint{8.307511in}{8.733146in}}%
\pgfpathquadraticcurveto{\pgfqpoint{8.300811in}{8.731724in}}{\pgfqpoint{8.293925in}{8.731724in}}%
\pgfpathquadraticcurveto{\pgfqpoint{8.276648in}{8.731724in}}{\pgfqpoint{8.266567in}{8.741764in}}%
\pgfpathquadraticcurveto{\pgfqpoint{8.256486in}{8.751825in}}{\pgfqpoint{8.256486in}{8.768978in}}%
\pgfpathquadraticcurveto{\pgfqpoint{8.256486in}{8.786687in}}{\pgfqpoint{8.266052in}{8.797078in}}%
\pgfpathquadraticcurveto{\pgfqpoint{8.275618in}{8.807489in}}{\pgfqpoint{8.291863in}{8.807489in}}%
\pgfpathquadraticcurveto{\pgfqpoint{8.306418in}{8.807489in}}{\pgfqpoint{8.314892in}{8.798108in}}%
\pgfpathquadraticcurveto{\pgfqpoint{8.323365in}{8.788749in}}{\pgfqpoint{8.323365in}{8.772647in}}%
\pgfpathlineto{\pgfqpoint{8.323365in}{8.772647in}}%
\pgfpathclose%
\pgfpathmoveto{\pgfqpoint{8.311511in}{8.776131in}}%
\pgfpathquadraticcurveto{\pgfqpoint{8.311387in}{8.785842in}}{\pgfqpoint{8.306068in}{8.791635in}}%
\pgfpathquadraticcurveto{\pgfqpoint{8.300749in}{8.797449in}}{\pgfqpoint{8.291987in}{8.797449in}}%
\pgfpathquadraticcurveto{\pgfqpoint{8.282071in}{8.797449in}}{\pgfqpoint{8.276112in}{8.791841in}}%
\pgfpathquadraticcurveto{\pgfqpoint{8.270154in}{8.786233in}}{\pgfqpoint{8.269247in}{8.776049in}}%
\pgfpathlineto{\pgfqpoint{8.311511in}{8.776131in}}%
\pgfpathlineto{\pgfqpoint{8.311511in}{8.776131in}}%
\pgfpathclose%
\pgfpathmoveto{\pgfqpoint{8.354557in}{8.826250in}}%
\pgfpathlineto{\pgfqpoint{8.354557in}{8.805757in}}%
\pgfpathlineto{\pgfqpoint{8.378967in}{8.805757in}}%
\pgfpathlineto{\pgfqpoint{8.378967in}{8.796542in}}%
\pgfpathlineto{\pgfqpoint{8.354557in}{8.796542in}}%
\pgfpathlineto{\pgfqpoint{8.354557in}{8.757371in}}%
\pgfpathquadraticcurveto{\pgfqpoint{8.354557in}{8.748547in}}{\pgfqpoint{8.356970in}{8.746032in}}%
\pgfpathquadraticcurveto{\pgfqpoint{8.359382in}{8.743516in}}{\pgfqpoint{8.366804in}{8.743516in}}%
\pgfpathlineto{\pgfqpoint{8.378967in}{8.743516in}}%
\pgfpathlineto{\pgfqpoint{8.378967in}{8.733600in}}%
\pgfpathlineto{\pgfqpoint{8.366804in}{8.733600in}}%
\pgfpathquadraticcurveto{\pgfqpoint{8.353073in}{8.733600in}}{\pgfqpoint{8.347857in}{8.738713in}}%
\pgfpathquadraticcurveto{\pgfqpoint{8.342641in}{8.743846in}}{\pgfqpoint{8.342641in}{8.757371in}}%
\pgfpathlineto{\pgfqpoint{8.342641in}{8.796542in}}%
\pgfpathlineto{\pgfqpoint{8.333941in}{8.796542in}}%
\pgfpathlineto{\pgfqpoint{8.333941in}{8.805757in}}%
\pgfpathlineto{\pgfqpoint{8.342641in}{8.805757in}}%
\pgfpathlineto{\pgfqpoint{8.342641in}{8.826250in}}%
\pgfpathlineto{\pgfqpoint{8.354557in}{8.826250in}}%
\pgfpathlineto{\pgfqpoint{8.354557in}{8.826250in}}%
\pgfpathclose%
\pgfpathmoveto{\pgfqpoint{8.427354in}{8.769864in}}%
\pgfpathquadraticcurveto{\pgfqpoint{8.412984in}{8.769864in}}{\pgfqpoint{8.407438in}{8.766586in}}%
\pgfpathquadraticcurveto{\pgfqpoint{8.401893in}{8.763308in}}{\pgfqpoint{8.401893in}{8.755371in}}%
\pgfpathquadraticcurveto{\pgfqpoint{8.401893in}{8.749062in}}{\pgfqpoint{8.406057in}{8.745351in}}%
\pgfpathquadraticcurveto{\pgfqpoint{8.410222in}{8.741661in}}{\pgfqpoint{8.417355in}{8.741661in}}%
\pgfpathquadraticcurveto{\pgfqpoint{8.427230in}{8.741661in}}{\pgfqpoint{8.433188in}{8.748650in}}%
\pgfpathquadraticcurveto{\pgfqpoint{8.439146in}{8.755639in}}{\pgfqpoint{8.439146in}{8.767225in}}%
\pgfpathlineto{\pgfqpoint{8.439146in}{8.769864in}}%
\pgfpathlineto{\pgfqpoint{8.427354in}{8.769864in}}%
\pgfpathlineto{\pgfqpoint{8.427354in}{8.769864in}}%
\pgfpathclose%
\pgfpathmoveto{\pgfqpoint{8.451001in}{8.774771in}}%
\pgfpathlineto{\pgfqpoint{8.451001in}{8.733600in}}%
\pgfpathlineto{\pgfqpoint{8.439146in}{8.733600in}}%
\pgfpathlineto{\pgfqpoint{8.439146in}{8.744547in}}%
\pgfpathquadraticcurveto{\pgfqpoint{8.435085in}{8.737991in}}{\pgfqpoint{8.429024in}{8.734858in}}%
\pgfpathquadraticcurveto{\pgfqpoint{8.422962in}{8.731724in}}{\pgfqpoint{8.414201in}{8.731724in}}%
\pgfpathquadraticcurveto{\pgfqpoint{8.403130in}{8.731724in}}{\pgfqpoint{8.396574in}{8.737950in}}%
\pgfpathquadraticcurveto{\pgfqpoint{8.390038in}{8.744176in}}{\pgfqpoint{8.390038in}{8.754608in}}%
\pgfpathquadraticcurveto{\pgfqpoint{8.390038in}{8.766772in}}{\pgfqpoint{8.398182in}{8.772957in}}%
\pgfpathquadraticcurveto{\pgfqpoint{8.406346in}{8.779141in}}{\pgfqpoint{8.422509in}{8.779141in}}%
\pgfpathlineto{\pgfqpoint{8.439146in}{8.779141in}}%
\pgfpathlineto{\pgfqpoint{8.439146in}{8.780317in}}%
\pgfpathquadraticcurveto{\pgfqpoint{8.439146in}{8.788501in}}{\pgfqpoint{8.433765in}{8.792975in}}%
\pgfpathquadraticcurveto{\pgfqpoint{8.428385in}{8.797449in}}{\pgfqpoint{8.418654in}{8.797449in}}%
\pgfpathquadraticcurveto{\pgfqpoint{8.412469in}{8.797449in}}{\pgfqpoint{8.406593in}{8.795964in}}%
\pgfpathquadraticcurveto{\pgfqpoint{8.400738in}{8.794480in}}{\pgfqpoint{8.395337in}{8.791511in}}%
\pgfpathlineto{\pgfqpoint{8.395337in}{8.802479in}}%
\pgfpathquadraticcurveto{\pgfqpoint{8.401831in}{8.804994in}}{\pgfqpoint{8.407954in}{8.806231in}}%
\pgfpathquadraticcurveto{\pgfqpoint{8.414077in}{8.807489in}}{\pgfqpoint{8.419870in}{8.807489in}}%
\pgfpathquadraticcurveto{\pgfqpoint{8.435538in}{8.807489in}}{\pgfqpoint{8.443270in}{8.799366in}}%
\pgfpathquadraticcurveto{\pgfqpoint{8.451001in}{8.791264in}}{\pgfqpoint{8.451001in}{8.774771in}}%
\pgfpathlineto{\pgfqpoint{8.451001in}{8.774771in}}%
\pgfpathclose%
\pgfpathmoveto{\pgfqpoint{8.475410in}{8.805757in}}%
\pgfpathlineto{\pgfqpoint{8.487265in}{8.805757in}}%
\pgfpathlineto{\pgfqpoint{8.487265in}{8.733600in}}%
\pgfpathlineto{\pgfqpoint{8.475410in}{8.733600in}}%
\pgfpathlineto{\pgfqpoint{8.475410in}{8.805757in}}%
\pgfpathlineto{\pgfqpoint{8.475410in}{8.805757in}}%
\pgfpathclose%
\pgfpathmoveto{\pgfqpoint{8.475410in}{8.833857in}}%
\pgfpathlineto{\pgfqpoint{8.487265in}{8.833857in}}%
\pgfpathlineto{\pgfqpoint{8.487265in}{8.818828in}}%
\pgfpathlineto{\pgfqpoint{8.475410in}{8.818828in}}%
\pgfpathlineto{\pgfqpoint{8.475410in}{8.833857in}}%
\pgfpathlineto{\pgfqpoint{8.475410in}{8.833857in}}%
\pgfpathclose%
\pgfpathmoveto{\pgfqpoint{8.572060in}{8.777162in}}%
\pgfpathlineto{\pgfqpoint{8.572060in}{8.733600in}}%
\pgfpathlineto{\pgfqpoint{8.560205in}{8.733600in}}%
\pgfpathlineto{\pgfqpoint{8.560205in}{8.776771in}}%
\pgfpathquadraticcurveto{\pgfqpoint{8.560205in}{8.787017in}}{\pgfqpoint{8.556206in}{8.792088in}}%
\pgfpathquadraticcurveto{\pgfqpoint{8.552206in}{8.797181in}}{\pgfqpoint{8.544228in}{8.797181in}}%
\pgfpathquadraticcurveto{\pgfqpoint{8.534620in}{8.797181in}}{\pgfqpoint{8.529075in}{8.791058in}}%
\pgfpathquadraticcurveto{\pgfqpoint{8.523529in}{8.784955in}}{\pgfqpoint{8.523529in}{8.774379in}}%
\pgfpathlineto{\pgfqpoint{8.523529in}{8.733600in}}%
\pgfpathlineto{\pgfqpoint{8.511613in}{8.733600in}}%
\pgfpathlineto{\pgfqpoint{8.511613in}{8.805757in}}%
\pgfpathlineto{\pgfqpoint{8.523529in}{8.805757in}}%
\pgfpathlineto{\pgfqpoint{8.523529in}{8.794542in}}%
\pgfpathquadraticcurveto{\pgfqpoint{8.527796in}{8.801057in}}{\pgfqpoint{8.533548in}{8.804273in}}%
\pgfpathquadraticcurveto{\pgfqpoint{8.539321in}{8.807489in}}{\pgfqpoint{8.546867in}{8.807489in}}%
\pgfpathquadraticcurveto{\pgfqpoint{8.559298in}{8.807489in}}{\pgfqpoint{8.565669in}{8.799799in}}%
\pgfpathquadraticcurveto{\pgfqpoint{8.572060in}{8.792109in}}{\pgfqpoint{8.572060in}{8.777162in}}%
\pgfpathlineto{\pgfqpoint{8.572060in}{8.777162in}}%
\pgfpathclose%
\pgfpathmoveto{\pgfqpoint{8.657411in}{8.772647in}}%
\pgfpathlineto{\pgfqpoint{8.657411in}{8.766854in}}%
\pgfpathlineto{\pgfqpoint{8.602902in}{8.766854in}}%
\pgfpathquadraticcurveto{\pgfqpoint{8.603685in}{8.754608in}}{\pgfqpoint{8.610282in}{8.748196in}}%
\pgfpathquadraticcurveto{\pgfqpoint{8.616880in}{8.741785in}}{\pgfqpoint{8.628672in}{8.741785in}}%
\pgfpathquadraticcurveto{\pgfqpoint{8.635496in}{8.741785in}}{\pgfqpoint{8.641908in}{8.743455in}}%
\pgfpathquadraticcurveto{\pgfqpoint{8.648319in}{8.745125in}}{\pgfqpoint{8.654649in}{8.748485in}}%
\pgfpathlineto{\pgfqpoint{8.654649in}{8.737270in}}%
\pgfpathquadraticcurveto{\pgfqpoint{8.648258in}{8.734569in}}{\pgfqpoint{8.641557in}{8.733146in}}%
\pgfpathquadraticcurveto{\pgfqpoint{8.634857in}{8.731724in}}{\pgfqpoint{8.627971in}{8.731724in}}%
\pgfpathquadraticcurveto{\pgfqpoint{8.610695in}{8.731724in}}{\pgfqpoint{8.600613in}{8.741764in}}%
\pgfpathquadraticcurveto{\pgfqpoint{8.590532in}{8.751825in}}{\pgfqpoint{8.590532in}{8.768978in}}%
\pgfpathquadraticcurveto{\pgfqpoint{8.590532in}{8.786687in}}{\pgfqpoint{8.600098in}{8.797078in}}%
\pgfpathquadraticcurveto{\pgfqpoint{8.609664in}{8.807489in}}{\pgfqpoint{8.625909in}{8.807489in}}%
\pgfpathquadraticcurveto{\pgfqpoint{8.640465in}{8.807489in}}{\pgfqpoint{8.648938in}{8.798108in}}%
\pgfpathquadraticcurveto{\pgfqpoint{8.657411in}{8.788749in}}{\pgfqpoint{8.657411in}{8.772647in}}%
\pgfpathlineto{\pgfqpoint{8.657411in}{8.772647in}}%
\pgfpathclose%
\pgfpathmoveto{\pgfqpoint{8.645557in}{8.776131in}}%
\pgfpathquadraticcurveto{\pgfqpoint{8.645433in}{8.785842in}}{\pgfqpoint{8.640114in}{8.791635in}}%
\pgfpathquadraticcurveto{\pgfqpoint{8.634795in}{8.797449in}}{\pgfqpoint{8.626033in}{8.797449in}}%
\pgfpathquadraticcurveto{\pgfqpoint{8.616117in}{8.797449in}}{\pgfqpoint{8.610159in}{8.791841in}}%
\pgfpathquadraticcurveto{\pgfqpoint{8.604201in}{8.786233in}}{\pgfqpoint{8.603293in}{8.776049in}}%
\pgfpathlineto{\pgfqpoint{8.645557in}{8.776131in}}%
\pgfpathlineto{\pgfqpoint{8.645557in}{8.776131in}}%
\pgfpathclose%
\pgfpathmoveto{\pgfqpoint{8.724352in}{8.794810in}}%
\pgfpathlineto{\pgfqpoint{8.724352in}{8.833857in}}%
\pgfpathlineto{\pgfqpoint{8.736207in}{8.833857in}}%
\pgfpathlineto{\pgfqpoint{8.736207in}{8.733600in}}%
\pgfpathlineto{\pgfqpoint{8.724352in}{8.733600in}}%
\pgfpathlineto{\pgfqpoint{8.724352in}{8.744424in}}%
\pgfpathquadraticcurveto{\pgfqpoint{8.720621in}{8.737991in}}{\pgfqpoint{8.714910in}{8.734858in}}%
\pgfpathquadraticcurveto{\pgfqpoint{8.709220in}{8.731724in}}{\pgfqpoint{8.701221in}{8.731724in}}%
\pgfpathquadraticcurveto{\pgfqpoint{8.688150in}{8.731724in}}{\pgfqpoint{8.679924in}{8.742156in}}%
\pgfpathquadraticcurveto{\pgfqpoint{8.671719in}{8.752608in}}{\pgfqpoint{8.671719in}{8.769617in}}%
\pgfpathquadraticcurveto{\pgfqpoint{8.671719in}{8.786625in}}{\pgfqpoint{8.679924in}{8.797057in}}%
\pgfpathquadraticcurveto{\pgfqpoint{8.688150in}{8.807489in}}{\pgfqpoint{8.701221in}{8.807489in}}%
\pgfpathquadraticcurveto{\pgfqpoint{8.709220in}{8.807489in}}{\pgfqpoint{8.714910in}{8.804355in}}%
\pgfpathquadraticcurveto{\pgfqpoint{8.720621in}{8.801242in}}{\pgfqpoint{8.724352in}{8.794810in}}%
\pgfpathlineto{\pgfqpoint{8.724352in}{8.794810in}}%
\pgfpathclose%
\pgfpathmoveto{\pgfqpoint{8.683965in}{8.769617in}}%
\pgfpathquadraticcurveto{\pgfqpoint{8.683965in}{8.756546in}}{\pgfqpoint{8.689346in}{8.749103in}}%
\pgfpathquadraticcurveto{\pgfqpoint{8.694727in}{8.741661in}}{\pgfqpoint{8.704128in}{8.741661in}}%
\pgfpathquadraticcurveto{\pgfqpoint{8.713529in}{8.741661in}}{\pgfqpoint{8.718930in}{8.749103in}}%
\pgfpathquadraticcurveto{\pgfqpoint{8.724352in}{8.756546in}}{\pgfqpoint{8.724352in}{8.769617in}}%
\pgfpathquadraticcurveto{\pgfqpoint{8.724352in}{8.782687in}}{\pgfqpoint{8.718930in}{8.790130in}}%
\pgfpathquadraticcurveto{\pgfqpoint{8.713529in}{8.797572in}}{\pgfqpoint{8.704128in}{8.797572in}}%
\pgfpathquadraticcurveto{\pgfqpoint{8.694727in}{8.797572in}}{\pgfqpoint{8.689346in}{8.790130in}}%
\pgfpathquadraticcurveto{\pgfqpoint{8.683965in}{8.782687in}}{\pgfqpoint{8.683965in}{8.769617in}}%
\pgfpathlineto{\pgfqpoint{8.683965in}{8.769617in}}%
\pgfpathclose%
\pgfpathmoveto{\pgfqpoint{8.823641in}{8.829796in}}%
\pgfpathlineto{\pgfqpoint{8.834588in}{8.829796in}}%
\pgfpathlineto{\pgfqpoint{8.801086in}{8.721354in}}%
\pgfpathlineto{\pgfqpoint{8.790139in}{8.721354in}}%
\pgfpathlineto{\pgfqpoint{8.823641in}{8.829796in}}%
\pgfpathlineto{\pgfqpoint{8.823641in}{8.829796in}}%
\pgfpathclose%
\pgfpathmoveto{\pgfqpoint{8.930062in}{8.785471in}}%
\pgfpathquadraticcurveto{\pgfqpoint{8.939401in}{8.783471in}}{\pgfqpoint{8.944638in}{8.777142in}}%
\pgfpathquadraticcurveto{\pgfqpoint{8.949895in}{8.770833in}}{\pgfqpoint{8.949895in}{8.761556in}}%
\pgfpathquadraticcurveto{\pgfqpoint{8.949895in}{8.747330in}}{\pgfqpoint{8.940102in}{8.739517in}}%
\pgfpathquadraticcurveto{\pgfqpoint{8.930309in}{8.731724in}}{\pgfqpoint{8.912270in}{8.731724in}}%
\pgfpathquadraticcurveto{\pgfqpoint{8.906230in}{8.731724in}}{\pgfqpoint{8.899818in}{8.732920in}}%
\pgfpathquadraticcurveto{\pgfqpoint{8.893406in}{8.734115in}}{\pgfqpoint{8.886582in}{8.736507in}}%
\pgfpathlineto{\pgfqpoint{8.886582in}{8.749062in}}%
\pgfpathquadraticcurveto{\pgfqpoint{8.891984in}{8.745908in}}{\pgfqpoint{8.898416in}{8.744300in}}%
\pgfpathquadraticcurveto{\pgfqpoint{8.904869in}{8.742692in}}{\pgfqpoint{8.911899in}{8.742692in}}%
\pgfpathquadraticcurveto{\pgfqpoint{8.924125in}{8.742692in}}{\pgfqpoint{8.930536in}{8.747516in}}%
\pgfpathquadraticcurveto{\pgfqpoint{8.936948in}{8.752340in}}{\pgfqpoint{8.936948in}{8.761556in}}%
\pgfpathquadraticcurveto{\pgfqpoint{8.936948in}{8.770070in}}{\pgfqpoint{8.930990in}{8.774853in}}%
\pgfpathquadraticcurveto{\pgfqpoint{8.925032in}{8.779657in}}{\pgfqpoint{8.914414in}{8.779657in}}%
\pgfpathlineto{\pgfqpoint{8.903199in}{8.779657in}}%
\pgfpathlineto{\pgfqpoint{8.903199in}{8.790357in}}%
\pgfpathlineto{\pgfqpoint{8.914930in}{8.790357in}}%
\pgfpathquadraticcurveto{\pgfqpoint{8.924516in}{8.790357in}}{\pgfqpoint{8.929608in}{8.794191in}}%
\pgfpathquadraticcurveto{\pgfqpoint{8.934701in}{8.798026in}}{\pgfqpoint{8.934701in}{8.805242in}}%
\pgfpathquadraticcurveto{\pgfqpoint{8.934701in}{8.812643in}}{\pgfqpoint{8.929444in}{8.816601in}}%
\pgfpathquadraticcurveto{\pgfqpoint{8.924207in}{8.820580in}}{\pgfqpoint{8.914414in}{8.820580in}}%
\pgfpathquadraticcurveto{\pgfqpoint{8.909054in}{8.820580in}}{\pgfqpoint{8.902931in}{8.819405in}}%
\pgfpathquadraticcurveto{\pgfqpoint{8.896808in}{8.818251in}}{\pgfqpoint{8.889469in}{8.815818in}}%
\pgfpathlineto{\pgfqpoint{8.889469in}{8.827404in}}%
\pgfpathquadraticcurveto{\pgfqpoint{8.896890in}{8.829466in}}{\pgfqpoint{8.903364in}{8.830497in}}%
\pgfpathquadraticcurveto{\pgfqpoint{8.909837in}{8.831528in}}{\pgfqpoint{8.915569in}{8.831528in}}%
\pgfpathquadraticcurveto{\pgfqpoint{8.930392in}{8.831528in}}{\pgfqpoint{8.939010in}{8.824786in}}%
\pgfpathquadraticcurveto{\pgfqpoint{8.947648in}{8.818065in}}{\pgfqpoint{8.947648in}{8.806602in}}%
\pgfpathquadraticcurveto{\pgfqpoint{8.947648in}{8.798603in}}{\pgfqpoint{8.943071in}{8.793099in}}%
\pgfpathquadraticcurveto{\pgfqpoint{8.938494in}{8.787594in}}{\pgfqpoint{8.930062in}{8.785471in}}%
\pgfpathlineto{\pgfqpoint{8.930062in}{8.785471in}}%
\pgfpathclose%
\pgfpathmoveto{\pgfqpoint{8.971295in}{8.829796in}}%
\pgfpathlineto{\pgfqpoint{9.033144in}{8.829796in}}%
\pgfpathlineto{\pgfqpoint{9.033144in}{8.824250in}}%
\pgfpathlineto{\pgfqpoint{8.998220in}{8.733600in}}%
\pgfpathlineto{\pgfqpoint{8.984633in}{8.733600in}}%
\pgfpathlineto{\pgfqpoint{9.017496in}{8.818828in}}%
\pgfpathlineto{\pgfqpoint{8.971295in}{8.818828in}}%
\pgfpathlineto{\pgfqpoint{8.971295in}{8.829796in}}%
\pgfpathlineto{\pgfqpoint{8.971295in}{8.829796in}}%
\pgfpathclose%
\pgfpathmoveto{\pgfqpoint{9.097961in}{8.785471in}}%
\pgfpathquadraticcurveto{\pgfqpoint{9.107301in}{8.783471in}}{\pgfqpoint{9.112537in}{8.777142in}}%
\pgfpathquadraticcurveto{\pgfqpoint{9.117794in}{8.770833in}}{\pgfqpoint{9.117794in}{8.761556in}}%
\pgfpathquadraticcurveto{\pgfqpoint{9.117794in}{8.747330in}}{\pgfqpoint{9.108001in}{8.739517in}}%
\pgfpathquadraticcurveto{\pgfqpoint{9.098209in}{8.731724in}}{\pgfqpoint{9.080169in}{8.731724in}}%
\pgfpathquadraticcurveto{\pgfqpoint{9.074129in}{8.731724in}}{\pgfqpoint{9.067717in}{8.732920in}}%
\pgfpathquadraticcurveto{\pgfqpoint{9.061306in}{8.734115in}}{\pgfqpoint{9.054482in}{8.736507in}}%
\pgfpathlineto{\pgfqpoint{9.054482in}{8.749062in}}%
\pgfpathquadraticcurveto{\pgfqpoint{9.059883in}{8.745908in}}{\pgfqpoint{9.066315in}{8.744300in}}%
\pgfpathquadraticcurveto{\pgfqpoint{9.072768in}{8.742692in}}{\pgfqpoint{9.079798in}{8.742692in}}%
\pgfpathquadraticcurveto{\pgfqpoint{9.092024in}{8.742692in}}{\pgfqpoint{9.098436in}{8.747516in}}%
\pgfpathquadraticcurveto{\pgfqpoint{9.104847in}{8.752340in}}{\pgfqpoint{9.104847in}{8.761556in}}%
\pgfpathquadraticcurveto{\pgfqpoint{9.104847in}{8.770070in}}{\pgfqpoint{9.098889in}{8.774853in}}%
\pgfpathquadraticcurveto{\pgfqpoint{9.092931in}{8.779657in}}{\pgfqpoint{9.082314in}{8.779657in}}%
\pgfpathlineto{\pgfqpoint{9.071098in}{8.779657in}}%
\pgfpathlineto{\pgfqpoint{9.071098in}{8.790357in}}%
\pgfpathlineto{\pgfqpoint{9.082829in}{8.790357in}}%
\pgfpathquadraticcurveto{\pgfqpoint{9.092416in}{8.790357in}}{\pgfqpoint{9.097508in}{8.794191in}}%
\pgfpathquadraticcurveto{\pgfqpoint{9.102600in}{8.798026in}}{\pgfqpoint{9.102600in}{8.805242in}}%
\pgfpathquadraticcurveto{\pgfqpoint{9.102600in}{8.812643in}}{\pgfqpoint{9.097343in}{8.816601in}}%
\pgfpathquadraticcurveto{\pgfqpoint{9.092106in}{8.820580in}}{\pgfqpoint{9.082314in}{8.820580in}}%
\pgfpathquadraticcurveto{\pgfqpoint{9.076953in}{8.820580in}}{\pgfqpoint{9.070830in}{8.819405in}}%
\pgfpathquadraticcurveto{\pgfqpoint{9.064707in}{8.818251in}}{\pgfqpoint{9.057368in}{8.815818in}}%
\pgfpathlineto{\pgfqpoint{9.057368in}{8.827404in}}%
\pgfpathquadraticcurveto{\pgfqpoint{9.064790in}{8.829466in}}{\pgfqpoint{9.071263in}{8.830497in}}%
\pgfpathquadraticcurveto{\pgfqpoint{9.077737in}{8.831528in}}{\pgfqpoint{9.083468in}{8.831528in}}%
\pgfpathquadraticcurveto{\pgfqpoint{9.098291in}{8.831528in}}{\pgfqpoint{9.106909in}{8.824786in}}%
\pgfpathquadraticcurveto{\pgfqpoint{9.115547in}{8.818065in}}{\pgfqpoint{9.115547in}{8.806602in}}%
\pgfpathquadraticcurveto{\pgfqpoint{9.115547in}{8.798603in}}{\pgfqpoint{9.110970in}{8.793099in}}%
\pgfpathquadraticcurveto{\pgfqpoint{9.106393in}{8.787594in}}{\pgfqpoint{9.097961in}{8.785471in}}%
\pgfpathlineto{\pgfqpoint{9.097961in}{8.785471in}}%
\pgfpathclose%
\pgfpathmoveto{\pgfqpoint{9.234668in}{8.802995in}}%
\pgfpathlineto{\pgfqpoint{9.234668in}{8.791903in}}%
\pgfpathquadraticcurveto{\pgfqpoint{9.229638in}{8.794686in}}{\pgfqpoint{9.224587in}{8.796067in}}%
\pgfpathquadraticcurveto{\pgfqpoint{9.219536in}{8.797449in}}{\pgfqpoint{9.214382in}{8.797449in}}%
\pgfpathquadraticcurveto{\pgfqpoint{9.202837in}{8.797449in}}{\pgfqpoint{9.196446in}{8.790130in}}%
\pgfpathquadraticcurveto{\pgfqpoint{9.190075in}{8.782832in}}{\pgfqpoint{9.190075in}{8.769617in}}%
\pgfpathquadraticcurveto{\pgfqpoint{9.190075in}{8.756402in}}{\pgfqpoint{9.196446in}{8.749083in}}%
\pgfpathquadraticcurveto{\pgfqpoint{9.202837in}{8.741785in}}{\pgfqpoint{9.214382in}{8.741785in}}%
\pgfpathquadraticcurveto{\pgfqpoint{9.219536in}{8.741785in}}{\pgfqpoint{9.224587in}{8.743166in}}%
\pgfpathquadraticcurveto{\pgfqpoint{9.229638in}{8.744547in}}{\pgfqpoint{9.234668in}{8.747330in}}%
\pgfpathlineto{\pgfqpoint{9.234668in}{8.736363in}}%
\pgfpathquadraticcurveto{\pgfqpoint{9.229700in}{8.734054in}}{\pgfqpoint{9.224381in}{8.732899in}}%
\pgfpathquadraticcurveto{\pgfqpoint{9.219082in}{8.731724in}}{\pgfqpoint{9.213083in}{8.731724in}}%
\pgfpathquadraticcurveto{\pgfqpoint{9.196775in}{8.731724in}}{\pgfqpoint{9.187168in}{8.741970in}}%
\pgfpathquadraticcurveto{\pgfqpoint{9.177582in}{8.752217in}}{\pgfqpoint{9.177582in}{8.769617in}}%
\pgfpathquadraticcurveto{\pgfqpoint{9.177582in}{8.787264in}}{\pgfqpoint{9.187271in}{8.797366in}}%
\pgfpathquadraticcurveto{\pgfqpoint{9.196982in}{8.807489in}}{\pgfqpoint{9.213866in}{8.807489in}}%
\pgfpathquadraticcurveto{\pgfqpoint{9.219330in}{8.807489in}}{\pgfqpoint{9.224546in}{8.806355in}}%
\pgfpathquadraticcurveto{\pgfqpoint{9.229761in}{8.805242in}}{\pgfqpoint{9.234668in}{8.802995in}}%
\pgfpathlineto{\pgfqpoint{9.234668in}{8.802995in}}%
\pgfpathclose%
\pgfpathmoveto{\pgfqpoint{9.283240in}{8.797449in}}%
\pgfpathquadraticcurveto{\pgfqpoint{9.273715in}{8.797449in}}{\pgfqpoint{9.268170in}{8.790006in}}%
\pgfpathquadraticcurveto{\pgfqpoint{9.262624in}{8.782564in}}{\pgfqpoint{9.262624in}{8.769617in}}%
\pgfpathquadraticcurveto{\pgfqpoint{9.262624in}{8.756670in}}{\pgfqpoint{9.268128in}{8.749227in}}%
\pgfpathquadraticcurveto{\pgfqpoint{9.273654in}{8.741785in}}{\pgfqpoint{9.283240in}{8.741785in}}%
\pgfpathquadraticcurveto{\pgfqpoint{9.292724in}{8.741785in}}{\pgfqpoint{9.298249in}{8.749248in}}%
\pgfpathquadraticcurveto{\pgfqpoint{9.303795in}{8.756732in}}{\pgfqpoint{9.303795in}{8.769617in}}%
\pgfpathquadraticcurveto{\pgfqpoint{9.303795in}{8.782440in}}{\pgfqpoint{9.298249in}{8.789944in}}%
\pgfpathquadraticcurveto{\pgfqpoint{9.292724in}{8.797449in}}{\pgfqpoint{9.283240in}{8.797449in}}%
\pgfpathlineto{\pgfqpoint{9.283240in}{8.797449in}}%
\pgfpathclose%
\pgfpathmoveto{\pgfqpoint{9.283240in}{8.807489in}}%
\pgfpathquadraticcurveto{\pgfqpoint{9.298702in}{8.807489in}}{\pgfqpoint{9.307526in}{8.797428in}}%
\pgfpathquadraticcurveto{\pgfqpoint{9.316371in}{8.787388in}}{\pgfqpoint{9.316371in}{8.769617in}}%
\pgfpathquadraticcurveto{\pgfqpoint{9.316371in}{8.751907in}}{\pgfqpoint{9.307526in}{8.741805in}}%
\pgfpathquadraticcurveto{\pgfqpoint{9.298702in}{8.731724in}}{\pgfqpoint{9.283240in}{8.731724in}}%
\pgfpathquadraticcurveto{\pgfqpoint{9.267716in}{8.731724in}}{\pgfqpoint{9.258913in}{8.741805in}}%
\pgfpathquadraticcurveto{\pgfqpoint{9.250130in}{8.751907in}}{\pgfqpoint{9.250130in}{8.769617in}}%
\pgfpathquadraticcurveto{\pgfqpoint{9.250130in}{8.787388in}}{\pgfqpoint{9.258913in}{8.797428in}}%
\pgfpathquadraticcurveto{\pgfqpoint{9.267716in}{8.807489in}}{\pgfqpoint{9.283240in}{8.807489in}}%
\pgfpathlineto{\pgfqpoint{9.283240in}{8.807489in}}%
\pgfpathclose%
\pgfpathmoveto{\pgfqpoint{9.392197in}{8.791903in}}%
\pgfpathquadraticcurveto{\pgfqpoint{9.396651in}{8.799902in}}{\pgfqpoint{9.402835in}{8.803695in}}%
\pgfpathquadraticcurveto{\pgfqpoint{9.409020in}{8.807489in}}{\pgfqpoint{9.417391in}{8.807489in}}%
\pgfpathquadraticcurveto{\pgfqpoint{9.428668in}{8.807489in}}{\pgfqpoint{9.434791in}{8.799593in}}%
\pgfpathquadraticcurveto{\pgfqpoint{9.440914in}{8.791717in}}{\pgfqpoint{9.440914in}{8.777162in}}%
\pgfpathlineto{\pgfqpoint{9.440914in}{8.733600in}}%
\pgfpathlineto{\pgfqpoint{9.428998in}{8.733600in}}%
\pgfpathlineto{\pgfqpoint{9.428998in}{8.776771in}}%
\pgfpathquadraticcurveto{\pgfqpoint{9.428998in}{8.787141in}}{\pgfqpoint{9.425307in}{8.792150in}}%
\pgfpathquadraticcurveto{\pgfqpoint{9.421638in}{8.797181in}}{\pgfqpoint{9.414113in}{8.797181in}}%
\pgfpathquadraticcurveto{\pgfqpoint{9.404897in}{8.797181in}}{\pgfqpoint{9.399537in}{8.791058in}}%
\pgfpathquadraticcurveto{\pgfqpoint{9.394197in}{8.784955in}}{\pgfqpoint{9.394197in}{8.774379in}}%
\pgfpathlineto{\pgfqpoint{9.394197in}{8.733600in}}%
\pgfpathlineto{\pgfqpoint{9.382281in}{8.733600in}}%
\pgfpathlineto{\pgfqpoint{9.382281in}{8.776771in}}%
\pgfpathquadraticcurveto{\pgfqpoint{9.382281in}{8.787202in}}{\pgfqpoint{9.378611in}{8.792192in}}%
\pgfpathquadraticcurveto{\pgfqpoint{9.374942in}{8.797181in}}{\pgfqpoint{9.367272in}{8.797181in}}%
\pgfpathquadraticcurveto{\pgfqpoint{9.358181in}{8.797181in}}{\pgfqpoint{9.352820in}{8.791037in}}%
\pgfpathquadraticcurveto{\pgfqpoint{9.347481in}{8.784893in}}{\pgfqpoint{9.347481in}{8.774379in}}%
\pgfpathlineto{\pgfqpoint{9.347481in}{8.733600in}}%
\pgfpathlineto{\pgfqpoint{9.335564in}{8.733600in}}%
\pgfpathlineto{\pgfqpoint{9.335564in}{8.805757in}}%
\pgfpathlineto{\pgfqpoint{9.347481in}{8.805757in}}%
\pgfpathlineto{\pgfqpoint{9.347481in}{8.794542in}}%
\pgfpathquadraticcurveto{\pgfqpoint{9.351542in}{8.801180in}}{\pgfqpoint{9.357212in}{8.804335in}}%
\pgfpathquadraticcurveto{\pgfqpoint{9.362881in}{8.807489in}}{\pgfqpoint{9.370674in}{8.807489in}}%
\pgfpathquadraticcurveto{\pgfqpoint{9.378549in}{8.807489in}}{\pgfqpoint{9.384054in}{8.803489in}}%
\pgfpathquadraticcurveto{\pgfqpoint{9.389559in}{8.799510in}}{\pgfqpoint{9.392197in}{8.791903in}}%
\pgfpathlineto{\pgfqpoint{9.392197in}{8.791903in}}%
\pgfpathclose%
\pgfpathmoveto{\pgfqpoint{9.476003in}{8.744424in}}%
\pgfpathlineto{\pgfqpoint{9.476003in}{8.706160in}}%
\pgfpathlineto{\pgfqpoint{9.464087in}{8.706160in}}%
\pgfpathlineto{\pgfqpoint{9.464087in}{8.805757in}}%
\pgfpathlineto{\pgfqpoint{9.476003in}{8.805757in}}%
\pgfpathlineto{\pgfqpoint{9.476003in}{8.794810in}}%
\pgfpathquadraticcurveto{\pgfqpoint{9.479755in}{8.801242in}}{\pgfqpoint{9.485445in}{8.804355in}}%
\pgfpathquadraticcurveto{\pgfqpoint{9.491156in}{8.807489in}}{\pgfqpoint{9.499072in}{8.807489in}}%
\pgfpathquadraticcurveto{\pgfqpoint{9.512226in}{8.807489in}}{\pgfqpoint{9.520431in}{8.797057in}}%
\pgfpathquadraticcurveto{\pgfqpoint{9.528657in}{8.786625in}}{\pgfqpoint{9.528657in}{8.769617in}}%
\pgfpathquadraticcurveto{\pgfqpoint{9.528657in}{8.752608in}}{\pgfqpoint{9.520431in}{8.742156in}}%
\pgfpathquadraticcurveto{\pgfqpoint{9.512226in}{8.731724in}}{\pgfqpoint{9.499072in}{8.731724in}}%
\pgfpathquadraticcurveto{\pgfqpoint{9.491156in}{8.731724in}}{\pgfqpoint{9.485445in}{8.734858in}}%
\pgfpathquadraticcurveto{\pgfqpoint{9.479755in}{8.737991in}}{\pgfqpoint{9.476003in}{8.744424in}}%
\pgfpathlineto{\pgfqpoint{9.476003in}{8.744424in}}%
\pgfpathclose%
\pgfpathmoveto{\pgfqpoint{9.516349in}{8.769617in}}%
\pgfpathquadraticcurveto{\pgfqpoint{9.516349in}{8.782687in}}{\pgfqpoint{9.510968in}{8.790130in}}%
\pgfpathquadraticcurveto{\pgfqpoint{9.505587in}{8.797572in}}{\pgfqpoint{9.496186in}{8.797572in}}%
\pgfpathquadraticcurveto{\pgfqpoint{9.486765in}{8.797572in}}{\pgfqpoint{9.481384in}{8.790130in}}%
\pgfpathquadraticcurveto{\pgfqpoint{9.476003in}{8.782687in}}{\pgfqpoint{9.476003in}{8.769617in}}%
\pgfpathquadraticcurveto{\pgfqpoint{9.476003in}{8.756546in}}{\pgfqpoint{9.481384in}{8.749103in}}%
\pgfpathquadraticcurveto{\pgfqpoint{9.486765in}{8.741661in}}{\pgfqpoint{9.496186in}{8.741661in}}%
\pgfpathquadraticcurveto{\pgfqpoint{9.505587in}{8.741661in}}{\pgfqpoint{9.510968in}{8.749103in}}%
\pgfpathquadraticcurveto{\pgfqpoint{9.516349in}{8.756546in}}{\pgfqpoint{9.516349in}{8.769617in}}%
\pgfpathlineto{\pgfqpoint{9.516349in}{8.769617in}}%
\pgfpathclose%
\pgfpathmoveto{\pgfqpoint{9.581105in}{8.769864in}}%
\pgfpathquadraticcurveto{\pgfqpoint{9.566735in}{8.769864in}}{\pgfqpoint{9.561189in}{8.766586in}}%
\pgfpathquadraticcurveto{\pgfqpoint{9.555644in}{8.763308in}}{\pgfqpoint{9.555644in}{8.755371in}}%
\pgfpathquadraticcurveto{\pgfqpoint{9.555644in}{8.749062in}}{\pgfqpoint{9.559808in}{8.745351in}}%
\pgfpathquadraticcurveto{\pgfqpoint{9.563973in}{8.741661in}}{\pgfqpoint{9.571106in}{8.741661in}}%
\pgfpathquadraticcurveto{\pgfqpoint{9.580981in}{8.741661in}}{\pgfqpoint{9.586939in}{8.748650in}}%
\pgfpathquadraticcurveto{\pgfqpoint{9.592897in}{8.755639in}}{\pgfqpoint{9.592897in}{8.767225in}}%
\pgfpathlineto{\pgfqpoint{9.592897in}{8.769864in}}%
\pgfpathlineto{\pgfqpoint{9.581105in}{8.769864in}}%
\pgfpathlineto{\pgfqpoint{9.581105in}{8.769864in}}%
\pgfpathclose%
\pgfpathmoveto{\pgfqpoint{9.604752in}{8.774771in}}%
\pgfpathlineto{\pgfqpoint{9.604752in}{8.733600in}}%
\pgfpathlineto{\pgfqpoint{9.592897in}{8.733600in}}%
\pgfpathlineto{\pgfqpoint{9.592897in}{8.744547in}}%
\pgfpathquadraticcurveto{\pgfqpoint{9.588836in}{8.737991in}}{\pgfqpoint{9.582775in}{8.734858in}}%
\pgfpathquadraticcurveto{\pgfqpoint{9.576714in}{8.731724in}}{\pgfqpoint{9.567952in}{8.731724in}}%
\pgfpathquadraticcurveto{\pgfqpoint{9.556881in}{8.731724in}}{\pgfqpoint{9.550325in}{8.737950in}}%
\pgfpathquadraticcurveto{\pgfqpoint{9.543789in}{8.744176in}}{\pgfqpoint{9.543789in}{8.754608in}}%
\pgfpathquadraticcurveto{\pgfqpoint{9.543789in}{8.766772in}}{\pgfqpoint{9.551933in}{8.772957in}}%
\pgfpathquadraticcurveto{\pgfqpoint{9.560097in}{8.779141in}}{\pgfqpoint{9.576260in}{8.779141in}}%
\pgfpathlineto{\pgfqpoint{9.592897in}{8.779141in}}%
\pgfpathlineto{\pgfqpoint{9.592897in}{8.780317in}}%
\pgfpathquadraticcurveto{\pgfqpoint{9.592897in}{8.788501in}}{\pgfqpoint{9.587516in}{8.792975in}}%
\pgfpathquadraticcurveto{\pgfqpoint{9.582136in}{8.797449in}}{\pgfqpoint{9.572405in}{8.797449in}}%
\pgfpathquadraticcurveto{\pgfqpoint{9.566220in}{8.797449in}}{\pgfqpoint{9.560344in}{8.795964in}}%
\pgfpathquadraticcurveto{\pgfqpoint{9.554489in}{8.794480in}}{\pgfqpoint{9.549088in}{8.791511in}}%
\pgfpathlineto{\pgfqpoint{9.549088in}{8.802479in}}%
\pgfpathquadraticcurveto{\pgfqpoint{9.555582in}{8.804994in}}{\pgfqpoint{9.561705in}{8.806231in}}%
\pgfpathquadraticcurveto{\pgfqpoint{9.567828in}{8.807489in}}{\pgfqpoint{9.573621in}{8.807489in}}%
\pgfpathquadraticcurveto{\pgfqpoint{9.589289in}{8.807489in}}{\pgfqpoint{9.597021in}{8.799366in}}%
\pgfpathquadraticcurveto{\pgfqpoint{9.604752in}{8.791264in}}{\pgfqpoint{9.604752in}{8.774771in}}%
\pgfpathlineto{\pgfqpoint{9.604752in}{8.774771in}}%
\pgfpathclose%
\pgfpathmoveto{\pgfqpoint{9.670971in}{8.794686in}}%
\pgfpathquadraticcurveto{\pgfqpoint{9.668972in}{8.795841in}}{\pgfqpoint{9.666621in}{8.796377in}}%
\pgfpathquadraticcurveto{\pgfqpoint{9.664271in}{8.796933in}}{\pgfqpoint{9.661447in}{8.796933in}}%
\pgfpathquadraticcurveto{\pgfqpoint{9.651386in}{8.796933in}}{\pgfqpoint{9.646005in}{8.790398in}}%
\pgfpathquadraticcurveto{\pgfqpoint{9.640624in}{8.783863in}}{\pgfqpoint{9.640624in}{8.771616in}}%
\pgfpathlineto{\pgfqpoint{9.640624in}{8.733600in}}%
\pgfpathlineto{\pgfqpoint{9.628708in}{8.733600in}}%
\pgfpathlineto{\pgfqpoint{9.628708in}{8.805757in}}%
\pgfpathlineto{\pgfqpoint{9.640624in}{8.805757in}}%
\pgfpathlineto{\pgfqpoint{9.640624in}{8.794542in}}%
\pgfpathquadraticcurveto{\pgfqpoint{9.644376in}{8.801118in}}{\pgfqpoint{9.650355in}{8.804293in}}%
\pgfpathquadraticcurveto{\pgfqpoint{9.656354in}{8.807489in}}{\pgfqpoint{9.664931in}{8.807489in}}%
\pgfpathquadraticcurveto{\pgfqpoint{9.666147in}{8.807489in}}{\pgfqpoint{9.667631in}{8.807324in}}%
\pgfpathquadraticcurveto{\pgfqpoint{9.669116in}{8.807180in}}{\pgfqpoint{9.670909in}{8.806850in}}%
\pgfpathlineto{\pgfqpoint{9.670971in}{8.794686in}}%
\pgfpathlineto{\pgfqpoint{9.670971in}{8.794686in}}%
\pgfpathclose%
\pgfpathmoveto{\pgfqpoint{9.683403in}{8.805757in}}%
\pgfpathlineto{\pgfqpoint{9.695257in}{8.805757in}}%
\pgfpathlineto{\pgfqpoint{9.695257in}{8.733600in}}%
\pgfpathlineto{\pgfqpoint{9.683403in}{8.733600in}}%
\pgfpathlineto{\pgfqpoint{9.683403in}{8.805757in}}%
\pgfpathlineto{\pgfqpoint{9.683403in}{8.805757in}}%
\pgfpathclose%
\pgfpathmoveto{\pgfqpoint{9.683403in}{8.833857in}}%
\pgfpathlineto{\pgfqpoint{9.695257in}{8.833857in}}%
\pgfpathlineto{\pgfqpoint{9.695257in}{8.818828in}}%
\pgfpathlineto{\pgfqpoint{9.683403in}{8.818828in}}%
\pgfpathlineto{\pgfqpoint{9.683403in}{8.833857in}}%
\pgfpathlineto{\pgfqpoint{9.683403in}{8.833857in}}%
\pgfpathclose%
\pgfpathmoveto{\pgfqpoint{9.766054in}{8.803634in}}%
\pgfpathlineto{\pgfqpoint{9.766054in}{8.792418in}}%
\pgfpathquadraticcurveto{\pgfqpoint{9.761044in}{8.794995in}}{\pgfqpoint{9.755622in}{8.796274in}}%
\pgfpathquadraticcurveto{\pgfqpoint{9.750220in}{8.797572in}}{\pgfqpoint{9.744407in}{8.797572in}}%
\pgfpathquadraticcurveto{\pgfqpoint{9.735583in}{8.797572in}}{\pgfqpoint{9.731171in}{8.794872in}}%
\pgfpathquadraticcurveto{\pgfqpoint{9.726759in}{8.792171in}}{\pgfqpoint{9.726759in}{8.786749in}}%
\pgfpathquadraticcurveto{\pgfqpoint{9.726759in}{8.782626in}}{\pgfqpoint{9.729913in}{8.780275in}}%
\pgfpathquadraticcurveto{\pgfqpoint{9.733068in}{8.777925in}}{\pgfqpoint{9.742613in}{8.775802in}}%
\pgfpathlineto{\pgfqpoint{9.746674in}{8.774894in}}%
\pgfpathquadraticcurveto{\pgfqpoint{9.759292in}{8.772194in}}{\pgfqpoint{9.764611in}{8.767266in}}%
\pgfpathquadraticcurveto{\pgfqpoint{9.769930in}{8.762339in}}{\pgfqpoint{9.769930in}{8.753515in}}%
\pgfpathquadraticcurveto{\pgfqpoint{9.769930in}{8.743455in}}{\pgfqpoint{9.761972in}{8.737579in}}%
\pgfpathquadraticcurveto{\pgfqpoint{9.754014in}{8.731724in}}{\pgfqpoint{9.740098in}{8.731724in}}%
\pgfpathquadraticcurveto{\pgfqpoint{9.734305in}{8.731724in}}{\pgfqpoint{9.728017in}{8.732858in}}%
\pgfpathquadraticcurveto{\pgfqpoint{9.721729in}{8.733992in}}{\pgfqpoint{9.714781in}{8.736239in}}%
\pgfpathlineto{\pgfqpoint{9.714781in}{8.748485in}}%
\pgfpathquadraticcurveto{\pgfqpoint{9.721358in}{8.745063in}}{\pgfqpoint{9.727728in}{8.743352in}}%
\pgfpathquadraticcurveto{\pgfqpoint{9.734099in}{8.741661in}}{\pgfqpoint{9.740366in}{8.741661in}}%
\pgfpathquadraticcurveto{\pgfqpoint{9.748736in}{8.741661in}}{\pgfqpoint{9.753230in}{8.744527in}}%
\pgfpathquadraticcurveto{\pgfqpoint{9.757745in}{8.747392in}}{\pgfqpoint{9.757745in}{8.752608in}}%
\pgfpathquadraticcurveto{\pgfqpoint{9.757745in}{8.757432in}}{\pgfqpoint{9.754488in}{8.760010in}}%
\pgfpathquadraticcurveto{\pgfqpoint{9.751251in}{8.762587in}}{\pgfqpoint{9.740222in}{8.764978in}}%
\pgfpathlineto{\pgfqpoint{9.736098in}{8.765947in}}%
\pgfpathquadraticcurveto{\pgfqpoint{9.725089in}{8.768256in}}{\pgfqpoint{9.720182in}{8.773060in}}%
\pgfpathquadraticcurveto{\pgfqpoint{9.715296in}{8.777863in}}{\pgfqpoint{9.715296in}{8.786233in}}%
\pgfpathquadraticcurveto{\pgfqpoint{9.715296in}{8.796418in}}{\pgfqpoint{9.722512in}{8.801943in}}%
\pgfpathquadraticcurveto{\pgfqpoint{9.729728in}{8.807489in}}{\pgfqpoint{9.743005in}{8.807489in}}%
\pgfpathquadraticcurveto{\pgfqpoint{9.749561in}{8.807489in}}{\pgfqpoint{9.755354in}{8.806520in}}%
\pgfpathquadraticcurveto{\pgfqpoint{9.761168in}{8.805572in}}{\pgfqpoint{9.766054in}{8.803634in}}%
\pgfpathlineto{\pgfqpoint{9.766054in}{8.803634in}}%
\pgfpathclose%
\pgfpathmoveto{\pgfqpoint{9.816749in}{8.797449in}}%
\pgfpathquadraticcurveto{\pgfqpoint{9.807225in}{8.797449in}}{\pgfqpoint{9.801679in}{8.790006in}}%
\pgfpathquadraticcurveto{\pgfqpoint{9.796133in}{8.782564in}}{\pgfqpoint{9.796133in}{8.769617in}}%
\pgfpathquadraticcurveto{\pgfqpoint{9.796133in}{8.756670in}}{\pgfqpoint{9.801638in}{8.749227in}}%
\pgfpathquadraticcurveto{\pgfqpoint{9.807163in}{8.741785in}}{\pgfqpoint{9.816749in}{8.741785in}}%
\pgfpathquadraticcurveto{\pgfqpoint{9.826233in}{8.741785in}}{\pgfqpoint{9.831758in}{8.749248in}}%
\pgfpathquadraticcurveto{\pgfqpoint{9.837304in}{8.756732in}}{\pgfqpoint{9.837304in}{8.769617in}}%
\pgfpathquadraticcurveto{\pgfqpoint{9.837304in}{8.782440in}}{\pgfqpoint{9.831758in}{8.789944in}}%
\pgfpathquadraticcurveto{\pgfqpoint{9.826233in}{8.797449in}}{\pgfqpoint{9.816749in}{8.797449in}}%
\pgfpathlineto{\pgfqpoint{9.816749in}{8.797449in}}%
\pgfpathclose%
\pgfpathmoveto{\pgfqpoint{9.816749in}{8.807489in}}%
\pgfpathquadraticcurveto{\pgfqpoint{9.832212in}{8.807489in}}{\pgfqpoint{9.841035in}{8.797428in}}%
\pgfpathquadraticcurveto{\pgfqpoint{9.849880in}{8.787388in}}{\pgfqpoint{9.849880in}{8.769617in}}%
\pgfpathquadraticcurveto{\pgfqpoint{9.849880in}{8.751907in}}{\pgfqpoint{9.841035in}{8.741805in}}%
\pgfpathquadraticcurveto{\pgfqpoint{9.832212in}{8.731724in}}{\pgfqpoint{9.816749in}{8.731724in}}%
\pgfpathquadraticcurveto{\pgfqpoint{9.801225in}{8.731724in}}{\pgfqpoint{9.792422in}{8.741805in}}%
\pgfpathquadraticcurveto{\pgfqpoint{9.783640in}{8.751907in}}{\pgfqpoint{9.783640in}{8.769617in}}%
\pgfpathquadraticcurveto{\pgfqpoint{9.783640in}{8.787388in}}{\pgfqpoint{9.792422in}{8.797428in}}%
\pgfpathquadraticcurveto{\pgfqpoint{9.801225in}{8.807489in}}{\pgfqpoint{9.816749in}{8.807489in}}%
\pgfpathlineto{\pgfqpoint{9.816749in}{8.807489in}}%
\pgfpathclose%
\pgfpathmoveto{\pgfqpoint{9.929521in}{8.777162in}}%
\pgfpathlineto{\pgfqpoint{9.929521in}{8.733600in}}%
\pgfpathlineto{\pgfqpoint{9.917666in}{8.733600in}}%
\pgfpathlineto{\pgfqpoint{9.917666in}{8.776771in}}%
\pgfpathquadraticcurveto{\pgfqpoint{9.917666in}{8.787017in}}{\pgfqpoint{9.913667in}{8.792088in}}%
\pgfpathquadraticcurveto{\pgfqpoint{9.909667in}{8.797181in}}{\pgfqpoint{9.901689in}{8.797181in}}%
\pgfpathquadraticcurveto{\pgfqpoint{9.892081in}{8.797181in}}{\pgfqpoint{9.886536in}{8.791058in}}%
\pgfpathquadraticcurveto{\pgfqpoint{9.880990in}{8.784955in}}{\pgfqpoint{9.880990in}{8.774379in}}%
\pgfpathlineto{\pgfqpoint{9.880990in}{8.733600in}}%
\pgfpathlineto{\pgfqpoint{9.869074in}{8.733600in}}%
\pgfpathlineto{\pgfqpoint{9.869074in}{8.805757in}}%
\pgfpathlineto{\pgfqpoint{9.880990in}{8.805757in}}%
\pgfpathlineto{\pgfqpoint{9.880990in}{8.794542in}}%
\pgfpathquadraticcurveto{\pgfqpoint{9.885257in}{8.801057in}}{\pgfqpoint{9.891009in}{8.804273in}}%
\pgfpathquadraticcurveto{\pgfqpoint{9.896782in}{8.807489in}}{\pgfqpoint{9.904327in}{8.807489in}}%
\pgfpathquadraticcurveto{\pgfqpoint{9.916759in}{8.807489in}}{\pgfqpoint{9.923130in}{8.799799in}}%
\pgfpathquadraticcurveto{\pgfqpoint{9.929521in}{8.792109in}}{\pgfqpoint{9.929521in}{8.777162in}}%
\pgfpathlineto{\pgfqpoint{9.929521in}{8.777162in}}%
\pgfpathclose%
\pgfpathmoveto{\pgfqpoint{9.999142in}{8.803634in}}%
\pgfpathlineto{\pgfqpoint{9.999142in}{8.792418in}}%
\pgfpathquadraticcurveto{\pgfqpoint{9.994132in}{8.794995in}}{\pgfqpoint{9.988710in}{8.796274in}}%
\pgfpathquadraticcurveto{\pgfqpoint{9.983309in}{8.797572in}}{\pgfqpoint{9.977495in}{8.797572in}}%
\pgfpathquadraticcurveto{\pgfqpoint{9.968671in}{8.797572in}}{\pgfqpoint{9.964259in}{8.794872in}}%
\pgfpathquadraticcurveto{\pgfqpoint{9.959847in}{8.792171in}}{\pgfqpoint{9.959847in}{8.786749in}}%
\pgfpathquadraticcurveto{\pgfqpoint{9.959847in}{8.782626in}}{\pgfqpoint{9.963001in}{8.780275in}}%
\pgfpathquadraticcurveto{\pgfqpoint{9.966156in}{8.777925in}}{\pgfqpoint{9.975701in}{8.775802in}}%
\pgfpathlineto{\pgfqpoint{9.979763in}{8.774894in}}%
\pgfpathquadraticcurveto{\pgfqpoint{9.992380in}{8.772194in}}{\pgfqpoint{9.997699in}{8.767266in}}%
\pgfpathquadraticcurveto{\pgfqpoint{10.003018in}{8.762339in}}{\pgfqpoint{10.003018in}{8.753515in}}%
\pgfpathquadraticcurveto{\pgfqpoint{10.003018in}{8.743455in}}{\pgfqpoint{9.995060in}{8.737579in}}%
\pgfpathquadraticcurveto{\pgfqpoint{9.987102in}{8.731724in}}{\pgfqpoint{9.973186in}{8.731724in}}%
\pgfpathquadraticcurveto{\pgfqpoint{9.967393in}{8.731724in}}{\pgfqpoint{9.961105in}{8.732858in}}%
\pgfpathquadraticcurveto{\pgfqpoint{9.954817in}{8.733992in}}{\pgfqpoint{9.947869in}{8.736239in}}%
\pgfpathlineto{\pgfqpoint{9.947869in}{8.748485in}}%
\pgfpathquadraticcurveto{\pgfqpoint{9.954446in}{8.745063in}}{\pgfqpoint{9.960816in}{8.743352in}}%
\pgfpathquadraticcurveto{\pgfqpoint{9.967187in}{8.741661in}}{\pgfqpoint{9.973454in}{8.741661in}}%
\pgfpathquadraticcurveto{\pgfqpoint{9.981824in}{8.741661in}}{\pgfqpoint{9.986319in}{8.744527in}}%
\pgfpathquadraticcurveto{\pgfqpoint{9.990834in}{8.747392in}}{\pgfqpoint{9.990834in}{8.752608in}}%
\pgfpathquadraticcurveto{\pgfqpoint{9.990834in}{8.757432in}}{\pgfqpoint{9.987576in}{8.760010in}}%
\pgfpathquadraticcurveto{\pgfqpoint{9.984339in}{8.762587in}}{\pgfqpoint{9.973310in}{8.764978in}}%
\pgfpathlineto{\pgfqpoint{9.969186in}{8.765947in}}%
\pgfpathquadraticcurveto{\pgfqpoint{9.958177in}{8.768256in}}{\pgfqpoint{9.953271in}{8.773060in}}%
\pgfpathquadraticcurveto{\pgfqpoint{9.948385in}{8.777863in}}{\pgfqpoint{9.948385in}{8.786233in}}%
\pgfpathquadraticcurveto{\pgfqpoint{9.948385in}{8.796418in}}{\pgfqpoint{9.955600in}{8.801943in}}%
\pgfpathquadraticcurveto{\pgfqpoint{9.962816in}{8.807489in}}{\pgfqpoint{9.976093in}{8.807489in}}%
\pgfpathquadraticcurveto{\pgfqpoint{9.982649in}{8.807489in}}{\pgfqpoint{9.988442in}{8.806520in}}%
\pgfpathquadraticcurveto{\pgfqpoint{9.994256in}{8.805572in}}{\pgfqpoint{9.999142in}{8.803634in}}%
\pgfpathlineto{\pgfqpoint{9.999142in}{8.803634in}}%
\pgfpathclose%
\pgfpathmoveto{\pgfqpoint{10.024912in}{8.801819in}}%
\pgfpathlineto{\pgfqpoint{10.038498in}{8.801819in}}%
\pgfpathlineto{\pgfqpoint{10.038498in}{8.785471in}}%
\pgfpathlineto{\pgfqpoint{10.024912in}{8.785471in}}%
\pgfpathlineto{\pgfqpoint{10.024912in}{8.801819in}}%
\pgfpathlineto{\pgfqpoint{10.024912in}{8.801819in}}%
\pgfpathclose%
\pgfpathmoveto{\pgfqpoint{10.024912in}{8.749969in}}%
\pgfpathlineto{\pgfqpoint{10.038498in}{8.749969in}}%
\pgfpathlineto{\pgfqpoint{10.038498in}{8.738878in}}%
\pgfpathlineto{\pgfqpoint{10.027943in}{8.718261in}}%
\pgfpathlineto{\pgfqpoint{10.019635in}{8.718261in}}%
\pgfpathlineto{\pgfqpoint{10.024912in}{8.738878in}}%
\pgfpathlineto{\pgfqpoint{10.024912in}{8.749969in}}%
\pgfpathlineto{\pgfqpoint{10.024912in}{8.749969in}}%
\pgfpathclose%
\pgfpathmoveto{\pgfqpoint{10.145703in}{8.818457in}}%
\pgfpathlineto{\pgfqpoint{10.112841in}{8.767102in}}%
\pgfpathlineto{\pgfqpoint{10.145703in}{8.767102in}}%
\pgfpathlineto{\pgfqpoint{10.145703in}{8.818457in}}%
\pgfpathlineto{\pgfqpoint{10.145703in}{8.818457in}}%
\pgfpathclose%
\pgfpathmoveto{\pgfqpoint{10.142281in}{8.829796in}}%
\pgfpathlineto{\pgfqpoint{10.158650in}{8.829796in}}%
\pgfpathlineto{\pgfqpoint{10.158650in}{8.767102in}}%
\pgfpathlineto{\pgfqpoint{10.172381in}{8.767102in}}%
\pgfpathlineto{\pgfqpoint{10.172381in}{8.756278in}}%
\pgfpathlineto{\pgfqpoint{10.158650in}{8.756278in}}%
\pgfpathlineto{\pgfqpoint{10.158650in}{8.733600in}}%
\pgfpathlineto{\pgfqpoint{10.145703in}{8.733600in}}%
\pgfpathlineto{\pgfqpoint{10.145703in}{8.756278in}}%
\pgfpathlineto{\pgfqpoint{10.102285in}{8.756278in}}%
\pgfpathlineto{\pgfqpoint{10.102285in}{8.768833in}}%
\pgfpathlineto{\pgfqpoint{10.142281in}{8.829796in}}%
\pgfpathlineto{\pgfqpoint{10.142281in}{8.829796in}}%
\pgfpathclose%
\pgfpathmoveto{\pgfqpoint{10.196151in}{8.744547in}}%
\pgfpathlineto{\pgfqpoint{10.217407in}{8.744547in}}%
\pgfpathlineto{\pgfqpoint{10.217407in}{8.817941in}}%
\pgfpathlineto{\pgfqpoint{10.194275in}{8.813303in}}%
\pgfpathlineto{\pgfqpoint{10.194275in}{8.825157in}}%
\pgfpathlineto{\pgfqpoint{10.217283in}{8.829796in}}%
\pgfpathlineto{\pgfqpoint{10.230292in}{8.829796in}}%
\pgfpathlineto{\pgfqpoint{10.230292in}{8.744547in}}%
\pgfpathlineto{\pgfqpoint{10.251548in}{8.744547in}}%
\pgfpathlineto{\pgfqpoint{10.251548in}{8.733600in}}%
\pgfpathlineto{\pgfqpoint{10.196151in}{8.733600in}}%
\pgfpathlineto{\pgfqpoint{10.196151in}{8.744547in}}%
\pgfpathlineto{\pgfqpoint{10.196151in}{8.744547in}}%
\pgfpathclose%
\pgfpathmoveto{\pgfqpoint{10.370029in}{8.802995in}}%
\pgfpathlineto{\pgfqpoint{10.370029in}{8.791903in}}%
\pgfpathquadraticcurveto{\pgfqpoint{10.364999in}{8.794686in}}{\pgfqpoint{10.359948in}{8.796067in}}%
\pgfpathquadraticcurveto{\pgfqpoint{10.354897in}{8.797449in}}{\pgfqpoint{10.349743in}{8.797449in}}%
\pgfpathquadraticcurveto{\pgfqpoint{10.338198in}{8.797449in}}{\pgfqpoint{10.331807in}{8.790130in}}%
\pgfpathquadraticcurveto{\pgfqpoint{10.325436in}{8.782832in}}{\pgfqpoint{10.325436in}{8.769617in}}%
\pgfpathquadraticcurveto{\pgfqpoint{10.325436in}{8.756402in}}{\pgfqpoint{10.331807in}{8.749083in}}%
\pgfpathquadraticcurveto{\pgfqpoint{10.338198in}{8.741785in}}{\pgfqpoint{10.349743in}{8.741785in}}%
\pgfpathquadraticcurveto{\pgfqpoint{10.354897in}{8.741785in}}{\pgfqpoint{10.359948in}{8.743166in}}%
\pgfpathquadraticcurveto{\pgfqpoint{10.364999in}{8.744547in}}{\pgfqpoint{10.370029in}{8.747330in}}%
\pgfpathlineto{\pgfqpoint{10.370029in}{8.736363in}}%
\pgfpathquadraticcurveto{\pgfqpoint{10.365061in}{8.734054in}}{\pgfqpoint{10.359742in}{8.732899in}}%
\pgfpathquadraticcurveto{\pgfqpoint{10.354444in}{8.731724in}}{\pgfqpoint{10.348444in}{8.731724in}}%
\pgfpathquadraticcurveto{\pgfqpoint{10.332137in}{8.731724in}}{\pgfqpoint{10.322529in}{8.741970in}}%
\pgfpathquadraticcurveto{\pgfqpoint{10.312943in}{8.752217in}}{\pgfqpoint{10.312943in}{8.769617in}}%
\pgfpathquadraticcurveto{\pgfqpoint{10.312943in}{8.787264in}}{\pgfqpoint{10.322633in}{8.797366in}}%
\pgfpathquadraticcurveto{\pgfqpoint{10.332343in}{8.807489in}}{\pgfqpoint{10.349228in}{8.807489in}}%
\pgfpathquadraticcurveto{\pgfqpoint{10.354691in}{8.807489in}}{\pgfqpoint{10.359907in}{8.806355in}}%
\pgfpathquadraticcurveto{\pgfqpoint{10.365123in}{8.805242in}}{\pgfqpoint{10.370029in}{8.802995in}}%
\pgfpathlineto{\pgfqpoint{10.370029in}{8.802995in}}%
\pgfpathclose%
\pgfpathmoveto{\pgfqpoint{10.423446in}{8.769864in}}%
\pgfpathquadraticcurveto{\pgfqpoint{10.409077in}{8.769864in}}{\pgfqpoint{10.403531in}{8.766586in}}%
\pgfpathquadraticcurveto{\pgfqpoint{10.397985in}{8.763308in}}{\pgfqpoint{10.397985in}{8.755371in}}%
\pgfpathquadraticcurveto{\pgfqpoint{10.397985in}{8.749062in}}{\pgfqpoint{10.402150in}{8.745351in}}%
\pgfpathquadraticcurveto{\pgfqpoint{10.406314in}{8.741661in}}{\pgfqpoint{10.413447in}{8.741661in}}%
\pgfpathquadraticcurveto{\pgfqpoint{10.423323in}{8.741661in}}{\pgfqpoint{10.429281in}{8.748650in}}%
\pgfpathquadraticcurveto{\pgfqpoint{10.435239in}{8.755639in}}{\pgfqpoint{10.435239in}{8.767225in}}%
\pgfpathlineto{\pgfqpoint{10.435239in}{8.769864in}}%
\pgfpathlineto{\pgfqpoint{10.423446in}{8.769864in}}%
\pgfpathlineto{\pgfqpoint{10.423446in}{8.769864in}}%
\pgfpathclose%
\pgfpathmoveto{\pgfqpoint{10.447093in}{8.774771in}}%
\pgfpathlineto{\pgfqpoint{10.447093in}{8.733600in}}%
\pgfpathlineto{\pgfqpoint{10.435239in}{8.733600in}}%
\pgfpathlineto{\pgfqpoint{10.435239in}{8.744547in}}%
\pgfpathquadraticcurveto{\pgfqpoint{10.431177in}{8.737991in}}{\pgfqpoint{10.425116in}{8.734858in}}%
\pgfpathquadraticcurveto{\pgfqpoint{10.419055in}{8.731724in}}{\pgfqpoint{10.410293in}{8.731724in}}%
\pgfpathquadraticcurveto{\pgfqpoint{10.399222in}{8.731724in}}{\pgfqpoint{10.392666in}{8.737950in}}%
\pgfpathquadraticcurveto{\pgfqpoint{10.386131in}{8.744176in}}{\pgfqpoint{10.386131in}{8.754608in}}%
\pgfpathquadraticcurveto{\pgfqpoint{10.386131in}{8.766772in}}{\pgfqpoint{10.394274in}{8.772957in}}%
\pgfpathquadraticcurveto{\pgfqpoint{10.402438in}{8.779141in}}{\pgfqpoint{10.418602in}{8.779141in}}%
\pgfpathlineto{\pgfqpoint{10.435239in}{8.779141in}}%
\pgfpathlineto{\pgfqpoint{10.435239in}{8.780317in}}%
\pgfpathquadraticcurveto{\pgfqpoint{10.435239in}{8.788501in}}{\pgfqpoint{10.429858in}{8.792975in}}%
\pgfpathquadraticcurveto{\pgfqpoint{10.424477in}{8.797449in}}{\pgfqpoint{10.414746in}{8.797449in}}%
\pgfpathquadraticcurveto{\pgfqpoint{10.408561in}{8.797449in}}{\pgfqpoint{10.402686in}{8.795964in}}%
\pgfpathquadraticcurveto{\pgfqpoint{10.396831in}{8.794480in}}{\pgfqpoint{10.391429in}{8.791511in}}%
\pgfpathlineto{\pgfqpoint{10.391429in}{8.802479in}}%
\pgfpathquadraticcurveto{\pgfqpoint{10.397923in}{8.804994in}}{\pgfqpoint{10.404046in}{8.806231in}}%
\pgfpathquadraticcurveto{\pgfqpoint{10.410169in}{8.807489in}}{\pgfqpoint{10.415963in}{8.807489in}}%
\pgfpathquadraticcurveto{\pgfqpoint{10.431631in}{8.807489in}}{\pgfqpoint{10.439362in}{8.799366in}}%
\pgfpathquadraticcurveto{\pgfqpoint{10.447093in}{8.791264in}}{\pgfqpoint{10.447093in}{8.774771in}}%
\pgfpathlineto{\pgfqpoint{10.447093in}{8.774771in}}%
\pgfpathclose%
\pgfpathmoveto{\pgfqpoint{10.482966in}{8.744424in}}%
\pgfpathlineto{\pgfqpoint{10.482966in}{8.706160in}}%
\pgfpathlineto{\pgfqpoint{10.471049in}{8.706160in}}%
\pgfpathlineto{\pgfqpoint{10.471049in}{8.805757in}}%
\pgfpathlineto{\pgfqpoint{10.482966in}{8.805757in}}%
\pgfpathlineto{\pgfqpoint{10.482966in}{8.794810in}}%
\pgfpathquadraticcurveto{\pgfqpoint{10.486718in}{8.801242in}}{\pgfqpoint{10.492408in}{8.804355in}}%
\pgfpathquadraticcurveto{\pgfqpoint{10.498119in}{8.807489in}}{\pgfqpoint{10.506035in}{8.807489in}}%
\pgfpathquadraticcurveto{\pgfqpoint{10.519189in}{8.807489in}}{\pgfqpoint{10.527394in}{8.797057in}}%
\pgfpathquadraticcurveto{\pgfqpoint{10.535620in}{8.786625in}}{\pgfqpoint{10.535620in}{8.769617in}}%
\pgfpathquadraticcurveto{\pgfqpoint{10.535620in}{8.752608in}}{\pgfqpoint{10.527394in}{8.742156in}}%
\pgfpathquadraticcurveto{\pgfqpoint{10.519189in}{8.731724in}}{\pgfqpoint{10.506035in}{8.731724in}}%
\pgfpathquadraticcurveto{\pgfqpoint{10.498119in}{8.731724in}}{\pgfqpoint{10.492408in}{8.734858in}}%
\pgfpathquadraticcurveto{\pgfqpoint{10.486718in}{8.737991in}}{\pgfqpoint{10.482966in}{8.744424in}}%
\pgfpathlineto{\pgfqpoint{10.482966in}{8.744424in}}%
\pgfpathclose%
\pgfpathmoveto{\pgfqpoint{10.523312in}{8.769617in}}%
\pgfpathquadraticcurveto{\pgfqpoint{10.523312in}{8.782687in}}{\pgfqpoint{10.517931in}{8.790130in}}%
\pgfpathquadraticcurveto{\pgfqpoint{10.512550in}{8.797572in}}{\pgfqpoint{10.503149in}{8.797572in}}%
\pgfpathquadraticcurveto{\pgfqpoint{10.493727in}{8.797572in}}{\pgfqpoint{10.488347in}{8.790130in}}%
\pgfpathquadraticcurveto{\pgfqpoint{10.482966in}{8.782687in}}{\pgfqpoint{10.482966in}{8.769617in}}%
\pgfpathquadraticcurveto{\pgfqpoint{10.482966in}{8.756546in}}{\pgfqpoint{10.488347in}{8.749103in}}%
\pgfpathquadraticcurveto{\pgfqpoint{10.493727in}{8.741661in}}{\pgfqpoint{10.503149in}{8.741661in}}%
\pgfpathquadraticcurveto{\pgfqpoint{10.512550in}{8.741661in}}{\pgfqpoint{10.517931in}{8.749103in}}%
\pgfpathquadraticcurveto{\pgfqpoint{10.523312in}{8.756546in}}{\pgfqpoint{10.523312in}{8.769617in}}%
\pgfpathlineto{\pgfqpoint{10.523312in}{8.769617in}}%
\pgfpathclose%
\pgfpathmoveto{\pgfqpoint{10.566730in}{8.744424in}}%
\pgfpathlineto{\pgfqpoint{10.566730in}{8.706160in}}%
\pgfpathlineto{\pgfqpoint{10.554814in}{8.706160in}}%
\pgfpathlineto{\pgfqpoint{10.554814in}{8.805757in}}%
\pgfpathlineto{\pgfqpoint{10.566730in}{8.805757in}}%
\pgfpathlineto{\pgfqpoint{10.566730in}{8.794810in}}%
\pgfpathquadraticcurveto{\pgfqpoint{10.570482in}{8.801242in}}{\pgfqpoint{10.576172in}{8.804355in}}%
\pgfpathquadraticcurveto{\pgfqpoint{10.581883in}{8.807489in}}{\pgfqpoint{10.589799in}{8.807489in}}%
\pgfpathquadraticcurveto{\pgfqpoint{10.602953in}{8.807489in}}{\pgfqpoint{10.611158in}{8.797057in}}%
\pgfpathquadraticcurveto{\pgfqpoint{10.619384in}{8.786625in}}{\pgfqpoint{10.619384in}{8.769617in}}%
\pgfpathquadraticcurveto{\pgfqpoint{10.619384in}{8.752608in}}{\pgfqpoint{10.611158in}{8.742156in}}%
\pgfpathquadraticcurveto{\pgfqpoint{10.602953in}{8.731724in}}{\pgfqpoint{10.589799in}{8.731724in}}%
\pgfpathquadraticcurveto{\pgfqpoint{10.581883in}{8.731724in}}{\pgfqpoint{10.576172in}{8.734858in}}%
\pgfpathquadraticcurveto{\pgfqpoint{10.570482in}{8.737991in}}{\pgfqpoint{10.566730in}{8.744424in}}%
\pgfpathlineto{\pgfqpoint{10.566730in}{8.744424in}}%
\pgfpathclose%
\pgfpathmoveto{\pgfqpoint{10.607076in}{8.769617in}}%
\pgfpathquadraticcurveto{\pgfqpoint{10.607076in}{8.782687in}}{\pgfqpoint{10.601695in}{8.790130in}}%
\pgfpathquadraticcurveto{\pgfqpoint{10.596314in}{8.797572in}}{\pgfqpoint{10.586913in}{8.797572in}}%
\pgfpathquadraticcurveto{\pgfqpoint{10.577492in}{8.797572in}}{\pgfqpoint{10.572111in}{8.790130in}}%
\pgfpathquadraticcurveto{\pgfqpoint{10.566730in}{8.782687in}}{\pgfqpoint{10.566730in}{8.769617in}}%
\pgfpathquadraticcurveto{\pgfqpoint{10.566730in}{8.756546in}}{\pgfqpoint{10.572111in}{8.749103in}}%
\pgfpathquadraticcurveto{\pgfqpoint{10.577492in}{8.741661in}}{\pgfqpoint{10.586913in}{8.741661in}}%
\pgfpathquadraticcurveto{\pgfqpoint{10.596314in}{8.741661in}}{\pgfqpoint{10.601695in}{8.749103in}}%
\pgfpathquadraticcurveto{\pgfqpoint{10.607076in}{8.756546in}}{\pgfqpoint{10.607076in}{8.769617in}}%
\pgfpathlineto{\pgfqpoint{10.607076in}{8.769617in}}%
\pgfpathclose%
\pgfpathmoveto{\pgfqpoint{10.700756in}{8.772647in}}%
\pgfpathlineto{\pgfqpoint{10.700756in}{8.766854in}}%
\pgfpathlineto{\pgfqpoint{10.646247in}{8.766854in}}%
\pgfpathquadraticcurveto{\pgfqpoint{10.647030in}{8.754608in}}{\pgfqpoint{10.653628in}{8.748196in}}%
\pgfpathquadraticcurveto{\pgfqpoint{10.660225in}{8.741785in}}{\pgfqpoint{10.672017in}{8.741785in}}%
\pgfpathquadraticcurveto{\pgfqpoint{10.678841in}{8.741785in}}{\pgfqpoint{10.685253in}{8.743455in}}%
\pgfpathquadraticcurveto{\pgfqpoint{10.691665in}{8.745125in}}{\pgfqpoint{10.697994in}{8.748485in}}%
\pgfpathlineto{\pgfqpoint{10.697994in}{8.737270in}}%
\pgfpathquadraticcurveto{\pgfqpoint{10.691603in}{8.734569in}}{\pgfqpoint{10.684903in}{8.733146in}}%
\pgfpathquadraticcurveto{\pgfqpoint{10.678202in}{8.731724in}}{\pgfqpoint{10.671316in}{8.731724in}}%
\pgfpathquadraticcurveto{\pgfqpoint{10.654040in}{8.731724in}}{\pgfqpoint{10.643959in}{8.741764in}}%
\pgfpathquadraticcurveto{\pgfqpoint{10.633877in}{8.751825in}}{\pgfqpoint{10.633877in}{8.768978in}}%
\pgfpathquadraticcurveto{\pgfqpoint{10.633877in}{8.786687in}}{\pgfqpoint{10.643443in}{8.797078in}}%
\pgfpathquadraticcurveto{\pgfqpoint{10.653009in}{8.807489in}}{\pgfqpoint{10.669255in}{8.807489in}}%
\pgfpathquadraticcurveto{\pgfqpoint{10.683810in}{8.807489in}}{\pgfqpoint{10.692283in}{8.798108in}}%
\pgfpathquadraticcurveto{\pgfqpoint{10.700756in}{8.788749in}}{\pgfqpoint{10.700756in}{8.772647in}}%
\pgfpathlineto{\pgfqpoint{10.700756in}{8.772647in}}%
\pgfpathclose%
\pgfpathmoveto{\pgfqpoint{10.688902in}{8.776131in}}%
\pgfpathquadraticcurveto{\pgfqpoint{10.688778in}{8.785842in}}{\pgfqpoint{10.683459in}{8.791635in}}%
\pgfpathquadraticcurveto{\pgfqpoint{10.678140in}{8.797449in}}{\pgfqpoint{10.669378in}{8.797449in}}%
\pgfpathquadraticcurveto{\pgfqpoint{10.659462in}{8.797449in}}{\pgfqpoint{10.653504in}{8.791841in}}%
\pgfpathquadraticcurveto{\pgfqpoint{10.647546in}{8.786233in}}{\pgfqpoint{10.646639in}{8.776049in}}%
\pgfpathlineto{\pgfqpoint{10.688902in}{8.776131in}}%
\pgfpathlineto{\pgfqpoint{10.688902in}{8.776131in}}%
\pgfpathclose%
\pgfpathmoveto{\pgfqpoint{10.767698in}{8.794810in}}%
\pgfpathlineto{\pgfqpoint{10.767698in}{8.833857in}}%
\pgfpathlineto{\pgfqpoint{10.779552in}{8.833857in}}%
\pgfpathlineto{\pgfqpoint{10.779552in}{8.733600in}}%
\pgfpathlineto{\pgfqpoint{10.767698in}{8.733600in}}%
\pgfpathlineto{\pgfqpoint{10.767698in}{8.744424in}}%
\pgfpathquadraticcurveto{\pgfqpoint{10.763966in}{8.737991in}}{\pgfqpoint{10.758255in}{8.734858in}}%
\pgfpathquadraticcurveto{\pgfqpoint{10.752565in}{8.731724in}}{\pgfqpoint{10.744566in}{8.731724in}}%
\pgfpathquadraticcurveto{\pgfqpoint{10.731495in}{8.731724in}}{\pgfqpoint{10.723270in}{8.742156in}}%
\pgfpathquadraticcurveto{\pgfqpoint{10.715064in}{8.752608in}}{\pgfqpoint{10.715064in}{8.769617in}}%
\pgfpathquadraticcurveto{\pgfqpoint{10.715064in}{8.786625in}}{\pgfqpoint{10.723270in}{8.797057in}}%
\pgfpathquadraticcurveto{\pgfqpoint{10.731495in}{8.807489in}}{\pgfqpoint{10.744566in}{8.807489in}}%
\pgfpathquadraticcurveto{\pgfqpoint{10.752565in}{8.807489in}}{\pgfqpoint{10.758255in}{8.804355in}}%
\pgfpathquadraticcurveto{\pgfqpoint{10.763966in}{8.801242in}}{\pgfqpoint{10.767698in}{8.794810in}}%
\pgfpathlineto{\pgfqpoint{10.767698in}{8.794810in}}%
\pgfpathclose%
\pgfpathmoveto{\pgfqpoint{10.727310in}{8.769617in}}%
\pgfpathquadraticcurveto{\pgfqpoint{10.727310in}{8.756546in}}{\pgfqpoint{10.732691in}{8.749103in}}%
\pgfpathquadraticcurveto{\pgfqpoint{10.738072in}{8.741661in}}{\pgfqpoint{10.747473in}{8.741661in}}%
\pgfpathquadraticcurveto{\pgfqpoint{10.756874in}{8.741661in}}{\pgfqpoint{10.762276in}{8.749103in}}%
\pgfpathquadraticcurveto{\pgfqpoint{10.767698in}{8.756546in}}{\pgfqpoint{10.767698in}{8.769617in}}%
\pgfpathquadraticcurveto{\pgfqpoint{10.767698in}{8.782687in}}{\pgfqpoint{10.762276in}{8.790130in}}%
\pgfpathquadraticcurveto{\pgfqpoint{10.756874in}{8.797572in}}{\pgfqpoint{10.747473in}{8.797572in}}%
\pgfpathquadraticcurveto{\pgfqpoint{10.738072in}{8.797572in}}{\pgfqpoint{10.732691in}{8.790130in}}%
\pgfpathquadraticcurveto{\pgfqpoint{10.727310in}{8.782687in}}{\pgfqpoint{10.727310in}{8.769617in}}%
\pgfpathlineto{\pgfqpoint{10.727310in}{8.769617in}}%
\pgfpathclose%
\pgfpathmoveto{\pgfqpoint{10.874387in}{8.833713in}}%
\pgfpathquadraticcurveto{\pgfqpoint{10.865770in}{8.818910in}}{\pgfqpoint{10.861564in}{8.804396in}}%
\pgfpathquadraticcurveto{\pgfqpoint{10.857379in}{8.789903in}}{\pgfqpoint{10.857379in}{8.775018in}}%
\pgfpathquadraticcurveto{\pgfqpoint{10.857379in}{8.760154in}}{\pgfqpoint{10.861605in}{8.745557in}}%
\pgfpathquadraticcurveto{\pgfqpoint{10.865831in}{8.730961in}}{\pgfqpoint{10.874387in}{8.716200in}}%
\pgfpathlineto{\pgfqpoint{10.864079in}{8.716200in}}%
\pgfpathquadraticcurveto{\pgfqpoint{10.854431in}{8.731353in}}{\pgfqpoint{10.849627in}{8.745970in}}%
\pgfpathquadraticcurveto{\pgfqpoint{10.844823in}{8.760587in}}{\pgfqpoint{10.844823in}{8.775018in}}%
\pgfpathquadraticcurveto{\pgfqpoint{10.844823in}{8.789388in}}{\pgfqpoint{10.849586in}{8.803943in}}%
\pgfpathquadraticcurveto{\pgfqpoint{10.854369in}{8.818519in}}{\pgfqpoint{10.864079in}{8.833713in}}%
\pgfpathlineto{\pgfqpoint{10.874387in}{8.833713in}}%
\pgfpathlineto{\pgfqpoint{10.874387in}{8.833713in}}%
\pgfpathclose%
\pgfpathmoveto{\pgfqpoint{10.944874in}{8.794810in}}%
\pgfpathlineto{\pgfqpoint{10.944874in}{8.833857in}}%
\pgfpathlineto{\pgfqpoint{10.956729in}{8.833857in}}%
\pgfpathlineto{\pgfqpoint{10.956729in}{8.733600in}}%
\pgfpathlineto{\pgfqpoint{10.944874in}{8.733600in}}%
\pgfpathlineto{\pgfqpoint{10.944874in}{8.744424in}}%
\pgfpathquadraticcurveto{\pgfqpoint{10.941143in}{8.737991in}}{\pgfqpoint{10.935432in}{8.734858in}}%
\pgfpathquadraticcurveto{\pgfqpoint{10.929742in}{8.731724in}}{\pgfqpoint{10.921743in}{8.731724in}}%
\pgfpathquadraticcurveto{\pgfqpoint{10.908672in}{8.731724in}}{\pgfqpoint{10.900446in}{8.742156in}}%
\pgfpathquadraticcurveto{\pgfqpoint{10.892241in}{8.752608in}}{\pgfqpoint{10.892241in}{8.769617in}}%
\pgfpathquadraticcurveto{\pgfqpoint{10.892241in}{8.786625in}}{\pgfqpoint{10.900446in}{8.797057in}}%
\pgfpathquadraticcurveto{\pgfqpoint{10.908672in}{8.807489in}}{\pgfqpoint{10.921743in}{8.807489in}}%
\pgfpathquadraticcurveto{\pgfqpoint{10.929742in}{8.807489in}}{\pgfqpoint{10.935432in}{8.804355in}}%
\pgfpathquadraticcurveto{\pgfqpoint{10.941143in}{8.801242in}}{\pgfqpoint{10.944874in}{8.794810in}}%
\pgfpathlineto{\pgfqpoint{10.944874in}{8.794810in}}%
\pgfpathclose%
\pgfpathmoveto{\pgfqpoint{10.904487in}{8.769617in}}%
\pgfpathquadraticcurveto{\pgfqpoint{10.904487in}{8.756546in}}{\pgfqpoint{10.909868in}{8.749103in}}%
\pgfpathquadraticcurveto{\pgfqpoint{10.915249in}{8.741661in}}{\pgfqpoint{10.924650in}{8.741661in}}%
\pgfpathquadraticcurveto{\pgfqpoint{10.934051in}{8.741661in}}{\pgfqpoint{10.939452in}{8.749103in}}%
\pgfpathquadraticcurveto{\pgfqpoint{10.944874in}{8.756546in}}{\pgfqpoint{10.944874in}{8.769617in}}%
\pgfpathquadraticcurveto{\pgfqpoint{10.944874in}{8.782687in}}{\pgfqpoint{10.939452in}{8.790130in}}%
\pgfpathquadraticcurveto{\pgfqpoint{10.934051in}{8.797572in}}{\pgfqpoint{10.924650in}{8.797572in}}%
\pgfpathquadraticcurveto{\pgfqpoint{10.915249in}{8.797572in}}{\pgfqpoint{10.909868in}{8.790130in}}%
\pgfpathquadraticcurveto{\pgfqpoint{10.904487in}{8.782687in}}{\pgfqpoint{10.904487in}{8.769617in}}%
\pgfpathlineto{\pgfqpoint{10.904487in}{8.769617in}}%
\pgfpathclose%
\pgfpathmoveto{\pgfqpoint{10.981159in}{8.805757in}}%
\pgfpathlineto{\pgfqpoint{10.993013in}{8.805757in}}%
\pgfpathlineto{\pgfqpoint{10.993013in}{8.733600in}}%
\pgfpathlineto{\pgfqpoint{10.981159in}{8.733600in}}%
\pgfpathlineto{\pgfqpoint{10.981159in}{8.805757in}}%
\pgfpathlineto{\pgfqpoint{10.981159in}{8.805757in}}%
\pgfpathclose%
\pgfpathmoveto{\pgfqpoint{10.981159in}{8.833857in}}%
\pgfpathlineto{\pgfqpoint{10.993013in}{8.833857in}}%
\pgfpathlineto{\pgfqpoint{10.993013in}{8.818828in}}%
\pgfpathlineto{\pgfqpoint{10.981159in}{8.818828in}}%
\pgfpathlineto{\pgfqpoint{10.981159in}{8.833857in}}%
\pgfpathlineto{\pgfqpoint{10.981159in}{8.833857in}}%
\pgfpathclose%
\pgfpathmoveto{\pgfqpoint{11.050615in}{8.769864in}}%
\pgfpathquadraticcurveto{\pgfqpoint{11.036246in}{8.769864in}}{\pgfqpoint{11.030700in}{8.766586in}}%
\pgfpathquadraticcurveto{\pgfqpoint{11.025154in}{8.763308in}}{\pgfqpoint{11.025154in}{8.755371in}}%
\pgfpathquadraticcurveto{\pgfqpoint{11.025154in}{8.749062in}}{\pgfqpoint{11.029319in}{8.745351in}}%
\pgfpathquadraticcurveto{\pgfqpoint{11.033483in}{8.741661in}}{\pgfqpoint{11.040617in}{8.741661in}}%
\pgfpathquadraticcurveto{\pgfqpoint{11.050492in}{8.741661in}}{\pgfqpoint{11.056450in}{8.748650in}}%
\pgfpathquadraticcurveto{\pgfqpoint{11.062408in}{8.755639in}}{\pgfqpoint{11.062408in}{8.767225in}}%
\pgfpathlineto{\pgfqpoint{11.062408in}{8.769864in}}%
\pgfpathlineto{\pgfqpoint{11.050615in}{8.769864in}}%
\pgfpathlineto{\pgfqpoint{11.050615in}{8.769864in}}%
\pgfpathclose%
\pgfpathmoveto{\pgfqpoint{11.074262in}{8.774771in}}%
\pgfpathlineto{\pgfqpoint{11.074262in}{8.733600in}}%
\pgfpathlineto{\pgfqpoint{11.062408in}{8.733600in}}%
\pgfpathlineto{\pgfqpoint{11.062408in}{8.744547in}}%
\pgfpathquadraticcurveto{\pgfqpoint{11.058347in}{8.737991in}}{\pgfqpoint{11.052285in}{8.734858in}}%
\pgfpathquadraticcurveto{\pgfqpoint{11.046224in}{8.731724in}}{\pgfqpoint{11.037462in}{8.731724in}}%
\pgfpathquadraticcurveto{\pgfqpoint{11.026391in}{8.731724in}}{\pgfqpoint{11.019835in}{8.737950in}}%
\pgfpathquadraticcurveto{\pgfqpoint{11.013300in}{8.744176in}}{\pgfqpoint{11.013300in}{8.754608in}}%
\pgfpathquadraticcurveto{\pgfqpoint{11.013300in}{8.766772in}}{\pgfqpoint{11.021443in}{8.772957in}}%
\pgfpathquadraticcurveto{\pgfqpoint{11.029607in}{8.779141in}}{\pgfqpoint{11.045771in}{8.779141in}}%
\pgfpathlineto{\pgfqpoint{11.062408in}{8.779141in}}%
\pgfpathlineto{\pgfqpoint{11.062408in}{8.780317in}}%
\pgfpathquadraticcurveto{\pgfqpoint{11.062408in}{8.788501in}}{\pgfqpoint{11.057027in}{8.792975in}}%
\pgfpathquadraticcurveto{\pgfqpoint{11.051646in}{8.797449in}}{\pgfqpoint{11.041915in}{8.797449in}}%
\pgfpathquadraticcurveto{\pgfqpoint{11.035730in}{8.797449in}}{\pgfqpoint{11.029855in}{8.795964in}}%
\pgfpathquadraticcurveto{\pgfqpoint{11.024000in}{8.794480in}}{\pgfqpoint{11.018598in}{8.791511in}}%
\pgfpathlineto{\pgfqpoint{11.018598in}{8.802479in}}%
\pgfpathquadraticcurveto{\pgfqpoint{11.025092in}{8.804994in}}{\pgfqpoint{11.031215in}{8.806231in}}%
\pgfpathquadraticcurveto{\pgfqpoint{11.037339in}{8.807489in}}{\pgfqpoint{11.043132in}{8.807489in}}%
\pgfpathquadraticcurveto{\pgfqpoint{11.058800in}{8.807489in}}{\pgfqpoint{11.066531in}{8.799366in}}%
\pgfpathquadraticcurveto{\pgfqpoint{11.074262in}{8.791264in}}{\pgfqpoint{11.074262in}{8.774771in}}%
\pgfpathlineto{\pgfqpoint{11.074262in}{8.774771in}}%
\pgfpathclose%
\pgfpathmoveto{\pgfqpoint{11.154852in}{8.791903in}}%
\pgfpathquadraticcurveto{\pgfqpoint{11.159305in}{8.799902in}}{\pgfqpoint{11.165490in}{8.803695in}}%
\pgfpathquadraticcurveto{\pgfqpoint{11.171674in}{8.807489in}}{\pgfqpoint{11.180045in}{8.807489in}}%
\pgfpathquadraticcurveto{\pgfqpoint{11.191322in}{8.807489in}}{\pgfqpoint{11.197445in}{8.799593in}}%
\pgfpathquadraticcurveto{\pgfqpoint{11.203568in}{8.791717in}}{\pgfqpoint{11.203568in}{8.777162in}}%
\pgfpathlineto{\pgfqpoint{11.203568in}{8.733600in}}%
\pgfpathlineto{\pgfqpoint{11.191652in}{8.733600in}}%
\pgfpathlineto{\pgfqpoint{11.191652in}{8.776771in}}%
\pgfpathquadraticcurveto{\pgfqpoint{11.191652in}{8.787141in}}{\pgfqpoint{11.187961in}{8.792150in}}%
\pgfpathquadraticcurveto{\pgfqpoint{11.184292in}{8.797181in}}{\pgfqpoint{11.176767in}{8.797181in}}%
\pgfpathquadraticcurveto{\pgfqpoint{11.167551in}{8.797181in}}{\pgfqpoint{11.162191in}{8.791058in}}%
\pgfpathquadraticcurveto{\pgfqpoint{11.156851in}{8.784955in}}{\pgfqpoint{11.156851in}{8.774379in}}%
\pgfpathlineto{\pgfqpoint{11.156851in}{8.733600in}}%
\pgfpathlineto{\pgfqpoint{11.144935in}{8.733600in}}%
\pgfpathlineto{\pgfqpoint{11.144935in}{8.776771in}}%
\pgfpathquadraticcurveto{\pgfqpoint{11.144935in}{8.787202in}}{\pgfqpoint{11.141265in}{8.792192in}}%
\pgfpathquadraticcurveto{\pgfqpoint{11.137596in}{8.797181in}}{\pgfqpoint{11.129926in}{8.797181in}}%
\pgfpathquadraticcurveto{\pgfqpoint{11.120835in}{8.797181in}}{\pgfqpoint{11.115474in}{8.791037in}}%
\pgfpathquadraticcurveto{\pgfqpoint{11.110135in}{8.784893in}}{\pgfqpoint{11.110135in}{8.774379in}}%
\pgfpathlineto{\pgfqpoint{11.110135in}{8.733600in}}%
\pgfpathlineto{\pgfqpoint{11.098219in}{8.733600in}}%
\pgfpathlineto{\pgfqpoint{11.098219in}{8.805757in}}%
\pgfpathlineto{\pgfqpoint{11.110135in}{8.805757in}}%
\pgfpathlineto{\pgfqpoint{11.110135in}{8.794542in}}%
\pgfpathquadraticcurveto{\pgfqpoint{11.114196in}{8.801180in}}{\pgfqpoint{11.119866in}{8.804335in}}%
\pgfpathquadraticcurveto{\pgfqpoint{11.125535in}{8.807489in}}{\pgfqpoint{11.133328in}{8.807489in}}%
\pgfpathquadraticcurveto{\pgfqpoint{11.141204in}{8.807489in}}{\pgfqpoint{11.146708in}{8.803489in}}%
\pgfpathquadraticcurveto{\pgfqpoint{11.152213in}{8.799510in}}{\pgfqpoint{11.154852in}{8.791903in}}%
\pgfpathlineto{\pgfqpoint{11.154852in}{8.791903in}}%
\pgfpathclose%
\pgfpathmoveto{\pgfqpoint{11.255150in}{8.797449in}}%
\pgfpathquadraticcurveto{\pgfqpoint{11.245625in}{8.797449in}}{\pgfqpoint{11.240079in}{8.790006in}}%
\pgfpathquadraticcurveto{\pgfqpoint{11.234534in}{8.782564in}}{\pgfqpoint{11.234534in}{8.769617in}}%
\pgfpathquadraticcurveto{\pgfqpoint{11.234534in}{8.756670in}}{\pgfqpoint{11.240038in}{8.749227in}}%
\pgfpathquadraticcurveto{\pgfqpoint{11.245563in}{8.741785in}}{\pgfqpoint{11.255150in}{8.741785in}}%
\pgfpathquadraticcurveto{\pgfqpoint{11.264633in}{8.741785in}}{\pgfqpoint{11.270159in}{8.749248in}}%
\pgfpathquadraticcurveto{\pgfqpoint{11.275704in}{8.756732in}}{\pgfqpoint{11.275704in}{8.769617in}}%
\pgfpathquadraticcurveto{\pgfqpoint{11.275704in}{8.782440in}}{\pgfqpoint{11.270159in}{8.789944in}}%
\pgfpathquadraticcurveto{\pgfqpoint{11.264633in}{8.797449in}}{\pgfqpoint{11.255150in}{8.797449in}}%
\pgfpathlineto{\pgfqpoint{11.255150in}{8.797449in}}%
\pgfpathclose%
\pgfpathmoveto{\pgfqpoint{11.255150in}{8.807489in}}%
\pgfpathquadraticcurveto{\pgfqpoint{11.270612in}{8.807489in}}{\pgfqpoint{11.279436in}{8.797428in}}%
\pgfpathquadraticcurveto{\pgfqpoint{11.288280in}{8.787388in}}{\pgfqpoint{11.288280in}{8.769617in}}%
\pgfpathquadraticcurveto{\pgfqpoint{11.288280in}{8.751907in}}{\pgfqpoint{11.279436in}{8.741805in}}%
\pgfpathquadraticcurveto{\pgfqpoint{11.270612in}{8.731724in}}{\pgfqpoint{11.255150in}{8.731724in}}%
\pgfpathquadraticcurveto{\pgfqpoint{11.239626in}{8.731724in}}{\pgfqpoint{11.230823in}{8.741805in}}%
\pgfpathquadraticcurveto{\pgfqpoint{11.222040in}{8.751907in}}{\pgfqpoint{11.222040in}{8.769617in}}%
\pgfpathquadraticcurveto{\pgfqpoint{11.222040in}{8.787388in}}{\pgfqpoint{11.230823in}{8.797428in}}%
\pgfpathquadraticcurveto{\pgfqpoint{11.239626in}{8.807489in}}{\pgfqpoint{11.255150in}{8.807489in}}%
\pgfpathlineto{\pgfqpoint{11.255150in}{8.807489in}}%
\pgfpathclose%
\pgfpathmoveto{\pgfqpoint{11.367921in}{8.777162in}}%
\pgfpathlineto{\pgfqpoint{11.367921in}{8.733600in}}%
\pgfpathlineto{\pgfqpoint{11.356067in}{8.733600in}}%
\pgfpathlineto{\pgfqpoint{11.356067in}{8.776771in}}%
\pgfpathquadraticcurveto{\pgfqpoint{11.356067in}{8.787017in}}{\pgfqpoint{11.352067in}{8.792088in}}%
\pgfpathquadraticcurveto{\pgfqpoint{11.348068in}{8.797181in}}{\pgfqpoint{11.340089in}{8.797181in}}%
\pgfpathquadraticcurveto{\pgfqpoint{11.330482in}{8.797181in}}{\pgfqpoint{11.324936in}{8.791058in}}%
\pgfpathquadraticcurveto{\pgfqpoint{11.319390in}{8.784955in}}{\pgfqpoint{11.319390in}{8.774379in}}%
\pgfpathlineto{\pgfqpoint{11.319390in}{8.733600in}}%
\pgfpathlineto{\pgfqpoint{11.307474in}{8.733600in}}%
\pgfpathlineto{\pgfqpoint{11.307474in}{8.805757in}}%
\pgfpathlineto{\pgfqpoint{11.319390in}{8.805757in}}%
\pgfpathlineto{\pgfqpoint{11.319390in}{8.794542in}}%
\pgfpathquadraticcurveto{\pgfqpoint{11.323658in}{8.801057in}}{\pgfqpoint{11.329410in}{8.804273in}}%
\pgfpathquadraticcurveto{\pgfqpoint{11.335182in}{8.807489in}}{\pgfqpoint{11.342728in}{8.807489in}}%
\pgfpathquadraticcurveto{\pgfqpoint{11.355160in}{8.807489in}}{\pgfqpoint{11.361530in}{8.799799in}}%
\pgfpathquadraticcurveto{\pgfqpoint{11.367921in}{8.792109in}}{\pgfqpoint{11.367921in}{8.777162in}}%
\pgfpathlineto{\pgfqpoint{11.367921in}{8.777162in}}%
\pgfpathclose%
\pgfpathmoveto{\pgfqpoint{11.439027in}{8.794810in}}%
\pgfpathlineto{\pgfqpoint{11.439027in}{8.833857in}}%
\pgfpathlineto{\pgfqpoint{11.450881in}{8.833857in}}%
\pgfpathlineto{\pgfqpoint{11.450881in}{8.733600in}}%
\pgfpathlineto{\pgfqpoint{11.439027in}{8.733600in}}%
\pgfpathlineto{\pgfqpoint{11.439027in}{8.744424in}}%
\pgfpathquadraticcurveto{\pgfqpoint{11.435295in}{8.737991in}}{\pgfqpoint{11.429585in}{8.734858in}}%
\pgfpathquadraticcurveto{\pgfqpoint{11.423895in}{8.731724in}}{\pgfqpoint{11.415895in}{8.731724in}}%
\pgfpathquadraticcurveto{\pgfqpoint{11.402825in}{8.731724in}}{\pgfqpoint{11.394599in}{8.742156in}}%
\pgfpathquadraticcurveto{\pgfqpoint{11.386393in}{8.752608in}}{\pgfqpoint{11.386393in}{8.769617in}}%
\pgfpathquadraticcurveto{\pgfqpoint{11.386393in}{8.786625in}}{\pgfqpoint{11.394599in}{8.797057in}}%
\pgfpathquadraticcurveto{\pgfqpoint{11.402825in}{8.807489in}}{\pgfqpoint{11.415895in}{8.807489in}}%
\pgfpathquadraticcurveto{\pgfqpoint{11.423895in}{8.807489in}}{\pgfqpoint{11.429585in}{8.804355in}}%
\pgfpathquadraticcurveto{\pgfqpoint{11.435295in}{8.801242in}}{\pgfqpoint{11.439027in}{8.794810in}}%
\pgfpathlineto{\pgfqpoint{11.439027in}{8.794810in}}%
\pgfpathclose%
\pgfpathmoveto{\pgfqpoint{11.398640in}{8.769617in}}%
\pgfpathquadraticcurveto{\pgfqpoint{11.398640in}{8.756546in}}{\pgfqpoint{11.404020in}{8.749103in}}%
\pgfpathquadraticcurveto{\pgfqpoint{11.409401in}{8.741661in}}{\pgfqpoint{11.418802in}{8.741661in}}%
\pgfpathquadraticcurveto{\pgfqpoint{11.428203in}{8.741661in}}{\pgfqpoint{11.433605in}{8.749103in}}%
\pgfpathquadraticcurveto{\pgfqpoint{11.439027in}{8.756546in}}{\pgfqpoint{11.439027in}{8.769617in}}%
\pgfpathquadraticcurveto{\pgfqpoint{11.439027in}{8.782687in}}{\pgfqpoint{11.433605in}{8.790130in}}%
\pgfpathquadraticcurveto{\pgfqpoint{11.428203in}{8.797572in}}{\pgfqpoint{11.418802in}{8.797572in}}%
\pgfpathquadraticcurveto{\pgfqpoint{11.409401in}{8.797572in}}{\pgfqpoint{11.404020in}{8.790130in}}%
\pgfpathquadraticcurveto{\pgfqpoint{11.398640in}{8.782687in}}{\pgfqpoint{11.398640in}{8.769617in}}%
\pgfpathlineto{\pgfqpoint{11.398640in}{8.769617in}}%
\pgfpathclose%
\pgfpathmoveto{\pgfqpoint{11.521307in}{8.803634in}}%
\pgfpathlineto{\pgfqpoint{11.521307in}{8.792418in}}%
\pgfpathquadraticcurveto{\pgfqpoint{11.516297in}{8.794995in}}{\pgfqpoint{11.510875in}{8.796274in}}%
\pgfpathquadraticcurveto{\pgfqpoint{11.505473in}{8.797572in}}{\pgfqpoint{11.499659in}{8.797572in}}%
\pgfpathquadraticcurveto{\pgfqpoint{11.490836in}{8.797572in}}{\pgfqpoint{11.486424in}{8.794872in}}%
\pgfpathquadraticcurveto{\pgfqpoint{11.482012in}{8.792171in}}{\pgfqpoint{11.482012in}{8.786749in}}%
\pgfpathquadraticcurveto{\pgfqpoint{11.482012in}{8.782626in}}{\pgfqpoint{11.485166in}{8.780275in}}%
\pgfpathquadraticcurveto{\pgfqpoint{11.488321in}{8.777925in}}{\pgfqpoint{11.497866in}{8.775802in}}%
\pgfpathlineto{\pgfqpoint{11.501927in}{8.774894in}}%
\pgfpathquadraticcurveto{\pgfqpoint{11.514544in}{8.772194in}}{\pgfqpoint{11.519863in}{8.767266in}}%
\pgfpathquadraticcurveto{\pgfqpoint{11.525182in}{8.762339in}}{\pgfqpoint{11.525182in}{8.753515in}}%
\pgfpathquadraticcurveto{\pgfqpoint{11.525182in}{8.743455in}}{\pgfqpoint{11.517225in}{8.737579in}}%
\pgfpathquadraticcurveto{\pgfqpoint{11.509267in}{8.731724in}}{\pgfqpoint{11.495351in}{8.731724in}}%
\pgfpathquadraticcurveto{\pgfqpoint{11.489557in}{8.731724in}}{\pgfqpoint{11.483270in}{8.732858in}}%
\pgfpathquadraticcurveto{\pgfqpoint{11.476982in}{8.733992in}}{\pgfqpoint{11.470034in}{8.736239in}}%
\pgfpathlineto{\pgfqpoint{11.470034in}{8.748485in}}%
\pgfpathquadraticcurveto{\pgfqpoint{11.476610in}{8.745063in}}{\pgfqpoint{11.482981in}{8.743352in}}%
\pgfpathquadraticcurveto{\pgfqpoint{11.489351in}{8.741661in}}{\pgfqpoint{11.495619in}{8.741661in}}%
\pgfpathquadraticcurveto{\pgfqpoint{11.503989in}{8.741661in}}{\pgfqpoint{11.508483in}{8.744527in}}%
\pgfpathquadraticcurveto{\pgfqpoint{11.512998in}{8.747392in}}{\pgfqpoint{11.512998in}{8.752608in}}%
\pgfpathquadraticcurveto{\pgfqpoint{11.512998in}{8.757432in}}{\pgfqpoint{11.509741in}{8.760010in}}%
\pgfpathquadraticcurveto{\pgfqpoint{11.506504in}{8.762587in}}{\pgfqpoint{11.495474in}{8.764978in}}%
\pgfpathlineto{\pgfqpoint{11.491351in}{8.765947in}}%
\pgfpathquadraticcurveto{\pgfqpoint{11.480342in}{8.768256in}}{\pgfqpoint{11.475435in}{8.773060in}}%
\pgfpathquadraticcurveto{\pgfqpoint{11.470549in}{8.777863in}}{\pgfqpoint{11.470549in}{8.786233in}}%
\pgfpathquadraticcurveto{\pgfqpoint{11.470549in}{8.796418in}}{\pgfqpoint{11.477765in}{8.801943in}}%
\pgfpathquadraticcurveto{\pgfqpoint{11.484981in}{8.807489in}}{\pgfqpoint{11.498258in}{8.807489in}}%
\pgfpathquadraticcurveto{\pgfqpoint{11.504814in}{8.807489in}}{\pgfqpoint{11.510607in}{8.806520in}}%
\pgfpathquadraticcurveto{\pgfqpoint{11.516421in}{8.805572in}}{\pgfqpoint{11.521307in}{8.803634in}}%
\pgfpathlineto{\pgfqpoint{11.521307in}{8.803634in}}%
\pgfpathclose%
\pgfpathmoveto{\pgfqpoint{11.542191in}{8.833713in}}%
\pgfpathlineto{\pgfqpoint{11.552499in}{8.833713in}}%
\pgfpathquadraticcurveto{\pgfqpoint{11.562148in}{8.818519in}}{\pgfqpoint{11.566951in}{8.803943in}}%
\pgfpathquadraticcurveto{\pgfqpoint{11.571755in}{8.789388in}}{\pgfqpoint{11.571755in}{8.775018in}}%
\pgfpathquadraticcurveto{\pgfqpoint{11.571755in}{8.760587in}}{\pgfqpoint{11.566951in}{8.745970in}}%
\pgfpathquadraticcurveto{\pgfqpoint{11.562148in}{8.731353in}}{\pgfqpoint{11.552499in}{8.716200in}}%
\pgfpathlineto{\pgfqpoint{11.542191in}{8.716200in}}%
\pgfpathquadraticcurveto{\pgfqpoint{11.550747in}{8.730961in}}{\pgfqpoint{11.554973in}{8.745557in}}%
\pgfpathquadraticcurveto{\pgfqpoint{11.559199in}{8.760154in}}{\pgfqpoint{11.559199in}{8.775018in}}%
\pgfpathquadraticcurveto{\pgfqpoint{11.559199in}{8.789903in}}{\pgfqpoint{11.554973in}{8.804396in}}%
\pgfpathquadraticcurveto{\pgfqpoint{11.550747in}{8.818910in}}{\pgfqpoint{11.542191in}{8.833713in}}%
\pgfpathlineto{\pgfqpoint{11.542191in}{8.833713in}}%
\pgfpathclose%
\pgfusepath{fill}%
\end{pgfscope}%
\begin{pgfscope}%
\pgfsetbuttcap%
\pgfsetroundjoin%
\definecolor{currentfill}{rgb}{0.090196,0.168627,0.227451}%
\pgfsetfillcolor{currentfill}%
\pgfsetlinewidth{0.000000pt}%
\definecolor{currentstroke}{rgb}{0.090196,0.168627,0.227451}%
\pgfsetstrokecolor{currentstroke}%
\pgfsetdash{}{0pt}%
\pgfpathmoveto{\pgfqpoint{7.978709in}{9.067751in}}%
\pgfpathquadraticcurveto{\pgfqpoint{7.987368in}{9.067751in}}{\pgfqpoint{7.991804in}{9.071549in}}%
\pgfpathquadraticcurveto{\pgfqpoint{7.996270in}{9.075346in}}{\pgfqpoint{7.996270in}{9.082760in}}%
\pgfpathquadraticcurveto{\pgfqpoint{7.996270in}{9.090082in}}{\pgfqpoint{7.991804in}{9.093910in}}%
\pgfpathquadraticcurveto{\pgfqpoint{7.987368in}{9.097768in}}{\pgfqpoint{7.978709in}{9.097768in}}%
\pgfpathlineto{\pgfqpoint{7.958505in}{9.097768in}}%
\pgfpathlineto{\pgfqpoint{7.958505in}{9.067751in}}%
\pgfpathlineto{\pgfqpoint{7.978709in}{9.067751in}}%
\pgfpathlineto{\pgfqpoint{7.978709in}{9.067751in}}%
\pgfpathclose%
\pgfpathmoveto{\pgfqpoint{7.979955in}{9.005772in}}%
\pgfpathquadraticcurveto{\pgfqpoint{7.990953in}{9.005772in}}{\pgfqpoint{7.996513in}{9.010420in}}%
\pgfpathquadraticcurveto{\pgfqpoint{8.002073in}{9.015069in}}{\pgfqpoint{8.002073in}{9.024457in}}%
\pgfpathquadraticcurveto{\pgfqpoint{8.002073in}{9.033693in}}{\pgfqpoint{7.996574in}{9.038280in}}%
\pgfpathquadraticcurveto{\pgfqpoint{7.991074in}{9.042898in}}{\pgfqpoint{7.979955in}{9.042898in}}%
\pgfpathlineto{\pgfqpoint{7.958505in}{9.042898in}}%
\pgfpathlineto{\pgfqpoint{7.958505in}{9.005772in}}%
\pgfpathlineto{\pgfqpoint{7.979955in}{9.005772in}}%
\pgfpathlineto{\pgfqpoint{7.979955in}{9.005772in}}%
\pgfpathclose%
\pgfpathmoveto{\pgfqpoint{8.013952in}{9.056753in}}%
\pgfpathquadraticcurveto{\pgfqpoint{8.025710in}{9.053319in}}{\pgfqpoint{8.032151in}{9.044114in}}%
\pgfpathquadraticcurveto{\pgfqpoint{8.038622in}{9.034908in}}{\pgfqpoint{8.038622in}{9.021540in}}%
\pgfpathquadraticcurveto{\pgfqpoint{8.038622in}{9.001032in}}{\pgfqpoint{8.024768in}{8.990945in}}%
\pgfpathquadraticcurveto{\pgfqpoint{8.010914in}{8.980889in}}{\pgfqpoint{7.982598in}{8.980889in}}%
\pgfpathlineto{\pgfqpoint{7.921955in}{8.980889in}}%
\pgfpathlineto{\pgfqpoint{7.921955in}{9.122651in}}%
\pgfpathlineto{\pgfqpoint{7.976825in}{9.122651in}}%
\pgfpathquadraticcurveto{\pgfqpoint{8.006357in}{9.122651in}}{\pgfqpoint{8.019603in}{9.113719in}}%
\pgfpathquadraticcurveto{\pgfqpoint{8.032850in}{9.104786in}}{\pgfqpoint{8.032850in}{9.085129in}}%
\pgfpathquadraticcurveto{\pgfqpoint{8.032850in}{9.074799in}}{\pgfqpoint{8.027988in}{9.067538in}}%
\pgfpathquadraticcurveto{\pgfqpoint{8.023158in}{9.060277in}}{\pgfqpoint{8.013952in}{9.056753in}}%
\pgfpathlineto{\pgfqpoint{8.013952in}{9.056753in}}%
\pgfpathclose%
\pgfpathmoveto{\pgfqpoint{8.070159in}{9.122651in}}%
\pgfpathlineto{\pgfqpoint{8.106708in}{9.122651in}}%
\pgfpathlineto{\pgfqpoint{8.106708in}{8.994652in}}%
\pgfpathquadraticcurveto{\pgfqpoint{8.106708in}{8.968159in}}{\pgfqpoint{8.092307in}{8.955064in}}%
\pgfpathquadraticcurveto{\pgfqpoint{8.077936in}{8.941970in}}{\pgfqpoint{8.048770in}{8.941970in}}%
\pgfpathlineto{\pgfqpoint{8.041387in}{8.941970in}}%
\pgfpathlineto{\pgfqpoint{8.041387in}{8.969587in}}%
\pgfpathlineto{\pgfqpoint{8.047068in}{8.969587in}}%
\pgfpathquadraticcurveto{\pgfqpoint{8.058462in}{8.969587in}}{\pgfqpoint{8.064295in}{8.975967in}}%
\pgfpathquadraticcurveto{\pgfqpoint{8.070159in}{8.982317in}}{\pgfqpoint{8.070159in}{8.994652in}}%
\pgfpathlineto{\pgfqpoint{8.070159in}{9.122651in}}%
\pgfpathlineto{\pgfqpoint{8.070159in}{9.122651in}}%
\pgfpathclose%
\pgfpathmoveto{\pgfqpoint{8.241118in}{9.118185in}}%
\pgfpathlineto{\pgfqpoint{8.241118in}{9.088168in}}%
\pgfpathquadraticcurveto{\pgfqpoint{8.229451in}{9.093393in}}{\pgfqpoint{8.218331in}{9.096036in}}%
\pgfpathquadraticcurveto{\pgfqpoint{8.207242in}{9.098710in}}{\pgfqpoint{8.197368in}{9.098710in}}%
\pgfpathquadraticcurveto{\pgfqpoint{8.184273in}{9.098710in}}{\pgfqpoint{8.177984in}{9.095095in}}%
\pgfpathquadraticcurveto{\pgfqpoint{8.171725in}{9.091510in}}{\pgfqpoint{8.171725in}{9.083914in}}%
\pgfpathquadraticcurveto{\pgfqpoint{8.171725in}{9.078202in}}{\pgfqpoint{8.175949in}{9.075012in}}%
\pgfpathquadraticcurveto{\pgfqpoint{8.180172in}{9.071852in}}{\pgfqpoint{8.191291in}{9.069574in}}%
\pgfpathlineto{\pgfqpoint{8.206847in}{9.066444in}}%
\pgfpathquadraticcurveto{\pgfqpoint{8.230484in}{9.061674in}}{\pgfqpoint{8.240449in}{9.051983in}}%
\pgfpathquadraticcurveto{\pgfqpoint{8.250445in}{9.042321in}}{\pgfqpoint{8.250445in}{9.024457in}}%
\pgfpathquadraticcurveto{\pgfqpoint{8.250445in}{9.001032in}}{\pgfqpoint{8.236530in}{8.989578in}}%
\pgfpathquadraticcurveto{\pgfqpoint{8.222615in}{8.978124in}}{\pgfqpoint{8.194026in}{8.978124in}}%
\pgfpathquadraticcurveto{\pgfqpoint{8.180567in}{8.978124in}}{\pgfqpoint{8.166986in}{8.980707in}}%
\pgfpathquadraticcurveto{\pgfqpoint{8.153405in}{8.983259in}}{\pgfqpoint{8.139824in}{8.988302in}}%
\pgfpathlineto{\pgfqpoint{8.139824in}{9.019140in}}%
\pgfpathquadraticcurveto{\pgfqpoint{8.153405in}{9.011939in}}{\pgfqpoint{8.166074in}{9.008263in}}%
\pgfpathquadraticcurveto{\pgfqpoint{8.178744in}{9.004617in}}{\pgfqpoint{8.190532in}{9.004617in}}%
\pgfpathquadraticcurveto{\pgfqpoint{8.202502in}{9.004617in}}{\pgfqpoint{8.208852in}{9.008597in}}%
\pgfpathquadraticcurveto{\pgfqpoint{8.215202in}{9.012608in}}{\pgfqpoint{8.215202in}{9.020021in}}%
\pgfpathquadraticcurveto{\pgfqpoint{8.215202in}{9.026644in}}{\pgfqpoint{8.210888in}{9.030260in}}%
\pgfpathquadraticcurveto{\pgfqpoint{8.206574in}{9.033875in}}{\pgfqpoint{8.193661in}{9.036731in}}%
\pgfpathlineto{\pgfqpoint{8.179503in}{9.039860in}}%
\pgfpathquadraticcurveto{\pgfqpoint{8.158236in}{9.044418in}}{\pgfqpoint{8.148392in}{9.054383in}}%
\pgfpathquadraticcurveto{\pgfqpoint{8.138579in}{9.064348in}}{\pgfqpoint{8.138579in}{9.081240in}}%
\pgfpathquadraticcurveto{\pgfqpoint{8.138579in}{9.102417in}}{\pgfqpoint{8.152251in}{9.113810in}}%
\pgfpathquadraticcurveto{\pgfqpoint{8.165923in}{9.125203in}}{\pgfqpoint{8.191565in}{9.125203in}}%
\pgfpathquadraticcurveto{\pgfqpoint{8.203262in}{9.125203in}}{\pgfqpoint{8.215597in}{9.123441in}}%
\pgfpathquadraticcurveto{\pgfqpoint{8.227932in}{9.121679in}}{\pgfqpoint{8.241118in}{9.118185in}}%
\pgfpathlineto{\pgfqpoint{8.241118in}{9.118185in}}%
\pgfpathclose%
\pgfpathmoveto{\pgfqpoint{8.335302in}{9.087226in}}%
\pgfpathlineto{\pgfqpoint{8.369299in}{9.087226in}}%
\pgfpathlineto{\pgfqpoint{8.395792in}{9.013732in}}%
\pgfpathlineto{\pgfqpoint{8.422164in}{9.087226in}}%
\pgfpathlineto{\pgfqpoint{8.456252in}{9.087226in}}%
\pgfpathlineto{\pgfqpoint{8.414386in}{8.980889in}}%
\pgfpathlineto{\pgfqpoint{8.377077in}{8.980889in}}%
\pgfpathlineto{\pgfqpoint{8.335302in}{9.087226in}}%
\pgfpathlineto{\pgfqpoint{8.335302in}{9.087226in}}%
\pgfpathclose%
\pgfpathmoveto{\pgfqpoint{8.558518in}{9.083914in}}%
\pgfpathlineto{\pgfqpoint{8.558518in}{9.058089in}}%
\pgfpathquadraticcurveto{\pgfqpoint{8.547611in}{9.062647in}}{\pgfqpoint{8.537433in}{9.064925in}}%
\pgfpathquadraticcurveto{\pgfqpoint{8.527285in}{9.067204in}}{\pgfqpoint{8.518262in}{9.067204in}}%
\pgfpathquadraticcurveto{\pgfqpoint{8.508570in}{9.067204in}}{\pgfqpoint{8.503861in}{9.064773in}}%
\pgfpathquadraticcurveto{\pgfqpoint{8.499182in}{9.062343in}}{\pgfqpoint{8.499182in}{9.057330in}}%
\pgfpathquadraticcurveto{\pgfqpoint{8.499182in}{9.053228in}}{\pgfqpoint{8.502737in}{9.051041in}}%
\pgfpathquadraticcurveto{\pgfqpoint{8.506291in}{9.048884in}}{\pgfqpoint{8.515497in}{9.047820in}}%
\pgfpathlineto{\pgfqpoint{8.521482in}{9.046970in}}%
\pgfpathquadraticcurveto{\pgfqpoint{8.547611in}{9.043658in}}{\pgfqpoint{8.556604in}{9.036062in}}%
\pgfpathquadraticcurveto{\pgfqpoint{8.565627in}{9.028467in}}{\pgfqpoint{8.565627in}{9.012213in}}%
\pgfpathquadraticcurveto{\pgfqpoint{8.565627in}{8.995229in}}{\pgfqpoint{8.553080in}{8.986661in}}%
\pgfpathquadraticcurveto{\pgfqpoint{8.540562in}{8.978124in}}{\pgfqpoint{8.515710in}{8.978124in}}%
\pgfpathquadraticcurveto{\pgfqpoint{8.505167in}{8.978124in}}{\pgfqpoint{8.493896in}{8.979795in}}%
\pgfpathquadraticcurveto{\pgfqpoint{8.482654in}{8.981466in}}{\pgfqpoint{8.470775in}{8.984778in}}%
\pgfpathlineto{\pgfqpoint{8.470775in}{9.010602in}}%
\pgfpathquadraticcurveto{\pgfqpoint{8.480953in}{9.005681in}}{\pgfqpoint{8.491617in}{9.003189in}}%
\pgfpathquadraticcurveto{\pgfqpoint{8.502311in}{9.000728in}}{\pgfqpoint{8.513310in}{9.000728in}}%
\pgfpathquadraticcurveto{\pgfqpoint{8.523305in}{9.000728in}}{\pgfqpoint{8.528318in}{9.003463in}}%
\pgfpathquadraticcurveto{\pgfqpoint{8.533362in}{9.006227in}}{\pgfqpoint{8.533362in}{9.011666in}}%
\pgfpathquadraticcurveto{\pgfqpoint{8.533362in}{9.016223in}}{\pgfqpoint{8.529898in}{9.018441in}}%
\pgfpathquadraticcurveto{\pgfqpoint{8.526435in}{9.020659in}}{\pgfqpoint{8.516074in}{9.021905in}}%
\pgfpathlineto{\pgfqpoint{8.510089in}{9.022664in}}%
\pgfpathquadraticcurveto{\pgfqpoint{8.487394in}{9.025520in}}{\pgfqpoint{8.478279in}{9.033207in}}%
\pgfpathquadraticcurveto{\pgfqpoint{8.469165in}{9.040893in}}{\pgfqpoint{8.469165in}{9.056570in}}%
\pgfpathquadraticcurveto{\pgfqpoint{8.469165in}{9.073463in}}{\pgfqpoint{8.480740in}{9.081605in}}%
\pgfpathquadraticcurveto{\pgfqpoint{8.492346in}{9.089778in}}{\pgfqpoint{8.516257in}{9.089778in}}%
\pgfpathquadraticcurveto{\pgfqpoint{8.525675in}{9.089778in}}{\pgfqpoint{8.536005in}{9.088350in}}%
\pgfpathquadraticcurveto{\pgfqpoint{8.546365in}{9.086952in}}{\pgfqpoint{8.558518in}{9.083914in}}%
\pgfpathlineto{\pgfqpoint{8.558518in}{9.083914in}}%
\pgfpathclose%
\pgfpathmoveto{\pgfqpoint{8.643466in}{9.122651in}}%
\pgfpathlineto{\pgfqpoint{8.774108in}{9.122651in}}%
\pgfpathlineto{\pgfqpoint{8.774108in}{9.095003in}}%
\pgfpathlineto{\pgfqpoint{8.727107in}{9.095003in}}%
\pgfpathlineto{\pgfqpoint{8.727107in}{8.980889in}}%
\pgfpathlineto{\pgfqpoint{8.690558in}{8.980889in}}%
\pgfpathlineto{\pgfqpoint{8.690558in}{9.095003in}}%
\pgfpathlineto{\pgfqpoint{8.643466in}{9.095003in}}%
\pgfpathlineto{\pgfqpoint{8.643466in}{9.122651in}}%
\pgfpathlineto{\pgfqpoint{8.643466in}{9.122651in}}%
\pgfpathclose%
\pgfpathmoveto{\pgfqpoint{8.780975in}{9.122651in}}%
\pgfpathlineto{\pgfqpoint{8.816005in}{9.122651in}}%
\pgfpathlineto{\pgfqpoint{8.840493in}{9.019626in}}%
\pgfpathlineto{\pgfqpoint{8.864798in}{9.122651in}}%
\pgfpathlineto{\pgfqpoint{8.900011in}{9.122651in}}%
\pgfpathlineto{\pgfqpoint{8.924317in}{9.019626in}}%
\pgfpathlineto{\pgfqpoint{8.948835in}{9.122651in}}%
\pgfpathlineto{\pgfqpoint{8.983561in}{9.122651in}}%
\pgfpathlineto{\pgfqpoint{8.950141in}{8.980889in}}%
\pgfpathlineto{\pgfqpoint{8.908001in}{8.980889in}}%
\pgfpathlineto{\pgfqpoint{8.882268in}{9.088654in}}%
\pgfpathlineto{\pgfqpoint{8.856838in}{8.980889in}}%
\pgfpathlineto{\pgfqpoint{8.814668in}{8.980889in}}%
\pgfpathlineto{\pgfqpoint{8.780975in}{9.122651in}}%
\pgfpathlineto{\pgfqpoint{8.780975in}{9.122651in}}%
\pgfpathclose%
\pgfpathmoveto{\pgfqpoint{9.007502in}{9.122651in}}%
\pgfpathlineto{\pgfqpoint{9.106122in}{9.122651in}}%
\pgfpathlineto{\pgfqpoint{9.106122in}{9.095003in}}%
\pgfpathlineto{\pgfqpoint{9.044052in}{9.095003in}}%
\pgfpathlineto{\pgfqpoint{9.044052in}{9.068632in}}%
\pgfpathlineto{\pgfqpoint{9.102446in}{9.068632in}}%
\pgfpathlineto{\pgfqpoint{9.102446in}{9.040984in}}%
\pgfpathlineto{\pgfqpoint{9.044052in}{9.040984in}}%
\pgfpathlineto{\pgfqpoint{9.044052in}{8.980889in}}%
\pgfpathlineto{\pgfqpoint{9.007502in}{8.980889in}}%
\pgfpathlineto{\pgfqpoint{9.007502in}{9.122651in}}%
\pgfpathlineto{\pgfqpoint{9.007502in}{9.122651in}}%
\pgfpathclose%
\pgfpathmoveto{\pgfqpoint{9.140332in}{9.122651in}}%
\pgfpathlineto{\pgfqpoint{9.238952in}{9.122651in}}%
\pgfpathlineto{\pgfqpoint{9.238952in}{9.095003in}}%
\pgfpathlineto{\pgfqpoint{9.176882in}{9.095003in}}%
\pgfpathlineto{\pgfqpoint{9.176882in}{9.068632in}}%
\pgfpathlineto{\pgfqpoint{9.235276in}{9.068632in}}%
\pgfpathlineto{\pgfqpoint{9.235276in}{9.040984in}}%
\pgfpathlineto{\pgfqpoint{9.176882in}{9.040984in}}%
\pgfpathlineto{\pgfqpoint{9.176882in}{9.008506in}}%
\pgfpathlineto{\pgfqpoint{9.241048in}{9.008506in}}%
\pgfpathlineto{\pgfqpoint{9.241048in}{8.980889in}}%
\pgfpathlineto{\pgfqpoint{9.140332in}{8.980889in}}%
\pgfpathlineto{\pgfqpoint{9.140332in}{9.122651in}}%
\pgfpathlineto{\pgfqpoint{9.140332in}{9.122651in}}%
\pgfpathclose%
\pgfusepath{fill}%
\end{pgfscope}%
\begin{pgfscope}%
\pgfpathrectangle{\pgfqpoint{7.904091in}{2.928000in}}{\pgfqpoint{5.664000in}{5.664000in}}%
\pgfusepath{clip}%
\pgfsetbuttcap%
\pgfsetroundjoin%
\definecolor{currentfill}{rgb}{0.541176,0.380392,0.658824}%
\pgfsetfillcolor{currentfill}%
\pgfsetfillopacity{0.900000}%
\pgfsetlinewidth{0.803000pt}%
\definecolor{currentstroke}{rgb}{0.541176,0.380392,0.658824}%
\pgfsetstrokecolor{currentstroke}%
\pgfsetstrokeopacity{0.900000}%
\pgfsetdash{}{0pt}%
\pgfsys@defobject{currentmarker}{\pgfqpoint{-0.038665in}{-0.038665in}}{\pgfqpoint{0.038665in}{0.038665in}}{%
\pgfpathmoveto{\pgfqpoint{0.000000in}{-0.038665in}}%
\pgfpathcurveto{\pgfqpoint{0.010254in}{-0.038665in}}{\pgfqpoint{0.020090in}{-0.034591in}}{\pgfqpoint{0.027340in}{-0.027340in}}%
\pgfpathcurveto{\pgfqpoint{0.034591in}{-0.020090in}}{\pgfqpoint{0.038665in}{-0.010254in}}{\pgfqpoint{0.038665in}{0.000000in}}%
\pgfpathcurveto{\pgfqpoint{0.038665in}{0.010254in}}{\pgfqpoint{0.034591in}{0.020090in}}{\pgfqpoint{0.027340in}{0.027340in}}%
\pgfpathcurveto{\pgfqpoint{0.020090in}{0.034591in}}{\pgfqpoint{0.010254in}{0.038665in}}{\pgfqpoint{0.000000in}{0.038665in}}%
\pgfpathcurveto{\pgfqpoint{-0.010254in}{0.038665in}}{\pgfqpoint{-0.020090in}{0.034591in}}{\pgfqpoint{-0.027340in}{0.027340in}}%
\pgfpathcurveto{\pgfqpoint{-0.034591in}{0.020090in}}{\pgfqpoint{-0.038665in}{0.010254in}}{\pgfqpoint{-0.038665in}{0.000000in}}%
\pgfpathcurveto{\pgfqpoint{-0.038665in}{-0.010254in}}{\pgfqpoint{-0.034591in}{-0.020090in}}{\pgfqpoint{-0.027340in}{-0.027340in}}%
\pgfpathcurveto{\pgfqpoint{-0.020090in}{-0.034591in}}{\pgfqpoint{-0.010254in}{-0.038665in}}{\pgfqpoint{0.000000in}{-0.038665in}}%
\pgfpathlineto{\pgfqpoint{0.000000in}{-0.038665in}}%
\pgfpathclose%
\pgfusepath{stroke,fill}%
}%
\begin{pgfscope}%
\pgfsys@transformshift{10.874722in}{4.869728in}%
\pgfsys@useobject{currentmarker}{}%
\end{pgfscope}%
\end{pgfscope}%
\begin{pgfscope}%
\pgfpathrectangle{\pgfqpoint{7.904091in}{2.928000in}}{\pgfqpoint{5.664000in}{5.664000in}}%
\pgfusepath{clip}%
\pgfsetbuttcap%
\pgfsetroundjoin%
\definecolor{currentfill}{rgb}{0.321569,0.321569,0.321569}%
\pgfsetfillcolor{currentfill}%
\pgfsetfillopacity{0.900000}%
\pgfsetlinewidth{0.803000pt}%
\definecolor{currentstroke}{rgb}{0.321569,0.321569,0.321569}%
\pgfsetstrokecolor{currentstroke}%
\pgfsetstrokeopacity{0.900000}%
\pgfsetdash{}{0pt}%
\pgfsys@defobject{currentmarker}{\pgfqpoint{-0.038665in}{-0.038665in}}{\pgfqpoint{0.038665in}{0.038665in}}{%
\pgfpathmoveto{\pgfqpoint{0.000000in}{-0.038665in}}%
\pgfpathcurveto{\pgfqpoint{0.010254in}{-0.038665in}}{\pgfqpoint{0.020090in}{-0.034591in}}{\pgfqpoint{0.027340in}{-0.027340in}}%
\pgfpathcurveto{\pgfqpoint{0.034591in}{-0.020090in}}{\pgfqpoint{0.038665in}{-0.010254in}}{\pgfqpoint{0.038665in}{0.000000in}}%
\pgfpathcurveto{\pgfqpoint{0.038665in}{0.010254in}}{\pgfqpoint{0.034591in}{0.020090in}}{\pgfqpoint{0.027340in}{0.027340in}}%
\pgfpathcurveto{\pgfqpoint{0.020090in}{0.034591in}}{\pgfqpoint{0.010254in}{0.038665in}}{\pgfqpoint{0.000000in}{0.038665in}}%
\pgfpathcurveto{\pgfqpoint{-0.010254in}{0.038665in}}{\pgfqpoint{-0.020090in}{0.034591in}}{\pgfqpoint{-0.027340in}{0.027340in}}%
\pgfpathcurveto{\pgfqpoint{-0.034591in}{0.020090in}}{\pgfqpoint{-0.038665in}{0.010254in}}{\pgfqpoint{-0.038665in}{0.000000in}}%
\pgfpathcurveto{\pgfqpoint{-0.038665in}{-0.010254in}}{\pgfqpoint{-0.034591in}{-0.020090in}}{\pgfqpoint{-0.027340in}{-0.027340in}}%
\pgfpathcurveto{\pgfqpoint{-0.020090in}{-0.034591in}}{\pgfqpoint{-0.010254in}{-0.038665in}}{\pgfqpoint{0.000000in}{-0.038665in}}%
\pgfpathlineto{\pgfqpoint{0.000000in}{-0.038665in}}%
\pgfpathclose%
\pgfusepath{stroke,fill}%
}%
\begin{pgfscope}%
\pgfsys@transformshift{9.475413in}{4.109736in}%
\pgfsys@useobject{currentmarker}{}%
\end{pgfscope}%
\begin{pgfscope}%
\pgfsys@transformshift{11.779332in}{7.180296in}%
\pgfsys@useobject{currentmarker}{}%
\end{pgfscope}%
\begin{pgfscope}%
\pgfsys@transformshift{9.528492in}{4.171503in}%
\pgfsys@useobject{currentmarker}{}%
\end{pgfscope}%
\begin{pgfscope}%
\pgfsys@transformshift{9.616740in}{4.236289in}%
\pgfsys@useobject{currentmarker}{}%
\end{pgfscope}%
\begin{pgfscope}%
\pgfsys@transformshift{9.189814in}{3.820815in}%
\pgfsys@useobject{currentmarker}{}%
\end{pgfscope}%
\begin{pgfscope}%
\pgfsys@transformshift{9.915081in}{4.627484in}%
\pgfsys@useobject{currentmarker}{}%
\end{pgfscope}%
\begin{pgfscope}%
\pgfsys@transformshift{11.459832in}{6.906902in}%
\pgfsys@useobject{currentmarker}{}%
\end{pgfscope}%
\begin{pgfscope}%
\pgfsys@transformshift{11.707179in}{7.324358in}%
\pgfsys@useobject{currentmarker}{}%
\end{pgfscope}%
\begin{pgfscope}%
\pgfsys@transformshift{11.303736in}{6.719453in}%
\pgfsys@useobject{currentmarker}{}%
\end{pgfscope}%
\end{pgfscope}%
\begin{pgfscope}%
\pgfpathrectangle{\pgfqpoint{7.904091in}{2.928000in}}{\pgfqpoint{5.664000in}{5.664000in}}%
\pgfusepath{clip}%
\pgfsetbuttcap%
\pgfsetroundjoin%
\definecolor{currentfill}{rgb}{0.200000,0.133333,0.533333}%
\pgfsetfillcolor{currentfill}%
\pgfsetfillopacity{0.900000}%
\pgfsetlinewidth{0.803000pt}%
\definecolor{currentstroke}{rgb}{0.200000,0.133333,0.533333}%
\pgfsetstrokecolor{currentstroke}%
\pgfsetstrokeopacity{0.900000}%
\pgfsetdash{}{0pt}%
\pgfsys@defobject{currentmarker}{\pgfqpoint{-0.038665in}{-0.038665in}}{\pgfqpoint{0.038665in}{0.038665in}}{%
\pgfpathmoveto{\pgfqpoint{0.000000in}{-0.038665in}}%
\pgfpathcurveto{\pgfqpoint{0.010254in}{-0.038665in}}{\pgfqpoint{0.020090in}{-0.034591in}}{\pgfqpoint{0.027340in}{-0.027340in}}%
\pgfpathcurveto{\pgfqpoint{0.034591in}{-0.020090in}}{\pgfqpoint{0.038665in}{-0.010254in}}{\pgfqpoint{0.038665in}{0.000000in}}%
\pgfpathcurveto{\pgfqpoint{0.038665in}{0.010254in}}{\pgfqpoint{0.034591in}{0.020090in}}{\pgfqpoint{0.027340in}{0.027340in}}%
\pgfpathcurveto{\pgfqpoint{0.020090in}{0.034591in}}{\pgfqpoint{0.010254in}{0.038665in}}{\pgfqpoint{0.000000in}{0.038665in}}%
\pgfpathcurveto{\pgfqpoint{-0.010254in}{0.038665in}}{\pgfqpoint{-0.020090in}{0.034591in}}{\pgfqpoint{-0.027340in}{0.027340in}}%
\pgfpathcurveto{\pgfqpoint{-0.034591in}{0.020090in}}{\pgfqpoint{-0.038665in}{0.010254in}}{\pgfqpoint{-0.038665in}{0.000000in}}%
\pgfpathcurveto{\pgfqpoint{-0.038665in}{-0.010254in}}{\pgfqpoint{-0.034591in}{-0.020090in}}{\pgfqpoint{-0.027340in}{-0.027340in}}%
\pgfpathcurveto{\pgfqpoint{-0.020090in}{-0.034591in}}{\pgfqpoint{-0.010254in}{-0.038665in}}{\pgfqpoint{0.000000in}{-0.038665in}}%
\pgfpathlineto{\pgfqpoint{0.000000in}{-0.038665in}}%
\pgfpathclose%
\pgfusepath{stroke,fill}%
}%
\begin{pgfscope}%
\pgfsys@transformshift{9.138135in}{3.493034in}%
\pgfsys@useobject{currentmarker}{}%
\end{pgfscope}%
\begin{pgfscope}%
\pgfsys@transformshift{9.746763in}{4.013126in}%
\pgfsys@useobject{currentmarker}{}%
\end{pgfscope}%
\begin{pgfscope}%
\pgfsys@transformshift{8.949745in}{3.206301in}%
\pgfsys@useobject{currentmarker}{}%
\end{pgfscope}%
\begin{pgfscope}%
\pgfsys@transformshift{9.633108in}{3.873435in}%
\pgfsys@useobject{currentmarker}{}%
\end{pgfscope}%
\end{pgfscope}%
\begin{pgfscope}%
\pgfsetbuttcap%
\pgfsetroundjoin%
\definecolor{currentfill}{rgb}{0.176471,0.713725,0.643137}%
\pgfsetfillcolor{currentfill}%
\pgfsetfillopacity{0.900000}%
\pgfsetlinewidth{0.803000pt}%
\definecolor{currentstroke}{rgb}{0.176471,0.713725,0.643137}%
\pgfsetstrokecolor{currentstroke}%
\pgfsetstrokeopacity{0.900000}%
\pgfsetdash{}{0pt}%
\pgfsys@defobject{currentmarker}{\pgfqpoint{-0.059738in}{-0.059738in}}{\pgfqpoint{0.059738in}{0.059738in}}{%
\pgfpathmoveto{\pgfqpoint{0.000000in}{-0.059738in}}%
\pgfpathlineto{\pgfqpoint{0.059738in}{0.000000in}}%
\pgfpathlineto{\pgfqpoint{0.000000in}{0.059738in}}%
\pgfpathlineto{\pgfqpoint{-0.059738in}{0.000000in}}%
\pgfpathlineto{\pgfqpoint{0.000000in}{-0.059738in}}%
\pgfpathclose%
\pgfusepath{stroke,fill}%
}%
\begin{pgfscope}%
\pgfsys@transformshift{11.515783in}{8.592000in}%
\pgfsys@useobject{currentmarker}{}%
\end{pgfscope}%
\end{pgfscope}%
\begin{pgfscope}%
\pgfpathrectangle{\pgfqpoint{7.904091in}{2.928000in}}{\pgfqpoint{5.664000in}{5.664000in}}%
\pgfusepath{clip}%
\pgfsetbuttcap%
\pgfsetroundjoin%
\definecolor{currentfill}{rgb}{0.000000,0.447059,0.698039}%
\pgfsetfillcolor{currentfill}%
\pgfsetfillopacity{0.900000}%
\pgfsetlinewidth{0.803000pt}%
\definecolor{currentstroke}{rgb}{0.000000,0.447059,0.698039}%
\pgfsetstrokecolor{currentstroke}%
\pgfsetstrokeopacity{0.900000}%
\pgfsetdash{}{0pt}%
\pgfsys@defobject{currentmarker}{\pgfqpoint{-0.038665in}{-0.038665in}}{\pgfqpoint{0.038665in}{0.038665in}}{%
\pgfpathmoveto{\pgfqpoint{0.000000in}{-0.038665in}}%
\pgfpathcurveto{\pgfqpoint{0.010254in}{-0.038665in}}{\pgfqpoint{0.020090in}{-0.034591in}}{\pgfqpoint{0.027340in}{-0.027340in}}%
\pgfpathcurveto{\pgfqpoint{0.034591in}{-0.020090in}}{\pgfqpoint{0.038665in}{-0.010254in}}{\pgfqpoint{0.038665in}{0.000000in}}%
\pgfpathcurveto{\pgfqpoint{0.038665in}{0.010254in}}{\pgfqpoint{0.034591in}{0.020090in}}{\pgfqpoint{0.027340in}{0.027340in}}%
\pgfpathcurveto{\pgfqpoint{0.020090in}{0.034591in}}{\pgfqpoint{0.010254in}{0.038665in}}{\pgfqpoint{0.000000in}{0.038665in}}%
\pgfpathcurveto{\pgfqpoint{-0.010254in}{0.038665in}}{\pgfqpoint{-0.020090in}{0.034591in}}{\pgfqpoint{-0.027340in}{0.027340in}}%
\pgfpathcurveto{\pgfqpoint{-0.034591in}{0.020090in}}{\pgfqpoint{-0.038665in}{0.010254in}}{\pgfqpoint{-0.038665in}{0.000000in}}%
\pgfpathcurveto{\pgfqpoint{-0.038665in}{-0.010254in}}{\pgfqpoint{-0.034591in}{-0.020090in}}{\pgfqpoint{-0.027340in}{-0.027340in}}%
\pgfpathcurveto{\pgfqpoint{-0.020090in}{-0.034591in}}{\pgfqpoint{-0.010254in}{-0.038665in}}{\pgfqpoint{0.000000in}{-0.038665in}}%
\pgfpathlineto{\pgfqpoint{0.000000in}{-0.038665in}}%
\pgfpathclose%
\pgfusepath{stroke,fill}%
}%
\begin{pgfscope}%
\pgfsys@transformshift{10.755143in}{8.067246in}%
\pgfsys@useobject{currentmarker}{}%
\end{pgfscope}%
\end{pgfscope}%
\begin{pgfscope}%
\pgfsetbuttcap%
\pgfsetroundjoin%
\definecolor{currentfill}{rgb}{0.000000,0.447059,0.698039}%
\pgfsetfillcolor{currentfill}%
\pgfsetfillopacity{0.900000}%
\pgfsetlinewidth{0.803000pt}%
\definecolor{currentstroke}{rgb}{0.000000,0.447059,0.698039}%
\pgfsetstrokecolor{currentstroke}%
\pgfsetstrokeopacity{0.900000}%
\pgfsetdash{}{0pt}%
\pgfsys@defobject{currentmarker}{\pgfqpoint{-0.059738in}{-0.059738in}}{\pgfqpoint{0.059738in}{0.059738in}}{%
\pgfpathmoveto{\pgfqpoint{0.000000in}{-0.059738in}}%
\pgfpathlineto{\pgfqpoint{0.059738in}{0.000000in}}%
\pgfpathlineto{\pgfqpoint{0.000000in}{0.059738in}}%
\pgfpathlineto{\pgfqpoint{-0.059738in}{0.000000in}}%
\pgfpathlineto{\pgfqpoint{0.000000in}{-0.059738in}}%
\pgfpathclose%
\pgfusepath{stroke,fill}%
}%
\begin{pgfscope}%
\pgfsys@transformshift{9.189749in}{2.928000in}%
\pgfsys@useobject{currentmarker}{}%
\end{pgfscope}%
\begin{pgfscope}%
\pgfsys@transformshift{9.211505in}{2.928000in}%
\pgfsys@useobject{currentmarker}{}%
\end{pgfscope}%
\begin{pgfscope}%
\pgfsys@transformshift{9.669630in}{2.928000in}%
\pgfsys@useobject{currentmarker}{}%
\end{pgfscope}%
\begin{pgfscope}%
\pgfsys@transformshift{9.732921in}{2.928000in}%
\pgfsys@useobject{currentmarker}{}%
\end{pgfscope}%
\end{pgfscope}%
\begin{pgfscope}%
\pgfsetbuttcap%
\pgfsetroundjoin%
\definecolor{currentfill}{rgb}{0.941176,0.584314,0.337255}%
\pgfsetfillcolor{currentfill}%
\pgfsetfillopacity{0.900000}%
\pgfsetlinewidth{0.803000pt}%
\definecolor{currentstroke}{rgb}{0.941176,0.584314,0.337255}%
\pgfsetstrokecolor{currentstroke}%
\pgfsetstrokeopacity{0.900000}%
\pgfsetdash{}{0pt}%
\pgfsys@defobject{currentmarker}{\pgfqpoint{-0.059738in}{-0.059738in}}{\pgfqpoint{0.059738in}{0.059738in}}{%
\pgfpathmoveto{\pgfqpoint{0.000000in}{-0.059738in}}%
\pgfpathlineto{\pgfqpoint{0.059738in}{0.000000in}}%
\pgfpathlineto{\pgfqpoint{0.000000in}{0.059738in}}%
\pgfpathlineto{\pgfqpoint{-0.059738in}{0.000000in}}%
\pgfpathlineto{\pgfqpoint{0.000000in}{-0.059738in}}%
\pgfpathclose%
\pgfusepath{stroke,fill}%
}%
\begin{pgfscope}%
\pgfsys@transformshift{13.568091in}{8.592000in}%
\pgfsys@useobject{currentmarker}{}%
\end{pgfscope}%
\begin{pgfscope}%
\pgfsys@transformshift{13.260919in}{8.592000in}%
\pgfsys@useobject{currentmarker}{}%
\end{pgfscope}%
\begin{pgfscope}%
\pgfsys@transformshift{13.568091in}{8.592000in}%
\pgfsys@useobject{currentmarker}{}%
\end{pgfscope}%
\end{pgfscope}%
\begin{pgfscope}%
\pgfpathrectangle{\pgfqpoint{7.904091in}{2.928000in}}{\pgfqpoint{5.664000in}{5.664000in}}%
\pgfusepath{clip}%
\pgfsetbuttcap%
\pgfsetroundjoin%
\definecolor{currentfill}{rgb}{0.549020,0.160784,0.505882}%
\pgfsetfillcolor{currentfill}%
\pgfsetfillopacity{0.900000}%
\pgfsetlinewidth{0.803000pt}%
\definecolor{currentstroke}{rgb}{0.549020,0.160784,0.505882}%
\pgfsetstrokecolor{currentstroke}%
\pgfsetstrokeopacity{0.900000}%
\pgfsetdash{}{0pt}%
\pgfsys@defobject{currentmarker}{\pgfqpoint{-0.038665in}{-0.038665in}}{\pgfqpoint{0.038665in}{0.038665in}}{%
\pgfpathmoveto{\pgfqpoint{0.000000in}{-0.038665in}}%
\pgfpathcurveto{\pgfqpoint{0.010254in}{-0.038665in}}{\pgfqpoint{0.020090in}{-0.034591in}}{\pgfqpoint{0.027340in}{-0.027340in}}%
\pgfpathcurveto{\pgfqpoint{0.034591in}{-0.020090in}}{\pgfqpoint{0.038665in}{-0.010254in}}{\pgfqpoint{0.038665in}{0.000000in}}%
\pgfpathcurveto{\pgfqpoint{0.038665in}{0.010254in}}{\pgfqpoint{0.034591in}{0.020090in}}{\pgfqpoint{0.027340in}{0.027340in}}%
\pgfpathcurveto{\pgfqpoint{0.020090in}{0.034591in}}{\pgfqpoint{0.010254in}{0.038665in}}{\pgfqpoint{0.000000in}{0.038665in}}%
\pgfpathcurveto{\pgfqpoint{-0.010254in}{0.038665in}}{\pgfqpoint{-0.020090in}{0.034591in}}{\pgfqpoint{-0.027340in}{0.027340in}}%
\pgfpathcurveto{\pgfqpoint{-0.034591in}{0.020090in}}{\pgfqpoint{-0.038665in}{0.010254in}}{\pgfqpoint{-0.038665in}{0.000000in}}%
\pgfpathcurveto{\pgfqpoint{-0.038665in}{-0.010254in}}{\pgfqpoint{-0.034591in}{-0.020090in}}{\pgfqpoint{-0.027340in}{-0.027340in}}%
\pgfpathcurveto{\pgfqpoint{-0.020090in}{-0.034591in}}{\pgfqpoint{-0.010254in}{-0.038665in}}{\pgfqpoint{0.000000in}{-0.038665in}}%
\pgfpathlineto{\pgfqpoint{0.000000in}{-0.038665in}}%
\pgfpathclose%
\pgfusepath{stroke,fill}%
}%
\begin{pgfscope}%
\pgfsys@transformshift{11.010582in}{5.248026in}%
\pgfsys@useobject{currentmarker}{}%
\end{pgfscope}%
\end{pgfscope}%
\begin{pgfscope}%
\pgfsetbuttcap%
\pgfsetroundjoin%
\definecolor{currentfill}{rgb}{0.549020,0.160784,0.505882}%
\pgfsetfillcolor{currentfill}%
\pgfsetfillopacity{0.900000}%
\pgfsetlinewidth{0.803000pt}%
\definecolor{currentstroke}{rgb}{0.549020,0.160784,0.505882}%
\pgfsetstrokecolor{currentstroke}%
\pgfsetstrokeopacity{0.900000}%
\pgfsetdash{}{0pt}%
\pgfsys@defobject{currentmarker}{\pgfqpoint{-0.059738in}{-0.059738in}}{\pgfqpoint{0.059738in}{0.059738in}}{%
\pgfpathmoveto{\pgfqpoint{0.000000in}{-0.059738in}}%
\pgfpathlineto{\pgfqpoint{0.059738in}{0.000000in}}%
\pgfpathlineto{\pgfqpoint{0.000000in}{0.059738in}}%
\pgfpathlineto{\pgfqpoint{-0.059738in}{0.000000in}}%
\pgfpathlineto{\pgfqpoint{0.000000in}{-0.059738in}}%
\pgfpathclose%
\pgfusepath{stroke,fill}%
}%
\begin{pgfscope}%
\pgfsys@transformshift{13.568091in}{4.007000in}%
\pgfsys@useobject{currentmarker}{}%
\end{pgfscope}%
\begin{pgfscope}%
\pgfsys@transformshift{13.568091in}{8.592000in}%
\pgfsys@useobject{currentmarker}{}%
\end{pgfscope}%
\end{pgfscope}%
\begin{pgfscope}%
\pgfpathrectangle{\pgfqpoint{7.904091in}{2.928000in}}{\pgfqpoint{5.664000in}{5.664000in}}%
\pgfusepath{clip}%
\pgfsetbuttcap%
\pgfsetroundjoin%
\definecolor{currentfill}{rgb}{0.000000,0.560784,0.611765}%
\pgfsetfillcolor{currentfill}%
\pgfsetfillopacity{0.900000}%
\pgfsetlinewidth{0.803000pt}%
\definecolor{currentstroke}{rgb}{0.000000,0.560784,0.611765}%
\pgfsetstrokecolor{currentstroke}%
\pgfsetstrokeopacity{0.900000}%
\pgfsetdash{}{0pt}%
\pgfsys@defobject{currentmarker}{\pgfqpoint{-0.038665in}{-0.038665in}}{\pgfqpoint{0.038665in}{0.038665in}}{%
\pgfpathmoveto{\pgfqpoint{0.000000in}{-0.038665in}}%
\pgfpathcurveto{\pgfqpoint{0.010254in}{-0.038665in}}{\pgfqpoint{0.020090in}{-0.034591in}}{\pgfqpoint{0.027340in}{-0.027340in}}%
\pgfpathcurveto{\pgfqpoint{0.034591in}{-0.020090in}}{\pgfqpoint{0.038665in}{-0.010254in}}{\pgfqpoint{0.038665in}{0.000000in}}%
\pgfpathcurveto{\pgfqpoint{0.038665in}{0.010254in}}{\pgfqpoint{0.034591in}{0.020090in}}{\pgfqpoint{0.027340in}{0.027340in}}%
\pgfpathcurveto{\pgfqpoint{0.020090in}{0.034591in}}{\pgfqpoint{0.010254in}{0.038665in}}{\pgfqpoint{0.000000in}{0.038665in}}%
\pgfpathcurveto{\pgfqpoint{-0.010254in}{0.038665in}}{\pgfqpoint{-0.020090in}{0.034591in}}{\pgfqpoint{-0.027340in}{0.027340in}}%
\pgfpathcurveto{\pgfqpoint{-0.034591in}{0.020090in}}{\pgfqpoint{-0.038665in}{0.010254in}}{\pgfqpoint{-0.038665in}{0.000000in}}%
\pgfpathcurveto{\pgfqpoint{-0.038665in}{-0.010254in}}{\pgfqpoint{-0.034591in}{-0.020090in}}{\pgfqpoint{-0.027340in}{-0.027340in}}%
\pgfpathcurveto{\pgfqpoint{-0.020090in}{-0.034591in}}{\pgfqpoint{-0.010254in}{-0.038665in}}{\pgfqpoint{0.000000in}{-0.038665in}}%
\pgfpathlineto{\pgfqpoint{0.000000in}{-0.038665in}}%
\pgfpathclose%
\pgfusepath{stroke,fill}%
}%
\begin{pgfscope}%
\pgfsys@transformshift{8.616431in}{4.215783in}%
\pgfsys@useobject{currentmarker}{}%
\end{pgfscope}%
\end{pgfscope}%
\begin{pgfscope}%
\pgfsetbuttcap%
\pgfsetroundjoin%
\definecolor{currentfill}{rgb}{0.000000,0.560784,0.611765}%
\pgfsetfillcolor{currentfill}%
\pgfsetfillopacity{0.900000}%
\pgfsetlinewidth{0.803000pt}%
\definecolor{currentstroke}{rgb}{0.000000,0.560784,0.611765}%
\pgfsetstrokecolor{currentstroke}%
\pgfsetstrokeopacity{0.900000}%
\pgfsetdash{}{0pt}%
\pgfsys@defobject{currentmarker}{\pgfqpoint{-0.059738in}{-0.059738in}}{\pgfqpoint{0.059738in}{0.059738in}}{%
\pgfpathmoveto{\pgfqpoint{0.000000in}{-0.059738in}}%
\pgfpathlineto{\pgfqpoint{0.059738in}{0.000000in}}%
\pgfpathlineto{\pgfqpoint{0.000000in}{0.059738in}}%
\pgfpathlineto{\pgfqpoint{-0.059738in}{0.000000in}}%
\pgfpathlineto{\pgfqpoint{0.000000in}{-0.059738in}}%
\pgfpathclose%
\pgfusepath{stroke,fill}%
}%
\begin{pgfscope}%
\pgfsys@transformshift{7.904091in}{2.928000in}%
\pgfsys@useobject{currentmarker}{}%
\end{pgfscope}%
\end{pgfscope}%
\begin{pgfscope}%
\pgfsetbuttcap%
\pgfsetroundjoin%
\pgfsetlinewidth{0.903375pt}%
\definecolor{currentstroke}{rgb}{0.321569,0.321569,0.321569}%
\pgfsetstrokecolor{currentstroke}%
\pgfsetdash{}{0pt}%
\pgfpathmoveto{\pgfqpoint{7.904091in}{2.855166in}}%
\pgfpathlineto{\pgfqpoint{7.976925in}{2.928000in}}%
\pgfpathlineto{\pgfqpoint{7.904091in}{3.000834in}}%
\pgfpathlineto{\pgfqpoint{7.831257in}{2.928000in}}%
\pgfpathlineto{\pgfqpoint{7.904091in}{2.855166in}}%
\pgfpathclose%
\pgfusepath{stroke}%
\end{pgfscope}%
\begin{pgfscope}%
\pgfsetbuttcap%
\pgfsetroundjoin%
\pgfsetlinewidth{0.903375pt}%
\definecolor{currentstroke}{rgb}{0.321569,0.321569,0.321569}%
\pgfsetstrokecolor{currentstroke}%
\pgfsetdash{}{0pt}%
\pgfpathmoveto{\pgfqpoint{7.904091in}{2.855166in}}%
\pgfpathlineto{\pgfqpoint{7.976925in}{2.928000in}}%
\pgfpathlineto{\pgfqpoint{7.904091in}{3.000834in}}%
\pgfpathlineto{\pgfqpoint{7.831257in}{2.928000in}}%
\pgfpathlineto{\pgfqpoint{7.904091in}{2.855166in}}%
\pgfpathclose%
\pgfusepath{stroke}%
\end{pgfscope}%
\begin{pgfscope}%
\pgfsetbuttcap%
\pgfsetroundjoin%
\pgfsetlinewidth{0.903375pt}%
\definecolor{currentstroke}{rgb}{0.321569,0.321569,0.321569}%
\pgfsetstrokecolor{currentstroke}%
\pgfsetdash{}{0pt}%
\pgfpathmoveto{\pgfqpoint{7.904091in}{2.855166in}}%
\pgfpathlineto{\pgfqpoint{7.976925in}{2.928000in}}%
\pgfpathlineto{\pgfqpoint{7.904091in}{3.000834in}}%
\pgfpathlineto{\pgfqpoint{7.831257in}{2.928000in}}%
\pgfpathlineto{\pgfqpoint{7.904091in}{2.855166in}}%
\pgfpathclose%
\pgfusepath{stroke}%
\end{pgfscope}%
\begin{pgfscope}%
\pgfsetbuttcap%
\pgfsetroundjoin%
\pgfsetlinewidth{0.903375pt}%
\definecolor{currentstroke}{rgb}{0.321569,0.321569,0.321569}%
\pgfsetstrokecolor{currentstroke}%
\pgfsetdash{}{0pt}%
\pgfpathmoveto{\pgfqpoint{7.904091in}{2.855166in}}%
\pgfpathlineto{\pgfqpoint{7.976925in}{2.928000in}}%
\pgfpathlineto{\pgfqpoint{7.904091in}{3.000834in}}%
\pgfpathlineto{\pgfqpoint{7.831257in}{2.928000in}}%
\pgfpathlineto{\pgfqpoint{7.904091in}{2.855166in}}%
\pgfpathclose%
\pgfusepath{stroke}%
\end{pgfscope}%
\begin{pgfscope}%
\pgfsetbuttcap%
\pgfsetroundjoin%
\pgfsetlinewidth{0.903375pt}%
\definecolor{currentstroke}{rgb}{0.321569,0.321569,0.321569}%
\pgfsetstrokecolor{currentstroke}%
\pgfsetdash{}{0pt}%
\pgfpathmoveto{\pgfqpoint{7.904091in}{2.855166in}}%
\pgfpathlineto{\pgfqpoint{7.976925in}{2.928000in}}%
\pgfpathlineto{\pgfqpoint{7.904091in}{3.000834in}}%
\pgfpathlineto{\pgfqpoint{7.831257in}{2.928000in}}%
\pgfpathlineto{\pgfqpoint{7.904091in}{2.855166in}}%
\pgfpathclose%
\pgfusepath{stroke}%
\end{pgfscope}%
\begin{pgfscope}%
\pgfsetbuttcap%
\pgfsetroundjoin%
\pgfsetlinewidth{0.903375pt}%
\definecolor{currentstroke}{rgb}{0.321569,0.321569,0.321569}%
\pgfsetstrokecolor{currentstroke}%
\pgfsetdash{}{0pt}%
\pgfpathmoveto{\pgfqpoint{7.904091in}{2.855166in}}%
\pgfpathlineto{\pgfqpoint{7.976925in}{2.928000in}}%
\pgfpathlineto{\pgfqpoint{7.904091in}{3.000834in}}%
\pgfpathlineto{\pgfqpoint{7.831257in}{2.928000in}}%
\pgfpathlineto{\pgfqpoint{7.904091in}{2.855166in}}%
\pgfpathclose%
\pgfusepath{stroke}%
\end{pgfscope}%
\begin{pgfscope}%
\pgfpathrectangle{\pgfqpoint{7.904091in}{2.928000in}}{\pgfqpoint{5.664000in}{5.664000in}}%
\pgfusepath{clip}%
\pgfsetbuttcap%
\pgfsetroundjoin%
\pgfsetlinewidth{0.803000pt}%
\definecolor{currentstroke}{rgb}{0.768627,0.603922,0.000000}%
\pgfsetstrokecolor{currentstroke}%
\pgfsetdash{}{0pt}%
\pgfpathmoveto{\pgfqpoint{11.348575in}{6.340870in}}%
\pgfpathlineto{\pgfqpoint{11.441745in}{6.340870in}}%
\pgfpathlineto{\pgfqpoint{11.441745in}{6.434039in}}%
\pgfpathlineto{\pgfqpoint{11.348575in}{6.434039in}}%
\pgfpathlineto{\pgfqpoint{11.348575in}{6.340870in}}%
\pgfpathclose%
\pgfusepath{stroke}%
\end{pgfscope}%
\begin{pgfscope}%
\pgfpathrectangle{\pgfqpoint{7.904091in}{2.928000in}}{\pgfqpoint{5.664000in}{5.664000in}}%
\pgfusepath{clip}%
\pgfsetbuttcap%
\pgfsetroundjoin%
\pgfsetlinewidth{0.803000pt}%
\definecolor{currentstroke}{rgb}{0.768627,0.603922,0.000000}%
\pgfsetstrokecolor{currentstroke}%
\pgfsetdash{}{0pt}%
\pgfpathmoveto{\pgfqpoint{11.612428in}{6.475710in}}%
\pgfpathlineto{\pgfqpoint{11.705597in}{6.475710in}}%
\pgfpathlineto{\pgfqpoint{11.705597in}{6.568880in}}%
\pgfpathlineto{\pgfqpoint{11.612428in}{6.568880in}}%
\pgfpathlineto{\pgfqpoint{11.612428in}{6.475710in}}%
\pgfpathclose%
\pgfusepath{stroke}%
\end{pgfscope}%
\begin{pgfscope}%
\pgfpathrectangle{\pgfqpoint{7.904091in}{2.928000in}}{\pgfqpoint{5.664000in}{5.664000in}}%
\pgfusepath{clip}%
\pgfsetbuttcap%
\pgfsetroundjoin%
\pgfsetlinewidth{0.803000pt}%
\definecolor{currentstroke}{rgb}{0.768627,0.603922,0.000000}%
\pgfsetstrokecolor{currentstroke}%
\pgfsetdash{}{0pt}%
\pgfpathmoveto{\pgfqpoint{10.785557in}{5.744203in}}%
\pgfpathlineto{\pgfqpoint{10.878726in}{5.744203in}}%
\pgfpathlineto{\pgfqpoint{10.878726in}{5.837373in}}%
\pgfpathlineto{\pgfqpoint{10.785557in}{5.837373in}}%
\pgfpathlineto{\pgfqpoint{10.785557in}{5.744203in}}%
\pgfpathclose%
\pgfusepath{stroke}%
\end{pgfscope}%
\begin{pgfscope}%
\pgfpathrectangle{\pgfqpoint{7.904091in}{2.928000in}}{\pgfqpoint{5.664000in}{5.664000in}}%
\pgfusepath{clip}%
\pgfsetbuttcap%
\pgfsetroundjoin%
\pgfsetlinewidth{0.803000pt}%
\definecolor{currentstroke}{rgb}{0.768627,0.603922,0.000000}%
\pgfsetstrokecolor{currentstroke}%
\pgfsetdash{}{0pt}%
\pgfpathmoveto{\pgfqpoint{10.441056in}{5.746140in}}%
\pgfpathlineto{\pgfqpoint{10.534226in}{5.746140in}}%
\pgfpathlineto{\pgfqpoint{10.534226in}{5.839309in}}%
\pgfpathlineto{\pgfqpoint{10.441056in}{5.839309in}}%
\pgfpathlineto{\pgfqpoint{10.441056in}{5.746140in}}%
\pgfpathclose%
\pgfusepath{stroke}%
\end{pgfscope}%
\begin{pgfscope}%
\pgfsetbuttcap%
\pgfsetroundjoin%
\pgfsetlinewidth{0.903375pt}%
\definecolor{currentstroke}{rgb}{0.768627,0.603922,0.000000}%
\pgfsetstrokecolor{currentstroke}%
\pgfsetdash{}{0pt}%
\pgfpathmoveto{\pgfqpoint{7.904091in}{2.855166in}}%
\pgfpathlineto{\pgfqpoint{7.976925in}{2.928000in}}%
\pgfpathlineto{\pgfqpoint{7.904091in}{3.000834in}}%
\pgfpathlineto{\pgfqpoint{7.831257in}{2.928000in}}%
\pgfpathlineto{\pgfqpoint{7.904091in}{2.855166in}}%
\pgfpathclose%
\pgfusepath{stroke}%
\end{pgfscope}%
\begin{pgfscope}%
\pgfsetbuttcap%
\pgfsetroundjoin%
\pgfsetlinewidth{0.903375pt}%
\definecolor{currentstroke}{rgb}{0.768627,0.603922,0.000000}%
\pgfsetstrokecolor{currentstroke}%
\pgfsetdash{}{0pt}%
\pgfpathmoveto{\pgfqpoint{12.935583in}{8.519166in}}%
\pgfpathlineto{\pgfqpoint{13.008417in}{8.592000in}}%
\pgfpathlineto{\pgfqpoint{12.935583in}{8.664834in}}%
\pgfpathlineto{\pgfqpoint{12.862749in}{8.592000in}}%
\pgfpathlineto{\pgfqpoint{12.935583in}{8.519166in}}%
\pgfpathclose%
\pgfusepath{stroke}%
\end{pgfscope}%
\begin{pgfscope}%
\pgfsetbuttcap%
\pgfsetroundjoin%
\pgfsetlinewidth{0.903375pt}%
\definecolor{currentstroke}{rgb}{0.768627,0.603922,0.000000}%
\pgfsetstrokecolor{currentstroke}%
\pgfsetdash{}{0pt}%
\pgfpathmoveto{\pgfqpoint{12.894696in}{8.519166in}}%
\pgfpathlineto{\pgfqpoint{12.967530in}{8.592000in}}%
\pgfpathlineto{\pgfqpoint{12.894696in}{8.664834in}}%
\pgfpathlineto{\pgfqpoint{12.821862in}{8.592000in}}%
\pgfpathlineto{\pgfqpoint{12.894696in}{8.519166in}}%
\pgfpathclose%
\pgfusepath{stroke}%
\end{pgfscope}%
\begin{pgfscope}%
\pgfsetbuttcap%
\pgfsetroundjoin%
\pgfsetlinewidth{0.903375pt}%
\definecolor{currentstroke}{rgb}{0.768627,0.603922,0.000000}%
\pgfsetstrokecolor{currentstroke}%
\pgfsetdash{}{0pt}%
\pgfpathmoveto{\pgfqpoint{12.723859in}{8.519166in}}%
\pgfpathlineto{\pgfqpoint{12.796693in}{8.592000in}}%
\pgfpathlineto{\pgfqpoint{12.723859in}{8.664834in}}%
\pgfpathlineto{\pgfqpoint{12.651025in}{8.592000in}}%
\pgfpathlineto{\pgfqpoint{12.723859in}{8.519166in}}%
\pgfpathclose%
\pgfusepath{stroke}%
\end{pgfscope}%
\begin{pgfscope}%
\pgfsetbuttcap%
\pgfsetroundjoin%
\pgfsetlinewidth{0.903375pt}%
\definecolor{currentstroke}{rgb}{0.768627,0.603922,0.000000}%
\pgfsetstrokecolor{currentstroke}%
\pgfsetdash{}{0pt}%
\pgfpathmoveto{\pgfqpoint{7.904091in}{2.855166in}}%
\pgfpathlineto{\pgfqpoint{7.976925in}{2.928000in}}%
\pgfpathlineto{\pgfqpoint{7.904091in}{3.000834in}}%
\pgfpathlineto{\pgfqpoint{7.831257in}{2.928000in}}%
\pgfpathlineto{\pgfqpoint{7.904091in}{2.855166in}}%
\pgfpathclose%
\pgfusepath{stroke}%
\end{pgfscope}%
\begin{pgfscope}%
\pgfsetbuttcap%
\pgfsetroundjoin%
\pgfsetlinewidth{0.903375pt}%
\definecolor{currentstroke}{rgb}{0.709804,0.741176,0.207843}%
\pgfsetstrokecolor{currentstroke}%
\pgfsetdash{}{0pt}%
\pgfpathmoveto{\pgfqpoint{8.497415in}{2.855166in}}%
\pgfpathlineto{\pgfqpoint{8.570249in}{2.928000in}}%
\pgfpathlineto{\pgfqpoint{8.497415in}{3.000834in}}%
\pgfpathlineto{\pgfqpoint{8.424581in}{2.928000in}}%
\pgfpathlineto{\pgfqpoint{8.497415in}{2.855166in}}%
\pgfpathclose%
\pgfusepath{stroke}%
\end{pgfscope}%
\begin{pgfscope}%
\pgfsetbuttcap%
\pgfsetroundjoin%
\pgfsetlinewidth{0.903375pt}%
\definecolor{currentstroke}{rgb}{0.709804,0.741176,0.207843}%
\pgfsetstrokecolor{currentstroke}%
\pgfsetdash{}{0pt}%
\pgfpathmoveto{\pgfqpoint{7.991605in}{2.855166in}}%
\pgfpathlineto{\pgfqpoint{8.064439in}{2.928000in}}%
\pgfpathlineto{\pgfqpoint{7.991605in}{3.000834in}}%
\pgfpathlineto{\pgfqpoint{7.918771in}{2.928000in}}%
\pgfpathlineto{\pgfqpoint{7.991605in}{2.855166in}}%
\pgfpathclose%
\pgfusepath{stroke}%
\end{pgfscope}%
\begin{pgfscope}%
\pgfsetbuttcap%
\pgfsetroundjoin%
\pgfsetlinewidth{0.903375pt}%
\definecolor{currentstroke}{rgb}{0.709804,0.741176,0.207843}%
\pgfsetstrokecolor{currentstroke}%
\pgfsetdash{}{0pt}%
\pgfpathmoveto{\pgfqpoint{8.248787in}{2.855166in}}%
\pgfpathlineto{\pgfqpoint{8.321621in}{2.928000in}}%
\pgfpathlineto{\pgfqpoint{8.248787in}{3.000834in}}%
\pgfpathlineto{\pgfqpoint{8.175953in}{2.928000in}}%
\pgfpathlineto{\pgfqpoint{8.248787in}{2.855166in}}%
\pgfpathclose%
\pgfusepath{stroke}%
\end{pgfscope}%
\begin{pgfscope}%
\pgfsetbuttcap%
\pgfsetroundjoin%
\pgfsetlinewidth{0.903375pt}%
\definecolor{currentstroke}{rgb}{0.709804,0.741176,0.207843}%
\pgfsetstrokecolor{currentstroke}%
\pgfsetdash{}{0pt}%
\pgfpathmoveto{\pgfqpoint{7.904091in}{2.855166in}}%
\pgfpathlineto{\pgfqpoint{7.976925in}{2.928000in}}%
\pgfpathlineto{\pgfqpoint{7.904091in}{3.000834in}}%
\pgfpathlineto{\pgfqpoint{7.831257in}{2.928000in}}%
\pgfpathlineto{\pgfqpoint{7.904091in}{2.855166in}}%
\pgfpathclose%
\pgfusepath{stroke}%
\end{pgfscope}%
\begin{pgfscope}%
\pgfsetbuttcap%
\pgfsetroundjoin%
\pgfsetlinewidth{0.903375pt}%
\definecolor{currentstroke}{rgb}{0.709804,0.741176,0.207843}%
\pgfsetstrokecolor{currentstroke}%
\pgfsetdash{}{0pt}%
\pgfpathmoveto{\pgfqpoint{7.904091in}{2.855166in}}%
\pgfpathlineto{\pgfqpoint{7.976925in}{2.928000in}}%
\pgfpathlineto{\pgfqpoint{7.904091in}{3.000834in}}%
\pgfpathlineto{\pgfqpoint{7.831257in}{2.928000in}}%
\pgfpathlineto{\pgfqpoint{7.904091in}{2.855166in}}%
\pgfpathclose%
\pgfusepath{stroke}%
\end{pgfscope}%
\begin{pgfscope}%
\pgfsetbuttcap%
\pgfsetroundjoin%
\pgfsetlinewidth{0.803000pt}%
\definecolor{currentstroke}{rgb}{0.901961,0.196078,0.156863}%
\pgfsetstrokecolor{currentstroke}%
\pgfsetdash{}{0pt}%
\pgfpathmoveto{\pgfqpoint{13.568091in}{8.498314in}}%
\pgfpathlineto{\pgfqpoint{13.661777in}{8.592000in}}%
\pgfpathlineto{\pgfqpoint{13.568091in}{8.685686in}}%
\pgfpathlineto{\pgfqpoint{13.474405in}{8.592000in}}%
\pgfpathlineto{\pgfqpoint{13.568091in}{8.498314in}}%
\pgfpathclose%
\pgfusepath{stroke}%
\end{pgfscope}%
\begin{pgfscope}%
\pgfsetbuttcap%
\pgfsetroundjoin%
\pgfsetlinewidth{0.803000pt}%
\definecolor{currentstroke}{rgb}{0.901961,0.196078,0.156863}%
\pgfsetstrokecolor{currentstroke}%
\pgfsetdash{}{0pt}%
\pgfpathmoveto{\pgfqpoint{13.521506in}{8.545415in}}%
\pgfpathlineto{\pgfqpoint{13.614676in}{8.545415in}}%
\pgfpathlineto{\pgfqpoint{13.614676in}{8.638585in}}%
\pgfpathlineto{\pgfqpoint{13.521506in}{8.638585in}}%
\pgfpathlineto{\pgfqpoint{13.521506in}{8.545415in}}%
\pgfpathclose%
\pgfusepath{stroke}%
\end{pgfscope}%
\begin{pgfscope}%
\pgfsetbuttcap%
\pgfsetroundjoin%
\pgfsetlinewidth{0.803000pt}%
\definecolor{currentstroke}{rgb}{0.498039,0.560784,0.678431}%
\pgfsetstrokecolor{currentstroke}%
\pgfsetdash{}{0pt}%
\pgfpathmoveto{\pgfqpoint{7.904091in}{2.834314in}}%
\pgfpathlineto{\pgfqpoint{7.997777in}{2.928000in}}%
\pgfpathlineto{\pgfqpoint{7.904091in}{3.021686in}}%
\pgfpathlineto{\pgfqpoint{7.810405in}{2.928000in}}%
\pgfpathlineto{\pgfqpoint{7.904091in}{2.834314in}}%
\pgfpathclose%
\pgfusepath{stroke}%
\end{pgfscope}%
\begin{pgfscope}%
\pgfsetbuttcap%
\pgfsetroundjoin%
\pgfsetlinewidth{0.803000pt}%
\definecolor{currentstroke}{rgb}{0.709804,0.121569,0.270588}%
\pgfsetstrokecolor{currentstroke}%
\pgfsetdash{}{0pt}%
\pgfpathmoveto{\pgfqpoint{13.313868in}{8.498314in}}%
\pgfpathlineto{\pgfqpoint{13.407554in}{8.592000in}}%
\pgfpathlineto{\pgfqpoint{13.313868in}{8.685686in}}%
\pgfpathlineto{\pgfqpoint{13.220183in}{8.592000in}}%
\pgfpathlineto{\pgfqpoint{13.313868in}{8.498314in}}%
\pgfpathclose%
\pgfusepath{stroke}%
\end{pgfscope}%
\begin{pgfscope}%
\pgfpathrectangle{\pgfqpoint{7.904091in}{2.928000in}}{\pgfqpoint{5.664000in}{5.664000in}}%
\pgfusepath{clip}%
\pgfsetbuttcap%
\pgfsetroundjoin%
\definecolor{currentfill}{rgb}{0.541176,0.380392,0.658824}%
\pgfsetfillcolor{currentfill}%
\pgfsetlinewidth{1.254687pt}%
\definecolor{currentstroke}{rgb}{0.541176,0.380392,0.658824}%
\pgfsetstrokecolor{currentstroke}%
\pgfsetdash{}{0pt}%
\pgfsys@defobject{currentmarker}{\pgfqpoint{-0.043368in}{-0.043368in}}{\pgfqpoint{0.043368in}{0.043368in}}{%
\pgfpathmoveto{\pgfqpoint{-0.043368in}{-0.043368in}}%
\pgfpathlineto{\pgfqpoint{0.043368in}{0.043368in}}%
\pgfpathmoveto{\pgfqpoint{-0.043368in}{0.043368in}}%
\pgfpathlineto{\pgfqpoint{0.043368in}{-0.043368in}}%
\pgfusepath{stroke,fill}%
}%
\begin{pgfscope}%
\pgfsys@transformshift{11.420804in}{7.863346in}%
\pgfsys@useobject{currentmarker}{}%
\end{pgfscope}%
\begin{pgfscope}%
\pgfsys@transformshift{12.242752in}{7.902909in}%
\pgfsys@useobject{currentmarker}{}%
\end{pgfscope}%
\begin{pgfscope}%
\pgfsys@transformshift{10.954415in}{8.258300in}%
\pgfsys@useobject{currentmarker}{}%
\end{pgfscope}%
\begin{pgfscope}%
\pgfsys@transformshift{10.355715in}{5.720483in}%
\pgfsys@useobject{currentmarker}{}%
\end{pgfscope}%
\begin{pgfscope}%
\pgfsys@transformshift{11.108777in}{5.485944in}%
\pgfsys@useobject{currentmarker}{}%
\end{pgfscope}%
\begin{pgfscope}%
\pgfsys@transformshift{12.045118in}{4.856815in}%
\pgfsys@useobject{currentmarker}{}%
\end{pgfscope}%
\begin{pgfscope}%
\pgfsys@transformshift{10.377284in}{3.951490in}%
\pgfsys@useobject{currentmarker}{}%
\end{pgfscope}%
\begin{pgfscope}%
\pgfsys@transformshift{10.380930in}{5.764121in}%
\pgfsys@useobject{currentmarker}{}%
\end{pgfscope}%
\begin{pgfscope}%
\pgfsys@transformshift{10.197066in}{4.550944in}%
\pgfsys@useobject{currentmarker}{}%
\end{pgfscope}%
\begin{pgfscope}%
\pgfsys@transformshift{10.834736in}{6.094635in}%
\pgfsys@useobject{currentmarker}{}%
\end{pgfscope}%
\begin{pgfscope}%
\pgfsys@transformshift{11.119437in}{6.161875in}%
\pgfsys@useobject{currentmarker}{}%
\end{pgfscope}%
\end{pgfscope}%
\begin{pgfscope}%
\pgfpathrectangle{\pgfqpoint{7.904091in}{2.928000in}}{\pgfqpoint{5.664000in}{5.664000in}}%
\pgfusepath{clip}%
\pgfsetbuttcap%
\pgfsetroundjoin%
\definecolor{currentfill}{rgb}{0.321569,0.321569,0.321569}%
\pgfsetfillcolor{currentfill}%
\pgfsetlinewidth{1.254687pt}%
\definecolor{currentstroke}{rgb}{0.321569,0.321569,0.321569}%
\pgfsetstrokecolor{currentstroke}%
\pgfsetdash{}{0pt}%
\pgfsys@defobject{currentmarker}{\pgfqpoint{-0.043368in}{-0.043368in}}{\pgfqpoint{0.043368in}{0.043368in}}{%
\pgfpathmoveto{\pgfqpoint{-0.043368in}{-0.043368in}}%
\pgfpathlineto{\pgfqpoint{0.043368in}{0.043368in}}%
\pgfpathmoveto{\pgfqpoint{-0.043368in}{0.043368in}}%
\pgfpathlineto{\pgfqpoint{0.043368in}{-0.043368in}}%
\pgfusepath{stroke,fill}%
}%
\begin{pgfscope}%
\pgfsys@transformshift{11.925480in}{7.353540in}%
\pgfsys@useobject{currentmarker}{}%
\end{pgfscope}%
\begin{pgfscope}%
\pgfsys@transformshift{11.863091in}{6.477847in}%
\pgfsys@useobject{currentmarker}{}%
\end{pgfscope}%
\begin{pgfscope}%
\pgfsys@transformshift{12.038330in}{6.610893in}%
\pgfsys@useobject{currentmarker}{}%
\end{pgfscope}%
\begin{pgfscope}%
\pgfsys@transformshift{12.232758in}{6.867248in}%
\pgfsys@useobject{currentmarker}{}%
\end{pgfscope}%
\begin{pgfscope}%
\pgfsys@transformshift{11.912148in}{6.414383in}%
\pgfsys@useobject{currentmarker}{}%
\end{pgfscope}%
\begin{pgfscope}%
\pgfsys@transformshift{9.410817in}{4.397013in}%
\pgfsys@useobject{currentmarker}{}%
\end{pgfscope}%
\begin{pgfscope}%
\pgfsys@transformshift{9.413841in}{4.570754in}%
\pgfsys@useobject{currentmarker}{}%
\end{pgfscope}%
\begin{pgfscope}%
\pgfsys@transformshift{9.368585in}{4.367735in}%
\pgfsys@useobject{currentmarker}{}%
\end{pgfscope}%
\begin{pgfscope}%
\pgfsys@transformshift{11.678865in}{6.880559in}%
\pgfsys@useobject{currentmarker}{}%
\end{pgfscope}%
\begin{pgfscope}%
\pgfsys@transformshift{9.221651in}{4.489192in}%
\pgfsys@useobject{currentmarker}{}%
\end{pgfscope}%
\begin{pgfscope}%
\pgfsys@transformshift{10.296765in}{5.364099in}%
\pgfsys@useobject{currentmarker}{}%
\end{pgfscope}%
\begin{pgfscope}%
\pgfsys@transformshift{8.993135in}{3.272937in}%
\pgfsys@useobject{currentmarker}{}%
\end{pgfscope}%
\begin{pgfscope}%
\pgfsys@transformshift{11.651610in}{5.815443in}%
\pgfsys@useobject{currentmarker}{}%
\end{pgfscope}%
\begin{pgfscope}%
\pgfsys@transformshift{13.549087in}{8.418024in}%
\pgfsys@useobject{currentmarker}{}%
\end{pgfscope}%
\end{pgfscope}%
\begin{pgfscope}%
\pgfsetbuttcap%
\pgfsetroundjoin%
\definecolor{currentfill}{rgb}{0.321569,0.321569,0.321569}%
\pgfsetfillcolor{currentfill}%
\pgfsetlinewidth{1.254687pt}%
\definecolor{currentstroke}{rgb}{0.321569,0.321569,0.321569}%
\pgfsetstrokecolor{currentstroke}%
\pgfsetdash{}{0pt}%
\pgfsys@defobject{currentmarker}{\pgfqpoint{-0.043368in}{-0.043368in}}{\pgfqpoint{0.043368in}{0.043368in}}{%
\pgfpathmoveto{\pgfqpoint{-0.043368in}{-0.043368in}}%
\pgfpathlineto{\pgfqpoint{0.043368in}{0.043368in}}%
\pgfpathmoveto{\pgfqpoint{-0.043368in}{0.043368in}}%
\pgfpathlineto{\pgfqpoint{0.043368in}{-0.043368in}}%
\pgfusepath{stroke,fill}%
}%
\begin{pgfscope}%
\pgfsys@transformshift{7.904091in}{2.928000in}%
\pgfsys@useobject{currentmarker}{}%
\end{pgfscope}%
\begin{pgfscope}%
\pgfsys@transformshift{7.904091in}{2.928000in}%
\pgfsys@useobject{currentmarker}{}%
\end{pgfscope}%
\begin{pgfscope}%
\pgfsys@transformshift{7.904091in}{2.928000in}%
\pgfsys@useobject{currentmarker}{}%
\end{pgfscope}%
\begin{pgfscope}%
\pgfsys@transformshift{7.904091in}{2.928000in}%
\pgfsys@useobject{currentmarker}{}%
\end{pgfscope}%
\begin{pgfscope}%
\pgfsys@transformshift{7.904091in}{2.928000in}%
\pgfsys@useobject{currentmarker}{}%
\end{pgfscope}%
\begin{pgfscope}%
\pgfsys@transformshift{7.904091in}{2.928000in}%
\pgfsys@useobject{currentmarker}{}%
\end{pgfscope}%
\end{pgfscope}%
\begin{pgfscope}%
\pgfpathrectangle{\pgfqpoint{7.904091in}{2.928000in}}{\pgfqpoint{5.664000in}{5.664000in}}%
\pgfusepath{clip}%
\pgfsetbuttcap%
\pgfsetroundjoin%
\definecolor{currentfill}{rgb}{0.768627,0.603922,0.000000}%
\pgfsetfillcolor{currentfill}%
\pgfsetlinewidth{1.254687pt}%
\definecolor{currentstroke}{rgb}{0.768627,0.603922,0.000000}%
\pgfsetstrokecolor{currentstroke}%
\pgfsetdash{}{0pt}%
\pgfsys@defobject{currentmarker}{\pgfqpoint{-0.043368in}{-0.043368in}}{\pgfqpoint{0.043368in}{0.043368in}}{%
\pgfpathmoveto{\pgfqpoint{-0.043368in}{-0.043368in}}%
\pgfpathlineto{\pgfqpoint{0.043368in}{0.043368in}}%
\pgfpathmoveto{\pgfqpoint{-0.043368in}{0.043368in}}%
\pgfpathlineto{\pgfqpoint{0.043368in}{-0.043368in}}%
\pgfusepath{stroke,fill}%
}%
\begin{pgfscope}%
\pgfsys@transformshift{11.395160in}{6.387455in}%
\pgfsys@useobject{currentmarker}{}%
\end{pgfscope}%
\begin{pgfscope}%
\pgfsys@transformshift{11.659012in}{6.522295in}%
\pgfsys@useobject{currentmarker}{}%
\end{pgfscope}%
\end{pgfscope}%
\begin{pgfscope}%
\pgfpathrectangle{\pgfqpoint{7.904091in}{2.928000in}}{\pgfqpoint{5.664000in}{5.664000in}}%
\pgfusepath{clip}%
\pgfsetbuttcap%
\pgfsetroundjoin%
\definecolor{currentfill}{rgb}{0.768627,0.603922,0.000000}%
\pgfsetfillcolor{currentfill}%
\pgfsetlinewidth{1.254687pt}%
\definecolor{currentstroke}{rgb}{0.768627,0.603922,0.000000}%
\pgfsetstrokecolor{currentstroke}%
\pgfsetdash{}{0pt}%
\pgfsys@defobject{currentmarker}{\pgfqpoint{-0.043368in}{-0.043368in}}{\pgfqpoint{0.043368in}{0.043368in}}{%
\pgfpathmoveto{\pgfqpoint{-0.043368in}{-0.043368in}}%
\pgfpathlineto{\pgfqpoint{0.043368in}{0.043368in}}%
\pgfpathmoveto{\pgfqpoint{-0.043368in}{0.043368in}}%
\pgfpathlineto{\pgfqpoint{0.043368in}{-0.043368in}}%
\pgfusepath{stroke,fill}%
}%
\begin{pgfscope}%
\pgfsys@transformshift{10.832141in}{5.790788in}%
\pgfsys@useobject{currentmarker}{}%
\end{pgfscope}%
\begin{pgfscope}%
\pgfsys@transformshift{10.487641in}{5.792724in}%
\pgfsys@useobject{currentmarker}{}%
\end{pgfscope}%
\end{pgfscope}%
\begin{pgfscope}%
\pgfpathrectangle{\pgfqpoint{7.904091in}{2.928000in}}{\pgfqpoint{5.664000in}{5.664000in}}%
\pgfusepath{clip}%
\pgfsetbuttcap%
\pgfsetroundjoin%
\definecolor{currentfill}{rgb}{0.768627,0.603922,0.000000}%
\pgfsetfillcolor{currentfill}%
\pgfsetlinewidth{1.254687pt}%
\definecolor{currentstroke}{rgb}{0.768627,0.603922,0.000000}%
\pgfsetstrokecolor{currentstroke}%
\pgfsetdash{}{0pt}%
\pgfsys@defobject{currentmarker}{\pgfqpoint{-0.043368in}{-0.043368in}}{\pgfqpoint{0.043368in}{0.043368in}}{%
\pgfpathmoveto{\pgfqpoint{-0.043368in}{-0.043368in}}%
\pgfpathlineto{\pgfqpoint{0.043368in}{0.043368in}}%
\pgfpathmoveto{\pgfqpoint{-0.043368in}{0.043368in}}%
\pgfpathlineto{\pgfqpoint{0.043368in}{-0.043368in}}%
\pgfusepath{stroke,fill}%
}%
\begin{pgfscope}%
\pgfsys@transformshift{11.470217in}{6.748279in}%
\pgfsys@useobject{currentmarker}{}%
\end{pgfscope}%
\begin{pgfscope}%
\pgfsys@transformshift{11.197077in}{6.915647in}%
\pgfsys@useobject{currentmarker}{}%
\end{pgfscope}%
\begin{pgfscope}%
\pgfsys@transformshift{8.449436in}{3.773569in}%
\pgfsys@useobject{currentmarker}{}%
\end{pgfscope}%
\begin{pgfscope}%
\pgfsys@transformshift{11.072545in}{5.632565in}%
\pgfsys@useobject{currentmarker}{}%
\end{pgfscope}%
\begin{pgfscope}%
\pgfsys@transformshift{12.433741in}{7.920518in}%
\pgfsys@useobject{currentmarker}{}%
\end{pgfscope}%
\end{pgfscope}%
\begin{pgfscope}%
\pgfsetbuttcap%
\pgfsetroundjoin%
\definecolor{currentfill}{rgb}{0.768627,0.603922,0.000000}%
\pgfsetfillcolor{currentfill}%
\pgfsetlinewidth{1.254687pt}%
\definecolor{currentstroke}{rgb}{0.768627,0.603922,0.000000}%
\pgfsetstrokecolor{currentstroke}%
\pgfsetdash{}{0pt}%
\pgfsys@defobject{currentmarker}{\pgfqpoint{-0.043368in}{-0.043368in}}{\pgfqpoint{0.043368in}{0.043368in}}{%
\pgfpathmoveto{\pgfqpoint{-0.043368in}{-0.043368in}}%
\pgfpathlineto{\pgfqpoint{0.043368in}{0.043368in}}%
\pgfpathmoveto{\pgfqpoint{-0.043368in}{0.043368in}}%
\pgfpathlineto{\pgfqpoint{0.043368in}{-0.043368in}}%
\pgfusepath{stroke,fill}%
}%
\begin{pgfscope}%
\pgfsys@transformshift{7.904091in}{2.928000in}%
\pgfsys@useobject{currentmarker}{}%
\end{pgfscope}%
\begin{pgfscope}%
\pgfsys@transformshift{12.935583in}{8.592000in}%
\pgfsys@useobject{currentmarker}{}%
\end{pgfscope}%
\begin{pgfscope}%
\pgfsys@transformshift{12.894696in}{8.592000in}%
\pgfsys@useobject{currentmarker}{}%
\end{pgfscope}%
\begin{pgfscope}%
\pgfsys@transformshift{12.723859in}{8.592000in}%
\pgfsys@useobject{currentmarker}{}%
\end{pgfscope}%
\begin{pgfscope}%
\pgfsys@transformshift{7.904091in}{2.928000in}%
\pgfsys@useobject{currentmarker}{}%
\end{pgfscope}%
\end{pgfscope}%
\begin{pgfscope}%
\pgfpathrectangle{\pgfqpoint{7.904091in}{2.928000in}}{\pgfqpoint{5.664000in}{5.664000in}}%
\pgfusepath{clip}%
\pgfsetbuttcap%
\pgfsetroundjoin%
\definecolor{currentfill}{rgb}{0.709804,0.741176,0.207843}%
\pgfsetfillcolor{currentfill}%
\pgfsetlinewidth{1.254687pt}%
\definecolor{currentstroke}{rgb}{0.709804,0.741176,0.207843}%
\pgfsetstrokecolor{currentstroke}%
\pgfsetdash{}{0pt}%
\pgfsys@defobject{currentmarker}{\pgfqpoint{-0.043368in}{-0.043368in}}{\pgfqpoint{0.043368in}{0.043368in}}{%
\pgfpathmoveto{\pgfqpoint{-0.043368in}{-0.043368in}}%
\pgfpathlineto{\pgfqpoint{0.043368in}{0.043368in}}%
\pgfpathmoveto{\pgfqpoint{-0.043368in}{0.043368in}}%
\pgfpathlineto{\pgfqpoint{0.043368in}{-0.043368in}}%
\pgfusepath{stroke,fill}%
}%
\begin{pgfscope}%
\pgfsys@transformshift{9.052780in}{3.710716in}%
\pgfsys@useobject{currentmarker}{}%
\end{pgfscope}%
\begin{pgfscope}%
\pgfsys@transformshift{9.361012in}{4.151034in}%
\pgfsys@useobject{currentmarker}{}%
\end{pgfscope}%
\begin{pgfscope}%
\pgfsys@transformshift{8.742359in}{3.219411in}%
\pgfsys@useobject{currentmarker}{}%
\end{pgfscope}%
\begin{pgfscope}%
\pgfsys@transformshift{8.859330in}{4.389232in}%
\pgfsys@useobject{currentmarker}{}%
\end{pgfscope}%
\begin{pgfscope}%
\pgfsys@transformshift{9.351925in}{4.285701in}%
\pgfsys@useobject{currentmarker}{}%
\end{pgfscope}%
\begin{pgfscope}%
\pgfsys@transformshift{9.230171in}{3.248140in}%
\pgfsys@useobject{currentmarker}{}%
\end{pgfscope}%
\begin{pgfscope}%
\pgfsys@transformshift{8.851857in}{4.427600in}%
\pgfsys@useobject{currentmarker}{}%
\end{pgfscope}%
\begin{pgfscope}%
\pgfsys@transformshift{9.465941in}{4.450731in}%
\pgfsys@useobject{currentmarker}{}%
\end{pgfscope}%
\begin{pgfscope}%
\pgfsys@transformshift{10.085561in}{4.603721in}%
\pgfsys@useobject{currentmarker}{}%
\end{pgfscope}%
\begin{pgfscope}%
\pgfsys@transformshift{9.712562in}{4.569786in}%
\pgfsys@useobject{currentmarker}{}%
\end{pgfscope}%
\begin{pgfscope}%
\pgfsys@transformshift{10.577246in}{5.286814in}%
\pgfsys@useobject{currentmarker}{}%
\end{pgfscope}%
\end{pgfscope}%
\begin{pgfscope}%
\pgfsetbuttcap%
\pgfsetroundjoin%
\definecolor{currentfill}{rgb}{0.709804,0.741176,0.207843}%
\pgfsetfillcolor{currentfill}%
\pgfsetlinewidth{1.254687pt}%
\definecolor{currentstroke}{rgb}{0.709804,0.741176,0.207843}%
\pgfsetstrokecolor{currentstroke}%
\pgfsetdash{}{0pt}%
\pgfsys@defobject{currentmarker}{\pgfqpoint{-0.043368in}{-0.043368in}}{\pgfqpoint{0.043368in}{0.043368in}}{%
\pgfpathmoveto{\pgfqpoint{-0.043368in}{-0.043368in}}%
\pgfpathlineto{\pgfqpoint{0.043368in}{0.043368in}}%
\pgfpathmoveto{\pgfqpoint{-0.043368in}{0.043368in}}%
\pgfpathlineto{\pgfqpoint{0.043368in}{-0.043368in}}%
\pgfusepath{stroke,fill}%
}%
\begin{pgfscope}%
\pgfsys@transformshift{8.497415in}{2.928000in}%
\pgfsys@useobject{currentmarker}{}%
\end{pgfscope}%
\begin{pgfscope}%
\pgfsys@transformshift{7.991605in}{2.928000in}%
\pgfsys@useobject{currentmarker}{}%
\end{pgfscope}%
\begin{pgfscope}%
\pgfsys@transformshift{8.248787in}{2.928000in}%
\pgfsys@useobject{currentmarker}{}%
\end{pgfscope}%
\begin{pgfscope}%
\pgfsys@transformshift{7.904091in}{2.928000in}%
\pgfsys@useobject{currentmarker}{}%
\end{pgfscope}%
\begin{pgfscope}%
\pgfsys@transformshift{7.904091in}{2.928000in}%
\pgfsys@useobject{currentmarker}{}%
\end{pgfscope}%
\end{pgfscope}%
\begin{pgfscope}%
\pgfpathrectangle{\pgfqpoint{7.904091in}{2.928000in}}{\pgfqpoint{5.664000in}{5.664000in}}%
\pgfusepath{clip}%
\pgfsetbuttcap%
\pgfsetroundjoin%
\definecolor{currentfill}{rgb}{0.200000,0.133333,0.533333}%
\pgfsetfillcolor{currentfill}%
\pgfsetlinewidth{1.254687pt}%
\definecolor{currentstroke}{rgb}{0.200000,0.133333,0.533333}%
\pgfsetstrokecolor{currentstroke}%
\pgfsetdash{}{0pt}%
\pgfsys@defobject{currentmarker}{\pgfqpoint{-0.043368in}{-0.043368in}}{\pgfqpoint{0.043368in}{0.043368in}}{%
\pgfpathmoveto{\pgfqpoint{-0.043368in}{-0.043368in}}%
\pgfpathlineto{\pgfqpoint{0.043368in}{0.043368in}}%
\pgfpathmoveto{\pgfqpoint{-0.043368in}{0.043368in}}%
\pgfpathlineto{\pgfqpoint{0.043368in}{-0.043368in}}%
\pgfusepath{stroke,fill}%
}%
\begin{pgfscope}%
\pgfsys@transformshift{8.896518in}{3.464148in}%
\pgfsys@useobject{currentmarker}{}%
\end{pgfscope}%
\begin{pgfscope}%
\pgfsys@transformshift{9.115204in}{3.667555in}%
\pgfsys@useobject{currentmarker}{}%
\end{pgfscope}%
\end{pgfscope}%
\begin{pgfscope}%
\pgfpathrectangle{\pgfqpoint{7.904091in}{2.928000in}}{\pgfqpoint{5.664000in}{5.664000in}}%
\pgfusepath{clip}%
\pgfsetbuttcap%
\pgfsetroundjoin%
\definecolor{currentfill}{rgb}{0.058824,0.541176,0.372549}%
\pgfsetfillcolor{currentfill}%
\pgfsetlinewidth{1.254687pt}%
\definecolor{currentstroke}{rgb}{0.058824,0.541176,0.372549}%
\pgfsetstrokecolor{currentstroke}%
\pgfsetdash{}{0pt}%
\pgfsys@defobject{currentmarker}{\pgfqpoint{-0.043368in}{-0.043368in}}{\pgfqpoint{0.043368in}{0.043368in}}{%
\pgfpathmoveto{\pgfqpoint{-0.043368in}{-0.043368in}}%
\pgfpathlineto{\pgfqpoint{0.043368in}{0.043368in}}%
\pgfpathmoveto{\pgfqpoint{-0.043368in}{0.043368in}}%
\pgfpathlineto{\pgfqpoint{0.043368in}{-0.043368in}}%
\pgfusepath{stroke,fill}%
}%
\begin{pgfscope}%
\pgfsys@transformshift{11.119171in}{6.339932in}%
\pgfsys@useobject{currentmarker}{}%
\end{pgfscope}%
\begin{pgfscope}%
\pgfsys@transformshift{10.806523in}{5.706437in}%
\pgfsys@useobject{currentmarker}{}%
\end{pgfscope}%
\begin{pgfscope}%
\pgfsys@transformshift{10.912578in}{5.367841in}%
\pgfsys@useobject{currentmarker}{}%
\end{pgfscope}%
\begin{pgfscope}%
\pgfsys@transformshift{11.423515in}{6.773615in}%
\pgfsys@useobject{currentmarker}{}%
\end{pgfscope}%
\begin{pgfscope}%
\pgfsys@transformshift{11.679922in}{6.568356in}%
\pgfsys@useobject{currentmarker}{}%
\end{pgfscope}%
\begin{pgfscope}%
\pgfsys@transformshift{10.845323in}{5.594985in}%
\pgfsys@useobject{currentmarker}{}%
\end{pgfscope}%
\begin{pgfscope}%
\pgfsys@transformshift{10.603520in}{5.057864in}%
\pgfsys@useobject{currentmarker}{}%
\end{pgfscope}%
\end{pgfscope}%
\begin{pgfscope}%
\pgfpathrectangle{\pgfqpoint{7.904091in}{2.928000in}}{\pgfqpoint{5.664000in}{5.664000in}}%
\pgfusepath{clip}%
\pgfsetbuttcap%
\pgfsetroundjoin%
\definecolor{currentfill}{rgb}{0.337255,0.705882,0.913725}%
\pgfsetfillcolor{currentfill}%
\pgfsetlinewidth{1.254687pt}%
\definecolor{currentstroke}{rgb}{0.337255,0.705882,0.913725}%
\pgfsetstrokecolor{currentstroke}%
\pgfsetdash{}{0pt}%
\pgfsys@defobject{currentmarker}{\pgfqpoint{-0.043368in}{-0.043368in}}{\pgfqpoint{0.043368in}{0.043368in}}{%
\pgfpathmoveto{\pgfqpoint{-0.043368in}{-0.043368in}}%
\pgfpathlineto{\pgfqpoint{0.043368in}{0.043368in}}%
\pgfpathmoveto{\pgfqpoint{-0.043368in}{0.043368in}}%
\pgfpathlineto{\pgfqpoint{0.043368in}{-0.043368in}}%
\pgfusepath{stroke,fill}%
}%
\begin{pgfscope}%
\pgfsys@transformshift{8.523168in}{4.615867in}%
\pgfsys@useobject{currentmarker}{}%
\end{pgfscope}%
\begin{pgfscope}%
\pgfsys@transformshift{8.680631in}{4.303285in}%
\pgfsys@useobject{currentmarker}{}%
\end{pgfscope}%
\end{pgfscope}%
\begin{pgfscope}%
\pgfpathrectangle{\pgfqpoint{7.904091in}{2.928000in}}{\pgfqpoint{5.664000in}{5.664000in}}%
\pgfusepath{clip}%
\pgfsetbuttcap%
\pgfsetroundjoin%
\definecolor{currentfill}{rgb}{0.901961,0.196078,0.156863}%
\pgfsetfillcolor{currentfill}%
\pgfsetlinewidth{1.254687pt}%
\definecolor{currentstroke}{rgb}{0.901961,0.196078,0.156863}%
\pgfsetstrokecolor{currentstroke}%
\pgfsetdash{}{0pt}%
\pgfsys@defobject{currentmarker}{\pgfqpoint{-0.043368in}{-0.043368in}}{\pgfqpoint{0.043368in}{0.043368in}}{%
\pgfpathmoveto{\pgfqpoint{-0.043368in}{-0.043368in}}%
\pgfpathlineto{\pgfqpoint{0.043368in}{0.043368in}}%
\pgfpathmoveto{\pgfqpoint{-0.043368in}{0.043368in}}%
\pgfpathlineto{\pgfqpoint{0.043368in}{-0.043368in}}%
\pgfusepath{stroke,fill}%
}%
\begin{pgfscope}%
\pgfsys@transformshift{13.323459in}{8.407523in}%
\pgfsys@useobject{currentmarker}{}%
\end{pgfscope}%
\end{pgfscope}%
\begin{pgfscope}%
\pgfsetbuttcap%
\pgfsetroundjoin%
\definecolor{currentfill}{rgb}{0.901961,0.196078,0.156863}%
\pgfsetfillcolor{currentfill}%
\pgfsetlinewidth{1.254687pt}%
\definecolor{currentstroke}{rgb}{0.901961,0.196078,0.156863}%
\pgfsetstrokecolor{currentstroke}%
\pgfsetdash{}{0pt}%
\pgfsys@defobject{currentmarker}{\pgfqpoint{-0.043368in}{-0.043368in}}{\pgfqpoint{0.043368in}{0.043368in}}{%
\pgfpathmoveto{\pgfqpoint{-0.043368in}{-0.043368in}}%
\pgfpathlineto{\pgfqpoint{0.043368in}{0.043368in}}%
\pgfpathmoveto{\pgfqpoint{-0.043368in}{0.043368in}}%
\pgfpathlineto{\pgfqpoint{0.043368in}{-0.043368in}}%
\pgfusepath{stroke,fill}%
}%
\begin{pgfscope}%
\pgfsys@transformshift{13.568091in}{8.592000in}%
\pgfsys@useobject{currentmarker}{}%
\end{pgfscope}%
\end{pgfscope}%
\begin{pgfscope}%
\pgfpathrectangle{\pgfqpoint{7.904091in}{2.928000in}}{\pgfqpoint{5.664000in}{5.664000in}}%
\pgfusepath{clip}%
\pgfsetbuttcap%
\pgfsetroundjoin%
\definecolor{currentfill}{rgb}{0.835294,0.368627,0.000000}%
\pgfsetfillcolor{currentfill}%
\pgfsetlinewidth{1.254687pt}%
\definecolor{currentstroke}{rgb}{0.835294,0.368627,0.000000}%
\pgfsetstrokecolor{currentstroke}%
\pgfsetdash{}{0pt}%
\pgfsys@defobject{currentmarker}{\pgfqpoint{-0.043368in}{-0.043368in}}{\pgfqpoint{0.043368in}{0.043368in}}{%
\pgfpathmoveto{\pgfqpoint{-0.043368in}{-0.043368in}}%
\pgfpathlineto{\pgfqpoint{0.043368in}{0.043368in}}%
\pgfpathmoveto{\pgfqpoint{-0.043368in}{0.043368in}}%
\pgfpathlineto{\pgfqpoint{0.043368in}{-0.043368in}}%
\pgfusepath{stroke,fill}%
}%
\begin{pgfscope}%
\pgfsys@transformshift{11.941953in}{7.389026in}%
\pgfsys@useobject{currentmarker}{}%
\end{pgfscope}%
\end{pgfscope}%
\begin{pgfscope}%
\pgfpathrectangle{\pgfqpoint{7.904091in}{2.928000in}}{\pgfqpoint{5.664000in}{5.664000in}}%
\pgfusepath{clip}%
\pgfsetbuttcap%
\pgfsetroundjoin%
\pgfsetlinewidth{1.003750pt}%
\definecolor{currentstroke}{rgb}{0.541176,0.380392,0.658824}%
\pgfsetstrokecolor{currentstroke}%
\pgfsetstrokeopacity{0.900000}%
\pgfsetdash{}{0pt}%
\pgfpathmoveto{\pgfqpoint{12.436992in}{7.167467in}}%
\pgfpathlineto{\pgfqpoint{12.521475in}{7.167467in}}%
\pgfpathlineto{\pgfqpoint{12.521475in}{7.251950in}}%
\pgfpathlineto{\pgfqpoint{12.436992in}{7.251950in}}%
\pgfpathlineto{\pgfqpoint{12.436992in}{7.167467in}}%
\pgfpathclose%
\pgfusepath{stroke}%
\end{pgfscope}%
\begin{pgfscope}%
\pgfpathrectangle{\pgfqpoint{7.904091in}{2.928000in}}{\pgfqpoint{5.664000in}{5.664000in}}%
\pgfusepath{clip}%
\pgfsetbuttcap%
\pgfsetroundjoin%
\pgfsetlinewidth{1.003750pt}%
\definecolor{currentstroke}{rgb}{0.541176,0.380392,0.658824}%
\pgfsetstrokecolor{currentstroke}%
\pgfsetstrokeopacity{0.900000}%
\pgfsetdash{}{0pt}%
\pgfpathmoveto{\pgfqpoint{11.706947in}{6.310886in}}%
\pgfpathlineto{\pgfqpoint{11.791430in}{6.310886in}}%
\pgfpathlineto{\pgfqpoint{11.791430in}{6.395369in}}%
\pgfpathlineto{\pgfqpoint{11.706947in}{6.395369in}}%
\pgfpathlineto{\pgfqpoint{11.706947in}{6.310886in}}%
\pgfpathclose%
\pgfusepath{stroke}%
\end{pgfscope}%
\begin{pgfscope}%
\pgfpathrectangle{\pgfqpoint{7.904091in}{2.928000in}}{\pgfqpoint{5.664000in}{5.664000in}}%
\pgfusepath{clip}%
\pgfsetbuttcap%
\pgfsetroundjoin%
\pgfsetlinewidth{1.003750pt}%
\definecolor{currentstroke}{rgb}{0.541176,0.380392,0.658824}%
\pgfsetstrokecolor{currentstroke}%
\pgfsetstrokeopacity{0.900000}%
\pgfsetdash{}{0pt}%
\pgfpathmoveto{\pgfqpoint{11.476017in}{6.248328in}}%
\pgfpathlineto{\pgfqpoint{11.560499in}{6.248328in}}%
\pgfpathlineto{\pgfqpoint{11.560499in}{6.332811in}}%
\pgfpathlineto{\pgfqpoint{11.476017in}{6.332811in}}%
\pgfpathlineto{\pgfqpoint{11.476017in}{6.248328in}}%
\pgfpathclose%
\pgfusepath{stroke}%
\end{pgfscope}%
\begin{pgfscope}%
\pgfpathrectangle{\pgfqpoint{7.904091in}{2.928000in}}{\pgfqpoint{5.664000in}{5.664000in}}%
\pgfusepath{clip}%
\pgfsetbuttcap%
\pgfsetroundjoin%
\pgfsetlinewidth{1.003750pt}%
\definecolor{currentstroke}{rgb}{0.541176,0.380392,0.658824}%
\pgfsetstrokecolor{currentstroke}%
\pgfsetstrokeopacity{0.900000}%
\pgfsetdash{}{0pt}%
\pgfpathmoveto{\pgfqpoint{11.739105in}{7.173351in}}%
\pgfpathlineto{\pgfqpoint{11.823588in}{7.173351in}}%
\pgfpathlineto{\pgfqpoint{11.823588in}{7.257834in}}%
\pgfpathlineto{\pgfqpoint{11.739105in}{7.257834in}}%
\pgfpathlineto{\pgfqpoint{11.739105in}{7.173351in}}%
\pgfpathclose%
\pgfusepath{stroke}%
\end{pgfscope}%
\begin{pgfscope}%
\pgfpathrectangle{\pgfqpoint{7.904091in}{2.928000in}}{\pgfqpoint{5.664000in}{5.664000in}}%
\pgfusepath{clip}%
\pgfsetbuttcap%
\pgfsetroundjoin%
\pgfsetlinewidth{1.003750pt}%
\definecolor{currentstroke}{rgb}{0.541176,0.380392,0.658824}%
\pgfsetstrokecolor{currentstroke}%
\pgfsetstrokeopacity{0.900000}%
\pgfsetdash{}{0pt}%
\pgfpathmoveto{\pgfqpoint{11.298195in}{6.562731in}}%
\pgfpathlineto{\pgfqpoint{11.382678in}{6.562731in}}%
\pgfpathlineto{\pgfqpoint{11.382678in}{6.647213in}}%
\pgfpathlineto{\pgfqpoint{11.298195in}{6.647213in}}%
\pgfpathlineto{\pgfqpoint{11.298195in}{6.562731in}}%
\pgfpathclose%
\pgfusepath{stroke}%
\end{pgfscope}%
\begin{pgfscope}%
\pgfpathrectangle{\pgfqpoint{7.904091in}{2.928000in}}{\pgfqpoint{5.664000in}{5.664000in}}%
\pgfusepath{clip}%
\pgfsetbuttcap%
\pgfsetroundjoin%
\pgfsetlinewidth{1.003750pt}%
\definecolor{currentstroke}{rgb}{0.541176,0.380392,0.658824}%
\pgfsetstrokecolor{currentstroke}%
\pgfsetstrokeopacity{0.900000}%
\pgfsetdash{}{0pt}%
\pgfpathmoveto{\pgfqpoint{11.352770in}{6.730880in}}%
\pgfpathlineto{\pgfqpoint{11.437253in}{6.730880in}}%
\pgfpathlineto{\pgfqpoint{11.437253in}{6.815363in}}%
\pgfpathlineto{\pgfqpoint{11.352770in}{6.815363in}}%
\pgfpathlineto{\pgfqpoint{11.352770in}{6.730880in}}%
\pgfpathclose%
\pgfusepath{stroke}%
\end{pgfscope}%
\begin{pgfscope}%
\pgfpathrectangle{\pgfqpoint{7.904091in}{2.928000in}}{\pgfqpoint{5.664000in}{5.664000in}}%
\pgfusepath{clip}%
\pgfsetbuttcap%
\pgfsetroundjoin%
\pgfsetlinewidth{1.003750pt}%
\definecolor{currentstroke}{rgb}{0.541176,0.380392,0.658824}%
\pgfsetstrokecolor{currentstroke}%
\pgfsetstrokeopacity{0.900000}%
\pgfsetdash{}{0pt}%
\pgfpathmoveto{\pgfqpoint{11.845406in}{6.466139in}}%
\pgfpathlineto{\pgfqpoint{11.929889in}{6.466139in}}%
\pgfpathlineto{\pgfqpoint{11.929889in}{6.550622in}}%
\pgfpathlineto{\pgfqpoint{11.845406in}{6.550622in}}%
\pgfpathlineto{\pgfqpoint{11.845406in}{6.466139in}}%
\pgfpathclose%
\pgfusepath{stroke}%
\end{pgfscope}%
\begin{pgfscope}%
\pgfpathrectangle{\pgfqpoint{7.904091in}{2.928000in}}{\pgfqpoint{5.664000in}{5.664000in}}%
\pgfusepath{clip}%
\pgfsetbuttcap%
\pgfsetroundjoin%
\pgfsetlinewidth{1.003750pt}%
\definecolor{currentstroke}{rgb}{0.541176,0.380392,0.658824}%
\pgfsetstrokecolor{currentstroke}%
\pgfsetstrokeopacity{0.900000}%
\pgfsetdash{}{0pt}%
\pgfpathmoveto{\pgfqpoint{10.995146in}{5.881744in}}%
\pgfpathlineto{\pgfqpoint{11.079629in}{5.881744in}}%
\pgfpathlineto{\pgfqpoint{11.079629in}{5.966227in}}%
\pgfpathlineto{\pgfqpoint{10.995146in}{5.966227in}}%
\pgfpathlineto{\pgfqpoint{10.995146in}{5.881744in}}%
\pgfpathclose%
\pgfusepath{stroke}%
\end{pgfscope}%
\begin{pgfscope}%
\pgfpathrectangle{\pgfqpoint{7.904091in}{2.928000in}}{\pgfqpoint{5.664000in}{5.664000in}}%
\pgfusepath{clip}%
\pgfsetbuttcap%
\pgfsetroundjoin%
\pgfsetlinewidth{1.003750pt}%
\definecolor{currentstroke}{rgb}{0.541176,0.380392,0.658824}%
\pgfsetstrokecolor{currentstroke}%
\pgfsetstrokeopacity{0.900000}%
\pgfsetdash{}{0pt}%
\pgfpathmoveto{\pgfqpoint{10.288383in}{4.930898in}}%
\pgfpathlineto{\pgfqpoint{10.372866in}{4.930898in}}%
\pgfpathlineto{\pgfqpoint{10.372866in}{5.015381in}}%
\pgfpathlineto{\pgfqpoint{10.288383in}{5.015381in}}%
\pgfpathlineto{\pgfqpoint{10.288383in}{4.930898in}}%
\pgfpathclose%
\pgfusepath{stroke}%
\end{pgfscope}%
\begin{pgfscope}%
\pgfpathrectangle{\pgfqpoint{7.904091in}{2.928000in}}{\pgfqpoint{5.664000in}{5.664000in}}%
\pgfusepath{clip}%
\pgfsetbuttcap%
\pgfsetroundjoin%
\pgfsetlinewidth{1.003750pt}%
\definecolor{currentstroke}{rgb}{0.541176,0.380392,0.658824}%
\pgfsetstrokecolor{currentstroke}%
\pgfsetstrokeopacity{0.900000}%
\pgfsetdash{}{0pt}%
\pgfpathmoveto{\pgfqpoint{11.869735in}{5.937140in}}%
\pgfpathlineto{\pgfqpoint{11.954218in}{5.937140in}}%
\pgfpathlineto{\pgfqpoint{11.954218in}{6.021623in}}%
\pgfpathlineto{\pgfqpoint{11.869735in}{6.021623in}}%
\pgfpathlineto{\pgfqpoint{11.869735in}{5.937140in}}%
\pgfpathclose%
\pgfusepath{stroke}%
\end{pgfscope}%
\begin{pgfscope}%
\pgfpathrectangle{\pgfqpoint{7.904091in}{2.928000in}}{\pgfqpoint{5.664000in}{5.664000in}}%
\pgfusepath{clip}%
\pgfsetbuttcap%
\pgfsetroundjoin%
\pgfsetlinewidth{1.003750pt}%
\definecolor{currentstroke}{rgb}{0.541176,0.380392,0.658824}%
\pgfsetstrokecolor{currentstroke}%
\pgfsetstrokeopacity{0.900000}%
\pgfsetdash{}{0pt}%
\pgfpathmoveto{\pgfqpoint{11.451595in}{6.315582in}}%
\pgfpathlineto{\pgfqpoint{11.536078in}{6.315582in}}%
\pgfpathlineto{\pgfqpoint{11.536078in}{6.400065in}}%
\pgfpathlineto{\pgfqpoint{11.451595in}{6.400065in}}%
\pgfpathlineto{\pgfqpoint{11.451595in}{6.315582in}}%
\pgfpathclose%
\pgfusepath{stroke}%
\end{pgfscope}%
\begin{pgfscope}%
\pgfpathrectangle{\pgfqpoint{7.904091in}{2.928000in}}{\pgfqpoint{5.664000in}{5.664000in}}%
\pgfusepath{clip}%
\pgfsetbuttcap%
\pgfsetroundjoin%
\pgfsetlinewidth{1.003750pt}%
\definecolor{currentstroke}{rgb}{0.541176,0.380392,0.658824}%
\pgfsetstrokecolor{currentstroke}%
\pgfsetstrokeopacity{0.900000}%
\pgfsetdash{}{0pt}%
\pgfpathmoveto{\pgfqpoint{10.262711in}{4.998625in}}%
\pgfpathlineto{\pgfqpoint{10.347194in}{4.998625in}}%
\pgfpathlineto{\pgfqpoint{10.347194in}{5.083108in}}%
\pgfpathlineto{\pgfqpoint{10.262711in}{5.083108in}}%
\pgfpathlineto{\pgfqpoint{10.262711in}{4.998625in}}%
\pgfpathclose%
\pgfusepath{stroke}%
\end{pgfscope}%
\begin{pgfscope}%
\pgfpathrectangle{\pgfqpoint{7.904091in}{2.928000in}}{\pgfqpoint{5.664000in}{5.664000in}}%
\pgfusepath{clip}%
\pgfsetbuttcap%
\pgfsetroundjoin%
\pgfsetlinewidth{1.003750pt}%
\definecolor{currentstroke}{rgb}{0.541176,0.380392,0.658824}%
\pgfsetstrokecolor{currentstroke}%
\pgfsetstrokeopacity{0.900000}%
\pgfsetdash{}{0pt}%
\pgfpathmoveto{\pgfqpoint{9.816303in}{4.840213in}}%
\pgfpathlineto{\pgfqpoint{9.900786in}{4.840213in}}%
\pgfpathlineto{\pgfqpoint{9.900786in}{4.924695in}}%
\pgfpathlineto{\pgfqpoint{9.816303in}{4.924695in}}%
\pgfpathlineto{\pgfqpoint{9.816303in}{4.840213in}}%
\pgfpathclose%
\pgfusepath{stroke}%
\end{pgfscope}%
\begin{pgfscope}%
\pgfpathrectangle{\pgfqpoint{7.904091in}{2.928000in}}{\pgfqpoint{5.664000in}{5.664000in}}%
\pgfusepath{clip}%
\pgfsetbuttcap%
\pgfsetroundjoin%
\pgfsetlinewidth{1.003750pt}%
\definecolor{currentstroke}{rgb}{0.541176,0.380392,0.658824}%
\pgfsetstrokecolor{currentstroke}%
\pgfsetstrokeopacity{0.900000}%
\pgfsetdash{}{0pt}%
\pgfpathmoveto{\pgfqpoint{10.186004in}{5.209913in}}%
\pgfpathlineto{\pgfqpoint{10.270487in}{5.209913in}}%
\pgfpathlineto{\pgfqpoint{10.270487in}{5.294396in}}%
\pgfpathlineto{\pgfqpoint{10.186004in}{5.294396in}}%
\pgfpathlineto{\pgfqpoint{10.186004in}{5.209913in}}%
\pgfpathclose%
\pgfusepath{stroke}%
\end{pgfscope}%
\begin{pgfscope}%
\pgfpathrectangle{\pgfqpoint{7.904091in}{2.928000in}}{\pgfqpoint{5.664000in}{5.664000in}}%
\pgfusepath{clip}%
\pgfsetbuttcap%
\pgfsetroundjoin%
\pgfsetlinewidth{1.003750pt}%
\definecolor{currentstroke}{rgb}{0.541176,0.380392,0.658824}%
\pgfsetstrokecolor{currentstroke}%
\pgfsetstrokeopacity{0.900000}%
\pgfsetdash{}{0pt}%
\pgfpathmoveto{\pgfqpoint{9.814665in}{4.838574in}}%
\pgfpathlineto{\pgfqpoint{9.899148in}{4.838574in}}%
\pgfpathlineto{\pgfqpoint{9.899148in}{4.923057in}}%
\pgfpathlineto{\pgfqpoint{9.814665in}{4.923057in}}%
\pgfpathlineto{\pgfqpoint{9.814665in}{4.838574in}}%
\pgfpathclose%
\pgfusepath{stroke}%
\end{pgfscope}%
\begin{pgfscope}%
\pgfpathrectangle{\pgfqpoint{7.904091in}{2.928000in}}{\pgfqpoint{5.664000in}{5.664000in}}%
\pgfusepath{clip}%
\pgfsetbuttcap%
\pgfsetroundjoin%
\pgfsetlinewidth{1.003750pt}%
\definecolor{currentstroke}{rgb}{0.541176,0.380392,0.658824}%
\pgfsetstrokecolor{currentstroke}%
\pgfsetstrokeopacity{0.900000}%
\pgfsetdash{}{0pt}%
\pgfpathmoveto{\pgfqpoint{11.327317in}{6.351226in}}%
\pgfpathlineto{\pgfqpoint{11.411800in}{6.351226in}}%
\pgfpathlineto{\pgfqpoint{11.411800in}{6.435709in}}%
\pgfpathlineto{\pgfqpoint{11.327317in}{6.435709in}}%
\pgfpathlineto{\pgfqpoint{11.327317in}{6.351226in}}%
\pgfpathclose%
\pgfusepath{stroke}%
\end{pgfscope}%
\begin{pgfscope}%
\pgfpathrectangle{\pgfqpoint{7.904091in}{2.928000in}}{\pgfqpoint{5.664000in}{5.664000in}}%
\pgfusepath{clip}%
\pgfsetbuttcap%
\pgfsetroundjoin%
\pgfsetlinewidth{1.003750pt}%
\definecolor{currentstroke}{rgb}{0.541176,0.380392,0.658824}%
\pgfsetstrokecolor{currentstroke}%
\pgfsetstrokeopacity{0.900000}%
\pgfsetdash{}{0pt}%
\pgfpathmoveto{\pgfqpoint{10.917626in}{5.941535in}}%
\pgfpathlineto{\pgfqpoint{11.002109in}{5.941535in}}%
\pgfpathlineto{\pgfqpoint{11.002109in}{6.026018in}}%
\pgfpathlineto{\pgfqpoint{10.917626in}{6.026018in}}%
\pgfpathlineto{\pgfqpoint{10.917626in}{5.941535in}}%
\pgfpathclose%
\pgfusepath{stroke}%
\end{pgfscope}%
\begin{pgfscope}%
\pgfpathrectangle{\pgfqpoint{7.904091in}{2.928000in}}{\pgfqpoint{5.664000in}{5.664000in}}%
\pgfusepath{clip}%
\pgfsetbuttcap%
\pgfsetroundjoin%
\pgfsetlinewidth{1.003750pt}%
\definecolor{currentstroke}{rgb}{0.541176,0.380392,0.658824}%
\pgfsetstrokecolor{currentstroke}%
\pgfsetstrokeopacity{0.900000}%
\pgfsetdash{}{0pt}%
\pgfpathmoveto{\pgfqpoint{11.747496in}{6.771405in}}%
\pgfpathlineto{\pgfqpoint{11.831979in}{6.771405in}}%
\pgfpathlineto{\pgfqpoint{11.831979in}{6.855888in}}%
\pgfpathlineto{\pgfqpoint{11.747496in}{6.855888in}}%
\pgfpathlineto{\pgfqpoint{11.747496in}{6.771405in}}%
\pgfpathclose%
\pgfusepath{stroke}%
\end{pgfscope}%
\begin{pgfscope}%
\pgfpathrectangle{\pgfqpoint{7.904091in}{2.928000in}}{\pgfqpoint{5.664000in}{5.664000in}}%
\pgfusepath{clip}%
\pgfsetbuttcap%
\pgfsetroundjoin%
\pgfsetlinewidth{1.003750pt}%
\definecolor{currentstroke}{rgb}{0.541176,0.380392,0.658824}%
\pgfsetstrokecolor{currentstroke}%
\pgfsetstrokeopacity{0.900000}%
\pgfsetdash{}{0pt}%
\pgfpathmoveto{\pgfqpoint{10.161076in}{5.184985in}}%
\pgfpathlineto{\pgfqpoint{10.245559in}{5.184985in}}%
\pgfpathlineto{\pgfqpoint{10.245559in}{5.269468in}}%
\pgfpathlineto{\pgfqpoint{10.161076in}{5.269468in}}%
\pgfpathlineto{\pgfqpoint{10.161076in}{5.184985in}}%
\pgfpathclose%
\pgfusepath{stroke}%
\end{pgfscope}%
\begin{pgfscope}%
\pgfpathrectangle{\pgfqpoint{7.904091in}{2.928000in}}{\pgfqpoint{5.664000in}{5.664000in}}%
\pgfusepath{clip}%
\pgfsetbuttcap%
\pgfsetroundjoin%
\pgfsetlinewidth{1.003750pt}%
\definecolor{currentstroke}{rgb}{0.541176,0.380392,0.658824}%
\pgfsetstrokecolor{currentstroke}%
\pgfsetstrokeopacity{0.900000}%
\pgfsetdash{}{0pt}%
\pgfpathmoveto{\pgfqpoint{9.523705in}{4.547615in}}%
\pgfpathlineto{\pgfqpoint{9.608188in}{4.547615in}}%
\pgfpathlineto{\pgfqpoint{9.608188in}{4.632097in}}%
\pgfpathlineto{\pgfqpoint{9.523705in}{4.632097in}}%
\pgfpathlineto{\pgfqpoint{9.523705in}{4.547615in}}%
\pgfpathclose%
\pgfusepath{stroke}%
\end{pgfscope}%
\begin{pgfscope}%
\pgfpathrectangle{\pgfqpoint{7.904091in}{2.928000in}}{\pgfqpoint{5.664000in}{5.664000in}}%
\pgfusepath{clip}%
\pgfsetbuttcap%
\pgfsetroundjoin%
\pgfsetlinewidth{1.003750pt}%
\definecolor{currentstroke}{rgb}{0.541176,0.380392,0.658824}%
\pgfsetstrokecolor{currentstroke}%
\pgfsetstrokeopacity{0.900000}%
\pgfsetdash{}{0pt}%
\pgfpathmoveto{\pgfqpoint{10.576035in}{5.599944in}}%
\pgfpathlineto{\pgfqpoint{10.660518in}{5.599944in}}%
\pgfpathlineto{\pgfqpoint{10.660518in}{5.684427in}}%
\pgfpathlineto{\pgfqpoint{10.576035in}{5.684427in}}%
\pgfpathlineto{\pgfqpoint{10.576035in}{5.599944in}}%
\pgfpathclose%
\pgfusepath{stroke}%
\end{pgfscope}%
\begin{pgfscope}%
\pgfpathrectangle{\pgfqpoint{7.904091in}{2.928000in}}{\pgfqpoint{5.664000in}{5.664000in}}%
\pgfusepath{clip}%
\pgfsetbuttcap%
\pgfsetroundjoin%
\pgfsetlinewidth{1.003750pt}%
\definecolor{currentstroke}{rgb}{0.541176,0.380392,0.658824}%
\pgfsetstrokecolor{currentstroke}%
\pgfsetstrokeopacity{0.900000}%
\pgfsetdash{}{0pt}%
\pgfpathmoveto{\pgfqpoint{9.538157in}{4.562067in}}%
\pgfpathlineto{\pgfqpoint{9.622640in}{4.562067in}}%
\pgfpathlineto{\pgfqpoint{9.622640in}{4.646549in}}%
\pgfpathlineto{\pgfqpoint{9.538157in}{4.646549in}}%
\pgfpathlineto{\pgfqpoint{9.538157in}{4.562067in}}%
\pgfpathclose%
\pgfusepath{stroke}%
\end{pgfscope}%
\begin{pgfscope}%
\pgfpathrectangle{\pgfqpoint{7.904091in}{2.928000in}}{\pgfqpoint{5.664000in}{5.664000in}}%
\pgfusepath{clip}%
\pgfsetbuttcap%
\pgfsetroundjoin%
\pgfsetlinewidth{1.003750pt}%
\definecolor{currentstroke}{rgb}{0.541176,0.380392,0.658824}%
\pgfsetstrokecolor{currentstroke}%
\pgfsetstrokeopacity{0.900000}%
\pgfsetdash{}{0pt}%
\pgfpathmoveto{\pgfqpoint{10.631772in}{5.655681in}}%
\pgfpathlineto{\pgfqpoint{10.716254in}{5.655681in}}%
\pgfpathlineto{\pgfqpoint{10.716254in}{5.740163in}}%
\pgfpathlineto{\pgfqpoint{10.631772in}{5.740163in}}%
\pgfpathlineto{\pgfqpoint{10.631772in}{5.655681in}}%
\pgfpathclose%
\pgfusepath{stroke}%
\end{pgfscope}%
\begin{pgfscope}%
\pgfpathrectangle{\pgfqpoint{7.904091in}{2.928000in}}{\pgfqpoint{5.664000in}{5.664000in}}%
\pgfusepath{clip}%
\pgfsetbuttcap%
\pgfsetroundjoin%
\pgfsetlinewidth{1.003750pt}%
\definecolor{currentstroke}{rgb}{0.541176,0.380392,0.658824}%
\pgfsetstrokecolor{currentstroke}%
\pgfsetstrokeopacity{0.900000}%
\pgfsetdash{}{0pt}%
\pgfpathmoveto{\pgfqpoint{10.738528in}{5.762438in}}%
\pgfpathlineto{\pgfqpoint{10.823011in}{5.762438in}}%
\pgfpathlineto{\pgfqpoint{10.823011in}{5.846920in}}%
\pgfpathlineto{\pgfqpoint{10.738528in}{5.846920in}}%
\pgfpathlineto{\pgfqpoint{10.738528in}{5.762438in}}%
\pgfpathclose%
\pgfusepath{stroke}%
\end{pgfscope}%
\begin{pgfscope}%
\pgfpathrectangle{\pgfqpoint{7.904091in}{2.928000in}}{\pgfqpoint{5.664000in}{5.664000in}}%
\pgfusepath{clip}%
\pgfsetbuttcap%
\pgfsetroundjoin%
\definecolor{currentfill}{rgb}{0.768627,0.603922,0.000000}%
\pgfsetfillcolor{currentfill}%
\pgfsetfillopacity{0.900000}%
\pgfsetlinewidth{0.803000pt}%
\definecolor{currentstroke}{rgb}{0.768627,0.603922,0.000000}%
\pgfsetstrokecolor{currentstroke}%
\pgfsetstrokeopacity{0.900000}%
\pgfsetdash{}{0pt}%
\pgfsys@defobject{currentmarker}{\pgfqpoint{-0.042241in}{-0.042241in}}{\pgfqpoint{0.042241in}{0.042241in}}{%
\pgfpathmoveto{\pgfqpoint{-0.042241in}{-0.042241in}}%
\pgfpathlineto{\pgfqpoint{0.042241in}{-0.042241in}}%
\pgfpathlineto{\pgfqpoint{0.042241in}{0.042241in}}%
\pgfpathlineto{\pgfqpoint{-0.042241in}{0.042241in}}%
\pgfpathlineto{\pgfqpoint{-0.042241in}{-0.042241in}}%
\pgfpathclose%
\pgfusepath{stroke,fill}%
}%
\begin{pgfscope}%
\pgfsys@transformshift{9.380097in}{4.824992in}%
\pgfsys@useobject{currentmarker}{}%
\end{pgfscope}%
\end{pgfscope}%
\begin{pgfscope}%
\pgfpathrectangle{\pgfqpoint{7.904091in}{2.928000in}}{\pgfqpoint{5.664000in}{5.664000in}}%
\pgfusepath{clip}%
\pgfsetbuttcap%
\pgfsetroundjoin%
\pgfsetlinewidth{1.003750pt}%
\definecolor{currentstroke}{rgb}{0.768627,0.603922,0.000000}%
\pgfsetstrokecolor{currentstroke}%
\pgfsetstrokeopacity{0.900000}%
\pgfsetdash{}{0pt}%
\pgfpathmoveto{\pgfqpoint{12.255775in}{7.360320in}}%
\pgfpathlineto{\pgfqpoint{12.340258in}{7.360320in}}%
\pgfpathlineto{\pgfqpoint{12.340258in}{7.444803in}}%
\pgfpathlineto{\pgfqpoint{12.255775in}{7.444803in}}%
\pgfpathlineto{\pgfqpoint{12.255775in}{7.360320in}}%
\pgfpathclose%
\pgfusepath{stroke}%
\end{pgfscope}%
\begin{pgfscope}%
\pgfpathrectangle{\pgfqpoint{7.904091in}{2.928000in}}{\pgfqpoint{5.664000in}{5.664000in}}%
\pgfusepath{clip}%
\pgfsetbuttcap%
\pgfsetroundjoin%
\pgfsetlinewidth{1.003750pt}%
\definecolor{currentstroke}{rgb}{0.768627,0.603922,0.000000}%
\pgfsetstrokecolor{currentstroke}%
\pgfsetstrokeopacity{0.900000}%
\pgfsetdash{}{0pt}%
\pgfpathmoveto{\pgfqpoint{12.403981in}{7.471891in}}%
\pgfpathlineto{\pgfqpoint{12.488464in}{7.471891in}}%
\pgfpathlineto{\pgfqpoint{12.488464in}{7.556374in}}%
\pgfpathlineto{\pgfqpoint{12.403981in}{7.556374in}}%
\pgfpathlineto{\pgfqpoint{12.403981in}{7.471891in}}%
\pgfpathclose%
\pgfusepath{stroke}%
\end{pgfscope}%
\begin{pgfscope}%
\pgfpathrectangle{\pgfqpoint{7.904091in}{2.928000in}}{\pgfqpoint{5.664000in}{5.664000in}}%
\pgfusepath{clip}%
\pgfsetbuttcap%
\pgfsetroundjoin%
\pgfsetlinewidth{1.003750pt}%
\definecolor{currentstroke}{rgb}{0.768627,0.603922,0.000000}%
\pgfsetstrokecolor{currentstroke}%
\pgfsetstrokeopacity{0.900000}%
\pgfsetdash{}{0pt}%
\pgfpathmoveto{\pgfqpoint{10.422089in}{5.606192in}}%
\pgfpathlineto{\pgfqpoint{10.506572in}{5.606192in}}%
\pgfpathlineto{\pgfqpoint{10.506572in}{5.690675in}}%
\pgfpathlineto{\pgfqpoint{10.422089in}{5.690675in}}%
\pgfpathlineto{\pgfqpoint{10.422089in}{5.606192in}}%
\pgfpathclose%
\pgfusepath{stroke}%
\end{pgfscope}%
\begin{pgfscope}%
\pgfpathrectangle{\pgfqpoint{7.904091in}{2.928000in}}{\pgfqpoint{5.664000in}{5.664000in}}%
\pgfusepath{clip}%
\pgfsetbuttcap%
\pgfsetroundjoin%
\pgfsetlinewidth{1.003750pt}%
\definecolor{currentstroke}{rgb}{0.200000,0.133333,0.533333}%
\pgfsetstrokecolor{currentstroke}%
\pgfsetstrokeopacity{0.900000}%
\pgfsetdash{}{0pt}%
\pgfpathmoveto{\pgfqpoint{9.471329in}{3.950912in}}%
\pgfpathlineto{\pgfqpoint{9.555812in}{3.950912in}}%
\pgfpathlineto{\pgfqpoint{9.555812in}{4.035394in}}%
\pgfpathlineto{\pgfqpoint{9.471329in}{4.035394in}}%
\pgfpathlineto{\pgfqpoint{9.471329in}{3.950912in}}%
\pgfpathclose%
\pgfusepath{stroke}%
\end{pgfscope}%
\begin{pgfscope}%
\pgfpathrectangle{\pgfqpoint{7.904091in}{2.928000in}}{\pgfqpoint{5.664000in}{5.664000in}}%
\pgfusepath{clip}%
\pgfsetbuttcap%
\pgfsetroundjoin%
\pgfsetlinewidth{1.003750pt}%
\definecolor{currentstroke}{rgb}{0.200000,0.133333,0.533333}%
\pgfsetstrokecolor{currentstroke}%
\pgfsetstrokeopacity{0.900000}%
\pgfsetdash{}{0pt}%
\pgfpathmoveto{\pgfqpoint{9.511877in}{4.044083in}}%
\pgfpathlineto{\pgfqpoint{9.596360in}{4.044083in}}%
\pgfpathlineto{\pgfqpoint{9.596360in}{4.128566in}}%
\pgfpathlineto{\pgfqpoint{9.511877in}{4.128566in}}%
\pgfpathlineto{\pgfqpoint{9.511877in}{4.044083in}}%
\pgfpathclose%
\pgfusepath{stroke}%
\end{pgfscope}%
\begin{pgfscope}%
\pgfpathrectangle{\pgfqpoint{7.904091in}{2.928000in}}{\pgfqpoint{5.664000in}{5.664000in}}%
\pgfusepath{clip}%
\pgfsetbuttcap%
\pgfsetroundjoin%
\pgfsetlinewidth{1.003750pt}%
\definecolor{currentstroke}{rgb}{0.200000,0.133333,0.533333}%
\pgfsetstrokecolor{currentstroke}%
\pgfsetstrokeopacity{0.900000}%
\pgfsetdash{}{0pt}%
\pgfpathmoveto{\pgfqpoint{9.938697in}{4.722009in}}%
\pgfpathlineto{\pgfqpoint{10.023180in}{4.722009in}}%
\pgfpathlineto{\pgfqpoint{10.023180in}{4.806492in}}%
\pgfpathlineto{\pgfqpoint{9.938697in}{4.806492in}}%
\pgfpathlineto{\pgfqpoint{9.938697in}{4.722009in}}%
\pgfpathclose%
\pgfusepath{stroke}%
\end{pgfscope}%
\begin{pgfscope}%
\pgfpathrectangle{\pgfqpoint{7.904091in}{2.928000in}}{\pgfqpoint{5.664000in}{5.664000in}}%
\pgfusepath{clip}%
\pgfsetbuttcap%
\pgfsetroundjoin%
\pgfsetlinewidth{1.003750pt}%
\definecolor{currentstroke}{rgb}{0.200000,0.133333,0.533333}%
\pgfsetstrokecolor{currentstroke}%
\pgfsetstrokeopacity{0.900000}%
\pgfsetdash{}{0pt}%
\pgfpathmoveto{\pgfqpoint{11.348076in}{6.371985in}}%
\pgfpathlineto{\pgfqpoint{11.432558in}{6.371985in}}%
\pgfpathlineto{\pgfqpoint{11.432558in}{6.456468in}}%
\pgfpathlineto{\pgfqpoint{11.348076in}{6.456468in}}%
\pgfpathlineto{\pgfqpoint{11.348076in}{6.371985in}}%
\pgfpathclose%
\pgfusepath{stroke}%
\end{pgfscope}%
\begin{pgfscope}%
\pgfpathrectangle{\pgfqpoint{7.904091in}{2.928000in}}{\pgfqpoint{5.664000in}{5.664000in}}%
\pgfusepath{clip}%
\pgfsetbuttcap%
\pgfsetroundjoin%
\pgfsetlinewidth{1.003750pt}%
\definecolor{currentstroke}{rgb}{0.200000,0.133333,0.533333}%
\pgfsetstrokecolor{currentstroke}%
\pgfsetstrokeopacity{0.900000}%
\pgfsetdash{}{0pt}%
\pgfpathmoveto{\pgfqpoint{11.019575in}{6.043485in}}%
\pgfpathlineto{\pgfqpoint{11.104058in}{6.043485in}}%
\pgfpathlineto{\pgfqpoint{11.104058in}{6.127967in}}%
\pgfpathlineto{\pgfqpoint{11.019575in}{6.127967in}}%
\pgfpathlineto{\pgfqpoint{11.019575in}{6.043485in}}%
\pgfpathclose%
\pgfusepath{stroke}%
\end{pgfscope}%
\begin{pgfscope}%
\pgfpathrectangle{\pgfqpoint{7.904091in}{2.928000in}}{\pgfqpoint{5.664000in}{5.664000in}}%
\pgfusepath{clip}%
\pgfsetbuttcap%
\pgfsetroundjoin%
\pgfsetlinewidth{1.003750pt}%
\definecolor{currentstroke}{rgb}{0.200000,0.133333,0.533333}%
\pgfsetstrokecolor{currentstroke}%
\pgfsetstrokeopacity{0.900000}%
\pgfsetdash{}{0pt}%
\pgfpathmoveto{\pgfqpoint{10.625596in}{5.649505in}}%
\pgfpathlineto{\pgfqpoint{10.710079in}{5.649505in}}%
\pgfpathlineto{\pgfqpoint{10.710079in}{5.733988in}}%
\pgfpathlineto{\pgfqpoint{10.625596in}{5.733988in}}%
\pgfpathlineto{\pgfqpoint{10.625596in}{5.649505in}}%
\pgfpathclose%
\pgfusepath{stroke}%
\end{pgfscope}%
\begin{pgfscope}%
\pgfpathrectangle{\pgfqpoint{7.904091in}{2.928000in}}{\pgfqpoint{5.664000in}{5.664000in}}%
\pgfusepath{clip}%
\pgfsetbuttcap%
\pgfsetroundjoin%
\pgfsetlinewidth{1.003750pt}%
\definecolor{currentstroke}{rgb}{0.176471,0.713725,0.643137}%
\pgfsetstrokecolor{currentstroke}%
\pgfsetstrokeopacity{0.900000}%
\pgfsetdash{}{0pt}%
\pgfpathmoveto{\pgfqpoint{12.134820in}{7.256770in}}%
\pgfpathlineto{\pgfqpoint{12.219303in}{7.256770in}}%
\pgfpathlineto{\pgfqpoint{12.219303in}{7.341253in}}%
\pgfpathlineto{\pgfqpoint{12.134820in}{7.341253in}}%
\pgfpathlineto{\pgfqpoint{12.134820in}{7.256770in}}%
\pgfpathclose%
\pgfusepath{stroke}%
\end{pgfscope}%
\begin{pgfscope}%
\pgfpathrectangle{\pgfqpoint{7.904091in}{2.928000in}}{\pgfqpoint{5.664000in}{5.664000in}}%
\pgfusepath{clip}%
\pgfsetbuttcap%
\pgfsetroundjoin%
\pgfsetlinewidth{1.003750pt}%
\definecolor{currentstroke}{rgb}{0.176471,0.713725,0.643137}%
\pgfsetstrokecolor{currentstroke}%
\pgfsetstrokeopacity{0.900000}%
\pgfsetdash{}{0pt}%
\pgfpathmoveto{\pgfqpoint{12.402625in}{7.916725in}}%
\pgfpathlineto{\pgfqpoint{12.487108in}{7.916725in}}%
\pgfpathlineto{\pgfqpoint{12.487108in}{8.001208in}}%
\pgfpathlineto{\pgfqpoint{12.402625in}{8.001208in}}%
\pgfpathlineto{\pgfqpoint{12.402625in}{7.916725in}}%
\pgfpathclose%
\pgfusepath{stroke}%
\end{pgfscope}%
\begin{pgfscope}%
\pgfpathrectangle{\pgfqpoint{7.904091in}{2.928000in}}{\pgfqpoint{5.664000in}{5.664000in}}%
\pgfusepath{clip}%
\pgfsetbuttcap%
\pgfsetroundjoin%
\pgfsetlinewidth{1.003750pt}%
\definecolor{currentstroke}{rgb}{0.176471,0.713725,0.643137}%
\pgfsetstrokecolor{currentstroke}%
\pgfsetstrokeopacity{0.900000}%
\pgfsetdash{}{0pt}%
\pgfpathmoveto{\pgfqpoint{11.929947in}{7.080785in}}%
\pgfpathlineto{\pgfqpoint{12.014430in}{7.080785in}}%
\pgfpathlineto{\pgfqpoint{12.014430in}{7.165268in}}%
\pgfpathlineto{\pgfqpoint{11.929947in}{7.165268in}}%
\pgfpathlineto{\pgfqpoint{11.929947in}{7.080785in}}%
\pgfpathclose%
\pgfusepath{stroke}%
\end{pgfscope}%
\begin{pgfscope}%
\pgfpathrectangle{\pgfqpoint{7.904091in}{2.928000in}}{\pgfqpoint{5.664000in}{5.664000in}}%
\pgfusepath{clip}%
\pgfsetbuttcap%
\pgfsetroundjoin%
\pgfsetlinewidth{1.003750pt}%
\definecolor{currentstroke}{rgb}{0.176471,0.713725,0.643137}%
\pgfsetstrokecolor{currentstroke}%
\pgfsetstrokeopacity{0.900000}%
\pgfsetdash{}{0pt}%
\pgfpathmoveto{\pgfqpoint{12.503619in}{7.557277in}}%
\pgfpathlineto{\pgfqpoint{12.588102in}{7.557277in}}%
\pgfpathlineto{\pgfqpoint{12.588102in}{7.641760in}}%
\pgfpathlineto{\pgfqpoint{12.503619in}{7.641760in}}%
\pgfpathlineto{\pgfqpoint{12.503619in}{7.557277in}}%
\pgfpathclose%
\pgfusepath{stroke}%
\end{pgfscope}%
\begin{pgfscope}%
\pgfpathrectangle{\pgfqpoint{7.904091in}{2.928000in}}{\pgfqpoint{5.664000in}{5.664000in}}%
\pgfusepath{clip}%
\pgfsetbuttcap%
\pgfsetroundjoin%
\pgfsetlinewidth{1.003750pt}%
\definecolor{currentstroke}{rgb}{0.176471,0.713725,0.643137}%
\pgfsetstrokecolor{currentstroke}%
\pgfsetstrokeopacity{0.900000}%
\pgfsetdash{}{0pt}%
\pgfpathmoveto{\pgfqpoint{10.084933in}{5.353212in}}%
\pgfpathlineto{\pgfqpoint{10.169416in}{5.353212in}}%
\pgfpathlineto{\pgfqpoint{10.169416in}{5.437694in}}%
\pgfpathlineto{\pgfqpoint{10.084933in}{5.437694in}}%
\pgfpathlineto{\pgfqpoint{10.084933in}{5.353212in}}%
\pgfpathclose%
\pgfusepath{stroke}%
\end{pgfscope}%
\begin{pgfscope}%
\pgfpathrectangle{\pgfqpoint{7.904091in}{2.928000in}}{\pgfqpoint{5.664000in}{5.664000in}}%
\pgfusepath{clip}%
\pgfsetbuttcap%
\pgfsetroundjoin%
\pgfsetlinewidth{1.003750pt}%
\definecolor{currentstroke}{rgb}{0.901961,0.196078,0.156863}%
\pgfsetstrokecolor{currentstroke}%
\pgfsetstrokeopacity{0.900000}%
\pgfsetdash{}{0pt}%
\pgfpathmoveto{\pgfqpoint{12.669591in}{7.718036in}}%
\pgfpathlineto{\pgfqpoint{12.754074in}{7.718036in}}%
\pgfpathlineto{\pgfqpoint{12.754074in}{7.802518in}}%
\pgfpathlineto{\pgfqpoint{12.669591in}{7.802518in}}%
\pgfpathlineto{\pgfqpoint{12.669591in}{7.718036in}}%
\pgfpathclose%
\pgfusepath{stroke}%
\end{pgfscope}%
\begin{pgfscope}%
\pgfpathrectangle{\pgfqpoint{7.904091in}{2.928000in}}{\pgfqpoint{5.664000in}{5.664000in}}%
\pgfusepath{clip}%
\pgfsetbuttcap%
\pgfsetroundjoin%
\pgfsetlinewidth{1.003750pt}%
\definecolor{currentstroke}{rgb}{0.901961,0.196078,0.156863}%
\pgfsetstrokecolor{currentstroke}%
\pgfsetstrokeopacity{0.900000}%
\pgfsetdash{}{0pt}%
\pgfpathmoveto{\pgfqpoint{11.752922in}{6.776831in}}%
\pgfpathlineto{\pgfqpoint{11.837405in}{6.776831in}}%
\pgfpathlineto{\pgfqpoint{11.837405in}{6.861314in}}%
\pgfpathlineto{\pgfqpoint{11.752922in}{6.861314in}}%
\pgfpathlineto{\pgfqpoint{11.752922in}{6.776831in}}%
\pgfpathclose%
\pgfusepath{stroke}%
\end{pgfscope}%
\begin{pgfscope}%
\pgfpathrectangle{\pgfqpoint{7.904091in}{2.928000in}}{\pgfqpoint{5.664000in}{5.664000in}}%
\pgfusepath{clip}%
\pgfsetbuttcap%
\pgfsetroundjoin%
\pgfsetlinewidth{1.003750pt}%
\definecolor{currentstroke}{rgb}{0.901961,0.196078,0.156863}%
\pgfsetstrokecolor{currentstroke}%
\pgfsetstrokeopacity{0.900000}%
\pgfsetdash{}{0pt}%
\pgfpathmoveto{\pgfqpoint{10.575373in}{5.599282in}}%
\pgfpathlineto{\pgfqpoint{10.659856in}{5.599282in}}%
\pgfpathlineto{\pgfqpoint{10.659856in}{5.683765in}}%
\pgfpathlineto{\pgfqpoint{10.575373in}{5.683765in}}%
\pgfpathlineto{\pgfqpoint{10.575373in}{5.599282in}}%
\pgfpathclose%
\pgfusepath{stroke}%
\end{pgfscope}%
\begin{pgfscope}%
\pgfpathrectangle{\pgfqpoint{7.904091in}{2.928000in}}{\pgfqpoint{5.664000in}{5.664000in}}%
\pgfusepath{clip}%
\pgfsetbuttcap%
\pgfsetroundjoin%
\pgfsetlinewidth{1.003750pt}%
\definecolor{currentstroke}{rgb}{0.498039,0.560784,0.678431}%
\pgfsetstrokecolor{currentstroke}%
\pgfsetstrokeopacity{0.900000}%
\pgfsetdash{}{0pt}%
\pgfpathmoveto{\pgfqpoint{11.285555in}{6.309465in}}%
\pgfpathlineto{\pgfqpoint{11.370038in}{6.309465in}}%
\pgfpathlineto{\pgfqpoint{11.370038in}{6.393947in}}%
\pgfpathlineto{\pgfqpoint{11.285555in}{6.393947in}}%
\pgfpathlineto{\pgfqpoint{11.285555in}{6.309465in}}%
\pgfpathclose%
\pgfusepath{stroke}%
\end{pgfscope}%
\begin{pgfscope}%
\pgfsetbuttcap%
\pgfsetroundjoin%
\pgfsetlinewidth{1.003750pt}%
\definecolor{currentstroke}{rgb}{0.498039,0.560784,0.678431}%
\pgfsetstrokecolor{currentstroke}%
\pgfsetstrokeopacity{0.900000}%
\pgfsetdash{}{0pt}%
\pgfpathmoveto{\pgfqpoint{7.861850in}{2.885759in}}%
\pgfpathlineto{\pgfqpoint{7.946332in}{2.885759in}}%
\pgfpathlineto{\pgfqpoint{7.946332in}{2.970241in}}%
\pgfpathlineto{\pgfqpoint{7.861850in}{2.970241in}}%
\pgfpathlineto{\pgfqpoint{7.861850in}{2.885759in}}%
\pgfpathclose%
\pgfusepath{stroke}%
\end{pgfscope}%
\begin{pgfscope}%
\pgfsetbuttcap%
\pgfsetroundjoin%
\definecolor{currentfill}{rgb}{0.709804,0.121569,0.270588}%
\pgfsetfillcolor{currentfill}%
\pgfsetfillopacity{0.900000}%
\pgfsetlinewidth{0.803000pt}%
\definecolor{currentstroke}{rgb}{0.709804,0.121569,0.270588}%
\pgfsetstrokecolor{currentstroke}%
\pgfsetstrokeopacity{0.900000}%
\pgfsetdash{}{0pt}%
\pgfsys@defobject{currentmarker}{\pgfqpoint{-0.042241in}{-0.042241in}}{\pgfqpoint{0.042241in}{0.042241in}}{%
\pgfpathmoveto{\pgfqpoint{-0.042241in}{-0.042241in}}%
\pgfpathlineto{\pgfqpoint{0.042241in}{-0.042241in}}%
\pgfpathlineto{\pgfqpoint{0.042241in}{0.042241in}}%
\pgfpathlineto{\pgfqpoint{-0.042241in}{0.042241in}}%
\pgfpathlineto{\pgfqpoint{-0.042241in}{-0.042241in}}%
\pgfpathclose%
\pgfusepath{stroke,fill}%
}%
\begin{pgfscope}%
\pgfsys@transformshift{13.313868in}{8.592000in}%
\pgfsys@useobject{currentmarker}{}%
\end{pgfscope}%
\end{pgfscope}%
\begin{pgfscope}%
\pgfsetbuttcap%
\pgfsetmiterjoin%
\definecolor{currentfill}{rgb}{1.000000,1.000000,1.000000}%
\pgfsetfillcolor{currentfill}%
\pgfsetfillopacity{0.970000}%
\pgfsetlinewidth{0.702625pt}%
\definecolor{currentstroke}{rgb}{0.796078,0.827451,0.850980}%
\pgfsetstrokecolor{currentstroke}%
\pgfsetstrokeopacity{0.970000}%
\pgfsetdash{}{0pt}%
\pgfpathmoveto{\pgfqpoint{8.102331in}{6.665893in}}%
\pgfpathlineto{\pgfqpoint{10.219206in}{6.665893in}}%
\pgfpathquadraticcurveto{\pgfqpoint{10.297123in}{6.665893in}}{\pgfqpoint{10.297123in}{6.743810in}}%
\pgfpathlineto{\pgfqpoint{10.297123in}{7.402560in}}%
\pgfpathquadraticcurveto{\pgfqpoint{10.297123in}{7.480477in}}{\pgfqpoint{10.219206in}{7.480477in}}%
\pgfpathlineto{\pgfqpoint{8.102331in}{7.480477in}}%
\pgfpathquadraticcurveto{\pgfqpoint{8.024414in}{7.480477in}}{\pgfqpoint{8.024414in}{7.402560in}}%
\pgfpathlineto{\pgfqpoint{8.024414in}{6.743810in}}%
\pgfpathquadraticcurveto{\pgfqpoint{8.024414in}{6.665893in}}{\pgfqpoint{8.102331in}{6.665893in}}%
\pgfpathlineto{\pgfqpoint{8.102331in}{6.665893in}}%
\pgfpathclose%
\pgfusepath{stroke,fill}%
\end{pgfscope}%
\begin{pgfscope}%
\pgfsetbuttcap%
\pgfsetroundjoin%
\definecolor{currentfill}{rgb}{0.090196,0.168627,0.227451}%
\pgfsetfillcolor{currentfill}%
\pgfsetlinewidth{0.000000pt}%
\definecolor{currentstroke}{rgb}{0.090196,0.168627,0.227451}%
\pgfsetstrokecolor{currentstroke}%
\pgfsetdash{}{0pt}%
\pgfpathmoveto{\pgfqpoint{8.178145in}{7.352665in}}%
\pgfpathlineto{\pgfqpoint{8.178145in}{7.339030in}}%
\pgfpathquadraticcurveto{\pgfqpoint{8.170198in}{7.342837in}}{\pgfqpoint{8.163137in}{7.344697in}}%
\pgfpathquadraticcurveto{\pgfqpoint{8.156076in}{7.346578in}}{\pgfqpoint{8.149501in}{7.346578in}}%
\pgfpathquadraticcurveto{\pgfqpoint{8.138102in}{7.346578in}}{\pgfqpoint{8.131904in}{7.342151in}}%
\pgfpathquadraticcurveto{\pgfqpoint{8.125706in}{7.337724in}}{\pgfqpoint{8.125706in}{7.329556in}}%
\pgfpathquadraticcurveto{\pgfqpoint{8.125706in}{7.322694in}}{\pgfqpoint{8.129823in}{7.319197in}}%
\pgfpathquadraticcurveto{\pgfqpoint{8.133940in}{7.315721in}}{\pgfqpoint{8.145429in}{7.313574in}}%
\pgfpathlineto{\pgfqpoint{8.153862in}{7.311848in}}%
\pgfpathquadraticcurveto{\pgfqpoint{8.169490in}{7.308859in}}{\pgfqpoint{8.176927in}{7.301356in}}%
\pgfpathquadraticcurveto{\pgfqpoint{8.184365in}{7.293852in}}{\pgfqpoint{8.184365in}{7.281279in}}%
\pgfpathquadraticcurveto{\pgfqpoint{8.184365in}{7.266249in}}{\pgfqpoint{8.174293in}{7.258501in}}%
\pgfpathquadraticcurveto{\pgfqpoint{8.164244in}{7.250754in}}{\pgfqpoint{8.144809in}{7.250754in}}%
\pgfpathquadraticcurveto{\pgfqpoint{8.137482in}{7.250754in}}{\pgfqpoint{8.129203in}{7.252414in}}%
\pgfpathquadraticcurveto{\pgfqpoint{8.120947in}{7.254074in}}{\pgfqpoint{8.112093in}{7.257328in}}%
\pgfpathlineto{\pgfqpoint{8.112093in}{7.271716in}}%
\pgfpathquadraticcurveto{\pgfqpoint{8.120593in}{7.266957in}}{\pgfqpoint{8.128761in}{7.264522in}}%
\pgfpathquadraticcurveto{\pgfqpoint{8.136929in}{7.262109in}}{\pgfqpoint{8.144809in}{7.262109in}}%
\pgfpathquadraticcurveto{\pgfqpoint{8.156762in}{7.262109in}}{\pgfqpoint{8.163270in}{7.266802in}}%
\pgfpathquadraticcurveto{\pgfqpoint{8.169778in}{7.271517in}}{\pgfqpoint{8.169778in}{7.280238in}}%
\pgfpathquadraticcurveto{\pgfqpoint{8.169778in}{7.287831in}}{\pgfqpoint{8.165107in}{7.292125in}}%
\pgfpathquadraticcurveto{\pgfqpoint{8.160436in}{7.296419in}}{\pgfqpoint{8.149789in}{7.298567in}}%
\pgfpathlineto{\pgfqpoint{8.141267in}{7.300227in}}%
\pgfpathquadraticcurveto{\pgfqpoint{8.125640in}{7.303326in}}{\pgfqpoint{8.118645in}{7.309966in}}%
\pgfpathquadraticcurveto{\pgfqpoint{8.111672in}{7.316607in}}{\pgfqpoint{8.111672in}{7.328449in}}%
\pgfpathquadraticcurveto{\pgfqpoint{8.111672in}{7.342151in}}{\pgfqpoint{8.121323in}{7.350031in}}%
\pgfpathquadraticcurveto{\pgfqpoint{8.130974in}{7.357912in}}{\pgfqpoint{8.147908in}{7.357912in}}%
\pgfpathquadraticcurveto{\pgfqpoint{8.155190in}{7.357912in}}{\pgfqpoint{8.162716in}{7.356583in}}%
\pgfpathquadraticcurveto{\pgfqpoint{8.170265in}{7.355277in}}{\pgfqpoint{8.178145in}{7.352665in}}%
\pgfpathlineto{\pgfqpoint{8.178145in}{7.352665in}}%
\pgfpathclose%
\pgfpathmoveto{\pgfqpoint{8.205615in}{7.330242in}}%
\pgfpathlineto{\pgfqpoint{8.218343in}{7.330242in}}%
\pgfpathlineto{\pgfqpoint{8.218343in}{7.252768in}}%
\pgfpathlineto{\pgfqpoint{8.205615in}{7.252768in}}%
\pgfpathlineto{\pgfqpoint{8.205615in}{7.330242in}}%
\pgfpathlineto{\pgfqpoint{8.205615in}{7.330242in}}%
\pgfpathclose%
\pgfpathmoveto{\pgfqpoint{8.205615in}{7.360413in}}%
\pgfpathlineto{\pgfqpoint{8.218343in}{7.360413in}}%
\pgfpathlineto{\pgfqpoint{8.218343in}{7.344276in}}%
\pgfpathlineto{\pgfqpoint{8.205615in}{7.344276in}}%
\pgfpathlineto{\pgfqpoint{8.205615in}{7.360413in}}%
\pgfpathlineto{\pgfqpoint{8.205615in}{7.360413in}}%
\pgfpathclose%
\pgfpathmoveto{\pgfqpoint{8.295949in}{7.292413in}}%
\pgfpathquadraticcurveto{\pgfqpoint{8.295949in}{7.306248in}}{\pgfqpoint{8.290238in}{7.313840in}}%
\pgfpathquadraticcurveto{\pgfqpoint{8.284550in}{7.321455in}}{\pgfqpoint{8.274235in}{7.321455in}}%
\pgfpathquadraticcurveto{\pgfqpoint{8.264008in}{7.321455in}}{\pgfqpoint{8.258297in}{7.313840in}}%
\pgfpathquadraticcurveto{\pgfqpoint{8.252586in}{7.306248in}}{\pgfqpoint{8.252586in}{7.292413in}}%
\pgfpathquadraticcurveto{\pgfqpoint{8.252586in}{7.278645in}}{\pgfqpoint{8.258297in}{7.271030in}}%
\pgfpathquadraticcurveto{\pgfqpoint{8.264008in}{7.263415in}}{\pgfqpoint{8.274235in}{7.263415in}}%
\pgfpathquadraticcurveto{\pgfqpoint{8.284550in}{7.263415in}}{\pgfqpoint{8.290238in}{7.271030in}}%
\pgfpathquadraticcurveto{\pgfqpoint{8.295949in}{7.278645in}}{\pgfqpoint{8.295949in}{7.292413in}}%
\pgfpathlineto{\pgfqpoint{8.295949in}{7.292413in}}%
\pgfpathclose%
\pgfpathmoveto{\pgfqpoint{8.308677in}{7.262375in}}%
\pgfpathquadraticcurveto{\pgfqpoint{8.308677in}{7.242608in}}{\pgfqpoint{8.299890in}{7.232957in}}%
\pgfpathquadraticcurveto{\pgfqpoint{8.291124in}{7.223306in}}{\pgfqpoint{8.272995in}{7.223306in}}%
\pgfpathquadraticcurveto{\pgfqpoint{8.266288in}{7.223306in}}{\pgfqpoint{8.260334in}{7.224302in}}%
\pgfpathquadraticcurveto{\pgfqpoint{8.254379in}{7.225298in}}{\pgfqpoint{8.248779in}{7.227379in}}%
\pgfpathlineto{\pgfqpoint{8.248779in}{7.239753in}}%
\pgfpathquadraticcurveto{\pgfqpoint{8.254379in}{7.236720in}}{\pgfqpoint{8.259847in}{7.235281in}}%
\pgfpathquadraticcurveto{\pgfqpoint{8.265314in}{7.233820in}}{\pgfqpoint{8.270981in}{7.233820in}}%
\pgfpathquadraticcurveto{\pgfqpoint{8.283509in}{7.233820in}}{\pgfqpoint{8.289729in}{7.240350in}}%
\pgfpathquadraticcurveto{\pgfqpoint{8.295949in}{7.246880in}}{\pgfqpoint{8.295949in}{7.260095in}}%
\pgfpathlineto{\pgfqpoint{8.295949in}{7.266404in}}%
\pgfpathquadraticcurveto{\pgfqpoint{8.292009in}{7.259542in}}{\pgfqpoint{8.285856in}{7.256155in}}%
\pgfpathquadraticcurveto{\pgfqpoint{8.279702in}{7.252768in}}{\pgfqpoint{8.271113in}{7.252768in}}%
\pgfpathquadraticcurveto{\pgfqpoint{8.256880in}{7.252768in}}{\pgfqpoint{8.248159in}{7.263615in}}%
\pgfpathquadraticcurveto{\pgfqpoint{8.239438in}{7.274483in}}{\pgfqpoint{8.239438in}{7.292413in}}%
\pgfpathquadraticcurveto{\pgfqpoint{8.239438in}{7.310387in}}{\pgfqpoint{8.248159in}{7.321233in}}%
\pgfpathquadraticcurveto{\pgfqpoint{8.256880in}{7.332102in}}{\pgfqpoint{8.271113in}{7.332102in}}%
\pgfpathquadraticcurveto{\pgfqpoint{8.279702in}{7.332102in}}{\pgfqpoint{8.285856in}{7.328715in}}%
\pgfpathquadraticcurveto{\pgfqpoint{8.292009in}{7.325328in}}{\pgfqpoint{8.295949in}{7.318488in}}%
\pgfpathlineto{\pgfqpoint{8.295949in}{7.330242in}}%
\pgfpathlineto{\pgfqpoint{8.308677in}{7.330242in}}%
\pgfpathlineto{\pgfqpoint{8.308677in}{7.262375in}}%
\pgfpathlineto{\pgfqpoint{8.308677in}{7.262375in}}%
\pgfpathclose%
\pgfpathmoveto{\pgfqpoint{8.399322in}{7.299540in}}%
\pgfpathlineto{\pgfqpoint{8.399322in}{7.252768in}}%
\pgfpathlineto{\pgfqpoint{8.386594in}{7.252768in}}%
\pgfpathlineto{\pgfqpoint{8.386594in}{7.299120in}}%
\pgfpathquadraticcurveto{\pgfqpoint{8.386594in}{7.310121in}}{\pgfqpoint{8.382300in}{7.315567in}}%
\pgfpathquadraticcurveto{\pgfqpoint{8.378005in}{7.321034in}}{\pgfqpoint{8.369439in}{7.321034in}}%
\pgfpathquadraticcurveto{\pgfqpoint{8.359124in}{7.321034in}}{\pgfqpoint{8.353169in}{7.314460in}}%
\pgfpathquadraticcurveto{\pgfqpoint{8.347215in}{7.307908in}}{\pgfqpoint{8.347215in}{7.296552in}}%
\pgfpathlineto{\pgfqpoint{8.347215in}{7.252768in}}%
\pgfpathlineto{\pgfqpoint{8.334421in}{7.252768in}}%
\pgfpathlineto{\pgfqpoint{8.334421in}{7.330242in}}%
\pgfpathlineto{\pgfqpoint{8.347215in}{7.330242in}}%
\pgfpathlineto{\pgfqpoint{8.347215in}{7.318201in}}%
\pgfpathquadraticcurveto{\pgfqpoint{8.351797in}{7.325195in}}{\pgfqpoint{8.357973in}{7.328649in}}%
\pgfpathquadraticcurveto{\pgfqpoint{8.364171in}{7.332102in}}{\pgfqpoint{8.372272in}{7.332102in}}%
\pgfpathquadraticcurveto{\pgfqpoint{8.385620in}{7.332102in}}{\pgfqpoint{8.392460in}{7.323845in}}%
\pgfpathquadraticcurveto{\pgfqpoint{8.399322in}{7.315589in}}{\pgfqpoint{8.399322in}{7.299540in}}%
\pgfpathlineto{\pgfqpoint{8.399322in}{7.299540in}}%
\pgfpathclose%
\pgfpathmoveto{\pgfqpoint{8.418270in}{7.297238in}}%
\pgfpathlineto{\pgfqpoint{8.455546in}{7.297238in}}%
\pgfpathlineto{\pgfqpoint{8.455546in}{7.285905in}}%
\pgfpathlineto{\pgfqpoint{8.418270in}{7.285905in}}%
\pgfpathlineto{\pgfqpoint{8.418270in}{7.297238in}}%
\pgfpathlineto{\pgfqpoint{8.418270in}{7.297238in}}%
\pgfpathclose%
\pgfpathmoveto{\pgfqpoint{8.520690in}{7.318356in}}%
\pgfpathquadraticcurveto{\pgfqpoint{8.518543in}{7.319595in}}{\pgfqpoint{8.516020in}{7.320171in}}%
\pgfpathquadraticcurveto{\pgfqpoint{8.513496in}{7.320768in}}{\pgfqpoint{8.510464in}{7.320768in}}%
\pgfpathquadraticcurveto{\pgfqpoint{8.499662in}{7.320768in}}{\pgfqpoint{8.493884in}{7.313751in}}%
\pgfpathquadraticcurveto{\pgfqpoint{8.488107in}{7.306734in}}{\pgfqpoint{8.488107in}{7.293586in}}%
\pgfpathlineto{\pgfqpoint{8.488107in}{7.252768in}}%
\pgfpathlineto{\pgfqpoint{8.475313in}{7.252768in}}%
\pgfpathlineto{\pgfqpoint{8.475313in}{7.330242in}}%
\pgfpathlineto{\pgfqpoint{8.488107in}{7.330242in}}%
\pgfpathlineto{\pgfqpoint{8.488107in}{7.318201in}}%
\pgfpathquadraticcurveto{\pgfqpoint{8.492136in}{7.325262in}}{\pgfqpoint{8.498555in}{7.328671in}}%
\pgfpathquadraticcurveto{\pgfqpoint{8.504996in}{7.332102in}}{\pgfqpoint{8.514205in}{7.332102in}}%
\pgfpathquadraticcurveto{\pgfqpoint{8.515511in}{7.332102in}}{\pgfqpoint{8.517104in}{7.331925in}}%
\pgfpathquadraticcurveto{\pgfqpoint{8.518698in}{7.331770in}}{\pgfqpoint{8.520624in}{7.331415in}}%
\pgfpathlineto{\pgfqpoint{8.520690in}{7.318356in}}%
\pgfpathlineto{\pgfqpoint{8.520690in}{7.318356in}}%
\pgfpathclose%
\pgfpathmoveto{\pgfqpoint{8.597190in}{7.294693in}}%
\pgfpathlineto{\pgfqpoint{8.597190in}{7.288473in}}%
\pgfpathlineto{\pgfqpoint{8.538664in}{7.288473in}}%
\pgfpathquadraticcurveto{\pgfqpoint{8.539505in}{7.275324in}}{\pgfqpoint{8.546589in}{7.268440in}}%
\pgfpathquadraticcurveto{\pgfqpoint{8.553672in}{7.261556in}}{\pgfqpoint{8.566334in}{7.261556in}}%
\pgfpathquadraticcurveto{\pgfqpoint{8.573660in}{7.261556in}}{\pgfqpoint{8.580544in}{7.263349in}}%
\pgfpathquadraticcurveto{\pgfqpoint{8.587429in}{7.265142in}}{\pgfqpoint{8.594224in}{7.268750in}}%
\pgfpathlineto{\pgfqpoint{8.594224in}{7.256708in}}%
\pgfpathquadraticcurveto{\pgfqpoint{8.587362in}{7.253809in}}{\pgfqpoint{8.580168in}{7.252281in}}%
\pgfpathquadraticcurveto{\pgfqpoint{8.572974in}{7.250754in}}{\pgfqpoint{8.565581in}{7.250754in}}%
\pgfpathquadraticcurveto{\pgfqpoint{8.547031in}{7.250754in}}{\pgfqpoint{8.536207in}{7.261534in}}%
\pgfpathquadraticcurveto{\pgfqpoint{8.525383in}{7.272336in}}{\pgfqpoint{8.525383in}{7.290753in}}%
\pgfpathquadraticcurveto{\pgfqpoint{8.525383in}{7.309767in}}{\pgfqpoint{8.535654in}{7.320923in}}%
\pgfpathquadraticcurveto{\pgfqpoint{8.545925in}{7.332102in}}{\pgfqpoint{8.563367in}{7.332102in}}%
\pgfpathquadraticcurveto{\pgfqpoint{8.578995in}{7.332102in}}{\pgfqpoint{8.588093in}{7.322030in}}%
\pgfpathquadraticcurveto{\pgfqpoint{8.597190in}{7.311981in}}{\pgfqpoint{8.597190in}{7.294693in}}%
\pgfpathlineto{\pgfqpoint{8.597190in}{7.294693in}}%
\pgfpathclose%
\pgfpathmoveto{\pgfqpoint{8.584462in}{7.298434in}}%
\pgfpathquadraticcurveto{\pgfqpoint{8.584330in}{7.308859in}}{\pgfqpoint{8.578619in}{7.315080in}}%
\pgfpathquadraticcurveto{\pgfqpoint{8.572908in}{7.321322in}}{\pgfqpoint{8.563500in}{7.321322in}}%
\pgfpathquadraticcurveto{\pgfqpoint{8.552853in}{7.321322in}}{\pgfqpoint{8.546456in}{7.315301in}}%
\pgfpathquadraticcurveto{\pgfqpoint{8.540059in}{7.309280in}}{\pgfqpoint{8.539085in}{7.298345in}}%
\pgfpathlineto{\pgfqpoint{8.584462in}{7.298434in}}%
\pgfpathlineto{\pgfqpoint{8.584462in}{7.298434in}}%
\pgfpathclose%
\pgfpathmoveto{\pgfqpoint{8.608966in}{7.330242in}}%
\pgfpathlineto{\pgfqpoint{8.622447in}{7.330242in}}%
\pgfpathlineto{\pgfqpoint{8.646663in}{7.265231in}}%
\pgfpathlineto{\pgfqpoint{8.670879in}{7.330242in}}%
\pgfpathlineto{\pgfqpoint{8.684360in}{7.330242in}}%
\pgfpathlineto{\pgfqpoint{8.655296in}{7.252768in}}%
\pgfpathlineto{\pgfqpoint{8.638008in}{7.252768in}}%
\pgfpathlineto{\pgfqpoint{8.608966in}{7.330242in}}%
\pgfpathlineto{\pgfqpoint{8.608966in}{7.330242in}}%
\pgfpathclose%
\pgfpathmoveto{\pgfqpoint{8.768209in}{7.294693in}}%
\pgfpathlineto{\pgfqpoint{8.768209in}{7.288473in}}%
\pgfpathlineto{\pgfqpoint{8.709682in}{7.288473in}}%
\pgfpathquadraticcurveto{\pgfqpoint{8.710524in}{7.275324in}}{\pgfqpoint{8.717607in}{7.268440in}}%
\pgfpathquadraticcurveto{\pgfqpoint{8.724690in}{7.261556in}}{\pgfqpoint{8.737352in}{7.261556in}}%
\pgfpathquadraticcurveto{\pgfqpoint{8.744679in}{7.261556in}}{\pgfqpoint{8.751563in}{7.263349in}}%
\pgfpathquadraticcurveto{\pgfqpoint{8.758447in}{7.265142in}}{\pgfqpoint{8.765242in}{7.268750in}}%
\pgfpathlineto{\pgfqpoint{8.765242in}{7.256708in}}%
\pgfpathquadraticcurveto{\pgfqpoint{8.758380in}{7.253809in}}{\pgfqpoint{8.751186in}{7.252281in}}%
\pgfpathquadraticcurveto{\pgfqpoint{8.743992in}{7.250754in}}{\pgfqpoint{8.736599in}{7.250754in}}%
\pgfpathquadraticcurveto{\pgfqpoint{8.718050in}{7.250754in}}{\pgfqpoint{8.707225in}{7.261534in}}%
\pgfpathquadraticcurveto{\pgfqpoint{8.696401in}{7.272336in}}{\pgfqpoint{8.696401in}{7.290753in}}%
\pgfpathquadraticcurveto{\pgfqpoint{8.696401in}{7.309767in}}{\pgfqpoint{8.706672in}{7.320923in}}%
\pgfpathquadraticcurveto{\pgfqpoint{8.716943in}{7.332102in}}{\pgfqpoint{8.734386in}{7.332102in}}%
\pgfpathquadraticcurveto{\pgfqpoint{8.750013in}{7.332102in}}{\pgfqpoint{8.759111in}{7.322030in}}%
\pgfpathquadraticcurveto{\pgfqpoint{8.768209in}{7.311981in}}{\pgfqpoint{8.768209in}{7.294693in}}%
\pgfpathlineto{\pgfqpoint{8.768209in}{7.294693in}}%
\pgfpathclose%
\pgfpathmoveto{\pgfqpoint{8.755481in}{7.298434in}}%
\pgfpathquadraticcurveto{\pgfqpoint{8.755348in}{7.308859in}}{\pgfqpoint{8.749637in}{7.315080in}}%
\pgfpathquadraticcurveto{\pgfqpoint{8.743926in}{7.321322in}}{\pgfqpoint{8.734518in}{7.321322in}}%
\pgfpathquadraticcurveto{\pgfqpoint{8.723871in}{7.321322in}}{\pgfqpoint{8.717474in}{7.315301in}}%
\pgfpathquadraticcurveto{\pgfqpoint{8.711077in}{7.309280in}}{\pgfqpoint{8.710103in}{7.298345in}}%
\pgfpathlineto{\pgfqpoint{8.755481in}{7.298434in}}%
\pgfpathlineto{\pgfqpoint{8.755481in}{7.298434in}}%
\pgfpathclose%
\pgfpathmoveto{\pgfqpoint{8.833995in}{7.318356in}}%
\pgfpathquadraticcurveto{\pgfqpoint{8.831848in}{7.319595in}}{\pgfqpoint{8.829324in}{7.320171in}}%
\pgfpathquadraticcurveto{\pgfqpoint{8.826801in}{7.320768in}}{\pgfqpoint{8.823768in}{7.320768in}}%
\pgfpathquadraticcurveto{\pgfqpoint{8.812966in}{7.320768in}}{\pgfqpoint{8.807189in}{7.313751in}}%
\pgfpathquadraticcurveto{\pgfqpoint{8.801412in}{7.306734in}}{\pgfqpoint{8.801412in}{7.293586in}}%
\pgfpathlineto{\pgfqpoint{8.801412in}{7.252768in}}%
\pgfpathlineto{\pgfqpoint{8.788617in}{7.252768in}}%
\pgfpathlineto{\pgfqpoint{8.788617in}{7.330242in}}%
\pgfpathlineto{\pgfqpoint{8.801412in}{7.330242in}}%
\pgfpathlineto{\pgfqpoint{8.801412in}{7.318201in}}%
\pgfpathquadraticcurveto{\pgfqpoint{8.805440in}{7.325262in}}{\pgfqpoint{8.811860in}{7.328671in}}%
\pgfpathquadraticcurveto{\pgfqpoint{8.818301in}{7.332102in}}{\pgfqpoint{8.827509in}{7.332102in}}%
\pgfpathquadraticcurveto{\pgfqpoint{8.828815in}{7.332102in}}{\pgfqpoint{8.830409in}{7.331925in}}%
\pgfpathquadraticcurveto{\pgfqpoint{8.832003in}{7.331770in}}{\pgfqpoint{8.833929in}{7.331415in}}%
\pgfpathlineto{\pgfqpoint{8.833995in}{7.318356in}}%
\pgfpathlineto{\pgfqpoint{8.833995in}{7.318356in}}%
\pgfpathclose%
\pgfpathmoveto{\pgfqpoint{8.896727in}{7.327962in}}%
\pgfpathlineto{\pgfqpoint{8.896727in}{7.315921in}}%
\pgfpathquadraticcurveto{\pgfqpoint{8.891348in}{7.318688in}}{\pgfqpoint{8.885526in}{7.320060in}}%
\pgfpathquadraticcurveto{\pgfqpoint{8.879727in}{7.321455in}}{\pgfqpoint{8.873485in}{7.321455in}}%
\pgfpathquadraticcurveto{\pgfqpoint{8.864011in}{7.321455in}}{\pgfqpoint{8.859274in}{7.318555in}}%
\pgfpathquadraticcurveto{\pgfqpoint{8.854537in}{7.315655in}}{\pgfqpoint{8.854537in}{7.309833in}}%
\pgfpathquadraticcurveto{\pgfqpoint{8.854537in}{7.305406in}}{\pgfqpoint{8.857923in}{7.302883in}}%
\pgfpathquadraticcurveto{\pgfqpoint{8.861310in}{7.300359in}}{\pgfqpoint{8.871559in}{7.298080in}}%
\pgfpathlineto{\pgfqpoint{8.875919in}{7.297106in}}%
\pgfpathquadraticcurveto{\pgfqpoint{8.889466in}{7.294206in}}{\pgfqpoint{8.895177in}{7.288915in}}%
\pgfpathquadraticcurveto{\pgfqpoint{8.900888in}{7.283625in}}{\pgfqpoint{8.900888in}{7.274151in}}%
\pgfpathquadraticcurveto{\pgfqpoint{8.900888in}{7.263349in}}{\pgfqpoint{8.892344in}{7.257040in}}%
\pgfpathquadraticcurveto{\pgfqpoint{8.883800in}{7.250754in}}{\pgfqpoint{8.868858in}{7.250754in}}%
\pgfpathquadraticcurveto{\pgfqpoint{8.862638in}{7.250754in}}{\pgfqpoint{8.855887in}{7.251971in}}%
\pgfpathquadraticcurveto{\pgfqpoint{8.849136in}{7.253189in}}{\pgfqpoint{8.841676in}{7.255602in}}%
\pgfpathlineto{\pgfqpoint{8.841676in}{7.268750in}}%
\pgfpathquadraticcurveto{\pgfqpoint{8.848737in}{7.265076in}}{\pgfqpoint{8.855577in}{7.263238in}}%
\pgfpathquadraticcurveto{\pgfqpoint{8.862417in}{7.261423in}}{\pgfqpoint{8.869146in}{7.261423in}}%
\pgfpathquadraticcurveto{\pgfqpoint{8.878133in}{7.261423in}}{\pgfqpoint{8.882959in}{7.264500in}}%
\pgfpathquadraticcurveto{\pgfqpoint{8.887806in}{7.267577in}}{\pgfqpoint{8.887806in}{7.273177in}}%
\pgfpathquadraticcurveto{\pgfqpoint{8.887806in}{7.278357in}}{\pgfqpoint{8.884309in}{7.281124in}}%
\pgfpathquadraticcurveto{\pgfqpoint{8.880834in}{7.283891in}}{\pgfqpoint{8.868991in}{7.286458in}}%
\pgfpathlineto{\pgfqpoint{8.864564in}{7.287499in}}%
\pgfpathquadraticcurveto{\pgfqpoint{8.852744in}{7.289978in}}{\pgfqpoint{8.847475in}{7.295136in}}%
\pgfpathquadraticcurveto{\pgfqpoint{8.842229in}{7.300293in}}{\pgfqpoint{8.842229in}{7.309280in}}%
\pgfpathquadraticcurveto{\pgfqpoint{8.842229in}{7.320215in}}{\pgfqpoint{8.849977in}{7.326147in}}%
\pgfpathquadraticcurveto{\pgfqpoint{8.857724in}{7.332102in}}{\pgfqpoint{8.871979in}{7.332102in}}%
\pgfpathquadraticcurveto{\pgfqpoint{8.879018in}{7.332102in}}{\pgfqpoint{8.885238in}{7.331061in}}%
\pgfpathquadraticcurveto{\pgfqpoint{8.891481in}{7.330043in}}{\pgfqpoint{8.896727in}{7.327962in}}%
\pgfpathlineto{\pgfqpoint{8.896727in}{7.327962in}}%
\pgfpathclose%
\pgfpathmoveto{\pgfqpoint{8.956360in}{7.291705in}}%
\pgfpathquadraticcurveto{\pgfqpoint{8.940931in}{7.291705in}}{\pgfqpoint{8.934977in}{7.288185in}}%
\pgfpathquadraticcurveto{\pgfqpoint{8.929022in}{7.284665in}}{\pgfqpoint{8.929022in}{7.276143in}}%
\pgfpathquadraticcurveto{\pgfqpoint{8.929022in}{7.269370in}}{\pgfqpoint{8.933494in}{7.265386in}}%
\pgfpathquadraticcurveto{\pgfqpoint{8.937965in}{7.261423in}}{\pgfqpoint{8.945624in}{7.261423in}}%
\pgfpathquadraticcurveto{\pgfqpoint{8.956227in}{7.261423in}}{\pgfqpoint{8.962624in}{7.268927in}}%
\pgfpathquadraticcurveto{\pgfqpoint{8.969021in}{7.276431in}}{\pgfqpoint{8.969021in}{7.288871in}}%
\pgfpathlineto{\pgfqpoint{8.969021in}{7.291705in}}%
\pgfpathlineto{\pgfqpoint{8.956360in}{7.291705in}}%
\pgfpathlineto{\pgfqpoint{8.956360in}{7.291705in}}%
\pgfpathclose%
\pgfpathmoveto{\pgfqpoint{8.981749in}{7.296973in}}%
\pgfpathlineto{\pgfqpoint{8.981749in}{7.252768in}}%
\pgfpathlineto{\pgfqpoint{8.969021in}{7.252768in}}%
\pgfpathlineto{\pgfqpoint{8.969021in}{7.264522in}}%
\pgfpathquadraticcurveto{\pgfqpoint{8.964660in}{7.257483in}}{\pgfqpoint{8.958153in}{7.254119in}}%
\pgfpathquadraticcurveto{\pgfqpoint{8.951645in}{7.250754in}}{\pgfqpoint{8.942237in}{7.250754in}}%
\pgfpathquadraticcurveto{\pgfqpoint{8.930350in}{7.250754in}}{\pgfqpoint{8.923311in}{7.257439in}}%
\pgfpathquadraticcurveto{\pgfqpoint{8.916294in}{7.264124in}}{\pgfqpoint{8.916294in}{7.275324in}}%
\pgfpathquadraticcurveto{\pgfqpoint{8.916294in}{7.288384in}}{\pgfqpoint{8.925038in}{7.295025in}}%
\pgfpathquadraticcurveto{\pgfqpoint{8.933804in}{7.301665in}}{\pgfqpoint{8.951158in}{7.301665in}}%
\pgfpathlineto{\pgfqpoint{8.969021in}{7.301665in}}%
\pgfpathlineto{\pgfqpoint{8.969021in}{7.302927in}}%
\pgfpathquadraticcurveto{\pgfqpoint{8.969021in}{7.311715in}}{\pgfqpoint{8.963244in}{7.316518in}}%
\pgfpathquadraticcurveto{\pgfqpoint{8.957466in}{7.321322in}}{\pgfqpoint{8.947018in}{7.321322in}}%
\pgfpathquadraticcurveto{\pgfqpoint{8.940378in}{7.321322in}}{\pgfqpoint{8.934069in}{7.319728in}}%
\pgfpathquadraticcurveto{\pgfqpoint{8.927783in}{7.318134in}}{\pgfqpoint{8.921983in}{7.314947in}}%
\pgfpathlineto{\pgfqpoint{8.921983in}{7.326723in}}%
\pgfpathquadraticcurveto{\pgfqpoint{8.928956in}{7.329423in}}{\pgfqpoint{8.935530in}{7.330751in}}%
\pgfpathquadraticcurveto{\pgfqpoint{8.942104in}{7.332102in}}{\pgfqpoint{8.948324in}{7.332102in}}%
\pgfpathquadraticcurveto{\pgfqpoint{8.965147in}{7.332102in}}{\pgfqpoint{8.973448in}{7.323380in}}%
\pgfpathquadraticcurveto{\pgfqpoint{8.981749in}{7.314681in}}{\pgfqpoint{8.981749in}{7.296973in}}%
\pgfpathlineto{\pgfqpoint{8.981749in}{7.296973in}}%
\pgfpathclose%
\pgfpathmoveto{\pgfqpoint{9.007957in}{7.360413in}}%
\pgfpathlineto{\pgfqpoint{9.020685in}{7.360413in}}%
\pgfpathlineto{\pgfqpoint{9.020685in}{7.252768in}}%
\pgfpathlineto{\pgfqpoint{9.007957in}{7.252768in}}%
\pgfpathlineto{\pgfqpoint{9.007957in}{7.360413in}}%
\pgfpathlineto{\pgfqpoint{9.007957in}{7.360413in}}%
\pgfpathclose%
\pgfpathmoveto{\pgfqpoint{9.096698in}{7.327962in}}%
\pgfpathlineto{\pgfqpoint{9.096698in}{7.315921in}}%
\pgfpathquadraticcurveto{\pgfqpoint{9.091319in}{7.318688in}}{\pgfqpoint{9.085498in}{7.320060in}}%
\pgfpathquadraticcurveto{\pgfqpoint{9.079698in}{7.321455in}}{\pgfqpoint{9.073456in}{7.321455in}}%
\pgfpathquadraticcurveto{\pgfqpoint{9.063982in}{7.321455in}}{\pgfqpoint{9.059245in}{7.318555in}}%
\pgfpathquadraticcurveto{\pgfqpoint{9.054508in}{7.315655in}}{\pgfqpoint{9.054508in}{7.309833in}}%
\pgfpathquadraticcurveto{\pgfqpoint{9.054508in}{7.305406in}}{\pgfqpoint{9.057895in}{7.302883in}}%
\pgfpathquadraticcurveto{\pgfqpoint{9.061281in}{7.300359in}}{\pgfqpoint{9.071530in}{7.298080in}}%
\pgfpathlineto{\pgfqpoint{9.075891in}{7.297106in}}%
\pgfpathquadraticcurveto{\pgfqpoint{9.089438in}{7.294206in}}{\pgfqpoint{9.095149in}{7.288915in}}%
\pgfpathquadraticcurveto{\pgfqpoint{9.100860in}{7.283625in}}{\pgfqpoint{9.100860in}{7.274151in}}%
\pgfpathquadraticcurveto{\pgfqpoint{9.100860in}{7.263349in}}{\pgfqpoint{9.092315in}{7.257040in}}%
\pgfpathquadraticcurveto{\pgfqpoint{9.083771in}{7.250754in}}{\pgfqpoint{9.068830in}{7.250754in}}%
\pgfpathquadraticcurveto{\pgfqpoint{9.062610in}{7.250754in}}{\pgfqpoint{9.055858in}{7.251971in}}%
\pgfpathquadraticcurveto{\pgfqpoint{9.049107in}{7.253189in}}{\pgfqpoint{9.041647in}{7.255602in}}%
\pgfpathlineto{\pgfqpoint{9.041647in}{7.268750in}}%
\pgfpathquadraticcurveto{\pgfqpoint{9.048709in}{7.265076in}}{\pgfqpoint{9.055548in}{7.263238in}}%
\pgfpathquadraticcurveto{\pgfqpoint{9.062388in}{7.261423in}}{\pgfqpoint{9.069117in}{7.261423in}}%
\pgfpathquadraticcurveto{\pgfqpoint{9.078104in}{7.261423in}}{\pgfqpoint{9.082930in}{7.264500in}}%
\pgfpathquadraticcurveto{\pgfqpoint{9.087778in}{7.267577in}}{\pgfqpoint{9.087778in}{7.273177in}}%
\pgfpathquadraticcurveto{\pgfqpoint{9.087778in}{7.278357in}}{\pgfqpoint{9.084280in}{7.281124in}}%
\pgfpathquadraticcurveto{\pgfqpoint{9.080805in}{7.283891in}}{\pgfqpoint{9.068962in}{7.286458in}}%
\pgfpathlineto{\pgfqpoint{9.064535in}{7.287499in}}%
\pgfpathquadraticcurveto{\pgfqpoint{9.052715in}{7.289978in}}{\pgfqpoint{9.047447in}{7.295136in}}%
\pgfpathquadraticcurveto{\pgfqpoint{9.042201in}{7.300293in}}{\pgfqpoint{9.042201in}{7.309280in}}%
\pgfpathquadraticcurveto{\pgfqpoint{9.042201in}{7.320215in}}{\pgfqpoint{9.049948in}{7.326147in}}%
\pgfpathquadraticcurveto{\pgfqpoint{9.057695in}{7.332102in}}{\pgfqpoint{9.071951in}{7.332102in}}%
\pgfpathquadraticcurveto{\pgfqpoint{9.078990in}{7.332102in}}{\pgfqpoint{9.085210in}{7.331061in}}%
\pgfpathquadraticcurveto{\pgfqpoint{9.091452in}{7.330043in}}{\pgfqpoint{9.096698in}{7.327962in}}%
\pgfpathlineto{\pgfqpoint{9.096698in}{7.327962in}}%
\pgfpathclose%
\pgfpathmoveto{\pgfqpoint{9.124367in}{7.270344in}}%
\pgfpathlineto{\pgfqpoint{9.138955in}{7.270344in}}%
\pgfpathlineto{\pgfqpoint{9.138955in}{7.252768in}}%
\pgfpathlineto{\pgfqpoint{9.124367in}{7.252768in}}%
\pgfpathlineto{\pgfqpoint{9.124367in}{7.270344in}}%
\pgfpathlineto{\pgfqpoint{9.124367in}{7.270344in}}%
\pgfpathclose%
\pgfpathmoveto{\pgfqpoint{9.124367in}{7.326014in}}%
\pgfpathlineto{\pgfqpoint{9.138955in}{7.326014in}}%
\pgfpathlineto{\pgfqpoint{9.138955in}{7.308461in}}%
\pgfpathlineto{\pgfqpoint{9.124367in}{7.308461in}}%
\pgfpathlineto{\pgfqpoint{9.124367in}{7.326014in}}%
\pgfpathlineto{\pgfqpoint{9.124367in}{7.326014in}}%
\pgfpathclose%
\pgfpathmoveto{\pgfqpoint{9.247285in}{7.309966in}}%
\pgfpathquadraticcurveto{\pgfqpoint{9.237878in}{7.309966in}}{\pgfqpoint{9.232366in}{7.303525in}}%
\pgfpathquadraticcurveto{\pgfqpoint{9.226876in}{7.297106in}}{\pgfqpoint{9.226876in}{7.285905in}}%
\pgfpathquadraticcurveto{\pgfqpoint{9.226876in}{7.274771in}}{\pgfqpoint{9.232366in}{7.268285in}}%
\pgfpathquadraticcurveto{\pgfqpoint{9.237878in}{7.261822in}}{\pgfqpoint{9.247285in}{7.261822in}}%
\pgfpathquadraticcurveto{\pgfqpoint{9.256693in}{7.261822in}}{\pgfqpoint{9.262182in}{7.268285in}}%
\pgfpathquadraticcurveto{\pgfqpoint{9.267672in}{7.274771in}}{\pgfqpoint{9.267672in}{7.285905in}}%
\pgfpathquadraticcurveto{\pgfqpoint{9.267672in}{7.297106in}}{\pgfqpoint{9.262182in}{7.303525in}}%
\pgfpathquadraticcurveto{\pgfqpoint{9.256693in}{7.309966in}}{\pgfqpoint{9.247285in}{7.309966in}}%
\pgfpathlineto{\pgfqpoint{9.247285in}{7.309966in}}%
\pgfpathclose%
\pgfpathmoveto{\pgfqpoint{9.275021in}{7.353772in}}%
\pgfpathlineto{\pgfqpoint{9.275021in}{7.341044in}}%
\pgfpathquadraticcurveto{\pgfqpoint{9.269753in}{7.343524in}}{\pgfqpoint{9.264396in}{7.344830in}}%
\pgfpathquadraticcurveto{\pgfqpoint{9.259039in}{7.346158in}}{\pgfqpoint{9.253771in}{7.346158in}}%
\pgfpathquadraticcurveto{\pgfqpoint{9.239936in}{7.346158in}}{\pgfqpoint{9.232632in}{7.336817in}}%
\pgfpathquadraticcurveto{\pgfqpoint{9.225349in}{7.327475in}}{\pgfqpoint{9.224309in}{7.308594in}}%
\pgfpathquadraticcurveto{\pgfqpoint{9.228382in}{7.314615in}}{\pgfqpoint{9.234535in}{7.317824in}}%
\pgfpathquadraticcurveto{\pgfqpoint{9.240711in}{7.321034in}}{\pgfqpoint{9.248104in}{7.321034in}}%
\pgfpathquadraticcurveto{\pgfqpoint{9.263666in}{7.321034in}}{\pgfqpoint{9.272697in}{7.311582in}}%
\pgfpathquadraticcurveto{\pgfqpoint{9.281728in}{7.302152in}}{\pgfqpoint{9.281728in}{7.285905in}}%
\pgfpathquadraticcurveto{\pgfqpoint{9.281728in}{7.269990in}}{\pgfqpoint{9.272320in}{7.260361in}}%
\pgfpathquadraticcurveto{\pgfqpoint{9.262913in}{7.250754in}}{\pgfqpoint{9.247285in}{7.250754in}}%
\pgfpathquadraticcurveto{\pgfqpoint{9.229356in}{7.250754in}}{\pgfqpoint{9.219882in}{7.264478in}}%
\pgfpathquadraticcurveto{\pgfqpoint{9.210408in}{7.278224in}}{\pgfqpoint{9.210408in}{7.304300in}}%
\pgfpathquadraticcurveto{\pgfqpoint{9.210408in}{7.328781in}}{\pgfqpoint{9.222029in}{7.343346in}}%
\pgfpathquadraticcurveto{\pgfqpoint{9.233650in}{7.357912in}}{\pgfqpoint{9.253218in}{7.357912in}}%
\pgfpathquadraticcurveto{\pgfqpoint{9.258486in}{7.357912in}}{\pgfqpoint{9.263843in}{7.356871in}}%
\pgfpathquadraticcurveto{\pgfqpoint{9.269199in}{7.355831in}}{\pgfqpoint{9.275021in}{7.353772in}}%
\pgfpathlineto{\pgfqpoint{9.275021in}{7.353772in}}%
\pgfpathclose%
\pgfpathmoveto{\pgfqpoint{9.305789in}{7.270344in}}%
\pgfpathlineto{\pgfqpoint{9.320399in}{7.270344in}}%
\pgfpathlineto{\pgfqpoint{9.320399in}{7.252768in}}%
\pgfpathlineto{\pgfqpoint{9.305789in}{7.252768in}}%
\pgfpathlineto{\pgfqpoint{9.305789in}{7.270344in}}%
\pgfpathlineto{\pgfqpoint{9.305789in}{7.270344in}}%
\pgfpathclose%
\pgfpathmoveto{\pgfqpoint{9.389218in}{7.343878in}}%
\pgfpathlineto{\pgfqpoint{9.353934in}{7.288738in}}%
\pgfpathlineto{\pgfqpoint{9.389218in}{7.288738in}}%
\pgfpathlineto{\pgfqpoint{9.389218in}{7.343878in}}%
\pgfpathlineto{\pgfqpoint{9.389218in}{7.343878in}}%
\pgfpathclose%
\pgfpathmoveto{\pgfqpoint{9.385543in}{7.356052in}}%
\pgfpathlineto{\pgfqpoint{9.403119in}{7.356052in}}%
\pgfpathlineto{\pgfqpoint{9.403119in}{7.288738in}}%
\pgfpathlineto{\pgfqpoint{9.417861in}{7.288738in}}%
\pgfpathlineto{\pgfqpoint{9.417861in}{7.277117in}}%
\pgfpathlineto{\pgfqpoint{9.403119in}{7.277117in}}%
\pgfpathlineto{\pgfqpoint{9.403119in}{7.252768in}}%
\pgfpathlineto{\pgfqpoint{9.389218in}{7.252768in}}%
\pgfpathlineto{\pgfqpoint{9.389218in}{7.277117in}}%
\pgfpathlineto{\pgfqpoint{9.342600in}{7.277117in}}%
\pgfpathlineto{\pgfqpoint{9.342600in}{7.290598in}}%
\pgfpathlineto{\pgfqpoint{9.385543in}{7.356052in}}%
\pgfpathlineto{\pgfqpoint{9.385543in}{7.356052in}}%
\pgfpathclose%
\pgfpathmoveto{\pgfqpoint{9.528804in}{7.298212in}}%
\pgfpathquadraticcurveto{\pgfqpoint{9.522783in}{7.298212in}}{\pgfqpoint{9.519352in}{7.293099in}}%
\pgfpathquadraticcurveto{\pgfqpoint{9.515943in}{7.287986in}}{\pgfqpoint{9.515943in}{7.278844in}}%
\pgfpathquadraticcurveto{\pgfqpoint{9.515943in}{7.269857in}}{\pgfqpoint{9.519352in}{7.264699in}}%
\pgfpathquadraticcurveto{\pgfqpoint{9.522783in}{7.259542in}}{\pgfqpoint{9.528804in}{7.259542in}}%
\pgfpathquadraticcurveto{\pgfqpoint{9.534692in}{7.259542in}}{\pgfqpoint{9.538100in}{7.264699in}}%
\pgfpathquadraticcurveto{\pgfqpoint{9.541531in}{7.269857in}}{\pgfqpoint{9.541531in}{7.278844in}}%
\pgfpathquadraticcurveto{\pgfqpoint{9.541531in}{7.287919in}}{\pgfqpoint{9.538100in}{7.293055in}}%
\pgfpathquadraticcurveto{\pgfqpoint{9.534692in}{7.298212in}}{\pgfqpoint{9.528804in}{7.298212in}}%
\pgfpathlineto{\pgfqpoint{9.528804in}{7.298212in}}%
\pgfpathclose%
\pgfpathmoveto{\pgfqpoint{9.528804in}{7.307000in}}%
\pgfpathquadraticcurveto{\pgfqpoint{9.539738in}{7.307000in}}{\pgfqpoint{9.546158in}{7.299386in}}%
\pgfpathquadraticcurveto{\pgfqpoint{9.552599in}{7.291793in}}{\pgfqpoint{9.552599in}{7.278844in}}%
\pgfpathquadraticcurveto{\pgfqpoint{9.552599in}{7.265917in}}{\pgfqpoint{9.546136in}{7.258324in}}%
\pgfpathquadraticcurveto{\pgfqpoint{9.539672in}{7.250754in}}{\pgfqpoint{9.528804in}{7.250754in}}%
\pgfpathquadraticcurveto{\pgfqpoint{9.517736in}{7.250754in}}{\pgfqpoint{9.511294in}{7.258324in}}%
\pgfpathquadraticcurveto{\pgfqpoint{9.504875in}{7.265917in}}{\pgfqpoint{9.504875in}{7.278844in}}%
\pgfpathquadraticcurveto{\pgfqpoint{9.504875in}{7.291859in}}{\pgfqpoint{9.511339in}{7.299430in}}%
\pgfpathquadraticcurveto{\pgfqpoint{9.517802in}{7.307000in}}{\pgfqpoint{9.528804in}{7.307000in}}%
\pgfpathlineto{\pgfqpoint{9.528804in}{7.307000in}}%
\pgfpathclose%
\pgfpathmoveto{\pgfqpoint{9.457417in}{7.349124in}}%
\pgfpathquadraticcurveto{\pgfqpoint{9.451462in}{7.349124in}}{\pgfqpoint{9.448031in}{7.343966in}}%
\pgfpathquadraticcurveto{\pgfqpoint{9.444623in}{7.338831in}}{\pgfqpoint{9.444623in}{7.329822in}}%
\pgfpathquadraticcurveto{\pgfqpoint{9.444623in}{7.320702in}}{\pgfqpoint{9.448009in}{7.315567in}}%
\pgfpathquadraticcurveto{\pgfqpoint{9.451396in}{7.310453in}}{\pgfqpoint{9.457417in}{7.310453in}}%
\pgfpathquadraticcurveto{\pgfqpoint{9.463438in}{7.310453in}}{\pgfqpoint{9.466847in}{7.315567in}}%
\pgfpathquadraticcurveto{\pgfqpoint{9.470278in}{7.320702in}}{\pgfqpoint{9.470278in}{7.329822in}}%
\pgfpathquadraticcurveto{\pgfqpoint{9.470278in}{7.338742in}}{\pgfqpoint{9.466824in}{7.343922in}}%
\pgfpathquadraticcurveto{\pgfqpoint{9.463371in}{7.349124in}}{\pgfqpoint{9.457417in}{7.349124in}}%
\pgfpathlineto{\pgfqpoint{9.457417in}{7.349124in}}%
\pgfpathclose%
\pgfpathmoveto{\pgfqpoint{9.519883in}{7.357912in}}%
\pgfpathlineto{\pgfqpoint{9.530951in}{7.357912in}}%
\pgfpathlineto{\pgfqpoint{9.466337in}{7.250754in}}%
\pgfpathlineto{\pgfqpoint{9.455270in}{7.250754in}}%
\pgfpathlineto{\pgfqpoint{9.519883in}{7.357912in}}%
\pgfpathlineto{\pgfqpoint{9.519883in}{7.357912in}}%
\pgfpathclose%
\pgfpathmoveto{\pgfqpoint{9.457417in}{7.357912in}}%
\pgfpathquadraticcurveto{\pgfqpoint{9.468352in}{7.357912in}}{\pgfqpoint{9.474837in}{7.350341in}}%
\pgfpathquadraticcurveto{\pgfqpoint{9.481345in}{7.342771in}}{\pgfqpoint{9.481345in}{7.329822in}}%
\pgfpathquadraticcurveto{\pgfqpoint{9.481345in}{7.316762in}}{\pgfqpoint{9.474882in}{7.309214in}}%
\pgfpathquadraticcurveto{\pgfqpoint{9.468418in}{7.301665in}}{\pgfqpoint{9.457417in}{7.301665in}}%
\pgfpathquadraticcurveto{\pgfqpoint{9.446416in}{7.301665in}}{\pgfqpoint{9.440018in}{7.309236in}}%
\pgfpathquadraticcurveto{\pgfqpoint{9.433621in}{7.316828in}}{\pgfqpoint{9.433621in}{7.329822in}}%
\pgfpathquadraticcurveto{\pgfqpoint{9.433621in}{7.342705in}}{\pgfqpoint{9.440041in}{7.350297in}}%
\pgfpathquadraticcurveto{\pgfqpoint{9.446482in}{7.357912in}}{\pgfqpoint{9.457417in}{7.357912in}}%
\pgfpathlineto{\pgfqpoint{9.457417in}{7.357912in}}%
\pgfpathclose%
\pgfusepath{fill}%
\end{pgfscope}%
\begin{pgfscope}%
\pgfsetbuttcap%
\pgfsetroundjoin%
\definecolor{currentfill}{rgb}{0.090196,0.168627,0.227451}%
\pgfsetfillcolor{currentfill}%
\pgfsetlinewidth{0.000000pt}%
\definecolor{currentstroke}{rgb}{0.090196,0.168627,0.227451}%
\pgfsetstrokecolor{currentstroke}%
\pgfsetdash{}{0pt}%
\pgfpathmoveto{\pgfqpoint{8.178145in}{7.133082in}}%
\pgfpathlineto{\pgfqpoint{8.178145in}{7.119447in}}%
\pgfpathquadraticcurveto{\pgfqpoint{8.170198in}{7.123254in}}{\pgfqpoint{8.163137in}{7.125113in}}%
\pgfpathquadraticcurveto{\pgfqpoint{8.156076in}{7.126995in}}{\pgfqpoint{8.149501in}{7.126995in}}%
\pgfpathquadraticcurveto{\pgfqpoint{8.138102in}{7.126995in}}{\pgfqpoint{8.131904in}{7.122568in}}%
\pgfpathquadraticcurveto{\pgfqpoint{8.125706in}{7.118141in}}{\pgfqpoint{8.125706in}{7.109973in}}%
\pgfpathquadraticcurveto{\pgfqpoint{8.125706in}{7.103111in}}{\pgfqpoint{8.129823in}{7.099613in}}%
\pgfpathquadraticcurveto{\pgfqpoint{8.133940in}{7.096138in}}{\pgfqpoint{8.145429in}{7.093991in}}%
\pgfpathlineto{\pgfqpoint{8.153862in}{7.092264in}}%
\pgfpathquadraticcurveto{\pgfqpoint{8.169490in}{7.089276in}}{\pgfqpoint{8.176927in}{7.081772in}}%
\pgfpathquadraticcurveto{\pgfqpoint{8.184365in}{7.074268in}}{\pgfqpoint{8.184365in}{7.061695in}}%
\pgfpathquadraticcurveto{\pgfqpoint{8.184365in}{7.046665in}}{\pgfqpoint{8.174293in}{7.038918in}}%
\pgfpathquadraticcurveto{\pgfqpoint{8.164244in}{7.031171in}}{\pgfqpoint{8.144809in}{7.031171in}}%
\pgfpathquadraticcurveto{\pgfqpoint{8.137482in}{7.031171in}}{\pgfqpoint{8.129203in}{7.032831in}}%
\pgfpathquadraticcurveto{\pgfqpoint{8.120947in}{7.034491in}}{\pgfqpoint{8.112093in}{7.037745in}}%
\pgfpathlineto{\pgfqpoint{8.112093in}{7.052133in}}%
\pgfpathquadraticcurveto{\pgfqpoint{8.120593in}{7.047374in}}{\pgfqpoint{8.128761in}{7.044939in}}%
\pgfpathquadraticcurveto{\pgfqpoint{8.136929in}{7.042526in}}{\pgfqpoint{8.144809in}{7.042526in}}%
\pgfpathquadraticcurveto{\pgfqpoint{8.156762in}{7.042526in}}{\pgfqpoint{8.163270in}{7.047219in}}%
\pgfpathquadraticcurveto{\pgfqpoint{8.169778in}{7.051934in}}{\pgfqpoint{8.169778in}{7.060655in}}%
\pgfpathquadraticcurveto{\pgfqpoint{8.169778in}{7.068248in}}{\pgfqpoint{8.165107in}{7.072542in}}%
\pgfpathquadraticcurveto{\pgfqpoint{8.160436in}{7.076836in}}{\pgfqpoint{8.149789in}{7.078983in}}%
\pgfpathlineto{\pgfqpoint{8.141267in}{7.080643in}}%
\pgfpathquadraticcurveto{\pgfqpoint{8.125640in}{7.083742in}}{\pgfqpoint{8.118645in}{7.090383in}}%
\pgfpathquadraticcurveto{\pgfqpoint{8.111672in}{7.097024in}}{\pgfqpoint{8.111672in}{7.108866in}}%
\pgfpathquadraticcurveto{\pgfqpoint{8.111672in}{7.122568in}}{\pgfqpoint{8.121323in}{7.130448in}}%
\pgfpathquadraticcurveto{\pgfqpoint{8.130974in}{7.138328in}}{\pgfqpoint{8.147908in}{7.138328in}}%
\pgfpathquadraticcurveto{\pgfqpoint{8.155190in}{7.138328in}}{\pgfqpoint{8.162716in}{7.137000in}}%
\pgfpathquadraticcurveto{\pgfqpoint{8.170265in}{7.135694in}}{\pgfqpoint{8.178145in}{7.133082in}}%
\pgfpathlineto{\pgfqpoint{8.178145in}{7.133082in}}%
\pgfpathclose%
\pgfpathmoveto{\pgfqpoint{8.205615in}{7.110659in}}%
\pgfpathlineto{\pgfqpoint{8.218343in}{7.110659in}}%
\pgfpathlineto{\pgfqpoint{8.218343in}{7.033185in}}%
\pgfpathlineto{\pgfqpoint{8.205615in}{7.033185in}}%
\pgfpathlineto{\pgfqpoint{8.205615in}{7.110659in}}%
\pgfpathlineto{\pgfqpoint{8.205615in}{7.110659in}}%
\pgfpathclose%
\pgfpathmoveto{\pgfqpoint{8.205615in}{7.140830in}}%
\pgfpathlineto{\pgfqpoint{8.218343in}{7.140830in}}%
\pgfpathlineto{\pgfqpoint{8.218343in}{7.124693in}}%
\pgfpathlineto{\pgfqpoint{8.205615in}{7.124693in}}%
\pgfpathlineto{\pgfqpoint{8.205615in}{7.140830in}}%
\pgfpathlineto{\pgfqpoint{8.205615in}{7.140830in}}%
\pgfpathclose%
\pgfpathmoveto{\pgfqpoint{8.295949in}{7.072830in}}%
\pgfpathquadraticcurveto{\pgfqpoint{8.295949in}{7.086664in}}{\pgfqpoint{8.290238in}{7.094257in}}%
\pgfpathquadraticcurveto{\pgfqpoint{8.284550in}{7.101871in}}{\pgfqpoint{8.274235in}{7.101871in}}%
\pgfpathquadraticcurveto{\pgfqpoint{8.264008in}{7.101871in}}{\pgfqpoint{8.258297in}{7.094257in}}%
\pgfpathquadraticcurveto{\pgfqpoint{8.252586in}{7.086664in}}{\pgfqpoint{8.252586in}{7.072830in}}%
\pgfpathquadraticcurveto{\pgfqpoint{8.252586in}{7.059061in}}{\pgfqpoint{8.258297in}{7.051447in}}%
\pgfpathquadraticcurveto{\pgfqpoint{8.264008in}{7.043832in}}{\pgfqpoint{8.274235in}{7.043832in}}%
\pgfpathquadraticcurveto{\pgfqpoint{8.284550in}{7.043832in}}{\pgfqpoint{8.290238in}{7.051447in}}%
\pgfpathquadraticcurveto{\pgfqpoint{8.295949in}{7.059061in}}{\pgfqpoint{8.295949in}{7.072830in}}%
\pgfpathlineto{\pgfqpoint{8.295949in}{7.072830in}}%
\pgfpathclose%
\pgfpathmoveto{\pgfqpoint{8.308677in}{7.042792in}}%
\pgfpathquadraticcurveto{\pgfqpoint{8.308677in}{7.023025in}}{\pgfqpoint{8.299890in}{7.013374in}}%
\pgfpathquadraticcurveto{\pgfqpoint{8.291124in}{7.003723in}}{\pgfqpoint{8.272995in}{7.003723in}}%
\pgfpathquadraticcurveto{\pgfqpoint{8.266288in}{7.003723in}}{\pgfqpoint{8.260334in}{7.004719in}}%
\pgfpathquadraticcurveto{\pgfqpoint{8.254379in}{7.005715in}}{\pgfqpoint{8.248779in}{7.007796in}}%
\pgfpathlineto{\pgfqpoint{8.248779in}{7.020169in}}%
\pgfpathquadraticcurveto{\pgfqpoint{8.254379in}{7.017137in}}{\pgfqpoint{8.259847in}{7.015698in}}%
\pgfpathquadraticcurveto{\pgfqpoint{8.265314in}{7.014237in}}{\pgfqpoint{8.270981in}{7.014237in}}%
\pgfpathquadraticcurveto{\pgfqpoint{8.283509in}{7.014237in}}{\pgfqpoint{8.289729in}{7.020767in}}%
\pgfpathquadraticcurveto{\pgfqpoint{8.295949in}{7.027297in}}{\pgfqpoint{8.295949in}{7.040512in}}%
\pgfpathlineto{\pgfqpoint{8.295949in}{7.046820in}}%
\pgfpathquadraticcurveto{\pgfqpoint{8.292009in}{7.039958in}}{\pgfqpoint{8.285856in}{7.036572in}}%
\pgfpathquadraticcurveto{\pgfqpoint{8.279702in}{7.033185in}}{\pgfqpoint{8.271113in}{7.033185in}}%
\pgfpathquadraticcurveto{\pgfqpoint{8.256880in}{7.033185in}}{\pgfqpoint{8.248159in}{7.044031in}}%
\pgfpathquadraticcurveto{\pgfqpoint{8.239438in}{7.054900in}}{\pgfqpoint{8.239438in}{7.072830in}}%
\pgfpathquadraticcurveto{\pgfqpoint{8.239438in}{7.090803in}}{\pgfqpoint{8.248159in}{7.101650in}}%
\pgfpathquadraticcurveto{\pgfqpoint{8.256880in}{7.112518in}}{\pgfqpoint{8.271113in}{7.112518in}}%
\pgfpathquadraticcurveto{\pgfqpoint{8.279702in}{7.112518in}}{\pgfqpoint{8.285856in}{7.109132in}}%
\pgfpathquadraticcurveto{\pgfqpoint{8.292009in}{7.105745in}}{\pgfqpoint{8.295949in}{7.098905in}}%
\pgfpathlineto{\pgfqpoint{8.295949in}{7.110659in}}%
\pgfpathlineto{\pgfqpoint{8.308677in}{7.110659in}}%
\pgfpathlineto{\pgfqpoint{8.308677in}{7.042792in}}%
\pgfpathlineto{\pgfqpoint{8.308677in}{7.042792in}}%
\pgfpathclose%
\pgfpathmoveto{\pgfqpoint{8.399322in}{7.079957in}}%
\pgfpathlineto{\pgfqpoint{8.399322in}{7.033185in}}%
\pgfpathlineto{\pgfqpoint{8.386594in}{7.033185in}}%
\pgfpathlineto{\pgfqpoint{8.386594in}{7.079537in}}%
\pgfpathquadraticcurveto{\pgfqpoint{8.386594in}{7.090538in}}{\pgfqpoint{8.382300in}{7.095983in}}%
\pgfpathquadraticcurveto{\pgfqpoint{8.378005in}{7.101451in}}{\pgfqpoint{8.369439in}{7.101451in}}%
\pgfpathquadraticcurveto{\pgfqpoint{8.359124in}{7.101451in}}{\pgfqpoint{8.353169in}{7.094876in}}%
\pgfpathquadraticcurveto{\pgfqpoint{8.347215in}{7.088324in}}{\pgfqpoint{8.347215in}{7.076969in}}%
\pgfpathlineto{\pgfqpoint{8.347215in}{7.033185in}}%
\pgfpathlineto{\pgfqpoint{8.334421in}{7.033185in}}%
\pgfpathlineto{\pgfqpoint{8.334421in}{7.110659in}}%
\pgfpathlineto{\pgfqpoint{8.347215in}{7.110659in}}%
\pgfpathlineto{\pgfqpoint{8.347215in}{7.098617in}}%
\pgfpathquadraticcurveto{\pgfqpoint{8.351797in}{7.105612in}}{\pgfqpoint{8.357973in}{7.109065in}}%
\pgfpathquadraticcurveto{\pgfqpoint{8.364171in}{7.112518in}}{\pgfqpoint{8.372272in}{7.112518in}}%
\pgfpathquadraticcurveto{\pgfqpoint{8.385620in}{7.112518in}}{\pgfqpoint{8.392460in}{7.104262in}}%
\pgfpathquadraticcurveto{\pgfqpoint{8.399322in}{7.096005in}}{\pgfqpoint{8.399322in}{7.079957in}}%
\pgfpathlineto{\pgfqpoint{8.399322in}{7.079957in}}%
\pgfpathclose%
\pgfpathmoveto{\pgfqpoint{8.424689in}{7.110659in}}%
\pgfpathlineto{\pgfqpoint{8.437417in}{7.110659in}}%
\pgfpathlineto{\pgfqpoint{8.437417in}{7.033185in}}%
\pgfpathlineto{\pgfqpoint{8.424689in}{7.033185in}}%
\pgfpathlineto{\pgfqpoint{8.424689in}{7.110659in}}%
\pgfpathlineto{\pgfqpoint{8.424689in}{7.110659in}}%
\pgfpathclose%
\pgfpathmoveto{\pgfqpoint{8.424689in}{7.140830in}}%
\pgfpathlineto{\pgfqpoint{8.437417in}{7.140830in}}%
\pgfpathlineto{\pgfqpoint{8.437417in}{7.124693in}}%
\pgfpathlineto{\pgfqpoint{8.424689in}{7.124693in}}%
\pgfpathlineto{\pgfqpoint{8.424689in}{7.140830in}}%
\pgfpathlineto{\pgfqpoint{8.424689in}{7.140830in}}%
\pgfpathclose%
\pgfpathmoveto{\pgfqpoint{8.526645in}{7.110659in}}%
\pgfpathlineto{\pgfqpoint{8.526645in}{7.033185in}}%
\pgfpathlineto{\pgfqpoint{8.513850in}{7.033185in}}%
\pgfpathlineto{\pgfqpoint{8.513850in}{7.100764in}}%
\pgfpathlineto{\pgfqpoint{8.478921in}{7.100764in}}%
\pgfpathlineto{\pgfqpoint{8.478921in}{7.033185in}}%
\pgfpathlineto{\pgfqpoint{8.466126in}{7.033185in}}%
\pgfpathlineto{\pgfqpoint{8.466126in}{7.100764in}}%
\pgfpathlineto{\pgfqpoint{8.453952in}{7.100764in}}%
\pgfpathlineto{\pgfqpoint{8.453952in}{7.110659in}}%
\pgfpathlineto{\pgfqpoint{8.466126in}{7.110659in}}%
\pgfpathlineto{\pgfqpoint{8.466126in}{7.116060in}}%
\pgfpathquadraticcurveto{\pgfqpoint{8.466126in}{7.128721in}}{\pgfqpoint{8.472103in}{7.134764in}}%
\pgfpathquadraticcurveto{\pgfqpoint{8.478102in}{7.140830in}}{\pgfqpoint{8.490475in}{7.140830in}}%
\pgfpathlineto{\pgfqpoint{8.503270in}{7.140830in}}%
\pgfpathlineto{\pgfqpoint{8.503270in}{7.130227in}}%
\pgfpathlineto{\pgfqpoint{8.491095in}{7.130227in}}%
\pgfpathquadraticcurveto{\pgfqpoint{8.484255in}{7.130227in}}{\pgfqpoint{8.481577in}{7.127460in}}%
\pgfpathquadraticcurveto{\pgfqpoint{8.478921in}{7.124693in}}{\pgfqpoint{8.478921in}{7.117499in}}%
\pgfpathlineto{\pgfqpoint{8.478921in}{7.110659in}}%
\pgfpathlineto{\pgfqpoint{8.526645in}{7.110659in}}%
\pgfpathlineto{\pgfqpoint{8.526645in}{7.110659in}}%
\pgfpathclose%
\pgfpathmoveto{\pgfqpoint{8.513850in}{7.140675in}}%
\pgfpathlineto{\pgfqpoint{8.526645in}{7.140675in}}%
\pgfpathlineto{\pgfqpoint{8.526645in}{7.124560in}}%
\pgfpathlineto{\pgfqpoint{8.513850in}{7.124560in}}%
\pgfpathlineto{\pgfqpoint{8.513850in}{7.140675in}}%
\pgfpathlineto{\pgfqpoint{8.513850in}{7.140675in}}%
\pgfpathclose%
\pgfpathmoveto{\pgfqpoint{8.609033in}{7.107693in}}%
\pgfpathlineto{\pgfqpoint{8.609033in}{7.095784in}}%
\pgfpathquadraticcurveto{\pgfqpoint{8.603632in}{7.098772in}}{\pgfqpoint{8.598209in}{7.100255in}}%
\pgfpathquadraticcurveto{\pgfqpoint{8.592785in}{7.101738in}}{\pgfqpoint{8.587251in}{7.101738in}}%
\pgfpathquadraticcurveto{\pgfqpoint{8.574856in}{7.101738in}}{\pgfqpoint{8.567994in}{7.093880in}}%
\pgfpathquadraticcurveto{\pgfqpoint{8.561154in}{7.086044in}}{\pgfqpoint{8.561154in}{7.071856in}}%
\pgfpathquadraticcurveto{\pgfqpoint{8.561154in}{7.057667in}}{\pgfqpoint{8.567994in}{7.049809in}}%
\pgfpathquadraticcurveto{\pgfqpoint{8.574856in}{7.041973in}}{\pgfqpoint{8.587251in}{7.041973in}}%
\pgfpathquadraticcurveto{\pgfqpoint{8.592785in}{7.041973in}}{\pgfqpoint{8.598209in}{7.043456in}}%
\pgfpathquadraticcurveto{\pgfqpoint{8.603632in}{7.044939in}}{\pgfqpoint{8.609033in}{7.047927in}}%
\pgfpathlineto{\pgfqpoint{8.609033in}{7.036151in}}%
\pgfpathquadraticcurveto{\pgfqpoint{8.603698in}{7.033672in}}{\pgfqpoint{8.597987in}{7.032432in}}%
\pgfpathquadraticcurveto{\pgfqpoint{8.592298in}{7.031171in}}{\pgfqpoint{8.585857in}{7.031171in}}%
\pgfpathquadraticcurveto{\pgfqpoint{8.568348in}{7.031171in}}{\pgfqpoint{8.558033in}{7.042172in}}%
\pgfpathquadraticcurveto{\pgfqpoint{8.547740in}{7.053173in}}{\pgfqpoint{8.547740in}{7.071856in}}%
\pgfpathquadraticcurveto{\pgfqpoint{8.547740in}{7.090803in}}{\pgfqpoint{8.558143in}{7.101650in}}%
\pgfpathquadraticcurveto{\pgfqpoint{8.568569in}{7.112518in}}{\pgfqpoint{8.586698in}{7.112518in}}%
\pgfpathquadraticcurveto{\pgfqpoint{8.592564in}{7.112518in}}{\pgfqpoint{8.598164in}{7.111301in}}%
\pgfpathquadraticcurveto{\pgfqpoint{8.603765in}{7.110106in}}{\pgfqpoint{8.609033in}{7.107693in}}%
\pgfpathlineto{\pgfqpoint{8.609033in}{7.107693in}}%
\pgfpathclose%
\pgfpathmoveto{\pgfqpoint{8.666386in}{7.072121in}}%
\pgfpathquadraticcurveto{\pgfqpoint{8.650957in}{7.072121in}}{\pgfqpoint{8.645003in}{7.068602in}}%
\pgfpathquadraticcurveto{\pgfqpoint{8.639048in}{7.065082in}}{\pgfqpoint{8.639048in}{7.056560in}}%
\pgfpathquadraticcurveto{\pgfqpoint{8.639048in}{7.049787in}}{\pgfqpoint{8.643520in}{7.045802in}}%
\pgfpathquadraticcurveto{\pgfqpoint{8.647991in}{7.041840in}}{\pgfqpoint{8.655650in}{7.041840in}}%
\pgfpathquadraticcurveto{\pgfqpoint{8.666253in}{7.041840in}}{\pgfqpoint{8.672650in}{7.049344in}}%
\pgfpathquadraticcurveto{\pgfqpoint{8.679047in}{7.056848in}}{\pgfqpoint{8.679047in}{7.069288in}}%
\pgfpathlineto{\pgfqpoint{8.679047in}{7.072121in}}%
\pgfpathlineto{\pgfqpoint{8.666386in}{7.072121in}}%
\pgfpathlineto{\pgfqpoint{8.666386in}{7.072121in}}%
\pgfpathclose%
\pgfpathmoveto{\pgfqpoint{8.691775in}{7.077389in}}%
\pgfpathlineto{\pgfqpoint{8.691775in}{7.033185in}}%
\pgfpathlineto{\pgfqpoint{8.679047in}{7.033185in}}%
\pgfpathlineto{\pgfqpoint{8.679047in}{7.044939in}}%
\pgfpathquadraticcurveto{\pgfqpoint{8.674686in}{7.037900in}}{\pgfqpoint{8.668179in}{7.034535in}}%
\pgfpathquadraticcurveto{\pgfqpoint{8.661671in}{7.031171in}}{\pgfqpoint{8.652263in}{7.031171in}}%
\pgfpathquadraticcurveto{\pgfqpoint{8.640376in}{7.031171in}}{\pgfqpoint{8.633337in}{7.037856in}}%
\pgfpathquadraticcurveto{\pgfqpoint{8.626320in}{7.044540in}}{\pgfqpoint{8.626320in}{7.055741in}}%
\pgfpathquadraticcurveto{\pgfqpoint{8.626320in}{7.068801in}}{\pgfqpoint{8.635064in}{7.075442in}}%
\pgfpathquadraticcurveto{\pgfqpoint{8.643830in}{7.082082in}}{\pgfqpoint{8.661184in}{7.082082in}}%
\pgfpathlineto{\pgfqpoint{8.679047in}{7.082082in}}%
\pgfpathlineto{\pgfqpoint{8.679047in}{7.083344in}}%
\pgfpathquadraticcurveto{\pgfqpoint{8.679047in}{7.092132in}}{\pgfqpoint{8.673270in}{7.096935in}}%
\pgfpathquadraticcurveto{\pgfqpoint{8.667492in}{7.101738in}}{\pgfqpoint{8.657044in}{7.101738in}}%
\pgfpathquadraticcurveto{\pgfqpoint{8.650404in}{7.101738in}}{\pgfqpoint{8.644095in}{7.100145in}}%
\pgfpathquadraticcurveto{\pgfqpoint{8.637809in}{7.098551in}}{\pgfqpoint{8.632009in}{7.095363in}}%
\pgfpathlineto{\pgfqpoint{8.632009in}{7.107139in}}%
\pgfpathquadraticcurveto{\pgfqpoint{8.638982in}{7.109840in}}{\pgfqpoint{8.645556in}{7.111168in}}%
\pgfpathquadraticcurveto{\pgfqpoint{8.652130in}{7.112518in}}{\pgfqpoint{8.658350in}{7.112518in}}%
\pgfpathquadraticcurveto{\pgfqpoint{8.675173in}{7.112518in}}{\pgfqpoint{8.683474in}{7.103797in}}%
\pgfpathquadraticcurveto{\pgfqpoint{8.691775in}{7.095098in}}{\pgfqpoint{8.691775in}{7.077389in}}%
\pgfpathlineto{\pgfqpoint{8.691775in}{7.077389in}}%
\pgfpathclose%
\pgfpathmoveto{\pgfqpoint{8.782397in}{7.079957in}}%
\pgfpathlineto{\pgfqpoint{8.782397in}{7.033185in}}%
\pgfpathlineto{\pgfqpoint{8.769669in}{7.033185in}}%
\pgfpathlineto{\pgfqpoint{8.769669in}{7.079537in}}%
\pgfpathquadraticcurveto{\pgfqpoint{8.769669in}{7.090538in}}{\pgfqpoint{8.765375in}{7.095983in}}%
\pgfpathquadraticcurveto{\pgfqpoint{8.761081in}{7.101451in}}{\pgfqpoint{8.752515in}{7.101451in}}%
\pgfpathquadraticcurveto{\pgfqpoint{8.742199in}{7.101451in}}{\pgfqpoint{8.736245in}{7.094876in}}%
\pgfpathquadraticcurveto{\pgfqpoint{8.730291in}{7.088324in}}{\pgfqpoint{8.730291in}{7.076969in}}%
\pgfpathlineto{\pgfqpoint{8.730291in}{7.033185in}}%
\pgfpathlineto{\pgfqpoint{8.717496in}{7.033185in}}%
\pgfpathlineto{\pgfqpoint{8.717496in}{7.110659in}}%
\pgfpathlineto{\pgfqpoint{8.730291in}{7.110659in}}%
\pgfpathlineto{\pgfqpoint{8.730291in}{7.098617in}}%
\pgfpathquadraticcurveto{\pgfqpoint{8.734873in}{7.105612in}}{\pgfqpoint{8.741048in}{7.109065in}}%
\pgfpathquadraticcurveto{\pgfqpoint{8.747246in}{7.112518in}}{\pgfqpoint{8.755348in}{7.112518in}}%
\pgfpathquadraticcurveto{\pgfqpoint{8.768695in}{7.112518in}}{\pgfqpoint{8.775535in}{7.104262in}}%
\pgfpathquadraticcurveto{\pgfqpoint{8.782397in}{7.096005in}}{\pgfqpoint{8.782397in}{7.079957in}}%
\pgfpathlineto{\pgfqpoint{8.782397in}{7.079957in}}%
\pgfpathclose%
\pgfpathmoveto{\pgfqpoint{8.820360in}{7.132662in}}%
\pgfpathlineto{\pgfqpoint{8.820360in}{7.110659in}}%
\pgfpathlineto{\pgfqpoint{8.846568in}{7.110659in}}%
\pgfpathlineto{\pgfqpoint{8.846568in}{7.100764in}}%
\pgfpathlineto{\pgfqpoint{8.820360in}{7.100764in}}%
\pgfpathlineto{\pgfqpoint{8.820360in}{7.058707in}}%
\pgfpathquadraticcurveto{\pgfqpoint{8.820360in}{7.049233in}}{\pgfqpoint{8.822949in}{7.046533in}}%
\pgfpathquadraticcurveto{\pgfqpoint{8.825539in}{7.043832in}}{\pgfqpoint{8.833508in}{7.043832in}}%
\pgfpathlineto{\pgfqpoint{8.846568in}{7.043832in}}%
\pgfpathlineto{\pgfqpoint{8.846568in}{7.033185in}}%
\pgfpathlineto{\pgfqpoint{8.833508in}{7.033185in}}%
\pgfpathquadraticcurveto{\pgfqpoint{8.818766in}{7.033185in}}{\pgfqpoint{8.813166in}{7.038675in}}%
\pgfpathquadraticcurveto{\pgfqpoint{8.807565in}{7.044186in}}{\pgfqpoint{8.807565in}{7.058707in}}%
\pgfpathlineto{\pgfqpoint{8.807565in}{7.100764in}}%
\pgfpathlineto{\pgfqpoint{8.798224in}{7.100764in}}%
\pgfpathlineto{\pgfqpoint{8.798224in}{7.110659in}}%
\pgfpathlineto{\pgfqpoint{8.807565in}{7.110659in}}%
\pgfpathlineto{\pgfqpoint{8.807565in}{7.132662in}}%
\pgfpathlineto{\pgfqpoint{8.820360in}{7.132662in}}%
\pgfpathlineto{\pgfqpoint{8.820360in}{7.132662in}}%
\pgfpathclose%
\pgfpathmoveto{\pgfqpoint{8.959304in}{7.098905in}}%
\pgfpathlineto{\pgfqpoint{8.959304in}{7.140830in}}%
\pgfpathlineto{\pgfqpoint{8.972031in}{7.140830in}}%
\pgfpathlineto{\pgfqpoint{8.972031in}{7.033185in}}%
\pgfpathlineto{\pgfqpoint{8.959304in}{7.033185in}}%
\pgfpathlineto{\pgfqpoint{8.959304in}{7.044806in}}%
\pgfpathquadraticcurveto{\pgfqpoint{8.955297in}{7.037900in}}{\pgfqpoint{8.949166in}{7.034535in}}%
\pgfpathquadraticcurveto{\pgfqpoint{8.943056in}{7.031171in}}{\pgfqpoint{8.934468in}{7.031171in}}%
\pgfpathquadraticcurveto{\pgfqpoint{8.920434in}{7.031171in}}{\pgfqpoint{8.911602in}{7.042371in}}%
\pgfpathquadraticcurveto{\pgfqpoint{8.902792in}{7.053594in}}{\pgfqpoint{8.902792in}{7.071856in}}%
\pgfpathquadraticcurveto{\pgfqpoint{8.902792in}{7.090117in}}{\pgfqpoint{8.911602in}{7.101318in}}%
\pgfpathquadraticcurveto{\pgfqpoint{8.920434in}{7.112518in}}{\pgfqpoint{8.934468in}{7.112518in}}%
\pgfpathquadraticcurveto{\pgfqpoint{8.943056in}{7.112518in}}{\pgfqpoint{8.949166in}{7.109154in}}%
\pgfpathquadraticcurveto{\pgfqpoint{8.955297in}{7.105811in}}{\pgfqpoint{8.959304in}{7.098905in}}%
\pgfpathlineto{\pgfqpoint{8.959304in}{7.098905in}}%
\pgfpathclose%
\pgfpathmoveto{\pgfqpoint{8.915940in}{7.071856in}}%
\pgfpathquadraticcurveto{\pgfqpoint{8.915940in}{7.057822in}}{\pgfqpoint{8.921718in}{7.049831in}}%
\pgfpathquadraticcurveto{\pgfqpoint{8.927495in}{7.041840in}}{\pgfqpoint{8.937589in}{7.041840in}}%
\pgfpathquadraticcurveto{\pgfqpoint{8.947682in}{7.041840in}}{\pgfqpoint{8.953482in}{7.049831in}}%
\pgfpathquadraticcurveto{\pgfqpoint{8.959304in}{7.057822in}}{\pgfqpoint{8.959304in}{7.071856in}}%
\pgfpathquadraticcurveto{\pgfqpoint{8.959304in}{7.085889in}}{\pgfqpoint{8.953482in}{7.093880in}}%
\pgfpathquadraticcurveto{\pgfqpoint{8.947682in}{7.101871in}}{\pgfqpoint{8.937589in}{7.101871in}}%
\pgfpathquadraticcurveto{\pgfqpoint{8.927495in}{7.101871in}}{\pgfqpoint{8.921718in}{7.093880in}}%
\pgfpathquadraticcurveto{\pgfqpoint{8.915940in}{7.085889in}}{\pgfqpoint{8.915940in}{7.071856in}}%
\pgfpathlineto{\pgfqpoint{8.915940in}{7.071856in}}%
\pgfpathclose%
\pgfpathmoveto{\pgfqpoint{8.998262in}{7.110659in}}%
\pgfpathlineto{\pgfqpoint{9.010990in}{7.110659in}}%
\pgfpathlineto{\pgfqpoint{9.010990in}{7.033185in}}%
\pgfpathlineto{\pgfqpoint{8.998262in}{7.033185in}}%
\pgfpathlineto{\pgfqpoint{8.998262in}{7.110659in}}%
\pgfpathlineto{\pgfqpoint{8.998262in}{7.110659in}}%
\pgfpathclose%
\pgfpathmoveto{\pgfqpoint{8.998262in}{7.140830in}}%
\pgfpathlineto{\pgfqpoint{9.010990in}{7.140830in}}%
\pgfpathlineto{\pgfqpoint{9.010990in}{7.124693in}}%
\pgfpathlineto{\pgfqpoint{8.998262in}{7.124693in}}%
\pgfpathlineto{\pgfqpoint{8.998262in}{7.140830in}}%
\pgfpathlineto{\pgfqpoint{8.998262in}{7.140830in}}%
\pgfpathclose%
\pgfpathmoveto{\pgfqpoint{9.124567in}{7.140830in}}%
\pgfpathlineto{\pgfqpoint{9.124567in}{7.130227in}}%
\pgfpathlineto{\pgfqpoint{9.112392in}{7.130227in}}%
\pgfpathquadraticcurveto{\pgfqpoint{9.105552in}{7.130227in}}{\pgfqpoint{9.102874in}{7.127460in}}%
\pgfpathquadraticcurveto{\pgfqpoint{9.100218in}{7.124693in}}{\pgfqpoint{9.100218in}{7.117499in}}%
\pgfpathlineto{\pgfqpoint{9.100218in}{7.110659in}}%
\pgfpathlineto{\pgfqpoint{9.121180in}{7.110659in}}%
\pgfpathlineto{\pgfqpoint{9.121180in}{7.100764in}}%
\pgfpathlineto{\pgfqpoint{9.100218in}{7.100764in}}%
\pgfpathlineto{\pgfqpoint{9.100218in}{7.033185in}}%
\pgfpathlineto{\pgfqpoint{9.087423in}{7.033185in}}%
\pgfpathlineto{\pgfqpoint{9.087423in}{7.100764in}}%
\pgfpathlineto{\pgfqpoint{9.052494in}{7.100764in}}%
\pgfpathlineto{\pgfqpoint{9.052494in}{7.033185in}}%
\pgfpathlineto{\pgfqpoint{9.039699in}{7.033185in}}%
\pgfpathlineto{\pgfqpoint{9.039699in}{7.100764in}}%
\pgfpathlineto{\pgfqpoint{9.027525in}{7.100764in}}%
\pgfpathlineto{\pgfqpoint{9.027525in}{7.110659in}}%
\pgfpathlineto{\pgfqpoint{9.039699in}{7.110659in}}%
\pgfpathlineto{\pgfqpoint{9.039699in}{7.116060in}}%
\pgfpathquadraticcurveto{\pgfqpoint{9.039699in}{7.128987in}}{\pgfqpoint{9.045720in}{7.134897in}}%
\pgfpathquadraticcurveto{\pgfqpoint{9.051741in}{7.140830in}}{\pgfqpoint{9.064801in}{7.140830in}}%
\pgfpathlineto{\pgfqpoint{9.076843in}{7.140830in}}%
\pgfpathlineto{\pgfqpoint{9.076843in}{7.130227in}}%
\pgfpathlineto{\pgfqpoint{9.064668in}{7.130227in}}%
\pgfpathquadraticcurveto{\pgfqpoint{9.057828in}{7.130227in}}{\pgfqpoint{9.055128in}{7.127460in}}%
\pgfpathquadraticcurveto{\pgfqpoint{9.052494in}{7.124693in}}{\pgfqpoint{9.052494in}{7.117499in}}%
\pgfpathlineto{\pgfqpoint{9.052494in}{7.110659in}}%
\pgfpathlineto{\pgfqpoint{9.087423in}{7.110659in}}%
\pgfpathlineto{\pgfqpoint{9.087423in}{7.116060in}}%
\pgfpathquadraticcurveto{\pgfqpoint{9.087423in}{7.128987in}}{\pgfqpoint{9.093444in}{7.134897in}}%
\pgfpathquadraticcurveto{\pgfqpoint{9.099465in}{7.140830in}}{\pgfqpoint{9.112547in}{7.140830in}}%
\pgfpathlineto{\pgfqpoint{9.124567in}{7.140830in}}%
\pgfpathlineto{\pgfqpoint{9.124567in}{7.140830in}}%
\pgfpathclose%
\pgfpathmoveto{\pgfqpoint{9.201487in}{7.075109in}}%
\pgfpathlineto{\pgfqpoint{9.201487in}{7.068889in}}%
\pgfpathlineto{\pgfqpoint{9.142961in}{7.068889in}}%
\pgfpathquadraticcurveto{\pgfqpoint{9.143802in}{7.055741in}}{\pgfqpoint{9.150886in}{7.048857in}}%
\pgfpathquadraticcurveto{\pgfqpoint{9.157969in}{7.041973in}}{\pgfqpoint{9.170630in}{7.041973in}}%
\pgfpathquadraticcurveto{\pgfqpoint{9.177957in}{7.041973in}}{\pgfqpoint{9.184841in}{7.043766in}}%
\pgfpathquadraticcurveto{\pgfqpoint{9.191725in}{7.045559in}}{\pgfqpoint{9.198521in}{7.049167in}}%
\pgfpathlineto{\pgfqpoint{9.198521in}{7.037125in}}%
\pgfpathquadraticcurveto{\pgfqpoint{9.191659in}{7.034225in}}{\pgfqpoint{9.184465in}{7.032698in}}%
\pgfpathquadraticcurveto{\pgfqpoint{9.177271in}{7.031171in}}{\pgfqpoint{9.169878in}{7.031171in}}%
\pgfpathquadraticcurveto{\pgfqpoint{9.151328in}{7.031171in}}{\pgfqpoint{9.140504in}{7.041951in}}%
\pgfpathquadraticcurveto{\pgfqpoint{9.129680in}{7.052753in}}{\pgfqpoint{9.129680in}{7.071169in}}%
\pgfpathquadraticcurveto{\pgfqpoint{9.129680in}{7.090184in}}{\pgfqpoint{9.139951in}{7.101340in}}%
\pgfpathquadraticcurveto{\pgfqpoint{9.150222in}{7.112518in}}{\pgfqpoint{9.167664in}{7.112518in}}%
\pgfpathquadraticcurveto{\pgfqpoint{9.183292in}{7.112518in}}{\pgfqpoint{9.192390in}{7.102447in}}%
\pgfpathquadraticcurveto{\pgfqpoint{9.201487in}{7.092397in}}{\pgfqpoint{9.201487in}{7.075109in}}%
\pgfpathlineto{\pgfqpoint{9.201487in}{7.075109in}}%
\pgfpathclose%
\pgfpathmoveto{\pgfqpoint{9.188759in}{7.078850in}}%
\pgfpathquadraticcurveto{\pgfqpoint{9.188626in}{7.089276in}}{\pgfqpoint{9.182916in}{7.095496in}}%
\pgfpathquadraticcurveto{\pgfqpoint{9.177205in}{7.101738in}}{\pgfqpoint{9.167797in}{7.101738in}}%
\pgfpathquadraticcurveto{\pgfqpoint{9.157150in}{7.101738in}}{\pgfqpoint{9.150753in}{7.095718in}}%
\pgfpathquadraticcurveto{\pgfqpoint{9.144356in}{7.089697in}}{\pgfqpoint{9.143382in}{7.078762in}}%
\pgfpathlineto{\pgfqpoint{9.188759in}{7.078850in}}%
\pgfpathlineto{\pgfqpoint{9.188759in}{7.078850in}}%
\pgfpathclose%
\pgfpathmoveto{\pgfqpoint{9.267274in}{7.098772in}}%
\pgfpathquadraticcurveto{\pgfqpoint{9.265126in}{7.100012in}}{\pgfqpoint{9.262603in}{7.100587in}}%
\pgfpathquadraticcurveto{\pgfqpoint{9.260080in}{7.101185in}}{\pgfqpoint{9.257047in}{7.101185in}}%
\pgfpathquadraticcurveto{\pgfqpoint{9.246245in}{7.101185in}}{\pgfqpoint{9.240468in}{7.094168in}}%
\pgfpathquadraticcurveto{\pgfqpoint{9.234690in}{7.087151in}}{\pgfqpoint{9.234690in}{7.074003in}}%
\pgfpathlineto{\pgfqpoint{9.234690in}{7.033185in}}%
\pgfpathlineto{\pgfqpoint{9.221896in}{7.033185in}}%
\pgfpathlineto{\pgfqpoint{9.221896in}{7.110659in}}%
\pgfpathlineto{\pgfqpoint{9.234690in}{7.110659in}}%
\pgfpathlineto{\pgfqpoint{9.234690in}{7.098617in}}%
\pgfpathquadraticcurveto{\pgfqpoint{9.238719in}{7.105678in}}{\pgfqpoint{9.245138in}{7.109087in}}%
\pgfpathquadraticcurveto{\pgfqpoint{9.251580in}{7.112518in}}{\pgfqpoint{9.260788in}{7.112518in}}%
\pgfpathquadraticcurveto{\pgfqpoint{9.262094in}{7.112518in}}{\pgfqpoint{9.263688in}{7.112341in}}%
\pgfpathquadraticcurveto{\pgfqpoint{9.265281in}{7.112186in}}{\pgfqpoint{9.267207in}{7.111832in}}%
\pgfpathlineto{\pgfqpoint{9.267274in}{7.098772in}}%
\pgfpathlineto{\pgfqpoint{9.267274in}{7.098772in}}%
\pgfpathclose%
\pgfpathmoveto{\pgfqpoint{9.343774in}{7.075109in}}%
\pgfpathlineto{\pgfqpoint{9.343774in}{7.068889in}}%
\pgfpathlineto{\pgfqpoint{9.285248in}{7.068889in}}%
\pgfpathquadraticcurveto{\pgfqpoint{9.286089in}{7.055741in}}{\pgfqpoint{9.293172in}{7.048857in}}%
\pgfpathquadraticcurveto{\pgfqpoint{9.300255in}{7.041973in}}{\pgfqpoint{9.312917in}{7.041973in}}%
\pgfpathquadraticcurveto{\pgfqpoint{9.320244in}{7.041973in}}{\pgfqpoint{9.327128in}{7.043766in}}%
\pgfpathquadraticcurveto{\pgfqpoint{9.334012in}{7.045559in}}{\pgfqpoint{9.340807in}{7.049167in}}%
\pgfpathlineto{\pgfqpoint{9.340807in}{7.037125in}}%
\pgfpathquadraticcurveto{\pgfqpoint{9.333945in}{7.034225in}}{\pgfqpoint{9.326751in}{7.032698in}}%
\pgfpathquadraticcurveto{\pgfqpoint{9.319557in}{7.031171in}}{\pgfqpoint{9.312164in}{7.031171in}}%
\pgfpathquadraticcurveto{\pgfqpoint{9.293615in}{7.031171in}}{\pgfqpoint{9.282791in}{7.041951in}}%
\pgfpathquadraticcurveto{\pgfqpoint{9.271966in}{7.052753in}}{\pgfqpoint{9.271966in}{7.071169in}}%
\pgfpathquadraticcurveto{\pgfqpoint{9.271966in}{7.090184in}}{\pgfqpoint{9.282237in}{7.101340in}}%
\pgfpathquadraticcurveto{\pgfqpoint{9.292508in}{7.112518in}}{\pgfqpoint{9.309951in}{7.112518in}}%
\pgfpathquadraticcurveto{\pgfqpoint{9.325578in}{7.112518in}}{\pgfqpoint{9.334676in}{7.102447in}}%
\pgfpathquadraticcurveto{\pgfqpoint{9.343774in}{7.092397in}}{\pgfqpoint{9.343774in}{7.075109in}}%
\pgfpathlineto{\pgfqpoint{9.343774in}{7.075109in}}%
\pgfpathclose%
\pgfpathmoveto{\pgfqpoint{9.331046in}{7.078850in}}%
\pgfpathquadraticcurveto{\pgfqpoint{9.330913in}{7.089276in}}{\pgfqpoint{9.325202in}{7.095496in}}%
\pgfpathquadraticcurveto{\pgfqpoint{9.319491in}{7.101738in}}{\pgfqpoint{9.310084in}{7.101738in}}%
\pgfpathquadraticcurveto{\pgfqpoint{9.299436in}{7.101738in}}{\pgfqpoint{9.293039in}{7.095718in}}%
\pgfpathquadraticcurveto{\pgfqpoint{9.286642in}{7.089697in}}{\pgfqpoint{9.285668in}{7.078762in}}%
\pgfpathlineto{\pgfqpoint{9.331046in}{7.078850in}}%
\pgfpathlineto{\pgfqpoint{9.331046in}{7.078850in}}%
\pgfpathclose%
\pgfpathmoveto{\pgfqpoint{9.429084in}{7.079957in}}%
\pgfpathlineto{\pgfqpoint{9.429084in}{7.033185in}}%
\pgfpathlineto{\pgfqpoint{9.416356in}{7.033185in}}%
\pgfpathlineto{\pgfqpoint{9.416356in}{7.079537in}}%
\pgfpathquadraticcurveto{\pgfqpoint{9.416356in}{7.090538in}}{\pgfqpoint{9.412061in}{7.095983in}}%
\pgfpathquadraticcurveto{\pgfqpoint{9.407767in}{7.101451in}}{\pgfqpoint{9.399201in}{7.101451in}}%
\pgfpathquadraticcurveto{\pgfqpoint{9.388886in}{7.101451in}}{\pgfqpoint{9.382931in}{7.094876in}}%
\pgfpathquadraticcurveto{\pgfqpoint{9.376977in}{7.088324in}}{\pgfqpoint{9.376977in}{7.076969in}}%
\pgfpathlineto{\pgfqpoint{9.376977in}{7.033185in}}%
\pgfpathlineto{\pgfqpoint{9.364182in}{7.033185in}}%
\pgfpathlineto{\pgfqpoint{9.364182in}{7.110659in}}%
\pgfpathlineto{\pgfqpoint{9.376977in}{7.110659in}}%
\pgfpathlineto{\pgfqpoint{9.376977in}{7.098617in}}%
\pgfpathquadraticcurveto{\pgfqpoint{9.381559in}{7.105612in}}{\pgfqpoint{9.387735in}{7.109065in}}%
\pgfpathquadraticcurveto{\pgfqpoint{9.393932in}{7.112518in}}{\pgfqpoint{9.402034in}{7.112518in}}%
\pgfpathquadraticcurveto{\pgfqpoint{9.415382in}{7.112518in}}{\pgfqpoint{9.422222in}{7.104262in}}%
\pgfpathquadraticcurveto{\pgfqpoint{9.429084in}{7.096005in}}{\pgfqpoint{9.429084in}{7.079957in}}%
\pgfpathlineto{\pgfqpoint{9.429084in}{7.079957in}}%
\pgfpathclose%
\pgfpathmoveto{\pgfqpoint{9.510210in}{7.107693in}}%
\pgfpathlineto{\pgfqpoint{9.510210in}{7.095784in}}%
\pgfpathquadraticcurveto{\pgfqpoint{9.504809in}{7.098772in}}{\pgfqpoint{9.499386in}{7.100255in}}%
\pgfpathquadraticcurveto{\pgfqpoint{9.493962in}{7.101738in}}{\pgfqpoint{9.488429in}{7.101738in}}%
\pgfpathquadraticcurveto{\pgfqpoint{9.476033in}{7.101738in}}{\pgfqpoint{9.469171in}{7.093880in}}%
\pgfpathquadraticcurveto{\pgfqpoint{9.462331in}{7.086044in}}{\pgfqpoint{9.462331in}{7.071856in}}%
\pgfpathquadraticcurveto{\pgfqpoint{9.462331in}{7.057667in}}{\pgfqpoint{9.469171in}{7.049809in}}%
\pgfpathquadraticcurveto{\pgfqpoint{9.476033in}{7.041973in}}{\pgfqpoint{9.488429in}{7.041973in}}%
\pgfpathquadraticcurveto{\pgfqpoint{9.493962in}{7.041973in}}{\pgfqpoint{9.499386in}{7.043456in}}%
\pgfpathquadraticcurveto{\pgfqpoint{9.504809in}{7.044939in}}{\pgfqpoint{9.510210in}{7.047927in}}%
\pgfpathlineto{\pgfqpoint{9.510210in}{7.036151in}}%
\pgfpathquadraticcurveto{\pgfqpoint{9.504875in}{7.033672in}}{\pgfqpoint{9.499164in}{7.032432in}}%
\pgfpathquadraticcurveto{\pgfqpoint{9.493475in}{7.031171in}}{\pgfqpoint{9.487034in}{7.031171in}}%
\pgfpathquadraticcurveto{\pgfqpoint{9.469525in}{7.031171in}}{\pgfqpoint{9.459210in}{7.042172in}}%
\pgfpathquadraticcurveto{\pgfqpoint{9.448917in}{7.053173in}}{\pgfqpoint{9.448917in}{7.071856in}}%
\pgfpathquadraticcurveto{\pgfqpoint{9.448917in}{7.090803in}}{\pgfqpoint{9.459320in}{7.101650in}}%
\pgfpathquadraticcurveto{\pgfqpoint{9.469746in}{7.112518in}}{\pgfqpoint{9.487875in}{7.112518in}}%
\pgfpathquadraticcurveto{\pgfqpoint{9.493741in}{7.112518in}}{\pgfqpoint{9.499341in}{7.111301in}}%
\pgfpathquadraticcurveto{\pgfqpoint{9.504942in}{7.110106in}}{\pgfqpoint{9.510210in}{7.107693in}}%
\pgfpathlineto{\pgfqpoint{9.510210in}{7.107693in}}%
\pgfpathclose%
\pgfpathmoveto{\pgfqpoint{9.598619in}{7.075109in}}%
\pgfpathlineto{\pgfqpoint{9.598619in}{7.068889in}}%
\pgfpathlineto{\pgfqpoint{9.540093in}{7.068889in}}%
\pgfpathquadraticcurveto{\pgfqpoint{9.540934in}{7.055741in}}{\pgfqpoint{9.548017in}{7.048857in}}%
\pgfpathquadraticcurveto{\pgfqpoint{9.555100in}{7.041973in}}{\pgfqpoint{9.567762in}{7.041973in}}%
\pgfpathquadraticcurveto{\pgfqpoint{9.575089in}{7.041973in}}{\pgfqpoint{9.581973in}{7.043766in}}%
\pgfpathquadraticcurveto{\pgfqpoint{9.588857in}{7.045559in}}{\pgfqpoint{9.595653in}{7.049167in}}%
\pgfpathlineto{\pgfqpoint{9.595653in}{7.037125in}}%
\pgfpathquadraticcurveto{\pgfqpoint{9.588791in}{7.034225in}}{\pgfqpoint{9.581597in}{7.032698in}}%
\pgfpathquadraticcurveto{\pgfqpoint{9.574403in}{7.031171in}}{\pgfqpoint{9.567009in}{7.031171in}}%
\pgfpathquadraticcurveto{\pgfqpoint{9.548460in}{7.031171in}}{\pgfqpoint{9.537636in}{7.041951in}}%
\pgfpathquadraticcurveto{\pgfqpoint{9.526811in}{7.052753in}}{\pgfqpoint{9.526811in}{7.071169in}}%
\pgfpathquadraticcurveto{\pgfqpoint{9.526811in}{7.090184in}}{\pgfqpoint{9.537082in}{7.101340in}}%
\pgfpathquadraticcurveto{\pgfqpoint{9.547353in}{7.112518in}}{\pgfqpoint{9.564796in}{7.112518in}}%
\pgfpathquadraticcurveto{\pgfqpoint{9.580423in}{7.112518in}}{\pgfqpoint{9.589521in}{7.102447in}}%
\pgfpathquadraticcurveto{\pgfqpoint{9.598619in}{7.092397in}}{\pgfqpoint{9.598619in}{7.075109in}}%
\pgfpathlineto{\pgfqpoint{9.598619in}{7.075109in}}%
\pgfpathclose%
\pgfpathmoveto{\pgfqpoint{9.585891in}{7.078850in}}%
\pgfpathquadraticcurveto{\pgfqpoint{9.585758in}{7.089276in}}{\pgfqpoint{9.580047in}{7.095496in}}%
\pgfpathquadraticcurveto{\pgfqpoint{9.574336in}{7.101738in}}{\pgfqpoint{9.564929in}{7.101738in}}%
\pgfpathquadraticcurveto{\pgfqpoint{9.554281in}{7.101738in}}{\pgfqpoint{9.547884in}{7.095718in}}%
\pgfpathquadraticcurveto{\pgfqpoint{9.541487in}{7.089697in}}{\pgfqpoint{9.540513in}{7.078762in}}%
\pgfpathlineto{\pgfqpoint{9.585891in}{7.078850in}}%
\pgfpathlineto{\pgfqpoint{9.585891in}{7.078850in}}%
\pgfpathclose%
\pgfpathmoveto{\pgfqpoint{9.668899in}{7.108379in}}%
\pgfpathlineto{\pgfqpoint{9.668899in}{7.096337in}}%
\pgfpathquadraticcurveto{\pgfqpoint{9.663520in}{7.099104in}}{\pgfqpoint{9.657698in}{7.100477in}}%
\pgfpathquadraticcurveto{\pgfqpoint{9.651899in}{7.101871in}}{\pgfqpoint{9.645656in}{7.101871in}}%
\pgfpathquadraticcurveto{\pgfqpoint{9.636182in}{7.101871in}}{\pgfqpoint{9.631445in}{7.098971in}}%
\pgfpathquadraticcurveto{\pgfqpoint{9.626709in}{7.096072in}}{\pgfqpoint{9.626709in}{7.090250in}}%
\pgfpathquadraticcurveto{\pgfqpoint{9.626709in}{7.085823in}}{\pgfqpoint{9.630095in}{7.083300in}}%
\pgfpathquadraticcurveto{\pgfqpoint{9.633482in}{7.080776in}}{\pgfqpoint{9.643731in}{7.078496in}}%
\pgfpathlineto{\pgfqpoint{9.648091in}{7.077522in}}%
\pgfpathquadraticcurveto{\pgfqpoint{9.661638in}{7.074623in}}{\pgfqpoint{9.667349in}{7.069332in}}%
\pgfpathquadraticcurveto{\pgfqpoint{9.673060in}{7.064042in}}{\pgfqpoint{9.673060in}{7.054568in}}%
\pgfpathquadraticcurveto{\pgfqpoint{9.673060in}{7.043766in}}{\pgfqpoint{9.664516in}{7.037457in}}%
\pgfpathquadraticcurveto{\pgfqpoint{9.655972in}{7.031171in}}{\pgfqpoint{9.641030in}{7.031171in}}%
\pgfpathquadraticcurveto{\pgfqpoint{9.634810in}{7.031171in}}{\pgfqpoint{9.628059in}{7.032388in}}%
\pgfpathquadraticcurveto{\pgfqpoint{9.621307in}{7.033606in}}{\pgfqpoint{9.613848in}{7.036018in}}%
\pgfpathlineto{\pgfqpoint{9.613848in}{7.049167in}}%
\pgfpathquadraticcurveto{\pgfqpoint{9.620909in}{7.045492in}}{\pgfqpoint{9.627749in}{7.043655in}}%
\pgfpathquadraticcurveto{\pgfqpoint{9.634589in}{7.041840in}}{\pgfqpoint{9.641318in}{7.041840in}}%
\pgfpathquadraticcurveto{\pgfqpoint{9.650305in}{7.041840in}}{\pgfqpoint{9.655130in}{7.044917in}}%
\pgfpathquadraticcurveto{\pgfqpoint{9.659978in}{7.047994in}}{\pgfqpoint{9.659978in}{7.053594in}}%
\pgfpathquadraticcurveto{\pgfqpoint{9.659978in}{7.058774in}}{\pgfqpoint{9.656481in}{7.061540in}}%
\pgfpathquadraticcurveto{\pgfqpoint{9.653005in}{7.064307in}}{\pgfqpoint{9.641163in}{7.066875in}}%
\pgfpathlineto{\pgfqpoint{9.636736in}{7.067915in}}%
\pgfpathquadraticcurveto{\pgfqpoint{9.624916in}{7.070395in}}{\pgfqpoint{9.619647in}{7.075552in}}%
\pgfpathquadraticcurveto{\pgfqpoint{9.614401in}{7.080710in}}{\pgfqpoint{9.614401in}{7.089697in}}%
\pgfpathquadraticcurveto{\pgfqpoint{9.614401in}{7.100632in}}{\pgfqpoint{9.622149in}{7.106564in}}%
\pgfpathquadraticcurveto{\pgfqpoint{9.629896in}{7.112518in}}{\pgfqpoint{9.644151in}{7.112518in}}%
\pgfpathquadraticcurveto{\pgfqpoint{9.651190in}{7.112518in}}{\pgfqpoint{9.657410in}{7.111478in}}%
\pgfpathquadraticcurveto{\pgfqpoint{9.663653in}{7.110460in}}{\pgfqpoint{9.668899in}{7.108379in}}%
\pgfpathlineto{\pgfqpoint{9.668899in}{7.108379in}}%
\pgfpathclose%
\pgfpathmoveto{\pgfqpoint{9.696568in}{7.050761in}}%
\pgfpathlineto{\pgfqpoint{9.711155in}{7.050761in}}%
\pgfpathlineto{\pgfqpoint{9.711155in}{7.033185in}}%
\pgfpathlineto{\pgfqpoint{9.696568in}{7.033185in}}%
\pgfpathlineto{\pgfqpoint{9.696568in}{7.050761in}}%
\pgfpathlineto{\pgfqpoint{9.696568in}{7.050761in}}%
\pgfpathclose%
\pgfpathmoveto{\pgfqpoint{9.696568in}{7.106431in}}%
\pgfpathlineto{\pgfqpoint{9.711155in}{7.106431in}}%
\pgfpathlineto{\pgfqpoint{9.711155in}{7.088878in}}%
\pgfpathlineto{\pgfqpoint{9.696568in}{7.088878in}}%
\pgfpathlineto{\pgfqpoint{9.696568in}{7.106431in}}%
\pgfpathlineto{\pgfqpoint{9.696568in}{7.106431in}}%
\pgfpathclose%
\pgfpathmoveto{\pgfqpoint{9.790289in}{7.044939in}}%
\pgfpathlineto{\pgfqpoint{9.813111in}{7.044939in}}%
\pgfpathlineto{\pgfqpoint{9.813111in}{7.123741in}}%
\pgfpathlineto{\pgfqpoint{9.788275in}{7.118761in}}%
\pgfpathlineto{\pgfqpoint{9.788275in}{7.131488in}}%
\pgfpathlineto{\pgfqpoint{9.812978in}{7.136469in}}%
\pgfpathlineto{\pgfqpoint{9.826945in}{7.136469in}}%
\pgfpathlineto{\pgfqpoint{9.826945in}{7.044939in}}%
\pgfpathlineto{\pgfqpoint{9.849767in}{7.044939in}}%
\pgfpathlineto{\pgfqpoint{9.849767in}{7.033185in}}%
\pgfpathlineto{\pgfqpoint{9.790289in}{7.033185in}}%
\pgfpathlineto{\pgfqpoint{9.790289in}{7.044939in}}%
\pgfpathlineto{\pgfqpoint{9.790289in}{7.044939in}}%
\pgfpathclose%
\pgfpathmoveto{\pgfqpoint{9.907873in}{7.127261in}}%
\pgfpathquadraticcurveto{\pgfqpoint{9.897093in}{7.127261in}}{\pgfqpoint{9.891647in}{7.116636in}}%
\pgfpathquadraticcurveto{\pgfqpoint{9.886224in}{7.106033in}}{\pgfqpoint{9.886224in}{7.084716in}}%
\pgfpathquadraticcurveto{\pgfqpoint{9.886224in}{7.063488in}}{\pgfqpoint{9.891647in}{7.052863in}}%
\pgfpathquadraticcurveto{\pgfqpoint{9.897093in}{7.042238in}}{\pgfqpoint{9.907873in}{7.042238in}}%
\pgfpathquadraticcurveto{\pgfqpoint{9.918741in}{7.042238in}}{\pgfqpoint{9.924164in}{7.052863in}}%
\pgfpathquadraticcurveto{\pgfqpoint{9.929610in}{7.063488in}}{\pgfqpoint{9.929610in}{7.084716in}}%
\pgfpathquadraticcurveto{\pgfqpoint{9.929610in}{7.106033in}}{\pgfqpoint{9.924164in}{7.116636in}}%
\pgfpathquadraticcurveto{\pgfqpoint{9.918741in}{7.127261in}}{\pgfqpoint{9.907873in}{7.127261in}}%
\pgfpathlineto{\pgfqpoint{9.907873in}{7.127261in}}%
\pgfpathclose%
\pgfpathmoveto{\pgfqpoint{9.907873in}{7.138328in}}%
\pgfpathquadraticcurveto{\pgfqpoint{9.925249in}{7.138328in}}{\pgfqpoint{9.934413in}{7.124582in}}%
\pgfpathquadraticcurveto{\pgfqpoint{9.943577in}{7.110858in}}{\pgfqpoint{9.943577in}{7.084716in}}%
\pgfpathquadraticcurveto{\pgfqpoint{9.943577in}{7.058641in}}{\pgfqpoint{9.934413in}{7.044895in}}%
\pgfpathquadraticcurveto{\pgfqpoint{9.925249in}{7.031171in}}{\pgfqpoint{9.907873in}{7.031171in}}%
\pgfpathquadraticcurveto{\pgfqpoint{9.890518in}{7.031171in}}{\pgfqpoint{9.881354in}{7.044895in}}%
\pgfpathquadraticcurveto{\pgfqpoint{9.872190in}{7.058641in}}{\pgfqpoint{9.872190in}{7.084716in}}%
\pgfpathquadraticcurveto{\pgfqpoint{9.872190in}{7.110858in}}{\pgfqpoint{9.881354in}{7.124582in}}%
\pgfpathquadraticcurveto{\pgfqpoint{9.890518in}{7.138328in}}{\pgfqpoint{9.907873in}{7.138328in}}%
\pgfpathlineto{\pgfqpoint{9.907873in}{7.138328in}}%
\pgfpathclose%
\pgfpathmoveto{\pgfqpoint{9.968125in}{7.050761in}}%
\pgfpathlineto{\pgfqpoint{9.982735in}{7.050761in}}%
\pgfpathlineto{\pgfqpoint{9.982735in}{7.033185in}}%
\pgfpathlineto{\pgfqpoint{9.968125in}{7.033185in}}%
\pgfpathlineto{\pgfqpoint{9.968125in}{7.050761in}}%
\pgfpathlineto{\pgfqpoint{9.968125in}{7.050761in}}%
\pgfpathclose%
\pgfpathmoveto{\pgfqpoint{10.015584in}{7.044939in}}%
\pgfpathlineto{\pgfqpoint{10.038405in}{7.044939in}}%
\pgfpathlineto{\pgfqpoint{10.038405in}{7.123741in}}%
\pgfpathlineto{\pgfqpoint{10.013569in}{7.118761in}}%
\pgfpathlineto{\pgfqpoint{10.013569in}{7.131488in}}%
\pgfpathlineto{\pgfqpoint{10.038272in}{7.136469in}}%
\pgfpathlineto{\pgfqpoint{10.052240in}{7.136469in}}%
\pgfpathlineto{\pgfqpoint{10.052240in}{7.044939in}}%
\pgfpathlineto{\pgfqpoint{10.075061in}{7.044939in}}%
\pgfpathlineto{\pgfqpoint{10.075061in}{7.033185in}}%
\pgfpathlineto{\pgfqpoint{10.015584in}{7.033185in}}%
\pgfpathlineto{\pgfqpoint{10.015584in}{7.044939in}}%
\pgfpathlineto{\pgfqpoint{10.015584in}{7.044939in}}%
\pgfpathclose%
\pgfpathmoveto{\pgfqpoint{10.191140in}{7.078629in}}%
\pgfpathquadraticcurveto{\pgfqpoint{10.185119in}{7.078629in}}{\pgfqpoint{10.181688in}{7.073516in}}%
\pgfpathquadraticcurveto{\pgfqpoint{10.178279in}{7.068402in}}{\pgfqpoint{10.178279in}{7.059261in}}%
\pgfpathquadraticcurveto{\pgfqpoint{10.178279in}{7.050274in}}{\pgfqpoint{10.181688in}{7.045116in}}%
\pgfpathquadraticcurveto{\pgfqpoint{10.185119in}{7.039958in}}{\pgfqpoint{10.191140in}{7.039958in}}%
\pgfpathquadraticcurveto{\pgfqpoint{10.197028in}{7.039958in}}{\pgfqpoint{10.200436in}{7.045116in}}%
\pgfpathquadraticcurveto{\pgfqpoint{10.203867in}{7.050274in}}{\pgfqpoint{10.203867in}{7.059261in}}%
\pgfpathquadraticcurveto{\pgfqpoint{10.203867in}{7.068336in}}{\pgfqpoint{10.200436in}{7.073471in}}%
\pgfpathquadraticcurveto{\pgfqpoint{10.197028in}{7.078629in}}{\pgfqpoint{10.191140in}{7.078629in}}%
\pgfpathlineto{\pgfqpoint{10.191140in}{7.078629in}}%
\pgfpathclose%
\pgfpathmoveto{\pgfqpoint{10.191140in}{7.087417in}}%
\pgfpathquadraticcurveto{\pgfqpoint{10.202074in}{7.087417in}}{\pgfqpoint{10.208494in}{7.079802in}}%
\pgfpathquadraticcurveto{\pgfqpoint{10.214935in}{7.072210in}}{\pgfqpoint{10.214935in}{7.059261in}}%
\pgfpathquadraticcurveto{\pgfqpoint{10.214935in}{7.046333in}}{\pgfqpoint{10.208472in}{7.038741in}}%
\pgfpathquadraticcurveto{\pgfqpoint{10.202008in}{7.031171in}}{\pgfqpoint{10.191140in}{7.031171in}}%
\pgfpathquadraticcurveto{\pgfqpoint{10.180072in}{7.031171in}}{\pgfqpoint{10.173630in}{7.038741in}}%
\pgfpathquadraticcurveto{\pgfqpoint{10.167211in}{7.046333in}}{\pgfqpoint{10.167211in}{7.059261in}}%
\pgfpathquadraticcurveto{\pgfqpoint{10.167211in}{7.072276in}}{\pgfqpoint{10.173675in}{7.079846in}}%
\pgfpathquadraticcurveto{\pgfqpoint{10.180138in}{7.087417in}}{\pgfqpoint{10.191140in}{7.087417in}}%
\pgfpathlineto{\pgfqpoint{10.191140in}{7.087417in}}%
\pgfpathclose%
\pgfpathmoveto{\pgfqpoint{10.119753in}{7.129540in}}%
\pgfpathquadraticcurveto{\pgfqpoint{10.113798in}{7.129540in}}{\pgfqpoint{10.110367in}{7.124383in}}%
\pgfpathquadraticcurveto{\pgfqpoint{10.106959in}{7.119248in}}{\pgfqpoint{10.106959in}{7.110238in}}%
\pgfpathquadraticcurveto{\pgfqpoint{10.106959in}{7.101119in}}{\pgfqpoint{10.110345in}{7.095983in}}%
\pgfpathquadraticcurveto{\pgfqpoint{10.113732in}{7.090870in}}{\pgfqpoint{10.119753in}{7.090870in}}%
\pgfpathquadraticcurveto{\pgfqpoint{10.125774in}{7.090870in}}{\pgfqpoint{10.129182in}{7.095983in}}%
\pgfpathquadraticcurveto{\pgfqpoint{10.132613in}{7.101119in}}{\pgfqpoint{10.132613in}{7.110238in}}%
\pgfpathquadraticcurveto{\pgfqpoint{10.132613in}{7.119159in}}{\pgfqpoint{10.129160in}{7.124339in}}%
\pgfpathquadraticcurveto{\pgfqpoint{10.125707in}{7.129540in}}{\pgfqpoint{10.119753in}{7.129540in}}%
\pgfpathlineto{\pgfqpoint{10.119753in}{7.129540in}}%
\pgfpathclose%
\pgfpathmoveto{\pgfqpoint{10.182219in}{7.138328in}}%
\pgfpathlineto{\pgfqpoint{10.193287in}{7.138328in}}%
\pgfpathlineto{\pgfqpoint{10.128673in}{7.031171in}}%
\pgfpathlineto{\pgfqpoint{10.117606in}{7.031171in}}%
\pgfpathlineto{\pgfqpoint{10.182219in}{7.138328in}}%
\pgfpathlineto{\pgfqpoint{10.182219in}{7.138328in}}%
\pgfpathclose%
\pgfpathmoveto{\pgfqpoint{10.119753in}{7.138328in}}%
\pgfpathquadraticcurveto{\pgfqpoint{10.130688in}{7.138328in}}{\pgfqpoint{10.137173in}{7.130758in}}%
\pgfpathquadraticcurveto{\pgfqpoint{10.143681in}{7.123188in}}{\pgfqpoint{10.143681in}{7.110238in}}%
\pgfpathquadraticcurveto{\pgfqpoint{10.143681in}{7.097178in}}{\pgfqpoint{10.137218in}{7.089630in}}%
\pgfpathquadraticcurveto{\pgfqpoint{10.130754in}{7.082082in}}{\pgfqpoint{10.119753in}{7.082082in}}%
\pgfpathquadraticcurveto{\pgfqpoint{10.108751in}{7.082082in}}{\pgfqpoint{10.102354in}{7.089652in}}%
\pgfpathquadraticcurveto{\pgfqpoint{10.095957in}{7.097245in}}{\pgfqpoint{10.095957in}{7.110238in}}%
\pgfpathquadraticcurveto{\pgfqpoint{10.095957in}{7.123121in}}{\pgfqpoint{10.102376in}{7.130714in}}%
\pgfpathquadraticcurveto{\pgfqpoint{10.108818in}{7.138328in}}{\pgfqpoint{10.119753in}{7.138328in}}%
\pgfpathlineto{\pgfqpoint{10.119753in}{7.138328in}}%
\pgfpathclose%
\pgfusepath{fill}%
\end{pgfscope}%
\begin{pgfscope}%
\pgfsetbuttcap%
\pgfsetroundjoin%
\definecolor{currentfill}{rgb}{0.090196,0.168627,0.227451}%
\pgfsetfillcolor{currentfill}%
\pgfsetlinewidth{0.000000pt}%
\definecolor{currentstroke}{rgb}{0.090196,0.168627,0.227451}%
\pgfsetstrokecolor{currentstroke}%
\pgfsetdash{}{0pt}%
\pgfpathmoveto{\pgfqpoint{8.178145in}{6.913499in}}%
\pgfpathlineto{\pgfqpoint{8.178145in}{6.899863in}}%
\pgfpathquadraticcurveto{\pgfqpoint{8.170198in}{6.903671in}}{\pgfqpoint{8.163137in}{6.905530in}}%
\pgfpathquadraticcurveto{\pgfqpoint{8.156076in}{6.907412in}}{\pgfqpoint{8.149501in}{6.907412in}}%
\pgfpathquadraticcurveto{\pgfqpoint{8.138102in}{6.907412in}}{\pgfqpoint{8.131904in}{6.902984in}}%
\pgfpathquadraticcurveto{\pgfqpoint{8.125706in}{6.898557in}}{\pgfqpoint{8.125706in}{6.890389in}}%
\pgfpathquadraticcurveto{\pgfqpoint{8.125706in}{6.883527in}}{\pgfqpoint{8.129823in}{6.880030in}}%
\pgfpathquadraticcurveto{\pgfqpoint{8.133940in}{6.876555in}}{\pgfqpoint{8.145429in}{6.874408in}}%
\pgfpathlineto{\pgfqpoint{8.153862in}{6.872681in}}%
\pgfpathquadraticcurveto{\pgfqpoint{8.169490in}{6.869693in}}{\pgfqpoint{8.176927in}{6.862189in}}%
\pgfpathquadraticcurveto{\pgfqpoint{8.184365in}{6.854685in}}{\pgfqpoint{8.184365in}{6.842112in}}%
\pgfpathquadraticcurveto{\pgfqpoint{8.184365in}{6.827082in}}{\pgfqpoint{8.174293in}{6.819335in}}%
\pgfpathquadraticcurveto{\pgfqpoint{8.164244in}{6.811587in}}{\pgfqpoint{8.144809in}{6.811587in}}%
\pgfpathquadraticcurveto{\pgfqpoint{8.137482in}{6.811587in}}{\pgfqpoint{8.129203in}{6.813248in}}%
\pgfpathquadraticcurveto{\pgfqpoint{8.120947in}{6.814908in}}{\pgfqpoint{8.112093in}{6.818162in}}%
\pgfpathlineto{\pgfqpoint{8.112093in}{6.832550in}}%
\pgfpathquadraticcurveto{\pgfqpoint{8.120593in}{6.827790in}}{\pgfqpoint{8.128761in}{6.825356in}}%
\pgfpathquadraticcurveto{\pgfqpoint{8.136929in}{6.822943in}}{\pgfqpoint{8.144809in}{6.822943in}}%
\pgfpathquadraticcurveto{\pgfqpoint{8.156762in}{6.822943in}}{\pgfqpoint{8.163270in}{6.827636in}}%
\pgfpathquadraticcurveto{\pgfqpoint{8.169778in}{6.832350in}}{\pgfqpoint{8.169778in}{6.841072in}}%
\pgfpathquadraticcurveto{\pgfqpoint{8.169778in}{6.848664in}}{\pgfqpoint{8.165107in}{6.852958in}}%
\pgfpathquadraticcurveto{\pgfqpoint{8.160436in}{6.857253in}}{\pgfqpoint{8.149789in}{6.859400in}}%
\pgfpathlineto{\pgfqpoint{8.141267in}{6.861060in}}%
\pgfpathquadraticcurveto{\pgfqpoint{8.125640in}{6.864159in}}{\pgfqpoint{8.118645in}{6.870800in}}%
\pgfpathquadraticcurveto{\pgfqpoint{8.111672in}{6.877440in}}{\pgfqpoint{8.111672in}{6.889283in}}%
\pgfpathquadraticcurveto{\pgfqpoint{8.111672in}{6.902984in}}{\pgfqpoint{8.121323in}{6.910865in}}%
\pgfpathquadraticcurveto{\pgfqpoint{8.130974in}{6.918745in}}{\pgfqpoint{8.147908in}{6.918745in}}%
\pgfpathquadraticcurveto{\pgfqpoint{8.155190in}{6.918745in}}{\pgfqpoint{8.162716in}{6.917417in}}%
\pgfpathquadraticcurveto{\pgfqpoint{8.170265in}{6.916111in}}{\pgfqpoint{8.178145in}{6.913499in}}%
\pgfpathlineto{\pgfqpoint{8.178145in}{6.913499in}}%
\pgfpathclose%
\pgfpathmoveto{\pgfqpoint{8.205615in}{6.891076in}}%
\pgfpathlineto{\pgfqpoint{8.218343in}{6.891076in}}%
\pgfpathlineto{\pgfqpoint{8.218343in}{6.813602in}}%
\pgfpathlineto{\pgfqpoint{8.205615in}{6.813602in}}%
\pgfpathlineto{\pgfqpoint{8.205615in}{6.891076in}}%
\pgfpathlineto{\pgfqpoint{8.205615in}{6.891076in}}%
\pgfpathclose%
\pgfpathmoveto{\pgfqpoint{8.205615in}{6.921246in}}%
\pgfpathlineto{\pgfqpoint{8.218343in}{6.921246in}}%
\pgfpathlineto{\pgfqpoint{8.218343in}{6.905109in}}%
\pgfpathlineto{\pgfqpoint{8.205615in}{6.905109in}}%
\pgfpathlineto{\pgfqpoint{8.205615in}{6.921246in}}%
\pgfpathlineto{\pgfqpoint{8.205615in}{6.921246in}}%
\pgfpathclose%
\pgfpathmoveto{\pgfqpoint{8.295949in}{6.853246in}}%
\pgfpathquadraticcurveto{\pgfqpoint{8.295949in}{6.867081in}}{\pgfqpoint{8.290238in}{6.874673in}}%
\pgfpathquadraticcurveto{\pgfqpoint{8.284550in}{6.882288in}}{\pgfqpoint{8.274235in}{6.882288in}}%
\pgfpathquadraticcurveto{\pgfqpoint{8.264008in}{6.882288in}}{\pgfqpoint{8.258297in}{6.874673in}}%
\pgfpathquadraticcurveto{\pgfqpoint{8.252586in}{6.867081in}}{\pgfqpoint{8.252586in}{6.853246in}}%
\pgfpathquadraticcurveto{\pgfqpoint{8.252586in}{6.839478in}}{\pgfqpoint{8.258297in}{6.831863in}}%
\pgfpathquadraticcurveto{\pgfqpoint{8.264008in}{6.824249in}}{\pgfqpoint{8.274235in}{6.824249in}}%
\pgfpathquadraticcurveto{\pgfqpoint{8.284550in}{6.824249in}}{\pgfqpoint{8.290238in}{6.831863in}}%
\pgfpathquadraticcurveto{\pgfqpoint{8.295949in}{6.839478in}}{\pgfqpoint{8.295949in}{6.853246in}}%
\pgfpathlineto{\pgfqpoint{8.295949in}{6.853246in}}%
\pgfpathclose%
\pgfpathmoveto{\pgfqpoint{8.308677in}{6.823208in}}%
\pgfpathquadraticcurveto{\pgfqpoint{8.308677in}{6.803442in}}{\pgfqpoint{8.299890in}{6.793790in}}%
\pgfpathquadraticcurveto{\pgfqpoint{8.291124in}{6.784139in}}{\pgfqpoint{8.272995in}{6.784139in}}%
\pgfpathquadraticcurveto{\pgfqpoint{8.266288in}{6.784139in}}{\pgfqpoint{8.260334in}{6.785136in}}%
\pgfpathquadraticcurveto{\pgfqpoint{8.254379in}{6.786132in}}{\pgfqpoint{8.248779in}{6.788212in}}%
\pgfpathlineto{\pgfqpoint{8.248779in}{6.800586in}}%
\pgfpathquadraticcurveto{\pgfqpoint{8.254379in}{6.797553in}}{\pgfqpoint{8.259847in}{6.796115in}}%
\pgfpathquadraticcurveto{\pgfqpoint{8.265314in}{6.794654in}}{\pgfqpoint{8.270981in}{6.794654in}}%
\pgfpathquadraticcurveto{\pgfqpoint{8.283509in}{6.794654in}}{\pgfqpoint{8.289729in}{6.801184in}}%
\pgfpathquadraticcurveto{\pgfqpoint{8.295949in}{6.807714in}}{\pgfqpoint{8.295949in}{6.820928in}}%
\pgfpathlineto{\pgfqpoint{8.295949in}{6.827237in}}%
\pgfpathquadraticcurveto{\pgfqpoint{8.292009in}{6.820375in}}{\pgfqpoint{8.285856in}{6.816988in}}%
\pgfpathquadraticcurveto{\pgfqpoint{8.279702in}{6.813602in}}{\pgfqpoint{8.271113in}{6.813602in}}%
\pgfpathquadraticcurveto{\pgfqpoint{8.256880in}{6.813602in}}{\pgfqpoint{8.248159in}{6.824448in}}%
\pgfpathquadraticcurveto{\pgfqpoint{8.239438in}{6.835317in}}{\pgfqpoint{8.239438in}{6.853246in}}%
\pgfpathquadraticcurveto{\pgfqpoint{8.239438in}{6.871220in}}{\pgfqpoint{8.248159in}{6.882067in}}%
\pgfpathquadraticcurveto{\pgfqpoint{8.256880in}{6.892935in}}{\pgfqpoint{8.271113in}{6.892935in}}%
\pgfpathquadraticcurveto{\pgfqpoint{8.279702in}{6.892935in}}{\pgfqpoint{8.285856in}{6.889548in}}%
\pgfpathquadraticcurveto{\pgfqpoint{8.292009in}{6.886162in}}{\pgfqpoint{8.295949in}{6.879322in}}%
\pgfpathlineto{\pgfqpoint{8.295949in}{6.891076in}}%
\pgfpathlineto{\pgfqpoint{8.308677in}{6.891076in}}%
\pgfpathlineto{\pgfqpoint{8.308677in}{6.823208in}}%
\pgfpathlineto{\pgfqpoint{8.308677in}{6.823208in}}%
\pgfpathclose%
\pgfpathmoveto{\pgfqpoint{8.399322in}{6.860374in}}%
\pgfpathlineto{\pgfqpoint{8.399322in}{6.813602in}}%
\pgfpathlineto{\pgfqpoint{8.386594in}{6.813602in}}%
\pgfpathlineto{\pgfqpoint{8.386594in}{6.859953in}}%
\pgfpathquadraticcurveto{\pgfqpoint{8.386594in}{6.870955in}}{\pgfqpoint{8.382300in}{6.876400in}}%
\pgfpathquadraticcurveto{\pgfqpoint{8.378005in}{6.881867in}}{\pgfqpoint{8.369439in}{6.881867in}}%
\pgfpathquadraticcurveto{\pgfqpoint{8.359124in}{6.881867in}}{\pgfqpoint{8.353169in}{6.875293in}}%
\pgfpathquadraticcurveto{\pgfqpoint{8.347215in}{6.868741in}}{\pgfqpoint{8.347215in}{6.857386in}}%
\pgfpathlineto{\pgfqpoint{8.347215in}{6.813602in}}%
\pgfpathlineto{\pgfqpoint{8.334421in}{6.813602in}}%
\pgfpathlineto{\pgfqpoint{8.334421in}{6.891076in}}%
\pgfpathlineto{\pgfqpoint{8.347215in}{6.891076in}}%
\pgfpathlineto{\pgfqpoint{8.347215in}{6.879034in}}%
\pgfpathquadraticcurveto{\pgfqpoint{8.351797in}{6.886029in}}{\pgfqpoint{8.357973in}{6.889482in}}%
\pgfpathquadraticcurveto{\pgfqpoint{8.364171in}{6.892935in}}{\pgfqpoint{8.372272in}{6.892935in}}%
\pgfpathquadraticcurveto{\pgfqpoint{8.385620in}{6.892935in}}{\pgfqpoint{8.392460in}{6.884678in}}%
\pgfpathquadraticcurveto{\pgfqpoint{8.399322in}{6.876422in}}{\pgfqpoint{8.399322in}{6.860374in}}%
\pgfpathlineto{\pgfqpoint{8.399322in}{6.860374in}}%
\pgfpathclose%
\pgfpathmoveto{\pgfqpoint{8.424689in}{6.891076in}}%
\pgfpathlineto{\pgfqpoint{8.437417in}{6.891076in}}%
\pgfpathlineto{\pgfqpoint{8.437417in}{6.813602in}}%
\pgfpathlineto{\pgfqpoint{8.424689in}{6.813602in}}%
\pgfpathlineto{\pgfqpoint{8.424689in}{6.891076in}}%
\pgfpathlineto{\pgfqpoint{8.424689in}{6.891076in}}%
\pgfpathclose%
\pgfpathmoveto{\pgfqpoint{8.424689in}{6.921246in}}%
\pgfpathlineto{\pgfqpoint{8.437417in}{6.921246in}}%
\pgfpathlineto{\pgfqpoint{8.437417in}{6.905109in}}%
\pgfpathlineto{\pgfqpoint{8.424689in}{6.905109in}}%
\pgfpathlineto{\pgfqpoint{8.424689in}{6.921246in}}%
\pgfpathlineto{\pgfqpoint{8.424689in}{6.921246in}}%
\pgfpathclose%
\pgfpathmoveto{\pgfqpoint{8.526645in}{6.891076in}}%
\pgfpathlineto{\pgfqpoint{8.526645in}{6.813602in}}%
\pgfpathlineto{\pgfqpoint{8.513850in}{6.813602in}}%
\pgfpathlineto{\pgfqpoint{8.513850in}{6.881181in}}%
\pgfpathlineto{\pgfqpoint{8.478921in}{6.881181in}}%
\pgfpathlineto{\pgfqpoint{8.478921in}{6.813602in}}%
\pgfpathlineto{\pgfqpoint{8.466126in}{6.813602in}}%
\pgfpathlineto{\pgfqpoint{8.466126in}{6.881181in}}%
\pgfpathlineto{\pgfqpoint{8.453952in}{6.881181in}}%
\pgfpathlineto{\pgfqpoint{8.453952in}{6.891076in}}%
\pgfpathlineto{\pgfqpoint{8.466126in}{6.891076in}}%
\pgfpathlineto{\pgfqpoint{8.466126in}{6.896477in}}%
\pgfpathquadraticcurveto{\pgfqpoint{8.466126in}{6.909138in}}{\pgfqpoint{8.472103in}{6.915181in}}%
\pgfpathquadraticcurveto{\pgfqpoint{8.478102in}{6.921246in}}{\pgfqpoint{8.490475in}{6.921246in}}%
\pgfpathlineto{\pgfqpoint{8.503270in}{6.921246in}}%
\pgfpathlineto{\pgfqpoint{8.503270in}{6.910643in}}%
\pgfpathlineto{\pgfqpoint{8.491095in}{6.910643in}}%
\pgfpathquadraticcurveto{\pgfqpoint{8.484255in}{6.910643in}}{\pgfqpoint{8.481577in}{6.907876in}}%
\pgfpathquadraticcurveto{\pgfqpoint{8.478921in}{6.905109in}}{\pgfqpoint{8.478921in}{6.897915in}}%
\pgfpathlineto{\pgfqpoint{8.478921in}{6.891076in}}%
\pgfpathlineto{\pgfqpoint{8.526645in}{6.891076in}}%
\pgfpathlineto{\pgfqpoint{8.526645in}{6.891076in}}%
\pgfpathclose%
\pgfpathmoveto{\pgfqpoint{8.513850in}{6.921091in}}%
\pgfpathlineto{\pgfqpoint{8.526645in}{6.921091in}}%
\pgfpathlineto{\pgfqpoint{8.526645in}{6.904977in}}%
\pgfpathlineto{\pgfqpoint{8.513850in}{6.904977in}}%
\pgfpathlineto{\pgfqpoint{8.513850in}{6.921091in}}%
\pgfpathlineto{\pgfqpoint{8.513850in}{6.921091in}}%
\pgfpathclose%
\pgfpathmoveto{\pgfqpoint{8.609033in}{6.888109in}}%
\pgfpathlineto{\pgfqpoint{8.609033in}{6.876201in}}%
\pgfpathquadraticcurveto{\pgfqpoint{8.603632in}{6.879189in}}{\pgfqpoint{8.598209in}{6.880672in}}%
\pgfpathquadraticcurveto{\pgfqpoint{8.592785in}{6.882155in}}{\pgfqpoint{8.587251in}{6.882155in}}%
\pgfpathquadraticcurveto{\pgfqpoint{8.574856in}{6.882155in}}{\pgfqpoint{8.567994in}{6.874297in}}%
\pgfpathquadraticcurveto{\pgfqpoint{8.561154in}{6.866461in}}{\pgfqpoint{8.561154in}{6.852272in}}%
\pgfpathquadraticcurveto{\pgfqpoint{8.561154in}{6.838083in}}{\pgfqpoint{8.567994in}{6.830225in}}%
\pgfpathquadraticcurveto{\pgfqpoint{8.574856in}{6.822389in}}{\pgfqpoint{8.587251in}{6.822389in}}%
\pgfpathquadraticcurveto{\pgfqpoint{8.592785in}{6.822389in}}{\pgfqpoint{8.598209in}{6.823873in}}%
\pgfpathquadraticcurveto{\pgfqpoint{8.603632in}{6.825356in}}{\pgfqpoint{8.609033in}{6.828344in}}%
\pgfpathlineto{\pgfqpoint{8.609033in}{6.816568in}}%
\pgfpathquadraticcurveto{\pgfqpoint{8.603698in}{6.814089in}}{\pgfqpoint{8.597987in}{6.812849in}}%
\pgfpathquadraticcurveto{\pgfqpoint{8.592298in}{6.811587in}}{\pgfqpoint{8.585857in}{6.811587in}}%
\pgfpathquadraticcurveto{\pgfqpoint{8.568348in}{6.811587in}}{\pgfqpoint{8.558033in}{6.822589in}}%
\pgfpathquadraticcurveto{\pgfqpoint{8.547740in}{6.833590in}}{\pgfqpoint{8.547740in}{6.852272in}}%
\pgfpathquadraticcurveto{\pgfqpoint{8.547740in}{6.871220in}}{\pgfqpoint{8.558143in}{6.882067in}}%
\pgfpathquadraticcurveto{\pgfqpoint{8.568569in}{6.892935in}}{\pgfqpoint{8.586698in}{6.892935in}}%
\pgfpathquadraticcurveto{\pgfqpoint{8.592564in}{6.892935in}}{\pgfqpoint{8.598164in}{6.891718in}}%
\pgfpathquadraticcurveto{\pgfqpoint{8.603765in}{6.890522in}}{\pgfqpoint{8.609033in}{6.888109in}}%
\pgfpathlineto{\pgfqpoint{8.609033in}{6.888109in}}%
\pgfpathclose%
\pgfpathmoveto{\pgfqpoint{8.666386in}{6.852538in}}%
\pgfpathquadraticcurveto{\pgfqpoint{8.650957in}{6.852538in}}{\pgfqpoint{8.645003in}{6.849018in}}%
\pgfpathquadraticcurveto{\pgfqpoint{8.639048in}{6.845499in}}{\pgfqpoint{8.639048in}{6.836977in}}%
\pgfpathquadraticcurveto{\pgfqpoint{8.639048in}{6.830203in}}{\pgfqpoint{8.643520in}{6.826219in}}%
\pgfpathquadraticcurveto{\pgfqpoint{8.647991in}{6.822257in}}{\pgfqpoint{8.655650in}{6.822257in}}%
\pgfpathquadraticcurveto{\pgfqpoint{8.666253in}{6.822257in}}{\pgfqpoint{8.672650in}{6.829761in}}%
\pgfpathquadraticcurveto{\pgfqpoint{8.679047in}{6.837264in}}{\pgfqpoint{8.679047in}{6.849705in}}%
\pgfpathlineto{\pgfqpoint{8.679047in}{6.852538in}}%
\pgfpathlineto{\pgfqpoint{8.666386in}{6.852538in}}%
\pgfpathlineto{\pgfqpoint{8.666386in}{6.852538in}}%
\pgfpathclose%
\pgfpathmoveto{\pgfqpoint{8.691775in}{6.857806in}}%
\pgfpathlineto{\pgfqpoint{8.691775in}{6.813602in}}%
\pgfpathlineto{\pgfqpoint{8.679047in}{6.813602in}}%
\pgfpathlineto{\pgfqpoint{8.679047in}{6.825356in}}%
\pgfpathquadraticcurveto{\pgfqpoint{8.674686in}{6.818317in}}{\pgfqpoint{8.668179in}{6.814952in}}%
\pgfpathquadraticcurveto{\pgfqpoint{8.661671in}{6.811587in}}{\pgfqpoint{8.652263in}{6.811587in}}%
\pgfpathquadraticcurveto{\pgfqpoint{8.640376in}{6.811587in}}{\pgfqpoint{8.633337in}{6.818272in}}%
\pgfpathquadraticcurveto{\pgfqpoint{8.626320in}{6.824957in}}{\pgfqpoint{8.626320in}{6.836158in}}%
\pgfpathquadraticcurveto{\pgfqpoint{8.626320in}{6.849218in}}{\pgfqpoint{8.635064in}{6.855858in}}%
\pgfpathquadraticcurveto{\pgfqpoint{8.643830in}{6.862499in}}{\pgfqpoint{8.661184in}{6.862499in}}%
\pgfpathlineto{\pgfqpoint{8.679047in}{6.862499in}}%
\pgfpathlineto{\pgfqpoint{8.679047in}{6.863761in}}%
\pgfpathquadraticcurveto{\pgfqpoint{8.679047in}{6.872548in}}{\pgfqpoint{8.673270in}{6.877352in}}%
\pgfpathquadraticcurveto{\pgfqpoint{8.667492in}{6.882155in}}{\pgfqpoint{8.657044in}{6.882155in}}%
\pgfpathquadraticcurveto{\pgfqpoint{8.650404in}{6.882155in}}{\pgfqpoint{8.644095in}{6.880561in}}%
\pgfpathquadraticcurveto{\pgfqpoint{8.637809in}{6.878968in}}{\pgfqpoint{8.632009in}{6.875780in}}%
\pgfpathlineto{\pgfqpoint{8.632009in}{6.887556in}}%
\pgfpathquadraticcurveto{\pgfqpoint{8.638982in}{6.890257in}}{\pgfqpoint{8.645556in}{6.891585in}}%
\pgfpathquadraticcurveto{\pgfqpoint{8.652130in}{6.892935in}}{\pgfqpoint{8.658350in}{6.892935in}}%
\pgfpathquadraticcurveto{\pgfqpoint{8.675173in}{6.892935in}}{\pgfqpoint{8.683474in}{6.884214in}}%
\pgfpathquadraticcurveto{\pgfqpoint{8.691775in}{6.875514in}}{\pgfqpoint{8.691775in}{6.857806in}}%
\pgfpathlineto{\pgfqpoint{8.691775in}{6.857806in}}%
\pgfpathclose%
\pgfpathmoveto{\pgfqpoint{8.782397in}{6.860374in}}%
\pgfpathlineto{\pgfqpoint{8.782397in}{6.813602in}}%
\pgfpathlineto{\pgfqpoint{8.769669in}{6.813602in}}%
\pgfpathlineto{\pgfqpoint{8.769669in}{6.859953in}}%
\pgfpathquadraticcurveto{\pgfqpoint{8.769669in}{6.870955in}}{\pgfqpoint{8.765375in}{6.876400in}}%
\pgfpathquadraticcurveto{\pgfqpoint{8.761081in}{6.881867in}}{\pgfqpoint{8.752515in}{6.881867in}}%
\pgfpathquadraticcurveto{\pgfqpoint{8.742199in}{6.881867in}}{\pgfqpoint{8.736245in}{6.875293in}}%
\pgfpathquadraticcurveto{\pgfqpoint{8.730291in}{6.868741in}}{\pgfqpoint{8.730291in}{6.857386in}}%
\pgfpathlineto{\pgfqpoint{8.730291in}{6.813602in}}%
\pgfpathlineto{\pgfqpoint{8.717496in}{6.813602in}}%
\pgfpathlineto{\pgfqpoint{8.717496in}{6.891076in}}%
\pgfpathlineto{\pgfqpoint{8.730291in}{6.891076in}}%
\pgfpathlineto{\pgfqpoint{8.730291in}{6.879034in}}%
\pgfpathquadraticcurveto{\pgfqpoint{8.734873in}{6.886029in}}{\pgfqpoint{8.741048in}{6.889482in}}%
\pgfpathquadraticcurveto{\pgfqpoint{8.747246in}{6.892935in}}{\pgfqpoint{8.755348in}{6.892935in}}%
\pgfpathquadraticcurveto{\pgfqpoint{8.768695in}{6.892935in}}{\pgfqpoint{8.775535in}{6.884678in}}%
\pgfpathquadraticcurveto{\pgfqpoint{8.782397in}{6.876422in}}{\pgfqpoint{8.782397in}{6.860374in}}%
\pgfpathlineto{\pgfqpoint{8.782397in}{6.860374in}}%
\pgfpathclose%
\pgfpathmoveto{\pgfqpoint{8.820360in}{6.913078in}}%
\pgfpathlineto{\pgfqpoint{8.820360in}{6.891076in}}%
\pgfpathlineto{\pgfqpoint{8.846568in}{6.891076in}}%
\pgfpathlineto{\pgfqpoint{8.846568in}{6.881181in}}%
\pgfpathlineto{\pgfqpoint{8.820360in}{6.881181in}}%
\pgfpathlineto{\pgfqpoint{8.820360in}{6.839124in}}%
\pgfpathquadraticcurveto{\pgfqpoint{8.820360in}{6.829650in}}{\pgfqpoint{8.822949in}{6.826949in}}%
\pgfpathquadraticcurveto{\pgfqpoint{8.825539in}{6.824249in}}{\pgfqpoint{8.833508in}{6.824249in}}%
\pgfpathlineto{\pgfqpoint{8.846568in}{6.824249in}}%
\pgfpathlineto{\pgfqpoint{8.846568in}{6.813602in}}%
\pgfpathlineto{\pgfqpoint{8.833508in}{6.813602in}}%
\pgfpathquadraticcurveto{\pgfqpoint{8.818766in}{6.813602in}}{\pgfqpoint{8.813166in}{6.819091in}}%
\pgfpathquadraticcurveto{\pgfqpoint{8.807565in}{6.824603in}}{\pgfqpoint{8.807565in}{6.839124in}}%
\pgfpathlineto{\pgfqpoint{8.807565in}{6.881181in}}%
\pgfpathlineto{\pgfqpoint{8.798224in}{6.881181in}}%
\pgfpathlineto{\pgfqpoint{8.798224in}{6.891076in}}%
\pgfpathlineto{\pgfqpoint{8.807565in}{6.891076in}}%
\pgfpathlineto{\pgfqpoint{8.807565in}{6.913078in}}%
\pgfpathlineto{\pgfqpoint{8.820360in}{6.913078in}}%
\pgfpathlineto{\pgfqpoint{8.820360in}{6.913078in}}%
\pgfpathclose%
\pgfpathmoveto{\pgfqpoint{8.960145in}{6.902431in}}%
\pgfpathlineto{\pgfqpoint{8.960145in}{6.863893in}}%
\pgfpathlineto{\pgfqpoint{8.998660in}{6.863893in}}%
\pgfpathlineto{\pgfqpoint{8.998660in}{6.852139in}}%
\pgfpathlineto{\pgfqpoint{8.960145in}{6.852139in}}%
\pgfpathlineto{\pgfqpoint{8.960145in}{6.813602in}}%
\pgfpathlineto{\pgfqpoint{8.948524in}{6.813602in}}%
\pgfpathlineto{\pgfqpoint{8.948524in}{6.852139in}}%
\pgfpathlineto{\pgfqpoint{8.909986in}{6.852139in}}%
\pgfpathlineto{\pgfqpoint{8.909986in}{6.863893in}}%
\pgfpathlineto{\pgfqpoint{8.948524in}{6.863893in}}%
\pgfpathlineto{\pgfqpoint{8.948524in}{6.902431in}}%
\pgfpathlineto{\pgfqpoint{8.960145in}{6.902431in}}%
\pgfpathlineto{\pgfqpoint{8.960145in}{6.902431in}}%
\pgfpathclose%
\pgfpathmoveto{\pgfqpoint{9.121600in}{6.862034in}}%
\pgfpathquadraticcurveto{\pgfqpoint{9.126094in}{6.860507in}}{\pgfqpoint{9.130344in}{6.855526in}}%
\pgfpathquadraticcurveto{\pgfqpoint{9.134594in}{6.850546in}}{\pgfqpoint{9.138888in}{6.841824in}}%
\pgfpathlineto{\pgfqpoint{9.153077in}{6.813602in}}%
\pgfpathlineto{\pgfqpoint{9.138047in}{6.813602in}}%
\pgfpathlineto{\pgfqpoint{9.124854in}{6.840098in}}%
\pgfpathquadraticcurveto{\pgfqpoint{9.119719in}{6.850479in}}{\pgfqpoint{9.114916in}{6.853866in}}%
\pgfpathquadraticcurveto{\pgfqpoint{9.110112in}{6.857253in}}{\pgfqpoint{9.101811in}{6.857253in}}%
\pgfpathlineto{\pgfqpoint{9.086582in}{6.857253in}}%
\pgfpathlineto{\pgfqpoint{9.086582in}{6.813602in}}%
\pgfpathlineto{\pgfqpoint{9.072615in}{6.813602in}}%
\pgfpathlineto{\pgfqpoint{9.072615in}{6.916886in}}%
\pgfpathlineto{\pgfqpoint{9.104158in}{6.916886in}}%
\pgfpathquadraticcurveto{\pgfqpoint{9.121866in}{6.916886in}}{\pgfqpoint{9.130587in}{6.909470in}}%
\pgfpathquadraticcurveto{\pgfqpoint{9.139309in}{6.902077in}}{\pgfqpoint{9.139309in}{6.887136in}}%
\pgfpathquadraticcurveto{\pgfqpoint{9.139309in}{6.877374in}}{\pgfqpoint{9.134771in}{6.870932in}}%
\pgfpathquadraticcurveto{\pgfqpoint{9.130233in}{6.864513in}}{\pgfqpoint{9.121600in}{6.862034in}}%
\pgfpathlineto{\pgfqpoint{9.121600in}{6.862034in}}%
\pgfpathclose%
\pgfpathmoveto{\pgfqpoint{9.086582in}{6.905397in}}%
\pgfpathlineto{\pgfqpoint{9.086582in}{6.868741in}}%
\pgfpathlineto{\pgfqpoint{9.104158in}{6.868741in}}%
\pgfpathquadraticcurveto{\pgfqpoint{9.114251in}{6.868741in}}{\pgfqpoint{9.119409in}{6.873412in}}%
\pgfpathquadraticcurveto{\pgfqpoint{9.124567in}{6.878082in}}{\pgfqpoint{9.124567in}{6.887136in}}%
\pgfpathquadraticcurveto{\pgfqpoint{9.124567in}{6.896189in}}{\pgfqpoint{9.119409in}{6.900793in}}%
\pgfpathquadraticcurveto{\pgfqpoint{9.114251in}{6.905397in}}{\pgfqpoint{9.104158in}{6.905397in}}%
\pgfpathlineto{\pgfqpoint{9.086582in}{6.905397in}}%
\pgfpathlineto{\pgfqpoint{9.086582in}{6.905397in}}%
\pgfpathclose%
\pgfpathmoveto{\pgfqpoint{9.230418in}{6.855526in}}%
\pgfpathlineto{\pgfqpoint{9.230418in}{6.849306in}}%
\pgfpathlineto{\pgfqpoint{9.171892in}{6.849306in}}%
\pgfpathquadraticcurveto{\pgfqpoint{9.172733in}{6.836158in}}{\pgfqpoint{9.179817in}{6.829274in}}%
\pgfpathquadraticcurveto{\pgfqpoint{9.186900in}{6.822389in}}{\pgfqpoint{9.199561in}{6.822389in}}%
\pgfpathquadraticcurveto{\pgfqpoint{9.206888in}{6.822389in}}{\pgfqpoint{9.213772in}{6.824182in}}%
\pgfpathquadraticcurveto{\pgfqpoint{9.220656in}{6.825975in}}{\pgfqpoint{9.227452in}{6.829583in}}%
\pgfpathlineto{\pgfqpoint{9.227452in}{6.817542in}}%
\pgfpathquadraticcurveto{\pgfqpoint{9.220590in}{6.814642in}}{\pgfqpoint{9.213396in}{6.813115in}}%
\pgfpathquadraticcurveto{\pgfqpoint{9.206202in}{6.811587in}}{\pgfqpoint{9.198809in}{6.811587in}}%
\pgfpathquadraticcurveto{\pgfqpoint{9.180259in}{6.811587in}}{\pgfqpoint{9.169435in}{6.822367in}}%
\pgfpathquadraticcurveto{\pgfqpoint{9.158611in}{6.833169in}}{\pgfqpoint{9.158611in}{6.851586in}}%
\pgfpathquadraticcurveto{\pgfqpoint{9.158611in}{6.870600in}}{\pgfqpoint{9.168882in}{6.881757in}}%
\pgfpathquadraticcurveto{\pgfqpoint{9.179153in}{6.892935in}}{\pgfqpoint{9.196595in}{6.892935in}}%
\pgfpathquadraticcurveto{\pgfqpoint{9.212223in}{6.892935in}}{\pgfqpoint{9.221320in}{6.882863in}}%
\pgfpathquadraticcurveto{\pgfqpoint{9.230418in}{6.872814in}}{\pgfqpoint{9.230418in}{6.855526in}}%
\pgfpathlineto{\pgfqpoint{9.230418in}{6.855526in}}%
\pgfpathclose%
\pgfpathmoveto{\pgfqpoint{9.217690in}{6.859267in}}%
\pgfpathquadraticcurveto{\pgfqpoint{9.217557in}{6.869693in}}{\pgfqpoint{9.211847in}{6.875913in}}%
\pgfpathquadraticcurveto{\pgfqpoint{9.206136in}{6.882155in}}{\pgfqpoint{9.196728in}{6.882155in}}%
\pgfpathquadraticcurveto{\pgfqpoint{9.186081in}{6.882155in}}{\pgfqpoint{9.179684in}{6.876134in}}%
\pgfpathquadraticcurveto{\pgfqpoint{9.173287in}{6.870113in}}{\pgfqpoint{9.172313in}{6.859178in}}%
\pgfpathlineto{\pgfqpoint{9.217690in}{6.859267in}}%
\pgfpathlineto{\pgfqpoint{9.217690in}{6.859267in}}%
\pgfpathclose%
\pgfpathmoveto{\pgfqpoint{9.242194in}{6.891076in}}%
\pgfpathlineto{\pgfqpoint{9.255675in}{6.891076in}}%
\pgfpathlineto{\pgfqpoint{9.279891in}{6.826064in}}%
\pgfpathlineto{\pgfqpoint{9.304107in}{6.891076in}}%
\pgfpathlineto{\pgfqpoint{9.317587in}{6.891076in}}%
\pgfpathlineto{\pgfqpoint{9.288524in}{6.813602in}}%
\pgfpathlineto{\pgfqpoint{9.271236in}{6.813602in}}%
\pgfpathlineto{\pgfqpoint{9.242194in}{6.891076in}}%
\pgfpathlineto{\pgfqpoint{9.242194in}{6.891076in}}%
\pgfpathclose%
\pgfpathmoveto{\pgfqpoint{9.401436in}{6.855526in}}%
\pgfpathlineto{\pgfqpoint{9.401436in}{6.849306in}}%
\pgfpathlineto{\pgfqpoint{9.342910in}{6.849306in}}%
\pgfpathquadraticcurveto{\pgfqpoint{9.343751in}{6.836158in}}{\pgfqpoint{9.350835in}{6.829274in}}%
\pgfpathquadraticcurveto{\pgfqpoint{9.357918in}{6.822389in}}{\pgfqpoint{9.370580in}{6.822389in}}%
\pgfpathquadraticcurveto{\pgfqpoint{9.377906in}{6.822389in}}{\pgfqpoint{9.384791in}{6.824182in}}%
\pgfpathquadraticcurveto{\pgfqpoint{9.391675in}{6.825975in}}{\pgfqpoint{9.398470in}{6.829583in}}%
\pgfpathlineto{\pgfqpoint{9.398470in}{6.817542in}}%
\pgfpathquadraticcurveto{\pgfqpoint{9.391608in}{6.814642in}}{\pgfqpoint{9.384414in}{6.813115in}}%
\pgfpathquadraticcurveto{\pgfqpoint{9.377220in}{6.811587in}}{\pgfqpoint{9.369827in}{6.811587in}}%
\pgfpathquadraticcurveto{\pgfqpoint{9.351278in}{6.811587in}}{\pgfqpoint{9.340453in}{6.822367in}}%
\pgfpathquadraticcurveto{\pgfqpoint{9.329629in}{6.833169in}}{\pgfqpoint{9.329629in}{6.851586in}}%
\pgfpathquadraticcurveto{\pgfqpoint{9.329629in}{6.870600in}}{\pgfqpoint{9.339900in}{6.881757in}}%
\pgfpathquadraticcurveto{\pgfqpoint{9.350171in}{6.892935in}}{\pgfqpoint{9.367613in}{6.892935in}}%
\pgfpathquadraticcurveto{\pgfqpoint{9.383241in}{6.892935in}}{\pgfqpoint{9.392339in}{6.882863in}}%
\pgfpathquadraticcurveto{\pgfqpoint{9.401436in}{6.872814in}}{\pgfqpoint{9.401436in}{6.855526in}}%
\pgfpathlineto{\pgfqpoint{9.401436in}{6.855526in}}%
\pgfpathclose%
\pgfpathmoveto{\pgfqpoint{9.388709in}{6.859267in}}%
\pgfpathquadraticcurveto{\pgfqpoint{9.388576in}{6.869693in}}{\pgfqpoint{9.382865in}{6.875913in}}%
\pgfpathquadraticcurveto{\pgfqpoint{9.377154in}{6.882155in}}{\pgfqpoint{9.367746in}{6.882155in}}%
\pgfpathquadraticcurveto{\pgfqpoint{9.357099in}{6.882155in}}{\pgfqpoint{9.350702in}{6.876134in}}%
\pgfpathquadraticcurveto{\pgfqpoint{9.344305in}{6.870113in}}{\pgfqpoint{9.343331in}{6.859178in}}%
\pgfpathlineto{\pgfqpoint{9.388709in}{6.859267in}}%
\pgfpathlineto{\pgfqpoint{9.388709in}{6.859267in}}%
\pgfpathclose%
\pgfpathmoveto{\pgfqpoint{9.467223in}{6.879189in}}%
\pgfpathquadraticcurveto{\pgfqpoint{9.465076in}{6.880428in}}{\pgfqpoint{9.462552in}{6.881004in}}%
\pgfpathquadraticcurveto{\pgfqpoint{9.460029in}{6.881602in}}{\pgfqpoint{9.456996in}{6.881602in}}%
\pgfpathquadraticcurveto{\pgfqpoint{9.446194in}{6.881602in}}{\pgfqpoint{9.440417in}{6.874585in}}%
\pgfpathquadraticcurveto{\pgfqpoint{9.434640in}{6.867568in}}{\pgfqpoint{9.434640in}{6.854419in}}%
\pgfpathlineto{\pgfqpoint{9.434640in}{6.813602in}}%
\pgfpathlineto{\pgfqpoint{9.421845in}{6.813602in}}%
\pgfpathlineto{\pgfqpoint{9.421845in}{6.891076in}}%
\pgfpathlineto{\pgfqpoint{9.434640in}{6.891076in}}%
\pgfpathlineto{\pgfqpoint{9.434640in}{6.879034in}}%
\pgfpathquadraticcurveto{\pgfqpoint{9.438668in}{6.886095in}}{\pgfqpoint{9.445087in}{6.889504in}}%
\pgfpathquadraticcurveto{\pgfqpoint{9.451529in}{6.892935in}}{\pgfqpoint{9.460737in}{6.892935in}}%
\pgfpathquadraticcurveto{\pgfqpoint{9.462043in}{6.892935in}}{\pgfqpoint{9.463637in}{6.892758in}}%
\pgfpathquadraticcurveto{\pgfqpoint{9.465231in}{6.892603in}}{\pgfqpoint{9.467156in}{6.892249in}}%
\pgfpathlineto{\pgfqpoint{9.467223in}{6.879189in}}%
\pgfpathlineto{\pgfqpoint{9.467223in}{6.879189in}}%
\pgfpathclose%
\pgfpathmoveto{\pgfqpoint{9.529955in}{6.888796in}}%
\pgfpathlineto{\pgfqpoint{9.529955in}{6.876754in}}%
\pgfpathquadraticcurveto{\pgfqpoint{9.524576in}{6.879521in}}{\pgfqpoint{9.518754in}{6.880893in}}%
\pgfpathquadraticcurveto{\pgfqpoint{9.512955in}{6.882288in}}{\pgfqpoint{9.506712in}{6.882288in}}%
\pgfpathquadraticcurveto{\pgfqpoint{9.497238in}{6.882288in}}{\pgfqpoint{9.492501in}{6.879388in}}%
\pgfpathquadraticcurveto{\pgfqpoint{9.487765in}{6.876488in}}{\pgfqpoint{9.487765in}{6.870667in}}%
\pgfpathquadraticcurveto{\pgfqpoint{9.487765in}{6.866240in}}{\pgfqpoint{9.491151in}{6.863716in}}%
\pgfpathquadraticcurveto{\pgfqpoint{9.494538in}{6.861193in}}{\pgfqpoint{9.504787in}{6.858913in}}%
\pgfpathlineto{\pgfqpoint{9.509147in}{6.857939in}}%
\pgfpathquadraticcurveto{\pgfqpoint{9.522694in}{6.855039in}}{\pgfqpoint{9.528405in}{6.849749in}}%
\pgfpathquadraticcurveto{\pgfqpoint{9.534116in}{6.844458in}}{\pgfqpoint{9.534116in}{6.834984in}}%
\pgfpathquadraticcurveto{\pgfqpoint{9.534116in}{6.824182in}}{\pgfqpoint{9.525572in}{6.817874in}}%
\pgfpathquadraticcurveto{\pgfqpoint{9.517028in}{6.811587in}}{\pgfqpoint{9.502086in}{6.811587in}}%
\pgfpathquadraticcurveto{\pgfqpoint{9.495866in}{6.811587in}}{\pgfqpoint{9.489115in}{6.812805in}}%
\pgfpathquadraticcurveto{\pgfqpoint{9.482363in}{6.814022in}}{\pgfqpoint{9.474904in}{6.816435in}}%
\pgfpathlineto{\pgfqpoint{9.474904in}{6.829583in}}%
\pgfpathquadraticcurveto{\pgfqpoint{9.481965in}{6.825909in}}{\pgfqpoint{9.488805in}{6.824072in}}%
\pgfpathquadraticcurveto{\pgfqpoint{9.495645in}{6.822257in}}{\pgfqpoint{9.502374in}{6.822257in}}%
\pgfpathquadraticcurveto{\pgfqpoint{9.511361in}{6.822257in}}{\pgfqpoint{9.516186in}{6.825333in}}%
\pgfpathquadraticcurveto{\pgfqpoint{9.521034in}{6.828410in}}{\pgfqpoint{9.521034in}{6.834011in}}%
\pgfpathquadraticcurveto{\pgfqpoint{9.521034in}{6.839190in}}{\pgfqpoint{9.517537in}{6.841957in}}%
\pgfpathquadraticcurveto{\pgfqpoint{9.514061in}{6.844724in}}{\pgfqpoint{9.502219in}{6.847292in}}%
\pgfpathlineto{\pgfqpoint{9.497792in}{6.848332in}}%
\pgfpathquadraticcurveto{\pgfqpoint{9.485972in}{6.850811in}}{\pgfqpoint{9.480703in}{6.855969in}}%
\pgfpathquadraticcurveto{\pgfqpoint{9.475457in}{6.861126in}}{\pgfqpoint{9.475457in}{6.870113in}}%
\pgfpathquadraticcurveto{\pgfqpoint{9.475457in}{6.881048in}}{\pgfqpoint{9.483205in}{6.886981in}}%
\pgfpathquadraticcurveto{\pgfqpoint{9.490952in}{6.892935in}}{\pgfqpoint{9.505207in}{6.892935in}}%
\pgfpathquadraticcurveto{\pgfqpoint{9.512246in}{6.892935in}}{\pgfqpoint{9.518466in}{6.891895in}}%
\pgfpathquadraticcurveto{\pgfqpoint{9.524709in}{6.890876in}}{\pgfqpoint{9.529955in}{6.888796in}}%
\pgfpathlineto{\pgfqpoint{9.529955in}{6.888796in}}%
\pgfpathclose%
\pgfpathmoveto{\pgfqpoint{9.589587in}{6.852538in}}%
\pgfpathquadraticcurveto{\pgfqpoint{9.574159in}{6.852538in}}{\pgfqpoint{9.568205in}{6.849018in}}%
\pgfpathquadraticcurveto{\pgfqpoint{9.562250in}{6.845499in}}{\pgfqpoint{9.562250in}{6.836977in}}%
\pgfpathquadraticcurveto{\pgfqpoint{9.562250in}{6.830203in}}{\pgfqpoint{9.566722in}{6.826219in}}%
\pgfpathquadraticcurveto{\pgfqpoint{9.571193in}{6.822257in}}{\pgfqpoint{9.578852in}{6.822257in}}%
\pgfpathquadraticcurveto{\pgfqpoint{9.589455in}{6.822257in}}{\pgfqpoint{9.595852in}{6.829761in}}%
\pgfpathquadraticcurveto{\pgfqpoint{9.602249in}{6.837264in}}{\pgfqpoint{9.602249in}{6.849705in}}%
\pgfpathlineto{\pgfqpoint{9.602249in}{6.852538in}}%
\pgfpathlineto{\pgfqpoint{9.589587in}{6.852538in}}%
\pgfpathlineto{\pgfqpoint{9.589587in}{6.852538in}}%
\pgfpathclose%
\pgfpathmoveto{\pgfqpoint{9.614977in}{6.857806in}}%
\pgfpathlineto{\pgfqpoint{9.614977in}{6.813602in}}%
\pgfpathlineto{\pgfqpoint{9.602249in}{6.813602in}}%
\pgfpathlineto{\pgfqpoint{9.602249in}{6.825356in}}%
\pgfpathquadraticcurveto{\pgfqpoint{9.597888in}{6.818317in}}{\pgfqpoint{9.591380in}{6.814952in}}%
\pgfpathquadraticcurveto{\pgfqpoint{9.584873in}{6.811587in}}{\pgfqpoint{9.575465in}{6.811587in}}%
\pgfpathquadraticcurveto{\pgfqpoint{9.563578in}{6.811587in}}{\pgfqpoint{9.556539in}{6.818272in}}%
\pgfpathquadraticcurveto{\pgfqpoint{9.549522in}{6.824957in}}{\pgfqpoint{9.549522in}{6.836158in}}%
\pgfpathquadraticcurveto{\pgfqpoint{9.549522in}{6.849218in}}{\pgfqpoint{9.558266in}{6.855858in}}%
\pgfpathquadraticcurveto{\pgfqpoint{9.567031in}{6.862499in}}{\pgfqpoint{9.584386in}{6.862499in}}%
\pgfpathlineto{\pgfqpoint{9.602249in}{6.862499in}}%
\pgfpathlineto{\pgfqpoint{9.602249in}{6.863761in}}%
\pgfpathquadraticcurveto{\pgfqpoint{9.602249in}{6.872548in}}{\pgfqpoint{9.596472in}{6.877352in}}%
\pgfpathquadraticcurveto{\pgfqpoint{9.590694in}{6.882155in}}{\pgfqpoint{9.580246in}{6.882155in}}%
\pgfpathquadraticcurveto{\pgfqpoint{9.573606in}{6.882155in}}{\pgfqpoint{9.567297in}{6.880561in}}%
\pgfpathquadraticcurveto{\pgfqpoint{9.561011in}{6.878968in}}{\pgfqpoint{9.555211in}{6.875780in}}%
\pgfpathlineto{\pgfqpoint{9.555211in}{6.887556in}}%
\pgfpathquadraticcurveto{\pgfqpoint{9.562184in}{6.890257in}}{\pgfqpoint{9.568758in}{6.891585in}}%
\pgfpathquadraticcurveto{\pgfqpoint{9.575332in}{6.892935in}}{\pgfqpoint{9.581552in}{6.892935in}}%
\pgfpathquadraticcurveto{\pgfqpoint{9.598375in}{6.892935in}}{\pgfqpoint{9.606676in}{6.884214in}}%
\pgfpathquadraticcurveto{\pgfqpoint{9.614977in}{6.875514in}}{\pgfqpoint{9.614977in}{6.857806in}}%
\pgfpathlineto{\pgfqpoint{9.614977in}{6.857806in}}%
\pgfpathclose%
\pgfpathmoveto{\pgfqpoint{9.641185in}{6.921246in}}%
\pgfpathlineto{\pgfqpoint{9.653913in}{6.921246in}}%
\pgfpathlineto{\pgfqpoint{9.653913in}{6.813602in}}%
\pgfpathlineto{\pgfqpoint{9.641185in}{6.813602in}}%
\pgfpathlineto{\pgfqpoint{9.641185in}{6.921246in}}%
\pgfpathlineto{\pgfqpoint{9.641185in}{6.921246in}}%
\pgfpathclose%
\pgfpathmoveto{\pgfqpoint{9.683796in}{6.831177in}}%
\pgfpathlineto{\pgfqpoint{9.698383in}{6.831177in}}%
\pgfpathlineto{\pgfqpoint{9.698383in}{6.813602in}}%
\pgfpathlineto{\pgfqpoint{9.683796in}{6.813602in}}%
\pgfpathlineto{\pgfqpoint{9.683796in}{6.831177in}}%
\pgfpathlineto{\pgfqpoint{9.683796in}{6.831177in}}%
\pgfpathclose%
\pgfpathmoveto{\pgfqpoint{9.683796in}{6.886848in}}%
\pgfpathlineto{\pgfqpoint{9.698383in}{6.886848in}}%
\pgfpathlineto{\pgfqpoint{9.698383in}{6.869294in}}%
\pgfpathlineto{\pgfqpoint{9.683796in}{6.869294in}}%
\pgfpathlineto{\pgfqpoint{9.683796in}{6.886848in}}%
\pgfpathlineto{\pgfqpoint{9.683796in}{6.886848in}}%
\pgfpathclose%
\pgfpathmoveto{\pgfqpoint{9.777517in}{6.825356in}}%
\pgfpathlineto{\pgfqpoint{9.800339in}{6.825356in}}%
\pgfpathlineto{\pgfqpoint{9.800339in}{6.904158in}}%
\pgfpathlineto{\pgfqpoint{9.775503in}{6.899177in}}%
\pgfpathlineto{\pgfqpoint{9.775503in}{6.911905in}}%
\pgfpathlineto{\pgfqpoint{9.800206in}{6.916886in}}%
\pgfpathlineto{\pgfqpoint{9.814173in}{6.916886in}}%
\pgfpathlineto{\pgfqpoint{9.814173in}{6.825356in}}%
\pgfpathlineto{\pgfqpoint{9.836995in}{6.825356in}}%
\pgfpathlineto{\pgfqpoint{9.836995in}{6.813602in}}%
\pgfpathlineto{\pgfqpoint{9.777517in}{6.813602in}}%
\pgfpathlineto{\pgfqpoint{9.777517in}{6.825356in}}%
\pgfpathlineto{\pgfqpoint{9.777517in}{6.825356in}}%
\pgfpathclose%
\pgfpathmoveto{\pgfqpoint{9.865218in}{6.831177in}}%
\pgfpathlineto{\pgfqpoint{9.879827in}{6.831177in}}%
\pgfpathlineto{\pgfqpoint{9.879827in}{6.813602in}}%
\pgfpathlineto{\pgfqpoint{9.865218in}{6.813602in}}%
\pgfpathlineto{\pgfqpoint{9.865218in}{6.831177in}}%
\pgfpathlineto{\pgfqpoint{9.865218in}{6.831177in}}%
\pgfpathclose%
\pgfpathmoveto{\pgfqpoint{9.940124in}{6.907677in}}%
\pgfpathquadraticcurveto{\pgfqpoint{9.929344in}{6.907677in}}{\pgfqpoint{9.923899in}{6.897052in}}%
\pgfpathquadraticcurveto{\pgfqpoint{9.918475in}{6.886449in}}{\pgfqpoint{9.918475in}{6.865133in}}%
\pgfpathquadraticcurveto{\pgfqpoint{9.918475in}{6.843905in}}{\pgfqpoint{9.923899in}{6.833280in}}%
\pgfpathquadraticcurveto{\pgfqpoint{9.929344in}{6.822655in}}{\pgfqpoint{9.940124in}{6.822655in}}%
\pgfpathquadraticcurveto{\pgfqpoint{9.950992in}{6.822655in}}{\pgfqpoint{9.956416in}{6.833280in}}%
\pgfpathquadraticcurveto{\pgfqpoint{9.961861in}{6.843905in}}{\pgfqpoint{9.961861in}{6.865133in}}%
\pgfpathquadraticcurveto{\pgfqpoint{9.961861in}{6.886449in}}{\pgfqpoint{9.956416in}{6.897052in}}%
\pgfpathquadraticcurveto{\pgfqpoint{9.950992in}{6.907677in}}{\pgfqpoint{9.940124in}{6.907677in}}%
\pgfpathlineto{\pgfqpoint{9.940124in}{6.907677in}}%
\pgfpathclose%
\pgfpathmoveto{\pgfqpoint{9.940124in}{6.918745in}}%
\pgfpathquadraticcurveto{\pgfqpoint{9.957500in}{6.918745in}}{\pgfqpoint{9.966664in}{6.904999in}}%
\pgfpathquadraticcurveto{\pgfqpoint{9.975828in}{6.891275in}}{\pgfqpoint{9.975828in}{6.865133in}}%
\pgfpathquadraticcurveto{\pgfqpoint{9.975828in}{6.839057in}}{\pgfqpoint{9.966664in}{6.825311in}}%
\pgfpathquadraticcurveto{\pgfqpoint{9.957500in}{6.811587in}}{\pgfqpoint{9.940124in}{6.811587in}}%
\pgfpathquadraticcurveto{\pgfqpoint{9.922770in}{6.811587in}}{\pgfqpoint{9.913606in}{6.825311in}}%
\pgfpathquadraticcurveto{\pgfqpoint{9.904442in}{6.839057in}}{\pgfqpoint{9.904442in}{6.865133in}}%
\pgfpathquadraticcurveto{\pgfqpoint{9.904442in}{6.891275in}}{\pgfqpoint{9.913606in}{6.904999in}}%
\pgfpathquadraticcurveto{\pgfqpoint{9.922770in}{6.918745in}}{\pgfqpoint{9.940124in}{6.918745in}}%
\pgfpathlineto{\pgfqpoint{9.940124in}{6.918745in}}%
\pgfpathclose%
\pgfpathmoveto{\pgfqpoint{10.088232in}{6.859046in}}%
\pgfpathquadraticcurveto{\pgfqpoint{10.082211in}{6.859046in}}{\pgfqpoint{10.078780in}{6.853932in}}%
\pgfpathquadraticcurveto{\pgfqpoint{10.075371in}{6.848819in}}{\pgfqpoint{10.075371in}{6.839677in}}%
\pgfpathquadraticcurveto{\pgfqpoint{10.075371in}{6.830690in}}{\pgfqpoint{10.078780in}{6.825533in}}%
\pgfpathquadraticcurveto{\pgfqpoint{10.082211in}{6.820375in}}{\pgfqpoint{10.088232in}{6.820375in}}%
\pgfpathquadraticcurveto{\pgfqpoint{10.094120in}{6.820375in}}{\pgfqpoint{10.097529in}{6.825533in}}%
\pgfpathquadraticcurveto{\pgfqpoint{10.100960in}{6.830690in}}{\pgfqpoint{10.100960in}{6.839677in}}%
\pgfpathquadraticcurveto{\pgfqpoint{10.100960in}{6.848753in}}{\pgfqpoint{10.097529in}{6.853888in}}%
\pgfpathquadraticcurveto{\pgfqpoint{10.094120in}{6.859046in}}{\pgfqpoint{10.088232in}{6.859046in}}%
\pgfpathlineto{\pgfqpoint{10.088232in}{6.859046in}}%
\pgfpathclose%
\pgfpathmoveto{\pgfqpoint{10.088232in}{6.867833in}}%
\pgfpathquadraticcurveto{\pgfqpoint{10.099167in}{6.867833in}}{\pgfqpoint{10.105586in}{6.860219in}}%
\pgfpathquadraticcurveto{\pgfqpoint{10.112028in}{6.852626in}}{\pgfqpoint{10.112028in}{6.839677in}}%
\pgfpathquadraticcurveto{\pgfqpoint{10.112028in}{6.826750in}}{\pgfqpoint{10.105564in}{6.819158in}}%
\pgfpathquadraticcurveto{\pgfqpoint{10.099100in}{6.811587in}}{\pgfqpoint{10.088232in}{6.811587in}}%
\pgfpathquadraticcurveto{\pgfqpoint{10.077164in}{6.811587in}}{\pgfqpoint{10.070723in}{6.819158in}}%
\pgfpathquadraticcurveto{\pgfqpoint{10.064304in}{6.826750in}}{\pgfqpoint{10.064304in}{6.839677in}}%
\pgfpathquadraticcurveto{\pgfqpoint{10.064304in}{6.852693in}}{\pgfqpoint{10.070767in}{6.860263in}}%
\pgfpathquadraticcurveto{\pgfqpoint{10.077231in}{6.867833in}}{\pgfqpoint{10.088232in}{6.867833in}}%
\pgfpathlineto{\pgfqpoint{10.088232in}{6.867833in}}%
\pgfpathclose%
\pgfpathmoveto{\pgfqpoint{10.016845in}{6.909957in}}%
\pgfpathquadraticcurveto{\pgfqpoint{10.010891in}{6.909957in}}{\pgfqpoint{10.007460in}{6.904800in}}%
\pgfpathquadraticcurveto{\pgfqpoint{10.004051in}{6.899664in}}{\pgfqpoint{10.004051in}{6.890655in}}%
\pgfpathquadraticcurveto{\pgfqpoint{10.004051in}{6.881535in}}{\pgfqpoint{10.007438in}{6.876400in}}%
\pgfpathquadraticcurveto{\pgfqpoint{10.010824in}{6.871287in}}{\pgfqpoint{10.016845in}{6.871287in}}%
\pgfpathquadraticcurveto{\pgfqpoint{10.022866in}{6.871287in}}{\pgfqpoint{10.026275in}{6.876400in}}%
\pgfpathquadraticcurveto{\pgfqpoint{10.029706in}{6.881535in}}{\pgfqpoint{10.029706in}{6.890655in}}%
\pgfpathquadraticcurveto{\pgfqpoint{10.029706in}{6.899576in}}{\pgfqpoint{10.026253in}{6.904755in}}%
\pgfpathquadraticcurveto{\pgfqpoint{10.022800in}{6.909957in}}{\pgfqpoint{10.016845in}{6.909957in}}%
\pgfpathlineto{\pgfqpoint{10.016845in}{6.909957in}}%
\pgfpathclose%
\pgfpathmoveto{\pgfqpoint{10.079311in}{6.918745in}}%
\pgfpathlineto{\pgfqpoint{10.090379in}{6.918745in}}%
\pgfpathlineto{\pgfqpoint{10.025766in}{6.811587in}}%
\pgfpathlineto{\pgfqpoint{10.014698in}{6.811587in}}%
\pgfpathlineto{\pgfqpoint{10.079311in}{6.918745in}}%
\pgfpathlineto{\pgfqpoint{10.079311in}{6.918745in}}%
\pgfpathclose%
\pgfpathmoveto{\pgfqpoint{10.016845in}{6.918745in}}%
\pgfpathquadraticcurveto{\pgfqpoint{10.027780in}{6.918745in}}{\pgfqpoint{10.034266in}{6.911175in}}%
\pgfpathquadraticcurveto{\pgfqpoint{10.040774in}{6.903604in}}{\pgfqpoint{10.040774in}{6.890655in}}%
\pgfpathquadraticcurveto{\pgfqpoint{10.040774in}{6.877595in}}{\pgfqpoint{10.034310in}{6.870047in}}%
\pgfpathquadraticcurveto{\pgfqpoint{10.027847in}{6.862499in}}{\pgfqpoint{10.016845in}{6.862499in}}%
\pgfpathquadraticcurveto{\pgfqpoint{10.005844in}{6.862499in}}{\pgfqpoint{9.999447in}{6.870069in}}%
\pgfpathquadraticcurveto{\pgfqpoint{9.993050in}{6.877662in}}{\pgfqpoint{9.993050in}{6.890655in}}%
\pgfpathquadraticcurveto{\pgfqpoint{9.993050in}{6.903538in}}{\pgfqpoint{9.999469in}{6.911130in}}%
\pgfpathquadraticcurveto{\pgfqpoint{10.005910in}{6.918745in}}{\pgfqpoint{10.016845in}{6.918745in}}%
\pgfpathlineto{\pgfqpoint{10.016845in}{6.918745in}}%
\pgfpathclose%
\pgfusepath{fill}%
\end{pgfscope}%
\begin{pgfscope}%
\pgfsetbuttcap%
\pgfsetroundjoin%
\definecolor{currentfill}{rgb}{0.541176,0.380392,0.658824}%
\pgfsetfillcolor{currentfill}%
\pgfsetlinewidth{1.003750pt}%
\definecolor{currentstroke}{rgb}{0.000000,0.000000,0.000000}%
\pgfsetstrokecolor{currentstroke}%
\pgfsetstrokeopacity{0.000000}%
\pgfsetdash{}{0pt}%
\pgfsys@defobject{currentmarker}{\pgfqpoint{-0.038194in}{-0.038194in}}{\pgfqpoint{0.038194in}{0.038194in}}{%
\pgfpathmoveto{\pgfqpoint{0.000000in}{-0.038194in}}%
\pgfpathcurveto{\pgfqpoint{0.010129in}{-0.038194in}}{\pgfqpoint{0.019845in}{-0.034170in}}{\pgfqpoint{0.027008in}{-0.027008in}}%
\pgfpathcurveto{\pgfqpoint{0.034170in}{-0.019845in}}{\pgfqpoint{0.038194in}{-0.010129in}}{\pgfqpoint{0.038194in}{0.000000in}}%
\pgfpathcurveto{\pgfqpoint{0.038194in}{0.010129in}}{\pgfqpoint{0.034170in}{0.019845in}}{\pgfqpoint{0.027008in}{0.027008in}}%
\pgfpathcurveto{\pgfqpoint{0.019845in}{0.034170in}}{\pgfqpoint{0.010129in}{0.038194in}}{\pgfqpoint{0.000000in}{0.038194in}}%
\pgfpathcurveto{\pgfqpoint{-0.010129in}{0.038194in}}{\pgfqpoint{-0.019845in}{0.034170in}}{\pgfqpoint{-0.027008in}{0.027008in}}%
\pgfpathcurveto{\pgfqpoint{-0.034170in}{0.019845in}}{\pgfqpoint{-0.038194in}{0.010129in}}{\pgfqpoint{-0.038194in}{0.000000in}}%
\pgfpathcurveto{\pgfqpoint{-0.038194in}{-0.010129in}}{\pgfqpoint{-0.034170in}{-0.019845in}}{\pgfqpoint{-0.027008in}{-0.027008in}}%
\pgfpathcurveto{\pgfqpoint{-0.019845in}{-0.034170in}}{\pgfqpoint{-0.010129in}{-0.038194in}}{\pgfqpoint{0.000000in}{-0.038194in}}%
\pgfpathlineto{\pgfqpoint{0.000000in}{-0.038194in}}%
\pgfpathclose%
\pgfusepath{stroke,fill}%
}%
\begin{pgfscope}%
\pgfsys@transformshift{1.078000in}{2.265600in}%
\pgfsys@useobject{currentmarker}{}%
\end{pgfscope}%
\end{pgfscope}%
\begin{pgfscope}%
\pgfsetbuttcap%
\pgfsetroundjoin%
\definecolor{currentfill}{rgb}{0.709804,0.741176,0.207843}%
\pgfsetfillcolor{currentfill}%
\pgfsetlinewidth{1.003750pt}%
\definecolor{currentstroke}{rgb}{0.000000,0.000000,0.000000}%
\pgfsetstrokecolor{currentstroke}%
\pgfsetstrokeopacity{0.000000}%
\pgfsetdash{}{0pt}%
\pgfsys@defobject{currentmarker}{\pgfqpoint{-0.038194in}{-0.038194in}}{\pgfqpoint{0.038194in}{0.038194in}}{%
\pgfpathmoveto{\pgfqpoint{0.000000in}{-0.038194in}}%
\pgfpathcurveto{\pgfqpoint{0.010129in}{-0.038194in}}{\pgfqpoint{0.019845in}{-0.034170in}}{\pgfqpoint{0.027008in}{-0.027008in}}%
\pgfpathcurveto{\pgfqpoint{0.034170in}{-0.019845in}}{\pgfqpoint{0.038194in}{-0.010129in}}{\pgfqpoint{0.038194in}{0.000000in}}%
\pgfpathcurveto{\pgfqpoint{0.038194in}{0.010129in}}{\pgfqpoint{0.034170in}{0.019845in}}{\pgfqpoint{0.027008in}{0.027008in}}%
\pgfpathcurveto{\pgfqpoint{0.019845in}{0.034170in}}{\pgfqpoint{0.010129in}{0.038194in}}{\pgfqpoint{0.000000in}{0.038194in}}%
\pgfpathcurveto{\pgfqpoint{-0.010129in}{0.038194in}}{\pgfqpoint{-0.019845in}{0.034170in}}{\pgfqpoint{-0.027008in}{0.027008in}}%
\pgfpathcurveto{\pgfqpoint{-0.034170in}{0.019845in}}{\pgfqpoint{-0.038194in}{0.010129in}}{\pgfqpoint{-0.038194in}{0.000000in}}%
\pgfpathcurveto{\pgfqpoint{-0.038194in}{-0.010129in}}{\pgfqpoint{-0.034170in}{-0.019845in}}{\pgfqpoint{-0.027008in}{-0.027008in}}%
\pgfpathcurveto{\pgfqpoint{-0.019845in}{-0.034170in}}{\pgfqpoint{-0.010129in}{-0.038194in}}{\pgfqpoint{0.000000in}{-0.038194in}}%
\pgfpathlineto{\pgfqpoint{0.000000in}{-0.038194in}}%
\pgfpathclose%
\pgfusepath{stroke,fill}%
}%
\begin{pgfscope}%
\pgfsys@transformshift{1.078000in}{2.088960in}%
\pgfsys@useobject{currentmarker}{}%
\end{pgfscope}%
\end{pgfscope}%
\begin{pgfscope}%
\pgfsetbuttcap%
\pgfsetroundjoin%
\definecolor{currentfill}{rgb}{0.000000,0.447059,0.698039}%
\pgfsetfillcolor{currentfill}%
\pgfsetlinewidth{1.003750pt}%
\definecolor{currentstroke}{rgb}{0.000000,0.000000,0.000000}%
\pgfsetstrokecolor{currentstroke}%
\pgfsetstrokeopacity{0.000000}%
\pgfsetdash{}{0pt}%
\pgfsys@defobject{currentmarker}{\pgfqpoint{-0.038194in}{-0.038194in}}{\pgfqpoint{0.038194in}{0.038194in}}{%
\pgfpathmoveto{\pgfqpoint{0.000000in}{-0.038194in}}%
\pgfpathcurveto{\pgfqpoint{0.010129in}{-0.038194in}}{\pgfqpoint{0.019845in}{-0.034170in}}{\pgfqpoint{0.027008in}{-0.027008in}}%
\pgfpathcurveto{\pgfqpoint{0.034170in}{-0.019845in}}{\pgfqpoint{0.038194in}{-0.010129in}}{\pgfqpoint{0.038194in}{0.000000in}}%
\pgfpathcurveto{\pgfqpoint{0.038194in}{0.010129in}}{\pgfqpoint{0.034170in}{0.019845in}}{\pgfqpoint{0.027008in}{0.027008in}}%
\pgfpathcurveto{\pgfqpoint{0.019845in}{0.034170in}}{\pgfqpoint{0.010129in}{0.038194in}}{\pgfqpoint{0.000000in}{0.038194in}}%
\pgfpathcurveto{\pgfqpoint{-0.010129in}{0.038194in}}{\pgfqpoint{-0.019845in}{0.034170in}}{\pgfqpoint{-0.027008in}{0.027008in}}%
\pgfpathcurveto{\pgfqpoint{-0.034170in}{0.019845in}}{\pgfqpoint{-0.038194in}{0.010129in}}{\pgfqpoint{-0.038194in}{0.000000in}}%
\pgfpathcurveto{\pgfqpoint{-0.038194in}{-0.010129in}}{\pgfqpoint{-0.034170in}{-0.019845in}}{\pgfqpoint{-0.027008in}{-0.027008in}}%
\pgfpathcurveto{\pgfqpoint{-0.019845in}{-0.034170in}}{\pgfqpoint{-0.010129in}{-0.038194in}}{\pgfqpoint{0.000000in}{-0.038194in}}%
\pgfpathlineto{\pgfqpoint{0.000000in}{-0.038194in}}%
\pgfpathclose%
\pgfusepath{stroke,fill}%
}%
\begin{pgfscope}%
\pgfsys@transformshift{1.078000in}{1.912320in}%
\pgfsys@useobject{currentmarker}{}%
\end{pgfscope}%
\end{pgfscope}%
\begin{pgfscope}%
\pgfsetbuttcap%
\pgfsetroundjoin%
\definecolor{currentfill}{rgb}{0.149020,0.274510,0.325490}%
\pgfsetfillcolor{currentfill}%
\pgfsetlinewidth{1.003750pt}%
\definecolor{currentstroke}{rgb}{0.000000,0.000000,0.000000}%
\pgfsetstrokecolor{currentstroke}%
\pgfsetstrokeopacity{0.000000}%
\pgfsetdash{}{0pt}%
\pgfsys@defobject{currentmarker}{\pgfqpoint{-0.038194in}{-0.038194in}}{\pgfqpoint{0.038194in}{0.038194in}}{%
\pgfpathmoveto{\pgfqpoint{0.000000in}{-0.038194in}}%
\pgfpathcurveto{\pgfqpoint{0.010129in}{-0.038194in}}{\pgfqpoint{0.019845in}{-0.034170in}}{\pgfqpoint{0.027008in}{-0.027008in}}%
\pgfpathcurveto{\pgfqpoint{0.034170in}{-0.019845in}}{\pgfqpoint{0.038194in}{-0.010129in}}{\pgfqpoint{0.038194in}{0.000000in}}%
\pgfpathcurveto{\pgfqpoint{0.038194in}{0.010129in}}{\pgfqpoint{0.034170in}{0.019845in}}{\pgfqpoint{0.027008in}{0.027008in}}%
\pgfpathcurveto{\pgfqpoint{0.019845in}{0.034170in}}{\pgfqpoint{0.010129in}{0.038194in}}{\pgfqpoint{0.000000in}{0.038194in}}%
\pgfpathcurveto{\pgfqpoint{-0.010129in}{0.038194in}}{\pgfqpoint{-0.019845in}{0.034170in}}{\pgfqpoint{-0.027008in}{0.027008in}}%
\pgfpathcurveto{\pgfqpoint{-0.034170in}{0.019845in}}{\pgfqpoint{-0.038194in}{0.010129in}}{\pgfqpoint{-0.038194in}{0.000000in}}%
\pgfpathcurveto{\pgfqpoint{-0.038194in}{-0.010129in}}{\pgfqpoint{-0.034170in}{-0.019845in}}{\pgfqpoint{-0.027008in}{-0.027008in}}%
\pgfpathcurveto{\pgfqpoint{-0.019845in}{-0.034170in}}{\pgfqpoint{-0.010129in}{-0.038194in}}{\pgfqpoint{0.000000in}{-0.038194in}}%
\pgfpathlineto{\pgfqpoint{0.000000in}{-0.038194in}}%
\pgfpathclose%
\pgfusepath{stroke,fill}%
}%
\begin{pgfscope}%
\pgfsys@transformshift{1.078000in}{1.735680in}%
\pgfsys@useobject{currentmarker}{}%
\end{pgfscope}%
\end{pgfscope}%
\begin{pgfscope}%
\pgfsetbuttcap%
\pgfsetroundjoin%
\definecolor{currentfill}{rgb}{0.549020,0.337255,0.294118}%
\pgfsetfillcolor{currentfill}%
\pgfsetlinewidth{1.003750pt}%
\definecolor{currentstroke}{rgb}{0.000000,0.000000,0.000000}%
\pgfsetstrokecolor{currentstroke}%
\pgfsetstrokeopacity{0.000000}%
\pgfsetdash{}{0pt}%
\pgfsys@defobject{currentmarker}{\pgfqpoint{-0.038194in}{-0.038194in}}{\pgfqpoint{0.038194in}{0.038194in}}{%
\pgfpathmoveto{\pgfqpoint{0.000000in}{-0.038194in}}%
\pgfpathcurveto{\pgfqpoint{0.010129in}{-0.038194in}}{\pgfqpoint{0.019845in}{-0.034170in}}{\pgfqpoint{0.027008in}{-0.027008in}}%
\pgfpathcurveto{\pgfqpoint{0.034170in}{-0.019845in}}{\pgfqpoint{0.038194in}{-0.010129in}}{\pgfqpoint{0.038194in}{0.000000in}}%
\pgfpathcurveto{\pgfqpoint{0.038194in}{0.010129in}}{\pgfqpoint{0.034170in}{0.019845in}}{\pgfqpoint{0.027008in}{0.027008in}}%
\pgfpathcurveto{\pgfqpoint{0.019845in}{0.034170in}}{\pgfqpoint{0.010129in}{0.038194in}}{\pgfqpoint{0.000000in}{0.038194in}}%
\pgfpathcurveto{\pgfqpoint{-0.010129in}{0.038194in}}{\pgfqpoint{-0.019845in}{0.034170in}}{\pgfqpoint{-0.027008in}{0.027008in}}%
\pgfpathcurveto{\pgfqpoint{-0.034170in}{0.019845in}}{\pgfqpoint{-0.038194in}{0.010129in}}{\pgfqpoint{-0.038194in}{0.000000in}}%
\pgfpathcurveto{\pgfqpoint{-0.038194in}{-0.010129in}}{\pgfqpoint{-0.034170in}{-0.019845in}}{\pgfqpoint{-0.027008in}{-0.027008in}}%
\pgfpathcurveto{\pgfqpoint{-0.019845in}{-0.034170in}}{\pgfqpoint{-0.010129in}{-0.038194in}}{\pgfqpoint{0.000000in}{-0.038194in}}%
\pgfpathlineto{\pgfqpoint{0.000000in}{-0.038194in}}%
\pgfpathclose%
\pgfusepath{stroke,fill}%
}%
\begin{pgfscope}%
\pgfsys@transformshift{1.078000in}{1.559040in}%
\pgfsys@useobject{currentmarker}{}%
\end{pgfscope}%
\end{pgfscope}%
\begin{pgfscope}%
\pgfsetbuttcap%
\pgfsetroundjoin%
\definecolor{currentfill}{rgb}{0.835294,0.368627,0.000000}%
\pgfsetfillcolor{currentfill}%
\pgfsetlinewidth{1.003750pt}%
\definecolor{currentstroke}{rgb}{0.000000,0.000000,0.000000}%
\pgfsetstrokecolor{currentstroke}%
\pgfsetstrokeopacity{0.000000}%
\pgfsetdash{}{0pt}%
\pgfsys@defobject{currentmarker}{\pgfqpoint{-0.038194in}{-0.038194in}}{\pgfqpoint{0.038194in}{0.038194in}}{%
\pgfpathmoveto{\pgfqpoint{0.000000in}{-0.038194in}}%
\pgfpathcurveto{\pgfqpoint{0.010129in}{-0.038194in}}{\pgfqpoint{0.019845in}{-0.034170in}}{\pgfqpoint{0.027008in}{-0.027008in}}%
\pgfpathcurveto{\pgfqpoint{0.034170in}{-0.019845in}}{\pgfqpoint{0.038194in}{-0.010129in}}{\pgfqpoint{0.038194in}{0.000000in}}%
\pgfpathcurveto{\pgfqpoint{0.038194in}{0.010129in}}{\pgfqpoint{0.034170in}{0.019845in}}{\pgfqpoint{0.027008in}{0.027008in}}%
\pgfpathcurveto{\pgfqpoint{0.019845in}{0.034170in}}{\pgfqpoint{0.010129in}{0.038194in}}{\pgfqpoint{0.000000in}{0.038194in}}%
\pgfpathcurveto{\pgfqpoint{-0.010129in}{0.038194in}}{\pgfqpoint{-0.019845in}{0.034170in}}{\pgfqpoint{-0.027008in}{0.027008in}}%
\pgfpathcurveto{\pgfqpoint{-0.034170in}{0.019845in}}{\pgfqpoint{-0.038194in}{0.010129in}}{\pgfqpoint{-0.038194in}{0.000000in}}%
\pgfpathcurveto{\pgfqpoint{-0.038194in}{-0.010129in}}{\pgfqpoint{-0.034170in}{-0.019845in}}{\pgfqpoint{-0.027008in}{-0.027008in}}%
\pgfpathcurveto{\pgfqpoint{-0.019845in}{-0.034170in}}{\pgfqpoint{-0.010129in}{-0.038194in}}{\pgfqpoint{0.000000in}{-0.038194in}}%
\pgfpathlineto{\pgfqpoint{0.000000in}{-0.038194in}}%
\pgfpathclose%
\pgfusepath{stroke,fill}%
}%
\begin{pgfscope}%
\pgfsys@transformshift{1.078000in}{1.382400in}%
\pgfsys@useobject{currentmarker}{}%
\end{pgfscope}%
\end{pgfscope}%
\begin{pgfscope}%
\pgfsetbuttcap%
\pgfsetroundjoin%
\definecolor{currentfill}{rgb}{0.662745,0.368627,0.000000}%
\pgfsetfillcolor{currentfill}%
\pgfsetlinewidth{1.003750pt}%
\definecolor{currentstroke}{rgb}{0.000000,0.000000,0.000000}%
\pgfsetstrokecolor{currentstroke}%
\pgfsetstrokeopacity{0.000000}%
\pgfsetdash{}{0pt}%
\pgfsys@defobject{currentmarker}{\pgfqpoint{-0.038194in}{-0.038194in}}{\pgfqpoint{0.038194in}{0.038194in}}{%
\pgfpathmoveto{\pgfqpoint{0.000000in}{-0.038194in}}%
\pgfpathcurveto{\pgfqpoint{0.010129in}{-0.038194in}}{\pgfqpoint{0.019845in}{-0.034170in}}{\pgfqpoint{0.027008in}{-0.027008in}}%
\pgfpathcurveto{\pgfqpoint{0.034170in}{-0.019845in}}{\pgfqpoint{0.038194in}{-0.010129in}}{\pgfqpoint{0.038194in}{0.000000in}}%
\pgfpathcurveto{\pgfqpoint{0.038194in}{0.010129in}}{\pgfqpoint{0.034170in}{0.019845in}}{\pgfqpoint{0.027008in}{0.027008in}}%
\pgfpathcurveto{\pgfqpoint{0.019845in}{0.034170in}}{\pgfqpoint{0.010129in}{0.038194in}}{\pgfqpoint{0.000000in}{0.038194in}}%
\pgfpathcurveto{\pgfqpoint{-0.010129in}{0.038194in}}{\pgfqpoint{-0.019845in}{0.034170in}}{\pgfqpoint{-0.027008in}{0.027008in}}%
\pgfpathcurveto{\pgfqpoint{-0.034170in}{0.019845in}}{\pgfqpoint{-0.038194in}{0.010129in}}{\pgfqpoint{-0.038194in}{0.000000in}}%
\pgfpathcurveto{\pgfqpoint{-0.038194in}{-0.010129in}}{\pgfqpoint{-0.034170in}{-0.019845in}}{\pgfqpoint{-0.027008in}{-0.027008in}}%
\pgfpathcurveto{\pgfqpoint{-0.019845in}{-0.034170in}}{\pgfqpoint{-0.010129in}{-0.038194in}}{\pgfqpoint{0.000000in}{-0.038194in}}%
\pgfpathlineto{\pgfqpoint{0.000000in}{-0.038194in}}%
\pgfpathclose%
\pgfusepath{stroke,fill}%
}%
\begin{pgfscope}%
\pgfsys@transformshift{4.312000in}{2.265600in}%
\pgfsys@useobject{currentmarker}{}%
\end{pgfscope}%
\end{pgfscope}%
\begin{pgfscope}%
\pgfsetbuttcap%
\pgfsetroundjoin%
\definecolor{currentfill}{rgb}{0.337255,0.705882,0.913725}%
\pgfsetfillcolor{currentfill}%
\pgfsetlinewidth{1.003750pt}%
\definecolor{currentstroke}{rgb}{0.000000,0.000000,0.000000}%
\pgfsetstrokecolor{currentstroke}%
\pgfsetstrokeopacity{0.000000}%
\pgfsetdash{}{0pt}%
\pgfsys@defobject{currentmarker}{\pgfqpoint{-0.038194in}{-0.038194in}}{\pgfqpoint{0.038194in}{0.038194in}}{%
\pgfpathmoveto{\pgfqpoint{0.000000in}{-0.038194in}}%
\pgfpathcurveto{\pgfqpoint{0.010129in}{-0.038194in}}{\pgfqpoint{0.019845in}{-0.034170in}}{\pgfqpoint{0.027008in}{-0.027008in}}%
\pgfpathcurveto{\pgfqpoint{0.034170in}{-0.019845in}}{\pgfqpoint{0.038194in}{-0.010129in}}{\pgfqpoint{0.038194in}{0.000000in}}%
\pgfpathcurveto{\pgfqpoint{0.038194in}{0.010129in}}{\pgfqpoint{0.034170in}{0.019845in}}{\pgfqpoint{0.027008in}{0.027008in}}%
\pgfpathcurveto{\pgfqpoint{0.019845in}{0.034170in}}{\pgfqpoint{0.010129in}{0.038194in}}{\pgfqpoint{0.000000in}{0.038194in}}%
\pgfpathcurveto{\pgfqpoint{-0.010129in}{0.038194in}}{\pgfqpoint{-0.019845in}{0.034170in}}{\pgfqpoint{-0.027008in}{0.027008in}}%
\pgfpathcurveto{\pgfqpoint{-0.034170in}{0.019845in}}{\pgfqpoint{-0.038194in}{0.010129in}}{\pgfqpoint{-0.038194in}{0.000000in}}%
\pgfpathcurveto{\pgfqpoint{-0.038194in}{-0.010129in}}{\pgfqpoint{-0.034170in}{-0.019845in}}{\pgfqpoint{-0.027008in}{-0.027008in}}%
\pgfpathcurveto{\pgfqpoint{-0.019845in}{-0.034170in}}{\pgfqpoint{-0.010129in}{-0.038194in}}{\pgfqpoint{0.000000in}{-0.038194in}}%
\pgfpathlineto{\pgfqpoint{0.000000in}{-0.038194in}}%
\pgfpathclose%
\pgfusepath{stroke,fill}%
}%
\begin{pgfscope}%
\pgfsys@transformshift{4.312000in}{2.088960in}%
\pgfsys@useobject{currentmarker}{}%
\end{pgfscope}%
\end{pgfscope}%
\begin{pgfscope}%
\pgfsetbuttcap%
\pgfsetroundjoin%
\definecolor{currentfill}{rgb}{0.000000,0.619608,0.450980}%
\pgfsetfillcolor{currentfill}%
\pgfsetlinewidth{1.003750pt}%
\definecolor{currentstroke}{rgb}{0.000000,0.000000,0.000000}%
\pgfsetstrokecolor{currentstroke}%
\pgfsetstrokeopacity{0.000000}%
\pgfsetdash{}{0pt}%
\pgfsys@defobject{currentmarker}{\pgfqpoint{-0.038194in}{-0.038194in}}{\pgfqpoint{0.038194in}{0.038194in}}{%
\pgfpathmoveto{\pgfqpoint{0.000000in}{-0.038194in}}%
\pgfpathcurveto{\pgfqpoint{0.010129in}{-0.038194in}}{\pgfqpoint{0.019845in}{-0.034170in}}{\pgfqpoint{0.027008in}{-0.027008in}}%
\pgfpathcurveto{\pgfqpoint{0.034170in}{-0.019845in}}{\pgfqpoint{0.038194in}{-0.010129in}}{\pgfqpoint{0.038194in}{0.000000in}}%
\pgfpathcurveto{\pgfqpoint{0.038194in}{0.010129in}}{\pgfqpoint{0.034170in}{0.019845in}}{\pgfqpoint{0.027008in}{0.027008in}}%
\pgfpathcurveto{\pgfqpoint{0.019845in}{0.034170in}}{\pgfqpoint{0.010129in}{0.038194in}}{\pgfqpoint{0.000000in}{0.038194in}}%
\pgfpathcurveto{\pgfqpoint{-0.010129in}{0.038194in}}{\pgfqpoint{-0.019845in}{0.034170in}}{\pgfqpoint{-0.027008in}{0.027008in}}%
\pgfpathcurveto{\pgfqpoint{-0.034170in}{0.019845in}}{\pgfqpoint{-0.038194in}{0.010129in}}{\pgfqpoint{-0.038194in}{0.000000in}}%
\pgfpathcurveto{\pgfqpoint{-0.038194in}{-0.010129in}}{\pgfqpoint{-0.034170in}{-0.019845in}}{\pgfqpoint{-0.027008in}{-0.027008in}}%
\pgfpathcurveto{\pgfqpoint{-0.019845in}{-0.034170in}}{\pgfqpoint{-0.010129in}{-0.038194in}}{\pgfqpoint{0.000000in}{-0.038194in}}%
\pgfpathlineto{\pgfqpoint{0.000000in}{-0.038194in}}%
\pgfpathclose%
\pgfusepath{stroke,fill}%
}%
\begin{pgfscope}%
\pgfsys@transformshift{4.312000in}{1.912320in}%
\pgfsys@useobject{currentmarker}{}%
\end{pgfscope}%
\end{pgfscope}%
\begin{pgfscope}%
\pgfsetbuttcap%
\pgfsetroundjoin%
\definecolor{currentfill}{rgb}{0.709804,0.121569,0.270588}%
\pgfsetfillcolor{currentfill}%
\pgfsetlinewidth{1.003750pt}%
\definecolor{currentstroke}{rgb}{0.000000,0.000000,0.000000}%
\pgfsetstrokecolor{currentstroke}%
\pgfsetstrokeopacity{0.000000}%
\pgfsetdash{}{0pt}%
\pgfsys@defobject{currentmarker}{\pgfqpoint{-0.038194in}{-0.038194in}}{\pgfqpoint{0.038194in}{0.038194in}}{%
\pgfpathmoveto{\pgfqpoint{0.000000in}{-0.038194in}}%
\pgfpathcurveto{\pgfqpoint{0.010129in}{-0.038194in}}{\pgfqpoint{0.019845in}{-0.034170in}}{\pgfqpoint{0.027008in}{-0.027008in}}%
\pgfpathcurveto{\pgfqpoint{0.034170in}{-0.019845in}}{\pgfqpoint{0.038194in}{-0.010129in}}{\pgfqpoint{0.038194in}{0.000000in}}%
\pgfpathcurveto{\pgfqpoint{0.038194in}{0.010129in}}{\pgfqpoint{0.034170in}{0.019845in}}{\pgfqpoint{0.027008in}{0.027008in}}%
\pgfpathcurveto{\pgfqpoint{0.019845in}{0.034170in}}{\pgfqpoint{0.010129in}{0.038194in}}{\pgfqpoint{0.000000in}{0.038194in}}%
\pgfpathcurveto{\pgfqpoint{-0.010129in}{0.038194in}}{\pgfqpoint{-0.019845in}{0.034170in}}{\pgfqpoint{-0.027008in}{0.027008in}}%
\pgfpathcurveto{\pgfqpoint{-0.034170in}{0.019845in}}{\pgfqpoint{-0.038194in}{0.010129in}}{\pgfqpoint{-0.038194in}{0.000000in}}%
\pgfpathcurveto{\pgfqpoint{-0.038194in}{-0.010129in}}{\pgfqpoint{-0.034170in}{-0.019845in}}{\pgfqpoint{-0.027008in}{-0.027008in}}%
\pgfpathcurveto{\pgfqpoint{-0.019845in}{-0.034170in}}{\pgfqpoint{-0.010129in}{-0.038194in}}{\pgfqpoint{0.000000in}{-0.038194in}}%
\pgfpathlineto{\pgfqpoint{0.000000in}{-0.038194in}}%
\pgfpathclose%
\pgfusepath{stroke,fill}%
}%
\begin{pgfscope}%
\pgfsys@transformshift{4.312000in}{1.735680in}%
\pgfsys@useobject{currentmarker}{}%
\end{pgfscope}%
\end{pgfscope}%
\begin{pgfscope}%
\pgfsetbuttcap%
\pgfsetroundjoin%
\definecolor{currentfill}{rgb}{0.058824,0.541176,0.372549}%
\pgfsetfillcolor{currentfill}%
\pgfsetlinewidth{1.003750pt}%
\definecolor{currentstroke}{rgb}{0.000000,0.000000,0.000000}%
\pgfsetstrokecolor{currentstroke}%
\pgfsetstrokeopacity{0.000000}%
\pgfsetdash{}{0pt}%
\pgfsys@defobject{currentmarker}{\pgfqpoint{-0.038194in}{-0.038194in}}{\pgfqpoint{0.038194in}{0.038194in}}{%
\pgfpathmoveto{\pgfqpoint{0.000000in}{-0.038194in}}%
\pgfpathcurveto{\pgfqpoint{0.010129in}{-0.038194in}}{\pgfqpoint{0.019845in}{-0.034170in}}{\pgfqpoint{0.027008in}{-0.027008in}}%
\pgfpathcurveto{\pgfqpoint{0.034170in}{-0.019845in}}{\pgfqpoint{0.038194in}{-0.010129in}}{\pgfqpoint{0.038194in}{0.000000in}}%
\pgfpathcurveto{\pgfqpoint{0.038194in}{0.010129in}}{\pgfqpoint{0.034170in}{0.019845in}}{\pgfqpoint{0.027008in}{0.027008in}}%
\pgfpathcurveto{\pgfqpoint{0.019845in}{0.034170in}}{\pgfqpoint{0.010129in}{0.038194in}}{\pgfqpoint{0.000000in}{0.038194in}}%
\pgfpathcurveto{\pgfqpoint{-0.010129in}{0.038194in}}{\pgfqpoint{-0.019845in}{0.034170in}}{\pgfqpoint{-0.027008in}{0.027008in}}%
\pgfpathcurveto{\pgfqpoint{-0.034170in}{0.019845in}}{\pgfqpoint{-0.038194in}{0.010129in}}{\pgfqpoint{-0.038194in}{0.000000in}}%
\pgfpathcurveto{\pgfqpoint{-0.038194in}{-0.010129in}}{\pgfqpoint{-0.034170in}{-0.019845in}}{\pgfqpoint{-0.027008in}{-0.027008in}}%
\pgfpathcurveto{\pgfqpoint{-0.019845in}{-0.034170in}}{\pgfqpoint{-0.010129in}{-0.038194in}}{\pgfqpoint{0.000000in}{-0.038194in}}%
\pgfpathlineto{\pgfqpoint{0.000000in}{-0.038194in}}%
\pgfpathclose%
\pgfusepath{stroke,fill}%
}%
\begin{pgfscope}%
\pgfsys@transformshift{4.312000in}{1.559040in}%
\pgfsys@useobject{currentmarker}{}%
\end{pgfscope}%
\end{pgfscope}%
\begin{pgfscope}%
\pgfsetbuttcap%
\pgfsetroundjoin%
\definecolor{currentfill}{rgb}{0.901961,0.403922,0.619608}%
\pgfsetfillcolor{currentfill}%
\pgfsetlinewidth{1.003750pt}%
\definecolor{currentstroke}{rgb}{0.000000,0.000000,0.000000}%
\pgfsetstrokecolor{currentstroke}%
\pgfsetstrokeopacity{0.000000}%
\pgfsetdash{}{0pt}%
\pgfsys@defobject{currentmarker}{\pgfqpoint{-0.038194in}{-0.038194in}}{\pgfqpoint{0.038194in}{0.038194in}}{%
\pgfpathmoveto{\pgfqpoint{0.000000in}{-0.038194in}}%
\pgfpathcurveto{\pgfqpoint{0.010129in}{-0.038194in}}{\pgfqpoint{0.019845in}{-0.034170in}}{\pgfqpoint{0.027008in}{-0.027008in}}%
\pgfpathcurveto{\pgfqpoint{0.034170in}{-0.019845in}}{\pgfqpoint{0.038194in}{-0.010129in}}{\pgfqpoint{0.038194in}{0.000000in}}%
\pgfpathcurveto{\pgfqpoint{0.038194in}{0.010129in}}{\pgfqpoint{0.034170in}{0.019845in}}{\pgfqpoint{0.027008in}{0.027008in}}%
\pgfpathcurveto{\pgfqpoint{0.019845in}{0.034170in}}{\pgfqpoint{0.010129in}{0.038194in}}{\pgfqpoint{0.000000in}{0.038194in}}%
\pgfpathcurveto{\pgfqpoint{-0.010129in}{0.038194in}}{\pgfqpoint{-0.019845in}{0.034170in}}{\pgfqpoint{-0.027008in}{0.027008in}}%
\pgfpathcurveto{\pgfqpoint{-0.034170in}{0.019845in}}{\pgfqpoint{-0.038194in}{0.010129in}}{\pgfqpoint{-0.038194in}{0.000000in}}%
\pgfpathcurveto{\pgfqpoint{-0.038194in}{-0.010129in}}{\pgfqpoint{-0.034170in}{-0.019845in}}{\pgfqpoint{-0.027008in}{-0.027008in}}%
\pgfpathcurveto{\pgfqpoint{-0.019845in}{-0.034170in}}{\pgfqpoint{-0.010129in}{-0.038194in}}{\pgfqpoint{0.000000in}{-0.038194in}}%
\pgfpathlineto{\pgfqpoint{0.000000in}{-0.038194in}}%
\pgfpathclose%
\pgfusepath{stroke,fill}%
}%
\begin{pgfscope}%
\pgfsys@transformshift{4.312000in}{1.382400in}%
\pgfsys@useobject{currentmarker}{}%
\end{pgfscope}%
\end{pgfscope}%
\begin{pgfscope}%
\pgfsetbuttcap%
\pgfsetroundjoin%
\definecolor{currentfill}{rgb}{0.478431,0.243137,0.615686}%
\pgfsetfillcolor{currentfill}%
\pgfsetlinewidth{1.003750pt}%
\definecolor{currentstroke}{rgb}{0.000000,0.000000,0.000000}%
\pgfsetstrokecolor{currentstroke}%
\pgfsetstrokeopacity{0.000000}%
\pgfsetdash{}{0pt}%
\pgfsys@defobject{currentmarker}{\pgfqpoint{-0.038194in}{-0.038194in}}{\pgfqpoint{0.038194in}{0.038194in}}{%
\pgfpathmoveto{\pgfqpoint{0.000000in}{-0.038194in}}%
\pgfpathcurveto{\pgfqpoint{0.010129in}{-0.038194in}}{\pgfqpoint{0.019845in}{-0.034170in}}{\pgfqpoint{0.027008in}{-0.027008in}}%
\pgfpathcurveto{\pgfqpoint{0.034170in}{-0.019845in}}{\pgfqpoint{0.038194in}{-0.010129in}}{\pgfqpoint{0.038194in}{0.000000in}}%
\pgfpathcurveto{\pgfqpoint{0.038194in}{0.010129in}}{\pgfqpoint{0.034170in}{0.019845in}}{\pgfqpoint{0.027008in}{0.027008in}}%
\pgfpathcurveto{\pgfqpoint{0.019845in}{0.034170in}}{\pgfqpoint{0.010129in}{0.038194in}}{\pgfqpoint{0.000000in}{0.038194in}}%
\pgfpathcurveto{\pgfqpoint{-0.010129in}{0.038194in}}{\pgfqpoint{-0.019845in}{0.034170in}}{\pgfqpoint{-0.027008in}{0.027008in}}%
\pgfpathcurveto{\pgfqpoint{-0.034170in}{0.019845in}}{\pgfqpoint{-0.038194in}{0.010129in}}{\pgfqpoint{-0.038194in}{0.000000in}}%
\pgfpathcurveto{\pgfqpoint{-0.038194in}{-0.010129in}}{\pgfqpoint{-0.034170in}{-0.019845in}}{\pgfqpoint{-0.027008in}{-0.027008in}}%
\pgfpathcurveto{\pgfqpoint{-0.019845in}{-0.034170in}}{\pgfqpoint{-0.010129in}{-0.038194in}}{\pgfqpoint{0.000000in}{-0.038194in}}%
\pgfpathlineto{\pgfqpoint{0.000000in}{-0.038194in}}%
\pgfpathclose%
\pgfusepath{stroke,fill}%
}%
\begin{pgfscope}%
\pgfsys@transformshift{7.546000in}{2.265600in}%
\pgfsys@useobject{currentmarker}{}%
\end{pgfscope}%
\end{pgfscope}%
\begin{pgfscope}%
\pgfsetbuttcap%
\pgfsetroundjoin%
\definecolor{currentfill}{rgb}{0.549020,0.160784,0.505882}%
\pgfsetfillcolor{currentfill}%
\pgfsetlinewidth{1.003750pt}%
\definecolor{currentstroke}{rgb}{0.000000,0.000000,0.000000}%
\pgfsetstrokecolor{currentstroke}%
\pgfsetstrokeopacity{0.000000}%
\pgfsetdash{}{0pt}%
\pgfsys@defobject{currentmarker}{\pgfqpoint{-0.038194in}{-0.038194in}}{\pgfqpoint{0.038194in}{0.038194in}}{%
\pgfpathmoveto{\pgfqpoint{0.000000in}{-0.038194in}}%
\pgfpathcurveto{\pgfqpoint{0.010129in}{-0.038194in}}{\pgfqpoint{0.019845in}{-0.034170in}}{\pgfqpoint{0.027008in}{-0.027008in}}%
\pgfpathcurveto{\pgfqpoint{0.034170in}{-0.019845in}}{\pgfqpoint{0.038194in}{-0.010129in}}{\pgfqpoint{0.038194in}{0.000000in}}%
\pgfpathcurveto{\pgfqpoint{0.038194in}{0.010129in}}{\pgfqpoint{0.034170in}{0.019845in}}{\pgfqpoint{0.027008in}{0.027008in}}%
\pgfpathcurveto{\pgfqpoint{0.019845in}{0.034170in}}{\pgfqpoint{0.010129in}{0.038194in}}{\pgfqpoint{0.000000in}{0.038194in}}%
\pgfpathcurveto{\pgfqpoint{-0.010129in}{0.038194in}}{\pgfqpoint{-0.019845in}{0.034170in}}{\pgfqpoint{-0.027008in}{0.027008in}}%
\pgfpathcurveto{\pgfqpoint{-0.034170in}{0.019845in}}{\pgfqpoint{-0.038194in}{0.010129in}}{\pgfqpoint{-0.038194in}{0.000000in}}%
\pgfpathcurveto{\pgfqpoint{-0.038194in}{-0.010129in}}{\pgfqpoint{-0.034170in}{-0.019845in}}{\pgfqpoint{-0.027008in}{-0.027008in}}%
\pgfpathcurveto{\pgfqpoint{-0.019845in}{-0.034170in}}{\pgfqpoint{-0.010129in}{-0.038194in}}{\pgfqpoint{0.000000in}{-0.038194in}}%
\pgfpathlineto{\pgfqpoint{0.000000in}{-0.038194in}}%
\pgfpathclose%
\pgfusepath{stroke,fill}%
}%
\begin{pgfscope}%
\pgfsys@transformshift{7.546000in}{2.088960in}%
\pgfsys@useobject{currentmarker}{}%
\end{pgfscope}%
\end{pgfscope}%
\begin{pgfscope}%
\pgfsetbuttcap%
\pgfsetroundjoin%
\definecolor{currentfill}{rgb}{0.941176,0.584314,0.337255}%
\pgfsetfillcolor{currentfill}%
\pgfsetlinewidth{1.003750pt}%
\definecolor{currentstroke}{rgb}{0.000000,0.000000,0.000000}%
\pgfsetstrokecolor{currentstroke}%
\pgfsetstrokeopacity{0.000000}%
\pgfsetdash{}{0pt}%
\pgfsys@defobject{currentmarker}{\pgfqpoint{-0.038194in}{-0.038194in}}{\pgfqpoint{0.038194in}{0.038194in}}{%
\pgfpathmoveto{\pgfqpoint{0.000000in}{-0.038194in}}%
\pgfpathcurveto{\pgfqpoint{0.010129in}{-0.038194in}}{\pgfqpoint{0.019845in}{-0.034170in}}{\pgfqpoint{0.027008in}{-0.027008in}}%
\pgfpathcurveto{\pgfqpoint{0.034170in}{-0.019845in}}{\pgfqpoint{0.038194in}{-0.010129in}}{\pgfqpoint{0.038194in}{0.000000in}}%
\pgfpathcurveto{\pgfqpoint{0.038194in}{0.010129in}}{\pgfqpoint{0.034170in}{0.019845in}}{\pgfqpoint{0.027008in}{0.027008in}}%
\pgfpathcurveto{\pgfqpoint{0.019845in}{0.034170in}}{\pgfqpoint{0.010129in}{0.038194in}}{\pgfqpoint{0.000000in}{0.038194in}}%
\pgfpathcurveto{\pgfqpoint{-0.010129in}{0.038194in}}{\pgfqpoint{-0.019845in}{0.034170in}}{\pgfqpoint{-0.027008in}{0.027008in}}%
\pgfpathcurveto{\pgfqpoint{-0.034170in}{0.019845in}}{\pgfqpoint{-0.038194in}{0.010129in}}{\pgfqpoint{-0.038194in}{0.000000in}}%
\pgfpathcurveto{\pgfqpoint{-0.038194in}{-0.010129in}}{\pgfqpoint{-0.034170in}{-0.019845in}}{\pgfqpoint{-0.027008in}{-0.027008in}}%
\pgfpathcurveto{\pgfqpoint{-0.019845in}{-0.034170in}}{\pgfqpoint{-0.010129in}{-0.038194in}}{\pgfqpoint{0.000000in}{-0.038194in}}%
\pgfpathlineto{\pgfqpoint{0.000000in}{-0.038194in}}%
\pgfpathclose%
\pgfusepath{stroke,fill}%
}%
\begin{pgfscope}%
\pgfsys@transformshift{7.546000in}{1.912320in}%
\pgfsys@useobject{currentmarker}{}%
\end{pgfscope}%
\end{pgfscope}%
\begin{pgfscope}%
\pgfsetbuttcap%
\pgfsetroundjoin%
\definecolor{currentfill}{rgb}{0.000000,0.560784,0.611765}%
\pgfsetfillcolor{currentfill}%
\pgfsetlinewidth{1.003750pt}%
\definecolor{currentstroke}{rgb}{0.000000,0.000000,0.000000}%
\pgfsetstrokecolor{currentstroke}%
\pgfsetstrokeopacity{0.000000}%
\pgfsetdash{}{0pt}%
\pgfsys@defobject{currentmarker}{\pgfqpoint{-0.038194in}{-0.038194in}}{\pgfqpoint{0.038194in}{0.038194in}}{%
\pgfpathmoveto{\pgfqpoint{0.000000in}{-0.038194in}}%
\pgfpathcurveto{\pgfqpoint{0.010129in}{-0.038194in}}{\pgfqpoint{0.019845in}{-0.034170in}}{\pgfqpoint{0.027008in}{-0.027008in}}%
\pgfpathcurveto{\pgfqpoint{0.034170in}{-0.019845in}}{\pgfqpoint{0.038194in}{-0.010129in}}{\pgfqpoint{0.038194in}{0.000000in}}%
\pgfpathcurveto{\pgfqpoint{0.038194in}{0.010129in}}{\pgfqpoint{0.034170in}{0.019845in}}{\pgfqpoint{0.027008in}{0.027008in}}%
\pgfpathcurveto{\pgfqpoint{0.019845in}{0.034170in}}{\pgfqpoint{0.010129in}{0.038194in}}{\pgfqpoint{0.000000in}{0.038194in}}%
\pgfpathcurveto{\pgfqpoint{-0.010129in}{0.038194in}}{\pgfqpoint{-0.019845in}{0.034170in}}{\pgfqpoint{-0.027008in}{0.027008in}}%
\pgfpathcurveto{\pgfqpoint{-0.034170in}{0.019845in}}{\pgfqpoint{-0.038194in}{0.010129in}}{\pgfqpoint{-0.038194in}{0.000000in}}%
\pgfpathcurveto{\pgfqpoint{-0.038194in}{-0.010129in}}{\pgfqpoint{-0.034170in}{-0.019845in}}{\pgfqpoint{-0.027008in}{-0.027008in}}%
\pgfpathcurveto{\pgfqpoint{-0.019845in}{-0.034170in}}{\pgfqpoint{-0.010129in}{-0.038194in}}{\pgfqpoint{0.000000in}{-0.038194in}}%
\pgfpathlineto{\pgfqpoint{0.000000in}{-0.038194in}}%
\pgfpathclose%
\pgfusepath{stroke,fill}%
}%
\begin{pgfscope}%
\pgfsys@transformshift{7.546000in}{1.735680in}%
\pgfsys@useobject{currentmarker}{}%
\end{pgfscope}%
\end{pgfscope}%
\begin{pgfscope}%
\pgfsetbuttcap%
\pgfsetroundjoin%
\definecolor{currentfill}{rgb}{0.200000,0.133333,0.533333}%
\pgfsetfillcolor{currentfill}%
\pgfsetlinewidth{1.003750pt}%
\definecolor{currentstroke}{rgb}{0.000000,0.000000,0.000000}%
\pgfsetstrokecolor{currentstroke}%
\pgfsetstrokeopacity{0.000000}%
\pgfsetdash{}{0pt}%
\pgfsys@defobject{currentmarker}{\pgfqpoint{-0.038194in}{-0.038194in}}{\pgfqpoint{0.038194in}{0.038194in}}{%
\pgfpathmoveto{\pgfqpoint{0.000000in}{-0.038194in}}%
\pgfpathcurveto{\pgfqpoint{0.010129in}{-0.038194in}}{\pgfqpoint{0.019845in}{-0.034170in}}{\pgfqpoint{0.027008in}{-0.027008in}}%
\pgfpathcurveto{\pgfqpoint{0.034170in}{-0.019845in}}{\pgfqpoint{0.038194in}{-0.010129in}}{\pgfqpoint{0.038194in}{0.000000in}}%
\pgfpathcurveto{\pgfqpoint{0.038194in}{0.010129in}}{\pgfqpoint{0.034170in}{0.019845in}}{\pgfqpoint{0.027008in}{0.027008in}}%
\pgfpathcurveto{\pgfqpoint{0.019845in}{0.034170in}}{\pgfqpoint{0.010129in}{0.038194in}}{\pgfqpoint{0.000000in}{0.038194in}}%
\pgfpathcurveto{\pgfqpoint{-0.010129in}{0.038194in}}{\pgfqpoint{-0.019845in}{0.034170in}}{\pgfqpoint{-0.027008in}{0.027008in}}%
\pgfpathcurveto{\pgfqpoint{-0.034170in}{0.019845in}}{\pgfqpoint{-0.038194in}{0.010129in}}{\pgfqpoint{-0.038194in}{0.000000in}}%
\pgfpathcurveto{\pgfqpoint{-0.038194in}{-0.010129in}}{\pgfqpoint{-0.034170in}{-0.019845in}}{\pgfqpoint{-0.027008in}{-0.027008in}}%
\pgfpathcurveto{\pgfqpoint{-0.019845in}{-0.034170in}}{\pgfqpoint{-0.010129in}{-0.038194in}}{\pgfqpoint{0.000000in}{-0.038194in}}%
\pgfpathlineto{\pgfqpoint{0.000000in}{-0.038194in}}%
\pgfpathclose%
\pgfusepath{stroke,fill}%
}%
\begin{pgfscope}%
\pgfsys@transformshift{7.546000in}{1.559040in}%
\pgfsys@useobject{currentmarker}{}%
\end{pgfscope}%
\end{pgfscope}%
\begin{pgfscope}%
\pgfsetbuttcap%
\pgfsetroundjoin%
\definecolor{currentfill}{rgb}{0.901961,0.196078,0.156863}%
\pgfsetfillcolor{currentfill}%
\pgfsetlinewidth{1.003750pt}%
\definecolor{currentstroke}{rgb}{0.000000,0.000000,0.000000}%
\pgfsetstrokecolor{currentstroke}%
\pgfsetstrokeopacity{0.000000}%
\pgfsetdash{}{0pt}%
\pgfsys@defobject{currentmarker}{\pgfqpoint{-0.038194in}{-0.038194in}}{\pgfqpoint{0.038194in}{0.038194in}}{%
\pgfpathmoveto{\pgfqpoint{0.000000in}{-0.038194in}}%
\pgfpathcurveto{\pgfqpoint{0.010129in}{-0.038194in}}{\pgfqpoint{0.019845in}{-0.034170in}}{\pgfqpoint{0.027008in}{-0.027008in}}%
\pgfpathcurveto{\pgfqpoint{0.034170in}{-0.019845in}}{\pgfqpoint{0.038194in}{-0.010129in}}{\pgfqpoint{0.038194in}{0.000000in}}%
\pgfpathcurveto{\pgfqpoint{0.038194in}{0.010129in}}{\pgfqpoint{0.034170in}{0.019845in}}{\pgfqpoint{0.027008in}{0.027008in}}%
\pgfpathcurveto{\pgfqpoint{0.019845in}{0.034170in}}{\pgfqpoint{0.010129in}{0.038194in}}{\pgfqpoint{0.000000in}{0.038194in}}%
\pgfpathcurveto{\pgfqpoint{-0.010129in}{0.038194in}}{\pgfqpoint{-0.019845in}{0.034170in}}{\pgfqpoint{-0.027008in}{0.027008in}}%
\pgfpathcurveto{\pgfqpoint{-0.034170in}{0.019845in}}{\pgfqpoint{-0.038194in}{0.010129in}}{\pgfqpoint{-0.038194in}{0.000000in}}%
\pgfpathcurveto{\pgfqpoint{-0.038194in}{-0.010129in}}{\pgfqpoint{-0.034170in}{-0.019845in}}{\pgfqpoint{-0.027008in}{-0.027008in}}%
\pgfpathcurveto{\pgfqpoint{-0.019845in}{-0.034170in}}{\pgfqpoint{-0.010129in}{-0.038194in}}{\pgfqpoint{0.000000in}{-0.038194in}}%
\pgfpathlineto{\pgfqpoint{0.000000in}{-0.038194in}}%
\pgfpathclose%
\pgfusepath{stroke,fill}%
}%
\begin{pgfscope}%
\pgfsys@transformshift{10.780000in}{2.265600in}%
\pgfsys@useobject{currentmarker}{}%
\end{pgfscope}%
\end{pgfscope}%
\begin{pgfscope}%
\pgfsetbuttcap%
\pgfsetroundjoin%
\definecolor{currentfill}{rgb}{0.768627,0.603922,0.000000}%
\pgfsetfillcolor{currentfill}%
\pgfsetlinewidth{1.003750pt}%
\definecolor{currentstroke}{rgb}{0.000000,0.000000,0.000000}%
\pgfsetstrokecolor{currentstroke}%
\pgfsetstrokeopacity{0.000000}%
\pgfsetdash{}{0pt}%
\pgfsys@defobject{currentmarker}{\pgfqpoint{-0.038194in}{-0.038194in}}{\pgfqpoint{0.038194in}{0.038194in}}{%
\pgfpathmoveto{\pgfqpoint{0.000000in}{-0.038194in}}%
\pgfpathcurveto{\pgfqpoint{0.010129in}{-0.038194in}}{\pgfqpoint{0.019845in}{-0.034170in}}{\pgfqpoint{0.027008in}{-0.027008in}}%
\pgfpathcurveto{\pgfqpoint{0.034170in}{-0.019845in}}{\pgfqpoint{0.038194in}{-0.010129in}}{\pgfqpoint{0.038194in}{0.000000in}}%
\pgfpathcurveto{\pgfqpoint{0.038194in}{0.010129in}}{\pgfqpoint{0.034170in}{0.019845in}}{\pgfqpoint{0.027008in}{0.027008in}}%
\pgfpathcurveto{\pgfqpoint{0.019845in}{0.034170in}}{\pgfqpoint{0.010129in}{0.038194in}}{\pgfqpoint{0.000000in}{0.038194in}}%
\pgfpathcurveto{\pgfqpoint{-0.010129in}{0.038194in}}{\pgfqpoint{-0.019845in}{0.034170in}}{\pgfqpoint{-0.027008in}{0.027008in}}%
\pgfpathcurveto{\pgfqpoint{-0.034170in}{0.019845in}}{\pgfqpoint{-0.038194in}{0.010129in}}{\pgfqpoint{-0.038194in}{0.000000in}}%
\pgfpathcurveto{\pgfqpoint{-0.038194in}{-0.010129in}}{\pgfqpoint{-0.034170in}{-0.019845in}}{\pgfqpoint{-0.027008in}{-0.027008in}}%
\pgfpathcurveto{\pgfqpoint{-0.019845in}{-0.034170in}}{\pgfqpoint{-0.010129in}{-0.038194in}}{\pgfqpoint{0.000000in}{-0.038194in}}%
\pgfpathlineto{\pgfqpoint{0.000000in}{-0.038194in}}%
\pgfpathclose%
\pgfusepath{stroke,fill}%
}%
\begin{pgfscope}%
\pgfsys@transformshift{10.780000in}{2.088960in}%
\pgfsys@useobject{currentmarker}{}%
\end{pgfscope}%
\end{pgfscope}%
\begin{pgfscope}%
\pgfsetbuttcap%
\pgfsetroundjoin%
\definecolor{currentfill}{rgb}{0.498039,0.560784,0.678431}%
\pgfsetfillcolor{currentfill}%
\pgfsetlinewidth{1.003750pt}%
\definecolor{currentstroke}{rgb}{0.000000,0.000000,0.000000}%
\pgfsetstrokecolor{currentstroke}%
\pgfsetstrokeopacity{0.000000}%
\pgfsetdash{}{0pt}%
\pgfsys@defobject{currentmarker}{\pgfqpoint{-0.038194in}{-0.038194in}}{\pgfqpoint{0.038194in}{0.038194in}}{%
\pgfpathmoveto{\pgfqpoint{0.000000in}{-0.038194in}}%
\pgfpathcurveto{\pgfqpoint{0.010129in}{-0.038194in}}{\pgfqpoint{0.019845in}{-0.034170in}}{\pgfqpoint{0.027008in}{-0.027008in}}%
\pgfpathcurveto{\pgfqpoint{0.034170in}{-0.019845in}}{\pgfqpoint{0.038194in}{-0.010129in}}{\pgfqpoint{0.038194in}{0.000000in}}%
\pgfpathcurveto{\pgfqpoint{0.038194in}{0.010129in}}{\pgfqpoint{0.034170in}{0.019845in}}{\pgfqpoint{0.027008in}{0.027008in}}%
\pgfpathcurveto{\pgfqpoint{0.019845in}{0.034170in}}{\pgfqpoint{0.010129in}{0.038194in}}{\pgfqpoint{0.000000in}{0.038194in}}%
\pgfpathcurveto{\pgfqpoint{-0.010129in}{0.038194in}}{\pgfqpoint{-0.019845in}{0.034170in}}{\pgfqpoint{-0.027008in}{0.027008in}}%
\pgfpathcurveto{\pgfqpoint{-0.034170in}{0.019845in}}{\pgfqpoint{-0.038194in}{0.010129in}}{\pgfqpoint{-0.038194in}{0.000000in}}%
\pgfpathcurveto{\pgfqpoint{-0.038194in}{-0.010129in}}{\pgfqpoint{-0.034170in}{-0.019845in}}{\pgfqpoint{-0.027008in}{-0.027008in}}%
\pgfpathcurveto{\pgfqpoint{-0.019845in}{-0.034170in}}{\pgfqpoint{-0.010129in}{-0.038194in}}{\pgfqpoint{0.000000in}{-0.038194in}}%
\pgfpathlineto{\pgfqpoint{0.000000in}{-0.038194in}}%
\pgfpathclose%
\pgfusepath{stroke,fill}%
}%
\begin{pgfscope}%
\pgfsys@transformshift{10.780000in}{1.912320in}%
\pgfsys@useobject{currentmarker}{}%
\end{pgfscope}%
\end{pgfscope}%
\begin{pgfscope}%
\pgfsetbuttcap%
\pgfsetroundjoin%
\definecolor{currentfill}{rgb}{0.176471,0.713725,0.643137}%
\pgfsetfillcolor{currentfill}%
\pgfsetlinewidth{1.003750pt}%
\definecolor{currentstroke}{rgb}{0.000000,0.000000,0.000000}%
\pgfsetstrokecolor{currentstroke}%
\pgfsetstrokeopacity{0.000000}%
\pgfsetdash{}{0pt}%
\pgfsys@defobject{currentmarker}{\pgfqpoint{-0.038194in}{-0.038194in}}{\pgfqpoint{0.038194in}{0.038194in}}{%
\pgfpathmoveto{\pgfqpoint{0.000000in}{-0.038194in}}%
\pgfpathcurveto{\pgfqpoint{0.010129in}{-0.038194in}}{\pgfqpoint{0.019845in}{-0.034170in}}{\pgfqpoint{0.027008in}{-0.027008in}}%
\pgfpathcurveto{\pgfqpoint{0.034170in}{-0.019845in}}{\pgfqpoint{0.038194in}{-0.010129in}}{\pgfqpoint{0.038194in}{0.000000in}}%
\pgfpathcurveto{\pgfqpoint{0.038194in}{0.010129in}}{\pgfqpoint{0.034170in}{0.019845in}}{\pgfqpoint{0.027008in}{0.027008in}}%
\pgfpathcurveto{\pgfqpoint{0.019845in}{0.034170in}}{\pgfqpoint{0.010129in}{0.038194in}}{\pgfqpoint{0.000000in}{0.038194in}}%
\pgfpathcurveto{\pgfqpoint{-0.010129in}{0.038194in}}{\pgfqpoint{-0.019845in}{0.034170in}}{\pgfqpoint{-0.027008in}{0.027008in}}%
\pgfpathcurveto{\pgfqpoint{-0.034170in}{0.019845in}}{\pgfqpoint{-0.038194in}{0.010129in}}{\pgfqpoint{-0.038194in}{0.000000in}}%
\pgfpathcurveto{\pgfqpoint{-0.038194in}{-0.010129in}}{\pgfqpoint{-0.034170in}{-0.019845in}}{\pgfqpoint{-0.027008in}{-0.027008in}}%
\pgfpathcurveto{\pgfqpoint{-0.019845in}{-0.034170in}}{\pgfqpoint{-0.010129in}{-0.038194in}}{\pgfqpoint{0.000000in}{-0.038194in}}%
\pgfpathlineto{\pgfqpoint{0.000000in}{-0.038194in}}%
\pgfpathclose%
\pgfusepath{stroke,fill}%
}%
\begin{pgfscope}%
\pgfsys@transformshift{10.780000in}{1.735680in}%
\pgfsys@useobject{currentmarker}{}%
\end{pgfscope}%
\end{pgfscope}%
\begin{pgfscope}%
\pgfsetbuttcap%
\pgfsetroundjoin%
\definecolor{currentfill}{rgb}{0.321569,0.321569,0.321569}%
\pgfsetfillcolor{currentfill}%
\pgfsetlinewidth{1.003750pt}%
\definecolor{currentstroke}{rgb}{0.000000,0.000000,0.000000}%
\pgfsetstrokecolor{currentstroke}%
\pgfsetstrokeopacity{0.000000}%
\pgfsetdash{}{0pt}%
\pgfsys@defobject{currentmarker}{\pgfqpoint{-0.038194in}{-0.038194in}}{\pgfqpoint{0.038194in}{0.038194in}}{%
\pgfpathmoveto{\pgfqpoint{0.000000in}{-0.038194in}}%
\pgfpathcurveto{\pgfqpoint{0.010129in}{-0.038194in}}{\pgfqpoint{0.019845in}{-0.034170in}}{\pgfqpoint{0.027008in}{-0.027008in}}%
\pgfpathcurveto{\pgfqpoint{0.034170in}{-0.019845in}}{\pgfqpoint{0.038194in}{-0.010129in}}{\pgfqpoint{0.038194in}{0.000000in}}%
\pgfpathcurveto{\pgfqpoint{0.038194in}{0.010129in}}{\pgfqpoint{0.034170in}{0.019845in}}{\pgfqpoint{0.027008in}{0.027008in}}%
\pgfpathcurveto{\pgfqpoint{0.019845in}{0.034170in}}{\pgfqpoint{0.010129in}{0.038194in}}{\pgfqpoint{0.000000in}{0.038194in}}%
\pgfpathcurveto{\pgfqpoint{-0.010129in}{0.038194in}}{\pgfqpoint{-0.019845in}{0.034170in}}{\pgfqpoint{-0.027008in}{0.027008in}}%
\pgfpathcurveto{\pgfqpoint{-0.034170in}{0.019845in}}{\pgfqpoint{-0.038194in}{0.010129in}}{\pgfqpoint{-0.038194in}{0.000000in}}%
\pgfpathcurveto{\pgfqpoint{-0.038194in}{-0.010129in}}{\pgfqpoint{-0.034170in}{-0.019845in}}{\pgfqpoint{-0.027008in}{-0.027008in}}%
\pgfpathcurveto{\pgfqpoint{-0.019845in}{-0.034170in}}{\pgfqpoint{-0.010129in}{-0.038194in}}{\pgfqpoint{0.000000in}{-0.038194in}}%
\pgfpathlineto{\pgfqpoint{0.000000in}{-0.038194in}}%
\pgfpathclose%
\pgfusepath{stroke,fill}%
}%
\begin{pgfscope}%
\pgfsys@transformshift{10.780000in}{1.559040in}%
\pgfsys@useobject{currentmarker}{}%
\end{pgfscope}%
\end{pgfscope}%
\begin{pgfscope}%
\pgfsetbuttcap%
\pgfsetroundjoin%
\definecolor{currentfill}{rgb}{0.309804,0.372549,0.419608}%
\pgfsetfillcolor{currentfill}%
\pgfsetlinewidth{0.000000pt}%
\definecolor{currentstroke}{rgb}{0.309804,0.372549,0.419608}%
\pgfsetstrokecolor{currentstroke}%
\pgfsetdash{}{0pt}%
\pgfpathmoveto{\pgfqpoint{0.953314in}{9.372676in}}%
\pgfpathlineto{\pgfqpoint{1.005903in}{9.372676in}}%
\pgfpathlineto{\pgfqpoint{1.005903in}{9.360000in}}%
\pgfpathlineto{\pgfqpoint{0.935196in}{9.360000in}}%
\pgfpathlineto{\pgfqpoint{0.935196in}{9.372676in}}%
\pgfpathquadraticcurveto{\pgfqpoint{0.943766in}{9.381556in}}{\pgfqpoint{0.958566in}{9.396500in}}%
\pgfpathquadraticcurveto{\pgfqpoint{0.973390in}{9.411467in}}{\pgfqpoint{0.977186in}{9.415812in}}%
\pgfpathquadraticcurveto{\pgfqpoint{0.984419in}{9.423928in}}{\pgfqpoint{0.987283in}{9.429562in}}%
\pgfpathquadraticcurveto{\pgfqpoint{0.990172in}{9.435195in}}{\pgfqpoint{0.990172in}{9.440638in}}%
\pgfpathquadraticcurveto{\pgfqpoint{0.990172in}{9.449518in}}{\pgfqpoint{0.983941in}{9.455104in}}%
\pgfpathquadraticcurveto{\pgfqpoint{0.977711in}{9.460714in}}{\pgfqpoint{0.967709in}{9.460714in}}%
\pgfpathquadraticcurveto{\pgfqpoint{0.960619in}{9.460714in}}{\pgfqpoint{0.952741in}{9.458255in}}%
\pgfpathquadraticcurveto{\pgfqpoint{0.944888in}{9.455796in}}{\pgfqpoint{0.935936in}{9.450783in}}%
\pgfpathlineto{\pgfqpoint{0.935936in}{9.466013in}}%
\pgfpathquadraticcurveto{\pgfqpoint{0.945031in}{9.469666in}}{\pgfqpoint{0.952932in}{9.471528in}}%
\pgfpathquadraticcurveto{\pgfqpoint{0.960858in}{9.473390in}}{\pgfqpoint{0.967422in}{9.473390in}}%
\pgfpathquadraticcurveto{\pgfqpoint{0.984729in}{9.473390in}}{\pgfqpoint{0.995018in}{9.464724in}}%
\pgfpathquadraticcurveto{\pgfqpoint{1.005306in}{9.456083in}}{\pgfqpoint{1.005306in}{9.441617in}}%
\pgfpathquadraticcurveto{\pgfqpoint{1.005306in}{9.434742in}}{\pgfqpoint{1.002728in}{9.428583in}}%
\pgfpathquadraticcurveto{\pgfqpoint{1.000174in}{9.422448in}}{\pgfqpoint{0.993371in}{9.414093in}}%
\pgfpathquadraticcurveto{\pgfqpoint{0.991509in}{9.411921in}}{\pgfqpoint{0.981507in}{9.401584in}}%
\pgfpathquadraticcurveto{\pgfqpoint{0.971528in}{9.391248in}}{\pgfqpoint{0.953314in}{9.372676in}}%
\pgfpathlineto{\pgfqpoint{0.953314in}{9.372676in}}%
\pgfpathclose%
\pgfpathmoveto{\pgfqpoint{1.050519in}{9.372676in}}%
\pgfpathlineto{\pgfqpoint{1.103108in}{9.372676in}}%
\pgfpathlineto{\pgfqpoint{1.103108in}{9.360000in}}%
\pgfpathlineto{\pgfqpoint{1.032401in}{9.360000in}}%
\pgfpathlineto{\pgfqpoint{1.032401in}{9.372676in}}%
\pgfpathquadraticcurveto{\pgfqpoint{1.040970in}{9.381556in}}{\pgfqpoint{1.055771in}{9.396500in}}%
\pgfpathquadraticcurveto{\pgfqpoint{1.070595in}{9.411467in}}{\pgfqpoint{1.074391in}{9.415812in}}%
\pgfpathquadraticcurveto{\pgfqpoint{1.081624in}{9.423928in}}{\pgfqpoint{1.084488in}{9.429562in}}%
\pgfpathquadraticcurveto{\pgfqpoint{1.087377in}{9.435195in}}{\pgfqpoint{1.087377in}{9.440638in}}%
\pgfpathquadraticcurveto{\pgfqpoint{1.087377in}{9.449518in}}{\pgfqpoint{1.081146in}{9.455104in}}%
\pgfpathquadraticcurveto{\pgfqpoint{1.074916in}{9.460714in}}{\pgfqpoint{1.064914in}{9.460714in}}%
\pgfpathquadraticcurveto{\pgfqpoint{1.057824in}{9.460714in}}{\pgfqpoint{1.049946in}{9.458255in}}%
\pgfpathquadraticcurveto{\pgfqpoint{1.042092in}{9.455796in}}{\pgfqpoint{1.033141in}{9.450783in}}%
\pgfpathlineto{\pgfqpoint{1.033141in}{9.466013in}}%
\pgfpathquadraticcurveto{\pgfqpoint{1.042236in}{9.469666in}}{\pgfqpoint{1.050137in}{9.471528in}}%
\pgfpathquadraticcurveto{\pgfqpoint{1.058063in}{9.473390in}}{\pgfqpoint{1.064627in}{9.473390in}}%
\pgfpathquadraticcurveto{\pgfqpoint{1.081934in}{9.473390in}}{\pgfqpoint{1.092223in}{9.464724in}}%
\pgfpathquadraticcurveto{\pgfqpoint{1.102511in}{9.456083in}}{\pgfqpoint{1.102511in}{9.441617in}}%
\pgfpathquadraticcurveto{\pgfqpoint{1.102511in}{9.434742in}}{\pgfqpoint{1.099933in}{9.428583in}}%
\pgfpathquadraticcurveto{\pgfqpoint{1.097379in}{9.422448in}}{\pgfqpoint{1.090576in}{9.414093in}}%
\pgfpathquadraticcurveto{\pgfqpoint{1.088714in}{9.411921in}}{\pgfqpoint{1.078711in}{9.401584in}}%
\pgfpathquadraticcurveto{\pgfqpoint{1.068733in}{9.391248in}}{\pgfqpoint{1.050519in}{9.372676in}}%
\pgfpathlineto{\pgfqpoint{1.050519in}{9.372676in}}%
\pgfpathclose%
\pgfpathmoveto{\pgfqpoint{1.234616in}{9.441092in}}%
\pgfpathlineto{\pgfqpoint{1.234616in}{9.428105in}}%
\pgfpathquadraticcurveto{\pgfqpoint{1.228816in}{9.431089in}}{\pgfqpoint{1.222537in}{9.432569in}}%
\pgfpathquadraticcurveto{\pgfqpoint{1.216283in}{9.434073in}}{\pgfqpoint{1.209551in}{9.434073in}}%
\pgfpathquadraticcurveto{\pgfqpoint{1.199334in}{9.434073in}}{\pgfqpoint{1.194226in}{9.430946in}}%
\pgfpathquadraticcurveto{\pgfqpoint{1.189117in}{9.427819in}}{\pgfqpoint{1.189117in}{9.421541in}}%
\pgfpathquadraticcurveto{\pgfqpoint{1.189117in}{9.416766in}}{\pgfqpoint{1.192770in}{9.414045in}}%
\pgfpathquadraticcurveto{\pgfqpoint{1.196422in}{9.411324in}}{\pgfqpoint{1.207474in}{9.408865in}}%
\pgfpathlineto{\pgfqpoint{1.212177in}{9.407815in}}%
\pgfpathquadraticcurveto{\pgfqpoint{1.226786in}{9.404687in}}{\pgfqpoint{1.232945in}{9.398982in}}%
\pgfpathquadraticcurveto{\pgfqpoint{1.239104in}{9.393277in}}{\pgfqpoint{1.239104in}{9.383060in}}%
\pgfpathquadraticcurveto{\pgfqpoint{1.239104in}{9.371411in}}{\pgfqpoint{1.229890in}{9.364607in}}%
\pgfpathquadraticcurveto{\pgfqpoint{1.220675in}{9.357828in}}{\pgfqpoint{1.204562in}{9.357828in}}%
\pgfpathquadraticcurveto{\pgfqpoint{1.197854in}{9.357828in}}{\pgfqpoint{1.190573in}{9.359141in}}%
\pgfpathquadraticcurveto{\pgfqpoint{1.183293in}{9.360454in}}{\pgfqpoint{1.175248in}{9.363056in}}%
\pgfpathlineto{\pgfqpoint{1.175248in}{9.377235in}}%
\pgfpathquadraticcurveto{\pgfqpoint{1.182863in}{9.373273in}}{\pgfqpoint{1.190239in}{9.371291in}}%
\pgfpathquadraticcurveto{\pgfqpoint{1.197615in}{9.369334in}}{\pgfqpoint{1.204872in}{9.369334in}}%
\pgfpathquadraticcurveto{\pgfqpoint{1.214564in}{9.369334in}}{\pgfqpoint{1.219768in}{9.372652in}}%
\pgfpathquadraticcurveto{\pgfqpoint{1.224996in}{9.375970in}}{\pgfqpoint{1.224996in}{9.382010in}}%
\pgfpathquadraticcurveto{\pgfqpoint{1.224996in}{9.387595in}}{\pgfqpoint{1.221224in}{9.390579in}}%
\pgfpathquadraticcurveto{\pgfqpoint{1.217477in}{9.393563in}}{\pgfqpoint{1.204705in}{9.396332in}}%
\pgfpathlineto{\pgfqpoint{1.199931in}{9.397454in}}%
\pgfpathquadraticcurveto{\pgfqpoint{1.187184in}{9.400128in}}{\pgfqpoint{1.181502in}{9.405690in}}%
\pgfpathquadraticcurveto{\pgfqpoint{1.175845in}{9.411252in}}{\pgfqpoint{1.175845in}{9.420944in}}%
\pgfpathquadraticcurveto{\pgfqpoint{1.175845in}{9.432737in}}{\pgfqpoint{1.184200in}{9.439134in}}%
\pgfpathquadraticcurveto{\pgfqpoint{1.192555in}{9.445556in}}{\pgfqpoint{1.207928in}{9.445556in}}%
\pgfpathquadraticcurveto{\pgfqpoint{1.215519in}{9.445556in}}{\pgfqpoint{1.222227in}{9.444434in}}%
\pgfpathquadraticcurveto{\pgfqpoint{1.228959in}{9.443336in}}{\pgfqpoint{1.234616in}{9.441092in}}%
\pgfpathlineto{\pgfqpoint{1.234616in}{9.441092in}}%
\pgfpathclose%
\pgfpathmoveto{\pgfqpoint{1.274530in}{9.467279in}}%
\pgfpathlineto{\pgfqpoint{1.274530in}{9.443550in}}%
\pgfpathlineto{\pgfqpoint{1.302793in}{9.443550in}}%
\pgfpathlineto{\pgfqpoint{1.302793in}{9.432880in}}%
\pgfpathlineto{\pgfqpoint{1.274530in}{9.432880in}}%
\pgfpathlineto{\pgfqpoint{1.274530in}{9.387524in}}%
\pgfpathquadraticcurveto{\pgfqpoint{1.274530in}{9.377307in}}{\pgfqpoint{1.277322in}{9.374395in}}%
\pgfpathquadraticcurveto{\pgfqpoint{1.280115in}{9.371482in}}{\pgfqpoint{1.288709in}{9.371482in}}%
\pgfpathlineto{\pgfqpoint{1.302793in}{9.371482in}}%
\pgfpathlineto{\pgfqpoint{1.302793in}{9.360000in}}%
\pgfpathlineto{\pgfqpoint{1.288709in}{9.360000in}}%
\pgfpathquadraticcurveto{\pgfqpoint{1.272811in}{9.360000in}}{\pgfqpoint{1.266771in}{9.365920in}}%
\pgfpathquadraticcurveto{\pgfqpoint{1.260732in}{9.371864in}}{\pgfqpoint{1.260732in}{9.387524in}}%
\pgfpathlineto{\pgfqpoint{1.260732in}{9.432880in}}%
\pgfpathlineto{\pgfqpoint{1.250658in}{9.432880in}}%
\pgfpathlineto{\pgfqpoint{1.250658in}{9.443550in}}%
\pgfpathlineto{\pgfqpoint{1.260732in}{9.443550in}}%
\pgfpathlineto{\pgfqpoint{1.260732in}{9.467279in}}%
\pgfpathlineto{\pgfqpoint{1.274530in}{9.467279in}}%
\pgfpathlineto{\pgfqpoint{1.274530in}{9.467279in}}%
\pgfpathclose%
\pgfpathmoveto{\pgfqpoint{1.319432in}{9.392967in}}%
\pgfpathlineto{\pgfqpoint{1.319432in}{9.443550in}}%
\pgfpathlineto{\pgfqpoint{1.333158in}{9.443550in}}%
\pgfpathlineto{\pgfqpoint{1.333158in}{9.393492in}}%
\pgfpathquadraticcurveto{\pgfqpoint{1.333158in}{9.381628in}}{\pgfqpoint{1.337765in}{9.375684in}}%
\pgfpathquadraticcurveto{\pgfqpoint{1.342396in}{9.369763in}}{\pgfqpoint{1.351658in}{9.369763in}}%
\pgfpathquadraticcurveto{\pgfqpoint{1.362759in}{9.369763in}}{\pgfqpoint{1.369204in}{9.376853in}}%
\pgfpathquadraticcurveto{\pgfqpoint{1.375673in}{9.383943in}}{\pgfqpoint{1.375673in}{9.396189in}}%
\pgfpathlineto{\pgfqpoint{1.375673in}{9.443550in}}%
\pgfpathlineto{\pgfqpoint{1.389399in}{9.443550in}}%
\pgfpathlineto{\pgfqpoint{1.389399in}{9.360000in}}%
\pgfpathlineto{\pgfqpoint{1.375673in}{9.360000in}}%
\pgfpathlineto{\pgfqpoint{1.375673in}{9.372843in}}%
\pgfpathquadraticcurveto{\pgfqpoint{1.370684in}{9.365228in}}{\pgfqpoint{1.364072in}{9.361528in}}%
\pgfpathquadraticcurveto{\pgfqpoint{1.357483in}{9.357828in}}{\pgfqpoint{1.348746in}{9.357828in}}%
\pgfpathquadraticcurveto{\pgfqpoint{1.334352in}{9.357828in}}{\pgfqpoint{1.326880in}{9.366780in}}%
\pgfpathquadraticcurveto{\pgfqpoint{1.319432in}{9.375731in}}{\pgfqpoint{1.319432in}{9.392967in}}%
\pgfpathlineto{\pgfqpoint{1.319432in}{9.392967in}}%
\pgfpathclose%
\pgfpathmoveto{\pgfqpoint{1.353974in}{9.445556in}}%
\pgfpathlineto{\pgfqpoint{1.353974in}{9.445556in}}%
\pgfpathlineto{\pgfqpoint{1.353974in}{9.445556in}}%
\pgfpathclose%
\pgfpathmoveto{\pgfqpoint{1.472639in}{9.430875in}}%
\pgfpathlineto{\pgfqpoint{1.472639in}{9.476087in}}%
\pgfpathlineto{\pgfqpoint{1.486365in}{9.476087in}}%
\pgfpathlineto{\pgfqpoint{1.486365in}{9.360000in}}%
\pgfpathlineto{\pgfqpoint{1.472639in}{9.360000in}}%
\pgfpathlineto{\pgfqpoint{1.472639in}{9.372533in}}%
\pgfpathquadraticcurveto{\pgfqpoint{1.468319in}{9.365085in}}{\pgfqpoint{1.461706in}{9.361456in}}%
\pgfpathquadraticcurveto{\pgfqpoint{1.455118in}{9.357828in}}{\pgfqpoint{1.445855in}{9.357828in}}%
\pgfpathquadraticcurveto{\pgfqpoint{1.430721in}{9.357828in}}{\pgfqpoint{1.421196in}{9.369907in}}%
\pgfpathquadraticcurveto{\pgfqpoint{1.411695in}{9.382010in}}{\pgfqpoint{1.411695in}{9.401704in}}%
\pgfpathquadraticcurveto{\pgfqpoint{1.411695in}{9.421398in}}{\pgfqpoint{1.421196in}{9.433477in}}%
\pgfpathquadraticcurveto{\pgfqpoint{1.430721in}{9.445556in}}{\pgfqpoint{1.445855in}{9.445556in}}%
\pgfpathquadraticcurveto{\pgfqpoint{1.455118in}{9.445556in}}{\pgfqpoint{1.461706in}{9.441927in}}%
\pgfpathquadraticcurveto{\pgfqpoint{1.468319in}{9.438322in}}{\pgfqpoint{1.472639in}{9.430875in}}%
\pgfpathlineto{\pgfqpoint{1.472639in}{9.430875in}}%
\pgfpathclose%
\pgfpathmoveto{\pgfqpoint{1.425875in}{9.401704in}}%
\pgfpathquadraticcurveto{\pgfqpoint{1.425875in}{9.386569in}}{\pgfqpoint{1.432105in}{9.377951in}}%
\pgfpathquadraticcurveto{\pgfqpoint{1.438336in}{9.369334in}}{\pgfqpoint{1.449221in}{9.369334in}}%
\pgfpathquadraticcurveto{\pgfqpoint{1.460107in}{9.369334in}}{\pgfqpoint{1.466361in}{9.377951in}}%
\pgfpathquadraticcurveto{\pgfqpoint{1.472639in}{9.386569in}}{\pgfqpoint{1.472639in}{9.401704in}}%
\pgfpathquadraticcurveto{\pgfqpoint{1.472639in}{9.416838in}}{\pgfqpoint{1.466361in}{9.425456in}}%
\pgfpathquadraticcurveto{\pgfqpoint{1.460107in}{9.434073in}}{\pgfqpoint{1.449221in}{9.434073in}}%
\pgfpathquadraticcurveto{\pgfqpoint{1.438336in}{9.434073in}}{\pgfqpoint{1.432105in}{9.425456in}}%
\pgfpathquadraticcurveto{\pgfqpoint{1.425875in}{9.416838in}}{\pgfqpoint{1.425875in}{9.401704in}}%
\pgfpathlineto{\pgfqpoint{1.425875in}{9.401704in}}%
\pgfpathclose%
\pgfpathmoveto{\pgfqpoint{1.514653in}{9.443550in}}%
\pgfpathlineto{\pgfqpoint{1.528379in}{9.443550in}}%
\pgfpathlineto{\pgfqpoint{1.528379in}{9.360000in}}%
\pgfpathlineto{\pgfqpoint{1.514653in}{9.360000in}}%
\pgfpathlineto{\pgfqpoint{1.514653in}{9.443550in}}%
\pgfpathlineto{\pgfqpoint{1.514653in}{9.443550in}}%
\pgfpathclose%
\pgfpathmoveto{\pgfqpoint{1.514653in}{9.476087in}}%
\pgfpathlineto{\pgfqpoint{1.528379in}{9.476087in}}%
\pgfpathlineto{\pgfqpoint{1.528379in}{9.458685in}}%
\pgfpathlineto{\pgfqpoint{1.514653in}{9.458685in}}%
\pgfpathlineto{\pgfqpoint{1.514653in}{9.476087in}}%
\pgfpathlineto{\pgfqpoint{1.514653in}{9.476087in}}%
\pgfpathclose%
\pgfpathmoveto{\pgfqpoint{1.628568in}{9.405213in}}%
\pgfpathlineto{\pgfqpoint{1.628568in}{9.398505in}}%
\pgfpathlineto{\pgfqpoint{1.565452in}{9.398505in}}%
\pgfpathquadraticcurveto{\pgfqpoint{1.566359in}{9.384325in}}{\pgfqpoint{1.573998in}{9.376901in}}%
\pgfpathquadraticcurveto{\pgfqpoint{1.581637in}{9.369477in}}{\pgfqpoint{1.595291in}{9.369477in}}%
\pgfpathquadraticcurveto{\pgfqpoint{1.603193in}{9.369477in}}{\pgfqpoint{1.610617in}{9.371411in}}%
\pgfpathquadraticcurveto{\pgfqpoint{1.618041in}{9.373344in}}{\pgfqpoint{1.625369in}{9.377235in}}%
\pgfpathlineto{\pgfqpoint{1.625369in}{9.364249in}}%
\pgfpathquadraticcurveto{\pgfqpoint{1.617969in}{9.361122in}}{\pgfqpoint{1.610211in}{9.359475in}}%
\pgfpathquadraticcurveto{\pgfqpoint{1.602453in}{9.357828in}}{\pgfqpoint{1.594480in}{9.357828in}}%
\pgfpathquadraticcurveto{\pgfqpoint{1.574475in}{9.357828in}}{\pgfqpoint{1.562802in}{9.369453in}}%
\pgfpathquadraticcurveto{\pgfqpoint{1.551129in}{9.381102in}}{\pgfqpoint{1.551129in}{9.400964in}}%
\pgfpathquadraticcurveto{\pgfqpoint{1.551129in}{9.421469in}}{\pgfqpoint{1.562205in}{9.433500in}}%
\pgfpathquadraticcurveto{\pgfqpoint{1.573282in}{9.445556in}}{\pgfqpoint{1.592092in}{9.445556in}}%
\pgfpathquadraticcurveto{\pgfqpoint{1.608946in}{9.445556in}}{\pgfqpoint{1.618757in}{9.434694in}}%
\pgfpathquadraticcurveto{\pgfqpoint{1.628568in}{9.423856in}}{\pgfqpoint{1.628568in}{9.405213in}}%
\pgfpathlineto{\pgfqpoint{1.628568in}{9.405213in}}%
\pgfpathclose%
\pgfpathmoveto{\pgfqpoint{1.614842in}{9.409247in}}%
\pgfpathquadraticcurveto{\pgfqpoint{1.614699in}{9.420490in}}{\pgfqpoint{1.608540in}{9.427198in}}%
\pgfpathquadraticcurveto{\pgfqpoint{1.602381in}{9.433930in}}{\pgfqpoint{1.592236in}{9.433930in}}%
\pgfpathquadraticcurveto{\pgfqpoint{1.580753in}{9.433930in}}{\pgfqpoint{1.573855in}{9.427437in}}%
\pgfpathquadraticcurveto{\pgfqpoint{1.566956in}{9.420944in}}{\pgfqpoint{1.565905in}{9.409151in}}%
\pgfpathlineto{\pgfqpoint{1.614842in}{9.409247in}}%
\pgfpathlineto{\pgfqpoint{1.614842in}{9.409247in}}%
\pgfpathclose%
\pgfpathmoveto{\pgfqpoint{1.704360in}{9.441092in}}%
\pgfpathlineto{\pgfqpoint{1.704360in}{9.428105in}}%
\pgfpathquadraticcurveto{\pgfqpoint{1.698559in}{9.431089in}}{\pgfqpoint{1.692281in}{9.432569in}}%
\pgfpathquadraticcurveto{\pgfqpoint{1.686027in}{9.434073in}}{\pgfqpoint{1.679295in}{9.434073in}}%
\pgfpathquadraticcurveto{\pgfqpoint{1.669078in}{9.434073in}}{\pgfqpoint{1.663970in}{9.430946in}}%
\pgfpathquadraticcurveto{\pgfqpoint{1.658861in}{9.427819in}}{\pgfqpoint{1.658861in}{9.421541in}}%
\pgfpathquadraticcurveto{\pgfqpoint{1.658861in}{9.416766in}}{\pgfqpoint{1.662513in}{9.414045in}}%
\pgfpathquadraticcurveto{\pgfqpoint{1.666166in}{9.411324in}}{\pgfqpoint{1.677218in}{9.408865in}}%
\pgfpathlineto{\pgfqpoint{1.681921in}{9.407815in}}%
\pgfpathquadraticcurveto{\pgfqpoint{1.696530in}{9.404687in}}{\pgfqpoint{1.702689in}{9.398982in}}%
\pgfpathquadraticcurveto{\pgfqpoint{1.708848in}{9.393277in}}{\pgfqpoint{1.708848in}{9.383060in}}%
\pgfpathquadraticcurveto{\pgfqpoint{1.708848in}{9.371411in}}{\pgfqpoint{1.699634in}{9.364607in}}%
\pgfpathquadraticcurveto{\pgfqpoint{1.690419in}{9.357828in}}{\pgfqpoint{1.674306in}{9.357828in}}%
\pgfpathquadraticcurveto{\pgfqpoint{1.667598in}{9.357828in}}{\pgfqpoint{1.660317in}{9.359141in}}%
\pgfpathquadraticcurveto{\pgfqpoint{1.653036in}{9.360454in}}{\pgfqpoint{1.644992in}{9.363056in}}%
\pgfpathlineto{\pgfqpoint{1.644992in}{9.377235in}}%
\pgfpathquadraticcurveto{\pgfqpoint{1.652607in}{9.373273in}}{\pgfqpoint{1.659983in}{9.371291in}}%
\pgfpathquadraticcurveto{\pgfqpoint{1.667359in}{9.369334in}}{\pgfqpoint{1.674616in}{9.369334in}}%
\pgfpathquadraticcurveto{\pgfqpoint{1.684308in}{9.369334in}}{\pgfqpoint{1.689512in}{9.372652in}}%
\pgfpathquadraticcurveto{\pgfqpoint{1.694740in}{9.375970in}}{\pgfqpoint{1.694740in}{9.382010in}}%
\pgfpathquadraticcurveto{\pgfqpoint{1.694740in}{9.387595in}}{\pgfqpoint{1.690968in}{9.390579in}}%
\pgfpathquadraticcurveto{\pgfqpoint{1.687220in}{9.393563in}}{\pgfqpoint{1.674449in}{9.396332in}}%
\pgfpathlineto{\pgfqpoint{1.669675in}{9.397454in}}%
\pgfpathquadraticcurveto{\pgfqpoint{1.656928in}{9.400128in}}{\pgfqpoint{1.651246in}{9.405690in}}%
\pgfpathquadraticcurveto{\pgfqpoint{1.645589in}{9.411252in}}{\pgfqpoint{1.645589in}{9.420944in}}%
\pgfpathquadraticcurveto{\pgfqpoint{1.645589in}{9.432737in}}{\pgfqpoint{1.653944in}{9.439134in}}%
\pgfpathquadraticcurveto{\pgfqpoint{1.662299in}{9.445556in}}{\pgfqpoint{1.677672in}{9.445556in}}%
\pgfpathquadraticcurveto{\pgfqpoint{1.685263in}{9.445556in}}{\pgfqpoint{1.691971in}{9.444434in}}%
\pgfpathquadraticcurveto{\pgfqpoint{1.698703in}{9.443336in}}{\pgfqpoint{1.704360in}{9.441092in}}%
\pgfpathlineto{\pgfqpoint{1.704360in}{9.441092in}}%
\pgfpathclose%
\pgfpathmoveto{\pgfqpoint{1.781179in}{9.422520in}}%
\pgfpathlineto{\pgfqpoint{1.796934in}{9.422520in}}%
\pgfpathlineto{\pgfqpoint{1.796934in}{9.403566in}}%
\pgfpathlineto{\pgfqpoint{1.781179in}{9.403566in}}%
\pgfpathlineto{\pgfqpoint{1.781179in}{9.422520in}}%
\pgfpathlineto{\pgfqpoint{1.781179in}{9.422520in}}%
\pgfpathclose%
\pgfpathmoveto{\pgfqpoint{1.876951in}{9.471385in}}%
\pgfpathlineto{\pgfqpoint{1.940951in}{9.471385in}}%
\pgfpathlineto{\pgfqpoint{1.940951in}{9.458685in}}%
\pgfpathlineto{\pgfqpoint{1.892014in}{9.458685in}}%
\pgfpathlineto{\pgfqpoint{1.892014in}{9.425862in}}%
\pgfpathlineto{\pgfqpoint{1.936177in}{9.425862in}}%
\pgfpathlineto{\pgfqpoint{1.936177in}{9.413186in}}%
\pgfpathlineto{\pgfqpoint{1.892014in}{9.413186in}}%
\pgfpathlineto{\pgfqpoint{1.892014in}{9.360000in}}%
\pgfpathlineto{\pgfqpoint{1.876951in}{9.360000in}}%
\pgfpathlineto{\pgfqpoint{1.876951in}{9.471385in}}%
\pgfpathlineto{\pgfqpoint{1.876951in}{9.471385in}}%
\pgfpathclose%
\pgfpathmoveto{\pgfqpoint{1.953102in}{9.443550in}}%
\pgfpathlineto{\pgfqpoint{1.966828in}{9.443550in}}%
\pgfpathlineto{\pgfqpoint{1.966828in}{9.360000in}}%
\pgfpathlineto{\pgfqpoint{1.953102in}{9.360000in}}%
\pgfpathlineto{\pgfqpoint{1.953102in}{9.443550in}}%
\pgfpathlineto{\pgfqpoint{1.953102in}{9.443550in}}%
\pgfpathclose%
\pgfpathmoveto{\pgfqpoint{1.953102in}{9.476087in}}%
\pgfpathlineto{\pgfqpoint{1.966828in}{9.476087in}}%
\pgfpathlineto{\pgfqpoint{1.966828in}{9.458685in}}%
\pgfpathlineto{\pgfqpoint{1.953102in}{9.458685in}}%
\pgfpathlineto{\pgfqpoint{1.953102in}{9.476087in}}%
\pgfpathlineto{\pgfqpoint{1.953102in}{9.476087in}}%
\pgfpathclose%
\pgfpathmoveto{\pgfqpoint{1.995545in}{9.476087in}}%
\pgfpathlineto{\pgfqpoint{2.009271in}{9.476087in}}%
\pgfpathlineto{\pgfqpoint{2.009271in}{9.360000in}}%
\pgfpathlineto{\pgfqpoint{1.995545in}{9.360000in}}%
\pgfpathlineto{\pgfqpoint{1.995545in}{9.476087in}}%
\pgfpathlineto{\pgfqpoint{1.995545in}{9.476087in}}%
\pgfpathclose%
\pgfpathmoveto{\pgfqpoint{2.037989in}{9.476087in}}%
\pgfpathlineto{\pgfqpoint{2.051715in}{9.476087in}}%
\pgfpathlineto{\pgfqpoint{2.051715in}{9.360000in}}%
\pgfpathlineto{\pgfqpoint{2.037989in}{9.360000in}}%
\pgfpathlineto{\pgfqpoint{2.037989in}{9.476087in}}%
\pgfpathlineto{\pgfqpoint{2.037989in}{9.476087in}}%
\pgfpathclose%
\pgfpathmoveto{\pgfqpoint{2.151904in}{9.405213in}}%
\pgfpathlineto{\pgfqpoint{2.151904in}{9.398505in}}%
\pgfpathlineto{\pgfqpoint{2.088787in}{9.398505in}}%
\pgfpathquadraticcurveto{\pgfqpoint{2.089694in}{9.384325in}}{\pgfqpoint{2.097333in}{9.376901in}}%
\pgfpathquadraticcurveto{\pgfqpoint{2.104972in}{9.369477in}}{\pgfqpoint{2.118627in}{9.369477in}}%
\pgfpathquadraticcurveto{\pgfqpoint{2.126528in}{9.369477in}}{\pgfqpoint{2.133952in}{9.371411in}}%
\pgfpathquadraticcurveto{\pgfqpoint{2.141376in}{9.373344in}}{\pgfqpoint{2.148705in}{9.377235in}}%
\pgfpathlineto{\pgfqpoint{2.148705in}{9.364249in}}%
\pgfpathquadraticcurveto{\pgfqpoint{2.141305in}{9.361122in}}{\pgfqpoint{2.133546in}{9.359475in}}%
\pgfpathquadraticcurveto{\pgfqpoint{2.125788in}{9.357828in}}{\pgfqpoint{2.117815in}{9.357828in}}%
\pgfpathquadraticcurveto{\pgfqpoint{2.097811in}{9.357828in}}{\pgfqpoint{2.086138in}{9.369453in}}%
\pgfpathquadraticcurveto{\pgfqpoint{2.074464in}{9.381102in}}{\pgfqpoint{2.074464in}{9.400964in}}%
\pgfpathquadraticcurveto{\pgfqpoint{2.074464in}{9.421469in}}{\pgfqpoint{2.085541in}{9.433500in}}%
\pgfpathquadraticcurveto{\pgfqpoint{2.096617in}{9.445556in}}{\pgfqpoint{2.115428in}{9.445556in}}%
\pgfpathquadraticcurveto{\pgfqpoint{2.132281in}{9.445556in}}{\pgfqpoint{2.142092in}{9.434694in}}%
\pgfpathquadraticcurveto{\pgfqpoint{2.151904in}{9.423856in}}{\pgfqpoint{2.151904in}{9.405213in}}%
\pgfpathlineto{\pgfqpoint{2.151904in}{9.405213in}}%
\pgfpathclose%
\pgfpathmoveto{\pgfqpoint{2.138178in}{9.409247in}}%
\pgfpathquadraticcurveto{\pgfqpoint{2.138034in}{9.420490in}}{\pgfqpoint{2.131875in}{9.427198in}}%
\pgfpathquadraticcurveto{\pgfqpoint{2.125717in}{9.433930in}}{\pgfqpoint{2.115571in}{9.433930in}}%
\pgfpathquadraticcurveto{\pgfqpoint{2.104089in}{9.433930in}}{\pgfqpoint{2.097190in}{9.427437in}}%
\pgfpathquadraticcurveto{\pgfqpoint{2.090291in}{9.420944in}}{\pgfqpoint{2.089241in}{9.409151in}}%
\pgfpathlineto{\pgfqpoint{2.138178in}{9.409247in}}%
\pgfpathlineto{\pgfqpoint{2.138178in}{9.409247in}}%
\pgfpathclose%
\pgfpathmoveto{\pgfqpoint{2.229414in}{9.430875in}}%
\pgfpathlineto{\pgfqpoint{2.229414in}{9.476087in}}%
\pgfpathlineto{\pgfqpoint{2.243141in}{9.476087in}}%
\pgfpathlineto{\pgfqpoint{2.243141in}{9.360000in}}%
\pgfpathlineto{\pgfqpoint{2.229414in}{9.360000in}}%
\pgfpathlineto{\pgfqpoint{2.229414in}{9.372533in}}%
\pgfpathquadraticcurveto{\pgfqpoint{2.225094in}{9.365085in}}{\pgfqpoint{2.218481in}{9.361456in}}%
\pgfpathquadraticcurveto{\pgfqpoint{2.211893in}{9.357828in}}{\pgfqpoint{2.202631in}{9.357828in}}%
\pgfpathquadraticcurveto{\pgfqpoint{2.187496in}{9.357828in}}{\pgfqpoint{2.177971in}{9.369907in}}%
\pgfpathquadraticcurveto{\pgfqpoint{2.168470in}{9.382010in}}{\pgfqpoint{2.168470in}{9.401704in}}%
\pgfpathquadraticcurveto{\pgfqpoint{2.168470in}{9.421398in}}{\pgfqpoint{2.177971in}{9.433477in}}%
\pgfpathquadraticcurveto{\pgfqpoint{2.187496in}{9.445556in}}{\pgfqpoint{2.202631in}{9.445556in}}%
\pgfpathquadraticcurveto{\pgfqpoint{2.211893in}{9.445556in}}{\pgfqpoint{2.218481in}{9.441927in}}%
\pgfpathquadraticcurveto{\pgfqpoint{2.225094in}{9.438322in}}{\pgfqpoint{2.229414in}{9.430875in}}%
\pgfpathlineto{\pgfqpoint{2.229414in}{9.430875in}}%
\pgfpathclose%
\pgfpathmoveto{\pgfqpoint{2.182650in}{9.401704in}}%
\pgfpathquadraticcurveto{\pgfqpoint{2.182650in}{9.386569in}}{\pgfqpoint{2.188881in}{9.377951in}}%
\pgfpathquadraticcurveto{\pgfqpoint{2.195111in}{9.369334in}}{\pgfqpoint{2.205997in}{9.369334in}}%
\pgfpathquadraticcurveto{\pgfqpoint{2.216882in}{9.369334in}}{\pgfqpoint{2.223136in}{9.377951in}}%
\pgfpathquadraticcurveto{\pgfqpoint{2.229414in}{9.386569in}}{\pgfqpoint{2.229414in}{9.401704in}}%
\pgfpathquadraticcurveto{\pgfqpoint{2.229414in}{9.416838in}}{\pgfqpoint{2.223136in}{9.425456in}}%
\pgfpathquadraticcurveto{\pgfqpoint{2.216882in}{9.434073in}}{\pgfqpoint{2.205997in}{9.434073in}}%
\pgfpathquadraticcurveto{\pgfqpoint{2.195111in}{9.434073in}}{\pgfqpoint{2.188881in}{9.425456in}}%
\pgfpathquadraticcurveto{\pgfqpoint{2.182650in}{9.416838in}}{\pgfqpoint{2.182650in}{9.401704in}}%
\pgfpathlineto{\pgfqpoint{2.182650in}{9.401704in}}%
\pgfpathclose%
\pgfpathmoveto{\pgfqpoint{2.274938in}{9.378954in}}%
\pgfpathlineto{\pgfqpoint{2.290669in}{9.378954in}}%
\pgfpathlineto{\pgfqpoint{2.290669in}{9.360000in}}%
\pgfpathlineto{\pgfqpoint{2.274938in}{9.360000in}}%
\pgfpathlineto{\pgfqpoint{2.274938in}{9.378954in}}%
\pgfpathlineto{\pgfqpoint{2.274938in}{9.378954in}}%
\pgfpathclose%
\pgfpathmoveto{\pgfqpoint{2.274938in}{9.438991in}}%
\pgfpathlineto{\pgfqpoint{2.290669in}{9.438991in}}%
\pgfpathlineto{\pgfqpoint{2.290669in}{9.420061in}}%
\pgfpathlineto{\pgfqpoint{2.274938in}{9.420061in}}%
\pgfpathlineto{\pgfqpoint{2.274938in}{9.438991in}}%
\pgfpathlineto{\pgfqpoint{2.274938in}{9.438991in}}%
\pgfpathclose%
\pgfpathmoveto{\pgfqpoint{2.384723in}{9.372533in}}%
\pgfpathlineto{\pgfqpoint{2.384723in}{9.328227in}}%
\pgfpathlineto{\pgfqpoint{2.370925in}{9.328227in}}%
\pgfpathlineto{\pgfqpoint{2.370925in}{9.443550in}}%
\pgfpathlineto{\pgfqpoint{2.384723in}{9.443550in}}%
\pgfpathlineto{\pgfqpoint{2.384723in}{9.430875in}}%
\pgfpathquadraticcurveto{\pgfqpoint{2.389067in}{9.438322in}}{\pgfqpoint{2.395656in}{9.441927in}}%
\pgfpathquadraticcurveto{\pgfqpoint{2.402268in}{9.445556in}}{\pgfqpoint{2.411435in}{9.445556in}}%
\pgfpathquadraticcurveto{\pgfqpoint{2.426665in}{9.445556in}}{\pgfqpoint{2.436166in}{9.433477in}}%
\pgfpathquadraticcurveto{\pgfqpoint{2.445691in}{9.421398in}}{\pgfqpoint{2.445691in}{9.401704in}}%
\pgfpathquadraticcurveto{\pgfqpoint{2.445691in}{9.382010in}}{\pgfqpoint{2.436166in}{9.369907in}}%
\pgfpathquadraticcurveto{\pgfqpoint{2.426665in}{9.357828in}}{\pgfqpoint{2.411435in}{9.357828in}}%
\pgfpathquadraticcurveto{\pgfqpoint{2.402268in}{9.357828in}}{\pgfqpoint{2.395656in}{9.361456in}}%
\pgfpathquadraticcurveto{\pgfqpoint{2.389067in}{9.365085in}}{\pgfqpoint{2.384723in}{9.372533in}}%
\pgfpathlineto{\pgfqpoint{2.384723in}{9.372533in}}%
\pgfpathclose%
\pgfpathmoveto{\pgfqpoint{2.431439in}{9.401704in}}%
\pgfpathquadraticcurveto{\pgfqpoint{2.431439in}{9.416838in}}{\pgfqpoint{2.425209in}{9.425456in}}%
\pgfpathquadraticcurveto{\pgfqpoint{2.418978in}{9.434073in}}{\pgfqpoint{2.408093in}{9.434073in}}%
\pgfpathquadraticcurveto{\pgfqpoint{2.397184in}{9.434073in}}{\pgfqpoint{2.390953in}{9.425456in}}%
\pgfpathquadraticcurveto{\pgfqpoint{2.384723in}{9.416838in}}{\pgfqpoint{2.384723in}{9.401704in}}%
\pgfpathquadraticcurveto{\pgfqpoint{2.384723in}{9.386569in}}{\pgfqpoint{2.390953in}{9.377951in}}%
\pgfpathquadraticcurveto{\pgfqpoint{2.397184in}{9.369334in}}{\pgfqpoint{2.408093in}{9.369334in}}%
\pgfpathquadraticcurveto{\pgfqpoint{2.418978in}{9.369334in}}{\pgfqpoint{2.425209in}{9.377951in}}%
\pgfpathquadraticcurveto{\pgfqpoint{2.431439in}{9.386569in}}{\pgfqpoint{2.431439in}{9.401704in}}%
\pgfpathlineto{\pgfqpoint{2.431439in}{9.401704in}}%
\pgfpathclose%
\pgfpathmoveto{\pgfqpoint{2.506420in}{9.401990in}}%
\pgfpathquadraticcurveto{\pgfqpoint{2.489781in}{9.401990in}}{\pgfqpoint{2.483360in}{9.398194in}}%
\pgfpathquadraticcurveto{\pgfqpoint{2.476938in}{9.394399in}}{\pgfqpoint{2.476938in}{9.385208in}}%
\pgfpathquadraticcurveto{\pgfqpoint{2.476938in}{9.377904in}}{\pgfqpoint{2.481760in}{9.373607in}}%
\pgfpathquadraticcurveto{\pgfqpoint{2.486582in}{9.369334in}}{\pgfqpoint{2.494842in}{9.369334in}}%
\pgfpathquadraticcurveto{\pgfqpoint{2.506276in}{9.369334in}}{\pgfqpoint{2.513175in}{9.377426in}}%
\pgfpathquadraticcurveto{\pgfqpoint{2.520074in}{9.385519in}}{\pgfqpoint{2.520074in}{9.398934in}}%
\pgfpathlineto{\pgfqpoint{2.520074in}{9.401990in}}%
\pgfpathlineto{\pgfqpoint{2.506420in}{9.401990in}}%
\pgfpathlineto{\pgfqpoint{2.506420in}{9.401990in}}%
\pgfpathclose%
\pgfpathmoveto{\pgfqpoint{2.533800in}{9.407671in}}%
\pgfpathlineto{\pgfqpoint{2.533800in}{9.360000in}}%
\pgfpathlineto{\pgfqpoint{2.520074in}{9.360000in}}%
\pgfpathlineto{\pgfqpoint{2.520074in}{9.372676in}}%
\pgfpathquadraticcurveto{\pgfqpoint{2.515372in}{9.365085in}}{\pgfqpoint{2.508353in}{9.361456in}}%
\pgfpathquadraticcurveto{\pgfqpoint{2.501335in}{9.357828in}}{\pgfqpoint{2.491190in}{9.357828in}}%
\pgfpathquadraticcurveto{\pgfqpoint{2.478371in}{9.357828in}}{\pgfqpoint{2.470780in}{9.365037in}}%
\pgfpathquadraticcurveto{\pgfqpoint{2.463212in}{9.372246in}}{\pgfqpoint{2.463212in}{9.384325in}}%
\pgfpathquadraticcurveto{\pgfqpoint{2.463212in}{9.398409in}}{\pgfqpoint{2.472641in}{9.405571in}}%
\pgfpathquadraticcurveto{\pgfqpoint{2.482095in}{9.412732in}}{\pgfqpoint{2.500810in}{9.412732in}}%
\pgfpathlineto{\pgfqpoint{2.520074in}{9.412732in}}%
\pgfpathlineto{\pgfqpoint{2.520074in}{9.414093in}}%
\pgfpathquadraticcurveto{\pgfqpoint{2.520074in}{9.423570in}}{\pgfqpoint{2.513844in}{9.428750in}}%
\pgfpathquadraticcurveto{\pgfqpoint{2.507613in}{9.433930in}}{\pgfqpoint{2.496346in}{9.433930in}}%
\pgfpathquadraticcurveto{\pgfqpoint{2.489184in}{9.433930in}}{\pgfqpoint{2.482381in}{9.432211in}}%
\pgfpathquadraticcurveto{\pgfqpoint{2.475602in}{9.430493in}}{\pgfqpoint{2.469347in}{9.427055in}}%
\pgfpathlineto{\pgfqpoint{2.469347in}{9.439755in}}%
\pgfpathquadraticcurveto{\pgfqpoint{2.476867in}{9.442667in}}{\pgfqpoint{2.483957in}{9.444099in}}%
\pgfpathquadraticcurveto{\pgfqpoint{2.491046in}{9.445556in}}{\pgfqpoint{2.497754in}{9.445556in}}%
\pgfpathquadraticcurveto{\pgfqpoint{2.515897in}{9.445556in}}{\pgfqpoint{2.524849in}{9.436150in}}%
\pgfpathquadraticcurveto{\pgfqpoint{2.533800in}{9.426769in}}{\pgfqpoint{2.533800in}{9.407671in}}%
\pgfpathlineto{\pgfqpoint{2.533800in}{9.407671in}}%
\pgfpathclose%
\pgfpathmoveto{\pgfqpoint{2.562064in}{9.443550in}}%
\pgfpathlineto{\pgfqpoint{2.575790in}{9.443550in}}%
\pgfpathlineto{\pgfqpoint{2.575790in}{9.360000in}}%
\pgfpathlineto{\pgfqpoint{2.562064in}{9.360000in}}%
\pgfpathlineto{\pgfqpoint{2.562064in}{9.443550in}}%
\pgfpathlineto{\pgfqpoint{2.562064in}{9.443550in}}%
\pgfpathclose%
\pgfpathmoveto{\pgfqpoint{2.562064in}{9.476087in}}%
\pgfpathlineto{\pgfqpoint{2.575790in}{9.476087in}}%
\pgfpathlineto{\pgfqpoint{2.575790in}{9.458685in}}%
\pgfpathlineto{\pgfqpoint{2.562064in}{9.458685in}}%
\pgfpathlineto{\pgfqpoint{2.562064in}{9.476087in}}%
\pgfpathlineto{\pgfqpoint{2.562064in}{9.476087in}}%
\pgfpathclose%
\pgfpathmoveto{\pgfqpoint{2.652919in}{9.430731in}}%
\pgfpathquadraticcurveto{\pgfqpoint{2.650604in}{9.432068in}}{\pgfqpoint{2.647882in}{9.432689in}}%
\pgfpathquadraticcurveto{\pgfqpoint{2.645161in}{9.433333in}}{\pgfqpoint{2.641891in}{9.433333in}}%
\pgfpathquadraticcurveto{\pgfqpoint{2.630241in}{9.433333in}}{\pgfqpoint{2.624011in}{9.425766in}}%
\pgfpathquadraticcurveto{\pgfqpoint{2.617780in}{9.418199in}}{\pgfqpoint{2.617780in}{9.404019in}}%
\pgfpathlineto{\pgfqpoint{2.617780in}{9.360000in}}%
\pgfpathlineto{\pgfqpoint{2.603983in}{9.360000in}}%
\pgfpathlineto{\pgfqpoint{2.603983in}{9.443550in}}%
\pgfpathlineto{\pgfqpoint{2.617780in}{9.443550in}}%
\pgfpathlineto{\pgfqpoint{2.617780in}{9.430564in}}%
\pgfpathquadraticcurveto{\pgfqpoint{2.622125in}{9.438179in}}{\pgfqpoint{2.629048in}{9.441855in}}%
\pgfpathquadraticcurveto{\pgfqpoint{2.635994in}{9.445556in}}{\pgfqpoint{2.645925in}{9.445556in}}%
\pgfpathquadraticcurveto{\pgfqpoint{2.647333in}{9.445556in}}{\pgfqpoint{2.649052in}{9.445365in}}%
\pgfpathquadraticcurveto{\pgfqpoint{2.650771in}{9.445197in}}{\pgfqpoint{2.652848in}{9.444816in}}%
\pgfpathlineto{\pgfqpoint{2.652919in}{9.430731in}}%
\pgfpathlineto{\pgfqpoint{2.652919in}{9.430731in}}%
\pgfpathclose%
\pgfpathmoveto{\pgfqpoint{2.735419in}{9.405213in}}%
\pgfpathlineto{\pgfqpoint{2.735419in}{9.398505in}}%
\pgfpathlineto{\pgfqpoint{2.672303in}{9.398505in}}%
\pgfpathquadraticcurveto{\pgfqpoint{2.673210in}{9.384325in}}{\pgfqpoint{2.680849in}{9.376901in}}%
\pgfpathquadraticcurveto{\pgfqpoint{2.688488in}{9.369477in}}{\pgfqpoint{2.702142in}{9.369477in}}%
\pgfpathquadraticcurveto{\pgfqpoint{2.710044in}{9.369477in}}{\pgfqpoint{2.717468in}{9.371411in}}%
\pgfpathquadraticcurveto{\pgfqpoint{2.724892in}{9.373344in}}{\pgfqpoint{2.732220in}{9.377235in}}%
\pgfpathlineto{\pgfqpoint{2.732220in}{9.364249in}}%
\pgfpathquadraticcurveto{\pgfqpoint{2.724820in}{9.361122in}}{\pgfqpoint{2.717062in}{9.359475in}}%
\pgfpathquadraticcurveto{\pgfqpoint{2.709304in}{9.357828in}}{\pgfqpoint{2.701331in}{9.357828in}}%
\pgfpathquadraticcurveto{\pgfqpoint{2.681326in}{9.357828in}}{\pgfqpoint{2.669653in}{9.369453in}}%
\pgfpathquadraticcurveto{\pgfqpoint{2.657980in}{9.381102in}}{\pgfqpoint{2.657980in}{9.400964in}}%
\pgfpathquadraticcurveto{\pgfqpoint{2.657980in}{9.421469in}}{\pgfqpoint{2.669056in}{9.433500in}}%
\pgfpathquadraticcurveto{\pgfqpoint{2.680133in}{9.445556in}}{\pgfqpoint{2.698944in}{9.445556in}}%
\pgfpathquadraticcurveto{\pgfqpoint{2.715797in}{9.445556in}}{\pgfqpoint{2.725608in}{9.434694in}}%
\pgfpathquadraticcurveto{\pgfqpoint{2.735419in}{9.423856in}}{\pgfqpoint{2.735419in}{9.405213in}}%
\pgfpathlineto{\pgfqpoint{2.735419in}{9.405213in}}%
\pgfpathclose%
\pgfpathmoveto{\pgfqpoint{2.721693in}{9.409247in}}%
\pgfpathquadraticcurveto{\pgfqpoint{2.721550in}{9.420490in}}{\pgfqpoint{2.715391in}{9.427198in}}%
\pgfpathquadraticcurveto{\pgfqpoint{2.709232in}{9.433930in}}{\pgfqpoint{2.699087in}{9.433930in}}%
\pgfpathquadraticcurveto{\pgfqpoint{2.687605in}{9.433930in}}{\pgfqpoint{2.680706in}{9.427437in}}%
\pgfpathquadraticcurveto{\pgfqpoint{2.673807in}{9.420944in}}{\pgfqpoint{2.672757in}{9.409151in}}%
\pgfpathlineto{\pgfqpoint{2.721693in}{9.409247in}}%
\pgfpathlineto{\pgfqpoint{2.721693in}{9.409247in}}%
\pgfpathclose%
\pgfpathmoveto{\pgfqpoint{2.812930in}{9.430875in}}%
\pgfpathlineto{\pgfqpoint{2.812930in}{9.476087in}}%
\pgfpathlineto{\pgfqpoint{2.826656in}{9.476087in}}%
\pgfpathlineto{\pgfqpoint{2.826656in}{9.360000in}}%
\pgfpathlineto{\pgfqpoint{2.812930in}{9.360000in}}%
\pgfpathlineto{\pgfqpoint{2.812930in}{9.372533in}}%
\pgfpathquadraticcurveto{\pgfqpoint{2.808609in}{9.365085in}}{\pgfqpoint{2.801997in}{9.361456in}}%
\pgfpathquadraticcurveto{\pgfqpoint{2.795408in}{9.357828in}}{\pgfqpoint{2.786146in}{9.357828in}}%
\pgfpathquadraticcurveto{\pgfqpoint{2.771012in}{9.357828in}}{\pgfqpoint{2.761487in}{9.369907in}}%
\pgfpathquadraticcurveto{\pgfqpoint{2.751986in}{9.382010in}}{\pgfqpoint{2.751986in}{9.401704in}}%
\pgfpathquadraticcurveto{\pgfqpoint{2.751986in}{9.421398in}}{\pgfqpoint{2.761487in}{9.433477in}}%
\pgfpathquadraticcurveto{\pgfqpoint{2.771012in}{9.445556in}}{\pgfqpoint{2.786146in}{9.445556in}}%
\pgfpathquadraticcurveto{\pgfqpoint{2.795408in}{9.445556in}}{\pgfqpoint{2.801997in}{9.441927in}}%
\pgfpathquadraticcurveto{\pgfqpoint{2.808609in}{9.438322in}}{\pgfqpoint{2.812930in}{9.430875in}}%
\pgfpathlineto{\pgfqpoint{2.812930in}{9.430875in}}%
\pgfpathclose%
\pgfpathmoveto{\pgfqpoint{2.766166in}{9.401704in}}%
\pgfpathquadraticcurveto{\pgfqpoint{2.766166in}{9.386569in}}{\pgfqpoint{2.772396in}{9.377951in}}%
\pgfpathquadraticcurveto{\pgfqpoint{2.778627in}{9.369334in}}{\pgfqpoint{2.789512in}{9.369334in}}%
\pgfpathquadraticcurveto{\pgfqpoint{2.800398in}{9.369334in}}{\pgfqpoint{2.806652in}{9.377951in}}%
\pgfpathquadraticcurveto{\pgfqpoint{2.812930in}{9.386569in}}{\pgfqpoint{2.812930in}{9.401704in}}%
\pgfpathquadraticcurveto{\pgfqpoint{2.812930in}{9.416838in}}{\pgfqpoint{2.806652in}{9.425456in}}%
\pgfpathquadraticcurveto{\pgfqpoint{2.800398in}{9.434073in}}{\pgfqpoint{2.789512in}{9.434073in}}%
\pgfpathquadraticcurveto{\pgfqpoint{2.778627in}{9.434073in}}{\pgfqpoint{2.772396in}{9.425456in}}%
\pgfpathquadraticcurveto{\pgfqpoint{2.766166in}{9.416838in}}{\pgfqpoint{2.766166in}{9.401704in}}%
\pgfpathlineto{\pgfqpoint{2.766166in}{9.401704in}}%
\pgfpathclose%
\pgfpathmoveto{\pgfqpoint{2.921187in}{9.476756in}}%
\pgfpathlineto{\pgfqpoint{2.921187in}{9.323978in}}%
\pgfpathlineto{\pgfqpoint{2.908512in}{9.323978in}}%
\pgfpathlineto{\pgfqpoint{2.908512in}{9.476756in}}%
\pgfpathlineto{\pgfqpoint{2.921187in}{9.476756in}}%
\pgfpathlineto{\pgfqpoint{2.921187in}{9.476756in}}%
\pgfpathclose%
\pgfpathmoveto{\pgfqpoint{2.968549in}{9.467279in}}%
\pgfpathlineto{\pgfqpoint{2.968549in}{9.443550in}}%
\pgfpathlineto{\pgfqpoint{2.996812in}{9.443550in}}%
\pgfpathlineto{\pgfqpoint{2.996812in}{9.432880in}}%
\pgfpathlineto{\pgfqpoint{2.968549in}{9.432880in}}%
\pgfpathlineto{\pgfqpoint{2.968549in}{9.387524in}}%
\pgfpathquadraticcurveto{\pgfqpoint{2.968549in}{9.377307in}}{\pgfqpoint{2.971342in}{9.374395in}}%
\pgfpathquadraticcurveto{\pgfqpoint{2.974135in}{9.371482in}}{\pgfqpoint{2.982728in}{9.371482in}}%
\pgfpathlineto{\pgfqpoint{2.996812in}{9.371482in}}%
\pgfpathlineto{\pgfqpoint{2.996812in}{9.360000in}}%
\pgfpathlineto{\pgfqpoint{2.982728in}{9.360000in}}%
\pgfpathquadraticcurveto{\pgfqpoint{2.966830in}{9.360000in}}{\pgfqpoint{2.960790in}{9.365920in}}%
\pgfpathquadraticcurveto{\pgfqpoint{2.954751in}{9.371864in}}{\pgfqpoint{2.954751in}{9.387524in}}%
\pgfpathlineto{\pgfqpoint{2.954751in}{9.432880in}}%
\pgfpathlineto{\pgfqpoint{2.944677in}{9.432880in}}%
\pgfpathlineto{\pgfqpoint{2.944677in}{9.443550in}}%
\pgfpathlineto{\pgfqpoint{2.954751in}{9.443550in}}%
\pgfpathlineto{\pgfqpoint{2.954751in}{9.467279in}}%
\pgfpathlineto{\pgfqpoint{2.968549in}{9.467279in}}%
\pgfpathlineto{\pgfqpoint{2.968549in}{9.467279in}}%
\pgfpathclose%
\pgfpathmoveto{\pgfqpoint{3.032548in}{9.476756in}}%
\pgfpathlineto{\pgfqpoint{3.032548in}{9.323978in}}%
\pgfpathlineto{\pgfqpoint{3.019872in}{9.323978in}}%
\pgfpathlineto{\pgfqpoint{3.019872in}{9.476756in}}%
\pgfpathlineto{\pgfqpoint{3.032548in}{9.476756in}}%
\pgfpathlineto{\pgfqpoint{3.032548in}{9.476756in}}%
\pgfpathclose%
\pgfpathmoveto{\pgfqpoint{3.116671in}{9.435195in}}%
\pgfpathlineto{\pgfqpoint{3.116671in}{9.448778in}}%
\pgfpathlineto{\pgfqpoint{3.212301in}{9.414093in}}%
\pgfpathlineto{\pgfqpoint{3.212301in}{9.401704in}}%
\pgfpathlineto{\pgfqpoint{3.116671in}{9.367018in}}%
\pgfpathlineto{\pgfqpoint{3.116671in}{9.380601in}}%
\pgfpathlineto{\pgfqpoint{3.193514in}{9.407815in}}%
\pgfpathlineto{\pgfqpoint{3.116671in}{9.435195in}}%
\pgfpathlineto{\pgfqpoint{3.116671in}{9.435195in}}%
\pgfpathclose%
\pgfpathmoveto{\pgfqpoint{3.296018in}{9.372676in}}%
\pgfpathlineto{\pgfqpoint{3.320630in}{9.372676in}}%
\pgfpathlineto{\pgfqpoint{3.320630in}{9.457658in}}%
\pgfpathlineto{\pgfqpoint{3.293846in}{9.452287in}}%
\pgfpathlineto{\pgfqpoint{3.293846in}{9.466013in}}%
\pgfpathlineto{\pgfqpoint{3.320487in}{9.471385in}}%
\pgfpathlineto{\pgfqpoint{3.335549in}{9.471385in}}%
\pgfpathlineto{\pgfqpoint{3.335549in}{9.372676in}}%
\pgfpathlineto{\pgfqpoint{3.360161in}{9.372676in}}%
\pgfpathlineto{\pgfqpoint{3.360161in}{9.360000in}}%
\pgfpathlineto{\pgfqpoint{3.296018in}{9.360000in}}%
\pgfpathlineto{\pgfqpoint{3.296018in}{9.372676in}}%
\pgfpathlineto{\pgfqpoint{3.296018in}{9.372676in}}%
\pgfpathclose%
\pgfpathmoveto{\pgfqpoint{3.390597in}{9.378954in}}%
\pgfpathlineto{\pgfqpoint{3.406352in}{9.378954in}}%
\pgfpathlineto{\pgfqpoint{3.406352in}{9.360000in}}%
\pgfpathlineto{\pgfqpoint{3.390597in}{9.360000in}}%
\pgfpathlineto{\pgfqpoint{3.390597in}{9.378954in}}%
\pgfpathlineto{\pgfqpoint{3.390597in}{9.378954in}}%
\pgfpathclose%
\pgfpathmoveto{\pgfqpoint{3.439605in}{9.362316in}}%
\pgfpathlineto{\pgfqpoint{3.439605in}{9.376042in}}%
\pgfpathquadraticcurveto{\pgfqpoint{3.445287in}{9.373344in}}{\pgfqpoint{3.451088in}{9.371936in}}%
\pgfpathquadraticcurveto{\pgfqpoint{3.456912in}{9.370527in}}{\pgfqpoint{3.462522in}{9.370527in}}%
\pgfpathquadraticcurveto{\pgfqpoint{3.477442in}{9.370527in}}{\pgfqpoint{3.485296in}{9.380553in}}%
\pgfpathquadraticcurveto{\pgfqpoint{3.493173in}{9.390579in}}{\pgfqpoint{3.494295in}{9.411037in}}%
\pgfpathquadraticcurveto{\pgfqpoint{3.489974in}{9.404616in}}{\pgfqpoint{3.483314in}{9.401178in}}%
\pgfpathquadraticcurveto{\pgfqpoint{3.476678in}{9.397741in}}{\pgfqpoint{3.468633in}{9.397741in}}%
\pgfpathquadraticcurveto{\pgfqpoint{3.451923in}{9.397741in}}{\pgfqpoint{3.442184in}{9.407839in}}%
\pgfpathquadraticcurveto{\pgfqpoint{3.432444in}{9.417960in}}{\pgfqpoint{3.432444in}{9.435506in}}%
\pgfpathquadraticcurveto{\pgfqpoint{3.432444in}{9.452645in}}{\pgfqpoint{3.442589in}{9.463006in}}%
\pgfpathquadraticcurveto{\pgfqpoint{3.452735in}{9.473390in}}{\pgfqpoint{3.469588in}{9.473390in}}%
\pgfpathquadraticcurveto{\pgfqpoint{3.488924in}{9.473390in}}{\pgfqpoint{3.499093in}{9.458566in}}%
\pgfpathquadraticcurveto{\pgfqpoint{3.509286in}{9.443765in}}{\pgfqpoint{3.509286in}{9.415573in}}%
\pgfpathquadraticcurveto{\pgfqpoint{3.509286in}{9.389243in}}{\pgfqpoint{3.496778in}{9.373535in}}%
\pgfpathquadraticcurveto{\pgfqpoint{3.484293in}{9.357828in}}{\pgfqpoint{3.463191in}{9.357828in}}%
\pgfpathquadraticcurveto{\pgfqpoint{3.457509in}{9.357828in}}{\pgfqpoint{3.451684in}{9.358950in}}%
\pgfpathquadraticcurveto{\pgfqpoint{3.445884in}{9.360072in}}{\pgfqpoint{3.439605in}{9.362316in}}%
\pgfpathlineto{\pgfqpoint{3.439605in}{9.362316in}}%
\pgfpathclose%
\pgfpathmoveto{\pgfqpoint{3.469588in}{9.409533in}}%
\pgfpathquadraticcurveto{\pgfqpoint{3.479734in}{9.409533in}}{\pgfqpoint{3.485654in}{9.416456in}}%
\pgfpathquadraticcurveto{\pgfqpoint{3.491598in}{9.423403in}}{\pgfqpoint{3.491598in}{9.435506in}}%
\pgfpathquadraticcurveto{\pgfqpoint{3.491598in}{9.447513in}}{\pgfqpoint{3.485654in}{9.454484in}}%
\pgfpathquadraticcurveto{\pgfqpoint{3.479734in}{9.461454in}}{\pgfqpoint{3.469588in}{9.461454in}}%
\pgfpathquadraticcurveto{\pgfqpoint{3.459443in}{9.461454in}}{\pgfqpoint{3.453523in}{9.454484in}}%
\pgfpathquadraticcurveto{\pgfqpoint{3.447602in}{9.447513in}}{\pgfqpoint{3.447602in}{9.435506in}}%
\pgfpathquadraticcurveto{\pgfqpoint{3.447602in}{9.423403in}}{\pgfqpoint{3.453523in}{9.416456in}}%
\pgfpathquadraticcurveto{\pgfqpoint{3.459443in}{9.409533in}}{\pgfqpoint{3.469588in}{9.409533in}}%
\pgfpathlineto{\pgfqpoint{3.469588in}{9.409533in}}%
\pgfpathclose%
\pgfpathmoveto{\pgfqpoint{3.570469in}{9.421684in}}%
\pgfpathquadraticcurveto{\pgfqpoint{3.560324in}{9.421684in}}{\pgfqpoint{3.554380in}{9.414737in}}%
\pgfpathquadraticcurveto{\pgfqpoint{3.548460in}{9.407815in}}{\pgfqpoint{3.548460in}{9.395736in}}%
\pgfpathquadraticcurveto{\pgfqpoint{3.548460in}{9.383728in}}{\pgfqpoint{3.554380in}{9.376734in}}%
\pgfpathquadraticcurveto{\pgfqpoint{3.560324in}{9.369763in}}{\pgfqpoint{3.570469in}{9.369763in}}%
\pgfpathquadraticcurveto{\pgfqpoint{3.580615in}{9.369763in}}{\pgfqpoint{3.586535in}{9.376734in}}%
\pgfpathquadraticcurveto{\pgfqpoint{3.592455in}{9.383728in}}{\pgfqpoint{3.592455in}{9.395736in}}%
\pgfpathquadraticcurveto{\pgfqpoint{3.592455in}{9.407815in}}{\pgfqpoint{3.586535in}{9.414737in}}%
\pgfpathquadraticcurveto{\pgfqpoint{3.580615in}{9.421684in}}{\pgfqpoint{3.570469in}{9.421684in}}%
\pgfpathlineto{\pgfqpoint{3.570469in}{9.421684in}}%
\pgfpathclose%
\pgfpathmoveto{\pgfqpoint{3.600380in}{9.468926in}}%
\pgfpathlineto{\pgfqpoint{3.600380in}{9.455200in}}%
\pgfpathquadraticcurveto{\pgfqpoint{3.594699in}{9.457873in}}{\pgfqpoint{3.588922in}{9.459282in}}%
\pgfpathquadraticcurveto{\pgfqpoint{3.583145in}{9.460714in}}{\pgfqpoint{3.577464in}{9.460714in}}%
\pgfpathquadraticcurveto{\pgfqpoint{3.562544in}{9.460714in}}{\pgfqpoint{3.554666in}{9.450640in}}%
\pgfpathquadraticcurveto{\pgfqpoint{3.546812in}{9.440566in}}{\pgfqpoint{3.545691in}{9.420204in}}%
\pgfpathquadraticcurveto{\pgfqpoint{3.550083in}{9.426697in}}{\pgfqpoint{3.556719in}{9.430158in}}%
\pgfpathquadraticcurveto{\pgfqpoint{3.563379in}{9.433620in}}{\pgfqpoint{3.571352in}{9.433620in}}%
\pgfpathquadraticcurveto{\pgfqpoint{3.588134in}{9.433620in}}{\pgfqpoint{3.597874in}{9.423427in}}%
\pgfpathquadraticcurveto{\pgfqpoint{3.607613in}{9.413257in}}{\pgfqpoint{3.607613in}{9.395736in}}%
\pgfpathquadraticcurveto{\pgfqpoint{3.607613in}{9.378572in}}{\pgfqpoint{3.597468in}{9.368188in}}%
\pgfpathquadraticcurveto{\pgfqpoint{3.587322in}{9.357828in}}{\pgfqpoint{3.570469in}{9.357828in}}%
\pgfpathquadraticcurveto{\pgfqpoint{3.551133in}{9.357828in}}{\pgfqpoint{3.540916in}{9.372628in}}%
\pgfpathquadraticcurveto{\pgfqpoint{3.530699in}{9.387452in}}{\pgfqpoint{3.530699in}{9.415573in}}%
\pgfpathquadraticcurveto{\pgfqpoint{3.530699in}{9.441975in}}{\pgfqpoint{3.543232in}{9.457682in}}%
\pgfpathquadraticcurveto{\pgfqpoint{3.555764in}{9.473390in}}{\pgfqpoint{3.576867in}{9.473390in}}%
\pgfpathquadraticcurveto{\pgfqpoint{3.582548in}{9.473390in}}{\pgfqpoint{3.588325in}{9.472268in}}%
\pgfpathquadraticcurveto{\pgfqpoint{3.594102in}{9.471146in}}{\pgfqpoint{3.600380in}{9.468926in}}%
\pgfpathlineto{\pgfqpoint{3.600380in}{9.468926in}}%
\pgfpathclose%
\pgfpathmoveto{\pgfqpoint{3.682116in}{9.422520in}}%
\pgfpathlineto{\pgfqpoint{3.697872in}{9.422520in}}%
\pgfpathlineto{\pgfqpoint{3.697872in}{9.403566in}}%
\pgfpathlineto{\pgfqpoint{3.682116in}{9.403566in}}%
\pgfpathlineto{\pgfqpoint{3.682116in}{9.422520in}}%
\pgfpathlineto{\pgfqpoint{3.682116in}{9.422520in}}%
\pgfpathclose%
\pgfpathmoveto{\pgfqpoint{3.777889in}{9.471385in}}%
\pgfpathlineto{\pgfqpoint{3.792952in}{9.471385in}}%
\pgfpathlineto{\pgfqpoint{3.792952in}{9.425718in}}%
\pgfpathlineto{\pgfqpoint{3.847713in}{9.425718in}}%
\pgfpathlineto{\pgfqpoint{3.847713in}{9.471385in}}%
\pgfpathlineto{\pgfqpoint{3.862776in}{9.471385in}}%
\pgfpathlineto{\pgfqpoint{3.862776in}{9.360000in}}%
\pgfpathlineto{\pgfqpoint{3.847713in}{9.360000in}}%
\pgfpathlineto{\pgfqpoint{3.847713in}{9.413043in}}%
\pgfpathlineto{\pgfqpoint{3.792952in}{9.413043in}}%
\pgfpathlineto{\pgfqpoint{3.792952in}{9.360000in}}%
\pgfpathlineto{\pgfqpoint{3.777889in}{9.360000in}}%
\pgfpathlineto{\pgfqpoint{3.777889in}{9.471385in}}%
\pgfpathlineto{\pgfqpoint{3.777889in}{9.471385in}}%
\pgfpathclose%
\pgfpathmoveto{\pgfqpoint{3.924556in}{9.433930in}}%
\pgfpathquadraticcurveto{\pgfqpoint{3.913527in}{9.433930in}}{\pgfqpoint{3.907105in}{9.425313in}}%
\pgfpathquadraticcurveto{\pgfqpoint{3.900684in}{9.416695in}}{\pgfqpoint{3.900684in}{9.401704in}}%
\pgfpathquadraticcurveto{\pgfqpoint{3.900684in}{9.386712in}}{\pgfqpoint{3.907058in}{9.378095in}}%
\pgfpathquadraticcurveto{\pgfqpoint{3.913455in}{9.369477in}}{\pgfqpoint{3.924556in}{9.369477in}}%
\pgfpathquadraticcurveto{\pgfqpoint{3.935536in}{9.369477in}}{\pgfqpoint{3.941934in}{9.378118in}}%
\pgfpathquadraticcurveto{\pgfqpoint{3.948355in}{9.386784in}}{\pgfqpoint{3.948355in}{9.401704in}}%
\pgfpathquadraticcurveto{\pgfqpoint{3.948355in}{9.416552in}}{\pgfqpoint{3.941934in}{9.425241in}}%
\pgfpathquadraticcurveto{\pgfqpoint{3.935536in}{9.433930in}}{\pgfqpoint{3.924556in}{9.433930in}}%
\pgfpathlineto{\pgfqpoint{3.924556in}{9.433930in}}%
\pgfpathclose%
\pgfpathmoveto{\pgfqpoint{3.924556in}{9.445556in}}%
\pgfpathquadraticcurveto{\pgfqpoint{3.942459in}{9.445556in}}{\pgfqpoint{3.952676in}{9.433906in}}%
\pgfpathquadraticcurveto{\pgfqpoint{3.962917in}{9.422281in}}{\pgfqpoint{3.962917in}{9.401704in}}%
\pgfpathquadraticcurveto{\pgfqpoint{3.962917in}{9.381198in}}{\pgfqpoint{3.952676in}{9.369501in}}%
\pgfpathquadraticcurveto{\pgfqpoint{3.942459in}{9.357828in}}{\pgfqpoint{3.924556in}{9.357828in}}%
\pgfpathquadraticcurveto{\pgfqpoint{3.906580in}{9.357828in}}{\pgfqpoint{3.896387in}{9.369501in}}%
\pgfpathquadraticcurveto{\pgfqpoint{3.886218in}{9.381198in}}{\pgfqpoint{3.886218in}{9.401704in}}%
\pgfpathquadraticcurveto{\pgfqpoint{3.886218in}{9.422281in}}{\pgfqpoint{3.896387in}{9.433906in}}%
\pgfpathquadraticcurveto{\pgfqpoint{3.906580in}{9.445556in}}{\pgfqpoint{3.924556in}{9.445556in}}%
\pgfpathlineto{\pgfqpoint{3.924556in}{9.445556in}}%
\pgfpathclose%
\pgfpathmoveto{\pgfqpoint{3.985667in}{9.476087in}}%
\pgfpathlineto{\pgfqpoint{3.999393in}{9.476087in}}%
\pgfpathlineto{\pgfqpoint{3.999393in}{9.360000in}}%
\pgfpathlineto{\pgfqpoint{3.985667in}{9.360000in}}%
\pgfpathlineto{\pgfqpoint{3.985667in}{9.476087in}}%
\pgfpathlineto{\pgfqpoint{3.985667in}{9.476087in}}%
\pgfpathclose%
\pgfpathmoveto{\pgfqpoint{4.028110in}{9.476087in}}%
\pgfpathlineto{\pgfqpoint{4.041836in}{9.476087in}}%
\pgfpathlineto{\pgfqpoint{4.041836in}{9.360000in}}%
\pgfpathlineto{\pgfqpoint{4.028110in}{9.360000in}}%
\pgfpathlineto{\pgfqpoint{4.028110in}{9.476087in}}%
\pgfpathlineto{\pgfqpoint{4.028110in}{9.476087in}}%
\pgfpathclose%
\pgfpathmoveto{\pgfqpoint{4.102924in}{9.433930in}}%
\pgfpathquadraticcurveto{\pgfqpoint{4.091895in}{9.433930in}}{\pgfqpoint{4.085474in}{9.425313in}}%
\pgfpathquadraticcurveto{\pgfqpoint{4.079052in}{9.416695in}}{\pgfqpoint{4.079052in}{9.401704in}}%
\pgfpathquadraticcurveto{\pgfqpoint{4.079052in}{9.386712in}}{\pgfqpoint{4.085426in}{9.378095in}}%
\pgfpathquadraticcurveto{\pgfqpoint{4.091823in}{9.369477in}}{\pgfqpoint{4.102924in}{9.369477in}}%
\pgfpathquadraticcurveto{\pgfqpoint{4.113905in}{9.369477in}}{\pgfqpoint{4.120302in}{9.378118in}}%
\pgfpathquadraticcurveto{\pgfqpoint{4.126724in}{9.386784in}}{\pgfqpoint{4.126724in}{9.401704in}}%
\pgfpathquadraticcurveto{\pgfqpoint{4.126724in}{9.416552in}}{\pgfqpoint{4.120302in}{9.425241in}}%
\pgfpathquadraticcurveto{\pgfqpoint{4.113905in}{9.433930in}}{\pgfqpoint{4.102924in}{9.433930in}}%
\pgfpathlineto{\pgfqpoint{4.102924in}{9.433930in}}%
\pgfpathclose%
\pgfpathmoveto{\pgfqpoint{4.102924in}{9.445556in}}%
\pgfpathquadraticcurveto{\pgfqpoint{4.120827in}{9.445556in}}{\pgfqpoint{4.131044in}{9.433906in}}%
\pgfpathquadraticcurveto{\pgfqpoint{4.141285in}{9.422281in}}{\pgfqpoint{4.141285in}{9.401704in}}%
\pgfpathquadraticcurveto{\pgfqpoint{4.141285in}{9.381198in}}{\pgfqpoint{4.131044in}{9.369501in}}%
\pgfpathquadraticcurveto{\pgfqpoint{4.120827in}{9.357828in}}{\pgfqpoint{4.102924in}{9.357828in}}%
\pgfpathquadraticcurveto{\pgfqpoint{4.084948in}{9.357828in}}{\pgfqpoint{4.074755in}{9.369501in}}%
\pgfpathquadraticcurveto{\pgfqpoint{4.064586in}{9.381198in}}{\pgfqpoint{4.064586in}{9.401704in}}%
\pgfpathquadraticcurveto{\pgfqpoint{4.064586in}{9.422281in}}{\pgfqpoint{4.074755in}{9.433906in}}%
\pgfpathquadraticcurveto{\pgfqpoint{4.084948in}{9.445556in}}{\pgfqpoint{4.102924in}{9.445556in}}%
\pgfpathlineto{\pgfqpoint{4.102924in}{9.445556in}}%
\pgfpathclose%
\pgfpathmoveto{\pgfqpoint{4.156062in}{9.443550in}}%
\pgfpathlineto{\pgfqpoint{4.169788in}{9.443550in}}%
\pgfpathlineto{\pgfqpoint{4.186951in}{9.378357in}}%
\pgfpathlineto{\pgfqpoint{4.204020in}{9.443550in}}%
\pgfpathlineto{\pgfqpoint{4.220204in}{9.443550in}}%
\pgfpathlineto{\pgfqpoint{4.237368in}{9.378357in}}%
\pgfpathlineto{\pgfqpoint{4.254460in}{9.443550in}}%
\pgfpathlineto{\pgfqpoint{4.268186in}{9.443550in}}%
\pgfpathlineto{\pgfqpoint{4.246320in}{9.360000in}}%
\pgfpathlineto{\pgfqpoint{4.230135in}{9.360000in}}%
\pgfpathlineto{\pgfqpoint{4.212160in}{9.428487in}}%
\pgfpathlineto{\pgfqpoint{4.194113in}{9.360000in}}%
\pgfpathlineto{\pgfqpoint{4.177904in}{9.360000in}}%
\pgfpathlineto{\pgfqpoint{4.156062in}{9.443550in}}%
\pgfpathlineto{\pgfqpoint{4.156062in}{9.443550in}}%
\pgfpathclose%
\pgfpathmoveto{\pgfqpoint{4.284132in}{9.378954in}}%
\pgfpathlineto{\pgfqpoint{4.299864in}{9.378954in}}%
\pgfpathlineto{\pgfqpoint{4.299864in}{9.360000in}}%
\pgfpathlineto{\pgfqpoint{4.284132in}{9.360000in}}%
\pgfpathlineto{\pgfqpoint{4.284132in}{9.378954in}}%
\pgfpathlineto{\pgfqpoint{4.284132in}{9.378954in}}%
\pgfpathclose%
\pgfpathmoveto{\pgfqpoint{4.284132in}{9.438991in}}%
\pgfpathlineto{\pgfqpoint{4.299864in}{9.438991in}}%
\pgfpathlineto{\pgfqpoint{4.299864in}{9.420061in}}%
\pgfpathlineto{\pgfqpoint{4.284132in}{9.420061in}}%
\pgfpathlineto{\pgfqpoint{4.284132in}{9.438991in}}%
\pgfpathlineto{\pgfqpoint{4.284132in}{9.438991in}}%
\pgfpathclose%
\pgfpathmoveto{\pgfqpoint{4.413015in}{9.433930in}}%
\pgfpathquadraticcurveto{\pgfqpoint{4.401986in}{9.433930in}}{\pgfqpoint{4.395565in}{9.425313in}}%
\pgfpathquadraticcurveto{\pgfqpoint{4.389143in}{9.416695in}}{\pgfqpoint{4.389143in}{9.401704in}}%
\pgfpathquadraticcurveto{\pgfqpoint{4.389143in}{9.386712in}}{\pgfqpoint{4.395517in}{9.378095in}}%
\pgfpathquadraticcurveto{\pgfqpoint{4.401914in}{9.369477in}}{\pgfqpoint{4.413015in}{9.369477in}}%
\pgfpathquadraticcurveto{\pgfqpoint{4.423996in}{9.369477in}}{\pgfqpoint{4.430393in}{9.378118in}}%
\pgfpathquadraticcurveto{\pgfqpoint{4.436815in}{9.386784in}}{\pgfqpoint{4.436815in}{9.401704in}}%
\pgfpathquadraticcurveto{\pgfqpoint{4.436815in}{9.416552in}}{\pgfqpoint{4.430393in}{9.425241in}}%
\pgfpathquadraticcurveto{\pgfqpoint{4.423996in}{9.433930in}}{\pgfqpoint{4.413015in}{9.433930in}}%
\pgfpathlineto{\pgfqpoint{4.413015in}{9.433930in}}%
\pgfpathclose%
\pgfpathmoveto{\pgfqpoint{4.413015in}{9.445556in}}%
\pgfpathquadraticcurveto{\pgfqpoint{4.430918in}{9.445556in}}{\pgfqpoint{4.441135in}{9.433906in}}%
\pgfpathquadraticcurveto{\pgfqpoint{4.451376in}{9.422281in}}{\pgfqpoint{4.451376in}{9.401704in}}%
\pgfpathquadraticcurveto{\pgfqpoint{4.451376in}{9.381198in}}{\pgfqpoint{4.441135in}{9.369501in}}%
\pgfpathquadraticcurveto{\pgfqpoint{4.430918in}{9.357828in}}{\pgfqpoint{4.413015in}{9.357828in}}%
\pgfpathquadraticcurveto{\pgfqpoint{4.395039in}{9.357828in}}{\pgfqpoint{4.384846in}{9.369501in}}%
\pgfpathquadraticcurveto{\pgfqpoint{4.374677in}{9.381198in}}{\pgfqpoint{4.374677in}{9.401704in}}%
\pgfpathquadraticcurveto{\pgfqpoint{4.374677in}{9.422281in}}{\pgfqpoint{4.384846in}{9.433906in}}%
\pgfpathquadraticcurveto{\pgfqpoint{4.395039in}{9.445556in}}{\pgfqpoint{4.413015in}{9.445556in}}%
\pgfpathlineto{\pgfqpoint{4.413015in}{9.445556in}}%
\pgfpathclose%
\pgfpathmoveto{\pgfqpoint{4.487709in}{9.467279in}}%
\pgfpathlineto{\pgfqpoint{4.487709in}{9.443550in}}%
\pgfpathlineto{\pgfqpoint{4.515973in}{9.443550in}}%
\pgfpathlineto{\pgfqpoint{4.515973in}{9.432880in}}%
\pgfpathlineto{\pgfqpoint{4.487709in}{9.432880in}}%
\pgfpathlineto{\pgfqpoint{4.487709in}{9.387524in}}%
\pgfpathquadraticcurveto{\pgfqpoint{4.487709in}{9.377307in}}{\pgfqpoint{4.490502in}{9.374395in}}%
\pgfpathquadraticcurveto{\pgfqpoint{4.493295in}{9.371482in}}{\pgfqpoint{4.501888in}{9.371482in}}%
\pgfpathlineto{\pgfqpoint{4.515973in}{9.371482in}}%
\pgfpathlineto{\pgfqpoint{4.515973in}{9.360000in}}%
\pgfpathlineto{\pgfqpoint{4.501888in}{9.360000in}}%
\pgfpathquadraticcurveto{\pgfqpoint{4.485990in}{9.360000in}}{\pgfqpoint{4.479951in}{9.365920in}}%
\pgfpathquadraticcurveto{\pgfqpoint{4.473911in}{9.371864in}}{\pgfqpoint{4.473911in}{9.387524in}}%
\pgfpathlineto{\pgfqpoint{4.473911in}{9.432880in}}%
\pgfpathlineto{\pgfqpoint{4.463837in}{9.432880in}}%
\pgfpathlineto{\pgfqpoint{4.463837in}{9.443550in}}%
\pgfpathlineto{\pgfqpoint{4.473911in}{9.443550in}}%
\pgfpathlineto{\pgfqpoint{4.473911in}{9.467279in}}%
\pgfpathlineto{\pgfqpoint{4.487709in}{9.467279in}}%
\pgfpathlineto{\pgfqpoint{4.487709in}{9.467279in}}%
\pgfpathclose%
\pgfpathmoveto{\pgfqpoint{4.603486in}{9.410441in}}%
\pgfpathlineto{\pgfqpoint{4.603486in}{9.360000in}}%
\pgfpathlineto{\pgfqpoint{4.589760in}{9.360000in}}%
\pgfpathlineto{\pgfqpoint{4.589760in}{9.409987in}}%
\pgfpathquadraticcurveto{\pgfqpoint{4.589760in}{9.421851in}}{\pgfqpoint{4.585128in}{9.427724in}}%
\pgfpathquadraticcurveto{\pgfqpoint{4.580497in}{9.433620in}}{\pgfqpoint{4.571259in}{9.433620in}}%
\pgfpathquadraticcurveto{\pgfqpoint{4.560135in}{9.433620in}}{\pgfqpoint{4.553714in}{9.426530in}}%
\pgfpathquadraticcurveto{\pgfqpoint{4.547292in}{9.419464in}}{\pgfqpoint{4.547292in}{9.407218in}}%
\pgfpathlineto{\pgfqpoint{4.547292in}{9.360000in}}%
\pgfpathlineto{\pgfqpoint{4.533494in}{9.360000in}}%
\pgfpathlineto{\pgfqpoint{4.533494in}{9.476087in}}%
\pgfpathlineto{\pgfqpoint{4.547292in}{9.476087in}}%
\pgfpathlineto{\pgfqpoint{4.547292in}{9.430564in}}%
\pgfpathquadraticcurveto{\pgfqpoint{4.552234in}{9.438108in}}{\pgfqpoint{4.558894in}{9.441832in}}%
\pgfpathquadraticcurveto{\pgfqpoint{4.565578in}{9.445556in}}{\pgfqpoint{4.574315in}{9.445556in}}%
\pgfpathquadraticcurveto{\pgfqpoint{4.588709in}{9.445556in}}{\pgfqpoint{4.596086in}{9.436651in}}%
\pgfpathquadraticcurveto{\pgfqpoint{4.603486in}{9.427747in}}{\pgfqpoint{4.603486in}{9.410441in}}%
\pgfpathlineto{\pgfqpoint{4.603486in}{9.410441in}}%
\pgfpathclose%
\pgfpathmoveto{\pgfqpoint{4.702314in}{9.405213in}}%
\pgfpathlineto{\pgfqpoint{4.702314in}{9.398505in}}%
\pgfpathlineto{\pgfqpoint{4.639197in}{9.398505in}}%
\pgfpathquadraticcurveto{\pgfqpoint{4.640105in}{9.384325in}}{\pgfqpoint{4.647743in}{9.376901in}}%
\pgfpathquadraticcurveto{\pgfqpoint{4.655382in}{9.369477in}}{\pgfqpoint{4.669037in}{9.369477in}}%
\pgfpathquadraticcurveto{\pgfqpoint{4.676938in}{9.369477in}}{\pgfqpoint{4.684362in}{9.371411in}}%
\pgfpathquadraticcurveto{\pgfqpoint{4.691786in}{9.373344in}}{\pgfqpoint{4.699115in}{9.377235in}}%
\pgfpathlineto{\pgfqpoint{4.699115in}{9.364249in}}%
\pgfpathquadraticcurveto{\pgfqpoint{4.691715in}{9.361122in}}{\pgfqpoint{4.683957in}{9.359475in}}%
\pgfpathquadraticcurveto{\pgfqpoint{4.676198in}{9.357828in}}{\pgfqpoint{4.668225in}{9.357828in}}%
\pgfpathquadraticcurveto{\pgfqpoint{4.648221in}{9.357828in}}{\pgfqpoint{4.636548in}{9.369453in}}%
\pgfpathquadraticcurveto{\pgfqpoint{4.624875in}{9.381102in}}{\pgfqpoint{4.624875in}{9.400964in}}%
\pgfpathquadraticcurveto{\pgfqpoint{4.624875in}{9.421469in}}{\pgfqpoint{4.635951in}{9.433500in}}%
\pgfpathquadraticcurveto{\pgfqpoint{4.647027in}{9.445556in}}{\pgfqpoint{4.665838in}{9.445556in}}%
\pgfpathquadraticcurveto{\pgfqpoint{4.682691in}{9.445556in}}{\pgfqpoint{4.692503in}{9.434694in}}%
\pgfpathquadraticcurveto{\pgfqpoint{4.702314in}{9.423856in}}{\pgfqpoint{4.702314in}{9.405213in}}%
\pgfpathlineto{\pgfqpoint{4.702314in}{9.405213in}}%
\pgfpathclose%
\pgfpathmoveto{\pgfqpoint{4.688588in}{9.409247in}}%
\pgfpathquadraticcurveto{\pgfqpoint{4.688444in}{9.420490in}}{\pgfqpoint{4.682286in}{9.427198in}}%
\pgfpathquadraticcurveto{\pgfqpoint{4.676127in}{9.433930in}}{\pgfqpoint{4.665981in}{9.433930in}}%
\pgfpathquadraticcurveto{\pgfqpoint{4.654499in}{9.433930in}}{\pgfqpoint{4.647600in}{9.427437in}}%
\pgfpathquadraticcurveto{\pgfqpoint{4.640701in}{9.420944in}}{\pgfqpoint{4.639651in}{9.409151in}}%
\pgfpathlineto{\pgfqpoint{4.688588in}{9.409247in}}%
\pgfpathlineto{\pgfqpoint{4.688588in}{9.409247in}}%
\pgfpathclose%
\pgfpathmoveto{\pgfqpoint{4.773260in}{9.430731in}}%
\pgfpathquadraticcurveto{\pgfqpoint{4.770944in}{9.432068in}}{\pgfqpoint{4.768223in}{9.432689in}}%
\pgfpathquadraticcurveto{\pgfqpoint{4.765502in}{9.433333in}}{\pgfqpoint{4.762231in}{9.433333in}}%
\pgfpathquadraticcurveto{\pgfqpoint{4.750582in}{9.433333in}}{\pgfqpoint{4.744352in}{9.425766in}}%
\pgfpathquadraticcurveto{\pgfqpoint{4.738121in}{9.418199in}}{\pgfqpoint{4.738121in}{9.404019in}}%
\pgfpathlineto{\pgfqpoint{4.738121in}{9.360000in}}%
\pgfpathlineto{\pgfqpoint{4.724323in}{9.360000in}}%
\pgfpathlineto{\pgfqpoint{4.724323in}{9.443550in}}%
\pgfpathlineto{\pgfqpoint{4.738121in}{9.443550in}}%
\pgfpathlineto{\pgfqpoint{4.738121in}{9.430564in}}%
\pgfpathquadraticcurveto{\pgfqpoint{4.742466in}{9.438179in}}{\pgfqpoint{4.749388in}{9.441855in}}%
\pgfpathquadraticcurveto{\pgfqpoint{4.756335in}{9.445556in}}{\pgfqpoint{4.766266in}{9.445556in}}%
\pgfpathquadraticcurveto{\pgfqpoint{4.767674in}{9.445556in}}{\pgfqpoint{4.769393in}{9.445365in}}%
\pgfpathquadraticcurveto{\pgfqpoint{4.771112in}{9.445197in}}{\pgfqpoint{4.773188in}{9.444816in}}%
\pgfpathlineto{\pgfqpoint{4.773260in}{9.430731in}}%
\pgfpathlineto{\pgfqpoint{4.773260in}{9.430731in}}%
\pgfpathclose%
\pgfpathmoveto{\pgfqpoint{4.779681in}{9.443550in}}%
\pgfpathlineto{\pgfqpoint{4.793408in}{9.443550in}}%
\pgfpathlineto{\pgfqpoint{4.810571in}{9.378357in}}%
\pgfpathlineto{\pgfqpoint{4.827639in}{9.443550in}}%
\pgfpathlineto{\pgfqpoint{4.843824in}{9.443550in}}%
\pgfpathlineto{\pgfqpoint{4.860988in}{9.378357in}}%
\pgfpathlineto{\pgfqpoint{4.878080in}{9.443550in}}%
\pgfpathlineto{\pgfqpoint{4.891806in}{9.443550in}}%
\pgfpathlineto{\pgfqpoint{4.869940in}{9.360000in}}%
\pgfpathlineto{\pgfqpoint{4.853755in}{9.360000in}}%
\pgfpathlineto{\pgfqpoint{4.835780in}{9.428487in}}%
\pgfpathlineto{\pgfqpoint{4.817733in}{9.360000in}}%
\pgfpathlineto{\pgfqpoint{4.801524in}{9.360000in}}%
\pgfpathlineto{\pgfqpoint{4.779681in}{9.443550in}}%
\pgfpathlineto{\pgfqpoint{4.779681in}{9.443550in}}%
\pgfpathclose%
\pgfpathmoveto{\pgfqpoint{4.912598in}{9.443550in}}%
\pgfpathlineto{\pgfqpoint{4.926324in}{9.443550in}}%
\pgfpathlineto{\pgfqpoint{4.926324in}{9.360000in}}%
\pgfpathlineto{\pgfqpoint{4.912598in}{9.360000in}}%
\pgfpathlineto{\pgfqpoint{4.912598in}{9.443550in}}%
\pgfpathlineto{\pgfqpoint{4.912598in}{9.443550in}}%
\pgfpathclose%
\pgfpathmoveto{\pgfqpoint{4.912598in}{9.476087in}}%
\pgfpathlineto{\pgfqpoint{4.926324in}{9.476087in}}%
\pgfpathlineto{\pgfqpoint{4.926324in}{9.458685in}}%
\pgfpathlineto{\pgfqpoint{4.912598in}{9.458685in}}%
\pgfpathlineto{\pgfqpoint{4.912598in}{9.476087in}}%
\pgfpathlineto{\pgfqpoint{4.912598in}{9.476087in}}%
\pgfpathclose%
\pgfpathmoveto{\pgfqpoint{5.008299in}{9.441092in}}%
\pgfpathlineto{\pgfqpoint{5.008299in}{9.428105in}}%
\pgfpathquadraticcurveto{\pgfqpoint{5.002498in}{9.431089in}}{\pgfqpoint{4.996220in}{9.432569in}}%
\pgfpathquadraticcurveto{\pgfqpoint{4.989966in}{9.434073in}}{\pgfqpoint{4.983234in}{9.434073in}}%
\pgfpathquadraticcurveto{\pgfqpoint{4.973017in}{9.434073in}}{\pgfqpoint{4.967908in}{9.430946in}}%
\pgfpathquadraticcurveto{\pgfqpoint{4.962800in}{9.427819in}}{\pgfqpoint{4.962800in}{9.421541in}}%
\pgfpathquadraticcurveto{\pgfqpoint{4.962800in}{9.416766in}}{\pgfqpoint{4.966452in}{9.414045in}}%
\pgfpathquadraticcurveto{\pgfqpoint{4.970105in}{9.411324in}}{\pgfqpoint{4.981157in}{9.408865in}}%
\pgfpathlineto{\pgfqpoint{4.985860in}{9.407815in}}%
\pgfpathquadraticcurveto{\pgfqpoint{5.000469in}{9.404687in}}{\pgfqpoint{5.006628in}{9.398982in}}%
\pgfpathquadraticcurveto{\pgfqpoint{5.012787in}{9.393277in}}{\pgfqpoint{5.012787in}{9.383060in}}%
\pgfpathquadraticcurveto{\pgfqpoint{5.012787in}{9.371411in}}{\pgfqpoint{5.003572in}{9.364607in}}%
\pgfpathquadraticcurveto{\pgfqpoint{4.994358in}{9.357828in}}{\pgfqpoint{4.978245in}{9.357828in}}%
\pgfpathquadraticcurveto{\pgfqpoint{4.971537in}{9.357828in}}{\pgfqpoint{4.964256in}{9.359141in}}%
\pgfpathquadraticcurveto{\pgfqpoint{4.956975in}{9.360454in}}{\pgfqpoint{4.948931in}{9.363056in}}%
\pgfpathlineto{\pgfqpoint{4.948931in}{9.377235in}}%
\pgfpathquadraticcurveto{\pgfqpoint{4.956546in}{9.373273in}}{\pgfqpoint{4.963922in}{9.371291in}}%
\pgfpathquadraticcurveto{\pgfqpoint{4.971298in}{9.369334in}}{\pgfqpoint{4.978555in}{9.369334in}}%
\pgfpathquadraticcurveto{\pgfqpoint{4.988247in}{9.369334in}}{\pgfqpoint{4.993451in}{9.372652in}}%
\pgfpathquadraticcurveto{\pgfqpoint{4.998679in}{9.375970in}}{\pgfqpoint{4.998679in}{9.382010in}}%
\pgfpathquadraticcurveto{\pgfqpoint{4.998679in}{9.387595in}}{\pgfqpoint{4.994907in}{9.390579in}}%
\pgfpathquadraticcurveto{\pgfqpoint{4.991159in}{9.393563in}}{\pgfqpoint{4.978388in}{9.396332in}}%
\pgfpathlineto{\pgfqpoint{4.973614in}{9.397454in}}%
\pgfpathquadraticcurveto{\pgfqpoint{4.960866in}{9.400128in}}{\pgfqpoint{4.955185in}{9.405690in}}%
\pgfpathquadraticcurveto{\pgfqpoint{4.949527in}{9.411252in}}{\pgfqpoint{4.949527in}{9.420944in}}%
\pgfpathquadraticcurveto{\pgfqpoint{4.949527in}{9.432737in}}{\pgfqpoint{4.957882in}{9.439134in}}%
\pgfpathquadraticcurveto{\pgfqpoint{4.966237in}{9.445556in}}{\pgfqpoint{4.981611in}{9.445556in}}%
\pgfpathquadraticcurveto{\pgfqpoint{4.989202in}{9.445556in}}{\pgfqpoint{4.995910in}{9.444434in}}%
\pgfpathquadraticcurveto{\pgfqpoint{5.002641in}{9.443336in}}{\pgfqpoint{5.008299in}{9.441092in}}%
\pgfpathlineto{\pgfqpoint{5.008299in}{9.441092in}}%
\pgfpathclose%
\pgfpathmoveto{\pgfqpoint{5.106101in}{9.405213in}}%
\pgfpathlineto{\pgfqpoint{5.106101in}{9.398505in}}%
\pgfpathlineto{\pgfqpoint{5.042984in}{9.398505in}}%
\pgfpathquadraticcurveto{\pgfqpoint{5.043891in}{9.384325in}}{\pgfqpoint{5.051530in}{9.376901in}}%
\pgfpathquadraticcurveto{\pgfqpoint{5.059169in}{9.369477in}}{\pgfqpoint{5.072824in}{9.369477in}}%
\pgfpathquadraticcurveto{\pgfqpoint{5.080725in}{9.369477in}}{\pgfqpoint{5.088149in}{9.371411in}}%
\pgfpathquadraticcurveto{\pgfqpoint{5.095573in}{9.373344in}}{\pgfqpoint{5.102902in}{9.377235in}}%
\pgfpathlineto{\pgfqpoint{5.102902in}{9.364249in}}%
\pgfpathquadraticcurveto{\pgfqpoint{5.095502in}{9.361122in}}{\pgfqpoint{5.087743in}{9.359475in}}%
\pgfpathquadraticcurveto{\pgfqpoint{5.079985in}{9.357828in}}{\pgfqpoint{5.072012in}{9.357828in}}%
\pgfpathquadraticcurveto{\pgfqpoint{5.052008in}{9.357828in}}{\pgfqpoint{5.040335in}{9.369453in}}%
\pgfpathquadraticcurveto{\pgfqpoint{5.028661in}{9.381102in}}{\pgfqpoint{5.028661in}{9.400964in}}%
\pgfpathquadraticcurveto{\pgfqpoint{5.028661in}{9.421469in}}{\pgfqpoint{5.039738in}{9.433500in}}%
\pgfpathquadraticcurveto{\pgfqpoint{5.050814in}{9.445556in}}{\pgfqpoint{5.069625in}{9.445556in}}%
\pgfpathquadraticcurveto{\pgfqpoint{5.086478in}{9.445556in}}{\pgfqpoint{5.096289in}{9.434694in}}%
\pgfpathquadraticcurveto{\pgfqpoint{5.106101in}{9.423856in}}{\pgfqpoint{5.106101in}{9.405213in}}%
\pgfpathlineto{\pgfqpoint{5.106101in}{9.405213in}}%
\pgfpathclose%
\pgfpathmoveto{\pgfqpoint{5.092375in}{9.409247in}}%
\pgfpathquadraticcurveto{\pgfqpoint{5.092231in}{9.420490in}}{\pgfqpoint{5.086072in}{9.427198in}}%
\pgfpathquadraticcurveto{\pgfqpoint{5.079914in}{9.433930in}}{\pgfqpoint{5.069768in}{9.433930in}}%
\pgfpathquadraticcurveto{\pgfqpoint{5.058286in}{9.433930in}}{\pgfqpoint{5.051387in}{9.427437in}}%
\pgfpathquadraticcurveto{\pgfqpoint{5.044488in}{9.420944in}}{\pgfqpoint{5.043438in}{9.409151in}}%
\pgfpathlineto{\pgfqpoint{5.092375in}{9.409247in}}%
\pgfpathlineto{\pgfqpoint{5.092375in}{9.409247in}}%
\pgfpathclose%
\pgfpathmoveto{\pgfqpoint{5.179124in}{9.422520in}}%
\pgfpathlineto{\pgfqpoint{5.194879in}{9.422520in}}%
\pgfpathlineto{\pgfqpoint{5.194879in}{9.403566in}}%
\pgfpathlineto{\pgfqpoint{5.179124in}{9.403566in}}%
\pgfpathlineto{\pgfqpoint{5.179124in}{9.422520in}}%
\pgfpathlineto{\pgfqpoint{5.179124in}{9.422520in}}%
\pgfpathclose%
\pgfpathmoveto{\pgfqpoint{5.358303in}{9.462791in}}%
\pgfpathlineto{\pgfqpoint{5.358303in}{9.446916in}}%
\pgfpathquadraticcurveto{\pgfqpoint{5.350688in}{9.454006in}}{\pgfqpoint{5.342071in}{9.457491in}}%
\pgfpathquadraticcurveto{\pgfqpoint{5.333453in}{9.461000in}}{\pgfqpoint{5.323761in}{9.461000in}}%
\pgfpathquadraticcurveto{\pgfqpoint{5.304664in}{9.461000in}}{\pgfqpoint{5.294519in}{9.449327in}}%
\pgfpathquadraticcurveto{\pgfqpoint{5.284373in}{9.437654in}}{\pgfqpoint{5.284373in}{9.415573in}}%
\pgfpathquadraticcurveto{\pgfqpoint{5.284373in}{9.393563in}}{\pgfqpoint{5.294519in}{9.381890in}}%
\pgfpathquadraticcurveto{\pgfqpoint{5.304664in}{9.370217in}}{\pgfqpoint{5.323761in}{9.370217in}}%
\pgfpathquadraticcurveto{\pgfqpoint{5.333453in}{9.370217in}}{\pgfqpoint{5.342071in}{9.373726in}}%
\pgfpathquadraticcurveto{\pgfqpoint{5.350688in}{9.377235in}}{\pgfqpoint{5.358303in}{9.384325in}}%
\pgfpathlineto{\pgfqpoint{5.358303in}{9.368570in}}%
\pgfpathquadraticcurveto{\pgfqpoint{5.350402in}{9.363199in}}{\pgfqpoint{5.341546in}{9.360501in}}%
\pgfpathquadraticcurveto{\pgfqpoint{5.332713in}{9.357828in}}{\pgfqpoint{5.322878in}{9.357828in}}%
\pgfpathquadraticcurveto{\pgfqpoint{5.297574in}{9.357828in}}{\pgfqpoint{5.283013in}{9.373296in}}%
\pgfpathquadraticcurveto{\pgfqpoint{5.268475in}{9.388789in}}{\pgfqpoint{5.268475in}{9.415573in}}%
\pgfpathquadraticcurveto{\pgfqpoint{5.268475in}{9.442428in}}{\pgfqpoint{5.283013in}{9.457897in}}%
\pgfpathquadraticcurveto{\pgfqpoint{5.297574in}{9.473390in}}{\pgfqpoint{5.322878in}{9.473390in}}%
\pgfpathquadraticcurveto{\pgfqpoint{5.332856in}{9.473390in}}{\pgfqpoint{5.341689in}{9.470740in}}%
\pgfpathquadraticcurveto{\pgfqpoint{5.350545in}{9.468090in}}{\pgfqpoint{5.358303in}{9.462791in}}%
\pgfpathlineto{\pgfqpoint{5.358303in}{9.462791in}}%
\pgfpathclose%
\pgfpathmoveto{\pgfqpoint{5.429393in}{9.430731in}}%
\pgfpathquadraticcurveto{\pgfqpoint{5.427077in}{9.432068in}}{\pgfqpoint{5.424356in}{9.432689in}}%
\pgfpathquadraticcurveto{\pgfqpoint{5.421635in}{9.433333in}}{\pgfqpoint{5.418364in}{9.433333in}}%
\pgfpathquadraticcurveto{\pgfqpoint{5.406715in}{9.433333in}}{\pgfqpoint{5.400484in}{9.425766in}}%
\pgfpathquadraticcurveto{\pgfqpoint{5.394254in}{9.418199in}}{\pgfqpoint{5.394254in}{9.404019in}}%
\pgfpathlineto{\pgfqpoint{5.394254in}{9.360000in}}%
\pgfpathlineto{\pgfqpoint{5.380456in}{9.360000in}}%
\pgfpathlineto{\pgfqpoint{5.380456in}{9.443550in}}%
\pgfpathlineto{\pgfqpoint{5.394254in}{9.443550in}}%
\pgfpathlineto{\pgfqpoint{5.394254in}{9.430564in}}%
\pgfpathquadraticcurveto{\pgfqpoint{5.398599in}{9.438179in}}{\pgfqpoint{5.405521in}{9.441855in}}%
\pgfpathquadraticcurveto{\pgfqpoint{5.412468in}{9.445556in}}{\pgfqpoint{5.422398in}{9.445556in}}%
\pgfpathquadraticcurveto{\pgfqpoint{5.423807in}{9.445556in}}{\pgfqpoint{5.425526in}{9.445365in}}%
\pgfpathquadraticcurveto{\pgfqpoint{5.427244in}{9.445197in}}{\pgfqpoint{5.429321in}{9.444816in}}%
\pgfpathlineto{\pgfqpoint{5.429393in}{9.430731in}}%
\pgfpathlineto{\pgfqpoint{5.429393in}{9.430731in}}%
\pgfpathclose%
\pgfpathmoveto{\pgfqpoint{5.472791in}{9.433930in}}%
\pgfpathquadraticcurveto{\pgfqpoint{5.461763in}{9.433930in}}{\pgfqpoint{5.455341in}{9.425313in}}%
\pgfpathquadraticcurveto{\pgfqpoint{5.448920in}{9.416695in}}{\pgfqpoint{5.448920in}{9.401704in}}%
\pgfpathquadraticcurveto{\pgfqpoint{5.448920in}{9.386712in}}{\pgfqpoint{5.455293in}{9.378095in}}%
\pgfpathquadraticcurveto{\pgfqpoint{5.461691in}{9.369477in}}{\pgfqpoint{5.472791in}{9.369477in}}%
\pgfpathquadraticcurveto{\pgfqpoint{5.483772in}{9.369477in}}{\pgfqpoint{5.490170in}{9.378118in}}%
\pgfpathquadraticcurveto{\pgfqpoint{5.496591in}{9.386784in}}{\pgfqpoint{5.496591in}{9.401704in}}%
\pgfpathquadraticcurveto{\pgfqpoint{5.496591in}{9.416552in}}{\pgfqpoint{5.490170in}{9.425241in}}%
\pgfpathquadraticcurveto{\pgfqpoint{5.483772in}{9.433930in}}{\pgfqpoint{5.472791in}{9.433930in}}%
\pgfpathlineto{\pgfqpoint{5.472791in}{9.433930in}}%
\pgfpathclose%
\pgfpathmoveto{\pgfqpoint{5.472791in}{9.445556in}}%
\pgfpathquadraticcurveto{\pgfqpoint{5.490695in}{9.445556in}}{\pgfqpoint{5.500912in}{9.433906in}}%
\pgfpathquadraticcurveto{\pgfqpoint{5.511153in}{9.422281in}}{\pgfqpoint{5.511153in}{9.401704in}}%
\pgfpathquadraticcurveto{\pgfqpoint{5.511153in}{9.381198in}}{\pgfqpoint{5.500912in}{9.369501in}}%
\pgfpathquadraticcurveto{\pgfqpoint{5.490695in}{9.357828in}}{\pgfqpoint{5.472791in}{9.357828in}}%
\pgfpathquadraticcurveto{\pgfqpoint{5.454816in}{9.357828in}}{\pgfqpoint{5.444623in}{9.369501in}}%
\pgfpathquadraticcurveto{\pgfqpoint{5.434454in}{9.381198in}}{\pgfqpoint{5.434454in}{9.401704in}}%
\pgfpathquadraticcurveto{\pgfqpoint{5.434454in}{9.422281in}}{\pgfqpoint{5.444623in}{9.433906in}}%
\pgfpathquadraticcurveto{\pgfqpoint{5.454816in}{9.445556in}}{\pgfqpoint{5.472791in}{9.445556in}}%
\pgfpathlineto{\pgfqpoint{5.472791in}{9.445556in}}%
\pgfpathclose%
\pgfpathmoveto{\pgfqpoint{5.587160in}{9.441092in}}%
\pgfpathlineto{\pgfqpoint{5.587160in}{9.428105in}}%
\pgfpathquadraticcurveto{\pgfqpoint{5.581359in}{9.431089in}}{\pgfqpoint{5.575081in}{9.432569in}}%
\pgfpathquadraticcurveto{\pgfqpoint{5.568826in}{9.434073in}}{\pgfqpoint{5.562095in}{9.434073in}}%
\pgfpathquadraticcurveto{\pgfqpoint{5.551878in}{9.434073in}}{\pgfqpoint{5.546769in}{9.430946in}}%
\pgfpathquadraticcurveto{\pgfqpoint{5.541661in}{9.427819in}}{\pgfqpoint{5.541661in}{9.421541in}}%
\pgfpathquadraticcurveto{\pgfqpoint{5.541661in}{9.416766in}}{\pgfqpoint{5.545313in}{9.414045in}}%
\pgfpathquadraticcurveto{\pgfqpoint{5.548965in}{9.411324in}}{\pgfqpoint{5.560018in}{9.408865in}}%
\pgfpathlineto{\pgfqpoint{5.564720in}{9.407815in}}%
\pgfpathquadraticcurveto{\pgfqpoint{5.579330in}{9.404687in}}{\pgfqpoint{5.585489in}{9.398982in}}%
\pgfpathquadraticcurveto{\pgfqpoint{5.591648in}{9.393277in}}{\pgfqpoint{5.591648in}{9.383060in}}%
\pgfpathquadraticcurveto{\pgfqpoint{5.591648in}{9.371411in}}{\pgfqpoint{5.582433in}{9.364607in}}%
\pgfpathquadraticcurveto{\pgfqpoint{5.573219in}{9.357828in}}{\pgfqpoint{5.557105in}{9.357828in}}%
\pgfpathquadraticcurveto{\pgfqpoint{5.550398in}{9.357828in}}{\pgfqpoint{5.543117in}{9.359141in}}%
\pgfpathquadraticcurveto{\pgfqpoint{5.535836in}{9.360454in}}{\pgfqpoint{5.527791in}{9.363056in}}%
\pgfpathlineto{\pgfqpoint{5.527791in}{9.377235in}}%
\pgfpathquadraticcurveto{\pgfqpoint{5.535406in}{9.373273in}}{\pgfqpoint{5.542783in}{9.371291in}}%
\pgfpathquadraticcurveto{\pgfqpoint{5.550159in}{9.369334in}}{\pgfqpoint{5.557416in}{9.369334in}}%
\pgfpathquadraticcurveto{\pgfqpoint{5.567108in}{9.369334in}}{\pgfqpoint{5.572312in}{9.372652in}}%
\pgfpathquadraticcurveto{\pgfqpoint{5.577539in}{9.375970in}}{\pgfqpoint{5.577539in}{9.382010in}}%
\pgfpathquadraticcurveto{\pgfqpoint{5.577539in}{9.387595in}}{\pgfqpoint{5.573768in}{9.390579in}}%
\pgfpathquadraticcurveto{\pgfqpoint{5.570020in}{9.393563in}}{\pgfqpoint{5.557249in}{9.396332in}}%
\pgfpathlineto{\pgfqpoint{5.552474in}{9.397454in}}%
\pgfpathquadraticcurveto{\pgfqpoint{5.539727in}{9.400128in}}{\pgfqpoint{5.534046in}{9.405690in}}%
\pgfpathquadraticcurveto{\pgfqpoint{5.528388in}{9.411252in}}{\pgfqpoint{5.528388in}{9.420944in}}%
\pgfpathquadraticcurveto{\pgfqpoint{5.528388in}{9.432737in}}{\pgfqpoint{5.536743in}{9.439134in}}%
\pgfpathquadraticcurveto{\pgfqpoint{5.545098in}{9.445556in}}{\pgfqpoint{5.560471in}{9.445556in}}%
\pgfpathquadraticcurveto{\pgfqpoint{5.568062in}{9.445556in}}{\pgfqpoint{5.574770in}{9.444434in}}%
\pgfpathquadraticcurveto{\pgfqpoint{5.581502in}{9.443336in}}{\pgfqpoint{5.587160in}{9.441092in}}%
\pgfpathlineto{\pgfqpoint{5.587160in}{9.441092in}}%
\pgfpathclose%
\pgfpathmoveto{\pgfqpoint{5.666747in}{9.441092in}}%
\pgfpathlineto{\pgfqpoint{5.666747in}{9.428105in}}%
\pgfpathquadraticcurveto{\pgfqpoint{5.660947in}{9.431089in}}{\pgfqpoint{5.654668in}{9.432569in}}%
\pgfpathquadraticcurveto{\pgfqpoint{5.648414in}{9.434073in}}{\pgfqpoint{5.641682in}{9.434073in}}%
\pgfpathquadraticcurveto{\pgfqpoint{5.631465in}{9.434073in}}{\pgfqpoint{5.626357in}{9.430946in}}%
\pgfpathquadraticcurveto{\pgfqpoint{5.621248in}{9.427819in}}{\pgfqpoint{5.621248in}{9.421541in}}%
\pgfpathquadraticcurveto{\pgfqpoint{5.621248in}{9.416766in}}{\pgfqpoint{5.624901in}{9.414045in}}%
\pgfpathquadraticcurveto{\pgfqpoint{5.628553in}{9.411324in}}{\pgfqpoint{5.639605in}{9.408865in}}%
\pgfpathlineto{\pgfqpoint{5.644308in}{9.407815in}}%
\pgfpathquadraticcurveto{\pgfqpoint{5.658918in}{9.404687in}}{\pgfqpoint{5.665076in}{9.398982in}}%
\pgfpathquadraticcurveto{\pgfqpoint{5.671235in}{9.393277in}}{\pgfqpoint{5.671235in}{9.383060in}}%
\pgfpathquadraticcurveto{\pgfqpoint{5.671235in}{9.371411in}}{\pgfqpoint{5.662021in}{9.364607in}}%
\pgfpathquadraticcurveto{\pgfqpoint{5.652806in}{9.357828in}}{\pgfqpoint{5.636693in}{9.357828in}}%
\pgfpathquadraticcurveto{\pgfqpoint{5.629985in}{9.357828in}}{\pgfqpoint{5.622704in}{9.359141in}}%
\pgfpathquadraticcurveto{\pgfqpoint{5.615424in}{9.360454in}}{\pgfqpoint{5.607379in}{9.363056in}}%
\pgfpathlineto{\pgfqpoint{5.607379in}{9.377235in}}%
\pgfpathquadraticcurveto{\pgfqpoint{5.614994in}{9.373273in}}{\pgfqpoint{5.622370in}{9.371291in}}%
\pgfpathquadraticcurveto{\pgfqpoint{5.629747in}{9.369334in}}{\pgfqpoint{5.637003in}{9.369334in}}%
\pgfpathquadraticcurveto{\pgfqpoint{5.646695in}{9.369334in}}{\pgfqpoint{5.651899in}{9.372652in}}%
\pgfpathquadraticcurveto{\pgfqpoint{5.657127in}{9.375970in}}{\pgfqpoint{5.657127in}{9.382010in}}%
\pgfpathquadraticcurveto{\pgfqpoint{5.657127in}{9.387595in}}{\pgfqpoint{5.653355in}{9.390579in}}%
\pgfpathquadraticcurveto{\pgfqpoint{5.649608in}{9.393563in}}{\pgfqpoint{5.636836in}{9.396332in}}%
\pgfpathlineto{\pgfqpoint{5.632062in}{9.397454in}}%
\pgfpathquadraticcurveto{\pgfqpoint{5.619315in}{9.400128in}}{\pgfqpoint{5.613633in}{9.405690in}}%
\pgfpathquadraticcurveto{\pgfqpoint{5.607976in}{9.411252in}}{\pgfqpoint{5.607976in}{9.420944in}}%
\pgfpathquadraticcurveto{\pgfqpoint{5.607976in}{9.432737in}}{\pgfqpoint{5.616331in}{9.439134in}}%
\pgfpathquadraticcurveto{\pgfqpoint{5.624686in}{9.445556in}}{\pgfqpoint{5.640059in}{9.445556in}}%
\pgfpathquadraticcurveto{\pgfqpoint{5.647650in}{9.445556in}}{\pgfqpoint{5.654358in}{9.444434in}}%
\pgfpathquadraticcurveto{\pgfqpoint{5.661090in}{9.443336in}}{\pgfqpoint{5.666747in}{9.441092in}}%
\pgfpathlineto{\pgfqpoint{5.666747in}{9.441092in}}%
\pgfpathclose%
\pgfpathmoveto{\pgfqpoint{5.696587in}{9.378954in}}%
\pgfpathlineto{\pgfqpoint{5.712318in}{9.378954in}}%
\pgfpathlineto{\pgfqpoint{5.712318in}{9.360000in}}%
\pgfpathlineto{\pgfqpoint{5.696587in}{9.360000in}}%
\pgfpathlineto{\pgfqpoint{5.696587in}{9.378954in}}%
\pgfpathlineto{\pgfqpoint{5.696587in}{9.378954in}}%
\pgfpathclose%
\pgfpathmoveto{\pgfqpoint{5.696587in}{9.438991in}}%
\pgfpathlineto{\pgfqpoint{5.712318in}{9.438991in}}%
\pgfpathlineto{\pgfqpoint{5.712318in}{9.420061in}}%
\pgfpathlineto{\pgfqpoint{5.696587in}{9.420061in}}%
\pgfpathlineto{\pgfqpoint{5.696587in}{9.438991in}}%
\pgfpathlineto{\pgfqpoint{5.696587in}{9.438991in}}%
\pgfpathclose%
\pgfpathmoveto{\pgfqpoint{5.835400in}{9.476087in}}%
\pgfpathlineto{\pgfqpoint{5.835400in}{9.464653in}}%
\pgfpathlineto{\pgfqpoint{5.822270in}{9.464653in}}%
\pgfpathquadraticcurveto{\pgfqpoint{5.814894in}{9.464653in}}{\pgfqpoint{5.812006in}{9.461669in}}%
\pgfpathquadraticcurveto{\pgfqpoint{5.809141in}{9.458685in}}{\pgfqpoint{5.809141in}{9.450927in}}%
\pgfpathlineto{\pgfqpoint{5.809141in}{9.443550in}}%
\pgfpathlineto{\pgfqpoint{5.831747in}{9.443550in}}%
\pgfpathlineto{\pgfqpoint{5.831747in}{9.432880in}}%
\pgfpathlineto{\pgfqpoint{5.809141in}{9.432880in}}%
\pgfpathlineto{\pgfqpoint{5.809141in}{9.360000in}}%
\pgfpathlineto{\pgfqpoint{5.795343in}{9.360000in}}%
\pgfpathlineto{\pgfqpoint{5.795343in}{9.432880in}}%
\pgfpathlineto{\pgfqpoint{5.782214in}{9.432880in}}%
\pgfpathlineto{\pgfqpoint{5.782214in}{9.443550in}}%
\pgfpathlineto{\pgfqpoint{5.795343in}{9.443550in}}%
\pgfpathlineto{\pgfqpoint{5.795343in}{9.449375in}}%
\pgfpathquadraticcurveto{\pgfqpoint{5.795343in}{9.463316in}}{\pgfqpoint{5.801836in}{9.469690in}}%
\pgfpathquadraticcurveto{\pgfqpoint{5.808329in}{9.476087in}}{\pgfqpoint{5.822414in}{9.476087in}}%
\pgfpathlineto{\pgfqpoint{5.835400in}{9.476087in}}%
\pgfpathlineto{\pgfqpoint{5.835400in}{9.476087in}}%
\pgfpathclose%
\pgfpathmoveto{\pgfqpoint{5.895293in}{9.430731in}}%
\pgfpathquadraticcurveto{\pgfqpoint{5.892978in}{9.432068in}}{\pgfqpoint{5.890257in}{9.432689in}}%
\pgfpathquadraticcurveto{\pgfqpoint{5.887535in}{9.433333in}}{\pgfqpoint{5.884265in}{9.433333in}}%
\pgfpathquadraticcurveto{\pgfqpoint{5.872615in}{9.433333in}}{\pgfqpoint{5.866385in}{9.425766in}}%
\pgfpathquadraticcurveto{\pgfqpoint{5.860155in}{9.418199in}}{\pgfqpoint{5.860155in}{9.404019in}}%
\pgfpathlineto{\pgfqpoint{5.860155in}{9.360000in}}%
\pgfpathlineto{\pgfqpoint{5.846357in}{9.360000in}}%
\pgfpathlineto{\pgfqpoint{5.846357in}{9.443550in}}%
\pgfpathlineto{\pgfqpoint{5.860155in}{9.443550in}}%
\pgfpathlineto{\pgfqpoint{5.860155in}{9.430564in}}%
\pgfpathquadraticcurveto{\pgfqpoint{5.864499in}{9.438179in}}{\pgfqpoint{5.871422in}{9.441855in}}%
\pgfpathquadraticcurveto{\pgfqpoint{5.878368in}{9.445556in}}{\pgfqpoint{5.888299in}{9.445556in}}%
\pgfpathquadraticcurveto{\pgfqpoint{5.889707in}{9.445556in}}{\pgfqpoint{5.891426in}{9.445365in}}%
\pgfpathquadraticcurveto{\pgfqpoint{5.893145in}{9.445197in}}{\pgfqpoint{5.895222in}{9.444816in}}%
\pgfpathlineto{\pgfqpoint{5.895293in}{9.430731in}}%
\pgfpathlineto{\pgfqpoint{5.895293in}{9.430731in}}%
\pgfpathclose%
\pgfpathmoveto{\pgfqpoint{5.947668in}{9.401990in}}%
\pgfpathquadraticcurveto{\pgfqpoint{5.931029in}{9.401990in}}{\pgfqpoint{5.924608in}{9.398194in}}%
\pgfpathquadraticcurveto{\pgfqpoint{5.918186in}{9.394399in}}{\pgfqpoint{5.918186in}{9.385208in}}%
\pgfpathquadraticcurveto{\pgfqpoint{5.918186in}{9.377904in}}{\pgfqpoint{5.923008in}{9.373607in}}%
\pgfpathquadraticcurveto{\pgfqpoint{5.927830in}{9.369334in}}{\pgfqpoint{5.936090in}{9.369334in}}%
\pgfpathquadraticcurveto{\pgfqpoint{5.947524in}{9.369334in}}{\pgfqpoint{5.954423in}{9.377426in}}%
\pgfpathquadraticcurveto{\pgfqpoint{5.961322in}{9.385519in}}{\pgfqpoint{5.961322in}{9.398934in}}%
\pgfpathlineto{\pgfqpoint{5.961322in}{9.401990in}}%
\pgfpathlineto{\pgfqpoint{5.947668in}{9.401990in}}%
\pgfpathlineto{\pgfqpoint{5.947668in}{9.401990in}}%
\pgfpathclose%
\pgfpathmoveto{\pgfqpoint{5.975048in}{9.407671in}}%
\pgfpathlineto{\pgfqpoint{5.975048in}{9.360000in}}%
\pgfpathlineto{\pgfqpoint{5.961322in}{9.360000in}}%
\pgfpathlineto{\pgfqpoint{5.961322in}{9.372676in}}%
\pgfpathquadraticcurveto{\pgfqpoint{5.956619in}{9.365085in}}{\pgfqpoint{5.949601in}{9.361456in}}%
\pgfpathquadraticcurveto{\pgfqpoint{5.942583in}{9.357828in}}{\pgfqpoint{5.932437in}{9.357828in}}%
\pgfpathquadraticcurveto{\pgfqpoint{5.919618in}{9.357828in}}{\pgfqpoint{5.912027in}{9.365037in}}%
\pgfpathquadraticcurveto{\pgfqpoint{5.904460in}{9.372246in}}{\pgfqpoint{5.904460in}{9.384325in}}%
\pgfpathquadraticcurveto{\pgfqpoint{5.904460in}{9.398409in}}{\pgfqpoint{5.913889in}{9.405571in}}%
\pgfpathquadraticcurveto{\pgfqpoint{5.923342in}{9.412732in}}{\pgfqpoint{5.942058in}{9.412732in}}%
\pgfpathlineto{\pgfqpoint{5.961322in}{9.412732in}}%
\pgfpathlineto{\pgfqpoint{5.961322in}{9.414093in}}%
\pgfpathquadraticcurveto{\pgfqpoint{5.961322in}{9.423570in}}{\pgfqpoint{5.955092in}{9.428750in}}%
\pgfpathquadraticcurveto{\pgfqpoint{5.948861in}{9.433930in}}{\pgfqpoint{5.937594in}{9.433930in}}%
\pgfpathquadraticcurveto{\pgfqpoint{5.930432in}{9.433930in}}{\pgfqpoint{5.923629in}{9.432211in}}%
\pgfpathquadraticcurveto{\pgfqpoint{5.916849in}{9.430493in}}{\pgfqpoint{5.910595in}{9.427055in}}%
\pgfpathlineto{\pgfqpoint{5.910595in}{9.439755in}}%
\pgfpathquadraticcurveto{\pgfqpoint{5.918115in}{9.442667in}}{\pgfqpoint{5.925204in}{9.444099in}}%
\pgfpathquadraticcurveto{\pgfqpoint{5.932294in}{9.445556in}}{\pgfqpoint{5.939002in}{9.445556in}}%
\pgfpathquadraticcurveto{\pgfqpoint{5.957145in}{9.445556in}}{\pgfqpoint{5.966096in}{9.436150in}}%
\pgfpathquadraticcurveto{\pgfqpoint{5.975048in}{9.426769in}}{\pgfqpoint{5.975048in}{9.407671in}}%
\pgfpathlineto{\pgfqpoint{5.975048in}{9.407671in}}%
\pgfpathclose%
\pgfpathmoveto{\pgfqpoint{6.058288in}{9.402754in}}%
\pgfpathquadraticcurveto{\pgfqpoint{6.058288in}{9.417674in}}{\pgfqpoint{6.052129in}{9.425862in}}%
\pgfpathquadraticcurveto{\pgfqpoint{6.045994in}{9.434073in}}{\pgfqpoint{6.034870in}{9.434073in}}%
\pgfpathquadraticcurveto{\pgfqpoint{6.023842in}{9.434073in}}{\pgfqpoint{6.017683in}{9.425862in}}%
\pgfpathquadraticcurveto{\pgfqpoint{6.011524in}{9.417674in}}{\pgfqpoint{6.011524in}{9.402754in}}%
\pgfpathquadraticcurveto{\pgfqpoint{6.011524in}{9.387906in}}{\pgfqpoint{6.017683in}{9.379694in}}%
\pgfpathquadraticcurveto{\pgfqpoint{6.023842in}{9.371482in}}{\pgfqpoint{6.034870in}{9.371482in}}%
\pgfpathquadraticcurveto{\pgfqpoint{6.045994in}{9.371482in}}{\pgfqpoint{6.052129in}{9.379694in}}%
\pgfpathquadraticcurveto{\pgfqpoint{6.058288in}{9.387906in}}{\pgfqpoint{6.058288in}{9.402754in}}%
\pgfpathlineto{\pgfqpoint{6.058288in}{9.402754in}}%
\pgfpathclose%
\pgfpathmoveto{\pgfqpoint{6.072014in}{9.370360in}}%
\pgfpathquadraticcurveto{\pgfqpoint{6.072014in}{9.349043in}}{\pgfqpoint{6.062537in}{9.338635in}}%
\pgfpathquadraticcurveto{\pgfqpoint{6.053084in}{9.328227in}}{\pgfqpoint{6.033533in}{9.328227in}}%
\pgfpathquadraticcurveto{\pgfqpoint{6.026300in}{9.328227in}}{\pgfqpoint{6.019879in}{9.329301in}}%
\pgfpathquadraticcurveto{\pgfqpoint{6.013457in}{9.330375in}}{\pgfqpoint{6.007418in}{9.332619in}}%
\pgfpathlineto{\pgfqpoint{6.007418in}{9.345964in}}%
\pgfpathquadraticcurveto{\pgfqpoint{6.013457in}{9.342693in}}{\pgfqpoint{6.019354in}{9.341141in}}%
\pgfpathquadraticcurveto{\pgfqpoint{6.025250in}{9.339566in}}{\pgfqpoint{6.031361in}{9.339566in}}%
\pgfpathquadraticcurveto{\pgfqpoint{6.044872in}{9.339566in}}{\pgfqpoint{6.051580in}{9.346608in}}%
\pgfpathquadraticcurveto{\pgfqpoint{6.058288in}{9.353650in}}{\pgfqpoint{6.058288in}{9.367901in}}%
\pgfpathlineto{\pgfqpoint{6.058288in}{9.374705in}}%
\pgfpathquadraticcurveto{\pgfqpoint{6.054039in}{9.367305in}}{\pgfqpoint{6.047403in}{9.363652in}}%
\pgfpathquadraticcurveto{\pgfqpoint{6.040766in}{9.360000in}}{\pgfqpoint{6.031504in}{9.360000in}}%
\pgfpathquadraticcurveto{\pgfqpoint{6.016155in}{9.360000in}}{\pgfqpoint{6.006750in}{9.371697in}}%
\pgfpathquadraticcurveto{\pgfqpoint{5.997344in}{9.383418in}}{\pgfqpoint{5.997344in}{9.402754in}}%
\pgfpathquadraticcurveto{\pgfqpoint{5.997344in}{9.422138in}}{\pgfqpoint{6.006750in}{9.433835in}}%
\pgfpathquadraticcurveto{\pgfqpoint{6.016155in}{9.445556in}}{\pgfqpoint{6.031504in}{9.445556in}}%
\pgfpathquadraticcurveto{\pgfqpoint{6.040766in}{9.445556in}}{\pgfqpoint{6.047403in}{9.441903in}}%
\pgfpathquadraticcurveto{\pgfqpoint{6.054039in}{9.438251in}}{\pgfqpoint{6.058288in}{9.430875in}}%
\pgfpathlineto{\pgfqpoint{6.058288in}{9.443550in}}%
\pgfpathlineto{\pgfqpoint{6.072014in}{9.443550in}}%
\pgfpathlineto{\pgfqpoint{6.072014in}{9.370360in}}%
\pgfpathlineto{\pgfqpoint{6.072014in}{9.370360in}}%
\pgfpathclose%
\pgfpathmoveto{\pgfqpoint{6.100302in}{9.443550in}}%
\pgfpathlineto{\pgfqpoint{6.114028in}{9.443550in}}%
\pgfpathlineto{\pgfqpoint{6.114028in}{9.360000in}}%
\pgfpathlineto{\pgfqpoint{6.100302in}{9.360000in}}%
\pgfpathlineto{\pgfqpoint{6.100302in}{9.443550in}}%
\pgfpathlineto{\pgfqpoint{6.100302in}{9.443550in}}%
\pgfpathclose%
\pgfpathmoveto{\pgfqpoint{6.100302in}{9.476087in}}%
\pgfpathlineto{\pgfqpoint{6.114028in}{9.476087in}}%
\pgfpathlineto{\pgfqpoint{6.114028in}{9.458685in}}%
\pgfpathlineto{\pgfqpoint{6.100302in}{9.458685in}}%
\pgfpathlineto{\pgfqpoint{6.100302in}{9.476087in}}%
\pgfpathlineto{\pgfqpoint{6.100302in}{9.476087in}}%
\pgfpathclose%
\pgfpathmoveto{\pgfqpoint{6.142746in}{9.476087in}}%
\pgfpathlineto{\pgfqpoint{6.156472in}{9.476087in}}%
\pgfpathlineto{\pgfqpoint{6.156472in}{9.360000in}}%
\pgfpathlineto{\pgfqpoint{6.142746in}{9.360000in}}%
\pgfpathlineto{\pgfqpoint{6.142746in}{9.476087in}}%
\pgfpathlineto{\pgfqpoint{6.142746in}{9.476087in}}%
\pgfpathclose%
\pgfpathmoveto{\pgfqpoint{6.256661in}{9.405213in}}%
\pgfpathlineto{\pgfqpoint{6.256661in}{9.398505in}}%
\pgfpathlineto{\pgfqpoint{6.193544in}{9.398505in}}%
\pgfpathquadraticcurveto{\pgfqpoint{6.194451in}{9.384325in}}{\pgfqpoint{6.202090in}{9.376901in}}%
\pgfpathquadraticcurveto{\pgfqpoint{6.209729in}{9.369477in}}{\pgfqpoint{6.223384in}{9.369477in}}%
\pgfpathquadraticcurveto{\pgfqpoint{6.231285in}{9.369477in}}{\pgfqpoint{6.238709in}{9.371411in}}%
\pgfpathquadraticcurveto{\pgfqpoint{6.246133in}{9.373344in}}{\pgfqpoint{6.253462in}{9.377235in}}%
\pgfpathlineto{\pgfqpoint{6.253462in}{9.364249in}}%
\pgfpathquadraticcurveto{\pgfqpoint{6.246062in}{9.361122in}}{\pgfqpoint{6.238303in}{9.359475in}}%
\pgfpathquadraticcurveto{\pgfqpoint{6.230545in}{9.357828in}}{\pgfqpoint{6.222572in}{9.357828in}}%
\pgfpathquadraticcurveto{\pgfqpoint{6.202568in}{9.357828in}}{\pgfqpoint{6.190895in}{9.369453in}}%
\pgfpathquadraticcurveto{\pgfqpoint{6.179221in}{9.381102in}}{\pgfqpoint{6.179221in}{9.400964in}}%
\pgfpathquadraticcurveto{\pgfqpoint{6.179221in}{9.421469in}}{\pgfqpoint{6.190298in}{9.433500in}}%
\pgfpathquadraticcurveto{\pgfqpoint{6.201374in}{9.445556in}}{\pgfqpoint{6.220185in}{9.445556in}}%
\pgfpathquadraticcurveto{\pgfqpoint{6.237038in}{9.445556in}}{\pgfqpoint{6.246849in}{9.434694in}}%
\pgfpathquadraticcurveto{\pgfqpoint{6.256661in}{9.423856in}}{\pgfqpoint{6.256661in}{9.405213in}}%
\pgfpathlineto{\pgfqpoint{6.256661in}{9.405213in}}%
\pgfpathclose%
\pgfpathmoveto{\pgfqpoint{6.242934in}{9.409247in}}%
\pgfpathquadraticcurveto{\pgfqpoint{6.242791in}{9.420490in}}{\pgfqpoint{6.236632in}{9.427198in}}%
\pgfpathquadraticcurveto{\pgfqpoint{6.230474in}{9.433930in}}{\pgfqpoint{6.220328in}{9.433930in}}%
\pgfpathquadraticcurveto{\pgfqpoint{6.208846in}{9.433930in}}{\pgfqpoint{6.201947in}{9.427437in}}%
\pgfpathquadraticcurveto{\pgfqpoint{6.195048in}{9.420944in}}{\pgfqpoint{6.193998in}{9.409151in}}%
\pgfpathlineto{\pgfqpoint{6.242934in}{9.409247in}}%
\pgfpathlineto{\pgfqpoint{6.242934in}{9.409247in}}%
\pgfpathclose%
\pgfpathmoveto{\pgfqpoint{6.327750in}{9.443550in}}%
\pgfpathlineto{\pgfqpoint{6.341476in}{9.443550in}}%
\pgfpathlineto{\pgfqpoint{6.341476in}{9.360000in}}%
\pgfpathlineto{\pgfqpoint{6.327750in}{9.360000in}}%
\pgfpathlineto{\pgfqpoint{6.327750in}{9.443550in}}%
\pgfpathlineto{\pgfqpoint{6.327750in}{9.443550in}}%
\pgfpathclose%
\pgfpathmoveto{\pgfqpoint{6.327750in}{9.476087in}}%
\pgfpathlineto{\pgfqpoint{6.341476in}{9.476087in}}%
\pgfpathlineto{\pgfqpoint{6.341476in}{9.458685in}}%
\pgfpathlineto{\pgfqpoint{6.327750in}{9.458685in}}%
\pgfpathlineto{\pgfqpoint{6.327750in}{9.476087in}}%
\pgfpathlineto{\pgfqpoint{6.327750in}{9.476087in}}%
\pgfpathclose%
\pgfpathmoveto{\pgfqpoint{6.439660in}{9.410441in}}%
\pgfpathlineto{\pgfqpoint{6.439660in}{9.360000in}}%
\pgfpathlineto{\pgfqpoint{6.425934in}{9.360000in}}%
\pgfpathlineto{\pgfqpoint{6.425934in}{9.409987in}}%
\pgfpathquadraticcurveto{\pgfqpoint{6.425934in}{9.421851in}}{\pgfqpoint{6.421303in}{9.427724in}}%
\pgfpathquadraticcurveto{\pgfqpoint{6.416671in}{9.433620in}}{\pgfqpoint{6.407433in}{9.433620in}}%
\pgfpathquadraticcurveto{\pgfqpoint{6.396309in}{9.433620in}}{\pgfqpoint{6.389888in}{9.426530in}}%
\pgfpathquadraticcurveto{\pgfqpoint{6.383466in}{9.419464in}}{\pgfqpoint{6.383466in}{9.407218in}}%
\pgfpathlineto{\pgfqpoint{6.383466in}{9.360000in}}%
\pgfpathlineto{\pgfqpoint{6.369668in}{9.360000in}}%
\pgfpathlineto{\pgfqpoint{6.369668in}{9.443550in}}%
\pgfpathlineto{\pgfqpoint{6.383466in}{9.443550in}}%
\pgfpathlineto{\pgfqpoint{6.383466in}{9.430564in}}%
\pgfpathquadraticcurveto{\pgfqpoint{6.388408in}{9.438108in}}{\pgfqpoint{6.395068in}{9.441832in}}%
\pgfpathquadraticcurveto{\pgfqpoint{6.401752in}{9.445556in}}{\pgfqpoint{6.410489in}{9.445556in}}%
\pgfpathquadraticcurveto{\pgfqpoint{6.424883in}{9.445556in}}{\pgfqpoint{6.432260in}{9.436651in}}%
\pgfpathquadraticcurveto{\pgfqpoint{6.439660in}{9.427747in}}{\pgfqpoint{6.439660in}{9.410441in}}%
\pgfpathlineto{\pgfqpoint{6.439660in}{9.410441in}}%
\pgfpathclose%
\pgfpathmoveto{\pgfqpoint{6.509317in}{9.476087in}}%
\pgfpathlineto{\pgfqpoint{6.509317in}{9.464653in}}%
\pgfpathlineto{\pgfqpoint{6.496187in}{9.464653in}}%
\pgfpathquadraticcurveto{\pgfqpoint{6.488811in}{9.464653in}}{\pgfqpoint{6.485923in}{9.461669in}}%
\pgfpathquadraticcurveto{\pgfqpoint{6.483058in}{9.458685in}}{\pgfqpoint{6.483058in}{9.450927in}}%
\pgfpathlineto{\pgfqpoint{6.483058in}{9.443550in}}%
\pgfpathlineto{\pgfqpoint{6.505664in}{9.443550in}}%
\pgfpathlineto{\pgfqpoint{6.505664in}{9.432880in}}%
\pgfpathlineto{\pgfqpoint{6.483058in}{9.432880in}}%
\pgfpathlineto{\pgfqpoint{6.483058in}{9.360000in}}%
\pgfpathlineto{\pgfqpoint{6.469260in}{9.360000in}}%
\pgfpathlineto{\pgfqpoint{6.469260in}{9.432880in}}%
\pgfpathlineto{\pgfqpoint{6.456131in}{9.432880in}}%
\pgfpathlineto{\pgfqpoint{6.456131in}{9.443550in}}%
\pgfpathlineto{\pgfqpoint{6.469260in}{9.443550in}}%
\pgfpathlineto{\pgfqpoint{6.469260in}{9.449375in}}%
\pgfpathquadraticcurveto{\pgfqpoint{6.469260in}{9.463316in}}{\pgfqpoint{6.475753in}{9.469690in}}%
\pgfpathquadraticcurveto{\pgfqpoint{6.482247in}{9.476087in}}{\pgfqpoint{6.496331in}{9.476087in}}%
\pgfpathlineto{\pgfqpoint{6.509317in}{9.476087in}}%
\pgfpathlineto{\pgfqpoint{6.509317in}{9.476087in}}%
\pgfpathclose%
\pgfpathmoveto{\pgfqpoint{6.592270in}{9.405213in}}%
\pgfpathlineto{\pgfqpoint{6.592270in}{9.398505in}}%
\pgfpathlineto{\pgfqpoint{6.529154in}{9.398505in}}%
\pgfpathquadraticcurveto{\pgfqpoint{6.530061in}{9.384325in}}{\pgfqpoint{6.537700in}{9.376901in}}%
\pgfpathquadraticcurveto{\pgfqpoint{6.545339in}{9.369477in}}{\pgfqpoint{6.558993in}{9.369477in}}%
\pgfpathquadraticcurveto{\pgfqpoint{6.566895in}{9.369477in}}{\pgfqpoint{6.574319in}{9.371411in}}%
\pgfpathquadraticcurveto{\pgfqpoint{6.581743in}{9.373344in}}{\pgfqpoint{6.589072in}{9.377235in}}%
\pgfpathlineto{\pgfqpoint{6.589072in}{9.364249in}}%
\pgfpathquadraticcurveto{\pgfqpoint{6.581671in}{9.361122in}}{\pgfqpoint{6.573913in}{9.359475in}}%
\pgfpathquadraticcurveto{\pgfqpoint{6.566155in}{9.357828in}}{\pgfqpoint{6.558182in}{9.357828in}}%
\pgfpathquadraticcurveto{\pgfqpoint{6.538178in}{9.357828in}}{\pgfqpoint{6.526504in}{9.369453in}}%
\pgfpathquadraticcurveto{\pgfqpoint{6.514831in}{9.381102in}}{\pgfqpoint{6.514831in}{9.400964in}}%
\pgfpathquadraticcurveto{\pgfqpoint{6.514831in}{9.421469in}}{\pgfqpoint{6.525908in}{9.433500in}}%
\pgfpathquadraticcurveto{\pgfqpoint{6.536984in}{9.445556in}}{\pgfqpoint{6.555795in}{9.445556in}}%
\pgfpathquadraticcurveto{\pgfqpoint{6.572648in}{9.445556in}}{\pgfqpoint{6.582459in}{9.434694in}}%
\pgfpathquadraticcurveto{\pgfqpoint{6.592270in}{9.423856in}}{\pgfqpoint{6.592270in}{9.405213in}}%
\pgfpathlineto{\pgfqpoint{6.592270in}{9.405213in}}%
\pgfpathclose%
\pgfpathmoveto{\pgfqpoint{6.578544in}{9.409247in}}%
\pgfpathquadraticcurveto{\pgfqpoint{6.578401in}{9.420490in}}{\pgfqpoint{6.572242in}{9.427198in}}%
\pgfpathquadraticcurveto{\pgfqpoint{6.566083in}{9.433930in}}{\pgfqpoint{6.555938in}{9.433930in}}%
\pgfpathquadraticcurveto{\pgfqpoint{6.544456in}{9.433930in}}{\pgfqpoint{6.537557in}{9.427437in}}%
\pgfpathquadraticcurveto{\pgfqpoint{6.530658in}{9.420944in}}{\pgfqpoint{6.529608in}{9.409151in}}%
\pgfpathlineto{\pgfqpoint{6.578544in}{9.409247in}}%
\pgfpathlineto{\pgfqpoint{6.578544in}{9.409247in}}%
\pgfpathclose%
\pgfpathmoveto{\pgfqpoint{6.663217in}{9.430731in}}%
\pgfpathquadraticcurveto{\pgfqpoint{6.660901in}{9.432068in}}{\pgfqpoint{6.658180in}{9.432689in}}%
\pgfpathquadraticcurveto{\pgfqpoint{6.655458in}{9.433333in}}{\pgfqpoint{6.652188in}{9.433333in}}%
\pgfpathquadraticcurveto{\pgfqpoint{6.640539in}{9.433333in}}{\pgfqpoint{6.634308in}{9.425766in}}%
\pgfpathquadraticcurveto{\pgfqpoint{6.628078in}{9.418199in}}{\pgfqpoint{6.628078in}{9.404019in}}%
\pgfpathlineto{\pgfqpoint{6.628078in}{9.360000in}}%
\pgfpathlineto{\pgfqpoint{6.614280in}{9.360000in}}%
\pgfpathlineto{\pgfqpoint{6.614280in}{9.443550in}}%
\pgfpathlineto{\pgfqpoint{6.628078in}{9.443550in}}%
\pgfpathlineto{\pgfqpoint{6.628078in}{9.430564in}}%
\pgfpathquadraticcurveto{\pgfqpoint{6.632422in}{9.438179in}}{\pgfqpoint{6.639345in}{9.441855in}}%
\pgfpathquadraticcurveto{\pgfqpoint{6.646292in}{9.445556in}}{\pgfqpoint{6.656222in}{9.445556in}}%
\pgfpathquadraticcurveto{\pgfqpoint{6.657631in}{9.445556in}}{\pgfqpoint{6.659349in}{9.445365in}}%
\pgfpathquadraticcurveto{\pgfqpoint{6.661068in}{9.445197in}}{\pgfqpoint{6.663145in}{9.444816in}}%
\pgfpathlineto{\pgfqpoint{6.663217in}{9.430731in}}%
\pgfpathlineto{\pgfqpoint{6.663217in}{9.430731in}}%
\pgfpathclose%
\pgfpathmoveto{\pgfqpoint{6.745717in}{9.405213in}}%
\pgfpathlineto{\pgfqpoint{6.745717in}{9.398505in}}%
\pgfpathlineto{\pgfqpoint{6.682600in}{9.398505in}}%
\pgfpathquadraticcurveto{\pgfqpoint{6.683507in}{9.384325in}}{\pgfqpoint{6.691146in}{9.376901in}}%
\pgfpathquadraticcurveto{\pgfqpoint{6.698785in}{9.369477in}}{\pgfqpoint{6.712440in}{9.369477in}}%
\pgfpathquadraticcurveto{\pgfqpoint{6.720341in}{9.369477in}}{\pgfqpoint{6.727765in}{9.371411in}}%
\pgfpathquadraticcurveto{\pgfqpoint{6.735189in}{9.373344in}}{\pgfqpoint{6.742518in}{9.377235in}}%
\pgfpathlineto{\pgfqpoint{6.742518in}{9.364249in}}%
\pgfpathquadraticcurveto{\pgfqpoint{6.735118in}{9.361122in}}{\pgfqpoint{6.727359in}{9.359475in}}%
\pgfpathquadraticcurveto{\pgfqpoint{6.719601in}{9.357828in}}{\pgfqpoint{6.711628in}{9.357828in}}%
\pgfpathquadraticcurveto{\pgfqpoint{6.691624in}{9.357828in}}{\pgfqpoint{6.679951in}{9.369453in}}%
\pgfpathquadraticcurveto{\pgfqpoint{6.668277in}{9.381102in}}{\pgfqpoint{6.668277in}{9.400964in}}%
\pgfpathquadraticcurveto{\pgfqpoint{6.668277in}{9.421469in}}{\pgfqpoint{6.679354in}{9.433500in}}%
\pgfpathquadraticcurveto{\pgfqpoint{6.690430in}{9.445556in}}{\pgfqpoint{6.709241in}{9.445556in}}%
\pgfpathquadraticcurveto{\pgfqpoint{6.726094in}{9.445556in}}{\pgfqpoint{6.735905in}{9.434694in}}%
\pgfpathquadraticcurveto{\pgfqpoint{6.745717in}{9.423856in}}{\pgfqpoint{6.745717in}{9.405213in}}%
\pgfpathlineto{\pgfqpoint{6.745717in}{9.405213in}}%
\pgfpathclose%
\pgfpathmoveto{\pgfqpoint{6.731990in}{9.409247in}}%
\pgfpathquadraticcurveto{\pgfqpoint{6.731847in}{9.420490in}}{\pgfqpoint{6.725688in}{9.427198in}}%
\pgfpathquadraticcurveto{\pgfqpoint{6.719530in}{9.433930in}}{\pgfqpoint{6.709384in}{9.433930in}}%
\pgfpathquadraticcurveto{\pgfqpoint{6.697902in}{9.433930in}}{\pgfqpoint{6.691003in}{9.427437in}}%
\pgfpathquadraticcurveto{\pgfqpoint{6.684104in}{9.420944in}}{\pgfqpoint{6.683054in}{9.409151in}}%
\pgfpathlineto{\pgfqpoint{6.731990in}{9.409247in}}%
\pgfpathlineto{\pgfqpoint{6.731990in}{9.409247in}}%
\pgfpathclose%
\pgfpathmoveto{\pgfqpoint{6.837717in}{9.410441in}}%
\pgfpathlineto{\pgfqpoint{6.837717in}{9.360000in}}%
\pgfpathlineto{\pgfqpoint{6.823991in}{9.360000in}}%
\pgfpathlineto{\pgfqpoint{6.823991in}{9.409987in}}%
\pgfpathquadraticcurveto{\pgfqpoint{6.823991in}{9.421851in}}{\pgfqpoint{6.819360in}{9.427724in}}%
\pgfpathquadraticcurveto{\pgfqpoint{6.814729in}{9.433620in}}{\pgfqpoint{6.805491in}{9.433620in}}%
\pgfpathquadraticcurveto{\pgfqpoint{6.794367in}{9.433620in}}{\pgfqpoint{6.787945in}{9.426530in}}%
\pgfpathquadraticcurveto{\pgfqpoint{6.781524in}{9.419464in}}{\pgfqpoint{6.781524in}{9.407218in}}%
\pgfpathlineto{\pgfqpoint{6.781524in}{9.360000in}}%
\pgfpathlineto{\pgfqpoint{6.767726in}{9.360000in}}%
\pgfpathlineto{\pgfqpoint{6.767726in}{9.443550in}}%
\pgfpathlineto{\pgfqpoint{6.781524in}{9.443550in}}%
\pgfpathlineto{\pgfqpoint{6.781524in}{9.430564in}}%
\pgfpathquadraticcurveto{\pgfqpoint{6.786465in}{9.438108in}}{\pgfqpoint{6.793125in}{9.441832in}}%
\pgfpathquadraticcurveto{\pgfqpoint{6.799809in}{9.445556in}}{\pgfqpoint{6.808546in}{9.445556in}}%
\pgfpathquadraticcurveto{\pgfqpoint{6.822941in}{9.445556in}}{\pgfqpoint{6.830317in}{9.436651in}}%
\pgfpathquadraticcurveto{\pgfqpoint{6.837717in}{9.427747in}}{\pgfqpoint{6.837717in}{9.410441in}}%
\pgfpathlineto{\pgfqpoint{6.837717in}{9.410441in}}%
\pgfpathclose%
\pgfpathmoveto{\pgfqpoint{6.925207in}{9.440352in}}%
\pgfpathlineto{\pgfqpoint{6.925207in}{9.427509in}}%
\pgfpathquadraticcurveto{\pgfqpoint{6.919382in}{9.430731in}}{\pgfqpoint{6.913533in}{9.432331in}}%
\pgfpathquadraticcurveto{\pgfqpoint{6.907685in}{9.433930in}}{\pgfqpoint{6.901717in}{9.433930in}}%
\pgfpathquadraticcurveto{\pgfqpoint{6.888349in}{9.433930in}}{\pgfqpoint{6.880949in}{9.425456in}}%
\pgfpathquadraticcurveto{\pgfqpoint{6.873572in}{9.417005in}}{\pgfqpoint{6.873572in}{9.401704in}}%
\pgfpathquadraticcurveto{\pgfqpoint{6.873572in}{9.386402in}}{\pgfqpoint{6.880949in}{9.377928in}}%
\pgfpathquadraticcurveto{\pgfqpoint{6.888349in}{9.369477in}}{\pgfqpoint{6.901717in}{9.369477in}}%
\pgfpathquadraticcurveto{\pgfqpoint{6.907685in}{9.369477in}}{\pgfqpoint{6.913533in}{9.371076in}}%
\pgfpathquadraticcurveto{\pgfqpoint{6.919382in}{9.372676in}}{\pgfqpoint{6.925207in}{9.375898in}}%
\pgfpathlineto{\pgfqpoint{6.925207in}{9.363199in}}%
\pgfpathquadraticcurveto{\pgfqpoint{6.919454in}{9.360525in}}{\pgfqpoint{6.913295in}{9.359188in}}%
\pgfpathquadraticcurveto{\pgfqpoint{6.907160in}{9.357828in}}{\pgfqpoint{6.900213in}{9.357828in}}%
\pgfpathquadraticcurveto{\pgfqpoint{6.881331in}{9.357828in}}{\pgfqpoint{6.870207in}{9.369692in}}%
\pgfpathquadraticcurveto{\pgfqpoint{6.859106in}{9.381556in}}{\pgfqpoint{6.859106in}{9.401704in}}%
\pgfpathquadraticcurveto{\pgfqpoint{6.859106in}{9.422138in}}{\pgfqpoint{6.870326in}{9.433835in}}%
\pgfpathquadraticcurveto{\pgfqpoint{6.881569in}{9.445556in}}{\pgfqpoint{6.901120in}{9.445556in}}%
\pgfpathquadraticcurveto{\pgfqpoint{6.907446in}{9.445556in}}{\pgfqpoint{6.913486in}{9.444243in}}%
\pgfpathquadraticcurveto{\pgfqpoint{6.919525in}{9.442954in}}{\pgfqpoint{6.925207in}{9.440352in}}%
\pgfpathlineto{\pgfqpoint{6.925207in}{9.440352in}}%
\pgfpathclose%
\pgfpathmoveto{\pgfqpoint{7.020549in}{9.405213in}}%
\pgfpathlineto{\pgfqpoint{7.020549in}{9.398505in}}%
\pgfpathlineto{\pgfqpoint{6.957433in}{9.398505in}}%
\pgfpathquadraticcurveto{\pgfqpoint{6.958340in}{9.384325in}}{\pgfqpoint{6.965979in}{9.376901in}}%
\pgfpathquadraticcurveto{\pgfqpoint{6.973618in}{9.369477in}}{\pgfqpoint{6.987273in}{9.369477in}}%
\pgfpathquadraticcurveto{\pgfqpoint{6.995174in}{9.369477in}}{\pgfqpoint{7.002598in}{9.371411in}}%
\pgfpathquadraticcurveto{\pgfqpoint{7.010022in}{9.373344in}}{\pgfqpoint{7.017351in}{9.377235in}}%
\pgfpathlineto{\pgfqpoint{7.017351in}{9.364249in}}%
\pgfpathquadraticcurveto{\pgfqpoint{7.009951in}{9.361122in}}{\pgfqpoint{7.002192in}{9.359475in}}%
\pgfpathquadraticcurveto{\pgfqpoint{6.994434in}{9.357828in}}{\pgfqpoint{6.986461in}{9.357828in}}%
\pgfpathquadraticcurveto{\pgfqpoint{6.966457in}{9.357828in}}{\pgfqpoint{6.954783in}{9.369453in}}%
\pgfpathquadraticcurveto{\pgfqpoint{6.943110in}{9.381102in}}{\pgfqpoint{6.943110in}{9.400964in}}%
\pgfpathquadraticcurveto{\pgfqpoint{6.943110in}{9.421469in}}{\pgfqpoint{6.954187in}{9.433500in}}%
\pgfpathquadraticcurveto{\pgfqpoint{6.965263in}{9.445556in}}{\pgfqpoint{6.984074in}{9.445556in}}%
\pgfpathquadraticcurveto{\pgfqpoint{7.000927in}{9.445556in}}{\pgfqpoint{7.010738in}{9.434694in}}%
\pgfpathquadraticcurveto{\pgfqpoint{7.020549in}{9.423856in}}{\pgfqpoint{7.020549in}{9.405213in}}%
\pgfpathlineto{\pgfqpoint{7.020549in}{9.405213in}}%
\pgfpathclose%
\pgfpathmoveto{\pgfqpoint{7.006823in}{9.409247in}}%
\pgfpathquadraticcurveto{\pgfqpoint{7.006680in}{9.420490in}}{\pgfqpoint{7.000521in}{9.427198in}}%
\pgfpathquadraticcurveto{\pgfqpoint{6.994362in}{9.433930in}}{\pgfqpoint{6.984217in}{9.433930in}}%
\pgfpathquadraticcurveto{\pgfqpoint{6.972735in}{9.433930in}}{\pgfqpoint{6.965836in}{9.427437in}}%
\pgfpathquadraticcurveto{\pgfqpoint{6.958937in}{9.420944in}}{\pgfqpoint{6.957887in}{9.409151in}}%
\pgfpathlineto{\pgfqpoint{7.006823in}{9.409247in}}%
\pgfpathlineto{\pgfqpoint{7.006823in}{9.409247in}}%
\pgfpathclose%
\pgfusepath{fill}%
\end{pgfscope}%
\begin{pgfscope}%
\pgfsetbuttcap%
\pgfsetroundjoin%
\definecolor{currentfill}{rgb}{0.090196,0.168627,0.227451}%
\pgfsetfillcolor{currentfill}%
\pgfsetlinewidth{0.000000pt}%
\definecolor{currentstroke}{rgb}{0.090196,0.168627,0.227451}%
\pgfsetstrokecolor{currentstroke}%
\pgfsetdash{}{0pt}%
\pgfpathmoveto{\pgfqpoint{1.316611in}{2.290690in}}%
\pgfpathlineto{\pgfqpoint{1.316611in}{2.277991in}}%
\pgfpathquadraticcurveto{\pgfqpoint{1.309210in}{2.281537in}}{\pgfqpoint{1.302633in}{2.283268in}}%
\pgfpathquadraticcurveto{\pgfqpoint{1.296056in}{2.285021in}}{\pgfqpoint{1.289933in}{2.285021in}}%
\pgfpathquadraticcurveto{\pgfqpoint{1.279316in}{2.285021in}}{\pgfqpoint{1.273543in}{2.280897in}}%
\pgfpathquadraticcurveto{\pgfqpoint{1.267771in}{2.276774in}}{\pgfqpoint{1.267771in}{2.269167in}}%
\pgfpathquadraticcurveto{\pgfqpoint{1.267771in}{2.262776in}}{\pgfqpoint{1.271605in}{2.259518in}}%
\pgfpathquadraticcurveto{\pgfqpoint{1.275440in}{2.256282in}}{\pgfqpoint{1.286140in}{2.254282in}}%
\pgfpathlineto{\pgfqpoint{1.293995in}{2.252674in}}%
\pgfpathquadraticcurveto{\pgfqpoint{1.308550in}{2.249891in}}{\pgfqpoint{1.315477in}{2.242902in}}%
\pgfpathquadraticcurveto{\pgfqpoint{1.322404in}{2.235913in}}{\pgfqpoint{1.322404in}{2.224203in}}%
\pgfpathquadraticcurveto{\pgfqpoint{1.322404in}{2.210204in}}{\pgfqpoint{1.313024in}{2.202988in}}%
\pgfpathquadraticcurveto{\pgfqpoint{1.303664in}{2.195773in}}{\pgfqpoint{1.285563in}{2.195773in}}%
\pgfpathquadraticcurveto{\pgfqpoint{1.278739in}{2.195773in}}{\pgfqpoint{1.271028in}{2.197319in}}%
\pgfpathquadraticcurveto{\pgfqpoint{1.263338in}{2.198865in}}{\pgfqpoint{1.255092in}{2.201896in}}%
\pgfpathlineto{\pgfqpoint{1.255092in}{2.215296in}}%
\pgfpathquadraticcurveto{\pgfqpoint{1.263008in}{2.210864in}}{\pgfqpoint{1.270616in}{2.208596in}}%
\pgfpathquadraticcurveto{\pgfqpoint{1.278223in}{2.206349in}}{\pgfqpoint{1.285563in}{2.206349in}}%
\pgfpathquadraticcurveto{\pgfqpoint{1.296696in}{2.206349in}}{\pgfqpoint{1.302757in}{2.210720in}}%
\pgfpathquadraticcurveto{\pgfqpoint{1.308818in}{2.215111in}}{\pgfqpoint{1.308818in}{2.223234in}}%
\pgfpathquadraticcurveto{\pgfqpoint{1.308818in}{2.230305in}}{\pgfqpoint{1.304468in}{2.234305in}}%
\pgfpathquadraticcurveto{\pgfqpoint{1.300118in}{2.238304in}}{\pgfqpoint{1.290201in}{2.240304in}}%
\pgfpathlineto{\pgfqpoint{1.282264in}{2.241850in}}%
\pgfpathquadraticcurveto{\pgfqpoint{1.267709in}{2.244736in}}{\pgfqpoint{1.261194in}{2.250921in}}%
\pgfpathquadraticcurveto{\pgfqpoint{1.254700in}{2.257106in}}{\pgfqpoint{1.254700in}{2.268136in}}%
\pgfpathquadraticcurveto{\pgfqpoint{1.254700in}{2.280897in}}{\pgfqpoint{1.263689in}{2.288237in}}%
\pgfpathquadraticcurveto{\pgfqpoint{1.272678in}{2.295576in}}{\pgfqpoint{1.288449in}{2.295576in}}%
\pgfpathquadraticcurveto{\pgfqpoint{1.295232in}{2.295576in}}{\pgfqpoint{1.302241in}{2.294339in}}%
\pgfpathquadraticcurveto{\pgfqpoint{1.309271in}{2.293123in}}{\pgfqpoint{1.316611in}{2.290690in}}%
\pgfpathlineto{\pgfqpoint{1.316611in}{2.290690in}}%
\pgfpathclose%
\pgfpathmoveto{\pgfqpoint{1.353926in}{2.290299in}}%
\pgfpathlineto{\pgfqpoint{1.353926in}{2.269806in}}%
\pgfpathlineto{\pgfqpoint{1.378336in}{2.269806in}}%
\pgfpathlineto{\pgfqpoint{1.378336in}{2.260590in}}%
\pgfpathlineto{\pgfqpoint{1.353926in}{2.260590in}}%
\pgfpathlineto{\pgfqpoint{1.353926in}{2.221419in}}%
\pgfpathquadraticcurveto{\pgfqpoint{1.353926in}{2.212596in}}{\pgfqpoint{1.356339in}{2.210080in}}%
\pgfpathquadraticcurveto{\pgfqpoint{1.358751in}{2.207565in}}{\pgfqpoint{1.366173in}{2.207565in}}%
\pgfpathlineto{\pgfqpoint{1.378336in}{2.207565in}}%
\pgfpathlineto{\pgfqpoint{1.378336in}{2.197649in}}%
\pgfpathlineto{\pgfqpoint{1.366173in}{2.197649in}}%
\pgfpathquadraticcurveto{\pgfqpoint{1.352442in}{2.197649in}}{\pgfqpoint{1.347226in}{2.202762in}}%
\pgfpathquadraticcurveto{\pgfqpoint{1.342010in}{2.207895in}}{\pgfqpoint{1.342010in}{2.221419in}}%
\pgfpathlineto{\pgfqpoint{1.342010in}{2.260590in}}%
\pgfpathlineto{\pgfqpoint{1.333310in}{2.260590in}}%
\pgfpathlineto{\pgfqpoint{1.333310in}{2.269806in}}%
\pgfpathlineto{\pgfqpoint{1.342010in}{2.269806in}}%
\pgfpathlineto{\pgfqpoint{1.342010in}{2.290299in}}%
\pgfpathlineto{\pgfqpoint{1.353926in}{2.290299in}}%
\pgfpathlineto{\pgfqpoint{1.353926in}{2.290299in}}%
\pgfpathclose%
\pgfpathmoveto{\pgfqpoint{1.392706in}{2.226120in}}%
\pgfpathlineto{\pgfqpoint{1.392706in}{2.269806in}}%
\pgfpathlineto{\pgfqpoint{1.404560in}{2.269806in}}%
\pgfpathlineto{\pgfqpoint{1.404560in}{2.226573in}}%
\pgfpathquadraticcurveto{\pgfqpoint{1.404560in}{2.216327in}}{\pgfqpoint{1.408539in}{2.211194in}}%
\pgfpathquadraticcurveto{\pgfqpoint{1.412539in}{2.206081in}}{\pgfqpoint{1.420538in}{2.206081in}}%
\pgfpathquadraticcurveto{\pgfqpoint{1.430124in}{2.206081in}}{\pgfqpoint{1.435691in}{2.212204in}}%
\pgfpathquadraticcurveto{\pgfqpoint{1.441278in}{2.218327in}}{\pgfqpoint{1.441278in}{2.228903in}}%
\pgfpathlineto{\pgfqpoint{1.441278in}{2.269806in}}%
\pgfpathlineto{\pgfqpoint{1.453132in}{2.269806in}}%
\pgfpathlineto{\pgfqpoint{1.453132in}{2.197649in}}%
\pgfpathlineto{\pgfqpoint{1.441278in}{2.197649in}}%
\pgfpathlineto{\pgfqpoint{1.441278in}{2.208740in}}%
\pgfpathquadraticcurveto{\pgfqpoint{1.436969in}{2.202164in}}{\pgfqpoint{1.431258in}{2.198968in}}%
\pgfpathquadraticcurveto{\pgfqpoint{1.425568in}{2.195773in}}{\pgfqpoint{1.418023in}{2.195773in}}%
\pgfpathquadraticcurveto{\pgfqpoint{1.405591in}{2.195773in}}{\pgfqpoint{1.399138in}{2.203504in}}%
\pgfpathquadraticcurveto{\pgfqpoint{1.392706in}{2.211235in}}{\pgfqpoint{1.392706in}{2.226120in}}%
\pgfpathlineto{\pgfqpoint{1.392706in}{2.226120in}}%
\pgfpathclose%
\pgfpathmoveto{\pgfqpoint{1.422538in}{2.271538in}}%
\pgfpathlineto{\pgfqpoint{1.422538in}{2.271538in}}%
\pgfpathlineto{\pgfqpoint{1.422538in}{2.271538in}}%
\pgfpathclose%
\pgfpathmoveto{\pgfqpoint{1.525021in}{2.258859in}}%
\pgfpathlineto{\pgfqpoint{1.525021in}{2.297906in}}%
\pgfpathlineto{\pgfqpoint{1.536876in}{2.297906in}}%
\pgfpathlineto{\pgfqpoint{1.536876in}{2.197649in}}%
\pgfpathlineto{\pgfqpoint{1.525021in}{2.197649in}}%
\pgfpathlineto{\pgfqpoint{1.525021in}{2.208472in}}%
\pgfpathquadraticcurveto{\pgfqpoint{1.521290in}{2.202040in}}{\pgfqpoint{1.515579in}{2.198906in}}%
\pgfpathquadraticcurveto{\pgfqpoint{1.509889in}{2.195773in}}{\pgfqpoint{1.501890in}{2.195773in}}%
\pgfpathquadraticcurveto{\pgfqpoint{1.488819in}{2.195773in}}{\pgfqpoint{1.480593in}{2.206205in}}%
\pgfpathquadraticcurveto{\pgfqpoint{1.472388in}{2.216657in}}{\pgfqpoint{1.472388in}{2.233666in}}%
\pgfpathquadraticcurveto{\pgfqpoint{1.472388in}{2.250674in}}{\pgfqpoint{1.480593in}{2.261106in}}%
\pgfpathquadraticcurveto{\pgfqpoint{1.488819in}{2.271538in}}{\pgfqpoint{1.501890in}{2.271538in}}%
\pgfpathquadraticcurveto{\pgfqpoint{1.509889in}{2.271538in}}{\pgfqpoint{1.515579in}{2.268404in}}%
\pgfpathquadraticcurveto{\pgfqpoint{1.521290in}{2.265291in}}{\pgfqpoint{1.525021in}{2.258859in}}%
\pgfpathlineto{\pgfqpoint{1.525021in}{2.258859in}}%
\pgfpathclose%
\pgfpathmoveto{\pgfqpoint{1.484634in}{2.233666in}}%
\pgfpathquadraticcurveto{\pgfqpoint{1.484634in}{2.220595in}}{\pgfqpoint{1.490015in}{2.213152in}}%
\pgfpathquadraticcurveto{\pgfqpoint{1.495396in}{2.205710in}}{\pgfqpoint{1.504797in}{2.205710in}}%
\pgfpathquadraticcurveto{\pgfqpoint{1.514198in}{2.205710in}}{\pgfqpoint{1.519599in}{2.213152in}}%
\pgfpathquadraticcurveto{\pgfqpoint{1.525021in}{2.220595in}}{\pgfqpoint{1.525021in}{2.233666in}}%
\pgfpathquadraticcurveto{\pgfqpoint{1.525021in}{2.246736in}}{\pgfqpoint{1.519599in}{2.254179in}}%
\pgfpathquadraticcurveto{\pgfqpoint{1.514198in}{2.261621in}}{\pgfqpoint{1.504797in}{2.261621in}}%
\pgfpathquadraticcurveto{\pgfqpoint{1.495396in}{2.261621in}}{\pgfqpoint{1.490015in}{2.254179in}}%
\pgfpathquadraticcurveto{\pgfqpoint{1.484634in}{2.246736in}}{\pgfqpoint{1.484634in}{2.233666in}}%
\pgfpathlineto{\pgfqpoint{1.484634in}{2.233666in}}%
\pgfpathclose%
\pgfpathmoveto{\pgfqpoint{1.591323in}{2.190948in}}%
\pgfpathquadraticcurveto{\pgfqpoint{1.586314in}{2.178063in}}{\pgfqpoint{1.581531in}{2.174146in}}%
\pgfpathquadraticcurveto{\pgfqpoint{1.576768in}{2.170208in}}{\pgfqpoint{1.568790in}{2.170208in}}%
\pgfpathlineto{\pgfqpoint{1.559306in}{2.170208in}}%
\pgfpathlineto{\pgfqpoint{1.559306in}{2.180125in}}%
\pgfpathlineto{\pgfqpoint{1.566275in}{2.180125in}}%
\pgfpathquadraticcurveto{\pgfqpoint{1.571161in}{2.180125in}}{\pgfqpoint{1.573861in}{2.182455in}}%
\pgfpathquadraticcurveto{\pgfqpoint{1.576583in}{2.184764in}}{\pgfqpoint{1.579861in}{2.193402in}}%
\pgfpathlineto{\pgfqpoint{1.581984in}{2.198803in}}%
\pgfpathlineto{\pgfqpoint{1.552812in}{2.269806in}}%
\pgfpathlineto{\pgfqpoint{1.565367in}{2.269806in}}%
\pgfpathlineto{\pgfqpoint{1.587922in}{2.213379in}}%
\pgfpathlineto{\pgfqpoint{1.610476in}{2.269806in}}%
\pgfpathlineto{\pgfqpoint{1.623031in}{2.269806in}}%
\pgfpathlineto{\pgfqpoint{1.591323in}{2.190948in}}%
\pgfpathlineto{\pgfqpoint{1.591323in}{2.190948in}}%
\pgfpathclose%
\pgfpathmoveto{\pgfqpoint{1.685272in}{2.208596in}}%
\pgfpathlineto{\pgfqpoint{1.706527in}{2.208596in}}%
\pgfpathlineto{\pgfqpoint{1.706527in}{2.281990in}}%
\pgfpathlineto{\pgfqpoint{1.683396in}{2.277351in}}%
\pgfpathlineto{\pgfqpoint{1.683396in}{2.289206in}}%
\pgfpathlineto{\pgfqpoint{1.706404in}{2.293845in}}%
\pgfpathlineto{\pgfqpoint{1.719413in}{2.293845in}}%
\pgfpathlineto{\pgfqpoint{1.719413in}{2.208596in}}%
\pgfpathlineto{\pgfqpoint{1.740668in}{2.208596in}}%
\pgfpathlineto{\pgfqpoint{1.740668in}{2.197649in}}%
\pgfpathlineto{\pgfqpoint{1.685272in}{2.197649in}}%
\pgfpathlineto{\pgfqpoint{1.685272in}{2.208596in}}%
\pgfpathlineto{\pgfqpoint{1.685272in}{2.208596in}}%
\pgfpathclose%
\pgfpathmoveto{\pgfqpoint{1.835689in}{2.297762in}}%
\pgfpathquadraticcurveto{\pgfqpoint{1.827071in}{2.282959in}}{\pgfqpoint{1.822865in}{2.268445in}}%
\pgfpathquadraticcurveto{\pgfqpoint{1.818680in}{2.253952in}}{\pgfqpoint{1.818680in}{2.239067in}}%
\pgfpathquadraticcurveto{\pgfqpoint{1.818680in}{2.224203in}}{\pgfqpoint{1.822906in}{2.209606in}}%
\pgfpathquadraticcurveto{\pgfqpoint{1.827133in}{2.195010in}}{\pgfqpoint{1.835689in}{2.180249in}}%
\pgfpathlineto{\pgfqpoint{1.825380in}{2.180249in}}%
\pgfpathquadraticcurveto{\pgfqpoint{1.815732in}{2.195402in}}{\pgfqpoint{1.810928in}{2.210019in}}%
\pgfpathquadraticcurveto{\pgfqpoint{1.806125in}{2.224636in}}{\pgfqpoint{1.806125in}{2.239067in}}%
\pgfpathquadraticcurveto{\pgfqpoint{1.806125in}{2.253437in}}{\pgfqpoint{1.810887in}{2.267992in}}%
\pgfpathquadraticcurveto{\pgfqpoint{1.815670in}{2.282567in}}{\pgfqpoint{1.825380in}{2.297762in}}%
\pgfpathlineto{\pgfqpoint{1.835689in}{2.297762in}}%
\pgfpathlineto{\pgfqpoint{1.835689in}{2.297762in}}%
\pgfpathclose%
\pgfpathmoveto{\pgfqpoint{1.871582in}{2.208596in}}%
\pgfpathlineto{\pgfqpoint{1.916999in}{2.208596in}}%
\pgfpathlineto{\pgfqpoint{1.916999in}{2.197649in}}%
\pgfpathlineto{\pgfqpoint{1.855934in}{2.197649in}}%
\pgfpathlineto{\pgfqpoint{1.855934in}{2.208596in}}%
\pgfpathquadraticcurveto{\pgfqpoint{1.863335in}{2.216265in}}{\pgfqpoint{1.876117in}{2.229171in}}%
\pgfpathquadraticcurveto{\pgfqpoint{1.888920in}{2.242098in}}{\pgfqpoint{1.892198in}{2.245850in}}%
\pgfpathquadraticcurveto{\pgfqpoint{1.898445in}{2.252859in}}{\pgfqpoint{1.900919in}{2.257725in}}%
\pgfpathquadraticcurveto{\pgfqpoint{1.903413in}{2.262590in}}{\pgfqpoint{1.903413in}{2.267291in}}%
\pgfpathquadraticcurveto{\pgfqpoint{1.903413in}{2.274960in}}{\pgfqpoint{1.898032in}{2.279784in}}%
\pgfpathquadraticcurveto{\pgfqpoint{1.892651in}{2.284629in}}{\pgfqpoint{1.884013in}{2.284629in}}%
\pgfpathquadraticcurveto{\pgfqpoint{1.877890in}{2.284629in}}{\pgfqpoint{1.871087in}{2.282506in}}%
\pgfpathquadraticcurveto{\pgfqpoint{1.864304in}{2.280382in}}{\pgfqpoint{1.856573in}{2.276053in}}%
\pgfpathlineto{\pgfqpoint{1.856573in}{2.289206in}}%
\pgfpathquadraticcurveto{\pgfqpoint{1.864428in}{2.292360in}}{\pgfqpoint{1.871252in}{2.293968in}}%
\pgfpathquadraticcurveto{\pgfqpoint{1.878096in}{2.295576in}}{\pgfqpoint{1.883766in}{2.295576in}}%
\pgfpathquadraticcurveto{\pgfqpoint{1.898713in}{2.295576in}}{\pgfqpoint{1.907598in}{2.288093in}}%
\pgfpathquadraticcurveto{\pgfqpoint{1.916484in}{2.280629in}}{\pgfqpoint{1.916484in}{2.268136in}}%
\pgfpathquadraticcurveto{\pgfqpoint{1.916484in}{2.262198in}}{\pgfqpoint{1.914257in}{2.256879in}}%
\pgfpathquadraticcurveto{\pgfqpoint{1.912051in}{2.251581in}}{\pgfqpoint{1.906176in}{2.244365in}}%
\pgfpathquadraticcurveto{\pgfqpoint{1.904568in}{2.242489in}}{\pgfqpoint{1.895929in}{2.233562in}}%
\pgfpathquadraticcurveto{\pgfqpoint{1.887312in}{2.224636in}}{\pgfqpoint{1.871582in}{2.208596in}}%
\pgfpathlineto{\pgfqpoint{1.871582in}{2.208596in}}%
\pgfpathclose%
\pgfpathmoveto{\pgfqpoint{1.972148in}{2.285268in}}%
\pgfpathquadraticcurveto{\pgfqpoint{1.962108in}{2.285268in}}{\pgfqpoint{1.957036in}{2.275372in}}%
\pgfpathquadraticcurveto{\pgfqpoint{1.951985in}{2.265497in}}{\pgfqpoint{1.951985in}{2.245644in}}%
\pgfpathquadraticcurveto{\pgfqpoint{1.951985in}{2.225873in}}{\pgfqpoint{1.957036in}{2.215977in}}%
\pgfpathquadraticcurveto{\pgfqpoint{1.962108in}{2.206081in}}{\pgfqpoint{1.972148in}{2.206081in}}%
\pgfpathquadraticcurveto{\pgfqpoint{1.982271in}{2.206081in}}{\pgfqpoint{1.987322in}{2.215977in}}%
\pgfpathquadraticcurveto{\pgfqpoint{1.992393in}{2.225873in}}{\pgfqpoint{1.992393in}{2.245644in}}%
\pgfpathquadraticcurveto{\pgfqpoint{1.992393in}{2.265497in}}{\pgfqpoint{1.987322in}{2.275372in}}%
\pgfpathquadraticcurveto{\pgfqpoint{1.982271in}{2.285268in}}{\pgfqpoint{1.972148in}{2.285268in}}%
\pgfpathlineto{\pgfqpoint{1.972148in}{2.285268in}}%
\pgfpathclose%
\pgfpathmoveto{\pgfqpoint{1.972148in}{2.295576in}}%
\pgfpathquadraticcurveto{\pgfqpoint{1.988332in}{2.295576in}}{\pgfqpoint{1.996867in}{2.282774in}}%
\pgfpathquadraticcurveto{\pgfqpoint{2.005402in}{2.269991in}}{\pgfqpoint{2.005402in}{2.245644in}}%
\pgfpathquadraticcurveto{\pgfqpoint{2.005402in}{2.221358in}}{\pgfqpoint{1.996867in}{2.208555in}}%
\pgfpathquadraticcurveto{\pgfqpoint{1.988332in}{2.195773in}}{\pgfqpoint{1.972148in}{2.195773in}}%
\pgfpathquadraticcurveto{\pgfqpoint{1.955985in}{2.195773in}}{\pgfqpoint{1.947450in}{2.208555in}}%
\pgfpathquadraticcurveto{\pgfqpoint{1.938914in}{2.221358in}}{\pgfqpoint{1.938914in}{2.245644in}}%
\pgfpathquadraticcurveto{\pgfqpoint{1.938914in}{2.269991in}}{\pgfqpoint{1.947450in}{2.282774in}}%
\pgfpathquadraticcurveto{\pgfqpoint{1.955985in}{2.295576in}}{\pgfqpoint{1.972148in}{2.295576in}}%
\pgfpathlineto{\pgfqpoint{1.972148in}{2.295576in}}%
\pgfpathclose%
\pgfpathmoveto{\pgfqpoint{2.030533in}{2.208596in}}%
\pgfpathlineto{\pgfqpoint{2.051789in}{2.208596in}}%
\pgfpathlineto{\pgfqpoint{2.051789in}{2.281990in}}%
\pgfpathlineto{\pgfqpoint{2.028657in}{2.277351in}}%
\pgfpathlineto{\pgfqpoint{2.028657in}{2.289206in}}%
\pgfpathlineto{\pgfqpoint{2.051665in}{2.293845in}}%
\pgfpathlineto{\pgfqpoint{2.064674in}{2.293845in}}%
\pgfpathlineto{\pgfqpoint{2.064674in}{2.208596in}}%
\pgfpathlineto{\pgfqpoint{2.085929in}{2.208596in}}%
\pgfpathlineto{\pgfqpoint{2.085929in}{2.197649in}}%
\pgfpathlineto{\pgfqpoint{2.030533in}{2.197649in}}%
\pgfpathlineto{\pgfqpoint{2.030533in}{2.208596in}}%
\pgfpathlineto{\pgfqpoint{2.030533in}{2.208596in}}%
\pgfpathclose%
\pgfpathmoveto{\pgfqpoint{2.108937in}{2.293845in}}%
\pgfpathlineto{\pgfqpoint{2.170786in}{2.293845in}}%
\pgfpathlineto{\pgfqpoint{2.170786in}{2.288299in}}%
\pgfpathlineto{\pgfqpoint{2.135862in}{2.197649in}}%
\pgfpathlineto{\pgfqpoint{2.122276in}{2.197649in}}%
\pgfpathlineto{\pgfqpoint{2.155138in}{2.282877in}}%
\pgfpathlineto{\pgfqpoint{2.108937in}{2.282877in}}%
\pgfpathlineto{\pgfqpoint{2.108937in}{2.293845in}}%
\pgfpathlineto{\pgfqpoint{2.108937in}{2.293845in}}%
\pgfpathclose%
\pgfpathmoveto{\pgfqpoint{2.192640in}{2.297762in}}%
\pgfpathlineto{\pgfqpoint{2.202948in}{2.297762in}}%
\pgfpathquadraticcurveto{\pgfqpoint{2.212596in}{2.282567in}}{\pgfqpoint{2.217400in}{2.267992in}}%
\pgfpathquadraticcurveto{\pgfqpoint{2.222203in}{2.253437in}}{\pgfqpoint{2.222203in}{2.239067in}}%
\pgfpathquadraticcurveto{\pgfqpoint{2.222203in}{2.224636in}}{\pgfqpoint{2.217400in}{2.210019in}}%
\pgfpathquadraticcurveto{\pgfqpoint{2.212596in}{2.195402in}}{\pgfqpoint{2.202948in}{2.180249in}}%
\pgfpathlineto{\pgfqpoint{2.192640in}{2.180249in}}%
\pgfpathquadraticcurveto{\pgfqpoint{2.201195in}{2.195010in}}{\pgfqpoint{2.205422in}{2.209606in}}%
\pgfpathquadraticcurveto{\pgfqpoint{2.209648in}{2.224203in}}{\pgfqpoint{2.209648in}{2.239067in}}%
\pgfpathquadraticcurveto{\pgfqpoint{2.209648in}{2.253952in}}{\pgfqpoint{2.205422in}{2.268445in}}%
\pgfpathquadraticcurveto{\pgfqpoint{2.201195in}{2.282959in}}{\pgfqpoint{2.192640in}{2.297762in}}%
\pgfpathlineto{\pgfqpoint{2.192640in}{2.297762in}}%
\pgfpathclose%
\pgfusepath{fill}%
\end{pgfscope}%
\begin{pgfscope}%
\pgfsetbuttcap%
\pgfsetroundjoin%
\definecolor{currentfill}{rgb}{0.090196,0.168627,0.227451}%
\pgfsetfillcolor{currentfill}%
\pgfsetlinewidth{0.000000pt}%
\definecolor{currentstroke}{rgb}{0.090196,0.168627,0.227451}%
\pgfsetstrokecolor{currentstroke}%
\pgfsetdash{}{0pt}%
\pgfpathmoveto{\pgfqpoint{1.316611in}{2.114050in}}%
\pgfpathlineto{\pgfqpoint{1.316611in}{2.101351in}}%
\pgfpathquadraticcurveto{\pgfqpoint{1.309210in}{2.104897in}}{\pgfqpoint{1.302633in}{2.106628in}}%
\pgfpathquadraticcurveto{\pgfqpoint{1.296056in}{2.108381in}}{\pgfqpoint{1.289933in}{2.108381in}}%
\pgfpathquadraticcurveto{\pgfqpoint{1.279316in}{2.108381in}}{\pgfqpoint{1.273543in}{2.104257in}}%
\pgfpathquadraticcurveto{\pgfqpoint{1.267771in}{2.100134in}}{\pgfqpoint{1.267771in}{2.092527in}}%
\pgfpathquadraticcurveto{\pgfqpoint{1.267771in}{2.086136in}}{\pgfqpoint{1.271605in}{2.082878in}}%
\pgfpathquadraticcurveto{\pgfqpoint{1.275440in}{2.079642in}}{\pgfqpoint{1.286140in}{2.077642in}}%
\pgfpathlineto{\pgfqpoint{1.293995in}{2.076034in}}%
\pgfpathquadraticcurveto{\pgfqpoint{1.308550in}{2.073251in}}{\pgfqpoint{1.315477in}{2.066262in}}%
\pgfpathquadraticcurveto{\pgfqpoint{1.322404in}{2.059273in}}{\pgfqpoint{1.322404in}{2.047563in}}%
\pgfpathquadraticcurveto{\pgfqpoint{1.322404in}{2.033564in}}{\pgfqpoint{1.313024in}{2.026348in}}%
\pgfpathquadraticcurveto{\pgfqpoint{1.303664in}{2.019133in}}{\pgfqpoint{1.285563in}{2.019133in}}%
\pgfpathquadraticcurveto{\pgfqpoint{1.278739in}{2.019133in}}{\pgfqpoint{1.271028in}{2.020679in}}%
\pgfpathquadraticcurveto{\pgfqpoint{1.263338in}{2.022225in}}{\pgfqpoint{1.255092in}{2.025256in}}%
\pgfpathlineto{\pgfqpoint{1.255092in}{2.038656in}}%
\pgfpathquadraticcurveto{\pgfqpoint{1.263008in}{2.034224in}}{\pgfqpoint{1.270616in}{2.031956in}}%
\pgfpathquadraticcurveto{\pgfqpoint{1.278223in}{2.029709in}}{\pgfqpoint{1.285563in}{2.029709in}}%
\pgfpathquadraticcurveto{\pgfqpoint{1.296696in}{2.029709in}}{\pgfqpoint{1.302757in}{2.034080in}}%
\pgfpathquadraticcurveto{\pgfqpoint{1.308818in}{2.038471in}}{\pgfqpoint{1.308818in}{2.046594in}}%
\pgfpathquadraticcurveto{\pgfqpoint{1.308818in}{2.053665in}}{\pgfqpoint{1.304468in}{2.057665in}}%
\pgfpathquadraticcurveto{\pgfqpoint{1.300118in}{2.061664in}}{\pgfqpoint{1.290201in}{2.063664in}}%
\pgfpathlineto{\pgfqpoint{1.282264in}{2.065210in}}%
\pgfpathquadraticcurveto{\pgfqpoint{1.267709in}{2.068096in}}{\pgfqpoint{1.261194in}{2.074281in}}%
\pgfpathquadraticcurveto{\pgfqpoint{1.254700in}{2.080466in}}{\pgfqpoint{1.254700in}{2.091496in}}%
\pgfpathquadraticcurveto{\pgfqpoint{1.254700in}{2.104257in}}{\pgfqpoint{1.263689in}{2.111597in}}%
\pgfpathquadraticcurveto{\pgfqpoint{1.272678in}{2.118936in}}{\pgfqpoint{1.288449in}{2.118936in}}%
\pgfpathquadraticcurveto{\pgfqpoint{1.295232in}{2.118936in}}{\pgfqpoint{1.302241in}{2.117699in}}%
\pgfpathquadraticcurveto{\pgfqpoint{1.309271in}{2.116483in}}{\pgfqpoint{1.316611in}{2.114050in}}%
\pgfpathlineto{\pgfqpoint{1.316611in}{2.114050in}}%
\pgfpathclose%
\pgfpathmoveto{\pgfqpoint{1.353926in}{2.113659in}}%
\pgfpathlineto{\pgfqpoint{1.353926in}{2.093166in}}%
\pgfpathlineto{\pgfqpoint{1.378336in}{2.093166in}}%
\pgfpathlineto{\pgfqpoint{1.378336in}{2.083950in}}%
\pgfpathlineto{\pgfqpoint{1.353926in}{2.083950in}}%
\pgfpathlineto{\pgfqpoint{1.353926in}{2.044779in}}%
\pgfpathquadraticcurveto{\pgfqpoint{1.353926in}{2.035956in}}{\pgfqpoint{1.356339in}{2.033440in}}%
\pgfpathquadraticcurveto{\pgfqpoint{1.358751in}{2.030925in}}{\pgfqpoint{1.366173in}{2.030925in}}%
\pgfpathlineto{\pgfqpoint{1.378336in}{2.030925in}}%
\pgfpathlineto{\pgfqpoint{1.378336in}{2.021009in}}%
\pgfpathlineto{\pgfqpoint{1.366173in}{2.021009in}}%
\pgfpathquadraticcurveto{\pgfqpoint{1.352442in}{2.021009in}}{\pgfqpoint{1.347226in}{2.026122in}}%
\pgfpathquadraticcurveto{\pgfqpoint{1.342010in}{2.031255in}}{\pgfqpoint{1.342010in}{2.044779in}}%
\pgfpathlineto{\pgfqpoint{1.342010in}{2.083950in}}%
\pgfpathlineto{\pgfqpoint{1.333310in}{2.083950in}}%
\pgfpathlineto{\pgfqpoint{1.333310in}{2.093166in}}%
\pgfpathlineto{\pgfqpoint{1.342010in}{2.093166in}}%
\pgfpathlineto{\pgfqpoint{1.342010in}{2.113659in}}%
\pgfpathlineto{\pgfqpoint{1.353926in}{2.113659in}}%
\pgfpathlineto{\pgfqpoint{1.353926in}{2.113659in}}%
\pgfpathclose%
\pgfpathmoveto{\pgfqpoint{1.392706in}{2.049480in}}%
\pgfpathlineto{\pgfqpoint{1.392706in}{2.093166in}}%
\pgfpathlineto{\pgfqpoint{1.404560in}{2.093166in}}%
\pgfpathlineto{\pgfqpoint{1.404560in}{2.049933in}}%
\pgfpathquadraticcurveto{\pgfqpoint{1.404560in}{2.039687in}}{\pgfqpoint{1.408539in}{2.034554in}}%
\pgfpathquadraticcurveto{\pgfqpoint{1.412539in}{2.029441in}}{\pgfqpoint{1.420538in}{2.029441in}}%
\pgfpathquadraticcurveto{\pgfqpoint{1.430124in}{2.029441in}}{\pgfqpoint{1.435691in}{2.035564in}}%
\pgfpathquadraticcurveto{\pgfqpoint{1.441278in}{2.041687in}}{\pgfqpoint{1.441278in}{2.052263in}}%
\pgfpathlineto{\pgfqpoint{1.441278in}{2.093166in}}%
\pgfpathlineto{\pgfqpoint{1.453132in}{2.093166in}}%
\pgfpathlineto{\pgfqpoint{1.453132in}{2.021009in}}%
\pgfpathlineto{\pgfqpoint{1.441278in}{2.021009in}}%
\pgfpathlineto{\pgfqpoint{1.441278in}{2.032100in}}%
\pgfpathquadraticcurveto{\pgfqpoint{1.436969in}{2.025524in}}{\pgfqpoint{1.431258in}{2.022328in}}%
\pgfpathquadraticcurveto{\pgfqpoint{1.425568in}{2.019133in}}{\pgfqpoint{1.418023in}{2.019133in}}%
\pgfpathquadraticcurveto{\pgfqpoint{1.405591in}{2.019133in}}{\pgfqpoint{1.399138in}{2.026864in}}%
\pgfpathquadraticcurveto{\pgfqpoint{1.392706in}{2.034595in}}{\pgfqpoint{1.392706in}{2.049480in}}%
\pgfpathlineto{\pgfqpoint{1.392706in}{2.049480in}}%
\pgfpathclose%
\pgfpathmoveto{\pgfqpoint{1.422538in}{2.094898in}}%
\pgfpathlineto{\pgfqpoint{1.422538in}{2.094898in}}%
\pgfpathlineto{\pgfqpoint{1.422538in}{2.094898in}}%
\pgfpathclose%
\pgfpathmoveto{\pgfqpoint{1.525021in}{2.082219in}}%
\pgfpathlineto{\pgfqpoint{1.525021in}{2.121266in}}%
\pgfpathlineto{\pgfqpoint{1.536876in}{2.121266in}}%
\pgfpathlineto{\pgfqpoint{1.536876in}{2.021009in}}%
\pgfpathlineto{\pgfqpoint{1.525021in}{2.021009in}}%
\pgfpathlineto{\pgfqpoint{1.525021in}{2.031832in}}%
\pgfpathquadraticcurveto{\pgfqpoint{1.521290in}{2.025400in}}{\pgfqpoint{1.515579in}{2.022266in}}%
\pgfpathquadraticcurveto{\pgfqpoint{1.509889in}{2.019133in}}{\pgfqpoint{1.501890in}{2.019133in}}%
\pgfpathquadraticcurveto{\pgfqpoint{1.488819in}{2.019133in}}{\pgfqpoint{1.480593in}{2.029565in}}%
\pgfpathquadraticcurveto{\pgfqpoint{1.472388in}{2.040017in}}{\pgfqpoint{1.472388in}{2.057026in}}%
\pgfpathquadraticcurveto{\pgfqpoint{1.472388in}{2.074034in}}{\pgfqpoint{1.480593in}{2.084466in}}%
\pgfpathquadraticcurveto{\pgfqpoint{1.488819in}{2.094898in}}{\pgfqpoint{1.501890in}{2.094898in}}%
\pgfpathquadraticcurveto{\pgfqpoint{1.509889in}{2.094898in}}{\pgfqpoint{1.515579in}{2.091764in}}%
\pgfpathquadraticcurveto{\pgfqpoint{1.521290in}{2.088651in}}{\pgfqpoint{1.525021in}{2.082219in}}%
\pgfpathlineto{\pgfqpoint{1.525021in}{2.082219in}}%
\pgfpathclose%
\pgfpathmoveto{\pgfqpoint{1.484634in}{2.057026in}}%
\pgfpathquadraticcurveto{\pgfqpoint{1.484634in}{2.043955in}}{\pgfqpoint{1.490015in}{2.036512in}}%
\pgfpathquadraticcurveto{\pgfqpoint{1.495396in}{2.029070in}}{\pgfqpoint{1.504797in}{2.029070in}}%
\pgfpathquadraticcurveto{\pgfqpoint{1.514198in}{2.029070in}}{\pgfqpoint{1.519599in}{2.036512in}}%
\pgfpathquadraticcurveto{\pgfqpoint{1.525021in}{2.043955in}}{\pgfqpoint{1.525021in}{2.057026in}}%
\pgfpathquadraticcurveto{\pgfqpoint{1.525021in}{2.070096in}}{\pgfqpoint{1.519599in}{2.077539in}}%
\pgfpathquadraticcurveto{\pgfqpoint{1.514198in}{2.084981in}}{\pgfqpoint{1.504797in}{2.084981in}}%
\pgfpathquadraticcurveto{\pgfqpoint{1.495396in}{2.084981in}}{\pgfqpoint{1.490015in}{2.077539in}}%
\pgfpathquadraticcurveto{\pgfqpoint{1.484634in}{2.070096in}}{\pgfqpoint{1.484634in}{2.057026in}}%
\pgfpathlineto{\pgfqpoint{1.484634in}{2.057026in}}%
\pgfpathclose%
\pgfpathmoveto{\pgfqpoint{1.591323in}{2.014308in}}%
\pgfpathquadraticcurveto{\pgfqpoint{1.586314in}{2.001423in}}{\pgfqpoint{1.581531in}{1.997506in}}%
\pgfpathquadraticcurveto{\pgfqpoint{1.576768in}{1.993568in}}{\pgfqpoint{1.568790in}{1.993568in}}%
\pgfpathlineto{\pgfqpoint{1.559306in}{1.993568in}}%
\pgfpathlineto{\pgfqpoint{1.559306in}{2.003485in}}%
\pgfpathlineto{\pgfqpoint{1.566275in}{2.003485in}}%
\pgfpathquadraticcurveto{\pgfqpoint{1.571161in}{2.003485in}}{\pgfqpoint{1.573861in}{2.005815in}}%
\pgfpathquadraticcurveto{\pgfqpoint{1.576583in}{2.008124in}}{\pgfqpoint{1.579861in}{2.016762in}}%
\pgfpathlineto{\pgfqpoint{1.581984in}{2.022163in}}%
\pgfpathlineto{\pgfqpoint{1.552812in}{2.093166in}}%
\pgfpathlineto{\pgfqpoint{1.565367in}{2.093166in}}%
\pgfpathlineto{\pgfqpoint{1.587922in}{2.036739in}}%
\pgfpathlineto{\pgfqpoint{1.610476in}{2.093166in}}%
\pgfpathlineto{\pgfqpoint{1.623031in}{2.093166in}}%
\pgfpathlineto{\pgfqpoint{1.591323in}{2.014308in}}%
\pgfpathlineto{\pgfqpoint{1.591323in}{2.014308in}}%
\pgfpathclose%
\pgfpathmoveto{\pgfqpoint{1.694219in}{2.031956in}}%
\pgfpathlineto{\pgfqpoint{1.739637in}{2.031956in}}%
\pgfpathlineto{\pgfqpoint{1.739637in}{2.021009in}}%
\pgfpathlineto{\pgfqpoint{1.678572in}{2.021009in}}%
\pgfpathlineto{\pgfqpoint{1.678572in}{2.031956in}}%
\pgfpathquadraticcurveto{\pgfqpoint{1.685973in}{2.039625in}}{\pgfqpoint{1.698755in}{2.052531in}}%
\pgfpathquadraticcurveto{\pgfqpoint{1.711558in}{2.065458in}}{\pgfqpoint{1.714836in}{2.069210in}}%
\pgfpathquadraticcurveto{\pgfqpoint{1.721082in}{2.076219in}}{\pgfqpoint{1.723556in}{2.081085in}}%
\pgfpathquadraticcurveto{\pgfqpoint{1.726051in}{2.085950in}}{\pgfqpoint{1.726051in}{2.090651in}}%
\pgfpathquadraticcurveto{\pgfqpoint{1.726051in}{2.098320in}}{\pgfqpoint{1.720670in}{2.103144in}}%
\pgfpathquadraticcurveto{\pgfqpoint{1.715289in}{2.107989in}}{\pgfqpoint{1.706651in}{2.107989in}}%
\pgfpathquadraticcurveto{\pgfqpoint{1.700528in}{2.107989in}}{\pgfqpoint{1.693725in}{2.105866in}}%
\pgfpathquadraticcurveto{\pgfqpoint{1.686942in}{2.103742in}}{\pgfqpoint{1.679211in}{2.099413in}}%
\pgfpathlineto{\pgfqpoint{1.679211in}{2.112566in}}%
\pgfpathquadraticcurveto{\pgfqpoint{1.687066in}{2.115720in}}{\pgfqpoint{1.693890in}{2.117328in}}%
\pgfpathquadraticcurveto{\pgfqpoint{1.700734in}{2.118936in}}{\pgfqpoint{1.706404in}{2.118936in}}%
\pgfpathquadraticcurveto{\pgfqpoint{1.721350in}{2.118936in}}{\pgfqpoint{1.730236in}{2.111453in}}%
\pgfpathquadraticcurveto{\pgfqpoint{1.739122in}{2.103989in}}{\pgfqpoint{1.739122in}{2.091496in}}%
\pgfpathquadraticcurveto{\pgfqpoint{1.739122in}{2.085558in}}{\pgfqpoint{1.736895in}{2.080239in}}%
\pgfpathquadraticcurveto{\pgfqpoint{1.734689in}{2.074941in}}{\pgfqpoint{1.728814in}{2.067725in}}%
\pgfpathquadraticcurveto{\pgfqpoint{1.727206in}{2.065849in}}{\pgfqpoint{1.718567in}{2.056922in}}%
\pgfpathquadraticcurveto{\pgfqpoint{1.709950in}{2.047996in}}{\pgfqpoint{1.694219in}{2.031956in}}%
\pgfpathlineto{\pgfqpoint{1.694219in}{2.031956in}}%
\pgfpathclose%
\pgfpathmoveto{\pgfqpoint{1.835689in}{2.121122in}}%
\pgfpathquadraticcurveto{\pgfqpoint{1.827071in}{2.106319in}}{\pgfqpoint{1.822865in}{2.091805in}}%
\pgfpathquadraticcurveto{\pgfqpoint{1.818680in}{2.077312in}}{\pgfqpoint{1.818680in}{2.062427in}}%
\pgfpathquadraticcurveto{\pgfqpoint{1.818680in}{2.047563in}}{\pgfqpoint{1.822906in}{2.032966in}}%
\pgfpathquadraticcurveto{\pgfqpoint{1.827133in}{2.018370in}}{\pgfqpoint{1.835689in}{2.003609in}}%
\pgfpathlineto{\pgfqpoint{1.825380in}{2.003609in}}%
\pgfpathquadraticcurveto{\pgfqpoint{1.815732in}{2.018762in}}{\pgfqpoint{1.810928in}{2.033379in}}%
\pgfpathquadraticcurveto{\pgfqpoint{1.806125in}{2.047996in}}{\pgfqpoint{1.806125in}{2.062427in}}%
\pgfpathquadraticcurveto{\pgfqpoint{1.806125in}{2.076797in}}{\pgfqpoint{1.810887in}{2.091352in}}%
\pgfpathquadraticcurveto{\pgfqpoint{1.815670in}{2.105927in}}{\pgfqpoint{1.825380in}{2.121122in}}%
\pgfpathlineto{\pgfqpoint{1.835689in}{2.121122in}}%
\pgfpathlineto{\pgfqpoint{1.835689in}{2.121122in}}%
\pgfpathclose%
\pgfpathmoveto{\pgfqpoint{1.871582in}{2.031956in}}%
\pgfpathlineto{\pgfqpoint{1.916999in}{2.031956in}}%
\pgfpathlineto{\pgfqpoint{1.916999in}{2.021009in}}%
\pgfpathlineto{\pgfqpoint{1.855934in}{2.021009in}}%
\pgfpathlineto{\pgfqpoint{1.855934in}{2.031956in}}%
\pgfpathquadraticcurveto{\pgfqpoint{1.863335in}{2.039625in}}{\pgfqpoint{1.876117in}{2.052531in}}%
\pgfpathquadraticcurveto{\pgfqpoint{1.888920in}{2.065458in}}{\pgfqpoint{1.892198in}{2.069210in}}%
\pgfpathquadraticcurveto{\pgfqpoint{1.898445in}{2.076219in}}{\pgfqpoint{1.900919in}{2.081085in}}%
\pgfpathquadraticcurveto{\pgfqpoint{1.903413in}{2.085950in}}{\pgfqpoint{1.903413in}{2.090651in}}%
\pgfpathquadraticcurveto{\pgfqpoint{1.903413in}{2.098320in}}{\pgfqpoint{1.898032in}{2.103144in}}%
\pgfpathquadraticcurveto{\pgfqpoint{1.892651in}{2.107989in}}{\pgfqpoint{1.884013in}{2.107989in}}%
\pgfpathquadraticcurveto{\pgfqpoint{1.877890in}{2.107989in}}{\pgfqpoint{1.871087in}{2.105866in}}%
\pgfpathquadraticcurveto{\pgfqpoint{1.864304in}{2.103742in}}{\pgfqpoint{1.856573in}{2.099413in}}%
\pgfpathlineto{\pgfqpoint{1.856573in}{2.112566in}}%
\pgfpathquadraticcurveto{\pgfqpoint{1.864428in}{2.115720in}}{\pgfqpoint{1.871252in}{2.117328in}}%
\pgfpathquadraticcurveto{\pgfqpoint{1.878096in}{2.118936in}}{\pgfqpoint{1.883766in}{2.118936in}}%
\pgfpathquadraticcurveto{\pgfqpoint{1.898713in}{2.118936in}}{\pgfqpoint{1.907598in}{2.111453in}}%
\pgfpathquadraticcurveto{\pgfqpoint{1.916484in}{2.103989in}}{\pgfqpoint{1.916484in}{2.091496in}}%
\pgfpathquadraticcurveto{\pgfqpoint{1.916484in}{2.085558in}}{\pgfqpoint{1.914257in}{2.080239in}}%
\pgfpathquadraticcurveto{\pgfqpoint{1.912051in}{2.074941in}}{\pgfqpoint{1.906176in}{2.067725in}}%
\pgfpathquadraticcurveto{\pgfqpoint{1.904568in}{2.065849in}}{\pgfqpoint{1.895929in}{2.056922in}}%
\pgfpathquadraticcurveto{\pgfqpoint{1.887312in}{2.047996in}}{\pgfqpoint{1.871582in}{2.031956in}}%
\pgfpathlineto{\pgfqpoint{1.871582in}{2.031956in}}%
\pgfpathclose%
\pgfpathmoveto{\pgfqpoint{1.972148in}{2.108628in}}%
\pgfpathquadraticcurveto{\pgfqpoint{1.962108in}{2.108628in}}{\pgfqpoint{1.957036in}{2.098732in}}%
\pgfpathquadraticcurveto{\pgfqpoint{1.951985in}{2.088857in}}{\pgfqpoint{1.951985in}{2.069004in}}%
\pgfpathquadraticcurveto{\pgfqpoint{1.951985in}{2.049233in}}{\pgfqpoint{1.957036in}{2.039337in}}%
\pgfpathquadraticcurveto{\pgfqpoint{1.962108in}{2.029441in}}{\pgfqpoint{1.972148in}{2.029441in}}%
\pgfpathquadraticcurveto{\pgfqpoint{1.982271in}{2.029441in}}{\pgfqpoint{1.987322in}{2.039337in}}%
\pgfpathquadraticcurveto{\pgfqpoint{1.992393in}{2.049233in}}{\pgfqpoint{1.992393in}{2.069004in}}%
\pgfpathquadraticcurveto{\pgfqpoint{1.992393in}{2.088857in}}{\pgfqpoint{1.987322in}{2.098732in}}%
\pgfpathquadraticcurveto{\pgfqpoint{1.982271in}{2.108628in}}{\pgfqpoint{1.972148in}{2.108628in}}%
\pgfpathlineto{\pgfqpoint{1.972148in}{2.108628in}}%
\pgfpathclose%
\pgfpathmoveto{\pgfqpoint{1.972148in}{2.118936in}}%
\pgfpathquadraticcurveto{\pgfqpoint{1.988332in}{2.118936in}}{\pgfqpoint{1.996867in}{2.106134in}}%
\pgfpathquadraticcurveto{\pgfqpoint{2.005402in}{2.093351in}}{\pgfqpoint{2.005402in}{2.069004in}}%
\pgfpathquadraticcurveto{\pgfqpoint{2.005402in}{2.044718in}}{\pgfqpoint{1.996867in}{2.031915in}}%
\pgfpathquadraticcurveto{\pgfqpoint{1.988332in}{2.019133in}}{\pgfqpoint{1.972148in}{2.019133in}}%
\pgfpathquadraticcurveto{\pgfqpoint{1.955985in}{2.019133in}}{\pgfqpoint{1.947450in}{2.031915in}}%
\pgfpathquadraticcurveto{\pgfqpoint{1.938914in}{2.044718in}}{\pgfqpoint{1.938914in}{2.069004in}}%
\pgfpathquadraticcurveto{\pgfqpoint{1.938914in}{2.093351in}}{\pgfqpoint{1.947450in}{2.106134in}}%
\pgfpathquadraticcurveto{\pgfqpoint{1.955985in}{2.118936in}}{\pgfqpoint{1.972148in}{2.118936in}}%
\pgfpathlineto{\pgfqpoint{1.972148in}{2.118936in}}%
\pgfpathclose%
\pgfpathmoveto{\pgfqpoint{2.030533in}{2.031956in}}%
\pgfpathlineto{\pgfqpoint{2.051789in}{2.031956in}}%
\pgfpathlineto{\pgfqpoint{2.051789in}{2.105350in}}%
\pgfpathlineto{\pgfqpoint{2.028657in}{2.100711in}}%
\pgfpathlineto{\pgfqpoint{2.028657in}{2.112566in}}%
\pgfpathlineto{\pgfqpoint{2.051665in}{2.117205in}}%
\pgfpathlineto{\pgfqpoint{2.064674in}{2.117205in}}%
\pgfpathlineto{\pgfqpoint{2.064674in}{2.031956in}}%
\pgfpathlineto{\pgfqpoint{2.085929in}{2.031956in}}%
\pgfpathlineto{\pgfqpoint{2.085929in}{2.021009in}}%
\pgfpathlineto{\pgfqpoint{2.030533in}{2.021009in}}%
\pgfpathlineto{\pgfqpoint{2.030533in}{2.031956in}}%
\pgfpathlineto{\pgfqpoint{2.030533in}{2.031956in}}%
\pgfpathclose%
\pgfpathmoveto{\pgfqpoint{2.108937in}{2.117205in}}%
\pgfpathlineto{\pgfqpoint{2.170786in}{2.117205in}}%
\pgfpathlineto{\pgfqpoint{2.170786in}{2.111659in}}%
\pgfpathlineto{\pgfqpoint{2.135862in}{2.021009in}}%
\pgfpathlineto{\pgfqpoint{2.122276in}{2.021009in}}%
\pgfpathlineto{\pgfqpoint{2.155138in}{2.106237in}}%
\pgfpathlineto{\pgfqpoint{2.108937in}{2.106237in}}%
\pgfpathlineto{\pgfqpoint{2.108937in}{2.117205in}}%
\pgfpathlineto{\pgfqpoint{2.108937in}{2.117205in}}%
\pgfpathclose%
\pgfpathmoveto{\pgfqpoint{2.192640in}{2.121122in}}%
\pgfpathlineto{\pgfqpoint{2.202948in}{2.121122in}}%
\pgfpathquadraticcurveto{\pgfqpoint{2.212596in}{2.105927in}}{\pgfqpoint{2.217400in}{2.091352in}}%
\pgfpathquadraticcurveto{\pgfqpoint{2.222203in}{2.076797in}}{\pgfqpoint{2.222203in}{2.062427in}}%
\pgfpathquadraticcurveto{\pgfqpoint{2.222203in}{2.047996in}}{\pgfqpoint{2.217400in}{2.033379in}}%
\pgfpathquadraticcurveto{\pgfqpoint{2.212596in}{2.018762in}}{\pgfqpoint{2.202948in}{2.003609in}}%
\pgfpathlineto{\pgfqpoint{2.192640in}{2.003609in}}%
\pgfpathquadraticcurveto{\pgfqpoint{2.201195in}{2.018370in}}{\pgfqpoint{2.205422in}{2.032966in}}%
\pgfpathquadraticcurveto{\pgfqpoint{2.209648in}{2.047563in}}{\pgfqpoint{2.209648in}{2.062427in}}%
\pgfpathquadraticcurveto{\pgfqpoint{2.209648in}{2.077312in}}{\pgfqpoint{2.205422in}{2.091805in}}%
\pgfpathquadraticcurveto{\pgfqpoint{2.201195in}{2.106319in}}{\pgfqpoint{2.192640in}{2.121122in}}%
\pgfpathlineto{\pgfqpoint{2.192640in}{2.121122in}}%
\pgfpathclose%
\pgfusepath{fill}%
\end{pgfscope}%
\begin{pgfscope}%
\pgfsetbuttcap%
\pgfsetroundjoin%
\definecolor{currentfill}{rgb}{0.090196,0.168627,0.227451}%
\pgfsetfillcolor{currentfill}%
\pgfsetlinewidth{0.000000pt}%
\definecolor{currentstroke}{rgb}{0.090196,0.168627,0.227451}%
\pgfsetstrokecolor{currentstroke}%
\pgfsetdash{}{0pt}%
\pgfpathmoveto{\pgfqpoint{1.316611in}{1.937410in}}%
\pgfpathlineto{\pgfqpoint{1.316611in}{1.924711in}}%
\pgfpathquadraticcurveto{\pgfqpoint{1.309210in}{1.928257in}}{\pgfqpoint{1.302633in}{1.929988in}}%
\pgfpathquadraticcurveto{\pgfqpoint{1.296056in}{1.931741in}}{\pgfqpoint{1.289933in}{1.931741in}}%
\pgfpathquadraticcurveto{\pgfqpoint{1.279316in}{1.931741in}}{\pgfqpoint{1.273543in}{1.927617in}}%
\pgfpathquadraticcurveto{\pgfqpoint{1.267771in}{1.923494in}}{\pgfqpoint{1.267771in}{1.915887in}}%
\pgfpathquadraticcurveto{\pgfqpoint{1.267771in}{1.909496in}}{\pgfqpoint{1.271605in}{1.906238in}}%
\pgfpathquadraticcurveto{\pgfqpoint{1.275440in}{1.903002in}}{\pgfqpoint{1.286140in}{1.901002in}}%
\pgfpathlineto{\pgfqpoint{1.293995in}{1.899394in}}%
\pgfpathquadraticcurveto{\pgfqpoint{1.308550in}{1.896611in}}{\pgfqpoint{1.315477in}{1.889622in}}%
\pgfpathquadraticcurveto{\pgfqpoint{1.322404in}{1.882633in}}{\pgfqpoint{1.322404in}{1.870923in}}%
\pgfpathquadraticcurveto{\pgfqpoint{1.322404in}{1.856924in}}{\pgfqpoint{1.313024in}{1.849708in}}%
\pgfpathquadraticcurveto{\pgfqpoint{1.303664in}{1.842493in}}{\pgfqpoint{1.285563in}{1.842493in}}%
\pgfpathquadraticcurveto{\pgfqpoint{1.278739in}{1.842493in}}{\pgfqpoint{1.271028in}{1.844039in}}%
\pgfpathquadraticcurveto{\pgfqpoint{1.263338in}{1.845585in}}{\pgfqpoint{1.255092in}{1.848616in}}%
\pgfpathlineto{\pgfqpoint{1.255092in}{1.862016in}}%
\pgfpathquadraticcurveto{\pgfqpoint{1.263008in}{1.857584in}}{\pgfqpoint{1.270616in}{1.855316in}}%
\pgfpathquadraticcurveto{\pgfqpoint{1.278223in}{1.853069in}}{\pgfqpoint{1.285563in}{1.853069in}}%
\pgfpathquadraticcurveto{\pgfqpoint{1.296696in}{1.853069in}}{\pgfqpoint{1.302757in}{1.857440in}}%
\pgfpathquadraticcurveto{\pgfqpoint{1.308818in}{1.861831in}}{\pgfqpoint{1.308818in}{1.869954in}}%
\pgfpathquadraticcurveto{\pgfqpoint{1.308818in}{1.877025in}}{\pgfqpoint{1.304468in}{1.881025in}}%
\pgfpathquadraticcurveto{\pgfqpoint{1.300118in}{1.885024in}}{\pgfqpoint{1.290201in}{1.887024in}}%
\pgfpathlineto{\pgfqpoint{1.282264in}{1.888570in}}%
\pgfpathquadraticcurveto{\pgfqpoint{1.267709in}{1.891456in}}{\pgfqpoint{1.261194in}{1.897641in}}%
\pgfpathquadraticcurveto{\pgfqpoint{1.254700in}{1.903826in}}{\pgfqpoint{1.254700in}{1.914856in}}%
\pgfpathquadraticcurveto{\pgfqpoint{1.254700in}{1.927617in}}{\pgfqpoint{1.263689in}{1.934957in}}%
\pgfpathquadraticcurveto{\pgfqpoint{1.272678in}{1.942296in}}{\pgfqpoint{1.288449in}{1.942296in}}%
\pgfpathquadraticcurveto{\pgfqpoint{1.295232in}{1.942296in}}{\pgfqpoint{1.302241in}{1.941059in}}%
\pgfpathquadraticcurveto{\pgfqpoint{1.309271in}{1.939843in}}{\pgfqpoint{1.316611in}{1.937410in}}%
\pgfpathlineto{\pgfqpoint{1.316611in}{1.937410in}}%
\pgfpathclose%
\pgfpathmoveto{\pgfqpoint{1.353926in}{1.937019in}}%
\pgfpathlineto{\pgfqpoint{1.353926in}{1.916526in}}%
\pgfpathlineto{\pgfqpoint{1.378336in}{1.916526in}}%
\pgfpathlineto{\pgfqpoint{1.378336in}{1.907310in}}%
\pgfpathlineto{\pgfqpoint{1.353926in}{1.907310in}}%
\pgfpathlineto{\pgfqpoint{1.353926in}{1.868139in}}%
\pgfpathquadraticcurveto{\pgfqpoint{1.353926in}{1.859316in}}{\pgfqpoint{1.356339in}{1.856800in}}%
\pgfpathquadraticcurveto{\pgfqpoint{1.358751in}{1.854285in}}{\pgfqpoint{1.366173in}{1.854285in}}%
\pgfpathlineto{\pgfqpoint{1.378336in}{1.854285in}}%
\pgfpathlineto{\pgfqpoint{1.378336in}{1.844369in}}%
\pgfpathlineto{\pgfqpoint{1.366173in}{1.844369in}}%
\pgfpathquadraticcurveto{\pgfqpoint{1.352442in}{1.844369in}}{\pgfqpoint{1.347226in}{1.849482in}}%
\pgfpathquadraticcurveto{\pgfqpoint{1.342010in}{1.854615in}}{\pgfqpoint{1.342010in}{1.868139in}}%
\pgfpathlineto{\pgfqpoint{1.342010in}{1.907310in}}%
\pgfpathlineto{\pgfqpoint{1.333310in}{1.907310in}}%
\pgfpathlineto{\pgfqpoint{1.333310in}{1.916526in}}%
\pgfpathlineto{\pgfqpoint{1.342010in}{1.916526in}}%
\pgfpathlineto{\pgfqpoint{1.342010in}{1.937019in}}%
\pgfpathlineto{\pgfqpoint{1.353926in}{1.937019in}}%
\pgfpathlineto{\pgfqpoint{1.353926in}{1.937019in}}%
\pgfpathclose%
\pgfpathmoveto{\pgfqpoint{1.392706in}{1.872840in}}%
\pgfpathlineto{\pgfqpoint{1.392706in}{1.916526in}}%
\pgfpathlineto{\pgfqpoint{1.404560in}{1.916526in}}%
\pgfpathlineto{\pgfqpoint{1.404560in}{1.873293in}}%
\pgfpathquadraticcurveto{\pgfqpoint{1.404560in}{1.863047in}}{\pgfqpoint{1.408539in}{1.857914in}}%
\pgfpathquadraticcurveto{\pgfqpoint{1.412539in}{1.852801in}}{\pgfqpoint{1.420538in}{1.852801in}}%
\pgfpathquadraticcurveto{\pgfqpoint{1.430124in}{1.852801in}}{\pgfqpoint{1.435691in}{1.858924in}}%
\pgfpathquadraticcurveto{\pgfqpoint{1.441278in}{1.865047in}}{\pgfqpoint{1.441278in}{1.875623in}}%
\pgfpathlineto{\pgfqpoint{1.441278in}{1.916526in}}%
\pgfpathlineto{\pgfqpoint{1.453132in}{1.916526in}}%
\pgfpathlineto{\pgfqpoint{1.453132in}{1.844369in}}%
\pgfpathlineto{\pgfqpoint{1.441278in}{1.844369in}}%
\pgfpathlineto{\pgfqpoint{1.441278in}{1.855460in}}%
\pgfpathquadraticcurveto{\pgfqpoint{1.436969in}{1.848884in}}{\pgfqpoint{1.431258in}{1.845688in}}%
\pgfpathquadraticcurveto{\pgfqpoint{1.425568in}{1.842493in}}{\pgfqpoint{1.418023in}{1.842493in}}%
\pgfpathquadraticcurveto{\pgfqpoint{1.405591in}{1.842493in}}{\pgfqpoint{1.399138in}{1.850224in}}%
\pgfpathquadraticcurveto{\pgfqpoint{1.392706in}{1.857955in}}{\pgfqpoint{1.392706in}{1.872840in}}%
\pgfpathlineto{\pgfqpoint{1.392706in}{1.872840in}}%
\pgfpathclose%
\pgfpathmoveto{\pgfqpoint{1.422538in}{1.918258in}}%
\pgfpathlineto{\pgfqpoint{1.422538in}{1.918258in}}%
\pgfpathlineto{\pgfqpoint{1.422538in}{1.918258in}}%
\pgfpathclose%
\pgfpathmoveto{\pgfqpoint{1.525021in}{1.905579in}}%
\pgfpathlineto{\pgfqpoint{1.525021in}{1.944626in}}%
\pgfpathlineto{\pgfqpoint{1.536876in}{1.944626in}}%
\pgfpathlineto{\pgfqpoint{1.536876in}{1.844369in}}%
\pgfpathlineto{\pgfqpoint{1.525021in}{1.844369in}}%
\pgfpathlineto{\pgfqpoint{1.525021in}{1.855192in}}%
\pgfpathquadraticcurveto{\pgfqpoint{1.521290in}{1.848760in}}{\pgfqpoint{1.515579in}{1.845626in}}%
\pgfpathquadraticcurveto{\pgfqpoint{1.509889in}{1.842493in}}{\pgfqpoint{1.501890in}{1.842493in}}%
\pgfpathquadraticcurveto{\pgfqpoint{1.488819in}{1.842493in}}{\pgfqpoint{1.480593in}{1.852925in}}%
\pgfpathquadraticcurveto{\pgfqpoint{1.472388in}{1.863377in}}{\pgfqpoint{1.472388in}{1.880386in}}%
\pgfpathquadraticcurveto{\pgfqpoint{1.472388in}{1.897394in}}{\pgfqpoint{1.480593in}{1.907826in}}%
\pgfpathquadraticcurveto{\pgfqpoint{1.488819in}{1.918258in}}{\pgfqpoint{1.501890in}{1.918258in}}%
\pgfpathquadraticcurveto{\pgfqpoint{1.509889in}{1.918258in}}{\pgfqpoint{1.515579in}{1.915124in}}%
\pgfpathquadraticcurveto{\pgfqpoint{1.521290in}{1.912011in}}{\pgfqpoint{1.525021in}{1.905579in}}%
\pgfpathlineto{\pgfqpoint{1.525021in}{1.905579in}}%
\pgfpathclose%
\pgfpathmoveto{\pgfqpoint{1.484634in}{1.880386in}}%
\pgfpathquadraticcurveto{\pgfqpoint{1.484634in}{1.867315in}}{\pgfqpoint{1.490015in}{1.859872in}}%
\pgfpathquadraticcurveto{\pgfqpoint{1.495396in}{1.852430in}}{\pgfqpoint{1.504797in}{1.852430in}}%
\pgfpathquadraticcurveto{\pgfqpoint{1.514198in}{1.852430in}}{\pgfqpoint{1.519599in}{1.859872in}}%
\pgfpathquadraticcurveto{\pgfqpoint{1.525021in}{1.867315in}}{\pgfqpoint{1.525021in}{1.880386in}}%
\pgfpathquadraticcurveto{\pgfqpoint{1.525021in}{1.893456in}}{\pgfqpoint{1.519599in}{1.900899in}}%
\pgfpathquadraticcurveto{\pgfqpoint{1.514198in}{1.908341in}}{\pgfqpoint{1.504797in}{1.908341in}}%
\pgfpathquadraticcurveto{\pgfqpoint{1.495396in}{1.908341in}}{\pgfqpoint{1.490015in}{1.900899in}}%
\pgfpathquadraticcurveto{\pgfqpoint{1.484634in}{1.893456in}}{\pgfqpoint{1.484634in}{1.880386in}}%
\pgfpathlineto{\pgfqpoint{1.484634in}{1.880386in}}%
\pgfpathclose%
\pgfpathmoveto{\pgfqpoint{1.591323in}{1.837668in}}%
\pgfpathquadraticcurveto{\pgfqpoint{1.586314in}{1.824783in}}{\pgfqpoint{1.581531in}{1.820866in}}%
\pgfpathquadraticcurveto{\pgfqpoint{1.576768in}{1.816928in}}{\pgfqpoint{1.568790in}{1.816928in}}%
\pgfpathlineto{\pgfqpoint{1.559306in}{1.816928in}}%
\pgfpathlineto{\pgfqpoint{1.559306in}{1.826845in}}%
\pgfpathlineto{\pgfqpoint{1.566275in}{1.826845in}}%
\pgfpathquadraticcurveto{\pgfqpoint{1.571161in}{1.826845in}}{\pgfqpoint{1.573861in}{1.829175in}}%
\pgfpathquadraticcurveto{\pgfqpoint{1.576583in}{1.831484in}}{\pgfqpoint{1.579861in}{1.840122in}}%
\pgfpathlineto{\pgfqpoint{1.581984in}{1.845523in}}%
\pgfpathlineto{\pgfqpoint{1.552812in}{1.916526in}}%
\pgfpathlineto{\pgfqpoint{1.565367in}{1.916526in}}%
\pgfpathlineto{\pgfqpoint{1.587922in}{1.860099in}}%
\pgfpathlineto{\pgfqpoint{1.610476in}{1.916526in}}%
\pgfpathlineto{\pgfqpoint{1.623031in}{1.916526in}}%
\pgfpathlineto{\pgfqpoint{1.591323in}{1.837668in}}%
\pgfpathlineto{\pgfqpoint{1.591323in}{1.837668in}}%
\pgfpathclose%
\pgfpathmoveto{\pgfqpoint{1.718773in}{1.929226in}}%
\pgfpathlineto{\pgfqpoint{1.685911in}{1.877870in}}%
\pgfpathlineto{\pgfqpoint{1.718773in}{1.877870in}}%
\pgfpathlineto{\pgfqpoint{1.718773in}{1.929226in}}%
\pgfpathlineto{\pgfqpoint{1.718773in}{1.929226in}}%
\pgfpathclose%
\pgfpathmoveto{\pgfqpoint{1.715351in}{1.940565in}}%
\pgfpathlineto{\pgfqpoint{1.731720in}{1.940565in}}%
\pgfpathlineto{\pgfqpoint{1.731720in}{1.877870in}}%
\pgfpathlineto{\pgfqpoint{1.745451in}{1.877870in}}%
\pgfpathlineto{\pgfqpoint{1.745451in}{1.867047in}}%
\pgfpathlineto{\pgfqpoint{1.731720in}{1.867047in}}%
\pgfpathlineto{\pgfqpoint{1.731720in}{1.844369in}}%
\pgfpathlineto{\pgfqpoint{1.718773in}{1.844369in}}%
\pgfpathlineto{\pgfqpoint{1.718773in}{1.867047in}}%
\pgfpathlineto{\pgfqpoint{1.675355in}{1.867047in}}%
\pgfpathlineto{\pgfqpoint{1.675355in}{1.879602in}}%
\pgfpathlineto{\pgfqpoint{1.715351in}{1.940565in}}%
\pgfpathlineto{\pgfqpoint{1.715351in}{1.940565in}}%
\pgfpathclose%
\pgfpathmoveto{\pgfqpoint{1.835689in}{1.944482in}}%
\pgfpathquadraticcurveto{\pgfqpoint{1.827071in}{1.929679in}}{\pgfqpoint{1.822865in}{1.915165in}}%
\pgfpathquadraticcurveto{\pgfqpoint{1.818680in}{1.900672in}}{\pgfqpoint{1.818680in}{1.885787in}}%
\pgfpathquadraticcurveto{\pgfqpoint{1.818680in}{1.870923in}}{\pgfqpoint{1.822906in}{1.856326in}}%
\pgfpathquadraticcurveto{\pgfqpoint{1.827133in}{1.841730in}}{\pgfqpoint{1.835689in}{1.826969in}}%
\pgfpathlineto{\pgfqpoint{1.825380in}{1.826969in}}%
\pgfpathquadraticcurveto{\pgfqpoint{1.815732in}{1.842122in}}{\pgfqpoint{1.810928in}{1.856739in}}%
\pgfpathquadraticcurveto{\pgfqpoint{1.806125in}{1.871356in}}{\pgfqpoint{1.806125in}{1.885787in}}%
\pgfpathquadraticcurveto{\pgfqpoint{1.806125in}{1.900157in}}{\pgfqpoint{1.810887in}{1.914712in}}%
\pgfpathquadraticcurveto{\pgfqpoint{1.815670in}{1.929287in}}{\pgfqpoint{1.825380in}{1.944482in}}%
\pgfpathlineto{\pgfqpoint{1.835689in}{1.944482in}}%
\pgfpathlineto{\pgfqpoint{1.835689in}{1.944482in}}%
\pgfpathclose%
\pgfpathmoveto{\pgfqpoint{1.871582in}{1.855316in}}%
\pgfpathlineto{\pgfqpoint{1.916999in}{1.855316in}}%
\pgfpathlineto{\pgfqpoint{1.916999in}{1.844369in}}%
\pgfpathlineto{\pgfqpoint{1.855934in}{1.844369in}}%
\pgfpathlineto{\pgfqpoint{1.855934in}{1.855316in}}%
\pgfpathquadraticcurveto{\pgfqpoint{1.863335in}{1.862985in}}{\pgfqpoint{1.876117in}{1.875891in}}%
\pgfpathquadraticcurveto{\pgfqpoint{1.888920in}{1.888818in}}{\pgfqpoint{1.892198in}{1.892570in}}%
\pgfpathquadraticcurveto{\pgfqpoint{1.898445in}{1.899579in}}{\pgfqpoint{1.900919in}{1.904445in}}%
\pgfpathquadraticcurveto{\pgfqpoint{1.903413in}{1.909310in}}{\pgfqpoint{1.903413in}{1.914011in}}%
\pgfpathquadraticcurveto{\pgfqpoint{1.903413in}{1.921680in}}{\pgfqpoint{1.898032in}{1.926504in}}%
\pgfpathquadraticcurveto{\pgfqpoint{1.892651in}{1.931349in}}{\pgfqpoint{1.884013in}{1.931349in}}%
\pgfpathquadraticcurveto{\pgfqpoint{1.877890in}{1.931349in}}{\pgfqpoint{1.871087in}{1.929226in}}%
\pgfpathquadraticcurveto{\pgfqpoint{1.864304in}{1.927102in}}{\pgfqpoint{1.856573in}{1.922773in}}%
\pgfpathlineto{\pgfqpoint{1.856573in}{1.935926in}}%
\pgfpathquadraticcurveto{\pgfqpoint{1.864428in}{1.939080in}}{\pgfqpoint{1.871252in}{1.940688in}}%
\pgfpathquadraticcurveto{\pgfqpoint{1.878096in}{1.942296in}}{\pgfqpoint{1.883766in}{1.942296in}}%
\pgfpathquadraticcurveto{\pgfqpoint{1.898713in}{1.942296in}}{\pgfqpoint{1.907598in}{1.934813in}}%
\pgfpathquadraticcurveto{\pgfqpoint{1.916484in}{1.927349in}}{\pgfqpoint{1.916484in}{1.914856in}}%
\pgfpathquadraticcurveto{\pgfqpoint{1.916484in}{1.908918in}}{\pgfqpoint{1.914257in}{1.903599in}}%
\pgfpathquadraticcurveto{\pgfqpoint{1.912051in}{1.898301in}}{\pgfqpoint{1.906176in}{1.891085in}}%
\pgfpathquadraticcurveto{\pgfqpoint{1.904568in}{1.889209in}}{\pgfqpoint{1.895929in}{1.880282in}}%
\pgfpathquadraticcurveto{\pgfqpoint{1.887312in}{1.871356in}}{\pgfqpoint{1.871582in}{1.855316in}}%
\pgfpathlineto{\pgfqpoint{1.871582in}{1.855316in}}%
\pgfpathclose%
\pgfpathmoveto{\pgfqpoint{1.972148in}{1.931988in}}%
\pgfpathquadraticcurveto{\pgfqpoint{1.962108in}{1.931988in}}{\pgfqpoint{1.957036in}{1.922092in}}%
\pgfpathquadraticcurveto{\pgfqpoint{1.951985in}{1.912217in}}{\pgfqpoint{1.951985in}{1.892364in}}%
\pgfpathquadraticcurveto{\pgfqpoint{1.951985in}{1.872593in}}{\pgfqpoint{1.957036in}{1.862697in}}%
\pgfpathquadraticcurveto{\pgfqpoint{1.962108in}{1.852801in}}{\pgfqpoint{1.972148in}{1.852801in}}%
\pgfpathquadraticcurveto{\pgfqpoint{1.982271in}{1.852801in}}{\pgfqpoint{1.987322in}{1.862697in}}%
\pgfpathquadraticcurveto{\pgfqpoint{1.992393in}{1.872593in}}{\pgfqpoint{1.992393in}{1.892364in}}%
\pgfpathquadraticcurveto{\pgfqpoint{1.992393in}{1.912217in}}{\pgfqpoint{1.987322in}{1.922092in}}%
\pgfpathquadraticcurveto{\pgfqpoint{1.982271in}{1.931988in}}{\pgfqpoint{1.972148in}{1.931988in}}%
\pgfpathlineto{\pgfqpoint{1.972148in}{1.931988in}}%
\pgfpathclose%
\pgfpathmoveto{\pgfqpoint{1.972148in}{1.942296in}}%
\pgfpathquadraticcurveto{\pgfqpoint{1.988332in}{1.942296in}}{\pgfqpoint{1.996867in}{1.929494in}}%
\pgfpathquadraticcurveto{\pgfqpoint{2.005402in}{1.916711in}}{\pgfqpoint{2.005402in}{1.892364in}}%
\pgfpathquadraticcurveto{\pgfqpoint{2.005402in}{1.868078in}}{\pgfqpoint{1.996867in}{1.855275in}}%
\pgfpathquadraticcurveto{\pgfqpoint{1.988332in}{1.842493in}}{\pgfqpoint{1.972148in}{1.842493in}}%
\pgfpathquadraticcurveto{\pgfqpoint{1.955985in}{1.842493in}}{\pgfqpoint{1.947450in}{1.855275in}}%
\pgfpathquadraticcurveto{\pgfqpoint{1.938914in}{1.868078in}}{\pgfqpoint{1.938914in}{1.892364in}}%
\pgfpathquadraticcurveto{\pgfqpoint{1.938914in}{1.916711in}}{\pgfqpoint{1.947450in}{1.929494in}}%
\pgfpathquadraticcurveto{\pgfqpoint{1.955985in}{1.942296in}}{\pgfqpoint{1.972148in}{1.942296in}}%
\pgfpathlineto{\pgfqpoint{1.972148in}{1.942296in}}%
\pgfpathclose%
\pgfpathmoveto{\pgfqpoint{2.030533in}{1.855316in}}%
\pgfpathlineto{\pgfqpoint{2.051789in}{1.855316in}}%
\pgfpathlineto{\pgfqpoint{2.051789in}{1.928710in}}%
\pgfpathlineto{\pgfqpoint{2.028657in}{1.924071in}}%
\pgfpathlineto{\pgfqpoint{2.028657in}{1.935926in}}%
\pgfpathlineto{\pgfqpoint{2.051665in}{1.940565in}}%
\pgfpathlineto{\pgfqpoint{2.064674in}{1.940565in}}%
\pgfpathlineto{\pgfqpoint{2.064674in}{1.855316in}}%
\pgfpathlineto{\pgfqpoint{2.085929in}{1.855316in}}%
\pgfpathlineto{\pgfqpoint{2.085929in}{1.844369in}}%
\pgfpathlineto{\pgfqpoint{2.030533in}{1.844369in}}%
\pgfpathlineto{\pgfqpoint{2.030533in}{1.855316in}}%
\pgfpathlineto{\pgfqpoint{2.030533in}{1.855316in}}%
\pgfpathclose%
\pgfpathmoveto{\pgfqpoint{2.140047in}{1.890055in}}%
\pgfpathquadraticcurveto{\pgfqpoint{2.130770in}{1.890055in}}{\pgfqpoint{2.125451in}{1.885086in}}%
\pgfpathquadraticcurveto{\pgfqpoint{2.120153in}{1.880117in}}{\pgfqpoint{2.120153in}{1.871438in}}%
\pgfpathquadraticcurveto{\pgfqpoint{2.120153in}{1.862738in}}{\pgfqpoint{2.125451in}{1.857769in}}%
\pgfpathquadraticcurveto{\pgfqpoint{2.130770in}{1.852801in}}{\pgfqpoint{2.140047in}{1.852801in}}%
\pgfpathquadraticcurveto{\pgfqpoint{2.149325in}{1.852801in}}{\pgfqpoint{2.154664in}{1.857790in}}%
\pgfpathquadraticcurveto{\pgfqpoint{2.160025in}{1.862800in}}{\pgfqpoint{2.160025in}{1.871438in}}%
\pgfpathquadraticcurveto{\pgfqpoint{2.160025in}{1.880117in}}{\pgfqpoint{2.154706in}{1.885086in}}%
\pgfpathquadraticcurveto{\pgfqpoint{2.149407in}{1.890055in}}{\pgfqpoint{2.140047in}{1.890055in}}%
\pgfpathlineto{\pgfqpoint{2.140047in}{1.890055in}}%
\pgfpathclose%
\pgfpathmoveto{\pgfqpoint{2.127038in}{1.895580in}}%
\pgfpathquadraticcurveto{\pgfqpoint{2.118668in}{1.897641in}}{\pgfqpoint{2.113988in}{1.903373in}}%
\pgfpathquadraticcurveto{\pgfqpoint{2.109329in}{1.909125in}}{\pgfqpoint{2.109329in}{1.917371in}}%
\pgfpathquadraticcurveto{\pgfqpoint{2.109329in}{1.928896in}}{\pgfqpoint{2.117534in}{1.935596in}}%
\pgfpathquadraticcurveto{\pgfqpoint{2.125760in}{1.942296in}}{\pgfqpoint{2.140047in}{1.942296in}}%
\pgfpathquadraticcurveto{\pgfqpoint{2.154417in}{1.942296in}}{\pgfqpoint{2.162602in}{1.935596in}}%
\pgfpathquadraticcurveto{\pgfqpoint{2.170786in}{1.928896in}}{\pgfqpoint{2.170786in}{1.917371in}}%
\pgfpathquadraticcurveto{\pgfqpoint{2.170786in}{1.909125in}}{\pgfqpoint{2.166106in}{1.903373in}}%
\pgfpathquadraticcurveto{\pgfqpoint{2.161447in}{1.897641in}}{\pgfqpoint{2.153139in}{1.895580in}}%
\pgfpathquadraticcurveto{\pgfqpoint{2.162540in}{1.893394in}}{\pgfqpoint{2.167776in}{1.887003in}}%
\pgfpathquadraticcurveto{\pgfqpoint{2.173033in}{1.880633in}}{\pgfqpoint{2.173033in}{1.871438in}}%
\pgfpathquadraticcurveto{\pgfqpoint{2.173033in}{1.857440in}}{\pgfqpoint{2.164498in}{1.849956in}}%
\pgfpathquadraticcurveto{\pgfqpoint{2.155963in}{1.842493in}}{\pgfqpoint{2.140047in}{1.842493in}}%
\pgfpathquadraticcurveto{\pgfqpoint{2.124152in}{1.842493in}}{\pgfqpoint{2.115596in}{1.849956in}}%
\pgfpathquadraticcurveto{\pgfqpoint{2.107061in}{1.857440in}}{\pgfqpoint{2.107061in}{1.871438in}}%
\pgfpathquadraticcurveto{\pgfqpoint{2.107061in}{1.880633in}}{\pgfqpoint{2.112339in}{1.887003in}}%
\pgfpathquadraticcurveto{\pgfqpoint{2.117637in}{1.893394in}}{\pgfqpoint{2.127038in}{1.895580in}}%
\pgfpathlineto{\pgfqpoint{2.127038in}{1.895580in}}%
\pgfpathclose%
\pgfpathmoveto{\pgfqpoint{2.122276in}{1.916134in}}%
\pgfpathquadraticcurveto{\pgfqpoint{2.122276in}{1.908671in}}{\pgfqpoint{2.126935in}{1.904486in}}%
\pgfpathquadraticcurveto{\pgfqpoint{2.131615in}{1.900301in}}{\pgfqpoint{2.140047in}{1.900301in}}%
\pgfpathquadraticcurveto{\pgfqpoint{2.148438in}{1.900301in}}{\pgfqpoint{2.153159in}{1.904486in}}%
\pgfpathquadraticcurveto{\pgfqpoint{2.157901in}{1.908671in}}{\pgfqpoint{2.157901in}{1.916134in}}%
\pgfpathquadraticcurveto{\pgfqpoint{2.157901in}{1.923618in}}{\pgfqpoint{2.153159in}{1.927803in}}%
\pgfpathquadraticcurveto{\pgfqpoint{2.148438in}{1.931988in}}{\pgfqpoint{2.140047in}{1.931988in}}%
\pgfpathquadraticcurveto{\pgfqpoint{2.131615in}{1.931988in}}{\pgfqpoint{2.126935in}{1.927803in}}%
\pgfpathquadraticcurveto{\pgfqpoint{2.122276in}{1.923618in}}{\pgfqpoint{2.122276in}{1.916134in}}%
\pgfpathlineto{\pgfqpoint{2.122276in}{1.916134in}}%
\pgfpathclose%
\pgfpathmoveto{\pgfqpoint{2.192640in}{1.944482in}}%
\pgfpathlineto{\pgfqpoint{2.202948in}{1.944482in}}%
\pgfpathquadraticcurveto{\pgfqpoint{2.212596in}{1.929287in}}{\pgfqpoint{2.217400in}{1.914712in}}%
\pgfpathquadraticcurveto{\pgfqpoint{2.222203in}{1.900157in}}{\pgfqpoint{2.222203in}{1.885787in}}%
\pgfpathquadraticcurveto{\pgfqpoint{2.222203in}{1.871356in}}{\pgfqpoint{2.217400in}{1.856739in}}%
\pgfpathquadraticcurveto{\pgfqpoint{2.212596in}{1.842122in}}{\pgfqpoint{2.202948in}{1.826969in}}%
\pgfpathlineto{\pgfqpoint{2.192640in}{1.826969in}}%
\pgfpathquadraticcurveto{\pgfqpoint{2.201195in}{1.841730in}}{\pgfqpoint{2.205422in}{1.856326in}}%
\pgfpathquadraticcurveto{\pgfqpoint{2.209648in}{1.870923in}}{\pgfqpoint{2.209648in}{1.885787in}}%
\pgfpathquadraticcurveto{\pgfqpoint{2.209648in}{1.900672in}}{\pgfqpoint{2.205422in}{1.915165in}}%
\pgfpathquadraticcurveto{\pgfqpoint{2.201195in}{1.929679in}}{\pgfqpoint{2.192640in}{1.944482in}}%
\pgfpathlineto{\pgfqpoint{2.192640in}{1.944482in}}%
\pgfpathclose%
\pgfusepath{fill}%
\end{pgfscope}%
\begin{pgfscope}%
\pgfsetbuttcap%
\pgfsetroundjoin%
\definecolor{currentfill}{rgb}{0.090196,0.168627,0.227451}%
\pgfsetfillcolor{currentfill}%
\pgfsetlinewidth{0.000000pt}%
\definecolor{currentstroke}{rgb}{0.090196,0.168627,0.227451}%
\pgfsetstrokecolor{currentstroke}%
\pgfsetdash{}{0pt}%
\pgfpathmoveto{\pgfqpoint{1.316611in}{1.760770in}}%
\pgfpathlineto{\pgfqpoint{1.316611in}{1.748071in}}%
\pgfpathquadraticcurveto{\pgfqpoint{1.309210in}{1.751617in}}{\pgfqpoint{1.302633in}{1.753348in}}%
\pgfpathquadraticcurveto{\pgfqpoint{1.296056in}{1.755101in}}{\pgfqpoint{1.289933in}{1.755101in}}%
\pgfpathquadraticcurveto{\pgfqpoint{1.279316in}{1.755101in}}{\pgfqpoint{1.273543in}{1.750977in}}%
\pgfpathquadraticcurveto{\pgfqpoint{1.267771in}{1.746854in}}{\pgfqpoint{1.267771in}{1.739247in}}%
\pgfpathquadraticcurveto{\pgfqpoint{1.267771in}{1.732856in}}{\pgfqpoint{1.271605in}{1.729598in}}%
\pgfpathquadraticcurveto{\pgfqpoint{1.275440in}{1.726362in}}{\pgfqpoint{1.286140in}{1.724362in}}%
\pgfpathlineto{\pgfqpoint{1.293995in}{1.722754in}}%
\pgfpathquadraticcurveto{\pgfqpoint{1.308550in}{1.719971in}}{\pgfqpoint{1.315477in}{1.712982in}}%
\pgfpathquadraticcurveto{\pgfqpoint{1.322404in}{1.705993in}}{\pgfqpoint{1.322404in}{1.694283in}}%
\pgfpathquadraticcurveto{\pgfqpoint{1.322404in}{1.680284in}}{\pgfqpoint{1.313024in}{1.673068in}}%
\pgfpathquadraticcurveto{\pgfqpoint{1.303664in}{1.665853in}}{\pgfqpoint{1.285563in}{1.665853in}}%
\pgfpathquadraticcurveto{\pgfqpoint{1.278739in}{1.665853in}}{\pgfqpoint{1.271028in}{1.667399in}}%
\pgfpathquadraticcurveto{\pgfqpoint{1.263338in}{1.668945in}}{\pgfqpoint{1.255092in}{1.671976in}}%
\pgfpathlineto{\pgfqpoint{1.255092in}{1.685376in}}%
\pgfpathquadraticcurveto{\pgfqpoint{1.263008in}{1.680944in}}{\pgfqpoint{1.270616in}{1.678676in}}%
\pgfpathquadraticcurveto{\pgfqpoint{1.278223in}{1.676429in}}{\pgfqpoint{1.285563in}{1.676429in}}%
\pgfpathquadraticcurveto{\pgfqpoint{1.296696in}{1.676429in}}{\pgfqpoint{1.302757in}{1.680800in}}%
\pgfpathquadraticcurveto{\pgfqpoint{1.308818in}{1.685191in}}{\pgfqpoint{1.308818in}{1.693314in}}%
\pgfpathquadraticcurveto{\pgfqpoint{1.308818in}{1.700385in}}{\pgfqpoint{1.304468in}{1.704385in}}%
\pgfpathquadraticcurveto{\pgfqpoint{1.300118in}{1.708384in}}{\pgfqpoint{1.290201in}{1.710384in}}%
\pgfpathlineto{\pgfqpoint{1.282264in}{1.711930in}}%
\pgfpathquadraticcurveto{\pgfqpoint{1.267709in}{1.714816in}}{\pgfqpoint{1.261194in}{1.721001in}}%
\pgfpathquadraticcurveto{\pgfqpoint{1.254700in}{1.727186in}}{\pgfqpoint{1.254700in}{1.738216in}}%
\pgfpathquadraticcurveto{\pgfqpoint{1.254700in}{1.750977in}}{\pgfqpoint{1.263689in}{1.758317in}}%
\pgfpathquadraticcurveto{\pgfqpoint{1.272678in}{1.765656in}}{\pgfqpoint{1.288449in}{1.765656in}}%
\pgfpathquadraticcurveto{\pgfqpoint{1.295232in}{1.765656in}}{\pgfqpoint{1.302241in}{1.764419in}}%
\pgfpathquadraticcurveto{\pgfqpoint{1.309271in}{1.763203in}}{\pgfqpoint{1.316611in}{1.760770in}}%
\pgfpathlineto{\pgfqpoint{1.316611in}{1.760770in}}%
\pgfpathclose%
\pgfpathmoveto{\pgfqpoint{1.353926in}{1.760379in}}%
\pgfpathlineto{\pgfqpoint{1.353926in}{1.739886in}}%
\pgfpathlineto{\pgfqpoint{1.378336in}{1.739886in}}%
\pgfpathlineto{\pgfqpoint{1.378336in}{1.730670in}}%
\pgfpathlineto{\pgfqpoint{1.353926in}{1.730670in}}%
\pgfpathlineto{\pgfqpoint{1.353926in}{1.691499in}}%
\pgfpathquadraticcurveto{\pgfqpoint{1.353926in}{1.682676in}}{\pgfqpoint{1.356339in}{1.680160in}}%
\pgfpathquadraticcurveto{\pgfqpoint{1.358751in}{1.677645in}}{\pgfqpoint{1.366173in}{1.677645in}}%
\pgfpathlineto{\pgfqpoint{1.378336in}{1.677645in}}%
\pgfpathlineto{\pgfqpoint{1.378336in}{1.667729in}}%
\pgfpathlineto{\pgfqpoint{1.366173in}{1.667729in}}%
\pgfpathquadraticcurveto{\pgfqpoint{1.352442in}{1.667729in}}{\pgfqpoint{1.347226in}{1.672842in}}%
\pgfpathquadraticcurveto{\pgfqpoint{1.342010in}{1.677975in}}{\pgfqpoint{1.342010in}{1.691499in}}%
\pgfpathlineto{\pgfqpoint{1.342010in}{1.730670in}}%
\pgfpathlineto{\pgfqpoint{1.333310in}{1.730670in}}%
\pgfpathlineto{\pgfqpoint{1.333310in}{1.739886in}}%
\pgfpathlineto{\pgfqpoint{1.342010in}{1.739886in}}%
\pgfpathlineto{\pgfqpoint{1.342010in}{1.760379in}}%
\pgfpathlineto{\pgfqpoint{1.353926in}{1.760379in}}%
\pgfpathlineto{\pgfqpoint{1.353926in}{1.760379in}}%
\pgfpathclose%
\pgfpathmoveto{\pgfqpoint{1.392706in}{1.696200in}}%
\pgfpathlineto{\pgfqpoint{1.392706in}{1.739886in}}%
\pgfpathlineto{\pgfqpoint{1.404560in}{1.739886in}}%
\pgfpathlineto{\pgfqpoint{1.404560in}{1.696653in}}%
\pgfpathquadraticcurveto{\pgfqpoint{1.404560in}{1.686407in}}{\pgfqpoint{1.408539in}{1.681274in}}%
\pgfpathquadraticcurveto{\pgfqpoint{1.412539in}{1.676161in}}{\pgfqpoint{1.420538in}{1.676161in}}%
\pgfpathquadraticcurveto{\pgfqpoint{1.430124in}{1.676161in}}{\pgfqpoint{1.435691in}{1.682284in}}%
\pgfpathquadraticcurveto{\pgfqpoint{1.441278in}{1.688407in}}{\pgfqpoint{1.441278in}{1.698983in}}%
\pgfpathlineto{\pgfqpoint{1.441278in}{1.739886in}}%
\pgfpathlineto{\pgfqpoint{1.453132in}{1.739886in}}%
\pgfpathlineto{\pgfqpoint{1.453132in}{1.667729in}}%
\pgfpathlineto{\pgfqpoint{1.441278in}{1.667729in}}%
\pgfpathlineto{\pgfqpoint{1.441278in}{1.678820in}}%
\pgfpathquadraticcurveto{\pgfqpoint{1.436969in}{1.672244in}}{\pgfqpoint{1.431258in}{1.669048in}}%
\pgfpathquadraticcurveto{\pgfqpoint{1.425568in}{1.665853in}}{\pgfqpoint{1.418023in}{1.665853in}}%
\pgfpathquadraticcurveto{\pgfqpoint{1.405591in}{1.665853in}}{\pgfqpoint{1.399138in}{1.673584in}}%
\pgfpathquadraticcurveto{\pgfqpoint{1.392706in}{1.681315in}}{\pgfqpoint{1.392706in}{1.696200in}}%
\pgfpathlineto{\pgfqpoint{1.392706in}{1.696200in}}%
\pgfpathclose%
\pgfpathmoveto{\pgfqpoint{1.422538in}{1.741618in}}%
\pgfpathlineto{\pgfqpoint{1.422538in}{1.741618in}}%
\pgfpathlineto{\pgfqpoint{1.422538in}{1.741618in}}%
\pgfpathclose%
\pgfpathmoveto{\pgfqpoint{1.525021in}{1.728939in}}%
\pgfpathlineto{\pgfqpoint{1.525021in}{1.767986in}}%
\pgfpathlineto{\pgfqpoint{1.536876in}{1.767986in}}%
\pgfpathlineto{\pgfqpoint{1.536876in}{1.667729in}}%
\pgfpathlineto{\pgfqpoint{1.525021in}{1.667729in}}%
\pgfpathlineto{\pgfqpoint{1.525021in}{1.678552in}}%
\pgfpathquadraticcurveto{\pgfqpoint{1.521290in}{1.672120in}}{\pgfqpoint{1.515579in}{1.668986in}}%
\pgfpathquadraticcurveto{\pgfqpoint{1.509889in}{1.665853in}}{\pgfqpoint{1.501890in}{1.665853in}}%
\pgfpathquadraticcurveto{\pgfqpoint{1.488819in}{1.665853in}}{\pgfqpoint{1.480593in}{1.676285in}}%
\pgfpathquadraticcurveto{\pgfqpoint{1.472388in}{1.686737in}}{\pgfqpoint{1.472388in}{1.703746in}}%
\pgfpathquadraticcurveto{\pgfqpoint{1.472388in}{1.720754in}}{\pgfqpoint{1.480593in}{1.731186in}}%
\pgfpathquadraticcurveto{\pgfqpoint{1.488819in}{1.741618in}}{\pgfqpoint{1.501890in}{1.741618in}}%
\pgfpathquadraticcurveto{\pgfqpoint{1.509889in}{1.741618in}}{\pgfqpoint{1.515579in}{1.738484in}}%
\pgfpathquadraticcurveto{\pgfqpoint{1.521290in}{1.735371in}}{\pgfqpoint{1.525021in}{1.728939in}}%
\pgfpathlineto{\pgfqpoint{1.525021in}{1.728939in}}%
\pgfpathclose%
\pgfpathmoveto{\pgfqpoint{1.484634in}{1.703746in}}%
\pgfpathquadraticcurveto{\pgfqpoint{1.484634in}{1.690675in}}{\pgfqpoint{1.490015in}{1.683232in}}%
\pgfpathquadraticcurveto{\pgfqpoint{1.495396in}{1.675790in}}{\pgfqpoint{1.504797in}{1.675790in}}%
\pgfpathquadraticcurveto{\pgfqpoint{1.514198in}{1.675790in}}{\pgfqpoint{1.519599in}{1.683232in}}%
\pgfpathquadraticcurveto{\pgfqpoint{1.525021in}{1.690675in}}{\pgfqpoint{1.525021in}{1.703746in}}%
\pgfpathquadraticcurveto{\pgfqpoint{1.525021in}{1.716816in}}{\pgfqpoint{1.519599in}{1.724259in}}%
\pgfpathquadraticcurveto{\pgfqpoint{1.514198in}{1.731701in}}{\pgfqpoint{1.504797in}{1.731701in}}%
\pgfpathquadraticcurveto{\pgfqpoint{1.495396in}{1.731701in}}{\pgfqpoint{1.490015in}{1.724259in}}%
\pgfpathquadraticcurveto{\pgfqpoint{1.484634in}{1.716816in}}{\pgfqpoint{1.484634in}{1.703746in}}%
\pgfpathlineto{\pgfqpoint{1.484634in}{1.703746in}}%
\pgfpathclose%
\pgfpathmoveto{\pgfqpoint{1.591323in}{1.661028in}}%
\pgfpathquadraticcurveto{\pgfqpoint{1.586314in}{1.648143in}}{\pgfqpoint{1.581531in}{1.644226in}}%
\pgfpathquadraticcurveto{\pgfqpoint{1.576768in}{1.640288in}}{\pgfqpoint{1.568790in}{1.640288in}}%
\pgfpathlineto{\pgfqpoint{1.559306in}{1.640288in}}%
\pgfpathlineto{\pgfqpoint{1.559306in}{1.650205in}}%
\pgfpathlineto{\pgfqpoint{1.566275in}{1.650205in}}%
\pgfpathquadraticcurveto{\pgfqpoint{1.571161in}{1.650205in}}{\pgfqpoint{1.573861in}{1.652535in}}%
\pgfpathquadraticcurveto{\pgfqpoint{1.576583in}{1.654844in}}{\pgfqpoint{1.579861in}{1.663482in}}%
\pgfpathlineto{\pgfqpoint{1.581984in}{1.668883in}}%
\pgfpathlineto{\pgfqpoint{1.552812in}{1.739886in}}%
\pgfpathlineto{\pgfqpoint{1.565367in}{1.739886in}}%
\pgfpathlineto{\pgfqpoint{1.587922in}{1.683459in}}%
\pgfpathlineto{\pgfqpoint{1.610476in}{1.739886in}}%
\pgfpathlineto{\pgfqpoint{1.623031in}{1.739886in}}%
\pgfpathlineto{\pgfqpoint{1.591323in}{1.661028in}}%
\pgfpathlineto{\pgfqpoint{1.591323in}{1.661028in}}%
\pgfpathclose%
\pgfpathmoveto{\pgfqpoint{1.683148in}{1.763925in}}%
\pgfpathlineto{\pgfqpoint{1.734236in}{1.763925in}}%
\pgfpathlineto{\pgfqpoint{1.734236in}{1.752957in}}%
\pgfpathlineto{\pgfqpoint{1.695065in}{1.752957in}}%
\pgfpathlineto{\pgfqpoint{1.695065in}{1.729392in}}%
\pgfpathquadraticcurveto{\pgfqpoint{1.697889in}{1.730361in}}{\pgfqpoint{1.700714in}{1.730835in}}%
\pgfpathquadraticcurveto{\pgfqpoint{1.703559in}{1.731310in}}{\pgfqpoint{1.706404in}{1.731310in}}%
\pgfpathquadraticcurveto{\pgfqpoint{1.722505in}{1.731310in}}{\pgfqpoint{1.731906in}{1.722486in}}%
\pgfpathquadraticcurveto{\pgfqpoint{1.741328in}{1.713662in}}{\pgfqpoint{1.741328in}{1.698591in}}%
\pgfpathquadraticcurveto{\pgfqpoint{1.741328in}{1.683067in}}{\pgfqpoint{1.731659in}{1.674450in}}%
\pgfpathquadraticcurveto{\pgfqpoint{1.721990in}{1.665853in}}{\pgfqpoint{1.704404in}{1.665853in}}%
\pgfpathquadraticcurveto{\pgfqpoint{1.698343in}{1.665853in}}{\pgfqpoint{1.692055in}{1.666884in}}%
\pgfpathquadraticcurveto{\pgfqpoint{1.685787in}{1.667914in}}{\pgfqpoint{1.679087in}{1.669976in}}%
\pgfpathlineto{\pgfqpoint{1.679087in}{1.683067in}}%
\pgfpathquadraticcurveto{\pgfqpoint{1.684880in}{1.679913in}}{\pgfqpoint{1.691065in}{1.678367in}}%
\pgfpathquadraticcurveto{\pgfqpoint{1.697250in}{1.676821in}}{\pgfqpoint{1.704136in}{1.676821in}}%
\pgfpathquadraticcurveto{\pgfqpoint{1.715289in}{1.676821in}}{\pgfqpoint{1.721783in}{1.682676in}}%
\pgfpathquadraticcurveto{\pgfqpoint{1.728298in}{1.688531in}}{\pgfqpoint{1.728298in}{1.698591in}}%
\pgfpathquadraticcurveto{\pgfqpoint{1.728298in}{1.708632in}}{\pgfqpoint{1.721783in}{1.714487in}}%
\pgfpathquadraticcurveto{\pgfqpoint{1.715289in}{1.720362in}}{\pgfqpoint{1.704136in}{1.720362in}}%
\pgfpathquadraticcurveto{\pgfqpoint{1.698920in}{1.720362in}}{\pgfqpoint{1.693725in}{1.719208in}}%
\pgfpathquadraticcurveto{\pgfqpoint{1.688550in}{1.718053in}}{\pgfqpoint{1.683148in}{1.715600in}}%
\pgfpathlineto{\pgfqpoint{1.683148in}{1.763925in}}%
\pgfpathlineto{\pgfqpoint{1.683148in}{1.763925in}}%
\pgfpathclose%
\pgfpathmoveto{\pgfqpoint{1.835689in}{1.767842in}}%
\pgfpathquadraticcurveto{\pgfqpoint{1.827071in}{1.753039in}}{\pgfqpoint{1.822865in}{1.738525in}}%
\pgfpathquadraticcurveto{\pgfqpoint{1.818680in}{1.724032in}}{\pgfqpoint{1.818680in}{1.709147in}}%
\pgfpathquadraticcurveto{\pgfqpoint{1.818680in}{1.694283in}}{\pgfqpoint{1.822906in}{1.679686in}}%
\pgfpathquadraticcurveto{\pgfqpoint{1.827133in}{1.665090in}}{\pgfqpoint{1.835689in}{1.650329in}}%
\pgfpathlineto{\pgfqpoint{1.825380in}{1.650329in}}%
\pgfpathquadraticcurveto{\pgfqpoint{1.815732in}{1.665482in}}{\pgfqpoint{1.810928in}{1.680099in}}%
\pgfpathquadraticcurveto{\pgfqpoint{1.806125in}{1.694716in}}{\pgfqpoint{1.806125in}{1.709147in}}%
\pgfpathquadraticcurveto{\pgfqpoint{1.806125in}{1.723517in}}{\pgfqpoint{1.810887in}{1.738072in}}%
\pgfpathquadraticcurveto{\pgfqpoint{1.815670in}{1.752647in}}{\pgfqpoint{1.825380in}{1.767842in}}%
\pgfpathlineto{\pgfqpoint{1.835689in}{1.767842in}}%
\pgfpathlineto{\pgfqpoint{1.835689in}{1.767842in}}%
\pgfpathclose%
\pgfpathmoveto{\pgfqpoint{1.871582in}{1.678676in}}%
\pgfpathlineto{\pgfqpoint{1.916999in}{1.678676in}}%
\pgfpathlineto{\pgfqpoint{1.916999in}{1.667729in}}%
\pgfpathlineto{\pgfqpoint{1.855934in}{1.667729in}}%
\pgfpathlineto{\pgfqpoint{1.855934in}{1.678676in}}%
\pgfpathquadraticcurveto{\pgfqpoint{1.863335in}{1.686345in}}{\pgfqpoint{1.876117in}{1.699251in}}%
\pgfpathquadraticcurveto{\pgfqpoint{1.888920in}{1.712178in}}{\pgfqpoint{1.892198in}{1.715930in}}%
\pgfpathquadraticcurveto{\pgfqpoint{1.898445in}{1.722939in}}{\pgfqpoint{1.900919in}{1.727805in}}%
\pgfpathquadraticcurveto{\pgfqpoint{1.903413in}{1.732670in}}{\pgfqpoint{1.903413in}{1.737371in}}%
\pgfpathquadraticcurveto{\pgfqpoint{1.903413in}{1.745040in}}{\pgfqpoint{1.898032in}{1.749864in}}%
\pgfpathquadraticcurveto{\pgfqpoint{1.892651in}{1.754709in}}{\pgfqpoint{1.884013in}{1.754709in}}%
\pgfpathquadraticcurveto{\pgfqpoint{1.877890in}{1.754709in}}{\pgfqpoint{1.871087in}{1.752586in}}%
\pgfpathquadraticcurveto{\pgfqpoint{1.864304in}{1.750462in}}{\pgfqpoint{1.856573in}{1.746133in}}%
\pgfpathlineto{\pgfqpoint{1.856573in}{1.759286in}}%
\pgfpathquadraticcurveto{\pgfqpoint{1.864428in}{1.762440in}}{\pgfqpoint{1.871252in}{1.764048in}}%
\pgfpathquadraticcurveto{\pgfqpoint{1.878096in}{1.765656in}}{\pgfqpoint{1.883766in}{1.765656in}}%
\pgfpathquadraticcurveto{\pgfqpoint{1.898713in}{1.765656in}}{\pgfqpoint{1.907598in}{1.758173in}}%
\pgfpathquadraticcurveto{\pgfqpoint{1.916484in}{1.750709in}}{\pgfqpoint{1.916484in}{1.738216in}}%
\pgfpathquadraticcurveto{\pgfqpoint{1.916484in}{1.732278in}}{\pgfqpoint{1.914257in}{1.726959in}}%
\pgfpathquadraticcurveto{\pgfqpoint{1.912051in}{1.721661in}}{\pgfqpoint{1.906176in}{1.714445in}}%
\pgfpathquadraticcurveto{\pgfqpoint{1.904568in}{1.712569in}}{\pgfqpoint{1.895929in}{1.703642in}}%
\pgfpathquadraticcurveto{\pgfqpoint{1.887312in}{1.694716in}}{\pgfqpoint{1.871582in}{1.678676in}}%
\pgfpathlineto{\pgfqpoint{1.871582in}{1.678676in}}%
\pgfpathclose%
\pgfpathmoveto{\pgfqpoint{1.972148in}{1.755348in}}%
\pgfpathquadraticcurveto{\pgfqpoint{1.962108in}{1.755348in}}{\pgfqpoint{1.957036in}{1.745452in}}%
\pgfpathquadraticcurveto{\pgfqpoint{1.951985in}{1.735577in}}{\pgfqpoint{1.951985in}{1.715724in}}%
\pgfpathquadraticcurveto{\pgfqpoint{1.951985in}{1.695953in}}{\pgfqpoint{1.957036in}{1.686057in}}%
\pgfpathquadraticcurveto{\pgfqpoint{1.962108in}{1.676161in}}{\pgfqpoint{1.972148in}{1.676161in}}%
\pgfpathquadraticcurveto{\pgfqpoint{1.982271in}{1.676161in}}{\pgfqpoint{1.987322in}{1.686057in}}%
\pgfpathquadraticcurveto{\pgfqpoint{1.992393in}{1.695953in}}{\pgfqpoint{1.992393in}{1.715724in}}%
\pgfpathquadraticcurveto{\pgfqpoint{1.992393in}{1.735577in}}{\pgfqpoint{1.987322in}{1.745452in}}%
\pgfpathquadraticcurveto{\pgfqpoint{1.982271in}{1.755348in}}{\pgfqpoint{1.972148in}{1.755348in}}%
\pgfpathlineto{\pgfqpoint{1.972148in}{1.755348in}}%
\pgfpathclose%
\pgfpathmoveto{\pgfqpoint{1.972148in}{1.765656in}}%
\pgfpathquadraticcurveto{\pgfqpoint{1.988332in}{1.765656in}}{\pgfqpoint{1.996867in}{1.752854in}}%
\pgfpathquadraticcurveto{\pgfqpoint{2.005402in}{1.740071in}}{\pgfqpoint{2.005402in}{1.715724in}}%
\pgfpathquadraticcurveto{\pgfqpoint{2.005402in}{1.691438in}}{\pgfqpoint{1.996867in}{1.678635in}}%
\pgfpathquadraticcurveto{\pgfqpoint{1.988332in}{1.665853in}}{\pgfqpoint{1.972148in}{1.665853in}}%
\pgfpathquadraticcurveto{\pgfqpoint{1.955985in}{1.665853in}}{\pgfqpoint{1.947450in}{1.678635in}}%
\pgfpathquadraticcurveto{\pgfqpoint{1.938914in}{1.691438in}}{\pgfqpoint{1.938914in}{1.715724in}}%
\pgfpathquadraticcurveto{\pgfqpoint{1.938914in}{1.740071in}}{\pgfqpoint{1.947450in}{1.752854in}}%
\pgfpathquadraticcurveto{\pgfqpoint{1.955985in}{1.765656in}}{\pgfqpoint{1.972148in}{1.765656in}}%
\pgfpathlineto{\pgfqpoint{1.972148in}{1.765656in}}%
\pgfpathclose%
\pgfpathmoveto{\pgfqpoint{2.030533in}{1.678676in}}%
\pgfpathlineto{\pgfqpoint{2.051789in}{1.678676in}}%
\pgfpathlineto{\pgfqpoint{2.051789in}{1.752070in}}%
\pgfpathlineto{\pgfqpoint{2.028657in}{1.747431in}}%
\pgfpathlineto{\pgfqpoint{2.028657in}{1.759286in}}%
\pgfpathlineto{\pgfqpoint{2.051665in}{1.763925in}}%
\pgfpathlineto{\pgfqpoint{2.064674in}{1.763925in}}%
\pgfpathlineto{\pgfqpoint{2.064674in}{1.678676in}}%
\pgfpathlineto{\pgfqpoint{2.085929in}{1.678676in}}%
\pgfpathlineto{\pgfqpoint{2.085929in}{1.667729in}}%
\pgfpathlineto{\pgfqpoint{2.030533in}{1.667729in}}%
\pgfpathlineto{\pgfqpoint{2.030533in}{1.678676in}}%
\pgfpathlineto{\pgfqpoint{2.030533in}{1.678676in}}%
\pgfpathclose%
\pgfpathmoveto{\pgfqpoint{2.140047in}{1.713415in}}%
\pgfpathquadraticcurveto{\pgfqpoint{2.130770in}{1.713415in}}{\pgfqpoint{2.125451in}{1.708446in}}%
\pgfpathquadraticcurveto{\pgfqpoint{2.120153in}{1.703477in}}{\pgfqpoint{2.120153in}{1.694798in}}%
\pgfpathquadraticcurveto{\pgfqpoint{2.120153in}{1.686098in}}{\pgfqpoint{2.125451in}{1.681129in}}%
\pgfpathquadraticcurveto{\pgfqpoint{2.130770in}{1.676161in}}{\pgfqpoint{2.140047in}{1.676161in}}%
\pgfpathquadraticcurveto{\pgfqpoint{2.149325in}{1.676161in}}{\pgfqpoint{2.154664in}{1.681150in}}%
\pgfpathquadraticcurveto{\pgfqpoint{2.160025in}{1.686160in}}{\pgfqpoint{2.160025in}{1.694798in}}%
\pgfpathquadraticcurveto{\pgfqpoint{2.160025in}{1.703477in}}{\pgfqpoint{2.154706in}{1.708446in}}%
\pgfpathquadraticcurveto{\pgfqpoint{2.149407in}{1.713415in}}{\pgfqpoint{2.140047in}{1.713415in}}%
\pgfpathlineto{\pgfqpoint{2.140047in}{1.713415in}}%
\pgfpathclose%
\pgfpathmoveto{\pgfqpoint{2.127038in}{1.718940in}}%
\pgfpathquadraticcurveto{\pgfqpoint{2.118668in}{1.721001in}}{\pgfqpoint{2.113988in}{1.726733in}}%
\pgfpathquadraticcurveto{\pgfqpoint{2.109329in}{1.732485in}}{\pgfqpoint{2.109329in}{1.740731in}}%
\pgfpathquadraticcurveto{\pgfqpoint{2.109329in}{1.752256in}}{\pgfqpoint{2.117534in}{1.758956in}}%
\pgfpathquadraticcurveto{\pgfqpoint{2.125760in}{1.765656in}}{\pgfqpoint{2.140047in}{1.765656in}}%
\pgfpathquadraticcurveto{\pgfqpoint{2.154417in}{1.765656in}}{\pgfqpoint{2.162602in}{1.758956in}}%
\pgfpathquadraticcurveto{\pgfqpoint{2.170786in}{1.752256in}}{\pgfqpoint{2.170786in}{1.740731in}}%
\pgfpathquadraticcurveto{\pgfqpoint{2.170786in}{1.732485in}}{\pgfqpoint{2.166106in}{1.726733in}}%
\pgfpathquadraticcurveto{\pgfqpoint{2.161447in}{1.721001in}}{\pgfqpoint{2.153139in}{1.718940in}}%
\pgfpathquadraticcurveto{\pgfqpoint{2.162540in}{1.716754in}}{\pgfqpoint{2.167776in}{1.710363in}}%
\pgfpathquadraticcurveto{\pgfqpoint{2.173033in}{1.703993in}}{\pgfqpoint{2.173033in}{1.694798in}}%
\pgfpathquadraticcurveto{\pgfqpoint{2.173033in}{1.680800in}}{\pgfqpoint{2.164498in}{1.673316in}}%
\pgfpathquadraticcurveto{\pgfqpoint{2.155963in}{1.665853in}}{\pgfqpoint{2.140047in}{1.665853in}}%
\pgfpathquadraticcurveto{\pgfqpoint{2.124152in}{1.665853in}}{\pgfqpoint{2.115596in}{1.673316in}}%
\pgfpathquadraticcurveto{\pgfqpoint{2.107061in}{1.680800in}}{\pgfqpoint{2.107061in}{1.694798in}}%
\pgfpathquadraticcurveto{\pgfqpoint{2.107061in}{1.703993in}}{\pgfqpoint{2.112339in}{1.710363in}}%
\pgfpathquadraticcurveto{\pgfqpoint{2.117637in}{1.716754in}}{\pgfqpoint{2.127038in}{1.718940in}}%
\pgfpathlineto{\pgfqpoint{2.127038in}{1.718940in}}%
\pgfpathclose%
\pgfpathmoveto{\pgfqpoint{2.122276in}{1.739494in}}%
\pgfpathquadraticcurveto{\pgfqpoint{2.122276in}{1.732031in}}{\pgfqpoint{2.126935in}{1.727846in}}%
\pgfpathquadraticcurveto{\pgfqpoint{2.131615in}{1.723661in}}{\pgfqpoint{2.140047in}{1.723661in}}%
\pgfpathquadraticcurveto{\pgfqpoint{2.148438in}{1.723661in}}{\pgfqpoint{2.153159in}{1.727846in}}%
\pgfpathquadraticcurveto{\pgfqpoint{2.157901in}{1.732031in}}{\pgfqpoint{2.157901in}{1.739494in}}%
\pgfpathquadraticcurveto{\pgfqpoint{2.157901in}{1.746978in}}{\pgfqpoint{2.153159in}{1.751163in}}%
\pgfpathquadraticcurveto{\pgfqpoint{2.148438in}{1.755348in}}{\pgfqpoint{2.140047in}{1.755348in}}%
\pgfpathquadraticcurveto{\pgfqpoint{2.131615in}{1.755348in}}{\pgfqpoint{2.126935in}{1.751163in}}%
\pgfpathquadraticcurveto{\pgfqpoint{2.122276in}{1.746978in}}{\pgfqpoint{2.122276in}{1.739494in}}%
\pgfpathlineto{\pgfqpoint{2.122276in}{1.739494in}}%
\pgfpathclose%
\pgfpathmoveto{\pgfqpoint{2.192640in}{1.767842in}}%
\pgfpathlineto{\pgfqpoint{2.202948in}{1.767842in}}%
\pgfpathquadraticcurveto{\pgfqpoint{2.212596in}{1.752647in}}{\pgfqpoint{2.217400in}{1.738072in}}%
\pgfpathquadraticcurveto{\pgfqpoint{2.222203in}{1.723517in}}{\pgfqpoint{2.222203in}{1.709147in}}%
\pgfpathquadraticcurveto{\pgfqpoint{2.222203in}{1.694716in}}{\pgfqpoint{2.217400in}{1.680099in}}%
\pgfpathquadraticcurveto{\pgfqpoint{2.212596in}{1.665482in}}{\pgfqpoint{2.202948in}{1.650329in}}%
\pgfpathlineto{\pgfqpoint{2.192640in}{1.650329in}}%
\pgfpathquadraticcurveto{\pgfqpoint{2.201195in}{1.665090in}}{\pgfqpoint{2.205422in}{1.679686in}}%
\pgfpathquadraticcurveto{\pgfqpoint{2.209648in}{1.694283in}}{\pgfqpoint{2.209648in}{1.709147in}}%
\pgfpathquadraticcurveto{\pgfqpoint{2.209648in}{1.724032in}}{\pgfqpoint{2.205422in}{1.738525in}}%
\pgfpathquadraticcurveto{\pgfqpoint{2.201195in}{1.753039in}}{\pgfqpoint{2.192640in}{1.767842in}}%
\pgfpathlineto{\pgfqpoint{2.192640in}{1.767842in}}%
\pgfpathclose%
\pgfusepath{fill}%
\end{pgfscope}%
\begin{pgfscope}%
\pgfsetbuttcap%
\pgfsetroundjoin%
\definecolor{currentfill}{rgb}{0.090196,0.168627,0.227451}%
\pgfsetfillcolor{currentfill}%
\pgfsetlinewidth{0.000000pt}%
\definecolor{currentstroke}{rgb}{0.090196,0.168627,0.227451}%
\pgfsetstrokecolor{currentstroke}%
\pgfsetdash{}{0pt}%
\pgfpathmoveto{\pgfqpoint{1.316611in}{1.584130in}}%
\pgfpathlineto{\pgfqpoint{1.316611in}{1.571431in}}%
\pgfpathquadraticcurveto{\pgfqpoint{1.309210in}{1.574977in}}{\pgfqpoint{1.302633in}{1.576708in}}%
\pgfpathquadraticcurveto{\pgfqpoint{1.296056in}{1.578461in}}{\pgfqpoint{1.289933in}{1.578461in}}%
\pgfpathquadraticcurveto{\pgfqpoint{1.279316in}{1.578461in}}{\pgfqpoint{1.273543in}{1.574337in}}%
\pgfpathquadraticcurveto{\pgfqpoint{1.267771in}{1.570214in}}{\pgfqpoint{1.267771in}{1.562607in}}%
\pgfpathquadraticcurveto{\pgfqpoint{1.267771in}{1.556216in}}{\pgfqpoint{1.271605in}{1.552958in}}%
\pgfpathquadraticcurveto{\pgfqpoint{1.275440in}{1.549722in}}{\pgfqpoint{1.286140in}{1.547722in}}%
\pgfpathlineto{\pgfqpoint{1.293995in}{1.546114in}}%
\pgfpathquadraticcurveto{\pgfqpoint{1.308550in}{1.543331in}}{\pgfqpoint{1.315477in}{1.536342in}}%
\pgfpathquadraticcurveto{\pgfqpoint{1.322404in}{1.529353in}}{\pgfqpoint{1.322404in}{1.517643in}}%
\pgfpathquadraticcurveto{\pgfqpoint{1.322404in}{1.503644in}}{\pgfqpoint{1.313024in}{1.496428in}}%
\pgfpathquadraticcurveto{\pgfqpoint{1.303664in}{1.489213in}}{\pgfqpoint{1.285563in}{1.489213in}}%
\pgfpathquadraticcurveto{\pgfqpoint{1.278739in}{1.489213in}}{\pgfqpoint{1.271028in}{1.490759in}}%
\pgfpathquadraticcurveto{\pgfqpoint{1.263338in}{1.492305in}}{\pgfqpoint{1.255092in}{1.495336in}}%
\pgfpathlineto{\pgfqpoint{1.255092in}{1.508736in}}%
\pgfpathquadraticcurveto{\pgfqpoint{1.263008in}{1.504304in}}{\pgfqpoint{1.270616in}{1.502036in}}%
\pgfpathquadraticcurveto{\pgfqpoint{1.278223in}{1.499789in}}{\pgfqpoint{1.285563in}{1.499789in}}%
\pgfpathquadraticcurveto{\pgfqpoint{1.296696in}{1.499789in}}{\pgfqpoint{1.302757in}{1.504160in}}%
\pgfpathquadraticcurveto{\pgfqpoint{1.308818in}{1.508551in}}{\pgfqpoint{1.308818in}{1.516674in}}%
\pgfpathquadraticcurveto{\pgfqpoint{1.308818in}{1.523745in}}{\pgfqpoint{1.304468in}{1.527745in}}%
\pgfpathquadraticcurveto{\pgfqpoint{1.300118in}{1.531744in}}{\pgfqpoint{1.290201in}{1.533744in}}%
\pgfpathlineto{\pgfqpoint{1.282264in}{1.535290in}}%
\pgfpathquadraticcurveto{\pgfqpoint{1.267709in}{1.538176in}}{\pgfqpoint{1.261194in}{1.544361in}}%
\pgfpathquadraticcurveto{\pgfqpoint{1.254700in}{1.550546in}}{\pgfqpoint{1.254700in}{1.561576in}}%
\pgfpathquadraticcurveto{\pgfqpoint{1.254700in}{1.574337in}}{\pgfqpoint{1.263689in}{1.581677in}}%
\pgfpathquadraticcurveto{\pgfqpoint{1.272678in}{1.589016in}}{\pgfqpoint{1.288449in}{1.589016in}}%
\pgfpathquadraticcurveto{\pgfqpoint{1.295232in}{1.589016in}}{\pgfqpoint{1.302241in}{1.587779in}}%
\pgfpathquadraticcurveto{\pgfqpoint{1.309271in}{1.586563in}}{\pgfqpoint{1.316611in}{1.584130in}}%
\pgfpathlineto{\pgfqpoint{1.316611in}{1.584130in}}%
\pgfpathclose%
\pgfpathmoveto{\pgfqpoint{1.353926in}{1.583739in}}%
\pgfpathlineto{\pgfqpoint{1.353926in}{1.563246in}}%
\pgfpathlineto{\pgfqpoint{1.378336in}{1.563246in}}%
\pgfpathlineto{\pgfqpoint{1.378336in}{1.554030in}}%
\pgfpathlineto{\pgfqpoint{1.353926in}{1.554030in}}%
\pgfpathlineto{\pgfqpoint{1.353926in}{1.514859in}}%
\pgfpathquadraticcurveto{\pgfqpoint{1.353926in}{1.506036in}}{\pgfqpoint{1.356339in}{1.503520in}}%
\pgfpathquadraticcurveto{\pgfqpoint{1.358751in}{1.501005in}}{\pgfqpoint{1.366173in}{1.501005in}}%
\pgfpathlineto{\pgfqpoint{1.378336in}{1.501005in}}%
\pgfpathlineto{\pgfqpoint{1.378336in}{1.491089in}}%
\pgfpathlineto{\pgfqpoint{1.366173in}{1.491089in}}%
\pgfpathquadraticcurveto{\pgfqpoint{1.352442in}{1.491089in}}{\pgfqpoint{1.347226in}{1.496202in}}%
\pgfpathquadraticcurveto{\pgfqpoint{1.342010in}{1.501335in}}{\pgfqpoint{1.342010in}{1.514859in}}%
\pgfpathlineto{\pgfqpoint{1.342010in}{1.554030in}}%
\pgfpathlineto{\pgfqpoint{1.333310in}{1.554030in}}%
\pgfpathlineto{\pgfqpoint{1.333310in}{1.563246in}}%
\pgfpathlineto{\pgfqpoint{1.342010in}{1.563246in}}%
\pgfpathlineto{\pgfqpoint{1.342010in}{1.583739in}}%
\pgfpathlineto{\pgfqpoint{1.353926in}{1.583739in}}%
\pgfpathlineto{\pgfqpoint{1.353926in}{1.583739in}}%
\pgfpathclose%
\pgfpathmoveto{\pgfqpoint{1.392706in}{1.519560in}}%
\pgfpathlineto{\pgfqpoint{1.392706in}{1.563246in}}%
\pgfpathlineto{\pgfqpoint{1.404560in}{1.563246in}}%
\pgfpathlineto{\pgfqpoint{1.404560in}{1.520013in}}%
\pgfpathquadraticcurveto{\pgfqpoint{1.404560in}{1.509767in}}{\pgfqpoint{1.408539in}{1.504634in}}%
\pgfpathquadraticcurveto{\pgfqpoint{1.412539in}{1.499521in}}{\pgfqpoint{1.420538in}{1.499521in}}%
\pgfpathquadraticcurveto{\pgfqpoint{1.430124in}{1.499521in}}{\pgfqpoint{1.435691in}{1.505644in}}%
\pgfpathquadraticcurveto{\pgfqpoint{1.441278in}{1.511767in}}{\pgfqpoint{1.441278in}{1.522343in}}%
\pgfpathlineto{\pgfqpoint{1.441278in}{1.563246in}}%
\pgfpathlineto{\pgfqpoint{1.453132in}{1.563246in}}%
\pgfpathlineto{\pgfqpoint{1.453132in}{1.491089in}}%
\pgfpathlineto{\pgfqpoint{1.441278in}{1.491089in}}%
\pgfpathlineto{\pgfqpoint{1.441278in}{1.502180in}}%
\pgfpathquadraticcurveto{\pgfqpoint{1.436969in}{1.495604in}}{\pgfqpoint{1.431258in}{1.492408in}}%
\pgfpathquadraticcurveto{\pgfqpoint{1.425568in}{1.489213in}}{\pgfqpoint{1.418023in}{1.489213in}}%
\pgfpathquadraticcurveto{\pgfqpoint{1.405591in}{1.489213in}}{\pgfqpoint{1.399138in}{1.496944in}}%
\pgfpathquadraticcurveto{\pgfqpoint{1.392706in}{1.504675in}}{\pgfqpoint{1.392706in}{1.519560in}}%
\pgfpathlineto{\pgfqpoint{1.392706in}{1.519560in}}%
\pgfpathclose%
\pgfpathmoveto{\pgfqpoint{1.422538in}{1.564978in}}%
\pgfpathlineto{\pgfqpoint{1.422538in}{1.564978in}}%
\pgfpathlineto{\pgfqpoint{1.422538in}{1.564978in}}%
\pgfpathclose%
\pgfpathmoveto{\pgfqpoint{1.525021in}{1.552299in}}%
\pgfpathlineto{\pgfqpoint{1.525021in}{1.591346in}}%
\pgfpathlineto{\pgfqpoint{1.536876in}{1.591346in}}%
\pgfpathlineto{\pgfqpoint{1.536876in}{1.491089in}}%
\pgfpathlineto{\pgfqpoint{1.525021in}{1.491089in}}%
\pgfpathlineto{\pgfqpoint{1.525021in}{1.501912in}}%
\pgfpathquadraticcurveto{\pgfqpoint{1.521290in}{1.495480in}}{\pgfqpoint{1.515579in}{1.492346in}}%
\pgfpathquadraticcurveto{\pgfqpoint{1.509889in}{1.489213in}}{\pgfqpoint{1.501890in}{1.489213in}}%
\pgfpathquadraticcurveto{\pgfqpoint{1.488819in}{1.489213in}}{\pgfqpoint{1.480593in}{1.499645in}}%
\pgfpathquadraticcurveto{\pgfqpoint{1.472388in}{1.510097in}}{\pgfqpoint{1.472388in}{1.527106in}}%
\pgfpathquadraticcurveto{\pgfqpoint{1.472388in}{1.544114in}}{\pgfqpoint{1.480593in}{1.554546in}}%
\pgfpathquadraticcurveto{\pgfqpoint{1.488819in}{1.564978in}}{\pgfqpoint{1.501890in}{1.564978in}}%
\pgfpathquadraticcurveto{\pgfqpoint{1.509889in}{1.564978in}}{\pgfqpoint{1.515579in}{1.561844in}}%
\pgfpathquadraticcurveto{\pgfqpoint{1.521290in}{1.558731in}}{\pgfqpoint{1.525021in}{1.552299in}}%
\pgfpathlineto{\pgfqpoint{1.525021in}{1.552299in}}%
\pgfpathclose%
\pgfpathmoveto{\pgfqpoint{1.484634in}{1.527106in}}%
\pgfpathquadraticcurveto{\pgfqpoint{1.484634in}{1.514035in}}{\pgfqpoint{1.490015in}{1.506592in}}%
\pgfpathquadraticcurveto{\pgfqpoint{1.495396in}{1.499150in}}{\pgfqpoint{1.504797in}{1.499150in}}%
\pgfpathquadraticcurveto{\pgfqpoint{1.514198in}{1.499150in}}{\pgfqpoint{1.519599in}{1.506592in}}%
\pgfpathquadraticcurveto{\pgfqpoint{1.525021in}{1.514035in}}{\pgfqpoint{1.525021in}{1.527106in}}%
\pgfpathquadraticcurveto{\pgfqpoint{1.525021in}{1.540176in}}{\pgfqpoint{1.519599in}{1.547619in}}%
\pgfpathquadraticcurveto{\pgfqpoint{1.514198in}{1.555061in}}{\pgfqpoint{1.504797in}{1.555061in}}%
\pgfpathquadraticcurveto{\pgfqpoint{1.495396in}{1.555061in}}{\pgfqpoint{1.490015in}{1.547619in}}%
\pgfpathquadraticcurveto{\pgfqpoint{1.484634in}{1.540176in}}{\pgfqpoint{1.484634in}{1.527106in}}%
\pgfpathlineto{\pgfqpoint{1.484634in}{1.527106in}}%
\pgfpathclose%
\pgfpathmoveto{\pgfqpoint{1.591323in}{1.484388in}}%
\pgfpathquadraticcurveto{\pgfqpoint{1.586314in}{1.471503in}}{\pgfqpoint{1.581531in}{1.467586in}}%
\pgfpathquadraticcurveto{\pgfqpoint{1.576768in}{1.463648in}}{\pgfqpoint{1.568790in}{1.463648in}}%
\pgfpathlineto{\pgfqpoint{1.559306in}{1.463648in}}%
\pgfpathlineto{\pgfqpoint{1.559306in}{1.473565in}}%
\pgfpathlineto{\pgfqpoint{1.566275in}{1.473565in}}%
\pgfpathquadraticcurveto{\pgfqpoint{1.571161in}{1.473565in}}{\pgfqpoint{1.573861in}{1.475895in}}%
\pgfpathquadraticcurveto{\pgfqpoint{1.576583in}{1.478204in}}{\pgfqpoint{1.579861in}{1.486842in}}%
\pgfpathlineto{\pgfqpoint{1.581984in}{1.492243in}}%
\pgfpathlineto{\pgfqpoint{1.552812in}{1.563246in}}%
\pgfpathlineto{\pgfqpoint{1.565367in}{1.563246in}}%
\pgfpathlineto{\pgfqpoint{1.587922in}{1.506819in}}%
\pgfpathlineto{\pgfqpoint{1.610476in}{1.563246in}}%
\pgfpathlineto{\pgfqpoint{1.623031in}{1.563246in}}%
\pgfpathlineto{\pgfqpoint{1.591323in}{1.484388in}}%
\pgfpathlineto{\pgfqpoint{1.591323in}{1.484388in}}%
\pgfpathclose%
\pgfpathmoveto{\pgfqpoint{1.712465in}{1.544361in}}%
\pgfpathquadraticcurveto{\pgfqpoint{1.703703in}{1.544361in}}{\pgfqpoint{1.698569in}{1.538362in}}%
\pgfpathquadraticcurveto{\pgfqpoint{1.693457in}{1.532383in}}{\pgfqpoint{1.693457in}{1.521951in}}%
\pgfpathquadraticcurveto{\pgfqpoint{1.693457in}{1.511581in}}{\pgfqpoint{1.698569in}{1.505541in}}%
\pgfpathquadraticcurveto{\pgfqpoint{1.703703in}{1.499521in}}{\pgfqpoint{1.712465in}{1.499521in}}%
\pgfpathquadraticcurveto{\pgfqpoint{1.721227in}{1.499521in}}{\pgfqpoint{1.726340in}{1.505541in}}%
\pgfpathquadraticcurveto{\pgfqpoint{1.731452in}{1.511581in}}{\pgfqpoint{1.731452in}{1.521951in}}%
\pgfpathquadraticcurveto{\pgfqpoint{1.731452in}{1.532383in}}{\pgfqpoint{1.726340in}{1.538362in}}%
\pgfpathquadraticcurveto{\pgfqpoint{1.721227in}{1.544361in}}{\pgfqpoint{1.712465in}{1.544361in}}%
\pgfpathlineto{\pgfqpoint{1.712465in}{1.544361in}}%
\pgfpathclose%
\pgfpathmoveto{\pgfqpoint{1.738297in}{1.585161in}}%
\pgfpathlineto{\pgfqpoint{1.738297in}{1.573307in}}%
\pgfpathquadraticcurveto{\pgfqpoint{1.733390in}{1.575616in}}{\pgfqpoint{1.728401in}{1.576832in}}%
\pgfpathquadraticcurveto{\pgfqpoint{1.723412in}{1.578069in}}{\pgfqpoint{1.718505in}{1.578069in}}%
\pgfpathquadraticcurveto{\pgfqpoint{1.705620in}{1.578069in}}{\pgfqpoint{1.698817in}{1.569369in}}%
\pgfpathquadraticcurveto{\pgfqpoint{1.692034in}{1.560669in}}{\pgfqpoint{1.691065in}{1.543083in}}%
\pgfpathquadraticcurveto{\pgfqpoint{1.694859in}{1.548691in}}{\pgfqpoint{1.700590in}{1.551680in}}%
\pgfpathquadraticcurveto{\pgfqpoint{1.706342in}{1.554670in}}{\pgfqpoint{1.713228in}{1.554670in}}%
\pgfpathquadraticcurveto{\pgfqpoint{1.727721in}{1.554670in}}{\pgfqpoint{1.736132in}{1.545866in}}%
\pgfpathquadraticcurveto{\pgfqpoint{1.744544in}{1.537084in}}{\pgfqpoint{1.744544in}{1.521951in}}%
\pgfpathquadraticcurveto{\pgfqpoint{1.744544in}{1.507128in}}{\pgfqpoint{1.735782in}{1.498160in}}%
\pgfpathquadraticcurveto{\pgfqpoint{1.727020in}{1.489213in}}{\pgfqpoint{1.712465in}{1.489213in}}%
\pgfpathquadraticcurveto{\pgfqpoint{1.695766in}{1.489213in}}{\pgfqpoint{1.686942in}{1.501995in}}%
\pgfpathquadraticcurveto{\pgfqpoint{1.678118in}{1.514798in}}{\pgfqpoint{1.678118in}{1.539084in}}%
\pgfpathquadraticcurveto{\pgfqpoint{1.678118in}{1.561885in}}{\pgfqpoint{1.688942in}{1.575451in}}%
\pgfpathquadraticcurveto{\pgfqpoint{1.699765in}{1.589016in}}{\pgfqpoint{1.717990in}{1.589016in}}%
\pgfpathquadraticcurveto{\pgfqpoint{1.722897in}{1.589016in}}{\pgfqpoint{1.727886in}{1.588047in}}%
\pgfpathquadraticcurveto{\pgfqpoint{1.732875in}{1.587078in}}{\pgfqpoint{1.738297in}{1.585161in}}%
\pgfpathlineto{\pgfqpoint{1.738297in}{1.585161in}}%
\pgfpathclose%
\pgfpathmoveto{\pgfqpoint{1.835689in}{1.591202in}}%
\pgfpathquadraticcurveto{\pgfqpoint{1.827071in}{1.576399in}}{\pgfqpoint{1.822865in}{1.561885in}}%
\pgfpathquadraticcurveto{\pgfqpoint{1.818680in}{1.547392in}}{\pgfqpoint{1.818680in}{1.532507in}}%
\pgfpathquadraticcurveto{\pgfqpoint{1.818680in}{1.517643in}}{\pgfqpoint{1.822906in}{1.503046in}}%
\pgfpathquadraticcurveto{\pgfqpoint{1.827133in}{1.488450in}}{\pgfqpoint{1.835689in}{1.473689in}}%
\pgfpathlineto{\pgfqpoint{1.825380in}{1.473689in}}%
\pgfpathquadraticcurveto{\pgfqpoint{1.815732in}{1.488842in}}{\pgfqpoint{1.810928in}{1.503459in}}%
\pgfpathquadraticcurveto{\pgfqpoint{1.806125in}{1.518076in}}{\pgfqpoint{1.806125in}{1.532507in}}%
\pgfpathquadraticcurveto{\pgfqpoint{1.806125in}{1.546877in}}{\pgfqpoint{1.810887in}{1.561432in}}%
\pgfpathquadraticcurveto{\pgfqpoint{1.815670in}{1.576007in}}{\pgfqpoint{1.825380in}{1.591202in}}%
\pgfpathlineto{\pgfqpoint{1.835689in}{1.591202in}}%
\pgfpathlineto{\pgfqpoint{1.835689in}{1.591202in}}%
\pgfpathclose%
\pgfpathmoveto{\pgfqpoint{1.871582in}{1.502036in}}%
\pgfpathlineto{\pgfqpoint{1.916999in}{1.502036in}}%
\pgfpathlineto{\pgfqpoint{1.916999in}{1.491089in}}%
\pgfpathlineto{\pgfqpoint{1.855934in}{1.491089in}}%
\pgfpathlineto{\pgfqpoint{1.855934in}{1.502036in}}%
\pgfpathquadraticcurveto{\pgfqpoint{1.863335in}{1.509705in}}{\pgfqpoint{1.876117in}{1.522611in}}%
\pgfpathquadraticcurveto{\pgfqpoint{1.888920in}{1.535538in}}{\pgfqpoint{1.892198in}{1.539290in}}%
\pgfpathquadraticcurveto{\pgfqpoint{1.898445in}{1.546299in}}{\pgfqpoint{1.900919in}{1.551165in}}%
\pgfpathquadraticcurveto{\pgfqpoint{1.903413in}{1.556030in}}{\pgfqpoint{1.903413in}{1.560731in}}%
\pgfpathquadraticcurveto{\pgfqpoint{1.903413in}{1.568400in}}{\pgfqpoint{1.898032in}{1.573224in}}%
\pgfpathquadraticcurveto{\pgfqpoint{1.892651in}{1.578069in}}{\pgfqpoint{1.884013in}{1.578069in}}%
\pgfpathquadraticcurveto{\pgfqpoint{1.877890in}{1.578069in}}{\pgfqpoint{1.871087in}{1.575946in}}%
\pgfpathquadraticcurveto{\pgfqpoint{1.864304in}{1.573822in}}{\pgfqpoint{1.856573in}{1.569493in}}%
\pgfpathlineto{\pgfqpoint{1.856573in}{1.582646in}}%
\pgfpathquadraticcurveto{\pgfqpoint{1.864428in}{1.585800in}}{\pgfqpoint{1.871252in}{1.587408in}}%
\pgfpathquadraticcurveto{\pgfqpoint{1.878096in}{1.589016in}}{\pgfqpoint{1.883766in}{1.589016in}}%
\pgfpathquadraticcurveto{\pgfqpoint{1.898713in}{1.589016in}}{\pgfqpoint{1.907598in}{1.581533in}}%
\pgfpathquadraticcurveto{\pgfqpoint{1.916484in}{1.574069in}}{\pgfqpoint{1.916484in}{1.561576in}}%
\pgfpathquadraticcurveto{\pgfqpoint{1.916484in}{1.555638in}}{\pgfqpoint{1.914257in}{1.550319in}}%
\pgfpathquadraticcurveto{\pgfqpoint{1.912051in}{1.545021in}}{\pgfqpoint{1.906176in}{1.537805in}}%
\pgfpathquadraticcurveto{\pgfqpoint{1.904568in}{1.535929in}}{\pgfqpoint{1.895929in}{1.527002in}}%
\pgfpathquadraticcurveto{\pgfqpoint{1.887312in}{1.518076in}}{\pgfqpoint{1.871582in}{1.502036in}}%
\pgfpathlineto{\pgfqpoint{1.871582in}{1.502036in}}%
\pgfpathclose%
\pgfpathmoveto{\pgfqpoint{1.972148in}{1.578708in}}%
\pgfpathquadraticcurveto{\pgfqpoint{1.962108in}{1.578708in}}{\pgfqpoint{1.957036in}{1.568812in}}%
\pgfpathquadraticcurveto{\pgfqpoint{1.951985in}{1.558937in}}{\pgfqpoint{1.951985in}{1.539084in}}%
\pgfpathquadraticcurveto{\pgfqpoint{1.951985in}{1.519313in}}{\pgfqpoint{1.957036in}{1.509417in}}%
\pgfpathquadraticcurveto{\pgfqpoint{1.962108in}{1.499521in}}{\pgfqpoint{1.972148in}{1.499521in}}%
\pgfpathquadraticcurveto{\pgfqpoint{1.982271in}{1.499521in}}{\pgfqpoint{1.987322in}{1.509417in}}%
\pgfpathquadraticcurveto{\pgfqpoint{1.992393in}{1.519313in}}{\pgfqpoint{1.992393in}{1.539084in}}%
\pgfpathquadraticcurveto{\pgfqpoint{1.992393in}{1.558937in}}{\pgfqpoint{1.987322in}{1.568812in}}%
\pgfpathquadraticcurveto{\pgfqpoint{1.982271in}{1.578708in}}{\pgfqpoint{1.972148in}{1.578708in}}%
\pgfpathlineto{\pgfqpoint{1.972148in}{1.578708in}}%
\pgfpathclose%
\pgfpathmoveto{\pgfqpoint{1.972148in}{1.589016in}}%
\pgfpathquadraticcurveto{\pgfqpoint{1.988332in}{1.589016in}}{\pgfqpoint{1.996867in}{1.576214in}}%
\pgfpathquadraticcurveto{\pgfqpoint{2.005402in}{1.563431in}}{\pgfqpoint{2.005402in}{1.539084in}}%
\pgfpathquadraticcurveto{\pgfqpoint{2.005402in}{1.514798in}}{\pgfqpoint{1.996867in}{1.501995in}}%
\pgfpathquadraticcurveto{\pgfqpoint{1.988332in}{1.489213in}}{\pgfqpoint{1.972148in}{1.489213in}}%
\pgfpathquadraticcurveto{\pgfqpoint{1.955985in}{1.489213in}}{\pgfqpoint{1.947450in}{1.501995in}}%
\pgfpathquadraticcurveto{\pgfqpoint{1.938914in}{1.514798in}}{\pgfqpoint{1.938914in}{1.539084in}}%
\pgfpathquadraticcurveto{\pgfqpoint{1.938914in}{1.563431in}}{\pgfqpoint{1.947450in}{1.576214in}}%
\pgfpathquadraticcurveto{\pgfqpoint{1.955985in}{1.589016in}}{\pgfqpoint{1.972148in}{1.589016in}}%
\pgfpathlineto{\pgfqpoint{1.972148in}{1.589016in}}%
\pgfpathclose%
\pgfpathmoveto{\pgfqpoint{2.030533in}{1.502036in}}%
\pgfpathlineto{\pgfqpoint{2.051789in}{1.502036in}}%
\pgfpathlineto{\pgfqpoint{2.051789in}{1.575430in}}%
\pgfpathlineto{\pgfqpoint{2.028657in}{1.570791in}}%
\pgfpathlineto{\pgfqpoint{2.028657in}{1.582646in}}%
\pgfpathlineto{\pgfqpoint{2.051665in}{1.587285in}}%
\pgfpathlineto{\pgfqpoint{2.064674in}{1.587285in}}%
\pgfpathlineto{\pgfqpoint{2.064674in}{1.502036in}}%
\pgfpathlineto{\pgfqpoint{2.085929in}{1.502036in}}%
\pgfpathlineto{\pgfqpoint{2.085929in}{1.491089in}}%
\pgfpathlineto{\pgfqpoint{2.030533in}{1.491089in}}%
\pgfpathlineto{\pgfqpoint{2.030533in}{1.502036in}}%
\pgfpathlineto{\pgfqpoint{2.030533in}{1.502036in}}%
\pgfpathclose%
\pgfpathmoveto{\pgfqpoint{2.140047in}{1.536775in}}%
\pgfpathquadraticcurveto{\pgfqpoint{2.130770in}{1.536775in}}{\pgfqpoint{2.125451in}{1.531806in}}%
\pgfpathquadraticcurveto{\pgfqpoint{2.120153in}{1.526837in}}{\pgfqpoint{2.120153in}{1.518158in}}%
\pgfpathquadraticcurveto{\pgfqpoint{2.120153in}{1.509458in}}{\pgfqpoint{2.125451in}{1.504489in}}%
\pgfpathquadraticcurveto{\pgfqpoint{2.130770in}{1.499521in}}{\pgfqpoint{2.140047in}{1.499521in}}%
\pgfpathquadraticcurveto{\pgfqpoint{2.149325in}{1.499521in}}{\pgfqpoint{2.154664in}{1.504510in}}%
\pgfpathquadraticcurveto{\pgfqpoint{2.160025in}{1.509520in}}{\pgfqpoint{2.160025in}{1.518158in}}%
\pgfpathquadraticcurveto{\pgfqpoint{2.160025in}{1.526837in}}{\pgfqpoint{2.154706in}{1.531806in}}%
\pgfpathquadraticcurveto{\pgfqpoint{2.149407in}{1.536775in}}{\pgfqpoint{2.140047in}{1.536775in}}%
\pgfpathlineto{\pgfqpoint{2.140047in}{1.536775in}}%
\pgfpathclose%
\pgfpathmoveto{\pgfqpoint{2.127038in}{1.542300in}}%
\pgfpathquadraticcurveto{\pgfqpoint{2.118668in}{1.544361in}}{\pgfqpoint{2.113988in}{1.550093in}}%
\pgfpathquadraticcurveto{\pgfqpoint{2.109329in}{1.555845in}}{\pgfqpoint{2.109329in}{1.564091in}}%
\pgfpathquadraticcurveto{\pgfqpoint{2.109329in}{1.575616in}}{\pgfqpoint{2.117534in}{1.582316in}}%
\pgfpathquadraticcurveto{\pgfqpoint{2.125760in}{1.589016in}}{\pgfqpoint{2.140047in}{1.589016in}}%
\pgfpathquadraticcurveto{\pgfqpoint{2.154417in}{1.589016in}}{\pgfqpoint{2.162602in}{1.582316in}}%
\pgfpathquadraticcurveto{\pgfqpoint{2.170786in}{1.575616in}}{\pgfqpoint{2.170786in}{1.564091in}}%
\pgfpathquadraticcurveto{\pgfqpoint{2.170786in}{1.555845in}}{\pgfqpoint{2.166106in}{1.550093in}}%
\pgfpathquadraticcurveto{\pgfqpoint{2.161447in}{1.544361in}}{\pgfqpoint{2.153139in}{1.542300in}}%
\pgfpathquadraticcurveto{\pgfqpoint{2.162540in}{1.540114in}}{\pgfqpoint{2.167776in}{1.533723in}}%
\pgfpathquadraticcurveto{\pgfqpoint{2.173033in}{1.527353in}}{\pgfqpoint{2.173033in}{1.518158in}}%
\pgfpathquadraticcurveto{\pgfqpoint{2.173033in}{1.504160in}}{\pgfqpoint{2.164498in}{1.496676in}}%
\pgfpathquadraticcurveto{\pgfqpoint{2.155963in}{1.489213in}}{\pgfqpoint{2.140047in}{1.489213in}}%
\pgfpathquadraticcurveto{\pgfqpoint{2.124152in}{1.489213in}}{\pgfqpoint{2.115596in}{1.496676in}}%
\pgfpathquadraticcurveto{\pgfqpoint{2.107061in}{1.504160in}}{\pgfqpoint{2.107061in}{1.518158in}}%
\pgfpathquadraticcurveto{\pgfqpoint{2.107061in}{1.527353in}}{\pgfqpoint{2.112339in}{1.533723in}}%
\pgfpathquadraticcurveto{\pgfqpoint{2.117637in}{1.540114in}}{\pgfqpoint{2.127038in}{1.542300in}}%
\pgfpathlineto{\pgfqpoint{2.127038in}{1.542300in}}%
\pgfpathclose%
\pgfpathmoveto{\pgfqpoint{2.122276in}{1.562854in}}%
\pgfpathquadraticcurveto{\pgfqpoint{2.122276in}{1.555391in}}{\pgfqpoint{2.126935in}{1.551206in}}%
\pgfpathquadraticcurveto{\pgfqpoint{2.131615in}{1.547021in}}{\pgfqpoint{2.140047in}{1.547021in}}%
\pgfpathquadraticcurveto{\pgfqpoint{2.148438in}{1.547021in}}{\pgfqpoint{2.153159in}{1.551206in}}%
\pgfpathquadraticcurveto{\pgfqpoint{2.157901in}{1.555391in}}{\pgfqpoint{2.157901in}{1.562854in}}%
\pgfpathquadraticcurveto{\pgfqpoint{2.157901in}{1.570338in}}{\pgfqpoint{2.153159in}{1.574523in}}%
\pgfpathquadraticcurveto{\pgfqpoint{2.148438in}{1.578708in}}{\pgfqpoint{2.140047in}{1.578708in}}%
\pgfpathquadraticcurveto{\pgfqpoint{2.131615in}{1.578708in}}{\pgfqpoint{2.126935in}{1.574523in}}%
\pgfpathquadraticcurveto{\pgfqpoint{2.122276in}{1.570338in}}{\pgfqpoint{2.122276in}{1.562854in}}%
\pgfpathlineto{\pgfqpoint{2.122276in}{1.562854in}}%
\pgfpathclose%
\pgfpathmoveto{\pgfqpoint{2.192640in}{1.591202in}}%
\pgfpathlineto{\pgfqpoint{2.202948in}{1.591202in}}%
\pgfpathquadraticcurveto{\pgfqpoint{2.212596in}{1.576007in}}{\pgfqpoint{2.217400in}{1.561432in}}%
\pgfpathquadraticcurveto{\pgfqpoint{2.222203in}{1.546877in}}{\pgfqpoint{2.222203in}{1.532507in}}%
\pgfpathquadraticcurveto{\pgfqpoint{2.222203in}{1.518076in}}{\pgfqpoint{2.217400in}{1.503459in}}%
\pgfpathquadraticcurveto{\pgfqpoint{2.212596in}{1.488842in}}{\pgfqpoint{2.202948in}{1.473689in}}%
\pgfpathlineto{\pgfqpoint{2.192640in}{1.473689in}}%
\pgfpathquadraticcurveto{\pgfqpoint{2.201195in}{1.488450in}}{\pgfqpoint{2.205422in}{1.503046in}}%
\pgfpathquadraticcurveto{\pgfqpoint{2.209648in}{1.517643in}}{\pgfqpoint{2.209648in}{1.532507in}}%
\pgfpathquadraticcurveto{\pgfqpoint{2.209648in}{1.547392in}}{\pgfqpoint{2.205422in}{1.561885in}}%
\pgfpathquadraticcurveto{\pgfqpoint{2.201195in}{1.576399in}}{\pgfqpoint{2.192640in}{1.591202in}}%
\pgfpathlineto{\pgfqpoint{2.192640in}{1.591202in}}%
\pgfpathclose%
\pgfusepath{fill}%
\end{pgfscope}%
\begin{pgfscope}%
\pgfsetbuttcap%
\pgfsetroundjoin%
\definecolor{currentfill}{rgb}{0.090196,0.168627,0.227451}%
\pgfsetfillcolor{currentfill}%
\pgfsetlinewidth{0.000000pt}%
\definecolor{currentstroke}{rgb}{0.090196,0.168627,0.227451}%
\pgfsetstrokecolor{currentstroke}%
\pgfsetdash{}{0pt}%
\pgfpathmoveto{\pgfqpoint{1.316611in}{1.407490in}}%
\pgfpathlineto{\pgfqpoint{1.316611in}{1.394791in}}%
\pgfpathquadraticcurveto{\pgfqpoint{1.309210in}{1.398337in}}{\pgfqpoint{1.302633in}{1.400068in}}%
\pgfpathquadraticcurveto{\pgfqpoint{1.296056in}{1.401821in}}{\pgfqpoint{1.289933in}{1.401821in}}%
\pgfpathquadraticcurveto{\pgfqpoint{1.279316in}{1.401821in}}{\pgfqpoint{1.273543in}{1.397697in}}%
\pgfpathquadraticcurveto{\pgfqpoint{1.267771in}{1.393574in}}{\pgfqpoint{1.267771in}{1.385967in}}%
\pgfpathquadraticcurveto{\pgfqpoint{1.267771in}{1.379576in}}{\pgfqpoint{1.271605in}{1.376318in}}%
\pgfpathquadraticcurveto{\pgfqpoint{1.275440in}{1.373082in}}{\pgfqpoint{1.286140in}{1.371082in}}%
\pgfpathlineto{\pgfqpoint{1.293995in}{1.369474in}}%
\pgfpathquadraticcurveto{\pgfqpoint{1.308550in}{1.366691in}}{\pgfqpoint{1.315477in}{1.359702in}}%
\pgfpathquadraticcurveto{\pgfqpoint{1.322404in}{1.352713in}}{\pgfqpoint{1.322404in}{1.341003in}}%
\pgfpathquadraticcurveto{\pgfqpoint{1.322404in}{1.327004in}}{\pgfqpoint{1.313024in}{1.319788in}}%
\pgfpathquadraticcurveto{\pgfqpoint{1.303664in}{1.312573in}}{\pgfqpoint{1.285563in}{1.312573in}}%
\pgfpathquadraticcurveto{\pgfqpoint{1.278739in}{1.312573in}}{\pgfqpoint{1.271028in}{1.314119in}}%
\pgfpathquadraticcurveto{\pgfqpoint{1.263338in}{1.315665in}}{\pgfqpoint{1.255092in}{1.318696in}}%
\pgfpathlineto{\pgfqpoint{1.255092in}{1.332096in}}%
\pgfpathquadraticcurveto{\pgfqpoint{1.263008in}{1.327664in}}{\pgfqpoint{1.270616in}{1.325396in}}%
\pgfpathquadraticcurveto{\pgfqpoint{1.278223in}{1.323149in}}{\pgfqpoint{1.285563in}{1.323149in}}%
\pgfpathquadraticcurveto{\pgfqpoint{1.296696in}{1.323149in}}{\pgfqpoint{1.302757in}{1.327520in}}%
\pgfpathquadraticcurveto{\pgfqpoint{1.308818in}{1.331911in}}{\pgfqpoint{1.308818in}{1.340034in}}%
\pgfpathquadraticcurveto{\pgfqpoint{1.308818in}{1.347105in}}{\pgfqpoint{1.304468in}{1.351105in}}%
\pgfpathquadraticcurveto{\pgfqpoint{1.300118in}{1.355104in}}{\pgfqpoint{1.290201in}{1.357104in}}%
\pgfpathlineto{\pgfqpoint{1.282264in}{1.358650in}}%
\pgfpathquadraticcurveto{\pgfqpoint{1.267709in}{1.361536in}}{\pgfqpoint{1.261194in}{1.367721in}}%
\pgfpathquadraticcurveto{\pgfqpoint{1.254700in}{1.373906in}}{\pgfqpoint{1.254700in}{1.384936in}}%
\pgfpathquadraticcurveto{\pgfqpoint{1.254700in}{1.397697in}}{\pgfqpoint{1.263689in}{1.405037in}}%
\pgfpathquadraticcurveto{\pgfqpoint{1.272678in}{1.412376in}}{\pgfqpoint{1.288449in}{1.412376in}}%
\pgfpathquadraticcurveto{\pgfqpoint{1.295232in}{1.412376in}}{\pgfqpoint{1.302241in}{1.411139in}}%
\pgfpathquadraticcurveto{\pgfqpoint{1.309271in}{1.409923in}}{\pgfqpoint{1.316611in}{1.407490in}}%
\pgfpathlineto{\pgfqpoint{1.316611in}{1.407490in}}%
\pgfpathclose%
\pgfpathmoveto{\pgfqpoint{1.353926in}{1.407099in}}%
\pgfpathlineto{\pgfqpoint{1.353926in}{1.386606in}}%
\pgfpathlineto{\pgfqpoint{1.378336in}{1.386606in}}%
\pgfpathlineto{\pgfqpoint{1.378336in}{1.377390in}}%
\pgfpathlineto{\pgfqpoint{1.353926in}{1.377390in}}%
\pgfpathlineto{\pgfqpoint{1.353926in}{1.338219in}}%
\pgfpathquadraticcurveto{\pgfqpoint{1.353926in}{1.329396in}}{\pgfqpoint{1.356339in}{1.326880in}}%
\pgfpathquadraticcurveto{\pgfqpoint{1.358751in}{1.324365in}}{\pgfqpoint{1.366173in}{1.324365in}}%
\pgfpathlineto{\pgfqpoint{1.378336in}{1.324365in}}%
\pgfpathlineto{\pgfqpoint{1.378336in}{1.314449in}}%
\pgfpathlineto{\pgfqpoint{1.366173in}{1.314449in}}%
\pgfpathquadraticcurveto{\pgfqpoint{1.352442in}{1.314449in}}{\pgfqpoint{1.347226in}{1.319562in}}%
\pgfpathquadraticcurveto{\pgfqpoint{1.342010in}{1.324695in}}{\pgfqpoint{1.342010in}{1.338219in}}%
\pgfpathlineto{\pgfqpoint{1.342010in}{1.377390in}}%
\pgfpathlineto{\pgfqpoint{1.333310in}{1.377390in}}%
\pgfpathlineto{\pgfqpoint{1.333310in}{1.386606in}}%
\pgfpathlineto{\pgfqpoint{1.342010in}{1.386606in}}%
\pgfpathlineto{\pgfqpoint{1.342010in}{1.407099in}}%
\pgfpathlineto{\pgfqpoint{1.353926in}{1.407099in}}%
\pgfpathlineto{\pgfqpoint{1.353926in}{1.407099in}}%
\pgfpathclose%
\pgfpathmoveto{\pgfqpoint{1.392706in}{1.342920in}}%
\pgfpathlineto{\pgfqpoint{1.392706in}{1.386606in}}%
\pgfpathlineto{\pgfqpoint{1.404560in}{1.386606in}}%
\pgfpathlineto{\pgfqpoint{1.404560in}{1.343373in}}%
\pgfpathquadraticcurveto{\pgfqpoint{1.404560in}{1.333127in}}{\pgfqpoint{1.408539in}{1.327994in}}%
\pgfpathquadraticcurveto{\pgfqpoint{1.412539in}{1.322881in}}{\pgfqpoint{1.420538in}{1.322881in}}%
\pgfpathquadraticcurveto{\pgfqpoint{1.430124in}{1.322881in}}{\pgfqpoint{1.435691in}{1.329004in}}%
\pgfpathquadraticcurveto{\pgfqpoint{1.441278in}{1.335127in}}{\pgfqpoint{1.441278in}{1.345703in}}%
\pgfpathlineto{\pgfqpoint{1.441278in}{1.386606in}}%
\pgfpathlineto{\pgfqpoint{1.453132in}{1.386606in}}%
\pgfpathlineto{\pgfqpoint{1.453132in}{1.314449in}}%
\pgfpathlineto{\pgfqpoint{1.441278in}{1.314449in}}%
\pgfpathlineto{\pgfqpoint{1.441278in}{1.325540in}}%
\pgfpathquadraticcurveto{\pgfqpoint{1.436969in}{1.318964in}}{\pgfqpoint{1.431258in}{1.315768in}}%
\pgfpathquadraticcurveto{\pgfqpoint{1.425568in}{1.312573in}}{\pgfqpoint{1.418023in}{1.312573in}}%
\pgfpathquadraticcurveto{\pgfqpoint{1.405591in}{1.312573in}}{\pgfqpoint{1.399138in}{1.320304in}}%
\pgfpathquadraticcurveto{\pgfqpoint{1.392706in}{1.328035in}}{\pgfqpoint{1.392706in}{1.342920in}}%
\pgfpathlineto{\pgfqpoint{1.392706in}{1.342920in}}%
\pgfpathclose%
\pgfpathmoveto{\pgfqpoint{1.422538in}{1.388338in}}%
\pgfpathlineto{\pgfqpoint{1.422538in}{1.388338in}}%
\pgfpathlineto{\pgfqpoint{1.422538in}{1.388338in}}%
\pgfpathclose%
\pgfpathmoveto{\pgfqpoint{1.525021in}{1.375659in}}%
\pgfpathlineto{\pgfqpoint{1.525021in}{1.414706in}}%
\pgfpathlineto{\pgfqpoint{1.536876in}{1.414706in}}%
\pgfpathlineto{\pgfqpoint{1.536876in}{1.314449in}}%
\pgfpathlineto{\pgfqpoint{1.525021in}{1.314449in}}%
\pgfpathlineto{\pgfqpoint{1.525021in}{1.325272in}}%
\pgfpathquadraticcurveto{\pgfqpoint{1.521290in}{1.318840in}}{\pgfqpoint{1.515579in}{1.315706in}}%
\pgfpathquadraticcurveto{\pgfqpoint{1.509889in}{1.312573in}}{\pgfqpoint{1.501890in}{1.312573in}}%
\pgfpathquadraticcurveto{\pgfqpoint{1.488819in}{1.312573in}}{\pgfqpoint{1.480593in}{1.323005in}}%
\pgfpathquadraticcurveto{\pgfqpoint{1.472388in}{1.333457in}}{\pgfqpoint{1.472388in}{1.350466in}}%
\pgfpathquadraticcurveto{\pgfqpoint{1.472388in}{1.367474in}}{\pgfqpoint{1.480593in}{1.377906in}}%
\pgfpathquadraticcurveto{\pgfqpoint{1.488819in}{1.388338in}}{\pgfqpoint{1.501890in}{1.388338in}}%
\pgfpathquadraticcurveto{\pgfqpoint{1.509889in}{1.388338in}}{\pgfqpoint{1.515579in}{1.385204in}}%
\pgfpathquadraticcurveto{\pgfqpoint{1.521290in}{1.382091in}}{\pgfqpoint{1.525021in}{1.375659in}}%
\pgfpathlineto{\pgfqpoint{1.525021in}{1.375659in}}%
\pgfpathclose%
\pgfpathmoveto{\pgfqpoint{1.484634in}{1.350466in}}%
\pgfpathquadraticcurveto{\pgfqpoint{1.484634in}{1.337395in}}{\pgfqpoint{1.490015in}{1.329952in}}%
\pgfpathquadraticcurveto{\pgfqpoint{1.495396in}{1.322510in}}{\pgfqpoint{1.504797in}{1.322510in}}%
\pgfpathquadraticcurveto{\pgfqpoint{1.514198in}{1.322510in}}{\pgfqpoint{1.519599in}{1.329952in}}%
\pgfpathquadraticcurveto{\pgfqpoint{1.525021in}{1.337395in}}{\pgfqpoint{1.525021in}{1.350466in}}%
\pgfpathquadraticcurveto{\pgfqpoint{1.525021in}{1.363536in}}{\pgfqpoint{1.519599in}{1.370979in}}%
\pgfpathquadraticcurveto{\pgfqpoint{1.514198in}{1.378421in}}{\pgfqpoint{1.504797in}{1.378421in}}%
\pgfpathquadraticcurveto{\pgfqpoint{1.495396in}{1.378421in}}{\pgfqpoint{1.490015in}{1.370979in}}%
\pgfpathquadraticcurveto{\pgfqpoint{1.484634in}{1.363536in}}{\pgfqpoint{1.484634in}{1.350466in}}%
\pgfpathlineto{\pgfqpoint{1.484634in}{1.350466in}}%
\pgfpathclose%
\pgfpathmoveto{\pgfqpoint{1.591323in}{1.307748in}}%
\pgfpathquadraticcurveto{\pgfqpoint{1.586314in}{1.294863in}}{\pgfqpoint{1.581531in}{1.290946in}}%
\pgfpathquadraticcurveto{\pgfqpoint{1.576768in}{1.287008in}}{\pgfqpoint{1.568790in}{1.287008in}}%
\pgfpathlineto{\pgfqpoint{1.559306in}{1.287008in}}%
\pgfpathlineto{\pgfqpoint{1.559306in}{1.296925in}}%
\pgfpathlineto{\pgfqpoint{1.566275in}{1.296925in}}%
\pgfpathquadraticcurveto{\pgfqpoint{1.571161in}{1.296925in}}{\pgfqpoint{1.573861in}{1.299255in}}%
\pgfpathquadraticcurveto{\pgfqpoint{1.576583in}{1.301564in}}{\pgfqpoint{1.579861in}{1.310202in}}%
\pgfpathlineto{\pgfqpoint{1.581984in}{1.315603in}}%
\pgfpathlineto{\pgfqpoint{1.552812in}{1.386606in}}%
\pgfpathlineto{\pgfqpoint{1.565367in}{1.386606in}}%
\pgfpathlineto{\pgfqpoint{1.587922in}{1.330179in}}%
\pgfpathlineto{\pgfqpoint{1.610476in}{1.386606in}}%
\pgfpathlineto{\pgfqpoint{1.623031in}{1.386606in}}%
\pgfpathlineto{\pgfqpoint{1.591323in}{1.307748in}}%
\pgfpathlineto{\pgfqpoint{1.591323in}{1.307748in}}%
\pgfpathclose%
\pgfpathmoveto{\pgfqpoint{1.679726in}{1.410645in}}%
\pgfpathlineto{\pgfqpoint{1.741575in}{1.410645in}}%
\pgfpathlineto{\pgfqpoint{1.741575in}{1.405099in}}%
\pgfpathlineto{\pgfqpoint{1.706651in}{1.314449in}}%
\pgfpathlineto{\pgfqpoint{1.693065in}{1.314449in}}%
\pgfpathlineto{\pgfqpoint{1.725927in}{1.399677in}}%
\pgfpathlineto{\pgfqpoint{1.679726in}{1.399677in}}%
\pgfpathlineto{\pgfqpoint{1.679726in}{1.410645in}}%
\pgfpathlineto{\pgfqpoint{1.679726in}{1.410645in}}%
\pgfpathclose%
\pgfpathmoveto{\pgfqpoint{1.835689in}{1.414562in}}%
\pgfpathquadraticcurveto{\pgfqpoint{1.827071in}{1.399759in}}{\pgfqpoint{1.822865in}{1.385245in}}%
\pgfpathquadraticcurveto{\pgfqpoint{1.818680in}{1.370752in}}{\pgfqpoint{1.818680in}{1.355867in}}%
\pgfpathquadraticcurveto{\pgfqpoint{1.818680in}{1.341003in}}{\pgfqpoint{1.822906in}{1.326406in}}%
\pgfpathquadraticcurveto{\pgfqpoint{1.827133in}{1.311810in}}{\pgfqpoint{1.835689in}{1.297049in}}%
\pgfpathlineto{\pgfqpoint{1.825380in}{1.297049in}}%
\pgfpathquadraticcurveto{\pgfqpoint{1.815732in}{1.312202in}}{\pgfqpoint{1.810928in}{1.326819in}}%
\pgfpathquadraticcurveto{\pgfqpoint{1.806125in}{1.341436in}}{\pgfqpoint{1.806125in}{1.355867in}}%
\pgfpathquadraticcurveto{\pgfqpoint{1.806125in}{1.370237in}}{\pgfqpoint{1.810887in}{1.384792in}}%
\pgfpathquadraticcurveto{\pgfqpoint{1.815670in}{1.399367in}}{\pgfqpoint{1.825380in}{1.414562in}}%
\pgfpathlineto{\pgfqpoint{1.835689in}{1.414562in}}%
\pgfpathlineto{\pgfqpoint{1.835689in}{1.414562in}}%
\pgfpathclose%
\pgfpathmoveto{\pgfqpoint{1.871582in}{1.325396in}}%
\pgfpathlineto{\pgfqpoint{1.916999in}{1.325396in}}%
\pgfpathlineto{\pgfqpoint{1.916999in}{1.314449in}}%
\pgfpathlineto{\pgfqpoint{1.855934in}{1.314449in}}%
\pgfpathlineto{\pgfqpoint{1.855934in}{1.325396in}}%
\pgfpathquadraticcurveto{\pgfqpoint{1.863335in}{1.333065in}}{\pgfqpoint{1.876117in}{1.345971in}}%
\pgfpathquadraticcurveto{\pgfqpoint{1.888920in}{1.358898in}}{\pgfqpoint{1.892198in}{1.362650in}}%
\pgfpathquadraticcurveto{\pgfqpoint{1.898445in}{1.369659in}}{\pgfqpoint{1.900919in}{1.374525in}}%
\pgfpathquadraticcurveto{\pgfqpoint{1.903413in}{1.379390in}}{\pgfqpoint{1.903413in}{1.384091in}}%
\pgfpathquadraticcurveto{\pgfqpoint{1.903413in}{1.391760in}}{\pgfqpoint{1.898032in}{1.396584in}}%
\pgfpathquadraticcurveto{\pgfqpoint{1.892651in}{1.401429in}}{\pgfqpoint{1.884013in}{1.401429in}}%
\pgfpathquadraticcurveto{\pgfqpoint{1.877890in}{1.401429in}}{\pgfqpoint{1.871087in}{1.399306in}}%
\pgfpathquadraticcurveto{\pgfqpoint{1.864304in}{1.397182in}}{\pgfqpoint{1.856573in}{1.392853in}}%
\pgfpathlineto{\pgfqpoint{1.856573in}{1.406006in}}%
\pgfpathquadraticcurveto{\pgfqpoint{1.864428in}{1.409160in}}{\pgfqpoint{1.871252in}{1.410768in}}%
\pgfpathquadraticcurveto{\pgfqpoint{1.878096in}{1.412376in}}{\pgfqpoint{1.883766in}{1.412376in}}%
\pgfpathquadraticcurveto{\pgfqpoint{1.898713in}{1.412376in}}{\pgfqpoint{1.907598in}{1.404893in}}%
\pgfpathquadraticcurveto{\pgfqpoint{1.916484in}{1.397429in}}{\pgfqpoint{1.916484in}{1.384936in}}%
\pgfpathquadraticcurveto{\pgfqpoint{1.916484in}{1.378998in}}{\pgfqpoint{1.914257in}{1.373679in}}%
\pgfpathquadraticcurveto{\pgfqpoint{1.912051in}{1.368381in}}{\pgfqpoint{1.906176in}{1.361165in}}%
\pgfpathquadraticcurveto{\pgfqpoint{1.904568in}{1.359289in}}{\pgfqpoint{1.895929in}{1.350362in}}%
\pgfpathquadraticcurveto{\pgfqpoint{1.887312in}{1.341436in}}{\pgfqpoint{1.871582in}{1.325396in}}%
\pgfpathlineto{\pgfqpoint{1.871582in}{1.325396in}}%
\pgfpathclose%
\pgfpathmoveto{\pgfqpoint{1.972148in}{1.402068in}}%
\pgfpathquadraticcurveto{\pgfqpoint{1.962108in}{1.402068in}}{\pgfqpoint{1.957036in}{1.392172in}}%
\pgfpathquadraticcurveto{\pgfqpoint{1.951985in}{1.382297in}}{\pgfqpoint{1.951985in}{1.362444in}}%
\pgfpathquadraticcurveto{\pgfqpoint{1.951985in}{1.342673in}}{\pgfqpoint{1.957036in}{1.332777in}}%
\pgfpathquadraticcurveto{\pgfqpoint{1.962108in}{1.322881in}}{\pgfqpoint{1.972148in}{1.322881in}}%
\pgfpathquadraticcurveto{\pgfqpoint{1.982271in}{1.322881in}}{\pgfqpoint{1.987322in}{1.332777in}}%
\pgfpathquadraticcurveto{\pgfqpoint{1.992393in}{1.342673in}}{\pgfqpoint{1.992393in}{1.362444in}}%
\pgfpathquadraticcurveto{\pgfqpoint{1.992393in}{1.382297in}}{\pgfqpoint{1.987322in}{1.392172in}}%
\pgfpathquadraticcurveto{\pgfqpoint{1.982271in}{1.402068in}}{\pgfqpoint{1.972148in}{1.402068in}}%
\pgfpathlineto{\pgfqpoint{1.972148in}{1.402068in}}%
\pgfpathclose%
\pgfpathmoveto{\pgfqpoint{1.972148in}{1.412376in}}%
\pgfpathquadraticcurveto{\pgfqpoint{1.988332in}{1.412376in}}{\pgfqpoint{1.996867in}{1.399574in}}%
\pgfpathquadraticcurveto{\pgfqpoint{2.005402in}{1.386791in}}{\pgfqpoint{2.005402in}{1.362444in}}%
\pgfpathquadraticcurveto{\pgfqpoint{2.005402in}{1.338158in}}{\pgfqpoint{1.996867in}{1.325355in}}%
\pgfpathquadraticcurveto{\pgfqpoint{1.988332in}{1.312573in}}{\pgfqpoint{1.972148in}{1.312573in}}%
\pgfpathquadraticcurveto{\pgfqpoint{1.955985in}{1.312573in}}{\pgfqpoint{1.947450in}{1.325355in}}%
\pgfpathquadraticcurveto{\pgfqpoint{1.938914in}{1.338158in}}{\pgfqpoint{1.938914in}{1.362444in}}%
\pgfpathquadraticcurveto{\pgfqpoint{1.938914in}{1.386791in}}{\pgfqpoint{1.947450in}{1.399574in}}%
\pgfpathquadraticcurveto{\pgfqpoint{1.955985in}{1.412376in}}{\pgfqpoint{1.972148in}{1.412376in}}%
\pgfpathlineto{\pgfqpoint{1.972148in}{1.412376in}}%
\pgfpathclose%
\pgfpathmoveto{\pgfqpoint{2.030533in}{1.325396in}}%
\pgfpathlineto{\pgfqpoint{2.051789in}{1.325396in}}%
\pgfpathlineto{\pgfqpoint{2.051789in}{1.398790in}}%
\pgfpathlineto{\pgfqpoint{2.028657in}{1.394151in}}%
\pgfpathlineto{\pgfqpoint{2.028657in}{1.406006in}}%
\pgfpathlineto{\pgfqpoint{2.051665in}{1.410645in}}%
\pgfpathlineto{\pgfqpoint{2.064674in}{1.410645in}}%
\pgfpathlineto{\pgfqpoint{2.064674in}{1.325396in}}%
\pgfpathlineto{\pgfqpoint{2.085929in}{1.325396in}}%
\pgfpathlineto{\pgfqpoint{2.085929in}{1.314449in}}%
\pgfpathlineto{\pgfqpoint{2.030533in}{1.314449in}}%
\pgfpathlineto{\pgfqpoint{2.030533in}{1.325396in}}%
\pgfpathlineto{\pgfqpoint{2.030533in}{1.325396in}}%
\pgfpathclose%
\pgfpathmoveto{\pgfqpoint{2.112607in}{1.316449in}}%
\pgfpathlineto{\pgfqpoint{2.112607in}{1.328303in}}%
\pgfpathquadraticcurveto{\pgfqpoint{2.117514in}{1.325973in}}{\pgfqpoint{2.122523in}{1.324757in}}%
\pgfpathquadraticcurveto{\pgfqpoint{2.127554in}{1.323541in}}{\pgfqpoint{2.132399in}{1.323541in}}%
\pgfpathquadraticcurveto{\pgfqpoint{2.145284in}{1.323541in}}{\pgfqpoint{2.152067in}{1.332199in}}%
\pgfpathquadraticcurveto{\pgfqpoint{2.158870in}{1.340858in}}{\pgfqpoint{2.159839in}{1.358526in}}%
\pgfpathquadraticcurveto{\pgfqpoint{2.156107in}{1.352981in}}{\pgfqpoint{2.150355in}{1.350012in}}%
\pgfpathquadraticcurveto{\pgfqpoint{2.144624in}{1.347043in}}{\pgfqpoint{2.137676in}{1.347043in}}%
\pgfpathquadraticcurveto{\pgfqpoint{2.123245in}{1.347043in}}{\pgfqpoint{2.114834in}{1.355764in}}%
\pgfpathquadraticcurveto{\pgfqpoint{2.106422in}{1.364505in}}{\pgfqpoint{2.106422in}{1.379658in}}%
\pgfpathquadraticcurveto{\pgfqpoint{2.106422in}{1.394461in}}{\pgfqpoint{2.115184in}{1.403408in}}%
\pgfpathquadraticcurveto{\pgfqpoint{2.123946in}{1.412376in}}{\pgfqpoint{2.138501in}{1.412376in}}%
\pgfpathquadraticcurveto{\pgfqpoint{2.155200in}{1.412376in}}{\pgfqpoint{2.163983in}{1.399574in}}%
\pgfpathquadraticcurveto{\pgfqpoint{2.172786in}{1.386791in}}{\pgfqpoint{2.172786in}{1.362444in}}%
\pgfpathquadraticcurveto{\pgfqpoint{2.172786in}{1.339704in}}{\pgfqpoint{2.161983in}{1.326138in}}%
\pgfpathquadraticcurveto{\pgfqpoint{2.151201in}{1.312573in}}{\pgfqpoint{2.132976in}{1.312573in}}%
\pgfpathquadraticcurveto{\pgfqpoint{2.128069in}{1.312573in}}{\pgfqpoint{2.123039in}{1.313542in}}%
\pgfpathquadraticcurveto{\pgfqpoint{2.118029in}{1.314511in}}{\pgfqpoint{2.112607in}{1.316449in}}%
\pgfpathlineto{\pgfqpoint{2.112607in}{1.316449in}}%
\pgfpathclose%
\pgfpathmoveto{\pgfqpoint{2.138501in}{1.357228in}}%
\pgfpathquadraticcurveto{\pgfqpoint{2.147263in}{1.357228in}}{\pgfqpoint{2.152376in}{1.363206in}}%
\pgfpathquadraticcurveto{\pgfqpoint{2.157509in}{1.369206in}}{\pgfqpoint{2.157509in}{1.379658in}}%
\pgfpathquadraticcurveto{\pgfqpoint{2.157509in}{1.390028in}}{\pgfqpoint{2.152376in}{1.396048in}}%
\pgfpathquadraticcurveto{\pgfqpoint{2.147263in}{1.402068in}}{\pgfqpoint{2.138501in}{1.402068in}}%
\pgfpathquadraticcurveto{\pgfqpoint{2.129739in}{1.402068in}}{\pgfqpoint{2.124626in}{1.396048in}}%
\pgfpathquadraticcurveto{\pgfqpoint{2.119513in}{1.390028in}}{\pgfqpoint{2.119513in}{1.379658in}}%
\pgfpathquadraticcurveto{\pgfqpoint{2.119513in}{1.369206in}}{\pgfqpoint{2.124626in}{1.363206in}}%
\pgfpathquadraticcurveto{\pgfqpoint{2.129739in}{1.357228in}}{\pgfqpoint{2.138501in}{1.357228in}}%
\pgfpathlineto{\pgfqpoint{2.138501in}{1.357228in}}%
\pgfpathclose%
\pgfpathmoveto{\pgfqpoint{2.192640in}{1.414562in}}%
\pgfpathlineto{\pgfqpoint{2.202948in}{1.414562in}}%
\pgfpathquadraticcurveto{\pgfqpoint{2.212596in}{1.399367in}}{\pgfqpoint{2.217400in}{1.384792in}}%
\pgfpathquadraticcurveto{\pgfqpoint{2.222203in}{1.370237in}}{\pgfqpoint{2.222203in}{1.355867in}}%
\pgfpathquadraticcurveto{\pgfqpoint{2.222203in}{1.341436in}}{\pgfqpoint{2.217400in}{1.326819in}}%
\pgfpathquadraticcurveto{\pgfqpoint{2.212596in}{1.312202in}}{\pgfqpoint{2.202948in}{1.297049in}}%
\pgfpathlineto{\pgfqpoint{2.192640in}{1.297049in}}%
\pgfpathquadraticcurveto{\pgfqpoint{2.201195in}{1.311810in}}{\pgfqpoint{2.205422in}{1.326406in}}%
\pgfpathquadraticcurveto{\pgfqpoint{2.209648in}{1.341003in}}{\pgfqpoint{2.209648in}{1.355867in}}%
\pgfpathquadraticcurveto{\pgfqpoint{2.209648in}{1.370752in}}{\pgfqpoint{2.205422in}{1.385245in}}%
\pgfpathquadraticcurveto{\pgfqpoint{2.201195in}{1.399759in}}{\pgfqpoint{2.192640in}{1.414562in}}%
\pgfpathlineto{\pgfqpoint{2.192640in}{1.414562in}}%
\pgfpathclose%
\pgfusepath{fill}%
\end{pgfscope}%
\begin{pgfscope}%
\pgfsetbuttcap%
\pgfsetroundjoin%
\definecolor{currentfill}{rgb}{0.090196,0.168627,0.227451}%
\pgfsetfillcolor{currentfill}%
\pgfsetlinewidth{0.000000pt}%
\definecolor{currentstroke}{rgb}{0.090196,0.168627,0.227451}%
\pgfsetstrokecolor{currentstroke}%
\pgfsetdash{}{0pt}%
\pgfpathmoveto{\pgfqpoint{4.550611in}{2.290690in}}%
\pgfpathlineto{\pgfqpoint{4.550611in}{2.277991in}}%
\pgfpathquadraticcurveto{\pgfqpoint{4.543210in}{2.281537in}}{\pgfqpoint{4.536633in}{2.283268in}}%
\pgfpathquadraticcurveto{\pgfqpoint{4.530056in}{2.285021in}}{\pgfqpoint{4.523933in}{2.285021in}}%
\pgfpathquadraticcurveto{\pgfqpoint{4.513316in}{2.285021in}}{\pgfqpoint{4.507543in}{2.280897in}}%
\pgfpathquadraticcurveto{\pgfqpoint{4.501771in}{2.276774in}}{\pgfqpoint{4.501771in}{2.269167in}}%
\pgfpathquadraticcurveto{\pgfqpoint{4.501771in}{2.262776in}}{\pgfqpoint{4.505605in}{2.259518in}}%
\pgfpathquadraticcurveto{\pgfqpoint{4.509440in}{2.256282in}}{\pgfqpoint{4.520140in}{2.254282in}}%
\pgfpathlineto{\pgfqpoint{4.527995in}{2.252674in}}%
\pgfpathquadraticcurveto{\pgfqpoint{4.542550in}{2.249891in}}{\pgfqpoint{4.549477in}{2.242902in}}%
\pgfpathquadraticcurveto{\pgfqpoint{4.556404in}{2.235913in}}{\pgfqpoint{4.556404in}{2.224203in}}%
\pgfpathquadraticcurveto{\pgfqpoint{4.556404in}{2.210204in}}{\pgfqpoint{4.547024in}{2.202988in}}%
\pgfpathquadraticcurveto{\pgfqpoint{4.537664in}{2.195773in}}{\pgfqpoint{4.519563in}{2.195773in}}%
\pgfpathquadraticcurveto{\pgfqpoint{4.512739in}{2.195773in}}{\pgfqpoint{4.505028in}{2.197319in}}%
\pgfpathquadraticcurveto{\pgfqpoint{4.497338in}{2.198865in}}{\pgfqpoint{4.489092in}{2.201896in}}%
\pgfpathlineto{\pgfqpoint{4.489092in}{2.215296in}}%
\pgfpathquadraticcurveto{\pgfqpoint{4.497008in}{2.210864in}}{\pgfqpoint{4.504616in}{2.208596in}}%
\pgfpathquadraticcurveto{\pgfqpoint{4.512223in}{2.206349in}}{\pgfqpoint{4.519563in}{2.206349in}}%
\pgfpathquadraticcurveto{\pgfqpoint{4.530696in}{2.206349in}}{\pgfqpoint{4.536757in}{2.210720in}}%
\pgfpathquadraticcurveto{\pgfqpoint{4.542818in}{2.215111in}}{\pgfqpoint{4.542818in}{2.223234in}}%
\pgfpathquadraticcurveto{\pgfqpoint{4.542818in}{2.230305in}}{\pgfqpoint{4.538468in}{2.234305in}}%
\pgfpathquadraticcurveto{\pgfqpoint{4.534118in}{2.238304in}}{\pgfqpoint{4.524201in}{2.240304in}}%
\pgfpathlineto{\pgfqpoint{4.516264in}{2.241850in}}%
\pgfpathquadraticcurveto{\pgfqpoint{4.501709in}{2.244736in}}{\pgfqpoint{4.495194in}{2.250921in}}%
\pgfpathquadraticcurveto{\pgfqpoint{4.488700in}{2.257106in}}{\pgfqpoint{4.488700in}{2.268136in}}%
\pgfpathquadraticcurveto{\pgfqpoint{4.488700in}{2.280897in}}{\pgfqpoint{4.497689in}{2.288237in}}%
\pgfpathquadraticcurveto{\pgfqpoint{4.506678in}{2.295576in}}{\pgfqpoint{4.522449in}{2.295576in}}%
\pgfpathquadraticcurveto{\pgfqpoint{4.529232in}{2.295576in}}{\pgfqpoint{4.536241in}{2.294339in}}%
\pgfpathquadraticcurveto{\pgfqpoint{4.543271in}{2.293123in}}{\pgfqpoint{4.550611in}{2.290690in}}%
\pgfpathlineto{\pgfqpoint{4.550611in}{2.290690in}}%
\pgfpathclose%
\pgfpathmoveto{\pgfqpoint{4.587926in}{2.290299in}}%
\pgfpathlineto{\pgfqpoint{4.587926in}{2.269806in}}%
\pgfpathlineto{\pgfqpoint{4.612336in}{2.269806in}}%
\pgfpathlineto{\pgfqpoint{4.612336in}{2.260590in}}%
\pgfpathlineto{\pgfqpoint{4.587926in}{2.260590in}}%
\pgfpathlineto{\pgfqpoint{4.587926in}{2.221419in}}%
\pgfpathquadraticcurveto{\pgfqpoint{4.587926in}{2.212596in}}{\pgfqpoint{4.590339in}{2.210080in}}%
\pgfpathquadraticcurveto{\pgfqpoint{4.592751in}{2.207565in}}{\pgfqpoint{4.600173in}{2.207565in}}%
\pgfpathlineto{\pgfqpoint{4.612336in}{2.207565in}}%
\pgfpathlineto{\pgfqpoint{4.612336in}{2.197649in}}%
\pgfpathlineto{\pgfqpoint{4.600173in}{2.197649in}}%
\pgfpathquadraticcurveto{\pgfqpoint{4.586442in}{2.197649in}}{\pgfqpoint{4.581226in}{2.202762in}}%
\pgfpathquadraticcurveto{\pgfqpoint{4.576010in}{2.207895in}}{\pgfqpoint{4.576010in}{2.221419in}}%
\pgfpathlineto{\pgfqpoint{4.576010in}{2.260590in}}%
\pgfpathlineto{\pgfqpoint{4.567310in}{2.260590in}}%
\pgfpathlineto{\pgfqpoint{4.567310in}{2.269806in}}%
\pgfpathlineto{\pgfqpoint{4.576010in}{2.269806in}}%
\pgfpathlineto{\pgfqpoint{4.576010in}{2.290299in}}%
\pgfpathlineto{\pgfqpoint{4.587926in}{2.290299in}}%
\pgfpathlineto{\pgfqpoint{4.587926in}{2.290299in}}%
\pgfpathclose%
\pgfpathmoveto{\pgfqpoint{4.626706in}{2.226120in}}%
\pgfpathlineto{\pgfqpoint{4.626706in}{2.269806in}}%
\pgfpathlineto{\pgfqpoint{4.638560in}{2.269806in}}%
\pgfpathlineto{\pgfqpoint{4.638560in}{2.226573in}}%
\pgfpathquadraticcurveto{\pgfqpoint{4.638560in}{2.216327in}}{\pgfqpoint{4.642539in}{2.211194in}}%
\pgfpathquadraticcurveto{\pgfqpoint{4.646539in}{2.206081in}}{\pgfqpoint{4.654538in}{2.206081in}}%
\pgfpathquadraticcurveto{\pgfqpoint{4.664124in}{2.206081in}}{\pgfqpoint{4.669691in}{2.212204in}}%
\pgfpathquadraticcurveto{\pgfqpoint{4.675278in}{2.218327in}}{\pgfqpoint{4.675278in}{2.228903in}}%
\pgfpathlineto{\pgfqpoint{4.675278in}{2.269806in}}%
\pgfpathlineto{\pgfqpoint{4.687132in}{2.269806in}}%
\pgfpathlineto{\pgfqpoint{4.687132in}{2.197649in}}%
\pgfpathlineto{\pgfqpoint{4.675278in}{2.197649in}}%
\pgfpathlineto{\pgfqpoint{4.675278in}{2.208740in}}%
\pgfpathquadraticcurveto{\pgfqpoint{4.670969in}{2.202164in}}{\pgfqpoint{4.665258in}{2.198968in}}%
\pgfpathquadraticcurveto{\pgfqpoint{4.659568in}{2.195773in}}{\pgfqpoint{4.652023in}{2.195773in}}%
\pgfpathquadraticcurveto{\pgfqpoint{4.639591in}{2.195773in}}{\pgfqpoint{4.633138in}{2.203504in}}%
\pgfpathquadraticcurveto{\pgfqpoint{4.626706in}{2.211235in}}{\pgfqpoint{4.626706in}{2.226120in}}%
\pgfpathlineto{\pgfqpoint{4.626706in}{2.226120in}}%
\pgfpathclose%
\pgfpathmoveto{\pgfqpoint{4.656538in}{2.271538in}}%
\pgfpathlineto{\pgfqpoint{4.656538in}{2.271538in}}%
\pgfpathlineto{\pgfqpoint{4.656538in}{2.271538in}}%
\pgfpathclose%
\pgfpathmoveto{\pgfqpoint{4.759021in}{2.258859in}}%
\pgfpathlineto{\pgfqpoint{4.759021in}{2.297906in}}%
\pgfpathlineto{\pgfqpoint{4.770876in}{2.297906in}}%
\pgfpathlineto{\pgfqpoint{4.770876in}{2.197649in}}%
\pgfpathlineto{\pgfqpoint{4.759021in}{2.197649in}}%
\pgfpathlineto{\pgfqpoint{4.759021in}{2.208472in}}%
\pgfpathquadraticcurveto{\pgfqpoint{4.755290in}{2.202040in}}{\pgfqpoint{4.749579in}{2.198906in}}%
\pgfpathquadraticcurveto{\pgfqpoint{4.743889in}{2.195773in}}{\pgfqpoint{4.735890in}{2.195773in}}%
\pgfpathquadraticcurveto{\pgfqpoint{4.722819in}{2.195773in}}{\pgfqpoint{4.714593in}{2.206205in}}%
\pgfpathquadraticcurveto{\pgfqpoint{4.706388in}{2.216657in}}{\pgfqpoint{4.706388in}{2.233666in}}%
\pgfpathquadraticcurveto{\pgfqpoint{4.706388in}{2.250674in}}{\pgfqpoint{4.714593in}{2.261106in}}%
\pgfpathquadraticcurveto{\pgfqpoint{4.722819in}{2.271538in}}{\pgfqpoint{4.735890in}{2.271538in}}%
\pgfpathquadraticcurveto{\pgfqpoint{4.743889in}{2.271538in}}{\pgfqpoint{4.749579in}{2.268404in}}%
\pgfpathquadraticcurveto{\pgfqpoint{4.755290in}{2.265291in}}{\pgfqpoint{4.759021in}{2.258859in}}%
\pgfpathlineto{\pgfqpoint{4.759021in}{2.258859in}}%
\pgfpathclose%
\pgfpathmoveto{\pgfqpoint{4.718634in}{2.233666in}}%
\pgfpathquadraticcurveto{\pgfqpoint{4.718634in}{2.220595in}}{\pgfqpoint{4.724015in}{2.213152in}}%
\pgfpathquadraticcurveto{\pgfqpoint{4.729396in}{2.205710in}}{\pgfqpoint{4.738797in}{2.205710in}}%
\pgfpathquadraticcurveto{\pgfqpoint{4.748198in}{2.205710in}}{\pgfqpoint{4.753599in}{2.213152in}}%
\pgfpathquadraticcurveto{\pgfqpoint{4.759021in}{2.220595in}}{\pgfqpoint{4.759021in}{2.233666in}}%
\pgfpathquadraticcurveto{\pgfqpoint{4.759021in}{2.246736in}}{\pgfqpoint{4.753599in}{2.254179in}}%
\pgfpathquadraticcurveto{\pgfqpoint{4.748198in}{2.261621in}}{\pgfqpoint{4.738797in}{2.261621in}}%
\pgfpathquadraticcurveto{\pgfqpoint{4.729396in}{2.261621in}}{\pgfqpoint{4.724015in}{2.254179in}}%
\pgfpathquadraticcurveto{\pgfqpoint{4.718634in}{2.246736in}}{\pgfqpoint{4.718634in}{2.233666in}}%
\pgfpathlineto{\pgfqpoint{4.718634in}{2.233666in}}%
\pgfpathclose%
\pgfpathmoveto{\pgfqpoint{4.825323in}{2.190948in}}%
\pgfpathquadraticcurveto{\pgfqpoint{4.820314in}{2.178063in}}{\pgfqpoint{4.815531in}{2.174146in}}%
\pgfpathquadraticcurveto{\pgfqpoint{4.810768in}{2.170208in}}{\pgfqpoint{4.802790in}{2.170208in}}%
\pgfpathlineto{\pgfqpoint{4.793306in}{2.170208in}}%
\pgfpathlineto{\pgfqpoint{4.793306in}{2.180125in}}%
\pgfpathlineto{\pgfqpoint{4.800275in}{2.180125in}}%
\pgfpathquadraticcurveto{\pgfqpoint{4.805161in}{2.180125in}}{\pgfqpoint{4.807861in}{2.182455in}}%
\pgfpathquadraticcurveto{\pgfqpoint{4.810583in}{2.184764in}}{\pgfqpoint{4.813861in}{2.193402in}}%
\pgfpathlineto{\pgfqpoint{4.815984in}{2.198803in}}%
\pgfpathlineto{\pgfqpoint{4.786812in}{2.269806in}}%
\pgfpathlineto{\pgfqpoint{4.799367in}{2.269806in}}%
\pgfpathlineto{\pgfqpoint{4.821922in}{2.213379in}}%
\pgfpathlineto{\pgfqpoint{4.844476in}{2.269806in}}%
\pgfpathlineto{\pgfqpoint{4.857031in}{2.269806in}}%
\pgfpathlineto{\pgfqpoint{4.825323in}{2.190948in}}%
\pgfpathlineto{\pgfqpoint{4.825323in}{2.190948in}}%
\pgfpathclose%
\pgfpathmoveto{\pgfqpoint{4.944836in}{2.243335in}}%
\pgfpathquadraticcurveto{\pgfqpoint{4.935559in}{2.243335in}}{\pgfqpoint{4.930240in}{2.238366in}}%
\pgfpathquadraticcurveto{\pgfqpoint{4.924941in}{2.233397in}}{\pgfqpoint{4.924941in}{2.224718in}}%
\pgfpathquadraticcurveto{\pgfqpoint{4.924941in}{2.216018in}}{\pgfqpoint{4.930240in}{2.211049in}}%
\pgfpathquadraticcurveto{\pgfqpoint{4.935559in}{2.206081in}}{\pgfqpoint{4.944836in}{2.206081in}}%
\pgfpathquadraticcurveto{\pgfqpoint{4.954113in}{2.206081in}}{\pgfqpoint{4.959453in}{2.211070in}}%
\pgfpathquadraticcurveto{\pgfqpoint{4.964813in}{2.216080in}}{\pgfqpoint{4.964813in}{2.224718in}}%
\pgfpathquadraticcurveto{\pgfqpoint{4.964813in}{2.233397in}}{\pgfqpoint{4.959494in}{2.238366in}}%
\pgfpathquadraticcurveto{\pgfqpoint{4.954196in}{2.243335in}}{\pgfqpoint{4.944836in}{2.243335in}}%
\pgfpathlineto{\pgfqpoint{4.944836in}{2.243335in}}%
\pgfpathclose%
\pgfpathmoveto{\pgfqpoint{4.931827in}{2.248860in}}%
\pgfpathquadraticcurveto{\pgfqpoint{4.923457in}{2.250921in}}{\pgfqpoint{4.918777in}{2.256653in}}%
\pgfpathquadraticcurveto{\pgfqpoint{4.914118in}{2.262405in}}{\pgfqpoint{4.914118in}{2.270651in}}%
\pgfpathquadraticcurveto{\pgfqpoint{4.914118in}{2.282176in}}{\pgfqpoint{4.922323in}{2.288876in}}%
\pgfpathquadraticcurveto{\pgfqpoint{4.930549in}{2.295576in}}{\pgfqpoint{4.944836in}{2.295576in}}%
\pgfpathquadraticcurveto{\pgfqpoint{4.959206in}{2.295576in}}{\pgfqpoint{4.967390in}{2.288876in}}%
\pgfpathquadraticcurveto{\pgfqpoint{4.975575in}{2.282176in}}{\pgfqpoint{4.975575in}{2.270651in}}%
\pgfpathquadraticcurveto{\pgfqpoint{4.975575in}{2.262405in}}{\pgfqpoint{4.970895in}{2.256653in}}%
\pgfpathquadraticcurveto{\pgfqpoint{4.966236in}{2.250921in}}{\pgfqpoint{4.957928in}{2.248860in}}%
\pgfpathquadraticcurveto{\pgfqpoint{4.967329in}{2.246674in}}{\pgfqpoint{4.972565in}{2.240283in}}%
\pgfpathquadraticcurveto{\pgfqpoint{4.977822in}{2.233913in}}{\pgfqpoint{4.977822in}{2.224718in}}%
\pgfpathquadraticcurveto{\pgfqpoint{4.977822in}{2.210720in}}{\pgfqpoint{4.969287in}{2.203236in}}%
\pgfpathquadraticcurveto{\pgfqpoint{4.960752in}{2.195773in}}{\pgfqpoint{4.944836in}{2.195773in}}%
\pgfpathquadraticcurveto{\pgfqpoint{4.928941in}{2.195773in}}{\pgfqpoint{4.920385in}{2.203236in}}%
\pgfpathquadraticcurveto{\pgfqpoint{4.911850in}{2.210720in}}{\pgfqpoint{4.911850in}{2.224718in}}%
\pgfpathquadraticcurveto{\pgfqpoint{4.911850in}{2.233913in}}{\pgfqpoint{4.917128in}{2.240283in}}%
\pgfpathquadraticcurveto{\pgfqpoint{4.922426in}{2.246674in}}{\pgfqpoint{4.931827in}{2.248860in}}%
\pgfpathlineto{\pgfqpoint{4.931827in}{2.248860in}}%
\pgfpathclose%
\pgfpathmoveto{\pgfqpoint{4.927065in}{2.269414in}}%
\pgfpathquadraticcurveto{\pgfqpoint{4.927065in}{2.261951in}}{\pgfqpoint{4.931724in}{2.257766in}}%
\pgfpathquadraticcurveto{\pgfqpoint{4.936404in}{2.253581in}}{\pgfqpoint{4.944836in}{2.253581in}}%
\pgfpathquadraticcurveto{\pgfqpoint{4.953227in}{2.253581in}}{\pgfqpoint{4.957948in}{2.257766in}}%
\pgfpathquadraticcurveto{\pgfqpoint{4.962690in}{2.261951in}}{\pgfqpoint{4.962690in}{2.269414in}}%
\pgfpathquadraticcurveto{\pgfqpoint{4.962690in}{2.276898in}}{\pgfqpoint{4.957948in}{2.281083in}}%
\pgfpathquadraticcurveto{\pgfqpoint{4.953227in}{2.285268in}}{\pgfqpoint{4.944836in}{2.285268in}}%
\pgfpathquadraticcurveto{\pgfqpoint{4.936404in}{2.285268in}}{\pgfqpoint{4.931724in}{2.281083in}}%
\pgfpathquadraticcurveto{\pgfqpoint{4.927065in}{2.276898in}}{\pgfqpoint{4.927065in}{2.269414in}}%
\pgfpathlineto{\pgfqpoint{4.927065in}{2.269414in}}%
\pgfpathclose%
\pgfpathmoveto{\pgfqpoint{5.069689in}{2.297762in}}%
\pgfpathquadraticcurveto{\pgfqpoint{5.061071in}{2.282959in}}{\pgfqpoint{5.056865in}{2.268445in}}%
\pgfpathquadraticcurveto{\pgfqpoint{5.052680in}{2.253952in}}{\pgfqpoint{5.052680in}{2.239067in}}%
\pgfpathquadraticcurveto{\pgfqpoint{5.052680in}{2.224203in}}{\pgfqpoint{5.056906in}{2.209606in}}%
\pgfpathquadraticcurveto{\pgfqpoint{5.061133in}{2.195010in}}{\pgfqpoint{5.069689in}{2.180249in}}%
\pgfpathlineto{\pgfqpoint{5.059380in}{2.180249in}}%
\pgfpathquadraticcurveto{\pgfqpoint{5.049732in}{2.195402in}}{\pgfqpoint{5.044928in}{2.210019in}}%
\pgfpathquadraticcurveto{\pgfqpoint{5.040125in}{2.224636in}}{\pgfqpoint{5.040125in}{2.239067in}}%
\pgfpathquadraticcurveto{\pgfqpoint{5.040125in}{2.253437in}}{\pgfqpoint{5.044887in}{2.267992in}}%
\pgfpathquadraticcurveto{\pgfqpoint{5.049670in}{2.282567in}}{\pgfqpoint{5.059380in}{2.297762in}}%
\pgfpathlineto{\pgfqpoint{5.069689in}{2.297762in}}%
\pgfpathlineto{\pgfqpoint{5.069689in}{2.297762in}}%
\pgfpathclose%
\pgfpathmoveto{\pgfqpoint{5.105582in}{2.208596in}}%
\pgfpathlineto{\pgfqpoint{5.150999in}{2.208596in}}%
\pgfpathlineto{\pgfqpoint{5.150999in}{2.197649in}}%
\pgfpathlineto{\pgfqpoint{5.089934in}{2.197649in}}%
\pgfpathlineto{\pgfqpoint{5.089934in}{2.208596in}}%
\pgfpathquadraticcurveto{\pgfqpoint{5.097335in}{2.216265in}}{\pgfqpoint{5.110117in}{2.229171in}}%
\pgfpathquadraticcurveto{\pgfqpoint{5.122920in}{2.242098in}}{\pgfqpoint{5.126198in}{2.245850in}}%
\pgfpathquadraticcurveto{\pgfqpoint{5.132445in}{2.252859in}}{\pgfqpoint{5.134919in}{2.257725in}}%
\pgfpathquadraticcurveto{\pgfqpoint{5.137413in}{2.262590in}}{\pgfqpoint{5.137413in}{2.267291in}}%
\pgfpathquadraticcurveto{\pgfqpoint{5.137413in}{2.274960in}}{\pgfqpoint{5.132032in}{2.279784in}}%
\pgfpathquadraticcurveto{\pgfqpoint{5.126651in}{2.284629in}}{\pgfqpoint{5.118013in}{2.284629in}}%
\pgfpathquadraticcurveto{\pgfqpoint{5.111890in}{2.284629in}}{\pgfqpoint{5.105087in}{2.282506in}}%
\pgfpathquadraticcurveto{\pgfqpoint{5.098304in}{2.280382in}}{\pgfqpoint{5.090573in}{2.276053in}}%
\pgfpathlineto{\pgfqpoint{5.090573in}{2.289206in}}%
\pgfpathquadraticcurveto{\pgfqpoint{5.098428in}{2.292360in}}{\pgfqpoint{5.105252in}{2.293968in}}%
\pgfpathquadraticcurveto{\pgfqpoint{5.112096in}{2.295576in}}{\pgfqpoint{5.117766in}{2.295576in}}%
\pgfpathquadraticcurveto{\pgfqpoint{5.132713in}{2.295576in}}{\pgfqpoint{5.141598in}{2.288093in}}%
\pgfpathquadraticcurveto{\pgfqpoint{5.150484in}{2.280629in}}{\pgfqpoint{5.150484in}{2.268136in}}%
\pgfpathquadraticcurveto{\pgfqpoint{5.150484in}{2.262198in}}{\pgfqpoint{5.148257in}{2.256879in}}%
\pgfpathquadraticcurveto{\pgfqpoint{5.146051in}{2.251581in}}{\pgfqpoint{5.140176in}{2.244365in}}%
\pgfpathquadraticcurveto{\pgfqpoint{5.138568in}{2.242489in}}{\pgfqpoint{5.129929in}{2.233562in}}%
\pgfpathquadraticcurveto{\pgfqpoint{5.121312in}{2.224636in}}{\pgfqpoint{5.105582in}{2.208596in}}%
\pgfpathlineto{\pgfqpoint{5.105582in}{2.208596in}}%
\pgfpathclose%
\pgfpathmoveto{\pgfqpoint{5.206148in}{2.285268in}}%
\pgfpathquadraticcurveto{\pgfqpoint{5.196108in}{2.285268in}}{\pgfqpoint{5.191036in}{2.275372in}}%
\pgfpathquadraticcurveto{\pgfqpoint{5.185985in}{2.265497in}}{\pgfqpoint{5.185985in}{2.245644in}}%
\pgfpathquadraticcurveto{\pgfqpoint{5.185985in}{2.225873in}}{\pgfqpoint{5.191036in}{2.215977in}}%
\pgfpathquadraticcurveto{\pgfqpoint{5.196108in}{2.206081in}}{\pgfqpoint{5.206148in}{2.206081in}}%
\pgfpathquadraticcurveto{\pgfqpoint{5.216271in}{2.206081in}}{\pgfqpoint{5.221322in}{2.215977in}}%
\pgfpathquadraticcurveto{\pgfqpoint{5.226393in}{2.225873in}}{\pgfqpoint{5.226393in}{2.245644in}}%
\pgfpathquadraticcurveto{\pgfqpoint{5.226393in}{2.265497in}}{\pgfqpoint{5.221322in}{2.275372in}}%
\pgfpathquadraticcurveto{\pgfqpoint{5.216271in}{2.285268in}}{\pgfqpoint{5.206148in}{2.285268in}}%
\pgfpathlineto{\pgfqpoint{5.206148in}{2.285268in}}%
\pgfpathclose%
\pgfpathmoveto{\pgfqpoint{5.206148in}{2.295576in}}%
\pgfpathquadraticcurveto{\pgfqpoint{5.222332in}{2.295576in}}{\pgfqpoint{5.230867in}{2.282774in}}%
\pgfpathquadraticcurveto{\pgfqpoint{5.239402in}{2.269991in}}{\pgfqpoint{5.239402in}{2.245644in}}%
\pgfpathquadraticcurveto{\pgfqpoint{5.239402in}{2.221358in}}{\pgfqpoint{5.230867in}{2.208555in}}%
\pgfpathquadraticcurveto{\pgfqpoint{5.222332in}{2.195773in}}{\pgfqpoint{5.206148in}{2.195773in}}%
\pgfpathquadraticcurveto{\pgfqpoint{5.189985in}{2.195773in}}{\pgfqpoint{5.181450in}{2.208555in}}%
\pgfpathquadraticcurveto{\pgfqpoint{5.172914in}{2.221358in}}{\pgfqpoint{5.172914in}{2.245644in}}%
\pgfpathquadraticcurveto{\pgfqpoint{5.172914in}{2.269991in}}{\pgfqpoint{5.181450in}{2.282774in}}%
\pgfpathquadraticcurveto{\pgfqpoint{5.189985in}{2.295576in}}{\pgfqpoint{5.206148in}{2.295576in}}%
\pgfpathlineto{\pgfqpoint{5.206148in}{2.295576in}}%
\pgfpathclose%
\pgfpathmoveto{\pgfqpoint{5.273481in}{2.208596in}}%
\pgfpathlineto{\pgfqpoint{5.318899in}{2.208596in}}%
\pgfpathlineto{\pgfqpoint{5.318899in}{2.197649in}}%
\pgfpathlineto{\pgfqpoint{5.257833in}{2.197649in}}%
\pgfpathlineto{\pgfqpoint{5.257833in}{2.208596in}}%
\pgfpathquadraticcurveto{\pgfqpoint{5.265234in}{2.216265in}}{\pgfqpoint{5.278016in}{2.229171in}}%
\pgfpathquadraticcurveto{\pgfqpoint{5.290819in}{2.242098in}}{\pgfqpoint{5.294097in}{2.245850in}}%
\pgfpathquadraticcurveto{\pgfqpoint{5.300344in}{2.252859in}}{\pgfqpoint{5.302818in}{2.257725in}}%
\pgfpathquadraticcurveto{\pgfqpoint{5.305313in}{2.262590in}}{\pgfqpoint{5.305313in}{2.267291in}}%
\pgfpathquadraticcurveto{\pgfqpoint{5.305313in}{2.274960in}}{\pgfqpoint{5.299932in}{2.279784in}}%
\pgfpathquadraticcurveto{\pgfqpoint{5.294551in}{2.284629in}}{\pgfqpoint{5.285913in}{2.284629in}}%
\pgfpathquadraticcurveto{\pgfqpoint{5.279789in}{2.284629in}}{\pgfqpoint{5.272986in}{2.282506in}}%
\pgfpathquadraticcurveto{\pgfqpoint{5.266203in}{2.280382in}}{\pgfqpoint{5.258472in}{2.276053in}}%
\pgfpathlineto{\pgfqpoint{5.258472in}{2.289206in}}%
\pgfpathquadraticcurveto{\pgfqpoint{5.266327in}{2.292360in}}{\pgfqpoint{5.273151in}{2.293968in}}%
\pgfpathquadraticcurveto{\pgfqpoint{5.279996in}{2.295576in}}{\pgfqpoint{5.285665in}{2.295576in}}%
\pgfpathquadraticcurveto{\pgfqpoint{5.300612in}{2.295576in}}{\pgfqpoint{5.309498in}{2.288093in}}%
\pgfpathquadraticcurveto{\pgfqpoint{5.318383in}{2.280629in}}{\pgfqpoint{5.318383in}{2.268136in}}%
\pgfpathquadraticcurveto{\pgfqpoint{5.318383in}{2.262198in}}{\pgfqpoint{5.316157in}{2.256879in}}%
\pgfpathquadraticcurveto{\pgfqpoint{5.313951in}{2.251581in}}{\pgfqpoint{5.308075in}{2.244365in}}%
\pgfpathquadraticcurveto{\pgfqpoint{5.306467in}{2.242489in}}{\pgfqpoint{5.297829in}{2.233562in}}%
\pgfpathquadraticcurveto{\pgfqpoint{5.289211in}{2.224636in}}{\pgfqpoint{5.273481in}{2.208596in}}%
\pgfpathlineto{\pgfqpoint{5.273481in}{2.208596in}}%
\pgfpathclose%
\pgfpathmoveto{\pgfqpoint{5.374047in}{2.285268in}}%
\pgfpathquadraticcurveto{\pgfqpoint{5.364007in}{2.285268in}}{\pgfqpoint{5.358936in}{2.275372in}}%
\pgfpathquadraticcurveto{\pgfqpoint{5.353885in}{2.265497in}}{\pgfqpoint{5.353885in}{2.245644in}}%
\pgfpathquadraticcurveto{\pgfqpoint{5.353885in}{2.225873in}}{\pgfqpoint{5.358936in}{2.215977in}}%
\pgfpathquadraticcurveto{\pgfqpoint{5.364007in}{2.206081in}}{\pgfqpoint{5.374047in}{2.206081in}}%
\pgfpathquadraticcurveto{\pgfqpoint{5.384170in}{2.206081in}}{\pgfqpoint{5.389221in}{2.215977in}}%
\pgfpathquadraticcurveto{\pgfqpoint{5.394293in}{2.225873in}}{\pgfqpoint{5.394293in}{2.245644in}}%
\pgfpathquadraticcurveto{\pgfqpoint{5.394293in}{2.265497in}}{\pgfqpoint{5.389221in}{2.275372in}}%
\pgfpathquadraticcurveto{\pgfqpoint{5.384170in}{2.285268in}}{\pgfqpoint{5.374047in}{2.285268in}}%
\pgfpathlineto{\pgfqpoint{5.374047in}{2.285268in}}%
\pgfpathclose%
\pgfpathmoveto{\pgfqpoint{5.374047in}{2.295576in}}%
\pgfpathquadraticcurveto{\pgfqpoint{5.390231in}{2.295576in}}{\pgfqpoint{5.398766in}{2.282774in}}%
\pgfpathquadraticcurveto{\pgfqpoint{5.407301in}{2.269991in}}{\pgfqpoint{5.407301in}{2.245644in}}%
\pgfpathquadraticcurveto{\pgfqpoint{5.407301in}{2.221358in}}{\pgfqpoint{5.398766in}{2.208555in}}%
\pgfpathquadraticcurveto{\pgfqpoint{5.390231in}{2.195773in}}{\pgfqpoint{5.374047in}{2.195773in}}%
\pgfpathquadraticcurveto{\pgfqpoint{5.357884in}{2.195773in}}{\pgfqpoint{5.349349in}{2.208555in}}%
\pgfpathquadraticcurveto{\pgfqpoint{5.340814in}{2.221358in}}{\pgfqpoint{5.340814in}{2.245644in}}%
\pgfpathquadraticcurveto{\pgfqpoint{5.340814in}{2.269991in}}{\pgfqpoint{5.349349in}{2.282774in}}%
\pgfpathquadraticcurveto{\pgfqpoint{5.357884in}{2.295576in}}{\pgfqpoint{5.374047in}{2.295576in}}%
\pgfpathlineto{\pgfqpoint{5.374047in}{2.295576in}}%
\pgfpathclose%
\pgfpathmoveto{\pgfqpoint{5.426640in}{2.297762in}}%
\pgfpathlineto{\pgfqpoint{5.436948in}{2.297762in}}%
\pgfpathquadraticcurveto{\pgfqpoint{5.446596in}{2.282567in}}{\pgfqpoint{5.451400in}{2.267992in}}%
\pgfpathquadraticcurveto{\pgfqpoint{5.456203in}{2.253437in}}{\pgfqpoint{5.456203in}{2.239067in}}%
\pgfpathquadraticcurveto{\pgfqpoint{5.456203in}{2.224636in}}{\pgfqpoint{5.451400in}{2.210019in}}%
\pgfpathquadraticcurveto{\pgfqpoint{5.446596in}{2.195402in}}{\pgfqpoint{5.436948in}{2.180249in}}%
\pgfpathlineto{\pgfqpoint{5.426640in}{2.180249in}}%
\pgfpathquadraticcurveto{\pgfqpoint{5.435195in}{2.195010in}}{\pgfqpoint{5.439422in}{2.209606in}}%
\pgfpathquadraticcurveto{\pgfqpoint{5.443648in}{2.224203in}}{\pgfqpoint{5.443648in}{2.239067in}}%
\pgfpathquadraticcurveto{\pgfqpoint{5.443648in}{2.253952in}}{\pgfqpoint{5.439422in}{2.268445in}}%
\pgfpathquadraticcurveto{\pgfqpoint{5.435195in}{2.282959in}}{\pgfqpoint{5.426640in}{2.297762in}}%
\pgfpathlineto{\pgfqpoint{5.426640in}{2.297762in}}%
\pgfpathclose%
\pgfusepath{fill}%
\end{pgfscope}%
\begin{pgfscope}%
\pgfsetbuttcap%
\pgfsetroundjoin%
\definecolor{currentfill}{rgb}{0.090196,0.168627,0.227451}%
\pgfsetfillcolor{currentfill}%
\pgfsetlinewidth{0.000000pt}%
\definecolor{currentstroke}{rgb}{0.090196,0.168627,0.227451}%
\pgfsetstrokecolor{currentstroke}%
\pgfsetdash{}{0pt}%
\pgfpathmoveto{\pgfqpoint{4.550611in}{2.114050in}}%
\pgfpathlineto{\pgfqpoint{4.550611in}{2.101351in}}%
\pgfpathquadraticcurveto{\pgfqpoint{4.543210in}{2.104897in}}{\pgfqpoint{4.536633in}{2.106628in}}%
\pgfpathquadraticcurveto{\pgfqpoint{4.530056in}{2.108381in}}{\pgfqpoint{4.523933in}{2.108381in}}%
\pgfpathquadraticcurveto{\pgfqpoint{4.513316in}{2.108381in}}{\pgfqpoint{4.507543in}{2.104257in}}%
\pgfpathquadraticcurveto{\pgfqpoint{4.501771in}{2.100134in}}{\pgfqpoint{4.501771in}{2.092527in}}%
\pgfpathquadraticcurveto{\pgfqpoint{4.501771in}{2.086136in}}{\pgfqpoint{4.505605in}{2.082878in}}%
\pgfpathquadraticcurveto{\pgfqpoint{4.509440in}{2.079642in}}{\pgfqpoint{4.520140in}{2.077642in}}%
\pgfpathlineto{\pgfqpoint{4.527995in}{2.076034in}}%
\pgfpathquadraticcurveto{\pgfqpoint{4.542550in}{2.073251in}}{\pgfqpoint{4.549477in}{2.066262in}}%
\pgfpathquadraticcurveto{\pgfqpoint{4.556404in}{2.059273in}}{\pgfqpoint{4.556404in}{2.047563in}}%
\pgfpathquadraticcurveto{\pgfqpoint{4.556404in}{2.033564in}}{\pgfqpoint{4.547024in}{2.026348in}}%
\pgfpathquadraticcurveto{\pgfqpoint{4.537664in}{2.019133in}}{\pgfqpoint{4.519563in}{2.019133in}}%
\pgfpathquadraticcurveto{\pgfqpoint{4.512739in}{2.019133in}}{\pgfqpoint{4.505028in}{2.020679in}}%
\pgfpathquadraticcurveto{\pgfqpoint{4.497338in}{2.022225in}}{\pgfqpoint{4.489092in}{2.025256in}}%
\pgfpathlineto{\pgfqpoint{4.489092in}{2.038656in}}%
\pgfpathquadraticcurveto{\pgfqpoint{4.497008in}{2.034224in}}{\pgfqpoint{4.504616in}{2.031956in}}%
\pgfpathquadraticcurveto{\pgfqpoint{4.512223in}{2.029709in}}{\pgfqpoint{4.519563in}{2.029709in}}%
\pgfpathquadraticcurveto{\pgfqpoint{4.530696in}{2.029709in}}{\pgfqpoint{4.536757in}{2.034080in}}%
\pgfpathquadraticcurveto{\pgfqpoint{4.542818in}{2.038471in}}{\pgfqpoint{4.542818in}{2.046594in}}%
\pgfpathquadraticcurveto{\pgfqpoint{4.542818in}{2.053665in}}{\pgfqpoint{4.538468in}{2.057665in}}%
\pgfpathquadraticcurveto{\pgfqpoint{4.534118in}{2.061664in}}{\pgfqpoint{4.524201in}{2.063664in}}%
\pgfpathlineto{\pgfqpoint{4.516264in}{2.065210in}}%
\pgfpathquadraticcurveto{\pgfqpoint{4.501709in}{2.068096in}}{\pgfqpoint{4.495194in}{2.074281in}}%
\pgfpathquadraticcurveto{\pgfqpoint{4.488700in}{2.080466in}}{\pgfqpoint{4.488700in}{2.091496in}}%
\pgfpathquadraticcurveto{\pgfqpoint{4.488700in}{2.104257in}}{\pgfqpoint{4.497689in}{2.111597in}}%
\pgfpathquadraticcurveto{\pgfqpoint{4.506678in}{2.118936in}}{\pgfqpoint{4.522449in}{2.118936in}}%
\pgfpathquadraticcurveto{\pgfqpoint{4.529232in}{2.118936in}}{\pgfqpoint{4.536241in}{2.117699in}}%
\pgfpathquadraticcurveto{\pgfqpoint{4.543271in}{2.116483in}}{\pgfqpoint{4.550611in}{2.114050in}}%
\pgfpathlineto{\pgfqpoint{4.550611in}{2.114050in}}%
\pgfpathclose%
\pgfpathmoveto{\pgfqpoint{4.587926in}{2.113659in}}%
\pgfpathlineto{\pgfqpoint{4.587926in}{2.093166in}}%
\pgfpathlineto{\pgfqpoint{4.612336in}{2.093166in}}%
\pgfpathlineto{\pgfqpoint{4.612336in}{2.083950in}}%
\pgfpathlineto{\pgfqpoint{4.587926in}{2.083950in}}%
\pgfpathlineto{\pgfqpoint{4.587926in}{2.044779in}}%
\pgfpathquadraticcurveto{\pgfqpoint{4.587926in}{2.035956in}}{\pgfqpoint{4.590339in}{2.033440in}}%
\pgfpathquadraticcurveto{\pgfqpoint{4.592751in}{2.030925in}}{\pgfqpoint{4.600173in}{2.030925in}}%
\pgfpathlineto{\pgfqpoint{4.612336in}{2.030925in}}%
\pgfpathlineto{\pgfqpoint{4.612336in}{2.021009in}}%
\pgfpathlineto{\pgfqpoint{4.600173in}{2.021009in}}%
\pgfpathquadraticcurveto{\pgfqpoint{4.586442in}{2.021009in}}{\pgfqpoint{4.581226in}{2.026122in}}%
\pgfpathquadraticcurveto{\pgfqpoint{4.576010in}{2.031255in}}{\pgfqpoint{4.576010in}{2.044779in}}%
\pgfpathlineto{\pgfqpoint{4.576010in}{2.083950in}}%
\pgfpathlineto{\pgfqpoint{4.567310in}{2.083950in}}%
\pgfpathlineto{\pgfqpoint{4.567310in}{2.093166in}}%
\pgfpathlineto{\pgfqpoint{4.576010in}{2.093166in}}%
\pgfpathlineto{\pgfqpoint{4.576010in}{2.113659in}}%
\pgfpathlineto{\pgfqpoint{4.587926in}{2.113659in}}%
\pgfpathlineto{\pgfqpoint{4.587926in}{2.113659in}}%
\pgfpathclose%
\pgfpathmoveto{\pgfqpoint{4.626706in}{2.049480in}}%
\pgfpathlineto{\pgfqpoint{4.626706in}{2.093166in}}%
\pgfpathlineto{\pgfqpoint{4.638560in}{2.093166in}}%
\pgfpathlineto{\pgfqpoint{4.638560in}{2.049933in}}%
\pgfpathquadraticcurveto{\pgfqpoint{4.638560in}{2.039687in}}{\pgfqpoint{4.642539in}{2.034554in}}%
\pgfpathquadraticcurveto{\pgfqpoint{4.646539in}{2.029441in}}{\pgfqpoint{4.654538in}{2.029441in}}%
\pgfpathquadraticcurveto{\pgfqpoint{4.664124in}{2.029441in}}{\pgfqpoint{4.669691in}{2.035564in}}%
\pgfpathquadraticcurveto{\pgfqpoint{4.675278in}{2.041687in}}{\pgfqpoint{4.675278in}{2.052263in}}%
\pgfpathlineto{\pgfqpoint{4.675278in}{2.093166in}}%
\pgfpathlineto{\pgfqpoint{4.687132in}{2.093166in}}%
\pgfpathlineto{\pgfqpoint{4.687132in}{2.021009in}}%
\pgfpathlineto{\pgfqpoint{4.675278in}{2.021009in}}%
\pgfpathlineto{\pgfqpoint{4.675278in}{2.032100in}}%
\pgfpathquadraticcurveto{\pgfqpoint{4.670969in}{2.025524in}}{\pgfqpoint{4.665258in}{2.022328in}}%
\pgfpathquadraticcurveto{\pgfqpoint{4.659568in}{2.019133in}}{\pgfqpoint{4.652023in}{2.019133in}}%
\pgfpathquadraticcurveto{\pgfqpoint{4.639591in}{2.019133in}}{\pgfqpoint{4.633138in}{2.026864in}}%
\pgfpathquadraticcurveto{\pgfqpoint{4.626706in}{2.034595in}}{\pgfqpoint{4.626706in}{2.049480in}}%
\pgfpathlineto{\pgfqpoint{4.626706in}{2.049480in}}%
\pgfpathclose%
\pgfpathmoveto{\pgfqpoint{4.656538in}{2.094898in}}%
\pgfpathlineto{\pgfqpoint{4.656538in}{2.094898in}}%
\pgfpathlineto{\pgfqpoint{4.656538in}{2.094898in}}%
\pgfpathclose%
\pgfpathmoveto{\pgfqpoint{4.759021in}{2.082219in}}%
\pgfpathlineto{\pgfqpoint{4.759021in}{2.121266in}}%
\pgfpathlineto{\pgfqpoint{4.770876in}{2.121266in}}%
\pgfpathlineto{\pgfqpoint{4.770876in}{2.021009in}}%
\pgfpathlineto{\pgfqpoint{4.759021in}{2.021009in}}%
\pgfpathlineto{\pgfqpoint{4.759021in}{2.031832in}}%
\pgfpathquadraticcurveto{\pgfqpoint{4.755290in}{2.025400in}}{\pgfqpoint{4.749579in}{2.022266in}}%
\pgfpathquadraticcurveto{\pgfqpoint{4.743889in}{2.019133in}}{\pgfqpoint{4.735890in}{2.019133in}}%
\pgfpathquadraticcurveto{\pgfqpoint{4.722819in}{2.019133in}}{\pgfqpoint{4.714593in}{2.029565in}}%
\pgfpathquadraticcurveto{\pgfqpoint{4.706388in}{2.040017in}}{\pgfqpoint{4.706388in}{2.057026in}}%
\pgfpathquadraticcurveto{\pgfqpoint{4.706388in}{2.074034in}}{\pgfqpoint{4.714593in}{2.084466in}}%
\pgfpathquadraticcurveto{\pgfqpoint{4.722819in}{2.094898in}}{\pgfqpoint{4.735890in}{2.094898in}}%
\pgfpathquadraticcurveto{\pgfqpoint{4.743889in}{2.094898in}}{\pgfqpoint{4.749579in}{2.091764in}}%
\pgfpathquadraticcurveto{\pgfqpoint{4.755290in}{2.088651in}}{\pgfqpoint{4.759021in}{2.082219in}}%
\pgfpathlineto{\pgfqpoint{4.759021in}{2.082219in}}%
\pgfpathclose%
\pgfpathmoveto{\pgfqpoint{4.718634in}{2.057026in}}%
\pgfpathquadraticcurveto{\pgfqpoint{4.718634in}{2.043955in}}{\pgfqpoint{4.724015in}{2.036512in}}%
\pgfpathquadraticcurveto{\pgfqpoint{4.729396in}{2.029070in}}{\pgfqpoint{4.738797in}{2.029070in}}%
\pgfpathquadraticcurveto{\pgfqpoint{4.748198in}{2.029070in}}{\pgfqpoint{4.753599in}{2.036512in}}%
\pgfpathquadraticcurveto{\pgfqpoint{4.759021in}{2.043955in}}{\pgfqpoint{4.759021in}{2.057026in}}%
\pgfpathquadraticcurveto{\pgfqpoint{4.759021in}{2.070096in}}{\pgfqpoint{4.753599in}{2.077539in}}%
\pgfpathquadraticcurveto{\pgfqpoint{4.748198in}{2.084981in}}{\pgfqpoint{4.738797in}{2.084981in}}%
\pgfpathquadraticcurveto{\pgfqpoint{4.729396in}{2.084981in}}{\pgfqpoint{4.724015in}{2.077539in}}%
\pgfpathquadraticcurveto{\pgfqpoint{4.718634in}{2.070096in}}{\pgfqpoint{4.718634in}{2.057026in}}%
\pgfpathlineto{\pgfqpoint{4.718634in}{2.057026in}}%
\pgfpathclose%
\pgfpathmoveto{\pgfqpoint{4.825323in}{2.014308in}}%
\pgfpathquadraticcurveto{\pgfqpoint{4.820314in}{2.001423in}}{\pgfqpoint{4.815531in}{1.997506in}}%
\pgfpathquadraticcurveto{\pgfqpoint{4.810768in}{1.993568in}}{\pgfqpoint{4.802790in}{1.993568in}}%
\pgfpathlineto{\pgfqpoint{4.793306in}{1.993568in}}%
\pgfpathlineto{\pgfqpoint{4.793306in}{2.003485in}}%
\pgfpathlineto{\pgfqpoint{4.800275in}{2.003485in}}%
\pgfpathquadraticcurveto{\pgfqpoint{4.805161in}{2.003485in}}{\pgfqpoint{4.807861in}{2.005815in}}%
\pgfpathquadraticcurveto{\pgfqpoint{4.810583in}{2.008124in}}{\pgfqpoint{4.813861in}{2.016762in}}%
\pgfpathlineto{\pgfqpoint{4.815984in}{2.022163in}}%
\pgfpathlineto{\pgfqpoint{4.786812in}{2.093166in}}%
\pgfpathlineto{\pgfqpoint{4.799367in}{2.093166in}}%
\pgfpathlineto{\pgfqpoint{4.821922in}{2.036739in}}%
\pgfpathlineto{\pgfqpoint{4.844476in}{2.093166in}}%
\pgfpathlineto{\pgfqpoint{4.857031in}{2.093166in}}%
\pgfpathlineto{\pgfqpoint{4.825323in}{2.014308in}}%
\pgfpathlineto{\pgfqpoint{4.825323in}{2.014308in}}%
\pgfpathclose%
\pgfpathmoveto{\pgfqpoint{4.917396in}{2.023009in}}%
\pgfpathlineto{\pgfqpoint{4.917396in}{2.034863in}}%
\pgfpathquadraticcurveto{\pgfqpoint{4.922303in}{2.032533in}}{\pgfqpoint{4.927312in}{2.031317in}}%
\pgfpathquadraticcurveto{\pgfqpoint{4.932343in}{2.030101in}}{\pgfqpoint{4.937188in}{2.030101in}}%
\pgfpathquadraticcurveto{\pgfqpoint{4.950073in}{2.030101in}}{\pgfqpoint{4.956855in}{2.038759in}}%
\pgfpathquadraticcurveto{\pgfqpoint{4.963659in}{2.047418in}}{\pgfqpoint{4.964628in}{2.065086in}}%
\pgfpathquadraticcurveto{\pgfqpoint{4.960896in}{2.059541in}}{\pgfqpoint{4.955144in}{2.056572in}}%
\pgfpathquadraticcurveto{\pgfqpoint{4.949413in}{2.053603in}}{\pgfqpoint{4.942465in}{2.053603in}}%
\pgfpathquadraticcurveto{\pgfqpoint{4.928034in}{2.053603in}}{\pgfqpoint{4.919622in}{2.062324in}}%
\pgfpathquadraticcurveto{\pgfqpoint{4.911211in}{2.071065in}}{\pgfqpoint{4.911211in}{2.086218in}}%
\pgfpathquadraticcurveto{\pgfqpoint{4.911211in}{2.101021in}}{\pgfqpoint{4.919973in}{2.109968in}}%
\pgfpathquadraticcurveto{\pgfqpoint{4.928735in}{2.118936in}}{\pgfqpoint{4.943290in}{2.118936in}}%
\pgfpathquadraticcurveto{\pgfqpoint{4.959989in}{2.118936in}}{\pgfqpoint{4.968772in}{2.106134in}}%
\pgfpathquadraticcurveto{\pgfqpoint{4.977575in}{2.093351in}}{\pgfqpoint{4.977575in}{2.069004in}}%
\pgfpathquadraticcurveto{\pgfqpoint{4.977575in}{2.046264in}}{\pgfqpoint{4.966772in}{2.032698in}}%
\pgfpathquadraticcurveto{\pgfqpoint{4.955990in}{2.019133in}}{\pgfqpoint{4.937765in}{2.019133in}}%
\pgfpathquadraticcurveto{\pgfqpoint{4.932858in}{2.019133in}}{\pgfqpoint{4.927828in}{2.020102in}}%
\pgfpathquadraticcurveto{\pgfqpoint{4.922818in}{2.021071in}}{\pgfqpoint{4.917396in}{2.023009in}}%
\pgfpathlineto{\pgfqpoint{4.917396in}{2.023009in}}%
\pgfpathclose%
\pgfpathmoveto{\pgfqpoint{4.943290in}{2.063788in}}%
\pgfpathquadraticcurveto{\pgfqpoint{4.952052in}{2.063788in}}{\pgfqpoint{4.957165in}{2.069766in}}%
\pgfpathquadraticcurveto{\pgfqpoint{4.962298in}{2.075766in}}{\pgfqpoint{4.962298in}{2.086218in}}%
\pgfpathquadraticcurveto{\pgfqpoint{4.962298in}{2.096588in}}{\pgfqpoint{4.957165in}{2.102608in}}%
\pgfpathquadraticcurveto{\pgfqpoint{4.952052in}{2.108628in}}{\pgfqpoint{4.943290in}{2.108628in}}%
\pgfpathquadraticcurveto{\pgfqpoint{4.934528in}{2.108628in}}{\pgfqpoint{4.929415in}{2.102608in}}%
\pgfpathquadraticcurveto{\pgfqpoint{4.924302in}{2.096588in}}{\pgfqpoint{4.924302in}{2.086218in}}%
\pgfpathquadraticcurveto{\pgfqpoint{4.924302in}{2.075766in}}{\pgfqpoint{4.929415in}{2.069766in}}%
\pgfpathquadraticcurveto{\pgfqpoint{4.934528in}{2.063788in}}{\pgfqpoint{4.943290in}{2.063788in}}%
\pgfpathlineto{\pgfqpoint{4.943290in}{2.063788in}}%
\pgfpathclose%
\pgfpathmoveto{\pgfqpoint{5.069689in}{2.121122in}}%
\pgfpathquadraticcurveto{\pgfqpoint{5.061071in}{2.106319in}}{\pgfqpoint{5.056865in}{2.091805in}}%
\pgfpathquadraticcurveto{\pgfqpoint{5.052680in}{2.077312in}}{\pgfqpoint{5.052680in}{2.062427in}}%
\pgfpathquadraticcurveto{\pgfqpoint{5.052680in}{2.047563in}}{\pgfqpoint{5.056906in}{2.032966in}}%
\pgfpathquadraticcurveto{\pgfqpoint{5.061133in}{2.018370in}}{\pgfqpoint{5.069689in}{2.003609in}}%
\pgfpathlineto{\pgfqpoint{5.059380in}{2.003609in}}%
\pgfpathquadraticcurveto{\pgfqpoint{5.049732in}{2.018762in}}{\pgfqpoint{5.044928in}{2.033379in}}%
\pgfpathquadraticcurveto{\pgfqpoint{5.040125in}{2.047996in}}{\pgfqpoint{5.040125in}{2.062427in}}%
\pgfpathquadraticcurveto{\pgfqpoint{5.040125in}{2.076797in}}{\pgfqpoint{5.044887in}{2.091352in}}%
\pgfpathquadraticcurveto{\pgfqpoint{5.049670in}{2.105927in}}{\pgfqpoint{5.059380in}{2.121122in}}%
\pgfpathlineto{\pgfqpoint{5.069689in}{2.121122in}}%
\pgfpathlineto{\pgfqpoint{5.069689in}{2.121122in}}%
\pgfpathclose%
\pgfpathmoveto{\pgfqpoint{5.105582in}{2.031956in}}%
\pgfpathlineto{\pgfqpoint{5.150999in}{2.031956in}}%
\pgfpathlineto{\pgfqpoint{5.150999in}{2.021009in}}%
\pgfpathlineto{\pgfqpoint{5.089934in}{2.021009in}}%
\pgfpathlineto{\pgfqpoint{5.089934in}{2.031956in}}%
\pgfpathquadraticcurveto{\pgfqpoint{5.097335in}{2.039625in}}{\pgfqpoint{5.110117in}{2.052531in}}%
\pgfpathquadraticcurveto{\pgfqpoint{5.122920in}{2.065458in}}{\pgfqpoint{5.126198in}{2.069210in}}%
\pgfpathquadraticcurveto{\pgfqpoint{5.132445in}{2.076219in}}{\pgfqpoint{5.134919in}{2.081085in}}%
\pgfpathquadraticcurveto{\pgfqpoint{5.137413in}{2.085950in}}{\pgfqpoint{5.137413in}{2.090651in}}%
\pgfpathquadraticcurveto{\pgfqpoint{5.137413in}{2.098320in}}{\pgfqpoint{5.132032in}{2.103144in}}%
\pgfpathquadraticcurveto{\pgfqpoint{5.126651in}{2.107989in}}{\pgfqpoint{5.118013in}{2.107989in}}%
\pgfpathquadraticcurveto{\pgfqpoint{5.111890in}{2.107989in}}{\pgfqpoint{5.105087in}{2.105866in}}%
\pgfpathquadraticcurveto{\pgfqpoint{5.098304in}{2.103742in}}{\pgfqpoint{5.090573in}{2.099413in}}%
\pgfpathlineto{\pgfqpoint{5.090573in}{2.112566in}}%
\pgfpathquadraticcurveto{\pgfqpoint{5.098428in}{2.115720in}}{\pgfqpoint{5.105252in}{2.117328in}}%
\pgfpathquadraticcurveto{\pgfqpoint{5.112096in}{2.118936in}}{\pgfqpoint{5.117766in}{2.118936in}}%
\pgfpathquadraticcurveto{\pgfqpoint{5.132713in}{2.118936in}}{\pgfqpoint{5.141598in}{2.111453in}}%
\pgfpathquadraticcurveto{\pgfqpoint{5.150484in}{2.103989in}}{\pgfqpoint{5.150484in}{2.091496in}}%
\pgfpathquadraticcurveto{\pgfqpoint{5.150484in}{2.085558in}}{\pgfqpoint{5.148257in}{2.080239in}}%
\pgfpathquadraticcurveto{\pgfqpoint{5.146051in}{2.074941in}}{\pgfqpoint{5.140176in}{2.067725in}}%
\pgfpathquadraticcurveto{\pgfqpoint{5.138568in}{2.065849in}}{\pgfqpoint{5.129929in}{2.056922in}}%
\pgfpathquadraticcurveto{\pgfqpoint{5.121312in}{2.047996in}}{\pgfqpoint{5.105582in}{2.031956in}}%
\pgfpathlineto{\pgfqpoint{5.105582in}{2.031956in}}%
\pgfpathclose%
\pgfpathmoveto{\pgfqpoint{5.206148in}{2.108628in}}%
\pgfpathquadraticcurveto{\pgfqpoint{5.196108in}{2.108628in}}{\pgfqpoint{5.191036in}{2.098732in}}%
\pgfpathquadraticcurveto{\pgfqpoint{5.185985in}{2.088857in}}{\pgfqpoint{5.185985in}{2.069004in}}%
\pgfpathquadraticcurveto{\pgfqpoint{5.185985in}{2.049233in}}{\pgfqpoint{5.191036in}{2.039337in}}%
\pgfpathquadraticcurveto{\pgfqpoint{5.196108in}{2.029441in}}{\pgfqpoint{5.206148in}{2.029441in}}%
\pgfpathquadraticcurveto{\pgfqpoint{5.216271in}{2.029441in}}{\pgfqpoint{5.221322in}{2.039337in}}%
\pgfpathquadraticcurveto{\pgfqpoint{5.226393in}{2.049233in}}{\pgfqpoint{5.226393in}{2.069004in}}%
\pgfpathquadraticcurveto{\pgfqpoint{5.226393in}{2.088857in}}{\pgfqpoint{5.221322in}{2.098732in}}%
\pgfpathquadraticcurveto{\pgfqpoint{5.216271in}{2.108628in}}{\pgfqpoint{5.206148in}{2.108628in}}%
\pgfpathlineto{\pgfqpoint{5.206148in}{2.108628in}}%
\pgfpathclose%
\pgfpathmoveto{\pgfqpoint{5.206148in}{2.118936in}}%
\pgfpathquadraticcurveto{\pgfqpoint{5.222332in}{2.118936in}}{\pgfqpoint{5.230867in}{2.106134in}}%
\pgfpathquadraticcurveto{\pgfqpoint{5.239402in}{2.093351in}}{\pgfqpoint{5.239402in}{2.069004in}}%
\pgfpathquadraticcurveto{\pgfqpoint{5.239402in}{2.044718in}}{\pgfqpoint{5.230867in}{2.031915in}}%
\pgfpathquadraticcurveto{\pgfqpoint{5.222332in}{2.019133in}}{\pgfqpoint{5.206148in}{2.019133in}}%
\pgfpathquadraticcurveto{\pgfqpoint{5.189985in}{2.019133in}}{\pgfqpoint{5.181450in}{2.031915in}}%
\pgfpathquadraticcurveto{\pgfqpoint{5.172914in}{2.044718in}}{\pgfqpoint{5.172914in}{2.069004in}}%
\pgfpathquadraticcurveto{\pgfqpoint{5.172914in}{2.093351in}}{\pgfqpoint{5.181450in}{2.106134in}}%
\pgfpathquadraticcurveto{\pgfqpoint{5.189985in}{2.118936in}}{\pgfqpoint{5.206148in}{2.118936in}}%
\pgfpathlineto{\pgfqpoint{5.206148in}{2.118936in}}%
\pgfpathclose%
\pgfpathmoveto{\pgfqpoint{5.273481in}{2.031956in}}%
\pgfpathlineto{\pgfqpoint{5.318899in}{2.031956in}}%
\pgfpathlineto{\pgfqpoint{5.318899in}{2.021009in}}%
\pgfpathlineto{\pgfqpoint{5.257833in}{2.021009in}}%
\pgfpathlineto{\pgfqpoint{5.257833in}{2.031956in}}%
\pgfpathquadraticcurveto{\pgfqpoint{5.265234in}{2.039625in}}{\pgfqpoint{5.278016in}{2.052531in}}%
\pgfpathquadraticcurveto{\pgfqpoint{5.290819in}{2.065458in}}{\pgfqpoint{5.294097in}{2.069210in}}%
\pgfpathquadraticcurveto{\pgfqpoint{5.300344in}{2.076219in}}{\pgfqpoint{5.302818in}{2.081085in}}%
\pgfpathquadraticcurveto{\pgfqpoint{5.305313in}{2.085950in}}{\pgfqpoint{5.305313in}{2.090651in}}%
\pgfpathquadraticcurveto{\pgfqpoint{5.305313in}{2.098320in}}{\pgfqpoint{5.299932in}{2.103144in}}%
\pgfpathquadraticcurveto{\pgfqpoint{5.294551in}{2.107989in}}{\pgfqpoint{5.285913in}{2.107989in}}%
\pgfpathquadraticcurveto{\pgfqpoint{5.279789in}{2.107989in}}{\pgfqpoint{5.272986in}{2.105866in}}%
\pgfpathquadraticcurveto{\pgfqpoint{5.266203in}{2.103742in}}{\pgfqpoint{5.258472in}{2.099413in}}%
\pgfpathlineto{\pgfqpoint{5.258472in}{2.112566in}}%
\pgfpathquadraticcurveto{\pgfqpoint{5.266327in}{2.115720in}}{\pgfqpoint{5.273151in}{2.117328in}}%
\pgfpathquadraticcurveto{\pgfqpoint{5.279996in}{2.118936in}}{\pgfqpoint{5.285665in}{2.118936in}}%
\pgfpathquadraticcurveto{\pgfqpoint{5.300612in}{2.118936in}}{\pgfqpoint{5.309498in}{2.111453in}}%
\pgfpathquadraticcurveto{\pgfqpoint{5.318383in}{2.103989in}}{\pgfqpoint{5.318383in}{2.091496in}}%
\pgfpathquadraticcurveto{\pgfqpoint{5.318383in}{2.085558in}}{\pgfqpoint{5.316157in}{2.080239in}}%
\pgfpathquadraticcurveto{\pgfqpoint{5.313951in}{2.074941in}}{\pgfqpoint{5.308075in}{2.067725in}}%
\pgfpathquadraticcurveto{\pgfqpoint{5.306467in}{2.065849in}}{\pgfqpoint{5.297829in}{2.056922in}}%
\pgfpathquadraticcurveto{\pgfqpoint{5.289211in}{2.047996in}}{\pgfqpoint{5.273481in}{2.031956in}}%
\pgfpathlineto{\pgfqpoint{5.273481in}{2.031956in}}%
\pgfpathclose%
\pgfpathmoveto{\pgfqpoint{5.374047in}{2.108628in}}%
\pgfpathquadraticcurveto{\pgfqpoint{5.364007in}{2.108628in}}{\pgfqpoint{5.358936in}{2.098732in}}%
\pgfpathquadraticcurveto{\pgfqpoint{5.353885in}{2.088857in}}{\pgfqpoint{5.353885in}{2.069004in}}%
\pgfpathquadraticcurveto{\pgfqpoint{5.353885in}{2.049233in}}{\pgfqpoint{5.358936in}{2.039337in}}%
\pgfpathquadraticcurveto{\pgfqpoint{5.364007in}{2.029441in}}{\pgfqpoint{5.374047in}{2.029441in}}%
\pgfpathquadraticcurveto{\pgfqpoint{5.384170in}{2.029441in}}{\pgfqpoint{5.389221in}{2.039337in}}%
\pgfpathquadraticcurveto{\pgfqpoint{5.394293in}{2.049233in}}{\pgfqpoint{5.394293in}{2.069004in}}%
\pgfpathquadraticcurveto{\pgfqpoint{5.394293in}{2.088857in}}{\pgfqpoint{5.389221in}{2.098732in}}%
\pgfpathquadraticcurveto{\pgfqpoint{5.384170in}{2.108628in}}{\pgfqpoint{5.374047in}{2.108628in}}%
\pgfpathlineto{\pgfqpoint{5.374047in}{2.108628in}}%
\pgfpathclose%
\pgfpathmoveto{\pgfqpoint{5.374047in}{2.118936in}}%
\pgfpathquadraticcurveto{\pgfqpoint{5.390231in}{2.118936in}}{\pgfqpoint{5.398766in}{2.106134in}}%
\pgfpathquadraticcurveto{\pgfqpoint{5.407301in}{2.093351in}}{\pgfqpoint{5.407301in}{2.069004in}}%
\pgfpathquadraticcurveto{\pgfqpoint{5.407301in}{2.044718in}}{\pgfqpoint{5.398766in}{2.031915in}}%
\pgfpathquadraticcurveto{\pgfqpoint{5.390231in}{2.019133in}}{\pgfqpoint{5.374047in}{2.019133in}}%
\pgfpathquadraticcurveto{\pgfqpoint{5.357884in}{2.019133in}}{\pgfqpoint{5.349349in}{2.031915in}}%
\pgfpathquadraticcurveto{\pgfqpoint{5.340814in}{2.044718in}}{\pgfqpoint{5.340814in}{2.069004in}}%
\pgfpathquadraticcurveto{\pgfqpoint{5.340814in}{2.093351in}}{\pgfqpoint{5.349349in}{2.106134in}}%
\pgfpathquadraticcurveto{\pgfqpoint{5.357884in}{2.118936in}}{\pgfqpoint{5.374047in}{2.118936in}}%
\pgfpathlineto{\pgfqpoint{5.374047in}{2.118936in}}%
\pgfpathclose%
\pgfpathmoveto{\pgfqpoint{5.426640in}{2.121122in}}%
\pgfpathlineto{\pgfqpoint{5.436948in}{2.121122in}}%
\pgfpathquadraticcurveto{\pgfqpoint{5.446596in}{2.105927in}}{\pgfqpoint{5.451400in}{2.091352in}}%
\pgfpathquadraticcurveto{\pgfqpoint{5.456203in}{2.076797in}}{\pgfqpoint{5.456203in}{2.062427in}}%
\pgfpathquadraticcurveto{\pgfqpoint{5.456203in}{2.047996in}}{\pgfqpoint{5.451400in}{2.033379in}}%
\pgfpathquadraticcurveto{\pgfqpoint{5.446596in}{2.018762in}}{\pgfqpoint{5.436948in}{2.003609in}}%
\pgfpathlineto{\pgfqpoint{5.426640in}{2.003609in}}%
\pgfpathquadraticcurveto{\pgfqpoint{5.435195in}{2.018370in}}{\pgfqpoint{5.439422in}{2.032966in}}%
\pgfpathquadraticcurveto{\pgfqpoint{5.443648in}{2.047563in}}{\pgfqpoint{5.443648in}{2.062427in}}%
\pgfpathquadraticcurveto{\pgfqpoint{5.443648in}{2.077312in}}{\pgfqpoint{5.439422in}{2.091805in}}%
\pgfpathquadraticcurveto{\pgfqpoint{5.435195in}{2.106319in}}{\pgfqpoint{5.426640in}{2.121122in}}%
\pgfpathlineto{\pgfqpoint{5.426640in}{2.121122in}}%
\pgfpathclose%
\pgfusepath{fill}%
\end{pgfscope}%
\begin{pgfscope}%
\pgfsetbuttcap%
\pgfsetroundjoin%
\definecolor{currentfill}{rgb}{0.090196,0.168627,0.227451}%
\pgfsetfillcolor{currentfill}%
\pgfsetlinewidth{0.000000pt}%
\definecolor{currentstroke}{rgb}{0.090196,0.168627,0.227451}%
\pgfsetstrokecolor{currentstroke}%
\pgfsetdash{}{0pt}%
\pgfpathmoveto{\pgfqpoint{4.550611in}{1.937410in}}%
\pgfpathlineto{\pgfqpoint{4.550611in}{1.924711in}}%
\pgfpathquadraticcurveto{\pgfqpoint{4.543210in}{1.928257in}}{\pgfqpoint{4.536633in}{1.929988in}}%
\pgfpathquadraticcurveto{\pgfqpoint{4.530056in}{1.931741in}}{\pgfqpoint{4.523933in}{1.931741in}}%
\pgfpathquadraticcurveto{\pgfqpoint{4.513316in}{1.931741in}}{\pgfqpoint{4.507543in}{1.927617in}}%
\pgfpathquadraticcurveto{\pgfqpoint{4.501771in}{1.923494in}}{\pgfqpoint{4.501771in}{1.915887in}}%
\pgfpathquadraticcurveto{\pgfqpoint{4.501771in}{1.909496in}}{\pgfqpoint{4.505605in}{1.906238in}}%
\pgfpathquadraticcurveto{\pgfqpoint{4.509440in}{1.903002in}}{\pgfqpoint{4.520140in}{1.901002in}}%
\pgfpathlineto{\pgfqpoint{4.527995in}{1.899394in}}%
\pgfpathquadraticcurveto{\pgfqpoint{4.542550in}{1.896611in}}{\pgfqpoint{4.549477in}{1.889622in}}%
\pgfpathquadraticcurveto{\pgfqpoint{4.556404in}{1.882633in}}{\pgfqpoint{4.556404in}{1.870923in}}%
\pgfpathquadraticcurveto{\pgfqpoint{4.556404in}{1.856924in}}{\pgfqpoint{4.547024in}{1.849708in}}%
\pgfpathquadraticcurveto{\pgfqpoint{4.537664in}{1.842493in}}{\pgfqpoint{4.519563in}{1.842493in}}%
\pgfpathquadraticcurveto{\pgfqpoint{4.512739in}{1.842493in}}{\pgfqpoint{4.505028in}{1.844039in}}%
\pgfpathquadraticcurveto{\pgfqpoint{4.497338in}{1.845585in}}{\pgfqpoint{4.489092in}{1.848616in}}%
\pgfpathlineto{\pgfqpoint{4.489092in}{1.862016in}}%
\pgfpathquadraticcurveto{\pgfqpoint{4.497008in}{1.857584in}}{\pgfqpoint{4.504616in}{1.855316in}}%
\pgfpathquadraticcurveto{\pgfqpoint{4.512223in}{1.853069in}}{\pgfqpoint{4.519563in}{1.853069in}}%
\pgfpathquadraticcurveto{\pgfqpoint{4.530696in}{1.853069in}}{\pgfqpoint{4.536757in}{1.857440in}}%
\pgfpathquadraticcurveto{\pgfqpoint{4.542818in}{1.861831in}}{\pgfqpoint{4.542818in}{1.869954in}}%
\pgfpathquadraticcurveto{\pgfqpoint{4.542818in}{1.877025in}}{\pgfqpoint{4.538468in}{1.881025in}}%
\pgfpathquadraticcurveto{\pgfqpoint{4.534118in}{1.885024in}}{\pgfqpoint{4.524201in}{1.887024in}}%
\pgfpathlineto{\pgfqpoint{4.516264in}{1.888570in}}%
\pgfpathquadraticcurveto{\pgfqpoint{4.501709in}{1.891456in}}{\pgfqpoint{4.495194in}{1.897641in}}%
\pgfpathquadraticcurveto{\pgfqpoint{4.488700in}{1.903826in}}{\pgfqpoint{4.488700in}{1.914856in}}%
\pgfpathquadraticcurveto{\pgfqpoint{4.488700in}{1.927617in}}{\pgfqpoint{4.497689in}{1.934957in}}%
\pgfpathquadraticcurveto{\pgfqpoint{4.506678in}{1.942296in}}{\pgfqpoint{4.522449in}{1.942296in}}%
\pgfpathquadraticcurveto{\pgfqpoint{4.529232in}{1.942296in}}{\pgfqpoint{4.536241in}{1.941059in}}%
\pgfpathquadraticcurveto{\pgfqpoint{4.543271in}{1.939843in}}{\pgfqpoint{4.550611in}{1.937410in}}%
\pgfpathlineto{\pgfqpoint{4.550611in}{1.937410in}}%
\pgfpathclose%
\pgfpathmoveto{\pgfqpoint{4.587926in}{1.937019in}}%
\pgfpathlineto{\pgfqpoint{4.587926in}{1.916526in}}%
\pgfpathlineto{\pgfqpoint{4.612336in}{1.916526in}}%
\pgfpathlineto{\pgfqpoint{4.612336in}{1.907310in}}%
\pgfpathlineto{\pgfqpoint{4.587926in}{1.907310in}}%
\pgfpathlineto{\pgfqpoint{4.587926in}{1.868139in}}%
\pgfpathquadraticcurveto{\pgfqpoint{4.587926in}{1.859316in}}{\pgfqpoint{4.590339in}{1.856800in}}%
\pgfpathquadraticcurveto{\pgfqpoint{4.592751in}{1.854285in}}{\pgfqpoint{4.600173in}{1.854285in}}%
\pgfpathlineto{\pgfqpoint{4.612336in}{1.854285in}}%
\pgfpathlineto{\pgfqpoint{4.612336in}{1.844369in}}%
\pgfpathlineto{\pgfqpoint{4.600173in}{1.844369in}}%
\pgfpathquadraticcurveto{\pgfqpoint{4.586442in}{1.844369in}}{\pgfqpoint{4.581226in}{1.849482in}}%
\pgfpathquadraticcurveto{\pgfqpoint{4.576010in}{1.854615in}}{\pgfqpoint{4.576010in}{1.868139in}}%
\pgfpathlineto{\pgfqpoint{4.576010in}{1.907310in}}%
\pgfpathlineto{\pgfqpoint{4.567310in}{1.907310in}}%
\pgfpathlineto{\pgfqpoint{4.567310in}{1.916526in}}%
\pgfpathlineto{\pgfqpoint{4.576010in}{1.916526in}}%
\pgfpathlineto{\pgfqpoint{4.576010in}{1.937019in}}%
\pgfpathlineto{\pgfqpoint{4.587926in}{1.937019in}}%
\pgfpathlineto{\pgfqpoint{4.587926in}{1.937019in}}%
\pgfpathclose%
\pgfpathmoveto{\pgfqpoint{4.626706in}{1.872840in}}%
\pgfpathlineto{\pgfqpoint{4.626706in}{1.916526in}}%
\pgfpathlineto{\pgfqpoint{4.638560in}{1.916526in}}%
\pgfpathlineto{\pgfqpoint{4.638560in}{1.873293in}}%
\pgfpathquadraticcurveto{\pgfqpoint{4.638560in}{1.863047in}}{\pgfqpoint{4.642539in}{1.857914in}}%
\pgfpathquadraticcurveto{\pgfqpoint{4.646539in}{1.852801in}}{\pgfqpoint{4.654538in}{1.852801in}}%
\pgfpathquadraticcurveto{\pgfqpoint{4.664124in}{1.852801in}}{\pgfqpoint{4.669691in}{1.858924in}}%
\pgfpathquadraticcurveto{\pgfqpoint{4.675278in}{1.865047in}}{\pgfqpoint{4.675278in}{1.875623in}}%
\pgfpathlineto{\pgfqpoint{4.675278in}{1.916526in}}%
\pgfpathlineto{\pgfqpoint{4.687132in}{1.916526in}}%
\pgfpathlineto{\pgfqpoint{4.687132in}{1.844369in}}%
\pgfpathlineto{\pgfqpoint{4.675278in}{1.844369in}}%
\pgfpathlineto{\pgfqpoint{4.675278in}{1.855460in}}%
\pgfpathquadraticcurveto{\pgfqpoint{4.670969in}{1.848884in}}{\pgfqpoint{4.665258in}{1.845688in}}%
\pgfpathquadraticcurveto{\pgfqpoint{4.659568in}{1.842493in}}{\pgfqpoint{4.652023in}{1.842493in}}%
\pgfpathquadraticcurveto{\pgfqpoint{4.639591in}{1.842493in}}{\pgfqpoint{4.633138in}{1.850224in}}%
\pgfpathquadraticcurveto{\pgfqpoint{4.626706in}{1.857955in}}{\pgfqpoint{4.626706in}{1.872840in}}%
\pgfpathlineto{\pgfqpoint{4.626706in}{1.872840in}}%
\pgfpathclose%
\pgfpathmoveto{\pgfqpoint{4.656538in}{1.918258in}}%
\pgfpathlineto{\pgfqpoint{4.656538in}{1.918258in}}%
\pgfpathlineto{\pgfqpoint{4.656538in}{1.918258in}}%
\pgfpathclose%
\pgfpathmoveto{\pgfqpoint{4.759021in}{1.905579in}}%
\pgfpathlineto{\pgfqpoint{4.759021in}{1.944626in}}%
\pgfpathlineto{\pgfqpoint{4.770876in}{1.944626in}}%
\pgfpathlineto{\pgfqpoint{4.770876in}{1.844369in}}%
\pgfpathlineto{\pgfqpoint{4.759021in}{1.844369in}}%
\pgfpathlineto{\pgfqpoint{4.759021in}{1.855192in}}%
\pgfpathquadraticcurveto{\pgfqpoint{4.755290in}{1.848760in}}{\pgfqpoint{4.749579in}{1.845626in}}%
\pgfpathquadraticcurveto{\pgfqpoint{4.743889in}{1.842493in}}{\pgfqpoint{4.735890in}{1.842493in}}%
\pgfpathquadraticcurveto{\pgfqpoint{4.722819in}{1.842493in}}{\pgfqpoint{4.714593in}{1.852925in}}%
\pgfpathquadraticcurveto{\pgfqpoint{4.706388in}{1.863377in}}{\pgfqpoint{4.706388in}{1.880386in}}%
\pgfpathquadraticcurveto{\pgfqpoint{4.706388in}{1.897394in}}{\pgfqpoint{4.714593in}{1.907826in}}%
\pgfpathquadraticcurveto{\pgfqpoint{4.722819in}{1.918258in}}{\pgfqpoint{4.735890in}{1.918258in}}%
\pgfpathquadraticcurveto{\pgfqpoint{4.743889in}{1.918258in}}{\pgfqpoint{4.749579in}{1.915124in}}%
\pgfpathquadraticcurveto{\pgfqpoint{4.755290in}{1.912011in}}{\pgfqpoint{4.759021in}{1.905579in}}%
\pgfpathlineto{\pgfqpoint{4.759021in}{1.905579in}}%
\pgfpathclose%
\pgfpathmoveto{\pgfqpoint{4.718634in}{1.880386in}}%
\pgfpathquadraticcurveto{\pgfqpoint{4.718634in}{1.867315in}}{\pgfqpoint{4.724015in}{1.859872in}}%
\pgfpathquadraticcurveto{\pgfqpoint{4.729396in}{1.852430in}}{\pgfqpoint{4.738797in}{1.852430in}}%
\pgfpathquadraticcurveto{\pgfqpoint{4.748198in}{1.852430in}}{\pgfqpoint{4.753599in}{1.859872in}}%
\pgfpathquadraticcurveto{\pgfqpoint{4.759021in}{1.867315in}}{\pgfqpoint{4.759021in}{1.880386in}}%
\pgfpathquadraticcurveto{\pgfqpoint{4.759021in}{1.893456in}}{\pgfqpoint{4.753599in}{1.900899in}}%
\pgfpathquadraticcurveto{\pgfqpoint{4.748198in}{1.908341in}}{\pgfqpoint{4.738797in}{1.908341in}}%
\pgfpathquadraticcurveto{\pgfqpoint{4.729396in}{1.908341in}}{\pgfqpoint{4.724015in}{1.900899in}}%
\pgfpathquadraticcurveto{\pgfqpoint{4.718634in}{1.893456in}}{\pgfqpoint{4.718634in}{1.880386in}}%
\pgfpathlineto{\pgfqpoint{4.718634in}{1.880386in}}%
\pgfpathclose%
\pgfpathmoveto{\pgfqpoint{4.825323in}{1.837668in}}%
\pgfpathquadraticcurveto{\pgfqpoint{4.820314in}{1.824783in}}{\pgfqpoint{4.815531in}{1.820866in}}%
\pgfpathquadraticcurveto{\pgfqpoint{4.810768in}{1.816928in}}{\pgfqpoint{4.802790in}{1.816928in}}%
\pgfpathlineto{\pgfqpoint{4.793306in}{1.816928in}}%
\pgfpathlineto{\pgfqpoint{4.793306in}{1.826845in}}%
\pgfpathlineto{\pgfqpoint{4.800275in}{1.826845in}}%
\pgfpathquadraticcurveto{\pgfqpoint{4.805161in}{1.826845in}}{\pgfqpoint{4.807861in}{1.829175in}}%
\pgfpathquadraticcurveto{\pgfqpoint{4.810583in}{1.831484in}}{\pgfqpoint{4.813861in}{1.840122in}}%
\pgfpathlineto{\pgfqpoint{4.815984in}{1.845523in}}%
\pgfpathlineto{\pgfqpoint{4.786812in}{1.916526in}}%
\pgfpathlineto{\pgfqpoint{4.799367in}{1.916526in}}%
\pgfpathlineto{\pgfqpoint{4.821922in}{1.860099in}}%
\pgfpathlineto{\pgfqpoint{4.844476in}{1.916526in}}%
\pgfpathlineto{\pgfqpoint{4.857031in}{1.916526in}}%
\pgfpathlineto{\pgfqpoint{4.825323in}{1.837668in}}%
\pgfpathlineto{\pgfqpoint{4.825323in}{1.837668in}}%
\pgfpathclose%
\pgfpathmoveto{\pgfqpoint{4.919272in}{1.855316in}}%
\pgfpathlineto{\pgfqpoint{4.940527in}{1.855316in}}%
\pgfpathlineto{\pgfqpoint{4.940527in}{1.928710in}}%
\pgfpathlineto{\pgfqpoint{4.917396in}{1.924071in}}%
\pgfpathlineto{\pgfqpoint{4.917396in}{1.935926in}}%
\pgfpathlineto{\pgfqpoint{4.940404in}{1.940565in}}%
\pgfpathlineto{\pgfqpoint{4.953413in}{1.940565in}}%
\pgfpathlineto{\pgfqpoint{4.953413in}{1.855316in}}%
\pgfpathlineto{\pgfqpoint{4.974668in}{1.855316in}}%
\pgfpathlineto{\pgfqpoint{4.974668in}{1.844369in}}%
\pgfpathlineto{\pgfqpoint{4.919272in}{1.844369in}}%
\pgfpathlineto{\pgfqpoint{4.919272in}{1.855316in}}%
\pgfpathlineto{\pgfqpoint{4.919272in}{1.855316in}}%
\pgfpathclose%
\pgfpathmoveto{\pgfqpoint{5.028786in}{1.931988in}}%
\pgfpathquadraticcurveto{\pgfqpoint{5.018746in}{1.931988in}}{\pgfqpoint{5.013674in}{1.922092in}}%
\pgfpathquadraticcurveto{\pgfqpoint{5.008623in}{1.912217in}}{\pgfqpoint{5.008623in}{1.892364in}}%
\pgfpathquadraticcurveto{\pgfqpoint{5.008623in}{1.872593in}}{\pgfqpoint{5.013674in}{1.862697in}}%
\pgfpathquadraticcurveto{\pgfqpoint{5.018746in}{1.852801in}}{\pgfqpoint{5.028786in}{1.852801in}}%
\pgfpathquadraticcurveto{\pgfqpoint{5.038908in}{1.852801in}}{\pgfqpoint{5.043959in}{1.862697in}}%
\pgfpathquadraticcurveto{\pgfqpoint{5.049031in}{1.872593in}}{\pgfqpoint{5.049031in}{1.892364in}}%
\pgfpathquadraticcurveto{\pgfqpoint{5.049031in}{1.912217in}}{\pgfqpoint{5.043959in}{1.922092in}}%
\pgfpathquadraticcurveto{\pgfqpoint{5.038908in}{1.931988in}}{\pgfqpoint{5.028786in}{1.931988in}}%
\pgfpathlineto{\pgfqpoint{5.028786in}{1.931988in}}%
\pgfpathclose%
\pgfpathmoveto{\pgfqpoint{5.028786in}{1.942296in}}%
\pgfpathquadraticcurveto{\pgfqpoint{5.044970in}{1.942296in}}{\pgfqpoint{5.053505in}{1.929494in}}%
\pgfpathquadraticcurveto{\pgfqpoint{5.062040in}{1.916711in}}{\pgfqpoint{5.062040in}{1.892364in}}%
\pgfpathquadraticcurveto{\pgfqpoint{5.062040in}{1.868078in}}{\pgfqpoint{5.053505in}{1.855275in}}%
\pgfpathquadraticcurveto{\pgfqpoint{5.044970in}{1.842493in}}{\pgfqpoint{5.028786in}{1.842493in}}%
\pgfpathquadraticcurveto{\pgfqpoint{5.012623in}{1.842493in}}{\pgfqpoint{5.004087in}{1.855275in}}%
\pgfpathquadraticcurveto{\pgfqpoint{4.995552in}{1.868078in}}{\pgfqpoint{4.995552in}{1.892364in}}%
\pgfpathquadraticcurveto{\pgfqpoint{4.995552in}{1.916711in}}{\pgfqpoint{5.004087in}{1.929494in}}%
\pgfpathquadraticcurveto{\pgfqpoint{5.012623in}{1.942296in}}{\pgfqpoint{5.028786in}{1.942296in}}%
\pgfpathlineto{\pgfqpoint{5.028786in}{1.942296in}}%
\pgfpathclose%
\pgfpathmoveto{\pgfqpoint{5.153638in}{1.944482in}}%
\pgfpathquadraticcurveto{\pgfqpoint{5.145021in}{1.929679in}}{\pgfqpoint{5.140815in}{1.915165in}}%
\pgfpathquadraticcurveto{\pgfqpoint{5.136630in}{1.900672in}}{\pgfqpoint{5.136630in}{1.885787in}}%
\pgfpathquadraticcurveto{\pgfqpoint{5.136630in}{1.870923in}}{\pgfqpoint{5.140856in}{1.856326in}}%
\pgfpathquadraticcurveto{\pgfqpoint{5.145082in}{1.841730in}}{\pgfqpoint{5.153638in}{1.826969in}}%
\pgfpathlineto{\pgfqpoint{5.143330in}{1.826969in}}%
\pgfpathquadraticcurveto{\pgfqpoint{5.133682in}{1.842122in}}{\pgfqpoint{5.128878in}{1.856739in}}%
\pgfpathquadraticcurveto{\pgfqpoint{5.124074in}{1.871356in}}{\pgfqpoint{5.124074in}{1.885787in}}%
\pgfpathquadraticcurveto{\pgfqpoint{5.124074in}{1.900157in}}{\pgfqpoint{5.128837in}{1.914712in}}%
\pgfpathquadraticcurveto{\pgfqpoint{5.133620in}{1.929287in}}{\pgfqpoint{5.143330in}{1.944482in}}%
\pgfpathlineto{\pgfqpoint{5.153638in}{1.944482in}}%
\pgfpathlineto{\pgfqpoint{5.153638in}{1.944482in}}%
\pgfpathclose%
\pgfpathmoveto{\pgfqpoint{5.189531in}{1.855316in}}%
\pgfpathlineto{\pgfqpoint{5.234949in}{1.855316in}}%
\pgfpathlineto{\pgfqpoint{5.234949in}{1.844369in}}%
\pgfpathlineto{\pgfqpoint{5.173883in}{1.844369in}}%
\pgfpathlineto{\pgfqpoint{5.173883in}{1.855316in}}%
\pgfpathquadraticcurveto{\pgfqpoint{5.181285in}{1.862985in}}{\pgfqpoint{5.194067in}{1.875891in}}%
\pgfpathquadraticcurveto{\pgfqpoint{5.206870in}{1.888818in}}{\pgfqpoint{5.210148in}{1.892570in}}%
\pgfpathquadraticcurveto{\pgfqpoint{5.216394in}{1.899579in}}{\pgfqpoint{5.218868in}{1.904445in}}%
\pgfpathquadraticcurveto{\pgfqpoint{5.221363in}{1.909310in}}{\pgfqpoint{5.221363in}{1.914011in}}%
\pgfpathquadraticcurveto{\pgfqpoint{5.221363in}{1.921680in}}{\pgfqpoint{5.215982in}{1.926504in}}%
\pgfpathquadraticcurveto{\pgfqpoint{5.210601in}{1.931349in}}{\pgfqpoint{5.201963in}{1.931349in}}%
\pgfpathquadraticcurveto{\pgfqpoint{5.195840in}{1.931349in}}{\pgfqpoint{5.189036in}{1.929226in}}%
\pgfpathquadraticcurveto{\pgfqpoint{5.182254in}{1.927102in}}{\pgfqpoint{5.174523in}{1.922773in}}%
\pgfpathlineto{\pgfqpoint{5.174523in}{1.935926in}}%
\pgfpathquadraticcurveto{\pgfqpoint{5.182377in}{1.939080in}}{\pgfqpoint{5.189201in}{1.940688in}}%
\pgfpathquadraticcurveto{\pgfqpoint{5.196046in}{1.942296in}}{\pgfqpoint{5.201715in}{1.942296in}}%
\pgfpathquadraticcurveto{\pgfqpoint{5.216662in}{1.942296in}}{\pgfqpoint{5.225548in}{1.934813in}}%
\pgfpathquadraticcurveto{\pgfqpoint{5.234434in}{1.927349in}}{\pgfqpoint{5.234434in}{1.914856in}}%
\pgfpathquadraticcurveto{\pgfqpoint{5.234434in}{1.908918in}}{\pgfqpoint{5.232207in}{1.903599in}}%
\pgfpathquadraticcurveto{\pgfqpoint{5.230001in}{1.898301in}}{\pgfqpoint{5.224125in}{1.891085in}}%
\pgfpathquadraticcurveto{\pgfqpoint{5.222517in}{1.889209in}}{\pgfqpoint{5.213879in}{1.880282in}}%
\pgfpathquadraticcurveto{\pgfqpoint{5.205262in}{1.871356in}}{\pgfqpoint{5.189531in}{1.855316in}}%
\pgfpathlineto{\pgfqpoint{5.189531in}{1.855316in}}%
\pgfpathclose%
\pgfpathmoveto{\pgfqpoint{5.290098in}{1.931988in}}%
\pgfpathquadraticcurveto{\pgfqpoint{5.280058in}{1.931988in}}{\pgfqpoint{5.274986in}{1.922092in}}%
\pgfpathquadraticcurveto{\pgfqpoint{5.269935in}{1.912217in}}{\pgfqpoint{5.269935in}{1.892364in}}%
\pgfpathquadraticcurveto{\pgfqpoint{5.269935in}{1.872593in}}{\pgfqpoint{5.274986in}{1.862697in}}%
\pgfpathquadraticcurveto{\pgfqpoint{5.280058in}{1.852801in}}{\pgfqpoint{5.290098in}{1.852801in}}%
\pgfpathquadraticcurveto{\pgfqpoint{5.300220in}{1.852801in}}{\pgfqpoint{5.305271in}{1.862697in}}%
\pgfpathquadraticcurveto{\pgfqpoint{5.310343in}{1.872593in}}{\pgfqpoint{5.310343in}{1.892364in}}%
\pgfpathquadraticcurveto{\pgfqpoint{5.310343in}{1.912217in}}{\pgfqpoint{5.305271in}{1.922092in}}%
\pgfpathquadraticcurveto{\pgfqpoint{5.300220in}{1.931988in}}{\pgfqpoint{5.290098in}{1.931988in}}%
\pgfpathlineto{\pgfqpoint{5.290098in}{1.931988in}}%
\pgfpathclose%
\pgfpathmoveto{\pgfqpoint{5.290098in}{1.942296in}}%
\pgfpathquadraticcurveto{\pgfqpoint{5.306281in}{1.942296in}}{\pgfqpoint{5.314817in}{1.929494in}}%
\pgfpathquadraticcurveto{\pgfqpoint{5.323352in}{1.916711in}}{\pgfqpoint{5.323352in}{1.892364in}}%
\pgfpathquadraticcurveto{\pgfqpoint{5.323352in}{1.868078in}}{\pgfqpoint{5.314817in}{1.855275in}}%
\pgfpathquadraticcurveto{\pgfqpoint{5.306281in}{1.842493in}}{\pgfqpoint{5.290098in}{1.842493in}}%
\pgfpathquadraticcurveto{\pgfqpoint{5.273934in}{1.842493in}}{\pgfqpoint{5.265399in}{1.855275in}}%
\pgfpathquadraticcurveto{\pgfqpoint{5.256864in}{1.868078in}}{\pgfqpoint{5.256864in}{1.892364in}}%
\pgfpathquadraticcurveto{\pgfqpoint{5.256864in}{1.916711in}}{\pgfqpoint{5.265399in}{1.929494in}}%
\pgfpathquadraticcurveto{\pgfqpoint{5.273934in}{1.942296in}}{\pgfqpoint{5.290098in}{1.942296in}}%
\pgfpathlineto{\pgfqpoint{5.290098in}{1.942296in}}%
\pgfpathclose%
\pgfpathmoveto{\pgfqpoint{5.357431in}{1.855316in}}%
\pgfpathlineto{\pgfqpoint{5.402848in}{1.855316in}}%
\pgfpathlineto{\pgfqpoint{5.402848in}{1.844369in}}%
\pgfpathlineto{\pgfqpoint{5.341783in}{1.844369in}}%
\pgfpathlineto{\pgfqpoint{5.341783in}{1.855316in}}%
\pgfpathquadraticcurveto{\pgfqpoint{5.349184in}{1.862985in}}{\pgfqpoint{5.361966in}{1.875891in}}%
\pgfpathquadraticcurveto{\pgfqpoint{5.374769in}{1.888818in}}{\pgfqpoint{5.378047in}{1.892570in}}%
\pgfpathquadraticcurveto{\pgfqpoint{5.384294in}{1.899579in}}{\pgfqpoint{5.386768in}{1.904445in}}%
\pgfpathquadraticcurveto{\pgfqpoint{5.389262in}{1.909310in}}{\pgfqpoint{5.389262in}{1.914011in}}%
\pgfpathquadraticcurveto{\pgfqpoint{5.389262in}{1.921680in}}{\pgfqpoint{5.383881in}{1.926504in}}%
\pgfpathquadraticcurveto{\pgfqpoint{5.378500in}{1.931349in}}{\pgfqpoint{5.369862in}{1.931349in}}%
\pgfpathquadraticcurveto{\pgfqpoint{5.363739in}{1.931349in}}{\pgfqpoint{5.356936in}{1.929226in}}%
\pgfpathquadraticcurveto{\pgfqpoint{5.350153in}{1.927102in}}{\pgfqpoint{5.342422in}{1.922773in}}%
\pgfpathlineto{\pgfqpoint{5.342422in}{1.935926in}}%
\pgfpathquadraticcurveto{\pgfqpoint{5.350277in}{1.939080in}}{\pgfqpoint{5.357101in}{1.940688in}}%
\pgfpathquadraticcurveto{\pgfqpoint{5.363945in}{1.942296in}}{\pgfqpoint{5.369615in}{1.942296in}}%
\pgfpathquadraticcurveto{\pgfqpoint{5.384562in}{1.942296in}}{\pgfqpoint{5.393447in}{1.934813in}}%
\pgfpathquadraticcurveto{\pgfqpoint{5.402333in}{1.927349in}}{\pgfqpoint{5.402333in}{1.914856in}}%
\pgfpathquadraticcurveto{\pgfqpoint{5.402333in}{1.908918in}}{\pgfqpoint{5.400106in}{1.903599in}}%
\pgfpathquadraticcurveto{\pgfqpoint{5.397900in}{1.898301in}}{\pgfqpoint{5.392025in}{1.891085in}}%
\pgfpathquadraticcurveto{\pgfqpoint{5.390417in}{1.889209in}}{\pgfqpoint{5.381778in}{1.880282in}}%
\pgfpathquadraticcurveto{\pgfqpoint{5.373161in}{1.871356in}}{\pgfqpoint{5.357431in}{1.855316in}}%
\pgfpathlineto{\pgfqpoint{5.357431in}{1.855316in}}%
\pgfpathclose%
\pgfpathmoveto{\pgfqpoint{5.432433in}{1.855316in}}%
\pgfpathlineto{\pgfqpoint{5.453688in}{1.855316in}}%
\pgfpathlineto{\pgfqpoint{5.453688in}{1.928710in}}%
\pgfpathlineto{\pgfqpoint{5.430557in}{1.924071in}}%
\pgfpathlineto{\pgfqpoint{5.430557in}{1.935926in}}%
\pgfpathlineto{\pgfqpoint{5.453564in}{1.940565in}}%
\pgfpathlineto{\pgfqpoint{5.466573in}{1.940565in}}%
\pgfpathlineto{\pgfqpoint{5.466573in}{1.855316in}}%
\pgfpathlineto{\pgfqpoint{5.487829in}{1.855316in}}%
\pgfpathlineto{\pgfqpoint{5.487829in}{1.844369in}}%
\pgfpathlineto{\pgfqpoint{5.432433in}{1.844369in}}%
\pgfpathlineto{\pgfqpoint{5.432433in}{1.855316in}}%
\pgfpathlineto{\pgfqpoint{5.432433in}{1.855316in}}%
\pgfpathclose%
\pgfpathmoveto{\pgfqpoint{5.510589in}{1.944482in}}%
\pgfpathlineto{\pgfqpoint{5.520897in}{1.944482in}}%
\pgfpathquadraticcurveto{\pgfqpoint{5.530546in}{1.929287in}}{\pgfqpoint{5.535349in}{1.914712in}}%
\pgfpathquadraticcurveto{\pgfqpoint{5.540153in}{1.900157in}}{\pgfqpoint{5.540153in}{1.885787in}}%
\pgfpathquadraticcurveto{\pgfqpoint{5.540153in}{1.871356in}}{\pgfqpoint{5.535349in}{1.856739in}}%
\pgfpathquadraticcurveto{\pgfqpoint{5.530546in}{1.842122in}}{\pgfqpoint{5.520897in}{1.826969in}}%
\pgfpathlineto{\pgfqpoint{5.510589in}{1.826969in}}%
\pgfpathquadraticcurveto{\pgfqpoint{5.519145in}{1.841730in}}{\pgfqpoint{5.523371in}{1.856326in}}%
\pgfpathquadraticcurveto{\pgfqpoint{5.527598in}{1.870923in}}{\pgfqpoint{5.527598in}{1.885787in}}%
\pgfpathquadraticcurveto{\pgfqpoint{5.527598in}{1.900672in}}{\pgfqpoint{5.523371in}{1.915165in}}%
\pgfpathquadraticcurveto{\pgfqpoint{5.519145in}{1.929679in}}{\pgfqpoint{5.510589in}{1.944482in}}%
\pgfpathlineto{\pgfqpoint{5.510589in}{1.944482in}}%
\pgfpathclose%
\pgfusepath{fill}%
\end{pgfscope}%
\begin{pgfscope}%
\pgfsetbuttcap%
\pgfsetroundjoin%
\definecolor{currentfill}{rgb}{0.090196,0.168627,0.227451}%
\pgfsetfillcolor{currentfill}%
\pgfsetlinewidth{0.000000pt}%
\definecolor{currentstroke}{rgb}{0.090196,0.168627,0.227451}%
\pgfsetstrokecolor{currentstroke}%
\pgfsetdash{}{0pt}%
\pgfpathmoveto{\pgfqpoint{4.550611in}{1.760770in}}%
\pgfpathlineto{\pgfqpoint{4.550611in}{1.748071in}}%
\pgfpathquadraticcurveto{\pgfqpoint{4.543210in}{1.751617in}}{\pgfqpoint{4.536633in}{1.753348in}}%
\pgfpathquadraticcurveto{\pgfqpoint{4.530056in}{1.755101in}}{\pgfqpoint{4.523933in}{1.755101in}}%
\pgfpathquadraticcurveto{\pgfqpoint{4.513316in}{1.755101in}}{\pgfqpoint{4.507543in}{1.750977in}}%
\pgfpathquadraticcurveto{\pgfqpoint{4.501771in}{1.746854in}}{\pgfqpoint{4.501771in}{1.739247in}}%
\pgfpathquadraticcurveto{\pgfqpoint{4.501771in}{1.732856in}}{\pgfqpoint{4.505605in}{1.729598in}}%
\pgfpathquadraticcurveto{\pgfqpoint{4.509440in}{1.726362in}}{\pgfqpoint{4.520140in}{1.724362in}}%
\pgfpathlineto{\pgfqpoint{4.527995in}{1.722754in}}%
\pgfpathquadraticcurveto{\pgfqpoint{4.542550in}{1.719971in}}{\pgfqpoint{4.549477in}{1.712982in}}%
\pgfpathquadraticcurveto{\pgfqpoint{4.556404in}{1.705993in}}{\pgfqpoint{4.556404in}{1.694283in}}%
\pgfpathquadraticcurveto{\pgfqpoint{4.556404in}{1.680284in}}{\pgfqpoint{4.547024in}{1.673068in}}%
\pgfpathquadraticcurveto{\pgfqpoint{4.537664in}{1.665853in}}{\pgfqpoint{4.519563in}{1.665853in}}%
\pgfpathquadraticcurveto{\pgfqpoint{4.512739in}{1.665853in}}{\pgfqpoint{4.505028in}{1.667399in}}%
\pgfpathquadraticcurveto{\pgfqpoint{4.497338in}{1.668945in}}{\pgfqpoint{4.489092in}{1.671976in}}%
\pgfpathlineto{\pgfqpoint{4.489092in}{1.685376in}}%
\pgfpathquadraticcurveto{\pgfqpoint{4.497008in}{1.680944in}}{\pgfqpoint{4.504616in}{1.678676in}}%
\pgfpathquadraticcurveto{\pgfqpoint{4.512223in}{1.676429in}}{\pgfqpoint{4.519563in}{1.676429in}}%
\pgfpathquadraticcurveto{\pgfqpoint{4.530696in}{1.676429in}}{\pgfqpoint{4.536757in}{1.680800in}}%
\pgfpathquadraticcurveto{\pgfqpoint{4.542818in}{1.685191in}}{\pgfqpoint{4.542818in}{1.693314in}}%
\pgfpathquadraticcurveto{\pgfqpoint{4.542818in}{1.700385in}}{\pgfqpoint{4.538468in}{1.704385in}}%
\pgfpathquadraticcurveto{\pgfqpoint{4.534118in}{1.708384in}}{\pgfqpoint{4.524201in}{1.710384in}}%
\pgfpathlineto{\pgfqpoint{4.516264in}{1.711930in}}%
\pgfpathquadraticcurveto{\pgfqpoint{4.501709in}{1.714816in}}{\pgfqpoint{4.495194in}{1.721001in}}%
\pgfpathquadraticcurveto{\pgfqpoint{4.488700in}{1.727186in}}{\pgfqpoint{4.488700in}{1.738216in}}%
\pgfpathquadraticcurveto{\pgfqpoint{4.488700in}{1.750977in}}{\pgfqpoint{4.497689in}{1.758317in}}%
\pgfpathquadraticcurveto{\pgfqpoint{4.506678in}{1.765656in}}{\pgfqpoint{4.522449in}{1.765656in}}%
\pgfpathquadraticcurveto{\pgfqpoint{4.529232in}{1.765656in}}{\pgfqpoint{4.536241in}{1.764419in}}%
\pgfpathquadraticcurveto{\pgfqpoint{4.543271in}{1.763203in}}{\pgfqpoint{4.550611in}{1.760770in}}%
\pgfpathlineto{\pgfqpoint{4.550611in}{1.760770in}}%
\pgfpathclose%
\pgfpathmoveto{\pgfqpoint{4.587926in}{1.760379in}}%
\pgfpathlineto{\pgfqpoint{4.587926in}{1.739886in}}%
\pgfpathlineto{\pgfqpoint{4.612336in}{1.739886in}}%
\pgfpathlineto{\pgfqpoint{4.612336in}{1.730670in}}%
\pgfpathlineto{\pgfqpoint{4.587926in}{1.730670in}}%
\pgfpathlineto{\pgfqpoint{4.587926in}{1.691499in}}%
\pgfpathquadraticcurveto{\pgfqpoint{4.587926in}{1.682676in}}{\pgfqpoint{4.590339in}{1.680160in}}%
\pgfpathquadraticcurveto{\pgfqpoint{4.592751in}{1.677645in}}{\pgfqpoint{4.600173in}{1.677645in}}%
\pgfpathlineto{\pgfqpoint{4.612336in}{1.677645in}}%
\pgfpathlineto{\pgfqpoint{4.612336in}{1.667729in}}%
\pgfpathlineto{\pgfqpoint{4.600173in}{1.667729in}}%
\pgfpathquadraticcurveto{\pgfqpoint{4.586442in}{1.667729in}}{\pgfqpoint{4.581226in}{1.672842in}}%
\pgfpathquadraticcurveto{\pgfqpoint{4.576010in}{1.677975in}}{\pgfqpoint{4.576010in}{1.691499in}}%
\pgfpathlineto{\pgfqpoint{4.576010in}{1.730670in}}%
\pgfpathlineto{\pgfqpoint{4.567310in}{1.730670in}}%
\pgfpathlineto{\pgfqpoint{4.567310in}{1.739886in}}%
\pgfpathlineto{\pgfqpoint{4.576010in}{1.739886in}}%
\pgfpathlineto{\pgfqpoint{4.576010in}{1.760379in}}%
\pgfpathlineto{\pgfqpoint{4.587926in}{1.760379in}}%
\pgfpathlineto{\pgfqpoint{4.587926in}{1.760379in}}%
\pgfpathclose%
\pgfpathmoveto{\pgfqpoint{4.626706in}{1.696200in}}%
\pgfpathlineto{\pgfqpoint{4.626706in}{1.739886in}}%
\pgfpathlineto{\pgfqpoint{4.638560in}{1.739886in}}%
\pgfpathlineto{\pgfqpoint{4.638560in}{1.696653in}}%
\pgfpathquadraticcurveto{\pgfqpoint{4.638560in}{1.686407in}}{\pgfqpoint{4.642539in}{1.681274in}}%
\pgfpathquadraticcurveto{\pgfqpoint{4.646539in}{1.676161in}}{\pgfqpoint{4.654538in}{1.676161in}}%
\pgfpathquadraticcurveto{\pgfqpoint{4.664124in}{1.676161in}}{\pgfqpoint{4.669691in}{1.682284in}}%
\pgfpathquadraticcurveto{\pgfqpoint{4.675278in}{1.688407in}}{\pgfqpoint{4.675278in}{1.698983in}}%
\pgfpathlineto{\pgfqpoint{4.675278in}{1.739886in}}%
\pgfpathlineto{\pgfqpoint{4.687132in}{1.739886in}}%
\pgfpathlineto{\pgfqpoint{4.687132in}{1.667729in}}%
\pgfpathlineto{\pgfqpoint{4.675278in}{1.667729in}}%
\pgfpathlineto{\pgfqpoint{4.675278in}{1.678820in}}%
\pgfpathquadraticcurveto{\pgfqpoint{4.670969in}{1.672244in}}{\pgfqpoint{4.665258in}{1.669048in}}%
\pgfpathquadraticcurveto{\pgfqpoint{4.659568in}{1.665853in}}{\pgfqpoint{4.652023in}{1.665853in}}%
\pgfpathquadraticcurveto{\pgfqpoint{4.639591in}{1.665853in}}{\pgfqpoint{4.633138in}{1.673584in}}%
\pgfpathquadraticcurveto{\pgfqpoint{4.626706in}{1.681315in}}{\pgfqpoint{4.626706in}{1.696200in}}%
\pgfpathlineto{\pgfqpoint{4.626706in}{1.696200in}}%
\pgfpathclose%
\pgfpathmoveto{\pgfqpoint{4.656538in}{1.741618in}}%
\pgfpathlineto{\pgfqpoint{4.656538in}{1.741618in}}%
\pgfpathlineto{\pgfqpoint{4.656538in}{1.741618in}}%
\pgfpathclose%
\pgfpathmoveto{\pgfqpoint{4.759021in}{1.728939in}}%
\pgfpathlineto{\pgfqpoint{4.759021in}{1.767986in}}%
\pgfpathlineto{\pgfqpoint{4.770876in}{1.767986in}}%
\pgfpathlineto{\pgfqpoint{4.770876in}{1.667729in}}%
\pgfpathlineto{\pgfqpoint{4.759021in}{1.667729in}}%
\pgfpathlineto{\pgfqpoint{4.759021in}{1.678552in}}%
\pgfpathquadraticcurveto{\pgfqpoint{4.755290in}{1.672120in}}{\pgfqpoint{4.749579in}{1.668986in}}%
\pgfpathquadraticcurveto{\pgfqpoint{4.743889in}{1.665853in}}{\pgfqpoint{4.735890in}{1.665853in}}%
\pgfpathquadraticcurveto{\pgfqpoint{4.722819in}{1.665853in}}{\pgfqpoint{4.714593in}{1.676285in}}%
\pgfpathquadraticcurveto{\pgfqpoint{4.706388in}{1.686737in}}{\pgfqpoint{4.706388in}{1.703746in}}%
\pgfpathquadraticcurveto{\pgfqpoint{4.706388in}{1.720754in}}{\pgfqpoint{4.714593in}{1.731186in}}%
\pgfpathquadraticcurveto{\pgfqpoint{4.722819in}{1.741618in}}{\pgfqpoint{4.735890in}{1.741618in}}%
\pgfpathquadraticcurveto{\pgfqpoint{4.743889in}{1.741618in}}{\pgfqpoint{4.749579in}{1.738484in}}%
\pgfpathquadraticcurveto{\pgfqpoint{4.755290in}{1.735371in}}{\pgfqpoint{4.759021in}{1.728939in}}%
\pgfpathlineto{\pgfqpoint{4.759021in}{1.728939in}}%
\pgfpathclose%
\pgfpathmoveto{\pgfqpoint{4.718634in}{1.703746in}}%
\pgfpathquadraticcurveto{\pgfqpoint{4.718634in}{1.690675in}}{\pgfqpoint{4.724015in}{1.683232in}}%
\pgfpathquadraticcurveto{\pgfqpoint{4.729396in}{1.675790in}}{\pgfqpoint{4.738797in}{1.675790in}}%
\pgfpathquadraticcurveto{\pgfqpoint{4.748198in}{1.675790in}}{\pgfqpoint{4.753599in}{1.683232in}}%
\pgfpathquadraticcurveto{\pgfqpoint{4.759021in}{1.690675in}}{\pgfqpoint{4.759021in}{1.703746in}}%
\pgfpathquadraticcurveto{\pgfqpoint{4.759021in}{1.716816in}}{\pgfqpoint{4.753599in}{1.724259in}}%
\pgfpathquadraticcurveto{\pgfqpoint{4.748198in}{1.731701in}}{\pgfqpoint{4.738797in}{1.731701in}}%
\pgfpathquadraticcurveto{\pgfqpoint{4.729396in}{1.731701in}}{\pgfqpoint{4.724015in}{1.724259in}}%
\pgfpathquadraticcurveto{\pgfqpoint{4.718634in}{1.716816in}}{\pgfqpoint{4.718634in}{1.703746in}}%
\pgfpathlineto{\pgfqpoint{4.718634in}{1.703746in}}%
\pgfpathclose%
\pgfpathmoveto{\pgfqpoint{4.825323in}{1.661028in}}%
\pgfpathquadraticcurveto{\pgfqpoint{4.820314in}{1.648143in}}{\pgfqpoint{4.815531in}{1.644226in}}%
\pgfpathquadraticcurveto{\pgfqpoint{4.810768in}{1.640288in}}{\pgfqpoint{4.802790in}{1.640288in}}%
\pgfpathlineto{\pgfqpoint{4.793306in}{1.640288in}}%
\pgfpathlineto{\pgfqpoint{4.793306in}{1.650205in}}%
\pgfpathlineto{\pgfqpoint{4.800275in}{1.650205in}}%
\pgfpathquadraticcurveto{\pgfqpoint{4.805161in}{1.650205in}}{\pgfqpoint{4.807861in}{1.652535in}}%
\pgfpathquadraticcurveto{\pgfqpoint{4.810583in}{1.654844in}}{\pgfqpoint{4.813861in}{1.663482in}}%
\pgfpathlineto{\pgfqpoint{4.815984in}{1.668883in}}%
\pgfpathlineto{\pgfqpoint{4.786812in}{1.739886in}}%
\pgfpathlineto{\pgfqpoint{4.799367in}{1.739886in}}%
\pgfpathlineto{\pgfqpoint{4.821922in}{1.683459in}}%
\pgfpathlineto{\pgfqpoint{4.844476in}{1.739886in}}%
\pgfpathlineto{\pgfqpoint{4.857031in}{1.739886in}}%
\pgfpathlineto{\pgfqpoint{4.825323in}{1.661028in}}%
\pgfpathlineto{\pgfqpoint{4.825323in}{1.661028in}}%
\pgfpathclose%
\pgfpathmoveto{\pgfqpoint{4.919272in}{1.678676in}}%
\pgfpathlineto{\pgfqpoint{4.940527in}{1.678676in}}%
\pgfpathlineto{\pgfqpoint{4.940527in}{1.752070in}}%
\pgfpathlineto{\pgfqpoint{4.917396in}{1.747431in}}%
\pgfpathlineto{\pgfqpoint{4.917396in}{1.759286in}}%
\pgfpathlineto{\pgfqpoint{4.940404in}{1.763925in}}%
\pgfpathlineto{\pgfqpoint{4.953413in}{1.763925in}}%
\pgfpathlineto{\pgfqpoint{4.953413in}{1.678676in}}%
\pgfpathlineto{\pgfqpoint{4.974668in}{1.678676in}}%
\pgfpathlineto{\pgfqpoint{4.974668in}{1.667729in}}%
\pgfpathlineto{\pgfqpoint{4.919272in}{1.667729in}}%
\pgfpathlineto{\pgfqpoint{4.919272in}{1.678676in}}%
\pgfpathlineto{\pgfqpoint{4.919272in}{1.678676in}}%
\pgfpathclose%
\pgfpathmoveto{\pgfqpoint{5.003222in}{1.678676in}}%
\pgfpathlineto{\pgfqpoint{5.024477in}{1.678676in}}%
\pgfpathlineto{\pgfqpoint{5.024477in}{1.752070in}}%
\pgfpathlineto{\pgfqpoint{5.001345in}{1.747431in}}%
\pgfpathlineto{\pgfqpoint{5.001345in}{1.759286in}}%
\pgfpathlineto{\pgfqpoint{5.024353in}{1.763925in}}%
\pgfpathlineto{\pgfqpoint{5.037362in}{1.763925in}}%
\pgfpathlineto{\pgfqpoint{5.037362in}{1.678676in}}%
\pgfpathlineto{\pgfqpoint{5.058618in}{1.678676in}}%
\pgfpathlineto{\pgfqpoint{5.058618in}{1.667729in}}%
\pgfpathlineto{\pgfqpoint{5.003222in}{1.667729in}}%
\pgfpathlineto{\pgfqpoint{5.003222in}{1.678676in}}%
\pgfpathlineto{\pgfqpoint{5.003222in}{1.678676in}}%
\pgfpathclose%
\pgfpathmoveto{\pgfqpoint{5.153638in}{1.767842in}}%
\pgfpathquadraticcurveto{\pgfqpoint{5.145021in}{1.753039in}}{\pgfqpoint{5.140815in}{1.738525in}}%
\pgfpathquadraticcurveto{\pgfqpoint{5.136630in}{1.724032in}}{\pgfqpoint{5.136630in}{1.709147in}}%
\pgfpathquadraticcurveto{\pgfqpoint{5.136630in}{1.694283in}}{\pgfqpoint{5.140856in}{1.679686in}}%
\pgfpathquadraticcurveto{\pgfqpoint{5.145082in}{1.665090in}}{\pgfqpoint{5.153638in}{1.650329in}}%
\pgfpathlineto{\pgfqpoint{5.143330in}{1.650329in}}%
\pgfpathquadraticcurveto{\pgfqpoint{5.133682in}{1.665482in}}{\pgfqpoint{5.128878in}{1.680099in}}%
\pgfpathquadraticcurveto{\pgfqpoint{5.124074in}{1.694716in}}{\pgfqpoint{5.124074in}{1.709147in}}%
\pgfpathquadraticcurveto{\pgfqpoint{5.124074in}{1.723517in}}{\pgfqpoint{5.128837in}{1.738072in}}%
\pgfpathquadraticcurveto{\pgfqpoint{5.133620in}{1.752647in}}{\pgfqpoint{5.143330in}{1.767842in}}%
\pgfpathlineto{\pgfqpoint{5.153638in}{1.767842in}}%
\pgfpathlineto{\pgfqpoint{5.153638in}{1.767842in}}%
\pgfpathclose%
\pgfpathmoveto{\pgfqpoint{5.189531in}{1.678676in}}%
\pgfpathlineto{\pgfqpoint{5.234949in}{1.678676in}}%
\pgfpathlineto{\pgfqpoint{5.234949in}{1.667729in}}%
\pgfpathlineto{\pgfqpoint{5.173883in}{1.667729in}}%
\pgfpathlineto{\pgfqpoint{5.173883in}{1.678676in}}%
\pgfpathquadraticcurveto{\pgfqpoint{5.181285in}{1.686345in}}{\pgfqpoint{5.194067in}{1.699251in}}%
\pgfpathquadraticcurveto{\pgfqpoint{5.206870in}{1.712178in}}{\pgfqpoint{5.210148in}{1.715930in}}%
\pgfpathquadraticcurveto{\pgfqpoint{5.216394in}{1.722939in}}{\pgfqpoint{5.218868in}{1.727805in}}%
\pgfpathquadraticcurveto{\pgfqpoint{5.221363in}{1.732670in}}{\pgfqpoint{5.221363in}{1.737371in}}%
\pgfpathquadraticcurveto{\pgfqpoint{5.221363in}{1.745040in}}{\pgfqpoint{5.215982in}{1.749864in}}%
\pgfpathquadraticcurveto{\pgfqpoint{5.210601in}{1.754709in}}{\pgfqpoint{5.201963in}{1.754709in}}%
\pgfpathquadraticcurveto{\pgfqpoint{5.195840in}{1.754709in}}{\pgfqpoint{5.189036in}{1.752586in}}%
\pgfpathquadraticcurveto{\pgfqpoint{5.182254in}{1.750462in}}{\pgfqpoint{5.174523in}{1.746133in}}%
\pgfpathlineto{\pgfqpoint{5.174523in}{1.759286in}}%
\pgfpathquadraticcurveto{\pgfqpoint{5.182377in}{1.762440in}}{\pgfqpoint{5.189201in}{1.764048in}}%
\pgfpathquadraticcurveto{\pgfqpoint{5.196046in}{1.765656in}}{\pgfqpoint{5.201715in}{1.765656in}}%
\pgfpathquadraticcurveto{\pgfqpoint{5.216662in}{1.765656in}}{\pgfqpoint{5.225548in}{1.758173in}}%
\pgfpathquadraticcurveto{\pgfqpoint{5.234434in}{1.750709in}}{\pgfqpoint{5.234434in}{1.738216in}}%
\pgfpathquadraticcurveto{\pgfqpoint{5.234434in}{1.732278in}}{\pgfqpoint{5.232207in}{1.726959in}}%
\pgfpathquadraticcurveto{\pgfqpoint{5.230001in}{1.721661in}}{\pgfqpoint{5.224125in}{1.714445in}}%
\pgfpathquadraticcurveto{\pgfqpoint{5.222517in}{1.712569in}}{\pgfqpoint{5.213879in}{1.703642in}}%
\pgfpathquadraticcurveto{\pgfqpoint{5.205262in}{1.694716in}}{\pgfqpoint{5.189531in}{1.678676in}}%
\pgfpathlineto{\pgfqpoint{5.189531in}{1.678676in}}%
\pgfpathclose%
\pgfpathmoveto{\pgfqpoint{5.290098in}{1.755348in}}%
\pgfpathquadraticcurveto{\pgfqpoint{5.280058in}{1.755348in}}{\pgfqpoint{5.274986in}{1.745452in}}%
\pgfpathquadraticcurveto{\pgfqpoint{5.269935in}{1.735577in}}{\pgfqpoint{5.269935in}{1.715724in}}%
\pgfpathquadraticcurveto{\pgfqpoint{5.269935in}{1.695953in}}{\pgfqpoint{5.274986in}{1.686057in}}%
\pgfpathquadraticcurveto{\pgfqpoint{5.280058in}{1.676161in}}{\pgfqpoint{5.290098in}{1.676161in}}%
\pgfpathquadraticcurveto{\pgfqpoint{5.300220in}{1.676161in}}{\pgfqpoint{5.305271in}{1.686057in}}%
\pgfpathquadraticcurveto{\pgfqpoint{5.310343in}{1.695953in}}{\pgfqpoint{5.310343in}{1.715724in}}%
\pgfpathquadraticcurveto{\pgfqpoint{5.310343in}{1.735577in}}{\pgfqpoint{5.305271in}{1.745452in}}%
\pgfpathquadraticcurveto{\pgfqpoint{5.300220in}{1.755348in}}{\pgfqpoint{5.290098in}{1.755348in}}%
\pgfpathlineto{\pgfqpoint{5.290098in}{1.755348in}}%
\pgfpathclose%
\pgfpathmoveto{\pgfqpoint{5.290098in}{1.765656in}}%
\pgfpathquadraticcurveto{\pgfqpoint{5.306281in}{1.765656in}}{\pgfqpoint{5.314817in}{1.752854in}}%
\pgfpathquadraticcurveto{\pgfqpoint{5.323352in}{1.740071in}}{\pgfqpoint{5.323352in}{1.715724in}}%
\pgfpathquadraticcurveto{\pgfqpoint{5.323352in}{1.691438in}}{\pgfqpoint{5.314817in}{1.678635in}}%
\pgfpathquadraticcurveto{\pgfqpoint{5.306281in}{1.665853in}}{\pgfqpoint{5.290098in}{1.665853in}}%
\pgfpathquadraticcurveto{\pgfqpoint{5.273934in}{1.665853in}}{\pgfqpoint{5.265399in}{1.678635in}}%
\pgfpathquadraticcurveto{\pgfqpoint{5.256864in}{1.691438in}}{\pgfqpoint{5.256864in}{1.715724in}}%
\pgfpathquadraticcurveto{\pgfqpoint{5.256864in}{1.740071in}}{\pgfqpoint{5.265399in}{1.752854in}}%
\pgfpathquadraticcurveto{\pgfqpoint{5.273934in}{1.765656in}}{\pgfqpoint{5.290098in}{1.765656in}}%
\pgfpathlineto{\pgfqpoint{5.290098in}{1.765656in}}%
\pgfpathclose%
\pgfpathmoveto{\pgfqpoint{5.357431in}{1.678676in}}%
\pgfpathlineto{\pgfqpoint{5.402848in}{1.678676in}}%
\pgfpathlineto{\pgfqpoint{5.402848in}{1.667729in}}%
\pgfpathlineto{\pgfqpoint{5.341783in}{1.667729in}}%
\pgfpathlineto{\pgfqpoint{5.341783in}{1.678676in}}%
\pgfpathquadraticcurveto{\pgfqpoint{5.349184in}{1.686345in}}{\pgfqpoint{5.361966in}{1.699251in}}%
\pgfpathquadraticcurveto{\pgfqpoint{5.374769in}{1.712178in}}{\pgfqpoint{5.378047in}{1.715930in}}%
\pgfpathquadraticcurveto{\pgfqpoint{5.384294in}{1.722939in}}{\pgfqpoint{5.386768in}{1.727805in}}%
\pgfpathquadraticcurveto{\pgfqpoint{5.389262in}{1.732670in}}{\pgfqpoint{5.389262in}{1.737371in}}%
\pgfpathquadraticcurveto{\pgfqpoint{5.389262in}{1.745040in}}{\pgfqpoint{5.383881in}{1.749864in}}%
\pgfpathquadraticcurveto{\pgfqpoint{5.378500in}{1.754709in}}{\pgfqpoint{5.369862in}{1.754709in}}%
\pgfpathquadraticcurveto{\pgfqpoint{5.363739in}{1.754709in}}{\pgfqpoint{5.356936in}{1.752586in}}%
\pgfpathquadraticcurveto{\pgfqpoint{5.350153in}{1.750462in}}{\pgfqpoint{5.342422in}{1.746133in}}%
\pgfpathlineto{\pgfqpoint{5.342422in}{1.759286in}}%
\pgfpathquadraticcurveto{\pgfqpoint{5.350277in}{1.762440in}}{\pgfqpoint{5.357101in}{1.764048in}}%
\pgfpathquadraticcurveto{\pgfqpoint{5.363945in}{1.765656in}}{\pgfqpoint{5.369615in}{1.765656in}}%
\pgfpathquadraticcurveto{\pgfqpoint{5.384562in}{1.765656in}}{\pgfqpoint{5.393447in}{1.758173in}}%
\pgfpathquadraticcurveto{\pgfqpoint{5.402333in}{1.750709in}}{\pgfqpoint{5.402333in}{1.738216in}}%
\pgfpathquadraticcurveto{\pgfqpoint{5.402333in}{1.732278in}}{\pgfqpoint{5.400106in}{1.726959in}}%
\pgfpathquadraticcurveto{\pgfqpoint{5.397900in}{1.721661in}}{\pgfqpoint{5.392025in}{1.714445in}}%
\pgfpathquadraticcurveto{\pgfqpoint{5.390417in}{1.712569in}}{\pgfqpoint{5.381778in}{1.703642in}}%
\pgfpathquadraticcurveto{\pgfqpoint{5.373161in}{1.694716in}}{\pgfqpoint{5.357431in}{1.678676in}}%
\pgfpathlineto{\pgfqpoint{5.357431in}{1.678676in}}%
\pgfpathclose%
\pgfpathmoveto{\pgfqpoint{5.432433in}{1.678676in}}%
\pgfpathlineto{\pgfqpoint{5.453688in}{1.678676in}}%
\pgfpathlineto{\pgfqpoint{5.453688in}{1.752070in}}%
\pgfpathlineto{\pgfqpoint{5.430557in}{1.747431in}}%
\pgfpathlineto{\pgfqpoint{5.430557in}{1.759286in}}%
\pgfpathlineto{\pgfqpoint{5.453564in}{1.763925in}}%
\pgfpathlineto{\pgfqpoint{5.466573in}{1.763925in}}%
\pgfpathlineto{\pgfqpoint{5.466573in}{1.678676in}}%
\pgfpathlineto{\pgfqpoint{5.487829in}{1.678676in}}%
\pgfpathlineto{\pgfqpoint{5.487829in}{1.667729in}}%
\pgfpathlineto{\pgfqpoint{5.432433in}{1.667729in}}%
\pgfpathlineto{\pgfqpoint{5.432433in}{1.678676in}}%
\pgfpathlineto{\pgfqpoint{5.432433in}{1.678676in}}%
\pgfpathclose%
\pgfpathmoveto{\pgfqpoint{5.510589in}{1.767842in}}%
\pgfpathlineto{\pgfqpoint{5.520897in}{1.767842in}}%
\pgfpathquadraticcurveto{\pgfqpoint{5.530546in}{1.752647in}}{\pgfqpoint{5.535349in}{1.738072in}}%
\pgfpathquadraticcurveto{\pgfqpoint{5.540153in}{1.723517in}}{\pgfqpoint{5.540153in}{1.709147in}}%
\pgfpathquadraticcurveto{\pgfqpoint{5.540153in}{1.694716in}}{\pgfqpoint{5.535349in}{1.680099in}}%
\pgfpathquadraticcurveto{\pgfqpoint{5.530546in}{1.665482in}}{\pgfqpoint{5.520897in}{1.650329in}}%
\pgfpathlineto{\pgfqpoint{5.510589in}{1.650329in}}%
\pgfpathquadraticcurveto{\pgfqpoint{5.519145in}{1.665090in}}{\pgfqpoint{5.523371in}{1.679686in}}%
\pgfpathquadraticcurveto{\pgfqpoint{5.527598in}{1.694283in}}{\pgfqpoint{5.527598in}{1.709147in}}%
\pgfpathquadraticcurveto{\pgfqpoint{5.527598in}{1.724032in}}{\pgfqpoint{5.523371in}{1.738525in}}%
\pgfpathquadraticcurveto{\pgfqpoint{5.519145in}{1.753039in}}{\pgfqpoint{5.510589in}{1.767842in}}%
\pgfpathlineto{\pgfqpoint{5.510589in}{1.767842in}}%
\pgfpathclose%
\pgfusepath{fill}%
\end{pgfscope}%
\begin{pgfscope}%
\pgfsetbuttcap%
\pgfsetroundjoin%
\definecolor{currentfill}{rgb}{0.090196,0.168627,0.227451}%
\pgfsetfillcolor{currentfill}%
\pgfsetlinewidth{0.000000pt}%
\definecolor{currentstroke}{rgb}{0.090196,0.168627,0.227451}%
\pgfsetstrokecolor{currentstroke}%
\pgfsetdash{}{0pt}%
\pgfpathmoveto{\pgfqpoint{4.550611in}{1.584130in}}%
\pgfpathlineto{\pgfqpoint{4.550611in}{1.571431in}}%
\pgfpathquadraticcurveto{\pgfqpoint{4.543210in}{1.574977in}}{\pgfqpoint{4.536633in}{1.576708in}}%
\pgfpathquadraticcurveto{\pgfqpoint{4.530056in}{1.578461in}}{\pgfqpoint{4.523933in}{1.578461in}}%
\pgfpathquadraticcurveto{\pgfqpoint{4.513316in}{1.578461in}}{\pgfqpoint{4.507543in}{1.574337in}}%
\pgfpathquadraticcurveto{\pgfqpoint{4.501771in}{1.570214in}}{\pgfqpoint{4.501771in}{1.562607in}}%
\pgfpathquadraticcurveto{\pgfqpoint{4.501771in}{1.556216in}}{\pgfqpoint{4.505605in}{1.552958in}}%
\pgfpathquadraticcurveto{\pgfqpoint{4.509440in}{1.549722in}}{\pgfqpoint{4.520140in}{1.547722in}}%
\pgfpathlineto{\pgfqpoint{4.527995in}{1.546114in}}%
\pgfpathquadraticcurveto{\pgfqpoint{4.542550in}{1.543331in}}{\pgfqpoint{4.549477in}{1.536342in}}%
\pgfpathquadraticcurveto{\pgfqpoint{4.556404in}{1.529353in}}{\pgfqpoint{4.556404in}{1.517643in}}%
\pgfpathquadraticcurveto{\pgfqpoint{4.556404in}{1.503644in}}{\pgfqpoint{4.547024in}{1.496428in}}%
\pgfpathquadraticcurveto{\pgfqpoint{4.537664in}{1.489213in}}{\pgfqpoint{4.519563in}{1.489213in}}%
\pgfpathquadraticcurveto{\pgfqpoint{4.512739in}{1.489213in}}{\pgfqpoint{4.505028in}{1.490759in}}%
\pgfpathquadraticcurveto{\pgfqpoint{4.497338in}{1.492305in}}{\pgfqpoint{4.489092in}{1.495336in}}%
\pgfpathlineto{\pgfqpoint{4.489092in}{1.508736in}}%
\pgfpathquadraticcurveto{\pgfqpoint{4.497008in}{1.504304in}}{\pgfqpoint{4.504616in}{1.502036in}}%
\pgfpathquadraticcurveto{\pgfqpoint{4.512223in}{1.499789in}}{\pgfqpoint{4.519563in}{1.499789in}}%
\pgfpathquadraticcurveto{\pgfqpoint{4.530696in}{1.499789in}}{\pgfqpoint{4.536757in}{1.504160in}}%
\pgfpathquadraticcurveto{\pgfqpoint{4.542818in}{1.508551in}}{\pgfqpoint{4.542818in}{1.516674in}}%
\pgfpathquadraticcurveto{\pgfqpoint{4.542818in}{1.523745in}}{\pgfqpoint{4.538468in}{1.527745in}}%
\pgfpathquadraticcurveto{\pgfqpoint{4.534118in}{1.531744in}}{\pgfqpoint{4.524201in}{1.533744in}}%
\pgfpathlineto{\pgfqpoint{4.516264in}{1.535290in}}%
\pgfpathquadraticcurveto{\pgfqpoint{4.501709in}{1.538176in}}{\pgfqpoint{4.495194in}{1.544361in}}%
\pgfpathquadraticcurveto{\pgfqpoint{4.488700in}{1.550546in}}{\pgfqpoint{4.488700in}{1.561576in}}%
\pgfpathquadraticcurveto{\pgfqpoint{4.488700in}{1.574337in}}{\pgfqpoint{4.497689in}{1.581677in}}%
\pgfpathquadraticcurveto{\pgfqpoint{4.506678in}{1.589016in}}{\pgfqpoint{4.522449in}{1.589016in}}%
\pgfpathquadraticcurveto{\pgfqpoint{4.529232in}{1.589016in}}{\pgfqpoint{4.536241in}{1.587779in}}%
\pgfpathquadraticcurveto{\pgfqpoint{4.543271in}{1.586563in}}{\pgfqpoint{4.550611in}{1.584130in}}%
\pgfpathlineto{\pgfqpoint{4.550611in}{1.584130in}}%
\pgfpathclose%
\pgfpathmoveto{\pgfqpoint{4.587926in}{1.583739in}}%
\pgfpathlineto{\pgfqpoint{4.587926in}{1.563246in}}%
\pgfpathlineto{\pgfqpoint{4.612336in}{1.563246in}}%
\pgfpathlineto{\pgfqpoint{4.612336in}{1.554030in}}%
\pgfpathlineto{\pgfqpoint{4.587926in}{1.554030in}}%
\pgfpathlineto{\pgfqpoint{4.587926in}{1.514859in}}%
\pgfpathquadraticcurveto{\pgfqpoint{4.587926in}{1.506036in}}{\pgfqpoint{4.590339in}{1.503520in}}%
\pgfpathquadraticcurveto{\pgfqpoint{4.592751in}{1.501005in}}{\pgfqpoint{4.600173in}{1.501005in}}%
\pgfpathlineto{\pgfqpoint{4.612336in}{1.501005in}}%
\pgfpathlineto{\pgfqpoint{4.612336in}{1.491089in}}%
\pgfpathlineto{\pgfqpoint{4.600173in}{1.491089in}}%
\pgfpathquadraticcurveto{\pgfqpoint{4.586442in}{1.491089in}}{\pgfqpoint{4.581226in}{1.496202in}}%
\pgfpathquadraticcurveto{\pgfqpoint{4.576010in}{1.501335in}}{\pgfqpoint{4.576010in}{1.514859in}}%
\pgfpathlineto{\pgfqpoint{4.576010in}{1.554030in}}%
\pgfpathlineto{\pgfqpoint{4.567310in}{1.554030in}}%
\pgfpathlineto{\pgfqpoint{4.567310in}{1.563246in}}%
\pgfpathlineto{\pgfqpoint{4.576010in}{1.563246in}}%
\pgfpathlineto{\pgfqpoint{4.576010in}{1.583739in}}%
\pgfpathlineto{\pgfqpoint{4.587926in}{1.583739in}}%
\pgfpathlineto{\pgfqpoint{4.587926in}{1.583739in}}%
\pgfpathclose%
\pgfpathmoveto{\pgfqpoint{4.626706in}{1.519560in}}%
\pgfpathlineto{\pgfqpoint{4.626706in}{1.563246in}}%
\pgfpathlineto{\pgfqpoint{4.638560in}{1.563246in}}%
\pgfpathlineto{\pgfqpoint{4.638560in}{1.520013in}}%
\pgfpathquadraticcurveto{\pgfqpoint{4.638560in}{1.509767in}}{\pgfqpoint{4.642539in}{1.504634in}}%
\pgfpathquadraticcurveto{\pgfqpoint{4.646539in}{1.499521in}}{\pgfqpoint{4.654538in}{1.499521in}}%
\pgfpathquadraticcurveto{\pgfqpoint{4.664124in}{1.499521in}}{\pgfqpoint{4.669691in}{1.505644in}}%
\pgfpathquadraticcurveto{\pgfqpoint{4.675278in}{1.511767in}}{\pgfqpoint{4.675278in}{1.522343in}}%
\pgfpathlineto{\pgfqpoint{4.675278in}{1.563246in}}%
\pgfpathlineto{\pgfqpoint{4.687132in}{1.563246in}}%
\pgfpathlineto{\pgfqpoint{4.687132in}{1.491089in}}%
\pgfpathlineto{\pgfqpoint{4.675278in}{1.491089in}}%
\pgfpathlineto{\pgfqpoint{4.675278in}{1.502180in}}%
\pgfpathquadraticcurveto{\pgfqpoint{4.670969in}{1.495604in}}{\pgfqpoint{4.665258in}{1.492408in}}%
\pgfpathquadraticcurveto{\pgfqpoint{4.659568in}{1.489213in}}{\pgfqpoint{4.652023in}{1.489213in}}%
\pgfpathquadraticcurveto{\pgfqpoint{4.639591in}{1.489213in}}{\pgfqpoint{4.633138in}{1.496944in}}%
\pgfpathquadraticcurveto{\pgfqpoint{4.626706in}{1.504675in}}{\pgfqpoint{4.626706in}{1.519560in}}%
\pgfpathlineto{\pgfqpoint{4.626706in}{1.519560in}}%
\pgfpathclose%
\pgfpathmoveto{\pgfqpoint{4.656538in}{1.564978in}}%
\pgfpathlineto{\pgfqpoint{4.656538in}{1.564978in}}%
\pgfpathlineto{\pgfqpoint{4.656538in}{1.564978in}}%
\pgfpathclose%
\pgfpathmoveto{\pgfqpoint{4.759021in}{1.552299in}}%
\pgfpathlineto{\pgfqpoint{4.759021in}{1.591346in}}%
\pgfpathlineto{\pgfqpoint{4.770876in}{1.591346in}}%
\pgfpathlineto{\pgfqpoint{4.770876in}{1.491089in}}%
\pgfpathlineto{\pgfqpoint{4.759021in}{1.491089in}}%
\pgfpathlineto{\pgfqpoint{4.759021in}{1.501912in}}%
\pgfpathquadraticcurveto{\pgfqpoint{4.755290in}{1.495480in}}{\pgfqpoint{4.749579in}{1.492346in}}%
\pgfpathquadraticcurveto{\pgfqpoint{4.743889in}{1.489213in}}{\pgfqpoint{4.735890in}{1.489213in}}%
\pgfpathquadraticcurveto{\pgfqpoint{4.722819in}{1.489213in}}{\pgfqpoint{4.714593in}{1.499645in}}%
\pgfpathquadraticcurveto{\pgfqpoint{4.706388in}{1.510097in}}{\pgfqpoint{4.706388in}{1.527106in}}%
\pgfpathquadraticcurveto{\pgfqpoint{4.706388in}{1.544114in}}{\pgfqpoint{4.714593in}{1.554546in}}%
\pgfpathquadraticcurveto{\pgfqpoint{4.722819in}{1.564978in}}{\pgfqpoint{4.735890in}{1.564978in}}%
\pgfpathquadraticcurveto{\pgfqpoint{4.743889in}{1.564978in}}{\pgfqpoint{4.749579in}{1.561844in}}%
\pgfpathquadraticcurveto{\pgfqpoint{4.755290in}{1.558731in}}{\pgfqpoint{4.759021in}{1.552299in}}%
\pgfpathlineto{\pgfqpoint{4.759021in}{1.552299in}}%
\pgfpathclose%
\pgfpathmoveto{\pgfqpoint{4.718634in}{1.527106in}}%
\pgfpathquadraticcurveto{\pgfqpoint{4.718634in}{1.514035in}}{\pgfqpoint{4.724015in}{1.506592in}}%
\pgfpathquadraticcurveto{\pgfqpoint{4.729396in}{1.499150in}}{\pgfqpoint{4.738797in}{1.499150in}}%
\pgfpathquadraticcurveto{\pgfqpoint{4.748198in}{1.499150in}}{\pgfqpoint{4.753599in}{1.506592in}}%
\pgfpathquadraticcurveto{\pgfqpoint{4.759021in}{1.514035in}}{\pgfqpoint{4.759021in}{1.527106in}}%
\pgfpathquadraticcurveto{\pgfqpoint{4.759021in}{1.540176in}}{\pgfqpoint{4.753599in}{1.547619in}}%
\pgfpathquadraticcurveto{\pgfqpoint{4.748198in}{1.555061in}}{\pgfqpoint{4.738797in}{1.555061in}}%
\pgfpathquadraticcurveto{\pgfqpoint{4.729396in}{1.555061in}}{\pgfqpoint{4.724015in}{1.547619in}}%
\pgfpathquadraticcurveto{\pgfqpoint{4.718634in}{1.540176in}}{\pgfqpoint{4.718634in}{1.527106in}}%
\pgfpathlineto{\pgfqpoint{4.718634in}{1.527106in}}%
\pgfpathclose%
\pgfpathmoveto{\pgfqpoint{4.825323in}{1.484388in}}%
\pgfpathquadraticcurveto{\pgfqpoint{4.820314in}{1.471503in}}{\pgfqpoint{4.815531in}{1.467586in}}%
\pgfpathquadraticcurveto{\pgfqpoint{4.810768in}{1.463648in}}{\pgfqpoint{4.802790in}{1.463648in}}%
\pgfpathlineto{\pgfqpoint{4.793306in}{1.463648in}}%
\pgfpathlineto{\pgfqpoint{4.793306in}{1.473565in}}%
\pgfpathlineto{\pgfqpoint{4.800275in}{1.473565in}}%
\pgfpathquadraticcurveto{\pgfqpoint{4.805161in}{1.473565in}}{\pgfqpoint{4.807861in}{1.475895in}}%
\pgfpathquadraticcurveto{\pgfqpoint{4.810583in}{1.478204in}}{\pgfqpoint{4.813861in}{1.486842in}}%
\pgfpathlineto{\pgfqpoint{4.815984in}{1.492243in}}%
\pgfpathlineto{\pgfqpoint{4.786812in}{1.563246in}}%
\pgfpathlineto{\pgfqpoint{4.799367in}{1.563246in}}%
\pgfpathlineto{\pgfqpoint{4.821922in}{1.506819in}}%
\pgfpathlineto{\pgfqpoint{4.844476in}{1.563246in}}%
\pgfpathlineto{\pgfqpoint{4.857031in}{1.563246in}}%
\pgfpathlineto{\pgfqpoint{4.825323in}{1.484388in}}%
\pgfpathlineto{\pgfqpoint{4.825323in}{1.484388in}}%
\pgfpathclose%
\pgfpathmoveto{\pgfqpoint{4.919272in}{1.502036in}}%
\pgfpathlineto{\pgfqpoint{4.940527in}{1.502036in}}%
\pgfpathlineto{\pgfqpoint{4.940527in}{1.575430in}}%
\pgfpathlineto{\pgfqpoint{4.917396in}{1.570791in}}%
\pgfpathlineto{\pgfqpoint{4.917396in}{1.582646in}}%
\pgfpathlineto{\pgfqpoint{4.940404in}{1.587285in}}%
\pgfpathlineto{\pgfqpoint{4.953413in}{1.587285in}}%
\pgfpathlineto{\pgfqpoint{4.953413in}{1.502036in}}%
\pgfpathlineto{\pgfqpoint{4.974668in}{1.502036in}}%
\pgfpathlineto{\pgfqpoint{4.974668in}{1.491089in}}%
\pgfpathlineto{\pgfqpoint{4.919272in}{1.491089in}}%
\pgfpathlineto{\pgfqpoint{4.919272in}{1.502036in}}%
\pgfpathlineto{\pgfqpoint{4.919272in}{1.502036in}}%
\pgfpathclose%
\pgfpathmoveto{\pgfqpoint{5.012169in}{1.502036in}}%
\pgfpathlineto{\pgfqpoint{5.057587in}{1.502036in}}%
\pgfpathlineto{\pgfqpoint{5.057587in}{1.491089in}}%
\pgfpathlineto{\pgfqpoint{4.996521in}{1.491089in}}%
\pgfpathlineto{\pgfqpoint{4.996521in}{1.502036in}}%
\pgfpathquadraticcurveto{\pgfqpoint{5.003923in}{1.509705in}}{\pgfqpoint{5.016705in}{1.522611in}}%
\pgfpathquadraticcurveto{\pgfqpoint{5.029507in}{1.535538in}}{\pgfqpoint{5.032785in}{1.539290in}}%
\pgfpathquadraticcurveto{\pgfqpoint{5.039032in}{1.546299in}}{\pgfqpoint{5.041506in}{1.551165in}}%
\pgfpathquadraticcurveto{\pgfqpoint{5.044001in}{1.556030in}}{\pgfqpoint{5.044001in}{1.560731in}}%
\pgfpathquadraticcurveto{\pgfqpoint{5.044001in}{1.568400in}}{\pgfqpoint{5.038620in}{1.573224in}}%
\pgfpathquadraticcurveto{\pgfqpoint{5.033239in}{1.578069in}}{\pgfqpoint{5.024601in}{1.578069in}}%
\pgfpathquadraticcurveto{\pgfqpoint{5.018478in}{1.578069in}}{\pgfqpoint{5.011674in}{1.575946in}}%
\pgfpathquadraticcurveto{\pgfqpoint{5.004891in}{1.573822in}}{\pgfqpoint{4.997160in}{1.569493in}}%
\pgfpathlineto{\pgfqpoint{4.997160in}{1.582646in}}%
\pgfpathquadraticcurveto{\pgfqpoint{5.005015in}{1.585800in}}{\pgfqpoint{5.011839in}{1.587408in}}%
\pgfpathquadraticcurveto{\pgfqpoint{5.018684in}{1.589016in}}{\pgfqpoint{5.024353in}{1.589016in}}%
\pgfpathquadraticcurveto{\pgfqpoint{5.039300in}{1.589016in}}{\pgfqpoint{5.048186in}{1.581533in}}%
\pgfpathquadraticcurveto{\pgfqpoint{5.057071in}{1.574069in}}{\pgfqpoint{5.057071in}{1.561576in}}%
\pgfpathquadraticcurveto{\pgfqpoint{5.057071in}{1.555638in}}{\pgfqpoint{5.054845in}{1.550319in}}%
\pgfpathquadraticcurveto{\pgfqpoint{5.052639in}{1.545021in}}{\pgfqpoint{5.046763in}{1.537805in}}%
\pgfpathquadraticcurveto{\pgfqpoint{5.045155in}{1.535929in}}{\pgfqpoint{5.036517in}{1.527002in}}%
\pgfpathquadraticcurveto{\pgfqpoint{5.027899in}{1.518076in}}{\pgfqpoint{5.012169in}{1.502036in}}%
\pgfpathlineto{\pgfqpoint{5.012169in}{1.502036in}}%
\pgfpathclose%
\pgfpathmoveto{\pgfqpoint{5.153638in}{1.591202in}}%
\pgfpathquadraticcurveto{\pgfqpoint{5.145021in}{1.576399in}}{\pgfqpoint{5.140815in}{1.561885in}}%
\pgfpathquadraticcurveto{\pgfqpoint{5.136630in}{1.547392in}}{\pgfqpoint{5.136630in}{1.532507in}}%
\pgfpathquadraticcurveto{\pgfqpoint{5.136630in}{1.517643in}}{\pgfqpoint{5.140856in}{1.503046in}}%
\pgfpathquadraticcurveto{\pgfqpoint{5.145082in}{1.488450in}}{\pgfqpoint{5.153638in}{1.473689in}}%
\pgfpathlineto{\pgfqpoint{5.143330in}{1.473689in}}%
\pgfpathquadraticcurveto{\pgfqpoint{5.133682in}{1.488842in}}{\pgfqpoint{5.128878in}{1.503459in}}%
\pgfpathquadraticcurveto{\pgfqpoint{5.124074in}{1.518076in}}{\pgfqpoint{5.124074in}{1.532507in}}%
\pgfpathquadraticcurveto{\pgfqpoint{5.124074in}{1.546877in}}{\pgfqpoint{5.128837in}{1.561432in}}%
\pgfpathquadraticcurveto{\pgfqpoint{5.133620in}{1.576007in}}{\pgfqpoint{5.143330in}{1.591202in}}%
\pgfpathlineto{\pgfqpoint{5.153638in}{1.591202in}}%
\pgfpathlineto{\pgfqpoint{5.153638in}{1.591202in}}%
\pgfpathclose%
\pgfpathmoveto{\pgfqpoint{5.189531in}{1.502036in}}%
\pgfpathlineto{\pgfqpoint{5.234949in}{1.502036in}}%
\pgfpathlineto{\pgfqpoint{5.234949in}{1.491089in}}%
\pgfpathlineto{\pgfqpoint{5.173883in}{1.491089in}}%
\pgfpathlineto{\pgfqpoint{5.173883in}{1.502036in}}%
\pgfpathquadraticcurveto{\pgfqpoint{5.181285in}{1.509705in}}{\pgfqpoint{5.194067in}{1.522611in}}%
\pgfpathquadraticcurveto{\pgfqpoint{5.206870in}{1.535538in}}{\pgfqpoint{5.210148in}{1.539290in}}%
\pgfpathquadraticcurveto{\pgfqpoint{5.216394in}{1.546299in}}{\pgfqpoint{5.218868in}{1.551165in}}%
\pgfpathquadraticcurveto{\pgfqpoint{5.221363in}{1.556030in}}{\pgfqpoint{5.221363in}{1.560731in}}%
\pgfpathquadraticcurveto{\pgfqpoint{5.221363in}{1.568400in}}{\pgfqpoint{5.215982in}{1.573224in}}%
\pgfpathquadraticcurveto{\pgfqpoint{5.210601in}{1.578069in}}{\pgfqpoint{5.201963in}{1.578069in}}%
\pgfpathquadraticcurveto{\pgfqpoint{5.195840in}{1.578069in}}{\pgfqpoint{5.189036in}{1.575946in}}%
\pgfpathquadraticcurveto{\pgfqpoint{5.182254in}{1.573822in}}{\pgfqpoint{5.174523in}{1.569493in}}%
\pgfpathlineto{\pgfqpoint{5.174523in}{1.582646in}}%
\pgfpathquadraticcurveto{\pgfqpoint{5.182377in}{1.585800in}}{\pgfqpoint{5.189201in}{1.587408in}}%
\pgfpathquadraticcurveto{\pgfqpoint{5.196046in}{1.589016in}}{\pgfqpoint{5.201715in}{1.589016in}}%
\pgfpathquadraticcurveto{\pgfqpoint{5.216662in}{1.589016in}}{\pgfqpoint{5.225548in}{1.581533in}}%
\pgfpathquadraticcurveto{\pgfqpoint{5.234434in}{1.574069in}}{\pgfqpoint{5.234434in}{1.561576in}}%
\pgfpathquadraticcurveto{\pgfqpoint{5.234434in}{1.555638in}}{\pgfqpoint{5.232207in}{1.550319in}}%
\pgfpathquadraticcurveto{\pgfqpoint{5.230001in}{1.545021in}}{\pgfqpoint{5.224125in}{1.537805in}}%
\pgfpathquadraticcurveto{\pgfqpoint{5.222517in}{1.535929in}}{\pgfqpoint{5.213879in}{1.527002in}}%
\pgfpathquadraticcurveto{\pgfqpoint{5.205262in}{1.518076in}}{\pgfqpoint{5.189531in}{1.502036in}}%
\pgfpathlineto{\pgfqpoint{5.189531in}{1.502036in}}%
\pgfpathclose%
\pgfpathmoveto{\pgfqpoint{5.290098in}{1.578708in}}%
\pgfpathquadraticcurveto{\pgfqpoint{5.280058in}{1.578708in}}{\pgfqpoint{5.274986in}{1.568812in}}%
\pgfpathquadraticcurveto{\pgfqpoint{5.269935in}{1.558937in}}{\pgfqpoint{5.269935in}{1.539084in}}%
\pgfpathquadraticcurveto{\pgfqpoint{5.269935in}{1.519313in}}{\pgfqpoint{5.274986in}{1.509417in}}%
\pgfpathquadraticcurveto{\pgfqpoint{5.280058in}{1.499521in}}{\pgfqpoint{5.290098in}{1.499521in}}%
\pgfpathquadraticcurveto{\pgfqpoint{5.300220in}{1.499521in}}{\pgfqpoint{5.305271in}{1.509417in}}%
\pgfpathquadraticcurveto{\pgfqpoint{5.310343in}{1.519313in}}{\pgfqpoint{5.310343in}{1.539084in}}%
\pgfpathquadraticcurveto{\pgfqpoint{5.310343in}{1.558937in}}{\pgfqpoint{5.305271in}{1.568812in}}%
\pgfpathquadraticcurveto{\pgfqpoint{5.300220in}{1.578708in}}{\pgfqpoint{5.290098in}{1.578708in}}%
\pgfpathlineto{\pgfqpoint{5.290098in}{1.578708in}}%
\pgfpathclose%
\pgfpathmoveto{\pgfqpoint{5.290098in}{1.589016in}}%
\pgfpathquadraticcurveto{\pgfqpoint{5.306281in}{1.589016in}}{\pgfqpoint{5.314817in}{1.576214in}}%
\pgfpathquadraticcurveto{\pgfqpoint{5.323352in}{1.563431in}}{\pgfqpoint{5.323352in}{1.539084in}}%
\pgfpathquadraticcurveto{\pgfqpoint{5.323352in}{1.514798in}}{\pgfqpoint{5.314817in}{1.501995in}}%
\pgfpathquadraticcurveto{\pgfqpoint{5.306281in}{1.489213in}}{\pgfqpoint{5.290098in}{1.489213in}}%
\pgfpathquadraticcurveto{\pgfqpoint{5.273934in}{1.489213in}}{\pgfqpoint{5.265399in}{1.501995in}}%
\pgfpathquadraticcurveto{\pgfqpoint{5.256864in}{1.514798in}}{\pgfqpoint{5.256864in}{1.539084in}}%
\pgfpathquadraticcurveto{\pgfqpoint{5.256864in}{1.563431in}}{\pgfqpoint{5.265399in}{1.576214in}}%
\pgfpathquadraticcurveto{\pgfqpoint{5.273934in}{1.589016in}}{\pgfqpoint{5.290098in}{1.589016in}}%
\pgfpathlineto{\pgfqpoint{5.290098in}{1.589016in}}%
\pgfpathclose%
\pgfpathmoveto{\pgfqpoint{5.357431in}{1.502036in}}%
\pgfpathlineto{\pgfqpoint{5.402848in}{1.502036in}}%
\pgfpathlineto{\pgfqpoint{5.402848in}{1.491089in}}%
\pgfpathlineto{\pgfqpoint{5.341783in}{1.491089in}}%
\pgfpathlineto{\pgfqpoint{5.341783in}{1.502036in}}%
\pgfpathquadraticcurveto{\pgfqpoint{5.349184in}{1.509705in}}{\pgfqpoint{5.361966in}{1.522611in}}%
\pgfpathquadraticcurveto{\pgfqpoint{5.374769in}{1.535538in}}{\pgfqpoint{5.378047in}{1.539290in}}%
\pgfpathquadraticcurveto{\pgfqpoint{5.384294in}{1.546299in}}{\pgfqpoint{5.386768in}{1.551165in}}%
\pgfpathquadraticcurveto{\pgfqpoint{5.389262in}{1.556030in}}{\pgfqpoint{5.389262in}{1.560731in}}%
\pgfpathquadraticcurveto{\pgfqpoint{5.389262in}{1.568400in}}{\pgfqpoint{5.383881in}{1.573224in}}%
\pgfpathquadraticcurveto{\pgfqpoint{5.378500in}{1.578069in}}{\pgfqpoint{5.369862in}{1.578069in}}%
\pgfpathquadraticcurveto{\pgfqpoint{5.363739in}{1.578069in}}{\pgfqpoint{5.356936in}{1.575946in}}%
\pgfpathquadraticcurveto{\pgfqpoint{5.350153in}{1.573822in}}{\pgfqpoint{5.342422in}{1.569493in}}%
\pgfpathlineto{\pgfqpoint{5.342422in}{1.582646in}}%
\pgfpathquadraticcurveto{\pgfqpoint{5.350277in}{1.585800in}}{\pgfqpoint{5.357101in}{1.587408in}}%
\pgfpathquadraticcurveto{\pgfqpoint{5.363945in}{1.589016in}}{\pgfqpoint{5.369615in}{1.589016in}}%
\pgfpathquadraticcurveto{\pgfqpoint{5.384562in}{1.589016in}}{\pgfqpoint{5.393447in}{1.581533in}}%
\pgfpathquadraticcurveto{\pgfqpoint{5.402333in}{1.574069in}}{\pgfqpoint{5.402333in}{1.561576in}}%
\pgfpathquadraticcurveto{\pgfqpoint{5.402333in}{1.555638in}}{\pgfqpoint{5.400106in}{1.550319in}}%
\pgfpathquadraticcurveto{\pgfqpoint{5.397900in}{1.545021in}}{\pgfqpoint{5.392025in}{1.537805in}}%
\pgfpathquadraticcurveto{\pgfqpoint{5.390417in}{1.535929in}}{\pgfqpoint{5.381778in}{1.527002in}}%
\pgfpathquadraticcurveto{\pgfqpoint{5.373161in}{1.518076in}}{\pgfqpoint{5.357431in}{1.502036in}}%
\pgfpathlineto{\pgfqpoint{5.357431in}{1.502036in}}%
\pgfpathclose%
\pgfpathmoveto{\pgfqpoint{5.432433in}{1.502036in}}%
\pgfpathlineto{\pgfqpoint{5.453688in}{1.502036in}}%
\pgfpathlineto{\pgfqpoint{5.453688in}{1.575430in}}%
\pgfpathlineto{\pgfqpoint{5.430557in}{1.570791in}}%
\pgfpathlineto{\pgfqpoint{5.430557in}{1.582646in}}%
\pgfpathlineto{\pgfqpoint{5.453564in}{1.587285in}}%
\pgfpathlineto{\pgfqpoint{5.466573in}{1.587285in}}%
\pgfpathlineto{\pgfqpoint{5.466573in}{1.502036in}}%
\pgfpathlineto{\pgfqpoint{5.487829in}{1.502036in}}%
\pgfpathlineto{\pgfqpoint{5.487829in}{1.491089in}}%
\pgfpathlineto{\pgfqpoint{5.432433in}{1.491089in}}%
\pgfpathlineto{\pgfqpoint{5.432433in}{1.502036in}}%
\pgfpathlineto{\pgfqpoint{5.432433in}{1.502036in}}%
\pgfpathclose%
\pgfpathmoveto{\pgfqpoint{5.510589in}{1.591202in}}%
\pgfpathlineto{\pgfqpoint{5.520897in}{1.591202in}}%
\pgfpathquadraticcurveto{\pgfqpoint{5.530546in}{1.576007in}}{\pgfqpoint{5.535349in}{1.561432in}}%
\pgfpathquadraticcurveto{\pgfqpoint{5.540153in}{1.546877in}}{\pgfqpoint{5.540153in}{1.532507in}}%
\pgfpathquadraticcurveto{\pgfqpoint{5.540153in}{1.518076in}}{\pgfqpoint{5.535349in}{1.503459in}}%
\pgfpathquadraticcurveto{\pgfqpoint{5.530546in}{1.488842in}}{\pgfqpoint{5.520897in}{1.473689in}}%
\pgfpathlineto{\pgfqpoint{5.510589in}{1.473689in}}%
\pgfpathquadraticcurveto{\pgfqpoint{5.519145in}{1.488450in}}{\pgfqpoint{5.523371in}{1.503046in}}%
\pgfpathquadraticcurveto{\pgfqpoint{5.527598in}{1.517643in}}{\pgfqpoint{5.527598in}{1.532507in}}%
\pgfpathquadraticcurveto{\pgfqpoint{5.527598in}{1.547392in}}{\pgfqpoint{5.523371in}{1.561885in}}%
\pgfpathquadraticcurveto{\pgfqpoint{5.519145in}{1.576399in}}{\pgfqpoint{5.510589in}{1.591202in}}%
\pgfpathlineto{\pgfqpoint{5.510589in}{1.591202in}}%
\pgfpathclose%
\pgfusepath{fill}%
\end{pgfscope}%
\begin{pgfscope}%
\pgfsetbuttcap%
\pgfsetroundjoin%
\definecolor{currentfill}{rgb}{0.090196,0.168627,0.227451}%
\pgfsetfillcolor{currentfill}%
\pgfsetlinewidth{0.000000pt}%
\definecolor{currentstroke}{rgb}{0.090196,0.168627,0.227451}%
\pgfsetstrokecolor{currentstroke}%
\pgfsetdash{}{0pt}%
\pgfpathmoveto{\pgfqpoint{4.550611in}{1.407490in}}%
\pgfpathlineto{\pgfqpoint{4.550611in}{1.394791in}}%
\pgfpathquadraticcurveto{\pgfqpoint{4.543210in}{1.398337in}}{\pgfqpoint{4.536633in}{1.400068in}}%
\pgfpathquadraticcurveto{\pgfqpoint{4.530056in}{1.401821in}}{\pgfqpoint{4.523933in}{1.401821in}}%
\pgfpathquadraticcurveto{\pgfqpoint{4.513316in}{1.401821in}}{\pgfqpoint{4.507543in}{1.397697in}}%
\pgfpathquadraticcurveto{\pgfqpoint{4.501771in}{1.393574in}}{\pgfqpoint{4.501771in}{1.385967in}}%
\pgfpathquadraticcurveto{\pgfqpoint{4.501771in}{1.379576in}}{\pgfqpoint{4.505605in}{1.376318in}}%
\pgfpathquadraticcurveto{\pgfqpoint{4.509440in}{1.373082in}}{\pgfqpoint{4.520140in}{1.371082in}}%
\pgfpathlineto{\pgfqpoint{4.527995in}{1.369474in}}%
\pgfpathquadraticcurveto{\pgfqpoint{4.542550in}{1.366691in}}{\pgfqpoint{4.549477in}{1.359702in}}%
\pgfpathquadraticcurveto{\pgfqpoint{4.556404in}{1.352713in}}{\pgfqpoint{4.556404in}{1.341003in}}%
\pgfpathquadraticcurveto{\pgfqpoint{4.556404in}{1.327004in}}{\pgfqpoint{4.547024in}{1.319788in}}%
\pgfpathquadraticcurveto{\pgfqpoint{4.537664in}{1.312573in}}{\pgfqpoint{4.519563in}{1.312573in}}%
\pgfpathquadraticcurveto{\pgfqpoint{4.512739in}{1.312573in}}{\pgfqpoint{4.505028in}{1.314119in}}%
\pgfpathquadraticcurveto{\pgfqpoint{4.497338in}{1.315665in}}{\pgfqpoint{4.489092in}{1.318696in}}%
\pgfpathlineto{\pgfqpoint{4.489092in}{1.332096in}}%
\pgfpathquadraticcurveto{\pgfqpoint{4.497008in}{1.327664in}}{\pgfqpoint{4.504616in}{1.325396in}}%
\pgfpathquadraticcurveto{\pgfqpoint{4.512223in}{1.323149in}}{\pgfqpoint{4.519563in}{1.323149in}}%
\pgfpathquadraticcurveto{\pgfqpoint{4.530696in}{1.323149in}}{\pgfqpoint{4.536757in}{1.327520in}}%
\pgfpathquadraticcurveto{\pgfqpoint{4.542818in}{1.331911in}}{\pgfqpoint{4.542818in}{1.340034in}}%
\pgfpathquadraticcurveto{\pgfqpoint{4.542818in}{1.347105in}}{\pgfqpoint{4.538468in}{1.351105in}}%
\pgfpathquadraticcurveto{\pgfqpoint{4.534118in}{1.355104in}}{\pgfqpoint{4.524201in}{1.357104in}}%
\pgfpathlineto{\pgfqpoint{4.516264in}{1.358650in}}%
\pgfpathquadraticcurveto{\pgfqpoint{4.501709in}{1.361536in}}{\pgfqpoint{4.495194in}{1.367721in}}%
\pgfpathquadraticcurveto{\pgfqpoint{4.488700in}{1.373906in}}{\pgfqpoint{4.488700in}{1.384936in}}%
\pgfpathquadraticcurveto{\pgfqpoint{4.488700in}{1.397697in}}{\pgfqpoint{4.497689in}{1.405037in}}%
\pgfpathquadraticcurveto{\pgfqpoint{4.506678in}{1.412376in}}{\pgfqpoint{4.522449in}{1.412376in}}%
\pgfpathquadraticcurveto{\pgfqpoint{4.529232in}{1.412376in}}{\pgfqpoint{4.536241in}{1.411139in}}%
\pgfpathquadraticcurveto{\pgfqpoint{4.543271in}{1.409923in}}{\pgfqpoint{4.550611in}{1.407490in}}%
\pgfpathlineto{\pgfqpoint{4.550611in}{1.407490in}}%
\pgfpathclose%
\pgfpathmoveto{\pgfqpoint{4.587926in}{1.407099in}}%
\pgfpathlineto{\pgfqpoint{4.587926in}{1.386606in}}%
\pgfpathlineto{\pgfqpoint{4.612336in}{1.386606in}}%
\pgfpathlineto{\pgfqpoint{4.612336in}{1.377390in}}%
\pgfpathlineto{\pgfqpoint{4.587926in}{1.377390in}}%
\pgfpathlineto{\pgfqpoint{4.587926in}{1.338219in}}%
\pgfpathquadraticcurveto{\pgfqpoint{4.587926in}{1.329396in}}{\pgfqpoint{4.590339in}{1.326880in}}%
\pgfpathquadraticcurveto{\pgfqpoint{4.592751in}{1.324365in}}{\pgfqpoint{4.600173in}{1.324365in}}%
\pgfpathlineto{\pgfqpoint{4.612336in}{1.324365in}}%
\pgfpathlineto{\pgfqpoint{4.612336in}{1.314449in}}%
\pgfpathlineto{\pgfqpoint{4.600173in}{1.314449in}}%
\pgfpathquadraticcurveto{\pgfqpoint{4.586442in}{1.314449in}}{\pgfqpoint{4.581226in}{1.319562in}}%
\pgfpathquadraticcurveto{\pgfqpoint{4.576010in}{1.324695in}}{\pgfqpoint{4.576010in}{1.338219in}}%
\pgfpathlineto{\pgfqpoint{4.576010in}{1.377390in}}%
\pgfpathlineto{\pgfqpoint{4.567310in}{1.377390in}}%
\pgfpathlineto{\pgfqpoint{4.567310in}{1.386606in}}%
\pgfpathlineto{\pgfqpoint{4.576010in}{1.386606in}}%
\pgfpathlineto{\pgfqpoint{4.576010in}{1.407099in}}%
\pgfpathlineto{\pgfqpoint{4.587926in}{1.407099in}}%
\pgfpathlineto{\pgfqpoint{4.587926in}{1.407099in}}%
\pgfpathclose%
\pgfpathmoveto{\pgfqpoint{4.626706in}{1.342920in}}%
\pgfpathlineto{\pgfqpoint{4.626706in}{1.386606in}}%
\pgfpathlineto{\pgfqpoint{4.638560in}{1.386606in}}%
\pgfpathlineto{\pgfqpoint{4.638560in}{1.343373in}}%
\pgfpathquadraticcurveto{\pgfqpoint{4.638560in}{1.333127in}}{\pgfqpoint{4.642539in}{1.327994in}}%
\pgfpathquadraticcurveto{\pgfqpoint{4.646539in}{1.322881in}}{\pgfqpoint{4.654538in}{1.322881in}}%
\pgfpathquadraticcurveto{\pgfqpoint{4.664124in}{1.322881in}}{\pgfqpoint{4.669691in}{1.329004in}}%
\pgfpathquadraticcurveto{\pgfqpoint{4.675278in}{1.335127in}}{\pgfqpoint{4.675278in}{1.345703in}}%
\pgfpathlineto{\pgfqpoint{4.675278in}{1.386606in}}%
\pgfpathlineto{\pgfqpoint{4.687132in}{1.386606in}}%
\pgfpathlineto{\pgfqpoint{4.687132in}{1.314449in}}%
\pgfpathlineto{\pgfqpoint{4.675278in}{1.314449in}}%
\pgfpathlineto{\pgfqpoint{4.675278in}{1.325540in}}%
\pgfpathquadraticcurveto{\pgfqpoint{4.670969in}{1.318964in}}{\pgfqpoint{4.665258in}{1.315768in}}%
\pgfpathquadraticcurveto{\pgfqpoint{4.659568in}{1.312573in}}{\pgfqpoint{4.652023in}{1.312573in}}%
\pgfpathquadraticcurveto{\pgfqpoint{4.639591in}{1.312573in}}{\pgfqpoint{4.633138in}{1.320304in}}%
\pgfpathquadraticcurveto{\pgfqpoint{4.626706in}{1.328035in}}{\pgfqpoint{4.626706in}{1.342920in}}%
\pgfpathlineto{\pgfqpoint{4.626706in}{1.342920in}}%
\pgfpathclose%
\pgfpathmoveto{\pgfqpoint{4.656538in}{1.388338in}}%
\pgfpathlineto{\pgfqpoint{4.656538in}{1.388338in}}%
\pgfpathlineto{\pgfqpoint{4.656538in}{1.388338in}}%
\pgfpathclose%
\pgfpathmoveto{\pgfqpoint{4.759021in}{1.375659in}}%
\pgfpathlineto{\pgfqpoint{4.759021in}{1.414706in}}%
\pgfpathlineto{\pgfqpoint{4.770876in}{1.414706in}}%
\pgfpathlineto{\pgfqpoint{4.770876in}{1.314449in}}%
\pgfpathlineto{\pgfqpoint{4.759021in}{1.314449in}}%
\pgfpathlineto{\pgfqpoint{4.759021in}{1.325272in}}%
\pgfpathquadraticcurveto{\pgfqpoint{4.755290in}{1.318840in}}{\pgfqpoint{4.749579in}{1.315706in}}%
\pgfpathquadraticcurveto{\pgfqpoint{4.743889in}{1.312573in}}{\pgfqpoint{4.735890in}{1.312573in}}%
\pgfpathquadraticcurveto{\pgfqpoint{4.722819in}{1.312573in}}{\pgfqpoint{4.714593in}{1.323005in}}%
\pgfpathquadraticcurveto{\pgfqpoint{4.706388in}{1.333457in}}{\pgfqpoint{4.706388in}{1.350466in}}%
\pgfpathquadraticcurveto{\pgfqpoint{4.706388in}{1.367474in}}{\pgfqpoint{4.714593in}{1.377906in}}%
\pgfpathquadraticcurveto{\pgfqpoint{4.722819in}{1.388338in}}{\pgfqpoint{4.735890in}{1.388338in}}%
\pgfpathquadraticcurveto{\pgfqpoint{4.743889in}{1.388338in}}{\pgfqpoint{4.749579in}{1.385204in}}%
\pgfpathquadraticcurveto{\pgfqpoint{4.755290in}{1.382091in}}{\pgfqpoint{4.759021in}{1.375659in}}%
\pgfpathlineto{\pgfqpoint{4.759021in}{1.375659in}}%
\pgfpathclose%
\pgfpathmoveto{\pgfqpoint{4.718634in}{1.350466in}}%
\pgfpathquadraticcurveto{\pgfqpoint{4.718634in}{1.337395in}}{\pgfqpoint{4.724015in}{1.329952in}}%
\pgfpathquadraticcurveto{\pgfqpoint{4.729396in}{1.322510in}}{\pgfqpoint{4.738797in}{1.322510in}}%
\pgfpathquadraticcurveto{\pgfqpoint{4.748198in}{1.322510in}}{\pgfqpoint{4.753599in}{1.329952in}}%
\pgfpathquadraticcurveto{\pgfqpoint{4.759021in}{1.337395in}}{\pgfqpoint{4.759021in}{1.350466in}}%
\pgfpathquadraticcurveto{\pgfqpoint{4.759021in}{1.363536in}}{\pgfqpoint{4.753599in}{1.370979in}}%
\pgfpathquadraticcurveto{\pgfqpoint{4.748198in}{1.378421in}}{\pgfqpoint{4.738797in}{1.378421in}}%
\pgfpathquadraticcurveto{\pgfqpoint{4.729396in}{1.378421in}}{\pgfqpoint{4.724015in}{1.370979in}}%
\pgfpathquadraticcurveto{\pgfqpoint{4.718634in}{1.363536in}}{\pgfqpoint{4.718634in}{1.350466in}}%
\pgfpathlineto{\pgfqpoint{4.718634in}{1.350466in}}%
\pgfpathclose%
\pgfpathmoveto{\pgfqpoint{4.825323in}{1.307748in}}%
\pgfpathquadraticcurveto{\pgfqpoint{4.820314in}{1.294863in}}{\pgfqpoint{4.815531in}{1.290946in}}%
\pgfpathquadraticcurveto{\pgfqpoint{4.810768in}{1.287008in}}{\pgfqpoint{4.802790in}{1.287008in}}%
\pgfpathlineto{\pgfqpoint{4.793306in}{1.287008in}}%
\pgfpathlineto{\pgfqpoint{4.793306in}{1.296925in}}%
\pgfpathlineto{\pgfqpoint{4.800275in}{1.296925in}}%
\pgfpathquadraticcurveto{\pgfqpoint{4.805161in}{1.296925in}}{\pgfqpoint{4.807861in}{1.299255in}}%
\pgfpathquadraticcurveto{\pgfqpoint{4.810583in}{1.301564in}}{\pgfqpoint{4.813861in}{1.310202in}}%
\pgfpathlineto{\pgfqpoint{4.815984in}{1.315603in}}%
\pgfpathlineto{\pgfqpoint{4.786812in}{1.386606in}}%
\pgfpathlineto{\pgfqpoint{4.799367in}{1.386606in}}%
\pgfpathlineto{\pgfqpoint{4.821922in}{1.330179in}}%
\pgfpathlineto{\pgfqpoint{4.844476in}{1.386606in}}%
\pgfpathlineto{\pgfqpoint{4.857031in}{1.386606in}}%
\pgfpathlineto{\pgfqpoint{4.825323in}{1.307748in}}%
\pgfpathlineto{\pgfqpoint{4.825323in}{1.307748in}}%
\pgfpathclose%
\pgfpathmoveto{\pgfqpoint{4.919272in}{1.325396in}}%
\pgfpathlineto{\pgfqpoint{4.940527in}{1.325396in}}%
\pgfpathlineto{\pgfqpoint{4.940527in}{1.398790in}}%
\pgfpathlineto{\pgfqpoint{4.917396in}{1.394151in}}%
\pgfpathlineto{\pgfqpoint{4.917396in}{1.406006in}}%
\pgfpathlineto{\pgfqpoint{4.940404in}{1.410645in}}%
\pgfpathlineto{\pgfqpoint{4.953413in}{1.410645in}}%
\pgfpathlineto{\pgfqpoint{4.953413in}{1.325396in}}%
\pgfpathlineto{\pgfqpoint{4.974668in}{1.325396in}}%
\pgfpathlineto{\pgfqpoint{4.974668in}{1.314449in}}%
\pgfpathlineto{\pgfqpoint{4.919272in}{1.314449in}}%
\pgfpathlineto{\pgfqpoint{4.919272in}{1.325396in}}%
\pgfpathlineto{\pgfqpoint{4.919272in}{1.325396in}}%
\pgfpathclose%
\pgfpathmoveto{\pgfqpoint{5.040393in}{1.366319in}}%
\pgfpathquadraticcurveto{\pgfqpoint{5.049732in}{1.364320in}}{\pgfqpoint{5.054969in}{1.357990in}}%
\pgfpathquadraticcurveto{\pgfqpoint{5.060226in}{1.351682in}}{\pgfqpoint{5.060226in}{1.342405in}}%
\pgfpathquadraticcurveto{\pgfqpoint{5.060226in}{1.328179in}}{\pgfqpoint{5.050433in}{1.320366in}}%
\pgfpathquadraticcurveto{\pgfqpoint{5.040640in}{1.312573in}}{\pgfqpoint{5.022601in}{1.312573in}}%
\pgfpathquadraticcurveto{\pgfqpoint{5.016560in}{1.312573in}}{\pgfqpoint{5.010149in}{1.313768in}}%
\pgfpathquadraticcurveto{\pgfqpoint{5.003737in}{1.314964in}}{\pgfqpoint{4.996913in}{1.317356in}}%
\pgfpathlineto{\pgfqpoint{4.996913in}{1.329911in}}%
\pgfpathquadraticcurveto{\pgfqpoint{5.002314in}{1.326757in}}{\pgfqpoint{5.008747in}{1.325149in}}%
\pgfpathquadraticcurveto{\pgfqpoint{5.015200in}{1.323541in}}{\pgfqpoint{5.022230in}{1.323541in}}%
\pgfpathquadraticcurveto{\pgfqpoint{5.034455in}{1.323541in}}{\pgfqpoint{5.040867in}{1.328365in}}%
\pgfpathquadraticcurveto{\pgfqpoint{5.047279in}{1.333189in}}{\pgfqpoint{5.047279in}{1.342405in}}%
\pgfpathquadraticcurveto{\pgfqpoint{5.047279in}{1.350919in}}{\pgfqpoint{5.041321in}{1.355702in}}%
\pgfpathquadraticcurveto{\pgfqpoint{5.035362in}{1.360506in}}{\pgfqpoint{5.024745in}{1.360506in}}%
\pgfpathlineto{\pgfqpoint{5.013530in}{1.360506in}}%
\pgfpathlineto{\pgfqpoint{5.013530in}{1.371206in}}%
\pgfpathlineto{\pgfqpoint{5.025260in}{1.371206in}}%
\pgfpathquadraticcurveto{\pgfqpoint{5.034847in}{1.371206in}}{\pgfqpoint{5.039939in}{1.375040in}}%
\pgfpathquadraticcurveto{\pgfqpoint{5.045031in}{1.378875in}}{\pgfqpoint{5.045031in}{1.386091in}}%
\pgfpathquadraticcurveto{\pgfqpoint{5.045031in}{1.393492in}}{\pgfqpoint{5.039774in}{1.397450in}}%
\pgfpathquadraticcurveto{\pgfqpoint{5.034538in}{1.401429in}}{\pgfqpoint{5.024745in}{1.401429in}}%
\pgfpathquadraticcurveto{\pgfqpoint{5.019385in}{1.401429in}}{\pgfqpoint{5.013262in}{1.400254in}}%
\pgfpathquadraticcurveto{\pgfqpoint{5.007139in}{1.399099in}}{\pgfqpoint{4.999799in}{1.396667in}}%
\pgfpathlineto{\pgfqpoint{4.999799in}{1.408253in}}%
\pgfpathquadraticcurveto{\pgfqpoint{5.007221in}{1.410315in}}{\pgfqpoint{5.013695in}{1.411345in}}%
\pgfpathquadraticcurveto{\pgfqpoint{5.020168in}{1.412376in}}{\pgfqpoint{5.025900in}{1.412376in}}%
\pgfpathquadraticcurveto{\pgfqpoint{5.040723in}{1.412376in}}{\pgfqpoint{5.049340in}{1.405635in}}%
\pgfpathquadraticcurveto{\pgfqpoint{5.057979in}{1.398914in}}{\pgfqpoint{5.057979in}{1.387451in}}%
\pgfpathquadraticcurveto{\pgfqpoint{5.057979in}{1.379452in}}{\pgfqpoint{5.053402in}{1.373947in}}%
\pgfpathquadraticcurveto{\pgfqpoint{5.048825in}{1.368443in}}{\pgfqpoint{5.040393in}{1.366319in}}%
\pgfpathlineto{\pgfqpoint{5.040393in}{1.366319in}}%
\pgfpathclose%
\pgfpathmoveto{\pgfqpoint{5.153638in}{1.414562in}}%
\pgfpathquadraticcurveto{\pgfqpoint{5.145021in}{1.399759in}}{\pgfqpoint{5.140815in}{1.385245in}}%
\pgfpathquadraticcurveto{\pgfqpoint{5.136630in}{1.370752in}}{\pgfqpoint{5.136630in}{1.355867in}}%
\pgfpathquadraticcurveto{\pgfqpoint{5.136630in}{1.341003in}}{\pgfqpoint{5.140856in}{1.326406in}}%
\pgfpathquadraticcurveto{\pgfqpoint{5.145082in}{1.311810in}}{\pgfqpoint{5.153638in}{1.297049in}}%
\pgfpathlineto{\pgfqpoint{5.143330in}{1.297049in}}%
\pgfpathquadraticcurveto{\pgfqpoint{5.133682in}{1.312202in}}{\pgfqpoint{5.128878in}{1.326819in}}%
\pgfpathquadraticcurveto{\pgfqpoint{5.124074in}{1.341436in}}{\pgfqpoint{5.124074in}{1.355867in}}%
\pgfpathquadraticcurveto{\pgfqpoint{5.124074in}{1.370237in}}{\pgfqpoint{5.128837in}{1.384792in}}%
\pgfpathquadraticcurveto{\pgfqpoint{5.133620in}{1.399367in}}{\pgfqpoint{5.143330in}{1.414562in}}%
\pgfpathlineto{\pgfqpoint{5.153638in}{1.414562in}}%
\pgfpathlineto{\pgfqpoint{5.153638in}{1.414562in}}%
\pgfpathclose%
\pgfpathmoveto{\pgfqpoint{5.189531in}{1.325396in}}%
\pgfpathlineto{\pgfqpoint{5.234949in}{1.325396in}}%
\pgfpathlineto{\pgfqpoint{5.234949in}{1.314449in}}%
\pgfpathlineto{\pgfqpoint{5.173883in}{1.314449in}}%
\pgfpathlineto{\pgfqpoint{5.173883in}{1.325396in}}%
\pgfpathquadraticcurveto{\pgfqpoint{5.181285in}{1.333065in}}{\pgfqpoint{5.194067in}{1.345971in}}%
\pgfpathquadraticcurveto{\pgfqpoint{5.206870in}{1.358898in}}{\pgfqpoint{5.210148in}{1.362650in}}%
\pgfpathquadraticcurveto{\pgfqpoint{5.216394in}{1.369659in}}{\pgfqpoint{5.218868in}{1.374525in}}%
\pgfpathquadraticcurveto{\pgfqpoint{5.221363in}{1.379390in}}{\pgfqpoint{5.221363in}{1.384091in}}%
\pgfpathquadraticcurveto{\pgfqpoint{5.221363in}{1.391760in}}{\pgfqpoint{5.215982in}{1.396584in}}%
\pgfpathquadraticcurveto{\pgfqpoint{5.210601in}{1.401429in}}{\pgfqpoint{5.201963in}{1.401429in}}%
\pgfpathquadraticcurveto{\pgfqpoint{5.195840in}{1.401429in}}{\pgfqpoint{5.189036in}{1.399306in}}%
\pgfpathquadraticcurveto{\pgfqpoint{5.182254in}{1.397182in}}{\pgfqpoint{5.174523in}{1.392853in}}%
\pgfpathlineto{\pgfqpoint{5.174523in}{1.406006in}}%
\pgfpathquadraticcurveto{\pgfqpoint{5.182377in}{1.409160in}}{\pgfqpoint{5.189201in}{1.410768in}}%
\pgfpathquadraticcurveto{\pgfqpoint{5.196046in}{1.412376in}}{\pgfqpoint{5.201715in}{1.412376in}}%
\pgfpathquadraticcurveto{\pgfqpoint{5.216662in}{1.412376in}}{\pgfqpoint{5.225548in}{1.404893in}}%
\pgfpathquadraticcurveto{\pgfqpoint{5.234434in}{1.397429in}}{\pgfqpoint{5.234434in}{1.384936in}}%
\pgfpathquadraticcurveto{\pgfqpoint{5.234434in}{1.378998in}}{\pgfqpoint{5.232207in}{1.373679in}}%
\pgfpathquadraticcurveto{\pgfqpoint{5.230001in}{1.368381in}}{\pgfqpoint{5.224125in}{1.361165in}}%
\pgfpathquadraticcurveto{\pgfqpoint{5.222517in}{1.359289in}}{\pgfqpoint{5.213879in}{1.350362in}}%
\pgfpathquadraticcurveto{\pgfqpoint{5.205262in}{1.341436in}}{\pgfqpoint{5.189531in}{1.325396in}}%
\pgfpathlineto{\pgfqpoint{5.189531in}{1.325396in}}%
\pgfpathclose%
\pgfpathmoveto{\pgfqpoint{5.290098in}{1.402068in}}%
\pgfpathquadraticcurveto{\pgfqpoint{5.280058in}{1.402068in}}{\pgfqpoint{5.274986in}{1.392172in}}%
\pgfpathquadraticcurveto{\pgfqpoint{5.269935in}{1.382297in}}{\pgfqpoint{5.269935in}{1.362444in}}%
\pgfpathquadraticcurveto{\pgfqpoint{5.269935in}{1.342673in}}{\pgfqpoint{5.274986in}{1.332777in}}%
\pgfpathquadraticcurveto{\pgfqpoint{5.280058in}{1.322881in}}{\pgfqpoint{5.290098in}{1.322881in}}%
\pgfpathquadraticcurveto{\pgfqpoint{5.300220in}{1.322881in}}{\pgfqpoint{5.305271in}{1.332777in}}%
\pgfpathquadraticcurveto{\pgfqpoint{5.310343in}{1.342673in}}{\pgfqpoint{5.310343in}{1.362444in}}%
\pgfpathquadraticcurveto{\pgfqpoint{5.310343in}{1.382297in}}{\pgfqpoint{5.305271in}{1.392172in}}%
\pgfpathquadraticcurveto{\pgfqpoint{5.300220in}{1.402068in}}{\pgfqpoint{5.290098in}{1.402068in}}%
\pgfpathlineto{\pgfqpoint{5.290098in}{1.402068in}}%
\pgfpathclose%
\pgfpathmoveto{\pgfqpoint{5.290098in}{1.412376in}}%
\pgfpathquadraticcurveto{\pgfqpoint{5.306281in}{1.412376in}}{\pgfqpoint{5.314817in}{1.399574in}}%
\pgfpathquadraticcurveto{\pgfqpoint{5.323352in}{1.386791in}}{\pgfqpoint{5.323352in}{1.362444in}}%
\pgfpathquadraticcurveto{\pgfqpoint{5.323352in}{1.338158in}}{\pgfqpoint{5.314817in}{1.325355in}}%
\pgfpathquadraticcurveto{\pgfqpoint{5.306281in}{1.312573in}}{\pgfqpoint{5.290098in}{1.312573in}}%
\pgfpathquadraticcurveto{\pgfqpoint{5.273934in}{1.312573in}}{\pgfqpoint{5.265399in}{1.325355in}}%
\pgfpathquadraticcurveto{\pgfqpoint{5.256864in}{1.338158in}}{\pgfqpoint{5.256864in}{1.362444in}}%
\pgfpathquadraticcurveto{\pgfqpoint{5.256864in}{1.386791in}}{\pgfqpoint{5.265399in}{1.399574in}}%
\pgfpathquadraticcurveto{\pgfqpoint{5.273934in}{1.412376in}}{\pgfqpoint{5.290098in}{1.412376in}}%
\pgfpathlineto{\pgfqpoint{5.290098in}{1.412376in}}%
\pgfpathclose%
\pgfpathmoveto{\pgfqpoint{5.357431in}{1.325396in}}%
\pgfpathlineto{\pgfqpoint{5.402848in}{1.325396in}}%
\pgfpathlineto{\pgfqpoint{5.402848in}{1.314449in}}%
\pgfpathlineto{\pgfqpoint{5.341783in}{1.314449in}}%
\pgfpathlineto{\pgfqpoint{5.341783in}{1.325396in}}%
\pgfpathquadraticcurveto{\pgfqpoint{5.349184in}{1.333065in}}{\pgfqpoint{5.361966in}{1.345971in}}%
\pgfpathquadraticcurveto{\pgfqpoint{5.374769in}{1.358898in}}{\pgfqpoint{5.378047in}{1.362650in}}%
\pgfpathquadraticcurveto{\pgfqpoint{5.384294in}{1.369659in}}{\pgfqpoint{5.386768in}{1.374525in}}%
\pgfpathquadraticcurveto{\pgfqpoint{5.389262in}{1.379390in}}{\pgfqpoint{5.389262in}{1.384091in}}%
\pgfpathquadraticcurveto{\pgfqpoint{5.389262in}{1.391760in}}{\pgfqpoint{5.383881in}{1.396584in}}%
\pgfpathquadraticcurveto{\pgfqpoint{5.378500in}{1.401429in}}{\pgfqpoint{5.369862in}{1.401429in}}%
\pgfpathquadraticcurveto{\pgfqpoint{5.363739in}{1.401429in}}{\pgfqpoint{5.356936in}{1.399306in}}%
\pgfpathquadraticcurveto{\pgfqpoint{5.350153in}{1.397182in}}{\pgfqpoint{5.342422in}{1.392853in}}%
\pgfpathlineto{\pgfqpoint{5.342422in}{1.406006in}}%
\pgfpathquadraticcurveto{\pgfqpoint{5.350277in}{1.409160in}}{\pgfqpoint{5.357101in}{1.410768in}}%
\pgfpathquadraticcurveto{\pgfqpoint{5.363945in}{1.412376in}}{\pgfqpoint{5.369615in}{1.412376in}}%
\pgfpathquadraticcurveto{\pgfqpoint{5.384562in}{1.412376in}}{\pgfqpoint{5.393447in}{1.404893in}}%
\pgfpathquadraticcurveto{\pgfqpoint{5.402333in}{1.397429in}}{\pgfqpoint{5.402333in}{1.384936in}}%
\pgfpathquadraticcurveto{\pgfqpoint{5.402333in}{1.378998in}}{\pgfqpoint{5.400106in}{1.373679in}}%
\pgfpathquadraticcurveto{\pgfqpoint{5.397900in}{1.368381in}}{\pgfqpoint{5.392025in}{1.361165in}}%
\pgfpathquadraticcurveto{\pgfqpoint{5.390417in}{1.359289in}}{\pgfqpoint{5.381778in}{1.350362in}}%
\pgfpathquadraticcurveto{\pgfqpoint{5.373161in}{1.341436in}}{\pgfqpoint{5.357431in}{1.325396in}}%
\pgfpathlineto{\pgfqpoint{5.357431in}{1.325396in}}%
\pgfpathclose%
\pgfpathmoveto{\pgfqpoint{5.441380in}{1.325396in}}%
\pgfpathlineto{\pgfqpoint{5.486798in}{1.325396in}}%
\pgfpathlineto{\pgfqpoint{5.486798in}{1.314449in}}%
\pgfpathlineto{\pgfqpoint{5.425732in}{1.314449in}}%
\pgfpathlineto{\pgfqpoint{5.425732in}{1.325396in}}%
\pgfpathquadraticcurveto{\pgfqpoint{5.433134in}{1.333065in}}{\pgfqpoint{5.445916in}{1.345971in}}%
\pgfpathquadraticcurveto{\pgfqpoint{5.458719in}{1.358898in}}{\pgfqpoint{5.461997in}{1.362650in}}%
\pgfpathquadraticcurveto{\pgfqpoint{5.468243in}{1.369659in}}{\pgfqpoint{5.470717in}{1.374525in}}%
\pgfpathquadraticcurveto{\pgfqpoint{5.473212in}{1.379390in}}{\pgfqpoint{5.473212in}{1.384091in}}%
\pgfpathquadraticcurveto{\pgfqpoint{5.473212in}{1.391760in}}{\pgfqpoint{5.467831in}{1.396584in}}%
\pgfpathquadraticcurveto{\pgfqpoint{5.462450in}{1.401429in}}{\pgfqpoint{5.453812in}{1.401429in}}%
\pgfpathquadraticcurveto{\pgfqpoint{5.447689in}{1.401429in}}{\pgfqpoint{5.440885in}{1.399306in}}%
\pgfpathquadraticcurveto{\pgfqpoint{5.434103in}{1.397182in}}{\pgfqpoint{5.426372in}{1.392853in}}%
\pgfpathlineto{\pgfqpoint{5.426372in}{1.406006in}}%
\pgfpathquadraticcurveto{\pgfqpoint{5.434226in}{1.409160in}}{\pgfqpoint{5.441050in}{1.410768in}}%
\pgfpathquadraticcurveto{\pgfqpoint{5.447895in}{1.412376in}}{\pgfqpoint{5.453564in}{1.412376in}}%
\pgfpathquadraticcurveto{\pgfqpoint{5.468511in}{1.412376in}}{\pgfqpoint{5.477397in}{1.404893in}}%
\pgfpathquadraticcurveto{\pgfqpoint{5.486283in}{1.397429in}}{\pgfqpoint{5.486283in}{1.384936in}}%
\pgfpathquadraticcurveto{\pgfqpoint{5.486283in}{1.378998in}}{\pgfqpoint{5.484056in}{1.373679in}}%
\pgfpathquadraticcurveto{\pgfqpoint{5.481850in}{1.368381in}}{\pgfqpoint{5.475974in}{1.361165in}}%
\pgfpathquadraticcurveto{\pgfqpoint{5.474366in}{1.359289in}}{\pgfqpoint{5.465728in}{1.350362in}}%
\pgfpathquadraticcurveto{\pgfqpoint{5.457110in}{1.341436in}}{\pgfqpoint{5.441380in}{1.325396in}}%
\pgfpathlineto{\pgfqpoint{5.441380in}{1.325396in}}%
\pgfpathclose%
\pgfpathmoveto{\pgfqpoint{5.510589in}{1.414562in}}%
\pgfpathlineto{\pgfqpoint{5.520897in}{1.414562in}}%
\pgfpathquadraticcurveto{\pgfqpoint{5.530546in}{1.399367in}}{\pgfqpoint{5.535349in}{1.384792in}}%
\pgfpathquadraticcurveto{\pgfqpoint{5.540153in}{1.370237in}}{\pgfqpoint{5.540153in}{1.355867in}}%
\pgfpathquadraticcurveto{\pgfqpoint{5.540153in}{1.341436in}}{\pgfqpoint{5.535349in}{1.326819in}}%
\pgfpathquadraticcurveto{\pgfqpoint{5.530546in}{1.312202in}}{\pgfqpoint{5.520897in}{1.297049in}}%
\pgfpathlineto{\pgfqpoint{5.510589in}{1.297049in}}%
\pgfpathquadraticcurveto{\pgfqpoint{5.519145in}{1.311810in}}{\pgfqpoint{5.523371in}{1.326406in}}%
\pgfpathquadraticcurveto{\pgfqpoint{5.527598in}{1.341003in}}{\pgfqpoint{5.527598in}{1.355867in}}%
\pgfpathquadraticcurveto{\pgfqpoint{5.527598in}{1.370752in}}{\pgfqpoint{5.523371in}{1.385245in}}%
\pgfpathquadraticcurveto{\pgfqpoint{5.519145in}{1.399759in}}{\pgfqpoint{5.510589in}{1.414562in}}%
\pgfpathlineto{\pgfqpoint{5.510589in}{1.414562in}}%
\pgfpathclose%
\pgfusepath{fill}%
\end{pgfscope}%
\begin{pgfscope}%
\pgfsetbuttcap%
\pgfsetroundjoin%
\definecolor{currentfill}{rgb}{0.090196,0.168627,0.227451}%
\pgfsetfillcolor{currentfill}%
\pgfsetlinewidth{0.000000pt}%
\definecolor{currentstroke}{rgb}{0.090196,0.168627,0.227451}%
\pgfsetstrokecolor{currentstroke}%
\pgfsetdash{}{0pt}%
\pgfpathmoveto{\pgfqpoint{7.784611in}{2.290690in}}%
\pgfpathlineto{\pgfqpoint{7.784611in}{2.277991in}}%
\pgfpathquadraticcurveto{\pgfqpoint{7.777210in}{2.281537in}}{\pgfqpoint{7.770633in}{2.283268in}}%
\pgfpathquadraticcurveto{\pgfqpoint{7.764056in}{2.285021in}}{\pgfqpoint{7.757933in}{2.285021in}}%
\pgfpathquadraticcurveto{\pgfqpoint{7.747316in}{2.285021in}}{\pgfqpoint{7.741543in}{2.280897in}}%
\pgfpathquadraticcurveto{\pgfqpoint{7.735771in}{2.276774in}}{\pgfqpoint{7.735771in}{2.269167in}}%
\pgfpathquadraticcurveto{\pgfqpoint{7.735771in}{2.262776in}}{\pgfqpoint{7.739605in}{2.259518in}}%
\pgfpathquadraticcurveto{\pgfqpoint{7.743440in}{2.256282in}}{\pgfqpoint{7.754140in}{2.254282in}}%
\pgfpathlineto{\pgfqpoint{7.761995in}{2.252674in}}%
\pgfpathquadraticcurveto{\pgfqpoint{7.776550in}{2.249891in}}{\pgfqpoint{7.783477in}{2.242902in}}%
\pgfpathquadraticcurveto{\pgfqpoint{7.790404in}{2.235913in}}{\pgfqpoint{7.790404in}{2.224203in}}%
\pgfpathquadraticcurveto{\pgfqpoint{7.790404in}{2.210204in}}{\pgfqpoint{7.781024in}{2.202988in}}%
\pgfpathquadraticcurveto{\pgfqpoint{7.771664in}{2.195773in}}{\pgfqpoint{7.753563in}{2.195773in}}%
\pgfpathquadraticcurveto{\pgfqpoint{7.746739in}{2.195773in}}{\pgfqpoint{7.739028in}{2.197319in}}%
\pgfpathquadraticcurveto{\pgfqpoint{7.731338in}{2.198865in}}{\pgfqpoint{7.723092in}{2.201896in}}%
\pgfpathlineto{\pgfqpoint{7.723092in}{2.215296in}}%
\pgfpathquadraticcurveto{\pgfqpoint{7.731008in}{2.210864in}}{\pgfqpoint{7.738616in}{2.208596in}}%
\pgfpathquadraticcurveto{\pgfqpoint{7.746223in}{2.206349in}}{\pgfqpoint{7.753563in}{2.206349in}}%
\pgfpathquadraticcurveto{\pgfqpoint{7.764696in}{2.206349in}}{\pgfqpoint{7.770757in}{2.210720in}}%
\pgfpathquadraticcurveto{\pgfqpoint{7.776818in}{2.215111in}}{\pgfqpoint{7.776818in}{2.223234in}}%
\pgfpathquadraticcurveto{\pgfqpoint{7.776818in}{2.230305in}}{\pgfqpoint{7.772468in}{2.234305in}}%
\pgfpathquadraticcurveto{\pgfqpoint{7.768118in}{2.238304in}}{\pgfqpoint{7.758201in}{2.240304in}}%
\pgfpathlineto{\pgfqpoint{7.750264in}{2.241850in}}%
\pgfpathquadraticcurveto{\pgfqpoint{7.735709in}{2.244736in}}{\pgfqpoint{7.729194in}{2.250921in}}%
\pgfpathquadraticcurveto{\pgfqpoint{7.722700in}{2.257106in}}{\pgfqpoint{7.722700in}{2.268136in}}%
\pgfpathquadraticcurveto{\pgfqpoint{7.722700in}{2.280897in}}{\pgfqpoint{7.731689in}{2.288237in}}%
\pgfpathquadraticcurveto{\pgfqpoint{7.740678in}{2.295576in}}{\pgfqpoint{7.756449in}{2.295576in}}%
\pgfpathquadraticcurveto{\pgfqpoint{7.763232in}{2.295576in}}{\pgfqpoint{7.770241in}{2.294339in}}%
\pgfpathquadraticcurveto{\pgfqpoint{7.777271in}{2.293123in}}{\pgfqpoint{7.784611in}{2.290690in}}%
\pgfpathlineto{\pgfqpoint{7.784611in}{2.290690in}}%
\pgfpathclose%
\pgfpathmoveto{\pgfqpoint{7.821926in}{2.290299in}}%
\pgfpathlineto{\pgfqpoint{7.821926in}{2.269806in}}%
\pgfpathlineto{\pgfqpoint{7.846336in}{2.269806in}}%
\pgfpathlineto{\pgfqpoint{7.846336in}{2.260590in}}%
\pgfpathlineto{\pgfqpoint{7.821926in}{2.260590in}}%
\pgfpathlineto{\pgfqpoint{7.821926in}{2.221419in}}%
\pgfpathquadraticcurveto{\pgfqpoint{7.821926in}{2.212596in}}{\pgfqpoint{7.824339in}{2.210080in}}%
\pgfpathquadraticcurveto{\pgfqpoint{7.826751in}{2.207565in}}{\pgfqpoint{7.834173in}{2.207565in}}%
\pgfpathlineto{\pgfqpoint{7.846336in}{2.207565in}}%
\pgfpathlineto{\pgfqpoint{7.846336in}{2.197649in}}%
\pgfpathlineto{\pgfqpoint{7.834173in}{2.197649in}}%
\pgfpathquadraticcurveto{\pgfqpoint{7.820442in}{2.197649in}}{\pgfqpoint{7.815226in}{2.202762in}}%
\pgfpathquadraticcurveto{\pgfqpoint{7.810010in}{2.207895in}}{\pgfqpoint{7.810010in}{2.221419in}}%
\pgfpathlineto{\pgfqpoint{7.810010in}{2.260590in}}%
\pgfpathlineto{\pgfqpoint{7.801310in}{2.260590in}}%
\pgfpathlineto{\pgfqpoint{7.801310in}{2.269806in}}%
\pgfpathlineto{\pgfqpoint{7.810010in}{2.269806in}}%
\pgfpathlineto{\pgfqpoint{7.810010in}{2.290299in}}%
\pgfpathlineto{\pgfqpoint{7.821926in}{2.290299in}}%
\pgfpathlineto{\pgfqpoint{7.821926in}{2.290299in}}%
\pgfpathclose%
\pgfpathmoveto{\pgfqpoint{7.860706in}{2.226120in}}%
\pgfpathlineto{\pgfqpoint{7.860706in}{2.269806in}}%
\pgfpathlineto{\pgfqpoint{7.872560in}{2.269806in}}%
\pgfpathlineto{\pgfqpoint{7.872560in}{2.226573in}}%
\pgfpathquadraticcurveto{\pgfqpoint{7.872560in}{2.216327in}}{\pgfqpoint{7.876539in}{2.211194in}}%
\pgfpathquadraticcurveto{\pgfqpoint{7.880539in}{2.206081in}}{\pgfqpoint{7.888538in}{2.206081in}}%
\pgfpathquadraticcurveto{\pgfqpoint{7.898124in}{2.206081in}}{\pgfqpoint{7.903691in}{2.212204in}}%
\pgfpathquadraticcurveto{\pgfqpoint{7.909278in}{2.218327in}}{\pgfqpoint{7.909278in}{2.228903in}}%
\pgfpathlineto{\pgfqpoint{7.909278in}{2.269806in}}%
\pgfpathlineto{\pgfqpoint{7.921132in}{2.269806in}}%
\pgfpathlineto{\pgfqpoint{7.921132in}{2.197649in}}%
\pgfpathlineto{\pgfqpoint{7.909278in}{2.197649in}}%
\pgfpathlineto{\pgfqpoint{7.909278in}{2.208740in}}%
\pgfpathquadraticcurveto{\pgfqpoint{7.904969in}{2.202164in}}{\pgfqpoint{7.899258in}{2.198968in}}%
\pgfpathquadraticcurveto{\pgfqpoint{7.893568in}{2.195773in}}{\pgfqpoint{7.886023in}{2.195773in}}%
\pgfpathquadraticcurveto{\pgfqpoint{7.873591in}{2.195773in}}{\pgfqpoint{7.867138in}{2.203504in}}%
\pgfpathquadraticcurveto{\pgfqpoint{7.860706in}{2.211235in}}{\pgfqpoint{7.860706in}{2.226120in}}%
\pgfpathlineto{\pgfqpoint{7.860706in}{2.226120in}}%
\pgfpathclose%
\pgfpathmoveto{\pgfqpoint{7.890538in}{2.271538in}}%
\pgfpathlineto{\pgfqpoint{7.890538in}{2.271538in}}%
\pgfpathlineto{\pgfqpoint{7.890538in}{2.271538in}}%
\pgfpathclose%
\pgfpathmoveto{\pgfqpoint{7.993021in}{2.258859in}}%
\pgfpathlineto{\pgfqpoint{7.993021in}{2.297906in}}%
\pgfpathlineto{\pgfqpoint{8.004876in}{2.297906in}}%
\pgfpathlineto{\pgfqpoint{8.004876in}{2.197649in}}%
\pgfpathlineto{\pgfqpoint{7.993021in}{2.197649in}}%
\pgfpathlineto{\pgfqpoint{7.993021in}{2.208472in}}%
\pgfpathquadraticcurveto{\pgfqpoint{7.989290in}{2.202040in}}{\pgfqpoint{7.983579in}{2.198906in}}%
\pgfpathquadraticcurveto{\pgfqpoint{7.977889in}{2.195773in}}{\pgfqpoint{7.969890in}{2.195773in}}%
\pgfpathquadraticcurveto{\pgfqpoint{7.956819in}{2.195773in}}{\pgfqpoint{7.948593in}{2.206205in}}%
\pgfpathquadraticcurveto{\pgfqpoint{7.940388in}{2.216657in}}{\pgfqpoint{7.940388in}{2.233666in}}%
\pgfpathquadraticcurveto{\pgfqpoint{7.940388in}{2.250674in}}{\pgfqpoint{7.948593in}{2.261106in}}%
\pgfpathquadraticcurveto{\pgfqpoint{7.956819in}{2.271538in}}{\pgfqpoint{7.969890in}{2.271538in}}%
\pgfpathquadraticcurveto{\pgfqpoint{7.977889in}{2.271538in}}{\pgfqpoint{7.983579in}{2.268404in}}%
\pgfpathquadraticcurveto{\pgfqpoint{7.989290in}{2.265291in}}{\pgfqpoint{7.993021in}{2.258859in}}%
\pgfpathlineto{\pgfqpoint{7.993021in}{2.258859in}}%
\pgfpathclose%
\pgfpathmoveto{\pgfqpoint{7.952634in}{2.233666in}}%
\pgfpathquadraticcurveto{\pgfqpoint{7.952634in}{2.220595in}}{\pgfqpoint{7.958015in}{2.213152in}}%
\pgfpathquadraticcurveto{\pgfqpoint{7.963396in}{2.205710in}}{\pgfqpoint{7.972797in}{2.205710in}}%
\pgfpathquadraticcurveto{\pgfqpoint{7.982198in}{2.205710in}}{\pgfqpoint{7.987599in}{2.213152in}}%
\pgfpathquadraticcurveto{\pgfqpoint{7.993021in}{2.220595in}}{\pgfqpoint{7.993021in}{2.233666in}}%
\pgfpathquadraticcurveto{\pgfqpoint{7.993021in}{2.246736in}}{\pgfqpoint{7.987599in}{2.254179in}}%
\pgfpathquadraticcurveto{\pgfqpoint{7.982198in}{2.261621in}}{\pgfqpoint{7.972797in}{2.261621in}}%
\pgfpathquadraticcurveto{\pgfqpoint{7.963396in}{2.261621in}}{\pgfqpoint{7.958015in}{2.254179in}}%
\pgfpathquadraticcurveto{\pgfqpoint{7.952634in}{2.246736in}}{\pgfqpoint{7.952634in}{2.233666in}}%
\pgfpathlineto{\pgfqpoint{7.952634in}{2.233666in}}%
\pgfpathclose%
\pgfpathmoveto{\pgfqpoint{8.059323in}{2.190948in}}%
\pgfpathquadraticcurveto{\pgfqpoint{8.054314in}{2.178063in}}{\pgfqpoint{8.049531in}{2.174146in}}%
\pgfpathquadraticcurveto{\pgfqpoint{8.044768in}{2.170208in}}{\pgfqpoint{8.036790in}{2.170208in}}%
\pgfpathlineto{\pgfqpoint{8.027306in}{2.170208in}}%
\pgfpathlineto{\pgfqpoint{8.027306in}{2.180125in}}%
\pgfpathlineto{\pgfqpoint{8.034275in}{2.180125in}}%
\pgfpathquadraticcurveto{\pgfqpoint{8.039161in}{2.180125in}}{\pgfqpoint{8.041861in}{2.182455in}}%
\pgfpathquadraticcurveto{\pgfqpoint{8.044583in}{2.184764in}}{\pgfqpoint{8.047861in}{2.193402in}}%
\pgfpathlineto{\pgfqpoint{8.049984in}{2.198803in}}%
\pgfpathlineto{\pgfqpoint{8.020812in}{2.269806in}}%
\pgfpathlineto{\pgfqpoint{8.033367in}{2.269806in}}%
\pgfpathlineto{\pgfqpoint{8.055922in}{2.213379in}}%
\pgfpathlineto{\pgfqpoint{8.078476in}{2.269806in}}%
\pgfpathlineto{\pgfqpoint{8.091031in}{2.269806in}}%
\pgfpathlineto{\pgfqpoint{8.059323in}{2.190948in}}%
\pgfpathlineto{\pgfqpoint{8.059323in}{2.190948in}}%
\pgfpathclose%
\pgfpathmoveto{\pgfqpoint{8.153272in}{2.208596in}}%
\pgfpathlineto{\pgfqpoint{8.174527in}{2.208596in}}%
\pgfpathlineto{\pgfqpoint{8.174527in}{2.281990in}}%
\pgfpathlineto{\pgfqpoint{8.151396in}{2.277351in}}%
\pgfpathlineto{\pgfqpoint{8.151396in}{2.289206in}}%
\pgfpathlineto{\pgfqpoint{8.174404in}{2.293845in}}%
\pgfpathlineto{\pgfqpoint{8.187413in}{2.293845in}}%
\pgfpathlineto{\pgfqpoint{8.187413in}{2.208596in}}%
\pgfpathlineto{\pgfqpoint{8.208668in}{2.208596in}}%
\pgfpathlineto{\pgfqpoint{8.208668in}{2.197649in}}%
\pgfpathlineto{\pgfqpoint{8.153272in}{2.197649in}}%
\pgfpathlineto{\pgfqpoint{8.153272in}{2.208596in}}%
\pgfpathlineto{\pgfqpoint{8.153272in}{2.208596in}}%
\pgfpathclose%
\pgfpathmoveto{\pgfqpoint{8.270723in}{2.282506in}}%
\pgfpathlineto{\pgfqpoint{8.237861in}{2.231150in}}%
\pgfpathlineto{\pgfqpoint{8.270723in}{2.231150in}}%
\pgfpathlineto{\pgfqpoint{8.270723in}{2.282506in}}%
\pgfpathlineto{\pgfqpoint{8.270723in}{2.282506in}}%
\pgfpathclose%
\pgfpathmoveto{\pgfqpoint{8.267301in}{2.293845in}}%
\pgfpathlineto{\pgfqpoint{8.283670in}{2.293845in}}%
\pgfpathlineto{\pgfqpoint{8.283670in}{2.231150in}}%
\pgfpathlineto{\pgfqpoint{8.297401in}{2.231150in}}%
\pgfpathlineto{\pgfqpoint{8.297401in}{2.220327in}}%
\pgfpathlineto{\pgfqpoint{8.283670in}{2.220327in}}%
\pgfpathlineto{\pgfqpoint{8.283670in}{2.197649in}}%
\pgfpathlineto{\pgfqpoint{8.270723in}{2.197649in}}%
\pgfpathlineto{\pgfqpoint{8.270723in}{2.220327in}}%
\pgfpathlineto{\pgfqpoint{8.227305in}{2.220327in}}%
\pgfpathlineto{\pgfqpoint{8.227305in}{2.232882in}}%
\pgfpathlineto{\pgfqpoint{8.267301in}{2.293845in}}%
\pgfpathlineto{\pgfqpoint{8.267301in}{2.293845in}}%
\pgfpathclose%
\pgfpathmoveto{\pgfqpoint{8.387638in}{2.297762in}}%
\pgfpathquadraticcurveto{\pgfqpoint{8.379021in}{2.282959in}}{\pgfqpoint{8.374815in}{2.268445in}}%
\pgfpathquadraticcurveto{\pgfqpoint{8.370630in}{2.253952in}}{\pgfqpoint{8.370630in}{2.239067in}}%
\pgfpathquadraticcurveto{\pgfqpoint{8.370630in}{2.224203in}}{\pgfqpoint{8.374856in}{2.209606in}}%
\pgfpathquadraticcurveto{\pgfqpoint{8.379082in}{2.195010in}}{\pgfqpoint{8.387638in}{2.180249in}}%
\pgfpathlineto{\pgfqpoint{8.377330in}{2.180249in}}%
\pgfpathquadraticcurveto{\pgfqpoint{8.367682in}{2.195402in}}{\pgfqpoint{8.362878in}{2.210019in}}%
\pgfpathquadraticcurveto{\pgfqpoint{8.358074in}{2.224636in}}{\pgfqpoint{8.358074in}{2.239067in}}%
\pgfpathquadraticcurveto{\pgfqpoint{8.358074in}{2.253437in}}{\pgfqpoint{8.362837in}{2.267992in}}%
\pgfpathquadraticcurveto{\pgfqpoint{8.367620in}{2.282567in}}{\pgfqpoint{8.377330in}{2.297762in}}%
\pgfpathlineto{\pgfqpoint{8.387638in}{2.297762in}}%
\pgfpathlineto{\pgfqpoint{8.387638in}{2.297762in}}%
\pgfpathclose%
\pgfpathmoveto{\pgfqpoint{8.423531in}{2.208596in}}%
\pgfpathlineto{\pgfqpoint{8.468949in}{2.208596in}}%
\pgfpathlineto{\pgfqpoint{8.468949in}{2.197649in}}%
\pgfpathlineto{\pgfqpoint{8.407883in}{2.197649in}}%
\pgfpathlineto{\pgfqpoint{8.407883in}{2.208596in}}%
\pgfpathquadraticcurveto{\pgfqpoint{8.415285in}{2.216265in}}{\pgfqpoint{8.428067in}{2.229171in}}%
\pgfpathquadraticcurveto{\pgfqpoint{8.440870in}{2.242098in}}{\pgfqpoint{8.444148in}{2.245850in}}%
\pgfpathquadraticcurveto{\pgfqpoint{8.450394in}{2.252859in}}{\pgfqpoint{8.452868in}{2.257725in}}%
\pgfpathquadraticcurveto{\pgfqpoint{8.455363in}{2.262590in}}{\pgfqpoint{8.455363in}{2.267291in}}%
\pgfpathquadraticcurveto{\pgfqpoint{8.455363in}{2.274960in}}{\pgfqpoint{8.449982in}{2.279784in}}%
\pgfpathquadraticcurveto{\pgfqpoint{8.444601in}{2.284629in}}{\pgfqpoint{8.435963in}{2.284629in}}%
\pgfpathquadraticcurveto{\pgfqpoint{8.429840in}{2.284629in}}{\pgfqpoint{8.423036in}{2.282506in}}%
\pgfpathquadraticcurveto{\pgfqpoint{8.416254in}{2.280382in}}{\pgfqpoint{8.408523in}{2.276053in}}%
\pgfpathlineto{\pgfqpoint{8.408523in}{2.289206in}}%
\pgfpathquadraticcurveto{\pgfqpoint{8.416377in}{2.292360in}}{\pgfqpoint{8.423201in}{2.293968in}}%
\pgfpathquadraticcurveto{\pgfqpoint{8.430046in}{2.295576in}}{\pgfqpoint{8.435715in}{2.295576in}}%
\pgfpathquadraticcurveto{\pgfqpoint{8.450662in}{2.295576in}}{\pgfqpoint{8.459548in}{2.288093in}}%
\pgfpathquadraticcurveto{\pgfqpoint{8.468434in}{2.280629in}}{\pgfqpoint{8.468434in}{2.268136in}}%
\pgfpathquadraticcurveto{\pgfqpoint{8.468434in}{2.262198in}}{\pgfqpoint{8.466207in}{2.256879in}}%
\pgfpathquadraticcurveto{\pgfqpoint{8.464001in}{2.251581in}}{\pgfqpoint{8.458125in}{2.244365in}}%
\pgfpathquadraticcurveto{\pgfqpoint{8.456517in}{2.242489in}}{\pgfqpoint{8.447879in}{2.233562in}}%
\pgfpathquadraticcurveto{\pgfqpoint{8.439262in}{2.224636in}}{\pgfqpoint{8.423531in}{2.208596in}}%
\pgfpathlineto{\pgfqpoint{8.423531in}{2.208596in}}%
\pgfpathclose%
\pgfpathmoveto{\pgfqpoint{8.524098in}{2.285268in}}%
\pgfpathquadraticcurveto{\pgfqpoint{8.514058in}{2.285268in}}{\pgfqpoint{8.508986in}{2.275372in}}%
\pgfpathquadraticcurveto{\pgfqpoint{8.503935in}{2.265497in}}{\pgfqpoint{8.503935in}{2.245644in}}%
\pgfpathquadraticcurveto{\pgfqpoint{8.503935in}{2.225873in}}{\pgfqpoint{8.508986in}{2.215977in}}%
\pgfpathquadraticcurveto{\pgfqpoint{8.514058in}{2.206081in}}{\pgfqpoint{8.524098in}{2.206081in}}%
\pgfpathquadraticcurveto{\pgfqpoint{8.534220in}{2.206081in}}{\pgfqpoint{8.539271in}{2.215977in}}%
\pgfpathquadraticcurveto{\pgfqpoint{8.544343in}{2.225873in}}{\pgfqpoint{8.544343in}{2.245644in}}%
\pgfpathquadraticcurveto{\pgfqpoint{8.544343in}{2.265497in}}{\pgfqpoint{8.539271in}{2.275372in}}%
\pgfpathquadraticcurveto{\pgfqpoint{8.534220in}{2.285268in}}{\pgfqpoint{8.524098in}{2.285268in}}%
\pgfpathlineto{\pgfqpoint{8.524098in}{2.285268in}}%
\pgfpathclose%
\pgfpathmoveto{\pgfqpoint{8.524098in}{2.295576in}}%
\pgfpathquadraticcurveto{\pgfqpoint{8.540281in}{2.295576in}}{\pgfqpoint{8.548817in}{2.282774in}}%
\pgfpathquadraticcurveto{\pgfqpoint{8.557352in}{2.269991in}}{\pgfqpoint{8.557352in}{2.245644in}}%
\pgfpathquadraticcurveto{\pgfqpoint{8.557352in}{2.221358in}}{\pgfqpoint{8.548817in}{2.208555in}}%
\pgfpathquadraticcurveto{\pgfqpoint{8.540281in}{2.195773in}}{\pgfqpoint{8.524098in}{2.195773in}}%
\pgfpathquadraticcurveto{\pgfqpoint{8.507934in}{2.195773in}}{\pgfqpoint{8.499399in}{2.208555in}}%
\pgfpathquadraticcurveto{\pgfqpoint{8.490864in}{2.221358in}}{\pgfqpoint{8.490864in}{2.245644in}}%
\pgfpathquadraticcurveto{\pgfqpoint{8.490864in}{2.269991in}}{\pgfqpoint{8.499399in}{2.282774in}}%
\pgfpathquadraticcurveto{\pgfqpoint{8.507934in}{2.295576in}}{\pgfqpoint{8.524098in}{2.295576in}}%
\pgfpathlineto{\pgfqpoint{8.524098in}{2.295576in}}%
\pgfpathclose%
\pgfpathmoveto{\pgfqpoint{8.591431in}{2.208596in}}%
\pgfpathlineto{\pgfqpoint{8.636848in}{2.208596in}}%
\pgfpathlineto{\pgfqpoint{8.636848in}{2.197649in}}%
\pgfpathlineto{\pgfqpoint{8.575783in}{2.197649in}}%
\pgfpathlineto{\pgfqpoint{8.575783in}{2.208596in}}%
\pgfpathquadraticcurveto{\pgfqpoint{8.583184in}{2.216265in}}{\pgfqpoint{8.595966in}{2.229171in}}%
\pgfpathquadraticcurveto{\pgfqpoint{8.608769in}{2.242098in}}{\pgfqpoint{8.612047in}{2.245850in}}%
\pgfpathquadraticcurveto{\pgfqpoint{8.618294in}{2.252859in}}{\pgfqpoint{8.620768in}{2.257725in}}%
\pgfpathquadraticcurveto{\pgfqpoint{8.623262in}{2.262590in}}{\pgfqpoint{8.623262in}{2.267291in}}%
\pgfpathquadraticcurveto{\pgfqpoint{8.623262in}{2.274960in}}{\pgfqpoint{8.617881in}{2.279784in}}%
\pgfpathquadraticcurveto{\pgfqpoint{8.612500in}{2.284629in}}{\pgfqpoint{8.603862in}{2.284629in}}%
\pgfpathquadraticcurveto{\pgfqpoint{8.597739in}{2.284629in}}{\pgfqpoint{8.590936in}{2.282506in}}%
\pgfpathquadraticcurveto{\pgfqpoint{8.584153in}{2.280382in}}{\pgfqpoint{8.576422in}{2.276053in}}%
\pgfpathlineto{\pgfqpoint{8.576422in}{2.289206in}}%
\pgfpathquadraticcurveto{\pgfqpoint{8.584277in}{2.292360in}}{\pgfqpoint{8.591101in}{2.293968in}}%
\pgfpathquadraticcurveto{\pgfqpoint{8.597945in}{2.295576in}}{\pgfqpoint{8.603615in}{2.295576in}}%
\pgfpathquadraticcurveto{\pgfqpoint{8.618562in}{2.295576in}}{\pgfqpoint{8.627447in}{2.288093in}}%
\pgfpathquadraticcurveto{\pgfqpoint{8.636333in}{2.280629in}}{\pgfqpoint{8.636333in}{2.268136in}}%
\pgfpathquadraticcurveto{\pgfqpoint{8.636333in}{2.262198in}}{\pgfqpoint{8.634106in}{2.256879in}}%
\pgfpathquadraticcurveto{\pgfqpoint{8.631900in}{2.251581in}}{\pgfqpoint{8.626025in}{2.244365in}}%
\pgfpathquadraticcurveto{\pgfqpoint{8.624417in}{2.242489in}}{\pgfqpoint{8.615778in}{2.233562in}}%
\pgfpathquadraticcurveto{\pgfqpoint{8.607161in}{2.224636in}}{\pgfqpoint{8.591431in}{2.208596in}}%
\pgfpathlineto{\pgfqpoint{8.591431in}{2.208596in}}%
\pgfpathclose%
\pgfpathmoveto{\pgfqpoint{8.675380in}{2.208596in}}%
\pgfpathlineto{\pgfqpoint{8.720798in}{2.208596in}}%
\pgfpathlineto{\pgfqpoint{8.720798in}{2.197649in}}%
\pgfpathlineto{\pgfqpoint{8.659732in}{2.197649in}}%
\pgfpathlineto{\pgfqpoint{8.659732in}{2.208596in}}%
\pgfpathquadraticcurveto{\pgfqpoint{8.667134in}{2.216265in}}{\pgfqpoint{8.679916in}{2.229171in}}%
\pgfpathquadraticcurveto{\pgfqpoint{8.692719in}{2.242098in}}{\pgfqpoint{8.695997in}{2.245850in}}%
\pgfpathquadraticcurveto{\pgfqpoint{8.702243in}{2.252859in}}{\pgfqpoint{8.704717in}{2.257725in}}%
\pgfpathquadraticcurveto{\pgfqpoint{8.707212in}{2.262590in}}{\pgfqpoint{8.707212in}{2.267291in}}%
\pgfpathquadraticcurveto{\pgfqpoint{8.707212in}{2.274960in}}{\pgfqpoint{8.701831in}{2.279784in}}%
\pgfpathquadraticcurveto{\pgfqpoint{8.696450in}{2.284629in}}{\pgfqpoint{8.687812in}{2.284629in}}%
\pgfpathquadraticcurveto{\pgfqpoint{8.681689in}{2.284629in}}{\pgfqpoint{8.674885in}{2.282506in}}%
\pgfpathquadraticcurveto{\pgfqpoint{8.668103in}{2.280382in}}{\pgfqpoint{8.660372in}{2.276053in}}%
\pgfpathlineto{\pgfqpoint{8.660372in}{2.289206in}}%
\pgfpathquadraticcurveto{\pgfqpoint{8.668226in}{2.292360in}}{\pgfqpoint{8.675050in}{2.293968in}}%
\pgfpathquadraticcurveto{\pgfqpoint{8.681895in}{2.295576in}}{\pgfqpoint{8.687564in}{2.295576in}}%
\pgfpathquadraticcurveto{\pgfqpoint{8.702511in}{2.295576in}}{\pgfqpoint{8.711397in}{2.288093in}}%
\pgfpathquadraticcurveto{\pgfqpoint{8.720283in}{2.280629in}}{\pgfqpoint{8.720283in}{2.268136in}}%
\pgfpathquadraticcurveto{\pgfqpoint{8.720283in}{2.262198in}}{\pgfqpoint{8.718056in}{2.256879in}}%
\pgfpathquadraticcurveto{\pgfqpoint{8.715850in}{2.251581in}}{\pgfqpoint{8.709974in}{2.244365in}}%
\pgfpathquadraticcurveto{\pgfqpoint{8.708366in}{2.242489in}}{\pgfqpoint{8.699728in}{2.233562in}}%
\pgfpathquadraticcurveto{\pgfqpoint{8.691110in}{2.224636in}}{\pgfqpoint{8.675380in}{2.208596in}}%
\pgfpathlineto{\pgfqpoint{8.675380in}{2.208596in}}%
\pgfpathclose%
\pgfpathmoveto{\pgfqpoint{8.744589in}{2.297762in}}%
\pgfpathlineto{\pgfqpoint{8.754897in}{2.297762in}}%
\pgfpathquadraticcurveto{\pgfqpoint{8.764546in}{2.282567in}}{\pgfqpoint{8.769349in}{2.267992in}}%
\pgfpathquadraticcurveto{\pgfqpoint{8.774153in}{2.253437in}}{\pgfqpoint{8.774153in}{2.239067in}}%
\pgfpathquadraticcurveto{\pgfqpoint{8.774153in}{2.224636in}}{\pgfqpoint{8.769349in}{2.210019in}}%
\pgfpathquadraticcurveto{\pgfqpoint{8.764546in}{2.195402in}}{\pgfqpoint{8.754897in}{2.180249in}}%
\pgfpathlineto{\pgfqpoint{8.744589in}{2.180249in}}%
\pgfpathquadraticcurveto{\pgfqpoint{8.753145in}{2.195010in}}{\pgfqpoint{8.757371in}{2.209606in}}%
\pgfpathquadraticcurveto{\pgfqpoint{8.761598in}{2.224203in}}{\pgfqpoint{8.761598in}{2.239067in}}%
\pgfpathquadraticcurveto{\pgfqpoint{8.761598in}{2.253952in}}{\pgfqpoint{8.757371in}{2.268445in}}%
\pgfpathquadraticcurveto{\pgfqpoint{8.753145in}{2.282959in}}{\pgfqpoint{8.744589in}{2.297762in}}%
\pgfpathlineto{\pgfqpoint{8.744589in}{2.297762in}}%
\pgfpathclose%
\pgfusepath{fill}%
\end{pgfscope}%
\begin{pgfscope}%
\pgfsetbuttcap%
\pgfsetroundjoin%
\definecolor{currentfill}{rgb}{0.090196,0.168627,0.227451}%
\pgfsetfillcolor{currentfill}%
\pgfsetlinewidth{0.000000pt}%
\definecolor{currentstroke}{rgb}{0.090196,0.168627,0.227451}%
\pgfsetstrokecolor{currentstroke}%
\pgfsetdash{}{0pt}%
\pgfpathmoveto{\pgfqpoint{7.784611in}{2.114050in}}%
\pgfpathlineto{\pgfqpoint{7.784611in}{2.101351in}}%
\pgfpathquadraticcurveto{\pgfqpoint{7.777210in}{2.104897in}}{\pgfqpoint{7.770633in}{2.106628in}}%
\pgfpathquadraticcurveto{\pgfqpoint{7.764056in}{2.108381in}}{\pgfqpoint{7.757933in}{2.108381in}}%
\pgfpathquadraticcurveto{\pgfqpoint{7.747316in}{2.108381in}}{\pgfqpoint{7.741543in}{2.104257in}}%
\pgfpathquadraticcurveto{\pgfqpoint{7.735771in}{2.100134in}}{\pgfqpoint{7.735771in}{2.092527in}}%
\pgfpathquadraticcurveto{\pgfqpoint{7.735771in}{2.086136in}}{\pgfqpoint{7.739605in}{2.082878in}}%
\pgfpathquadraticcurveto{\pgfqpoint{7.743440in}{2.079642in}}{\pgfqpoint{7.754140in}{2.077642in}}%
\pgfpathlineto{\pgfqpoint{7.761995in}{2.076034in}}%
\pgfpathquadraticcurveto{\pgfqpoint{7.776550in}{2.073251in}}{\pgfqpoint{7.783477in}{2.066262in}}%
\pgfpathquadraticcurveto{\pgfqpoint{7.790404in}{2.059273in}}{\pgfqpoint{7.790404in}{2.047563in}}%
\pgfpathquadraticcurveto{\pgfqpoint{7.790404in}{2.033564in}}{\pgfqpoint{7.781024in}{2.026348in}}%
\pgfpathquadraticcurveto{\pgfqpoint{7.771664in}{2.019133in}}{\pgfqpoint{7.753563in}{2.019133in}}%
\pgfpathquadraticcurveto{\pgfqpoint{7.746739in}{2.019133in}}{\pgfqpoint{7.739028in}{2.020679in}}%
\pgfpathquadraticcurveto{\pgfqpoint{7.731338in}{2.022225in}}{\pgfqpoint{7.723092in}{2.025256in}}%
\pgfpathlineto{\pgfqpoint{7.723092in}{2.038656in}}%
\pgfpathquadraticcurveto{\pgfqpoint{7.731008in}{2.034224in}}{\pgfqpoint{7.738616in}{2.031956in}}%
\pgfpathquadraticcurveto{\pgfqpoint{7.746223in}{2.029709in}}{\pgfqpoint{7.753563in}{2.029709in}}%
\pgfpathquadraticcurveto{\pgfqpoint{7.764696in}{2.029709in}}{\pgfqpoint{7.770757in}{2.034080in}}%
\pgfpathquadraticcurveto{\pgfqpoint{7.776818in}{2.038471in}}{\pgfqpoint{7.776818in}{2.046594in}}%
\pgfpathquadraticcurveto{\pgfqpoint{7.776818in}{2.053665in}}{\pgfqpoint{7.772468in}{2.057665in}}%
\pgfpathquadraticcurveto{\pgfqpoint{7.768118in}{2.061664in}}{\pgfqpoint{7.758201in}{2.063664in}}%
\pgfpathlineto{\pgfqpoint{7.750264in}{2.065210in}}%
\pgfpathquadraticcurveto{\pgfqpoint{7.735709in}{2.068096in}}{\pgfqpoint{7.729194in}{2.074281in}}%
\pgfpathquadraticcurveto{\pgfqpoint{7.722700in}{2.080466in}}{\pgfqpoint{7.722700in}{2.091496in}}%
\pgfpathquadraticcurveto{\pgfqpoint{7.722700in}{2.104257in}}{\pgfqpoint{7.731689in}{2.111597in}}%
\pgfpathquadraticcurveto{\pgfqpoint{7.740678in}{2.118936in}}{\pgfqpoint{7.756449in}{2.118936in}}%
\pgfpathquadraticcurveto{\pgfqpoint{7.763232in}{2.118936in}}{\pgfqpoint{7.770241in}{2.117699in}}%
\pgfpathquadraticcurveto{\pgfqpoint{7.777271in}{2.116483in}}{\pgfqpoint{7.784611in}{2.114050in}}%
\pgfpathlineto{\pgfqpoint{7.784611in}{2.114050in}}%
\pgfpathclose%
\pgfpathmoveto{\pgfqpoint{7.821926in}{2.113659in}}%
\pgfpathlineto{\pgfqpoint{7.821926in}{2.093166in}}%
\pgfpathlineto{\pgfqpoint{7.846336in}{2.093166in}}%
\pgfpathlineto{\pgfqpoint{7.846336in}{2.083950in}}%
\pgfpathlineto{\pgfqpoint{7.821926in}{2.083950in}}%
\pgfpathlineto{\pgfqpoint{7.821926in}{2.044779in}}%
\pgfpathquadraticcurveto{\pgfqpoint{7.821926in}{2.035956in}}{\pgfqpoint{7.824339in}{2.033440in}}%
\pgfpathquadraticcurveto{\pgfqpoint{7.826751in}{2.030925in}}{\pgfqpoint{7.834173in}{2.030925in}}%
\pgfpathlineto{\pgfqpoint{7.846336in}{2.030925in}}%
\pgfpathlineto{\pgfqpoint{7.846336in}{2.021009in}}%
\pgfpathlineto{\pgfqpoint{7.834173in}{2.021009in}}%
\pgfpathquadraticcurveto{\pgfqpoint{7.820442in}{2.021009in}}{\pgfqpoint{7.815226in}{2.026122in}}%
\pgfpathquadraticcurveto{\pgfqpoint{7.810010in}{2.031255in}}{\pgfqpoint{7.810010in}{2.044779in}}%
\pgfpathlineto{\pgfqpoint{7.810010in}{2.083950in}}%
\pgfpathlineto{\pgfqpoint{7.801310in}{2.083950in}}%
\pgfpathlineto{\pgfqpoint{7.801310in}{2.093166in}}%
\pgfpathlineto{\pgfqpoint{7.810010in}{2.093166in}}%
\pgfpathlineto{\pgfqpoint{7.810010in}{2.113659in}}%
\pgfpathlineto{\pgfqpoint{7.821926in}{2.113659in}}%
\pgfpathlineto{\pgfqpoint{7.821926in}{2.113659in}}%
\pgfpathclose%
\pgfpathmoveto{\pgfqpoint{7.860706in}{2.049480in}}%
\pgfpathlineto{\pgfqpoint{7.860706in}{2.093166in}}%
\pgfpathlineto{\pgfqpoint{7.872560in}{2.093166in}}%
\pgfpathlineto{\pgfqpoint{7.872560in}{2.049933in}}%
\pgfpathquadraticcurveto{\pgfqpoint{7.872560in}{2.039687in}}{\pgfqpoint{7.876539in}{2.034554in}}%
\pgfpathquadraticcurveto{\pgfqpoint{7.880539in}{2.029441in}}{\pgfqpoint{7.888538in}{2.029441in}}%
\pgfpathquadraticcurveto{\pgfqpoint{7.898124in}{2.029441in}}{\pgfqpoint{7.903691in}{2.035564in}}%
\pgfpathquadraticcurveto{\pgfqpoint{7.909278in}{2.041687in}}{\pgfqpoint{7.909278in}{2.052263in}}%
\pgfpathlineto{\pgfqpoint{7.909278in}{2.093166in}}%
\pgfpathlineto{\pgfqpoint{7.921132in}{2.093166in}}%
\pgfpathlineto{\pgfqpoint{7.921132in}{2.021009in}}%
\pgfpathlineto{\pgfqpoint{7.909278in}{2.021009in}}%
\pgfpathlineto{\pgfqpoint{7.909278in}{2.032100in}}%
\pgfpathquadraticcurveto{\pgfqpoint{7.904969in}{2.025524in}}{\pgfqpoint{7.899258in}{2.022328in}}%
\pgfpathquadraticcurveto{\pgfqpoint{7.893568in}{2.019133in}}{\pgfqpoint{7.886023in}{2.019133in}}%
\pgfpathquadraticcurveto{\pgfqpoint{7.873591in}{2.019133in}}{\pgfqpoint{7.867138in}{2.026864in}}%
\pgfpathquadraticcurveto{\pgfqpoint{7.860706in}{2.034595in}}{\pgfqpoint{7.860706in}{2.049480in}}%
\pgfpathlineto{\pgfqpoint{7.860706in}{2.049480in}}%
\pgfpathclose%
\pgfpathmoveto{\pgfqpoint{7.890538in}{2.094898in}}%
\pgfpathlineto{\pgfqpoint{7.890538in}{2.094898in}}%
\pgfpathlineto{\pgfqpoint{7.890538in}{2.094898in}}%
\pgfpathclose%
\pgfpathmoveto{\pgfqpoint{7.993021in}{2.082219in}}%
\pgfpathlineto{\pgfqpoint{7.993021in}{2.121266in}}%
\pgfpathlineto{\pgfqpoint{8.004876in}{2.121266in}}%
\pgfpathlineto{\pgfqpoint{8.004876in}{2.021009in}}%
\pgfpathlineto{\pgfqpoint{7.993021in}{2.021009in}}%
\pgfpathlineto{\pgfqpoint{7.993021in}{2.031832in}}%
\pgfpathquadraticcurveto{\pgfqpoint{7.989290in}{2.025400in}}{\pgfqpoint{7.983579in}{2.022266in}}%
\pgfpathquadraticcurveto{\pgfqpoint{7.977889in}{2.019133in}}{\pgfqpoint{7.969890in}{2.019133in}}%
\pgfpathquadraticcurveto{\pgfqpoint{7.956819in}{2.019133in}}{\pgfqpoint{7.948593in}{2.029565in}}%
\pgfpathquadraticcurveto{\pgfqpoint{7.940388in}{2.040017in}}{\pgfqpoint{7.940388in}{2.057026in}}%
\pgfpathquadraticcurveto{\pgfqpoint{7.940388in}{2.074034in}}{\pgfqpoint{7.948593in}{2.084466in}}%
\pgfpathquadraticcurveto{\pgfqpoint{7.956819in}{2.094898in}}{\pgfqpoint{7.969890in}{2.094898in}}%
\pgfpathquadraticcurveto{\pgfqpoint{7.977889in}{2.094898in}}{\pgfqpoint{7.983579in}{2.091764in}}%
\pgfpathquadraticcurveto{\pgfqpoint{7.989290in}{2.088651in}}{\pgfqpoint{7.993021in}{2.082219in}}%
\pgfpathlineto{\pgfqpoint{7.993021in}{2.082219in}}%
\pgfpathclose%
\pgfpathmoveto{\pgfqpoint{7.952634in}{2.057026in}}%
\pgfpathquadraticcurveto{\pgfqpoint{7.952634in}{2.043955in}}{\pgfqpoint{7.958015in}{2.036512in}}%
\pgfpathquadraticcurveto{\pgfqpoint{7.963396in}{2.029070in}}{\pgfqpoint{7.972797in}{2.029070in}}%
\pgfpathquadraticcurveto{\pgfqpoint{7.982198in}{2.029070in}}{\pgfqpoint{7.987599in}{2.036512in}}%
\pgfpathquadraticcurveto{\pgfqpoint{7.993021in}{2.043955in}}{\pgfqpoint{7.993021in}{2.057026in}}%
\pgfpathquadraticcurveto{\pgfqpoint{7.993021in}{2.070096in}}{\pgfqpoint{7.987599in}{2.077539in}}%
\pgfpathquadraticcurveto{\pgfqpoint{7.982198in}{2.084981in}}{\pgfqpoint{7.972797in}{2.084981in}}%
\pgfpathquadraticcurveto{\pgfqpoint{7.963396in}{2.084981in}}{\pgfqpoint{7.958015in}{2.077539in}}%
\pgfpathquadraticcurveto{\pgfqpoint{7.952634in}{2.070096in}}{\pgfqpoint{7.952634in}{2.057026in}}%
\pgfpathlineto{\pgfqpoint{7.952634in}{2.057026in}}%
\pgfpathclose%
\pgfpathmoveto{\pgfqpoint{8.059323in}{2.014308in}}%
\pgfpathquadraticcurveto{\pgfqpoint{8.054314in}{2.001423in}}{\pgfqpoint{8.049531in}{1.997506in}}%
\pgfpathquadraticcurveto{\pgfqpoint{8.044768in}{1.993568in}}{\pgfqpoint{8.036790in}{1.993568in}}%
\pgfpathlineto{\pgfqpoint{8.027306in}{1.993568in}}%
\pgfpathlineto{\pgfqpoint{8.027306in}{2.003485in}}%
\pgfpathlineto{\pgfqpoint{8.034275in}{2.003485in}}%
\pgfpathquadraticcurveto{\pgfqpoint{8.039161in}{2.003485in}}{\pgfqpoint{8.041861in}{2.005815in}}%
\pgfpathquadraticcurveto{\pgfqpoint{8.044583in}{2.008124in}}{\pgfqpoint{8.047861in}{2.016762in}}%
\pgfpathlineto{\pgfqpoint{8.049984in}{2.022163in}}%
\pgfpathlineto{\pgfqpoint{8.020812in}{2.093166in}}%
\pgfpathlineto{\pgfqpoint{8.033367in}{2.093166in}}%
\pgfpathlineto{\pgfqpoint{8.055922in}{2.036739in}}%
\pgfpathlineto{\pgfqpoint{8.078476in}{2.093166in}}%
\pgfpathlineto{\pgfqpoint{8.091031in}{2.093166in}}%
\pgfpathlineto{\pgfqpoint{8.059323in}{2.014308in}}%
\pgfpathlineto{\pgfqpoint{8.059323in}{2.014308in}}%
\pgfpathclose%
\pgfpathmoveto{\pgfqpoint{8.153272in}{2.031956in}}%
\pgfpathlineto{\pgfqpoint{8.174527in}{2.031956in}}%
\pgfpathlineto{\pgfqpoint{8.174527in}{2.105350in}}%
\pgfpathlineto{\pgfqpoint{8.151396in}{2.100711in}}%
\pgfpathlineto{\pgfqpoint{8.151396in}{2.112566in}}%
\pgfpathlineto{\pgfqpoint{8.174404in}{2.117205in}}%
\pgfpathlineto{\pgfqpoint{8.187413in}{2.117205in}}%
\pgfpathlineto{\pgfqpoint{8.187413in}{2.031956in}}%
\pgfpathlineto{\pgfqpoint{8.208668in}{2.031956in}}%
\pgfpathlineto{\pgfqpoint{8.208668in}{2.021009in}}%
\pgfpathlineto{\pgfqpoint{8.153272in}{2.021009in}}%
\pgfpathlineto{\pgfqpoint{8.153272in}{2.031956in}}%
\pgfpathlineto{\pgfqpoint{8.153272in}{2.031956in}}%
\pgfpathclose%
\pgfpathmoveto{\pgfqpoint{8.235098in}{2.117205in}}%
\pgfpathlineto{\pgfqpoint{8.286185in}{2.117205in}}%
\pgfpathlineto{\pgfqpoint{8.286185in}{2.106237in}}%
\pgfpathlineto{\pgfqpoint{8.247014in}{2.106237in}}%
\pgfpathlineto{\pgfqpoint{8.247014in}{2.082672in}}%
\pgfpathquadraticcurveto{\pgfqpoint{8.249839in}{2.083641in}}{\pgfqpoint{8.252663in}{2.084115in}}%
\pgfpathquadraticcurveto{\pgfqpoint{8.255508in}{2.084590in}}{\pgfqpoint{8.258353in}{2.084590in}}%
\pgfpathquadraticcurveto{\pgfqpoint{8.274455in}{2.084590in}}{\pgfqpoint{8.283856in}{2.075766in}}%
\pgfpathquadraticcurveto{\pgfqpoint{8.293277in}{2.066942in}}{\pgfqpoint{8.293277in}{2.051871in}}%
\pgfpathquadraticcurveto{\pgfqpoint{8.293277in}{2.036347in}}{\pgfqpoint{8.283608in}{2.027730in}}%
\pgfpathquadraticcurveto{\pgfqpoint{8.273939in}{2.019133in}}{\pgfqpoint{8.256354in}{2.019133in}}%
\pgfpathquadraticcurveto{\pgfqpoint{8.250292in}{2.019133in}}{\pgfqpoint{8.244004in}{2.020164in}}%
\pgfpathquadraticcurveto{\pgfqpoint{8.237737in}{2.021194in}}{\pgfqpoint{8.231037in}{2.023256in}}%
\pgfpathlineto{\pgfqpoint{8.231037in}{2.036347in}}%
\pgfpathquadraticcurveto{\pgfqpoint{8.236830in}{2.033193in}}{\pgfqpoint{8.243015in}{2.031647in}}%
\pgfpathquadraticcurveto{\pgfqpoint{8.249200in}{2.030101in}}{\pgfqpoint{8.256086in}{2.030101in}}%
\pgfpathquadraticcurveto{\pgfqpoint{8.267239in}{2.030101in}}{\pgfqpoint{8.273733in}{2.035956in}}%
\pgfpathquadraticcurveto{\pgfqpoint{8.280248in}{2.041811in}}{\pgfqpoint{8.280248in}{2.051871in}}%
\pgfpathquadraticcurveto{\pgfqpoint{8.280248in}{2.061912in}}{\pgfqpoint{8.273733in}{2.067767in}}%
\pgfpathquadraticcurveto{\pgfqpoint{8.267239in}{2.073642in}}{\pgfqpoint{8.256086in}{2.073642in}}%
\pgfpathquadraticcurveto{\pgfqpoint{8.250870in}{2.073642in}}{\pgfqpoint{8.245674in}{2.072488in}}%
\pgfpathquadraticcurveto{\pgfqpoint{8.240500in}{2.071333in}}{\pgfqpoint{8.235098in}{2.068880in}}%
\pgfpathlineto{\pgfqpoint{8.235098in}{2.117205in}}%
\pgfpathlineto{\pgfqpoint{8.235098in}{2.117205in}}%
\pgfpathclose%
\pgfpathmoveto{\pgfqpoint{8.387638in}{2.121122in}}%
\pgfpathquadraticcurveto{\pgfqpoint{8.379021in}{2.106319in}}{\pgfqpoint{8.374815in}{2.091805in}}%
\pgfpathquadraticcurveto{\pgfqpoint{8.370630in}{2.077312in}}{\pgfqpoint{8.370630in}{2.062427in}}%
\pgfpathquadraticcurveto{\pgfqpoint{8.370630in}{2.047563in}}{\pgfqpoint{8.374856in}{2.032966in}}%
\pgfpathquadraticcurveto{\pgfqpoint{8.379082in}{2.018370in}}{\pgfqpoint{8.387638in}{2.003609in}}%
\pgfpathlineto{\pgfqpoint{8.377330in}{2.003609in}}%
\pgfpathquadraticcurveto{\pgfqpoint{8.367682in}{2.018762in}}{\pgfqpoint{8.362878in}{2.033379in}}%
\pgfpathquadraticcurveto{\pgfqpoint{8.358074in}{2.047996in}}{\pgfqpoint{8.358074in}{2.062427in}}%
\pgfpathquadraticcurveto{\pgfqpoint{8.358074in}{2.076797in}}{\pgfqpoint{8.362837in}{2.091352in}}%
\pgfpathquadraticcurveto{\pgfqpoint{8.367620in}{2.105927in}}{\pgfqpoint{8.377330in}{2.121122in}}%
\pgfpathlineto{\pgfqpoint{8.387638in}{2.121122in}}%
\pgfpathlineto{\pgfqpoint{8.387638in}{2.121122in}}%
\pgfpathclose%
\pgfpathmoveto{\pgfqpoint{8.423531in}{2.031956in}}%
\pgfpathlineto{\pgfqpoint{8.468949in}{2.031956in}}%
\pgfpathlineto{\pgfqpoint{8.468949in}{2.021009in}}%
\pgfpathlineto{\pgfqpoint{8.407883in}{2.021009in}}%
\pgfpathlineto{\pgfqpoint{8.407883in}{2.031956in}}%
\pgfpathquadraticcurveto{\pgfqpoint{8.415285in}{2.039625in}}{\pgfqpoint{8.428067in}{2.052531in}}%
\pgfpathquadraticcurveto{\pgfqpoint{8.440870in}{2.065458in}}{\pgfqpoint{8.444148in}{2.069210in}}%
\pgfpathquadraticcurveto{\pgfqpoint{8.450394in}{2.076219in}}{\pgfqpoint{8.452868in}{2.081085in}}%
\pgfpathquadraticcurveto{\pgfqpoint{8.455363in}{2.085950in}}{\pgfqpoint{8.455363in}{2.090651in}}%
\pgfpathquadraticcurveto{\pgfqpoint{8.455363in}{2.098320in}}{\pgfqpoint{8.449982in}{2.103144in}}%
\pgfpathquadraticcurveto{\pgfqpoint{8.444601in}{2.107989in}}{\pgfqpoint{8.435963in}{2.107989in}}%
\pgfpathquadraticcurveto{\pgfqpoint{8.429840in}{2.107989in}}{\pgfqpoint{8.423036in}{2.105866in}}%
\pgfpathquadraticcurveto{\pgfqpoint{8.416254in}{2.103742in}}{\pgfqpoint{8.408523in}{2.099413in}}%
\pgfpathlineto{\pgfqpoint{8.408523in}{2.112566in}}%
\pgfpathquadraticcurveto{\pgfqpoint{8.416377in}{2.115720in}}{\pgfqpoint{8.423201in}{2.117328in}}%
\pgfpathquadraticcurveto{\pgfqpoint{8.430046in}{2.118936in}}{\pgfqpoint{8.435715in}{2.118936in}}%
\pgfpathquadraticcurveto{\pgfqpoint{8.450662in}{2.118936in}}{\pgfqpoint{8.459548in}{2.111453in}}%
\pgfpathquadraticcurveto{\pgfqpoint{8.468434in}{2.103989in}}{\pgfqpoint{8.468434in}{2.091496in}}%
\pgfpathquadraticcurveto{\pgfqpoint{8.468434in}{2.085558in}}{\pgfqpoint{8.466207in}{2.080239in}}%
\pgfpathquadraticcurveto{\pgfqpoint{8.464001in}{2.074941in}}{\pgfqpoint{8.458125in}{2.067725in}}%
\pgfpathquadraticcurveto{\pgfqpoint{8.456517in}{2.065849in}}{\pgfqpoint{8.447879in}{2.056922in}}%
\pgfpathquadraticcurveto{\pgfqpoint{8.439262in}{2.047996in}}{\pgfqpoint{8.423531in}{2.031956in}}%
\pgfpathlineto{\pgfqpoint{8.423531in}{2.031956in}}%
\pgfpathclose%
\pgfpathmoveto{\pgfqpoint{8.524098in}{2.108628in}}%
\pgfpathquadraticcurveto{\pgfqpoint{8.514058in}{2.108628in}}{\pgfqpoint{8.508986in}{2.098732in}}%
\pgfpathquadraticcurveto{\pgfqpoint{8.503935in}{2.088857in}}{\pgfqpoint{8.503935in}{2.069004in}}%
\pgfpathquadraticcurveto{\pgfqpoint{8.503935in}{2.049233in}}{\pgfqpoint{8.508986in}{2.039337in}}%
\pgfpathquadraticcurveto{\pgfqpoint{8.514058in}{2.029441in}}{\pgfqpoint{8.524098in}{2.029441in}}%
\pgfpathquadraticcurveto{\pgfqpoint{8.534220in}{2.029441in}}{\pgfqpoint{8.539271in}{2.039337in}}%
\pgfpathquadraticcurveto{\pgfqpoint{8.544343in}{2.049233in}}{\pgfqpoint{8.544343in}{2.069004in}}%
\pgfpathquadraticcurveto{\pgfqpoint{8.544343in}{2.088857in}}{\pgfqpoint{8.539271in}{2.098732in}}%
\pgfpathquadraticcurveto{\pgfqpoint{8.534220in}{2.108628in}}{\pgfqpoint{8.524098in}{2.108628in}}%
\pgfpathlineto{\pgfqpoint{8.524098in}{2.108628in}}%
\pgfpathclose%
\pgfpathmoveto{\pgfqpoint{8.524098in}{2.118936in}}%
\pgfpathquadraticcurveto{\pgfqpoint{8.540281in}{2.118936in}}{\pgfqpoint{8.548817in}{2.106134in}}%
\pgfpathquadraticcurveto{\pgfqpoint{8.557352in}{2.093351in}}{\pgfqpoint{8.557352in}{2.069004in}}%
\pgfpathquadraticcurveto{\pgfqpoint{8.557352in}{2.044718in}}{\pgfqpoint{8.548817in}{2.031915in}}%
\pgfpathquadraticcurveto{\pgfqpoint{8.540281in}{2.019133in}}{\pgfqpoint{8.524098in}{2.019133in}}%
\pgfpathquadraticcurveto{\pgfqpoint{8.507934in}{2.019133in}}{\pgfqpoint{8.499399in}{2.031915in}}%
\pgfpathquadraticcurveto{\pgfqpoint{8.490864in}{2.044718in}}{\pgfqpoint{8.490864in}{2.069004in}}%
\pgfpathquadraticcurveto{\pgfqpoint{8.490864in}{2.093351in}}{\pgfqpoint{8.499399in}{2.106134in}}%
\pgfpathquadraticcurveto{\pgfqpoint{8.507934in}{2.118936in}}{\pgfqpoint{8.524098in}{2.118936in}}%
\pgfpathlineto{\pgfqpoint{8.524098in}{2.118936in}}%
\pgfpathclose%
\pgfpathmoveto{\pgfqpoint{8.591431in}{2.031956in}}%
\pgfpathlineto{\pgfqpoint{8.636848in}{2.031956in}}%
\pgfpathlineto{\pgfqpoint{8.636848in}{2.021009in}}%
\pgfpathlineto{\pgfqpoint{8.575783in}{2.021009in}}%
\pgfpathlineto{\pgfqpoint{8.575783in}{2.031956in}}%
\pgfpathquadraticcurveto{\pgfqpoint{8.583184in}{2.039625in}}{\pgfqpoint{8.595966in}{2.052531in}}%
\pgfpathquadraticcurveto{\pgfqpoint{8.608769in}{2.065458in}}{\pgfqpoint{8.612047in}{2.069210in}}%
\pgfpathquadraticcurveto{\pgfqpoint{8.618294in}{2.076219in}}{\pgfqpoint{8.620768in}{2.081085in}}%
\pgfpathquadraticcurveto{\pgfqpoint{8.623262in}{2.085950in}}{\pgfqpoint{8.623262in}{2.090651in}}%
\pgfpathquadraticcurveto{\pgfqpoint{8.623262in}{2.098320in}}{\pgfqpoint{8.617881in}{2.103144in}}%
\pgfpathquadraticcurveto{\pgfqpoint{8.612500in}{2.107989in}}{\pgfqpoint{8.603862in}{2.107989in}}%
\pgfpathquadraticcurveto{\pgfqpoint{8.597739in}{2.107989in}}{\pgfqpoint{8.590936in}{2.105866in}}%
\pgfpathquadraticcurveto{\pgfqpoint{8.584153in}{2.103742in}}{\pgfqpoint{8.576422in}{2.099413in}}%
\pgfpathlineto{\pgfqpoint{8.576422in}{2.112566in}}%
\pgfpathquadraticcurveto{\pgfqpoint{8.584277in}{2.115720in}}{\pgfqpoint{8.591101in}{2.117328in}}%
\pgfpathquadraticcurveto{\pgfqpoint{8.597945in}{2.118936in}}{\pgfqpoint{8.603615in}{2.118936in}}%
\pgfpathquadraticcurveto{\pgfqpoint{8.618562in}{2.118936in}}{\pgfqpoint{8.627447in}{2.111453in}}%
\pgfpathquadraticcurveto{\pgfqpoint{8.636333in}{2.103989in}}{\pgfqpoint{8.636333in}{2.091496in}}%
\pgfpathquadraticcurveto{\pgfqpoint{8.636333in}{2.085558in}}{\pgfqpoint{8.634106in}{2.080239in}}%
\pgfpathquadraticcurveto{\pgfqpoint{8.631900in}{2.074941in}}{\pgfqpoint{8.626025in}{2.067725in}}%
\pgfpathquadraticcurveto{\pgfqpoint{8.624417in}{2.065849in}}{\pgfqpoint{8.615778in}{2.056922in}}%
\pgfpathquadraticcurveto{\pgfqpoint{8.607161in}{2.047996in}}{\pgfqpoint{8.591431in}{2.031956in}}%
\pgfpathlineto{\pgfqpoint{8.591431in}{2.031956in}}%
\pgfpathclose%
\pgfpathmoveto{\pgfqpoint{8.703604in}{2.072879in}}%
\pgfpathquadraticcurveto{\pgfqpoint{8.712943in}{2.070880in}}{\pgfqpoint{8.718180in}{2.064550in}}%
\pgfpathquadraticcurveto{\pgfqpoint{8.723437in}{2.058242in}}{\pgfqpoint{8.723437in}{2.048965in}}%
\pgfpathquadraticcurveto{\pgfqpoint{8.723437in}{2.034739in}}{\pgfqpoint{8.713644in}{2.026926in}}%
\pgfpathquadraticcurveto{\pgfqpoint{8.703851in}{2.019133in}}{\pgfqpoint{8.685812in}{2.019133in}}%
\pgfpathquadraticcurveto{\pgfqpoint{8.679771in}{2.019133in}}{\pgfqpoint{8.673360in}{2.020328in}}%
\pgfpathquadraticcurveto{\pgfqpoint{8.666948in}{2.021524in}}{\pgfqpoint{8.660124in}{2.023916in}}%
\pgfpathlineto{\pgfqpoint{8.660124in}{2.036471in}}%
\pgfpathquadraticcurveto{\pgfqpoint{8.665526in}{2.033317in}}{\pgfqpoint{8.671958in}{2.031709in}}%
\pgfpathquadraticcurveto{\pgfqpoint{8.678411in}{2.030101in}}{\pgfqpoint{8.685441in}{2.030101in}}%
\pgfpathquadraticcurveto{\pgfqpoint{8.697666in}{2.030101in}}{\pgfqpoint{8.704078in}{2.034925in}}%
\pgfpathquadraticcurveto{\pgfqpoint{8.710490in}{2.039749in}}{\pgfqpoint{8.710490in}{2.048965in}}%
\pgfpathquadraticcurveto{\pgfqpoint{8.710490in}{2.057479in}}{\pgfqpoint{8.704532in}{2.062262in}}%
\pgfpathquadraticcurveto{\pgfqpoint{8.698574in}{2.067066in}}{\pgfqpoint{8.687956in}{2.067066in}}%
\pgfpathlineto{\pgfqpoint{8.676741in}{2.067066in}}%
\pgfpathlineto{\pgfqpoint{8.676741in}{2.077766in}}%
\pgfpathlineto{\pgfqpoint{8.688472in}{2.077766in}}%
\pgfpathquadraticcurveto{\pgfqpoint{8.698058in}{2.077766in}}{\pgfqpoint{8.703150in}{2.081600in}}%
\pgfpathquadraticcurveto{\pgfqpoint{8.708243in}{2.085435in}}{\pgfqpoint{8.708243in}{2.092651in}}%
\pgfpathquadraticcurveto{\pgfqpoint{8.708243in}{2.100052in}}{\pgfqpoint{8.702985in}{2.104010in}}%
\pgfpathquadraticcurveto{\pgfqpoint{8.697749in}{2.107989in}}{\pgfqpoint{8.687956in}{2.107989in}}%
\pgfpathquadraticcurveto{\pgfqpoint{8.682596in}{2.107989in}}{\pgfqpoint{8.676473in}{2.106814in}}%
\pgfpathquadraticcurveto{\pgfqpoint{8.670350in}{2.105659in}}{\pgfqpoint{8.663010in}{2.103227in}}%
\pgfpathlineto{\pgfqpoint{8.663010in}{2.114813in}}%
\pgfpathquadraticcurveto{\pgfqpoint{8.670432in}{2.116875in}}{\pgfqpoint{8.676906in}{2.117905in}}%
\pgfpathquadraticcurveto{\pgfqpoint{8.683379in}{2.118936in}}{\pgfqpoint{8.689111in}{2.118936in}}%
\pgfpathquadraticcurveto{\pgfqpoint{8.703934in}{2.118936in}}{\pgfqpoint{8.712551in}{2.112195in}}%
\pgfpathquadraticcurveto{\pgfqpoint{8.721190in}{2.105474in}}{\pgfqpoint{8.721190in}{2.094011in}}%
\pgfpathquadraticcurveto{\pgfqpoint{8.721190in}{2.086012in}}{\pgfqpoint{8.716613in}{2.080507in}}%
\pgfpathquadraticcurveto{\pgfqpoint{8.712036in}{2.075003in}}{\pgfqpoint{8.703604in}{2.072879in}}%
\pgfpathlineto{\pgfqpoint{8.703604in}{2.072879in}}%
\pgfpathclose%
\pgfpathmoveto{\pgfqpoint{8.744589in}{2.121122in}}%
\pgfpathlineto{\pgfqpoint{8.754897in}{2.121122in}}%
\pgfpathquadraticcurveto{\pgfqpoint{8.764546in}{2.105927in}}{\pgfqpoint{8.769349in}{2.091352in}}%
\pgfpathquadraticcurveto{\pgfqpoint{8.774153in}{2.076797in}}{\pgfqpoint{8.774153in}{2.062427in}}%
\pgfpathquadraticcurveto{\pgfqpoint{8.774153in}{2.047996in}}{\pgfqpoint{8.769349in}{2.033379in}}%
\pgfpathquadraticcurveto{\pgfqpoint{8.764546in}{2.018762in}}{\pgfqpoint{8.754897in}{2.003609in}}%
\pgfpathlineto{\pgfqpoint{8.744589in}{2.003609in}}%
\pgfpathquadraticcurveto{\pgfqpoint{8.753145in}{2.018370in}}{\pgfqpoint{8.757371in}{2.032966in}}%
\pgfpathquadraticcurveto{\pgfqpoint{8.761598in}{2.047563in}}{\pgfqpoint{8.761598in}{2.062427in}}%
\pgfpathquadraticcurveto{\pgfqpoint{8.761598in}{2.077312in}}{\pgfqpoint{8.757371in}{2.091805in}}%
\pgfpathquadraticcurveto{\pgfqpoint{8.753145in}{2.106319in}}{\pgfqpoint{8.744589in}{2.121122in}}%
\pgfpathlineto{\pgfqpoint{8.744589in}{2.121122in}}%
\pgfpathclose%
\pgfusepath{fill}%
\end{pgfscope}%
\begin{pgfscope}%
\pgfsetbuttcap%
\pgfsetroundjoin%
\definecolor{currentfill}{rgb}{0.090196,0.168627,0.227451}%
\pgfsetfillcolor{currentfill}%
\pgfsetlinewidth{0.000000pt}%
\definecolor{currentstroke}{rgb}{0.090196,0.168627,0.227451}%
\pgfsetstrokecolor{currentstroke}%
\pgfsetdash{}{0pt}%
\pgfpathmoveto{\pgfqpoint{7.784611in}{1.937410in}}%
\pgfpathlineto{\pgfqpoint{7.784611in}{1.924711in}}%
\pgfpathquadraticcurveto{\pgfqpoint{7.777210in}{1.928257in}}{\pgfqpoint{7.770633in}{1.929988in}}%
\pgfpathquadraticcurveto{\pgfqpoint{7.764056in}{1.931741in}}{\pgfqpoint{7.757933in}{1.931741in}}%
\pgfpathquadraticcurveto{\pgfqpoint{7.747316in}{1.931741in}}{\pgfqpoint{7.741543in}{1.927617in}}%
\pgfpathquadraticcurveto{\pgfqpoint{7.735771in}{1.923494in}}{\pgfqpoint{7.735771in}{1.915887in}}%
\pgfpathquadraticcurveto{\pgfqpoint{7.735771in}{1.909496in}}{\pgfqpoint{7.739605in}{1.906238in}}%
\pgfpathquadraticcurveto{\pgfqpoint{7.743440in}{1.903002in}}{\pgfqpoint{7.754140in}{1.901002in}}%
\pgfpathlineto{\pgfqpoint{7.761995in}{1.899394in}}%
\pgfpathquadraticcurveto{\pgfqpoint{7.776550in}{1.896611in}}{\pgfqpoint{7.783477in}{1.889622in}}%
\pgfpathquadraticcurveto{\pgfqpoint{7.790404in}{1.882633in}}{\pgfqpoint{7.790404in}{1.870923in}}%
\pgfpathquadraticcurveto{\pgfqpoint{7.790404in}{1.856924in}}{\pgfqpoint{7.781024in}{1.849708in}}%
\pgfpathquadraticcurveto{\pgfqpoint{7.771664in}{1.842493in}}{\pgfqpoint{7.753563in}{1.842493in}}%
\pgfpathquadraticcurveto{\pgfqpoint{7.746739in}{1.842493in}}{\pgfqpoint{7.739028in}{1.844039in}}%
\pgfpathquadraticcurveto{\pgfqpoint{7.731338in}{1.845585in}}{\pgfqpoint{7.723092in}{1.848616in}}%
\pgfpathlineto{\pgfqpoint{7.723092in}{1.862016in}}%
\pgfpathquadraticcurveto{\pgfqpoint{7.731008in}{1.857584in}}{\pgfqpoint{7.738616in}{1.855316in}}%
\pgfpathquadraticcurveto{\pgfqpoint{7.746223in}{1.853069in}}{\pgfqpoint{7.753563in}{1.853069in}}%
\pgfpathquadraticcurveto{\pgfqpoint{7.764696in}{1.853069in}}{\pgfqpoint{7.770757in}{1.857440in}}%
\pgfpathquadraticcurveto{\pgfqpoint{7.776818in}{1.861831in}}{\pgfqpoint{7.776818in}{1.869954in}}%
\pgfpathquadraticcurveto{\pgfqpoint{7.776818in}{1.877025in}}{\pgfqpoint{7.772468in}{1.881025in}}%
\pgfpathquadraticcurveto{\pgfqpoint{7.768118in}{1.885024in}}{\pgfqpoint{7.758201in}{1.887024in}}%
\pgfpathlineto{\pgfqpoint{7.750264in}{1.888570in}}%
\pgfpathquadraticcurveto{\pgfqpoint{7.735709in}{1.891456in}}{\pgfqpoint{7.729194in}{1.897641in}}%
\pgfpathquadraticcurveto{\pgfqpoint{7.722700in}{1.903826in}}{\pgfqpoint{7.722700in}{1.914856in}}%
\pgfpathquadraticcurveto{\pgfqpoint{7.722700in}{1.927617in}}{\pgfqpoint{7.731689in}{1.934957in}}%
\pgfpathquadraticcurveto{\pgfqpoint{7.740678in}{1.942296in}}{\pgfqpoint{7.756449in}{1.942296in}}%
\pgfpathquadraticcurveto{\pgfqpoint{7.763232in}{1.942296in}}{\pgfqpoint{7.770241in}{1.941059in}}%
\pgfpathquadraticcurveto{\pgfqpoint{7.777271in}{1.939843in}}{\pgfqpoint{7.784611in}{1.937410in}}%
\pgfpathlineto{\pgfqpoint{7.784611in}{1.937410in}}%
\pgfpathclose%
\pgfpathmoveto{\pgfqpoint{7.821926in}{1.937019in}}%
\pgfpathlineto{\pgfqpoint{7.821926in}{1.916526in}}%
\pgfpathlineto{\pgfqpoint{7.846336in}{1.916526in}}%
\pgfpathlineto{\pgfqpoint{7.846336in}{1.907310in}}%
\pgfpathlineto{\pgfqpoint{7.821926in}{1.907310in}}%
\pgfpathlineto{\pgfqpoint{7.821926in}{1.868139in}}%
\pgfpathquadraticcurveto{\pgfqpoint{7.821926in}{1.859316in}}{\pgfqpoint{7.824339in}{1.856800in}}%
\pgfpathquadraticcurveto{\pgfqpoint{7.826751in}{1.854285in}}{\pgfqpoint{7.834173in}{1.854285in}}%
\pgfpathlineto{\pgfqpoint{7.846336in}{1.854285in}}%
\pgfpathlineto{\pgfqpoint{7.846336in}{1.844369in}}%
\pgfpathlineto{\pgfqpoint{7.834173in}{1.844369in}}%
\pgfpathquadraticcurveto{\pgfqpoint{7.820442in}{1.844369in}}{\pgfqpoint{7.815226in}{1.849482in}}%
\pgfpathquadraticcurveto{\pgfqpoint{7.810010in}{1.854615in}}{\pgfqpoint{7.810010in}{1.868139in}}%
\pgfpathlineto{\pgfqpoint{7.810010in}{1.907310in}}%
\pgfpathlineto{\pgfqpoint{7.801310in}{1.907310in}}%
\pgfpathlineto{\pgfqpoint{7.801310in}{1.916526in}}%
\pgfpathlineto{\pgfqpoint{7.810010in}{1.916526in}}%
\pgfpathlineto{\pgfqpoint{7.810010in}{1.937019in}}%
\pgfpathlineto{\pgfqpoint{7.821926in}{1.937019in}}%
\pgfpathlineto{\pgfqpoint{7.821926in}{1.937019in}}%
\pgfpathclose%
\pgfpathmoveto{\pgfqpoint{7.860706in}{1.872840in}}%
\pgfpathlineto{\pgfqpoint{7.860706in}{1.916526in}}%
\pgfpathlineto{\pgfqpoint{7.872560in}{1.916526in}}%
\pgfpathlineto{\pgfqpoint{7.872560in}{1.873293in}}%
\pgfpathquadraticcurveto{\pgfqpoint{7.872560in}{1.863047in}}{\pgfqpoint{7.876539in}{1.857914in}}%
\pgfpathquadraticcurveto{\pgfqpoint{7.880539in}{1.852801in}}{\pgfqpoint{7.888538in}{1.852801in}}%
\pgfpathquadraticcurveto{\pgfqpoint{7.898124in}{1.852801in}}{\pgfqpoint{7.903691in}{1.858924in}}%
\pgfpathquadraticcurveto{\pgfqpoint{7.909278in}{1.865047in}}{\pgfqpoint{7.909278in}{1.875623in}}%
\pgfpathlineto{\pgfqpoint{7.909278in}{1.916526in}}%
\pgfpathlineto{\pgfqpoint{7.921132in}{1.916526in}}%
\pgfpathlineto{\pgfqpoint{7.921132in}{1.844369in}}%
\pgfpathlineto{\pgfqpoint{7.909278in}{1.844369in}}%
\pgfpathlineto{\pgfqpoint{7.909278in}{1.855460in}}%
\pgfpathquadraticcurveto{\pgfqpoint{7.904969in}{1.848884in}}{\pgfqpoint{7.899258in}{1.845688in}}%
\pgfpathquadraticcurveto{\pgfqpoint{7.893568in}{1.842493in}}{\pgfqpoint{7.886023in}{1.842493in}}%
\pgfpathquadraticcurveto{\pgfqpoint{7.873591in}{1.842493in}}{\pgfqpoint{7.867138in}{1.850224in}}%
\pgfpathquadraticcurveto{\pgfqpoint{7.860706in}{1.857955in}}{\pgfqpoint{7.860706in}{1.872840in}}%
\pgfpathlineto{\pgfqpoint{7.860706in}{1.872840in}}%
\pgfpathclose%
\pgfpathmoveto{\pgfqpoint{7.890538in}{1.918258in}}%
\pgfpathlineto{\pgfqpoint{7.890538in}{1.918258in}}%
\pgfpathlineto{\pgfqpoint{7.890538in}{1.918258in}}%
\pgfpathclose%
\pgfpathmoveto{\pgfqpoint{7.993021in}{1.905579in}}%
\pgfpathlineto{\pgfqpoint{7.993021in}{1.944626in}}%
\pgfpathlineto{\pgfqpoint{8.004876in}{1.944626in}}%
\pgfpathlineto{\pgfqpoint{8.004876in}{1.844369in}}%
\pgfpathlineto{\pgfqpoint{7.993021in}{1.844369in}}%
\pgfpathlineto{\pgfqpoint{7.993021in}{1.855192in}}%
\pgfpathquadraticcurveto{\pgfqpoint{7.989290in}{1.848760in}}{\pgfqpoint{7.983579in}{1.845626in}}%
\pgfpathquadraticcurveto{\pgfqpoint{7.977889in}{1.842493in}}{\pgfqpoint{7.969890in}{1.842493in}}%
\pgfpathquadraticcurveto{\pgfqpoint{7.956819in}{1.842493in}}{\pgfqpoint{7.948593in}{1.852925in}}%
\pgfpathquadraticcurveto{\pgfqpoint{7.940388in}{1.863377in}}{\pgfqpoint{7.940388in}{1.880386in}}%
\pgfpathquadraticcurveto{\pgfqpoint{7.940388in}{1.897394in}}{\pgfqpoint{7.948593in}{1.907826in}}%
\pgfpathquadraticcurveto{\pgfqpoint{7.956819in}{1.918258in}}{\pgfqpoint{7.969890in}{1.918258in}}%
\pgfpathquadraticcurveto{\pgfqpoint{7.977889in}{1.918258in}}{\pgfqpoint{7.983579in}{1.915124in}}%
\pgfpathquadraticcurveto{\pgfqpoint{7.989290in}{1.912011in}}{\pgfqpoint{7.993021in}{1.905579in}}%
\pgfpathlineto{\pgfqpoint{7.993021in}{1.905579in}}%
\pgfpathclose%
\pgfpathmoveto{\pgfqpoint{7.952634in}{1.880386in}}%
\pgfpathquadraticcurveto{\pgfqpoint{7.952634in}{1.867315in}}{\pgfqpoint{7.958015in}{1.859872in}}%
\pgfpathquadraticcurveto{\pgfqpoint{7.963396in}{1.852430in}}{\pgfqpoint{7.972797in}{1.852430in}}%
\pgfpathquadraticcurveto{\pgfqpoint{7.982198in}{1.852430in}}{\pgfqpoint{7.987599in}{1.859872in}}%
\pgfpathquadraticcurveto{\pgfqpoint{7.993021in}{1.867315in}}{\pgfqpoint{7.993021in}{1.880386in}}%
\pgfpathquadraticcurveto{\pgfqpoint{7.993021in}{1.893456in}}{\pgfqpoint{7.987599in}{1.900899in}}%
\pgfpathquadraticcurveto{\pgfqpoint{7.982198in}{1.908341in}}{\pgfqpoint{7.972797in}{1.908341in}}%
\pgfpathquadraticcurveto{\pgfqpoint{7.963396in}{1.908341in}}{\pgfqpoint{7.958015in}{1.900899in}}%
\pgfpathquadraticcurveto{\pgfqpoint{7.952634in}{1.893456in}}{\pgfqpoint{7.952634in}{1.880386in}}%
\pgfpathlineto{\pgfqpoint{7.952634in}{1.880386in}}%
\pgfpathclose%
\pgfpathmoveto{\pgfqpoint{8.059323in}{1.837668in}}%
\pgfpathquadraticcurveto{\pgfqpoint{8.054314in}{1.824783in}}{\pgfqpoint{8.049531in}{1.820866in}}%
\pgfpathquadraticcurveto{\pgfqpoint{8.044768in}{1.816928in}}{\pgfqpoint{8.036790in}{1.816928in}}%
\pgfpathlineto{\pgfqpoint{8.027306in}{1.816928in}}%
\pgfpathlineto{\pgfqpoint{8.027306in}{1.826845in}}%
\pgfpathlineto{\pgfqpoint{8.034275in}{1.826845in}}%
\pgfpathquadraticcurveto{\pgfqpoint{8.039161in}{1.826845in}}{\pgfqpoint{8.041861in}{1.829175in}}%
\pgfpathquadraticcurveto{\pgfqpoint{8.044583in}{1.831484in}}{\pgfqpoint{8.047861in}{1.840122in}}%
\pgfpathlineto{\pgfqpoint{8.049984in}{1.845523in}}%
\pgfpathlineto{\pgfqpoint{8.020812in}{1.916526in}}%
\pgfpathlineto{\pgfqpoint{8.033367in}{1.916526in}}%
\pgfpathlineto{\pgfqpoint{8.055922in}{1.860099in}}%
\pgfpathlineto{\pgfqpoint{8.078476in}{1.916526in}}%
\pgfpathlineto{\pgfqpoint{8.091031in}{1.916526in}}%
\pgfpathlineto{\pgfqpoint{8.059323in}{1.837668in}}%
\pgfpathlineto{\pgfqpoint{8.059323in}{1.837668in}}%
\pgfpathclose%
\pgfpathmoveto{\pgfqpoint{8.153272in}{1.855316in}}%
\pgfpathlineto{\pgfqpoint{8.174527in}{1.855316in}}%
\pgfpathlineto{\pgfqpoint{8.174527in}{1.928710in}}%
\pgfpathlineto{\pgfqpoint{8.151396in}{1.924071in}}%
\pgfpathlineto{\pgfqpoint{8.151396in}{1.935926in}}%
\pgfpathlineto{\pgfqpoint{8.174404in}{1.940565in}}%
\pgfpathlineto{\pgfqpoint{8.187413in}{1.940565in}}%
\pgfpathlineto{\pgfqpoint{8.187413in}{1.855316in}}%
\pgfpathlineto{\pgfqpoint{8.208668in}{1.855316in}}%
\pgfpathlineto{\pgfqpoint{8.208668in}{1.844369in}}%
\pgfpathlineto{\pgfqpoint{8.153272in}{1.844369in}}%
\pgfpathlineto{\pgfqpoint{8.153272in}{1.855316in}}%
\pgfpathlineto{\pgfqpoint{8.153272in}{1.855316in}}%
\pgfpathclose%
\pgfpathmoveto{\pgfqpoint{8.264414in}{1.897641in}}%
\pgfpathquadraticcurveto{\pgfqpoint{8.255653in}{1.897641in}}{\pgfqpoint{8.250519in}{1.891642in}}%
\pgfpathquadraticcurveto{\pgfqpoint{8.245406in}{1.885663in}}{\pgfqpoint{8.245406in}{1.875231in}}%
\pgfpathquadraticcurveto{\pgfqpoint{8.245406in}{1.864861in}}{\pgfqpoint{8.250519in}{1.858821in}}%
\pgfpathquadraticcurveto{\pgfqpoint{8.255653in}{1.852801in}}{\pgfqpoint{8.264414in}{1.852801in}}%
\pgfpathquadraticcurveto{\pgfqpoint{8.273176in}{1.852801in}}{\pgfqpoint{8.278289in}{1.858821in}}%
\pgfpathquadraticcurveto{\pgfqpoint{8.283402in}{1.864861in}}{\pgfqpoint{8.283402in}{1.875231in}}%
\pgfpathquadraticcurveto{\pgfqpoint{8.283402in}{1.885663in}}{\pgfqpoint{8.278289in}{1.891642in}}%
\pgfpathquadraticcurveto{\pgfqpoint{8.273176in}{1.897641in}}{\pgfqpoint{8.264414in}{1.897641in}}%
\pgfpathlineto{\pgfqpoint{8.264414in}{1.897641in}}%
\pgfpathclose%
\pgfpathmoveto{\pgfqpoint{8.290247in}{1.938441in}}%
\pgfpathlineto{\pgfqpoint{8.290247in}{1.926587in}}%
\pgfpathquadraticcurveto{\pgfqpoint{8.285340in}{1.928896in}}{\pgfqpoint{8.280351in}{1.930112in}}%
\pgfpathquadraticcurveto{\pgfqpoint{8.275362in}{1.931349in}}{\pgfqpoint{8.270455in}{1.931349in}}%
\pgfpathquadraticcurveto{\pgfqpoint{8.257570in}{1.931349in}}{\pgfqpoint{8.250766in}{1.922649in}}%
\pgfpathquadraticcurveto{\pgfqpoint{8.243984in}{1.913949in}}{\pgfqpoint{8.243015in}{1.896363in}}%
\pgfpathquadraticcurveto{\pgfqpoint{8.246808in}{1.901971in}}{\pgfqpoint{8.252539in}{1.904960in}}%
\pgfpathquadraticcurveto{\pgfqpoint{8.258291in}{1.907950in}}{\pgfqpoint{8.265177in}{1.907950in}}%
\pgfpathquadraticcurveto{\pgfqpoint{8.279671in}{1.907950in}}{\pgfqpoint{8.288082in}{1.899146in}}%
\pgfpathquadraticcurveto{\pgfqpoint{8.296493in}{1.890364in}}{\pgfqpoint{8.296493in}{1.875231in}}%
\pgfpathquadraticcurveto{\pgfqpoint{8.296493in}{1.860408in}}{\pgfqpoint{8.287732in}{1.851440in}}%
\pgfpathquadraticcurveto{\pgfqpoint{8.278970in}{1.842493in}}{\pgfqpoint{8.264414in}{1.842493in}}%
\pgfpathquadraticcurveto{\pgfqpoint{8.247715in}{1.842493in}}{\pgfqpoint{8.238891in}{1.855275in}}%
\pgfpathquadraticcurveto{\pgfqpoint{8.230068in}{1.868078in}}{\pgfqpoint{8.230068in}{1.892364in}}%
\pgfpathquadraticcurveto{\pgfqpoint{8.230068in}{1.915165in}}{\pgfqpoint{8.240891in}{1.928731in}}%
\pgfpathquadraticcurveto{\pgfqpoint{8.251715in}{1.942296in}}{\pgfqpoint{8.269940in}{1.942296in}}%
\pgfpathquadraticcurveto{\pgfqpoint{8.274846in}{1.942296in}}{\pgfqpoint{8.279836in}{1.941327in}}%
\pgfpathquadraticcurveto{\pgfqpoint{8.284825in}{1.940358in}}{\pgfqpoint{8.290247in}{1.938441in}}%
\pgfpathlineto{\pgfqpoint{8.290247in}{1.938441in}}%
\pgfpathclose%
\pgfpathmoveto{\pgfqpoint{8.387638in}{1.944482in}}%
\pgfpathquadraticcurveto{\pgfqpoint{8.379021in}{1.929679in}}{\pgfqpoint{8.374815in}{1.915165in}}%
\pgfpathquadraticcurveto{\pgfqpoint{8.370630in}{1.900672in}}{\pgfqpoint{8.370630in}{1.885787in}}%
\pgfpathquadraticcurveto{\pgfqpoint{8.370630in}{1.870923in}}{\pgfqpoint{8.374856in}{1.856326in}}%
\pgfpathquadraticcurveto{\pgfqpoint{8.379082in}{1.841730in}}{\pgfqpoint{8.387638in}{1.826969in}}%
\pgfpathlineto{\pgfqpoint{8.377330in}{1.826969in}}%
\pgfpathquadraticcurveto{\pgfqpoint{8.367682in}{1.842122in}}{\pgfqpoint{8.362878in}{1.856739in}}%
\pgfpathquadraticcurveto{\pgfqpoint{8.358074in}{1.871356in}}{\pgfqpoint{8.358074in}{1.885787in}}%
\pgfpathquadraticcurveto{\pgfqpoint{8.358074in}{1.900157in}}{\pgfqpoint{8.362837in}{1.914712in}}%
\pgfpathquadraticcurveto{\pgfqpoint{8.367620in}{1.929287in}}{\pgfqpoint{8.377330in}{1.944482in}}%
\pgfpathlineto{\pgfqpoint{8.387638in}{1.944482in}}%
\pgfpathlineto{\pgfqpoint{8.387638in}{1.944482in}}%
\pgfpathclose%
\pgfpathmoveto{\pgfqpoint{8.423531in}{1.855316in}}%
\pgfpathlineto{\pgfqpoint{8.468949in}{1.855316in}}%
\pgfpathlineto{\pgfqpoint{8.468949in}{1.844369in}}%
\pgfpathlineto{\pgfqpoint{8.407883in}{1.844369in}}%
\pgfpathlineto{\pgfqpoint{8.407883in}{1.855316in}}%
\pgfpathquadraticcurveto{\pgfqpoint{8.415285in}{1.862985in}}{\pgfqpoint{8.428067in}{1.875891in}}%
\pgfpathquadraticcurveto{\pgfqpoint{8.440870in}{1.888818in}}{\pgfqpoint{8.444148in}{1.892570in}}%
\pgfpathquadraticcurveto{\pgfqpoint{8.450394in}{1.899579in}}{\pgfqpoint{8.452868in}{1.904445in}}%
\pgfpathquadraticcurveto{\pgfqpoint{8.455363in}{1.909310in}}{\pgfqpoint{8.455363in}{1.914011in}}%
\pgfpathquadraticcurveto{\pgfqpoint{8.455363in}{1.921680in}}{\pgfqpoint{8.449982in}{1.926504in}}%
\pgfpathquadraticcurveto{\pgfqpoint{8.444601in}{1.931349in}}{\pgfqpoint{8.435963in}{1.931349in}}%
\pgfpathquadraticcurveto{\pgfqpoint{8.429840in}{1.931349in}}{\pgfqpoint{8.423036in}{1.929226in}}%
\pgfpathquadraticcurveto{\pgfqpoint{8.416254in}{1.927102in}}{\pgfqpoint{8.408523in}{1.922773in}}%
\pgfpathlineto{\pgfqpoint{8.408523in}{1.935926in}}%
\pgfpathquadraticcurveto{\pgfqpoint{8.416377in}{1.939080in}}{\pgfqpoint{8.423201in}{1.940688in}}%
\pgfpathquadraticcurveto{\pgfqpoint{8.430046in}{1.942296in}}{\pgfqpoint{8.435715in}{1.942296in}}%
\pgfpathquadraticcurveto{\pgfqpoint{8.450662in}{1.942296in}}{\pgfqpoint{8.459548in}{1.934813in}}%
\pgfpathquadraticcurveto{\pgfqpoint{8.468434in}{1.927349in}}{\pgfqpoint{8.468434in}{1.914856in}}%
\pgfpathquadraticcurveto{\pgfqpoint{8.468434in}{1.908918in}}{\pgfqpoint{8.466207in}{1.903599in}}%
\pgfpathquadraticcurveto{\pgfqpoint{8.464001in}{1.898301in}}{\pgfqpoint{8.458125in}{1.891085in}}%
\pgfpathquadraticcurveto{\pgfqpoint{8.456517in}{1.889209in}}{\pgfqpoint{8.447879in}{1.880282in}}%
\pgfpathquadraticcurveto{\pgfqpoint{8.439262in}{1.871356in}}{\pgfqpoint{8.423531in}{1.855316in}}%
\pgfpathlineto{\pgfqpoint{8.423531in}{1.855316in}}%
\pgfpathclose%
\pgfpathmoveto{\pgfqpoint{8.524098in}{1.931988in}}%
\pgfpathquadraticcurveto{\pgfqpoint{8.514058in}{1.931988in}}{\pgfqpoint{8.508986in}{1.922092in}}%
\pgfpathquadraticcurveto{\pgfqpoint{8.503935in}{1.912217in}}{\pgfqpoint{8.503935in}{1.892364in}}%
\pgfpathquadraticcurveto{\pgfqpoint{8.503935in}{1.872593in}}{\pgfqpoint{8.508986in}{1.862697in}}%
\pgfpathquadraticcurveto{\pgfqpoint{8.514058in}{1.852801in}}{\pgfqpoint{8.524098in}{1.852801in}}%
\pgfpathquadraticcurveto{\pgfqpoint{8.534220in}{1.852801in}}{\pgfqpoint{8.539271in}{1.862697in}}%
\pgfpathquadraticcurveto{\pgfqpoint{8.544343in}{1.872593in}}{\pgfqpoint{8.544343in}{1.892364in}}%
\pgfpathquadraticcurveto{\pgfqpoint{8.544343in}{1.912217in}}{\pgfqpoint{8.539271in}{1.922092in}}%
\pgfpathquadraticcurveto{\pgfqpoint{8.534220in}{1.931988in}}{\pgfqpoint{8.524098in}{1.931988in}}%
\pgfpathlineto{\pgfqpoint{8.524098in}{1.931988in}}%
\pgfpathclose%
\pgfpathmoveto{\pgfqpoint{8.524098in}{1.942296in}}%
\pgfpathquadraticcurveto{\pgfqpoint{8.540281in}{1.942296in}}{\pgfqpoint{8.548817in}{1.929494in}}%
\pgfpathquadraticcurveto{\pgfqpoint{8.557352in}{1.916711in}}{\pgfqpoint{8.557352in}{1.892364in}}%
\pgfpathquadraticcurveto{\pgfqpoint{8.557352in}{1.868078in}}{\pgfqpoint{8.548817in}{1.855275in}}%
\pgfpathquadraticcurveto{\pgfqpoint{8.540281in}{1.842493in}}{\pgfqpoint{8.524098in}{1.842493in}}%
\pgfpathquadraticcurveto{\pgfqpoint{8.507934in}{1.842493in}}{\pgfqpoint{8.499399in}{1.855275in}}%
\pgfpathquadraticcurveto{\pgfqpoint{8.490864in}{1.868078in}}{\pgfqpoint{8.490864in}{1.892364in}}%
\pgfpathquadraticcurveto{\pgfqpoint{8.490864in}{1.916711in}}{\pgfqpoint{8.499399in}{1.929494in}}%
\pgfpathquadraticcurveto{\pgfqpoint{8.507934in}{1.942296in}}{\pgfqpoint{8.524098in}{1.942296in}}%
\pgfpathlineto{\pgfqpoint{8.524098in}{1.942296in}}%
\pgfpathclose%
\pgfpathmoveto{\pgfqpoint{8.591431in}{1.855316in}}%
\pgfpathlineto{\pgfqpoint{8.636848in}{1.855316in}}%
\pgfpathlineto{\pgfqpoint{8.636848in}{1.844369in}}%
\pgfpathlineto{\pgfqpoint{8.575783in}{1.844369in}}%
\pgfpathlineto{\pgfqpoint{8.575783in}{1.855316in}}%
\pgfpathquadraticcurveto{\pgfqpoint{8.583184in}{1.862985in}}{\pgfqpoint{8.595966in}{1.875891in}}%
\pgfpathquadraticcurveto{\pgfqpoint{8.608769in}{1.888818in}}{\pgfqpoint{8.612047in}{1.892570in}}%
\pgfpathquadraticcurveto{\pgfqpoint{8.618294in}{1.899579in}}{\pgfqpoint{8.620768in}{1.904445in}}%
\pgfpathquadraticcurveto{\pgfqpoint{8.623262in}{1.909310in}}{\pgfqpoint{8.623262in}{1.914011in}}%
\pgfpathquadraticcurveto{\pgfqpoint{8.623262in}{1.921680in}}{\pgfqpoint{8.617881in}{1.926504in}}%
\pgfpathquadraticcurveto{\pgfqpoint{8.612500in}{1.931349in}}{\pgfqpoint{8.603862in}{1.931349in}}%
\pgfpathquadraticcurveto{\pgfqpoint{8.597739in}{1.931349in}}{\pgfqpoint{8.590936in}{1.929226in}}%
\pgfpathquadraticcurveto{\pgfqpoint{8.584153in}{1.927102in}}{\pgfqpoint{8.576422in}{1.922773in}}%
\pgfpathlineto{\pgfqpoint{8.576422in}{1.935926in}}%
\pgfpathquadraticcurveto{\pgfqpoint{8.584277in}{1.939080in}}{\pgfqpoint{8.591101in}{1.940688in}}%
\pgfpathquadraticcurveto{\pgfqpoint{8.597945in}{1.942296in}}{\pgfqpoint{8.603615in}{1.942296in}}%
\pgfpathquadraticcurveto{\pgfqpoint{8.618562in}{1.942296in}}{\pgfqpoint{8.627447in}{1.934813in}}%
\pgfpathquadraticcurveto{\pgfqpoint{8.636333in}{1.927349in}}{\pgfqpoint{8.636333in}{1.914856in}}%
\pgfpathquadraticcurveto{\pgfqpoint{8.636333in}{1.908918in}}{\pgfqpoint{8.634106in}{1.903599in}}%
\pgfpathquadraticcurveto{\pgfqpoint{8.631900in}{1.898301in}}{\pgfqpoint{8.626025in}{1.891085in}}%
\pgfpathquadraticcurveto{\pgfqpoint{8.624417in}{1.889209in}}{\pgfqpoint{8.615778in}{1.880282in}}%
\pgfpathquadraticcurveto{\pgfqpoint{8.607161in}{1.871356in}}{\pgfqpoint{8.591431in}{1.855316in}}%
\pgfpathlineto{\pgfqpoint{8.591431in}{1.855316in}}%
\pgfpathclose%
\pgfpathmoveto{\pgfqpoint{8.703604in}{1.896239in}}%
\pgfpathquadraticcurveto{\pgfqpoint{8.712943in}{1.894240in}}{\pgfqpoint{8.718180in}{1.887910in}}%
\pgfpathquadraticcurveto{\pgfqpoint{8.723437in}{1.881602in}}{\pgfqpoint{8.723437in}{1.872325in}}%
\pgfpathquadraticcurveto{\pgfqpoint{8.723437in}{1.858099in}}{\pgfqpoint{8.713644in}{1.850286in}}%
\pgfpathquadraticcurveto{\pgfqpoint{8.703851in}{1.842493in}}{\pgfqpoint{8.685812in}{1.842493in}}%
\pgfpathquadraticcurveto{\pgfqpoint{8.679771in}{1.842493in}}{\pgfqpoint{8.673360in}{1.843688in}}%
\pgfpathquadraticcurveto{\pgfqpoint{8.666948in}{1.844884in}}{\pgfqpoint{8.660124in}{1.847276in}}%
\pgfpathlineto{\pgfqpoint{8.660124in}{1.859831in}}%
\pgfpathquadraticcurveto{\pgfqpoint{8.665526in}{1.856677in}}{\pgfqpoint{8.671958in}{1.855069in}}%
\pgfpathquadraticcurveto{\pgfqpoint{8.678411in}{1.853461in}}{\pgfqpoint{8.685441in}{1.853461in}}%
\pgfpathquadraticcurveto{\pgfqpoint{8.697666in}{1.853461in}}{\pgfqpoint{8.704078in}{1.858285in}}%
\pgfpathquadraticcurveto{\pgfqpoint{8.710490in}{1.863109in}}{\pgfqpoint{8.710490in}{1.872325in}}%
\pgfpathquadraticcurveto{\pgfqpoint{8.710490in}{1.880839in}}{\pgfqpoint{8.704532in}{1.885622in}}%
\pgfpathquadraticcurveto{\pgfqpoint{8.698574in}{1.890426in}}{\pgfqpoint{8.687956in}{1.890426in}}%
\pgfpathlineto{\pgfqpoint{8.676741in}{1.890426in}}%
\pgfpathlineto{\pgfqpoint{8.676741in}{1.901126in}}%
\pgfpathlineto{\pgfqpoint{8.688472in}{1.901126in}}%
\pgfpathquadraticcurveto{\pgfqpoint{8.698058in}{1.901126in}}{\pgfqpoint{8.703150in}{1.904960in}}%
\pgfpathquadraticcurveto{\pgfqpoint{8.708243in}{1.908795in}}{\pgfqpoint{8.708243in}{1.916011in}}%
\pgfpathquadraticcurveto{\pgfqpoint{8.708243in}{1.923412in}}{\pgfqpoint{8.702985in}{1.927370in}}%
\pgfpathquadraticcurveto{\pgfqpoint{8.697749in}{1.931349in}}{\pgfqpoint{8.687956in}{1.931349in}}%
\pgfpathquadraticcurveto{\pgfqpoint{8.682596in}{1.931349in}}{\pgfqpoint{8.676473in}{1.930174in}}%
\pgfpathquadraticcurveto{\pgfqpoint{8.670350in}{1.929019in}}{\pgfqpoint{8.663010in}{1.926587in}}%
\pgfpathlineto{\pgfqpoint{8.663010in}{1.938173in}}%
\pgfpathquadraticcurveto{\pgfqpoint{8.670432in}{1.940235in}}{\pgfqpoint{8.676906in}{1.941265in}}%
\pgfpathquadraticcurveto{\pgfqpoint{8.683379in}{1.942296in}}{\pgfqpoint{8.689111in}{1.942296in}}%
\pgfpathquadraticcurveto{\pgfqpoint{8.703934in}{1.942296in}}{\pgfqpoint{8.712551in}{1.935555in}}%
\pgfpathquadraticcurveto{\pgfqpoint{8.721190in}{1.928834in}}{\pgfqpoint{8.721190in}{1.917371in}}%
\pgfpathquadraticcurveto{\pgfqpoint{8.721190in}{1.909372in}}{\pgfqpoint{8.716613in}{1.903867in}}%
\pgfpathquadraticcurveto{\pgfqpoint{8.712036in}{1.898363in}}{\pgfqpoint{8.703604in}{1.896239in}}%
\pgfpathlineto{\pgfqpoint{8.703604in}{1.896239in}}%
\pgfpathclose%
\pgfpathmoveto{\pgfqpoint{8.744589in}{1.944482in}}%
\pgfpathlineto{\pgfqpoint{8.754897in}{1.944482in}}%
\pgfpathquadraticcurveto{\pgfqpoint{8.764546in}{1.929287in}}{\pgfqpoint{8.769349in}{1.914712in}}%
\pgfpathquadraticcurveto{\pgfqpoint{8.774153in}{1.900157in}}{\pgfqpoint{8.774153in}{1.885787in}}%
\pgfpathquadraticcurveto{\pgfqpoint{8.774153in}{1.871356in}}{\pgfqpoint{8.769349in}{1.856739in}}%
\pgfpathquadraticcurveto{\pgfqpoint{8.764546in}{1.842122in}}{\pgfqpoint{8.754897in}{1.826969in}}%
\pgfpathlineto{\pgfqpoint{8.744589in}{1.826969in}}%
\pgfpathquadraticcurveto{\pgfqpoint{8.753145in}{1.841730in}}{\pgfqpoint{8.757371in}{1.856326in}}%
\pgfpathquadraticcurveto{\pgfqpoint{8.761598in}{1.870923in}}{\pgfqpoint{8.761598in}{1.885787in}}%
\pgfpathquadraticcurveto{\pgfqpoint{8.761598in}{1.900672in}}{\pgfqpoint{8.757371in}{1.915165in}}%
\pgfpathquadraticcurveto{\pgfqpoint{8.753145in}{1.929679in}}{\pgfqpoint{8.744589in}{1.944482in}}%
\pgfpathlineto{\pgfqpoint{8.744589in}{1.944482in}}%
\pgfpathclose%
\pgfusepath{fill}%
\end{pgfscope}%
\begin{pgfscope}%
\pgfsetbuttcap%
\pgfsetroundjoin%
\definecolor{currentfill}{rgb}{0.090196,0.168627,0.227451}%
\pgfsetfillcolor{currentfill}%
\pgfsetlinewidth{0.000000pt}%
\definecolor{currentstroke}{rgb}{0.090196,0.168627,0.227451}%
\pgfsetstrokecolor{currentstroke}%
\pgfsetdash{}{0pt}%
\pgfpathmoveto{\pgfqpoint{7.784611in}{1.760770in}}%
\pgfpathlineto{\pgfqpoint{7.784611in}{1.748071in}}%
\pgfpathquadraticcurveto{\pgfqpoint{7.777210in}{1.751617in}}{\pgfqpoint{7.770633in}{1.753348in}}%
\pgfpathquadraticcurveto{\pgfqpoint{7.764056in}{1.755101in}}{\pgfqpoint{7.757933in}{1.755101in}}%
\pgfpathquadraticcurveto{\pgfqpoint{7.747316in}{1.755101in}}{\pgfqpoint{7.741543in}{1.750977in}}%
\pgfpathquadraticcurveto{\pgfqpoint{7.735771in}{1.746854in}}{\pgfqpoint{7.735771in}{1.739247in}}%
\pgfpathquadraticcurveto{\pgfqpoint{7.735771in}{1.732856in}}{\pgfqpoint{7.739605in}{1.729598in}}%
\pgfpathquadraticcurveto{\pgfqpoint{7.743440in}{1.726362in}}{\pgfqpoint{7.754140in}{1.724362in}}%
\pgfpathlineto{\pgfqpoint{7.761995in}{1.722754in}}%
\pgfpathquadraticcurveto{\pgfqpoint{7.776550in}{1.719971in}}{\pgfqpoint{7.783477in}{1.712982in}}%
\pgfpathquadraticcurveto{\pgfqpoint{7.790404in}{1.705993in}}{\pgfqpoint{7.790404in}{1.694283in}}%
\pgfpathquadraticcurveto{\pgfqpoint{7.790404in}{1.680284in}}{\pgfqpoint{7.781024in}{1.673068in}}%
\pgfpathquadraticcurveto{\pgfqpoint{7.771664in}{1.665853in}}{\pgfqpoint{7.753563in}{1.665853in}}%
\pgfpathquadraticcurveto{\pgfqpoint{7.746739in}{1.665853in}}{\pgfqpoint{7.739028in}{1.667399in}}%
\pgfpathquadraticcurveto{\pgfqpoint{7.731338in}{1.668945in}}{\pgfqpoint{7.723092in}{1.671976in}}%
\pgfpathlineto{\pgfqpoint{7.723092in}{1.685376in}}%
\pgfpathquadraticcurveto{\pgfqpoint{7.731008in}{1.680944in}}{\pgfqpoint{7.738616in}{1.678676in}}%
\pgfpathquadraticcurveto{\pgfqpoint{7.746223in}{1.676429in}}{\pgfqpoint{7.753563in}{1.676429in}}%
\pgfpathquadraticcurveto{\pgfqpoint{7.764696in}{1.676429in}}{\pgfqpoint{7.770757in}{1.680800in}}%
\pgfpathquadraticcurveto{\pgfqpoint{7.776818in}{1.685191in}}{\pgfqpoint{7.776818in}{1.693314in}}%
\pgfpathquadraticcurveto{\pgfqpoint{7.776818in}{1.700385in}}{\pgfqpoint{7.772468in}{1.704385in}}%
\pgfpathquadraticcurveto{\pgfqpoint{7.768118in}{1.708384in}}{\pgfqpoint{7.758201in}{1.710384in}}%
\pgfpathlineto{\pgfqpoint{7.750264in}{1.711930in}}%
\pgfpathquadraticcurveto{\pgfqpoint{7.735709in}{1.714816in}}{\pgfqpoint{7.729194in}{1.721001in}}%
\pgfpathquadraticcurveto{\pgfqpoint{7.722700in}{1.727186in}}{\pgfqpoint{7.722700in}{1.738216in}}%
\pgfpathquadraticcurveto{\pgfqpoint{7.722700in}{1.750977in}}{\pgfqpoint{7.731689in}{1.758317in}}%
\pgfpathquadraticcurveto{\pgfqpoint{7.740678in}{1.765656in}}{\pgfqpoint{7.756449in}{1.765656in}}%
\pgfpathquadraticcurveto{\pgfqpoint{7.763232in}{1.765656in}}{\pgfqpoint{7.770241in}{1.764419in}}%
\pgfpathquadraticcurveto{\pgfqpoint{7.777271in}{1.763203in}}{\pgfqpoint{7.784611in}{1.760770in}}%
\pgfpathlineto{\pgfqpoint{7.784611in}{1.760770in}}%
\pgfpathclose%
\pgfpathmoveto{\pgfqpoint{7.821926in}{1.760379in}}%
\pgfpathlineto{\pgfqpoint{7.821926in}{1.739886in}}%
\pgfpathlineto{\pgfqpoint{7.846336in}{1.739886in}}%
\pgfpathlineto{\pgfqpoint{7.846336in}{1.730670in}}%
\pgfpathlineto{\pgfqpoint{7.821926in}{1.730670in}}%
\pgfpathlineto{\pgfqpoint{7.821926in}{1.691499in}}%
\pgfpathquadraticcurveto{\pgfqpoint{7.821926in}{1.682676in}}{\pgfqpoint{7.824339in}{1.680160in}}%
\pgfpathquadraticcurveto{\pgfqpoint{7.826751in}{1.677645in}}{\pgfqpoint{7.834173in}{1.677645in}}%
\pgfpathlineto{\pgfqpoint{7.846336in}{1.677645in}}%
\pgfpathlineto{\pgfqpoint{7.846336in}{1.667729in}}%
\pgfpathlineto{\pgfqpoint{7.834173in}{1.667729in}}%
\pgfpathquadraticcurveto{\pgfqpoint{7.820442in}{1.667729in}}{\pgfqpoint{7.815226in}{1.672842in}}%
\pgfpathquadraticcurveto{\pgfqpoint{7.810010in}{1.677975in}}{\pgfqpoint{7.810010in}{1.691499in}}%
\pgfpathlineto{\pgfqpoint{7.810010in}{1.730670in}}%
\pgfpathlineto{\pgfqpoint{7.801310in}{1.730670in}}%
\pgfpathlineto{\pgfqpoint{7.801310in}{1.739886in}}%
\pgfpathlineto{\pgfqpoint{7.810010in}{1.739886in}}%
\pgfpathlineto{\pgfqpoint{7.810010in}{1.760379in}}%
\pgfpathlineto{\pgfqpoint{7.821926in}{1.760379in}}%
\pgfpathlineto{\pgfqpoint{7.821926in}{1.760379in}}%
\pgfpathclose%
\pgfpathmoveto{\pgfqpoint{7.860706in}{1.696200in}}%
\pgfpathlineto{\pgfqpoint{7.860706in}{1.739886in}}%
\pgfpathlineto{\pgfqpoint{7.872560in}{1.739886in}}%
\pgfpathlineto{\pgfqpoint{7.872560in}{1.696653in}}%
\pgfpathquadraticcurveto{\pgfqpoint{7.872560in}{1.686407in}}{\pgfqpoint{7.876539in}{1.681274in}}%
\pgfpathquadraticcurveto{\pgfqpoint{7.880539in}{1.676161in}}{\pgfqpoint{7.888538in}{1.676161in}}%
\pgfpathquadraticcurveto{\pgfqpoint{7.898124in}{1.676161in}}{\pgfqpoint{7.903691in}{1.682284in}}%
\pgfpathquadraticcurveto{\pgfqpoint{7.909278in}{1.688407in}}{\pgfqpoint{7.909278in}{1.698983in}}%
\pgfpathlineto{\pgfqpoint{7.909278in}{1.739886in}}%
\pgfpathlineto{\pgfqpoint{7.921132in}{1.739886in}}%
\pgfpathlineto{\pgfqpoint{7.921132in}{1.667729in}}%
\pgfpathlineto{\pgfqpoint{7.909278in}{1.667729in}}%
\pgfpathlineto{\pgfqpoint{7.909278in}{1.678820in}}%
\pgfpathquadraticcurveto{\pgfqpoint{7.904969in}{1.672244in}}{\pgfqpoint{7.899258in}{1.669048in}}%
\pgfpathquadraticcurveto{\pgfqpoint{7.893568in}{1.665853in}}{\pgfqpoint{7.886023in}{1.665853in}}%
\pgfpathquadraticcurveto{\pgfqpoint{7.873591in}{1.665853in}}{\pgfqpoint{7.867138in}{1.673584in}}%
\pgfpathquadraticcurveto{\pgfqpoint{7.860706in}{1.681315in}}{\pgfqpoint{7.860706in}{1.696200in}}%
\pgfpathlineto{\pgfqpoint{7.860706in}{1.696200in}}%
\pgfpathclose%
\pgfpathmoveto{\pgfqpoint{7.890538in}{1.741618in}}%
\pgfpathlineto{\pgfqpoint{7.890538in}{1.741618in}}%
\pgfpathlineto{\pgfqpoint{7.890538in}{1.741618in}}%
\pgfpathclose%
\pgfpathmoveto{\pgfqpoint{7.993021in}{1.728939in}}%
\pgfpathlineto{\pgfqpoint{7.993021in}{1.767986in}}%
\pgfpathlineto{\pgfqpoint{8.004876in}{1.767986in}}%
\pgfpathlineto{\pgfqpoint{8.004876in}{1.667729in}}%
\pgfpathlineto{\pgfqpoint{7.993021in}{1.667729in}}%
\pgfpathlineto{\pgfqpoint{7.993021in}{1.678552in}}%
\pgfpathquadraticcurveto{\pgfqpoint{7.989290in}{1.672120in}}{\pgfqpoint{7.983579in}{1.668986in}}%
\pgfpathquadraticcurveto{\pgfqpoint{7.977889in}{1.665853in}}{\pgfqpoint{7.969890in}{1.665853in}}%
\pgfpathquadraticcurveto{\pgfqpoint{7.956819in}{1.665853in}}{\pgfqpoint{7.948593in}{1.676285in}}%
\pgfpathquadraticcurveto{\pgfqpoint{7.940388in}{1.686737in}}{\pgfqpoint{7.940388in}{1.703746in}}%
\pgfpathquadraticcurveto{\pgfqpoint{7.940388in}{1.720754in}}{\pgfqpoint{7.948593in}{1.731186in}}%
\pgfpathquadraticcurveto{\pgfqpoint{7.956819in}{1.741618in}}{\pgfqpoint{7.969890in}{1.741618in}}%
\pgfpathquadraticcurveto{\pgfqpoint{7.977889in}{1.741618in}}{\pgfqpoint{7.983579in}{1.738484in}}%
\pgfpathquadraticcurveto{\pgfqpoint{7.989290in}{1.735371in}}{\pgfqpoint{7.993021in}{1.728939in}}%
\pgfpathlineto{\pgfqpoint{7.993021in}{1.728939in}}%
\pgfpathclose%
\pgfpathmoveto{\pgfqpoint{7.952634in}{1.703746in}}%
\pgfpathquadraticcurveto{\pgfqpoint{7.952634in}{1.690675in}}{\pgfqpoint{7.958015in}{1.683232in}}%
\pgfpathquadraticcurveto{\pgfqpoint{7.963396in}{1.675790in}}{\pgfqpoint{7.972797in}{1.675790in}}%
\pgfpathquadraticcurveto{\pgfqpoint{7.982198in}{1.675790in}}{\pgfqpoint{7.987599in}{1.683232in}}%
\pgfpathquadraticcurveto{\pgfqpoint{7.993021in}{1.690675in}}{\pgfqpoint{7.993021in}{1.703746in}}%
\pgfpathquadraticcurveto{\pgfqpoint{7.993021in}{1.716816in}}{\pgfqpoint{7.987599in}{1.724259in}}%
\pgfpathquadraticcurveto{\pgfqpoint{7.982198in}{1.731701in}}{\pgfqpoint{7.972797in}{1.731701in}}%
\pgfpathquadraticcurveto{\pgfqpoint{7.963396in}{1.731701in}}{\pgfqpoint{7.958015in}{1.724259in}}%
\pgfpathquadraticcurveto{\pgfqpoint{7.952634in}{1.716816in}}{\pgfqpoint{7.952634in}{1.703746in}}%
\pgfpathlineto{\pgfqpoint{7.952634in}{1.703746in}}%
\pgfpathclose%
\pgfpathmoveto{\pgfqpoint{8.059323in}{1.661028in}}%
\pgfpathquadraticcurveto{\pgfqpoint{8.054314in}{1.648143in}}{\pgfqpoint{8.049531in}{1.644226in}}%
\pgfpathquadraticcurveto{\pgfqpoint{8.044768in}{1.640288in}}{\pgfqpoint{8.036790in}{1.640288in}}%
\pgfpathlineto{\pgfqpoint{8.027306in}{1.640288in}}%
\pgfpathlineto{\pgfqpoint{8.027306in}{1.650205in}}%
\pgfpathlineto{\pgfqpoint{8.034275in}{1.650205in}}%
\pgfpathquadraticcurveto{\pgfqpoint{8.039161in}{1.650205in}}{\pgfqpoint{8.041861in}{1.652535in}}%
\pgfpathquadraticcurveto{\pgfqpoint{8.044583in}{1.654844in}}{\pgfqpoint{8.047861in}{1.663482in}}%
\pgfpathlineto{\pgfqpoint{8.049984in}{1.668883in}}%
\pgfpathlineto{\pgfqpoint{8.020812in}{1.739886in}}%
\pgfpathlineto{\pgfqpoint{8.033367in}{1.739886in}}%
\pgfpathlineto{\pgfqpoint{8.055922in}{1.683459in}}%
\pgfpathlineto{\pgfqpoint{8.078476in}{1.739886in}}%
\pgfpathlineto{\pgfqpoint{8.091031in}{1.739886in}}%
\pgfpathlineto{\pgfqpoint{8.059323in}{1.661028in}}%
\pgfpathlineto{\pgfqpoint{8.059323in}{1.661028in}}%
\pgfpathclose%
\pgfpathmoveto{\pgfqpoint{8.153272in}{1.678676in}}%
\pgfpathlineto{\pgfqpoint{8.174527in}{1.678676in}}%
\pgfpathlineto{\pgfqpoint{8.174527in}{1.752070in}}%
\pgfpathlineto{\pgfqpoint{8.151396in}{1.747431in}}%
\pgfpathlineto{\pgfqpoint{8.151396in}{1.759286in}}%
\pgfpathlineto{\pgfqpoint{8.174404in}{1.763925in}}%
\pgfpathlineto{\pgfqpoint{8.187413in}{1.763925in}}%
\pgfpathlineto{\pgfqpoint{8.187413in}{1.678676in}}%
\pgfpathlineto{\pgfqpoint{8.208668in}{1.678676in}}%
\pgfpathlineto{\pgfqpoint{8.208668in}{1.667729in}}%
\pgfpathlineto{\pgfqpoint{8.153272in}{1.667729in}}%
\pgfpathlineto{\pgfqpoint{8.153272in}{1.678676in}}%
\pgfpathlineto{\pgfqpoint{8.153272in}{1.678676in}}%
\pgfpathclose%
\pgfpathmoveto{\pgfqpoint{8.231676in}{1.763925in}}%
\pgfpathlineto{\pgfqpoint{8.293525in}{1.763925in}}%
\pgfpathlineto{\pgfqpoint{8.293525in}{1.758379in}}%
\pgfpathlineto{\pgfqpoint{8.258601in}{1.667729in}}%
\pgfpathlineto{\pgfqpoint{8.245015in}{1.667729in}}%
\pgfpathlineto{\pgfqpoint{8.277877in}{1.752957in}}%
\pgfpathlineto{\pgfqpoint{8.231676in}{1.752957in}}%
\pgfpathlineto{\pgfqpoint{8.231676in}{1.763925in}}%
\pgfpathlineto{\pgfqpoint{8.231676in}{1.763925in}}%
\pgfpathclose%
\pgfpathmoveto{\pgfqpoint{8.387638in}{1.767842in}}%
\pgfpathquadraticcurveto{\pgfqpoint{8.379021in}{1.753039in}}{\pgfqpoint{8.374815in}{1.738525in}}%
\pgfpathquadraticcurveto{\pgfqpoint{8.370630in}{1.724032in}}{\pgfqpoint{8.370630in}{1.709147in}}%
\pgfpathquadraticcurveto{\pgfqpoint{8.370630in}{1.694283in}}{\pgfqpoint{8.374856in}{1.679686in}}%
\pgfpathquadraticcurveto{\pgfqpoint{8.379082in}{1.665090in}}{\pgfqpoint{8.387638in}{1.650329in}}%
\pgfpathlineto{\pgfqpoint{8.377330in}{1.650329in}}%
\pgfpathquadraticcurveto{\pgfqpoint{8.367682in}{1.665482in}}{\pgfqpoint{8.362878in}{1.680099in}}%
\pgfpathquadraticcurveto{\pgfqpoint{8.358074in}{1.694716in}}{\pgfqpoint{8.358074in}{1.709147in}}%
\pgfpathquadraticcurveto{\pgfqpoint{8.358074in}{1.723517in}}{\pgfqpoint{8.362837in}{1.738072in}}%
\pgfpathquadraticcurveto{\pgfqpoint{8.367620in}{1.752647in}}{\pgfqpoint{8.377330in}{1.767842in}}%
\pgfpathlineto{\pgfqpoint{8.387638in}{1.767842in}}%
\pgfpathlineto{\pgfqpoint{8.387638in}{1.767842in}}%
\pgfpathclose%
\pgfpathmoveto{\pgfqpoint{8.423531in}{1.678676in}}%
\pgfpathlineto{\pgfqpoint{8.468949in}{1.678676in}}%
\pgfpathlineto{\pgfqpoint{8.468949in}{1.667729in}}%
\pgfpathlineto{\pgfqpoint{8.407883in}{1.667729in}}%
\pgfpathlineto{\pgfqpoint{8.407883in}{1.678676in}}%
\pgfpathquadraticcurveto{\pgfqpoint{8.415285in}{1.686345in}}{\pgfqpoint{8.428067in}{1.699251in}}%
\pgfpathquadraticcurveto{\pgfqpoint{8.440870in}{1.712178in}}{\pgfqpoint{8.444148in}{1.715930in}}%
\pgfpathquadraticcurveto{\pgfqpoint{8.450394in}{1.722939in}}{\pgfqpoint{8.452868in}{1.727805in}}%
\pgfpathquadraticcurveto{\pgfqpoint{8.455363in}{1.732670in}}{\pgfqpoint{8.455363in}{1.737371in}}%
\pgfpathquadraticcurveto{\pgfqpoint{8.455363in}{1.745040in}}{\pgfqpoint{8.449982in}{1.749864in}}%
\pgfpathquadraticcurveto{\pgfqpoint{8.444601in}{1.754709in}}{\pgfqpoint{8.435963in}{1.754709in}}%
\pgfpathquadraticcurveto{\pgfqpoint{8.429840in}{1.754709in}}{\pgfqpoint{8.423036in}{1.752586in}}%
\pgfpathquadraticcurveto{\pgfqpoint{8.416254in}{1.750462in}}{\pgfqpoint{8.408523in}{1.746133in}}%
\pgfpathlineto{\pgfqpoint{8.408523in}{1.759286in}}%
\pgfpathquadraticcurveto{\pgfqpoint{8.416377in}{1.762440in}}{\pgfqpoint{8.423201in}{1.764048in}}%
\pgfpathquadraticcurveto{\pgfqpoint{8.430046in}{1.765656in}}{\pgfqpoint{8.435715in}{1.765656in}}%
\pgfpathquadraticcurveto{\pgfqpoint{8.450662in}{1.765656in}}{\pgfqpoint{8.459548in}{1.758173in}}%
\pgfpathquadraticcurveto{\pgfqpoint{8.468434in}{1.750709in}}{\pgfqpoint{8.468434in}{1.738216in}}%
\pgfpathquadraticcurveto{\pgfqpoint{8.468434in}{1.732278in}}{\pgfqpoint{8.466207in}{1.726959in}}%
\pgfpathquadraticcurveto{\pgfqpoint{8.464001in}{1.721661in}}{\pgfqpoint{8.458125in}{1.714445in}}%
\pgfpathquadraticcurveto{\pgfqpoint{8.456517in}{1.712569in}}{\pgfqpoint{8.447879in}{1.703642in}}%
\pgfpathquadraticcurveto{\pgfqpoint{8.439262in}{1.694716in}}{\pgfqpoint{8.423531in}{1.678676in}}%
\pgfpathlineto{\pgfqpoint{8.423531in}{1.678676in}}%
\pgfpathclose%
\pgfpathmoveto{\pgfqpoint{8.524098in}{1.755348in}}%
\pgfpathquadraticcurveto{\pgfqpoint{8.514058in}{1.755348in}}{\pgfqpoint{8.508986in}{1.745452in}}%
\pgfpathquadraticcurveto{\pgfqpoint{8.503935in}{1.735577in}}{\pgfqpoint{8.503935in}{1.715724in}}%
\pgfpathquadraticcurveto{\pgfqpoint{8.503935in}{1.695953in}}{\pgfqpoint{8.508986in}{1.686057in}}%
\pgfpathquadraticcurveto{\pgfqpoint{8.514058in}{1.676161in}}{\pgfqpoint{8.524098in}{1.676161in}}%
\pgfpathquadraticcurveto{\pgfqpoint{8.534220in}{1.676161in}}{\pgfqpoint{8.539271in}{1.686057in}}%
\pgfpathquadraticcurveto{\pgfqpoint{8.544343in}{1.695953in}}{\pgfqpoint{8.544343in}{1.715724in}}%
\pgfpathquadraticcurveto{\pgfqpoint{8.544343in}{1.735577in}}{\pgfqpoint{8.539271in}{1.745452in}}%
\pgfpathquadraticcurveto{\pgfqpoint{8.534220in}{1.755348in}}{\pgfqpoint{8.524098in}{1.755348in}}%
\pgfpathlineto{\pgfqpoint{8.524098in}{1.755348in}}%
\pgfpathclose%
\pgfpathmoveto{\pgfqpoint{8.524098in}{1.765656in}}%
\pgfpathquadraticcurveto{\pgfqpoint{8.540281in}{1.765656in}}{\pgfqpoint{8.548817in}{1.752854in}}%
\pgfpathquadraticcurveto{\pgfqpoint{8.557352in}{1.740071in}}{\pgfqpoint{8.557352in}{1.715724in}}%
\pgfpathquadraticcurveto{\pgfqpoint{8.557352in}{1.691438in}}{\pgfqpoint{8.548817in}{1.678635in}}%
\pgfpathquadraticcurveto{\pgfqpoint{8.540281in}{1.665853in}}{\pgfqpoint{8.524098in}{1.665853in}}%
\pgfpathquadraticcurveto{\pgfqpoint{8.507934in}{1.665853in}}{\pgfqpoint{8.499399in}{1.678635in}}%
\pgfpathquadraticcurveto{\pgfqpoint{8.490864in}{1.691438in}}{\pgfqpoint{8.490864in}{1.715724in}}%
\pgfpathquadraticcurveto{\pgfqpoint{8.490864in}{1.740071in}}{\pgfqpoint{8.499399in}{1.752854in}}%
\pgfpathquadraticcurveto{\pgfqpoint{8.507934in}{1.765656in}}{\pgfqpoint{8.524098in}{1.765656in}}%
\pgfpathlineto{\pgfqpoint{8.524098in}{1.765656in}}%
\pgfpathclose%
\pgfpathmoveto{\pgfqpoint{8.591431in}{1.678676in}}%
\pgfpathlineto{\pgfqpoint{8.636848in}{1.678676in}}%
\pgfpathlineto{\pgfqpoint{8.636848in}{1.667729in}}%
\pgfpathlineto{\pgfqpoint{8.575783in}{1.667729in}}%
\pgfpathlineto{\pgfqpoint{8.575783in}{1.678676in}}%
\pgfpathquadraticcurveto{\pgfqpoint{8.583184in}{1.686345in}}{\pgfqpoint{8.595966in}{1.699251in}}%
\pgfpathquadraticcurveto{\pgfqpoint{8.608769in}{1.712178in}}{\pgfqpoint{8.612047in}{1.715930in}}%
\pgfpathquadraticcurveto{\pgfqpoint{8.618294in}{1.722939in}}{\pgfqpoint{8.620768in}{1.727805in}}%
\pgfpathquadraticcurveto{\pgfqpoint{8.623262in}{1.732670in}}{\pgfqpoint{8.623262in}{1.737371in}}%
\pgfpathquadraticcurveto{\pgfqpoint{8.623262in}{1.745040in}}{\pgfqpoint{8.617881in}{1.749864in}}%
\pgfpathquadraticcurveto{\pgfqpoint{8.612500in}{1.754709in}}{\pgfqpoint{8.603862in}{1.754709in}}%
\pgfpathquadraticcurveto{\pgfqpoint{8.597739in}{1.754709in}}{\pgfqpoint{8.590936in}{1.752586in}}%
\pgfpathquadraticcurveto{\pgfqpoint{8.584153in}{1.750462in}}{\pgfqpoint{8.576422in}{1.746133in}}%
\pgfpathlineto{\pgfqpoint{8.576422in}{1.759286in}}%
\pgfpathquadraticcurveto{\pgfqpoint{8.584277in}{1.762440in}}{\pgfqpoint{8.591101in}{1.764048in}}%
\pgfpathquadraticcurveto{\pgfqpoint{8.597945in}{1.765656in}}{\pgfqpoint{8.603615in}{1.765656in}}%
\pgfpathquadraticcurveto{\pgfqpoint{8.618562in}{1.765656in}}{\pgfqpoint{8.627447in}{1.758173in}}%
\pgfpathquadraticcurveto{\pgfqpoint{8.636333in}{1.750709in}}{\pgfqpoint{8.636333in}{1.738216in}}%
\pgfpathquadraticcurveto{\pgfqpoint{8.636333in}{1.732278in}}{\pgfqpoint{8.634106in}{1.726959in}}%
\pgfpathquadraticcurveto{\pgfqpoint{8.631900in}{1.721661in}}{\pgfqpoint{8.626025in}{1.714445in}}%
\pgfpathquadraticcurveto{\pgfqpoint{8.624417in}{1.712569in}}{\pgfqpoint{8.615778in}{1.703642in}}%
\pgfpathquadraticcurveto{\pgfqpoint{8.607161in}{1.694716in}}{\pgfqpoint{8.591431in}{1.678676in}}%
\pgfpathlineto{\pgfqpoint{8.591431in}{1.678676in}}%
\pgfpathclose%
\pgfpathmoveto{\pgfqpoint{8.703604in}{1.719599in}}%
\pgfpathquadraticcurveto{\pgfqpoint{8.712943in}{1.717600in}}{\pgfqpoint{8.718180in}{1.711270in}}%
\pgfpathquadraticcurveto{\pgfqpoint{8.723437in}{1.704962in}}{\pgfqpoint{8.723437in}{1.695685in}}%
\pgfpathquadraticcurveto{\pgfqpoint{8.723437in}{1.681459in}}{\pgfqpoint{8.713644in}{1.673646in}}%
\pgfpathquadraticcurveto{\pgfqpoint{8.703851in}{1.665853in}}{\pgfqpoint{8.685812in}{1.665853in}}%
\pgfpathquadraticcurveto{\pgfqpoint{8.679771in}{1.665853in}}{\pgfqpoint{8.673360in}{1.667048in}}%
\pgfpathquadraticcurveto{\pgfqpoint{8.666948in}{1.668244in}}{\pgfqpoint{8.660124in}{1.670636in}}%
\pgfpathlineto{\pgfqpoint{8.660124in}{1.683191in}}%
\pgfpathquadraticcurveto{\pgfqpoint{8.665526in}{1.680037in}}{\pgfqpoint{8.671958in}{1.678429in}}%
\pgfpathquadraticcurveto{\pgfqpoint{8.678411in}{1.676821in}}{\pgfqpoint{8.685441in}{1.676821in}}%
\pgfpathquadraticcurveto{\pgfqpoint{8.697666in}{1.676821in}}{\pgfqpoint{8.704078in}{1.681645in}}%
\pgfpathquadraticcurveto{\pgfqpoint{8.710490in}{1.686469in}}{\pgfqpoint{8.710490in}{1.695685in}}%
\pgfpathquadraticcurveto{\pgfqpoint{8.710490in}{1.704199in}}{\pgfqpoint{8.704532in}{1.708982in}}%
\pgfpathquadraticcurveto{\pgfqpoint{8.698574in}{1.713786in}}{\pgfqpoint{8.687956in}{1.713786in}}%
\pgfpathlineto{\pgfqpoint{8.676741in}{1.713786in}}%
\pgfpathlineto{\pgfqpoint{8.676741in}{1.724486in}}%
\pgfpathlineto{\pgfqpoint{8.688472in}{1.724486in}}%
\pgfpathquadraticcurveto{\pgfqpoint{8.698058in}{1.724486in}}{\pgfqpoint{8.703150in}{1.728320in}}%
\pgfpathquadraticcurveto{\pgfqpoint{8.708243in}{1.732155in}}{\pgfqpoint{8.708243in}{1.739371in}}%
\pgfpathquadraticcurveto{\pgfqpoint{8.708243in}{1.746772in}}{\pgfqpoint{8.702985in}{1.750730in}}%
\pgfpathquadraticcurveto{\pgfqpoint{8.697749in}{1.754709in}}{\pgfqpoint{8.687956in}{1.754709in}}%
\pgfpathquadraticcurveto{\pgfqpoint{8.682596in}{1.754709in}}{\pgfqpoint{8.676473in}{1.753534in}}%
\pgfpathquadraticcurveto{\pgfqpoint{8.670350in}{1.752379in}}{\pgfqpoint{8.663010in}{1.749947in}}%
\pgfpathlineto{\pgfqpoint{8.663010in}{1.761533in}}%
\pgfpathquadraticcurveto{\pgfqpoint{8.670432in}{1.763595in}}{\pgfqpoint{8.676906in}{1.764625in}}%
\pgfpathquadraticcurveto{\pgfqpoint{8.683379in}{1.765656in}}{\pgfqpoint{8.689111in}{1.765656in}}%
\pgfpathquadraticcurveto{\pgfqpoint{8.703934in}{1.765656in}}{\pgfqpoint{8.712551in}{1.758915in}}%
\pgfpathquadraticcurveto{\pgfqpoint{8.721190in}{1.752194in}}{\pgfqpoint{8.721190in}{1.740731in}}%
\pgfpathquadraticcurveto{\pgfqpoint{8.721190in}{1.732732in}}{\pgfqpoint{8.716613in}{1.727227in}}%
\pgfpathquadraticcurveto{\pgfqpoint{8.712036in}{1.721723in}}{\pgfqpoint{8.703604in}{1.719599in}}%
\pgfpathlineto{\pgfqpoint{8.703604in}{1.719599in}}%
\pgfpathclose%
\pgfpathmoveto{\pgfqpoint{8.744589in}{1.767842in}}%
\pgfpathlineto{\pgfqpoint{8.754897in}{1.767842in}}%
\pgfpathquadraticcurveto{\pgfqpoint{8.764546in}{1.752647in}}{\pgfqpoint{8.769349in}{1.738072in}}%
\pgfpathquadraticcurveto{\pgfqpoint{8.774153in}{1.723517in}}{\pgfqpoint{8.774153in}{1.709147in}}%
\pgfpathquadraticcurveto{\pgfqpoint{8.774153in}{1.694716in}}{\pgfqpoint{8.769349in}{1.680099in}}%
\pgfpathquadraticcurveto{\pgfqpoint{8.764546in}{1.665482in}}{\pgfqpoint{8.754897in}{1.650329in}}%
\pgfpathlineto{\pgfqpoint{8.744589in}{1.650329in}}%
\pgfpathquadraticcurveto{\pgfqpoint{8.753145in}{1.665090in}}{\pgfqpoint{8.757371in}{1.679686in}}%
\pgfpathquadraticcurveto{\pgfqpoint{8.761598in}{1.694283in}}{\pgfqpoint{8.761598in}{1.709147in}}%
\pgfpathquadraticcurveto{\pgfqpoint{8.761598in}{1.724032in}}{\pgfqpoint{8.757371in}{1.738525in}}%
\pgfpathquadraticcurveto{\pgfqpoint{8.753145in}{1.753039in}}{\pgfqpoint{8.744589in}{1.767842in}}%
\pgfpathlineto{\pgfqpoint{8.744589in}{1.767842in}}%
\pgfpathclose%
\pgfusepath{fill}%
\end{pgfscope}%
\begin{pgfscope}%
\pgfsetbuttcap%
\pgfsetroundjoin%
\definecolor{currentfill}{rgb}{0.090196,0.168627,0.227451}%
\pgfsetfillcolor{currentfill}%
\pgfsetlinewidth{0.000000pt}%
\definecolor{currentstroke}{rgb}{0.090196,0.168627,0.227451}%
\pgfsetstrokecolor{currentstroke}%
\pgfsetdash{}{0pt}%
\pgfpathmoveto{\pgfqpoint{7.784611in}{1.584130in}}%
\pgfpathlineto{\pgfqpoint{7.784611in}{1.571431in}}%
\pgfpathquadraticcurveto{\pgfqpoint{7.777210in}{1.574977in}}{\pgfqpoint{7.770633in}{1.576708in}}%
\pgfpathquadraticcurveto{\pgfqpoint{7.764056in}{1.578461in}}{\pgfqpoint{7.757933in}{1.578461in}}%
\pgfpathquadraticcurveto{\pgfqpoint{7.747316in}{1.578461in}}{\pgfqpoint{7.741543in}{1.574337in}}%
\pgfpathquadraticcurveto{\pgfqpoint{7.735771in}{1.570214in}}{\pgfqpoint{7.735771in}{1.562607in}}%
\pgfpathquadraticcurveto{\pgfqpoint{7.735771in}{1.556216in}}{\pgfqpoint{7.739605in}{1.552958in}}%
\pgfpathquadraticcurveto{\pgfqpoint{7.743440in}{1.549722in}}{\pgfqpoint{7.754140in}{1.547722in}}%
\pgfpathlineto{\pgfqpoint{7.761995in}{1.546114in}}%
\pgfpathquadraticcurveto{\pgfqpoint{7.776550in}{1.543331in}}{\pgfqpoint{7.783477in}{1.536342in}}%
\pgfpathquadraticcurveto{\pgfqpoint{7.790404in}{1.529353in}}{\pgfqpoint{7.790404in}{1.517643in}}%
\pgfpathquadraticcurveto{\pgfqpoint{7.790404in}{1.503644in}}{\pgfqpoint{7.781024in}{1.496428in}}%
\pgfpathquadraticcurveto{\pgfqpoint{7.771664in}{1.489213in}}{\pgfqpoint{7.753563in}{1.489213in}}%
\pgfpathquadraticcurveto{\pgfqpoint{7.746739in}{1.489213in}}{\pgfqpoint{7.739028in}{1.490759in}}%
\pgfpathquadraticcurveto{\pgfqpoint{7.731338in}{1.492305in}}{\pgfqpoint{7.723092in}{1.495336in}}%
\pgfpathlineto{\pgfqpoint{7.723092in}{1.508736in}}%
\pgfpathquadraticcurveto{\pgfqpoint{7.731008in}{1.504304in}}{\pgfqpoint{7.738616in}{1.502036in}}%
\pgfpathquadraticcurveto{\pgfqpoint{7.746223in}{1.499789in}}{\pgfqpoint{7.753563in}{1.499789in}}%
\pgfpathquadraticcurveto{\pgfqpoint{7.764696in}{1.499789in}}{\pgfqpoint{7.770757in}{1.504160in}}%
\pgfpathquadraticcurveto{\pgfqpoint{7.776818in}{1.508551in}}{\pgfqpoint{7.776818in}{1.516674in}}%
\pgfpathquadraticcurveto{\pgfqpoint{7.776818in}{1.523745in}}{\pgfqpoint{7.772468in}{1.527745in}}%
\pgfpathquadraticcurveto{\pgfqpoint{7.768118in}{1.531744in}}{\pgfqpoint{7.758201in}{1.533744in}}%
\pgfpathlineto{\pgfqpoint{7.750264in}{1.535290in}}%
\pgfpathquadraticcurveto{\pgfqpoint{7.735709in}{1.538176in}}{\pgfqpoint{7.729194in}{1.544361in}}%
\pgfpathquadraticcurveto{\pgfqpoint{7.722700in}{1.550546in}}{\pgfqpoint{7.722700in}{1.561576in}}%
\pgfpathquadraticcurveto{\pgfqpoint{7.722700in}{1.574337in}}{\pgfqpoint{7.731689in}{1.581677in}}%
\pgfpathquadraticcurveto{\pgfqpoint{7.740678in}{1.589016in}}{\pgfqpoint{7.756449in}{1.589016in}}%
\pgfpathquadraticcurveto{\pgfqpoint{7.763232in}{1.589016in}}{\pgfqpoint{7.770241in}{1.587779in}}%
\pgfpathquadraticcurveto{\pgfqpoint{7.777271in}{1.586563in}}{\pgfqpoint{7.784611in}{1.584130in}}%
\pgfpathlineto{\pgfqpoint{7.784611in}{1.584130in}}%
\pgfpathclose%
\pgfpathmoveto{\pgfqpoint{7.821926in}{1.583739in}}%
\pgfpathlineto{\pgfqpoint{7.821926in}{1.563246in}}%
\pgfpathlineto{\pgfqpoint{7.846336in}{1.563246in}}%
\pgfpathlineto{\pgfqpoint{7.846336in}{1.554030in}}%
\pgfpathlineto{\pgfqpoint{7.821926in}{1.554030in}}%
\pgfpathlineto{\pgfqpoint{7.821926in}{1.514859in}}%
\pgfpathquadraticcurveto{\pgfqpoint{7.821926in}{1.506036in}}{\pgfqpoint{7.824339in}{1.503520in}}%
\pgfpathquadraticcurveto{\pgfqpoint{7.826751in}{1.501005in}}{\pgfqpoint{7.834173in}{1.501005in}}%
\pgfpathlineto{\pgfqpoint{7.846336in}{1.501005in}}%
\pgfpathlineto{\pgfqpoint{7.846336in}{1.491089in}}%
\pgfpathlineto{\pgfqpoint{7.834173in}{1.491089in}}%
\pgfpathquadraticcurveto{\pgfqpoint{7.820442in}{1.491089in}}{\pgfqpoint{7.815226in}{1.496202in}}%
\pgfpathquadraticcurveto{\pgfqpoint{7.810010in}{1.501335in}}{\pgfqpoint{7.810010in}{1.514859in}}%
\pgfpathlineto{\pgfqpoint{7.810010in}{1.554030in}}%
\pgfpathlineto{\pgfqpoint{7.801310in}{1.554030in}}%
\pgfpathlineto{\pgfqpoint{7.801310in}{1.563246in}}%
\pgfpathlineto{\pgfqpoint{7.810010in}{1.563246in}}%
\pgfpathlineto{\pgfqpoint{7.810010in}{1.583739in}}%
\pgfpathlineto{\pgfqpoint{7.821926in}{1.583739in}}%
\pgfpathlineto{\pgfqpoint{7.821926in}{1.583739in}}%
\pgfpathclose%
\pgfpathmoveto{\pgfqpoint{7.860706in}{1.519560in}}%
\pgfpathlineto{\pgfqpoint{7.860706in}{1.563246in}}%
\pgfpathlineto{\pgfqpoint{7.872560in}{1.563246in}}%
\pgfpathlineto{\pgfqpoint{7.872560in}{1.520013in}}%
\pgfpathquadraticcurveto{\pgfqpoint{7.872560in}{1.509767in}}{\pgfqpoint{7.876539in}{1.504634in}}%
\pgfpathquadraticcurveto{\pgfqpoint{7.880539in}{1.499521in}}{\pgfqpoint{7.888538in}{1.499521in}}%
\pgfpathquadraticcurveto{\pgfqpoint{7.898124in}{1.499521in}}{\pgfqpoint{7.903691in}{1.505644in}}%
\pgfpathquadraticcurveto{\pgfqpoint{7.909278in}{1.511767in}}{\pgfqpoint{7.909278in}{1.522343in}}%
\pgfpathlineto{\pgfqpoint{7.909278in}{1.563246in}}%
\pgfpathlineto{\pgfqpoint{7.921132in}{1.563246in}}%
\pgfpathlineto{\pgfqpoint{7.921132in}{1.491089in}}%
\pgfpathlineto{\pgfqpoint{7.909278in}{1.491089in}}%
\pgfpathlineto{\pgfqpoint{7.909278in}{1.502180in}}%
\pgfpathquadraticcurveto{\pgfqpoint{7.904969in}{1.495604in}}{\pgfqpoint{7.899258in}{1.492408in}}%
\pgfpathquadraticcurveto{\pgfqpoint{7.893568in}{1.489213in}}{\pgfqpoint{7.886023in}{1.489213in}}%
\pgfpathquadraticcurveto{\pgfqpoint{7.873591in}{1.489213in}}{\pgfqpoint{7.867138in}{1.496944in}}%
\pgfpathquadraticcurveto{\pgfqpoint{7.860706in}{1.504675in}}{\pgfqpoint{7.860706in}{1.519560in}}%
\pgfpathlineto{\pgfqpoint{7.860706in}{1.519560in}}%
\pgfpathclose%
\pgfpathmoveto{\pgfqpoint{7.890538in}{1.564978in}}%
\pgfpathlineto{\pgfqpoint{7.890538in}{1.564978in}}%
\pgfpathlineto{\pgfqpoint{7.890538in}{1.564978in}}%
\pgfpathclose%
\pgfpathmoveto{\pgfqpoint{7.993021in}{1.552299in}}%
\pgfpathlineto{\pgfqpoint{7.993021in}{1.591346in}}%
\pgfpathlineto{\pgfqpoint{8.004876in}{1.591346in}}%
\pgfpathlineto{\pgfqpoint{8.004876in}{1.491089in}}%
\pgfpathlineto{\pgfqpoint{7.993021in}{1.491089in}}%
\pgfpathlineto{\pgfqpoint{7.993021in}{1.501912in}}%
\pgfpathquadraticcurveto{\pgfqpoint{7.989290in}{1.495480in}}{\pgfqpoint{7.983579in}{1.492346in}}%
\pgfpathquadraticcurveto{\pgfqpoint{7.977889in}{1.489213in}}{\pgfqpoint{7.969890in}{1.489213in}}%
\pgfpathquadraticcurveto{\pgfqpoint{7.956819in}{1.489213in}}{\pgfqpoint{7.948593in}{1.499645in}}%
\pgfpathquadraticcurveto{\pgfqpoint{7.940388in}{1.510097in}}{\pgfqpoint{7.940388in}{1.527106in}}%
\pgfpathquadraticcurveto{\pgfqpoint{7.940388in}{1.544114in}}{\pgfqpoint{7.948593in}{1.554546in}}%
\pgfpathquadraticcurveto{\pgfqpoint{7.956819in}{1.564978in}}{\pgfqpoint{7.969890in}{1.564978in}}%
\pgfpathquadraticcurveto{\pgfqpoint{7.977889in}{1.564978in}}{\pgfqpoint{7.983579in}{1.561844in}}%
\pgfpathquadraticcurveto{\pgfqpoint{7.989290in}{1.558731in}}{\pgfqpoint{7.993021in}{1.552299in}}%
\pgfpathlineto{\pgfqpoint{7.993021in}{1.552299in}}%
\pgfpathclose%
\pgfpathmoveto{\pgfqpoint{7.952634in}{1.527106in}}%
\pgfpathquadraticcurveto{\pgfqpoint{7.952634in}{1.514035in}}{\pgfqpoint{7.958015in}{1.506592in}}%
\pgfpathquadraticcurveto{\pgfqpoint{7.963396in}{1.499150in}}{\pgfqpoint{7.972797in}{1.499150in}}%
\pgfpathquadraticcurveto{\pgfqpoint{7.982198in}{1.499150in}}{\pgfqpoint{7.987599in}{1.506592in}}%
\pgfpathquadraticcurveto{\pgfqpoint{7.993021in}{1.514035in}}{\pgfqpoint{7.993021in}{1.527106in}}%
\pgfpathquadraticcurveto{\pgfqpoint{7.993021in}{1.540176in}}{\pgfqpoint{7.987599in}{1.547619in}}%
\pgfpathquadraticcurveto{\pgfqpoint{7.982198in}{1.555061in}}{\pgfqpoint{7.972797in}{1.555061in}}%
\pgfpathquadraticcurveto{\pgfqpoint{7.963396in}{1.555061in}}{\pgfqpoint{7.958015in}{1.547619in}}%
\pgfpathquadraticcurveto{\pgfqpoint{7.952634in}{1.540176in}}{\pgfqpoint{7.952634in}{1.527106in}}%
\pgfpathlineto{\pgfqpoint{7.952634in}{1.527106in}}%
\pgfpathclose%
\pgfpathmoveto{\pgfqpoint{8.059323in}{1.484388in}}%
\pgfpathquadraticcurveto{\pgfqpoint{8.054314in}{1.471503in}}{\pgfqpoint{8.049531in}{1.467586in}}%
\pgfpathquadraticcurveto{\pgfqpoint{8.044768in}{1.463648in}}{\pgfqpoint{8.036790in}{1.463648in}}%
\pgfpathlineto{\pgfqpoint{8.027306in}{1.463648in}}%
\pgfpathlineto{\pgfqpoint{8.027306in}{1.473565in}}%
\pgfpathlineto{\pgfqpoint{8.034275in}{1.473565in}}%
\pgfpathquadraticcurveto{\pgfqpoint{8.039161in}{1.473565in}}{\pgfqpoint{8.041861in}{1.475895in}}%
\pgfpathquadraticcurveto{\pgfqpoint{8.044583in}{1.478204in}}{\pgfqpoint{8.047861in}{1.486842in}}%
\pgfpathlineto{\pgfqpoint{8.049984in}{1.492243in}}%
\pgfpathlineto{\pgfqpoint{8.020812in}{1.563246in}}%
\pgfpathlineto{\pgfqpoint{8.033367in}{1.563246in}}%
\pgfpathlineto{\pgfqpoint{8.055922in}{1.506819in}}%
\pgfpathlineto{\pgfqpoint{8.078476in}{1.563246in}}%
\pgfpathlineto{\pgfqpoint{8.091031in}{1.563246in}}%
\pgfpathlineto{\pgfqpoint{8.059323in}{1.484388in}}%
\pgfpathlineto{\pgfqpoint{8.059323in}{1.484388in}}%
\pgfpathclose%
\pgfpathmoveto{\pgfqpoint{8.153272in}{1.502036in}}%
\pgfpathlineto{\pgfqpoint{8.174527in}{1.502036in}}%
\pgfpathlineto{\pgfqpoint{8.174527in}{1.575430in}}%
\pgfpathlineto{\pgfqpoint{8.151396in}{1.570791in}}%
\pgfpathlineto{\pgfqpoint{8.151396in}{1.582646in}}%
\pgfpathlineto{\pgfqpoint{8.174404in}{1.587285in}}%
\pgfpathlineto{\pgfqpoint{8.187413in}{1.587285in}}%
\pgfpathlineto{\pgfqpoint{8.187413in}{1.502036in}}%
\pgfpathlineto{\pgfqpoint{8.208668in}{1.502036in}}%
\pgfpathlineto{\pgfqpoint{8.208668in}{1.491089in}}%
\pgfpathlineto{\pgfqpoint{8.153272in}{1.491089in}}%
\pgfpathlineto{\pgfqpoint{8.153272in}{1.502036in}}%
\pgfpathlineto{\pgfqpoint{8.153272in}{1.502036in}}%
\pgfpathclose%
\pgfpathmoveto{\pgfqpoint{8.262786in}{1.536775in}}%
\pgfpathquadraticcurveto{\pgfqpoint{8.253508in}{1.536775in}}{\pgfqpoint{8.248189in}{1.531806in}}%
\pgfpathquadraticcurveto{\pgfqpoint{8.242891in}{1.526837in}}{\pgfqpoint{8.242891in}{1.518158in}}%
\pgfpathquadraticcurveto{\pgfqpoint{8.242891in}{1.509458in}}{\pgfqpoint{8.248189in}{1.504489in}}%
\pgfpathquadraticcurveto{\pgfqpoint{8.253508in}{1.499521in}}{\pgfqpoint{8.262786in}{1.499521in}}%
\pgfpathquadraticcurveto{\pgfqpoint{8.272063in}{1.499521in}}{\pgfqpoint{8.277403in}{1.504510in}}%
\pgfpathquadraticcurveto{\pgfqpoint{8.282763in}{1.509520in}}{\pgfqpoint{8.282763in}{1.518158in}}%
\pgfpathquadraticcurveto{\pgfqpoint{8.282763in}{1.526837in}}{\pgfqpoint{8.277444in}{1.531806in}}%
\pgfpathquadraticcurveto{\pgfqpoint{8.272146in}{1.536775in}}{\pgfqpoint{8.262786in}{1.536775in}}%
\pgfpathlineto{\pgfqpoint{8.262786in}{1.536775in}}%
\pgfpathclose%
\pgfpathmoveto{\pgfqpoint{8.249777in}{1.542300in}}%
\pgfpathquadraticcurveto{\pgfqpoint{8.241407in}{1.544361in}}{\pgfqpoint{8.236727in}{1.550093in}}%
\pgfpathquadraticcurveto{\pgfqpoint{8.232067in}{1.555845in}}{\pgfqpoint{8.232067in}{1.564091in}}%
\pgfpathquadraticcurveto{\pgfqpoint{8.232067in}{1.575616in}}{\pgfqpoint{8.240273in}{1.582316in}}%
\pgfpathquadraticcurveto{\pgfqpoint{8.248499in}{1.589016in}}{\pgfqpoint{8.262786in}{1.589016in}}%
\pgfpathquadraticcurveto{\pgfqpoint{8.277155in}{1.589016in}}{\pgfqpoint{8.285340in}{1.582316in}}%
\pgfpathquadraticcurveto{\pgfqpoint{8.293525in}{1.575616in}}{\pgfqpoint{8.293525in}{1.564091in}}%
\pgfpathquadraticcurveto{\pgfqpoint{8.293525in}{1.555845in}}{\pgfqpoint{8.288845in}{1.550093in}}%
\pgfpathquadraticcurveto{\pgfqpoint{8.284186in}{1.544361in}}{\pgfqpoint{8.275877in}{1.542300in}}%
\pgfpathquadraticcurveto{\pgfqpoint{8.285278in}{1.540114in}}{\pgfqpoint{8.290515in}{1.533723in}}%
\pgfpathquadraticcurveto{\pgfqpoint{8.295772in}{1.527353in}}{\pgfqpoint{8.295772in}{1.518158in}}%
\pgfpathquadraticcurveto{\pgfqpoint{8.295772in}{1.504160in}}{\pgfqpoint{8.287237in}{1.496676in}}%
\pgfpathquadraticcurveto{\pgfqpoint{8.278702in}{1.489213in}}{\pgfqpoint{8.262786in}{1.489213in}}%
\pgfpathquadraticcurveto{\pgfqpoint{8.246891in}{1.489213in}}{\pgfqpoint{8.238335in}{1.496676in}}%
\pgfpathquadraticcurveto{\pgfqpoint{8.229800in}{1.504160in}}{\pgfqpoint{8.229800in}{1.518158in}}%
\pgfpathquadraticcurveto{\pgfqpoint{8.229800in}{1.527353in}}{\pgfqpoint{8.235077in}{1.533723in}}%
\pgfpathquadraticcurveto{\pgfqpoint{8.240376in}{1.540114in}}{\pgfqpoint{8.249777in}{1.542300in}}%
\pgfpathlineto{\pgfqpoint{8.249777in}{1.542300in}}%
\pgfpathclose%
\pgfpathmoveto{\pgfqpoint{8.245015in}{1.562854in}}%
\pgfpathquadraticcurveto{\pgfqpoint{8.245015in}{1.555391in}}{\pgfqpoint{8.249674in}{1.551206in}}%
\pgfpathquadraticcurveto{\pgfqpoint{8.254354in}{1.547021in}}{\pgfqpoint{8.262786in}{1.547021in}}%
\pgfpathquadraticcurveto{\pgfqpoint{8.271177in}{1.547021in}}{\pgfqpoint{8.275898in}{1.551206in}}%
\pgfpathquadraticcurveto{\pgfqpoint{8.280640in}{1.555391in}}{\pgfqpoint{8.280640in}{1.562854in}}%
\pgfpathquadraticcurveto{\pgfqpoint{8.280640in}{1.570338in}}{\pgfqpoint{8.275898in}{1.574523in}}%
\pgfpathquadraticcurveto{\pgfqpoint{8.271177in}{1.578708in}}{\pgfqpoint{8.262786in}{1.578708in}}%
\pgfpathquadraticcurveto{\pgfqpoint{8.254354in}{1.578708in}}{\pgfqpoint{8.249674in}{1.574523in}}%
\pgfpathquadraticcurveto{\pgfqpoint{8.245015in}{1.570338in}}{\pgfqpoint{8.245015in}{1.562854in}}%
\pgfpathlineto{\pgfqpoint{8.245015in}{1.562854in}}%
\pgfpathclose%
\pgfpathmoveto{\pgfqpoint{8.387638in}{1.591202in}}%
\pgfpathquadraticcurveto{\pgfqpoint{8.379021in}{1.576399in}}{\pgfqpoint{8.374815in}{1.561885in}}%
\pgfpathquadraticcurveto{\pgfqpoint{8.370630in}{1.547392in}}{\pgfqpoint{8.370630in}{1.532507in}}%
\pgfpathquadraticcurveto{\pgfqpoint{8.370630in}{1.517643in}}{\pgfqpoint{8.374856in}{1.503046in}}%
\pgfpathquadraticcurveto{\pgfqpoint{8.379082in}{1.488450in}}{\pgfqpoint{8.387638in}{1.473689in}}%
\pgfpathlineto{\pgfqpoint{8.377330in}{1.473689in}}%
\pgfpathquadraticcurveto{\pgfqpoint{8.367682in}{1.488842in}}{\pgfqpoint{8.362878in}{1.503459in}}%
\pgfpathquadraticcurveto{\pgfqpoint{8.358074in}{1.518076in}}{\pgfqpoint{8.358074in}{1.532507in}}%
\pgfpathquadraticcurveto{\pgfqpoint{8.358074in}{1.546877in}}{\pgfqpoint{8.362837in}{1.561432in}}%
\pgfpathquadraticcurveto{\pgfqpoint{8.367620in}{1.576007in}}{\pgfqpoint{8.377330in}{1.591202in}}%
\pgfpathlineto{\pgfqpoint{8.387638in}{1.591202in}}%
\pgfpathlineto{\pgfqpoint{8.387638in}{1.591202in}}%
\pgfpathclose%
\pgfpathmoveto{\pgfqpoint{8.423531in}{1.502036in}}%
\pgfpathlineto{\pgfqpoint{8.468949in}{1.502036in}}%
\pgfpathlineto{\pgfqpoint{8.468949in}{1.491089in}}%
\pgfpathlineto{\pgfqpoint{8.407883in}{1.491089in}}%
\pgfpathlineto{\pgfqpoint{8.407883in}{1.502036in}}%
\pgfpathquadraticcurveto{\pgfqpoint{8.415285in}{1.509705in}}{\pgfqpoint{8.428067in}{1.522611in}}%
\pgfpathquadraticcurveto{\pgfqpoint{8.440870in}{1.535538in}}{\pgfqpoint{8.444148in}{1.539290in}}%
\pgfpathquadraticcurveto{\pgfqpoint{8.450394in}{1.546299in}}{\pgfqpoint{8.452868in}{1.551165in}}%
\pgfpathquadraticcurveto{\pgfqpoint{8.455363in}{1.556030in}}{\pgfqpoint{8.455363in}{1.560731in}}%
\pgfpathquadraticcurveto{\pgfqpoint{8.455363in}{1.568400in}}{\pgfqpoint{8.449982in}{1.573224in}}%
\pgfpathquadraticcurveto{\pgfqpoint{8.444601in}{1.578069in}}{\pgfqpoint{8.435963in}{1.578069in}}%
\pgfpathquadraticcurveto{\pgfqpoint{8.429840in}{1.578069in}}{\pgfqpoint{8.423036in}{1.575946in}}%
\pgfpathquadraticcurveto{\pgfqpoint{8.416254in}{1.573822in}}{\pgfqpoint{8.408523in}{1.569493in}}%
\pgfpathlineto{\pgfqpoint{8.408523in}{1.582646in}}%
\pgfpathquadraticcurveto{\pgfqpoint{8.416377in}{1.585800in}}{\pgfqpoint{8.423201in}{1.587408in}}%
\pgfpathquadraticcurveto{\pgfqpoint{8.430046in}{1.589016in}}{\pgfqpoint{8.435715in}{1.589016in}}%
\pgfpathquadraticcurveto{\pgfqpoint{8.450662in}{1.589016in}}{\pgfqpoint{8.459548in}{1.581533in}}%
\pgfpathquadraticcurveto{\pgfqpoint{8.468434in}{1.574069in}}{\pgfqpoint{8.468434in}{1.561576in}}%
\pgfpathquadraticcurveto{\pgfqpoint{8.468434in}{1.555638in}}{\pgfqpoint{8.466207in}{1.550319in}}%
\pgfpathquadraticcurveto{\pgfqpoint{8.464001in}{1.545021in}}{\pgfqpoint{8.458125in}{1.537805in}}%
\pgfpathquadraticcurveto{\pgfqpoint{8.456517in}{1.535929in}}{\pgfqpoint{8.447879in}{1.527002in}}%
\pgfpathquadraticcurveto{\pgfqpoint{8.439262in}{1.518076in}}{\pgfqpoint{8.423531in}{1.502036in}}%
\pgfpathlineto{\pgfqpoint{8.423531in}{1.502036in}}%
\pgfpathclose%
\pgfpathmoveto{\pgfqpoint{8.524098in}{1.578708in}}%
\pgfpathquadraticcurveto{\pgfqpoint{8.514058in}{1.578708in}}{\pgfqpoint{8.508986in}{1.568812in}}%
\pgfpathquadraticcurveto{\pgfqpoint{8.503935in}{1.558937in}}{\pgfqpoint{8.503935in}{1.539084in}}%
\pgfpathquadraticcurveto{\pgfqpoint{8.503935in}{1.519313in}}{\pgfqpoint{8.508986in}{1.509417in}}%
\pgfpathquadraticcurveto{\pgfqpoint{8.514058in}{1.499521in}}{\pgfqpoint{8.524098in}{1.499521in}}%
\pgfpathquadraticcurveto{\pgfqpoint{8.534220in}{1.499521in}}{\pgfqpoint{8.539271in}{1.509417in}}%
\pgfpathquadraticcurveto{\pgfqpoint{8.544343in}{1.519313in}}{\pgfqpoint{8.544343in}{1.539084in}}%
\pgfpathquadraticcurveto{\pgfqpoint{8.544343in}{1.558937in}}{\pgfqpoint{8.539271in}{1.568812in}}%
\pgfpathquadraticcurveto{\pgfqpoint{8.534220in}{1.578708in}}{\pgfqpoint{8.524098in}{1.578708in}}%
\pgfpathlineto{\pgfqpoint{8.524098in}{1.578708in}}%
\pgfpathclose%
\pgfpathmoveto{\pgfqpoint{8.524098in}{1.589016in}}%
\pgfpathquadraticcurveto{\pgfqpoint{8.540281in}{1.589016in}}{\pgfqpoint{8.548817in}{1.576214in}}%
\pgfpathquadraticcurveto{\pgfqpoint{8.557352in}{1.563431in}}{\pgfqpoint{8.557352in}{1.539084in}}%
\pgfpathquadraticcurveto{\pgfqpoint{8.557352in}{1.514798in}}{\pgfqpoint{8.548817in}{1.501995in}}%
\pgfpathquadraticcurveto{\pgfqpoint{8.540281in}{1.489213in}}{\pgfqpoint{8.524098in}{1.489213in}}%
\pgfpathquadraticcurveto{\pgfqpoint{8.507934in}{1.489213in}}{\pgfqpoint{8.499399in}{1.501995in}}%
\pgfpathquadraticcurveto{\pgfqpoint{8.490864in}{1.514798in}}{\pgfqpoint{8.490864in}{1.539084in}}%
\pgfpathquadraticcurveto{\pgfqpoint{8.490864in}{1.563431in}}{\pgfqpoint{8.499399in}{1.576214in}}%
\pgfpathquadraticcurveto{\pgfqpoint{8.507934in}{1.589016in}}{\pgfqpoint{8.524098in}{1.589016in}}%
\pgfpathlineto{\pgfqpoint{8.524098in}{1.589016in}}%
\pgfpathclose%
\pgfpathmoveto{\pgfqpoint{8.591431in}{1.502036in}}%
\pgfpathlineto{\pgfqpoint{8.636848in}{1.502036in}}%
\pgfpathlineto{\pgfqpoint{8.636848in}{1.491089in}}%
\pgfpathlineto{\pgfqpoint{8.575783in}{1.491089in}}%
\pgfpathlineto{\pgfqpoint{8.575783in}{1.502036in}}%
\pgfpathquadraticcurveto{\pgfqpoint{8.583184in}{1.509705in}}{\pgfqpoint{8.595966in}{1.522611in}}%
\pgfpathquadraticcurveto{\pgfqpoint{8.608769in}{1.535538in}}{\pgfqpoint{8.612047in}{1.539290in}}%
\pgfpathquadraticcurveto{\pgfqpoint{8.618294in}{1.546299in}}{\pgfqpoint{8.620768in}{1.551165in}}%
\pgfpathquadraticcurveto{\pgfqpoint{8.623262in}{1.556030in}}{\pgfqpoint{8.623262in}{1.560731in}}%
\pgfpathquadraticcurveto{\pgfqpoint{8.623262in}{1.568400in}}{\pgfqpoint{8.617881in}{1.573224in}}%
\pgfpathquadraticcurveto{\pgfqpoint{8.612500in}{1.578069in}}{\pgfqpoint{8.603862in}{1.578069in}}%
\pgfpathquadraticcurveto{\pgfqpoint{8.597739in}{1.578069in}}{\pgfqpoint{8.590936in}{1.575946in}}%
\pgfpathquadraticcurveto{\pgfqpoint{8.584153in}{1.573822in}}{\pgfqpoint{8.576422in}{1.569493in}}%
\pgfpathlineto{\pgfqpoint{8.576422in}{1.582646in}}%
\pgfpathquadraticcurveto{\pgfqpoint{8.584277in}{1.585800in}}{\pgfqpoint{8.591101in}{1.587408in}}%
\pgfpathquadraticcurveto{\pgfqpoint{8.597945in}{1.589016in}}{\pgfqpoint{8.603615in}{1.589016in}}%
\pgfpathquadraticcurveto{\pgfqpoint{8.618562in}{1.589016in}}{\pgfqpoint{8.627447in}{1.581533in}}%
\pgfpathquadraticcurveto{\pgfqpoint{8.636333in}{1.574069in}}{\pgfqpoint{8.636333in}{1.561576in}}%
\pgfpathquadraticcurveto{\pgfqpoint{8.636333in}{1.555638in}}{\pgfqpoint{8.634106in}{1.550319in}}%
\pgfpathquadraticcurveto{\pgfqpoint{8.631900in}{1.545021in}}{\pgfqpoint{8.626025in}{1.537805in}}%
\pgfpathquadraticcurveto{\pgfqpoint{8.624417in}{1.535929in}}{\pgfqpoint{8.615778in}{1.527002in}}%
\pgfpathquadraticcurveto{\pgfqpoint{8.607161in}{1.518076in}}{\pgfqpoint{8.591431in}{1.502036in}}%
\pgfpathlineto{\pgfqpoint{8.591431in}{1.502036in}}%
\pgfpathclose%
\pgfpathmoveto{\pgfqpoint{8.699934in}{1.575946in}}%
\pgfpathlineto{\pgfqpoint{8.667072in}{1.524590in}}%
\pgfpathlineto{\pgfqpoint{8.699934in}{1.524590in}}%
\pgfpathlineto{\pgfqpoint{8.699934in}{1.575946in}}%
\pgfpathlineto{\pgfqpoint{8.699934in}{1.575946in}}%
\pgfpathclose%
\pgfpathmoveto{\pgfqpoint{8.696512in}{1.587285in}}%
\pgfpathlineto{\pgfqpoint{8.712881in}{1.587285in}}%
\pgfpathlineto{\pgfqpoint{8.712881in}{1.524590in}}%
\pgfpathlineto{\pgfqpoint{8.726612in}{1.524590in}}%
\pgfpathlineto{\pgfqpoint{8.726612in}{1.513767in}}%
\pgfpathlineto{\pgfqpoint{8.712881in}{1.513767in}}%
\pgfpathlineto{\pgfqpoint{8.712881in}{1.491089in}}%
\pgfpathlineto{\pgfqpoint{8.699934in}{1.491089in}}%
\pgfpathlineto{\pgfqpoint{8.699934in}{1.513767in}}%
\pgfpathlineto{\pgfqpoint{8.656516in}{1.513767in}}%
\pgfpathlineto{\pgfqpoint{8.656516in}{1.526322in}}%
\pgfpathlineto{\pgfqpoint{8.696512in}{1.587285in}}%
\pgfpathlineto{\pgfqpoint{8.696512in}{1.587285in}}%
\pgfpathclose%
\pgfpathmoveto{\pgfqpoint{8.744589in}{1.591202in}}%
\pgfpathlineto{\pgfqpoint{8.754897in}{1.591202in}}%
\pgfpathquadraticcurveto{\pgfqpoint{8.764546in}{1.576007in}}{\pgfqpoint{8.769349in}{1.561432in}}%
\pgfpathquadraticcurveto{\pgfqpoint{8.774153in}{1.546877in}}{\pgfqpoint{8.774153in}{1.532507in}}%
\pgfpathquadraticcurveto{\pgfqpoint{8.774153in}{1.518076in}}{\pgfqpoint{8.769349in}{1.503459in}}%
\pgfpathquadraticcurveto{\pgfqpoint{8.764546in}{1.488842in}}{\pgfqpoint{8.754897in}{1.473689in}}%
\pgfpathlineto{\pgfqpoint{8.744589in}{1.473689in}}%
\pgfpathquadraticcurveto{\pgfqpoint{8.753145in}{1.488450in}}{\pgfqpoint{8.757371in}{1.503046in}}%
\pgfpathquadraticcurveto{\pgfqpoint{8.761598in}{1.517643in}}{\pgfqpoint{8.761598in}{1.532507in}}%
\pgfpathquadraticcurveto{\pgfqpoint{8.761598in}{1.547392in}}{\pgfqpoint{8.757371in}{1.561885in}}%
\pgfpathquadraticcurveto{\pgfqpoint{8.753145in}{1.576399in}}{\pgfqpoint{8.744589in}{1.591202in}}%
\pgfpathlineto{\pgfqpoint{8.744589in}{1.591202in}}%
\pgfpathclose%
\pgfusepath{fill}%
\end{pgfscope}%
\begin{pgfscope}%
\pgfsetbuttcap%
\pgfsetroundjoin%
\definecolor{currentfill}{rgb}{0.090196,0.168627,0.227451}%
\pgfsetfillcolor{currentfill}%
\pgfsetlinewidth{0.000000pt}%
\definecolor{currentstroke}{rgb}{0.090196,0.168627,0.227451}%
\pgfsetstrokecolor{currentstroke}%
\pgfsetdash{}{0pt}%
\pgfpathmoveto{\pgfqpoint{11.018611in}{2.290690in}}%
\pgfpathlineto{\pgfqpoint{11.018611in}{2.277991in}}%
\pgfpathquadraticcurveto{\pgfqpoint{11.011210in}{2.281537in}}{\pgfqpoint{11.004633in}{2.283268in}}%
\pgfpathquadraticcurveto{\pgfqpoint{10.998056in}{2.285021in}}{\pgfqpoint{10.991933in}{2.285021in}}%
\pgfpathquadraticcurveto{\pgfqpoint{10.981316in}{2.285021in}}{\pgfqpoint{10.975543in}{2.280897in}}%
\pgfpathquadraticcurveto{\pgfqpoint{10.969771in}{2.276774in}}{\pgfqpoint{10.969771in}{2.269167in}}%
\pgfpathquadraticcurveto{\pgfqpoint{10.969771in}{2.262776in}}{\pgfqpoint{10.973605in}{2.259518in}}%
\pgfpathquadraticcurveto{\pgfqpoint{10.977440in}{2.256282in}}{\pgfqpoint{10.988140in}{2.254282in}}%
\pgfpathlineto{\pgfqpoint{10.995995in}{2.252674in}}%
\pgfpathquadraticcurveto{\pgfqpoint{11.010550in}{2.249891in}}{\pgfqpoint{11.017477in}{2.242902in}}%
\pgfpathquadraticcurveto{\pgfqpoint{11.024404in}{2.235913in}}{\pgfqpoint{11.024404in}{2.224203in}}%
\pgfpathquadraticcurveto{\pgfqpoint{11.024404in}{2.210204in}}{\pgfqpoint{11.015024in}{2.202988in}}%
\pgfpathquadraticcurveto{\pgfqpoint{11.005664in}{2.195773in}}{\pgfqpoint{10.987563in}{2.195773in}}%
\pgfpathquadraticcurveto{\pgfqpoint{10.980739in}{2.195773in}}{\pgfqpoint{10.973028in}{2.197319in}}%
\pgfpathquadraticcurveto{\pgfqpoint{10.965338in}{2.198865in}}{\pgfqpoint{10.957092in}{2.201896in}}%
\pgfpathlineto{\pgfqpoint{10.957092in}{2.215296in}}%
\pgfpathquadraticcurveto{\pgfqpoint{10.965008in}{2.210864in}}{\pgfqpoint{10.972616in}{2.208596in}}%
\pgfpathquadraticcurveto{\pgfqpoint{10.980223in}{2.206349in}}{\pgfqpoint{10.987563in}{2.206349in}}%
\pgfpathquadraticcurveto{\pgfqpoint{10.998696in}{2.206349in}}{\pgfqpoint{11.004757in}{2.210720in}}%
\pgfpathquadraticcurveto{\pgfqpoint{11.010818in}{2.215111in}}{\pgfqpoint{11.010818in}{2.223234in}}%
\pgfpathquadraticcurveto{\pgfqpoint{11.010818in}{2.230305in}}{\pgfqpoint{11.006468in}{2.234305in}}%
\pgfpathquadraticcurveto{\pgfqpoint{11.002118in}{2.238304in}}{\pgfqpoint{10.992201in}{2.240304in}}%
\pgfpathlineto{\pgfqpoint{10.984264in}{2.241850in}}%
\pgfpathquadraticcurveto{\pgfqpoint{10.969709in}{2.244736in}}{\pgfqpoint{10.963194in}{2.250921in}}%
\pgfpathquadraticcurveto{\pgfqpoint{10.956700in}{2.257106in}}{\pgfqpoint{10.956700in}{2.268136in}}%
\pgfpathquadraticcurveto{\pgfqpoint{10.956700in}{2.280897in}}{\pgfqpoint{10.965689in}{2.288237in}}%
\pgfpathquadraticcurveto{\pgfqpoint{10.974678in}{2.295576in}}{\pgfqpoint{10.990449in}{2.295576in}}%
\pgfpathquadraticcurveto{\pgfqpoint{10.997232in}{2.295576in}}{\pgfqpoint{11.004241in}{2.294339in}}%
\pgfpathquadraticcurveto{\pgfqpoint{11.011271in}{2.293123in}}{\pgfqpoint{11.018611in}{2.290690in}}%
\pgfpathlineto{\pgfqpoint{11.018611in}{2.290690in}}%
\pgfpathclose%
\pgfpathmoveto{\pgfqpoint{11.055926in}{2.290299in}}%
\pgfpathlineto{\pgfqpoint{11.055926in}{2.269806in}}%
\pgfpathlineto{\pgfqpoint{11.080336in}{2.269806in}}%
\pgfpathlineto{\pgfqpoint{11.080336in}{2.260590in}}%
\pgfpathlineto{\pgfqpoint{11.055926in}{2.260590in}}%
\pgfpathlineto{\pgfqpoint{11.055926in}{2.221419in}}%
\pgfpathquadraticcurveto{\pgfqpoint{11.055926in}{2.212596in}}{\pgfqpoint{11.058339in}{2.210080in}}%
\pgfpathquadraticcurveto{\pgfqpoint{11.060751in}{2.207565in}}{\pgfqpoint{11.068173in}{2.207565in}}%
\pgfpathlineto{\pgfqpoint{11.080336in}{2.207565in}}%
\pgfpathlineto{\pgfqpoint{11.080336in}{2.197649in}}%
\pgfpathlineto{\pgfqpoint{11.068173in}{2.197649in}}%
\pgfpathquadraticcurveto{\pgfqpoint{11.054442in}{2.197649in}}{\pgfqpoint{11.049226in}{2.202762in}}%
\pgfpathquadraticcurveto{\pgfqpoint{11.044010in}{2.207895in}}{\pgfqpoint{11.044010in}{2.221419in}}%
\pgfpathlineto{\pgfqpoint{11.044010in}{2.260590in}}%
\pgfpathlineto{\pgfqpoint{11.035310in}{2.260590in}}%
\pgfpathlineto{\pgfqpoint{11.035310in}{2.269806in}}%
\pgfpathlineto{\pgfqpoint{11.044010in}{2.269806in}}%
\pgfpathlineto{\pgfqpoint{11.044010in}{2.290299in}}%
\pgfpathlineto{\pgfqpoint{11.055926in}{2.290299in}}%
\pgfpathlineto{\pgfqpoint{11.055926in}{2.290299in}}%
\pgfpathclose%
\pgfpathmoveto{\pgfqpoint{11.094706in}{2.226120in}}%
\pgfpathlineto{\pgfqpoint{11.094706in}{2.269806in}}%
\pgfpathlineto{\pgfqpoint{11.106560in}{2.269806in}}%
\pgfpathlineto{\pgfqpoint{11.106560in}{2.226573in}}%
\pgfpathquadraticcurveto{\pgfqpoint{11.106560in}{2.216327in}}{\pgfqpoint{11.110539in}{2.211194in}}%
\pgfpathquadraticcurveto{\pgfqpoint{11.114539in}{2.206081in}}{\pgfqpoint{11.122538in}{2.206081in}}%
\pgfpathquadraticcurveto{\pgfqpoint{11.132124in}{2.206081in}}{\pgfqpoint{11.137691in}{2.212204in}}%
\pgfpathquadraticcurveto{\pgfqpoint{11.143278in}{2.218327in}}{\pgfqpoint{11.143278in}{2.228903in}}%
\pgfpathlineto{\pgfqpoint{11.143278in}{2.269806in}}%
\pgfpathlineto{\pgfqpoint{11.155132in}{2.269806in}}%
\pgfpathlineto{\pgfqpoint{11.155132in}{2.197649in}}%
\pgfpathlineto{\pgfqpoint{11.143278in}{2.197649in}}%
\pgfpathlineto{\pgfqpoint{11.143278in}{2.208740in}}%
\pgfpathquadraticcurveto{\pgfqpoint{11.138969in}{2.202164in}}{\pgfqpoint{11.133258in}{2.198968in}}%
\pgfpathquadraticcurveto{\pgfqpoint{11.127568in}{2.195773in}}{\pgfqpoint{11.120023in}{2.195773in}}%
\pgfpathquadraticcurveto{\pgfqpoint{11.107591in}{2.195773in}}{\pgfqpoint{11.101138in}{2.203504in}}%
\pgfpathquadraticcurveto{\pgfqpoint{11.094706in}{2.211235in}}{\pgfqpoint{11.094706in}{2.226120in}}%
\pgfpathlineto{\pgfqpoint{11.094706in}{2.226120in}}%
\pgfpathclose%
\pgfpathmoveto{\pgfqpoint{11.124538in}{2.271538in}}%
\pgfpathlineto{\pgfqpoint{11.124538in}{2.271538in}}%
\pgfpathlineto{\pgfqpoint{11.124538in}{2.271538in}}%
\pgfpathclose%
\pgfpathmoveto{\pgfqpoint{11.227021in}{2.258859in}}%
\pgfpathlineto{\pgfqpoint{11.227021in}{2.297906in}}%
\pgfpathlineto{\pgfqpoint{11.238876in}{2.297906in}}%
\pgfpathlineto{\pgfqpoint{11.238876in}{2.197649in}}%
\pgfpathlineto{\pgfqpoint{11.227021in}{2.197649in}}%
\pgfpathlineto{\pgfqpoint{11.227021in}{2.208472in}}%
\pgfpathquadraticcurveto{\pgfqpoint{11.223290in}{2.202040in}}{\pgfqpoint{11.217579in}{2.198906in}}%
\pgfpathquadraticcurveto{\pgfqpoint{11.211889in}{2.195773in}}{\pgfqpoint{11.203890in}{2.195773in}}%
\pgfpathquadraticcurveto{\pgfqpoint{11.190819in}{2.195773in}}{\pgfqpoint{11.182593in}{2.206205in}}%
\pgfpathquadraticcurveto{\pgfqpoint{11.174388in}{2.216657in}}{\pgfqpoint{11.174388in}{2.233666in}}%
\pgfpathquadraticcurveto{\pgfqpoint{11.174388in}{2.250674in}}{\pgfqpoint{11.182593in}{2.261106in}}%
\pgfpathquadraticcurveto{\pgfqpoint{11.190819in}{2.271538in}}{\pgfqpoint{11.203890in}{2.271538in}}%
\pgfpathquadraticcurveto{\pgfqpoint{11.211889in}{2.271538in}}{\pgfqpoint{11.217579in}{2.268404in}}%
\pgfpathquadraticcurveto{\pgfqpoint{11.223290in}{2.265291in}}{\pgfqpoint{11.227021in}{2.258859in}}%
\pgfpathlineto{\pgfqpoint{11.227021in}{2.258859in}}%
\pgfpathclose%
\pgfpathmoveto{\pgfqpoint{11.186634in}{2.233666in}}%
\pgfpathquadraticcurveto{\pgfqpoint{11.186634in}{2.220595in}}{\pgfqpoint{11.192015in}{2.213152in}}%
\pgfpathquadraticcurveto{\pgfqpoint{11.197396in}{2.205710in}}{\pgfqpoint{11.206797in}{2.205710in}}%
\pgfpathquadraticcurveto{\pgfqpoint{11.216198in}{2.205710in}}{\pgfqpoint{11.221599in}{2.213152in}}%
\pgfpathquadraticcurveto{\pgfqpoint{11.227021in}{2.220595in}}{\pgfqpoint{11.227021in}{2.233666in}}%
\pgfpathquadraticcurveto{\pgfqpoint{11.227021in}{2.246736in}}{\pgfqpoint{11.221599in}{2.254179in}}%
\pgfpathquadraticcurveto{\pgfqpoint{11.216198in}{2.261621in}}{\pgfqpoint{11.206797in}{2.261621in}}%
\pgfpathquadraticcurveto{\pgfqpoint{11.197396in}{2.261621in}}{\pgfqpoint{11.192015in}{2.254179in}}%
\pgfpathquadraticcurveto{\pgfqpoint{11.186634in}{2.246736in}}{\pgfqpoint{11.186634in}{2.233666in}}%
\pgfpathlineto{\pgfqpoint{11.186634in}{2.233666in}}%
\pgfpathclose%
\pgfpathmoveto{\pgfqpoint{11.293323in}{2.190948in}}%
\pgfpathquadraticcurveto{\pgfqpoint{11.288314in}{2.178063in}}{\pgfqpoint{11.283531in}{2.174146in}}%
\pgfpathquadraticcurveto{\pgfqpoint{11.278768in}{2.170208in}}{\pgfqpoint{11.270790in}{2.170208in}}%
\pgfpathlineto{\pgfqpoint{11.261306in}{2.170208in}}%
\pgfpathlineto{\pgfqpoint{11.261306in}{2.180125in}}%
\pgfpathlineto{\pgfqpoint{11.268275in}{2.180125in}}%
\pgfpathquadraticcurveto{\pgfqpoint{11.273161in}{2.180125in}}{\pgfqpoint{11.275861in}{2.182455in}}%
\pgfpathquadraticcurveto{\pgfqpoint{11.278583in}{2.184764in}}{\pgfqpoint{11.281861in}{2.193402in}}%
\pgfpathlineto{\pgfqpoint{11.283984in}{2.198803in}}%
\pgfpathlineto{\pgfqpoint{11.254812in}{2.269806in}}%
\pgfpathlineto{\pgfqpoint{11.267367in}{2.269806in}}%
\pgfpathlineto{\pgfqpoint{11.289922in}{2.213379in}}%
\pgfpathlineto{\pgfqpoint{11.312476in}{2.269806in}}%
\pgfpathlineto{\pgfqpoint{11.325031in}{2.269806in}}%
\pgfpathlineto{\pgfqpoint{11.293323in}{2.190948in}}%
\pgfpathlineto{\pgfqpoint{11.293323in}{2.190948in}}%
\pgfpathclose%
\pgfpathmoveto{\pgfqpoint{11.387272in}{2.208596in}}%
\pgfpathlineto{\pgfqpoint{11.408527in}{2.208596in}}%
\pgfpathlineto{\pgfqpoint{11.408527in}{2.281990in}}%
\pgfpathlineto{\pgfqpoint{11.385396in}{2.277351in}}%
\pgfpathlineto{\pgfqpoint{11.385396in}{2.289206in}}%
\pgfpathlineto{\pgfqpoint{11.408404in}{2.293845in}}%
\pgfpathlineto{\pgfqpoint{11.421413in}{2.293845in}}%
\pgfpathlineto{\pgfqpoint{11.421413in}{2.208596in}}%
\pgfpathlineto{\pgfqpoint{11.442668in}{2.208596in}}%
\pgfpathlineto{\pgfqpoint{11.442668in}{2.197649in}}%
\pgfpathlineto{\pgfqpoint{11.387272in}{2.197649in}}%
\pgfpathlineto{\pgfqpoint{11.387272in}{2.208596in}}%
\pgfpathlineto{\pgfqpoint{11.387272in}{2.208596in}}%
\pgfpathclose%
\pgfpathmoveto{\pgfqpoint{11.469345in}{2.199649in}}%
\pgfpathlineto{\pgfqpoint{11.469345in}{2.211503in}}%
\pgfpathquadraticcurveto{\pgfqpoint{11.474252in}{2.209173in}}{\pgfqpoint{11.479262in}{2.207957in}}%
\pgfpathquadraticcurveto{\pgfqpoint{11.484292in}{2.206741in}}{\pgfqpoint{11.489137in}{2.206741in}}%
\pgfpathquadraticcurveto{\pgfqpoint{11.502022in}{2.206741in}}{\pgfqpoint{11.508805in}{2.215399in}}%
\pgfpathquadraticcurveto{\pgfqpoint{11.515609in}{2.224058in}}{\pgfqpoint{11.516577in}{2.241726in}}%
\pgfpathquadraticcurveto{\pgfqpoint{11.512846in}{2.236181in}}{\pgfqpoint{11.507094in}{2.233212in}}%
\pgfpathquadraticcurveto{\pgfqpoint{11.501363in}{2.230243in}}{\pgfqpoint{11.494415in}{2.230243in}}%
\pgfpathquadraticcurveto{\pgfqpoint{11.479984in}{2.230243in}}{\pgfqpoint{11.471572in}{2.238964in}}%
\pgfpathquadraticcurveto{\pgfqpoint{11.463161in}{2.247705in}}{\pgfqpoint{11.463161in}{2.262858in}}%
\pgfpathquadraticcurveto{\pgfqpoint{11.463161in}{2.277661in}}{\pgfqpoint{11.471923in}{2.286608in}}%
\pgfpathquadraticcurveto{\pgfqpoint{11.480684in}{2.295576in}}{\pgfqpoint{11.495240in}{2.295576in}}%
\pgfpathquadraticcurveto{\pgfqpoint{11.511939in}{2.295576in}}{\pgfqpoint{11.520721in}{2.282774in}}%
\pgfpathquadraticcurveto{\pgfqpoint{11.529525in}{2.269991in}}{\pgfqpoint{11.529525in}{2.245644in}}%
\pgfpathquadraticcurveto{\pgfqpoint{11.529525in}{2.222904in}}{\pgfqpoint{11.518722in}{2.209338in}}%
\pgfpathquadraticcurveto{\pgfqpoint{11.507939in}{2.195773in}}{\pgfqpoint{11.489714in}{2.195773in}}%
\pgfpathquadraticcurveto{\pgfqpoint{11.484808in}{2.195773in}}{\pgfqpoint{11.479777in}{2.196742in}}%
\pgfpathquadraticcurveto{\pgfqpoint{11.474768in}{2.197711in}}{\pgfqpoint{11.469345in}{2.199649in}}%
\pgfpathlineto{\pgfqpoint{11.469345in}{2.199649in}}%
\pgfpathclose%
\pgfpathmoveto{\pgfqpoint{11.495240in}{2.240428in}}%
\pgfpathquadraticcurveto{\pgfqpoint{11.504002in}{2.240428in}}{\pgfqpoint{11.509114in}{2.246406in}}%
\pgfpathquadraticcurveto{\pgfqpoint{11.514248in}{2.252406in}}{\pgfqpoint{11.514248in}{2.262858in}}%
\pgfpathquadraticcurveto{\pgfqpoint{11.514248in}{2.273228in}}{\pgfqpoint{11.509114in}{2.279248in}}%
\pgfpathquadraticcurveto{\pgfqpoint{11.504002in}{2.285268in}}{\pgfqpoint{11.495240in}{2.285268in}}%
\pgfpathquadraticcurveto{\pgfqpoint{11.486478in}{2.285268in}}{\pgfqpoint{11.481365in}{2.279248in}}%
\pgfpathquadraticcurveto{\pgfqpoint{11.476252in}{2.273228in}}{\pgfqpoint{11.476252in}{2.262858in}}%
\pgfpathquadraticcurveto{\pgfqpoint{11.476252in}{2.252406in}}{\pgfqpoint{11.481365in}{2.246406in}}%
\pgfpathquadraticcurveto{\pgfqpoint{11.486478in}{2.240428in}}{\pgfqpoint{11.495240in}{2.240428in}}%
\pgfpathlineto{\pgfqpoint{11.495240in}{2.240428in}}%
\pgfpathclose%
\pgfpathmoveto{\pgfqpoint{11.621638in}{2.297762in}}%
\pgfpathquadraticcurveto{\pgfqpoint{11.613021in}{2.282959in}}{\pgfqpoint{11.608815in}{2.268445in}}%
\pgfpathquadraticcurveto{\pgfqpoint{11.604630in}{2.253952in}}{\pgfqpoint{11.604630in}{2.239067in}}%
\pgfpathquadraticcurveto{\pgfqpoint{11.604630in}{2.224203in}}{\pgfqpoint{11.608856in}{2.209606in}}%
\pgfpathquadraticcurveto{\pgfqpoint{11.613082in}{2.195010in}}{\pgfqpoint{11.621638in}{2.180249in}}%
\pgfpathlineto{\pgfqpoint{11.611330in}{2.180249in}}%
\pgfpathquadraticcurveto{\pgfqpoint{11.601682in}{2.195402in}}{\pgfqpoint{11.596878in}{2.210019in}}%
\pgfpathquadraticcurveto{\pgfqpoint{11.592074in}{2.224636in}}{\pgfqpoint{11.592074in}{2.239067in}}%
\pgfpathquadraticcurveto{\pgfqpoint{11.592074in}{2.253437in}}{\pgfqpoint{11.596837in}{2.267992in}}%
\pgfpathquadraticcurveto{\pgfqpoint{11.601620in}{2.282567in}}{\pgfqpoint{11.611330in}{2.297762in}}%
\pgfpathlineto{\pgfqpoint{11.621638in}{2.297762in}}%
\pgfpathlineto{\pgfqpoint{11.621638in}{2.297762in}}%
\pgfpathclose%
\pgfpathmoveto{\pgfqpoint{11.657531in}{2.208596in}}%
\pgfpathlineto{\pgfqpoint{11.702949in}{2.208596in}}%
\pgfpathlineto{\pgfqpoint{11.702949in}{2.197649in}}%
\pgfpathlineto{\pgfqpoint{11.641883in}{2.197649in}}%
\pgfpathlineto{\pgfqpoint{11.641883in}{2.208596in}}%
\pgfpathquadraticcurveto{\pgfqpoint{11.649285in}{2.216265in}}{\pgfqpoint{11.662067in}{2.229171in}}%
\pgfpathquadraticcurveto{\pgfqpoint{11.674870in}{2.242098in}}{\pgfqpoint{11.678148in}{2.245850in}}%
\pgfpathquadraticcurveto{\pgfqpoint{11.684394in}{2.252859in}}{\pgfqpoint{11.686868in}{2.257725in}}%
\pgfpathquadraticcurveto{\pgfqpoint{11.689363in}{2.262590in}}{\pgfqpoint{11.689363in}{2.267291in}}%
\pgfpathquadraticcurveto{\pgfqpoint{11.689363in}{2.274960in}}{\pgfqpoint{11.683982in}{2.279784in}}%
\pgfpathquadraticcurveto{\pgfqpoint{11.678601in}{2.284629in}}{\pgfqpoint{11.669963in}{2.284629in}}%
\pgfpathquadraticcurveto{\pgfqpoint{11.663840in}{2.284629in}}{\pgfqpoint{11.657036in}{2.282506in}}%
\pgfpathquadraticcurveto{\pgfqpoint{11.650254in}{2.280382in}}{\pgfqpoint{11.642523in}{2.276053in}}%
\pgfpathlineto{\pgfqpoint{11.642523in}{2.289206in}}%
\pgfpathquadraticcurveto{\pgfqpoint{11.650377in}{2.292360in}}{\pgfqpoint{11.657201in}{2.293968in}}%
\pgfpathquadraticcurveto{\pgfqpoint{11.664046in}{2.295576in}}{\pgfqpoint{11.669715in}{2.295576in}}%
\pgfpathquadraticcurveto{\pgfqpoint{11.684662in}{2.295576in}}{\pgfqpoint{11.693548in}{2.288093in}}%
\pgfpathquadraticcurveto{\pgfqpoint{11.702434in}{2.280629in}}{\pgfqpoint{11.702434in}{2.268136in}}%
\pgfpathquadraticcurveto{\pgfqpoint{11.702434in}{2.262198in}}{\pgfqpoint{11.700207in}{2.256879in}}%
\pgfpathquadraticcurveto{\pgfqpoint{11.698001in}{2.251581in}}{\pgfqpoint{11.692125in}{2.244365in}}%
\pgfpathquadraticcurveto{\pgfqpoint{11.690517in}{2.242489in}}{\pgfqpoint{11.681879in}{2.233562in}}%
\pgfpathquadraticcurveto{\pgfqpoint{11.673262in}{2.224636in}}{\pgfqpoint{11.657531in}{2.208596in}}%
\pgfpathlineto{\pgfqpoint{11.657531in}{2.208596in}}%
\pgfpathclose%
\pgfpathmoveto{\pgfqpoint{11.758098in}{2.285268in}}%
\pgfpathquadraticcurveto{\pgfqpoint{11.748058in}{2.285268in}}{\pgfqpoint{11.742986in}{2.275372in}}%
\pgfpathquadraticcurveto{\pgfqpoint{11.737935in}{2.265497in}}{\pgfqpoint{11.737935in}{2.245644in}}%
\pgfpathquadraticcurveto{\pgfqpoint{11.737935in}{2.225873in}}{\pgfqpoint{11.742986in}{2.215977in}}%
\pgfpathquadraticcurveto{\pgfqpoint{11.748058in}{2.206081in}}{\pgfqpoint{11.758098in}{2.206081in}}%
\pgfpathquadraticcurveto{\pgfqpoint{11.768220in}{2.206081in}}{\pgfqpoint{11.773271in}{2.215977in}}%
\pgfpathquadraticcurveto{\pgfqpoint{11.778343in}{2.225873in}}{\pgfqpoint{11.778343in}{2.245644in}}%
\pgfpathquadraticcurveto{\pgfqpoint{11.778343in}{2.265497in}}{\pgfqpoint{11.773271in}{2.275372in}}%
\pgfpathquadraticcurveto{\pgfqpoint{11.768220in}{2.285268in}}{\pgfqpoint{11.758098in}{2.285268in}}%
\pgfpathlineto{\pgfqpoint{11.758098in}{2.285268in}}%
\pgfpathclose%
\pgfpathmoveto{\pgfqpoint{11.758098in}{2.295576in}}%
\pgfpathquadraticcurveto{\pgfqpoint{11.774281in}{2.295576in}}{\pgfqpoint{11.782817in}{2.282774in}}%
\pgfpathquadraticcurveto{\pgfqpoint{11.791352in}{2.269991in}}{\pgfqpoint{11.791352in}{2.245644in}}%
\pgfpathquadraticcurveto{\pgfqpoint{11.791352in}{2.221358in}}{\pgfqpoint{11.782817in}{2.208555in}}%
\pgfpathquadraticcurveto{\pgfqpoint{11.774281in}{2.195773in}}{\pgfqpoint{11.758098in}{2.195773in}}%
\pgfpathquadraticcurveto{\pgfqpoint{11.741934in}{2.195773in}}{\pgfqpoint{11.733399in}{2.208555in}}%
\pgfpathquadraticcurveto{\pgfqpoint{11.724864in}{2.221358in}}{\pgfqpoint{11.724864in}{2.245644in}}%
\pgfpathquadraticcurveto{\pgfqpoint{11.724864in}{2.269991in}}{\pgfqpoint{11.733399in}{2.282774in}}%
\pgfpathquadraticcurveto{\pgfqpoint{11.741934in}{2.295576in}}{\pgfqpoint{11.758098in}{2.295576in}}%
\pgfpathlineto{\pgfqpoint{11.758098in}{2.295576in}}%
\pgfpathclose%
\pgfpathmoveto{\pgfqpoint{11.825431in}{2.208596in}}%
\pgfpathlineto{\pgfqpoint{11.870848in}{2.208596in}}%
\pgfpathlineto{\pgfqpoint{11.870848in}{2.197649in}}%
\pgfpathlineto{\pgfqpoint{11.809783in}{2.197649in}}%
\pgfpathlineto{\pgfqpoint{11.809783in}{2.208596in}}%
\pgfpathquadraticcurveto{\pgfqpoint{11.817184in}{2.216265in}}{\pgfqpoint{11.829966in}{2.229171in}}%
\pgfpathquadraticcurveto{\pgfqpoint{11.842769in}{2.242098in}}{\pgfqpoint{11.846047in}{2.245850in}}%
\pgfpathquadraticcurveto{\pgfqpoint{11.852294in}{2.252859in}}{\pgfqpoint{11.854768in}{2.257725in}}%
\pgfpathquadraticcurveto{\pgfqpoint{11.857262in}{2.262590in}}{\pgfqpoint{11.857262in}{2.267291in}}%
\pgfpathquadraticcurveto{\pgfqpoint{11.857262in}{2.274960in}}{\pgfqpoint{11.851881in}{2.279784in}}%
\pgfpathquadraticcurveto{\pgfqpoint{11.846500in}{2.284629in}}{\pgfqpoint{11.837862in}{2.284629in}}%
\pgfpathquadraticcurveto{\pgfqpoint{11.831739in}{2.284629in}}{\pgfqpoint{11.824936in}{2.282506in}}%
\pgfpathquadraticcurveto{\pgfqpoint{11.818153in}{2.280382in}}{\pgfqpoint{11.810422in}{2.276053in}}%
\pgfpathlineto{\pgfqpoint{11.810422in}{2.289206in}}%
\pgfpathquadraticcurveto{\pgfqpoint{11.818277in}{2.292360in}}{\pgfqpoint{11.825101in}{2.293968in}}%
\pgfpathquadraticcurveto{\pgfqpoint{11.831945in}{2.295576in}}{\pgfqpoint{11.837615in}{2.295576in}}%
\pgfpathquadraticcurveto{\pgfqpoint{11.852562in}{2.295576in}}{\pgfqpoint{11.861447in}{2.288093in}}%
\pgfpathquadraticcurveto{\pgfqpoint{11.870333in}{2.280629in}}{\pgfqpoint{11.870333in}{2.268136in}}%
\pgfpathquadraticcurveto{\pgfqpoint{11.870333in}{2.262198in}}{\pgfqpoint{11.868106in}{2.256879in}}%
\pgfpathquadraticcurveto{\pgfqpoint{11.865900in}{2.251581in}}{\pgfqpoint{11.860025in}{2.244365in}}%
\pgfpathquadraticcurveto{\pgfqpoint{11.858417in}{2.242489in}}{\pgfqpoint{11.849778in}{2.233562in}}%
\pgfpathquadraticcurveto{\pgfqpoint{11.841161in}{2.224636in}}{\pgfqpoint{11.825431in}{2.208596in}}%
\pgfpathlineto{\pgfqpoint{11.825431in}{2.208596in}}%
\pgfpathclose%
\pgfpathmoveto{\pgfqpoint{11.933934in}{2.282506in}}%
\pgfpathlineto{\pgfqpoint{11.901072in}{2.231150in}}%
\pgfpathlineto{\pgfqpoint{11.933934in}{2.231150in}}%
\pgfpathlineto{\pgfqpoint{11.933934in}{2.282506in}}%
\pgfpathlineto{\pgfqpoint{11.933934in}{2.282506in}}%
\pgfpathclose%
\pgfpathmoveto{\pgfqpoint{11.930512in}{2.293845in}}%
\pgfpathlineto{\pgfqpoint{11.946881in}{2.293845in}}%
\pgfpathlineto{\pgfqpoint{11.946881in}{2.231150in}}%
\pgfpathlineto{\pgfqpoint{11.960612in}{2.231150in}}%
\pgfpathlineto{\pgfqpoint{11.960612in}{2.220327in}}%
\pgfpathlineto{\pgfqpoint{11.946881in}{2.220327in}}%
\pgfpathlineto{\pgfqpoint{11.946881in}{2.197649in}}%
\pgfpathlineto{\pgfqpoint{11.933934in}{2.197649in}}%
\pgfpathlineto{\pgfqpoint{11.933934in}{2.220327in}}%
\pgfpathlineto{\pgfqpoint{11.890516in}{2.220327in}}%
\pgfpathlineto{\pgfqpoint{11.890516in}{2.232882in}}%
\pgfpathlineto{\pgfqpoint{11.930512in}{2.293845in}}%
\pgfpathlineto{\pgfqpoint{11.930512in}{2.293845in}}%
\pgfpathclose%
\pgfpathmoveto{\pgfqpoint{11.978589in}{2.297762in}}%
\pgfpathlineto{\pgfqpoint{11.988897in}{2.297762in}}%
\pgfpathquadraticcurveto{\pgfqpoint{11.998546in}{2.282567in}}{\pgfqpoint{12.003349in}{2.267992in}}%
\pgfpathquadraticcurveto{\pgfqpoint{12.008153in}{2.253437in}}{\pgfqpoint{12.008153in}{2.239067in}}%
\pgfpathquadraticcurveto{\pgfqpoint{12.008153in}{2.224636in}}{\pgfqpoint{12.003349in}{2.210019in}}%
\pgfpathquadraticcurveto{\pgfqpoint{11.998546in}{2.195402in}}{\pgfqpoint{11.988897in}{2.180249in}}%
\pgfpathlineto{\pgfqpoint{11.978589in}{2.180249in}}%
\pgfpathquadraticcurveto{\pgfqpoint{11.987145in}{2.195010in}}{\pgfqpoint{11.991371in}{2.209606in}}%
\pgfpathquadraticcurveto{\pgfqpoint{11.995598in}{2.224203in}}{\pgfqpoint{11.995598in}{2.239067in}}%
\pgfpathquadraticcurveto{\pgfqpoint{11.995598in}{2.253952in}}{\pgfqpoint{11.991371in}{2.268445in}}%
\pgfpathquadraticcurveto{\pgfqpoint{11.987145in}{2.282959in}}{\pgfqpoint{11.978589in}{2.297762in}}%
\pgfpathlineto{\pgfqpoint{11.978589in}{2.297762in}}%
\pgfpathclose%
\pgfusepath{fill}%
\end{pgfscope}%
\begin{pgfscope}%
\pgfsetbuttcap%
\pgfsetroundjoin%
\definecolor{currentfill}{rgb}{0.090196,0.168627,0.227451}%
\pgfsetfillcolor{currentfill}%
\pgfsetlinewidth{0.000000pt}%
\definecolor{currentstroke}{rgb}{0.090196,0.168627,0.227451}%
\pgfsetstrokecolor{currentstroke}%
\pgfsetdash{}{0pt}%
\pgfpathmoveto{\pgfqpoint{11.018611in}{2.114050in}}%
\pgfpathlineto{\pgfqpoint{11.018611in}{2.101351in}}%
\pgfpathquadraticcurveto{\pgfqpoint{11.011210in}{2.104897in}}{\pgfqpoint{11.004633in}{2.106628in}}%
\pgfpathquadraticcurveto{\pgfqpoint{10.998056in}{2.108381in}}{\pgfqpoint{10.991933in}{2.108381in}}%
\pgfpathquadraticcurveto{\pgfqpoint{10.981316in}{2.108381in}}{\pgfqpoint{10.975543in}{2.104257in}}%
\pgfpathquadraticcurveto{\pgfqpoint{10.969771in}{2.100134in}}{\pgfqpoint{10.969771in}{2.092527in}}%
\pgfpathquadraticcurveto{\pgfqpoint{10.969771in}{2.086136in}}{\pgfqpoint{10.973605in}{2.082878in}}%
\pgfpathquadraticcurveto{\pgfqpoint{10.977440in}{2.079642in}}{\pgfqpoint{10.988140in}{2.077642in}}%
\pgfpathlineto{\pgfqpoint{10.995995in}{2.076034in}}%
\pgfpathquadraticcurveto{\pgfqpoint{11.010550in}{2.073251in}}{\pgfqpoint{11.017477in}{2.066262in}}%
\pgfpathquadraticcurveto{\pgfqpoint{11.024404in}{2.059273in}}{\pgfqpoint{11.024404in}{2.047563in}}%
\pgfpathquadraticcurveto{\pgfqpoint{11.024404in}{2.033564in}}{\pgfqpoint{11.015024in}{2.026348in}}%
\pgfpathquadraticcurveto{\pgfqpoint{11.005664in}{2.019133in}}{\pgfqpoint{10.987563in}{2.019133in}}%
\pgfpathquadraticcurveto{\pgfqpoint{10.980739in}{2.019133in}}{\pgfqpoint{10.973028in}{2.020679in}}%
\pgfpathquadraticcurveto{\pgfqpoint{10.965338in}{2.022225in}}{\pgfqpoint{10.957092in}{2.025256in}}%
\pgfpathlineto{\pgfqpoint{10.957092in}{2.038656in}}%
\pgfpathquadraticcurveto{\pgfqpoint{10.965008in}{2.034224in}}{\pgfqpoint{10.972616in}{2.031956in}}%
\pgfpathquadraticcurveto{\pgfqpoint{10.980223in}{2.029709in}}{\pgfqpoint{10.987563in}{2.029709in}}%
\pgfpathquadraticcurveto{\pgfqpoint{10.998696in}{2.029709in}}{\pgfqpoint{11.004757in}{2.034080in}}%
\pgfpathquadraticcurveto{\pgfqpoint{11.010818in}{2.038471in}}{\pgfqpoint{11.010818in}{2.046594in}}%
\pgfpathquadraticcurveto{\pgfqpoint{11.010818in}{2.053665in}}{\pgfqpoint{11.006468in}{2.057665in}}%
\pgfpathquadraticcurveto{\pgfqpoint{11.002118in}{2.061664in}}{\pgfqpoint{10.992201in}{2.063664in}}%
\pgfpathlineto{\pgfqpoint{10.984264in}{2.065210in}}%
\pgfpathquadraticcurveto{\pgfqpoint{10.969709in}{2.068096in}}{\pgfqpoint{10.963194in}{2.074281in}}%
\pgfpathquadraticcurveto{\pgfqpoint{10.956700in}{2.080466in}}{\pgfqpoint{10.956700in}{2.091496in}}%
\pgfpathquadraticcurveto{\pgfqpoint{10.956700in}{2.104257in}}{\pgfqpoint{10.965689in}{2.111597in}}%
\pgfpathquadraticcurveto{\pgfqpoint{10.974678in}{2.118936in}}{\pgfqpoint{10.990449in}{2.118936in}}%
\pgfpathquadraticcurveto{\pgfqpoint{10.997232in}{2.118936in}}{\pgfqpoint{11.004241in}{2.117699in}}%
\pgfpathquadraticcurveto{\pgfqpoint{11.011271in}{2.116483in}}{\pgfqpoint{11.018611in}{2.114050in}}%
\pgfpathlineto{\pgfqpoint{11.018611in}{2.114050in}}%
\pgfpathclose%
\pgfpathmoveto{\pgfqpoint{11.055926in}{2.113659in}}%
\pgfpathlineto{\pgfqpoint{11.055926in}{2.093166in}}%
\pgfpathlineto{\pgfqpoint{11.080336in}{2.093166in}}%
\pgfpathlineto{\pgfqpoint{11.080336in}{2.083950in}}%
\pgfpathlineto{\pgfqpoint{11.055926in}{2.083950in}}%
\pgfpathlineto{\pgfqpoint{11.055926in}{2.044779in}}%
\pgfpathquadraticcurveto{\pgfqpoint{11.055926in}{2.035956in}}{\pgfqpoint{11.058339in}{2.033440in}}%
\pgfpathquadraticcurveto{\pgfqpoint{11.060751in}{2.030925in}}{\pgfqpoint{11.068173in}{2.030925in}}%
\pgfpathlineto{\pgfqpoint{11.080336in}{2.030925in}}%
\pgfpathlineto{\pgfqpoint{11.080336in}{2.021009in}}%
\pgfpathlineto{\pgfqpoint{11.068173in}{2.021009in}}%
\pgfpathquadraticcurveto{\pgfqpoint{11.054442in}{2.021009in}}{\pgfqpoint{11.049226in}{2.026122in}}%
\pgfpathquadraticcurveto{\pgfqpoint{11.044010in}{2.031255in}}{\pgfqpoint{11.044010in}{2.044779in}}%
\pgfpathlineto{\pgfqpoint{11.044010in}{2.083950in}}%
\pgfpathlineto{\pgfqpoint{11.035310in}{2.083950in}}%
\pgfpathlineto{\pgfqpoint{11.035310in}{2.093166in}}%
\pgfpathlineto{\pgfqpoint{11.044010in}{2.093166in}}%
\pgfpathlineto{\pgfqpoint{11.044010in}{2.113659in}}%
\pgfpathlineto{\pgfqpoint{11.055926in}{2.113659in}}%
\pgfpathlineto{\pgfqpoint{11.055926in}{2.113659in}}%
\pgfpathclose%
\pgfpathmoveto{\pgfqpoint{11.094706in}{2.049480in}}%
\pgfpathlineto{\pgfqpoint{11.094706in}{2.093166in}}%
\pgfpathlineto{\pgfqpoint{11.106560in}{2.093166in}}%
\pgfpathlineto{\pgfqpoint{11.106560in}{2.049933in}}%
\pgfpathquadraticcurveto{\pgfqpoint{11.106560in}{2.039687in}}{\pgfqpoint{11.110539in}{2.034554in}}%
\pgfpathquadraticcurveto{\pgfqpoint{11.114539in}{2.029441in}}{\pgfqpoint{11.122538in}{2.029441in}}%
\pgfpathquadraticcurveto{\pgfqpoint{11.132124in}{2.029441in}}{\pgfqpoint{11.137691in}{2.035564in}}%
\pgfpathquadraticcurveto{\pgfqpoint{11.143278in}{2.041687in}}{\pgfqpoint{11.143278in}{2.052263in}}%
\pgfpathlineto{\pgfqpoint{11.143278in}{2.093166in}}%
\pgfpathlineto{\pgfqpoint{11.155132in}{2.093166in}}%
\pgfpathlineto{\pgfqpoint{11.155132in}{2.021009in}}%
\pgfpathlineto{\pgfqpoint{11.143278in}{2.021009in}}%
\pgfpathlineto{\pgfqpoint{11.143278in}{2.032100in}}%
\pgfpathquadraticcurveto{\pgfqpoint{11.138969in}{2.025524in}}{\pgfqpoint{11.133258in}{2.022328in}}%
\pgfpathquadraticcurveto{\pgfqpoint{11.127568in}{2.019133in}}{\pgfqpoint{11.120023in}{2.019133in}}%
\pgfpathquadraticcurveto{\pgfqpoint{11.107591in}{2.019133in}}{\pgfqpoint{11.101138in}{2.026864in}}%
\pgfpathquadraticcurveto{\pgfqpoint{11.094706in}{2.034595in}}{\pgfqpoint{11.094706in}{2.049480in}}%
\pgfpathlineto{\pgfqpoint{11.094706in}{2.049480in}}%
\pgfpathclose%
\pgfpathmoveto{\pgfqpoint{11.124538in}{2.094898in}}%
\pgfpathlineto{\pgfqpoint{11.124538in}{2.094898in}}%
\pgfpathlineto{\pgfqpoint{11.124538in}{2.094898in}}%
\pgfpathclose%
\pgfpathmoveto{\pgfqpoint{11.227021in}{2.082219in}}%
\pgfpathlineto{\pgfqpoint{11.227021in}{2.121266in}}%
\pgfpathlineto{\pgfqpoint{11.238876in}{2.121266in}}%
\pgfpathlineto{\pgfqpoint{11.238876in}{2.021009in}}%
\pgfpathlineto{\pgfqpoint{11.227021in}{2.021009in}}%
\pgfpathlineto{\pgfqpoint{11.227021in}{2.031832in}}%
\pgfpathquadraticcurveto{\pgfqpoint{11.223290in}{2.025400in}}{\pgfqpoint{11.217579in}{2.022266in}}%
\pgfpathquadraticcurveto{\pgfqpoint{11.211889in}{2.019133in}}{\pgfqpoint{11.203890in}{2.019133in}}%
\pgfpathquadraticcurveto{\pgfqpoint{11.190819in}{2.019133in}}{\pgfqpoint{11.182593in}{2.029565in}}%
\pgfpathquadraticcurveto{\pgfqpoint{11.174388in}{2.040017in}}{\pgfqpoint{11.174388in}{2.057026in}}%
\pgfpathquadraticcurveto{\pgfqpoint{11.174388in}{2.074034in}}{\pgfqpoint{11.182593in}{2.084466in}}%
\pgfpathquadraticcurveto{\pgfqpoint{11.190819in}{2.094898in}}{\pgfqpoint{11.203890in}{2.094898in}}%
\pgfpathquadraticcurveto{\pgfqpoint{11.211889in}{2.094898in}}{\pgfqpoint{11.217579in}{2.091764in}}%
\pgfpathquadraticcurveto{\pgfqpoint{11.223290in}{2.088651in}}{\pgfqpoint{11.227021in}{2.082219in}}%
\pgfpathlineto{\pgfqpoint{11.227021in}{2.082219in}}%
\pgfpathclose%
\pgfpathmoveto{\pgfqpoint{11.186634in}{2.057026in}}%
\pgfpathquadraticcurveto{\pgfqpoint{11.186634in}{2.043955in}}{\pgfqpoint{11.192015in}{2.036512in}}%
\pgfpathquadraticcurveto{\pgfqpoint{11.197396in}{2.029070in}}{\pgfqpoint{11.206797in}{2.029070in}}%
\pgfpathquadraticcurveto{\pgfqpoint{11.216198in}{2.029070in}}{\pgfqpoint{11.221599in}{2.036512in}}%
\pgfpathquadraticcurveto{\pgfqpoint{11.227021in}{2.043955in}}{\pgfqpoint{11.227021in}{2.057026in}}%
\pgfpathquadraticcurveto{\pgfqpoint{11.227021in}{2.070096in}}{\pgfqpoint{11.221599in}{2.077539in}}%
\pgfpathquadraticcurveto{\pgfqpoint{11.216198in}{2.084981in}}{\pgfqpoint{11.206797in}{2.084981in}}%
\pgfpathquadraticcurveto{\pgfqpoint{11.197396in}{2.084981in}}{\pgfqpoint{11.192015in}{2.077539in}}%
\pgfpathquadraticcurveto{\pgfqpoint{11.186634in}{2.070096in}}{\pgfqpoint{11.186634in}{2.057026in}}%
\pgfpathlineto{\pgfqpoint{11.186634in}{2.057026in}}%
\pgfpathclose%
\pgfpathmoveto{\pgfqpoint{11.293323in}{2.014308in}}%
\pgfpathquadraticcurveto{\pgfqpoint{11.288314in}{2.001423in}}{\pgfqpoint{11.283531in}{1.997506in}}%
\pgfpathquadraticcurveto{\pgfqpoint{11.278768in}{1.993568in}}{\pgfqpoint{11.270790in}{1.993568in}}%
\pgfpathlineto{\pgfqpoint{11.261306in}{1.993568in}}%
\pgfpathlineto{\pgfqpoint{11.261306in}{2.003485in}}%
\pgfpathlineto{\pgfqpoint{11.268275in}{2.003485in}}%
\pgfpathquadraticcurveto{\pgfqpoint{11.273161in}{2.003485in}}{\pgfqpoint{11.275861in}{2.005815in}}%
\pgfpathquadraticcurveto{\pgfqpoint{11.278583in}{2.008124in}}{\pgfqpoint{11.281861in}{2.016762in}}%
\pgfpathlineto{\pgfqpoint{11.283984in}{2.022163in}}%
\pgfpathlineto{\pgfqpoint{11.254812in}{2.093166in}}%
\pgfpathlineto{\pgfqpoint{11.267367in}{2.093166in}}%
\pgfpathlineto{\pgfqpoint{11.289922in}{2.036739in}}%
\pgfpathlineto{\pgfqpoint{11.312476in}{2.093166in}}%
\pgfpathlineto{\pgfqpoint{11.325031in}{2.093166in}}%
\pgfpathlineto{\pgfqpoint{11.293323in}{2.014308in}}%
\pgfpathlineto{\pgfqpoint{11.293323in}{2.014308in}}%
\pgfpathclose%
\pgfpathmoveto{\pgfqpoint{11.396219in}{2.031956in}}%
\pgfpathlineto{\pgfqpoint{11.441637in}{2.031956in}}%
\pgfpathlineto{\pgfqpoint{11.441637in}{2.021009in}}%
\pgfpathlineto{\pgfqpoint{11.380572in}{2.021009in}}%
\pgfpathlineto{\pgfqpoint{11.380572in}{2.031956in}}%
\pgfpathquadraticcurveto{\pgfqpoint{11.387973in}{2.039625in}}{\pgfqpoint{11.400755in}{2.052531in}}%
\pgfpathquadraticcurveto{\pgfqpoint{11.413558in}{2.065458in}}{\pgfqpoint{11.416836in}{2.069210in}}%
\pgfpathquadraticcurveto{\pgfqpoint{11.423082in}{2.076219in}}{\pgfqpoint{11.425556in}{2.081085in}}%
\pgfpathquadraticcurveto{\pgfqpoint{11.428051in}{2.085950in}}{\pgfqpoint{11.428051in}{2.090651in}}%
\pgfpathquadraticcurveto{\pgfqpoint{11.428051in}{2.098320in}}{\pgfqpoint{11.422670in}{2.103144in}}%
\pgfpathquadraticcurveto{\pgfqpoint{11.417289in}{2.107989in}}{\pgfqpoint{11.408651in}{2.107989in}}%
\pgfpathquadraticcurveto{\pgfqpoint{11.402528in}{2.107989in}}{\pgfqpoint{11.395725in}{2.105866in}}%
\pgfpathquadraticcurveto{\pgfqpoint{11.388942in}{2.103742in}}{\pgfqpoint{11.381211in}{2.099413in}}%
\pgfpathlineto{\pgfqpoint{11.381211in}{2.112566in}}%
\pgfpathquadraticcurveto{\pgfqpoint{11.389066in}{2.115720in}}{\pgfqpoint{11.395890in}{2.117328in}}%
\pgfpathquadraticcurveto{\pgfqpoint{11.402734in}{2.118936in}}{\pgfqpoint{11.408404in}{2.118936in}}%
\pgfpathquadraticcurveto{\pgfqpoint{11.423350in}{2.118936in}}{\pgfqpoint{11.432236in}{2.111453in}}%
\pgfpathquadraticcurveto{\pgfqpoint{11.441122in}{2.103989in}}{\pgfqpoint{11.441122in}{2.091496in}}%
\pgfpathquadraticcurveto{\pgfqpoint{11.441122in}{2.085558in}}{\pgfqpoint{11.438895in}{2.080239in}}%
\pgfpathquadraticcurveto{\pgfqpoint{11.436689in}{2.074941in}}{\pgfqpoint{11.430814in}{2.067725in}}%
\pgfpathquadraticcurveto{\pgfqpoint{11.429206in}{2.065849in}}{\pgfqpoint{11.420567in}{2.056922in}}%
\pgfpathquadraticcurveto{\pgfqpoint{11.411950in}{2.047996in}}{\pgfqpoint{11.396219in}{2.031956in}}%
\pgfpathlineto{\pgfqpoint{11.396219in}{2.031956in}}%
\pgfpathclose%
\pgfpathmoveto{\pgfqpoint{11.496786in}{2.108628in}}%
\pgfpathquadraticcurveto{\pgfqpoint{11.486746in}{2.108628in}}{\pgfqpoint{11.481674in}{2.098732in}}%
\pgfpathquadraticcurveto{\pgfqpoint{11.476623in}{2.088857in}}{\pgfqpoint{11.476623in}{2.069004in}}%
\pgfpathquadraticcurveto{\pgfqpoint{11.476623in}{2.049233in}}{\pgfqpoint{11.481674in}{2.039337in}}%
\pgfpathquadraticcurveto{\pgfqpoint{11.486746in}{2.029441in}}{\pgfqpoint{11.496786in}{2.029441in}}%
\pgfpathquadraticcurveto{\pgfqpoint{11.506908in}{2.029441in}}{\pgfqpoint{11.511959in}{2.039337in}}%
\pgfpathquadraticcurveto{\pgfqpoint{11.517031in}{2.049233in}}{\pgfqpoint{11.517031in}{2.069004in}}%
\pgfpathquadraticcurveto{\pgfqpoint{11.517031in}{2.088857in}}{\pgfqpoint{11.511959in}{2.098732in}}%
\pgfpathquadraticcurveto{\pgfqpoint{11.506908in}{2.108628in}}{\pgfqpoint{11.496786in}{2.108628in}}%
\pgfpathlineto{\pgfqpoint{11.496786in}{2.108628in}}%
\pgfpathclose%
\pgfpathmoveto{\pgfqpoint{11.496786in}{2.118936in}}%
\pgfpathquadraticcurveto{\pgfqpoint{11.512970in}{2.118936in}}{\pgfqpoint{11.521505in}{2.106134in}}%
\pgfpathquadraticcurveto{\pgfqpoint{11.530040in}{2.093351in}}{\pgfqpoint{11.530040in}{2.069004in}}%
\pgfpathquadraticcurveto{\pgfqpoint{11.530040in}{2.044718in}}{\pgfqpoint{11.521505in}{2.031915in}}%
\pgfpathquadraticcurveto{\pgfqpoint{11.512970in}{2.019133in}}{\pgfqpoint{11.496786in}{2.019133in}}%
\pgfpathquadraticcurveto{\pgfqpoint{11.480623in}{2.019133in}}{\pgfqpoint{11.472087in}{2.031915in}}%
\pgfpathquadraticcurveto{\pgfqpoint{11.463552in}{2.044718in}}{\pgfqpoint{11.463552in}{2.069004in}}%
\pgfpathquadraticcurveto{\pgfqpoint{11.463552in}{2.093351in}}{\pgfqpoint{11.472087in}{2.106134in}}%
\pgfpathquadraticcurveto{\pgfqpoint{11.480623in}{2.118936in}}{\pgfqpoint{11.496786in}{2.118936in}}%
\pgfpathlineto{\pgfqpoint{11.496786in}{2.118936in}}%
\pgfpathclose%
\pgfpathmoveto{\pgfqpoint{11.621638in}{2.121122in}}%
\pgfpathquadraticcurveto{\pgfqpoint{11.613021in}{2.106319in}}{\pgfqpoint{11.608815in}{2.091805in}}%
\pgfpathquadraticcurveto{\pgfqpoint{11.604630in}{2.077312in}}{\pgfqpoint{11.604630in}{2.062427in}}%
\pgfpathquadraticcurveto{\pgfqpoint{11.604630in}{2.047563in}}{\pgfqpoint{11.608856in}{2.032966in}}%
\pgfpathquadraticcurveto{\pgfqpoint{11.613082in}{2.018370in}}{\pgfqpoint{11.621638in}{2.003609in}}%
\pgfpathlineto{\pgfqpoint{11.611330in}{2.003609in}}%
\pgfpathquadraticcurveto{\pgfqpoint{11.601682in}{2.018762in}}{\pgfqpoint{11.596878in}{2.033379in}}%
\pgfpathquadraticcurveto{\pgfqpoint{11.592074in}{2.047996in}}{\pgfqpoint{11.592074in}{2.062427in}}%
\pgfpathquadraticcurveto{\pgfqpoint{11.592074in}{2.076797in}}{\pgfqpoint{11.596837in}{2.091352in}}%
\pgfpathquadraticcurveto{\pgfqpoint{11.601620in}{2.105927in}}{\pgfqpoint{11.611330in}{2.121122in}}%
\pgfpathlineto{\pgfqpoint{11.621638in}{2.121122in}}%
\pgfpathlineto{\pgfqpoint{11.621638in}{2.121122in}}%
\pgfpathclose%
\pgfpathmoveto{\pgfqpoint{11.657531in}{2.031956in}}%
\pgfpathlineto{\pgfqpoint{11.702949in}{2.031956in}}%
\pgfpathlineto{\pgfqpoint{11.702949in}{2.021009in}}%
\pgfpathlineto{\pgfqpoint{11.641883in}{2.021009in}}%
\pgfpathlineto{\pgfqpoint{11.641883in}{2.031956in}}%
\pgfpathquadraticcurveto{\pgfqpoint{11.649285in}{2.039625in}}{\pgfqpoint{11.662067in}{2.052531in}}%
\pgfpathquadraticcurveto{\pgfqpoint{11.674870in}{2.065458in}}{\pgfqpoint{11.678148in}{2.069210in}}%
\pgfpathquadraticcurveto{\pgfqpoint{11.684394in}{2.076219in}}{\pgfqpoint{11.686868in}{2.081085in}}%
\pgfpathquadraticcurveto{\pgfqpoint{11.689363in}{2.085950in}}{\pgfqpoint{11.689363in}{2.090651in}}%
\pgfpathquadraticcurveto{\pgfqpoint{11.689363in}{2.098320in}}{\pgfqpoint{11.683982in}{2.103144in}}%
\pgfpathquadraticcurveto{\pgfqpoint{11.678601in}{2.107989in}}{\pgfqpoint{11.669963in}{2.107989in}}%
\pgfpathquadraticcurveto{\pgfqpoint{11.663840in}{2.107989in}}{\pgfqpoint{11.657036in}{2.105866in}}%
\pgfpathquadraticcurveto{\pgfqpoint{11.650254in}{2.103742in}}{\pgfqpoint{11.642523in}{2.099413in}}%
\pgfpathlineto{\pgfqpoint{11.642523in}{2.112566in}}%
\pgfpathquadraticcurveto{\pgfqpoint{11.650377in}{2.115720in}}{\pgfqpoint{11.657201in}{2.117328in}}%
\pgfpathquadraticcurveto{\pgfqpoint{11.664046in}{2.118936in}}{\pgfqpoint{11.669715in}{2.118936in}}%
\pgfpathquadraticcurveto{\pgfqpoint{11.684662in}{2.118936in}}{\pgfqpoint{11.693548in}{2.111453in}}%
\pgfpathquadraticcurveto{\pgfqpoint{11.702434in}{2.103989in}}{\pgfqpoint{11.702434in}{2.091496in}}%
\pgfpathquadraticcurveto{\pgfqpoint{11.702434in}{2.085558in}}{\pgfqpoint{11.700207in}{2.080239in}}%
\pgfpathquadraticcurveto{\pgfqpoint{11.698001in}{2.074941in}}{\pgfqpoint{11.692125in}{2.067725in}}%
\pgfpathquadraticcurveto{\pgfqpoint{11.690517in}{2.065849in}}{\pgfqpoint{11.681879in}{2.056922in}}%
\pgfpathquadraticcurveto{\pgfqpoint{11.673262in}{2.047996in}}{\pgfqpoint{11.657531in}{2.031956in}}%
\pgfpathlineto{\pgfqpoint{11.657531in}{2.031956in}}%
\pgfpathclose%
\pgfpathmoveto{\pgfqpoint{11.758098in}{2.108628in}}%
\pgfpathquadraticcurveto{\pgfqpoint{11.748058in}{2.108628in}}{\pgfqpoint{11.742986in}{2.098732in}}%
\pgfpathquadraticcurveto{\pgfqpoint{11.737935in}{2.088857in}}{\pgfqpoint{11.737935in}{2.069004in}}%
\pgfpathquadraticcurveto{\pgfqpoint{11.737935in}{2.049233in}}{\pgfqpoint{11.742986in}{2.039337in}}%
\pgfpathquadraticcurveto{\pgfqpoint{11.748058in}{2.029441in}}{\pgfqpoint{11.758098in}{2.029441in}}%
\pgfpathquadraticcurveto{\pgfqpoint{11.768220in}{2.029441in}}{\pgfqpoint{11.773271in}{2.039337in}}%
\pgfpathquadraticcurveto{\pgfqpoint{11.778343in}{2.049233in}}{\pgfqpoint{11.778343in}{2.069004in}}%
\pgfpathquadraticcurveto{\pgfqpoint{11.778343in}{2.088857in}}{\pgfqpoint{11.773271in}{2.098732in}}%
\pgfpathquadraticcurveto{\pgfqpoint{11.768220in}{2.108628in}}{\pgfqpoint{11.758098in}{2.108628in}}%
\pgfpathlineto{\pgfqpoint{11.758098in}{2.108628in}}%
\pgfpathclose%
\pgfpathmoveto{\pgfqpoint{11.758098in}{2.118936in}}%
\pgfpathquadraticcurveto{\pgfqpoint{11.774281in}{2.118936in}}{\pgfqpoint{11.782817in}{2.106134in}}%
\pgfpathquadraticcurveto{\pgfqpoint{11.791352in}{2.093351in}}{\pgfqpoint{11.791352in}{2.069004in}}%
\pgfpathquadraticcurveto{\pgfqpoint{11.791352in}{2.044718in}}{\pgfqpoint{11.782817in}{2.031915in}}%
\pgfpathquadraticcurveto{\pgfqpoint{11.774281in}{2.019133in}}{\pgfqpoint{11.758098in}{2.019133in}}%
\pgfpathquadraticcurveto{\pgfqpoint{11.741934in}{2.019133in}}{\pgfqpoint{11.733399in}{2.031915in}}%
\pgfpathquadraticcurveto{\pgfqpoint{11.724864in}{2.044718in}}{\pgfqpoint{11.724864in}{2.069004in}}%
\pgfpathquadraticcurveto{\pgfqpoint{11.724864in}{2.093351in}}{\pgfqpoint{11.733399in}{2.106134in}}%
\pgfpathquadraticcurveto{\pgfqpoint{11.741934in}{2.118936in}}{\pgfqpoint{11.758098in}{2.118936in}}%
\pgfpathlineto{\pgfqpoint{11.758098in}{2.118936in}}%
\pgfpathclose%
\pgfpathmoveto{\pgfqpoint{11.825431in}{2.031956in}}%
\pgfpathlineto{\pgfqpoint{11.870848in}{2.031956in}}%
\pgfpathlineto{\pgfqpoint{11.870848in}{2.021009in}}%
\pgfpathlineto{\pgfqpoint{11.809783in}{2.021009in}}%
\pgfpathlineto{\pgfqpoint{11.809783in}{2.031956in}}%
\pgfpathquadraticcurveto{\pgfqpoint{11.817184in}{2.039625in}}{\pgfqpoint{11.829966in}{2.052531in}}%
\pgfpathquadraticcurveto{\pgfqpoint{11.842769in}{2.065458in}}{\pgfqpoint{11.846047in}{2.069210in}}%
\pgfpathquadraticcurveto{\pgfqpoint{11.852294in}{2.076219in}}{\pgfqpoint{11.854768in}{2.081085in}}%
\pgfpathquadraticcurveto{\pgfqpoint{11.857262in}{2.085950in}}{\pgfqpoint{11.857262in}{2.090651in}}%
\pgfpathquadraticcurveto{\pgfqpoint{11.857262in}{2.098320in}}{\pgfqpoint{11.851881in}{2.103144in}}%
\pgfpathquadraticcurveto{\pgfqpoint{11.846500in}{2.107989in}}{\pgfqpoint{11.837862in}{2.107989in}}%
\pgfpathquadraticcurveto{\pgfqpoint{11.831739in}{2.107989in}}{\pgfqpoint{11.824936in}{2.105866in}}%
\pgfpathquadraticcurveto{\pgfqpoint{11.818153in}{2.103742in}}{\pgfqpoint{11.810422in}{2.099413in}}%
\pgfpathlineto{\pgfqpoint{11.810422in}{2.112566in}}%
\pgfpathquadraticcurveto{\pgfqpoint{11.818277in}{2.115720in}}{\pgfqpoint{11.825101in}{2.117328in}}%
\pgfpathquadraticcurveto{\pgfqpoint{11.831945in}{2.118936in}}{\pgfqpoint{11.837615in}{2.118936in}}%
\pgfpathquadraticcurveto{\pgfqpoint{11.852562in}{2.118936in}}{\pgfqpoint{11.861447in}{2.111453in}}%
\pgfpathquadraticcurveto{\pgfqpoint{11.870333in}{2.103989in}}{\pgfqpoint{11.870333in}{2.091496in}}%
\pgfpathquadraticcurveto{\pgfqpoint{11.870333in}{2.085558in}}{\pgfqpoint{11.868106in}{2.080239in}}%
\pgfpathquadraticcurveto{\pgfqpoint{11.865900in}{2.074941in}}{\pgfqpoint{11.860025in}{2.067725in}}%
\pgfpathquadraticcurveto{\pgfqpoint{11.858417in}{2.065849in}}{\pgfqpoint{11.849778in}{2.056922in}}%
\pgfpathquadraticcurveto{\pgfqpoint{11.841161in}{2.047996in}}{\pgfqpoint{11.825431in}{2.031956in}}%
\pgfpathlineto{\pgfqpoint{11.825431in}{2.031956in}}%
\pgfpathclose%
\pgfpathmoveto{\pgfqpoint{11.933934in}{2.105866in}}%
\pgfpathlineto{\pgfqpoint{11.901072in}{2.054510in}}%
\pgfpathlineto{\pgfqpoint{11.933934in}{2.054510in}}%
\pgfpathlineto{\pgfqpoint{11.933934in}{2.105866in}}%
\pgfpathlineto{\pgfqpoint{11.933934in}{2.105866in}}%
\pgfpathclose%
\pgfpathmoveto{\pgfqpoint{11.930512in}{2.117205in}}%
\pgfpathlineto{\pgfqpoint{11.946881in}{2.117205in}}%
\pgfpathlineto{\pgfqpoint{11.946881in}{2.054510in}}%
\pgfpathlineto{\pgfqpoint{11.960612in}{2.054510in}}%
\pgfpathlineto{\pgfqpoint{11.960612in}{2.043687in}}%
\pgfpathlineto{\pgfqpoint{11.946881in}{2.043687in}}%
\pgfpathlineto{\pgfqpoint{11.946881in}{2.021009in}}%
\pgfpathlineto{\pgfqpoint{11.933934in}{2.021009in}}%
\pgfpathlineto{\pgfqpoint{11.933934in}{2.043687in}}%
\pgfpathlineto{\pgfqpoint{11.890516in}{2.043687in}}%
\pgfpathlineto{\pgfqpoint{11.890516in}{2.056242in}}%
\pgfpathlineto{\pgfqpoint{11.930512in}{2.117205in}}%
\pgfpathlineto{\pgfqpoint{11.930512in}{2.117205in}}%
\pgfpathclose%
\pgfpathmoveto{\pgfqpoint{11.978589in}{2.121122in}}%
\pgfpathlineto{\pgfqpoint{11.988897in}{2.121122in}}%
\pgfpathquadraticcurveto{\pgfqpoint{11.998546in}{2.105927in}}{\pgfqpoint{12.003349in}{2.091352in}}%
\pgfpathquadraticcurveto{\pgfqpoint{12.008153in}{2.076797in}}{\pgfqpoint{12.008153in}{2.062427in}}%
\pgfpathquadraticcurveto{\pgfqpoint{12.008153in}{2.047996in}}{\pgfqpoint{12.003349in}{2.033379in}}%
\pgfpathquadraticcurveto{\pgfqpoint{11.998546in}{2.018762in}}{\pgfqpoint{11.988897in}{2.003609in}}%
\pgfpathlineto{\pgfqpoint{11.978589in}{2.003609in}}%
\pgfpathquadraticcurveto{\pgfqpoint{11.987145in}{2.018370in}}{\pgfqpoint{11.991371in}{2.032966in}}%
\pgfpathquadraticcurveto{\pgfqpoint{11.995598in}{2.047563in}}{\pgfqpoint{11.995598in}{2.062427in}}%
\pgfpathquadraticcurveto{\pgfqpoint{11.995598in}{2.077312in}}{\pgfqpoint{11.991371in}{2.091805in}}%
\pgfpathquadraticcurveto{\pgfqpoint{11.987145in}{2.106319in}}{\pgfqpoint{11.978589in}{2.121122in}}%
\pgfpathlineto{\pgfqpoint{11.978589in}{2.121122in}}%
\pgfpathclose%
\pgfusepath{fill}%
\end{pgfscope}%
\begin{pgfscope}%
\pgfsetbuttcap%
\pgfsetroundjoin%
\definecolor{currentfill}{rgb}{0.090196,0.168627,0.227451}%
\pgfsetfillcolor{currentfill}%
\pgfsetlinewidth{0.000000pt}%
\definecolor{currentstroke}{rgb}{0.090196,0.168627,0.227451}%
\pgfsetstrokecolor{currentstroke}%
\pgfsetdash{}{0pt}%
\pgfpathmoveto{\pgfqpoint{11.018611in}{1.937410in}}%
\pgfpathlineto{\pgfqpoint{11.018611in}{1.924711in}}%
\pgfpathquadraticcurveto{\pgfqpoint{11.011210in}{1.928257in}}{\pgfqpoint{11.004633in}{1.929988in}}%
\pgfpathquadraticcurveto{\pgfqpoint{10.998056in}{1.931741in}}{\pgfqpoint{10.991933in}{1.931741in}}%
\pgfpathquadraticcurveto{\pgfqpoint{10.981316in}{1.931741in}}{\pgfqpoint{10.975543in}{1.927617in}}%
\pgfpathquadraticcurveto{\pgfqpoint{10.969771in}{1.923494in}}{\pgfqpoint{10.969771in}{1.915887in}}%
\pgfpathquadraticcurveto{\pgfqpoint{10.969771in}{1.909496in}}{\pgfqpoint{10.973605in}{1.906238in}}%
\pgfpathquadraticcurveto{\pgfqpoint{10.977440in}{1.903002in}}{\pgfqpoint{10.988140in}{1.901002in}}%
\pgfpathlineto{\pgfqpoint{10.995995in}{1.899394in}}%
\pgfpathquadraticcurveto{\pgfqpoint{11.010550in}{1.896611in}}{\pgfqpoint{11.017477in}{1.889622in}}%
\pgfpathquadraticcurveto{\pgfqpoint{11.024404in}{1.882633in}}{\pgfqpoint{11.024404in}{1.870923in}}%
\pgfpathquadraticcurveto{\pgfqpoint{11.024404in}{1.856924in}}{\pgfqpoint{11.015024in}{1.849708in}}%
\pgfpathquadraticcurveto{\pgfqpoint{11.005664in}{1.842493in}}{\pgfqpoint{10.987563in}{1.842493in}}%
\pgfpathquadraticcurveto{\pgfqpoint{10.980739in}{1.842493in}}{\pgfqpoint{10.973028in}{1.844039in}}%
\pgfpathquadraticcurveto{\pgfqpoint{10.965338in}{1.845585in}}{\pgfqpoint{10.957092in}{1.848616in}}%
\pgfpathlineto{\pgfqpoint{10.957092in}{1.862016in}}%
\pgfpathquadraticcurveto{\pgfqpoint{10.965008in}{1.857584in}}{\pgfqpoint{10.972616in}{1.855316in}}%
\pgfpathquadraticcurveto{\pgfqpoint{10.980223in}{1.853069in}}{\pgfqpoint{10.987563in}{1.853069in}}%
\pgfpathquadraticcurveto{\pgfqpoint{10.998696in}{1.853069in}}{\pgfqpoint{11.004757in}{1.857440in}}%
\pgfpathquadraticcurveto{\pgfqpoint{11.010818in}{1.861831in}}{\pgfqpoint{11.010818in}{1.869954in}}%
\pgfpathquadraticcurveto{\pgfqpoint{11.010818in}{1.877025in}}{\pgfqpoint{11.006468in}{1.881025in}}%
\pgfpathquadraticcurveto{\pgfqpoint{11.002118in}{1.885024in}}{\pgfqpoint{10.992201in}{1.887024in}}%
\pgfpathlineto{\pgfqpoint{10.984264in}{1.888570in}}%
\pgfpathquadraticcurveto{\pgfqpoint{10.969709in}{1.891456in}}{\pgfqpoint{10.963194in}{1.897641in}}%
\pgfpathquadraticcurveto{\pgfqpoint{10.956700in}{1.903826in}}{\pgfqpoint{10.956700in}{1.914856in}}%
\pgfpathquadraticcurveto{\pgfqpoint{10.956700in}{1.927617in}}{\pgfqpoint{10.965689in}{1.934957in}}%
\pgfpathquadraticcurveto{\pgfqpoint{10.974678in}{1.942296in}}{\pgfqpoint{10.990449in}{1.942296in}}%
\pgfpathquadraticcurveto{\pgfqpoint{10.997232in}{1.942296in}}{\pgfqpoint{11.004241in}{1.941059in}}%
\pgfpathquadraticcurveto{\pgfqpoint{11.011271in}{1.939843in}}{\pgfqpoint{11.018611in}{1.937410in}}%
\pgfpathlineto{\pgfqpoint{11.018611in}{1.937410in}}%
\pgfpathclose%
\pgfpathmoveto{\pgfqpoint{11.055926in}{1.937019in}}%
\pgfpathlineto{\pgfqpoint{11.055926in}{1.916526in}}%
\pgfpathlineto{\pgfqpoint{11.080336in}{1.916526in}}%
\pgfpathlineto{\pgfqpoint{11.080336in}{1.907310in}}%
\pgfpathlineto{\pgfqpoint{11.055926in}{1.907310in}}%
\pgfpathlineto{\pgfqpoint{11.055926in}{1.868139in}}%
\pgfpathquadraticcurveto{\pgfqpoint{11.055926in}{1.859316in}}{\pgfqpoint{11.058339in}{1.856800in}}%
\pgfpathquadraticcurveto{\pgfqpoint{11.060751in}{1.854285in}}{\pgfqpoint{11.068173in}{1.854285in}}%
\pgfpathlineto{\pgfqpoint{11.080336in}{1.854285in}}%
\pgfpathlineto{\pgfqpoint{11.080336in}{1.844369in}}%
\pgfpathlineto{\pgfqpoint{11.068173in}{1.844369in}}%
\pgfpathquadraticcurveto{\pgfqpoint{11.054442in}{1.844369in}}{\pgfqpoint{11.049226in}{1.849482in}}%
\pgfpathquadraticcurveto{\pgfqpoint{11.044010in}{1.854615in}}{\pgfqpoint{11.044010in}{1.868139in}}%
\pgfpathlineto{\pgfqpoint{11.044010in}{1.907310in}}%
\pgfpathlineto{\pgfqpoint{11.035310in}{1.907310in}}%
\pgfpathlineto{\pgfqpoint{11.035310in}{1.916526in}}%
\pgfpathlineto{\pgfqpoint{11.044010in}{1.916526in}}%
\pgfpathlineto{\pgfqpoint{11.044010in}{1.937019in}}%
\pgfpathlineto{\pgfqpoint{11.055926in}{1.937019in}}%
\pgfpathlineto{\pgfqpoint{11.055926in}{1.937019in}}%
\pgfpathclose%
\pgfpathmoveto{\pgfqpoint{11.094706in}{1.872840in}}%
\pgfpathlineto{\pgfqpoint{11.094706in}{1.916526in}}%
\pgfpathlineto{\pgfqpoint{11.106560in}{1.916526in}}%
\pgfpathlineto{\pgfqpoint{11.106560in}{1.873293in}}%
\pgfpathquadraticcurveto{\pgfqpoint{11.106560in}{1.863047in}}{\pgfqpoint{11.110539in}{1.857914in}}%
\pgfpathquadraticcurveto{\pgfqpoint{11.114539in}{1.852801in}}{\pgfqpoint{11.122538in}{1.852801in}}%
\pgfpathquadraticcurveto{\pgfqpoint{11.132124in}{1.852801in}}{\pgfqpoint{11.137691in}{1.858924in}}%
\pgfpathquadraticcurveto{\pgfqpoint{11.143278in}{1.865047in}}{\pgfqpoint{11.143278in}{1.875623in}}%
\pgfpathlineto{\pgfqpoint{11.143278in}{1.916526in}}%
\pgfpathlineto{\pgfqpoint{11.155132in}{1.916526in}}%
\pgfpathlineto{\pgfqpoint{11.155132in}{1.844369in}}%
\pgfpathlineto{\pgfqpoint{11.143278in}{1.844369in}}%
\pgfpathlineto{\pgfqpoint{11.143278in}{1.855460in}}%
\pgfpathquadraticcurveto{\pgfqpoint{11.138969in}{1.848884in}}{\pgfqpoint{11.133258in}{1.845688in}}%
\pgfpathquadraticcurveto{\pgfqpoint{11.127568in}{1.842493in}}{\pgfqpoint{11.120023in}{1.842493in}}%
\pgfpathquadraticcurveto{\pgfqpoint{11.107591in}{1.842493in}}{\pgfqpoint{11.101138in}{1.850224in}}%
\pgfpathquadraticcurveto{\pgfqpoint{11.094706in}{1.857955in}}{\pgfqpoint{11.094706in}{1.872840in}}%
\pgfpathlineto{\pgfqpoint{11.094706in}{1.872840in}}%
\pgfpathclose%
\pgfpathmoveto{\pgfqpoint{11.124538in}{1.918258in}}%
\pgfpathlineto{\pgfqpoint{11.124538in}{1.918258in}}%
\pgfpathlineto{\pgfqpoint{11.124538in}{1.918258in}}%
\pgfpathclose%
\pgfpathmoveto{\pgfqpoint{11.227021in}{1.905579in}}%
\pgfpathlineto{\pgfqpoint{11.227021in}{1.944626in}}%
\pgfpathlineto{\pgfqpoint{11.238876in}{1.944626in}}%
\pgfpathlineto{\pgfqpoint{11.238876in}{1.844369in}}%
\pgfpathlineto{\pgfqpoint{11.227021in}{1.844369in}}%
\pgfpathlineto{\pgfqpoint{11.227021in}{1.855192in}}%
\pgfpathquadraticcurveto{\pgfqpoint{11.223290in}{1.848760in}}{\pgfqpoint{11.217579in}{1.845626in}}%
\pgfpathquadraticcurveto{\pgfqpoint{11.211889in}{1.842493in}}{\pgfqpoint{11.203890in}{1.842493in}}%
\pgfpathquadraticcurveto{\pgfqpoint{11.190819in}{1.842493in}}{\pgfqpoint{11.182593in}{1.852925in}}%
\pgfpathquadraticcurveto{\pgfqpoint{11.174388in}{1.863377in}}{\pgfqpoint{11.174388in}{1.880386in}}%
\pgfpathquadraticcurveto{\pgfqpoint{11.174388in}{1.897394in}}{\pgfqpoint{11.182593in}{1.907826in}}%
\pgfpathquadraticcurveto{\pgfqpoint{11.190819in}{1.918258in}}{\pgfqpoint{11.203890in}{1.918258in}}%
\pgfpathquadraticcurveto{\pgfqpoint{11.211889in}{1.918258in}}{\pgfqpoint{11.217579in}{1.915124in}}%
\pgfpathquadraticcurveto{\pgfqpoint{11.223290in}{1.912011in}}{\pgfqpoint{11.227021in}{1.905579in}}%
\pgfpathlineto{\pgfqpoint{11.227021in}{1.905579in}}%
\pgfpathclose%
\pgfpathmoveto{\pgfqpoint{11.186634in}{1.880386in}}%
\pgfpathquadraticcurveto{\pgfqpoint{11.186634in}{1.867315in}}{\pgfqpoint{11.192015in}{1.859872in}}%
\pgfpathquadraticcurveto{\pgfqpoint{11.197396in}{1.852430in}}{\pgfqpoint{11.206797in}{1.852430in}}%
\pgfpathquadraticcurveto{\pgfqpoint{11.216198in}{1.852430in}}{\pgfqpoint{11.221599in}{1.859872in}}%
\pgfpathquadraticcurveto{\pgfqpoint{11.227021in}{1.867315in}}{\pgfqpoint{11.227021in}{1.880386in}}%
\pgfpathquadraticcurveto{\pgfqpoint{11.227021in}{1.893456in}}{\pgfqpoint{11.221599in}{1.900899in}}%
\pgfpathquadraticcurveto{\pgfqpoint{11.216198in}{1.908341in}}{\pgfqpoint{11.206797in}{1.908341in}}%
\pgfpathquadraticcurveto{\pgfqpoint{11.197396in}{1.908341in}}{\pgfqpoint{11.192015in}{1.900899in}}%
\pgfpathquadraticcurveto{\pgfqpoint{11.186634in}{1.893456in}}{\pgfqpoint{11.186634in}{1.880386in}}%
\pgfpathlineto{\pgfqpoint{11.186634in}{1.880386in}}%
\pgfpathclose%
\pgfpathmoveto{\pgfqpoint{11.293323in}{1.837668in}}%
\pgfpathquadraticcurveto{\pgfqpoint{11.288314in}{1.824783in}}{\pgfqpoint{11.283531in}{1.820866in}}%
\pgfpathquadraticcurveto{\pgfqpoint{11.278768in}{1.816928in}}{\pgfqpoint{11.270790in}{1.816928in}}%
\pgfpathlineto{\pgfqpoint{11.261306in}{1.816928in}}%
\pgfpathlineto{\pgfqpoint{11.261306in}{1.826845in}}%
\pgfpathlineto{\pgfqpoint{11.268275in}{1.826845in}}%
\pgfpathquadraticcurveto{\pgfqpoint{11.273161in}{1.826845in}}{\pgfqpoint{11.275861in}{1.829175in}}%
\pgfpathquadraticcurveto{\pgfqpoint{11.278583in}{1.831484in}}{\pgfqpoint{11.281861in}{1.840122in}}%
\pgfpathlineto{\pgfqpoint{11.283984in}{1.845523in}}%
\pgfpathlineto{\pgfqpoint{11.254812in}{1.916526in}}%
\pgfpathlineto{\pgfqpoint{11.267367in}{1.916526in}}%
\pgfpathlineto{\pgfqpoint{11.289922in}{1.860099in}}%
\pgfpathlineto{\pgfqpoint{11.312476in}{1.916526in}}%
\pgfpathlineto{\pgfqpoint{11.325031in}{1.916526in}}%
\pgfpathlineto{\pgfqpoint{11.293323in}{1.837668in}}%
\pgfpathlineto{\pgfqpoint{11.293323in}{1.837668in}}%
\pgfpathclose%
\pgfpathmoveto{\pgfqpoint{11.396219in}{1.855316in}}%
\pgfpathlineto{\pgfqpoint{11.441637in}{1.855316in}}%
\pgfpathlineto{\pgfqpoint{11.441637in}{1.844369in}}%
\pgfpathlineto{\pgfqpoint{11.380572in}{1.844369in}}%
\pgfpathlineto{\pgfqpoint{11.380572in}{1.855316in}}%
\pgfpathquadraticcurveto{\pgfqpoint{11.387973in}{1.862985in}}{\pgfqpoint{11.400755in}{1.875891in}}%
\pgfpathquadraticcurveto{\pgfqpoint{11.413558in}{1.888818in}}{\pgfqpoint{11.416836in}{1.892570in}}%
\pgfpathquadraticcurveto{\pgfqpoint{11.423082in}{1.899579in}}{\pgfqpoint{11.425556in}{1.904445in}}%
\pgfpathquadraticcurveto{\pgfqpoint{11.428051in}{1.909310in}}{\pgfqpoint{11.428051in}{1.914011in}}%
\pgfpathquadraticcurveto{\pgfqpoint{11.428051in}{1.921680in}}{\pgfqpoint{11.422670in}{1.926504in}}%
\pgfpathquadraticcurveto{\pgfqpoint{11.417289in}{1.931349in}}{\pgfqpoint{11.408651in}{1.931349in}}%
\pgfpathquadraticcurveto{\pgfqpoint{11.402528in}{1.931349in}}{\pgfqpoint{11.395725in}{1.929226in}}%
\pgfpathquadraticcurveto{\pgfqpoint{11.388942in}{1.927102in}}{\pgfqpoint{11.381211in}{1.922773in}}%
\pgfpathlineto{\pgfqpoint{11.381211in}{1.935926in}}%
\pgfpathquadraticcurveto{\pgfqpoint{11.389066in}{1.939080in}}{\pgfqpoint{11.395890in}{1.940688in}}%
\pgfpathquadraticcurveto{\pgfqpoint{11.402734in}{1.942296in}}{\pgfqpoint{11.408404in}{1.942296in}}%
\pgfpathquadraticcurveto{\pgfqpoint{11.423350in}{1.942296in}}{\pgfqpoint{11.432236in}{1.934813in}}%
\pgfpathquadraticcurveto{\pgfqpoint{11.441122in}{1.927349in}}{\pgfqpoint{11.441122in}{1.914856in}}%
\pgfpathquadraticcurveto{\pgfqpoint{11.441122in}{1.908918in}}{\pgfqpoint{11.438895in}{1.903599in}}%
\pgfpathquadraticcurveto{\pgfqpoint{11.436689in}{1.898301in}}{\pgfqpoint{11.430814in}{1.891085in}}%
\pgfpathquadraticcurveto{\pgfqpoint{11.429206in}{1.889209in}}{\pgfqpoint{11.420567in}{1.880282in}}%
\pgfpathquadraticcurveto{\pgfqpoint{11.411950in}{1.871356in}}{\pgfqpoint{11.396219in}{1.855316in}}%
\pgfpathlineto{\pgfqpoint{11.396219in}{1.855316in}}%
\pgfpathclose%
\pgfpathmoveto{\pgfqpoint{11.480169in}{1.855316in}}%
\pgfpathlineto{\pgfqpoint{11.525587in}{1.855316in}}%
\pgfpathlineto{\pgfqpoint{11.525587in}{1.844369in}}%
\pgfpathlineto{\pgfqpoint{11.464521in}{1.844369in}}%
\pgfpathlineto{\pgfqpoint{11.464521in}{1.855316in}}%
\pgfpathquadraticcurveto{\pgfqpoint{11.471923in}{1.862985in}}{\pgfqpoint{11.484705in}{1.875891in}}%
\pgfpathquadraticcurveto{\pgfqpoint{11.497507in}{1.888818in}}{\pgfqpoint{11.500785in}{1.892570in}}%
\pgfpathquadraticcurveto{\pgfqpoint{11.507032in}{1.899579in}}{\pgfqpoint{11.509506in}{1.904445in}}%
\pgfpathquadraticcurveto{\pgfqpoint{11.512001in}{1.909310in}}{\pgfqpoint{11.512001in}{1.914011in}}%
\pgfpathquadraticcurveto{\pgfqpoint{11.512001in}{1.921680in}}{\pgfqpoint{11.506620in}{1.926504in}}%
\pgfpathquadraticcurveto{\pgfqpoint{11.501239in}{1.931349in}}{\pgfqpoint{11.492601in}{1.931349in}}%
\pgfpathquadraticcurveto{\pgfqpoint{11.486478in}{1.931349in}}{\pgfqpoint{11.479674in}{1.929226in}}%
\pgfpathquadraticcurveto{\pgfqpoint{11.472891in}{1.927102in}}{\pgfqpoint{11.465160in}{1.922773in}}%
\pgfpathlineto{\pgfqpoint{11.465160in}{1.935926in}}%
\pgfpathquadraticcurveto{\pgfqpoint{11.473015in}{1.939080in}}{\pgfqpoint{11.479839in}{1.940688in}}%
\pgfpathquadraticcurveto{\pgfqpoint{11.486684in}{1.942296in}}{\pgfqpoint{11.492353in}{1.942296in}}%
\pgfpathquadraticcurveto{\pgfqpoint{11.507300in}{1.942296in}}{\pgfqpoint{11.516186in}{1.934813in}}%
\pgfpathquadraticcurveto{\pgfqpoint{11.525071in}{1.927349in}}{\pgfqpoint{11.525071in}{1.914856in}}%
\pgfpathquadraticcurveto{\pgfqpoint{11.525071in}{1.908918in}}{\pgfqpoint{11.522845in}{1.903599in}}%
\pgfpathquadraticcurveto{\pgfqpoint{11.520639in}{1.898301in}}{\pgfqpoint{11.514763in}{1.891085in}}%
\pgfpathquadraticcurveto{\pgfqpoint{11.513155in}{1.889209in}}{\pgfqpoint{11.504517in}{1.880282in}}%
\pgfpathquadraticcurveto{\pgfqpoint{11.495899in}{1.871356in}}{\pgfqpoint{11.480169in}{1.855316in}}%
\pgfpathlineto{\pgfqpoint{11.480169in}{1.855316in}}%
\pgfpathclose%
\pgfpathmoveto{\pgfqpoint{11.621638in}{1.944482in}}%
\pgfpathquadraticcurveto{\pgfqpoint{11.613021in}{1.929679in}}{\pgfqpoint{11.608815in}{1.915165in}}%
\pgfpathquadraticcurveto{\pgfqpoint{11.604630in}{1.900672in}}{\pgfqpoint{11.604630in}{1.885787in}}%
\pgfpathquadraticcurveto{\pgfqpoint{11.604630in}{1.870923in}}{\pgfqpoint{11.608856in}{1.856326in}}%
\pgfpathquadraticcurveto{\pgfqpoint{11.613082in}{1.841730in}}{\pgfqpoint{11.621638in}{1.826969in}}%
\pgfpathlineto{\pgfqpoint{11.611330in}{1.826969in}}%
\pgfpathquadraticcurveto{\pgfqpoint{11.601682in}{1.842122in}}{\pgfqpoint{11.596878in}{1.856739in}}%
\pgfpathquadraticcurveto{\pgfqpoint{11.592074in}{1.871356in}}{\pgfqpoint{11.592074in}{1.885787in}}%
\pgfpathquadraticcurveto{\pgfqpoint{11.592074in}{1.900157in}}{\pgfqpoint{11.596837in}{1.914712in}}%
\pgfpathquadraticcurveto{\pgfqpoint{11.601620in}{1.929287in}}{\pgfqpoint{11.611330in}{1.944482in}}%
\pgfpathlineto{\pgfqpoint{11.621638in}{1.944482in}}%
\pgfpathlineto{\pgfqpoint{11.621638in}{1.944482in}}%
\pgfpathclose%
\pgfpathmoveto{\pgfqpoint{11.657531in}{1.855316in}}%
\pgfpathlineto{\pgfqpoint{11.702949in}{1.855316in}}%
\pgfpathlineto{\pgfqpoint{11.702949in}{1.844369in}}%
\pgfpathlineto{\pgfqpoint{11.641883in}{1.844369in}}%
\pgfpathlineto{\pgfqpoint{11.641883in}{1.855316in}}%
\pgfpathquadraticcurveto{\pgfqpoint{11.649285in}{1.862985in}}{\pgfqpoint{11.662067in}{1.875891in}}%
\pgfpathquadraticcurveto{\pgfqpoint{11.674870in}{1.888818in}}{\pgfqpoint{11.678148in}{1.892570in}}%
\pgfpathquadraticcurveto{\pgfqpoint{11.684394in}{1.899579in}}{\pgfqpoint{11.686868in}{1.904445in}}%
\pgfpathquadraticcurveto{\pgfqpoint{11.689363in}{1.909310in}}{\pgfqpoint{11.689363in}{1.914011in}}%
\pgfpathquadraticcurveto{\pgfqpoint{11.689363in}{1.921680in}}{\pgfqpoint{11.683982in}{1.926504in}}%
\pgfpathquadraticcurveto{\pgfqpoint{11.678601in}{1.931349in}}{\pgfqpoint{11.669963in}{1.931349in}}%
\pgfpathquadraticcurveto{\pgfqpoint{11.663840in}{1.931349in}}{\pgfqpoint{11.657036in}{1.929226in}}%
\pgfpathquadraticcurveto{\pgfqpoint{11.650254in}{1.927102in}}{\pgfqpoint{11.642523in}{1.922773in}}%
\pgfpathlineto{\pgfqpoint{11.642523in}{1.935926in}}%
\pgfpathquadraticcurveto{\pgfqpoint{11.650377in}{1.939080in}}{\pgfqpoint{11.657201in}{1.940688in}}%
\pgfpathquadraticcurveto{\pgfqpoint{11.664046in}{1.942296in}}{\pgfqpoint{11.669715in}{1.942296in}}%
\pgfpathquadraticcurveto{\pgfqpoint{11.684662in}{1.942296in}}{\pgfqpoint{11.693548in}{1.934813in}}%
\pgfpathquadraticcurveto{\pgfqpoint{11.702434in}{1.927349in}}{\pgfqpoint{11.702434in}{1.914856in}}%
\pgfpathquadraticcurveto{\pgfqpoint{11.702434in}{1.908918in}}{\pgfqpoint{11.700207in}{1.903599in}}%
\pgfpathquadraticcurveto{\pgfqpoint{11.698001in}{1.898301in}}{\pgfqpoint{11.692125in}{1.891085in}}%
\pgfpathquadraticcurveto{\pgfqpoint{11.690517in}{1.889209in}}{\pgfqpoint{11.681879in}{1.880282in}}%
\pgfpathquadraticcurveto{\pgfqpoint{11.673262in}{1.871356in}}{\pgfqpoint{11.657531in}{1.855316in}}%
\pgfpathlineto{\pgfqpoint{11.657531in}{1.855316in}}%
\pgfpathclose%
\pgfpathmoveto{\pgfqpoint{11.758098in}{1.931988in}}%
\pgfpathquadraticcurveto{\pgfqpoint{11.748058in}{1.931988in}}{\pgfqpoint{11.742986in}{1.922092in}}%
\pgfpathquadraticcurveto{\pgfqpoint{11.737935in}{1.912217in}}{\pgfqpoint{11.737935in}{1.892364in}}%
\pgfpathquadraticcurveto{\pgfqpoint{11.737935in}{1.872593in}}{\pgfqpoint{11.742986in}{1.862697in}}%
\pgfpathquadraticcurveto{\pgfqpoint{11.748058in}{1.852801in}}{\pgfqpoint{11.758098in}{1.852801in}}%
\pgfpathquadraticcurveto{\pgfqpoint{11.768220in}{1.852801in}}{\pgfqpoint{11.773271in}{1.862697in}}%
\pgfpathquadraticcurveto{\pgfqpoint{11.778343in}{1.872593in}}{\pgfqpoint{11.778343in}{1.892364in}}%
\pgfpathquadraticcurveto{\pgfqpoint{11.778343in}{1.912217in}}{\pgfqpoint{11.773271in}{1.922092in}}%
\pgfpathquadraticcurveto{\pgfqpoint{11.768220in}{1.931988in}}{\pgfqpoint{11.758098in}{1.931988in}}%
\pgfpathlineto{\pgfqpoint{11.758098in}{1.931988in}}%
\pgfpathclose%
\pgfpathmoveto{\pgfqpoint{11.758098in}{1.942296in}}%
\pgfpathquadraticcurveto{\pgfqpoint{11.774281in}{1.942296in}}{\pgfqpoint{11.782817in}{1.929494in}}%
\pgfpathquadraticcurveto{\pgfqpoint{11.791352in}{1.916711in}}{\pgfqpoint{11.791352in}{1.892364in}}%
\pgfpathquadraticcurveto{\pgfqpoint{11.791352in}{1.868078in}}{\pgfqpoint{11.782817in}{1.855275in}}%
\pgfpathquadraticcurveto{\pgfqpoint{11.774281in}{1.842493in}}{\pgfqpoint{11.758098in}{1.842493in}}%
\pgfpathquadraticcurveto{\pgfqpoint{11.741934in}{1.842493in}}{\pgfqpoint{11.733399in}{1.855275in}}%
\pgfpathquadraticcurveto{\pgfqpoint{11.724864in}{1.868078in}}{\pgfqpoint{11.724864in}{1.892364in}}%
\pgfpathquadraticcurveto{\pgfqpoint{11.724864in}{1.916711in}}{\pgfqpoint{11.733399in}{1.929494in}}%
\pgfpathquadraticcurveto{\pgfqpoint{11.741934in}{1.942296in}}{\pgfqpoint{11.758098in}{1.942296in}}%
\pgfpathlineto{\pgfqpoint{11.758098in}{1.942296in}}%
\pgfpathclose%
\pgfpathmoveto{\pgfqpoint{11.825431in}{1.855316in}}%
\pgfpathlineto{\pgfqpoint{11.870848in}{1.855316in}}%
\pgfpathlineto{\pgfqpoint{11.870848in}{1.844369in}}%
\pgfpathlineto{\pgfqpoint{11.809783in}{1.844369in}}%
\pgfpathlineto{\pgfqpoint{11.809783in}{1.855316in}}%
\pgfpathquadraticcurveto{\pgfqpoint{11.817184in}{1.862985in}}{\pgfqpoint{11.829966in}{1.875891in}}%
\pgfpathquadraticcurveto{\pgfqpoint{11.842769in}{1.888818in}}{\pgfqpoint{11.846047in}{1.892570in}}%
\pgfpathquadraticcurveto{\pgfqpoint{11.852294in}{1.899579in}}{\pgfqpoint{11.854768in}{1.904445in}}%
\pgfpathquadraticcurveto{\pgfqpoint{11.857262in}{1.909310in}}{\pgfqpoint{11.857262in}{1.914011in}}%
\pgfpathquadraticcurveto{\pgfqpoint{11.857262in}{1.921680in}}{\pgfqpoint{11.851881in}{1.926504in}}%
\pgfpathquadraticcurveto{\pgfqpoint{11.846500in}{1.931349in}}{\pgfqpoint{11.837862in}{1.931349in}}%
\pgfpathquadraticcurveto{\pgfqpoint{11.831739in}{1.931349in}}{\pgfqpoint{11.824936in}{1.929226in}}%
\pgfpathquadraticcurveto{\pgfqpoint{11.818153in}{1.927102in}}{\pgfqpoint{11.810422in}{1.922773in}}%
\pgfpathlineto{\pgfqpoint{11.810422in}{1.935926in}}%
\pgfpathquadraticcurveto{\pgfqpoint{11.818277in}{1.939080in}}{\pgfqpoint{11.825101in}{1.940688in}}%
\pgfpathquadraticcurveto{\pgfqpoint{11.831945in}{1.942296in}}{\pgfqpoint{11.837615in}{1.942296in}}%
\pgfpathquadraticcurveto{\pgfqpoint{11.852562in}{1.942296in}}{\pgfqpoint{11.861447in}{1.934813in}}%
\pgfpathquadraticcurveto{\pgfqpoint{11.870333in}{1.927349in}}{\pgfqpoint{11.870333in}{1.914856in}}%
\pgfpathquadraticcurveto{\pgfqpoint{11.870333in}{1.908918in}}{\pgfqpoint{11.868106in}{1.903599in}}%
\pgfpathquadraticcurveto{\pgfqpoint{11.865900in}{1.898301in}}{\pgfqpoint{11.860025in}{1.891085in}}%
\pgfpathquadraticcurveto{\pgfqpoint{11.858417in}{1.889209in}}{\pgfqpoint{11.849778in}{1.880282in}}%
\pgfpathquadraticcurveto{\pgfqpoint{11.841161in}{1.871356in}}{\pgfqpoint{11.825431in}{1.855316in}}%
\pgfpathlineto{\pgfqpoint{11.825431in}{1.855316in}}%
\pgfpathclose%
\pgfpathmoveto{\pgfqpoint{11.933934in}{1.929226in}}%
\pgfpathlineto{\pgfqpoint{11.901072in}{1.877870in}}%
\pgfpathlineto{\pgfqpoint{11.933934in}{1.877870in}}%
\pgfpathlineto{\pgfqpoint{11.933934in}{1.929226in}}%
\pgfpathlineto{\pgfqpoint{11.933934in}{1.929226in}}%
\pgfpathclose%
\pgfpathmoveto{\pgfqpoint{11.930512in}{1.940565in}}%
\pgfpathlineto{\pgfqpoint{11.946881in}{1.940565in}}%
\pgfpathlineto{\pgfqpoint{11.946881in}{1.877870in}}%
\pgfpathlineto{\pgfqpoint{11.960612in}{1.877870in}}%
\pgfpathlineto{\pgfqpoint{11.960612in}{1.867047in}}%
\pgfpathlineto{\pgfqpoint{11.946881in}{1.867047in}}%
\pgfpathlineto{\pgfqpoint{11.946881in}{1.844369in}}%
\pgfpathlineto{\pgfqpoint{11.933934in}{1.844369in}}%
\pgfpathlineto{\pgfqpoint{11.933934in}{1.867047in}}%
\pgfpathlineto{\pgfqpoint{11.890516in}{1.867047in}}%
\pgfpathlineto{\pgfqpoint{11.890516in}{1.879602in}}%
\pgfpathlineto{\pgfqpoint{11.930512in}{1.940565in}}%
\pgfpathlineto{\pgfqpoint{11.930512in}{1.940565in}}%
\pgfpathclose%
\pgfpathmoveto{\pgfqpoint{11.978589in}{1.944482in}}%
\pgfpathlineto{\pgfqpoint{11.988897in}{1.944482in}}%
\pgfpathquadraticcurveto{\pgfqpoint{11.998546in}{1.929287in}}{\pgfqpoint{12.003349in}{1.914712in}}%
\pgfpathquadraticcurveto{\pgfqpoint{12.008153in}{1.900157in}}{\pgfqpoint{12.008153in}{1.885787in}}%
\pgfpathquadraticcurveto{\pgfqpoint{12.008153in}{1.871356in}}{\pgfqpoint{12.003349in}{1.856739in}}%
\pgfpathquadraticcurveto{\pgfqpoint{11.998546in}{1.842122in}}{\pgfqpoint{11.988897in}{1.826969in}}%
\pgfpathlineto{\pgfqpoint{11.978589in}{1.826969in}}%
\pgfpathquadraticcurveto{\pgfqpoint{11.987145in}{1.841730in}}{\pgfqpoint{11.991371in}{1.856326in}}%
\pgfpathquadraticcurveto{\pgfqpoint{11.995598in}{1.870923in}}{\pgfqpoint{11.995598in}{1.885787in}}%
\pgfpathquadraticcurveto{\pgfqpoint{11.995598in}{1.900672in}}{\pgfqpoint{11.991371in}{1.915165in}}%
\pgfpathquadraticcurveto{\pgfqpoint{11.987145in}{1.929679in}}{\pgfqpoint{11.978589in}{1.944482in}}%
\pgfpathlineto{\pgfqpoint{11.978589in}{1.944482in}}%
\pgfpathclose%
\pgfusepath{fill}%
\end{pgfscope}%
\begin{pgfscope}%
\pgfsetbuttcap%
\pgfsetroundjoin%
\definecolor{currentfill}{rgb}{0.090196,0.168627,0.227451}%
\pgfsetfillcolor{currentfill}%
\pgfsetlinewidth{0.000000pt}%
\definecolor{currentstroke}{rgb}{0.090196,0.168627,0.227451}%
\pgfsetstrokecolor{currentstroke}%
\pgfsetdash{}{0pt}%
\pgfpathmoveto{\pgfqpoint{11.018611in}{1.760770in}}%
\pgfpathlineto{\pgfqpoint{11.018611in}{1.748071in}}%
\pgfpathquadraticcurveto{\pgfqpoint{11.011210in}{1.751617in}}{\pgfqpoint{11.004633in}{1.753348in}}%
\pgfpathquadraticcurveto{\pgfqpoint{10.998056in}{1.755101in}}{\pgfqpoint{10.991933in}{1.755101in}}%
\pgfpathquadraticcurveto{\pgfqpoint{10.981316in}{1.755101in}}{\pgfqpoint{10.975543in}{1.750977in}}%
\pgfpathquadraticcurveto{\pgfqpoint{10.969771in}{1.746854in}}{\pgfqpoint{10.969771in}{1.739247in}}%
\pgfpathquadraticcurveto{\pgfqpoint{10.969771in}{1.732856in}}{\pgfqpoint{10.973605in}{1.729598in}}%
\pgfpathquadraticcurveto{\pgfqpoint{10.977440in}{1.726362in}}{\pgfqpoint{10.988140in}{1.724362in}}%
\pgfpathlineto{\pgfqpoint{10.995995in}{1.722754in}}%
\pgfpathquadraticcurveto{\pgfqpoint{11.010550in}{1.719971in}}{\pgfqpoint{11.017477in}{1.712982in}}%
\pgfpathquadraticcurveto{\pgfqpoint{11.024404in}{1.705993in}}{\pgfqpoint{11.024404in}{1.694283in}}%
\pgfpathquadraticcurveto{\pgfqpoint{11.024404in}{1.680284in}}{\pgfqpoint{11.015024in}{1.673068in}}%
\pgfpathquadraticcurveto{\pgfqpoint{11.005664in}{1.665853in}}{\pgfqpoint{10.987563in}{1.665853in}}%
\pgfpathquadraticcurveto{\pgfqpoint{10.980739in}{1.665853in}}{\pgfqpoint{10.973028in}{1.667399in}}%
\pgfpathquadraticcurveto{\pgfqpoint{10.965338in}{1.668945in}}{\pgfqpoint{10.957092in}{1.671976in}}%
\pgfpathlineto{\pgfqpoint{10.957092in}{1.685376in}}%
\pgfpathquadraticcurveto{\pgfqpoint{10.965008in}{1.680944in}}{\pgfqpoint{10.972616in}{1.678676in}}%
\pgfpathquadraticcurveto{\pgfqpoint{10.980223in}{1.676429in}}{\pgfqpoint{10.987563in}{1.676429in}}%
\pgfpathquadraticcurveto{\pgfqpoint{10.998696in}{1.676429in}}{\pgfqpoint{11.004757in}{1.680800in}}%
\pgfpathquadraticcurveto{\pgfqpoint{11.010818in}{1.685191in}}{\pgfqpoint{11.010818in}{1.693314in}}%
\pgfpathquadraticcurveto{\pgfqpoint{11.010818in}{1.700385in}}{\pgfqpoint{11.006468in}{1.704385in}}%
\pgfpathquadraticcurveto{\pgfqpoint{11.002118in}{1.708384in}}{\pgfqpoint{10.992201in}{1.710384in}}%
\pgfpathlineto{\pgfqpoint{10.984264in}{1.711930in}}%
\pgfpathquadraticcurveto{\pgfqpoint{10.969709in}{1.714816in}}{\pgfqpoint{10.963194in}{1.721001in}}%
\pgfpathquadraticcurveto{\pgfqpoint{10.956700in}{1.727186in}}{\pgfqpoint{10.956700in}{1.738216in}}%
\pgfpathquadraticcurveto{\pgfqpoint{10.956700in}{1.750977in}}{\pgfqpoint{10.965689in}{1.758317in}}%
\pgfpathquadraticcurveto{\pgfqpoint{10.974678in}{1.765656in}}{\pgfqpoint{10.990449in}{1.765656in}}%
\pgfpathquadraticcurveto{\pgfqpoint{10.997232in}{1.765656in}}{\pgfqpoint{11.004241in}{1.764419in}}%
\pgfpathquadraticcurveto{\pgfqpoint{11.011271in}{1.763203in}}{\pgfqpoint{11.018611in}{1.760770in}}%
\pgfpathlineto{\pgfqpoint{11.018611in}{1.760770in}}%
\pgfpathclose%
\pgfpathmoveto{\pgfqpoint{11.055926in}{1.760379in}}%
\pgfpathlineto{\pgfqpoint{11.055926in}{1.739886in}}%
\pgfpathlineto{\pgfqpoint{11.080336in}{1.739886in}}%
\pgfpathlineto{\pgfqpoint{11.080336in}{1.730670in}}%
\pgfpathlineto{\pgfqpoint{11.055926in}{1.730670in}}%
\pgfpathlineto{\pgfqpoint{11.055926in}{1.691499in}}%
\pgfpathquadraticcurveto{\pgfqpoint{11.055926in}{1.682676in}}{\pgfqpoint{11.058339in}{1.680160in}}%
\pgfpathquadraticcurveto{\pgfqpoint{11.060751in}{1.677645in}}{\pgfqpoint{11.068173in}{1.677645in}}%
\pgfpathlineto{\pgfqpoint{11.080336in}{1.677645in}}%
\pgfpathlineto{\pgfqpoint{11.080336in}{1.667729in}}%
\pgfpathlineto{\pgfqpoint{11.068173in}{1.667729in}}%
\pgfpathquadraticcurveto{\pgfqpoint{11.054442in}{1.667729in}}{\pgfqpoint{11.049226in}{1.672842in}}%
\pgfpathquadraticcurveto{\pgfqpoint{11.044010in}{1.677975in}}{\pgfqpoint{11.044010in}{1.691499in}}%
\pgfpathlineto{\pgfqpoint{11.044010in}{1.730670in}}%
\pgfpathlineto{\pgfqpoint{11.035310in}{1.730670in}}%
\pgfpathlineto{\pgfqpoint{11.035310in}{1.739886in}}%
\pgfpathlineto{\pgfqpoint{11.044010in}{1.739886in}}%
\pgfpathlineto{\pgfqpoint{11.044010in}{1.760379in}}%
\pgfpathlineto{\pgfqpoint{11.055926in}{1.760379in}}%
\pgfpathlineto{\pgfqpoint{11.055926in}{1.760379in}}%
\pgfpathclose%
\pgfpathmoveto{\pgfqpoint{11.094706in}{1.696200in}}%
\pgfpathlineto{\pgfqpoint{11.094706in}{1.739886in}}%
\pgfpathlineto{\pgfqpoint{11.106560in}{1.739886in}}%
\pgfpathlineto{\pgfqpoint{11.106560in}{1.696653in}}%
\pgfpathquadraticcurveto{\pgfqpoint{11.106560in}{1.686407in}}{\pgfqpoint{11.110539in}{1.681274in}}%
\pgfpathquadraticcurveto{\pgfqpoint{11.114539in}{1.676161in}}{\pgfqpoint{11.122538in}{1.676161in}}%
\pgfpathquadraticcurveto{\pgfqpoint{11.132124in}{1.676161in}}{\pgfqpoint{11.137691in}{1.682284in}}%
\pgfpathquadraticcurveto{\pgfqpoint{11.143278in}{1.688407in}}{\pgfqpoint{11.143278in}{1.698983in}}%
\pgfpathlineto{\pgfqpoint{11.143278in}{1.739886in}}%
\pgfpathlineto{\pgfqpoint{11.155132in}{1.739886in}}%
\pgfpathlineto{\pgfqpoint{11.155132in}{1.667729in}}%
\pgfpathlineto{\pgfqpoint{11.143278in}{1.667729in}}%
\pgfpathlineto{\pgfqpoint{11.143278in}{1.678820in}}%
\pgfpathquadraticcurveto{\pgfqpoint{11.138969in}{1.672244in}}{\pgfqpoint{11.133258in}{1.669048in}}%
\pgfpathquadraticcurveto{\pgfqpoint{11.127568in}{1.665853in}}{\pgfqpoint{11.120023in}{1.665853in}}%
\pgfpathquadraticcurveto{\pgfqpoint{11.107591in}{1.665853in}}{\pgfqpoint{11.101138in}{1.673584in}}%
\pgfpathquadraticcurveto{\pgfqpoint{11.094706in}{1.681315in}}{\pgfqpoint{11.094706in}{1.696200in}}%
\pgfpathlineto{\pgfqpoint{11.094706in}{1.696200in}}%
\pgfpathclose%
\pgfpathmoveto{\pgfqpoint{11.124538in}{1.741618in}}%
\pgfpathlineto{\pgfqpoint{11.124538in}{1.741618in}}%
\pgfpathlineto{\pgfqpoint{11.124538in}{1.741618in}}%
\pgfpathclose%
\pgfpathmoveto{\pgfqpoint{11.227021in}{1.728939in}}%
\pgfpathlineto{\pgfqpoint{11.227021in}{1.767986in}}%
\pgfpathlineto{\pgfqpoint{11.238876in}{1.767986in}}%
\pgfpathlineto{\pgfqpoint{11.238876in}{1.667729in}}%
\pgfpathlineto{\pgfqpoint{11.227021in}{1.667729in}}%
\pgfpathlineto{\pgfqpoint{11.227021in}{1.678552in}}%
\pgfpathquadraticcurveto{\pgfqpoint{11.223290in}{1.672120in}}{\pgfqpoint{11.217579in}{1.668986in}}%
\pgfpathquadraticcurveto{\pgfqpoint{11.211889in}{1.665853in}}{\pgfqpoint{11.203890in}{1.665853in}}%
\pgfpathquadraticcurveto{\pgfqpoint{11.190819in}{1.665853in}}{\pgfqpoint{11.182593in}{1.676285in}}%
\pgfpathquadraticcurveto{\pgfqpoint{11.174388in}{1.686737in}}{\pgfqpoint{11.174388in}{1.703746in}}%
\pgfpathquadraticcurveto{\pgfqpoint{11.174388in}{1.720754in}}{\pgfqpoint{11.182593in}{1.731186in}}%
\pgfpathquadraticcurveto{\pgfqpoint{11.190819in}{1.741618in}}{\pgfqpoint{11.203890in}{1.741618in}}%
\pgfpathquadraticcurveto{\pgfqpoint{11.211889in}{1.741618in}}{\pgfqpoint{11.217579in}{1.738484in}}%
\pgfpathquadraticcurveto{\pgfqpoint{11.223290in}{1.735371in}}{\pgfqpoint{11.227021in}{1.728939in}}%
\pgfpathlineto{\pgfqpoint{11.227021in}{1.728939in}}%
\pgfpathclose%
\pgfpathmoveto{\pgfqpoint{11.186634in}{1.703746in}}%
\pgfpathquadraticcurveto{\pgfqpoint{11.186634in}{1.690675in}}{\pgfqpoint{11.192015in}{1.683232in}}%
\pgfpathquadraticcurveto{\pgfqpoint{11.197396in}{1.675790in}}{\pgfqpoint{11.206797in}{1.675790in}}%
\pgfpathquadraticcurveto{\pgfqpoint{11.216198in}{1.675790in}}{\pgfqpoint{11.221599in}{1.683232in}}%
\pgfpathquadraticcurveto{\pgfqpoint{11.227021in}{1.690675in}}{\pgfqpoint{11.227021in}{1.703746in}}%
\pgfpathquadraticcurveto{\pgfqpoint{11.227021in}{1.716816in}}{\pgfqpoint{11.221599in}{1.724259in}}%
\pgfpathquadraticcurveto{\pgfqpoint{11.216198in}{1.731701in}}{\pgfqpoint{11.206797in}{1.731701in}}%
\pgfpathquadraticcurveto{\pgfqpoint{11.197396in}{1.731701in}}{\pgfqpoint{11.192015in}{1.724259in}}%
\pgfpathquadraticcurveto{\pgfqpoint{11.186634in}{1.716816in}}{\pgfqpoint{11.186634in}{1.703746in}}%
\pgfpathlineto{\pgfqpoint{11.186634in}{1.703746in}}%
\pgfpathclose%
\pgfpathmoveto{\pgfqpoint{11.293323in}{1.661028in}}%
\pgfpathquadraticcurveto{\pgfqpoint{11.288314in}{1.648143in}}{\pgfqpoint{11.283531in}{1.644226in}}%
\pgfpathquadraticcurveto{\pgfqpoint{11.278768in}{1.640288in}}{\pgfqpoint{11.270790in}{1.640288in}}%
\pgfpathlineto{\pgfqpoint{11.261306in}{1.640288in}}%
\pgfpathlineto{\pgfqpoint{11.261306in}{1.650205in}}%
\pgfpathlineto{\pgfqpoint{11.268275in}{1.650205in}}%
\pgfpathquadraticcurveto{\pgfqpoint{11.273161in}{1.650205in}}{\pgfqpoint{11.275861in}{1.652535in}}%
\pgfpathquadraticcurveto{\pgfqpoint{11.278583in}{1.654844in}}{\pgfqpoint{11.281861in}{1.663482in}}%
\pgfpathlineto{\pgfqpoint{11.283984in}{1.668883in}}%
\pgfpathlineto{\pgfqpoint{11.254812in}{1.739886in}}%
\pgfpathlineto{\pgfqpoint{11.267367in}{1.739886in}}%
\pgfpathlineto{\pgfqpoint{11.289922in}{1.683459in}}%
\pgfpathlineto{\pgfqpoint{11.312476in}{1.739886in}}%
\pgfpathlineto{\pgfqpoint{11.325031in}{1.739886in}}%
\pgfpathlineto{\pgfqpoint{11.293323in}{1.661028in}}%
\pgfpathlineto{\pgfqpoint{11.293323in}{1.661028in}}%
\pgfpathclose%
\pgfpathmoveto{\pgfqpoint{11.396219in}{1.678676in}}%
\pgfpathlineto{\pgfqpoint{11.441637in}{1.678676in}}%
\pgfpathlineto{\pgfqpoint{11.441637in}{1.667729in}}%
\pgfpathlineto{\pgfqpoint{11.380572in}{1.667729in}}%
\pgfpathlineto{\pgfqpoint{11.380572in}{1.678676in}}%
\pgfpathquadraticcurveto{\pgfqpoint{11.387973in}{1.686345in}}{\pgfqpoint{11.400755in}{1.699251in}}%
\pgfpathquadraticcurveto{\pgfqpoint{11.413558in}{1.712178in}}{\pgfqpoint{11.416836in}{1.715930in}}%
\pgfpathquadraticcurveto{\pgfqpoint{11.423082in}{1.722939in}}{\pgfqpoint{11.425556in}{1.727805in}}%
\pgfpathquadraticcurveto{\pgfqpoint{11.428051in}{1.732670in}}{\pgfqpoint{11.428051in}{1.737371in}}%
\pgfpathquadraticcurveto{\pgfqpoint{11.428051in}{1.745040in}}{\pgfqpoint{11.422670in}{1.749864in}}%
\pgfpathquadraticcurveto{\pgfqpoint{11.417289in}{1.754709in}}{\pgfqpoint{11.408651in}{1.754709in}}%
\pgfpathquadraticcurveto{\pgfqpoint{11.402528in}{1.754709in}}{\pgfqpoint{11.395725in}{1.752586in}}%
\pgfpathquadraticcurveto{\pgfqpoint{11.388942in}{1.750462in}}{\pgfqpoint{11.381211in}{1.746133in}}%
\pgfpathlineto{\pgfqpoint{11.381211in}{1.759286in}}%
\pgfpathquadraticcurveto{\pgfqpoint{11.389066in}{1.762440in}}{\pgfqpoint{11.395890in}{1.764048in}}%
\pgfpathquadraticcurveto{\pgfqpoint{11.402734in}{1.765656in}}{\pgfqpoint{11.408404in}{1.765656in}}%
\pgfpathquadraticcurveto{\pgfqpoint{11.423350in}{1.765656in}}{\pgfqpoint{11.432236in}{1.758173in}}%
\pgfpathquadraticcurveto{\pgfqpoint{11.441122in}{1.750709in}}{\pgfqpoint{11.441122in}{1.738216in}}%
\pgfpathquadraticcurveto{\pgfqpoint{11.441122in}{1.732278in}}{\pgfqpoint{11.438895in}{1.726959in}}%
\pgfpathquadraticcurveto{\pgfqpoint{11.436689in}{1.721661in}}{\pgfqpoint{11.430814in}{1.714445in}}%
\pgfpathquadraticcurveto{\pgfqpoint{11.429206in}{1.712569in}}{\pgfqpoint{11.420567in}{1.703642in}}%
\pgfpathquadraticcurveto{\pgfqpoint{11.411950in}{1.694716in}}{\pgfqpoint{11.396219in}{1.678676in}}%
\pgfpathlineto{\pgfqpoint{11.396219in}{1.678676in}}%
\pgfpathclose%
\pgfpathmoveto{\pgfqpoint{11.508393in}{1.719599in}}%
\pgfpathquadraticcurveto{\pgfqpoint{11.517732in}{1.717600in}}{\pgfqpoint{11.522969in}{1.711270in}}%
\pgfpathquadraticcurveto{\pgfqpoint{11.528226in}{1.704962in}}{\pgfqpoint{11.528226in}{1.695685in}}%
\pgfpathquadraticcurveto{\pgfqpoint{11.528226in}{1.681459in}}{\pgfqpoint{11.518433in}{1.673646in}}%
\pgfpathquadraticcurveto{\pgfqpoint{11.508640in}{1.665853in}}{\pgfqpoint{11.490601in}{1.665853in}}%
\pgfpathquadraticcurveto{\pgfqpoint{11.484560in}{1.665853in}}{\pgfqpoint{11.478149in}{1.667048in}}%
\pgfpathquadraticcurveto{\pgfqpoint{11.471737in}{1.668244in}}{\pgfqpoint{11.464913in}{1.670636in}}%
\pgfpathlineto{\pgfqpoint{11.464913in}{1.683191in}}%
\pgfpathquadraticcurveto{\pgfqpoint{11.470314in}{1.680037in}}{\pgfqpoint{11.476747in}{1.678429in}}%
\pgfpathquadraticcurveto{\pgfqpoint{11.483200in}{1.676821in}}{\pgfqpoint{11.490230in}{1.676821in}}%
\pgfpathquadraticcurveto{\pgfqpoint{11.502455in}{1.676821in}}{\pgfqpoint{11.508867in}{1.681645in}}%
\pgfpathquadraticcurveto{\pgfqpoint{11.515279in}{1.686469in}}{\pgfqpoint{11.515279in}{1.695685in}}%
\pgfpathquadraticcurveto{\pgfqpoint{11.515279in}{1.704199in}}{\pgfqpoint{11.509321in}{1.708982in}}%
\pgfpathquadraticcurveto{\pgfqpoint{11.503362in}{1.713786in}}{\pgfqpoint{11.492745in}{1.713786in}}%
\pgfpathlineto{\pgfqpoint{11.481530in}{1.713786in}}%
\pgfpathlineto{\pgfqpoint{11.481530in}{1.724486in}}%
\pgfpathlineto{\pgfqpoint{11.493260in}{1.724486in}}%
\pgfpathquadraticcurveto{\pgfqpoint{11.502847in}{1.724486in}}{\pgfqpoint{11.507939in}{1.728320in}}%
\pgfpathquadraticcurveto{\pgfqpoint{11.513031in}{1.732155in}}{\pgfqpoint{11.513031in}{1.739371in}}%
\pgfpathquadraticcurveto{\pgfqpoint{11.513031in}{1.746772in}}{\pgfqpoint{11.507774in}{1.750730in}}%
\pgfpathquadraticcurveto{\pgfqpoint{11.502538in}{1.754709in}}{\pgfqpoint{11.492745in}{1.754709in}}%
\pgfpathquadraticcurveto{\pgfqpoint{11.487385in}{1.754709in}}{\pgfqpoint{11.481262in}{1.753534in}}%
\pgfpathquadraticcurveto{\pgfqpoint{11.475139in}{1.752379in}}{\pgfqpoint{11.467799in}{1.749947in}}%
\pgfpathlineto{\pgfqpoint{11.467799in}{1.761533in}}%
\pgfpathquadraticcurveto{\pgfqpoint{11.475221in}{1.763595in}}{\pgfqpoint{11.481695in}{1.764625in}}%
\pgfpathquadraticcurveto{\pgfqpoint{11.488168in}{1.765656in}}{\pgfqpoint{11.493900in}{1.765656in}}%
\pgfpathquadraticcurveto{\pgfqpoint{11.508723in}{1.765656in}}{\pgfqpoint{11.517340in}{1.758915in}}%
\pgfpathquadraticcurveto{\pgfqpoint{11.525979in}{1.752194in}}{\pgfqpoint{11.525979in}{1.740731in}}%
\pgfpathquadraticcurveto{\pgfqpoint{11.525979in}{1.732732in}}{\pgfqpoint{11.521402in}{1.727227in}}%
\pgfpathquadraticcurveto{\pgfqpoint{11.516825in}{1.721723in}}{\pgfqpoint{11.508393in}{1.719599in}}%
\pgfpathlineto{\pgfqpoint{11.508393in}{1.719599in}}%
\pgfpathclose%
\pgfpathmoveto{\pgfqpoint{11.621638in}{1.767842in}}%
\pgfpathquadraticcurveto{\pgfqpoint{11.613021in}{1.753039in}}{\pgfqpoint{11.608815in}{1.738525in}}%
\pgfpathquadraticcurveto{\pgfqpoint{11.604630in}{1.724032in}}{\pgfqpoint{11.604630in}{1.709147in}}%
\pgfpathquadraticcurveto{\pgfqpoint{11.604630in}{1.694283in}}{\pgfqpoint{11.608856in}{1.679686in}}%
\pgfpathquadraticcurveto{\pgfqpoint{11.613082in}{1.665090in}}{\pgfqpoint{11.621638in}{1.650329in}}%
\pgfpathlineto{\pgfqpoint{11.611330in}{1.650329in}}%
\pgfpathquadraticcurveto{\pgfqpoint{11.601682in}{1.665482in}}{\pgfqpoint{11.596878in}{1.680099in}}%
\pgfpathquadraticcurveto{\pgfqpoint{11.592074in}{1.694716in}}{\pgfqpoint{11.592074in}{1.709147in}}%
\pgfpathquadraticcurveto{\pgfqpoint{11.592074in}{1.723517in}}{\pgfqpoint{11.596837in}{1.738072in}}%
\pgfpathquadraticcurveto{\pgfqpoint{11.601620in}{1.752647in}}{\pgfqpoint{11.611330in}{1.767842in}}%
\pgfpathlineto{\pgfqpoint{11.621638in}{1.767842in}}%
\pgfpathlineto{\pgfqpoint{11.621638in}{1.767842in}}%
\pgfpathclose%
\pgfpathmoveto{\pgfqpoint{11.657531in}{1.678676in}}%
\pgfpathlineto{\pgfqpoint{11.702949in}{1.678676in}}%
\pgfpathlineto{\pgfqpoint{11.702949in}{1.667729in}}%
\pgfpathlineto{\pgfqpoint{11.641883in}{1.667729in}}%
\pgfpathlineto{\pgfqpoint{11.641883in}{1.678676in}}%
\pgfpathquadraticcurveto{\pgfqpoint{11.649285in}{1.686345in}}{\pgfqpoint{11.662067in}{1.699251in}}%
\pgfpathquadraticcurveto{\pgfqpoint{11.674870in}{1.712178in}}{\pgfqpoint{11.678148in}{1.715930in}}%
\pgfpathquadraticcurveto{\pgfqpoint{11.684394in}{1.722939in}}{\pgfqpoint{11.686868in}{1.727805in}}%
\pgfpathquadraticcurveto{\pgfqpoint{11.689363in}{1.732670in}}{\pgfqpoint{11.689363in}{1.737371in}}%
\pgfpathquadraticcurveto{\pgfqpoint{11.689363in}{1.745040in}}{\pgfqpoint{11.683982in}{1.749864in}}%
\pgfpathquadraticcurveto{\pgfqpoint{11.678601in}{1.754709in}}{\pgfqpoint{11.669963in}{1.754709in}}%
\pgfpathquadraticcurveto{\pgfqpoint{11.663840in}{1.754709in}}{\pgfqpoint{11.657036in}{1.752586in}}%
\pgfpathquadraticcurveto{\pgfqpoint{11.650254in}{1.750462in}}{\pgfqpoint{11.642523in}{1.746133in}}%
\pgfpathlineto{\pgfqpoint{11.642523in}{1.759286in}}%
\pgfpathquadraticcurveto{\pgfqpoint{11.650377in}{1.762440in}}{\pgfqpoint{11.657201in}{1.764048in}}%
\pgfpathquadraticcurveto{\pgfqpoint{11.664046in}{1.765656in}}{\pgfqpoint{11.669715in}{1.765656in}}%
\pgfpathquadraticcurveto{\pgfqpoint{11.684662in}{1.765656in}}{\pgfqpoint{11.693548in}{1.758173in}}%
\pgfpathquadraticcurveto{\pgfqpoint{11.702434in}{1.750709in}}{\pgfqpoint{11.702434in}{1.738216in}}%
\pgfpathquadraticcurveto{\pgfqpoint{11.702434in}{1.732278in}}{\pgfqpoint{11.700207in}{1.726959in}}%
\pgfpathquadraticcurveto{\pgfqpoint{11.698001in}{1.721661in}}{\pgfqpoint{11.692125in}{1.714445in}}%
\pgfpathquadraticcurveto{\pgfqpoint{11.690517in}{1.712569in}}{\pgfqpoint{11.681879in}{1.703642in}}%
\pgfpathquadraticcurveto{\pgfqpoint{11.673262in}{1.694716in}}{\pgfqpoint{11.657531in}{1.678676in}}%
\pgfpathlineto{\pgfqpoint{11.657531in}{1.678676in}}%
\pgfpathclose%
\pgfpathmoveto{\pgfqpoint{11.758098in}{1.755348in}}%
\pgfpathquadraticcurveto{\pgfqpoint{11.748058in}{1.755348in}}{\pgfqpoint{11.742986in}{1.745452in}}%
\pgfpathquadraticcurveto{\pgfqpoint{11.737935in}{1.735577in}}{\pgfqpoint{11.737935in}{1.715724in}}%
\pgfpathquadraticcurveto{\pgfqpoint{11.737935in}{1.695953in}}{\pgfqpoint{11.742986in}{1.686057in}}%
\pgfpathquadraticcurveto{\pgfqpoint{11.748058in}{1.676161in}}{\pgfqpoint{11.758098in}{1.676161in}}%
\pgfpathquadraticcurveto{\pgfqpoint{11.768220in}{1.676161in}}{\pgfqpoint{11.773271in}{1.686057in}}%
\pgfpathquadraticcurveto{\pgfqpoint{11.778343in}{1.695953in}}{\pgfqpoint{11.778343in}{1.715724in}}%
\pgfpathquadraticcurveto{\pgfqpoint{11.778343in}{1.735577in}}{\pgfqpoint{11.773271in}{1.745452in}}%
\pgfpathquadraticcurveto{\pgfqpoint{11.768220in}{1.755348in}}{\pgfqpoint{11.758098in}{1.755348in}}%
\pgfpathlineto{\pgfqpoint{11.758098in}{1.755348in}}%
\pgfpathclose%
\pgfpathmoveto{\pgfqpoint{11.758098in}{1.765656in}}%
\pgfpathquadraticcurveto{\pgfqpoint{11.774281in}{1.765656in}}{\pgfqpoint{11.782817in}{1.752854in}}%
\pgfpathquadraticcurveto{\pgfqpoint{11.791352in}{1.740071in}}{\pgfqpoint{11.791352in}{1.715724in}}%
\pgfpathquadraticcurveto{\pgfqpoint{11.791352in}{1.691438in}}{\pgfqpoint{11.782817in}{1.678635in}}%
\pgfpathquadraticcurveto{\pgfqpoint{11.774281in}{1.665853in}}{\pgfqpoint{11.758098in}{1.665853in}}%
\pgfpathquadraticcurveto{\pgfqpoint{11.741934in}{1.665853in}}{\pgfqpoint{11.733399in}{1.678635in}}%
\pgfpathquadraticcurveto{\pgfqpoint{11.724864in}{1.691438in}}{\pgfqpoint{11.724864in}{1.715724in}}%
\pgfpathquadraticcurveto{\pgfqpoint{11.724864in}{1.740071in}}{\pgfqpoint{11.733399in}{1.752854in}}%
\pgfpathquadraticcurveto{\pgfqpoint{11.741934in}{1.765656in}}{\pgfqpoint{11.758098in}{1.765656in}}%
\pgfpathlineto{\pgfqpoint{11.758098in}{1.765656in}}%
\pgfpathclose%
\pgfpathmoveto{\pgfqpoint{11.825431in}{1.678676in}}%
\pgfpathlineto{\pgfqpoint{11.870848in}{1.678676in}}%
\pgfpathlineto{\pgfqpoint{11.870848in}{1.667729in}}%
\pgfpathlineto{\pgfqpoint{11.809783in}{1.667729in}}%
\pgfpathlineto{\pgfqpoint{11.809783in}{1.678676in}}%
\pgfpathquadraticcurveto{\pgfqpoint{11.817184in}{1.686345in}}{\pgfqpoint{11.829966in}{1.699251in}}%
\pgfpathquadraticcurveto{\pgfqpoint{11.842769in}{1.712178in}}{\pgfqpoint{11.846047in}{1.715930in}}%
\pgfpathquadraticcurveto{\pgfqpoint{11.852294in}{1.722939in}}{\pgfqpoint{11.854768in}{1.727805in}}%
\pgfpathquadraticcurveto{\pgfqpoint{11.857262in}{1.732670in}}{\pgfqpoint{11.857262in}{1.737371in}}%
\pgfpathquadraticcurveto{\pgfqpoint{11.857262in}{1.745040in}}{\pgfqpoint{11.851881in}{1.749864in}}%
\pgfpathquadraticcurveto{\pgfqpoint{11.846500in}{1.754709in}}{\pgfqpoint{11.837862in}{1.754709in}}%
\pgfpathquadraticcurveto{\pgfqpoint{11.831739in}{1.754709in}}{\pgfqpoint{11.824936in}{1.752586in}}%
\pgfpathquadraticcurveto{\pgfqpoint{11.818153in}{1.750462in}}{\pgfqpoint{11.810422in}{1.746133in}}%
\pgfpathlineto{\pgfqpoint{11.810422in}{1.759286in}}%
\pgfpathquadraticcurveto{\pgfqpoint{11.818277in}{1.762440in}}{\pgfqpoint{11.825101in}{1.764048in}}%
\pgfpathquadraticcurveto{\pgfqpoint{11.831945in}{1.765656in}}{\pgfqpoint{11.837615in}{1.765656in}}%
\pgfpathquadraticcurveto{\pgfqpoint{11.852562in}{1.765656in}}{\pgfqpoint{11.861447in}{1.758173in}}%
\pgfpathquadraticcurveto{\pgfqpoint{11.870333in}{1.750709in}}{\pgfqpoint{11.870333in}{1.738216in}}%
\pgfpathquadraticcurveto{\pgfqpoint{11.870333in}{1.732278in}}{\pgfqpoint{11.868106in}{1.726959in}}%
\pgfpathquadraticcurveto{\pgfqpoint{11.865900in}{1.721661in}}{\pgfqpoint{11.860025in}{1.714445in}}%
\pgfpathquadraticcurveto{\pgfqpoint{11.858417in}{1.712569in}}{\pgfqpoint{11.849778in}{1.703642in}}%
\pgfpathquadraticcurveto{\pgfqpoint{11.841161in}{1.694716in}}{\pgfqpoint{11.825431in}{1.678676in}}%
\pgfpathlineto{\pgfqpoint{11.825431in}{1.678676in}}%
\pgfpathclose%
\pgfpathmoveto{\pgfqpoint{11.898309in}{1.763925in}}%
\pgfpathlineto{\pgfqpoint{11.949396in}{1.763925in}}%
\pgfpathlineto{\pgfqpoint{11.949396in}{1.752957in}}%
\pgfpathlineto{\pgfqpoint{11.910225in}{1.752957in}}%
\pgfpathlineto{\pgfqpoint{11.910225in}{1.729392in}}%
\pgfpathquadraticcurveto{\pgfqpoint{11.913050in}{1.730361in}}{\pgfqpoint{11.915874in}{1.730835in}}%
\pgfpathquadraticcurveto{\pgfqpoint{11.918719in}{1.731310in}}{\pgfqpoint{11.921564in}{1.731310in}}%
\pgfpathquadraticcurveto{\pgfqpoint{11.937666in}{1.731310in}}{\pgfqpoint{11.947067in}{1.722486in}}%
\pgfpathquadraticcurveto{\pgfqpoint{11.956488in}{1.713662in}}{\pgfqpoint{11.956488in}{1.698591in}}%
\pgfpathquadraticcurveto{\pgfqpoint{11.956488in}{1.683067in}}{\pgfqpoint{11.946819in}{1.674450in}}%
\pgfpathquadraticcurveto{\pgfqpoint{11.937150in}{1.665853in}}{\pgfqpoint{11.919565in}{1.665853in}}%
\pgfpathquadraticcurveto{\pgfqpoint{11.913503in}{1.665853in}}{\pgfqpoint{11.907215in}{1.666884in}}%
\pgfpathquadraticcurveto{\pgfqpoint{11.900948in}{1.667914in}}{\pgfqpoint{11.894248in}{1.669976in}}%
\pgfpathlineto{\pgfqpoint{11.894248in}{1.683067in}}%
\pgfpathquadraticcurveto{\pgfqpoint{11.900041in}{1.679913in}}{\pgfqpoint{11.906226in}{1.678367in}}%
\pgfpathquadraticcurveto{\pgfqpoint{11.912411in}{1.676821in}}{\pgfqpoint{11.919297in}{1.676821in}}%
\pgfpathquadraticcurveto{\pgfqpoint{11.930450in}{1.676821in}}{\pgfqpoint{11.936944in}{1.682676in}}%
\pgfpathquadraticcurveto{\pgfqpoint{11.943459in}{1.688531in}}{\pgfqpoint{11.943459in}{1.698591in}}%
\pgfpathquadraticcurveto{\pgfqpoint{11.943459in}{1.708632in}}{\pgfqpoint{11.936944in}{1.714487in}}%
\pgfpathquadraticcurveto{\pgfqpoint{11.930450in}{1.720362in}}{\pgfqpoint{11.919297in}{1.720362in}}%
\pgfpathquadraticcurveto{\pgfqpoint{11.914081in}{1.720362in}}{\pgfqpoint{11.908885in}{1.719208in}}%
\pgfpathquadraticcurveto{\pgfqpoint{11.903711in}{1.718053in}}{\pgfqpoint{11.898309in}{1.715600in}}%
\pgfpathlineto{\pgfqpoint{11.898309in}{1.763925in}}%
\pgfpathlineto{\pgfqpoint{11.898309in}{1.763925in}}%
\pgfpathclose%
\pgfpathmoveto{\pgfqpoint{11.978589in}{1.767842in}}%
\pgfpathlineto{\pgfqpoint{11.988897in}{1.767842in}}%
\pgfpathquadraticcurveto{\pgfqpoint{11.998546in}{1.752647in}}{\pgfqpoint{12.003349in}{1.738072in}}%
\pgfpathquadraticcurveto{\pgfqpoint{12.008153in}{1.723517in}}{\pgfqpoint{12.008153in}{1.709147in}}%
\pgfpathquadraticcurveto{\pgfqpoint{12.008153in}{1.694716in}}{\pgfqpoint{12.003349in}{1.680099in}}%
\pgfpathquadraticcurveto{\pgfqpoint{11.998546in}{1.665482in}}{\pgfqpoint{11.988897in}{1.650329in}}%
\pgfpathlineto{\pgfqpoint{11.978589in}{1.650329in}}%
\pgfpathquadraticcurveto{\pgfqpoint{11.987145in}{1.665090in}}{\pgfqpoint{11.991371in}{1.679686in}}%
\pgfpathquadraticcurveto{\pgfqpoint{11.995598in}{1.694283in}}{\pgfqpoint{11.995598in}{1.709147in}}%
\pgfpathquadraticcurveto{\pgfqpoint{11.995598in}{1.724032in}}{\pgfqpoint{11.991371in}{1.738525in}}%
\pgfpathquadraticcurveto{\pgfqpoint{11.987145in}{1.753039in}}{\pgfqpoint{11.978589in}{1.767842in}}%
\pgfpathlineto{\pgfqpoint{11.978589in}{1.767842in}}%
\pgfpathclose%
\pgfusepath{fill}%
\end{pgfscope}%
\begin{pgfscope}%
\pgfsetbuttcap%
\pgfsetroundjoin%
\definecolor{currentfill}{rgb}{0.090196,0.168627,0.227451}%
\pgfsetfillcolor{currentfill}%
\pgfsetlinewidth{0.000000pt}%
\definecolor{currentstroke}{rgb}{0.090196,0.168627,0.227451}%
\pgfsetstrokecolor{currentstroke}%
\pgfsetdash{}{0pt}%
\pgfpathmoveto{\pgfqpoint{11.018611in}{1.584130in}}%
\pgfpathlineto{\pgfqpoint{11.018611in}{1.571431in}}%
\pgfpathquadraticcurveto{\pgfqpoint{11.011210in}{1.574977in}}{\pgfqpoint{11.004633in}{1.576708in}}%
\pgfpathquadraticcurveto{\pgfqpoint{10.998056in}{1.578461in}}{\pgfqpoint{10.991933in}{1.578461in}}%
\pgfpathquadraticcurveto{\pgfqpoint{10.981316in}{1.578461in}}{\pgfqpoint{10.975543in}{1.574337in}}%
\pgfpathquadraticcurveto{\pgfqpoint{10.969771in}{1.570214in}}{\pgfqpoint{10.969771in}{1.562607in}}%
\pgfpathquadraticcurveto{\pgfqpoint{10.969771in}{1.556216in}}{\pgfqpoint{10.973605in}{1.552958in}}%
\pgfpathquadraticcurveto{\pgfqpoint{10.977440in}{1.549722in}}{\pgfqpoint{10.988140in}{1.547722in}}%
\pgfpathlineto{\pgfqpoint{10.995995in}{1.546114in}}%
\pgfpathquadraticcurveto{\pgfqpoint{11.010550in}{1.543331in}}{\pgfqpoint{11.017477in}{1.536342in}}%
\pgfpathquadraticcurveto{\pgfqpoint{11.024404in}{1.529353in}}{\pgfqpoint{11.024404in}{1.517643in}}%
\pgfpathquadraticcurveto{\pgfqpoint{11.024404in}{1.503644in}}{\pgfqpoint{11.015024in}{1.496428in}}%
\pgfpathquadraticcurveto{\pgfqpoint{11.005664in}{1.489213in}}{\pgfqpoint{10.987563in}{1.489213in}}%
\pgfpathquadraticcurveto{\pgfqpoint{10.980739in}{1.489213in}}{\pgfqpoint{10.973028in}{1.490759in}}%
\pgfpathquadraticcurveto{\pgfqpoint{10.965338in}{1.492305in}}{\pgfqpoint{10.957092in}{1.495336in}}%
\pgfpathlineto{\pgfqpoint{10.957092in}{1.508736in}}%
\pgfpathquadraticcurveto{\pgfqpoint{10.965008in}{1.504304in}}{\pgfqpoint{10.972616in}{1.502036in}}%
\pgfpathquadraticcurveto{\pgfqpoint{10.980223in}{1.499789in}}{\pgfqpoint{10.987563in}{1.499789in}}%
\pgfpathquadraticcurveto{\pgfqpoint{10.998696in}{1.499789in}}{\pgfqpoint{11.004757in}{1.504160in}}%
\pgfpathquadraticcurveto{\pgfqpoint{11.010818in}{1.508551in}}{\pgfqpoint{11.010818in}{1.516674in}}%
\pgfpathquadraticcurveto{\pgfqpoint{11.010818in}{1.523745in}}{\pgfqpoint{11.006468in}{1.527745in}}%
\pgfpathquadraticcurveto{\pgfqpoint{11.002118in}{1.531744in}}{\pgfqpoint{10.992201in}{1.533744in}}%
\pgfpathlineto{\pgfqpoint{10.984264in}{1.535290in}}%
\pgfpathquadraticcurveto{\pgfqpoint{10.969709in}{1.538176in}}{\pgfqpoint{10.963194in}{1.544361in}}%
\pgfpathquadraticcurveto{\pgfqpoint{10.956700in}{1.550546in}}{\pgfqpoint{10.956700in}{1.561576in}}%
\pgfpathquadraticcurveto{\pgfqpoint{10.956700in}{1.574337in}}{\pgfqpoint{10.965689in}{1.581677in}}%
\pgfpathquadraticcurveto{\pgfqpoint{10.974678in}{1.589016in}}{\pgfqpoint{10.990449in}{1.589016in}}%
\pgfpathquadraticcurveto{\pgfqpoint{10.997232in}{1.589016in}}{\pgfqpoint{11.004241in}{1.587779in}}%
\pgfpathquadraticcurveto{\pgfqpoint{11.011271in}{1.586563in}}{\pgfqpoint{11.018611in}{1.584130in}}%
\pgfpathlineto{\pgfqpoint{11.018611in}{1.584130in}}%
\pgfpathclose%
\pgfpathmoveto{\pgfqpoint{11.055926in}{1.583739in}}%
\pgfpathlineto{\pgfqpoint{11.055926in}{1.563246in}}%
\pgfpathlineto{\pgfqpoint{11.080336in}{1.563246in}}%
\pgfpathlineto{\pgfqpoint{11.080336in}{1.554030in}}%
\pgfpathlineto{\pgfqpoint{11.055926in}{1.554030in}}%
\pgfpathlineto{\pgfqpoint{11.055926in}{1.514859in}}%
\pgfpathquadraticcurveto{\pgfqpoint{11.055926in}{1.506036in}}{\pgfqpoint{11.058339in}{1.503520in}}%
\pgfpathquadraticcurveto{\pgfqpoint{11.060751in}{1.501005in}}{\pgfqpoint{11.068173in}{1.501005in}}%
\pgfpathlineto{\pgfqpoint{11.080336in}{1.501005in}}%
\pgfpathlineto{\pgfqpoint{11.080336in}{1.491089in}}%
\pgfpathlineto{\pgfqpoint{11.068173in}{1.491089in}}%
\pgfpathquadraticcurveto{\pgfqpoint{11.054442in}{1.491089in}}{\pgfqpoint{11.049226in}{1.496202in}}%
\pgfpathquadraticcurveto{\pgfqpoint{11.044010in}{1.501335in}}{\pgfqpoint{11.044010in}{1.514859in}}%
\pgfpathlineto{\pgfqpoint{11.044010in}{1.554030in}}%
\pgfpathlineto{\pgfqpoint{11.035310in}{1.554030in}}%
\pgfpathlineto{\pgfqpoint{11.035310in}{1.563246in}}%
\pgfpathlineto{\pgfqpoint{11.044010in}{1.563246in}}%
\pgfpathlineto{\pgfqpoint{11.044010in}{1.583739in}}%
\pgfpathlineto{\pgfqpoint{11.055926in}{1.583739in}}%
\pgfpathlineto{\pgfqpoint{11.055926in}{1.583739in}}%
\pgfpathclose%
\pgfpathmoveto{\pgfqpoint{11.094706in}{1.519560in}}%
\pgfpathlineto{\pgfqpoint{11.094706in}{1.563246in}}%
\pgfpathlineto{\pgfqpoint{11.106560in}{1.563246in}}%
\pgfpathlineto{\pgfqpoint{11.106560in}{1.520013in}}%
\pgfpathquadraticcurveto{\pgfqpoint{11.106560in}{1.509767in}}{\pgfqpoint{11.110539in}{1.504634in}}%
\pgfpathquadraticcurveto{\pgfqpoint{11.114539in}{1.499521in}}{\pgfqpoint{11.122538in}{1.499521in}}%
\pgfpathquadraticcurveto{\pgfqpoint{11.132124in}{1.499521in}}{\pgfqpoint{11.137691in}{1.505644in}}%
\pgfpathquadraticcurveto{\pgfqpoint{11.143278in}{1.511767in}}{\pgfqpoint{11.143278in}{1.522343in}}%
\pgfpathlineto{\pgfqpoint{11.143278in}{1.563246in}}%
\pgfpathlineto{\pgfqpoint{11.155132in}{1.563246in}}%
\pgfpathlineto{\pgfqpoint{11.155132in}{1.491089in}}%
\pgfpathlineto{\pgfqpoint{11.143278in}{1.491089in}}%
\pgfpathlineto{\pgfqpoint{11.143278in}{1.502180in}}%
\pgfpathquadraticcurveto{\pgfqpoint{11.138969in}{1.495604in}}{\pgfqpoint{11.133258in}{1.492408in}}%
\pgfpathquadraticcurveto{\pgfqpoint{11.127568in}{1.489213in}}{\pgfqpoint{11.120023in}{1.489213in}}%
\pgfpathquadraticcurveto{\pgfqpoint{11.107591in}{1.489213in}}{\pgfqpoint{11.101138in}{1.496944in}}%
\pgfpathquadraticcurveto{\pgfqpoint{11.094706in}{1.504675in}}{\pgfqpoint{11.094706in}{1.519560in}}%
\pgfpathlineto{\pgfqpoint{11.094706in}{1.519560in}}%
\pgfpathclose%
\pgfpathmoveto{\pgfqpoint{11.124538in}{1.564978in}}%
\pgfpathlineto{\pgfqpoint{11.124538in}{1.564978in}}%
\pgfpathlineto{\pgfqpoint{11.124538in}{1.564978in}}%
\pgfpathclose%
\pgfpathmoveto{\pgfqpoint{11.227021in}{1.552299in}}%
\pgfpathlineto{\pgfqpoint{11.227021in}{1.591346in}}%
\pgfpathlineto{\pgfqpoint{11.238876in}{1.591346in}}%
\pgfpathlineto{\pgfqpoint{11.238876in}{1.491089in}}%
\pgfpathlineto{\pgfqpoint{11.227021in}{1.491089in}}%
\pgfpathlineto{\pgfqpoint{11.227021in}{1.501912in}}%
\pgfpathquadraticcurveto{\pgfqpoint{11.223290in}{1.495480in}}{\pgfqpoint{11.217579in}{1.492346in}}%
\pgfpathquadraticcurveto{\pgfqpoint{11.211889in}{1.489213in}}{\pgfqpoint{11.203890in}{1.489213in}}%
\pgfpathquadraticcurveto{\pgfqpoint{11.190819in}{1.489213in}}{\pgfqpoint{11.182593in}{1.499645in}}%
\pgfpathquadraticcurveto{\pgfqpoint{11.174388in}{1.510097in}}{\pgfqpoint{11.174388in}{1.527106in}}%
\pgfpathquadraticcurveto{\pgfqpoint{11.174388in}{1.544114in}}{\pgfqpoint{11.182593in}{1.554546in}}%
\pgfpathquadraticcurveto{\pgfqpoint{11.190819in}{1.564978in}}{\pgfqpoint{11.203890in}{1.564978in}}%
\pgfpathquadraticcurveto{\pgfqpoint{11.211889in}{1.564978in}}{\pgfqpoint{11.217579in}{1.561844in}}%
\pgfpathquadraticcurveto{\pgfqpoint{11.223290in}{1.558731in}}{\pgfqpoint{11.227021in}{1.552299in}}%
\pgfpathlineto{\pgfqpoint{11.227021in}{1.552299in}}%
\pgfpathclose%
\pgfpathmoveto{\pgfqpoint{11.186634in}{1.527106in}}%
\pgfpathquadraticcurveto{\pgfqpoint{11.186634in}{1.514035in}}{\pgfqpoint{11.192015in}{1.506592in}}%
\pgfpathquadraticcurveto{\pgfqpoint{11.197396in}{1.499150in}}{\pgfqpoint{11.206797in}{1.499150in}}%
\pgfpathquadraticcurveto{\pgfqpoint{11.216198in}{1.499150in}}{\pgfqpoint{11.221599in}{1.506592in}}%
\pgfpathquadraticcurveto{\pgfqpoint{11.227021in}{1.514035in}}{\pgfqpoint{11.227021in}{1.527106in}}%
\pgfpathquadraticcurveto{\pgfqpoint{11.227021in}{1.540176in}}{\pgfqpoint{11.221599in}{1.547619in}}%
\pgfpathquadraticcurveto{\pgfqpoint{11.216198in}{1.555061in}}{\pgfqpoint{11.206797in}{1.555061in}}%
\pgfpathquadraticcurveto{\pgfqpoint{11.197396in}{1.555061in}}{\pgfqpoint{11.192015in}{1.547619in}}%
\pgfpathquadraticcurveto{\pgfqpoint{11.186634in}{1.540176in}}{\pgfqpoint{11.186634in}{1.527106in}}%
\pgfpathlineto{\pgfqpoint{11.186634in}{1.527106in}}%
\pgfpathclose%
\pgfpathmoveto{\pgfqpoint{11.293323in}{1.484388in}}%
\pgfpathquadraticcurveto{\pgfqpoint{11.288314in}{1.471503in}}{\pgfqpoint{11.283531in}{1.467586in}}%
\pgfpathquadraticcurveto{\pgfqpoint{11.278768in}{1.463648in}}{\pgfqpoint{11.270790in}{1.463648in}}%
\pgfpathlineto{\pgfqpoint{11.261306in}{1.463648in}}%
\pgfpathlineto{\pgfqpoint{11.261306in}{1.473565in}}%
\pgfpathlineto{\pgfqpoint{11.268275in}{1.473565in}}%
\pgfpathquadraticcurveto{\pgfqpoint{11.273161in}{1.473565in}}{\pgfqpoint{11.275861in}{1.475895in}}%
\pgfpathquadraticcurveto{\pgfqpoint{11.278583in}{1.478204in}}{\pgfqpoint{11.281861in}{1.486842in}}%
\pgfpathlineto{\pgfqpoint{11.283984in}{1.492243in}}%
\pgfpathlineto{\pgfqpoint{11.254812in}{1.563246in}}%
\pgfpathlineto{\pgfqpoint{11.267367in}{1.563246in}}%
\pgfpathlineto{\pgfqpoint{11.289922in}{1.506819in}}%
\pgfpathlineto{\pgfqpoint{11.312476in}{1.563246in}}%
\pgfpathlineto{\pgfqpoint{11.325031in}{1.563246in}}%
\pgfpathlineto{\pgfqpoint{11.293323in}{1.484388in}}%
\pgfpathlineto{\pgfqpoint{11.293323in}{1.484388in}}%
\pgfpathclose%
\pgfpathmoveto{\pgfqpoint{11.396219in}{1.502036in}}%
\pgfpathlineto{\pgfqpoint{11.441637in}{1.502036in}}%
\pgfpathlineto{\pgfqpoint{11.441637in}{1.491089in}}%
\pgfpathlineto{\pgfqpoint{11.380572in}{1.491089in}}%
\pgfpathlineto{\pgfqpoint{11.380572in}{1.502036in}}%
\pgfpathquadraticcurveto{\pgfqpoint{11.387973in}{1.509705in}}{\pgfqpoint{11.400755in}{1.522611in}}%
\pgfpathquadraticcurveto{\pgfqpoint{11.413558in}{1.535538in}}{\pgfqpoint{11.416836in}{1.539290in}}%
\pgfpathquadraticcurveto{\pgfqpoint{11.423082in}{1.546299in}}{\pgfqpoint{11.425556in}{1.551165in}}%
\pgfpathquadraticcurveto{\pgfqpoint{11.428051in}{1.556030in}}{\pgfqpoint{11.428051in}{1.560731in}}%
\pgfpathquadraticcurveto{\pgfqpoint{11.428051in}{1.568400in}}{\pgfqpoint{11.422670in}{1.573224in}}%
\pgfpathquadraticcurveto{\pgfqpoint{11.417289in}{1.578069in}}{\pgfqpoint{11.408651in}{1.578069in}}%
\pgfpathquadraticcurveto{\pgfqpoint{11.402528in}{1.578069in}}{\pgfqpoint{11.395725in}{1.575946in}}%
\pgfpathquadraticcurveto{\pgfqpoint{11.388942in}{1.573822in}}{\pgfqpoint{11.381211in}{1.569493in}}%
\pgfpathlineto{\pgfqpoint{11.381211in}{1.582646in}}%
\pgfpathquadraticcurveto{\pgfqpoint{11.389066in}{1.585800in}}{\pgfqpoint{11.395890in}{1.587408in}}%
\pgfpathquadraticcurveto{\pgfqpoint{11.402734in}{1.589016in}}{\pgfqpoint{11.408404in}{1.589016in}}%
\pgfpathquadraticcurveto{\pgfqpoint{11.423350in}{1.589016in}}{\pgfqpoint{11.432236in}{1.581533in}}%
\pgfpathquadraticcurveto{\pgfqpoint{11.441122in}{1.574069in}}{\pgfqpoint{11.441122in}{1.561576in}}%
\pgfpathquadraticcurveto{\pgfqpoint{11.441122in}{1.555638in}}{\pgfqpoint{11.438895in}{1.550319in}}%
\pgfpathquadraticcurveto{\pgfqpoint{11.436689in}{1.545021in}}{\pgfqpoint{11.430814in}{1.537805in}}%
\pgfpathquadraticcurveto{\pgfqpoint{11.429206in}{1.535929in}}{\pgfqpoint{11.420567in}{1.527002in}}%
\pgfpathquadraticcurveto{\pgfqpoint{11.411950in}{1.518076in}}{\pgfqpoint{11.396219in}{1.502036in}}%
\pgfpathlineto{\pgfqpoint{11.396219in}{1.502036in}}%
\pgfpathclose%
\pgfpathmoveto{\pgfqpoint{11.504723in}{1.575946in}}%
\pgfpathlineto{\pgfqpoint{11.471861in}{1.524590in}}%
\pgfpathlineto{\pgfqpoint{11.504723in}{1.524590in}}%
\pgfpathlineto{\pgfqpoint{11.504723in}{1.575946in}}%
\pgfpathlineto{\pgfqpoint{11.504723in}{1.575946in}}%
\pgfpathclose%
\pgfpathmoveto{\pgfqpoint{11.501301in}{1.587285in}}%
\pgfpathlineto{\pgfqpoint{11.517670in}{1.587285in}}%
\pgfpathlineto{\pgfqpoint{11.517670in}{1.524590in}}%
\pgfpathlineto{\pgfqpoint{11.531401in}{1.524590in}}%
\pgfpathlineto{\pgfqpoint{11.531401in}{1.513767in}}%
\pgfpathlineto{\pgfqpoint{11.517670in}{1.513767in}}%
\pgfpathlineto{\pgfqpoint{11.517670in}{1.491089in}}%
\pgfpathlineto{\pgfqpoint{11.504723in}{1.491089in}}%
\pgfpathlineto{\pgfqpoint{11.504723in}{1.513767in}}%
\pgfpathlineto{\pgfqpoint{11.461305in}{1.513767in}}%
\pgfpathlineto{\pgfqpoint{11.461305in}{1.526322in}}%
\pgfpathlineto{\pgfqpoint{11.501301in}{1.587285in}}%
\pgfpathlineto{\pgfqpoint{11.501301in}{1.587285in}}%
\pgfpathclose%
\pgfpathmoveto{\pgfqpoint{11.621638in}{1.591202in}}%
\pgfpathquadraticcurveto{\pgfqpoint{11.613021in}{1.576399in}}{\pgfqpoint{11.608815in}{1.561885in}}%
\pgfpathquadraticcurveto{\pgfqpoint{11.604630in}{1.547392in}}{\pgfqpoint{11.604630in}{1.532507in}}%
\pgfpathquadraticcurveto{\pgfqpoint{11.604630in}{1.517643in}}{\pgfqpoint{11.608856in}{1.503046in}}%
\pgfpathquadraticcurveto{\pgfqpoint{11.613082in}{1.488450in}}{\pgfqpoint{11.621638in}{1.473689in}}%
\pgfpathlineto{\pgfqpoint{11.611330in}{1.473689in}}%
\pgfpathquadraticcurveto{\pgfqpoint{11.601682in}{1.488842in}}{\pgfqpoint{11.596878in}{1.503459in}}%
\pgfpathquadraticcurveto{\pgfqpoint{11.592074in}{1.518076in}}{\pgfqpoint{11.592074in}{1.532507in}}%
\pgfpathquadraticcurveto{\pgfqpoint{11.592074in}{1.546877in}}{\pgfqpoint{11.596837in}{1.561432in}}%
\pgfpathquadraticcurveto{\pgfqpoint{11.601620in}{1.576007in}}{\pgfqpoint{11.611330in}{1.591202in}}%
\pgfpathlineto{\pgfqpoint{11.621638in}{1.591202in}}%
\pgfpathlineto{\pgfqpoint{11.621638in}{1.591202in}}%
\pgfpathclose%
\pgfpathmoveto{\pgfqpoint{11.657531in}{1.502036in}}%
\pgfpathlineto{\pgfqpoint{11.702949in}{1.502036in}}%
\pgfpathlineto{\pgfqpoint{11.702949in}{1.491089in}}%
\pgfpathlineto{\pgfqpoint{11.641883in}{1.491089in}}%
\pgfpathlineto{\pgfqpoint{11.641883in}{1.502036in}}%
\pgfpathquadraticcurveto{\pgfqpoint{11.649285in}{1.509705in}}{\pgfqpoint{11.662067in}{1.522611in}}%
\pgfpathquadraticcurveto{\pgfqpoint{11.674870in}{1.535538in}}{\pgfqpoint{11.678148in}{1.539290in}}%
\pgfpathquadraticcurveto{\pgfqpoint{11.684394in}{1.546299in}}{\pgfqpoint{11.686868in}{1.551165in}}%
\pgfpathquadraticcurveto{\pgfqpoint{11.689363in}{1.556030in}}{\pgfqpoint{11.689363in}{1.560731in}}%
\pgfpathquadraticcurveto{\pgfqpoint{11.689363in}{1.568400in}}{\pgfqpoint{11.683982in}{1.573224in}}%
\pgfpathquadraticcurveto{\pgfqpoint{11.678601in}{1.578069in}}{\pgfqpoint{11.669963in}{1.578069in}}%
\pgfpathquadraticcurveto{\pgfqpoint{11.663840in}{1.578069in}}{\pgfqpoint{11.657036in}{1.575946in}}%
\pgfpathquadraticcurveto{\pgfqpoint{11.650254in}{1.573822in}}{\pgfqpoint{11.642523in}{1.569493in}}%
\pgfpathlineto{\pgfqpoint{11.642523in}{1.582646in}}%
\pgfpathquadraticcurveto{\pgfqpoint{11.650377in}{1.585800in}}{\pgfqpoint{11.657201in}{1.587408in}}%
\pgfpathquadraticcurveto{\pgfqpoint{11.664046in}{1.589016in}}{\pgfqpoint{11.669715in}{1.589016in}}%
\pgfpathquadraticcurveto{\pgfqpoint{11.684662in}{1.589016in}}{\pgfqpoint{11.693548in}{1.581533in}}%
\pgfpathquadraticcurveto{\pgfqpoint{11.702434in}{1.574069in}}{\pgfqpoint{11.702434in}{1.561576in}}%
\pgfpathquadraticcurveto{\pgfqpoint{11.702434in}{1.555638in}}{\pgfqpoint{11.700207in}{1.550319in}}%
\pgfpathquadraticcurveto{\pgfqpoint{11.698001in}{1.545021in}}{\pgfqpoint{11.692125in}{1.537805in}}%
\pgfpathquadraticcurveto{\pgfqpoint{11.690517in}{1.535929in}}{\pgfqpoint{11.681879in}{1.527002in}}%
\pgfpathquadraticcurveto{\pgfqpoint{11.673262in}{1.518076in}}{\pgfqpoint{11.657531in}{1.502036in}}%
\pgfpathlineto{\pgfqpoint{11.657531in}{1.502036in}}%
\pgfpathclose%
\pgfpathmoveto{\pgfqpoint{11.758098in}{1.578708in}}%
\pgfpathquadraticcurveto{\pgfqpoint{11.748058in}{1.578708in}}{\pgfqpoint{11.742986in}{1.568812in}}%
\pgfpathquadraticcurveto{\pgfqpoint{11.737935in}{1.558937in}}{\pgfqpoint{11.737935in}{1.539084in}}%
\pgfpathquadraticcurveto{\pgfqpoint{11.737935in}{1.519313in}}{\pgfqpoint{11.742986in}{1.509417in}}%
\pgfpathquadraticcurveto{\pgfqpoint{11.748058in}{1.499521in}}{\pgfqpoint{11.758098in}{1.499521in}}%
\pgfpathquadraticcurveto{\pgfqpoint{11.768220in}{1.499521in}}{\pgfqpoint{11.773271in}{1.509417in}}%
\pgfpathquadraticcurveto{\pgfqpoint{11.778343in}{1.519313in}}{\pgfqpoint{11.778343in}{1.539084in}}%
\pgfpathquadraticcurveto{\pgfqpoint{11.778343in}{1.558937in}}{\pgfqpoint{11.773271in}{1.568812in}}%
\pgfpathquadraticcurveto{\pgfqpoint{11.768220in}{1.578708in}}{\pgfqpoint{11.758098in}{1.578708in}}%
\pgfpathlineto{\pgfqpoint{11.758098in}{1.578708in}}%
\pgfpathclose%
\pgfpathmoveto{\pgfqpoint{11.758098in}{1.589016in}}%
\pgfpathquadraticcurveto{\pgfqpoint{11.774281in}{1.589016in}}{\pgfqpoint{11.782817in}{1.576214in}}%
\pgfpathquadraticcurveto{\pgfqpoint{11.791352in}{1.563431in}}{\pgfqpoint{11.791352in}{1.539084in}}%
\pgfpathquadraticcurveto{\pgfqpoint{11.791352in}{1.514798in}}{\pgfqpoint{11.782817in}{1.501995in}}%
\pgfpathquadraticcurveto{\pgfqpoint{11.774281in}{1.489213in}}{\pgfqpoint{11.758098in}{1.489213in}}%
\pgfpathquadraticcurveto{\pgfqpoint{11.741934in}{1.489213in}}{\pgfqpoint{11.733399in}{1.501995in}}%
\pgfpathquadraticcurveto{\pgfqpoint{11.724864in}{1.514798in}}{\pgfqpoint{11.724864in}{1.539084in}}%
\pgfpathquadraticcurveto{\pgfqpoint{11.724864in}{1.563431in}}{\pgfqpoint{11.733399in}{1.576214in}}%
\pgfpathquadraticcurveto{\pgfqpoint{11.741934in}{1.589016in}}{\pgfqpoint{11.758098in}{1.589016in}}%
\pgfpathlineto{\pgfqpoint{11.758098in}{1.589016in}}%
\pgfpathclose%
\pgfpathmoveto{\pgfqpoint{11.825431in}{1.502036in}}%
\pgfpathlineto{\pgfqpoint{11.870848in}{1.502036in}}%
\pgfpathlineto{\pgfqpoint{11.870848in}{1.491089in}}%
\pgfpathlineto{\pgfqpoint{11.809783in}{1.491089in}}%
\pgfpathlineto{\pgfqpoint{11.809783in}{1.502036in}}%
\pgfpathquadraticcurveto{\pgfqpoint{11.817184in}{1.509705in}}{\pgfqpoint{11.829966in}{1.522611in}}%
\pgfpathquadraticcurveto{\pgfqpoint{11.842769in}{1.535538in}}{\pgfqpoint{11.846047in}{1.539290in}}%
\pgfpathquadraticcurveto{\pgfqpoint{11.852294in}{1.546299in}}{\pgfqpoint{11.854768in}{1.551165in}}%
\pgfpathquadraticcurveto{\pgfqpoint{11.857262in}{1.556030in}}{\pgfqpoint{11.857262in}{1.560731in}}%
\pgfpathquadraticcurveto{\pgfqpoint{11.857262in}{1.568400in}}{\pgfqpoint{11.851881in}{1.573224in}}%
\pgfpathquadraticcurveto{\pgfqpoint{11.846500in}{1.578069in}}{\pgfqpoint{11.837862in}{1.578069in}}%
\pgfpathquadraticcurveto{\pgfqpoint{11.831739in}{1.578069in}}{\pgfqpoint{11.824936in}{1.575946in}}%
\pgfpathquadraticcurveto{\pgfqpoint{11.818153in}{1.573822in}}{\pgfqpoint{11.810422in}{1.569493in}}%
\pgfpathlineto{\pgfqpoint{11.810422in}{1.582646in}}%
\pgfpathquadraticcurveto{\pgfqpoint{11.818277in}{1.585800in}}{\pgfqpoint{11.825101in}{1.587408in}}%
\pgfpathquadraticcurveto{\pgfqpoint{11.831945in}{1.589016in}}{\pgfqpoint{11.837615in}{1.589016in}}%
\pgfpathquadraticcurveto{\pgfqpoint{11.852562in}{1.589016in}}{\pgfqpoint{11.861447in}{1.581533in}}%
\pgfpathquadraticcurveto{\pgfqpoint{11.870333in}{1.574069in}}{\pgfqpoint{11.870333in}{1.561576in}}%
\pgfpathquadraticcurveto{\pgfqpoint{11.870333in}{1.555638in}}{\pgfqpoint{11.868106in}{1.550319in}}%
\pgfpathquadraticcurveto{\pgfqpoint{11.865900in}{1.545021in}}{\pgfqpoint{11.860025in}{1.537805in}}%
\pgfpathquadraticcurveto{\pgfqpoint{11.858417in}{1.535929in}}{\pgfqpoint{11.849778in}{1.527002in}}%
\pgfpathquadraticcurveto{\pgfqpoint{11.841161in}{1.518076in}}{\pgfqpoint{11.825431in}{1.502036in}}%
\pgfpathlineto{\pgfqpoint{11.825431in}{1.502036in}}%
\pgfpathclose%
\pgfpathmoveto{\pgfqpoint{11.898309in}{1.587285in}}%
\pgfpathlineto{\pgfqpoint{11.949396in}{1.587285in}}%
\pgfpathlineto{\pgfqpoint{11.949396in}{1.576317in}}%
\pgfpathlineto{\pgfqpoint{11.910225in}{1.576317in}}%
\pgfpathlineto{\pgfqpoint{11.910225in}{1.552752in}}%
\pgfpathquadraticcurveto{\pgfqpoint{11.913050in}{1.553721in}}{\pgfqpoint{11.915874in}{1.554195in}}%
\pgfpathquadraticcurveto{\pgfqpoint{11.918719in}{1.554670in}}{\pgfqpoint{11.921564in}{1.554670in}}%
\pgfpathquadraticcurveto{\pgfqpoint{11.937666in}{1.554670in}}{\pgfqpoint{11.947067in}{1.545846in}}%
\pgfpathquadraticcurveto{\pgfqpoint{11.956488in}{1.537022in}}{\pgfqpoint{11.956488in}{1.521951in}}%
\pgfpathquadraticcurveto{\pgfqpoint{11.956488in}{1.506427in}}{\pgfqpoint{11.946819in}{1.497810in}}%
\pgfpathquadraticcurveto{\pgfqpoint{11.937150in}{1.489213in}}{\pgfqpoint{11.919565in}{1.489213in}}%
\pgfpathquadraticcurveto{\pgfqpoint{11.913503in}{1.489213in}}{\pgfqpoint{11.907215in}{1.490244in}}%
\pgfpathquadraticcurveto{\pgfqpoint{11.900948in}{1.491274in}}{\pgfqpoint{11.894248in}{1.493336in}}%
\pgfpathlineto{\pgfqpoint{11.894248in}{1.506427in}}%
\pgfpathquadraticcurveto{\pgfqpoint{11.900041in}{1.503273in}}{\pgfqpoint{11.906226in}{1.501727in}}%
\pgfpathquadraticcurveto{\pgfqpoint{11.912411in}{1.500181in}}{\pgfqpoint{11.919297in}{1.500181in}}%
\pgfpathquadraticcurveto{\pgfqpoint{11.930450in}{1.500181in}}{\pgfqpoint{11.936944in}{1.506036in}}%
\pgfpathquadraticcurveto{\pgfqpoint{11.943459in}{1.511891in}}{\pgfqpoint{11.943459in}{1.521951in}}%
\pgfpathquadraticcurveto{\pgfqpoint{11.943459in}{1.531992in}}{\pgfqpoint{11.936944in}{1.537847in}}%
\pgfpathquadraticcurveto{\pgfqpoint{11.930450in}{1.543722in}}{\pgfqpoint{11.919297in}{1.543722in}}%
\pgfpathquadraticcurveto{\pgfqpoint{11.914081in}{1.543722in}}{\pgfqpoint{11.908885in}{1.542568in}}%
\pgfpathquadraticcurveto{\pgfqpoint{11.903711in}{1.541413in}}{\pgfqpoint{11.898309in}{1.538960in}}%
\pgfpathlineto{\pgfqpoint{11.898309in}{1.587285in}}%
\pgfpathlineto{\pgfqpoint{11.898309in}{1.587285in}}%
\pgfpathclose%
\pgfpathmoveto{\pgfqpoint{11.978589in}{1.591202in}}%
\pgfpathlineto{\pgfqpoint{11.988897in}{1.591202in}}%
\pgfpathquadraticcurveto{\pgfqpoint{11.998546in}{1.576007in}}{\pgfqpoint{12.003349in}{1.561432in}}%
\pgfpathquadraticcurveto{\pgfqpoint{12.008153in}{1.546877in}}{\pgfqpoint{12.008153in}{1.532507in}}%
\pgfpathquadraticcurveto{\pgfqpoint{12.008153in}{1.518076in}}{\pgfqpoint{12.003349in}{1.503459in}}%
\pgfpathquadraticcurveto{\pgfqpoint{11.998546in}{1.488842in}}{\pgfqpoint{11.988897in}{1.473689in}}%
\pgfpathlineto{\pgfqpoint{11.978589in}{1.473689in}}%
\pgfpathquadraticcurveto{\pgfqpoint{11.987145in}{1.488450in}}{\pgfqpoint{11.991371in}{1.503046in}}%
\pgfpathquadraticcurveto{\pgfqpoint{11.995598in}{1.517643in}}{\pgfqpoint{11.995598in}{1.532507in}}%
\pgfpathquadraticcurveto{\pgfqpoint{11.995598in}{1.547392in}}{\pgfqpoint{11.991371in}{1.561885in}}%
\pgfpathquadraticcurveto{\pgfqpoint{11.987145in}{1.576399in}}{\pgfqpoint{11.978589in}{1.591202in}}%
\pgfpathlineto{\pgfqpoint{11.978589in}{1.591202in}}%
\pgfpathclose%
\pgfusepath{fill}%
\end{pgfscope}%
\begin{pgfscope}%
\pgfsetbuttcap%
\pgfsetroundjoin%
\definecolor{currentfill}{rgb}{0.309804,0.372549,0.419608}%
\pgfsetfillcolor{currentfill}%
\pgfsetlinewidth{0.000000pt}%
\definecolor{currentstroke}{rgb}{0.309804,0.372549,0.419608}%
\pgfsetstrokecolor{currentstroke}%
\pgfsetdash{}{0pt}%
\pgfpathmoveto{\pgfqpoint{0.935312in}{1.207245in}}%
\pgfpathlineto{\pgfqpoint{0.950622in}{1.207245in}}%
\pgfpathlineto{\pgfqpoint{0.987889in}{1.136943in}}%
\pgfpathlineto{\pgfqpoint{0.987889in}{1.207245in}}%
\pgfpathlineto{\pgfqpoint{0.998913in}{1.207245in}}%
\pgfpathlineto{\pgfqpoint{0.998913in}{1.123200in}}%
\pgfpathlineto{\pgfqpoint{0.983602in}{1.123200in}}%
\pgfpathlineto{\pgfqpoint{0.946353in}{1.193501in}}%
\pgfpathlineto{\pgfqpoint{0.946353in}{1.123200in}}%
\pgfpathlineto{\pgfqpoint{0.935312in}{1.123200in}}%
\pgfpathlineto{\pgfqpoint{0.935312in}{1.207245in}}%
\pgfpathlineto{\pgfqpoint{0.935312in}{1.207245in}}%
\pgfpathclose%
\pgfpathmoveto{\pgfqpoint{1.045528in}{1.178984in}}%
\pgfpathquadraticcurveto{\pgfqpoint{1.037206in}{1.178984in}}{\pgfqpoint{1.032361in}{1.172481in}}%
\pgfpathquadraticcurveto{\pgfqpoint{1.027516in}{1.165979in}}{\pgfqpoint{1.027516in}{1.154667in}}%
\pgfpathquadraticcurveto{\pgfqpoint{1.027516in}{1.143356in}}{\pgfqpoint{1.032325in}{1.136853in}}%
\pgfpathquadraticcurveto{\pgfqpoint{1.037152in}{1.130351in}}{\pgfqpoint{1.045528in}{1.130351in}}%
\pgfpathquadraticcurveto{\pgfqpoint{1.053814in}{1.130351in}}{\pgfqpoint{1.058641in}{1.136871in}}%
\pgfpathquadraticcurveto{\pgfqpoint{1.063486in}{1.143410in}}{\pgfqpoint{1.063486in}{1.154667in}}%
\pgfpathquadraticcurveto{\pgfqpoint{1.063486in}{1.165871in}}{\pgfqpoint{1.058641in}{1.172427in}}%
\pgfpathquadraticcurveto{\pgfqpoint{1.053814in}{1.178984in}}{\pgfqpoint{1.045528in}{1.178984in}}%
\pgfpathlineto{\pgfqpoint{1.045528in}{1.178984in}}%
\pgfpathclose%
\pgfpathmoveto{\pgfqpoint{1.045528in}{1.187756in}}%
\pgfpathquadraticcurveto{\pgfqpoint{1.059037in}{1.187756in}}{\pgfqpoint{1.066746in}{1.178966in}}%
\pgfpathquadraticcurveto{\pgfqpoint{1.074474in}{1.170194in}}{\pgfqpoint{1.074474in}{1.154667in}}%
\pgfpathquadraticcurveto{\pgfqpoint{1.074474in}{1.139195in}}{\pgfqpoint{1.066746in}{1.130369in}}%
\pgfpathquadraticcurveto{\pgfqpoint{1.059037in}{1.121561in}}{\pgfqpoint{1.045528in}{1.121561in}}%
\pgfpathquadraticcurveto{\pgfqpoint{1.031965in}{1.121561in}}{\pgfqpoint{1.024274in}{1.130369in}}%
\pgfpathquadraticcurveto{\pgfqpoint{1.016600in}{1.139195in}}{\pgfqpoint{1.016600in}{1.154667in}}%
\pgfpathquadraticcurveto{\pgfqpoint{1.016600in}{1.170194in}}{\pgfqpoint{1.024274in}{1.178966in}}%
\pgfpathquadraticcurveto{\pgfqpoint{1.031965in}{1.187756in}}{\pgfqpoint{1.045528in}{1.187756in}}%
\pgfpathlineto{\pgfqpoint{1.045528in}{1.187756in}}%
\pgfpathclose%
\pgfpathmoveto{\pgfqpoint{1.101888in}{1.204147in}}%
\pgfpathlineto{\pgfqpoint{1.101888in}{1.186243in}}%
\pgfpathlineto{\pgfqpoint{1.123214in}{1.186243in}}%
\pgfpathlineto{\pgfqpoint{1.123214in}{1.178191in}}%
\pgfpathlineto{\pgfqpoint{1.101888in}{1.178191in}}%
\pgfpathlineto{\pgfqpoint{1.101888in}{1.143968in}}%
\pgfpathquadraticcurveto{\pgfqpoint{1.101888in}{1.136259in}}{\pgfqpoint{1.103995in}{1.134061in}}%
\pgfpathquadraticcurveto{\pgfqpoint{1.106103in}{1.131864in}}{\pgfqpoint{1.112587in}{1.131864in}}%
\pgfpathlineto{\pgfqpoint{1.123214in}{1.131864in}}%
\pgfpathlineto{\pgfqpoint{1.123214in}{1.123200in}}%
\pgfpathlineto{\pgfqpoint{1.112587in}{1.123200in}}%
\pgfpathquadraticcurveto{\pgfqpoint{1.100591in}{1.123200in}}{\pgfqpoint{1.096034in}{1.127667in}}%
\pgfpathquadraticcurveto{\pgfqpoint{1.091477in}{1.132152in}}{\pgfqpoint{1.091477in}{1.143968in}}%
\pgfpathlineto{\pgfqpoint{1.091477in}{1.178191in}}%
\pgfpathlineto{\pgfqpoint{1.083876in}{1.178191in}}%
\pgfpathlineto{\pgfqpoint{1.083876in}{1.186243in}}%
\pgfpathlineto{\pgfqpoint{1.091477in}{1.186243in}}%
\pgfpathlineto{\pgfqpoint{1.091477in}{1.204147in}}%
\pgfpathlineto{\pgfqpoint{1.101888in}{1.204147in}}%
\pgfpathlineto{\pgfqpoint{1.101888in}{1.204147in}}%
\pgfpathclose%
\pgfpathmoveto{\pgfqpoint{1.190760in}{1.157315in}}%
\pgfpathlineto{\pgfqpoint{1.190760in}{1.152254in}}%
\pgfpathlineto{\pgfqpoint{1.143136in}{1.152254in}}%
\pgfpathquadraticcurveto{\pgfqpoint{1.143820in}{1.141554in}}{\pgfqpoint{1.149584in}{1.135953in}}%
\pgfpathquadraticcurveto{\pgfqpoint{1.155348in}{1.130351in}}{\pgfqpoint{1.165651in}{1.130351in}}%
\pgfpathquadraticcurveto{\pgfqpoint{1.171613in}{1.130351in}}{\pgfqpoint{1.177215in}{1.131810in}}%
\pgfpathquadraticcurveto{\pgfqpoint{1.182817in}{1.133269in}}{\pgfqpoint{1.188346in}{1.136205in}}%
\pgfpathlineto{\pgfqpoint{1.188346in}{1.126406in}}%
\pgfpathquadraticcurveto{\pgfqpoint{1.182763in}{1.124047in}}{\pgfqpoint{1.176909in}{1.122804in}}%
\pgfpathquadraticcurveto{\pgfqpoint{1.171055in}{1.121561in}}{\pgfqpoint{1.165039in}{1.121561in}}%
\pgfpathquadraticcurveto{\pgfqpoint{1.149944in}{1.121561in}}{\pgfqpoint{1.141137in}{1.130333in}}%
\pgfpathquadraticcurveto{\pgfqpoint{1.132329in}{1.139123in}}{\pgfqpoint{1.132329in}{1.154109in}}%
\pgfpathquadraticcurveto{\pgfqpoint{1.132329in}{1.169581in}}{\pgfqpoint{1.140686in}{1.178659in}}%
\pgfpathquadraticcurveto{\pgfqpoint{1.149044in}{1.187756in}}{\pgfqpoint{1.163237in}{1.187756in}}%
\pgfpathquadraticcurveto{\pgfqpoint{1.175954in}{1.187756in}}{\pgfqpoint{1.183357in}{1.179560in}}%
\pgfpathquadraticcurveto{\pgfqpoint{1.190760in}{1.171383in}}{\pgfqpoint{1.190760in}{1.157315in}}%
\pgfpathlineto{\pgfqpoint{1.190760in}{1.157315in}}%
\pgfpathclose%
\pgfpathmoveto{\pgfqpoint{1.180403in}{1.160359in}}%
\pgfpathquadraticcurveto{\pgfqpoint{1.180295in}{1.168843in}}{\pgfqpoint{1.175648in}{1.173904in}}%
\pgfpathquadraticcurveto{\pgfqpoint{1.171001in}{1.178984in}}{\pgfqpoint{1.163345in}{1.178984in}}%
\pgfpathquadraticcurveto{\pgfqpoint{1.154682in}{1.178984in}}{\pgfqpoint{1.149476in}{1.174084in}}%
\pgfpathquadraticcurveto{\pgfqpoint{1.144271in}{1.169185in}}{\pgfqpoint{1.143478in}{1.160287in}}%
\pgfpathlineto{\pgfqpoint{1.180403in}{1.160359in}}%
\pgfpathlineto{\pgfqpoint{1.180403in}{1.160359in}}%
\pgfpathclose%
\pgfpathmoveto{\pgfqpoint{1.247949in}{1.184387in}}%
\pgfpathlineto{\pgfqpoint{1.247949in}{1.174589in}}%
\pgfpathquadraticcurveto{\pgfqpoint{1.243572in}{1.176840in}}{\pgfqpoint{1.238834in}{1.177957in}}%
\pgfpathquadraticcurveto{\pgfqpoint{1.234115in}{1.179092in}}{\pgfqpoint{1.229036in}{1.179092in}}%
\pgfpathquadraticcurveto{\pgfqpoint{1.221327in}{1.179092in}}{\pgfqpoint{1.217472in}{1.176732in}}%
\pgfpathquadraticcurveto{\pgfqpoint{1.213617in}{1.174373in}}{\pgfqpoint{1.213617in}{1.169635in}}%
\pgfpathquadraticcurveto{\pgfqpoint{1.213617in}{1.166033in}}{\pgfqpoint{1.216373in}{1.163980in}}%
\pgfpathquadraticcurveto{\pgfqpoint{1.219129in}{1.161926in}}{\pgfqpoint{1.227469in}{1.160071in}}%
\pgfpathlineto{\pgfqpoint{1.231017in}{1.159278in}}%
\pgfpathquadraticcurveto{\pgfqpoint{1.242041in}{1.156919in}}{\pgfqpoint{1.246688in}{1.152614in}}%
\pgfpathquadraticcurveto{\pgfqpoint{1.251335in}{1.148309in}}{\pgfqpoint{1.251335in}{1.140600in}}%
\pgfpathquadraticcurveto{\pgfqpoint{1.251335in}{1.131810in}}{\pgfqpoint{1.244382in}{1.126676in}}%
\pgfpathquadraticcurveto{\pgfqpoint{1.237429in}{1.121561in}}{\pgfqpoint{1.225271in}{1.121561in}}%
\pgfpathquadraticcurveto{\pgfqpoint{1.220210in}{1.121561in}}{\pgfqpoint{1.214716in}{1.122552in}}%
\pgfpathquadraticcurveto{\pgfqpoint{1.209222in}{1.123542in}}{\pgfqpoint{1.203152in}{1.125506in}}%
\pgfpathlineto{\pgfqpoint{1.203152in}{1.136205in}}%
\pgfpathquadraticcurveto{\pgfqpoint{1.208898in}{1.133215in}}{\pgfqpoint{1.214464in}{1.131720in}}%
\pgfpathquadraticcurveto{\pgfqpoint{1.220030in}{1.130243in}}{\pgfqpoint{1.225505in}{1.130243in}}%
\pgfpathquadraticcurveto{\pgfqpoint{1.232818in}{1.130243in}}{\pgfqpoint{1.236745in}{1.132746in}}%
\pgfpathquadraticcurveto{\pgfqpoint{1.240690in}{1.135250in}}{\pgfqpoint{1.240690in}{1.139807in}}%
\pgfpathquadraticcurveto{\pgfqpoint{1.240690in}{1.144022in}}{\pgfqpoint{1.237844in}{1.146274in}}%
\pgfpathquadraticcurveto{\pgfqpoint{1.235016in}{1.148525in}}{\pgfqpoint{1.225379in}{1.150614in}}%
\pgfpathlineto{\pgfqpoint{1.221777in}{1.151461in}}%
\pgfpathquadraticcurveto{\pgfqpoint{1.212158in}{1.153478in}}{\pgfqpoint{1.207872in}{1.157675in}}%
\pgfpathquadraticcurveto{\pgfqpoint{1.203603in}{1.161872in}}{\pgfqpoint{1.203603in}{1.169185in}}%
\pgfpathquadraticcurveto{\pgfqpoint{1.203603in}{1.178083in}}{\pgfqpoint{1.209907in}{1.182910in}}%
\pgfpathquadraticcurveto{\pgfqpoint{1.216211in}{1.187756in}}{\pgfqpoint{1.227811in}{1.187756in}}%
\pgfpathquadraticcurveto{\pgfqpoint{1.233539in}{1.187756in}}{\pgfqpoint{1.238600in}{1.186909in}}%
\pgfpathquadraticcurveto{\pgfqpoint{1.243680in}{1.186080in}}{\pgfqpoint{1.247949in}{1.184387in}}%
\pgfpathlineto{\pgfqpoint{1.247949in}{1.184387in}}%
\pgfpathclose%
\pgfpathmoveto{\pgfqpoint{1.270464in}{1.137502in}}%
\pgfpathlineto{\pgfqpoint{1.282334in}{1.137502in}}%
\pgfpathlineto{\pgfqpoint{1.282334in}{1.123200in}}%
\pgfpathlineto{\pgfqpoint{1.270464in}{1.123200in}}%
\pgfpathlineto{\pgfqpoint{1.270464in}{1.137502in}}%
\pgfpathlineto{\pgfqpoint{1.270464in}{1.137502in}}%
\pgfpathclose%
\pgfpathmoveto{\pgfqpoint{1.270464in}{1.182802in}}%
\pgfpathlineto{\pgfqpoint{1.282334in}{1.182802in}}%
\pgfpathlineto{\pgfqpoint{1.282334in}{1.168519in}}%
\pgfpathlineto{\pgfqpoint{1.270464in}{1.168519in}}%
\pgfpathlineto{\pgfqpoint{1.270464in}{1.182802in}}%
\pgfpathlineto{\pgfqpoint{1.270464in}{1.182802in}}%
\pgfpathclose%
\pgfusepath{fill}%
\end{pgfscope}%
\begin{pgfscope}%
\pgfsetbuttcap%
\pgfsetroundjoin%
\definecolor{currentfill}{rgb}{0.309804,0.372549,0.419608}%
\pgfsetfillcolor{currentfill}%
\pgfsetlinewidth{0.000000pt}%
\definecolor{currentstroke}{rgb}{0.309804,0.372549,0.419608}%
\pgfsetstrokecolor{currentstroke}%
\pgfsetdash{}{0pt}%
\pgfpathmoveto{\pgfqpoint{1.437264in}{1.200760in}}%
\pgfpathlineto{\pgfqpoint{1.437264in}{1.188782in}}%
\pgfpathquadraticcurveto{\pgfqpoint{1.431518in}{1.194132in}}{\pgfqpoint{1.425016in}{1.196762in}}%
\pgfpathquadraticcurveto{\pgfqpoint{1.418514in}{1.199409in}}{\pgfqpoint{1.411201in}{1.199409in}}%
\pgfpathquadraticcurveto{\pgfqpoint{1.396791in}{1.199409in}}{\pgfqpoint{1.389136in}{1.190601in}}%
\pgfpathquadraticcurveto{\pgfqpoint{1.381481in}{1.181794in}}{\pgfqpoint{1.381481in}{1.165132in}}%
\pgfpathquadraticcurveto{\pgfqpoint{1.381481in}{1.148525in}}{\pgfqpoint{1.389136in}{1.139717in}}%
\pgfpathquadraticcurveto{\pgfqpoint{1.396791in}{1.130909in}}{\pgfqpoint{1.411201in}{1.130909in}}%
\pgfpathquadraticcurveto{\pgfqpoint{1.418514in}{1.130909in}}{\pgfqpoint{1.425016in}{1.133557in}}%
\pgfpathquadraticcurveto{\pgfqpoint{1.431518in}{1.136205in}}{\pgfqpoint{1.437264in}{1.141554in}}%
\pgfpathlineto{\pgfqpoint{1.437264in}{1.129666in}}%
\pgfpathquadraticcurveto{\pgfqpoint{1.431302in}{1.125614in}}{\pgfqpoint{1.424620in}{1.123578in}}%
\pgfpathquadraticcurveto{\pgfqpoint{1.417955in}{1.121561in}}{\pgfqpoint{1.410534in}{1.121561in}}%
\pgfpathquadraticcurveto{\pgfqpoint{1.391441in}{1.121561in}}{\pgfqpoint{1.380454in}{1.133233in}}%
\pgfpathquadraticcurveto{\pgfqpoint{1.369485in}{1.144923in}}{\pgfqpoint{1.369485in}{1.165132in}}%
\pgfpathquadraticcurveto{\pgfqpoint{1.369485in}{1.185396in}}{\pgfqpoint{1.380454in}{1.197068in}}%
\pgfpathquadraticcurveto{\pgfqpoint{1.391441in}{1.208758in}}{\pgfqpoint{1.410534in}{1.208758in}}%
\pgfpathquadraticcurveto{\pgfqpoint{1.418063in}{1.208758in}}{\pgfqpoint{1.424728in}{1.206758in}}%
\pgfpathquadraticcurveto{\pgfqpoint{1.431410in}{1.204759in}}{\pgfqpoint{1.437264in}{1.200760in}}%
\pgfpathlineto{\pgfqpoint{1.437264in}{1.200760in}}%
\pgfpathclose%
\pgfpathmoveto{\pgfqpoint{1.454376in}{1.186243in}}%
\pgfpathlineto{\pgfqpoint{1.464733in}{1.186243in}}%
\pgfpathlineto{\pgfqpoint{1.464733in}{1.123200in}}%
\pgfpathlineto{\pgfqpoint{1.454376in}{1.123200in}}%
\pgfpathlineto{\pgfqpoint{1.454376in}{1.186243in}}%
\pgfpathlineto{\pgfqpoint{1.454376in}{1.186243in}}%
\pgfpathclose%
\pgfpathmoveto{\pgfqpoint{1.454376in}{1.210793in}}%
\pgfpathlineto{\pgfqpoint{1.464733in}{1.210793in}}%
\pgfpathlineto{\pgfqpoint{1.464733in}{1.197662in}}%
\pgfpathlineto{\pgfqpoint{1.454376in}{1.197662in}}%
\pgfpathlineto{\pgfqpoint{1.454376in}{1.210793in}}%
\pgfpathlineto{\pgfqpoint{1.454376in}{1.210793in}}%
\pgfpathclose%
\pgfpathmoveto{\pgfqpoint{1.522930in}{1.176570in}}%
\pgfpathquadraticcurveto{\pgfqpoint{1.521183in}{1.177579in}}{\pgfqpoint{1.519130in}{1.178047in}}%
\pgfpathquadraticcurveto{\pgfqpoint{1.517076in}{1.178533in}}{\pgfqpoint{1.514609in}{1.178533in}}%
\pgfpathquadraticcurveto{\pgfqpoint{1.505819in}{1.178533in}}{\pgfqpoint{1.501117in}{1.172823in}}%
\pgfpathquadraticcurveto{\pgfqpoint{1.496416in}{1.167114in}}{\pgfqpoint{1.496416in}{1.156414in}}%
\pgfpathlineto{\pgfqpoint{1.496416in}{1.123200in}}%
\pgfpathlineto{\pgfqpoint{1.486005in}{1.123200in}}%
\pgfpathlineto{\pgfqpoint{1.486005in}{1.186243in}}%
\pgfpathlineto{\pgfqpoint{1.496416in}{1.186243in}}%
\pgfpathlineto{\pgfqpoint{1.496416in}{1.176444in}}%
\pgfpathquadraticcurveto{\pgfqpoint{1.499694in}{1.182190in}}{\pgfqpoint{1.504918in}{1.184964in}}%
\pgfpathquadraticcurveto{\pgfqpoint{1.510160in}{1.187756in}}{\pgfqpoint{1.517653in}{1.187756in}}%
\pgfpathquadraticcurveto{\pgfqpoint{1.518715in}{1.187756in}}{\pgfqpoint{1.520012in}{1.187611in}}%
\pgfpathquadraticcurveto{\pgfqpoint{1.521309in}{1.187485in}}{\pgfqpoint{1.522876in}{1.187197in}}%
\pgfpathlineto{\pgfqpoint{1.522930in}{1.176570in}}%
\pgfpathlineto{\pgfqpoint{1.522930in}{1.176570in}}%
\pgfpathclose%
\pgfpathmoveto{\pgfqpoint{1.576624in}{1.183829in}}%
\pgfpathlineto{\pgfqpoint{1.576624in}{1.174138in}}%
\pgfpathquadraticcurveto{\pgfqpoint{1.572229in}{1.176570in}}{\pgfqpoint{1.567816in}{1.177777in}}%
\pgfpathquadraticcurveto{\pgfqpoint{1.563403in}{1.178984in}}{\pgfqpoint{1.558900in}{1.178984in}}%
\pgfpathquadraticcurveto{\pgfqpoint{1.548814in}{1.178984in}}{\pgfqpoint{1.543230in}{1.172589in}}%
\pgfpathquadraticcurveto{\pgfqpoint{1.537664in}{1.166213in}}{\pgfqpoint{1.537664in}{1.154667in}}%
\pgfpathquadraticcurveto{\pgfqpoint{1.537664in}{1.143121in}}{\pgfqpoint{1.543230in}{1.136727in}}%
\pgfpathquadraticcurveto{\pgfqpoint{1.548814in}{1.130351in}}{\pgfqpoint{1.558900in}{1.130351in}}%
\pgfpathquadraticcurveto{\pgfqpoint{1.563403in}{1.130351in}}{\pgfqpoint{1.567816in}{1.131558in}}%
\pgfpathquadraticcurveto{\pgfqpoint{1.572229in}{1.132764in}}{\pgfqpoint{1.576624in}{1.135196in}}%
\pgfpathlineto{\pgfqpoint{1.576624in}{1.125614in}}%
\pgfpathquadraticcurveto{\pgfqpoint{1.572283in}{1.123596in}}{\pgfqpoint{1.567636in}{1.122588in}}%
\pgfpathquadraticcurveto{\pgfqpoint{1.563007in}{1.121561in}}{\pgfqpoint{1.557766in}{1.121561in}}%
\pgfpathquadraticcurveto{\pgfqpoint{1.543518in}{1.121561in}}{\pgfqpoint{1.535124in}{1.130513in}}%
\pgfpathquadraticcurveto{\pgfqpoint{1.526749in}{1.139465in}}{\pgfqpoint{1.526749in}{1.154667in}}%
\pgfpathquadraticcurveto{\pgfqpoint{1.526749in}{1.170086in}}{\pgfqpoint{1.535214in}{1.178912in}}%
\pgfpathquadraticcurveto{\pgfqpoint{1.543698in}{1.187756in}}{\pgfqpoint{1.558450in}{1.187756in}}%
\pgfpathquadraticcurveto{\pgfqpoint{1.563223in}{1.187756in}}{\pgfqpoint{1.567780in}{1.186765in}}%
\pgfpathquadraticcurveto{\pgfqpoint{1.572337in}{1.185792in}}{\pgfqpoint{1.576624in}{1.183829in}}%
\pgfpathlineto{\pgfqpoint{1.576624in}{1.183829in}}%
\pgfpathclose%
\pgfpathmoveto{\pgfqpoint{1.594637in}{1.210793in}}%
\pgfpathlineto{\pgfqpoint{1.604993in}{1.210793in}}%
\pgfpathlineto{\pgfqpoint{1.604993in}{1.123200in}}%
\pgfpathlineto{\pgfqpoint{1.594637in}{1.123200in}}%
\pgfpathlineto{\pgfqpoint{1.594637in}{1.210793in}}%
\pgfpathlineto{\pgfqpoint{1.594637in}{1.210793in}}%
\pgfpathclose%
\pgfpathmoveto{\pgfqpoint{1.680590in}{1.157315in}}%
\pgfpathlineto{\pgfqpoint{1.680590in}{1.152254in}}%
\pgfpathlineto{\pgfqpoint{1.632966in}{1.152254in}}%
\pgfpathquadraticcurveto{\pgfqpoint{1.633651in}{1.141554in}}{\pgfqpoint{1.639415in}{1.135953in}}%
\pgfpathquadraticcurveto{\pgfqpoint{1.645179in}{1.130351in}}{\pgfqpoint{1.655482in}{1.130351in}}%
\pgfpathquadraticcurveto{\pgfqpoint{1.661444in}{1.130351in}}{\pgfqpoint{1.667045in}{1.131810in}}%
\pgfpathquadraticcurveto{\pgfqpoint{1.672647in}{1.133269in}}{\pgfqpoint{1.678177in}{1.136205in}}%
\pgfpathlineto{\pgfqpoint{1.678177in}{1.126406in}}%
\pgfpathquadraticcurveto{\pgfqpoint{1.672593in}{1.124047in}}{\pgfqpoint{1.666739in}{1.122804in}}%
\pgfpathquadraticcurveto{\pgfqpoint{1.660885in}{1.121561in}}{\pgfqpoint{1.654869in}{1.121561in}}%
\pgfpathquadraticcurveto{\pgfqpoint{1.639775in}{1.121561in}}{\pgfqpoint{1.630967in}{1.130333in}}%
\pgfpathquadraticcurveto{\pgfqpoint{1.622159in}{1.139123in}}{\pgfqpoint{1.622159in}{1.154109in}}%
\pgfpathquadraticcurveto{\pgfqpoint{1.622159in}{1.169581in}}{\pgfqpoint{1.630517in}{1.178659in}}%
\pgfpathquadraticcurveto{\pgfqpoint{1.638874in}{1.187756in}}{\pgfqpoint{1.653068in}{1.187756in}}%
\pgfpathquadraticcurveto{\pgfqpoint{1.665785in}{1.187756in}}{\pgfqpoint{1.673188in}{1.179560in}}%
\pgfpathquadraticcurveto{\pgfqpoint{1.680590in}{1.171383in}}{\pgfqpoint{1.680590in}{1.157315in}}%
\pgfpathlineto{\pgfqpoint{1.680590in}{1.157315in}}%
\pgfpathclose%
\pgfpathmoveto{\pgfqpoint{1.670234in}{1.160359in}}%
\pgfpathquadraticcurveto{\pgfqpoint{1.670125in}{1.168843in}}{\pgfqpoint{1.665478in}{1.173904in}}%
\pgfpathquadraticcurveto{\pgfqpoint{1.660831in}{1.178984in}}{\pgfqpoint{1.653176in}{1.178984in}}%
\pgfpathquadraticcurveto{\pgfqpoint{1.644512in}{1.178984in}}{\pgfqpoint{1.639307in}{1.174084in}}%
\pgfpathquadraticcurveto{\pgfqpoint{1.634101in}{1.169185in}}{\pgfqpoint{1.633309in}{1.160287in}}%
\pgfpathlineto{\pgfqpoint{1.670234in}{1.160359in}}%
\pgfpathlineto{\pgfqpoint{1.670234in}{1.160359in}}%
\pgfpathclose%
\pgfpathmoveto{\pgfqpoint{1.737779in}{1.184387in}}%
\pgfpathlineto{\pgfqpoint{1.737779in}{1.174589in}}%
\pgfpathquadraticcurveto{\pgfqpoint{1.733402in}{1.176840in}}{\pgfqpoint{1.728665in}{1.177957in}}%
\pgfpathquadraticcurveto{\pgfqpoint{1.723946in}{1.179092in}}{\pgfqpoint{1.718866in}{1.179092in}}%
\pgfpathquadraticcurveto{\pgfqpoint{1.711157in}{1.179092in}}{\pgfqpoint{1.707303in}{1.176732in}}%
\pgfpathquadraticcurveto{\pgfqpoint{1.703448in}{1.174373in}}{\pgfqpoint{1.703448in}{1.169635in}}%
\pgfpathquadraticcurveto{\pgfqpoint{1.703448in}{1.166033in}}{\pgfqpoint{1.706204in}{1.163980in}}%
\pgfpathquadraticcurveto{\pgfqpoint{1.708960in}{1.161926in}}{\pgfqpoint{1.717299in}{1.160071in}}%
\pgfpathlineto{\pgfqpoint{1.720848in}{1.159278in}}%
\pgfpathquadraticcurveto{\pgfqpoint{1.731871in}{1.156919in}}{\pgfqpoint{1.736518in}{1.152614in}}%
\pgfpathquadraticcurveto{\pgfqpoint{1.741165in}{1.148309in}}{\pgfqpoint{1.741165in}{1.140600in}}%
\pgfpathquadraticcurveto{\pgfqpoint{1.741165in}{1.131810in}}{\pgfqpoint{1.734213in}{1.126676in}}%
\pgfpathquadraticcurveto{\pgfqpoint{1.727260in}{1.121561in}}{\pgfqpoint{1.715102in}{1.121561in}}%
\pgfpathquadraticcurveto{\pgfqpoint{1.710040in}{1.121561in}}{\pgfqpoint{1.704547in}{1.122552in}}%
\pgfpathquadraticcurveto{\pgfqpoint{1.699053in}{1.123542in}}{\pgfqpoint{1.692983in}{1.125506in}}%
\pgfpathlineto{\pgfqpoint{1.692983in}{1.136205in}}%
\pgfpathquadraticcurveto{\pgfqpoint{1.698729in}{1.133215in}}{\pgfqpoint{1.704294in}{1.131720in}}%
\pgfpathquadraticcurveto{\pgfqpoint{1.709860in}{1.130243in}}{\pgfqpoint{1.715336in}{1.130243in}}%
\pgfpathquadraticcurveto{\pgfqpoint{1.722649in}{1.130243in}}{\pgfqpoint{1.726576in}{1.132746in}}%
\pgfpathquadraticcurveto{\pgfqpoint{1.730520in}{1.135250in}}{\pgfqpoint{1.730520in}{1.139807in}}%
\pgfpathquadraticcurveto{\pgfqpoint{1.730520in}{1.144022in}}{\pgfqpoint{1.727674in}{1.146274in}}%
\pgfpathquadraticcurveto{\pgfqpoint{1.724846in}{1.148525in}}{\pgfqpoint{1.715210in}{1.150614in}}%
\pgfpathlineto{\pgfqpoint{1.711607in}{1.151461in}}%
\pgfpathquadraticcurveto{\pgfqpoint{1.701989in}{1.153478in}}{\pgfqpoint{1.697702in}{1.157675in}}%
\pgfpathquadraticcurveto{\pgfqpoint{1.693433in}{1.161872in}}{\pgfqpoint{1.693433in}{1.169185in}}%
\pgfpathquadraticcurveto{\pgfqpoint{1.693433in}{1.178083in}}{\pgfqpoint{1.699737in}{1.182910in}}%
\pgfpathquadraticcurveto{\pgfqpoint{1.706042in}{1.187756in}}{\pgfqpoint{1.717641in}{1.187756in}}%
\pgfpathquadraticcurveto{\pgfqpoint{1.723369in}{1.187756in}}{\pgfqpoint{1.728431in}{1.186909in}}%
\pgfpathquadraticcurveto{\pgfqpoint{1.733510in}{1.186080in}}{\pgfqpoint{1.737779in}{1.184387in}}%
\pgfpathlineto{\pgfqpoint{1.737779in}{1.184387in}}%
\pgfpathclose%
\pgfpathmoveto{\pgfqpoint{1.760294in}{1.137502in}}%
\pgfpathlineto{\pgfqpoint{1.772164in}{1.137502in}}%
\pgfpathlineto{\pgfqpoint{1.772164in}{1.123200in}}%
\pgfpathlineto{\pgfqpoint{1.760294in}{1.123200in}}%
\pgfpathlineto{\pgfqpoint{1.760294in}{1.137502in}}%
\pgfpathlineto{\pgfqpoint{1.760294in}{1.137502in}}%
\pgfpathclose%
\pgfpathmoveto{\pgfqpoint{1.760294in}{1.182802in}}%
\pgfpathlineto{\pgfqpoint{1.772164in}{1.182802in}}%
\pgfpathlineto{\pgfqpoint{1.772164in}{1.168519in}}%
\pgfpathlineto{\pgfqpoint{1.760294in}{1.168519in}}%
\pgfpathlineto{\pgfqpoint{1.760294in}{1.182802in}}%
\pgfpathlineto{\pgfqpoint{1.760294in}{1.182802in}}%
\pgfpathclose%
\pgfusepath{fill}%
\end{pgfscope}%
\begin{pgfscope}%
\pgfsetbuttcap%
\pgfsetroundjoin%
\definecolor{currentfill}{rgb}{0.309804,0.372549,0.419608}%
\pgfsetfillcolor{currentfill}%
\pgfsetlinewidth{0.000000pt}%
\definecolor{currentstroke}{rgb}{0.309804,0.372549,0.419608}%
\pgfsetstrokecolor{currentstroke}%
\pgfsetdash{}{0pt}%
\pgfpathmoveto{\pgfqpoint{1.900583in}{1.184387in}}%
\pgfpathlineto{\pgfqpoint{1.900583in}{1.174589in}}%
\pgfpathquadraticcurveto{\pgfqpoint{1.896206in}{1.176840in}}{\pgfqpoint{1.891469in}{1.177957in}}%
\pgfpathquadraticcurveto{\pgfqpoint{1.886750in}{1.179092in}}{\pgfqpoint{1.881670in}{1.179092in}}%
\pgfpathquadraticcurveto{\pgfqpoint{1.873961in}{1.179092in}}{\pgfqpoint{1.870106in}{1.176732in}}%
\pgfpathquadraticcurveto{\pgfqpoint{1.866252in}{1.174373in}}{\pgfqpoint{1.866252in}{1.169635in}}%
\pgfpathquadraticcurveto{\pgfqpoint{1.866252in}{1.166033in}}{\pgfqpoint{1.869008in}{1.163980in}}%
\pgfpathquadraticcurveto{\pgfqpoint{1.871763in}{1.161926in}}{\pgfqpoint{1.880103in}{1.160071in}}%
\pgfpathlineto{\pgfqpoint{1.883651in}{1.159278in}}%
\pgfpathquadraticcurveto{\pgfqpoint{1.894675in}{1.156919in}}{\pgfqpoint{1.899322in}{1.152614in}}%
\pgfpathquadraticcurveto{\pgfqpoint{1.903969in}{1.148309in}}{\pgfqpoint{1.903969in}{1.140600in}}%
\pgfpathquadraticcurveto{\pgfqpoint{1.903969in}{1.131810in}}{\pgfqpoint{1.897016in}{1.126676in}}%
\pgfpathquadraticcurveto{\pgfqpoint{1.890064in}{1.121561in}}{\pgfqpoint{1.877906in}{1.121561in}}%
\pgfpathquadraticcurveto{\pgfqpoint{1.872844in}{1.121561in}}{\pgfqpoint{1.867350in}{1.122552in}}%
\pgfpathquadraticcurveto{\pgfqpoint{1.861857in}{1.123542in}}{\pgfqpoint{1.855787in}{1.125506in}}%
\pgfpathlineto{\pgfqpoint{1.855787in}{1.136205in}}%
\pgfpathquadraticcurveto{\pgfqpoint{1.861533in}{1.133215in}}{\pgfqpoint{1.867098in}{1.131720in}}%
\pgfpathquadraticcurveto{\pgfqpoint{1.872664in}{1.130243in}}{\pgfqpoint{1.878140in}{1.130243in}}%
\pgfpathquadraticcurveto{\pgfqpoint{1.885453in}{1.130243in}}{\pgfqpoint{1.889379in}{1.132746in}}%
\pgfpathquadraticcurveto{\pgfqpoint{1.893324in}{1.135250in}}{\pgfqpoint{1.893324in}{1.139807in}}%
\pgfpathquadraticcurveto{\pgfqpoint{1.893324in}{1.144022in}}{\pgfqpoint{1.890478in}{1.146274in}}%
\pgfpathquadraticcurveto{\pgfqpoint{1.887650in}{1.148525in}}{\pgfqpoint{1.878014in}{1.150614in}}%
\pgfpathlineto{\pgfqpoint{1.874411in}{1.151461in}}%
\pgfpathquadraticcurveto{\pgfqpoint{1.864793in}{1.153478in}}{\pgfqpoint{1.860506in}{1.157675in}}%
\pgfpathquadraticcurveto{\pgfqpoint{1.856237in}{1.161872in}}{\pgfqpoint{1.856237in}{1.169185in}}%
\pgfpathquadraticcurveto{\pgfqpoint{1.856237in}{1.178083in}}{\pgfqpoint{1.862541in}{1.182910in}}%
\pgfpathquadraticcurveto{\pgfqpoint{1.868845in}{1.187756in}}{\pgfqpoint{1.880445in}{1.187756in}}%
\pgfpathquadraticcurveto{\pgfqpoint{1.886173in}{1.187756in}}{\pgfqpoint{1.891235in}{1.186909in}}%
\pgfpathquadraticcurveto{\pgfqpoint{1.896314in}{1.186080in}}{\pgfqpoint{1.900583in}{1.184387in}}%
\pgfpathlineto{\pgfqpoint{1.900583in}{1.184387in}}%
\pgfpathclose%
\pgfpathmoveto{\pgfqpoint{1.930699in}{1.204147in}}%
\pgfpathlineto{\pgfqpoint{1.930699in}{1.186243in}}%
\pgfpathlineto{\pgfqpoint{1.952026in}{1.186243in}}%
\pgfpathlineto{\pgfqpoint{1.952026in}{1.178191in}}%
\pgfpathlineto{\pgfqpoint{1.930699in}{1.178191in}}%
\pgfpathlineto{\pgfqpoint{1.930699in}{1.143968in}}%
\pgfpathquadraticcurveto{\pgfqpoint{1.930699in}{1.136259in}}{\pgfqpoint{1.932807in}{1.134061in}}%
\pgfpathquadraticcurveto{\pgfqpoint{1.934914in}{1.131864in}}{\pgfqpoint{1.941398in}{1.131864in}}%
\pgfpathlineto{\pgfqpoint{1.952026in}{1.131864in}}%
\pgfpathlineto{\pgfqpoint{1.952026in}{1.123200in}}%
\pgfpathlineto{\pgfqpoint{1.941398in}{1.123200in}}%
\pgfpathquadraticcurveto{\pgfqpoint{1.929402in}{1.123200in}}{\pgfqpoint{1.924845in}{1.127667in}}%
\pgfpathquadraticcurveto{\pgfqpoint{1.920288in}{1.132152in}}{\pgfqpoint{1.920288in}{1.143968in}}%
\pgfpathlineto{\pgfqpoint{1.920288in}{1.178191in}}%
\pgfpathlineto{\pgfqpoint{1.912687in}{1.178191in}}%
\pgfpathlineto{\pgfqpoint{1.912687in}{1.186243in}}%
\pgfpathlineto{\pgfqpoint{1.920288in}{1.186243in}}%
\pgfpathlineto{\pgfqpoint{1.920288in}{1.204147in}}%
\pgfpathlineto{\pgfqpoint{1.930699in}{1.204147in}}%
\pgfpathlineto{\pgfqpoint{1.930699in}{1.204147in}}%
\pgfpathclose%
\pgfpathmoveto{\pgfqpoint{1.994300in}{1.154883in}}%
\pgfpathquadraticcurveto{\pgfqpoint{1.981746in}{1.154883in}}{\pgfqpoint{1.976900in}{1.152019in}}%
\pgfpathquadraticcurveto{\pgfqpoint{1.972055in}{1.149156in}}{\pgfqpoint{1.972055in}{1.142221in}}%
\pgfpathquadraticcurveto{\pgfqpoint{1.972055in}{1.136709in}}{\pgfqpoint{1.975694in}{1.133467in}}%
\pgfpathquadraticcurveto{\pgfqpoint{1.979332in}{1.130243in}}{\pgfqpoint{1.985564in}{1.130243in}}%
\pgfpathquadraticcurveto{\pgfqpoint{1.994192in}{1.130243in}}{\pgfqpoint{1.999398in}{1.136349in}}%
\pgfpathquadraticcurveto{\pgfqpoint{2.004603in}{1.142455in}}{\pgfqpoint{2.004603in}{1.152578in}}%
\pgfpathlineto{\pgfqpoint{2.004603in}{1.154883in}}%
\pgfpathlineto{\pgfqpoint{1.994300in}{1.154883in}}%
\pgfpathlineto{\pgfqpoint{1.994300in}{1.154883in}}%
\pgfpathclose%
\pgfpathmoveto{\pgfqpoint{2.014960in}{1.159170in}}%
\pgfpathlineto{\pgfqpoint{2.014960in}{1.123200in}}%
\pgfpathlineto{\pgfqpoint{2.004603in}{1.123200in}}%
\pgfpathlineto{\pgfqpoint{2.004603in}{1.132764in}}%
\pgfpathquadraticcurveto{\pgfqpoint{2.001055in}{1.127037in}}{\pgfqpoint{1.995759in}{1.124299in}}%
\pgfpathquadraticcurveto{\pgfqpoint{1.990464in}{1.121561in}}{\pgfqpoint{1.982808in}{1.121561in}}%
\pgfpathquadraticcurveto{\pgfqpoint{1.973136in}{1.121561in}}{\pgfqpoint{1.967408in}{1.127001in}}%
\pgfpathquadraticcurveto{\pgfqpoint{1.961698in}{1.132440in}}{\pgfqpoint{1.961698in}{1.141554in}}%
\pgfpathquadraticcurveto{\pgfqpoint{1.961698in}{1.152182in}}{\pgfqpoint{1.968813in}{1.157585in}}%
\pgfpathquadraticcurveto{\pgfqpoint{1.975946in}{1.162989in}}{\pgfqpoint{1.990067in}{1.162989in}}%
\pgfpathlineto{\pgfqpoint{2.004603in}{1.162989in}}%
\pgfpathlineto{\pgfqpoint{2.004603in}{1.164016in}}%
\pgfpathquadraticcurveto{\pgfqpoint{2.004603in}{1.171166in}}{\pgfqpoint{1.999902in}{1.175075in}}%
\pgfpathquadraticcurveto{\pgfqpoint{1.995201in}{1.178984in}}{\pgfqpoint{1.986699in}{1.178984in}}%
\pgfpathquadraticcurveto{\pgfqpoint{1.981295in}{1.178984in}}{\pgfqpoint{1.976162in}{1.177687in}}%
\pgfpathquadraticcurveto{\pgfqpoint{1.971046in}{1.176390in}}{\pgfqpoint{1.966327in}{1.173796in}}%
\pgfpathlineto{\pgfqpoint{1.966327in}{1.183379in}}%
\pgfpathquadraticcurveto{\pgfqpoint{1.972001in}{1.185576in}}{\pgfqpoint{1.977351in}{1.186657in}}%
\pgfpathquadraticcurveto{\pgfqpoint{1.982700in}{1.187756in}}{\pgfqpoint{1.987762in}{1.187756in}}%
\pgfpathquadraticcurveto{\pgfqpoint{2.001451in}{1.187756in}}{\pgfqpoint{2.008206in}{1.180659in}}%
\pgfpathquadraticcurveto{\pgfqpoint{2.014960in}{1.173580in}}{\pgfqpoint{2.014960in}{1.159170in}}%
\pgfpathlineto{\pgfqpoint{2.014960in}{1.159170in}}%
\pgfpathclose%
\pgfpathmoveto{\pgfqpoint{2.046535in}{1.204147in}}%
\pgfpathlineto{\pgfqpoint{2.046535in}{1.186243in}}%
\pgfpathlineto{\pgfqpoint{2.067862in}{1.186243in}}%
\pgfpathlineto{\pgfqpoint{2.067862in}{1.178191in}}%
\pgfpathlineto{\pgfqpoint{2.046535in}{1.178191in}}%
\pgfpathlineto{\pgfqpoint{2.046535in}{1.143968in}}%
\pgfpathquadraticcurveto{\pgfqpoint{2.046535in}{1.136259in}}{\pgfqpoint{2.048643in}{1.134061in}}%
\pgfpathquadraticcurveto{\pgfqpoint{2.050750in}{1.131864in}}{\pgfqpoint{2.057235in}{1.131864in}}%
\pgfpathlineto{\pgfqpoint{2.067862in}{1.131864in}}%
\pgfpathlineto{\pgfqpoint{2.067862in}{1.123200in}}%
\pgfpathlineto{\pgfqpoint{2.057235in}{1.123200in}}%
\pgfpathquadraticcurveto{\pgfqpoint{2.045238in}{1.123200in}}{\pgfqpoint{2.040681in}{1.127667in}}%
\pgfpathquadraticcurveto{\pgfqpoint{2.036124in}{1.132152in}}{\pgfqpoint{2.036124in}{1.143968in}}%
\pgfpathlineto{\pgfqpoint{2.036124in}{1.178191in}}%
\pgfpathlineto{\pgfqpoint{2.028523in}{1.178191in}}%
\pgfpathlineto{\pgfqpoint{2.028523in}{1.186243in}}%
\pgfpathlineto{\pgfqpoint{2.036124in}{1.186243in}}%
\pgfpathlineto{\pgfqpoint{2.036124in}{1.204147in}}%
\pgfpathlineto{\pgfqpoint{2.046535in}{1.204147in}}%
\pgfpathlineto{\pgfqpoint{2.046535in}{1.204147in}}%
\pgfpathclose%
\pgfpathmoveto{\pgfqpoint{2.081479in}{1.186243in}}%
\pgfpathlineto{\pgfqpoint{2.091836in}{1.186243in}}%
\pgfpathlineto{\pgfqpoint{2.091836in}{1.123200in}}%
\pgfpathlineto{\pgfqpoint{2.081479in}{1.123200in}}%
\pgfpathlineto{\pgfqpoint{2.081479in}{1.186243in}}%
\pgfpathlineto{\pgfqpoint{2.081479in}{1.186243in}}%
\pgfpathclose%
\pgfpathmoveto{\pgfqpoint{2.081479in}{1.210793in}}%
\pgfpathlineto{\pgfqpoint{2.091836in}{1.210793in}}%
\pgfpathlineto{\pgfqpoint{2.091836in}{1.197662in}}%
\pgfpathlineto{\pgfqpoint{2.081479in}{1.197662in}}%
\pgfpathlineto{\pgfqpoint{2.081479in}{1.210793in}}%
\pgfpathlineto{\pgfqpoint{2.081479in}{1.210793in}}%
\pgfpathclose%
\pgfpathmoveto{\pgfqpoint{2.158877in}{1.183829in}}%
\pgfpathlineto{\pgfqpoint{2.158877in}{1.174138in}}%
\pgfpathquadraticcurveto{\pgfqpoint{2.154482in}{1.176570in}}{\pgfqpoint{2.150069in}{1.177777in}}%
\pgfpathquadraticcurveto{\pgfqpoint{2.145656in}{1.178984in}}{\pgfqpoint{2.141153in}{1.178984in}}%
\pgfpathquadraticcurveto{\pgfqpoint{2.131066in}{1.178984in}}{\pgfqpoint{2.125483in}{1.172589in}}%
\pgfpathquadraticcurveto{\pgfqpoint{2.119917in}{1.166213in}}{\pgfqpoint{2.119917in}{1.154667in}}%
\pgfpathquadraticcurveto{\pgfqpoint{2.119917in}{1.143121in}}{\pgfqpoint{2.125483in}{1.136727in}}%
\pgfpathquadraticcurveto{\pgfqpoint{2.131066in}{1.130351in}}{\pgfqpoint{2.141153in}{1.130351in}}%
\pgfpathquadraticcurveto{\pgfqpoint{2.145656in}{1.130351in}}{\pgfqpoint{2.150069in}{1.131558in}}%
\pgfpathquadraticcurveto{\pgfqpoint{2.154482in}{1.132764in}}{\pgfqpoint{2.158877in}{1.135196in}}%
\pgfpathlineto{\pgfqpoint{2.158877in}{1.125614in}}%
\pgfpathquadraticcurveto{\pgfqpoint{2.154536in}{1.123596in}}{\pgfqpoint{2.149889in}{1.122588in}}%
\pgfpathquadraticcurveto{\pgfqpoint{2.145260in}{1.121561in}}{\pgfqpoint{2.140018in}{1.121561in}}%
\pgfpathquadraticcurveto{\pgfqpoint{2.125771in}{1.121561in}}{\pgfqpoint{2.117377in}{1.130513in}}%
\pgfpathquadraticcurveto{\pgfqpoint{2.109002in}{1.139465in}}{\pgfqpoint{2.109002in}{1.154667in}}%
\pgfpathquadraticcurveto{\pgfqpoint{2.109002in}{1.170086in}}{\pgfqpoint{2.117467in}{1.178912in}}%
\pgfpathquadraticcurveto{\pgfqpoint{2.125951in}{1.187756in}}{\pgfqpoint{2.140703in}{1.187756in}}%
\pgfpathquadraticcurveto{\pgfqpoint{2.145476in}{1.187756in}}{\pgfqpoint{2.150033in}{1.186765in}}%
\pgfpathquadraticcurveto{\pgfqpoint{2.154590in}{1.185792in}}{\pgfqpoint{2.158877in}{1.183829in}}%
\pgfpathlineto{\pgfqpoint{2.158877in}{1.183829in}}%
\pgfpathclose%
\pgfusepath{fill}%
\end{pgfscope}%
\begin{pgfscope}%
\pgfsetbuttcap%
\pgfsetroundjoin%
\definecolor{currentfill}{rgb}{0.309804,0.372549,0.419608}%
\pgfsetfillcolor{currentfill}%
\pgfsetlinewidth{0.000000pt}%
\definecolor{currentstroke}{rgb}{0.309804,0.372549,0.419608}%
\pgfsetstrokecolor{currentstroke}%
\pgfsetdash{}{0pt}%
\pgfpathmoveto{\pgfqpoint{2.289601in}{1.184387in}}%
\pgfpathlineto{\pgfqpoint{2.289601in}{1.174589in}}%
\pgfpathquadraticcurveto{\pgfqpoint{2.285224in}{1.176840in}}{\pgfqpoint{2.280487in}{1.177957in}}%
\pgfpathquadraticcurveto{\pgfqpoint{2.275768in}{1.179092in}}{\pgfqpoint{2.270688in}{1.179092in}}%
\pgfpathquadraticcurveto{\pgfqpoint{2.262979in}{1.179092in}}{\pgfqpoint{2.259125in}{1.176732in}}%
\pgfpathquadraticcurveto{\pgfqpoint{2.255270in}{1.174373in}}{\pgfqpoint{2.255270in}{1.169635in}}%
\pgfpathquadraticcurveto{\pgfqpoint{2.255270in}{1.166033in}}{\pgfqpoint{2.258026in}{1.163980in}}%
\pgfpathquadraticcurveto{\pgfqpoint{2.260782in}{1.161926in}}{\pgfqpoint{2.269121in}{1.160071in}}%
\pgfpathlineto{\pgfqpoint{2.272670in}{1.159278in}}%
\pgfpathquadraticcurveto{\pgfqpoint{2.283693in}{1.156919in}}{\pgfqpoint{2.288340in}{1.152614in}}%
\pgfpathquadraticcurveto{\pgfqpoint{2.292987in}{1.148309in}}{\pgfqpoint{2.292987in}{1.140600in}}%
\pgfpathquadraticcurveto{\pgfqpoint{2.292987in}{1.131810in}}{\pgfqpoint{2.286035in}{1.126676in}}%
\pgfpathquadraticcurveto{\pgfqpoint{2.279082in}{1.121561in}}{\pgfqpoint{2.266924in}{1.121561in}}%
\pgfpathquadraticcurveto{\pgfqpoint{2.261862in}{1.121561in}}{\pgfqpoint{2.256369in}{1.122552in}}%
\pgfpathquadraticcurveto{\pgfqpoint{2.250875in}{1.123542in}}{\pgfqpoint{2.244805in}{1.125506in}}%
\pgfpathlineto{\pgfqpoint{2.244805in}{1.136205in}}%
\pgfpathquadraticcurveto{\pgfqpoint{2.250551in}{1.133215in}}{\pgfqpoint{2.256117in}{1.131720in}}%
\pgfpathquadraticcurveto{\pgfqpoint{2.261682in}{1.130243in}}{\pgfqpoint{2.267158in}{1.130243in}}%
\pgfpathquadraticcurveto{\pgfqpoint{2.274471in}{1.130243in}}{\pgfqpoint{2.278398in}{1.132746in}}%
\pgfpathquadraticcurveto{\pgfqpoint{2.282342in}{1.135250in}}{\pgfqpoint{2.282342in}{1.139807in}}%
\pgfpathquadraticcurveto{\pgfqpoint{2.282342in}{1.144022in}}{\pgfqpoint{2.279496in}{1.146274in}}%
\pgfpathquadraticcurveto{\pgfqpoint{2.276668in}{1.148525in}}{\pgfqpoint{2.267032in}{1.150614in}}%
\pgfpathlineto{\pgfqpoint{2.263429in}{1.151461in}}%
\pgfpathquadraticcurveto{\pgfqpoint{2.253811in}{1.153478in}}{\pgfqpoint{2.249524in}{1.157675in}}%
\pgfpathquadraticcurveto{\pgfqpoint{2.245255in}{1.161872in}}{\pgfqpoint{2.245255in}{1.169185in}}%
\pgfpathquadraticcurveto{\pgfqpoint{2.245255in}{1.178083in}}{\pgfqpoint{2.251559in}{1.182910in}}%
\pgfpathquadraticcurveto{\pgfqpoint{2.257864in}{1.187756in}}{\pgfqpoint{2.269464in}{1.187756in}}%
\pgfpathquadraticcurveto{\pgfqpoint{2.275191in}{1.187756in}}{\pgfqpoint{2.280253in}{1.186909in}}%
\pgfpathquadraticcurveto{\pgfqpoint{2.285332in}{1.186080in}}{\pgfqpoint{2.289601in}{1.184387in}}%
\pgfpathlineto{\pgfqpoint{2.289601in}{1.184387in}}%
\pgfpathclose%
\pgfpathmoveto{\pgfqpoint{2.319483in}{1.132656in}}%
\pgfpathlineto{\pgfqpoint{2.319483in}{1.099226in}}%
\pgfpathlineto{\pgfqpoint{2.309072in}{1.099226in}}%
\pgfpathlineto{\pgfqpoint{2.309072in}{1.186243in}}%
\pgfpathlineto{\pgfqpoint{2.319483in}{1.186243in}}%
\pgfpathlineto{\pgfqpoint{2.319483in}{1.176678in}}%
\pgfpathquadraticcurveto{\pgfqpoint{2.322762in}{1.182298in}}{\pgfqpoint{2.327733in}{1.185018in}}%
\pgfpathquadraticcurveto{\pgfqpoint{2.332722in}{1.187756in}}{\pgfqpoint{2.339639in}{1.187756in}}%
\pgfpathquadraticcurveto{\pgfqpoint{2.351131in}{1.187756in}}{\pgfqpoint{2.358299in}{1.178641in}}%
\pgfpathquadraticcurveto{\pgfqpoint{2.365486in}{1.169527in}}{\pgfqpoint{2.365486in}{1.154667in}}%
\pgfpathquadraticcurveto{\pgfqpoint{2.365486in}{1.139807in}}{\pgfqpoint{2.358299in}{1.130675in}}%
\pgfpathquadraticcurveto{\pgfqpoint{2.351131in}{1.121561in}}{\pgfqpoint{2.339639in}{1.121561in}}%
\pgfpathquadraticcurveto{\pgfqpoint{2.332722in}{1.121561in}}{\pgfqpoint{2.327733in}{1.124299in}}%
\pgfpathquadraticcurveto{\pgfqpoint{2.322762in}{1.127037in}}{\pgfqpoint{2.319483in}{1.132656in}}%
\pgfpathlineto{\pgfqpoint{2.319483in}{1.132656in}}%
\pgfpathclose%
\pgfpathmoveto{\pgfqpoint{2.354733in}{1.154667in}}%
\pgfpathquadraticcurveto{\pgfqpoint{2.354733in}{1.166087in}}{\pgfqpoint{2.350032in}{1.172589in}}%
\pgfpathquadraticcurveto{\pgfqpoint{2.345331in}{1.179092in}}{\pgfqpoint{2.337117in}{1.179092in}}%
\pgfpathquadraticcurveto{\pgfqpoint{2.328886in}{1.179092in}}{\pgfqpoint{2.324184in}{1.172589in}}%
\pgfpathquadraticcurveto{\pgfqpoint{2.319483in}{1.166087in}}{\pgfqpoint{2.319483in}{1.154667in}}%
\pgfpathquadraticcurveto{\pgfqpoint{2.319483in}{1.143248in}}{\pgfqpoint{2.324184in}{1.136745in}}%
\pgfpathquadraticcurveto{\pgfqpoint{2.328886in}{1.130243in}}{\pgfqpoint{2.337117in}{1.130243in}}%
\pgfpathquadraticcurveto{\pgfqpoint{2.345331in}{1.130243in}}{\pgfqpoint{2.350032in}{1.136745in}}%
\pgfpathquadraticcurveto{\pgfqpoint{2.354733in}{1.143248in}}{\pgfqpoint{2.354733in}{1.154667in}}%
\pgfpathlineto{\pgfqpoint{2.354733in}{1.154667in}}%
\pgfpathclose%
\pgfpathmoveto{\pgfqpoint{2.436580in}{1.157315in}}%
\pgfpathlineto{\pgfqpoint{2.436580in}{1.152254in}}%
\pgfpathlineto{\pgfqpoint{2.388956in}{1.152254in}}%
\pgfpathquadraticcurveto{\pgfqpoint{2.389641in}{1.141554in}}{\pgfqpoint{2.395405in}{1.135953in}}%
\pgfpathquadraticcurveto{\pgfqpoint{2.401168in}{1.130351in}}{\pgfqpoint{2.411471in}{1.130351in}}%
\pgfpathquadraticcurveto{\pgfqpoint{2.417433in}{1.130351in}}{\pgfqpoint{2.423035in}{1.131810in}}%
\pgfpathquadraticcurveto{\pgfqpoint{2.428637in}{1.133269in}}{\pgfqpoint{2.434167in}{1.136205in}}%
\pgfpathlineto{\pgfqpoint{2.434167in}{1.126406in}}%
\pgfpathquadraticcurveto{\pgfqpoint{2.428583in}{1.124047in}}{\pgfqpoint{2.422729in}{1.122804in}}%
\pgfpathquadraticcurveto{\pgfqpoint{2.416875in}{1.121561in}}{\pgfqpoint{2.410859in}{1.121561in}}%
\pgfpathquadraticcurveto{\pgfqpoint{2.395765in}{1.121561in}}{\pgfqpoint{2.386957in}{1.130333in}}%
\pgfpathquadraticcurveto{\pgfqpoint{2.378149in}{1.139123in}}{\pgfqpoint{2.378149in}{1.154109in}}%
\pgfpathquadraticcurveto{\pgfqpoint{2.378149in}{1.169581in}}{\pgfqpoint{2.386507in}{1.178659in}}%
\pgfpathquadraticcurveto{\pgfqpoint{2.394864in}{1.187756in}}{\pgfqpoint{2.409058in}{1.187756in}}%
\pgfpathquadraticcurveto{\pgfqpoint{2.421774in}{1.187756in}}{\pgfqpoint{2.429177in}{1.179560in}}%
\pgfpathquadraticcurveto{\pgfqpoint{2.436580in}{1.171383in}}{\pgfqpoint{2.436580in}{1.157315in}}%
\pgfpathlineto{\pgfqpoint{2.436580in}{1.157315in}}%
\pgfpathclose%
\pgfpathmoveto{\pgfqpoint{2.426223in}{1.160359in}}%
\pgfpathquadraticcurveto{\pgfqpoint{2.426115in}{1.168843in}}{\pgfqpoint{2.421468in}{1.173904in}}%
\pgfpathquadraticcurveto{\pgfqpoint{2.416821in}{1.178984in}}{\pgfqpoint{2.409166in}{1.178984in}}%
\pgfpathquadraticcurveto{\pgfqpoint{2.400502in}{1.178984in}}{\pgfqpoint{2.395296in}{1.174084in}}%
\pgfpathquadraticcurveto{\pgfqpoint{2.390091in}{1.169185in}}{\pgfqpoint{2.389298in}{1.160287in}}%
\pgfpathlineto{\pgfqpoint{2.426223in}{1.160359in}}%
\pgfpathlineto{\pgfqpoint{2.426223in}{1.160359in}}%
\pgfpathclose%
\pgfpathmoveto{\pgfqpoint{2.498956in}{1.183829in}}%
\pgfpathlineto{\pgfqpoint{2.498956in}{1.174138in}}%
\pgfpathquadraticcurveto{\pgfqpoint{2.494561in}{1.176570in}}{\pgfqpoint{2.490148in}{1.177777in}}%
\pgfpathquadraticcurveto{\pgfqpoint{2.485735in}{1.178984in}}{\pgfqpoint{2.481232in}{1.178984in}}%
\pgfpathquadraticcurveto{\pgfqpoint{2.471146in}{1.178984in}}{\pgfqpoint{2.465562in}{1.172589in}}%
\pgfpathquadraticcurveto{\pgfqpoint{2.459996in}{1.166213in}}{\pgfqpoint{2.459996in}{1.154667in}}%
\pgfpathquadraticcurveto{\pgfqpoint{2.459996in}{1.143121in}}{\pgfqpoint{2.465562in}{1.136727in}}%
\pgfpathquadraticcurveto{\pgfqpoint{2.471146in}{1.130351in}}{\pgfqpoint{2.481232in}{1.130351in}}%
\pgfpathquadraticcurveto{\pgfqpoint{2.485735in}{1.130351in}}{\pgfqpoint{2.490148in}{1.131558in}}%
\pgfpathquadraticcurveto{\pgfqpoint{2.494561in}{1.132764in}}{\pgfqpoint{2.498956in}{1.135196in}}%
\pgfpathlineto{\pgfqpoint{2.498956in}{1.125614in}}%
\pgfpathquadraticcurveto{\pgfqpoint{2.494615in}{1.123596in}}{\pgfqpoint{2.489968in}{1.122588in}}%
\pgfpathquadraticcurveto{\pgfqpoint{2.485339in}{1.121561in}}{\pgfqpoint{2.480098in}{1.121561in}}%
\pgfpathquadraticcurveto{\pgfqpoint{2.465850in}{1.121561in}}{\pgfqpoint{2.457456in}{1.130513in}}%
\pgfpathquadraticcurveto{\pgfqpoint{2.449081in}{1.139465in}}{\pgfqpoint{2.449081in}{1.154667in}}%
\pgfpathquadraticcurveto{\pgfqpoint{2.449081in}{1.170086in}}{\pgfqpoint{2.457546in}{1.178912in}}%
\pgfpathquadraticcurveto{\pgfqpoint{2.466030in}{1.187756in}}{\pgfqpoint{2.480782in}{1.187756in}}%
\pgfpathquadraticcurveto{\pgfqpoint{2.485555in}{1.187756in}}{\pgfqpoint{2.490112in}{1.186765in}}%
\pgfpathquadraticcurveto{\pgfqpoint{2.494669in}{1.185792in}}{\pgfqpoint{2.498956in}{1.183829in}}%
\pgfpathlineto{\pgfqpoint{2.498956in}{1.183829in}}%
\pgfpathclose%
\pgfpathmoveto{\pgfqpoint{2.516969in}{1.186243in}}%
\pgfpathlineto{\pgfqpoint{2.527326in}{1.186243in}}%
\pgfpathlineto{\pgfqpoint{2.527326in}{1.123200in}}%
\pgfpathlineto{\pgfqpoint{2.516969in}{1.123200in}}%
\pgfpathlineto{\pgfqpoint{2.516969in}{1.186243in}}%
\pgfpathlineto{\pgfqpoint{2.516969in}{1.186243in}}%
\pgfpathclose%
\pgfpathmoveto{\pgfqpoint{2.516969in}{1.210793in}}%
\pgfpathlineto{\pgfqpoint{2.527326in}{1.210793in}}%
\pgfpathlineto{\pgfqpoint{2.527326in}{1.197662in}}%
\pgfpathlineto{\pgfqpoint{2.516969in}{1.197662in}}%
\pgfpathlineto{\pgfqpoint{2.516969in}{1.210793in}}%
\pgfpathlineto{\pgfqpoint{2.516969in}{1.210793in}}%
\pgfpathclose%
\pgfpathmoveto{\pgfqpoint{2.599933in}{1.186243in}}%
\pgfpathlineto{\pgfqpoint{2.599933in}{1.123200in}}%
\pgfpathlineto{\pgfqpoint{2.589521in}{1.123200in}}%
\pgfpathlineto{\pgfqpoint{2.589521in}{1.178191in}}%
\pgfpathlineto{\pgfqpoint{2.561098in}{1.178191in}}%
\pgfpathlineto{\pgfqpoint{2.561098in}{1.123200in}}%
\pgfpathlineto{\pgfqpoint{2.550687in}{1.123200in}}%
\pgfpathlineto{\pgfqpoint{2.550687in}{1.178191in}}%
\pgfpathlineto{\pgfqpoint{2.540781in}{1.178191in}}%
\pgfpathlineto{\pgfqpoint{2.540781in}{1.186243in}}%
\pgfpathlineto{\pgfqpoint{2.550687in}{1.186243in}}%
\pgfpathlineto{\pgfqpoint{2.550687in}{1.190638in}}%
\pgfpathquadraticcurveto{\pgfqpoint{2.550687in}{1.200940in}}{\pgfqpoint{2.555551in}{1.205858in}}%
\pgfpathquadraticcurveto{\pgfqpoint{2.560432in}{1.210793in}}{\pgfqpoint{2.570501in}{1.210793in}}%
\pgfpathlineto{\pgfqpoint{2.580912in}{1.210793in}}%
\pgfpathlineto{\pgfqpoint{2.580912in}{1.202165in}}%
\pgfpathlineto{\pgfqpoint{2.571005in}{1.202165in}}%
\pgfpathquadraticcurveto{\pgfqpoint{2.565439in}{1.202165in}}{\pgfqpoint{2.563260in}{1.199914in}}%
\pgfpathquadraticcurveto{\pgfqpoint{2.561098in}{1.197662in}}{\pgfqpoint{2.561098in}{1.191808in}}%
\pgfpathlineto{\pgfqpoint{2.561098in}{1.186243in}}%
\pgfpathlineto{\pgfqpoint{2.599933in}{1.186243in}}%
\pgfpathlineto{\pgfqpoint{2.599933in}{1.186243in}}%
\pgfpathclose%
\pgfpathmoveto{\pgfqpoint{2.589521in}{1.210667in}}%
\pgfpathlineto{\pgfqpoint{2.599933in}{1.210667in}}%
\pgfpathlineto{\pgfqpoint{2.599933in}{1.197554in}}%
\pgfpathlineto{\pgfqpoint{2.589521in}{1.197554in}}%
\pgfpathlineto{\pgfqpoint{2.589521in}{1.210667in}}%
\pgfpathlineto{\pgfqpoint{2.589521in}{1.210667in}}%
\pgfpathclose%
\pgfpathmoveto{\pgfqpoint{2.666974in}{1.183829in}}%
\pgfpathlineto{\pgfqpoint{2.666974in}{1.174138in}}%
\pgfpathquadraticcurveto{\pgfqpoint{2.662579in}{1.176570in}}{\pgfqpoint{2.658166in}{1.177777in}}%
\pgfpathquadraticcurveto{\pgfqpoint{2.653753in}{1.178984in}}{\pgfqpoint{2.649250in}{1.178984in}}%
\pgfpathquadraticcurveto{\pgfqpoint{2.639163in}{1.178984in}}{\pgfqpoint{2.633579in}{1.172589in}}%
\pgfpathquadraticcurveto{\pgfqpoint{2.628013in}{1.166213in}}{\pgfqpoint{2.628013in}{1.154667in}}%
\pgfpathquadraticcurveto{\pgfqpoint{2.628013in}{1.143121in}}{\pgfqpoint{2.633579in}{1.136727in}}%
\pgfpathquadraticcurveto{\pgfqpoint{2.639163in}{1.130351in}}{\pgfqpoint{2.649250in}{1.130351in}}%
\pgfpathquadraticcurveto{\pgfqpoint{2.653753in}{1.130351in}}{\pgfqpoint{2.658166in}{1.131558in}}%
\pgfpathquadraticcurveto{\pgfqpoint{2.662579in}{1.132764in}}{\pgfqpoint{2.666974in}{1.135196in}}%
\pgfpathlineto{\pgfqpoint{2.666974in}{1.125614in}}%
\pgfpathquadraticcurveto{\pgfqpoint{2.662633in}{1.123596in}}{\pgfqpoint{2.657986in}{1.122588in}}%
\pgfpathquadraticcurveto{\pgfqpoint{2.653357in}{1.121561in}}{\pgfqpoint{2.648115in}{1.121561in}}%
\pgfpathquadraticcurveto{\pgfqpoint{2.633867in}{1.121561in}}{\pgfqpoint{2.625474in}{1.130513in}}%
\pgfpathquadraticcurveto{\pgfqpoint{2.617098in}{1.139465in}}{\pgfqpoint{2.617098in}{1.154667in}}%
\pgfpathquadraticcurveto{\pgfqpoint{2.617098in}{1.170086in}}{\pgfqpoint{2.625564in}{1.178912in}}%
\pgfpathquadraticcurveto{\pgfqpoint{2.634048in}{1.187756in}}{\pgfqpoint{2.648799in}{1.187756in}}%
\pgfpathquadraticcurveto{\pgfqpoint{2.653573in}{1.187756in}}{\pgfqpoint{2.658130in}{1.186765in}}%
\pgfpathquadraticcurveto{\pgfqpoint{2.662687in}{1.185792in}}{\pgfqpoint{2.666974in}{1.183829in}}%
\pgfpathlineto{\pgfqpoint{2.666974in}{1.183829in}}%
\pgfpathclose%
\pgfpathmoveto{\pgfqpoint{2.713643in}{1.154883in}}%
\pgfpathquadraticcurveto{\pgfqpoint{2.701089in}{1.154883in}}{\pgfqpoint{2.696243in}{1.152019in}}%
\pgfpathquadraticcurveto{\pgfqpoint{2.691398in}{1.149156in}}{\pgfqpoint{2.691398in}{1.142221in}}%
\pgfpathquadraticcurveto{\pgfqpoint{2.691398in}{1.136709in}}{\pgfqpoint{2.695037in}{1.133467in}}%
\pgfpathquadraticcurveto{\pgfqpoint{2.698675in}{1.130243in}}{\pgfqpoint{2.704907in}{1.130243in}}%
\pgfpathquadraticcurveto{\pgfqpoint{2.713535in}{1.130243in}}{\pgfqpoint{2.718741in}{1.136349in}}%
\pgfpathquadraticcurveto{\pgfqpoint{2.723946in}{1.142455in}}{\pgfqpoint{2.723946in}{1.152578in}}%
\pgfpathlineto{\pgfqpoint{2.723946in}{1.154883in}}%
\pgfpathlineto{\pgfqpoint{2.713643in}{1.154883in}}%
\pgfpathlineto{\pgfqpoint{2.713643in}{1.154883in}}%
\pgfpathclose%
\pgfpathmoveto{\pgfqpoint{2.734303in}{1.159170in}}%
\pgfpathlineto{\pgfqpoint{2.734303in}{1.123200in}}%
\pgfpathlineto{\pgfqpoint{2.723946in}{1.123200in}}%
\pgfpathlineto{\pgfqpoint{2.723946in}{1.132764in}}%
\pgfpathquadraticcurveto{\pgfqpoint{2.720398in}{1.127037in}}{\pgfqpoint{2.715102in}{1.124299in}}%
\pgfpathquadraticcurveto{\pgfqpoint{2.709807in}{1.121561in}}{\pgfqpoint{2.702151in}{1.121561in}}%
\pgfpathquadraticcurveto{\pgfqpoint{2.692479in}{1.121561in}}{\pgfqpoint{2.686751in}{1.127001in}}%
\pgfpathquadraticcurveto{\pgfqpoint{2.681041in}{1.132440in}}{\pgfqpoint{2.681041in}{1.141554in}}%
\pgfpathquadraticcurveto{\pgfqpoint{2.681041in}{1.152182in}}{\pgfqpoint{2.688156in}{1.157585in}}%
\pgfpathquadraticcurveto{\pgfqpoint{2.695289in}{1.162989in}}{\pgfqpoint{2.709410in}{1.162989in}}%
\pgfpathlineto{\pgfqpoint{2.723946in}{1.162989in}}%
\pgfpathlineto{\pgfqpoint{2.723946in}{1.164016in}}%
\pgfpathquadraticcurveto{\pgfqpoint{2.723946in}{1.171166in}}{\pgfqpoint{2.719245in}{1.175075in}}%
\pgfpathquadraticcurveto{\pgfqpoint{2.714544in}{1.178984in}}{\pgfqpoint{2.706042in}{1.178984in}}%
\pgfpathquadraticcurveto{\pgfqpoint{2.700638in}{1.178984in}}{\pgfqpoint{2.695505in}{1.177687in}}%
\pgfpathquadraticcurveto{\pgfqpoint{2.690390in}{1.176390in}}{\pgfqpoint{2.685670in}{1.173796in}}%
\pgfpathlineto{\pgfqpoint{2.685670in}{1.183379in}}%
\pgfpathquadraticcurveto{\pgfqpoint{2.691344in}{1.185576in}}{\pgfqpoint{2.696694in}{1.186657in}}%
\pgfpathquadraticcurveto{\pgfqpoint{2.702043in}{1.187756in}}{\pgfqpoint{2.707105in}{1.187756in}}%
\pgfpathquadraticcurveto{\pgfqpoint{2.720794in}{1.187756in}}{\pgfqpoint{2.727549in}{1.180659in}}%
\pgfpathquadraticcurveto{\pgfqpoint{2.734303in}{1.173580in}}{\pgfqpoint{2.734303in}{1.159170in}}%
\pgfpathlineto{\pgfqpoint{2.734303in}{1.159170in}}%
\pgfpathclose%
\pgfpathmoveto{\pgfqpoint{2.765878in}{1.204147in}}%
\pgfpathlineto{\pgfqpoint{2.765878in}{1.186243in}}%
\pgfpathlineto{\pgfqpoint{2.787205in}{1.186243in}}%
\pgfpathlineto{\pgfqpoint{2.787205in}{1.178191in}}%
\pgfpathlineto{\pgfqpoint{2.765878in}{1.178191in}}%
\pgfpathlineto{\pgfqpoint{2.765878in}{1.143968in}}%
\pgfpathquadraticcurveto{\pgfqpoint{2.765878in}{1.136259in}}{\pgfqpoint{2.767986in}{1.134061in}}%
\pgfpathquadraticcurveto{\pgfqpoint{2.770093in}{1.131864in}}{\pgfqpoint{2.776578in}{1.131864in}}%
\pgfpathlineto{\pgfqpoint{2.787205in}{1.131864in}}%
\pgfpathlineto{\pgfqpoint{2.787205in}{1.123200in}}%
\pgfpathlineto{\pgfqpoint{2.776578in}{1.123200in}}%
\pgfpathquadraticcurveto{\pgfqpoint{2.764582in}{1.123200in}}{\pgfqpoint{2.760025in}{1.127667in}}%
\pgfpathquadraticcurveto{\pgfqpoint{2.755467in}{1.132152in}}{\pgfqpoint{2.755467in}{1.143968in}}%
\pgfpathlineto{\pgfqpoint{2.755467in}{1.178191in}}%
\pgfpathlineto{\pgfqpoint{2.747866in}{1.178191in}}%
\pgfpathlineto{\pgfqpoint{2.747866in}{1.186243in}}%
\pgfpathlineto{\pgfqpoint{2.755467in}{1.186243in}}%
\pgfpathlineto{\pgfqpoint{2.755467in}{1.204147in}}%
\pgfpathlineto{\pgfqpoint{2.765878in}{1.204147in}}%
\pgfpathlineto{\pgfqpoint{2.765878in}{1.204147in}}%
\pgfpathclose%
\pgfpathmoveto{\pgfqpoint{2.800822in}{1.186243in}}%
\pgfpathlineto{\pgfqpoint{2.811179in}{1.186243in}}%
\pgfpathlineto{\pgfqpoint{2.811179in}{1.123200in}}%
\pgfpathlineto{\pgfqpoint{2.800822in}{1.123200in}}%
\pgfpathlineto{\pgfqpoint{2.800822in}{1.186243in}}%
\pgfpathlineto{\pgfqpoint{2.800822in}{1.186243in}}%
\pgfpathclose%
\pgfpathmoveto{\pgfqpoint{2.800822in}{1.210793in}}%
\pgfpathlineto{\pgfqpoint{2.811179in}{1.210793in}}%
\pgfpathlineto{\pgfqpoint{2.811179in}{1.197662in}}%
\pgfpathlineto{\pgfqpoint{2.800822in}{1.197662in}}%
\pgfpathlineto{\pgfqpoint{2.800822in}{1.210793in}}%
\pgfpathlineto{\pgfqpoint{2.800822in}{1.210793in}}%
\pgfpathclose%
\pgfpathmoveto{\pgfqpoint{2.857272in}{1.178984in}}%
\pgfpathquadraticcurveto{\pgfqpoint{2.848951in}{1.178984in}}{\pgfqpoint{2.844105in}{1.172481in}}%
\pgfpathquadraticcurveto{\pgfqpoint{2.839260in}{1.165979in}}{\pgfqpoint{2.839260in}{1.154667in}}%
\pgfpathquadraticcurveto{\pgfqpoint{2.839260in}{1.143356in}}{\pgfqpoint{2.844069in}{1.136853in}}%
\pgfpathquadraticcurveto{\pgfqpoint{2.848896in}{1.130351in}}{\pgfqpoint{2.857272in}{1.130351in}}%
\pgfpathquadraticcurveto{\pgfqpoint{2.865558in}{1.130351in}}{\pgfqpoint{2.870385in}{1.136871in}}%
\pgfpathquadraticcurveto{\pgfqpoint{2.875230in}{1.143410in}}{\pgfqpoint{2.875230in}{1.154667in}}%
\pgfpathquadraticcurveto{\pgfqpoint{2.875230in}{1.165871in}}{\pgfqpoint{2.870385in}{1.172427in}}%
\pgfpathquadraticcurveto{\pgfqpoint{2.865558in}{1.178984in}}{\pgfqpoint{2.857272in}{1.178984in}}%
\pgfpathlineto{\pgfqpoint{2.857272in}{1.178984in}}%
\pgfpathclose%
\pgfpathmoveto{\pgfqpoint{2.857272in}{1.187756in}}%
\pgfpathquadraticcurveto{\pgfqpoint{2.870781in}{1.187756in}}{\pgfqpoint{2.878490in}{1.178966in}}%
\pgfpathquadraticcurveto{\pgfqpoint{2.886218in}{1.170194in}}{\pgfqpoint{2.886218in}{1.154667in}}%
\pgfpathquadraticcurveto{\pgfqpoint{2.886218in}{1.139195in}}{\pgfqpoint{2.878490in}{1.130369in}}%
\pgfpathquadraticcurveto{\pgfqpoint{2.870781in}{1.121561in}}{\pgfqpoint{2.857272in}{1.121561in}}%
\pgfpathquadraticcurveto{\pgfqpoint{2.843709in}{1.121561in}}{\pgfqpoint{2.836018in}{1.130369in}}%
\pgfpathquadraticcurveto{\pgfqpoint{2.828345in}{1.139195in}}{\pgfqpoint{2.828345in}{1.154667in}}%
\pgfpathquadraticcurveto{\pgfqpoint{2.828345in}{1.170194in}}{\pgfqpoint{2.836018in}{1.178966in}}%
\pgfpathquadraticcurveto{\pgfqpoint{2.843709in}{1.187756in}}{\pgfqpoint{2.857272in}{1.187756in}}%
\pgfpathlineto{\pgfqpoint{2.857272in}{1.187756in}}%
\pgfpathclose%
\pgfpathmoveto{\pgfqpoint{2.955799in}{1.161260in}}%
\pgfpathlineto{\pgfqpoint{2.955799in}{1.123200in}}%
\pgfpathlineto{\pgfqpoint{2.945442in}{1.123200in}}%
\pgfpathlineto{\pgfqpoint{2.945442in}{1.160917in}}%
\pgfpathquadraticcurveto{\pgfqpoint{2.945442in}{1.169869in}}{\pgfqpoint{2.941947in}{1.174300in}}%
\pgfpathquadraticcurveto{\pgfqpoint{2.938453in}{1.178749in}}{\pgfqpoint{2.931482in}{1.178749in}}%
\pgfpathquadraticcurveto{\pgfqpoint{2.923089in}{1.178749in}}{\pgfqpoint{2.918243in}{1.173400in}}%
\pgfpathquadraticcurveto{\pgfqpoint{2.913398in}{1.168068in}}{\pgfqpoint{2.913398in}{1.158828in}}%
\pgfpathlineto{\pgfqpoint{2.913398in}{1.123200in}}%
\pgfpathlineto{\pgfqpoint{2.902987in}{1.123200in}}%
\pgfpathlineto{\pgfqpoint{2.902987in}{1.186243in}}%
\pgfpathlineto{\pgfqpoint{2.913398in}{1.186243in}}%
\pgfpathlineto{\pgfqpoint{2.913398in}{1.176444in}}%
\pgfpathquadraticcurveto{\pgfqpoint{2.917127in}{1.182136in}}{\pgfqpoint{2.922152in}{1.184946in}}%
\pgfpathquadraticcurveto{\pgfqpoint{2.927195in}{1.187756in}}{\pgfqpoint{2.933788in}{1.187756in}}%
\pgfpathquadraticcurveto{\pgfqpoint{2.944649in}{1.187756in}}{\pgfqpoint{2.950215in}{1.181037in}}%
\pgfpathquadraticcurveto{\pgfqpoint{2.955799in}{1.174318in}}{\pgfqpoint{2.955799in}{1.161260in}}%
\pgfpathlineto{\pgfqpoint{2.955799in}{1.161260in}}%
\pgfpathclose%
\pgfpathmoveto{\pgfqpoint{3.016626in}{1.184387in}}%
\pgfpathlineto{\pgfqpoint{3.016626in}{1.174589in}}%
\pgfpathquadraticcurveto{\pgfqpoint{3.012249in}{1.176840in}}{\pgfqpoint{3.007512in}{1.177957in}}%
\pgfpathquadraticcurveto{\pgfqpoint{3.002792in}{1.179092in}}{\pgfqpoint{2.997713in}{1.179092in}}%
\pgfpathquadraticcurveto{\pgfqpoint{2.990004in}{1.179092in}}{\pgfqpoint{2.986149in}{1.176732in}}%
\pgfpathquadraticcurveto{\pgfqpoint{2.982294in}{1.174373in}}{\pgfqpoint{2.982294in}{1.169635in}}%
\pgfpathquadraticcurveto{\pgfqpoint{2.982294in}{1.166033in}}{\pgfqpoint{2.985050in}{1.163980in}}%
\pgfpathquadraticcurveto{\pgfqpoint{2.987806in}{1.161926in}}{\pgfqpoint{2.996146in}{1.160071in}}%
\pgfpathlineto{\pgfqpoint{2.999694in}{1.159278in}}%
\pgfpathquadraticcurveto{\pgfqpoint{3.010718in}{1.156919in}}{\pgfqpoint{3.015365in}{1.152614in}}%
\pgfpathquadraticcurveto{\pgfqpoint{3.020012in}{1.148309in}}{\pgfqpoint{3.020012in}{1.140600in}}%
\pgfpathquadraticcurveto{\pgfqpoint{3.020012in}{1.131810in}}{\pgfqpoint{3.013059in}{1.126676in}}%
\pgfpathquadraticcurveto{\pgfqpoint{3.006107in}{1.121561in}}{\pgfqpoint{2.993948in}{1.121561in}}%
\pgfpathquadraticcurveto{\pgfqpoint{2.988887in}{1.121561in}}{\pgfqpoint{2.983393in}{1.122552in}}%
\pgfpathquadraticcurveto{\pgfqpoint{2.977900in}{1.123542in}}{\pgfqpoint{2.971829in}{1.125506in}}%
\pgfpathlineto{\pgfqpoint{2.971829in}{1.136205in}}%
\pgfpathquadraticcurveto{\pgfqpoint{2.977575in}{1.133215in}}{\pgfqpoint{2.983141in}{1.131720in}}%
\pgfpathquadraticcurveto{\pgfqpoint{2.988707in}{1.130243in}}{\pgfqpoint{2.994183in}{1.130243in}}%
\pgfpathquadraticcurveto{\pgfqpoint{3.001495in}{1.130243in}}{\pgfqpoint{3.005422in}{1.132746in}}%
\pgfpathquadraticcurveto{\pgfqpoint{3.009367in}{1.135250in}}{\pgfqpoint{3.009367in}{1.139807in}}%
\pgfpathquadraticcurveto{\pgfqpoint{3.009367in}{1.144022in}}{\pgfqpoint{3.006521in}{1.146274in}}%
\pgfpathquadraticcurveto{\pgfqpoint{3.003693in}{1.148525in}}{\pgfqpoint{2.994056in}{1.150614in}}%
\pgfpathlineto{\pgfqpoint{2.990454in}{1.151461in}}%
\pgfpathquadraticcurveto{\pgfqpoint{2.980836in}{1.153478in}}{\pgfqpoint{2.976549in}{1.157675in}}%
\pgfpathquadraticcurveto{\pgfqpoint{2.972280in}{1.161872in}}{\pgfqpoint{2.972280in}{1.169185in}}%
\pgfpathquadraticcurveto{\pgfqpoint{2.972280in}{1.178083in}}{\pgfqpoint{2.978584in}{1.182910in}}%
\pgfpathquadraticcurveto{\pgfqpoint{2.984888in}{1.187756in}}{\pgfqpoint{2.996488in}{1.187756in}}%
\pgfpathquadraticcurveto{\pgfqpoint{3.002216in}{1.187756in}}{\pgfqpoint{3.007277in}{1.186909in}}%
\pgfpathquadraticcurveto{\pgfqpoint{3.012357in}{1.186080in}}{\pgfqpoint{3.016626in}{1.184387in}}%
\pgfpathlineto{\pgfqpoint{3.016626in}{1.184387in}}%
\pgfpathclose%
\pgfpathmoveto{\pgfqpoint{3.037952in}{1.137502in}}%
\pgfpathlineto{\pgfqpoint{3.049840in}{1.137502in}}%
\pgfpathlineto{\pgfqpoint{3.049840in}{1.123200in}}%
\pgfpathlineto{\pgfqpoint{3.037952in}{1.123200in}}%
\pgfpathlineto{\pgfqpoint{3.037952in}{1.137502in}}%
\pgfpathlineto{\pgfqpoint{3.037952in}{1.137502in}}%
\pgfpathclose%
\pgfusepath{fill}%
\end{pgfscope}%
\begin{pgfscope}%
\pgfsetbuttcap%
\pgfsetroundjoin%
\definecolor{currentfill}{rgb}{0.309804,0.372549,0.419608}%
\pgfsetfillcolor{currentfill}%
\pgfsetlinewidth{0.000000pt}%
\definecolor{currentstroke}{rgb}{0.309804,0.372549,0.419608}%
\pgfsetstrokecolor{currentstroke}%
\pgfsetdash{}{0pt}%
\pgfpathmoveto{\pgfqpoint{3.219265in}{1.204489in}}%
\pgfpathlineto{\pgfqpoint{3.219265in}{1.193393in}}%
\pgfpathquadraticcurveto{\pgfqpoint{3.212798in}{1.196491in}}{\pgfqpoint{3.207052in}{1.198004in}}%
\pgfpathquadraticcurveto{\pgfqpoint{3.201306in}{1.199536in}}{\pgfqpoint{3.195957in}{1.199536in}}%
\pgfpathquadraticcurveto{\pgfqpoint{3.186681in}{1.199536in}}{\pgfqpoint{3.181637in}{1.195933in}}%
\pgfpathquadraticcurveto{\pgfqpoint{3.176594in}{1.192331in}}{\pgfqpoint{3.176594in}{1.185684in}}%
\pgfpathquadraticcurveto{\pgfqpoint{3.176594in}{1.180100in}}{\pgfqpoint{3.179944in}{1.177254in}}%
\pgfpathquadraticcurveto{\pgfqpoint{3.183294in}{1.174427in}}{\pgfqpoint{3.192643in}{1.172679in}}%
\pgfpathlineto{\pgfqpoint{3.199505in}{1.171274in}}%
\pgfpathquadraticcurveto{\pgfqpoint{3.212222in}{1.168843in}}{\pgfqpoint{3.218274in}{1.162737in}}%
\pgfpathquadraticcurveto{\pgfqpoint{3.224326in}{1.156631in}}{\pgfqpoint{3.224326in}{1.146400in}}%
\pgfpathquadraticcurveto{\pgfqpoint{3.224326in}{1.134169in}}{\pgfqpoint{3.216130in}{1.127865in}}%
\pgfpathquadraticcurveto{\pgfqpoint{3.207953in}{1.121561in}}{\pgfqpoint{3.192138in}{1.121561in}}%
\pgfpathquadraticcurveto{\pgfqpoint{3.186176in}{1.121561in}}{\pgfqpoint{3.179440in}{1.122912in}}%
\pgfpathquadraticcurveto{\pgfqpoint{3.172721in}{1.124263in}}{\pgfqpoint{3.165516in}{1.126911in}}%
\pgfpathlineto{\pgfqpoint{3.165516in}{1.138618in}}%
\pgfpathquadraticcurveto{\pgfqpoint{3.172433in}{1.134746in}}{\pgfqpoint{3.179079in}{1.132764in}}%
\pgfpathquadraticcurveto{\pgfqpoint{3.185726in}{1.130801in}}{\pgfqpoint{3.192138in}{1.130801in}}%
\pgfpathquadraticcurveto{\pgfqpoint{3.201865in}{1.130801in}}{\pgfqpoint{3.207160in}{1.134620in}}%
\pgfpathquadraticcurveto{\pgfqpoint{3.212456in}{1.138456in}}{\pgfqpoint{3.212456in}{1.145553in}}%
\pgfpathquadraticcurveto{\pgfqpoint{3.212456in}{1.151731in}}{\pgfqpoint{3.208655in}{1.155226in}}%
\pgfpathquadraticcurveto{\pgfqpoint{3.204855in}{1.158720in}}{\pgfqpoint{3.196191in}{1.160467in}}%
\pgfpathlineto{\pgfqpoint{3.189256in}{1.161818in}}%
\pgfpathquadraticcurveto{\pgfqpoint{3.176540in}{1.164340in}}{\pgfqpoint{3.170848in}{1.169743in}}%
\pgfpathquadraticcurveto{\pgfqpoint{3.165174in}{1.175147in}}{\pgfqpoint{3.165174in}{1.184784in}}%
\pgfpathquadraticcurveto{\pgfqpoint{3.165174in}{1.195933in}}{\pgfqpoint{3.173027in}{1.202345in}}%
\pgfpathquadraticcurveto{\pgfqpoint{3.180881in}{1.208758in}}{\pgfqpoint{3.194660in}{1.208758in}}%
\pgfpathquadraticcurveto{\pgfqpoint{3.200586in}{1.208758in}}{\pgfqpoint{3.206710in}{1.207677in}}%
\pgfpathquadraticcurveto{\pgfqpoint{3.212852in}{1.206614in}}{\pgfqpoint{3.219265in}{1.204489in}}%
\pgfpathlineto{\pgfqpoint{3.219265in}{1.204489in}}%
\pgfpathclose%
\pgfpathmoveto{\pgfqpoint{3.247814in}{1.154667in}}%
\pgfpathquadraticcurveto{\pgfqpoint{3.247814in}{1.143248in}}{\pgfqpoint{3.252515in}{1.136745in}}%
\pgfpathquadraticcurveto{\pgfqpoint{3.257216in}{1.130243in}}{\pgfqpoint{3.265430in}{1.130243in}}%
\pgfpathquadraticcurveto{\pgfqpoint{3.273643in}{1.130243in}}{\pgfqpoint{3.278362in}{1.136745in}}%
\pgfpathquadraticcurveto{\pgfqpoint{3.283100in}{1.143248in}}{\pgfqpoint{3.283100in}{1.154667in}}%
\pgfpathquadraticcurveto{\pgfqpoint{3.283100in}{1.166087in}}{\pgfqpoint{3.278362in}{1.172589in}}%
\pgfpathquadraticcurveto{\pgfqpoint{3.273643in}{1.179092in}}{\pgfqpoint{3.265430in}{1.179092in}}%
\pgfpathquadraticcurveto{\pgfqpoint{3.257216in}{1.179092in}}{\pgfqpoint{3.252515in}{1.172589in}}%
\pgfpathquadraticcurveto{\pgfqpoint{3.247814in}{1.166087in}}{\pgfqpoint{3.247814in}{1.154667in}}%
\pgfpathlineto{\pgfqpoint{3.247814in}{1.154667in}}%
\pgfpathclose%
\pgfpathmoveto{\pgfqpoint{3.283100in}{1.132656in}}%
\pgfpathquadraticcurveto{\pgfqpoint{3.279839in}{1.127037in}}{\pgfqpoint{3.274850in}{1.124299in}}%
\pgfpathquadraticcurveto{\pgfqpoint{3.269879in}{1.121561in}}{\pgfqpoint{3.262890in}{1.121561in}}%
\pgfpathquadraticcurveto{\pgfqpoint{3.251470in}{1.121561in}}{\pgfqpoint{3.244283in}{1.130675in}}%
\pgfpathquadraticcurveto{\pgfqpoint{3.237115in}{1.139807in}}{\pgfqpoint{3.237115in}{1.154667in}}%
\pgfpathquadraticcurveto{\pgfqpoint{3.237115in}{1.169527in}}{\pgfqpoint{3.244283in}{1.178641in}}%
\pgfpathquadraticcurveto{\pgfqpoint{3.251470in}{1.187756in}}{\pgfqpoint{3.262890in}{1.187756in}}%
\pgfpathquadraticcurveto{\pgfqpoint{3.269879in}{1.187756in}}{\pgfqpoint{3.274850in}{1.185018in}}%
\pgfpathquadraticcurveto{\pgfqpoint{3.279839in}{1.182298in}}{\pgfqpoint{3.283100in}{1.176678in}}%
\pgfpathlineto{\pgfqpoint{3.283100in}{1.186243in}}%
\pgfpathlineto{\pgfqpoint{3.293457in}{1.186243in}}%
\pgfpathlineto{\pgfqpoint{3.293457in}{1.099226in}}%
\pgfpathlineto{\pgfqpoint{3.283100in}{1.099226in}}%
\pgfpathlineto{\pgfqpoint{3.283100in}{1.132656in}}%
\pgfpathlineto{\pgfqpoint{3.283100in}{1.132656in}}%
\pgfpathclose%
\pgfpathmoveto{\pgfqpoint{3.313738in}{1.148075in}}%
\pgfpathlineto{\pgfqpoint{3.313738in}{1.186243in}}%
\pgfpathlineto{\pgfqpoint{3.324095in}{1.186243in}}%
\pgfpathlineto{\pgfqpoint{3.324095in}{1.148471in}}%
\pgfpathquadraticcurveto{\pgfqpoint{3.324095in}{1.139519in}}{\pgfqpoint{3.327572in}{1.135034in}}%
\pgfpathquadraticcurveto{\pgfqpoint{3.331066in}{1.130567in}}{\pgfqpoint{3.338055in}{1.130567in}}%
\pgfpathquadraticcurveto{\pgfqpoint{3.346430in}{1.130567in}}{\pgfqpoint{3.351294in}{1.135917in}}%
\pgfpathquadraticcurveto{\pgfqpoint{3.356175in}{1.141266in}}{\pgfqpoint{3.356175in}{1.150506in}}%
\pgfpathlineto{\pgfqpoint{3.356175in}{1.186243in}}%
\pgfpathlineto{\pgfqpoint{3.366532in}{1.186243in}}%
\pgfpathlineto{\pgfqpoint{3.366532in}{1.123200in}}%
\pgfpathlineto{\pgfqpoint{3.356175in}{1.123200in}}%
\pgfpathlineto{\pgfqpoint{3.356175in}{1.132891in}}%
\pgfpathquadraticcurveto{\pgfqpoint{3.352410in}{1.127145in}}{\pgfqpoint{3.347421in}{1.124353in}}%
\pgfpathquadraticcurveto{\pgfqpoint{3.342450in}{1.121561in}}{\pgfqpoint{3.335857in}{1.121561in}}%
\pgfpathquadraticcurveto{\pgfqpoint{3.324996in}{1.121561in}}{\pgfqpoint{3.319358in}{1.128315in}}%
\pgfpathquadraticcurveto{\pgfqpoint{3.313738in}{1.135070in}}{\pgfqpoint{3.313738in}{1.148075in}}%
\pgfpathlineto{\pgfqpoint{3.313738in}{1.148075in}}%
\pgfpathclose%
\pgfpathmoveto{\pgfqpoint{3.339802in}{1.187756in}}%
\pgfpathlineto{\pgfqpoint{3.339802in}{1.187756in}}%
\pgfpathlineto{\pgfqpoint{3.339802in}{1.187756in}}%
\pgfpathclose%
\pgfpathmoveto{\pgfqpoint{3.416516in}{1.154883in}}%
\pgfpathquadraticcurveto{\pgfqpoint{3.403961in}{1.154883in}}{\pgfqpoint{3.399116in}{1.152019in}}%
\pgfpathquadraticcurveto{\pgfqpoint{3.394271in}{1.149156in}}{\pgfqpoint{3.394271in}{1.142221in}}%
\pgfpathquadraticcurveto{\pgfqpoint{3.394271in}{1.136709in}}{\pgfqpoint{3.397909in}{1.133467in}}%
\pgfpathquadraticcurveto{\pgfqpoint{3.401548in}{1.130243in}}{\pgfqpoint{3.407780in}{1.130243in}}%
\pgfpathquadraticcurveto{\pgfqpoint{3.416408in}{1.130243in}}{\pgfqpoint{3.421613in}{1.136349in}}%
\pgfpathquadraticcurveto{\pgfqpoint{3.426819in}{1.142455in}}{\pgfqpoint{3.426819in}{1.152578in}}%
\pgfpathlineto{\pgfqpoint{3.426819in}{1.154883in}}%
\pgfpathlineto{\pgfqpoint{3.416516in}{1.154883in}}%
\pgfpathlineto{\pgfqpoint{3.416516in}{1.154883in}}%
\pgfpathclose%
\pgfpathmoveto{\pgfqpoint{3.437176in}{1.159170in}}%
\pgfpathlineto{\pgfqpoint{3.437176in}{1.123200in}}%
\pgfpathlineto{\pgfqpoint{3.426819in}{1.123200in}}%
\pgfpathlineto{\pgfqpoint{3.426819in}{1.132764in}}%
\pgfpathquadraticcurveto{\pgfqpoint{3.423270in}{1.127037in}}{\pgfqpoint{3.417975in}{1.124299in}}%
\pgfpathquadraticcurveto{\pgfqpoint{3.412679in}{1.121561in}}{\pgfqpoint{3.405024in}{1.121561in}}%
\pgfpathquadraticcurveto{\pgfqpoint{3.395351in}{1.121561in}}{\pgfqpoint{3.389623in}{1.127001in}}%
\pgfpathquadraticcurveto{\pgfqpoint{3.383914in}{1.132440in}}{\pgfqpoint{3.383914in}{1.141554in}}%
\pgfpathquadraticcurveto{\pgfqpoint{3.383914in}{1.152182in}}{\pgfqpoint{3.391028in}{1.157585in}}%
\pgfpathquadraticcurveto{\pgfqpoint{3.398161in}{1.162989in}}{\pgfqpoint{3.412283in}{1.162989in}}%
\pgfpathlineto{\pgfqpoint{3.426819in}{1.162989in}}%
\pgfpathlineto{\pgfqpoint{3.426819in}{1.164016in}}%
\pgfpathquadraticcurveto{\pgfqpoint{3.426819in}{1.171166in}}{\pgfqpoint{3.422117in}{1.175075in}}%
\pgfpathquadraticcurveto{\pgfqpoint{3.417416in}{1.178984in}}{\pgfqpoint{3.408914in}{1.178984in}}%
\pgfpathquadraticcurveto{\pgfqpoint{3.403511in}{1.178984in}}{\pgfqpoint{3.398377in}{1.177687in}}%
\pgfpathquadraticcurveto{\pgfqpoint{3.393262in}{1.176390in}}{\pgfqpoint{3.388543in}{1.173796in}}%
\pgfpathlineto{\pgfqpoint{3.388543in}{1.183379in}}%
\pgfpathquadraticcurveto{\pgfqpoint{3.394217in}{1.185576in}}{\pgfqpoint{3.399566in}{1.186657in}}%
\pgfpathquadraticcurveto{\pgfqpoint{3.404916in}{1.187756in}}{\pgfqpoint{3.409977in}{1.187756in}}%
\pgfpathquadraticcurveto{\pgfqpoint{3.423666in}{1.187756in}}{\pgfqpoint{3.430421in}{1.180659in}}%
\pgfpathquadraticcurveto{\pgfqpoint{3.437176in}{1.173580in}}{\pgfqpoint{3.437176in}{1.159170in}}%
\pgfpathlineto{\pgfqpoint{3.437176in}{1.159170in}}%
\pgfpathclose%
\pgfpathmoveto{\pgfqpoint{3.495031in}{1.176570in}}%
\pgfpathquadraticcurveto{\pgfqpoint{3.493283in}{1.177579in}}{\pgfqpoint{3.491230in}{1.178047in}}%
\pgfpathquadraticcurveto{\pgfqpoint{3.489177in}{1.178533in}}{\pgfqpoint{3.486709in}{1.178533in}}%
\pgfpathquadraticcurveto{\pgfqpoint{3.477919in}{1.178533in}}{\pgfqpoint{3.473218in}{1.172823in}}%
\pgfpathquadraticcurveto{\pgfqpoint{3.468517in}{1.167114in}}{\pgfqpoint{3.468517in}{1.156414in}}%
\pgfpathlineto{\pgfqpoint{3.468517in}{1.123200in}}%
\pgfpathlineto{\pgfqpoint{3.458106in}{1.123200in}}%
\pgfpathlineto{\pgfqpoint{3.458106in}{1.186243in}}%
\pgfpathlineto{\pgfqpoint{3.468517in}{1.186243in}}%
\pgfpathlineto{\pgfqpoint{3.468517in}{1.176444in}}%
\pgfpathquadraticcurveto{\pgfqpoint{3.471795in}{1.182190in}}{\pgfqpoint{3.477018in}{1.184964in}}%
\pgfpathquadraticcurveto{\pgfqpoint{3.482260in}{1.187756in}}{\pgfqpoint{3.489753in}{1.187756in}}%
\pgfpathquadraticcurveto{\pgfqpoint{3.490816in}{1.187756in}}{\pgfqpoint{3.492113in}{1.187611in}}%
\pgfpathquadraticcurveto{\pgfqpoint{3.493410in}{1.187485in}}{\pgfqpoint{3.494977in}{1.187197in}}%
\pgfpathlineto{\pgfqpoint{3.495031in}{1.176570in}}%
\pgfpathlineto{\pgfqpoint{3.495031in}{1.176570in}}%
\pgfpathclose%
\pgfpathmoveto{\pgfqpoint{3.557281in}{1.157315in}}%
\pgfpathlineto{\pgfqpoint{3.557281in}{1.152254in}}%
\pgfpathlineto{\pgfqpoint{3.509656in}{1.152254in}}%
\pgfpathquadraticcurveto{\pgfqpoint{3.510341in}{1.141554in}}{\pgfqpoint{3.516105in}{1.135953in}}%
\pgfpathquadraticcurveto{\pgfqpoint{3.521869in}{1.130351in}}{\pgfqpoint{3.532172in}{1.130351in}}%
\pgfpathquadraticcurveto{\pgfqpoint{3.538134in}{1.130351in}}{\pgfqpoint{3.543735in}{1.131810in}}%
\pgfpathquadraticcurveto{\pgfqpoint{3.549337in}{1.133269in}}{\pgfqpoint{3.554867in}{1.136205in}}%
\pgfpathlineto{\pgfqpoint{3.554867in}{1.126406in}}%
\pgfpathquadraticcurveto{\pgfqpoint{3.549283in}{1.124047in}}{\pgfqpoint{3.543429in}{1.122804in}}%
\pgfpathquadraticcurveto{\pgfqpoint{3.537575in}{1.121561in}}{\pgfqpoint{3.531559in}{1.121561in}}%
\pgfpathquadraticcurveto{\pgfqpoint{3.516465in}{1.121561in}}{\pgfqpoint{3.507657in}{1.130333in}}%
\pgfpathquadraticcurveto{\pgfqpoint{3.498849in}{1.139123in}}{\pgfqpoint{3.498849in}{1.154109in}}%
\pgfpathquadraticcurveto{\pgfqpoint{3.498849in}{1.169581in}}{\pgfqpoint{3.507207in}{1.178659in}}%
\pgfpathquadraticcurveto{\pgfqpoint{3.515564in}{1.187756in}}{\pgfqpoint{3.529758in}{1.187756in}}%
\pgfpathquadraticcurveto{\pgfqpoint{3.542475in}{1.187756in}}{\pgfqpoint{3.549878in}{1.179560in}}%
\pgfpathquadraticcurveto{\pgfqpoint{3.557281in}{1.171383in}}{\pgfqpoint{3.557281in}{1.157315in}}%
\pgfpathlineto{\pgfqpoint{3.557281in}{1.157315in}}%
\pgfpathclose%
\pgfpathmoveto{\pgfqpoint{3.546924in}{1.160359in}}%
\pgfpathquadraticcurveto{\pgfqpoint{3.546816in}{1.168843in}}{\pgfqpoint{3.542168in}{1.173904in}}%
\pgfpathquadraticcurveto{\pgfqpoint{3.537521in}{1.178984in}}{\pgfqpoint{3.529866in}{1.178984in}}%
\pgfpathquadraticcurveto{\pgfqpoint{3.521202in}{1.178984in}}{\pgfqpoint{3.515997in}{1.174084in}}%
\pgfpathquadraticcurveto{\pgfqpoint{3.510791in}{1.169185in}}{\pgfqpoint{3.509999in}{1.160287in}}%
\pgfpathlineto{\pgfqpoint{3.546924in}{1.160359in}}%
\pgfpathlineto{\pgfqpoint{3.546924in}{1.160359in}}%
\pgfpathclose%
\pgfpathmoveto{\pgfqpoint{3.614469in}{1.184387in}}%
\pgfpathlineto{\pgfqpoint{3.614469in}{1.174589in}}%
\pgfpathquadraticcurveto{\pgfqpoint{3.610092in}{1.176840in}}{\pgfqpoint{3.605355in}{1.177957in}}%
\pgfpathquadraticcurveto{\pgfqpoint{3.600636in}{1.179092in}}{\pgfqpoint{3.595556in}{1.179092in}}%
\pgfpathquadraticcurveto{\pgfqpoint{3.587847in}{1.179092in}}{\pgfqpoint{3.583993in}{1.176732in}}%
\pgfpathquadraticcurveto{\pgfqpoint{3.580138in}{1.174373in}}{\pgfqpoint{3.580138in}{1.169635in}}%
\pgfpathquadraticcurveto{\pgfqpoint{3.580138in}{1.166033in}}{\pgfqpoint{3.582894in}{1.163980in}}%
\pgfpathquadraticcurveto{\pgfqpoint{3.585650in}{1.161926in}}{\pgfqpoint{3.593989in}{1.160071in}}%
\pgfpathlineto{\pgfqpoint{3.597538in}{1.159278in}}%
\pgfpathquadraticcurveto{\pgfqpoint{3.608561in}{1.156919in}}{\pgfqpoint{3.613208in}{1.152614in}}%
\pgfpathquadraticcurveto{\pgfqpoint{3.617855in}{1.148309in}}{\pgfqpoint{3.617855in}{1.140600in}}%
\pgfpathquadraticcurveto{\pgfqpoint{3.617855in}{1.131810in}}{\pgfqpoint{3.610903in}{1.126676in}}%
\pgfpathquadraticcurveto{\pgfqpoint{3.603950in}{1.121561in}}{\pgfqpoint{3.591792in}{1.121561in}}%
\pgfpathquadraticcurveto{\pgfqpoint{3.586730in}{1.121561in}}{\pgfqpoint{3.581237in}{1.122552in}}%
\pgfpathquadraticcurveto{\pgfqpoint{3.575743in}{1.123542in}}{\pgfqpoint{3.569673in}{1.125506in}}%
\pgfpathlineto{\pgfqpoint{3.569673in}{1.136205in}}%
\pgfpathquadraticcurveto{\pgfqpoint{3.575419in}{1.133215in}}{\pgfqpoint{3.580985in}{1.131720in}}%
\pgfpathquadraticcurveto{\pgfqpoint{3.586550in}{1.130243in}}{\pgfqpoint{3.592026in}{1.130243in}}%
\pgfpathquadraticcurveto{\pgfqpoint{3.599339in}{1.130243in}}{\pgfqpoint{3.603266in}{1.132746in}}%
\pgfpathquadraticcurveto{\pgfqpoint{3.607210in}{1.135250in}}{\pgfqpoint{3.607210in}{1.139807in}}%
\pgfpathquadraticcurveto{\pgfqpoint{3.607210in}{1.144022in}}{\pgfqpoint{3.604364in}{1.146274in}}%
\pgfpathquadraticcurveto{\pgfqpoint{3.601536in}{1.148525in}}{\pgfqpoint{3.591900in}{1.150614in}}%
\pgfpathlineto{\pgfqpoint{3.588298in}{1.151461in}}%
\pgfpathquadraticcurveto{\pgfqpoint{3.578679in}{1.153478in}}{\pgfqpoint{3.574392in}{1.157675in}}%
\pgfpathquadraticcurveto{\pgfqpoint{3.570123in}{1.161872in}}{\pgfqpoint{3.570123in}{1.169185in}}%
\pgfpathquadraticcurveto{\pgfqpoint{3.570123in}{1.178083in}}{\pgfqpoint{3.576428in}{1.182910in}}%
\pgfpathquadraticcurveto{\pgfqpoint{3.582732in}{1.187756in}}{\pgfqpoint{3.594332in}{1.187756in}}%
\pgfpathquadraticcurveto{\pgfqpoint{3.600059in}{1.187756in}}{\pgfqpoint{3.605121in}{1.186909in}}%
\pgfpathquadraticcurveto{\pgfqpoint{3.610200in}{1.186080in}}{\pgfqpoint{3.614469in}{1.184387in}}%
\pgfpathlineto{\pgfqpoint{3.614469in}{1.184387in}}%
\pgfpathclose%
\pgfpathmoveto{\pgfqpoint{3.636984in}{1.137502in}}%
\pgfpathlineto{\pgfqpoint{3.648854in}{1.137502in}}%
\pgfpathlineto{\pgfqpoint{3.648854in}{1.123200in}}%
\pgfpathlineto{\pgfqpoint{3.636984in}{1.123200in}}%
\pgfpathlineto{\pgfqpoint{3.636984in}{1.137502in}}%
\pgfpathlineto{\pgfqpoint{3.636984in}{1.137502in}}%
\pgfpathclose%
\pgfpathmoveto{\pgfqpoint{3.636984in}{1.182802in}}%
\pgfpathlineto{\pgfqpoint{3.648854in}{1.182802in}}%
\pgfpathlineto{\pgfqpoint{3.648854in}{1.168519in}}%
\pgfpathlineto{\pgfqpoint{3.636984in}{1.168519in}}%
\pgfpathlineto{\pgfqpoint{3.636984in}{1.182802in}}%
\pgfpathlineto{\pgfqpoint{3.636984in}{1.182802in}}%
\pgfpathclose%
\pgfusepath{fill}%
\end{pgfscope}%
\begin{pgfscope}%
\pgfsetbuttcap%
\pgfsetroundjoin%
\definecolor{currentfill}{rgb}{0.309804,0.372549,0.419608}%
\pgfsetfillcolor{currentfill}%
\pgfsetlinewidth{0.000000pt}%
\definecolor{currentstroke}{rgb}{0.309804,0.372549,0.419608}%
\pgfsetstrokecolor{currentstroke}%
\pgfsetdash{}{0pt}%
\pgfpathmoveto{\pgfqpoint{3.808881in}{1.157315in}}%
\pgfpathlineto{\pgfqpoint{3.808881in}{1.152254in}}%
\pgfpathlineto{\pgfqpoint{3.761257in}{1.152254in}}%
\pgfpathquadraticcurveto{\pgfqpoint{3.761941in}{1.141554in}}{\pgfqpoint{3.767705in}{1.135953in}}%
\pgfpathquadraticcurveto{\pgfqpoint{3.773469in}{1.130351in}}{\pgfqpoint{3.783772in}{1.130351in}}%
\pgfpathquadraticcurveto{\pgfqpoint{3.789734in}{1.130351in}}{\pgfqpoint{3.795336in}{1.131810in}}%
\pgfpathquadraticcurveto{\pgfqpoint{3.800937in}{1.133269in}}{\pgfqpoint{3.806467in}{1.136205in}}%
\pgfpathlineto{\pgfqpoint{3.806467in}{1.126406in}}%
\pgfpathquadraticcurveto{\pgfqpoint{3.800883in}{1.124047in}}{\pgfqpoint{3.795030in}{1.122804in}}%
\pgfpathquadraticcurveto{\pgfqpoint{3.789176in}{1.121561in}}{\pgfqpoint{3.783160in}{1.121561in}}%
\pgfpathquadraticcurveto{\pgfqpoint{3.768065in}{1.121561in}}{\pgfqpoint{3.759257in}{1.130333in}}%
\pgfpathquadraticcurveto{\pgfqpoint{3.750449in}{1.139123in}}{\pgfqpoint{3.750449in}{1.154109in}}%
\pgfpathquadraticcurveto{\pgfqpoint{3.750449in}{1.169581in}}{\pgfqpoint{3.758807in}{1.178659in}}%
\pgfpathquadraticcurveto{\pgfqpoint{3.767165in}{1.187756in}}{\pgfqpoint{3.781358in}{1.187756in}}%
\pgfpathquadraticcurveto{\pgfqpoint{3.794075in}{1.187756in}}{\pgfqpoint{3.801478in}{1.179560in}}%
\pgfpathquadraticcurveto{\pgfqpoint{3.808881in}{1.171383in}}{\pgfqpoint{3.808881in}{1.157315in}}%
\pgfpathlineto{\pgfqpoint{3.808881in}{1.157315in}}%
\pgfpathclose%
\pgfpathmoveto{\pgfqpoint{3.798524in}{1.160359in}}%
\pgfpathquadraticcurveto{\pgfqpoint{3.798416in}{1.168843in}}{\pgfqpoint{3.793769in}{1.173904in}}%
\pgfpathquadraticcurveto{\pgfqpoint{3.789122in}{1.178984in}}{\pgfqpoint{3.781466in}{1.178984in}}%
\pgfpathquadraticcurveto{\pgfqpoint{3.772803in}{1.178984in}}{\pgfqpoint{3.767597in}{1.174084in}}%
\pgfpathquadraticcurveto{\pgfqpoint{3.762391in}{1.169185in}}{\pgfqpoint{3.761599in}{1.160287in}}%
\pgfpathlineto{\pgfqpoint{3.798524in}{1.160359in}}%
\pgfpathlineto{\pgfqpoint{3.798524in}{1.160359in}}%
\pgfpathclose%
\pgfpathmoveto{\pgfqpoint{3.818463in}{1.186243in}}%
\pgfpathlineto{\pgfqpoint{3.829433in}{1.186243in}}%
\pgfpathlineto{\pgfqpoint{3.849138in}{1.133341in}}%
\pgfpathlineto{\pgfqpoint{3.868843in}{1.186243in}}%
\pgfpathlineto{\pgfqpoint{3.879813in}{1.186243in}}%
\pgfpathlineto{\pgfqpoint{3.856163in}{1.123200in}}%
\pgfpathlineto{\pgfqpoint{3.842095in}{1.123200in}}%
\pgfpathlineto{\pgfqpoint{3.818463in}{1.186243in}}%
\pgfpathlineto{\pgfqpoint{3.818463in}{1.186243in}}%
\pgfpathclose%
\pgfpathmoveto{\pgfqpoint{3.948043in}{1.157315in}}%
\pgfpathlineto{\pgfqpoint{3.948043in}{1.152254in}}%
\pgfpathlineto{\pgfqpoint{3.900419in}{1.152254in}}%
\pgfpathquadraticcurveto{\pgfqpoint{3.901103in}{1.141554in}}{\pgfqpoint{3.906867in}{1.135953in}}%
\pgfpathquadraticcurveto{\pgfqpoint{3.912631in}{1.130351in}}{\pgfqpoint{3.922934in}{1.130351in}}%
\pgfpathquadraticcurveto{\pgfqpoint{3.928896in}{1.130351in}}{\pgfqpoint{3.934498in}{1.131810in}}%
\pgfpathquadraticcurveto{\pgfqpoint{3.940099in}{1.133269in}}{\pgfqpoint{3.945629in}{1.136205in}}%
\pgfpathlineto{\pgfqpoint{3.945629in}{1.126406in}}%
\pgfpathquadraticcurveto{\pgfqpoint{3.940045in}{1.124047in}}{\pgfqpoint{3.934191in}{1.122804in}}%
\pgfpathquadraticcurveto{\pgfqpoint{3.928337in}{1.121561in}}{\pgfqpoint{3.922321in}{1.121561in}}%
\pgfpathquadraticcurveto{\pgfqpoint{3.907227in}{1.121561in}}{\pgfqpoint{3.898419in}{1.130333in}}%
\pgfpathquadraticcurveto{\pgfqpoint{3.889611in}{1.139123in}}{\pgfqpoint{3.889611in}{1.154109in}}%
\pgfpathquadraticcurveto{\pgfqpoint{3.889611in}{1.169581in}}{\pgfqpoint{3.897969in}{1.178659in}}%
\pgfpathquadraticcurveto{\pgfqpoint{3.906327in}{1.187756in}}{\pgfqpoint{3.920520in}{1.187756in}}%
\pgfpathquadraticcurveto{\pgfqpoint{3.933237in}{1.187756in}}{\pgfqpoint{3.940640in}{1.179560in}}%
\pgfpathquadraticcurveto{\pgfqpoint{3.948043in}{1.171383in}}{\pgfqpoint{3.948043in}{1.157315in}}%
\pgfpathlineto{\pgfqpoint{3.948043in}{1.157315in}}%
\pgfpathclose%
\pgfpathmoveto{\pgfqpoint{3.937686in}{1.160359in}}%
\pgfpathquadraticcurveto{\pgfqpoint{3.937578in}{1.168843in}}{\pgfqpoint{3.932931in}{1.173904in}}%
\pgfpathquadraticcurveto{\pgfqpoint{3.928283in}{1.178984in}}{\pgfqpoint{3.920628in}{1.178984in}}%
\pgfpathquadraticcurveto{\pgfqpoint{3.911964in}{1.178984in}}{\pgfqpoint{3.906759in}{1.174084in}}%
\pgfpathquadraticcurveto{\pgfqpoint{3.901553in}{1.169185in}}{\pgfqpoint{3.900761in}{1.160287in}}%
\pgfpathlineto{\pgfqpoint{3.937686in}{1.160359in}}%
\pgfpathlineto{\pgfqpoint{3.937686in}{1.160359in}}%
\pgfpathclose%
\pgfpathmoveto{\pgfqpoint{4.017462in}{1.161260in}}%
\pgfpathlineto{\pgfqpoint{4.017462in}{1.123200in}}%
\pgfpathlineto{\pgfqpoint{4.007105in}{1.123200in}}%
\pgfpathlineto{\pgfqpoint{4.007105in}{1.160917in}}%
\pgfpathquadraticcurveto{\pgfqpoint{4.007105in}{1.169869in}}{\pgfqpoint{4.003610in}{1.174300in}}%
\pgfpathquadraticcurveto{\pgfqpoint{4.000116in}{1.178749in}}{\pgfqpoint{3.993145in}{1.178749in}}%
\pgfpathquadraticcurveto{\pgfqpoint{3.984752in}{1.178749in}}{\pgfqpoint{3.979906in}{1.173400in}}%
\pgfpathquadraticcurveto{\pgfqpoint{3.975061in}{1.168068in}}{\pgfqpoint{3.975061in}{1.158828in}}%
\pgfpathlineto{\pgfqpoint{3.975061in}{1.123200in}}%
\pgfpathlineto{\pgfqpoint{3.964650in}{1.123200in}}%
\pgfpathlineto{\pgfqpoint{3.964650in}{1.186243in}}%
\pgfpathlineto{\pgfqpoint{3.975061in}{1.186243in}}%
\pgfpathlineto{\pgfqpoint{3.975061in}{1.176444in}}%
\pgfpathquadraticcurveto{\pgfqpoint{3.978789in}{1.182136in}}{\pgfqpoint{3.983815in}{1.184946in}}%
\pgfpathquadraticcurveto{\pgfqpoint{3.988858in}{1.187756in}}{\pgfqpoint{3.995451in}{1.187756in}}%
\pgfpathquadraticcurveto{\pgfqpoint{4.006312in}{1.187756in}}{\pgfqpoint{4.011878in}{1.181037in}}%
\pgfpathquadraticcurveto{\pgfqpoint{4.017462in}{1.174318in}}{\pgfqpoint{4.017462in}{1.161260in}}%
\pgfpathlineto{\pgfqpoint{4.017462in}{1.161260in}}%
\pgfpathclose%
\pgfpathmoveto{\pgfqpoint{4.048352in}{1.204147in}}%
\pgfpathlineto{\pgfqpoint{4.048352in}{1.186243in}}%
\pgfpathlineto{\pgfqpoint{4.069679in}{1.186243in}}%
\pgfpathlineto{\pgfqpoint{4.069679in}{1.178191in}}%
\pgfpathlineto{\pgfqpoint{4.048352in}{1.178191in}}%
\pgfpathlineto{\pgfqpoint{4.048352in}{1.143968in}}%
\pgfpathquadraticcurveto{\pgfqpoint{4.048352in}{1.136259in}}{\pgfqpoint{4.050460in}{1.134061in}}%
\pgfpathquadraticcurveto{\pgfqpoint{4.052567in}{1.131864in}}{\pgfqpoint{4.059052in}{1.131864in}}%
\pgfpathlineto{\pgfqpoint{4.069679in}{1.131864in}}%
\pgfpathlineto{\pgfqpoint{4.069679in}{1.123200in}}%
\pgfpathlineto{\pgfqpoint{4.059052in}{1.123200in}}%
\pgfpathquadraticcurveto{\pgfqpoint{4.047056in}{1.123200in}}{\pgfqpoint{4.042498in}{1.127667in}}%
\pgfpathquadraticcurveto{\pgfqpoint{4.037941in}{1.132152in}}{\pgfqpoint{4.037941in}{1.143968in}}%
\pgfpathlineto{\pgfqpoint{4.037941in}{1.178191in}}%
\pgfpathlineto{\pgfqpoint{4.030340in}{1.178191in}}%
\pgfpathlineto{\pgfqpoint{4.030340in}{1.186243in}}%
\pgfpathlineto{\pgfqpoint{4.037941in}{1.186243in}}%
\pgfpathlineto{\pgfqpoint{4.037941in}{1.204147in}}%
\pgfpathlineto{\pgfqpoint{4.048352in}{1.204147in}}%
\pgfpathlineto{\pgfqpoint{4.048352in}{1.204147in}}%
\pgfpathclose%
\pgfpathmoveto{\pgfqpoint{4.078072in}{1.159386in}}%
\pgfpathlineto{\pgfqpoint{4.108405in}{1.159386in}}%
\pgfpathlineto{\pgfqpoint{4.108405in}{1.150164in}}%
\pgfpathlineto{\pgfqpoint{4.078072in}{1.150164in}}%
\pgfpathlineto{\pgfqpoint{4.078072in}{1.159386in}}%
\pgfpathlineto{\pgfqpoint{4.078072in}{1.159386in}}%
\pgfpathclose%
\pgfpathmoveto{\pgfqpoint{4.165071in}{1.184387in}}%
\pgfpathlineto{\pgfqpoint{4.165071in}{1.174589in}}%
\pgfpathquadraticcurveto{\pgfqpoint{4.160694in}{1.176840in}}{\pgfqpoint{4.155957in}{1.177957in}}%
\pgfpathquadraticcurveto{\pgfqpoint{4.151238in}{1.179092in}}{\pgfqpoint{4.146158in}{1.179092in}}%
\pgfpathquadraticcurveto{\pgfqpoint{4.138449in}{1.179092in}}{\pgfqpoint{4.134595in}{1.176732in}}%
\pgfpathquadraticcurveto{\pgfqpoint{4.130740in}{1.174373in}}{\pgfqpoint{4.130740in}{1.169635in}}%
\pgfpathquadraticcurveto{\pgfqpoint{4.130740in}{1.166033in}}{\pgfqpoint{4.133496in}{1.163980in}}%
\pgfpathquadraticcurveto{\pgfqpoint{4.136252in}{1.161926in}}{\pgfqpoint{4.144591in}{1.160071in}}%
\pgfpathlineto{\pgfqpoint{4.148140in}{1.159278in}}%
\pgfpathquadraticcurveto{\pgfqpoint{4.159163in}{1.156919in}}{\pgfqpoint{4.163810in}{1.152614in}}%
\pgfpathquadraticcurveto{\pgfqpoint{4.168457in}{1.148309in}}{\pgfqpoint{4.168457in}{1.140600in}}%
\pgfpathquadraticcurveto{\pgfqpoint{4.168457in}{1.131810in}}{\pgfqpoint{4.161505in}{1.126676in}}%
\pgfpathquadraticcurveto{\pgfqpoint{4.154552in}{1.121561in}}{\pgfqpoint{4.142394in}{1.121561in}}%
\pgfpathquadraticcurveto{\pgfqpoint{4.137332in}{1.121561in}}{\pgfqpoint{4.131839in}{1.122552in}}%
\pgfpathquadraticcurveto{\pgfqpoint{4.126345in}{1.123542in}}{\pgfqpoint{4.120275in}{1.125506in}}%
\pgfpathlineto{\pgfqpoint{4.120275in}{1.136205in}}%
\pgfpathquadraticcurveto{\pgfqpoint{4.126021in}{1.133215in}}{\pgfqpoint{4.131587in}{1.131720in}}%
\pgfpathquadraticcurveto{\pgfqpoint{4.137152in}{1.130243in}}{\pgfqpoint{4.142628in}{1.130243in}}%
\pgfpathquadraticcurveto{\pgfqpoint{4.149941in}{1.130243in}}{\pgfqpoint{4.153868in}{1.132746in}}%
\pgfpathquadraticcurveto{\pgfqpoint{4.157812in}{1.135250in}}{\pgfqpoint{4.157812in}{1.139807in}}%
\pgfpathquadraticcurveto{\pgfqpoint{4.157812in}{1.144022in}}{\pgfqpoint{4.154966in}{1.146274in}}%
\pgfpathquadraticcurveto{\pgfqpoint{4.152138in}{1.148525in}}{\pgfqpoint{4.142502in}{1.150614in}}%
\pgfpathlineto{\pgfqpoint{4.138900in}{1.151461in}}%
\pgfpathquadraticcurveto{\pgfqpoint{4.129281in}{1.153478in}}{\pgfqpoint{4.124994in}{1.157675in}}%
\pgfpathquadraticcurveto{\pgfqpoint{4.120725in}{1.161872in}}{\pgfqpoint{4.120725in}{1.169185in}}%
\pgfpathquadraticcurveto{\pgfqpoint{4.120725in}{1.178083in}}{\pgfqpoint{4.127030in}{1.182910in}}%
\pgfpathquadraticcurveto{\pgfqpoint{4.133334in}{1.187756in}}{\pgfqpoint{4.144934in}{1.187756in}}%
\pgfpathquadraticcurveto{\pgfqpoint{4.150661in}{1.187756in}}{\pgfqpoint{4.155723in}{1.186909in}}%
\pgfpathquadraticcurveto{\pgfqpoint{4.160802in}{1.186080in}}{\pgfqpoint{4.165071in}{1.184387in}}%
\pgfpathlineto{\pgfqpoint{4.165071in}{1.184387in}}%
\pgfpathclose%
\pgfpathmoveto{\pgfqpoint{4.195188in}{1.204147in}}%
\pgfpathlineto{\pgfqpoint{4.195188in}{1.186243in}}%
\pgfpathlineto{\pgfqpoint{4.216514in}{1.186243in}}%
\pgfpathlineto{\pgfqpoint{4.216514in}{1.178191in}}%
\pgfpathlineto{\pgfqpoint{4.195188in}{1.178191in}}%
\pgfpathlineto{\pgfqpoint{4.195188in}{1.143968in}}%
\pgfpathquadraticcurveto{\pgfqpoint{4.195188in}{1.136259in}}{\pgfqpoint{4.197295in}{1.134061in}}%
\pgfpathquadraticcurveto{\pgfqpoint{4.199402in}{1.131864in}}{\pgfqpoint{4.205887in}{1.131864in}}%
\pgfpathlineto{\pgfqpoint{4.216514in}{1.131864in}}%
\pgfpathlineto{\pgfqpoint{4.216514in}{1.123200in}}%
\pgfpathlineto{\pgfqpoint{4.205887in}{1.123200in}}%
\pgfpathquadraticcurveto{\pgfqpoint{4.193891in}{1.123200in}}{\pgfqpoint{4.189334in}{1.127667in}}%
\pgfpathquadraticcurveto{\pgfqpoint{4.184776in}{1.132152in}}{\pgfqpoint{4.184776in}{1.143968in}}%
\pgfpathlineto{\pgfqpoint{4.184776in}{1.178191in}}%
\pgfpathlineto{\pgfqpoint{4.177175in}{1.178191in}}%
\pgfpathlineto{\pgfqpoint{4.177175in}{1.186243in}}%
\pgfpathlineto{\pgfqpoint{4.184776in}{1.186243in}}%
\pgfpathlineto{\pgfqpoint{4.184776in}{1.204147in}}%
\pgfpathlineto{\pgfqpoint{4.195188in}{1.204147in}}%
\pgfpathlineto{\pgfqpoint{4.195188in}{1.204147in}}%
\pgfpathclose%
\pgfpathmoveto{\pgfqpoint{4.229068in}{1.148075in}}%
\pgfpathlineto{\pgfqpoint{4.229068in}{1.186243in}}%
\pgfpathlineto{\pgfqpoint{4.239425in}{1.186243in}}%
\pgfpathlineto{\pgfqpoint{4.239425in}{1.148471in}}%
\pgfpathquadraticcurveto{\pgfqpoint{4.239425in}{1.139519in}}{\pgfqpoint{4.242902in}{1.135034in}}%
\pgfpathquadraticcurveto{\pgfqpoint{4.246396in}{1.130567in}}{\pgfqpoint{4.253385in}{1.130567in}}%
\pgfpathquadraticcurveto{\pgfqpoint{4.261760in}{1.130567in}}{\pgfqpoint{4.266624in}{1.135917in}}%
\pgfpathquadraticcurveto{\pgfqpoint{4.271505in}{1.141266in}}{\pgfqpoint{4.271505in}{1.150506in}}%
\pgfpathlineto{\pgfqpoint{4.271505in}{1.186243in}}%
\pgfpathlineto{\pgfqpoint{4.281862in}{1.186243in}}%
\pgfpathlineto{\pgfqpoint{4.281862in}{1.123200in}}%
\pgfpathlineto{\pgfqpoint{4.271505in}{1.123200in}}%
\pgfpathlineto{\pgfqpoint{4.271505in}{1.132891in}}%
\pgfpathquadraticcurveto{\pgfqpoint{4.267740in}{1.127145in}}{\pgfqpoint{4.262751in}{1.124353in}}%
\pgfpathquadraticcurveto{\pgfqpoint{4.257780in}{1.121561in}}{\pgfqpoint{4.251187in}{1.121561in}}%
\pgfpathquadraticcurveto{\pgfqpoint{4.240326in}{1.121561in}}{\pgfqpoint{4.234688in}{1.128315in}}%
\pgfpathquadraticcurveto{\pgfqpoint{4.229068in}{1.135070in}}{\pgfqpoint{4.229068in}{1.148075in}}%
\pgfpathlineto{\pgfqpoint{4.229068in}{1.148075in}}%
\pgfpathclose%
\pgfpathmoveto{\pgfqpoint{4.255132in}{1.187756in}}%
\pgfpathlineto{\pgfqpoint{4.255132in}{1.187756in}}%
\pgfpathlineto{\pgfqpoint{4.255132in}{1.187756in}}%
\pgfpathclose%
\pgfpathmoveto{\pgfqpoint{4.344670in}{1.176678in}}%
\pgfpathlineto{\pgfqpoint{4.344670in}{1.210793in}}%
\pgfpathlineto{\pgfqpoint{4.355027in}{1.210793in}}%
\pgfpathlineto{\pgfqpoint{4.355027in}{1.123200in}}%
\pgfpathlineto{\pgfqpoint{4.344670in}{1.123200in}}%
\pgfpathlineto{\pgfqpoint{4.344670in}{1.132656in}}%
\pgfpathquadraticcurveto{\pgfqpoint{4.341410in}{1.127037in}}{\pgfqpoint{4.336421in}{1.124299in}}%
\pgfpathquadraticcurveto{\pgfqpoint{4.331449in}{1.121561in}}{\pgfqpoint{4.324461in}{1.121561in}}%
\pgfpathquadraticcurveto{\pgfqpoint{4.313041in}{1.121561in}}{\pgfqpoint{4.305854in}{1.130675in}}%
\pgfpathquadraticcurveto{\pgfqpoint{4.298685in}{1.139807in}}{\pgfqpoint{4.298685in}{1.154667in}}%
\pgfpathquadraticcurveto{\pgfqpoint{4.298685in}{1.169527in}}{\pgfqpoint{4.305854in}{1.178641in}}%
\pgfpathquadraticcurveto{\pgfqpoint{4.313041in}{1.187756in}}{\pgfqpoint{4.324461in}{1.187756in}}%
\pgfpathquadraticcurveto{\pgfqpoint{4.331449in}{1.187756in}}{\pgfqpoint{4.336421in}{1.185018in}}%
\pgfpathquadraticcurveto{\pgfqpoint{4.341410in}{1.182298in}}{\pgfqpoint{4.344670in}{1.176678in}}%
\pgfpathlineto{\pgfqpoint{4.344670in}{1.176678in}}%
\pgfpathclose%
\pgfpathmoveto{\pgfqpoint{4.309385in}{1.154667in}}%
\pgfpathquadraticcurveto{\pgfqpoint{4.309385in}{1.143248in}}{\pgfqpoint{4.314086in}{1.136745in}}%
\pgfpathquadraticcurveto{\pgfqpoint{4.318787in}{1.130243in}}{\pgfqpoint{4.327000in}{1.130243in}}%
\pgfpathquadraticcurveto{\pgfqpoint{4.335214in}{1.130243in}}{\pgfqpoint{4.339933in}{1.136745in}}%
\pgfpathquadraticcurveto{\pgfqpoint{4.344670in}{1.143248in}}{\pgfqpoint{4.344670in}{1.154667in}}%
\pgfpathquadraticcurveto{\pgfqpoint{4.344670in}{1.166087in}}{\pgfqpoint{4.339933in}{1.172589in}}%
\pgfpathquadraticcurveto{\pgfqpoint{4.335214in}{1.179092in}}{\pgfqpoint{4.327000in}{1.179092in}}%
\pgfpathquadraticcurveto{\pgfqpoint{4.318787in}{1.179092in}}{\pgfqpoint{4.314086in}{1.172589in}}%
\pgfpathquadraticcurveto{\pgfqpoint{4.309385in}{1.166087in}}{\pgfqpoint{4.309385in}{1.154667in}}%
\pgfpathlineto{\pgfqpoint{4.309385in}{1.154667in}}%
\pgfpathclose%
\pgfpathmoveto{\pgfqpoint{4.402597in}{1.117346in}}%
\pgfpathquadraticcurveto{\pgfqpoint{4.398220in}{1.106088in}}{\pgfqpoint{4.394042in}{1.102666in}}%
\pgfpathquadraticcurveto{\pgfqpoint{4.389881in}{1.099226in}}{\pgfqpoint{4.382910in}{1.099226in}}%
\pgfpathlineto{\pgfqpoint{4.374625in}{1.099226in}}%
\pgfpathlineto{\pgfqpoint{4.374625in}{1.107890in}}%
\pgfpathlineto{\pgfqpoint{4.380713in}{1.107890in}}%
\pgfpathquadraticcurveto{\pgfqpoint{4.384982in}{1.107890in}}{\pgfqpoint{4.387341in}{1.109925in}}%
\pgfpathquadraticcurveto{\pgfqpoint{4.389719in}{1.111942in}}{\pgfqpoint{4.392583in}{1.119489in}}%
\pgfpathlineto{\pgfqpoint{4.394438in}{1.124209in}}%
\pgfpathlineto{\pgfqpoint{4.368951in}{1.186243in}}%
\pgfpathlineto{\pgfqpoint{4.379920in}{1.186243in}}%
\pgfpathlineto{\pgfqpoint{4.399625in}{1.136943in}}%
\pgfpathlineto{\pgfqpoint{4.419331in}{1.186243in}}%
\pgfpathlineto{\pgfqpoint{4.430300in}{1.186243in}}%
\pgfpathlineto{\pgfqpoint{4.402597in}{1.117346in}}%
\pgfpathlineto{\pgfqpoint{4.402597in}{1.117346in}}%
\pgfpathclose%
\pgfusepath{fill}%
\end{pgfscope}%
\begin{pgfscope}%
\pgfsetbuttcap%
\pgfsetroundjoin%
\definecolor{currentfill}{rgb}{0.309804,0.372549,0.419608}%
\pgfsetfillcolor{currentfill}%
\pgfsetlinewidth{0.000000pt}%
\definecolor{currentstroke}{rgb}{0.309804,0.372549,0.419608}%
\pgfsetstrokecolor{currentstroke}%
\pgfsetdash{}{0pt}%
\pgfpathmoveto{\pgfqpoint{4.562628in}{1.154883in}}%
\pgfpathquadraticcurveto{\pgfqpoint{4.550074in}{1.154883in}}{\pgfqpoint{4.545228in}{1.152019in}}%
\pgfpathquadraticcurveto{\pgfqpoint{4.540383in}{1.149156in}}{\pgfqpoint{4.540383in}{1.142221in}}%
\pgfpathquadraticcurveto{\pgfqpoint{4.540383in}{1.136709in}}{\pgfqpoint{4.544021in}{1.133467in}}%
\pgfpathquadraticcurveto{\pgfqpoint{4.547660in}{1.130243in}}{\pgfqpoint{4.553892in}{1.130243in}}%
\pgfpathquadraticcurveto{\pgfqpoint{4.562520in}{1.130243in}}{\pgfqpoint{4.567725in}{1.136349in}}%
\pgfpathquadraticcurveto{\pgfqpoint{4.572931in}{1.142455in}}{\pgfqpoint{4.572931in}{1.152578in}}%
\pgfpathlineto{\pgfqpoint{4.572931in}{1.154883in}}%
\pgfpathlineto{\pgfqpoint{4.562628in}{1.154883in}}%
\pgfpathlineto{\pgfqpoint{4.562628in}{1.154883in}}%
\pgfpathclose%
\pgfpathmoveto{\pgfqpoint{4.583288in}{1.159170in}}%
\pgfpathlineto{\pgfqpoint{4.583288in}{1.123200in}}%
\pgfpathlineto{\pgfqpoint{4.572931in}{1.123200in}}%
\pgfpathlineto{\pgfqpoint{4.572931in}{1.132764in}}%
\pgfpathquadraticcurveto{\pgfqpoint{4.569383in}{1.127037in}}{\pgfqpoint{4.564087in}{1.124299in}}%
\pgfpathquadraticcurveto{\pgfqpoint{4.558791in}{1.121561in}}{\pgfqpoint{4.551136in}{1.121561in}}%
\pgfpathquadraticcurveto{\pgfqpoint{4.541464in}{1.121561in}}{\pgfqpoint{4.535736in}{1.127001in}}%
\pgfpathquadraticcurveto{\pgfqpoint{4.530026in}{1.132440in}}{\pgfqpoint{4.530026in}{1.141554in}}%
\pgfpathquadraticcurveto{\pgfqpoint{4.530026in}{1.152182in}}{\pgfqpoint{4.537141in}{1.157585in}}%
\pgfpathquadraticcurveto{\pgfqpoint{4.544274in}{1.162989in}}{\pgfqpoint{4.558395in}{1.162989in}}%
\pgfpathlineto{\pgfqpoint{4.572931in}{1.162989in}}%
\pgfpathlineto{\pgfqpoint{4.572931in}{1.164016in}}%
\pgfpathquadraticcurveto{\pgfqpoint{4.572931in}{1.171166in}}{\pgfqpoint{4.568230in}{1.175075in}}%
\pgfpathquadraticcurveto{\pgfqpoint{4.563529in}{1.178984in}}{\pgfqpoint{4.555027in}{1.178984in}}%
\pgfpathquadraticcurveto{\pgfqpoint{4.549623in}{1.178984in}}{\pgfqpoint{4.544490in}{1.177687in}}%
\pgfpathquadraticcurveto{\pgfqpoint{4.539374in}{1.176390in}}{\pgfqpoint{4.534655in}{1.173796in}}%
\pgfpathlineto{\pgfqpoint{4.534655in}{1.183379in}}%
\pgfpathquadraticcurveto{\pgfqpoint{4.540329in}{1.185576in}}{\pgfqpoint{4.545679in}{1.186657in}}%
\pgfpathquadraticcurveto{\pgfqpoint{4.551028in}{1.187756in}}{\pgfqpoint{4.556090in}{1.187756in}}%
\pgfpathquadraticcurveto{\pgfqpoint{4.569779in}{1.187756in}}{\pgfqpoint{4.576533in}{1.180659in}}%
\pgfpathquadraticcurveto{\pgfqpoint{4.583288in}{1.173580in}}{\pgfqpoint{4.583288in}{1.159170in}}%
\pgfpathlineto{\pgfqpoint{4.583288in}{1.159170in}}%
\pgfpathclose%
\pgfpathmoveto{\pgfqpoint{4.597193in}{1.186243in}}%
\pgfpathlineto{\pgfqpoint{4.608163in}{1.186243in}}%
\pgfpathlineto{\pgfqpoint{4.627868in}{1.133341in}}%
\pgfpathlineto{\pgfqpoint{4.647573in}{1.186243in}}%
\pgfpathlineto{\pgfqpoint{4.658543in}{1.186243in}}%
\pgfpathlineto{\pgfqpoint{4.634893in}{1.123200in}}%
\pgfpathlineto{\pgfqpoint{4.620825in}{1.123200in}}%
\pgfpathlineto{\pgfqpoint{4.597193in}{1.186243in}}%
\pgfpathlineto{\pgfqpoint{4.597193in}{1.186243in}}%
\pgfpathclose%
\pgfpathmoveto{\pgfqpoint{4.726773in}{1.157315in}}%
\pgfpathlineto{\pgfqpoint{4.726773in}{1.152254in}}%
\pgfpathlineto{\pgfqpoint{4.679149in}{1.152254in}}%
\pgfpathquadraticcurveto{\pgfqpoint{4.679833in}{1.141554in}}{\pgfqpoint{4.685597in}{1.135953in}}%
\pgfpathquadraticcurveto{\pgfqpoint{4.691361in}{1.130351in}}{\pgfqpoint{4.701664in}{1.130351in}}%
\pgfpathquadraticcurveto{\pgfqpoint{4.707626in}{1.130351in}}{\pgfqpoint{4.713228in}{1.131810in}}%
\pgfpathquadraticcurveto{\pgfqpoint{4.718829in}{1.133269in}}{\pgfqpoint{4.724359in}{1.136205in}}%
\pgfpathlineto{\pgfqpoint{4.724359in}{1.126406in}}%
\pgfpathquadraticcurveto{\pgfqpoint{4.718775in}{1.124047in}}{\pgfqpoint{4.712921in}{1.122804in}}%
\pgfpathquadraticcurveto{\pgfqpoint{4.707067in}{1.121561in}}{\pgfqpoint{4.701051in}{1.121561in}}%
\pgfpathquadraticcurveto{\pgfqpoint{4.685957in}{1.121561in}}{\pgfqpoint{4.677149in}{1.130333in}}%
\pgfpathquadraticcurveto{\pgfqpoint{4.668341in}{1.139123in}}{\pgfqpoint{4.668341in}{1.154109in}}%
\pgfpathquadraticcurveto{\pgfqpoint{4.668341in}{1.169581in}}{\pgfqpoint{4.676699in}{1.178659in}}%
\pgfpathquadraticcurveto{\pgfqpoint{4.685057in}{1.187756in}}{\pgfqpoint{4.699250in}{1.187756in}}%
\pgfpathquadraticcurveto{\pgfqpoint{4.711967in}{1.187756in}}{\pgfqpoint{4.719370in}{1.179560in}}%
\pgfpathquadraticcurveto{\pgfqpoint{4.726773in}{1.171383in}}{\pgfqpoint{4.726773in}{1.157315in}}%
\pgfpathlineto{\pgfqpoint{4.726773in}{1.157315in}}%
\pgfpathclose%
\pgfpathmoveto{\pgfqpoint{4.716416in}{1.160359in}}%
\pgfpathquadraticcurveto{\pgfqpoint{4.716308in}{1.168843in}}{\pgfqpoint{4.711661in}{1.173904in}}%
\pgfpathquadraticcurveto{\pgfqpoint{4.707013in}{1.178984in}}{\pgfqpoint{4.699358in}{1.178984in}}%
\pgfpathquadraticcurveto{\pgfqpoint{4.690694in}{1.178984in}}{\pgfqpoint{4.685489in}{1.174084in}}%
\pgfpathquadraticcurveto{\pgfqpoint{4.680283in}{1.169185in}}{\pgfqpoint{4.679491in}{1.160287in}}%
\pgfpathlineto{\pgfqpoint{4.716416in}{1.160359in}}%
\pgfpathlineto{\pgfqpoint{4.716416in}{1.160359in}}%
\pgfpathclose%
\pgfpathmoveto{\pgfqpoint{4.780305in}{1.176570in}}%
\pgfpathquadraticcurveto{\pgfqpoint{4.778558in}{1.177579in}}{\pgfqpoint{4.776504in}{1.178047in}}%
\pgfpathquadraticcurveto{\pgfqpoint{4.774451in}{1.178533in}}{\pgfqpoint{4.771983in}{1.178533in}}%
\pgfpathquadraticcurveto{\pgfqpoint{4.763193in}{1.178533in}}{\pgfqpoint{4.758492in}{1.172823in}}%
\pgfpathquadraticcurveto{\pgfqpoint{4.753791in}{1.167114in}}{\pgfqpoint{4.753791in}{1.156414in}}%
\pgfpathlineto{\pgfqpoint{4.753791in}{1.123200in}}%
\pgfpathlineto{\pgfqpoint{4.743380in}{1.123200in}}%
\pgfpathlineto{\pgfqpoint{4.743380in}{1.186243in}}%
\pgfpathlineto{\pgfqpoint{4.753791in}{1.186243in}}%
\pgfpathlineto{\pgfqpoint{4.753791in}{1.176444in}}%
\pgfpathquadraticcurveto{\pgfqpoint{4.757069in}{1.182190in}}{\pgfqpoint{4.762293in}{1.184964in}}%
\pgfpathquadraticcurveto{\pgfqpoint{4.767534in}{1.187756in}}{\pgfqpoint{4.775027in}{1.187756in}}%
\pgfpathquadraticcurveto{\pgfqpoint{4.776090in}{1.187756in}}{\pgfqpoint{4.777387in}{1.187611in}}%
\pgfpathquadraticcurveto{\pgfqpoint{4.778684in}{1.187485in}}{\pgfqpoint{4.780251in}{1.187197in}}%
\pgfpathlineto{\pgfqpoint{4.780305in}{1.176570in}}%
\pgfpathlineto{\pgfqpoint{4.780305in}{1.176570in}}%
\pgfpathclose%
\pgfpathmoveto{\pgfqpoint{4.819824in}{1.154883in}}%
\pgfpathquadraticcurveto{\pgfqpoint{4.807269in}{1.154883in}}{\pgfqpoint{4.802424in}{1.152019in}}%
\pgfpathquadraticcurveto{\pgfqpoint{4.797579in}{1.149156in}}{\pgfqpoint{4.797579in}{1.142221in}}%
\pgfpathquadraticcurveto{\pgfqpoint{4.797579in}{1.136709in}}{\pgfqpoint{4.801217in}{1.133467in}}%
\pgfpathquadraticcurveto{\pgfqpoint{4.804855in}{1.130243in}}{\pgfqpoint{4.811088in}{1.130243in}}%
\pgfpathquadraticcurveto{\pgfqpoint{4.819715in}{1.130243in}}{\pgfqpoint{4.824921in}{1.136349in}}%
\pgfpathquadraticcurveto{\pgfqpoint{4.830127in}{1.142455in}}{\pgfqpoint{4.830127in}{1.152578in}}%
\pgfpathlineto{\pgfqpoint{4.830127in}{1.154883in}}%
\pgfpathlineto{\pgfqpoint{4.819824in}{1.154883in}}%
\pgfpathlineto{\pgfqpoint{4.819824in}{1.154883in}}%
\pgfpathclose%
\pgfpathmoveto{\pgfqpoint{4.840484in}{1.159170in}}%
\pgfpathlineto{\pgfqpoint{4.840484in}{1.123200in}}%
\pgfpathlineto{\pgfqpoint{4.830127in}{1.123200in}}%
\pgfpathlineto{\pgfqpoint{4.830127in}{1.132764in}}%
\pgfpathquadraticcurveto{\pgfqpoint{4.826578in}{1.127037in}}{\pgfqpoint{4.821283in}{1.124299in}}%
\pgfpathquadraticcurveto{\pgfqpoint{4.815987in}{1.121561in}}{\pgfqpoint{4.808332in}{1.121561in}}%
\pgfpathquadraticcurveto{\pgfqpoint{4.798659in}{1.121561in}}{\pgfqpoint{4.792931in}{1.127001in}}%
\pgfpathquadraticcurveto{\pgfqpoint{4.787222in}{1.132440in}}{\pgfqpoint{4.787222in}{1.141554in}}%
\pgfpathquadraticcurveto{\pgfqpoint{4.787222in}{1.152182in}}{\pgfqpoint{4.794336in}{1.157585in}}%
\pgfpathquadraticcurveto{\pgfqpoint{4.801469in}{1.162989in}}{\pgfqpoint{4.815591in}{1.162989in}}%
\pgfpathlineto{\pgfqpoint{4.830127in}{1.162989in}}%
\pgfpathlineto{\pgfqpoint{4.830127in}{1.164016in}}%
\pgfpathquadraticcurveto{\pgfqpoint{4.830127in}{1.171166in}}{\pgfqpoint{4.825425in}{1.175075in}}%
\pgfpathquadraticcurveto{\pgfqpoint{4.820724in}{1.178984in}}{\pgfqpoint{4.812222in}{1.178984in}}%
\pgfpathquadraticcurveto{\pgfqpoint{4.806819in}{1.178984in}}{\pgfqpoint{4.801685in}{1.177687in}}%
\pgfpathquadraticcurveto{\pgfqpoint{4.796570in}{1.176390in}}{\pgfqpoint{4.791851in}{1.173796in}}%
\pgfpathlineto{\pgfqpoint{4.791851in}{1.183379in}}%
\pgfpathquadraticcurveto{\pgfqpoint{4.797525in}{1.185576in}}{\pgfqpoint{4.802874in}{1.186657in}}%
\pgfpathquadraticcurveto{\pgfqpoint{4.808224in}{1.187756in}}{\pgfqpoint{4.813285in}{1.187756in}}%
\pgfpathquadraticcurveto{\pgfqpoint{4.826974in}{1.187756in}}{\pgfqpoint{4.833729in}{1.180659in}}%
\pgfpathquadraticcurveto{\pgfqpoint{4.840484in}{1.173580in}}{\pgfqpoint{4.840484in}{1.159170in}}%
\pgfpathlineto{\pgfqpoint{4.840484in}{1.159170in}}%
\pgfpathclose%
\pgfpathmoveto{\pgfqpoint{4.903292in}{1.155460in}}%
\pgfpathquadraticcurveto{\pgfqpoint{4.903292in}{1.166717in}}{\pgfqpoint{4.898645in}{1.172896in}}%
\pgfpathquadraticcurveto{\pgfqpoint{4.894016in}{1.179092in}}{\pgfqpoint{4.885622in}{1.179092in}}%
\pgfpathquadraticcurveto{\pgfqpoint{4.877300in}{1.179092in}}{\pgfqpoint{4.872653in}{1.172896in}}%
\pgfpathquadraticcurveto{\pgfqpoint{4.868006in}{1.166717in}}{\pgfqpoint{4.868006in}{1.155460in}}%
\pgfpathquadraticcurveto{\pgfqpoint{4.868006in}{1.144256in}}{\pgfqpoint{4.872653in}{1.138060in}}%
\pgfpathquadraticcurveto{\pgfqpoint{4.877300in}{1.131864in}}{\pgfqpoint{4.885622in}{1.131864in}}%
\pgfpathquadraticcurveto{\pgfqpoint{4.894016in}{1.131864in}}{\pgfqpoint{4.898645in}{1.138060in}}%
\pgfpathquadraticcurveto{\pgfqpoint{4.903292in}{1.144256in}}{\pgfqpoint{4.903292in}{1.155460in}}%
\pgfpathlineto{\pgfqpoint{4.903292in}{1.155460in}}%
\pgfpathclose%
\pgfpathmoveto{\pgfqpoint{4.913649in}{1.131017in}}%
\pgfpathquadraticcurveto{\pgfqpoint{4.913649in}{1.114932in}}{\pgfqpoint{4.906498in}{1.107079in}}%
\pgfpathquadraticcurveto{\pgfqpoint{4.899365in}{1.099226in}}{\pgfqpoint{4.884613in}{1.099226in}}%
\pgfpathquadraticcurveto{\pgfqpoint{4.879156in}{1.099226in}}{\pgfqpoint{4.874310in}{1.100036in}}%
\pgfpathquadraticcurveto{\pgfqpoint{4.869465in}{1.100847in}}{\pgfqpoint{4.864908in}{1.102540in}}%
\pgfpathlineto{\pgfqpoint{4.864908in}{1.112609in}}%
\pgfpathquadraticcurveto{\pgfqpoint{4.869465in}{1.110141in}}{\pgfqpoint{4.873914in}{1.108970in}}%
\pgfpathquadraticcurveto{\pgfqpoint{4.878363in}{1.107782in}}{\pgfqpoint{4.882974in}{1.107782in}}%
\pgfpathquadraticcurveto{\pgfqpoint{4.893169in}{1.107782in}}{\pgfqpoint{4.898230in}{1.113095in}}%
\pgfpathquadraticcurveto{\pgfqpoint{4.903292in}{1.118409in}}{\pgfqpoint{4.903292in}{1.129162in}}%
\pgfpathlineto{\pgfqpoint{4.903292in}{1.134295in}}%
\pgfpathquadraticcurveto{\pgfqpoint{4.900086in}{1.128712in}}{\pgfqpoint{4.895078in}{1.125956in}}%
\pgfpathquadraticcurveto{\pgfqpoint{4.890071in}{1.123200in}}{\pgfqpoint{4.883082in}{1.123200in}}%
\pgfpathquadraticcurveto{\pgfqpoint{4.871500in}{1.123200in}}{\pgfqpoint{4.864404in}{1.132026in}}%
\pgfpathquadraticcurveto{\pgfqpoint{4.857307in}{1.140870in}}{\pgfqpoint{4.857307in}{1.155460in}}%
\pgfpathquadraticcurveto{\pgfqpoint{4.857307in}{1.170086in}}{\pgfqpoint{4.864404in}{1.178912in}}%
\pgfpathquadraticcurveto{\pgfqpoint{4.871500in}{1.187756in}}{\pgfqpoint{4.883082in}{1.187756in}}%
\pgfpathquadraticcurveto{\pgfqpoint{4.890071in}{1.187756in}}{\pgfqpoint{4.895078in}{1.185000in}}%
\pgfpathquadraticcurveto{\pgfqpoint{4.900086in}{1.182244in}}{\pgfqpoint{4.903292in}{1.176678in}}%
\pgfpathlineto{\pgfqpoint{4.903292in}{1.186243in}}%
\pgfpathlineto{\pgfqpoint{4.913649in}{1.186243in}}%
\pgfpathlineto{\pgfqpoint{4.913649in}{1.131017in}}%
\pgfpathlineto{\pgfqpoint{4.913649in}{1.131017in}}%
\pgfpathclose%
\pgfpathmoveto{\pgfqpoint{4.988922in}{1.157315in}}%
\pgfpathlineto{\pgfqpoint{4.988922in}{1.152254in}}%
\pgfpathlineto{\pgfqpoint{4.941298in}{1.152254in}}%
\pgfpathquadraticcurveto{\pgfqpoint{4.941982in}{1.141554in}}{\pgfqpoint{4.947746in}{1.135953in}}%
\pgfpathquadraticcurveto{\pgfqpoint{4.953510in}{1.130351in}}{\pgfqpoint{4.963813in}{1.130351in}}%
\pgfpathquadraticcurveto{\pgfqpoint{4.969775in}{1.130351in}}{\pgfqpoint{4.975377in}{1.131810in}}%
\pgfpathquadraticcurveto{\pgfqpoint{4.980978in}{1.133269in}}{\pgfqpoint{4.986508in}{1.136205in}}%
\pgfpathlineto{\pgfqpoint{4.986508in}{1.126406in}}%
\pgfpathquadraticcurveto{\pgfqpoint{4.980924in}{1.124047in}}{\pgfqpoint{4.975070in}{1.122804in}}%
\pgfpathquadraticcurveto{\pgfqpoint{4.969216in}{1.121561in}}{\pgfqpoint{4.963200in}{1.121561in}}%
\pgfpathquadraticcurveto{\pgfqpoint{4.948106in}{1.121561in}}{\pgfqpoint{4.939298in}{1.130333in}}%
\pgfpathquadraticcurveto{\pgfqpoint{4.930490in}{1.139123in}}{\pgfqpoint{4.930490in}{1.154109in}}%
\pgfpathquadraticcurveto{\pgfqpoint{4.930490in}{1.169581in}}{\pgfqpoint{4.938848in}{1.178659in}}%
\pgfpathquadraticcurveto{\pgfqpoint{4.947206in}{1.187756in}}{\pgfqpoint{4.961399in}{1.187756in}}%
\pgfpathquadraticcurveto{\pgfqpoint{4.974116in}{1.187756in}}{\pgfqpoint{4.981519in}{1.179560in}}%
\pgfpathquadraticcurveto{\pgfqpoint{4.988922in}{1.171383in}}{\pgfqpoint{4.988922in}{1.157315in}}%
\pgfpathlineto{\pgfqpoint{4.988922in}{1.157315in}}%
\pgfpathclose%
\pgfpathmoveto{\pgfqpoint{4.978565in}{1.160359in}}%
\pgfpathquadraticcurveto{\pgfqpoint{4.978457in}{1.168843in}}{\pgfqpoint{4.973809in}{1.173904in}}%
\pgfpathquadraticcurveto{\pgfqpoint{4.969162in}{1.178984in}}{\pgfqpoint{4.961507in}{1.178984in}}%
\pgfpathquadraticcurveto{\pgfqpoint{4.952843in}{1.178984in}}{\pgfqpoint{4.947638in}{1.174084in}}%
\pgfpathquadraticcurveto{\pgfqpoint{4.942432in}{1.169185in}}{\pgfqpoint{4.941640in}{1.160287in}}%
\pgfpathlineto{\pgfqpoint{4.978565in}{1.160359in}}%
\pgfpathlineto{\pgfqpoint{4.978565in}{1.160359in}}%
\pgfpathclose%
\pgfpathmoveto{\pgfqpoint{5.046110in}{1.184387in}}%
\pgfpathlineto{\pgfqpoint{5.046110in}{1.174589in}}%
\pgfpathquadraticcurveto{\pgfqpoint{5.041733in}{1.176840in}}{\pgfqpoint{5.036996in}{1.177957in}}%
\pgfpathquadraticcurveto{\pgfqpoint{5.032277in}{1.179092in}}{\pgfqpoint{5.027197in}{1.179092in}}%
\pgfpathquadraticcurveto{\pgfqpoint{5.019488in}{1.179092in}}{\pgfqpoint{5.015634in}{1.176732in}}%
\pgfpathquadraticcurveto{\pgfqpoint{5.011779in}{1.174373in}}{\pgfqpoint{5.011779in}{1.169635in}}%
\pgfpathquadraticcurveto{\pgfqpoint{5.011779in}{1.166033in}}{\pgfqpoint{5.014535in}{1.163980in}}%
\pgfpathquadraticcurveto{\pgfqpoint{5.017291in}{1.161926in}}{\pgfqpoint{5.025630in}{1.160071in}}%
\pgfpathlineto{\pgfqpoint{5.029179in}{1.159278in}}%
\pgfpathquadraticcurveto{\pgfqpoint{5.040202in}{1.156919in}}{\pgfqpoint{5.044849in}{1.152614in}}%
\pgfpathquadraticcurveto{\pgfqpoint{5.049497in}{1.148309in}}{\pgfqpoint{5.049497in}{1.140600in}}%
\pgfpathquadraticcurveto{\pgfqpoint{5.049497in}{1.131810in}}{\pgfqpoint{5.042544in}{1.126676in}}%
\pgfpathquadraticcurveto{\pgfqpoint{5.035591in}{1.121561in}}{\pgfqpoint{5.023433in}{1.121561in}}%
\pgfpathquadraticcurveto{\pgfqpoint{5.018372in}{1.121561in}}{\pgfqpoint{5.012878in}{1.122552in}}%
\pgfpathquadraticcurveto{\pgfqpoint{5.007384in}{1.123542in}}{\pgfqpoint{5.001314in}{1.125506in}}%
\pgfpathlineto{\pgfqpoint{5.001314in}{1.136205in}}%
\pgfpathquadraticcurveto{\pgfqpoint{5.007060in}{1.133215in}}{\pgfqpoint{5.012626in}{1.131720in}}%
\pgfpathquadraticcurveto{\pgfqpoint{5.018191in}{1.130243in}}{\pgfqpoint{5.023667in}{1.130243in}}%
\pgfpathquadraticcurveto{\pgfqpoint{5.030980in}{1.130243in}}{\pgfqpoint{5.034907in}{1.132746in}}%
\pgfpathquadraticcurveto{\pgfqpoint{5.038851in}{1.135250in}}{\pgfqpoint{5.038851in}{1.139807in}}%
\pgfpathquadraticcurveto{\pgfqpoint{5.038851in}{1.144022in}}{\pgfqpoint{5.036005in}{1.146274in}}%
\pgfpathquadraticcurveto{\pgfqpoint{5.033178in}{1.148525in}}{\pgfqpoint{5.023541in}{1.150614in}}%
\pgfpathlineto{\pgfqpoint{5.019939in}{1.151461in}}%
\pgfpathquadraticcurveto{\pgfqpoint{5.010320in}{1.153478in}}{\pgfqpoint{5.006033in}{1.157675in}}%
\pgfpathquadraticcurveto{\pgfqpoint{5.001764in}{1.161872in}}{\pgfqpoint{5.001764in}{1.169185in}}%
\pgfpathquadraticcurveto{\pgfqpoint{5.001764in}{1.178083in}}{\pgfqpoint{5.008069in}{1.182910in}}%
\pgfpathquadraticcurveto{\pgfqpoint{5.014373in}{1.187756in}}{\pgfqpoint{5.025973in}{1.187756in}}%
\pgfpathquadraticcurveto{\pgfqpoint{5.031701in}{1.187756in}}{\pgfqpoint{5.036762in}{1.186909in}}%
\pgfpathquadraticcurveto{\pgfqpoint{5.041841in}{1.186080in}}{\pgfqpoint{5.046110in}{1.184387in}}%
\pgfpathlineto{\pgfqpoint{5.046110in}{1.184387in}}%
\pgfpathclose%
\pgfpathmoveto{\pgfqpoint{5.068625in}{1.137502in}}%
\pgfpathlineto{\pgfqpoint{5.080495in}{1.137502in}}%
\pgfpathlineto{\pgfqpoint{5.080495in}{1.127811in}}%
\pgfpathlineto{\pgfqpoint{5.071273in}{1.109799in}}%
\pgfpathlineto{\pgfqpoint{5.064014in}{1.109799in}}%
\pgfpathlineto{\pgfqpoint{5.068625in}{1.127811in}}%
\pgfpathlineto{\pgfqpoint{5.068625in}{1.137502in}}%
\pgfpathlineto{\pgfqpoint{5.068625in}{1.137502in}}%
\pgfpathclose%
\pgfusepath{fill}%
\end{pgfscope}%
\begin{pgfscope}%
\pgfsetbuttcap%
\pgfsetroundjoin%
\definecolor{currentfill}{rgb}{0.309804,0.372549,0.419608}%
\pgfsetfillcolor{currentfill}%
\pgfsetlinewidth{0.000000pt}%
\definecolor{currentstroke}{rgb}{0.309804,0.372549,0.419608}%
\pgfsetstrokecolor{currentstroke}%
\pgfsetdash{}{0pt}%
\pgfpathmoveto{\pgfqpoint{5.166973in}{1.186243in}}%
\pgfpathlineto{\pgfqpoint{5.177330in}{1.186243in}}%
\pgfpathlineto{\pgfqpoint{5.190281in}{1.137051in}}%
\pgfpathlineto{\pgfqpoint{5.203159in}{1.186243in}}%
\pgfpathlineto{\pgfqpoint{5.215372in}{1.186243in}}%
\pgfpathlineto{\pgfqpoint{5.228322in}{1.137051in}}%
\pgfpathlineto{\pgfqpoint{5.241219in}{1.186243in}}%
\pgfpathlineto{\pgfqpoint{5.251576in}{1.186243in}}%
\pgfpathlineto{\pgfqpoint{5.235077in}{1.123200in}}%
\pgfpathlineto{\pgfqpoint{5.222865in}{1.123200in}}%
\pgfpathlineto{\pgfqpoint{5.209301in}{1.174877in}}%
\pgfpathlineto{\pgfqpoint{5.195684in}{1.123200in}}%
\pgfpathlineto{\pgfqpoint{5.183454in}{1.123200in}}%
\pgfpathlineto{\pgfqpoint{5.166973in}{1.186243in}}%
\pgfpathlineto{\pgfqpoint{5.166973in}{1.186243in}}%
\pgfpathclose%
\pgfpathmoveto{\pgfqpoint{5.267265in}{1.186243in}}%
\pgfpathlineto{\pgfqpoint{5.277622in}{1.186243in}}%
\pgfpathlineto{\pgfqpoint{5.277622in}{1.123200in}}%
\pgfpathlineto{\pgfqpoint{5.267265in}{1.123200in}}%
\pgfpathlineto{\pgfqpoint{5.267265in}{1.186243in}}%
\pgfpathlineto{\pgfqpoint{5.267265in}{1.186243in}}%
\pgfpathclose%
\pgfpathmoveto{\pgfqpoint{5.267265in}{1.210793in}}%
\pgfpathlineto{\pgfqpoint{5.277622in}{1.210793in}}%
\pgfpathlineto{\pgfqpoint{5.277622in}{1.197662in}}%
\pgfpathlineto{\pgfqpoint{5.267265in}{1.197662in}}%
\pgfpathlineto{\pgfqpoint{5.267265in}{1.210793in}}%
\pgfpathlineto{\pgfqpoint{5.267265in}{1.210793in}}%
\pgfpathclose%
\pgfpathmoveto{\pgfqpoint{5.309539in}{1.204147in}}%
\pgfpathlineto{\pgfqpoint{5.309539in}{1.186243in}}%
\pgfpathlineto{\pgfqpoint{5.330865in}{1.186243in}}%
\pgfpathlineto{\pgfqpoint{5.330865in}{1.178191in}}%
\pgfpathlineto{\pgfqpoint{5.309539in}{1.178191in}}%
\pgfpathlineto{\pgfqpoint{5.309539in}{1.143968in}}%
\pgfpathquadraticcurveto{\pgfqpoint{5.309539in}{1.136259in}}{\pgfqpoint{5.311646in}{1.134061in}}%
\pgfpathquadraticcurveto{\pgfqpoint{5.313754in}{1.131864in}}{\pgfqpoint{5.320238in}{1.131864in}}%
\pgfpathlineto{\pgfqpoint{5.330865in}{1.131864in}}%
\pgfpathlineto{\pgfqpoint{5.330865in}{1.123200in}}%
\pgfpathlineto{\pgfqpoint{5.320238in}{1.123200in}}%
\pgfpathquadraticcurveto{\pgfqpoint{5.308242in}{1.123200in}}{\pgfqpoint{5.303685in}{1.127667in}}%
\pgfpathquadraticcurveto{\pgfqpoint{5.299128in}{1.132152in}}{\pgfqpoint{5.299128in}{1.143968in}}%
\pgfpathlineto{\pgfqpoint{5.299128in}{1.178191in}}%
\pgfpathlineto{\pgfqpoint{5.291527in}{1.178191in}}%
\pgfpathlineto{\pgfqpoint{5.291527in}{1.186243in}}%
\pgfpathlineto{\pgfqpoint{5.299128in}{1.186243in}}%
\pgfpathlineto{\pgfqpoint{5.299128in}{1.204147in}}%
\pgfpathlineto{\pgfqpoint{5.309539in}{1.204147in}}%
\pgfpathlineto{\pgfqpoint{5.309539in}{1.204147in}}%
\pgfpathclose%
\pgfpathmoveto{\pgfqpoint{5.396898in}{1.161260in}}%
\pgfpathlineto{\pgfqpoint{5.396898in}{1.123200in}}%
\pgfpathlineto{\pgfqpoint{5.386541in}{1.123200in}}%
\pgfpathlineto{\pgfqpoint{5.386541in}{1.160917in}}%
\pgfpathquadraticcurveto{\pgfqpoint{5.386541in}{1.169869in}}{\pgfqpoint{5.383047in}{1.174300in}}%
\pgfpathquadraticcurveto{\pgfqpoint{5.379552in}{1.178749in}}{\pgfqpoint{5.372582in}{1.178749in}}%
\pgfpathquadraticcurveto{\pgfqpoint{5.364188in}{1.178749in}}{\pgfqpoint{5.359343in}{1.173400in}}%
\pgfpathquadraticcurveto{\pgfqpoint{5.354497in}{1.168068in}}{\pgfqpoint{5.354497in}{1.158828in}}%
\pgfpathlineto{\pgfqpoint{5.354497in}{1.123200in}}%
\pgfpathlineto{\pgfqpoint{5.344086in}{1.123200in}}%
\pgfpathlineto{\pgfqpoint{5.344086in}{1.210793in}}%
\pgfpathlineto{\pgfqpoint{5.354497in}{1.210793in}}%
\pgfpathlineto{\pgfqpoint{5.354497in}{1.176444in}}%
\pgfpathquadraticcurveto{\pgfqpoint{5.358226in}{1.182136in}}{\pgfqpoint{5.363251in}{1.184946in}}%
\pgfpathquadraticcurveto{\pgfqpoint{5.368295in}{1.187756in}}{\pgfqpoint{5.374887in}{1.187756in}}%
\pgfpathquadraticcurveto{\pgfqpoint{5.385748in}{1.187756in}}{\pgfqpoint{5.391314in}{1.181037in}}%
\pgfpathquadraticcurveto{\pgfqpoint{5.396898in}{1.174318in}}{\pgfqpoint{5.396898in}{1.161260in}}%
\pgfpathlineto{\pgfqpoint{5.396898in}{1.161260in}}%
\pgfpathclose%
\pgfusepath{fill}%
\end{pgfscope}%
\begin{pgfscope}%
\pgfsetbuttcap%
\pgfsetroundjoin%
\definecolor{currentfill}{rgb}{0.309804,0.372549,0.419608}%
\pgfsetfillcolor{currentfill}%
\pgfsetlinewidth{0.000000pt}%
\definecolor{currentstroke}{rgb}{0.309804,0.372549,0.419608}%
\pgfsetstrokecolor{currentstroke}%
\pgfsetdash{}{0pt}%
\pgfpathmoveto{\pgfqpoint{5.502022in}{1.132656in}}%
\pgfpathlineto{\pgfqpoint{5.502022in}{1.099226in}}%
\pgfpathlineto{\pgfqpoint{5.491611in}{1.099226in}}%
\pgfpathlineto{\pgfqpoint{5.491611in}{1.186243in}}%
\pgfpathlineto{\pgfqpoint{5.502022in}{1.186243in}}%
\pgfpathlineto{\pgfqpoint{5.502022in}{1.176678in}}%
\pgfpathquadraticcurveto{\pgfqpoint{5.505300in}{1.182298in}}{\pgfqpoint{5.510271in}{1.185018in}}%
\pgfpathquadraticcurveto{\pgfqpoint{5.515261in}{1.187756in}}{\pgfqpoint{5.522178in}{1.187756in}}%
\pgfpathquadraticcurveto{\pgfqpoint{5.533669in}{1.187756in}}{\pgfqpoint{5.540838in}{1.178641in}}%
\pgfpathquadraticcurveto{\pgfqpoint{5.548025in}{1.169527in}}{\pgfqpoint{5.548025in}{1.154667in}}%
\pgfpathquadraticcurveto{\pgfqpoint{5.548025in}{1.139807in}}{\pgfqpoint{5.540838in}{1.130675in}}%
\pgfpathquadraticcurveto{\pgfqpoint{5.533669in}{1.121561in}}{\pgfqpoint{5.522178in}{1.121561in}}%
\pgfpathquadraticcurveto{\pgfqpoint{5.515261in}{1.121561in}}{\pgfqpoint{5.510271in}{1.124299in}}%
\pgfpathquadraticcurveto{\pgfqpoint{5.505300in}{1.127037in}}{\pgfqpoint{5.502022in}{1.132656in}}%
\pgfpathlineto{\pgfqpoint{5.502022in}{1.132656in}}%
\pgfpathclose%
\pgfpathmoveto{\pgfqpoint{5.537272in}{1.154667in}}%
\pgfpathquadraticcurveto{\pgfqpoint{5.537272in}{1.166087in}}{\pgfqpoint{5.532571in}{1.172589in}}%
\pgfpathquadraticcurveto{\pgfqpoint{5.527869in}{1.179092in}}{\pgfqpoint{5.519656in}{1.179092in}}%
\pgfpathquadraticcurveto{\pgfqpoint{5.511424in}{1.179092in}}{\pgfqpoint{5.506723in}{1.172589in}}%
\pgfpathquadraticcurveto{\pgfqpoint{5.502022in}{1.166087in}}{\pgfqpoint{5.502022in}{1.154667in}}%
\pgfpathquadraticcurveto{\pgfqpoint{5.502022in}{1.143248in}}{\pgfqpoint{5.506723in}{1.136745in}}%
\pgfpathquadraticcurveto{\pgfqpoint{5.511424in}{1.130243in}}{\pgfqpoint{5.519656in}{1.130243in}}%
\pgfpathquadraticcurveto{\pgfqpoint{5.527869in}{1.130243in}}{\pgfqpoint{5.532571in}{1.136745in}}%
\pgfpathquadraticcurveto{\pgfqpoint{5.537272in}{1.143248in}}{\pgfqpoint{5.537272in}{1.154667in}}%
\pgfpathlineto{\pgfqpoint{5.537272in}{1.154667in}}%
\pgfpathclose%
\pgfpathmoveto{\pgfqpoint{5.601719in}{1.176570in}}%
\pgfpathquadraticcurveto{\pgfqpoint{5.599972in}{1.177579in}}{\pgfqpoint{5.597919in}{1.178047in}}%
\pgfpathquadraticcurveto{\pgfqpoint{5.595865in}{1.178533in}}{\pgfqpoint{5.593398in}{1.178533in}}%
\pgfpathquadraticcurveto{\pgfqpoint{5.584608in}{1.178533in}}{\pgfqpoint{5.579906in}{1.172823in}}%
\pgfpathquadraticcurveto{\pgfqpoint{5.575205in}{1.167114in}}{\pgfqpoint{5.575205in}{1.156414in}}%
\pgfpathlineto{\pgfqpoint{5.575205in}{1.123200in}}%
\pgfpathlineto{\pgfqpoint{5.564794in}{1.123200in}}%
\pgfpathlineto{\pgfqpoint{5.564794in}{1.186243in}}%
\pgfpathlineto{\pgfqpoint{5.575205in}{1.186243in}}%
\pgfpathlineto{\pgfqpoint{5.575205in}{1.176444in}}%
\pgfpathquadraticcurveto{\pgfqpoint{5.578484in}{1.182190in}}{\pgfqpoint{5.583707in}{1.184964in}}%
\pgfpathquadraticcurveto{\pgfqpoint{5.588949in}{1.187756in}}{\pgfqpoint{5.596442in}{1.187756in}}%
\pgfpathquadraticcurveto{\pgfqpoint{5.597504in}{1.187756in}}{\pgfqpoint{5.598801in}{1.187611in}}%
\pgfpathquadraticcurveto{\pgfqpoint{5.600098in}{1.187485in}}{\pgfqpoint{5.601665in}{1.187197in}}%
\pgfpathlineto{\pgfqpoint{5.601719in}{1.176570in}}%
\pgfpathlineto{\pgfqpoint{5.601719in}{1.176570in}}%
\pgfpathclose%
\pgfpathmoveto{\pgfqpoint{5.663969in}{1.157315in}}%
\pgfpathlineto{\pgfqpoint{5.663969in}{1.152254in}}%
\pgfpathlineto{\pgfqpoint{5.616345in}{1.152254in}}%
\pgfpathquadraticcurveto{\pgfqpoint{5.617030in}{1.141554in}}{\pgfqpoint{5.622793in}{1.135953in}}%
\pgfpathquadraticcurveto{\pgfqpoint{5.628557in}{1.130351in}}{\pgfqpoint{5.638860in}{1.130351in}}%
\pgfpathquadraticcurveto{\pgfqpoint{5.644822in}{1.130351in}}{\pgfqpoint{5.650424in}{1.131810in}}%
\pgfpathquadraticcurveto{\pgfqpoint{5.656026in}{1.133269in}}{\pgfqpoint{5.661556in}{1.136205in}}%
\pgfpathlineto{\pgfqpoint{5.661556in}{1.126406in}}%
\pgfpathquadraticcurveto{\pgfqpoint{5.655972in}{1.124047in}}{\pgfqpoint{5.650118in}{1.122804in}}%
\pgfpathquadraticcurveto{\pgfqpoint{5.644264in}{1.121561in}}{\pgfqpoint{5.638248in}{1.121561in}}%
\pgfpathquadraticcurveto{\pgfqpoint{5.623154in}{1.121561in}}{\pgfqpoint{5.614346in}{1.130333in}}%
\pgfpathquadraticcurveto{\pgfqpoint{5.605538in}{1.139123in}}{\pgfqpoint{5.605538in}{1.154109in}}%
\pgfpathquadraticcurveto{\pgfqpoint{5.605538in}{1.169581in}}{\pgfqpoint{5.613895in}{1.178659in}}%
\pgfpathquadraticcurveto{\pgfqpoint{5.622253in}{1.187756in}}{\pgfqpoint{5.636447in}{1.187756in}}%
\pgfpathquadraticcurveto{\pgfqpoint{5.649163in}{1.187756in}}{\pgfqpoint{5.656566in}{1.179560in}}%
\pgfpathquadraticcurveto{\pgfqpoint{5.663969in}{1.171383in}}{\pgfqpoint{5.663969in}{1.157315in}}%
\pgfpathlineto{\pgfqpoint{5.663969in}{1.157315in}}%
\pgfpathclose%
\pgfpathmoveto{\pgfqpoint{5.653612in}{1.160359in}}%
\pgfpathquadraticcurveto{\pgfqpoint{5.653504in}{1.168843in}}{\pgfqpoint{5.648857in}{1.173904in}}%
\pgfpathquadraticcurveto{\pgfqpoint{5.644210in}{1.178984in}}{\pgfqpoint{5.636555in}{1.178984in}}%
\pgfpathquadraticcurveto{\pgfqpoint{5.627891in}{1.178984in}}{\pgfqpoint{5.622685in}{1.174084in}}%
\pgfpathquadraticcurveto{\pgfqpoint{5.617480in}{1.169185in}}{\pgfqpoint{5.616687in}{1.160287in}}%
\pgfpathlineto{\pgfqpoint{5.653612in}{1.160359in}}%
\pgfpathlineto{\pgfqpoint{5.653612in}{1.160359in}}%
\pgfpathclose%
\pgfusepath{fill}%
\end{pgfscope}%
\begin{pgfscope}%
\pgfsetbuttcap%
\pgfsetroundjoin%
\definecolor{currentfill}{rgb}{0.309804,0.372549,0.419608}%
\pgfsetfillcolor{currentfill}%
\pgfsetlinewidth{0.000000pt}%
\definecolor{currentstroke}{rgb}{0.309804,0.372549,0.419608}%
\pgfsetstrokecolor{currentstroke}%
\pgfsetdash{}{0pt}%
\pgfpathmoveto{\pgfqpoint{5.777183in}{1.154883in}}%
\pgfpathquadraticcurveto{\pgfqpoint{5.764628in}{1.154883in}}{\pgfqpoint{5.759783in}{1.152019in}}%
\pgfpathquadraticcurveto{\pgfqpoint{5.754938in}{1.149156in}}{\pgfqpoint{5.754938in}{1.142221in}}%
\pgfpathquadraticcurveto{\pgfqpoint{5.754938in}{1.136709in}}{\pgfqpoint{5.758576in}{1.133467in}}%
\pgfpathquadraticcurveto{\pgfqpoint{5.762215in}{1.130243in}}{\pgfqpoint{5.768447in}{1.130243in}}%
\pgfpathquadraticcurveto{\pgfqpoint{5.777075in}{1.130243in}}{\pgfqpoint{5.782280in}{1.136349in}}%
\pgfpathquadraticcurveto{\pgfqpoint{5.787486in}{1.142455in}}{\pgfqpoint{5.787486in}{1.152578in}}%
\pgfpathlineto{\pgfqpoint{5.787486in}{1.154883in}}%
\pgfpathlineto{\pgfqpoint{5.777183in}{1.154883in}}%
\pgfpathlineto{\pgfqpoint{5.777183in}{1.154883in}}%
\pgfpathclose%
\pgfpathmoveto{\pgfqpoint{5.797843in}{1.159170in}}%
\pgfpathlineto{\pgfqpoint{5.797843in}{1.123200in}}%
\pgfpathlineto{\pgfqpoint{5.787486in}{1.123200in}}%
\pgfpathlineto{\pgfqpoint{5.787486in}{1.132764in}}%
\pgfpathquadraticcurveto{\pgfqpoint{5.783937in}{1.127037in}}{\pgfqpoint{5.778642in}{1.124299in}}%
\pgfpathquadraticcurveto{\pgfqpoint{5.773346in}{1.121561in}}{\pgfqpoint{5.765691in}{1.121561in}}%
\pgfpathquadraticcurveto{\pgfqpoint{5.756018in}{1.121561in}}{\pgfqpoint{5.750291in}{1.127001in}}%
\pgfpathquadraticcurveto{\pgfqpoint{5.744581in}{1.132440in}}{\pgfqpoint{5.744581in}{1.141554in}}%
\pgfpathquadraticcurveto{\pgfqpoint{5.744581in}{1.152182in}}{\pgfqpoint{5.751696in}{1.157585in}}%
\pgfpathquadraticcurveto{\pgfqpoint{5.758828in}{1.162989in}}{\pgfqpoint{5.772950in}{1.162989in}}%
\pgfpathlineto{\pgfqpoint{5.787486in}{1.162989in}}%
\pgfpathlineto{\pgfqpoint{5.787486in}{1.164016in}}%
\pgfpathquadraticcurveto{\pgfqpoint{5.787486in}{1.171166in}}{\pgfqpoint{5.782785in}{1.175075in}}%
\pgfpathquadraticcurveto{\pgfqpoint{5.778083in}{1.178984in}}{\pgfqpoint{5.769582in}{1.178984in}}%
\pgfpathquadraticcurveto{\pgfqpoint{5.764178in}{1.178984in}}{\pgfqpoint{5.759044in}{1.177687in}}%
\pgfpathquadraticcurveto{\pgfqpoint{5.753929in}{1.176390in}}{\pgfqpoint{5.749210in}{1.173796in}}%
\pgfpathlineto{\pgfqpoint{5.749210in}{1.183379in}}%
\pgfpathquadraticcurveto{\pgfqpoint{5.754884in}{1.185576in}}{\pgfqpoint{5.760233in}{1.186657in}}%
\pgfpathquadraticcurveto{\pgfqpoint{5.765583in}{1.187756in}}{\pgfqpoint{5.770644in}{1.187756in}}%
\pgfpathquadraticcurveto{\pgfqpoint{5.784334in}{1.187756in}}{\pgfqpoint{5.791088in}{1.180659in}}%
\pgfpathquadraticcurveto{\pgfqpoint{5.797843in}{1.173580in}}{\pgfqpoint{5.797843in}{1.159170in}}%
\pgfpathlineto{\pgfqpoint{5.797843in}{1.159170in}}%
\pgfpathclose%
\pgfpathmoveto{\pgfqpoint{5.871584in}{1.161260in}}%
\pgfpathlineto{\pgfqpoint{5.871584in}{1.123200in}}%
\pgfpathlineto{\pgfqpoint{5.861227in}{1.123200in}}%
\pgfpathlineto{\pgfqpoint{5.861227in}{1.160917in}}%
\pgfpathquadraticcurveto{\pgfqpoint{5.861227in}{1.169869in}}{\pgfqpoint{5.857733in}{1.174300in}}%
\pgfpathquadraticcurveto{\pgfqpoint{5.854239in}{1.178749in}}{\pgfqpoint{5.847268in}{1.178749in}}%
\pgfpathquadraticcurveto{\pgfqpoint{5.838874in}{1.178749in}}{\pgfqpoint{5.834029in}{1.173400in}}%
\pgfpathquadraticcurveto{\pgfqpoint{5.829184in}{1.168068in}}{\pgfqpoint{5.829184in}{1.158828in}}%
\pgfpathlineto{\pgfqpoint{5.829184in}{1.123200in}}%
\pgfpathlineto{\pgfqpoint{5.818773in}{1.123200in}}%
\pgfpathlineto{\pgfqpoint{5.818773in}{1.186243in}}%
\pgfpathlineto{\pgfqpoint{5.829184in}{1.186243in}}%
\pgfpathlineto{\pgfqpoint{5.829184in}{1.176444in}}%
\pgfpathquadraticcurveto{\pgfqpoint{5.832912in}{1.182136in}}{\pgfqpoint{5.837938in}{1.184946in}}%
\pgfpathquadraticcurveto{\pgfqpoint{5.842981in}{1.187756in}}{\pgfqpoint{5.849574in}{1.187756in}}%
\pgfpathquadraticcurveto{\pgfqpoint{5.860435in}{1.187756in}}{\pgfqpoint{5.866001in}{1.181037in}}%
\pgfpathquadraticcurveto{\pgfqpoint{5.871584in}{1.174318in}}{\pgfqpoint{5.871584in}{1.161260in}}%
\pgfpathlineto{\pgfqpoint{5.871584in}{1.161260in}}%
\pgfpathclose%
\pgfpathmoveto{\pgfqpoint{5.933708in}{1.176678in}}%
\pgfpathlineto{\pgfqpoint{5.933708in}{1.210793in}}%
\pgfpathlineto{\pgfqpoint{5.944065in}{1.210793in}}%
\pgfpathlineto{\pgfqpoint{5.944065in}{1.123200in}}%
\pgfpathlineto{\pgfqpoint{5.933708in}{1.123200in}}%
\pgfpathlineto{\pgfqpoint{5.933708in}{1.132656in}}%
\pgfpathquadraticcurveto{\pgfqpoint{5.930448in}{1.127037in}}{\pgfqpoint{5.925459in}{1.124299in}}%
\pgfpathquadraticcurveto{\pgfqpoint{5.920487in}{1.121561in}}{\pgfqpoint{5.913499in}{1.121561in}}%
\pgfpathquadraticcurveto{\pgfqpoint{5.902079in}{1.121561in}}{\pgfqpoint{5.894892in}{1.130675in}}%
\pgfpathquadraticcurveto{\pgfqpoint{5.887723in}{1.139807in}}{\pgfqpoint{5.887723in}{1.154667in}}%
\pgfpathquadraticcurveto{\pgfqpoint{5.887723in}{1.169527in}}{\pgfqpoint{5.894892in}{1.178641in}}%
\pgfpathquadraticcurveto{\pgfqpoint{5.902079in}{1.187756in}}{\pgfqpoint{5.913499in}{1.187756in}}%
\pgfpathquadraticcurveto{\pgfqpoint{5.920487in}{1.187756in}}{\pgfqpoint{5.925459in}{1.185018in}}%
\pgfpathquadraticcurveto{\pgfqpoint{5.930448in}{1.182298in}}{\pgfqpoint{5.933708in}{1.176678in}}%
\pgfpathlineto{\pgfqpoint{5.933708in}{1.176678in}}%
\pgfpathclose%
\pgfpathmoveto{\pgfqpoint{5.898423in}{1.154667in}}%
\pgfpathquadraticcurveto{\pgfqpoint{5.898423in}{1.143248in}}{\pgfqpoint{5.903124in}{1.136745in}}%
\pgfpathquadraticcurveto{\pgfqpoint{5.907825in}{1.130243in}}{\pgfqpoint{5.916038in}{1.130243in}}%
\pgfpathquadraticcurveto{\pgfqpoint{5.924252in}{1.130243in}}{\pgfqpoint{5.928971in}{1.136745in}}%
\pgfpathquadraticcurveto{\pgfqpoint{5.933708in}{1.143248in}}{\pgfqpoint{5.933708in}{1.154667in}}%
\pgfpathquadraticcurveto{\pgfqpoint{5.933708in}{1.166087in}}{\pgfqpoint{5.928971in}{1.172589in}}%
\pgfpathquadraticcurveto{\pgfqpoint{5.924252in}{1.179092in}}{\pgfqpoint{5.916038in}{1.179092in}}%
\pgfpathquadraticcurveto{\pgfqpoint{5.907825in}{1.179092in}}{\pgfqpoint{5.903124in}{1.172589in}}%
\pgfpathquadraticcurveto{\pgfqpoint{5.898423in}{1.166087in}}{\pgfqpoint{5.898423in}{1.154667in}}%
\pgfpathlineto{\pgfqpoint{5.898423in}{1.154667in}}%
\pgfpathclose%
\pgfusepath{fill}%
\end{pgfscope}%
\begin{pgfscope}%
\pgfsetbuttcap%
\pgfsetroundjoin%
\definecolor{currentfill}{rgb}{0.309804,0.372549,0.419608}%
\pgfsetfillcolor{currentfill}%
\pgfsetlinewidth{0.000000pt}%
\definecolor{currentstroke}{rgb}{0.309804,0.372549,0.419608}%
\pgfsetstrokecolor{currentstroke}%
\pgfsetdash{}{0pt}%
\pgfpathmoveto{\pgfqpoint{6.047558in}{1.132656in}}%
\pgfpathlineto{\pgfqpoint{6.047558in}{1.099226in}}%
\pgfpathlineto{\pgfqpoint{6.037147in}{1.099226in}}%
\pgfpathlineto{\pgfqpoint{6.037147in}{1.186243in}}%
\pgfpathlineto{\pgfqpoint{6.047558in}{1.186243in}}%
\pgfpathlineto{\pgfqpoint{6.047558in}{1.176678in}}%
\pgfpathquadraticcurveto{\pgfqpoint{6.050837in}{1.182298in}}{\pgfqpoint{6.055808in}{1.185018in}}%
\pgfpathquadraticcurveto{\pgfqpoint{6.060797in}{1.187756in}}{\pgfqpoint{6.067714in}{1.187756in}}%
\pgfpathquadraticcurveto{\pgfqpoint{6.079206in}{1.187756in}}{\pgfqpoint{6.086375in}{1.178641in}}%
\pgfpathquadraticcurveto{\pgfqpoint{6.093561in}{1.169527in}}{\pgfqpoint{6.093561in}{1.154667in}}%
\pgfpathquadraticcurveto{\pgfqpoint{6.093561in}{1.139807in}}{\pgfqpoint{6.086375in}{1.130675in}}%
\pgfpathquadraticcurveto{\pgfqpoint{6.079206in}{1.121561in}}{\pgfqpoint{6.067714in}{1.121561in}}%
\pgfpathquadraticcurveto{\pgfqpoint{6.060797in}{1.121561in}}{\pgfqpoint{6.055808in}{1.124299in}}%
\pgfpathquadraticcurveto{\pgfqpoint{6.050837in}{1.127037in}}{\pgfqpoint{6.047558in}{1.132656in}}%
\pgfpathlineto{\pgfqpoint{6.047558in}{1.132656in}}%
\pgfpathclose%
\pgfpathmoveto{\pgfqpoint{6.082808in}{1.154667in}}%
\pgfpathquadraticcurveto{\pgfqpoint{6.082808in}{1.166087in}}{\pgfqpoint{6.078107in}{1.172589in}}%
\pgfpathquadraticcurveto{\pgfqpoint{6.073406in}{1.179092in}}{\pgfqpoint{6.065192in}{1.179092in}}%
\pgfpathquadraticcurveto{\pgfqpoint{6.056961in}{1.179092in}}{\pgfqpoint{6.052260in}{1.172589in}}%
\pgfpathquadraticcurveto{\pgfqpoint{6.047558in}{1.166087in}}{\pgfqpoint{6.047558in}{1.154667in}}%
\pgfpathquadraticcurveto{\pgfqpoint{6.047558in}{1.143248in}}{\pgfqpoint{6.052260in}{1.136745in}}%
\pgfpathquadraticcurveto{\pgfqpoint{6.056961in}{1.130243in}}{\pgfqpoint{6.065192in}{1.130243in}}%
\pgfpathquadraticcurveto{\pgfqpoint{6.073406in}{1.130243in}}{\pgfqpoint{6.078107in}{1.136745in}}%
\pgfpathquadraticcurveto{\pgfqpoint{6.082808in}{1.143248in}}{\pgfqpoint{6.082808in}{1.154667in}}%
\pgfpathlineto{\pgfqpoint{6.082808in}{1.154667in}}%
\pgfpathclose%
\pgfpathmoveto{\pgfqpoint{6.135151in}{1.178984in}}%
\pgfpathquadraticcurveto{\pgfqpoint{6.126830in}{1.178984in}}{\pgfqpoint{6.121985in}{1.172481in}}%
\pgfpathquadraticcurveto{\pgfqpoint{6.117139in}{1.165979in}}{\pgfqpoint{6.117139in}{1.154667in}}%
\pgfpathquadraticcurveto{\pgfqpoint{6.117139in}{1.143356in}}{\pgfqpoint{6.121949in}{1.136853in}}%
\pgfpathquadraticcurveto{\pgfqpoint{6.126776in}{1.130351in}}{\pgfqpoint{6.135151in}{1.130351in}}%
\pgfpathquadraticcurveto{\pgfqpoint{6.143437in}{1.130351in}}{\pgfqpoint{6.148264in}{1.136871in}}%
\pgfpathquadraticcurveto{\pgfqpoint{6.153110in}{1.143410in}}{\pgfqpoint{6.153110in}{1.154667in}}%
\pgfpathquadraticcurveto{\pgfqpoint{6.153110in}{1.165871in}}{\pgfqpoint{6.148264in}{1.172427in}}%
\pgfpathquadraticcurveto{\pgfqpoint{6.143437in}{1.178984in}}{\pgfqpoint{6.135151in}{1.178984in}}%
\pgfpathlineto{\pgfqpoint{6.135151in}{1.178984in}}%
\pgfpathclose%
\pgfpathmoveto{\pgfqpoint{6.135151in}{1.187756in}}%
\pgfpathquadraticcurveto{\pgfqpoint{6.148661in}{1.187756in}}{\pgfqpoint{6.156370in}{1.178966in}}%
\pgfpathquadraticcurveto{\pgfqpoint{6.164097in}{1.170194in}}{\pgfqpoint{6.164097in}{1.154667in}}%
\pgfpathquadraticcurveto{\pgfqpoint{6.164097in}{1.139195in}}{\pgfqpoint{6.156370in}{1.130369in}}%
\pgfpathquadraticcurveto{\pgfqpoint{6.148661in}{1.121561in}}{\pgfqpoint{6.135151in}{1.121561in}}%
\pgfpathquadraticcurveto{\pgfqpoint{6.121588in}{1.121561in}}{\pgfqpoint{6.113897in}{1.130369in}}%
\pgfpathquadraticcurveto{\pgfqpoint{6.106224in}{1.139195in}}{\pgfqpoint{6.106224in}{1.154667in}}%
\pgfpathquadraticcurveto{\pgfqpoint{6.106224in}{1.170194in}}{\pgfqpoint{6.113897in}{1.178966in}}%
\pgfpathquadraticcurveto{\pgfqpoint{6.121588in}{1.187756in}}{\pgfqpoint{6.135151in}{1.187756in}}%
\pgfpathlineto{\pgfqpoint{6.135151in}{1.187756in}}%
\pgfpathclose%
\pgfpathmoveto{\pgfqpoint{6.221448in}{1.184387in}}%
\pgfpathlineto{\pgfqpoint{6.221448in}{1.174589in}}%
\pgfpathquadraticcurveto{\pgfqpoint{6.217071in}{1.176840in}}{\pgfqpoint{6.212334in}{1.177957in}}%
\pgfpathquadraticcurveto{\pgfqpoint{6.207614in}{1.179092in}}{\pgfqpoint{6.202535in}{1.179092in}}%
\pgfpathquadraticcurveto{\pgfqpoint{6.194826in}{1.179092in}}{\pgfqpoint{6.190971in}{1.176732in}}%
\pgfpathquadraticcurveto{\pgfqpoint{6.187117in}{1.174373in}}{\pgfqpoint{6.187117in}{1.169635in}}%
\pgfpathquadraticcurveto{\pgfqpoint{6.187117in}{1.166033in}}{\pgfqpoint{6.189872in}{1.163980in}}%
\pgfpathquadraticcurveto{\pgfqpoint{6.192628in}{1.161926in}}{\pgfqpoint{6.200968in}{1.160071in}}%
\pgfpathlineto{\pgfqpoint{6.204516in}{1.159278in}}%
\pgfpathquadraticcurveto{\pgfqpoint{6.215540in}{1.156919in}}{\pgfqpoint{6.220187in}{1.152614in}}%
\pgfpathquadraticcurveto{\pgfqpoint{6.224834in}{1.148309in}}{\pgfqpoint{6.224834in}{1.140600in}}%
\pgfpathquadraticcurveto{\pgfqpoint{6.224834in}{1.131810in}}{\pgfqpoint{6.217881in}{1.126676in}}%
\pgfpathquadraticcurveto{\pgfqpoint{6.210929in}{1.121561in}}{\pgfqpoint{6.198770in}{1.121561in}}%
\pgfpathquadraticcurveto{\pgfqpoint{6.193709in}{1.121561in}}{\pgfqpoint{6.188215in}{1.122552in}}%
\pgfpathquadraticcurveto{\pgfqpoint{6.182722in}{1.123542in}}{\pgfqpoint{6.176651in}{1.125506in}}%
\pgfpathlineto{\pgfqpoint{6.176651in}{1.136205in}}%
\pgfpathquadraticcurveto{\pgfqpoint{6.182397in}{1.133215in}}{\pgfqpoint{6.187963in}{1.131720in}}%
\pgfpathquadraticcurveto{\pgfqpoint{6.193529in}{1.130243in}}{\pgfqpoint{6.199005in}{1.130243in}}%
\pgfpathquadraticcurveto{\pgfqpoint{6.206317in}{1.130243in}}{\pgfqpoint{6.210244in}{1.132746in}}%
\pgfpathquadraticcurveto{\pgfqpoint{6.214189in}{1.135250in}}{\pgfqpoint{6.214189in}{1.139807in}}%
\pgfpathquadraticcurveto{\pgfqpoint{6.214189in}{1.144022in}}{\pgfqpoint{6.211343in}{1.146274in}}%
\pgfpathquadraticcurveto{\pgfqpoint{6.208515in}{1.148525in}}{\pgfqpoint{6.198878in}{1.150614in}}%
\pgfpathlineto{\pgfqpoint{6.195276in}{1.151461in}}%
\pgfpathquadraticcurveto{\pgfqpoint{6.185658in}{1.153478in}}{\pgfqpoint{6.181371in}{1.157675in}}%
\pgfpathquadraticcurveto{\pgfqpoint{6.177102in}{1.161872in}}{\pgfqpoint{6.177102in}{1.169185in}}%
\pgfpathquadraticcurveto{\pgfqpoint{6.177102in}{1.178083in}}{\pgfqpoint{6.183406in}{1.182910in}}%
\pgfpathquadraticcurveto{\pgfqpoint{6.189710in}{1.187756in}}{\pgfqpoint{6.201310in}{1.187756in}}%
\pgfpathquadraticcurveto{\pgfqpoint{6.207038in}{1.187756in}}{\pgfqpoint{6.212099in}{1.186909in}}%
\pgfpathquadraticcurveto{\pgfqpoint{6.217179in}{1.186080in}}{\pgfqpoint{6.221448in}{1.184387in}}%
\pgfpathlineto{\pgfqpoint{6.221448in}{1.184387in}}%
\pgfpathclose%
\pgfpathmoveto{\pgfqpoint{6.251564in}{1.204147in}}%
\pgfpathlineto{\pgfqpoint{6.251564in}{1.186243in}}%
\pgfpathlineto{\pgfqpoint{6.272890in}{1.186243in}}%
\pgfpathlineto{\pgfqpoint{6.272890in}{1.178191in}}%
\pgfpathlineto{\pgfqpoint{6.251564in}{1.178191in}}%
\pgfpathlineto{\pgfqpoint{6.251564in}{1.143968in}}%
\pgfpathquadraticcurveto{\pgfqpoint{6.251564in}{1.136259in}}{\pgfqpoint{6.253671in}{1.134061in}}%
\pgfpathquadraticcurveto{\pgfqpoint{6.255779in}{1.131864in}}{\pgfqpoint{6.262263in}{1.131864in}}%
\pgfpathlineto{\pgfqpoint{6.272890in}{1.131864in}}%
\pgfpathlineto{\pgfqpoint{6.272890in}{1.123200in}}%
\pgfpathlineto{\pgfqpoint{6.262263in}{1.123200in}}%
\pgfpathquadraticcurveto{\pgfqpoint{6.250267in}{1.123200in}}{\pgfqpoint{6.245710in}{1.127667in}}%
\pgfpathquadraticcurveto{\pgfqpoint{6.241153in}{1.132152in}}{\pgfqpoint{6.241153in}{1.143968in}}%
\pgfpathlineto{\pgfqpoint{6.241153in}{1.178191in}}%
\pgfpathlineto{\pgfqpoint{6.233552in}{1.178191in}}%
\pgfpathlineto{\pgfqpoint{6.233552in}{1.186243in}}%
\pgfpathlineto{\pgfqpoint{6.241153in}{1.186243in}}%
\pgfpathlineto{\pgfqpoint{6.241153in}{1.204147in}}%
\pgfpathlineto{\pgfqpoint{6.251564in}{1.204147in}}%
\pgfpathlineto{\pgfqpoint{6.251564in}{1.204147in}}%
\pgfpathclose%
\pgfusepath{fill}%
\end{pgfscope}%
\begin{pgfscope}%
\pgfsetbuttcap%
\pgfsetroundjoin%
\definecolor{currentfill}{rgb}{0.309804,0.372549,0.419608}%
\pgfsetfillcolor{currentfill}%
\pgfsetlinewidth{0.000000pt}%
\definecolor{currentstroke}{rgb}{0.309804,0.372549,0.419608}%
\pgfsetstrokecolor{currentstroke}%
\pgfsetdash{}{0pt}%
\pgfpathmoveto{\pgfqpoint{6.410490in}{1.157315in}}%
\pgfpathlineto{\pgfqpoint{6.410490in}{1.152254in}}%
\pgfpathlineto{\pgfqpoint{6.362866in}{1.152254in}}%
\pgfpathquadraticcurveto{\pgfqpoint{6.363551in}{1.141554in}}{\pgfqpoint{6.369314in}{1.135953in}}%
\pgfpathquadraticcurveto{\pgfqpoint{6.375078in}{1.130351in}}{\pgfqpoint{6.385381in}{1.130351in}}%
\pgfpathquadraticcurveto{\pgfqpoint{6.391343in}{1.130351in}}{\pgfqpoint{6.396945in}{1.131810in}}%
\pgfpathquadraticcurveto{\pgfqpoint{6.402547in}{1.133269in}}{\pgfqpoint{6.408077in}{1.136205in}}%
\pgfpathlineto{\pgfqpoint{6.408077in}{1.126406in}}%
\pgfpathquadraticcurveto{\pgfqpoint{6.402493in}{1.124047in}}{\pgfqpoint{6.396639in}{1.122804in}}%
\pgfpathquadraticcurveto{\pgfqpoint{6.390785in}{1.121561in}}{\pgfqpoint{6.384769in}{1.121561in}}%
\pgfpathquadraticcurveto{\pgfqpoint{6.369675in}{1.121561in}}{\pgfqpoint{6.360867in}{1.130333in}}%
\pgfpathquadraticcurveto{\pgfqpoint{6.352059in}{1.139123in}}{\pgfqpoint{6.352059in}{1.154109in}}%
\pgfpathquadraticcurveto{\pgfqpoint{6.352059in}{1.169581in}}{\pgfqpoint{6.360416in}{1.178659in}}%
\pgfpathquadraticcurveto{\pgfqpoint{6.368774in}{1.187756in}}{\pgfqpoint{6.382968in}{1.187756in}}%
\pgfpathquadraticcurveto{\pgfqpoint{6.395684in}{1.187756in}}{\pgfqpoint{6.403087in}{1.179560in}}%
\pgfpathquadraticcurveto{\pgfqpoint{6.410490in}{1.171383in}}{\pgfqpoint{6.410490in}{1.157315in}}%
\pgfpathlineto{\pgfqpoint{6.410490in}{1.157315in}}%
\pgfpathclose%
\pgfpathmoveto{\pgfqpoint{6.400133in}{1.160359in}}%
\pgfpathquadraticcurveto{\pgfqpoint{6.400025in}{1.168843in}}{\pgfqpoint{6.395378in}{1.173904in}}%
\pgfpathquadraticcurveto{\pgfqpoint{6.390731in}{1.178984in}}{\pgfqpoint{6.383076in}{1.178984in}}%
\pgfpathquadraticcurveto{\pgfqpoint{6.374412in}{1.178984in}}{\pgfqpoint{6.369206in}{1.174084in}}%
\pgfpathquadraticcurveto{\pgfqpoint{6.364001in}{1.169185in}}{\pgfqpoint{6.363208in}{1.160287in}}%
\pgfpathlineto{\pgfqpoint{6.400133in}{1.160359in}}%
\pgfpathlineto{\pgfqpoint{6.400133in}{1.160359in}}%
\pgfpathclose%
\pgfpathmoveto{\pgfqpoint{6.479909in}{1.161260in}}%
\pgfpathlineto{\pgfqpoint{6.479909in}{1.123200in}}%
\pgfpathlineto{\pgfqpoint{6.469552in}{1.123200in}}%
\pgfpathlineto{\pgfqpoint{6.469552in}{1.160917in}}%
\pgfpathquadraticcurveto{\pgfqpoint{6.469552in}{1.169869in}}{\pgfqpoint{6.466058in}{1.174300in}}%
\pgfpathquadraticcurveto{\pgfqpoint{6.462563in}{1.178749in}}{\pgfqpoint{6.455593in}{1.178749in}}%
\pgfpathquadraticcurveto{\pgfqpoint{6.447199in}{1.178749in}}{\pgfqpoint{6.442354in}{1.173400in}}%
\pgfpathquadraticcurveto{\pgfqpoint{6.437508in}{1.168068in}}{\pgfqpoint{6.437508in}{1.158828in}}%
\pgfpathlineto{\pgfqpoint{6.437508in}{1.123200in}}%
\pgfpathlineto{\pgfqpoint{6.427097in}{1.123200in}}%
\pgfpathlineto{\pgfqpoint{6.427097in}{1.186243in}}%
\pgfpathlineto{\pgfqpoint{6.437508in}{1.186243in}}%
\pgfpathlineto{\pgfqpoint{6.437508in}{1.176444in}}%
\pgfpathquadraticcurveto{\pgfqpoint{6.441237in}{1.182136in}}{\pgfqpoint{6.446262in}{1.184946in}}%
\pgfpathquadraticcurveto{\pgfqpoint{6.451306in}{1.187756in}}{\pgfqpoint{6.457898in}{1.187756in}}%
\pgfpathquadraticcurveto{\pgfqpoint{6.468760in}{1.187756in}}{\pgfqpoint{6.474325in}{1.181037in}}%
\pgfpathquadraticcurveto{\pgfqpoint{6.479909in}{1.174318in}}{\pgfqpoint{6.479909in}{1.161260in}}%
\pgfpathlineto{\pgfqpoint{6.479909in}{1.161260in}}%
\pgfpathclose%
\pgfpathmoveto{\pgfqpoint{6.510800in}{1.204147in}}%
\pgfpathlineto{\pgfqpoint{6.510800in}{1.186243in}}%
\pgfpathlineto{\pgfqpoint{6.532126in}{1.186243in}}%
\pgfpathlineto{\pgfqpoint{6.532126in}{1.178191in}}%
\pgfpathlineto{\pgfqpoint{6.510800in}{1.178191in}}%
\pgfpathlineto{\pgfqpoint{6.510800in}{1.143968in}}%
\pgfpathquadraticcurveto{\pgfqpoint{6.510800in}{1.136259in}}{\pgfqpoint{6.512907in}{1.134061in}}%
\pgfpathquadraticcurveto{\pgfqpoint{6.515015in}{1.131864in}}{\pgfqpoint{6.521499in}{1.131864in}}%
\pgfpathlineto{\pgfqpoint{6.532126in}{1.131864in}}%
\pgfpathlineto{\pgfqpoint{6.532126in}{1.123200in}}%
\pgfpathlineto{\pgfqpoint{6.521499in}{1.123200in}}%
\pgfpathquadraticcurveto{\pgfqpoint{6.509503in}{1.123200in}}{\pgfqpoint{6.504946in}{1.127667in}}%
\pgfpathquadraticcurveto{\pgfqpoint{6.500389in}{1.132152in}}{\pgfqpoint{6.500389in}{1.143968in}}%
\pgfpathlineto{\pgfqpoint{6.500389in}{1.178191in}}%
\pgfpathlineto{\pgfqpoint{6.492788in}{1.178191in}}%
\pgfpathlineto{\pgfqpoint{6.492788in}{1.186243in}}%
\pgfpathlineto{\pgfqpoint{6.500389in}{1.186243in}}%
\pgfpathlineto{\pgfqpoint{6.500389in}{1.204147in}}%
\pgfpathlineto{\pgfqpoint{6.510800in}{1.204147in}}%
\pgfpathlineto{\pgfqpoint{6.510800in}{1.204147in}}%
\pgfpathclose%
\pgfpathmoveto{\pgfqpoint{6.599672in}{1.157315in}}%
\pgfpathlineto{\pgfqpoint{6.599672in}{1.152254in}}%
\pgfpathlineto{\pgfqpoint{6.552048in}{1.152254in}}%
\pgfpathquadraticcurveto{\pgfqpoint{6.552732in}{1.141554in}}{\pgfqpoint{6.558496in}{1.135953in}}%
\pgfpathquadraticcurveto{\pgfqpoint{6.564260in}{1.130351in}}{\pgfqpoint{6.574563in}{1.130351in}}%
\pgfpathquadraticcurveto{\pgfqpoint{6.580525in}{1.130351in}}{\pgfqpoint{6.586127in}{1.131810in}}%
\pgfpathquadraticcurveto{\pgfqpoint{6.591729in}{1.133269in}}{\pgfqpoint{6.597258in}{1.136205in}}%
\pgfpathlineto{\pgfqpoint{6.597258in}{1.126406in}}%
\pgfpathquadraticcurveto{\pgfqpoint{6.591674in}{1.124047in}}{\pgfqpoint{6.585821in}{1.122804in}}%
\pgfpathquadraticcurveto{\pgfqpoint{6.579967in}{1.121561in}}{\pgfqpoint{6.573951in}{1.121561in}}%
\pgfpathquadraticcurveto{\pgfqpoint{6.558856in}{1.121561in}}{\pgfqpoint{6.550048in}{1.130333in}}%
\pgfpathquadraticcurveto{\pgfqpoint{6.541240in}{1.139123in}}{\pgfqpoint{6.541240in}{1.154109in}}%
\pgfpathquadraticcurveto{\pgfqpoint{6.541240in}{1.169581in}}{\pgfqpoint{6.549598in}{1.178659in}}%
\pgfpathquadraticcurveto{\pgfqpoint{6.557956in}{1.187756in}}{\pgfqpoint{6.572149in}{1.187756in}}%
\pgfpathquadraticcurveto{\pgfqpoint{6.584866in}{1.187756in}}{\pgfqpoint{6.592269in}{1.179560in}}%
\pgfpathquadraticcurveto{\pgfqpoint{6.599672in}{1.171383in}}{\pgfqpoint{6.599672in}{1.157315in}}%
\pgfpathlineto{\pgfqpoint{6.599672in}{1.157315in}}%
\pgfpathclose%
\pgfpathmoveto{\pgfqpoint{6.589315in}{1.160359in}}%
\pgfpathquadraticcurveto{\pgfqpoint{6.589207in}{1.168843in}}{\pgfqpoint{6.584560in}{1.173904in}}%
\pgfpathquadraticcurveto{\pgfqpoint{6.579913in}{1.178984in}}{\pgfqpoint{6.572257in}{1.178984in}}%
\pgfpathquadraticcurveto{\pgfqpoint{6.563594in}{1.178984in}}{\pgfqpoint{6.558388in}{1.174084in}}%
\pgfpathquadraticcurveto{\pgfqpoint{6.553183in}{1.169185in}}{\pgfqpoint{6.552390in}{1.160287in}}%
\pgfpathlineto{\pgfqpoint{6.589315in}{1.160359in}}%
\pgfpathlineto{\pgfqpoint{6.589315in}{1.160359in}}%
\pgfpathclose%
\pgfpathmoveto{\pgfqpoint{6.653204in}{1.176570in}}%
\pgfpathquadraticcurveto{\pgfqpoint{6.651457in}{1.177579in}}{\pgfqpoint{6.649403in}{1.178047in}}%
\pgfpathquadraticcurveto{\pgfqpoint{6.647350in}{1.178533in}}{\pgfqpoint{6.644882in}{1.178533in}}%
\pgfpathquadraticcurveto{\pgfqpoint{6.636092in}{1.178533in}}{\pgfqpoint{6.631391in}{1.172823in}}%
\pgfpathquadraticcurveto{\pgfqpoint{6.626690in}{1.167114in}}{\pgfqpoint{6.626690in}{1.156414in}}%
\pgfpathlineto{\pgfqpoint{6.626690in}{1.123200in}}%
\pgfpathlineto{\pgfqpoint{6.616279in}{1.123200in}}%
\pgfpathlineto{\pgfqpoint{6.616279in}{1.186243in}}%
\pgfpathlineto{\pgfqpoint{6.626690in}{1.186243in}}%
\pgfpathlineto{\pgfqpoint{6.626690in}{1.176444in}}%
\pgfpathquadraticcurveto{\pgfqpoint{6.629968in}{1.182190in}}{\pgfqpoint{6.635192in}{1.184964in}}%
\pgfpathquadraticcurveto{\pgfqpoint{6.640433in}{1.187756in}}{\pgfqpoint{6.647926in}{1.187756in}}%
\pgfpathquadraticcurveto{\pgfqpoint{6.648989in}{1.187756in}}{\pgfqpoint{6.650286in}{1.187611in}}%
\pgfpathquadraticcurveto{\pgfqpoint{6.651583in}{1.187485in}}{\pgfqpoint{6.653150in}{1.187197in}}%
\pgfpathlineto{\pgfqpoint{6.653204in}{1.176570in}}%
\pgfpathlineto{\pgfqpoint{6.653204in}{1.176570in}}%
\pgfpathclose%
\pgfpathmoveto{\pgfqpoint{6.664065in}{1.186243in}}%
\pgfpathlineto{\pgfqpoint{6.674422in}{1.186243in}}%
\pgfpathlineto{\pgfqpoint{6.674422in}{1.123200in}}%
\pgfpathlineto{\pgfqpoint{6.664065in}{1.123200in}}%
\pgfpathlineto{\pgfqpoint{6.664065in}{1.186243in}}%
\pgfpathlineto{\pgfqpoint{6.664065in}{1.186243in}}%
\pgfpathclose%
\pgfpathmoveto{\pgfqpoint{6.664065in}{1.210793in}}%
\pgfpathlineto{\pgfqpoint{6.674422in}{1.210793in}}%
\pgfpathlineto{\pgfqpoint{6.674422in}{1.197662in}}%
\pgfpathlineto{\pgfqpoint{6.664065in}{1.197662in}}%
\pgfpathlineto{\pgfqpoint{6.664065in}{1.210793in}}%
\pgfpathlineto{\pgfqpoint{6.664065in}{1.210793in}}%
\pgfpathclose%
\pgfpathmoveto{\pgfqpoint{6.748506in}{1.161260in}}%
\pgfpathlineto{\pgfqpoint{6.748506in}{1.123200in}}%
\pgfpathlineto{\pgfqpoint{6.738149in}{1.123200in}}%
\pgfpathlineto{\pgfqpoint{6.738149in}{1.160917in}}%
\pgfpathquadraticcurveto{\pgfqpoint{6.738149in}{1.169869in}}{\pgfqpoint{6.734655in}{1.174300in}}%
\pgfpathquadraticcurveto{\pgfqpoint{6.731161in}{1.178749in}}{\pgfqpoint{6.724190in}{1.178749in}}%
\pgfpathquadraticcurveto{\pgfqpoint{6.715796in}{1.178749in}}{\pgfqpoint{6.710951in}{1.173400in}}%
\pgfpathquadraticcurveto{\pgfqpoint{6.706106in}{1.168068in}}{\pgfqpoint{6.706106in}{1.158828in}}%
\pgfpathlineto{\pgfqpoint{6.706106in}{1.123200in}}%
\pgfpathlineto{\pgfqpoint{6.695695in}{1.123200in}}%
\pgfpathlineto{\pgfqpoint{6.695695in}{1.186243in}}%
\pgfpathlineto{\pgfqpoint{6.706106in}{1.186243in}}%
\pgfpathlineto{\pgfqpoint{6.706106in}{1.176444in}}%
\pgfpathquadraticcurveto{\pgfqpoint{6.709834in}{1.182136in}}{\pgfqpoint{6.714860in}{1.184946in}}%
\pgfpathquadraticcurveto{\pgfqpoint{6.719903in}{1.187756in}}{\pgfqpoint{6.726495in}{1.187756in}}%
\pgfpathquadraticcurveto{\pgfqpoint{6.737357in}{1.187756in}}{\pgfqpoint{6.742923in}{1.181037in}}%
\pgfpathquadraticcurveto{\pgfqpoint{6.748506in}{1.174318in}}{\pgfqpoint{6.748506in}{1.161260in}}%
\pgfpathlineto{\pgfqpoint{6.748506in}{1.161260in}}%
\pgfpathclose%
\pgfpathmoveto{\pgfqpoint{6.810630in}{1.155460in}}%
\pgfpathquadraticcurveto{\pgfqpoint{6.810630in}{1.166717in}}{\pgfqpoint{6.805983in}{1.172896in}}%
\pgfpathquadraticcurveto{\pgfqpoint{6.801354in}{1.179092in}}{\pgfqpoint{6.792960in}{1.179092in}}%
\pgfpathquadraticcurveto{\pgfqpoint{6.784639in}{1.179092in}}{\pgfqpoint{6.779992in}{1.172896in}}%
\pgfpathquadraticcurveto{\pgfqpoint{6.775344in}{1.166717in}}{\pgfqpoint{6.775344in}{1.155460in}}%
\pgfpathquadraticcurveto{\pgfqpoint{6.775344in}{1.144256in}}{\pgfqpoint{6.779992in}{1.138060in}}%
\pgfpathquadraticcurveto{\pgfqpoint{6.784639in}{1.131864in}}{\pgfqpoint{6.792960in}{1.131864in}}%
\pgfpathquadraticcurveto{\pgfqpoint{6.801354in}{1.131864in}}{\pgfqpoint{6.805983in}{1.138060in}}%
\pgfpathquadraticcurveto{\pgfqpoint{6.810630in}{1.144256in}}{\pgfqpoint{6.810630in}{1.155460in}}%
\pgfpathlineto{\pgfqpoint{6.810630in}{1.155460in}}%
\pgfpathclose%
\pgfpathmoveto{\pgfqpoint{6.820987in}{1.131017in}}%
\pgfpathquadraticcurveto{\pgfqpoint{6.820987in}{1.114932in}}{\pgfqpoint{6.813836in}{1.107079in}}%
\pgfpathquadraticcurveto{\pgfqpoint{6.806704in}{1.099226in}}{\pgfqpoint{6.791952in}{1.099226in}}%
\pgfpathquadraticcurveto{\pgfqpoint{6.786494in}{1.099226in}}{\pgfqpoint{6.781649in}{1.100036in}}%
\pgfpathquadraticcurveto{\pgfqpoint{6.776803in}{1.100847in}}{\pgfqpoint{6.772246in}{1.102540in}}%
\pgfpathlineto{\pgfqpoint{6.772246in}{1.112609in}}%
\pgfpathquadraticcurveto{\pgfqpoint{6.776803in}{1.110141in}}{\pgfqpoint{6.781252in}{1.108970in}}%
\pgfpathquadraticcurveto{\pgfqpoint{6.785701in}{1.107782in}}{\pgfqpoint{6.790313in}{1.107782in}}%
\pgfpathquadraticcurveto{\pgfqpoint{6.800507in}{1.107782in}}{\pgfqpoint{6.805569in}{1.113095in}}%
\pgfpathquadraticcurveto{\pgfqpoint{6.810630in}{1.118409in}}{\pgfqpoint{6.810630in}{1.129162in}}%
\pgfpathlineto{\pgfqpoint{6.810630in}{1.134295in}}%
\pgfpathquadraticcurveto{\pgfqpoint{6.807424in}{1.128712in}}{\pgfqpoint{6.802417in}{1.125956in}}%
\pgfpathquadraticcurveto{\pgfqpoint{6.797409in}{1.123200in}}{\pgfqpoint{6.790421in}{1.123200in}}%
\pgfpathquadraticcurveto{\pgfqpoint{6.778839in}{1.123200in}}{\pgfqpoint{6.771742in}{1.132026in}}%
\pgfpathquadraticcurveto{\pgfqpoint{6.764645in}{1.140870in}}{\pgfqpoint{6.764645in}{1.155460in}}%
\pgfpathquadraticcurveto{\pgfqpoint{6.764645in}{1.170086in}}{\pgfqpoint{6.771742in}{1.178912in}}%
\pgfpathquadraticcurveto{\pgfqpoint{6.778839in}{1.187756in}}{\pgfqpoint{6.790421in}{1.187756in}}%
\pgfpathquadraticcurveto{\pgfqpoint{6.797409in}{1.187756in}}{\pgfqpoint{6.802417in}{1.185000in}}%
\pgfpathquadraticcurveto{\pgfqpoint{6.807424in}{1.182244in}}{\pgfqpoint{6.810630in}{1.176678in}}%
\pgfpathlineto{\pgfqpoint{6.810630in}{1.186243in}}%
\pgfpathlineto{\pgfqpoint{6.820987in}{1.186243in}}%
\pgfpathlineto{\pgfqpoint{6.820987in}{1.131017in}}%
\pgfpathlineto{\pgfqpoint{6.820987in}{1.131017in}}%
\pgfpathclose%
\pgfusepath{fill}%
\end{pgfscope}%
\begin{pgfscope}%
\pgfsetbuttcap%
\pgfsetroundjoin%
\definecolor{currentfill}{rgb}{0.309804,0.372549,0.419608}%
\pgfsetfillcolor{currentfill}%
\pgfsetlinewidth{0.000000pt}%
\definecolor{currentstroke}{rgb}{0.309804,0.372549,0.419608}%
\pgfsetstrokecolor{currentstroke}%
\pgfsetdash{}{0pt}%
\pgfpathmoveto{\pgfqpoint{6.965765in}{1.184387in}}%
\pgfpathlineto{\pgfqpoint{6.965765in}{1.174589in}}%
\pgfpathquadraticcurveto{\pgfqpoint{6.961388in}{1.176840in}}{\pgfqpoint{6.956651in}{1.177957in}}%
\pgfpathquadraticcurveto{\pgfqpoint{6.951932in}{1.179092in}}{\pgfqpoint{6.946852in}{1.179092in}}%
\pgfpathquadraticcurveto{\pgfqpoint{6.939143in}{1.179092in}}{\pgfqpoint{6.935289in}{1.176732in}}%
\pgfpathquadraticcurveto{\pgfqpoint{6.931434in}{1.174373in}}{\pgfqpoint{6.931434in}{1.169635in}}%
\pgfpathquadraticcurveto{\pgfqpoint{6.931434in}{1.166033in}}{\pgfqpoint{6.934190in}{1.163980in}}%
\pgfpathquadraticcurveto{\pgfqpoint{6.936946in}{1.161926in}}{\pgfqpoint{6.945285in}{1.160071in}}%
\pgfpathlineto{\pgfqpoint{6.948834in}{1.159278in}}%
\pgfpathquadraticcurveto{\pgfqpoint{6.959857in}{1.156919in}}{\pgfqpoint{6.964504in}{1.152614in}}%
\pgfpathquadraticcurveto{\pgfqpoint{6.969151in}{1.148309in}}{\pgfqpoint{6.969151in}{1.140600in}}%
\pgfpathquadraticcurveto{\pgfqpoint{6.969151in}{1.131810in}}{\pgfqpoint{6.962199in}{1.126676in}}%
\pgfpathquadraticcurveto{\pgfqpoint{6.955246in}{1.121561in}}{\pgfqpoint{6.943088in}{1.121561in}}%
\pgfpathquadraticcurveto{\pgfqpoint{6.938026in}{1.121561in}}{\pgfqpoint{6.932533in}{1.122552in}}%
\pgfpathquadraticcurveto{\pgfqpoint{6.927039in}{1.123542in}}{\pgfqpoint{6.920969in}{1.125506in}}%
\pgfpathlineto{\pgfqpoint{6.920969in}{1.136205in}}%
\pgfpathquadraticcurveto{\pgfqpoint{6.926715in}{1.133215in}}{\pgfqpoint{6.932281in}{1.131720in}}%
\pgfpathquadraticcurveto{\pgfqpoint{6.937846in}{1.130243in}}{\pgfqpoint{6.943322in}{1.130243in}}%
\pgfpathquadraticcurveto{\pgfqpoint{6.950635in}{1.130243in}}{\pgfqpoint{6.954562in}{1.132746in}}%
\pgfpathquadraticcurveto{\pgfqpoint{6.958506in}{1.135250in}}{\pgfqpoint{6.958506in}{1.139807in}}%
\pgfpathquadraticcurveto{\pgfqpoint{6.958506in}{1.144022in}}{\pgfqpoint{6.955660in}{1.146274in}}%
\pgfpathquadraticcurveto{\pgfqpoint{6.952832in}{1.148525in}}{\pgfqpoint{6.943196in}{1.150614in}}%
\pgfpathlineto{\pgfqpoint{6.939594in}{1.151461in}}%
\pgfpathquadraticcurveto{\pgfqpoint{6.929975in}{1.153478in}}{\pgfqpoint{6.925688in}{1.157675in}}%
\pgfpathquadraticcurveto{\pgfqpoint{6.921419in}{1.161872in}}{\pgfqpoint{6.921419in}{1.169185in}}%
\pgfpathquadraticcurveto{\pgfqpoint{6.921419in}{1.178083in}}{\pgfqpoint{6.927724in}{1.182910in}}%
\pgfpathquadraticcurveto{\pgfqpoint{6.934028in}{1.187756in}}{\pgfqpoint{6.945628in}{1.187756in}}%
\pgfpathquadraticcurveto{\pgfqpoint{6.951355in}{1.187756in}}{\pgfqpoint{6.956417in}{1.186909in}}%
\pgfpathquadraticcurveto{\pgfqpoint{6.961496in}{1.186080in}}{\pgfqpoint{6.965765in}{1.184387in}}%
\pgfpathlineto{\pgfqpoint{6.965765in}{1.184387in}}%
\pgfpathclose%
\pgfpathmoveto{\pgfqpoint{7.039561in}{1.157315in}}%
\pgfpathlineto{\pgfqpoint{7.039561in}{1.152254in}}%
\pgfpathlineto{\pgfqpoint{6.991937in}{1.152254in}}%
\pgfpathquadraticcurveto{\pgfqpoint{6.992621in}{1.141554in}}{\pgfqpoint{6.998385in}{1.135953in}}%
\pgfpathquadraticcurveto{\pgfqpoint{7.004149in}{1.130351in}}{\pgfqpoint{7.014452in}{1.130351in}}%
\pgfpathquadraticcurveto{\pgfqpoint{7.020414in}{1.130351in}}{\pgfqpoint{7.026016in}{1.131810in}}%
\pgfpathquadraticcurveto{\pgfqpoint{7.031618in}{1.133269in}}{\pgfqpoint{7.037147in}{1.136205in}}%
\pgfpathlineto{\pgfqpoint{7.037147in}{1.126406in}}%
\pgfpathquadraticcurveto{\pgfqpoint{7.031564in}{1.124047in}}{\pgfqpoint{7.025710in}{1.122804in}}%
\pgfpathquadraticcurveto{\pgfqpoint{7.019856in}{1.121561in}}{\pgfqpoint{7.013840in}{1.121561in}}%
\pgfpathquadraticcurveto{\pgfqpoint{6.998745in}{1.121561in}}{\pgfqpoint{6.989938in}{1.130333in}}%
\pgfpathquadraticcurveto{\pgfqpoint{6.981130in}{1.139123in}}{\pgfqpoint{6.981130in}{1.154109in}}%
\pgfpathquadraticcurveto{\pgfqpoint{6.981130in}{1.169581in}}{\pgfqpoint{6.989487in}{1.178659in}}%
\pgfpathquadraticcurveto{\pgfqpoint{6.997845in}{1.187756in}}{\pgfqpoint{7.012038in}{1.187756in}}%
\pgfpathquadraticcurveto{\pgfqpoint{7.024755in}{1.187756in}}{\pgfqpoint{7.032158in}{1.179560in}}%
\pgfpathquadraticcurveto{\pgfqpoint{7.039561in}{1.171383in}}{\pgfqpoint{7.039561in}{1.157315in}}%
\pgfpathlineto{\pgfqpoint{7.039561in}{1.157315in}}%
\pgfpathclose%
\pgfpathmoveto{\pgfqpoint{7.029204in}{1.160359in}}%
\pgfpathquadraticcurveto{\pgfqpoint{7.029096in}{1.168843in}}{\pgfqpoint{7.024449in}{1.173904in}}%
\pgfpathquadraticcurveto{\pgfqpoint{7.019802in}{1.178984in}}{\pgfqpoint{7.012146in}{1.178984in}}%
\pgfpathquadraticcurveto{\pgfqpoint{7.003483in}{1.178984in}}{\pgfqpoint{6.998277in}{1.174084in}}%
\pgfpathquadraticcurveto{\pgfqpoint{6.993072in}{1.169185in}}{\pgfqpoint{6.992279in}{1.160287in}}%
\pgfpathlineto{\pgfqpoint{7.029204in}{1.160359in}}%
\pgfpathlineto{\pgfqpoint{7.029204in}{1.160359in}}%
\pgfpathclose%
\pgfpathmoveto{\pgfqpoint{7.066579in}{1.132656in}}%
\pgfpathlineto{\pgfqpoint{7.066579in}{1.099226in}}%
\pgfpathlineto{\pgfqpoint{7.056168in}{1.099226in}}%
\pgfpathlineto{\pgfqpoint{7.056168in}{1.186243in}}%
\pgfpathlineto{\pgfqpoint{7.066579in}{1.186243in}}%
\pgfpathlineto{\pgfqpoint{7.066579in}{1.176678in}}%
\pgfpathquadraticcurveto{\pgfqpoint{7.069857in}{1.182298in}}{\pgfqpoint{7.074829in}{1.185018in}}%
\pgfpathquadraticcurveto{\pgfqpoint{7.079818in}{1.187756in}}{\pgfqpoint{7.086735in}{1.187756in}}%
\pgfpathquadraticcurveto{\pgfqpoint{7.098227in}{1.187756in}}{\pgfqpoint{7.105395in}{1.178641in}}%
\pgfpathquadraticcurveto{\pgfqpoint{7.112582in}{1.169527in}}{\pgfqpoint{7.112582in}{1.154667in}}%
\pgfpathquadraticcurveto{\pgfqpoint{7.112582in}{1.139807in}}{\pgfqpoint{7.105395in}{1.130675in}}%
\pgfpathquadraticcurveto{\pgfqpoint{7.098227in}{1.121561in}}{\pgfqpoint{7.086735in}{1.121561in}}%
\pgfpathquadraticcurveto{\pgfqpoint{7.079818in}{1.121561in}}{\pgfqpoint{7.074829in}{1.124299in}}%
\pgfpathquadraticcurveto{\pgfqpoint{7.069857in}{1.127037in}}{\pgfqpoint{7.066579in}{1.132656in}}%
\pgfpathlineto{\pgfqpoint{7.066579in}{1.132656in}}%
\pgfpathclose%
\pgfpathmoveto{\pgfqpoint{7.101829in}{1.154667in}}%
\pgfpathquadraticcurveto{\pgfqpoint{7.101829in}{1.166087in}}{\pgfqpoint{7.097128in}{1.172589in}}%
\pgfpathquadraticcurveto{\pgfqpoint{7.092427in}{1.179092in}}{\pgfqpoint{7.084213in}{1.179092in}}%
\pgfpathquadraticcurveto{\pgfqpoint{7.075982in}{1.179092in}}{\pgfqpoint{7.071280in}{1.172589in}}%
\pgfpathquadraticcurveto{\pgfqpoint{7.066579in}{1.166087in}}{\pgfqpoint{7.066579in}{1.154667in}}%
\pgfpathquadraticcurveto{\pgfqpoint{7.066579in}{1.143248in}}{\pgfqpoint{7.071280in}{1.136745in}}%
\pgfpathquadraticcurveto{\pgfqpoint{7.075982in}{1.130243in}}{\pgfqpoint{7.084213in}{1.130243in}}%
\pgfpathquadraticcurveto{\pgfqpoint{7.092427in}{1.130243in}}{\pgfqpoint{7.097128in}{1.136745in}}%
\pgfpathquadraticcurveto{\pgfqpoint{7.101829in}{1.143248in}}{\pgfqpoint{7.101829in}{1.154667in}}%
\pgfpathlineto{\pgfqpoint{7.101829in}{1.154667in}}%
\pgfpathclose%
\pgfpathmoveto{\pgfqpoint{7.158405in}{1.154883in}}%
\pgfpathquadraticcurveto{\pgfqpoint{7.145851in}{1.154883in}}{\pgfqpoint{7.141005in}{1.152019in}}%
\pgfpathquadraticcurveto{\pgfqpoint{7.136160in}{1.149156in}}{\pgfqpoint{7.136160in}{1.142221in}}%
\pgfpathquadraticcurveto{\pgfqpoint{7.136160in}{1.136709in}}{\pgfqpoint{7.139799in}{1.133467in}}%
\pgfpathquadraticcurveto{\pgfqpoint{7.143437in}{1.130243in}}{\pgfqpoint{7.149669in}{1.130243in}}%
\pgfpathquadraticcurveto{\pgfqpoint{7.158297in}{1.130243in}}{\pgfqpoint{7.163503in}{1.136349in}}%
\pgfpathquadraticcurveto{\pgfqpoint{7.168708in}{1.142455in}}{\pgfqpoint{7.168708in}{1.152578in}}%
\pgfpathlineto{\pgfqpoint{7.168708in}{1.154883in}}%
\pgfpathlineto{\pgfqpoint{7.158405in}{1.154883in}}%
\pgfpathlineto{\pgfqpoint{7.158405in}{1.154883in}}%
\pgfpathclose%
\pgfpathmoveto{\pgfqpoint{7.179065in}{1.159170in}}%
\pgfpathlineto{\pgfqpoint{7.179065in}{1.123200in}}%
\pgfpathlineto{\pgfqpoint{7.168708in}{1.123200in}}%
\pgfpathlineto{\pgfqpoint{7.168708in}{1.132764in}}%
\pgfpathquadraticcurveto{\pgfqpoint{7.165160in}{1.127037in}}{\pgfqpoint{7.159864in}{1.124299in}}%
\pgfpathquadraticcurveto{\pgfqpoint{7.154569in}{1.121561in}}{\pgfqpoint{7.146913in}{1.121561in}}%
\pgfpathquadraticcurveto{\pgfqpoint{7.137241in}{1.121561in}}{\pgfqpoint{7.131513in}{1.127001in}}%
\pgfpathquadraticcurveto{\pgfqpoint{7.125803in}{1.132440in}}{\pgfqpoint{7.125803in}{1.141554in}}%
\pgfpathquadraticcurveto{\pgfqpoint{7.125803in}{1.152182in}}{\pgfqpoint{7.132918in}{1.157585in}}%
\pgfpathquadraticcurveto{\pgfqpoint{7.140051in}{1.162989in}}{\pgfqpoint{7.154172in}{1.162989in}}%
\pgfpathlineto{\pgfqpoint{7.168708in}{1.162989in}}%
\pgfpathlineto{\pgfqpoint{7.168708in}{1.164016in}}%
\pgfpathquadraticcurveto{\pgfqpoint{7.168708in}{1.171166in}}{\pgfqpoint{7.164007in}{1.175075in}}%
\pgfpathquadraticcurveto{\pgfqpoint{7.159306in}{1.178984in}}{\pgfqpoint{7.150804in}{1.178984in}}%
\pgfpathquadraticcurveto{\pgfqpoint{7.145400in}{1.178984in}}{\pgfqpoint{7.140267in}{1.177687in}}%
\pgfpathquadraticcurveto{\pgfqpoint{7.135151in}{1.176390in}}{\pgfqpoint{7.130432in}{1.173796in}}%
\pgfpathlineto{\pgfqpoint{7.130432in}{1.183379in}}%
\pgfpathquadraticcurveto{\pgfqpoint{7.136106in}{1.185576in}}{\pgfqpoint{7.141456in}{1.186657in}}%
\pgfpathquadraticcurveto{\pgfqpoint{7.146805in}{1.187756in}}{\pgfqpoint{7.151867in}{1.187756in}}%
\pgfpathquadraticcurveto{\pgfqpoint{7.165556in}{1.187756in}}{\pgfqpoint{7.172311in}{1.180659in}}%
\pgfpathquadraticcurveto{\pgfqpoint{7.179065in}{1.173580in}}{\pgfqpoint{7.179065in}{1.159170in}}%
\pgfpathlineto{\pgfqpoint{7.179065in}{1.159170in}}%
\pgfpathclose%
\pgfpathmoveto{\pgfqpoint{7.236920in}{1.176570in}}%
\pgfpathquadraticcurveto{\pgfqpoint{7.235173in}{1.177579in}}{\pgfqpoint{7.233120in}{1.178047in}}%
\pgfpathquadraticcurveto{\pgfqpoint{7.231066in}{1.178533in}}{\pgfqpoint{7.228599in}{1.178533in}}%
\pgfpathquadraticcurveto{\pgfqpoint{7.219809in}{1.178533in}}{\pgfqpoint{7.215107in}{1.172823in}}%
\pgfpathquadraticcurveto{\pgfqpoint{7.210406in}{1.167114in}}{\pgfqpoint{7.210406in}{1.156414in}}%
\pgfpathlineto{\pgfqpoint{7.210406in}{1.123200in}}%
\pgfpathlineto{\pgfqpoint{7.199995in}{1.123200in}}%
\pgfpathlineto{\pgfqpoint{7.199995in}{1.186243in}}%
\pgfpathlineto{\pgfqpoint{7.210406in}{1.186243in}}%
\pgfpathlineto{\pgfqpoint{7.210406in}{1.176444in}}%
\pgfpathquadraticcurveto{\pgfqpoint{7.213684in}{1.182190in}}{\pgfqpoint{7.218908in}{1.184964in}}%
\pgfpathquadraticcurveto{\pgfqpoint{7.224150in}{1.187756in}}{\pgfqpoint{7.231643in}{1.187756in}}%
\pgfpathquadraticcurveto{\pgfqpoint{7.232705in}{1.187756in}}{\pgfqpoint{7.234002in}{1.187611in}}%
\pgfpathquadraticcurveto{\pgfqpoint{7.235299in}{1.187485in}}{\pgfqpoint{7.236866in}{1.187197in}}%
\pgfpathlineto{\pgfqpoint{7.236920in}{1.176570in}}%
\pgfpathlineto{\pgfqpoint{7.236920in}{1.176570in}}%
\pgfpathclose%
\pgfpathmoveto{\pgfqpoint{7.276439in}{1.154883in}}%
\pgfpathquadraticcurveto{\pgfqpoint{7.263884in}{1.154883in}}{\pgfqpoint{7.259039in}{1.152019in}}%
\pgfpathquadraticcurveto{\pgfqpoint{7.254194in}{1.149156in}}{\pgfqpoint{7.254194in}{1.142221in}}%
\pgfpathquadraticcurveto{\pgfqpoint{7.254194in}{1.136709in}}{\pgfqpoint{7.257832in}{1.133467in}}%
\pgfpathquadraticcurveto{\pgfqpoint{7.261471in}{1.130243in}}{\pgfqpoint{7.267703in}{1.130243in}}%
\pgfpathquadraticcurveto{\pgfqpoint{7.276331in}{1.130243in}}{\pgfqpoint{7.281536in}{1.136349in}}%
\pgfpathquadraticcurveto{\pgfqpoint{7.286742in}{1.142455in}}{\pgfqpoint{7.286742in}{1.152578in}}%
\pgfpathlineto{\pgfqpoint{7.286742in}{1.154883in}}%
\pgfpathlineto{\pgfqpoint{7.276439in}{1.154883in}}%
\pgfpathlineto{\pgfqpoint{7.276439in}{1.154883in}}%
\pgfpathclose%
\pgfpathmoveto{\pgfqpoint{7.297099in}{1.159170in}}%
\pgfpathlineto{\pgfqpoint{7.297099in}{1.123200in}}%
\pgfpathlineto{\pgfqpoint{7.286742in}{1.123200in}}%
\pgfpathlineto{\pgfqpoint{7.286742in}{1.132764in}}%
\pgfpathquadraticcurveto{\pgfqpoint{7.283193in}{1.127037in}}{\pgfqpoint{7.277898in}{1.124299in}}%
\pgfpathquadraticcurveto{\pgfqpoint{7.272602in}{1.121561in}}{\pgfqpoint{7.264947in}{1.121561in}}%
\pgfpathquadraticcurveto{\pgfqpoint{7.255275in}{1.121561in}}{\pgfqpoint{7.249547in}{1.127001in}}%
\pgfpathquadraticcurveto{\pgfqpoint{7.243837in}{1.132440in}}{\pgfqpoint{7.243837in}{1.141554in}}%
\pgfpathquadraticcurveto{\pgfqpoint{7.243837in}{1.152182in}}{\pgfqpoint{7.250952in}{1.157585in}}%
\pgfpathquadraticcurveto{\pgfqpoint{7.258084in}{1.162989in}}{\pgfqpoint{7.272206in}{1.162989in}}%
\pgfpathlineto{\pgfqpoint{7.286742in}{1.162989in}}%
\pgfpathlineto{\pgfqpoint{7.286742in}{1.164016in}}%
\pgfpathquadraticcurveto{\pgfqpoint{7.286742in}{1.171166in}}{\pgfqpoint{7.282041in}{1.175075in}}%
\pgfpathquadraticcurveto{\pgfqpoint{7.277339in}{1.178984in}}{\pgfqpoint{7.268838in}{1.178984in}}%
\pgfpathquadraticcurveto{\pgfqpoint{7.263434in}{1.178984in}}{\pgfqpoint{7.258301in}{1.177687in}}%
\pgfpathquadraticcurveto{\pgfqpoint{7.253185in}{1.176390in}}{\pgfqpoint{7.248466in}{1.173796in}}%
\pgfpathlineto{\pgfqpoint{7.248466in}{1.183379in}}%
\pgfpathquadraticcurveto{\pgfqpoint{7.254140in}{1.185576in}}{\pgfqpoint{7.259489in}{1.186657in}}%
\pgfpathquadraticcurveto{\pgfqpoint{7.264839in}{1.187756in}}{\pgfqpoint{7.269900in}{1.187756in}}%
\pgfpathquadraticcurveto{\pgfqpoint{7.283590in}{1.187756in}}{\pgfqpoint{7.290344in}{1.180659in}}%
\pgfpathquadraticcurveto{\pgfqpoint{7.297099in}{1.173580in}}{\pgfqpoint{7.297099in}{1.159170in}}%
\pgfpathlineto{\pgfqpoint{7.297099in}{1.159170in}}%
\pgfpathclose%
\pgfpathmoveto{\pgfqpoint{7.328674in}{1.204147in}}%
\pgfpathlineto{\pgfqpoint{7.328674in}{1.186243in}}%
\pgfpathlineto{\pgfqpoint{7.350000in}{1.186243in}}%
\pgfpathlineto{\pgfqpoint{7.350000in}{1.178191in}}%
\pgfpathlineto{\pgfqpoint{7.328674in}{1.178191in}}%
\pgfpathlineto{\pgfqpoint{7.328674in}{1.143968in}}%
\pgfpathquadraticcurveto{\pgfqpoint{7.328674in}{1.136259in}}{\pgfqpoint{7.330781in}{1.134061in}}%
\pgfpathquadraticcurveto{\pgfqpoint{7.332889in}{1.131864in}}{\pgfqpoint{7.339373in}{1.131864in}}%
\pgfpathlineto{\pgfqpoint{7.350000in}{1.131864in}}%
\pgfpathlineto{\pgfqpoint{7.350000in}{1.123200in}}%
\pgfpathlineto{\pgfqpoint{7.339373in}{1.123200in}}%
\pgfpathquadraticcurveto{\pgfqpoint{7.327377in}{1.123200in}}{\pgfqpoint{7.322820in}{1.127667in}}%
\pgfpathquadraticcurveto{\pgfqpoint{7.318263in}{1.132152in}}{\pgfqpoint{7.318263in}{1.143968in}}%
\pgfpathlineto{\pgfqpoint{7.318263in}{1.178191in}}%
\pgfpathlineto{\pgfqpoint{7.310662in}{1.178191in}}%
\pgfpathlineto{\pgfqpoint{7.310662in}{1.186243in}}%
\pgfpathlineto{\pgfqpoint{7.318263in}{1.186243in}}%
\pgfpathlineto{\pgfqpoint{7.318263in}{1.204147in}}%
\pgfpathlineto{\pgfqpoint{7.328674in}{1.204147in}}%
\pgfpathlineto{\pgfqpoint{7.328674in}{1.204147in}}%
\pgfpathclose%
\pgfpathmoveto{\pgfqpoint{7.417546in}{1.157315in}}%
\pgfpathlineto{\pgfqpoint{7.417546in}{1.152254in}}%
\pgfpathlineto{\pgfqpoint{7.369922in}{1.152254in}}%
\pgfpathquadraticcurveto{\pgfqpoint{7.370606in}{1.141554in}}{\pgfqpoint{7.376370in}{1.135953in}}%
\pgfpathquadraticcurveto{\pgfqpoint{7.382134in}{1.130351in}}{\pgfqpoint{7.392437in}{1.130351in}}%
\pgfpathquadraticcurveto{\pgfqpoint{7.398399in}{1.130351in}}{\pgfqpoint{7.404001in}{1.131810in}}%
\pgfpathquadraticcurveto{\pgfqpoint{7.409603in}{1.133269in}}{\pgfqpoint{7.415132in}{1.136205in}}%
\pgfpathlineto{\pgfqpoint{7.415132in}{1.126406in}}%
\pgfpathquadraticcurveto{\pgfqpoint{7.409549in}{1.124047in}}{\pgfqpoint{7.403695in}{1.122804in}}%
\pgfpathquadraticcurveto{\pgfqpoint{7.397841in}{1.121561in}}{\pgfqpoint{7.391825in}{1.121561in}}%
\pgfpathquadraticcurveto{\pgfqpoint{7.376730in}{1.121561in}}{\pgfqpoint{7.367923in}{1.130333in}}%
\pgfpathquadraticcurveto{\pgfqpoint{7.359115in}{1.139123in}}{\pgfqpoint{7.359115in}{1.154109in}}%
\pgfpathquadraticcurveto{\pgfqpoint{7.359115in}{1.169581in}}{\pgfqpoint{7.367472in}{1.178659in}}%
\pgfpathquadraticcurveto{\pgfqpoint{7.375830in}{1.187756in}}{\pgfqpoint{7.390023in}{1.187756in}}%
\pgfpathquadraticcurveto{\pgfqpoint{7.402740in}{1.187756in}}{\pgfqpoint{7.410143in}{1.179560in}}%
\pgfpathquadraticcurveto{\pgfqpoint{7.417546in}{1.171383in}}{\pgfqpoint{7.417546in}{1.157315in}}%
\pgfpathlineto{\pgfqpoint{7.417546in}{1.157315in}}%
\pgfpathclose%
\pgfpathmoveto{\pgfqpoint{7.407189in}{1.160359in}}%
\pgfpathquadraticcurveto{\pgfqpoint{7.407081in}{1.168843in}}{\pgfqpoint{7.402434in}{1.173904in}}%
\pgfpathquadraticcurveto{\pgfqpoint{7.397787in}{1.178984in}}{\pgfqpoint{7.390132in}{1.178984in}}%
\pgfpathquadraticcurveto{\pgfqpoint{7.381468in}{1.178984in}}{\pgfqpoint{7.376262in}{1.174084in}}%
\pgfpathquadraticcurveto{\pgfqpoint{7.371057in}{1.169185in}}{\pgfqpoint{7.370264in}{1.160287in}}%
\pgfpathlineto{\pgfqpoint{7.407189in}{1.160359in}}%
\pgfpathlineto{\pgfqpoint{7.407189in}{1.160359in}}%
\pgfpathclose%
\pgfpathmoveto{\pgfqpoint{7.434549in}{1.210793in}}%
\pgfpathlineto{\pgfqpoint{7.444906in}{1.210793in}}%
\pgfpathlineto{\pgfqpoint{7.444906in}{1.123200in}}%
\pgfpathlineto{\pgfqpoint{7.434549in}{1.123200in}}%
\pgfpathlineto{\pgfqpoint{7.434549in}{1.210793in}}%
\pgfpathlineto{\pgfqpoint{7.434549in}{1.210793in}}%
\pgfpathclose%
\pgfpathmoveto{\pgfqpoint{7.492801in}{1.117346in}}%
\pgfpathquadraticcurveto{\pgfqpoint{7.488424in}{1.106088in}}{\pgfqpoint{7.484245in}{1.102666in}}%
\pgfpathquadraticcurveto{\pgfqpoint{7.480084in}{1.099226in}}{\pgfqpoint{7.473113in}{1.099226in}}%
\pgfpathlineto{\pgfqpoint{7.464828in}{1.099226in}}%
\pgfpathlineto{\pgfqpoint{7.464828in}{1.107890in}}%
\pgfpathlineto{\pgfqpoint{7.470916in}{1.107890in}}%
\pgfpathquadraticcurveto{\pgfqpoint{7.475185in}{1.107890in}}{\pgfqpoint{7.477544in}{1.109925in}}%
\pgfpathquadraticcurveto{\pgfqpoint{7.479922in}{1.111942in}}{\pgfqpoint{7.482786in}{1.119489in}}%
\pgfpathlineto{\pgfqpoint{7.484641in}{1.124209in}}%
\pgfpathlineto{\pgfqpoint{7.459154in}{1.186243in}}%
\pgfpathlineto{\pgfqpoint{7.470123in}{1.186243in}}%
\pgfpathlineto{\pgfqpoint{7.489829in}{1.136943in}}%
\pgfpathlineto{\pgfqpoint{7.509534in}{1.186243in}}%
\pgfpathlineto{\pgfqpoint{7.520503in}{1.186243in}}%
\pgfpathlineto{\pgfqpoint{7.492801in}{1.117346in}}%
\pgfpathlineto{\pgfqpoint{7.492801in}{1.117346in}}%
\pgfpathclose%
\pgfpathmoveto{\pgfqpoint{7.519837in}{1.137502in}}%
\pgfpathlineto{\pgfqpoint{7.531725in}{1.137502in}}%
\pgfpathlineto{\pgfqpoint{7.531725in}{1.123200in}}%
\pgfpathlineto{\pgfqpoint{7.519837in}{1.123200in}}%
\pgfpathlineto{\pgfqpoint{7.519837in}{1.137502in}}%
\pgfpathlineto{\pgfqpoint{7.519837in}{1.137502in}}%
\pgfpathclose%
\pgfusepath{fill}%
\end{pgfscope}%
\begin{pgfscope}%
\pgfsetbuttcap%
\pgfsetroundjoin%
\definecolor{currentfill}{rgb}{0.309804,0.372549,0.419608}%
\pgfsetfillcolor{currentfill}%
\pgfsetlinewidth{0.000000pt}%
\definecolor{currentstroke}{rgb}{0.309804,0.372549,0.419608}%
\pgfsetstrokecolor{currentstroke}%
\pgfsetdash{}{0pt}%
\pgfpathmoveto{\pgfqpoint{7.640008in}{1.197896in}}%
\pgfpathlineto{\pgfqpoint{7.640008in}{1.166321in}}%
\pgfpathlineto{\pgfqpoint{7.654310in}{1.166321in}}%
\pgfpathquadraticcurveto{\pgfqpoint{7.662253in}{1.166321in}}{\pgfqpoint{7.666576in}{1.170428in}}%
\pgfpathquadraticcurveto{\pgfqpoint{7.670917in}{1.174535in}}{\pgfqpoint{7.670917in}{1.182136in}}%
\pgfpathquadraticcurveto{\pgfqpoint{7.670917in}{1.189683in}}{\pgfqpoint{7.666576in}{1.193790in}}%
\pgfpathquadraticcurveto{\pgfqpoint{7.662253in}{1.197896in}}{\pgfqpoint{7.654310in}{1.197896in}}%
\pgfpathlineto{\pgfqpoint{7.640008in}{1.197896in}}%
\pgfpathlineto{\pgfqpoint{7.640008in}{1.197896in}}%
\pgfpathclose%
\pgfpathmoveto{\pgfqpoint{7.628642in}{1.207245in}}%
\pgfpathlineto{\pgfqpoint{7.654310in}{1.207245in}}%
\pgfpathquadraticcurveto{\pgfqpoint{7.668449in}{1.207245in}}{\pgfqpoint{7.675672in}{1.200850in}}%
\pgfpathquadraticcurveto{\pgfqpoint{7.682913in}{1.194456in}}{\pgfqpoint{7.682913in}{1.182136in}}%
\pgfpathquadraticcurveto{\pgfqpoint{7.682913in}{1.169689in}}{\pgfqpoint{7.675672in}{1.163331in}}%
\pgfpathquadraticcurveto{\pgfqpoint{7.668449in}{1.156973in}}{\pgfqpoint{7.654310in}{1.156973in}}%
\pgfpathlineto{\pgfqpoint{7.640008in}{1.156973in}}%
\pgfpathlineto{\pgfqpoint{7.640008in}{1.123200in}}%
\pgfpathlineto{\pgfqpoint{7.628642in}{1.123200in}}%
\pgfpathlineto{\pgfqpoint{7.628642in}{1.207245in}}%
\pgfpathlineto{\pgfqpoint{7.628642in}{1.207245in}}%
\pgfpathclose%
\pgfpathmoveto{\pgfqpoint{7.732212in}{1.176570in}}%
\pgfpathquadraticcurveto{\pgfqpoint{7.730465in}{1.177579in}}{\pgfqpoint{7.728412in}{1.178047in}}%
\pgfpathquadraticcurveto{\pgfqpoint{7.726358in}{1.178533in}}{\pgfqpoint{7.723891in}{1.178533in}}%
\pgfpathquadraticcurveto{\pgfqpoint{7.715101in}{1.178533in}}{\pgfqpoint{7.710400in}{1.172823in}}%
\pgfpathquadraticcurveto{\pgfqpoint{7.705698in}{1.167114in}}{\pgfqpoint{7.705698in}{1.156414in}}%
\pgfpathlineto{\pgfqpoint{7.705698in}{1.123200in}}%
\pgfpathlineto{\pgfqpoint{7.695287in}{1.123200in}}%
\pgfpathlineto{\pgfqpoint{7.695287in}{1.186243in}}%
\pgfpathlineto{\pgfqpoint{7.705698in}{1.186243in}}%
\pgfpathlineto{\pgfqpoint{7.705698in}{1.176444in}}%
\pgfpathquadraticcurveto{\pgfqpoint{7.708977in}{1.182190in}}{\pgfqpoint{7.714200in}{1.184964in}}%
\pgfpathquadraticcurveto{\pgfqpoint{7.719442in}{1.187756in}}{\pgfqpoint{7.726935in}{1.187756in}}%
\pgfpathquadraticcurveto{\pgfqpoint{7.727997in}{1.187756in}}{\pgfqpoint{7.729294in}{1.187611in}}%
\pgfpathquadraticcurveto{\pgfqpoint{7.730591in}{1.187485in}}{\pgfqpoint{7.732158in}{1.187197in}}%
\pgfpathlineto{\pgfqpoint{7.732212in}{1.176570in}}%
\pgfpathlineto{\pgfqpoint{7.732212in}{1.176570in}}%
\pgfpathclose%
\pgfpathmoveto{\pgfqpoint{7.794462in}{1.157315in}}%
\pgfpathlineto{\pgfqpoint{7.794462in}{1.152254in}}%
\pgfpathlineto{\pgfqpoint{7.746838in}{1.152254in}}%
\pgfpathquadraticcurveto{\pgfqpoint{7.747523in}{1.141554in}}{\pgfqpoint{7.753286in}{1.135953in}}%
\pgfpathquadraticcurveto{\pgfqpoint{7.759050in}{1.130351in}}{\pgfqpoint{7.769353in}{1.130351in}}%
\pgfpathquadraticcurveto{\pgfqpoint{7.775315in}{1.130351in}}{\pgfqpoint{7.780917in}{1.131810in}}%
\pgfpathquadraticcurveto{\pgfqpoint{7.786519in}{1.133269in}}{\pgfqpoint{7.792049in}{1.136205in}}%
\pgfpathlineto{\pgfqpoint{7.792049in}{1.126406in}}%
\pgfpathquadraticcurveto{\pgfqpoint{7.786465in}{1.124047in}}{\pgfqpoint{7.780611in}{1.122804in}}%
\pgfpathquadraticcurveto{\pgfqpoint{7.774757in}{1.121561in}}{\pgfqpoint{7.768741in}{1.121561in}}%
\pgfpathquadraticcurveto{\pgfqpoint{7.753647in}{1.121561in}}{\pgfqpoint{7.744839in}{1.130333in}}%
\pgfpathquadraticcurveto{\pgfqpoint{7.736031in}{1.139123in}}{\pgfqpoint{7.736031in}{1.154109in}}%
\pgfpathquadraticcurveto{\pgfqpoint{7.736031in}{1.169581in}}{\pgfqpoint{7.744388in}{1.178659in}}%
\pgfpathquadraticcurveto{\pgfqpoint{7.752746in}{1.187756in}}{\pgfqpoint{7.766940in}{1.187756in}}%
\pgfpathquadraticcurveto{\pgfqpoint{7.779656in}{1.187756in}}{\pgfqpoint{7.787059in}{1.179560in}}%
\pgfpathquadraticcurveto{\pgfqpoint{7.794462in}{1.171383in}}{\pgfqpoint{7.794462in}{1.157315in}}%
\pgfpathlineto{\pgfqpoint{7.794462in}{1.157315in}}%
\pgfpathclose%
\pgfpathmoveto{\pgfqpoint{7.784105in}{1.160359in}}%
\pgfpathquadraticcurveto{\pgfqpoint{7.783997in}{1.168843in}}{\pgfqpoint{7.779350in}{1.173904in}}%
\pgfpathquadraticcurveto{\pgfqpoint{7.774703in}{1.178984in}}{\pgfqpoint{7.767048in}{1.178984in}}%
\pgfpathquadraticcurveto{\pgfqpoint{7.758384in}{1.178984in}}{\pgfqpoint{7.753178in}{1.174084in}}%
\pgfpathquadraticcurveto{\pgfqpoint{7.747973in}{1.169185in}}{\pgfqpoint{7.747180in}{1.160287in}}%
\pgfpathlineto{\pgfqpoint{7.784105in}{1.160359in}}%
\pgfpathlineto{\pgfqpoint{7.784105in}{1.160359in}}%
\pgfpathclose%
\pgfusepath{fill}%
\end{pgfscope}%
\begin{pgfscope}%
\pgfsetbuttcap%
\pgfsetroundjoin%
\definecolor{currentfill}{rgb}{0.309804,0.372549,0.419608}%
\pgfsetfillcolor{currentfill}%
\pgfsetlinewidth{0.000000pt}%
\definecolor{currentstroke}{rgb}{0.309804,0.372549,0.419608}%
\pgfsetstrokecolor{currentstroke}%
\pgfsetdash{}{0pt}%
\pgfpathmoveto{\pgfqpoint{7.901336in}{1.154883in}}%
\pgfpathquadraticcurveto{\pgfqpoint{7.888782in}{1.154883in}}{\pgfqpoint{7.883937in}{1.152019in}}%
\pgfpathquadraticcurveto{\pgfqpoint{7.879091in}{1.149156in}}{\pgfqpoint{7.879091in}{1.142221in}}%
\pgfpathquadraticcurveto{\pgfqpoint{7.879091in}{1.136709in}}{\pgfqpoint{7.882730in}{1.133467in}}%
\pgfpathquadraticcurveto{\pgfqpoint{7.886368in}{1.130243in}}{\pgfqpoint{7.892600in}{1.130243in}}%
\pgfpathquadraticcurveto{\pgfqpoint{7.901228in}{1.130243in}}{\pgfqpoint{7.906434in}{1.136349in}}%
\pgfpathquadraticcurveto{\pgfqpoint{7.911639in}{1.142455in}}{\pgfqpoint{7.911639in}{1.152578in}}%
\pgfpathlineto{\pgfqpoint{7.911639in}{1.154883in}}%
\pgfpathlineto{\pgfqpoint{7.901336in}{1.154883in}}%
\pgfpathlineto{\pgfqpoint{7.901336in}{1.154883in}}%
\pgfpathclose%
\pgfpathmoveto{\pgfqpoint{7.921996in}{1.159170in}}%
\pgfpathlineto{\pgfqpoint{7.921996in}{1.123200in}}%
\pgfpathlineto{\pgfqpoint{7.911639in}{1.123200in}}%
\pgfpathlineto{\pgfqpoint{7.911639in}{1.132764in}}%
\pgfpathquadraticcurveto{\pgfqpoint{7.908091in}{1.127037in}}{\pgfqpoint{7.902795in}{1.124299in}}%
\pgfpathquadraticcurveto{\pgfqpoint{7.897500in}{1.121561in}}{\pgfqpoint{7.889845in}{1.121561in}}%
\pgfpathquadraticcurveto{\pgfqpoint{7.880172in}{1.121561in}}{\pgfqpoint{7.874444in}{1.127001in}}%
\pgfpathquadraticcurveto{\pgfqpoint{7.868734in}{1.132440in}}{\pgfqpoint{7.868734in}{1.141554in}}%
\pgfpathquadraticcurveto{\pgfqpoint{7.868734in}{1.152182in}}{\pgfqpoint{7.875849in}{1.157585in}}%
\pgfpathquadraticcurveto{\pgfqpoint{7.882982in}{1.162989in}}{\pgfqpoint{7.897104in}{1.162989in}}%
\pgfpathlineto{\pgfqpoint{7.911639in}{1.162989in}}%
\pgfpathlineto{\pgfqpoint{7.911639in}{1.164016in}}%
\pgfpathquadraticcurveto{\pgfqpoint{7.911639in}{1.171166in}}{\pgfqpoint{7.906938in}{1.175075in}}%
\pgfpathquadraticcurveto{\pgfqpoint{7.902237in}{1.178984in}}{\pgfqpoint{7.893735in}{1.178984in}}%
\pgfpathquadraticcurveto{\pgfqpoint{7.888332in}{1.178984in}}{\pgfqpoint{7.883198in}{1.177687in}}%
\pgfpathquadraticcurveto{\pgfqpoint{7.878083in}{1.176390in}}{\pgfqpoint{7.873363in}{1.173796in}}%
\pgfpathlineto{\pgfqpoint{7.873363in}{1.183379in}}%
\pgfpathquadraticcurveto{\pgfqpoint{7.879037in}{1.185576in}}{\pgfqpoint{7.884387in}{1.186657in}}%
\pgfpathquadraticcurveto{\pgfqpoint{7.889737in}{1.187756in}}{\pgfqpoint{7.894798in}{1.187756in}}%
\pgfpathquadraticcurveto{\pgfqpoint{7.908487in}{1.187756in}}{\pgfqpoint{7.915242in}{1.180659in}}%
\pgfpathquadraticcurveto{\pgfqpoint{7.921996in}{1.173580in}}{\pgfqpoint{7.921996in}{1.159170in}}%
\pgfpathlineto{\pgfqpoint{7.921996in}{1.159170in}}%
\pgfpathclose%
\pgfpathmoveto{\pgfqpoint{7.935902in}{1.186243in}}%
\pgfpathlineto{\pgfqpoint{7.946871in}{1.186243in}}%
\pgfpathlineto{\pgfqpoint{7.966576in}{1.133341in}}%
\pgfpathlineto{\pgfqpoint{7.986282in}{1.186243in}}%
\pgfpathlineto{\pgfqpoint{7.997251in}{1.186243in}}%
\pgfpathlineto{\pgfqpoint{7.973601in}{1.123200in}}%
\pgfpathlineto{\pgfqpoint{7.959534in}{1.123200in}}%
\pgfpathlineto{\pgfqpoint{7.935902in}{1.186243in}}%
\pgfpathlineto{\pgfqpoint{7.935902in}{1.186243in}}%
\pgfpathclose%
\pgfpathmoveto{\pgfqpoint{8.065481in}{1.157315in}}%
\pgfpathlineto{\pgfqpoint{8.065481in}{1.152254in}}%
\pgfpathlineto{\pgfqpoint{8.017857in}{1.152254in}}%
\pgfpathquadraticcurveto{\pgfqpoint{8.018541in}{1.141554in}}{\pgfqpoint{8.024305in}{1.135953in}}%
\pgfpathquadraticcurveto{\pgfqpoint{8.030069in}{1.130351in}}{\pgfqpoint{8.040372in}{1.130351in}}%
\pgfpathquadraticcurveto{\pgfqpoint{8.046334in}{1.130351in}}{\pgfqpoint{8.051936in}{1.131810in}}%
\pgfpathquadraticcurveto{\pgfqpoint{8.057538in}{1.133269in}}{\pgfqpoint{8.063067in}{1.136205in}}%
\pgfpathlineto{\pgfqpoint{8.063067in}{1.126406in}}%
\pgfpathquadraticcurveto{\pgfqpoint{8.057484in}{1.124047in}}{\pgfqpoint{8.051630in}{1.122804in}}%
\pgfpathquadraticcurveto{\pgfqpoint{8.045776in}{1.121561in}}{\pgfqpoint{8.039760in}{1.121561in}}%
\pgfpathquadraticcurveto{\pgfqpoint{8.024666in}{1.121561in}}{\pgfqpoint{8.015858in}{1.130333in}}%
\pgfpathquadraticcurveto{\pgfqpoint{8.007050in}{1.139123in}}{\pgfqpoint{8.007050in}{1.154109in}}%
\pgfpathquadraticcurveto{\pgfqpoint{8.007050in}{1.169581in}}{\pgfqpoint{8.015407in}{1.178659in}}%
\pgfpathquadraticcurveto{\pgfqpoint{8.023765in}{1.187756in}}{\pgfqpoint{8.037959in}{1.187756in}}%
\pgfpathquadraticcurveto{\pgfqpoint{8.050675in}{1.187756in}}{\pgfqpoint{8.058078in}{1.179560in}}%
\pgfpathquadraticcurveto{\pgfqpoint{8.065481in}{1.171383in}}{\pgfqpoint{8.065481in}{1.157315in}}%
\pgfpathlineto{\pgfqpoint{8.065481in}{1.157315in}}%
\pgfpathclose%
\pgfpathmoveto{\pgfqpoint{8.055124in}{1.160359in}}%
\pgfpathquadraticcurveto{\pgfqpoint{8.055016in}{1.168843in}}{\pgfqpoint{8.050369in}{1.173904in}}%
\pgfpathquadraticcurveto{\pgfqpoint{8.045722in}{1.178984in}}{\pgfqpoint{8.038067in}{1.178984in}}%
\pgfpathquadraticcurveto{\pgfqpoint{8.029403in}{1.178984in}}{\pgfqpoint{8.024197in}{1.174084in}}%
\pgfpathquadraticcurveto{\pgfqpoint{8.018992in}{1.169185in}}{\pgfqpoint{8.018199in}{1.160287in}}%
\pgfpathlineto{\pgfqpoint{8.055124in}{1.160359in}}%
\pgfpathlineto{\pgfqpoint{8.055124in}{1.160359in}}%
\pgfpathclose%
\pgfpathmoveto{\pgfqpoint{8.119013in}{1.176570in}}%
\pgfpathquadraticcurveto{\pgfqpoint{8.117266in}{1.177579in}}{\pgfqpoint{8.115213in}{1.178047in}}%
\pgfpathquadraticcurveto{\pgfqpoint{8.113159in}{1.178533in}}{\pgfqpoint{8.110692in}{1.178533in}}%
\pgfpathquadraticcurveto{\pgfqpoint{8.101902in}{1.178533in}}{\pgfqpoint{8.097201in}{1.172823in}}%
\pgfpathquadraticcurveto{\pgfqpoint{8.092499in}{1.167114in}}{\pgfqpoint{8.092499in}{1.156414in}}%
\pgfpathlineto{\pgfqpoint{8.092499in}{1.123200in}}%
\pgfpathlineto{\pgfqpoint{8.082088in}{1.123200in}}%
\pgfpathlineto{\pgfqpoint{8.082088in}{1.186243in}}%
\pgfpathlineto{\pgfqpoint{8.092499in}{1.186243in}}%
\pgfpathlineto{\pgfqpoint{8.092499in}{1.176444in}}%
\pgfpathquadraticcurveto{\pgfqpoint{8.095778in}{1.182190in}}{\pgfqpoint{8.101001in}{1.184964in}}%
\pgfpathquadraticcurveto{\pgfqpoint{8.106243in}{1.187756in}}{\pgfqpoint{8.113736in}{1.187756in}}%
\pgfpathquadraticcurveto{\pgfqpoint{8.114798in}{1.187756in}}{\pgfqpoint{8.116095in}{1.187611in}}%
\pgfpathquadraticcurveto{\pgfqpoint{8.117392in}{1.187485in}}{\pgfqpoint{8.118959in}{1.187197in}}%
\pgfpathlineto{\pgfqpoint{8.119013in}{1.176570in}}%
\pgfpathlineto{\pgfqpoint{8.119013in}{1.176570in}}%
\pgfpathclose%
\pgfpathmoveto{\pgfqpoint{8.158532in}{1.154883in}}%
\pgfpathquadraticcurveto{\pgfqpoint{8.145977in}{1.154883in}}{\pgfqpoint{8.141132in}{1.152019in}}%
\pgfpathquadraticcurveto{\pgfqpoint{8.136287in}{1.149156in}}{\pgfqpoint{8.136287in}{1.142221in}}%
\pgfpathquadraticcurveto{\pgfqpoint{8.136287in}{1.136709in}}{\pgfqpoint{8.139925in}{1.133467in}}%
\pgfpathquadraticcurveto{\pgfqpoint{8.143564in}{1.130243in}}{\pgfqpoint{8.149796in}{1.130243in}}%
\pgfpathquadraticcurveto{\pgfqpoint{8.158424in}{1.130243in}}{\pgfqpoint{8.163629in}{1.136349in}}%
\pgfpathquadraticcurveto{\pgfqpoint{8.168835in}{1.142455in}}{\pgfqpoint{8.168835in}{1.152578in}}%
\pgfpathlineto{\pgfqpoint{8.168835in}{1.154883in}}%
\pgfpathlineto{\pgfqpoint{8.158532in}{1.154883in}}%
\pgfpathlineto{\pgfqpoint{8.158532in}{1.154883in}}%
\pgfpathclose%
\pgfpathmoveto{\pgfqpoint{8.179192in}{1.159170in}}%
\pgfpathlineto{\pgfqpoint{8.179192in}{1.123200in}}%
\pgfpathlineto{\pgfqpoint{8.168835in}{1.123200in}}%
\pgfpathlineto{\pgfqpoint{8.168835in}{1.132764in}}%
\pgfpathquadraticcurveto{\pgfqpoint{8.165286in}{1.127037in}}{\pgfqpoint{8.159991in}{1.124299in}}%
\pgfpathquadraticcurveto{\pgfqpoint{8.154695in}{1.121561in}}{\pgfqpoint{8.147040in}{1.121561in}}%
\pgfpathquadraticcurveto{\pgfqpoint{8.137368in}{1.121561in}}{\pgfqpoint{8.131640in}{1.127001in}}%
\pgfpathquadraticcurveto{\pgfqpoint{8.125930in}{1.132440in}}{\pgfqpoint{8.125930in}{1.141554in}}%
\pgfpathquadraticcurveto{\pgfqpoint{8.125930in}{1.152182in}}{\pgfqpoint{8.133045in}{1.157585in}}%
\pgfpathquadraticcurveto{\pgfqpoint{8.140178in}{1.162989in}}{\pgfqpoint{8.154299in}{1.162989in}}%
\pgfpathlineto{\pgfqpoint{8.168835in}{1.162989in}}%
\pgfpathlineto{\pgfqpoint{8.168835in}{1.164016in}}%
\pgfpathquadraticcurveto{\pgfqpoint{8.168835in}{1.171166in}}{\pgfqpoint{8.164134in}{1.175075in}}%
\pgfpathquadraticcurveto{\pgfqpoint{8.159433in}{1.178984in}}{\pgfqpoint{8.150931in}{1.178984in}}%
\pgfpathquadraticcurveto{\pgfqpoint{8.145527in}{1.178984in}}{\pgfqpoint{8.140394in}{1.177687in}}%
\pgfpathquadraticcurveto{\pgfqpoint{8.135278in}{1.176390in}}{\pgfqpoint{8.130559in}{1.173796in}}%
\pgfpathlineto{\pgfqpoint{8.130559in}{1.183379in}}%
\pgfpathquadraticcurveto{\pgfqpoint{8.136233in}{1.185576in}}{\pgfqpoint{8.141582in}{1.186657in}}%
\pgfpathquadraticcurveto{\pgfqpoint{8.146932in}{1.187756in}}{\pgfqpoint{8.151993in}{1.187756in}}%
\pgfpathquadraticcurveto{\pgfqpoint{8.165683in}{1.187756in}}{\pgfqpoint{8.172437in}{1.180659in}}%
\pgfpathquadraticcurveto{\pgfqpoint{8.179192in}{1.173580in}}{\pgfqpoint{8.179192in}{1.159170in}}%
\pgfpathlineto{\pgfqpoint{8.179192in}{1.159170in}}%
\pgfpathclose%
\pgfpathmoveto{\pgfqpoint{8.242000in}{1.155460in}}%
\pgfpathquadraticcurveto{\pgfqpoint{8.242000in}{1.166717in}}{\pgfqpoint{8.237353in}{1.172896in}}%
\pgfpathquadraticcurveto{\pgfqpoint{8.232724in}{1.179092in}}{\pgfqpoint{8.224330in}{1.179092in}}%
\pgfpathquadraticcurveto{\pgfqpoint{8.216009in}{1.179092in}}{\pgfqpoint{8.211362in}{1.172896in}}%
\pgfpathquadraticcurveto{\pgfqpoint{8.206714in}{1.166717in}}{\pgfqpoint{8.206714in}{1.155460in}}%
\pgfpathquadraticcurveto{\pgfqpoint{8.206714in}{1.144256in}}{\pgfqpoint{8.211362in}{1.138060in}}%
\pgfpathquadraticcurveto{\pgfqpoint{8.216009in}{1.131864in}}{\pgfqpoint{8.224330in}{1.131864in}}%
\pgfpathquadraticcurveto{\pgfqpoint{8.232724in}{1.131864in}}{\pgfqpoint{8.237353in}{1.138060in}}%
\pgfpathquadraticcurveto{\pgfqpoint{8.242000in}{1.144256in}}{\pgfqpoint{8.242000in}{1.155460in}}%
\pgfpathlineto{\pgfqpoint{8.242000in}{1.155460in}}%
\pgfpathclose%
\pgfpathmoveto{\pgfqpoint{8.252357in}{1.131017in}}%
\pgfpathquadraticcurveto{\pgfqpoint{8.252357in}{1.114932in}}{\pgfqpoint{8.245206in}{1.107079in}}%
\pgfpathquadraticcurveto{\pgfqpoint{8.238074in}{1.099226in}}{\pgfqpoint{8.223322in}{1.099226in}}%
\pgfpathquadraticcurveto{\pgfqpoint{8.217864in}{1.099226in}}{\pgfqpoint{8.213019in}{1.100036in}}%
\pgfpathquadraticcurveto{\pgfqpoint{8.208173in}{1.100847in}}{\pgfqpoint{8.203616in}{1.102540in}}%
\pgfpathlineto{\pgfqpoint{8.203616in}{1.112609in}}%
\pgfpathquadraticcurveto{\pgfqpoint{8.208173in}{1.110141in}}{\pgfqpoint{8.212622in}{1.108970in}}%
\pgfpathquadraticcurveto{\pgfqpoint{8.217071in}{1.107782in}}{\pgfqpoint{8.221683in}{1.107782in}}%
\pgfpathquadraticcurveto{\pgfqpoint{8.231877in}{1.107782in}}{\pgfqpoint{8.236939in}{1.113095in}}%
\pgfpathquadraticcurveto{\pgfqpoint{8.242000in}{1.118409in}}{\pgfqpoint{8.242000in}{1.129162in}}%
\pgfpathlineto{\pgfqpoint{8.242000in}{1.134295in}}%
\pgfpathquadraticcurveto{\pgfqpoint{8.238794in}{1.128712in}}{\pgfqpoint{8.233787in}{1.125956in}}%
\pgfpathquadraticcurveto{\pgfqpoint{8.228779in}{1.123200in}}{\pgfqpoint{8.221791in}{1.123200in}}%
\pgfpathquadraticcurveto{\pgfqpoint{8.210209in}{1.123200in}}{\pgfqpoint{8.203112in}{1.132026in}}%
\pgfpathquadraticcurveto{\pgfqpoint{8.196015in}{1.140870in}}{\pgfqpoint{8.196015in}{1.155460in}}%
\pgfpathquadraticcurveto{\pgfqpoint{8.196015in}{1.170086in}}{\pgfqpoint{8.203112in}{1.178912in}}%
\pgfpathquadraticcurveto{\pgfqpoint{8.210209in}{1.187756in}}{\pgfqpoint{8.221791in}{1.187756in}}%
\pgfpathquadraticcurveto{\pgfqpoint{8.228779in}{1.187756in}}{\pgfqpoint{8.233787in}{1.185000in}}%
\pgfpathquadraticcurveto{\pgfqpoint{8.238794in}{1.182244in}}{\pgfqpoint{8.242000in}{1.176678in}}%
\pgfpathlineto{\pgfqpoint{8.242000in}{1.186243in}}%
\pgfpathlineto{\pgfqpoint{8.252357in}{1.186243in}}%
\pgfpathlineto{\pgfqpoint{8.252357in}{1.131017in}}%
\pgfpathlineto{\pgfqpoint{8.252357in}{1.131017in}}%
\pgfpathclose%
\pgfpathmoveto{\pgfqpoint{8.327630in}{1.157315in}}%
\pgfpathlineto{\pgfqpoint{8.327630in}{1.152254in}}%
\pgfpathlineto{\pgfqpoint{8.280006in}{1.152254in}}%
\pgfpathquadraticcurveto{\pgfqpoint{8.280690in}{1.141554in}}{\pgfqpoint{8.286454in}{1.135953in}}%
\pgfpathquadraticcurveto{\pgfqpoint{8.292218in}{1.130351in}}{\pgfqpoint{8.302521in}{1.130351in}}%
\pgfpathquadraticcurveto{\pgfqpoint{8.308483in}{1.130351in}}{\pgfqpoint{8.314085in}{1.131810in}}%
\pgfpathquadraticcurveto{\pgfqpoint{8.319687in}{1.133269in}}{\pgfqpoint{8.325216in}{1.136205in}}%
\pgfpathlineto{\pgfqpoint{8.325216in}{1.126406in}}%
\pgfpathquadraticcurveto{\pgfqpoint{8.319633in}{1.124047in}}{\pgfqpoint{8.313779in}{1.122804in}}%
\pgfpathquadraticcurveto{\pgfqpoint{8.307925in}{1.121561in}}{\pgfqpoint{8.301909in}{1.121561in}}%
\pgfpathquadraticcurveto{\pgfqpoint{8.286814in}{1.121561in}}{\pgfqpoint{8.278007in}{1.130333in}}%
\pgfpathquadraticcurveto{\pgfqpoint{8.269199in}{1.139123in}}{\pgfqpoint{8.269199in}{1.154109in}}%
\pgfpathquadraticcurveto{\pgfqpoint{8.269199in}{1.169581in}}{\pgfqpoint{8.277556in}{1.178659in}}%
\pgfpathquadraticcurveto{\pgfqpoint{8.285914in}{1.187756in}}{\pgfqpoint{8.300107in}{1.187756in}}%
\pgfpathquadraticcurveto{\pgfqpoint{8.312824in}{1.187756in}}{\pgfqpoint{8.320227in}{1.179560in}}%
\pgfpathquadraticcurveto{\pgfqpoint{8.327630in}{1.171383in}}{\pgfqpoint{8.327630in}{1.157315in}}%
\pgfpathlineto{\pgfqpoint{8.327630in}{1.157315in}}%
\pgfpathclose%
\pgfpathmoveto{\pgfqpoint{8.317273in}{1.160359in}}%
\pgfpathquadraticcurveto{\pgfqpoint{8.317165in}{1.168843in}}{\pgfqpoint{8.312518in}{1.173904in}}%
\pgfpathquadraticcurveto{\pgfqpoint{8.307871in}{1.178984in}}{\pgfqpoint{8.300215in}{1.178984in}}%
\pgfpathquadraticcurveto{\pgfqpoint{8.291552in}{1.178984in}}{\pgfqpoint{8.286346in}{1.174084in}}%
\pgfpathquadraticcurveto{\pgfqpoint{8.281141in}{1.169185in}}{\pgfqpoint{8.280348in}{1.160287in}}%
\pgfpathlineto{\pgfqpoint{8.317273in}{1.160359in}}%
\pgfpathlineto{\pgfqpoint{8.317273in}{1.160359in}}%
\pgfpathclose%
\pgfpathmoveto{\pgfqpoint{8.384819in}{1.184387in}}%
\pgfpathlineto{\pgfqpoint{8.384819in}{1.174589in}}%
\pgfpathquadraticcurveto{\pgfqpoint{8.380442in}{1.176840in}}{\pgfqpoint{8.375704in}{1.177957in}}%
\pgfpathquadraticcurveto{\pgfqpoint{8.370985in}{1.179092in}}{\pgfqpoint{8.365906in}{1.179092in}}%
\pgfpathquadraticcurveto{\pgfqpoint{8.358197in}{1.179092in}}{\pgfqpoint{8.354342in}{1.176732in}}%
\pgfpathquadraticcurveto{\pgfqpoint{8.350487in}{1.174373in}}{\pgfqpoint{8.350487in}{1.169635in}}%
\pgfpathquadraticcurveto{\pgfqpoint{8.350487in}{1.166033in}}{\pgfqpoint{8.353243in}{1.163980in}}%
\pgfpathquadraticcurveto{\pgfqpoint{8.355999in}{1.161926in}}{\pgfqpoint{8.364339in}{1.160071in}}%
\pgfpathlineto{\pgfqpoint{8.367887in}{1.159278in}}%
\pgfpathquadraticcurveto{\pgfqpoint{8.378911in}{1.156919in}}{\pgfqpoint{8.383558in}{1.152614in}}%
\pgfpathquadraticcurveto{\pgfqpoint{8.388205in}{1.148309in}}{\pgfqpoint{8.388205in}{1.140600in}}%
\pgfpathquadraticcurveto{\pgfqpoint{8.388205in}{1.131810in}}{\pgfqpoint{8.381252in}{1.126676in}}%
\pgfpathquadraticcurveto{\pgfqpoint{8.374299in}{1.121561in}}{\pgfqpoint{8.362141in}{1.121561in}}%
\pgfpathquadraticcurveto{\pgfqpoint{8.357080in}{1.121561in}}{\pgfqpoint{8.351586in}{1.122552in}}%
\pgfpathquadraticcurveto{\pgfqpoint{8.346092in}{1.123542in}}{\pgfqpoint{8.340022in}{1.125506in}}%
\pgfpathlineto{\pgfqpoint{8.340022in}{1.136205in}}%
\pgfpathquadraticcurveto{\pgfqpoint{8.345768in}{1.133215in}}{\pgfqpoint{8.351334in}{1.131720in}}%
\pgfpathquadraticcurveto{\pgfqpoint{8.356900in}{1.130243in}}{\pgfqpoint{8.362375in}{1.130243in}}%
\pgfpathquadraticcurveto{\pgfqpoint{8.369688in}{1.130243in}}{\pgfqpoint{8.373615in}{1.132746in}}%
\pgfpathquadraticcurveto{\pgfqpoint{8.377560in}{1.135250in}}{\pgfqpoint{8.377560in}{1.139807in}}%
\pgfpathquadraticcurveto{\pgfqpoint{8.377560in}{1.144022in}}{\pgfqpoint{8.374714in}{1.146274in}}%
\pgfpathquadraticcurveto{\pgfqpoint{8.371886in}{1.148525in}}{\pgfqpoint{8.362249in}{1.150614in}}%
\pgfpathlineto{\pgfqpoint{8.358647in}{1.151461in}}%
\pgfpathquadraticcurveto{\pgfqpoint{8.349028in}{1.153478in}}{\pgfqpoint{8.344742in}{1.157675in}}%
\pgfpathquadraticcurveto{\pgfqpoint{8.340473in}{1.161872in}}{\pgfqpoint{8.340473in}{1.169185in}}%
\pgfpathquadraticcurveto{\pgfqpoint{8.340473in}{1.178083in}}{\pgfqpoint{8.346777in}{1.182910in}}%
\pgfpathquadraticcurveto{\pgfqpoint{8.353081in}{1.187756in}}{\pgfqpoint{8.364681in}{1.187756in}}%
\pgfpathquadraticcurveto{\pgfqpoint{8.370409in}{1.187756in}}{\pgfqpoint{8.375470in}{1.186909in}}%
\pgfpathquadraticcurveto{\pgfqpoint{8.380550in}{1.186080in}}{\pgfqpoint{8.384819in}{1.184387in}}%
\pgfpathlineto{\pgfqpoint{8.384819in}{1.184387in}}%
\pgfpathclose%
\pgfusepath{fill}%
\end{pgfscope}%
\begin{pgfscope}%
\pgfsetbuttcap%
\pgfsetroundjoin%
\definecolor{currentfill}{rgb}{0.309804,0.372549,0.419608}%
\pgfsetfillcolor{currentfill}%
\pgfsetlinewidth{0.000000pt}%
\definecolor{currentstroke}{rgb}{0.309804,0.372549,0.419608}%
\pgfsetstrokecolor{currentstroke}%
\pgfsetdash{}{0pt}%
\pgfpathmoveto{\pgfqpoint{8.500355in}{1.154883in}}%
\pgfpathquadraticcurveto{\pgfqpoint{8.487800in}{1.154883in}}{\pgfqpoint{8.482955in}{1.152019in}}%
\pgfpathquadraticcurveto{\pgfqpoint{8.478110in}{1.149156in}}{\pgfqpoint{8.478110in}{1.142221in}}%
\pgfpathquadraticcurveto{\pgfqpoint{8.478110in}{1.136709in}}{\pgfqpoint{8.481748in}{1.133467in}}%
\pgfpathquadraticcurveto{\pgfqpoint{8.485387in}{1.130243in}}{\pgfqpoint{8.491619in}{1.130243in}}%
\pgfpathquadraticcurveto{\pgfqpoint{8.500247in}{1.130243in}}{\pgfqpoint{8.505452in}{1.136349in}}%
\pgfpathquadraticcurveto{\pgfqpoint{8.510658in}{1.142455in}}{\pgfqpoint{8.510658in}{1.152578in}}%
\pgfpathlineto{\pgfqpoint{8.510658in}{1.154883in}}%
\pgfpathlineto{\pgfqpoint{8.500355in}{1.154883in}}%
\pgfpathlineto{\pgfqpoint{8.500355in}{1.154883in}}%
\pgfpathclose%
\pgfpathmoveto{\pgfqpoint{8.521015in}{1.159170in}}%
\pgfpathlineto{\pgfqpoint{8.521015in}{1.123200in}}%
\pgfpathlineto{\pgfqpoint{8.510658in}{1.123200in}}%
\pgfpathlineto{\pgfqpoint{8.510658in}{1.132764in}}%
\pgfpathquadraticcurveto{\pgfqpoint{8.507109in}{1.127037in}}{\pgfqpoint{8.501814in}{1.124299in}}%
\pgfpathquadraticcurveto{\pgfqpoint{8.496518in}{1.121561in}}{\pgfqpoint{8.488863in}{1.121561in}}%
\pgfpathquadraticcurveto{\pgfqpoint{8.479190in}{1.121561in}}{\pgfqpoint{8.473462in}{1.127001in}}%
\pgfpathquadraticcurveto{\pgfqpoint{8.467753in}{1.132440in}}{\pgfqpoint{8.467753in}{1.141554in}}%
\pgfpathquadraticcurveto{\pgfqpoint{8.467753in}{1.152182in}}{\pgfqpoint{8.474867in}{1.157585in}}%
\pgfpathquadraticcurveto{\pgfqpoint{8.482000in}{1.162989in}}{\pgfqpoint{8.496122in}{1.162989in}}%
\pgfpathlineto{\pgfqpoint{8.510658in}{1.162989in}}%
\pgfpathlineto{\pgfqpoint{8.510658in}{1.164016in}}%
\pgfpathquadraticcurveto{\pgfqpoint{8.510658in}{1.171166in}}{\pgfqpoint{8.505956in}{1.175075in}}%
\pgfpathquadraticcurveto{\pgfqpoint{8.501255in}{1.178984in}}{\pgfqpoint{8.492753in}{1.178984in}}%
\pgfpathquadraticcurveto{\pgfqpoint{8.487350in}{1.178984in}}{\pgfqpoint{8.482216in}{1.177687in}}%
\pgfpathquadraticcurveto{\pgfqpoint{8.477101in}{1.176390in}}{\pgfqpoint{8.472382in}{1.173796in}}%
\pgfpathlineto{\pgfqpoint{8.472382in}{1.183379in}}%
\pgfpathquadraticcurveto{\pgfqpoint{8.478056in}{1.185576in}}{\pgfqpoint{8.483405in}{1.186657in}}%
\pgfpathquadraticcurveto{\pgfqpoint{8.488755in}{1.187756in}}{\pgfqpoint{8.493816in}{1.187756in}}%
\pgfpathquadraticcurveto{\pgfqpoint{8.507505in}{1.187756in}}{\pgfqpoint{8.514260in}{1.180659in}}%
\pgfpathquadraticcurveto{\pgfqpoint{8.521015in}{1.173580in}}{\pgfqpoint{8.521015in}{1.159170in}}%
\pgfpathlineto{\pgfqpoint{8.521015in}{1.159170in}}%
\pgfpathclose%
\pgfpathmoveto{\pgfqpoint{8.578870in}{1.176570in}}%
\pgfpathquadraticcurveto{\pgfqpoint{8.577122in}{1.177579in}}{\pgfqpoint{8.575069in}{1.178047in}}%
\pgfpathquadraticcurveto{\pgfqpoint{8.573016in}{1.178533in}}{\pgfqpoint{8.570548in}{1.178533in}}%
\pgfpathquadraticcurveto{\pgfqpoint{8.561758in}{1.178533in}}{\pgfqpoint{8.557057in}{1.172823in}}%
\pgfpathquadraticcurveto{\pgfqpoint{8.552356in}{1.167114in}}{\pgfqpoint{8.552356in}{1.156414in}}%
\pgfpathlineto{\pgfqpoint{8.552356in}{1.123200in}}%
\pgfpathlineto{\pgfqpoint{8.541945in}{1.123200in}}%
\pgfpathlineto{\pgfqpoint{8.541945in}{1.186243in}}%
\pgfpathlineto{\pgfqpoint{8.552356in}{1.186243in}}%
\pgfpathlineto{\pgfqpoint{8.552356in}{1.176444in}}%
\pgfpathquadraticcurveto{\pgfqpoint{8.555634in}{1.182190in}}{\pgfqpoint{8.560857in}{1.184964in}}%
\pgfpathquadraticcurveto{\pgfqpoint{8.566099in}{1.187756in}}{\pgfqpoint{8.573592in}{1.187756in}}%
\pgfpathquadraticcurveto{\pgfqpoint{8.574655in}{1.187756in}}{\pgfqpoint{8.575952in}{1.187611in}}%
\pgfpathquadraticcurveto{\pgfqpoint{8.577248in}{1.187485in}}{\pgfqpoint{8.578816in}{1.187197in}}%
\pgfpathlineto{\pgfqpoint{8.578870in}{1.176570in}}%
\pgfpathlineto{\pgfqpoint{8.578870in}{1.176570in}}%
\pgfpathclose%
\pgfpathmoveto{\pgfqpoint{8.641120in}{1.157315in}}%
\pgfpathlineto{\pgfqpoint{8.641120in}{1.152254in}}%
\pgfpathlineto{\pgfqpoint{8.593495in}{1.152254in}}%
\pgfpathquadraticcurveto{\pgfqpoint{8.594180in}{1.141554in}}{\pgfqpoint{8.599944in}{1.135953in}}%
\pgfpathquadraticcurveto{\pgfqpoint{8.605708in}{1.130351in}}{\pgfqpoint{8.616011in}{1.130351in}}%
\pgfpathquadraticcurveto{\pgfqpoint{8.621973in}{1.130351in}}{\pgfqpoint{8.627574in}{1.131810in}}%
\pgfpathquadraticcurveto{\pgfqpoint{8.633176in}{1.133269in}}{\pgfqpoint{8.638706in}{1.136205in}}%
\pgfpathlineto{\pgfqpoint{8.638706in}{1.126406in}}%
\pgfpathquadraticcurveto{\pgfqpoint{8.633122in}{1.124047in}}{\pgfqpoint{8.627268in}{1.122804in}}%
\pgfpathquadraticcurveto{\pgfqpoint{8.621414in}{1.121561in}}{\pgfqpoint{8.615398in}{1.121561in}}%
\pgfpathquadraticcurveto{\pgfqpoint{8.600304in}{1.121561in}}{\pgfqpoint{8.591496in}{1.130333in}}%
\pgfpathquadraticcurveto{\pgfqpoint{8.582688in}{1.139123in}}{\pgfqpoint{8.582688in}{1.154109in}}%
\pgfpathquadraticcurveto{\pgfqpoint{8.582688in}{1.169581in}}{\pgfqpoint{8.591046in}{1.178659in}}%
\pgfpathquadraticcurveto{\pgfqpoint{8.599403in}{1.187756in}}{\pgfqpoint{8.613597in}{1.187756in}}%
\pgfpathquadraticcurveto{\pgfqpoint{8.626314in}{1.187756in}}{\pgfqpoint{8.633717in}{1.179560in}}%
\pgfpathquadraticcurveto{\pgfqpoint{8.641120in}{1.171383in}}{\pgfqpoint{8.641120in}{1.157315in}}%
\pgfpathlineto{\pgfqpoint{8.641120in}{1.157315in}}%
\pgfpathclose%
\pgfpathmoveto{\pgfqpoint{8.630763in}{1.160359in}}%
\pgfpathquadraticcurveto{\pgfqpoint{8.630655in}{1.168843in}}{\pgfqpoint{8.626007in}{1.173904in}}%
\pgfpathquadraticcurveto{\pgfqpoint{8.621360in}{1.178984in}}{\pgfqpoint{8.613705in}{1.178984in}}%
\pgfpathquadraticcurveto{\pgfqpoint{8.605041in}{1.178984in}}{\pgfqpoint{8.599836in}{1.174084in}}%
\pgfpathquadraticcurveto{\pgfqpoint{8.594630in}{1.169185in}}{\pgfqpoint{8.593838in}{1.160287in}}%
\pgfpathlineto{\pgfqpoint{8.630763in}{1.160359in}}%
\pgfpathlineto{\pgfqpoint{8.630763in}{1.160359in}}%
\pgfpathclose%
\pgfusepath{fill}%
\end{pgfscope}%
\begin{pgfscope}%
\pgfsetbuttcap%
\pgfsetroundjoin%
\definecolor{currentfill}{rgb}{0.309804,0.372549,0.419608}%
\pgfsetfillcolor{currentfill}%
\pgfsetlinewidth{0.000000pt}%
\definecolor{currentstroke}{rgb}{0.309804,0.372549,0.419608}%
\pgfsetstrokecolor{currentstroke}%
\pgfsetdash{}{0pt}%
\pgfpathmoveto{\pgfqpoint{8.759697in}{1.176678in}}%
\pgfpathlineto{\pgfqpoint{8.759697in}{1.210793in}}%
\pgfpathlineto{\pgfqpoint{8.770054in}{1.210793in}}%
\pgfpathlineto{\pgfqpoint{8.770054in}{1.123200in}}%
\pgfpathlineto{\pgfqpoint{8.759697in}{1.123200in}}%
\pgfpathlineto{\pgfqpoint{8.759697in}{1.132656in}}%
\pgfpathquadraticcurveto{\pgfqpoint{8.756437in}{1.127037in}}{\pgfqpoint{8.751448in}{1.124299in}}%
\pgfpathquadraticcurveto{\pgfqpoint{8.746477in}{1.121561in}}{\pgfqpoint{8.739488in}{1.121561in}}%
\pgfpathquadraticcurveto{\pgfqpoint{8.728068in}{1.121561in}}{\pgfqpoint{8.720881in}{1.130675in}}%
\pgfpathquadraticcurveto{\pgfqpoint{8.713712in}{1.139807in}}{\pgfqpoint{8.713712in}{1.154667in}}%
\pgfpathquadraticcurveto{\pgfqpoint{8.713712in}{1.169527in}}{\pgfqpoint{8.720881in}{1.178641in}}%
\pgfpathquadraticcurveto{\pgfqpoint{8.728068in}{1.187756in}}{\pgfqpoint{8.739488in}{1.187756in}}%
\pgfpathquadraticcurveto{\pgfqpoint{8.746477in}{1.187756in}}{\pgfqpoint{8.751448in}{1.185018in}}%
\pgfpathquadraticcurveto{\pgfqpoint{8.756437in}{1.182298in}}{\pgfqpoint{8.759697in}{1.176678in}}%
\pgfpathlineto{\pgfqpoint{8.759697in}{1.176678in}}%
\pgfpathclose%
\pgfpathmoveto{\pgfqpoint{8.724412in}{1.154667in}}%
\pgfpathquadraticcurveto{\pgfqpoint{8.724412in}{1.143248in}}{\pgfqpoint{8.729113in}{1.136745in}}%
\pgfpathquadraticcurveto{\pgfqpoint{8.733814in}{1.130243in}}{\pgfqpoint{8.742028in}{1.130243in}}%
\pgfpathquadraticcurveto{\pgfqpoint{8.750241in}{1.130243in}}{\pgfqpoint{8.754960in}{1.136745in}}%
\pgfpathquadraticcurveto{\pgfqpoint{8.759697in}{1.143248in}}{\pgfqpoint{8.759697in}{1.154667in}}%
\pgfpathquadraticcurveto{\pgfqpoint{8.759697in}{1.166087in}}{\pgfqpoint{8.754960in}{1.172589in}}%
\pgfpathquadraticcurveto{\pgfqpoint{8.750241in}{1.179092in}}{\pgfqpoint{8.742028in}{1.179092in}}%
\pgfpathquadraticcurveto{\pgfqpoint{8.733814in}{1.179092in}}{\pgfqpoint{8.729113in}{1.172589in}}%
\pgfpathquadraticcurveto{\pgfqpoint{8.724412in}{1.166087in}}{\pgfqpoint{8.724412in}{1.154667in}}%
\pgfpathlineto{\pgfqpoint{8.724412in}{1.154667in}}%
\pgfpathclose%
\pgfpathmoveto{\pgfqpoint{8.791399in}{1.186243in}}%
\pgfpathlineto{\pgfqpoint{8.801756in}{1.186243in}}%
\pgfpathlineto{\pgfqpoint{8.801756in}{1.123200in}}%
\pgfpathlineto{\pgfqpoint{8.791399in}{1.123200in}}%
\pgfpathlineto{\pgfqpoint{8.791399in}{1.186243in}}%
\pgfpathlineto{\pgfqpoint{8.791399in}{1.186243in}}%
\pgfpathclose%
\pgfpathmoveto{\pgfqpoint{8.791399in}{1.210793in}}%
\pgfpathlineto{\pgfqpoint{8.801756in}{1.210793in}}%
\pgfpathlineto{\pgfqpoint{8.801756in}{1.197662in}}%
\pgfpathlineto{\pgfqpoint{8.791399in}{1.197662in}}%
\pgfpathlineto{\pgfqpoint{8.791399in}{1.210793in}}%
\pgfpathlineto{\pgfqpoint{8.791399in}{1.210793in}}%
\pgfpathclose%
\pgfpathmoveto{\pgfqpoint{8.852082in}{1.154883in}}%
\pgfpathquadraticcurveto{\pgfqpoint{8.839527in}{1.154883in}}{\pgfqpoint{8.834682in}{1.152019in}}%
\pgfpathquadraticcurveto{\pgfqpoint{8.829837in}{1.149156in}}{\pgfqpoint{8.829837in}{1.142221in}}%
\pgfpathquadraticcurveto{\pgfqpoint{8.829837in}{1.136709in}}{\pgfqpoint{8.833475in}{1.133467in}}%
\pgfpathquadraticcurveto{\pgfqpoint{8.837114in}{1.130243in}}{\pgfqpoint{8.843346in}{1.130243in}}%
\pgfpathquadraticcurveto{\pgfqpoint{8.851974in}{1.130243in}}{\pgfqpoint{8.857179in}{1.136349in}}%
\pgfpathquadraticcurveto{\pgfqpoint{8.862385in}{1.142455in}}{\pgfqpoint{8.862385in}{1.152578in}}%
\pgfpathlineto{\pgfqpoint{8.862385in}{1.154883in}}%
\pgfpathlineto{\pgfqpoint{8.852082in}{1.154883in}}%
\pgfpathlineto{\pgfqpoint{8.852082in}{1.154883in}}%
\pgfpathclose%
\pgfpathmoveto{\pgfqpoint{8.872742in}{1.159170in}}%
\pgfpathlineto{\pgfqpoint{8.872742in}{1.123200in}}%
\pgfpathlineto{\pgfqpoint{8.862385in}{1.123200in}}%
\pgfpathlineto{\pgfqpoint{8.862385in}{1.132764in}}%
\pgfpathquadraticcurveto{\pgfqpoint{8.858836in}{1.127037in}}{\pgfqpoint{8.853541in}{1.124299in}}%
\pgfpathquadraticcurveto{\pgfqpoint{8.848245in}{1.121561in}}{\pgfqpoint{8.840590in}{1.121561in}}%
\pgfpathquadraticcurveto{\pgfqpoint{8.830918in}{1.121561in}}{\pgfqpoint{8.825190in}{1.127001in}}%
\pgfpathquadraticcurveto{\pgfqpoint{8.819480in}{1.132440in}}{\pgfqpoint{8.819480in}{1.141554in}}%
\pgfpathquadraticcurveto{\pgfqpoint{8.819480in}{1.152182in}}{\pgfqpoint{8.826595in}{1.157585in}}%
\pgfpathquadraticcurveto{\pgfqpoint{8.833727in}{1.162989in}}{\pgfqpoint{8.847849in}{1.162989in}}%
\pgfpathlineto{\pgfqpoint{8.862385in}{1.162989in}}%
\pgfpathlineto{\pgfqpoint{8.862385in}{1.164016in}}%
\pgfpathquadraticcurveto{\pgfqpoint{8.862385in}{1.171166in}}{\pgfqpoint{8.857684in}{1.175075in}}%
\pgfpathquadraticcurveto{\pgfqpoint{8.852982in}{1.178984in}}{\pgfqpoint{8.844481in}{1.178984in}}%
\pgfpathquadraticcurveto{\pgfqpoint{8.839077in}{1.178984in}}{\pgfqpoint{8.833944in}{1.177687in}}%
\pgfpathquadraticcurveto{\pgfqpoint{8.828828in}{1.176390in}}{\pgfqpoint{8.824109in}{1.173796in}}%
\pgfpathlineto{\pgfqpoint{8.824109in}{1.183379in}}%
\pgfpathquadraticcurveto{\pgfqpoint{8.829783in}{1.185576in}}{\pgfqpoint{8.835132in}{1.186657in}}%
\pgfpathquadraticcurveto{\pgfqpoint{8.840482in}{1.187756in}}{\pgfqpoint{8.845543in}{1.187756in}}%
\pgfpathquadraticcurveto{\pgfqpoint{8.859233in}{1.187756in}}{\pgfqpoint{8.865987in}{1.180659in}}%
\pgfpathquadraticcurveto{\pgfqpoint{8.872742in}{1.173580in}}{\pgfqpoint{8.872742in}{1.159170in}}%
\pgfpathlineto{\pgfqpoint{8.872742in}{1.159170in}}%
\pgfpathclose%
\pgfpathmoveto{\pgfqpoint{8.935550in}{1.155460in}}%
\pgfpathquadraticcurveto{\pgfqpoint{8.935550in}{1.166717in}}{\pgfqpoint{8.930903in}{1.172896in}}%
\pgfpathquadraticcurveto{\pgfqpoint{8.926274in}{1.179092in}}{\pgfqpoint{8.917880in}{1.179092in}}%
\pgfpathquadraticcurveto{\pgfqpoint{8.909559in}{1.179092in}}{\pgfqpoint{8.904911in}{1.172896in}}%
\pgfpathquadraticcurveto{\pgfqpoint{8.900264in}{1.166717in}}{\pgfqpoint{8.900264in}{1.155460in}}%
\pgfpathquadraticcurveto{\pgfqpoint{8.900264in}{1.144256in}}{\pgfqpoint{8.904911in}{1.138060in}}%
\pgfpathquadraticcurveto{\pgfqpoint{8.909559in}{1.131864in}}{\pgfqpoint{8.917880in}{1.131864in}}%
\pgfpathquadraticcurveto{\pgfqpoint{8.926274in}{1.131864in}}{\pgfqpoint{8.930903in}{1.138060in}}%
\pgfpathquadraticcurveto{\pgfqpoint{8.935550in}{1.144256in}}{\pgfqpoint{8.935550in}{1.155460in}}%
\pgfpathlineto{\pgfqpoint{8.935550in}{1.155460in}}%
\pgfpathclose%
\pgfpathmoveto{\pgfqpoint{8.945907in}{1.131017in}}%
\pgfpathquadraticcurveto{\pgfqpoint{8.945907in}{1.114932in}}{\pgfqpoint{8.938756in}{1.107079in}}%
\pgfpathquadraticcurveto{\pgfqpoint{8.931623in}{1.099226in}}{\pgfqpoint{8.916872in}{1.099226in}}%
\pgfpathquadraticcurveto{\pgfqpoint{8.911414in}{1.099226in}}{\pgfqpoint{8.906569in}{1.100036in}}%
\pgfpathquadraticcurveto{\pgfqpoint{8.901723in}{1.100847in}}{\pgfqpoint{8.897166in}{1.102540in}}%
\pgfpathlineto{\pgfqpoint{8.897166in}{1.112609in}}%
\pgfpathquadraticcurveto{\pgfqpoint{8.901723in}{1.110141in}}{\pgfqpoint{8.906172in}{1.108970in}}%
\pgfpathquadraticcurveto{\pgfqpoint{8.910621in}{1.107782in}}{\pgfqpoint{8.915232in}{1.107782in}}%
\pgfpathquadraticcurveto{\pgfqpoint{8.925427in}{1.107782in}}{\pgfqpoint{8.930489in}{1.113095in}}%
\pgfpathquadraticcurveto{\pgfqpoint{8.935550in}{1.118409in}}{\pgfqpoint{8.935550in}{1.129162in}}%
\pgfpathlineto{\pgfqpoint{8.935550in}{1.134295in}}%
\pgfpathquadraticcurveto{\pgfqpoint{8.932344in}{1.128712in}}{\pgfqpoint{8.927337in}{1.125956in}}%
\pgfpathquadraticcurveto{\pgfqpoint{8.922329in}{1.123200in}}{\pgfqpoint{8.915340in}{1.123200in}}%
\pgfpathquadraticcurveto{\pgfqpoint{8.903759in}{1.123200in}}{\pgfqpoint{8.896662in}{1.132026in}}%
\pgfpathquadraticcurveto{\pgfqpoint{8.889565in}{1.140870in}}{\pgfqpoint{8.889565in}{1.155460in}}%
\pgfpathquadraticcurveto{\pgfqpoint{8.889565in}{1.170086in}}{\pgfqpoint{8.896662in}{1.178912in}}%
\pgfpathquadraticcurveto{\pgfqpoint{8.903759in}{1.187756in}}{\pgfqpoint{8.915340in}{1.187756in}}%
\pgfpathquadraticcurveto{\pgfqpoint{8.922329in}{1.187756in}}{\pgfqpoint{8.927337in}{1.185000in}}%
\pgfpathquadraticcurveto{\pgfqpoint{8.932344in}{1.182244in}}{\pgfqpoint{8.935550in}{1.176678in}}%
\pgfpathlineto{\pgfqpoint{8.935550in}{1.186243in}}%
\pgfpathlineto{\pgfqpoint{8.945907in}{1.186243in}}%
\pgfpathlineto{\pgfqpoint{8.945907in}{1.131017in}}%
\pgfpathlineto{\pgfqpoint{8.945907in}{1.131017in}}%
\pgfpathclose%
\pgfpathmoveto{\pgfqpoint{9.019667in}{1.161260in}}%
\pgfpathlineto{\pgfqpoint{9.019667in}{1.123200in}}%
\pgfpathlineto{\pgfqpoint{9.009310in}{1.123200in}}%
\pgfpathlineto{\pgfqpoint{9.009310in}{1.160917in}}%
\pgfpathquadraticcurveto{\pgfqpoint{9.009310in}{1.169869in}}{\pgfqpoint{9.005816in}{1.174300in}}%
\pgfpathquadraticcurveto{\pgfqpoint{9.002321in}{1.178749in}}{\pgfqpoint{8.995350in}{1.178749in}}%
\pgfpathquadraticcurveto{\pgfqpoint{8.986957in}{1.178749in}}{\pgfqpoint{8.982112in}{1.173400in}}%
\pgfpathquadraticcurveto{\pgfqpoint{8.977266in}{1.168068in}}{\pgfqpoint{8.977266in}{1.158828in}}%
\pgfpathlineto{\pgfqpoint{8.977266in}{1.123200in}}%
\pgfpathlineto{\pgfqpoint{8.966855in}{1.123200in}}%
\pgfpathlineto{\pgfqpoint{8.966855in}{1.186243in}}%
\pgfpathlineto{\pgfqpoint{8.977266in}{1.186243in}}%
\pgfpathlineto{\pgfqpoint{8.977266in}{1.176444in}}%
\pgfpathquadraticcurveto{\pgfqpoint{8.980995in}{1.182136in}}{\pgfqpoint{8.986020in}{1.184946in}}%
\pgfpathquadraticcurveto{\pgfqpoint{8.991064in}{1.187756in}}{\pgfqpoint{8.997656in}{1.187756in}}%
\pgfpathquadraticcurveto{\pgfqpoint{9.008517in}{1.187756in}}{\pgfqpoint{9.014083in}{1.181037in}}%
\pgfpathquadraticcurveto{\pgfqpoint{9.019667in}{1.174318in}}{\pgfqpoint{9.019667in}{1.161260in}}%
\pgfpathlineto{\pgfqpoint{9.019667in}{1.161260in}}%
\pgfpathclose%
\pgfpathmoveto{\pgfqpoint{9.064733in}{1.178984in}}%
\pgfpathquadraticcurveto{\pgfqpoint{9.056412in}{1.178984in}}{\pgfqpoint{9.051566in}{1.172481in}}%
\pgfpathquadraticcurveto{\pgfqpoint{9.046721in}{1.165979in}}{\pgfqpoint{9.046721in}{1.154667in}}%
\pgfpathquadraticcurveto{\pgfqpoint{9.046721in}{1.143356in}}{\pgfqpoint{9.051530in}{1.136853in}}%
\pgfpathquadraticcurveto{\pgfqpoint{9.056358in}{1.130351in}}{\pgfqpoint{9.064733in}{1.130351in}}%
\pgfpathquadraticcurveto{\pgfqpoint{9.073019in}{1.130351in}}{\pgfqpoint{9.077846in}{1.136871in}}%
\pgfpathquadraticcurveto{\pgfqpoint{9.082691in}{1.143410in}}{\pgfqpoint{9.082691in}{1.154667in}}%
\pgfpathquadraticcurveto{\pgfqpoint{9.082691in}{1.165871in}}{\pgfqpoint{9.077846in}{1.172427in}}%
\pgfpathquadraticcurveto{\pgfqpoint{9.073019in}{1.178984in}}{\pgfqpoint{9.064733in}{1.178984in}}%
\pgfpathlineto{\pgfqpoint{9.064733in}{1.178984in}}%
\pgfpathclose%
\pgfpathmoveto{\pgfqpoint{9.064733in}{1.187756in}}%
\pgfpathquadraticcurveto{\pgfqpoint{9.078242in}{1.187756in}}{\pgfqpoint{9.085952in}{1.178966in}}%
\pgfpathquadraticcurveto{\pgfqpoint{9.093679in}{1.170194in}}{\pgfqpoint{9.093679in}{1.154667in}}%
\pgfpathquadraticcurveto{\pgfqpoint{9.093679in}{1.139195in}}{\pgfqpoint{9.085952in}{1.130369in}}%
\pgfpathquadraticcurveto{\pgfqpoint{9.078242in}{1.121561in}}{\pgfqpoint{9.064733in}{1.121561in}}%
\pgfpathquadraticcurveto{\pgfqpoint{9.051170in}{1.121561in}}{\pgfqpoint{9.043479in}{1.130369in}}%
\pgfpathquadraticcurveto{\pgfqpoint{9.035806in}{1.139195in}}{\pgfqpoint{9.035806in}{1.154667in}}%
\pgfpathquadraticcurveto{\pgfqpoint{9.035806in}{1.170194in}}{\pgfqpoint{9.043479in}{1.178966in}}%
\pgfpathquadraticcurveto{\pgfqpoint{9.051170in}{1.187756in}}{\pgfqpoint{9.064733in}{1.187756in}}%
\pgfpathlineto{\pgfqpoint{9.064733in}{1.187756in}}%
\pgfpathclose%
\pgfpathmoveto{\pgfqpoint{9.151030in}{1.184387in}}%
\pgfpathlineto{\pgfqpoint{9.151030in}{1.174589in}}%
\pgfpathquadraticcurveto{\pgfqpoint{9.146653in}{1.176840in}}{\pgfqpoint{9.141915in}{1.177957in}}%
\pgfpathquadraticcurveto{\pgfqpoint{9.137196in}{1.179092in}}{\pgfqpoint{9.132117in}{1.179092in}}%
\pgfpathquadraticcurveto{\pgfqpoint{9.124408in}{1.179092in}}{\pgfqpoint{9.120553in}{1.176732in}}%
\pgfpathquadraticcurveto{\pgfqpoint{9.116698in}{1.174373in}}{\pgfqpoint{9.116698in}{1.169635in}}%
\pgfpathquadraticcurveto{\pgfqpoint{9.116698in}{1.166033in}}{\pgfqpoint{9.119454in}{1.163980in}}%
\pgfpathquadraticcurveto{\pgfqpoint{9.122210in}{1.161926in}}{\pgfqpoint{9.130550in}{1.160071in}}%
\pgfpathlineto{\pgfqpoint{9.134098in}{1.159278in}}%
\pgfpathquadraticcurveto{\pgfqpoint{9.145122in}{1.156919in}}{\pgfqpoint{9.149769in}{1.152614in}}%
\pgfpathquadraticcurveto{\pgfqpoint{9.154416in}{1.148309in}}{\pgfqpoint{9.154416in}{1.140600in}}%
\pgfpathquadraticcurveto{\pgfqpoint{9.154416in}{1.131810in}}{\pgfqpoint{9.147463in}{1.126676in}}%
\pgfpathquadraticcurveto{\pgfqpoint{9.140510in}{1.121561in}}{\pgfqpoint{9.128352in}{1.121561in}}%
\pgfpathquadraticcurveto{\pgfqpoint{9.123291in}{1.121561in}}{\pgfqpoint{9.117797in}{1.122552in}}%
\pgfpathquadraticcurveto{\pgfqpoint{9.112303in}{1.123542in}}{\pgfqpoint{9.106233in}{1.125506in}}%
\pgfpathlineto{\pgfqpoint{9.106233in}{1.136205in}}%
\pgfpathquadraticcurveto{\pgfqpoint{9.111979in}{1.133215in}}{\pgfqpoint{9.117545in}{1.131720in}}%
\pgfpathquadraticcurveto{\pgfqpoint{9.123111in}{1.130243in}}{\pgfqpoint{9.128586in}{1.130243in}}%
\pgfpathquadraticcurveto{\pgfqpoint{9.135899in}{1.130243in}}{\pgfqpoint{9.139826in}{1.132746in}}%
\pgfpathquadraticcurveto{\pgfqpoint{9.143771in}{1.135250in}}{\pgfqpoint{9.143771in}{1.139807in}}%
\pgfpathquadraticcurveto{\pgfqpoint{9.143771in}{1.144022in}}{\pgfqpoint{9.140925in}{1.146274in}}%
\pgfpathquadraticcurveto{\pgfqpoint{9.138097in}{1.148525in}}{\pgfqpoint{9.128460in}{1.150614in}}%
\pgfpathlineto{\pgfqpoint{9.124858in}{1.151461in}}%
\pgfpathquadraticcurveto{\pgfqpoint{9.115239in}{1.153478in}}{\pgfqpoint{9.110952in}{1.157675in}}%
\pgfpathquadraticcurveto{\pgfqpoint{9.106684in}{1.161872in}}{\pgfqpoint{9.106684in}{1.169185in}}%
\pgfpathquadraticcurveto{\pgfqpoint{9.106684in}{1.178083in}}{\pgfqpoint{9.112988in}{1.182910in}}%
\pgfpathquadraticcurveto{\pgfqpoint{9.119292in}{1.187756in}}{\pgfqpoint{9.130892in}{1.187756in}}%
\pgfpathquadraticcurveto{\pgfqpoint{9.136620in}{1.187756in}}{\pgfqpoint{9.141681in}{1.186909in}}%
\pgfpathquadraticcurveto{\pgfqpoint{9.146761in}{1.186080in}}{\pgfqpoint{9.151030in}{1.184387in}}%
\pgfpathlineto{\pgfqpoint{9.151030in}{1.184387in}}%
\pgfpathclose%
\pgfpathmoveto{\pgfqpoint{9.181146in}{1.204147in}}%
\pgfpathlineto{\pgfqpoint{9.181146in}{1.186243in}}%
\pgfpathlineto{\pgfqpoint{9.202472in}{1.186243in}}%
\pgfpathlineto{\pgfqpoint{9.202472in}{1.178191in}}%
\pgfpathlineto{\pgfqpoint{9.181146in}{1.178191in}}%
\pgfpathlineto{\pgfqpoint{9.181146in}{1.143968in}}%
\pgfpathquadraticcurveto{\pgfqpoint{9.181146in}{1.136259in}}{\pgfqpoint{9.183253in}{1.134061in}}%
\pgfpathquadraticcurveto{\pgfqpoint{9.185361in}{1.131864in}}{\pgfqpoint{9.191845in}{1.131864in}}%
\pgfpathlineto{\pgfqpoint{9.202472in}{1.131864in}}%
\pgfpathlineto{\pgfqpoint{9.202472in}{1.123200in}}%
\pgfpathlineto{\pgfqpoint{9.191845in}{1.123200in}}%
\pgfpathquadraticcurveto{\pgfqpoint{9.179849in}{1.123200in}}{\pgfqpoint{9.175292in}{1.127667in}}%
\pgfpathquadraticcurveto{\pgfqpoint{9.170735in}{1.132152in}}{\pgfqpoint{9.170735in}{1.143968in}}%
\pgfpathlineto{\pgfqpoint{9.170735in}{1.178191in}}%
\pgfpathlineto{\pgfqpoint{9.163134in}{1.178191in}}%
\pgfpathlineto{\pgfqpoint{9.163134in}{1.186243in}}%
\pgfpathlineto{\pgfqpoint{9.170735in}{1.186243in}}%
\pgfpathlineto{\pgfqpoint{9.170735in}{1.204147in}}%
\pgfpathlineto{\pgfqpoint{9.181146in}{1.204147in}}%
\pgfpathlineto{\pgfqpoint{9.181146in}{1.204147in}}%
\pgfpathclose%
\pgfpathmoveto{\pgfqpoint{9.216089in}{1.186243in}}%
\pgfpathlineto{\pgfqpoint{9.226446in}{1.186243in}}%
\pgfpathlineto{\pgfqpoint{9.226446in}{1.123200in}}%
\pgfpathlineto{\pgfqpoint{9.216089in}{1.123200in}}%
\pgfpathlineto{\pgfqpoint{9.216089in}{1.186243in}}%
\pgfpathlineto{\pgfqpoint{9.216089in}{1.186243in}}%
\pgfpathclose%
\pgfpathmoveto{\pgfqpoint{9.216089in}{1.210793in}}%
\pgfpathlineto{\pgfqpoint{9.226446in}{1.210793in}}%
\pgfpathlineto{\pgfqpoint{9.226446in}{1.197662in}}%
\pgfpathlineto{\pgfqpoint{9.216089in}{1.197662in}}%
\pgfpathlineto{\pgfqpoint{9.216089in}{1.210793in}}%
\pgfpathlineto{\pgfqpoint{9.216089in}{1.210793in}}%
\pgfpathclose%
\pgfpathmoveto{\pgfqpoint{9.293488in}{1.183829in}}%
\pgfpathlineto{\pgfqpoint{9.293488in}{1.174138in}}%
\pgfpathquadraticcurveto{\pgfqpoint{9.289093in}{1.176570in}}{\pgfqpoint{9.284680in}{1.177777in}}%
\pgfpathquadraticcurveto{\pgfqpoint{9.280267in}{1.178984in}}{\pgfqpoint{9.275764in}{1.178984in}}%
\pgfpathquadraticcurveto{\pgfqpoint{9.265677in}{1.178984in}}{\pgfqpoint{9.260093in}{1.172589in}}%
\pgfpathquadraticcurveto{\pgfqpoint{9.254527in}{1.166213in}}{\pgfqpoint{9.254527in}{1.154667in}}%
\pgfpathquadraticcurveto{\pgfqpoint{9.254527in}{1.143121in}}{\pgfqpoint{9.260093in}{1.136727in}}%
\pgfpathquadraticcurveto{\pgfqpoint{9.265677in}{1.130351in}}{\pgfqpoint{9.275764in}{1.130351in}}%
\pgfpathquadraticcurveto{\pgfqpoint{9.280267in}{1.130351in}}{\pgfqpoint{9.284680in}{1.131558in}}%
\pgfpathquadraticcurveto{\pgfqpoint{9.289093in}{1.132764in}}{\pgfqpoint{9.293488in}{1.135196in}}%
\pgfpathlineto{\pgfqpoint{9.293488in}{1.125614in}}%
\pgfpathquadraticcurveto{\pgfqpoint{9.289147in}{1.123596in}}{\pgfqpoint{9.284500in}{1.122588in}}%
\pgfpathquadraticcurveto{\pgfqpoint{9.279870in}{1.121561in}}{\pgfqpoint{9.274629in}{1.121561in}}%
\pgfpathquadraticcurveto{\pgfqpoint{9.260381in}{1.121561in}}{\pgfqpoint{9.251988in}{1.130513in}}%
\pgfpathquadraticcurveto{\pgfqpoint{9.243612in}{1.139465in}}{\pgfqpoint{9.243612in}{1.154667in}}%
\pgfpathquadraticcurveto{\pgfqpoint{9.243612in}{1.170086in}}{\pgfqpoint{9.252078in}{1.178912in}}%
\pgfpathquadraticcurveto{\pgfqpoint{9.260561in}{1.187756in}}{\pgfqpoint{9.275313in}{1.187756in}}%
\pgfpathquadraticcurveto{\pgfqpoint{9.280087in}{1.187756in}}{\pgfqpoint{9.284644in}{1.186765in}}%
\pgfpathquadraticcurveto{\pgfqpoint{9.289201in}{1.185792in}}{\pgfqpoint{9.293488in}{1.183829in}}%
\pgfpathlineto{\pgfqpoint{9.293488in}{1.183829in}}%
\pgfpathclose%
\pgfpathmoveto{\pgfqpoint{9.351685in}{1.184387in}}%
\pgfpathlineto{\pgfqpoint{9.351685in}{1.174589in}}%
\pgfpathquadraticcurveto{\pgfqpoint{9.347308in}{1.176840in}}{\pgfqpoint{9.342571in}{1.177957in}}%
\pgfpathquadraticcurveto{\pgfqpoint{9.337852in}{1.179092in}}{\pgfqpoint{9.332772in}{1.179092in}}%
\pgfpathquadraticcurveto{\pgfqpoint{9.325063in}{1.179092in}}{\pgfqpoint{9.321208in}{1.176732in}}%
\pgfpathquadraticcurveto{\pgfqpoint{9.317354in}{1.174373in}}{\pgfqpoint{9.317354in}{1.169635in}}%
\pgfpathquadraticcurveto{\pgfqpoint{9.317354in}{1.166033in}}{\pgfqpoint{9.320110in}{1.163980in}}%
\pgfpathquadraticcurveto{\pgfqpoint{9.322865in}{1.161926in}}{\pgfqpoint{9.331205in}{1.160071in}}%
\pgfpathlineto{\pgfqpoint{9.334753in}{1.159278in}}%
\pgfpathquadraticcurveto{\pgfqpoint{9.345777in}{1.156919in}}{\pgfqpoint{9.350424in}{1.152614in}}%
\pgfpathquadraticcurveto{\pgfqpoint{9.355071in}{1.148309in}}{\pgfqpoint{9.355071in}{1.140600in}}%
\pgfpathquadraticcurveto{\pgfqpoint{9.355071in}{1.131810in}}{\pgfqpoint{9.348118in}{1.126676in}}%
\pgfpathquadraticcurveto{\pgfqpoint{9.341166in}{1.121561in}}{\pgfqpoint{9.329008in}{1.121561in}}%
\pgfpathquadraticcurveto{\pgfqpoint{9.323946in}{1.121561in}}{\pgfqpoint{9.318452in}{1.122552in}}%
\pgfpathquadraticcurveto{\pgfqpoint{9.312959in}{1.123542in}}{\pgfqpoint{9.306889in}{1.125506in}}%
\pgfpathlineto{\pgfqpoint{9.306889in}{1.136205in}}%
\pgfpathquadraticcurveto{\pgfqpoint{9.312635in}{1.133215in}}{\pgfqpoint{9.318200in}{1.131720in}}%
\pgfpathquadraticcurveto{\pgfqpoint{9.323766in}{1.130243in}}{\pgfqpoint{9.329242in}{1.130243in}}%
\pgfpathquadraticcurveto{\pgfqpoint{9.336555in}{1.130243in}}{\pgfqpoint{9.340481in}{1.132746in}}%
\pgfpathquadraticcurveto{\pgfqpoint{9.344426in}{1.135250in}}{\pgfqpoint{9.344426in}{1.139807in}}%
\pgfpathquadraticcurveto{\pgfqpoint{9.344426in}{1.144022in}}{\pgfqpoint{9.341580in}{1.146274in}}%
\pgfpathquadraticcurveto{\pgfqpoint{9.338752in}{1.148525in}}{\pgfqpoint{9.329116in}{1.150614in}}%
\pgfpathlineto{\pgfqpoint{9.325513in}{1.151461in}}%
\pgfpathquadraticcurveto{\pgfqpoint{9.315895in}{1.153478in}}{\pgfqpoint{9.311608in}{1.157675in}}%
\pgfpathquadraticcurveto{\pgfqpoint{9.307339in}{1.161872in}}{\pgfqpoint{9.307339in}{1.169185in}}%
\pgfpathquadraticcurveto{\pgfqpoint{9.307339in}{1.178083in}}{\pgfqpoint{9.313643in}{1.182910in}}%
\pgfpathquadraticcurveto{\pgfqpoint{9.319947in}{1.187756in}}{\pgfqpoint{9.331547in}{1.187756in}}%
\pgfpathquadraticcurveto{\pgfqpoint{9.337275in}{1.187756in}}{\pgfqpoint{9.342337in}{1.186909in}}%
\pgfpathquadraticcurveto{\pgfqpoint{9.347416in}{1.186080in}}{\pgfqpoint{9.351685in}{1.184387in}}%
\pgfpathlineto{\pgfqpoint{9.351685in}{1.184387in}}%
\pgfpathclose%
\pgfpathmoveto{\pgfqpoint{9.374200in}{1.137502in}}%
\pgfpathlineto{\pgfqpoint{9.386070in}{1.137502in}}%
\pgfpathlineto{\pgfqpoint{9.386070in}{1.127811in}}%
\pgfpathlineto{\pgfqpoint{9.376848in}{1.109799in}}%
\pgfpathlineto{\pgfqpoint{9.369589in}{1.109799in}}%
\pgfpathlineto{\pgfqpoint{9.374200in}{1.127811in}}%
\pgfpathlineto{\pgfqpoint{9.374200in}{1.137502in}}%
\pgfpathlineto{\pgfqpoint{9.374200in}{1.137502in}}%
\pgfpathclose%
\pgfusepath{fill}%
\end{pgfscope}%
\begin{pgfscope}%
\pgfsetbuttcap%
\pgfsetroundjoin%
\definecolor{currentfill}{rgb}{0.309804,0.372549,0.419608}%
\pgfsetfillcolor{currentfill}%
\pgfsetlinewidth{0.000000pt}%
\definecolor{currentstroke}{rgb}{0.309804,0.372549,0.419608}%
\pgfsetstrokecolor{currentstroke}%
\pgfsetdash{}{0pt}%
\pgfpathmoveto{\pgfqpoint{9.549649in}{1.161260in}}%
\pgfpathlineto{\pgfqpoint{9.549649in}{1.123200in}}%
\pgfpathlineto{\pgfqpoint{9.539292in}{1.123200in}}%
\pgfpathlineto{\pgfqpoint{9.539292in}{1.160917in}}%
\pgfpathquadraticcurveto{\pgfqpoint{9.539292in}{1.169869in}}{\pgfqpoint{9.535798in}{1.174300in}}%
\pgfpathquadraticcurveto{\pgfqpoint{9.532303in}{1.178749in}}{\pgfqpoint{9.525333in}{1.178749in}}%
\pgfpathquadraticcurveto{\pgfqpoint{9.516939in}{1.178749in}}{\pgfqpoint{9.512094in}{1.173400in}}%
\pgfpathquadraticcurveto{\pgfqpoint{9.507248in}{1.168068in}}{\pgfqpoint{9.507248in}{1.158828in}}%
\pgfpathlineto{\pgfqpoint{9.507248in}{1.123200in}}%
\pgfpathlineto{\pgfqpoint{9.496837in}{1.123200in}}%
\pgfpathlineto{\pgfqpoint{9.496837in}{1.186243in}}%
\pgfpathlineto{\pgfqpoint{9.507248in}{1.186243in}}%
\pgfpathlineto{\pgfqpoint{9.507248in}{1.176444in}}%
\pgfpathquadraticcurveto{\pgfqpoint{9.510977in}{1.182136in}}{\pgfqpoint{9.516002in}{1.184946in}}%
\pgfpathquadraticcurveto{\pgfqpoint{9.521046in}{1.187756in}}{\pgfqpoint{9.527638in}{1.187756in}}%
\pgfpathquadraticcurveto{\pgfqpoint{9.538500in}{1.187756in}}{\pgfqpoint{9.544065in}{1.181037in}}%
\pgfpathquadraticcurveto{\pgfqpoint{9.549649in}{1.174318in}}{\pgfqpoint{9.549649in}{1.161260in}}%
\pgfpathlineto{\pgfqpoint{9.549649in}{1.161260in}}%
\pgfpathclose%
\pgfpathmoveto{\pgfqpoint{9.594715in}{1.178984in}}%
\pgfpathquadraticcurveto{\pgfqpoint{9.586394in}{1.178984in}}{\pgfqpoint{9.581549in}{1.172481in}}%
\pgfpathquadraticcurveto{\pgfqpoint{9.576703in}{1.165979in}}{\pgfqpoint{9.576703in}{1.154667in}}%
\pgfpathquadraticcurveto{\pgfqpoint{9.576703in}{1.143356in}}{\pgfqpoint{9.581513in}{1.136853in}}%
\pgfpathquadraticcurveto{\pgfqpoint{9.586340in}{1.130351in}}{\pgfqpoint{9.594715in}{1.130351in}}%
\pgfpathquadraticcurveto{\pgfqpoint{9.603001in}{1.130351in}}{\pgfqpoint{9.607828in}{1.136871in}}%
\pgfpathquadraticcurveto{\pgfqpoint{9.612674in}{1.143410in}}{\pgfqpoint{9.612674in}{1.154667in}}%
\pgfpathquadraticcurveto{\pgfqpoint{9.612674in}{1.165871in}}{\pgfqpoint{9.607828in}{1.172427in}}%
\pgfpathquadraticcurveto{\pgfqpoint{9.603001in}{1.178984in}}{\pgfqpoint{9.594715in}{1.178984in}}%
\pgfpathlineto{\pgfqpoint{9.594715in}{1.178984in}}%
\pgfpathclose%
\pgfpathmoveto{\pgfqpoint{9.594715in}{1.187756in}}%
\pgfpathquadraticcurveto{\pgfqpoint{9.608225in}{1.187756in}}{\pgfqpoint{9.615934in}{1.178966in}}%
\pgfpathquadraticcurveto{\pgfqpoint{9.623661in}{1.170194in}}{\pgfqpoint{9.623661in}{1.154667in}}%
\pgfpathquadraticcurveto{\pgfqpoint{9.623661in}{1.139195in}}{\pgfqpoint{9.615934in}{1.130369in}}%
\pgfpathquadraticcurveto{\pgfqpoint{9.608225in}{1.121561in}}{\pgfqpoint{9.594715in}{1.121561in}}%
\pgfpathquadraticcurveto{\pgfqpoint{9.581152in}{1.121561in}}{\pgfqpoint{9.573461in}{1.130369in}}%
\pgfpathquadraticcurveto{\pgfqpoint{9.565788in}{1.139195in}}{\pgfqpoint{9.565788in}{1.154667in}}%
\pgfpathquadraticcurveto{\pgfqpoint{9.565788in}{1.170194in}}{\pgfqpoint{9.573461in}{1.178966in}}%
\pgfpathquadraticcurveto{\pgfqpoint{9.581152in}{1.187756in}}{\pgfqpoint{9.594715in}{1.187756in}}%
\pgfpathlineto{\pgfqpoint{9.594715in}{1.187756in}}%
\pgfpathclose%
\pgfpathmoveto{\pgfqpoint{9.651076in}{1.204147in}}%
\pgfpathlineto{\pgfqpoint{9.651076in}{1.186243in}}%
\pgfpathlineto{\pgfqpoint{9.672402in}{1.186243in}}%
\pgfpathlineto{\pgfqpoint{9.672402in}{1.178191in}}%
\pgfpathlineto{\pgfqpoint{9.651076in}{1.178191in}}%
\pgfpathlineto{\pgfqpoint{9.651076in}{1.143968in}}%
\pgfpathquadraticcurveto{\pgfqpoint{9.651076in}{1.136259in}}{\pgfqpoint{9.653183in}{1.134061in}}%
\pgfpathquadraticcurveto{\pgfqpoint{9.655290in}{1.131864in}}{\pgfqpoint{9.661775in}{1.131864in}}%
\pgfpathlineto{\pgfqpoint{9.672402in}{1.131864in}}%
\pgfpathlineto{\pgfqpoint{9.672402in}{1.123200in}}%
\pgfpathlineto{\pgfqpoint{9.661775in}{1.123200in}}%
\pgfpathquadraticcurveto{\pgfqpoint{9.649779in}{1.123200in}}{\pgfqpoint{9.645222in}{1.127667in}}%
\pgfpathquadraticcurveto{\pgfqpoint{9.640664in}{1.132152in}}{\pgfqpoint{9.640664in}{1.143968in}}%
\pgfpathlineto{\pgfqpoint{9.640664in}{1.178191in}}%
\pgfpathlineto{\pgfqpoint{9.633063in}{1.178191in}}%
\pgfpathlineto{\pgfqpoint{9.633063in}{1.186243in}}%
\pgfpathlineto{\pgfqpoint{9.640664in}{1.186243in}}%
\pgfpathlineto{\pgfqpoint{9.640664in}{1.204147in}}%
\pgfpathlineto{\pgfqpoint{9.651076in}{1.204147in}}%
\pgfpathlineto{\pgfqpoint{9.651076in}{1.204147in}}%
\pgfpathclose%
\pgfusepath{fill}%
\end{pgfscope}%
\begin{pgfscope}%
\pgfsetbuttcap%
\pgfsetroundjoin%
\definecolor{currentfill}{rgb}{0.309804,0.372549,0.419608}%
\pgfsetfillcolor{currentfill}%
\pgfsetlinewidth{0.000000pt}%
\definecolor{currentstroke}{rgb}{0.309804,0.372549,0.419608}%
\pgfsetstrokecolor{currentstroke}%
\pgfsetdash{}{0pt}%
\pgfpathmoveto{\pgfqpoint{9.766501in}{1.204147in}}%
\pgfpathlineto{\pgfqpoint{9.766501in}{1.186243in}}%
\pgfpathlineto{\pgfqpoint{9.787827in}{1.186243in}}%
\pgfpathlineto{\pgfqpoint{9.787827in}{1.178191in}}%
\pgfpathlineto{\pgfqpoint{9.766501in}{1.178191in}}%
\pgfpathlineto{\pgfqpoint{9.766501in}{1.143968in}}%
\pgfpathquadraticcurveto{\pgfqpoint{9.766501in}{1.136259in}}{\pgfqpoint{9.768608in}{1.134061in}}%
\pgfpathquadraticcurveto{\pgfqpoint{9.770716in}{1.131864in}}{\pgfqpoint{9.777200in}{1.131864in}}%
\pgfpathlineto{\pgfqpoint{9.787827in}{1.131864in}}%
\pgfpathlineto{\pgfqpoint{9.787827in}{1.123200in}}%
\pgfpathlineto{\pgfqpoint{9.777200in}{1.123200in}}%
\pgfpathquadraticcurveto{\pgfqpoint{9.765204in}{1.123200in}}{\pgfqpoint{9.760647in}{1.127667in}}%
\pgfpathquadraticcurveto{\pgfqpoint{9.756090in}{1.132152in}}{\pgfqpoint{9.756090in}{1.143968in}}%
\pgfpathlineto{\pgfqpoint{9.756090in}{1.178191in}}%
\pgfpathlineto{\pgfqpoint{9.748489in}{1.178191in}}%
\pgfpathlineto{\pgfqpoint{9.748489in}{1.186243in}}%
\pgfpathlineto{\pgfqpoint{9.756090in}{1.186243in}}%
\pgfpathlineto{\pgfqpoint{9.756090in}{1.204147in}}%
\pgfpathlineto{\pgfqpoint{9.766501in}{1.204147in}}%
\pgfpathlineto{\pgfqpoint{9.766501in}{1.204147in}}%
\pgfpathclose%
\pgfpathmoveto{\pgfqpoint{9.837973in}{1.176570in}}%
\pgfpathquadraticcurveto{\pgfqpoint{9.836226in}{1.177579in}}{\pgfqpoint{9.834173in}{1.178047in}}%
\pgfpathquadraticcurveto{\pgfqpoint{9.832119in}{1.178533in}}{\pgfqpoint{9.829651in}{1.178533in}}%
\pgfpathquadraticcurveto{\pgfqpoint{9.820862in}{1.178533in}}{\pgfqpoint{9.816160in}{1.172823in}}%
\pgfpathquadraticcurveto{\pgfqpoint{9.811459in}{1.167114in}}{\pgfqpoint{9.811459in}{1.156414in}}%
\pgfpathlineto{\pgfqpoint{9.811459in}{1.123200in}}%
\pgfpathlineto{\pgfqpoint{9.801048in}{1.123200in}}%
\pgfpathlineto{\pgfqpoint{9.801048in}{1.186243in}}%
\pgfpathlineto{\pgfqpoint{9.811459in}{1.186243in}}%
\pgfpathlineto{\pgfqpoint{9.811459in}{1.176444in}}%
\pgfpathquadraticcurveto{\pgfqpoint{9.814737in}{1.182190in}}{\pgfqpoint{9.819961in}{1.184964in}}%
\pgfpathquadraticcurveto{\pgfqpoint{9.825202in}{1.187756in}}{\pgfqpoint{9.832696in}{1.187756in}}%
\pgfpathquadraticcurveto{\pgfqpoint{9.833758in}{1.187756in}}{\pgfqpoint{9.835055in}{1.187611in}}%
\pgfpathquadraticcurveto{\pgfqpoint{9.836352in}{1.187485in}}{\pgfqpoint{9.837919in}{1.187197in}}%
\pgfpathlineto{\pgfqpoint{9.837973in}{1.176570in}}%
\pgfpathlineto{\pgfqpoint{9.837973in}{1.176570in}}%
\pgfpathclose%
\pgfpathmoveto{\pgfqpoint{9.900223in}{1.157315in}}%
\pgfpathlineto{\pgfqpoint{9.900223in}{1.152254in}}%
\pgfpathlineto{\pgfqpoint{9.852599in}{1.152254in}}%
\pgfpathquadraticcurveto{\pgfqpoint{9.853283in}{1.141554in}}{\pgfqpoint{9.859047in}{1.135953in}}%
\pgfpathquadraticcurveto{\pgfqpoint{9.864811in}{1.130351in}}{\pgfqpoint{9.875114in}{1.130351in}}%
\pgfpathquadraticcurveto{\pgfqpoint{9.881076in}{1.130351in}}{\pgfqpoint{9.886678in}{1.131810in}}%
\pgfpathquadraticcurveto{\pgfqpoint{9.892280in}{1.133269in}}{\pgfqpoint{9.897809in}{1.136205in}}%
\pgfpathlineto{\pgfqpoint{9.897809in}{1.126406in}}%
\pgfpathquadraticcurveto{\pgfqpoint{9.892226in}{1.124047in}}{\pgfqpoint{9.886372in}{1.122804in}}%
\pgfpathquadraticcurveto{\pgfqpoint{9.880518in}{1.121561in}}{\pgfqpoint{9.874502in}{1.121561in}}%
\pgfpathquadraticcurveto{\pgfqpoint{9.859408in}{1.121561in}}{\pgfqpoint{9.850600in}{1.130333in}}%
\pgfpathquadraticcurveto{\pgfqpoint{9.841792in}{1.139123in}}{\pgfqpoint{9.841792in}{1.154109in}}%
\pgfpathquadraticcurveto{\pgfqpoint{9.841792in}{1.169581in}}{\pgfqpoint{9.850149in}{1.178659in}}%
\pgfpathquadraticcurveto{\pgfqpoint{9.858507in}{1.187756in}}{\pgfqpoint{9.872701in}{1.187756in}}%
\pgfpathquadraticcurveto{\pgfqpoint{9.885417in}{1.187756in}}{\pgfqpoint{9.892820in}{1.179560in}}%
\pgfpathquadraticcurveto{\pgfqpoint{9.900223in}{1.171383in}}{\pgfqpoint{9.900223in}{1.157315in}}%
\pgfpathlineto{\pgfqpoint{9.900223in}{1.157315in}}%
\pgfpathclose%
\pgfpathmoveto{\pgfqpoint{9.889866in}{1.160359in}}%
\pgfpathquadraticcurveto{\pgfqpoint{9.889758in}{1.168843in}}{\pgfqpoint{9.885111in}{1.173904in}}%
\pgfpathquadraticcurveto{\pgfqpoint{9.880464in}{1.178984in}}{\pgfqpoint{9.872809in}{1.178984in}}%
\pgfpathquadraticcurveto{\pgfqpoint{9.864145in}{1.178984in}}{\pgfqpoint{9.858939in}{1.174084in}}%
\pgfpathquadraticcurveto{\pgfqpoint{9.853734in}{1.169185in}}{\pgfqpoint{9.852941in}{1.160287in}}%
\pgfpathlineto{\pgfqpoint{9.889866in}{1.160359in}}%
\pgfpathlineto{\pgfqpoint{9.889866in}{1.160359in}}%
\pgfpathclose%
\pgfpathmoveto{\pgfqpoint{9.945884in}{1.154883in}}%
\pgfpathquadraticcurveto{\pgfqpoint{9.933329in}{1.154883in}}{\pgfqpoint{9.928484in}{1.152019in}}%
\pgfpathquadraticcurveto{\pgfqpoint{9.923639in}{1.149156in}}{\pgfqpoint{9.923639in}{1.142221in}}%
\pgfpathquadraticcurveto{\pgfqpoint{9.923639in}{1.136709in}}{\pgfqpoint{9.927277in}{1.133467in}}%
\pgfpathquadraticcurveto{\pgfqpoint{9.930916in}{1.130243in}}{\pgfqpoint{9.937148in}{1.130243in}}%
\pgfpathquadraticcurveto{\pgfqpoint{9.945776in}{1.130243in}}{\pgfqpoint{9.950981in}{1.136349in}}%
\pgfpathquadraticcurveto{\pgfqpoint{9.956187in}{1.142455in}}{\pgfqpoint{9.956187in}{1.152578in}}%
\pgfpathlineto{\pgfqpoint{9.956187in}{1.154883in}}%
\pgfpathlineto{\pgfqpoint{9.945884in}{1.154883in}}%
\pgfpathlineto{\pgfqpoint{9.945884in}{1.154883in}}%
\pgfpathclose%
\pgfpathmoveto{\pgfqpoint{9.966544in}{1.159170in}}%
\pgfpathlineto{\pgfqpoint{9.966544in}{1.123200in}}%
\pgfpathlineto{\pgfqpoint{9.956187in}{1.123200in}}%
\pgfpathlineto{\pgfqpoint{9.956187in}{1.132764in}}%
\pgfpathquadraticcurveto{\pgfqpoint{9.952638in}{1.127037in}}{\pgfqpoint{9.947343in}{1.124299in}}%
\pgfpathquadraticcurveto{\pgfqpoint{9.942047in}{1.121561in}}{\pgfqpoint{9.934392in}{1.121561in}}%
\pgfpathquadraticcurveto{\pgfqpoint{9.924720in}{1.121561in}}{\pgfqpoint{9.918992in}{1.127001in}}%
\pgfpathquadraticcurveto{\pgfqpoint{9.913282in}{1.132440in}}{\pgfqpoint{9.913282in}{1.141554in}}%
\pgfpathquadraticcurveto{\pgfqpoint{9.913282in}{1.152182in}}{\pgfqpoint{9.920397in}{1.157585in}}%
\pgfpathquadraticcurveto{\pgfqpoint{9.927530in}{1.162989in}}{\pgfqpoint{9.941651in}{1.162989in}}%
\pgfpathlineto{\pgfqpoint{9.956187in}{1.162989in}}%
\pgfpathlineto{\pgfqpoint{9.956187in}{1.164016in}}%
\pgfpathquadraticcurveto{\pgfqpoint{9.956187in}{1.171166in}}{\pgfqpoint{9.951486in}{1.175075in}}%
\pgfpathquadraticcurveto{\pgfqpoint{9.946785in}{1.178984in}}{\pgfqpoint{9.938283in}{1.178984in}}%
\pgfpathquadraticcurveto{\pgfqpoint{9.932879in}{1.178984in}}{\pgfqpoint{9.927746in}{1.177687in}}%
\pgfpathquadraticcurveto{\pgfqpoint{9.922630in}{1.176390in}}{\pgfqpoint{9.917911in}{1.173796in}}%
\pgfpathlineto{\pgfqpoint{9.917911in}{1.183379in}}%
\pgfpathquadraticcurveto{\pgfqpoint{9.923585in}{1.185576in}}{\pgfqpoint{9.928934in}{1.186657in}}%
\pgfpathquadraticcurveto{\pgfqpoint{9.934284in}{1.187756in}}{\pgfqpoint{9.939345in}{1.187756in}}%
\pgfpathquadraticcurveto{\pgfqpoint{9.953035in}{1.187756in}}{\pgfqpoint{9.959789in}{1.180659in}}%
\pgfpathquadraticcurveto{\pgfqpoint{9.966544in}{1.173580in}}{\pgfqpoint{9.966544in}{1.159170in}}%
\pgfpathlineto{\pgfqpoint{9.966544in}{1.159170in}}%
\pgfpathclose%
\pgfpathmoveto{\pgfqpoint{9.998119in}{1.204147in}}%
\pgfpathlineto{\pgfqpoint{9.998119in}{1.186243in}}%
\pgfpathlineto{\pgfqpoint{10.019446in}{1.186243in}}%
\pgfpathlineto{\pgfqpoint{10.019446in}{1.178191in}}%
\pgfpathlineto{\pgfqpoint{9.998119in}{1.178191in}}%
\pgfpathlineto{\pgfqpoint{9.998119in}{1.143968in}}%
\pgfpathquadraticcurveto{\pgfqpoint{9.998119in}{1.136259in}}{\pgfqpoint{10.000227in}{1.134061in}}%
\pgfpathquadraticcurveto{\pgfqpoint{10.002334in}{1.131864in}}{\pgfqpoint{10.008818in}{1.131864in}}%
\pgfpathlineto{\pgfqpoint{10.019446in}{1.131864in}}%
\pgfpathlineto{\pgfqpoint{10.019446in}{1.123200in}}%
\pgfpathlineto{\pgfqpoint{10.008818in}{1.123200in}}%
\pgfpathquadraticcurveto{\pgfqpoint{9.996822in}{1.123200in}}{\pgfqpoint{9.992265in}{1.127667in}}%
\pgfpathquadraticcurveto{\pgfqpoint{9.987708in}{1.132152in}}{\pgfqpoint{9.987708in}{1.143968in}}%
\pgfpathlineto{\pgfqpoint{9.987708in}{1.178191in}}%
\pgfpathlineto{\pgfqpoint{9.980107in}{1.178191in}}%
\pgfpathlineto{\pgfqpoint{9.980107in}{1.186243in}}%
\pgfpathlineto{\pgfqpoint{9.987708in}{1.186243in}}%
\pgfpathlineto{\pgfqpoint{9.987708in}{1.204147in}}%
\pgfpathlineto{\pgfqpoint{9.998119in}{1.204147in}}%
\pgfpathlineto{\pgfqpoint{9.998119in}{1.204147in}}%
\pgfpathclose%
\pgfpathmoveto{\pgfqpoint{10.082146in}{1.174138in}}%
\pgfpathquadraticcurveto{\pgfqpoint{10.086036in}{1.181127in}}{\pgfqpoint{10.091440in}{1.184441in}}%
\pgfpathquadraticcurveto{\pgfqpoint{10.096844in}{1.187756in}}{\pgfqpoint{10.104157in}{1.187756in}}%
\pgfpathquadraticcurveto{\pgfqpoint{10.114009in}{1.187756in}}{\pgfqpoint{10.119359in}{1.180857in}}%
\pgfpathquadraticcurveto{\pgfqpoint{10.124709in}{1.173976in}}{\pgfqpoint{10.124709in}{1.161260in}}%
\pgfpathlineto{\pgfqpoint{10.124709in}{1.123200in}}%
\pgfpathlineto{\pgfqpoint{10.114298in}{1.123200in}}%
\pgfpathlineto{\pgfqpoint{10.114298in}{1.160917in}}%
\pgfpathquadraticcurveto{\pgfqpoint{10.114298in}{1.169978in}}{\pgfqpoint{10.111073in}{1.174355in}}%
\pgfpathquadraticcurveto{\pgfqpoint{10.107867in}{1.178749in}}{\pgfqpoint{10.101293in}{1.178749in}}%
\pgfpathquadraticcurveto{\pgfqpoint{10.093241in}{1.178749in}}{\pgfqpoint{10.088558in}{1.173400in}}%
\pgfpathquadraticcurveto{\pgfqpoint{10.083893in}{1.168068in}}{\pgfqpoint{10.083893in}{1.158828in}}%
\pgfpathlineto{\pgfqpoint{10.083893in}{1.123200in}}%
\pgfpathlineto{\pgfqpoint{10.073482in}{1.123200in}}%
\pgfpathlineto{\pgfqpoint{10.073482in}{1.160917in}}%
\pgfpathquadraticcurveto{\pgfqpoint{10.073482in}{1.170032in}}{\pgfqpoint{10.070276in}{1.174391in}}%
\pgfpathquadraticcurveto{\pgfqpoint{10.067070in}{1.178749in}}{\pgfqpoint{10.060369in}{1.178749in}}%
\pgfpathquadraticcurveto{\pgfqpoint{10.052426in}{1.178749in}}{\pgfqpoint{10.047743in}{1.173382in}}%
\pgfpathquadraticcurveto{\pgfqpoint{10.043077in}{1.168014in}}{\pgfqpoint{10.043077in}{1.158828in}}%
\pgfpathlineto{\pgfqpoint{10.043077in}{1.123200in}}%
\pgfpathlineto{\pgfqpoint{10.032666in}{1.123200in}}%
\pgfpathlineto{\pgfqpoint{10.032666in}{1.186243in}}%
\pgfpathlineto{\pgfqpoint{10.043077in}{1.186243in}}%
\pgfpathlineto{\pgfqpoint{10.043077in}{1.176444in}}%
\pgfpathquadraticcurveto{\pgfqpoint{10.046626in}{1.182244in}}{\pgfqpoint{10.051579in}{1.185000in}}%
\pgfpathquadraticcurveto{\pgfqpoint{10.056533in}{1.187756in}}{\pgfqpoint{10.063341in}{1.187756in}}%
\pgfpathquadraticcurveto{\pgfqpoint{10.070222in}{1.187756in}}{\pgfqpoint{10.075031in}{1.184261in}}%
\pgfpathquadraticcurveto{\pgfqpoint{10.079840in}{1.180785in}}{\pgfqpoint{10.082146in}{1.174138in}}%
\pgfpathlineto{\pgfqpoint{10.082146in}{1.174138in}}%
\pgfpathclose%
\pgfpathmoveto{\pgfqpoint{10.199279in}{1.157315in}}%
\pgfpathlineto{\pgfqpoint{10.199279in}{1.152254in}}%
\pgfpathlineto{\pgfqpoint{10.151655in}{1.152254in}}%
\pgfpathquadraticcurveto{\pgfqpoint{10.152339in}{1.141554in}}{\pgfqpoint{10.158103in}{1.135953in}}%
\pgfpathquadraticcurveto{\pgfqpoint{10.163867in}{1.130351in}}{\pgfqpoint{10.174170in}{1.130351in}}%
\pgfpathquadraticcurveto{\pgfqpoint{10.180132in}{1.130351in}}{\pgfqpoint{10.185734in}{1.131810in}}%
\pgfpathquadraticcurveto{\pgfqpoint{10.191336in}{1.133269in}}{\pgfqpoint{10.196865in}{1.136205in}}%
\pgfpathlineto{\pgfqpoint{10.196865in}{1.126406in}}%
\pgfpathquadraticcurveto{\pgfqpoint{10.191281in}{1.124047in}}{\pgfqpoint{10.185428in}{1.122804in}}%
\pgfpathquadraticcurveto{\pgfqpoint{10.179574in}{1.121561in}}{\pgfqpoint{10.173558in}{1.121561in}}%
\pgfpathquadraticcurveto{\pgfqpoint{10.158463in}{1.121561in}}{\pgfqpoint{10.149655in}{1.130333in}}%
\pgfpathquadraticcurveto{\pgfqpoint{10.140847in}{1.139123in}}{\pgfqpoint{10.140847in}{1.154109in}}%
\pgfpathquadraticcurveto{\pgfqpoint{10.140847in}{1.169581in}}{\pgfqpoint{10.149205in}{1.178659in}}%
\pgfpathquadraticcurveto{\pgfqpoint{10.157563in}{1.187756in}}{\pgfqpoint{10.171756in}{1.187756in}}%
\pgfpathquadraticcurveto{\pgfqpoint{10.184473in}{1.187756in}}{\pgfqpoint{10.191876in}{1.179560in}}%
\pgfpathquadraticcurveto{\pgfqpoint{10.199279in}{1.171383in}}{\pgfqpoint{10.199279in}{1.157315in}}%
\pgfpathlineto{\pgfqpoint{10.199279in}{1.157315in}}%
\pgfpathclose%
\pgfpathmoveto{\pgfqpoint{10.188922in}{1.160359in}}%
\pgfpathquadraticcurveto{\pgfqpoint{10.188814in}{1.168843in}}{\pgfqpoint{10.184167in}{1.173904in}}%
\pgfpathquadraticcurveto{\pgfqpoint{10.179520in}{1.178984in}}{\pgfqpoint{10.171864in}{1.178984in}}%
\pgfpathquadraticcurveto{\pgfqpoint{10.163201in}{1.178984in}}{\pgfqpoint{10.157995in}{1.174084in}}%
\pgfpathquadraticcurveto{\pgfqpoint{10.152789in}{1.169185in}}{\pgfqpoint{10.151997in}{1.160287in}}%
\pgfpathlineto{\pgfqpoint{10.188922in}{1.160359in}}%
\pgfpathlineto{\pgfqpoint{10.188922in}{1.160359in}}%
\pgfpathclose%
\pgfpathmoveto{\pgfqpoint{10.268698in}{1.161260in}}%
\pgfpathlineto{\pgfqpoint{10.268698in}{1.123200in}}%
\pgfpathlineto{\pgfqpoint{10.258341in}{1.123200in}}%
\pgfpathlineto{\pgfqpoint{10.258341in}{1.160917in}}%
\pgfpathquadraticcurveto{\pgfqpoint{10.258341in}{1.169869in}}{\pgfqpoint{10.254846in}{1.174300in}}%
\pgfpathquadraticcurveto{\pgfqpoint{10.251352in}{1.178749in}}{\pgfqpoint{10.244381in}{1.178749in}}%
\pgfpathquadraticcurveto{\pgfqpoint{10.235988in}{1.178749in}}{\pgfqpoint{10.231142in}{1.173400in}}%
\pgfpathquadraticcurveto{\pgfqpoint{10.226297in}{1.168068in}}{\pgfqpoint{10.226297in}{1.158828in}}%
\pgfpathlineto{\pgfqpoint{10.226297in}{1.123200in}}%
\pgfpathlineto{\pgfqpoint{10.215886in}{1.123200in}}%
\pgfpathlineto{\pgfqpoint{10.215886in}{1.186243in}}%
\pgfpathlineto{\pgfqpoint{10.226297in}{1.186243in}}%
\pgfpathlineto{\pgfqpoint{10.226297in}{1.176444in}}%
\pgfpathquadraticcurveto{\pgfqpoint{10.230026in}{1.182136in}}{\pgfqpoint{10.235051in}{1.184946in}}%
\pgfpathquadraticcurveto{\pgfqpoint{10.240094in}{1.187756in}}{\pgfqpoint{10.246687in}{1.187756in}}%
\pgfpathquadraticcurveto{\pgfqpoint{10.257548in}{1.187756in}}{\pgfqpoint{10.263114in}{1.181037in}}%
\pgfpathquadraticcurveto{\pgfqpoint{10.268698in}{1.174318in}}{\pgfqpoint{10.268698in}{1.161260in}}%
\pgfpathlineto{\pgfqpoint{10.268698in}{1.161260in}}%
\pgfpathclose%
\pgfpathmoveto{\pgfqpoint{10.299589in}{1.204147in}}%
\pgfpathlineto{\pgfqpoint{10.299589in}{1.186243in}}%
\pgfpathlineto{\pgfqpoint{10.320915in}{1.186243in}}%
\pgfpathlineto{\pgfqpoint{10.320915in}{1.178191in}}%
\pgfpathlineto{\pgfqpoint{10.299589in}{1.178191in}}%
\pgfpathlineto{\pgfqpoint{10.299589in}{1.143968in}}%
\pgfpathquadraticcurveto{\pgfqpoint{10.299589in}{1.136259in}}{\pgfqpoint{10.301696in}{1.134061in}}%
\pgfpathquadraticcurveto{\pgfqpoint{10.303803in}{1.131864in}}{\pgfqpoint{10.310288in}{1.131864in}}%
\pgfpathlineto{\pgfqpoint{10.320915in}{1.131864in}}%
\pgfpathlineto{\pgfqpoint{10.320915in}{1.123200in}}%
\pgfpathlineto{\pgfqpoint{10.310288in}{1.123200in}}%
\pgfpathquadraticcurveto{\pgfqpoint{10.298292in}{1.123200in}}{\pgfqpoint{10.293735in}{1.127667in}}%
\pgfpathquadraticcurveto{\pgfqpoint{10.289178in}{1.132152in}}{\pgfqpoint{10.289178in}{1.143968in}}%
\pgfpathlineto{\pgfqpoint{10.289178in}{1.178191in}}%
\pgfpathlineto{\pgfqpoint{10.281576in}{1.178191in}}%
\pgfpathlineto{\pgfqpoint{10.281576in}{1.186243in}}%
\pgfpathlineto{\pgfqpoint{10.289178in}{1.186243in}}%
\pgfpathlineto{\pgfqpoint{10.289178in}{1.204147in}}%
\pgfpathlineto{\pgfqpoint{10.299589in}{1.204147in}}%
\pgfpathlineto{\pgfqpoint{10.299589in}{1.204147in}}%
\pgfpathclose%
\pgfusepath{fill}%
\end{pgfscope}%
\begin{pgfscope}%
\pgfsetbuttcap%
\pgfsetroundjoin%
\definecolor{currentfill}{rgb}{0.309804,0.372549,0.419608}%
\pgfsetfillcolor{currentfill}%
\pgfsetlinewidth{0.000000pt}%
\definecolor{currentstroke}{rgb}{0.309804,0.372549,0.419608}%
\pgfsetstrokecolor{currentstroke}%
\pgfsetdash{}{0pt}%
\pgfpathmoveto{\pgfqpoint{10.476699in}{1.157315in}}%
\pgfpathlineto{\pgfqpoint{10.476699in}{1.152254in}}%
\pgfpathlineto{\pgfqpoint{10.429074in}{1.152254in}}%
\pgfpathquadraticcurveto{\pgfqpoint{10.429759in}{1.141554in}}{\pgfqpoint{10.435523in}{1.135953in}}%
\pgfpathquadraticcurveto{\pgfqpoint{10.441287in}{1.130351in}}{\pgfqpoint{10.451590in}{1.130351in}}%
\pgfpathquadraticcurveto{\pgfqpoint{10.457552in}{1.130351in}}{\pgfqpoint{10.463153in}{1.131810in}}%
\pgfpathquadraticcurveto{\pgfqpoint{10.468755in}{1.133269in}}{\pgfqpoint{10.474285in}{1.136205in}}%
\pgfpathlineto{\pgfqpoint{10.474285in}{1.126406in}}%
\pgfpathquadraticcurveto{\pgfqpoint{10.468701in}{1.124047in}}{\pgfqpoint{10.462847in}{1.122804in}}%
\pgfpathquadraticcurveto{\pgfqpoint{10.456993in}{1.121561in}}{\pgfqpoint{10.450977in}{1.121561in}}%
\pgfpathquadraticcurveto{\pgfqpoint{10.435883in}{1.121561in}}{\pgfqpoint{10.427075in}{1.130333in}}%
\pgfpathquadraticcurveto{\pgfqpoint{10.418267in}{1.139123in}}{\pgfqpoint{10.418267in}{1.154109in}}%
\pgfpathquadraticcurveto{\pgfqpoint{10.418267in}{1.169581in}}{\pgfqpoint{10.426625in}{1.178659in}}%
\pgfpathquadraticcurveto{\pgfqpoint{10.434982in}{1.187756in}}{\pgfqpoint{10.449176in}{1.187756in}}%
\pgfpathquadraticcurveto{\pgfqpoint{10.461893in}{1.187756in}}{\pgfqpoint{10.469296in}{1.179560in}}%
\pgfpathquadraticcurveto{\pgfqpoint{10.476699in}{1.171383in}}{\pgfqpoint{10.476699in}{1.157315in}}%
\pgfpathlineto{\pgfqpoint{10.476699in}{1.157315in}}%
\pgfpathclose%
\pgfpathmoveto{\pgfqpoint{10.466342in}{1.160359in}}%
\pgfpathquadraticcurveto{\pgfqpoint{10.466234in}{1.168843in}}{\pgfqpoint{10.461586in}{1.173904in}}%
\pgfpathquadraticcurveto{\pgfqpoint{10.456939in}{1.178984in}}{\pgfqpoint{10.449284in}{1.178984in}}%
\pgfpathquadraticcurveto{\pgfqpoint{10.440620in}{1.178984in}}{\pgfqpoint{10.435415in}{1.174084in}}%
\pgfpathquadraticcurveto{\pgfqpoint{10.430209in}{1.169185in}}{\pgfqpoint{10.429417in}{1.160287in}}%
\pgfpathlineto{\pgfqpoint{10.466342in}{1.160359in}}%
\pgfpathlineto{\pgfqpoint{10.466342in}{1.160359in}}%
\pgfpathclose%
\pgfpathmoveto{\pgfqpoint{10.564454in}{1.210793in}}%
\pgfpathlineto{\pgfqpoint{10.564454in}{1.202165in}}%
\pgfpathlineto{\pgfqpoint{10.554547in}{1.202165in}}%
\pgfpathquadraticcurveto{\pgfqpoint{10.548981in}{1.202165in}}{\pgfqpoint{10.546802in}{1.199914in}}%
\pgfpathquadraticcurveto{\pgfqpoint{10.544640in}{1.197662in}}{\pgfqpoint{10.544640in}{1.191808in}}%
\pgfpathlineto{\pgfqpoint{10.544640in}{1.186243in}}%
\pgfpathlineto{\pgfqpoint{10.561698in}{1.186243in}}%
\pgfpathlineto{\pgfqpoint{10.561698in}{1.178191in}}%
\pgfpathlineto{\pgfqpoint{10.544640in}{1.178191in}}%
\pgfpathlineto{\pgfqpoint{10.544640in}{1.123200in}}%
\pgfpathlineto{\pgfqpoint{10.534229in}{1.123200in}}%
\pgfpathlineto{\pgfqpoint{10.534229in}{1.178191in}}%
\pgfpathlineto{\pgfqpoint{10.505806in}{1.178191in}}%
\pgfpathlineto{\pgfqpoint{10.505806in}{1.123200in}}%
\pgfpathlineto{\pgfqpoint{10.495395in}{1.123200in}}%
\pgfpathlineto{\pgfqpoint{10.495395in}{1.178191in}}%
\pgfpathlineto{\pgfqpoint{10.485488in}{1.178191in}}%
\pgfpathlineto{\pgfqpoint{10.485488in}{1.186243in}}%
\pgfpathlineto{\pgfqpoint{10.495395in}{1.186243in}}%
\pgfpathlineto{\pgfqpoint{10.495395in}{1.190638in}}%
\pgfpathquadraticcurveto{\pgfqpoint{10.495395in}{1.201157in}}{\pgfqpoint{10.500294in}{1.205966in}}%
\pgfpathquadraticcurveto{\pgfqpoint{10.505194in}{1.210793in}}{\pgfqpoint{10.515821in}{1.210793in}}%
\pgfpathlineto{\pgfqpoint{10.525620in}{1.210793in}}%
\pgfpathlineto{\pgfqpoint{10.525620in}{1.202165in}}%
\pgfpathlineto{\pgfqpoint{10.515713in}{1.202165in}}%
\pgfpathquadraticcurveto{\pgfqpoint{10.510147in}{1.202165in}}{\pgfqpoint{10.507950in}{1.199914in}}%
\pgfpathquadraticcurveto{\pgfqpoint{10.505806in}{1.197662in}}{\pgfqpoint{10.505806in}{1.191808in}}%
\pgfpathlineto{\pgfqpoint{10.505806in}{1.186243in}}%
\pgfpathlineto{\pgfqpoint{10.534229in}{1.186243in}}%
\pgfpathlineto{\pgfqpoint{10.534229in}{1.190638in}}%
\pgfpathquadraticcurveto{\pgfqpoint{10.534229in}{1.201157in}}{\pgfqpoint{10.539129in}{1.205966in}}%
\pgfpathquadraticcurveto{\pgfqpoint{10.544028in}{1.210793in}}{\pgfqpoint{10.554673in}{1.210793in}}%
\pgfpathlineto{\pgfqpoint{10.564454in}{1.210793in}}%
\pgfpathlineto{\pgfqpoint{10.564454in}{1.210793in}}%
\pgfpathclose%
\pgfpathmoveto{\pgfqpoint{10.627046in}{1.157315in}}%
\pgfpathlineto{\pgfqpoint{10.627046in}{1.152254in}}%
\pgfpathlineto{\pgfqpoint{10.579422in}{1.152254in}}%
\pgfpathquadraticcurveto{\pgfqpoint{10.580106in}{1.141554in}}{\pgfqpoint{10.585870in}{1.135953in}}%
\pgfpathquadraticcurveto{\pgfqpoint{10.591634in}{1.130351in}}{\pgfqpoint{10.601937in}{1.130351in}}%
\pgfpathquadraticcurveto{\pgfqpoint{10.607899in}{1.130351in}}{\pgfqpoint{10.613501in}{1.131810in}}%
\pgfpathquadraticcurveto{\pgfqpoint{10.619103in}{1.133269in}}{\pgfqpoint{10.624632in}{1.136205in}}%
\pgfpathlineto{\pgfqpoint{10.624632in}{1.126406in}}%
\pgfpathquadraticcurveto{\pgfqpoint{10.619049in}{1.124047in}}{\pgfqpoint{10.613195in}{1.122804in}}%
\pgfpathquadraticcurveto{\pgfqpoint{10.607341in}{1.121561in}}{\pgfqpoint{10.601325in}{1.121561in}}%
\pgfpathquadraticcurveto{\pgfqpoint{10.586230in}{1.121561in}}{\pgfqpoint{10.577423in}{1.130333in}}%
\pgfpathquadraticcurveto{\pgfqpoint{10.568615in}{1.139123in}}{\pgfqpoint{10.568615in}{1.154109in}}%
\pgfpathquadraticcurveto{\pgfqpoint{10.568615in}{1.169581in}}{\pgfqpoint{10.576972in}{1.178659in}}%
\pgfpathquadraticcurveto{\pgfqpoint{10.585330in}{1.187756in}}{\pgfqpoint{10.599523in}{1.187756in}}%
\pgfpathquadraticcurveto{\pgfqpoint{10.612240in}{1.187756in}}{\pgfqpoint{10.619643in}{1.179560in}}%
\pgfpathquadraticcurveto{\pgfqpoint{10.627046in}{1.171383in}}{\pgfqpoint{10.627046in}{1.157315in}}%
\pgfpathlineto{\pgfqpoint{10.627046in}{1.157315in}}%
\pgfpathclose%
\pgfpathmoveto{\pgfqpoint{10.616689in}{1.160359in}}%
\pgfpathquadraticcurveto{\pgfqpoint{10.616581in}{1.168843in}}{\pgfqpoint{10.611934in}{1.173904in}}%
\pgfpathquadraticcurveto{\pgfqpoint{10.607287in}{1.178984in}}{\pgfqpoint{10.599632in}{1.178984in}}%
\pgfpathquadraticcurveto{\pgfqpoint{10.590968in}{1.178984in}}{\pgfqpoint{10.585762in}{1.174084in}}%
\pgfpathquadraticcurveto{\pgfqpoint{10.580557in}{1.169185in}}{\pgfqpoint{10.579764in}{1.160287in}}%
\pgfpathlineto{\pgfqpoint{10.616689in}{1.160359in}}%
\pgfpathlineto{\pgfqpoint{10.616689in}{1.160359in}}%
\pgfpathclose%
\pgfpathmoveto{\pgfqpoint{10.689422in}{1.183829in}}%
\pgfpathlineto{\pgfqpoint{10.689422in}{1.174138in}}%
\pgfpathquadraticcurveto{\pgfqpoint{10.685027in}{1.176570in}}{\pgfqpoint{10.680614in}{1.177777in}}%
\pgfpathquadraticcurveto{\pgfqpoint{10.676201in}{1.178984in}}{\pgfqpoint{10.671698in}{1.178984in}}%
\pgfpathquadraticcurveto{\pgfqpoint{10.661611in}{1.178984in}}{\pgfqpoint{10.656028in}{1.172589in}}%
\pgfpathquadraticcurveto{\pgfqpoint{10.650462in}{1.166213in}}{\pgfqpoint{10.650462in}{1.154667in}}%
\pgfpathquadraticcurveto{\pgfqpoint{10.650462in}{1.143121in}}{\pgfqpoint{10.656028in}{1.136727in}}%
\pgfpathquadraticcurveto{\pgfqpoint{10.661611in}{1.130351in}}{\pgfqpoint{10.671698in}{1.130351in}}%
\pgfpathquadraticcurveto{\pgfqpoint{10.676201in}{1.130351in}}{\pgfqpoint{10.680614in}{1.131558in}}%
\pgfpathquadraticcurveto{\pgfqpoint{10.685027in}{1.132764in}}{\pgfqpoint{10.689422in}{1.135196in}}%
\pgfpathlineto{\pgfqpoint{10.689422in}{1.125614in}}%
\pgfpathquadraticcurveto{\pgfqpoint{10.685081in}{1.123596in}}{\pgfqpoint{10.680434in}{1.122588in}}%
\pgfpathquadraticcurveto{\pgfqpoint{10.675805in}{1.121561in}}{\pgfqpoint{10.670563in}{1.121561in}}%
\pgfpathquadraticcurveto{\pgfqpoint{10.656316in}{1.121561in}}{\pgfqpoint{10.647922in}{1.130513in}}%
\pgfpathquadraticcurveto{\pgfqpoint{10.639546in}{1.139465in}}{\pgfqpoint{10.639546in}{1.154667in}}%
\pgfpathquadraticcurveto{\pgfqpoint{10.639546in}{1.170086in}}{\pgfqpoint{10.648012in}{1.178912in}}%
\pgfpathquadraticcurveto{\pgfqpoint{10.656496in}{1.187756in}}{\pgfqpoint{10.671248in}{1.187756in}}%
\pgfpathquadraticcurveto{\pgfqpoint{10.676021in}{1.187756in}}{\pgfqpoint{10.680578in}{1.186765in}}%
\pgfpathquadraticcurveto{\pgfqpoint{10.685135in}{1.185792in}}{\pgfqpoint{10.689422in}{1.183829in}}%
\pgfpathlineto{\pgfqpoint{10.689422in}{1.183829in}}%
\pgfpathclose%
\pgfpathmoveto{\pgfqpoint{10.717683in}{1.204147in}}%
\pgfpathlineto{\pgfqpoint{10.717683in}{1.186243in}}%
\pgfpathlineto{\pgfqpoint{10.739010in}{1.186243in}}%
\pgfpathlineto{\pgfqpoint{10.739010in}{1.178191in}}%
\pgfpathlineto{\pgfqpoint{10.717683in}{1.178191in}}%
\pgfpathlineto{\pgfqpoint{10.717683in}{1.143968in}}%
\pgfpathquadraticcurveto{\pgfqpoint{10.717683in}{1.136259in}}{\pgfqpoint{10.719791in}{1.134061in}}%
\pgfpathquadraticcurveto{\pgfqpoint{10.721898in}{1.131864in}}{\pgfqpoint{10.728382in}{1.131864in}}%
\pgfpathlineto{\pgfqpoint{10.739010in}{1.131864in}}%
\pgfpathlineto{\pgfqpoint{10.739010in}{1.123200in}}%
\pgfpathlineto{\pgfqpoint{10.728382in}{1.123200in}}%
\pgfpathquadraticcurveto{\pgfqpoint{10.716386in}{1.123200in}}{\pgfqpoint{10.711829in}{1.127667in}}%
\pgfpathquadraticcurveto{\pgfqpoint{10.707272in}{1.132152in}}{\pgfqpoint{10.707272in}{1.143968in}}%
\pgfpathlineto{\pgfqpoint{10.707272in}{1.178191in}}%
\pgfpathlineto{\pgfqpoint{10.699671in}{1.178191in}}%
\pgfpathlineto{\pgfqpoint{10.699671in}{1.186243in}}%
\pgfpathlineto{\pgfqpoint{10.707272in}{1.186243in}}%
\pgfpathlineto{\pgfqpoint{10.707272in}{1.204147in}}%
\pgfpathlineto{\pgfqpoint{10.717683in}{1.204147in}}%
\pgfpathlineto{\pgfqpoint{10.717683in}{1.204147in}}%
\pgfpathclose%
\pgfpathmoveto{\pgfqpoint{10.792812in}{1.184387in}}%
\pgfpathlineto{\pgfqpoint{10.792812in}{1.174589in}}%
\pgfpathquadraticcurveto{\pgfqpoint{10.788435in}{1.176840in}}{\pgfqpoint{10.783698in}{1.177957in}}%
\pgfpathquadraticcurveto{\pgfqpoint{10.778979in}{1.179092in}}{\pgfqpoint{10.773899in}{1.179092in}}%
\pgfpathquadraticcurveto{\pgfqpoint{10.766190in}{1.179092in}}{\pgfqpoint{10.762335in}{1.176732in}}%
\pgfpathquadraticcurveto{\pgfqpoint{10.758481in}{1.174373in}}{\pgfqpoint{10.758481in}{1.169635in}}%
\pgfpathquadraticcurveto{\pgfqpoint{10.758481in}{1.166033in}}{\pgfqpoint{10.761237in}{1.163980in}}%
\pgfpathquadraticcurveto{\pgfqpoint{10.763992in}{1.161926in}}{\pgfqpoint{10.772332in}{1.160071in}}%
\pgfpathlineto{\pgfqpoint{10.775880in}{1.159278in}}%
\pgfpathquadraticcurveto{\pgfqpoint{10.786904in}{1.156919in}}{\pgfqpoint{10.791551in}{1.152614in}}%
\pgfpathquadraticcurveto{\pgfqpoint{10.796198in}{1.148309in}}{\pgfqpoint{10.796198in}{1.140600in}}%
\pgfpathquadraticcurveto{\pgfqpoint{10.796198in}{1.131810in}}{\pgfqpoint{10.789245in}{1.126676in}}%
\pgfpathquadraticcurveto{\pgfqpoint{10.782293in}{1.121561in}}{\pgfqpoint{10.770135in}{1.121561in}}%
\pgfpathquadraticcurveto{\pgfqpoint{10.765073in}{1.121561in}}{\pgfqpoint{10.759579in}{1.122552in}}%
\pgfpathquadraticcurveto{\pgfqpoint{10.754086in}{1.123542in}}{\pgfqpoint{10.748016in}{1.125506in}}%
\pgfpathlineto{\pgfqpoint{10.748016in}{1.136205in}}%
\pgfpathquadraticcurveto{\pgfqpoint{10.753762in}{1.133215in}}{\pgfqpoint{10.759327in}{1.131720in}}%
\pgfpathquadraticcurveto{\pgfqpoint{10.764893in}{1.130243in}}{\pgfqpoint{10.770369in}{1.130243in}}%
\pgfpathquadraticcurveto{\pgfqpoint{10.777682in}{1.130243in}}{\pgfqpoint{10.781608in}{1.132746in}}%
\pgfpathquadraticcurveto{\pgfqpoint{10.785553in}{1.135250in}}{\pgfqpoint{10.785553in}{1.139807in}}%
\pgfpathquadraticcurveto{\pgfqpoint{10.785553in}{1.144022in}}{\pgfqpoint{10.782707in}{1.146274in}}%
\pgfpathquadraticcurveto{\pgfqpoint{10.779879in}{1.148525in}}{\pgfqpoint{10.770243in}{1.150614in}}%
\pgfpathlineto{\pgfqpoint{10.766640in}{1.151461in}}%
\pgfpathquadraticcurveto{\pgfqpoint{10.757022in}{1.153478in}}{\pgfqpoint{10.752735in}{1.157675in}}%
\pgfpathquadraticcurveto{\pgfqpoint{10.748466in}{1.161872in}}{\pgfqpoint{10.748466in}{1.169185in}}%
\pgfpathquadraticcurveto{\pgfqpoint{10.748466in}{1.178083in}}{\pgfqpoint{10.754770in}{1.182910in}}%
\pgfpathquadraticcurveto{\pgfqpoint{10.761074in}{1.187756in}}{\pgfqpoint{10.772674in}{1.187756in}}%
\pgfpathquadraticcurveto{\pgfqpoint{10.778402in}{1.187756in}}{\pgfqpoint{10.783464in}{1.186909in}}%
\pgfpathquadraticcurveto{\pgfqpoint{10.788543in}{1.186080in}}{\pgfqpoint{10.792812in}{1.184387in}}%
\pgfpathlineto{\pgfqpoint{10.792812in}{1.184387in}}%
\pgfpathclose%
\pgfpathmoveto{\pgfqpoint{10.814138in}{1.137502in}}%
\pgfpathlineto{\pgfqpoint{10.826026in}{1.137502in}}%
\pgfpathlineto{\pgfqpoint{10.826026in}{1.123200in}}%
\pgfpathlineto{\pgfqpoint{10.814138in}{1.123200in}}%
\pgfpathlineto{\pgfqpoint{10.814138in}{1.137502in}}%
\pgfpathlineto{\pgfqpoint{10.814138in}{1.137502in}}%
\pgfpathclose%
\pgfusepath{fill}%
\end{pgfscope}%
\begin{pgfscope}%
\pgfsetbuttcap%
\pgfsetroundjoin%
\definecolor{currentfill}{rgb}{0.309804,0.372549,0.419608}%
\pgfsetfillcolor{currentfill}%
\pgfsetlinewidth{0.000000pt}%
\definecolor{currentstroke}{rgb}{0.309804,0.372549,0.419608}%
\pgfsetstrokecolor{currentstroke}%
\pgfsetdash{}{0pt}%
\pgfpathmoveto{\pgfqpoint{10.972619in}{1.204489in}}%
\pgfpathlineto{\pgfqpoint{10.972619in}{1.193393in}}%
\pgfpathquadraticcurveto{\pgfqpoint{10.966152in}{1.196491in}}{\pgfqpoint{10.960406in}{1.198004in}}%
\pgfpathquadraticcurveto{\pgfqpoint{10.954661in}{1.199536in}}{\pgfqpoint{10.949311in}{1.199536in}}%
\pgfpathquadraticcurveto{\pgfqpoint{10.940035in}{1.199536in}}{\pgfqpoint{10.934991in}{1.195933in}}%
\pgfpathquadraticcurveto{\pgfqpoint{10.929948in}{1.192331in}}{\pgfqpoint{10.929948in}{1.185684in}}%
\pgfpathquadraticcurveto{\pgfqpoint{10.929948in}{1.180100in}}{\pgfqpoint{10.933298in}{1.177254in}}%
\pgfpathquadraticcurveto{\pgfqpoint{10.936648in}{1.174427in}}{\pgfqpoint{10.945997in}{1.172679in}}%
\pgfpathlineto{\pgfqpoint{10.952859in}{1.171274in}}%
\pgfpathquadraticcurveto{\pgfqpoint{10.965576in}{1.168843in}}{\pgfqpoint{10.971628in}{1.162737in}}%
\pgfpathquadraticcurveto{\pgfqpoint{10.977680in}{1.156631in}}{\pgfqpoint{10.977680in}{1.146400in}}%
\pgfpathquadraticcurveto{\pgfqpoint{10.977680in}{1.134169in}}{\pgfqpoint{10.969485in}{1.127865in}}%
\pgfpathquadraticcurveto{\pgfqpoint{10.961307in}{1.121561in}}{\pgfqpoint{10.945492in}{1.121561in}}%
\pgfpathquadraticcurveto{\pgfqpoint{10.939530in}{1.121561in}}{\pgfqpoint{10.932794in}{1.122912in}}%
\pgfpathquadraticcurveto{\pgfqpoint{10.926075in}{1.124263in}}{\pgfqpoint{10.918870in}{1.126911in}}%
\pgfpathlineto{\pgfqpoint{10.918870in}{1.138618in}}%
\pgfpathquadraticcurveto{\pgfqpoint{10.925787in}{1.134746in}}{\pgfqpoint{10.932434in}{1.132764in}}%
\pgfpathquadraticcurveto{\pgfqpoint{10.939080in}{1.130801in}}{\pgfqpoint{10.945492in}{1.130801in}}%
\pgfpathquadraticcurveto{\pgfqpoint{10.955219in}{1.130801in}}{\pgfqpoint{10.960515in}{1.134620in}}%
\pgfpathquadraticcurveto{\pgfqpoint{10.965810in}{1.138456in}}{\pgfqpoint{10.965810in}{1.145553in}}%
\pgfpathquadraticcurveto{\pgfqpoint{10.965810in}{1.151731in}}{\pgfqpoint{10.962010in}{1.155226in}}%
\pgfpathquadraticcurveto{\pgfqpoint{10.958209in}{1.158720in}}{\pgfqpoint{10.949545in}{1.160467in}}%
\pgfpathlineto{\pgfqpoint{10.942610in}{1.161818in}}%
\pgfpathquadraticcurveto{\pgfqpoint{10.929894in}{1.164340in}}{\pgfqpoint{10.924202in}{1.169743in}}%
\pgfpathquadraticcurveto{\pgfqpoint{10.918528in}{1.175147in}}{\pgfqpoint{10.918528in}{1.184784in}}%
\pgfpathquadraticcurveto{\pgfqpoint{10.918528in}{1.195933in}}{\pgfqpoint{10.926382in}{1.202345in}}%
\pgfpathquadraticcurveto{\pgfqpoint{10.934235in}{1.208758in}}{\pgfqpoint{10.948014in}{1.208758in}}%
\pgfpathquadraticcurveto{\pgfqpoint{10.953940in}{1.208758in}}{\pgfqpoint{10.960064in}{1.207677in}}%
\pgfpathquadraticcurveto{\pgfqpoint{10.966206in}{1.206614in}}{\pgfqpoint{10.972619in}{1.204489in}}%
\pgfpathlineto{\pgfqpoint{10.972619in}{1.204489in}}%
\pgfpathclose%
\pgfpathmoveto{\pgfqpoint{11.047387in}{1.161260in}}%
\pgfpathlineto{\pgfqpoint{11.047387in}{1.123200in}}%
\pgfpathlineto{\pgfqpoint{11.037030in}{1.123200in}}%
\pgfpathlineto{\pgfqpoint{11.037030in}{1.160917in}}%
\pgfpathquadraticcurveto{\pgfqpoint{11.037030in}{1.169869in}}{\pgfqpoint{11.033536in}{1.174300in}}%
\pgfpathquadraticcurveto{\pgfqpoint{11.030041in}{1.178749in}}{\pgfqpoint{11.023071in}{1.178749in}}%
\pgfpathquadraticcurveto{\pgfqpoint{11.014677in}{1.178749in}}{\pgfqpoint{11.009832in}{1.173400in}}%
\pgfpathquadraticcurveto{\pgfqpoint{11.004987in}{1.168068in}}{\pgfqpoint{11.004987in}{1.158828in}}%
\pgfpathlineto{\pgfqpoint{11.004987in}{1.123200in}}%
\pgfpathlineto{\pgfqpoint{10.994576in}{1.123200in}}%
\pgfpathlineto{\pgfqpoint{10.994576in}{1.210793in}}%
\pgfpathlineto{\pgfqpoint{11.004987in}{1.210793in}}%
\pgfpathlineto{\pgfqpoint{11.004987in}{1.176444in}}%
\pgfpathquadraticcurveto{\pgfqpoint{11.008715in}{1.182136in}}{\pgfqpoint{11.013740in}{1.184946in}}%
\pgfpathquadraticcurveto{\pgfqpoint{11.018784in}{1.187756in}}{\pgfqpoint{11.025376in}{1.187756in}}%
\pgfpathquadraticcurveto{\pgfqpoint{11.036238in}{1.187756in}}{\pgfqpoint{11.041803in}{1.181037in}}%
\pgfpathquadraticcurveto{\pgfqpoint{11.047387in}{1.174318in}}{\pgfqpoint{11.047387in}{1.161260in}}%
\pgfpathlineto{\pgfqpoint{11.047387in}{1.161260in}}%
\pgfpathclose%
\pgfpathmoveto{\pgfqpoint{11.096686in}{1.154883in}}%
\pgfpathquadraticcurveto{\pgfqpoint{11.084132in}{1.154883in}}{\pgfqpoint{11.079287in}{1.152019in}}%
\pgfpathquadraticcurveto{\pgfqpoint{11.074441in}{1.149156in}}{\pgfqpoint{11.074441in}{1.142221in}}%
\pgfpathquadraticcurveto{\pgfqpoint{11.074441in}{1.136709in}}{\pgfqpoint{11.078080in}{1.133467in}}%
\pgfpathquadraticcurveto{\pgfqpoint{11.081718in}{1.130243in}}{\pgfqpoint{11.087951in}{1.130243in}}%
\pgfpathquadraticcurveto{\pgfqpoint{11.096578in}{1.130243in}}{\pgfqpoint{11.101784in}{1.136349in}}%
\pgfpathquadraticcurveto{\pgfqpoint{11.106989in}{1.142455in}}{\pgfqpoint{11.106989in}{1.152578in}}%
\pgfpathlineto{\pgfqpoint{11.106989in}{1.154883in}}%
\pgfpathlineto{\pgfqpoint{11.096686in}{1.154883in}}%
\pgfpathlineto{\pgfqpoint{11.096686in}{1.154883in}}%
\pgfpathclose%
\pgfpathmoveto{\pgfqpoint{11.117346in}{1.159170in}}%
\pgfpathlineto{\pgfqpoint{11.117346in}{1.123200in}}%
\pgfpathlineto{\pgfqpoint{11.106989in}{1.123200in}}%
\pgfpathlineto{\pgfqpoint{11.106989in}{1.132764in}}%
\pgfpathquadraticcurveto{\pgfqpoint{11.103441in}{1.127037in}}{\pgfqpoint{11.098145in}{1.124299in}}%
\pgfpathquadraticcurveto{\pgfqpoint{11.092850in}{1.121561in}}{\pgfqpoint{11.085195in}{1.121561in}}%
\pgfpathquadraticcurveto{\pgfqpoint{11.075522in}{1.121561in}}{\pgfqpoint{11.069794in}{1.127001in}}%
\pgfpathquadraticcurveto{\pgfqpoint{11.064084in}{1.132440in}}{\pgfqpoint{11.064084in}{1.141554in}}%
\pgfpathquadraticcurveto{\pgfqpoint{11.064084in}{1.152182in}}{\pgfqpoint{11.071199in}{1.157585in}}%
\pgfpathquadraticcurveto{\pgfqpoint{11.078332in}{1.162989in}}{\pgfqpoint{11.092454in}{1.162989in}}%
\pgfpathlineto{\pgfqpoint{11.106989in}{1.162989in}}%
\pgfpathlineto{\pgfqpoint{11.106989in}{1.164016in}}%
\pgfpathquadraticcurveto{\pgfqpoint{11.106989in}{1.171166in}}{\pgfqpoint{11.102288in}{1.175075in}}%
\pgfpathquadraticcurveto{\pgfqpoint{11.097587in}{1.178984in}}{\pgfqpoint{11.089085in}{1.178984in}}%
\pgfpathquadraticcurveto{\pgfqpoint{11.083682in}{1.178984in}}{\pgfqpoint{11.078548in}{1.177687in}}%
\pgfpathquadraticcurveto{\pgfqpoint{11.073433in}{1.176390in}}{\pgfqpoint{11.068714in}{1.173796in}}%
\pgfpathlineto{\pgfqpoint{11.068714in}{1.183379in}}%
\pgfpathquadraticcurveto{\pgfqpoint{11.074387in}{1.185576in}}{\pgfqpoint{11.079737in}{1.186657in}}%
\pgfpathquadraticcurveto{\pgfqpoint{11.085087in}{1.187756in}}{\pgfqpoint{11.090148in}{1.187756in}}%
\pgfpathquadraticcurveto{\pgfqpoint{11.103837in}{1.187756in}}{\pgfqpoint{11.110592in}{1.180659in}}%
\pgfpathquadraticcurveto{\pgfqpoint{11.117346in}{1.173580in}}{\pgfqpoint{11.117346in}{1.159170in}}%
\pgfpathlineto{\pgfqpoint{11.117346in}{1.159170in}}%
\pgfpathclose%
\pgfpathmoveto{\pgfqpoint{11.175201in}{1.176570in}}%
\pgfpathquadraticcurveto{\pgfqpoint{11.173454in}{1.177579in}}{\pgfqpoint{11.171401in}{1.178047in}}%
\pgfpathquadraticcurveto{\pgfqpoint{11.169347in}{1.178533in}}{\pgfqpoint{11.166880in}{1.178533in}}%
\pgfpathquadraticcurveto{\pgfqpoint{11.158090in}{1.178533in}}{\pgfqpoint{11.153389in}{1.172823in}}%
\pgfpathquadraticcurveto{\pgfqpoint{11.148688in}{1.167114in}}{\pgfqpoint{11.148688in}{1.156414in}}%
\pgfpathlineto{\pgfqpoint{11.148688in}{1.123200in}}%
\pgfpathlineto{\pgfqpoint{11.138276in}{1.123200in}}%
\pgfpathlineto{\pgfqpoint{11.138276in}{1.186243in}}%
\pgfpathlineto{\pgfqpoint{11.148688in}{1.186243in}}%
\pgfpathlineto{\pgfqpoint{11.148688in}{1.176444in}}%
\pgfpathquadraticcurveto{\pgfqpoint{11.151966in}{1.182190in}}{\pgfqpoint{11.157189in}{1.184964in}}%
\pgfpathquadraticcurveto{\pgfqpoint{11.162431in}{1.187756in}}{\pgfqpoint{11.169924in}{1.187756in}}%
\pgfpathquadraticcurveto{\pgfqpoint{11.170987in}{1.187756in}}{\pgfqpoint{11.172283in}{1.187611in}}%
\pgfpathquadraticcurveto{\pgfqpoint{11.173580in}{1.187485in}}{\pgfqpoint{11.175147in}{1.187197in}}%
\pgfpathlineto{\pgfqpoint{11.175201in}{1.176570in}}%
\pgfpathlineto{\pgfqpoint{11.175201in}{1.176570in}}%
\pgfpathclose%
\pgfpathmoveto{\pgfqpoint{11.237451in}{1.157315in}}%
\pgfpathlineto{\pgfqpoint{11.237451in}{1.152254in}}%
\pgfpathlineto{\pgfqpoint{11.189827in}{1.152254in}}%
\pgfpathquadraticcurveto{\pgfqpoint{11.190512in}{1.141554in}}{\pgfqpoint{11.196276in}{1.135953in}}%
\pgfpathquadraticcurveto{\pgfqpoint{11.202039in}{1.130351in}}{\pgfqpoint{11.212342in}{1.130351in}}%
\pgfpathquadraticcurveto{\pgfqpoint{11.218304in}{1.130351in}}{\pgfqpoint{11.223906in}{1.131810in}}%
\pgfpathquadraticcurveto{\pgfqpoint{11.229508in}{1.133269in}}{\pgfqpoint{11.235038in}{1.136205in}}%
\pgfpathlineto{\pgfqpoint{11.235038in}{1.126406in}}%
\pgfpathquadraticcurveto{\pgfqpoint{11.229454in}{1.124047in}}{\pgfqpoint{11.223600in}{1.122804in}}%
\pgfpathquadraticcurveto{\pgfqpoint{11.217746in}{1.121561in}}{\pgfqpoint{11.211730in}{1.121561in}}%
\pgfpathquadraticcurveto{\pgfqpoint{11.196636in}{1.121561in}}{\pgfqpoint{11.187828in}{1.130333in}}%
\pgfpathquadraticcurveto{\pgfqpoint{11.179020in}{1.139123in}}{\pgfqpoint{11.179020in}{1.154109in}}%
\pgfpathquadraticcurveto{\pgfqpoint{11.179020in}{1.169581in}}{\pgfqpoint{11.187378in}{1.178659in}}%
\pgfpathquadraticcurveto{\pgfqpoint{11.195735in}{1.187756in}}{\pgfqpoint{11.209929in}{1.187756in}}%
\pgfpathquadraticcurveto{\pgfqpoint{11.222645in}{1.187756in}}{\pgfqpoint{11.230048in}{1.179560in}}%
\pgfpathquadraticcurveto{\pgfqpoint{11.237451in}{1.171383in}}{\pgfqpoint{11.237451in}{1.157315in}}%
\pgfpathlineto{\pgfqpoint{11.237451in}{1.157315in}}%
\pgfpathclose%
\pgfpathmoveto{\pgfqpoint{11.227094in}{1.160359in}}%
\pgfpathquadraticcurveto{\pgfqpoint{11.226986in}{1.168843in}}{\pgfqpoint{11.222339in}{1.173904in}}%
\pgfpathquadraticcurveto{\pgfqpoint{11.217692in}{1.178984in}}{\pgfqpoint{11.210037in}{1.178984in}}%
\pgfpathquadraticcurveto{\pgfqpoint{11.201373in}{1.178984in}}{\pgfqpoint{11.196168in}{1.174084in}}%
\pgfpathquadraticcurveto{\pgfqpoint{11.190962in}{1.169185in}}{\pgfqpoint{11.190169in}{1.160287in}}%
\pgfpathlineto{\pgfqpoint{11.227094in}{1.160359in}}%
\pgfpathlineto{\pgfqpoint{11.227094in}{1.160359in}}%
\pgfpathclose%
\pgfpathmoveto{\pgfqpoint{11.294640in}{1.184387in}}%
\pgfpathlineto{\pgfqpoint{11.294640in}{1.174589in}}%
\pgfpathquadraticcurveto{\pgfqpoint{11.290263in}{1.176840in}}{\pgfqpoint{11.285526in}{1.177957in}}%
\pgfpathquadraticcurveto{\pgfqpoint{11.280807in}{1.179092in}}{\pgfqpoint{11.275727in}{1.179092in}}%
\pgfpathquadraticcurveto{\pgfqpoint{11.268018in}{1.179092in}}{\pgfqpoint{11.264163in}{1.176732in}}%
\pgfpathquadraticcurveto{\pgfqpoint{11.260309in}{1.174373in}}{\pgfqpoint{11.260309in}{1.169635in}}%
\pgfpathquadraticcurveto{\pgfqpoint{11.260309in}{1.166033in}}{\pgfqpoint{11.263065in}{1.163980in}}%
\pgfpathquadraticcurveto{\pgfqpoint{11.265821in}{1.161926in}}{\pgfqpoint{11.274160in}{1.160071in}}%
\pgfpathlineto{\pgfqpoint{11.277709in}{1.159278in}}%
\pgfpathquadraticcurveto{\pgfqpoint{11.288732in}{1.156919in}}{\pgfqpoint{11.293379in}{1.152614in}}%
\pgfpathquadraticcurveto{\pgfqpoint{11.298026in}{1.148309in}}{\pgfqpoint{11.298026in}{1.140600in}}%
\pgfpathquadraticcurveto{\pgfqpoint{11.298026in}{1.131810in}}{\pgfqpoint{11.291074in}{1.126676in}}%
\pgfpathquadraticcurveto{\pgfqpoint{11.284121in}{1.121561in}}{\pgfqpoint{11.271963in}{1.121561in}}%
\pgfpathquadraticcurveto{\pgfqpoint{11.266901in}{1.121561in}}{\pgfqpoint{11.261408in}{1.122552in}}%
\pgfpathquadraticcurveto{\pgfqpoint{11.255914in}{1.123542in}}{\pgfqpoint{11.249844in}{1.125506in}}%
\pgfpathlineto{\pgfqpoint{11.249844in}{1.136205in}}%
\pgfpathquadraticcurveto{\pgfqpoint{11.255590in}{1.133215in}}{\pgfqpoint{11.261155in}{1.131720in}}%
\pgfpathquadraticcurveto{\pgfqpoint{11.266721in}{1.130243in}}{\pgfqpoint{11.272197in}{1.130243in}}%
\pgfpathquadraticcurveto{\pgfqpoint{11.279510in}{1.130243in}}{\pgfqpoint{11.283436in}{1.132746in}}%
\pgfpathquadraticcurveto{\pgfqpoint{11.287381in}{1.135250in}}{\pgfqpoint{11.287381in}{1.139807in}}%
\pgfpathquadraticcurveto{\pgfqpoint{11.287381in}{1.144022in}}{\pgfqpoint{11.284535in}{1.146274in}}%
\pgfpathquadraticcurveto{\pgfqpoint{11.281707in}{1.148525in}}{\pgfqpoint{11.272071in}{1.150614in}}%
\pgfpathlineto{\pgfqpoint{11.268468in}{1.151461in}}%
\pgfpathquadraticcurveto{\pgfqpoint{11.258850in}{1.153478in}}{\pgfqpoint{11.254563in}{1.157675in}}%
\pgfpathquadraticcurveto{\pgfqpoint{11.250294in}{1.161872in}}{\pgfqpoint{11.250294in}{1.169185in}}%
\pgfpathquadraticcurveto{\pgfqpoint{11.250294in}{1.178083in}}{\pgfqpoint{11.256598in}{1.182910in}}%
\pgfpathquadraticcurveto{\pgfqpoint{11.262903in}{1.187756in}}{\pgfqpoint{11.274502in}{1.187756in}}%
\pgfpathquadraticcurveto{\pgfqpoint{11.280230in}{1.187756in}}{\pgfqpoint{11.285292in}{1.186909in}}%
\pgfpathquadraticcurveto{\pgfqpoint{11.290371in}{1.186080in}}{\pgfqpoint{11.294640in}{1.184387in}}%
\pgfpathlineto{\pgfqpoint{11.294640in}{1.184387in}}%
\pgfpathclose%
\pgfusepath{fill}%
\end{pgfscope}%
\begin{pgfscope}%
\pgfsetbuttcap%
\pgfsetroundjoin%
\definecolor{currentfill}{rgb}{0.309804,0.372549,0.419608}%
\pgfsetfillcolor{currentfill}%
\pgfsetlinewidth{0.000000pt}%
\definecolor{currentstroke}{rgb}{0.309804,0.372549,0.419608}%
\pgfsetstrokecolor{currentstroke}%
\pgfsetdash{}{0pt}%
\pgfpathmoveto{\pgfqpoint{11.387244in}{1.148075in}}%
\pgfpathlineto{\pgfqpoint{11.387244in}{1.186243in}}%
\pgfpathlineto{\pgfqpoint{11.397601in}{1.186243in}}%
\pgfpathlineto{\pgfqpoint{11.397601in}{1.148471in}}%
\pgfpathquadraticcurveto{\pgfqpoint{11.397601in}{1.139519in}}{\pgfqpoint{11.401077in}{1.135034in}}%
\pgfpathquadraticcurveto{\pgfqpoint{11.404572in}{1.130567in}}{\pgfqpoint{11.411560in}{1.130567in}}%
\pgfpathquadraticcurveto{\pgfqpoint{11.419936in}{1.130567in}}{\pgfqpoint{11.424799in}{1.135917in}}%
\pgfpathquadraticcurveto{\pgfqpoint{11.429681in}{1.141266in}}{\pgfqpoint{11.429681in}{1.150506in}}%
\pgfpathlineto{\pgfqpoint{11.429681in}{1.186243in}}%
\pgfpathlineto{\pgfqpoint{11.440038in}{1.186243in}}%
\pgfpathlineto{\pgfqpoint{11.440038in}{1.123200in}}%
\pgfpathlineto{\pgfqpoint{11.429681in}{1.123200in}}%
\pgfpathlineto{\pgfqpoint{11.429681in}{1.132891in}}%
\pgfpathquadraticcurveto{\pgfqpoint{11.425916in}{1.127145in}}{\pgfqpoint{11.420927in}{1.124353in}}%
\pgfpathquadraticcurveto{\pgfqpoint{11.415955in}{1.121561in}}{\pgfqpoint{11.409363in}{1.121561in}}%
\pgfpathquadraticcurveto{\pgfqpoint{11.398502in}{1.121561in}}{\pgfqpoint{11.392864in}{1.128315in}}%
\pgfpathquadraticcurveto{\pgfqpoint{11.387244in}{1.135070in}}{\pgfqpoint{11.387244in}{1.148075in}}%
\pgfpathlineto{\pgfqpoint{11.387244in}{1.148075in}}%
\pgfpathclose%
\pgfpathmoveto{\pgfqpoint{11.413308in}{1.187756in}}%
\pgfpathlineto{\pgfqpoint{11.413308in}{1.187756in}}%
\pgfpathlineto{\pgfqpoint{11.413308in}{1.187756in}}%
\pgfpathclose%
\pgfpathmoveto{\pgfqpoint{11.501549in}{1.184387in}}%
\pgfpathlineto{\pgfqpoint{11.501549in}{1.174589in}}%
\pgfpathquadraticcurveto{\pgfqpoint{11.497172in}{1.176840in}}{\pgfqpoint{11.492435in}{1.177957in}}%
\pgfpathquadraticcurveto{\pgfqpoint{11.487716in}{1.179092in}}{\pgfqpoint{11.482636in}{1.179092in}}%
\pgfpathquadraticcurveto{\pgfqpoint{11.474927in}{1.179092in}}{\pgfqpoint{11.471072in}{1.176732in}}%
\pgfpathquadraticcurveto{\pgfqpoint{11.467218in}{1.174373in}}{\pgfqpoint{11.467218in}{1.169635in}}%
\pgfpathquadraticcurveto{\pgfqpoint{11.467218in}{1.166033in}}{\pgfqpoint{11.469974in}{1.163980in}}%
\pgfpathquadraticcurveto{\pgfqpoint{11.472730in}{1.161926in}}{\pgfqpoint{11.481069in}{1.160071in}}%
\pgfpathlineto{\pgfqpoint{11.484618in}{1.159278in}}%
\pgfpathquadraticcurveto{\pgfqpoint{11.495641in}{1.156919in}}{\pgfqpoint{11.500288in}{1.152614in}}%
\pgfpathquadraticcurveto{\pgfqpoint{11.504935in}{1.148309in}}{\pgfqpoint{11.504935in}{1.140600in}}%
\pgfpathquadraticcurveto{\pgfqpoint{11.504935in}{1.131810in}}{\pgfqpoint{11.497983in}{1.126676in}}%
\pgfpathquadraticcurveto{\pgfqpoint{11.491030in}{1.121561in}}{\pgfqpoint{11.478872in}{1.121561in}}%
\pgfpathquadraticcurveto{\pgfqpoint{11.473810in}{1.121561in}}{\pgfqpoint{11.468317in}{1.122552in}}%
\pgfpathquadraticcurveto{\pgfqpoint{11.462823in}{1.123542in}}{\pgfqpoint{11.456753in}{1.125506in}}%
\pgfpathlineto{\pgfqpoint{11.456753in}{1.136205in}}%
\pgfpathquadraticcurveto{\pgfqpoint{11.462499in}{1.133215in}}{\pgfqpoint{11.468064in}{1.131720in}}%
\pgfpathquadraticcurveto{\pgfqpoint{11.473630in}{1.130243in}}{\pgfqpoint{11.479106in}{1.130243in}}%
\pgfpathquadraticcurveto{\pgfqpoint{11.486419in}{1.130243in}}{\pgfqpoint{11.490345in}{1.132746in}}%
\pgfpathquadraticcurveto{\pgfqpoint{11.494290in}{1.135250in}}{\pgfqpoint{11.494290in}{1.139807in}}%
\pgfpathquadraticcurveto{\pgfqpoint{11.494290in}{1.144022in}}{\pgfqpoint{11.491444in}{1.146274in}}%
\pgfpathquadraticcurveto{\pgfqpoint{11.488616in}{1.148525in}}{\pgfqpoint{11.478980in}{1.150614in}}%
\pgfpathlineto{\pgfqpoint{11.475377in}{1.151461in}}%
\pgfpathquadraticcurveto{\pgfqpoint{11.465759in}{1.153478in}}{\pgfqpoint{11.461472in}{1.157675in}}%
\pgfpathquadraticcurveto{\pgfqpoint{11.457203in}{1.161872in}}{\pgfqpoint{11.457203in}{1.169185in}}%
\pgfpathquadraticcurveto{\pgfqpoint{11.457203in}{1.178083in}}{\pgfqpoint{11.463507in}{1.182910in}}%
\pgfpathquadraticcurveto{\pgfqpoint{11.469812in}{1.187756in}}{\pgfqpoint{11.481411in}{1.187756in}}%
\pgfpathquadraticcurveto{\pgfqpoint{11.487139in}{1.187756in}}{\pgfqpoint{11.492201in}{1.186909in}}%
\pgfpathquadraticcurveto{\pgfqpoint{11.497280in}{1.186080in}}{\pgfqpoint{11.501549in}{1.184387in}}%
\pgfpathlineto{\pgfqpoint{11.501549in}{1.184387in}}%
\pgfpathclose%
\pgfpathmoveto{\pgfqpoint{11.575345in}{1.157315in}}%
\pgfpathlineto{\pgfqpoint{11.575345in}{1.152254in}}%
\pgfpathlineto{\pgfqpoint{11.527721in}{1.152254in}}%
\pgfpathquadraticcurveto{\pgfqpoint{11.528405in}{1.141554in}}{\pgfqpoint{11.534169in}{1.135953in}}%
\pgfpathquadraticcurveto{\pgfqpoint{11.539933in}{1.130351in}}{\pgfqpoint{11.550236in}{1.130351in}}%
\pgfpathquadraticcurveto{\pgfqpoint{11.556198in}{1.130351in}}{\pgfqpoint{11.561800in}{1.131810in}}%
\pgfpathquadraticcurveto{\pgfqpoint{11.567401in}{1.133269in}}{\pgfqpoint{11.572931in}{1.136205in}}%
\pgfpathlineto{\pgfqpoint{11.572931in}{1.126406in}}%
\pgfpathquadraticcurveto{\pgfqpoint{11.567347in}{1.124047in}}{\pgfqpoint{11.561493in}{1.122804in}}%
\pgfpathquadraticcurveto{\pgfqpoint{11.555640in}{1.121561in}}{\pgfqpoint{11.549623in}{1.121561in}}%
\pgfpathquadraticcurveto{\pgfqpoint{11.534529in}{1.121561in}}{\pgfqpoint{11.525721in}{1.130333in}}%
\pgfpathquadraticcurveto{\pgfqpoint{11.516913in}{1.139123in}}{\pgfqpoint{11.516913in}{1.154109in}}%
\pgfpathquadraticcurveto{\pgfqpoint{11.516913in}{1.169581in}}{\pgfqpoint{11.525271in}{1.178659in}}%
\pgfpathquadraticcurveto{\pgfqpoint{11.533629in}{1.187756in}}{\pgfqpoint{11.547822in}{1.187756in}}%
\pgfpathquadraticcurveto{\pgfqpoint{11.560539in}{1.187756in}}{\pgfqpoint{11.567942in}{1.179560in}}%
\pgfpathquadraticcurveto{\pgfqpoint{11.575345in}{1.171383in}}{\pgfqpoint{11.575345in}{1.157315in}}%
\pgfpathlineto{\pgfqpoint{11.575345in}{1.157315in}}%
\pgfpathclose%
\pgfpathmoveto{\pgfqpoint{11.564988in}{1.160359in}}%
\pgfpathquadraticcurveto{\pgfqpoint{11.564880in}{1.168843in}}{\pgfqpoint{11.560233in}{1.173904in}}%
\pgfpathquadraticcurveto{\pgfqpoint{11.555586in}{1.178984in}}{\pgfqpoint{11.547930in}{1.178984in}}%
\pgfpathquadraticcurveto{\pgfqpoint{11.539266in}{1.178984in}}{\pgfqpoint{11.534061in}{1.174084in}}%
\pgfpathquadraticcurveto{\pgfqpoint{11.528855in}{1.169185in}}{\pgfqpoint{11.528063in}{1.160287in}}%
\pgfpathlineto{\pgfqpoint{11.564988in}{1.160359in}}%
\pgfpathlineto{\pgfqpoint{11.564988in}{1.160359in}}%
\pgfpathclose%
\pgfusepath{fill}%
\end{pgfscope}%
\begin{pgfscope}%
\pgfsetbuttcap%
\pgfsetroundjoin%
\definecolor{currentfill}{rgb}{0.309804,0.372549,0.419608}%
\pgfsetfillcolor{currentfill}%
\pgfsetlinewidth{0.000000pt}%
\definecolor{currentstroke}{rgb}{0.309804,0.372549,0.419608}%
\pgfsetstrokecolor{currentstroke}%
\pgfsetdash{}{0pt}%
\pgfpathmoveto{\pgfqpoint{11.693854in}{1.176570in}}%
\pgfpathquadraticcurveto{\pgfqpoint{11.692106in}{1.177579in}}{\pgfqpoint{11.690053in}{1.178047in}}%
\pgfpathquadraticcurveto{\pgfqpoint{11.688000in}{1.178533in}}{\pgfqpoint{11.685532in}{1.178533in}}%
\pgfpathquadraticcurveto{\pgfqpoint{11.676742in}{1.178533in}}{\pgfqpoint{11.672041in}{1.172823in}}%
\pgfpathquadraticcurveto{\pgfqpoint{11.667340in}{1.167114in}}{\pgfqpoint{11.667340in}{1.156414in}}%
\pgfpathlineto{\pgfqpoint{11.667340in}{1.123200in}}%
\pgfpathlineto{\pgfqpoint{11.656929in}{1.123200in}}%
\pgfpathlineto{\pgfqpoint{11.656929in}{1.186243in}}%
\pgfpathlineto{\pgfqpoint{11.667340in}{1.186243in}}%
\pgfpathlineto{\pgfqpoint{11.667340in}{1.176444in}}%
\pgfpathquadraticcurveto{\pgfqpoint{11.670618in}{1.182190in}}{\pgfqpoint{11.675841in}{1.184964in}}%
\pgfpathquadraticcurveto{\pgfqpoint{11.681083in}{1.187756in}}{\pgfqpoint{11.688576in}{1.187756in}}%
\pgfpathquadraticcurveto{\pgfqpoint{11.689639in}{1.187756in}}{\pgfqpoint{11.690936in}{1.187611in}}%
\pgfpathquadraticcurveto{\pgfqpoint{11.692232in}{1.187485in}}{\pgfqpoint{11.693799in}{1.187197in}}%
\pgfpathlineto{\pgfqpoint{11.693854in}{1.176570in}}%
\pgfpathlineto{\pgfqpoint{11.693854in}{1.176570in}}%
\pgfpathclose%
\pgfpathmoveto{\pgfqpoint{11.756104in}{1.157315in}}%
\pgfpathlineto{\pgfqpoint{11.756104in}{1.152254in}}%
\pgfpathlineto{\pgfqpoint{11.708479in}{1.152254in}}%
\pgfpathquadraticcurveto{\pgfqpoint{11.709164in}{1.141554in}}{\pgfqpoint{11.714928in}{1.135953in}}%
\pgfpathquadraticcurveto{\pgfqpoint{11.720692in}{1.130351in}}{\pgfqpoint{11.730995in}{1.130351in}}%
\pgfpathquadraticcurveto{\pgfqpoint{11.736957in}{1.130351in}}{\pgfqpoint{11.742558in}{1.131810in}}%
\pgfpathquadraticcurveto{\pgfqpoint{11.748160in}{1.133269in}}{\pgfqpoint{11.753690in}{1.136205in}}%
\pgfpathlineto{\pgfqpoint{11.753690in}{1.126406in}}%
\pgfpathquadraticcurveto{\pgfqpoint{11.748106in}{1.124047in}}{\pgfqpoint{11.742252in}{1.122804in}}%
\pgfpathquadraticcurveto{\pgfqpoint{11.736398in}{1.121561in}}{\pgfqpoint{11.730382in}{1.121561in}}%
\pgfpathquadraticcurveto{\pgfqpoint{11.715288in}{1.121561in}}{\pgfqpoint{11.706480in}{1.130333in}}%
\pgfpathquadraticcurveto{\pgfqpoint{11.697672in}{1.139123in}}{\pgfqpoint{11.697672in}{1.154109in}}%
\pgfpathquadraticcurveto{\pgfqpoint{11.697672in}{1.169581in}}{\pgfqpoint{11.706030in}{1.178659in}}%
\pgfpathquadraticcurveto{\pgfqpoint{11.714387in}{1.187756in}}{\pgfqpoint{11.728581in}{1.187756in}}%
\pgfpathquadraticcurveto{\pgfqpoint{11.741298in}{1.187756in}}{\pgfqpoint{11.748701in}{1.179560in}}%
\pgfpathquadraticcurveto{\pgfqpoint{11.756104in}{1.171383in}}{\pgfqpoint{11.756104in}{1.157315in}}%
\pgfpathlineto{\pgfqpoint{11.756104in}{1.157315in}}%
\pgfpathclose%
\pgfpathmoveto{\pgfqpoint{11.745747in}{1.160359in}}%
\pgfpathquadraticcurveto{\pgfqpoint{11.745638in}{1.168843in}}{\pgfqpoint{11.740991in}{1.173904in}}%
\pgfpathquadraticcurveto{\pgfqpoint{11.736344in}{1.178984in}}{\pgfqpoint{11.728689in}{1.178984in}}%
\pgfpathquadraticcurveto{\pgfqpoint{11.720025in}{1.178984in}}{\pgfqpoint{11.714820in}{1.174084in}}%
\pgfpathquadraticcurveto{\pgfqpoint{11.709614in}{1.169185in}}{\pgfqpoint{11.708822in}{1.160287in}}%
\pgfpathlineto{\pgfqpoint{11.745747in}{1.160359in}}%
\pgfpathlineto{\pgfqpoint{11.745747in}{1.160359in}}%
\pgfpathclose%
\pgfpathmoveto{\pgfqpoint{11.783356in}{1.204147in}}%
\pgfpathlineto{\pgfqpoint{11.783356in}{1.186243in}}%
\pgfpathlineto{\pgfqpoint{11.804682in}{1.186243in}}%
\pgfpathlineto{\pgfqpoint{11.804682in}{1.178191in}}%
\pgfpathlineto{\pgfqpoint{11.783356in}{1.178191in}}%
\pgfpathlineto{\pgfqpoint{11.783356in}{1.143968in}}%
\pgfpathquadraticcurveto{\pgfqpoint{11.783356in}{1.136259in}}{\pgfqpoint{11.785463in}{1.134061in}}%
\pgfpathquadraticcurveto{\pgfqpoint{11.787571in}{1.131864in}}{\pgfqpoint{11.794055in}{1.131864in}}%
\pgfpathlineto{\pgfqpoint{11.804682in}{1.131864in}}%
\pgfpathlineto{\pgfqpoint{11.804682in}{1.123200in}}%
\pgfpathlineto{\pgfqpoint{11.794055in}{1.123200in}}%
\pgfpathquadraticcurveto{\pgfqpoint{11.782059in}{1.123200in}}{\pgfqpoint{11.777502in}{1.127667in}}%
\pgfpathquadraticcurveto{\pgfqpoint{11.772945in}{1.132152in}}{\pgfqpoint{11.772945in}{1.143968in}}%
\pgfpathlineto{\pgfqpoint{11.772945in}{1.178191in}}%
\pgfpathlineto{\pgfqpoint{11.765344in}{1.178191in}}%
\pgfpathlineto{\pgfqpoint{11.765344in}{1.186243in}}%
\pgfpathlineto{\pgfqpoint{11.772945in}{1.186243in}}%
\pgfpathlineto{\pgfqpoint{11.772945in}{1.204147in}}%
\pgfpathlineto{\pgfqpoint{11.783356in}{1.204147in}}%
\pgfpathlineto{\pgfqpoint{11.783356in}{1.204147in}}%
\pgfpathclose%
\pgfpathmoveto{\pgfqpoint{11.846957in}{1.154883in}}%
\pgfpathquadraticcurveto{\pgfqpoint{11.834402in}{1.154883in}}{\pgfqpoint{11.829557in}{1.152019in}}%
\pgfpathquadraticcurveto{\pgfqpoint{11.824712in}{1.149156in}}{\pgfqpoint{11.824712in}{1.142221in}}%
\pgfpathquadraticcurveto{\pgfqpoint{11.824712in}{1.136709in}}{\pgfqpoint{11.828350in}{1.133467in}}%
\pgfpathquadraticcurveto{\pgfqpoint{11.831989in}{1.130243in}}{\pgfqpoint{11.838221in}{1.130243in}}%
\pgfpathquadraticcurveto{\pgfqpoint{11.846849in}{1.130243in}}{\pgfqpoint{11.852054in}{1.136349in}}%
\pgfpathquadraticcurveto{\pgfqpoint{11.857260in}{1.142455in}}{\pgfqpoint{11.857260in}{1.152578in}}%
\pgfpathlineto{\pgfqpoint{11.857260in}{1.154883in}}%
\pgfpathlineto{\pgfqpoint{11.846957in}{1.154883in}}%
\pgfpathlineto{\pgfqpoint{11.846957in}{1.154883in}}%
\pgfpathclose%
\pgfpathmoveto{\pgfqpoint{11.867617in}{1.159170in}}%
\pgfpathlineto{\pgfqpoint{11.867617in}{1.123200in}}%
\pgfpathlineto{\pgfqpoint{11.857260in}{1.123200in}}%
\pgfpathlineto{\pgfqpoint{11.857260in}{1.132764in}}%
\pgfpathquadraticcurveto{\pgfqpoint{11.853711in}{1.127037in}}{\pgfqpoint{11.848416in}{1.124299in}}%
\pgfpathquadraticcurveto{\pgfqpoint{11.843120in}{1.121561in}}{\pgfqpoint{11.835465in}{1.121561in}}%
\pgfpathquadraticcurveto{\pgfqpoint{11.825793in}{1.121561in}}{\pgfqpoint{11.820065in}{1.127001in}}%
\pgfpathquadraticcurveto{\pgfqpoint{11.814355in}{1.132440in}}{\pgfqpoint{11.814355in}{1.141554in}}%
\pgfpathquadraticcurveto{\pgfqpoint{11.814355in}{1.152182in}}{\pgfqpoint{11.821470in}{1.157585in}}%
\pgfpathquadraticcurveto{\pgfqpoint{11.828602in}{1.162989in}}{\pgfqpoint{11.842724in}{1.162989in}}%
\pgfpathlineto{\pgfqpoint{11.857260in}{1.162989in}}%
\pgfpathlineto{\pgfqpoint{11.857260in}{1.164016in}}%
\pgfpathquadraticcurveto{\pgfqpoint{11.857260in}{1.171166in}}{\pgfqpoint{11.852559in}{1.175075in}}%
\pgfpathquadraticcurveto{\pgfqpoint{11.847857in}{1.178984in}}{\pgfqpoint{11.839356in}{1.178984in}}%
\pgfpathquadraticcurveto{\pgfqpoint{11.833952in}{1.178984in}}{\pgfqpoint{11.828819in}{1.177687in}}%
\pgfpathquadraticcurveto{\pgfqpoint{11.823703in}{1.176390in}}{\pgfqpoint{11.818984in}{1.173796in}}%
\pgfpathlineto{\pgfqpoint{11.818984in}{1.183379in}}%
\pgfpathquadraticcurveto{\pgfqpoint{11.824658in}{1.185576in}}{\pgfqpoint{11.830007in}{1.186657in}}%
\pgfpathquadraticcurveto{\pgfqpoint{11.835357in}{1.187756in}}{\pgfqpoint{11.840418in}{1.187756in}}%
\pgfpathquadraticcurveto{\pgfqpoint{11.854108in}{1.187756in}}{\pgfqpoint{11.860862in}{1.180659in}}%
\pgfpathquadraticcurveto{\pgfqpoint{11.867617in}{1.173580in}}{\pgfqpoint{11.867617in}{1.159170in}}%
\pgfpathlineto{\pgfqpoint{11.867617in}{1.159170in}}%
\pgfpathclose%
\pgfpathmoveto{\pgfqpoint{11.888943in}{1.186243in}}%
\pgfpathlineto{\pgfqpoint{11.899300in}{1.186243in}}%
\pgfpathlineto{\pgfqpoint{11.899300in}{1.123200in}}%
\pgfpathlineto{\pgfqpoint{11.888943in}{1.123200in}}%
\pgfpathlineto{\pgfqpoint{11.888943in}{1.186243in}}%
\pgfpathlineto{\pgfqpoint{11.888943in}{1.186243in}}%
\pgfpathclose%
\pgfpathmoveto{\pgfqpoint{11.888943in}{1.210793in}}%
\pgfpathlineto{\pgfqpoint{11.899300in}{1.210793in}}%
\pgfpathlineto{\pgfqpoint{11.899300in}{1.197662in}}%
\pgfpathlineto{\pgfqpoint{11.888943in}{1.197662in}}%
\pgfpathlineto{\pgfqpoint{11.888943in}{1.210793in}}%
\pgfpathlineto{\pgfqpoint{11.888943in}{1.210793in}}%
\pgfpathclose%
\pgfpathmoveto{\pgfqpoint{11.973384in}{1.161260in}}%
\pgfpathlineto{\pgfqpoint{11.973384in}{1.123200in}}%
\pgfpathlineto{\pgfqpoint{11.963027in}{1.123200in}}%
\pgfpathlineto{\pgfqpoint{11.963027in}{1.160917in}}%
\pgfpathquadraticcurveto{\pgfqpoint{11.963027in}{1.169869in}}{\pgfqpoint{11.959533in}{1.174300in}}%
\pgfpathquadraticcurveto{\pgfqpoint{11.956038in}{1.178749in}}{\pgfqpoint{11.949068in}{1.178749in}}%
\pgfpathquadraticcurveto{\pgfqpoint{11.940674in}{1.178749in}}{\pgfqpoint{11.935829in}{1.173400in}}%
\pgfpathquadraticcurveto{\pgfqpoint{11.930984in}{1.168068in}}{\pgfqpoint{11.930984in}{1.158828in}}%
\pgfpathlineto{\pgfqpoint{11.930984in}{1.123200in}}%
\pgfpathlineto{\pgfqpoint{11.920572in}{1.123200in}}%
\pgfpathlineto{\pgfqpoint{11.920572in}{1.186243in}}%
\pgfpathlineto{\pgfqpoint{11.930984in}{1.186243in}}%
\pgfpathlineto{\pgfqpoint{11.930984in}{1.176444in}}%
\pgfpathquadraticcurveto{\pgfqpoint{11.934712in}{1.182136in}}{\pgfqpoint{11.939737in}{1.184946in}}%
\pgfpathquadraticcurveto{\pgfqpoint{11.944781in}{1.187756in}}{\pgfqpoint{11.951373in}{1.187756in}}%
\pgfpathquadraticcurveto{\pgfqpoint{11.962235in}{1.187756in}}{\pgfqpoint{11.967800in}{1.181037in}}%
\pgfpathquadraticcurveto{\pgfqpoint{11.973384in}{1.174318in}}{\pgfqpoint{11.973384in}{1.161260in}}%
\pgfpathlineto{\pgfqpoint{11.973384in}{1.161260in}}%
\pgfpathclose%
\pgfpathmoveto{\pgfqpoint{12.047954in}{1.157315in}}%
\pgfpathlineto{\pgfqpoint{12.047954in}{1.152254in}}%
\pgfpathlineto{\pgfqpoint{12.000330in}{1.152254in}}%
\pgfpathquadraticcurveto{\pgfqpoint{12.001015in}{1.141554in}}{\pgfqpoint{12.006779in}{1.135953in}}%
\pgfpathquadraticcurveto{\pgfqpoint{12.012543in}{1.130351in}}{\pgfqpoint{12.022845in}{1.130351in}}%
\pgfpathquadraticcurveto{\pgfqpoint{12.028808in}{1.130351in}}{\pgfqpoint{12.034409in}{1.131810in}}%
\pgfpathquadraticcurveto{\pgfqpoint{12.040011in}{1.133269in}}{\pgfqpoint{12.045541in}{1.136205in}}%
\pgfpathlineto{\pgfqpoint{12.045541in}{1.126406in}}%
\pgfpathquadraticcurveto{\pgfqpoint{12.039957in}{1.124047in}}{\pgfqpoint{12.034103in}{1.122804in}}%
\pgfpathquadraticcurveto{\pgfqpoint{12.028249in}{1.121561in}}{\pgfqpoint{12.022233in}{1.121561in}}%
\pgfpathquadraticcurveto{\pgfqpoint{12.007139in}{1.121561in}}{\pgfqpoint{11.998331in}{1.130333in}}%
\pgfpathquadraticcurveto{\pgfqpoint{11.989523in}{1.139123in}}{\pgfqpoint{11.989523in}{1.154109in}}%
\pgfpathquadraticcurveto{\pgfqpoint{11.989523in}{1.169581in}}{\pgfqpoint{11.997881in}{1.178659in}}%
\pgfpathquadraticcurveto{\pgfqpoint{12.006238in}{1.187756in}}{\pgfqpoint{12.020432in}{1.187756in}}%
\pgfpathquadraticcurveto{\pgfqpoint{12.033148in}{1.187756in}}{\pgfqpoint{12.040551in}{1.179560in}}%
\pgfpathquadraticcurveto{\pgfqpoint{12.047954in}{1.171383in}}{\pgfqpoint{12.047954in}{1.157315in}}%
\pgfpathlineto{\pgfqpoint{12.047954in}{1.157315in}}%
\pgfpathclose%
\pgfpathmoveto{\pgfqpoint{12.037597in}{1.160359in}}%
\pgfpathquadraticcurveto{\pgfqpoint{12.037489in}{1.168843in}}{\pgfqpoint{12.032842in}{1.173904in}}%
\pgfpathquadraticcurveto{\pgfqpoint{12.028195in}{1.178984in}}{\pgfqpoint{12.020540in}{1.178984in}}%
\pgfpathquadraticcurveto{\pgfqpoint{12.011876in}{1.178984in}}{\pgfqpoint{12.006671in}{1.174084in}}%
\pgfpathquadraticcurveto{\pgfqpoint{12.001465in}{1.169185in}}{\pgfqpoint{12.000673in}{1.160287in}}%
\pgfpathlineto{\pgfqpoint{12.037597in}{1.160359in}}%
\pgfpathlineto{\pgfqpoint{12.037597in}{1.160359in}}%
\pgfpathclose%
\pgfpathmoveto{\pgfqpoint{12.106440in}{1.176678in}}%
\pgfpathlineto{\pgfqpoint{12.106440in}{1.210793in}}%
\pgfpathlineto{\pgfqpoint{12.116797in}{1.210793in}}%
\pgfpathlineto{\pgfqpoint{12.116797in}{1.123200in}}%
\pgfpathlineto{\pgfqpoint{12.106440in}{1.123200in}}%
\pgfpathlineto{\pgfqpoint{12.106440in}{1.132656in}}%
\pgfpathquadraticcurveto{\pgfqpoint{12.103180in}{1.127037in}}{\pgfqpoint{12.098190in}{1.124299in}}%
\pgfpathquadraticcurveto{\pgfqpoint{12.093219in}{1.121561in}}{\pgfqpoint{12.086230in}{1.121561in}}%
\pgfpathquadraticcurveto{\pgfqpoint{12.074811in}{1.121561in}}{\pgfqpoint{12.067624in}{1.130675in}}%
\pgfpathquadraticcurveto{\pgfqpoint{12.060455in}{1.139807in}}{\pgfqpoint{12.060455in}{1.154667in}}%
\pgfpathquadraticcurveto{\pgfqpoint{12.060455in}{1.169527in}}{\pgfqpoint{12.067624in}{1.178641in}}%
\pgfpathquadraticcurveto{\pgfqpoint{12.074811in}{1.187756in}}{\pgfqpoint{12.086230in}{1.187756in}}%
\pgfpathquadraticcurveto{\pgfqpoint{12.093219in}{1.187756in}}{\pgfqpoint{12.098190in}{1.185018in}}%
\pgfpathquadraticcurveto{\pgfqpoint{12.103180in}{1.182298in}}{\pgfqpoint{12.106440in}{1.176678in}}%
\pgfpathlineto{\pgfqpoint{12.106440in}{1.176678in}}%
\pgfpathclose%
\pgfpathmoveto{\pgfqpoint{12.071154in}{1.154667in}}%
\pgfpathquadraticcurveto{\pgfqpoint{12.071154in}{1.143248in}}{\pgfqpoint{12.075855in}{1.136745in}}%
\pgfpathquadraticcurveto{\pgfqpoint{12.080556in}{1.130243in}}{\pgfqpoint{12.088770in}{1.130243in}}%
\pgfpathquadraticcurveto{\pgfqpoint{12.096984in}{1.130243in}}{\pgfqpoint{12.101703in}{1.136745in}}%
\pgfpathquadraticcurveto{\pgfqpoint{12.106440in}{1.143248in}}{\pgfqpoint{12.106440in}{1.154667in}}%
\pgfpathquadraticcurveto{\pgfqpoint{12.106440in}{1.166087in}}{\pgfqpoint{12.101703in}{1.172589in}}%
\pgfpathquadraticcurveto{\pgfqpoint{12.096984in}{1.179092in}}{\pgfqpoint{12.088770in}{1.179092in}}%
\pgfpathquadraticcurveto{\pgfqpoint{12.080556in}{1.179092in}}{\pgfqpoint{12.075855in}{1.172589in}}%
\pgfpathquadraticcurveto{\pgfqpoint{12.071154in}{1.166087in}}{\pgfqpoint{12.071154in}{1.154667in}}%
\pgfpathlineto{\pgfqpoint{12.071154in}{1.154667in}}%
\pgfpathclose%
\pgfusepath{fill}%
\end{pgfscope}%
\begin{pgfscope}%
\pgfsetbuttcap%
\pgfsetroundjoin%
\definecolor{currentfill}{rgb}{0.309804,0.372549,0.419608}%
\pgfsetfillcolor{currentfill}%
\pgfsetlinewidth{0.000000pt}%
\definecolor{currentstroke}{rgb}{0.309804,0.372549,0.419608}%
\pgfsetstrokecolor{currentstroke}%
\pgfsetdash{}{0pt}%
\pgfpathmoveto{\pgfqpoint{12.259216in}{1.183829in}}%
\pgfpathlineto{\pgfqpoint{12.259216in}{1.174138in}}%
\pgfpathquadraticcurveto{\pgfqpoint{12.254821in}{1.176570in}}{\pgfqpoint{12.250408in}{1.177777in}}%
\pgfpathquadraticcurveto{\pgfqpoint{12.245995in}{1.178984in}}{\pgfqpoint{12.241492in}{1.178984in}}%
\pgfpathquadraticcurveto{\pgfqpoint{12.231405in}{1.178984in}}{\pgfqpoint{12.225821in}{1.172589in}}%
\pgfpathquadraticcurveto{\pgfqpoint{12.220255in}{1.166213in}}{\pgfqpoint{12.220255in}{1.154667in}}%
\pgfpathquadraticcurveto{\pgfqpoint{12.220255in}{1.143121in}}{\pgfqpoint{12.225821in}{1.136727in}}%
\pgfpathquadraticcurveto{\pgfqpoint{12.231405in}{1.130351in}}{\pgfqpoint{12.241492in}{1.130351in}}%
\pgfpathquadraticcurveto{\pgfqpoint{12.245995in}{1.130351in}}{\pgfqpoint{12.250408in}{1.131558in}}%
\pgfpathquadraticcurveto{\pgfqpoint{12.254821in}{1.132764in}}{\pgfqpoint{12.259216in}{1.135196in}}%
\pgfpathlineto{\pgfqpoint{12.259216in}{1.125614in}}%
\pgfpathquadraticcurveto{\pgfqpoint{12.254875in}{1.123596in}}{\pgfqpoint{12.250228in}{1.122588in}}%
\pgfpathquadraticcurveto{\pgfqpoint{12.245599in}{1.121561in}}{\pgfqpoint{12.240357in}{1.121561in}}%
\pgfpathquadraticcurveto{\pgfqpoint{12.226109in}{1.121561in}}{\pgfqpoint{12.217716in}{1.130513in}}%
\pgfpathquadraticcurveto{\pgfqpoint{12.209340in}{1.139465in}}{\pgfqpoint{12.209340in}{1.154667in}}%
\pgfpathquadraticcurveto{\pgfqpoint{12.209340in}{1.170086in}}{\pgfqpoint{12.217806in}{1.178912in}}%
\pgfpathquadraticcurveto{\pgfqpoint{12.226289in}{1.187756in}}{\pgfqpoint{12.241041in}{1.187756in}}%
\pgfpathquadraticcurveto{\pgfqpoint{12.245815in}{1.187756in}}{\pgfqpoint{12.250372in}{1.186765in}}%
\pgfpathquadraticcurveto{\pgfqpoint{12.254929in}{1.185792in}}{\pgfqpoint{12.259216in}{1.183829in}}%
\pgfpathlineto{\pgfqpoint{12.259216in}{1.183829in}}%
\pgfpathclose%
\pgfpathmoveto{\pgfqpoint{12.301652in}{1.178984in}}%
\pgfpathquadraticcurveto{\pgfqpoint{12.293331in}{1.178984in}}{\pgfqpoint{12.288485in}{1.172481in}}%
\pgfpathquadraticcurveto{\pgfqpoint{12.283640in}{1.165979in}}{\pgfqpoint{12.283640in}{1.154667in}}%
\pgfpathquadraticcurveto{\pgfqpoint{12.283640in}{1.143356in}}{\pgfqpoint{12.288449in}{1.136853in}}%
\pgfpathquadraticcurveto{\pgfqpoint{12.293277in}{1.130351in}}{\pgfqpoint{12.301652in}{1.130351in}}%
\pgfpathquadraticcurveto{\pgfqpoint{12.309938in}{1.130351in}}{\pgfqpoint{12.314765in}{1.136871in}}%
\pgfpathquadraticcurveto{\pgfqpoint{12.319610in}{1.143410in}}{\pgfqpoint{12.319610in}{1.154667in}}%
\pgfpathquadraticcurveto{\pgfqpoint{12.319610in}{1.165871in}}{\pgfqpoint{12.314765in}{1.172427in}}%
\pgfpathquadraticcurveto{\pgfqpoint{12.309938in}{1.178984in}}{\pgfqpoint{12.301652in}{1.178984in}}%
\pgfpathlineto{\pgfqpoint{12.301652in}{1.178984in}}%
\pgfpathclose%
\pgfpathmoveto{\pgfqpoint{12.301652in}{1.187756in}}%
\pgfpathquadraticcurveto{\pgfqpoint{12.315161in}{1.187756in}}{\pgfqpoint{12.322871in}{1.178966in}}%
\pgfpathquadraticcurveto{\pgfqpoint{12.330598in}{1.170194in}}{\pgfqpoint{12.330598in}{1.154667in}}%
\pgfpathquadraticcurveto{\pgfqpoint{12.330598in}{1.139195in}}{\pgfqpoint{12.322871in}{1.130369in}}%
\pgfpathquadraticcurveto{\pgfqpoint{12.315161in}{1.121561in}}{\pgfqpoint{12.301652in}{1.121561in}}%
\pgfpathquadraticcurveto{\pgfqpoint{12.288089in}{1.121561in}}{\pgfqpoint{12.280398in}{1.130369in}}%
\pgfpathquadraticcurveto{\pgfqpoint{12.272725in}{1.139195in}}{\pgfqpoint{12.272725in}{1.154667in}}%
\pgfpathquadraticcurveto{\pgfqpoint{12.272725in}{1.170194in}}{\pgfqpoint{12.280398in}{1.178966in}}%
\pgfpathquadraticcurveto{\pgfqpoint{12.288089in}{1.187756in}}{\pgfqpoint{12.301652in}{1.187756in}}%
\pgfpathlineto{\pgfqpoint{12.301652in}{1.187756in}}%
\pgfpathclose%
\pgfpathmoveto{\pgfqpoint{12.396847in}{1.174138in}}%
\pgfpathquadraticcurveto{\pgfqpoint{12.400737in}{1.181127in}}{\pgfqpoint{12.406141in}{1.184441in}}%
\pgfpathquadraticcurveto{\pgfqpoint{12.411544in}{1.187756in}}{\pgfqpoint{12.418857in}{1.187756in}}%
\pgfpathquadraticcurveto{\pgfqpoint{12.428710in}{1.187756in}}{\pgfqpoint{12.434060in}{1.180857in}}%
\pgfpathquadraticcurveto{\pgfqpoint{12.439409in}{1.173976in}}{\pgfqpoint{12.439409in}{1.161260in}}%
\pgfpathlineto{\pgfqpoint{12.439409in}{1.123200in}}%
\pgfpathlineto{\pgfqpoint{12.428998in}{1.123200in}}%
\pgfpathlineto{\pgfqpoint{12.428998in}{1.160917in}}%
\pgfpathquadraticcurveto{\pgfqpoint{12.428998in}{1.169978in}}{\pgfqpoint{12.425774in}{1.174355in}}%
\pgfpathquadraticcurveto{\pgfqpoint{12.422568in}{1.178749in}}{\pgfqpoint{12.415993in}{1.178749in}}%
\pgfpathquadraticcurveto{\pgfqpoint{12.407942in}{1.178749in}}{\pgfqpoint{12.403259in}{1.173400in}}%
\pgfpathquadraticcurveto{\pgfqpoint{12.398594in}{1.168068in}}{\pgfqpoint{12.398594in}{1.158828in}}%
\pgfpathlineto{\pgfqpoint{12.398594in}{1.123200in}}%
\pgfpathlineto{\pgfqpoint{12.388183in}{1.123200in}}%
\pgfpathlineto{\pgfqpoint{12.388183in}{1.160917in}}%
\pgfpathquadraticcurveto{\pgfqpoint{12.388183in}{1.170032in}}{\pgfqpoint{12.384977in}{1.174391in}}%
\pgfpathquadraticcurveto{\pgfqpoint{12.381770in}{1.178749in}}{\pgfqpoint{12.375070in}{1.178749in}}%
\pgfpathquadraticcurveto{\pgfqpoint{12.367127in}{1.178749in}}{\pgfqpoint{12.362443in}{1.173382in}}%
\pgfpathquadraticcurveto{\pgfqpoint{12.357778in}{1.168014in}}{\pgfqpoint{12.357778in}{1.158828in}}%
\pgfpathlineto{\pgfqpoint{12.357778in}{1.123200in}}%
\pgfpathlineto{\pgfqpoint{12.347367in}{1.123200in}}%
\pgfpathlineto{\pgfqpoint{12.347367in}{1.186243in}}%
\pgfpathlineto{\pgfqpoint{12.357778in}{1.186243in}}%
\pgfpathlineto{\pgfqpoint{12.357778in}{1.176444in}}%
\pgfpathquadraticcurveto{\pgfqpoint{12.361327in}{1.182244in}}{\pgfqpoint{12.366280in}{1.185000in}}%
\pgfpathquadraticcurveto{\pgfqpoint{12.371233in}{1.187756in}}{\pgfqpoint{12.378042in}{1.187756in}}%
\pgfpathquadraticcurveto{\pgfqpoint{12.384923in}{1.187756in}}{\pgfqpoint{12.389732in}{1.184261in}}%
\pgfpathquadraticcurveto{\pgfqpoint{12.394541in}{1.180785in}}{\pgfqpoint{12.396847in}{1.174138in}}%
\pgfpathlineto{\pgfqpoint{12.396847in}{1.174138in}}%
\pgfpathclose%
\pgfpathmoveto{\pgfqpoint{12.470066in}{1.132656in}}%
\pgfpathlineto{\pgfqpoint{12.470066in}{1.099226in}}%
\pgfpathlineto{\pgfqpoint{12.459655in}{1.099226in}}%
\pgfpathlineto{\pgfqpoint{12.459655in}{1.186243in}}%
\pgfpathlineto{\pgfqpoint{12.470066in}{1.186243in}}%
\pgfpathlineto{\pgfqpoint{12.470066in}{1.176678in}}%
\pgfpathquadraticcurveto{\pgfqpoint{12.473344in}{1.182298in}}{\pgfqpoint{12.478316in}{1.185018in}}%
\pgfpathquadraticcurveto{\pgfqpoint{12.483305in}{1.187756in}}{\pgfqpoint{12.490222in}{1.187756in}}%
\pgfpathquadraticcurveto{\pgfqpoint{12.501713in}{1.187756in}}{\pgfqpoint{12.508882in}{1.178641in}}%
\pgfpathquadraticcurveto{\pgfqpoint{12.516069in}{1.169527in}}{\pgfqpoint{12.516069in}{1.154667in}}%
\pgfpathquadraticcurveto{\pgfqpoint{12.516069in}{1.139807in}}{\pgfqpoint{12.508882in}{1.130675in}}%
\pgfpathquadraticcurveto{\pgfqpoint{12.501713in}{1.121561in}}{\pgfqpoint{12.490222in}{1.121561in}}%
\pgfpathquadraticcurveto{\pgfqpoint{12.483305in}{1.121561in}}{\pgfqpoint{12.478316in}{1.124299in}}%
\pgfpathquadraticcurveto{\pgfqpoint{12.473344in}{1.127037in}}{\pgfqpoint{12.470066in}{1.132656in}}%
\pgfpathlineto{\pgfqpoint{12.470066in}{1.132656in}}%
\pgfpathclose%
\pgfpathmoveto{\pgfqpoint{12.505316in}{1.154667in}}%
\pgfpathquadraticcurveto{\pgfqpoint{12.505316in}{1.166087in}}{\pgfqpoint{12.500615in}{1.172589in}}%
\pgfpathquadraticcurveto{\pgfqpoint{12.495913in}{1.179092in}}{\pgfqpoint{12.487700in}{1.179092in}}%
\pgfpathquadraticcurveto{\pgfqpoint{12.479468in}{1.179092in}}{\pgfqpoint{12.474767in}{1.172589in}}%
\pgfpathquadraticcurveto{\pgfqpoint{12.470066in}{1.166087in}}{\pgfqpoint{12.470066in}{1.154667in}}%
\pgfpathquadraticcurveto{\pgfqpoint{12.470066in}{1.143248in}}{\pgfqpoint{12.474767in}{1.136745in}}%
\pgfpathquadraticcurveto{\pgfqpoint{12.479468in}{1.130243in}}{\pgfqpoint{12.487700in}{1.130243in}}%
\pgfpathquadraticcurveto{\pgfqpoint{12.495913in}{1.130243in}}{\pgfqpoint{12.500615in}{1.136745in}}%
\pgfpathquadraticcurveto{\pgfqpoint{12.505316in}{1.143248in}}{\pgfqpoint{12.505316in}{1.154667in}}%
\pgfpathlineto{\pgfqpoint{12.505316in}{1.154667in}}%
\pgfpathclose%
\pgfpathmoveto{\pgfqpoint{12.561892in}{1.154883in}}%
\pgfpathquadraticcurveto{\pgfqpoint{12.549337in}{1.154883in}}{\pgfqpoint{12.544492in}{1.152019in}}%
\pgfpathquadraticcurveto{\pgfqpoint{12.539647in}{1.149156in}}{\pgfqpoint{12.539647in}{1.142221in}}%
\pgfpathquadraticcurveto{\pgfqpoint{12.539647in}{1.136709in}}{\pgfqpoint{12.543285in}{1.133467in}}%
\pgfpathquadraticcurveto{\pgfqpoint{12.546924in}{1.130243in}}{\pgfqpoint{12.553156in}{1.130243in}}%
\pgfpathquadraticcurveto{\pgfqpoint{12.561784in}{1.130243in}}{\pgfqpoint{12.566989in}{1.136349in}}%
\pgfpathquadraticcurveto{\pgfqpoint{12.572195in}{1.142455in}}{\pgfqpoint{12.572195in}{1.152578in}}%
\pgfpathlineto{\pgfqpoint{12.572195in}{1.154883in}}%
\pgfpathlineto{\pgfqpoint{12.561892in}{1.154883in}}%
\pgfpathlineto{\pgfqpoint{12.561892in}{1.154883in}}%
\pgfpathclose%
\pgfpathmoveto{\pgfqpoint{12.582552in}{1.159170in}}%
\pgfpathlineto{\pgfqpoint{12.582552in}{1.123200in}}%
\pgfpathlineto{\pgfqpoint{12.572195in}{1.123200in}}%
\pgfpathlineto{\pgfqpoint{12.572195in}{1.132764in}}%
\pgfpathquadraticcurveto{\pgfqpoint{12.568646in}{1.127037in}}{\pgfqpoint{12.563351in}{1.124299in}}%
\pgfpathquadraticcurveto{\pgfqpoint{12.558055in}{1.121561in}}{\pgfqpoint{12.550400in}{1.121561in}}%
\pgfpathquadraticcurveto{\pgfqpoint{12.540728in}{1.121561in}}{\pgfqpoint{12.535000in}{1.127001in}}%
\pgfpathquadraticcurveto{\pgfqpoint{12.529290in}{1.132440in}}{\pgfqpoint{12.529290in}{1.141554in}}%
\pgfpathquadraticcurveto{\pgfqpoint{12.529290in}{1.152182in}}{\pgfqpoint{12.536405in}{1.157585in}}%
\pgfpathquadraticcurveto{\pgfqpoint{12.543538in}{1.162989in}}{\pgfqpoint{12.557659in}{1.162989in}}%
\pgfpathlineto{\pgfqpoint{12.572195in}{1.162989in}}%
\pgfpathlineto{\pgfqpoint{12.572195in}{1.164016in}}%
\pgfpathquadraticcurveto{\pgfqpoint{12.572195in}{1.171166in}}{\pgfqpoint{12.567494in}{1.175075in}}%
\pgfpathquadraticcurveto{\pgfqpoint{12.562793in}{1.178984in}}{\pgfqpoint{12.554291in}{1.178984in}}%
\pgfpathquadraticcurveto{\pgfqpoint{12.548887in}{1.178984in}}{\pgfqpoint{12.543754in}{1.177687in}}%
\pgfpathquadraticcurveto{\pgfqpoint{12.538638in}{1.176390in}}{\pgfqpoint{12.533919in}{1.173796in}}%
\pgfpathlineto{\pgfqpoint{12.533919in}{1.183379in}}%
\pgfpathquadraticcurveto{\pgfqpoint{12.539593in}{1.185576in}}{\pgfqpoint{12.544942in}{1.186657in}}%
\pgfpathquadraticcurveto{\pgfqpoint{12.550292in}{1.187756in}}{\pgfqpoint{12.555354in}{1.187756in}}%
\pgfpathquadraticcurveto{\pgfqpoint{12.569043in}{1.187756in}}{\pgfqpoint{12.575797in}{1.180659in}}%
\pgfpathquadraticcurveto{\pgfqpoint{12.582552in}{1.173580in}}{\pgfqpoint{12.582552in}{1.159170in}}%
\pgfpathlineto{\pgfqpoint{12.582552in}{1.159170in}}%
\pgfpathclose%
\pgfpathmoveto{\pgfqpoint{12.640407in}{1.176570in}}%
\pgfpathquadraticcurveto{\pgfqpoint{12.638660in}{1.177579in}}{\pgfqpoint{12.636606in}{1.178047in}}%
\pgfpathquadraticcurveto{\pgfqpoint{12.634553in}{1.178533in}}{\pgfqpoint{12.632085in}{1.178533in}}%
\pgfpathquadraticcurveto{\pgfqpoint{12.623295in}{1.178533in}}{\pgfqpoint{12.618594in}{1.172823in}}%
\pgfpathquadraticcurveto{\pgfqpoint{12.613893in}{1.167114in}}{\pgfqpoint{12.613893in}{1.156414in}}%
\pgfpathlineto{\pgfqpoint{12.613893in}{1.123200in}}%
\pgfpathlineto{\pgfqpoint{12.603482in}{1.123200in}}%
\pgfpathlineto{\pgfqpoint{12.603482in}{1.186243in}}%
\pgfpathlineto{\pgfqpoint{12.613893in}{1.186243in}}%
\pgfpathlineto{\pgfqpoint{12.613893in}{1.176444in}}%
\pgfpathquadraticcurveto{\pgfqpoint{12.617171in}{1.182190in}}{\pgfqpoint{12.622395in}{1.184964in}}%
\pgfpathquadraticcurveto{\pgfqpoint{12.627636in}{1.187756in}}{\pgfqpoint{12.635129in}{1.187756in}}%
\pgfpathquadraticcurveto{\pgfqpoint{12.636192in}{1.187756in}}{\pgfqpoint{12.637489in}{1.187611in}}%
\pgfpathquadraticcurveto{\pgfqpoint{12.638786in}{1.187485in}}{\pgfqpoint{12.640353in}{1.187197in}}%
\pgfpathlineto{\pgfqpoint{12.640407in}{1.176570in}}%
\pgfpathlineto{\pgfqpoint{12.640407in}{1.176570in}}%
\pgfpathclose%
\pgfpathmoveto{\pgfqpoint{12.651268in}{1.186243in}}%
\pgfpathlineto{\pgfqpoint{12.661625in}{1.186243in}}%
\pgfpathlineto{\pgfqpoint{12.661625in}{1.123200in}}%
\pgfpathlineto{\pgfqpoint{12.651268in}{1.123200in}}%
\pgfpathlineto{\pgfqpoint{12.651268in}{1.186243in}}%
\pgfpathlineto{\pgfqpoint{12.651268in}{1.186243in}}%
\pgfpathclose%
\pgfpathmoveto{\pgfqpoint{12.651268in}{1.210793in}}%
\pgfpathlineto{\pgfqpoint{12.661625in}{1.210793in}}%
\pgfpathlineto{\pgfqpoint{12.661625in}{1.197662in}}%
\pgfpathlineto{\pgfqpoint{12.651268in}{1.197662in}}%
\pgfpathlineto{\pgfqpoint{12.651268in}{1.210793in}}%
\pgfpathlineto{\pgfqpoint{12.651268in}{1.210793in}}%
\pgfpathclose%
\pgfpathmoveto{\pgfqpoint{12.723479in}{1.184387in}}%
\pgfpathlineto{\pgfqpoint{12.723479in}{1.174589in}}%
\pgfpathquadraticcurveto{\pgfqpoint{12.719102in}{1.176840in}}{\pgfqpoint{12.714365in}{1.177957in}}%
\pgfpathquadraticcurveto{\pgfqpoint{12.709646in}{1.179092in}}{\pgfqpoint{12.704566in}{1.179092in}}%
\pgfpathquadraticcurveto{\pgfqpoint{12.696857in}{1.179092in}}{\pgfqpoint{12.693002in}{1.176732in}}%
\pgfpathquadraticcurveto{\pgfqpoint{12.689148in}{1.174373in}}{\pgfqpoint{12.689148in}{1.169635in}}%
\pgfpathquadraticcurveto{\pgfqpoint{12.689148in}{1.166033in}}{\pgfqpoint{12.691904in}{1.163980in}}%
\pgfpathquadraticcurveto{\pgfqpoint{12.694660in}{1.161926in}}{\pgfqpoint{12.702999in}{1.160071in}}%
\pgfpathlineto{\pgfqpoint{12.706548in}{1.159278in}}%
\pgfpathquadraticcurveto{\pgfqpoint{12.717571in}{1.156919in}}{\pgfqpoint{12.722218in}{1.152614in}}%
\pgfpathquadraticcurveto{\pgfqpoint{12.726865in}{1.148309in}}{\pgfqpoint{12.726865in}{1.140600in}}%
\pgfpathquadraticcurveto{\pgfqpoint{12.726865in}{1.131810in}}{\pgfqpoint{12.719913in}{1.126676in}}%
\pgfpathquadraticcurveto{\pgfqpoint{12.712960in}{1.121561in}}{\pgfqpoint{12.700802in}{1.121561in}}%
\pgfpathquadraticcurveto{\pgfqpoint{12.695740in}{1.121561in}}{\pgfqpoint{12.690247in}{1.122552in}}%
\pgfpathquadraticcurveto{\pgfqpoint{12.684753in}{1.123542in}}{\pgfqpoint{12.678683in}{1.125506in}}%
\pgfpathlineto{\pgfqpoint{12.678683in}{1.136205in}}%
\pgfpathquadraticcurveto{\pgfqpoint{12.684429in}{1.133215in}}{\pgfqpoint{12.689994in}{1.131720in}}%
\pgfpathquadraticcurveto{\pgfqpoint{12.695560in}{1.130243in}}{\pgfqpoint{12.701036in}{1.130243in}}%
\pgfpathquadraticcurveto{\pgfqpoint{12.708349in}{1.130243in}}{\pgfqpoint{12.712275in}{1.132746in}}%
\pgfpathquadraticcurveto{\pgfqpoint{12.716220in}{1.135250in}}{\pgfqpoint{12.716220in}{1.139807in}}%
\pgfpathquadraticcurveto{\pgfqpoint{12.716220in}{1.144022in}}{\pgfqpoint{12.713374in}{1.146274in}}%
\pgfpathquadraticcurveto{\pgfqpoint{12.710546in}{1.148525in}}{\pgfqpoint{12.700910in}{1.150614in}}%
\pgfpathlineto{\pgfqpoint{12.697307in}{1.151461in}}%
\pgfpathquadraticcurveto{\pgfqpoint{12.687689in}{1.153478in}}{\pgfqpoint{12.683402in}{1.157675in}}%
\pgfpathquadraticcurveto{\pgfqpoint{12.679133in}{1.161872in}}{\pgfqpoint{12.679133in}{1.169185in}}%
\pgfpathquadraticcurveto{\pgfqpoint{12.679133in}{1.178083in}}{\pgfqpoint{12.685437in}{1.182910in}}%
\pgfpathquadraticcurveto{\pgfqpoint{12.691742in}{1.187756in}}{\pgfqpoint{12.703341in}{1.187756in}}%
\pgfpathquadraticcurveto{\pgfqpoint{12.709069in}{1.187756in}}{\pgfqpoint{12.714131in}{1.186909in}}%
\pgfpathquadraticcurveto{\pgfqpoint{12.719210in}{1.186080in}}{\pgfqpoint{12.723479in}{1.184387in}}%
\pgfpathlineto{\pgfqpoint{12.723479in}{1.184387in}}%
\pgfpathclose%
\pgfpathmoveto{\pgfqpoint{12.767771in}{1.178984in}}%
\pgfpathquadraticcurveto{\pgfqpoint{12.759449in}{1.178984in}}{\pgfqpoint{12.754604in}{1.172481in}}%
\pgfpathquadraticcurveto{\pgfqpoint{12.749759in}{1.165979in}}{\pgfqpoint{12.749759in}{1.154667in}}%
\pgfpathquadraticcurveto{\pgfqpoint{12.749759in}{1.143356in}}{\pgfqpoint{12.754568in}{1.136853in}}%
\pgfpathquadraticcurveto{\pgfqpoint{12.759395in}{1.130351in}}{\pgfqpoint{12.767771in}{1.130351in}}%
\pgfpathquadraticcurveto{\pgfqpoint{12.776056in}{1.130351in}}{\pgfqpoint{12.780884in}{1.136871in}}%
\pgfpathquadraticcurveto{\pgfqpoint{12.785729in}{1.143410in}}{\pgfqpoint{12.785729in}{1.154667in}}%
\pgfpathquadraticcurveto{\pgfqpoint{12.785729in}{1.165871in}}{\pgfqpoint{12.780884in}{1.172427in}}%
\pgfpathquadraticcurveto{\pgfqpoint{12.776056in}{1.178984in}}{\pgfqpoint{12.767771in}{1.178984in}}%
\pgfpathlineto{\pgfqpoint{12.767771in}{1.178984in}}%
\pgfpathclose%
\pgfpathmoveto{\pgfqpoint{12.767771in}{1.187756in}}%
\pgfpathquadraticcurveto{\pgfqpoint{12.781280in}{1.187756in}}{\pgfqpoint{12.788989in}{1.178966in}}%
\pgfpathquadraticcurveto{\pgfqpoint{12.796716in}{1.170194in}}{\pgfqpoint{12.796716in}{1.154667in}}%
\pgfpathquadraticcurveto{\pgfqpoint{12.796716in}{1.139195in}}{\pgfqpoint{12.788989in}{1.130369in}}%
\pgfpathquadraticcurveto{\pgfqpoint{12.781280in}{1.121561in}}{\pgfqpoint{12.767771in}{1.121561in}}%
\pgfpathquadraticcurveto{\pgfqpoint{12.754208in}{1.121561in}}{\pgfqpoint{12.746516in}{1.130369in}}%
\pgfpathquadraticcurveto{\pgfqpoint{12.738843in}{1.139195in}}{\pgfqpoint{12.738843in}{1.154667in}}%
\pgfpathquadraticcurveto{\pgfqpoint{12.738843in}{1.170194in}}{\pgfqpoint{12.746516in}{1.178966in}}%
\pgfpathquadraticcurveto{\pgfqpoint{12.754208in}{1.187756in}}{\pgfqpoint{12.767771in}{1.187756in}}%
\pgfpathlineto{\pgfqpoint{12.767771in}{1.187756in}}%
\pgfpathclose%
\pgfpathmoveto{\pgfqpoint{12.866297in}{1.161260in}}%
\pgfpathlineto{\pgfqpoint{12.866297in}{1.123200in}}%
\pgfpathlineto{\pgfqpoint{12.855940in}{1.123200in}}%
\pgfpathlineto{\pgfqpoint{12.855940in}{1.160917in}}%
\pgfpathquadraticcurveto{\pgfqpoint{12.855940in}{1.169869in}}{\pgfqpoint{12.852446in}{1.174300in}}%
\pgfpathquadraticcurveto{\pgfqpoint{12.848952in}{1.178749in}}{\pgfqpoint{12.841981in}{1.178749in}}%
\pgfpathquadraticcurveto{\pgfqpoint{12.833587in}{1.178749in}}{\pgfqpoint{12.828742in}{1.173400in}}%
\pgfpathquadraticcurveto{\pgfqpoint{12.823897in}{1.168068in}}{\pgfqpoint{12.823897in}{1.158828in}}%
\pgfpathlineto{\pgfqpoint{12.823897in}{1.123200in}}%
\pgfpathlineto{\pgfqpoint{12.813486in}{1.123200in}}%
\pgfpathlineto{\pgfqpoint{12.813486in}{1.186243in}}%
\pgfpathlineto{\pgfqpoint{12.823897in}{1.186243in}}%
\pgfpathlineto{\pgfqpoint{12.823897in}{1.176444in}}%
\pgfpathquadraticcurveto{\pgfqpoint{12.827625in}{1.182136in}}{\pgfqpoint{12.832651in}{1.184946in}}%
\pgfpathquadraticcurveto{\pgfqpoint{12.837694in}{1.187756in}}{\pgfqpoint{12.844286in}{1.187756in}}%
\pgfpathquadraticcurveto{\pgfqpoint{12.855148in}{1.187756in}}{\pgfqpoint{12.860714in}{1.181037in}}%
\pgfpathquadraticcurveto{\pgfqpoint{12.866297in}{1.174318in}}{\pgfqpoint{12.866297in}{1.161260in}}%
\pgfpathlineto{\pgfqpoint{12.866297in}{1.161260in}}%
\pgfpathclose%
\pgfpathmoveto{\pgfqpoint{12.927124in}{1.184387in}}%
\pgfpathlineto{\pgfqpoint{12.927124in}{1.174589in}}%
\pgfpathquadraticcurveto{\pgfqpoint{12.922747in}{1.176840in}}{\pgfqpoint{12.918010in}{1.177957in}}%
\pgfpathquadraticcurveto{\pgfqpoint{12.913291in}{1.179092in}}{\pgfqpoint{12.908212in}{1.179092in}}%
\pgfpathquadraticcurveto{\pgfqpoint{12.900502in}{1.179092in}}{\pgfqpoint{12.896648in}{1.176732in}}%
\pgfpathquadraticcurveto{\pgfqpoint{12.892793in}{1.174373in}}{\pgfqpoint{12.892793in}{1.169635in}}%
\pgfpathquadraticcurveto{\pgfqpoint{12.892793in}{1.166033in}}{\pgfqpoint{12.895549in}{1.163980in}}%
\pgfpathquadraticcurveto{\pgfqpoint{12.898305in}{1.161926in}}{\pgfqpoint{12.906645in}{1.160071in}}%
\pgfpathlineto{\pgfqpoint{12.910193in}{1.159278in}}%
\pgfpathquadraticcurveto{\pgfqpoint{12.921216in}{1.156919in}}{\pgfqpoint{12.925863in}{1.152614in}}%
\pgfpathquadraticcurveto{\pgfqpoint{12.930511in}{1.148309in}}{\pgfqpoint{12.930511in}{1.140600in}}%
\pgfpathquadraticcurveto{\pgfqpoint{12.930511in}{1.131810in}}{\pgfqpoint{12.923558in}{1.126676in}}%
\pgfpathquadraticcurveto{\pgfqpoint{12.916605in}{1.121561in}}{\pgfqpoint{12.904447in}{1.121561in}}%
\pgfpathquadraticcurveto{\pgfqpoint{12.899386in}{1.121561in}}{\pgfqpoint{12.893892in}{1.122552in}}%
\pgfpathquadraticcurveto{\pgfqpoint{12.888398in}{1.123542in}}{\pgfqpoint{12.882328in}{1.125506in}}%
\pgfpathlineto{\pgfqpoint{12.882328in}{1.136205in}}%
\pgfpathquadraticcurveto{\pgfqpoint{12.888074in}{1.133215in}}{\pgfqpoint{12.893640in}{1.131720in}}%
\pgfpathquadraticcurveto{\pgfqpoint{12.899206in}{1.130243in}}{\pgfqpoint{12.904681in}{1.130243in}}%
\pgfpathquadraticcurveto{\pgfqpoint{12.911994in}{1.130243in}}{\pgfqpoint{12.915921in}{1.132746in}}%
\pgfpathquadraticcurveto{\pgfqpoint{12.919865in}{1.135250in}}{\pgfqpoint{12.919865in}{1.139807in}}%
\pgfpathquadraticcurveto{\pgfqpoint{12.919865in}{1.144022in}}{\pgfqpoint{12.917020in}{1.146274in}}%
\pgfpathquadraticcurveto{\pgfqpoint{12.914192in}{1.148525in}}{\pgfqpoint{12.904555in}{1.150614in}}%
\pgfpathlineto{\pgfqpoint{12.900953in}{1.151461in}}%
\pgfpathquadraticcurveto{\pgfqpoint{12.891334in}{1.153478in}}{\pgfqpoint{12.887047in}{1.157675in}}%
\pgfpathquadraticcurveto{\pgfqpoint{12.882778in}{1.161872in}}{\pgfqpoint{12.882778in}{1.169185in}}%
\pgfpathquadraticcurveto{\pgfqpoint{12.882778in}{1.178083in}}{\pgfqpoint{12.889083in}{1.182910in}}%
\pgfpathquadraticcurveto{\pgfqpoint{12.895387in}{1.187756in}}{\pgfqpoint{12.906987in}{1.187756in}}%
\pgfpathquadraticcurveto{\pgfqpoint{12.912715in}{1.187756in}}{\pgfqpoint{12.917776in}{1.186909in}}%
\pgfpathquadraticcurveto{\pgfqpoint{12.922855in}{1.186080in}}{\pgfqpoint{12.927124in}{1.184387in}}%
\pgfpathlineto{\pgfqpoint{12.927124in}{1.184387in}}%
\pgfpathclose%
\pgfpathmoveto{\pgfqpoint{12.949640in}{1.137502in}}%
\pgfpathlineto{\pgfqpoint{12.961510in}{1.137502in}}%
\pgfpathlineto{\pgfqpoint{12.961510in}{1.127811in}}%
\pgfpathlineto{\pgfqpoint{12.952287in}{1.109799in}}%
\pgfpathlineto{\pgfqpoint{12.945028in}{1.109799in}}%
\pgfpathlineto{\pgfqpoint{12.949640in}{1.127811in}}%
\pgfpathlineto{\pgfqpoint{12.949640in}{1.137502in}}%
\pgfpathlineto{\pgfqpoint{12.949640in}{1.137502in}}%
\pgfpathclose%
\pgfusepath{fill}%
\end{pgfscope}%
\begin{pgfscope}%
\pgfsetbuttcap%
\pgfsetroundjoin%
\definecolor{currentfill}{rgb}{0.309804,0.372549,0.419608}%
\pgfsetfillcolor{currentfill}%
\pgfsetlinewidth{0.000000pt}%
\definecolor{currentstroke}{rgb}{0.309804,0.372549,0.419608}%
\pgfsetstrokecolor{currentstroke}%
\pgfsetdash{}{0pt}%
\pgfpathmoveto{\pgfqpoint{13.072861in}{1.186243in}}%
\pgfpathlineto{\pgfqpoint{13.083218in}{1.186243in}}%
\pgfpathlineto{\pgfqpoint{13.083218in}{1.123200in}}%
\pgfpathlineto{\pgfqpoint{13.072861in}{1.123200in}}%
\pgfpathlineto{\pgfqpoint{13.072861in}{1.186243in}}%
\pgfpathlineto{\pgfqpoint{13.072861in}{1.186243in}}%
\pgfpathclose%
\pgfpathmoveto{\pgfqpoint{13.072861in}{1.210793in}}%
\pgfpathlineto{\pgfqpoint{13.083218in}{1.210793in}}%
\pgfpathlineto{\pgfqpoint{13.083218in}{1.197662in}}%
\pgfpathlineto{\pgfqpoint{13.072861in}{1.197662in}}%
\pgfpathlineto{\pgfqpoint{13.072861in}{1.210793in}}%
\pgfpathlineto{\pgfqpoint{13.072861in}{1.210793in}}%
\pgfpathclose%
\pgfpathmoveto{\pgfqpoint{13.157302in}{1.161260in}}%
\pgfpathlineto{\pgfqpoint{13.157302in}{1.123200in}}%
\pgfpathlineto{\pgfqpoint{13.146945in}{1.123200in}}%
\pgfpathlineto{\pgfqpoint{13.146945in}{1.160917in}}%
\pgfpathquadraticcurveto{\pgfqpoint{13.146945in}{1.169869in}}{\pgfqpoint{13.143451in}{1.174300in}}%
\pgfpathquadraticcurveto{\pgfqpoint{13.139957in}{1.178749in}}{\pgfqpoint{13.132986in}{1.178749in}}%
\pgfpathquadraticcurveto{\pgfqpoint{13.124592in}{1.178749in}}{\pgfqpoint{13.119747in}{1.173400in}}%
\pgfpathquadraticcurveto{\pgfqpoint{13.114902in}{1.168068in}}{\pgfqpoint{13.114902in}{1.158828in}}%
\pgfpathlineto{\pgfqpoint{13.114902in}{1.123200in}}%
\pgfpathlineto{\pgfqpoint{13.104491in}{1.123200in}}%
\pgfpathlineto{\pgfqpoint{13.104491in}{1.186243in}}%
\pgfpathlineto{\pgfqpoint{13.114902in}{1.186243in}}%
\pgfpathlineto{\pgfqpoint{13.114902in}{1.176444in}}%
\pgfpathquadraticcurveto{\pgfqpoint{13.118630in}{1.182136in}}{\pgfqpoint{13.123656in}{1.184946in}}%
\pgfpathquadraticcurveto{\pgfqpoint{13.128699in}{1.187756in}}{\pgfqpoint{13.135291in}{1.187756in}}%
\pgfpathquadraticcurveto{\pgfqpoint{13.146153in}{1.187756in}}{\pgfqpoint{13.151719in}{1.181037in}}%
\pgfpathquadraticcurveto{\pgfqpoint{13.157302in}{1.174318in}}{\pgfqpoint{13.157302in}{1.161260in}}%
\pgfpathlineto{\pgfqpoint{13.157302in}{1.161260in}}%
\pgfpathclose%
\pgfpathmoveto{\pgfqpoint{13.223317in}{1.183829in}}%
\pgfpathlineto{\pgfqpoint{13.223317in}{1.174138in}}%
\pgfpathquadraticcurveto{\pgfqpoint{13.218922in}{1.176570in}}{\pgfqpoint{13.214509in}{1.177777in}}%
\pgfpathquadraticcurveto{\pgfqpoint{13.210096in}{1.178984in}}{\pgfqpoint{13.205593in}{1.178984in}}%
\pgfpathquadraticcurveto{\pgfqpoint{13.195506in}{1.178984in}}{\pgfqpoint{13.189922in}{1.172589in}}%
\pgfpathquadraticcurveto{\pgfqpoint{13.184357in}{1.166213in}}{\pgfqpoint{13.184357in}{1.154667in}}%
\pgfpathquadraticcurveto{\pgfqpoint{13.184357in}{1.143121in}}{\pgfqpoint{13.189922in}{1.136727in}}%
\pgfpathquadraticcurveto{\pgfqpoint{13.195506in}{1.130351in}}{\pgfqpoint{13.205593in}{1.130351in}}%
\pgfpathquadraticcurveto{\pgfqpoint{13.210096in}{1.130351in}}{\pgfqpoint{13.214509in}{1.131558in}}%
\pgfpathquadraticcurveto{\pgfqpoint{13.218922in}{1.132764in}}{\pgfqpoint{13.223317in}{1.135196in}}%
\pgfpathlineto{\pgfqpoint{13.223317in}{1.125614in}}%
\pgfpathquadraticcurveto{\pgfqpoint{13.218976in}{1.123596in}}{\pgfqpoint{13.214329in}{1.122588in}}%
\pgfpathquadraticcurveto{\pgfqpoint{13.209700in}{1.121561in}}{\pgfqpoint{13.204458in}{1.121561in}}%
\pgfpathquadraticcurveto{\pgfqpoint{13.190211in}{1.121561in}}{\pgfqpoint{13.181817in}{1.130513in}}%
\pgfpathquadraticcurveto{\pgfqpoint{13.173441in}{1.139465in}}{\pgfqpoint{13.173441in}{1.154667in}}%
\pgfpathquadraticcurveto{\pgfqpoint{13.173441in}{1.170086in}}{\pgfqpoint{13.181907in}{1.178912in}}%
\pgfpathquadraticcurveto{\pgfqpoint{13.190391in}{1.187756in}}{\pgfqpoint{13.205143in}{1.187756in}}%
\pgfpathquadraticcurveto{\pgfqpoint{13.209916in}{1.187756in}}{\pgfqpoint{13.214473in}{1.186765in}}%
\pgfpathquadraticcurveto{\pgfqpoint{13.219030in}{1.185792in}}{\pgfqpoint{13.223317in}{1.183829in}}%
\pgfpathlineto{\pgfqpoint{13.223317in}{1.183829in}}%
\pgfpathclose%
\pgfpathmoveto{\pgfqpoint{13.241329in}{1.210793in}}%
\pgfpathlineto{\pgfqpoint{13.251686in}{1.210793in}}%
\pgfpathlineto{\pgfqpoint{13.251686in}{1.123200in}}%
\pgfpathlineto{\pgfqpoint{13.241329in}{1.123200in}}%
\pgfpathlineto{\pgfqpoint{13.241329in}{1.210793in}}%
\pgfpathlineto{\pgfqpoint{13.241329in}{1.210793in}}%
\pgfpathclose%
\pgfpathmoveto{\pgfqpoint{13.272292in}{1.148075in}}%
\pgfpathlineto{\pgfqpoint{13.272292in}{1.186243in}}%
\pgfpathlineto{\pgfqpoint{13.282649in}{1.186243in}}%
\pgfpathlineto{\pgfqpoint{13.282649in}{1.148471in}}%
\pgfpathquadraticcurveto{\pgfqpoint{13.282649in}{1.139519in}}{\pgfqpoint{13.286125in}{1.135034in}}%
\pgfpathquadraticcurveto{\pgfqpoint{13.289620in}{1.130567in}}{\pgfqpoint{13.296608in}{1.130567in}}%
\pgfpathquadraticcurveto{\pgfqpoint{13.304984in}{1.130567in}}{\pgfqpoint{13.309847in}{1.135917in}}%
\pgfpathquadraticcurveto{\pgfqpoint{13.314729in}{1.141266in}}{\pgfqpoint{13.314729in}{1.150506in}}%
\pgfpathlineto{\pgfqpoint{13.314729in}{1.186243in}}%
\pgfpathlineto{\pgfqpoint{13.325086in}{1.186243in}}%
\pgfpathlineto{\pgfqpoint{13.325086in}{1.123200in}}%
\pgfpathlineto{\pgfqpoint{13.314729in}{1.123200in}}%
\pgfpathlineto{\pgfqpoint{13.314729in}{1.132891in}}%
\pgfpathquadraticcurveto{\pgfqpoint{13.310964in}{1.127145in}}{\pgfqpoint{13.305975in}{1.124353in}}%
\pgfpathquadraticcurveto{\pgfqpoint{13.301003in}{1.121561in}}{\pgfqpoint{13.294411in}{1.121561in}}%
\pgfpathquadraticcurveto{\pgfqpoint{13.283549in}{1.121561in}}{\pgfqpoint{13.277912in}{1.128315in}}%
\pgfpathquadraticcurveto{\pgfqpoint{13.272292in}{1.135070in}}{\pgfqpoint{13.272292in}{1.148075in}}%
\pgfpathlineto{\pgfqpoint{13.272292in}{1.148075in}}%
\pgfpathclose%
\pgfpathmoveto{\pgfqpoint{13.298355in}{1.187756in}}%
\pgfpathlineto{\pgfqpoint{13.298355in}{1.187756in}}%
\pgfpathlineto{\pgfqpoint{13.298355in}{1.187756in}}%
\pgfpathclose%
\pgfpathmoveto{\pgfqpoint{13.387894in}{1.176678in}}%
\pgfpathlineto{\pgfqpoint{13.387894in}{1.210793in}}%
\pgfpathlineto{\pgfqpoint{13.398251in}{1.210793in}}%
\pgfpathlineto{\pgfqpoint{13.398251in}{1.123200in}}%
\pgfpathlineto{\pgfqpoint{13.387894in}{1.123200in}}%
\pgfpathlineto{\pgfqpoint{13.387894in}{1.132656in}}%
\pgfpathquadraticcurveto{\pgfqpoint{13.384634in}{1.127037in}}{\pgfqpoint{13.379644in}{1.124299in}}%
\pgfpathquadraticcurveto{\pgfqpoint{13.374673in}{1.121561in}}{\pgfqpoint{13.367684in}{1.121561in}}%
\pgfpathquadraticcurveto{\pgfqpoint{13.356265in}{1.121561in}}{\pgfqpoint{13.349078in}{1.130675in}}%
\pgfpathquadraticcurveto{\pgfqpoint{13.341909in}{1.139807in}}{\pgfqpoint{13.341909in}{1.154667in}}%
\pgfpathquadraticcurveto{\pgfqpoint{13.341909in}{1.169527in}}{\pgfqpoint{13.349078in}{1.178641in}}%
\pgfpathquadraticcurveto{\pgfqpoint{13.356265in}{1.187756in}}{\pgfqpoint{13.367684in}{1.187756in}}%
\pgfpathquadraticcurveto{\pgfqpoint{13.374673in}{1.187756in}}{\pgfqpoint{13.379644in}{1.185018in}}%
\pgfpathquadraticcurveto{\pgfqpoint{13.384634in}{1.182298in}}{\pgfqpoint{13.387894in}{1.176678in}}%
\pgfpathlineto{\pgfqpoint{13.387894in}{1.176678in}}%
\pgfpathclose%
\pgfpathmoveto{\pgfqpoint{13.352608in}{1.154667in}}%
\pgfpathquadraticcurveto{\pgfqpoint{13.352608in}{1.143248in}}{\pgfqpoint{13.357309in}{1.136745in}}%
\pgfpathquadraticcurveto{\pgfqpoint{13.362010in}{1.130243in}}{\pgfqpoint{13.370224in}{1.130243in}}%
\pgfpathquadraticcurveto{\pgfqpoint{13.378437in}{1.130243in}}{\pgfqpoint{13.383157in}{1.136745in}}%
\pgfpathquadraticcurveto{\pgfqpoint{13.387894in}{1.143248in}}{\pgfqpoint{13.387894in}{1.154667in}}%
\pgfpathquadraticcurveto{\pgfqpoint{13.387894in}{1.166087in}}{\pgfqpoint{13.383157in}{1.172589in}}%
\pgfpathquadraticcurveto{\pgfqpoint{13.378437in}{1.179092in}}{\pgfqpoint{13.370224in}{1.179092in}}%
\pgfpathquadraticcurveto{\pgfqpoint{13.362010in}{1.179092in}}{\pgfqpoint{13.357309in}{1.172589in}}%
\pgfpathquadraticcurveto{\pgfqpoint{13.352608in}{1.166087in}}{\pgfqpoint{13.352608in}{1.154667in}}%
\pgfpathlineto{\pgfqpoint{13.352608in}{1.154667in}}%
\pgfpathclose%
\pgfpathmoveto{\pgfqpoint{13.419595in}{1.186243in}}%
\pgfpathlineto{\pgfqpoint{13.429952in}{1.186243in}}%
\pgfpathlineto{\pgfqpoint{13.429952in}{1.123200in}}%
\pgfpathlineto{\pgfqpoint{13.419595in}{1.123200in}}%
\pgfpathlineto{\pgfqpoint{13.419595in}{1.186243in}}%
\pgfpathlineto{\pgfqpoint{13.419595in}{1.186243in}}%
\pgfpathclose%
\pgfpathmoveto{\pgfqpoint{13.419595in}{1.210793in}}%
\pgfpathlineto{\pgfqpoint{13.429952in}{1.210793in}}%
\pgfpathlineto{\pgfqpoint{13.429952in}{1.197662in}}%
\pgfpathlineto{\pgfqpoint{13.419595in}{1.197662in}}%
\pgfpathlineto{\pgfqpoint{13.419595in}{1.210793in}}%
\pgfpathlineto{\pgfqpoint{13.419595in}{1.210793in}}%
\pgfpathclose%
\pgfpathmoveto{\pgfqpoint{13.504036in}{1.161260in}}%
\pgfpathlineto{\pgfqpoint{13.504036in}{1.123200in}}%
\pgfpathlineto{\pgfqpoint{13.493679in}{1.123200in}}%
\pgfpathlineto{\pgfqpoint{13.493679in}{1.160917in}}%
\pgfpathquadraticcurveto{\pgfqpoint{13.493679in}{1.169869in}}{\pgfqpoint{13.490185in}{1.174300in}}%
\pgfpathquadraticcurveto{\pgfqpoint{13.486691in}{1.178749in}}{\pgfqpoint{13.479720in}{1.178749in}}%
\pgfpathquadraticcurveto{\pgfqpoint{13.471326in}{1.178749in}}{\pgfqpoint{13.466481in}{1.173400in}}%
\pgfpathquadraticcurveto{\pgfqpoint{13.461636in}{1.168068in}}{\pgfqpoint{13.461636in}{1.158828in}}%
\pgfpathlineto{\pgfqpoint{13.461636in}{1.123200in}}%
\pgfpathlineto{\pgfqpoint{13.451225in}{1.123200in}}%
\pgfpathlineto{\pgfqpoint{13.451225in}{1.186243in}}%
\pgfpathlineto{\pgfqpoint{13.461636in}{1.186243in}}%
\pgfpathlineto{\pgfqpoint{13.461636in}{1.176444in}}%
\pgfpathquadraticcurveto{\pgfqpoint{13.465364in}{1.182136in}}{\pgfqpoint{13.470390in}{1.184946in}}%
\pgfpathquadraticcurveto{\pgfqpoint{13.475433in}{1.187756in}}{\pgfqpoint{13.482025in}{1.187756in}}%
\pgfpathquadraticcurveto{\pgfqpoint{13.492887in}{1.187756in}}{\pgfqpoint{13.498452in}{1.181037in}}%
\pgfpathquadraticcurveto{\pgfqpoint{13.504036in}{1.174318in}}{\pgfqpoint{13.504036in}{1.161260in}}%
\pgfpathlineto{\pgfqpoint{13.504036in}{1.161260in}}%
\pgfpathclose%
\pgfpathmoveto{\pgfqpoint{13.566160in}{1.155460in}}%
\pgfpathquadraticcurveto{\pgfqpoint{13.566160in}{1.166717in}}{\pgfqpoint{13.561513in}{1.172896in}}%
\pgfpathquadraticcurveto{\pgfqpoint{13.556884in}{1.179092in}}{\pgfqpoint{13.548490in}{1.179092in}}%
\pgfpathquadraticcurveto{\pgfqpoint{13.540169in}{1.179092in}}{\pgfqpoint{13.535521in}{1.172896in}}%
\pgfpathquadraticcurveto{\pgfqpoint{13.530874in}{1.166717in}}{\pgfqpoint{13.530874in}{1.155460in}}%
\pgfpathquadraticcurveto{\pgfqpoint{13.530874in}{1.144256in}}{\pgfqpoint{13.535521in}{1.138060in}}%
\pgfpathquadraticcurveto{\pgfqpoint{13.540169in}{1.131864in}}{\pgfqpoint{13.548490in}{1.131864in}}%
\pgfpathquadraticcurveto{\pgfqpoint{13.556884in}{1.131864in}}{\pgfqpoint{13.561513in}{1.138060in}}%
\pgfpathquadraticcurveto{\pgfqpoint{13.566160in}{1.144256in}}{\pgfqpoint{13.566160in}{1.155460in}}%
\pgfpathlineto{\pgfqpoint{13.566160in}{1.155460in}}%
\pgfpathclose%
\pgfpathmoveto{\pgfqpoint{13.576517in}{1.131017in}}%
\pgfpathquadraticcurveto{\pgfqpoint{13.576517in}{1.114932in}}{\pgfqpoint{13.569366in}{1.107079in}}%
\pgfpathquadraticcurveto{\pgfqpoint{13.562234in}{1.099226in}}{\pgfqpoint{13.547482in}{1.099226in}}%
\pgfpathquadraticcurveto{\pgfqpoint{13.542024in}{1.099226in}}{\pgfqpoint{13.537179in}{1.100036in}}%
\pgfpathquadraticcurveto{\pgfqpoint{13.532333in}{1.100847in}}{\pgfqpoint{13.527776in}{1.102540in}}%
\pgfpathlineto{\pgfqpoint{13.527776in}{1.112609in}}%
\pgfpathquadraticcurveto{\pgfqpoint{13.532333in}{1.110141in}}{\pgfqpoint{13.536782in}{1.108970in}}%
\pgfpathquadraticcurveto{\pgfqpoint{13.541231in}{1.107782in}}{\pgfqpoint{13.545842in}{1.107782in}}%
\pgfpathquadraticcurveto{\pgfqpoint{13.556037in}{1.107782in}}{\pgfqpoint{13.561099in}{1.113095in}}%
\pgfpathquadraticcurveto{\pgfqpoint{13.566160in}{1.118409in}}{\pgfqpoint{13.566160in}{1.129162in}}%
\pgfpathlineto{\pgfqpoint{13.566160in}{1.134295in}}%
\pgfpathquadraticcurveto{\pgfqpoint{13.562954in}{1.128712in}}{\pgfqpoint{13.557947in}{1.125956in}}%
\pgfpathquadraticcurveto{\pgfqpoint{13.552939in}{1.123200in}}{\pgfqpoint{13.545951in}{1.123200in}}%
\pgfpathquadraticcurveto{\pgfqpoint{13.534369in}{1.123200in}}{\pgfqpoint{13.527272in}{1.132026in}}%
\pgfpathquadraticcurveto{\pgfqpoint{13.520175in}{1.140870in}}{\pgfqpoint{13.520175in}{1.155460in}}%
\pgfpathquadraticcurveto{\pgfqpoint{13.520175in}{1.170086in}}{\pgfqpoint{13.527272in}{1.178912in}}%
\pgfpathquadraticcurveto{\pgfqpoint{13.534369in}{1.187756in}}{\pgfqpoint{13.545951in}{1.187756in}}%
\pgfpathquadraticcurveto{\pgfqpoint{13.552939in}{1.187756in}}{\pgfqpoint{13.557947in}{1.185000in}}%
\pgfpathquadraticcurveto{\pgfqpoint{13.562954in}{1.182244in}}{\pgfqpoint{13.566160in}{1.176678in}}%
\pgfpathlineto{\pgfqpoint{13.566160in}{1.186243in}}%
\pgfpathlineto{\pgfqpoint{13.576517in}{1.186243in}}%
\pgfpathlineto{\pgfqpoint{13.576517in}{1.131017in}}%
\pgfpathlineto{\pgfqpoint{13.576517in}{1.131017in}}%
\pgfpathclose%
\pgfusepath{fill}%
\end{pgfscope}%
\begin{pgfscope}%
\pgfsetbuttcap%
\pgfsetroundjoin%
\definecolor{currentfill}{rgb}{0.309804,0.372549,0.419608}%
\pgfsetfillcolor{currentfill}%
\pgfsetlinewidth{0.000000pt}%
\definecolor{currentstroke}{rgb}{0.309804,0.372549,0.419608}%
\pgfsetstrokecolor{currentstroke}%
\pgfsetdash{}{0pt}%
\pgfpathmoveto{\pgfqpoint{0.956980in}{1.047593in}}%
\pgfpathlineto{\pgfqpoint{0.985800in}{1.047593in}}%
\pgfpathlineto{\pgfqpoint{0.985800in}{0.960000in}}%
\pgfpathlineto{\pgfqpoint{0.975389in}{0.960000in}}%
\pgfpathlineto{\pgfqpoint{0.975389in}{1.038965in}}%
\pgfpathlineto{\pgfqpoint{0.956872in}{1.038965in}}%
\pgfpathquadraticcurveto{\pgfqpoint{0.951306in}{1.038965in}}{\pgfqpoint{0.949127in}{1.036714in}}%
\pgfpathquadraticcurveto{\pgfqpoint{0.946965in}{1.034462in}}{\pgfqpoint{0.946965in}{1.028608in}}%
\pgfpathlineto{\pgfqpoint{0.946965in}{1.023043in}}%
\pgfpathlineto{\pgfqpoint{0.964023in}{1.023043in}}%
\pgfpathlineto{\pgfqpoint{0.964023in}{1.014991in}}%
\pgfpathlineto{\pgfqpoint{0.946965in}{1.014991in}}%
\pgfpathlineto{\pgfqpoint{0.946965in}{0.960000in}}%
\pgfpathlineto{\pgfqpoint{0.936554in}{0.960000in}}%
\pgfpathlineto{\pgfqpoint{0.936554in}{1.014991in}}%
\pgfpathlineto{\pgfqpoint{0.926648in}{1.014991in}}%
\pgfpathlineto{\pgfqpoint{0.926648in}{1.023043in}}%
\pgfpathlineto{\pgfqpoint{0.936554in}{1.023043in}}%
\pgfpathlineto{\pgfqpoint{0.936554in}{1.027438in}}%
\pgfpathquadraticcurveto{\pgfqpoint{0.936554in}{1.037957in}}{\pgfqpoint{0.941454in}{1.042766in}}%
\pgfpathquadraticcurveto{\pgfqpoint{0.946353in}{1.047593in}}{\pgfqpoint{0.956980in}{1.047593in}}%
\pgfpathlineto{\pgfqpoint{0.956980in}{1.047593in}}%
\pgfpathclose%
\pgfpathmoveto{\pgfqpoint{1.036126in}{0.991683in}}%
\pgfpathquadraticcurveto{\pgfqpoint{1.023571in}{0.991683in}}{\pgfqpoint{1.018726in}{0.988819in}}%
\pgfpathquadraticcurveto{\pgfqpoint{1.013881in}{0.985956in}}{\pgfqpoint{1.013881in}{0.979021in}}%
\pgfpathquadraticcurveto{\pgfqpoint{1.013881in}{0.973509in}}{\pgfqpoint{1.017519in}{0.970267in}}%
\pgfpathquadraticcurveto{\pgfqpoint{1.021158in}{0.967043in}}{\pgfqpoint{1.027390in}{0.967043in}}%
\pgfpathquadraticcurveto{\pgfqpoint{1.036018in}{0.967043in}}{\pgfqpoint{1.041223in}{0.973149in}}%
\pgfpathquadraticcurveto{\pgfqpoint{1.046429in}{0.979255in}}{\pgfqpoint{1.046429in}{0.989378in}}%
\pgfpathlineto{\pgfqpoint{1.046429in}{0.991683in}}%
\pgfpathlineto{\pgfqpoint{1.036126in}{0.991683in}}%
\pgfpathlineto{\pgfqpoint{1.036126in}{0.991683in}}%
\pgfpathclose%
\pgfpathmoveto{\pgfqpoint{1.056786in}{0.995970in}}%
\pgfpathlineto{\pgfqpoint{1.056786in}{0.960000in}}%
\pgfpathlineto{\pgfqpoint{1.046429in}{0.960000in}}%
\pgfpathlineto{\pgfqpoint{1.046429in}{0.969564in}}%
\pgfpathquadraticcurveto{\pgfqpoint{1.042880in}{0.963837in}}{\pgfqpoint{1.037585in}{0.961099in}}%
\pgfpathquadraticcurveto{\pgfqpoint{1.032289in}{0.958361in}}{\pgfqpoint{1.024634in}{0.958361in}}%
\pgfpathquadraticcurveto{\pgfqpoint{1.014961in}{0.958361in}}{\pgfqpoint{1.009234in}{0.963801in}}%
\pgfpathquadraticcurveto{\pgfqpoint{1.003524in}{0.969240in}}{\pgfqpoint{1.003524in}{0.978354in}}%
\pgfpathquadraticcurveto{\pgfqpoint{1.003524in}{0.988982in}}{\pgfqpoint{1.010638in}{0.994385in}}%
\pgfpathquadraticcurveto{\pgfqpoint{1.017771in}{0.999789in}}{\pgfqpoint{1.031893in}{0.999789in}}%
\pgfpathlineto{\pgfqpoint{1.046429in}{0.999789in}}%
\pgfpathlineto{\pgfqpoint{1.046429in}{1.000816in}}%
\pgfpathquadraticcurveto{\pgfqpoint{1.046429in}{1.007966in}}{\pgfqpoint{1.041727in}{1.011875in}}%
\pgfpathquadraticcurveto{\pgfqpoint{1.037026in}{1.015784in}}{\pgfqpoint{1.028525in}{1.015784in}}%
\pgfpathquadraticcurveto{\pgfqpoint{1.023121in}{1.015784in}}{\pgfqpoint{1.017987in}{1.014487in}}%
\pgfpathquadraticcurveto{\pgfqpoint{1.012872in}{1.013190in}}{\pgfqpoint{1.008153in}{1.010596in}}%
\pgfpathlineto{\pgfqpoint{1.008153in}{1.020179in}}%
\pgfpathquadraticcurveto{\pgfqpoint{1.013827in}{1.022376in}}{\pgfqpoint{1.019176in}{1.023457in}}%
\pgfpathquadraticcurveto{\pgfqpoint{1.024526in}{1.024556in}}{\pgfqpoint{1.029587in}{1.024556in}}%
\pgfpathquadraticcurveto{\pgfqpoint{1.043276in}{1.024556in}}{\pgfqpoint{1.050031in}{1.017459in}}%
\pgfpathquadraticcurveto{\pgfqpoint{1.056786in}{1.010380in}}{\pgfqpoint{1.056786in}{0.995970in}}%
\pgfpathlineto{\pgfqpoint{1.056786in}{0.995970in}}%
\pgfpathclose%
\pgfpathmoveto{\pgfqpoint{1.119594in}{0.992260in}}%
\pgfpathquadraticcurveto{\pgfqpoint{1.119594in}{1.003517in}}{\pgfqpoint{1.114947in}{1.009696in}}%
\pgfpathquadraticcurveto{\pgfqpoint{1.110318in}{1.015892in}}{\pgfqpoint{1.101924in}{1.015892in}}%
\pgfpathquadraticcurveto{\pgfqpoint{1.093602in}{1.015892in}}{\pgfqpoint{1.088955in}{1.009696in}}%
\pgfpathquadraticcurveto{\pgfqpoint{1.084308in}{1.003517in}}{\pgfqpoint{1.084308in}{0.992260in}}%
\pgfpathquadraticcurveto{\pgfqpoint{1.084308in}{0.981056in}}{\pgfqpoint{1.088955in}{0.974860in}}%
\pgfpathquadraticcurveto{\pgfqpoint{1.093602in}{0.968664in}}{\pgfqpoint{1.101924in}{0.968664in}}%
\pgfpathquadraticcurveto{\pgfqpoint{1.110318in}{0.968664in}}{\pgfqpoint{1.114947in}{0.974860in}}%
\pgfpathquadraticcurveto{\pgfqpoint{1.119594in}{0.981056in}}{\pgfqpoint{1.119594in}{0.992260in}}%
\pgfpathlineto{\pgfqpoint{1.119594in}{0.992260in}}%
\pgfpathclose%
\pgfpathmoveto{\pgfqpoint{1.129951in}{0.967817in}}%
\pgfpathquadraticcurveto{\pgfqpoint{1.129951in}{0.951732in}}{\pgfqpoint{1.122800in}{0.943879in}}%
\pgfpathquadraticcurveto{\pgfqpoint{1.115667in}{0.936026in}}{\pgfqpoint{1.100915in}{0.936026in}}%
\pgfpathquadraticcurveto{\pgfqpoint{1.095458in}{0.936026in}}{\pgfqpoint{1.090612in}{0.936836in}}%
\pgfpathquadraticcurveto{\pgfqpoint{1.085767in}{0.937647in}}{\pgfqpoint{1.081210in}{0.939340in}}%
\pgfpathlineto{\pgfqpoint{1.081210in}{0.949409in}}%
\pgfpathquadraticcurveto{\pgfqpoint{1.085767in}{0.946941in}}{\pgfqpoint{1.090216in}{0.945770in}}%
\pgfpathquadraticcurveto{\pgfqpoint{1.094665in}{0.944582in}}{\pgfqpoint{1.099276in}{0.944582in}}%
\pgfpathquadraticcurveto{\pgfqpoint{1.109471in}{0.944582in}}{\pgfqpoint{1.114533in}{0.949895in}}%
\pgfpathquadraticcurveto{\pgfqpoint{1.119594in}{0.955209in}}{\pgfqpoint{1.119594in}{0.965962in}}%
\pgfpathlineto{\pgfqpoint{1.119594in}{0.971095in}}%
\pgfpathquadraticcurveto{\pgfqpoint{1.116388in}{0.965512in}}{\pgfqpoint{1.111380in}{0.962756in}}%
\pgfpathquadraticcurveto{\pgfqpoint{1.106373in}{0.960000in}}{\pgfqpoint{1.099384in}{0.960000in}}%
\pgfpathquadraticcurveto{\pgfqpoint{1.087803in}{0.960000in}}{\pgfqpoint{1.080706in}{0.968826in}}%
\pgfpathquadraticcurveto{\pgfqpoint{1.073609in}{0.977670in}}{\pgfqpoint{1.073609in}{0.992260in}}%
\pgfpathquadraticcurveto{\pgfqpoint{1.073609in}{1.006886in}}{\pgfqpoint{1.080706in}{1.015712in}}%
\pgfpathquadraticcurveto{\pgfqpoint{1.087803in}{1.024556in}}{\pgfqpoint{1.099384in}{1.024556in}}%
\pgfpathquadraticcurveto{\pgfqpoint{1.106373in}{1.024556in}}{\pgfqpoint{1.111380in}{1.021800in}}%
\pgfpathquadraticcurveto{\pgfqpoint{1.116388in}{1.019044in}}{\pgfqpoint{1.119594in}{1.013478in}}%
\pgfpathlineto{\pgfqpoint{1.119594in}{1.023043in}}%
\pgfpathlineto{\pgfqpoint{1.129951in}{1.023043in}}%
\pgfpathlineto{\pgfqpoint{1.129951in}{0.967817in}}%
\pgfpathlineto{\pgfqpoint{1.129951in}{0.967817in}}%
\pgfpathclose%
\pgfpathmoveto{\pgfqpoint{1.192777in}{0.992260in}}%
\pgfpathquadraticcurveto{\pgfqpoint{1.192777in}{1.003517in}}{\pgfqpoint{1.188130in}{1.009696in}}%
\pgfpathquadraticcurveto{\pgfqpoint{1.183501in}{1.015892in}}{\pgfqpoint{1.175107in}{1.015892in}}%
\pgfpathquadraticcurveto{\pgfqpoint{1.166786in}{1.015892in}}{\pgfqpoint{1.162139in}{1.009696in}}%
\pgfpathquadraticcurveto{\pgfqpoint{1.157492in}{1.003517in}}{\pgfqpoint{1.157492in}{0.992260in}}%
\pgfpathquadraticcurveto{\pgfqpoint{1.157492in}{0.981056in}}{\pgfqpoint{1.162139in}{0.974860in}}%
\pgfpathquadraticcurveto{\pgfqpoint{1.166786in}{0.968664in}}{\pgfqpoint{1.175107in}{0.968664in}}%
\pgfpathquadraticcurveto{\pgfqpoint{1.183501in}{0.968664in}}{\pgfqpoint{1.188130in}{0.974860in}}%
\pgfpathquadraticcurveto{\pgfqpoint{1.192777in}{0.981056in}}{\pgfqpoint{1.192777in}{0.992260in}}%
\pgfpathlineto{\pgfqpoint{1.192777in}{0.992260in}}%
\pgfpathclose%
\pgfpathmoveto{\pgfqpoint{1.203134in}{0.967817in}}%
\pgfpathquadraticcurveto{\pgfqpoint{1.203134in}{0.951732in}}{\pgfqpoint{1.195984in}{0.943879in}}%
\pgfpathquadraticcurveto{\pgfqpoint{1.188851in}{0.936026in}}{\pgfqpoint{1.174099in}{0.936026in}}%
\pgfpathquadraticcurveto{\pgfqpoint{1.168641in}{0.936026in}}{\pgfqpoint{1.163796in}{0.936836in}}%
\pgfpathquadraticcurveto{\pgfqpoint{1.158951in}{0.937647in}}{\pgfqpoint{1.154393in}{0.939340in}}%
\pgfpathlineto{\pgfqpoint{1.154393in}{0.949409in}}%
\pgfpathquadraticcurveto{\pgfqpoint{1.158951in}{0.946941in}}{\pgfqpoint{1.163400in}{0.945770in}}%
\pgfpathquadraticcurveto{\pgfqpoint{1.167849in}{0.944582in}}{\pgfqpoint{1.172460in}{0.944582in}}%
\pgfpathquadraticcurveto{\pgfqpoint{1.182655in}{0.944582in}}{\pgfqpoint{1.187716in}{0.949895in}}%
\pgfpathquadraticcurveto{\pgfqpoint{1.192777in}{0.955209in}}{\pgfqpoint{1.192777in}{0.965962in}}%
\pgfpathlineto{\pgfqpoint{1.192777in}{0.971095in}}%
\pgfpathquadraticcurveto{\pgfqpoint{1.189571in}{0.965512in}}{\pgfqpoint{1.184564in}{0.962756in}}%
\pgfpathquadraticcurveto{\pgfqpoint{1.179556in}{0.960000in}}{\pgfqpoint{1.172568in}{0.960000in}}%
\pgfpathquadraticcurveto{\pgfqpoint{1.160986in}{0.960000in}}{\pgfqpoint{1.153889in}{0.968826in}}%
\pgfpathquadraticcurveto{\pgfqpoint{1.146792in}{0.977670in}}{\pgfqpoint{1.146792in}{0.992260in}}%
\pgfpathquadraticcurveto{\pgfqpoint{1.146792in}{1.006886in}}{\pgfqpoint{1.153889in}{1.015712in}}%
\pgfpathquadraticcurveto{\pgfqpoint{1.160986in}{1.024556in}}{\pgfqpoint{1.172568in}{1.024556in}}%
\pgfpathquadraticcurveto{\pgfqpoint{1.179556in}{1.024556in}}{\pgfqpoint{1.184564in}{1.021800in}}%
\pgfpathquadraticcurveto{\pgfqpoint{1.189571in}{1.019044in}}{\pgfqpoint{1.192777in}{1.013478in}}%
\pgfpathlineto{\pgfqpoint{1.192777in}{1.023043in}}%
\pgfpathlineto{\pgfqpoint{1.203134in}{1.023043in}}%
\pgfpathlineto{\pgfqpoint{1.203134in}{0.967817in}}%
\pgfpathlineto{\pgfqpoint{1.203134in}{0.967817in}}%
\pgfpathclose%
\pgfpathmoveto{\pgfqpoint{1.278407in}{0.994115in}}%
\pgfpathlineto{\pgfqpoint{1.278407in}{0.989054in}}%
\pgfpathlineto{\pgfqpoint{1.230783in}{0.989054in}}%
\pgfpathquadraticcurveto{\pgfqpoint{1.231467in}{0.978354in}}{\pgfqpoint{1.237231in}{0.972753in}}%
\pgfpathquadraticcurveto{\pgfqpoint{1.242995in}{0.967151in}}{\pgfqpoint{1.253298in}{0.967151in}}%
\pgfpathquadraticcurveto{\pgfqpoint{1.259260in}{0.967151in}}{\pgfqpoint{1.264862in}{0.968610in}}%
\pgfpathquadraticcurveto{\pgfqpoint{1.270464in}{0.970069in}}{\pgfqpoint{1.275993in}{0.973005in}}%
\pgfpathlineto{\pgfqpoint{1.275993in}{0.963206in}}%
\pgfpathquadraticcurveto{\pgfqpoint{1.270410in}{0.960847in}}{\pgfqpoint{1.264556in}{0.959604in}}%
\pgfpathquadraticcurveto{\pgfqpoint{1.258702in}{0.958361in}}{\pgfqpoint{1.252686in}{0.958361in}}%
\pgfpathquadraticcurveto{\pgfqpoint{1.237592in}{0.958361in}}{\pgfqpoint{1.228784in}{0.967133in}}%
\pgfpathquadraticcurveto{\pgfqpoint{1.219976in}{0.975923in}}{\pgfqpoint{1.219976in}{0.990909in}}%
\pgfpathquadraticcurveto{\pgfqpoint{1.219976in}{1.006381in}}{\pgfqpoint{1.228333in}{1.015459in}}%
\pgfpathquadraticcurveto{\pgfqpoint{1.236691in}{1.024556in}}{\pgfqpoint{1.250885in}{1.024556in}}%
\pgfpathquadraticcurveto{\pgfqpoint{1.263601in}{1.024556in}}{\pgfqpoint{1.271004in}{1.016360in}}%
\pgfpathquadraticcurveto{\pgfqpoint{1.278407in}{1.008183in}}{\pgfqpoint{1.278407in}{0.994115in}}%
\pgfpathlineto{\pgfqpoint{1.278407in}{0.994115in}}%
\pgfpathclose%
\pgfpathmoveto{\pgfqpoint{1.268050in}{0.997159in}}%
\pgfpathquadraticcurveto{\pgfqpoint{1.267942in}{1.005643in}}{\pgfqpoint{1.263295in}{1.010704in}}%
\pgfpathquadraticcurveto{\pgfqpoint{1.258648in}{1.015784in}}{\pgfqpoint{1.250993in}{1.015784in}}%
\pgfpathquadraticcurveto{\pgfqpoint{1.242329in}{1.015784in}}{\pgfqpoint{1.237123in}{1.010884in}}%
\pgfpathquadraticcurveto{\pgfqpoint{1.231918in}{1.005985in}}{\pgfqpoint{1.231125in}{0.997087in}}%
\pgfpathlineto{\pgfqpoint{1.268050in}{0.997159in}}%
\pgfpathlineto{\pgfqpoint{1.268050in}{0.997159in}}%
\pgfpathclose%
\pgfpathmoveto{\pgfqpoint{1.336893in}{1.013478in}}%
\pgfpathlineto{\pgfqpoint{1.336893in}{1.047593in}}%
\pgfpathlineto{\pgfqpoint{1.347250in}{1.047593in}}%
\pgfpathlineto{\pgfqpoint{1.347250in}{0.960000in}}%
\pgfpathlineto{\pgfqpoint{1.336893in}{0.960000in}}%
\pgfpathlineto{\pgfqpoint{1.336893in}{0.969456in}}%
\pgfpathquadraticcurveto{\pgfqpoint{1.333632in}{0.963837in}}{\pgfqpoint{1.328643in}{0.961099in}}%
\pgfpathquadraticcurveto{\pgfqpoint{1.323672in}{0.958361in}}{\pgfqpoint{1.316683in}{0.958361in}}%
\pgfpathquadraticcurveto{\pgfqpoint{1.305263in}{0.958361in}}{\pgfqpoint{1.298076in}{0.967475in}}%
\pgfpathquadraticcurveto{\pgfqpoint{1.290908in}{0.976607in}}{\pgfqpoint{1.290908in}{0.991467in}}%
\pgfpathquadraticcurveto{\pgfqpoint{1.290908in}{1.006327in}}{\pgfqpoint{1.298076in}{1.015441in}}%
\pgfpathquadraticcurveto{\pgfqpoint{1.305263in}{1.024556in}}{\pgfqpoint{1.316683in}{1.024556in}}%
\pgfpathquadraticcurveto{\pgfqpoint{1.323672in}{1.024556in}}{\pgfqpoint{1.328643in}{1.021818in}}%
\pgfpathquadraticcurveto{\pgfqpoint{1.333632in}{1.019098in}}{\pgfqpoint{1.336893in}{1.013478in}}%
\pgfpathlineto{\pgfqpoint{1.336893in}{1.013478in}}%
\pgfpathclose%
\pgfpathmoveto{\pgfqpoint{1.301607in}{0.991467in}}%
\pgfpathquadraticcurveto{\pgfqpoint{1.301607in}{0.980048in}}{\pgfqpoint{1.306308in}{0.973545in}}%
\pgfpathquadraticcurveto{\pgfqpoint{1.311009in}{0.967043in}}{\pgfqpoint{1.319223in}{0.967043in}}%
\pgfpathquadraticcurveto{\pgfqpoint{1.327436in}{0.967043in}}{\pgfqpoint{1.332155in}{0.973545in}}%
\pgfpathquadraticcurveto{\pgfqpoint{1.336893in}{0.980048in}}{\pgfqpoint{1.336893in}{0.991467in}}%
\pgfpathquadraticcurveto{\pgfqpoint{1.336893in}{1.002887in}}{\pgfqpoint{1.332155in}{1.009389in}}%
\pgfpathquadraticcurveto{\pgfqpoint{1.327436in}{1.015892in}}{\pgfqpoint{1.319223in}{1.015892in}}%
\pgfpathquadraticcurveto{\pgfqpoint{1.311009in}{1.015892in}}{\pgfqpoint{1.306308in}{1.009389in}}%
\pgfpathquadraticcurveto{\pgfqpoint{1.301607in}{1.002887in}}{\pgfqpoint{1.301607in}{0.991467in}}%
\pgfpathlineto{\pgfqpoint{1.301607in}{0.991467in}}%
\pgfpathclose%
\pgfusepath{fill}%
\end{pgfscope}%
\begin{pgfscope}%
\pgfsetbuttcap%
\pgfsetroundjoin%
\definecolor{currentfill}{rgb}{0.309804,0.372549,0.419608}%
\pgfsetfillcolor{currentfill}%
\pgfsetlinewidth{0.000000pt}%
\definecolor{currentstroke}{rgb}{0.309804,0.372549,0.419608}%
\pgfsetstrokecolor{currentstroke}%
\pgfsetdash{}{0pt}%
\pgfpathmoveto{\pgfqpoint{1.454769in}{0.969456in}}%
\pgfpathlineto{\pgfqpoint{1.454769in}{0.936026in}}%
\pgfpathlineto{\pgfqpoint{1.444358in}{0.936026in}}%
\pgfpathlineto{\pgfqpoint{1.444358in}{1.023043in}}%
\pgfpathlineto{\pgfqpoint{1.454769in}{1.023043in}}%
\pgfpathlineto{\pgfqpoint{1.454769in}{1.013478in}}%
\pgfpathquadraticcurveto{\pgfqpoint{1.458048in}{1.019098in}}{\pgfqpoint{1.463019in}{1.021818in}}%
\pgfpathquadraticcurveto{\pgfqpoint{1.468008in}{1.024556in}}{\pgfqpoint{1.474925in}{1.024556in}}%
\pgfpathquadraticcurveto{\pgfqpoint{1.486417in}{1.024556in}}{\pgfqpoint{1.493586in}{1.015441in}}%
\pgfpathquadraticcurveto{\pgfqpoint{1.500772in}{1.006327in}}{\pgfqpoint{1.500772in}{0.991467in}}%
\pgfpathquadraticcurveto{\pgfqpoint{1.500772in}{0.976607in}}{\pgfqpoint{1.493586in}{0.967475in}}%
\pgfpathquadraticcurveto{\pgfqpoint{1.486417in}{0.958361in}}{\pgfqpoint{1.474925in}{0.958361in}}%
\pgfpathquadraticcurveto{\pgfqpoint{1.468008in}{0.958361in}}{\pgfqpoint{1.463019in}{0.961099in}}%
\pgfpathquadraticcurveto{\pgfqpoint{1.458048in}{0.963837in}}{\pgfqpoint{1.454769in}{0.969456in}}%
\pgfpathlineto{\pgfqpoint{1.454769in}{0.969456in}}%
\pgfpathclose%
\pgfpathmoveto{\pgfqpoint{1.490019in}{0.991467in}}%
\pgfpathquadraticcurveto{\pgfqpoint{1.490019in}{1.002887in}}{\pgfqpoint{1.485318in}{1.009389in}}%
\pgfpathquadraticcurveto{\pgfqpoint{1.480617in}{1.015892in}}{\pgfqpoint{1.472403in}{1.015892in}}%
\pgfpathquadraticcurveto{\pgfqpoint{1.464172in}{1.015892in}}{\pgfqpoint{1.459471in}{1.009389in}}%
\pgfpathquadraticcurveto{\pgfqpoint{1.454769in}{1.002887in}}{\pgfqpoint{1.454769in}{0.991467in}}%
\pgfpathquadraticcurveto{\pgfqpoint{1.454769in}{0.980048in}}{\pgfqpoint{1.459471in}{0.973545in}}%
\pgfpathquadraticcurveto{\pgfqpoint{1.464172in}{0.967043in}}{\pgfqpoint{1.472403in}{0.967043in}}%
\pgfpathquadraticcurveto{\pgfqpoint{1.480617in}{0.967043in}}{\pgfqpoint{1.485318in}{0.973545in}}%
\pgfpathquadraticcurveto{\pgfqpoint{1.490019in}{0.980048in}}{\pgfqpoint{1.490019in}{0.991467in}}%
\pgfpathlineto{\pgfqpoint{1.490019in}{0.991467in}}%
\pgfpathclose%
\pgfpathmoveto{\pgfqpoint{1.542363in}{1.015784in}}%
\pgfpathquadraticcurveto{\pgfqpoint{1.534041in}{1.015784in}}{\pgfqpoint{1.529196in}{1.009281in}}%
\pgfpathquadraticcurveto{\pgfqpoint{1.524350in}{1.002779in}}{\pgfqpoint{1.524350in}{0.991467in}}%
\pgfpathquadraticcurveto{\pgfqpoint{1.524350in}{0.980156in}}{\pgfqpoint{1.529160in}{0.973653in}}%
\pgfpathquadraticcurveto{\pgfqpoint{1.533987in}{0.967151in}}{\pgfqpoint{1.542363in}{0.967151in}}%
\pgfpathquadraticcurveto{\pgfqpoint{1.550648in}{0.967151in}}{\pgfqpoint{1.555475in}{0.973671in}}%
\pgfpathquadraticcurveto{\pgfqpoint{1.560321in}{0.980210in}}{\pgfqpoint{1.560321in}{0.991467in}}%
\pgfpathquadraticcurveto{\pgfqpoint{1.560321in}{1.002671in}}{\pgfqpoint{1.555475in}{1.009227in}}%
\pgfpathquadraticcurveto{\pgfqpoint{1.550648in}{1.015784in}}{\pgfqpoint{1.542363in}{1.015784in}}%
\pgfpathlineto{\pgfqpoint{1.542363in}{1.015784in}}%
\pgfpathclose%
\pgfpathmoveto{\pgfqpoint{1.542363in}{1.024556in}}%
\pgfpathquadraticcurveto{\pgfqpoint{1.555872in}{1.024556in}}{\pgfqpoint{1.563581in}{1.015766in}}%
\pgfpathquadraticcurveto{\pgfqpoint{1.571308in}{1.006994in}}{\pgfqpoint{1.571308in}{0.991467in}}%
\pgfpathquadraticcurveto{\pgfqpoint{1.571308in}{0.975995in}}{\pgfqpoint{1.563581in}{0.967169in}}%
\pgfpathquadraticcurveto{\pgfqpoint{1.555872in}{0.958361in}}{\pgfqpoint{1.542363in}{0.958361in}}%
\pgfpathquadraticcurveto{\pgfqpoint{1.528799in}{0.958361in}}{\pgfqpoint{1.521108in}{0.967169in}}%
\pgfpathquadraticcurveto{\pgfqpoint{1.513435in}{0.975995in}}{\pgfqpoint{1.513435in}{0.991467in}}%
\pgfpathquadraticcurveto{\pgfqpoint{1.513435in}{1.006994in}}{\pgfqpoint{1.521108in}{1.015766in}}%
\pgfpathquadraticcurveto{\pgfqpoint{1.528799in}{1.024556in}}{\pgfqpoint{1.542363in}{1.024556in}}%
\pgfpathlineto{\pgfqpoint{1.542363in}{1.024556in}}%
\pgfpathclose%
\pgfpathmoveto{\pgfqpoint{1.588474in}{1.023043in}}%
\pgfpathlineto{\pgfqpoint{1.598831in}{1.023043in}}%
\pgfpathlineto{\pgfqpoint{1.598831in}{0.960000in}}%
\pgfpathlineto{\pgfqpoint{1.588474in}{0.960000in}}%
\pgfpathlineto{\pgfqpoint{1.588474in}{1.023043in}}%
\pgfpathlineto{\pgfqpoint{1.588474in}{1.023043in}}%
\pgfpathclose%
\pgfpathmoveto{\pgfqpoint{1.588474in}{1.047593in}}%
\pgfpathlineto{\pgfqpoint{1.598831in}{1.047593in}}%
\pgfpathlineto{\pgfqpoint{1.598831in}{1.034462in}}%
\pgfpathlineto{\pgfqpoint{1.588474in}{1.034462in}}%
\pgfpathlineto{\pgfqpoint{1.588474in}{1.047593in}}%
\pgfpathlineto{\pgfqpoint{1.588474in}{1.047593in}}%
\pgfpathclose%
\pgfpathmoveto{\pgfqpoint{1.672915in}{0.998060in}}%
\pgfpathlineto{\pgfqpoint{1.672915in}{0.960000in}}%
\pgfpathlineto{\pgfqpoint{1.662558in}{0.960000in}}%
\pgfpathlineto{\pgfqpoint{1.662558in}{0.997717in}}%
\pgfpathquadraticcurveto{\pgfqpoint{1.662558in}{1.006669in}}{\pgfqpoint{1.659063in}{1.011100in}}%
\pgfpathquadraticcurveto{\pgfqpoint{1.655569in}{1.015549in}}{\pgfqpoint{1.648598in}{1.015549in}}%
\pgfpathquadraticcurveto{\pgfqpoint{1.640205in}{1.015549in}}{\pgfqpoint{1.635359in}{1.010200in}}%
\pgfpathquadraticcurveto{\pgfqpoint{1.630514in}{1.004868in}}{\pgfqpoint{1.630514in}{0.995628in}}%
\pgfpathlineto{\pgfqpoint{1.630514in}{0.960000in}}%
\pgfpathlineto{\pgfqpoint{1.620103in}{0.960000in}}%
\pgfpathlineto{\pgfqpoint{1.620103in}{1.023043in}}%
\pgfpathlineto{\pgfqpoint{1.630514in}{1.023043in}}%
\pgfpathlineto{\pgfqpoint{1.630514in}{1.013244in}}%
\pgfpathquadraticcurveto{\pgfqpoint{1.634242in}{1.018936in}}{\pgfqpoint{1.639268in}{1.021746in}}%
\pgfpathquadraticcurveto{\pgfqpoint{1.644311in}{1.024556in}}{\pgfqpoint{1.650904in}{1.024556in}}%
\pgfpathquadraticcurveto{\pgfqpoint{1.661765in}{1.024556in}}{\pgfqpoint{1.667331in}{1.017837in}}%
\pgfpathquadraticcurveto{\pgfqpoint{1.672915in}{1.011118in}}{\pgfqpoint{1.672915in}{0.998060in}}%
\pgfpathlineto{\pgfqpoint{1.672915in}{0.998060in}}%
\pgfpathclose%
\pgfpathmoveto{\pgfqpoint{1.703805in}{1.040947in}}%
\pgfpathlineto{\pgfqpoint{1.703805in}{1.023043in}}%
\pgfpathlineto{\pgfqpoint{1.725132in}{1.023043in}}%
\pgfpathlineto{\pgfqpoint{1.725132in}{1.014991in}}%
\pgfpathlineto{\pgfqpoint{1.703805in}{1.014991in}}%
\pgfpathlineto{\pgfqpoint{1.703805in}{0.980768in}}%
\pgfpathquadraticcurveto{\pgfqpoint{1.703805in}{0.973059in}}{\pgfqpoint{1.705913in}{0.970861in}}%
\pgfpathquadraticcurveto{\pgfqpoint{1.708020in}{0.968664in}}{\pgfqpoint{1.714505in}{0.968664in}}%
\pgfpathlineto{\pgfqpoint{1.725132in}{0.968664in}}%
\pgfpathlineto{\pgfqpoint{1.725132in}{0.960000in}}%
\pgfpathlineto{\pgfqpoint{1.714505in}{0.960000in}}%
\pgfpathquadraticcurveto{\pgfqpoint{1.702509in}{0.960000in}}{\pgfqpoint{1.697951in}{0.964467in}}%
\pgfpathquadraticcurveto{\pgfqpoint{1.693394in}{0.968952in}}{\pgfqpoint{1.693394in}{0.980768in}}%
\pgfpathlineto{\pgfqpoint{1.693394in}{1.014991in}}%
\pgfpathlineto{\pgfqpoint{1.685793in}{1.014991in}}%
\pgfpathlineto{\pgfqpoint{1.685793in}{1.023043in}}%
\pgfpathlineto{\pgfqpoint{1.693394in}{1.023043in}}%
\pgfpathlineto{\pgfqpoint{1.693394in}{1.040947in}}%
\pgfpathlineto{\pgfqpoint{1.703805in}{1.040947in}}%
\pgfpathlineto{\pgfqpoint{1.703805in}{1.040947in}}%
\pgfpathclose%
\pgfpathmoveto{\pgfqpoint{1.778934in}{1.021187in}}%
\pgfpathlineto{\pgfqpoint{1.778934in}{1.011389in}}%
\pgfpathquadraticcurveto{\pgfqpoint{1.774557in}{1.013640in}}{\pgfqpoint{1.769820in}{1.014757in}}%
\pgfpathquadraticcurveto{\pgfqpoint{1.765101in}{1.015892in}}{\pgfqpoint{1.760021in}{1.015892in}}%
\pgfpathquadraticcurveto{\pgfqpoint{1.752312in}{1.015892in}}{\pgfqpoint{1.748458in}{1.013532in}}%
\pgfpathquadraticcurveto{\pgfqpoint{1.744603in}{1.011173in}}{\pgfqpoint{1.744603in}{1.006435in}}%
\pgfpathquadraticcurveto{\pgfqpoint{1.744603in}{1.002833in}}{\pgfqpoint{1.747359in}{1.000780in}}%
\pgfpathquadraticcurveto{\pgfqpoint{1.750115in}{0.998726in}}{\pgfqpoint{1.758454in}{0.996871in}}%
\pgfpathlineto{\pgfqpoint{1.762003in}{0.996078in}}%
\pgfpathquadraticcurveto{\pgfqpoint{1.773026in}{0.993719in}}{\pgfqpoint{1.777673in}{0.989414in}}%
\pgfpathquadraticcurveto{\pgfqpoint{1.782320in}{0.985109in}}{\pgfqpoint{1.782320in}{0.977400in}}%
\pgfpathquadraticcurveto{\pgfqpoint{1.782320in}{0.968610in}}{\pgfqpoint{1.775368in}{0.963476in}}%
\pgfpathquadraticcurveto{\pgfqpoint{1.768415in}{0.958361in}}{\pgfqpoint{1.756257in}{0.958361in}}%
\pgfpathquadraticcurveto{\pgfqpoint{1.751195in}{0.958361in}}{\pgfqpoint{1.745702in}{0.959352in}}%
\pgfpathquadraticcurveto{\pgfqpoint{1.740208in}{0.960342in}}{\pgfqpoint{1.734138in}{0.962306in}}%
\pgfpathlineto{\pgfqpoint{1.734138in}{0.973005in}}%
\pgfpathquadraticcurveto{\pgfqpoint{1.739884in}{0.970015in}}{\pgfqpoint{1.745450in}{0.968520in}}%
\pgfpathquadraticcurveto{\pgfqpoint{1.751015in}{0.967043in}}{\pgfqpoint{1.756491in}{0.967043in}}%
\pgfpathquadraticcurveto{\pgfqpoint{1.763804in}{0.967043in}}{\pgfqpoint{1.767731in}{0.969546in}}%
\pgfpathquadraticcurveto{\pgfqpoint{1.771675in}{0.972050in}}{\pgfqpoint{1.771675in}{0.976607in}}%
\pgfpathquadraticcurveto{\pgfqpoint{1.771675in}{0.980822in}}{\pgfqpoint{1.768829in}{0.983074in}}%
\pgfpathquadraticcurveto{\pgfqpoint{1.766001in}{0.985325in}}{\pgfqpoint{1.756365in}{0.987414in}}%
\pgfpathlineto{\pgfqpoint{1.752762in}{0.988261in}}%
\pgfpathquadraticcurveto{\pgfqpoint{1.743144in}{0.990278in}}{\pgfqpoint{1.738857in}{0.994475in}}%
\pgfpathquadraticcurveto{\pgfqpoint{1.734588in}{0.998672in}}{\pgfqpoint{1.734588in}{1.005985in}}%
\pgfpathquadraticcurveto{\pgfqpoint{1.734588in}{1.014883in}}{\pgfqpoint{1.740892in}{1.019710in}}%
\pgfpathquadraticcurveto{\pgfqpoint{1.747197in}{1.024556in}}{\pgfqpoint{1.758797in}{1.024556in}}%
\pgfpathquadraticcurveto{\pgfqpoint{1.764524in}{1.024556in}}{\pgfqpoint{1.769586in}{1.023709in}}%
\pgfpathquadraticcurveto{\pgfqpoint{1.774665in}{1.022880in}}{\pgfqpoint{1.778934in}{1.021187in}}%
\pgfpathlineto{\pgfqpoint{1.778934in}{1.021187in}}%
\pgfpathclose%
\pgfpathmoveto{\pgfqpoint{1.800261in}{0.974302in}}%
\pgfpathlineto{\pgfqpoint{1.812149in}{0.974302in}}%
\pgfpathlineto{\pgfqpoint{1.812149in}{0.960000in}}%
\pgfpathlineto{\pgfqpoint{1.800261in}{0.960000in}}%
\pgfpathlineto{\pgfqpoint{1.800261in}{0.974302in}}%
\pgfpathlineto{\pgfqpoint{1.800261in}{0.974302in}}%
\pgfpathclose%
\pgfusepath{fill}%
\end{pgfscope}%
\begin{pgfscope}%
\pgfsetbuttcap%
\pgfsetroundjoin%
\definecolor{currentfill}{rgb}{0.309804,0.372549,0.419608}%
\pgfsetfillcolor{currentfill}%
\pgfsetlinewidth{0.000000pt}%
\definecolor{currentstroke}{rgb}{0.309804,0.372549,0.419608}%
\pgfsetstrokecolor{currentstroke}%
\pgfsetdash{}{0pt}%
\pgfpathmoveto{\pgfqpoint{1.955478in}{1.041289in}}%
\pgfpathlineto{\pgfqpoint{1.955478in}{1.030193in}}%
\pgfpathquadraticcurveto{\pgfqpoint{1.949012in}{1.033291in}}{\pgfqpoint{1.943266in}{1.034804in}}%
\pgfpathquadraticcurveto{\pgfqpoint{1.937520in}{1.036336in}}{\pgfqpoint{1.932171in}{1.036336in}}%
\pgfpathquadraticcurveto{\pgfqpoint{1.922894in}{1.036336in}}{\pgfqpoint{1.917851in}{1.032733in}}%
\pgfpathquadraticcurveto{\pgfqpoint{1.912807in}{1.029131in}}{\pgfqpoint{1.912807in}{1.022484in}}%
\pgfpathquadraticcurveto{\pgfqpoint{1.912807in}{1.016900in}}{\pgfqpoint{1.916158in}{1.014054in}}%
\pgfpathquadraticcurveto{\pgfqpoint{1.919508in}{1.011227in}}{\pgfqpoint{1.928856in}{1.009479in}}%
\pgfpathlineto{\pgfqpoint{1.935719in}{1.008074in}}%
\pgfpathquadraticcurveto{\pgfqpoint{1.948436in}{1.005643in}}{\pgfqpoint{1.954488in}{0.999537in}}%
\pgfpathquadraticcurveto{\pgfqpoint{1.960540in}{0.993431in}}{\pgfqpoint{1.960540in}{0.983200in}}%
\pgfpathquadraticcurveto{\pgfqpoint{1.960540in}{0.970969in}}{\pgfqpoint{1.952344in}{0.964665in}}%
\pgfpathquadraticcurveto{\pgfqpoint{1.944167in}{0.958361in}}{\pgfqpoint{1.928352in}{0.958361in}}%
\pgfpathquadraticcurveto{\pgfqpoint{1.922390in}{0.958361in}}{\pgfqpoint{1.915653in}{0.959712in}}%
\pgfpathquadraticcurveto{\pgfqpoint{1.908935in}{0.961063in}}{\pgfqpoint{1.901730in}{0.963711in}}%
\pgfpathlineto{\pgfqpoint{1.901730in}{0.975418in}}%
\pgfpathquadraticcurveto{\pgfqpoint{1.908647in}{0.971546in}}{\pgfqpoint{1.915293in}{0.969564in}}%
\pgfpathquadraticcurveto{\pgfqpoint{1.921940in}{0.967601in}}{\pgfqpoint{1.928352in}{0.967601in}}%
\pgfpathquadraticcurveto{\pgfqpoint{1.938079in}{0.967601in}}{\pgfqpoint{1.943374in}{0.971420in}}%
\pgfpathquadraticcurveto{\pgfqpoint{1.948670in}{0.975256in}}{\pgfqpoint{1.948670in}{0.982353in}}%
\pgfpathquadraticcurveto{\pgfqpoint{1.948670in}{0.988531in}}{\pgfqpoint{1.944869in}{0.992026in}}%
\pgfpathquadraticcurveto{\pgfqpoint{1.941069in}{0.995520in}}{\pgfqpoint{1.932405in}{0.997267in}}%
\pgfpathlineto{\pgfqpoint{1.925470in}{0.998618in}}%
\pgfpathquadraticcurveto{\pgfqpoint{1.912753in}{1.001140in}}{\pgfqpoint{1.907062in}{1.006543in}}%
\pgfpathquadraticcurveto{\pgfqpoint{1.901388in}{1.011947in}}{\pgfqpoint{1.901388in}{1.021584in}}%
\pgfpathquadraticcurveto{\pgfqpoint{1.901388in}{1.032733in}}{\pgfqpoint{1.909241in}{1.039145in}}%
\pgfpathquadraticcurveto{\pgfqpoint{1.917094in}{1.045558in}}{\pgfqpoint{1.930874in}{1.045558in}}%
\pgfpathquadraticcurveto{\pgfqpoint{1.936800in}{1.045558in}}{\pgfqpoint{1.942924in}{1.044477in}}%
\pgfpathquadraticcurveto{\pgfqpoint{1.949066in}{1.043414in}}{\pgfqpoint{1.955478in}{1.041289in}}%
\pgfpathlineto{\pgfqpoint{1.955478in}{1.041289in}}%
\pgfpathclose%
\pgfpathmoveto{\pgfqpoint{1.977831in}{1.023043in}}%
\pgfpathlineto{\pgfqpoint{1.988188in}{1.023043in}}%
\pgfpathlineto{\pgfqpoint{1.988188in}{0.960000in}}%
\pgfpathlineto{\pgfqpoint{1.977831in}{0.960000in}}%
\pgfpathlineto{\pgfqpoint{1.977831in}{1.023043in}}%
\pgfpathlineto{\pgfqpoint{1.977831in}{1.023043in}}%
\pgfpathclose%
\pgfpathmoveto{\pgfqpoint{1.977831in}{1.047593in}}%
\pgfpathlineto{\pgfqpoint{1.988188in}{1.047593in}}%
\pgfpathlineto{\pgfqpoint{1.988188in}{1.034462in}}%
\pgfpathlineto{\pgfqpoint{1.977831in}{1.034462in}}%
\pgfpathlineto{\pgfqpoint{1.977831in}{1.047593in}}%
\pgfpathlineto{\pgfqpoint{1.977831in}{1.047593in}}%
\pgfpathclose%
\pgfpathmoveto{\pgfqpoint{2.051339in}{0.992260in}}%
\pgfpathquadraticcurveto{\pgfqpoint{2.051339in}{1.003517in}}{\pgfqpoint{2.046692in}{1.009696in}}%
\pgfpathquadraticcurveto{\pgfqpoint{2.042063in}{1.015892in}}{\pgfqpoint{2.033669in}{1.015892in}}%
\pgfpathquadraticcurveto{\pgfqpoint{2.025347in}{1.015892in}}{\pgfqpoint{2.020700in}{1.009696in}}%
\pgfpathquadraticcurveto{\pgfqpoint{2.016053in}{1.003517in}}{\pgfqpoint{2.016053in}{0.992260in}}%
\pgfpathquadraticcurveto{\pgfqpoint{2.016053in}{0.981056in}}{\pgfqpoint{2.020700in}{0.974860in}}%
\pgfpathquadraticcurveto{\pgfqpoint{2.025347in}{0.968664in}}{\pgfqpoint{2.033669in}{0.968664in}}%
\pgfpathquadraticcurveto{\pgfqpoint{2.042063in}{0.968664in}}{\pgfqpoint{2.046692in}{0.974860in}}%
\pgfpathquadraticcurveto{\pgfqpoint{2.051339in}{0.981056in}}{\pgfqpoint{2.051339in}{0.992260in}}%
\pgfpathlineto{\pgfqpoint{2.051339in}{0.992260in}}%
\pgfpathclose%
\pgfpathmoveto{\pgfqpoint{2.061696in}{0.967817in}}%
\pgfpathquadraticcurveto{\pgfqpoint{2.061696in}{0.951732in}}{\pgfqpoint{2.054545in}{0.943879in}}%
\pgfpathquadraticcurveto{\pgfqpoint{2.047412in}{0.936026in}}{\pgfqpoint{2.032660in}{0.936026in}}%
\pgfpathquadraticcurveto{\pgfqpoint{2.027203in}{0.936026in}}{\pgfqpoint{2.022357in}{0.936836in}}%
\pgfpathquadraticcurveto{\pgfqpoint{2.017512in}{0.937647in}}{\pgfqpoint{2.012955in}{0.939340in}}%
\pgfpathlineto{\pgfqpoint{2.012955in}{0.949409in}}%
\pgfpathquadraticcurveto{\pgfqpoint{2.017512in}{0.946941in}}{\pgfqpoint{2.021961in}{0.945770in}}%
\pgfpathquadraticcurveto{\pgfqpoint{2.026410in}{0.944582in}}{\pgfqpoint{2.031021in}{0.944582in}}%
\pgfpathquadraticcurveto{\pgfqpoint{2.041216in}{0.944582in}}{\pgfqpoint{2.046278in}{0.949895in}}%
\pgfpathquadraticcurveto{\pgfqpoint{2.051339in}{0.955209in}}{\pgfqpoint{2.051339in}{0.965962in}}%
\pgfpathlineto{\pgfqpoint{2.051339in}{0.971095in}}%
\pgfpathquadraticcurveto{\pgfqpoint{2.048133in}{0.965512in}}{\pgfqpoint{2.043125in}{0.962756in}}%
\pgfpathquadraticcurveto{\pgfqpoint{2.038118in}{0.960000in}}{\pgfqpoint{2.031129in}{0.960000in}}%
\pgfpathquadraticcurveto{\pgfqpoint{2.019547in}{0.960000in}}{\pgfqpoint{2.012451in}{0.968826in}}%
\pgfpathquadraticcurveto{\pgfqpoint{2.005354in}{0.977670in}}{\pgfqpoint{2.005354in}{0.992260in}}%
\pgfpathquadraticcurveto{\pgfqpoint{2.005354in}{1.006886in}}{\pgfqpoint{2.012451in}{1.015712in}}%
\pgfpathquadraticcurveto{\pgfqpoint{2.019547in}{1.024556in}}{\pgfqpoint{2.031129in}{1.024556in}}%
\pgfpathquadraticcurveto{\pgfqpoint{2.038118in}{1.024556in}}{\pgfqpoint{2.043125in}{1.021800in}}%
\pgfpathquadraticcurveto{\pgfqpoint{2.048133in}{1.019044in}}{\pgfqpoint{2.051339in}{1.013478in}}%
\pgfpathlineto{\pgfqpoint{2.051339in}{1.023043in}}%
\pgfpathlineto{\pgfqpoint{2.061696in}{1.023043in}}%
\pgfpathlineto{\pgfqpoint{2.061696in}{0.967817in}}%
\pgfpathlineto{\pgfqpoint{2.061696in}{0.967817in}}%
\pgfpathclose%
\pgfpathmoveto{\pgfqpoint{2.135456in}{0.998060in}}%
\pgfpathlineto{\pgfqpoint{2.135456in}{0.960000in}}%
\pgfpathlineto{\pgfqpoint{2.125099in}{0.960000in}}%
\pgfpathlineto{\pgfqpoint{2.125099in}{0.997717in}}%
\pgfpathquadraticcurveto{\pgfqpoint{2.125099in}{1.006669in}}{\pgfqpoint{2.121604in}{1.011100in}}%
\pgfpathquadraticcurveto{\pgfqpoint{2.118110in}{1.015549in}}{\pgfqpoint{2.111139in}{1.015549in}}%
\pgfpathquadraticcurveto{\pgfqpoint{2.102746in}{1.015549in}}{\pgfqpoint{2.097900in}{1.010200in}}%
\pgfpathquadraticcurveto{\pgfqpoint{2.093055in}{1.004868in}}{\pgfqpoint{2.093055in}{0.995628in}}%
\pgfpathlineto{\pgfqpoint{2.093055in}{0.960000in}}%
\pgfpathlineto{\pgfqpoint{2.082644in}{0.960000in}}%
\pgfpathlineto{\pgfqpoint{2.082644in}{1.023043in}}%
\pgfpathlineto{\pgfqpoint{2.093055in}{1.023043in}}%
\pgfpathlineto{\pgfqpoint{2.093055in}{1.013244in}}%
\pgfpathquadraticcurveto{\pgfqpoint{2.096784in}{1.018936in}}{\pgfqpoint{2.101809in}{1.021746in}}%
\pgfpathquadraticcurveto{\pgfqpoint{2.106852in}{1.024556in}}{\pgfqpoint{2.113445in}{1.024556in}}%
\pgfpathquadraticcurveto{\pgfqpoint{2.124306in}{1.024556in}}{\pgfqpoint{2.129872in}{1.017837in}}%
\pgfpathquadraticcurveto{\pgfqpoint{2.135456in}{1.011118in}}{\pgfqpoint{2.135456in}{0.998060in}}%
\pgfpathlineto{\pgfqpoint{2.135456in}{0.998060in}}%
\pgfpathclose%
\pgfpathmoveto{\pgfqpoint{2.156098in}{1.023043in}}%
\pgfpathlineto{\pgfqpoint{2.166455in}{1.023043in}}%
\pgfpathlineto{\pgfqpoint{2.166455in}{0.960000in}}%
\pgfpathlineto{\pgfqpoint{2.156098in}{0.960000in}}%
\pgfpathlineto{\pgfqpoint{2.156098in}{1.023043in}}%
\pgfpathlineto{\pgfqpoint{2.156098in}{1.023043in}}%
\pgfpathclose%
\pgfpathmoveto{\pgfqpoint{2.156098in}{1.047593in}}%
\pgfpathlineto{\pgfqpoint{2.166455in}{1.047593in}}%
\pgfpathlineto{\pgfqpoint{2.166455in}{1.034462in}}%
\pgfpathlineto{\pgfqpoint{2.156098in}{1.034462in}}%
\pgfpathlineto{\pgfqpoint{2.156098in}{1.047593in}}%
\pgfpathlineto{\pgfqpoint{2.156098in}{1.047593in}}%
\pgfpathclose%
\pgfpathmoveto{\pgfqpoint{2.239062in}{1.023043in}}%
\pgfpathlineto{\pgfqpoint{2.239062in}{0.960000in}}%
\pgfpathlineto{\pgfqpoint{2.228651in}{0.960000in}}%
\pgfpathlineto{\pgfqpoint{2.228651in}{1.014991in}}%
\pgfpathlineto{\pgfqpoint{2.200227in}{1.014991in}}%
\pgfpathlineto{\pgfqpoint{2.200227in}{0.960000in}}%
\pgfpathlineto{\pgfqpoint{2.189816in}{0.960000in}}%
\pgfpathlineto{\pgfqpoint{2.189816in}{1.014991in}}%
\pgfpathlineto{\pgfqpoint{2.179910in}{1.014991in}}%
\pgfpathlineto{\pgfqpoint{2.179910in}{1.023043in}}%
\pgfpathlineto{\pgfqpoint{2.189816in}{1.023043in}}%
\pgfpathlineto{\pgfqpoint{2.189816in}{1.027438in}}%
\pgfpathquadraticcurveto{\pgfqpoint{2.189816in}{1.037740in}}{\pgfqpoint{2.194680in}{1.042658in}}%
\pgfpathquadraticcurveto{\pgfqpoint{2.199561in}{1.047593in}}{\pgfqpoint{2.209630in}{1.047593in}}%
\pgfpathlineto{\pgfqpoint{2.220041in}{1.047593in}}%
\pgfpathlineto{\pgfqpoint{2.220041in}{1.038965in}}%
\pgfpathlineto{\pgfqpoint{2.210134in}{1.038965in}}%
\pgfpathquadraticcurveto{\pgfqpoint{2.204568in}{1.038965in}}{\pgfqpoint{2.202389in}{1.036714in}}%
\pgfpathquadraticcurveto{\pgfqpoint{2.200227in}{1.034462in}}{\pgfqpoint{2.200227in}{1.028608in}}%
\pgfpathlineto{\pgfqpoint{2.200227in}{1.023043in}}%
\pgfpathlineto{\pgfqpoint{2.239062in}{1.023043in}}%
\pgfpathlineto{\pgfqpoint{2.239062in}{1.023043in}}%
\pgfpathclose%
\pgfpathmoveto{\pgfqpoint{2.228651in}{1.047467in}}%
\pgfpathlineto{\pgfqpoint{2.239062in}{1.047467in}}%
\pgfpathlineto{\pgfqpoint{2.239062in}{1.034354in}}%
\pgfpathlineto{\pgfqpoint{2.228651in}{1.034354in}}%
\pgfpathlineto{\pgfqpoint{2.228651in}{1.047467in}}%
\pgfpathlineto{\pgfqpoint{2.228651in}{1.047467in}}%
\pgfpathclose%
\pgfpathmoveto{\pgfqpoint{2.306103in}{1.020629in}}%
\pgfpathlineto{\pgfqpoint{2.306103in}{1.010938in}}%
\pgfpathquadraticcurveto{\pgfqpoint{2.301708in}{1.013370in}}{\pgfqpoint{2.297295in}{1.014577in}}%
\pgfpathquadraticcurveto{\pgfqpoint{2.292882in}{1.015784in}}{\pgfqpoint{2.288379in}{1.015784in}}%
\pgfpathquadraticcurveto{\pgfqpoint{2.278292in}{1.015784in}}{\pgfqpoint{2.272708in}{1.009389in}}%
\pgfpathquadraticcurveto{\pgfqpoint{2.267143in}{1.003013in}}{\pgfqpoint{2.267143in}{0.991467in}}%
\pgfpathquadraticcurveto{\pgfqpoint{2.267143in}{0.979921in}}{\pgfqpoint{2.272708in}{0.973527in}}%
\pgfpathquadraticcurveto{\pgfqpoint{2.278292in}{0.967151in}}{\pgfqpoint{2.288379in}{0.967151in}}%
\pgfpathquadraticcurveto{\pgfqpoint{2.292882in}{0.967151in}}{\pgfqpoint{2.297295in}{0.968358in}}%
\pgfpathquadraticcurveto{\pgfqpoint{2.301708in}{0.969564in}}{\pgfqpoint{2.306103in}{0.971996in}}%
\pgfpathlineto{\pgfqpoint{2.306103in}{0.962414in}}%
\pgfpathquadraticcurveto{\pgfqpoint{2.301762in}{0.960396in}}{\pgfqpoint{2.297115in}{0.959388in}}%
\pgfpathquadraticcurveto{\pgfqpoint{2.292486in}{0.958361in}}{\pgfqpoint{2.287244in}{0.958361in}}%
\pgfpathquadraticcurveto{\pgfqpoint{2.272996in}{0.958361in}}{\pgfqpoint{2.264603in}{0.967313in}}%
\pgfpathquadraticcurveto{\pgfqpoint{2.256227in}{0.976265in}}{\pgfqpoint{2.256227in}{0.991467in}}%
\pgfpathquadraticcurveto{\pgfqpoint{2.256227in}{1.006886in}}{\pgfqpoint{2.264693in}{1.015712in}}%
\pgfpathquadraticcurveto{\pgfqpoint{2.273177in}{1.024556in}}{\pgfqpoint{2.287929in}{1.024556in}}%
\pgfpathquadraticcurveto{\pgfqpoint{2.292702in}{1.024556in}}{\pgfqpoint{2.297259in}{1.023565in}}%
\pgfpathquadraticcurveto{\pgfqpoint{2.301816in}{1.022592in}}{\pgfqpoint{2.306103in}{1.020629in}}%
\pgfpathlineto{\pgfqpoint{2.306103in}{1.020629in}}%
\pgfpathclose%
\pgfpathmoveto{\pgfqpoint{2.352772in}{0.991683in}}%
\pgfpathquadraticcurveto{\pgfqpoint{2.340218in}{0.991683in}}{\pgfqpoint{2.335373in}{0.988819in}}%
\pgfpathquadraticcurveto{\pgfqpoint{2.330527in}{0.985956in}}{\pgfqpoint{2.330527in}{0.979021in}}%
\pgfpathquadraticcurveto{\pgfqpoint{2.330527in}{0.973509in}}{\pgfqpoint{2.334166in}{0.970267in}}%
\pgfpathquadraticcurveto{\pgfqpoint{2.337804in}{0.967043in}}{\pgfqpoint{2.344036in}{0.967043in}}%
\pgfpathquadraticcurveto{\pgfqpoint{2.352664in}{0.967043in}}{\pgfqpoint{2.357870in}{0.973149in}}%
\pgfpathquadraticcurveto{\pgfqpoint{2.363075in}{0.979255in}}{\pgfqpoint{2.363075in}{0.989378in}}%
\pgfpathlineto{\pgfqpoint{2.363075in}{0.991683in}}%
\pgfpathlineto{\pgfqpoint{2.352772in}{0.991683in}}%
\pgfpathlineto{\pgfqpoint{2.352772in}{0.991683in}}%
\pgfpathclose%
\pgfpathmoveto{\pgfqpoint{2.373432in}{0.995970in}}%
\pgfpathlineto{\pgfqpoint{2.373432in}{0.960000in}}%
\pgfpathlineto{\pgfqpoint{2.363075in}{0.960000in}}%
\pgfpathlineto{\pgfqpoint{2.363075in}{0.969564in}}%
\pgfpathquadraticcurveto{\pgfqpoint{2.359527in}{0.963837in}}{\pgfqpoint{2.354231in}{0.961099in}}%
\pgfpathquadraticcurveto{\pgfqpoint{2.348936in}{0.958361in}}{\pgfqpoint{2.341281in}{0.958361in}}%
\pgfpathquadraticcurveto{\pgfqpoint{2.331608in}{0.958361in}}{\pgfqpoint{2.325880in}{0.963801in}}%
\pgfpathquadraticcurveto{\pgfqpoint{2.320170in}{0.969240in}}{\pgfqpoint{2.320170in}{0.978354in}}%
\pgfpathquadraticcurveto{\pgfqpoint{2.320170in}{0.988982in}}{\pgfqpoint{2.327285in}{0.994385in}}%
\pgfpathquadraticcurveto{\pgfqpoint{2.334418in}{0.999789in}}{\pgfqpoint{2.348539in}{0.999789in}}%
\pgfpathlineto{\pgfqpoint{2.363075in}{0.999789in}}%
\pgfpathlineto{\pgfqpoint{2.363075in}{1.000816in}}%
\pgfpathquadraticcurveto{\pgfqpoint{2.363075in}{1.007966in}}{\pgfqpoint{2.358374in}{1.011875in}}%
\pgfpathquadraticcurveto{\pgfqpoint{2.353673in}{1.015784in}}{\pgfqpoint{2.345171in}{1.015784in}}%
\pgfpathquadraticcurveto{\pgfqpoint{2.339768in}{1.015784in}}{\pgfqpoint{2.334634in}{1.014487in}}%
\pgfpathquadraticcurveto{\pgfqpoint{2.329519in}{1.013190in}}{\pgfqpoint{2.324799in}{1.010596in}}%
\pgfpathlineto{\pgfqpoint{2.324799in}{1.020179in}}%
\pgfpathquadraticcurveto{\pgfqpoint{2.330473in}{1.022376in}}{\pgfqpoint{2.335823in}{1.023457in}}%
\pgfpathquadraticcurveto{\pgfqpoint{2.341172in}{1.024556in}}{\pgfqpoint{2.346234in}{1.024556in}}%
\pgfpathquadraticcurveto{\pgfqpoint{2.359923in}{1.024556in}}{\pgfqpoint{2.366678in}{1.017459in}}%
\pgfpathquadraticcurveto{\pgfqpoint{2.373432in}{1.010380in}}{\pgfqpoint{2.373432in}{0.995970in}}%
\pgfpathlineto{\pgfqpoint{2.373432in}{0.995970in}}%
\pgfpathclose%
\pgfpathmoveto{\pgfqpoint{2.447174in}{0.998060in}}%
\pgfpathlineto{\pgfqpoint{2.447174in}{0.960000in}}%
\pgfpathlineto{\pgfqpoint{2.436817in}{0.960000in}}%
\pgfpathlineto{\pgfqpoint{2.436817in}{0.997717in}}%
\pgfpathquadraticcurveto{\pgfqpoint{2.436817in}{1.006669in}}{\pgfqpoint{2.433323in}{1.011100in}}%
\pgfpathquadraticcurveto{\pgfqpoint{2.429828in}{1.015549in}}{\pgfqpoint{2.422858in}{1.015549in}}%
\pgfpathquadraticcurveto{\pgfqpoint{2.414464in}{1.015549in}}{\pgfqpoint{2.409619in}{1.010200in}}%
\pgfpathquadraticcurveto{\pgfqpoint{2.404773in}{1.004868in}}{\pgfqpoint{2.404773in}{0.995628in}}%
\pgfpathlineto{\pgfqpoint{2.404773in}{0.960000in}}%
\pgfpathlineto{\pgfqpoint{2.394362in}{0.960000in}}%
\pgfpathlineto{\pgfqpoint{2.394362in}{1.023043in}}%
\pgfpathlineto{\pgfqpoint{2.404773in}{1.023043in}}%
\pgfpathlineto{\pgfqpoint{2.404773in}{1.013244in}}%
\pgfpathquadraticcurveto{\pgfqpoint{2.408502in}{1.018936in}}{\pgfqpoint{2.413527in}{1.021746in}}%
\pgfpathquadraticcurveto{\pgfqpoint{2.418571in}{1.024556in}}{\pgfqpoint{2.425163in}{1.024556in}}%
\pgfpathquadraticcurveto{\pgfqpoint{2.436024in}{1.024556in}}{\pgfqpoint{2.441590in}{1.017837in}}%
\pgfpathquadraticcurveto{\pgfqpoint{2.447174in}{1.011118in}}{\pgfqpoint{2.447174in}{0.998060in}}%
\pgfpathlineto{\pgfqpoint{2.447174in}{0.998060in}}%
\pgfpathclose%
\pgfpathmoveto{\pgfqpoint{2.478065in}{1.040947in}}%
\pgfpathlineto{\pgfqpoint{2.478065in}{1.023043in}}%
\pgfpathlineto{\pgfqpoint{2.499391in}{1.023043in}}%
\pgfpathlineto{\pgfqpoint{2.499391in}{1.014991in}}%
\pgfpathlineto{\pgfqpoint{2.478065in}{1.014991in}}%
\pgfpathlineto{\pgfqpoint{2.478065in}{0.980768in}}%
\pgfpathquadraticcurveto{\pgfqpoint{2.478065in}{0.973059in}}{\pgfqpoint{2.480172in}{0.970861in}}%
\pgfpathquadraticcurveto{\pgfqpoint{2.482280in}{0.968664in}}{\pgfqpoint{2.488764in}{0.968664in}}%
\pgfpathlineto{\pgfqpoint{2.499391in}{0.968664in}}%
\pgfpathlineto{\pgfqpoint{2.499391in}{0.960000in}}%
\pgfpathlineto{\pgfqpoint{2.488764in}{0.960000in}}%
\pgfpathquadraticcurveto{\pgfqpoint{2.476768in}{0.960000in}}{\pgfqpoint{2.472211in}{0.964467in}}%
\pgfpathquadraticcurveto{\pgfqpoint{2.467654in}{0.968952in}}{\pgfqpoint{2.467654in}{0.980768in}}%
\pgfpathlineto{\pgfqpoint{2.467654in}{1.014991in}}%
\pgfpathlineto{\pgfqpoint{2.460053in}{1.014991in}}%
\pgfpathlineto{\pgfqpoint{2.460053in}{1.023043in}}%
\pgfpathlineto{\pgfqpoint{2.467654in}{1.023043in}}%
\pgfpathlineto{\pgfqpoint{2.467654in}{1.040947in}}%
\pgfpathlineto{\pgfqpoint{2.478065in}{1.040947in}}%
\pgfpathlineto{\pgfqpoint{2.478065in}{1.040947in}}%
\pgfpathclose%
\pgfusepath{fill}%
\end{pgfscope}%
\begin{pgfscope}%
\pgfsetbuttcap%
\pgfsetroundjoin%
\definecolor{currentfill}{rgb}{0.309804,0.372549,0.419608}%
\pgfsetfillcolor{currentfill}%
\pgfsetlinewidth{0.000000pt}%
\definecolor{currentstroke}{rgb}{0.309804,0.372549,0.419608}%
\pgfsetstrokecolor{currentstroke}%
\pgfsetdash{}{0pt}%
\pgfpathmoveto{\pgfqpoint{2.646708in}{1.032283in}}%
\pgfpathlineto{\pgfqpoint{2.646708in}{1.000924in}}%
\pgfpathlineto{\pgfqpoint{2.678049in}{1.000924in}}%
\pgfpathlineto{\pgfqpoint{2.678049in}{0.991359in}}%
\pgfpathlineto{\pgfqpoint{2.646708in}{0.991359in}}%
\pgfpathlineto{\pgfqpoint{2.646708in}{0.960000in}}%
\pgfpathlineto{\pgfqpoint{2.637251in}{0.960000in}}%
\pgfpathlineto{\pgfqpoint{2.637251in}{0.991359in}}%
\pgfpathlineto{\pgfqpoint{2.605892in}{0.991359in}}%
\pgfpathlineto{\pgfqpoint{2.605892in}{1.000924in}}%
\pgfpathlineto{\pgfqpoint{2.637251in}{1.000924in}}%
\pgfpathlineto{\pgfqpoint{2.637251in}{1.032283in}}%
\pgfpathlineto{\pgfqpoint{2.646708in}{1.032283in}}%
\pgfpathlineto{\pgfqpoint{2.646708in}{1.032283in}}%
\pgfpathclose%
\pgfusepath{fill}%
\end{pgfscope}%
\begin{pgfscope}%
\pgfsetbuttcap%
\pgfsetroundjoin%
\definecolor{currentfill}{rgb}{0.309804,0.372549,0.419608}%
\pgfsetfillcolor{currentfill}%
\pgfsetlinewidth{0.000000pt}%
\definecolor{currentstroke}{rgb}{0.309804,0.372549,0.419608}%
\pgfsetstrokecolor{currentstroke}%
\pgfsetdash{}{0pt}%
\pgfpathmoveto{\pgfqpoint{2.794746in}{0.999411in}}%
\pgfpathquadraticcurveto{\pgfqpoint{2.798402in}{0.998168in}}{\pgfqpoint{2.801861in}{0.994115in}}%
\pgfpathquadraticcurveto{\pgfqpoint{2.805319in}{0.990062in}}{\pgfqpoint{2.808813in}{0.982965in}}%
\pgfpathlineto{\pgfqpoint{2.820359in}{0.960000in}}%
\pgfpathlineto{\pgfqpoint{2.808129in}{0.960000in}}%
\pgfpathlineto{\pgfqpoint{2.797394in}{0.981561in}}%
\pgfpathquadraticcurveto{\pgfqpoint{2.793215in}{0.990008in}}{\pgfqpoint{2.789306in}{0.992764in}}%
\pgfpathquadraticcurveto{\pgfqpoint{2.785397in}{0.995520in}}{\pgfqpoint{2.778643in}{0.995520in}}%
\pgfpathlineto{\pgfqpoint{2.766251in}{0.995520in}}%
\pgfpathlineto{\pgfqpoint{2.766251in}{0.960000in}}%
\pgfpathlineto{\pgfqpoint{2.754885in}{0.960000in}}%
\pgfpathlineto{\pgfqpoint{2.754885in}{1.044045in}}%
\pgfpathlineto{\pgfqpoint{2.780552in}{1.044045in}}%
\pgfpathquadraticcurveto{\pgfqpoint{2.794962in}{1.044045in}}{\pgfqpoint{2.802059in}{1.038011in}}%
\pgfpathquadraticcurveto{\pgfqpoint{2.809156in}{1.031995in}}{\pgfqpoint{2.809156in}{1.019836in}}%
\pgfpathquadraticcurveto{\pgfqpoint{2.809156in}{1.011893in}}{\pgfqpoint{2.805463in}{1.006651in}}%
\pgfpathquadraticcurveto{\pgfqpoint{2.801771in}{1.001428in}}{\pgfqpoint{2.794746in}{0.999411in}}%
\pgfpathlineto{\pgfqpoint{2.794746in}{0.999411in}}%
\pgfpathclose%
\pgfpathmoveto{\pgfqpoint{2.766251in}{1.034696in}}%
\pgfpathlineto{\pgfqpoint{2.766251in}{1.004868in}}%
\pgfpathlineto{\pgfqpoint{2.780552in}{1.004868in}}%
\pgfpathquadraticcurveto{\pgfqpoint{2.788766in}{1.004868in}}{\pgfqpoint{2.792963in}{1.008669in}}%
\pgfpathquadraticcurveto{\pgfqpoint{2.797159in}{1.012469in}}{\pgfqpoint{2.797159in}{1.019836in}}%
\pgfpathquadraticcurveto{\pgfqpoint{2.797159in}{1.027203in}}{\pgfqpoint{2.792963in}{1.030950in}}%
\pgfpathquadraticcurveto{\pgfqpoint{2.788766in}{1.034696in}}{\pgfqpoint{2.780552in}{1.034696in}}%
\pgfpathlineto{\pgfqpoint{2.766251in}{1.034696in}}%
\pgfpathlineto{\pgfqpoint{2.766251in}{1.034696in}}%
\pgfpathclose%
\pgfpathmoveto{\pgfqpoint{2.883294in}{0.994115in}}%
\pgfpathlineto{\pgfqpoint{2.883294in}{0.989054in}}%
\pgfpathlineto{\pgfqpoint{2.835669in}{0.989054in}}%
\pgfpathquadraticcurveto{\pgfqpoint{2.836354in}{0.978354in}}{\pgfqpoint{2.842118in}{0.972753in}}%
\pgfpathquadraticcurveto{\pgfqpoint{2.847882in}{0.967151in}}{\pgfqpoint{2.858185in}{0.967151in}}%
\pgfpathquadraticcurveto{\pgfqpoint{2.864147in}{0.967151in}}{\pgfqpoint{2.869748in}{0.968610in}}%
\pgfpathquadraticcurveto{\pgfqpoint{2.875350in}{0.970069in}}{\pgfqpoint{2.880880in}{0.973005in}}%
\pgfpathlineto{\pgfqpoint{2.880880in}{0.963206in}}%
\pgfpathquadraticcurveto{\pgfqpoint{2.875296in}{0.960847in}}{\pgfqpoint{2.869442in}{0.959604in}}%
\pgfpathquadraticcurveto{\pgfqpoint{2.863588in}{0.958361in}}{\pgfqpoint{2.857572in}{0.958361in}}%
\pgfpathquadraticcurveto{\pgfqpoint{2.842478in}{0.958361in}}{\pgfqpoint{2.833670in}{0.967133in}}%
\pgfpathquadraticcurveto{\pgfqpoint{2.824862in}{0.975923in}}{\pgfqpoint{2.824862in}{0.990909in}}%
\pgfpathquadraticcurveto{\pgfqpoint{2.824862in}{1.006381in}}{\pgfqpoint{2.833220in}{1.015459in}}%
\pgfpathquadraticcurveto{\pgfqpoint{2.841577in}{1.024556in}}{\pgfqpoint{2.855771in}{1.024556in}}%
\pgfpathquadraticcurveto{\pgfqpoint{2.868488in}{1.024556in}}{\pgfqpoint{2.875891in}{1.016360in}}%
\pgfpathquadraticcurveto{\pgfqpoint{2.883294in}{1.008183in}}{\pgfqpoint{2.883294in}{0.994115in}}%
\pgfpathlineto{\pgfqpoint{2.883294in}{0.994115in}}%
\pgfpathclose%
\pgfpathmoveto{\pgfqpoint{2.872937in}{0.997159in}}%
\pgfpathquadraticcurveto{\pgfqpoint{2.872828in}{1.005643in}}{\pgfqpoint{2.868181in}{1.010704in}}%
\pgfpathquadraticcurveto{\pgfqpoint{2.863534in}{1.015784in}}{\pgfqpoint{2.855879in}{1.015784in}}%
\pgfpathquadraticcurveto{\pgfqpoint{2.847215in}{1.015784in}}{\pgfqpoint{2.842010in}{1.010884in}}%
\pgfpathquadraticcurveto{\pgfqpoint{2.836804in}{1.005985in}}{\pgfqpoint{2.836012in}{0.997087in}}%
\pgfpathlineto{\pgfqpoint{2.872937in}{0.997159in}}%
\pgfpathlineto{\pgfqpoint{2.872937in}{0.997159in}}%
\pgfpathclose%
\pgfpathmoveto{\pgfqpoint{2.892876in}{1.023043in}}%
\pgfpathlineto{\pgfqpoint{2.903845in}{1.023043in}}%
\pgfpathlineto{\pgfqpoint{2.923551in}{0.970141in}}%
\pgfpathlineto{\pgfqpoint{2.943256in}{1.023043in}}%
\pgfpathlineto{\pgfqpoint{2.954225in}{1.023043in}}%
\pgfpathlineto{\pgfqpoint{2.930575in}{0.960000in}}%
\pgfpathlineto{\pgfqpoint{2.916508in}{0.960000in}}%
\pgfpathlineto{\pgfqpoint{2.892876in}{1.023043in}}%
\pgfpathlineto{\pgfqpoint{2.892876in}{1.023043in}}%
\pgfpathclose%
\pgfpathmoveto{\pgfqpoint{3.022455in}{0.994115in}}%
\pgfpathlineto{\pgfqpoint{3.022455in}{0.989054in}}%
\pgfpathlineto{\pgfqpoint{2.974831in}{0.989054in}}%
\pgfpathquadraticcurveto{\pgfqpoint{2.975516in}{0.978354in}}{\pgfqpoint{2.981280in}{0.972753in}}%
\pgfpathquadraticcurveto{\pgfqpoint{2.987044in}{0.967151in}}{\pgfqpoint{2.997346in}{0.967151in}}%
\pgfpathquadraticcurveto{\pgfqpoint{3.003309in}{0.967151in}}{\pgfqpoint{3.008910in}{0.968610in}}%
\pgfpathquadraticcurveto{\pgfqpoint{3.014512in}{0.970069in}}{\pgfqpoint{3.020042in}{0.973005in}}%
\pgfpathlineto{\pgfqpoint{3.020042in}{0.963206in}}%
\pgfpathquadraticcurveto{\pgfqpoint{3.014458in}{0.960847in}}{\pgfqpoint{3.008604in}{0.959604in}}%
\pgfpathquadraticcurveto{\pgfqpoint{3.002750in}{0.958361in}}{\pgfqpoint{2.996734in}{0.958361in}}%
\pgfpathquadraticcurveto{\pgfqpoint{2.981640in}{0.958361in}}{\pgfqpoint{2.972832in}{0.967133in}}%
\pgfpathquadraticcurveto{\pgfqpoint{2.964024in}{0.975923in}}{\pgfqpoint{2.964024in}{0.990909in}}%
\pgfpathquadraticcurveto{\pgfqpoint{2.964024in}{1.006381in}}{\pgfqpoint{2.972382in}{1.015459in}}%
\pgfpathquadraticcurveto{\pgfqpoint{2.980739in}{1.024556in}}{\pgfqpoint{2.994933in}{1.024556in}}%
\pgfpathquadraticcurveto{\pgfqpoint{3.007649in}{1.024556in}}{\pgfqpoint{3.015052in}{1.016360in}}%
\pgfpathquadraticcurveto{\pgfqpoint{3.022455in}{1.008183in}}{\pgfqpoint{3.022455in}{0.994115in}}%
\pgfpathlineto{\pgfqpoint{3.022455in}{0.994115in}}%
\pgfpathclose%
\pgfpathmoveto{\pgfqpoint{3.012098in}{0.997159in}}%
\pgfpathquadraticcurveto{\pgfqpoint{3.011990in}{1.005643in}}{\pgfqpoint{3.007343in}{1.010704in}}%
\pgfpathquadraticcurveto{\pgfqpoint{3.002696in}{1.015784in}}{\pgfqpoint{2.995041in}{1.015784in}}%
\pgfpathquadraticcurveto{\pgfqpoint{2.986377in}{1.015784in}}{\pgfqpoint{2.981172in}{1.010884in}}%
\pgfpathquadraticcurveto{\pgfqpoint{2.975966in}{1.005985in}}{\pgfqpoint{2.975174in}{0.997087in}}%
\pgfpathlineto{\pgfqpoint{3.012098in}{0.997159in}}%
\pgfpathlineto{\pgfqpoint{3.012098in}{0.997159in}}%
\pgfpathclose%
\pgfpathmoveto{\pgfqpoint{3.075988in}{1.013370in}}%
\pgfpathquadraticcurveto{\pgfqpoint{3.074240in}{1.014379in}}{\pgfqpoint{3.072187in}{1.014847in}}%
\pgfpathquadraticcurveto{\pgfqpoint{3.070134in}{1.015333in}}{\pgfqpoint{3.067666in}{1.015333in}}%
\pgfpathquadraticcurveto{\pgfqpoint{3.058876in}{1.015333in}}{\pgfqpoint{3.054175in}{1.009623in}}%
\pgfpathquadraticcurveto{\pgfqpoint{3.049474in}{1.003914in}}{\pgfqpoint{3.049474in}{0.993214in}}%
\pgfpathlineto{\pgfqpoint{3.049474in}{0.960000in}}%
\pgfpathlineto{\pgfqpoint{3.039063in}{0.960000in}}%
\pgfpathlineto{\pgfqpoint{3.039063in}{1.023043in}}%
\pgfpathlineto{\pgfqpoint{3.049474in}{1.023043in}}%
\pgfpathlineto{\pgfqpoint{3.049474in}{1.013244in}}%
\pgfpathquadraticcurveto{\pgfqpoint{3.052752in}{1.018990in}}{\pgfqpoint{3.057975in}{1.021764in}}%
\pgfpathquadraticcurveto{\pgfqpoint{3.063217in}{1.024556in}}{\pgfqpoint{3.070710in}{1.024556in}}%
\pgfpathquadraticcurveto{\pgfqpoint{3.071773in}{1.024556in}}{\pgfqpoint{3.073070in}{1.024411in}}%
\pgfpathquadraticcurveto{\pgfqpoint{3.074366in}{1.024285in}}{\pgfqpoint{3.075934in}{1.023997in}}%
\pgfpathlineto{\pgfqpoint{3.075988in}{1.013370in}}%
\pgfpathlineto{\pgfqpoint{3.075988in}{1.013370in}}%
\pgfpathclose%
\pgfpathmoveto{\pgfqpoint{3.127034in}{1.021187in}}%
\pgfpathlineto{\pgfqpoint{3.127034in}{1.011389in}}%
\pgfpathquadraticcurveto{\pgfqpoint{3.122657in}{1.013640in}}{\pgfqpoint{3.117920in}{1.014757in}}%
\pgfpathquadraticcurveto{\pgfqpoint{3.113201in}{1.015892in}}{\pgfqpoint{3.108121in}{1.015892in}}%
\pgfpathquadraticcurveto{\pgfqpoint{3.100412in}{1.015892in}}{\pgfqpoint{3.096557in}{1.013532in}}%
\pgfpathquadraticcurveto{\pgfqpoint{3.092703in}{1.011173in}}{\pgfqpoint{3.092703in}{1.006435in}}%
\pgfpathquadraticcurveto{\pgfqpoint{3.092703in}{1.002833in}}{\pgfqpoint{3.095459in}{1.000780in}}%
\pgfpathquadraticcurveto{\pgfqpoint{3.098215in}{0.998726in}}{\pgfqpoint{3.106554in}{0.996871in}}%
\pgfpathlineto{\pgfqpoint{3.110103in}{0.996078in}}%
\pgfpathquadraticcurveto{\pgfqpoint{3.121126in}{0.993719in}}{\pgfqpoint{3.125773in}{0.989414in}}%
\pgfpathquadraticcurveto{\pgfqpoint{3.130420in}{0.985109in}}{\pgfqpoint{3.130420in}{0.977400in}}%
\pgfpathquadraticcurveto{\pgfqpoint{3.130420in}{0.968610in}}{\pgfqpoint{3.123468in}{0.963476in}}%
\pgfpathquadraticcurveto{\pgfqpoint{3.116515in}{0.958361in}}{\pgfqpoint{3.104357in}{0.958361in}}%
\pgfpathquadraticcurveto{\pgfqpoint{3.099295in}{0.958361in}}{\pgfqpoint{3.093802in}{0.959352in}}%
\pgfpathquadraticcurveto{\pgfqpoint{3.088308in}{0.960342in}}{\pgfqpoint{3.082238in}{0.962306in}}%
\pgfpathlineto{\pgfqpoint{3.082238in}{0.973005in}}%
\pgfpathquadraticcurveto{\pgfqpoint{3.087984in}{0.970015in}}{\pgfqpoint{3.093549in}{0.968520in}}%
\pgfpathquadraticcurveto{\pgfqpoint{3.099115in}{0.967043in}}{\pgfqpoint{3.104591in}{0.967043in}}%
\pgfpathquadraticcurveto{\pgfqpoint{3.111904in}{0.967043in}}{\pgfqpoint{3.115830in}{0.969546in}}%
\pgfpathquadraticcurveto{\pgfqpoint{3.119775in}{0.972050in}}{\pgfqpoint{3.119775in}{0.976607in}}%
\pgfpathquadraticcurveto{\pgfqpoint{3.119775in}{0.980822in}}{\pgfqpoint{3.116929in}{0.983074in}}%
\pgfpathquadraticcurveto{\pgfqpoint{3.114101in}{0.985325in}}{\pgfqpoint{3.104465in}{0.987414in}}%
\pgfpathlineto{\pgfqpoint{3.100862in}{0.988261in}}%
\pgfpathquadraticcurveto{\pgfqpoint{3.091244in}{0.990278in}}{\pgfqpoint{3.086957in}{0.994475in}}%
\pgfpathquadraticcurveto{\pgfqpoint{3.082688in}{0.998672in}}{\pgfqpoint{3.082688in}{1.005985in}}%
\pgfpathquadraticcurveto{\pgfqpoint{3.082688in}{1.014883in}}{\pgfqpoint{3.088992in}{1.019710in}}%
\pgfpathquadraticcurveto{\pgfqpoint{3.095297in}{1.024556in}}{\pgfqpoint{3.106896in}{1.024556in}}%
\pgfpathquadraticcurveto{\pgfqpoint{3.112624in}{1.024556in}}{\pgfqpoint{3.117686in}{1.023709in}}%
\pgfpathquadraticcurveto{\pgfqpoint{3.122765in}{1.022880in}}{\pgfqpoint{3.127034in}{1.021187in}}%
\pgfpathlineto{\pgfqpoint{3.127034in}{1.021187in}}%
\pgfpathclose%
\pgfpathmoveto{\pgfqpoint{3.175559in}{0.991683in}}%
\pgfpathquadraticcurveto{\pgfqpoint{3.163004in}{0.991683in}}{\pgfqpoint{3.158159in}{0.988819in}}%
\pgfpathquadraticcurveto{\pgfqpoint{3.153314in}{0.985956in}}{\pgfqpoint{3.153314in}{0.979021in}}%
\pgfpathquadraticcurveto{\pgfqpoint{3.153314in}{0.973509in}}{\pgfqpoint{3.156952in}{0.970267in}}%
\pgfpathquadraticcurveto{\pgfqpoint{3.160591in}{0.967043in}}{\pgfqpoint{3.166823in}{0.967043in}}%
\pgfpathquadraticcurveto{\pgfqpoint{3.175451in}{0.967043in}}{\pgfqpoint{3.180656in}{0.973149in}}%
\pgfpathquadraticcurveto{\pgfqpoint{3.185862in}{0.979255in}}{\pgfqpoint{3.185862in}{0.989378in}}%
\pgfpathlineto{\pgfqpoint{3.185862in}{0.991683in}}%
\pgfpathlineto{\pgfqpoint{3.175559in}{0.991683in}}%
\pgfpathlineto{\pgfqpoint{3.175559in}{0.991683in}}%
\pgfpathclose%
\pgfpathmoveto{\pgfqpoint{3.196219in}{0.995970in}}%
\pgfpathlineto{\pgfqpoint{3.196219in}{0.960000in}}%
\pgfpathlineto{\pgfqpoint{3.185862in}{0.960000in}}%
\pgfpathlineto{\pgfqpoint{3.185862in}{0.969564in}}%
\pgfpathquadraticcurveto{\pgfqpoint{3.182313in}{0.963837in}}{\pgfqpoint{3.177018in}{0.961099in}}%
\pgfpathquadraticcurveto{\pgfqpoint{3.171722in}{0.958361in}}{\pgfqpoint{3.164067in}{0.958361in}}%
\pgfpathquadraticcurveto{\pgfqpoint{3.154394in}{0.958361in}}{\pgfqpoint{3.148667in}{0.963801in}}%
\pgfpathquadraticcurveto{\pgfqpoint{3.142957in}{0.969240in}}{\pgfqpoint{3.142957in}{0.978354in}}%
\pgfpathquadraticcurveto{\pgfqpoint{3.142957in}{0.988982in}}{\pgfqpoint{3.150072in}{0.994385in}}%
\pgfpathquadraticcurveto{\pgfqpoint{3.157204in}{0.999789in}}{\pgfqpoint{3.171326in}{0.999789in}}%
\pgfpathlineto{\pgfqpoint{3.185862in}{0.999789in}}%
\pgfpathlineto{\pgfqpoint{3.185862in}{1.000816in}}%
\pgfpathquadraticcurveto{\pgfqpoint{3.185862in}{1.007966in}}{\pgfqpoint{3.181161in}{1.011875in}}%
\pgfpathquadraticcurveto{\pgfqpoint{3.176459in}{1.015784in}}{\pgfqpoint{3.167958in}{1.015784in}}%
\pgfpathquadraticcurveto{\pgfqpoint{3.162554in}{1.015784in}}{\pgfqpoint{3.157420in}{1.014487in}}%
\pgfpathquadraticcurveto{\pgfqpoint{3.152305in}{1.013190in}}{\pgfqpoint{3.147586in}{1.010596in}}%
\pgfpathlineto{\pgfqpoint{3.147586in}{1.020179in}}%
\pgfpathquadraticcurveto{\pgfqpoint{3.153260in}{1.022376in}}{\pgfqpoint{3.158609in}{1.023457in}}%
\pgfpathquadraticcurveto{\pgfqpoint{3.163959in}{1.024556in}}{\pgfqpoint{3.169020in}{1.024556in}}%
\pgfpathquadraticcurveto{\pgfqpoint{3.182710in}{1.024556in}}{\pgfqpoint{3.189464in}{1.017459in}}%
\pgfpathquadraticcurveto{\pgfqpoint{3.196219in}{1.010380in}}{\pgfqpoint{3.196219in}{0.995970in}}%
\pgfpathlineto{\pgfqpoint{3.196219in}{0.995970in}}%
\pgfpathclose%
\pgfpathmoveto{\pgfqpoint{3.217545in}{1.047593in}}%
\pgfpathlineto{\pgfqpoint{3.227902in}{1.047593in}}%
\pgfpathlineto{\pgfqpoint{3.227902in}{0.960000in}}%
\pgfpathlineto{\pgfqpoint{3.217545in}{0.960000in}}%
\pgfpathlineto{\pgfqpoint{3.217545in}{1.047593in}}%
\pgfpathlineto{\pgfqpoint{3.217545in}{1.047593in}}%
\pgfpathclose%
\pgfpathmoveto{\pgfqpoint{3.252218in}{0.974302in}}%
\pgfpathlineto{\pgfqpoint{3.264088in}{0.974302in}}%
\pgfpathlineto{\pgfqpoint{3.264088in}{0.960000in}}%
\pgfpathlineto{\pgfqpoint{3.252218in}{0.960000in}}%
\pgfpathlineto{\pgfqpoint{3.252218in}{0.974302in}}%
\pgfpathlineto{\pgfqpoint{3.252218in}{0.974302in}}%
\pgfpathclose%
\pgfpathmoveto{\pgfqpoint{3.252218in}{1.019602in}}%
\pgfpathlineto{\pgfqpoint{3.264088in}{1.019602in}}%
\pgfpathlineto{\pgfqpoint{3.264088in}{1.005319in}}%
\pgfpathlineto{\pgfqpoint{3.252218in}{1.005319in}}%
\pgfpathlineto{\pgfqpoint{3.252218in}{1.019602in}}%
\pgfpathlineto{\pgfqpoint{3.252218in}{1.019602in}}%
\pgfpathclose%
\pgfusepath{fill}%
\end{pgfscope}%
\begin{pgfscope}%
\pgfsetbuttcap%
\pgfsetroundjoin%
\definecolor{currentfill}{rgb}{0.309804,0.372549,0.419608}%
\pgfsetfillcolor{currentfill}%
\pgfsetlinewidth{0.000000pt}%
\definecolor{currentstroke}{rgb}{0.309804,0.372549,0.419608}%
\pgfsetstrokecolor{currentstroke}%
\pgfsetdash{}{0pt}%
\pgfpathmoveto{\pgfqpoint{3.363596in}{1.015784in}}%
\pgfpathquadraticcurveto{\pgfqpoint{3.355275in}{1.015784in}}{\pgfqpoint{3.350429in}{1.009281in}}%
\pgfpathquadraticcurveto{\pgfqpoint{3.345584in}{1.002779in}}{\pgfqpoint{3.345584in}{0.991467in}}%
\pgfpathquadraticcurveto{\pgfqpoint{3.345584in}{0.980156in}}{\pgfqpoint{3.350393in}{0.973653in}}%
\pgfpathquadraticcurveto{\pgfqpoint{3.355221in}{0.967151in}}{\pgfqpoint{3.363596in}{0.967151in}}%
\pgfpathquadraticcurveto{\pgfqpoint{3.371882in}{0.967151in}}{\pgfqpoint{3.376709in}{0.973671in}}%
\pgfpathquadraticcurveto{\pgfqpoint{3.381554in}{0.980210in}}{\pgfqpoint{3.381554in}{0.991467in}}%
\pgfpathquadraticcurveto{\pgfqpoint{3.381554in}{1.002671in}}{\pgfqpoint{3.376709in}{1.009227in}}%
\pgfpathquadraticcurveto{\pgfqpoint{3.371882in}{1.015784in}}{\pgfqpoint{3.363596in}{1.015784in}}%
\pgfpathlineto{\pgfqpoint{3.363596in}{1.015784in}}%
\pgfpathclose%
\pgfpathmoveto{\pgfqpoint{3.363596in}{1.024556in}}%
\pgfpathquadraticcurveto{\pgfqpoint{3.377105in}{1.024556in}}{\pgfqpoint{3.384814in}{1.015766in}}%
\pgfpathquadraticcurveto{\pgfqpoint{3.392542in}{1.006994in}}{\pgfqpoint{3.392542in}{0.991467in}}%
\pgfpathquadraticcurveto{\pgfqpoint{3.392542in}{0.975995in}}{\pgfqpoint{3.384814in}{0.967169in}}%
\pgfpathquadraticcurveto{\pgfqpoint{3.377105in}{0.958361in}}{\pgfqpoint{3.363596in}{0.958361in}}%
\pgfpathquadraticcurveto{\pgfqpoint{3.350033in}{0.958361in}}{\pgfqpoint{3.342342in}{0.967169in}}%
\pgfpathquadraticcurveto{\pgfqpoint{3.334669in}{0.975995in}}{\pgfqpoint{3.334669in}{0.991467in}}%
\pgfpathquadraticcurveto{\pgfqpoint{3.334669in}{1.006994in}}{\pgfqpoint{3.342342in}{1.015766in}}%
\pgfpathquadraticcurveto{\pgfqpoint{3.350033in}{1.024556in}}{\pgfqpoint{3.363596in}{1.024556in}}%
\pgfpathlineto{\pgfqpoint{3.363596in}{1.024556in}}%
\pgfpathclose%
\pgfpathmoveto{\pgfqpoint{3.419722in}{0.969456in}}%
\pgfpathlineto{\pgfqpoint{3.419722in}{0.936026in}}%
\pgfpathlineto{\pgfqpoint{3.409311in}{0.936026in}}%
\pgfpathlineto{\pgfqpoint{3.409311in}{1.023043in}}%
\pgfpathlineto{\pgfqpoint{3.419722in}{1.023043in}}%
\pgfpathlineto{\pgfqpoint{3.419722in}{1.013478in}}%
\pgfpathquadraticcurveto{\pgfqpoint{3.423000in}{1.019098in}}{\pgfqpoint{3.427972in}{1.021818in}}%
\pgfpathquadraticcurveto{\pgfqpoint{3.432961in}{1.024556in}}{\pgfqpoint{3.439878in}{1.024556in}}%
\pgfpathquadraticcurveto{\pgfqpoint{3.451369in}{1.024556in}}{\pgfqpoint{3.458538in}{1.015441in}}%
\pgfpathquadraticcurveto{\pgfqpoint{3.465725in}{1.006327in}}{\pgfqpoint{3.465725in}{0.991467in}}%
\pgfpathquadraticcurveto{\pgfqpoint{3.465725in}{0.976607in}}{\pgfqpoint{3.458538in}{0.967475in}}%
\pgfpathquadraticcurveto{\pgfqpoint{3.451369in}{0.958361in}}{\pgfqpoint{3.439878in}{0.958361in}}%
\pgfpathquadraticcurveto{\pgfqpoint{3.432961in}{0.958361in}}{\pgfqpoint{3.427972in}{0.961099in}}%
\pgfpathquadraticcurveto{\pgfqpoint{3.423000in}{0.963837in}}{\pgfqpoint{3.419722in}{0.969456in}}%
\pgfpathlineto{\pgfqpoint{3.419722in}{0.969456in}}%
\pgfpathclose%
\pgfpathmoveto{\pgfqpoint{3.454972in}{0.991467in}}%
\pgfpathquadraticcurveto{\pgfqpoint{3.454972in}{1.002887in}}{\pgfqpoint{3.450271in}{1.009389in}}%
\pgfpathquadraticcurveto{\pgfqpoint{3.445569in}{1.015892in}}{\pgfqpoint{3.437356in}{1.015892in}}%
\pgfpathquadraticcurveto{\pgfqpoint{3.429124in}{1.015892in}}{\pgfqpoint{3.424423in}{1.009389in}}%
\pgfpathquadraticcurveto{\pgfqpoint{3.419722in}{1.002887in}}{\pgfqpoint{3.419722in}{0.991467in}}%
\pgfpathquadraticcurveto{\pgfqpoint{3.419722in}{0.980048in}}{\pgfqpoint{3.424423in}{0.973545in}}%
\pgfpathquadraticcurveto{\pgfqpoint{3.429124in}{0.967043in}}{\pgfqpoint{3.437356in}{0.967043in}}%
\pgfpathquadraticcurveto{\pgfqpoint{3.445569in}{0.967043in}}{\pgfqpoint{3.450271in}{0.973545in}}%
\pgfpathquadraticcurveto{\pgfqpoint{3.454972in}{0.980048in}}{\pgfqpoint{3.454972in}{0.991467in}}%
\pgfpathlineto{\pgfqpoint{3.454972in}{0.991467in}}%
\pgfpathclose%
\pgfpathmoveto{\pgfqpoint{3.492905in}{0.969456in}}%
\pgfpathlineto{\pgfqpoint{3.492905in}{0.936026in}}%
\pgfpathlineto{\pgfqpoint{3.482494in}{0.936026in}}%
\pgfpathlineto{\pgfqpoint{3.482494in}{1.023043in}}%
\pgfpathlineto{\pgfqpoint{3.492905in}{1.023043in}}%
\pgfpathlineto{\pgfqpoint{3.492905in}{1.013478in}}%
\pgfpathquadraticcurveto{\pgfqpoint{3.496184in}{1.019098in}}{\pgfqpoint{3.501155in}{1.021818in}}%
\pgfpathquadraticcurveto{\pgfqpoint{3.506144in}{1.024556in}}{\pgfqpoint{3.513061in}{1.024556in}}%
\pgfpathquadraticcurveto{\pgfqpoint{3.524553in}{1.024556in}}{\pgfqpoint{3.531722in}{1.015441in}}%
\pgfpathquadraticcurveto{\pgfqpoint{3.538908in}{1.006327in}}{\pgfqpoint{3.538908in}{0.991467in}}%
\pgfpathquadraticcurveto{\pgfqpoint{3.538908in}{0.976607in}}{\pgfqpoint{3.531722in}{0.967475in}}%
\pgfpathquadraticcurveto{\pgfqpoint{3.524553in}{0.958361in}}{\pgfqpoint{3.513061in}{0.958361in}}%
\pgfpathquadraticcurveto{\pgfqpoint{3.506144in}{0.958361in}}{\pgfqpoint{3.501155in}{0.961099in}}%
\pgfpathquadraticcurveto{\pgfqpoint{3.496184in}{0.963837in}}{\pgfqpoint{3.492905in}{0.969456in}}%
\pgfpathlineto{\pgfqpoint{3.492905in}{0.969456in}}%
\pgfpathclose%
\pgfpathmoveto{\pgfqpoint{3.528155in}{0.991467in}}%
\pgfpathquadraticcurveto{\pgfqpoint{3.528155in}{1.002887in}}{\pgfqpoint{3.523454in}{1.009389in}}%
\pgfpathquadraticcurveto{\pgfqpoint{3.518753in}{1.015892in}}{\pgfqpoint{3.510539in}{1.015892in}}%
\pgfpathquadraticcurveto{\pgfqpoint{3.502308in}{1.015892in}}{\pgfqpoint{3.497607in}{1.009389in}}%
\pgfpathquadraticcurveto{\pgfqpoint{3.492905in}{1.002887in}}{\pgfqpoint{3.492905in}{0.991467in}}%
\pgfpathquadraticcurveto{\pgfqpoint{3.492905in}{0.980048in}}{\pgfqpoint{3.497607in}{0.973545in}}%
\pgfpathquadraticcurveto{\pgfqpoint{3.502308in}{0.967043in}}{\pgfqpoint{3.510539in}{0.967043in}}%
\pgfpathquadraticcurveto{\pgfqpoint{3.518753in}{0.967043in}}{\pgfqpoint{3.523454in}{0.973545in}}%
\pgfpathquadraticcurveto{\pgfqpoint{3.528155in}{0.980048in}}{\pgfqpoint{3.528155in}{0.991467in}}%
\pgfpathlineto{\pgfqpoint{3.528155in}{0.991467in}}%
\pgfpathclose%
\pgfpathmoveto{\pgfqpoint{3.580498in}{1.015784in}}%
\pgfpathquadraticcurveto{\pgfqpoint{3.572177in}{1.015784in}}{\pgfqpoint{3.567332in}{1.009281in}}%
\pgfpathquadraticcurveto{\pgfqpoint{3.562486in}{1.002779in}}{\pgfqpoint{3.562486in}{0.991467in}}%
\pgfpathquadraticcurveto{\pgfqpoint{3.562486in}{0.980156in}}{\pgfqpoint{3.567296in}{0.973653in}}%
\pgfpathquadraticcurveto{\pgfqpoint{3.572123in}{0.967151in}}{\pgfqpoint{3.580498in}{0.967151in}}%
\pgfpathquadraticcurveto{\pgfqpoint{3.588784in}{0.967151in}}{\pgfqpoint{3.593611in}{0.973671in}}%
\pgfpathquadraticcurveto{\pgfqpoint{3.598457in}{0.980210in}}{\pgfqpoint{3.598457in}{0.991467in}}%
\pgfpathquadraticcurveto{\pgfqpoint{3.598457in}{1.002671in}}{\pgfqpoint{3.593611in}{1.009227in}}%
\pgfpathquadraticcurveto{\pgfqpoint{3.588784in}{1.015784in}}{\pgfqpoint{3.580498in}{1.015784in}}%
\pgfpathlineto{\pgfqpoint{3.580498in}{1.015784in}}%
\pgfpathclose%
\pgfpathmoveto{\pgfqpoint{3.580498in}{1.024556in}}%
\pgfpathquadraticcurveto{\pgfqpoint{3.594008in}{1.024556in}}{\pgfqpoint{3.601717in}{1.015766in}}%
\pgfpathquadraticcurveto{\pgfqpoint{3.609444in}{1.006994in}}{\pgfqpoint{3.609444in}{0.991467in}}%
\pgfpathquadraticcurveto{\pgfqpoint{3.609444in}{0.975995in}}{\pgfqpoint{3.601717in}{0.967169in}}%
\pgfpathquadraticcurveto{\pgfqpoint{3.594008in}{0.958361in}}{\pgfqpoint{3.580498in}{0.958361in}}%
\pgfpathquadraticcurveto{\pgfqpoint{3.566935in}{0.958361in}}{\pgfqpoint{3.559244in}{0.967169in}}%
\pgfpathquadraticcurveto{\pgfqpoint{3.551571in}{0.975995in}}{\pgfqpoint{3.551571in}{0.991467in}}%
\pgfpathquadraticcurveto{\pgfqpoint{3.551571in}{1.006994in}}{\pgfqpoint{3.559244in}{1.015766in}}%
\pgfpathquadraticcurveto{\pgfqpoint{3.566935in}{1.024556in}}{\pgfqpoint{3.580498in}{1.024556in}}%
\pgfpathlineto{\pgfqpoint{3.580498in}{1.024556in}}%
\pgfpathclose%
\pgfpathmoveto{\pgfqpoint{3.666795in}{1.021187in}}%
\pgfpathlineto{\pgfqpoint{3.666795in}{1.011389in}}%
\pgfpathquadraticcurveto{\pgfqpoint{3.662418in}{1.013640in}}{\pgfqpoint{3.657681in}{1.014757in}}%
\pgfpathquadraticcurveto{\pgfqpoint{3.652961in}{1.015892in}}{\pgfqpoint{3.647882in}{1.015892in}}%
\pgfpathquadraticcurveto{\pgfqpoint{3.640173in}{1.015892in}}{\pgfqpoint{3.636318in}{1.013532in}}%
\pgfpathquadraticcurveto{\pgfqpoint{3.632464in}{1.011173in}}{\pgfqpoint{3.632464in}{1.006435in}}%
\pgfpathquadraticcurveto{\pgfqpoint{3.632464in}{1.002833in}}{\pgfqpoint{3.635219in}{1.000780in}}%
\pgfpathquadraticcurveto{\pgfqpoint{3.637975in}{0.998726in}}{\pgfqpoint{3.646315in}{0.996871in}}%
\pgfpathlineto{\pgfqpoint{3.649863in}{0.996078in}}%
\pgfpathquadraticcurveto{\pgfqpoint{3.660887in}{0.993719in}}{\pgfqpoint{3.665534in}{0.989414in}}%
\pgfpathquadraticcurveto{\pgfqpoint{3.670181in}{0.985109in}}{\pgfqpoint{3.670181in}{0.977400in}}%
\pgfpathquadraticcurveto{\pgfqpoint{3.670181in}{0.968610in}}{\pgfqpoint{3.663228in}{0.963476in}}%
\pgfpathquadraticcurveto{\pgfqpoint{3.656276in}{0.958361in}}{\pgfqpoint{3.644117in}{0.958361in}}%
\pgfpathquadraticcurveto{\pgfqpoint{3.639056in}{0.958361in}}{\pgfqpoint{3.633562in}{0.959352in}}%
\pgfpathquadraticcurveto{\pgfqpoint{3.628069in}{0.960342in}}{\pgfqpoint{3.621998in}{0.962306in}}%
\pgfpathlineto{\pgfqpoint{3.621998in}{0.973005in}}%
\pgfpathquadraticcurveto{\pgfqpoint{3.627744in}{0.970015in}}{\pgfqpoint{3.633310in}{0.968520in}}%
\pgfpathquadraticcurveto{\pgfqpoint{3.638876in}{0.967043in}}{\pgfqpoint{3.644352in}{0.967043in}}%
\pgfpathquadraticcurveto{\pgfqpoint{3.651665in}{0.967043in}}{\pgfqpoint{3.655591in}{0.969546in}}%
\pgfpathquadraticcurveto{\pgfqpoint{3.659536in}{0.972050in}}{\pgfqpoint{3.659536in}{0.976607in}}%
\pgfpathquadraticcurveto{\pgfqpoint{3.659536in}{0.980822in}}{\pgfqpoint{3.656690in}{0.983074in}}%
\pgfpathquadraticcurveto{\pgfqpoint{3.653862in}{0.985325in}}{\pgfqpoint{3.644225in}{0.987414in}}%
\pgfpathlineto{\pgfqpoint{3.640623in}{0.988261in}}%
\pgfpathquadraticcurveto{\pgfqpoint{3.631005in}{0.990278in}}{\pgfqpoint{3.626718in}{0.994475in}}%
\pgfpathquadraticcurveto{\pgfqpoint{3.622449in}{0.998672in}}{\pgfqpoint{3.622449in}{1.005985in}}%
\pgfpathquadraticcurveto{\pgfqpoint{3.622449in}{1.014883in}}{\pgfqpoint{3.628753in}{1.019710in}}%
\pgfpathquadraticcurveto{\pgfqpoint{3.635057in}{1.024556in}}{\pgfqpoint{3.646657in}{1.024556in}}%
\pgfpathquadraticcurveto{\pgfqpoint{3.652385in}{1.024556in}}{\pgfqpoint{3.657446in}{1.023709in}}%
\pgfpathquadraticcurveto{\pgfqpoint{3.662526in}{1.022880in}}{\pgfqpoint{3.666795in}{1.021187in}}%
\pgfpathlineto{\pgfqpoint{3.666795in}{1.021187in}}%
\pgfpathclose%
\pgfpathmoveto{\pgfqpoint{3.686662in}{1.023043in}}%
\pgfpathlineto{\pgfqpoint{3.697019in}{1.023043in}}%
\pgfpathlineto{\pgfqpoint{3.697019in}{0.960000in}}%
\pgfpathlineto{\pgfqpoint{3.686662in}{0.960000in}}%
\pgfpathlineto{\pgfqpoint{3.686662in}{1.023043in}}%
\pgfpathlineto{\pgfqpoint{3.686662in}{1.023043in}}%
\pgfpathclose%
\pgfpathmoveto{\pgfqpoint{3.686662in}{1.047593in}}%
\pgfpathlineto{\pgfqpoint{3.697019in}{1.047593in}}%
\pgfpathlineto{\pgfqpoint{3.697019in}{1.034462in}}%
\pgfpathlineto{\pgfqpoint{3.686662in}{1.034462in}}%
\pgfpathlineto{\pgfqpoint{3.686662in}{1.047593in}}%
\pgfpathlineto{\pgfqpoint{3.686662in}{1.047593in}}%
\pgfpathclose%
\pgfpathmoveto{\pgfqpoint{3.728937in}{1.040947in}}%
\pgfpathlineto{\pgfqpoint{3.728937in}{1.023043in}}%
\pgfpathlineto{\pgfqpoint{3.750263in}{1.023043in}}%
\pgfpathlineto{\pgfqpoint{3.750263in}{1.014991in}}%
\pgfpathlineto{\pgfqpoint{3.728937in}{1.014991in}}%
\pgfpathlineto{\pgfqpoint{3.728937in}{0.980768in}}%
\pgfpathquadraticcurveto{\pgfqpoint{3.728937in}{0.973059in}}{\pgfqpoint{3.731044in}{0.970861in}}%
\pgfpathquadraticcurveto{\pgfqpoint{3.733151in}{0.968664in}}{\pgfqpoint{3.739636in}{0.968664in}}%
\pgfpathlineto{\pgfqpoint{3.750263in}{0.968664in}}%
\pgfpathlineto{\pgfqpoint{3.750263in}{0.960000in}}%
\pgfpathlineto{\pgfqpoint{3.739636in}{0.960000in}}%
\pgfpathquadraticcurveto{\pgfqpoint{3.727640in}{0.960000in}}{\pgfqpoint{3.723083in}{0.964467in}}%
\pgfpathquadraticcurveto{\pgfqpoint{3.718526in}{0.968952in}}{\pgfqpoint{3.718526in}{0.980768in}}%
\pgfpathlineto{\pgfqpoint{3.718526in}{1.014991in}}%
\pgfpathlineto{\pgfqpoint{3.710924in}{1.014991in}}%
\pgfpathlineto{\pgfqpoint{3.710924in}{1.023043in}}%
\pgfpathlineto{\pgfqpoint{3.718526in}{1.023043in}}%
\pgfpathlineto{\pgfqpoint{3.718526in}{1.040947in}}%
\pgfpathlineto{\pgfqpoint{3.728937in}{1.040947in}}%
\pgfpathlineto{\pgfqpoint{3.728937in}{1.040947in}}%
\pgfpathclose%
\pgfpathmoveto{\pgfqpoint{3.817809in}{0.994115in}}%
\pgfpathlineto{\pgfqpoint{3.817809in}{0.989054in}}%
\pgfpathlineto{\pgfqpoint{3.770184in}{0.989054in}}%
\pgfpathquadraticcurveto{\pgfqpoint{3.770869in}{0.978354in}}{\pgfqpoint{3.776633in}{0.972753in}}%
\pgfpathquadraticcurveto{\pgfqpoint{3.782397in}{0.967151in}}{\pgfqpoint{3.792700in}{0.967151in}}%
\pgfpathquadraticcurveto{\pgfqpoint{3.798662in}{0.967151in}}{\pgfqpoint{3.804263in}{0.968610in}}%
\pgfpathquadraticcurveto{\pgfqpoint{3.809865in}{0.970069in}}{\pgfqpoint{3.815395in}{0.973005in}}%
\pgfpathlineto{\pgfqpoint{3.815395in}{0.963206in}}%
\pgfpathquadraticcurveto{\pgfqpoint{3.809811in}{0.960847in}}{\pgfqpoint{3.803957in}{0.959604in}}%
\pgfpathquadraticcurveto{\pgfqpoint{3.798103in}{0.958361in}}{\pgfqpoint{3.792087in}{0.958361in}}%
\pgfpathquadraticcurveto{\pgfqpoint{3.776993in}{0.958361in}}{\pgfqpoint{3.768185in}{0.967133in}}%
\pgfpathquadraticcurveto{\pgfqpoint{3.759377in}{0.975923in}}{\pgfqpoint{3.759377in}{0.990909in}}%
\pgfpathquadraticcurveto{\pgfqpoint{3.759377in}{1.006381in}}{\pgfqpoint{3.767735in}{1.015459in}}%
\pgfpathquadraticcurveto{\pgfqpoint{3.776092in}{1.024556in}}{\pgfqpoint{3.790286in}{1.024556in}}%
\pgfpathquadraticcurveto{\pgfqpoint{3.803003in}{1.024556in}}{\pgfqpoint{3.810406in}{1.016360in}}%
\pgfpathquadraticcurveto{\pgfqpoint{3.817809in}{1.008183in}}{\pgfqpoint{3.817809in}{0.994115in}}%
\pgfpathlineto{\pgfqpoint{3.817809in}{0.994115in}}%
\pgfpathclose%
\pgfpathmoveto{\pgfqpoint{3.807452in}{0.997159in}}%
\pgfpathquadraticcurveto{\pgfqpoint{3.807344in}{1.005643in}}{\pgfqpoint{3.802696in}{1.010704in}}%
\pgfpathquadraticcurveto{\pgfqpoint{3.798049in}{1.015784in}}{\pgfqpoint{3.790394in}{1.015784in}}%
\pgfpathquadraticcurveto{\pgfqpoint{3.781730in}{1.015784in}}{\pgfqpoint{3.776525in}{1.010884in}}%
\pgfpathquadraticcurveto{\pgfqpoint{3.771319in}{1.005985in}}{\pgfqpoint{3.770527in}{0.997087in}}%
\pgfpathlineto{\pgfqpoint{3.807452in}{0.997159in}}%
\pgfpathlineto{\pgfqpoint{3.807452in}{0.997159in}}%
\pgfpathclose%
\pgfusepath{fill}%
\end{pgfscope}%
\begin{pgfscope}%
\pgfsetbuttcap%
\pgfsetroundjoin%
\definecolor{currentfill}{rgb}{0.309804,0.372549,0.419608}%
\pgfsetfillcolor{currentfill}%
\pgfsetlinewidth{0.000000pt}%
\definecolor{currentstroke}{rgb}{0.309804,0.372549,0.419608}%
\pgfsetstrokecolor{currentstroke}%
\pgfsetdash{}{0pt}%
\pgfpathmoveto{\pgfqpoint{3.951480in}{0.998060in}}%
\pgfpathlineto{\pgfqpoint{3.951480in}{0.960000in}}%
\pgfpathlineto{\pgfqpoint{3.941123in}{0.960000in}}%
\pgfpathlineto{\pgfqpoint{3.941123in}{0.997717in}}%
\pgfpathquadraticcurveto{\pgfqpoint{3.941123in}{1.006669in}}{\pgfqpoint{3.937629in}{1.011100in}}%
\pgfpathquadraticcurveto{\pgfqpoint{3.934135in}{1.015549in}}{\pgfqpoint{3.927164in}{1.015549in}}%
\pgfpathquadraticcurveto{\pgfqpoint{3.918770in}{1.015549in}}{\pgfqpoint{3.913925in}{1.010200in}}%
\pgfpathquadraticcurveto{\pgfqpoint{3.909080in}{1.004868in}}{\pgfqpoint{3.909080in}{0.995628in}}%
\pgfpathlineto{\pgfqpoint{3.909080in}{0.960000in}}%
\pgfpathlineto{\pgfqpoint{3.898669in}{0.960000in}}%
\pgfpathlineto{\pgfqpoint{3.898669in}{1.023043in}}%
\pgfpathlineto{\pgfqpoint{3.909080in}{1.023043in}}%
\pgfpathlineto{\pgfqpoint{3.909080in}{1.013244in}}%
\pgfpathquadraticcurveto{\pgfqpoint{3.912808in}{1.018936in}}{\pgfqpoint{3.917834in}{1.021746in}}%
\pgfpathquadraticcurveto{\pgfqpoint{3.922877in}{1.024556in}}{\pgfqpoint{3.929470in}{1.024556in}}%
\pgfpathquadraticcurveto{\pgfqpoint{3.940331in}{1.024556in}}{\pgfqpoint{3.945897in}{1.017837in}}%
\pgfpathquadraticcurveto{\pgfqpoint{3.951480in}{1.011118in}}{\pgfqpoint{3.951480in}{0.998060in}}%
\pgfpathlineto{\pgfqpoint{3.951480in}{0.998060in}}%
\pgfpathclose%
\pgfpathmoveto{\pgfqpoint{3.996547in}{1.015784in}}%
\pgfpathquadraticcurveto{\pgfqpoint{3.988225in}{1.015784in}}{\pgfqpoint{3.983380in}{1.009281in}}%
\pgfpathquadraticcurveto{\pgfqpoint{3.978535in}{1.002779in}}{\pgfqpoint{3.978535in}{0.991467in}}%
\pgfpathquadraticcurveto{\pgfqpoint{3.978535in}{0.980156in}}{\pgfqpoint{3.983344in}{0.973653in}}%
\pgfpathquadraticcurveto{\pgfqpoint{3.988171in}{0.967151in}}{\pgfqpoint{3.996547in}{0.967151in}}%
\pgfpathquadraticcurveto{\pgfqpoint{4.004832in}{0.967151in}}{\pgfqpoint{4.009660in}{0.973671in}}%
\pgfpathquadraticcurveto{\pgfqpoint{4.014505in}{0.980210in}}{\pgfqpoint{4.014505in}{0.991467in}}%
\pgfpathquadraticcurveto{\pgfqpoint{4.014505in}{1.002671in}}{\pgfqpoint{4.009660in}{1.009227in}}%
\pgfpathquadraticcurveto{\pgfqpoint{4.004832in}{1.015784in}}{\pgfqpoint{3.996547in}{1.015784in}}%
\pgfpathlineto{\pgfqpoint{3.996547in}{1.015784in}}%
\pgfpathclose%
\pgfpathmoveto{\pgfqpoint{3.996547in}{1.024556in}}%
\pgfpathquadraticcurveto{\pgfqpoint{4.010056in}{1.024556in}}{\pgfqpoint{4.017765in}{1.015766in}}%
\pgfpathquadraticcurveto{\pgfqpoint{4.025492in}{1.006994in}}{\pgfqpoint{4.025492in}{0.991467in}}%
\pgfpathquadraticcurveto{\pgfqpoint{4.025492in}{0.975995in}}{\pgfqpoint{4.017765in}{0.967169in}}%
\pgfpathquadraticcurveto{\pgfqpoint{4.010056in}{0.958361in}}{\pgfqpoint{3.996547in}{0.958361in}}%
\pgfpathquadraticcurveto{\pgfqpoint{3.982984in}{0.958361in}}{\pgfqpoint{3.975292in}{0.967169in}}%
\pgfpathquadraticcurveto{\pgfqpoint{3.967619in}{0.975995in}}{\pgfqpoint{3.967619in}{0.991467in}}%
\pgfpathquadraticcurveto{\pgfqpoint{3.967619in}{1.006994in}}{\pgfqpoint{3.975292in}{1.015766in}}%
\pgfpathquadraticcurveto{\pgfqpoint{3.982984in}{1.024556in}}{\pgfqpoint{3.996547in}{1.024556in}}%
\pgfpathlineto{\pgfqpoint{3.996547in}{1.024556in}}%
\pgfpathclose%
\pgfpathmoveto{\pgfqpoint{4.095073in}{0.998060in}}%
\pgfpathlineto{\pgfqpoint{4.095073in}{0.960000in}}%
\pgfpathlineto{\pgfqpoint{4.084716in}{0.960000in}}%
\pgfpathlineto{\pgfqpoint{4.084716in}{0.997717in}}%
\pgfpathquadraticcurveto{\pgfqpoint{4.084716in}{1.006669in}}{\pgfqpoint{4.081222in}{1.011100in}}%
\pgfpathquadraticcurveto{\pgfqpoint{4.077728in}{1.015549in}}{\pgfqpoint{4.070757in}{1.015549in}}%
\pgfpathquadraticcurveto{\pgfqpoint{4.062363in}{1.015549in}}{\pgfqpoint{4.057518in}{1.010200in}}%
\pgfpathquadraticcurveto{\pgfqpoint{4.052673in}{1.004868in}}{\pgfqpoint{4.052673in}{0.995628in}}%
\pgfpathlineto{\pgfqpoint{4.052673in}{0.960000in}}%
\pgfpathlineto{\pgfqpoint{4.042262in}{0.960000in}}%
\pgfpathlineto{\pgfqpoint{4.042262in}{1.023043in}}%
\pgfpathlineto{\pgfqpoint{4.052673in}{1.023043in}}%
\pgfpathlineto{\pgfqpoint{4.052673in}{1.013244in}}%
\pgfpathquadraticcurveto{\pgfqpoint{4.056401in}{1.018936in}}{\pgfqpoint{4.061427in}{1.021746in}}%
\pgfpathquadraticcurveto{\pgfqpoint{4.066470in}{1.024556in}}{\pgfqpoint{4.073062in}{1.024556in}}%
\pgfpathquadraticcurveto{\pgfqpoint{4.083924in}{1.024556in}}{\pgfqpoint{4.089489in}{1.017837in}}%
\pgfpathquadraticcurveto{\pgfqpoint{4.095073in}{1.011118in}}{\pgfqpoint{4.095073in}{0.998060in}}%
\pgfpathlineto{\pgfqpoint{4.095073in}{0.998060in}}%
\pgfpathclose%
\pgfpathmoveto{\pgfqpoint{4.111212in}{1.023043in}}%
\pgfpathlineto{\pgfqpoint{4.160403in}{1.023043in}}%
\pgfpathlineto{\pgfqpoint{4.160403in}{1.013586in}}%
\pgfpathlineto{\pgfqpoint{4.121461in}{0.968268in}}%
\pgfpathlineto{\pgfqpoint{4.160403in}{0.968268in}}%
\pgfpathlineto{\pgfqpoint{4.160403in}{0.960000in}}%
\pgfpathlineto{\pgfqpoint{4.109807in}{0.960000in}}%
\pgfpathlineto{\pgfqpoint{4.109807in}{0.969456in}}%
\pgfpathlineto{\pgfqpoint{4.148767in}{1.014775in}}%
\pgfpathlineto{\pgfqpoint{4.111212in}{1.014775in}}%
\pgfpathlineto{\pgfqpoint{4.111212in}{1.023043in}}%
\pgfpathlineto{\pgfqpoint{4.111212in}{1.023043in}}%
\pgfpathclose%
\pgfpathmoveto{\pgfqpoint{4.230146in}{0.994115in}}%
\pgfpathlineto{\pgfqpoint{4.230146in}{0.989054in}}%
\pgfpathlineto{\pgfqpoint{4.182522in}{0.989054in}}%
\pgfpathquadraticcurveto{\pgfqpoint{4.183207in}{0.978354in}}{\pgfqpoint{4.188971in}{0.972753in}}%
\pgfpathquadraticcurveto{\pgfqpoint{4.194734in}{0.967151in}}{\pgfqpoint{4.205037in}{0.967151in}}%
\pgfpathquadraticcurveto{\pgfqpoint{4.210999in}{0.967151in}}{\pgfqpoint{4.216601in}{0.968610in}}%
\pgfpathquadraticcurveto{\pgfqpoint{4.222203in}{0.970069in}}{\pgfqpoint{4.227733in}{0.973005in}}%
\pgfpathlineto{\pgfqpoint{4.227733in}{0.963206in}}%
\pgfpathquadraticcurveto{\pgfqpoint{4.222149in}{0.960847in}}{\pgfqpoint{4.216295in}{0.959604in}}%
\pgfpathquadraticcurveto{\pgfqpoint{4.210441in}{0.958361in}}{\pgfqpoint{4.204425in}{0.958361in}}%
\pgfpathquadraticcurveto{\pgfqpoint{4.189331in}{0.958361in}}{\pgfqpoint{4.180523in}{0.967133in}}%
\pgfpathquadraticcurveto{\pgfqpoint{4.171715in}{0.975923in}}{\pgfqpoint{4.171715in}{0.990909in}}%
\pgfpathquadraticcurveto{\pgfqpoint{4.171715in}{1.006381in}}{\pgfqpoint{4.180073in}{1.015459in}}%
\pgfpathquadraticcurveto{\pgfqpoint{4.188430in}{1.024556in}}{\pgfqpoint{4.202624in}{1.024556in}}%
\pgfpathquadraticcurveto{\pgfqpoint{4.215340in}{1.024556in}}{\pgfqpoint{4.222743in}{1.016360in}}%
\pgfpathquadraticcurveto{\pgfqpoint{4.230146in}{1.008183in}}{\pgfqpoint{4.230146in}{0.994115in}}%
\pgfpathlineto{\pgfqpoint{4.230146in}{0.994115in}}%
\pgfpathclose%
\pgfpathmoveto{\pgfqpoint{4.219789in}{0.997159in}}%
\pgfpathquadraticcurveto{\pgfqpoint{4.219681in}{1.005643in}}{\pgfqpoint{4.215034in}{1.010704in}}%
\pgfpathquadraticcurveto{\pgfqpoint{4.210387in}{1.015784in}}{\pgfqpoint{4.202732in}{1.015784in}}%
\pgfpathquadraticcurveto{\pgfqpoint{4.194068in}{1.015784in}}{\pgfqpoint{4.188863in}{1.010884in}}%
\pgfpathquadraticcurveto{\pgfqpoint{4.183657in}{1.005985in}}{\pgfqpoint{4.182864in}{0.997087in}}%
\pgfpathlineto{\pgfqpoint{4.219789in}{0.997159in}}%
\pgfpathlineto{\pgfqpoint{4.219789in}{0.997159in}}%
\pgfpathclose%
\pgfpathmoveto{\pgfqpoint{4.283678in}{1.013370in}}%
\pgfpathquadraticcurveto{\pgfqpoint{4.281931in}{1.014379in}}{\pgfqpoint{4.279878in}{1.014847in}}%
\pgfpathquadraticcurveto{\pgfqpoint{4.277825in}{1.015333in}}{\pgfqpoint{4.275357in}{1.015333in}}%
\pgfpathquadraticcurveto{\pgfqpoint{4.266567in}{1.015333in}}{\pgfqpoint{4.261866in}{1.009623in}}%
\pgfpathquadraticcurveto{\pgfqpoint{4.257165in}{1.003914in}}{\pgfqpoint{4.257165in}{0.993214in}}%
\pgfpathlineto{\pgfqpoint{4.257165in}{0.960000in}}%
\pgfpathlineto{\pgfqpoint{4.246754in}{0.960000in}}%
\pgfpathlineto{\pgfqpoint{4.246754in}{1.023043in}}%
\pgfpathlineto{\pgfqpoint{4.257165in}{1.023043in}}%
\pgfpathlineto{\pgfqpoint{4.257165in}{1.013244in}}%
\pgfpathquadraticcurveto{\pgfqpoint{4.260443in}{1.018990in}}{\pgfqpoint{4.265666in}{1.021764in}}%
\pgfpathquadraticcurveto{\pgfqpoint{4.270908in}{1.024556in}}{\pgfqpoint{4.278401in}{1.024556in}}%
\pgfpathquadraticcurveto{\pgfqpoint{4.279464in}{1.024556in}}{\pgfqpoint{4.280761in}{1.024411in}}%
\pgfpathquadraticcurveto{\pgfqpoint{4.282057in}{1.024285in}}{\pgfqpoint{4.283624in}{1.023997in}}%
\pgfpathlineto{\pgfqpoint{4.283678in}{1.013370in}}%
\pgfpathlineto{\pgfqpoint{4.283678in}{1.013370in}}%
\pgfpathclose%
\pgfpathmoveto{\pgfqpoint{4.316425in}{1.015784in}}%
\pgfpathquadraticcurveto{\pgfqpoint{4.308103in}{1.015784in}}{\pgfqpoint{4.303258in}{1.009281in}}%
\pgfpathquadraticcurveto{\pgfqpoint{4.298412in}{1.002779in}}{\pgfqpoint{4.298412in}{0.991467in}}%
\pgfpathquadraticcurveto{\pgfqpoint{4.298412in}{0.980156in}}{\pgfqpoint{4.303222in}{0.973653in}}%
\pgfpathquadraticcurveto{\pgfqpoint{4.308049in}{0.967151in}}{\pgfqpoint{4.316425in}{0.967151in}}%
\pgfpathquadraticcurveto{\pgfqpoint{4.324710in}{0.967151in}}{\pgfqpoint{4.329537in}{0.973671in}}%
\pgfpathquadraticcurveto{\pgfqpoint{4.334383in}{0.980210in}}{\pgfqpoint{4.334383in}{0.991467in}}%
\pgfpathquadraticcurveto{\pgfqpoint{4.334383in}{1.002671in}}{\pgfqpoint{4.329537in}{1.009227in}}%
\pgfpathquadraticcurveto{\pgfqpoint{4.324710in}{1.015784in}}{\pgfqpoint{4.316425in}{1.015784in}}%
\pgfpathlineto{\pgfqpoint{4.316425in}{1.015784in}}%
\pgfpathclose%
\pgfpathmoveto{\pgfqpoint{4.316425in}{1.024556in}}%
\pgfpathquadraticcurveto{\pgfqpoint{4.329934in}{1.024556in}}{\pgfqpoint{4.337643in}{1.015766in}}%
\pgfpathquadraticcurveto{\pgfqpoint{4.345370in}{1.006994in}}{\pgfqpoint{4.345370in}{0.991467in}}%
\pgfpathquadraticcurveto{\pgfqpoint{4.345370in}{0.975995in}}{\pgfqpoint{4.337643in}{0.967169in}}%
\pgfpathquadraticcurveto{\pgfqpoint{4.329934in}{0.958361in}}{\pgfqpoint{4.316425in}{0.958361in}}%
\pgfpathquadraticcurveto{\pgfqpoint{4.302861in}{0.958361in}}{\pgfqpoint{4.295170in}{0.967169in}}%
\pgfpathquadraticcurveto{\pgfqpoint{4.287497in}{0.975995in}}{\pgfqpoint{4.287497in}{0.991467in}}%
\pgfpathquadraticcurveto{\pgfqpoint{4.287497in}{1.006994in}}{\pgfqpoint{4.295170in}{1.015766in}}%
\pgfpathquadraticcurveto{\pgfqpoint{4.302861in}{1.024556in}}{\pgfqpoint{4.316425in}{1.024556in}}%
\pgfpathlineto{\pgfqpoint{4.316425in}{1.024556in}}%
\pgfpathclose%
\pgfusepath{fill}%
\end{pgfscope}%
\begin{pgfscope}%
\pgfsetbuttcap%
\pgfsetroundjoin%
\definecolor{currentfill}{rgb}{0.309804,0.372549,0.419608}%
\pgfsetfillcolor{currentfill}%
\pgfsetlinewidth{0.000000pt}%
\definecolor{currentstroke}{rgb}{0.309804,0.372549,0.419608}%
\pgfsetstrokecolor{currentstroke}%
\pgfsetdash{}{0pt}%
\pgfpathmoveto{\pgfqpoint{4.466643in}{1.021187in}}%
\pgfpathlineto{\pgfqpoint{4.466643in}{1.011389in}}%
\pgfpathquadraticcurveto{\pgfqpoint{4.462266in}{1.013640in}}{\pgfqpoint{4.457529in}{1.014757in}}%
\pgfpathquadraticcurveto{\pgfqpoint{4.452810in}{1.015892in}}{\pgfqpoint{4.447731in}{1.015892in}}%
\pgfpathquadraticcurveto{\pgfqpoint{4.440021in}{1.015892in}}{\pgfqpoint{4.436167in}{1.013532in}}%
\pgfpathquadraticcurveto{\pgfqpoint{4.432312in}{1.011173in}}{\pgfqpoint{4.432312in}{1.006435in}}%
\pgfpathquadraticcurveto{\pgfqpoint{4.432312in}{1.002833in}}{\pgfqpoint{4.435068in}{1.000780in}}%
\pgfpathquadraticcurveto{\pgfqpoint{4.437824in}{0.998726in}}{\pgfqpoint{4.446164in}{0.996871in}}%
\pgfpathlineto{\pgfqpoint{4.449712in}{0.996078in}}%
\pgfpathquadraticcurveto{\pgfqpoint{4.460735in}{0.993719in}}{\pgfqpoint{4.465383in}{0.989414in}}%
\pgfpathquadraticcurveto{\pgfqpoint{4.470030in}{0.985109in}}{\pgfqpoint{4.470030in}{0.977400in}}%
\pgfpathquadraticcurveto{\pgfqpoint{4.470030in}{0.968610in}}{\pgfqpoint{4.463077in}{0.963476in}}%
\pgfpathquadraticcurveto{\pgfqpoint{4.456124in}{0.958361in}}{\pgfqpoint{4.443966in}{0.958361in}}%
\pgfpathquadraticcurveto{\pgfqpoint{4.438905in}{0.958361in}}{\pgfqpoint{4.433411in}{0.959352in}}%
\pgfpathquadraticcurveto{\pgfqpoint{4.427917in}{0.960342in}}{\pgfqpoint{4.421847in}{0.962306in}}%
\pgfpathlineto{\pgfqpoint{4.421847in}{0.973005in}}%
\pgfpathquadraticcurveto{\pgfqpoint{4.427593in}{0.970015in}}{\pgfqpoint{4.433159in}{0.968520in}}%
\pgfpathquadraticcurveto{\pgfqpoint{4.438725in}{0.967043in}}{\pgfqpoint{4.444200in}{0.967043in}}%
\pgfpathquadraticcurveto{\pgfqpoint{4.451513in}{0.967043in}}{\pgfqpoint{4.455440in}{0.969546in}}%
\pgfpathquadraticcurveto{\pgfqpoint{4.459385in}{0.972050in}}{\pgfqpoint{4.459385in}{0.976607in}}%
\pgfpathquadraticcurveto{\pgfqpoint{4.459385in}{0.980822in}}{\pgfqpoint{4.456539in}{0.983074in}}%
\pgfpathquadraticcurveto{\pgfqpoint{4.453711in}{0.985325in}}{\pgfqpoint{4.444074in}{0.987414in}}%
\pgfpathlineto{\pgfqpoint{4.440472in}{0.988261in}}%
\pgfpathquadraticcurveto{\pgfqpoint{4.430853in}{0.990278in}}{\pgfqpoint{4.426566in}{0.994475in}}%
\pgfpathquadraticcurveto{\pgfqpoint{4.422298in}{0.998672in}}{\pgfqpoint{4.422298in}{1.005985in}}%
\pgfpathquadraticcurveto{\pgfqpoint{4.422298in}{1.014883in}}{\pgfqpoint{4.428602in}{1.019710in}}%
\pgfpathquadraticcurveto{\pgfqpoint{4.434906in}{1.024556in}}{\pgfqpoint{4.446506in}{1.024556in}}%
\pgfpathquadraticcurveto{\pgfqpoint{4.452234in}{1.024556in}}{\pgfqpoint{4.457295in}{1.023709in}}%
\pgfpathquadraticcurveto{\pgfqpoint{4.462375in}{1.022880in}}{\pgfqpoint{4.466643in}{1.021187in}}%
\pgfpathlineto{\pgfqpoint{4.466643in}{1.021187in}}%
\pgfpathclose%
\pgfpathmoveto{\pgfqpoint{4.486511in}{1.023043in}}%
\pgfpathlineto{\pgfqpoint{4.496868in}{1.023043in}}%
\pgfpathlineto{\pgfqpoint{4.496868in}{0.960000in}}%
\pgfpathlineto{\pgfqpoint{4.486511in}{0.960000in}}%
\pgfpathlineto{\pgfqpoint{4.486511in}{1.023043in}}%
\pgfpathlineto{\pgfqpoint{4.486511in}{1.023043in}}%
\pgfpathclose%
\pgfpathmoveto{\pgfqpoint{4.486511in}{1.047593in}}%
\pgfpathlineto{\pgfqpoint{4.496868in}{1.047593in}}%
\pgfpathlineto{\pgfqpoint{4.496868in}{1.034462in}}%
\pgfpathlineto{\pgfqpoint{4.486511in}{1.034462in}}%
\pgfpathlineto{\pgfqpoint{4.486511in}{1.047593in}}%
\pgfpathlineto{\pgfqpoint{4.486511in}{1.047593in}}%
\pgfpathclose%
\pgfpathmoveto{\pgfqpoint{4.560018in}{0.992260in}}%
\pgfpathquadraticcurveto{\pgfqpoint{4.560018in}{1.003517in}}{\pgfqpoint{4.555371in}{1.009696in}}%
\pgfpathquadraticcurveto{\pgfqpoint{4.550742in}{1.015892in}}{\pgfqpoint{4.542349in}{1.015892in}}%
\pgfpathquadraticcurveto{\pgfqpoint{4.534027in}{1.015892in}}{\pgfqpoint{4.529380in}{1.009696in}}%
\pgfpathquadraticcurveto{\pgfqpoint{4.524733in}{1.003517in}}{\pgfqpoint{4.524733in}{0.992260in}}%
\pgfpathquadraticcurveto{\pgfqpoint{4.524733in}{0.981056in}}{\pgfqpoint{4.529380in}{0.974860in}}%
\pgfpathquadraticcurveto{\pgfqpoint{4.534027in}{0.968664in}}{\pgfqpoint{4.542349in}{0.968664in}}%
\pgfpathquadraticcurveto{\pgfqpoint{4.550742in}{0.968664in}}{\pgfqpoint{4.555371in}{0.974860in}}%
\pgfpathquadraticcurveto{\pgfqpoint{4.560018in}{0.981056in}}{\pgfqpoint{4.560018in}{0.992260in}}%
\pgfpathlineto{\pgfqpoint{4.560018in}{0.992260in}}%
\pgfpathclose%
\pgfpathmoveto{\pgfqpoint{4.570375in}{0.967817in}}%
\pgfpathquadraticcurveto{\pgfqpoint{4.570375in}{0.951732in}}{\pgfqpoint{4.563225in}{0.943879in}}%
\pgfpathquadraticcurveto{\pgfqpoint{4.556092in}{0.936026in}}{\pgfqpoint{4.541340in}{0.936026in}}%
\pgfpathquadraticcurveto{\pgfqpoint{4.535882in}{0.936026in}}{\pgfqpoint{4.531037in}{0.936836in}}%
\pgfpathquadraticcurveto{\pgfqpoint{4.526192in}{0.937647in}}{\pgfqpoint{4.521635in}{0.939340in}}%
\pgfpathlineto{\pgfqpoint{4.521635in}{0.949409in}}%
\pgfpathquadraticcurveto{\pgfqpoint{4.526192in}{0.946941in}}{\pgfqpoint{4.530641in}{0.945770in}}%
\pgfpathquadraticcurveto{\pgfqpoint{4.535090in}{0.944582in}}{\pgfqpoint{4.539701in}{0.944582in}}%
\pgfpathquadraticcurveto{\pgfqpoint{4.549896in}{0.944582in}}{\pgfqpoint{4.554957in}{0.949895in}}%
\pgfpathquadraticcurveto{\pgfqpoint{4.560018in}{0.955209in}}{\pgfqpoint{4.560018in}{0.965962in}}%
\pgfpathlineto{\pgfqpoint{4.560018in}{0.971095in}}%
\pgfpathquadraticcurveto{\pgfqpoint{4.556812in}{0.965512in}}{\pgfqpoint{4.551805in}{0.962756in}}%
\pgfpathquadraticcurveto{\pgfqpoint{4.546798in}{0.960000in}}{\pgfqpoint{4.539809in}{0.960000in}}%
\pgfpathquadraticcurveto{\pgfqpoint{4.528227in}{0.960000in}}{\pgfqpoint{4.521130in}{0.968826in}}%
\pgfpathquadraticcurveto{\pgfqpoint{4.514033in}{0.977670in}}{\pgfqpoint{4.514033in}{0.992260in}}%
\pgfpathquadraticcurveto{\pgfqpoint{4.514033in}{1.006886in}}{\pgfqpoint{4.521130in}{1.015712in}}%
\pgfpathquadraticcurveto{\pgfqpoint{4.528227in}{1.024556in}}{\pgfqpoint{4.539809in}{1.024556in}}%
\pgfpathquadraticcurveto{\pgfqpoint{4.546798in}{1.024556in}}{\pgfqpoint{4.551805in}{1.021800in}}%
\pgfpathquadraticcurveto{\pgfqpoint{4.556812in}{1.019044in}}{\pgfqpoint{4.560018in}{1.013478in}}%
\pgfpathlineto{\pgfqpoint{4.560018in}{1.023043in}}%
\pgfpathlineto{\pgfqpoint{4.570375in}{1.023043in}}%
\pgfpathlineto{\pgfqpoint{4.570375in}{0.967817in}}%
\pgfpathlineto{\pgfqpoint{4.570375in}{0.967817in}}%
\pgfpathclose%
\pgfpathmoveto{\pgfqpoint{4.644135in}{0.998060in}}%
\pgfpathlineto{\pgfqpoint{4.644135in}{0.960000in}}%
\pgfpathlineto{\pgfqpoint{4.633778in}{0.960000in}}%
\pgfpathlineto{\pgfqpoint{4.633778in}{0.997717in}}%
\pgfpathquadraticcurveto{\pgfqpoint{4.633778in}{1.006669in}}{\pgfqpoint{4.630284in}{1.011100in}}%
\pgfpathquadraticcurveto{\pgfqpoint{4.626789in}{1.015549in}}{\pgfqpoint{4.619819in}{1.015549in}}%
\pgfpathquadraticcurveto{\pgfqpoint{4.611425in}{1.015549in}}{\pgfqpoint{4.606580in}{1.010200in}}%
\pgfpathquadraticcurveto{\pgfqpoint{4.601735in}{1.004868in}}{\pgfqpoint{4.601735in}{0.995628in}}%
\pgfpathlineto{\pgfqpoint{4.601735in}{0.960000in}}%
\pgfpathlineto{\pgfqpoint{4.591324in}{0.960000in}}%
\pgfpathlineto{\pgfqpoint{4.591324in}{1.023043in}}%
\pgfpathlineto{\pgfqpoint{4.601735in}{1.023043in}}%
\pgfpathlineto{\pgfqpoint{4.601735in}{1.013244in}}%
\pgfpathquadraticcurveto{\pgfqpoint{4.605463in}{1.018936in}}{\pgfqpoint{4.610488in}{1.021746in}}%
\pgfpathquadraticcurveto{\pgfqpoint{4.615532in}{1.024556in}}{\pgfqpoint{4.622124in}{1.024556in}}%
\pgfpathquadraticcurveto{\pgfqpoint{4.632986in}{1.024556in}}{\pgfqpoint{4.638551in}{1.017837in}}%
\pgfpathquadraticcurveto{\pgfqpoint{4.644135in}{1.011118in}}{\pgfqpoint{4.644135in}{0.998060in}}%
\pgfpathlineto{\pgfqpoint{4.644135in}{0.998060in}}%
\pgfpathclose%
\pgfpathmoveto{\pgfqpoint{4.704962in}{1.021187in}}%
\pgfpathlineto{\pgfqpoint{4.704962in}{1.011389in}}%
\pgfpathquadraticcurveto{\pgfqpoint{4.700585in}{1.013640in}}{\pgfqpoint{4.695848in}{1.014757in}}%
\pgfpathquadraticcurveto{\pgfqpoint{4.691129in}{1.015892in}}{\pgfqpoint{4.686049in}{1.015892in}}%
\pgfpathquadraticcurveto{\pgfqpoint{4.678340in}{1.015892in}}{\pgfqpoint{4.674486in}{1.013532in}}%
\pgfpathquadraticcurveto{\pgfqpoint{4.670631in}{1.011173in}}{\pgfqpoint{4.670631in}{1.006435in}}%
\pgfpathquadraticcurveto{\pgfqpoint{4.670631in}{1.002833in}}{\pgfqpoint{4.673387in}{1.000780in}}%
\pgfpathquadraticcurveto{\pgfqpoint{4.676143in}{0.998726in}}{\pgfqpoint{4.684482in}{0.996871in}}%
\pgfpathlineto{\pgfqpoint{4.688031in}{0.996078in}}%
\pgfpathquadraticcurveto{\pgfqpoint{4.699054in}{0.993719in}}{\pgfqpoint{4.703701in}{0.989414in}}%
\pgfpathquadraticcurveto{\pgfqpoint{4.708349in}{0.985109in}}{\pgfqpoint{4.708349in}{0.977400in}}%
\pgfpathquadraticcurveto{\pgfqpoint{4.708349in}{0.968610in}}{\pgfqpoint{4.701396in}{0.963476in}}%
\pgfpathquadraticcurveto{\pgfqpoint{4.694443in}{0.958361in}}{\pgfqpoint{4.682285in}{0.958361in}}%
\pgfpathquadraticcurveto{\pgfqpoint{4.677224in}{0.958361in}}{\pgfqpoint{4.671730in}{0.959352in}}%
\pgfpathquadraticcurveto{\pgfqpoint{4.666236in}{0.960342in}}{\pgfqpoint{4.660166in}{0.962306in}}%
\pgfpathlineto{\pgfqpoint{4.660166in}{0.973005in}}%
\pgfpathquadraticcurveto{\pgfqpoint{4.665912in}{0.970015in}}{\pgfqpoint{4.671478in}{0.968520in}}%
\pgfpathquadraticcurveto{\pgfqpoint{4.677043in}{0.967043in}}{\pgfqpoint{4.682519in}{0.967043in}}%
\pgfpathquadraticcurveto{\pgfqpoint{4.689832in}{0.967043in}}{\pgfqpoint{4.693759in}{0.969546in}}%
\pgfpathquadraticcurveto{\pgfqpoint{4.697703in}{0.972050in}}{\pgfqpoint{4.697703in}{0.976607in}}%
\pgfpathquadraticcurveto{\pgfqpoint{4.697703in}{0.980822in}}{\pgfqpoint{4.694857in}{0.983074in}}%
\pgfpathquadraticcurveto{\pgfqpoint{4.692029in}{0.985325in}}{\pgfqpoint{4.682393in}{0.987414in}}%
\pgfpathlineto{\pgfqpoint{4.678791in}{0.988261in}}%
\pgfpathquadraticcurveto{\pgfqpoint{4.669172in}{0.990278in}}{\pgfqpoint{4.664885in}{0.994475in}}%
\pgfpathquadraticcurveto{\pgfqpoint{4.660616in}{0.998672in}}{\pgfqpoint{4.660616in}{1.005985in}}%
\pgfpathquadraticcurveto{\pgfqpoint{4.660616in}{1.014883in}}{\pgfqpoint{4.666921in}{1.019710in}}%
\pgfpathquadraticcurveto{\pgfqpoint{4.673225in}{1.024556in}}{\pgfqpoint{4.684825in}{1.024556in}}%
\pgfpathquadraticcurveto{\pgfqpoint{4.690552in}{1.024556in}}{\pgfqpoint{4.695614in}{1.023709in}}%
\pgfpathquadraticcurveto{\pgfqpoint{4.700693in}{1.022880in}}{\pgfqpoint{4.704962in}{1.021187in}}%
\pgfpathlineto{\pgfqpoint{4.704962in}{1.021187in}}%
\pgfpathclose%
\pgfusepath{fill}%
\end{pgfscope}%
\begin{pgfscope}%
\pgfsetbuttcap%
\pgfsetroundjoin%
\definecolor{currentfill}{rgb}{0.309804,0.372549,0.419608}%
\pgfsetfillcolor{currentfill}%
\pgfsetlinewidth{0.000000pt}%
\definecolor{currentstroke}{rgb}{0.309804,0.372549,0.419608}%
\pgfsetstrokecolor{currentstroke}%
\pgfsetdash{}{0pt}%
\pgfpathmoveto{\pgfqpoint{4.815009in}{0.991683in}}%
\pgfpathquadraticcurveto{\pgfqpoint{4.802454in}{0.991683in}}{\pgfqpoint{4.797609in}{0.988819in}}%
\pgfpathquadraticcurveto{\pgfqpoint{4.792764in}{0.985956in}}{\pgfqpoint{4.792764in}{0.979021in}}%
\pgfpathquadraticcurveto{\pgfqpoint{4.792764in}{0.973509in}}{\pgfqpoint{4.796402in}{0.970267in}}%
\pgfpathquadraticcurveto{\pgfqpoint{4.800041in}{0.967043in}}{\pgfqpoint{4.806273in}{0.967043in}}%
\pgfpathquadraticcurveto{\pgfqpoint{4.814901in}{0.967043in}}{\pgfqpoint{4.820106in}{0.973149in}}%
\pgfpathquadraticcurveto{\pgfqpoint{4.825312in}{0.979255in}}{\pgfqpoint{4.825312in}{0.989378in}}%
\pgfpathlineto{\pgfqpoint{4.825312in}{0.991683in}}%
\pgfpathlineto{\pgfqpoint{4.815009in}{0.991683in}}%
\pgfpathlineto{\pgfqpoint{4.815009in}{0.991683in}}%
\pgfpathclose%
\pgfpathmoveto{\pgfqpoint{4.835669in}{0.995970in}}%
\pgfpathlineto{\pgfqpoint{4.835669in}{0.960000in}}%
\pgfpathlineto{\pgfqpoint{4.825312in}{0.960000in}}%
\pgfpathlineto{\pgfqpoint{4.825312in}{0.969564in}}%
\pgfpathquadraticcurveto{\pgfqpoint{4.821764in}{0.963837in}}{\pgfqpoint{4.816468in}{0.961099in}}%
\pgfpathquadraticcurveto{\pgfqpoint{4.811172in}{0.958361in}}{\pgfqpoint{4.803517in}{0.958361in}}%
\pgfpathquadraticcurveto{\pgfqpoint{4.793845in}{0.958361in}}{\pgfqpoint{4.788117in}{0.963801in}}%
\pgfpathquadraticcurveto{\pgfqpoint{4.782407in}{0.969240in}}{\pgfqpoint{4.782407in}{0.978354in}}%
\pgfpathquadraticcurveto{\pgfqpoint{4.782407in}{0.988982in}}{\pgfqpoint{4.789522in}{0.994385in}}%
\pgfpathquadraticcurveto{\pgfqpoint{4.796655in}{0.999789in}}{\pgfqpoint{4.810776in}{0.999789in}}%
\pgfpathlineto{\pgfqpoint{4.825312in}{0.999789in}}%
\pgfpathlineto{\pgfqpoint{4.825312in}{1.000816in}}%
\pgfpathquadraticcurveto{\pgfqpoint{4.825312in}{1.007966in}}{\pgfqpoint{4.820611in}{1.011875in}}%
\pgfpathquadraticcurveto{\pgfqpoint{4.815910in}{1.015784in}}{\pgfqpoint{4.807408in}{1.015784in}}%
\pgfpathquadraticcurveto{\pgfqpoint{4.802004in}{1.015784in}}{\pgfqpoint{4.796871in}{1.014487in}}%
\pgfpathquadraticcurveto{\pgfqpoint{4.791755in}{1.013190in}}{\pgfqpoint{4.787036in}{1.010596in}}%
\pgfpathlineto{\pgfqpoint{4.787036in}{1.020179in}}%
\pgfpathquadraticcurveto{\pgfqpoint{4.792710in}{1.022376in}}{\pgfqpoint{4.798060in}{1.023457in}}%
\pgfpathquadraticcurveto{\pgfqpoint{4.803409in}{1.024556in}}{\pgfqpoint{4.808471in}{1.024556in}}%
\pgfpathquadraticcurveto{\pgfqpoint{4.822160in}{1.024556in}}{\pgfqpoint{4.828914in}{1.017459in}}%
\pgfpathquadraticcurveto{\pgfqpoint{4.835669in}{1.010380in}}{\pgfqpoint{4.835669in}{0.995970in}}%
\pgfpathlineto{\pgfqpoint{4.835669in}{0.995970in}}%
\pgfpathclose%
\pgfpathmoveto{\pgfqpoint{4.909411in}{0.998060in}}%
\pgfpathlineto{\pgfqpoint{4.909411in}{0.960000in}}%
\pgfpathlineto{\pgfqpoint{4.899054in}{0.960000in}}%
\pgfpathlineto{\pgfqpoint{4.899054in}{0.997717in}}%
\pgfpathquadraticcurveto{\pgfqpoint{4.899054in}{1.006669in}}{\pgfqpoint{4.895559in}{1.011100in}}%
\pgfpathquadraticcurveto{\pgfqpoint{4.892065in}{1.015549in}}{\pgfqpoint{4.885094in}{1.015549in}}%
\pgfpathquadraticcurveto{\pgfqpoint{4.876701in}{1.015549in}}{\pgfqpoint{4.871855in}{1.010200in}}%
\pgfpathquadraticcurveto{\pgfqpoint{4.867010in}{1.004868in}}{\pgfqpoint{4.867010in}{0.995628in}}%
\pgfpathlineto{\pgfqpoint{4.867010in}{0.960000in}}%
\pgfpathlineto{\pgfqpoint{4.856599in}{0.960000in}}%
\pgfpathlineto{\pgfqpoint{4.856599in}{1.023043in}}%
\pgfpathlineto{\pgfqpoint{4.867010in}{1.023043in}}%
\pgfpathlineto{\pgfqpoint{4.867010in}{1.013244in}}%
\pgfpathquadraticcurveto{\pgfqpoint{4.870739in}{1.018936in}}{\pgfqpoint{4.875764in}{1.021746in}}%
\pgfpathquadraticcurveto{\pgfqpoint{4.880807in}{1.024556in}}{\pgfqpoint{4.887400in}{1.024556in}}%
\pgfpathquadraticcurveto{\pgfqpoint{4.898261in}{1.024556in}}{\pgfqpoint{4.903827in}{1.017837in}}%
\pgfpathquadraticcurveto{\pgfqpoint{4.909411in}{1.011118in}}{\pgfqpoint{4.909411in}{0.998060in}}%
\pgfpathlineto{\pgfqpoint{4.909411in}{0.998060in}}%
\pgfpathclose%
\pgfpathmoveto{\pgfqpoint{4.971535in}{1.013478in}}%
\pgfpathlineto{\pgfqpoint{4.971535in}{1.047593in}}%
\pgfpathlineto{\pgfqpoint{4.981892in}{1.047593in}}%
\pgfpathlineto{\pgfqpoint{4.981892in}{0.960000in}}%
\pgfpathlineto{\pgfqpoint{4.971535in}{0.960000in}}%
\pgfpathlineto{\pgfqpoint{4.971535in}{0.969456in}}%
\pgfpathquadraticcurveto{\pgfqpoint{4.968274in}{0.963837in}}{\pgfqpoint{4.963285in}{0.961099in}}%
\pgfpathquadraticcurveto{\pgfqpoint{4.958314in}{0.958361in}}{\pgfqpoint{4.951325in}{0.958361in}}%
\pgfpathquadraticcurveto{\pgfqpoint{4.939905in}{0.958361in}}{\pgfqpoint{4.932718in}{0.967475in}}%
\pgfpathquadraticcurveto{\pgfqpoint{4.925550in}{0.976607in}}{\pgfqpoint{4.925550in}{0.991467in}}%
\pgfpathquadraticcurveto{\pgfqpoint{4.925550in}{1.006327in}}{\pgfqpoint{4.932718in}{1.015441in}}%
\pgfpathquadraticcurveto{\pgfqpoint{4.939905in}{1.024556in}}{\pgfqpoint{4.951325in}{1.024556in}}%
\pgfpathquadraticcurveto{\pgfqpoint{4.958314in}{1.024556in}}{\pgfqpoint{4.963285in}{1.021818in}}%
\pgfpathquadraticcurveto{\pgfqpoint{4.968274in}{1.019098in}}{\pgfqpoint{4.971535in}{1.013478in}}%
\pgfpathlineto{\pgfqpoint{4.971535in}{1.013478in}}%
\pgfpathclose%
\pgfpathmoveto{\pgfqpoint{4.936249in}{0.991467in}}%
\pgfpathquadraticcurveto{\pgfqpoint{4.936249in}{0.980048in}}{\pgfqpoint{4.940950in}{0.973545in}}%
\pgfpathquadraticcurveto{\pgfqpoint{4.945651in}{0.967043in}}{\pgfqpoint{4.953865in}{0.967043in}}%
\pgfpathquadraticcurveto{\pgfqpoint{4.962078in}{0.967043in}}{\pgfqpoint{4.966797in}{0.973545in}}%
\pgfpathquadraticcurveto{\pgfqpoint{4.971535in}{0.980048in}}{\pgfqpoint{4.971535in}{0.991467in}}%
\pgfpathquadraticcurveto{\pgfqpoint{4.971535in}{1.002887in}}{\pgfqpoint{4.966797in}{1.009389in}}%
\pgfpathquadraticcurveto{\pgfqpoint{4.962078in}{1.015892in}}{\pgfqpoint{4.953865in}{1.015892in}}%
\pgfpathquadraticcurveto{\pgfqpoint{4.945651in}{1.015892in}}{\pgfqpoint{4.940950in}{1.009389in}}%
\pgfpathquadraticcurveto{\pgfqpoint{4.936249in}{1.002887in}}{\pgfqpoint{4.936249in}{0.991467in}}%
\pgfpathlineto{\pgfqpoint{4.936249in}{0.991467in}}%
\pgfpathclose%
\pgfusepath{fill}%
\end{pgfscope}%
\begin{pgfscope}%
\pgfsetbuttcap%
\pgfsetroundjoin%
\definecolor{currentfill}{rgb}{0.309804,0.372549,0.419608}%
\pgfsetfillcolor{currentfill}%
\pgfsetlinewidth{0.000000pt}%
\definecolor{currentstroke}{rgb}{0.309804,0.372549,0.419608}%
\pgfsetstrokecolor{currentstroke}%
\pgfsetdash{}{0pt}%
\pgfpathmoveto{\pgfqpoint{5.094902in}{0.991683in}}%
\pgfpathquadraticcurveto{\pgfqpoint{5.082348in}{0.991683in}}{\pgfqpoint{5.077503in}{0.988819in}}%
\pgfpathquadraticcurveto{\pgfqpoint{5.072657in}{0.985956in}}{\pgfqpoint{5.072657in}{0.979021in}}%
\pgfpathquadraticcurveto{\pgfqpoint{5.072657in}{0.973509in}}{\pgfqpoint{5.076296in}{0.970267in}}%
\pgfpathquadraticcurveto{\pgfqpoint{5.079934in}{0.967043in}}{\pgfqpoint{5.086166in}{0.967043in}}%
\pgfpathquadraticcurveto{\pgfqpoint{5.094794in}{0.967043in}}{\pgfqpoint{5.100000in}{0.973149in}}%
\pgfpathquadraticcurveto{\pgfqpoint{5.105205in}{0.979255in}}{\pgfqpoint{5.105205in}{0.989378in}}%
\pgfpathlineto{\pgfqpoint{5.105205in}{0.991683in}}%
\pgfpathlineto{\pgfqpoint{5.094902in}{0.991683in}}%
\pgfpathlineto{\pgfqpoint{5.094902in}{0.991683in}}%
\pgfpathclose%
\pgfpathmoveto{\pgfqpoint{5.115562in}{0.995970in}}%
\pgfpathlineto{\pgfqpoint{5.115562in}{0.960000in}}%
\pgfpathlineto{\pgfqpoint{5.105205in}{0.960000in}}%
\pgfpathlineto{\pgfqpoint{5.105205in}{0.969564in}}%
\pgfpathquadraticcurveto{\pgfqpoint{5.101657in}{0.963837in}}{\pgfqpoint{5.096361in}{0.961099in}}%
\pgfpathquadraticcurveto{\pgfqpoint{5.091066in}{0.958361in}}{\pgfqpoint{5.083411in}{0.958361in}}%
\pgfpathquadraticcurveto{\pgfqpoint{5.073738in}{0.958361in}}{\pgfqpoint{5.068010in}{0.963801in}}%
\pgfpathquadraticcurveto{\pgfqpoint{5.062300in}{0.969240in}}{\pgfqpoint{5.062300in}{0.978354in}}%
\pgfpathquadraticcurveto{\pgfqpoint{5.062300in}{0.988982in}}{\pgfqpoint{5.069415in}{0.994385in}}%
\pgfpathquadraticcurveto{\pgfqpoint{5.076548in}{0.999789in}}{\pgfqpoint{5.090669in}{0.999789in}}%
\pgfpathlineto{\pgfqpoint{5.105205in}{0.999789in}}%
\pgfpathlineto{\pgfqpoint{5.105205in}{1.000816in}}%
\pgfpathquadraticcurveto{\pgfqpoint{5.105205in}{1.007966in}}{\pgfqpoint{5.100504in}{1.011875in}}%
\pgfpathquadraticcurveto{\pgfqpoint{5.095803in}{1.015784in}}{\pgfqpoint{5.087301in}{1.015784in}}%
\pgfpathquadraticcurveto{\pgfqpoint{5.081898in}{1.015784in}}{\pgfqpoint{5.076764in}{1.014487in}}%
\pgfpathquadraticcurveto{\pgfqpoint{5.071649in}{1.013190in}}{\pgfqpoint{5.066929in}{1.010596in}}%
\pgfpathlineto{\pgfqpoint{5.066929in}{1.020179in}}%
\pgfpathquadraticcurveto{\pgfqpoint{5.072603in}{1.022376in}}{\pgfqpoint{5.077953in}{1.023457in}}%
\pgfpathquadraticcurveto{\pgfqpoint{5.083302in}{1.024556in}}{\pgfqpoint{5.088364in}{1.024556in}}%
\pgfpathquadraticcurveto{\pgfqpoint{5.102053in}{1.024556in}}{\pgfqpoint{5.108808in}{1.017459in}}%
\pgfpathquadraticcurveto{\pgfqpoint{5.115562in}{1.010380in}}{\pgfqpoint{5.115562in}{0.995970in}}%
\pgfpathlineto{\pgfqpoint{5.115562in}{0.995970in}}%
\pgfpathclose%
\pgfusepath{fill}%
\end{pgfscope}%
\begin{pgfscope}%
\pgfsetbuttcap%
\pgfsetroundjoin%
\definecolor{currentfill}{rgb}{0.309804,0.372549,0.419608}%
\pgfsetfillcolor{currentfill}%
\pgfsetlinewidth{0.000000pt}%
\definecolor{currentstroke}{rgb}{0.309804,0.372549,0.419608}%
\pgfsetstrokecolor{currentstroke}%
\pgfsetdash{}{0pt}%
\pgfpathmoveto{\pgfqpoint{5.237077in}{1.023043in}}%
\pgfpathlineto{\pgfqpoint{5.237077in}{0.960000in}}%
\pgfpathlineto{\pgfqpoint{5.226666in}{0.960000in}}%
\pgfpathlineto{\pgfqpoint{5.226666in}{1.014991in}}%
\pgfpathlineto{\pgfqpoint{5.198242in}{1.014991in}}%
\pgfpathlineto{\pgfqpoint{5.198242in}{0.960000in}}%
\pgfpathlineto{\pgfqpoint{5.187831in}{0.960000in}}%
\pgfpathlineto{\pgfqpoint{5.187831in}{1.014991in}}%
\pgfpathlineto{\pgfqpoint{5.177925in}{1.014991in}}%
\pgfpathlineto{\pgfqpoint{5.177925in}{1.023043in}}%
\pgfpathlineto{\pgfqpoint{5.187831in}{1.023043in}}%
\pgfpathlineto{\pgfqpoint{5.187831in}{1.027438in}}%
\pgfpathquadraticcurveto{\pgfqpoint{5.187831in}{1.037740in}}{\pgfqpoint{5.192695in}{1.042658in}}%
\pgfpathquadraticcurveto{\pgfqpoint{5.197576in}{1.047593in}}{\pgfqpoint{5.207645in}{1.047593in}}%
\pgfpathlineto{\pgfqpoint{5.218056in}{1.047593in}}%
\pgfpathlineto{\pgfqpoint{5.218056in}{1.038965in}}%
\pgfpathlineto{\pgfqpoint{5.208149in}{1.038965in}}%
\pgfpathquadraticcurveto{\pgfqpoint{5.202583in}{1.038965in}}{\pgfqpoint{5.200404in}{1.036714in}}%
\pgfpathquadraticcurveto{\pgfqpoint{5.198242in}{1.034462in}}{\pgfqpoint{5.198242in}{1.028608in}}%
\pgfpathlineto{\pgfqpoint{5.198242in}{1.023043in}}%
\pgfpathlineto{\pgfqpoint{5.237077in}{1.023043in}}%
\pgfpathlineto{\pgfqpoint{5.237077in}{1.023043in}}%
\pgfpathclose%
\pgfpathmoveto{\pgfqpoint{5.226666in}{1.047467in}}%
\pgfpathlineto{\pgfqpoint{5.237077in}{1.047467in}}%
\pgfpathlineto{\pgfqpoint{5.237077in}{1.034354in}}%
\pgfpathlineto{\pgfqpoint{5.226666in}{1.034354in}}%
\pgfpathlineto{\pgfqpoint{5.226666in}{1.047467in}}%
\pgfpathlineto{\pgfqpoint{5.226666in}{1.047467in}}%
\pgfpathclose%
\pgfpathmoveto{\pgfqpoint{5.258745in}{1.047593in}}%
\pgfpathlineto{\pgfqpoint{5.269102in}{1.047593in}}%
\pgfpathlineto{\pgfqpoint{5.269102in}{0.960000in}}%
\pgfpathlineto{\pgfqpoint{5.258745in}{0.960000in}}%
\pgfpathlineto{\pgfqpoint{5.258745in}{1.047593in}}%
\pgfpathlineto{\pgfqpoint{5.258745in}{1.047593in}}%
\pgfpathclose%
\pgfpathmoveto{\pgfqpoint{5.290771in}{1.047593in}}%
\pgfpathlineto{\pgfqpoint{5.301128in}{1.047593in}}%
\pgfpathlineto{\pgfqpoint{5.301128in}{0.960000in}}%
\pgfpathlineto{\pgfqpoint{5.290771in}{0.960000in}}%
\pgfpathlineto{\pgfqpoint{5.290771in}{1.047593in}}%
\pgfpathlineto{\pgfqpoint{5.290771in}{1.047593in}}%
\pgfpathclose%
\pgfpathmoveto{\pgfqpoint{5.376725in}{0.994115in}}%
\pgfpathlineto{\pgfqpoint{5.376725in}{0.989054in}}%
\pgfpathlineto{\pgfqpoint{5.329101in}{0.989054in}}%
\pgfpathquadraticcurveto{\pgfqpoint{5.329785in}{0.978354in}}{\pgfqpoint{5.335549in}{0.972753in}}%
\pgfpathquadraticcurveto{\pgfqpoint{5.341313in}{0.967151in}}{\pgfqpoint{5.351616in}{0.967151in}}%
\pgfpathquadraticcurveto{\pgfqpoint{5.357578in}{0.967151in}}{\pgfqpoint{5.363180in}{0.968610in}}%
\pgfpathquadraticcurveto{\pgfqpoint{5.368781in}{0.970069in}}{\pgfqpoint{5.374311in}{0.973005in}}%
\pgfpathlineto{\pgfqpoint{5.374311in}{0.963206in}}%
\pgfpathquadraticcurveto{\pgfqpoint{5.368727in}{0.960847in}}{\pgfqpoint{5.362874in}{0.959604in}}%
\pgfpathquadraticcurveto{\pgfqpoint{5.357020in}{0.958361in}}{\pgfqpoint{5.351004in}{0.958361in}}%
\pgfpathquadraticcurveto{\pgfqpoint{5.335909in}{0.958361in}}{\pgfqpoint{5.327101in}{0.967133in}}%
\pgfpathquadraticcurveto{\pgfqpoint{5.318293in}{0.975923in}}{\pgfqpoint{5.318293in}{0.990909in}}%
\pgfpathquadraticcurveto{\pgfqpoint{5.318293in}{1.006381in}}{\pgfqpoint{5.326651in}{1.015459in}}%
\pgfpathquadraticcurveto{\pgfqpoint{5.335009in}{1.024556in}}{\pgfqpoint{5.349202in}{1.024556in}}%
\pgfpathquadraticcurveto{\pgfqpoint{5.361919in}{1.024556in}}{\pgfqpoint{5.369322in}{1.016360in}}%
\pgfpathquadraticcurveto{\pgfqpoint{5.376725in}{1.008183in}}{\pgfqpoint{5.376725in}{0.994115in}}%
\pgfpathlineto{\pgfqpoint{5.376725in}{0.994115in}}%
\pgfpathclose%
\pgfpathmoveto{\pgfqpoint{5.366368in}{0.997159in}}%
\pgfpathquadraticcurveto{\pgfqpoint{5.366260in}{1.005643in}}{\pgfqpoint{5.361613in}{1.010704in}}%
\pgfpathquadraticcurveto{\pgfqpoint{5.356966in}{1.015784in}}{\pgfqpoint{5.349310in}{1.015784in}}%
\pgfpathquadraticcurveto{\pgfqpoint{5.340647in}{1.015784in}}{\pgfqpoint{5.335441in}{1.010884in}}%
\pgfpathquadraticcurveto{\pgfqpoint{5.330235in}{1.005985in}}{\pgfqpoint{5.329443in}{0.997087in}}%
\pgfpathlineto{\pgfqpoint{5.366368in}{0.997159in}}%
\pgfpathlineto{\pgfqpoint{5.366368in}{0.997159in}}%
\pgfpathclose%
\pgfpathmoveto{\pgfqpoint{5.435210in}{1.013478in}}%
\pgfpathlineto{\pgfqpoint{5.435210in}{1.047593in}}%
\pgfpathlineto{\pgfqpoint{5.445567in}{1.047593in}}%
\pgfpathlineto{\pgfqpoint{5.445567in}{0.960000in}}%
\pgfpathlineto{\pgfqpoint{5.435210in}{0.960000in}}%
\pgfpathlineto{\pgfqpoint{5.435210in}{0.969456in}}%
\pgfpathquadraticcurveto{\pgfqpoint{5.431950in}{0.963837in}}{\pgfqpoint{5.426961in}{0.961099in}}%
\pgfpathquadraticcurveto{\pgfqpoint{5.421989in}{0.958361in}}{\pgfqpoint{5.415001in}{0.958361in}}%
\pgfpathquadraticcurveto{\pgfqpoint{5.403581in}{0.958361in}}{\pgfqpoint{5.396394in}{0.967475in}}%
\pgfpathquadraticcurveto{\pgfqpoint{5.389225in}{0.976607in}}{\pgfqpoint{5.389225in}{0.991467in}}%
\pgfpathquadraticcurveto{\pgfqpoint{5.389225in}{1.006327in}}{\pgfqpoint{5.396394in}{1.015441in}}%
\pgfpathquadraticcurveto{\pgfqpoint{5.403581in}{1.024556in}}{\pgfqpoint{5.415001in}{1.024556in}}%
\pgfpathquadraticcurveto{\pgfqpoint{5.421989in}{1.024556in}}{\pgfqpoint{5.426961in}{1.021818in}}%
\pgfpathquadraticcurveto{\pgfqpoint{5.431950in}{1.019098in}}{\pgfqpoint{5.435210in}{1.013478in}}%
\pgfpathlineto{\pgfqpoint{5.435210in}{1.013478in}}%
\pgfpathclose%
\pgfpathmoveto{\pgfqpoint{5.399925in}{0.991467in}}%
\pgfpathquadraticcurveto{\pgfqpoint{5.399925in}{0.980048in}}{\pgfqpoint{5.404626in}{0.973545in}}%
\pgfpathquadraticcurveto{\pgfqpoint{5.409327in}{0.967043in}}{\pgfqpoint{5.417540in}{0.967043in}}%
\pgfpathquadraticcurveto{\pgfqpoint{5.425754in}{0.967043in}}{\pgfqpoint{5.430473in}{0.973545in}}%
\pgfpathquadraticcurveto{\pgfqpoint{5.435210in}{0.980048in}}{\pgfqpoint{5.435210in}{0.991467in}}%
\pgfpathquadraticcurveto{\pgfqpoint{5.435210in}{1.002887in}}{\pgfqpoint{5.430473in}{1.009389in}}%
\pgfpathquadraticcurveto{\pgfqpoint{5.425754in}{1.015892in}}{\pgfqpoint{5.417540in}{1.015892in}}%
\pgfpathquadraticcurveto{\pgfqpoint{5.409327in}{1.015892in}}{\pgfqpoint{5.404626in}{1.009389in}}%
\pgfpathquadraticcurveto{\pgfqpoint{5.399925in}{1.002887in}}{\pgfqpoint{5.399925in}{0.991467in}}%
\pgfpathlineto{\pgfqpoint{5.399925in}{0.991467in}}%
\pgfpathclose%
\pgfusepath{fill}%
\end{pgfscope}%
\begin{pgfscope}%
\pgfsetbuttcap%
\pgfsetroundjoin%
\definecolor{currentfill}{rgb}{0.309804,0.372549,0.419608}%
\pgfsetfillcolor{currentfill}%
\pgfsetlinewidth{0.000000pt}%
\definecolor{currentstroke}{rgb}{0.309804,0.372549,0.419608}%
\pgfsetstrokecolor{currentstroke}%
\pgfsetdash{}{0pt}%
\pgfpathmoveto{\pgfqpoint{5.575115in}{1.010938in}}%
\pgfpathquadraticcurveto{\pgfqpoint{5.579005in}{1.017927in}}{\pgfqpoint{5.584409in}{1.021241in}}%
\pgfpathquadraticcurveto{\pgfqpoint{5.589813in}{1.024556in}}{\pgfqpoint{5.597126in}{1.024556in}}%
\pgfpathquadraticcurveto{\pgfqpoint{5.606978in}{1.024556in}}{\pgfqpoint{5.612328in}{1.017657in}}%
\pgfpathquadraticcurveto{\pgfqpoint{5.617677in}{1.010776in}}{\pgfqpoint{5.617677in}{0.998060in}}%
\pgfpathlineto{\pgfqpoint{5.617677in}{0.960000in}}%
\pgfpathlineto{\pgfqpoint{5.607266in}{0.960000in}}%
\pgfpathlineto{\pgfqpoint{5.607266in}{0.997717in}}%
\pgfpathquadraticcurveto{\pgfqpoint{5.607266in}{1.006778in}}{\pgfqpoint{5.604042in}{1.011155in}}%
\pgfpathquadraticcurveto{\pgfqpoint{5.600836in}{1.015549in}}{\pgfqpoint{5.594262in}{1.015549in}}%
\pgfpathquadraticcurveto{\pgfqpoint{5.586210in}{1.015549in}}{\pgfqpoint{5.581527in}{1.010200in}}%
\pgfpathquadraticcurveto{\pgfqpoint{5.576862in}{1.004868in}}{\pgfqpoint{5.576862in}{0.995628in}}%
\pgfpathlineto{\pgfqpoint{5.576862in}{0.960000in}}%
\pgfpathlineto{\pgfqpoint{5.566451in}{0.960000in}}%
\pgfpathlineto{\pgfqpoint{5.566451in}{0.997717in}}%
\pgfpathquadraticcurveto{\pgfqpoint{5.566451in}{1.006832in}}{\pgfqpoint{5.563245in}{1.011191in}}%
\pgfpathquadraticcurveto{\pgfqpoint{5.560039in}{1.015549in}}{\pgfqpoint{5.553338in}{1.015549in}}%
\pgfpathquadraticcurveto{\pgfqpoint{5.545395in}{1.015549in}}{\pgfqpoint{5.540711in}{1.010182in}}%
\pgfpathquadraticcurveto{\pgfqpoint{5.536046in}{1.004814in}}{\pgfqpoint{5.536046in}{0.995628in}}%
\pgfpathlineto{\pgfqpoint{5.536046in}{0.960000in}}%
\pgfpathlineto{\pgfqpoint{5.525635in}{0.960000in}}%
\pgfpathlineto{\pgfqpoint{5.525635in}{1.023043in}}%
\pgfpathlineto{\pgfqpoint{5.536046in}{1.023043in}}%
\pgfpathlineto{\pgfqpoint{5.536046in}{1.013244in}}%
\pgfpathquadraticcurveto{\pgfqpoint{5.539595in}{1.019044in}}{\pgfqpoint{5.544548in}{1.021800in}}%
\pgfpathquadraticcurveto{\pgfqpoint{5.549501in}{1.024556in}}{\pgfqpoint{5.556310in}{1.024556in}}%
\pgfpathquadraticcurveto{\pgfqpoint{5.563191in}{1.024556in}}{\pgfqpoint{5.568000in}{1.021061in}}%
\pgfpathquadraticcurveto{\pgfqpoint{5.572809in}{1.017585in}}{\pgfqpoint{5.575115in}{1.010938in}}%
\pgfpathlineto{\pgfqpoint{5.575115in}{1.010938in}}%
\pgfpathclose%
\pgfpathmoveto{\pgfqpoint{5.666977in}{0.991683in}}%
\pgfpathquadraticcurveto{\pgfqpoint{5.654422in}{0.991683in}}{\pgfqpoint{5.649577in}{0.988819in}}%
\pgfpathquadraticcurveto{\pgfqpoint{5.644732in}{0.985956in}}{\pgfqpoint{5.644732in}{0.979021in}}%
\pgfpathquadraticcurveto{\pgfqpoint{5.644732in}{0.973509in}}{\pgfqpoint{5.648370in}{0.970267in}}%
\pgfpathquadraticcurveto{\pgfqpoint{5.652009in}{0.967043in}}{\pgfqpoint{5.658241in}{0.967043in}}%
\pgfpathquadraticcurveto{\pgfqpoint{5.666869in}{0.967043in}}{\pgfqpoint{5.672074in}{0.973149in}}%
\pgfpathquadraticcurveto{\pgfqpoint{5.677280in}{0.979255in}}{\pgfqpoint{5.677280in}{0.989378in}}%
\pgfpathlineto{\pgfqpoint{5.677280in}{0.991683in}}%
\pgfpathlineto{\pgfqpoint{5.666977in}{0.991683in}}%
\pgfpathlineto{\pgfqpoint{5.666977in}{0.991683in}}%
\pgfpathclose%
\pgfpathmoveto{\pgfqpoint{5.687637in}{0.995970in}}%
\pgfpathlineto{\pgfqpoint{5.687637in}{0.960000in}}%
\pgfpathlineto{\pgfqpoint{5.677280in}{0.960000in}}%
\pgfpathlineto{\pgfqpoint{5.677280in}{0.969564in}}%
\pgfpathquadraticcurveto{\pgfqpoint{5.673731in}{0.963837in}}{\pgfqpoint{5.668436in}{0.961099in}}%
\pgfpathquadraticcurveto{\pgfqpoint{5.663140in}{0.958361in}}{\pgfqpoint{5.655485in}{0.958361in}}%
\pgfpathquadraticcurveto{\pgfqpoint{5.645812in}{0.958361in}}{\pgfqpoint{5.640085in}{0.963801in}}%
\pgfpathquadraticcurveto{\pgfqpoint{5.634375in}{0.969240in}}{\pgfqpoint{5.634375in}{0.978354in}}%
\pgfpathquadraticcurveto{\pgfqpoint{5.634375in}{0.988982in}}{\pgfqpoint{5.641489in}{0.994385in}}%
\pgfpathquadraticcurveto{\pgfqpoint{5.648622in}{0.999789in}}{\pgfqpoint{5.662744in}{0.999789in}}%
\pgfpathlineto{\pgfqpoint{5.677280in}{0.999789in}}%
\pgfpathlineto{\pgfqpoint{5.677280in}{1.000816in}}%
\pgfpathquadraticcurveto{\pgfqpoint{5.677280in}{1.007966in}}{\pgfqpoint{5.672578in}{1.011875in}}%
\pgfpathquadraticcurveto{\pgfqpoint{5.667877in}{1.015784in}}{\pgfqpoint{5.659376in}{1.015784in}}%
\pgfpathquadraticcurveto{\pgfqpoint{5.653972in}{1.015784in}}{\pgfqpoint{5.648838in}{1.014487in}}%
\pgfpathquadraticcurveto{\pgfqpoint{5.643723in}{1.013190in}}{\pgfqpoint{5.639004in}{1.010596in}}%
\pgfpathlineto{\pgfqpoint{5.639004in}{1.020179in}}%
\pgfpathquadraticcurveto{\pgfqpoint{5.644678in}{1.022376in}}{\pgfqpoint{5.650027in}{1.023457in}}%
\pgfpathquadraticcurveto{\pgfqpoint{5.655377in}{1.024556in}}{\pgfqpoint{5.660438in}{1.024556in}}%
\pgfpathquadraticcurveto{\pgfqpoint{5.674128in}{1.024556in}}{\pgfqpoint{5.680882in}{1.017459in}}%
\pgfpathquadraticcurveto{\pgfqpoint{5.687637in}{1.010380in}}{\pgfqpoint{5.687637in}{0.995970in}}%
\pgfpathlineto{\pgfqpoint{5.687637in}{0.995970in}}%
\pgfpathclose%
\pgfpathmoveto{\pgfqpoint{5.745492in}{1.013370in}}%
\pgfpathquadraticcurveto{\pgfqpoint{5.743744in}{1.014379in}}{\pgfqpoint{5.741691in}{1.014847in}}%
\pgfpathquadraticcurveto{\pgfqpoint{5.739638in}{1.015333in}}{\pgfqpoint{5.737170in}{1.015333in}}%
\pgfpathquadraticcurveto{\pgfqpoint{5.728380in}{1.015333in}}{\pgfqpoint{5.723679in}{1.009623in}}%
\pgfpathquadraticcurveto{\pgfqpoint{5.718978in}{1.003914in}}{\pgfqpoint{5.718978in}{0.993214in}}%
\pgfpathlineto{\pgfqpoint{5.718978in}{0.960000in}}%
\pgfpathlineto{\pgfqpoint{5.708567in}{0.960000in}}%
\pgfpathlineto{\pgfqpoint{5.708567in}{1.023043in}}%
\pgfpathlineto{\pgfqpoint{5.718978in}{1.023043in}}%
\pgfpathlineto{\pgfqpoint{5.718978in}{1.013244in}}%
\pgfpathquadraticcurveto{\pgfqpoint{5.722256in}{1.018990in}}{\pgfqpoint{5.727480in}{1.021764in}}%
\pgfpathquadraticcurveto{\pgfqpoint{5.732721in}{1.024556in}}{\pgfqpoint{5.740214in}{1.024556in}}%
\pgfpathquadraticcurveto{\pgfqpoint{5.741277in}{1.024556in}}{\pgfqpoint{5.742574in}{1.024411in}}%
\pgfpathquadraticcurveto{\pgfqpoint{5.743871in}{1.024285in}}{\pgfqpoint{5.745438in}{1.023997in}}%
\pgfpathlineto{\pgfqpoint{5.745492in}{1.013370in}}%
\pgfpathlineto{\pgfqpoint{5.745492in}{1.013370in}}%
\pgfpathclose%
\pgfpathmoveto{\pgfqpoint{5.755957in}{1.047593in}}%
\pgfpathlineto{\pgfqpoint{5.766368in}{1.047593in}}%
\pgfpathlineto{\pgfqpoint{5.766368in}{0.995862in}}%
\pgfpathlineto{\pgfqpoint{5.797277in}{1.023043in}}%
\pgfpathlineto{\pgfqpoint{5.810498in}{1.023043in}}%
\pgfpathlineto{\pgfqpoint{5.777067in}{0.993557in}}%
\pgfpathlineto{\pgfqpoint{5.811920in}{0.960000in}}%
\pgfpathlineto{\pgfqpoint{5.798411in}{0.960000in}}%
\pgfpathlineto{\pgfqpoint{5.766368in}{0.990783in}}%
\pgfpathlineto{\pgfqpoint{5.766368in}{0.960000in}}%
\pgfpathlineto{\pgfqpoint{5.755957in}{0.960000in}}%
\pgfpathlineto{\pgfqpoint{5.755957in}{1.047593in}}%
\pgfpathlineto{\pgfqpoint{5.755957in}{1.047593in}}%
\pgfpathclose%
\pgfpathmoveto{\pgfqpoint{5.872928in}{0.994115in}}%
\pgfpathlineto{\pgfqpoint{5.872928in}{0.989054in}}%
\pgfpathlineto{\pgfqpoint{5.825304in}{0.989054in}}%
\pgfpathquadraticcurveto{\pgfqpoint{5.825988in}{0.978354in}}{\pgfqpoint{5.831752in}{0.972753in}}%
\pgfpathquadraticcurveto{\pgfqpoint{5.837516in}{0.967151in}}{\pgfqpoint{5.847819in}{0.967151in}}%
\pgfpathquadraticcurveto{\pgfqpoint{5.853781in}{0.967151in}}{\pgfqpoint{5.859382in}{0.968610in}}%
\pgfpathquadraticcurveto{\pgfqpoint{5.864984in}{0.970069in}}{\pgfqpoint{5.870514in}{0.973005in}}%
\pgfpathlineto{\pgfqpoint{5.870514in}{0.963206in}}%
\pgfpathquadraticcurveto{\pgfqpoint{5.864930in}{0.960847in}}{\pgfqpoint{5.859076in}{0.959604in}}%
\pgfpathquadraticcurveto{\pgfqpoint{5.853222in}{0.958361in}}{\pgfqpoint{5.847206in}{0.958361in}}%
\pgfpathquadraticcurveto{\pgfqpoint{5.832112in}{0.958361in}}{\pgfqpoint{5.823304in}{0.967133in}}%
\pgfpathquadraticcurveto{\pgfqpoint{5.814496in}{0.975923in}}{\pgfqpoint{5.814496in}{0.990909in}}%
\pgfpathquadraticcurveto{\pgfqpoint{5.814496in}{1.006381in}}{\pgfqpoint{5.822854in}{1.015459in}}%
\pgfpathquadraticcurveto{\pgfqpoint{5.831211in}{1.024556in}}{\pgfqpoint{5.845405in}{1.024556in}}%
\pgfpathquadraticcurveto{\pgfqpoint{5.858122in}{1.024556in}}{\pgfqpoint{5.865525in}{1.016360in}}%
\pgfpathquadraticcurveto{\pgfqpoint{5.872928in}{1.008183in}}{\pgfqpoint{5.872928in}{0.994115in}}%
\pgfpathlineto{\pgfqpoint{5.872928in}{0.994115in}}%
\pgfpathclose%
\pgfpathmoveto{\pgfqpoint{5.862571in}{0.997159in}}%
\pgfpathquadraticcurveto{\pgfqpoint{5.862463in}{1.005643in}}{\pgfqpoint{5.857815in}{1.010704in}}%
\pgfpathquadraticcurveto{\pgfqpoint{5.853168in}{1.015784in}}{\pgfqpoint{5.845513in}{1.015784in}}%
\pgfpathquadraticcurveto{\pgfqpoint{5.836849in}{1.015784in}}{\pgfqpoint{5.831644in}{1.010884in}}%
\pgfpathquadraticcurveto{\pgfqpoint{5.826438in}{1.005985in}}{\pgfqpoint{5.825646in}{0.997087in}}%
\pgfpathlineto{\pgfqpoint{5.862571in}{0.997159in}}%
\pgfpathlineto{\pgfqpoint{5.862571in}{0.997159in}}%
\pgfpathclose%
\pgfpathmoveto{\pgfqpoint{5.926460in}{1.013370in}}%
\pgfpathquadraticcurveto{\pgfqpoint{5.924713in}{1.014379in}}{\pgfqpoint{5.922659in}{1.014847in}}%
\pgfpathquadraticcurveto{\pgfqpoint{5.920606in}{1.015333in}}{\pgfqpoint{5.918138in}{1.015333in}}%
\pgfpathquadraticcurveto{\pgfqpoint{5.909348in}{1.015333in}}{\pgfqpoint{5.904647in}{1.009623in}}%
\pgfpathquadraticcurveto{\pgfqpoint{5.899946in}{1.003914in}}{\pgfqpoint{5.899946in}{0.993214in}}%
\pgfpathlineto{\pgfqpoint{5.899946in}{0.960000in}}%
\pgfpathlineto{\pgfqpoint{5.889535in}{0.960000in}}%
\pgfpathlineto{\pgfqpoint{5.889535in}{1.023043in}}%
\pgfpathlineto{\pgfqpoint{5.899946in}{1.023043in}}%
\pgfpathlineto{\pgfqpoint{5.899946in}{1.013244in}}%
\pgfpathquadraticcurveto{\pgfqpoint{5.903224in}{1.018990in}}{\pgfqpoint{5.908448in}{1.021764in}}%
\pgfpathquadraticcurveto{\pgfqpoint{5.913689in}{1.024556in}}{\pgfqpoint{5.921182in}{1.024556in}}%
\pgfpathquadraticcurveto{\pgfqpoint{5.922245in}{1.024556in}}{\pgfqpoint{5.923542in}{1.024411in}}%
\pgfpathquadraticcurveto{\pgfqpoint{5.924839in}{1.024285in}}{\pgfqpoint{5.926406in}{1.023997in}}%
\pgfpathlineto{\pgfqpoint{5.926460in}{1.013370in}}%
\pgfpathlineto{\pgfqpoint{5.926460in}{1.013370in}}%
\pgfpathclose%
\pgfpathmoveto{\pgfqpoint{5.928207in}{0.974302in}}%
\pgfpathlineto{\pgfqpoint{5.940095in}{0.974302in}}%
\pgfpathlineto{\pgfqpoint{5.940095in}{0.960000in}}%
\pgfpathlineto{\pgfqpoint{5.928207in}{0.960000in}}%
\pgfpathlineto{\pgfqpoint{5.928207in}{0.974302in}}%
\pgfpathlineto{\pgfqpoint{5.928207in}{0.974302in}}%
\pgfpathclose%
\pgfusepath{fill}%
\end{pgfscope}%
\begin{pgfscope}%
\pgfsetbuttcap%
\pgfsetroundjoin%
\definecolor{currentfill}{rgb}{0.309804,0.372549,0.419608}%
\pgfsetfillcolor{currentfill}%
\pgfsetlinewidth{0.000000pt}%
\definecolor{currentstroke}{rgb}{0.309804,0.372549,0.419608}%
\pgfsetstrokecolor{currentstroke}%
\pgfsetdash{}{0pt}%
\pgfpathmoveto{\pgfqpoint{6.043053in}{1.034696in}}%
\pgfpathlineto{\pgfqpoint{6.043053in}{0.969348in}}%
\pgfpathlineto{\pgfqpoint{6.056797in}{0.969348in}}%
\pgfpathquadraticcurveto{\pgfqpoint{6.074196in}{0.969348in}}{\pgfqpoint{6.082266in}{0.977220in}}%
\pgfpathquadraticcurveto{\pgfqpoint{6.090335in}{0.985109in}}{\pgfqpoint{6.090335in}{1.002112in}}%
\pgfpathquadraticcurveto{\pgfqpoint{6.090335in}{1.018990in}}{\pgfqpoint{6.082266in}{1.026843in}}%
\pgfpathquadraticcurveto{\pgfqpoint{6.074196in}{1.034696in}}{\pgfqpoint{6.056797in}{1.034696in}}%
\pgfpathlineto{\pgfqpoint{6.043053in}{1.034696in}}%
\pgfpathlineto{\pgfqpoint{6.043053in}{1.034696in}}%
\pgfpathclose%
\pgfpathmoveto{\pgfqpoint{6.031688in}{1.044045in}}%
\pgfpathlineto{\pgfqpoint{6.055049in}{1.044045in}}%
\pgfpathquadraticcurveto{\pgfqpoint{6.079474in}{1.044045in}}{\pgfqpoint{6.090894in}{1.033886in}}%
\pgfpathquadraticcurveto{\pgfqpoint{6.102331in}{1.023727in}}{\pgfqpoint{6.102331in}{1.002112in}}%
\pgfpathquadraticcurveto{\pgfqpoint{6.102331in}{0.980372in}}{\pgfqpoint{6.090840in}{0.970177in}}%
\pgfpathquadraticcurveto{\pgfqpoint{6.079366in}{0.960000in}}{\pgfqpoint{6.055049in}{0.960000in}}%
\pgfpathlineto{\pgfqpoint{6.031688in}{0.960000in}}%
\pgfpathlineto{\pgfqpoint{6.031688in}{1.044045in}}%
\pgfpathlineto{\pgfqpoint{6.031688in}{1.044045in}}%
\pgfpathclose%
\pgfpathmoveto{\pgfqpoint{6.120001in}{1.023043in}}%
\pgfpathlineto{\pgfqpoint{6.130358in}{1.023043in}}%
\pgfpathlineto{\pgfqpoint{6.130358in}{0.960000in}}%
\pgfpathlineto{\pgfqpoint{6.120001in}{0.960000in}}%
\pgfpathlineto{\pgfqpoint{6.120001in}{1.023043in}}%
\pgfpathlineto{\pgfqpoint{6.120001in}{1.023043in}}%
\pgfpathclose%
\pgfpathmoveto{\pgfqpoint{6.120001in}{1.047593in}}%
\pgfpathlineto{\pgfqpoint{6.130358in}{1.047593in}}%
\pgfpathlineto{\pgfqpoint{6.130358in}{1.034462in}}%
\pgfpathlineto{\pgfqpoint{6.120001in}{1.034462in}}%
\pgfpathlineto{\pgfqpoint{6.120001in}{1.047593in}}%
\pgfpathlineto{\pgfqpoint{6.120001in}{1.047593in}}%
\pgfpathclose%
\pgfpathmoveto{\pgfqpoint{6.180684in}{0.991683in}}%
\pgfpathquadraticcurveto{\pgfqpoint{6.168130in}{0.991683in}}{\pgfqpoint{6.163284in}{0.988819in}}%
\pgfpathquadraticcurveto{\pgfqpoint{6.158439in}{0.985956in}}{\pgfqpoint{6.158439in}{0.979021in}}%
\pgfpathquadraticcurveto{\pgfqpoint{6.158439in}{0.973509in}}{\pgfqpoint{6.162078in}{0.970267in}}%
\pgfpathquadraticcurveto{\pgfqpoint{6.165716in}{0.967043in}}{\pgfqpoint{6.171948in}{0.967043in}}%
\pgfpathquadraticcurveto{\pgfqpoint{6.180576in}{0.967043in}}{\pgfqpoint{6.185782in}{0.973149in}}%
\pgfpathquadraticcurveto{\pgfqpoint{6.190987in}{0.979255in}}{\pgfqpoint{6.190987in}{0.989378in}}%
\pgfpathlineto{\pgfqpoint{6.190987in}{0.991683in}}%
\pgfpathlineto{\pgfqpoint{6.180684in}{0.991683in}}%
\pgfpathlineto{\pgfqpoint{6.180684in}{0.991683in}}%
\pgfpathclose%
\pgfpathmoveto{\pgfqpoint{6.201344in}{0.995970in}}%
\pgfpathlineto{\pgfqpoint{6.201344in}{0.960000in}}%
\pgfpathlineto{\pgfqpoint{6.190987in}{0.960000in}}%
\pgfpathlineto{\pgfqpoint{6.190987in}{0.969564in}}%
\pgfpathquadraticcurveto{\pgfqpoint{6.187439in}{0.963837in}}{\pgfqpoint{6.182143in}{0.961099in}}%
\pgfpathquadraticcurveto{\pgfqpoint{6.176848in}{0.958361in}}{\pgfqpoint{6.169192in}{0.958361in}}%
\pgfpathquadraticcurveto{\pgfqpoint{6.159520in}{0.958361in}}{\pgfqpoint{6.153792in}{0.963801in}}%
\pgfpathquadraticcurveto{\pgfqpoint{6.148082in}{0.969240in}}{\pgfqpoint{6.148082in}{0.978354in}}%
\pgfpathquadraticcurveto{\pgfqpoint{6.148082in}{0.988982in}}{\pgfqpoint{6.155197in}{0.994385in}}%
\pgfpathquadraticcurveto{\pgfqpoint{6.162330in}{0.999789in}}{\pgfqpoint{6.176451in}{0.999789in}}%
\pgfpathlineto{\pgfqpoint{6.190987in}{0.999789in}}%
\pgfpathlineto{\pgfqpoint{6.190987in}{1.000816in}}%
\pgfpathquadraticcurveto{\pgfqpoint{6.190987in}{1.007966in}}{\pgfqpoint{6.186286in}{1.011875in}}%
\pgfpathquadraticcurveto{\pgfqpoint{6.181585in}{1.015784in}}{\pgfqpoint{6.173083in}{1.015784in}}%
\pgfpathquadraticcurveto{\pgfqpoint{6.167679in}{1.015784in}}{\pgfqpoint{6.162546in}{1.014487in}}%
\pgfpathquadraticcurveto{\pgfqpoint{6.157431in}{1.013190in}}{\pgfqpoint{6.152711in}{1.010596in}}%
\pgfpathlineto{\pgfqpoint{6.152711in}{1.020179in}}%
\pgfpathquadraticcurveto{\pgfqpoint{6.158385in}{1.022376in}}{\pgfqpoint{6.163735in}{1.023457in}}%
\pgfpathquadraticcurveto{\pgfqpoint{6.169084in}{1.024556in}}{\pgfqpoint{6.174146in}{1.024556in}}%
\pgfpathquadraticcurveto{\pgfqpoint{6.187835in}{1.024556in}}{\pgfqpoint{6.194590in}{1.017459in}}%
\pgfpathquadraticcurveto{\pgfqpoint{6.201344in}{1.010380in}}{\pgfqpoint{6.201344in}{0.995970in}}%
\pgfpathlineto{\pgfqpoint{6.201344in}{0.995970in}}%
\pgfpathclose%
\pgfpathmoveto{\pgfqpoint{6.271754in}{1.010938in}}%
\pgfpathquadraticcurveto{\pgfqpoint{6.275644in}{1.017927in}}{\pgfqpoint{6.281048in}{1.021241in}}%
\pgfpathquadraticcurveto{\pgfqpoint{6.286452in}{1.024556in}}{\pgfqpoint{6.293765in}{1.024556in}}%
\pgfpathquadraticcurveto{\pgfqpoint{6.303617in}{1.024556in}}{\pgfqpoint{6.308967in}{1.017657in}}%
\pgfpathquadraticcurveto{\pgfqpoint{6.314316in}{1.010776in}}{\pgfqpoint{6.314316in}{0.998060in}}%
\pgfpathlineto{\pgfqpoint{6.314316in}{0.960000in}}%
\pgfpathlineto{\pgfqpoint{6.303905in}{0.960000in}}%
\pgfpathlineto{\pgfqpoint{6.303905in}{0.997717in}}%
\pgfpathquadraticcurveto{\pgfqpoint{6.303905in}{1.006778in}}{\pgfqpoint{6.300681in}{1.011155in}}%
\pgfpathquadraticcurveto{\pgfqpoint{6.297475in}{1.015549in}}{\pgfqpoint{6.290901in}{1.015549in}}%
\pgfpathquadraticcurveto{\pgfqpoint{6.282849in}{1.015549in}}{\pgfqpoint{6.278166in}{1.010200in}}%
\pgfpathquadraticcurveto{\pgfqpoint{6.273501in}{1.004868in}}{\pgfqpoint{6.273501in}{0.995628in}}%
\pgfpathlineto{\pgfqpoint{6.273501in}{0.960000in}}%
\pgfpathlineto{\pgfqpoint{6.263090in}{0.960000in}}%
\pgfpathlineto{\pgfqpoint{6.263090in}{0.997717in}}%
\pgfpathquadraticcurveto{\pgfqpoint{6.263090in}{1.006832in}}{\pgfqpoint{6.259884in}{1.011191in}}%
\pgfpathquadraticcurveto{\pgfqpoint{6.256678in}{1.015549in}}{\pgfqpoint{6.249977in}{1.015549in}}%
\pgfpathquadraticcurveto{\pgfqpoint{6.242034in}{1.015549in}}{\pgfqpoint{6.237350in}{1.010182in}}%
\pgfpathquadraticcurveto{\pgfqpoint{6.232685in}{1.004814in}}{\pgfqpoint{6.232685in}{0.995628in}}%
\pgfpathlineto{\pgfqpoint{6.232685in}{0.960000in}}%
\pgfpathlineto{\pgfqpoint{6.222274in}{0.960000in}}%
\pgfpathlineto{\pgfqpoint{6.222274in}{1.023043in}}%
\pgfpathlineto{\pgfqpoint{6.232685in}{1.023043in}}%
\pgfpathlineto{\pgfqpoint{6.232685in}{1.013244in}}%
\pgfpathquadraticcurveto{\pgfqpoint{6.236234in}{1.019044in}}{\pgfqpoint{6.241187in}{1.021800in}}%
\pgfpathquadraticcurveto{\pgfqpoint{6.246140in}{1.024556in}}{\pgfqpoint{6.252949in}{1.024556in}}%
\pgfpathquadraticcurveto{\pgfqpoint{6.259830in}{1.024556in}}{\pgfqpoint{6.264639in}{1.021061in}}%
\pgfpathquadraticcurveto{\pgfqpoint{6.269448in}{1.017585in}}{\pgfqpoint{6.271754in}{1.010938in}}%
\pgfpathlineto{\pgfqpoint{6.271754in}{1.010938in}}%
\pgfpathclose%
\pgfpathmoveto{\pgfqpoint{6.359383in}{1.015784in}}%
\pgfpathquadraticcurveto{\pgfqpoint{6.351061in}{1.015784in}}{\pgfqpoint{6.346216in}{1.009281in}}%
\pgfpathquadraticcurveto{\pgfqpoint{6.341371in}{1.002779in}}{\pgfqpoint{6.341371in}{0.991467in}}%
\pgfpathquadraticcurveto{\pgfqpoint{6.341371in}{0.980156in}}{\pgfqpoint{6.346180in}{0.973653in}}%
\pgfpathquadraticcurveto{\pgfqpoint{6.351007in}{0.967151in}}{\pgfqpoint{6.359383in}{0.967151in}}%
\pgfpathquadraticcurveto{\pgfqpoint{6.367668in}{0.967151in}}{\pgfqpoint{6.372496in}{0.973671in}}%
\pgfpathquadraticcurveto{\pgfqpoint{6.377341in}{0.980210in}}{\pgfqpoint{6.377341in}{0.991467in}}%
\pgfpathquadraticcurveto{\pgfqpoint{6.377341in}{1.002671in}}{\pgfqpoint{6.372496in}{1.009227in}}%
\pgfpathquadraticcurveto{\pgfqpoint{6.367668in}{1.015784in}}{\pgfqpoint{6.359383in}{1.015784in}}%
\pgfpathlineto{\pgfqpoint{6.359383in}{1.015784in}}%
\pgfpathclose%
\pgfpathmoveto{\pgfqpoint{6.359383in}{1.024556in}}%
\pgfpathquadraticcurveto{\pgfqpoint{6.372892in}{1.024556in}}{\pgfqpoint{6.380601in}{1.015766in}}%
\pgfpathquadraticcurveto{\pgfqpoint{6.388328in}{1.006994in}}{\pgfqpoint{6.388328in}{0.991467in}}%
\pgfpathquadraticcurveto{\pgfqpoint{6.388328in}{0.975995in}}{\pgfqpoint{6.380601in}{0.967169in}}%
\pgfpathquadraticcurveto{\pgfqpoint{6.372892in}{0.958361in}}{\pgfqpoint{6.359383in}{0.958361in}}%
\pgfpathquadraticcurveto{\pgfqpoint{6.345820in}{0.958361in}}{\pgfqpoint{6.338128in}{0.967169in}}%
\pgfpathquadraticcurveto{\pgfqpoint{6.330455in}{0.975995in}}{\pgfqpoint{6.330455in}{0.991467in}}%
\pgfpathquadraticcurveto{\pgfqpoint{6.330455in}{1.006994in}}{\pgfqpoint{6.338128in}{1.015766in}}%
\pgfpathquadraticcurveto{\pgfqpoint{6.345820in}{1.024556in}}{\pgfqpoint{6.359383in}{1.024556in}}%
\pgfpathlineto{\pgfqpoint{6.359383in}{1.024556in}}%
\pgfpathclose%
\pgfpathmoveto{\pgfqpoint{6.457909in}{0.998060in}}%
\pgfpathlineto{\pgfqpoint{6.457909in}{0.960000in}}%
\pgfpathlineto{\pgfqpoint{6.447552in}{0.960000in}}%
\pgfpathlineto{\pgfqpoint{6.447552in}{0.997717in}}%
\pgfpathquadraticcurveto{\pgfqpoint{6.447552in}{1.006669in}}{\pgfqpoint{6.444058in}{1.011100in}}%
\pgfpathquadraticcurveto{\pgfqpoint{6.440564in}{1.015549in}}{\pgfqpoint{6.433593in}{1.015549in}}%
\pgfpathquadraticcurveto{\pgfqpoint{6.425199in}{1.015549in}}{\pgfqpoint{6.420354in}{1.010200in}}%
\pgfpathquadraticcurveto{\pgfqpoint{6.415509in}{1.004868in}}{\pgfqpoint{6.415509in}{0.995628in}}%
\pgfpathlineto{\pgfqpoint{6.415509in}{0.960000in}}%
\pgfpathlineto{\pgfqpoint{6.405098in}{0.960000in}}%
\pgfpathlineto{\pgfqpoint{6.405098in}{1.023043in}}%
\pgfpathlineto{\pgfqpoint{6.415509in}{1.023043in}}%
\pgfpathlineto{\pgfqpoint{6.415509in}{1.013244in}}%
\pgfpathquadraticcurveto{\pgfqpoint{6.419237in}{1.018936in}}{\pgfqpoint{6.424263in}{1.021746in}}%
\pgfpathquadraticcurveto{\pgfqpoint{6.429306in}{1.024556in}}{\pgfqpoint{6.435898in}{1.024556in}}%
\pgfpathquadraticcurveto{\pgfqpoint{6.446760in}{1.024556in}}{\pgfqpoint{6.452326in}{1.017837in}}%
\pgfpathquadraticcurveto{\pgfqpoint{6.457909in}{1.011118in}}{\pgfqpoint{6.457909in}{0.998060in}}%
\pgfpathlineto{\pgfqpoint{6.457909in}{0.998060in}}%
\pgfpathclose%
\pgfpathmoveto{\pgfqpoint{6.520033in}{1.013478in}}%
\pgfpathlineto{\pgfqpoint{6.520033in}{1.047593in}}%
\pgfpathlineto{\pgfqpoint{6.530390in}{1.047593in}}%
\pgfpathlineto{\pgfqpoint{6.530390in}{0.960000in}}%
\pgfpathlineto{\pgfqpoint{6.520033in}{0.960000in}}%
\pgfpathlineto{\pgfqpoint{6.520033in}{0.969456in}}%
\pgfpathquadraticcurveto{\pgfqpoint{6.516773in}{0.963837in}}{\pgfqpoint{6.511784in}{0.961099in}}%
\pgfpathquadraticcurveto{\pgfqpoint{6.506812in}{0.958361in}}{\pgfqpoint{6.499824in}{0.958361in}}%
\pgfpathquadraticcurveto{\pgfqpoint{6.488404in}{0.958361in}}{\pgfqpoint{6.481217in}{0.967475in}}%
\pgfpathquadraticcurveto{\pgfqpoint{6.474048in}{0.976607in}}{\pgfqpoint{6.474048in}{0.991467in}}%
\pgfpathquadraticcurveto{\pgfqpoint{6.474048in}{1.006327in}}{\pgfqpoint{6.481217in}{1.015441in}}%
\pgfpathquadraticcurveto{\pgfqpoint{6.488404in}{1.024556in}}{\pgfqpoint{6.499824in}{1.024556in}}%
\pgfpathquadraticcurveto{\pgfqpoint{6.506812in}{1.024556in}}{\pgfqpoint{6.511784in}{1.021818in}}%
\pgfpathquadraticcurveto{\pgfqpoint{6.516773in}{1.019098in}}{\pgfqpoint{6.520033in}{1.013478in}}%
\pgfpathlineto{\pgfqpoint{6.520033in}{1.013478in}}%
\pgfpathclose%
\pgfpathmoveto{\pgfqpoint{6.484747in}{0.991467in}}%
\pgfpathquadraticcurveto{\pgfqpoint{6.484747in}{0.980048in}}{\pgfqpoint{6.489449in}{0.973545in}}%
\pgfpathquadraticcurveto{\pgfqpoint{6.494150in}{0.967043in}}{\pgfqpoint{6.502363in}{0.967043in}}%
\pgfpathquadraticcurveto{\pgfqpoint{6.510577in}{0.967043in}}{\pgfqpoint{6.515296in}{0.973545in}}%
\pgfpathquadraticcurveto{\pgfqpoint{6.520033in}{0.980048in}}{\pgfqpoint{6.520033in}{0.991467in}}%
\pgfpathquadraticcurveto{\pgfqpoint{6.520033in}{1.002887in}}{\pgfqpoint{6.515296in}{1.009389in}}%
\pgfpathquadraticcurveto{\pgfqpoint{6.510577in}{1.015892in}}{\pgfqpoint{6.502363in}{1.015892in}}%
\pgfpathquadraticcurveto{\pgfqpoint{6.494150in}{1.015892in}}{\pgfqpoint{6.489449in}{1.009389in}}%
\pgfpathquadraticcurveto{\pgfqpoint{6.484747in}{1.002887in}}{\pgfqpoint{6.484747in}{0.991467in}}%
\pgfpathlineto{\pgfqpoint{6.484747in}{0.991467in}}%
\pgfpathclose%
\pgfpathmoveto{\pgfqpoint{6.591920in}{1.021187in}}%
\pgfpathlineto{\pgfqpoint{6.591920in}{1.011389in}}%
\pgfpathquadraticcurveto{\pgfqpoint{6.587543in}{1.013640in}}{\pgfqpoint{6.582806in}{1.014757in}}%
\pgfpathquadraticcurveto{\pgfqpoint{6.578086in}{1.015892in}}{\pgfqpoint{6.573007in}{1.015892in}}%
\pgfpathquadraticcurveto{\pgfqpoint{6.565298in}{1.015892in}}{\pgfqpoint{6.561443in}{1.013532in}}%
\pgfpathquadraticcurveto{\pgfqpoint{6.557589in}{1.011173in}}{\pgfqpoint{6.557589in}{1.006435in}}%
\pgfpathquadraticcurveto{\pgfqpoint{6.557589in}{1.002833in}}{\pgfqpoint{6.560344in}{1.000780in}}%
\pgfpathquadraticcurveto{\pgfqpoint{6.563100in}{0.998726in}}{\pgfqpoint{6.571440in}{0.996871in}}%
\pgfpathlineto{\pgfqpoint{6.574988in}{0.996078in}}%
\pgfpathquadraticcurveto{\pgfqpoint{6.586012in}{0.993719in}}{\pgfqpoint{6.590659in}{0.989414in}}%
\pgfpathquadraticcurveto{\pgfqpoint{6.595306in}{0.985109in}}{\pgfqpoint{6.595306in}{0.977400in}}%
\pgfpathquadraticcurveto{\pgfqpoint{6.595306in}{0.968610in}}{\pgfqpoint{6.588353in}{0.963476in}}%
\pgfpathquadraticcurveto{\pgfqpoint{6.581401in}{0.958361in}}{\pgfqpoint{6.569242in}{0.958361in}}%
\pgfpathquadraticcurveto{\pgfqpoint{6.564181in}{0.958361in}}{\pgfqpoint{6.558687in}{0.959352in}}%
\pgfpathquadraticcurveto{\pgfqpoint{6.553194in}{0.960342in}}{\pgfqpoint{6.547123in}{0.962306in}}%
\pgfpathlineto{\pgfqpoint{6.547123in}{0.973005in}}%
\pgfpathquadraticcurveto{\pgfqpoint{6.552869in}{0.970015in}}{\pgfqpoint{6.558435in}{0.968520in}}%
\pgfpathquadraticcurveto{\pgfqpoint{6.564001in}{0.967043in}}{\pgfqpoint{6.569477in}{0.967043in}}%
\pgfpathquadraticcurveto{\pgfqpoint{6.576789in}{0.967043in}}{\pgfqpoint{6.580716in}{0.969546in}}%
\pgfpathquadraticcurveto{\pgfqpoint{6.584661in}{0.972050in}}{\pgfqpoint{6.584661in}{0.976607in}}%
\pgfpathquadraticcurveto{\pgfqpoint{6.584661in}{0.980822in}}{\pgfqpoint{6.581815in}{0.983074in}}%
\pgfpathquadraticcurveto{\pgfqpoint{6.578987in}{0.985325in}}{\pgfqpoint{6.569350in}{0.987414in}}%
\pgfpathlineto{\pgfqpoint{6.565748in}{0.988261in}}%
\pgfpathquadraticcurveto{\pgfqpoint{6.556130in}{0.990278in}}{\pgfqpoint{6.551843in}{0.994475in}}%
\pgfpathquadraticcurveto{\pgfqpoint{6.547574in}{0.998672in}}{\pgfqpoint{6.547574in}{1.005985in}}%
\pgfpathquadraticcurveto{\pgfqpoint{6.547574in}{1.014883in}}{\pgfqpoint{6.553878in}{1.019710in}}%
\pgfpathquadraticcurveto{\pgfqpoint{6.560182in}{1.024556in}}{\pgfqpoint{6.571782in}{1.024556in}}%
\pgfpathquadraticcurveto{\pgfqpoint{6.577510in}{1.024556in}}{\pgfqpoint{6.582571in}{1.023709in}}%
\pgfpathquadraticcurveto{\pgfqpoint{6.587651in}{1.022880in}}{\pgfqpoint{6.591920in}{1.021187in}}%
\pgfpathlineto{\pgfqpoint{6.591920in}{1.021187in}}%
\pgfpathclose%
\pgfpathmoveto{\pgfqpoint{6.614435in}{0.974302in}}%
\pgfpathlineto{\pgfqpoint{6.626305in}{0.974302in}}%
\pgfpathlineto{\pgfqpoint{6.626305in}{0.960000in}}%
\pgfpathlineto{\pgfqpoint{6.614435in}{0.960000in}}%
\pgfpathlineto{\pgfqpoint{6.614435in}{0.974302in}}%
\pgfpathlineto{\pgfqpoint{6.614435in}{0.974302in}}%
\pgfpathclose%
\pgfpathmoveto{\pgfqpoint{6.614435in}{1.019602in}}%
\pgfpathlineto{\pgfqpoint{6.626305in}{1.019602in}}%
\pgfpathlineto{\pgfqpoint{6.626305in}{1.005319in}}%
\pgfpathlineto{\pgfqpoint{6.614435in}{1.005319in}}%
\pgfpathlineto{\pgfqpoint{6.614435in}{1.019602in}}%
\pgfpathlineto{\pgfqpoint{6.614435in}{1.019602in}}%
\pgfpathclose%
\pgfusepath{fill}%
\end{pgfscope}%
\begin{pgfscope}%
\pgfsetbuttcap%
\pgfsetroundjoin%
\definecolor{currentfill}{rgb}{0.309804,0.372549,0.419608}%
\pgfsetfillcolor{currentfill}%
\pgfsetlinewidth{0.000000pt}%
\definecolor{currentstroke}{rgb}{0.309804,0.372549,0.419608}%
\pgfsetstrokecolor{currentstroke}%
\pgfsetdash{}{0pt}%
\pgfpathmoveto{\pgfqpoint{6.766503in}{1.020629in}}%
\pgfpathlineto{\pgfqpoint{6.766503in}{1.010938in}}%
\pgfpathquadraticcurveto{\pgfqpoint{6.762108in}{1.013370in}}{\pgfqpoint{6.757695in}{1.014577in}}%
\pgfpathquadraticcurveto{\pgfqpoint{6.753282in}{1.015784in}}{\pgfqpoint{6.748779in}{1.015784in}}%
\pgfpathquadraticcurveto{\pgfqpoint{6.738693in}{1.015784in}}{\pgfqpoint{6.733109in}{1.009389in}}%
\pgfpathquadraticcurveto{\pgfqpoint{6.727543in}{1.003013in}}{\pgfqpoint{6.727543in}{0.991467in}}%
\pgfpathquadraticcurveto{\pgfqpoint{6.727543in}{0.979921in}}{\pgfqpoint{6.733109in}{0.973527in}}%
\pgfpathquadraticcurveto{\pgfqpoint{6.738693in}{0.967151in}}{\pgfqpoint{6.748779in}{0.967151in}}%
\pgfpathquadraticcurveto{\pgfqpoint{6.753282in}{0.967151in}}{\pgfqpoint{6.757695in}{0.968358in}}%
\pgfpathquadraticcurveto{\pgfqpoint{6.762108in}{0.969564in}}{\pgfqpoint{6.766503in}{0.971996in}}%
\pgfpathlineto{\pgfqpoint{6.766503in}{0.962414in}}%
\pgfpathquadraticcurveto{\pgfqpoint{6.762162in}{0.960396in}}{\pgfqpoint{6.757515in}{0.959388in}}%
\pgfpathquadraticcurveto{\pgfqpoint{6.752886in}{0.958361in}}{\pgfqpoint{6.747645in}{0.958361in}}%
\pgfpathquadraticcurveto{\pgfqpoint{6.733397in}{0.958361in}}{\pgfqpoint{6.725003in}{0.967313in}}%
\pgfpathquadraticcurveto{\pgfqpoint{6.716628in}{0.976265in}}{\pgfqpoint{6.716628in}{0.991467in}}%
\pgfpathquadraticcurveto{\pgfqpoint{6.716628in}{1.006886in}}{\pgfqpoint{6.725093in}{1.015712in}}%
\pgfpathquadraticcurveto{\pgfqpoint{6.733577in}{1.024556in}}{\pgfqpoint{6.748329in}{1.024556in}}%
\pgfpathquadraticcurveto{\pgfqpoint{6.753102in}{1.024556in}}{\pgfqpoint{6.757659in}{1.023565in}}%
\pgfpathquadraticcurveto{\pgfqpoint{6.762216in}{1.022592in}}{\pgfqpoint{6.766503in}{1.020629in}}%
\pgfpathlineto{\pgfqpoint{6.766503in}{1.020629in}}%
\pgfpathclose%
\pgfpathmoveto{\pgfqpoint{6.808940in}{1.015784in}}%
\pgfpathquadraticcurveto{\pgfqpoint{6.800618in}{1.015784in}}{\pgfqpoint{6.795773in}{1.009281in}}%
\pgfpathquadraticcurveto{\pgfqpoint{6.790928in}{1.002779in}}{\pgfqpoint{6.790928in}{0.991467in}}%
\pgfpathquadraticcurveto{\pgfqpoint{6.790928in}{0.980156in}}{\pgfqpoint{6.795737in}{0.973653in}}%
\pgfpathquadraticcurveto{\pgfqpoint{6.800564in}{0.967151in}}{\pgfqpoint{6.808940in}{0.967151in}}%
\pgfpathquadraticcurveto{\pgfqpoint{6.817226in}{0.967151in}}{\pgfqpoint{6.822053in}{0.973671in}}%
\pgfpathquadraticcurveto{\pgfqpoint{6.826898in}{0.980210in}}{\pgfqpoint{6.826898in}{0.991467in}}%
\pgfpathquadraticcurveto{\pgfqpoint{6.826898in}{1.002671in}}{\pgfqpoint{6.822053in}{1.009227in}}%
\pgfpathquadraticcurveto{\pgfqpoint{6.817226in}{1.015784in}}{\pgfqpoint{6.808940in}{1.015784in}}%
\pgfpathlineto{\pgfqpoint{6.808940in}{1.015784in}}%
\pgfpathclose%
\pgfpathmoveto{\pgfqpoint{6.808940in}{1.024556in}}%
\pgfpathquadraticcurveto{\pgfqpoint{6.822449in}{1.024556in}}{\pgfqpoint{6.830158in}{1.015766in}}%
\pgfpathquadraticcurveto{\pgfqpoint{6.837885in}{1.006994in}}{\pgfqpoint{6.837885in}{0.991467in}}%
\pgfpathquadraticcurveto{\pgfqpoint{6.837885in}{0.975995in}}{\pgfqpoint{6.830158in}{0.967169in}}%
\pgfpathquadraticcurveto{\pgfqpoint{6.822449in}{0.958361in}}{\pgfqpoint{6.808940in}{0.958361in}}%
\pgfpathquadraticcurveto{\pgfqpoint{6.795377in}{0.958361in}}{\pgfqpoint{6.787686in}{0.967169in}}%
\pgfpathquadraticcurveto{\pgfqpoint{6.780012in}{0.975995in}}{\pgfqpoint{6.780012in}{0.991467in}}%
\pgfpathquadraticcurveto{\pgfqpoint{6.780012in}{1.006994in}}{\pgfqpoint{6.787686in}{1.015766in}}%
\pgfpathquadraticcurveto{\pgfqpoint{6.795377in}{1.024556in}}{\pgfqpoint{6.808940in}{1.024556in}}%
\pgfpathlineto{\pgfqpoint{6.808940in}{1.024556in}}%
\pgfpathclose%
\pgfpathmoveto{\pgfqpoint{6.879476in}{1.015784in}}%
\pgfpathquadraticcurveto{\pgfqpoint{6.871154in}{1.015784in}}{\pgfqpoint{6.866309in}{1.009281in}}%
\pgfpathquadraticcurveto{\pgfqpoint{6.861463in}{1.002779in}}{\pgfqpoint{6.861463in}{0.991467in}}%
\pgfpathquadraticcurveto{\pgfqpoint{6.861463in}{0.980156in}}{\pgfqpoint{6.866273in}{0.973653in}}%
\pgfpathquadraticcurveto{\pgfqpoint{6.871100in}{0.967151in}}{\pgfqpoint{6.879476in}{0.967151in}}%
\pgfpathquadraticcurveto{\pgfqpoint{6.887761in}{0.967151in}}{\pgfqpoint{6.892588in}{0.973671in}}%
\pgfpathquadraticcurveto{\pgfqpoint{6.897434in}{0.980210in}}{\pgfqpoint{6.897434in}{0.991467in}}%
\pgfpathquadraticcurveto{\pgfqpoint{6.897434in}{1.002671in}}{\pgfqpoint{6.892588in}{1.009227in}}%
\pgfpathquadraticcurveto{\pgfqpoint{6.887761in}{1.015784in}}{\pgfqpoint{6.879476in}{1.015784in}}%
\pgfpathlineto{\pgfqpoint{6.879476in}{1.015784in}}%
\pgfpathclose%
\pgfpathmoveto{\pgfqpoint{6.879476in}{1.024556in}}%
\pgfpathquadraticcurveto{\pgfqpoint{6.892985in}{1.024556in}}{\pgfqpoint{6.900694in}{1.015766in}}%
\pgfpathquadraticcurveto{\pgfqpoint{6.908421in}{1.006994in}}{\pgfqpoint{6.908421in}{0.991467in}}%
\pgfpathquadraticcurveto{\pgfqpoint{6.908421in}{0.975995in}}{\pgfqpoint{6.900694in}{0.967169in}}%
\pgfpathquadraticcurveto{\pgfqpoint{6.892985in}{0.958361in}}{\pgfqpoint{6.879476in}{0.958361in}}%
\pgfpathquadraticcurveto{\pgfqpoint{6.865912in}{0.958361in}}{\pgfqpoint{6.858221in}{0.967169in}}%
\pgfpathquadraticcurveto{\pgfqpoint{6.850548in}{0.975995in}}{\pgfqpoint{6.850548in}{0.991467in}}%
\pgfpathquadraticcurveto{\pgfqpoint{6.850548in}{1.006994in}}{\pgfqpoint{6.858221in}{1.015766in}}%
\pgfpathquadraticcurveto{\pgfqpoint{6.865912in}{1.024556in}}{\pgfqpoint{6.879476in}{1.024556in}}%
\pgfpathlineto{\pgfqpoint{6.879476in}{1.024556in}}%
\pgfpathclose%
\pgfpathmoveto{\pgfqpoint{6.962115in}{1.013370in}}%
\pgfpathquadraticcurveto{\pgfqpoint{6.960368in}{1.014379in}}{\pgfqpoint{6.958315in}{1.014847in}}%
\pgfpathquadraticcurveto{\pgfqpoint{6.956261in}{1.015333in}}{\pgfqpoint{6.953794in}{1.015333in}}%
\pgfpathquadraticcurveto{\pgfqpoint{6.945004in}{1.015333in}}{\pgfqpoint{6.940303in}{1.009623in}}%
\pgfpathquadraticcurveto{\pgfqpoint{6.935601in}{1.003914in}}{\pgfqpoint{6.935601in}{0.993214in}}%
\pgfpathlineto{\pgfqpoint{6.935601in}{0.960000in}}%
\pgfpathlineto{\pgfqpoint{6.925190in}{0.960000in}}%
\pgfpathlineto{\pgfqpoint{6.925190in}{1.023043in}}%
\pgfpathlineto{\pgfqpoint{6.935601in}{1.023043in}}%
\pgfpathlineto{\pgfqpoint{6.935601in}{1.013244in}}%
\pgfpathquadraticcurveto{\pgfqpoint{6.938880in}{1.018990in}}{\pgfqpoint{6.944103in}{1.021764in}}%
\pgfpathquadraticcurveto{\pgfqpoint{6.949345in}{1.024556in}}{\pgfqpoint{6.956838in}{1.024556in}}%
\pgfpathquadraticcurveto{\pgfqpoint{6.957900in}{1.024556in}}{\pgfqpoint{6.959197in}{1.024411in}}%
\pgfpathquadraticcurveto{\pgfqpoint{6.960494in}{1.024285in}}{\pgfqpoint{6.962061in}{1.023997in}}%
\pgfpathlineto{\pgfqpoint{6.962115in}{1.013370in}}%
\pgfpathlineto{\pgfqpoint{6.962115in}{1.013370in}}%
\pgfpathclose%
\pgfpathmoveto{\pgfqpoint{7.012441in}{1.013478in}}%
\pgfpathlineto{\pgfqpoint{7.012441in}{1.047593in}}%
\pgfpathlineto{\pgfqpoint{7.022798in}{1.047593in}}%
\pgfpathlineto{\pgfqpoint{7.022798in}{0.960000in}}%
\pgfpathlineto{\pgfqpoint{7.012441in}{0.960000in}}%
\pgfpathlineto{\pgfqpoint{7.012441in}{0.969456in}}%
\pgfpathquadraticcurveto{\pgfqpoint{7.009181in}{0.963837in}}{\pgfqpoint{7.004192in}{0.961099in}}%
\pgfpathquadraticcurveto{\pgfqpoint{6.999220in}{0.958361in}}{\pgfqpoint{6.992232in}{0.958361in}}%
\pgfpathquadraticcurveto{\pgfqpoint{6.980812in}{0.958361in}}{\pgfqpoint{6.973625in}{0.967475in}}%
\pgfpathquadraticcurveto{\pgfqpoint{6.966456in}{0.976607in}}{\pgfqpoint{6.966456in}{0.991467in}}%
\pgfpathquadraticcurveto{\pgfqpoint{6.966456in}{1.006327in}}{\pgfqpoint{6.973625in}{1.015441in}}%
\pgfpathquadraticcurveto{\pgfqpoint{6.980812in}{1.024556in}}{\pgfqpoint{6.992232in}{1.024556in}}%
\pgfpathquadraticcurveto{\pgfqpoint{6.999220in}{1.024556in}}{\pgfqpoint{7.004192in}{1.021818in}}%
\pgfpathquadraticcurveto{\pgfqpoint{7.009181in}{1.019098in}}{\pgfqpoint{7.012441in}{1.013478in}}%
\pgfpathlineto{\pgfqpoint{7.012441in}{1.013478in}}%
\pgfpathclose%
\pgfpathmoveto{\pgfqpoint{6.977155in}{0.991467in}}%
\pgfpathquadraticcurveto{\pgfqpoint{6.977155in}{0.980048in}}{\pgfqpoint{6.981857in}{0.973545in}}%
\pgfpathquadraticcurveto{\pgfqpoint{6.986558in}{0.967043in}}{\pgfqpoint{6.994771in}{0.967043in}}%
\pgfpathquadraticcurveto{\pgfqpoint{7.002985in}{0.967043in}}{\pgfqpoint{7.007704in}{0.973545in}}%
\pgfpathquadraticcurveto{\pgfqpoint{7.012441in}{0.980048in}}{\pgfqpoint{7.012441in}{0.991467in}}%
\pgfpathquadraticcurveto{\pgfqpoint{7.012441in}{1.002887in}}{\pgfqpoint{7.007704in}{1.009389in}}%
\pgfpathquadraticcurveto{\pgfqpoint{7.002985in}{1.015892in}}{\pgfqpoint{6.994771in}{1.015892in}}%
\pgfpathquadraticcurveto{\pgfqpoint{6.986558in}{1.015892in}}{\pgfqpoint{6.981857in}{1.009389in}}%
\pgfpathquadraticcurveto{\pgfqpoint{6.977155in}{1.002887in}}{\pgfqpoint{6.977155in}{0.991467in}}%
\pgfpathlineto{\pgfqpoint{6.977155in}{0.991467in}}%
\pgfpathclose%
\pgfpathmoveto{\pgfqpoint{7.044143in}{1.023043in}}%
\pgfpathlineto{\pgfqpoint{7.054500in}{1.023043in}}%
\pgfpathlineto{\pgfqpoint{7.054500in}{0.960000in}}%
\pgfpathlineto{\pgfqpoint{7.044143in}{0.960000in}}%
\pgfpathlineto{\pgfqpoint{7.044143in}{1.023043in}}%
\pgfpathlineto{\pgfqpoint{7.044143in}{1.023043in}}%
\pgfpathclose%
\pgfpathmoveto{\pgfqpoint{7.044143in}{1.047593in}}%
\pgfpathlineto{\pgfqpoint{7.054500in}{1.047593in}}%
\pgfpathlineto{\pgfqpoint{7.054500in}{1.034462in}}%
\pgfpathlineto{\pgfqpoint{7.044143in}{1.034462in}}%
\pgfpathlineto{\pgfqpoint{7.044143in}{1.047593in}}%
\pgfpathlineto{\pgfqpoint{7.044143in}{1.047593in}}%
\pgfpathclose%
\pgfpathmoveto{\pgfqpoint{7.128584in}{0.998060in}}%
\pgfpathlineto{\pgfqpoint{7.128584in}{0.960000in}}%
\pgfpathlineto{\pgfqpoint{7.118227in}{0.960000in}}%
\pgfpathlineto{\pgfqpoint{7.118227in}{0.997717in}}%
\pgfpathquadraticcurveto{\pgfqpoint{7.118227in}{1.006669in}}{\pgfqpoint{7.114732in}{1.011100in}}%
\pgfpathquadraticcurveto{\pgfqpoint{7.111238in}{1.015549in}}{\pgfqpoint{7.104267in}{1.015549in}}%
\pgfpathquadraticcurveto{\pgfqpoint{7.095874in}{1.015549in}}{\pgfqpoint{7.091028in}{1.010200in}}%
\pgfpathquadraticcurveto{\pgfqpoint{7.086183in}{1.004868in}}{\pgfqpoint{7.086183in}{0.995628in}}%
\pgfpathlineto{\pgfqpoint{7.086183in}{0.960000in}}%
\pgfpathlineto{\pgfqpoint{7.075772in}{0.960000in}}%
\pgfpathlineto{\pgfqpoint{7.075772in}{1.023043in}}%
\pgfpathlineto{\pgfqpoint{7.086183in}{1.023043in}}%
\pgfpathlineto{\pgfqpoint{7.086183in}{1.013244in}}%
\pgfpathquadraticcurveto{\pgfqpoint{7.089912in}{1.018936in}}{\pgfqpoint{7.094937in}{1.021746in}}%
\pgfpathquadraticcurveto{\pgfqpoint{7.099980in}{1.024556in}}{\pgfqpoint{7.106573in}{1.024556in}}%
\pgfpathquadraticcurveto{\pgfqpoint{7.117434in}{1.024556in}}{\pgfqpoint{7.123000in}{1.017837in}}%
\pgfpathquadraticcurveto{\pgfqpoint{7.128584in}{1.011118in}}{\pgfqpoint{7.128584in}{0.998060in}}%
\pgfpathlineto{\pgfqpoint{7.128584in}{0.998060in}}%
\pgfpathclose%
\pgfpathmoveto{\pgfqpoint{7.177883in}{0.991683in}}%
\pgfpathquadraticcurveto{\pgfqpoint{7.165328in}{0.991683in}}{\pgfqpoint{7.160483in}{0.988819in}}%
\pgfpathquadraticcurveto{\pgfqpoint{7.155638in}{0.985956in}}{\pgfqpoint{7.155638in}{0.979021in}}%
\pgfpathquadraticcurveto{\pgfqpoint{7.155638in}{0.973509in}}{\pgfqpoint{7.159276in}{0.970267in}}%
\pgfpathquadraticcurveto{\pgfqpoint{7.162915in}{0.967043in}}{\pgfqpoint{7.169147in}{0.967043in}}%
\pgfpathquadraticcurveto{\pgfqpoint{7.177775in}{0.967043in}}{\pgfqpoint{7.182980in}{0.973149in}}%
\pgfpathquadraticcurveto{\pgfqpoint{7.188186in}{0.979255in}}{\pgfqpoint{7.188186in}{0.989378in}}%
\pgfpathlineto{\pgfqpoint{7.188186in}{0.991683in}}%
\pgfpathlineto{\pgfqpoint{7.177883in}{0.991683in}}%
\pgfpathlineto{\pgfqpoint{7.177883in}{0.991683in}}%
\pgfpathclose%
\pgfpathmoveto{\pgfqpoint{7.198543in}{0.995970in}}%
\pgfpathlineto{\pgfqpoint{7.198543in}{0.960000in}}%
\pgfpathlineto{\pgfqpoint{7.188186in}{0.960000in}}%
\pgfpathlineto{\pgfqpoint{7.188186in}{0.969564in}}%
\pgfpathquadraticcurveto{\pgfqpoint{7.184637in}{0.963837in}}{\pgfqpoint{7.179342in}{0.961099in}}%
\pgfpathquadraticcurveto{\pgfqpoint{7.174046in}{0.958361in}}{\pgfqpoint{7.166391in}{0.958361in}}%
\pgfpathquadraticcurveto{\pgfqpoint{7.156719in}{0.958361in}}{\pgfqpoint{7.150991in}{0.963801in}}%
\pgfpathquadraticcurveto{\pgfqpoint{7.145281in}{0.969240in}}{\pgfqpoint{7.145281in}{0.978354in}}%
\pgfpathquadraticcurveto{\pgfqpoint{7.145281in}{0.988982in}}{\pgfqpoint{7.152396in}{0.994385in}}%
\pgfpathquadraticcurveto{\pgfqpoint{7.159529in}{0.999789in}}{\pgfqpoint{7.173650in}{0.999789in}}%
\pgfpathlineto{\pgfqpoint{7.188186in}{0.999789in}}%
\pgfpathlineto{\pgfqpoint{7.188186in}{1.000816in}}%
\pgfpathquadraticcurveto{\pgfqpoint{7.188186in}{1.007966in}}{\pgfqpoint{7.183485in}{1.011875in}}%
\pgfpathquadraticcurveto{\pgfqpoint{7.178784in}{1.015784in}}{\pgfqpoint{7.170282in}{1.015784in}}%
\pgfpathquadraticcurveto{\pgfqpoint{7.164878in}{1.015784in}}{\pgfqpoint{7.159745in}{1.014487in}}%
\pgfpathquadraticcurveto{\pgfqpoint{7.154629in}{1.013190in}}{\pgfqpoint{7.149910in}{1.010596in}}%
\pgfpathlineto{\pgfqpoint{7.149910in}{1.020179in}}%
\pgfpathquadraticcurveto{\pgfqpoint{7.155584in}{1.022376in}}{\pgfqpoint{7.160933in}{1.023457in}}%
\pgfpathquadraticcurveto{\pgfqpoint{7.166283in}{1.024556in}}{\pgfqpoint{7.171344in}{1.024556in}}%
\pgfpathquadraticcurveto{\pgfqpoint{7.185034in}{1.024556in}}{\pgfqpoint{7.191788in}{1.017459in}}%
\pgfpathquadraticcurveto{\pgfqpoint{7.198543in}{1.010380in}}{\pgfqpoint{7.198543in}{0.995970in}}%
\pgfpathlineto{\pgfqpoint{7.198543in}{0.995970in}}%
\pgfpathclose%
\pgfpathmoveto{\pgfqpoint{7.230118in}{1.040947in}}%
\pgfpathlineto{\pgfqpoint{7.230118in}{1.023043in}}%
\pgfpathlineto{\pgfqpoint{7.251445in}{1.023043in}}%
\pgfpathlineto{\pgfqpoint{7.251445in}{1.014991in}}%
\pgfpathlineto{\pgfqpoint{7.230118in}{1.014991in}}%
\pgfpathlineto{\pgfqpoint{7.230118in}{0.980768in}}%
\pgfpathquadraticcurveto{\pgfqpoint{7.230118in}{0.973059in}}{\pgfqpoint{7.232226in}{0.970861in}}%
\pgfpathquadraticcurveto{\pgfqpoint{7.234333in}{0.968664in}}{\pgfqpoint{7.240817in}{0.968664in}}%
\pgfpathlineto{\pgfqpoint{7.251445in}{0.968664in}}%
\pgfpathlineto{\pgfqpoint{7.251445in}{0.960000in}}%
\pgfpathlineto{\pgfqpoint{7.240817in}{0.960000in}}%
\pgfpathquadraticcurveto{\pgfqpoint{7.228821in}{0.960000in}}{\pgfqpoint{7.224264in}{0.964467in}}%
\pgfpathquadraticcurveto{\pgfqpoint{7.219707in}{0.968952in}}{\pgfqpoint{7.219707in}{0.980768in}}%
\pgfpathlineto{\pgfqpoint{7.219707in}{1.014991in}}%
\pgfpathlineto{\pgfqpoint{7.212106in}{1.014991in}}%
\pgfpathlineto{\pgfqpoint{7.212106in}{1.023043in}}%
\pgfpathlineto{\pgfqpoint{7.219707in}{1.023043in}}%
\pgfpathlineto{\pgfqpoint{7.219707in}{1.040947in}}%
\pgfpathlineto{\pgfqpoint{7.230118in}{1.040947in}}%
\pgfpathlineto{\pgfqpoint{7.230118in}{1.040947in}}%
\pgfpathclose%
\pgfpathmoveto{\pgfqpoint{7.318990in}{0.994115in}}%
\pgfpathlineto{\pgfqpoint{7.318990in}{0.989054in}}%
\pgfpathlineto{\pgfqpoint{7.271366in}{0.989054in}}%
\pgfpathquadraticcurveto{\pgfqpoint{7.272050in}{0.978354in}}{\pgfqpoint{7.277814in}{0.972753in}}%
\pgfpathquadraticcurveto{\pgfqpoint{7.283578in}{0.967151in}}{\pgfqpoint{7.293881in}{0.967151in}}%
\pgfpathquadraticcurveto{\pgfqpoint{7.299843in}{0.967151in}}{\pgfqpoint{7.305445in}{0.968610in}}%
\pgfpathquadraticcurveto{\pgfqpoint{7.311047in}{0.970069in}}{\pgfqpoint{7.316576in}{0.973005in}}%
\pgfpathlineto{\pgfqpoint{7.316576in}{0.963206in}}%
\pgfpathquadraticcurveto{\pgfqpoint{7.310993in}{0.960847in}}{\pgfqpoint{7.305139in}{0.959604in}}%
\pgfpathquadraticcurveto{\pgfqpoint{7.299285in}{0.958361in}}{\pgfqpoint{7.293269in}{0.958361in}}%
\pgfpathquadraticcurveto{\pgfqpoint{7.278175in}{0.958361in}}{\pgfqpoint{7.269367in}{0.967133in}}%
\pgfpathquadraticcurveto{\pgfqpoint{7.260559in}{0.975923in}}{\pgfqpoint{7.260559in}{0.990909in}}%
\pgfpathquadraticcurveto{\pgfqpoint{7.260559in}{1.006381in}}{\pgfqpoint{7.268916in}{1.015459in}}%
\pgfpathquadraticcurveto{\pgfqpoint{7.277274in}{1.024556in}}{\pgfqpoint{7.291468in}{1.024556in}}%
\pgfpathquadraticcurveto{\pgfqpoint{7.304184in}{1.024556in}}{\pgfqpoint{7.311587in}{1.016360in}}%
\pgfpathquadraticcurveto{\pgfqpoint{7.318990in}{1.008183in}}{\pgfqpoint{7.318990in}{0.994115in}}%
\pgfpathlineto{\pgfqpoint{7.318990in}{0.994115in}}%
\pgfpathclose%
\pgfpathmoveto{\pgfqpoint{7.308633in}{0.997159in}}%
\pgfpathquadraticcurveto{\pgfqpoint{7.308525in}{1.005643in}}{\pgfqpoint{7.303878in}{1.010704in}}%
\pgfpathquadraticcurveto{\pgfqpoint{7.299231in}{1.015784in}}{\pgfqpoint{7.291576in}{1.015784in}}%
\pgfpathquadraticcurveto{\pgfqpoint{7.282912in}{1.015784in}}{\pgfqpoint{7.277706in}{1.010884in}}%
\pgfpathquadraticcurveto{\pgfqpoint{7.272501in}{1.005985in}}{\pgfqpoint{7.271708in}{0.997087in}}%
\pgfpathlineto{\pgfqpoint{7.308633in}{0.997159in}}%
\pgfpathlineto{\pgfqpoint{7.308633in}{0.997159in}}%
\pgfpathclose%
\pgfpathmoveto{\pgfqpoint{7.376179in}{1.021187in}}%
\pgfpathlineto{\pgfqpoint{7.376179in}{1.011389in}}%
\pgfpathquadraticcurveto{\pgfqpoint{7.371802in}{1.013640in}}{\pgfqpoint{7.367065in}{1.014757in}}%
\pgfpathquadraticcurveto{\pgfqpoint{7.362345in}{1.015892in}}{\pgfqpoint{7.357266in}{1.015892in}}%
\pgfpathquadraticcurveto{\pgfqpoint{7.349557in}{1.015892in}}{\pgfqpoint{7.345702in}{1.013532in}}%
\pgfpathquadraticcurveto{\pgfqpoint{7.341848in}{1.011173in}}{\pgfqpoint{7.341848in}{1.006435in}}%
\pgfpathquadraticcurveto{\pgfqpoint{7.341848in}{1.002833in}}{\pgfqpoint{7.344603in}{1.000780in}}%
\pgfpathquadraticcurveto{\pgfqpoint{7.347359in}{0.998726in}}{\pgfqpoint{7.355699in}{0.996871in}}%
\pgfpathlineto{\pgfqpoint{7.359247in}{0.996078in}}%
\pgfpathquadraticcurveto{\pgfqpoint{7.370271in}{0.993719in}}{\pgfqpoint{7.374918in}{0.989414in}}%
\pgfpathquadraticcurveto{\pgfqpoint{7.379565in}{0.985109in}}{\pgfqpoint{7.379565in}{0.977400in}}%
\pgfpathquadraticcurveto{\pgfqpoint{7.379565in}{0.968610in}}{\pgfqpoint{7.372612in}{0.963476in}}%
\pgfpathquadraticcurveto{\pgfqpoint{7.365660in}{0.958361in}}{\pgfqpoint{7.353501in}{0.958361in}}%
\pgfpathquadraticcurveto{\pgfqpoint{7.348440in}{0.958361in}}{\pgfqpoint{7.342946in}{0.959352in}}%
\pgfpathquadraticcurveto{\pgfqpoint{7.337453in}{0.960342in}}{\pgfqpoint{7.331382in}{0.962306in}}%
\pgfpathlineto{\pgfqpoint{7.331382in}{0.973005in}}%
\pgfpathquadraticcurveto{\pgfqpoint{7.337128in}{0.970015in}}{\pgfqpoint{7.342694in}{0.968520in}}%
\pgfpathquadraticcurveto{\pgfqpoint{7.348260in}{0.967043in}}{\pgfqpoint{7.353736in}{0.967043in}}%
\pgfpathquadraticcurveto{\pgfqpoint{7.361048in}{0.967043in}}{\pgfqpoint{7.364975in}{0.969546in}}%
\pgfpathquadraticcurveto{\pgfqpoint{7.368920in}{0.972050in}}{\pgfqpoint{7.368920in}{0.976607in}}%
\pgfpathquadraticcurveto{\pgfqpoint{7.368920in}{0.980822in}}{\pgfqpoint{7.366074in}{0.983074in}}%
\pgfpathquadraticcurveto{\pgfqpoint{7.363246in}{0.985325in}}{\pgfqpoint{7.353609in}{0.987414in}}%
\pgfpathlineto{\pgfqpoint{7.350007in}{0.988261in}}%
\pgfpathquadraticcurveto{\pgfqpoint{7.340389in}{0.990278in}}{\pgfqpoint{7.336102in}{0.994475in}}%
\pgfpathquadraticcurveto{\pgfqpoint{7.331833in}{0.998672in}}{\pgfqpoint{7.331833in}{1.005985in}}%
\pgfpathquadraticcurveto{\pgfqpoint{7.331833in}{1.014883in}}{\pgfqpoint{7.338137in}{1.019710in}}%
\pgfpathquadraticcurveto{\pgfqpoint{7.344441in}{1.024556in}}{\pgfqpoint{7.356041in}{1.024556in}}%
\pgfpathquadraticcurveto{\pgfqpoint{7.361769in}{1.024556in}}{\pgfqpoint{7.366830in}{1.023709in}}%
\pgfpathquadraticcurveto{\pgfqpoint{7.371910in}{1.022880in}}{\pgfqpoint{7.376179in}{1.021187in}}%
\pgfpathlineto{\pgfqpoint{7.376179in}{1.021187in}}%
\pgfpathclose%
\pgfusepath{fill}%
\end{pgfscope}%
\begin{pgfscope}%
\pgfsetbuttcap%
\pgfsetroundjoin%
\definecolor{currentfill}{rgb}{0.309804,0.372549,0.419608}%
\pgfsetfillcolor{currentfill}%
\pgfsetlinewidth{0.000000pt}%
\definecolor{currentstroke}{rgb}{0.309804,0.372549,0.419608}%
\pgfsetstrokecolor{currentstroke}%
\pgfsetdash{}{0pt}%
\pgfpathmoveto{\pgfqpoint{7.514365in}{1.020629in}}%
\pgfpathlineto{\pgfqpoint{7.514365in}{1.010938in}}%
\pgfpathquadraticcurveto{\pgfqpoint{7.509970in}{1.013370in}}{\pgfqpoint{7.505557in}{1.014577in}}%
\pgfpathquadraticcurveto{\pgfqpoint{7.501144in}{1.015784in}}{\pgfqpoint{7.496641in}{1.015784in}}%
\pgfpathquadraticcurveto{\pgfqpoint{7.486555in}{1.015784in}}{\pgfqpoint{7.480971in}{1.009389in}}%
\pgfpathquadraticcurveto{\pgfqpoint{7.475405in}{1.003013in}}{\pgfqpoint{7.475405in}{0.991467in}}%
\pgfpathquadraticcurveto{\pgfqpoint{7.475405in}{0.979921in}}{\pgfqpoint{7.480971in}{0.973527in}}%
\pgfpathquadraticcurveto{\pgfqpoint{7.486555in}{0.967151in}}{\pgfqpoint{7.496641in}{0.967151in}}%
\pgfpathquadraticcurveto{\pgfqpoint{7.501144in}{0.967151in}}{\pgfqpoint{7.505557in}{0.968358in}}%
\pgfpathquadraticcurveto{\pgfqpoint{7.509970in}{0.969564in}}{\pgfqpoint{7.514365in}{0.971996in}}%
\pgfpathlineto{\pgfqpoint{7.514365in}{0.962414in}}%
\pgfpathquadraticcurveto{\pgfqpoint{7.510024in}{0.960396in}}{\pgfqpoint{7.505377in}{0.959388in}}%
\pgfpathquadraticcurveto{\pgfqpoint{7.500748in}{0.958361in}}{\pgfqpoint{7.495507in}{0.958361in}}%
\pgfpathquadraticcurveto{\pgfqpoint{7.481259in}{0.958361in}}{\pgfqpoint{7.472865in}{0.967313in}}%
\pgfpathquadraticcurveto{\pgfqpoint{7.464490in}{0.976265in}}{\pgfqpoint{7.464490in}{0.991467in}}%
\pgfpathquadraticcurveto{\pgfqpoint{7.464490in}{1.006886in}}{\pgfqpoint{7.472955in}{1.015712in}}%
\pgfpathquadraticcurveto{\pgfqpoint{7.481439in}{1.024556in}}{\pgfqpoint{7.496191in}{1.024556in}}%
\pgfpathquadraticcurveto{\pgfqpoint{7.500964in}{1.024556in}}{\pgfqpoint{7.505521in}{1.023565in}}%
\pgfpathquadraticcurveto{\pgfqpoint{7.510079in}{1.022592in}}{\pgfqpoint{7.514365in}{1.020629in}}%
\pgfpathlineto{\pgfqpoint{7.514365in}{1.020629in}}%
\pgfpathclose%
\pgfpathmoveto{\pgfqpoint{7.561035in}{0.991683in}}%
\pgfpathquadraticcurveto{\pgfqpoint{7.548480in}{0.991683in}}{\pgfqpoint{7.543635in}{0.988819in}}%
\pgfpathquadraticcurveto{\pgfqpoint{7.538790in}{0.985956in}}{\pgfqpoint{7.538790in}{0.979021in}}%
\pgfpathquadraticcurveto{\pgfqpoint{7.538790in}{0.973509in}}{\pgfqpoint{7.542428in}{0.970267in}}%
\pgfpathquadraticcurveto{\pgfqpoint{7.546067in}{0.967043in}}{\pgfqpoint{7.552299in}{0.967043in}}%
\pgfpathquadraticcurveto{\pgfqpoint{7.560927in}{0.967043in}}{\pgfqpoint{7.566132in}{0.973149in}}%
\pgfpathquadraticcurveto{\pgfqpoint{7.571338in}{0.979255in}}{\pgfqpoint{7.571338in}{0.989378in}}%
\pgfpathlineto{\pgfqpoint{7.571338in}{0.991683in}}%
\pgfpathlineto{\pgfqpoint{7.561035in}{0.991683in}}%
\pgfpathlineto{\pgfqpoint{7.561035in}{0.991683in}}%
\pgfpathclose%
\pgfpathmoveto{\pgfqpoint{7.581695in}{0.995970in}}%
\pgfpathlineto{\pgfqpoint{7.581695in}{0.960000in}}%
\pgfpathlineto{\pgfqpoint{7.571338in}{0.960000in}}%
\pgfpathlineto{\pgfqpoint{7.571338in}{0.969564in}}%
\pgfpathquadraticcurveto{\pgfqpoint{7.567789in}{0.963837in}}{\pgfqpoint{7.562494in}{0.961099in}}%
\pgfpathquadraticcurveto{\pgfqpoint{7.557198in}{0.958361in}}{\pgfqpoint{7.549543in}{0.958361in}}%
\pgfpathquadraticcurveto{\pgfqpoint{7.539871in}{0.958361in}}{\pgfqpoint{7.534143in}{0.963801in}}%
\pgfpathquadraticcurveto{\pgfqpoint{7.528433in}{0.969240in}}{\pgfqpoint{7.528433in}{0.978354in}}%
\pgfpathquadraticcurveto{\pgfqpoint{7.528433in}{0.988982in}}{\pgfqpoint{7.535548in}{0.994385in}}%
\pgfpathquadraticcurveto{\pgfqpoint{7.542681in}{0.999789in}}{\pgfqpoint{7.556802in}{0.999789in}}%
\pgfpathlineto{\pgfqpoint{7.571338in}{0.999789in}}%
\pgfpathlineto{\pgfqpoint{7.571338in}{1.000816in}}%
\pgfpathquadraticcurveto{\pgfqpoint{7.571338in}{1.007966in}}{\pgfqpoint{7.566637in}{1.011875in}}%
\pgfpathquadraticcurveto{\pgfqpoint{7.561936in}{1.015784in}}{\pgfqpoint{7.553434in}{1.015784in}}%
\pgfpathquadraticcurveto{\pgfqpoint{7.548030in}{1.015784in}}{\pgfqpoint{7.542897in}{1.014487in}}%
\pgfpathquadraticcurveto{\pgfqpoint{7.537781in}{1.013190in}}{\pgfqpoint{7.533062in}{1.010596in}}%
\pgfpathlineto{\pgfqpoint{7.533062in}{1.020179in}}%
\pgfpathquadraticcurveto{\pgfqpoint{7.538736in}{1.022376in}}{\pgfqpoint{7.544085in}{1.023457in}}%
\pgfpathquadraticcurveto{\pgfqpoint{7.549435in}{1.024556in}}{\pgfqpoint{7.554496in}{1.024556in}}%
\pgfpathquadraticcurveto{\pgfqpoint{7.568186in}{1.024556in}}{\pgfqpoint{7.574940in}{1.017459in}}%
\pgfpathquadraticcurveto{\pgfqpoint{7.581695in}{1.010380in}}{\pgfqpoint{7.581695in}{0.995970in}}%
\pgfpathlineto{\pgfqpoint{7.581695in}{0.995970in}}%
\pgfpathclose%
\pgfpathmoveto{\pgfqpoint{7.613036in}{0.969456in}}%
\pgfpathlineto{\pgfqpoint{7.613036in}{0.936026in}}%
\pgfpathlineto{\pgfqpoint{7.602625in}{0.936026in}}%
\pgfpathlineto{\pgfqpoint{7.602625in}{1.023043in}}%
\pgfpathlineto{\pgfqpoint{7.613036in}{1.023043in}}%
\pgfpathlineto{\pgfqpoint{7.613036in}{1.013478in}}%
\pgfpathquadraticcurveto{\pgfqpoint{7.616314in}{1.019098in}}{\pgfqpoint{7.621286in}{1.021818in}}%
\pgfpathquadraticcurveto{\pgfqpoint{7.626275in}{1.024556in}}{\pgfqpoint{7.633192in}{1.024556in}}%
\pgfpathquadraticcurveto{\pgfqpoint{7.644683in}{1.024556in}}{\pgfqpoint{7.651852in}{1.015441in}}%
\pgfpathquadraticcurveto{\pgfqpoint{7.659039in}{1.006327in}}{\pgfqpoint{7.659039in}{0.991467in}}%
\pgfpathquadraticcurveto{\pgfqpoint{7.659039in}{0.976607in}}{\pgfqpoint{7.651852in}{0.967475in}}%
\pgfpathquadraticcurveto{\pgfqpoint{7.644683in}{0.958361in}}{\pgfqpoint{7.633192in}{0.958361in}}%
\pgfpathquadraticcurveto{\pgfqpoint{7.626275in}{0.958361in}}{\pgfqpoint{7.621286in}{0.961099in}}%
\pgfpathquadraticcurveto{\pgfqpoint{7.616314in}{0.963837in}}{\pgfqpoint{7.613036in}{0.969456in}}%
\pgfpathlineto{\pgfqpoint{7.613036in}{0.969456in}}%
\pgfpathclose%
\pgfpathmoveto{\pgfqpoint{7.648286in}{0.991467in}}%
\pgfpathquadraticcurveto{\pgfqpoint{7.648286in}{1.002887in}}{\pgfqpoint{7.643585in}{1.009389in}}%
\pgfpathquadraticcurveto{\pgfqpoint{7.638883in}{1.015892in}}{\pgfqpoint{7.630670in}{1.015892in}}%
\pgfpathquadraticcurveto{\pgfqpoint{7.622438in}{1.015892in}}{\pgfqpoint{7.617737in}{1.009389in}}%
\pgfpathquadraticcurveto{\pgfqpoint{7.613036in}{1.002887in}}{\pgfqpoint{7.613036in}{0.991467in}}%
\pgfpathquadraticcurveto{\pgfqpoint{7.613036in}{0.980048in}}{\pgfqpoint{7.617737in}{0.973545in}}%
\pgfpathquadraticcurveto{\pgfqpoint{7.622438in}{0.967043in}}{\pgfqpoint{7.630670in}{0.967043in}}%
\pgfpathquadraticcurveto{\pgfqpoint{7.638883in}{0.967043in}}{\pgfqpoint{7.643585in}{0.973545in}}%
\pgfpathquadraticcurveto{\pgfqpoint{7.648286in}{0.980048in}}{\pgfqpoint{7.648286in}{0.991467in}}%
\pgfpathlineto{\pgfqpoint{7.648286in}{0.991467in}}%
\pgfpathclose%
\pgfpathmoveto{\pgfqpoint{7.686219in}{0.969456in}}%
\pgfpathlineto{\pgfqpoint{7.686219in}{0.936026in}}%
\pgfpathlineto{\pgfqpoint{7.675808in}{0.936026in}}%
\pgfpathlineto{\pgfqpoint{7.675808in}{1.023043in}}%
\pgfpathlineto{\pgfqpoint{7.686219in}{1.023043in}}%
\pgfpathlineto{\pgfqpoint{7.686219in}{1.013478in}}%
\pgfpathquadraticcurveto{\pgfqpoint{7.689498in}{1.019098in}}{\pgfqpoint{7.694469in}{1.021818in}}%
\pgfpathquadraticcurveto{\pgfqpoint{7.699458in}{1.024556in}}{\pgfqpoint{7.706375in}{1.024556in}}%
\pgfpathquadraticcurveto{\pgfqpoint{7.717867in}{1.024556in}}{\pgfqpoint{7.725036in}{1.015441in}}%
\pgfpathquadraticcurveto{\pgfqpoint{7.732222in}{1.006327in}}{\pgfqpoint{7.732222in}{0.991467in}}%
\pgfpathquadraticcurveto{\pgfqpoint{7.732222in}{0.976607in}}{\pgfqpoint{7.725036in}{0.967475in}}%
\pgfpathquadraticcurveto{\pgfqpoint{7.717867in}{0.958361in}}{\pgfqpoint{7.706375in}{0.958361in}}%
\pgfpathquadraticcurveto{\pgfqpoint{7.699458in}{0.958361in}}{\pgfqpoint{7.694469in}{0.961099in}}%
\pgfpathquadraticcurveto{\pgfqpoint{7.689498in}{0.963837in}}{\pgfqpoint{7.686219in}{0.969456in}}%
\pgfpathlineto{\pgfqpoint{7.686219in}{0.969456in}}%
\pgfpathclose%
\pgfpathmoveto{\pgfqpoint{7.721469in}{0.991467in}}%
\pgfpathquadraticcurveto{\pgfqpoint{7.721469in}{1.002887in}}{\pgfqpoint{7.716768in}{1.009389in}}%
\pgfpathquadraticcurveto{\pgfqpoint{7.712067in}{1.015892in}}{\pgfqpoint{7.703853in}{1.015892in}}%
\pgfpathquadraticcurveto{\pgfqpoint{7.695622in}{1.015892in}}{\pgfqpoint{7.690921in}{1.009389in}}%
\pgfpathquadraticcurveto{\pgfqpoint{7.686219in}{1.002887in}}{\pgfqpoint{7.686219in}{0.991467in}}%
\pgfpathquadraticcurveto{\pgfqpoint{7.686219in}{0.980048in}}{\pgfqpoint{7.690921in}{0.973545in}}%
\pgfpathquadraticcurveto{\pgfqpoint{7.695622in}{0.967043in}}{\pgfqpoint{7.703853in}{0.967043in}}%
\pgfpathquadraticcurveto{\pgfqpoint{7.712067in}{0.967043in}}{\pgfqpoint{7.716768in}{0.973545in}}%
\pgfpathquadraticcurveto{\pgfqpoint{7.721469in}{0.980048in}}{\pgfqpoint{7.721469in}{0.991467in}}%
\pgfpathlineto{\pgfqpoint{7.721469in}{0.991467in}}%
\pgfpathclose%
\pgfpathmoveto{\pgfqpoint{7.803316in}{0.994115in}}%
\pgfpathlineto{\pgfqpoint{7.803316in}{0.989054in}}%
\pgfpathlineto{\pgfqpoint{7.755692in}{0.989054in}}%
\pgfpathquadraticcurveto{\pgfqpoint{7.756377in}{0.978354in}}{\pgfqpoint{7.762141in}{0.972753in}}%
\pgfpathquadraticcurveto{\pgfqpoint{7.767904in}{0.967151in}}{\pgfqpoint{7.778207in}{0.967151in}}%
\pgfpathquadraticcurveto{\pgfqpoint{7.784169in}{0.967151in}}{\pgfqpoint{7.789771in}{0.968610in}}%
\pgfpathquadraticcurveto{\pgfqpoint{7.795373in}{0.970069in}}{\pgfqpoint{7.800903in}{0.973005in}}%
\pgfpathlineto{\pgfqpoint{7.800903in}{0.963206in}}%
\pgfpathquadraticcurveto{\pgfqpoint{7.795319in}{0.960847in}}{\pgfqpoint{7.789465in}{0.959604in}}%
\pgfpathquadraticcurveto{\pgfqpoint{7.783611in}{0.958361in}}{\pgfqpoint{7.777595in}{0.958361in}}%
\pgfpathquadraticcurveto{\pgfqpoint{7.762501in}{0.958361in}}{\pgfqpoint{7.753693in}{0.967133in}}%
\pgfpathquadraticcurveto{\pgfqpoint{7.744885in}{0.975923in}}{\pgfqpoint{7.744885in}{0.990909in}}%
\pgfpathquadraticcurveto{\pgfqpoint{7.744885in}{1.006381in}}{\pgfqpoint{7.753243in}{1.015459in}}%
\pgfpathquadraticcurveto{\pgfqpoint{7.761600in}{1.024556in}}{\pgfqpoint{7.775794in}{1.024556in}}%
\pgfpathquadraticcurveto{\pgfqpoint{7.788510in}{1.024556in}}{\pgfqpoint{7.795913in}{1.016360in}}%
\pgfpathquadraticcurveto{\pgfqpoint{7.803316in}{1.008183in}}{\pgfqpoint{7.803316in}{0.994115in}}%
\pgfpathlineto{\pgfqpoint{7.803316in}{0.994115in}}%
\pgfpathclose%
\pgfpathmoveto{\pgfqpoint{7.792959in}{0.997159in}}%
\pgfpathquadraticcurveto{\pgfqpoint{7.792851in}{1.005643in}}{\pgfqpoint{7.788204in}{1.010704in}}%
\pgfpathquadraticcurveto{\pgfqpoint{7.783557in}{1.015784in}}{\pgfqpoint{7.775902in}{1.015784in}}%
\pgfpathquadraticcurveto{\pgfqpoint{7.767238in}{1.015784in}}{\pgfqpoint{7.762033in}{1.010884in}}%
\pgfpathquadraticcurveto{\pgfqpoint{7.756827in}{1.005985in}}{\pgfqpoint{7.756034in}{0.997087in}}%
\pgfpathlineto{\pgfqpoint{7.792959in}{0.997159in}}%
\pgfpathlineto{\pgfqpoint{7.792959in}{0.997159in}}%
\pgfpathclose%
\pgfpathmoveto{\pgfqpoint{7.861802in}{1.013478in}}%
\pgfpathlineto{\pgfqpoint{7.861802in}{1.047593in}}%
\pgfpathlineto{\pgfqpoint{7.872159in}{1.047593in}}%
\pgfpathlineto{\pgfqpoint{7.872159in}{0.960000in}}%
\pgfpathlineto{\pgfqpoint{7.861802in}{0.960000in}}%
\pgfpathlineto{\pgfqpoint{7.861802in}{0.969456in}}%
\pgfpathquadraticcurveto{\pgfqpoint{7.858542in}{0.963837in}}{\pgfqpoint{7.853552in}{0.961099in}}%
\pgfpathquadraticcurveto{\pgfqpoint{7.848581in}{0.958361in}}{\pgfqpoint{7.841592in}{0.958361in}}%
\pgfpathquadraticcurveto{\pgfqpoint{7.830172in}{0.958361in}}{\pgfqpoint{7.822986in}{0.967475in}}%
\pgfpathquadraticcurveto{\pgfqpoint{7.815817in}{0.976607in}}{\pgfqpoint{7.815817in}{0.991467in}}%
\pgfpathquadraticcurveto{\pgfqpoint{7.815817in}{1.006327in}}{\pgfqpoint{7.822986in}{1.015441in}}%
\pgfpathquadraticcurveto{\pgfqpoint{7.830172in}{1.024556in}}{\pgfqpoint{7.841592in}{1.024556in}}%
\pgfpathquadraticcurveto{\pgfqpoint{7.848581in}{1.024556in}}{\pgfqpoint{7.853552in}{1.021818in}}%
\pgfpathquadraticcurveto{\pgfqpoint{7.858542in}{1.019098in}}{\pgfqpoint{7.861802in}{1.013478in}}%
\pgfpathlineto{\pgfqpoint{7.861802in}{1.013478in}}%
\pgfpathclose%
\pgfpathmoveto{\pgfqpoint{7.826516in}{0.991467in}}%
\pgfpathquadraticcurveto{\pgfqpoint{7.826516in}{0.980048in}}{\pgfqpoint{7.831217in}{0.973545in}}%
\pgfpathquadraticcurveto{\pgfqpoint{7.835918in}{0.967043in}}{\pgfqpoint{7.844132in}{0.967043in}}%
\pgfpathquadraticcurveto{\pgfqpoint{7.852345in}{0.967043in}}{\pgfqpoint{7.857065in}{0.973545in}}%
\pgfpathquadraticcurveto{\pgfqpoint{7.861802in}{0.980048in}}{\pgfqpoint{7.861802in}{0.991467in}}%
\pgfpathquadraticcurveto{\pgfqpoint{7.861802in}{1.002887in}}{\pgfqpoint{7.857065in}{1.009389in}}%
\pgfpathquadraticcurveto{\pgfqpoint{7.852345in}{1.015892in}}{\pgfqpoint{7.844132in}{1.015892in}}%
\pgfpathquadraticcurveto{\pgfqpoint{7.835918in}{1.015892in}}{\pgfqpoint{7.831217in}{1.009389in}}%
\pgfpathquadraticcurveto{\pgfqpoint{7.826516in}{1.002887in}}{\pgfqpoint{7.826516in}{0.991467in}}%
\pgfpathlineto{\pgfqpoint{7.826516in}{0.991467in}}%
\pgfpathclose%
\pgfusepath{fill}%
\end{pgfscope}%
\begin{pgfscope}%
\pgfsetbuttcap%
\pgfsetroundjoin%
\definecolor{currentfill}{rgb}{0.309804,0.372549,0.419608}%
\pgfsetfillcolor{currentfill}%
\pgfsetlinewidth{0.000000pt}%
\definecolor{currentstroke}{rgb}{0.309804,0.372549,0.419608}%
\pgfsetstrokecolor{currentstroke}%
\pgfsetdash{}{0pt}%
\pgfpathmoveto{\pgfqpoint{7.997543in}{0.991683in}}%
\pgfpathquadraticcurveto{\pgfqpoint{7.984989in}{0.991683in}}{\pgfqpoint{7.980144in}{0.988819in}}%
\pgfpathquadraticcurveto{\pgfqpoint{7.975298in}{0.985956in}}{\pgfqpoint{7.975298in}{0.979021in}}%
\pgfpathquadraticcurveto{\pgfqpoint{7.975298in}{0.973509in}}{\pgfqpoint{7.978937in}{0.970267in}}%
\pgfpathquadraticcurveto{\pgfqpoint{7.982575in}{0.967043in}}{\pgfqpoint{7.988808in}{0.967043in}}%
\pgfpathquadraticcurveto{\pgfqpoint{7.997435in}{0.967043in}}{\pgfqpoint{8.002641in}{0.973149in}}%
\pgfpathquadraticcurveto{\pgfqpoint{8.007846in}{0.979255in}}{\pgfqpoint{8.007846in}{0.989378in}}%
\pgfpathlineto{\pgfqpoint{8.007846in}{0.991683in}}%
\pgfpathlineto{\pgfqpoint{7.997543in}{0.991683in}}%
\pgfpathlineto{\pgfqpoint{7.997543in}{0.991683in}}%
\pgfpathclose%
\pgfpathmoveto{\pgfqpoint{8.018203in}{0.995970in}}%
\pgfpathlineto{\pgfqpoint{8.018203in}{0.960000in}}%
\pgfpathlineto{\pgfqpoint{8.007846in}{0.960000in}}%
\pgfpathlineto{\pgfqpoint{8.007846in}{0.969564in}}%
\pgfpathquadraticcurveto{\pgfqpoint{8.004298in}{0.963837in}}{\pgfqpoint{7.999002in}{0.961099in}}%
\pgfpathquadraticcurveto{\pgfqpoint{7.993707in}{0.958361in}}{\pgfqpoint{7.986052in}{0.958361in}}%
\pgfpathquadraticcurveto{\pgfqpoint{7.976379in}{0.958361in}}{\pgfqpoint{7.970651in}{0.963801in}}%
\pgfpathquadraticcurveto{\pgfqpoint{7.964941in}{0.969240in}}{\pgfqpoint{7.964941in}{0.978354in}}%
\pgfpathquadraticcurveto{\pgfqpoint{7.964941in}{0.988982in}}{\pgfqpoint{7.972056in}{0.994385in}}%
\pgfpathquadraticcurveto{\pgfqpoint{7.979189in}{0.999789in}}{\pgfqpoint{7.993311in}{0.999789in}}%
\pgfpathlineto{\pgfqpoint{8.007846in}{0.999789in}}%
\pgfpathlineto{\pgfqpoint{8.007846in}{1.000816in}}%
\pgfpathquadraticcurveto{\pgfqpoint{8.007846in}{1.007966in}}{\pgfqpoint{8.003145in}{1.011875in}}%
\pgfpathquadraticcurveto{\pgfqpoint{7.998444in}{1.015784in}}{\pgfqpoint{7.989942in}{1.015784in}}%
\pgfpathquadraticcurveto{\pgfqpoint{7.984539in}{1.015784in}}{\pgfqpoint{7.979405in}{1.014487in}}%
\pgfpathquadraticcurveto{\pgfqpoint{7.974290in}{1.013190in}}{\pgfqpoint{7.969571in}{1.010596in}}%
\pgfpathlineto{\pgfqpoint{7.969571in}{1.020179in}}%
\pgfpathquadraticcurveto{\pgfqpoint{7.975244in}{1.022376in}}{\pgfqpoint{7.980594in}{1.023457in}}%
\pgfpathquadraticcurveto{\pgfqpoint{7.985944in}{1.024556in}}{\pgfqpoint{7.991005in}{1.024556in}}%
\pgfpathquadraticcurveto{\pgfqpoint{8.004694in}{1.024556in}}{\pgfqpoint{8.011449in}{1.017459in}}%
\pgfpathquadraticcurveto{\pgfqpoint{8.018203in}{1.010380in}}{\pgfqpoint{8.018203in}{0.995970in}}%
\pgfpathlineto{\pgfqpoint{8.018203in}{0.995970in}}%
\pgfpathclose%
\pgfpathmoveto{\pgfqpoint{8.049779in}{1.040947in}}%
\pgfpathlineto{\pgfqpoint{8.049779in}{1.023043in}}%
\pgfpathlineto{\pgfqpoint{8.071105in}{1.023043in}}%
\pgfpathlineto{\pgfqpoint{8.071105in}{1.014991in}}%
\pgfpathlineto{\pgfqpoint{8.049779in}{1.014991in}}%
\pgfpathlineto{\pgfqpoint{8.049779in}{0.980768in}}%
\pgfpathquadraticcurveto{\pgfqpoint{8.049779in}{0.973059in}}{\pgfqpoint{8.051886in}{0.970861in}}%
\pgfpathquadraticcurveto{\pgfqpoint{8.053994in}{0.968664in}}{\pgfqpoint{8.060478in}{0.968664in}}%
\pgfpathlineto{\pgfqpoint{8.071105in}{0.968664in}}%
\pgfpathlineto{\pgfqpoint{8.071105in}{0.960000in}}%
\pgfpathlineto{\pgfqpoint{8.060478in}{0.960000in}}%
\pgfpathquadraticcurveto{\pgfqpoint{8.048482in}{0.960000in}}{\pgfqpoint{8.043925in}{0.964467in}}%
\pgfpathquadraticcurveto{\pgfqpoint{8.039368in}{0.968952in}}{\pgfqpoint{8.039368in}{0.980768in}}%
\pgfpathlineto{\pgfqpoint{8.039368in}{1.014991in}}%
\pgfpathlineto{\pgfqpoint{8.031767in}{1.014991in}}%
\pgfpathlineto{\pgfqpoint{8.031767in}{1.023043in}}%
\pgfpathlineto{\pgfqpoint{8.039368in}{1.023043in}}%
\pgfpathlineto{\pgfqpoint{8.039368in}{1.040947in}}%
\pgfpathlineto{\pgfqpoint{8.049779in}{1.040947in}}%
\pgfpathlineto{\pgfqpoint{8.049779in}{1.040947in}}%
\pgfpathclose%
\pgfusepath{fill}%
\end{pgfscope}%
\begin{pgfscope}%
\pgfsetbuttcap%
\pgfsetroundjoin%
\definecolor{currentfill}{rgb}{0.309804,0.372549,0.419608}%
\pgfsetfillcolor{currentfill}%
\pgfsetlinewidth{0.000000pt}%
\definecolor{currentstroke}{rgb}{0.309804,0.372549,0.419608}%
\pgfsetstrokecolor{currentstroke}%
\pgfsetdash{}{0pt}%
\pgfpathmoveto{\pgfqpoint{8.180946in}{1.032283in}}%
\pgfpathlineto{\pgfqpoint{8.180946in}{1.010488in}}%
\pgfpathlineto{\pgfqpoint{8.212287in}{1.010488in}}%
\pgfpathlineto{\pgfqpoint{8.212287in}{1.000924in}}%
\pgfpathlineto{\pgfqpoint{8.180946in}{1.000924in}}%
\pgfpathlineto{\pgfqpoint{8.180946in}{0.979147in}}%
\pgfpathlineto{\pgfqpoint{8.171490in}{0.979147in}}%
\pgfpathlineto{\pgfqpoint{8.171490in}{1.000924in}}%
\pgfpathlineto{\pgfqpoint{8.140130in}{1.000924in}}%
\pgfpathlineto{\pgfqpoint{8.140130in}{1.010488in}}%
\pgfpathlineto{\pgfqpoint{8.171490in}{1.010488in}}%
\pgfpathlineto{\pgfqpoint{8.171490in}{1.032283in}}%
\pgfpathlineto{\pgfqpoint{8.180946in}{1.032283in}}%
\pgfpathlineto{\pgfqpoint{8.180946in}{1.032283in}}%
\pgfpathclose%
\pgfpathmoveto{\pgfqpoint{8.140130in}{0.969564in}}%
\pgfpathlineto{\pgfqpoint{8.212287in}{0.969564in}}%
\pgfpathlineto{\pgfqpoint{8.212287in}{0.960000in}}%
\pgfpathlineto{\pgfqpoint{8.140130in}{0.960000in}}%
\pgfpathlineto{\pgfqpoint{8.140130in}{0.969564in}}%
\pgfpathlineto{\pgfqpoint{8.140130in}{0.969564in}}%
\pgfpathclose%
\pgfpathmoveto{\pgfqpoint{8.236964in}{1.044045in}}%
\pgfpathlineto{\pgfqpoint{8.281598in}{1.044045in}}%
\pgfpathlineto{\pgfqpoint{8.281598in}{1.034462in}}%
\pgfpathlineto{\pgfqpoint{8.247375in}{1.034462in}}%
\pgfpathlineto{\pgfqpoint{8.247375in}{1.013874in}}%
\pgfpathquadraticcurveto{\pgfqpoint{8.249842in}{1.014721in}}{\pgfqpoint{8.252310in}{1.015135in}}%
\pgfpathquadraticcurveto{\pgfqpoint{8.254796in}{1.015549in}}{\pgfqpoint{8.257281in}{1.015549in}}%
\pgfpathquadraticcurveto{\pgfqpoint{8.271349in}{1.015549in}}{\pgfqpoint{8.279562in}{1.007840in}}%
\pgfpathquadraticcurveto{\pgfqpoint{8.287794in}{1.000131in}}{\pgfqpoint{8.287794in}{0.986964in}}%
\pgfpathquadraticcurveto{\pgfqpoint{8.287794in}{0.973401in}}{\pgfqpoint{8.279346in}{0.965872in}}%
\pgfpathquadraticcurveto{\pgfqpoint{8.270899in}{0.958361in}}{\pgfqpoint{8.255534in}{0.958361in}}%
\pgfpathquadraticcurveto{\pgfqpoint{8.250239in}{0.958361in}}{\pgfqpoint{8.244745in}{0.959262in}}%
\pgfpathquadraticcurveto{\pgfqpoint{8.239269in}{0.960162in}}{\pgfqpoint{8.233415in}{0.961963in}}%
\pgfpathlineto{\pgfqpoint{8.233415in}{0.973401in}}%
\pgfpathquadraticcurveto{\pgfqpoint{8.238477in}{0.970645in}}{\pgfqpoint{8.243880in}{0.969294in}}%
\pgfpathquadraticcurveto{\pgfqpoint{8.249284in}{0.967943in}}{\pgfqpoint{8.255300in}{0.967943in}}%
\pgfpathquadraticcurveto{\pgfqpoint{8.265045in}{0.967943in}}{\pgfqpoint{8.270718in}{0.973059in}}%
\pgfpathquadraticcurveto{\pgfqpoint{8.276410in}{0.978174in}}{\pgfqpoint{8.276410in}{0.986964in}}%
\pgfpathquadraticcurveto{\pgfqpoint{8.276410in}{0.995736in}}{\pgfqpoint{8.270718in}{1.000852in}}%
\pgfpathquadraticcurveto{\pgfqpoint{8.265045in}{1.005985in}}{\pgfqpoint{8.255300in}{1.005985in}}%
\pgfpathquadraticcurveto{\pgfqpoint{8.250743in}{1.005985in}}{\pgfqpoint{8.246204in}{1.004976in}}%
\pgfpathquadraticcurveto{\pgfqpoint{8.241683in}{1.003968in}}{\pgfqpoint{8.236964in}{1.001824in}}%
\pgfpathlineto{\pgfqpoint{8.236964in}{1.044045in}}%
\pgfpathlineto{\pgfqpoint{8.236964in}{1.044045in}}%
\pgfpathclose%
\pgfpathmoveto{\pgfqpoint{8.311372in}{1.019602in}}%
\pgfpathlineto{\pgfqpoint{8.323242in}{1.019602in}}%
\pgfpathlineto{\pgfqpoint{8.323242in}{1.005319in}}%
\pgfpathlineto{\pgfqpoint{8.311372in}{1.005319in}}%
\pgfpathlineto{\pgfqpoint{8.311372in}{1.019602in}}%
\pgfpathlineto{\pgfqpoint{8.311372in}{1.019602in}}%
\pgfpathclose%
\pgfpathmoveto{\pgfqpoint{8.311372in}{0.974302in}}%
\pgfpathlineto{\pgfqpoint{8.323242in}{0.974302in}}%
\pgfpathlineto{\pgfqpoint{8.323242in}{0.964611in}}%
\pgfpathlineto{\pgfqpoint{8.314020in}{0.946599in}}%
\pgfpathlineto{\pgfqpoint{8.306761in}{0.946599in}}%
\pgfpathlineto{\pgfqpoint{8.311372in}{0.964611in}}%
\pgfpathlineto{\pgfqpoint{8.311372in}{0.974302in}}%
\pgfpathlineto{\pgfqpoint{8.311372in}{0.974302in}}%
\pgfpathclose%
\pgfusepath{fill}%
\end{pgfscope}%
\begin{pgfscope}%
\pgfsetbuttcap%
\pgfsetroundjoin%
\definecolor{currentfill}{rgb}{0.309804,0.372549,0.419608}%
\pgfsetfillcolor{currentfill}%
\pgfsetlinewidth{0.000000pt}%
\definecolor{currentstroke}{rgb}{0.309804,0.372549,0.419608}%
\pgfsetstrokecolor{currentstroke}%
\pgfsetdash{}{0pt}%
\pgfpathmoveto{\pgfqpoint{8.450155in}{1.013478in}}%
\pgfpathlineto{\pgfqpoint{8.450155in}{1.047593in}}%
\pgfpathlineto{\pgfqpoint{8.460512in}{1.047593in}}%
\pgfpathlineto{\pgfqpoint{8.460512in}{0.960000in}}%
\pgfpathlineto{\pgfqpoint{8.450155in}{0.960000in}}%
\pgfpathlineto{\pgfqpoint{8.450155in}{0.969456in}}%
\pgfpathquadraticcurveto{\pgfqpoint{8.446895in}{0.963837in}}{\pgfqpoint{8.441905in}{0.961099in}}%
\pgfpathquadraticcurveto{\pgfqpoint{8.436934in}{0.958361in}}{\pgfqpoint{8.429945in}{0.958361in}}%
\pgfpathquadraticcurveto{\pgfqpoint{8.418525in}{0.958361in}}{\pgfqpoint{8.411339in}{0.967475in}}%
\pgfpathquadraticcurveto{\pgfqpoint{8.404170in}{0.976607in}}{\pgfqpoint{8.404170in}{0.991467in}}%
\pgfpathquadraticcurveto{\pgfqpoint{8.404170in}{1.006327in}}{\pgfqpoint{8.411339in}{1.015441in}}%
\pgfpathquadraticcurveto{\pgfqpoint{8.418525in}{1.024556in}}{\pgfqpoint{8.429945in}{1.024556in}}%
\pgfpathquadraticcurveto{\pgfqpoint{8.436934in}{1.024556in}}{\pgfqpoint{8.441905in}{1.021818in}}%
\pgfpathquadraticcurveto{\pgfqpoint{8.446895in}{1.019098in}}{\pgfqpoint{8.450155in}{1.013478in}}%
\pgfpathlineto{\pgfqpoint{8.450155in}{1.013478in}}%
\pgfpathclose%
\pgfpathmoveto{\pgfqpoint{8.414869in}{0.991467in}}%
\pgfpathquadraticcurveto{\pgfqpoint{8.414869in}{0.980048in}}{\pgfqpoint{8.419570in}{0.973545in}}%
\pgfpathquadraticcurveto{\pgfqpoint{8.424271in}{0.967043in}}{\pgfqpoint{8.432485in}{0.967043in}}%
\pgfpathquadraticcurveto{\pgfqpoint{8.440698in}{0.967043in}}{\pgfqpoint{8.445418in}{0.973545in}}%
\pgfpathquadraticcurveto{\pgfqpoint{8.450155in}{0.980048in}}{\pgfqpoint{8.450155in}{0.991467in}}%
\pgfpathquadraticcurveto{\pgfqpoint{8.450155in}{1.002887in}}{\pgfqpoint{8.445418in}{1.009389in}}%
\pgfpathquadraticcurveto{\pgfqpoint{8.440698in}{1.015892in}}{\pgfqpoint{8.432485in}{1.015892in}}%
\pgfpathquadraticcurveto{\pgfqpoint{8.424271in}{1.015892in}}{\pgfqpoint{8.419570in}{1.009389in}}%
\pgfpathquadraticcurveto{\pgfqpoint{8.414869in}{1.002887in}}{\pgfqpoint{8.414869in}{0.991467in}}%
\pgfpathlineto{\pgfqpoint{8.414869in}{0.991467in}}%
\pgfpathclose%
\pgfpathmoveto{\pgfqpoint{8.481856in}{1.023043in}}%
\pgfpathlineto{\pgfqpoint{8.492213in}{1.023043in}}%
\pgfpathlineto{\pgfqpoint{8.492213in}{0.960000in}}%
\pgfpathlineto{\pgfqpoint{8.481856in}{0.960000in}}%
\pgfpathlineto{\pgfqpoint{8.481856in}{1.023043in}}%
\pgfpathlineto{\pgfqpoint{8.481856in}{1.023043in}}%
\pgfpathclose%
\pgfpathmoveto{\pgfqpoint{8.481856in}{1.047593in}}%
\pgfpathlineto{\pgfqpoint{8.492213in}{1.047593in}}%
\pgfpathlineto{\pgfqpoint{8.492213in}{1.034462in}}%
\pgfpathlineto{\pgfqpoint{8.481856in}{1.034462in}}%
\pgfpathlineto{\pgfqpoint{8.481856in}{1.047593in}}%
\pgfpathlineto{\pgfqpoint{8.481856in}{1.047593in}}%
\pgfpathclose%
\pgfpathmoveto{\pgfqpoint{8.554067in}{1.021187in}}%
\pgfpathlineto{\pgfqpoint{8.554067in}{1.011389in}}%
\pgfpathquadraticcurveto{\pgfqpoint{8.549690in}{1.013640in}}{\pgfqpoint{8.544953in}{1.014757in}}%
\pgfpathquadraticcurveto{\pgfqpoint{8.540234in}{1.015892in}}{\pgfqpoint{8.535154in}{1.015892in}}%
\pgfpathquadraticcurveto{\pgfqpoint{8.527445in}{1.015892in}}{\pgfqpoint{8.523590in}{1.013532in}}%
\pgfpathquadraticcurveto{\pgfqpoint{8.519736in}{1.011173in}}{\pgfqpoint{8.519736in}{1.006435in}}%
\pgfpathquadraticcurveto{\pgfqpoint{8.519736in}{1.002833in}}{\pgfqpoint{8.522492in}{1.000780in}}%
\pgfpathquadraticcurveto{\pgfqpoint{8.525247in}{0.998726in}}{\pgfqpoint{8.533587in}{0.996871in}}%
\pgfpathlineto{\pgfqpoint{8.537135in}{0.996078in}}%
\pgfpathquadraticcurveto{\pgfqpoint{8.548159in}{0.993719in}}{\pgfqpoint{8.552806in}{0.989414in}}%
\pgfpathquadraticcurveto{\pgfqpoint{8.557453in}{0.985109in}}{\pgfqpoint{8.557453in}{0.977400in}}%
\pgfpathquadraticcurveto{\pgfqpoint{8.557453in}{0.968610in}}{\pgfqpoint{8.550500in}{0.963476in}}%
\pgfpathquadraticcurveto{\pgfqpoint{8.543548in}{0.958361in}}{\pgfqpoint{8.531390in}{0.958361in}}%
\pgfpathquadraticcurveto{\pgfqpoint{8.526328in}{0.958361in}}{\pgfqpoint{8.520834in}{0.959352in}}%
\pgfpathquadraticcurveto{\pgfqpoint{8.515341in}{0.960342in}}{\pgfqpoint{8.509271in}{0.962306in}}%
\pgfpathlineto{\pgfqpoint{8.509271in}{0.973005in}}%
\pgfpathquadraticcurveto{\pgfqpoint{8.515017in}{0.970015in}}{\pgfqpoint{8.520582in}{0.968520in}}%
\pgfpathquadraticcurveto{\pgfqpoint{8.526148in}{0.967043in}}{\pgfqpoint{8.531624in}{0.967043in}}%
\pgfpathquadraticcurveto{\pgfqpoint{8.538937in}{0.967043in}}{\pgfqpoint{8.542863in}{0.969546in}}%
\pgfpathquadraticcurveto{\pgfqpoint{8.546808in}{0.972050in}}{\pgfqpoint{8.546808in}{0.976607in}}%
\pgfpathquadraticcurveto{\pgfqpoint{8.546808in}{0.980822in}}{\pgfqpoint{8.543962in}{0.983074in}}%
\pgfpathquadraticcurveto{\pgfqpoint{8.541134in}{0.985325in}}{\pgfqpoint{8.531498in}{0.987414in}}%
\pgfpathlineto{\pgfqpoint{8.527895in}{0.988261in}}%
\pgfpathquadraticcurveto{\pgfqpoint{8.518277in}{0.990278in}}{\pgfqpoint{8.513990in}{0.994475in}}%
\pgfpathquadraticcurveto{\pgfqpoint{8.509721in}{0.998672in}}{\pgfqpoint{8.509721in}{1.005985in}}%
\pgfpathquadraticcurveto{\pgfqpoint{8.509721in}{1.014883in}}{\pgfqpoint{8.516025in}{1.019710in}}%
\pgfpathquadraticcurveto{\pgfqpoint{8.522329in}{1.024556in}}{\pgfqpoint{8.533929in}{1.024556in}}%
\pgfpathquadraticcurveto{\pgfqpoint{8.539657in}{1.024556in}}{\pgfqpoint{8.544719in}{1.023709in}}%
\pgfpathquadraticcurveto{\pgfqpoint{8.549798in}{1.022880in}}{\pgfqpoint{8.554067in}{1.021187in}}%
\pgfpathlineto{\pgfqpoint{8.554067in}{1.021187in}}%
\pgfpathclose%
\pgfpathmoveto{\pgfqpoint{8.583949in}{0.969456in}}%
\pgfpathlineto{\pgfqpoint{8.583949in}{0.936026in}}%
\pgfpathlineto{\pgfqpoint{8.573538in}{0.936026in}}%
\pgfpathlineto{\pgfqpoint{8.573538in}{1.023043in}}%
\pgfpathlineto{\pgfqpoint{8.583949in}{1.023043in}}%
\pgfpathlineto{\pgfqpoint{8.583949in}{1.013478in}}%
\pgfpathquadraticcurveto{\pgfqpoint{8.587227in}{1.019098in}}{\pgfqpoint{8.592199in}{1.021818in}}%
\pgfpathquadraticcurveto{\pgfqpoint{8.597188in}{1.024556in}}{\pgfqpoint{8.604105in}{1.024556in}}%
\pgfpathquadraticcurveto{\pgfqpoint{8.615596in}{1.024556in}}{\pgfqpoint{8.622765in}{1.015441in}}%
\pgfpathquadraticcurveto{\pgfqpoint{8.629952in}{1.006327in}}{\pgfqpoint{8.629952in}{0.991467in}}%
\pgfpathquadraticcurveto{\pgfqpoint{8.629952in}{0.976607in}}{\pgfqpoint{8.622765in}{0.967475in}}%
\pgfpathquadraticcurveto{\pgfqpoint{8.615596in}{0.958361in}}{\pgfqpoint{8.604105in}{0.958361in}}%
\pgfpathquadraticcurveto{\pgfqpoint{8.597188in}{0.958361in}}{\pgfqpoint{8.592199in}{0.961099in}}%
\pgfpathquadraticcurveto{\pgfqpoint{8.587227in}{0.963837in}}{\pgfqpoint{8.583949in}{0.969456in}}%
\pgfpathlineto{\pgfqpoint{8.583949in}{0.969456in}}%
\pgfpathclose%
\pgfpathmoveto{\pgfqpoint{8.619199in}{0.991467in}}%
\pgfpathquadraticcurveto{\pgfqpoint{8.619199in}{1.002887in}}{\pgfqpoint{8.614498in}{1.009389in}}%
\pgfpathquadraticcurveto{\pgfqpoint{8.609796in}{1.015892in}}{\pgfqpoint{8.601583in}{1.015892in}}%
\pgfpathquadraticcurveto{\pgfqpoint{8.593351in}{1.015892in}}{\pgfqpoint{8.588650in}{1.009389in}}%
\pgfpathquadraticcurveto{\pgfqpoint{8.583949in}{1.002887in}}{\pgfqpoint{8.583949in}{0.991467in}}%
\pgfpathquadraticcurveto{\pgfqpoint{8.583949in}{0.980048in}}{\pgfqpoint{8.588650in}{0.973545in}}%
\pgfpathquadraticcurveto{\pgfqpoint{8.593351in}{0.967043in}}{\pgfqpoint{8.601583in}{0.967043in}}%
\pgfpathquadraticcurveto{\pgfqpoint{8.609796in}{0.967043in}}{\pgfqpoint{8.614498in}{0.973545in}}%
\pgfpathquadraticcurveto{\pgfqpoint{8.619199in}{0.980048in}}{\pgfqpoint{8.619199in}{0.991467in}}%
\pgfpathlineto{\pgfqpoint{8.619199in}{0.991467in}}%
\pgfpathclose%
\pgfpathmoveto{\pgfqpoint{8.647118in}{1.047593in}}%
\pgfpathlineto{\pgfqpoint{8.657475in}{1.047593in}}%
\pgfpathlineto{\pgfqpoint{8.657475in}{0.960000in}}%
\pgfpathlineto{\pgfqpoint{8.647118in}{0.960000in}}%
\pgfpathlineto{\pgfqpoint{8.647118in}{1.047593in}}%
\pgfpathlineto{\pgfqpoint{8.647118in}{1.047593in}}%
\pgfpathclose%
\pgfpathmoveto{\pgfqpoint{8.707801in}{0.991683in}}%
\pgfpathquadraticcurveto{\pgfqpoint{8.695246in}{0.991683in}}{\pgfqpoint{8.690401in}{0.988819in}}%
\pgfpathquadraticcurveto{\pgfqpoint{8.685556in}{0.985956in}}{\pgfqpoint{8.685556in}{0.979021in}}%
\pgfpathquadraticcurveto{\pgfqpoint{8.685556in}{0.973509in}}{\pgfqpoint{8.689194in}{0.970267in}}%
\pgfpathquadraticcurveto{\pgfqpoint{8.692832in}{0.967043in}}{\pgfqpoint{8.699065in}{0.967043in}}%
\pgfpathquadraticcurveto{\pgfqpoint{8.707692in}{0.967043in}}{\pgfqpoint{8.712898in}{0.973149in}}%
\pgfpathquadraticcurveto{\pgfqpoint{8.718104in}{0.979255in}}{\pgfqpoint{8.718104in}{0.989378in}}%
\pgfpathlineto{\pgfqpoint{8.718104in}{0.991683in}}%
\pgfpathlineto{\pgfqpoint{8.707801in}{0.991683in}}%
\pgfpathlineto{\pgfqpoint{8.707801in}{0.991683in}}%
\pgfpathclose%
\pgfpathmoveto{\pgfqpoint{8.728461in}{0.995970in}}%
\pgfpathlineto{\pgfqpoint{8.728461in}{0.960000in}}%
\pgfpathlineto{\pgfqpoint{8.718104in}{0.960000in}}%
\pgfpathlineto{\pgfqpoint{8.718104in}{0.969564in}}%
\pgfpathquadraticcurveto{\pgfqpoint{8.714555in}{0.963837in}}{\pgfqpoint{8.709260in}{0.961099in}}%
\pgfpathquadraticcurveto{\pgfqpoint{8.703964in}{0.958361in}}{\pgfqpoint{8.696309in}{0.958361in}}%
\pgfpathquadraticcurveto{\pgfqpoint{8.686636in}{0.958361in}}{\pgfqpoint{8.680908in}{0.963801in}}%
\pgfpathquadraticcurveto{\pgfqpoint{8.675199in}{0.969240in}}{\pgfqpoint{8.675199in}{0.978354in}}%
\pgfpathquadraticcurveto{\pgfqpoint{8.675199in}{0.988982in}}{\pgfqpoint{8.682313in}{0.994385in}}%
\pgfpathquadraticcurveto{\pgfqpoint{8.689446in}{0.999789in}}{\pgfqpoint{8.703568in}{0.999789in}}%
\pgfpathlineto{\pgfqpoint{8.718104in}{0.999789in}}%
\pgfpathlineto{\pgfqpoint{8.718104in}{1.000816in}}%
\pgfpathquadraticcurveto{\pgfqpoint{8.718104in}{1.007966in}}{\pgfqpoint{8.713402in}{1.011875in}}%
\pgfpathquadraticcurveto{\pgfqpoint{8.708701in}{1.015784in}}{\pgfqpoint{8.700199in}{1.015784in}}%
\pgfpathquadraticcurveto{\pgfqpoint{8.694796in}{1.015784in}}{\pgfqpoint{8.689662in}{1.014487in}}%
\pgfpathquadraticcurveto{\pgfqpoint{8.684547in}{1.013190in}}{\pgfqpoint{8.679828in}{1.010596in}}%
\pgfpathlineto{\pgfqpoint{8.679828in}{1.020179in}}%
\pgfpathquadraticcurveto{\pgfqpoint{8.685502in}{1.022376in}}{\pgfqpoint{8.690851in}{1.023457in}}%
\pgfpathquadraticcurveto{\pgfqpoint{8.696201in}{1.024556in}}{\pgfqpoint{8.701262in}{1.024556in}}%
\pgfpathquadraticcurveto{\pgfqpoint{8.714951in}{1.024556in}}{\pgfqpoint{8.721706in}{1.017459in}}%
\pgfpathquadraticcurveto{\pgfqpoint{8.728461in}{1.010380in}}{\pgfqpoint{8.728461in}{0.995970in}}%
\pgfpathlineto{\pgfqpoint{8.728461in}{0.995970in}}%
\pgfpathclose%
\pgfpathmoveto{\pgfqpoint{8.776013in}{0.954146in}}%
\pgfpathquadraticcurveto{\pgfqpoint{8.771636in}{0.942888in}}{\pgfqpoint{8.767457in}{0.939466in}}%
\pgfpathquadraticcurveto{\pgfqpoint{8.763296in}{0.936026in}}{\pgfqpoint{8.756325in}{0.936026in}}%
\pgfpathlineto{\pgfqpoint{8.748040in}{0.936026in}}%
\pgfpathlineto{\pgfqpoint{8.748040in}{0.944690in}}%
\pgfpathlineto{\pgfqpoint{8.754128in}{0.944690in}}%
\pgfpathquadraticcurveto{\pgfqpoint{8.758397in}{0.944690in}}{\pgfqpoint{8.760756in}{0.946725in}}%
\pgfpathquadraticcurveto{\pgfqpoint{8.763134in}{0.948742in}}{\pgfqpoint{8.765998in}{0.956289in}}%
\pgfpathlineto{\pgfqpoint{8.767853in}{0.961009in}}%
\pgfpathlineto{\pgfqpoint{8.742366in}{1.023043in}}%
\pgfpathlineto{\pgfqpoint{8.753335in}{1.023043in}}%
\pgfpathlineto{\pgfqpoint{8.773041in}{0.973743in}}%
\pgfpathlineto{\pgfqpoint{8.792746in}{1.023043in}}%
\pgfpathlineto{\pgfqpoint{8.803715in}{1.023043in}}%
\pgfpathlineto{\pgfqpoint{8.776013in}{0.954146in}}%
\pgfpathlineto{\pgfqpoint{8.776013in}{0.954146in}}%
\pgfpathclose%
\pgfpathmoveto{\pgfqpoint{8.871945in}{0.994115in}}%
\pgfpathlineto{\pgfqpoint{8.871945in}{0.989054in}}%
\pgfpathlineto{\pgfqpoint{8.824321in}{0.989054in}}%
\pgfpathquadraticcurveto{\pgfqpoint{8.825006in}{0.978354in}}{\pgfqpoint{8.830770in}{0.972753in}}%
\pgfpathquadraticcurveto{\pgfqpoint{8.836533in}{0.967151in}}{\pgfqpoint{8.846836in}{0.967151in}}%
\pgfpathquadraticcurveto{\pgfqpoint{8.852798in}{0.967151in}}{\pgfqpoint{8.858400in}{0.968610in}}%
\pgfpathquadraticcurveto{\pgfqpoint{8.864002in}{0.970069in}}{\pgfqpoint{8.869532in}{0.973005in}}%
\pgfpathlineto{\pgfqpoint{8.869532in}{0.963206in}}%
\pgfpathquadraticcurveto{\pgfqpoint{8.863948in}{0.960847in}}{\pgfqpoint{8.858094in}{0.959604in}}%
\pgfpathquadraticcurveto{\pgfqpoint{8.852240in}{0.958361in}}{\pgfqpoint{8.846224in}{0.958361in}}%
\pgfpathquadraticcurveto{\pgfqpoint{8.831130in}{0.958361in}}{\pgfqpoint{8.822322in}{0.967133in}}%
\pgfpathquadraticcurveto{\pgfqpoint{8.813514in}{0.975923in}}{\pgfqpoint{8.813514in}{0.990909in}}%
\pgfpathquadraticcurveto{\pgfqpoint{8.813514in}{1.006381in}}{\pgfqpoint{8.821872in}{1.015459in}}%
\pgfpathquadraticcurveto{\pgfqpoint{8.830229in}{1.024556in}}{\pgfqpoint{8.844423in}{1.024556in}}%
\pgfpathquadraticcurveto{\pgfqpoint{8.857139in}{1.024556in}}{\pgfqpoint{8.864542in}{1.016360in}}%
\pgfpathquadraticcurveto{\pgfqpoint{8.871945in}{1.008183in}}{\pgfqpoint{8.871945in}{0.994115in}}%
\pgfpathlineto{\pgfqpoint{8.871945in}{0.994115in}}%
\pgfpathclose%
\pgfpathmoveto{\pgfqpoint{8.861588in}{0.997159in}}%
\pgfpathquadraticcurveto{\pgfqpoint{8.861480in}{1.005643in}}{\pgfqpoint{8.856833in}{1.010704in}}%
\pgfpathquadraticcurveto{\pgfqpoint{8.852186in}{1.015784in}}{\pgfqpoint{8.844531in}{1.015784in}}%
\pgfpathquadraticcurveto{\pgfqpoint{8.835867in}{1.015784in}}{\pgfqpoint{8.830661in}{1.010884in}}%
\pgfpathquadraticcurveto{\pgfqpoint{8.825456in}{1.005985in}}{\pgfqpoint{8.824663in}{0.997087in}}%
\pgfpathlineto{\pgfqpoint{8.861588in}{0.997159in}}%
\pgfpathlineto{\pgfqpoint{8.861588in}{0.997159in}}%
\pgfpathclose%
\pgfpathmoveto{\pgfqpoint{8.930431in}{1.013478in}}%
\pgfpathlineto{\pgfqpoint{8.930431in}{1.047593in}}%
\pgfpathlineto{\pgfqpoint{8.940788in}{1.047593in}}%
\pgfpathlineto{\pgfqpoint{8.940788in}{0.960000in}}%
\pgfpathlineto{\pgfqpoint{8.930431in}{0.960000in}}%
\pgfpathlineto{\pgfqpoint{8.930431in}{0.969456in}}%
\pgfpathquadraticcurveto{\pgfqpoint{8.927171in}{0.963837in}}{\pgfqpoint{8.922181in}{0.961099in}}%
\pgfpathquadraticcurveto{\pgfqpoint{8.917210in}{0.958361in}}{\pgfqpoint{8.910221in}{0.958361in}}%
\pgfpathquadraticcurveto{\pgfqpoint{8.898801in}{0.958361in}}{\pgfqpoint{8.891615in}{0.967475in}}%
\pgfpathquadraticcurveto{\pgfqpoint{8.884446in}{0.976607in}}{\pgfqpoint{8.884446in}{0.991467in}}%
\pgfpathquadraticcurveto{\pgfqpoint{8.884446in}{1.006327in}}{\pgfqpoint{8.891615in}{1.015441in}}%
\pgfpathquadraticcurveto{\pgfqpoint{8.898801in}{1.024556in}}{\pgfqpoint{8.910221in}{1.024556in}}%
\pgfpathquadraticcurveto{\pgfqpoint{8.917210in}{1.024556in}}{\pgfqpoint{8.922181in}{1.021818in}}%
\pgfpathquadraticcurveto{\pgfqpoint{8.927171in}{1.019098in}}{\pgfqpoint{8.930431in}{1.013478in}}%
\pgfpathlineto{\pgfqpoint{8.930431in}{1.013478in}}%
\pgfpathclose%
\pgfpathmoveto{\pgfqpoint{8.895145in}{0.991467in}}%
\pgfpathquadraticcurveto{\pgfqpoint{8.895145in}{0.980048in}}{\pgfqpoint{8.899846in}{0.973545in}}%
\pgfpathquadraticcurveto{\pgfqpoint{8.904547in}{0.967043in}}{\pgfqpoint{8.912761in}{0.967043in}}%
\pgfpathquadraticcurveto{\pgfqpoint{8.920974in}{0.967043in}}{\pgfqpoint{8.925694in}{0.973545in}}%
\pgfpathquadraticcurveto{\pgfqpoint{8.930431in}{0.980048in}}{\pgfqpoint{8.930431in}{0.991467in}}%
\pgfpathquadraticcurveto{\pgfqpoint{8.930431in}{1.002887in}}{\pgfqpoint{8.925694in}{1.009389in}}%
\pgfpathquadraticcurveto{\pgfqpoint{8.920974in}{1.015892in}}{\pgfqpoint{8.912761in}{1.015892in}}%
\pgfpathquadraticcurveto{\pgfqpoint{8.904547in}{1.015892in}}{\pgfqpoint{8.899846in}{1.009389in}}%
\pgfpathquadraticcurveto{\pgfqpoint{8.895145in}{1.002887in}}{\pgfqpoint{8.895145in}{0.991467in}}%
\pgfpathlineto{\pgfqpoint{8.895145in}{0.991467in}}%
\pgfpathclose%
\pgfusepath{fill}%
\end{pgfscope}%
\begin{pgfscope}%
\pgfsetbuttcap%
\pgfsetroundjoin%
\definecolor{currentfill}{rgb}{0.309804,0.372549,0.419608}%
\pgfsetfillcolor{currentfill}%
\pgfsetlinewidth{0.000000pt}%
\definecolor{currentstroke}{rgb}{0.309804,0.372549,0.419608}%
\pgfsetstrokecolor{currentstroke}%
\pgfsetdash{}{0pt}%
\pgfpathmoveto{\pgfqpoint{9.070048in}{0.992260in}}%
\pgfpathquadraticcurveto{\pgfqpoint{9.070048in}{1.003517in}}{\pgfqpoint{9.065401in}{1.009696in}}%
\pgfpathquadraticcurveto{\pgfqpoint{9.060772in}{1.015892in}}{\pgfqpoint{9.052378in}{1.015892in}}%
\pgfpathquadraticcurveto{\pgfqpoint{9.044057in}{1.015892in}}{\pgfqpoint{9.039409in}{1.009696in}}%
\pgfpathquadraticcurveto{\pgfqpoint{9.034762in}{1.003517in}}{\pgfqpoint{9.034762in}{0.992260in}}%
\pgfpathquadraticcurveto{\pgfqpoint{9.034762in}{0.981056in}}{\pgfqpoint{9.039409in}{0.974860in}}%
\pgfpathquadraticcurveto{\pgfqpoint{9.044057in}{0.968664in}}{\pgfqpoint{9.052378in}{0.968664in}}%
\pgfpathquadraticcurveto{\pgfqpoint{9.060772in}{0.968664in}}{\pgfqpoint{9.065401in}{0.974860in}}%
\pgfpathquadraticcurveto{\pgfqpoint{9.070048in}{0.981056in}}{\pgfqpoint{9.070048in}{0.992260in}}%
\pgfpathlineto{\pgfqpoint{9.070048in}{0.992260in}}%
\pgfpathclose%
\pgfpathmoveto{\pgfqpoint{9.080405in}{0.967817in}}%
\pgfpathquadraticcurveto{\pgfqpoint{9.080405in}{0.951732in}}{\pgfqpoint{9.073254in}{0.943879in}}%
\pgfpathquadraticcurveto{\pgfqpoint{9.066121in}{0.936026in}}{\pgfqpoint{9.051369in}{0.936026in}}%
\pgfpathquadraticcurveto{\pgfqpoint{9.045912in}{0.936026in}}{\pgfqpoint{9.041067in}{0.936836in}}%
\pgfpathquadraticcurveto{\pgfqpoint{9.036221in}{0.937647in}}{\pgfqpoint{9.031664in}{0.939340in}}%
\pgfpathlineto{\pgfqpoint{9.031664in}{0.949409in}}%
\pgfpathquadraticcurveto{\pgfqpoint{9.036221in}{0.946941in}}{\pgfqpoint{9.040670in}{0.945770in}}%
\pgfpathquadraticcurveto{\pgfqpoint{9.045119in}{0.944582in}}{\pgfqpoint{9.049730in}{0.944582in}}%
\pgfpathquadraticcurveto{\pgfqpoint{9.059925in}{0.944582in}}{\pgfqpoint{9.064987in}{0.949895in}}%
\pgfpathquadraticcurveto{\pgfqpoint{9.070048in}{0.955209in}}{\pgfqpoint{9.070048in}{0.965962in}}%
\pgfpathlineto{\pgfqpoint{9.070048in}{0.971095in}}%
\pgfpathquadraticcurveto{\pgfqpoint{9.066842in}{0.965512in}}{\pgfqpoint{9.061835in}{0.962756in}}%
\pgfpathquadraticcurveto{\pgfqpoint{9.056827in}{0.960000in}}{\pgfqpoint{9.049838in}{0.960000in}}%
\pgfpathquadraticcurveto{\pgfqpoint{9.038257in}{0.960000in}}{\pgfqpoint{9.031160in}{0.968826in}}%
\pgfpathquadraticcurveto{\pgfqpoint{9.024063in}{0.977670in}}{\pgfqpoint{9.024063in}{0.992260in}}%
\pgfpathquadraticcurveto{\pgfqpoint{9.024063in}{1.006886in}}{\pgfqpoint{9.031160in}{1.015712in}}%
\pgfpathquadraticcurveto{\pgfqpoint{9.038257in}{1.024556in}}{\pgfqpoint{9.049838in}{1.024556in}}%
\pgfpathquadraticcurveto{\pgfqpoint{9.056827in}{1.024556in}}{\pgfqpoint{9.061835in}{1.021800in}}%
\pgfpathquadraticcurveto{\pgfqpoint{9.066842in}{1.019044in}}{\pgfqpoint{9.070048in}{1.013478in}}%
\pgfpathlineto{\pgfqpoint{9.070048in}{1.023043in}}%
\pgfpathlineto{\pgfqpoint{9.080405in}{1.023043in}}%
\pgfpathlineto{\pgfqpoint{9.080405in}{0.967817in}}%
\pgfpathlineto{\pgfqpoint{9.080405in}{0.967817in}}%
\pgfpathclose%
\pgfpathmoveto{\pgfqpoint{9.130407in}{0.991683in}}%
\pgfpathquadraticcurveto{\pgfqpoint{9.117852in}{0.991683in}}{\pgfqpoint{9.113007in}{0.988819in}}%
\pgfpathquadraticcurveto{\pgfqpoint{9.108162in}{0.985956in}}{\pgfqpoint{9.108162in}{0.979021in}}%
\pgfpathquadraticcurveto{\pgfqpoint{9.108162in}{0.973509in}}{\pgfqpoint{9.111800in}{0.970267in}}%
\pgfpathquadraticcurveto{\pgfqpoint{9.115439in}{0.967043in}}{\pgfqpoint{9.121671in}{0.967043in}}%
\pgfpathquadraticcurveto{\pgfqpoint{9.130299in}{0.967043in}}{\pgfqpoint{9.135504in}{0.973149in}}%
\pgfpathquadraticcurveto{\pgfqpoint{9.140710in}{0.979255in}}{\pgfqpoint{9.140710in}{0.989378in}}%
\pgfpathlineto{\pgfqpoint{9.140710in}{0.991683in}}%
\pgfpathlineto{\pgfqpoint{9.130407in}{0.991683in}}%
\pgfpathlineto{\pgfqpoint{9.130407in}{0.991683in}}%
\pgfpathclose%
\pgfpathmoveto{\pgfqpoint{9.151067in}{0.995970in}}%
\pgfpathlineto{\pgfqpoint{9.151067in}{0.960000in}}%
\pgfpathlineto{\pgfqpoint{9.140710in}{0.960000in}}%
\pgfpathlineto{\pgfqpoint{9.140710in}{0.969564in}}%
\pgfpathquadraticcurveto{\pgfqpoint{9.137161in}{0.963837in}}{\pgfqpoint{9.131866in}{0.961099in}}%
\pgfpathquadraticcurveto{\pgfqpoint{9.126570in}{0.958361in}}{\pgfqpoint{9.118915in}{0.958361in}}%
\pgfpathquadraticcurveto{\pgfqpoint{9.109243in}{0.958361in}}{\pgfqpoint{9.103515in}{0.963801in}}%
\pgfpathquadraticcurveto{\pgfqpoint{9.097805in}{0.969240in}}{\pgfqpoint{9.097805in}{0.978354in}}%
\pgfpathquadraticcurveto{\pgfqpoint{9.097805in}{0.988982in}}{\pgfqpoint{9.104920in}{0.994385in}}%
\pgfpathquadraticcurveto{\pgfqpoint{9.112052in}{0.999789in}}{\pgfqpoint{9.126174in}{0.999789in}}%
\pgfpathlineto{\pgfqpoint{9.140710in}{0.999789in}}%
\pgfpathlineto{\pgfqpoint{9.140710in}{1.000816in}}%
\pgfpathquadraticcurveto{\pgfqpoint{9.140710in}{1.007966in}}{\pgfqpoint{9.136009in}{1.011875in}}%
\pgfpathquadraticcurveto{\pgfqpoint{9.131307in}{1.015784in}}{\pgfqpoint{9.122806in}{1.015784in}}%
\pgfpathquadraticcurveto{\pgfqpoint{9.117402in}{1.015784in}}{\pgfqpoint{9.112269in}{1.014487in}}%
\pgfpathquadraticcurveto{\pgfqpoint{9.107153in}{1.013190in}}{\pgfqpoint{9.102434in}{1.010596in}}%
\pgfpathlineto{\pgfqpoint{9.102434in}{1.020179in}}%
\pgfpathquadraticcurveto{\pgfqpoint{9.108108in}{1.022376in}}{\pgfqpoint{9.113457in}{1.023457in}}%
\pgfpathquadraticcurveto{\pgfqpoint{9.118807in}{1.024556in}}{\pgfqpoint{9.123868in}{1.024556in}}%
\pgfpathquadraticcurveto{\pgfqpoint{9.137558in}{1.024556in}}{\pgfqpoint{9.144312in}{1.017459in}}%
\pgfpathquadraticcurveto{\pgfqpoint{9.151067in}{1.010380in}}{\pgfqpoint{9.151067in}{0.995970in}}%
\pgfpathlineto{\pgfqpoint{9.151067in}{0.995970in}}%
\pgfpathclose%
\pgfpathmoveto{\pgfqpoint{9.182408in}{0.969456in}}%
\pgfpathlineto{\pgfqpoint{9.182408in}{0.936026in}}%
\pgfpathlineto{\pgfqpoint{9.171997in}{0.936026in}}%
\pgfpathlineto{\pgfqpoint{9.171997in}{1.023043in}}%
\pgfpathlineto{\pgfqpoint{9.182408in}{1.023043in}}%
\pgfpathlineto{\pgfqpoint{9.182408in}{1.013478in}}%
\pgfpathquadraticcurveto{\pgfqpoint{9.185686in}{1.019098in}}{\pgfqpoint{9.190657in}{1.021818in}}%
\pgfpathquadraticcurveto{\pgfqpoint{9.195647in}{1.024556in}}{\pgfqpoint{9.202563in}{1.024556in}}%
\pgfpathquadraticcurveto{\pgfqpoint{9.214055in}{1.024556in}}{\pgfqpoint{9.221224in}{1.015441in}}%
\pgfpathquadraticcurveto{\pgfqpoint{9.228411in}{1.006327in}}{\pgfqpoint{9.228411in}{0.991467in}}%
\pgfpathquadraticcurveto{\pgfqpoint{9.228411in}{0.976607in}}{\pgfqpoint{9.221224in}{0.967475in}}%
\pgfpathquadraticcurveto{\pgfqpoint{9.214055in}{0.958361in}}{\pgfqpoint{9.202563in}{0.958361in}}%
\pgfpathquadraticcurveto{\pgfqpoint{9.195647in}{0.958361in}}{\pgfqpoint{9.190657in}{0.961099in}}%
\pgfpathquadraticcurveto{\pgfqpoint{9.185686in}{0.963837in}}{\pgfqpoint{9.182408in}{0.969456in}}%
\pgfpathlineto{\pgfqpoint{9.182408in}{0.969456in}}%
\pgfpathclose%
\pgfpathmoveto{\pgfqpoint{9.217658in}{0.991467in}}%
\pgfpathquadraticcurveto{\pgfqpoint{9.217658in}{1.002887in}}{\pgfqpoint{9.212956in}{1.009389in}}%
\pgfpathquadraticcurveto{\pgfqpoint{9.208255in}{1.015892in}}{\pgfqpoint{9.200042in}{1.015892in}}%
\pgfpathquadraticcurveto{\pgfqpoint{9.191810in}{1.015892in}}{\pgfqpoint{9.187109in}{1.009389in}}%
\pgfpathquadraticcurveto{\pgfqpoint{9.182408in}{1.002887in}}{\pgfqpoint{9.182408in}{0.991467in}}%
\pgfpathquadraticcurveto{\pgfqpoint{9.182408in}{0.980048in}}{\pgfqpoint{9.187109in}{0.973545in}}%
\pgfpathquadraticcurveto{\pgfqpoint{9.191810in}{0.967043in}}{\pgfqpoint{9.200042in}{0.967043in}}%
\pgfpathquadraticcurveto{\pgfqpoint{9.208255in}{0.967043in}}{\pgfqpoint{9.212956in}{0.973545in}}%
\pgfpathquadraticcurveto{\pgfqpoint{9.217658in}{0.980048in}}{\pgfqpoint{9.217658in}{0.991467in}}%
\pgfpathlineto{\pgfqpoint{9.217658in}{0.991467in}}%
\pgfpathclose%
\pgfpathmoveto{\pgfqpoint{9.285762in}{1.021187in}}%
\pgfpathlineto{\pgfqpoint{9.285762in}{1.011389in}}%
\pgfpathquadraticcurveto{\pgfqpoint{9.281385in}{1.013640in}}{\pgfqpoint{9.276647in}{1.014757in}}%
\pgfpathquadraticcurveto{\pgfqpoint{9.271928in}{1.015892in}}{\pgfqpoint{9.266849in}{1.015892in}}%
\pgfpathquadraticcurveto{\pgfqpoint{9.259140in}{1.015892in}}{\pgfqpoint{9.255285in}{1.013532in}}%
\pgfpathquadraticcurveto{\pgfqpoint{9.251430in}{1.011173in}}{\pgfqpoint{9.251430in}{1.006435in}}%
\pgfpathquadraticcurveto{\pgfqpoint{9.251430in}{1.002833in}}{\pgfqpoint{9.254186in}{1.000780in}}%
\pgfpathquadraticcurveto{\pgfqpoint{9.256942in}{0.998726in}}{\pgfqpoint{9.265282in}{0.996871in}}%
\pgfpathlineto{\pgfqpoint{9.268830in}{0.996078in}}%
\pgfpathquadraticcurveto{\pgfqpoint{9.279854in}{0.993719in}}{\pgfqpoint{9.284501in}{0.989414in}}%
\pgfpathquadraticcurveto{\pgfqpoint{9.289148in}{0.985109in}}{\pgfqpoint{9.289148in}{0.977400in}}%
\pgfpathquadraticcurveto{\pgfqpoint{9.289148in}{0.968610in}}{\pgfqpoint{9.282195in}{0.963476in}}%
\pgfpathquadraticcurveto{\pgfqpoint{9.275243in}{0.958361in}}{\pgfqpoint{9.263084in}{0.958361in}}%
\pgfpathquadraticcurveto{\pgfqpoint{9.258023in}{0.958361in}}{\pgfqpoint{9.252529in}{0.959352in}}%
\pgfpathquadraticcurveto{\pgfqpoint{9.247035in}{0.960342in}}{\pgfqpoint{9.240965in}{0.962306in}}%
\pgfpathlineto{\pgfqpoint{9.240965in}{0.973005in}}%
\pgfpathquadraticcurveto{\pgfqpoint{9.246711in}{0.970015in}}{\pgfqpoint{9.252277in}{0.968520in}}%
\pgfpathquadraticcurveto{\pgfqpoint{9.257843in}{0.967043in}}{\pgfqpoint{9.263318in}{0.967043in}}%
\pgfpathquadraticcurveto{\pgfqpoint{9.270631in}{0.967043in}}{\pgfqpoint{9.274558in}{0.969546in}}%
\pgfpathquadraticcurveto{\pgfqpoint{9.278503in}{0.972050in}}{\pgfqpoint{9.278503in}{0.976607in}}%
\pgfpathquadraticcurveto{\pgfqpoint{9.278503in}{0.980822in}}{\pgfqpoint{9.275657in}{0.983074in}}%
\pgfpathquadraticcurveto{\pgfqpoint{9.272829in}{0.985325in}}{\pgfqpoint{9.263192in}{0.987414in}}%
\pgfpathlineto{\pgfqpoint{9.259590in}{0.988261in}}%
\pgfpathquadraticcurveto{\pgfqpoint{9.249971in}{0.990278in}}{\pgfqpoint{9.245685in}{0.994475in}}%
\pgfpathquadraticcurveto{\pgfqpoint{9.241416in}{0.998672in}}{\pgfqpoint{9.241416in}{1.005985in}}%
\pgfpathquadraticcurveto{\pgfqpoint{9.241416in}{1.014883in}}{\pgfqpoint{9.247720in}{1.019710in}}%
\pgfpathquadraticcurveto{\pgfqpoint{9.254024in}{1.024556in}}{\pgfqpoint{9.265624in}{1.024556in}}%
\pgfpathquadraticcurveto{\pgfqpoint{9.271352in}{1.024556in}}{\pgfqpoint{9.276413in}{1.023709in}}%
\pgfpathquadraticcurveto{\pgfqpoint{9.281493in}{1.022880in}}{\pgfqpoint{9.285762in}{1.021187in}}%
\pgfpathlineto{\pgfqpoint{9.285762in}{1.021187in}}%
\pgfpathclose%
\pgfusepath{fill}%
\end{pgfscope}%
\begin{pgfscope}%
\pgfsetbuttcap%
\pgfsetroundjoin%
\definecolor{currentfill}{rgb}{0.309804,0.372549,0.419608}%
\pgfsetfillcolor{currentfill}%
\pgfsetlinewidth{0.000000pt}%
\definecolor{currentstroke}{rgb}{0.309804,0.372549,0.419608}%
\pgfsetstrokecolor{currentstroke}%
\pgfsetdash{}{0pt}%
\pgfpathmoveto{\pgfqpoint{9.417543in}{1.010938in}}%
\pgfpathquadraticcurveto{\pgfqpoint{9.421433in}{1.017927in}}{\pgfqpoint{9.426837in}{1.021241in}}%
\pgfpathquadraticcurveto{\pgfqpoint{9.432240in}{1.024556in}}{\pgfqpoint{9.439553in}{1.024556in}}%
\pgfpathquadraticcurveto{\pgfqpoint{9.449406in}{1.024556in}}{\pgfqpoint{9.454756in}{1.017657in}}%
\pgfpathquadraticcurveto{\pgfqpoint{9.460105in}{1.010776in}}{\pgfqpoint{9.460105in}{0.998060in}}%
\pgfpathlineto{\pgfqpoint{9.460105in}{0.960000in}}%
\pgfpathlineto{\pgfqpoint{9.449694in}{0.960000in}}%
\pgfpathlineto{\pgfqpoint{9.449694in}{0.997717in}}%
\pgfpathquadraticcurveto{\pgfqpoint{9.449694in}{1.006778in}}{\pgfqpoint{9.446470in}{1.011155in}}%
\pgfpathquadraticcurveto{\pgfqpoint{9.443264in}{1.015549in}}{\pgfqpoint{9.436689in}{1.015549in}}%
\pgfpathquadraticcurveto{\pgfqpoint{9.428638in}{1.015549in}}{\pgfqpoint{9.423955in}{1.010200in}}%
\pgfpathquadraticcurveto{\pgfqpoint{9.419290in}{1.004868in}}{\pgfqpoint{9.419290in}{0.995628in}}%
\pgfpathlineto{\pgfqpoint{9.419290in}{0.960000in}}%
\pgfpathlineto{\pgfqpoint{9.408879in}{0.960000in}}%
\pgfpathlineto{\pgfqpoint{9.408879in}{0.997717in}}%
\pgfpathquadraticcurveto{\pgfqpoint{9.408879in}{1.006832in}}{\pgfqpoint{9.405672in}{1.011191in}}%
\pgfpathquadraticcurveto{\pgfqpoint{9.402466in}{1.015549in}}{\pgfqpoint{9.395766in}{1.015549in}}%
\pgfpathquadraticcurveto{\pgfqpoint{9.387822in}{1.015549in}}{\pgfqpoint{9.383139in}{1.010182in}}%
\pgfpathquadraticcurveto{\pgfqpoint{9.378474in}{1.004814in}}{\pgfqpoint{9.378474in}{0.995628in}}%
\pgfpathlineto{\pgfqpoint{9.378474in}{0.960000in}}%
\pgfpathlineto{\pgfqpoint{9.368063in}{0.960000in}}%
\pgfpathlineto{\pgfqpoint{9.368063in}{1.023043in}}%
\pgfpathlineto{\pgfqpoint{9.378474in}{1.023043in}}%
\pgfpathlineto{\pgfqpoint{9.378474in}{1.013244in}}%
\pgfpathquadraticcurveto{\pgfqpoint{9.382023in}{1.019044in}}{\pgfqpoint{9.386976in}{1.021800in}}%
\pgfpathquadraticcurveto{\pgfqpoint{9.391929in}{1.024556in}}{\pgfqpoint{9.398738in}{1.024556in}}%
\pgfpathquadraticcurveto{\pgfqpoint{9.405618in}{1.024556in}}{\pgfqpoint{9.410428in}{1.021061in}}%
\pgfpathquadraticcurveto{\pgfqpoint{9.415237in}{1.017585in}}{\pgfqpoint{9.417543in}{1.010938in}}%
\pgfpathlineto{\pgfqpoint{9.417543in}{1.010938in}}%
\pgfpathclose%
\pgfpathmoveto{\pgfqpoint{9.509404in}{0.991683in}}%
\pgfpathquadraticcurveto{\pgfqpoint{9.496850in}{0.991683in}}{\pgfqpoint{9.492005in}{0.988819in}}%
\pgfpathquadraticcurveto{\pgfqpoint{9.487159in}{0.985956in}}{\pgfqpoint{9.487159in}{0.979021in}}%
\pgfpathquadraticcurveto{\pgfqpoint{9.487159in}{0.973509in}}{\pgfqpoint{9.490798in}{0.970267in}}%
\pgfpathquadraticcurveto{\pgfqpoint{9.494436in}{0.967043in}}{\pgfqpoint{9.500669in}{0.967043in}}%
\pgfpathquadraticcurveto{\pgfqpoint{9.509296in}{0.967043in}}{\pgfqpoint{9.514502in}{0.973149in}}%
\pgfpathquadraticcurveto{\pgfqpoint{9.519707in}{0.979255in}}{\pgfqpoint{9.519707in}{0.989378in}}%
\pgfpathlineto{\pgfqpoint{9.519707in}{0.991683in}}%
\pgfpathlineto{\pgfqpoint{9.509404in}{0.991683in}}%
\pgfpathlineto{\pgfqpoint{9.509404in}{0.991683in}}%
\pgfpathclose%
\pgfpathmoveto{\pgfqpoint{9.530064in}{0.995970in}}%
\pgfpathlineto{\pgfqpoint{9.530064in}{0.960000in}}%
\pgfpathlineto{\pgfqpoint{9.519707in}{0.960000in}}%
\pgfpathlineto{\pgfqpoint{9.519707in}{0.969564in}}%
\pgfpathquadraticcurveto{\pgfqpoint{9.516159in}{0.963837in}}{\pgfqpoint{9.510863in}{0.961099in}}%
\pgfpathquadraticcurveto{\pgfqpoint{9.505568in}{0.958361in}}{\pgfqpoint{9.497913in}{0.958361in}}%
\pgfpathquadraticcurveto{\pgfqpoint{9.488240in}{0.958361in}}{\pgfqpoint{9.482512in}{0.963801in}}%
\pgfpathquadraticcurveto{\pgfqpoint{9.476802in}{0.969240in}}{\pgfqpoint{9.476802in}{0.978354in}}%
\pgfpathquadraticcurveto{\pgfqpoint{9.476802in}{0.988982in}}{\pgfqpoint{9.483917in}{0.994385in}}%
\pgfpathquadraticcurveto{\pgfqpoint{9.491050in}{0.999789in}}{\pgfqpoint{9.505172in}{0.999789in}}%
\pgfpathlineto{\pgfqpoint{9.519707in}{0.999789in}}%
\pgfpathlineto{\pgfqpoint{9.519707in}{1.000816in}}%
\pgfpathquadraticcurveto{\pgfqpoint{9.519707in}{1.007966in}}{\pgfqpoint{9.515006in}{1.011875in}}%
\pgfpathquadraticcurveto{\pgfqpoint{9.510305in}{1.015784in}}{\pgfqpoint{9.501803in}{1.015784in}}%
\pgfpathquadraticcurveto{\pgfqpoint{9.496400in}{1.015784in}}{\pgfqpoint{9.491266in}{1.014487in}}%
\pgfpathquadraticcurveto{\pgfqpoint{9.486151in}{1.013190in}}{\pgfqpoint{9.481432in}{1.010596in}}%
\pgfpathlineto{\pgfqpoint{9.481432in}{1.020179in}}%
\pgfpathquadraticcurveto{\pgfqpoint{9.487105in}{1.022376in}}{\pgfqpoint{9.492455in}{1.023457in}}%
\pgfpathquadraticcurveto{\pgfqpoint{9.497805in}{1.024556in}}{\pgfqpoint{9.502866in}{1.024556in}}%
\pgfpathquadraticcurveto{\pgfqpoint{9.516555in}{1.024556in}}{\pgfqpoint{9.523310in}{1.017459in}}%
\pgfpathquadraticcurveto{\pgfqpoint{9.530064in}{1.010380in}}{\pgfqpoint{9.530064in}{0.995970in}}%
\pgfpathlineto{\pgfqpoint{9.530064in}{0.995970in}}%
\pgfpathclose%
\pgfpathmoveto{\pgfqpoint{9.577617in}{0.954146in}}%
\pgfpathquadraticcurveto{\pgfqpoint{9.573240in}{0.942888in}}{\pgfqpoint{9.569061in}{0.939466in}}%
\pgfpathquadraticcurveto{\pgfqpoint{9.564900in}{0.936026in}}{\pgfqpoint{9.557929in}{0.936026in}}%
\pgfpathlineto{\pgfqpoint{9.549644in}{0.936026in}}%
\pgfpathlineto{\pgfqpoint{9.549644in}{0.944690in}}%
\pgfpathlineto{\pgfqpoint{9.555732in}{0.944690in}}%
\pgfpathquadraticcurveto{\pgfqpoint{9.560001in}{0.944690in}}{\pgfqpoint{9.562360in}{0.946725in}}%
\pgfpathquadraticcurveto{\pgfqpoint{9.564738in}{0.948742in}}{\pgfqpoint{9.567602in}{0.956289in}}%
\pgfpathlineto{\pgfqpoint{9.569457in}{0.961009in}}%
\pgfpathlineto{\pgfqpoint{9.543970in}{1.023043in}}%
\pgfpathlineto{\pgfqpoint{9.554939in}{1.023043in}}%
\pgfpathlineto{\pgfqpoint{9.574645in}{0.973743in}}%
\pgfpathlineto{\pgfqpoint{9.594350in}{1.023043in}}%
\pgfpathlineto{\pgfqpoint{9.605319in}{1.023043in}}%
\pgfpathlineto{\pgfqpoint{9.577617in}{0.954146in}}%
\pgfpathlineto{\pgfqpoint{9.577617in}{0.954146in}}%
\pgfpathclose%
\pgfusepath{fill}%
\end{pgfscope}%
\begin{pgfscope}%
\pgfsetbuttcap%
\pgfsetroundjoin%
\definecolor{currentfill}{rgb}{0.309804,0.372549,0.419608}%
\pgfsetfillcolor{currentfill}%
\pgfsetlinewidth{0.000000pt}%
\definecolor{currentstroke}{rgb}{0.309804,0.372549,0.419608}%
\pgfsetstrokecolor{currentstroke}%
\pgfsetdash{}{0pt}%
\pgfpathmoveto{\pgfqpoint{9.719835in}{1.013478in}}%
\pgfpathlineto{\pgfqpoint{9.719835in}{1.047593in}}%
\pgfpathlineto{\pgfqpoint{9.730192in}{1.047593in}}%
\pgfpathlineto{\pgfqpoint{9.730192in}{0.960000in}}%
\pgfpathlineto{\pgfqpoint{9.719835in}{0.960000in}}%
\pgfpathlineto{\pgfqpoint{9.719835in}{0.969456in}}%
\pgfpathquadraticcurveto{\pgfqpoint{9.716574in}{0.963837in}}{\pgfqpoint{9.711585in}{0.961099in}}%
\pgfpathquadraticcurveto{\pgfqpoint{9.706614in}{0.958361in}}{\pgfqpoint{9.699625in}{0.958361in}}%
\pgfpathquadraticcurveto{\pgfqpoint{9.688205in}{0.958361in}}{\pgfqpoint{9.681019in}{0.967475in}}%
\pgfpathquadraticcurveto{\pgfqpoint{9.673850in}{0.976607in}}{\pgfqpoint{9.673850in}{0.991467in}}%
\pgfpathquadraticcurveto{\pgfqpoint{9.673850in}{1.006327in}}{\pgfqpoint{9.681019in}{1.015441in}}%
\pgfpathquadraticcurveto{\pgfqpoint{9.688205in}{1.024556in}}{\pgfqpoint{9.699625in}{1.024556in}}%
\pgfpathquadraticcurveto{\pgfqpoint{9.706614in}{1.024556in}}{\pgfqpoint{9.711585in}{1.021818in}}%
\pgfpathquadraticcurveto{\pgfqpoint{9.716574in}{1.019098in}}{\pgfqpoint{9.719835in}{1.013478in}}%
\pgfpathlineto{\pgfqpoint{9.719835in}{1.013478in}}%
\pgfpathclose%
\pgfpathmoveto{\pgfqpoint{9.684549in}{0.991467in}}%
\pgfpathquadraticcurveto{\pgfqpoint{9.684549in}{0.980048in}}{\pgfqpoint{9.689250in}{0.973545in}}%
\pgfpathquadraticcurveto{\pgfqpoint{9.693951in}{0.967043in}}{\pgfqpoint{9.702165in}{0.967043in}}%
\pgfpathquadraticcurveto{\pgfqpoint{9.710378in}{0.967043in}}{\pgfqpoint{9.715097in}{0.973545in}}%
\pgfpathquadraticcurveto{\pgfqpoint{9.719835in}{0.980048in}}{\pgfqpoint{9.719835in}{0.991467in}}%
\pgfpathquadraticcurveto{\pgfqpoint{9.719835in}{1.002887in}}{\pgfqpoint{9.715097in}{1.009389in}}%
\pgfpathquadraticcurveto{\pgfqpoint{9.710378in}{1.015892in}}{\pgfqpoint{9.702165in}{1.015892in}}%
\pgfpathquadraticcurveto{\pgfqpoint{9.693951in}{1.015892in}}{\pgfqpoint{9.689250in}{1.009389in}}%
\pgfpathquadraticcurveto{\pgfqpoint{9.684549in}{1.002887in}}{\pgfqpoint{9.684549in}{0.991467in}}%
\pgfpathlineto{\pgfqpoint{9.684549in}{0.991467in}}%
\pgfpathclose%
\pgfpathmoveto{\pgfqpoint{9.751536in}{1.023043in}}%
\pgfpathlineto{\pgfqpoint{9.761893in}{1.023043in}}%
\pgfpathlineto{\pgfqpoint{9.761893in}{0.960000in}}%
\pgfpathlineto{\pgfqpoint{9.751536in}{0.960000in}}%
\pgfpathlineto{\pgfqpoint{9.751536in}{1.023043in}}%
\pgfpathlineto{\pgfqpoint{9.751536in}{1.023043in}}%
\pgfpathclose%
\pgfpathmoveto{\pgfqpoint{9.751536in}{1.047593in}}%
\pgfpathlineto{\pgfqpoint{9.761893in}{1.047593in}}%
\pgfpathlineto{\pgfqpoint{9.761893in}{1.034462in}}%
\pgfpathlineto{\pgfqpoint{9.751536in}{1.034462in}}%
\pgfpathlineto{\pgfqpoint{9.751536in}{1.047593in}}%
\pgfpathlineto{\pgfqpoint{9.751536in}{1.047593in}}%
\pgfpathclose%
\pgfpathmoveto{\pgfqpoint{9.854313in}{1.047593in}}%
\pgfpathlineto{\pgfqpoint{9.854313in}{1.038965in}}%
\pgfpathlineto{\pgfqpoint{9.844407in}{1.038965in}}%
\pgfpathquadraticcurveto{\pgfqpoint{9.838841in}{1.038965in}}{\pgfqpoint{9.836662in}{1.036714in}}%
\pgfpathquadraticcurveto{\pgfqpoint{9.834500in}{1.034462in}}{\pgfqpoint{9.834500in}{1.028608in}}%
\pgfpathlineto{\pgfqpoint{9.834500in}{1.023043in}}%
\pgfpathlineto{\pgfqpoint{9.851558in}{1.023043in}}%
\pgfpathlineto{\pgfqpoint{9.851558in}{1.014991in}}%
\pgfpathlineto{\pgfqpoint{9.834500in}{1.014991in}}%
\pgfpathlineto{\pgfqpoint{9.834500in}{0.960000in}}%
\pgfpathlineto{\pgfqpoint{9.824089in}{0.960000in}}%
\pgfpathlineto{\pgfqpoint{9.824089in}{1.014991in}}%
\pgfpathlineto{\pgfqpoint{9.795666in}{1.014991in}}%
\pgfpathlineto{\pgfqpoint{9.795666in}{0.960000in}}%
\pgfpathlineto{\pgfqpoint{9.785255in}{0.960000in}}%
\pgfpathlineto{\pgfqpoint{9.785255in}{1.014991in}}%
\pgfpathlineto{\pgfqpoint{9.775348in}{1.014991in}}%
\pgfpathlineto{\pgfqpoint{9.775348in}{1.023043in}}%
\pgfpathlineto{\pgfqpoint{9.785255in}{1.023043in}}%
\pgfpathlineto{\pgfqpoint{9.785255in}{1.027438in}}%
\pgfpathquadraticcurveto{\pgfqpoint{9.785255in}{1.037957in}}{\pgfqpoint{9.790154in}{1.042766in}}%
\pgfpathquadraticcurveto{\pgfqpoint{9.795053in}{1.047593in}}{\pgfqpoint{9.805681in}{1.047593in}}%
\pgfpathlineto{\pgfqpoint{9.815479in}{1.047593in}}%
\pgfpathlineto{\pgfqpoint{9.815479in}{1.038965in}}%
\pgfpathlineto{\pgfqpoint{9.805573in}{1.038965in}}%
\pgfpathquadraticcurveto{\pgfqpoint{9.800007in}{1.038965in}}{\pgfqpoint{9.797809in}{1.036714in}}%
\pgfpathquadraticcurveto{\pgfqpoint{9.795666in}{1.034462in}}{\pgfqpoint{9.795666in}{1.028608in}}%
\pgfpathlineto{\pgfqpoint{9.795666in}{1.023043in}}%
\pgfpathlineto{\pgfqpoint{9.824089in}{1.023043in}}%
\pgfpathlineto{\pgfqpoint{9.824089in}{1.027438in}}%
\pgfpathquadraticcurveto{\pgfqpoint{9.824089in}{1.037957in}}{\pgfqpoint{9.828988in}{1.042766in}}%
\pgfpathquadraticcurveto{\pgfqpoint{9.833888in}{1.047593in}}{\pgfqpoint{9.844533in}{1.047593in}}%
\pgfpathlineto{\pgfqpoint{9.854313in}{1.047593in}}%
\pgfpathlineto{\pgfqpoint{9.854313in}{1.047593in}}%
\pgfpathclose%
\pgfpathmoveto{\pgfqpoint{9.916906in}{0.994115in}}%
\pgfpathlineto{\pgfqpoint{9.916906in}{0.989054in}}%
\pgfpathlineto{\pgfqpoint{9.869282in}{0.989054in}}%
\pgfpathquadraticcurveto{\pgfqpoint{9.869966in}{0.978354in}}{\pgfqpoint{9.875730in}{0.972753in}}%
\pgfpathquadraticcurveto{\pgfqpoint{9.881494in}{0.967151in}}{\pgfqpoint{9.891797in}{0.967151in}}%
\pgfpathquadraticcurveto{\pgfqpoint{9.897759in}{0.967151in}}{\pgfqpoint{9.903361in}{0.968610in}}%
\pgfpathquadraticcurveto{\pgfqpoint{9.908962in}{0.970069in}}{\pgfqpoint{9.914492in}{0.973005in}}%
\pgfpathlineto{\pgfqpoint{9.914492in}{0.963206in}}%
\pgfpathquadraticcurveto{\pgfqpoint{9.908908in}{0.960847in}}{\pgfqpoint{9.903054in}{0.959604in}}%
\pgfpathquadraticcurveto{\pgfqpoint{9.897200in}{0.958361in}}{\pgfqpoint{9.891184in}{0.958361in}}%
\pgfpathquadraticcurveto{\pgfqpoint{9.876090in}{0.958361in}}{\pgfqpoint{9.867282in}{0.967133in}}%
\pgfpathquadraticcurveto{\pgfqpoint{9.858474in}{0.975923in}}{\pgfqpoint{9.858474in}{0.990909in}}%
\pgfpathquadraticcurveto{\pgfqpoint{9.858474in}{1.006381in}}{\pgfqpoint{9.866832in}{1.015459in}}%
\pgfpathquadraticcurveto{\pgfqpoint{9.875190in}{1.024556in}}{\pgfqpoint{9.889383in}{1.024556in}}%
\pgfpathquadraticcurveto{\pgfqpoint{9.902100in}{1.024556in}}{\pgfqpoint{9.909503in}{1.016360in}}%
\pgfpathquadraticcurveto{\pgfqpoint{9.916906in}{1.008183in}}{\pgfqpoint{9.916906in}{0.994115in}}%
\pgfpathlineto{\pgfqpoint{9.916906in}{0.994115in}}%
\pgfpathclose%
\pgfpathmoveto{\pgfqpoint{9.906549in}{0.997159in}}%
\pgfpathquadraticcurveto{\pgfqpoint{9.906441in}{1.005643in}}{\pgfqpoint{9.901793in}{1.010704in}}%
\pgfpathquadraticcurveto{\pgfqpoint{9.897146in}{1.015784in}}{\pgfqpoint{9.889491in}{1.015784in}}%
\pgfpathquadraticcurveto{\pgfqpoint{9.880827in}{1.015784in}}{\pgfqpoint{9.875622in}{1.010884in}}%
\pgfpathquadraticcurveto{\pgfqpoint{9.870416in}{1.005985in}}{\pgfqpoint{9.869624in}{0.997087in}}%
\pgfpathlineto{\pgfqpoint{9.906549in}{0.997159in}}%
\pgfpathlineto{\pgfqpoint{9.906549in}{0.997159in}}%
\pgfpathclose%
\pgfpathmoveto{\pgfqpoint{9.970438in}{1.013370in}}%
\pgfpathquadraticcurveto{\pgfqpoint{9.968691in}{1.014379in}}{\pgfqpoint{9.966637in}{1.014847in}}%
\pgfpathquadraticcurveto{\pgfqpoint{9.964584in}{1.015333in}}{\pgfqpoint{9.962116in}{1.015333in}}%
\pgfpathquadraticcurveto{\pgfqpoint{9.953326in}{1.015333in}}{\pgfqpoint{9.948625in}{1.009623in}}%
\pgfpathquadraticcurveto{\pgfqpoint{9.943924in}{1.003914in}}{\pgfqpoint{9.943924in}{0.993214in}}%
\pgfpathlineto{\pgfqpoint{9.943924in}{0.960000in}}%
\pgfpathlineto{\pgfqpoint{9.933513in}{0.960000in}}%
\pgfpathlineto{\pgfqpoint{9.933513in}{1.023043in}}%
\pgfpathlineto{\pgfqpoint{9.943924in}{1.023043in}}%
\pgfpathlineto{\pgfqpoint{9.943924in}{1.013244in}}%
\pgfpathquadraticcurveto{\pgfqpoint{9.947202in}{1.018990in}}{\pgfqpoint{9.952426in}{1.021764in}}%
\pgfpathquadraticcurveto{\pgfqpoint{9.957667in}{1.024556in}}{\pgfqpoint{9.965160in}{1.024556in}}%
\pgfpathquadraticcurveto{\pgfqpoint{9.966223in}{1.024556in}}{\pgfqpoint{9.967520in}{1.024411in}}%
\pgfpathquadraticcurveto{\pgfqpoint{9.968817in}{1.024285in}}{\pgfqpoint{9.970384in}{1.023997in}}%
\pgfpathlineto{\pgfqpoint{9.970438in}{1.013370in}}%
\pgfpathlineto{\pgfqpoint{9.970438in}{1.013370in}}%
\pgfpathclose%
\pgfusepath{fill}%
\end{pgfscope}%
\begin{pgfscope}%
\pgfsetbuttcap%
\pgfsetroundjoin%
\definecolor{currentfill}{rgb}{0.309804,0.372549,0.419608}%
\pgfsetfillcolor{currentfill}%
\pgfsetlinewidth{0.000000pt}%
\definecolor{currentstroke}{rgb}{0.309804,0.372549,0.419608}%
\pgfsetstrokecolor{currentstroke}%
\pgfsetdash{}{0pt}%
\pgfpathmoveto{\pgfqpoint{10.070164in}{1.047593in}}%
\pgfpathlineto{\pgfqpoint{10.070164in}{1.038965in}}%
\pgfpathlineto{\pgfqpoint{10.060257in}{1.038965in}}%
\pgfpathquadraticcurveto{\pgfqpoint{10.054691in}{1.038965in}}{\pgfqpoint{10.052512in}{1.036714in}}%
\pgfpathquadraticcurveto{\pgfqpoint{10.050350in}{1.034462in}}{\pgfqpoint{10.050350in}{1.028608in}}%
\pgfpathlineto{\pgfqpoint{10.050350in}{1.023043in}}%
\pgfpathlineto{\pgfqpoint{10.067408in}{1.023043in}}%
\pgfpathlineto{\pgfqpoint{10.067408in}{1.014991in}}%
\pgfpathlineto{\pgfqpoint{10.050350in}{1.014991in}}%
\pgfpathlineto{\pgfqpoint{10.050350in}{0.960000in}}%
\pgfpathlineto{\pgfqpoint{10.039939in}{0.960000in}}%
\pgfpathlineto{\pgfqpoint{10.039939in}{1.014991in}}%
\pgfpathlineto{\pgfqpoint{10.030032in}{1.014991in}}%
\pgfpathlineto{\pgfqpoint{10.030032in}{1.023043in}}%
\pgfpathlineto{\pgfqpoint{10.039939in}{1.023043in}}%
\pgfpathlineto{\pgfqpoint{10.039939in}{1.027438in}}%
\pgfpathquadraticcurveto{\pgfqpoint{10.039939in}{1.037957in}}{\pgfqpoint{10.044838in}{1.042766in}}%
\pgfpathquadraticcurveto{\pgfqpoint{10.049738in}{1.047593in}}{\pgfqpoint{10.060365in}{1.047593in}}%
\pgfpathlineto{\pgfqpoint{10.070164in}{1.047593in}}%
\pgfpathlineto{\pgfqpoint{10.070164in}{1.047593in}}%
\pgfpathclose%
\pgfpathmoveto{\pgfqpoint{10.115356in}{1.013370in}}%
\pgfpathquadraticcurveto{\pgfqpoint{10.113609in}{1.014379in}}{\pgfqpoint{10.111555in}{1.014847in}}%
\pgfpathquadraticcurveto{\pgfqpoint{10.109502in}{1.015333in}}{\pgfqpoint{10.107034in}{1.015333in}}%
\pgfpathquadraticcurveto{\pgfqpoint{10.098245in}{1.015333in}}{\pgfqpoint{10.093543in}{1.009623in}}%
\pgfpathquadraticcurveto{\pgfqpoint{10.088842in}{1.003914in}}{\pgfqpoint{10.088842in}{0.993214in}}%
\pgfpathlineto{\pgfqpoint{10.088842in}{0.960000in}}%
\pgfpathlineto{\pgfqpoint{10.078431in}{0.960000in}}%
\pgfpathlineto{\pgfqpoint{10.078431in}{1.023043in}}%
\pgfpathlineto{\pgfqpoint{10.088842in}{1.023043in}}%
\pgfpathlineto{\pgfqpoint{10.088842in}{1.013244in}}%
\pgfpathquadraticcurveto{\pgfqpoint{10.092120in}{1.018990in}}{\pgfqpoint{10.097344in}{1.021764in}}%
\pgfpathquadraticcurveto{\pgfqpoint{10.102585in}{1.024556in}}{\pgfqpoint{10.110078in}{1.024556in}}%
\pgfpathquadraticcurveto{\pgfqpoint{10.111141in}{1.024556in}}{\pgfqpoint{10.112438in}{1.024411in}}%
\pgfpathquadraticcurveto{\pgfqpoint{10.113735in}{1.024285in}}{\pgfqpoint{10.115302in}{1.023997in}}%
\pgfpathlineto{\pgfqpoint{10.115356in}{1.013370in}}%
\pgfpathlineto{\pgfqpoint{10.115356in}{1.013370in}}%
\pgfpathclose%
\pgfpathmoveto{\pgfqpoint{10.148102in}{1.015784in}}%
\pgfpathquadraticcurveto{\pgfqpoint{10.139781in}{1.015784in}}{\pgfqpoint{10.134935in}{1.009281in}}%
\pgfpathquadraticcurveto{\pgfqpoint{10.130090in}{1.002779in}}{\pgfqpoint{10.130090in}{0.991467in}}%
\pgfpathquadraticcurveto{\pgfqpoint{10.130090in}{0.980156in}}{\pgfqpoint{10.134899in}{0.973653in}}%
\pgfpathquadraticcurveto{\pgfqpoint{10.139726in}{0.967151in}}{\pgfqpoint{10.148102in}{0.967151in}}%
\pgfpathquadraticcurveto{\pgfqpoint{10.156388in}{0.967151in}}{\pgfqpoint{10.161215in}{0.973671in}}%
\pgfpathquadraticcurveto{\pgfqpoint{10.166060in}{0.980210in}}{\pgfqpoint{10.166060in}{0.991467in}}%
\pgfpathquadraticcurveto{\pgfqpoint{10.166060in}{1.002671in}}{\pgfqpoint{10.161215in}{1.009227in}}%
\pgfpathquadraticcurveto{\pgfqpoint{10.156388in}{1.015784in}}{\pgfqpoint{10.148102in}{1.015784in}}%
\pgfpathlineto{\pgfqpoint{10.148102in}{1.015784in}}%
\pgfpathclose%
\pgfpathmoveto{\pgfqpoint{10.148102in}{1.024556in}}%
\pgfpathquadraticcurveto{\pgfqpoint{10.161611in}{1.024556in}}{\pgfqpoint{10.169320in}{1.015766in}}%
\pgfpathquadraticcurveto{\pgfqpoint{10.177048in}{1.006994in}}{\pgfqpoint{10.177048in}{0.991467in}}%
\pgfpathquadraticcurveto{\pgfqpoint{10.177048in}{0.975995in}}{\pgfqpoint{10.169320in}{0.967169in}}%
\pgfpathquadraticcurveto{\pgfqpoint{10.161611in}{0.958361in}}{\pgfqpoint{10.148102in}{0.958361in}}%
\pgfpathquadraticcurveto{\pgfqpoint{10.134539in}{0.958361in}}{\pgfqpoint{10.126848in}{0.967169in}}%
\pgfpathquadraticcurveto{\pgfqpoint{10.119175in}{0.975995in}}{\pgfqpoint{10.119175in}{0.991467in}}%
\pgfpathquadraticcurveto{\pgfqpoint{10.119175in}{1.006994in}}{\pgfqpoint{10.126848in}{1.015766in}}%
\pgfpathquadraticcurveto{\pgfqpoint{10.134539in}{1.024556in}}{\pgfqpoint{10.148102in}{1.024556in}}%
\pgfpathlineto{\pgfqpoint{10.148102in}{1.024556in}}%
\pgfpathclose%
\pgfpathmoveto{\pgfqpoint{10.243296in}{1.010938in}}%
\pgfpathquadraticcurveto{\pgfqpoint{10.247187in}{1.017927in}}{\pgfqpoint{10.252591in}{1.021241in}}%
\pgfpathquadraticcurveto{\pgfqpoint{10.257994in}{1.024556in}}{\pgfqpoint{10.265307in}{1.024556in}}%
\pgfpathquadraticcurveto{\pgfqpoint{10.275160in}{1.024556in}}{\pgfqpoint{10.280509in}{1.017657in}}%
\pgfpathquadraticcurveto{\pgfqpoint{10.285859in}{1.010776in}}{\pgfqpoint{10.285859in}{0.998060in}}%
\pgfpathlineto{\pgfqpoint{10.285859in}{0.960000in}}%
\pgfpathlineto{\pgfqpoint{10.275448in}{0.960000in}}%
\pgfpathlineto{\pgfqpoint{10.275448in}{0.997717in}}%
\pgfpathquadraticcurveto{\pgfqpoint{10.275448in}{1.006778in}}{\pgfqpoint{10.272224in}{1.011155in}}%
\pgfpathquadraticcurveto{\pgfqpoint{10.269018in}{1.015549in}}{\pgfqpoint{10.262443in}{1.015549in}}%
\pgfpathquadraticcurveto{\pgfqpoint{10.254392in}{1.015549in}}{\pgfqpoint{10.249709in}{1.010200in}}%
\pgfpathquadraticcurveto{\pgfqpoint{10.245044in}{1.004868in}}{\pgfqpoint{10.245044in}{0.995628in}}%
\pgfpathlineto{\pgfqpoint{10.245044in}{0.960000in}}%
\pgfpathlineto{\pgfqpoint{10.234633in}{0.960000in}}%
\pgfpathlineto{\pgfqpoint{10.234633in}{0.997717in}}%
\pgfpathquadraticcurveto{\pgfqpoint{10.234633in}{1.006832in}}{\pgfqpoint{10.231426in}{1.011191in}}%
\pgfpathquadraticcurveto{\pgfqpoint{10.228220in}{1.015549in}}{\pgfqpoint{10.221520in}{1.015549in}}%
\pgfpathquadraticcurveto{\pgfqpoint{10.213576in}{1.015549in}}{\pgfqpoint{10.208893in}{1.010182in}}%
\pgfpathquadraticcurveto{\pgfqpoint{10.204228in}{1.004814in}}{\pgfqpoint{10.204228in}{0.995628in}}%
\pgfpathlineto{\pgfqpoint{10.204228in}{0.960000in}}%
\pgfpathlineto{\pgfqpoint{10.193817in}{0.960000in}}%
\pgfpathlineto{\pgfqpoint{10.193817in}{1.023043in}}%
\pgfpathlineto{\pgfqpoint{10.204228in}{1.023043in}}%
\pgfpathlineto{\pgfqpoint{10.204228in}{1.013244in}}%
\pgfpathquadraticcurveto{\pgfqpoint{10.207776in}{1.019044in}}{\pgfqpoint{10.212730in}{1.021800in}}%
\pgfpathquadraticcurveto{\pgfqpoint{10.217683in}{1.024556in}}{\pgfqpoint{10.224492in}{1.024556in}}%
\pgfpathquadraticcurveto{\pgfqpoint{10.231372in}{1.024556in}}{\pgfqpoint{10.236182in}{1.021061in}}%
\pgfpathquadraticcurveto{\pgfqpoint{10.240991in}{1.017585in}}{\pgfqpoint{10.243296in}{1.010938in}}%
\pgfpathlineto{\pgfqpoint{10.243296in}{1.010938in}}%
\pgfpathclose%
\pgfusepath{fill}%
\end{pgfscope}%
\begin{pgfscope}%
\pgfsetbuttcap%
\pgfsetroundjoin%
\definecolor{currentfill}{rgb}{0.309804,0.372549,0.419608}%
\pgfsetfillcolor{currentfill}%
\pgfsetlinewidth{0.000000pt}%
\definecolor{currentstroke}{rgb}{0.309804,0.372549,0.419608}%
\pgfsetstrokecolor{currentstroke}%
\pgfsetdash{}{0pt}%
\pgfpathmoveto{\pgfqpoint{10.394297in}{0.991683in}}%
\pgfpathquadraticcurveto{\pgfqpoint{10.381742in}{0.991683in}}{\pgfqpoint{10.376897in}{0.988819in}}%
\pgfpathquadraticcurveto{\pgfqpoint{10.372052in}{0.985956in}}{\pgfqpoint{10.372052in}{0.979021in}}%
\pgfpathquadraticcurveto{\pgfqpoint{10.372052in}{0.973509in}}{\pgfqpoint{10.375690in}{0.970267in}}%
\pgfpathquadraticcurveto{\pgfqpoint{10.379329in}{0.967043in}}{\pgfqpoint{10.385561in}{0.967043in}}%
\pgfpathquadraticcurveto{\pgfqpoint{10.394189in}{0.967043in}}{\pgfqpoint{10.399394in}{0.973149in}}%
\pgfpathquadraticcurveto{\pgfqpoint{10.404600in}{0.979255in}}{\pgfqpoint{10.404600in}{0.989378in}}%
\pgfpathlineto{\pgfqpoint{10.404600in}{0.991683in}}%
\pgfpathlineto{\pgfqpoint{10.394297in}{0.991683in}}%
\pgfpathlineto{\pgfqpoint{10.394297in}{0.991683in}}%
\pgfpathclose%
\pgfpathmoveto{\pgfqpoint{10.414957in}{0.995970in}}%
\pgfpathlineto{\pgfqpoint{10.414957in}{0.960000in}}%
\pgfpathlineto{\pgfqpoint{10.404600in}{0.960000in}}%
\pgfpathlineto{\pgfqpoint{10.404600in}{0.969564in}}%
\pgfpathquadraticcurveto{\pgfqpoint{10.401051in}{0.963837in}}{\pgfqpoint{10.395756in}{0.961099in}}%
\pgfpathquadraticcurveto{\pgfqpoint{10.390460in}{0.958361in}}{\pgfqpoint{10.382805in}{0.958361in}}%
\pgfpathquadraticcurveto{\pgfqpoint{10.373132in}{0.958361in}}{\pgfqpoint{10.367405in}{0.963801in}}%
\pgfpathquadraticcurveto{\pgfqpoint{10.361695in}{0.969240in}}{\pgfqpoint{10.361695in}{0.978354in}}%
\pgfpathquadraticcurveto{\pgfqpoint{10.361695in}{0.988982in}}{\pgfqpoint{10.368809in}{0.994385in}}%
\pgfpathquadraticcurveto{\pgfqpoint{10.375942in}{0.999789in}}{\pgfqpoint{10.390064in}{0.999789in}}%
\pgfpathlineto{\pgfqpoint{10.404600in}{0.999789in}}%
\pgfpathlineto{\pgfqpoint{10.404600in}{1.000816in}}%
\pgfpathquadraticcurveto{\pgfqpoint{10.404600in}{1.007966in}}{\pgfqpoint{10.399898in}{1.011875in}}%
\pgfpathquadraticcurveto{\pgfqpoint{10.395197in}{1.015784in}}{\pgfqpoint{10.386696in}{1.015784in}}%
\pgfpathquadraticcurveto{\pgfqpoint{10.381292in}{1.015784in}}{\pgfqpoint{10.376158in}{1.014487in}}%
\pgfpathquadraticcurveto{\pgfqpoint{10.371043in}{1.013190in}}{\pgfqpoint{10.366324in}{1.010596in}}%
\pgfpathlineto{\pgfqpoint{10.366324in}{1.020179in}}%
\pgfpathquadraticcurveto{\pgfqpoint{10.371998in}{1.022376in}}{\pgfqpoint{10.377347in}{1.023457in}}%
\pgfpathquadraticcurveto{\pgfqpoint{10.382697in}{1.024556in}}{\pgfqpoint{10.387758in}{1.024556in}}%
\pgfpathquadraticcurveto{\pgfqpoint{10.401448in}{1.024556in}}{\pgfqpoint{10.408202in}{1.017459in}}%
\pgfpathquadraticcurveto{\pgfqpoint{10.414957in}{1.010380in}}{\pgfqpoint{10.414957in}{0.995970in}}%
\pgfpathlineto{\pgfqpoint{10.414957in}{0.995970in}}%
\pgfpathclose%
\pgfpathmoveto{\pgfqpoint{10.481656in}{1.020629in}}%
\pgfpathlineto{\pgfqpoint{10.481656in}{1.010938in}}%
\pgfpathquadraticcurveto{\pgfqpoint{10.477261in}{1.013370in}}{\pgfqpoint{10.472848in}{1.014577in}}%
\pgfpathquadraticcurveto{\pgfqpoint{10.468435in}{1.015784in}}{\pgfqpoint{10.463932in}{1.015784in}}%
\pgfpathquadraticcurveto{\pgfqpoint{10.453845in}{1.015784in}}{\pgfqpoint{10.448261in}{1.009389in}}%
\pgfpathquadraticcurveto{\pgfqpoint{10.442695in}{1.003013in}}{\pgfqpoint{10.442695in}{0.991467in}}%
\pgfpathquadraticcurveto{\pgfqpoint{10.442695in}{0.979921in}}{\pgfqpoint{10.448261in}{0.973527in}}%
\pgfpathquadraticcurveto{\pgfqpoint{10.453845in}{0.967151in}}{\pgfqpoint{10.463932in}{0.967151in}}%
\pgfpathquadraticcurveto{\pgfqpoint{10.468435in}{0.967151in}}{\pgfqpoint{10.472848in}{0.968358in}}%
\pgfpathquadraticcurveto{\pgfqpoint{10.477261in}{0.969564in}}{\pgfqpoint{10.481656in}{0.971996in}}%
\pgfpathlineto{\pgfqpoint{10.481656in}{0.962414in}}%
\pgfpathquadraticcurveto{\pgfqpoint{10.477315in}{0.960396in}}{\pgfqpoint{10.472668in}{0.959388in}}%
\pgfpathquadraticcurveto{\pgfqpoint{10.468038in}{0.958361in}}{\pgfqpoint{10.462797in}{0.958361in}}%
\pgfpathquadraticcurveto{\pgfqpoint{10.448549in}{0.958361in}}{\pgfqpoint{10.440156in}{0.967313in}}%
\pgfpathquadraticcurveto{\pgfqpoint{10.431780in}{0.976265in}}{\pgfqpoint{10.431780in}{0.991467in}}%
\pgfpathquadraticcurveto{\pgfqpoint{10.431780in}{1.006886in}}{\pgfqpoint{10.440246in}{1.015712in}}%
\pgfpathquadraticcurveto{\pgfqpoint{10.448729in}{1.024556in}}{\pgfqpoint{10.463481in}{1.024556in}}%
\pgfpathquadraticcurveto{\pgfqpoint{10.468255in}{1.024556in}}{\pgfqpoint{10.472812in}{1.023565in}}%
\pgfpathquadraticcurveto{\pgfqpoint{10.477369in}{1.022592in}}{\pgfqpoint{10.481656in}{1.020629in}}%
\pgfpathlineto{\pgfqpoint{10.481656in}{1.020629in}}%
\pgfpathclose%
\pgfpathmoveto{\pgfqpoint{10.509917in}{1.040947in}}%
\pgfpathlineto{\pgfqpoint{10.509917in}{1.023043in}}%
\pgfpathlineto{\pgfqpoint{10.531243in}{1.023043in}}%
\pgfpathlineto{\pgfqpoint{10.531243in}{1.014991in}}%
\pgfpathlineto{\pgfqpoint{10.509917in}{1.014991in}}%
\pgfpathlineto{\pgfqpoint{10.509917in}{0.980768in}}%
\pgfpathquadraticcurveto{\pgfqpoint{10.509917in}{0.973059in}}{\pgfqpoint{10.512024in}{0.970861in}}%
\pgfpathquadraticcurveto{\pgfqpoint{10.514132in}{0.968664in}}{\pgfqpoint{10.520616in}{0.968664in}}%
\pgfpathlineto{\pgfqpoint{10.531243in}{0.968664in}}%
\pgfpathlineto{\pgfqpoint{10.531243in}{0.960000in}}%
\pgfpathlineto{\pgfqpoint{10.520616in}{0.960000in}}%
\pgfpathquadraticcurveto{\pgfqpoint{10.508620in}{0.960000in}}{\pgfqpoint{10.504063in}{0.964467in}}%
\pgfpathquadraticcurveto{\pgfqpoint{10.499506in}{0.968952in}}{\pgfqpoint{10.499506in}{0.980768in}}%
\pgfpathlineto{\pgfqpoint{10.499506in}{1.014991in}}%
\pgfpathlineto{\pgfqpoint{10.491905in}{1.014991in}}%
\pgfpathlineto{\pgfqpoint{10.491905in}{1.023043in}}%
\pgfpathlineto{\pgfqpoint{10.499506in}{1.023043in}}%
\pgfpathlineto{\pgfqpoint{10.499506in}{1.040947in}}%
\pgfpathlineto{\pgfqpoint{10.509917in}{1.040947in}}%
\pgfpathlineto{\pgfqpoint{10.509917in}{1.040947in}}%
\pgfpathclose%
\pgfpathmoveto{\pgfqpoint{10.543798in}{0.984875in}}%
\pgfpathlineto{\pgfqpoint{10.543798in}{1.023043in}}%
\pgfpathlineto{\pgfqpoint{10.554155in}{1.023043in}}%
\pgfpathlineto{\pgfqpoint{10.554155in}{0.985271in}}%
\pgfpathquadraticcurveto{\pgfqpoint{10.554155in}{0.976319in}}{\pgfqpoint{10.557631in}{0.971834in}}%
\pgfpathquadraticcurveto{\pgfqpoint{10.561125in}{0.967367in}}{\pgfqpoint{10.568114in}{0.967367in}}%
\pgfpathquadraticcurveto{\pgfqpoint{10.576490in}{0.967367in}}{\pgfqpoint{10.581353in}{0.972717in}}%
\pgfpathquadraticcurveto{\pgfqpoint{10.586234in}{0.978066in}}{\pgfqpoint{10.586234in}{0.987306in}}%
\pgfpathlineto{\pgfqpoint{10.586234in}{1.023043in}}%
\pgfpathlineto{\pgfqpoint{10.596591in}{1.023043in}}%
\pgfpathlineto{\pgfqpoint{10.596591in}{0.960000in}}%
\pgfpathlineto{\pgfqpoint{10.586234in}{0.960000in}}%
\pgfpathlineto{\pgfqpoint{10.586234in}{0.969691in}}%
\pgfpathquadraticcurveto{\pgfqpoint{10.582470in}{0.963945in}}{\pgfqpoint{10.577480in}{0.961153in}}%
\pgfpathquadraticcurveto{\pgfqpoint{10.572509in}{0.958361in}}{\pgfqpoint{10.565916in}{0.958361in}}%
\pgfpathquadraticcurveto{\pgfqpoint{10.555055in}{0.958361in}}{\pgfqpoint{10.549417in}{0.965115in}}%
\pgfpathquadraticcurveto{\pgfqpoint{10.543798in}{0.971870in}}{\pgfqpoint{10.543798in}{0.984875in}}%
\pgfpathlineto{\pgfqpoint{10.543798in}{0.984875in}}%
\pgfpathclose%
\pgfpathmoveto{\pgfqpoint{10.569861in}{1.024556in}}%
\pgfpathlineto{\pgfqpoint{10.569861in}{1.024556in}}%
\pgfpathlineto{\pgfqpoint{10.569861in}{1.024556in}}%
\pgfpathclose%
\pgfpathmoveto{\pgfqpoint{10.646575in}{0.991683in}}%
\pgfpathquadraticcurveto{\pgfqpoint{10.634020in}{0.991683in}}{\pgfqpoint{10.629175in}{0.988819in}}%
\pgfpathquadraticcurveto{\pgfqpoint{10.624330in}{0.985956in}}{\pgfqpoint{10.624330in}{0.979021in}}%
\pgfpathquadraticcurveto{\pgfqpoint{10.624330in}{0.973509in}}{\pgfqpoint{10.627968in}{0.970267in}}%
\pgfpathquadraticcurveto{\pgfqpoint{10.631607in}{0.967043in}}{\pgfqpoint{10.637839in}{0.967043in}}%
\pgfpathquadraticcurveto{\pgfqpoint{10.646467in}{0.967043in}}{\pgfqpoint{10.651672in}{0.973149in}}%
\pgfpathquadraticcurveto{\pgfqpoint{10.656878in}{0.979255in}}{\pgfqpoint{10.656878in}{0.989378in}}%
\pgfpathlineto{\pgfqpoint{10.656878in}{0.991683in}}%
\pgfpathlineto{\pgfqpoint{10.646575in}{0.991683in}}%
\pgfpathlineto{\pgfqpoint{10.646575in}{0.991683in}}%
\pgfpathclose%
\pgfpathmoveto{\pgfqpoint{10.667235in}{0.995970in}}%
\pgfpathlineto{\pgfqpoint{10.667235in}{0.960000in}}%
\pgfpathlineto{\pgfqpoint{10.656878in}{0.960000in}}%
\pgfpathlineto{\pgfqpoint{10.656878in}{0.969564in}}%
\pgfpathquadraticcurveto{\pgfqpoint{10.653329in}{0.963837in}}{\pgfqpoint{10.648034in}{0.961099in}}%
\pgfpathquadraticcurveto{\pgfqpoint{10.642738in}{0.958361in}}{\pgfqpoint{10.635083in}{0.958361in}}%
\pgfpathquadraticcurveto{\pgfqpoint{10.625411in}{0.958361in}}{\pgfqpoint{10.619683in}{0.963801in}}%
\pgfpathquadraticcurveto{\pgfqpoint{10.613973in}{0.969240in}}{\pgfqpoint{10.613973in}{0.978354in}}%
\pgfpathquadraticcurveto{\pgfqpoint{10.613973in}{0.988982in}}{\pgfqpoint{10.621088in}{0.994385in}}%
\pgfpathquadraticcurveto{\pgfqpoint{10.628221in}{0.999789in}}{\pgfqpoint{10.642342in}{0.999789in}}%
\pgfpathlineto{\pgfqpoint{10.656878in}{0.999789in}}%
\pgfpathlineto{\pgfqpoint{10.656878in}{1.000816in}}%
\pgfpathquadraticcurveto{\pgfqpoint{10.656878in}{1.007966in}}{\pgfqpoint{10.652177in}{1.011875in}}%
\pgfpathquadraticcurveto{\pgfqpoint{10.647475in}{1.015784in}}{\pgfqpoint{10.638974in}{1.015784in}}%
\pgfpathquadraticcurveto{\pgfqpoint{10.633570in}{1.015784in}}{\pgfqpoint{10.628437in}{1.014487in}}%
\pgfpathquadraticcurveto{\pgfqpoint{10.623321in}{1.013190in}}{\pgfqpoint{10.618602in}{1.010596in}}%
\pgfpathlineto{\pgfqpoint{10.618602in}{1.020179in}}%
\pgfpathquadraticcurveto{\pgfqpoint{10.624276in}{1.022376in}}{\pgfqpoint{10.629625in}{1.023457in}}%
\pgfpathquadraticcurveto{\pgfqpoint{10.634975in}{1.024556in}}{\pgfqpoint{10.640036in}{1.024556in}}%
\pgfpathquadraticcurveto{\pgfqpoint{10.653726in}{1.024556in}}{\pgfqpoint{10.660480in}{1.017459in}}%
\pgfpathquadraticcurveto{\pgfqpoint{10.667235in}{1.010380in}}{\pgfqpoint{10.667235in}{0.995970in}}%
\pgfpathlineto{\pgfqpoint{10.667235in}{0.995970in}}%
\pgfpathclose%
\pgfpathmoveto{\pgfqpoint{10.688561in}{1.047593in}}%
\pgfpathlineto{\pgfqpoint{10.698918in}{1.047593in}}%
\pgfpathlineto{\pgfqpoint{10.698918in}{0.960000in}}%
\pgfpathlineto{\pgfqpoint{10.688561in}{0.960000in}}%
\pgfpathlineto{\pgfqpoint{10.688561in}{1.047593in}}%
\pgfpathlineto{\pgfqpoint{10.688561in}{1.047593in}}%
\pgfpathclose%
\pgfusepath{fill}%
\end{pgfscope}%
\begin{pgfscope}%
\pgfsetbuttcap%
\pgfsetroundjoin%
\definecolor{currentfill}{rgb}{0.309804,0.372549,0.419608}%
\pgfsetfillcolor{currentfill}%
\pgfsetlinewidth{0.000000pt}%
\definecolor{currentstroke}{rgb}{0.309804,0.372549,0.419608}%
\pgfsetstrokecolor{currentstroke}%
\pgfsetdash{}{0pt}%
\pgfpathmoveto{\pgfqpoint{10.827015in}{0.992260in}}%
\pgfpathquadraticcurveto{\pgfqpoint{10.827015in}{1.003517in}}{\pgfqpoint{10.822368in}{1.009696in}}%
\pgfpathquadraticcurveto{\pgfqpoint{10.817738in}{1.015892in}}{\pgfqpoint{10.809345in}{1.015892in}}%
\pgfpathquadraticcurveto{\pgfqpoint{10.801023in}{1.015892in}}{\pgfqpoint{10.796376in}{1.009696in}}%
\pgfpathquadraticcurveto{\pgfqpoint{10.791729in}{1.003517in}}{\pgfqpoint{10.791729in}{0.992260in}}%
\pgfpathquadraticcurveto{\pgfqpoint{10.791729in}{0.981056in}}{\pgfqpoint{10.796376in}{0.974860in}}%
\pgfpathquadraticcurveto{\pgfqpoint{10.801023in}{0.968664in}}{\pgfqpoint{10.809345in}{0.968664in}}%
\pgfpathquadraticcurveto{\pgfqpoint{10.817738in}{0.968664in}}{\pgfqpoint{10.822368in}{0.974860in}}%
\pgfpathquadraticcurveto{\pgfqpoint{10.827015in}{0.981056in}}{\pgfqpoint{10.827015in}{0.992260in}}%
\pgfpathlineto{\pgfqpoint{10.827015in}{0.992260in}}%
\pgfpathclose%
\pgfpathmoveto{\pgfqpoint{10.837372in}{0.967817in}}%
\pgfpathquadraticcurveto{\pgfqpoint{10.837372in}{0.951732in}}{\pgfqpoint{10.830221in}{0.943879in}}%
\pgfpathquadraticcurveto{\pgfqpoint{10.823088in}{0.936026in}}{\pgfqpoint{10.808336in}{0.936026in}}%
\pgfpathquadraticcurveto{\pgfqpoint{10.802878in}{0.936026in}}{\pgfqpoint{10.798033in}{0.936836in}}%
\pgfpathquadraticcurveto{\pgfqpoint{10.793188in}{0.937647in}}{\pgfqpoint{10.788631in}{0.939340in}}%
\pgfpathlineto{\pgfqpoint{10.788631in}{0.949409in}}%
\pgfpathquadraticcurveto{\pgfqpoint{10.793188in}{0.946941in}}{\pgfqpoint{10.797637in}{0.945770in}}%
\pgfpathquadraticcurveto{\pgfqpoint{10.802086in}{0.944582in}}{\pgfqpoint{10.806697in}{0.944582in}}%
\pgfpathquadraticcurveto{\pgfqpoint{10.816892in}{0.944582in}}{\pgfqpoint{10.821953in}{0.949895in}}%
\pgfpathquadraticcurveto{\pgfqpoint{10.827015in}{0.955209in}}{\pgfqpoint{10.827015in}{0.965962in}}%
\pgfpathlineto{\pgfqpoint{10.827015in}{0.971095in}}%
\pgfpathquadraticcurveto{\pgfqpoint{10.823808in}{0.965512in}}{\pgfqpoint{10.818801in}{0.962756in}}%
\pgfpathquadraticcurveto{\pgfqpoint{10.813794in}{0.960000in}}{\pgfqpoint{10.806805in}{0.960000in}}%
\pgfpathquadraticcurveto{\pgfqpoint{10.795223in}{0.960000in}}{\pgfqpoint{10.788126in}{0.968826in}}%
\pgfpathquadraticcurveto{\pgfqpoint{10.781030in}{0.977670in}}{\pgfqpoint{10.781030in}{0.992260in}}%
\pgfpathquadraticcurveto{\pgfqpoint{10.781030in}{1.006886in}}{\pgfqpoint{10.788126in}{1.015712in}}%
\pgfpathquadraticcurveto{\pgfqpoint{10.795223in}{1.024556in}}{\pgfqpoint{10.806805in}{1.024556in}}%
\pgfpathquadraticcurveto{\pgfqpoint{10.813794in}{1.024556in}}{\pgfqpoint{10.818801in}{1.021800in}}%
\pgfpathquadraticcurveto{\pgfqpoint{10.823808in}{1.019044in}}{\pgfqpoint{10.827015in}{1.013478in}}%
\pgfpathlineto{\pgfqpoint{10.827015in}{1.023043in}}%
\pgfpathlineto{\pgfqpoint{10.837372in}{1.023043in}}%
\pgfpathlineto{\pgfqpoint{10.837372in}{0.967817in}}%
\pgfpathlineto{\pgfqpoint{10.837372in}{0.967817in}}%
\pgfpathclose%
\pgfpathmoveto{\pgfqpoint{10.887373in}{0.991683in}}%
\pgfpathquadraticcurveto{\pgfqpoint{10.874819in}{0.991683in}}{\pgfqpoint{10.869974in}{0.988819in}}%
\pgfpathquadraticcurveto{\pgfqpoint{10.865128in}{0.985956in}}{\pgfqpoint{10.865128in}{0.979021in}}%
\pgfpathquadraticcurveto{\pgfqpoint{10.865128in}{0.973509in}}{\pgfqpoint{10.868767in}{0.970267in}}%
\pgfpathquadraticcurveto{\pgfqpoint{10.872405in}{0.967043in}}{\pgfqpoint{10.878637in}{0.967043in}}%
\pgfpathquadraticcurveto{\pgfqpoint{10.887265in}{0.967043in}}{\pgfqpoint{10.892471in}{0.973149in}}%
\pgfpathquadraticcurveto{\pgfqpoint{10.897676in}{0.979255in}}{\pgfqpoint{10.897676in}{0.989378in}}%
\pgfpathlineto{\pgfqpoint{10.897676in}{0.991683in}}%
\pgfpathlineto{\pgfqpoint{10.887373in}{0.991683in}}%
\pgfpathlineto{\pgfqpoint{10.887373in}{0.991683in}}%
\pgfpathclose%
\pgfpathmoveto{\pgfqpoint{10.908033in}{0.995970in}}%
\pgfpathlineto{\pgfqpoint{10.908033in}{0.960000in}}%
\pgfpathlineto{\pgfqpoint{10.897676in}{0.960000in}}%
\pgfpathlineto{\pgfqpoint{10.897676in}{0.969564in}}%
\pgfpathquadraticcurveto{\pgfqpoint{10.894128in}{0.963837in}}{\pgfqpoint{10.888832in}{0.961099in}}%
\pgfpathquadraticcurveto{\pgfqpoint{10.883537in}{0.958361in}}{\pgfqpoint{10.875882in}{0.958361in}}%
\pgfpathquadraticcurveto{\pgfqpoint{10.866209in}{0.958361in}}{\pgfqpoint{10.860481in}{0.963801in}}%
\pgfpathquadraticcurveto{\pgfqpoint{10.854771in}{0.969240in}}{\pgfqpoint{10.854771in}{0.978354in}}%
\pgfpathquadraticcurveto{\pgfqpoint{10.854771in}{0.988982in}}{\pgfqpoint{10.861886in}{0.994385in}}%
\pgfpathquadraticcurveto{\pgfqpoint{10.869019in}{0.999789in}}{\pgfqpoint{10.883141in}{0.999789in}}%
\pgfpathlineto{\pgfqpoint{10.897676in}{0.999789in}}%
\pgfpathlineto{\pgfqpoint{10.897676in}{1.000816in}}%
\pgfpathquadraticcurveto{\pgfqpoint{10.897676in}{1.007966in}}{\pgfqpoint{10.892975in}{1.011875in}}%
\pgfpathquadraticcurveto{\pgfqpoint{10.888274in}{1.015784in}}{\pgfqpoint{10.879772in}{1.015784in}}%
\pgfpathquadraticcurveto{\pgfqpoint{10.874369in}{1.015784in}}{\pgfqpoint{10.869235in}{1.014487in}}%
\pgfpathquadraticcurveto{\pgfqpoint{10.864120in}{1.013190in}}{\pgfqpoint{10.859401in}{1.010596in}}%
\pgfpathlineto{\pgfqpoint{10.859401in}{1.020179in}}%
\pgfpathquadraticcurveto{\pgfqpoint{10.865074in}{1.022376in}}{\pgfqpoint{10.870424in}{1.023457in}}%
\pgfpathquadraticcurveto{\pgfqpoint{10.875774in}{1.024556in}}{\pgfqpoint{10.880835in}{1.024556in}}%
\pgfpathquadraticcurveto{\pgfqpoint{10.894524in}{1.024556in}}{\pgfqpoint{10.901279in}{1.017459in}}%
\pgfpathquadraticcurveto{\pgfqpoint{10.908033in}{1.010380in}}{\pgfqpoint{10.908033in}{0.995970in}}%
\pgfpathlineto{\pgfqpoint{10.908033in}{0.995970in}}%
\pgfpathclose%
\pgfpathmoveto{\pgfqpoint{10.939374in}{0.969456in}}%
\pgfpathlineto{\pgfqpoint{10.939374in}{0.936026in}}%
\pgfpathlineto{\pgfqpoint{10.928963in}{0.936026in}}%
\pgfpathlineto{\pgfqpoint{10.928963in}{1.023043in}}%
\pgfpathlineto{\pgfqpoint{10.939374in}{1.023043in}}%
\pgfpathlineto{\pgfqpoint{10.939374in}{1.013478in}}%
\pgfpathquadraticcurveto{\pgfqpoint{10.942653in}{1.019098in}}{\pgfqpoint{10.947624in}{1.021818in}}%
\pgfpathquadraticcurveto{\pgfqpoint{10.952613in}{1.024556in}}{\pgfqpoint{10.959530in}{1.024556in}}%
\pgfpathquadraticcurveto{\pgfqpoint{10.971022in}{1.024556in}}{\pgfqpoint{10.978191in}{1.015441in}}%
\pgfpathquadraticcurveto{\pgfqpoint{10.985377in}{1.006327in}}{\pgfqpoint{10.985377in}{0.991467in}}%
\pgfpathquadraticcurveto{\pgfqpoint{10.985377in}{0.976607in}}{\pgfqpoint{10.978191in}{0.967475in}}%
\pgfpathquadraticcurveto{\pgfqpoint{10.971022in}{0.958361in}}{\pgfqpoint{10.959530in}{0.958361in}}%
\pgfpathquadraticcurveto{\pgfqpoint{10.952613in}{0.958361in}}{\pgfqpoint{10.947624in}{0.961099in}}%
\pgfpathquadraticcurveto{\pgfqpoint{10.942653in}{0.963837in}}{\pgfqpoint{10.939374in}{0.969456in}}%
\pgfpathlineto{\pgfqpoint{10.939374in}{0.969456in}}%
\pgfpathclose%
\pgfpathmoveto{\pgfqpoint{10.974624in}{0.991467in}}%
\pgfpathquadraticcurveto{\pgfqpoint{10.974624in}{1.002887in}}{\pgfqpoint{10.969923in}{1.009389in}}%
\pgfpathquadraticcurveto{\pgfqpoint{10.965222in}{1.015892in}}{\pgfqpoint{10.957008in}{1.015892in}}%
\pgfpathquadraticcurveto{\pgfqpoint{10.948777in}{1.015892in}}{\pgfqpoint{10.944076in}{1.009389in}}%
\pgfpathquadraticcurveto{\pgfqpoint{10.939374in}{1.002887in}}{\pgfqpoint{10.939374in}{0.991467in}}%
\pgfpathquadraticcurveto{\pgfqpoint{10.939374in}{0.980048in}}{\pgfqpoint{10.944076in}{0.973545in}}%
\pgfpathquadraticcurveto{\pgfqpoint{10.948777in}{0.967043in}}{\pgfqpoint{10.957008in}{0.967043in}}%
\pgfpathquadraticcurveto{\pgfqpoint{10.965222in}{0.967043in}}{\pgfqpoint{10.969923in}{0.973545in}}%
\pgfpathquadraticcurveto{\pgfqpoint{10.974624in}{0.980048in}}{\pgfqpoint{10.974624in}{0.991467in}}%
\pgfpathlineto{\pgfqpoint{10.974624in}{0.991467in}}%
\pgfpathclose%
\pgfpathmoveto{\pgfqpoint{11.042728in}{1.021187in}}%
\pgfpathlineto{\pgfqpoint{11.042728in}{1.011389in}}%
\pgfpathquadraticcurveto{\pgfqpoint{11.038351in}{1.013640in}}{\pgfqpoint{11.033614in}{1.014757in}}%
\pgfpathquadraticcurveto{\pgfqpoint{11.028895in}{1.015892in}}{\pgfqpoint{11.023815in}{1.015892in}}%
\pgfpathquadraticcurveto{\pgfqpoint{11.016106in}{1.015892in}}{\pgfqpoint{11.012252in}{1.013532in}}%
\pgfpathquadraticcurveto{\pgfqpoint{11.008397in}{1.011173in}}{\pgfqpoint{11.008397in}{1.006435in}}%
\pgfpathquadraticcurveto{\pgfqpoint{11.008397in}{1.002833in}}{\pgfqpoint{11.011153in}{1.000780in}}%
\pgfpathquadraticcurveto{\pgfqpoint{11.013909in}{0.998726in}}{\pgfqpoint{11.022248in}{0.996871in}}%
\pgfpathlineto{\pgfqpoint{11.025797in}{0.996078in}}%
\pgfpathquadraticcurveto{\pgfqpoint{11.036820in}{0.993719in}}{\pgfqpoint{11.041467in}{0.989414in}}%
\pgfpathquadraticcurveto{\pgfqpoint{11.046114in}{0.985109in}}{\pgfqpoint{11.046114in}{0.977400in}}%
\pgfpathquadraticcurveto{\pgfqpoint{11.046114in}{0.968610in}}{\pgfqpoint{11.039162in}{0.963476in}}%
\pgfpathquadraticcurveto{\pgfqpoint{11.032209in}{0.958361in}}{\pgfqpoint{11.020051in}{0.958361in}}%
\pgfpathquadraticcurveto{\pgfqpoint{11.014989in}{0.958361in}}{\pgfqpoint{11.009496in}{0.959352in}}%
\pgfpathquadraticcurveto{\pgfqpoint{11.004002in}{0.960342in}}{\pgfqpoint{10.997932in}{0.962306in}}%
\pgfpathlineto{\pgfqpoint{10.997932in}{0.973005in}}%
\pgfpathquadraticcurveto{\pgfqpoint{11.003678in}{0.970015in}}{\pgfqpoint{11.009244in}{0.968520in}}%
\pgfpathquadraticcurveto{\pgfqpoint{11.014809in}{0.967043in}}{\pgfqpoint{11.020285in}{0.967043in}}%
\pgfpathquadraticcurveto{\pgfqpoint{11.027598in}{0.967043in}}{\pgfqpoint{11.031525in}{0.969546in}}%
\pgfpathquadraticcurveto{\pgfqpoint{11.035469in}{0.972050in}}{\pgfqpoint{11.035469in}{0.976607in}}%
\pgfpathquadraticcurveto{\pgfqpoint{11.035469in}{0.980822in}}{\pgfqpoint{11.032623in}{0.983074in}}%
\pgfpathquadraticcurveto{\pgfqpoint{11.029795in}{0.985325in}}{\pgfqpoint{11.020159in}{0.987414in}}%
\pgfpathlineto{\pgfqpoint{11.016557in}{0.988261in}}%
\pgfpathquadraticcurveto{\pgfqpoint{11.006938in}{0.990278in}}{\pgfqpoint{11.002651in}{0.994475in}}%
\pgfpathquadraticcurveto{\pgfqpoint{10.998382in}{0.998672in}}{\pgfqpoint{10.998382in}{1.005985in}}%
\pgfpathquadraticcurveto{\pgfqpoint{10.998382in}{1.014883in}}{\pgfqpoint{11.004687in}{1.019710in}}%
\pgfpathquadraticcurveto{\pgfqpoint{11.010991in}{1.024556in}}{\pgfqpoint{11.022591in}{1.024556in}}%
\pgfpathquadraticcurveto{\pgfqpoint{11.028318in}{1.024556in}}{\pgfqpoint{11.033380in}{1.023709in}}%
\pgfpathquadraticcurveto{\pgfqpoint{11.038459in}{1.022880in}}{\pgfqpoint{11.042728in}{1.021187in}}%
\pgfpathlineto{\pgfqpoint{11.042728in}{1.021187in}}%
\pgfpathclose%
\pgfpathmoveto{\pgfqpoint{11.064055in}{0.974302in}}%
\pgfpathlineto{\pgfqpoint{11.075943in}{0.974302in}}%
\pgfpathlineto{\pgfqpoint{11.075943in}{0.960000in}}%
\pgfpathlineto{\pgfqpoint{11.064055in}{0.960000in}}%
\pgfpathlineto{\pgfqpoint{11.064055in}{0.974302in}}%
\pgfpathlineto{\pgfqpoint{11.064055in}{0.974302in}}%
\pgfpathclose%
\pgfusepath{fill}%
\end{pgfscope}%
\begin{pgfscope}%
\pgfsetbuttcap%
\pgfsetroundjoin%
\definecolor{currentfill}{rgb}{0.309804,0.372549,0.419608}%
\pgfsetfillcolor{currentfill}%
\pgfsetlinewidth{0.000000pt}%
\definecolor{currentstroke}{rgb}{0.309804,0.372549,0.419608}%
\pgfsetstrokecolor{currentstroke}%
\pgfsetdash{}{0pt}%
\pgfpathmoveto{\pgfqpoint{11.165876in}{1.044045in}}%
\pgfpathlineto{\pgfqpoint{11.219012in}{1.044045in}}%
\pgfpathlineto{\pgfqpoint{11.219012in}{1.034462in}}%
\pgfpathlineto{\pgfqpoint{11.177242in}{1.034462in}}%
\pgfpathlineto{\pgfqpoint{11.177242in}{1.009587in}}%
\pgfpathlineto{\pgfqpoint{11.217265in}{1.009587in}}%
\pgfpathlineto{\pgfqpoint{11.217265in}{1.000023in}}%
\pgfpathlineto{\pgfqpoint{11.177242in}{1.000023in}}%
\pgfpathlineto{\pgfqpoint{11.177242in}{0.969564in}}%
\pgfpathlineto{\pgfqpoint{11.220021in}{0.969564in}}%
\pgfpathlineto{\pgfqpoint{11.220021in}{0.960000in}}%
\pgfpathlineto{\pgfqpoint{11.165876in}{0.960000in}}%
\pgfpathlineto{\pgfqpoint{11.165876in}{1.044045in}}%
\pgfpathlineto{\pgfqpoint{11.165876in}{1.044045in}}%
\pgfpathclose%
\pgfpathmoveto{\pgfqpoint{11.290682in}{1.023043in}}%
\pgfpathlineto{\pgfqpoint{11.267879in}{0.992368in}}%
\pgfpathlineto{\pgfqpoint{11.291853in}{0.960000in}}%
\pgfpathlineto{\pgfqpoint{11.279641in}{0.960000in}}%
\pgfpathlineto{\pgfqpoint{11.261287in}{0.984767in}}%
\pgfpathlineto{\pgfqpoint{11.242950in}{0.960000in}}%
\pgfpathlineto{\pgfqpoint{11.230720in}{0.960000in}}%
\pgfpathlineto{\pgfqpoint{11.255217in}{0.992980in}}%
\pgfpathlineto{\pgfqpoint{11.232809in}{1.023043in}}%
\pgfpathlineto{\pgfqpoint{11.245022in}{1.023043in}}%
\pgfpathlineto{\pgfqpoint{11.261737in}{1.000581in}}%
\pgfpathlineto{\pgfqpoint{11.278452in}{1.023043in}}%
\pgfpathlineto{\pgfqpoint{11.290682in}{1.023043in}}%
\pgfpathlineto{\pgfqpoint{11.290682in}{1.023043in}}%
\pgfpathclose%
\pgfpathmoveto{\pgfqpoint{11.349852in}{1.020629in}}%
\pgfpathlineto{\pgfqpoint{11.349852in}{1.010938in}}%
\pgfpathquadraticcurveto{\pgfqpoint{11.345457in}{1.013370in}}{\pgfqpoint{11.341044in}{1.014577in}}%
\pgfpathquadraticcurveto{\pgfqpoint{11.336631in}{1.015784in}}{\pgfqpoint{11.332128in}{1.015784in}}%
\pgfpathquadraticcurveto{\pgfqpoint{11.322042in}{1.015784in}}{\pgfqpoint{11.316458in}{1.009389in}}%
\pgfpathquadraticcurveto{\pgfqpoint{11.310892in}{1.003013in}}{\pgfqpoint{11.310892in}{0.991467in}}%
\pgfpathquadraticcurveto{\pgfqpoint{11.310892in}{0.979921in}}{\pgfqpoint{11.316458in}{0.973527in}}%
\pgfpathquadraticcurveto{\pgfqpoint{11.322042in}{0.967151in}}{\pgfqpoint{11.332128in}{0.967151in}}%
\pgfpathquadraticcurveto{\pgfqpoint{11.336631in}{0.967151in}}{\pgfqpoint{11.341044in}{0.968358in}}%
\pgfpathquadraticcurveto{\pgfqpoint{11.345457in}{0.969564in}}{\pgfqpoint{11.349852in}{0.971996in}}%
\pgfpathlineto{\pgfqpoint{11.349852in}{0.962414in}}%
\pgfpathquadraticcurveto{\pgfqpoint{11.345511in}{0.960396in}}{\pgfqpoint{11.340864in}{0.959388in}}%
\pgfpathquadraticcurveto{\pgfqpoint{11.336235in}{0.958361in}}{\pgfqpoint{11.330994in}{0.958361in}}%
\pgfpathquadraticcurveto{\pgfqpoint{11.316746in}{0.958361in}}{\pgfqpoint{11.308352in}{0.967313in}}%
\pgfpathquadraticcurveto{\pgfqpoint{11.299977in}{0.976265in}}{\pgfqpoint{11.299977in}{0.991467in}}%
\pgfpathquadraticcurveto{\pgfqpoint{11.299977in}{1.006886in}}{\pgfqpoint{11.308442in}{1.015712in}}%
\pgfpathquadraticcurveto{\pgfqpoint{11.316926in}{1.024556in}}{\pgfqpoint{11.331678in}{1.024556in}}%
\pgfpathquadraticcurveto{\pgfqpoint{11.336451in}{1.024556in}}{\pgfqpoint{11.341008in}{1.023565in}}%
\pgfpathquadraticcurveto{\pgfqpoint{11.345566in}{1.022592in}}{\pgfqpoint{11.349852in}{1.020629in}}%
\pgfpathlineto{\pgfqpoint{11.349852in}{1.020629in}}%
\pgfpathclose%
\pgfpathmoveto{\pgfqpoint{11.367865in}{1.047593in}}%
\pgfpathlineto{\pgfqpoint{11.378222in}{1.047593in}}%
\pgfpathlineto{\pgfqpoint{11.378222in}{0.960000in}}%
\pgfpathlineto{\pgfqpoint{11.367865in}{0.960000in}}%
\pgfpathlineto{\pgfqpoint{11.367865in}{1.047593in}}%
\pgfpathlineto{\pgfqpoint{11.367865in}{1.047593in}}%
\pgfpathclose%
\pgfpathmoveto{\pgfqpoint{11.398827in}{0.984875in}}%
\pgfpathlineto{\pgfqpoint{11.398827in}{1.023043in}}%
\pgfpathlineto{\pgfqpoint{11.409184in}{1.023043in}}%
\pgfpathlineto{\pgfqpoint{11.409184in}{0.985271in}}%
\pgfpathquadraticcurveto{\pgfqpoint{11.409184in}{0.976319in}}{\pgfqpoint{11.412661in}{0.971834in}}%
\pgfpathquadraticcurveto{\pgfqpoint{11.416155in}{0.967367in}}{\pgfqpoint{11.423144in}{0.967367in}}%
\pgfpathquadraticcurveto{\pgfqpoint{11.431520in}{0.967367in}}{\pgfqpoint{11.436383in}{0.972717in}}%
\pgfpathquadraticcurveto{\pgfqpoint{11.441264in}{0.978066in}}{\pgfqpoint{11.441264in}{0.987306in}}%
\pgfpathlineto{\pgfqpoint{11.441264in}{1.023043in}}%
\pgfpathlineto{\pgfqpoint{11.451621in}{1.023043in}}%
\pgfpathlineto{\pgfqpoint{11.451621in}{0.960000in}}%
\pgfpathlineto{\pgfqpoint{11.441264in}{0.960000in}}%
\pgfpathlineto{\pgfqpoint{11.441264in}{0.969691in}}%
\pgfpathquadraticcurveto{\pgfqpoint{11.437500in}{0.963945in}}{\pgfqpoint{11.432510in}{0.961153in}}%
\pgfpathquadraticcurveto{\pgfqpoint{11.427539in}{0.958361in}}{\pgfqpoint{11.420946in}{0.958361in}}%
\pgfpathquadraticcurveto{\pgfqpoint{11.410085in}{0.958361in}}{\pgfqpoint{11.404447in}{0.965115in}}%
\pgfpathquadraticcurveto{\pgfqpoint{11.398827in}{0.971870in}}{\pgfqpoint{11.398827in}{0.984875in}}%
\pgfpathlineto{\pgfqpoint{11.398827in}{0.984875in}}%
\pgfpathclose%
\pgfpathmoveto{\pgfqpoint{11.424891in}{1.024556in}}%
\pgfpathlineto{\pgfqpoint{11.424891in}{1.024556in}}%
\pgfpathlineto{\pgfqpoint{11.424891in}{1.024556in}}%
\pgfpathclose%
\pgfpathmoveto{\pgfqpoint{11.514429in}{1.013478in}}%
\pgfpathlineto{\pgfqpoint{11.514429in}{1.047593in}}%
\pgfpathlineto{\pgfqpoint{11.524786in}{1.047593in}}%
\pgfpathlineto{\pgfqpoint{11.524786in}{0.960000in}}%
\pgfpathlineto{\pgfqpoint{11.514429in}{0.960000in}}%
\pgfpathlineto{\pgfqpoint{11.514429in}{0.969456in}}%
\pgfpathquadraticcurveto{\pgfqpoint{11.511169in}{0.963837in}}{\pgfqpoint{11.506180in}{0.961099in}}%
\pgfpathquadraticcurveto{\pgfqpoint{11.501209in}{0.958361in}}{\pgfqpoint{11.494220in}{0.958361in}}%
\pgfpathquadraticcurveto{\pgfqpoint{11.482800in}{0.958361in}}{\pgfqpoint{11.475613in}{0.967475in}}%
\pgfpathquadraticcurveto{\pgfqpoint{11.468444in}{0.976607in}}{\pgfqpoint{11.468444in}{0.991467in}}%
\pgfpathquadraticcurveto{\pgfqpoint{11.468444in}{1.006327in}}{\pgfqpoint{11.475613in}{1.015441in}}%
\pgfpathquadraticcurveto{\pgfqpoint{11.482800in}{1.024556in}}{\pgfqpoint{11.494220in}{1.024556in}}%
\pgfpathquadraticcurveto{\pgfqpoint{11.501209in}{1.024556in}}{\pgfqpoint{11.506180in}{1.021818in}}%
\pgfpathquadraticcurveto{\pgfqpoint{11.511169in}{1.019098in}}{\pgfqpoint{11.514429in}{1.013478in}}%
\pgfpathlineto{\pgfqpoint{11.514429in}{1.013478in}}%
\pgfpathclose%
\pgfpathmoveto{\pgfqpoint{11.479144in}{0.991467in}}%
\pgfpathquadraticcurveto{\pgfqpoint{11.479144in}{0.980048in}}{\pgfqpoint{11.483845in}{0.973545in}}%
\pgfpathquadraticcurveto{\pgfqpoint{11.488546in}{0.967043in}}{\pgfqpoint{11.496760in}{0.967043in}}%
\pgfpathquadraticcurveto{\pgfqpoint{11.504973in}{0.967043in}}{\pgfqpoint{11.509692in}{0.973545in}}%
\pgfpathquadraticcurveto{\pgfqpoint{11.514429in}{0.980048in}}{\pgfqpoint{11.514429in}{0.991467in}}%
\pgfpathquadraticcurveto{\pgfqpoint{11.514429in}{1.002887in}}{\pgfqpoint{11.509692in}{1.009389in}}%
\pgfpathquadraticcurveto{\pgfqpoint{11.504973in}{1.015892in}}{\pgfqpoint{11.496760in}{1.015892in}}%
\pgfpathquadraticcurveto{\pgfqpoint{11.488546in}{1.015892in}}{\pgfqpoint{11.483845in}{1.009389in}}%
\pgfpathquadraticcurveto{\pgfqpoint{11.479144in}{1.002887in}}{\pgfqpoint{11.479144in}{0.991467in}}%
\pgfpathlineto{\pgfqpoint{11.479144in}{0.991467in}}%
\pgfpathclose%
\pgfpathmoveto{\pgfqpoint{11.600059in}{0.994115in}}%
\pgfpathlineto{\pgfqpoint{11.600059in}{0.989054in}}%
\pgfpathlineto{\pgfqpoint{11.552435in}{0.989054in}}%
\pgfpathquadraticcurveto{\pgfqpoint{11.553120in}{0.978354in}}{\pgfqpoint{11.558883in}{0.972753in}}%
\pgfpathquadraticcurveto{\pgfqpoint{11.564647in}{0.967151in}}{\pgfqpoint{11.574950in}{0.967151in}}%
\pgfpathquadraticcurveto{\pgfqpoint{11.580912in}{0.967151in}}{\pgfqpoint{11.586514in}{0.968610in}}%
\pgfpathquadraticcurveto{\pgfqpoint{11.592116in}{0.970069in}}{\pgfqpoint{11.597646in}{0.973005in}}%
\pgfpathlineto{\pgfqpoint{11.597646in}{0.963206in}}%
\pgfpathquadraticcurveto{\pgfqpoint{11.592062in}{0.960847in}}{\pgfqpoint{11.586208in}{0.959604in}}%
\pgfpathquadraticcurveto{\pgfqpoint{11.580354in}{0.958361in}}{\pgfqpoint{11.574338in}{0.958361in}}%
\pgfpathquadraticcurveto{\pgfqpoint{11.559244in}{0.958361in}}{\pgfqpoint{11.550436in}{0.967133in}}%
\pgfpathquadraticcurveto{\pgfqpoint{11.541628in}{0.975923in}}{\pgfqpoint{11.541628in}{0.990909in}}%
\pgfpathquadraticcurveto{\pgfqpoint{11.541628in}{1.006381in}}{\pgfqpoint{11.549985in}{1.015459in}}%
\pgfpathquadraticcurveto{\pgfqpoint{11.558343in}{1.024556in}}{\pgfqpoint{11.572537in}{1.024556in}}%
\pgfpathquadraticcurveto{\pgfqpoint{11.585253in}{1.024556in}}{\pgfqpoint{11.592656in}{1.016360in}}%
\pgfpathquadraticcurveto{\pgfqpoint{11.600059in}{1.008183in}}{\pgfqpoint{11.600059in}{0.994115in}}%
\pgfpathlineto{\pgfqpoint{11.600059in}{0.994115in}}%
\pgfpathclose%
\pgfpathmoveto{\pgfqpoint{11.589702in}{0.997159in}}%
\pgfpathquadraticcurveto{\pgfqpoint{11.589594in}{1.005643in}}{\pgfqpoint{11.584947in}{1.010704in}}%
\pgfpathquadraticcurveto{\pgfqpoint{11.580300in}{1.015784in}}{\pgfqpoint{11.572645in}{1.015784in}}%
\pgfpathquadraticcurveto{\pgfqpoint{11.563981in}{1.015784in}}{\pgfqpoint{11.558775in}{1.010884in}}%
\pgfpathquadraticcurveto{\pgfqpoint{11.553570in}{1.005985in}}{\pgfqpoint{11.552777in}{0.997087in}}%
\pgfpathlineto{\pgfqpoint{11.589702in}{0.997159in}}%
\pgfpathlineto{\pgfqpoint{11.589702in}{0.997159in}}%
\pgfpathclose%
\pgfpathmoveto{\pgfqpoint{11.658545in}{1.013478in}}%
\pgfpathlineto{\pgfqpoint{11.658545in}{1.047593in}}%
\pgfpathlineto{\pgfqpoint{11.668902in}{1.047593in}}%
\pgfpathlineto{\pgfqpoint{11.668902in}{0.960000in}}%
\pgfpathlineto{\pgfqpoint{11.658545in}{0.960000in}}%
\pgfpathlineto{\pgfqpoint{11.658545in}{0.969456in}}%
\pgfpathquadraticcurveto{\pgfqpoint{11.655284in}{0.963837in}}{\pgfqpoint{11.650295in}{0.961099in}}%
\pgfpathquadraticcurveto{\pgfqpoint{11.645324in}{0.958361in}}{\pgfqpoint{11.638335in}{0.958361in}}%
\pgfpathquadraticcurveto{\pgfqpoint{11.626915in}{0.958361in}}{\pgfqpoint{11.619729in}{0.967475in}}%
\pgfpathquadraticcurveto{\pgfqpoint{11.612560in}{0.976607in}}{\pgfqpoint{11.612560in}{0.991467in}}%
\pgfpathquadraticcurveto{\pgfqpoint{11.612560in}{1.006327in}}{\pgfqpoint{11.619729in}{1.015441in}}%
\pgfpathquadraticcurveto{\pgfqpoint{11.626915in}{1.024556in}}{\pgfqpoint{11.638335in}{1.024556in}}%
\pgfpathquadraticcurveto{\pgfqpoint{11.645324in}{1.024556in}}{\pgfqpoint{11.650295in}{1.021818in}}%
\pgfpathquadraticcurveto{\pgfqpoint{11.655284in}{1.019098in}}{\pgfqpoint{11.658545in}{1.013478in}}%
\pgfpathlineto{\pgfqpoint{11.658545in}{1.013478in}}%
\pgfpathclose%
\pgfpathmoveto{\pgfqpoint{11.623259in}{0.991467in}}%
\pgfpathquadraticcurveto{\pgfqpoint{11.623259in}{0.980048in}}{\pgfqpoint{11.627960in}{0.973545in}}%
\pgfpathquadraticcurveto{\pgfqpoint{11.632661in}{0.967043in}}{\pgfqpoint{11.640875in}{0.967043in}}%
\pgfpathquadraticcurveto{\pgfqpoint{11.649088in}{0.967043in}}{\pgfqpoint{11.653807in}{0.973545in}}%
\pgfpathquadraticcurveto{\pgfqpoint{11.658545in}{0.980048in}}{\pgfqpoint{11.658545in}{0.991467in}}%
\pgfpathquadraticcurveto{\pgfqpoint{11.658545in}{1.002887in}}{\pgfqpoint{11.653807in}{1.009389in}}%
\pgfpathquadraticcurveto{\pgfqpoint{11.649088in}{1.015892in}}{\pgfqpoint{11.640875in}{1.015892in}}%
\pgfpathquadraticcurveto{\pgfqpoint{11.632661in}{1.015892in}}{\pgfqpoint{11.627960in}{1.009389in}}%
\pgfpathquadraticcurveto{\pgfqpoint{11.623259in}{1.002887in}}{\pgfqpoint{11.623259in}{0.991467in}}%
\pgfpathlineto{\pgfqpoint{11.623259in}{0.991467in}}%
\pgfpathclose%
\pgfusepath{fill}%
\end{pgfscope}%
\begin{pgfscope}%
\pgfsetbuttcap%
\pgfsetroundjoin%
\definecolor{currentfill}{rgb}{0.309804,0.372549,0.419608}%
\pgfsetfillcolor{currentfill}%
\pgfsetlinewidth{0.000000pt}%
\definecolor{currentstroke}{rgb}{0.309804,0.372549,0.419608}%
\pgfsetstrokecolor{currentstroke}%
\pgfsetdash{}{0pt}%
\pgfpathmoveto{\pgfqpoint{11.773288in}{1.023043in}}%
\pgfpathlineto{\pgfqpoint{11.783645in}{1.023043in}}%
\pgfpathlineto{\pgfqpoint{11.783645in}{0.960000in}}%
\pgfpathlineto{\pgfqpoint{11.773288in}{0.960000in}}%
\pgfpathlineto{\pgfqpoint{11.773288in}{1.023043in}}%
\pgfpathlineto{\pgfqpoint{11.773288in}{1.023043in}}%
\pgfpathclose%
\pgfpathmoveto{\pgfqpoint{11.773288in}{1.047593in}}%
\pgfpathlineto{\pgfqpoint{11.783645in}{1.047593in}}%
\pgfpathlineto{\pgfqpoint{11.783645in}{1.034462in}}%
\pgfpathlineto{\pgfqpoint{11.773288in}{1.034462in}}%
\pgfpathlineto{\pgfqpoint{11.773288in}{1.047593in}}%
\pgfpathlineto{\pgfqpoint{11.773288in}{1.047593in}}%
\pgfpathclose%
\pgfpathmoveto{\pgfqpoint{11.837231in}{1.047593in}}%
\pgfpathlineto{\pgfqpoint{11.837231in}{1.038965in}}%
\pgfpathlineto{\pgfqpoint{11.827325in}{1.038965in}}%
\pgfpathquadraticcurveto{\pgfqpoint{11.821759in}{1.038965in}}{\pgfqpoint{11.819579in}{1.036714in}}%
\pgfpathquadraticcurveto{\pgfqpoint{11.817418in}{1.034462in}}{\pgfqpoint{11.817418in}{1.028608in}}%
\pgfpathlineto{\pgfqpoint{11.817418in}{1.023043in}}%
\pgfpathlineto{\pgfqpoint{11.834475in}{1.023043in}}%
\pgfpathlineto{\pgfqpoint{11.834475in}{1.014991in}}%
\pgfpathlineto{\pgfqpoint{11.817418in}{1.014991in}}%
\pgfpathlineto{\pgfqpoint{11.817418in}{0.960000in}}%
\pgfpathlineto{\pgfqpoint{11.807007in}{0.960000in}}%
\pgfpathlineto{\pgfqpoint{11.807007in}{1.014991in}}%
\pgfpathlineto{\pgfqpoint{11.797100in}{1.014991in}}%
\pgfpathlineto{\pgfqpoint{11.797100in}{1.023043in}}%
\pgfpathlineto{\pgfqpoint{11.807007in}{1.023043in}}%
\pgfpathlineto{\pgfqpoint{11.807007in}{1.027438in}}%
\pgfpathquadraticcurveto{\pgfqpoint{11.807007in}{1.037957in}}{\pgfqpoint{11.811906in}{1.042766in}}%
\pgfpathquadraticcurveto{\pgfqpoint{11.816805in}{1.047593in}}{\pgfqpoint{11.827433in}{1.047593in}}%
\pgfpathlineto{\pgfqpoint{11.837231in}{1.047593in}}%
\pgfpathlineto{\pgfqpoint{11.837231in}{1.047593in}}%
\pgfpathclose%
\pgfusepath{fill}%
\end{pgfscope}%
\begin{pgfscope}%
\pgfsetbuttcap%
\pgfsetroundjoin%
\definecolor{currentfill}{rgb}{0.309804,0.372549,0.419608}%
\pgfsetfillcolor{currentfill}%
\pgfsetlinewidth{0.000000pt}%
\definecolor{currentstroke}{rgb}{0.309804,0.372549,0.419608}%
\pgfsetstrokecolor{currentstroke}%
\pgfsetdash{}{0pt}%
\pgfpathmoveto{\pgfqpoint{11.947110in}{0.994115in}}%
\pgfpathlineto{\pgfqpoint{11.947110in}{0.989054in}}%
\pgfpathlineto{\pgfqpoint{11.899486in}{0.989054in}}%
\pgfpathquadraticcurveto{\pgfqpoint{11.900170in}{0.978354in}}{\pgfqpoint{11.905934in}{0.972753in}}%
\pgfpathquadraticcurveto{\pgfqpoint{11.911698in}{0.967151in}}{\pgfqpoint{11.922001in}{0.967151in}}%
\pgfpathquadraticcurveto{\pgfqpoint{11.927963in}{0.967151in}}{\pgfqpoint{11.933565in}{0.968610in}}%
\pgfpathquadraticcurveto{\pgfqpoint{11.939166in}{0.970069in}}{\pgfqpoint{11.944696in}{0.973005in}}%
\pgfpathlineto{\pgfqpoint{11.944696in}{0.963206in}}%
\pgfpathquadraticcurveto{\pgfqpoint{11.939112in}{0.960847in}}{\pgfqpoint{11.933258in}{0.959604in}}%
\pgfpathquadraticcurveto{\pgfqpoint{11.927404in}{0.958361in}}{\pgfqpoint{11.921388in}{0.958361in}}%
\pgfpathquadraticcurveto{\pgfqpoint{11.906294in}{0.958361in}}{\pgfqpoint{11.897486in}{0.967133in}}%
\pgfpathquadraticcurveto{\pgfqpoint{11.888678in}{0.975923in}}{\pgfqpoint{11.888678in}{0.990909in}}%
\pgfpathquadraticcurveto{\pgfqpoint{11.888678in}{1.006381in}}{\pgfqpoint{11.897036in}{1.015459in}}%
\pgfpathquadraticcurveto{\pgfqpoint{11.905394in}{1.024556in}}{\pgfqpoint{11.919587in}{1.024556in}}%
\pgfpathquadraticcurveto{\pgfqpoint{11.932304in}{1.024556in}}{\pgfqpoint{11.939707in}{1.016360in}}%
\pgfpathquadraticcurveto{\pgfqpoint{11.947110in}{1.008183in}}{\pgfqpoint{11.947110in}{0.994115in}}%
\pgfpathlineto{\pgfqpoint{11.947110in}{0.994115in}}%
\pgfpathclose%
\pgfpathmoveto{\pgfqpoint{11.936753in}{0.997159in}}%
\pgfpathquadraticcurveto{\pgfqpoint{11.936645in}{1.005643in}}{\pgfqpoint{11.931998in}{1.010704in}}%
\pgfpathquadraticcurveto{\pgfqpoint{11.927350in}{1.015784in}}{\pgfqpoint{11.919695in}{1.015784in}}%
\pgfpathquadraticcurveto{\pgfqpoint{11.911031in}{1.015784in}}{\pgfqpoint{11.905826in}{1.010884in}}%
\pgfpathquadraticcurveto{\pgfqpoint{11.900620in}{1.005985in}}{\pgfqpoint{11.899828in}{0.997087in}}%
\pgfpathlineto{\pgfqpoint{11.936753in}{0.997159in}}%
\pgfpathlineto{\pgfqpoint{11.936753in}{0.997159in}}%
\pgfpathclose%
\pgfpathmoveto{\pgfqpoint{11.964113in}{1.023043in}}%
\pgfpathlineto{\pgfqpoint{11.974470in}{1.023043in}}%
\pgfpathlineto{\pgfqpoint{11.974470in}{0.960000in}}%
\pgfpathlineto{\pgfqpoint{11.964113in}{0.960000in}}%
\pgfpathlineto{\pgfqpoint{11.964113in}{1.023043in}}%
\pgfpathlineto{\pgfqpoint{11.964113in}{1.023043in}}%
\pgfpathclose%
\pgfpathmoveto{\pgfqpoint{11.964113in}{1.047593in}}%
\pgfpathlineto{\pgfqpoint{11.974470in}{1.047593in}}%
\pgfpathlineto{\pgfqpoint{11.974470in}{1.034462in}}%
\pgfpathlineto{\pgfqpoint{11.964113in}{1.034462in}}%
\pgfpathlineto{\pgfqpoint{11.964113in}{1.047593in}}%
\pgfpathlineto{\pgfqpoint{11.964113in}{1.047593in}}%
\pgfpathclose%
\pgfpathmoveto{\pgfqpoint{12.006388in}{1.040947in}}%
\pgfpathlineto{\pgfqpoint{12.006388in}{1.023043in}}%
\pgfpathlineto{\pgfqpoint{12.027714in}{1.023043in}}%
\pgfpathlineto{\pgfqpoint{12.027714in}{1.014991in}}%
\pgfpathlineto{\pgfqpoint{12.006388in}{1.014991in}}%
\pgfpathlineto{\pgfqpoint{12.006388in}{0.980768in}}%
\pgfpathquadraticcurveto{\pgfqpoint{12.006388in}{0.973059in}}{\pgfqpoint{12.008495in}{0.970861in}}%
\pgfpathquadraticcurveto{\pgfqpoint{12.010603in}{0.968664in}}{\pgfqpoint{12.017087in}{0.968664in}}%
\pgfpathlineto{\pgfqpoint{12.027714in}{0.968664in}}%
\pgfpathlineto{\pgfqpoint{12.027714in}{0.960000in}}%
\pgfpathlineto{\pgfqpoint{12.017087in}{0.960000in}}%
\pgfpathquadraticcurveto{\pgfqpoint{12.005091in}{0.960000in}}{\pgfqpoint{12.000534in}{0.964467in}}%
\pgfpathquadraticcurveto{\pgfqpoint{11.995977in}{0.968952in}}{\pgfqpoint{11.995977in}{0.980768in}}%
\pgfpathlineto{\pgfqpoint{11.995977in}{1.014991in}}%
\pgfpathlineto{\pgfqpoint{11.988376in}{1.014991in}}%
\pgfpathlineto{\pgfqpoint{11.988376in}{1.023043in}}%
\pgfpathlineto{\pgfqpoint{11.995977in}{1.023043in}}%
\pgfpathlineto{\pgfqpoint{11.995977in}{1.040947in}}%
\pgfpathlineto{\pgfqpoint{12.006388in}{1.040947in}}%
\pgfpathlineto{\pgfqpoint{12.006388in}{1.040947in}}%
\pgfpathclose%
\pgfpathmoveto{\pgfqpoint{12.093747in}{0.998060in}}%
\pgfpathlineto{\pgfqpoint{12.093747in}{0.960000in}}%
\pgfpathlineto{\pgfqpoint{12.083390in}{0.960000in}}%
\pgfpathlineto{\pgfqpoint{12.083390in}{0.997717in}}%
\pgfpathquadraticcurveto{\pgfqpoint{12.083390in}{1.006669in}}{\pgfqpoint{12.079895in}{1.011100in}}%
\pgfpathquadraticcurveto{\pgfqpoint{12.076401in}{1.015549in}}{\pgfqpoint{12.069430in}{1.015549in}}%
\pgfpathquadraticcurveto{\pgfqpoint{12.061037in}{1.015549in}}{\pgfqpoint{12.056191in}{1.010200in}}%
\pgfpathquadraticcurveto{\pgfqpoint{12.051346in}{1.004868in}}{\pgfqpoint{12.051346in}{0.995628in}}%
\pgfpathlineto{\pgfqpoint{12.051346in}{0.960000in}}%
\pgfpathlineto{\pgfqpoint{12.040935in}{0.960000in}}%
\pgfpathlineto{\pgfqpoint{12.040935in}{1.047593in}}%
\pgfpathlineto{\pgfqpoint{12.051346in}{1.047593in}}%
\pgfpathlineto{\pgfqpoint{12.051346in}{1.013244in}}%
\pgfpathquadraticcurveto{\pgfqpoint{12.055075in}{1.018936in}}{\pgfqpoint{12.060100in}{1.021746in}}%
\pgfpathquadraticcurveto{\pgfqpoint{12.065143in}{1.024556in}}{\pgfqpoint{12.071736in}{1.024556in}}%
\pgfpathquadraticcurveto{\pgfqpoint{12.082597in}{1.024556in}}{\pgfqpoint{12.088163in}{1.017837in}}%
\pgfpathquadraticcurveto{\pgfqpoint{12.093747in}{1.011118in}}{\pgfqpoint{12.093747in}{0.998060in}}%
\pgfpathlineto{\pgfqpoint{12.093747in}{0.998060in}}%
\pgfpathclose%
\pgfpathmoveto{\pgfqpoint{12.168317in}{0.994115in}}%
\pgfpathlineto{\pgfqpoint{12.168317in}{0.989054in}}%
\pgfpathlineto{\pgfqpoint{12.120693in}{0.989054in}}%
\pgfpathquadraticcurveto{\pgfqpoint{12.121377in}{0.978354in}}{\pgfqpoint{12.127141in}{0.972753in}}%
\pgfpathquadraticcurveto{\pgfqpoint{12.132905in}{0.967151in}}{\pgfqpoint{12.143208in}{0.967151in}}%
\pgfpathquadraticcurveto{\pgfqpoint{12.149170in}{0.967151in}}{\pgfqpoint{12.154772in}{0.968610in}}%
\pgfpathquadraticcurveto{\pgfqpoint{12.160374in}{0.970069in}}{\pgfqpoint{12.165903in}{0.973005in}}%
\pgfpathlineto{\pgfqpoint{12.165903in}{0.963206in}}%
\pgfpathquadraticcurveto{\pgfqpoint{12.160320in}{0.960847in}}{\pgfqpoint{12.154466in}{0.959604in}}%
\pgfpathquadraticcurveto{\pgfqpoint{12.148612in}{0.958361in}}{\pgfqpoint{12.142596in}{0.958361in}}%
\pgfpathquadraticcurveto{\pgfqpoint{12.127501in}{0.958361in}}{\pgfqpoint{12.118694in}{0.967133in}}%
\pgfpathquadraticcurveto{\pgfqpoint{12.109886in}{0.975923in}}{\pgfqpoint{12.109886in}{0.990909in}}%
\pgfpathquadraticcurveto{\pgfqpoint{12.109886in}{1.006381in}}{\pgfqpoint{12.118243in}{1.015459in}}%
\pgfpathquadraticcurveto{\pgfqpoint{12.126601in}{1.024556in}}{\pgfqpoint{12.140794in}{1.024556in}}%
\pgfpathquadraticcurveto{\pgfqpoint{12.153511in}{1.024556in}}{\pgfqpoint{12.160914in}{1.016360in}}%
\pgfpathquadraticcurveto{\pgfqpoint{12.168317in}{1.008183in}}{\pgfqpoint{12.168317in}{0.994115in}}%
\pgfpathlineto{\pgfqpoint{12.168317in}{0.994115in}}%
\pgfpathclose%
\pgfpathmoveto{\pgfqpoint{12.157960in}{0.997159in}}%
\pgfpathquadraticcurveto{\pgfqpoint{12.157852in}{1.005643in}}{\pgfqpoint{12.153205in}{1.010704in}}%
\pgfpathquadraticcurveto{\pgfqpoint{12.148558in}{1.015784in}}{\pgfqpoint{12.140903in}{1.015784in}}%
\pgfpathquadraticcurveto{\pgfqpoint{12.132239in}{1.015784in}}{\pgfqpoint{12.127033in}{1.010884in}}%
\pgfpathquadraticcurveto{\pgfqpoint{12.121828in}{1.005985in}}{\pgfqpoint{12.121035in}{0.997087in}}%
\pgfpathlineto{\pgfqpoint{12.157960in}{0.997159in}}%
\pgfpathlineto{\pgfqpoint{12.157960in}{0.997159in}}%
\pgfpathclose%
\pgfpathmoveto{\pgfqpoint{12.221849in}{1.013370in}}%
\pgfpathquadraticcurveto{\pgfqpoint{12.220102in}{1.014379in}}{\pgfqpoint{12.218049in}{1.014847in}}%
\pgfpathquadraticcurveto{\pgfqpoint{12.215995in}{1.015333in}}{\pgfqpoint{12.213528in}{1.015333in}}%
\pgfpathquadraticcurveto{\pgfqpoint{12.204738in}{1.015333in}}{\pgfqpoint{12.200036in}{1.009623in}}%
\pgfpathquadraticcurveto{\pgfqpoint{12.195335in}{1.003914in}}{\pgfqpoint{12.195335in}{0.993214in}}%
\pgfpathlineto{\pgfqpoint{12.195335in}{0.960000in}}%
\pgfpathlineto{\pgfqpoint{12.184924in}{0.960000in}}%
\pgfpathlineto{\pgfqpoint{12.184924in}{1.023043in}}%
\pgfpathlineto{\pgfqpoint{12.195335in}{1.023043in}}%
\pgfpathlineto{\pgfqpoint{12.195335in}{1.013244in}}%
\pgfpathquadraticcurveto{\pgfqpoint{12.198613in}{1.018990in}}{\pgfqpoint{12.203837in}{1.021764in}}%
\pgfpathquadraticcurveto{\pgfqpoint{12.209079in}{1.024556in}}{\pgfqpoint{12.216572in}{1.024556in}}%
\pgfpathquadraticcurveto{\pgfqpoint{12.217634in}{1.024556in}}{\pgfqpoint{12.218931in}{1.024411in}}%
\pgfpathquadraticcurveto{\pgfqpoint{12.220228in}{1.024285in}}{\pgfqpoint{12.221795in}{1.023997in}}%
\pgfpathlineto{\pgfqpoint{12.221849in}{1.013370in}}%
\pgfpathlineto{\pgfqpoint{12.221849in}{1.013370in}}%
\pgfpathclose%
\pgfusepath{fill}%
\end{pgfscope}%
\begin{pgfscope}%
\pgfsetbuttcap%
\pgfsetroundjoin%
\definecolor{currentfill}{rgb}{0.309804,0.372549,0.419608}%
\pgfsetfillcolor{currentfill}%
\pgfsetlinewidth{0.000000pt}%
\definecolor{currentstroke}{rgb}{0.309804,0.372549,0.419608}%
\pgfsetstrokecolor{currentstroke}%
\pgfsetdash{}{0pt}%
\pgfpathmoveto{\pgfqpoint{12.347003in}{0.994115in}}%
\pgfpathlineto{\pgfqpoint{12.347003in}{0.989054in}}%
\pgfpathlineto{\pgfqpoint{12.299379in}{0.989054in}}%
\pgfpathquadraticcurveto{\pgfqpoint{12.300063in}{0.978354in}}{\pgfqpoint{12.305827in}{0.972753in}}%
\pgfpathquadraticcurveto{\pgfqpoint{12.311591in}{0.967151in}}{\pgfqpoint{12.321894in}{0.967151in}}%
\pgfpathquadraticcurveto{\pgfqpoint{12.327856in}{0.967151in}}{\pgfqpoint{12.333458in}{0.968610in}}%
\pgfpathquadraticcurveto{\pgfqpoint{12.339060in}{0.970069in}}{\pgfqpoint{12.344589in}{0.973005in}}%
\pgfpathlineto{\pgfqpoint{12.344589in}{0.963206in}}%
\pgfpathquadraticcurveto{\pgfqpoint{12.339006in}{0.960847in}}{\pgfqpoint{12.333152in}{0.959604in}}%
\pgfpathquadraticcurveto{\pgfqpoint{12.327298in}{0.958361in}}{\pgfqpoint{12.321282in}{0.958361in}}%
\pgfpathquadraticcurveto{\pgfqpoint{12.306188in}{0.958361in}}{\pgfqpoint{12.297380in}{0.967133in}}%
\pgfpathquadraticcurveto{\pgfqpoint{12.288572in}{0.975923in}}{\pgfqpoint{12.288572in}{0.990909in}}%
\pgfpathquadraticcurveto{\pgfqpoint{12.288572in}{1.006381in}}{\pgfqpoint{12.296929in}{1.015459in}}%
\pgfpathquadraticcurveto{\pgfqpoint{12.305287in}{1.024556in}}{\pgfqpoint{12.319481in}{1.024556in}}%
\pgfpathquadraticcurveto{\pgfqpoint{12.332197in}{1.024556in}}{\pgfqpoint{12.339600in}{1.016360in}}%
\pgfpathquadraticcurveto{\pgfqpoint{12.347003in}{1.008183in}}{\pgfqpoint{12.347003in}{0.994115in}}%
\pgfpathlineto{\pgfqpoint{12.347003in}{0.994115in}}%
\pgfpathclose%
\pgfpathmoveto{\pgfqpoint{12.336646in}{0.997159in}}%
\pgfpathquadraticcurveto{\pgfqpoint{12.336538in}{1.005643in}}{\pgfqpoint{12.331891in}{1.010704in}}%
\pgfpathquadraticcurveto{\pgfqpoint{12.327244in}{1.015784in}}{\pgfqpoint{12.319589in}{1.015784in}}%
\pgfpathquadraticcurveto{\pgfqpoint{12.310925in}{1.015784in}}{\pgfqpoint{12.305719in}{1.010884in}}%
\pgfpathquadraticcurveto{\pgfqpoint{12.300514in}{1.005985in}}{\pgfqpoint{12.299721in}{0.997087in}}%
\pgfpathlineto{\pgfqpoint{12.336646in}{0.997159in}}%
\pgfpathlineto{\pgfqpoint{12.336646in}{0.997159in}}%
\pgfpathclose%
\pgfpathmoveto{\pgfqpoint{12.404192in}{1.021187in}}%
\pgfpathlineto{\pgfqpoint{12.404192in}{1.011389in}}%
\pgfpathquadraticcurveto{\pgfqpoint{12.399815in}{1.013640in}}{\pgfqpoint{12.395078in}{1.014757in}}%
\pgfpathquadraticcurveto{\pgfqpoint{12.390358in}{1.015892in}}{\pgfqpoint{12.385279in}{1.015892in}}%
\pgfpathquadraticcurveto{\pgfqpoint{12.377570in}{1.015892in}}{\pgfqpoint{12.373715in}{1.013532in}}%
\pgfpathquadraticcurveto{\pgfqpoint{12.369860in}{1.011173in}}{\pgfqpoint{12.369860in}{1.006435in}}%
\pgfpathquadraticcurveto{\pgfqpoint{12.369860in}{1.002833in}}{\pgfqpoint{12.372616in}{1.000780in}}%
\pgfpathquadraticcurveto{\pgfqpoint{12.375372in}{0.998726in}}{\pgfqpoint{12.383712in}{0.996871in}}%
\pgfpathlineto{\pgfqpoint{12.387260in}{0.996078in}}%
\pgfpathquadraticcurveto{\pgfqpoint{12.398284in}{0.993719in}}{\pgfqpoint{12.402931in}{0.989414in}}%
\pgfpathquadraticcurveto{\pgfqpoint{12.407578in}{0.985109in}}{\pgfqpoint{12.407578in}{0.977400in}}%
\pgfpathquadraticcurveto{\pgfqpoint{12.407578in}{0.968610in}}{\pgfqpoint{12.400625in}{0.963476in}}%
\pgfpathquadraticcurveto{\pgfqpoint{12.393673in}{0.958361in}}{\pgfqpoint{12.381514in}{0.958361in}}%
\pgfpathquadraticcurveto{\pgfqpoint{12.376453in}{0.958361in}}{\pgfqpoint{12.370959in}{0.959352in}}%
\pgfpathquadraticcurveto{\pgfqpoint{12.365466in}{0.960342in}}{\pgfqpoint{12.359395in}{0.962306in}}%
\pgfpathlineto{\pgfqpoint{12.359395in}{0.973005in}}%
\pgfpathquadraticcurveto{\pgfqpoint{12.365141in}{0.970015in}}{\pgfqpoint{12.370707in}{0.968520in}}%
\pgfpathquadraticcurveto{\pgfqpoint{12.376273in}{0.967043in}}{\pgfqpoint{12.381749in}{0.967043in}}%
\pgfpathquadraticcurveto{\pgfqpoint{12.389061in}{0.967043in}}{\pgfqpoint{12.392988in}{0.969546in}}%
\pgfpathquadraticcurveto{\pgfqpoint{12.396933in}{0.972050in}}{\pgfqpoint{12.396933in}{0.976607in}}%
\pgfpathquadraticcurveto{\pgfqpoint{12.396933in}{0.980822in}}{\pgfqpoint{12.394087in}{0.983074in}}%
\pgfpathquadraticcurveto{\pgfqpoint{12.391259in}{0.985325in}}{\pgfqpoint{12.381622in}{0.987414in}}%
\pgfpathlineto{\pgfqpoint{12.378020in}{0.988261in}}%
\pgfpathquadraticcurveto{\pgfqpoint{12.368402in}{0.990278in}}{\pgfqpoint{12.364115in}{0.994475in}}%
\pgfpathquadraticcurveto{\pgfqpoint{12.359846in}{0.998672in}}{\pgfqpoint{12.359846in}{1.005985in}}%
\pgfpathquadraticcurveto{\pgfqpoint{12.359846in}{1.014883in}}{\pgfqpoint{12.366150in}{1.019710in}}%
\pgfpathquadraticcurveto{\pgfqpoint{12.372454in}{1.024556in}}{\pgfqpoint{12.384054in}{1.024556in}}%
\pgfpathquadraticcurveto{\pgfqpoint{12.389782in}{1.024556in}}{\pgfqpoint{12.394843in}{1.023709in}}%
\pgfpathquadraticcurveto{\pgfqpoint{12.399923in}{1.022880in}}{\pgfqpoint{12.404192in}{1.021187in}}%
\pgfpathlineto{\pgfqpoint{12.404192in}{1.021187in}}%
\pgfpathclose%
\pgfpathmoveto{\pgfqpoint{12.434308in}{1.040947in}}%
\pgfpathlineto{\pgfqpoint{12.434308in}{1.023043in}}%
\pgfpathlineto{\pgfqpoint{12.455634in}{1.023043in}}%
\pgfpathlineto{\pgfqpoint{12.455634in}{1.014991in}}%
\pgfpathlineto{\pgfqpoint{12.434308in}{1.014991in}}%
\pgfpathlineto{\pgfqpoint{12.434308in}{0.980768in}}%
\pgfpathquadraticcurveto{\pgfqpoint{12.434308in}{0.973059in}}{\pgfqpoint{12.436415in}{0.970861in}}%
\pgfpathquadraticcurveto{\pgfqpoint{12.438523in}{0.968664in}}{\pgfqpoint{12.445007in}{0.968664in}}%
\pgfpathlineto{\pgfqpoint{12.455634in}{0.968664in}}%
\pgfpathlineto{\pgfqpoint{12.455634in}{0.960000in}}%
\pgfpathlineto{\pgfqpoint{12.445007in}{0.960000in}}%
\pgfpathquadraticcurveto{\pgfqpoint{12.433011in}{0.960000in}}{\pgfqpoint{12.428454in}{0.964467in}}%
\pgfpathquadraticcurveto{\pgfqpoint{12.423897in}{0.968952in}}{\pgfqpoint{12.423897in}{0.980768in}}%
\pgfpathlineto{\pgfqpoint{12.423897in}{1.014991in}}%
\pgfpathlineto{\pgfqpoint{12.416296in}{1.014991in}}%
\pgfpathlineto{\pgfqpoint{12.416296in}{1.023043in}}%
\pgfpathlineto{\pgfqpoint{12.423897in}{1.023043in}}%
\pgfpathlineto{\pgfqpoint{12.423897in}{1.040947in}}%
\pgfpathlineto{\pgfqpoint{12.434308in}{1.040947in}}%
\pgfpathlineto{\pgfqpoint{12.434308in}{1.040947in}}%
\pgfpathclose%
\pgfpathmoveto{\pgfqpoint{12.469252in}{1.023043in}}%
\pgfpathlineto{\pgfqpoint{12.479609in}{1.023043in}}%
\pgfpathlineto{\pgfqpoint{12.479609in}{0.960000in}}%
\pgfpathlineto{\pgfqpoint{12.469252in}{0.960000in}}%
\pgfpathlineto{\pgfqpoint{12.469252in}{1.023043in}}%
\pgfpathlineto{\pgfqpoint{12.469252in}{1.023043in}}%
\pgfpathclose%
\pgfpathmoveto{\pgfqpoint{12.469252in}{1.047593in}}%
\pgfpathlineto{\pgfqpoint{12.479609in}{1.047593in}}%
\pgfpathlineto{\pgfqpoint{12.479609in}{1.034462in}}%
\pgfpathlineto{\pgfqpoint{12.469252in}{1.034462in}}%
\pgfpathlineto{\pgfqpoint{12.469252in}{1.047593in}}%
\pgfpathlineto{\pgfqpoint{12.469252in}{1.047593in}}%
\pgfpathclose%
\pgfpathmoveto{\pgfqpoint{12.550360in}{1.010938in}}%
\pgfpathquadraticcurveto{\pgfqpoint{12.554251in}{1.017927in}}{\pgfqpoint{12.559655in}{1.021241in}}%
\pgfpathquadraticcurveto{\pgfqpoint{12.565058in}{1.024556in}}{\pgfqpoint{12.572371in}{1.024556in}}%
\pgfpathquadraticcurveto{\pgfqpoint{12.582224in}{1.024556in}}{\pgfqpoint{12.587573in}{1.017657in}}%
\pgfpathquadraticcurveto{\pgfqpoint{12.592923in}{1.010776in}}{\pgfqpoint{12.592923in}{0.998060in}}%
\pgfpathlineto{\pgfqpoint{12.592923in}{0.960000in}}%
\pgfpathlineto{\pgfqpoint{12.582512in}{0.960000in}}%
\pgfpathlineto{\pgfqpoint{12.582512in}{0.997717in}}%
\pgfpathquadraticcurveto{\pgfqpoint{12.582512in}{1.006778in}}{\pgfqpoint{12.579288in}{1.011155in}}%
\pgfpathquadraticcurveto{\pgfqpoint{12.576082in}{1.015549in}}{\pgfqpoint{12.569507in}{1.015549in}}%
\pgfpathquadraticcurveto{\pgfqpoint{12.561456in}{1.015549in}}{\pgfqpoint{12.556773in}{1.010200in}}%
\pgfpathquadraticcurveto{\pgfqpoint{12.552107in}{1.004868in}}{\pgfqpoint{12.552107in}{0.995628in}}%
\pgfpathlineto{\pgfqpoint{12.552107in}{0.960000in}}%
\pgfpathlineto{\pgfqpoint{12.541696in}{0.960000in}}%
\pgfpathlineto{\pgfqpoint{12.541696in}{0.997717in}}%
\pgfpathquadraticcurveto{\pgfqpoint{12.541696in}{1.006832in}}{\pgfqpoint{12.538490in}{1.011191in}}%
\pgfpathquadraticcurveto{\pgfqpoint{12.535284in}{1.015549in}}{\pgfqpoint{12.528584in}{1.015549in}}%
\pgfpathquadraticcurveto{\pgfqpoint{12.520640in}{1.015549in}}{\pgfqpoint{12.515957in}{1.010182in}}%
\pgfpathquadraticcurveto{\pgfqpoint{12.511292in}{1.004814in}}{\pgfqpoint{12.511292in}{0.995628in}}%
\pgfpathlineto{\pgfqpoint{12.511292in}{0.960000in}}%
\pgfpathlineto{\pgfqpoint{12.500881in}{0.960000in}}%
\pgfpathlineto{\pgfqpoint{12.500881in}{1.023043in}}%
\pgfpathlineto{\pgfqpoint{12.511292in}{1.023043in}}%
\pgfpathlineto{\pgfqpoint{12.511292in}{1.013244in}}%
\pgfpathquadraticcurveto{\pgfqpoint{12.514840in}{1.019044in}}{\pgfqpoint{12.519794in}{1.021800in}}%
\pgfpathquadraticcurveto{\pgfqpoint{12.524747in}{1.024556in}}{\pgfqpoint{12.531556in}{1.024556in}}%
\pgfpathquadraticcurveto{\pgfqpoint{12.538436in}{1.024556in}}{\pgfqpoint{12.543245in}{1.021061in}}%
\pgfpathquadraticcurveto{\pgfqpoint{12.548055in}{1.017585in}}{\pgfqpoint{12.550360in}{1.010938in}}%
\pgfpathlineto{\pgfqpoint{12.550360in}{1.010938in}}%
\pgfpathclose%
\pgfpathmoveto{\pgfqpoint{12.642222in}{0.991683in}}%
\pgfpathquadraticcurveto{\pgfqpoint{12.629668in}{0.991683in}}{\pgfqpoint{12.624823in}{0.988819in}}%
\pgfpathquadraticcurveto{\pgfqpoint{12.619977in}{0.985956in}}{\pgfqpoint{12.619977in}{0.979021in}}%
\pgfpathquadraticcurveto{\pgfqpoint{12.619977in}{0.973509in}}{\pgfqpoint{12.623616in}{0.970267in}}%
\pgfpathquadraticcurveto{\pgfqpoint{12.627254in}{0.967043in}}{\pgfqpoint{12.633486in}{0.967043in}}%
\pgfpathquadraticcurveto{\pgfqpoint{12.642114in}{0.967043in}}{\pgfqpoint{12.647320in}{0.973149in}}%
\pgfpathquadraticcurveto{\pgfqpoint{12.652525in}{0.979255in}}{\pgfqpoint{12.652525in}{0.989378in}}%
\pgfpathlineto{\pgfqpoint{12.652525in}{0.991683in}}%
\pgfpathlineto{\pgfqpoint{12.642222in}{0.991683in}}%
\pgfpathlineto{\pgfqpoint{12.642222in}{0.991683in}}%
\pgfpathclose%
\pgfpathmoveto{\pgfqpoint{12.662882in}{0.995970in}}%
\pgfpathlineto{\pgfqpoint{12.662882in}{0.960000in}}%
\pgfpathlineto{\pgfqpoint{12.652525in}{0.960000in}}%
\pgfpathlineto{\pgfqpoint{12.652525in}{0.969564in}}%
\pgfpathquadraticcurveto{\pgfqpoint{12.648977in}{0.963837in}}{\pgfqpoint{12.643681in}{0.961099in}}%
\pgfpathquadraticcurveto{\pgfqpoint{12.638386in}{0.958361in}}{\pgfqpoint{12.630731in}{0.958361in}}%
\pgfpathquadraticcurveto{\pgfqpoint{12.621058in}{0.958361in}}{\pgfqpoint{12.615330in}{0.963801in}}%
\pgfpathquadraticcurveto{\pgfqpoint{12.609620in}{0.969240in}}{\pgfqpoint{12.609620in}{0.978354in}}%
\pgfpathquadraticcurveto{\pgfqpoint{12.609620in}{0.988982in}}{\pgfqpoint{12.616735in}{0.994385in}}%
\pgfpathquadraticcurveto{\pgfqpoint{12.623868in}{0.999789in}}{\pgfqpoint{12.637989in}{0.999789in}}%
\pgfpathlineto{\pgfqpoint{12.652525in}{0.999789in}}%
\pgfpathlineto{\pgfqpoint{12.652525in}{1.000816in}}%
\pgfpathquadraticcurveto{\pgfqpoint{12.652525in}{1.007966in}}{\pgfqpoint{12.647824in}{1.011875in}}%
\pgfpathquadraticcurveto{\pgfqpoint{12.643123in}{1.015784in}}{\pgfqpoint{12.634621in}{1.015784in}}%
\pgfpathquadraticcurveto{\pgfqpoint{12.629217in}{1.015784in}}{\pgfqpoint{12.624084in}{1.014487in}}%
\pgfpathquadraticcurveto{\pgfqpoint{12.618969in}{1.013190in}}{\pgfqpoint{12.614249in}{1.010596in}}%
\pgfpathlineto{\pgfqpoint{12.614249in}{1.020179in}}%
\pgfpathquadraticcurveto{\pgfqpoint{12.619923in}{1.022376in}}{\pgfqpoint{12.625273in}{1.023457in}}%
\pgfpathquadraticcurveto{\pgfqpoint{12.630622in}{1.024556in}}{\pgfqpoint{12.635684in}{1.024556in}}%
\pgfpathquadraticcurveto{\pgfqpoint{12.649373in}{1.024556in}}{\pgfqpoint{12.656128in}{1.017459in}}%
\pgfpathquadraticcurveto{\pgfqpoint{12.662882in}{1.010380in}}{\pgfqpoint{12.662882in}{0.995970in}}%
\pgfpathlineto{\pgfqpoint{12.662882in}{0.995970in}}%
\pgfpathclose%
\pgfpathmoveto{\pgfqpoint{12.694458in}{1.040947in}}%
\pgfpathlineto{\pgfqpoint{12.694458in}{1.023043in}}%
\pgfpathlineto{\pgfqpoint{12.715784in}{1.023043in}}%
\pgfpathlineto{\pgfqpoint{12.715784in}{1.014991in}}%
\pgfpathlineto{\pgfqpoint{12.694458in}{1.014991in}}%
\pgfpathlineto{\pgfqpoint{12.694458in}{0.980768in}}%
\pgfpathquadraticcurveto{\pgfqpoint{12.694458in}{0.973059in}}{\pgfqpoint{12.696565in}{0.970861in}}%
\pgfpathquadraticcurveto{\pgfqpoint{12.698672in}{0.968664in}}{\pgfqpoint{12.705157in}{0.968664in}}%
\pgfpathlineto{\pgfqpoint{12.715784in}{0.968664in}}%
\pgfpathlineto{\pgfqpoint{12.715784in}{0.960000in}}%
\pgfpathlineto{\pgfqpoint{12.705157in}{0.960000in}}%
\pgfpathquadraticcurveto{\pgfqpoint{12.693161in}{0.960000in}}{\pgfqpoint{12.688604in}{0.964467in}}%
\pgfpathquadraticcurveto{\pgfqpoint{12.684046in}{0.968952in}}{\pgfqpoint{12.684046in}{0.980768in}}%
\pgfpathlineto{\pgfqpoint{12.684046in}{1.014991in}}%
\pgfpathlineto{\pgfqpoint{12.676445in}{1.014991in}}%
\pgfpathlineto{\pgfqpoint{12.676445in}{1.023043in}}%
\pgfpathlineto{\pgfqpoint{12.684046in}{1.023043in}}%
\pgfpathlineto{\pgfqpoint{12.684046in}{1.040947in}}%
\pgfpathlineto{\pgfqpoint{12.694458in}{1.040947in}}%
\pgfpathlineto{\pgfqpoint{12.694458in}{1.040947in}}%
\pgfpathclose%
\pgfpathmoveto{\pgfqpoint{12.783329in}{0.994115in}}%
\pgfpathlineto{\pgfqpoint{12.783329in}{0.989054in}}%
\pgfpathlineto{\pgfqpoint{12.735705in}{0.989054in}}%
\pgfpathquadraticcurveto{\pgfqpoint{12.736390in}{0.978354in}}{\pgfqpoint{12.742154in}{0.972753in}}%
\pgfpathquadraticcurveto{\pgfqpoint{12.747918in}{0.967151in}}{\pgfqpoint{12.758221in}{0.967151in}}%
\pgfpathquadraticcurveto{\pgfqpoint{12.764183in}{0.967151in}}{\pgfqpoint{12.769784in}{0.968610in}}%
\pgfpathquadraticcurveto{\pgfqpoint{12.775386in}{0.970069in}}{\pgfqpoint{12.780916in}{0.973005in}}%
\pgfpathlineto{\pgfqpoint{12.780916in}{0.963206in}}%
\pgfpathquadraticcurveto{\pgfqpoint{12.775332in}{0.960847in}}{\pgfqpoint{12.769478in}{0.959604in}}%
\pgfpathquadraticcurveto{\pgfqpoint{12.763624in}{0.958361in}}{\pgfqpoint{12.757608in}{0.958361in}}%
\pgfpathquadraticcurveto{\pgfqpoint{12.742514in}{0.958361in}}{\pgfqpoint{12.733706in}{0.967133in}}%
\pgfpathquadraticcurveto{\pgfqpoint{12.724898in}{0.975923in}}{\pgfqpoint{12.724898in}{0.990909in}}%
\pgfpathquadraticcurveto{\pgfqpoint{12.724898in}{1.006381in}}{\pgfqpoint{12.733256in}{1.015459in}}%
\pgfpathquadraticcurveto{\pgfqpoint{12.741613in}{1.024556in}}{\pgfqpoint{12.755807in}{1.024556in}}%
\pgfpathquadraticcurveto{\pgfqpoint{12.768523in}{1.024556in}}{\pgfqpoint{12.775926in}{1.016360in}}%
\pgfpathquadraticcurveto{\pgfqpoint{12.783329in}{1.008183in}}{\pgfqpoint{12.783329in}{0.994115in}}%
\pgfpathlineto{\pgfqpoint{12.783329in}{0.994115in}}%
\pgfpathclose%
\pgfpathmoveto{\pgfqpoint{12.772972in}{0.997159in}}%
\pgfpathquadraticcurveto{\pgfqpoint{12.772864in}{1.005643in}}{\pgfqpoint{12.768217in}{1.010704in}}%
\pgfpathquadraticcurveto{\pgfqpoint{12.763570in}{1.015784in}}{\pgfqpoint{12.755915in}{1.015784in}}%
\pgfpathquadraticcurveto{\pgfqpoint{12.747251in}{1.015784in}}{\pgfqpoint{12.742046in}{1.010884in}}%
\pgfpathquadraticcurveto{\pgfqpoint{12.736840in}{1.005985in}}{\pgfqpoint{12.736048in}{0.997087in}}%
\pgfpathlineto{\pgfqpoint{12.772972in}{0.997159in}}%
\pgfpathlineto{\pgfqpoint{12.772972in}{0.997159in}}%
\pgfpathclose%
\pgfusepath{fill}%
\end{pgfscope}%
\begin{pgfscope}%
\pgfsetbuttcap%
\pgfsetroundjoin%
\definecolor{currentfill}{rgb}{0.309804,0.372549,0.419608}%
\pgfsetfillcolor{currentfill}%
\pgfsetlinewidth{0.000000pt}%
\definecolor{currentstroke}{rgb}{0.309804,0.372549,0.419608}%
\pgfsetstrokecolor{currentstroke}%
\pgfsetdash{}{0pt}%
\pgfpathmoveto{\pgfqpoint{12.915383in}{0.998060in}}%
\pgfpathlineto{\pgfqpoint{12.915383in}{0.960000in}}%
\pgfpathlineto{\pgfqpoint{12.905026in}{0.960000in}}%
\pgfpathlineto{\pgfqpoint{12.905026in}{0.997717in}}%
\pgfpathquadraticcurveto{\pgfqpoint{12.905026in}{1.006669in}}{\pgfqpoint{12.901532in}{1.011100in}}%
\pgfpathquadraticcurveto{\pgfqpoint{12.898038in}{1.015549in}}{\pgfqpoint{12.891067in}{1.015549in}}%
\pgfpathquadraticcurveto{\pgfqpoint{12.882673in}{1.015549in}}{\pgfqpoint{12.877828in}{1.010200in}}%
\pgfpathquadraticcurveto{\pgfqpoint{12.872983in}{1.004868in}}{\pgfqpoint{12.872983in}{0.995628in}}%
\pgfpathlineto{\pgfqpoint{12.872983in}{0.960000in}}%
\pgfpathlineto{\pgfqpoint{12.862572in}{0.960000in}}%
\pgfpathlineto{\pgfqpoint{12.862572in}{1.047593in}}%
\pgfpathlineto{\pgfqpoint{12.872983in}{1.047593in}}%
\pgfpathlineto{\pgfqpoint{12.872983in}{1.013244in}}%
\pgfpathquadraticcurveto{\pgfqpoint{12.876711in}{1.018936in}}{\pgfqpoint{12.881737in}{1.021746in}}%
\pgfpathquadraticcurveto{\pgfqpoint{12.886780in}{1.024556in}}{\pgfqpoint{12.893373in}{1.024556in}}%
\pgfpathquadraticcurveto{\pgfqpoint{12.904234in}{1.024556in}}{\pgfqpoint{12.909800in}{1.017837in}}%
\pgfpathquadraticcurveto{\pgfqpoint{12.915383in}{1.011118in}}{\pgfqpoint{12.915383in}{0.998060in}}%
\pgfpathlineto{\pgfqpoint{12.915383in}{0.998060in}}%
\pgfpathclose%
\pgfpathmoveto{\pgfqpoint{12.964683in}{0.991683in}}%
\pgfpathquadraticcurveto{\pgfqpoint{12.952128in}{0.991683in}}{\pgfqpoint{12.947283in}{0.988819in}}%
\pgfpathquadraticcurveto{\pgfqpoint{12.942438in}{0.985956in}}{\pgfqpoint{12.942438in}{0.979021in}}%
\pgfpathquadraticcurveto{\pgfqpoint{12.942438in}{0.973509in}}{\pgfqpoint{12.946076in}{0.970267in}}%
\pgfpathquadraticcurveto{\pgfqpoint{12.949715in}{0.967043in}}{\pgfqpoint{12.955947in}{0.967043in}}%
\pgfpathquadraticcurveto{\pgfqpoint{12.964575in}{0.967043in}}{\pgfqpoint{12.969780in}{0.973149in}}%
\pgfpathquadraticcurveto{\pgfqpoint{12.974986in}{0.979255in}}{\pgfqpoint{12.974986in}{0.989378in}}%
\pgfpathlineto{\pgfqpoint{12.974986in}{0.991683in}}%
\pgfpathlineto{\pgfqpoint{12.964683in}{0.991683in}}%
\pgfpathlineto{\pgfqpoint{12.964683in}{0.991683in}}%
\pgfpathclose%
\pgfpathmoveto{\pgfqpoint{12.985343in}{0.995970in}}%
\pgfpathlineto{\pgfqpoint{12.985343in}{0.960000in}}%
\pgfpathlineto{\pgfqpoint{12.974986in}{0.960000in}}%
\pgfpathlineto{\pgfqpoint{12.974986in}{0.969564in}}%
\pgfpathquadraticcurveto{\pgfqpoint{12.971437in}{0.963837in}}{\pgfqpoint{12.966142in}{0.961099in}}%
\pgfpathquadraticcurveto{\pgfqpoint{12.960846in}{0.958361in}}{\pgfqpoint{12.953191in}{0.958361in}}%
\pgfpathquadraticcurveto{\pgfqpoint{12.943518in}{0.958361in}}{\pgfqpoint{12.937790in}{0.963801in}}%
\pgfpathquadraticcurveto{\pgfqpoint{12.932081in}{0.969240in}}{\pgfqpoint{12.932081in}{0.978354in}}%
\pgfpathquadraticcurveto{\pgfqpoint{12.932081in}{0.988982in}}{\pgfqpoint{12.939195in}{0.994385in}}%
\pgfpathquadraticcurveto{\pgfqpoint{12.946328in}{0.999789in}}{\pgfqpoint{12.960450in}{0.999789in}}%
\pgfpathlineto{\pgfqpoint{12.974986in}{0.999789in}}%
\pgfpathlineto{\pgfqpoint{12.974986in}{1.000816in}}%
\pgfpathquadraticcurveto{\pgfqpoint{12.974986in}{1.007966in}}{\pgfqpoint{12.970284in}{1.011875in}}%
\pgfpathquadraticcurveto{\pgfqpoint{12.965583in}{1.015784in}}{\pgfqpoint{12.957082in}{1.015784in}}%
\pgfpathquadraticcurveto{\pgfqpoint{12.951678in}{1.015784in}}{\pgfqpoint{12.946544in}{1.014487in}}%
\pgfpathquadraticcurveto{\pgfqpoint{12.941429in}{1.013190in}}{\pgfqpoint{12.936710in}{1.010596in}}%
\pgfpathlineto{\pgfqpoint{12.936710in}{1.020179in}}%
\pgfpathquadraticcurveto{\pgfqpoint{12.942384in}{1.022376in}}{\pgfqpoint{12.947733in}{1.023457in}}%
\pgfpathquadraticcurveto{\pgfqpoint{12.953083in}{1.024556in}}{\pgfqpoint{12.958144in}{1.024556in}}%
\pgfpathquadraticcurveto{\pgfqpoint{12.971833in}{1.024556in}}{\pgfqpoint{12.978588in}{1.017459in}}%
\pgfpathquadraticcurveto{\pgfqpoint{12.985343in}{1.010380in}}{\pgfqpoint{12.985343in}{0.995970in}}%
\pgfpathlineto{\pgfqpoint{12.985343in}{0.995970in}}%
\pgfpathclose%
\pgfpathmoveto{\pgfqpoint{13.046854in}{1.021187in}}%
\pgfpathlineto{\pgfqpoint{13.046854in}{1.011389in}}%
\pgfpathquadraticcurveto{\pgfqpoint{13.042477in}{1.013640in}}{\pgfqpoint{13.037740in}{1.014757in}}%
\pgfpathquadraticcurveto{\pgfqpoint{13.033021in}{1.015892in}}{\pgfqpoint{13.027941in}{1.015892in}}%
\pgfpathquadraticcurveto{\pgfqpoint{13.020232in}{1.015892in}}{\pgfqpoint{13.016378in}{1.013532in}}%
\pgfpathquadraticcurveto{\pgfqpoint{13.012523in}{1.011173in}}{\pgfqpoint{13.012523in}{1.006435in}}%
\pgfpathquadraticcurveto{\pgfqpoint{13.012523in}{1.002833in}}{\pgfqpoint{13.015279in}{1.000780in}}%
\pgfpathquadraticcurveto{\pgfqpoint{13.018035in}{0.998726in}}{\pgfqpoint{13.026374in}{0.996871in}}%
\pgfpathlineto{\pgfqpoint{13.029923in}{0.996078in}}%
\pgfpathquadraticcurveto{\pgfqpoint{13.040946in}{0.993719in}}{\pgfqpoint{13.045593in}{0.989414in}}%
\pgfpathquadraticcurveto{\pgfqpoint{13.050240in}{0.985109in}}{\pgfqpoint{13.050240in}{0.977400in}}%
\pgfpathquadraticcurveto{\pgfqpoint{13.050240in}{0.968610in}}{\pgfqpoint{13.043288in}{0.963476in}}%
\pgfpathquadraticcurveto{\pgfqpoint{13.036335in}{0.958361in}}{\pgfqpoint{13.024177in}{0.958361in}}%
\pgfpathquadraticcurveto{\pgfqpoint{13.019115in}{0.958361in}}{\pgfqpoint{13.013622in}{0.959352in}}%
\pgfpathquadraticcurveto{\pgfqpoint{13.008128in}{0.960342in}}{\pgfqpoint{13.002058in}{0.962306in}}%
\pgfpathlineto{\pgfqpoint{13.002058in}{0.973005in}}%
\pgfpathquadraticcurveto{\pgfqpoint{13.007804in}{0.970015in}}{\pgfqpoint{13.013369in}{0.968520in}}%
\pgfpathquadraticcurveto{\pgfqpoint{13.018935in}{0.967043in}}{\pgfqpoint{13.024411in}{0.967043in}}%
\pgfpathquadraticcurveto{\pgfqpoint{13.031724in}{0.967043in}}{\pgfqpoint{13.035651in}{0.969546in}}%
\pgfpathquadraticcurveto{\pgfqpoint{13.039595in}{0.972050in}}{\pgfqpoint{13.039595in}{0.976607in}}%
\pgfpathquadraticcurveto{\pgfqpoint{13.039595in}{0.980822in}}{\pgfqpoint{13.036749in}{0.983074in}}%
\pgfpathquadraticcurveto{\pgfqpoint{13.033921in}{0.985325in}}{\pgfqpoint{13.024285in}{0.987414in}}%
\pgfpathlineto{\pgfqpoint{13.020682in}{0.988261in}}%
\pgfpathquadraticcurveto{\pgfqpoint{13.011064in}{0.990278in}}{\pgfqpoint{13.006777in}{0.994475in}}%
\pgfpathquadraticcurveto{\pgfqpoint{13.002508in}{0.998672in}}{\pgfqpoint{13.002508in}{1.005985in}}%
\pgfpathquadraticcurveto{\pgfqpoint{13.002508in}{1.014883in}}{\pgfqpoint{13.008812in}{1.019710in}}%
\pgfpathquadraticcurveto{\pgfqpoint{13.015117in}{1.024556in}}{\pgfqpoint{13.026716in}{1.024556in}}%
\pgfpathquadraticcurveto{\pgfqpoint{13.032444in}{1.024556in}}{\pgfqpoint{13.037506in}{1.023709in}}%
\pgfpathquadraticcurveto{\pgfqpoint{13.042585in}{1.022880in}}{\pgfqpoint{13.046854in}{1.021187in}}%
\pgfpathlineto{\pgfqpoint{13.046854in}{1.021187in}}%
\pgfpathclose%
\pgfusepath{fill}%
\end{pgfscope}%
\begin{pgfscope}%
\pgfsetbuttcap%
\pgfsetroundjoin%
\definecolor{currentfill}{rgb}{0.309804,0.372549,0.419608}%
\pgfsetfillcolor{currentfill}%
\pgfsetlinewidth{0.000000pt}%
\definecolor{currentstroke}{rgb}{0.309804,0.372549,0.419608}%
\pgfsetstrokecolor{currentstroke}%
\pgfsetdash{}{0pt}%
\pgfpathmoveto{\pgfqpoint{13.176790in}{0.994115in}}%
\pgfpathlineto{\pgfqpoint{13.176790in}{0.989054in}}%
\pgfpathlineto{\pgfqpoint{13.129166in}{0.989054in}}%
\pgfpathquadraticcurveto{\pgfqpoint{13.129850in}{0.978354in}}{\pgfqpoint{13.135614in}{0.972753in}}%
\pgfpathquadraticcurveto{\pgfqpoint{13.141378in}{0.967151in}}{\pgfqpoint{13.151681in}{0.967151in}}%
\pgfpathquadraticcurveto{\pgfqpoint{13.157643in}{0.967151in}}{\pgfqpoint{13.163245in}{0.968610in}}%
\pgfpathquadraticcurveto{\pgfqpoint{13.168846in}{0.970069in}}{\pgfqpoint{13.174376in}{0.973005in}}%
\pgfpathlineto{\pgfqpoint{13.174376in}{0.963206in}}%
\pgfpathquadraticcurveto{\pgfqpoint{13.168792in}{0.960847in}}{\pgfqpoint{13.162938in}{0.959604in}}%
\pgfpathquadraticcurveto{\pgfqpoint{13.157084in}{0.958361in}}{\pgfqpoint{13.151068in}{0.958361in}}%
\pgfpathquadraticcurveto{\pgfqpoint{13.135974in}{0.958361in}}{\pgfqpoint{13.127166in}{0.967133in}}%
\pgfpathquadraticcurveto{\pgfqpoint{13.118358in}{0.975923in}}{\pgfqpoint{13.118358in}{0.990909in}}%
\pgfpathquadraticcurveto{\pgfqpoint{13.118358in}{1.006381in}}{\pgfqpoint{13.126716in}{1.015459in}}%
\pgfpathquadraticcurveto{\pgfqpoint{13.135074in}{1.024556in}}{\pgfqpoint{13.149267in}{1.024556in}}%
\pgfpathquadraticcurveto{\pgfqpoint{13.161984in}{1.024556in}}{\pgfqpoint{13.169387in}{1.016360in}}%
\pgfpathquadraticcurveto{\pgfqpoint{13.176790in}{1.008183in}}{\pgfqpoint{13.176790in}{0.994115in}}%
\pgfpathlineto{\pgfqpoint{13.176790in}{0.994115in}}%
\pgfpathclose%
\pgfpathmoveto{\pgfqpoint{13.166433in}{0.997159in}}%
\pgfpathquadraticcurveto{\pgfqpoint{13.166325in}{1.005643in}}{\pgfqpoint{13.161678in}{1.010704in}}%
\pgfpathquadraticcurveto{\pgfqpoint{13.157030in}{1.015784in}}{\pgfqpoint{13.149375in}{1.015784in}}%
\pgfpathquadraticcurveto{\pgfqpoint{13.140711in}{1.015784in}}{\pgfqpoint{13.135506in}{1.010884in}}%
\pgfpathquadraticcurveto{\pgfqpoint{13.130300in}{1.005985in}}{\pgfqpoint{13.129508in}{0.997087in}}%
\pgfpathlineto{\pgfqpoint{13.166433in}{0.997159in}}%
\pgfpathlineto{\pgfqpoint{13.166433in}{0.997159in}}%
\pgfpathclose%
\pgfpathmoveto{\pgfqpoint{13.264545in}{1.047593in}}%
\pgfpathlineto{\pgfqpoint{13.264545in}{1.038965in}}%
\pgfpathlineto{\pgfqpoint{13.254638in}{1.038965in}}%
\pgfpathquadraticcurveto{\pgfqpoint{13.249072in}{1.038965in}}{\pgfqpoint{13.246893in}{1.036714in}}%
\pgfpathquadraticcurveto{\pgfqpoint{13.244732in}{1.034462in}}{\pgfqpoint{13.244732in}{1.028608in}}%
\pgfpathlineto{\pgfqpoint{13.244732in}{1.023043in}}%
\pgfpathlineto{\pgfqpoint{13.261789in}{1.023043in}}%
\pgfpathlineto{\pgfqpoint{13.261789in}{1.014991in}}%
\pgfpathlineto{\pgfqpoint{13.244732in}{1.014991in}}%
\pgfpathlineto{\pgfqpoint{13.244732in}{0.960000in}}%
\pgfpathlineto{\pgfqpoint{13.234321in}{0.960000in}}%
\pgfpathlineto{\pgfqpoint{13.234321in}{1.014991in}}%
\pgfpathlineto{\pgfqpoint{13.205897in}{1.014991in}}%
\pgfpathlineto{\pgfqpoint{13.205897in}{0.960000in}}%
\pgfpathlineto{\pgfqpoint{13.195486in}{0.960000in}}%
\pgfpathlineto{\pgfqpoint{13.195486in}{1.014991in}}%
\pgfpathlineto{\pgfqpoint{13.185580in}{1.014991in}}%
\pgfpathlineto{\pgfqpoint{13.185580in}{1.023043in}}%
\pgfpathlineto{\pgfqpoint{13.195486in}{1.023043in}}%
\pgfpathlineto{\pgfqpoint{13.195486in}{1.027438in}}%
\pgfpathquadraticcurveto{\pgfqpoint{13.195486in}{1.037957in}}{\pgfqpoint{13.200386in}{1.042766in}}%
\pgfpathquadraticcurveto{\pgfqpoint{13.205285in}{1.047593in}}{\pgfqpoint{13.215912in}{1.047593in}}%
\pgfpathlineto{\pgfqpoint{13.225711in}{1.047593in}}%
\pgfpathlineto{\pgfqpoint{13.225711in}{1.038965in}}%
\pgfpathlineto{\pgfqpoint{13.215804in}{1.038965in}}%
\pgfpathquadraticcurveto{\pgfqpoint{13.210238in}{1.038965in}}{\pgfqpoint{13.208041in}{1.036714in}}%
\pgfpathquadraticcurveto{\pgfqpoint{13.205897in}{1.034462in}}{\pgfqpoint{13.205897in}{1.028608in}}%
\pgfpathlineto{\pgfqpoint{13.205897in}{1.023043in}}%
\pgfpathlineto{\pgfqpoint{13.234321in}{1.023043in}}%
\pgfpathlineto{\pgfqpoint{13.234321in}{1.027438in}}%
\pgfpathquadraticcurveto{\pgfqpoint{13.234321in}{1.037957in}}{\pgfqpoint{13.239220in}{1.042766in}}%
\pgfpathquadraticcurveto{\pgfqpoint{13.244119in}{1.047593in}}{\pgfqpoint{13.254764in}{1.047593in}}%
\pgfpathlineto{\pgfqpoint{13.264545in}{1.047593in}}%
\pgfpathlineto{\pgfqpoint{13.264545in}{1.047593in}}%
\pgfpathclose%
\pgfpathmoveto{\pgfqpoint{13.327137in}{0.994115in}}%
\pgfpathlineto{\pgfqpoint{13.327137in}{0.989054in}}%
\pgfpathlineto{\pgfqpoint{13.279513in}{0.989054in}}%
\pgfpathquadraticcurveto{\pgfqpoint{13.280197in}{0.978354in}}{\pgfqpoint{13.285961in}{0.972753in}}%
\pgfpathquadraticcurveto{\pgfqpoint{13.291725in}{0.967151in}}{\pgfqpoint{13.302028in}{0.967151in}}%
\pgfpathquadraticcurveto{\pgfqpoint{13.307990in}{0.967151in}}{\pgfqpoint{13.313592in}{0.968610in}}%
\pgfpathquadraticcurveto{\pgfqpoint{13.319194in}{0.970069in}}{\pgfqpoint{13.324724in}{0.973005in}}%
\pgfpathlineto{\pgfqpoint{13.324724in}{0.963206in}}%
\pgfpathquadraticcurveto{\pgfqpoint{13.319140in}{0.960847in}}{\pgfqpoint{13.313286in}{0.959604in}}%
\pgfpathquadraticcurveto{\pgfqpoint{13.307432in}{0.958361in}}{\pgfqpoint{13.301416in}{0.958361in}}%
\pgfpathquadraticcurveto{\pgfqpoint{13.286322in}{0.958361in}}{\pgfqpoint{13.277514in}{0.967133in}}%
\pgfpathquadraticcurveto{\pgfqpoint{13.268706in}{0.975923in}}{\pgfqpoint{13.268706in}{0.990909in}}%
\pgfpathquadraticcurveto{\pgfqpoint{13.268706in}{1.006381in}}{\pgfqpoint{13.277063in}{1.015459in}}%
\pgfpathquadraticcurveto{\pgfqpoint{13.285421in}{1.024556in}}{\pgfqpoint{13.299615in}{1.024556in}}%
\pgfpathquadraticcurveto{\pgfqpoint{13.312331in}{1.024556in}}{\pgfqpoint{13.319734in}{1.016360in}}%
\pgfpathquadraticcurveto{\pgfqpoint{13.327137in}{1.008183in}}{\pgfqpoint{13.327137in}{0.994115in}}%
\pgfpathlineto{\pgfqpoint{13.327137in}{0.994115in}}%
\pgfpathclose%
\pgfpathmoveto{\pgfqpoint{13.316780in}{0.997159in}}%
\pgfpathquadraticcurveto{\pgfqpoint{13.316672in}{1.005643in}}{\pgfqpoint{13.312025in}{1.010704in}}%
\pgfpathquadraticcurveto{\pgfqpoint{13.307378in}{1.015784in}}{\pgfqpoint{13.299723in}{1.015784in}}%
\pgfpathquadraticcurveto{\pgfqpoint{13.291059in}{1.015784in}}{\pgfqpoint{13.285853in}{1.010884in}}%
\pgfpathquadraticcurveto{\pgfqpoint{13.280648in}{1.005985in}}{\pgfqpoint{13.279855in}{0.997087in}}%
\pgfpathlineto{\pgfqpoint{13.316780in}{0.997159in}}%
\pgfpathlineto{\pgfqpoint{13.316780in}{0.997159in}}%
\pgfpathclose%
\pgfpathmoveto{\pgfqpoint{13.389513in}{1.020629in}}%
\pgfpathlineto{\pgfqpoint{13.389513in}{1.010938in}}%
\pgfpathquadraticcurveto{\pgfqpoint{13.385118in}{1.013370in}}{\pgfqpoint{13.380705in}{1.014577in}}%
\pgfpathquadraticcurveto{\pgfqpoint{13.376292in}{1.015784in}}{\pgfqpoint{13.371789in}{1.015784in}}%
\pgfpathquadraticcurveto{\pgfqpoint{13.361702in}{1.015784in}}{\pgfqpoint{13.356119in}{1.009389in}}%
\pgfpathquadraticcurveto{\pgfqpoint{13.350553in}{1.003013in}}{\pgfqpoint{13.350553in}{0.991467in}}%
\pgfpathquadraticcurveto{\pgfqpoint{13.350553in}{0.979921in}}{\pgfqpoint{13.356119in}{0.973527in}}%
\pgfpathquadraticcurveto{\pgfqpoint{13.361702in}{0.967151in}}{\pgfqpoint{13.371789in}{0.967151in}}%
\pgfpathquadraticcurveto{\pgfqpoint{13.376292in}{0.967151in}}{\pgfqpoint{13.380705in}{0.968358in}}%
\pgfpathquadraticcurveto{\pgfqpoint{13.385118in}{0.969564in}}{\pgfqpoint{13.389513in}{0.971996in}}%
\pgfpathlineto{\pgfqpoint{13.389513in}{0.962414in}}%
\pgfpathquadraticcurveto{\pgfqpoint{13.385172in}{0.960396in}}{\pgfqpoint{13.380525in}{0.959388in}}%
\pgfpathquadraticcurveto{\pgfqpoint{13.375896in}{0.958361in}}{\pgfqpoint{13.370655in}{0.958361in}}%
\pgfpathquadraticcurveto{\pgfqpoint{13.356407in}{0.958361in}}{\pgfqpoint{13.348013in}{0.967313in}}%
\pgfpathquadraticcurveto{\pgfqpoint{13.339638in}{0.976265in}}{\pgfqpoint{13.339638in}{0.991467in}}%
\pgfpathquadraticcurveto{\pgfqpoint{13.339638in}{1.006886in}}{\pgfqpoint{13.348103in}{1.015712in}}%
\pgfpathquadraticcurveto{\pgfqpoint{13.356587in}{1.024556in}}{\pgfqpoint{13.371339in}{1.024556in}}%
\pgfpathquadraticcurveto{\pgfqpoint{13.376112in}{1.024556in}}{\pgfqpoint{13.380669in}{1.023565in}}%
\pgfpathquadraticcurveto{\pgfqpoint{13.385226in}{1.022592in}}{\pgfqpoint{13.389513in}{1.020629in}}%
\pgfpathlineto{\pgfqpoint{13.389513in}{1.020629in}}%
\pgfpathclose%
\pgfpathmoveto{\pgfqpoint{13.417774in}{1.040947in}}%
\pgfpathlineto{\pgfqpoint{13.417774in}{1.023043in}}%
\pgfpathlineto{\pgfqpoint{13.439101in}{1.023043in}}%
\pgfpathlineto{\pgfqpoint{13.439101in}{1.014991in}}%
\pgfpathlineto{\pgfqpoint{13.417774in}{1.014991in}}%
\pgfpathlineto{\pgfqpoint{13.417774in}{0.980768in}}%
\pgfpathquadraticcurveto{\pgfqpoint{13.417774in}{0.973059in}}{\pgfqpoint{13.419882in}{0.970861in}}%
\pgfpathquadraticcurveto{\pgfqpoint{13.421989in}{0.968664in}}{\pgfqpoint{13.428474in}{0.968664in}}%
\pgfpathlineto{\pgfqpoint{13.439101in}{0.968664in}}%
\pgfpathlineto{\pgfqpoint{13.439101in}{0.960000in}}%
\pgfpathlineto{\pgfqpoint{13.428474in}{0.960000in}}%
\pgfpathquadraticcurveto{\pgfqpoint{13.416477in}{0.960000in}}{\pgfqpoint{13.411920in}{0.964467in}}%
\pgfpathquadraticcurveto{\pgfqpoint{13.407363in}{0.968952in}}{\pgfqpoint{13.407363in}{0.980768in}}%
\pgfpathlineto{\pgfqpoint{13.407363in}{1.014991in}}%
\pgfpathlineto{\pgfqpoint{13.399762in}{1.014991in}}%
\pgfpathlineto{\pgfqpoint{13.399762in}{1.023043in}}%
\pgfpathlineto{\pgfqpoint{13.407363in}{1.023043in}}%
\pgfpathlineto{\pgfqpoint{13.407363in}{1.040947in}}%
\pgfpathlineto{\pgfqpoint{13.417774in}{1.040947in}}%
\pgfpathlineto{\pgfqpoint{13.417774in}{1.040947in}}%
\pgfpathclose%
\pgfpathmoveto{\pgfqpoint{13.452718in}{1.023043in}}%
\pgfpathlineto{\pgfqpoint{13.463075in}{1.023043in}}%
\pgfpathlineto{\pgfqpoint{13.463075in}{0.960000in}}%
\pgfpathlineto{\pgfqpoint{13.452718in}{0.960000in}}%
\pgfpathlineto{\pgfqpoint{13.452718in}{1.023043in}}%
\pgfpathlineto{\pgfqpoint{13.452718in}{1.023043in}}%
\pgfpathclose%
\pgfpathmoveto{\pgfqpoint{13.452718in}{1.047593in}}%
\pgfpathlineto{\pgfqpoint{13.463075in}{1.047593in}}%
\pgfpathlineto{\pgfqpoint{13.463075in}{1.034462in}}%
\pgfpathlineto{\pgfqpoint{13.452718in}{1.034462in}}%
\pgfpathlineto{\pgfqpoint{13.452718in}{1.047593in}}%
\pgfpathlineto{\pgfqpoint{13.452718in}{1.047593in}}%
\pgfpathclose%
\pgfpathmoveto{\pgfqpoint{13.477322in}{1.023043in}}%
\pgfpathlineto{\pgfqpoint{13.488292in}{1.023043in}}%
\pgfpathlineto{\pgfqpoint{13.507997in}{0.970141in}}%
\pgfpathlineto{\pgfqpoint{13.527702in}{1.023043in}}%
\pgfpathlineto{\pgfqpoint{13.538672in}{1.023043in}}%
\pgfpathlineto{\pgfqpoint{13.515022in}{0.960000in}}%
\pgfpathlineto{\pgfqpoint{13.500954in}{0.960000in}}%
\pgfpathlineto{\pgfqpoint{13.477322in}{1.023043in}}%
\pgfpathlineto{\pgfqpoint{13.477322in}{1.023043in}}%
\pgfpathclose%
\pgfpathmoveto{\pgfqpoint{13.606902in}{0.994115in}}%
\pgfpathlineto{\pgfqpoint{13.606902in}{0.989054in}}%
\pgfpathlineto{\pgfqpoint{13.559278in}{0.989054in}}%
\pgfpathquadraticcurveto{\pgfqpoint{13.559962in}{0.978354in}}{\pgfqpoint{13.565726in}{0.972753in}}%
\pgfpathquadraticcurveto{\pgfqpoint{13.571490in}{0.967151in}}{\pgfqpoint{13.581793in}{0.967151in}}%
\pgfpathquadraticcurveto{\pgfqpoint{13.587755in}{0.967151in}}{\pgfqpoint{13.593357in}{0.968610in}}%
\pgfpathquadraticcurveto{\pgfqpoint{13.598959in}{0.970069in}}{\pgfqpoint{13.604488in}{0.973005in}}%
\pgfpathlineto{\pgfqpoint{13.604488in}{0.963206in}}%
\pgfpathquadraticcurveto{\pgfqpoint{13.598905in}{0.960847in}}{\pgfqpoint{13.593051in}{0.959604in}}%
\pgfpathquadraticcurveto{\pgfqpoint{13.587197in}{0.958361in}}{\pgfqpoint{13.581181in}{0.958361in}}%
\pgfpathquadraticcurveto{\pgfqpoint{13.566086in}{0.958361in}}{\pgfqpoint{13.557278in}{0.967133in}}%
\pgfpathquadraticcurveto{\pgfqpoint{13.548470in}{0.975923in}}{\pgfqpoint{13.548470in}{0.990909in}}%
\pgfpathquadraticcurveto{\pgfqpoint{13.548470in}{1.006381in}}{\pgfqpoint{13.556828in}{1.015459in}}%
\pgfpathquadraticcurveto{\pgfqpoint{13.565186in}{1.024556in}}{\pgfqpoint{13.579379in}{1.024556in}}%
\pgfpathquadraticcurveto{\pgfqpoint{13.592096in}{1.024556in}}{\pgfqpoint{13.599499in}{1.016360in}}%
\pgfpathquadraticcurveto{\pgfqpoint{13.606902in}{1.008183in}}{\pgfqpoint{13.606902in}{0.994115in}}%
\pgfpathlineto{\pgfqpoint{13.606902in}{0.994115in}}%
\pgfpathclose%
\pgfpathmoveto{\pgfqpoint{13.596545in}{0.997159in}}%
\pgfpathquadraticcurveto{\pgfqpoint{13.596437in}{1.005643in}}{\pgfqpoint{13.591790in}{1.010704in}}%
\pgfpathquadraticcurveto{\pgfqpoint{13.587143in}{1.015784in}}{\pgfqpoint{13.579487in}{1.015784in}}%
\pgfpathquadraticcurveto{\pgfqpoint{13.570824in}{1.015784in}}{\pgfqpoint{13.565618in}{1.010884in}}%
\pgfpathquadraticcurveto{\pgfqpoint{13.560413in}{1.005985in}}{\pgfqpoint{13.559620in}{0.997087in}}%
\pgfpathlineto{\pgfqpoint{13.596545in}{0.997159in}}%
\pgfpathlineto{\pgfqpoint{13.596545in}{0.997159in}}%
\pgfpathclose%
\pgfusepath{fill}%
\end{pgfscope}%
\begin{pgfscope}%
\pgfsetbuttcap%
\pgfsetroundjoin%
\definecolor{currentfill}{rgb}{0.309804,0.372549,0.419608}%
\pgfsetfillcolor{currentfill}%
\pgfsetlinewidth{0.000000pt}%
\definecolor{currentstroke}{rgb}{0.309804,0.372549,0.419608}%
\pgfsetstrokecolor{currentstroke}%
\pgfsetdash{}{0pt}%
\pgfpathmoveto{\pgfqpoint{0.980234in}{0.857429in}}%
\pgfpathlineto{\pgfqpoint{0.980234in}{0.847738in}}%
\pgfpathquadraticcurveto{\pgfqpoint{0.975839in}{0.850170in}}{\pgfqpoint{0.971426in}{0.851377in}}%
\pgfpathquadraticcurveto{\pgfqpoint{0.967013in}{0.852584in}}{\pgfqpoint{0.962510in}{0.852584in}}%
\pgfpathquadraticcurveto{\pgfqpoint{0.952423in}{0.852584in}}{\pgfqpoint{0.946839in}{0.846189in}}%
\pgfpathquadraticcurveto{\pgfqpoint{0.941274in}{0.839813in}}{\pgfqpoint{0.941274in}{0.828267in}}%
\pgfpathquadraticcurveto{\pgfqpoint{0.941274in}{0.816721in}}{\pgfqpoint{0.946839in}{0.810327in}}%
\pgfpathquadraticcurveto{\pgfqpoint{0.952423in}{0.803951in}}{\pgfqpoint{0.962510in}{0.803951in}}%
\pgfpathquadraticcurveto{\pgfqpoint{0.967013in}{0.803951in}}{\pgfqpoint{0.971426in}{0.805158in}}%
\pgfpathquadraticcurveto{\pgfqpoint{0.975839in}{0.806364in}}{\pgfqpoint{0.980234in}{0.808796in}}%
\pgfpathlineto{\pgfqpoint{0.980234in}{0.799214in}}%
\pgfpathquadraticcurveto{\pgfqpoint{0.975893in}{0.797196in}}{\pgfqpoint{0.971246in}{0.796188in}}%
\pgfpathquadraticcurveto{\pgfqpoint{0.966617in}{0.795161in}}{\pgfqpoint{0.961375in}{0.795161in}}%
\pgfpathquadraticcurveto{\pgfqpoint{0.947128in}{0.795161in}}{\pgfqpoint{0.938734in}{0.804113in}}%
\pgfpathquadraticcurveto{\pgfqpoint{0.930358in}{0.813065in}}{\pgfqpoint{0.930358in}{0.828267in}}%
\pgfpathquadraticcurveto{\pgfqpoint{0.930358in}{0.843686in}}{\pgfqpoint{0.938824in}{0.852512in}}%
\pgfpathquadraticcurveto{\pgfqpoint{0.947308in}{0.861356in}}{\pgfqpoint{0.962060in}{0.861356in}}%
\pgfpathquadraticcurveto{\pgfqpoint{0.966833in}{0.861356in}}{\pgfqpoint{0.971390in}{0.860365in}}%
\pgfpathquadraticcurveto{\pgfqpoint{0.975947in}{0.859392in}}{\pgfqpoint{0.980234in}{0.857429in}}%
\pgfpathlineto{\pgfqpoint{0.980234in}{0.857429in}}%
\pgfpathclose%
\pgfpathmoveto{\pgfqpoint{0.998246in}{0.884393in}}%
\pgfpathlineto{\pgfqpoint{1.008603in}{0.884393in}}%
\pgfpathlineto{\pgfqpoint{1.008603in}{0.796800in}}%
\pgfpathlineto{\pgfqpoint{0.998246in}{0.796800in}}%
\pgfpathlineto{\pgfqpoint{0.998246in}{0.884393in}}%
\pgfpathlineto{\pgfqpoint{0.998246in}{0.884393in}}%
\pgfpathclose%
\pgfpathmoveto{\pgfqpoint{1.029209in}{0.821675in}}%
\pgfpathlineto{\pgfqpoint{1.029209in}{0.859843in}}%
\pgfpathlineto{\pgfqpoint{1.039566in}{0.859843in}}%
\pgfpathlineto{\pgfqpoint{1.039566in}{0.822071in}}%
\pgfpathquadraticcurveto{\pgfqpoint{1.039566in}{0.813119in}}{\pgfqpoint{1.043042in}{0.808634in}}%
\pgfpathquadraticcurveto{\pgfqpoint{1.046537in}{0.804167in}}{\pgfqpoint{1.053525in}{0.804167in}}%
\pgfpathquadraticcurveto{\pgfqpoint{1.061901in}{0.804167in}}{\pgfqpoint{1.066764in}{0.809517in}}%
\pgfpathquadraticcurveto{\pgfqpoint{1.071646in}{0.814866in}}{\pgfqpoint{1.071646in}{0.824106in}}%
\pgfpathlineto{\pgfqpoint{1.071646in}{0.859843in}}%
\pgfpathlineto{\pgfqpoint{1.082003in}{0.859843in}}%
\pgfpathlineto{\pgfqpoint{1.082003in}{0.796800in}}%
\pgfpathlineto{\pgfqpoint{1.071646in}{0.796800in}}%
\pgfpathlineto{\pgfqpoint{1.071646in}{0.806491in}}%
\pgfpathquadraticcurveto{\pgfqpoint{1.067881in}{0.800745in}}{\pgfqpoint{1.062892in}{0.797953in}}%
\pgfpathquadraticcurveto{\pgfqpoint{1.057920in}{0.795161in}}{\pgfqpoint{1.051328in}{0.795161in}}%
\pgfpathquadraticcurveto{\pgfqpoint{1.040467in}{0.795161in}}{\pgfqpoint{1.034829in}{0.801915in}}%
\pgfpathquadraticcurveto{\pgfqpoint{1.029209in}{0.808670in}}{\pgfqpoint{1.029209in}{0.821675in}}%
\pgfpathlineto{\pgfqpoint{1.029209in}{0.821675in}}%
\pgfpathclose%
\pgfpathmoveto{\pgfqpoint{1.055273in}{0.861356in}}%
\pgfpathlineto{\pgfqpoint{1.055273in}{0.861356in}}%
\pgfpathlineto{\pgfqpoint{1.055273in}{0.861356in}}%
\pgfpathclose%
\pgfpathmoveto{\pgfqpoint{1.143514in}{0.857987in}}%
\pgfpathlineto{\pgfqpoint{1.143514in}{0.848189in}}%
\pgfpathquadraticcurveto{\pgfqpoint{1.139137in}{0.850440in}}{\pgfqpoint{1.134400in}{0.851557in}}%
\pgfpathquadraticcurveto{\pgfqpoint{1.129681in}{0.852692in}}{\pgfqpoint{1.124601in}{0.852692in}}%
\pgfpathquadraticcurveto{\pgfqpoint{1.116892in}{0.852692in}}{\pgfqpoint{1.113038in}{0.850332in}}%
\pgfpathquadraticcurveto{\pgfqpoint{1.109183in}{0.847973in}}{\pgfqpoint{1.109183in}{0.843235in}}%
\pgfpathquadraticcurveto{\pgfqpoint{1.109183in}{0.839633in}}{\pgfqpoint{1.111939in}{0.837580in}}%
\pgfpathquadraticcurveto{\pgfqpoint{1.114695in}{0.835526in}}{\pgfqpoint{1.123034in}{0.833671in}}%
\pgfpathlineto{\pgfqpoint{1.126583in}{0.832878in}}%
\pgfpathquadraticcurveto{\pgfqpoint{1.137606in}{0.830519in}}{\pgfqpoint{1.142253in}{0.826214in}}%
\pgfpathquadraticcurveto{\pgfqpoint{1.146900in}{0.821909in}}{\pgfqpoint{1.146900in}{0.814200in}}%
\pgfpathquadraticcurveto{\pgfqpoint{1.146900in}{0.805410in}}{\pgfqpoint{1.139948in}{0.800276in}}%
\pgfpathquadraticcurveto{\pgfqpoint{1.132995in}{0.795161in}}{\pgfqpoint{1.120837in}{0.795161in}}%
\pgfpathquadraticcurveto{\pgfqpoint{1.115775in}{0.795161in}}{\pgfqpoint{1.110282in}{0.796152in}}%
\pgfpathquadraticcurveto{\pgfqpoint{1.104788in}{0.797142in}}{\pgfqpoint{1.098718in}{0.799106in}}%
\pgfpathlineto{\pgfqpoint{1.098718in}{0.809805in}}%
\pgfpathquadraticcurveto{\pgfqpoint{1.104464in}{0.806815in}}{\pgfqpoint{1.110030in}{0.805320in}}%
\pgfpathquadraticcurveto{\pgfqpoint{1.115595in}{0.803843in}}{\pgfqpoint{1.121071in}{0.803843in}}%
\pgfpathquadraticcurveto{\pgfqpoint{1.128384in}{0.803843in}}{\pgfqpoint{1.132311in}{0.806346in}}%
\pgfpathquadraticcurveto{\pgfqpoint{1.136255in}{0.808850in}}{\pgfqpoint{1.136255in}{0.813407in}}%
\pgfpathquadraticcurveto{\pgfqpoint{1.136255in}{0.817622in}}{\pgfqpoint{1.133409in}{0.819874in}}%
\pgfpathquadraticcurveto{\pgfqpoint{1.130581in}{0.822125in}}{\pgfqpoint{1.120945in}{0.824214in}}%
\pgfpathlineto{\pgfqpoint{1.117342in}{0.825061in}}%
\pgfpathquadraticcurveto{\pgfqpoint{1.107724in}{0.827078in}}{\pgfqpoint{1.103437in}{0.831275in}}%
\pgfpathquadraticcurveto{\pgfqpoint{1.099168in}{0.835472in}}{\pgfqpoint{1.099168in}{0.842785in}}%
\pgfpathquadraticcurveto{\pgfqpoint{1.099168in}{0.851683in}}{\pgfqpoint{1.105472in}{0.856510in}}%
\pgfpathquadraticcurveto{\pgfqpoint{1.111777in}{0.861356in}}{\pgfqpoint{1.123377in}{0.861356in}}%
\pgfpathquadraticcurveto{\pgfqpoint{1.129104in}{0.861356in}}{\pgfqpoint{1.134166in}{0.860509in}}%
\pgfpathquadraticcurveto{\pgfqpoint{1.139245in}{0.859680in}}{\pgfqpoint{1.143514in}{0.857987in}}%
\pgfpathlineto{\pgfqpoint{1.143514in}{0.857987in}}%
\pgfpathclose%
\pgfpathmoveto{\pgfqpoint{1.173630in}{0.877747in}}%
\pgfpathlineto{\pgfqpoint{1.173630in}{0.859843in}}%
\pgfpathlineto{\pgfqpoint{1.194957in}{0.859843in}}%
\pgfpathlineto{\pgfqpoint{1.194957in}{0.851791in}}%
\pgfpathlineto{\pgfqpoint{1.173630in}{0.851791in}}%
\pgfpathlineto{\pgfqpoint{1.173630in}{0.817568in}}%
\pgfpathquadraticcurveto{\pgfqpoint{1.173630in}{0.809859in}}{\pgfqpoint{1.175738in}{0.807661in}}%
\pgfpathquadraticcurveto{\pgfqpoint{1.177845in}{0.805464in}}{\pgfqpoint{1.184330in}{0.805464in}}%
\pgfpathlineto{\pgfqpoint{1.194957in}{0.805464in}}%
\pgfpathlineto{\pgfqpoint{1.194957in}{0.796800in}}%
\pgfpathlineto{\pgfqpoint{1.184330in}{0.796800in}}%
\pgfpathquadraticcurveto{\pgfqpoint{1.172334in}{0.796800in}}{\pgfqpoint{1.167776in}{0.801267in}}%
\pgfpathquadraticcurveto{\pgfqpoint{1.163219in}{0.805752in}}{\pgfqpoint{1.163219in}{0.817568in}}%
\pgfpathlineto{\pgfqpoint{1.163219in}{0.851791in}}%
\pgfpathlineto{\pgfqpoint{1.155618in}{0.851791in}}%
\pgfpathlineto{\pgfqpoint{1.155618in}{0.859843in}}%
\pgfpathlineto{\pgfqpoint{1.163219in}{0.859843in}}%
\pgfpathlineto{\pgfqpoint{1.163219in}{0.877747in}}%
\pgfpathlineto{\pgfqpoint{1.173630in}{0.877747in}}%
\pgfpathlineto{\pgfqpoint{1.173630in}{0.877747in}}%
\pgfpathclose%
\pgfpathmoveto{\pgfqpoint{1.262502in}{0.830915in}}%
\pgfpathlineto{\pgfqpoint{1.262502in}{0.825854in}}%
\pgfpathlineto{\pgfqpoint{1.214878in}{0.825854in}}%
\pgfpathquadraticcurveto{\pgfqpoint{1.215563in}{0.815154in}}{\pgfqpoint{1.221327in}{0.809553in}}%
\pgfpathquadraticcurveto{\pgfqpoint{1.227090in}{0.803951in}}{\pgfqpoint{1.237393in}{0.803951in}}%
\pgfpathquadraticcurveto{\pgfqpoint{1.243355in}{0.803951in}}{\pgfqpoint{1.248957in}{0.805410in}}%
\pgfpathquadraticcurveto{\pgfqpoint{1.254559in}{0.806869in}}{\pgfqpoint{1.260089in}{0.809805in}}%
\pgfpathlineto{\pgfqpoint{1.260089in}{0.800006in}}%
\pgfpathquadraticcurveto{\pgfqpoint{1.254505in}{0.797647in}}{\pgfqpoint{1.248651in}{0.796404in}}%
\pgfpathquadraticcurveto{\pgfqpoint{1.242797in}{0.795161in}}{\pgfqpoint{1.236781in}{0.795161in}}%
\pgfpathquadraticcurveto{\pgfqpoint{1.221687in}{0.795161in}}{\pgfqpoint{1.212879in}{0.803933in}}%
\pgfpathquadraticcurveto{\pgfqpoint{1.204071in}{0.812723in}}{\pgfqpoint{1.204071in}{0.827709in}}%
\pgfpathquadraticcurveto{\pgfqpoint{1.204071in}{0.843181in}}{\pgfqpoint{1.212429in}{0.852259in}}%
\pgfpathquadraticcurveto{\pgfqpoint{1.220786in}{0.861356in}}{\pgfqpoint{1.234980in}{0.861356in}}%
\pgfpathquadraticcurveto{\pgfqpoint{1.247696in}{0.861356in}}{\pgfqpoint{1.255099in}{0.853160in}}%
\pgfpathquadraticcurveto{\pgfqpoint{1.262502in}{0.844983in}}{\pgfqpoint{1.262502in}{0.830915in}}%
\pgfpathlineto{\pgfqpoint{1.262502in}{0.830915in}}%
\pgfpathclose%
\pgfpathmoveto{\pgfqpoint{1.252145in}{0.833959in}}%
\pgfpathquadraticcurveto{\pgfqpoint{1.252037in}{0.842443in}}{\pgfqpoint{1.247390in}{0.847504in}}%
\pgfpathquadraticcurveto{\pgfqpoint{1.242743in}{0.852584in}}{\pgfqpoint{1.235088in}{0.852584in}}%
\pgfpathquadraticcurveto{\pgfqpoint{1.226424in}{0.852584in}}{\pgfqpoint{1.221219in}{0.847684in}}%
\pgfpathquadraticcurveto{\pgfqpoint{1.216013in}{0.842785in}}{\pgfqpoint{1.215220in}{0.833887in}}%
\pgfpathlineto{\pgfqpoint{1.252145in}{0.833959in}}%
\pgfpathlineto{\pgfqpoint{1.252145in}{0.833959in}}%
\pgfpathclose%
\pgfpathmoveto{\pgfqpoint{1.316035in}{0.850170in}}%
\pgfpathquadraticcurveto{\pgfqpoint{1.314287in}{0.851179in}}{\pgfqpoint{1.312234in}{0.851647in}}%
\pgfpathquadraticcurveto{\pgfqpoint{1.310181in}{0.852133in}}{\pgfqpoint{1.307713in}{0.852133in}}%
\pgfpathquadraticcurveto{\pgfqpoint{1.298923in}{0.852133in}}{\pgfqpoint{1.294222in}{0.846423in}}%
\pgfpathquadraticcurveto{\pgfqpoint{1.289521in}{0.840714in}}{\pgfqpoint{1.289521in}{0.830014in}}%
\pgfpathlineto{\pgfqpoint{1.289521in}{0.796800in}}%
\pgfpathlineto{\pgfqpoint{1.279110in}{0.796800in}}%
\pgfpathlineto{\pgfqpoint{1.279110in}{0.859843in}}%
\pgfpathlineto{\pgfqpoint{1.289521in}{0.859843in}}%
\pgfpathlineto{\pgfqpoint{1.289521in}{0.850044in}}%
\pgfpathquadraticcurveto{\pgfqpoint{1.292799in}{0.855790in}}{\pgfqpoint{1.298022in}{0.858564in}}%
\pgfpathquadraticcurveto{\pgfqpoint{1.303264in}{0.861356in}}{\pgfqpoint{1.310757in}{0.861356in}}%
\pgfpathquadraticcurveto{\pgfqpoint{1.311820in}{0.861356in}}{\pgfqpoint{1.313117in}{0.861211in}}%
\pgfpathquadraticcurveto{\pgfqpoint{1.314413in}{0.861085in}}{\pgfqpoint{1.315980in}{0.860797in}}%
\pgfpathlineto{\pgfqpoint{1.316035in}{0.850170in}}%
\pgfpathlineto{\pgfqpoint{1.316035in}{0.850170in}}%
\pgfpathclose%
\pgfpathmoveto{\pgfqpoint{1.367081in}{0.857987in}}%
\pgfpathlineto{\pgfqpoint{1.367081in}{0.848189in}}%
\pgfpathquadraticcurveto{\pgfqpoint{1.362704in}{0.850440in}}{\pgfqpoint{1.357967in}{0.851557in}}%
\pgfpathquadraticcurveto{\pgfqpoint{1.353248in}{0.852692in}}{\pgfqpoint{1.348168in}{0.852692in}}%
\pgfpathquadraticcurveto{\pgfqpoint{1.340459in}{0.852692in}}{\pgfqpoint{1.336604in}{0.850332in}}%
\pgfpathquadraticcurveto{\pgfqpoint{1.332750in}{0.847973in}}{\pgfqpoint{1.332750in}{0.843235in}}%
\pgfpathquadraticcurveto{\pgfqpoint{1.332750in}{0.839633in}}{\pgfqpoint{1.335506in}{0.837580in}}%
\pgfpathquadraticcurveto{\pgfqpoint{1.338262in}{0.835526in}}{\pgfqpoint{1.346601in}{0.833671in}}%
\pgfpathlineto{\pgfqpoint{1.350150in}{0.832878in}}%
\pgfpathquadraticcurveto{\pgfqpoint{1.361173in}{0.830519in}}{\pgfqpoint{1.365820in}{0.826214in}}%
\pgfpathquadraticcurveto{\pgfqpoint{1.370467in}{0.821909in}}{\pgfqpoint{1.370467in}{0.814200in}}%
\pgfpathquadraticcurveto{\pgfqpoint{1.370467in}{0.805410in}}{\pgfqpoint{1.363515in}{0.800276in}}%
\pgfpathquadraticcurveto{\pgfqpoint{1.356562in}{0.795161in}}{\pgfqpoint{1.344404in}{0.795161in}}%
\pgfpathquadraticcurveto{\pgfqpoint{1.339342in}{0.795161in}}{\pgfqpoint{1.333849in}{0.796152in}}%
\pgfpathquadraticcurveto{\pgfqpoint{1.328355in}{0.797142in}}{\pgfqpoint{1.322285in}{0.799106in}}%
\pgfpathlineto{\pgfqpoint{1.322285in}{0.809805in}}%
\pgfpathquadraticcurveto{\pgfqpoint{1.328031in}{0.806815in}}{\pgfqpoint{1.333596in}{0.805320in}}%
\pgfpathquadraticcurveto{\pgfqpoint{1.339162in}{0.803843in}}{\pgfqpoint{1.344638in}{0.803843in}}%
\pgfpathquadraticcurveto{\pgfqpoint{1.351951in}{0.803843in}}{\pgfqpoint{1.355877in}{0.806346in}}%
\pgfpathquadraticcurveto{\pgfqpoint{1.359822in}{0.808850in}}{\pgfqpoint{1.359822in}{0.813407in}}%
\pgfpathquadraticcurveto{\pgfqpoint{1.359822in}{0.817622in}}{\pgfqpoint{1.356976in}{0.819874in}}%
\pgfpathquadraticcurveto{\pgfqpoint{1.354148in}{0.822125in}}{\pgfqpoint{1.344512in}{0.824214in}}%
\pgfpathlineto{\pgfqpoint{1.340909in}{0.825061in}}%
\pgfpathquadraticcurveto{\pgfqpoint{1.331291in}{0.827078in}}{\pgfqpoint{1.327004in}{0.831275in}}%
\pgfpathquadraticcurveto{\pgfqpoint{1.322735in}{0.835472in}}{\pgfqpoint{1.322735in}{0.842785in}}%
\pgfpathquadraticcurveto{\pgfqpoint{1.322735in}{0.851683in}}{\pgfqpoint{1.329039in}{0.856510in}}%
\pgfpathquadraticcurveto{\pgfqpoint{1.335344in}{0.861356in}}{\pgfqpoint{1.346943in}{0.861356in}}%
\pgfpathquadraticcurveto{\pgfqpoint{1.352671in}{0.861356in}}{\pgfqpoint{1.357733in}{0.860509in}}%
\pgfpathquadraticcurveto{\pgfqpoint{1.362812in}{0.859680in}}{\pgfqpoint{1.367081in}{0.857987in}}%
\pgfpathlineto{\pgfqpoint{1.367081in}{0.857987in}}%
\pgfpathclose%
\pgfusepath{fill}%
\end{pgfscope}%
\begin{pgfscope}%
\pgfsetbuttcap%
\pgfsetroundjoin%
\definecolor{currentfill}{rgb}{0.309804,0.372549,0.419608}%
\pgfsetfillcolor{currentfill}%
\pgfsetlinewidth{0.000000pt}%
\definecolor{currentstroke}{rgb}{0.309804,0.372549,0.419608}%
\pgfsetstrokecolor{currentstroke}%
\pgfsetdash{}{0pt}%
\pgfpathmoveto{\pgfqpoint{1.532074in}{0.853538in}}%
\pgfpathlineto{\pgfqpoint{1.473985in}{0.832878in}}%
\pgfpathlineto{\pgfqpoint{1.532074in}{0.812344in}}%
\pgfpathlineto{\pgfqpoint{1.532074in}{0.802096in}}%
\pgfpathlineto{\pgfqpoint{1.459917in}{0.828267in}}%
\pgfpathlineto{\pgfqpoint{1.459917in}{0.837616in}}%
\pgfpathlineto{\pgfqpoint{1.532074in}{0.863787in}}%
\pgfpathlineto{\pgfqpoint{1.532074in}{0.853538in}}%
\pgfpathlineto{\pgfqpoint{1.532074in}{0.853538in}}%
\pgfpathclose%
\pgfpathmoveto{\pgfqpoint{1.556750in}{0.880845in}}%
\pgfpathlineto{\pgfqpoint{1.601385in}{0.880845in}}%
\pgfpathlineto{\pgfqpoint{1.601385in}{0.871262in}}%
\pgfpathlineto{\pgfqpoint{1.567161in}{0.871262in}}%
\pgfpathlineto{\pgfqpoint{1.567161in}{0.850674in}}%
\pgfpathquadraticcurveto{\pgfqpoint{1.569629in}{0.851521in}}{\pgfqpoint{1.572097in}{0.851935in}}%
\pgfpathquadraticcurveto{\pgfqpoint{1.574582in}{0.852349in}}{\pgfqpoint{1.577068in}{0.852349in}}%
\pgfpathquadraticcurveto{\pgfqpoint{1.591136in}{0.852349in}}{\pgfqpoint{1.599349in}{0.844640in}}%
\pgfpathquadraticcurveto{\pgfqpoint{1.607581in}{0.836931in}}{\pgfqpoint{1.607581in}{0.823764in}}%
\pgfpathquadraticcurveto{\pgfqpoint{1.607581in}{0.810201in}}{\pgfqpoint{1.599133in}{0.802672in}}%
\pgfpathquadraticcurveto{\pgfqpoint{1.590685in}{0.795161in}}{\pgfqpoint{1.575321in}{0.795161in}}%
\pgfpathquadraticcurveto{\pgfqpoint{1.570025in}{0.795161in}}{\pgfqpoint{1.564532in}{0.796062in}}%
\pgfpathquadraticcurveto{\pgfqpoint{1.559056in}{0.796962in}}{\pgfqpoint{1.553202in}{0.798763in}}%
\pgfpathlineto{\pgfqpoint{1.553202in}{0.810201in}}%
\pgfpathquadraticcurveto{\pgfqpoint{1.558263in}{0.807445in}}{\pgfqpoint{1.563667in}{0.806094in}}%
\pgfpathquadraticcurveto{\pgfqpoint{1.569071in}{0.804743in}}{\pgfqpoint{1.575087in}{0.804743in}}%
\pgfpathquadraticcurveto{\pgfqpoint{1.584831in}{0.804743in}}{\pgfqpoint{1.590505in}{0.809859in}}%
\pgfpathquadraticcurveto{\pgfqpoint{1.596197in}{0.814974in}}{\pgfqpoint{1.596197in}{0.823764in}}%
\pgfpathquadraticcurveto{\pgfqpoint{1.596197in}{0.832536in}}{\pgfqpoint{1.590505in}{0.837652in}}%
\pgfpathquadraticcurveto{\pgfqpoint{1.584831in}{0.842785in}}{\pgfqpoint{1.575087in}{0.842785in}}%
\pgfpathquadraticcurveto{\pgfqpoint{1.570530in}{0.842785in}}{\pgfqpoint{1.565991in}{0.841776in}}%
\pgfpathquadraticcurveto{\pgfqpoint{1.561470in}{0.840768in}}{\pgfqpoint{1.556750in}{0.838624in}}%
\pgfpathlineto{\pgfqpoint{1.556750in}{0.880845in}}%
\pgfpathlineto{\pgfqpoint{1.556750in}{0.880845in}}%
\pgfpathclose%
\pgfusepath{fill}%
\end{pgfscope}%
\begin{pgfscope}%
\pgfsetbuttcap%
\pgfsetroundjoin%
\definecolor{currentfill}{rgb}{0.309804,0.372549,0.419608}%
\pgfsetfillcolor{currentfill}%
\pgfsetlinewidth{0.000000pt}%
\definecolor{currentstroke}{rgb}{0.309804,0.372549,0.419608}%
\pgfsetstrokecolor{currentstroke}%
\pgfsetdash{}{0pt}%
\pgfpathmoveto{\pgfqpoint{1.716696in}{0.852584in}}%
\pgfpathquadraticcurveto{\pgfqpoint{1.708374in}{0.852584in}}{\pgfqpoint{1.703529in}{0.846081in}}%
\pgfpathquadraticcurveto{\pgfqpoint{1.698683in}{0.839579in}}{\pgfqpoint{1.698683in}{0.828267in}}%
\pgfpathquadraticcurveto{\pgfqpoint{1.698683in}{0.816956in}}{\pgfqpoint{1.703493in}{0.810453in}}%
\pgfpathquadraticcurveto{\pgfqpoint{1.708320in}{0.803951in}}{\pgfqpoint{1.716696in}{0.803951in}}%
\pgfpathquadraticcurveto{\pgfqpoint{1.724981in}{0.803951in}}{\pgfqpoint{1.729808in}{0.810471in}}%
\pgfpathquadraticcurveto{\pgfqpoint{1.734654in}{0.817010in}}{\pgfqpoint{1.734654in}{0.828267in}}%
\pgfpathquadraticcurveto{\pgfqpoint{1.734654in}{0.839471in}}{\pgfqpoint{1.729808in}{0.846027in}}%
\pgfpathquadraticcurveto{\pgfqpoint{1.724981in}{0.852584in}}{\pgfqpoint{1.716696in}{0.852584in}}%
\pgfpathlineto{\pgfqpoint{1.716696in}{0.852584in}}%
\pgfpathclose%
\pgfpathmoveto{\pgfqpoint{1.716696in}{0.861356in}}%
\pgfpathquadraticcurveto{\pgfqpoint{1.730205in}{0.861356in}}{\pgfqpoint{1.737914in}{0.852566in}}%
\pgfpathquadraticcurveto{\pgfqpoint{1.745641in}{0.843794in}}{\pgfqpoint{1.745641in}{0.828267in}}%
\pgfpathquadraticcurveto{\pgfqpoint{1.745641in}{0.812795in}}{\pgfqpoint{1.737914in}{0.803969in}}%
\pgfpathquadraticcurveto{\pgfqpoint{1.730205in}{0.795161in}}{\pgfqpoint{1.716696in}{0.795161in}}%
\pgfpathquadraticcurveto{\pgfqpoint{1.703132in}{0.795161in}}{\pgfqpoint{1.695441in}{0.803969in}}%
\pgfpathquadraticcurveto{\pgfqpoint{1.687768in}{0.812795in}}{\pgfqpoint{1.687768in}{0.828267in}}%
\pgfpathquadraticcurveto{\pgfqpoint{1.687768in}{0.843794in}}{\pgfqpoint{1.695441in}{0.852566in}}%
\pgfpathquadraticcurveto{\pgfqpoint{1.703132in}{0.861356in}}{\pgfqpoint{1.716696in}{0.861356in}}%
\pgfpathlineto{\pgfqpoint{1.716696in}{0.861356in}}%
\pgfpathclose%
\pgfpathmoveto{\pgfqpoint{1.799335in}{0.850170in}}%
\pgfpathquadraticcurveto{\pgfqpoint{1.797588in}{0.851179in}}{\pgfqpoint{1.795535in}{0.851647in}}%
\pgfpathquadraticcurveto{\pgfqpoint{1.793481in}{0.852133in}}{\pgfqpoint{1.791014in}{0.852133in}}%
\pgfpathquadraticcurveto{\pgfqpoint{1.782224in}{0.852133in}}{\pgfqpoint{1.777523in}{0.846423in}}%
\pgfpathquadraticcurveto{\pgfqpoint{1.772821in}{0.840714in}}{\pgfqpoint{1.772821in}{0.830014in}}%
\pgfpathlineto{\pgfqpoint{1.772821in}{0.796800in}}%
\pgfpathlineto{\pgfqpoint{1.762410in}{0.796800in}}%
\pgfpathlineto{\pgfqpoint{1.762410in}{0.859843in}}%
\pgfpathlineto{\pgfqpoint{1.772821in}{0.859843in}}%
\pgfpathlineto{\pgfqpoint{1.772821in}{0.850044in}}%
\pgfpathquadraticcurveto{\pgfqpoint{1.776100in}{0.855790in}}{\pgfqpoint{1.781323in}{0.858564in}}%
\pgfpathquadraticcurveto{\pgfqpoint{1.786565in}{0.861356in}}{\pgfqpoint{1.794058in}{0.861356in}}%
\pgfpathquadraticcurveto{\pgfqpoint{1.795120in}{0.861356in}}{\pgfqpoint{1.796417in}{0.861211in}}%
\pgfpathquadraticcurveto{\pgfqpoint{1.797714in}{0.861085in}}{\pgfqpoint{1.799281in}{0.860797in}}%
\pgfpathlineto{\pgfqpoint{1.799335in}{0.850170in}}%
\pgfpathlineto{\pgfqpoint{1.799335in}{0.850170in}}%
\pgfpathclose%
\pgfusepath{fill}%
\end{pgfscope}%
\begin{pgfscope}%
\pgfsetbuttcap%
\pgfsetroundjoin%
\definecolor{currentfill}{rgb}{0.309804,0.372549,0.419608}%
\pgfsetfillcolor{currentfill}%
\pgfsetlinewidth{0.000000pt}%
\definecolor{currentstroke}{rgb}{0.309804,0.372549,0.419608}%
\pgfsetstrokecolor{currentstroke}%
\pgfsetdash{}{0pt}%
\pgfpathmoveto{\pgfqpoint{1.915059in}{0.847738in}}%
\pgfpathquadraticcurveto{\pgfqpoint{1.918950in}{0.854727in}}{\pgfqpoint{1.924353in}{0.858041in}}%
\pgfpathquadraticcurveto{\pgfqpoint{1.929757in}{0.861356in}}{\pgfqpoint{1.937070in}{0.861356in}}%
\pgfpathquadraticcurveto{\pgfqpoint{1.946923in}{0.861356in}}{\pgfqpoint{1.952272in}{0.854457in}}%
\pgfpathquadraticcurveto{\pgfqpoint{1.957622in}{0.847576in}}{\pgfqpoint{1.957622in}{0.834860in}}%
\pgfpathlineto{\pgfqpoint{1.957622in}{0.796800in}}%
\pgfpathlineto{\pgfqpoint{1.947211in}{0.796800in}}%
\pgfpathlineto{\pgfqpoint{1.947211in}{0.834517in}}%
\pgfpathquadraticcurveto{\pgfqpoint{1.947211in}{0.843578in}}{\pgfqpoint{1.943987in}{0.847955in}}%
\pgfpathquadraticcurveto{\pgfqpoint{1.940780in}{0.852349in}}{\pgfqpoint{1.934206in}{0.852349in}}%
\pgfpathquadraticcurveto{\pgfqpoint{1.926155in}{0.852349in}}{\pgfqpoint{1.921471in}{0.847000in}}%
\pgfpathquadraticcurveto{\pgfqpoint{1.916806in}{0.841668in}}{\pgfqpoint{1.916806in}{0.832428in}}%
\pgfpathlineto{\pgfqpoint{1.916806in}{0.796800in}}%
\pgfpathlineto{\pgfqpoint{1.906395in}{0.796800in}}%
\pgfpathlineto{\pgfqpoint{1.906395in}{0.834517in}}%
\pgfpathquadraticcurveto{\pgfqpoint{1.906395in}{0.843632in}}{\pgfqpoint{1.903189in}{0.847991in}}%
\pgfpathquadraticcurveto{\pgfqpoint{1.899983in}{0.852349in}}{\pgfqpoint{1.893282in}{0.852349in}}%
\pgfpathquadraticcurveto{\pgfqpoint{1.885339in}{0.852349in}}{\pgfqpoint{1.880656in}{0.846982in}}%
\pgfpathquadraticcurveto{\pgfqpoint{1.875991in}{0.841614in}}{\pgfqpoint{1.875991in}{0.832428in}}%
\pgfpathlineto{\pgfqpoint{1.875991in}{0.796800in}}%
\pgfpathlineto{\pgfqpoint{1.865580in}{0.796800in}}%
\pgfpathlineto{\pgfqpoint{1.865580in}{0.859843in}}%
\pgfpathlineto{\pgfqpoint{1.875991in}{0.859843in}}%
\pgfpathlineto{\pgfqpoint{1.875991in}{0.850044in}}%
\pgfpathquadraticcurveto{\pgfqpoint{1.879539in}{0.855844in}}{\pgfqpoint{1.884492in}{0.858600in}}%
\pgfpathquadraticcurveto{\pgfqpoint{1.889446in}{0.861356in}}{\pgfqpoint{1.896254in}{0.861356in}}%
\pgfpathquadraticcurveto{\pgfqpoint{1.903135in}{0.861356in}}{\pgfqpoint{1.907944in}{0.857861in}}%
\pgfpathquadraticcurveto{\pgfqpoint{1.912753in}{0.854385in}}{\pgfqpoint{1.915059in}{0.847738in}}%
\pgfpathlineto{\pgfqpoint{1.915059in}{0.847738in}}%
\pgfpathclose%
\pgfpathmoveto{\pgfqpoint{2.006921in}{0.828483in}}%
\pgfpathquadraticcurveto{\pgfqpoint{1.994367in}{0.828483in}}{\pgfqpoint{1.989521in}{0.825619in}}%
\pgfpathquadraticcurveto{\pgfqpoint{1.984676in}{0.822756in}}{\pgfqpoint{1.984676in}{0.815821in}}%
\pgfpathquadraticcurveto{\pgfqpoint{1.984676in}{0.810309in}}{\pgfqpoint{1.988314in}{0.807067in}}%
\pgfpathquadraticcurveto{\pgfqpoint{1.991953in}{0.803843in}}{\pgfqpoint{1.998185in}{0.803843in}}%
\pgfpathquadraticcurveto{\pgfqpoint{2.006813in}{0.803843in}}{\pgfqpoint{2.012018in}{0.809949in}}%
\pgfpathquadraticcurveto{\pgfqpoint{2.017224in}{0.816055in}}{\pgfqpoint{2.017224in}{0.826178in}}%
\pgfpathlineto{\pgfqpoint{2.017224in}{0.828483in}}%
\pgfpathlineto{\pgfqpoint{2.006921in}{0.828483in}}%
\pgfpathlineto{\pgfqpoint{2.006921in}{0.828483in}}%
\pgfpathclose%
\pgfpathmoveto{\pgfqpoint{2.027581in}{0.832770in}}%
\pgfpathlineto{\pgfqpoint{2.027581in}{0.796800in}}%
\pgfpathlineto{\pgfqpoint{2.017224in}{0.796800in}}%
\pgfpathlineto{\pgfqpoint{2.017224in}{0.806364in}}%
\pgfpathquadraticcurveto{\pgfqpoint{2.013676in}{0.800637in}}{\pgfqpoint{2.008380in}{0.797899in}}%
\pgfpathquadraticcurveto{\pgfqpoint{2.003084in}{0.795161in}}{\pgfqpoint{1.995429in}{0.795161in}}%
\pgfpathquadraticcurveto{\pgfqpoint{1.985757in}{0.795161in}}{\pgfqpoint{1.980029in}{0.800601in}}%
\pgfpathquadraticcurveto{\pgfqpoint{1.974319in}{0.806040in}}{\pgfqpoint{1.974319in}{0.815154in}}%
\pgfpathquadraticcurveto{\pgfqpoint{1.974319in}{0.825782in}}{\pgfqpoint{1.981434in}{0.831185in}}%
\pgfpathquadraticcurveto{\pgfqpoint{1.988567in}{0.836589in}}{\pgfqpoint{2.002688in}{0.836589in}}%
\pgfpathlineto{\pgfqpoint{2.017224in}{0.836589in}}%
\pgfpathlineto{\pgfqpoint{2.017224in}{0.837616in}}%
\pgfpathquadraticcurveto{\pgfqpoint{2.017224in}{0.844766in}}{\pgfqpoint{2.012523in}{0.848675in}}%
\pgfpathquadraticcurveto{\pgfqpoint{2.007822in}{0.852584in}}{\pgfqpoint{1.999320in}{0.852584in}}%
\pgfpathquadraticcurveto{\pgfqpoint{1.993916in}{0.852584in}}{\pgfqpoint{1.988783in}{0.851287in}}%
\pgfpathquadraticcurveto{\pgfqpoint{1.983667in}{0.849990in}}{\pgfqpoint{1.978948in}{0.847396in}}%
\pgfpathlineto{\pgfqpoint{1.978948in}{0.856979in}}%
\pgfpathquadraticcurveto{\pgfqpoint{1.984622in}{0.859176in}}{\pgfqpoint{1.989972in}{0.860257in}}%
\pgfpathquadraticcurveto{\pgfqpoint{1.995321in}{0.861356in}}{\pgfqpoint{2.000383in}{0.861356in}}%
\pgfpathquadraticcurveto{\pgfqpoint{2.014072in}{0.861356in}}{\pgfqpoint{2.020826in}{0.854259in}}%
\pgfpathquadraticcurveto{\pgfqpoint{2.027581in}{0.847180in}}{\pgfqpoint{2.027581in}{0.832770in}}%
\pgfpathlineto{\pgfqpoint{2.027581in}{0.832770in}}%
\pgfpathclose%
\pgfpathmoveto{\pgfqpoint{2.101323in}{0.859843in}}%
\pgfpathlineto{\pgfqpoint{2.078519in}{0.829168in}}%
\pgfpathlineto{\pgfqpoint{2.102493in}{0.796800in}}%
\pgfpathlineto{\pgfqpoint{2.090281in}{0.796800in}}%
\pgfpathlineto{\pgfqpoint{2.071927in}{0.821567in}}%
\pgfpathlineto{\pgfqpoint{2.053590in}{0.796800in}}%
\pgfpathlineto{\pgfqpoint{2.041360in}{0.796800in}}%
\pgfpathlineto{\pgfqpoint{2.065857in}{0.829780in}}%
\pgfpathlineto{\pgfqpoint{2.043450in}{0.859843in}}%
\pgfpathlineto{\pgfqpoint{2.055662in}{0.859843in}}%
\pgfpathlineto{\pgfqpoint{2.072377in}{0.837381in}}%
\pgfpathlineto{\pgfqpoint{2.089092in}{0.859843in}}%
\pgfpathlineto{\pgfqpoint{2.101323in}{0.859843in}}%
\pgfpathlineto{\pgfqpoint{2.101323in}{0.859843in}}%
\pgfpathclose%
\pgfpathmoveto{\pgfqpoint{2.117137in}{0.859843in}}%
\pgfpathlineto{\pgfqpoint{2.127494in}{0.859843in}}%
\pgfpathlineto{\pgfqpoint{2.127494in}{0.796800in}}%
\pgfpathlineto{\pgfqpoint{2.117137in}{0.796800in}}%
\pgfpathlineto{\pgfqpoint{2.117137in}{0.859843in}}%
\pgfpathlineto{\pgfqpoint{2.117137in}{0.859843in}}%
\pgfpathclose%
\pgfpathmoveto{\pgfqpoint{2.117137in}{0.884393in}}%
\pgfpathlineto{\pgfqpoint{2.127494in}{0.884393in}}%
\pgfpathlineto{\pgfqpoint{2.127494in}{0.871262in}}%
\pgfpathlineto{\pgfqpoint{2.117137in}{0.871262in}}%
\pgfpathlineto{\pgfqpoint{2.117137in}{0.884393in}}%
\pgfpathlineto{\pgfqpoint{2.117137in}{0.884393in}}%
\pgfpathclose%
\pgfpathmoveto{\pgfqpoint{2.198246in}{0.847738in}}%
\pgfpathquadraticcurveto{\pgfqpoint{2.202137in}{0.854727in}}{\pgfqpoint{2.207540in}{0.858041in}}%
\pgfpathquadraticcurveto{\pgfqpoint{2.212944in}{0.861356in}}{\pgfqpoint{2.220257in}{0.861356in}}%
\pgfpathquadraticcurveto{\pgfqpoint{2.230110in}{0.861356in}}{\pgfqpoint{2.235459in}{0.854457in}}%
\pgfpathquadraticcurveto{\pgfqpoint{2.240809in}{0.847576in}}{\pgfqpoint{2.240809in}{0.834860in}}%
\pgfpathlineto{\pgfqpoint{2.240809in}{0.796800in}}%
\pgfpathlineto{\pgfqpoint{2.230398in}{0.796800in}}%
\pgfpathlineto{\pgfqpoint{2.230398in}{0.834517in}}%
\pgfpathquadraticcurveto{\pgfqpoint{2.230398in}{0.843578in}}{\pgfqpoint{2.227174in}{0.847955in}}%
\pgfpathquadraticcurveto{\pgfqpoint{2.223967in}{0.852349in}}{\pgfqpoint{2.217393in}{0.852349in}}%
\pgfpathquadraticcurveto{\pgfqpoint{2.209342in}{0.852349in}}{\pgfqpoint{2.204658in}{0.847000in}}%
\pgfpathquadraticcurveto{\pgfqpoint{2.199993in}{0.841668in}}{\pgfqpoint{2.199993in}{0.832428in}}%
\pgfpathlineto{\pgfqpoint{2.199993in}{0.796800in}}%
\pgfpathlineto{\pgfqpoint{2.189582in}{0.796800in}}%
\pgfpathlineto{\pgfqpoint{2.189582in}{0.834517in}}%
\pgfpathquadraticcurveto{\pgfqpoint{2.189582in}{0.843632in}}{\pgfqpoint{2.186376in}{0.847991in}}%
\pgfpathquadraticcurveto{\pgfqpoint{2.183170in}{0.852349in}}{\pgfqpoint{2.176469in}{0.852349in}}%
\pgfpathquadraticcurveto{\pgfqpoint{2.168526in}{0.852349in}}{\pgfqpoint{2.163843in}{0.846982in}}%
\pgfpathquadraticcurveto{\pgfqpoint{2.159178in}{0.841614in}}{\pgfqpoint{2.159178in}{0.832428in}}%
\pgfpathlineto{\pgfqpoint{2.159178in}{0.796800in}}%
\pgfpathlineto{\pgfqpoint{2.148767in}{0.796800in}}%
\pgfpathlineto{\pgfqpoint{2.148767in}{0.859843in}}%
\pgfpathlineto{\pgfqpoint{2.159178in}{0.859843in}}%
\pgfpathlineto{\pgfqpoint{2.159178in}{0.850044in}}%
\pgfpathquadraticcurveto{\pgfqpoint{2.162726in}{0.855844in}}{\pgfqpoint{2.167679in}{0.858600in}}%
\pgfpathquadraticcurveto{\pgfqpoint{2.172633in}{0.861356in}}{\pgfqpoint{2.179441in}{0.861356in}}%
\pgfpathquadraticcurveto{\pgfqpoint{2.186322in}{0.861356in}}{\pgfqpoint{2.191131in}{0.857861in}}%
\pgfpathquadraticcurveto{\pgfqpoint{2.195941in}{0.854385in}}{\pgfqpoint{2.198246in}{0.847738in}}%
\pgfpathlineto{\pgfqpoint{2.198246in}{0.847738in}}%
\pgfpathclose%
\pgfpathmoveto{\pgfqpoint{2.260388in}{0.821675in}}%
\pgfpathlineto{\pgfqpoint{2.260388in}{0.859843in}}%
\pgfpathlineto{\pgfqpoint{2.270745in}{0.859843in}}%
\pgfpathlineto{\pgfqpoint{2.270745in}{0.822071in}}%
\pgfpathquadraticcurveto{\pgfqpoint{2.270745in}{0.813119in}}{\pgfqpoint{2.274221in}{0.808634in}}%
\pgfpathquadraticcurveto{\pgfqpoint{2.277716in}{0.804167in}}{\pgfqpoint{2.284704in}{0.804167in}}%
\pgfpathquadraticcurveto{\pgfqpoint{2.293080in}{0.804167in}}{\pgfqpoint{2.297943in}{0.809517in}}%
\pgfpathquadraticcurveto{\pgfqpoint{2.302825in}{0.814866in}}{\pgfqpoint{2.302825in}{0.824106in}}%
\pgfpathlineto{\pgfqpoint{2.302825in}{0.859843in}}%
\pgfpathlineto{\pgfqpoint{2.313182in}{0.859843in}}%
\pgfpathlineto{\pgfqpoint{2.313182in}{0.796800in}}%
\pgfpathlineto{\pgfqpoint{2.302825in}{0.796800in}}%
\pgfpathlineto{\pgfqpoint{2.302825in}{0.806491in}}%
\pgfpathquadraticcurveto{\pgfqpoint{2.299060in}{0.800745in}}{\pgfqpoint{2.294071in}{0.797953in}}%
\pgfpathquadraticcurveto{\pgfqpoint{2.289099in}{0.795161in}}{\pgfqpoint{2.282507in}{0.795161in}}%
\pgfpathquadraticcurveto{\pgfqpoint{2.271646in}{0.795161in}}{\pgfqpoint{2.266008in}{0.801915in}}%
\pgfpathquadraticcurveto{\pgfqpoint{2.260388in}{0.808670in}}{\pgfqpoint{2.260388in}{0.821675in}}%
\pgfpathlineto{\pgfqpoint{2.260388in}{0.821675in}}%
\pgfpathclose%
\pgfpathmoveto{\pgfqpoint{2.286452in}{0.861356in}}%
\pgfpathlineto{\pgfqpoint{2.286452in}{0.861356in}}%
\pgfpathlineto{\pgfqpoint{2.286452in}{0.861356in}}%
\pgfpathclose%
\pgfpathmoveto{\pgfqpoint{2.383591in}{0.847738in}}%
\pgfpathquadraticcurveto{\pgfqpoint{2.387482in}{0.854727in}}{\pgfqpoint{2.392885in}{0.858041in}}%
\pgfpathquadraticcurveto{\pgfqpoint{2.398289in}{0.861356in}}{\pgfqpoint{2.405602in}{0.861356in}}%
\pgfpathquadraticcurveto{\pgfqpoint{2.415455in}{0.861356in}}{\pgfqpoint{2.420804in}{0.854457in}}%
\pgfpathquadraticcurveto{\pgfqpoint{2.426154in}{0.847576in}}{\pgfqpoint{2.426154in}{0.834860in}}%
\pgfpathlineto{\pgfqpoint{2.426154in}{0.796800in}}%
\pgfpathlineto{\pgfqpoint{2.415743in}{0.796800in}}%
\pgfpathlineto{\pgfqpoint{2.415743in}{0.834517in}}%
\pgfpathquadraticcurveto{\pgfqpoint{2.415743in}{0.843578in}}{\pgfqpoint{2.412519in}{0.847955in}}%
\pgfpathquadraticcurveto{\pgfqpoint{2.409312in}{0.852349in}}{\pgfqpoint{2.402738in}{0.852349in}}%
\pgfpathquadraticcurveto{\pgfqpoint{2.394687in}{0.852349in}}{\pgfqpoint{2.390003in}{0.847000in}}%
\pgfpathquadraticcurveto{\pgfqpoint{2.385338in}{0.841668in}}{\pgfqpoint{2.385338in}{0.832428in}}%
\pgfpathlineto{\pgfqpoint{2.385338in}{0.796800in}}%
\pgfpathlineto{\pgfqpoint{2.374927in}{0.796800in}}%
\pgfpathlineto{\pgfqpoint{2.374927in}{0.834517in}}%
\pgfpathquadraticcurveto{\pgfqpoint{2.374927in}{0.843632in}}{\pgfqpoint{2.371721in}{0.847991in}}%
\pgfpathquadraticcurveto{\pgfqpoint{2.368515in}{0.852349in}}{\pgfqpoint{2.361814in}{0.852349in}}%
\pgfpathquadraticcurveto{\pgfqpoint{2.353871in}{0.852349in}}{\pgfqpoint{2.349188in}{0.846982in}}%
\pgfpathquadraticcurveto{\pgfqpoint{2.344523in}{0.841614in}}{\pgfqpoint{2.344523in}{0.832428in}}%
\pgfpathlineto{\pgfqpoint{2.344523in}{0.796800in}}%
\pgfpathlineto{\pgfqpoint{2.334112in}{0.796800in}}%
\pgfpathlineto{\pgfqpoint{2.334112in}{0.859843in}}%
\pgfpathlineto{\pgfqpoint{2.344523in}{0.859843in}}%
\pgfpathlineto{\pgfqpoint{2.344523in}{0.850044in}}%
\pgfpathquadraticcurveto{\pgfqpoint{2.348071in}{0.855844in}}{\pgfqpoint{2.353025in}{0.858600in}}%
\pgfpathquadraticcurveto{\pgfqpoint{2.357978in}{0.861356in}}{\pgfqpoint{2.364786in}{0.861356in}}%
\pgfpathquadraticcurveto{\pgfqpoint{2.371667in}{0.861356in}}{\pgfqpoint{2.376476in}{0.857861in}}%
\pgfpathquadraticcurveto{\pgfqpoint{2.381286in}{0.854385in}}{\pgfqpoint{2.383591in}{0.847738in}}%
\pgfpathlineto{\pgfqpoint{2.383591in}{0.847738in}}%
\pgfpathclose%
\pgfusepath{fill}%
\end{pgfscope}%
\begin{pgfscope}%
\pgfsetbuttcap%
\pgfsetroundjoin%
\definecolor{currentfill}{rgb}{0.309804,0.372549,0.419608}%
\pgfsetfillcolor{currentfill}%
\pgfsetlinewidth{0.000000pt}%
\definecolor{currentstroke}{rgb}{0.309804,0.372549,0.419608}%
\pgfsetstrokecolor{currentstroke}%
\pgfsetdash{}{0pt}%
\pgfpathmoveto{\pgfqpoint{2.575053in}{0.857429in}}%
\pgfpathlineto{\pgfqpoint{2.575053in}{0.847738in}}%
\pgfpathquadraticcurveto{\pgfqpoint{2.570658in}{0.850170in}}{\pgfqpoint{2.566245in}{0.851377in}}%
\pgfpathquadraticcurveto{\pgfqpoint{2.561832in}{0.852584in}}{\pgfqpoint{2.557329in}{0.852584in}}%
\pgfpathquadraticcurveto{\pgfqpoint{2.547243in}{0.852584in}}{\pgfqpoint{2.541659in}{0.846189in}}%
\pgfpathquadraticcurveto{\pgfqpoint{2.536093in}{0.839813in}}{\pgfqpoint{2.536093in}{0.828267in}}%
\pgfpathquadraticcurveto{\pgfqpoint{2.536093in}{0.816721in}}{\pgfqpoint{2.541659in}{0.810327in}}%
\pgfpathquadraticcurveto{\pgfqpoint{2.547243in}{0.803951in}}{\pgfqpoint{2.557329in}{0.803951in}}%
\pgfpathquadraticcurveto{\pgfqpoint{2.561832in}{0.803951in}}{\pgfqpoint{2.566245in}{0.805158in}}%
\pgfpathquadraticcurveto{\pgfqpoint{2.570658in}{0.806364in}}{\pgfqpoint{2.575053in}{0.808796in}}%
\pgfpathlineto{\pgfqpoint{2.575053in}{0.799214in}}%
\pgfpathquadraticcurveto{\pgfqpoint{2.570712in}{0.797196in}}{\pgfqpoint{2.566065in}{0.796188in}}%
\pgfpathquadraticcurveto{\pgfqpoint{2.561436in}{0.795161in}}{\pgfqpoint{2.556195in}{0.795161in}}%
\pgfpathquadraticcurveto{\pgfqpoint{2.541947in}{0.795161in}}{\pgfqpoint{2.533553in}{0.804113in}}%
\pgfpathquadraticcurveto{\pgfqpoint{2.525178in}{0.813065in}}{\pgfqpoint{2.525178in}{0.828267in}}%
\pgfpathquadraticcurveto{\pgfqpoint{2.525178in}{0.843686in}}{\pgfqpoint{2.533643in}{0.852512in}}%
\pgfpathquadraticcurveto{\pgfqpoint{2.542127in}{0.861356in}}{\pgfqpoint{2.556879in}{0.861356in}}%
\pgfpathquadraticcurveto{\pgfqpoint{2.561652in}{0.861356in}}{\pgfqpoint{2.566209in}{0.860365in}}%
\pgfpathquadraticcurveto{\pgfqpoint{2.570766in}{0.859392in}}{\pgfqpoint{2.575053in}{0.857429in}}%
\pgfpathlineto{\pgfqpoint{2.575053in}{0.857429in}}%
\pgfpathclose%
\pgfpathmoveto{\pgfqpoint{2.593066in}{0.884393in}}%
\pgfpathlineto{\pgfqpoint{2.603423in}{0.884393in}}%
\pgfpathlineto{\pgfqpoint{2.603423in}{0.796800in}}%
\pgfpathlineto{\pgfqpoint{2.593066in}{0.796800in}}%
\pgfpathlineto{\pgfqpoint{2.593066in}{0.884393in}}%
\pgfpathlineto{\pgfqpoint{2.593066in}{0.884393in}}%
\pgfpathclose%
\pgfpathmoveto{\pgfqpoint{2.624028in}{0.821675in}}%
\pgfpathlineto{\pgfqpoint{2.624028in}{0.859843in}}%
\pgfpathlineto{\pgfqpoint{2.634385in}{0.859843in}}%
\pgfpathlineto{\pgfqpoint{2.634385in}{0.822071in}}%
\pgfpathquadraticcurveto{\pgfqpoint{2.634385in}{0.813119in}}{\pgfqpoint{2.637862in}{0.808634in}}%
\pgfpathquadraticcurveto{\pgfqpoint{2.641356in}{0.804167in}}{\pgfqpoint{2.648345in}{0.804167in}}%
\pgfpathquadraticcurveto{\pgfqpoint{2.656720in}{0.804167in}}{\pgfqpoint{2.661584in}{0.809517in}}%
\pgfpathquadraticcurveto{\pgfqpoint{2.666465in}{0.814866in}}{\pgfqpoint{2.666465in}{0.824106in}}%
\pgfpathlineto{\pgfqpoint{2.666465in}{0.859843in}}%
\pgfpathlineto{\pgfqpoint{2.676822in}{0.859843in}}%
\pgfpathlineto{\pgfqpoint{2.676822in}{0.796800in}}%
\pgfpathlineto{\pgfqpoint{2.666465in}{0.796800in}}%
\pgfpathlineto{\pgfqpoint{2.666465in}{0.806491in}}%
\pgfpathquadraticcurveto{\pgfqpoint{2.662701in}{0.800745in}}{\pgfqpoint{2.657711in}{0.797953in}}%
\pgfpathquadraticcurveto{\pgfqpoint{2.652740in}{0.795161in}}{\pgfqpoint{2.646147in}{0.795161in}}%
\pgfpathquadraticcurveto{\pgfqpoint{2.635286in}{0.795161in}}{\pgfqpoint{2.629648in}{0.801915in}}%
\pgfpathquadraticcurveto{\pgfqpoint{2.624028in}{0.808670in}}{\pgfqpoint{2.624028in}{0.821675in}}%
\pgfpathlineto{\pgfqpoint{2.624028in}{0.821675in}}%
\pgfpathclose%
\pgfpathmoveto{\pgfqpoint{2.650092in}{0.861356in}}%
\pgfpathlineto{\pgfqpoint{2.650092in}{0.861356in}}%
\pgfpathlineto{\pgfqpoint{2.650092in}{0.861356in}}%
\pgfpathclose%
\pgfpathmoveto{\pgfqpoint{2.738334in}{0.857987in}}%
\pgfpathlineto{\pgfqpoint{2.738334in}{0.848189in}}%
\pgfpathquadraticcurveto{\pgfqpoint{2.733957in}{0.850440in}}{\pgfqpoint{2.729219in}{0.851557in}}%
\pgfpathquadraticcurveto{\pgfqpoint{2.724500in}{0.852692in}}{\pgfqpoint{2.719421in}{0.852692in}}%
\pgfpathquadraticcurveto{\pgfqpoint{2.711712in}{0.852692in}}{\pgfqpoint{2.707857in}{0.850332in}}%
\pgfpathquadraticcurveto{\pgfqpoint{2.704002in}{0.847973in}}{\pgfqpoint{2.704002in}{0.843235in}}%
\pgfpathquadraticcurveto{\pgfqpoint{2.704002in}{0.839633in}}{\pgfqpoint{2.706758in}{0.837580in}}%
\pgfpathquadraticcurveto{\pgfqpoint{2.709514in}{0.835526in}}{\pgfqpoint{2.717854in}{0.833671in}}%
\pgfpathlineto{\pgfqpoint{2.721402in}{0.832878in}}%
\pgfpathquadraticcurveto{\pgfqpoint{2.732426in}{0.830519in}}{\pgfqpoint{2.737073in}{0.826214in}}%
\pgfpathquadraticcurveto{\pgfqpoint{2.741720in}{0.821909in}}{\pgfqpoint{2.741720in}{0.814200in}}%
\pgfpathquadraticcurveto{\pgfqpoint{2.741720in}{0.805410in}}{\pgfqpoint{2.734767in}{0.800276in}}%
\pgfpathquadraticcurveto{\pgfqpoint{2.727814in}{0.795161in}}{\pgfqpoint{2.715656in}{0.795161in}}%
\pgfpathquadraticcurveto{\pgfqpoint{2.710595in}{0.795161in}}{\pgfqpoint{2.705101in}{0.796152in}}%
\pgfpathquadraticcurveto{\pgfqpoint{2.699607in}{0.797142in}}{\pgfqpoint{2.693537in}{0.799106in}}%
\pgfpathlineto{\pgfqpoint{2.693537in}{0.809805in}}%
\pgfpathquadraticcurveto{\pgfqpoint{2.699283in}{0.806815in}}{\pgfqpoint{2.704849in}{0.805320in}}%
\pgfpathquadraticcurveto{\pgfqpoint{2.710415in}{0.803843in}}{\pgfqpoint{2.715890in}{0.803843in}}%
\pgfpathquadraticcurveto{\pgfqpoint{2.723203in}{0.803843in}}{\pgfqpoint{2.727130in}{0.806346in}}%
\pgfpathquadraticcurveto{\pgfqpoint{2.731075in}{0.808850in}}{\pgfqpoint{2.731075in}{0.813407in}}%
\pgfpathquadraticcurveto{\pgfqpoint{2.731075in}{0.817622in}}{\pgfqpoint{2.728229in}{0.819874in}}%
\pgfpathquadraticcurveto{\pgfqpoint{2.725401in}{0.822125in}}{\pgfqpoint{2.715764in}{0.824214in}}%
\pgfpathlineto{\pgfqpoint{2.712162in}{0.825061in}}%
\pgfpathquadraticcurveto{\pgfqpoint{2.702543in}{0.827078in}}{\pgfqpoint{2.698257in}{0.831275in}}%
\pgfpathquadraticcurveto{\pgfqpoint{2.693988in}{0.835472in}}{\pgfqpoint{2.693988in}{0.842785in}}%
\pgfpathquadraticcurveto{\pgfqpoint{2.693988in}{0.851683in}}{\pgfqpoint{2.700292in}{0.856510in}}%
\pgfpathquadraticcurveto{\pgfqpoint{2.706596in}{0.861356in}}{\pgfqpoint{2.718196in}{0.861356in}}%
\pgfpathquadraticcurveto{\pgfqpoint{2.723924in}{0.861356in}}{\pgfqpoint{2.728985in}{0.860509in}}%
\pgfpathquadraticcurveto{\pgfqpoint{2.734065in}{0.859680in}}{\pgfqpoint{2.738334in}{0.857987in}}%
\pgfpathlineto{\pgfqpoint{2.738334in}{0.857987in}}%
\pgfpathclose%
\pgfpathmoveto{\pgfqpoint{2.768450in}{0.877747in}}%
\pgfpathlineto{\pgfqpoint{2.768450in}{0.859843in}}%
\pgfpathlineto{\pgfqpoint{2.789776in}{0.859843in}}%
\pgfpathlineto{\pgfqpoint{2.789776in}{0.851791in}}%
\pgfpathlineto{\pgfqpoint{2.768450in}{0.851791in}}%
\pgfpathlineto{\pgfqpoint{2.768450in}{0.817568in}}%
\pgfpathquadraticcurveto{\pgfqpoint{2.768450in}{0.809859in}}{\pgfqpoint{2.770557in}{0.807661in}}%
\pgfpathquadraticcurveto{\pgfqpoint{2.772665in}{0.805464in}}{\pgfqpoint{2.779149in}{0.805464in}}%
\pgfpathlineto{\pgfqpoint{2.789776in}{0.805464in}}%
\pgfpathlineto{\pgfqpoint{2.789776in}{0.796800in}}%
\pgfpathlineto{\pgfqpoint{2.779149in}{0.796800in}}%
\pgfpathquadraticcurveto{\pgfqpoint{2.767153in}{0.796800in}}{\pgfqpoint{2.762596in}{0.801267in}}%
\pgfpathquadraticcurveto{\pgfqpoint{2.758039in}{0.805752in}}{\pgfqpoint{2.758039in}{0.817568in}}%
\pgfpathlineto{\pgfqpoint{2.758039in}{0.851791in}}%
\pgfpathlineto{\pgfqpoint{2.750438in}{0.851791in}}%
\pgfpathlineto{\pgfqpoint{2.750438in}{0.859843in}}%
\pgfpathlineto{\pgfqpoint{2.758039in}{0.859843in}}%
\pgfpathlineto{\pgfqpoint{2.758039in}{0.877747in}}%
\pgfpathlineto{\pgfqpoint{2.768450in}{0.877747in}}%
\pgfpathlineto{\pgfqpoint{2.768450in}{0.877747in}}%
\pgfpathclose%
\pgfpathmoveto{\pgfqpoint{2.857322in}{0.830915in}}%
\pgfpathlineto{\pgfqpoint{2.857322in}{0.825854in}}%
\pgfpathlineto{\pgfqpoint{2.809698in}{0.825854in}}%
\pgfpathquadraticcurveto{\pgfqpoint{2.810382in}{0.815154in}}{\pgfqpoint{2.816146in}{0.809553in}}%
\pgfpathquadraticcurveto{\pgfqpoint{2.821910in}{0.803951in}}{\pgfqpoint{2.832213in}{0.803951in}}%
\pgfpathquadraticcurveto{\pgfqpoint{2.838175in}{0.803951in}}{\pgfqpoint{2.843777in}{0.805410in}}%
\pgfpathquadraticcurveto{\pgfqpoint{2.849378in}{0.806869in}}{\pgfqpoint{2.854908in}{0.809805in}}%
\pgfpathlineto{\pgfqpoint{2.854908in}{0.800006in}}%
\pgfpathquadraticcurveto{\pgfqpoint{2.849324in}{0.797647in}}{\pgfqpoint{2.843470in}{0.796404in}}%
\pgfpathquadraticcurveto{\pgfqpoint{2.837617in}{0.795161in}}{\pgfqpoint{2.831600in}{0.795161in}}%
\pgfpathquadraticcurveto{\pgfqpoint{2.816506in}{0.795161in}}{\pgfqpoint{2.807698in}{0.803933in}}%
\pgfpathquadraticcurveto{\pgfqpoint{2.798890in}{0.812723in}}{\pgfqpoint{2.798890in}{0.827709in}}%
\pgfpathquadraticcurveto{\pgfqpoint{2.798890in}{0.843181in}}{\pgfqpoint{2.807248in}{0.852259in}}%
\pgfpathquadraticcurveto{\pgfqpoint{2.815606in}{0.861356in}}{\pgfqpoint{2.829799in}{0.861356in}}%
\pgfpathquadraticcurveto{\pgfqpoint{2.842516in}{0.861356in}}{\pgfqpoint{2.849919in}{0.853160in}}%
\pgfpathquadraticcurveto{\pgfqpoint{2.857322in}{0.844983in}}{\pgfqpoint{2.857322in}{0.830915in}}%
\pgfpathlineto{\pgfqpoint{2.857322in}{0.830915in}}%
\pgfpathclose%
\pgfpathmoveto{\pgfqpoint{2.846965in}{0.833959in}}%
\pgfpathquadraticcurveto{\pgfqpoint{2.846857in}{0.842443in}}{\pgfqpoint{2.842210in}{0.847504in}}%
\pgfpathquadraticcurveto{\pgfqpoint{2.837563in}{0.852584in}}{\pgfqpoint{2.829907in}{0.852584in}}%
\pgfpathquadraticcurveto{\pgfqpoint{2.821243in}{0.852584in}}{\pgfqpoint{2.816038in}{0.847684in}}%
\pgfpathquadraticcurveto{\pgfqpoint{2.810832in}{0.842785in}}{\pgfqpoint{2.810040in}{0.833887in}}%
\pgfpathlineto{\pgfqpoint{2.846965in}{0.833959in}}%
\pgfpathlineto{\pgfqpoint{2.846965in}{0.833959in}}%
\pgfpathclose%
\pgfpathmoveto{\pgfqpoint{2.910854in}{0.850170in}}%
\pgfpathquadraticcurveto{\pgfqpoint{2.909107in}{0.851179in}}{\pgfqpoint{2.907053in}{0.851647in}}%
\pgfpathquadraticcurveto{\pgfqpoint{2.905000in}{0.852133in}}{\pgfqpoint{2.902532in}{0.852133in}}%
\pgfpathquadraticcurveto{\pgfqpoint{2.893742in}{0.852133in}}{\pgfqpoint{2.889041in}{0.846423in}}%
\pgfpathquadraticcurveto{\pgfqpoint{2.884340in}{0.840714in}}{\pgfqpoint{2.884340in}{0.830014in}}%
\pgfpathlineto{\pgfqpoint{2.884340in}{0.796800in}}%
\pgfpathlineto{\pgfqpoint{2.873929in}{0.796800in}}%
\pgfpathlineto{\pgfqpoint{2.873929in}{0.859843in}}%
\pgfpathlineto{\pgfqpoint{2.884340in}{0.859843in}}%
\pgfpathlineto{\pgfqpoint{2.884340in}{0.850044in}}%
\pgfpathquadraticcurveto{\pgfqpoint{2.887618in}{0.855790in}}{\pgfqpoint{2.892842in}{0.858564in}}%
\pgfpathquadraticcurveto{\pgfqpoint{2.898083in}{0.861356in}}{\pgfqpoint{2.905576in}{0.861356in}}%
\pgfpathquadraticcurveto{\pgfqpoint{2.906639in}{0.861356in}}{\pgfqpoint{2.907936in}{0.861211in}}%
\pgfpathquadraticcurveto{\pgfqpoint{2.909233in}{0.861085in}}{\pgfqpoint{2.910800in}{0.860797in}}%
\pgfpathlineto{\pgfqpoint{2.910854in}{0.850170in}}%
\pgfpathlineto{\pgfqpoint{2.910854in}{0.850170in}}%
\pgfpathclose%
\pgfusepath{fill}%
\end{pgfscope}%
\begin{pgfscope}%
\pgfsetbuttcap%
\pgfsetroundjoin%
\definecolor{currentfill}{rgb}{0.309804,0.372549,0.419608}%
\pgfsetfillcolor{currentfill}%
\pgfsetlinewidth{0.000000pt}%
\definecolor{currentstroke}{rgb}{0.309804,0.372549,0.419608}%
\pgfsetstrokecolor{currentstroke}%
\pgfsetdash{}{0pt}%
\pgfpathmoveto{\pgfqpoint{2.993386in}{0.884393in}}%
\pgfpathlineto{\pgfqpoint{3.003743in}{0.884393in}}%
\pgfpathlineto{\pgfqpoint{3.003743in}{0.796800in}}%
\pgfpathlineto{\pgfqpoint{2.993386in}{0.796800in}}%
\pgfpathlineto{\pgfqpoint{2.993386in}{0.884393in}}%
\pgfpathlineto{\pgfqpoint{2.993386in}{0.884393in}}%
\pgfpathclose%
\pgfpathmoveto{\pgfqpoint{3.079340in}{0.830915in}}%
\pgfpathlineto{\pgfqpoint{3.079340in}{0.825854in}}%
\pgfpathlineto{\pgfqpoint{3.031715in}{0.825854in}}%
\pgfpathquadraticcurveto{\pgfqpoint{3.032400in}{0.815154in}}{\pgfqpoint{3.038164in}{0.809553in}}%
\pgfpathquadraticcurveto{\pgfqpoint{3.043928in}{0.803951in}}{\pgfqpoint{3.054231in}{0.803951in}}%
\pgfpathquadraticcurveto{\pgfqpoint{3.060193in}{0.803951in}}{\pgfqpoint{3.065794in}{0.805410in}}%
\pgfpathquadraticcurveto{\pgfqpoint{3.071396in}{0.806869in}}{\pgfqpoint{3.076926in}{0.809805in}}%
\pgfpathlineto{\pgfqpoint{3.076926in}{0.800006in}}%
\pgfpathquadraticcurveto{\pgfqpoint{3.071342in}{0.797647in}}{\pgfqpoint{3.065488in}{0.796404in}}%
\pgfpathquadraticcurveto{\pgfqpoint{3.059634in}{0.795161in}}{\pgfqpoint{3.053618in}{0.795161in}}%
\pgfpathquadraticcurveto{\pgfqpoint{3.038524in}{0.795161in}}{\pgfqpoint{3.029716in}{0.803933in}}%
\pgfpathquadraticcurveto{\pgfqpoint{3.020908in}{0.812723in}}{\pgfqpoint{3.020908in}{0.827709in}}%
\pgfpathquadraticcurveto{\pgfqpoint{3.020908in}{0.843181in}}{\pgfqpoint{3.029266in}{0.852259in}}%
\pgfpathquadraticcurveto{\pgfqpoint{3.037623in}{0.861356in}}{\pgfqpoint{3.051817in}{0.861356in}}%
\pgfpathquadraticcurveto{\pgfqpoint{3.064534in}{0.861356in}}{\pgfqpoint{3.071937in}{0.853160in}}%
\pgfpathquadraticcurveto{\pgfqpoint{3.079340in}{0.844983in}}{\pgfqpoint{3.079340in}{0.830915in}}%
\pgfpathlineto{\pgfqpoint{3.079340in}{0.830915in}}%
\pgfpathclose%
\pgfpathmoveto{\pgfqpoint{3.068983in}{0.833959in}}%
\pgfpathquadraticcurveto{\pgfqpoint{3.068875in}{0.842443in}}{\pgfqpoint{3.064227in}{0.847504in}}%
\pgfpathquadraticcurveto{\pgfqpoint{3.059580in}{0.852584in}}{\pgfqpoint{3.051925in}{0.852584in}}%
\pgfpathquadraticcurveto{\pgfqpoint{3.043261in}{0.852584in}}{\pgfqpoint{3.038056in}{0.847684in}}%
\pgfpathquadraticcurveto{\pgfqpoint{3.032850in}{0.842785in}}{\pgfqpoint{3.032058in}{0.833887in}}%
\pgfpathlineto{\pgfqpoint{3.068983in}{0.833959in}}%
\pgfpathlineto{\pgfqpoint{3.068983in}{0.833959in}}%
\pgfpathclose%
\pgfpathmoveto{\pgfqpoint{3.088922in}{0.859843in}}%
\pgfpathlineto{\pgfqpoint{3.099891in}{0.859843in}}%
\pgfpathlineto{\pgfqpoint{3.119597in}{0.806941in}}%
\pgfpathlineto{\pgfqpoint{3.139302in}{0.859843in}}%
\pgfpathlineto{\pgfqpoint{3.150271in}{0.859843in}}%
\pgfpathlineto{\pgfqpoint{3.126622in}{0.796800in}}%
\pgfpathlineto{\pgfqpoint{3.112554in}{0.796800in}}%
\pgfpathlineto{\pgfqpoint{3.088922in}{0.859843in}}%
\pgfpathlineto{\pgfqpoint{3.088922in}{0.859843in}}%
\pgfpathclose%
\pgfpathmoveto{\pgfqpoint{3.218502in}{0.830915in}}%
\pgfpathlineto{\pgfqpoint{3.218502in}{0.825854in}}%
\pgfpathlineto{\pgfqpoint{3.170877in}{0.825854in}}%
\pgfpathquadraticcurveto{\pgfqpoint{3.171562in}{0.815154in}}{\pgfqpoint{3.177326in}{0.809553in}}%
\pgfpathquadraticcurveto{\pgfqpoint{3.183090in}{0.803951in}}{\pgfqpoint{3.193393in}{0.803951in}}%
\pgfpathquadraticcurveto{\pgfqpoint{3.199355in}{0.803951in}}{\pgfqpoint{3.204956in}{0.805410in}}%
\pgfpathquadraticcurveto{\pgfqpoint{3.210558in}{0.806869in}}{\pgfqpoint{3.216088in}{0.809805in}}%
\pgfpathlineto{\pgfqpoint{3.216088in}{0.800006in}}%
\pgfpathquadraticcurveto{\pgfqpoint{3.210504in}{0.797647in}}{\pgfqpoint{3.204650in}{0.796404in}}%
\pgfpathquadraticcurveto{\pgfqpoint{3.198796in}{0.795161in}}{\pgfqpoint{3.192780in}{0.795161in}}%
\pgfpathquadraticcurveto{\pgfqpoint{3.177686in}{0.795161in}}{\pgfqpoint{3.168878in}{0.803933in}}%
\pgfpathquadraticcurveto{\pgfqpoint{3.160070in}{0.812723in}}{\pgfqpoint{3.160070in}{0.827709in}}%
\pgfpathquadraticcurveto{\pgfqpoint{3.160070in}{0.843181in}}{\pgfqpoint{3.168428in}{0.852259in}}%
\pgfpathquadraticcurveto{\pgfqpoint{3.176785in}{0.861356in}}{\pgfqpoint{3.190979in}{0.861356in}}%
\pgfpathquadraticcurveto{\pgfqpoint{3.203696in}{0.861356in}}{\pgfqpoint{3.211099in}{0.853160in}}%
\pgfpathquadraticcurveto{\pgfqpoint{3.218502in}{0.844983in}}{\pgfqpoint{3.218502in}{0.830915in}}%
\pgfpathlineto{\pgfqpoint{3.218502in}{0.830915in}}%
\pgfpathclose%
\pgfpathmoveto{\pgfqpoint{3.208145in}{0.833959in}}%
\pgfpathquadraticcurveto{\pgfqpoint{3.208036in}{0.842443in}}{\pgfqpoint{3.203389in}{0.847504in}}%
\pgfpathquadraticcurveto{\pgfqpoint{3.198742in}{0.852584in}}{\pgfqpoint{3.191087in}{0.852584in}}%
\pgfpathquadraticcurveto{\pgfqpoint{3.182423in}{0.852584in}}{\pgfqpoint{3.177218in}{0.847684in}}%
\pgfpathquadraticcurveto{\pgfqpoint{3.172012in}{0.842785in}}{\pgfqpoint{3.171220in}{0.833887in}}%
\pgfpathlineto{\pgfqpoint{3.208145in}{0.833959in}}%
\pgfpathlineto{\pgfqpoint{3.208145in}{0.833959in}}%
\pgfpathclose%
\pgfpathmoveto{\pgfqpoint{3.272034in}{0.850170in}}%
\pgfpathquadraticcurveto{\pgfqpoint{3.270286in}{0.851179in}}{\pgfqpoint{3.268233in}{0.851647in}}%
\pgfpathquadraticcurveto{\pgfqpoint{3.266180in}{0.852133in}}{\pgfqpoint{3.263712in}{0.852133in}}%
\pgfpathquadraticcurveto{\pgfqpoint{3.254922in}{0.852133in}}{\pgfqpoint{3.250221in}{0.846423in}}%
\pgfpathquadraticcurveto{\pgfqpoint{3.245520in}{0.840714in}}{\pgfqpoint{3.245520in}{0.830014in}}%
\pgfpathlineto{\pgfqpoint{3.245520in}{0.796800in}}%
\pgfpathlineto{\pgfqpoint{3.235109in}{0.796800in}}%
\pgfpathlineto{\pgfqpoint{3.235109in}{0.859843in}}%
\pgfpathlineto{\pgfqpoint{3.245520in}{0.859843in}}%
\pgfpathlineto{\pgfqpoint{3.245520in}{0.850044in}}%
\pgfpathquadraticcurveto{\pgfqpoint{3.248798in}{0.855790in}}{\pgfqpoint{3.254021in}{0.858564in}}%
\pgfpathquadraticcurveto{\pgfqpoint{3.259263in}{0.861356in}}{\pgfqpoint{3.266756in}{0.861356in}}%
\pgfpathquadraticcurveto{\pgfqpoint{3.267819in}{0.861356in}}{\pgfqpoint{3.269116in}{0.861211in}}%
\pgfpathquadraticcurveto{\pgfqpoint{3.270413in}{0.861085in}}{\pgfqpoint{3.271980in}{0.860797in}}%
\pgfpathlineto{\pgfqpoint{3.272034in}{0.850170in}}%
\pgfpathlineto{\pgfqpoint{3.272034in}{0.850170in}}%
\pgfpathclose%
\pgfpathmoveto{\pgfqpoint{3.311552in}{0.828483in}}%
\pgfpathquadraticcurveto{\pgfqpoint{3.298998in}{0.828483in}}{\pgfqpoint{3.294153in}{0.825619in}}%
\pgfpathquadraticcurveto{\pgfqpoint{3.289307in}{0.822756in}}{\pgfqpoint{3.289307in}{0.815821in}}%
\pgfpathquadraticcurveto{\pgfqpoint{3.289307in}{0.810309in}}{\pgfqpoint{3.292946in}{0.807067in}}%
\pgfpathquadraticcurveto{\pgfqpoint{3.296584in}{0.803843in}}{\pgfqpoint{3.302816in}{0.803843in}}%
\pgfpathquadraticcurveto{\pgfqpoint{3.311444in}{0.803843in}}{\pgfqpoint{3.316650in}{0.809949in}}%
\pgfpathquadraticcurveto{\pgfqpoint{3.321855in}{0.816055in}}{\pgfqpoint{3.321855in}{0.826178in}}%
\pgfpathlineto{\pgfqpoint{3.321855in}{0.828483in}}%
\pgfpathlineto{\pgfqpoint{3.311552in}{0.828483in}}%
\pgfpathlineto{\pgfqpoint{3.311552in}{0.828483in}}%
\pgfpathclose%
\pgfpathmoveto{\pgfqpoint{3.332212in}{0.832770in}}%
\pgfpathlineto{\pgfqpoint{3.332212in}{0.796800in}}%
\pgfpathlineto{\pgfqpoint{3.321855in}{0.796800in}}%
\pgfpathlineto{\pgfqpoint{3.321855in}{0.806364in}}%
\pgfpathquadraticcurveto{\pgfqpoint{3.318307in}{0.800637in}}{\pgfqpoint{3.313011in}{0.797899in}}%
\pgfpathquadraticcurveto{\pgfqpoint{3.307716in}{0.795161in}}{\pgfqpoint{3.300061in}{0.795161in}}%
\pgfpathquadraticcurveto{\pgfqpoint{3.290388in}{0.795161in}}{\pgfqpoint{3.284660in}{0.800601in}}%
\pgfpathquadraticcurveto{\pgfqpoint{3.278950in}{0.806040in}}{\pgfqpoint{3.278950in}{0.815154in}}%
\pgfpathquadraticcurveto{\pgfqpoint{3.278950in}{0.825782in}}{\pgfqpoint{3.286065in}{0.831185in}}%
\pgfpathquadraticcurveto{\pgfqpoint{3.293198in}{0.836589in}}{\pgfqpoint{3.307319in}{0.836589in}}%
\pgfpathlineto{\pgfqpoint{3.321855in}{0.836589in}}%
\pgfpathlineto{\pgfqpoint{3.321855in}{0.837616in}}%
\pgfpathquadraticcurveto{\pgfqpoint{3.321855in}{0.844766in}}{\pgfqpoint{3.317154in}{0.848675in}}%
\pgfpathquadraticcurveto{\pgfqpoint{3.312453in}{0.852584in}}{\pgfqpoint{3.303951in}{0.852584in}}%
\pgfpathquadraticcurveto{\pgfqpoint{3.298548in}{0.852584in}}{\pgfqpoint{3.293414in}{0.851287in}}%
\pgfpathquadraticcurveto{\pgfqpoint{3.288299in}{0.849990in}}{\pgfqpoint{3.283579in}{0.847396in}}%
\pgfpathlineto{\pgfqpoint{3.283579in}{0.856979in}}%
\pgfpathquadraticcurveto{\pgfqpoint{3.289253in}{0.859176in}}{\pgfqpoint{3.294603in}{0.860257in}}%
\pgfpathquadraticcurveto{\pgfqpoint{3.299952in}{0.861356in}}{\pgfqpoint{3.305014in}{0.861356in}}%
\pgfpathquadraticcurveto{\pgfqpoint{3.318703in}{0.861356in}}{\pgfqpoint{3.325458in}{0.854259in}}%
\pgfpathquadraticcurveto{\pgfqpoint{3.332212in}{0.847180in}}{\pgfqpoint{3.332212in}{0.832770in}}%
\pgfpathlineto{\pgfqpoint{3.332212in}{0.832770in}}%
\pgfpathclose%
\pgfpathmoveto{\pgfqpoint{3.395021in}{0.829060in}}%
\pgfpathquadraticcurveto{\pgfqpoint{3.395021in}{0.840317in}}{\pgfqpoint{3.390373in}{0.846496in}}%
\pgfpathquadraticcurveto{\pgfqpoint{3.385744in}{0.852692in}}{\pgfqpoint{3.377351in}{0.852692in}}%
\pgfpathquadraticcurveto{\pgfqpoint{3.369029in}{0.852692in}}{\pgfqpoint{3.364382in}{0.846496in}}%
\pgfpathquadraticcurveto{\pgfqpoint{3.359735in}{0.840317in}}{\pgfqpoint{3.359735in}{0.829060in}}%
\pgfpathquadraticcurveto{\pgfqpoint{3.359735in}{0.817856in}}{\pgfqpoint{3.364382in}{0.811660in}}%
\pgfpathquadraticcurveto{\pgfqpoint{3.369029in}{0.805464in}}{\pgfqpoint{3.377351in}{0.805464in}}%
\pgfpathquadraticcurveto{\pgfqpoint{3.385744in}{0.805464in}}{\pgfqpoint{3.390373in}{0.811660in}}%
\pgfpathquadraticcurveto{\pgfqpoint{3.395021in}{0.817856in}}{\pgfqpoint{3.395021in}{0.829060in}}%
\pgfpathlineto{\pgfqpoint{3.395021in}{0.829060in}}%
\pgfpathclose%
\pgfpathmoveto{\pgfqpoint{3.405378in}{0.804617in}}%
\pgfpathquadraticcurveto{\pgfqpoint{3.405378in}{0.788532in}}{\pgfqpoint{3.398227in}{0.780679in}}%
\pgfpathquadraticcurveto{\pgfqpoint{3.391094in}{0.772826in}}{\pgfqpoint{3.376342in}{0.772826in}}%
\pgfpathquadraticcurveto{\pgfqpoint{3.370884in}{0.772826in}}{\pgfqpoint{3.366039in}{0.773636in}}%
\pgfpathquadraticcurveto{\pgfqpoint{3.361194in}{0.774447in}}{\pgfqpoint{3.356637in}{0.776140in}}%
\pgfpathlineto{\pgfqpoint{3.356637in}{0.786209in}}%
\pgfpathquadraticcurveto{\pgfqpoint{3.361194in}{0.783741in}}{\pgfqpoint{3.365643in}{0.782570in}}%
\pgfpathquadraticcurveto{\pgfqpoint{3.370092in}{0.781382in}}{\pgfqpoint{3.374703in}{0.781382in}}%
\pgfpathquadraticcurveto{\pgfqpoint{3.384898in}{0.781382in}}{\pgfqpoint{3.389959in}{0.786695in}}%
\pgfpathquadraticcurveto{\pgfqpoint{3.395021in}{0.792009in}}{\pgfqpoint{3.395021in}{0.802762in}}%
\pgfpathlineto{\pgfqpoint{3.395021in}{0.807895in}}%
\pgfpathquadraticcurveto{\pgfqpoint{3.391814in}{0.802312in}}{\pgfqpoint{3.386807in}{0.799556in}}%
\pgfpathquadraticcurveto{\pgfqpoint{3.381800in}{0.796800in}}{\pgfqpoint{3.374811in}{0.796800in}}%
\pgfpathquadraticcurveto{\pgfqpoint{3.363229in}{0.796800in}}{\pgfqpoint{3.356132in}{0.805626in}}%
\pgfpathquadraticcurveto{\pgfqpoint{3.349036in}{0.814470in}}{\pgfqpoint{3.349036in}{0.829060in}}%
\pgfpathquadraticcurveto{\pgfqpoint{3.349036in}{0.843686in}}{\pgfqpoint{3.356132in}{0.852512in}}%
\pgfpathquadraticcurveto{\pgfqpoint{3.363229in}{0.861356in}}{\pgfqpoint{3.374811in}{0.861356in}}%
\pgfpathquadraticcurveto{\pgfqpoint{3.381800in}{0.861356in}}{\pgfqpoint{3.386807in}{0.858600in}}%
\pgfpathquadraticcurveto{\pgfqpoint{3.391814in}{0.855844in}}{\pgfqpoint{3.395021in}{0.850278in}}%
\pgfpathlineto{\pgfqpoint{3.395021in}{0.859843in}}%
\pgfpathlineto{\pgfqpoint{3.405378in}{0.859843in}}%
\pgfpathlineto{\pgfqpoint{3.405378in}{0.804617in}}%
\pgfpathlineto{\pgfqpoint{3.405378in}{0.804617in}}%
\pgfpathclose%
\pgfpathmoveto{\pgfqpoint{3.480650in}{0.830915in}}%
\pgfpathlineto{\pgfqpoint{3.480650in}{0.825854in}}%
\pgfpathlineto{\pgfqpoint{3.433026in}{0.825854in}}%
\pgfpathquadraticcurveto{\pgfqpoint{3.433711in}{0.815154in}}{\pgfqpoint{3.439475in}{0.809553in}}%
\pgfpathquadraticcurveto{\pgfqpoint{3.445238in}{0.803951in}}{\pgfqpoint{3.455541in}{0.803951in}}%
\pgfpathquadraticcurveto{\pgfqpoint{3.461503in}{0.803951in}}{\pgfqpoint{3.467105in}{0.805410in}}%
\pgfpathquadraticcurveto{\pgfqpoint{3.472707in}{0.806869in}}{\pgfqpoint{3.478237in}{0.809805in}}%
\pgfpathlineto{\pgfqpoint{3.478237in}{0.800006in}}%
\pgfpathquadraticcurveto{\pgfqpoint{3.472653in}{0.797647in}}{\pgfqpoint{3.466799in}{0.796404in}}%
\pgfpathquadraticcurveto{\pgfqpoint{3.460945in}{0.795161in}}{\pgfqpoint{3.454929in}{0.795161in}}%
\pgfpathquadraticcurveto{\pgfqpoint{3.439835in}{0.795161in}}{\pgfqpoint{3.431027in}{0.803933in}}%
\pgfpathquadraticcurveto{\pgfqpoint{3.422219in}{0.812723in}}{\pgfqpoint{3.422219in}{0.827709in}}%
\pgfpathquadraticcurveto{\pgfqpoint{3.422219in}{0.843181in}}{\pgfqpoint{3.430577in}{0.852259in}}%
\pgfpathquadraticcurveto{\pgfqpoint{3.438934in}{0.861356in}}{\pgfqpoint{3.453128in}{0.861356in}}%
\pgfpathquadraticcurveto{\pgfqpoint{3.465844in}{0.861356in}}{\pgfqpoint{3.473247in}{0.853160in}}%
\pgfpathquadraticcurveto{\pgfqpoint{3.480650in}{0.844983in}}{\pgfqpoint{3.480650in}{0.830915in}}%
\pgfpathlineto{\pgfqpoint{3.480650in}{0.830915in}}%
\pgfpathclose%
\pgfpathmoveto{\pgfqpoint{3.470293in}{0.833959in}}%
\pgfpathquadraticcurveto{\pgfqpoint{3.470185in}{0.842443in}}{\pgfqpoint{3.465538in}{0.847504in}}%
\pgfpathquadraticcurveto{\pgfqpoint{3.460891in}{0.852584in}}{\pgfqpoint{3.453236in}{0.852584in}}%
\pgfpathquadraticcurveto{\pgfqpoint{3.444572in}{0.852584in}}{\pgfqpoint{3.439367in}{0.847684in}}%
\pgfpathquadraticcurveto{\pgfqpoint{3.434161in}{0.842785in}}{\pgfqpoint{3.433368in}{0.833887in}}%
\pgfpathlineto{\pgfqpoint{3.470293in}{0.833959in}}%
\pgfpathlineto{\pgfqpoint{3.470293in}{0.833959in}}%
\pgfpathclose%
\pgfusepath{fill}%
\end{pgfscope}%
\begin{pgfscope}%
\pgfsetbuttcap%
\pgfsetroundjoin%
\definecolor{currentfill}{rgb}{0.309804,0.372549,0.419608}%
\pgfsetfillcolor{currentfill}%
\pgfsetlinewidth{0.000000pt}%
\definecolor{currentstroke}{rgb}{0.309804,0.372549,0.419608}%
\pgfsetstrokecolor{currentstroke}%
\pgfsetdash{}{0pt}%
\pgfpathmoveto{\pgfqpoint{3.558441in}{0.853538in}}%
\pgfpathlineto{\pgfqpoint{3.558441in}{0.863787in}}%
\pgfpathlineto{\pgfqpoint{3.630598in}{0.837616in}}%
\pgfpathlineto{\pgfqpoint{3.630598in}{0.828267in}}%
\pgfpathlineto{\pgfqpoint{3.558441in}{0.802096in}}%
\pgfpathlineto{\pgfqpoint{3.558441in}{0.812344in}}%
\pgfpathlineto{\pgfqpoint{3.616423in}{0.832878in}}%
\pgfpathlineto{\pgfqpoint{3.558441in}{0.853538in}}%
\pgfpathlineto{\pgfqpoint{3.558441in}{0.853538in}}%
\pgfpathclose%
\pgfpathmoveto{\pgfqpoint{3.655275in}{0.880845in}}%
\pgfpathlineto{\pgfqpoint{3.699909in}{0.880845in}}%
\pgfpathlineto{\pgfqpoint{3.699909in}{0.871262in}}%
\pgfpathlineto{\pgfqpoint{3.665686in}{0.871262in}}%
\pgfpathlineto{\pgfqpoint{3.665686in}{0.850674in}}%
\pgfpathquadraticcurveto{\pgfqpoint{3.668153in}{0.851521in}}{\pgfqpoint{3.670621in}{0.851935in}}%
\pgfpathquadraticcurveto{\pgfqpoint{3.673107in}{0.852349in}}{\pgfqpoint{3.675592in}{0.852349in}}%
\pgfpathquadraticcurveto{\pgfqpoint{3.689660in}{0.852349in}}{\pgfqpoint{3.697873in}{0.844640in}}%
\pgfpathquadraticcurveto{\pgfqpoint{3.706105in}{0.836931in}}{\pgfqpoint{3.706105in}{0.823764in}}%
\pgfpathquadraticcurveto{\pgfqpoint{3.706105in}{0.810201in}}{\pgfqpoint{3.697657in}{0.802672in}}%
\pgfpathquadraticcurveto{\pgfqpoint{3.689210in}{0.795161in}}{\pgfqpoint{3.673845in}{0.795161in}}%
\pgfpathquadraticcurveto{\pgfqpoint{3.668550in}{0.795161in}}{\pgfqpoint{3.663056in}{0.796062in}}%
\pgfpathquadraticcurveto{\pgfqpoint{3.657580in}{0.796962in}}{\pgfqpoint{3.651726in}{0.798763in}}%
\pgfpathlineto{\pgfqpoint{3.651726in}{0.810201in}}%
\pgfpathquadraticcurveto{\pgfqpoint{3.656788in}{0.807445in}}{\pgfqpoint{3.662191in}{0.806094in}}%
\pgfpathquadraticcurveto{\pgfqpoint{3.667595in}{0.804743in}}{\pgfqpoint{3.673611in}{0.804743in}}%
\pgfpathquadraticcurveto{\pgfqpoint{3.683356in}{0.804743in}}{\pgfqpoint{3.689030in}{0.809859in}}%
\pgfpathquadraticcurveto{\pgfqpoint{3.694721in}{0.814974in}}{\pgfqpoint{3.694721in}{0.823764in}}%
\pgfpathquadraticcurveto{\pgfqpoint{3.694721in}{0.832536in}}{\pgfqpoint{3.689030in}{0.837652in}}%
\pgfpathquadraticcurveto{\pgfqpoint{3.683356in}{0.842785in}}{\pgfqpoint{3.673611in}{0.842785in}}%
\pgfpathquadraticcurveto{\pgfqpoint{3.669054in}{0.842785in}}{\pgfqpoint{3.664515in}{0.841776in}}%
\pgfpathquadraticcurveto{\pgfqpoint{3.659994in}{0.840768in}}{\pgfqpoint{3.655275in}{0.838624in}}%
\pgfpathlineto{\pgfqpoint{3.655275in}{0.880845in}}%
\pgfpathlineto{\pgfqpoint{3.655275in}{0.880845in}}%
\pgfpathclose%
\pgfpathmoveto{\pgfqpoint{3.752811in}{0.873352in}}%
\pgfpathquadraticcurveto{\pgfqpoint{3.744039in}{0.873352in}}{\pgfqpoint{3.739608in}{0.864706in}}%
\pgfpathquadraticcurveto{\pgfqpoint{3.735195in}{0.856078in}}{\pgfqpoint{3.735195in}{0.838732in}}%
\pgfpathquadraticcurveto{\pgfqpoint{3.735195in}{0.821459in}}{\pgfqpoint{3.739608in}{0.812813in}}%
\pgfpathquadraticcurveto{\pgfqpoint{3.744039in}{0.804167in}}{\pgfqpoint{3.752811in}{0.804167in}}%
\pgfpathquadraticcurveto{\pgfqpoint{3.761655in}{0.804167in}}{\pgfqpoint{3.766067in}{0.812813in}}%
\pgfpathquadraticcurveto{\pgfqpoint{3.770498in}{0.821459in}}{\pgfqpoint{3.770498in}{0.838732in}}%
\pgfpathquadraticcurveto{\pgfqpoint{3.770498in}{0.856078in}}{\pgfqpoint{3.766067in}{0.864706in}}%
\pgfpathquadraticcurveto{\pgfqpoint{3.761655in}{0.873352in}}{\pgfqpoint{3.752811in}{0.873352in}}%
\pgfpathlineto{\pgfqpoint{3.752811in}{0.873352in}}%
\pgfpathclose%
\pgfpathmoveto{\pgfqpoint{3.752811in}{0.882358in}}%
\pgfpathquadraticcurveto{\pgfqpoint{3.766950in}{0.882358in}}{\pgfqpoint{3.774407in}{0.871172in}}%
\pgfpathquadraticcurveto{\pgfqpoint{3.781864in}{0.860005in}}{\pgfqpoint{3.781864in}{0.838732in}}%
\pgfpathquadraticcurveto{\pgfqpoint{3.781864in}{0.817514in}}{\pgfqpoint{3.774407in}{0.806328in}}%
\pgfpathquadraticcurveto{\pgfqpoint{3.766950in}{0.795161in}}{\pgfqpoint{3.752811in}{0.795161in}}%
\pgfpathquadraticcurveto{\pgfqpoint{3.738689in}{0.795161in}}{\pgfqpoint{3.731232in}{0.806328in}}%
\pgfpathquadraticcurveto{\pgfqpoint{3.723775in}{0.817514in}}{\pgfqpoint{3.723775in}{0.838732in}}%
\pgfpathquadraticcurveto{\pgfqpoint{3.723775in}{0.860005in}}{\pgfqpoint{3.731232in}{0.871172in}}%
\pgfpathquadraticcurveto{\pgfqpoint{3.738689in}{0.882358in}}{\pgfqpoint{3.752811in}{0.882358in}}%
\pgfpathlineto{\pgfqpoint{3.752811in}{0.882358in}}%
\pgfpathclose%
\pgfpathmoveto{\pgfqpoint{3.873330in}{0.833779in}}%
\pgfpathquadraticcurveto{\pgfqpoint{3.868431in}{0.833779in}}{\pgfqpoint{3.865639in}{0.829618in}}%
\pgfpathquadraticcurveto{\pgfqpoint{3.862865in}{0.825457in}}{\pgfqpoint{3.862865in}{0.818018in}}%
\pgfpathquadraticcurveto{\pgfqpoint{3.862865in}{0.810705in}}{\pgfqpoint{3.865639in}{0.806509in}}%
\pgfpathquadraticcurveto{\pgfqpoint{3.868431in}{0.802312in}}{\pgfqpoint{3.873330in}{0.802312in}}%
\pgfpathquadraticcurveto{\pgfqpoint{3.878121in}{0.802312in}}{\pgfqpoint{3.880895in}{0.806509in}}%
\pgfpathquadraticcurveto{\pgfqpoint{3.883687in}{0.810705in}}{\pgfqpoint{3.883687in}{0.818018in}}%
\pgfpathquadraticcurveto{\pgfqpoint{3.883687in}{0.825403in}}{\pgfqpoint{3.880895in}{0.829582in}}%
\pgfpathquadraticcurveto{\pgfqpoint{3.878121in}{0.833779in}}{\pgfqpoint{3.873330in}{0.833779in}}%
\pgfpathlineto{\pgfqpoint{3.873330in}{0.833779in}}%
\pgfpathclose%
\pgfpathmoveto{\pgfqpoint{3.873330in}{0.840930in}}%
\pgfpathquadraticcurveto{\pgfqpoint{3.882228in}{0.840930in}}{\pgfqpoint{3.887451in}{0.834734in}}%
\pgfpathquadraticcurveto{\pgfqpoint{3.892693in}{0.828555in}}{\pgfqpoint{3.892693in}{0.818018in}}%
\pgfpathquadraticcurveto{\pgfqpoint{3.892693in}{0.807499in}}{\pgfqpoint{3.887433in}{0.801321in}}%
\pgfpathquadraticcurveto{\pgfqpoint{3.882174in}{0.795161in}}{\pgfqpoint{3.873330in}{0.795161in}}%
\pgfpathquadraticcurveto{\pgfqpoint{3.864324in}{0.795161in}}{\pgfqpoint{3.859082in}{0.801321in}}%
\pgfpathquadraticcurveto{\pgfqpoint{3.853859in}{0.807499in}}{\pgfqpoint{3.853859in}{0.818018in}}%
\pgfpathquadraticcurveto{\pgfqpoint{3.853859in}{0.828609in}}{\pgfqpoint{3.859118in}{0.834770in}}%
\pgfpathquadraticcurveto{\pgfqpoint{3.864378in}{0.840930in}}{\pgfqpoint{3.873330in}{0.840930in}}%
\pgfpathlineto{\pgfqpoint{3.873330in}{0.840930in}}%
\pgfpathclose%
\pgfpathmoveto{\pgfqpoint{3.815241in}{0.875207in}}%
\pgfpathquadraticcurveto{\pgfqpoint{3.810395in}{0.875207in}}{\pgfqpoint{3.807604in}{0.871010in}}%
\pgfpathquadraticcurveto{\pgfqpoint{3.804830in}{0.866831in}}{\pgfqpoint{3.804830in}{0.859500in}}%
\pgfpathquadraticcurveto{\pgfqpoint{3.804830in}{0.852079in}}{\pgfqpoint{3.807586in}{0.847900in}}%
\pgfpathquadraticcurveto{\pgfqpoint{3.810341in}{0.843740in}}{\pgfqpoint{3.815241in}{0.843740in}}%
\pgfpathquadraticcurveto{\pgfqpoint{3.820140in}{0.843740in}}{\pgfqpoint{3.822914in}{0.847900in}}%
\pgfpathquadraticcurveto{\pgfqpoint{3.825706in}{0.852079in}}{\pgfqpoint{3.825706in}{0.859500in}}%
\pgfpathquadraticcurveto{\pgfqpoint{3.825706in}{0.866759in}}{\pgfqpoint{3.822896in}{0.870974in}}%
\pgfpathquadraticcurveto{\pgfqpoint{3.820086in}{0.875207in}}{\pgfqpoint{3.815241in}{0.875207in}}%
\pgfpathlineto{\pgfqpoint{3.815241in}{0.875207in}}%
\pgfpathclose%
\pgfpathmoveto{\pgfqpoint{3.866071in}{0.882358in}}%
\pgfpathlineto{\pgfqpoint{3.875077in}{0.882358in}}%
\pgfpathlineto{\pgfqpoint{3.822500in}{0.795161in}}%
\pgfpathlineto{\pgfqpoint{3.813493in}{0.795161in}}%
\pgfpathlineto{\pgfqpoint{3.866071in}{0.882358in}}%
\pgfpathlineto{\pgfqpoint{3.866071in}{0.882358in}}%
\pgfpathclose%
\pgfpathmoveto{\pgfqpoint{3.815241in}{0.882358in}}%
\pgfpathquadraticcurveto{\pgfqpoint{3.824139in}{0.882358in}}{\pgfqpoint{3.829416in}{0.876198in}}%
\pgfpathquadraticcurveto{\pgfqpoint{3.834712in}{0.870037in}}{\pgfqpoint{3.834712in}{0.859500in}}%
\pgfpathquadraticcurveto{\pgfqpoint{3.834712in}{0.848873in}}{\pgfqpoint{3.829452in}{0.842731in}}%
\pgfpathquadraticcurveto{\pgfqpoint{3.824193in}{0.836589in}}{\pgfqpoint{3.815241in}{0.836589in}}%
\pgfpathquadraticcurveto{\pgfqpoint{3.806289in}{0.836589in}}{\pgfqpoint{3.801083in}{0.842749in}}%
\pgfpathquadraticcurveto{\pgfqpoint{3.795878in}{0.848927in}}{\pgfqpoint{3.795878in}{0.859500in}}%
\pgfpathquadraticcurveto{\pgfqpoint{3.795878in}{0.869983in}}{\pgfqpoint{3.801101in}{0.876162in}}%
\pgfpathquadraticcurveto{\pgfqpoint{3.806343in}{0.882358in}}{\pgfqpoint{3.815241in}{0.882358in}}%
\pgfpathlineto{\pgfqpoint{3.815241in}{0.882358in}}%
\pgfpathclose%
\pgfpathmoveto{\pgfqpoint{3.911372in}{0.811102in}}%
\pgfpathlineto{\pgfqpoint{3.923260in}{0.811102in}}%
\pgfpathlineto{\pgfqpoint{3.923260in}{0.796800in}}%
\pgfpathlineto{\pgfqpoint{3.911372in}{0.796800in}}%
\pgfpathlineto{\pgfqpoint{3.911372in}{0.811102in}}%
\pgfpathlineto{\pgfqpoint{3.911372in}{0.811102in}}%
\pgfpathclose%
\pgfusepath{fill}%
\end{pgfscope}%
\begin{pgfscope}%
\pgfsetbuttcap%
\pgfsetroundjoin%
\definecolor{currentfill}{rgb}{0.309804,0.372549,0.419608}%
\pgfsetfillcolor{currentfill}%
\pgfsetlinewidth{0.000000pt}%
\definecolor{currentstroke}{rgb}{0.309804,0.372549,0.419608}%
\pgfsetstrokecolor{currentstroke}%
\pgfsetdash{}{0pt}%
\pgfpathmoveto{\pgfqpoint{4.084180in}{0.874360in}}%
\pgfpathlineto{\pgfqpoint{4.084180in}{0.862382in}}%
\pgfpathquadraticcurveto{\pgfqpoint{4.078434in}{0.867732in}}{\pgfqpoint{4.071932in}{0.870362in}}%
\pgfpathquadraticcurveto{\pgfqpoint{4.065429in}{0.873009in}}{\pgfqpoint{4.058117in}{0.873009in}}%
\pgfpathquadraticcurveto{\pgfqpoint{4.043707in}{0.873009in}}{\pgfqpoint{4.036052in}{0.864201in}}%
\pgfpathquadraticcurveto{\pgfqpoint{4.028396in}{0.855394in}}{\pgfqpoint{4.028396in}{0.838732in}}%
\pgfpathquadraticcurveto{\pgfqpoint{4.028396in}{0.822125in}}{\pgfqpoint{4.036052in}{0.813317in}}%
\pgfpathquadraticcurveto{\pgfqpoint{4.043707in}{0.804509in}}{\pgfqpoint{4.058117in}{0.804509in}}%
\pgfpathquadraticcurveto{\pgfqpoint{4.065429in}{0.804509in}}{\pgfqpoint{4.071932in}{0.807157in}}%
\pgfpathquadraticcurveto{\pgfqpoint{4.078434in}{0.809805in}}{\pgfqpoint{4.084180in}{0.815154in}}%
\pgfpathlineto{\pgfqpoint{4.084180in}{0.803266in}}%
\pgfpathquadraticcurveto{\pgfqpoint{4.078218in}{0.799214in}}{\pgfqpoint{4.071536in}{0.797178in}}%
\pgfpathquadraticcurveto{\pgfqpoint{4.064871in}{0.795161in}}{\pgfqpoint{4.057450in}{0.795161in}}%
\pgfpathquadraticcurveto{\pgfqpoint{4.038357in}{0.795161in}}{\pgfqpoint{4.027370in}{0.806833in}}%
\pgfpathquadraticcurveto{\pgfqpoint{4.016400in}{0.818523in}}{\pgfqpoint{4.016400in}{0.838732in}}%
\pgfpathquadraticcurveto{\pgfqpoint{4.016400in}{0.858996in}}{\pgfqpoint{4.027370in}{0.870668in}}%
\pgfpathquadraticcurveto{\pgfqpoint{4.038357in}{0.882358in}}{\pgfqpoint{4.057450in}{0.882358in}}%
\pgfpathquadraticcurveto{\pgfqpoint{4.064979in}{0.882358in}}{\pgfqpoint{4.071644in}{0.880358in}}%
\pgfpathquadraticcurveto{\pgfqpoint{4.078326in}{0.878359in}}{\pgfqpoint{4.084180in}{0.874360in}}%
\pgfpathlineto{\pgfqpoint{4.084180in}{0.874360in}}%
\pgfpathclose%
\pgfpathmoveto{\pgfqpoint{4.137820in}{0.850170in}}%
\pgfpathquadraticcurveto{\pgfqpoint{4.136073in}{0.851179in}}{\pgfqpoint{4.134020in}{0.851647in}}%
\pgfpathquadraticcurveto{\pgfqpoint{4.131966in}{0.852133in}}{\pgfqpoint{4.129499in}{0.852133in}}%
\pgfpathquadraticcurveto{\pgfqpoint{4.120709in}{0.852133in}}{\pgfqpoint{4.116008in}{0.846423in}}%
\pgfpathquadraticcurveto{\pgfqpoint{4.111306in}{0.840714in}}{\pgfqpoint{4.111306in}{0.830014in}}%
\pgfpathlineto{\pgfqpoint{4.111306in}{0.796800in}}%
\pgfpathlineto{\pgfqpoint{4.100895in}{0.796800in}}%
\pgfpathlineto{\pgfqpoint{4.100895in}{0.859843in}}%
\pgfpathlineto{\pgfqpoint{4.111306in}{0.859843in}}%
\pgfpathlineto{\pgfqpoint{4.111306in}{0.850044in}}%
\pgfpathquadraticcurveto{\pgfqpoint{4.114585in}{0.855790in}}{\pgfqpoint{4.119808in}{0.858564in}}%
\pgfpathquadraticcurveto{\pgfqpoint{4.125050in}{0.861356in}}{\pgfqpoint{4.132543in}{0.861356in}}%
\pgfpathquadraticcurveto{\pgfqpoint{4.133605in}{0.861356in}}{\pgfqpoint{4.134902in}{0.861211in}}%
\pgfpathquadraticcurveto{\pgfqpoint{4.136199in}{0.861085in}}{\pgfqpoint{4.137766in}{0.860797in}}%
\pgfpathlineto{\pgfqpoint{4.137820in}{0.850170in}}%
\pgfpathlineto{\pgfqpoint{4.137820in}{0.850170in}}%
\pgfpathclose%
\pgfpathmoveto{\pgfqpoint{4.170566in}{0.852584in}}%
\pgfpathquadraticcurveto{\pgfqpoint{4.162245in}{0.852584in}}{\pgfqpoint{4.157400in}{0.846081in}}%
\pgfpathquadraticcurveto{\pgfqpoint{4.152554in}{0.839579in}}{\pgfqpoint{4.152554in}{0.828267in}}%
\pgfpathquadraticcurveto{\pgfqpoint{4.152554in}{0.816956in}}{\pgfqpoint{4.157363in}{0.810453in}}%
\pgfpathquadraticcurveto{\pgfqpoint{4.162191in}{0.803951in}}{\pgfqpoint{4.170566in}{0.803951in}}%
\pgfpathquadraticcurveto{\pgfqpoint{4.178852in}{0.803951in}}{\pgfqpoint{4.183679in}{0.810471in}}%
\pgfpathquadraticcurveto{\pgfqpoint{4.188525in}{0.817010in}}{\pgfqpoint{4.188525in}{0.828267in}}%
\pgfpathquadraticcurveto{\pgfqpoint{4.188525in}{0.839471in}}{\pgfqpoint{4.183679in}{0.846027in}}%
\pgfpathquadraticcurveto{\pgfqpoint{4.178852in}{0.852584in}}{\pgfqpoint{4.170566in}{0.852584in}}%
\pgfpathlineto{\pgfqpoint{4.170566in}{0.852584in}}%
\pgfpathclose%
\pgfpathmoveto{\pgfqpoint{4.170566in}{0.861356in}}%
\pgfpathquadraticcurveto{\pgfqpoint{4.184076in}{0.861356in}}{\pgfqpoint{4.191785in}{0.852566in}}%
\pgfpathquadraticcurveto{\pgfqpoint{4.199512in}{0.843794in}}{\pgfqpoint{4.199512in}{0.828267in}}%
\pgfpathquadraticcurveto{\pgfqpoint{4.199512in}{0.812795in}}{\pgfqpoint{4.191785in}{0.803969in}}%
\pgfpathquadraticcurveto{\pgfqpoint{4.184076in}{0.795161in}}{\pgfqpoint{4.170566in}{0.795161in}}%
\pgfpathquadraticcurveto{\pgfqpoint{4.157003in}{0.795161in}}{\pgfqpoint{4.149312in}{0.803969in}}%
\pgfpathquadraticcurveto{\pgfqpoint{4.141639in}{0.812795in}}{\pgfqpoint{4.141639in}{0.828267in}}%
\pgfpathquadraticcurveto{\pgfqpoint{4.141639in}{0.843794in}}{\pgfqpoint{4.149312in}{0.852566in}}%
\pgfpathquadraticcurveto{\pgfqpoint{4.157003in}{0.861356in}}{\pgfqpoint{4.170566in}{0.861356in}}%
\pgfpathlineto{\pgfqpoint{4.170566in}{0.861356in}}%
\pgfpathclose%
\pgfpathmoveto{\pgfqpoint{4.256863in}{0.857987in}}%
\pgfpathlineto{\pgfqpoint{4.256863in}{0.848189in}}%
\pgfpathquadraticcurveto{\pgfqpoint{4.252486in}{0.850440in}}{\pgfqpoint{4.247748in}{0.851557in}}%
\pgfpathquadraticcurveto{\pgfqpoint{4.243029in}{0.852692in}}{\pgfqpoint{4.237950in}{0.852692in}}%
\pgfpathquadraticcurveto{\pgfqpoint{4.230241in}{0.852692in}}{\pgfqpoint{4.226386in}{0.850332in}}%
\pgfpathquadraticcurveto{\pgfqpoint{4.222531in}{0.847973in}}{\pgfqpoint{4.222531in}{0.843235in}}%
\pgfpathquadraticcurveto{\pgfqpoint{4.222531in}{0.839633in}}{\pgfqpoint{4.225287in}{0.837580in}}%
\pgfpathquadraticcurveto{\pgfqpoint{4.228043in}{0.835526in}}{\pgfqpoint{4.236383in}{0.833671in}}%
\pgfpathlineto{\pgfqpoint{4.239931in}{0.832878in}}%
\pgfpathquadraticcurveto{\pgfqpoint{4.250955in}{0.830519in}}{\pgfqpoint{4.255602in}{0.826214in}}%
\pgfpathquadraticcurveto{\pgfqpoint{4.260249in}{0.821909in}}{\pgfqpoint{4.260249in}{0.814200in}}%
\pgfpathquadraticcurveto{\pgfqpoint{4.260249in}{0.805410in}}{\pgfqpoint{4.253296in}{0.800276in}}%
\pgfpathquadraticcurveto{\pgfqpoint{4.246344in}{0.795161in}}{\pgfqpoint{4.234185in}{0.795161in}}%
\pgfpathquadraticcurveto{\pgfqpoint{4.229124in}{0.795161in}}{\pgfqpoint{4.223630in}{0.796152in}}%
\pgfpathquadraticcurveto{\pgfqpoint{4.218137in}{0.797142in}}{\pgfqpoint{4.212066in}{0.799106in}}%
\pgfpathlineto{\pgfqpoint{4.212066in}{0.809805in}}%
\pgfpathquadraticcurveto{\pgfqpoint{4.217812in}{0.806815in}}{\pgfqpoint{4.223378in}{0.805320in}}%
\pgfpathquadraticcurveto{\pgfqpoint{4.228944in}{0.803843in}}{\pgfqpoint{4.234419in}{0.803843in}}%
\pgfpathquadraticcurveto{\pgfqpoint{4.241732in}{0.803843in}}{\pgfqpoint{4.245659in}{0.806346in}}%
\pgfpathquadraticcurveto{\pgfqpoint{4.249604in}{0.808850in}}{\pgfqpoint{4.249604in}{0.813407in}}%
\pgfpathquadraticcurveto{\pgfqpoint{4.249604in}{0.817622in}}{\pgfqpoint{4.246758in}{0.819874in}}%
\pgfpathquadraticcurveto{\pgfqpoint{4.243930in}{0.822125in}}{\pgfqpoint{4.234293in}{0.824214in}}%
\pgfpathlineto{\pgfqpoint{4.230691in}{0.825061in}}%
\pgfpathquadraticcurveto{\pgfqpoint{4.221072in}{0.827078in}}{\pgfqpoint{4.216786in}{0.831275in}}%
\pgfpathquadraticcurveto{\pgfqpoint{4.212517in}{0.835472in}}{\pgfqpoint{4.212517in}{0.842785in}}%
\pgfpathquadraticcurveto{\pgfqpoint{4.212517in}{0.851683in}}{\pgfqpoint{4.218821in}{0.856510in}}%
\pgfpathquadraticcurveto{\pgfqpoint{4.225125in}{0.861356in}}{\pgfqpoint{4.236725in}{0.861356in}}%
\pgfpathquadraticcurveto{\pgfqpoint{4.242453in}{0.861356in}}{\pgfqpoint{4.247514in}{0.860509in}}%
\pgfpathquadraticcurveto{\pgfqpoint{4.252594in}{0.859680in}}{\pgfqpoint{4.256863in}{0.857987in}}%
\pgfpathlineto{\pgfqpoint{4.256863in}{0.857987in}}%
\pgfpathclose%
\pgfpathmoveto{\pgfqpoint{4.316915in}{0.857987in}}%
\pgfpathlineto{\pgfqpoint{4.316915in}{0.848189in}}%
\pgfpathquadraticcurveto{\pgfqpoint{4.312538in}{0.850440in}}{\pgfqpoint{4.307801in}{0.851557in}}%
\pgfpathquadraticcurveto{\pgfqpoint{4.303082in}{0.852692in}}{\pgfqpoint{4.298002in}{0.852692in}}%
\pgfpathquadraticcurveto{\pgfqpoint{4.290293in}{0.852692in}}{\pgfqpoint{4.286439in}{0.850332in}}%
\pgfpathquadraticcurveto{\pgfqpoint{4.282584in}{0.847973in}}{\pgfqpoint{4.282584in}{0.843235in}}%
\pgfpathquadraticcurveto{\pgfqpoint{4.282584in}{0.839633in}}{\pgfqpoint{4.285340in}{0.837580in}}%
\pgfpathquadraticcurveto{\pgfqpoint{4.288096in}{0.835526in}}{\pgfqpoint{4.296435in}{0.833671in}}%
\pgfpathlineto{\pgfqpoint{4.299984in}{0.832878in}}%
\pgfpathquadraticcurveto{\pgfqpoint{4.311007in}{0.830519in}}{\pgfqpoint{4.315654in}{0.826214in}}%
\pgfpathquadraticcurveto{\pgfqpoint{4.320301in}{0.821909in}}{\pgfqpoint{4.320301in}{0.814200in}}%
\pgfpathquadraticcurveto{\pgfqpoint{4.320301in}{0.805410in}}{\pgfqpoint{4.313349in}{0.800276in}}%
\pgfpathquadraticcurveto{\pgfqpoint{4.306396in}{0.795161in}}{\pgfqpoint{4.294238in}{0.795161in}}%
\pgfpathquadraticcurveto{\pgfqpoint{4.289176in}{0.795161in}}{\pgfqpoint{4.283683in}{0.796152in}}%
\pgfpathquadraticcurveto{\pgfqpoint{4.278189in}{0.797142in}}{\pgfqpoint{4.272119in}{0.799106in}}%
\pgfpathlineto{\pgfqpoint{4.272119in}{0.809805in}}%
\pgfpathquadraticcurveto{\pgfqpoint{4.277865in}{0.806815in}}{\pgfqpoint{4.283431in}{0.805320in}}%
\pgfpathquadraticcurveto{\pgfqpoint{4.288996in}{0.803843in}}{\pgfqpoint{4.294472in}{0.803843in}}%
\pgfpathquadraticcurveto{\pgfqpoint{4.301785in}{0.803843in}}{\pgfqpoint{4.305712in}{0.806346in}}%
\pgfpathquadraticcurveto{\pgfqpoint{4.309656in}{0.808850in}}{\pgfqpoint{4.309656in}{0.813407in}}%
\pgfpathquadraticcurveto{\pgfqpoint{4.309656in}{0.817622in}}{\pgfqpoint{4.306810in}{0.819874in}}%
\pgfpathquadraticcurveto{\pgfqpoint{4.303982in}{0.822125in}}{\pgfqpoint{4.294346in}{0.824214in}}%
\pgfpathlineto{\pgfqpoint{4.290743in}{0.825061in}}%
\pgfpathquadraticcurveto{\pgfqpoint{4.281125in}{0.827078in}}{\pgfqpoint{4.276838in}{0.831275in}}%
\pgfpathquadraticcurveto{\pgfqpoint{4.272569in}{0.835472in}}{\pgfqpoint{4.272569in}{0.842785in}}%
\pgfpathquadraticcurveto{\pgfqpoint{4.272569in}{0.851683in}}{\pgfqpoint{4.278873in}{0.856510in}}%
\pgfpathquadraticcurveto{\pgfqpoint{4.285178in}{0.861356in}}{\pgfqpoint{4.296778in}{0.861356in}}%
\pgfpathquadraticcurveto{\pgfqpoint{4.302505in}{0.861356in}}{\pgfqpoint{4.307567in}{0.860509in}}%
\pgfpathquadraticcurveto{\pgfqpoint{4.312646in}{0.859680in}}{\pgfqpoint{4.316915in}{0.857987in}}%
\pgfpathlineto{\pgfqpoint{4.316915in}{0.857987in}}%
\pgfpathclose%
\pgfpathmoveto{\pgfqpoint{4.390711in}{0.830915in}}%
\pgfpathlineto{\pgfqpoint{4.390711in}{0.825854in}}%
\pgfpathlineto{\pgfqpoint{4.343087in}{0.825854in}}%
\pgfpathquadraticcurveto{\pgfqpoint{4.343771in}{0.815154in}}{\pgfqpoint{4.349535in}{0.809553in}}%
\pgfpathquadraticcurveto{\pgfqpoint{4.355299in}{0.803951in}}{\pgfqpoint{4.365602in}{0.803951in}}%
\pgfpathquadraticcurveto{\pgfqpoint{4.371564in}{0.803951in}}{\pgfqpoint{4.377166in}{0.805410in}}%
\pgfpathquadraticcurveto{\pgfqpoint{4.382768in}{0.806869in}}{\pgfqpoint{4.388297in}{0.809805in}}%
\pgfpathlineto{\pgfqpoint{4.388297in}{0.800006in}}%
\pgfpathquadraticcurveto{\pgfqpoint{4.382714in}{0.797647in}}{\pgfqpoint{4.376860in}{0.796404in}}%
\pgfpathquadraticcurveto{\pgfqpoint{4.371006in}{0.795161in}}{\pgfqpoint{4.364990in}{0.795161in}}%
\pgfpathquadraticcurveto{\pgfqpoint{4.349895in}{0.795161in}}{\pgfqpoint{4.341087in}{0.803933in}}%
\pgfpathquadraticcurveto{\pgfqpoint{4.332280in}{0.812723in}}{\pgfqpoint{4.332280in}{0.827709in}}%
\pgfpathquadraticcurveto{\pgfqpoint{4.332280in}{0.843181in}}{\pgfqpoint{4.340637in}{0.852259in}}%
\pgfpathquadraticcurveto{\pgfqpoint{4.348995in}{0.861356in}}{\pgfqpoint{4.363188in}{0.861356in}}%
\pgfpathquadraticcurveto{\pgfqpoint{4.375905in}{0.861356in}}{\pgfqpoint{4.383308in}{0.853160in}}%
\pgfpathquadraticcurveto{\pgfqpoint{4.390711in}{0.844983in}}{\pgfqpoint{4.390711in}{0.830915in}}%
\pgfpathlineto{\pgfqpoint{4.390711in}{0.830915in}}%
\pgfpathclose%
\pgfpathmoveto{\pgfqpoint{4.380354in}{0.833959in}}%
\pgfpathquadraticcurveto{\pgfqpoint{4.380246in}{0.842443in}}{\pgfqpoint{4.375599in}{0.847504in}}%
\pgfpathquadraticcurveto{\pgfqpoint{4.370952in}{0.852584in}}{\pgfqpoint{4.363296in}{0.852584in}}%
\pgfpathquadraticcurveto{\pgfqpoint{4.354633in}{0.852584in}}{\pgfqpoint{4.349427in}{0.847684in}}%
\pgfpathquadraticcurveto{\pgfqpoint{4.344222in}{0.842785in}}{\pgfqpoint{4.343429in}{0.833887in}}%
\pgfpathlineto{\pgfqpoint{4.380354in}{0.833959in}}%
\pgfpathlineto{\pgfqpoint{4.380354in}{0.833959in}}%
\pgfpathclose%
\pgfpathmoveto{\pgfqpoint{4.447900in}{0.857987in}}%
\pgfpathlineto{\pgfqpoint{4.447900in}{0.848189in}}%
\pgfpathquadraticcurveto{\pgfqpoint{4.443523in}{0.850440in}}{\pgfqpoint{4.438785in}{0.851557in}}%
\pgfpathquadraticcurveto{\pgfqpoint{4.434066in}{0.852692in}}{\pgfqpoint{4.428987in}{0.852692in}}%
\pgfpathquadraticcurveto{\pgfqpoint{4.421278in}{0.852692in}}{\pgfqpoint{4.417423in}{0.850332in}}%
\pgfpathquadraticcurveto{\pgfqpoint{4.413568in}{0.847973in}}{\pgfqpoint{4.413568in}{0.843235in}}%
\pgfpathquadraticcurveto{\pgfqpoint{4.413568in}{0.839633in}}{\pgfqpoint{4.416324in}{0.837580in}}%
\pgfpathquadraticcurveto{\pgfqpoint{4.419080in}{0.835526in}}{\pgfqpoint{4.427420in}{0.833671in}}%
\pgfpathlineto{\pgfqpoint{4.430968in}{0.832878in}}%
\pgfpathquadraticcurveto{\pgfqpoint{4.441992in}{0.830519in}}{\pgfqpoint{4.446639in}{0.826214in}}%
\pgfpathquadraticcurveto{\pgfqpoint{4.451286in}{0.821909in}}{\pgfqpoint{4.451286in}{0.814200in}}%
\pgfpathquadraticcurveto{\pgfqpoint{4.451286in}{0.805410in}}{\pgfqpoint{4.444333in}{0.800276in}}%
\pgfpathquadraticcurveto{\pgfqpoint{4.437380in}{0.795161in}}{\pgfqpoint{4.425222in}{0.795161in}}%
\pgfpathquadraticcurveto{\pgfqpoint{4.420161in}{0.795161in}}{\pgfqpoint{4.414667in}{0.796152in}}%
\pgfpathquadraticcurveto{\pgfqpoint{4.409173in}{0.797142in}}{\pgfqpoint{4.403103in}{0.799106in}}%
\pgfpathlineto{\pgfqpoint{4.403103in}{0.809805in}}%
\pgfpathquadraticcurveto{\pgfqpoint{4.408849in}{0.806815in}}{\pgfqpoint{4.414415in}{0.805320in}}%
\pgfpathquadraticcurveto{\pgfqpoint{4.419981in}{0.803843in}}{\pgfqpoint{4.425456in}{0.803843in}}%
\pgfpathquadraticcurveto{\pgfqpoint{4.432769in}{0.803843in}}{\pgfqpoint{4.436696in}{0.806346in}}%
\pgfpathquadraticcurveto{\pgfqpoint{4.440641in}{0.808850in}}{\pgfqpoint{4.440641in}{0.813407in}}%
\pgfpathquadraticcurveto{\pgfqpoint{4.440641in}{0.817622in}}{\pgfqpoint{4.437795in}{0.819874in}}%
\pgfpathquadraticcurveto{\pgfqpoint{4.434967in}{0.822125in}}{\pgfqpoint{4.425330in}{0.824214in}}%
\pgfpathlineto{\pgfqpoint{4.421728in}{0.825061in}}%
\pgfpathquadraticcurveto{\pgfqpoint{4.412109in}{0.827078in}}{\pgfqpoint{4.407822in}{0.831275in}}%
\pgfpathquadraticcurveto{\pgfqpoint{4.403554in}{0.835472in}}{\pgfqpoint{4.403554in}{0.842785in}}%
\pgfpathquadraticcurveto{\pgfqpoint{4.403554in}{0.851683in}}{\pgfqpoint{4.409858in}{0.856510in}}%
\pgfpathquadraticcurveto{\pgfqpoint{4.416162in}{0.861356in}}{\pgfqpoint{4.427762in}{0.861356in}}%
\pgfpathquadraticcurveto{\pgfqpoint{4.433490in}{0.861356in}}{\pgfqpoint{4.438551in}{0.860509in}}%
\pgfpathquadraticcurveto{\pgfqpoint{4.443631in}{0.859680in}}{\pgfqpoint{4.447900in}{0.857987in}}%
\pgfpathlineto{\pgfqpoint{4.447900in}{0.857987in}}%
\pgfpathclose%
\pgfpathmoveto{\pgfqpoint{4.470415in}{0.811102in}}%
\pgfpathlineto{\pgfqpoint{4.482285in}{0.811102in}}%
\pgfpathlineto{\pgfqpoint{4.482285in}{0.796800in}}%
\pgfpathlineto{\pgfqpoint{4.470415in}{0.796800in}}%
\pgfpathlineto{\pgfqpoint{4.470415in}{0.811102in}}%
\pgfpathlineto{\pgfqpoint{4.470415in}{0.811102in}}%
\pgfpathclose%
\pgfpathmoveto{\pgfqpoint{4.470415in}{0.856402in}}%
\pgfpathlineto{\pgfqpoint{4.482285in}{0.856402in}}%
\pgfpathlineto{\pgfqpoint{4.482285in}{0.842119in}}%
\pgfpathlineto{\pgfqpoint{4.470415in}{0.842119in}}%
\pgfpathlineto{\pgfqpoint{4.470415in}{0.856402in}}%
\pgfpathlineto{\pgfqpoint{4.470415in}{0.856402in}}%
\pgfpathclose%
\pgfusepath{fill}%
\end{pgfscope}%
\begin{pgfscope}%
\pgfsetbuttcap%
\pgfsetroundjoin%
\definecolor{currentfill}{rgb}{0.309804,0.372549,0.419608}%
\pgfsetfillcolor{currentfill}%
\pgfsetlinewidth{0.000000pt}%
\definecolor{currentstroke}{rgb}{0.309804,0.372549,0.419608}%
\pgfsetstrokecolor{currentstroke}%
\pgfsetdash{}{0pt}%
\pgfpathmoveto{\pgfqpoint{4.615929in}{0.830915in}}%
\pgfpathlineto{\pgfqpoint{4.615929in}{0.825854in}}%
\pgfpathlineto{\pgfqpoint{4.568304in}{0.825854in}}%
\pgfpathquadraticcurveto{\pgfqpoint{4.568989in}{0.815154in}}{\pgfqpoint{4.574753in}{0.809553in}}%
\pgfpathquadraticcurveto{\pgfqpoint{4.580517in}{0.803951in}}{\pgfqpoint{4.590820in}{0.803951in}}%
\pgfpathquadraticcurveto{\pgfqpoint{4.596782in}{0.803951in}}{\pgfqpoint{4.602383in}{0.805410in}}%
\pgfpathquadraticcurveto{\pgfqpoint{4.607985in}{0.806869in}}{\pgfqpoint{4.613515in}{0.809805in}}%
\pgfpathlineto{\pgfqpoint{4.613515in}{0.800006in}}%
\pgfpathquadraticcurveto{\pgfqpoint{4.607931in}{0.797647in}}{\pgfqpoint{4.602077in}{0.796404in}}%
\pgfpathquadraticcurveto{\pgfqpoint{4.596223in}{0.795161in}}{\pgfqpoint{4.590207in}{0.795161in}}%
\pgfpathquadraticcurveto{\pgfqpoint{4.575113in}{0.795161in}}{\pgfqpoint{4.566305in}{0.803933in}}%
\pgfpathquadraticcurveto{\pgfqpoint{4.557497in}{0.812723in}}{\pgfqpoint{4.557497in}{0.827709in}}%
\pgfpathquadraticcurveto{\pgfqpoint{4.557497in}{0.843181in}}{\pgfqpoint{4.565855in}{0.852259in}}%
\pgfpathquadraticcurveto{\pgfqpoint{4.574212in}{0.861356in}}{\pgfqpoint{4.588406in}{0.861356in}}%
\pgfpathquadraticcurveto{\pgfqpoint{4.601123in}{0.861356in}}{\pgfqpoint{4.608526in}{0.853160in}}%
\pgfpathquadraticcurveto{\pgfqpoint{4.615929in}{0.844983in}}{\pgfqpoint{4.615929in}{0.830915in}}%
\pgfpathlineto{\pgfqpoint{4.615929in}{0.830915in}}%
\pgfpathclose%
\pgfpathmoveto{\pgfqpoint{4.605572in}{0.833959in}}%
\pgfpathquadraticcurveto{\pgfqpoint{4.605464in}{0.842443in}}{\pgfqpoint{4.600816in}{0.847504in}}%
\pgfpathquadraticcurveto{\pgfqpoint{4.596169in}{0.852584in}}{\pgfqpoint{4.588514in}{0.852584in}}%
\pgfpathquadraticcurveto{\pgfqpoint{4.579850in}{0.852584in}}{\pgfqpoint{4.574645in}{0.847684in}}%
\pgfpathquadraticcurveto{\pgfqpoint{4.569439in}{0.842785in}}{\pgfqpoint{4.568647in}{0.833887in}}%
\pgfpathlineto{\pgfqpoint{4.605572in}{0.833959in}}%
\pgfpathlineto{\pgfqpoint{4.605572in}{0.833959in}}%
\pgfpathclose%
\pgfpathmoveto{\pgfqpoint{4.703684in}{0.884393in}}%
\pgfpathlineto{\pgfqpoint{4.703684in}{0.875765in}}%
\pgfpathlineto{\pgfqpoint{4.693777in}{0.875765in}}%
\pgfpathquadraticcurveto{\pgfqpoint{4.688211in}{0.875765in}}{\pgfqpoint{4.686032in}{0.873514in}}%
\pgfpathquadraticcurveto{\pgfqpoint{4.683870in}{0.871262in}}{\pgfqpoint{4.683870in}{0.865408in}}%
\pgfpathlineto{\pgfqpoint{4.683870in}{0.859843in}}%
\pgfpathlineto{\pgfqpoint{4.700928in}{0.859843in}}%
\pgfpathlineto{\pgfqpoint{4.700928in}{0.851791in}}%
\pgfpathlineto{\pgfqpoint{4.683870in}{0.851791in}}%
\pgfpathlineto{\pgfqpoint{4.683870in}{0.796800in}}%
\pgfpathlineto{\pgfqpoint{4.673459in}{0.796800in}}%
\pgfpathlineto{\pgfqpoint{4.673459in}{0.851791in}}%
\pgfpathlineto{\pgfqpoint{4.645036in}{0.851791in}}%
\pgfpathlineto{\pgfqpoint{4.645036in}{0.796800in}}%
\pgfpathlineto{\pgfqpoint{4.634625in}{0.796800in}}%
\pgfpathlineto{\pgfqpoint{4.634625in}{0.851791in}}%
\pgfpathlineto{\pgfqpoint{4.624719in}{0.851791in}}%
\pgfpathlineto{\pgfqpoint{4.624719in}{0.859843in}}%
\pgfpathlineto{\pgfqpoint{4.634625in}{0.859843in}}%
\pgfpathlineto{\pgfqpoint{4.634625in}{0.864238in}}%
\pgfpathquadraticcurveto{\pgfqpoint{4.634625in}{0.874757in}}{\pgfqpoint{4.639525in}{0.879566in}}%
\pgfpathquadraticcurveto{\pgfqpoint{4.644424in}{0.884393in}}{\pgfqpoint{4.655051in}{0.884393in}}%
\pgfpathlineto{\pgfqpoint{4.664850in}{0.884393in}}%
\pgfpathlineto{\pgfqpoint{4.664850in}{0.875765in}}%
\pgfpathlineto{\pgfqpoint{4.654943in}{0.875765in}}%
\pgfpathquadraticcurveto{\pgfqpoint{4.649377in}{0.875765in}}{\pgfqpoint{4.647180in}{0.873514in}}%
\pgfpathquadraticcurveto{\pgfqpoint{4.645036in}{0.871262in}}{\pgfqpoint{4.645036in}{0.865408in}}%
\pgfpathlineto{\pgfqpoint{4.645036in}{0.859843in}}%
\pgfpathlineto{\pgfqpoint{4.673459in}{0.859843in}}%
\pgfpathlineto{\pgfqpoint{4.673459in}{0.864238in}}%
\pgfpathquadraticcurveto{\pgfqpoint{4.673459in}{0.874757in}}{\pgfqpoint{4.678359in}{0.879566in}}%
\pgfpathquadraticcurveto{\pgfqpoint{4.683258in}{0.884393in}}{\pgfqpoint{4.693903in}{0.884393in}}%
\pgfpathlineto{\pgfqpoint{4.703684in}{0.884393in}}%
\pgfpathlineto{\pgfqpoint{4.703684in}{0.884393in}}%
\pgfpathclose%
\pgfpathmoveto{\pgfqpoint{4.766276in}{0.830915in}}%
\pgfpathlineto{\pgfqpoint{4.766276in}{0.825854in}}%
\pgfpathlineto{\pgfqpoint{4.718652in}{0.825854in}}%
\pgfpathquadraticcurveto{\pgfqpoint{4.719336in}{0.815154in}}{\pgfqpoint{4.725100in}{0.809553in}}%
\pgfpathquadraticcurveto{\pgfqpoint{4.730864in}{0.803951in}}{\pgfqpoint{4.741167in}{0.803951in}}%
\pgfpathquadraticcurveto{\pgfqpoint{4.747129in}{0.803951in}}{\pgfqpoint{4.752731in}{0.805410in}}%
\pgfpathquadraticcurveto{\pgfqpoint{4.758333in}{0.806869in}}{\pgfqpoint{4.763862in}{0.809805in}}%
\pgfpathlineto{\pgfqpoint{4.763862in}{0.800006in}}%
\pgfpathquadraticcurveto{\pgfqpoint{4.758279in}{0.797647in}}{\pgfqpoint{4.752425in}{0.796404in}}%
\pgfpathquadraticcurveto{\pgfqpoint{4.746571in}{0.795161in}}{\pgfqpoint{4.740555in}{0.795161in}}%
\pgfpathquadraticcurveto{\pgfqpoint{4.725461in}{0.795161in}}{\pgfqpoint{4.716653in}{0.803933in}}%
\pgfpathquadraticcurveto{\pgfqpoint{4.707845in}{0.812723in}}{\pgfqpoint{4.707845in}{0.827709in}}%
\pgfpathquadraticcurveto{\pgfqpoint{4.707845in}{0.843181in}}{\pgfqpoint{4.716202in}{0.852259in}}%
\pgfpathquadraticcurveto{\pgfqpoint{4.724560in}{0.861356in}}{\pgfqpoint{4.738753in}{0.861356in}}%
\pgfpathquadraticcurveto{\pgfqpoint{4.751470in}{0.861356in}}{\pgfqpoint{4.758873in}{0.853160in}}%
\pgfpathquadraticcurveto{\pgfqpoint{4.766276in}{0.844983in}}{\pgfqpoint{4.766276in}{0.830915in}}%
\pgfpathlineto{\pgfqpoint{4.766276in}{0.830915in}}%
\pgfpathclose%
\pgfpathmoveto{\pgfqpoint{4.755919in}{0.833959in}}%
\pgfpathquadraticcurveto{\pgfqpoint{4.755811in}{0.842443in}}{\pgfqpoint{4.751164in}{0.847504in}}%
\pgfpathquadraticcurveto{\pgfqpoint{4.746517in}{0.852584in}}{\pgfqpoint{4.738862in}{0.852584in}}%
\pgfpathquadraticcurveto{\pgfqpoint{4.730198in}{0.852584in}}{\pgfqpoint{4.724992in}{0.847684in}}%
\pgfpathquadraticcurveto{\pgfqpoint{4.719787in}{0.842785in}}{\pgfqpoint{4.718994in}{0.833887in}}%
\pgfpathlineto{\pgfqpoint{4.755919in}{0.833959in}}%
\pgfpathlineto{\pgfqpoint{4.755919in}{0.833959in}}%
\pgfpathclose%
\pgfpathmoveto{\pgfqpoint{4.828652in}{0.857429in}}%
\pgfpathlineto{\pgfqpoint{4.828652in}{0.847738in}}%
\pgfpathquadraticcurveto{\pgfqpoint{4.824257in}{0.850170in}}{\pgfqpoint{4.819844in}{0.851377in}}%
\pgfpathquadraticcurveto{\pgfqpoint{4.815431in}{0.852584in}}{\pgfqpoint{4.810928in}{0.852584in}}%
\pgfpathquadraticcurveto{\pgfqpoint{4.800841in}{0.852584in}}{\pgfqpoint{4.795258in}{0.846189in}}%
\pgfpathquadraticcurveto{\pgfqpoint{4.789692in}{0.839813in}}{\pgfqpoint{4.789692in}{0.828267in}}%
\pgfpathquadraticcurveto{\pgfqpoint{4.789692in}{0.816721in}}{\pgfqpoint{4.795258in}{0.810327in}}%
\pgfpathquadraticcurveto{\pgfqpoint{4.800841in}{0.803951in}}{\pgfqpoint{4.810928in}{0.803951in}}%
\pgfpathquadraticcurveto{\pgfqpoint{4.815431in}{0.803951in}}{\pgfqpoint{4.819844in}{0.805158in}}%
\pgfpathquadraticcurveto{\pgfqpoint{4.824257in}{0.806364in}}{\pgfqpoint{4.828652in}{0.808796in}}%
\pgfpathlineto{\pgfqpoint{4.828652in}{0.799214in}}%
\pgfpathquadraticcurveto{\pgfqpoint{4.824311in}{0.797196in}}{\pgfqpoint{4.819664in}{0.796188in}}%
\pgfpathquadraticcurveto{\pgfqpoint{4.815035in}{0.795161in}}{\pgfqpoint{4.809793in}{0.795161in}}%
\pgfpathquadraticcurveto{\pgfqpoint{4.795546in}{0.795161in}}{\pgfqpoint{4.787152in}{0.804113in}}%
\pgfpathquadraticcurveto{\pgfqpoint{4.778776in}{0.813065in}}{\pgfqpoint{4.778776in}{0.828267in}}%
\pgfpathquadraticcurveto{\pgfqpoint{4.778776in}{0.843686in}}{\pgfqpoint{4.787242in}{0.852512in}}%
\pgfpathquadraticcurveto{\pgfqpoint{4.795726in}{0.861356in}}{\pgfqpoint{4.810478in}{0.861356in}}%
\pgfpathquadraticcurveto{\pgfqpoint{4.815251in}{0.861356in}}{\pgfqpoint{4.819808in}{0.860365in}}%
\pgfpathquadraticcurveto{\pgfqpoint{4.824365in}{0.859392in}}{\pgfqpoint{4.828652in}{0.857429in}}%
\pgfpathlineto{\pgfqpoint{4.828652in}{0.857429in}}%
\pgfpathclose%
\pgfpathmoveto{\pgfqpoint{4.856913in}{0.877747in}}%
\pgfpathlineto{\pgfqpoint{4.856913in}{0.859843in}}%
\pgfpathlineto{\pgfqpoint{4.878240in}{0.859843in}}%
\pgfpathlineto{\pgfqpoint{4.878240in}{0.851791in}}%
\pgfpathlineto{\pgfqpoint{4.856913in}{0.851791in}}%
\pgfpathlineto{\pgfqpoint{4.856913in}{0.817568in}}%
\pgfpathquadraticcurveto{\pgfqpoint{4.856913in}{0.809859in}}{\pgfqpoint{4.859021in}{0.807661in}}%
\pgfpathquadraticcurveto{\pgfqpoint{4.861128in}{0.805464in}}{\pgfqpoint{4.867612in}{0.805464in}}%
\pgfpathlineto{\pgfqpoint{4.878240in}{0.805464in}}%
\pgfpathlineto{\pgfqpoint{4.878240in}{0.796800in}}%
\pgfpathlineto{\pgfqpoint{4.867612in}{0.796800in}}%
\pgfpathquadraticcurveto{\pgfqpoint{4.855616in}{0.796800in}}{\pgfqpoint{4.851059in}{0.801267in}}%
\pgfpathquadraticcurveto{\pgfqpoint{4.846502in}{0.805752in}}{\pgfqpoint{4.846502in}{0.817568in}}%
\pgfpathlineto{\pgfqpoint{4.846502in}{0.851791in}}%
\pgfpathlineto{\pgfqpoint{4.838901in}{0.851791in}}%
\pgfpathlineto{\pgfqpoint{4.838901in}{0.859843in}}%
\pgfpathlineto{\pgfqpoint{4.846502in}{0.859843in}}%
\pgfpathlineto{\pgfqpoint{4.846502in}{0.877747in}}%
\pgfpathlineto{\pgfqpoint{4.856913in}{0.877747in}}%
\pgfpathlineto{\pgfqpoint{4.856913in}{0.877747in}}%
\pgfpathclose%
\pgfpathmoveto{\pgfqpoint{4.891857in}{0.859843in}}%
\pgfpathlineto{\pgfqpoint{4.902214in}{0.859843in}}%
\pgfpathlineto{\pgfqpoint{4.902214in}{0.796800in}}%
\pgfpathlineto{\pgfqpoint{4.891857in}{0.796800in}}%
\pgfpathlineto{\pgfqpoint{4.891857in}{0.859843in}}%
\pgfpathlineto{\pgfqpoint{4.891857in}{0.859843in}}%
\pgfpathclose%
\pgfpathmoveto{\pgfqpoint{4.891857in}{0.884393in}}%
\pgfpathlineto{\pgfqpoint{4.902214in}{0.884393in}}%
\pgfpathlineto{\pgfqpoint{4.902214in}{0.871262in}}%
\pgfpathlineto{\pgfqpoint{4.891857in}{0.871262in}}%
\pgfpathlineto{\pgfqpoint{4.891857in}{0.884393in}}%
\pgfpathlineto{\pgfqpoint{4.891857in}{0.884393in}}%
\pgfpathclose%
\pgfpathmoveto{\pgfqpoint{4.916461in}{0.859843in}}%
\pgfpathlineto{\pgfqpoint{4.927431in}{0.859843in}}%
\pgfpathlineto{\pgfqpoint{4.947136in}{0.806941in}}%
\pgfpathlineto{\pgfqpoint{4.966841in}{0.859843in}}%
\pgfpathlineto{\pgfqpoint{4.977811in}{0.859843in}}%
\pgfpathlineto{\pgfqpoint{4.954161in}{0.796800in}}%
\pgfpathlineto{\pgfqpoint{4.940093in}{0.796800in}}%
\pgfpathlineto{\pgfqpoint{4.916461in}{0.859843in}}%
\pgfpathlineto{\pgfqpoint{4.916461in}{0.859843in}}%
\pgfpathclose%
\pgfpathmoveto{\pgfqpoint{5.046041in}{0.830915in}}%
\pgfpathlineto{\pgfqpoint{5.046041in}{0.825854in}}%
\pgfpathlineto{\pgfqpoint{4.998417in}{0.825854in}}%
\pgfpathquadraticcurveto{\pgfqpoint{4.999101in}{0.815154in}}{\pgfqpoint{5.004865in}{0.809553in}}%
\pgfpathquadraticcurveto{\pgfqpoint{5.010629in}{0.803951in}}{\pgfqpoint{5.020932in}{0.803951in}}%
\pgfpathquadraticcurveto{\pgfqpoint{5.026894in}{0.803951in}}{\pgfqpoint{5.032496in}{0.805410in}}%
\pgfpathquadraticcurveto{\pgfqpoint{5.038097in}{0.806869in}}{\pgfqpoint{5.043627in}{0.809805in}}%
\pgfpathlineto{\pgfqpoint{5.043627in}{0.800006in}}%
\pgfpathquadraticcurveto{\pgfqpoint{5.038043in}{0.797647in}}{\pgfqpoint{5.032189in}{0.796404in}}%
\pgfpathquadraticcurveto{\pgfqpoint{5.026336in}{0.795161in}}{\pgfqpoint{5.020319in}{0.795161in}}%
\pgfpathquadraticcurveto{\pgfqpoint{5.005225in}{0.795161in}}{\pgfqpoint{4.996417in}{0.803933in}}%
\pgfpathquadraticcurveto{\pgfqpoint{4.987609in}{0.812723in}}{\pgfqpoint{4.987609in}{0.827709in}}%
\pgfpathquadraticcurveto{\pgfqpoint{4.987609in}{0.843181in}}{\pgfqpoint{4.995967in}{0.852259in}}%
\pgfpathquadraticcurveto{\pgfqpoint{5.004325in}{0.861356in}}{\pgfqpoint{5.018518in}{0.861356in}}%
\pgfpathquadraticcurveto{\pgfqpoint{5.031235in}{0.861356in}}{\pgfqpoint{5.038638in}{0.853160in}}%
\pgfpathquadraticcurveto{\pgfqpoint{5.046041in}{0.844983in}}{\pgfqpoint{5.046041in}{0.830915in}}%
\pgfpathlineto{\pgfqpoint{5.046041in}{0.830915in}}%
\pgfpathclose%
\pgfpathmoveto{\pgfqpoint{5.035684in}{0.833959in}}%
\pgfpathquadraticcurveto{\pgfqpoint{5.035576in}{0.842443in}}{\pgfqpoint{5.030929in}{0.847504in}}%
\pgfpathquadraticcurveto{\pgfqpoint{5.026281in}{0.852584in}}{\pgfqpoint{5.018626in}{0.852584in}}%
\pgfpathquadraticcurveto{\pgfqpoint{5.009962in}{0.852584in}}{\pgfqpoint{5.004757in}{0.847684in}}%
\pgfpathquadraticcurveto{\pgfqpoint{4.999551in}{0.842785in}}{\pgfqpoint{4.998759in}{0.833887in}}%
\pgfpathlineto{\pgfqpoint{5.035684in}{0.833959in}}%
\pgfpathlineto{\pgfqpoint{5.035684in}{0.833959in}}%
\pgfpathclose%
\pgfusepath{fill}%
\end{pgfscope}%
\begin{pgfscope}%
\pgfsetbuttcap%
\pgfsetroundjoin%
\definecolor{currentfill}{rgb}{0.309804,0.372549,0.419608}%
\pgfsetfillcolor{currentfill}%
\pgfsetlinewidth{0.000000pt}%
\definecolor{currentstroke}{rgb}{0.309804,0.372549,0.419608}%
\pgfsetstrokecolor{currentstroke}%
\pgfsetdash{}{0pt}%
\pgfpathmoveto{\pgfqpoint{5.171078in}{0.857429in}}%
\pgfpathlineto{\pgfqpoint{5.171078in}{0.847738in}}%
\pgfpathquadraticcurveto{\pgfqpoint{5.166683in}{0.850170in}}{\pgfqpoint{5.162270in}{0.851377in}}%
\pgfpathquadraticcurveto{\pgfqpoint{5.157857in}{0.852584in}}{\pgfqpoint{5.153354in}{0.852584in}}%
\pgfpathquadraticcurveto{\pgfqpoint{5.143267in}{0.852584in}}{\pgfqpoint{5.137683in}{0.846189in}}%
\pgfpathquadraticcurveto{\pgfqpoint{5.132117in}{0.839813in}}{\pgfqpoint{5.132117in}{0.828267in}}%
\pgfpathquadraticcurveto{\pgfqpoint{5.132117in}{0.816721in}}{\pgfqpoint{5.137683in}{0.810327in}}%
\pgfpathquadraticcurveto{\pgfqpoint{5.143267in}{0.803951in}}{\pgfqpoint{5.153354in}{0.803951in}}%
\pgfpathquadraticcurveto{\pgfqpoint{5.157857in}{0.803951in}}{\pgfqpoint{5.162270in}{0.805158in}}%
\pgfpathquadraticcurveto{\pgfqpoint{5.166683in}{0.806364in}}{\pgfqpoint{5.171078in}{0.808796in}}%
\pgfpathlineto{\pgfqpoint{5.171078in}{0.799214in}}%
\pgfpathquadraticcurveto{\pgfqpoint{5.166737in}{0.797196in}}{\pgfqpoint{5.162090in}{0.796188in}}%
\pgfpathquadraticcurveto{\pgfqpoint{5.157461in}{0.795161in}}{\pgfqpoint{5.152219in}{0.795161in}}%
\pgfpathquadraticcurveto{\pgfqpoint{5.137971in}{0.795161in}}{\pgfqpoint{5.129578in}{0.804113in}}%
\pgfpathquadraticcurveto{\pgfqpoint{5.121202in}{0.813065in}}{\pgfqpoint{5.121202in}{0.828267in}}%
\pgfpathquadraticcurveto{\pgfqpoint{5.121202in}{0.843686in}}{\pgfqpoint{5.129668in}{0.852512in}}%
\pgfpathquadraticcurveto{\pgfqpoint{5.138151in}{0.861356in}}{\pgfqpoint{5.152903in}{0.861356in}}%
\pgfpathquadraticcurveto{\pgfqpoint{5.157677in}{0.861356in}}{\pgfqpoint{5.162234in}{0.860365in}}%
\pgfpathquadraticcurveto{\pgfqpoint{5.166791in}{0.859392in}}{\pgfqpoint{5.171078in}{0.857429in}}%
\pgfpathlineto{\pgfqpoint{5.171078in}{0.857429in}}%
\pgfpathclose%
\pgfpathmoveto{\pgfqpoint{5.189090in}{0.884393in}}%
\pgfpathlineto{\pgfqpoint{5.199447in}{0.884393in}}%
\pgfpathlineto{\pgfqpoint{5.199447in}{0.796800in}}%
\pgfpathlineto{\pgfqpoint{5.189090in}{0.796800in}}%
\pgfpathlineto{\pgfqpoint{5.189090in}{0.884393in}}%
\pgfpathlineto{\pgfqpoint{5.189090in}{0.884393in}}%
\pgfpathclose%
\pgfpathmoveto{\pgfqpoint{5.220053in}{0.821675in}}%
\pgfpathlineto{\pgfqpoint{5.220053in}{0.859843in}}%
\pgfpathlineto{\pgfqpoint{5.230410in}{0.859843in}}%
\pgfpathlineto{\pgfqpoint{5.230410in}{0.822071in}}%
\pgfpathquadraticcurveto{\pgfqpoint{5.230410in}{0.813119in}}{\pgfqpoint{5.233886in}{0.808634in}}%
\pgfpathquadraticcurveto{\pgfqpoint{5.237380in}{0.804167in}}{\pgfqpoint{5.244369in}{0.804167in}}%
\pgfpathquadraticcurveto{\pgfqpoint{5.252745in}{0.804167in}}{\pgfqpoint{5.257608in}{0.809517in}}%
\pgfpathquadraticcurveto{\pgfqpoint{5.262489in}{0.814866in}}{\pgfqpoint{5.262489in}{0.824106in}}%
\pgfpathlineto{\pgfqpoint{5.262489in}{0.859843in}}%
\pgfpathlineto{\pgfqpoint{5.272846in}{0.859843in}}%
\pgfpathlineto{\pgfqpoint{5.272846in}{0.796800in}}%
\pgfpathlineto{\pgfqpoint{5.262489in}{0.796800in}}%
\pgfpathlineto{\pgfqpoint{5.262489in}{0.806491in}}%
\pgfpathquadraticcurveto{\pgfqpoint{5.258725in}{0.800745in}}{\pgfqpoint{5.253735in}{0.797953in}}%
\pgfpathquadraticcurveto{\pgfqpoint{5.248764in}{0.795161in}}{\pgfqpoint{5.242172in}{0.795161in}}%
\pgfpathquadraticcurveto{\pgfqpoint{5.231310in}{0.795161in}}{\pgfqpoint{5.225673in}{0.801915in}}%
\pgfpathquadraticcurveto{\pgfqpoint{5.220053in}{0.808670in}}{\pgfqpoint{5.220053in}{0.821675in}}%
\pgfpathlineto{\pgfqpoint{5.220053in}{0.821675in}}%
\pgfpathclose%
\pgfpathmoveto{\pgfqpoint{5.246116in}{0.861356in}}%
\pgfpathlineto{\pgfqpoint{5.246116in}{0.861356in}}%
\pgfpathlineto{\pgfqpoint{5.246116in}{0.861356in}}%
\pgfpathclose%
\pgfpathmoveto{\pgfqpoint{5.334358in}{0.857987in}}%
\pgfpathlineto{\pgfqpoint{5.334358in}{0.848189in}}%
\pgfpathquadraticcurveto{\pgfqpoint{5.329981in}{0.850440in}}{\pgfqpoint{5.325244in}{0.851557in}}%
\pgfpathquadraticcurveto{\pgfqpoint{5.320525in}{0.852692in}}{\pgfqpoint{5.315445in}{0.852692in}}%
\pgfpathquadraticcurveto{\pgfqpoint{5.307736in}{0.852692in}}{\pgfqpoint{5.303881in}{0.850332in}}%
\pgfpathquadraticcurveto{\pgfqpoint{5.300027in}{0.847973in}}{\pgfqpoint{5.300027in}{0.843235in}}%
\pgfpathquadraticcurveto{\pgfqpoint{5.300027in}{0.839633in}}{\pgfqpoint{5.302783in}{0.837580in}}%
\pgfpathquadraticcurveto{\pgfqpoint{5.305538in}{0.835526in}}{\pgfqpoint{5.313878in}{0.833671in}}%
\pgfpathlineto{\pgfqpoint{5.317426in}{0.832878in}}%
\pgfpathquadraticcurveto{\pgfqpoint{5.328450in}{0.830519in}}{\pgfqpoint{5.333097in}{0.826214in}}%
\pgfpathquadraticcurveto{\pgfqpoint{5.337744in}{0.821909in}}{\pgfqpoint{5.337744in}{0.814200in}}%
\pgfpathquadraticcurveto{\pgfqpoint{5.337744in}{0.805410in}}{\pgfqpoint{5.330791in}{0.800276in}}%
\pgfpathquadraticcurveto{\pgfqpoint{5.323839in}{0.795161in}}{\pgfqpoint{5.311681in}{0.795161in}}%
\pgfpathquadraticcurveto{\pgfqpoint{5.306619in}{0.795161in}}{\pgfqpoint{5.301125in}{0.796152in}}%
\pgfpathquadraticcurveto{\pgfqpoint{5.295632in}{0.797142in}}{\pgfqpoint{5.289562in}{0.799106in}}%
\pgfpathlineto{\pgfqpoint{5.289562in}{0.809805in}}%
\pgfpathquadraticcurveto{\pgfqpoint{5.295308in}{0.806815in}}{\pgfqpoint{5.300873in}{0.805320in}}%
\pgfpathquadraticcurveto{\pgfqpoint{5.306439in}{0.803843in}}{\pgfqpoint{5.311915in}{0.803843in}}%
\pgfpathquadraticcurveto{\pgfqpoint{5.319228in}{0.803843in}}{\pgfqpoint{5.323154in}{0.806346in}}%
\pgfpathquadraticcurveto{\pgfqpoint{5.327099in}{0.808850in}}{\pgfqpoint{5.327099in}{0.813407in}}%
\pgfpathquadraticcurveto{\pgfqpoint{5.327099in}{0.817622in}}{\pgfqpoint{5.324253in}{0.819874in}}%
\pgfpathquadraticcurveto{\pgfqpoint{5.321425in}{0.822125in}}{\pgfqpoint{5.311789in}{0.824214in}}%
\pgfpathlineto{\pgfqpoint{5.308186in}{0.825061in}}%
\pgfpathquadraticcurveto{\pgfqpoint{5.298568in}{0.827078in}}{\pgfqpoint{5.294281in}{0.831275in}}%
\pgfpathquadraticcurveto{\pgfqpoint{5.290012in}{0.835472in}}{\pgfqpoint{5.290012in}{0.842785in}}%
\pgfpathquadraticcurveto{\pgfqpoint{5.290012in}{0.851683in}}{\pgfqpoint{5.296316in}{0.856510in}}%
\pgfpathquadraticcurveto{\pgfqpoint{5.302620in}{0.861356in}}{\pgfqpoint{5.314220in}{0.861356in}}%
\pgfpathquadraticcurveto{\pgfqpoint{5.319948in}{0.861356in}}{\pgfqpoint{5.325010in}{0.860509in}}%
\pgfpathquadraticcurveto{\pgfqpoint{5.330089in}{0.859680in}}{\pgfqpoint{5.334358in}{0.857987in}}%
\pgfpathlineto{\pgfqpoint{5.334358in}{0.857987in}}%
\pgfpathclose%
\pgfpathmoveto{\pgfqpoint{5.364474in}{0.877747in}}%
\pgfpathlineto{\pgfqpoint{5.364474in}{0.859843in}}%
\pgfpathlineto{\pgfqpoint{5.385801in}{0.859843in}}%
\pgfpathlineto{\pgfqpoint{5.385801in}{0.851791in}}%
\pgfpathlineto{\pgfqpoint{5.364474in}{0.851791in}}%
\pgfpathlineto{\pgfqpoint{5.364474in}{0.817568in}}%
\pgfpathquadraticcurveto{\pgfqpoint{5.364474in}{0.809859in}}{\pgfqpoint{5.366582in}{0.807661in}}%
\pgfpathquadraticcurveto{\pgfqpoint{5.368689in}{0.805464in}}{\pgfqpoint{5.375173in}{0.805464in}}%
\pgfpathlineto{\pgfqpoint{5.385801in}{0.805464in}}%
\pgfpathlineto{\pgfqpoint{5.385801in}{0.796800in}}%
\pgfpathlineto{\pgfqpoint{5.375173in}{0.796800in}}%
\pgfpathquadraticcurveto{\pgfqpoint{5.363177in}{0.796800in}}{\pgfqpoint{5.358620in}{0.801267in}}%
\pgfpathquadraticcurveto{\pgfqpoint{5.354063in}{0.805752in}}{\pgfqpoint{5.354063in}{0.817568in}}%
\pgfpathlineto{\pgfqpoint{5.354063in}{0.851791in}}%
\pgfpathlineto{\pgfqpoint{5.346462in}{0.851791in}}%
\pgfpathlineto{\pgfqpoint{5.346462in}{0.859843in}}%
\pgfpathlineto{\pgfqpoint{5.354063in}{0.859843in}}%
\pgfpathlineto{\pgfqpoint{5.354063in}{0.877747in}}%
\pgfpathlineto{\pgfqpoint{5.364474in}{0.877747in}}%
\pgfpathlineto{\pgfqpoint{5.364474in}{0.877747in}}%
\pgfpathclose%
\pgfpathmoveto{\pgfqpoint{5.453346in}{0.830915in}}%
\pgfpathlineto{\pgfqpoint{5.453346in}{0.825854in}}%
\pgfpathlineto{\pgfqpoint{5.405722in}{0.825854in}}%
\pgfpathquadraticcurveto{\pgfqpoint{5.406406in}{0.815154in}}{\pgfqpoint{5.412170in}{0.809553in}}%
\pgfpathquadraticcurveto{\pgfqpoint{5.417934in}{0.803951in}}{\pgfqpoint{5.428237in}{0.803951in}}%
\pgfpathquadraticcurveto{\pgfqpoint{5.434199in}{0.803951in}}{\pgfqpoint{5.439801in}{0.805410in}}%
\pgfpathquadraticcurveto{\pgfqpoint{5.445403in}{0.806869in}}{\pgfqpoint{5.450933in}{0.809805in}}%
\pgfpathlineto{\pgfqpoint{5.450933in}{0.800006in}}%
\pgfpathquadraticcurveto{\pgfqpoint{5.445349in}{0.797647in}}{\pgfqpoint{5.439495in}{0.796404in}}%
\pgfpathquadraticcurveto{\pgfqpoint{5.433641in}{0.795161in}}{\pgfqpoint{5.427625in}{0.795161in}}%
\pgfpathquadraticcurveto{\pgfqpoint{5.412531in}{0.795161in}}{\pgfqpoint{5.403723in}{0.803933in}}%
\pgfpathquadraticcurveto{\pgfqpoint{5.394915in}{0.812723in}}{\pgfqpoint{5.394915in}{0.827709in}}%
\pgfpathquadraticcurveto{\pgfqpoint{5.394915in}{0.843181in}}{\pgfqpoint{5.403272in}{0.852259in}}%
\pgfpathquadraticcurveto{\pgfqpoint{5.411630in}{0.861356in}}{\pgfqpoint{5.425824in}{0.861356in}}%
\pgfpathquadraticcurveto{\pgfqpoint{5.438540in}{0.861356in}}{\pgfqpoint{5.445943in}{0.853160in}}%
\pgfpathquadraticcurveto{\pgfqpoint{5.453346in}{0.844983in}}{\pgfqpoint{5.453346in}{0.830915in}}%
\pgfpathlineto{\pgfqpoint{5.453346in}{0.830915in}}%
\pgfpathclose%
\pgfpathmoveto{\pgfqpoint{5.442989in}{0.833959in}}%
\pgfpathquadraticcurveto{\pgfqpoint{5.442881in}{0.842443in}}{\pgfqpoint{5.438234in}{0.847504in}}%
\pgfpathquadraticcurveto{\pgfqpoint{5.433587in}{0.852584in}}{\pgfqpoint{5.425932in}{0.852584in}}%
\pgfpathquadraticcurveto{\pgfqpoint{5.417268in}{0.852584in}}{\pgfqpoint{5.412062in}{0.847684in}}%
\pgfpathquadraticcurveto{\pgfqpoint{5.406857in}{0.842785in}}{\pgfqpoint{5.406064in}{0.833887in}}%
\pgfpathlineto{\pgfqpoint{5.442989in}{0.833959in}}%
\pgfpathlineto{\pgfqpoint{5.442989in}{0.833959in}}%
\pgfpathclose%
\pgfpathmoveto{\pgfqpoint{5.506878in}{0.850170in}}%
\pgfpathquadraticcurveto{\pgfqpoint{5.505131in}{0.851179in}}{\pgfqpoint{5.503078in}{0.851647in}}%
\pgfpathquadraticcurveto{\pgfqpoint{5.501024in}{0.852133in}}{\pgfqpoint{5.498557in}{0.852133in}}%
\pgfpathquadraticcurveto{\pgfqpoint{5.489767in}{0.852133in}}{\pgfqpoint{5.485066in}{0.846423in}}%
\pgfpathquadraticcurveto{\pgfqpoint{5.480364in}{0.840714in}}{\pgfqpoint{5.480364in}{0.830014in}}%
\pgfpathlineto{\pgfqpoint{5.480364in}{0.796800in}}%
\pgfpathlineto{\pgfqpoint{5.469953in}{0.796800in}}%
\pgfpathlineto{\pgfqpoint{5.469953in}{0.859843in}}%
\pgfpathlineto{\pgfqpoint{5.480364in}{0.859843in}}%
\pgfpathlineto{\pgfqpoint{5.480364in}{0.850044in}}%
\pgfpathquadraticcurveto{\pgfqpoint{5.483643in}{0.855790in}}{\pgfqpoint{5.488866in}{0.858564in}}%
\pgfpathquadraticcurveto{\pgfqpoint{5.494108in}{0.861356in}}{\pgfqpoint{5.501601in}{0.861356in}}%
\pgfpathquadraticcurveto{\pgfqpoint{5.502663in}{0.861356in}}{\pgfqpoint{5.503960in}{0.861211in}}%
\pgfpathquadraticcurveto{\pgfqpoint{5.505257in}{0.861085in}}{\pgfqpoint{5.506824in}{0.860797in}}%
\pgfpathlineto{\pgfqpoint{5.506878in}{0.850170in}}%
\pgfpathlineto{\pgfqpoint{5.506878in}{0.850170in}}%
\pgfpathclose%
\pgfpathmoveto{\pgfqpoint{5.557925in}{0.857987in}}%
\pgfpathlineto{\pgfqpoint{5.557925in}{0.848189in}}%
\pgfpathquadraticcurveto{\pgfqpoint{5.553548in}{0.850440in}}{\pgfqpoint{5.548811in}{0.851557in}}%
\pgfpathquadraticcurveto{\pgfqpoint{5.544091in}{0.852692in}}{\pgfqpoint{5.539012in}{0.852692in}}%
\pgfpathquadraticcurveto{\pgfqpoint{5.531303in}{0.852692in}}{\pgfqpoint{5.527448in}{0.850332in}}%
\pgfpathquadraticcurveto{\pgfqpoint{5.523594in}{0.847973in}}{\pgfqpoint{5.523594in}{0.843235in}}%
\pgfpathquadraticcurveto{\pgfqpoint{5.523594in}{0.839633in}}{\pgfqpoint{5.526349in}{0.837580in}}%
\pgfpathquadraticcurveto{\pgfqpoint{5.529105in}{0.835526in}}{\pgfqpoint{5.537445in}{0.833671in}}%
\pgfpathlineto{\pgfqpoint{5.540993in}{0.832878in}}%
\pgfpathquadraticcurveto{\pgfqpoint{5.552017in}{0.830519in}}{\pgfqpoint{5.556664in}{0.826214in}}%
\pgfpathquadraticcurveto{\pgfqpoint{5.561311in}{0.821909in}}{\pgfqpoint{5.561311in}{0.814200in}}%
\pgfpathquadraticcurveto{\pgfqpoint{5.561311in}{0.805410in}}{\pgfqpoint{5.554358in}{0.800276in}}%
\pgfpathquadraticcurveto{\pgfqpoint{5.547406in}{0.795161in}}{\pgfqpoint{5.535247in}{0.795161in}}%
\pgfpathquadraticcurveto{\pgfqpoint{5.530186in}{0.795161in}}{\pgfqpoint{5.524692in}{0.796152in}}%
\pgfpathquadraticcurveto{\pgfqpoint{5.519199in}{0.797142in}}{\pgfqpoint{5.513128in}{0.799106in}}%
\pgfpathlineto{\pgfqpoint{5.513128in}{0.809805in}}%
\pgfpathquadraticcurveto{\pgfqpoint{5.518874in}{0.806815in}}{\pgfqpoint{5.524440in}{0.805320in}}%
\pgfpathquadraticcurveto{\pgfqpoint{5.530006in}{0.803843in}}{\pgfqpoint{5.535482in}{0.803843in}}%
\pgfpathquadraticcurveto{\pgfqpoint{5.542794in}{0.803843in}}{\pgfqpoint{5.546721in}{0.806346in}}%
\pgfpathquadraticcurveto{\pgfqpoint{5.550666in}{0.808850in}}{\pgfqpoint{5.550666in}{0.813407in}}%
\pgfpathquadraticcurveto{\pgfqpoint{5.550666in}{0.817622in}}{\pgfqpoint{5.547820in}{0.819874in}}%
\pgfpathquadraticcurveto{\pgfqpoint{5.544992in}{0.822125in}}{\pgfqpoint{5.535355in}{0.824214in}}%
\pgfpathlineto{\pgfqpoint{5.531753in}{0.825061in}}%
\pgfpathquadraticcurveto{\pgfqpoint{5.522135in}{0.827078in}}{\pgfqpoint{5.517848in}{0.831275in}}%
\pgfpathquadraticcurveto{\pgfqpoint{5.513579in}{0.835472in}}{\pgfqpoint{5.513579in}{0.842785in}}%
\pgfpathquadraticcurveto{\pgfqpoint{5.513579in}{0.851683in}}{\pgfqpoint{5.519883in}{0.856510in}}%
\pgfpathquadraticcurveto{\pgfqpoint{5.526187in}{0.861356in}}{\pgfqpoint{5.537787in}{0.861356in}}%
\pgfpathquadraticcurveto{\pgfqpoint{5.543515in}{0.861356in}}{\pgfqpoint{5.548576in}{0.860509in}}%
\pgfpathquadraticcurveto{\pgfqpoint{5.553656in}{0.859680in}}{\pgfqpoint{5.557925in}{0.857987in}}%
\pgfpathlineto{\pgfqpoint{5.557925in}{0.857987in}}%
\pgfpathclose%
\pgfusepath{fill}%
\end{pgfscope}%
\begin{pgfscope}%
\pgfsetbuttcap%
\pgfsetroundjoin%
\definecolor{currentfill}{rgb}{0.309804,0.372549,0.419608}%
\pgfsetfillcolor{currentfill}%
\pgfsetlinewidth{0.000000pt}%
\definecolor{currentstroke}{rgb}{0.309804,0.372549,0.419608}%
\pgfsetstrokecolor{currentstroke}%
\pgfsetdash{}{0pt}%
\pgfpathmoveto{\pgfqpoint{5.722918in}{0.853538in}}%
\pgfpathlineto{\pgfqpoint{5.664828in}{0.832878in}}%
\pgfpathlineto{\pgfqpoint{5.722918in}{0.812344in}}%
\pgfpathlineto{\pgfqpoint{5.722918in}{0.802096in}}%
\pgfpathlineto{\pgfqpoint{5.650761in}{0.828267in}}%
\pgfpathlineto{\pgfqpoint{5.650761in}{0.837616in}}%
\pgfpathlineto{\pgfqpoint{5.722918in}{0.863787in}}%
\pgfpathlineto{\pgfqpoint{5.722918in}{0.853538in}}%
\pgfpathlineto{\pgfqpoint{5.722918in}{0.853538in}}%
\pgfpathclose%
\pgfpathmoveto{\pgfqpoint{5.749449in}{0.806364in}}%
\pgfpathlineto{\pgfqpoint{5.768020in}{0.806364in}}%
\pgfpathlineto{\pgfqpoint{5.768020in}{0.870488in}}%
\pgfpathlineto{\pgfqpoint{5.747810in}{0.866435in}}%
\pgfpathlineto{\pgfqpoint{5.747810in}{0.876792in}}%
\pgfpathlineto{\pgfqpoint{5.767912in}{0.880845in}}%
\pgfpathlineto{\pgfqpoint{5.779278in}{0.880845in}}%
\pgfpathlineto{\pgfqpoint{5.779278in}{0.806364in}}%
\pgfpathlineto{\pgfqpoint{5.797848in}{0.806364in}}%
\pgfpathlineto{\pgfqpoint{5.797848in}{0.796800in}}%
\pgfpathlineto{\pgfqpoint{5.749449in}{0.796800in}}%
\pgfpathlineto{\pgfqpoint{5.749449in}{0.806364in}}%
\pgfpathlineto{\pgfqpoint{5.749449in}{0.806364in}}%
\pgfpathclose%
\pgfpathmoveto{\pgfqpoint{5.845130in}{0.873352in}}%
\pgfpathquadraticcurveto{\pgfqpoint{5.836358in}{0.873352in}}{\pgfqpoint{5.831927in}{0.864706in}}%
\pgfpathquadraticcurveto{\pgfqpoint{5.827514in}{0.856078in}}{\pgfqpoint{5.827514in}{0.838732in}}%
\pgfpathquadraticcurveto{\pgfqpoint{5.827514in}{0.821459in}}{\pgfqpoint{5.831927in}{0.812813in}}%
\pgfpathquadraticcurveto{\pgfqpoint{5.836358in}{0.804167in}}{\pgfqpoint{5.845130in}{0.804167in}}%
\pgfpathquadraticcurveto{\pgfqpoint{5.853974in}{0.804167in}}{\pgfqpoint{5.858387in}{0.812813in}}%
\pgfpathquadraticcurveto{\pgfqpoint{5.862818in}{0.821459in}}{\pgfqpoint{5.862818in}{0.838732in}}%
\pgfpathquadraticcurveto{\pgfqpoint{5.862818in}{0.856078in}}{\pgfqpoint{5.858387in}{0.864706in}}%
\pgfpathquadraticcurveto{\pgfqpoint{5.853974in}{0.873352in}}{\pgfqpoint{5.845130in}{0.873352in}}%
\pgfpathlineto{\pgfqpoint{5.845130in}{0.873352in}}%
\pgfpathclose%
\pgfpathmoveto{\pgfqpoint{5.845130in}{0.882358in}}%
\pgfpathquadraticcurveto{\pgfqpoint{5.859270in}{0.882358in}}{\pgfqpoint{5.866727in}{0.871172in}}%
\pgfpathquadraticcurveto{\pgfqpoint{5.874184in}{0.860005in}}{\pgfqpoint{5.874184in}{0.838732in}}%
\pgfpathquadraticcurveto{\pgfqpoint{5.874184in}{0.817514in}}{\pgfqpoint{5.866727in}{0.806328in}}%
\pgfpathquadraticcurveto{\pgfqpoint{5.859270in}{0.795161in}}{\pgfqpoint{5.845130in}{0.795161in}}%
\pgfpathquadraticcurveto{\pgfqpoint{5.831008in}{0.795161in}}{\pgfqpoint{5.823551in}{0.806328in}}%
\pgfpathquadraticcurveto{\pgfqpoint{5.816094in}{0.817514in}}{\pgfqpoint{5.816094in}{0.838732in}}%
\pgfpathquadraticcurveto{\pgfqpoint{5.816094in}{0.860005in}}{\pgfqpoint{5.823551in}{0.871172in}}%
\pgfpathquadraticcurveto{\pgfqpoint{5.831008in}{0.882358in}}{\pgfqpoint{5.845130in}{0.882358in}}%
\pgfpathlineto{\pgfqpoint{5.845130in}{0.882358in}}%
\pgfpathclose%
\pgfusepath{fill}%
\end{pgfscope}%
\begin{pgfscope}%
\pgfsetbuttcap%
\pgfsetroundjoin%
\definecolor{currentfill}{rgb}{0.309804,0.372549,0.419608}%
\pgfsetfillcolor{currentfill}%
\pgfsetlinewidth{0.000000pt}%
\definecolor{currentstroke}{rgb}{0.309804,0.372549,0.419608}%
\pgfsetstrokecolor{currentstroke}%
\pgfsetdash{}{0pt}%
\pgfpathmoveto{\pgfqpoint{5.987539in}{0.852584in}}%
\pgfpathquadraticcurveto{\pgfqpoint{5.979218in}{0.852584in}}{\pgfqpoint{5.974372in}{0.846081in}}%
\pgfpathquadraticcurveto{\pgfqpoint{5.969527in}{0.839579in}}{\pgfqpoint{5.969527in}{0.828267in}}%
\pgfpathquadraticcurveto{\pgfqpoint{5.969527in}{0.816956in}}{\pgfqpoint{5.974336in}{0.810453in}}%
\pgfpathquadraticcurveto{\pgfqpoint{5.979164in}{0.803951in}}{\pgfqpoint{5.987539in}{0.803951in}}%
\pgfpathquadraticcurveto{\pgfqpoint{5.995825in}{0.803951in}}{\pgfqpoint{6.000652in}{0.810471in}}%
\pgfpathquadraticcurveto{\pgfqpoint{6.005497in}{0.817010in}}{\pgfqpoint{6.005497in}{0.828267in}}%
\pgfpathquadraticcurveto{\pgfqpoint{6.005497in}{0.839471in}}{\pgfqpoint{6.000652in}{0.846027in}}%
\pgfpathquadraticcurveto{\pgfqpoint{5.995825in}{0.852584in}}{\pgfqpoint{5.987539in}{0.852584in}}%
\pgfpathlineto{\pgfqpoint{5.987539in}{0.852584in}}%
\pgfpathclose%
\pgfpathmoveto{\pgfqpoint{5.987539in}{0.861356in}}%
\pgfpathquadraticcurveto{\pgfqpoint{6.001048in}{0.861356in}}{\pgfqpoint{6.008758in}{0.852566in}}%
\pgfpathquadraticcurveto{\pgfqpoint{6.016485in}{0.843794in}}{\pgfqpoint{6.016485in}{0.828267in}}%
\pgfpathquadraticcurveto{\pgfqpoint{6.016485in}{0.812795in}}{\pgfqpoint{6.008758in}{0.803969in}}%
\pgfpathquadraticcurveto{\pgfqpoint{6.001048in}{0.795161in}}{\pgfqpoint{5.987539in}{0.795161in}}%
\pgfpathquadraticcurveto{\pgfqpoint{5.973976in}{0.795161in}}{\pgfqpoint{5.966285in}{0.803969in}}%
\pgfpathquadraticcurveto{\pgfqpoint{5.958612in}{0.812795in}}{\pgfqpoint{5.958612in}{0.828267in}}%
\pgfpathquadraticcurveto{\pgfqpoint{5.958612in}{0.843794in}}{\pgfqpoint{5.966285in}{0.852566in}}%
\pgfpathquadraticcurveto{\pgfqpoint{5.973976in}{0.861356in}}{\pgfqpoint{5.987539in}{0.861356in}}%
\pgfpathlineto{\pgfqpoint{5.987539in}{0.861356in}}%
\pgfpathclose%
\pgfpathmoveto{\pgfqpoint{6.070179in}{0.850170in}}%
\pgfpathquadraticcurveto{\pgfqpoint{6.068432in}{0.851179in}}{\pgfqpoint{6.066378in}{0.851647in}}%
\pgfpathquadraticcurveto{\pgfqpoint{6.064325in}{0.852133in}}{\pgfqpoint{6.061857in}{0.852133in}}%
\pgfpathquadraticcurveto{\pgfqpoint{6.053067in}{0.852133in}}{\pgfqpoint{6.048366in}{0.846423in}}%
\pgfpathquadraticcurveto{\pgfqpoint{6.043665in}{0.840714in}}{\pgfqpoint{6.043665in}{0.830014in}}%
\pgfpathlineto{\pgfqpoint{6.043665in}{0.796800in}}%
\pgfpathlineto{\pgfqpoint{6.033254in}{0.796800in}}%
\pgfpathlineto{\pgfqpoint{6.033254in}{0.859843in}}%
\pgfpathlineto{\pgfqpoint{6.043665in}{0.859843in}}%
\pgfpathlineto{\pgfqpoint{6.043665in}{0.850044in}}%
\pgfpathquadraticcurveto{\pgfqpoint{6.046943in}{0.855790in}}{\pgfqpoint{6.052167in}{0.858564in}}%
\pgfpathquadraticcurveto{\pgfqpoint{6.057408in}{0.861356in}}{\pgfqpoint{6.064901in}{0.861356in}}%
\pgfpathquadraticcurveto{\pgfqpoint{6.065964in}{0.861356in}}{\pgfqpoint{6.067261in}{0.861211in}}%
\pgfpathquadraticcurveto{\pgfqpoint{6.068558in}{0.861085in}}{\pgfqpoint{6.070125in}{0.860797in}}%
\pgfpathlineto{\pgfqpoint{6.070179in}{0.850170in}}%
\pgfpathlineto{\pgfqpoint{6.070179in}{0.850170in}}%
\pgfpathclose%
\pgfusepath{fill}%
\end{pgfscope}%
\begin{pgfscope}%
\pgfsetbuttcap%
\pgfsetroundjoin%
\definecolor{currentfill}{rgb}{0.309804,0.372549,0.419608}%
\pgfsetfillcolor{currentfill}%
\pgfsetlinewidth{0.000000pt}%
\definecolor{currentstroke}{rgb}{0.309804,0.372549,0.419608}%
\pgfsetstrokecolor{currentstroke}%
\pgfsetdash{}{0pt}%
\pgfpathmoveto{\pgfqpoint{6.185903in}{0.847738in}}%
\pgfpathquadraticcurveto{\pgfqpoint{6.189793in}{0.854727in}}{\pgfqpoint{6.195197in}{0.858041in}}%
\pgfpathquadraticcurveto{\pgfqpoint{6.200601in}{0.861356in}}{\pgfqpoint{6.207914in}{0.861356in}}%
\pgfpathquadraticcurveto{\pgfqpoint{6.217766in}{0.861356in}}{\pgfqpoint{6.223116in}{0.854457in}}%
\pgfpathquadraticcurveto{\pgfqpoint{6.228465in}{0.847576in}}{\pgfqpoint{6.228465in}{0.834860in}}%
\pgfpathlineto{\pgfqpoint{6.228465in}{0.796800in}}%
\pgfpathlineto{\pgfqpoint{6.218054in}{0.796800in}}%
\pgfpathlineto{\pgfqpoint{6.218054in}{0.834517in}}%
\pgfpathquadraticcurveto{\pgfqpoint{6.218054in}{0.843578in}}{\pgfqpoint{6.214830in}{0.847955in}}%
\pgfpathquadraticcurveto{\pgfqpoint{6.211624in}{0.852349in}}{\pgfqpoint{6.205050in}{0.852349in}}%
\pgfpathquadraticcurveto{\pgfqpoint{6.196998in}{0.852349in}}{\pgfqpoint{6.192315in}{0.847000in}}%
\pgfpathquadraticcurveto{\pgfqpoint{6.187650in}{0.841668in}}{\pgfqpoint{6.187650in}{0.832428in}}%
\pgfpathlineto{\pgfqpoint{6.187650in}{0.796800in}}%
\pgfpathlineto{\pgfqpoint{6.177239in}{0.796800in}}%
\pgfpathlineto{\pgfqpoint{6.177239in}{0.834517in}}%
\pgfpathquadraticcurveto{\pgfqpoint{6.177239in}{0.843632in}}{\pgfqpoint{6.174033in}{0.847991in}}%
\pgfpathquadraticcurveto{\pgfqpoint{6.170827in}{0.852349in}}{\pgfqpoint{6.164126in}{0.852349in}}%
\pgfpathquadraticcurveto{\pgfqpoint{6.156183in}{0.852349in}}{\pgfqpoint{6.151500in}{0.846982in}}%
\pgfpathquadraticcurveto{\pgfqpoint{6.146834in}{0.841614in}}{\pgfqpoint{6.146834in}{0.832428in}}%
\pgfpathlineto{\pgfqpoint{6.146834in}{0.796800in}}%
\pgfpathlineto{\pgfqpoint{6.136423in}{0.796800in}}%
\pgfpathlineto{\pgfqpoint{6.136423in}{0.859843in}}%
\pgfpathlineto{\pgfqpoint{6.146834in}{0.859843in}}%
\pgfpathlineto{\pgfqpoint{6.146834in}{0.850044in}}%
\pgfpathquadraticcurveto{\pgfqpoint{6.150383in}{0.855844in}}{\pgfqpoint{6.155336in}{0.858600in}}%
\pgfpathquadraticcurveto{\pgfqpoint{6.160289in}{0.861356in}}{\pgfqpoint{6.167098in}{0.861356in}}%
\pgfpathquadraticcurveto{\pgfqpoint{6.173979in}{0.861356in}}{\pgfqpoint{6.178788in}{0.857861in}}%
\pgfpathquadraticcurveto{\pgfqpoint{6.183597in}{0.854385in}}{\pgfqpoint{6.185903in}{0.847738in}}%
\pgfpathlineto{\pgfqpoint{6.185903in}{0.847738in}}%
\pgfpathclose%
\pgfpathmoveto{\pgfqpoint{6.277765in}{0.828483in}}%
\pgfpathquadraticcurveto{\pgfqpoint{6.265210in}{0.828483in}}{\pgfqpoint{6.260365in}{0.825619in}}%
\pgfpathquadraticcurveto{\pgfqpoint{6.255520in}{0.822756in}}{\pgfqpoint{6.255520in}{0.815821in}}%
\pgfpathquadraticcurveto{\pgfqpoint{6.255520in}{0.810309in}}{\pgfqpoint{6.259158in}{0.807067in}}%
\pgfpathquadraticcurveto{\pgfqpoint{6.262797in}{0.803843in}}{\pgfqpoint{6.269029in}{0.803843in}}%
\pgfpathquadraticcurveto{\pgfqpoint{6.277657in}{0.803843in}}{\pgfqpoint{6.282862in}{0.809949in}}%
\pgfpathquadraticcurveto{\pgfqpoint{6.288068in}{0.816055in}}{\pgfqpoint{6.288068in}{0.826178in}}%
\pgfpathlineto{\pgfqpoint{6.288068in}{0.828483in}}%
\pgfpathlineto{\pgfqpoint{6.277765in}{0.828483in}}%
\pgfpathlineto{\pgfqpoint{6.277765in}{0.828483in}}%
\pgfpathclose%
\pgfpathmoveto{\pgfqpoint{6.298425in}{0.832770in}}%
\pgfpathlineto{\pgfqpoint{6.298425in}{0.796800in}}%
\pgfpathlineto{\pgfqpoint{6.288068in}{0.796800in}}%
\pgfpathlineto{\pgfqpoint{6.288068in}{0.806364in}}%
\pgfpathquadraticcurveto{\pgfqpoint{6.284519in}{0.800637in}}{\pgfqpoint{6.279224in}{0.797899in}}%
\pgfpathquadraticcurveto{\pgfqpoint{6.273928in}{0.795161in}}{\pgfqpoint{6.266273in}{0.795161in}}%
\pgfpathquadraticcurveto{\pgfqpoint{6.256600in}{0.795161in}}{\pgfqpoint{6.250873in}{0.800601in}}%
\pgfpathquadraticcurveto{\pgfqpoint{6.245163in}{0.806040in}}{\pgfqpoint{6.245163in}{0.815154in}}%
\pgfpathquadraticcurveto{\pgfqpoint{6.245163in}{0.825782in}}{\pgfqpoint{6.252278in}{0.831185in}}%
\pgfpathquadraticcurveto{\pgfqpoint{6.259410in}{0.836589in}}{\pgfqpoint{6.273532in}{0.836589in}}%
\pgfpathlineto{\pgfqpoint{6.288068in}{0.836589in}}%
\pgfpathlineto{\pgfqpoint{6.288068in}{0.837616in}}%
\pgfpathquadraticcurveto{\pgfqpoint{6.288068in}{0.844766in}}{\pgfqpoint{6.283367in}{0.848675in}}%
\pgfpathquadraticcurveto{\pgfqpoint{6.278665in}{0.852584in}}{\pgfqpoint{6.270164in}{0.852584in}}%
\pgfpathquadraticcurveto{\pgfqpoint{6.264760in}{0.852584in}}{\pgfqpoint{6.259627in}{0.851287in}}%
\pgfpathquadraticcurveto{\pgfqpoint{6.254511in}{0.849990in}}{\pgfqpoint{6.249792in}{0.847396in}}%
\pgfpathlineto{\pgfqpoint{6.249792in}{0.856979in}}%
\pgfpathquadraticcurveto{\pgfqpoint{6.255466in}{0.859176in}}{\pgfqpoint{6.260815in}{0.860257in}}%
\pgfpathquadraticcurveto{\pgfqpoint{6.266165in}{0.861356in}}{\pgfqpoint{6.271226in}{0.861356in}}%
\pgfpathquadraticcurveto{\pgfqpoint{6.284916in}{0.861356in}}{\pgfqpoint{6.291670in}{0.854259in}}%
\pgfpathquadraticcurveto{\pgfqpoint{6.298425in}{0.847180in}}{\pgfqpoint{6.298425in}{0.832770in}}%
\pgfpathlineto{\pgfqpoint{6.298425in}{0.832770in}}%
\pgfpathclose%
\pgfpathmoveto{\pgfqpoint{6.372166in}{0.859843in}}%
\pgfpathlineto{\pgfqpoint{6.349363in}{0.829168in}}%
\pgfpathlineto{\pgfqpoint{6.373337in}{0.796800in}}%
\pgfpathlineto{\pgfqpoint{6.361125in}{0.796800in}}%
\pgfpathlineto{\pgfqpoint{6.342771in}{0.821567in}}%
\pgfpathlineto{\pgfqpoint{6.324434in}{0.796800in}}%
\pgfpathlineto{\pgfqpoint{6.312204in}{0.796800in}}%
\pgfpathlineto{\pgfqpoint{6.336701in}{0.829780in}}%
\pgfpathlineto{\pgfqpoint{6.314293in}{0.859843in}}%
\pgfpathlineto{\pgfqpoint{6.326506in}{0.859843in}}%
\pgfpathlineto{\pgfqpoint{6.343221in}{0.837381in}}%
\pgfpathlineto{\pgfqpoint{6.359936in}{0.859843in}}%
\pgfpathlineto{\pgfqpoint{6.372166in}{0.859843in}}%
\pgfpathlineto{\pgfqpoint{6.372166in}{0.859843in}}%
\pgfpathclose%
\pgfpathmoveto{\pgfqpoint{6.387981in}{0.859843in}}%
\pgfpathlineto{\pgfqpoint{6.398338in}{0.859843in}}%
\pgfpathlineto{\pgfqpoint{6.398338in}{0.796800in}}%
\pgfpathlineto{\pgfqpoint{6.387981in}{0.796800in}}%
\pgfpathlineto{\pgfqpoint{6.387981in}{0.859843in}}%
\pgfpathlineto{\pgfqpoint{6.387981in}{0.859843in}}%
\pgfpathclose%
\pgfpathmoveto{\pgfqpoint{6.387981in}{0.884393in}}%
\pgfpathlineto{\pgfqpoint{6.398338in}{0.884393in}}%
\pgfpathlineto{\pgfqpoint{6.398338in}{0.871262in}}%
\pgfpathlineto{\pgfqpoint{6.387981in}{0.871262in}}%
\pgfpathlineto{\pgfqpoint{6.387981in}{0.884393in}}%
\pgfpathlineto{\pgfqpoint{6.387981in}{0.884393in}}%
\pgfpathclose%
\pgfpathmoveto{\pgfqpoint{6.469090in}{0.847738in}}%
\pgfpathquadraticcurveto{\pgfqpoint{6.472980in}{0.854727in}}{\pgfqpoint{6.478384in}{0.858041in}}%
\pgfpathquadraticcurveto{\pgfqpoint{6.483788in}{0.861356in}}{\pgfqpoint{6.491101in}{0.861356in}}%
\pgfpathquadraticcurveto{\pgfqpoint{6.500953in}{0.861356in}}{\pgfqpoint{6.506303in}{0.854457in}}%
\pgfpathquadraticcurveto{\pgfqpoint{6.511653in}{0.847576in}}{\pgfqpoint{6.511653in}{0.834860in}}%
\pgfpathlineto{\pgfqpoint{6.511653in}{0.796800in}}%
\pgfpathlineto{\pgfqpoint{6.501242in}{0.796800in}}%
\pgfpathlineto{\pgfqpoint{6.501242in}{0.834517in}}%
\pgfpathquadraticcurveto{\pgfqpoint{6.501242in}{0.843578in}}{\pgfqpoint{6.498017in}{0.847955in}}%
\pgfpathquadraticcurveto{\pgfqpoint{6.494811in}{0.852349in}}{\pgfqpoint{6.488237in}{0.852349in}}%
\pgfpathquadraticcurveto{\pgfqpoint{6.480185in}{0.852349in}}{\pgfqpoint{6.475502in}{0.847000in}}%
\pgfpathquadraticcurveto{\pgfqpoint{6.470837in}{0.841668in}}{\pgfqpoint{6.470837in}{0.832428in}}%
\pgfpathlineto{\pgfqpoint{6.470837in}{0.796800in}}%
\pgfpathlineto{\pgfqpoint{6.460426in}{0.796800in}}%
\pgfpathlineto{\pgfqpoint{6.460426in}{0.834517in}}%
\pgfpathquadraticcurveto{\pgfqpoint{6.460426in}{0.843632in}}{\pgfqpoint{6.457220in}{0.847991in}}%
\pgfpathquadraticcurveto{\pgfqpoint{6.454014in}{0.852349in}}{\pgfqpoint{6.447313in}{0.852349in}}%
\pgfpathquadraticcurveto{\pgfqpoint{6.439370in}{0.852349in}}{\pgfqpoint{6.434687in}{0.846982in}}%
\pgfpathquadraticcurveto{\pgfqpoint{6.430021in}{0.841614in}}{\pgfqpoint{6.430021in}{0.832428in}}%
\pgfpathlineto{\pgfqpoint{6.430021in}{0.796800in}}%
\pgfpathlineto{\pgfqpoint{6.419610in}{0.796800in}}%
\pgfpathlineto{\pgfqpoint{6.419610in}{0.859843in}}%
\pgfpathlineto{\pgfqpoint{6.430021in}{0.859843in}}%
\pgfpathlineto{\pgfqpoint{6.430021in}{0.850044in}}%
\pgfpathquadraticcurveto{\pgfqpoint{6.433570in}{0.855844in}}{\pgfqpoint{6.438523in}{0.858600in}}%
\pgfpathquadraticcurveto{\pgfqpoint{6.443477in}{0.861356in}}{\pgfqpoint{6.450285in}{0.861356in}}%
\pgfpathquadraticcurveto{\pgfqpoint{6.457166in}{0.861356in}}{\pgfqpoint{6.461975in}{0.857861in}}%
\pgfpathquadraticcurveto{\pgfqpoint{6.466784in}{0.854385in}}{\pgfqpoint{6.469090in}{0.847738in}}%
\pgfpathlineto{\pgfqpoint{6.469090in}{0.847738in}}%
\pgfpathclose%
\pgfpathmoveto{\pgfqpoint{6.531232in}{0.821675in}}%
\pgfpathlineto{\pgfqpoint{6.531232in}{0.859843in}}%
\pgfpathlineto{\pgfqpoint{6.541589in}{0.859843in}}%
\pgfpathlineto{\pgfqpoint{6.541589in}{0.822071in}}%
\pgfpathquadraticcurveto{\pgfqpoint{6.541589in}{0.813119in}}{\pgfqpoint{6.545065in}{0.808634in}}%
\pgfpathquadraticcurveto{\pgfqpoint{6.548559in}{0.804167in}}{\pgfqpoint{6.555548in}{0.804167in}}%
\pgfpathquadraticcurveto{\pgfqpoint{6.563924in}{0.804167in}}{\pgfqpoint{6.568787in}{0.809517in}}%
\pgfpathquadraticcurveto{\pgfqpoint{6.573668in}{0.814866in}}{\pgfqpoint{6.573668in}{0.824106in}}%
\pgfpathlineto{\pgfqpoint{6.573668in}{0.859843in}}%
\pgfpathlineto{\pgfqpoint{6.584025in}{0.859843in}}%
\pgfpathlineto{\pgfqpoint{6.584025in}{0.796800in}}%
\pgfpathlineto{\pgfqpoint{6.573668in}{0.796800in}}%
\pgfpathlineto{\pgfqpoint{6.573668in}{0.806491in}}%
\pgfpathquadraticcurveto{\pgfqpoint{6.569904in}{0.800745in}}{\pgfqpoint{6.564914in}{0.797953in}}%
\pgfpathquadraticcurveto{\pgfqpoint{6.559943in}{0.795161in}}{\pgfqpoint{6.553351in}{0.795161in}}%
\pgfpathquadraticcurveto{\pgfqpoint{6.542489in}{0.795161in}}{\pgfqpoint{6.536852in}{0.801915in}}%
\pgfpathquadraticcurveto{\pgfqpoint{6.531232in}{0.808670in}}{\pgfqpoint{6.531232in}{0.821675in}}%
\pgfpathlineto{\pgfqpoint{6.531232in}{0.821675in}}%
\pgfpathclose%
\pgfpathmoveto{\pgfqpoint{6.557295in}{0.861356in}}%
\pgfpathlineto{\pgfqpoint{6.557295in}{0.861356in}}%
\pgfpathlineto{\pgfqpoint{6.557295in}{0.861356in}}%
\pgfpathclose%
\pgfpathmoveto{\pgfqpoint{6.654435in}{0.847738in}}%
\pgfpathquadraticcurveto{\pgfqpoint{6.658326in}{0.854727in}}{\pgfqpoint{6.663729in}{0.858041in}}%
\pgfpathquadraticcurveto{\pgfqpoint{6.669133in}{0.861356in}}{\pgfqpoint{6.676446in}{0.861356in}}%
\pgfpathquadraticcurveto{\pgfqpoint{6.686298in}{0.861356in}}{\pgfqpoint{6.691648in}{0.854457in}}%
\pgfpathquadraticcurveto{\pgfqpoint{6.696998in}{0.847576in}}{\pgfqpoint{6.696998in}{0.834860in}}%
\pgfpathlineto{\pgfqpoint{6.696998in}{0.796800in}}%
\pgfpathlineto{\pgfqpoint{6.686587in}{0.796800in}}%
\pgfpathlineto{\pgfqpoint{6.686587in}{0.834517in}}%
\pgfpathquadraticcurveto{\pgfqpoint{6.686587in}{0.843578in}}{\pgfqpoint{6.683362in}{0.847955in}}%
\pgfpathquadraticcurveto{\pgfqpoint{6.680156in}{0.852349in}}{\pgfqpoint{6.673582in}{0.852349in}}%
\pgfpathquadraticcurveto{\pgfqpoint{6.665530in}{0.852349in}}{\pgfqpoint{6.660847in}{0.847000in}}%
\pgfpathquadraticcurveto{\pgfqpoint{6.656182in}{0.841668in}}{\pgfqpoint{6.656182in}{0.832428in}}%
\pgfpathlineto{\pgfqpoint{6.656182in}{0.796800in}}%
\pgfpathlineto{\pgfqpoint{6.645771in}{0.796800in}}%
\pgfpathlineto{\pgfqpoint{6.645771in}{0.834517in}}%
\pgfpathquadraticcurveto{\pgfqpoint{6.645771in}{0.843632in}}{\pgfqpoint{6.642565in}{0.847991in}}%
\pgfpathquadraticcurveto{\pgfqpoint{6.639359in}{0.852349in}}{\pgfqpoint{6.632658in}{0.852349in}}%
\pgfpathquadraticcurveto{\pgfqpoint{6.624715in}{0.852349in}}{\pgfqpoint{6.620032in}{0.846982in}}%
\pgfpathquadraticcurveto{\pgfqpoint{6.615367in}{0.841614in}}{\pgfqpoint{6.615367in}{0.832428in}}%
\pgfpathlineto{\pgfqpoint{6.615367in}{0.796800in}}%
\pgfpathlineto{\pgfqpoint{6.604956in}{0.796800in}}%
\pgfpathlineto{\pgfqpoint{6.604956in}{0.859843in}}%
\pgfpathlineto{\pgfqpoint{6.615367in}{0.859843in}}%
\pgfpathlineto{\pgfqpoint{6.615367in}{0.850044in}}%
\pgfpathquadraticcurveto{\pgfqpoint{6.618915in}{0.855844in}}{\pgfqpoint{6.623868in}{0.858600in}}%
\pgfpathquadraticcurveto{\pgfqpoint{6.628822in}{0.861356in}}{\pgfqpoint{6.635630in}{0.861356in}}%
\pgfpathquadraticcurveto{\pgfqpoint{6.642511in}{0.861356in}}{\pgfqpoint{6.647320in}{0.857861in}}%
\pgfpathquadraticcurveto{\pgfqpoint{6.652129in}{0.854385in}}{\pgfqpoint{6.654435in}{0.847738in}}%
\pgfpathlineto{\pgfqpoint{6.654435in}{0.847738in}}%
\pgfpathclose%
\pgfusepath{fill}%
\end{pgfscope}%
\begin{pgfscope}%
\pgfsetbuttcap%
\pgfsetroundjoin%
\definecolor{currentfill}{rgb}{0.309804,0.372549,0.419608}%
\pgfsetfillcolor{currentfill}%
\pgfsetlinewidth{0.000000pt}%
\definecolor{currentstroke}{rgb}{0.309804,0.372549,0.419608}%
\pgfsetstrokecolor{currentstroke}%
\pgfsetdash{}{0pt}%
\pgfpathmoveto{\pgfqpoint{6.800525in}{0.884393in}}%
\pgfpathlineto{\pgfqpoint{6.810882in}{0.884393in}}%
\pgfpathlineto{\pgfqpoint{6.810882in}{0.796800in}}%
\pgfpathlineto{\pgfqpoint{6.800525in}{0.796800in}}%
\pgfpathlineto{\pgfqpoint{6.800525in}{0.884393in}}%
\pgfpathlineto{\pgfqpoint{6.800525in}{0.884393in}}%
\pgfpathclose%
\pgfpathmoveto{\pgfqpoint{6.886479in}{0.830915in}}%
\pgfpathlineto{\pgfqpoint{6.886479in}{0.825854in}}%
\pgfpathlineto{\pgfqpoint{6.838854in}{0.825854in}}%
\pgfpathquadraticcurveto{\pgfqpoint{6.839539in}{0.815154in}}{\pgfqpoint{6.845303in}{0.809553in}}%
\pgfpathquadraticcurveto{\pgfqpoint{6.851067in}{0.803951in}}{\pgfqpoint{6.861370in}{0.803951in}}%
\pgfpathquadraticcurveto{\pgfqpoint{6.867332in}{0.803951in}}{\pgfqpoint{6.872933in}{0.805410in}}%
\pgfpathquadraticcurveto{\pgfqpoint{6.878535in}{0.806869in}}{\pgfqpoint{6.884065in}{0.809805in}}%
\pgfpathlineto{\pgfqpoint{6.884065in}{0.800006in}}%
\pgfpathquadraticcurveto{\pgfqpoint{6.878481in}{0.797647in}}{\pgfqpoint{6.872627in}{0.796404in}}%
\pgfpathquadraticcurveto{\pgfqpoint{6.866773in}{0.795161in}}{\pgfqpoint{6.860757in}{0.795161in}}%
\pgfpathquadraticcurveto{\pgfqpoint{6.845663in}{0.795161in}}{\pgfqpoint{6.836855in}{0.803933in}}%
\pgfpathquadraticcurveto{\pgfqpoint{6.828047in}{0.812723in}}{\pgfqpoint{6.828047in}{0.827709in}}%
\pgfpathquadraticcurveto{\pgfqpoint{6.828047in}{0.843181in}}{\pgfqpoint{6.836405in}{0.852259in}}%
\pgfpathquadraticcurveto{\pgfqpoint{6.844762in}{0.861356in}}{\pgfqpoint{6.858956in}{0.861356in}}%
\pgfpathquadraticcurveto{\pgfqpoint{6.871673in}{0.861356in}}{\pgfqpoint{6.879076in}{0.853160in}}%
\pgfpathquadraticcurveto{\pgfqpoint{6.886479in}{0.844983in}}{\pgfqpoint{6.886479in}{0.830915in}}%
\pgfpathlineto{\pgfqpoint{6.886479in}{0.830915in}}%
\pgfpathclose%
\pgfpathmoveto{\pgfqpoint{6.876122in}{0.833959in}}%
\pgfpathquadraticcurveto{\pgfqpoint{6.876013in}{0.842443in}}{\pgfqpoint{6.871366in}{0.847504in}}%
\pgfpathquadraticcurveto{\pgfqpoint{6.866719in}{0.852584in}}{\pgfqpoint{6.859064in}{0.852584in}}%
\pgfpathquadraticcurveto{\pgfqpoint{6.850400in}{0.852584in}}{\pgfqpoint{6.845195in}{0.847684in}}%
\pgfpathquadraticcurveto{\pgfqpoint{6.839989in}{0.842785in}}{\pgfqpoint{6.839197in}{0.833887in}}%
\pgfpathlineto{\pgfqpoint{6.876122in}{0.833959in}}%
\pgfpathlineto{\pgfqpoint{6.876122in}{0.833959in}}%
\pgfpathclose%
\pgfpathmoveto{\pgfqpoint{6.896061in}{0.859843in}}%
\pgfpathlineto{\pgfqpoint{6.907030in}{0.859843in}}%
\pgfpathlineto{\pgfqpoint{6.926736in}{0.806941in}}%
\pgfpathlineto{\pgfqpoint{6.946441in}{0.859843in}}%
\pgfpathlineto{\pgfqpoint{6.957410in}{0.859843in}}%
\pgfpathlineto{\pgfqpoint{6.933760in}{0.796800in}}%
\pgfpathlineto{\pgfqpoint{6.919693in}{0.796800in}}%
\pgfpathlineto{\pgfqpoint{6.896061in}{0.859843in}}%
\pgfpathlineto{\pgfqpoint{6.896061in}{0.859843in}}%
\pgfpathclose%
\pgfpathmoveto{\pgfqpoint{7.025640in}{0.830915in}}%
\pgfpathlineto{\pgfqpoint{7.025640in}{0.825854in}}%
\pgfpathlineto{\pgfqpoint{6.978016in}{0.825854in}}%
\pgfpathquadraticcurveto{\pgfqpoint{6.978701in}{0.815154in}}{\pgfqpoint{6.984465in}{0.809553in}}%
\pgfpathquadraticcurveto{\pgfqpoint{6.990229in}{0.803951in}}{\pgfqpoint{7.000531in}{0.803951in}}%
\pgfpathquadraticcurveto{\pgfqpoint{7.006493in}{0.803951in}}{\pgfqpoint{7.012095in}{0.805410in}}%
\pgfpathquadraticcurveto{\pgfqpoint{7.017697in}{0.806869in}}{\pgfqpoint{7.023227in}{0.809805in}}%
\pgfpathlineto{\pgfqpoint{7.023227in}{0.800006in}}%
\pgfpathquadraticcurveto{\pgfqpoint{7.017643in}{0.797647in}}{\pgfqpoint{7.011789in}{0.796404in}}%
\pgfpathquadraticcurveto{\pgfqpoint{7.005935in}{0.795161in}}{\pgfqpoint{6.999919in}{0.795161in}}%
\pgfpathquadraticcurveto{\pgfqpoint{6.984825in}{0.795161in}}{\pgfqpoint{6.976017in}{0.803933in}}%
\pgfpathquadraticcurveto{\pgfqpoint{6.967209in}{0.812723in}}{\pgfqpoint{6.967209in}{0.827709in}}%
\pgfpathquadraticcurveto{\pgfqpoint{6.967209in}{0.843181in}}{\pgfqpoint{6.975567in}{0.852259in}}%
\pgfpathquadraticcurveto{\pgfqpoint{6.983924in}{0.861356in}}{\pgfqpoint{6.998118in}{0.861356in}}%
\pgfpathquadraticcurveto{\pgfqpoint{7.010834in}{0.861356in}}{\pgfqpoint{7.018237in}{0.853160in}}%
\pgfpathquadraticcurveto{\pgfqpoint{7.025640in}{0.844983in}}{\pgfqpoint{7.025640in}{0.830915in}}%
\pgfpathlineto{\pgfqpoint{7.025640in}{0.830915in}}%
\pgfpathclose%
\pgfpathmoveto{\pgfqpoint{7.015283in}{0.833959in}}%
\pgfpathquadraticcurveto{\pgfqpoint{7.015175in}{0.842443in}}{\pgfqpoint{7.010528in}{0.847504in}}%
\pgfpathquadraticcurveto{\pgfqpoint{7.005881in}{0.852584in}}{\pgfqpoint{6.998226in}{0.852584in}}%
\pgfpathquadraticcurveto{\pgfqpoint{6.989562in}{0.852584in}}{\pgfqpoint{6.984357in}{0.847684in}}%
\pgfpathquadraticcurveto{\pgfqpoint{6.979151in}{0.842785in}}{\pgfqpoint{6.978359in}{0.833887in}}%
\pgfpathlineto{\pgfqpoint{7.015283in}{0.833959in}}%
\pgfpathlineto{\pgfqpoint{7.015283in}{0.833959in}}%
\pgfpathclose%
\pgfpathmoveto{\pgfqpoint{7.079173in}{0.850170in}}%
\pgfpathquadraticcurveto{\pgfqpoint{7.077425in}{0.851179in}}{\pgfqpoint{7.075372in}{0.851647in}}%
\pgfpathquadraticcurveto{\pgfqpoint{7.073319in}{0.852133in}}{\pgfqpoint{7.070851in}{0.852133in}}%
\pgfpathquadraticcurveto{\pgfqpoint{7.062061in}{0.852133in}}{\pgfqpoint{7.057360in}{0.846423in}}%
\pgfpathquadraticcurveto{\pgfqpoint{7.052659in}{0.840714in}}{\pgfqpoint{7.052659in}{0.830014in}}%
\pgfpathlineto{\pgfqpoint{7.052659in}{0.796800in}}%
\pgfpathlineto{\pgfqpoint{7.042248in}{0.796800in}}%
\pgfpathlineto{\pgfqpoint{7.042248in}{0.859843in}}%
\pgfpathlineto{\pgfqpoint{7.052659in}{0.859843in}}%
\pgfpathlineto{\pgfqpoint{7.052659in}{0.850044in}}%
\pgfpathquadraticcurveto{\pgfqpoint{7.055937in}{0.855790in}}{\pgfqpoint{7.061160in}{0.858564in}}%
\pgfpathquadraticcurveto{\pgfqpoint{7.066402in}{0.861356in}}{\pgfqpoint{7.073895in}{0.861356in}}%
\pgfpathquadraticcurveto{\pgfqpoint{7.074958in}{0.861356in}}{\pgfqpoint{7.076255in}{0.861211in}}%
\pgfpathquadraticcurveto{\pgfqpoint{7.077551in}{0.861085in}}{\pgfqpoint{7.079118in}{0.860797in}}%
\pgfpathlineto{\pgfqpoint{7.079173in}{0.850170in}}%
\pgfpathlineto{\pgfqpoint{7.079173in}{0.850170in}}%
\pgfpathclose%
\pgfpathmoveto{\pgfqpoint{7.118691in}{0.828483in}}%
\pgfpathquadraticcurveto{\pgfqpoint{7.106137in}{0.828483in}}{\pgfqpoint{7.101291in}{0.825619in}}%
\pgfpathquadraticcurveto{\pgfqpoint{7.096446in}{0.822756in}}{\pgfqpoint{7.096446in}{0.815821in}}%
\pgfpathquadraticcurveto{\pgfqpoint{7.096446in}{0.810309in}}{\pgfqpoint{7.100085in}{0.807067in}}%
\pgfpathquadraticcurveto{\pgfqpoint{7.103723in}{0.803843in}}{\pgfqpoint{7.109955in}{0.803843in}}%
\pgfpathquadraticcurveto{\pgfqpoint{7.118583in}{0.803843in}}{\pgfqpoint{7.123789in}{0.809949in}}%
\pgfpathquadraticcurveto{\pgfqpoint{7.128994in}{0.816055in}}{\pgfqpoint{7.128994in}{0.826178in}}%
\pgfpathlineto{\pgfqpoint{7.128994in}{0.828483in}}%
\pgfpathlineto{\pgfqpoint{7.118691in}{0.828483in}}%
\pgfpathlineto{\pgfqpoint{7.118691in}{0.828483in}}%
\pgfpathclose%
\pgfpathmoveto{\pgfqpoint{7.139351in}{0.832770in}}%
\pgfpathlineto{\pgfqpoint{7.139351in}{0.796800in}}%
\pgfpathlineto{\pgfqpoint{7.128994in}{0.796800in}}%
\pgfpathlineto{\pgfqpoint{7.128994in}{0.806364in}}%
\pgfpathquadraticcurveto{\pgfqpoint{7.125446in}{0.800637in}}{\pgfqpoint{7.120150in}{0.797899in}}%
\pgfpathquadraticcurveto{\pgfqpoint{7.114855in}{0.795161in}}{\pgfqpoint{7.107199in}{0.795161in}}%
\pgfpathquadraticcurveto{\pgfqpoint{7.097527in}{0.795161in}}{\pgfqpoint{7.091799in}{0.800601in}}%
\pgfpathquadraticcurveto{\pgfqpoint{7.086089in}{0.806040in}}{\pgfqpoint{7.086089in}{0.815154in}}%
\pgfpathquadraticcurveto{\pgfqpoint{7.086089in}{0.825782in}}{\pgfqpoint{7.093204in}{0.831185in}}%
\pgfpathquadraticcurveto{\pgfqpoint{7.100337in}{0.836589in}}{\pgfqpoint{7.114458in}{0.836589in}}%
\pgfpathlineto{\pgfqpoint{7.128994in}{0.836589in}}%
\pgfpathlineto{\pgfqpoint{7.128994in}{0.837616in}}%
\pgfpathquadraticcurveto{\pgfqpoint{7.128994in}{0.844766in}}{\pgfqpoint{7.124293in}{0.848675in}}%
\pgfpathquadraticcurveto{\pgfqpoint{7.119592in}{0.852584in}}{\pgfqpoint{7.111090in}{0.852584in}}%
\pgfpathquadraticcurveto{\pgfqpoint{7.105686in}{0.852584in}}{\pgfqpoint{7.100553in}{0.851287in}}%
\pgfpathquadraticcurveto{\pgfqpoint{7.095437in}{0.849990in}}{\pgfqpoint{7.090718in}{0.847396in}}%
\pgfpathlineto{\pgfqpoint{7.090718in}{0.856979in}}%
\pgfpathquadraticcurveto{\pgfqpoint{7.096392in}{0.859176in}}{\pgfqpoint{7.101742in}{0.860257in}}%
\pgfpathquadraticcurveto{\pgfqpoint{7.107091in}{0.861356in}}{\pgfqpoint{7.112153in}{0.861356in}}%
\pgfpathquadraticcurveto{\pgfqpoint{7.125842in}{0.861356in}}{\pgfqpoint{7.132597in}{0.854259in}}%
\pgfpathquadraticcurveto{\pgfqpoint{7.139351in}{0.847180in}}{\pgfqpoint{7.139351in}{0.832770in}}%
\pgfpathlineto{\pgfqpoint{7.139351in}{0.832770in}}%
\pgfpathclose%
\pgfpathmoveto{\pgfqpoint{7.202160in}{0.829060in}}%
\pgfpathquadraticcurveto{\pgfqpoint{7.202160in}{0.840317in}}{\pgfqpoint{7.197512in}{0.846496in}}%
\pgfpathquadraticcurveto{\pgfqpoint{7.192883in}{0.852692in}}{\pgfqpoint{7.184490in}{0.852692in}}%
\pgfpathquadraticcurveto{\pgfqpoint{7.176168in}{0.852692in}}{\pgfqpoint{7.171521in}{0.846496in}}%
\pgfpathquadraticcurveto{\pgfqpoint{7.166874in}{0.840317in}}{\pgfqpoint{7.166874in}{0.829060in}}%
\pgfpathquadraticcurveto{\pgfqpoint{7.166874in}{0.817856in}}{\pgfqpoint{7.171521in}{0.811660in}}%
\pgfpathquadraticcurveto{\pgfqpoint{7.176168in}{0.805464in}}{\pgfqpoint{7.184490in}{0.805464in}}%
\pgfpathquadraticcurveto{\pgfqpoint{7.192883in}{0.805464in}}{\pgfqpoint{7.197512in}{0.811660in}}%
\pgfpathquadraticcurveto{\pgfqpoint{7.202160in}{0.817856in}}{\pgfqpoint{7.202160in}{0.829060in}}%
\pgfpathlineto{\pgfqpoint{7.202160in}{0.829060in}}%
\pgfpathclose%
\pgfpathmoveto{\pgfqpoint{7.212516in}{0.804617in}}%
\pgfpathquadraticcurveto{\pgfqpoint{7.212516in}{0.788532in}}{\pgfqpoint{7.205366in}{0.780679in}}%
\pgfpathquadraticcurveto{\pgfqpoint{7.198233in}{0.772826in}}{\pgfqpoint{7.183481in}{0.772826in}}%
\pgfpathquadraticcurveto{\pgfqpoint{7.178023in}{0.772826in}}{\pgfqpoint{7.173178in}{0.773636in}}%
\pgfpathquadraticcurveto{\pgfqpoint{7.168333in}{0.774447in}}{\pgfqpoint{7.163776in}{0.776140in}}%
\pgfpathlineto{\pgfqpoint{7.163776in}{0.786209in}}%
\pgfpathquadraticcurveto{\pgfqpoint{7.168333in}{0.783741in}}{\pgfqpoint{7.172782in}{0.782570in}}%
\pgfpathquadraticcurveto{\pgfqpoint{7.177231in}{0.781382in}}{\pgfqpoint{7.181842in}{0.781382in}}%
\pgfpathquadraticcurveto{\pgfqpoint{7.192037in}{0.781382in}}{\pgfqpoint{7.197098in}{0.786695in}}%
\pgfpathquadraticcurveto{\pgfqpoint{7.202160in}{0.792009in}}{\pgfqpoint{7.202160in}{0.802762in}}%
\pgfpathlineto{\pgfqpoint{7.202160in}{0.807895in}}%
\pgfpathquadraticcurveto{\pgfqpoint{7.198953in}{0.802312in}}{\pgfqpoint{7.193946in}{0.799556in}}%
\pgfpathquadraticcurveto{\pgfqpoint{7.188939in}{0.796800in}}{\pgfqpoint{7.181950in}{0.796800in}}%
\pgfpathquadraticcurveto{\pgfqpoint{7.170368in}{0.796800in}}{\pgfqpoint{7.163271in}{0.805626in}}%
\pgfpathquadraticcurveto{\pgfqpoint{7.156174in}{0.814470in}}{\pgfqpoint{7.156174in}{0.829060in}}%
\pgfpathquadraticcurveto{\pgfqpoint{7.156174in}{0.843686in}}{\pgfqpoint{7.163271in}{0.852512in}}%
\pgfpathquadraticcurveto{\pgfqpoint{7.170368in}{0.861356in}}{\pgfqpoint{7.181950in}{0.861356in}}%
\pgfpathquadraticcurveto{\pgfqpoint{7.188939in}{0.861356in}}{\pgfqpoint{7.193946in}{0.858600in}}%
\pgfpathquadraticcurveto{\pgfqpoint{7.198953in}{0.855844in}}{\pgfqpoint{7.202160in}{0.850278in}}%
\pgfpathlineto{\pgfqpoint{7.202160in}{0.859843in}}%
\pgfpathlineto{\pgfqpoint{7.212516in}{0.859843in}}%
\pgfpathlineto{\pgfqpoint{7.212516in}{0.804617in}}%
\pgfpathlineto{\pgfqpoint{7.212516in}{0.804617in}}%
\pgfpathclose%
\pgfpathmoveto{\pgfqpoint{7.287789in}{0.830915in}}%
\pgfpathlineto{\pgfqpoint{7.287789in}{0.825854in}}%
\pgfpathlineto{\pgfqpoint{7.240165in}{0.825854in}}%
\pgfpathquadraticcurveto{\pgfqpoint{7.240850in}{0.815154in}}{\pgfqpoint{7.246613in}{0.809553in}}%
\pgfpathquadraticcurveto{\pgfqpoint{7.252377in}{0.803951in}}{\pgfqpoint{7.262680in}{0.803951in}}%
\pgfpathquadraticcurveto{\pgfqpoint{7.268642in}{0.803951in}}{\pgfqpoint{7.274244in}{0.805410in}}%
\pgfpathquadraticcurveto{\pgfqpoint{7.279846in}{0.806869in}}{\pgfqpoint{7.285376in}{0.809805in}}%
\pgfpathlineto{\pgfqpoint{7.285376in}{0.800006in}}%
\pgfpathquadraticcurveto{\pgfqpoint{7.279792in}{0.797647in}}{\pgfqpoint{7.273938in}{0.796404in}}%
\pgfpathquadraticcurveto{\pgfqpoint{7.268084in}{0.795161in}}{\pgfqpoint{7.262068in}{0.795161in}}%
\pgfpathquadraticcurveto{\pgfqpoint{7.246974in}{0.795161in}}{\pgfqpoint{7.238166in}{0.803933in}}%
\pgfpathquadraticcurveto{\pgfqpoint{7.229358in}{0.812723in}}{\pgfqpoint{7.229358in}{0.827709in}}%
\pgfpathquadraticcurveto{\pgfqpoint{7.229358in}{0.843181in}}{\pgfqpoint{7.237715in}{0.852259in}}%
\pgfpathquadraticcurveto{\pgfqpoint{7.246073in}{0.861356in}}{\pgfqpoint{7.260267in}{0.861356in}}%
\pgfpathquadraticcurveto{\pgfqpoint{7.272983in}{0.861356in}}{\pgfqpoint{7.280386in}{0.853160in}}%
\pgfpathquadraticcurveto{\pgfqpoint{7.287789in}{0.844983in}}{\pgfqpoint{7.287789in}{0.830915in}}%
\pgfpathlineto{\pgfqpoint{7.287789in}{0.830915in}}%
\pgfpathclose%
\pgfpathmoveto{\pgfqpoint{7.277432in}{0.833959in}}%
\pgfpathquadraticcurveto{\pgfqpoint{7.277324in}{0.842443in}}{\pgfqpoint{7.272677in}{0.847504in}}%
\pgfpathquadraticcurveto{\pgfqpoint{7.268030in}{0.852584in}}{\pgfqpoint{7.260375in}{0.852584in}}%
\pgfpathquadraticcurveto{\pgfqpoint{7.251711in}{0.852584in}}{\pgfqpoint{7.246505in}{0.847684in}}%
\pgfpathquadraticcurveto{\pgfqpoint{7.241300in}{0.842785in}}{\pgfqpoint{7.240507in}{0.833887in}}%
\pgfpathlineto{\pgfqpoint{7.277432in}{0.833959in}}%
\pgfpathlineto{\pgfqpoint{7.277432in}{0.833959in}}%
\pgfpathclose%
\pgfusepath{fill}%
\end{pgfscope}%
\begin{pgfscope}%
\pgfsetbuttcap%
\pgfsetroundjoin%
\definecolor{currentfill}{rgb}{0.309804,0.372549,0.419608}%
\pgfsetfillcolor{currentfill}%
\pgfsetlinewidth{0.000000pt}%
\definecolor{currentstroke}{rgb}{0.309804,0.372549,0.419608}%
\pgfsetstrokecolor{currentstroke}%
\pgfsetdash{}{0pt}%
\pgfpathmoveto{\pgfqpoint{7.365580in}{0.853538in}}%
\pgfpathlineto{\pgfqpoint{7.365580in}{0.863787in}}%
\pgfpathlineto{\pgfqpoint{7.437737in}{0.837616in}}%
\pgfpathlineto{\pgfqpoint{7.437737in}{0.828267in}}%
\pgfpathlineto{\pgfqpoint{7.365580in}{0.802096in}}%
\pgfpathlineto{\pgfqpoint{7.365580in}{0.812344in}}%
\pgfpathlineto{\pgfqpoint{7.423561in}{0.832878in}}%
\pgfpathlineto{\pgfqpoint{7.365580in}{0.853538in}}%
\pgfpathlineto{\pgfqpoint{7.365580in}{0.853538in}}%
\pgfpathclose%
\pgfpathmoveto{\pgfqpoint{7.472086in}{0.806364in}}%
\pgfpathlineto{\pgfqpoint{7.511767in}{0.806364in}}%
\pgfpathlineto{\pgfqpoint{7.511767in}{0.796800in}}%
\pgfpathlineto{\pgfqpoint{7.458415in}{0.796800in}}%
\pgfpathlineto{\pgfqpoint{7.458415in}{0.806364in}}%
\pgfpathquadraticcurveto{\pgfqpoint{7.464881in}{0.813065in}}{\pgfqpoint{7.476049in}{0.824341in}}%
\pgfpathquadraticcurveto{\pgfqpoint{7.487234in}{0.835634in}}{\pgfqpoint{7.490098in}{0.838912in}}%
\pgfpathquadraticcurveto{\pgfqpoint{7.495556in}{0.845037in}}{\pgfqpoint{7.497717in}{0.849287in}}%
\pgfpathquadraticcurveto{\pgfqpoint{7.499897in}{0.853538in}}{\pgfqpoint{7.499897in}{0.857645in}}%
\pgfpathquadraticcurveto{\pgfqpoint{7.499897in}{0.864346in}}{\pgfqpoint{7.495196in}{0.868560in}}%
\pgfpathquadraticcurveto{\pgfqpoint{7.490495in}{0.872793in}}{\pgfqpoint{7.482947in}{0.872793in}}%
\pgfpathquadraticcurveto{\pgfqpoint{7.477598in}{0.872793in}}{\pgfqpoint{7.471654in}{0.870938in}}%
\pgfpathquadraticcurveto{\pgfqpoint{7.465728in}{0.869083in}}{\pgfqpoint{7.458973in}{0.865300in}}%
\pgfpathlineto{\pgfqpoint{7.458973in}{0.876792in}}%
\pgfpathquadraticcurveto{\pgfqpoint{7.465836in}{0.879548in}}{\pgfqpoint{7.471798in}{0.880953in}}%
\pgfpathquadraticcurveto{\pgfqpoint{7.477778in}{0.882358in}}{\pgfqpoint{7.482731in}{0.882358in}}%
\pgfpathquadraticcurveto{\pgfqpoint{7.495790in}{0.882358in}}{\pgfqpoint{7.503553in}{0.875819in}}%
\pgfpathquadraticcurveto{\pgfqpoint{7.511317in}{0.869299in}}{\pgfqpoint{7.511317in}{0.858384in}}%
\pgfpathquadraticcurveto{\pgfqpoint{7.511317in}{0.853196in}}{\pgfqpoint{7.509371in}{0.848549in}}%
\pgfpathquadraticcurveto{\pgfqpoint{7.507444in}{0.843920in}}{\pgfqpoint{7.502311in}{0.837616in}}%
\pgfpathquadraticcurveto{\pgfqpoint{7.500906in}{0.835976in}}{\pgfqpoint{7.493359in}{0.828177in}}%
\pgfpathquadraticcurveto{\pgfqpoint{7.485829in}{0.820378in}}{\pgfqpoint{7.472086in}{0.806364in}}%
\pgfpathlineto{\pgfqpoint{7.472086in}{0.806364in}}%
\pgfpathclose%
\pgfpathmoveto{\pgfqpoint{7.535759in}{0.880845in}}%
\pgfpathlineto{\pgfqpoint{7.580393in}{0.880845in}}%
\pgfpathlineto{\pgfqpoint{7.580393in}{0.871262in}}%
\pgfpathlineto{\pgfqpoint{7.546170in}{0.871262in}}%
\pgfpathlineto{\pgfqpoint{7.546170in}{0.850674in}}%
\pgfpathquadraticcurveto{\pgfqpoint{7.548638in}{0.851521in}}{\pgfqpoint{7.551105in}{0.851935in}}%
\pgfpathquadraticcurveto{\pgfqpoint{7.553591in}{0.852349in}}{\pgfqpoint{7.556077in}{0.852349in}}%
\pgfpathquadraticcurveto{\pgfqpoint{7.570144in}{0.852349in}}{\pgfqpoint{7.578358in}{0.844640in}}%
\pgfpathquadraticcurveto{\pgfqpoint{7.586589in}{0.836931in}}{\pgfqpoint{7.586589in}{0.823764in}}%
\pgfpathquadraticcurveto{\pgfqpoint{7.586589in}{0.810201in}}{\pgfqpoint{7.578142in}{0.802672in}}%
\pgfpathquadraticcurveto{\pgfqpoint{7.569694in}{0.795161in}}{\pgfqpoint{7.554330in}{0.795161in}}%
\pgfpathquadraticcurveto{\pgfqpoint{7.549034in}{0.795161in}}{\pgfqpoint{7.543540in}{0.796062in}}%
\pgfpathquadraticcurveto{\pgfqpoint{7.538065in}{0.796962in}}{\pgfqpoint{7.532211in}{0.798763in}}%
\pgfpathlineto{\pgfqpoint{7.532211in}{0.810201in}}%
\pgfpathquadraticcurveto{\pgfqpoint{7.537272in}{0.807445in}}{\pgfqpoint{7.542676in}{0.806094in}}%
\pgfpathquadraticcurveto{\pgfqpoint{7.548079in}{0.804743in}}{\pgfqpoint{7.554095in}{0.804743in}}%
\pgfpathquadraticcurveto{\pgfqpoint{7.563840in}{0.804743in}}{\pgfqpoint{7.569514in}{0.809859in}}%
\pgfpathquadraticcurveto{\pgfqpoint{7.575206in}{0.814974in}}{\pgfqpoint{7.575206in}{0.823764in}}%
\pgfpathquadraticcurveto{\pgfqpoint{7.575206in}{0.832536in}}{\pgfqpoint{7.569514in}{0.837652in}}%
\pgfpathquadraticcurveto{\pgfqpoint{7.563840in}{0.842785in}}{\pgfqpoint{7.554095in}{0.842785in}}%
\pgfpathquadraticcurveto{\pgfqpoint{7.549538in}{0.842785in}}{\pgfqpoint{7.544999in}{0.841776in}}%
\pgfpathquadraticcurveto{\pgfqpoint{7.540478in}{0.840768in}}{\pgfqpoint{7.535759in}{0.838624in}}%
\pgfpathlineto{\pgfqpoint{7.535759in}{0.880845in}}%
\pgfpathlineto{\pgfqpoint{7.535759in}{0.880845in}}%
\pgfpathclose%
\pgfpathmoveto{\pgfqpoint{7.680469in}{0.833779in}}%
\pgfpathquadraticcurveto{\pgfqpoint{7.675569in}{0.833779in}}{\pgfqpoint{7.672778in}{0.829618in}}%
\pgfpathquadraticcurveto{\pgfqpoint{7.670004in}{0.825457in}}{\pgfqpoint{7.670004in}{0.818018in}}%
\pgfpathquadraticcurveto{\pgfqpoint{7.670004in}{0.810705in}}{\pgfqpoint{7.672778in}{0.806509in}}%
\pgfpathquadraticcurveto{\pgfqpoint{7.675569in}{0.802312in}}{\pgfqpoint{7.680469in}{0.802312in}}%
\pgfpathquadraticcurveto{\pgfqpoint{7.685260in}{0.802312in}}{\pgfqpoint{7.688034in}{0.806509in}}%
\pgfpathquadraticcurveto{\pgfqpoint{7.690826in}{0.810705in}}{\pgfqpoint{7.690826in}{0.818018in}}%
\pgfpathquadraticcurveto{\pgfqpoint{7.690826in}{0.825403in}}{\pgfqpoint{7.688034in}{0.829582in}}%
\pgfpathquadraticcurveto{\pgfqpoint{7.685260in}{0.833779in}}{\pgfqpoint{7.680469in}{0.833779in}}%
\pgfpathlineto{\pgfqpoint{7.680469in}{0.833779in}}%
\pgfpathclose%
\pgfpathmoveto{\pgfqpoint{7.680469in}{0.840930in}}%
\pgfpathquadraticcurveto{\pgfqpoint{7.689367in}{0.840930in}}{\pgfqpoint{7.694590in}{0.834734in}}%
\pgfpathquadraticcurveto{\pgfqpoint{7.699832in}{0.828555in}}{\pgfqpoint{7.699832in}{0.818018in}}%
\pgfpathquadraticcurveto{\pgfqpoint{7.699832in}{0.807499in}}{\pgfqpoint{7.694572in}{0.801321in}}%
\pgfpathquadraticcurveto{\pgfqpoint{7.689313in}{0.795161in}}{\pgfqpoint{7.680469in}{0.795161in}}%
\pgfpathquadraticcurveto{\pgfqpoint{7.671463in}{0.795161in}}{\pgfqpoint{7.666221in}{0.801321in}}%
\pgfpathquadraticcurveto{\pgfqpoint{7.660998in}{0.807499in}}{\pgfqpoint{7.660998in}{0.818018in}}%
\pgfpathquadraticcurveto{\pgfqpoint{7.660998in}{0.828609in}}{\pgfqpoint{7.666257in}{0.834770in}}%
\pgfpathquadraticcurveto{\pgfqpoint{7.671517in}{0.840930in}}{\pgfqpoint{7.680469in}{0.840930in}}%
\pgfpathlineto{\pgfqpoint{7.680469in}{0.840930in}}%
\pgfpathclose%
\pgfpathmoveto{\pgfqpoint{7.622380in}{0.875207in}}%
\pgfpathquadraticcurveto{\pgfqpoint{7.617534in}{0.875207in}}{\pgfqpoint{7.614742in}{0.871010in}}%
\pgfpathquadraticcurveto{\pgfqpoint{7.611969in}{0.866831in}}{\pgfqpoint{7.611969in}{0.859500in}}%
\pgfpathquadraticcurveto{\pgfqpoint{7.611969in}{0.852079in}}{\pgfqpoint{7.614724in}{0.847900in}}%
\pgfpathquadraticcurveto{\pgfqpoint{7.617480in}{0.843740in}}{\pgfqpoint{7.622380in}{0.843740in}}%
\pgfpathquadraticcurveto{\pgfqpoint{7.627279in}{0.843740in}}{\pgfqpoint{7.630053in}{0.847900in}}%
\pgfpathquadraticcurveto{\pgfqpoint{7.632845in}{0.852079in}}{\pgfqpoint{7.632845in}{0.859500in}}%
\pgfpathquadraticcurveto{\pgfqpoint{7.632845in}{0.866759in}}{\pgfqpoint{7.630035in}{0.870974in}}%
\pgfpathquadraticcurveto{\pgfqpoint{7.627225in}{0.875207in}}{\pgfqpoint{7.622380in}{0.875207in}}%
\pgfpathlineto{\pgfqpoint{7.622380in}{0.875207in}}%
\pgfpathclose%
\pgfpathmoveto{\pgfqpoint{7.673210in}{0.882358in}}%
\pgfpathlineto{\pgfqpoint{7.682216in}{0.882358in}}%
\pgfpathlineto{\pgfqpoint{7.629638in}{0.795161in}}%
\pgfpathlineto{\pgfqpoint{7.620632in}{0.795161in}}%
\pgfpathlineto{\pgfqpoint{7.673210in}{0.882358in}}%
\pgfpathlineto{\pgfqpoint{7.673210in}{0.882358in}}%
\pgfpathclose%
\pgfpathmoveto{\pgfqpoint{7.622380in}{0.882358in}}%
\pgfpathquadraticcurveto{\pgfqpoint{7.631278in}{0.882358in}}{\pgfqpoint{7.636555in}{0.876198in}}%
\pgfpathquadraticcurveto{\pgfqpoint{7.641851in}{0.870037in}}{\pgfqpoint{7.641851in}{0.859500in}}%
\pgfpathquadraticcurveto{\pgfqpoint{7.641851in}{0.848873in}}{\pgfqpoint{7.636591in}{0.842731in}}%
\pgfpathquadraticcurveto{\pgfqpoint{7.631332in}{0.836589in}}{\pgfqpoint{7.622380in}{0.836589in}}%
\pgfpathquadraticcurveto{\pgfqpoint{7.613428in}{0.836589in}}{\pgfqpoint{7.608222in}{0.842749in}}%
\pgfpathquadraticcurveto{\pgfqpoint{7.603016in}{0.848927in}}{\pgfqpoint{7.603016in}{0.859500in}}%
\pgfpathquadraticcurveto{\pgfqpoint{7.603016in}{0.869983in}}{\pgfqpoint{7.608240in}{0.876162in}}%
\pgfpathquadraticcurveto{\pgfqpoint{7.613482in}{0.882358in}}{\pgfqpoint{7.622380in}{0.882358in}}%
\pgfpathlineto{\pgfqpoint{7.622380in}{0.882358in}}%
\pgfpathclose%
\pgfusepath{fill}%
\end{pgfscope}%
\begin{pgfscope}%
\pgfsetbuttcap%
\pgfsetroundjoin%
\definecolor{currentfill}{rgb}{0.309804,0.372549,0.419608}%
\pgfsetfillcolor{currentfill}%
\pgfsetlinewidth{0.000000pt}%
\definecolor{currentstroke}{rgb}{0.309804,0.372549,0.419608}%
\pgfsetstrokecolor{currentstroke}%
\pgfsetdash{}{0pt}%
\pgfpathmoveto{\pgfqpoint{7.819852in}{0.884393in}}%
\pgfpathlineto{\pgfqpoint{7.819852in}{0.875765in}}%
\pgfpathlineto{\pgfqpoint{7.809945in}{0.875765in}}%
\pgfpathquadraticcurveto{\pgfqpoint{7.804379in}{0.875765in}}{\pgfqpoint{7.802200in}{0.873514in}}%
\pgfpathquadraticcurveto{\pgfqpoint{7.800038in}{0.871262in}}{\pgfqpoint{7.800038in}{0.865408in}}%
\pgfpathlineto{\pgfqpoint{7.800038in}{0.859843in}}%
\pgfpathlineto{\pgfqpoint{7.817096in}{0.859843in}}%
\pgfpathlineto{\pgfqpoint{7.817096in}{0.851791in}}%
\pgfpathlineto{\pgfqpoint{7.800038in}{0.851791in}}%
\pgfpathlineto{\pgfqpoint{7.800038in}{0.796800in}}%
\pgfpathlineto{\pgfqpoint{7.789627in}{0.796800in}}%
\pgfpathlineto{\pgfqpoint{7.789627in}{0.851791in}}%
\pgfpathlineto{\pgfqpoint{7.779721in}{0.851791in}}%
\pgfpathlineto{\pgfqpoint{7.779721in}{0.859843in}}%
\pgfpathlineto{\pgfqpoint{7.789627in}{0.859843in}}%
\pgfpathlineto{\pgfqpoint{7.789627in}{0.864238in}}%
\pgfpathquadraticcurveto{\pgfqpoint{7.789627in}{0.874757in}}{\pgfqpoint{7.794527in}{0.879566in}}%
\pgfpathquadraticcurveto{\pgfqpoint{7.799426in}{0.884393in}}{\pgfqpoint{7.810053in}{0.884393in}}%
\pgfpathlineto{\pgfqpoint{7.819852in}{0.884393in}}%
\pgfpathlineto{\pgfqpoint{7.819852in}{0.884393in}}%
\pgfpathclose%
\pgfpathmoveto{\pgfqpoint{7.852940in}{0.852584in}}%
\pgfpathquadraticcurveto{\pgfqpoint{7.844618in}{0.852584in}}{\pgfqpoint{7.839773in}{0.846081in}}%
\pgfpathquadraticcurveto{\pgfqpoint{7.834928in}{0.839579in}}{\pgfqpoint{7.834928in}{0.828267in}}%
\pgfpathquadraticcurveto{\pgfqpoint{7.834928in}{0.816956in}}{\pgfqpoint{7.839737in}{0.810453in}}%
\pgfpathquadraticcurveto{\pgfqpoint{7.844564in}{0.803951in}}{\pgfqpoint{7.852940in}{0.803951in}}%
\pgfpathquadraticcurveto{\pgfqpoint{7.861226in}{0.803951in}}{\pgfqpoint{7.866053in}{0.810471in}}%
\pgfpathquadraticcurveto{\pgfqpoint{7.870898in}{0.817010in}}{\pgfqpoint{7.870898in}{0.828267in}}%
\pgfpathquadraticcurveto{\pgfqpoint{7.870898in}{0.839471in}}{\pgfqpoint{7.866053in}{0.846027in}}%
\pgfpathquadraticcurveto{\pgfqpoint{7.861226in}{0.852584in}}{\pgfqpoint{7.852940in}{0.852584in}}%
\pgfpathlineto{\pgfqpoint{7.852940in}{0.852584in}}%
\pgfpathclose%
\pgfpathmoveto{\pgfqpoint{7.852940in}{0.861356in}}%
\pgfpathquadraticcurveto{\pgfqpoint{7.866449in}{0.861356in}}{\pgfqpoint{7.874158in}{0.852566in}}%
\pgfpathquadraticcurveto{\pgfqpoint{7.881886in}{0.843794in}}{\pgfqpoint{7.881886in}{0.828267in}}%
\pgfpathquadraticcurveto{\pgfqpoint{7.881886in}{0.812795in}}{\pgfqpoint{7.874158in}{0.803969in}}%
\pgfpathquadraticcurveto{\pgfqpoint{7.866449in}{0.795161in}}{\pgfqpoint{7.852940in}{0.795161in}}%
\pgfpathquadraticcurveto{\pgfqpoint{7.839377in}{0.795161in}}{\pgfqpoint{7.831686in}{0.803969in}}%
\pgfpathquadraticcurveto{\pgfqpoint{7.824013in}{0.812795in}}{\pgfqpoint{7.824013in}{0.828267in}}%
\pgfpathquadraticcurveto{\pgfqpoint{7.824013in}{0.843794in}}{\pgfqpoint{7.831686in}{0.852566in}}%
\pgfpathquadraticcurveto{\pgfqpoint{7.839377in}{0.861356in}}{\pgfqpoint{7.852940in}{0.861356in}}%
\pgfpathlineto{\pgfqpoint{7.852940in}{0.861356in}}%
\pgfpathclose%
\pgfpathmoveto{\pgfqpoint{7.935580in}{0.850170in}}%
\pgfpathquadraticcurveto{\pgfqpoint{7.933833in}{0.851179in}}{\pgfqpoint{7.931779in}{0.851647in}}%
\pgfpathquadraticcurveto{\pgfqpoint{7.929726in}{0.852133in}}{\pgfqpoint{7.927258in}{0.852133in}}%
\pgfpathquadraticcurveto{\pgfqpoint{7.918468in}{0.852133in}}{\pgfqpoint{7.913767in}{0.846423in}}%
\pgfpathquadraticcurveto{\pgfqpoint{7.909066in}{0.840714in}}{\pgfqpoint{7.909066in}{0.830014in}}%
\pgfpathlineto{\pgfqpoint{7.909066in}{0.796800in}}%
\pgfpathlineto{\pgfqpoint{7.898655in}{0.796800in}}%
\pgfpathlineto{\pgfqpoint{7.898655in}{0.859843in}}%
\pgfpathlineto{\pgfqpoint{7.909066in}{0.859843in}}%
\pgfpathlineto{\pgfqpoint{7.909066in}{0.850044in}}%
\pgfpathquadraticcurveto{\pgfqpoint{7.912344in}{0.855790in}}{\pgfqpoint{7.917568in}{0.858564in}}%
\pgfpathquadraticcurveto{\pgfqpoint{7.922809in}{0.861356in}}{\pgfqpoint{7.930302in}{0.861356in}}%
\pgfpathquadraticcurveto{\pgfqpoint{7.931365in}{0.861356in}}{\pgfqpoint{7.932662in}{0.861211in}}%
\pgfpathquadraticcurveto{\pgfqpoint{7.933959in}{0.861085in}}{\pgfqpoint{7.935526in}{0.860797in}}%
\pgfpathlineto{\pgfqpoint{7.935580in}{0.850170in}}%
\pgfpathlineto{\pgfqpoint{7.935580in}{0.850170in}}%
\pgfpathclose%
\pgfusepath{fill}%
\end{pgfscope}%
\begin{pgfscope}%
\pgfsetbuttcap%
\pgfsetroundjoin%
\definecolor{currentfill}{rgb}{0.309804,0.372549,0.419608}%
\pgfsetfillcolor{currentfill}%
\pgfsetlinewidth{0.000000pt}%
\definecolor{currentstroke}{rgb}{0.309804,0.372549,0.419608}%
\pgfsetstrokecolor{currentstroke}%
\pgfsetdash{}{0pt}%
\pgfpathmoveto{\pgfqpoint{8.055567in}{0.830915in}}%
\pgfpathlineto{\pgfqpoint{8.055567in}{0.825854in}}%
\pgfpathlineto{\pgfqpoint{8.007943in}{0.825854in}}%
\pgfpathquadraticcurveto{\pgfqpoint{8.008628in}{0.815154in}}{\pgfqpoint{8.014392in}{0.809553in}}%
\pgfpathquadraticcurveto{\pgfqpoint{8.020156in}{0.803951in}}{\pgfqpoint{8.030459in}{0.803951in}}%
\pgfpathquadraticcurveto{\pgfqpoint{8.036421in}{0.803951in}}{\pgfqpoint{8.042022in}{0.805410in}}%
\pgfpathquadraticcurveto{\pgfqpoint{8.047624in}{0.806869in}}{\pgfqpoint{8.053154in}{0.809805in}}%
\pgfpathlineto{\pgfqpoint{8.053154in}{0.800006in}}%
\pgfpathquadraticcurveto{\pgfqpoint{8.047570in}{0.797647in}}{\pgfqpoint{8.041716in}{0.796404in}}%
\pgfpathquadraticcurveto{\pgfqpoint{8.035862in}{0.795161in}}{\pgfqpoint{8.029846in}{0.795161in}}%
\pgfpathquadraticcurveto{\pgfqpoint{8.014752in}{0.795161in}}{\pgfqpoint{8.005944in}{0.803933in}}%
\pgfpathquadraticcurveto{\pgfqpoint{7.997136in}{0.812723in}}{\pgfqpoint{7.997136in}{0.827709in}}%
\pgfpathquadraticcurveto{\pgfqpoint{7.997136in}{0.843181in}}{\pgfqpoint{8.005494in}{0.852259in}}%
\pgfpathquadraticcurveto{\pgfqpoint{8.013851in}{0.861356in}}{\pgfqpoint{8.028045in}{0.861356in}}%
\pgfpathquadraticcurveto{\pgfqpoint{8.040762in}{0.861356in}}{\pgfqpoint{8.048164in}{0.853160in}}%
\pgfpathquadraticcurveto{\pgfqpoint{8.055567in}{0.844983in}}{\pgfqpoint{8.055567in}{0.830915in}}%
\pgfpathlineto{\pgfqpoint{8.055567in}{0.830915in}}%
\pgfpathclose%
\pgfpathmoveto{\pgfqpoint{8.045211in}{0.833959in}}%
\pgfpathquadraticcurveto{\pgfqpoint{8.045102in}{0.842443in}}{\pgfqpoint{8.040455in}{0.847504in}}%
\pgfpathquadraticcurveto{\pgfqpoint{8.035808in}{0.852584in}}{\pgfqpoint{8.028153in}{0.852584in}}%
\pgfpathquadraticcurveto{\pgfqpoint{8.019489in}{0.852584in}}{\pgfqpoint{8.014284in}{0.847684in}}%
\pgfpathquadraticcurveto{\pgfqpoint{8.009078in}{0.842785in}}{\pgfqpoint{8.008286in}{0.833887in}}%
\pgfpathlineto{\pgfqpoint{8.045211in}{0.833959in}}%
\pgfpathlineto{\pgfqpoint{8.045211in}{0.833959in}}%
\pgfpathclose%
\pgfpathmoveto{\pgfqpoint{8.072571in}{0.859843in}}%
\pgfpathlineto{\pgfqpoint{8.082928in}{0.859843in}}%
\pgfpathlineto{\pgfqpoint{8.082928in}{0.796800in}}%
\pgfpathlineto{\pgfqpoint{8.072571in}{0.796800in}}%
\pgfpathlineto{\pgfqpoint{8.072571in}{0.859843in}}%
\pgfpathlineto{\pgfqpoint{8.072571in}{0.859843in}}%
\pgfpathclose%
\pgfpathmoveto{\pgfqpoint{8.072571in}{0.884393in}}%
\pgfpathlineto{\pgfqpoint{8.082928in}{0.884393in}}%
\pgfpathlineto{\pgfqpoint{8.082928in}{0.871262in}}%
\pgfpathlineto{\pgfqpoint{8.072571in}{0.871262in}}%
\pgfpathlineto{\pgfqpoint{8.072571in}{0.884393in}}%
\pgfpathlineto{\pgfqpoint{8.072571in}{0.884393in}}%
\pgfpathclose%
\pgfpathmoveto{\pgfqpoint{8.114845in}{0.877747in}}%
\pgfpathlineto{\pgfqpoint{8.114845in}{0.859843in}}%
\pgfpathlineto{\pgfqpoint{8.136172in}{0.859843in}}%
\pgfpathlineto{\pgfqpoint{8.136172in}{0.851791in}}%
\pgfpathlineto{\pgfqpoint{8.114845in}{0.851791in}}%
\pgfpathlineto{\pgfqpoint{8.114845in}{0.817568in}}%
\pgfpathquadraticcurveto{\pgfqpoint{8.114845in}{0.809859in}}{\pgfqpoint{8.116953in}{0.807661in}}%
\pgfpathquadraticcurveto{\pgfqpoint{8.119060in}{0.805464in}}{\pgfqpoint{8.125545in}{0.805464in}}%
\pgfpathlineto{\pgfqpoint{8.136172in}{0.805464in}}%
\pgfpathlineto{\pgfqpoint{8.136172in}{0.796800in}}%
\pgfpathlineto{\pgfqpoint{8.125545in}{0.796800in}}%
\pgfpathquadraticcurveto{\pgfqpoint{8.113549in}{0.796800in}}{\pgfqpoint{8.108992in}{0.801267in}}%
\pgfpathquadraticcurveto{\pgfqpoint{8.104434in}{0.805752in}}{\pgfqpoint{8.104434in}{0.817568in}}%
\pgfpathlineto{\pgfqpoint{8.104434in}{0.851791in}}%
\pgfpathlineto{\pgfqpoint{8.096833in}{0.851791in}}%
\pgfpathlineto{\pgfqpoint{8.096833in}{0.859843in}}%
\pgfpathlineto{\pgfqpoint{8.104434in}{0.859843in}}%
\pgfpathlineto{\pgfqpoint{8.104434in}{0.877747in}}%
\pgfpathlineto{\pgfqpoint{8.114845in}{0.877747in}}%
\pgfpathlineto{\pgfqpoint{8.114845in}{0.877747in}}%
\pgfpathclose%
\pgfpathmoveto{\pgfqpoint{8.202204in}{0.834860in}}%
\pgfpathlineto{\pgfqpoint{8.202204in}{0.796800in}}%
\pgfpathlineto{\pgfqpoint{8.191847in}{0.796800in}}%
\pgfpathlineto{\pgfqpoint{8.191847in}{0.834517in}}%
\pgfpathquadraticcurveto{\pgfqpoint{8.191847in}{0.843469in}}{\pgfqpoint{8.188353in}{0.847900in}}%
\pgfpathquadraticcurveto{\pgfqpoint{8.184859in}{0.852349in}}{\pgfqpoint{8.177888in}{0.852349in}}%
\pgfpathquadraticcurveto{\pgfqpoint{8.169494in}{0.852349in}}{\pgfqpoint{8.164649in}{0.847000in}}%
\pgfpathquadraticcurveto{\pgfqpoint{8.159804in}{0.841668in}}{\pgfqpoint{8.159804in}{0.832428in}}%
\pgfpathlineto{\pgfqpoint{8.159804in}{0.796800in}}%
\pgfpathlineto{\pgfqpoint{8.149393in}{0.796800in}}%
\pgfpathlineto{\pgfqpoint{8.149393in}{0.884393in}}%
\pgfpathlineto{\pgfqpoint{8.159804in}{0.884393in}}%
\pgfpathlineto{\pgfqpoint{8.159804in}{0.850044in}}%
\pgfpathquadraticcurveto{\pgfqpoint{8.163532in}{0.855736in}}{\pgfqpoint{8.168558in}{0.858546in}}%
\pgfpathquadraticcurveto{\pgfqpoint{8.173601in}{0.861356in}}{\pgfqpoint{8.180194in}{0.861356in}}%
\pgfpathquadraticcurveto{\pgfqpoint{8.191055in}{0.861356in}}{\pgfqpoint{8.196621in}{0.854637in}}%
\pgfpathquadraticcurveto{\pgfqpoint{8.202204in}{0.847918in}}{\pgfqpoint{8.202204in}{0.834860in}}%
\pgfpathlineto{\pgfqpoint{8.202204in}{0.834860in}}%
\pgfpathclose%
\pgfpathmoveto{\pgfqpoint{8.276775in}{0.830915in}}%
\pgfpathlineto{\pgfqpoint{8.276775in}{0.825854in}}%
\pgfpathlineto{\pgfqpoint{8.229151in}{0.825854in}}%
\pgfpathquadraticcurveto{\pgfqpoint{8.229835in}{0.815154in}}{\pgfqpoint{8.235599in}{0.809553in}}%
\pgfpathquadraticcurveto{\pgfqpoint{8.241363in}{0.803951in}}{\pgfqpoint{8.251666in}{0.803951in}}%
\pgfpathquadraticcurveto{\pgfqpoint{8.257628in}{0.803951in}}{\pgfqpoint{8.263230in}{0.805410in}}%
\pgfpathquadraticcurveto{\pgfqpoint{8.268831in}{0.806869in}}{\pgfqpoint{8.274361in}{0.809805in}}%
\pgfpathlineto{\pgfqpoint{8.274361in}{0.800006in}}%
\pgfpathquadraticcurveto{\pgfqpoint{8.268777in}{0.797647in}}{\pgfqpoint{8.262923in}{0.796404in}}%
\pgfpathquadraticcurveto{\pgfqpoint{8.257069in}{0.795161in}}{\pgfqpoint{8.251053in}{0.795161in}}%
\pgfpathquadraticcurveto{\pgfqpoint{8.235959in}{0.795161in}}{\pgfqpoint{8.227151in}{0.803933in}}%
\pgfpathquadraticcurveto{\pgfqpoint{8.218343in}{0.812723in}}{\pgfqpoint{8.218343in}{0.827709in}}%
\pgfpathquadraticcurveto{\pgfqpoint{8.218343in}{0.843181in}}{\pgfqpoint{8.226701in}{0.852259in}}%
\pgfpathquadraticcurveto{\pgfqpoint{8.235059in}{0.861356in}}{\pgfqpoint{8.249252in}{0.861356in}}%
\pgfpathquadraticcurveto{\pgfqpoint{8.261969in}{0.861356in}}{\pgfqpoint{8.269372in}{0.853160in}}%
\pgfpathquadraticcurveto{\pgfqpoint{8.276775in}{0.844983in}}{\pgfqpoint{8.276775in}{0.830915in}}%
\pgfpathlineto{\pgfqpoint{8.276775in}{0.830915in}}%
\pgfpathclose%
\pgfpathmoveto{\pgfqpoint{8.266418in}{0.833959in}}%
\pgfpathquadraticcurveto{\pgfqpoint{8.266310in}{0.842443in}}{\pgfqpoint{8.261663in}{0.847504in}}%
\pgfpathquadraticcurveto{\pgfqpoint{8.257015in}{0.852584in}}{\pgfqpoint{8.249360in}{0.852584in}}%
\pgfpathquadraticcurveto{\pgfqpoint{8.240696in}{0.852584in}}{\pgfqpoint{8.235491in}{0.847684in}}%
\pgfpathquadraticcurveto{\pgfqpoint{8.230285in}{0.842785in}}{\pgfqpoint{8.229493in}{0.833887in}}%
\pgfpathlineto{\pgfqpoint{8.266418in}{0.833959in}}%
\pgfpathlineto{\pgfqpoint{8.266418in}{0.833959in}}%
\pgfpathclose%
\pgfpathmoveto{\pgfqpoint{8.330307in}{0.850170in}}%
\pgfpathquadraticcurveto{\pgfqpoint{8.328560in}{0.851179in}}{\pgfqpoint{8.326506in}{0.851647in}}%
\pgfpathquadraticcurveto{\pgfqpoint{8.324453in}{0.852133in}}{\pgfqpoint{8.321985in}{0.852133in}}%
\pgfpathquadraticcurveto{\pgfqpoint{8.313195in}{0.852133in}}{\pgfqpoint{8.308494in}{0.846423in}}%
\pgfpathquadraticcurveto{\pgfqpoint{8.303793in}{0.840714in}}{\pgfqpoint{8.303793in}{0.830014in}}%
\pgfpathlineto{\pgfqpoint{8.303793in}{0.796800in}}%
\pgfpathlineto{\pgfqpoint{8.293382in}{0.796800in}}%
\pgfpathlineto{\pgfqpoint{8.293382in}{0.859843in}}%
\pgfpathlineto{\pgfqpoint{8.303793in}{0.859843in}}%
\pgfpathlineto{\pgfqpoint{8.303793in}{0.850044in}}%
\pgfpathquadraticcurveto{\pgfqpoint{8.307071in}{0.855790in}}{\pgfqpoint{8.312295in}{0.858564in}}%
\pgfpathquadraticcurveto{\pgfqpoint{8.317536in}{0.861356in}}{\pgfqpoint{8.325029in}{0.861356in}}%
\pgfpathquadraticcurveto{\pgfqpoint{8.326092in}{0.861356in}}{\pgfqpoint{8.327389in}{0.861211in}}%
\pgfpathquadraticcurveto{\pgfqpoint{8.328686in}{0.861085in}}{\pgfqpoint{8.330253in}{0.860797in}}%
\pgfpathlineto{\pgfqpoint{8.330307in}{0.850170in}}%
\pgfpathlineto{\pgfqpoint{8.330307in}{0.850170in}}%
\pgfpathclose%
\pgfusepath{fill}%
\end{pgfscope}%
\begin{pgfscope}%
\pgfsetbuttcap%
\pgfsetroundjoin%
\definecolor{currentfill}{rgb}{0.309804,0.372549,0.419608}%
\pgfsetfillcolor{currentfill}%
\pgfsetlinewidth{0.000000pt}%
\definecolor{currentstroke}{rgb}{0.309804,0.372549,0.419608}%
\pgfsetstrokecolor{currentstroke}%
\pgfsetdash{}{0pt}%
\pgfpathmoveto{\pgfqpoint{8.459272in}{0.830915in}}%
\pgfpathlineto{\pgfqpoint{8.459272in}{0.825854in}}%
\pgfpathlineto{\pgfqpoint{8.411648in}{0.825854in}}%
\pgfpathquadraticcurveto{\pgfqpoint{8.412333in}{0.815154in}}{\pgfqpoint{8.418097in}{0.809553in}}%
\pgfpathquadraticcurveto{\pgfqpoint{8.423860in}{0.803951in}}{\pgfqpoint{8.434163in}{0.803951in}}%
\pgfpathquadraticcurveto{\pgfqpoint{8.440125in}{0.803951in}}{\pgfqpoint{8.445727in}{0.805410in}}%
\pgfpathquadraticcurveto{\pgfqpoint{8.451329in}{0.806869in}}{\pgfqpoint{8.456859in}{0.809805in}}%
\pgfpathlineto{\pgfqpoint{8.456859in}{0.800006in}}%
\pgfpathquadraticcurveto{\pgfqpoint{8.451275in}{0.797647in}}{\pgfqpoint{8.445421in}{0.796404in}}%
\pgfpathquadraticcurveto{\pgfqpoint{8.439567in}{0.795161in}}{\pgfqpoint{8.433551in}{0.795161in}}%
\pgfpathquadraticcurveto{\pgfqpoint{8.418457in}{0.795161in}}{\pgfqpoint{8.409649in}{0.803933in}}%
\pgfpathquadraticcurveto{\pgfqpoint{8.400841in}{0.812723in}}{\pgfqpoint{8.400841in}{0.827709in}}%
\pgfpathquadraticcurveto{\pgfqpoint{8.400841in}{0.843181in}}{\pgfqpoint{8.409199in}{0.852259in}}%
\pgfpathquadraticcurveto{\pgfqpoint{8.417556in}{0.861356in}}{\pgfqpoint{8.431750in}{0.861356in}}%
\pgfpathquadraticcurveto{\pgfqpoint{8.444466in}{0.861356in}}{\pgfqpoint{8.451869in}{0.853160in}}%
\pgfpathquadraticcurveto{\pgfqpoint{8.459272in}{0.844983in}}{\pgfqpoint{8.459272in}{0.830915in}}%
\pgfpathlineto{\pgfqpoint{8.459272in}{0.830915in}}%
\pgfpathclose%
\pgfpathmoveto{\pgfqpoint{8.448915in}{0.833959in}}%
\pgfpathquadraticcurveto{\pgfqpoint{8.448807in}{0.842443in}}{\pgfqpoint{8.444160in}{0.847504in}}%
\pgfpathquadraticcurveto{\pgfqpoint{8.439513in}{0.852584in}}{\pgfqpoint{8.431858in}{0.852584in}}%
\pgfpathquadraticcurveto{\pgfqpoint{8.423194in}{0.852584in}}{\pgfqpoint{8.417988in}{0.847684in}}%
\pgfpathquadraticcurveto{\pgfqpoint{8.412783in}{0.842785in}}{\pgfqpoint{8.411990in}{0.833887in}}%
\pgfpathlineto{\pgfqpoint{8.448915in}{0.833959in}}%
\pgfpathlineto{\pgfqpoint{8.448915in}{0.833959in}}%
\pgfpathclose%
\pgfpathmoveto{\pgfqpoint{8.516461in}{0.857987in}}%
\pgfpathlineto{\pgfqpoint{8.516461in}{0.848189in}}%
\pgfpathquadraticcurveto{\pgfqpoint{8.512084in}{0.850440in}}{\pgfqpoint{8.507347in}{0.851557in}}%
\pgfpathquadraticcurveto{\pgfqpoint{8.502628in}{0.852692in}}{\pgfqpoint{8.497548in}{0.852692in}}%
\pgfpathquadraticcurveto{\pgfqpoint{8.489839in}{0.852692in}}{\pgfqpoint{8.485984in}{0.850332in}}%
\pgfpathquadraticcurveto{\pgfqpoint{8.482130in}{0.847973in}}{\pgfqpoint{8.482130in}{0.843235in}}%
\pgfpathquadraticcurveto{\pgfqpoint{8.482130in}{0.839633in}}{\pgfqpoint{8.484886in}{0.837580in}}%
\pgfpathquadraticcurveto{\pgfqpoint{8.487641in}{0.835526in}}{\pgfqpoint{8.495981in}{0.833671in}}%
\pgfpathlineto{\pgfqpoint{8.499530in}{0.832878in}}%
\pgfpathquadraticcurveto{\pgfqpoint{8.510553in}{0.830519in}}{\pgfqpoint{8.515200in}{0.826214in}}%
\pgfpathquadraticcurveto{\pgfqpoint{8.519847in}{0.821909in}}{\pgfqpoint{8.519847in}{0.814200in}}%
\pgfpathquadraticcurveto{\pgfqpoint{8.519847in}{0.805410in}}{\pgfqpoint{8.512895in}{0.800276in}}%
\pgfpathquadraticcurveto{\pgfqpoint{8.505942in}{0.795161in}}{\pgfqpoint{8.493784in}{0.795161in}}%
\pgfpathquadraticcurveto{\pgfqpoint{8.488722in}{0.795161in}}{\pgfqpoint{8.483229in}{0.796152in}}%
\pgfpathquadraticcurveto{\pgfqpoint{8.477735in}{0.797142in}}{\pgfqpoint{8.471665in}{0.799106in}}%
\pgfpathlineto{\pgfqpoint{8.471665in}{0.809805in}}%
\pgfpathquadraticcurveto{\pgfqpoint{8.477411in}{0.806815in}}{\pgfqpoint{8.482976in}{0.805320in}}%
\pgfpathquadraticcurveto{\pgfqpoint{8.488542in}{0.803843in}}{\pgfqpoint{8.494018in}{0.803843in}}%
\pgfpathquadraticcurveto{\pgfqpoint{8.501331in}{0.803843in}}{\pgfqpoint{8.505257in}{0.806346in}}%
\pgfpathquadraticcurveto{\pgfqpoint{8.509202in}{0.808850in}}{\pgfqpoint{8.509202in}{0.813407in}}%
\pgfpathquadraticcurveto{\pgfqpoint{8.509202in}{0.817622in}}{\pgfqpoint{8.506356in}{0.819874in}}%
\pgfpathquadraticcurveto{\pgfqpoint{8.503528in}{0.822125in}}{\pgfqpoint{8.493892in}{0.824214in}}%
\pgfpathlineto{\pgfqpoint{8.490289in}{0.825061in}}%
\pgfpathquadraticcurveto{\pgfqpoint{8.480671in}{0.827078in}}{\pgfqpoint{8.476384in}{0.831275in}}%
\pgfpathquadraticcurveto{\pgfqpoint{8.472115in}{0.835472in}}{\pgfqpoint{8.472115in}{0.842785in}}%
\pgfpathquadraticcurveto{\pgfqpoint{8.472115in}{0.851683in}}{\pgfqpoint{8.478419in}{0.856510in}}%
\pgfpathquadraticcurveto{\pgfqpoint{8.484724in}{0.861356in}}{\pgfqpoint{8.496323in}{0.861356in}}%
\pgfpathquadraticcurveto{\pgfqpoint{8.502051in}{0.861356in}}{\pgfqpoint{8.507113in}{0.860509in}}%
\pgfpathquadraticcurveto{\pgfqpoint{8.512192in}{0.859680in}}{\pgfqpoint{8.516461in}{0.857987in}}%
\pgfpathlineto{\pgfqpoint{8.516461in}{0.857987in}}%
\pgfpathclose%
\pgfpathmoveto{\pgfqpoint{8.546577in}{0.877747in}}%
\pgfpathlineto{\pgfqpoint{8.546577in}{0.859843in}}%
\pgfpathlineto{\pgfqpoint{8.567904in}{0.859843in}}%
\pgfpathlineto{\pgfqpoint{8.567904in}{0.851791in}}%
\pgfpathlineto{\pgfqpoint{8.546577in}{0.851791in}}%
\pgfpathlineto{\pgfqpoint{8.546577in}{0.817568in}}%
\pgfpathquadraticcurveto{\pgfqpoint{8.546577in}{0.809859in}}{\pgfqpoint{8.548685in}{0.807661in}}%
\pgfpathquadraticcurveto{\pgfqpoint{8.550792in}{0.805464in}}{\pgfqpoint{8.557276in}{0.805464in}}%
\pgfpathlineto{\pgfqpoint{8.567904in}{0.805464in}}%
\pgfpathlineto{\pgfqpoint{8.567904in}{0.796800in}}%
\pgfpathlineto{\pgfqpoint{8.557276in}{0.796800in}}%
\pgfpathquadraticcurveto{\pgfqpoint{8.545280in}{0.796800in}}{\pgfqpoint{8.540723in}{0.801267in}}%
\pgfpathquadraticcurveto{\pgfqpoint{8.536166in}{0.805752in}}{\pgfqpoint{8.536166in}{0.817568in}}%
\pgfpathlineto{\pgfqpoint{8.536166in}{0.851791in}}%
\pgfpathlineto{\pgfqpoint{8.528565in}{0.851791in}}%
\pgfpathlineto{\pgfqpoint{8.528565in}{0.859843in}}%
\pgfpathlineto{\pgfqpoint{8.536166in}{0.859843in}}%
\pgfpathlineto{\pgfqpoint{8.536166in}{0.877747in}}%
\pgfpathlineto{\pgfqpoint{8.546577in}{0.877747in}}%
\pgfpathlineto{\pgfqpoint{8.546577in}{0.877747in}}%
\pgfpathclose%
\pgfpathmoveto{\pgfqpoint{8.581521in}{0.859843in}}%
\pgfpathlineto{\pgfqpoint{8.591878in}{0.859843in}}%
\pgfpathlineto{\pgfqpoint{8.591878in}{0.796800in}}%
\pgfpathlineto{\pgfqpoint{8.581521in}{0.796800in}}%
\pgfpathlineto{\pgfqpoint{8.581521in}{0.859843in}}%
\pgfpathlineto{\pgfqpoint{8.581521in}{0.859843in}}%
\pgfpathclose%
\pgfpathmoveto{\pgfqpoint{8.581521in}{0.884393in}}%
\pgfpathlineto{\pgfqpoint{8.591878in}{0.884393in}}%
\pgfpathlineto{\pgfqpoint{8.591878in}{0.871262in}}%
\pgfpathlineto{\pgfqpoint{8.581521in}{0.871262in}}%
\pgfpathlineto{\pgfqpoint{8.581521in}{0.884393in}}%
\pgfpathlineto{\pgfqpoint{8.581521in}{0.884393in}}%
\pgfpathclose%
\pgfpathmoveto{\pgfqpoint{8.662630in}{0.847738in}}%
\pgfpathquadraticcurveto{\pgfqpoint{8.666520in}{0.854727in}}{\pgfqpoint{8.671924in}{0.858041in}}%
\pgfpathquadraticcurveto{\pgfqpoint{8.677327in}{0.861356in}}{\pgfqpoint{8.684640in}{0.861356in}}%
\pgfpathquadraticcurveto{\pgfqpoint{8.694493in}{0.861356in}}{\pgfqpoint{8.699843in}{0.854457in}}%
\pgfpathquadraticcurveto{\pgfqpoint{8.705192in}{0.847576in}}{\pgfqpoint{8.705192in}{0.834860in}}%
\pgfpathlineto{\pgfqpoint{8.705192in}{0.796800in}}%
\pgfpathlineto{\pgfqpoint{8.694781in}{0.796800in}}%
\pgfpathlineto{\pgfqpoint{8.694781in}{0.834517in}}%
\pgfpathquadraticcurveto{\pgfqpoint{8.694781in}{0.843578in}}{\pgfqpoint{8.691557in}{0.847955in}}%
\pgfpathquadraticcurveto{\pgfqpoint{8.688351in}{0.852349in}}{\pgfqpoint{8.681776in}{0.852349in}}%
\pgfpathquadraticcurveto{\pgfqpoint{8.673725in}{0.852349in}}{\pgfqpoint{8.669042in}{0.847000in}}%
\pgfpathquadraticcurveto{\pgfqpoint{8.664377in}{0.841668in}}{\pgfqpoint{8.664377in}{0.832428in}}%
\pgfpathlineto{\pgfqpoint{8.664377in}{0.796800in}}%
\pgfpathlineto{\pgfqpoint{8.653966in}{0.796800in}}%
\pgfpathlineto{\pgfqpoint{8.653966in}{0.834517in}}%
\pgfpathquadraticcurveto{\pgfqpoint{8.653966in}{0.843632in}}{\pgfqpoint{8.650760in}{0.847991in}}%
\pgfpathquadraticcurveto{\pgfqpoint{8.647553in}{0.852349in}}{\pgfqpoint{8.640853in}{0.852349in}}%
\pgfpathquadraticcurveto{\pgfqpoint{8.632910in}{0.852349in}}{\pgfqpoint{8.628226in}{0.846982in}}%
\pgfpathquadraticcurveto{\pgfqpoint{8.623561in}{0.841614in}}{\pgfqpoint{8.623561in}{0.832428in}}%
\pgfpathlineto{\pgfqpoint{8.623561in}{0.796800in}}%
\pgfpathlineto{\pgfqpoint{8.613150in}{0.796800in}}%
\pgfpathlineto{\pgfqpoint{8.613150in}{0.859843in}}%
\pgfpathlineto{\pgfqpoint{8.623561in}{0.859843in}}%
\pgfpathlineto{\pgfqpoint{8.623561in}{0.850044in}}%
\pgfpathquadraticcurveto{\pgfqpoint{8.627110in}{0.855844in}}{\pgfqpoint{8.632063in}{0.858600in}}%
\pgfpathquadraticcurveto{\pgfqpoint{8.637016in}{0.861356in}}{\pgfqpoint{8.643825in}{0.861356in}}%
\pgfpathquadraticcurveto{\pgfqpoint{8.650706in}{0.861356in}}{\pgfqpoint{8.655515in}{0.857861in}}%
\pgfpathquadraticcurveto{\pgfqpoint{8.660324in}{0.854385in}}{\pgfqpoint{8.662630in}{0.847738in}}%
\pgfpathlineto{\pgfqpoint{8.662630in}{0.847738in}}%
\pgfpathclose%
\pgfpathmoveto{\pgfqpoint{8.754492in}{0.828483in}}%
\pgfpathquadraticcurveto{\pgfqpoint{8.741937in}{0.828483in}}{\pgfqpoint{8.737092in}{0.825619in}}%
\pgfpathquadraticcurveto{\pgfqpoint{8.732247in}{0.822756in}}{\pgfqpoint{8.732247in}{0.815821in}}%
\pgfpathquadraticcurveto{\pgfqpoint{8.732247in}{0.810309in}}{\pgfqpoint{8.735885in}{0.807067in}}%
\pgfpathquadraticcurveto{\pgfqpoint{8.739523in}{0.803843in}}{\pgfqpoint{8.745756in}{0.803843in}}%
\pgfpathquadraticcurveto{\pgfqpoint{8.754383in}{0.803843in}}{\pgfqpoint{8.759589in}{0.809949in}}%
\pgfpathquadraticcurveto{\pgfqpoint{8.764794in}{0.816055in}}{\pgfqpoint{8.764794in}{0.826178in}}%
\pgfpathlineto{\pgfqpoint{8.764794in}{0.828483in}}%
\pgfpathlineto{\pgfqpoint{8.754492in}{0.828483in}}%
\pgfpathlineto{\pgfqpoint{8.754492in}{0.828483in}}%
\pgfpathclose%
\pgfpathmoveto{\pgfqpoint{8.775151in}{0.832770in}}%
\pgfpathlineto{\pgfqpoint{8.775151in}{0.796800in}}%
\pgfpathlineto{\pgfqpoint{8.764794in}{0.796800in}}%
\pgfpathlineto{\pgfqpoint{8.764794in}{0.806364in}}%
\pgfpathquadraticcurveto{\pgfqpoint{8.761246in}{0.800637in}}{\pgfqpoint{8.755951in}{0.797899in}}%
\pgfpathquadraticcurveto{\pgfqpoint{8.750655in}{0.795161in}}{\pgfqpoint{8.743000in}{0.795161in}}%
\pgfpathquadraticcurveto{\pgfqpoint{8.733327in}{0.795161in}}{\pgfqpoint{8.727599in}{0.800601in}}%
\pgfpathquadraticcurveto{\pgfqpoint{8.721890in}{0.806040in}}{\pgfqpoint{8.721890in}{0.815154in}}%
\pgfpathquadraticcurveto{\pgfqpoint{8.721890in}{0.825782in}}{\pgfqpoint{8.729004in}{0.831185in}}%
\pgfpathquadraticcurveto{\pgfqpoint{8.736137in}{0.836589in}}{\pgfqpoint{8.750259in}{0.836589in}}%
\pgfpathlineto{\pgfqpoint{8.764794in}{0.836589in}}%
\pgfpathlineto{\pgfqpoint{8.764794in}{0.837616in}}%
\pgfpathquadraticcurveto{\pgfqpoint{8.764794in}{0.844766in}}{\pgfqpoint{8.760093in}{0.848675in}}%
\pgfpathquadraticcurveto{\pgfqpoint{8.755392in}{0.852584in}}{\pgfqpoint{8.746890in}{0.852584in}}%
\pgfpathquadraticcurveto{\pgfqpoint{8.741487in}{0.852584in}}{\pgfqpoint{8.736353in}{0.851287in}}%
\pgfpathquadraticcurveto{\pgfqpoint{8.731238in}{0.849990in}}{\pgfqpoint{8.726519in}{0.847396in}}%
\pgfpathlineto{\pgfqpoint{8.726519in}{0.856979in}}%
\pgfpathquadraticcurveto{\pgfqpoint{8.732192in}{0.859176in}}{\pgfqpoint{8.737542in}{0.860257in}}%
\pgfpathquadraticcurveto{\pgfqpoint{8.742892in}{0.861356in}}{\pgfqpoint{8.747953in}{0.861356in}}%
\pgfpathquadraticcurveto{\pgfqpoint{8.761642in}{0.861356in}}{\pgfqpoint{8.768397in}{0.854259in}}%
\pgfpathquadraticcurveto{\pgfqpoint{8.775151in}{0.847180in}}{\pgfqpoint{8.775151in}{0.832770in}}%
\pgfpathlineto{\pgfqpoint{8.775151in}{0.832770in}}%
\pgfpathclose%
\pgfpathmoveto{\pgfqpoint{8.806727in}{0.877747in}}%
\pgfpathlineto{\pgfqpoint{8.806727in}{0.859843in}}%
\pgfpathlineto{\pgfqpoint{8.828053in}{0.859843in}}%
\pgfpathlineto{\pgfqpoint{8.828053in}{0.851791in}}%
\pgfpathlineto{\pgfqpoint{8.806727in}{0.851791in}}%
\pgfpathlineto{\pgfqpoint{8.806727in}{0.817568in}}%
\pgfpathquadraticcurveto{\pgfqpoint{8.806727in}{0.809859in}}{\pgfqpoint{8.808834in}{0.807661in}}%
\pgfpathquadraticcurveto{\pgfqpoint{8.810942in}{0.805464in}}{\pgfqpoint{8.817426in}{0.805464in}}%
\pgfpathlineto{\pgfqpoint{8.828053in}{0.805464in}}%
\pgfpathlineto{\pgfqpoint{8.828053in}{0.796800in}}%
\pgfpathlineto{\pgfqpoint{8.817426in}{0.796800in}}%
\pgfpathquadraticcurveto{\pgfqpoint{8.805430in}{0.796800in}}{\pgfqpoint{8.800873in}{0.801267in}}%
\pgfpathquadraticcurveto{\pgfqpoint{8.796316in}{0.805752in}}{\pgfqpoint{8.796316in}{0.817568in}}%
\pgfpathlineto{\pgfqpoint{8.796316in}{0.851791in}}%
\pgfpathlineto{\pgfqpoint{8.788715in}{0.851791in}}%
\pgfpathlineto{\pgfqpoint{8.788715in}{0.859843in}}%
\pgfpathlineto{\pgfqpoint{8.796316in}{0.859843in}}%
\pgfpathlineto{\pgfqpoint{8.796316in}{0.877747in}}%
\pgfpathlineto{\pgfqpoint{8.806727in}{0.877747in}}%
\pgfpathlineto{\pgfqpoint{8.806727in}{0.877747in}}%
\pgfpathclose%
\pgfpathmoveto{\pgfqpoint{8.895599in}{0.830915in}}%
\pgfpathlineto{\pgfqpoint{8.895599in}{0.825854in}}%
\pgfpathlineto{\pgfqpoint{8.847975in}{0.825854in}}%
\pgfpathquadraticcurveto{\pgfqpoint{8.848659in}{0.815154in}}{\pgfqpoint{8.854423in}{0.809553in}}%
\pgfpathquadraticcurveto{\pgfqpoint{8.860187in}{0.803951in}}{\pgfqpoint{8.870490in}{0.803951in}}%
\pgfpathquadraticcurveto{\pgfqpoint{8.876452in}{0.803951in}}{\pgfqpoint{8.882054in}{0.805410in}}%
\pgfpathquadraticcurveto{\pgfqpoint{8.887655in}{0.806869in}}{\pgfqpoint{8.893185in}{0.809805in}}%
\pgfpathlineto{\pgfqpoint{8.893185in}{0.800006in}}%
\pgfpathquadraticcurveto{\pgfqpoint{8.887601in}{0.797647in}}{\pgfqpoint{8.881747in}{0.796404in}}%
\pgfpathquadraticcurveto{\pgfqpoint{8.875893in}{0.795161in}}{\pgfqpoint{8.869877in}{0.795161in}}%
\pgfpathquadraticcurveto{\pgfqpoint{8.854783in}{0.795161in}}{\pgfqpoint{8.845975in}{0.803933in}}%
\pgfpathquadraticcurveto{\pgfqpoint{8.837167in}{0.812723in}}{\pgfqpoint{8.837167in}{0.827709in}}%
\pgfpathquadraticcurveto{\pgfqpoint{8.837167in}{0.843181in}}{\pgfqpoint{8.845525in}{0.852259in}}%
\pgfpathquadraticcurveto{\pgfqpoint{8.853883in}{0.861356in}}{\pgfqpoint{8.868076in}{0.861356in}}%
\pgfpathquadraticcurveto{\pgfqpoint{8.880793in}{0.861356in}}{\pgfqpoint{8.888196in}{0.853160in}}%
\pgfpathquadraticcurveto{\pgfqpoint{8.895599in}{0.844983in}}{\pgfqpoint{8.895599in}{0.830915in}}%
\pgfpathlineto{\pgfqpoint{8.895599in}{0.830915in}}%
\pgfpathclose%
\pgfpathmoveto{\pgfqpoint{8.885242in}{0.833959in}}%
\pgfpathquadraticcurveto{\pgfqpoint{8.885134in}{0.842443in}}{\pgfqpoint{8.880487in}{0.847504in}}%
\pgfpathquadraticcurveto{\pgfqpoint{8.875839in}{0.852584in}}{\pgfqpoint{8.868184in}{0.852584in}}%
\pgfpathquadraticcurveto{\pgfqpoint{8.859520in}{0.852584in}}{\pgfqpoint{8.854315in}{0.847684in}}%
\pgfpathquadraticcurveto{\pgfqpoint{8.849109in}{0.842785in}}{\pgfqpoint{8.848317in}{0.833887in}}%
\pgfpathlineto{\pgfqpoint{8.885242in}{0.833959in}}%
\pgfpathlineto{\pgfqpoint{8.885242in}{0.833959in}}%
\pgfpathclose%
\pgfusepath{fill}%
\end{pgfscope}%
\begin{pgfscope}%
\pgfsetbuttcap%
\pgfsetroundjoin%
\definecolor{currentfill}{rgb}{0.309804,0.372549,0.419608}%
\pgfsetfillcolor{currentfill}%
\pgfsetlinewidth{0.000000pt}%
\definecolor{currentstroke}{rgb}{0.309804,0.372549,0.419608}%
\pgfsetstrokecolor{currentstroke}%
\pgfsetdash{}{0pt}%
\pgfpathmoveto{\pgfqpoint{9.003473in}{0.852584in}}%
\pgfpathquadraticcurveto{\pgfqpoint{8.995152in}{0.852584in}}{\pgfqpoint{8.990306in}{0.846081in}}%
\pgfpathquadraticcurveto{\pgfqpoint{8.985461in}{0.839579in}}{\pgfqpoint{8.985461in}{0.828267in}}%
\pgfpathquadraticcurveto{\pgfqpoint{8.985461in}{0.816956in}}{\pgfqpoint{8.990270in}{0.810453in}}%
\pgfpathquadraticcurveto{\pgfqpoint{8.995098in}{0.803951in}}{\pgfqpoint{9.003473in}{0.803951in}}%
\pgfpathquadraticcurveto{\pgfqpoint{9.011759in}{0.803951in}}{\pgfqpoint{9.016586in}{0.810471in}}%
\pgfpathquadraticcurveto{\pgfqpoint{9.021431in}{0.817010in}}{\pgfqpoint{9.021431in}{0.828267in}}%
\pgfpathquadraticcurveto{\pgfqpoint{9.021431in}{0.839471in}}{\pgfqpoint{9.016586in}{0.846027in}}%
\pgfpathquadraticcurveto{\pgfqpoint{9.011759in}{0.852584in}}{\pgfqpoint{9.003473in}{0.852584in}}%
\pgfpathlineto{\pgfqpoint{9.003473in}{0.852584in}}%
\pgfpathclose%
\pgfpathmoveto{\pgfqpoint{9.003473in}{0.861356in}}%
\pgfpathquadraticcurveto{\pgfqpoint{9.016982in}{0.861356in}}{\pgfqpoint{9.024692in}{0.852566in}}%
\pgfpathquadraticcurveto{\pgfqpoint{9.032419in}{0.843794in}}{\pgfqpoint{9.032419in}{0.828267in}}%
\pgfpathquadraticcurveto{\pgfqpoint{9.032419in}{0.812795in}}{\pgfqpoint{9.024692in}{0.803969in}}%
\pgfpathquadraticcurveto{\pgfqpoint{9.016982in}{0.795161in}}{\pgfqpoint{9.003473in}{0.795161in}}%
\pgfpathquadraticcurveto{\pgfqpoint{8.989910in}{0.795161in}}{\pgfqpoint{8.982219in}{0.803969in}}%
\pgfpathquadraticcurveto{\pgfqpoint{8.974546in}{0.812795in}}{\pgfqpoint{8.974546in}{0.828267in}}%
\pgfpathquadraticcurveto{\pgfqpoint{8.974546in}{0.843794in}}{\pgfqpoint{8.982219in}{0.852566in}}%
\pgfpathquadraticcurveto{\pgfqpoint{8.989910in}{0.861356in}}{\pgfqpoint{9.003473in}{0.861356in}}%
\pgfpathlineto{\pgfqpoint{9.003473in}{0.861356in}}%
\pgfpathclose%
\pgfpathmoveto{\pgfqpoint{9.086113in}{0.850170in}}%
\pgfpathquadraticcurveto{\pgfqpoint{9.084366in}{0.851179in}}{\pgfqpoint{9.082312in}{0.851647in}}%
\pgfpathquadraticcurveto{\pgfqpoint{9.080259in}{0.852133in}}{\pgfqpoint{9.077791in}{0.852133in}}%
\pgfpathquadraticcurveto{\pgfqpoint{9.069002in}{0.852133in}}{\pgfqpoint{9.064300in}{0.846423in}}%
\pgfpathquadraticcurveto{\pgfqpoint{9.059599in}{0.840714in}}{\pgfqpoint{9.059599in}{0.830014in}}%
\pgfpathlineto{\pgfqpoint{9.059599in}{0.796800in}}%
\pgfpathlineto{\pgfqpoint{9.049188in}{0.796800in}}%
\pgfpathlineto{\pgfqpoint{9.049188in}{0.859843in}}%
\pgfpathlineto{\pgfqpoint{9.059599in}{0.859843in}}%
\pgfpathlineto{\pgfqpoint{9.059599in}{0.850044in}}%
\pgfpathquadraticcurveto{\pgfqpoint{9.062877in}{0.855790in}}{\pgfqpoint{9.068101in}{0.858564in}}%
\pgfpathquadraticcurveto{\pgfqpoint{9.073342in}{0.861356in}}{\pgfqpoint{9.080836in}{0.861356in}}%
\pgfpathquadraticcurveto{\pgfqpoint{9.081898in}{0.861356in}}{\pgfqpoint{9.083195in}{0.861211in}}%
\pgfpathquadraticcurveto{\pgfqpoint{9.084492in}{0.861085in}}{\pgfqpoint{9.086059in}{0.860797in}}%
\pgfpathlineto{\pgfqpoint{9.086113in}{0.850170in}}%
\pgfpathlineto{\pgfqpoint{9.086113in}{0.850170in}}%
\pgfpathclose%
\pgfusepath{fill}%
\end{pgfscope}%
\begin{pgfscope}%
\pgfsetbuttcap%
\pgfsetroundjoin%
\definecolor{currentfill}{rgb}{0.309804,0.372549,0.419608}%
\pgfsetfillcolor{currentfill}%
\pgfsetlinewidth{0.000000pt}%
\definecolor{currentstroke}{rgb}{0.309804,0.372549,0.419608}%
\pgfsetstrokecolor{currentstroke}%
\pgfsetdash{}{0pt}%
\pgfpathmoveto{\pgfqpoint{9.163003in}{0.877747in}}%
\pgfpathlineto{\pgfqpoint{9.163003in}{0.859843in}}%
\pgfpathlineto{\pgfqpoint{9.184329in}{0.859843in}}%
\pgfpathlineto{\pgfqpoint{9.184329in}{0.851791in}}%
\pgfpathlineto{\pgfqpoint{9.163003in}{0.851791in}}%
\pgfpathlineto{\pgfqpoint{9.163003in}{0.817568in}}%
\pgfpathquadraticcurveto{\pgfqpoint{9.163003in}{0.809859in}}{\pgfqpoint{9.165110in}{0.807661in}}%
\pgfpathquadraticcurveto{\pgfqpoint{9.167217in}{0.805464in}}{\pgfqpoint{9.173702in}{0.805464in}}%
\pgfpathlineto{\pgfqpoint{9.184329in}{0.805464in}}%
\pgfpathlineto{\pgfqpoint{9.184329in}{0.796800in}}%
\pgfpathlineto{\pgfqpoint{9.173702in}{0.796800in}}%
\pgfpathquadraticcurveto{\pgfqpoint{9.161706in}{0.796800in}}{\pgfqpoint{9.157149in}{0.801267in}}%
\pgfpathquadraticcurveto{\pgfqpoint{9.152592in}{0.805752in}}{\pgfqpoint{9.152592in}{0.817568in}}%
\pgfpathlineto{\pgfqpoint{9.152592in}{0.851791in}}%
\pgfpathlineto{\pgfqpoint{9.144990in}{0.851791in}}%
\pgfpathlineto{\pgfqpoint{9.144990in}{0.859843in}}%
\pgfpathlineto{\pgfqpoint{9.152592in}{0.859843in}}%
\pgfpathlineto{\pgfqpoint{9.152592in}{0.877747in}}%
\pgfpathlineto{\pgfqpoint{9.163003in}{0.877747in}}%
\pgfpathlineto{\pgfqpoint{9.163003in}{0.877747in}}%
\pgfpathclose%
\pgfpathmoveto{\pgfqpoint{9.250362in}{0.834860in}}%
\pgfpathlineto{\pgfqpoint{9.250362in}{0.796800in}}%
\pgfpathlineto{\pgfqpoint{9.240005in}{0.796800in}}%
\pgfpathlineto{\pgfqpoint{9.240005in}{0.834517in}}%
\pgfpathquadraticcurveto{\pgfqpoint{9.240005in}{0.843469in}}{\pgfqpoint{9.236510in}{0.847900in}}%
\pgfpathquadraticcurveto{\pgfqpoint{9.233016in}{0.852349in}}{\pgfqpoint{9.226045in}{0.852349in}}%
\pgfpathquadraticcurveto{\pgfqpoint{9.217651in}{0.852349in}}{\pgfqpoint{9.212806in}{0.847000in}}%
\pgfpathquadraticcurveto{\pgfqpoint{9.207961in}{0.841668in}}{\pgfqpoint{9.207961in}{0.832428in}}%
\pgfpathlineto{\pgfqpoint{9.207961in}{0.796800in}}%
\pgfpathlineto{\pgfqpoint{9.197550in}{0.796800in}}%
\pgfpathlineto{\pgfqpoint{9.197550in}{0.884393in}}%
\pgfpathlineto{\pgfqpoint{9.207961in}{0.884393in}}%
\pgfpathlineto{\pgfqpoint{9.207961in}{0.850044in}}%
\pgfpathquadraticcurveto{\pgfqpoint{9.211689in}{0.855736in}}{\pgfqpoint{9.216715in}{0.858546in}}%
\pgfpathquadraticcurveto{\pgfqpoint{9.221758in}{0.861356in}}{\pgfqpoint{9.228351in}{0.861356in}}%
\pgfpathquadraticcurveto{\pgfqpoint{9.239212in}{0.861356in}}{\pgfqpoint{9.244778in}{0.854637in}}%
\pgfpathquadraticcurveto{\pgfqpoint{9.250362in}{0.847918in}}{\pgfqpoint{9.250362in}{0.834860in}}%
\pgfpathlineto{\pgfqpoint{9.250362in}{0.834860in}}%
\pgfpathclose%
\pgfpathmoveto{\pgfqpoint{9.324932in}{0.830915in}}%
\pgfpathlineto{\pgfqpoint{9.324932in}{0.825854in}}%
\pgfpathlineto{\pgfqpoint{9.277308in}{0.825854in}}%
\pgfpathquadraticcurveto{\pgfqpoint{9.277992in}{0.815154in}}{\pgfqpoint{9.283756in}{0.809553in}}%
\pgfpathquadraticcurveto{\pgfqpoint{9.289520in}{0.803951in}}{\pgfqpoint{9.299823in}{0.803951in}}%
\pgfpathquadraticcurveto{\pgfqpoint{9.305785in}{0.803951in}}{\pgfqpoint{9.311387in}{0.805410in}}%
\pgfpathquadraticcurveto{\pgfqpoint{9.316988in}{0.806869in}}{\pgfqpoint{9.322518in}{0.809805in}}%
\pgfpathlineto{\pgfqpoint{9.322518in}{0.800006in}}%
\pgfpathquadraticcurveto{\pgfqpoint{9.316934in}{0.797647in}}{\pgfqpoint{9.311081in}{0.796404in}}%
\pgfpathquadraticcurveto{\pgfqpoint{9.305227in}{0.795161in}}{\pgfqpoint{9.299211in}{0.795161in}}%
\pgfpathquadraticcurveto{\pgfqpoint{9.284116in}{0.795161in}}{\pgfqpoint{9.275308in}{0.803933in}}%
\pgfpathquadraticcurveto{\pgfqpoint{9.266500in}{0.812723in}}{\pgfqpoint{9.266500in}{0.827709in}}%
\pgfpathquadraticcurveto{\pgfqpoint{9.266500in}{0.843181in}}{\pgfqpoint{9.274858in}{0.852259in}}%
\pgfpathquadraticcurveto{\pgfqpoint{9.283216in}{0.861356in}}{\pgfqpoint{9.297409in}{0.861356in}}%
\pgfpathquadraticcurveto{\pgfqpoint{9.310126in}{0.861356in}}{\pgfqpoint{9.317529in}{0.853160in}}%
\pgfpathquadraticcurveto{\pgfqpoint{9.324932in}{0.844983in}}{\pgfqpoint{9.324932in}{0.830915in}}%
\pgfpathlineto{\pgfqpoint{9.324932in}{0.830915in}}%
\pgfpathclose%
\pgfpathmoveto{\pgfqpoint{9.314575in}{0.833959in}}%
\pgfpathquadraticcurveto{\pgfqpoint{9.314467in}{0.842443in}}{\pgfqpoint{9.309820in}{0.847504in}}%
\pgfpathquadraticcurveto{\pgfqpoint{9.305173in}{0.852584in}}{\pgfqpoint{9.297517in}{0.852584in}}%
\pgfpathquadraticcurveto{\pgfqpoint{9.288854in}{0.852584in}}{\pgfqpoint{9.283648in}{0.847684in}}%
\pgfpathquadraticcurveto{\pgfqpoint{9.278442in}{0.842785in}}{\pgfqpoint{9.277650in}{0.833887in}}%
\pgfpathlineto{\pgfqpoint{9.314575in}{0.833959in}}%
\pgfpathlineto{\pgfqpoint{9.314575in}{0.833959in}}%
\pgfpathclose%
\pgfpathmoveto{\pgfqpoint{9.341935in}{0.859843in}}%
\pgfpathlineto{\pgfqpoint{9.352292in}{0.859843in}}%
\pgfpathlineto{\pgfqpoint{9.352292in}{0.796800in}}%
\pgfpathlineto{\pgfqpoint{9.341935in}{0.796800in}}%
\pgfpathlineto{\pgfqpoint{9.341935in}{0.859843in}}%
\pgfpathlineto{\pgfqpoint{9.341935in}{0.859843in}}%
\pgfpathclose%
\pgfpathmoveto{\pgfqpoint{9.341935in}{0.884393in}}%
\pgfpathlineto{\pgfqpoint{9.352292in}{0.884393in}}%
\pgfpathlineto{\pgfqpoint{9.352292in}{0.871262in}}%
\pgfpathlineto{\pgfqpoint{9.341935in}{0.871262in}}%
\pgfpathlineto{\pgfqpoint{9.341935in}{0.884393in}}%
\pgfpathlineto{\pgfqpoint{9.341935in}{0.884393in}}%
\pgfpathclose%
\pgfpathmoveto{\pgfqpoint{9.410490in}{0.850170in}}%
\pgfpathquadraticcurveto{\pgfqpoint{9.408742in}{0.851179in}}{\pgfqpoint{9.406689in}{0.851647in}}%
\pgfpathquadraticcurveto{\pgfqpoint{9.404636in}{0.852133in}}{\pgfqpoint{9.402168in}{0.852133in}}%
\pgfpathquadraticcurveto{\pgfqpoint{9.393378in}{0.852133in}}{\pgfqpoint{9.388677in}{0.846423in}}%
\pgfpathquadraticcurveto{\pgfqpoint{9.383976in}{0.840714in}}{\pgfqpoint{9.383976in}{0.830014in}}%
\pgfpathlineto{\pgfqpoint{9.383976in}{0.796800in}}%
\pgfpathlineto{\pgfqpoint{9.373565in}{0.796800in}}%
\pgfpathlineto{\pgfqpoint{9.373565in}{0.859843in}}%
\pgfpathlineto{\pgfqpoint{9.383976in}{0.859843in}}%
\pgfpathlineto{\pgfqpoint{9.383976in}{0.850044in}}%
\pgfpathquadraticcurveto{\pgfqpoint{9.387254in}{0.855790in}}{\pgfqpoint{9.392477in}{0.858564in}}%
\pgfpathquadraticcurveto{\pgfqpoint{9.397719in}{0.861356in}}{\pgfqpoint{9.405212in}{0.861356in}}%
\pgfpathquadraticcurveto{\pgfqpoint{9.406275in}{0.861356in}}{\pgfqpoint{9.407572in}{0.861211in}}%
\pgfpathquadraticcurveto{\pgfqpoint{9.408868in}{0.861085in}}{\pgfqpoint{9.410436in}{0.860797in}}%
\pgfpathlineto{\pgfqpoint{9.410490in}{0.850170in}}%
\pgfpathlineto{\pgfqpoint{9.410490in}{0.850170in}}%
\pgfpathclose%
\pgfusepath{fill}%
\end{pgfscope}%
\begin{pgfscope}%
\pgfsetbuttcap%
\pgfsetroundjoin%
\definecolor{currentfill}{rgb}{0.309804,0.372549,0.419608}%
\pgfsetfillcolor{currentfill}%
\pgfsetlinewidth{0.000000pt}%
\definecolor{currentstroke}{rgb}{0.309804,0.372549,0.419608}%
\pgfsetstrokecolor{currentstroke}%
\pgfsetdash{}{0pt}%
\pgfpathmoveto{\pgfqpoint{9.527941in}{0.850278in}}%
\pgfpathlineto{\pgfqpoint{9.527941in}{0.884393in}}%
\pgfpathlineto{\pgfqpoint{9.538298in}{0.884393in}}%
\pgfpathlineto{\pgfqpoint{9.538298in}{0.796800in}}%
\pgfpathlineto{\pgfqpoint{9.527941in}{0.796800in}}%
\pgfpathlineto{\pgfqpoint{9.527941in}{0.806256in}}%
\pgfpathquadraticcurveto{\pgfqpoint{9.524680in}{0.800637in}}{\pgfqpoint{9.519691in}{0.797899in}}%
\pgfpathquadraticcurveto{\pgfqpoint{9.514720in}{0.795161in}}{\pgfqpoint{9.507731in}{0.795161in}}%
\pgfpathquadraticcurveto{\pgfqpoint{9.496311in}{0.795161in}}{\pgfqpoint{9.489124in}{0.804275in}}%
\pgfpathquadraticcurveto{\pgfqpoint{9.481956in}{0.813407in}}{\pgfqpoint{9.481956in}{0.828267in}}%
\pgfpathquadraticcurveto{\pgfqpoint{9.481956in}{0.843127in}}{\pgfqpoint{9.489124in}{0.852241in}}%
\pgfpathquadraticcurveto{\pgfqpoint{9.496311in}{0.861356in}}{\pgfqpoint{9.507731in}{0.861356in}}%
\pgfpathquadraticcurveto{\pgfqpoint{9.514720in}{0.861356in}}{\pgfqpoint{9.519691in}{0.858618in}}%
\pgfpathquadraticcurveto{\pgfqpoint{9.524680in}{0.855898in}}{\pgfqpoint{9.527941in}{0.850278in}}%
\pgfpathlineto{\pgfqpoint{9.527941in}{0.850278in}}%
\pgfpathclose%
\pgfpathmoveto{\pgfqpoint{9.492655in}{0.828267in}}%
\pgfpathquadraticcurveto{\pgfqpoint{9.492655in}{0.816848in}}{\pgfqpoint{9.497356in}{0.810345in}}%
\pgfpathquadraticcurveto{\pgfqpoint{9.502057in}{0.803843in}}{\pgfqpoint{9.510271in}{0.803843in}}%
\pgfpathquadraticcurveto{\pgfqpoint{9.518484in}{0.803843in}}{\pgfqpoint{9.523203in}{0.810345in}}%
\pgfpathquadraticcurveto{\pgfqpoint{9.527941in}{0.816848in}}{\pgfqpoint{9.527941in}{0.828267in}}%
\pgfpathquadraticcurveto{\pgfqpoint{9.527941in}{0.839687in}}{\pgfqpoint{9.523203in}{0.846189in}}%
\pgfpathquadraticcurveto{\pgfqpoint{9.518484in}{0.852692in}}{\pgfqpoint{9.510271in}{0.852692in}}%
\pgfpathquadraticcurveto{\pgfqpoint{9.502057in}{0.852692in}}{\pgfqpoint{9.497356in}{0.846189in}}%
\pgfpathquadraticcurveto{\pgfqpoint{9.492655in}{0.839687in}}{\pgfqpoint{9.492655in}{0.828267in}}%
\pgfpathlineto{\pgfqpoint{9.492655in}{0.828267in}}%
\pgfpathclose%
\pgfpathmoveto{\pgfqpoint{9.559642in}{0.859843in}}%
\pgfpathlineto{\pgfqpoint{9.569999in}{0.859843in}}%
\pgfpathlineto{\pgfqpoint{9.569999in}{0.796800in}}%
\pgfpathlineto{\pgfqpoint{9.559642in}{0.796800in}}%
\pgfpathlineto{\pgfqpoint{9.559642in}{0.859843in}}%
\pgfpathlineto{\pgfqpoint{9.559642in}{0.859843in}}%
\pgfpathclose%
\pgfpathmoveto{\pgfqpoint{9.559642in}{0.884393in}}%
\pgfpathlineto{\pgfqpoint{9.569999in}{0.884393in}}%
\pgfpathlineto{\pgfqpoint{9.569999in}{0.871262in}}%
\pgfpathlineto{\pgfqpoint{9.559642in}{0.871262in}}%
\pgfpathlineto{\pgfqpoint{9.559642in}{0.884393in}}%
\pgfpathlineto{\pgfqpoint{9.559642in}{0.884393in}}%
\pgfpathclose%
\pgfpathmoveto{\pgfqpoint{9.662419in}{0.884393in}}%
\pgfpathlineto{\pgfqpoint{9.662419in}{0.875765in}}%
\pgfpathlineto{\pgfqpoint{9.652513in}{0.875765in}}%
\pgfpathquadraticcurveto{\pgfqpoint{9.646947in}{0.875765in}}{\pgfqpoint{9.644767in}{0.873514in}}%
\pgfpathquadraticcurveto{\pgfqpoint{9.642606in}{0.871262in}}{\pgfqpoint{9.642606in}{0.865408in}}%
\pgfpathlineto{\pgfqpoint{9.642606in}{0.859843in}}%
\pgfpathlineto{\pgfqpoint{9.659663in}{0.859843in}}%
\pgfpathlineto{\pgfqpoint{9.659663in}{0.851791in}}%
\pgfpathlineto{\pgfqpoint{9.642606in}{0.851791in}}%
\pgfpathlineto{\pgfqpoint{9.642606in}{0.796800in}}%
\pgfpathlineto{\pgfqpoint{9.632195in}{0.796800in}}%
\pgfpathlineto{\pgfqpoint{9.632195in}{0.851791in}}%
\pgfpathlineto{\pgfqpoint{9.603772in}{0.851791in}}%
\pgfpathlineto{\pgfqpoint{9.603772in}{0.796800in}}%
\pgfpathlineto{\pgfqpoint{9.593361in}{0.796800in}}%
\pgfpathlineto{\pgfqpoint{9.593361in}{0.851791in}}%
\pgfpathlineto{\pgfqpoint{9.583454in}{0.851791in}}%
\pgfpathlineto{\pgfqpoint{9.583454in}{0.859843in}}%
\pgfpathlineto{\pgfqpoint{9.593361in}{0.859843in}}%
\pgfpathlineto{\pgfqpoint{9.593361in}{0.864238in}}%
\pgfpathquadraticcurveto{\pgfqpoint{9.593361in}{0.874757in}}{\pgfqpoint{9.598260in}{0.879566in}}%
\pgfpathquadraticcurveto{\pgfqpoint{9.603159in}{0.884393in}}{\pgfqpoint{9.613786in}{0.884393in}}%
\pgfpathlineto{\pgfqpoint{9.623585in}{0.884393in}}%
\pgfpathlineto{\pgfqpoint{9.623585in}{0.875765in}}%
\pgfpathlineto{\pgfqpoint{9.613678in}{0.875765in}}%
\pgfpathquadraticcurveto{\pgfqpoint{9.608113in}{0.875765in}}{\pgfqpoint{9.605915in}{0.873514in}}%
\pgfpathquadraticcurveto{\pgfqpoint{9.603772in}{0.871262in}}{\pgfqpoint{9.603772in}{0.865408in}}%
\pgfpathlineto{\pgfqpoint{9.603772in}{0.859843in}}%
\pgfpathlineto{\pgfqpoint{9.632195in}{0.859843in}}%
\pgfpathlineto{\pgfqpoint{9.632195in}{0.864238in}}%
\pgfpathquadraticcurveto{\pgfqpoint{9.632195in}{0.874757in}}{\pgfqpoint{9.637094in}{0.879566in}}%
\pgfpathquadraticcurveto{\pgfqpoint{9.641993in}{0.884393in}}{\pgfqpoint{9.652639in}{0.884393in}}%
\pgfpathlineto{\pgfqpoint{9.662419in}{0.884393in}}%
\pgfpathlineto{\pgfqpoint{9.662419in}{0.884393in}}%
\pgfpathclose%
\pgfpathmoveto{\pgfqpoint{9.725012in}{0.830915in}}%
\pgfpathlineto{\pgfqpoint{9.725012in}{0.825854in}}%
\pgfpathlineto{\pgfqpoint{9.677387in}{0.825854in}}%
\pgfpathquadraticcurveto{\pgfqpoint{9.678072in}{0.815154in}}{\pgfqpoint{9.683836in}{0.809553in}}%
\pgfpathquadraticcurveto{\pgfqpoint{9.689600in}{0.803951in}}{\pgfqpoint{9.699903in}{0.803951in}}%
\pgfpathquadraticcurveto{\pgfqpoint{9.705865in}{0.803951in}}{\pgfqpoint{9.711466in}{0.805410in}}%
\pgfpathquadraticcurveto{\pgfqpoint{9.717068in}{0.806869in}}{\pgfqpoint{9.722598in}{0.809805in}}%
\pgfpathlineto{\pgfqpoint{9.722598in}{0.800006in}}%
\pgfpathquadraticcurveto{\pgfqpoint{9.717014in}{0.797647in}}{\pgfqpoint{9.711160in}{0.796404in}}%
\pgfpathquadraticcurveto{\pgfqpoint{9.705306in}{0.795161in}}{\pgfqpoint{9.699290in}{0.795161in}}%
\pgfpathquadraticcurveto{\pgfqpoint{9.684196in}{0.795161in}}{\pgfqpoint{9.675388in}{0.803933in}}%
\pgfpathquadraticcurveto{\pgfqpoint{9.666580in}{0.812723in}}{\pgfqpoint{9.666580in}{0.827709in}}%
\pgfpathquadraticcurveto{\pgfqpoint{9.666580in}{0.843181in}}{\pgfqpoint{9.674938in}{0.852259in}}%
\pgfpathquadraticcurveto{\pgfqpoint{9.683295in}{0.861356in}}{\pgfqpoint{9.697489in}{0.861356in}}%
\pgfpathquadraticcurveto{\pgfqpoint{9.710206in}{0.861356in}}{\pgfqpoint{9.717609in}{0.853160in}}%
\pgfpathquadraticcurveto{\pgfqpoint{9.725012in}{0.844983in}}{\pgfqpoint{9.725012in}{0.830915in}}%
\pgfpathlineto{\pgfqpoint{9.725012in}{0.830915in}}%
\pgfpathclose%
\pgfpathmoveto{\pgfqpoint{9.714655in}{0.833959in}}%
\pgfpathquadraticcurveto{\pgfqpoint{9.714546in}{0.842443in}}{\pgfqpoint{9.709899in}{0.847504in}}%
\pgfpathquadraticcurveto{\pgfqpoint{9.705252in}{0.852584in}}{\pgfqpoint{9.697597in}{0.852584in}}%
\pgfpathquadraticcurveto{\pgfqpoint{9.688933in}{0.852584in}}{\pgfqpoint{9.683728in}{0.847684in}}%
\pgfpathquadraticcurveto{\pgfqpoint{9.678522in}{0.842785in}}{\pgfqpoint{9.677730in}{0.833887in}}%
\pgfpathlineto{\pgfqpoint{9.714655in}{0.833959in}}%
\pgfpathlineto{\pgfqpoint{9.714655in}{0.833959in}}%
\pgfpathclose%
\pgfpathmoveto{\pgfqpoint{9.778544in}{0.850170in}}%
\pgfpathquadraticcurveto{\pgfqpoint{9.776796in}{0.851179in}}{\pgfqpoint{9.774743in}{0.851647in}}%
\pgfpathquadraticcurveto{\pgfqpoint{9.772690in}{0.852133in}}{\pgfqpoint{9.770222in}{0.852133in}}%
\pgfpathquadraticcurveto{\pgfqpoint{9.761432in}{0.852133in}}{\pgfqpoint{9.756731in}{0.846423in}}%
\pgfpathquadraticcurveto{\pgfqpoint{9.752030in}{0.840714in}}{\pgfqpoint{9.752030in}{0.830014in}}%
\pgfpathlineto{\pgfqpoint{9.752030in}{0.796800in}}%
\pgfpathlineto{\pgfqpoint{9.741619in}{0.796800in}}%
\pgfpathlineto{\pgfqpoint{9.741619in}{0.859843in}}%
\pgfpathlineto{\pgfqpoint{9.752030in}{0.859843in}}%
\pgfpathlineto{\pgfqpoint{9.752030in}{0.850044in}}%
\pgfpathquadraticcurveto{\pgfqpoint{9.755308in}{0.855790in}}{\pgfqpoint{9.760531in}{0.858564in}}%
\pgfpathquadraticcurveto{\pgfqpoint{9.765773in}{0.861356in}}{\pgfqpoint{9.773266in}{0.861356in}}%
\pgfpathquadraticcurveto{\pgfqpoint{9.774329in}{0.861356in}}{\pgfqpoint{9.775626in}{0.861211in}}%
\pgfpathquadraticcurveto{\pgfqpoint{9.776923in}{0.861085in}}{\pgfqpoint{9.778490in}{0.860797in}}%
\pgfpathlineto{\pgfqpoint{9.778544in}{0.850170in}}%
\pgfpathlineto{\pgfqpoint{9.778544in}{0.850170in}}%
\pgfpathclose%
\pgfpathmoveto{\pgfqpoint{9.840794in}{0.830915in}}%
\pgfpathlineto{\pgfqpoint{9.840794in}{0.825854in}}%
\pgfpathlineto{\pgfqpoint{9.793169in}{0.825854in}}%
\pgfpathquadraticcurveto{\pgfqpoint{9.793854in}{0.815154in}}{\pgfqpoint{9.799618in}{0.809553in}}%
\pgfpathquadraticcurveto{\pgfqpoint{9.805382in}{0.803951in}}{\pgfqpoint{9.815685in}{0.803951in}}%
\pgfpathquadraticcurveto{\pgfqpoint{9.821647in}{0.803951in}}{\pgfqpoint{9.827248in}{0.805410in}}%
\pgfpathquadraticcurveto{\pgfqpoint{9.832850in}{0.806869in}}{\pgfqpoint{9.838380in}{0.809805in}}%
\pgfpathlineto{\pgfqpoint{9.838380in}{0.800006in}}%
\pgfpathquadraticcurveto{\pgfqpoint{9.832796in}{0.797647in}}{\pgfqpoint{9.826942in}{0.796404in}}%
\pgfpathquadraticcurveto{\pgfqpoint{9.821088in}{0.795161in}}{\pgfqpoint{9.815072in}{0.795161in}}%
\pgfpathquadraticcurveto{\pgfqpoint{9.799978in}{0.795161in}}{\pgfqpoint{9.791170in}{0.803933in}}%
\pgfpathquadraticcurveto{\pgfqpoint{9.782362in}{0.812723in}}{\pgfqpoint{9.782362in}{0.827709in}}%
\pgfpathquadraticcurveto{\pgfqpoint{9.782362in}{0.843181in}}{\pgfqpoint{9.790720in}{0.852259in}}%
\pgfpathquadraticcurveto{\pgfqpoint{9.799077in}{0.861356in}}{\pgfqpoint{9.813271in}{0.861356in}}%
\pgfpathquadraticcurveto{\pgfqpoint{9.825988in}{0.861356in}}{\pgfqpoint{9.833391in}{0.853160in}}%
\pgfpathquadraticcurveto{\pgfqpoint{9.840794in}{0.844983in}}{\pgfqpoint{9.840794in}{0.830915in}}%
\pgfpathlineto{\pgfqpoint{9.840794in}{0.830915in}}%
\pgfpathclose%
\pgfpathmoveto{\pgfqpoint{9.830437in}{0.833959in}}%
\pgfpathquadraticcurveto{\pgfqpoint{9.830329in}{0.842443in}}{\pgfqpoint{9.825681in}{0.847504in}}%
\pgfpathquadraticcurveto{\pgfqpoint{9.821034in}{0.852584in}}{\pgfqpoint{9.813379in}{0.852584in}}%
\pgfpathquadraticcurveto{\pgfqpoint{9.804715in}{0.852584in}}{\pgfqpoint{9.799510in}{0.847684in}}%
\pgfpathquadraticcurveto{\pgfqpoint{9.794304in}{0.842785in}}{\pgfqpoint{9.793512in}{0.833887in}}%
\pgfpathlineto{\pgfqpoint{9.830437in}{0.833959in}}%
\pgfpathlineto{\pgfqpoint{9.830437in}{0.833959in}}%
\pgfpathclose%
\pgfpathmoveto{\pgfqpoint{9.910212in}{0.834860in}}%
\pgfpathlineto{\pgfqpoint{9.910212in}{0.796800in}}%
\pgfpathlineto{\pgfqpoint{9.899855in}{0.796800in}}%
\pgfpathlineto{\pgfqpoint{9.899855in}{0.834517in}}%
\pgfpathquadraticcurveto{\pgfqpoint{9.899855in}{0.843469in}}{\pgfqpoint{9.896361in}{0.847900in}}%
\pgfpathquadraticcurveto{\pgfqpoint{9.892867in}{0.852349in}}{\pgfqpoint{9.885896in}{0.852349in}}%
\pgfpathquadraticcurveto{\pgfqpoint{9.877502in}{0.852349in}}{\pgfqpoint{9.872657in}{0.847000in}}%
\pgfpathquadraticcurveto{\pgfqpoint{9.867812in}{0.841668in}}{\pgfqpoint{9.867812in}{0.832428in}}%
\pgfpathlineto{\pgfqpoint{9.867812in}{0.796800in}}%
\pgfpathlineto{\pgfqpoint{9.857401in}{0.796800in}}%
\pgfpathlineto{\pgfqpoint{9.857401in}{0.859843in}}%
\pgfpathlineto{\pgfqpoint{9.867812in}{0.859843in}}%
\pgfpathlineto{\pgfqpoint{9.867812in}{0.850044in}}%
\pgfpathquadraticcurveto{\pgfqpoint{9.871540in}{0.855736in}}{\pgfqpoint{9.876566in}{0.858546in}}%
\pgfpathquadraticcurveto{\pgfqpoint{9.881609in}{0.861356in}}{\pgfqpoint{9.888202in}{0.861356in}}%
\pgfpathquadraticcurveto{\pgfqpoint{9.899063in}{0.861356in}}{\pgfqpoint{9.904629in}{0.854637in}}%
\pgfpathquadraticcurveto{\pgfqpoint{9.910212in}{0.847918in}}{\pgfqpoint{9.910212in}{0.834860in}}%
\pgfpathlineto{\pgfqpoint{9.910212in}{0.834860in}}%
\pgfpathclose%
\pgfpathmoveto{\pgfqpoint{9.976227in}{0.857429in}}%
\pgfpathlineto{\pgfqpoint{9.976227in}{0.847738in}}%
\pgfpathquadraticcurveto{\pgfqpoint{9.971832in}{0.850170in}}{\pgfqpoint{9.967419in}{0.851377in}}%
\pgfpathquadraticcurveto{\pgfqpoint{9.963006in}{0.852584in}}{\pgfqpoint{9.958503in}{0.852584in}}%
\pgfpathquadraticcurveto{\pgfqpoint{9.948416in}{0.852584in}}{\pgfqpoint{9.942832in}{0.846189in}}%
\pgfpathquadraticcurveto{\pgfqpoint{9.937267in}{0.839813in}}{\pgfqpoint{9.937267in}{0.828267in}}%
\pgfpathquadraticcurveto{\pgfqpoint{9.937267in}{0.816721in}}{\pgfqpoint{9.942832in}{0.810327in}}%
\pgfpathquadraticcurveto{\pgfqpoint{9.948416in}{0.803951in}}{\pgfqpoint{9.958503in}{0.803951in}}%
\pgfpathquadraticcurveto{\pgfqpoint{9.963006in}{0.803951in}}{\pgfqpoint{9.967419in}{0.805158in}}%
\pgfpathquadraticcurveto{\pgfqpoint{9.971832in}{0.806364in}}{\pgfqpoint{9.976227in}{0.808796in}}%
\pgfpathlineto{\pgfqpoint{9.976227in}{0.799214in}}%
\pgfpathquadraticcurveto{\pgfqpoint{9.971886in}{0.797196in}}{\pgfqpoint{9.967239in}{0.796188in}}%
\pgfpathquadraticcurveto{\pgfqpoint{9.962610in}{0.795161in}}{\pgfqpoint{9.957368in}{0.795161in}}%
\pgfpathquadraticcurveto{\pgfqpoint{9.943121in}{0.795161in}}{\pgfqpoint{9.934727in}{0.804113in}}%
\pgfpathquadraticcurveto{\pgfqpoint{9.926351in}{0.813065in}}{\pgfqpoint{9.926351in}{0.828267in}}%
\pgfpathquadraticcurveto{\pgfqpoint{9.926351in}{0.843686in}}{\pgfqpoint{9.934817in}{0.852512in}}%
\pgfpathquadraticcurveto{\pgfqpoint{9.943301in}{0.861356in}}{\pgfqpoint{9.958053in}{0.861356in}}%
\pgfpathquadraticcurveto{\pgfqpoint{9.962826in}{0.861356in}}{\pgfqpoint{9.967383in}{0.860365in}}%
\pgfpathquadraticcurveto{\pgfqpoint{9.971940in}{0.859392in}}{\pgfqpoint{9.976227in}{0.857429in}}%
\pgfpathlineto{\pgfqpoint{9.976227in}{0.857429in}}%
\pgfpathclose%
\pgfpathmoveto{\pgfqpoint{10.048168in}{0.830915in}}%
\pgfpathlineto{\pgfqpoint{10.048168in}{0.825854in}}%
\pgfpathlineto{\pgfqpoint{10.000543in}{0.825854in}}%
\pgfpathquadraticcurveto{\pgfqpoint{10.001228in}{0.815154in}}{\pgfqpoint{10.006992in}{0.809553in}}%
\pgfpathquadraticcurveto{\pgfqpoint{10.012756in}{0.803951in}}{\pgfqpoint{10.023059in}{0.803951in}}%
\pgfpathquadraticcurveto{\pgfqpoint{10.029021in}{0.803951in}}{\pgfqpoint{10.034622in}{0.805410in}}%
\pgfpathquadraticcurveto{\pgfqpoint{10.040224in}{0.806869in}}{\pgfqpoint{10.045754in}{0.809805in}}%
\pgfpathlineto{\pgfqpoint{10.045754in}{0.800006in}}%
\pgfpathquadraticcurveto{\pgfqpoint{10.040170in}{0.797647in}}{\pgfqpoint{10.034316in}{0.796404in}}%
\pgfpathquadraticcurveto{\pgfqpoint{10.028462in}{0.795161in}}{\pgfqpoint{10.022446in}{0.795161in}}%
\pgfpathquadraticcurveto{\pgfqpoint{10.007352in}{0.795161in}}{\pgfqpoint{9.998544in}{0.803933in}}%
\pgfpathquadraticcurveto{\pgfqpoint{9.989736in}{0.812723in}}{\pgfqpoint{9.989736in}{0.827709in}}%
\pgfpathquadraticcurveto{\pgfqpoint{9.989736in}{0.843181in}}{\pgfqpoint{9.998094in}{0.852259in}}%
\pgfpathquadraticcurveto{\pgfqpoint{10.006451in}{0.861356in}}{\pgfqpoint{10.020645in}{0.861356in}}%
\pgfpathquadraticcurveto{\pgfqpoint{10.033362in}{0.861356in}}{\pgfqpoint{10.040765in}{0.853160in}}%
\pgfpathquadraticcurveto{\pgfqpoint{10.048168in}{0.844983in}}{\pgfqpoint{10.048168in}{0.830915in}}%
\pgfpathlineto{\pgfqpoint{10.048168in}{0.830915in}}%
\pgfpathclose%
\pgfpathmoveto{\pgfqpoint{10.037811in}{0.833959in}}%
\pgfpathquadraticcurveto{\pgfqpoint{10.037702in}{0.842443in}}{\pgfqpoint{10.033055in}{0.847504in}}%
\pgfpathquadraticcurveto{\pgfqpoint{10.028408in}{0.852584in}}{\pgfqpoint{10.020753in}{0.852584in}}%
\pgfpathquadraticcurveto{\pgfqpoint{10.012089in}{0.852584in}}{\pgfqpoint{10.006884in}{0.847684in}}%
\pgfpathquadraticcurveto{\pgfqpoint{10.001678in}{0.842785in}}{\pgfqpoint{10.000886in}{0.833887in}}%
\pgfpathlineto{\pgfqpoint{10.037811in}{0.833959in}}%
\pgfpathlineto{\pgfqpoint{10.037811in}{0.833959in}}%
\pgfpathclose%
\pgfpathmoveto{\pgfqpoint{10.067819in}{0.811102in}}%
\pgfpathlineto{\pgfqpoint{10.079689in}{0.811102in}}%
\pgfpathlineto{\pgfqpoint{10.079689in}{0.801411in}}%
\pgfpathlineto{\pgfqpoint{10.070467in}{0.783399in}}%
\pgfpathlineto{\pgfqpoint{10.063208in}{0.783399in}}%
\pgfpathlineto{\pgfqpoint{10.067819in}{0.801411in}}%
\pgfpathlineto{\pgfqpoint{10.067819in}{0.811102in}}%
\pgfpathlineto{\pgfqpoint{10.067819in}{0.811102in}}%
\pgfpathclose%
\pgfusepath{fill}%
\end{pgfscope}%
\begin{pgfscope}%
\pgfsetbuttcap%
\pgfsetroundjoin%
\definecolor{currentfill}{rgb}{0.309804,0.372549,0.419608}%
\pgfsetfillcolor{currentfill}%
\pgfsetlinewidth{0.000000pt}%
\definecolor{currentstroke}{rgb}{0.309804,0.372549,0.419608}%
\pgfsetstrokecolor{currentstroke}%
\pgfsetdash{}{0pt}%
\pgfpathmoveto{\pgfqpoint{10.202088in}{0.852584in}}%
\pgfpathquadraticcurveto{\pgfqpoint{10.193766in}{0.852584in}}{\pgfqpoint{10.188921in}{0.846081in}}%
\pgfpathquadraticcurveto{\pgfqpoint{10.184076in}{0.839579in}}{\pgfqpoint{10.184076in}{0.828267in}}%
\pgfpathquadraticcurveto{\pgfqpoint{10.184076in}{0.816956in}}{\pgfqpoint{10.188885in}{0.810453in}}%
\pgfpathquadraticcurveto{\pgfqpoint{10.193712in}{0.803951in}}{\pgfqpoint{10.202088in}{0.803951in}}%
\pgfpathquadraticcurveto{\pgfqpoint{10.210373in}{0.803951in}}{\pgfqpoint{10.215201in}{0.810471in}}%
\pgfpathquadraticcurveto{\pgfqpoint{10.220046in}{0.817010in}}{\pgfqpoint{10.220046in}{0.828267in}}%
\pgfpathquadraticcurveto{\pgfqpoint{10.220046in}{0.839471in}}{\pgfqpoint{10.215201in}{0.846027in}}%
\pgfpathquadraticcurveto{\pgfqpoint{10.210373in}{0.852584in}}{\pgfqpoint{10.202088in}{0.852584in}}%
\pgfpathlineto{\pgfqpoint{10.202088in}{0.852584in}}%
\pgfpathclose%
\pgfpathmoveto{\pgfqpoint{10.202088in}{0.861356in}}%
\pgfpathquadraticcurveto{\pgfqpoint{10.215597in}{0.861356in}}{\pgfqpoint{10.223306in}{0.852566in}}%
\pgfpathquadraticcurveto{\pgfqpoint{10.231033in}{0.843794in}}{\pgfqpoint{10.231033in}{0.828267in}}%
\pgfpathquadraticcurveto{\pgfqpoint{10.231033in}{0.812795in}}{\pgfqpoint{10.223306in}{0.803969in}}%
\pgfpathquadraticcurveto{\pgfqpoint{10.215597in}{0.795161in}}{\pgfqpoint{10.202088in}{0.795161in}}%
\pgfpathquadraticcurveto{\pgfqpoint{10.188525in}{0.795161in}}{\pgfqpoint{10.180834in}{0.803969in}}%
\pgfpathquadraticcurveto{\pgfqpoint{10.173160in}{0.812795in}}{\pgfqpoint{10.173160in}{0.828267in}}%
\pgfpathquadraticcurveto{\pgfqpoint{10.173160in}{0.843794in}}{\pgfqpoint{10.180834in}{0.852566in}}%
\pgfpathquadraticcurveto{\pgfqpoint{10.188525in}{0.861356in}}{\pgfqpoint{10.202088in}{0.861356in}}%
\pgfpathlineto{\pgfqpoint{10.202088in}{0.861356in}}%
\pgfpathclose%
\pgfpathmoveto{\pgfqpoint{10.284728in}{0.850170in}}%
\pgfpathquadraticcurveto{\pgfqpoint{10.282980in}{0.851179in}}{\pgfqpoint{10.280927in}{0.851647in}}%
\pgfpathquadraticcurveto{\pgfqpoint{10.278874in}{0.852133in}}{\pgfqpoint{10.276406in}{0.852133in}}%
\pgfpathquadraticcurveto{\pgfqpoint{10.267616in}{0.852133in}}{\pgfqpoint{10.262915in}{0.846423in}}%
\pgfpathquadraticcurveto{\pgfqpoint{10.258214in}{0.840714in}}{\pgfqpoint{10.258214in}{0.830014in}}%
\pgfpathlineto{\pgfqpoint{10.258214in}{0.796800in}}%
\pgfpathlineto{\pgfqpoint{10.247803in}{0.796800in}}%
\pgfpathlineto{\pgfqpoint{10.247803in}{0.859843in}}%
\pgfpathlineto{\pgfqpoint{10.258214in}{0.859843in}}%
\pgfpathlineto{\pgfqpoint{10.258214in}{0.850044in}}%
\pgfpathquadraticcurveto{\pgfqpoint{10.261492in}{0.855790in}}{\pgfqpoint{10.266715in}{0.858564in}}%
\pgfpathquadraticcurveto{\pgfqpoint{10.271957in}{0.861356in}}{\pgfqpoint{10.279450in}{0.861356in}}%
\pgfpathquadraticcurveto{\pgfqpoint{10.280513in}{0.861356in}}{\pgfqpoint{10.281810in}{0.861211in}}%
\pgfpathquadraticcurveto{\pgfqpoint{10.283107in}{0.861085in}}{\pgfqpoint{10.284674in}{0.860797in}}%
\pgfpathlineto{\pgfqpoint{10.284728in}{0.850170in}}%
\pgfpathlineto{\pgfqpoint{10.284728in}{0.850170in}}%
\pgfpathclose%
\pgfusepath{fill}%
\end{pgfscope}%
\begin{pgfscope}%
\pgfsetbuttcap%
\pgfsetroundjoin%
\definecolor{currentfill}{rgb}{0.309804,0.372549,0.419608}%
\pgfsetfillcolor{currentfill}%
\pgfsetlinewidth{0.000000pt}%
\definecolor{currentstroke}{rgb}{0.309804,0.372549,0.419608}%
\pgfsetstrokecolor{currentstroke}%
\pgfsetdash{}{0pt}%
\pgfpathmoveto{\pgfqpoint{10.361383in}{0.806256in}}%
\pgfpathlineto{\pgfqpoint{10.361383in}{0.772826in}}%
\pgfpathlineto{\pgfqpoint{10.350972in}{0.772826in}}%
\pgfpathlineto{\pgfqpoint{10.350972in}{0.859843in}}%
\pgfpathlineto{\pgfqpoint{10.361383in}{0.859843in}}%
\pgfpathlineto{\pgfqpoint{10.361383in}{0.850278in}}%
\pgfpathquadraticcurveto{\pgfqpoint{10.364661in}{0.855898in}}{\pgfqpoint{10.369633in}{0.858618in}}%
\pgfpathquadraticcurveto{\pgfqpoint{10.374622in}{0.861356in}}{\pgfqpoint{10.381539in}{0.861356in}}%
\pgfpathquadraticcurveto{\pgfqpoint{10.393030in}{0.861356in}}{\pgfqpoint{10.400199in}{0.852241in}}%
\pgfpathquadraticcurveto{\pgfqpoint{10.407386in}{0.843127in}}{\pgfqpoint{10.407386in}{0.828267in}}%
\pgfpathquadraticcurveto{\pgfqpoint{10.407386in}{0.813407in}}{\pgfqpoint{10.400199in}{0.804275in}}%
\pgfpathquadraticcurveto{\pgfqpoint{10.393030in}{0.795161in}}{\pgfqpoint{10.381539in}{0.795161in}}%
\pgfpathquadraticcurveto{\pgfqpoint{10.374622in}{0.795161in}}{\pgfqpoint{10.369633in}{0.797899in}}%
\pgfpathquadraticcurveto{\pgfqpoint{10.364661in}{0.800637in}}{\pgfqpoint{10.361383in}{0.806256in}}%
\pgfpathlineto{\pgfqpoint{10.361383in}{0.806256in}}%
\pgfpathclose%
\pgfpathmoveto{\pgfqpoint{10.396633in}{0.828267in}}%
\pgfpathquadraticcurveto{\pgfqpoint{10.396633in}{0.839687in}}{\pgfqpoint{10.391932in}{0.846189in}}%
\pgfpathquadraticcurveto{\pgfqpoint{10.387230in}{0.852692in}}{\pgfqpoint{10.379017in}{0.852692in}}%
\pgfpathquadraticcurveto{\pgfqpoint{10.370785in}{0.852692in}}{\pgfqpoint{10.366084in}{0.846189in}}%
\pgfpathquadraticcurveto{\pgfqpoint{10.361383in}{0.839687in}}{\pgfqpoint{10.361383in}{0.828267in}}%
\pgfpathquadraticcurveto{\pgfqpoint{10.361383in}{0.816848in}}{\pgfqpoint{10.366084in}{0.810345in}}%
\pgfpathquadraticcurveto{\pgfqpoint{10.370785in}{0.803843in}}{\pgfqpoint{10.379017in}{0.803843in}}%
\pgfpathquadraticcurveto{\pgfqpoint{10.387230in}{0.803843in}}{\pgfqpoint{10.391932in}{0.810345in}}%
\pgfpathquadraticcurveto{\pgfqpoint{10.396633in}{0.816848in}}{\pgfqpoint{10.396633in}{0.828267in}}%
\pgfpathlineto{\pgfqpoint{10.396633in}{0.828267in}}%
\pgfpathclose%
\pgfpathmoveto{\pgfqpoint{10.461080in}{0.850170in}}%
\pgfpathquadraticcurveto{\pgfqpoint{10.459333in}{0.851179in}}{\pgfqpoint{10.457280in}{0.851647in}}%
\pgfpathquadraticcurveto{\pgfqpoint{10.455226in}{0.852133in}}{\pgfqpoint{10.452759in}{0.852133in}}%
\pgfpathquadraticcurveto{\pgfqpoint{10.443969in}{0.852133in}}{\pgfqpoint{10.439268in}{0.846423in}}%
\pgfpathquadraticcurveto{\pgfqpoint{10.434566in}{0.840714in}}{\pgfqpoint{10.434566in}{0.830014in}}%
\pgfpathlineto{\pgfqpoint{10.434566in}{0.796800in}}%
\pgfpathlineto{\pgfqpoint{10.424155in}{0.796800in}}%
\pgfpathlineto{\pgfqpoint{10.424155in}{0.859843in}}%
\pgfpathlineto{\pgfqpoint{10.434566in}{0.859843in}}%
\pgfpathlineto{\pgfqpoint{10.434566in}{0.850044in}}%
\pgfpathquadraticcurveto{\pgfqpoint{10.437845in}{0.855790in}}{\pgfqpoint{10.443068in}{0.858564in}}%
\pgfpathquadraticcurveto{\pgfqpoint{10.448310in}{0.861356in}}{\pgfqpoint{10.455803in}{0.861356in}}%
\pgfpathquadraticcurveto{\pgfqpoint{10.456865in}{0.861356in}}{\pgfqpoint{10.458162in}{0.861211in}}%
\pgfpathquadraticcurveto{\pgfqpoint{10.459459in}{0.861085in}}{\pgfqpoint{10.461026in}{0.860797in}}%
\pgfpathlineto{\pgfqpoint{10.461080in}{0.850170in}}%
\pgfpathlineto{\pgfqpoint{10.461080in}{0.850170in}}%
\pgfpathclose%
\pgfpathmoveto{\pgfqpoint{10.471942in}{0.859843in}}%
\pgfpathlineto{\pgfqpoint{10.482299in}{0.859843in}}%
\pgfpathlineto{\pgfqpoint{10.482299in}{0.796800in}}%
\pgfpathlineto{\pgfqpoint{10.471942in}{0.796800in}}%
\pgfpathlineto{\pgfqpoint{10.471942in}{0.859843in}}%
\pgfpathlineto{\pgfqpoint{10.471942in}{0.859843in}}%
\pgfpathclose%
\pgfpathmoveto{\pgfqpoint{10.471942in}{0.884393in}}%
\pgfpathlineto{\pgfqpoint{10.482299in}{0.884393in}}%
\pgfpathlineto{\pgfqpoint{10.482299in}{0.871262in}}%
\pgfpathlineto{\pgfqpoint{10.471942in}{0.871262in}}%
\pgfpathlineto{\pgfqpoint{10.471942in}{0.884393in}}%
\pgfpathlineto{\pgfqpoint{10.471942in}{0.884393in}}%
\pgfpathclose%
\pgfpathmoveto{\pgfqpoint{10.528392in}{0.852584in}}%
\pgfpathquadraticcurveto{\pgfqpoint{10.520070in}{0.852584in}}{\pgfqpoint{10.515225in}{0.846081in}}%
\pgfpathquadraticcurveto{\pgfqpoint{10.510380in}{0.839579in}}{\pgfqpoint{10.510380in}{0.828267in}}%
\pgfpathquadraticcurveto{\pgfqpoint{10.510380in}{0.816956in}}{\pgfqpoint{10.515189in}{0.810453in}}%
\pgfpathquadraticcurveto{\pgfqpoint{10.520016in}{0.803951in}}{\pgfqpoint{10.528392in}{0.803951in}}%
\pgfpathquadraticcurveto{\pgfqpoint{10.536677in}{0.803951in}}{\pgfqpoint{10.541505in}{0.810471in}}%
\pgfpathquadraticcurveto{\pgfqpoint{10.546350in}{0.817010in}}{\pgfqpoint{10.546350in}{0.828267in}}%
\pgfpathquadraticcurveto{\pgfqpoint{10.546350in}{0.839471in}}{\pgfqpoint{10.541505in}{0.846027in}}%
\pgfpathquadraticcurveto{\pgfqpoint{10.536677in}{0.852584in}}{\pgfqpoint{10.528392in}{0.852584in}}%
\pgfpathlineto{\pgfqpoint{10.528392in}{0.852584in}}%
\pgfpathclose%
\pgfpathmoveto{\pgfqpoint{10.528392in}{0.861356in}}%
\pgfpathquadraticcurveto{\pgfqpoint{10.541901in}{0.861356in}}{\pgfqpoint{10.549610in}{0.852566in}}%
\pgfpathquadraticcurveto{\pgfqpoint{10.557337in}{0.843794in}}{\pgfqpoint{10.557337in}{0.828267in}}%
\pgfpathquadraticcurveto{\pgfqpoint{10.557337in}{0.812795in}}{\pgfqpoint{10.549610in}{0.803969in}}%
\pgfpathquadraticcurveto{\pgfqpoint{10.541901in}{0.795161in}}{\pgfqpoint{10.528392in}{0.795161in}}%
\pgfpathquadraticcurveto{\pgfqpoint{10.514829in}{0.795161in}}{\pgfqpoint{10.507137in}{0.803969in}}%
\pgfpathquadraticcurveto{\pgfqpoint{10.499464in}{0.812795in}}{\pgfqpoint{10.499464in}{0.828267in}}%
\pgfpathquadraticcurveto{\pgfqpoint{10.499464in}{0.843794in}}{\pgfqpoint{10.507137in}{0.852566in}}%
\pgfpathquadraticcurveto{\pgfqpoint{10.514829in}{0.861356in}}{\pgfqpoint{10.528392in}{0.861356in}}%
\pgfpathlineto{\pgfqpoint{10.528392in}{0.861356in}}%
\pgfpathclose%
\pgfpathmoveto{\pgfqpoint{10.611031in}{0.850170in}}%
\pgfpathquadraticcurveto{\pgfqpoint{10.609284in}{0.851179in}}{\pgfqpoint{10.607231in}{0.851647in}}%
\pgfpathquadraticcurveto{\pgfqpoint{10.605178in}{0.852133in}}{\pgfqpoint{10.602710in}{0.852133in}}%
\pgfpathquadraticcurveto{\pgfqpoint{10.593920in}{0.852133in}}{\pgfqpoint{10.589219in}{0.846423in}}%
\pgfpathquadraticcurveto{\pgfqpoint{10.584518in}{0.840714in}}{\pgfqpoint{10.584518in}{0.830014in}}%
\pgfpathlineto{\pgfqpoint{10.584518in}{0.796800in}}%
\pgfpathlineto{\pgfqpoint{10.574107in}{0.796800in}}%
\pgfpathlineto{\pgfqpoint{10.574107in}{0.859843in}}%
\pgfpathlineto{\pgfqpoint{10.584518in}{0.859843in}}%
\pgfpathlineto{\pgfqpoint{10.584518in}{0.850044in}}%
\pgfpathquadraticcurveto{\pgfqpoint{10.587796in}{0.855790in}}{\pgfqpoint{10.593019in}{0.858564in}}%
\pgfpathquadraticcurveto{\pgfqpoint{10.598261in}{0.861356in}}{\pgfqpoint{10.605754in}{0.861356in}}%
\pgfpathquadraticcurveto{\pgfqpoint{10.606817in}{0.861356in}}{\pgfqpoint{10.608113in}{0.861211in}}%
\pgfpathquadraticcurveto{\pgfqpoint{10.609410in}{0.861085in}}{\pgfqpoint{10.610977in}{0.860797in}}%
\pgfpathlineto{\pgfqpoint{10.611031in}{0.850170in}}%
\pgfpathlineto{\pgfqpoint{10.611031in}{0.850170in}}%
\pgfpathclose%
\pgfusepath{fill}%
\end{pgfscope}%
\begin{pgfscope}%
\pgfsetbuttcap%
\pgfsetroundjoin%
\definecolor{currentfill}{rgb}{0.309804,0.372549,0.419608}%
\pgfsetfillcolor{currentfill}%
\pgfsetlinewidth{0.000000pt}%
\definecolor{currentstroke}{rgb}{0.309804,0.372549,0.419608}%
\pgfsetstrokecolor{currentstroke}%
\pgfsetdash{}{0pt}%
\pgfpathmoveto{\pgfqpoint{10.685073in}{0.859843in}}%
\pgfpathlineto{\pgfqpoint{10.695430in}{0.859843in}}%
\pgfpathlineto{\pgfqpoint{10.695430in}{0.796800in}}%
\pgfpathlineto{\pgfqpoint{10.685073in}{0.796800in}}%
\pgfpathlineto{\pgfqpoint{10.685073in}{0.859843in}}%
\pgfpathlineto{\pgfqpoint{10.685073in}{0.859843in}}%
\pgfpathclose%
\pgfpathmoveto{\pgfqpoint{10.685073in}{0.884393in}}%
\pgfpathlineto{\pgfqpoint{10.695430in}{0.884393in}}%
\pgfpathlineto{\pgfqpoint{10.695430in}{0.871262in}}%
\pgfpathlineto{\pgfqpoint{10.685073in}{0.871262in}}%
\pgfpathlineto{\pgfqpoint{10.685073in}{0.884393in}}%
\pgfpathlineto{\pgfqpoint{10.685073in}{0.884393in}}%
\pgfpathclose%
\pgfpathmoveto{\pgfqpoint{10.769514in}{0.834860in}}%
\pgfpathlineto{\pgfqpoint{10.769514in}{0.796800in}}%
\pgfpathlineto{\pgfqpoint{10.759157in}{0.796800in}}%
\pgfpathlineto{\pgfqpoint{10.759157in}{0.834517in}}%
\pgfpathquadraticcurveto{\pgfqpoint{10.759157in}{0.843469in}}{\pgfqpoint{10.755663in}{0.847900in}}%
\pgfpathquadraticcurveto{\pgfqpoint{10.752168in}{0.852349in}}{\pgfqpoint{10.745198in}{0.852349in}}%
\pgfpathquadraticcurveto{\pgfqpoint{10.736804in}{0.852349in}}{\pgfqpoint{10.731959in}{0.847000in}}%
\pgfpathquadraticcurveto{\pgfqpoint{10.727113in}{0.841668in}}{\pgfqpoint{10.727113in}{0.832428in}}%
\pgfpathlineto{\pgfqpoint{10.727113in}{0.796800in}}%
\pgfpathlineto{\pgfqpoint{10.716702in}{0.796800in}}%
\pgfpathlineto{\pgfqpoint{10.716702in}{0.859843in}}%
\pgfpathlineto{\pgfqpoint{10.727113in}{0.859843in}}%
\pgfpathlineto{\pgfqpoint{10.727113in}{0.850044in}}%
\pgfpathquadraticcurveto{\pgfqpoint{10.730842in}{0.855736in}}{\pgfqpoint{10.735867in}{0.858546in}}%
\pgfpathquadraticcurveto{\pgfqpoint{10.740911in}{0.861356in}}{\pgfqpoint{10.747503in}{0.861356in}}%
\pgfpathquadraticcurveto{\pgfqpoint{10.758365in}{0.861356in}}{\pgfqpoint{10.763930in}{0.854637in}}%
\pgfpathquadraticcurveto{\pgfqpoint{10.769514in}{0.847918in}}{\pgfqpoint{10.769514in}{0.834860in}}%
\pgfpathlineto{\pgfqpoint{10.769514in}{0.834860in}}%
\pgfpathclose%
\pgfpathmoveto{\pgfqpoint{10.822074in}{0.884393in}}%
\pgfpathlineto{\pgfqpoint{10.822074in}{0.875765in}}%
\pgfpathlineto{\pgfqpoint{10.812167in}{0.875765in}}%
\pgfpathquadraticcurveto{\pgfqpoint{10.806601in}{0.875765in}}{\pgfqpoint{10.804422in}{0.873514in}}%
\pgfpathquadraticcurveto{\pgfqpoint{10.802260in}{0.871262in}}{\pgfqpoint{10.802260in}{0.865408in}}%
\pgfpathlineto{\pgfqpoint{10.802260in}{0.859843in}}%
\pgfpathlineto{\pgfqpoint{10.819318in}{0.859843in}}%
\pgfpathlineto{\pgfqpoint{10.819318in}{0.851791in}}%
\pgfpathlineto{\pgfqpoint{10.802260in}{0.851791in}}%
\pgfpathlineto{\pgfqpoint{10.802260in}{0.796800in}}%
\pgfpathlineto{\pgfqpoint{10.791849in}{0.796800in}}%
\pgfpathlineto{\pgfqpoint{10.791849in}{0.851791in}}%
\pgfpathlineto{\pgfqpoint{10.781942in}{0.851791in}}%
\pgfpathlineto{\pgfqpoint{10.781942in}{0.859843in}}%
\pgfpathlineto{\pgfqpoint{10.791849in}{0.859843in}}%
\pgfpathlineto{\pgfqpoint{10.791849in}{0.864238in}}%
\pgfpathquadraticcurveto{\pgfqpoint{10.791849in}{0.874757in}}{\pgfqpoint{10.796748in}{0.879566in}}%
\pgfpathquadraticcurveto{\pgfqpoint{10.801648in}{0.884393in}}{\pgfqpoint{10.812275in}{0.884393in}}%
\pgfpathlineto{\pgfqpoint{10.822074in}{0.884393in}}%
\pgfpathlineto{\pgfqpoint{10.822074in}{0.884393in}}%
\pgfpathclose%
\pgfpathmoveto{\pgfqpoint{10.884666in}{0.830915in}}%
\pgfpathlineto{\pgfqpoint{10.884666in}{0.825854in}}%
\pgfpathlineto{\pgfqpoint{10.837042in}{0.825854in}}%
\pgfpathquadraticcurveto{\pgfqpoint{10.837726in}{0.815154in}}{\pgfqpoint{10.843490in}{0.809553in}}%
\pgfpathquadraticcurveto{\pgfqpoint{10.849254in}{0.803951in}}{\pgfqpoint{10.859557in}{0.803951in}}%
\pgfpathquadraticcurveto{\pgfqpoint{10.865519in}{0.803951in}}{\pgfqpoint{10.871121in}{0.805410in}}%
\pgfpathquadraticcurveto{\pgfqpoint{10.876722in}{0.806869in}}{\pgfqpoint{10.882252in}{0.809805in}}%
\pgfpathlineto{\pgfqpoint{10.882252in}{0.800006in}}%
\pgfpathquadraticcurveto{\pgfqpoint{10.876668in}{0.797647in}}{\pgfqpoint{10.870814in}{0.796404in}}%
\pgfpathquadraticcurveto{\pgfqpoint{10.864961in}{0.795161in}}{\pgfqpoint{10.858944in}{0.795161in}}%
\pgfpathquadraticcurveto{\pgfqpoint{10.843850in}{0.795161in}}{\pgfqpoint{10.835042in}{0.803933in}}%
\pgfpathquadraticcurveto{\pgfqpoint{10.826234in}{0.812723in}}{\pgfqpoint{10.826234in}{0.827709in}}%
\pgfpathquadraticcurveto{\pgfqpoint{10.826234in}{0.843181in}}{\pgfqpoint{10.834592in}{0.852259in}}%
\pgfpathquadraticcurveto{\pgfqpoint{10.842950in}{0.861356in}}{\pgfqpoint{10.857143in}{0.861356in}}%
\pgfpathquadraticcurveto{\pgfqpoint{10.869860in}{0.861356in}}{\pgfqpoint{10.877263in}{0.853160in}}%
\pgfpathquadraticcurveto{\pgfqpoint{10.884666in}{0.844983in}}{\pgfqpoint{10.884666in}{0.830915in}}%
\pgfpathlineto{\pgfqpoint{10.884666in}{0.830915in}}%
\pgfpathclose%
\pgfpathmoveto{\pgfqpoint{10.874309in}{0.833959in}}%
\pgfpathquadraticcurveto{\pgfqpoint{10.874201in}{0.842443in}}{\pgfqpoint{10.869554in}{0.847504in}}%
\pgfpathquadraticcurveto{\pgfqpoint{10.864906in}{0.852584in}}{\pgfqpoint{10.857251in}{0.852584in}}%
\pgfpathquadraticcurveto{\pgfqpoint{10.848587in}{0.852584in}}{\pgfqpoint{10.843382in}{0.847684in}}%
\pgfpathquadraticcurveto{\pgfqpoint{10.838176in}{0.842785in}}{\pgfqpoint{10.837384in}{0.833887in}}%
\pgfpathlineto{\pgfqpoint{10.874309in}{0.833959in}}%
\pgfpathlineto{\pgfqpoint{10.874309in}{0.833959in}}%
\pgfpathclose%
\pgfpathmoveto{\pgfqpoint{10.938198in}{0.850170in}}%
\pgfpathquadraticcurveto{\pgfqpoint{10.936451in}{0.851179in}}{\pgfqpoint{10.934397in}{0.851647in}}%
\pgfpathquadraticcurveto{\pgfqpoint{10.932344in}{0.852133in}}{\pgfqpoint{10.929876in}{0.852133in}}%
\pgfpathquadraticcurveto{\pgfqpoint{10.921086in}{0.852133in}}{\pgfqpoint{10.916385in}{0.846423in}}%
\pgfpathquadraticcurveto{\pgfqpoint{10.911684in}{0.840714in}}{\pgfqpoint{10.911684in}{0.830014in}}%
\pgfpathlineto{\pgfqpoint{10.911684in}{0.796800in}}%
\pgfpathlineto{\pgfqpoint{10.901273in}{0.796800in}}%
\pgfpathlineto{\pgfqpoint{10.901273in}{0.859843in}}%
\pgfpathlineto{\pgfqpoint{10.911684in}{0.859843in}}%
\pgfpathlineto{\pgfqpoint{10.911684in}{0.850044in}}%
\pgfpathquadraticcurveto{\pgfqpoint{10.914962in}{0.855790in}}{\pgfqpoint{10.920186in}{0.858564in}}%
\pgfpathquadraticcurveto{\pgfqpoint{10.925427in}{0.861356in}}{\pgfqpoint{10.932920in}{0.861356in}}%
\pgfpathquadraticcurveto{\pgfqpoint{10.933983in}{0.861356in}}{\pgfqpoint{10.935280in}{0.861211in}}%
\pgfpathquadraticcurveto{\pgfqpoint{10.936577in}{0.861085in}}{\pgfqpoint{10.938144in}{0.860797in}}%
\pgfpathlineto{\pgfqpoint{10.938198in}{0.850170in}}%
\pgfpathlineto{\pgfqpoint{10.938198in}{0.850170in}}%
\pgfpathclose%
\pgfpathmoveto{\pgfqpoint{11.000448in}{0.830915in}}%
\pgfpathlineto{\pgfqpoint{11.000448in}{0.825854in}}%
\pgfpathlineto{\pgfqpoint{10.952824in}{0.825854in}}%
\pgfpathquadraticcurveto{\pgfqpoint{10.953508in}{0.815154in}}{\pgfqpoint{10.959272in}{0.809553in}}%
\pgfpathquadraticcurveto{\pgfqpoint{10.965036in}{0.803951in}}{\pgfqpoint{10.975339in}{0.803951in}}%
\pgfpathquadraticcurveto{\pgfqpoint{10.981301in}{0.803951in}}{\pgfqpoint{10.986903in}{0.805410in}}%
\pgfpathquadraticcurveto{\pgfqpoint{10.992505in}{0.806869in}}{\pgfqpoint{10.998034in}{0.809805in}}%
\pgfpathlineto{\pgfqpoint{10.998034in}{0.800006in}}%
\pgfpathquadraticcurveto{\pgfqpoint{10.992451in}{0.797647in}}{\pgfqpoint{10.986597in}{0.796404in}}%
\pgfpathquadraticcurveto{\pgfqpoint{10.980743in}{0.795161in}}{\pgfqpoint{10.974727in}{0.795161in}}%
\pgfpathquadraticcurveto{\pgfqpoint{10.959632in}{0.795161in}}{\pgfqpoint{10.950824in}{0.803933in}}%
\pgfpathquadraticcurveto{\pgfqpoint{10.942016in}{0.812723in}}{\pgfqpoint{10.942016in}{0.827709in}}%
\pgfpathquadraticcurveto{\pgfqpoint{10.942016in}{0.843181in}}{\pgfqpoint{10.950374in}{0.852259in}}%
\pgfpathquadraticcurveto{\pgfqpoint{10.958732in}{0.861356in}}{\pgfqpoint{10.972925in}{0.861356in}}%
\pgfpathquadraticcurveto{\pgfqpoint{10.985642in}{0.861356in}}{\pgfqpoint{10.993045in}{0.853160in}}%
\pgfpathquadraticcurveto{\pgfqpoint{11.000448in}{0.844983in}}{\pgfqpoint{11.000448in}{0.830915in}}%
\pgfpathlineto{\pgfqpoint{11.000448in}{0.830915in}}%
\pgfpathclose%
\pgfpathmoveto{\pgfqpoint{10.990091in}{0.833959in}}%
\pgfpathquadraticcurveto{\pgfqpoint{10.989983in}{0.842443in}}{\pgfqpoint{10.985336in}{0.847504in}}%
\pgfpathquadraticcurveto{\pgfqpoint{10.980689in}{0.852584in}}{\pgfqpoint{10.973033in}{0.852584in}}%
\pgfpathquadraticcurveto{\pgfqpoint{10.964370in}{0.852584in}}{\pgfqpoint{10.959164in}{0.847684in}}%
\pgfpathquadraticcurveto{\pgfqpoint{10.953959in}{0.842785in}}{\pgfqpoint{10.953166in}{0.833887in}}%
\pgfpathlineto{\pgfqpoint{10.990091in}{0.833959in}}%
\pgfpathlineto{\pgfqpoint{10.990091in}{0.833959in}}%
\pgfpathclose%
\pgfpathmoveto{\pgfqpoint{11.069867in}{0.834860in}}%
\pgfpathlineto{\pgfqpoint{11.069867in}{0.796800in}}%
\pgfpathlineto{\pgfqpoint{11.059510in}{0.796800in}}%
\pgfpathlineto{\pgfqpoint{11.059510in}{0.834517in}}%
\pgfpathquadraticcurveto{\pgfqpoint{11.059510in}{0.843469in}}{\pgfqpoint{11.056015in}{0.847900in}}%
\pgfpathquadraticcurveto{\pgfqpoint{11.052521in}{0.852349in}}{\pgfqpoint{11.045550in}{0.852349in}}%
\pgfpathquadraticcurveto{\pgfqpoint{11.037157in}{0.852349in}}{\pgfqpoint{11.032311in}{0.847000in}}%
\pgfpathquadraticcurveto{\pgfqpoint{11.027466in}{0.841668in}}{\pgfqpoint{11.027466in}{0.832428in}}%
\pgfpathlineto{\pgfqpoint{11.027466in}{0.796800in}}%
\pgfpathlineto{\pgfqpoint{11.017055in}{0.796800in}}%
\pgfpathlineto{\pgfqpoint{11.017055in}{0.859843in}}%
\pgfpathlineto{\pgfqpoint{11.027466in}{0.859843in}}%
\pgfpathlineto{\pgfqpoint{11.027466in}{0.850044in}}%
\pgfpathquadraticcurveto{\pgfqpoint{11.031195in}{0.855736in}}{\pgfqpoint{11.036220in}{0.858546in}}%
\pgfpathquadraticcurveto{\pgfqpoint{11.041263in}{0.861356in}}{\pgfqpoint{11.047856in}{0.861356in}}%
\pgfpathquadraticcurveto{\pgfqpoint{11.058717in}{0.861356in}}{\pgfqpoint{11.064283in}{0.854637in}}%
\pgfpathquadraticcurveto{\pgfqpoint{11.069867in}{0.847918in}}{\pgfqpoint{11.069867in}{0.834860in}}%
\pgfpathlineto{\pgfqpoint{11.069867in}{0.834860in}}%
\pgfpathclose%
\pgfpathmoveto{\pgfqpoint{11.135881in}{0.857429in}}%
\pgfpathlineto{\pgfqpoint{11.135881in}{0.847738in}}%
\pgfpathquadraticcurveto{\pgfqpoint{11.131486in}{0.850170in}}{\pgfqpoint{11.127073in}{0.851377in}}%
\pgfpathquadraticcurveto{\pgfqpoint{11.122660in}{0.852584in}}{\pgfqpoint{11.118157in}{0.852584in}}%
\pgfpathquadraticcurveto{\pgfqpoint{11.108071in}{0.852584in}}{\pgfqpoint{11.102487in}{0.846189in}}%
\pgfpathquadraticcurveto{\pgfqpoint{11.096921in}{0.839813in}}{\pgfqpoint{11.096921in}{0.828267in}}%
\pgfpathquadraticcurveto{\pgfqpoint{11.096921in}{0.816721in}}{\pgfqpoint{11.102487in}{0.810327in}}%
\pgfpathquadraticcurveto{\pgfqpoint{11.108071in}{0.803951in}}{\pgfqpoint{11.118157in}{0.803951in}}%
\pgfpathquadraticcurveto{\pgfqpoint{11.122660in}{0.803951in}}{\pgfqpoint{11.127073in}{0.805158in}}%
\pgfpathquadraticcurveto{\pgfqpoint{11.131486in}{0.806364in}}{\pgfqpoint{11.135881in}{0.808796in}}%
\pgfpathlineto{\pgfqpoint{11.135881in}{0.799214in}}%
\pgfpathquadraticcurveto{\pgfqpoint{11.131540in}{0.797196in}}{\pgfqpoint{11.126893in}{0.796188in}}%
\pgfpathquadraticcurveto{\pgfqpoint{11.122264in}{0.795161in}}{\pgfqpoint{11.117023in}{0.795161in}}%
\pgfpathquadraticcurveto{\pgfqpoint{11.102775in}{0.795161in}}{\pgfqpoint{11.094381in}{0.804113in}}%
\pgfpathquadraticcurveto{\pgfqpoint{11.086006in}{0.813065in}}{\pgfqpoint{11.086006in}{0.828267in}}%
\pgfpathquadraticcurveto{\pgfqpoint{11.086006in}{0.843686in}}{\pgfqpoint{11.094471in}{0.852512in}}%
\pgfpathquadraticcurveto{\pgfqpoint{11.102955in}{0.861356in}}{\pgfqpoint{11.117707in}{0.861356in}}%
\pgfpathquadraticcurveto{\pgfqpoint{11.122480in}{0.861356in}}{\pgfqpoint{11.127037in}{0.860365in}}%
\pgfpathquadraticcurveto{\pgfqpoint{11.131594in}{0.859392in}}{\pgfqpoint{11.135881in}{0.857429in}}%
\pgfpathlineto{\pgfqpoint{11.135881in}{0.857429in}}%
\pgfpathclose%
\pgfpathmoveto{\pgfqpoint{11.207822in}{0.830915in}}%
\pgfpathlineto{\pgfqpoint{11.207822in}{0.825854in}}%
\pgfpathlineto{\pgfqpoint{11.160198in}{0.825854in}}%
\pgfpathquadraticcurveto{\pgfqpoint{11.160882in}{0.815154in}}{\pgfqpoint{11.166646in}{0.809553in}}%
\pgfpathquadraticcurveto{\pgfqpoint{11.172410in}{0.803951in}}{\pgfqpoint{11.182713in}{0.803951in}}%
\pgfpathquadraticcurveto{\pgfqpoint{11.188675in}{0.803951in}}{\pgfqpoint{11.194277in}{0.805410in}}%
\pgfpathquadraticcurveto{\pgfqpoint{11.199878in}{0.806869in}}{\pgfqpoint{11.205408in}{0.809805in}}%
\pgfpathlineto{\pgfqpoint{11.205408in}{0.800006in}}%
\pgfpathquadraticcurveto{\pgfqpoint{11.199824in}{0.797647in}}{\pgfqpoint{11.193970in}{0.796404in}}%
\pgfpathquadraticcurveto{\pgfqpoint{11.188117in}{0.795161in}}{\pgfqpoint{11.182100in}{0.795161in}}%
\pgfpathquadraticcurveto{\pgfqpoint{11.167006in}{0.795161in}}{\pgfqpoint{11.158198in}{0.803933in}}%
\pgfpathquadraticcurveto{\pgfqpoint{11.149390in}{0.812723in}}{\pgfqpoint{11.149390in}{0.827709in}}%
\pgfpathquadraticcurveto{\pgfqpoint{11.149390in}{0.843181in}}{\pgfqpoint{11.157748in}{0.852259in}}%
\pgfpathquadraticcurveto{\pgfqpoint{11.166106in}{0.861356in}}{\pgfqpoint{11.180299in}{0.861356in}}%
\pgfpathquadraticcurveto{\pgfqpoint{11.193016in}{0.861356in}}{\pgfqpoint{11.200419in}{0.853160in}}%
\pgfpathquadraticcurveto{\pgfqpoint{11.207822in}{0.844983in}}{\pgfqpoint{11.207822in}{0.830915in}}%
\pgfpathlineto{\pgfqpoint{11.207822in}{0.830915in}}%
\pgfpathclose%
\pgfpathmoveto{\pgfqpoint{11.197465in}{0.833959in}}%
\pgfpathquadraticcurveto{\pgfqpoint{11.197357in}{0.842443in}}{\pgfqpoint{11.192710in}{0.847504in}}%
\pgfpathquadraticcurveto{\pgfqpoint{11.188062in}{0.852584in}}{\pgfqpoint{11.180407in}{0.852584in}}%
\pgfpathquadraticcurveto{\pgfqpoint{11.171743in}{0.852584in}}{\pgfqpoint{11.166538in}{0.847684in}}%
\pgfpathquadraticcurveto{\pgfqpoint{11.161332in}{0.842785in}}{\pgfqpoint{11.160540in}{0.833887in}}%
\pgfpathlineto{\pgfqpoint{11.197465in}{0.833959in}}%
\pgfpathlineto{\pgfqpoint{11.197465in}{0.833959in}}%
\pgfpathclose%
\pgfusepath{fill}%
\end{pgfscope}%
\begin{pgfscope}%
\pgfsetbuttcap%
\pgfsetroundjoin%
\definecolor{currentfill}{rgb}{0.309804,0.372549,0.419608}%
\pgfsetfillcolor{currentfill}%
\pgfsetlinewidth{0.000000pt}%
\definecolor{currentstroke}{rgb}{0.309804,0.372549,0.419608}%
\pgfsetstrokecolor{currentstroke}%
\pgfsetdash{}{0pt}%
\pgfpathmoveto{\pgfqpoint{11.290262in}{0.859843in}}%
\pgfpathlineto{\pgfqpoint{11.300619in}{0.859843in}}%
\pgfpathlineto{\pgfqpoint{11.313570in}{0.810651in}}%
\pgfpathlineto{\pgfqpoint{11.326448in}{0.859843in}}%
\pgfpathlineto{\pgfqpoint{11.338661in}{0.859843in}}%
\pgfpathlineto{\pgfqpoint{11.351611in}{0.810651in}}%
\pgfpathlineto{\pgfqpoint{11.364508in}{0.859843in}}%
\pgfpathlineto{\pgfqpoint{11.374865in}{0.859843in}}%
\pgfpathlineto{\pgfqpoint{11.358366in}{0.796800in}}%
\pgfpathlineto{\pgfqpoint{11.346154in}{0.796800in}}%
\pgfpathlineto{\pgfqpoint{11.332590in}{0.848477in}}%
\pgfpathlineto{\pgfqpoint{11.318973in}{0.796800in}}%
\pgfpathlineto{\pgfqpoint{11.306743in}{0.796800in}}%
\pgfpathlineto{\pgfqpoint{11.290262in}{0.859843in}}%
\pgfpathlineto{\pgfqpoint{11.290262in}{0.859843in}}%
\pgfpathclose%
\pgfpathmoveto{\pgfqpoint{11.419211in}{0.828483in}}%
\pgfpathquadraticcurveto{\pgfqpoint{11.406656in}{0.828483in}}{\pgfqpoint{11.401811in}{0.825619in}}%
\pgfpathquadraticcurveto{\pgfqpoint{11.396966in}{0.822756in}}{\pgfqpoint{11.396966in}{0.815821in}}%
\pgfpathquadraticcurveto{\pgfqpoint{11.396966in}{0.810309in}}{\pgfqpoint{11.400604in}{0.807067in}}%
\pgfpathquadraticcurveto{\pgfqpoint{11.404243in}{0.803843in}}{\pgfqpoint{11.410475in}{0.803843in}}%
\pgfpathquadraticcurveto{\pgfqpoint{11.419103in}{0.803843in}}{\pgfqpoint{11.424308in}{0.809949in}}%
\pgfpathquadraticcurveto{\pgfqpoint{11.429514in}{0.816055in}}{\pgfqpoint{11.429514in}{0.826178in}}%
\pgfpathlineto{\pgfqpoint{11.429514in}{0.828483in}}%
\pgfpathlineto{\pgfqpoint{11.419211in}{0.828483in}}%
\pgfpathlineto{\pgfqpoint{11.419211in}{0.828483in}}%
\pgfpathclose%
\pgfpathmoveto{\pgfqpoint{11.439871in}{0.832770in}}%
\pgfpathlineto{\pgfqpoint{11.439871in}{0.796800in}}%
\pgfpathlineto{\pgfqpoint{11.429514in}{0.796800in}}%
\pgfpathlineto{\pgfqpoint{11.429514in}{0.806364in}}%
\pgfpathquadraticcurveto{\pgfqpoint{11.425965in}{0.800637in}}{\pgfqpoint{11.420670in}{0.797899in}}%
\pgfpathquadraticcurveto{\pgfqpoint{11.415374in}{0.795161in}}{\pgfqpoint{11.407719in}{0.795161in}}%
\pgfpathquadraticcurveto{\pgfqpoint{11.398047in}{0.795161in}}{\pgfqpoint{11.392319in}{0.800601in}}%
\pgfpathquadraticcurveto{\pgfqpoint{11.386609in}{0.806040in}}{\pgfqpoint{11.386609in}{0.815154in}}%
\pgfpathquadraticcurveto{\pgfqpoint{11.386609in}{0.825782in}}{\pgfqpoint{11.393724in}{0.831185in}}%
\pgfpathquadraticcurveto{\pgfqpoint{11.400857in}{0.836589in}}{\pgfqpoint{11.414978in}{0.836589in}}%
\pgfpathlineto{\pgfqpoint{11.429514in}{0.836589in}}%
\pgfpathlineto{\pgfqpoint{11.429514in}{0.837616in}}%
\pgfpathquadraticcurveto{\pgfqpoint{11.429514in}{0.844766in}}{\pgfqpoint{11.424813in}{0.848675in}}%
\pgfpathquadraticcurveto{\pgfqpoint{11.420112in}{0.852584in}}{\pgfqpoint{11.411610in}{0.852584in}}%
\pgfpathquadraticcurveto{\pgfqpoint{11.406206in}{0.852584in}}{\pgfqpoint{11.401073in}{0.851287in}}%
\pgfpathquadraticcurveto{\pgfqpoint{11.395957in}{0.849990in}}{\pgfqpoint{11.391238in}{0.847396in}}%
\pgfpathlineto{\pgfqpoint{11.391238in}{0.856979in}}%
\pgfpathquadraticcurveto{\pgfqpoint{11.396912in}{0.859176in}}{\pgfqpoint{11.402262in}{0.860257in}}%
\pgfpathquadraticcurveto{\pgfqpoint{11.407611in}{0.861356in}}{\pgfqpoint{11.412673in}{0.861356in}}%
\pgfpathquadraticcurveto{\pgfqpoint{11.426362in}{0.861356in}}{\pgfqpoint{11.433116in}{0.854259in}}%
\pgfpathquadraticcurveto{\pgfqpoint{11.439871in}{0.847180in}}{\pgfqpoint{11.439871in}{0.832770in}}%
\pgfpathlineto{\pgfqpoint{11.439871in}{0.832770in}}%
\pgfpathclose%
\pgfpathmoveto{\pgfqpoint{11.497726in}{0.850170in}}%
\pgfpathquadraticcurveto{\pgfqpoint{11.495979in}{0.851179in}}{\pgfqpoint{11.493925in}{0.851647in}}%
\pgfpathquadraticcurveto{\pgfqpoint{11.491872in}{0.852133in}}{\pgfqpoint{11.489404in}{0.852133in}}%
\pgfpathquadraticcurveto{\pgfqpoint{11.480614in}{0.852133in}}{\pgfqpoint{11.475913in}{0.846423in}}%
\pgfpathquadraticcurveto{\pgfqpoint{11.471212in}{0.840714in}}{\pgfqpoint{11.471212in}{0.830014in}}%
\pgfpathlineto{\pgfqpoint{11.471212in}{0.796800in}}%
\pgfpathlineto{\pgfqpoint{11.460801in}{0.796800in}}%
\pgfpathlineto{\pgfqpoint{11.460801in}{0.859843in}}%
\pgfpathlineto{\pgfqpoint{11.471212in}{0.859843in}}%
\pgfpathlineto{\pgfqpoint{11.471212in}{0.850044in}}%
\pgfpathquadraticcurveto{\pgfqpoint{11.474490in}{0.855790in}}{\pgfqpoint{11.479714in}{0.858564in}}%
\pgfpathquadraticcurveto{\pgfqpoint{11.484955in}{0.861356in}}{\pgfqpoint{11.492448in}{0.861356in}}%
\pgfpathquadraticcurveto{\pgfqpoint{11.493511in}{0.861356in}}{\pgfqpoint{11.494808in}{0.861211in}}%
\pgfpathquadraticcurveto{\pgfqpoint{11.496105in}{0.861085in}}{\pgfqpoint{11.497672in}{0.860797in}}%
\pgfpathlineto{\pgfqpoint{11.497726in}{0.850170in}}%
\pgfpathlineto{\pgfqpoint{11.497726in}{0.850170in}}%
\pgfpathclose%
\pgfpathmoveto{\pgfqpoint{11.558985in}{0.834860in}}%
\pgfpathlineto{\pgfqpoint{11.558985in}{0.796800in}}%
\pgfpathlineto{\pgfqpoint{11.548628in}{0.796800in}}%
\pgfpathlineto{\pgfqpoint{11.548628in}{0.834517in}}%
\pgfpathquadraticcurveto{\pgfqpoint{11.548628in}{0.843469in}}{\pgfqpoint{11.545134in}{0.847900in}}%
\pgfpathquadraticcurveto{\pgfqpoint{11.541640in}{0.852349in}}{\pgfqpoint{11.534669in}{0.852349in}}%
\pgfpathquadraticcurveto{\pgfqpoint{11.526275in}{0.852349in}}{\pgfqpoint{11.521430in}{0.847000in}}%
\pgfpathquadraticcurveto{\pgfqpoint{11.516585in}{0.841668in}}{\pgfqpoint{11.516585in}{0.832428in}}%
\pgfpathlineto{\pgfqpoint{11.516585in}{0.796800in}}%
\pgfpathlineto{\pgfqpoint{11.506174in}{0.796800in}}%
\pgfpathlineto{\pgfqpoint{11.506174in}{0.859843in}}%
\pgfpathlineto{\pgfqpoint{11.516585in}{0.859843in}}%
\pgfpathlineto{\pgfqpoint{11.516585in}{0.850044in}}%
\pgfpathquadraticcurveto{\pgfqpoint{11.520313in}{0.855736in}}{\pgfqpoint{11.525339in}{0.858546in}}%
\pgfpathquadraticcurveto{\pgfqpoint{11.530382in}{0.861356in}}{\pgfqpoint{11.536974in}{0.861356in}}%
\pgfpathquadraticcurveto{\pgfqpoint{11.547836in}{0.861356in}}{\pgfqpoint{11.553401in}{0.854637in}}%
\pgfpathquadraticcurveto{\pgfqpoint{11.558985in}{0.847918in}}{\pgfqpoint{11.558985in}{0.834860in}}%
\pgfpathlineto{\pgfqpoint{11.558985in}{0.834860in}}%
\pgfpathclose%
\pgfpathmoveto{\pgfqpoint{11.579627in}{0.859843in}}%
\pgfpathlineto{\pgfqpoint{11.589984in}{0.859843in}}%
\pgfpathlineto{\pgfqpoint{11.589984in}{0.796800in}}%
\pgfpathlineto{\pgfqpoint{11.579627in}{0.796800in}}%
\pgfpathlineto{\pgfqpoint{11.579627in}{0.859843in}}%
\pgfpathlineto{\pgfqpoint{11.579627in}{0.859843in}}%
\pgfpathclose%
\pgfpathmoveto{\pgfqpoint{11.579627in}{0.884393in}}%
\pgfpathlineto{\pgfqpoint{11.589984in}{0.884393in}}%
\pgfpathlineto{\pgfqpoint{11.589984in}{0.871262in}}%
\pgfpathlineto{\pgfqpoint{11.579627in}{0.871262in}}%
\pgfpathlineto{\pgfqpoint{11.579627in}{0.884393in}}%
\pgfpathlineto{\pgfqpoint{11.579627in}{0.884393in}}%
\pgfpathclose%
\pgfpathmoveto{\pgfqpoint{11.664068in}{0.834860in}}%
\pgfpathlineto{\pgfqpoint{11.664068in}{0.796800in}}%
\pgfpathlineto{\pgfqpoint{11.653711in}{0.796800in}}%
\pgfpathlineto{\pgfqpoint{11.653711in}{0.834517in}}%
\pgfpathquadraticcurveto{\pgfqpoint{11.653711in}{0.843469in}}{\pgfqpoint{11.650217in}{0.847900in}}%
\pgfpathquadraticcurveto{\pgfqpoint{11.646722in}{0.852349in}}{\pgfqpoint{11.639752in}{0.852349in}}%
\pgfpathquadraticcurveto{\pgfqpoint{11.631358in}{0.852349in}}{\pgfqpoint{11.626513in}{0.847000in}}%
\pgfpathquadraticcurveto{\pgfqpoint{11.621668in}{0.841668in}}{\pgfqpoint{11.621668in}{0.832428in}}%
\pgfpathlineto{\pgfqpoint{11.621668in}{0.796800in}}%
\pgfpathlineto{\pgfqpoint{11.611257in}{0.796800in}}%
\pgfpathlineto{\pgfqpoint{11.611257in}{0.859843in}}%
\pgfpathlineto{\pgfqpoint{11.621668in}{0.859843in}}%
\pgfpathlineto{\pgfqpoint{11.621668in}{0.850044in}}%
\pgfpathquadraticcurveto{\pgfqpoint{11.625396in}{0.855736in}}{\pgfqpoint{11.630421in}{0.858546in}}%
\pgfpathquadraticcurveto{\pgfqpoint{11.635465in}{0.861356in}}{\pgfqpoint{11.642057in}{0.861356in}}%
\pgfpathquadraticcurveto{\pgfqpoint{11.652919in}{0.861356in}}{\pgfqpoint{11.658484in}{0.854637in}}%
\pgfpathquadraticcurveto{\pgfqpoint{11.664068in}{0.847918in}}{\pgfqpoint{11.664068in}{0.834860in}}%
\pgfpathlineto{\pgfqpoint{11.664068in}{0.834860in}}%
\pgfpathclose%
\pgfpathmoveto{\pgfqpoint{11.726192in}{0.829060in}}%
\pgfpathquadraticcurveto{\pgfqpoint{11.726192in}{0.840317in}}{\pgfqpoint{11.721545in}{0.846496in}}%
\pgfpathquadraticcurveto{\pgfqpoint{11.716916in}{0.852692in}}{\pgfqpoint{11.708522in}{0.852692in}}%
\pgfpathquadraticcurveto{\pgfqpoint{11.700201in}{0.852692in}}{\pgfqpoint{11.695553in}{0.846496in}}%
\pgfpathquadraticcurveto{\pgfqpoint{11.690906in}{0.840317in}}{\pgfqpoint{11.690906in}{0.829060in}}%
\pgfpathquadraticcurveto{\pgfqpoint{11.690906in}{0.817856in}}{\pgfqpoint{11.695553in}{0.811660in}}%
\pgfpathquadraticcurveto{\pgfqpoint{11.700201in}{0.805464in}}{\pgfqpoint{11.708522in}{0.805464in}}%
\pgfpathquadraticcurveto{\pgfqpoint{11.716916in}{0.805464in}}{\pgfqpoint{11.721545in}{0.811660in}}%
\pgfpathquadraticcurveto{\pgfqpoint{11.726192in}{0.817856in}}{\pgfqpoint{11.726192in}{0.829060in}}%
\pgfpathlineto{\pgfqpoint{11.726192in}{0.829060in}}%
\pgfpathclose%
\pgfpathmoveto{\pgfqpoint{11.736549in}{0.804617in}}%
\pgfpathquadraticcurveto{\pgfqpoint{11.736549in}{0.788532in}}{\pgfqpoint{11.729398in}{0.780679in}}%
\pgfpathquadraticcurveto{\pgfqpoint{11.722265in}{0.772826in}}{\pgfqpoint{11.707513in}{0.772826in}}%
\pgfpathquadraticcurveto{\pgfqpoint{11.702056in}{0.772826in}}{\pgfqpoint{11.697211in}{0.773636in}}%
\pgfpathquadraticcurveto{\pgfqpoint{11.692365in}{0.774447in}}{\pgfqpoint{11.687808in}{0.776140in}}%
\pgfpathlineto{\pgfqpoint{11.687808in}{0.786209in}}%
\pgfpathquadraticcurveto{\pgfqpoint{11.692365in}{0.783741in}}{\pgfqpoint{11.696814in}{0.782570in}}%
\pgfpathquadraticcurveto{\pgfqpoint{11.701263in}{0.781382in}}{\pgfqpoint{11.705874in}{0.781382in}}%
\pgfpathquadraticcurveto{\pgfqpoint{11.716069in}{0.781382in}}{\pgfqpoint{11.721131in}{0.786695in}}%
\pgfpathquadraticcurveto{\pgfqpoint{11.726192in}{0.792009in}}{\pgfqpoint{11.726192in}{0.802762in}}%
\pgfpathlineto{\pgfqpoint{11.726192in}{0.807895in}}%
\pgfpathquadraticcurveto{\pgfqpoint{11.722986in}{0.802312in}}{\pgfqpoint{11.717979in}{0.799556in}}%
\pgfpathquadraticcurveto{\pgfqpoint{11.712971in}{0.796800in}}{\pgfqpoint{11.705982in}{0.796800in}}%
\pgfpathquadraticcurveto{\pgfqpoint{11.694401in}{0.796800in}}{\pgfqpoint{11.687304in}{0.805626in}}%
\pgfpathquadraticcurveto{\pgfqpoint{11.680207in}{0.814470in}}{\pgfqpoint{11.680207in}{0.829060in}}%
\pgfpathquadraticcurveto{\pgfqpoint{11.680207in}{0.843686in}}{\pgfqpoint{11.687304in}{0.852512in}}%
\pgfpathquadraticcurveto{\pgfqpoint{11.694401in}{0.861356in}}{\pgfqpoint{11.705982in}{0.861356in}}%
\pgfpathquadraticcurveto{\pgfqpoint{11.712971in}{0.861356in}}{\pgfqpoint{11.717979in}{0.858600in}}%
\pgfpathquadraticcurveto{\pgfqpoint{11.722986in}{0.855844in}}{\pgfqpoint{11.726192in}{0.850278in}}%
\pgfpathlineto{\pgfqpoint{11.726192in}{0.859843in}}%
\pgfpathlineto{\pgfqpoint{11.736549in}{0.859843in}}%
\pgfpathlineto{\pgfqpoint{11.736549in}{0.804617in}}%
\pgfpathlineto{\pgfqpoint{11.736549in}{0.804617in}}%
\pgfpathclose%
\pgfpathmoveto{\pgfqpoint{11.798079in}{0.857987in}}%
\pgfpathlineto{\pgfqpoint{11.798079in}{0.848189in}}%
\pgfpathquadraticcurveto{\pgfqpoint{11.793702in}{0.850440in}}{\pgfqpoint{11.788964in}{0.851557in}}%
\pgfpathquadraticcurveto{\pgfqpoint{11.784245in}{0.852692in}}{\pgfqpoint{11.779166in}{0.852692in}}%
\pgfpathquadraticcurveto{\pgfqpoint{11.771457in}{0.852692in}}{\pgfqpoint{11.767602in}{0.850332in}}%
\pgfpathquadraticcurveto{\pgfqpoint{11.763747in}{0.847973in}}{\pgfqpoint{11.763747in}{0.843235in}}%
\pgfpathquadraticcurveto{\pgfqpoint{11.763747in}{0.839633in}}{\pgfqpoint{11.766503in}{0.837580in}}%
\pgfpathquadraticcurveto{\pgfqpoint{11.769259in}{0.835526in}}{\pgfqpoint{11.777599in}{0.833671in}}%
\pgfpathlineto{\pgfqpoint{11.781147in}{0.832878in}}%
\pgfpathquadraticcurveto{\pgfqpoint{11.792171in}{0.830519in}}{\pgfqpoint{11.796818in}{0.826214in}}%
\pgfpathquadraticcurveto{\pgfqpoint{11.801465in}{0.821909in}}{\pgfqpoint{11.801465in}{0.814200in}}%
\pgfpathquadraticcurveto{\pgfqpoint{11.801465in}{0.805410in}}{\pgfqpoint{11.794512in}{0.800276in}}%
\pgfpathquadraticcurveto{\pgfqpoint{11.787559in}{0.795161in}}{\pgfqpoint{11.775401in}{0.795161in}}%
\pgfpathquadraticcurveto{\pgfqpoint{11.770340in}{0.795161in}}{\pgfqpoint{11.764846in}{0.796152in}}%
\pgfpathquadraticcurveto{\pgfqpoint{11.759352in}{0.797142in}}{\pgfqpoint{11.753282in}{0.799106in}}%
\pgfpathlineto{\pgfqpoint{11.753282in}{0.809805in}}%
\pgfpathquadraticcurveto{\pgfqpoint{11.759028in}{0.806815in}}{\pgfqpoint{11.764594in}{0.805320in}}%
\pgfpathquadraticcurveto{\pgfqpoint{11.770160in}{0.803843in}}{\pgfqpoint{11.775635in}{0.803843in}}%
\pgfpathquadraticcurveto{\pgfqpoint{11.782948in}{0.803843in}}{\pgfqpoint{11.786875in}{0.806346in}}%
\pgfpathquadraticcurveto{\pgfqpoint{11.790820in}{0.808850in}}{\pgfqpoint{11.790820in}{0.813407in}}%
\pgfpathquadraticcurveto{\pgfqpoint{11.790820in}{0.817622in}}{\pgfqpoint{11.787974in}{0.819874in}}%
\pgfpathquadraticcurveto{\pgfqpoint{11.785146in}{0.822125in}}{\pgfqpoint{11.775509in}{0.824214in}}%
\pgfpathlineto{\pgfqpoint{11.771907in}{0.825061in}}%
\pgfpathquadraticcurveto{\pgfqpoint{11.762288in}{0.827078in}}{\pgfqpoint{11.758002in}{0.831275in}}%
\pgfpathquadraticcurveto{\pgfqpoint{11.753733in}{0.835472in}}{\pgfqpoint{11.753733in}{0.842785in}}%
\pgfpathquadraticcurveto{\pgfqpoint{11.753733in}{0.851683in}}{\pgfqpoint{11.760037in}{0.856510in}}%
\pgfpathquadraticcurveto{\pgfqpoint{11.766341in}{0.861356in}}{\pgfqpoint{11.777941in}{0.861356in}}%
\pgfpathquadraticcurveto{\pgfqpoint{11.783669in}{0.861356in}}{\pgfqpoint{11.788730in}{0.860509in}}%
\pgfpathquadraticcurveto{\pgfqpoint{11.793810in}{0.859680in}}{\pgfqpoint{11.798079in}{0.857987in}}%
\pgfpathlineto{\pgfqpoint{11.798079in}{0.857987in}}%
\pgfpathclose%
\pgfpathmoveto{\pgfqpoint{11.820594in}{0.856402in}}%
\pgfpathlineto{\pgfqpoint{11.832464in}{0.856402in}}%
\pgfpathlineto{\pgfqpoint{11.832464in}{0.842119in}}%
\pgfpathlineto{\pgfqpoint{11.820594in}{0.842119in}}%
\pgfpathlineto{\pgfqpoint{11.820594in}{0.856402in}}%
\pgfpathlineto{\pgfqpoint{11.820594in}{0.856402in}}%
\pgfpathclose%
\pgfpathmoveto{\pgfqpoint{11.820594in}{0.811102in}}%
\pgfpathlineto{\pgfqpoint{11.832464in}{0.811102in}}%
\pgfpathlineto{\pgfqpoint{11.832464in}{0.801411in}}%
\pgfpathlineto{\pgfqpoint{11.823242in}{0.783399in}}%
\pgfpathlineto{\pgfqpoint{11.815983in}{0.783399in}}%
\pgfpathlineto{\pgfqpoint{11.820594in}{0.801411in}}%
\pgfpathlineto{\pgfqpoint{11.820594in}{0.811102in}}%
\pgfpathlineto{\pgfqpoint{11.820594in}{0.811102in}}%
\pgfpathclose%
\pgfusepath{fill}%
\end{pgfscope}%
\begin{pgfscope}%
\pgfsetbuttcap%
\pgfsetroundjoin%
\definecolor{currentfill}{rgb}{0.309804,0.372549,0.419608}%
\pgfsetfillcolor{currentfill}%
\pgfsetlinewidth{0.000000pt}%
\definecolor{currentstroke}{rgb}{0.309804,0.372549,0.419608}%
\pgfsetstrokecolor{currentstroke}%
\pgfsetdash{}{0pt}%
\pgfpathmoveto{\pgfqpoint{11.978137in}{0.857987in}}%
\pgfpathlineto{\pgfqpoint{11.978137in}{0.848189in}}%
\pgfpathquadraticcurveto{\pgfqpoint{11.973760in}{0.850440in}}{\pgfqpoint{11.969023in}{0.851557in}}%
\pgfpathquadraticcurveto{\pgfqpoint{11.964303in}{0.852692in}}{\pgfqpoint{11.959224in}{0.852692in}}%
\pgfpathquadraticcurveto{\pgfqpoint{11.951515in}{0.852692in}}{\pgfqpoint{11.947660in}{0.850332in}}%
\pgfpathquadraticcurveto{\pgfqpoint{11.943806in}{0.847973in}}{\pgfqpoint{11.943806in}{0.843235in}}%
\pgfpathquadraticcurveto{\pgfqpoint{11.943806in}{0.839633in}}{\pgfqpoint{11.946561in}{0.837580in}}%
\pgfpathquadraticcurveto{\pgfqpoint{11.949317in}{0.835526in}}{\pgfqpoint{11.957657in}{0.833671in}}%
\pgfpathlineto{\pgfqpoint{11.961205in}{0.832878in}}%
\pgfpathquadraticcurveto{\pgfqpoint{11.972229in}{0.830519in}}{\pgfqpoint{11.976876in}{0.826214in}}%
\pgfpathquadraticcurveto{\pgfqpoint{11.981523in}{0.821909in}}{\pgfqpoint{11.981523in}{0.814200in}}%
\pgfpathquadraticcurveto{\pgfqpoint{11.981523in}{0.805410in}}{\pgfqpoint{11.974570in}{0.800276in}}%
\pgfpathquadraticcurveto{\pgfqpoint{11.967618in}{0.795161in}}{\pgfqpoint{11.955459in}{0.795161in}}%
\pgfpathquadraticcurveto{\pgfqpoint{11.950398in}{0.795161in}}{\pgfqpoint{11.944904in}{0.796152in}}%
\pgfpathquadraticcurveto{\pgfqpoint{11.939411in}{0.797142in}}{\pgfqpoint{11.933340in}{0.799106in}}%
\pgfpathlineto{\pgfqpoint{11.933340in}{0.809805in}}%
\pgfpathquadraticcurveto{\pgfqpoint{11.939086in}{0.806815in}}{\pgfqpoint{11.944652in}{0.805320in}}%
\pgfpathquadraticcurveto{\pgfqpoint{11.950218in}{0.803843in}}{\pgfqpoint{11.955694in}{0.803843in}}%
\pgfpathquadraticcurveto{\pgfqpoint{11.963007in}{0.803843in}}{\pgfqpoint{11.966933in}{0.806346in}}%
\pgfpathquadraticcurveto{\pgfqpoint{11.970878in}{0.808850in}}{\pgfqpoint{11.970878in}{0.813407in}}%
\pgfpathquadraticcurveto{\pgfqpoint{11.970878in}{0.817622in}}{\pgfqpoint{11.968032in}{0.819874in}}%
\pgfpathquadraticcurveto{\pgfqpoint{11.965204in}{0.822125in}}{\pgfqpoint{11.955567in}{0.824214in}}%
\pgfpathlineto{\pgfqpoint{11.951965in}{0.825061in}}%
\pgfpathquadraticcurveto{\pgfqpoint{11.942347in}{0.827078in}}{\pgfqpoint{11.938060in}{0.831275in}}%
\pgfpathquadraticcurveto{\pgfqpoint{11.933791in}{0.835472in}}{\pgfqpoint{11.933791in}{0.842785in}}%
\pgfpathquadraticcurveto{\pgfqpoint{11.933791in}{0.851683in}}{\pgfqpoint{11.940095in}{0.856510in}}%
\pgfpathquadraticcurveto{\pgfqpoint{11.946399in}{0.861356in}}{\pgfqpoint{11.957999in}{0.861356in}}%
\pgfpathquadraticcurveto{\pgfqpoint{11.963727in}{0.861356in}}{\pgfqpoint{11.968788in}{0.860509in}}%
\pgfpathquadraticcurveto{\pgfqpoint{11.973868in}{0.859680in}}{\pgfqpoint{11.978137in}{0.857987in}}%
\pgfpathlineto{\pgfqpoint{11.978137in}{0.857987in}}%
\pgfpathclose%
\pgfpathmoveto{\pgfqpoint{11.998004in}{0.859843in}}%
\pgfpathlineto{\pgfqpoint{12.008361in}{0.859843in}}%
\pgfpathlineto{\pgfqpoint{12.008361in}{0.796800in}}%
\pgfpathlineto{\pgfqpoint{11.998004in}{0.796800in}}%
\pgfpathlineto{\pgfqpoint{11.998004in}{0.859843in}}%
\pgfpathlineto{\pgfqpoint{11.998004in}{0.859843in}}%
\pgfpathclose%
\pgfpathmoveto{\pgfqpoint{11.998004in}{0.884393in}}%
\pgfpathlineto{\pgfqpoint{12.008361in}{0.884393in}}%
\pgfpathlineto{\pgfqpoint{12.008361in}{0.871262in}}%
\pgfpathlineto{\pgfqpoint{11.998004in}{0.871262in}}%
\pgfpathlineto{\pgfqpoint{11.998004in}{0.884393in}}%
\pgfpathlineto{\pgfqpoint{11.998004in}{0.884393in}}%
\pgfpathclose%
\pgfpathmoveto{\pgfqpoint{12.071512in}{0.829060in}}%
\pgfpathquadraticcurveto{\pgfqpoint{12.071512in}{0.840317in}}{\pgfqpoint{12.066865in}{0.846496in}}%
\pgfpathquadraticcurveto{\pgfqpoint{12.062235in}{0.852692in}}{\pgfqpoint{12.053842in}{0.852692in}}%
\pgfpathquadraticcurveto{\pgfqpoint{12.045520in}{0.852692in}}{\pgfqpoint{12.040873in}{0.846496in}}%
\pgfpathquadraticcurveto{\pgfqpoint{12.036226in}{0.840317in}}{\pgfqpoint{12.036226in}{0.829060in}}%
\pgfpathquadraticcurveto{\pgfqpoint{12.036226in}{0.817856in}}{\pgfqpoint{12.040873in}{0.811660in}}%
\pgfpathquadraticcurveto{\pgfqpoint{12.045520in}{0.805464in}}{\pgfqpoint{12.053842in}{0.805464in}}%
\pgfpathquadraticcurveto{\pgfqpoint{12.062235in}{0.805464in}}{\pgfqpoint{12.066865in}{0.811660in}}%
\pgfpathquadraticcurveto{\pgfqpoint{12.071512in}{0.817856in}}{\pgfqpoint{12.071512in}{0.829060in}}%
\pgfpathlineto{\pgfqpoint{12.071512in}{0.829060in}}%
\pgfpathclose%
\pgfpathmoveto{\pgfqpoint{12.081869in}{0.804617in}}%
\pgfpathquadraticcurveto{\pgfqpoint{12.081869in}{0.788532in}}{\pgfqpoint{12.074718in}{0.780679in}}%
\pgfpathquadraticcurveto{\pgfqpoint{12.067585in}{0.772826in}}{\pgfqpoint{12.052833in}{0.772826in}}%
\pgfpathquadraticcurveto{\pgfqpoint{12.047375in}{0.772826in}}{\pgfqpoint{12.042530in}{0.773636in}}%
\pgfpathquadraticcurveto{\pgfqpoint{12.037685in}{0.774447in}}{\pgfqpoint{12.033128in}{0.776140in}}%
\pgfpathlineto{\pgfqpoint{12.033128in}{0.786209in}}%
\pgfpathquadraticcurveto{\pgfqpoint{12.037685in}{0.783741in}}{\pgfqpoint{12.042134in}{0.782570in}}%
\pgfpathquadraticcurveto{\pgfqpoint{12.046583in}{0.781382in}}{\pgfqpoint{12.051194in}{0.781382in}}%
\pgfpathquadraticcurveto{\pgfqpoint{12.061389in}{0.781382in}}{\pgfqpoint{12.066450in}{0.786695in}}%
\pgfpathquadraticcurveto{\pgfqpoint{12.071512in}{0.792009in}}{\pgfqpoint{12.071512in}{0.802762in}}%
\pgfpathlineto{\pgfqpoint{12.071512in}{0.807895in}}%
\pgfpathquadraticcurveto{\pgfqpoint{12.068306in}{0.802312in}}{\pgfqpoint{12.063298in}{0.799556in}}%
\pgfpathquadraticcurveto{\pgfqpoint{12.058291in}{0.796800in}}{\pgfqpoint{12.051302in}{0.796800in}}%
\pgfpathquadraticcurveto{\pgfqpoint{12.039720in}{0.796800in}}{\pgfqpoint{12.032623in}{0.805626in}}%
\pgfpathquadraticcurveto{\pgfqpoint{12.025527in}{0.814470in}}{\pgfqpoint{12.025527in}{0.829060in}}%
\pgfpathquadraticcurveto{\pgfqpoint{12.025527in}{0.843686in}}{\pgfqpoint{12.032623in}{0.852512in}}%
\pgfpathquadraticcurveto{\pgfqpoint{12.039720in}{0.861356in}}{\pgfqpoint{12.051302in}{0.861356in}}%
\pgfpathquadraticcurveto{\pgfqpoint{12.058291in}{0.861356in}}{\pgfqpoint{12.063298in}{0.858600in}}%
\pgfpathquadraticcurveto{\pgfqpoint{12.068306in}{0.855844in}}{\pgfqpoint{12.071512in}{0.850278in}}%
\pgfpathlineto{\pgfqpoint{12.071512in}{0.859843in}}%
\pgfpathlineto{\pgfqpoint{12.081869in}{0.859843in}}%
\pgfpathlineto{\pgfqpoint{12.081869in}{0.804617in}}%
\pgfpathlineto{\pgfqpoint{12.081869in}{0.804617in}}%
\pgfpathclose%
\pgfpathmoveto{\pgfqpoint{12.155628in}{0.834860in}}%
\pgfpathlineto{\pgfqpoint{12.155628in}{0.796800in}}%
\pgfpathlineto{\pgfqpoint{12.145271in}{0.796800in}}%
\pgfpathlineto{\pgfqpoint{12.145271in}{0.834517in}}%
\pgfpathquadraticcurveto{\pgfqpoint{12.145271in}{0.843469in}}{\pgfqpoint{12.141777in}{0.847900in}}%
\pgfpathquadraticcurveto{\pgfqpoint{12.138283in}{0.852349in}}{\pgfqpoint{12.131312in}{0.852349in}}%
\pgfpathquadraticcurveto{\pgfqpoint{12.122918in}{0.852349in}}{\pgfqpoint{12.118073in}{0.847000in}}%
\pgfpathquadraticcurveto{\pgfqpoint{12.113228in}{0.841668in}}{\pgfqpoint{12.113228in}{0.832428in}}%
\pgfpathlineto{\pgfqpoint{12.113228in}{0.796800in}}%
\pgfpathlineto{\pgfqpoint{12.102817in}{0.796800in}}%
\pgfpathlineto{\pgfqpoint{12.102817in}{0.859843in}}%
\pgfpathlineto{\pgfqpoint{12.113228in}{0.859843in}}%
\pgfpathlineto{\pgfqpoint{12.113228in}{0.850044in}}%
\pgfpathquadraticcurveto{\pgfqpoint{12.116956in}{0.855736in}}{\pgfqpoint{12.121982in}{0.858546in}}%
\pgfpathquadraticcurveto{\pgfqpoint{12.127025in}{0.861356in}}{\pgfqpoint{12.133618in}{0.861356in}}%
\pgfpathquadraticcurveto{\pgfqpoint{12.144479in}{0.861356in}}{\pgfqpoint{12.150045in}{0.854637in}}%
\pgfpathquadraticcurveto{\pgfqpoint{12.155628in}{0.847918in}}{\pgfqpoint{12.155628in}{0.834860in}}%
\pgfpathlineto{\pgfqpoint{12.155628in}{0.834860in}}%
\pgfpathclose%
\pgfpathmoveto{\pgfqpoint{12.176270in}{0.859843in}}%
\pgfpathlineto{\pgfqpoint{12.186627in}{0.859843in}}%
\pgfpathlineto{\pgfqpoint{12.186627in}{0.796800in}}%
\pgfpathlineto{\pgfqpoint{12.176270in}{0.796800in}}%
\pgfpathlineto{\pgfqpoint{12.176270in}{0.859843in}}%
\pgfpathlineto{\pgfqpoint{12.176270in}{0.859843in}}%
\pgfpathclose%
\pgfpathmoveto{\pgfqpoint{12.176270in}{0.884393in}}%
\pgfpathlineto{\pgfqpoint{12.186627in}{0.884393in}}%
\pgfpathlineto{\pgfqpoint{12.186627in}{0.871262in}}%
\pgfpathlineto{\pgfqpoint{12.176270in}{0.871262in}}%
\pgfpathlineto{\pgfqpoint{12.176270in}{0.884393in}}%
\pgfpathlineto{\pgfqpoint{12.176270in}{0.884393in}}%
\pgfpathclose%
\pgfpathmoveto{\pgfqpoint{12.259234in}{0.859843in}}%
\pgfpathlineto{\pgfqpoint{12.259234in}{0.796800in}}%
\pgfpathlineto{\pgfqpoint{12.248823in}{0.796800in}}%
\pgfpathlineto{\pgfqpoint{12.248823in}{0.851791in}}%
\pgfpathlineto{\pgfqpoint{12.220400in}{0.851791in}}%
\pgfpathlineto{\pgfqpoint{12.220400in}{0.796800in}}%
\pgfpathlineto{\pgfqpoint{12.209989in}{0.796800in}}%
\pgfpathlineto{\pgfqpoint{12.209989in}{0.851791in}}%
\pgfpathlineto{\pgfqpoint{12.200082in}{0.851791in}}%
\pgfpathlineto{\pgfqpoint{12.200082in}{0.859843in}}%
\pgfpathlineto{\pgfqpoint{12.209989in}{0.859843in}}%
\pgfpathlineto{\pgfqpoint{12.209989in}{0.864238in}}%
\pgfpathquadraticcurveto{\pgfqpoint{12.209989in}{0.874540in}}{\pgfqpoint{12.214852in}{0.879458in}}%
\pgfpathquadraticcurveto{\pgfqpoint{12.219734in}{0.884393in}}{\pgfqpoint{12.229803in}{0.884393in}}%
\pgfpathlineto{\pgfqpoint{12.240214in}{0.884393in}}%
\pgfpathlineto{\pgfqpoint{12.240214in}{0.875765in}}%
\pgfpathlineto{\pgfqpoint{12.230307in}{0.875765in}}%
\pgfpathquadraticcurveto{\pgfqpoint{12.224741in}{0.875765in}}{\pgfqpoint{12.222562in}{0.873514in}}%
\pgfpathquadraticcurveto{\pgfqpoint{12.220400in}{0.871262in}}{\pgfqpoint{12.220400in}{0.865408in}}%
\pgfpathlineto{\pgfqpoint{12.220400in}{0.859843in}}%
\pgfpathlineto{\pgfqpoint{12.259234in}{0.859843in}}%
\pgfpathlineto{\pgfqpoint{12.259234in}{0.859843in}}%
\pgfpathclose%
\pgfpathmoveto{\pgfqpoint{12.248823in}{0.884267in}}%
\pgfpathlineto{\pgfqpoint{12.259234in}{0.884267in}}%
\pgfpathlineto{\pgfqpoint{12.259234in}{0.871154in}}%
\pgfpathlineto{\pgfqpoint{12.248823in}{0.871154in}}%
\pgfpathlineto{\pgfqpoint{12.248823in}{0.884267in}}%
\pgfpathlineto{\pgfqpoint{12.248823in}{0.884267in}}%
\pgfpathclose%
\pgfpathmoveto{\pgfqpoint{12.326276in}{0.857429in}}%
\pgfpathlineto{\pgfqpoint{12.326276in}{0.847738in}}%
\pgfpathquadraticcurveto{\pgfqpoint{12.321881in}{0.850170in}}{\pgfqpoint{12.317468in}{0.851377in}}%
\pgfpathquadraticcurveto{\pgfqpoint{12.313055in}{0.852584in}}{\pgfqpoint{12.308552in}{0.852584in}}%
\pgfpathquadraticcurveto{\pgfqpoint{12.298465in}{0.852584in}}{\pgfqpoint{12.292881in}{0.846189in}}%
\pgfpathquadraticcurveto{\pgfqpoint{12.287315in}{0.839813in}}{\pgfqpoint{12.287315in}{0.828267in}}%
\pgfpathquadraticcurveto{\pgfqpoint{12.287315in}{0.816721in}}{\pgfqpoint{12.292881in}{0.810327in}}%
\pgfpathquadraticcurveto{\pgfqpoint{12.298465in}{0.803951in}}{\pgfqpoint{12.308552in}{0.803951in}}%
\pgfpathquadraticcurveto{\pgfqpoint{12.313055in}{0.803951in}}{\pgfqpoint{12.317468in}{0.805158in}}%
\pgfpathquadraticcurveto{\pgfqpoint{12.321881in}{0.806364in}}{\pgfqpoint{12.326276in}{0.808796in}}%
\pgfpathlineto{\pgfqpoint{12.326276in}{0.799214in}}%
\pgfpathquadraticcurveto{\pgfqpoint{12.321935in}{0.797196in}}{\pgfqpoint{12.317288in}{0.796188in}}%
\pgfpathquadraticcurveto{\pgfqpoint{12.312658in}{0.795161in}}{\pgfqpoint{12.307417in}{0.795161in}}%
\pgfpathquadraticcurveto{\pgfqpoint{12.293169in}{0.795161in}}{\pgfqpoint{12.284776in}{0.804113in}}%
\pgfpathquadraticcurveto{\pgfqpoint{12.276400in}{0.813065in}}{\pgfqpoint{12.276400in}{0.828267in}}%
\pgfpathquadraticcurveto{\pgfqpoint{12.276400in}{0.843686in}}{\pgfqpoint{12.284866in}{0.852512in}}%
\pgfpathquadraticcurveto{\pgfqpoint{12.293349in}{0.861356in}}{\pgfqpoint{12.308101in}{0.861356in}}%
\pgfpathquadraticcurveto{\pgfqpoint{12.312875in}{0.861356in}}{\pgfqpoint{12.317432in}{0.860365in}}%
\pgfpathquadraticcurveto{\pgfqpoint{12.321989in}{0.859392in}}{\pgfqpoint{12.326276in}{0.857429in}}%
\pgfpathlineto{\pgfqpoint{12.326276in}{0.857429in}}%
\pgfpathclose%
\pgfpathmoveto{\pgfqpoint{12.372945in}{0.828483in}}%
\pgfpathquadraticcurveto{\pgfqpoint{12.360391in}{0.828483in}}{\pgfqpoint{12.355545in}{0.825619in}}%
\pgfpathquadraticcurveto{\pgfqpoint{12.350700in}{0.822756in}}{\pgfqpoint{12.350700in}{0.815821in}}%
\pgfpathquadraticcurveto{\pgfqpoint{12.350700in}{0.810309in}}{\pgfqpoint{12.354339in}{0.807067in}}%
\pgfpathquadraticcurveto{\pgfqpoint{12.357977in}{0.803843in}}{\pgfqpoint{12.364209in}{0.803843in}}%
\pgfpathquadraticcurveto{\pgfqpoint{12.372837in}{0.803843in}}{\pgfqpoint{12.378043in}{0.809949in}}%
\pgfpathquadraticcurveto{\pgfqpoint{12.383248in}{0.816055in}}{\pgfqpoint{12.383248in}{0.826178in}}%
\pgfpathlineto{\pgfqpoint{12.383248in}{0.828483in}}%
\pgfpathlineto{\pgfqpoint{12.372945in}{0.828483in}}%
\pgfpathlineto{\pgfqpoint{12.372945in}{0.828483in}}%
\pgfpathclose%
\pgfpathmoveto{\pgfqpoint{12.393605in}{0.832770in}}%
\pgfpathlineto{\pgfqpoint{12.393605in}{0.796800in}}%
\pgfpathlineto{\pgfqpoint{12.383248in}{0.796800in}}%
\pgfpathlineto{\pgfqpoint{12.383248in}{0.806364in}}%
\pgfpathquadraticcurveto{\pgfqpoint{12.379700in}{0.800637in}}{\pgfqpoint{12.374404in}{0.797899in}}%
\pgfpathquadraticcurveto{\pgfqpoint{12.369109in}{0.795161in}}{\pgfqpoint{12.361453in}{0.795161in}}%
\pgfpathquadraticcurveto{\pgfqpoint{12.351781in}{0.795161in}}{\pgfqpoint{12.346053in}{0.800601in}}%
\pgfpathquadraticcurveto{\pgfqpoint{12.340343in}{0.806040in}}{\pgfqpoint{12.340343in}{0.815154in}}%
\pgfpathquadraticcurveto{\pgfqpoint{12.340343in}{0.825782in}}{\pgfqpoint{12.347458in}{0.831185in}}%
\pgfpathquadraticcurveto{\pgfqpoint{12.354591in}{0.836589in}}{\pgfqpoint{12.368712in}{0.836589in}}%
\pgfpathlineto{\pgfqpoint{12.383248in}{0.836589in}}%
\pgfpathlineto{\pgfqpoint{12.383248in}{0.837616in}}%
\pgfpathquadraticcurveto{\pgfqpoint{12.383248in}{0.844766in}}{\pgfqpoint{12.378547in}{0.848675in}}%
\pgfpathquadraticcurveto{\pgfqpoint{12.373846in}{0.852584in}}{\pgfqpoint{12.365344in}{0.852584in}}%
\pgfpathquadraticcurveto{\pgfqpoint{12.359940in}{0.852584in}}{\pgfqpoint{12.354807in}{0.851287in}}%
\pgfpathquadraticcurveto{\pgfqpoint{12.349691in}{0.849990in}}{\pgfqpoint{12.344972in}{0.847396in}}%
\pgfpathlineto{\pgfqpoint{12.344972in}{0.856979in}}%
\pgfpathquadraticcurveto{\pgfqpoint{12.350646in}{0.859176in}}{\pgfqpoint{12.355996in}{0.860257in}}%
\pgfpathquadraticcurveto{\pgfqpoint{12.361345in}{0.861356in}}{\pgfqpoint{12.366407in}{0.861356in}}%
\pgfpathquadraticcurveto{\pgfqpoint{12.380096in}{0.861356in}}{\pgfqpoint{12.386850in}{0.854259in}}%
\pgfpathquadraticcurveto{\pgfqpoint{12.393605in}{0.847180in}}{\pgfqpoint{12.393605in}{0.832770in}}%
\pgfpathlineto{\pgfqpoint{12.393605in}{0.832770in}}%
\pgfpathclose%
\pgfpathmoveto{\pgfqpoint{12.467347in}{0.834860in}}%
\pgfpathlineto{\pgfqpoint{12.467347in}{0.796800in}}%
\pgfpathlineto{\pgfqpoint{12.456990in}{0.796800in}}%
\pgfpathlineto{\pgfqpoint{12.456990in}{0.834517in}}%
\pgfpathquadraticcurveto{\pgfqpoint{12.456990in}{0.843469in}}{\pgfqpoint{12.453495in}{0.847900in}}%
\pgfpathquadraticcurveto{\pgfqpoint{12.450001in}{0.852349in}}{\pgfqpoint{12.443030in}{0.852349in}}%
\pgfpathquadraticcurveto{\pgfqpoint{12.434637in}{0.852349in}}{\pgfqpoint{12.429791in}{0.847000in}}%
\pgfpathquadraticcurveto{\pgfqpoint{12.424946in}{0.841668in}}{\pgfqpoint{12.424946in}{0.832428in}}%
\pgfpathlineto{\pgfqpoint{12.424946in}{0.796800in}}%
\pgfpathlineto{\pgfqpoint{12.414535in}{0.796800in}}%
\pgfpathlineto{\pgfqpoint{12.414535in}{0.859843in}}%
\pgfpathlineto{\pgfqpoint{12.424946in}{0.859843in}}%
\pgfpathlineto{\pgfqpoint{12.424946in}{0.850044in}}%
\pgfpathquadraticcurveto{\pgfqpoint{12.428675in}{0.855736in}}{\pgfqpoint{12.433700in}{0.858546in}}%
\pgfpathquadraticcurveto{\pgfqpoint{12.438743in}{0.861356in}}{\pgfqpoint{12.445336in}{0.861356in}}%
\pgfpathquadraticcurveto{\pgfqpoint{12.456197in}{0.861356in}}{\pgfqpoint{12.461763in}{0.854637in}}%
\pgfpathquadraticcurveto{\pgfqpoint{12.467347in}{0.847918in}}{\pgfqpoint{12.467347in}{0.834860in}}%
\pgfpathlineto{\pgfqpoint{12.467347in}{0.834860in}}%
\pgfpathclose%
\pgfpathmoveto{\pgfqpoint{12.533361in}{0.857429in}}%
\pgfpathlineto{\pgfqpoint{12.533361in}{0.847738in}}%
\pgfpathquadraticcurveto{\pgfqpoint{12.528966in}{0.850170in}}{\pgfqpoint{12.524553in}{0.851377in}}%
\pgfpathquadraticcurveto{\pgfqpoint{12.520140in}{0.852584in}}{\pgfqpoint{12.515637in}{0.852584in}}%
\pgfpathquadraticcurveto{\pgfqpoint{12.505551in}{0.852584in}}{\pgfqpoint{12.499967in}{0.846189in}}%
\pgfpathquadraticcurveto{\pgfqpoint{12.494401in}{0.839813in}}{\pgfqpoint{12.494401in}{0.828267in}}%
\pgfpathquadraticcurveto{\pgfqpoint{12.494401in}{0.816721in}}{\pgfqpoint{12.499967in}{0.810327in}}%
\pgfpathquadraticcurveto{\pgfqpoint{12.505551in}{0.803951in}}{\pgfqpoint{12.515637in}{0.803951in}}%
\pgfpathquadraticcurveto{\pgfqpoint{12.520140in}{0.803951in}}{\pgfqpoint{12.524553in}{0.805158in}}%
\pgfpathquadraticcurveto{\pgfqpoint{12.528966in}{0.806364in}}{\pgfqpoint{12.533361in}{0.808796in}}%
\pgfpathlineto{\pgfqpoint{12.533361in}{0.799214in}}%
\pgfpathquadraticcurveto{\pgfqpoint{12.529020in}{0.797196in}}{\pgfqpoint{12.524373in}{0.796188in}}%
\pgfpathquadraticcurveto{\pgfqpoint{12.519744in}{0.795161in}}{\pgfqpoint{12.514503in}{0.795161in}}%
\pgfpathquadraticcurveto{\pgfqpoint{12.500255in}{0.795161in}}{\pgfqpoint{12.491861in}{0.804113in}}%
\pgfpathquadraticcurveto{\pgfqpoint{12.483486in}{0.813065in}}{\pgfqpoint{12.483486in}{0.828267in}}%
\pgfpathquadraticcurveto{\pgfqpoint{12.483486in}{0.843686in}}{\pgfqpoint{12.491951in}{0.852512in}}%
\pgfpathquadraticcurveto{\pgfqpoint{12.500435in}{0.861356in}}{\pgfqpoint{12.515187in}{0.861356in}}%
\pgfpathquadraticcurveto{\pgfqpoint{12.519960in}{0.861356in}}{\pgfqpoint{12.524517in}{0.860365in}}%
\pgfpathquadraticcurveto{\pgfqpoint{12.529074in}{0.859392in}}{\pgfqpoint{12.533361in}{0.857429in}}%
\pgfpathlineto{\pgfqpoint{12.533361in}{0.857429in}}%
\pgfpathclose%
\pgfpathmoveto{\pgfqpoint{12.605302in}{0.830915in}}%
\pgfpathlineto{\pgfqpoint{12.605302in}{0.825854in}}%
\pgfpathlineto{\pgfqpoint{12.557678in}{0.825854in}}%
\pgfpathquadraticcurveto{\pgfqpoint{12.558362in}{0.815154in}}{\pgfqpoint{12.564126in}{0.809553in}}%
\pgfpathquadraticcurveto{\pgfqpoint{12.569890in}{0.803951in}}{\pgfqpoint{12.580193in}{0.803951in}}%
\pgfpathquadraticcurveto{\pgfqpoint{12.586155in}{0.803951in}}{\pgfqpoint{12.591757in}{0.805410in}}%
\pgfpathquadraticcurveto{\pgfqpoint{12.597359in}{0.806869in}}{\pgfqpoint{12.602888in}{0.809805in}}%
\pgfpathlineto{\pgfqpoint{12.602888in}{0.800006in}}%
\pgfpathquadraticcurveto{\pgfqpoint{12.597304in}{0.797647in}}{\pgfqpoint{12.591451in}{0.796404in}}%
\pgfpathquadraticcurveto{\pgfqpoint{12.585597in}{0.795161in}}{\pgfqpoint{12.579581in}{0.795161in}}%
\pgfpathquadraticcurveto{\pgfqpoint{12.564486in}{0.795161in}}{\pgfqpoint{12.555678in}{0.803933in}}%
\pgfpathquadraticcurveto{\pgfqpoint{12.546870in}{0.812723in}}{\pgfqpoint{12.546870in}{0.827709in}}%
\pgfpathquadraticcurveto{\pgfqpoint{12.546870in}{0.843181in}}{\pgfqpoint{12.555228in}{0.852259in}}%
\pgfpathquadraticcurveto{\pgfqpoint{12.563586in}{0.861356in}}{\pgfqpoint{12.577779in}{0.861356in}}%
\pgfpathquadraticcurveto{\pgfqpoint{12.590496in}{0.861356in}}{\pgfqpoint{12.597899in}{0.853160in}}%
\pgfpathquadraticcurveto{\pgfqpoint{12.605302in}{0.844983in}}{\pgfqpoint{12.605302in}{0.830915in}}%
\pgfpathlineto{\pgfqpoint{12.605302in}{0.830915in}}%
\pgfpathclose%
\pgfpathmoveto{\pgfqpoint{12.594945in}{0.833959in}}%
\pgfpathquadraticcurveto{\pgfqpoint{12.594837in}{0.842443in}}{\pgfqpoint{12.590190in}{0.847504in}}%
\pgfpathquadraticcurveto{\pgfqpoint{12.585543in}{0.852584in}}{\pgfqpoint{12.577887in}{0.852584in}}%
\pgfpathquadraticcurveto{\pgfqpoint{12.569224in}{0.852584in}}{\pgfqpoint{12.564018in}{0.847684in}}%
\pgfpathquadraticcurveto{\pgfqpoint{12.558812in}{0.842785in}}{\pgfqpoint{12.558020in}{0.833887in}}%
\pgfpathlineto{\pgfqpoint{12.594945in}{0.833959in}}%
\pgfpathlineto{\pgfqpoint{12.594945in}{0.833959in}}%
\pgfpathclose%
\pgfusepath{fill}%
\end{pgfscope}%
\begin{pgfscope}%
\pgfsetbuttcap%
\pgfsetroundjoin%
\definecolor{currentfill}{rgb}{0.309804,0.372549,0.419608}%
\pgfsetfillcolor{currentfill}%
\pgfsetlinewidth{0.000000pt}%
\definecolor{currentstroke}{rgb}{0.309804,0.372549,0.419608}%
\pgfsetstrokecolor{currentstroke}%
\pgfsetdash{}{0pt}%
\pgfpathmoveto{\pgfqpoint{12.762595in}{0.859843in}}%
\pgfpathlineto{\pgfqpoint{12.762595in}{0.796800in}}%
\pgfpathlineto{\pgfqpoint{12.752184in}{0.796800in}}%
\pgfpathlineto{\pgfqpoint{12.752184in}{0.851791in}}%
\pgfpathlineto{\pgfqpoint{12.723761in}{0.851791in}}%
\pgfpathlineto{\pgfqpoint{12.723761in}{0.796800in}}%
\pgfpathlineto{\pgfqpoint{12.713350in}{0.796800in}}%
\pgfpathlineto{\pgfqpoint{12.713350in}{0.851791in}}%
\pgfpathlineto{\pgfqpoint{12.703443in}{0.851791in}}%
\pgfpathlineto{\pgfqpoint{12.703443in}{0.859843in}}%
\pgfpathlineto{\pgfqpoint{12.713350in}{0.859843in}}%
\pgfpathlineto{\pgfqpoint{12.713350in}{0.864238in}}%
\pgfpathquadraticcurveto{\pgfqpoint{12.713350in}{0.874540in}}{\pgfqpoint{12.718213in}{0.879458in}}%
\pgfpathquadraticcurveto{\pgfqpoint{12.723094in}{0.884393in}}{\pgfqpoint{12.733163in}{0.884393in}}%
\pgfpathlineto{\pgfqpoint{12.743574in}{0.884393in}}%
\pgfpathlineto{\pgfqpoint{12.743574in}{0.875765in}}%
\pgfpathlineto{\pgfqpoint{12.733667in}{0.875765in}}%
\pgfpathquadraticcurveto{\pgfqpoint{12.728102in}{0.875765in}}{\pgfqpoint{12.725922in}{0.873514in}}%
\pgfpathquadraticcurveto{\pgfqpoint{12.723761in}{0.871262in}}{\pgfqpoint{12.723761in}{0.865408in}}%
\pgfpathlineto{\pgfqpoint{12.723761in}{0.859843in}}%
\pgfpathlineto{\pgfqpoint{12.762595in}{0.859843in}}%
\pgfpathlineto{\pgfqpoint{12.762595in}{0.859843in}}%
\pgfpathclose%
\pgfpathmoveto{\pgfqpoint{12.752184in}{0.884267in}}%
\pgfpathlineto{\pgfqpoint{12.762595in}{0.884267in}}%
\pgfpathlineto{\pgfqpoint{12.762595in}{0.871154in}}%
\pgfpathlineto{\pgfqpoint{12.752184in}{0.871154in}}%
\pgfpathlineto{\pgfqpoint{12.752184in}{0.884267in}}%
\pgfpathlineto{\pgfqpoint{12.752184in}{0.884267in}}%
\pgfpathclose%
\pgfpathmoveto{\pgfqpoint{12.784263in}{0.884393in}}%
\pgfpathlineto{\pgfqpoint{12.794620in}{0.884393in}}%
\pgfpathlineto{\pgfqpoint{12.794620in}{0.796800in}}%
\pgfpathlineto{\pgfqpoint{12.784263in}{0.796800in}}%
\pgfpathlineto{\pgfqpoint{12.784263in}{0.884393in}}%
\pgfpathlineto{\pgfqpoint{12.784263in}{0.884393in}}%
\pgfpathclose%
\pgfpathmoveto{\pgfqpoint{12.816289in}{0.884393in}}%
\pgfpathlineto{\pgfqpoint{12.826646in}{0.884393in}}%
\pgfpathlineto{\pgfqpoint{12.826646in}{0.796800in}}%
\pgfpathlineto{\pgfqpoint{12.816289in}{0.796800in}}%
\pgfpathlineto{\pgfqpoint{12.816289in}{0.884393in}}%
\pgfpathlineto{\pgfqpoint{12.816289in}{0.884393in}}%
\pgfpathclose%
\pgfusepath{fill}%
\end{pgfscope}%
\begin{pgfscope}%
\pgfsetbuttcap%
\pgfsetroundjoin%
\definecolor{currentfill}{rgb}{0.309804,0.372549,0.419608}%
\pgfsetfillcolor{currentfill}%
\pgfsetlinewidth{0.000000pt}%
\definecolor{currentstroke}{rgb}{0.309804,0.372549,0.419608}%
\pgfsetstrokecolor{currentstroke}%
\pgfsetdash{}{0pt}%
\pgfpathmoveto{\pgfqpoint{12.945546in}{0.857987in}}%
\pgfpathlineto{\pgfqpoint{12.945546in}{0.848189in}}%
\pgfpathquadraticcurveto{\pgfqpoint{12.941169in}{0.850440in}}{\pgfqpoint{12.936432in}{0.851557in}}%
\pgfpathquadraticcurveto{\pgfqpoint{12.931713in}{0.852692in}}{\pgfqpoint{12.926634in}{0.852692in}}%
\pgfpathquadraticcurveto{\pgfqpoint{12.918924in}{0.852692in}}{\pgfqpoint{12.915070in}{0.850332in}}%
\pgfpathquadraticcurveto{\pgfqpoint{12.911215in}{0.847973in}}{\pgfqpoint{12.911215in}{0.843235in}}%
\pgfpathquadraticcurveto{\pgfqpoint{12.911215in}{0.839633in}}{\pgfqpoint{12.913971in}{0.837580in}}%
\pgfpathquadraticcurveto{\pgfqpoint{12.916727in}{0.835526in}}{\pgfqpoint{12.925067in}{0.833671in}}%
\pgfpathlineto{\pgfqpoint{12.928615in}{0.832878in}}%
\pgfpathquadraticcurveto{\pgfqpoint{12.939638in}{0.830519in}}{\pgfqpoint{12.944286in}{0.826214in}}%
\pgfpathquadraticcurveto{\pgfqpoint{12.948933in}{0.821909in}}{\pgfqpoint{12.948933in}{0.814200in}}%
\pgfpathquadraticcurveto{\pgfqpoint{12.948933in}{0.805410in}}{\pgfqpoint{12.941980in}{0.800276in}}%
\pgfpathquadraticcurveto{\pgfqpoint{12.935027in}{0.795161in}}{\pgfqpoint{12.922869in}{0.795161in}}%
\pgfpathquadraticcurveto{\pgfqpoint{12.917808in}{0.795161in}}{\pgfqpoint{12.912314in}{0.796152in}}%
\pgfpathquadraticcurveto{\pgfqpoint{12.906820in}{0.797142in}}{\pgfqpoint{12.900750in}{0.799106in}}%
\pgfpathlineto{\pgfqpoint{12.900750in}{0.809805in}}%
\pgfpathquadraticcurveto{\pgfqpoint{12.906496in}{0.806815in}}{\pgfqpoint{12.912062in}{0.805320in}}%
\pgfpathquadraticcurveto{\pgfqpoint{12.917628in}{0.803843in}}{\pgfqpoint{12.923103in}{0.803843in}}%
\pgfpathquadraticcurveto{\pgfqpoint{12.930416in}{0.803843in}}{\pgfqpoint{12.934343in}{0.806346in}}%
\pgfpathquadraticcurveto{\pgfqpoint{12.938288in}{0.808850in}}{\pgfqpoint{12.938288in}{0.813407in}}%
\pgfpathquadraticcurveto{\pgfqpoint{12.938288in}{0.817622in}}{\pgfqpoint{12.935442in}{0.819874in}}%
\pgfpathquadraticcurveto{\pgfqpoint{12.932614in}{0.822125in}}{\pgfqpoint{12.922977in}{0.824214in}}%
\pgfpathlineto{\pgfqpoint{12.919375in}{0.825061in}}%
\pgfpathquadraticcurveto{\pgfqpoint{12.909756in}{0.827078in}}{\pgfqpoint{12.905469in}{0.831275in}}%
\pgfpathquadraticcurveto{\pgfqpoint{12.901201in}{0.835472in}}{\pgfqpoint{12.901201in}{0.842785in}}%
\pgfpathquadraticcurveto{\pgfqpoint{12.901201in}{0.851683in}}{\pgfqpoint{12.907505in}{0.856510in}}%
\pgfpathquadraticcurveto{\pgfqpoint{12.913809in}{0.861356in}}{\pgfqpoint{12.925409in}{0.861356in}}%
\pgfpathquadraticcurveto{\pgfqpoint{12.931137in}{0.861356in}}{\pgfqpoint{12.936198in}{0.860509in}}%
\pgfpathquadraticcurveto{\pgfqpoint{12.941278in}{0.859680in}}{\pgfqpoint{12.945546in}{0.857987in}}%
\pgfpathlineto{\pgfqpoint{12.945546in}{0.857987in}}%
\pgfpathclose%
\pgfpathmoveto{\pgfqpoint{12.964351in}{0.821675in}}%
\pgfpathlineto{\pgfqpoint{12.964351in}{0.859843in}}%
\pgfpathlineto{\pgfqpoint{12.974708in}{0.859843in}}%
\pgfpathlineto{\pgfqpoint{12.974708in}{0.822071in}}%
\pgfpathquadraticcurveto{\pgfqpoint{12.974708in}{0.813119in}}{\pgfqpoint{12.978184in}{0.808634in}}%
\pgfpathquadraticcurveto{\pgfqpoint{12.981679in}{0.804167in}}{\pgfqpoint{12.988668in}{0.804167in}}%
\pgfpathquadraticcurveto{\pgfqpoint{12.997043in}{0.804167in}}{\pgfqpoint{13.001906in}{0.809517in}}%
\pgfpathquadraticcurveto{\pgfqpoint{13.006788in}{0.814866in}}{\pgfqpoint{13.006788in}{0.824106in}}%
\pgfpathlineto{\pgfqpoint{13.006788in}{0.859843in}}%
\pgfpathlineto{\pgfqpoint{13.017145in}{0.859843in}}%
\pgfpathlineto{\pgfqpoint{13.017145in}{0.796800in}}%
\pgfpathlineto{\pgfqpoint{13.006788in}{0.796800in}}%
\pgfpathlineto{\pgfqpoint{13.006788in}{0.806491in}}%
\pgfpathquadraticcurveto{\pgfqpoint{13.003023in}{0.800745in}}{\pgfqpoint{12.998034in}{0.797953in}}%
\pgfpathquadraticcurveto{\pgfqpoint{12.993062in}{0.795161in}}{\pgfqpoint{12.986470in}{0.795161in}}%
\pgfpathquadraticcurveto{\pgfqpoint{12.975609in}{0.795161in}}{\pgfqpoint{12.969971in}{0.801915in}}%
\pgfpathquadraticcurveto{\pgfqpoint{12.964351in}{0.808670in}}{\pgfqpoint{12.964351in}{0.821675in}}%
\pgfpathlineto{\pgfqpoint{12.964351in}{0.821675in}}%
\pgfpathclose%
\pgfpathmoveto{\pgfqpoint{12.990415in}{0.861356in}}%
\pgfpathlineto{\pgfqpoint{12.990415in}{0.861356in}}%
\pgfpathlineto{\pgfqpoint{12.990415in}{0.861356in}}%
\pgfpathclose%
\pgfpathmoveto{\pgfqpoint{13.048486in}{0.806256in}}%
\pgfpathlineto{\pgfqpoint{13.048486in}{0.772826in}}%
\pgfpathlineto{\pgfqpoint{13.038075in}{0.772826in}}%
\pgfpathlineto{\pgfqpoint{13.038075in}{0.859843in}}%
\pgfpathlineto{\pgfqpoint{13.048486in}{0.859843in}}%
\pgfpathlineto{\pgfqpoint{13.048486in}{0.850278in}}%
\pgfpathquadraticcurveto{\pgfqpoint{13.051764in}{0.855898in}}{\pgfqpoint{13.056735in}{0.858618in}}%
\pgfpathquadraticcurveto{\pgfqpoint{13.061725in}{0.861356in}}{\pgfqpoint{13.068641in}{0.861356in}}%
\pgfpathquadraticcurveto{\pgfqpoint{13.080133in}{0.861356in}}{\pgfqpoint{13.087302in}{0.852241in}}%
\pgfpathquadraticcurveto{\pgfqpoint{13.094489in}{0.843127in}}{\pgfqpoint{13.094489in}{0.828267in}}%
\pgfpathquadraticcurveto{\pgfqpoint{13.094489in}{0.813407in}}{\pgfqpoint{13.087302in}{0.804275in}}%
\pgfpathquadraticcurveto{\pgfqpoint{13.080133in}{0.795161in}}{\pgfqpoint{13.068641in}{0.795161in}}%
\pgfpathquadraticcurveto{\pgfqpoint{13.061725in}{0.795161in}}{\pgfqpoint{13.056735in}{0.797899in}}%
\pgfpathquadraticcurveto{\pgfqpoint{13.051764in}{0.800637in}}{\pgfqpoint{13.048486in}{0.806256in}}%
\pgfpathlineto{\pgfqpoint{13.048486in}{0.806256in}}%
\pgfpathclose%
\pgfpathmoveto{\pgfqpoint{13.083736in}{0.828267in}}%
\pgfpathquadraticcurveto{\pgfqpoint{13.083736in}{0.839687in}}{\pgfqpoint{13.079035in}{0.846189in}}%
\pgfpathquadraticcurveto{\pgfqpoint{13.074333in}{0.852692in}}{\pgfqpoint{13.066120in}{0.852692in}}%
\pgfpathquadraticcurveto{\pgfqpoint{13.057888in}{0.852692in}}{\pgfqpoint{13.053187in}{0.846189in}}%
\pgfpathquadraticcurveto{\pgfqpoint{13.048486in}{0.839687in}}{\pgfqpoint{13.048486in}{0.828267in}}%
\pgfpathquadraticcurveto{\pgfqpoint{13.048486in}{0.816848in}}{\pgfqpoint{13.053187in}{0.810345in}}%
\pgfpathquadraticcurveto{\pgfqpoint{13.057888in}{0.803843in}}{\pgfqpoint{13.066120in}{0.803843in}}%
\pgfpathquadraticcurveto{\pgfqpoint{13.074333in}{0.803843in}}{\pgfqpoint{13.079035in}{0.810345in}}%
\pgfpathquadraticcurveto{\pgfqpoint{13.083736in}{0.816848in}}{\pgfqpoint{13.083736in}{0.828267in}}%
\pgfpathlineto{\pgfqpoint{13.083736in}{0.828267in}}%
\pgfpathclose%
\pgfpathmoveto{\pgfqpoint{13.121669in}{0.806256in}}%
\pgfpathlineto{\pgfqpoint{13.121669in}{0.772826in}}%
\pgfpathlineto{\pgfqpoint{13.111258in}{0.772826in}}%
\pgfpathlineto{\pgfqpoint{13.111258in}{0.859843in}}%
\pgfpathlineto{\pgfqpoint{13.121669in}{0.859843in}}%
\pgfpathlineto{\pgfqpoint{13.121669in}{0.850278in}}%
\pgfpathquadraticcurveto{\pgfqpoint{13.124947in}{0.855898in}}{\pgfqpoint{13.129919in}{0.858618in}}%
\pgfpathquadraticcurveto{\pgfqpoint{13.134908in}{0.861356in}}{\pgfqpoint{13.141825in}{0.861356in}}%
\pgfpathquadraticcurveto{\pgfqpoint{13.153317in}{0.861356in}}{\pgfqpoint{13.160485in}{0.852241in}}%
\pgfpathquadraticcurveto{\pgfqpoint{13.167672in}{0.843127in}}{\pgfqpoint{13.167672in}{0.828267in}}%
\pgfpathquadraticcurveto{\pgfqpoint{13.167672in}{0.813407in}}{\pgfqpoint{13.160485in}{0.804275in}}%
\pgfpathquadraticcurveto{\pgfqpoint{13.153317in}{0.795161in}}{\pgfqpoint{13.141825in}{0.795161in}}%
\pgfpathquadraticcurveto{\pgfqpoint{13.134908in}{0.795161in}}{\pgfqpoint{13.129919in}{0.797899in}}%
\pgfpathquadraticcurveto{\pgfqpoint{13.124947in}{0.800637in}}{\pgfqpoint{13.121669in}{0.806256in}}%
\pgfpathlineto{\pgfqpoint{13.121669in}{0.806256in}}%
\pgfpathclose%
\pgfpathmoveto{\pgfqpoint{13.156919in}{0.828267in}}%
\pgfpathquadraticcurveto{\pgfqpoint{13.156919in}{0.839687in}}{\pgfqpoint{13.152218in}{0.846189in}}%
\pgfpathquadraticcurveto{\pgfqpoint{13.147517in}{0.852692in}}{\pgfqpoint{13.139303in}{0.852692in}}%
\pgfpathquadraticcurveto{\pgfqpoint{13.131072in}{0.852692in}}{\pgfqpoint{13.126370in}{0.846189in}}%
\pgfpathquadraticcurveto{\pgfqpoint{13.121669in}{0.839687in}}{\pgfqpoint{13.121669in}{0.828267in}}%
\pgfpathquadraticcurveto{\pgfqpoint{13.121669in}{0.816848in}}{\pgfqpoint{13.126370in}{0.810345in}}%
\pgfpathquadraticcurveto{\pgfqpoint{13.131072in}{0.803843in}}{\pgfqpoint{13.139303in}{0.803843in}}%
\pgfpathquadraticcurveto{\pgfqpoint{13.147517in}{0.803843in}}{\pgfqpoint{13.152218in}{0.810345in}}%
\pgfpathquadraticcurveto{\pgfqpoint{13.156919in}{0.816848in}}{\pgfqpoint{13.156919in}{0.828267in}}%
\pgfpathlineto{\pgfqpoint{13.156919in}{0.828267in}}%
\pgfpathclose%
\pgfpathmoveto{\pgfqpoint{13.221367in}{0.850170in}}%
\pgfpathquadraticcurveto{\pgfqpoint{13.219619in}{0.851179in}}{\pgfqpoint{13.217566in}{0.851647in}}%
\pgfpathquadraticcurveto{\pgfqpoint{13.215513in}{0.852133in}}{\pgfqpoint{13.213045in}{0.852133in}}%
\pgfpathquadraticcurveto{\pgfqpoint{13.204255in}{0.852133in}}{\pgfqpoint{13.199554in}{0.846423in}}%
\pgfpathquadraticcurveto{\pgfqpoint{13.194853in}{0.840714in}}{\pgfqpoint{13.194853in}{0.830014in}}%
\pgfpathlineto{\pgfqpoint{13.194853in}{0.796800in}}%
\pgfpathlineto{\pgfqpoint{13.184442in}{0.796800in}}%
\pgfpathlineto{\pgfqpoint{13.184442in}{0.859843in}}%
\pgfpathlineto{\pgfqpoint{13.194853in}{0.859843in}}%
\pgfpathlineto{\pgfqpoint{13.194853in}{0.850044in}}%
\pgfpathquadraticcurveto{\pgfqpoint{13.198131in}{0.855790in}}{\pgfqpoint{13.203354in}{0.858564in}}%
\pgfpathquadraticcurveto{\pgfqpoint{13.208596in}{0.861356in}}{\pgfqpoint{13.216089in}{0.861356in}}%
\pgfpathquadraticcurveto{\pgfqpoint{13.217152in}{0.861356in}}{\pgfqpoint{13.218449in}{0.861211in}}%
\pgfpathquadraticcurveto{\pgfqpoint{13.219745in}{0.861085in}}{\pgfqpoint{13.221312in}{0.860797in}}%
\pgfpathlineto{\pgfqpoint{13.221367in}{0.850170in}}%
\pgfpathlineto{\pgfqpoint{13.221367in}{0.850170in}}%
\pgfpathclose%
\pgfpathmoveto{\pgfqpoint{13.283617in}{0.830915in}}%
\pgfpathlineto{\pgfqpoint{13.283617in}{0.825854in}}%
\pgfpathlineto{\pgfqpoint{13.235992in}{0.825854in}}%
\pgfpathquadraticcurveto{\pgfqpoint{13.236677in}{0.815154in}}{\pgfqpoint{13.242441in}{0.809553in}}%
\pgfpathquadraticcurveto{\pgfqpoint{13.248205in}{0.803951in}}{\pgfqpoint{13.258508in}{0.803951in}}%
\pgfpathquadraticcurveto{\pgfqpoint{13.264470in}{0.803951in}}{\pgfqpoint{13.270071in}{0.805410in}}%
\pgfpathquadraticcurveto{\pgfqpoint{13.275673in}{0.806869in}}{\pgfqpoint{13.281203in}{0.809805in}}%
\pgfpathlineto{\pgfqpoint{13.281203in}{0.800006in}}%
\pgfpathquadraticcurveto{\pgfqpoint{13.275619in}{0.797647in}}{\pgfqpoint{13.269765in}{0.796404in}}%
\pgfpathquadraticcurveto{\pgfqpoint{13.263911in}{0.795161in}}{\pgfqpoint{13.257895in}{0.795161in}}%
\pgfpathquadraticcurveto{\pgfqpoint{13.242801in}{0.795161in}}{\pgfqpoint{13.233993in}{0.803933in}}%
\pgfpathquadraticcurveto{\pgfqpoint{13.225185in}{0.812723in}}{\pgfqpoint{13.225185in}{0.827709in}}%
\pgfpathquadraticcurveto{\pgfqpoint{13.225185in}{0.843181in}}{\pgfqpoint{13.233543in}{0.852259in}}%
\pgfpathquadraticcurveto{\pgfqpoint{13.241900in}{0.861356in}}{\pgfqpoint{13.256094in}{0.861356in}}%
\pgfpathquadraticcurveto{\pgfqpoint{13.268811in}{0.861356in}}{\pgfqpoint{13.276214in}{0.853160in}}%
\pgfpathquadraticcurveto{\pgfqpoint{13.283617in}{0.844983in}}{\pgfqpoint{13.283617in}{0.830915in}}%
\pgfpathlineto{\pgfqpoint{13.283617in}{0.830915in}}%
\pgfpathclose%
\pgfpathmoveto{\pgfqpoint{13.273260in}{0.833959in}}%
\pgfpathquadraticcurveto{\pgfqpoint{13.273151in}{0.842443in}}{\pgfqpoint{13.268504in}{0.847504in}}%
\pgfpathquadraticcurveto{\pgfqpoint{13.263857in}{0.852584in}}{\pgfqpoint{13.256202in}{0.852584in}}%
\pgfpathquadraticcurveto{\pgfqpoint{13.247538in}{0.852584in}}{\pgfqpoint{13.242333in}{0.847684in}}%
\pgfpathquadraticcurveto{\pgfqpoint{13.237127in}{0.842785in}}{\pgfqpoint{13.236335in}{0.833887in}}%
\pgfpathlineto{\pgfqpoint{13.273260in}{0.833959in}}%
\pgfpathlineto{\pgfqpoint{13.273260in}{0.833959in}}%
\pgfpathclose%
\pgfpathmoveto{\pgfqpoint{13.340805in}{0.857987in}}%
\pgfpathlineto{\pgfqpoint{13.340805in}{0.848189in}}%
\pgfpathquadraticcurveto{\pgfqpoint{13.336428in}{0.850440in}}{\pgfqpoint{13.331691in}{0.851557in}}%
\pgfpathquadraticcurveto{\pgfqpoint{13.326972in}{0.852692in}}{\pgfqpoint{13.321892in}{0.852692in}}%
\pgfpathquadraticcurveto{\pgfqpoint{13.314183in}{0.852692in}}{\pgfqpoint{13.310329in}{0.850332in}}%
\pgfpathquadraticcurveto{\pgfqpoint{13.306474in}{0.847973in}}{\pgfqpoint{13.306474in}{0.843235in}}%
\pgfpathquadraticcurveto{\pgfqpoint{13.306474in}{0.839633in}}{\pgfqpoint{13.309230in}{0.837580in}}%
\pgfpathquadraticcurveto{\pgfqpoint{13.311986in}{0.835526in}}{\pgfqpoint{13.320325in}{0.833671in}}%
\pgfpathlineto{\pgfqpoint{13.323874in}{0.832878in}}%
\pgfpathquadraticcurveto{\pgfqpoint{13.334897in}{0.830519in}}{\pgfqpoint{13.339544in}{0.826214in}}%
\pgfpathquadraticcurveto{\pgfqpoint{13.344191in}{0.821909in}}{\pgfqpoint{13.344191in}{0.814200in}}%
\pgfpathquadraticcurveto{\pgfqpoint{13.344191in}{0.805410in}}{\pgfqpoint{13.337239in}{0.800276in}}%
\pgfpathquadraticcurveto{\pgfqpoint{13.330286in}{0.795161in}}{\pgfqpoint{13.318128in}{0.795161in}}%
\pgfpathquadraticcurveto{\pgfqpoint{13.313066in}{0.795161in}}{\pgfqpoint{13.307573in}{0.796152in}}%
\pgfpathquadraticcurveto{\pgfqpoint{13.302079in}{0.797142in}}{\pgfqpoint{13.296009in}{0.799106in}}%
\pgfpathlineto{\pgfqpoint{13.296009in}{0.809805in}}%
\pgfpathquadraticcurveto{\pgfqpoint{13.301755in}{0.806815in}}{\pgfqpoint{13.307321in}{0.805320in}}%
\pgfpathquadraticcurveto{\pgfqpoint{13.312886in}{0.803843in}}{\pgfqpoint{13.318362in}{0.803843in}}%
\pgfpathquadraticcurveto{\pgfqpoint{13.325675in}{0.803843in}}{\pgfqpoint{13.329602in}{0.806346in}}%
\pgfpathquadraticcurveto{\pgfqpoint{13.333546in}{0.808850in}}{\pgfqpoint{13.333546in}{0.813407in}}%
\pgfpathquadraticcurveto{\pgfqpoint{13.333546in}{0.817622in}}{\pgfqpoint{13.330700in}{0.819874in}}%
\pgfpathquadraticcurveto{\pgfqpoint{13.327872in}{0.822125in}}{\pgfqpoint{13.318236in}{0.824214in}}%
\pgfpathlineto{\pgfqpoint{13.314633in}{0.825061in}}%
\pgfpathquadraticcurveto{\pgfqpoint{13.305015in}{0.827078in}}{\pgfqpoint{13.300728in}{0.831275in}}%
\pgfpathquadraticcurveto{\pgfqpoint{13.296459in}{0.835472in}}{\pgfqpoint{13.296459in}{0.842785in}}%
\pgfpathquadraticcurveto{\pgfqpoint{13.296459in}{0.851683in}}{\pgfqpoint{13.302763in}{0.856510in}}%
\pgfpathquadraticcurveto{\pgfqpoint{13.309068in}{0.861356in}}{\pgfqpoint{13.320668in}{0.861356in}}%
\pgfpathquadraticcurveto{\pgfqpoint{13.326395in}{0.861356in}}{\pgfqpoint{13.331457in}{0.860509in}}%
\pgfpathquadraticcurveto{\pgfqpoint{13.336536in}{0.859680in}}{\pgfqpoint{13.340805in}{0.857987in}}%
\pgfpathlineto{\pgfqpoint{13.340805in}{0.857987in}}%
\pgfpathclose%
\pgfpathmoveto{\pgfqpoint{13.400858in}{0.857987in}}%
\pgfpathlineto{\pgfqpoint{13.400858in}{0.848189in}}%
\pgfpathquadraticcurveto{\pgfqpoint{13.396481in}{0.850440in}}{\pgfqpoint{13.391743in}{0.851557in}}%
\pgfpathquadraticcurveto{\pgfqpoint{13.387024in}{0.852692in}}{\pgfqpoint{13.381945in}{0.852692in}}%
\pgfpathquadraticcurveto{\pgfqpoint{13.374236in}{0.852692in}}{\pgfqpoint{13.370381in}{0.850332in}}%
\pgfpathquadraticcurveto{\pgfqpoint{13.366526in}{0.847973in}}{\pgfqpoint{13.366526in}{0.843235in}}%
\pgfpathquadraticcurveto{\pgfqpoint{13.366526in}{0.839633in}}{\pgfqpoint{13.369282in}{0.837580in}}%
\pgfpathquadraticcurveto{\pgfqpoint{13.372038in}{0.835526in}}{\pgfqpoint{13.380378in}{0.833671in}}%
\pgfpathlineto{\pgfqpoint{13.383926in}{0.832878in}}%
\pgfpathquadraticcurveto{\pgfqpoint{13.394950in}{0.830519in}}{\pgfqpoint{13.399597in}{0.826214in}}%
\pgfpathquadraticcurveto{\pgfqpoint{13.404244in}{0.821909in}}{\pgfqpoint{13.404244in}{0.814200in}}%
\pgfpathquadraticcurveto{\pgfqpoint{13.404244in}{0.805410in}}{\pgfqpoint{13.397291in}{0.800276in}}%
\pgfpathquadraticcurveto{\pgfqpoint{13.390339in}{0.795161in}}{\pgfqpoint{13.378180in}{0.795161in}}%
\pgfpathquadraticcurveto{\pgfqpoint{13.373119in}{0.795161in}}{\pgfqpoint{13.367625in}{0.796152in}}%
\pgfpathquadraticcurveto{\pgfqpoint{13.362132in}{0.797142in}}{\pgfqpoint{13.356061in}{0.799106in}}%
\pgfpathlineto{\pgfqpoint{13.356061in}{0.809805in}}%
\pgfpathquadraticcurveto{\pgfqpoint{13.361807in}{0.806815in}}{\pgfqpoint{13.367373in}{0.805320in}}%
\pgfpathquadraticcurveto{\pgfqpoint{13.372939in}{0.803843in}}{\pgfqpoint{13.378414in}{0.803843in}}%
\pgfpathquadraticcurveto{\pgfqpoint{13.385727in}{0.803843in}}{\pgfqpoint{13.389654in}{0.806346in}}%
\pgfpathquadraticcurveto{\pgfqpoint{13.393599in}{0.808850in}}{\pgfqpoint{13.393599in}{0.813407in}}%
\pgfpathquadraticcurveto{\pgfqpoint{13.393599in}{0.817622in}}{\pgfqpoint{13.390753in}{0.819874in}}%
\pgfpathquadraticcurveto{\pgfqpoint{13.387925in}{0.822125in}}{\pgfqpoint{13.378288in}{0.824214in}}%
\pgfpathlineto{\pgfqpoint{13.374686in}{0.825061in}}%
\pgfpathquadraticcurveto{\pgfqpoint{13.365067in}{0.827078in}}{\pgfqpoint{13.360781in}{0.831275in}}%
\pgfpathquadraticcurveto{\pgfqpoint{13.356512in}{0.835472in}}{\pgfqpoint{13.356512in}{0.842785in}}%
\pgfpathquadraticcurveto{\pgfqpoint{13.356512in}{0.851683in}}{\pgfqpoint{13.362816in}{0.856510in}}%
\pgfpathquadraticcurveto{\pgfqpoint{13.369120in}{0.861356in}}{\pgfqpoint{13.380720in}{0.861356in}}%
\pgfpathquadraticcurveto{\pgfqpoint{13.386448in}{0.861356in}}{\pgfqpoint{13.391509in}{0.860509in}}%
\pgfpathquadraticcurveto{\pgfqpoint{13.396589in}{0.859680in}}{\pgfqpoint{13.400858in}{0.857987in}}%
\pgfpathlineto{\pgfqpoint{13.400858in}{0.857987in}}%
\pgfpathclose%
\pgfpathmoveto{\pgfqpoint{13.474653in}{0.830915in}}%
\pgfpathlineto{\pgfqpoint{13.474653in}{0.825854in}}%
\pgfpathlineto{\pgfqpoint{13.427029in}{0.825854in}}%
\pgfpathquadraticcurveto{\pgfqpoint{13.427714in}{0.815154in}}{\pgfqpoint{13.433478in}{0.809553in}}%
\pgfpathquadraticcurveto{\pgfqpoint{13.439242in}{0.803951in}}{\pgfqpoint{13.449544in}{0.803951in}}%
\pgfpathquadraticcurveto{\pgfqpoint{13.455507in}{0.803951in}}{\pgfqpoint{13.461108in}{0.805410in}}%
\pgfpathquadraticcurveto{\pgfqpoint{13.466710in}{0.806869in}}{\pgfqpoint{13.472240in}{0.809805in}}%
\pgfpathlineto{\pgfqpoint{13.472240in}{0.800006in}}%
\pgfpathquadraticcurveto{\pgfqpoint{13.466656in}{0.797647in}}{\pgfqpoint{13.460802in}{0.796404in}}%
\pgfpathquadraticcurveto{\pgfqpoint{13.454948in}{0.795161in}}{\pgfqpoint{13.448932in}{0.795161in}}%
\pgfpathquadraticcurveto{\pgfqpoint{13.433838in}{0.795161in}}{\pgfqpoint{13.425030in}{0.803933in}}%
\pgfpathquadraticcurveto{\pgfqpoint{13.416222in}{0.812723in}}{\pgfqpoint{13.416222in}{0.827709in}}%
\pgfpathquadraticcurveto{\pgfqpoint{13.416222in}{0.843181in}}{\pgfqpoint{13.424580in}{0.852259in}}%
\pgfpathquadraticcurveto{\pgfqpoint{13.432937in}{0.861356in}}{\pgfqpoint{13.447131in}{0.861356in}}%
\pgfpathquadraticcurveto{\pgfqpoint{13.459847in}{0.861356in}}{\pgfqpoint{13.467250in}{0.853160in}}%
\pgfpathquadraticcurveto{\pgfqpoint{13.474653in}{0.844983in}}{\pgfqpoint{13.474653in}{0.830915in}}%
\pgfpathlineto{\pgfqpoint{13.474653in}{0.830915in}}%
\pgfpathclose%
\pgfpathmoveto{\pgfqpoint{13.464296in}{0.833959in}}%
\pgfpathquadraticcurveto{\pgfqpoint{13.464188in}{0.842443in}}{\pgfqpoint{13.459541in}{0.847504in}}%
\pgfpathquadraticcurveto{\pgfqpoint{13.454894in}{0.852584in}}{\pgfqpoint{13.447239in}{0.852584in}}%
\pgfpathquadraticcurveto{\pgfqpoint{13.438575in}{0.852584in}}{\pgfqpoint{13.433370in}{0.847684in}}%
\pgfpathquadraticcurveto{\pgfqpoint{13.428164in}{0.842785in}}{\pgfqpoint{13.427372in}{0.833887in}}%
\pgfpathlineto{\pgfqpoint{13.464296in}{0.833959in}}%
\pgfpathlineto{\pgfqpoint{13.464296in}{0.833959in}}%
\pgfpathclose%
\pgfpathmoveto{\pgfqpoint{13.533139in}{0.850278in}}%
\pgfpathlineto{\pgfqpoint{13.533139in}{0.884393in}}%
\pgfpathlineto{\pgfqpoint{13.543496in}{0.884393in}}%
\pgfpathlineto{\pgfqpoint{13.543496in}{0.796800in}}%
\pgfpathlineto{\pgfqpoint{13.533139in}{0.796800in}}%
\pgfpathlineto{\pgfqpoint{13.533139in}{0.806256in}}%
\pgfpathquadraticcurveto{\pgfqpoint{13.529879in}{0.800637in}}{\pgfqpoint{13.524889in}{0.797899in}}%
\pgfpathquadraticcurveto{\pgfqpoint{13.519918in}{0.795161in}}{\pgfqpoint{13.512929in}{0.795161in}}%
\pgfpathquadraticcurveto{\pgfqpoint{13.501510in}{0.795161in}}{\pgfqpoint{13.494323in}{0.804275in}}%
\pgfpathquadraticcurveto{\pgfqpoint{13.487154in}{0.813407in}}{\pgfqpoint{13.487154in}{0.828267in}}%
\pgfpathquadraticcurveto{\pgfqpoint{13.487154in}{0.843127in}}{\pgfqpoint{13.494323in}{0.852241in}}%
\pgfpathquadraticcurveto{\pgfqpoint{13.501510in}{0.861356in}}{\pgfqpoint{13.512929in}{0.861356in}}%
\pgfpathquadraticcurveto{\pgfqpoint{13.519918in}{0.861356in}}{\pgfqpoint{13.524889in}{0.858618in}}%
\pgfpathquadraticcurveto{\pgfqpoint{13.529879in}{0.855898in}}{\pgfqpoint{13.533139in}{0.850278in}}%
\pgfpathlineto{\pgfqpoint{13.533139in}{0.850278in}}%
\pgfpathclose%
\pgfpathmoveto{\pgfqpoint{13.497853in}{0.828267in}}%
\pgfpathquadraticcurveto{\pgfqpoint{13.497853in}{0.816848in}}{\pgfqpoint{13.502554in}{0.810345in}}%
\pgfpathquadraticcurveto{\pgfqpoint{13.507255in}{0.803843in}}{\pgfqpoint{13.515469in}{0.803843in}}%
\pgfpathquadraticcurveto{\pgfqpoint{13.523683in}{0.803843in}}{\pgfqpoint{13.528402in}{0.810345in}}%
\pgfpathquadraticcurveto{\pgfqpoint{13.533139in}{0.816848in}}{\pgfqpoint{13.533139in}{0.828267in}}%
\pgfpathquadraticcurveto{\pgfqpoint{13.533139in}{0.839687in}}{\pgfqpoint{13.528402in}{0.846189in}}%
\pgfpathquadraticcurveto{\pgfqpoint{13.523683in}{0.852692in}}{\pgfqpoint{13.515469in}{0.852692in}}%
\pgfpathquadraticcurveto{\pgfqpoint{13.507255in}{0.852692in}}{\pgfqpoint{13.502554in}{0.846189in}}%
\pgfpathquadraticcurveto{\pgfqpoint{13.497853in}{0.839687in}}{\pgfqpoint{13.497853in}{0.828267in}}%
\pgfpathlineto{\pgfqpoint{13.497853in}{0.828267in}}%
\pgfpathclose%
\pgfpathmoveto{\pgfqpoint{13.566299in}{0.811102in}}%
\pgfpathlineto{\pgfqpoint{13.578187in}{0.811102in}}%
\pgfpathlineto{\pgfqpoint{13.578187in}{0.796800in}}%
\pgfpathlineto{\pgfqpoint{13.566299in}{0.796800in}}%
\pgfpathlineto{\pgfqpoint{13.566299in}{0.811102in}}%
\pgfpathlineto{\pgfqpoint{13.566299in}{0.811102in}}%
\pgfpathclose%
\pgfusepath{fill}%
\end{pgfscope}%
\begin{pgfscope}%
\pgfsetbuttcap%
\pgfsetroundjoin%
\definecolor{currentfill}{rgb}{0.309804,0.372549,0.419608}%
\pgfsetfillcolor{currentfill}%
\pgfsetlinewidth{0.000000pt}%
\definecolor{currentstroke}{rgb}{0.309804,0.372549,0.419608}%
\pgfsetstrokecolor{currentstroke}%
\pgfsetdash{}{0pt}%
\pgfpathmoveto{\pgfqpoint{0.935312in}{0.717645in}}%
\pgfpathlineto{\pgfqpoint{0.946677in}{0.717645in}}%
\pgfpathlineto{\pgfqpoint{0.946677in}{0.643164in}}%
\pgfpathlineto{\pgfqpoint{0.987601in}{0.643164in}}%
\pgfpathlineto{\pgfqpoint{0.987601in}{0.633600in}}%
\pgfpathlineto{\pgfqpoint{0.935312in}{0.633600in}}%
\pgfpathlineto{\pgfqpoint{0.935312in}{0.717645in}}%
\pgfpathlineto{\pgfqpoint{0.935312in}{0.717645in}}%
\pgfpathclose%
\pgfpathmoveto{\pgfqpoint{1.051004in}{0.667715in}}%
\pgfpathlineto{\pgfqpoint{1.051004in}{0.662654in}}%
\pgfpathlineto{\pgfqpoint{1.003380in}{0.662654in}}%
\pgfpathquadraticcurveto{\pgfqpoint{1.004064in}{0.651954in}}{\pgfqpoint{1.009828in}{0.646353in}}%
\pgfpathquadraticcurveto{\pgfqpoint{1.015592in}{0.640751in}}{\pgfqpoint{1.025895in}{0.640751in}}%
\pgfpathquadraticcurveto{\pgfqpoint{1.031857in}{0.640751in}}{\pgfqpoint{1.037459in}{0.642210in}}%
\pgfpathquadraticcurveto{\pgfqpoint{1.043060in}{0.643669in}}{\pgfqpoint{1.048590in}{0.646605in}}%
\pgfpathlineto{\pgfqpoint{1.048590in}{0.636806in}}%
\pgfpathquadraticcurveto{\pgfqpoint{1.043006in}{0.634447in}}{\pgfqpoint{1.037152in}{0.633204in}}%
\pgfpathquadraticcurveto{\pgfqpoint{1.031298in}{0.631961in}}{\pgfqpoint{1.025282in}{0.631961in}}%
\pgfpathquadraticcurveto{\pgfqpoint{1.010188in}{0.631961in}}{\pgfqpoint{1.001380in}{0.640733in}}%
\pgfpathquadraticcurveto{\pgfqpoint{0.992572in}{0.649523in}}{\pgfqpoint{0.992572in}{0.664509in}}%
\pgfpathquadraticcurveto{\pgfqpoint{0.992572in}{0.679981in}}{\pgfqpoint{1.000930in}{0.689059in}}%
\pgfpathquadraticcurveto{\pgfqpoint{1.009288in}{0.698156in}}{\pgfqpoint{1.023481in}{0.698156in}}%
\pgfpathquadraticcurveto{\pgfqpoint{1.036198in}{0.698156in}}{\pgfqpoint{1.043601in}{0.689960in}}%
\pgfpathquadraticcurveto{\pgfqpoint{1.051004in}{0.681783in}}{\pgfqpoint{1.051004in}{0.667715in}}%
\pgfpathlineto{\pgfqpoint{1.051004in}{0.667715in}}%
\pgfpathclose%
\pgfpathmoveto{\pgfqpoint{1.040647in}{0.670759in}}%
\pgfpathquadraticcurveto{\pgfqpoint{1.040539in}{0.679243in}}{\pgfqpoint{1.035891in}{0.684304in}}%
\pgfpathquadraticcurveto{\pgfqpoint{1.031244in}{0.689384in}}{\pgfqpoint{1.023589in}{0.689384in}}%
\pgfpathquadraticcurveto{\pgfqpoint{1.014925in}{0.689384in}}{\pgfqpoint{1.009720in}{0.684484in}}%
\pgfpathquadraticcurveto{\pgfqpoint{1.004514in}{0.679585in}}{\pgfqpoint{1.003722in}{0.670687in}}%
\pgfpathlineto{\pgfqpoint{1.040647in}{0.670759in}}%
\pgfpathlineto{\pgfqpoint{1.040647in}{0.670759in}}%
\pgfpathclose%
\pgfpathmoveto{\pgfqpoint{1.060586in}{0.696643in}}%
\pgfpathlineto{\pgfqpoint{1.071556in}{0.696643in}}%
\pgfpathlineto{\pgfqpoint{1.091261in}{0.643741in}}%
\pgfpathlineto{\pgfqpoint{1.110966in}{0.696643in}}%
\pgfpathlineto{\pgfqpoint{1.121936in}{0.696643in}}%
\pgfpathlineto{\pgfqpoint{1.098286in}{0.633600in}}%
\pgfpathlineto{\pgfqpoint{1.084218in}{0.633600in}}%
\pgfpathlineto{\pgfqpoint{1.060586in}{0.696643in}}%
\pgfpathlineto{\pgfqpoint{1.060586in}{0.696643in}}%
\pgfpathclose%
\pgfpathmoveto{\pgfqpoint{1.190166in}{0.667715in}}%
\pgfpathlineto{\pgfqpoint{1.190166in}{0.662654in}}%
\pgfpathlineto{\pgfqpoint{1.142541in}{0.662654in}}%
\pgfpathquadraticcurveto{\pgfqpoint{1.143226in}{0.651954in}}{\pgfqpoint{1.148990in}{0.646353in}}%
\pgfpathquadraticcurveto{\pgfqpoint{1.154754in}{0.640751in}}{\pgfqpoint{1.165057in}{0.640751in}}%
\pgfpathquadraticcurveto{\pgfqpoint{1.171019in}{0.640751in}}{\pgfqpoint{1.176620in}{0.642210in}}%
\pgfpathquadraticcurveto{\pgfqpoint{1.182222in}{0.643669in}}{\pgfqpoint{1.187752in}{0.646605in}}%
\pgfpathlineto{\pgfqpoint{1.187752in}{0.636806in}}%
\pgfpathquadraticcurveto{\pgfqpoint{1.182168in}{0.634447in}}{\pgfqpoint{1.176314in}{0.633204in}}%
\pgfpathquadraticcurveto{\pgfqpoint{1.170460in}{0.631961in}}{\pgfqpoint{1.164444in}{0.631961in}}%
\pgfpathquadraticcurveto{\pgfqpoint{1.149350in}{0.631961in}}{\pgfqpoint{1.140542in}{0.640733in}}%
\pgfpathquadraticcurveto{\pgfqpoint{1.131734in}{0.649523in}}{\pgfqpoint{1.131734in}{0.664509in}}%
\pgfpathquadraticcurveto{\pgfqpoint{1.131734in}{0.679981in}}{\pgfqpoint{1.140092in}{0.689059in}}%
\pgfpathquadraticcurveto{\pgfqpoint{1.148449in}{0.698156in}}{\pgfqpoint{1.162643in}{0.698156in}}%
\pgfpathquadraticcurveto{\pgfqpoint{1.175360in}{0.698156in}}{\pgfqpoint{1.182763in}{0.689960in}}%
\pgfpathquadraticcurveto{\pgfqpoint{1.190166in}{0.681783in}}{\pgfqpoint{1.190166in}{0.667715in}}%
\pgfpathlineto{\pgfqpoint{1.190166in}{0.667715in}}%
\pgfpathclose%
\pgfpathmoveto{\pgfqpoint{1.179809in}{0.670759in}}%
\pgfpathquadraticcurveto{\pgfqpoint{1.179701in}{0.679243in}}{\pgfqpoint{1.175053in}{0.684304in}}%
\pgfpathquadraticcurveto{\pgfqpoint{1.170406in}{0.689384in}}{\pgfqpoint{1.162751in}{0.689384in}}%
\pgfpathquadraticcurveto{\pgfqpoint{1.154087in}{0.689384in}}{\pgfqpoint{1.148882in}{0.684484in}}%
\pgfpathquadraticcurveto{\pgfqpoint{1.143676in}{0.679585in}}{\pgfqpoint{1.142884in}{0.670687in}}%
\pgfpathlineto{\pgfqpoint{1.179809in}{0.670759in}}%
\pgfpathlineto{\pgfqpoint{1.179809in}{0.670759in}}%
\pgfpathclose%
\pgfpathmoveto{\pgfqpoint{1.243698in}{0.686970in}}%
\pgfpathquadraticcurveto{\pgfqpoint{1.241951in}{0.687979in}}{\pgfqpoint{1.239897in}{0.688447in}}%
\pgfpathquadraticcurveto{\pgfqpoint{1.237844in}{0.688933in}}{\pgfqpoint{1.235376in}{0.688933in}}%
\pgfpathquadraticcurveto{\pgfqpoint{1.226586in}{0.688933in}}{\pgfqpoint{1.221885in}{0.683223in}}%
\pgfpathquadraticcurveto{\pgfqpoint{1.217184in}{0.677514in}}{\pgfqpoint{1.217184in}{0.666814in}}%
\pgfpathlineto{\pgfqpoint{1.217184in}{0.633600in}}%
\pgfpathlineto{\pgfqpoint{1.206773in}{0.633600in}}%
\pgfpathlineto{\pgfqpoint{1.206773in}{0.696643in}}%
\pgfpathlineto{\pgfqpoint{1.217184in}{0.696643in}}%
\pgfpathlineto{\pgfqpoint{1.217184in}{0.686844in}}%
\pgfpathquadraticcurveto{\pgfqpoint{1.220462in}{0.692590in}}{\pgfqpoint{1.225686in}{0.695364in}}%
\pgfpathquadraticcurveto{\pgfqpoint{1.230927in}{0.698156in}}{\pgfqpoint{1.238420in}{0.698156in}}%
\pgfpathquadraticcurveto{\pgfqpoint{1.239483in}{0.698156in}}{\pgfqpoint{1.240780in}{0.698011in}}%
\pgfpathquadraticcurveto{\pgfqpoint{1.242077in}{0.697885in}}{\pgfqpoint{1.243644in}{0.697597in}}%
\pgfpathlineto{\pgfqpoint{1.243698in}{0.686970in}}%
\pgfpathlineto{\pgfqpoint{1.243698in}{0.686970in}}%
\pgfpathclose%
\pgfpathmoveto{\pgfqpoint{1.283216in}{0.665283in}}%
\pgfpathquadraticcurveto{\pgfqpoint{1.270662in}{0.665283in}}{\pgfqpoint{1.265817in}{0.662419in}}%
\pgfpathquadraticcurveto{\pgfqpoint{1.260971in}{0.659556in}}{\pgfqpoint{1.260971in}{0.652621in}}%
\pgfpathquadraticcurveto{\pgfqpoint{1.260971in}{0.647109in}}{\pgfqpoint{1.264610in}{0.643867in}}%
\pgfpathquadraticcurveto{\pgfqpoint{1.268248in}{0.640643in}}{\pgfqpoint{1.274480in}{0.640643in}}%
\pgfpathquadraticcurveto{\pgfqpoint{1.283108in}{0.640643in}}{\pgfqpoint{1.288314in}{0.646749in}}%
\pgfpathquadraticcurveto{\pgfqpoint{1.293519in}{0.652855in}}{\pgfqpoint{1.293519in}{0.662978in}}%
\pgfpathlineto{\pgfqpoint{1.293519in}{0.665283in}}%
\pgfpathlineto{\pgfqpoint{1.283216in}{0.665283in}}%
\pgfpathlineto{\pgfqpoint{1.283216in}{0.665283in}}%
\pgfpathclose%
\pgfpathmoveto{\pgfqpoint{1.303876in}{0.669570in}}%
\pgfpathlineto{\pgfqpoint{1.303876in}{0.633600in}}%
\pgfpathlineto{\pgfqpoint{1.293519in}{0.633600in}}%
\pgfpathlineto{\pgfqpoint{1.293519in}{0.643164in}}%
\pgfpathquadraticcurveto{\pgfqpoint{1.289971in}{0.637437in}}{\pgfqpoint{1.284675in}{0.634699in}}%
\pgfpathquadraticcurveto{\pgfqpoint{1.279380in}{0.631961in}}{\pgfqpoint{1.271725in}{0.631961in}}%
\pgfpathquadraticcurveto{\pgfqpoint{1.262052in}{0.631961in}}{\pgfqpoint{1.256324in}{0.637401in}}%
\pgfpathquadraticcurveto{\pgfqpoint{1.250614in}{0.642840in}}{\pgfqpoint{1.250614in}{0.651954in}}%
\pgfpathquadraticcurveto{\pgfqpoint{1.250614in}{0.662582in}}{\pgfqpoint{1.257729in}{0.667985in}}%
\pgfpathquadraticcurveto{\pgfqpoint{1.264862in}{0.673389in}}{\pgfqpoint{1.278984in}{0.673389in}}%
\pgfpathlineto{\pgfqpoint{1.293519in}{0.673389in}}%
\pgfpathlineto{\pgfqpoint{1.293519in}{0.674416in}}%
\pgfpathquadraticcurveto{\pgfqpoint{1.293519in}{0.681566in}}{\pgfqpoint{1.288818in}{0.685475in}}%
\pgfpathquadraticcurveto{\pgfqpoint{1.284117in}{0.689384in}}{\pgfqpoint{1.275615in}{0.689384in}}%
\pgfpathquadraticcurveto{\pgfqpoint{1.270212in}{0.689384in}}{\pgfqpoint{1.265078in}{0.688087in}}%
\pgfpathquadraticcurveto{\pgfqpoint{1.259963in}{0.686790in}}{\pgfqpoint{1.255243in}{0.684196in}}%
\pgfpathlineto{\pgfqpoint{1.255243in}{0.693779in}}%
\pgfpathquadraticcurveto{\pgfqpoint{1.260917in}{0.695976in}}{\pgfqpoint{1.266267in}{0.697057in}}%
\pgfpathquadraticcurveto{\pgfqpoint{1.271617in}{0.698156in}}{\pgfqpoint{1.276678in}{0.698156in}}%
\pgfpathquadraticcurveto{\pgfqpoint{1.290367in}{0.698156in}}{\pgfqpoint{1.297122in}{0.691059in}}%
\pgfpathquadraticcurveto{\pgfqpoint{1.303876in}{0.683980in}}{\pgfqpoint{1.303876in}{0.669570in}}%
\pgfpathlineto{\pgfqpoint{1.303876in}{0.669570in}}%
\pgfpathclose%
\pgfpathmoveto{\pgfqpoint{1.366685in}{0.665860in}}%
\pgfpathquadraticcurveto{\pgfqpoint{1.366685in}{0.677117in}}{\pgfqpoint{1.362038in}{0.683296in}}%
\pgfpathquadraticcurveto{\pgfqpoint{1.357408in}{0.689492in}}{\pgfqpoint{1.349015in}{0.689492in}}%
\pgfpathquadraticcurveto{\pgfqpoint{1.340693in}{0.689492in}}{\pgfqpoint{1.336046in}{0.683296in}}%
\pgfpathquadraticcurveto{\pgfqpoint{1.331399in}{0.677117in}}{\pgfqpoint{1.331399in}{0.665860in}}%
\pgfpathquadraticcurveto{\pgfqpoint{1.331399in}{0.654656in}}{\pgfqpoint{1.336046in}{0.648460in}}%
\pgfpathquadraticcurveto{\pgfqpoint{1.340693in}{0.642264in}}{\pgfqpoint{1.349015in}{0.642264in}}%
\pgfpathquadraticcurveto{\pgfqpoint{1.357408in}{0.642264in}}{\pgfqpoint{1.362038in}{0.648460in}}%
\pgfpathquadraticcurveto{\pgfqpoint{1.366685in}{0.654656in}}{\pgfqpoint{1.366685in}{0.665860in}}%
\pgfpathlineto{\pgfqpoint{1.366685in}{0.665860in}}%
\pgfpathclose%
\pgfpathmoveto{\pgfqpoint{1.377042in}{0.641417in}}%
\pgfpathquadraticcurveto{\pgfqpoint{1.377042in}{0.625332in}}{\pgfqpoint{1.369891in}{0.617479in}}%
\pgfpathquadraticcurveto{\pgfqpoint{1.362758in}{0.609626in}}{\pgfqpoint{1.348006in}{0.609626in}}%
\pgfpathquadraticcurveto{\pgfqpoint{1.342548in}{0.609626in}}{\pgfqpoint{1.337703in}{0.610436in}}%
\pgfpathquadraticcurveto{\pgfqpoint{1.332858in}{0.611247in}}{\pgfqpoint{1.328301in}{0.612940in}}%
\pgfpathlineto{\pgfqpoint{1.328301in}{0.623009in}}%
\pgfpathquadraticcurveto{\pgfqpoint{1.332858in}{0.620541in}}{\pgfqpoint{1.337307in}{0.619370in}}%
\pgfpathquadraticcurveto{\pgfqpoint{1.341756in}{0.618182in}}{\pgfqpoint{1.346367in}{0.618182in}}%
\pgfpathquadraticcurveto{\pgfqpoint{1.356562in}{0.618182in}}{\pgfqpoint{1.361623in}{0.623495in}}%
\pgfpathquadraticcurveto{\pgfqpoint{1.366685in}{0.628809in}}{\pgfqpoint{1.366685in}{0.639562in}}%
\pgfpathlineto{\pgfqpoint{1.366685in}{0.644695in}}%
\pgfpathquadraticcurveto{\pgfqpoint{1.363479in}{0.639112in}}{\pgfqpoint{1.358471in}{0.636356in}}%
\pgfpathquadraticcurveto{\pgfqpoint{1.353464in}{0.633600in}}{\pgfqpoint{1.346475in}{0.633600in}}%
\pgfpathquadraticcurveto{\pgfqpoint{1.334893in}{0.633600in}}{\pgfqpoint{1.327796in}{0.642426in}}%
\pgfpathquadraticcurveto{\pgfqpoint{1.320700in}{0.651270in}}{\pgfqpoint{1.320700in}{0.665860in}}%
\pgfpathquadraticcurveto{\pgfqpoint{1.320700in}{0.680486in}}{\pgfqpoint{1.327796in}{0.689312in}}%
\pgfpathquadraticcurveto{\pgfqpoint{1.334893in}{0.698156in}}{\pgfqpoint{1.346475in}{0.698156in}}%
\pgfpathquadraticcurveto{\pgfqpoint{1.353464in}{0.698156in}}{\pgfqpoint{1.358471in}{0.695400in}}%
\pgfpathquadraticcurveto{\pgfqpoint{1.363479in}{0.692644in}}{\pgfqpoint{1.366685in}{0.687078in}}%
\pgfpathlineto{\pgfqpoint{1.366685in}{0.696643in}}%
\pgfpathlineto{\pgfqpoint{1.377042in}{0.696643in}}%
\pgfpathlineto{\pgfqpoint{1.377042in}{0.641417in}}%
\pgfpathlineto{\pgfqpoint{1.377042in}{0.641417in}}%
\pgfpathclose%
\pgfpathmoveto{\pgfqpoint{1.452314in}{0.667715in}}%
\pgfpathlineto{\pgfqpoint{1.452314in}{0.662654in}}%
\pgfpathlineto{\pgfqpoint{1.404690in}{0.662654in}}%
\pgfpathquadraticcurveto{\pgfqpoint{1.405375in}{0.651954in}}{\pgfqpoint{1.411139in}{0.646353in}}%
\pgfpathquadraticcurveto{\pgfqpoint{1.416903in}{0.640751in}}{\pgfqpoint{1.427206in}{0.640751in}}%
\pgfpathquadraticcurveto{\pgfqpoint{1.433168in}{0.640751in}}{\pgfqpoint{1.438769in}{0.642210in}}%
\pgfpathquadraticcurveto{\pgfqpoint{1.444371in}{0.643669in}}{\pgfqpoint{1.449901in}{0.646605in}}%
\pgfpathlineto{\pgfqpoint{1.449901in}{0.636806in}}%
\pgfpathquadraticcurveto{\pgfqpoint{1.444317in}{0.634447in}}{\pgfqpoint{1.438463in}{0.633204in}}%
\pgfpathquadraticcurveto{\pgfqpoint{1.432609in}{0.631961in}}{\pgfqpoint{1.426593in}{0.631961in}}%
\pgfpathquadraticcurveto{\pgfqpoint{1.411499in}{0.631961in}}{\pgfqpoint{1.402691in}{0.640733in}}%
\pgfpathquadraticcurveto{\pgfqpoint{1.393883in}{0.649523in}}{\pgfqpoint{1.393883in}{0.664509in}}%
\pgfpathquadraticcurveto{\pgfqpoint{1.393883in}{0.679981in}}{\pgfqpoint{1.402241in}{0.689059in}}%
\pgfpathquadraticcurveto{\pgfqpoint{1.410598in}{0.698156in}}{\pgfqpoint{1.424792in}{0.698156in}}%
\pgfpathquadraticcurveto{\pgfqpoint{1.437508in}{0.698156in}}{\pgfqpoint{1.444911in}{0.689960in}}%
\pgfpathquadraticcurveto{\pgfqpoint{1.452314in}{0.681783in}}{\pgfqpoint{1.452314in}{0.667715in}}%
\pgfpathlineto{\pgfqpoint{1.452314in}{0.667715in}}%
\pgfpathclose%
\pgfpathmoveto{\pgfqpoint{1.441957in}{0.670759in}}%
\pgfpathquadraticcurveto{\pgfqpoint{1.441849in}{0.679243in}}{\pgfqpoint{1.437202in}{0.684304in}}%
\pgfpathquadraticcurveto{\pgfqpoint{1.432555in}{0.689384in}}{\pgfqpoint{1.424900in}{0.689384in}}%
\pgfpathquadraticcurveto{\pgfqpoint{1.416236in}{0.689384in}}{\pgfqpoint{1.411031in}{0.684484in}}%
\pgfpathquadraticcurveto{\pgfqpoint{1.405825in}{0.679585in}}{\pgfqpoint{1.405033in}{0.670687in}}%
\pgfpathlineto{\pgfqpoint{1.441957in}{0.670759in}}%
\pgfpathlineto{\pgfqpoint{1.441957in}{0.670759in}}%
\pgfpathclose%
\pgfusepath{fill}%
\end{pgfscope}%
\begin{pgfscope}%
\pgfsetbuttcap%
\pgfsetroundjoin%
\definecolor{currentfill}{rgb}{0.309804,0.372549,0.419608}%
\pgfsetfillcolor{currentfill}%
\pgfsetlinewidth{0.000000pt}%
\definecolor{currentstroke}{rgb}{0.309804,0.372549,0.419608}%
\pgfsetstrokecolor{currentstroke}%
\pgfsetdash{}{0pt}%
\pgfpathmoveto{\pgfqpoint{1.539344in}{0.658475in}}%
\pgfpathlineto{\pgfqpoint{1.539344in}{0.696643in}}%
\pgfpathlineto{\pgfqpoint{1.549701in}{0.696643in}}%
\pgfpathlineto{\pgfqpoint{1.549701in}{0.658871in}}%
\pgfpathquadraticcurveto{\pgfqpoint{1.549701in}{0.649919in}}{\pgfqpoint{1.553177in}{0.645434in}}%
\pgfpathquadraticcurveto{\pgfqpoint{1.556671in}{0.640967in}}{\pgfqpoint{1.563660in}{0.640967in}}%
\pgfpathquadraticcurveto{\pgfqpoint{1.572036in}{0.640967in}}{\pgfqpoint{1.576899in}{0.646317in}}%
\pgfpathquadraticcurveto{\pgfqpoint{1.581780in}{0.651666in}}{\pgfqpoint{1.581780in}{0.660906in}}%
\pgfpathlineto{\pgfqpoint{1.581780in}{0.696643in}}%
\pgfpathlineto{\pgfqpoint{1.592137in}{0.696643in}}%
\pgfpathlineto{\pgfqpoint{1.592137in}{0.633600in}}%
\pgfpathlineto{\pgfqpoint{1.581780in}{0.633600in}}%
\pgfpathlineto{\pgfqpoint{1.581780in}{0.643291in}}%
\pgfpathquadraticcurveto{\pgfqpoint{1.578016in}{0.637545in}}{\pgfqpoint{1.573026in}{0.634753in}}%
\pgfpathquadraticcurveto{\pgfqpoint{1.568055in}{0.631961in}}{\pgfqpoint{1.561463in}{0.631961in}}%
\pgfpathquadraticcurveto{\pgfqpoint{1.550601in}{0.631961in}}{\pgfqpoint{1.544964in}{0.638715in}}%
\pgfpathquadraticcurveto{\pgfqpoint{1.539344in}{0.645470in}}{\pgfqpoint{1.539344in}{0.658475in}}%
\pgfpathlineto{\pgfqpoint{1.539344in}{0.658475in}}%
\pgfpathclose%
\pgfpathmoveto{\pgfqpoint{1.565407in}{0.698156in}}%
\pgfpathlineto{\pgfqpoint{1.565407in}{0.698156in}}%
\pgfpathlineto{\pgfqpoint{1.565407in}{0.698156in}}%
\pgfpathclose%
\pgfpathmoveto{\pgfqpoint{1.653649in}{0.694787in}}%
\pgfpathlineto{\pgfqpoint{1.653649in}{0.684989in}}%
\pgfpathquadraticcurveto{\pgfqpoint{1.649272in}{0.687240in}}{\pgfqpoint{1.644535in}{0.688357in}}%
\pgfpathquadraticcurveto{\pgfqpoint{1.639816in}{0.689492in}}{\pgfqpoint{1.634736in}{0.689492in}}%
\pgfpathquadraticcurveto{\pgfqpoint{1.627027in}{0.689492in}}{\pgfqpoint{1.623172in}{0.687132in}}%
\pgfpathquadraticcurveto{\pgfqpoint{1.619318in}{0.684773in}}{\pgfqpoint{1.619318in}{0.680035in}}%
\pgfpathquadraticcurveto{\pgfqpoint{1.619318in}{0.676433in}}{\pgfqpoint{1.622074in}{0.674380in}}%
\pgfpathquadraticcurveto{\pgfqpoint{1.624829in}{0.672326in}}{\pgfqpoint{1.633169in}{0.670471in}}%
\pgfpathlineto{\pgfqpoint{1.636717in}{0.669678in}}%
\pgfpathquadraticcurveto{\pgfqpoint{1.647741in}{0.667319in}}{\pgfqpoint{1.652388in}{0.663014in}}%
\pgfpathquadraticcurveto{\pgfqpoint{1.657035in}{0.658709in}}{\pgfqpoint{1.657035in}{0.651000in}}%
\pgfpathquadraticcurveto{\pgfqpoint{1.657035in}{0.642210in}}{\pgfqpoint{1.650082in}{0.637076in}}%
\pgfpathquadraticcurveto{\pgfqpoint{1.643130in}{0.631961in}}{\pgfqpoint{1.630972in}{0.631961in}}%
\pgfpathquadraticcurveto{\pgfqpoint{1.625910in}{0.631961in}}{\pgfqpoint{1.620416in}{0.632952in}}%
\pgfpathquadraticcurveto{\pgfqpoint{1.614923in}{0.633942in}}{\pgfqpoint{1.608853in}{0.635906in}}%
\pgfpathlineto{\pgfqpoint{1.608853in}{0.646605in}}%
\pgfpathquadraticcurveto{\pgfqpoint{1.614599in}{0.643615in}}{\pgfqpoint{1.620164in}{0.642120in}}%
\pgfpathquadraticcurveto{\pgfqpoint{1.625730in}{0.640643in}}{\pgfqpoint{1.631206in}{0.640643in}}%
\pgfpathquadraticcurveto{\pgfqpoint{1.638519in}{0.640643in}}{\pgfqpoint{1.642445in}{0.643146in}}%
\pgfpathquadraticcurveto{\pgfqpoint{1.646390in}{0.645650in}}{\pgfqpoint{1.646390in}{0.650207in}}%
\pgfpathquadraticcurveto{\pgfqpoint{1.646390in}{0.654422in}}{\pgfqpoint{1.643544in}{0.656674in}}%
\pgfpathquadraticcurveto{\pgfqpoint{1.640716in}{0.658925in}}{\pgfqpoint{1.631080in}{0.661014in}}%
\pgfpathlineto{\pgfqpoint{1.627477in}{0.661861in}}%
\pgfpathquadraticcurveto{\pgfqpoint{1.617859in}{0.663878in}}{\pgfqpoint{1.613572in}{0.668075in}}%
\pgfpathquadraticcurveto{\pgfqpoint{1.609303in}{0.672272in}}{\pgfqpoint{1.609303in}{0.679585in}}%
\pgfpathquadraticcurveto{\pgfqpoint{1.609303in}{0.688483in}}{\pgfqpoint{1.615607in}{0.693310in}}%
\pgfpathquadraticcurveto{\pgfqpoint{1.621911in}{0.698156in}}{\pgfqpoint{1.633511in}{0.698156in}}%
\pgfpathquadraticcurveto{\pgfqpoint{1.639239in}{0.698156in}}{\pgfqpoint{1.644301in}{0.697309in}}%
\pgfpathquadraticcurveto{\pgfqpoint{1.649380in}{0.696480in}}{\pgfqpoint{1.653649in}{0.694787in}}%
\pgfpathlineto{\pgfqpoint{1.653649in}{0.694787in}}%
\pgfpathclose%
\pgfpathmoveto{\pgfqpoint{1.727445in}{0.667715in}}%
\pgfpathlineto{\pgfqpoint{1.727445in}{0.662654in}}%
\pgfpathlineto{\pgfqpoint{1.679821in}{0.662654in}}%
\pgfpathquadraticcurveto{\pgfqpoint{1.680505in}{0.651954in}}{\pgfqpoint{1.686269in}{0.646353in}}%
\pgfpathquadraticcurveto{\pgfqpoint{1.692033in}{0.640751in}}{\pgfqpoint{1.702336in}{0.640751in}}%
\pgfpathquadraticcurveto{\pgfqpoint{1.708298in}{0.640751in}}{\pgfqpoint{1.713900in}{0.642210in}}%
\pgfpathquadraticcurveto{\pgfqpoint{1.719501in}{0.643669in}}{\pgfqpoint{1.725031in}{0.646605in}}%
\pgfpathlineto{\pgfqpoint{1.725031in}{0.636806in}}%
\pgfpathquadraticcurveto{\pgfqpoint{1.719447in}{0.634447in}}{\pgfqpoint{1.713593in}{0.633204in}}%
\pgfpathquadraticcurveto{\pgfqpoint{1.707739in}{0.631961in}}{\pgfqpoint{1.701723in}{0.631961in}}%
\pgfpathquadraticcurveto{\pgfqpoint{1.686629in}{0.631961in}}{\pgfqpoint{1.677821in}{0.640733in}}%
\pgfpathquadraticcurveto{\pgfqpoint{1.669013in}{0.649523in}}{\pgfqpoint{1.669013in}{0.664509in}}%
\pgfpathquadraticcurveto{\pgfqpoint{1.669013in}{0.679981in}}{\pgfqpoint{1.677371in}{0.689059in}}%
\pgfpathquadraticcurveto{\pgfqpoint{1.685729in}{0.698156in}}{\pgfqpoint{1.699922in}{0.698156in}}%
\pgfpathquadraticcurveto{\pgfqpoint{1.712639in}{0.698156in}}{\pgfqpoint{1.720042in}{0.689960in}}%
\pgfpathquadraticcurveto{\pgfqpoint{1.727445in}{0.681783in}}{\pgfqpoint{1.727445in}{0.667715in}}%
\pgfpathlineto{\pgfqpoint{1.727445in}{0.667715in}}%
\pgfpathclose%
\pgfpathmoveto{\pgfqpoint{1.717088in}{0.670759in}}%
\pgfpathquadraticcurveto{\pgfqpoint{1.716980in}{0.679243in}}{\pgfqpoint{1.712332in}{0.684304in}}%
\pgfpathquadraticcurveto{\pgfqpoint{1.707685in}{0.689384in}}{\pgfqpoint{1.700030in}{0.689384in}}%
\pgfpathquadraticcurveto{\pgfqpoint{1.691366in}{0.689384in}}{\pgfqpoint{1.686161in}{0.684484in}}%
\pgfpathquadraticcurveto{\pgfqpoint{1.680955in}{0.679585in}}{\pgfqpoint{1.680163in}{0.670687in}}%
\pgfpathlineto{\pgfqpoint{1.717088in}{0.670759in}}%
\pgfpathlineto{\pgfqpoint{1.717088in}{0.670759in}}%
\pgfpathclose%
\pgfpathmoveto{\pgfqpoint{1.784633in}{0.694787in}}%
\pgfpathlineto{\pgfqpoint{1.784633in}{0.684989in}}%
\pgfpathquadraticcurveto{\pgfqpoint{1.780256in}{0.687240in}}{\pgfqpoint{1.775519in}{0.688357in}}%
\pgfpathquadraticcurveto{\pgfqpoint{1.770800in}{0.689492in}}{\pgfqpoint{1.765720in}{0.689492in}}%
\pgfpathquadraticcurveto{\pgfqpoint{1.758011in}{0.689492in}}{\pgfqpoint{1.754157in}{0.687132in}}%
\pgfpathquadraticcurveto{\pgfqpoint{1.750302in}{0.684773in}}{\pgfqpoint{1.750302in}{0.680035in}}%
\pgfpathquadraticcurveto{\pgfqpoint{1.750302in}{0.676433in}}{\pgfqpoint{1.753058in}{0.674380in}}%
\pgfpathquadraticcurveto{\pgfqpoint{1.755814in}{0.672326in}}{\pgfqpoint{1.764153in}{0.670471in}}%
\pgfpathlineto{\pgfqpoint{1.767702in}{0.669678in}}%
\pgfpathquadraticcurveto{\pgfqpoint{1.778725in}{0.667319in}}{\pgfqpoint{1.783372in}{0.663014in}}%
\pgfpathquadraticcurveto{\pgfqpoint{1.788020in}{0.658709in}}{\pgfqpoint{1.788020in}{0.651000in}}%
\pgfpathquadraticcurveto{\pgfqpoint{1.788020in}{0.642210in}}{\pgfqpoint{1.781067in}{0.637076in}}%
\pgfpathquadraticcurveto{\pgfqpoint{1.774114in}{0.631961in}}{\pgfqpoint{1.761956in}{0.631961in}}%
\pgfpathquadraticcurveto{\pgfqpoint{1.756895in}{0.631961in}}{\pgfqpoint{1.751401in}{0.632952in}}%
\pgfpathquadraticcurveto{\pgfqpoint{1.745907in}{0.633942in}}{\pgfqpoint{1.739837in}{0.635906in}}%
\pgfpathlineto{\pgfqpoint{1.739837in}{0.646605in}}%
\pgfpathquadraticcurveto{\pgfqpoint{1.745583in}{0.643615in}}{\pgfqpoint{1.751149in}{0.642120in}}%
\pgfpathquadraticcurveto{\pgfqpoint{1.756714in}{0.640643in}}{\pgfqpoint{1.762190in}{0.640643in}}%
\pgfpathquadraticcurveto{\pgfqpoint{1.769503in}{0.640643in}}{\pgfqpoint{1.773430in}{0.643146in}}%
\pgfpathquadraticcurveto{\pgfqpoint{1.777374in}{0.645650in}}{\pgfqpoint{1.777374in}{0.650207in}}%
\pgfpathquadraticcurveto{\pgfqpoint{1.777374in}{0.654422in}}{\pgfqpoint{1.774528in}{0.656674in}}%
\pgfpathquadraticcurveto{\pgfqpoint{1.771701in}{0.658925in}}{\pgfqpoint{1.762064in}{0.661014in}}%
\pgfpathlineto{\pgfqpoint{1.758462in}{0.661861in}}%
\pgfpathquadraticcurveto{\pgfqpoint{1.748843in}{0.663878in}}{\pgfqpoint{1.744556in}{0.668075in}}%
\pgfpathquadraticcurveto{\pgfqpoint{1.740287in}{0.672272in}}{\pgfqpoint{1.740287in}{0.679585in}}%
\pgfpathquadraticcurveto{\pgfqpoint{1.740287in}{0.688483in}}{\pgfqpoint{1.746592in}{0.693310in}}%
\pgfpathquadraticcurveto{\pgfqpoint{1.752896in}{0.698156in}}{\pgfqpoint{1.764496in}{0.698156in}}%
\pgfpathquadraticcurveto{\pgfqpoint{1.770224in}{0.698156in}}{\pgfqpoint{1.775285in}{0.697309in}}%
\pgfpathquadraticcurveto{\pgfqpoint{1.780364in}{0.696480in}}{\pgfqpoint{1.784633in}{0.694787in}}%
\pgfpathlineto{\pgfqpoint{1.784633in}{0.694787in}}%
\pgfpathclose%
\pgfusepath{fill}%
\end{pgfscope}%
\begin{pgfscope}%
\pgfsetbuttcap%
\pgfsetroundjoin%
\definecolor{currentfill}{rgb}{0.309804,0.372549,0.419608}%
\pgfsetfillcolor{currentfill}%
\pgfsetlinewidth{0.000000pt}%
\definecolor{currentstroke}{rgb}{0.309804,0.372549,0.419608}%
\pgfsetstrokecolor{currentstroke}%
\pgfsetdash{}{0pt}%
\pgfpathmoveto{\pgfqpoint{1.908168in}{0.694787in}}%
\pgfpathlineto{\pgfqpoint{1.908168in}{0.684989in}}%
\pgfpathquadraticcurveto{\pgfqpoint{1.903791in}{0.687240in}}{\pgfqpoint{1.899054in}{0.688357in}}%
\pgfpathquadraticcurveto{\pgfqpoint{1.894335in}{0.689492in}}{\pgfqpoint{1.889255in}{0.689492in}}%
\pgfpathquadraticcurveto{\pgfqpoint{1.881546in}{0.689492in}}{\pgfqpoint{1.877691in}{0.687132in}}%
\pgfpathquadraticcurveto{\pgfqpoint{1.873837in}{0.684773in}}{\pgfqpoint{1.873837in}{0.680035in}}%
\pgfpathquadraticcurveto{\pgfqpoint{1.873837in}{0.676433in}}{\pgfqpoint{1.876593in}{0.674380in}}%
\pgfpathquadraticcurveto{\pgfqpoint{1.879349in}{0.672326in}}{\pgfqpoint{1.887688in}{0.670471in}}%
\pgfpathlineto{\pgfqpoint{1.891237in}{0.669678in}}%
\pgfpathquadraticcurveto{\pgfqpoint{1.902260in}{0.667319in}}{\pgfqpoint{1.906907in}{0.663014in}}%
\pgfpathquadraticcurveto{\pgfqpoint{1.911554in}{0.658709in}}{\pgfqpoint{1.911554in}{0.651000in}}%
\pgfpathquadraticcurveto{\pgfqpoint{1.911554in}{0.642210in}}{\pgfqpoint{1.904602in}{0.637076in}}%
\pgfpathquadraticcurveto{\pgfqpoint{1.897649in}{0.631961in}}{\pgfqpoint{1.885491in}{0.631961in}}%
\pgfpathquadraticcurveto{\pgfqpoint{1.880429in}{0.631961in}}{\pgfqpoint{1.874936in}{0.632952in}}%
\pgfpathquadraticcurveto{\pgfqpoint{1.869442in}{0.633942in}}{\pgfqpoint{1.863372in}{0.635906in}}%
\pgfpathlineto{\pgfqpoint{1.863372in}{0.646605in}}%
\pgfpathquadraticcurveto{\pgfqpoint{1.869118in}{0.643615in}}{\pgfqpoint{1.874683in}{0.642120in}}%
\pgfpathquadraticcurveto{\pgfqpoint{1.880249in}{0.640643in}}{\pgfqpoint{1.885725in}{0.640643in}}%
\pgfpathquadraticcurveto{\pgfqpoint{1.893038in}{0.640643in}}{\pgfqpoint{1.896964in}{0.643146in}}%
\pgfpathquadraticcurveto{\pgfqpoint{1.900909in}{0.645650in}}{\pgfqpoint{1.900909in}{0.650207in}}%
\pgfpathquadraticcurveto{\pgfqpoint{1.900909in}{0.654422in}}{\pgfqpoint{1.898063in}{0.656674in}}%
\pgfpathquadraticcurveto{\pgfqpoint{1.895235in}{0.658925in}}{\pgfqpoint{1.885599in}{0.661014in}}%
\pgfpathlineto{\pgfqpoint{1.881996in}{0.661861in}}%
\pgfpathquadraticcurveto{\pgfqpoint{1.872378in}{0.663878in}}{\pgfqpoint{1.868091in}{0.668075in}}%
\pgfpathquadraticcurveto{\pgfqpoint{1.863822in}{0.672272in}}{\pgfqpoint{1.863822in}{0.679585in}}%
\pgfpathquadraticcurveto{\pgfqpoint{1.863822in}{0.688483in}}{\pgfqpoint{1.870126in}{0.693310in}}%
\pgfpathquadraticcurveto{\pgfqpoint{1.876431in}{0.698156in}}{\pgfqpoint{1.888030in}{0.698156in}}%
\pgfpathquadraticcurveto{\pgfqpoint{1.893758in}{0.698156in}}{\pgfqpoint{1.898820in}{0.697309in}}%
\pgfpathquadraticcurveto{\pgfqpoint{1.903899in}{0.696480in}}{\pgfqpoint{1.908168in}{0.694787in}}%
\pgfpathlineto{\pgfqpoint{1.908168in}{0.694787in}}%
\pgfpathclose%
\pgfpathmoveto{\pgfqpoint{1.934232in}{0.665067in}}%
\pgfpathquadraticcurveto{\pgfqpoint{1.934232in}{0.653648in}}{\pgfqpoint{1.938933in}{0.647145in}}%
\pgfpathquadraticcurveto{\pgfqpoint{1.943634in}{0.640643in}}{\pgfqpoint{1.951847in}{0.640643in}}%
\pgfpathquadraticcurveto{\pgfqpoint{1.960061in}{0.640643in}}{\pgfqpoint{1.964780in}{0.647145in}}%
\pgfpathquadraticcurveto{\pgfqpoint{1.969517in}{0.653648in}}{\pgfqpoint{1.969517in}{0.665067in}}%
\pgfpathquadraticcurveto{\pgfqpoint{1.969517in}{0.676487in}}{\pgfqpoint{1.964780in}{0.682989in}}%
\pgfpathquadraticcurveto{\pgfqpoint{1.960061in}{0.689492in}}{\pgfqpoint{1.951847in}{0.689492in}}%
\pgfpathquadraticcurveto{\pgfqpoint{1.943634in}{0.689492in}}{\pgfqpoint{1.938933in}{0.682989in}}%
\pgfpathquadraticcurveto{\pgfqpoint{1.934232in}{0.676487in}}{\pgfqpoint{1.934232in}{0.665067in}}%
\pgfpathlineto{\pgfqpoint{1.934232in}{0.665067in}}%
\pgfpathclose%
\pgfpathmoveto{\pgfqpoint{1.969517in}{0.643056in}}%
\pgfpathquadraticcurveto{\pgfqpoint{1.966257in}{0.637437in}}{\pgfqpoint{1.961268in}{0.634699in}}%
\pgfpathquadraticcurveto{\pgfqpoint{1.956296in}{0.631961in}}{\pgfqpoint{1.949308in}{0.631961in}}%
\pgfpathquadraticcurveto{\pgfqpoint{1.937888in}{0.631961in}}{\pgfqpoint{1.930701in}{0.641075in}}%
\pgfpathquadraticcurveto{\pgfqpoint{1.923532in}{0.650207in}}{\pgfqpoint{1.923532in}{0.665067in}}%
\pgfpathquadraticcurveto{\pgfqpoint{1.923532in}{0.679927in}}{\pgfqpoint{1.930701in}{0.689041in}}%
\pgfpathquadraticcurveto{\pgfqpoint{1.937888in}{0.698156in}}{\pgfqpoint{1.949308in}{0.698156in}}%
\pgfpathquadraticcurveto{\pgfqpoint{1.956296in}{0.698156in}}{\pgfqpoint{1.961268in}{0.695418in}}%
\pgfpathquadraticcurveto{\pgfqpoint{1.966257in}{0.692698in}}{\pgfqpoint{1.969517in}{0.687078in}}%
\pgfpathlineto{\pgfqpoint{1.969517in}{0.696643in}}%
\pgfpathlineto{\pgfqpoint{1.979874in}{0.696643in}}%
\pgfpathlineto{\pgfqpoint{1.979874in}{0.609626in}}%
\pgfpathlineto{\pgfqpoint{1.969517in}{0.609626in}}%
\pgfpathlineto{\pgfqpoint{1.969517in}{0.643056in}}%
\pgfpathlineto{\pgfqpoint{1.969517in}{0.643056in}}%
\pgfpathclose%
\pgfpathmoveto{\pgfqpoint{2.000156in}{0.658475in}}%
\pgfpathlineto{\pgfqpoint{2.000156in}{0.696643in}}%
\pgfpathlineto{\pgfqpoint{2.010513in}{0.696643in}}%
\pgfpathlineto{\pgfqpoint{2.010513in}{0.658871in}}%
\pgfpathquadraticcurveto{\pgfqpoint{2.010513in}{0.649919in}}{\pgfqpoint{2.013989in}{0.645434in}}%
\pgfpathquadraticcurveto{\pgfqpoint{2.017484in}{0.640967in}}{\pgfqpoint{2.024472in}{0.640967in}}%
\pgfpathquadraticcurveto{\pgfqpoint{2.032848in}{0.640967in}}{\pgfqpoint{2.037711in}{0.646317in}}%
\pgfpathquadraticcurveto{\pgfqpoint{2.042593in}{0.651666in}}{\pgfqpoint{2.042593in}{0.660906in}}%
\pgfpathlineto{\pgfqpoint{2.042593in}{0.696643in}}%
\pgfpathlineto{\pgfqpoint{2.052950in}{0.696643in}}%
\pgfpathlineto{\pgfqpoint{2.052950in}{0.633600in}}%
\pgfpathlineto{\pgfqpoint{2.042593in}{0.633600in}}%
\pgfpathlineto{\pgfqpoint{2.042593in}{0.643291in}}%
\pgfpathquadraticcurveto{\pgfqpoint{2.038828in}{0.637545in}}{\pgfqpoint{2.033839in}{0.634753in}}%
\pgfpathquadraticcurveto{\pgfqpoint{2.028867in}{0.631961in}}{\pgfqpoint{2.022275in}{0.631961in}}%
\pgfpathquadraticcurveto{\pgfqpoint{2.011414in}{0.631961in}}{\pgfqpoint{2.005776in}{0.638715in}}%
\pgfpathquadraticcurveto{\pgfqpoint{2.000156in}{0.645470in}}{\pgfqpoint{2.000156in}{0.658475in}}%
\pgfpathlineto{\pgfqpoint{2.000156in}{0.658475in}}%
\pgfpathclose%
\pgfpathmoveto{\pgfqpoint{2.026220in}{0.698156in}}%
\pgfpathlineto{\pgfqpoint{2.026220in}{0.698156in}}%
\pgfpathlineto{\pgfqpoint{2.026220in}{0.698156in}}%
\pgfpathclose%
\pgfpathmoveto{\pgfqpoint{2.102933in}{0.665283in}}%
\pgfpathquadraticcurveto{\pgfqpoint{2.090379in}{0.665283in}}{\pgfqpoint{2.085534in}{0.662419in}}%
\pgfpathquadraticcurveto{\pgfqpoint{2.080688in}{0.659556in}}{\pgfqpoint{2.080688in}{0.652621in}}%
\pgfpathquadraticcurveto{\pgfqpoint{2.080688in}{0.647109in}}{\pgfqpoint{2.084327in}{0.643867in}}%
\pgfpathquadraticcurveto{\pgfqpoint{2.087965in}{0.640643in}}{\pgfqpoint{2.094197in}{0.640643in}}%
\pgfpathquadraticcurveto{\pgfqpoint{2.102825in}{0.640643in}}{\pgfqpoint{2.108031in}{0.646749in}}%
\pgfpathquadraticcurveto{\pgfqpoint{2.113236in}{0.652855in}}{\pgfqpoint{2.113236in}{0.662978in}}%
\pgfpathlineto{\pgfqpoint{2.113236in}{0.665283in}}%
\pgfpathlineto{\pgfqpoint{2.102933in}{0.665283in}}%
\pgfpathlineto{\pgfqpoint{2.102933in}{0.665283in}}%
\pgfpathclose%
\pgfpathmoveto{\pgfqpoint{2.123593in}{0.669570in}}%
\pgfpathlineto{\pgfqpoint{2.123593in}{0.633600in}}%
\pgfpathlineto{\pgfqpoint{2.113236in}{0.633600in}}%
\pgfpathlineto{\pgfqpoint{2.113236in}{0.643164in}}%
\pgfpathquadraticcurveto{\pgfqpoint{2.109688in}{0.637437in}}{\pgfqpoint{2.104392in}{0.634699in}}%
\pgfpathquadraticcurveto{\pgfqpoint{2.099097in}{0.631961in}}{\pgfqpoint{2.091442in}{0.631961in}}%
\pgfpathquadraticcurveto{\pgfqpoint{2.081769in}{0.631961in}}{\pgfqpoint{2.076041in}{0.637401in}}%
\pgfpathquadraticcurveto{\pgfqpoint{2.070331in}{0.642840in}}{\pgfqpoint{2.070331in}{0.651954in}}%
\pgfpathquadraticcurveto{\pgfqpoint{2.070331in}{0.662582in}}{\pgfqpoint{2.077446in}{0.667985in}}%
\pgfpathquadraticcurveto{\pgfqpoint{2.084579in}{0.673389in}}{\pgfqpoint{2.098701in}{0.673389in}}%
\pgfpathlineto{\pgfqpoint{2.113236in}{0.673389in}}%
\pgfpathlineto{\pgfqpoint{2.113236in}{0.674416in}}%
\pgfpathquadraticcurveto{\pgfqpoint{2.113236in}{0.681566in}}{\pgfqpoint{2.108535in}{0.685475in}}%
\pgfpathquadraticcurveto{\pgfqpoint{2.103834in}{0.689384in}}{\pgfqpoint{2.095332in}{0.689384in}}%
\pgfpathquadraticcurveto{\pgfqpoint{2.089929in}{0.689384in}}{\pgfqpoint{2.084795in}{0.688087in}}%
\pgfpathquadraticcurveto{\pgfqpoint{2.079680in}{0.686790in}}{\pgfqpoint{2.074961in}{0.684196in}}%
\pgfpathlineto{\pgfqpoint{2.074961in}{0.693779in}}%
\pgfpathquadraticcurveto{\pgfqpoint{2.080634in}{0.695976in}}{\pgfqpoint{2.085984in}{0.697057in}}%
\pgfpathquadraticcurveto{\pgfqpoint{2.091334in}{0.698156in}}{\pgfqpoint{2.096395in}{0.698156in}}%
\pgfpathquadraticcurveto{\pgfqpoint{2.110084in}{0.698156in}}{\pgfqpoint{2.116839in}{0.691059in}}%
\pgfpathquadraticcurveto{\pgfqpoint{2.123593in}{0.683980in}}{\pgfqpoint{2.123593in}{0.669570in}}%
\pgfpathlineto{\pgfqpoint{2.123593in}{0.669570in}}%
\pgfpathclose%
\pgfpathmoveto{\pgfqpoint{2.181448in}{0.686970in}}%
\pgfpathquadraticcurveto{\pgfqpoint{2.179701in}{0.687979in}}{\pgfqpoint{2.177648in}{0.688447in}}%
\pgfpathquadraticcurveto{\pgfqpoint{2.175594in}{0.688933in}}{\pgfqpoint{2.173127in}{0.688933in}}%
\pgfpathquadraticcurveto{\pgfqpoint{2.164337in}{0.688933in}}{\pgfqpoint{2.159636in}{0.683223in}}%
\pgfpathquadraticcurveto{\pgfqpoint{2.154934in}{0.677514in}}{\pgfqpoint{2.154934in}{0.666814in}}%
\pgfpathlineto{\pgfqpoint{2.154934in}{0.633600in}}%
\pgfpathlineto{\pgfqpoint{2.144523in}{0.633600in}}%
\pgfpathlineto{\pgfqpoint{2.144523in}{0.696643in}}%
\pgfpathlineto{\pgfqpoint{2.154934in}{0.696643in}}%
\pgfpathlineto{\pgfqpoint{2.154934in}{0.686844in}}%
\pgfpathquadraticcurveto{\pgfqpoint{2.158213in}{0.692590in}}{\pgfqpoint{2.163436in}{0.695364in}}%
\pgfpathquadraticcurveto{\pgfqpoint{2.168678in}{0.698156in}}{\pgfqpoint{2.176171in}{0.698156in}}%
\pgfpathquadraticcurveto{\pgfqpoint{2.177234in}{0.698156in}}{\pgfqpoint{2.178530in}{0.698011in}}%
\pgfpathquadraticcurveto{\pgfqpoint{2.179827in}{0.697885in}}{\pgfqpoint{2.181394in}{0.697597in}}%
\pgfpathlineto{\pgfqpoint{2.181448in}{0.686970in}}%
\pgfpathlineto{\pgfqpoint{2.181448in}{0.686970in}}%
\pgfpathclose%
\pgfpathmoveto{\pgfqpoint{2.243698in}{0.667715in}}%
\pgfpathlineto{\pgfqpoint{2.243698in}{0.662654in}}%
\pgfpathlineto{\pgfqpoint{2.196074in}{0.662654in}}%
\pgfpathquadraticcurveto{\pgfqpoint{2.196759in}{0.651954in}}{\pgfqpoint{2.202523in}{0.646353in}}%
\pgfpathquadraticcurveto{\pgfqpoint{2.208286in}{0.640751in}}{\pgfqpoint{2.218589in}{0.640751in}}%
\pgfpathquadraticcurveto{\pgfqpoint{2.224551in}{0.640751in}}{\pgfqpoint{2.230153in}{0.642210in}}%
\pgfpathquadraticcurveto{\pgfqpoint{2.235755in}{0.643669in}}{\pgfqpoint{2.241285in}{0.646605in}}%
\pgfpathlineto{\pgfqpoint{2.241285in}{0.636806in}}%
\pgfpathquadraticcurveto{\pgfqpoint{2.235701in}{0.634447in}}{\pgfqpoint{2.229847in}{0.633204in}}%
\pgfpathquadraticcurveto{\pgfqpoint{2.223993in}{0.631961in}}{\pgfqpoint{2.217977in}{0.631961in}}%
\pgfpathquadraticcurveto{\pgfqpoint{2.202883in}{0.631961in}}{\pgfqpoint{2.194075in}{0.640733in}}%
\pgfpathquadraticcurveto{\pgfqpoint{2.185267in}{0.649523in}}{\pgfqpoint{2.185267in}{0.664509in}}%
\pgfpathquadraticcurveto{\pgfqpoint{2.185267in}{0.679981in}}{\pgfqpoint{2.193625in}{0.689059in}}%
\pgfpathquadraticcurveto{\pgfqpoint{2.201982in}{0.698156in}}{\pgfqpoint{2.216176in}{0.698156in}}%
\pgfpathquadraticcurveto{\pgfqpoint{2.228892in}{0.698156in}}{\pgfqpoint{2.236295in}{0.689960in}}%
\pgfpathquadraticcurveto{\pgfqpoint{2.243698in}{0.681783in}}{\pgfqpoint{2.243698in}{0.667715in}}%
\pgfpathlineto{\pgfqpoint{2.243698in}{0.667715in}}%
\pgfpathclose%
\pgfpathmoveto{\pgfqpoint{2.233341in}{0.670759in}}%
\pgfpathquadraticcurveto{\pgfqpoint{2.233233in}{0.679243in}}{\pgfqpoint{2.228586in}{0.684304in}}%
\pgfpathquadraticcurveto{\pgfqpoint{2.223939in}{0.689384in}}{\pgfqpoint{2.216284in}{0.689384in}}%
\pgfpathquadraticcurveto{\pgfqpoint{2.207620in}{0.689384in}}{\pgfqpoint{2.202414in}{0.684484in}}%
\pgfpathquadraticcurveto{\pgfqpoint{2.197209in}{0.679585in}}{\pgfqpoint{2.196416in}{0.670687in}}%
\pgfpathlineto{\pgfqpoint{2.233341in}{0.670759in}}%
\pgfpathlineto{\pgfqpoint{2.233341in}{0.670759in}}%
\pgfpathclose%
\pgfpathmoveto{\pgfqpoint{2.302184in}{0.687078in}}%
\pgfpathlineto{\pgfqpoint{2.302184in}{0.721193in}}%
\pgfpathlineto{\pgfqpoint{2.312541in}{0.721193in}}%
\pgfpathlineto{\pgfqpoint{2.312541in}{0.633600in}}%
\pgfpathlineto{\pgfqpoint{2.302184in}{0.633600in}}%
\pgfpathlineto{\pgfqpoint{2.302184in}{0.643056in}}%
\pgfpathquadraticcurveto{\pgfqpoint{2.298924in}{0.637437in}}{\pgfqpoint{2.293934in}{0.634699in}}%
\pgfpathquadraticcurveto{\pgfqpoint{2.288963in}{0.631961in}}{\pgfqpoint{2.281974in}{0.631961in}}%
\pgfpathquadraticcurveto{\pgfqpoint{2.270554in}{0.631961in}}{\pgfqpoint{2.263368in}{0.641075in}}%
\pgfpathquadraticcurveto{\pgfqpoint{2.256199in}{0.650207in}}{\pgfqpoint{2.256199in}{0.665067in}}%
\pgfpathquadraticcurveto{\pgfqpoint{2.256199in}{0.679927in}}{\pgfqpoint{2.263368in}{0.689041in}}%
\pgfpathquadraticcurveto{\pgfqpoint{2.270554in}{0.698156in}}{\pgfqpoint{2.281974in}{0.698156in}}%
\pgfpathquadraticcurveto{\pgfqpoint{2.288963in}{0.698156in}}{\pgfqpoint{2.293934in}{0.695418in}}%
\pgfpathquadraticcurveto{\pgfqpoint{2.298924in}{0.692698in}}{\pgfqpoint{2.302184in}{0.687078in}}%
\pgfpathlineto{\pgfqpoint{2.302184in}{0.687078in}}%
\pgfpathclose%
\pgfpathmoveto{\pgfqpoint{2.266898in}{0.665067in}}%
\pgfpathquadraticcurveto{\pgfqpoint{2.266898in}{0.653648in}}{\pgfqpoint{2.271599in}{0.647145in}}%
\pgfpathquadraticcurveto{\pgfqpoint{2.276300in}{0.640643in}}{\pgfqpoint{2.284514in}{0.640643in}}%
\pgfpathquadraticcurveto{\pgfqpoint{2.292727in}{0.640643in}}{\pgfqpoint{2.297447in}{0.647145in}}%
\pgfpathquadraticcurveto{\pgfqpoint{2.302184in}{0.653648in}}{\pgfqpoint{2.302184in}{0.665067in}}%
\pgfpathquadraticcurveto{\pgfqpoint{2.302184in}{0.676487in}}{\pgfqpoint{2.297447in}{0.682989in}}%
\pgfpathquadraticcurveto{\pgfqpoint{2.292727in}{0.689492in}}{\pgfqpoint{2.284514in}{0.689492in}}%
\pgfpathquadraticcurveto{\pgfqpoint{2.276300in}{0.689492in}}{\pgfqpoint{2.271599in}{0.682989in}}%
\pgfpathquadraticcurveto{\pgfqpoint{2.266898in}{0.676487in}}{\pgfqpoint{2.266898in}{0.665067in}}%
\pgfpathlineto{\pgfqpoint{2.266898in}{0.665067in}}%
\pgfpathclose%
\pgfusepath{fill}%
\end{pgfscope}%
\begin{pgfscope}%
\pgfsetbuttcap%
\pgfsetroundjoin%
\definecolor{currentfill}{rgb}{0.309804,0.372549,0.419608}%
\pgfsetfillcolor{currentfill}%
\pgfsetlinewidth{0.000000pt}%
\definecolor{currentstroke}{rgb}{0.309804,0.372549,0.419608}%
\pgfsetstrokecolor{currentstroke}%
\pgfsetdash{}{0pt}%
\pgfpathmoveto{\pgfqpoint{2.437484in}{0.689384in}}%
\pgfpathquadraticcurveto{\pgfqpoint{2.429162in}{0.689384in}}{\pgfqpoint{2.424317in}{0.682881in}}%
\pgfpathquadraticcurveto{\pgfqpoint{2.419472in}{0.676379in}}{\pgfqpoint{2.419472in}{0.665067in}}%
\pgfpathquadraticcurveto{\pgfqpoint{2.419472in}{0.653756in}}{\pgfqpoint{2.424281in}{0.647253in}}%
\pgfpathquadraticcurveto{\pgfqpoint{2.429108in}{0.640751in}}{\pgfqpoint{2.437484in}{0.640751in}}%
\pgfpathquadraticcurveto{\pgfqpoint{2.445769in}{0.640751in}}{\pgfqpoint{2.450597in}{0.647271in}}%
\pgfpathquadraticcurveto{\pgfqpoint{2.455442in}{0.653810in}}{\pgfqpoint{2.455442in}{0.665067in}}%
\pgfpathquadraticcurveto{\pgfqpoint{2.455442in}{0.676271in}}{\pgfqpoint{2.450597in}{0.682827in}}%
\pgfpathquadraticcurveto{\pgfqpoint{2.445769in}{0.689384in}}{\pgfqpoint{2.437484in}{0.689384in}}%
\pgfpathlineto{\pgfqpoint{2.437484in}{0.689384in}}%
\pgfpathclose%
\pgfpathmoveto{\pgfqpoint{2.437484in}{0.698156in}}%
\pgfpathquadraticcurveto{\pgfqpoint{2.450993in}{0.698156in}}{\pgfqpoint{2.458702in}{0.689366in}}%
\pgfpathquadraticcurveto{\pgfqpoint{2.466429in}{0.680594in}}{\pgfqpoint{2.466429in}{0.665067in}}%
\pgfpathquadraticcurveto{\pgfqpoint{2.466429in}{0.649595in}}{\pgfqpoint{2.458702in}{0.640769in}}%
\pgfpathquadraticcurveto{\pgfqpoint{2.450993in}{0.631961in}}{\pgfqpoint{2.437484in}{0.631961in}}%
\pgfpathquadraticcurveto{\pgfqpoint{2.423921in}{0.631961in}}{\pgfqpoint{2.416229in}{0.640769in}}%
\pgfpathquadraticcurveto{\pgfqpoint{2.408556in}{0.649595in}}{\pgfqpoint{2.408556in}{0.665067in}}%
\pgfpathquadraticcurveto{\pgfqpoint{2.408556in}{0.680594in}}{\pgfqpoint{2.416229in}{0.689366in}}%
\pgfpathquadraticcurveto{\pgfqpoint{2.423921in}{0.698156in}}{\pgfqpoint{2.437484in}{0.698156in}}%
\pgfpathlineto{\pgfqpoint{2.437484in}{0.698156in}}%
\pgfpathclose%
\pgfpathmoveto{\pgfqpoint{2.482532in}{0.658475in}}%
\pgfpathlineto{\pgfqpoint{2.482532in}{0.696643in}}%
\pgfpathlineto{\pgfqpoint{2.492889in}{0.696643in}}%
\pgfpathlineto{\pgfqpoint{2.492889in}{0.658871in}}%
\pgfpathquadraticcurveto{\pgfqpoint{2.492889in}{0.649919in}}{\pgfqpoint{2.496365in}{0.645434in}}%
\pgfpathquadraticcurveto{\pgfqpoint{2.499860in}{0.640967in}}{\pgfqpoint{2.506849in}{0.640967in}}%
\pgfpathquadraticcurveto{\pgfqpoint{2.515224in}{0.640967in}}{\pgfqpoint{2.520087in}{0.646317in}}%
\pgfpathquadraticcurveto{\pgfqpoint{2.524969in}{0.651666in}}{\pgfqpoint{2.524969in}{0.660906in}}%
\pgfpathlineto{\pgfqpoint{2.524969in}{0.696643in}}%
\pgfpathlineto{\pgfqpoint{2.535326in}{0.696643in}}%
\pgfpathlineto{\pgfqpoint{2.535326in}{0.633600in}}%
\pgfpathlineto{\pgfqpoint{2.524969in}{0.633600in}}%
\pgfpathlineto{\pgfqpoint{2.524969in}{0.643291in}}%
\pgfpathquadraticcurveto{\pgfqpoint{2.521204in}{0.637545in}}{\pgfqpoint{2.516215in}{0.634753in}}%
\pgfpathquadraticcurveto{\pgfqpoint{2.511243in}{0.631961in}}{\pgfqpoint{2.504651in}{0.631961in}}%
\pgfpathquadraticcurveto{\pgfqpoint{2.493790in}{0.631961in}}{\pgfqpoint{2.488152in}{0.638715in}}%
\pgfpathquadraticcurveto{\pgfqpoint{2.482532in}{0.645470in}}{\pgfqpoint{2.482532in}{0.658475in}}%
\pgfpathlineto{\pgfqpoint{2.482532in}{0.658475in}}%
\pgfpathclose%
\pgfpathmoveto{\pgfqpoint{2.508596in}{0.698156in}}%
\pgfpathlineto{\pgfqpoint{2.508596in}{0.698156in}}%
\pgfpathlineto{\pgfqpoint{2.508596in}{0.698156in}}%
\pgfpathclose%
\pgfpathmoveto{\pgfqpoint{2.566901in}{0.714547in}}%
\pgfpathlineto{\pgfqpoint{2.566901in}{0.696643in}}%
\pgfpathlineto{\pgfqpoint{2.588227in}{0.696643in}}%
\pgfpathlineto{\pgfqpoint{2.588227in}{0.688591in}}%
\pgfpathlineto{\pgfqpoint{2.566901in}{0.688591in}}%
\pgfpathlineto{\pgfqpoint{2.566901in}{0.654368in}}%
\pgfpathquadraticcurveto{\pgfqpoint{2.566901in}{0.646659in}}{\pgfqpoint{2.569008in}{0.644461in}}%
\pgfpathquadraticcurveto{\pgfqpoint{2.571116in}{0.642264in}}{\pgfqpoint{2.577600in}{0.642264in}}%
\pgfpathlineto{\pgfqpoint{2.588227in}{0.642264in}}%
\pgfpathlineto{\pgfqpoint{2.588227in}{0.633600in}}%
\pgfpathlineto{\pgfqpoint{2.577600in}{0.633600in}}%
\pgfpathquadraticcurveto{\pgfqpoint{2.565604in}{0.633600in}}{\pgfqpoint{2.561047in}{0.638067in}}%
\pgfpathquadraticcurveto{\pgfqpoint{2.556490in}{0.642552in}}{\pgfqpoint{2.556490in}{0.654368in}}%
\pgfpathlineto{\pgfqpoint{2.556490in}{0.688591in}}%
\pgfpathlineto{\pgfqpoint{2.548889in}{0.688591in}}%
\pgfpathlineto{\pgfqpoint{2.548889in}{0.696643in}}%
\pgfpathlineto{\pgfqpoint{2.556490in}{0.696643in}}%
\pgfpathlineto{\pgfqpoint{2.556490in}{0.714547in}}%
\pgfpathlineto{\pgfqpoint{2.566901in}{0.714547in}}%
\pgfpathlineto{\pgfqpoint{2.566901in}{0.714547in}}%
\pgfpathclose%
\pgfpathmoveto{\pgfqpoint{2.647217in}{0.694229in}}%
\pgfpathlineto{\pgfqpoint{2.647217in}{0.684538in}}%
\pgfpathquadraticcurveto{\pgfqpoint{2.642822in}{0.686970in}}{\pgfqpoint{2.638409in}{0.688177in}}%
\pgfpathquadraticcurveto{\pgfqpoint{2.633996in}{0.689384in}}{\pgfqpoint{2.629493in}{0.689384in}}%
\pgfpathquadraticcurveto{\pgfqpoint{2.619406in}{0.689384in}}{\pgfqpoint{2.613823in}{0.682989in}}%
\pgfpathquadraticcurveto{\pgfqpoint{2.608257in}{0.676613in}}{\pgfqpoint{2.608257in}{0.665067in}}%
\pgfpathquadraticcurveto{\pgfqpoint{2.608257in}{0.653521in}}{\pgfqpoint{2.613823in}{0.647127in}}%
\pgfpathquadraticcurveto{\pgfqpoint{2.619406in}{0.640751in}}{\pgfqpoint{2.629493in}{0.640751in}}%
\pgfpathquadraticcurveto{\pgfqpoint{2.633996in}{0.640751in}}{\pgfqpoint{2.638409in}{0.641958in}}%
\pgfpathquadraticcurveto{\pgfqpoint{2.642822in}{0.643164in}}{\pgfqpoint{2.647217in}{0.645596in}}%
\pgfpathlineto{\pgfqpoint{2.647217in}{0.636014in}}%
\pgfpathquadraticcurveto{\pgfqpoint{2.642876in}{0.633996in}}{\pgfqpoint{2.638229in}{0.632988in}}%
\pgfpathquadraticcurveto{\pgfqpoint{2.633600in}{0.631961in}}{\pgfqpoint{2.628359in}{0.631961in}}%
\pgfpathquadraticcurveto{\pgfqpoint{2.614111in}{0.631961in}}{\pgfqpoint{2.605717in}{0.640913in}}%
\pgfpathquadraticcurveto{\pgfqpoint{2.597342in}{0.649865in}}{\pgfqpoint{2.597342in}{0.665067in}}%
\pgfpathquadraticcurveto{\pgfqpoint{2.597342in}{0.680486in}}{\pgfqpoint{2.605807in}{0.689312in}}%
\pgfpathquadraticcurveto{\pgfqpoint{2.614291in}{0.698156in}}{\pgfqpoint{2.629043in}{0.698156in}}%
\pgfpathquadraticcurveto{\pgfqpoint{2.633816in}{0.698156in}}{\pgfqpoint{2.638373in}{0.697165in}}%
\pgfpathquadraticcurveto{\pgfqpoint{2.642930in}{0.696192in}}{\pgfqpoint{2.647217in}{0.694229in}}%
\pgfpathlineto{\pgfqpoint{2.647217in}{0.694229in}}%
\pgfpathclose%
\pgfpathmoveto{\pgfqpoint{2.689654in}{0.689384in}}%
\pgfpathquadraticcurveto{\pgfqpoint{2.681332in}{0.689384in}}{\pgfqpoint{2.676487in}{0.682881in}}%
\pgfpathquadraticcurveto{\pgfqpoint{2.671642in}{0.676379in}}{\pgfqpoint{2.671642in}{0.665067in}}%
\pgfpathquadraticcurveto{\pgfqpoint{2.671642in}{0.653756in}}{\pgfqpoint{2.676451in}{0.647253in}}%
\pgfpathquadraticcurveto{\pgfqpoint{2.681278in}{0.640751in}}{\pgfqpoint{2.689654in}{0.640751in}}%
\pgfpathquadraticcurveto{\pgfqpoint{2.697939in}{0.640751in}}{\pgfqpoint{2.702767in}{0.647271in}}%
\pgfpathquadraticcurveto{\pgfqpoint{2.707612in}{0.653810in}}{\pgfqpoint{2.707612in}{0.665067in}}%
\pgfpathquadraticcurveto{\pgfqpoint{2.707612in}{0.676271in}}{\pgfqpoint{2.702767in}{0.682827in}}%
\pgfpathquadraticcurveto{\pgfqpoint{2.697939in}{0.689384in}}{\pgfqpoint{2.689654in}{0.689384in}}%
\pgfpathlineto{\pgfqpoint{2.689654in}{0.689384in}}%
\pgfpathclose%
\pgfpathmoveto{\pgfqpoint{2.689654in}{0.698156in}}%
\pgfpathquadraticcurveto{\pgfqpoint{2.703163in}{0.698156in}}{\pgfqpoint{2.710872in}{0.689366in}}%
\pgfpathquadraticcurveto{\pgfqpoint{2.718599in}{0.680594in}}{\pgfqpoint{2.718599in}{0.665067in}}%
\pgfpathquadraticcurveto{\pgfqpoint{2.718599in}{0.649595in}}{\pgfqpoint{2.710872in}{0.640769in}}%
\pgfpathquadraticcurveto{\pgfqpoint{2.703163in}{0.631961in}}{\pgfqpoint{2.689654in}{0.631961in}}%
\pgfpathquadraticcurveto{\pgfqpoint{2.676091in}{0.631961in}}{\pgfqpoint{2.668400in}{0.640769in}}%
\pgfpathquadraticcurveto{\pgfqpoint{2.660726in}{0.649595in}}{\pgfqpoint{2.660726in}{0.665067in}}%
\pgfpathquadraticcurveto{\pgfqpoint{2.660726in}{0.680594in}}{\pgfqpoint{2.668400in}{0.689366in}}%
\pgfpathquadraticcurveto{\pgfqpoint{2.676091in}{0.698156in}}{\pgfqpoint{2.689654in}{0.698156in}}%
\pgfpathlineto{\pgfqpoint{2.689654in}{0.698156in}}%
\pgfpathclose%
\pgfpathmoveto{\pgfqpoint{2.784848in}{0.684538in}}%
\pgfpathquadraticcurveto{\pgfqpoint{2.788739in}{0.691527in}}{\pgfqpoint{2.794142in}{0.694841in}}%
\pgfpathquadraticcurveto{\pgfqpoint{2.799546in}{0.698156in}}{\pgfqpoint{2.806859in}{0.698156in}}%
\pgfpathquadraticcurveto{\pgfqpoint{2.816712in}{0.698156in}}{\pgfqpoint{2.822061in}{0.691257in}}%
\pgfpathquadraticcurveto{\pgfqpoint{2.827411in}{0.684376in}}{\pgfqpoint{2.827411in}{0.671660in}}%
\pgfpathlineto{\pgfqpoint{2.827411in}{0.633600in}}%
\pgfpathlineto{\pgfqpoint{2.817000in}{0.633600in}}%
\pgfpathlineto{\pgfqpoint{2.817000in}{0.671317in}}%
\pgfpathquadraticcurveto{\pgfqpoint{2.817000in}{0.680378in}}{\pgfqpoint{2.813776in}{0.684755in}}%
\pgfpathquadraticcurveto{\pgfqpoint{2.810569in}{0.689149in}}{\pgfqpoint{2.803995in}{0.689149in}}%
\pgfpathquadraticcurveto{\pgfqpoint{2.795944in}{0.689149in}}{\pgfqpoint{2.791260in}{0.683800in}}%
\pgfpathquadraticcurveto{\pgfqpoint{2.786595in}{0.678468in}}{\pgfqpoint{2.786595in}{0.669228in}}%
\pgfpathlineto{\pgfqpoint{2.786595in}{0.633600in}}%
\pgfpathlineto{\pgfqpoint{2.776184in}{0.633600in}}%
\pgfpathlineto{\pgfqpoint{2.776184in}{0.671317in}}%
\pgfpathquadraticcurveto{\pgfqpoint{2.776184in}{0.680432in}}{\pgfqpoint{2.772978in}{0.684791in}}%
\pgfpathquadraticcurveto{\pgfqpoint{2.769772in}{0.689149in}}{\pgfqpoint{2.763071in}{0.689149in}}%
\pgfpathquadraticcurveto{\pgfqpoint{2.755128in}{0.689149in}}{\pgfqpoint{2.750445in}{0.683782in}}%
\pgfpathquadraticcurveto{\pgfqpoint{2.745780in}{0.678414in}}{\pgfqpoint{2.745780in}{0.669228in}}%
\pgfpathlineto{\pgfqpoint{2.745780in}{0.633600in}}%
\pgfpathlineto{\pgfqpoint{2.735369in}{0.633600in}}%
\pgfpathlineto{\pgfqpoint{2.735369in}{0.696643in}}%
\pgfpathlineto{\pgfqpoint{2.745780in}{0.696643in}}%
\pgfpathlineto{\pgfqpoint{2.745780in}{0.686844in}}%
\pgfpathquadraticcurveto{\pgfqpoint{2.749328in}{0.692644in}}{\pgfqpoint{2.754281in}{0.695400in}}%
\pgfpathquadraticcurveto{\pgfqpoint{2.759235in}{0.698156in}}{\pgfqpoint{2.766043in}{0.698156in}}%
\pgfpathquadraticcurveto{\pgfqpoint{2.772924in}{0.698156in}}{\pgfqpoint{2.777733in}{0.694661in}}%
\pgfpathquadraticcurveto{\pgfqpoint{2.782543in}{0.691185in}}{\pgfqpoint{2.784848in}{0.684538in}}%
\pgfpathlineto{\pgfqpoint{2.784848in}{0.684538in}}%
\pgfpathclose%
\pgfpathmoveto{\pgfqpoint{2.901981in}{0.667715in}}%
\pgfpathlineto{\pgfqpoint{2.901981in}{0.662654in}}%
\pgfpathlineto{\pgfqpoint{2.854357in}{0.662654in}}%
\pgfpathquadraticcurveto{\pgfqpoint{2.855041in}{0.651954in}}{\pgfqpoint{2.860805in}{0.646353in}}%
\pgfpathquadraticcurveto{\pgfqpoint{2.866569in}{0.640751in}}{\pgfqpoint{2.876872in}{0.640751in}}%
\pgfpathquadraticcurveto{\pgfqpoint{2.882834in}{0.640751in}}{\pgfqpoint{2.888436in}{0.642210in}}%
\pgfpathquadraticcurveto{\pgfqpoint{2.894038in}{0.643669in}}{\pgfqpoint{2.899567in}{0.646605in}}%
\pgfpathlineto{\pgfqpoint{2.899567in}{0.636806in}}%
\pgfpathquadraticcurveto{\pgfqpoint{2.893984in}{0.634447in}}{\pgfqpoint{2.888130in}{0.633204in}}%
\pgfpathquadraticcurveto{\pgfqpoint{2.882276in}{0.631961in}}{\pgfqpoint{2.876260in}{0.631961in}}%
\pgfpathquadraticcurveto{\pgfqpoint{2.861166in}{0.631961in}}{\pgfqpoint{2.852358in}{0.640733in}}%
\pgfpathquadraticcurveto{\pgfqpoint{2.843550in}{0.649523in}}{\pgfqpoint{2.843550in}{0.664509in}}%
\pgfpathquadraticcurveto{\pgfqpoint{2.843550in}{0.679981in}}{\pgfqpoint{2.851907in}{0.689059in}}%
\pgfpathquadraticcurveto{\pgfqpoint{2.860265in}{0.698156in}}{\pgfqpoint{2.874459in}{0.698156in}}%
\pgfpathquadraticcurveto{\pgfqpoint{2.887175in}{0.698156in}}{\pgfqpoint{2.894578in}{0.689960in}}%
\pgfpathquadraticcurveto{\pgfqpoint{2.901981in}{0.681783in}}{\pgfqpoint{2.901981in}{0.667715in}}%
\pgfpathlineto{\pgfqpoint{2.901981in}{0.667715in}}%
\pgfpathclose%
\pgfpathmoveto{\pgfqpoint{2.891624in}{0.670759in}}%
\pgfpathquadraticcurveto{\pgfqpoint{2.891516in}{0.679243in}}{\pgfqpoint{2.886869in}{0.684304in}}%
\pgfpathquadraticcurveto{\pgfqpoint{2.882222in}{0.689384in}}{\pgfqpoint{2.874567in}{0.689384in}}%
\pgfpathquadraticcurveto{\pgfqpoint{2.865903in}{0.689384in}}{\pgfqpoint{2.860697in}{0.684484in}}%
\pgfpathquadraticcurveto{\pgfqpoint{2.855492in}{0.679585in}}{\pgfqpoint{2.854699in}{0.670687in}}%
\pgfpathlineto{\pgfqpoint{2.891624in}{0.670759in}}%
\pgfpathlineto{\pgfqpoint{2.891624in}{0.670759in}}%
\pgfpathclose%
\pgfpathmoveto{\pgfqpoint{2.913761in}{0.669786in}}%
\pgfpathlineto{\pgfqpoint{2.944094in}{0.669786in}}%
\pgfpathlineto{\pgfqpoint{2.944094in}{0.660564in}}%
\pgfpathlineto{\pgfqpoint{2.913761in}{0.660564in}}%
\pgfpathlineto{\pgfqpoint{2.913761in}{0.669786in}}%
\pgfpathlineto{\pgfqpoint{2.913761in}{0.669786in}}%
\pgfpathclose%
\pgfpathmoveto{\pgfqpoint{2.960575in}{0.721193in}}%
\pgfpathlineto{\pgfqpoint{2.970932in}{0.721193in}}%
\pgfpathlineto{\pgfqpoint{2.970932in}{0.633600in}}%
\pgfpathlineto{\pgfqpoint{2.960575in}{0.633600in}}%
\pgfpathlineto{\pgfqpoint{2.960575in}{0.721193in}}%
\pgfpathlineto{\pgfqpoint{2.960575in}{0.721193in}}%
\pgfpathclose%
\pgfpathmoveto{\pgfqpoint{2.992600in}{0.696643in}}%
\pgfpathlineto{\pgfqpoint{3.002957in}{0.696643in}}%
\pgfpathlineto{\pgfqpoint{3.002957in}{0.633600in}}%
\pgfpathlineto{\pgfqpoint{2.992600in}{0.633600in}}%
\pgfpathlineto{\pgfqpoint{2.992600in}{0.696643in}}%
\pgfpathlineto{\pgfqpoint{2.992600in}{0.696643in}}%
\pgfpathclose%
\pgfpathmoveto{\pgfqpoint{2.992600in}{0.721193in}}%
\pgfpathlineto{\pgfqpoint{3.002957in}{0.721193in}}%
\pgfpathlineto{\pgfqpoint{3.002957in}{0.708062in}}%
\pgfpathlineto{\pgfqpoint{2.992600in}{0.708062in}}%
\pgfpathlineto{\pgfqpoint{2.992600in}{0.721193in}}%
\pgfpathlineto{\pgfqpoint{2.992600in}{0.721193in}}%
\pgfpathclose%
\pgfpathmoveto{\pgfqpoint{3.077041in}{0.671660in}}%
\pgfpathlineto{\pgfqpoint{3.077041in}{0.633600in}}%
\pgfpathlineto{\pgfqpoint{3.066684in}{0.633600in}}%
\pgfpathlineto{\pgfqpoint{3.066684in}{0.671317in}}%
\pgfpathquadraticcurveto{\pgfqpoint{3.066684in}{0.680269in}}{\pgfqpoint{3.063190in}{0.684700in}}%
\pgfpathquadraticcurveto{\pgfqpoint{3.059696in}{0.689149in}}{\pgfqpoint{3.052725in}{0.689149in}}%
\pgfpathquadraticcurveto{\pgfqpoint{3.044331in}{0.689149in}}{\pgfqpoint{3.039486in}{0.683800in}}%
\pgfpathquadraticcurveto{\pgfqpoint{3.034641in}{0.678468in}}{\pgfqpoint{3.034641in}{0.669228in}}%
\pgfpathlineto{\pgfqpoint{3.034641in}{0.633600in}}%
\pgfpathlineto{\pgfqpoint{3.024230in}{0.633600in}}%
\pgfpathlineto{\pgfqpoint{3.024230in}{0.696643in}}%
\pgfpathlineto{\pgfqpoint{3.034641in}{0.696643in}}%
\pgfpathlineto{\pgfqpoint{3.034641in}{0.686844in}}%
\pgfpathquadraticcurveto{\pgfqpoint{3.038369in}{0.692536in}}{\pgfqpoint{3.043395in}{0.695346in}}%
\pgfpathquadraticcurveto{\pgfqpoint{3.048438in}{0.698156in}}{\pgfqpoint{3.055030in}{0.698156in}}%
\pgfpathquadraticcurveto{\pgfqpoint{3.065892in}{0.698156in}}{\pgfqpoint{3.071457in}{0.691437in}}%
\pgfpathquadraticcurveto{\pgfqpoint{3.077041in}{0.684718in}}{\pgfqpoint{3.077041in}{0.671660in}}%
\pgfpathlineto{\pgfqpoint{3.077041in}{0.671660in}}%
\pgfpathclose%
\pgfpathmoveto{\pgfqpoint{3.151612in}{0.667715in}}%
\pgfpathlineto{\pgfqpoint{3.151612in}{0.662654in}}%
\pgfpathlineto{\pgfqpoint{3.103987in}{0.662654in}}%
\pgfpathquadraticcurveto{\pgfqpoint{3.104672in}{0.651954in}}{\pgfqpoint{3.110436in}{0.646353in}}%
\pgfpathquadraticcurveto{\pgfqpoint{3.116200in}{0.640751in}}{\pgfqpoint{3.126503in}{0.640751in}}%
\pgfpathquadraticcurveto{\pgfqpoint{3.132465in}{0.640751in}}{\pgfqpoint{3.138066in}{0.642210in}}%
\pgfpathquadraticcurveto{\pgfqpoint{3.143668in}{0.643669in}}{\pgfqpoint{3.149198in}{0.646605in}}%
\pgfpathlineto{\pgfqpoint{3.149198in}{0.636806in}}%
\pgfpathquadraticcurveto{\pgfqpoint{3.143614in}{0.634447in}}{\pgfqpoint{3.137760in}{0.633204in}}%
\pgfpathquadraticcurveto{\pgfqpoint{3.131906in}{0.631961in}}{\pgfqpoint{3.125890in}{0.631961in}}%
\pgfpathquadraticcurveto{\pgfqpoint{3.110796in}{0.631961in}}{\pgfqpoint{3.101988in}{0.640733in}}%
\pgfpathquadraticcurveto{\pgfqpoint{3.093180in}{0.649523in}}{\pgfqpoint{3.093180in}{0.664509in}}%
\pgfpathquadraticcurveto{\pgfqpoint{3.093180in}{0.679981in}}{\pgfqpoint{3.101538in}{0.689059in}}%
\pgfpathquadraticcurveto{\pgfqpoint{3.109895in}{0.698156in}}{\pgfqpoint{3.124089in}{0.698156in}}%
\pgfpathquadraticcurveto{\pgfqpoint{3.136806in}{0.698156in}}{\pgfqpoint{3.144209in}{0.689960in}}%
\pgfpathquadraticcurveto{\pgfqpoint{3.151612in}{0.681783in}}{\pgfqpoint{3.151612in}{0.667715in}}%
\pgfpathlineto{\pgfqpoint{3.151612in}{0.667715in}}%
\pgfpathclose%
\pgfpathmoveto{\pgfqpoint{3.141255in}{0.670759in}}%
\pgfpathquadraticcurveto{\pgfqpoint{3.141146in}{0.679243in}}{\pgfqpoint{3.136499in}{0.684304in}}%
\pgfpathquadraticcurveto{\pgfqpoint{3.131852in}{0.689384in}}{\pgfqpoint{3.124197in}{0.689384in}}%
\pgfpathquadraticcurveto{\pgfqpoint{3.115533in}{0.689384in}}{\pgfqpoint{3.110328in}{0.684484in}}%
\pgfpathquadraticcurveto{\pgfqpoint{3.105122in}{0.679585in}}{\pgfqpoint{3.104330in}{0.670687in}}%
\pgfpathlineto{\pgfqpoint{3.141255in}{0.670759in}}%
\pgfpathlineto{\pgfqpoint{3.141255in}{0.670759in}}%
\pgfpathclose%
\pgfpathmoveto{\pgfqpoint{3.197272in}{0.665283in}}%
\pgfpathquadraticcurveto{\pgfqpoint{3.184718in}{0.665283in}}{\pgfqpoint{3.179873in}{0.662419in}}%
\pgfpathquadraticcurveto{\pgfqpoint{3.175027in}{0.659556in}}{\pgfqpoint{3.175027in}{0.652621in}}%
\pgfpathquadraticcurveto{\pgfqpoint{3.175027in}{0.647109in}}{\pgfqpoint{3.178666in}{0.643867in}}%
\pgfpathquadraticcurveto{\pgfqpoint{3.182304in}{0.640643in}}{\pgfqpoint{3.188536in}{0.640643in}}%
\pgfpathquadraticcurveto{\pgfqpoint{3.197164in}{0.640643in}}{\pgfqpoint{3.202370in}{0.646749in}}%
\pgfpathquadraticcurveto{\pgfqpoint{3.207575in}{0.652855in}}{\pgfqpoint{3.207575in}{0.662978in}}%
\pgfpathlineto{\pgfqpoint{3.207575in}{0.665283in}}%
\pgfpathlineto{\pgfqpoint{3.197272in}{0.665283in}}%
\pgfpathlineto{\pgfqpoint{3.197272in}{0.665283in}}%
\pgfpathclose%
\pgfpathmoveto{\pgfqpoint{3.217932in}{0.669570in}}%
\pgfpathlineto{\pgfqpoint{3.217932in}{0.633600in}}%
\pgfpathlineto{\pgfqpoint{3.207575in}{0.633600in}}%
\pgfpathlineto{\pgfqpoint{3.207575in}{0.643164in}}%
\pgfpathquadraticcurveto{\pgfqpoint{3.204027in}{0.637437in}}{\pgfqpoint{3.198731in}{0.634699in}}%
\pgfpathquadraticcurveto{\pgfqpoint{3.193436in}{0.631961in}}{\pgfqpoint{3.185781in}{0.631961in}}%
\pgfpathquadraticcurveto{\pgfqpoint{3.176108in}{0.631961in}}{\pgfqpoint{3.170380in}{0.637401in}}%
\pgfpathquadraticcurveto{\pgfqpoint{3.164670in}{0.642840in}}{\pgfqpoint{3.164670in}{0.651954in}}%
\pgfpathquadraticcurveto{\pgfqpoint{3.164670in}{0.662582in}}{\pgfqpoint{3.171785in}{0.667985in}}%
\pgfpathquadraticcurveto{\pgfqpoint{3.178918in}{0.673389in}}{\pgfqpoint{3.193039in}{0.673389in}}%
\pgfpathlineto{\pgfqpoint{3.207575in}{0.673389in}}%
\pgfpathlineto{\pgfqpoint{3.207575in}{0.674416in}}%
\pgfpathquadraticcurveto{\pgfqpoint{3.207575in}{0.681566in}}{\pgfqpoint{3.202874in}{0.685475in}}%
\pgfpathquadraticcurveto{\pgfqpoint{3.198173in}{0.689384in}}{\pgfqpoint{3.189671in}{0.689384in}}%
\pgfpathquadraticcurveto{\pgfqpoint{3.184268in}{0.689384in}}{\pgfqpoint{3.179134in}{0.688087in}}%
\pgfpathquadraticcurveto{\pgfqpoint{3.174019in}{0.686790in}}{\pgfqpoint{3.169299in}{0.684196in}}%
\pgfpathlineto{\pgfqpoint{3.169299in}{0.693779in}}%
\pgfpathquadraticcurveto{\pgfqpoint{3.174973in}{0.695976in}}{\pgfqpoint{3.180323in}{0.697057in}}%
\pgfpathquadraticcurveto{\pgfqpoint{3.185673in}{0.698156in}}{\pgfqpoint{3.190734in}{0.698156in}}%
\pgfpathquadraticcurveto{\pgfqpoint{3.204423in}{0.698156in}}{\pgfqpoint{3.211178in}{0.691059in}}%
\pgfpathquadraticcurveto{\pgfqpoint{3.217932in}{0.683980in}}{\pgfqpoint{3.217932in}{0.669570in}}%
\pgfpathlineto{\pgfqpoint{3.217932in}{0.669570in}}%
\pgfpathclose%
\pgfpathmoveto{\pgfqpoint{3.275787in}{0.686970in}}%
\pgfpathquadraticcurveto{\pgfqpoint{3.274040in}{0.687979in}}{\pgfqpoint{3.271987in}{0.688447in}}%
\pgfpathquadraticcurveto{\pgfqpoint{3.269933in}{0.688933in}}{\pgfqpoint{3.267466in}{0.688933in}}%
\pgfpathquadraticcurveto{\pgfqpoint{3.258676in}{0.688933in}}{\pgfqpoint{3.253975in}{0.683223in}}%
\pgfpathquadraticcurveto{\pgfqpoint{3.249273in}{0.677514in}}{\pgfqpoint{3.249273in}{0.666814in}}%
\pgfpathlineto{\pgfqpoint{3.249273in}{0.633600in}}%
\pgfpathlineto{\pgfqpoint{3.238862in}{0.633600in}}%
\pgfpathlineto{\pgfqpoint{3.238862in}{0.696643in}}%
\pgfpathlineto{\pgfqpoint{3.249273in}{0.696643in}}%
\pgfpathlineto{\pgfqpoint{3.249273in}{0.686844in}}%
\pgfpathquadraticcurveto{\pgfqpoint{3.252552in}{0.692590in}}{\pgfqpoint{3.257775in}{0.695364in}}%
\pgfpathquadraticcurveto{\pgfqpoint{3.263017in}{0.698156in}}{\pgfqpoint{3.270510in}{0.698156in}}%
\pgfpathquadraticcurveto{\pgfqpoint{3.271572in}{0.698156in}}{\pgfqpoint{3.272869in}{0.698011in}}%
\pgfpathquadraticcurveto{\pgfqpoint{3.274166in}{0.697885in}}{\pgfqpoint{3.275733in}{0.697597in}}%
\pgfpathlineto{\pgfqpoint{3.275787in}{0.686970in}}%
\pgfpathlineto{\pgfqpoint{3.275787in}{0.686970in}}%
\pgfpathclose%
\pgfusepath{fill}%
\end{pgfscope}%
\begin{pgfscope}%
\pgfsetbuttcap%
\pgfsetroundjoin%
\definecolor{currentfill}{rgb}{0.309804,0.372549,0.419608}%
\pgfsetfillcolor{currentfill}%
\pgfsetlinewidth{0.000000pt}%
\definecolor{currentstroke}{rgb}{0.309804,0.372549,0.419608}%
\pgfsetstrokecolor{currentstroke}%
\pgfsetdash{}{0pt}%
\pgfpathmoveto{\pgfqpoint{3.364620in}{0.696643in}}%
\pgfpathlineto{\pgfqpoint{3.374977in}{0.696643in}}%
\pgfpathlineto{\pgfqpoint{3.387927in}{0.647451in}}%
\pgfpathlineto{\pgfqpoint{3.400806in}{0.696643in}}%
\pgfpathlineto{\pgfqpoint{3.413018in}{0.696643in}}%
\pgfpathlineto{\pgfqpoint{3.425969in}{0.647451in}}%
\pgfpathlineto{\pgfqpoint{3.438866in}{0.696643in}}%
\pgfpathlineto{\pgfqpoint{3.449223in}{0.696643in}}%
\pgfpathlineto{\pgfqpoint{3.432724in}{0.633600in}}%
\pgfpathlineto{\pgfqpoint{3.420511in}{0.633600in}}%
\pgfpathlineto{\pgfqpoint{3.406948in}{0.685277in}}%
\pgfpathlineto{\pgfqpoint{3.393331in}{0.633600in}}%
\pgfpathlineto{\pgfqpoint{3.381101in}{0.633600in}}%
\pgfpathlineto{\pgfqpoint{3.364620in}{0.696643in}}%
\pgfpathlineto{\pgfqpoint{3.364620in}{0.696643in}}%
\pgfpathclose%
\pgfpathmoveto{\pgfqpoint{3.518840in}{0.667715in}}%
\pgfpathlineto{\pgfqpoint{3.518840in}{0.662654in}}%
\pgfpathlineto{\pgfqpoint{3.471215in}{0.662654in}}%
\pgfpathquadraticcurveto{\pgfqpoint{3.471900in}{0.651954in}}{\pgfqpoint{3.477664in}{0.646353in}}%
\pgfpathquadraticcurveto{\pgfqpoint{3.483428in}{0.640751in}}{\pgfqpoint{3.493731in}{0.640751in}}%
\pgfpathquadraticcurveto{\pgfqpoint{3.499693in}{0.640751in}}{\pgfqpoint{3.505294in}{0.642210in}}%
\pgfpathquadraticcurveto{\pgfqpoint{3.510896in}{0.643669in}}{\pgfqpoint{3.516426in}{0.646605in}}%
\pgfpathlineto{\pgfqpoint{3.516426in}{0.636806in}}%
\pgfpathquadraticcurveto{\pgfqpoint{3.510842in}{0.634447in}}{\pgfqpoint{3.504988in}{0.633204in}}%
\pgfpathquadraticcurveto{\pgfqpoint{3.499134in}{0.631961in}}{\pgfqpoint{3.493118in}{0.631961in}}%
\pgfpathquadraticcurveto{\pgfqpoint{3.478024in}{0.631961in}}{\pgfqpoint{3.469216in}{0.640733in}}%
\pgfpathquadraticcurveto{\pgfqpoint{3.460408in}{0.649523in}}{\pgfqpoint{3.460408in}{0.664509in}}%
\pgfpathquadraticcurveto{\pgfqpoint{3.460408in}{0.679981in}}{\pgfqpoint{3.468766in}{0.689059in}}%
\pgfpathquadraticcurveto{\pgfqpoint{3.477123in}{0.698156in}}{\pgfqpoint{3.491317in}{0.698156in}}%
\pgfpathquadraticcurveto{\pgfqpoint{3.504034in}{0.698156in}}{\pgfqpoint{3.511437in}{0.689960in}}%
\pgfpathquadraticcurveto{\pgfqpoint{3.518840in}{0.681783in}}{\pgfqpoint{3.518840in}{0.667715in}}%
\pgfpathlineto{\pgfqpoint{3.518840in}{0.667715in}}%
\pgfpathclose%
\pgfpathmoveto{\pgfqpoint{3.508483in}{0.670759in}}%
\pgfpathquadraticcurveto{\pgfqpoint{3.508375in}{0.679243in}}{\pgfqpoint{3.503727in}{0.684304in}}%
\pgfpathquadraticcurveto{\pgfqpoint{3.499080in}{0.689384in}}{\pgfqpoint{3.491425in}{0.689384in}}%
\pgfpathquadraticcurveto{\pgfqpoint{3.482761in}{0.689384in}}{\pgfqpoint{3.477556in}{0.684484in}}%
\pgfpathquadraticcurveto{\pgfqpoint{3.472350in}{0.679585in}}{\pgfqpoint{3.471558in}{0.670687in}}%
\pgfpathlineto{\pgfqpoint{3.508483in}{0.670759in}}%
\pgfpathlineto{\pgfqpoint{3.508483in}{0.670759in}}%
\pgfpathclose%
\pgfpathmoveto{\pgfqpoint{3.535843in}{0.696643in}}%
\pgfpathlineto{\pgfqpoint{3.546200in}{0.696643in}}%
\pgfpathlineto{\pgfqpoint{3.546200in}{0.633600in}}%
\pgfpathlineto{\pgfqpoint{3.535843in}{0.633600in}}%
\pgfpathlineto{\pgfqpoint{3.535843in}{0.696643in}}%
\pgfpathlineto{\pgfqpoint{3.535843in}{0.696643in}}%
\pgfpathclose%
\pgfpathmoveto{\pgfqpoint{3.535843in}{0.721193in}}%
\pgfpathlineto{\pgfqpoint{3.546200in}{0.721193in}}%
\pgfpathlineto{\pgfqpoint{3.546200in}{0.708062in}}%
\pgfpathlineto{\pgfqpoint{3.535843in}{0.708062in}}%
\pgfpathlineto{\pgfqpoint{3.535843in}{0.721193in}}%
\pgfpathlineto{\pgfqpoint{3.535843in}{0.721193in}}%
\pgfpathclose%
\pgfpathmoveto{\pgfqpoint{3.609351in}{0.665860in}}%
\pgfpathquadraticcurveto{\pgfqpoint{3.609351in}{0.677117in}}{\pgfqpoint{3.604704in}{0.683296in}}%
\pgfpathquadraticcurveto{\pgfqpoint{3.600074in}{0.689492in}}{\pgfqpoint{3.591681in}{0.689492in}}%
\pgfpathquadraticcurveto{\pgfqpoint{3.583359in}{0.689492in}}{\pgfqpoint{3.578712in}{0.683296in}}%
\pgfpathquadraticcurveto{\pgfqpoint{3.574065in}{0.677117in}}{\pgfqpoint{3.574065in}{0.665860in}}%
\pgfpathquadraticcurveto{\pgfqpoint{3.574065in}{0.654656in}}{\pgfqpoint{3.578712in}{0.648460in}}%
\pgfpathquadraticcurveto{\pgfqpoint{3.583359in}{0.642264in}}{\pgfqpoint{3.591681in}{0.642264in}}%
\pgfpathquadraticcurveto{\pgfqpoint{3.600074in}{0.642264in}}{\pgfqpoint{3.604704in}{0.648460in}}%
\pgfpathquadraticcurveto{\pgfqpoint{3.609351in}{0.654656in}}{\pgfqpoint{3.609351in}{0.665860in}}%
\pgfpathlineto{\pgfqpoint{3.609351in}{0.665860in}}%
\pgfpathclose%
\pgfpathmoveto{\pgfqpoint{3.619708in}{0.641417in}}%
\pgfpathquadraticcurveto{\pgfqpoint{3.619708in}{0.625332in}}{\pgfqpoint{3.612557in}{0.617479in}}%
\pgfpathquadraticcurveto{\pgfqpoint{3.605424in}{0.609626in}}{\pgfqpoint{3.590672in}{0.609626in}}%
\pgfpathquadraticcurveto{\pgfqpoint{3.585214in}{0.609626in}}{\pgfqpoint{3.580369in}{0.610436in}}%
\pgfpathquadraticcurveto{\pgfqpoint{3.575524in}{0.611247in}}{\pgfqpoint{3.570967in}{0.612940in}}%
\pgfpathlineto{\pgfqpoint{3.570967in}{0.623009in}}%
\pgfpathquadraticcurveto{\pgfqpoint{3.575524in}{0.620541in}}{\pgfqpoint{3.579973in}{0.619370in}}%
\pgfpathquadraticcurveto{\pgfqpoint{3.584422in}{0.618182in}}{\pgfqpoint{3.589033in}{0.618182in}}%
\pgfpathquadraticcurveto{\pgfqpoint{3.599228in}{0.618182in}}{\pgfqpoint{3.604289in}{0.623495in}}%
\pgfpathquadraticcurveto{\pgfqpoint{3.609351in}{0.628809in}}{\pgfqpoint{3.609351in}{0.639562in}}%
\pgfpathlineto{\pgfqpoint{3.609351in}{0.644695in}}%
\pgfpathquadraticcurveto{\pgfqpoint{3.606145in}{0.639112in}}{\pgfqpoint{3.601137in}{0.636356in}}%
\pgfpathquadraticcurveto{\pgfqpoint{3.596130in}{0.633600in}}{\pgfqpoint{3.589141in}{0.633600in}}%
\pgfpathquadraticcurveto{\pgfqpoint{3.577559in}{0.633600in}}{\pgfqpoint{3.570462in}{0.642426in}}%
\pgfpathquadraticcurveto{\pgfqpoint{3.563366in}{0.651270in}}{\pgfqpoint{3.563366in}{0.665860in}}%
\pgfpathquadraticcurveto{\pgfqpoint{3.563366in}{0.680486in}}{\pgfqpoint{3.570462in}{0.689312in}}%
\pgfpathquadraticcurveto{\pgfqpoint{3.577559in}{0.698156in}}{\pgfqpoint{3.589141in}{0.698156in}}%
\pgfpathquadraticcurveto{\pgfqpoint{3.596130in}{0.698156in}}{\pgfqpoint{3.601137in}{0.695400in}}%
\pgfpathquadraticcurveto{\pgfqpoint{3.606145in}{0.692644in}}{\pgfqpoint{3.609351in}{0.687078in}}%
\pgfpathlineto{\pgfqpoint{3.609351in}{0.696643in}}%
\pgfpathlineto{\pgfqpoint{3.619708in}{0.696643in}}%
\pgfpathlineto{\pgfqpoint{3.619708in}{0.641417in}}%
\pgfpathlineto{\pgfqpoint{3.619708in}{0.641417in}}%
\pgfpathclose%
\pgfpathmoveto{\pgfqpoint{3.693467in}{0.671660in}}%
\pgfpathlineto{\pgfqpoint{3.693467in}{0.633600in}}%
\pgfpathlineto{\pgfqpoint{3.683110in}{0.633600in}}%
\pgfpathlineto{\pgfqpoint{3.683110in}{0.671317in}}%
\pgfpathquadraticcurveto{\pgfqpoint{3.683110in}{0.680269in}}{\pgfqpoint{3.679616in}{0.684700in}}%
\pgfpathquadraticcurveto{\pgfqpoint{3.676122in}{0.689149in}}{\pgfqpoint{3.669151in}{0.689149in}}%
\pgfpathquadraticcurveto{\pgfqpoint{3.660757in}{0.689149in}}{\pgfqpoint{3.655912in}{0.683800in}}%
\pgfpathquadraticcurveto{\pgfqpoint{3.651067in}{0.678468in}}{\pgfqpoint{3.651067in}{0.669228in}}%
\pgfpathlineto{\pgfqpoint{3.651067in}{0.633600in}}%
\pgfpathlineto{\pgfqpoint{3.640656in}{0.633600in}}%
\pgfpathlineto{\pgfqpoint{3.640656in}{0.721193in}}%
\pgfpathlineto{\pgfqpoint{3.651067in}{0.721193in}}%
\pgfpathlineto{\pgfqpoint{3.651067in}{0.686844in}}%
\pgfpathquadraticcurveto{\pgfqpoint{3.654795in}{0.692536in}}{\pgfqpoint{3.659821in}{0.695346in}}%
\pgfpathquadraticcurveto{\pgfqpoint{3.664864in}{0.698156in}}{\pgfqpoint{3.671457in}{0.698156in}}%
\pgfpathquadraticcurveto{\pgfqpoint{3.682318in}{0.698156in}}{\pgfqpoint{3.687884in}{0.691437in}}%
\pgfpathquadraticcurveto{\pgfqpoint{3.693467in}{0.684718in}}{\pgfqpoint{3.693467in}{0.671660in}}%
\pgfpathlineto{\pgfqpoint{3.693467in}{0.671660in}}%
\pgfpathclose%
\pgfpathmoveto{\pgfqpoint{3.724358in}{0.714547in}}%
\pgfpathlineto{\pgfqpoint{3.724358in}{0.696643in}}%
\pgfpathlineto{\pgfqpoint{3.745685in}{0.696643in}}%
\pgfpathlineto{\pgfqpoint{3.745685in}{0.688591in}}%
\pgfpathlineto{\pgfqpoint{3.724358in}{0.688591in}}%
\pgfpathlineto{\pgfqpoint{3.724358in}{0.654368in}}%
\pgfpathquadraticcurveto{\pgfqpoint{3.724358in}{0.646659in}}{\pgfqpoint{3.726466in}{0.644461in}}%
\pgfpathquadraticcurveto{\pgfqpoint{3.728573in}{0.642264in}}{\pgfqpoint{3.735058in}{0.642264in}}%
\pgfpathlineto{\pgfqpoint{3.745685in}{0.642264in}}%
\pgfpathlineto{\pgfqpoint{3.745685in}{0.633600in}}%
\pgfpathlineto{\pgfqpoint{3.735058in}{0.633600in}}%
\pgfpathquadraticcurveto{\pgfqpoint{3.723061in}{0.633600in}}{\pgfqpoint{3.718504in}{0.638067in}}%
\pgfpathquadraticcurveto{\pgfqpoint{3.713947in}{0.642552in}}{\pgfqpoint{3.713947in}{0.654368in}}%
\pgfpathlineto{\pgfqpoint{3.713947in}{0.688591in}}%
\pgfpathlineto{\pgfqpoint{3.706346in}{0.688591in}}%
\pgfpathlineto{\pgfqpoint{3.706346in}{0.696643in}}%
\pgfpathlineto{\pgfqpoint{3.713947in}{0.696643in}}%
\pgfpathlineto{\pgfqpoint{3.713947in}{0.714547in}}%
\pgfpathlineto{\pgfqpoint{3.724358in}{0.714547in}}%
\pgfpathlineto{\pgfqpoint{3.724358in}{0.714547in}}%
\pgfpathclose%
\pgfpathmoveto{\pgfqpoint{3.799487in}{0.694787in}}%
\pgfpathlineto{\pgfqpoint{3.799487in}{0.684989in}}%
\pgfpathquadraticcurveto{\pgfqpoint{3.795110in}{0.687240in}}{\pgfqpoint{3.790373in}{0.688357in}}%
\pgfpathquadraticcurveto{\pgfqpoint{3.785654in}{0.689492in}}{\pgfqpoint{3.780574in}{0.689492in}}%
\pgfpathquadraticcurveto{\pgfqpoint{3.772865in}{0.689492in}}{\pgfqpoint{3.769010in}{0.687132in}}%
\pgfpathquadraticcurveto{\pgfqpoint{3.765156in}{0.684773in}}{\pgfqpoint{3.765156in}{0.680035in}}%
\pgfpathquadraticcurveto{\pgfqpoint{3.765156in}{0.676433in}}{\pgfqpoint{3.767912in}{0.674380in}}%
\pgfpathquadraticcurveto{\pgfqpoint{3.770668in}{0.672326in}}{\pgfqpoint{3.779007in}{0.670471in}}%
\pgfpathlineto{\pgfqpoint{3.782556in}{0.669678in}}%
\pgfpathquadraticcurveto{\pgfqpoint{3.793579in}{0.667319in}}{\pgfqpoint{3.798226in}{0.663014in}}%
\pgfpathquadraticcurveto{\pgfqpoint{3.802873in}{0.658709in}}{\pgfqpoint{3.802873in}{0.651000in}}%
\pgfpathquadraticcurveto{\pgfqpoint{3.802873in}{0.642210in}}{\pgfqpoint{3.795921in}{0.637076in}}%
\pgfpathquadraticcurveto{\pgfqpoint{3.788968in}{0.631961in}}{\pgfqpoint{3.776810in}{0.631961in}}%
\pgfpathquadraticcurveto{\pgfqpoint{3.771748in}{0.631961in}}{\pgfqpoint{3.766255in}{0.632952in}}%
\pgfpathquadraticcurveto{\pgfqpoint{3.760761in}{0.633942in}}{\pgfqpoint{3.754691in}{0.635906in}}%
\pgfpathlineto{\pgfqpoint{3.754691in}{0.646605in}}%
\pgfpathquadraticcurveto{\pgfqpoint{3.760437in}{0.643615in}}{\pgfqpoint{3.766002in}{0.642120in}}%
\pgfpathquadraticcurveto{\pgfqpoint{3.771568in}{0.640643in}}{\pgfqpoint{3.777044in}{0.640643in}}%
\pgfpathquadraticcurveto{\pgfqpoint{3.784357in}{0.640643in}}{\pgfqpoint{3.788283in}{0.643146in}}%
\pgfpathquadraticcurveto{\pgfqpoint{3.792228in}{0.645650in}}{\pgfqpoint{3.792228in}{0.650207in}}%
\pgfpathquadraticcurveto{\pgfqpoint{3.792228in}{0.654422in}}{\pgfqpoint{3.789382in}{0.656674in}}%
\pgfpathquadraticcurveto{\pgfqpoint{3.786554in}{0.658925in}}{\pgfqpoint{3.776918in}{0.661014in}}%
\pgfpathlineto{\pgfqpoint{3.773315in}{0.661861in}}%
\pgfpathquadraticcurveto{\pgfqpoint{3.763697in}{0.663878in}}{\pgfqpoint{3.759410in}{0.668075in}}%
\pgfpathquadraticcurveto{\pgfqpoint{3.755141in}{0.672272in}}{\pgfqpoint{3.755141in}{0.679585in}}%
\pgfpathquadraticcurveto{\pgfqpoint{3.755141in}{0.688483in}}{\pgfqpoint{3.761445in}{0.693310in}}%
\pgfpathquadraticcurveto{\pgfqpoint{3.767750in}{0.698156in}}{\pgfqpoint{3.779349in}{0.698156in}}%
\pgfpathquadraticcurveto{\pgfqpoint{3.785077in}{0.698156in}}{\pgfqpoint{3.790139in}{0.697309in}}%
\pgfpathquadraticcurveto{\pgfqpoint{3.795218in}{0.696480in}}{\pgfqpoint{3.799487in}{0.694787in}}%
\pgfpathlineto{\pgfqpoint{3.799487in}{0.694787in}}%
\pgfpathclose%
\pgfusepath{fill}%
\end{pgfscope}%
\begin{pgfscope}%
\pgfsetbuttcap%
\pgfsetroundjoin%
\definecolor{currentfill}{rgb}{0.309804,0.372549,0.419608}%
\pgfsetfillcolor{currentfill}%
\pgfsetlinewidth{0.000000pt}%
\definecolor{currentstroke}{rgb}{0.309804,0.372549,0.419608}%
\pgfsetstrokecolor{currentstroke}%
\pgfsetdash{}{0pt}%
\pgfpathmoveto{\pgfqpoint{3.943477in}{0.665067in}}%
\pgfpathquadraticcurveto{\pgfqpoint{3.943477in}{0.676487in}}{\pgfqpoint{3.938775in}{0.682989in}}%
\pgfpathquadraticcurveto{\pgfqpoint{3.934074in}{0.689492in}}{\pgfqpoint{3.925861in}{0.689492in}}%
\pgfpathquadraticcurveto{\pgfqpoint{3.917629in}{0.689492in}}{\pgfqpoint{3.912928in}{0.682989in}}%
\pgfpathquadraticcurveto{\pgfqpoint{3.908227in}{0.676487in}}{\pgfqpoint{3.908227in}{0.665067in}}%
\pgfpathquadraticcurveto{\pgfqpoint{3.908227in}{0.653648in}}{\pgfqpoint{3.912928in}{0.647145in}}%
\pgfpathquadraticcurveto{\pgfqpoint{3.917629in}{0.640643in}}{\pgfqpoint{3.925861in}{0.640643in}}%
\pgfpathquadraticcurveto{\pgfqpoint{3.934074in}{0.640643in}}{\pgfqpoint{3.938775in}{0.647145in}}%
\pgfpathquadraticcurveto{\pgfqpoint{3.943477in}{0.653648in}}{\pgfqpoint{3.943477in}{0.665067in}}%
\pgfpathlineto{\pgfqpoint{3.943477in}{0.665067in}}%
\pgfpathclose%
\pgfpathmoveto{\pgfqpoint{3.908227in}{0.687078in}}%
\pgfpathquadraticcurveto{\pgfqpoint{3.911505in}{0.692698in}}{\pgfqpoint{3.916476in}{0.695418in}}%
\pgfpathquadraticcurveto{\pgfqpoint{3.921466in}{0.698156in}}{\pgfqpoint{3.928382in}{0.698156in}}%
\pgfpathquadraticcurveto{\pgfqpoint{3.939874in}{0.698156in}}{\pgfqpoint{3.947043in}{0.689041in}}%
\pgfpathquadraticcurveto{\pgfqpoint{3.954230in}{0.679927in}}{\pgfqpoint{3.954230in}{0.665067in}}%
\pgfpathquadraticcurveto{\pgfqpoint{3.954230in}{0.650207in}}{\pgfqpoint{3.947043in}{0.641075in}}%
\pgfpathquadraticcurveto{\pgfqpoint{3.939874in}{0.631961in}}{\pgfqpoint{3.928382in}{0.631961in}}%
\pgfpathquadraticcurveto{\pgfqpoint{3.921466in}{0.631961in}}{\pgfqpoint{3.916476in}{0.634699in}}%
\pgfpathquadraticcurveto{\pgfqpoint{3.911505in}{0.637437in}}{\pgfqpoint{3.908227in}{0.643056in}}%
\pgfpathlineto{\pgfqpoint{3.908227in}{0.633600in}}%
\pgfpathlineto{\pgfqpoint{3.897816in}{0.633600in}}%
\pgfpathlineto{\pgfqpoint{3.897816in}{0.721193in}}%
\pgfpathlineto{\pgfqpoint{3.908227in}{0.721193in}}%
\pgfpathlineto{\pgfqpoint{3.908227in}{0.687078in}}%
\pgfpathlineto{\pgfqpoint{3.908227in}{0.687078in}}%
\pgfpathclose%
\pgfpathmoveto{\pgfqpoint{3.997621in}{0.627746in}}%
\pgfpathquadraticcurveto{\pgfqpoint{3.993244in}{0.616488in}}{\pgfqpoint{3.989065in}{0.613066in}}%
\pgfpathquadraticcurveto{\pgfqpoint{3.984905in}{0.609626in}}{\pgfqpoint{3.977934in}{0.609626in}}%
\pgfpathlineto{\pgfqpoint{3.969648in}{0.609626in}}%
\pgfpathlineto{\pgfqpoint{3.969648in}{0.618290in}}%
\pgfpathlineto{\pgfqpoint{3.975736in}{0.618290in}}%
\pgfpathquadraticcurveto{\pgfqpoint{3.980005in}{0.618290in}}{\pgfqpoint{3.982365in}{0.620325in}}%
\pgfpathquadraticcurveto{\pgfqpoint{3.984742in}{0.622342in}}{\pgfqpoint{3.987606in}{0.629889in}}%
\pgfpathlineto{\pgfqpoint{3.989462in}{0.634609in}}%
\pgfpathlineto{\pgfqpoint{3.963974in}{0.696643in}}%
\pgfpathlineto{\pgfqpoint{3.974944in}{0.696643in}}%
\pgfpathlineto{\pgfqpoint{3.994649in}{0.647343in}}%
\pgfpathlineto{\pgfqpoint{4.014354in}{0.696643in}}%
\pgfpathlineto{\pgfqpoint{4.025324in}{0.696643in}}%
\pgfpathlineto{\pgfqpoint{3.997621in}{0.627746in}}%
\pgfpathlineto{\pgfqpoint{3.997621in}{0.627746in}}%
\pgfpathclose%
\pgfusepath{fill}%
\end{pgfscope}%
\begin{pgfscope}%
\pgfsetbuttcap%
\pgfsetroundjoin%
\definecolor{currentfill}{rgb}{0.309804,0.372549,0.419608}%
\pgfsetfillcolor{currentfill}%
\pgfsetlinewidth{0.000000pt}%
\definecolor{currentstroke}{rgb}{0.309804,0.372549,0.419608}%
\pgfsetstrokecolor{currentstroke}%
\pgfsetdash{}{0pt}%
\pgfpathmoveto{\pgfqpoint{4.116037in}{0.714547in}}%
\pgfpathlineto{\pgfqpoint{4.116037in}{0.696643in}}%
\pgfpathlineto{\pgfqpoint{4.137364in}{0.696643in}}%
\pgfpathlineto{\pgfqpoint{4.137364in}{0.688591in}}%
\pgfpathlineto{\pgfqpoint{4.116037in}{0.688591in}}%
\pgfpathlineto{\pgfqpoint{4.116037in}{0.654368in}}%
\pgfpathquadraticcurveto{\pgfqpoint{4.116037in}{0.646659in}}{\pgfqpoint{4.118145in}{0.644461in}}%
\pgfpathquadraticcurveto{\pgfqpoint{4.120252in}{0.642264in}}{\pgfqpoint{4.126737in}{0.642264in}}%
\pgfpathlineto{\pgfqpoint{4.137364in}{0.642264in}}%
\pgfpathlineto{\pgfqpoint{4.137364in}{0.633600in}}%
\pgfpathlineto{\pgfqpoint{4.126737in}{0.633600in}}%
\pgfpathquadraticcurveto{\pgfqpoint{4.114740in}{0.633600in}}{\pgfqpoint{4.110183in}{0.638067in}}%
\pgfpathquadraticcurveto{\pgfqpoint{4.105626in}{0.642552in}}{\pgfqpoint{4.105626in}{0.654368in}}%
\pgfpathlineto{\pgfqpoint{4.105626in}{0.688591in}}%
\pgfpathlineto{\pgfqpoint{4.098025in}{0.688591in}}%
\pgfpathlineto{\pgfqpoint{4.098025in}{0.696643in}}%
\pgfpathlineto{\pgfqpoint{4.105626in}{0.696643in}}%
\pgfpathlineto{\pgfqpoint{4.105626in}{0.714547in}}%
\pgfpathlineto{\pgfqpoint{4.116037in}{0.714547in}}%
\pgfpathlineto{\pgfqpoint{4.116037in}{0.714547in}}%
\pgfpathclose%
\pgfpathmoveto{\pgfqpoint{4.203396in}{0.671660in}}%
\pgfpathlineto{\pgfqpoint{4.203396in}{0.633600in}}%
\pgfpathlineto{\pgfqpoint{4.193039in}{0.633600in}}%
\pgfpathlineto{\pgfqpoint{4.193039in}{0.671317in}}%
\pgfpathquadraticcurveto{\pgfqpoint{4.193039in}{0.680269in}}{\pgfqpoint{4.189545in}{0.684700in}}%
\pgfpathquadraticcurveto{\pgfqpoint{4.186051in}{0.689149in}}{\pgfqpoint{4.179080in}{0.689149in}}%
\pgfpathquadraticcurveto{\pgfqpoint{4.170686in}{0.689149in}}{\pgfqpoint{4.165841in}{0.683800in}}%
\pgfpathquadraticcurveto{\pgfqpoint{4.160996in}{0.678468in}}{\pgfqpoint{4.160996in}{0.669228in}}%
\pgfpathlineto{\pgfqpoint{4.160996in}{0.633600in}}%
\pgfpathlineto{\pgfqpoint{4.150585in}{0.633600in}}%
\pgfpathlineto{\pgfqpoint{4.150585in}{0.721193in}}%
\pgfpathlineto{\pgfqpoint{4.160996in}{0.721193in}}%
\pgfpathlineto{\pgfqpoint{4.160996in}{0.686844in}}%
\pgfpathquadraticcurveto{\pgfqpoint{4.164724in}{0.692536in}}{\pgfqpoint{4.169750in}{0.695346in}}%
\pgfpathquadraticcurveto{\pgfqpoint{4.174793in}{0.698156in}}{\pgfqpoint{4.181385in}{0.698156in}}%
\pgfpathquadraticcurveto{\pgfqpoint{4.192247in}{0.698156in}}{\pgfqpoint{4.197813in}{0.691437in}}%
\pgfpathquadraticcurveto{\pgfqpoint{4.203396in}{0.684718in}}{\pgfqpoint{4.203396in}{0.671660in}}%
\pgfpathlineto{\pgfqpoint{4.203396in}{0.671660in}}%
\pgfpathclose%
\pgfpathmoveto{\pgfqpoint{4.277967in}{0.667715in}}%
\pgfpathlineto{\pgfqpoint{4.277967in}{0.662654in}}%
\pgfpathlineto{\pgfqpoint{4.230342in}{0.662654in}}%
\pgfpathquadraticcurveto{\pgfqpoint{4.231027in}{0.651954in}}{\pgfqpoint{4.236791in}{0.646353in}}%
\pgfpathquadraticcurveto{\pgfqpoint{4.242555in}{0.640751in}}{\pgfqpoint{4.252858in}{0.640751in}}%
\pgfpathquadraticcurveto{\pgfqpoint{4.258820in}{0.640751in}}{\pgfqpoint{4.264421in}{0.642210in}}%
\pgfpathquadraticcurveto{\pgfqpoint{4.270023in}{0.643669in}}{\pgfqpoint{4.275553in}{0.646605in}}%
\pgfpathlineto{\pgfqpoint{4.275553in}{0.636806in}}%
\pgfpathquadraticcurveto{\pgfqpoint{4.269969in}{0.634447in}}{\pgfqpoint{4.264115in}{0.633204in}}%
\pgfpathquadraticcurveto{\pgfqpoint{4.258261in}{0.631961in}}{\pgfqpoint{4.252245in}{0.631961in}}%
\pgfpathquadraticcurveto{\pgfqpoint{4.237151in}{0.631961in}}{\pgfqpoint{4.228343in}{0.640733in}}%
\pgfpathquadraticcurveto{\pgfqpoint{4.219535in}{0.649523in}}{\pgfqpoint{4.219535in}{0.664509in}}%
\pgfpathquadraticcurveto{\pgfqpoint{4.219535in}{0.679981in}}{\pgfqpoint{4.227893in}{0.689059in}}%
\pgfpathquadraticcurveto{\pgfqpoint{4.236250in}{0.698156in}}{\pgfqpoint{4.250444in}{0.698156in}}%
\pgfpathquadraticcurveto{\pgfqpoint{4.263161in}{0.698156in}}{\pgfqpoint{4.270564in}{0.689960in}}%
\pgfpathquadraticcurveto{\pgfqpoint{4.277967in}{0.681783in}}{\pgfqpoint{4.277967in}{0.667715in}}%
\pgfpathlineto{\pgfqpoint{4.277967in}{0.667715in}}%
\pgfpathclose%
\pgfpathmoveto{\pgfqpoint{4.267610in}{0.670759in}}%
\pgfpathquadraticcurveto{\pgfqpoint{4.267502in}{0.679243in}}{\pgfqpoint{4.262854in}{0.684304in}}%
\pgfpathquadraticcurveto{\pgfqpoint{4.258207in}{0.689384in}}{\pgfqpoint{4.250552in}{0.689384in}}%
\pgfpathquadraticcurveto{\pgfqpoint{4.241888in}{0.689384in}}{\pgfqpoint{4.236683in}{0.684484in}}%
\pgfpathquadraticcurveto{\pgfqpoint{4.231477in}{0.679585in}}{\pgfqpoint{4.230685in}{0.670687in}}%
\pgfpathlineto{\pgfqpoint{4.267610in}{0.670759in}}%
\pgfpathlineto{\pgfqpoint{4.267610in}{0.670759in}}%
\pgfpathclose%
\pgfusepath{fill}%
\end{pgfscope}%
\begin{pgfscope}%
\pgfsetbuttcap%
\pgfsetroundjoin%
\definecolor{currentfill}{rgb}{0.309804,0.372549,0.419608}%
\pgfsetfillcolor{currentfill}%
\pgfsetlinewidth{0.000000pt}%
\definecolor{currentstroke}{rgb}{0.309804,0.372549,0.419608}%
\pgfsetstrokecolor{currentstroke}%
\pgfsetdash{}{0pt}%
\pgfpathmoveto{\pgfqpoint{4.417293in}{0.667715in}}%
\pgfpathlineto{\pgfqpoint{4.417293in}{0.662654in}}%
\pgfpathlineto{\pgfqpoint{4.369669in}{0.662654in}}%
\pgfpathquadraticcurveto{\pgfqpoint{4.370354in}{0.651954in}}{\pgfqpoint{4.376117in}{0.646353in}}%
\pgfpathquadraticcurveto{\pgfqpoint{4.381881in}{0.640751in}}{\pgfqpoint{4.392184in}{0.640751in}}%
\pgfpathquadraticcurveto{\pgfqpoint{4.398146in}{0.640751in}}{\pgfqpoint{4.403748in}{0.642210in}}%
\pgfpathquadraticcurveto{\pgfqpoint{4.409350in}{0.643669in}}{\pgfqpoint{4.414880in}{0.646605in}}%
\pgfpathlineto{\pgfqpoint{4.414880in}{0.636806in}}%
\pgfpathquadraticcurveto{\pgfqpoint{4.409296in}{0.634447in}}{\pgfqpoint{4.403442in}{0.633204in}}%
\pgfpathquadraticcurveto{\pgfqpoint{4.397588in}{0.631961in}}{\pgfqpoint{4.391572in}{0.631961in}}%
\pgfpathquadraticcurveto{\pgfqpoint{4.376478in}{0.631961in}}{\pgfqpoint{4.367670in}{0.640733in}}%
\pgfpathquadraticcurveto{\pgfqpoint{4.358862in}{0.649523in}}{\pgfqpoint{4.358862in}{0.664509in}}%
\pgfpathquadraticcurveto{\pgfqpoint{4.358862in}{0.679981in}}{\pgfqpoint{4.367219in}{0.689059in}}%
\pgfpathquadraticcurveto{\pgfqpoint{4.375577in}{0.698156in}}{\pgfqpoint{4.389771in}{0.698156in}}%
\pgfpathquadraticcurveto{\pgfqpoint{4.402487in}{0.698156in}}{\pgfqpoint{4.409890in}{0.689960in}}%
\pgfpathquadraticcurveto{\pgfqpoint{4.417293in}{0.681783in}}{\pgfqpoint{4.417293in}{0.667715in}}%
\pgfpathlineto{\pgfqpoint{4.417293in}{0.667715in}}%
\pgfpathclose%
\pgfpathmoveto{\pgfqpoint{4.406936in}{0.670759in}}%
\pgfpathquadraticcurveto{\pgfqpoint{4.406828in}{0.679243in}}{\pgfqpoint{4.402181in}{0.684304in}}%
\pgfpathquadraticcurveto{\pgfqpoint{4.397534in}{0.689384in}}{\pgfqpoint{4.389879in}{0.689384in}}%
\pgfpathquadraticcurveto{\pgfqpoint{4.381215in}{0.689384in}}{\pgfqpoint{4.376009in}{0.684484in}}%
\pgfpathquadraticcurveto{\pgfqpoint{4.370804in}{0.679585in}}{\pgfqpoint{4.370011in}{0.670687in}}%
\pgfpathlineto{\pgfqpoint{4.406936in}{0.670759in}}%
\pgfpathlineto{\pgfqpoint{4.406936in}{0.670759in}}%
\pgfpathclose%
\pgfpathmoveto{\pgfqpoint{4.484695in}{0.696643in}}%
\pgfpathlineto{\pgfqpoint{4.461891in}{0.665968in}}%
\pgfpathlineto{\pgfqpoint{4.485865in}{0.633600in}}%
\pgfpathlineto{\pgfqpoint{4.473653in}{0.633600in}}%
\pgfpathlineto{\pgfqpoint{4.455299in}{0.658367in}}%
\pgfpathlineto{\pgfqpoint{4.436962in}{0.633600in}}%
\pgfpathlineto{\pgfqpoint{4.424732in}{0.633600in}}%
\pgfpathlineto{\pgfqpoint{4.449229in}{0.666580in}}%
\pgfpathlineto{\pgfqpoint{4.426822in}{0.696643in}}%
\pgfpathlineto{\pgfqpoint{4.439034in}{0.696643in}}%
\pgfpathlineto{\pgfqpoint{4.455749in}{0.674181in}}%
\pgfpathlineto{\pgfqpoint{4.472464in}{0.696643in}}%
\pgfpathlineto{\pgfqpoint{4.484695in}{0.696643in}}%
\pgfpathlineto{\pgfqpoint{4.484695in}{0.696643in}}%
\pgfpathclose%
\pgfpathmoveto{\pgfqpoint{4.500509in}{0.696643in}}%
\pgfpathlineto{\pgfqpoint{4.510866in}{0.696643in}}%
\pgfpathlineto{\pgfqpoint{4.510866in}{0.633600in}}%
\pgfpathlineto{\pgfqpoint{4.500509in}{0.633600in}}%
\pgfpathlineto{\pgfqpoint{4.500509in}{0.696643in}}%
\pgfpathlineto{\pgfqpoint{4.500509in}{0.696643in}}%
\pgfpathclose%
\pgfpathmoveto{\pgfqpoint{4.500509in}{0.721193in}}%
\pgfpathlineto{\pgfqpoint{4.510866in}{0.721193in}}%
\pgfpathlineto{\pgfqpoint{4.510866in}{0.708062in}}%
\pgfpathlineto{\pgfqpoint{4.500509in}{0.708062in}}%
\pgfpathlineto{\pgfqpoint{4.500509in}{0.721193in}}%
\pgfpathlineto{\pgfqpoint{4.500509in}{0.721193in}}%
\pgfpathclose%
\pgfpathmoveto{\pgfqpoint{4.572720in}{0.694787in}}%
\pgfpathlineto{\pgfqpoint{4.572720in}{0.684989in}}%
\pgfpathquadraticcurveto{\pgfqpoint{4.568343in}{0.687240in}}{\pgfqpoint{4.563606in}{0.688357in}}%
\pgfpathquadraticcurveto{\pgfqpoint{4.558887in}{0.689492in}}{\pgfqpoint{4.553807in}{0.689492in}}%
\pgfpathquadraticcurveto{\pgfqpoint{4.546098in}{0.689492in}}{\pgfqpoint{4.542243in}{0.687132in}}%
\pgfpathquadraticcurveto{\pgfqpoint{4.538389in}{0.684773in}}{\pgfqpoint{4.538389in}{0.680035in}}%
\pgfpathquadraticcurveto{\pgfqpoint{4.538389in}{0.676433in}}{\pgfqpoint{4.541145in}{0.674380in}}%
\pgfpathquadraticcurveto{\pgfqpoint{4.543901in}{0.672326in}}{\pgfqpoint{4.552240in}{0.670471in}}%
\pgfpathlineto{\pgfqpoint{4.555789in}{0.669678in}}%
\pgfpathquadraticcurveto{\pgfqpoint{4.566812in}{0.667319in}}{\pgfqpoint{4.571459in}{0.663014in}}%
\pgfpathquadraticcurveto{\pgfqpoint{4.576106in}{0.658709in}}{\pgfqpoint{4.576106in}{0.651000in}}%
\pgfpathquadraticcurveto{\pgfqpoint{4.576106in}{0.642210in}}{\pgfqpoint{4.569154in}{0.637076in}}%
\pgfpathquadraticcurveto{\pgfqpoint{4.562201in}{0.631961in}}{\pgfqpoint{4.550043in}{0.631961in}}%
\pgfpathquadraticcurveto{\pgfqpoint{4.544981in}{0.631961in}}{\pgfqpoint{4.539488in}{0.632952in}}%
\pgfpathquadraticcurveto{\pgfqpoint{4.533994in}{0.633942in}}{\pgfqpoint{4.527924in}{0.635906in}}%
\pgfpathlineto{\pgfqpoint{4.527924in}{0.646605in}}%
\pgfpathquadraticcurveto{\pgfqpoint{4.533670in}{0.643615in}}{\pgfqpoint{4.539235in}{0.642120in}}%
\pgfpathquadraticcurveto{\pgfqpoint{4.544801in}{0.640643in}}{\pgfqpoint{4.550277in}{0.640643in}}%
\pgfpathquadraticcurveto{\pgfqpoint{4.557590in}{0.640643in}}{\pgfqpoint{4.561516in}{0.643146in}}%
\pgfpathquadraticcurveto{\pgfqpoint{4.565461in}{0.645650in}}{\pgfqpoint{4.565461in}{0.650207in}}%
\pgfpathquadraticcurveto{\pgfqpoint{4.565461in}{0.654422in}}{\pgfqpoint{4.562615in}{0.656674in}}%
\pgfpathquadraticcurveto{\pgfqpoint{4.559787in}{0.658925in}}{\pgfqpoint{4.550151in}{0.661014in}}%
\pgfpathlineto{\pgfqpoint{4.546548in}{0.661861in}}%
\pgfpathquadraticcurveto{\pgfqpoint{4.536930in}{0.663878in}}{\pgfqpoint{4.532643in}{0.668075in}}%
\pgfpathquadraticcurveto{\pgfqpoint{4.528374in}{0.672272in}}{\pgfqpoint{4.528374in}{0.679585in}}%
\pgfpathquadraticcurveto{\pgfqpoint{4.528374in}{0.688483in}}{\pgfqpoint{4.534678in}{0.693310in}}%
\pgfpathquadraticcurveto{\pgfqpoint{4.540983in}{0.698156in}}{\pgfqpoint{4.552582in}{0.698156in}}%
\pgfpathquadraticcurveto{\pgfqpoint{4.558310in}{0.698156in}}{\pgfqpoint{4.563372in}{0.697309in}}%
\pgfpathquadraticcurveto{\pgfqpoint{4.568451in}{0.696480in}}{\pgfqpoint{4.572720in}{0.694787in}}%
\pgfpathlineto{\pgfqpoint{4.572720in}{0.694787in}}%
\pgfpathclose%
\pgfpathmoveto{\pgfqpoint{4.602836in}{0.714547in}}%
\pgfpathlineto{\pgfqpoint{4.602836in}{0.696643in}}%
\pgfpathlineto{\pgfqpoint{4.624163in}{0.696643in}}%
\pgfpathlineto{\pgfqpoint{4.624163in}{0.688591in}}%
\pgfpathlineto{\pgfqpoint{4.602836in}{0.688591in}}%
\pgfpathlineto{\pgfqpoint{4.602836in}{0.654368in}}%
\pgfpathquadraticcurveto{\pgfqpoint{4.602836in}{0.646659in}}{\pgfqpoint{4.604944in}{0.644461in}}%
\pgfpathquadraticcurveto{\pgfqpoint{4.607051in}{0.642264in}}{\pgfqpoint{4.613536in}{0.642264in}}%
\pgfpathlineto{\pgfqpoint{4.624163in}{0.642264in}}%
\pgfpathlineto{\pgfqpoint{4.624163in}{0.633600in}}%
\pgfpathlineto{\pgfqpoint{4.613536in}{0.633600in}}%
\pgfpathquadraticcurveto{\pgfqpoint{4.601539in}{0.633600in}}{\pgfqpoint{4.596982in}{0.638067in}}%
\pgfpathquadraticcurveto{\pgfqpoint{4.592425in}{0.642552in}}{\pgfqpoint{4.592425in}{0.654368in}}%
\pgfpathlineto{\pgfqpoint{4.592425in}{0.688591in}}%
\pgfpathlineto{\pgfqpoint{4.584824in}{0.688591in}}%
\pgfpathlineto{\pgfqpoint{4.584824in}{0.696643in}}%
\pgfpathlineto{\pgfqpoint{4.592425in}{0.696643in}}%
\pgfpathlineto{\pgfqpoint{4.592425in}{0.714547in}}%
\pgfpathlineto{\pgfqpoint{4.602836in}{0.714547in}}%
\pgfpathlineto{\pgfqpoint{4.602836in}{0.714547in}}%
\pgfpathclose%
\pgfpathmoveto{\pgfqpoint{4.637780in}{0.696643in}}%
\pgfpathlineto{\pgfqpoint{4.648137in}{0.696643in}}%
\pgfpathlineto{\pgfqpoint{4.648137in}{0.633600in}}%
\pgfpathlineto{\pgfqpoint{4.637780in}{0.633600in}}%
\pgfpathlineto{\pgfqpoint{4.637780in}{0.696643in}}%
\pgfpathlineto{\pgfqpoint{4.637780in}{0.696643in}}%
\pgfpathclose%
\pgfpathmoveto{\pgfqpoint{4.637780in}{0.721193in}}%
\pgfpathlineto{\pgfqpoint{4.648137in}{0.721193in}}%
\pgfpathlineto{\pgfqpoint{4.648137in}{0.708062in}}%
\pgfpathlineto{\pgfqpoint{4.637780in}{0.708062in}}%
\pgfpathlineto{\pgfqpoint{4.637780in}{0.721193in}}%
\pgfpathlineto{\pgfqpoint{4.637780in}{0.721193in}}%
\pgfpathclose%
\pgfpathmoveto{\pgfqpoint{4.722221in}{0.671660in}}%
\pgfpathlineto{\pgfqpoint{4.722221in}{0.633600in}}%
\pgfpathlineto{\pgfqpoint{4.711864in}{0.633600in}}%
\pgfpathlineto{\pgfqpoint{4.711864in}{0.671317in}}%
\pgfpathquadraticcurveto{\pgfqpoint{4.711864in}{0.680269in}}{\pgfqpoint{4.708370in}{0.684700in}}%
\pgfpathquadraticcurveto{\pgfqpoint{4.704875in}{0.689149in}}{\pgfqpoint{4.697905in}{0.689149in}}%
\pgfpathquadraticcurveto{\pgfqpoint{4.689511in}{0.689149in}}{\pgfqpoint{4.684666in}{0.683800in}}%
\pgfpathquadraticcurveto{\pgfqpoint{4.679820in}{0.678468in}}{\pgfqpoint{4.679820in}{0.669228in}}%
\pgfpathlineto{\pgfqpoint{4.679820in}{0.633600in}}%
\pgfpathlineto{\pgfqpoint{4.669409in}{0.633600in}}%
\pgfpathlineto{\pgfqpoint{4.669409in}{0.696643in}}%
\pgfpathlineto{\pgfqpoint{4.679820in}{0.696643in}}%
\pgfpathlineto{\pgfqpoint{4.679820in}{0.686844in}}%
\pgfpathquadraticcurveto{\pgfqpoint{4.683549in}{0.692536in}}{\pgfqpoint{4.688574in}{0.695346in}}%
\pgfpathquadraticcurveto{\pgfqpoint{4.693618in}{0.698156in}}{\pgfqpoint{4.700210in}{0.698156in}}%
\pgfpathquadraticcurveto{\pgfqpoint{4.711071in}{0.698156in}}{\pgfqpoint{4.716637in}{0.691437in}}%
\pgfpathquadraticcurveto{\pgfqpoint{4.722221in}{0.684718in}}{\pgfqpoint{4.722221in}{0.671660in}}%
\pgfpathlineto{\pgfqpoint{4.722221in}{0.671660in}}%
\pgfpathclose%
\pgfpathmoveto{\pgfqpoint{4.784345in}{0.665860in}}%
\pgfpathquadraticcurveto{\pgfqpoint{4.784345in}{0.677117in}}{\pgfqpoint{4.779698in}{0.683296in}}%
\pgfpathquadraticcurveto{\pgfqpoint{4.775069in}{0.689492in}}{\pgfqpoint{4.766675in}{0.689492in}}%
\pgfpathquadraticcurveto{\pgfqpoint{4.758353in}{0.689492in}}{\pgfqpoint{4.753706in}{0.683296in}}%
\pgfpathquadraticcurveto{\pgfqpoint{4.749059in}{0.677117in}}{\pgfqpoint{4.749059in}{0.665860in}}%
\pgfpathquadraticcurveto{\pgfqpoint{4.749059in}{0.654656in}}{\pgfqpoint{4.753706in}{0.648460in}}%
\pgfpathquadraticcurveto{\pgfqpoint{4.758353in}{0.642264in}}{\pgfqpoint{4.766675in}{0.642264in}}%
\pgfpathquadraticcurveto{\pgfqpoint{4.775069in}{0.642264in}}{\pgfqpoint{4.779698in}{0.648460in}}%
\pgfpathquadraticcurveto{\pgfqpoint{4.784345in}{0.654656in}}{\pgfqpoint{4.784345in}{0.665860in}}%
\pgfpathlineto{\pgfqpoint{4.784345in}{0.665860in}}%
\pgfpathclose%
\pgfpathmoveto{\pgfqpoint{4.794702in}{0.641417in}}%
\pgfpathquadraticcurveto{\pgfqpoint{4.794702in}{0.625332in}}{\pgfqpoint{4.787551in}{0.617479in}}%
\pgfpathquadraticcurveto{\pgfqpoint{4.780418in}{0.609626in}}{\pgfqpoint{4.765666in}{0.609626in}}%
\pgfpathquadraticcurveto{\pgfqpoint{4.760209in}{0.609626in}}{\pgfqpoint{4.755363in}{0.610436in}}%
\pgfpathquadraticcurveto{\pgfqpoint{4.750518in}{0.611247in}}{\pgfqpoint{4.745961in}{0.612940in}}%
\pgfpathlineto{\pgfqpoint{4.745961in}{0.623009in}}%
\pgfpathquadraticcurveto{\pgfqpoint{4.750518in}{0.620541in}}{\pgfqpoint{4.754967in}{0.619370in}}%
\pgfpathquadraticcurveto{\pgfqpoint{4.759416in}{0.618182in}}{\pgfqpoint{4.764027in}{0.618182in}}%
\pgfpathquadraticcurveto{\pgfqpoint{4.774222in}{0.618182in}}{\pgfqpoint{4.779283in}{0.623495in}}%
\pgfpathquadraticcurveto{\pgfqpoint{4.784345in}{0.628809in}}{\pgfqpoint{4.784345in}{0.639562in}}%
\pgfpathlineto{\pgfqpoint{4.784345in}{0.644695in}}%
\pgfpathquadraticcurveto{\pgfqpoint{4.781139in}{0.639112in}}{\pgfqpoint{4.776131in}{0.636356in}}%
\pgfpathquadraticcurveto{\pgfqpoint{4.771124in}{0.633600in}}{\pgfqpoint{4.764135in}{0.633600in}}%
\pgfpathquadraticcurveto{\pgfqpoint{4.752553in}{0.633600in}}{\pgfqpoint{4.745457in}{0.642426in}}%
\pgfpathquadraticcurveto{\pgfqpoint{4.738360in}{0.651270in}}{\pgfqpoint{4.738360in}{0.665860in}}%
\pgfpathquadraticcurveto{\pgfqpoint{4.738360in}{0.680486in}}{\pgfqpoint{4.745457in}{0.689312in}}%
\pgfpathquadraticcurveto{\pgfqpoint{4.752553in}{0.698156in}}{\pgfqpoint{4.764135in}{0.698156in}}%
\pgfpathquadraticcurveto{\pgfqpoint{4.771124in}{0.698156in}}{\pgfqpoint{4.776131in}{0.695400in}}%
\pgfpathquadraticcurveto{\pgfqpoint{4.781139in}{0.692644in}}{\pgfqpoint{4.784345in}{0.687078in}}%
\pgfpathlineto{\pgfqpoint{4.784345in}{0.696643in}}%
\pgfpathlineto{\pgfqpoint{4.794702in}{0.696643in}}%
\pgfpathlineto{\pgfqpoint{4.794702in}{0.641417in}}%
\pgfpathlineto{\pgfqpoint{4.794702in}{0.641417in}}%
\pgfpathclose%
\pgfusepath{fill}%
\end{pgfscope}%
\begin{pgfscope}%
\pgfsetbuttcap%
\pgfsetroundjoin%
\definecolor{currentfill}{rgb}{0.309804,0.372549,0.419608}%
\pgfsetfillcolor{currentfill}%
\pgfsetlinewidth{0.000000pt}%
\definecolor{currentstroke}{rgb}{0.309804,0.372549,0.419608}%
\pgfsetstrokecolor{currentstroke}%
\pgfsetdash{}{0pt}%
\pgfpathmoveto{\pgfqpoint{4.934174in}{0.665067in}}%
\pgfpathquadraticcurveto{\pgfqpoint{4.934174in}{0.676487in}}{\pgfqpoint{4.929473in}{0.682989in}}%
\pgfpathquadraticcurveto{\pgfqpoint{4.924772in}{0.689492in}}{\pgfqpoint{4.916559in}{0.689492in}}%
\pgfpathquadraticcurveto{\pgfqpoint{4.908327in}{0.689492in}}{\pgfqpoint{4.903626in}{0.682989in}}%
\pgfpathquadraticcurveto{\pgfqpoint{4.898925in}{0.676487in}}{\pgfqpoint{4.898925in}{0.665067in}}%
\pgfpathquadraticcurveto{\pgfqpoint{4.898925in}{0.653648in}}{\pgfqpoint{4.903626in}{0.647145in}}%
\pgfpathquadraticcurveto{\pgfqpoint{4.908327in}{0.640643in}}{\pgfqpoint{4.916559in}{0.640643in}}%
\pgfpathquadraticcurveto{\pgfqpoint{4.924772in}{0.640643in}}{\pgfqpoint{4.929473in}{0.647145in}}%
\pgfpathquadraticcurveto{\pgfqpoint{4.934174in}{0.653648in}}{\pgfqpoint{4.934174in}{0.665067in}}%
\pgfpathlineto{\pgfqpoint{4.934174in}{0.665067in}}%
\pgfpathclose%
\pgfpathmoveto{\pgfqpoint{4.898925in}{0.687078in}}%
\pgfpathquadraticcurveto{\pgfqpoint{4.902203in}{0.692698in}}{\pgfqpoint{4.907174in}{0.695418in}}%
\pgfpathquadraticcurveto{\pgfqpoint{4.912164in}{0.698156in}}{\pgfqpoint{4.919080in}{0.698156in}}%
\pgfpathquadraticcurveto{\pgfqpoint{4.930572in}{0.698156in}}{\pgfqpoint{4.937741in}{0.689041in}}%
\pgfpathquadraticcurveto{\pgfqpoint{4.944928in}{0.679927in}}{\pgfqpoint{4.944928in}{0.665067in}}%
\pgfpathquadraticcurveto{\pgfqpoint{4.944928in}{0.650207in}}{\pgfqpoint{4.937741in}{0.641075in}}%
\pgfpathquadraticcurveto{\pgfqpoint{4.930572in}{0.631961in}}{\pgfqpoint{4.919080in}{0.631961in}}%
\pgfpathquadraticcurveto{\pgfqpoint{4.912164in}{0.631961in}}{\pgfqpoint{4.907174in}{0.634699in}}%
\pgfpathquadraticcurveto{\pgfqpoint{4.902203in}{0.637437in}}{\pgfqpoint{4.898925in}{0.643056in}}%
\pgfpathlineto{\pgfqpoint{4.898925in}{0.633600in}}%
\pgfpathlineto{\pgfqpoint{4.888514in}{0.633600in}}%
\pgfpathlineto{\pgfqpoint{4.888514in}{0.721193in}}%
\pgfpathlineto{\pgfqpoint{4.898925in}{0.721193in}}%
\pgfpathlineto{\pgfqpoint{4.898925in}{0.687078in}}%
\pgfpathlineto{\pgfqpoint{4.898925in}{0.687078in}}%
\pgfpathclose%
\pgfpathmoveto{\pgfqpoint{4.986518in}{0.689384in}}%
\pgfpathquadraticcurveto{\pgfqpoint{4.978196in}{0.689384in}}{\pgfqpoint{4.973351in}{0.682881in}}%
\pgfpathquadraticcurveto{\pgfqpoint{4.968506in}{0.676379in}}{\pgfqpoint{4.968506in}{0.665067in}}%
\pgfpathquadraticcurveto{\pgfqpoint{4.968506in}{0.653756in}}{\pgfqpoint{4.973315in}{0.647253in}}%
\pgfpathquadraticcurveto{\pgfqpoint{4.978142in}{0.640751in}}{\pgfqpoint{4.986518in}{0.640751in}}%
\pgfpathquadraticcurveto{\pgfqpoint{4.994803in}{0.640751in}}{\pgfqpoint{4.999631in}{0.647271in}}%
\pgfpathquadraticcurveto{\pgfqpoint{5.004476in}{0.653810in}}{\pgfqpoint{5.004476in}{0.665067in}}%
\pgfpathquadraticcurveto{\pgfqpoint{5.004476in}{0.676271in}}{\pgfqpoint{4.999631in}{0.682827in}}%
\pgfpathquadraticcurveto{\pgfqpoint{4.994803in}{0.689384in}}{\pgfqpoint{4.986518in}{0.689384in}}%
\pgfpathlineto{\pgfqpoint{4.986518in}{0.689384in}}%
\pgfpathclose%
\pgfpathmoveto{\pgfqpoint{4.986518in}{0.698156in}}%
\pgfpathquadraticcurveto{\pgfqpoint{5.000027in}{0.698156in}}{\pgfqpoint{5.007736in}{0.689366in}}%
\pgfpathquadraticcurveto{\pgfqpoint{5.015463in}{0.680594in}}{\pgfqpoint{5.015463in}{0.665067in}}%
\pgfpathquadraticcurveto{\pgfqpoint{5.015463in}{0.649595in}}{\pgfqpoint{5.007736in}{0.640769in}}%
\pgfpathquadraticcurveto{\pgfqpoint{5.000027in}{0.631961in}}{\pgfqpoint{4.986518in}{0.631961in}}%
\pgfpathquadraticcurveto{\pgfqpoint{4.972955in}{0.631961in}}{\pgfqpoint{4.965263in}{0.640769in}}%
\pgfpathquadraticcurveto{\pgfqpoint{4.957590in}{0.649595in}}{\pgfqpoint{4.957590in}{0.665067in}}%
\pgfpathquadraticcurveto{\pgfqpoint{4.957590in}{0.680594in}}{\pgfqpoint{4.965263in}{0.689366in}}%
\pgfpathquadraticcurveto{\pgfqpoint{4.972955in}{0.698156in}}{\pgfqpoint{4.986518in}{0.698156in}}%
\pgfpathlineto{\pgfqpoint{4.986518in}{0.698156in}}%
\pgfpathclose%
\pgfpathmoveto{\pgfqpoint{5.057053in}{0.689384in}}%
\pgfpathquadraticcurveto{\pgfqpoint{5.048732in}{0.689384in}}{\pgfqpoint{5.043887in}{0.682881in}}%
\pgfpathquadraticcurveto{\pgfqpoint{5.039041in}{0.676379in}}{\pgfqpoint{5.039041in}{0.665067in}}%
\pgfpathquadraticcurveto{\pgfqpoint{5.039041in}{0.653756in}}{\pgfqpoint{5.043850in}{0.647253in}}%
\pgfpathquadraticcurveto{\pgfqpoint{5.048678in}{0.640751in}}{\pgfqpoint{5.057053in}{0.640751in}}%
\pgfpathquadraticcurveto{\pgfqpoint{5.065339in}{0.640751in}}{\pgfqpoint{5.070166in}{0.647271in}}%
\pgfpathquadraticcurveto{\pgfqpoint{5.075012in}{0.653810in}}{\pgfqpoint{5.075012in}{0.665067in}}%
\pgfpathquadraticcurveto{\pgfqpoint{5.075012in}{0.676271in}}{\pgfqpoint{5.070166in}{0.682827in}}%
\pgfpathquadraticcurveto{\pgfqpoint{5.065339in}{0.689384in}}{\pgfqpoint{5.057053in}{0.689384in}}%
\pgfpathlineto{\pgfqpoint{5.057053in}{0.689384in}}%
\pgfpathclose%
\pgfpathmoveto{\pgfqpoint{5.057053in}{0.698156in}}%
\pgfpathquadraticcurveto{\pgfqpoint{5.070563in}{0.698156in}}{\pgfqpoint{5.078272in}{0.689366in}}%
\pgfpathquadraticcurveto{\pgfqpoint{5.085999in}{0.680594in}}{\pgfqpoint{5.085999in}{0.665067in}}%
\pgfpathquadraticcurveto{\pgfqpoint{5.085999in}{0.649595in}}{\pgfqpoint{5.078272in}{0.640769in}}%
\pgfpathquadraticcurveto{\pgfqpoint{5.070563in}{0.631961in}}{\pgfqpoint{5.057053in}{0.631961in}}%
\pgfpathquadraticcurveto{\pgfqpoint{5.043490in}{0.631961in}}{\pgfqpoint{5.035799in}{0.640769in}}%
\pgfpathquadraticcurveto{\pgfqpoint{5.028126in}{0.649595in}}{\pgfqpoint{5.028126in}{0.665067in}}%
\pgfpathquadraticcurveto{\pgfqpoint{5.028126in}{0.680594in}}{\pgfqpoint{5.035799in}{0.689366in}}%
\pgfpathquadraticcurveto{\pgfqpoint{5.043490in}{0.698156in}}{\pgfqpoint{5.057053in}{0.698156in}}%
\pgfpathlineto{\pgfqpoint{5.057053in}{0.698156in}}%
\pgfpathclose%
\pgfpathmoveto{\pgfqpoint{5.113413in}{0.714547in}}%
\pgfpathlineto{\pgfqpoint{5.113413in}{0.696643in}}%
\pgfpathlineto{\pgfqpoint{5.134740in}{0.696643in}}%
\pgfpathlineto{\pgfqpoint{5.134740in}{0.688591in}}%
\pgfpathlineto{\pgfqpoint{5.113413in}{0.688591in}}%
\pgfpathlineto{\pgfqpoint{5.113413in}{0.654368in}}%
\pgfpathquadraticcurveto{\pgfqpoint{5.113413in}{0.646659in}}{\pgfqpoint{5.115521in}{0.644461in}}%
\pgfpathquadraticcurveto{\pgfqpoint{5.117628in}{0.642264in}}{\pgfqpoint{5.124113in}{0.642264in}}%
\pgfpathlineto{\pgfqpoint{5.134740in}{0.642264in}}%
\pgfpathlineto{\pgfqpoint{5.134740in}{0.633600in}}%
\pgfpathlineto{\pgfqpoint{5.124113in}{0.633600in}}%
\pgfpathquadraticcurveto{\pgfqpoint{5.112117in}{0.633600in}}{\pgfqpoint{5.107559in}{0.638067in}}%
\pgfpathquadraticcurveto{\pgfqpoint{5.103002in}{0.642552in}}{\pgfqpoint{5.103002in}{0.654368in}}%
\pgfpathlineto{\pgfqpoint{5.103002in}{0.688591in}}%
\pgfpathlineto{\pgfqpoint{5.095401in}{0.688591in}}%
\pgfpathlineto{\pgfqpoint{5.095401in}{0.696643in}}%
\pgfpathlineto{\pgfqpoint{5.103002in}{0.696643in}}%
\pgfpathlineto{\pgfqpoint{5.103002in}{0.714547in}}%
\pgfpathlineto{\pgfqpoint{5.113413in}{0.714547in}}%
\pgfpathlineto{\pgfqpoint{5.113413in}{0.714547in}}%
\pgfpathclose%
\pgfpathmoveto{\pgfqpoint{5.188542in}{0.694787in}}%
\pgfpathlineto{\pgfqpoint{5.188542in}{0.684989in}}%
\pgfpathquadraticcurveto{\pgfqpoint{5.184165in}{0.687240in}}{\pgfqpoint{5.179428in}{0.688357in}}%
\pgfpathquadraticcurveto{\pgfqpoint{5.174709in}{0.689492in}}{\pgfqpoint{5.169629in}{0.689492in}}%
\pgfpathquadraticcurveto{\pgfqpoint{5.161920in}{0.689492in}}{\pgfqpoint{5.158066in}{0.687132in}}%
\pgfpathquadraticcurveto{\pgfqpoint{5.154211in}{0.684773in}}{\pgfqpoint{5.154211in}{0.680035in}}%
\pgfpathquadraticcurveto{\pgfqpoint{5.154211in}{0.676433in}}{\pgfqpoint{5.156967in}{0.674380in}}%
\pgfpathquadraticcurveto{\pgfqpoint{5.159723in}{0.672326in}}{\pgfqpoint{5.168062in}{0.670471in}}%
\pgfpathlineto{\pgfqpoint{5.171611in}{0.669678in}}%
\pgfpathquadraticcurveto{\pgfqpoint{5.182634in}{0.667319in}}{\pgfqpoint{5.187281in}{0.663014in}}%
\pgfpathquadraticcurveto{\pgfqpoint{5.191928in}{0.658709in}}{\pgfqpoint{5.191928in}{0.651000in}}%
\pgfpathquadraticcurveto{\pgfqpoint{5.191928in}{0.642210in}}{\pgfqpoint{5.184976in}{0.637076in}}%
\pgfpathquadraticcurveto{\pgfqpoint{5.178023in}{0.631961in}}{\pgfqpoint{5.165865in}{0.631961in}}%
\pgfpathquadraticcurveto{\pgfqpoint{5.160803in}{0.631961in}}{\pgfqpoint{5.155310in}{0.632952in}}%
\pgfpathquadraticcurveto{\pgfqpoint{5.149816in}{0.633942in}}{\pgfqpoint{5.143746in}{0.635906in}}%
\pgfpathlineto{\pgfqpoint{5.143746in}{0.646605in}}%
\pgfpathquadraticcurveto{\pgfqpoint{5.149492in}{0.643615in}}{\pgfqpoint{5.155058in}{0.642120in}}%
\pgfpathquadraticcurveto{\pgfqpoint{5.160623in}{0.640643in}}{\pgfqpoint{5.166099in}{0.640643in}}%
\pgfpathquadraticcurveto{\pgfqpoint{5.173412in}{0.640643in}}{\pgfqpoint{5.177339in}{0.643146in}}%
\pgfpathquadraticcurveto{\pgfqpoint{5.181283in}{0.645650in}}{\pgfqpoint{5.181283in}{0.650207in}}%
\pgfpathquadraticcurveto{\pgfqpoint{5.181283in}{0.654422in}}{\pgfqpoint{5.178437in}{0.656674in}}%
\pgfpathquadraticcurveto{\pgfqpoint{5.175609in}{0.658925in}}{\pgfqpoint{5.165973in}{0.661014in}}%
\pgfpathlineto{\pgfqpoint{5.162370in}{0.661861in}}%
\pgfpathquadraticcurveto{\pgfqpoint{5.152752in}{0.663878in}}{\pgfqpoint{5.148465in}{0.668075in}}%
\pgfpathquadraticcurveto{\pgfqpoint{5.144196in}{0.672272in}}{\pgfqpoint{5.144196in}{0.679585in}}%
\pgfpathquadraticcurveto{\pgfqpoint{5.144196in}{0.688483in}}{\pgfqpoint{5.150500in}{0.693310in}}%
\pgfpathquadraticcurveto{\pgfqpoint{5.156805in}{0.698156in}}{\pgfqpoint{5.168405in}{0.698156in}}%
\pgfpathquadraticcurveto{\pgfqpoint{5.174132in}{0.698156in}}{\pgfqpoint{5.179194in}{0.697309in}}%
\pgfpathquadraticcurveto{\pgfqpoint{5.184273in}{0.696480in}}{\pgfqpoint{5.188542in}{0.694787in}}%
\pgfpathlineto{\pgfqpoint{5.188542in}{0.694787in}}%
\pgfpathclose%
\pgfpathmoveto{\pgfqpoint{5.218658in}{0.714547in}}%
\pgfpathlineto{\pgfqpoint{5.218658in}{0.696643in}}%
\pgfpathlineto{\pgfqpoint{5.239985in}{0.696643in}}%
\pgfpathlineto{\pgfqpoint{5.239985in}{0.688591in}}%
\pgfpathlineto{\pgfqpoint{5.218658in}{0.688591in}}%
\pgfpathlineto{\pgfqpoint{5.218658in}{0.654368in}}%
\pgfpathquadraticcurveto{\pgfqpoint{5.218658in}{0.646659in}}{\pgfqpoint{5.220766in}{0.644461in}}%
\pgfpathquadraticcurveto{\pgfqpoint{5.222873in}{0.642264in}}{\pgfqpoint{5.229358in}{0.642264in}}%
\pgfpathlineto{\pgfqpoint{5.239985in}{0.642264in}}%
\pgfpathlineto{\pgfqpoint{5.239985in}{0.633600in}}%
\pgfpathlineto{\pgfqpoint{5.229358in}{0.633600in}}%
\pgfpathquadraticcurveto{\pgfqpoint{5.217362in}{0.633600in}}{\pgfqpoint{5.212804in}{0.638067in}}%
\pgfpathquadraticcurveto{\pgfqpoint{5.208247in}{0.642552in}}{\pgfqpoint{5.208247in}{0.654368in}}%
\pgfpathlineto{\pgfqpoint{5.208247in}{0.688591in}}%
\pgfpathlineto{\pgfqpoint{5.200646in}{0.688591in}}%
\pgfpathlineto{\pgfqpoint{5.200646in}{0.696643in}}%
\pgfpathlineto{\pgfqpoint{5.208247in}{0.696643in}}%
\pgfpathlineto{\pgfqpoint{5.208247in}{0.714547in}}%
\pgfpathlineto{\pgfqpoint{5.218658in}{0.714547in}}%
\pgfpathlineto{\pgfqpoint{5.218658in}{0.714547in}}%
\pgfpathclose%
\pgfpathmoveto{\pgfqpoint{5.290131in}{0.686970in}}%
\pgfpathquadraticcurveto{\pgfqpoint{5.288383in}{0.687979in}}{\pgfqpoint{5.286330in}{0.688447in}}%
\pgfpathquadraticcurveto{\pgfqpoint{5.284277in}{0.688933in}}{\pgfqpoint{5.281809in}{0.688933in}}%
\pgfpathquadraticcurveto{\pgfqpoint{5.273019in}{0.688933in}}{\pgfqpoint{5.268318in}{0.683223in}}%
\pgfpathquadraticcurveto{\pgfqpoint{5.263617in}{0.677514in}}{\pgfqpoint{5.263617in}{0.666814in}}%
\pgfpathlineto{\pgfqpoint{5.263617in}{0.633600in}}%
\pgfpathlineto{\pgfqpoint{5.253206in}{0.633600in}}%
\pgfpathlineto{\pgfqpoint{5.253206in}{0.696643in}}%
\pgfpathlineto{\pgfqpoint{5.263617in}{0.696643in}}%
\pgfpathlineto{\pgfqpoint{5.263617in}{0.686844in}}%
\pgfpathquadraticcurveto{\pgfqpoint{5.266895in}{0.692590in}}{\pgfqpoint{5.272118in}{0.695364in}}%
\pgfpathquadraticcurveto{\pgfqpoint{5.277360in}{0.698156in}}{\pgfqpoint{5.284853in}{0.698156in}}%
\pgfpathquadraticcurveto{\pgfqpoint{5.285916in}{0.698156in}}{\pgfqpoint{5.287213in}{0.698011in}}%
\pgfpathquadraticcurveto{\pgfqpoint{5.288510in}{0.697885in}}{\pgfqpoint{5.290077in}{0.697597in}}%
\pgfpathlineto{\pgfqpoint{5.290131in}{0.686970in}}%
\pgfpathlineto{\pgfqpoint{5.290131in}{0.686970in}}%
\pgfpathclose%
\pgfpathmoveto{\pgfqpoint{5.329649in}{0.665283in}}%
\pgfpathquadraticcurveto{\pgfqpoint{5.317095in}{0.665283in}}{\pgfqpoint{5.312250in}{0.662419in}}%
\pgfpathquadraticcurveto{\pgfqpoint{5.307404in}{0.659556in}}{\pgfqpoint{5.307404in}{0.652621in}}%
\pgfpathquadraticcurveto{\pgfqpoint{5.307404in}{0.647109in}}{\pgfqpoint{5.311043in}{0.643867in}}%
\pgfpathquadraticcurveto{\pgfqpoint{5.314681in}{0.640643in}}{\pgfqpoint{5.320913in}{0.640643in}}%
\pgfpathquadraticcurveto{\pgfqpoint{5.329541in}{0.640643in}}{\pgfqpoint{5.334747in}{0.646749in}}%
\pgfpathquadraticcurveto{\pgfqpoint{5.339952in}{0.652855in}}{\pgfqpoint{5.339952in}{0.662978in}}%
\pgfpathlineto{\pgfqpoint{5.339952in}{0.665283in}}%
\pgfpathlineto{\pgfqpoint{5.329649in}{0.665283in}}%
\pgfpathlineto{\pgfqpoint{5.329649in}{0.665283in}}%
\pgfpathclose%
\pgfpathmoveto{\pgfqpoint{5.350309in}{0.669570in}}%
\pgfpathlineto{\pgfqpoint{5.350309in}{0.633600in}}%
\pgfpathlineto{\pgfqpoint{5.339952in}{0.633600in}}%
\pgfpathlineto{\pgfqpoint{5.339952in}{0.643164in}}%
\pgfpathquadraticcurveto{\pgfqpoint{5.336404in}{0.637437in}}{\pgfqpoint{5.331108in}{0.634699in}}%
\pgfpathquadraticcurveto{\pgfqpoint{5.325813in}{0.631961in}}{\pgfqpoint{5.318158in}{0.631961in}}%
\pgfpathquadraticcurveto{\pgfqpoint{5.308485in}{0.631961in}}{\pgfqpoint{5.302757in}{0.637401in}}%
\pgfpathquadraticcurveto{\pgfqpoint{5.297047in}{0.642840in}}{\pgfqpoint{5.297047in}{0.651954in}}%
\pgfpathquadraticcurveto{\pgfqpoint{5.297047in}{0.662582in}}{\pgfqpoint{5.304162in}{0.667985in}}%
\pgfpathquadraticcurveto{\pgfqpoint{5.311295in}{0.673389in}}{\pgfqpoint{5.325416in}{0.673389in}}%
\pgfpathlineto{\pgfqpoint{5.339952in}{0.673389in}}%
\pgfpathlineto{\pgfqpoint{5.339952in}{0.674416in}}%
\pgfpathquadraticcurveto{\pgfqpoint{5.339952in}{0.681566in}}{\pgfqpoint{5.335251in}{0.685475in}}%
\pgfpathquadraticcurveto{\pgfqpoint{5.330550in}{0.689384in}}{\pgfqpoint{5.322048in}{0.689384in}}%
\pgfpathquadraticcurveto{\pgfqpoint{5.316645in}{0.689384in}}{\pgfqpoint{5.311511in}{0.688087in}}%
\pgfpathquadraticcurveto{\pgfqpoint{5.306396in}{0.686790in}}{\pgfqpoint{5.301676in}{0.684196in}}%
\pgfpathlineto{\pgfqpoint{5.301676in}{0.693779in}}%
\pgfpathquadraticcurveto{\pgfqpoint{5.307350in}{0.695976in}}{\pgfqpoint{5.312700in}{0.697057in}}%
\pgfpathquadraticcurveto{\pgfqpoint{5.318049in}{0.698156in}}{\pgfqpoint{5.323111in}{0.698156in}}%
\pgfpathquadraticcurveto{\pgfqpoint{5.336800in}{0.698156in}}{\pgfqpoint{5.343555in}{0.691059in}}%
\pgfpathquadraticcurveto{\pgfqpoint{5.350309in}{0.683980in}}{\pgfqpoint{5.350309in}{0.669570in}}%
\pgfpathlineto{\pgfqpoint{5.350309in}{0.669570in}}%
\pgfpathclose%
\pgfpathmoveto{\pgfqpoint{5.381650in}{0.643056in}}%
\pgfpathlineto{\pgfqpoint{5.381650in}{0.609626in}}%
\pgfpathlineto{\pgfqpoint{5.371239in}{0.609626in}}%
\pgfpathlineto{\pgfqpoint{5.371239in}{0.696643in}}%
\pgfpathlineto{\pgfqpoint{5.381650in}{0.696643in}}%
\pgfpathlineto{\pgfqpoint{5.381650in}{0.687078in}}%
\pgfpathquadraticcurveto{\pgfqpoint{5.384929in}{0.692698in}}{\pgfqpoint{5.389900in}{0.695418in}}%
\pgfpathquadraticcurveto{\pgfqpoint{5.394889in}{0.698156in}}{\pgfqpoint{5.401806in}{0.698156in}}%
\pgfpathquadraticcurveto{\pgfqpoint{5.413298in}{0.698156in}}{\pgfqpoint{5.420467in}{0.689041in}}%
\pgfpathquadraticcurveto{\pgfqpoint{5.427653in}{0.679927in}}{\pgfqpoint{5.427653in}{0.665067in}}%
\pgfpathquadraticcurveto{\pgfqpoint{5.427653in}{0.650207in}}{\pgfqpoint{5.420467in}{0.641075in}}%
\pgfpathquadraticcurveto{\pgfqpoint{5.413298in}{0.631961in}}{\pgfqpoint{5.401806in}{0.631961in}}%
\pgfpathquadraticcurveto{\pgfqpoint{5.394889in}{0.631961in}}{\pgfqpoint{5.389900in}{0.634699in}}%
\pgfpathquadraticcurveto{\pgfqpoint{5.384929in}{0.637437in}}{\pgfqpoint{5.381650in}{0.643056in}}%
\pgfpathlineto{\pgfqpoint{5.381650in}{0.643056in}}%
\pgfpathclose%
\pgfpathmoveto{\pgfqpoint{5.416900in}{0.665067in}}%
\pgfpathquadraticcurveto{\pgfqpoint{5.416900in}{0.676487in}}{\pgfqpoint{5.412199in}{0.682989in}}%
\pgfpathquadraticcurveto{\pgfqpoint{5.407498in}{0.689492in}}{\pgfqpoint{5.399284in}{0.689492in}}%
\pgfpathquadraticcurveto{\pgfqpoint{5.391053in}{0.689492in}}{\pgfqpoint{5.386352in}{0.682989in}}%
\pgfpathquadraticcurveto{\pgfqpoint{5.381650in}{0.676487in}}{\pgfqpoint{5.381650in}{0.665067in}}%
\pgfpathquadraticcurveto{\pgfqpoint{5.381650in}{0.653648in}}{\pgfqpoint{5.386352in}{0.647145in}}%
\pgfpathquadraticcurveto{\pgfqpoint{5.391053in}{0.640643in}}{\pgfqpoint{5.399284in}{0.640643in}}%
\pgfpathquadraticcurveto{\pgfqpoint{5.407498in}{0.640643in}}{\pgfqpoint{5.412199in}{0.647145in}}%
\pgfpathquadraticcurveto{\pgfqpoint{5.416900in}{0.653648in}}{\pgfqpoint{5.416900in}{0.665067in}}%
\pgfpathlineto{\pgfqpoint{5.416900in}{0.665067in}}%
\pgfpathclose%
\pgfusepath{fill}%
\end{pgfscope}%
\begin{pgfscope}%
\pgfsetbuttcap%
\pgfsetroundjoin%
\definecolor{currentfill}{rgb}{0.309804,0.372549,0.419608}%
\pgfsetfillcolor{currentfill}%
\pgfsetlinewidth{0.000000pt}%
\definecolor{currentstroke}{rgb}{0.309804,0.372549,0.419608}%
\pgfsetstrokecolor{currentstroke}%
\pgfsetdash{}{0pt}%
\pgfpathmoveto{\pgfqpoint{5.571859in}{0.694229in}}%
\pgfpathlineto{\pgfqpoint{5.571859in}{0.684538in}}%
\pgfpathquadraticcurveto{\pgfqpoint{5.567464in}{0.686970in}}{\pgfqpoint{5.563051in}{0.688177in}}%
\pgfpathquadraticcurveto{\pgfqpoint{5.558638in}{0.689384in}}{\pgfqpoint{5.554135in}{0.689384in}}%
\pgfpathquadraticcurveto{\pgfqpoint{5.544048in}{0.689384in}}{\pgfqpoint{5.538464in}{0.682989in}}%
\pgfpathquadraticcurveto{\pgfqpoint{5.532899in}{0.676613in}}{\pgfqpoint{5.532899in}{0.665067in}}%
\pgfpathquadraticcurveto{\pgfqpoint{5.532899in}{0.653521in}}{\pgfqpoint{5.538464in}{0.647127in}}%
\pgfpathquadraticcurveto{\pgfqpoint{5.544048in}{0.640751in}}{\pgfqpoint{5.554135in}{0.640751in}}%
\pgfpathquadraticcurveto{\pgfqpoint{5.558638in}{0.640751in}}{\pgfqpoint{5.563051in}{0.641958in}}%
\pgfpathquadraticcurveto{\pgfqpoint{5.567464in}{0.643164in}}{\pgfqpoint{5.571859in}{0.645596in}}%
\pgfpathlineto{\pgfqpoint{5.571859in}{0.636014in}}%
\pgfpathquadraticcurveto{\pgfqpoint{5.567518in}{0.633996in}}{\pgfqpoint{5.562871in}{0.632988in}}%
\pgfpathquadraticcurveto{\pgfqpoint{5.558242in}{0.631961in}}{\pgfqpoint{5.553000in}{0.631961in}}%
\pgfpathquadraticcurveto{\pgfqpoint{5.538753in}{0.631961in}}{\pgfqpoint{5.530359in}{0.640913in}}%
\pgfpathquadraticcurveto{\pgfqpoint{5.521983in}{0.649865in}}{\pgfqpoint{5.521983in}{0.665067in}}%
\pgfpathquadraticcurveto{\pgfqpoint{5.521983in}{0.680486in}}{\pgfqpoint{5.530449in}{0.689312in}}%
\pgfpathquadraticcurveto{\pgfqpoint{5.538933in}{0.698156in}}{\pgfqpoint{5.553685in}{0.698156in}}%
\pgfpathquadraticcurveto{\pgfqpoint{5.558458in}{0.698156in}}{\pgfqpoint{5.563015in}{0.697165in}}%
\pgfpathquadraticcurveto{\pgfqpoint{5.567572in}{0.696192in}}{\pgfqpoint{5.571859in}{0.694229in}}%
\pgfpathlineto{\pgfqpoint{5.571859in}{0.694229in}}%
\pgfpathclose%
\pgfpathmoveto{\pgfqpoint{5.589871in}{0.721193in}}%
\pgfpathlineto{\pgfqpoint{5.600228in}{0.721193in}}%
\pgfpathlineto{\pgfqpoint{5.600228in}{0.633600in}}%
\pgfpathlineto{\pgfqpoint{5.589871in}{0.633600in}}%
\pgfpathlineto{\pgfqpoint{5.589871in}{0.721193in}}%
\pgfpathlineto{\pgfqpoint{5.589871in}{0.721193in}}%
\pgfpathclose%
\pgfpathmoveto{\pgfqpoint{5.620834in}{0.658475in}}%
\pgfpathlineto{\pgfqpoint{5.620834in}{0.696643in}}%
\pgfpathlineto{\pgfqpoint{5.631191in}{0.696643in}}%
\pgfpathlineto{\pgfqpoint{5.631191in}{0.658871in}}%
\pgfpathquadraticcurveto{\pgfqpoint{5.631191in}{0.649919in}}{\pgfqpoint{5.634667in}{0.645434in}}%
\pgfpathquadraticcurveto{\pgfqpoint{5.638162in}{0.640967in}}{\pgfqpoint{5.645150in}{0.640967in}}%
\pgfpathquadraticcurveto{\pgfqpoint{5.653526in}{0.640967in}}{\pgfqpoint{5.658389in}{0.646317in}}%
\pgfpathquadraticcurveto{\pgfqpoint{5.663271in}{0.651666in}}{\pgfqpoint{5.663271in}{0.660906in}}%
\pgfpathlineto{\pgfqpoint{5.663271in}{0.696643in}}%
\pgfpathlineto{\pgfqpoint{5.673628in}{0.696643in}}%
\pgfpathlineto{\pgfqpoint{5.673628in}{0.633600in}}%
\pgfpathlineto{\pgfqpoint{5.663271in}{0.633600in}}%
\pgfpathlineto{\pgfqpoint{5.663271in}{0.643291in}}%
\pgfpathquadraticcurveto{\pgfqpoint{5.659506in}{0.637545in}}{\pgfqpoint{5.654517in}{0.634753in}}%
\pgfpathquadraticcurveto{\pgfqpoint{5.649545in}{0.631961in}}{\pgfqpoint{5.642953in}{0.631961in}}%
\pgfpathquadraticcurveto{\pgfqpoint{5.632092in}{0.631961in}}{\pgfqpoint{5.626454in}{0.638715in}}%
\pgfpathquadraticcurveto{\pgfqpoint{5.620834in}{0.645470in}}{\pgfqpoint{5.620834in}{0.658475in}}%
\pgfpathlineto{\pgfqpoint{5.620834in}{0.658475in}}%
\pgfpathclose%
\pgfpathmoveto{\pgfqpoint{5.646898in}{0.698156in}}%
\pgfpathlineto{\pgfqpoint{5.646898in}{0.698156in}}%
\pgfpathlineto{\pgfqpoint{5.646898in}{0.698156in}}%
\pgfpathclose%
\pgfpathmoveto{\pgfqpoint{5.735139in}{0.694787in}}%
\pgfpathlineto{\pgfqpoint{5.735139in}{0.684989in}}%
\pgfpathquadraticcurveto{\pgfqpoint{5.730762in}{0.687240in}}{\pgfqpoint{5.726025in}{0.688357in}}%
\pgfpathquadraticcurveto{\pgfqpoint{5.721306in}{0.689492in}}{\pgfqpoint{5.716226in}{0.689492in}}%
\pgfpathquadraticcurveto{\pgfqpoint{5.708517in}{0.689492in}}{\pgfqpoint{5.704663in}{0.687132in}}%
\pgfpathquadraticcurveto{\pgfqpoint{5.700808in}{0.684773in}}{\pgfqpoint{5.700808in}{0.680035in}}%
\pgfpathquadraticcurveto{\pgfqpoint{5.700808in}{0.676433in}}{\pgfqpoint{5.703564in}{0.674380in}}%
\pgfpathquadraticcurveto{\pgfqpoint{5.706320in}{0.672326in}}{\pgfqpoint{5.714659in}{0.670471in}}%
\pgfpathlineto{\pgfqpoint{5.718208in}{0.669678in}}%
\pgfpathquadraticcurveto{\pgfqpoint{5.729231in}{0.667319in}}{\pgfqpoint{5.733878in}{0.663014in}}%
\pgfpathquadraticcurveto{\pgfqpoint{5.738525in}{0.658709in}}{\pgfqpoint{5.738525in}{0.651000in}}%
\pgfpathquadraticcurveto{\pgfqpoint{5.738525in}{0.642210in}}{\pgfqpoint{5.731573in}{0.637076in}}%
\pgfpathquadraticcurveto{\pgfqpoint{5.724620in}{0.631961in}}{\pgfqpoint{5.712462in}{0.631961in}}%
\pgfpathquadraticcurveto{\pgfqpoint{5.707400in}{0.631961in}}{\pgfqpoint{5.701907in}{0.632952in}}%
\pgfpathquadraticcurveto{\pgfqpoint{5.696413in}{0.633942in}}{\pgfqpoint{5.690343in}{0.635906in}}%
\pgfpathlineto{\pgfqpoint{5.690343in}{0.646605in}}%
\pgfpathquadraticcurveto{\pgfqpoint{5.696089in}{0.643615in}}{\pgfqpoint{5.701655in}{0.642120in}}%
\pgfpathquadraticcurveto{\pgfqpoint{5.707220in}{0.640643in}}{\pgfqpoint{5.712696in}{0.640643in}}%
\pgfpathquadraticcurveto{\pgfqpoint{5.720009in}{0.640643in}}{\pgfqpoint{5.723936in}{0.643146in}}%
\pgfpathquadraticcurveto{\pgfqpoint{5.727880in}{0.645650in}}{\pgfqpoint{5.727880in}{0.650207in}}%
\pgfpathquadraticcurveto{\pgfqpoint{5.727880in}{0.654422in}}{\pgfqpoint{5.725034in}{0.656674in}}%
\pgfpathquadraticcurveto{\pgfqpoint{5.722206in}{0.658925in}}{\pgfqpoint{5.712570in}{0.661014in}}%
\pgfpathlineto{\pgfqpoint{5.708967in}{0.661861in}}%
\pgfpathquadraticcurveto{\pgfqpoint{5.699349in}{0.663878in}}{\pgfqpoint{5.695062in}{0.668075in}}%
\pgfpathquadraticcurveto{\pgfqpoint{5.690793in}{0.672272in}}{\pgfqpoint{5.690793in}{0.679585in}}%
\pgfpathquadraticcurveto{\pgfqpoint{5.690793in}{0.688483in}}{\pgfqpoint{5.697097in}{0.693310in}}%
\pgfpathquadraticcurveto{\pgfqpoint{5.703402in}{0.698156in}}{\pgfqpoint{5.715002in}{0.698156in}}%
\pgfpathquadraticcurveto{\pgfqpoint{5.720729in}{0.698156in}}{\pgfqpoint{5.725791in}{0.697309in}}%
\pgfpathquadraticcurveto{\pgfqpoint{5.730870in}{0.696480in}}{\pgfqpoint{5.735139in}{0.694787in}}%
\pgfpathlineto{\pgfqpoint{5.735139in}{0.694787in}}%
\pgfpathclose%
\pgfpathmoveto{\pgfqpoint{5.765255in}{0.714547in}}%
\pgfpathlineto{\pgfqpoint{5.765255in}{0.696643in}}%
\pgfpathlineto{\pgfqpoint{5.786582in}{0.696643in}}%
\pgfpathlineto{\pgfqpoint{5.786582in}{0.688591in}}%
\pgfpathlineto{\pgfqpoint{5.765255in}{0.688591in}}%
\pgfpathlineto{\pgfqpoint{5.765255in}{0.654368in}}%
\pgfpathquadraticcurveto{\pgfqpoint{5.765255in}{0.646659in}}{\pgfqpoint{5.767363in}{0.644461in}}%
\pgfpathquadraticcurveto{\pgfqpoint{5.769470in}{0.642264in}}{\pgfqpoint{5.775955in}{0.642264in}}%
\pgfpathlineto{\pgfqpoint{5.786582in}{0.642264in}}%
\pgfpathlineto{\pgfqpoint{5.786582in}{0.633600in}}%
\pgfpathlineto{\pgfqpoint{5.775955in}{0.633600in}}%
\pgfpathquadraticcurveto{\pgfqpoint{5.763959in}{0.633600in}}{\pgfqpoint{5.759401in}{0.638067in}}%
\pgfpathquadraticcurveto{\pgfqpoint{5.754844in}{0.642552in}}{\pgfqpoint{5.754844in}{0.654368in}}%
\pgfpathlineto{\pgfqpoint{5.754844in}{0.688591in}}%
\pgfpathlineto{\pgfqpoint{5.747243in}{0.688591in}}%
\pgfpathlineto{\pgfqpoint{5.747243in}{0.696643in}}%
\pgfpathlineto{\pgfqpoint{5.754844in}{0.696643in}}%
\pgfpathlineto{\pgfqpoint{5.754844in}{0.714547in}}%
\pgfpathlineto{\pgfqpoint{5.765255in}{0.714547in}}%
\pgfpathlineto{\pgfqpoint{5.765255in}{0.714547in}}%
\pgfpathclose%
\pgfpathmoveto{\pgfqpoint{5.854127in}{0.667715in}}%
\pgfpathlineto{\pgfqpoint{5.854127in}{0.662654in}}%
\pgfpathlineto{\pgfqpoint{5.806503in}{0.662654in}}%
\pgfpathquadraticcurveto{\pgfqpoint{5.807188in}{0.651954in}}{\pgfqpoint{5.812952in}{0.646353in}}%
\pgfpathquadraticcurveto{\pgfqpoint{5.818715in}{0.640751in}}{\pgfqpoint{5.829018in}{0.640751in}}%
\pgfpathquadraticcurveto{\pgfqpoint{5.834980in}{0.640751in}}{\pgfqpoint{5.840582in}{0.642210in}}%
\pgfpathquadraticcurveto{\pgfqpoint{5.846184in}{0.643669in}}{\pgfqpoint{5.851714in}{0.646605in}}%
\pgfpathlineto{\pgfqpoint{5.851714in}{0.636806in}}%
\pgfpathquadraticcurveto{\pgfqpoint{5.846130in}{0.634447in}}{\pgfqpoint{5.840276in}{0.633204in}}%
\pgfpathquadraticcurveto{\pgfqpoint{5.834422in}{0.631961in}}{\pgfqpoint{5.828406in}{0.631961in}}%
\pgfpathquadraticcurveto{\pgfqpoint{5.813312in}{0.631961in}}{\pgfqpoint{5.804504in}{0.640733in}}%
\pgfpathquadraticcurveto{\pgfqpoint{5.795696in}{0.649523in}}{\pgfqpoint{5.795696in}{0.664509in}}%
\pgfpathquadraticcurveto{\pgfqpoint{5.795696in}{0.679981in}}{\pgfqpoint{5.804054in}{0.689059in}}%
\pgfpathquadraticcurveto{\pgfqpoint{5.812411in}{0.698156in}}{\pgfqpoint{5.826605in}{0.698156in}}%
\pgfpathquadraticcurveto{\pgfqpoint{5.839321in}{0.698156in}}{\pgfqpoint{5.846724in}{0.689960in}}%
\pgfpathquadraticcurveto{\pgfqpoint{5.854127in}{0.681783in}}{\pgfqpoint{5.854127in}{0.667715in}}%
\pgfpathlineto{\pgfqpoint{5.854127in}{0.667715in}}%
\pgfpathclose%
\pgfpathmoveto{\pgfqpoint{5.843770in}{0.670759in}}%
\pgfpathquadraticcurveto{\pgfqpoint{5.843662in}{0.679243in}}{\pgfqpoint{5.839015in}{0.684304in}}%
\pgfpathquadraticcurveto{\pgfqpoint{5.834368in}{0.689384in}}{\pgfqpoint{5.826713in}{0.689384in}}%
\pgfpathquadraticcurveto{\pgfqpoint{5.818049in}{0.689384in}}{\pgfqpoint{5.812844in}{0.684484in}}%
\pgfpathquadraticcurveto{\pgfqpoint{5.807638in}{0.679585in}}{\pgfqpoint{5.806845in}{0.670687in}}%
\pgfpathlineto{\pgfqpoint{5.843770in}{0.670759in}}%
\pgfpathlineto{\pgfqpoint{5.843770in}{0.670759in}}%
\pgfpathclose%
\pgfpathmoveto{\pgfqpoint{5.907660in}{0.686970in}}%
\pgfpathquadraticcurveto{\pgfqpoint{5.905912in}{0.687979in}}{\pgfqpoint{5.903859in}{0.688447in}}%
\pgfpathquadraticcurveto{\pgfqpoint{5.901806in}{0.688933in}}{\pgfqpoint{5.899338in}{0.688933in}}%
\pgfpathquadraticcurveto{\pgfqpoint{5.890548in}{0.688933in}}{\pgfqpoint{5.885847in}{0.683223in}}%
\pgfpathquadraticcurveto{\pgfqpoint{5.881146in}{0.677514in}}{\pgfqpoint{5.881146in}{0.666814in}}%
\pgfpathlineto{\pgfqpoint{5.881146in}{0.633600in}}%
\pgfpathlineto{\pgfqpoint{5.870735in}{0.633600in}}%
\pgfpathlineto{\pgfqpoint{5.870735in}{0.696643in}}%
\pgfpathlineto{\pgfqpoint{5.881146in}{0.696643in}}%
\pgfpathlineto{\pgfqpoint{5.881146in}{0.686844in}}%
\pgfpathquadraticcurveto{\pgfqpoint{5.884424in}{0.692590in}}{\pgfqpoint{5.889647in}{0.695364in}}%
\pgfpathquadraticcurveto{\pgfqpoint{5.894889in}{0.698156in}}{\pgfqpoint{5.902382in}{0.698156in}}%
\pgfpathquadraticcurveto{\pgfqpoint{5.903445in}{0.698156in}}{\pgfqpoint{5.904742in}{0.698011in}}%
\pgfpathquadraticcurveto{\pgfqpoint{5.906038in}{0.697885in}}{\pgfqpoint{5.907605in}{0.697597in}}%
\pgfpathlineto{\pgfqpoint{5.907660in}{0.686970in}}%
\pgfpathlineto{\pgfqpoint{5.907660in}{0.686970in}}%
\pgfpathclose%
\pgfpathmoveto{\pgfqpoint{5.958706in}{0.694787in}}%
\pgfpathlineto{\pgfqpoint{5.958706in}{0.684989in}}%
\pgfpathquadraticcurveto{\pgfqpoint{5.954329in}{0.687240in}}{\pgfqpoint{5.949592in}{0.688357in}}%
\pgfpathquadraticcurveto{\pgfqpoint{5.944873in}{0.689492in}}{\pgfqpoint{5.939793in}{0.689492in}}%
\pgfpathquadraticcurveto{\pgfqpoint{5.932084in}{0.689492in}}{\pgfqpoint{5.928229in}{0.687132in}}%
\pgfpathquadraticcurveto{\pgfqpoint{5.924375in}{0.684773in}}{\pgfqpoint{5.924375in}{0.680035in}}%
\pgfpathquadraticcurveto{\pgfqpoint{5.924375in}{0.676433in}}{\pgfqpoint{5.927131in}{0.674380in}}%
\pgfpathquadraticcurveto{\pgfqpoint{5.929887in}{0.672326in}}{\pgfqpoint{5.938226in}{0.670471in}}%
\pgfpathlineto{\pgfqpoint{5.941775in}{0.669678in}}%
\pgfpathquadraticcurveto{\pgfqpoint{5.952798in}{0.667319in}}{\pgfqpoint{5.957445in}{0.663014in}}%
\pgfpathquadraticcurveto{\pgfqpoint{5.962092in}{0.658709in}}{\pgfqpoint{5.962092in}{0.651000in}}%
\pgfpathquadraticcurveto{\pgfqpoint{5.962092in}{0.642210in}}{\pgfqpoint{5.955140in}{0.637076in}}%
\pgfpathquadraticcurveto{\pgfqpoint{5.948187in}{0.631961in}}{\pgfqpoint{5.936029in}{0.631961in}}%
\pgfpathquadraticcurveto{\pgfqpoint{5.930967in}{0.631961in}}{\pgfqpoint{5.925474in}{0.632952in}}%
\pgfpathquadraticcurveto{\pgfqpoint{5.919980in}{0.633942in}}{\pgfqpoint{5.913910in}{0.635906in}}%
\pgfpathlineto{\pgfqpoint{5.913910in}{0.646605in}}%
\pgfpathquadraticcurveto{\pgfqpoint{5.919656in}{0.643615in}}{\pgfqpoint{5.925221in}{0.642120in}}%
\pgfpathquadraticcurveto{\pgfqpoint{5.930787in}{0.640643in}}{\pgfqpoint{5.936263in}{0.640643in}}%
\pgfpathquadraticcurveto{\pgfqpoint{5.943576in}{0.640643in}}{\pgfqpoint{5.947502in}{0.643146in}}%
\pgfpathquadraticcurveto{\pgfqpoint{5.951447in}{0.645650in}}{\pgfqpoint{5.951447in}{0.650207in}}%
\pgfpathquadraticcurveto{\pgfqpoint{5.951447in}{0.654422in}}{\pgfqpoint{5.948601in}{0.656674in}}%
\pgfpathquadraticcurveto{\pgfqpoint{5.945773in}{0.658925in}}{\pgfqpoint{5.936137in}{0.661014in}}%
\pgfpathlineto{\pgfqpoint{5.932534in}{0.661861in}}%
\pgfpathquadraticcurveto{\pgfqpoint{5.922916in}{0.663878in}}{\pgfqpoint{5.918629in}{0.668075in}}%
\pgfpathquadraticcurveto{\pgfqpoint{5.914360in}{0.672272in}}{\pgfqpoint{5.914360in}{0.679585in}}%
\pgfpathquadraticcurveto{\pgfqpoint{5.914360in}{0.688483in}}{\pgfqpoint{5.920664in}{0.693310in}}%
\pgfpathquadraticcurveto{\pgfqpoint{5.926969in}{0.698156in}}{\pgfqpoint{5.938568in}{0.698156in}}%
\pgfpathquadraticcurveto{\pgfqpoint{5.944296in}{0.698156in}}{\pgfqpoint{5.949358in}{0.697309in}}%
\pgfpathquadraticcurveto{\pgfqpoint{5.954437in}{0.696480in}}{\pgfqpoint{5.958706in}{0.694787in}}%
\pgfpathlineto{\pgfqpoint{5.958706in}{0.694787in}}%
\pgfpathclose%
\pgfpathmoveto{\pgfqpoint{5.981221in}{0.693202in}}%
\pgfpathlineto{\pgfqpoint{5.993091in}{0.693202in}}%
\pgfpathlineto{\pgfqpoint{5.993091in}{0.678919in}}%
\pgfpathlineto{\pgfqpoint{5.981221in}{0.678919in}}%
\pgfpathlineto{\pgfqpoint{5.981221in}{0.693202in}}%
\pgfpathlineto{\pgfqpoint{5.981221in}{0.693202in}}%
\pgfpathclose%
\pgfpathmoveto{\pgfqpoint{5.981221in}{0.647902in}}%
\pgfpathlineto{\pgfqpoint{5.993091in}{0.647902in}}%
\pgfpathlineto{\pgfqpoint{5.993091in}{0.638211in}}%
\pgfpathlineto{\pgfqpoint{5.983869in}{0.620199in}}%
\pgfpathlineto{\pgfqpoint{5.976610in}{0.620199in}}%
\pgfpathlineto{\pgfqpoint{5.981221in}{0.638211in}}%
\pgfpathlineto{\pgfqpoint{5.981221in}{0.647902in}}%
\pgfpathlineto{\pgfqpoint{5.981221in}{0.647902in}}%
\pgfpathclose%
\pgfusepath{fill}%
\end{pgfscope}%
\begin{pgfscope}%
\pgfsetbuttcap%
\pgfsetroundjoin%
\definecolor{currentfill}{rgb}{0.309804,0.372549,0.419608}%
\pgfsetfillcolor{currentfill}%
\pgfsetlinewidth{0.000000pt}%
\definecolor{currentstroke}{rgb}{0.309804,0.372549,0.419608}%
\pgfsetstrokecolor{currentstroke}%
\pgfsetdash{}{0pt}%
\pgfpathmoveto{\pgfqpoint{6.094063in}{0.696643in}}%
\pgfpathlineto{\pgfqpoint{6.104420in}{0.696643in}}%
\pgfpathlineto{\pgfqpoint{6.104420in}{0.633600in}}%
\pgfpathlineto{\pgfqpoint{6.094063in}{0.633600in}}%
\pgfpathlineto{\pgfqpoint{6.094063in}{0.696643in}}%
\pgfpathlineto{\pgfqpoint{6.094063in}{0.696643in}}%
\pgfpathclose%
\pgfpathmoveto{\pgfqpoint{6.094063in}{0.721193in}}%
\pgfpathlineto{\pgfqpoint{6.104420in}{0.721193in}}%
\pgfpathlineto{\pgfqpoint{6.104420in}{0.708062in}}%
\pgfpathlineto{\pgfqpoint{6.094063in}{0.708062in}}%
\pgfpathlineto{\pgfqpoint{6.094063in}{0.721193in}}%
\pgfpathlineto{\pgfqpoint{6.094063in}{0.721193in}}%
\pgfpathclose%
\pgfpathmoveto{\pgfqpoint{6.136337in}{0.714547in}}%
\pgfpathlineto{\pgfqpoint{6.136337in}{0.696643in}}%
\pgfpathlineto{\pgfqpoint{6.157664in}{0.696643in}}%
\pgfpathlineto{\pgfqpoint{6.157664in}{0.688591in}}%
\pgfpathlineto{\pgfqpoint{6.136337in}{0.688591in}}%
\pgfpathlineto{\pgfqpoint{6.136337in}{0.654368in}}%
\pgfpathquadraticcurveto{\pgfqpoint{6.136337in}{0.646659in}}{\pgfqpoint{6.138445in}{0.644461in}}%
\pgfpathquadraticcurveto{\pgfqpoint{6.140552in}{0.642264in}}{\pgfqpoint{6.147036in}{0.642264in}}%
\pgfpathlineto{\pgfqpoint{6.157664in}{0.642264in}}%
\pgfpathlineto{\pgfqpoint{6.157664in}{0.633600in}}%
\pgfpathlineto{\pgfqpoint{6.147036in}{0.633600in}}%
\pgfpathquadraticcurveto{\pgfqpoint{6.135040in}{0.633600in}}{\pgfqpoint{6.130483in}{0.638067in}}%
\pgfpathquadraticcurveto{\pgfqpoint{6.125926in}{0.642552in}}{\pgfqpoint{6.125926in}{0.654368in}}%
\pgfpathlineto{\pgfqpoint{6.125926in}{0.688591in}}%
\pgfpathlineto{\pgfqpoint{6.118325in}{0.688591in}}%
\pgfpathlineto{\pgfqpoint{6.118325in}{0.696643in}}%
\pgfpathlineto{\pgfqpoint{6.125926in}{0.696643in}}%
\pgfpathlineto{\pgfqpoint{6.125926in}{0.714547in}}%
\pgfpathlineto{\pgfqpoint{6.136337in}{0.714547in}}%
\pgfpathlineto{\pgfqpoint{6.136337in}{0.714547in}}%
\pgfpathclose%
\pgfusepath{fill}%
\end{pgfscope}%
\begin{pgfscope}%
\pgfsetbuttcap%
\pgfsetroundjoin%
\definecolor{currentfill}{rgb}{0.309804,0.372549,0.419608}%
\pgfsetfillcolor{currentfill}%
\pgfsetlinewidth{0.000000pt}%
\definecolor{currentstroke}{rgb}{0.309804,0.372549,0.419608}%
\pgfsetstrokecolor{currentstroke}%
\pgfsetdash{}{0pt}%
\pgfpathmoveto{\pgfqpoint{6.231639in}{0.696643in}}%
\pgfpathlineto{\pgfqpoint{6.241996in}{0.696643in}}%
\pgfpathlineto{\pgfqpoint{6.241996in}{0.633600in}}%
\pgfpathlineto{\pgfqpoint{6.231639in}{0.633600in}}%
\pgfpathlineto{\pgfqpoint{6.231639in}{0.696643in}}%
\pgfpathlineto{\pgfqpoint{6.231639in}{0.696643in}}%
\pgfpathclose%
\pgfpathmoveto{\pgfqpoint{6.231639in}{0.721193in}}%
\pgfpathlineto{\pgfqpoint{6.241996in}{0.721193in}}%
\pgfpathlineto{\pgfqpoint{6.241996in}{0.708062in}}%
\pgfpathlineto{\pgfqpoint{6.231639in}{0.708062in}}%
\pgfpathlineto{\pgfqpoint{6.231639in}{0.721193in}}%
\pgfpathlineto{\pgfqpoint{6.231639in}{0.721193in}}%
\pgfpathclose%
\pgfpathmoveto{\pgfqpoint{6.303850in}{0.694787in}}%
\pgfpathlineto{\pgfqpoint{6.303850in}{0.684989in}}%
\pgfpathquadraticcurveto{\pgfqpoint{6.299473in}{0.687240in}}{\pgfqpoint{6.294736in}{0.688357in}}%
\pgfpathquadraticcurveto{\pgfqpoint{6.290016in}{0.689492in}}{\pgfqpoint{6.284937in}{0.689492in}}%
\pgfpathquadraticcurveto{\pgfqpoint{6.277228in}{0.689492in}}{\pgfqpoint{6.273373in}{0.687132in}}%
\pgfpathquadraticcurveto{\pgfqpoint{6.269519in}{0.684773in}}{\pgfqpoint{6.269519in}{0.680035in}}%
\pgfpathquadraticcurveto{\pgfqpoint{6.269519in}{0.676433in}}{\pgfqpoint{6.272275in}{0.674380in}}%
\pgfpathquadraticcurveto{\pgfqpoint{6.275030in}{0.672326in}}{\pgfqpoint{6.283370in}{0.670471in}}%
\pgfpathlineto{\pgfqpoint{6.286918in}{0.669678in}}%
\pgfpathquadraticcurveto{\pgfqpoint{6.297942in}{0.667319in}}{\pgfqpoint{6.302589in}{0.663014in}}%
\pgfpathquadraticcurveto{\pgfqpoint{6.307236in}{0.658709in}}{\pgfqpoint{6.307236in}{0.651000in}}%
\pgfpathquadraticcurveto{\pgfqpoint{6.307236in}{0.642210in}}{\pgfqpoint{6.300283in}{0.637076in}}%
\pgfpathquadraticcurveto{\pgfqpoint{6.293331in}{0.631961in}}{\pgfqpoint{6.281173in}{0.631961in}}%
\pgfpathquadraticcurveto{\pgfqpoint{6.276111in}{0.631961in}}{\pgfqpoint{6.270617in}{0.632952in}}%
\pgfpathquadraticcurveto{\pgfqpoint{6.265124in}{0.633942in}}{\pgfqpoint{6.259054in}{0.635906in}}%
\pgfpathlineto{\pgfqpoint{6.259054in}{0.646605in}}%
\pgfpathquadraticcurveto{\pgfqpoint{6.264799in}{0.643615in}}{\pgfqpoint{6.270365in}{0.642120in}}%
\pgfpathquadraticcurveto{\pgfqpoint{6.275931in}{0.640643in}}{\pgfqpoint{6.281407in}{0.640643in}}%
\pgfpathquadraticcurveto{\pgfqpoint{6.288720in}{0.640643in}}{\pgfqpoint{6.292646in}{0.643146in}}%
\pgfpathquadraticcurveto{\pgfqpoint{6.296591in}{0.645650in}}{\pgfqpoint{6.296591in}{0.650207in}}%
\pgfpathquadraticcurveto{\pgfqpoint{6.296591in}{0.654422in}}{\pgfqpoint{6.293745in}{0.656674in}}%
\pgfpathquadraticcurveto{\pgfqpoint{6.290917in}{0.658925in}}{\pgfqpoint{6.281281in}{0.661014in}}%
\pgfpathlineto{\pgfqpoint{6.277678in}{0.661861in}}%
\pgfpathquadraticcurveto{\pgfqpoint{6.268060in}{0.663878in}}{\pgfqpoint{6.263773in}{0.668075in}}%
\pgfpathquadraticcurveto{\pgfqpoint{6.259504in}{0.672272in}}{\pgfqpoint{6.259504in}{0.679585in}}%
\pgfpathquadraticcurveto{\pgfqpoint{6.259504in}{0.688483in}}{\pgfqpoint{6.265808in}{0.693310in}}%
\pgfpathquadraticcurveto{\pgfqpoint{6.272112in}{0.698156in}}{\pgfqpoint{6.283712in}{0.698156in}}%
\pgfpathquadraticcurveto{\pgfqpoint{6.289440in}{0.698156in}}{\pgfqpoint{6.294502in}{0.697309in}}%
\pgfpathquadraticcurveto{\pgfqpoint{6.299581in}{0.696480in}}{\pgfqpoint{6.303850in}{0.694787in}}%
\pgfpathlineto{\pgfqpoint{6.303850in}{0.694787in}}%
\pgfpathclose%
\pgfusepath{fill}%
\end{pgfscope}%
\begin{pgfscope}%
\pgfsetbuttcap%
\pgfsetroundjoin%
\definecolor{currentfill}{rgb}{0.309804,0.372549,0.419608}%
\pgfsetfillcolor{currentfill}%
\pgfsetlinewidth{0.000000pt}%
\definecolor{currentstroke}{rgb}{0.309804,0.372549,0.419608}%
\pgfsetstrokecolor{currentstroke}%
\pgfsetdash{}{0pt}%
\pgfpathmoveto{\pgfqpoint{6.407873in}{0.665283in}}%
\pgfpathquadraticcurveto{\pgfqpoint{6.395318in}{0.665283in}}{\pgfqpoint{6.390473in}{0.662419in}}%
\pgfpathquadraticcurveto{\pgfqpoint{6.385628in}{0.659556in}}{\pgfqpoint{6.385628in}{0.652621in}}%
\pgfpathquadraticcurveto{\pgfqpoint{6.385628in}{0.647109in}}{\pgfqpoint{6.389266in}{0.643867in}}%
\pgfpathquadraticcurveto{\pgfqpoint{6.392905in}{0.640643in}}{\pgfqpoint{6.399137in}{0.640643in}}%
\pgfpathquadraticcurveto{\pgfqpoint{6.407765in}{0.640643in}}{\pgfqpoint{6.412970in}{0.646749in}}%
\pgfpathquadraticcurveto{\pgfqpoint{6.418176in}{0.652855in}}{\pgfqpoint{6.418176in}{0.662978in}}%
\pgfpathlineto{\pgfqpoint{6.418176in}{0.665283in}}%
\pgfpathlineto{\pgfqpoint{6.407873in}{0.665283in}}%
\pgfpathlineto{\pgfqpoint{6.407873in}{0.665283in}}%
\pgfpathclose%
\pgfpathmoveto{\pgfqpoint{6.428533in}{0.669570in}}%
\pgfpathlineto{\pgfqpoint{6.428533in}{0.633600in}}%
\pgfpathlineto{\pgfqpoint{6.418176in}{0.633600in}}%
\pgfpathlineto{\pgfqpoint{6.418176in}{0.643164in}}%
\pgfpathquadraticcurveto{\pgfqpoint{6.414627in}{0.637437in}}{\pgfqpoint{6.409332in}{0.634699in}}%
\pgfpathquadraticcurveto{\pgfqpoint{6.404036in}{0.631961in}}{\pgfqpoint{6.396381in}{0.631961in}}%
\pgfpathquadraticcurveto{\pgfqpoint{6.386709in}{0.631961in}}{\pgfqpoint{6.380981in}{0.637401in}}%
\pgfpathquadraticcurveto{\pgfqpoint{6.375271in}{0.642840in}}{\pgfqpoint{6.375271in}{0.651954in}}%
\pgfpathquadraticcurveto{\pgfqpoint{6.375271in}{0.662582in}}{\pgfqpoint{6.382386in}{0.667985in}}%
\pgfpathquadraticcurveto{\pgfqpoint{6.389518in}{0.673389in}}{\pgfqpoint{6.403640in}{0.673389in}}%
\pgfpathlineto{\pgfqpoint{6.418176in}{0.673389in}}%
\pgfpathlineto{\pgfqpoint{6.418176in}{0.674416in}}%
\pgfpathquadraticcurveto{\pgfqpoint{6.418176in}{0.681566in}}{\pgfqpoint{6.413475in}{0.685475in}}%
\pgfpathquadraticcurveto{\pgfqpoint{6.408773in}{0.689384in}}{\pgfqpoint{6.400272in}{0.689384in}}%
\pgfpathquadraticcurveto{\pgfqpoint{6.394868in}{0.689384in}}{\pgfqpoint{6.389735in}{0.688087in}}%
\pgfpathquadraticcurveto{\pgfqpoint{6.384619in}{0.686790in}}{\pgfqpoint{6.379900in}{0.684196in}}%
\pgfpathlineto{\pgfqpoint{6.379900in}{0.693779in}}%
\pgfpathquadraticcurveto{\pgfqpoint{6.385574in}{0.695976in}}{\pgfqpoint{6.390923in}{0.697057in}}%
\pgfpathquadraticcurveto{\pgfqpoint{6.396273in}{0.698156in}}{\pgfqpoint{6.401334in}{0.698156in}}%
\pgfpathquadraticcurveto{\pgfqpoint{6.415024in}{0.698156in}}{\pgfqpoint{6.421778in}{0.691059in}}%
\pgfpathquadraticcurveto{\pgfqpoint{6.428533in}{0.683980in}}{\pgfqpoint{6.428533in}{0.669570in}}%
\pgfpathlineto{\pgfqpoint{6.428533in}{0.669570in}}%
\pgfpathclose%
\pgfusepath{fill}%
\end{pgfscope}%
\begin{pgfscope}%
\pgfsetbuttcap%
\pgfsetroundjoin%
\definecolor{currentfill}{rgb}{0.309804,0.372549,0.419608}%
\pgfsetfillcolor{currentfill}%
\pgfsetlinewidth{0.000000pt}%
\definecolor{currentstroke}{rgb}{0.309804,0.372549,0.419608}%
\pgfsetstrokecolor{currentstroke}%
\pgfsetdash{}{0pt}%
\pgfpathmoveto{\pgfqpoint{6.548274in}{0.687078in}}%
\pgfpathlineto{\pgfqpoint{6.548274in}{0.721193in}}%
\pgfpathlineto{\pgfqpoint{6.558631in}{0.721193in}}%
\pgfpathlineto{\pgfqpoint{6.558631in}{0.633600in}}%
\pgfpathlineto{\pgfqpoint{6.548274in}{0.633600in}}%
\pgfpathlineto{\pgfqpoint{6.548274in}{0.643056in}}%
\pgfpathquadraticcurveto{\pgfqpoint{6.545014in}{0.637437in}}{\pgfqpoint{6.540024in}{0.634699in}}%
\pgfpathquadraticcurveto{\pgfqpoint{6.535053in}{0.631961in}}{\pgfqpoint{6.528064in}{0.631961in}}%
\pgfpathquadraticcurveto{\pgfqpoint{6.516645in}{0.631961in}}{\pgfqpoint{6.509458in}{0.641075in}}%
\pgfpathquadraticcurveto{\pgfqpoint{6.502289in}{0.650207in}}{\pgfqpoint{6.502289in}{0.665067in}}%
\pgfpathquadraticcurveto{\pgfqpoint{6.502289in}{0.679927in}}{\pgfqpoint{6.509458in}{0.689041in}}%
\pgfpathquadraticcurveto{\pgfqpoint{6.516645in}{0.698156in}}{\pgfqpoint{6.528064in}{0.698156in}}%
\pgfpathquadraticcurveto{\pgfqpoint{6.535053in}{0.698156in}}{\pgfqpoint{6.540024in}{0.695418in}}%
\pgfpathquadraticcurveto{\pgfqpoint{6.545014in}{0.692698in}}{\pgfqpoint{6.548274in}{0.687078in}}%
\pgfpathlineto{\pgfqpoint{6.548274in}{0.687078in}}%
\pgfpathclose%
\pgfpathmoveto{\pgfqpoint{6.512988in}{0.665067in}}%
\pgfpathquadraticcurveto{\pgfqpoint{6.512988in}{0.653648in}}{\pgfqpoint{6.517689in}{0.647145in}}%
\pgfpathquadraticcurveto{\pgfqpoint{6.522390in}{0.640643in}}{\pgfqpoint{6.530604in}{0.640643in}}%
\pgfpathquadraticcurveto{\pgfqpoint{6.538817in}{0.640643in}}{\pgfqpoint{6.543537in}{0.647145in}}%
\pgfpathquadraticcurveto{\pgfqpoint{6.548274in}{0.653648in}}{\pgfqpoint{6.548274in}{0.665067in}}%
\pgfpathquadraticcurveto{\pgfqpoint{6.548274in}{0.676487in}}{\pgfqpoint{6.543537in}{0.682989in}}%
\pgfpathquadraticcurveto{\pgfqpoint{6.538817in}{0.689492in}}{\pgfqpoint{6.530604in}{0.689492in}}%
\pgfpathquadraticcurveto{\pgfqpoint{6.522390in}{0.689492in}}{\pgfqpoint{6.517689in}{0.682989in}}%
\pgfpathquadraticcurveto{\pgfqpoint{6.512988in}{0.676487in}}{\pgfqpoint{6.512988in}{0.665067in}}%
\pgfpathlineto{\pgfqpoint{6.512988in}{0.665067in}}%
\pgfpathclose%
\pgfpathmoveto{\pgfqpoint{6.633904in}{0.667715in}}%
\pgfpathlineto{\pgfqpoint{6.633904in}{0.662654in}}%
\pgfpathlineto{\pgfqpoint{6.586280in}{0.662654in}}%
\pgfpathquadraticcurveto{\pgfqpoint{6.586964in}{0.651954in}}{\pgfqpoint{6.592728in}{0.646353in}}%
\pgfpathquadraticcurveto{\pgfqpoint{6.598492in}{0.640751in}}{\pgfqpoint{6.608795in}{0.640751in}}%
\pgfpathquadraticcurveto{\pgfqpoint{6.614757in}{0.640751in}}{\pgfqpoint{6.620359in}{0.642210in}}%
\pgfpathquadraticcurveto{\pgfqpoint{6.625960in}{0.643669in}}{\pgfqpoint{6.631490in}{0.646605in}}%
\pgfpathlineto{\pgfqpoint{6.631490in}{0.636806in}}%
\pgfpathquadraticcurveto{\pgfqpoint{6.625906in}{0.634447in}}{\pgfqpoint{6.620052in}{0.633204in}}%
\pgfpathquadraticcurveto{\pgfqpoint{6.614198in}{0.631961in}}{\pgfqpoint{6.608182in}{0.631961in}}%
\pgfpathquadraticcurveto{\pgfqpoint{6.593088in}{0.631961in}}{\pgfqpoint{6.584280in}{0.640733in}}%
\pgfpathquadraticcurveto{\pgfqpoint{6.575472in}{0.649523in}}{\pgfqpoint{6.575472in}{0.664509in}}%
\pgfpathquadraticcurveto{\pgfqpoint{6.575472in}{0.679981in}}{\pgfqpoint{6.583830in}{0.689059in}}%
\pgfpathquadraticcurveto{\pgfqpoint{6.592188in}{0.698156in}}{\pgfqpoint{6.606381in}{0.698156in}}%
\pgfpathquadraticcurveto{\pgfqpoint{6.619098in}{0.698156in}}{\pgfqpoint{6.626501in}{0.689960in}}%
\pgfpathquadraticcurveto{\pgfqpoint{6.633904in}{0.681783in}}{\pgfqpoint{6.633904in}{0.667715in}}%
\pgfpathlineto{\pgfqpoint{6.633904in}{0.667715in}}%
\pgfpathclose%
\pgfpathmoveto{\pgfqpoint{6.623547in}{0.670759in}}%
\pgfpathquadraticcurveto{\pgfqpoint{6.623439in}{0.679243in}}{\pgfqpoint{6.618791in}{0.684304in}}%
\pgfpathquadraticcurveto{\pgfqpoint{6.614144in}{0.689384in}}{\pgfqpoint{6.606489in}{0.689384in}}%
\pgfpathquadraticcurveto{\pgfqpoint{6.597825in}{0.689384in}}{\pgfqpoint{6.592620in}{0.684484in}}%
\pgfpathquadraticcurveto{\pgfqpoint{6.587414in}{0.679585in}}{\pgfqpoint{6.586622in}{0.670687in}}%
\pgfpathlineto{\pgfqpoint{6.623547in}{0.670759in}}%
\pgfpathlineto{\pgfqpoint{6.623547in}{0.670759in}}%
\pgfpathclose%
\pgfpathmoveto{\pgfqpoint{6.691092in}{0.694787in}}%
\pgfpathlineto{\pgfqpoint{6.691092in}{0.684989in}}%
\pgfpathquadraticcurveto{\pgfqpoint{6.686715in}{0.687240in}}{\pgfqpoint{6.681978in}{0.688357in}}%
\pgfpathquadraticcurveto{\pgfqpoint{6.677259in}{0.689492in}}{\pgfqpoint{6.672179in}{0.689492in}}%
\pgfpathquadraticcurveto{\pgfqpoint{6.664470in}{0.689492in}}{\pgfqpoint{6.660616in}{0.687132in}}%
\pgfpathquadraticcurveto{\pgfqpoint{6.656761in}{0.684773in}}{\pgfqpoint{6.656761in}{0.680035in}}%
\pgfpathquadraticcurveto{\pgfqpoint{6.656761in}{0.676433in}}{\pgfqpoint{6.659517in}{0.674380in}}%
\pgfpathquadraticcurveto{\pgfqpoint{6.662273in}{0.672326in}}{\pgfqpoint{6.670612in}{0.670471in}}%
\pgfpathlineto{\pgfqpoint{6.674161in}{0.669678in}}%
\pgfpathquadraticcurveto{\pgfqpoint{6.685184in}{0.667319in}}{\pgfqpoint{6.689831in}{0.663014in}}%
\pgfpathquadraticcurveto{\pgfqpoint{6.694479in}{0.658709in}}{\pgfqpoint{6.694479in}{0.651000in}}%
\pgfpathquadraticcurveto{\pgfqpoint{6.694479in}{0.642210in}}{\pgfqpoint{6.687526in}{0.637076in}}%
\pgfpathquadraticcurveto{\pgfqpoint{6.680573in}{0.631961in}}{\pgfqpoint{6.668415in}{0.631961in}}%
\pgfpathquadraticcurveto{\pgfqpoint{6.663354in}{0.631961in}}{\pgfqpoint{6.657860in}{0.632952in}}%
\pgfpathquadraticcurveto{\pgfqpoint{6.652366in}{0.633942in}}{\pgfqpoint{6.646296in}{0.635906in}}%
\pgfpathlineto{\pgfqpoint{6.646296in}{0.646605in}}%
\pgfpathquadraticcurveto{\pgfqpoint{6.652042in}{0.643615in}}{\pgfqpoint{6.657608in}{0.642120in}}%
\pgfpathquadraticcurveto{\pgfqpoint{6.663173in}{0.640643in}}{\pgfqpoint{6.668649in}{0.640643in}}%
\pgfpathquadraticcurveto{\pgfqpoint{6.675962in}{0.640643in}}{\pgfqpoint{6.679889in}{0.643146in}}%
\pgfpathquadraticcurveto{\pgfqpoint{6.683833in}{0.645650in}}{\pgfqpoint{6.683833in}{0.650207in}}%
\pgfpathquadraticcurveto{\pgfqpoint{6.683833in}{0.654422in}}{\pgfqpoint{6.680987in}{0.656674in}}%
\pgfpathquadraticcurveto{\pgfqpoint{6.678160in}{0.658925in}}{\pgfqpoint{6.668523in}{0.661014in}}%
\pgfpathlineto{\pgfqpoint{6.664921in}{0.661861in}}%
\pgfpathquadraticcurveto{\pgfqpoint{6.655302in}{0.663878in}}{\pgfqpoint{6.651015in}{0.668075in}}%
\pgfpathquadraticcurveto{\pgfqpoint{6.646746in}{0.672272in}}{\pgfqpoint{6.646746in}{0.679585in}}%
\pgfpathquadraticcurveto{\pgfqpoint{6.646746in}{0.688483in}}{\pgfqpoint{6.653051in}{0.693310in}}%
\pgfpathquadraticcurveto{\pgfqpoint{6.659355in}{0.698156in}}{\pgfqpoint{6.670955in}{0.698156in}}%
\pgfpathquadraticcurveto{\pgfqpoint{6.676683in}{0.698156in}}{\pgfqpoint{6.681744in}{0.697309in}}%
\pgfpathquadraticcurveto{\pgfqpoint{6.686823in}{0.696480in}}{\pgfqpoint{6.691092in}{0.694787in}}%
\pgfpathlineto{\pgfqpoint{6.691092in}{0.694787in}}%
\pgfpathclose%
\pgfpathmoveto{\pgfqpoint{6.710960in}{0.696643in}}%
\pgfpathlineto{\pgfqpoint{6.721317in}{0.696643in}}%
\pgfpathlineto{\pgfqpoint{6.721317in}{0.633600in}}%
\pgfpathlineto{\pgfqpoint{6.710960in}{0.633600in}}%
\pgfpathlineto{\pgfqpoint{6.710960in}{0.696643in}}%
\pgfpathlineto{\pgfqpoint{6.710960in}{0.696643in}}%
\pgfpathclose%
\pgfpathmoveto{\pgfqpoint{6.710960in}{0.721193in}}%
\pgfpathlineto{\pgfqpoint{6.721317in}{0.721193in}}%
\pgfpathlineto{\pgfqpoint{6.721317in}{0.708062in}}%
\pgfpathlineto{\pgfqpoint{6.710960in}{0.708062in}}%
\pgfpathlineto{\pgfqpoint{6.710960in}{0.721193in}}%
\pgfpathlineto{\pgfqpoint{6.710960in}{0.721193in}}%
\pgfpathclose%
\pgfpathmoveto{\pgfqpoint{6.784467in}{0.665860in}}%
\pgfpathquadraticcurveto{\pgfqpoint{6.784467in}{0.677117in}}{\pgfqpoint{6.779820in}{0.683296in}}%
\pgfpathquadraticcurveto{\pgfqpoint{6.775191in}{0.689492in}}{\pgfqpoint{6.766797in}{0.689492in}}%
\pgfpathquadraticcurveto{\pgfqpoint{6.758476in}{0.689492in}}{\pgfqpoint{6.753829in}{0.683296in}}%
\pgfpathquadraticcurveto{\pgfqpoint{6.749181in}{0.677117in}}{\pgfqpoint{6.749181in}{0.665860in}}%
\pgfpathquadraticcurveto{\pgfqpoint{6.749181in}{0.654656in}}{\pgfqpoint{6.753829in}{0.648460in}}%
\pgfpathquadraticcurveto{\pgfqpoint{6.758476in}{0.642264in}}{\pgfqpoint{6.766797in}{0.642264in}}%
\pgfpathquadraticcurveto{\pgfqpoint{6.775191in}{0.642264in}}{\pgfqpoint{6.779820in}{0.648460in}}%
\pgfpathquadraticcurveto{\pgfqpoint{6.784467in}{0.654656in}}{\pgfqpoint{6.784467in}{0.665860in}}%
\pgfpathlineto{\pgfqpoint{6.784467in}{0.665860in}}%
\pgfpathclose%
\pgfpathmoveto{\pgfqpoint{6.794824in}{0.641417in}}%
\pgfpathquadraticcurveto{\pgfqpoint{6.794824in}{0.625332in}}{\pgfqpoint{6.787673in}{0.617479in}}%
\pgfpathquadraticcurveto{\pgfqpoint{6.780541in}{0.609626in}}{\pgfqpoint{6.765789in}{0.609626in}}%
\pgfpathquadraticcurveto{\pgfqpoint{6.760331in}{0.609626in}}{\pgfqpoint{6.755486in}{0.610436in}}%
\pgfpathquadraticcurveto{\pgfqpoint{6.750640in}{0.611247in}}{\pgfqpoint{6.746083in}{0.612940in}}%
\pgfpathlineto{\pgfqpoint{6.746083in}{0.623009in}}%
\pgfpathquadraticcurveto{\pgfqpoint{6.750640in}{0.620541in}}{\pgfqpoint{6.755089in}{0.619370in}}%
\pgfpathquadraticcurveto{\pgfqpoint{6.759538in}{0.618182in}}{\pgfqpoint{6.764150in}{0.618182in}}%
\pgfpathquadraticcurveto{\pgfqpoint{6.774344in}{0.618182in}}{\pgfqpoint{6.779406in}{0.623495in}}%
\pgfpathquadraticcurveto{\pgfqpoint{6.784467in}{0.628809in}}{\pgfqpoint{6.784467in}{0.639562in}}%
\pgfpathlineto{\pgfqpoint{6.784467in}{0.644695in}}%
\pgfpathquadraticcurveto{\pgfqpoint{6.781261in}{0.639112in}}{\pgfqpoint{6.776254in}{0.636356in}}%
\pgfpathquadraticcurveto{\pgfqpoint{6.771246in}{0.633600in}}{\pgfqpoint{6.764258in}{0.633600in}}%
\pgfpathquadraticcurveto{\pgfqpoint{6.752676in}{0.633600in}}{\pgfqpoint{6.745579in}{0.642426in}}%
\pgfpathquadraticcurveto{\pgfqpoint{6.738482in}{0.651270in}}{\pgfqpoint{6.738482in}{0.665860in}}%
\pgfpathquadraticcurveto{\pgfqpoint{6.738482in}{0.680486in}}{\pgfqpoint{6.745579in}{0.689312in}}%
\pgfpathquadraticcurveto{\pgfqpoint{6.752676in}{0.698156in}}{\pgfqpoint{6.764258in}{0.698156in}}%
\pgfpathquadraticcurveto{\pgfqpoint{6.771246in}{0.698156in}}{\pgfqpoint{6.776254in}{0.695400in}}%
\pgfpathquadraticcurveto{\pgfqpoint{6.781261in}{0.692644in}}{\pgfqpoint{6.784467in}{0.687078in}}%
\pgfpathlineto{\pgfqpoint{6.784467in}{0.696643in}}%
\pgfpathlineto{\pgfqpoint{6.794824in}{0.696643in}}%
\pgfpathlineto{\pgfqpoint{6.794824in}{0.641417in}}%
\pgfpathlineto{\pgfqpoint{6.794824in}{0.641417in}}%
\pgfpathclose%
\pgfpathmoveto{\pgfqpoint{6.868584in}{0.671660in}}%
\pgfpathlineto{\pgfqpoint{6.868584in}{0.633600in}}%
\pgfpathlineto{\pgfqpoint{6.858227in}{0.633600in}}%
\pgfpathlineto{\pgfqpoint{6.858227in}{0.671317in}}%
\pgfpathquadraticcurveto{\pgfqpoint{6.858227in}{0.680269in}}{\pgfqpoint{6.854733in}{0.684700in}}%
\pgfpathquadraticcurveto{\pgfqpoint{6.851238in}{0.689149in}}{\pgfqpoint{6.844268in}{0.689149in}}%
\pgfpathquadraticcurveto{\pgfqpoint{6.835874in}{0.689149in}}{\pgfqpoint{6.831029in}{0.683800in}}%
\pgfpathquadraticcurveto{\pgfqpoint{6.826183in}{0.678468in}}{\pgfqpoint{6.826183in}{0.669228in}}%
\pgfpathlineto{\pgfqpoint{6.826183in}{0.633600in}}%
\pgfpathlineto{\pgfqpoint{6.815772in}{0.633600in}}%
\pgfpathlineto{\pgfqpoint{6.815772in}{0.696643in}}%
\pgfpathlineto{\pgfqpoint{6.826183in}{0.696643in}}%
\pgfpathlineto{\pgfqpoint{6.826183in}{0.686844in}}%
\pgfpathquadraticcurveto{\pgfqpoint{6.829912in}{0.692536in}}{\pgfqpoint{6.834937in}{0.695346in}}%
\pgfpathquadraticcurveto{\pgfqpoint{6.839981in}{0.698156in}}{\pgfqpoint{6.846573in}{0.698156in}}%
\pgfpathquadraticcurveto{\pgfqpoint{6.857434in}{0.698156in}}{\pgfqpoint{6.863000in}{0.691437in}}%
\pgfpathquadraticcurveto{\pgfqpoint{6.868584in}{0.684718in}}{\pgfqpoint{6.868584in}{0.671660in}}%
\pgfpathlineto{\pgfqpoint{6.868584in}{0.671660in}}%
\pgfpathclose%
\pgfusepath{fill}%
\end{pgfscope}%
\begin{pgfscope}%
\pgfsetbuttcap%
\pgfsetroundjoin%
\definecolor{currentfill}{rgb}{0.309804,0.372549,0.419608}%
\pgfsetfillcolor{currentfill}%
\pgfsetlinewidth{0.000000pt}%
\definecolor{currentstroke}{rgb}{0.309804,0.372549,0.419608}%
\pgfsetstrokecolor{currentstroke}%
\pgfsetdash{}{0pt}%
\pgfpathmoveto{\pgfqpoint{7.005850in}{0.687078in}}%
\pgfpathlineto{\pgfqpoint{7.005850in}{0.721193in}}%
\pgfpathlineto{\pgfqpoint{7.016207in}{0.721193in}}%
\pgfpathlineto{\pgfqpoint{7.016207in}{0.633600in}}%
\pgfpathlineto{\pgfqpoint{7.005850in}{0.633600in}}%
\pgfpathlineto{\pgfqpoint{7.005850in}{0.643056in}}%
\pgfpathquadraticcurveto{\pgfqpoint{7.002590in}{0.637437in}}{\pgfqpoint{6.997601in}{0.634699in}}%
\pgfpathquadraticcurveto{\pgfqpoint{6.992629in}{0.631961in}}{\pgfqpoint{6.985641in}{0.631961in}}%
\pgfpathquadraticcurveto{\pgfqpoint{6.974221in}{0.631961in}}{\pgfqpoint{6.967034in}{0.641075in}}%
\pgfpathquadraticcurveto{\pgfqpoint{6.959865in}{0.650207in}}{\pgfqpoint{6.959865in}{0.665067in}}%
\pgfpathquadraticcurveto{\pgfqpoint{6.959865in}{0.679927in}}{\pgfqpoint{6.967034in}{0.689041in}}%
\pgfpathquadraticcurveto{\pgfqpoint{6.974221in}{0.698156in}}{\pgfqpoint{6.985641in}{0.698156in}}%
\pgfpathquadraticcurveto{\pgfqpoint{6.992629in}{0.698156in}}{\pgfqpoint{6.997601in}{0.695418in}}%
\pgfpathquadraticcurveto{\pgfqpoint{7.002590in}{0.692698in}}{\pgfqpoint{7.005850in}{0.687078in}}%
\pgfpathlineto{\pgfqpoint{7.005850in}{0.687078in}}%
\pgfpathclose%
\pgfpathmoveto{\pgfqpoint{6.970564in}{0.665067in}}%
\pgfpathquadraticcurveto{\pgfqpoint{6.970564in}{0.653648in}}{\pgfqpoint{6.975266in}{0.647145in}}%
\pgfpathquadraticcurveto{\pgfqpoint{6.979967in}{0.640643in}}{\pgfqpoint{6.988180in}{0.640643in}}%
\pgfpathquadraticcurveto{\pgfqpoint{6.996394in}{0.640643in}}{\pgfqpoint{7.001113in}{0.647145in}}%
\pgfpathquadraticcurveto{\pgfqpoint{7.005850in}{0.653648in}}{\pgfqpoint{7.005850in}{0.665067in}}%
\pgfpathquadraticcurveto{\pgfqpoint{7.005850in}{0.676487in}}{\pgfqpoint{7.001113in}{0.682989in}}%
\pgfpathquadraticcurveto{\pgfqpoint{6.996394in}{0.689492in}}{\pgfqpoint{6.988180in}{0.689492in}}%
\pgfpathquadraticcurveto{\pgfqpoint{6.979967in}{0.689492in}}{\pgfqpoint{6.975266in}{0.682989in}}%
\pgfpathquadraticcurveto{\pgfqpoint{6.970564in}{0.676487in}}{\pgfqpoint{6.970564in}{0.665067in}}%
\pgfpathlineto{\pgfqpoint{6.970564in}{0.665067in}}%
\pgfpathclose%
\pgfpathmoveto{\pgfqpoint{7.037552in}{0.696643in}}%
\pgfpathlineto{\pgfqpoint{7.047909in}{0.696643in}}%
\pgfpathlineto{\pgfqpoint{7.047909in}{0.633600in}}%
\pgfpathlineto{\pgfqpoint{7.037552in}{0.633600in}}%
\pgfpathlineto{\pgfqpoint{7.037552in}{0.696643in}}%
\pgfpathlineto{\pgfqpoint{7.037552in}{0.696643in}}%
\pgfpathclose%
\pgfpathmoveto{\pgfqpoint{7.037552in}{0.721193in}}%
\pgfpathlineto{\pgfqpoint{7.047909in}{0.721193in}}%
\pgfpathlineto{\pgfqpoint{7.047909in}{0.708062in}}%
\pgfpathlineto{\pgfqpoint{7.037552in}{0.708062in}}%
\pgfpathlineto{\pgfqpoint{7.037552in}{0.721193in}}%
\pgfpathlineto{\pgfqpoint{7.037552in}{0.721193in}}%
\pgfpathclose%
\pgfpathmoveto{\pgfqpoint{7.098235in}{0.665283in}}%
\pgfpathquadraticcurveto{\pgfqpoint{7.085680in}{0.665283in}}{\pgfqpoint{7.080835in}{0.662419in}}%
\pgfpathquadraticcurveto{\pgfqpoint{7.075990in}{0.659556in}}{\pgfqpoint{7.075990in}{0.652621in}}%
\pgfpathquadraticcurveto{\pgfqpoint{7.075990in}{0.647109in}}{\pgfqpoint{7.079628in}{0.643867in}}%
\pgfpathquadraticcurveto{\pgfqpoint{7.083266in}{0.640643in}}{\pgfqpoint{7.089499in}{0.640643in}}%
\pgfpathquadraticcurveto{\pgfqpoint{7.098127in}{0.640643in}}{\pgfqpoint{7.103332in}{0.646749in}}%
\pgfpathquadraticcurveto{\pgfqpoint{7.108538in}{0.652855in}}{\pgfqpoint{7.108538in}{0.662978in}}%
\pgfpathlineto{\pgfqpoint{7.108538in}{0.665283in}}%
\pgfpathlineto{\pgfqpoint{7.098235in}{0.665283in}}%
\pgfpathlineto{\pgfqpoint{7.098235in}{0.665283in}}%
\pgfpathclose%
\pgfpathmoveto{\pgfqpoint{7.118895in}{0.669570in}}%
\pgfpathlineto{\pgfqpoint{7.118895in}{0.633600in}}%
\pgfpathlineto{\pgfqpoint{7.108538in}{0.633600in}}%
\pgfpathlineto{\pgfqpoint{7.108538in}{0.643164in}}%
\pgfpathquadraticcurveto{\pgfqpoint{7.104989in}{0.637437in}}{\pgfqpoint{7.099694in}{0.634699in}}%
\pgfpathquadraticcurveto{\pgfqpoint{7.094398in}{0.631961in}}{\pgfqpoint{7.086743in}{0.631961in}}%
\pgfpathquadraticcurveto{\pgfqpoint{7.077070in}{0.631961in}}{\pgfqpoint{7.071342in}{0.637401in}}%
\pgfpathquadraticcurveto{\pgfqpoint{7.065633in}{0.642840in}}{\pgfqpoint{7.065633in}{0.651954in}}%
\pgfpathquadraticcurveto{\pgfqpoint{7.065633in}{0.662582in}}{\pgfqpoint{7.072747in}{0.667985in}}%
\pgfpathquadraticcurveto{\pgfqpoint{7.079880in}{0.673389in}}{\pgfqpoint{7.094002in}{0.673389in}}%
\pgfpathlineto{\pgfqpoint{7.108538in}{0.673389in}}%
\pgfpathlineto{\pgfqpoint{7.108538in}{0.674416in}}%
\pgfpathquadraticcurveto{\pgfqpoint{7.108538in}{0.681566in}}{\pgfqpoint{7.103836in}{0.685475in}}%
\pgfpathquadraticcurveto{\pgfqpoint{7.099135in}{0.689384in}}{\pgfqpoint{7.090633in}{0.689384in}}%
\pgfpathquadraticcurveto{\pgfqpoint{7.085230in}{0.689384in}}{\pgfqpoint{7.080096in}{0.688087in}}%
\pgfpathquadraticcurveto{\pgfqpoint{7.074981in}{0.686790in}}{\pgfqpoint{7.070262in}{0.684196in}}%
\pgfpathlineto{\pgfqpoint{7.070262in}{0.693779in}}%
\pgfpathquadraticcurveto{\pgfqpoint{7.075936in}{0.695976in}}{\pgfqpoint{7.081285in}{0.697057in}}%
\pgfpathquadraticcurveto{\pgfqpoint{7.086635in}{0.698156in}}{\pgfqpoint{7.091696in}{0.698156in}}%
\pgfpathquadraticcurveto{\pgfqpoint{7.105385in}{0.698156in}}{\pgfqpoint{7.112140in}{0.691059in}}%
\pgfpathquadraticcurveto{\pgfqpoint{7.118895in}{0.683980in}}{\pgfqpoint{7.118895in}{0.669570in}}%
\pgfpathlineto{\pgfqpoint{7.118895in}{0.669570in}}%
\pgfpathclose%
\pgfpathmoveto{\pgfqpoint{7.181703in}{0.665860in}}%
\pgfpathquadraticcurveto{\pgfqpoint{7.181703in}{0.677117in}}{\pgfqpoint{7.177056in}{0.683296in}}%
\pgfpathquadraticcurveto{\pgfqpoint{7.172427in}{0.689492in}}{\pgfqpoint{7.164033in}{0.689492in}}%
\pgfpathquadraticcurveto{\pgfqpoint{7.155711in}{0.689492in}}{\pgfqpoint{7.151064in}{0.683296in}}%
\pgfpathquadraticcurveto{\pgfqpoint{7.146417in}{0.677117in}}{\pgfqpoint{7.146417in}{0.665860in}}%
\pgfpathquadraticcurveto{\pgfqpoint{7.146417in}{0.654656in}}{\pgfqpoint{7.151064in}{0.648460in}}%
\pgfpathquadraticcurveto{\pgfqpoint{7.155711in}{0.642264in}}{\pgfqpoint{7.164033in}{0.642264in}}%
\pgfpathquadraticcurveto{\pgfqpoint{7.172427in}{0.642264in}}{\pgfqpoint{7.177056in}{0.648460in}}%
\pgfpathquadraticcurveto{\pgfqpoint{7.181703in}{0.654656in}}{\pgfqpoint{7.181703in}{0.665860in}}%
\pgfpathlineto{\pgfqpoint{7.181703in}{0.665860in}}%
\pgfpathclose%
\pgfpathmoveto{\pgfqpoint{7.192060in}{0.641417in}}%
\pgfpathquadraticcurveto{\pgfqpoint{7.192060in}{0.625332in}}{\pgfqpoint{7.184909in}{0.617479in}}%
\pgfpathquadraticcurveto{\pgfqpoint{7.177776in}{0.609626in}}{\pgfqpoint{7.163024in}{0.609626in}}%
\pgfpathquadraticcurveto{\pgfqpoint{7.157567in}{0.609626in}}{\pgfqpoint{7.152721in}{0.610436in}}%
\pgfpathquadraticcurveto{\pgfqpoint{7.147876in}{0.611247in}}{\pgfqpoint{7.143319in}{0.612940in}}%
\pgfpathlineto{\pgfqpoint{7.143319in}{0.623009in}}%
\pgfpathquadraticcurveto{\pgfqpoint{7.147876in}{0.620541in}}{\pgfqpoint{7.152325in}{0.619370in}}%
\pgfpathquadraticcurveto{\pgfqpoint{7.156774in}{0.618182in}}{\pgfqpoint{7.161385in}{0.618182in}}%
\pgfpathquadraticcurveto{\pgfqpoint{7.171580in}{0.618182in}}{\pgfqpoint{7.176641in}{0.623495in}}%
\pgfpathquadraticcurveto{\pgfqpoint{7.181703in}{0.628809in}}{\pgfqpoint{7.181703in}{0.639562in}}%
\pgfpathlineto{\pgfqpoint{7.181703in}{0.644695in}}%
\pgfpathquadraticcurveto{\pgfqpoint{7.178497in}{0.639112in}}{\pgfqpoint{7.173489in}{0.636356in}}%
\pgfpathquadraticcurveto{\pgfqpoint{7.168482in}{0.633600in}}{\pgfqpoint{7.161493in}{0.633600in}}%
\pgfpathquadraticcurveto{\pgfqpoint{7.149911in}{0.633600in}}{\pgfqpoint{7.142815in}{0.642426in}}%
\pgfpathquadraticcurveto{\pgfqpoint{7.135718in}{0.651270in}}{\pgfqpoint{7.135718in}{0.665860in}}%
\pgfpathquadraticcurveto{\pgfqpoint{7.135718in}{0.680486in}}{\pgfqpoint{7.142815in}{0.689312in}}%
\pgfpathquadraticcurveto{\pgfqpoint{7.149911in}{0.698156in}}{\pgfqpoint{7.161493in}{0.698156in}}%
\pgfpathquadraticcurveto{\pgfqpoint{7.168482in}{0.698156in}}{\pgfqpoint{7.173489in}{0.695400in}}%
\pgfpathquadraticcurveto{\pgfqpoint{7.178497in}{0.692644in}}{\pgfqpoint{7.181703in}{0.687078in}}%
\pgfpathlineto{\pgfqpoint{7.181703in}{0.696643in}}%
\pgfpathlineto{\pgfqpoint{7.192060in}{0.696643in}}%
\pgfpathlineto{\pgfqpoint{7.192060in}{0.641417in}}%
\pgfpathlineto{\pgfqpoint{7.192060in}{0.641417in}}%
\pgfpathclose%
\pgfpathmoveto{\pgfqpoint{7.265820in}{0.671660in}}%
\pgfpathlineto{\pgfqpoint{7.265820in}{0.633600in}}%
\pgfpathlineto{\pgfqpoint{7.255463in}{0.633600in}}%
\pgfpathlineto{\pgfqpoint{7.255463in}{0.671317in}}%
\pgfpathquadraticcurveto{\pgfqpoint{7.255463in}{0.680269in}}{\pgfqpoint{7.251968in}{0.684700in}}%
\pgfpathquadraticcurveto{\pgfqpoint{7.248474in}{0.689149in}}{\pgfqpoint{7.241503in}{0.689149in}}%
\pgfpathquadraticcurveto{\pgfqpoint{7.233110in}{0.689149in}}{\pgfqpoint{7.228264in}{0.683800in}}%
\pgfpathquadraticcurveto{\pgfqpoint{7.223419in}{0.678468in}}{\pgfqpoint{7.223419in}{0.669228in}}%
\pgfpathlineto{\pgfqpoint{7.223419in}{0.633600in}}%
\pgfpathlineto{\pgfqpoint{7.213008in}{0.633600in}}%
\pgfpathlineto{\pgfqpoint{7.213008in}{0.696643in}}%
\pgfpathlineto{\pgfqpoint{7.223419in}{0.696643in}}%
\pgfpathlineto{\pgfqpoint{7.223419in}{0.686844in}}%
\pgfpathquadraticcurveto{\pgfqpoint{7.227148in}{0.692536in}}{\pgfqpoint{7.232173in}{0.695346in}}%
\pgfpathquadraticcurveto{\pgfqpoint{7.237216in}{0.698156in}}{\pgfqpoint{7.243809in}{0.698156in}}%
\pgfpathquadraticcurveto{\pgfqpoint{7.254670in}{0.698156in}}{\pgfqpoint{7.260236in}{0.691437in}}%
\pgfpathquadraticcurveto{\pgfqpoint{7.265820in}{0.684718in}}{\pgfqpoint{7.265820in}{0.671660in}}%
\pgfpathlineto{\pgfqpoint{7.265820in}{0.671660in}}%
\pgfpathclose%
\pgfpathmoveto{\pgfqpoint{7.310886in}{0.689384in}}%
\pgfpathquadraticcurveto{\pgfqpoint{7.302564in}{0.689384in}}{\pgfqpoint{7.297719in}{0.682881in}}%
\pgfpathquadraticcurveto{\pgfqpoint{7.292874in}{0.676379in}}{\pgfqpoint{7.292874in}{0.665067in}}%
\pgfpathquadraticcurveto{\pgfqpoint{7.292874in}{0.653756in}}{\pgfqpoint{7.297683in}{0.647253in}}%
\pgfpathquadraticcurveto{\pgfqpoint{7.302510in}{0.640751in}}{\pgfqpoint{7.310886in}{0.640751in}}%
\pgfpathquadraticcurveto{\pgfqpoint{7.319172in}{0.640751in}}{\pgfqpoint{7.323999in}{0.647271in}}%
\pgfpathquadraticcurveto{\pgfqpoint{7.328844in}{0.653810in}}{\pgfqpoint{7.328844in}{0.665067in}}%
\pgfpathquadraticcurveto{\pgfqpoint{7.328844in}{0.676271in}}{\pgfqpoint{7.323999in}{0.682827in}}%
\pgfpathquadraticcurveto{\pgfqpoint{7.319172in}{0.689384in}}{\pgfqpoint{7.310886in}{0.689384in}}%
\pgfpathlineto{\pgfqpoint{7.310886in}{0.689384in}}%
\pgfpathclose%
\pgfpathmoveto{\pgfqpoint{7.310886in}{0.698156in}}%
\pgfpathquadraticcurveto{\pgfqpoint{7.324395in}{0.698156in}}{\pgfqpoint{7.332104in}{0.689366in}}%
\pgfpathquadraticcurveto{\pgfqpoint{7.339832in}{0.680594in}}{\pgfqpoint{7.339832in}{0.665067in}}%
\pgfpathquadraticcurveto{\pgfqpoint{7.339832in}{0.649595in}}{\pgfqpoint{7.332104in}{0.640769in}}%
\pgfpathquadraticcurveto{\pgfqpoint{7.324395in}{0.631961in}}{\pgfqpoint{7.310886in}{0.631961in}}%
\pgfpathquadraticcurveto{\pgfqpoint{7.297323in}{0.631961in}}{\pgfqpoint{7.289632in}{0.640769in}}%
\pgfpathquadraticcurveto{\pgfqpoint{7.281959in}{0.649595in}}{\pgfqpoint{7.281959in}{0.665067in}}%
\pgfpathquadraticcurveto{\pgfqpoint{7.281959in}{0.680594in}}{\pgfqpoint{7.289632in}{0.689366in}}%
\pgfpathquadraticcurveto{\pgfqpoint{7.297323in}{0.698156in}}{\pgfqpoint{7.310886in}{0.698156in}}%
\pgfpathlineto{\pgfqpoint{7.310886in}{0.698156in}}%
\pgfpathclose%
\pgfpathmoveto{\pgfqpoint{7.397182in}{0.694787in}}%
\pgfpathlineto{\pgfqpoint{7.397182in}{0.684989in}}%
\pgfpathquadraticcurveto{\pgfqpoint{7.392805in}{0.687240in}}{\pgfqpoint{7.388068in}{0.688357in}}%
\pgfpathquadraticcurveto{\pgfqpoint{7.383349in}{0.689492in}}{\pgfqpoint{7.378270in}{0.689492in}}%
\pgfpathquadraticcurveto{\pgfqpoint{7.370560in}{0.689492in}}{\pgfqpoint{7.366706in}{0.687132in}}%
\pgfpathquadraticcurveto{\pgfqpoint{7.362851in}{0.684773in}}{\pgfqpoint{7.362851in}{0.680035in}}%
\pgfpathquadraticcurveto{\pgfqpoint{7.362851in}{0.676433in}}{\pgfqpoint{7.365607in}{0.674380in}}%
\pgfpathquadraticcurveto{\pgfqpoint{7.368363in}{0.672326in}}{\pgfqpoint{7.376702in}{0.670471in}}%
\pgfpathlineto{\pgfqpoint{7.380251in}{0.669678in}}%
\pgfpathquadraticcurveto{\pgfqpoint{7.391274in}{0.667319in}}{\pgfqpoint{7.395921in}{0.663014in}}%
\pgfpathquadraticcurveto{\pgfqpoint{7.400569in}{0.658709in}}{\pgfqpoint{7.400569in}{0.651000in}}%
\pgfpathquadraticcurveto{\pgfqpoint{7.400569in}{0.642210in}}{\pgfqpoint{7.393616in}{0.637076in}}%
\pgfpathquadraticcurveto{\pgfqpoint{7.386663in}{0.631961in}}{\pgfqpoint{7.374505in}{0.631961in}}%
\pgfpathquadraticcurveto{\pgfqpoint{7.369444in}{0.631961in}}{\pgfqpoint{7.363950in}{0.632952in}}%
\pgfpathquadraticcurveto{\pgfqpoint{7.358456in}{0.633942in}}{\pgfqpoint{7.352386in}{0.635906in}}%
\pgfpathlineto{\pgfqpoint{7.352386in}{0.646605in}}%
\pgfpathquadraticcurveto{\pgfqpoint{7.358132in}{0.643615in}}{\pgfqpoint{7.363698in}{0.642120in}}%
\pgfpathquadraticcurveto{\pgfqpoint{7.369263in}{0.640643in}}{\pgfqpoint{7.374739in}{0.640643in}}%
\pgfpathquadraticcurveto{\pgfqpoint{7.382052in}{0.640643in}}{\pgfqpoint{7.385979in}{0.643146in}}%
\pgfpathquadraticcurveto{\pgfqpoint{7.389923in}{0.645650in}}{\pgfqpoint{7.389923in}{0.650207in}}%
\pgfpathquadraticcurveto{\pgfqpoint{7.389923in}{0.654422in}}{\pgfqpoint{7.387077in}{0.656674in}}%
\pgfpathquadraticcurveto{\pgfqpoint{7.384250in}{0.658925in}}{\pgfqpoint{7.374613in}{0.661014in}}%
\pgfpathlineto{\pgfqpoint{7.371011in}{0.661861in}}%
\pgfpathquadraticcurveto{\pgfqpoint{7.361392in}{0.663878in}}{\pgfqpoint{7.357105in}{0.668075in}}%
\pgfpathquadraticcurveto{\pgfqpoint{7.352836in}{0.672272in}}{\pgfqpoint{7.352836in}{0.679585in}}%
\pgfpathquadraticcurveto{\pgfqpoint{7.352836in}{0.688483in}}{\pgfqpoint{7.359141in}{0.693310in}}%
\pgfpathquadraticcurveto{\pgfqpoint{7.365445in}{0.698156in}}{\pgfqpoint{7.377045in}{0.698156in}}%
\pgfpathquadraticcurveto{\pgfqpoint{7.382773in}{0.698156in}}{\pgfqpoint{7.387834in}{0.697309in}}%
\pgfpathquadraticcurveto{\pgfqpoint{7.392913in}{0.696480in}}{\pgfqpoint{7.397182in}{0.694787in}}%
\pgfpathlineto{\pgfqpoint{7.397182in}{0.694787in}}%
\pgfpathclose%
\pgfpathmoveto{\pgfqpoint{7.427299in}{0.714547in}}%
\pgfpathlineto{\pgfqpoint{7.427299in}{0.696643in}}%
\pgfpathlineto{\pgfqpoint{7.448625in}{0.696643in}}%
\pgfpathlineto{\pgfqpoint{7.448625in}{0.688591in}}%
\pgfpathlineto{\pgfqpoint{7.427299in}{0.688591in}}%
\pgfpathlineto{\pgfqpoint{7.427299in}{0.654368in}}%
\pgfpathquadraticcurveto{\pgfqpoint{7.427299in}{0.646659in}}{\pgfqpoint{7.429406in}{0.644461in}}%
\pgfpathquadraticcurveto{\pgfqpoint{7.431513in}{0.642264in}}{\pgfqpoint{7.437998in}{0.642264in}}%
\pgfpathlineto{\pgfqpoint{7.448625in}{0.642264in}}%
\pgfpathlineto{\pgfqpoint{7.448625in}{0.633600in}}%
\pgfpathlineto{\pgfqpoint{7.437998in}{0.633600in}}%
\pgfpathquadraticcurveto{\pgfqpoint{7.426002in}{0.633600in}}{\pgfqpoint{7.421445in}{0.638067in}}%
\pgfpathquadraticcurveto{\pgfqpoint{7.416888in}{0.642552in}}{\pgfqpoint{7.416888in}{0.654368in}}%
\pgfpathlineto{\pgfqpoint{7.416888in}{0.688591in}}%
\pgfpathlineto{\pgfqpoint{7.409286in}{0.688591in}}%
\pgfpathlineto{\pgfqpoint{7.409286in}{0.696643in}}%
\pgfpathlineto{\pgfqpoint{7.416888in}{0.696643in}}%
\pgfpathlineto{\pgfqpoint{7.416888in}{0.714547in}}%
\pgfpathlineto{\pgfqpoint{7.427299in}{0.714547in}}%
\pgfpathlineto{\pgfqpoint{7.427299in}{0.714547in}}%
\pgfpathclose%
\pgfpathmoveto{\pgfqpoint{7.462242in}{0.696643in}}%
\pgfpathlineto{\pgfqpoint{7.472599in}{0.696643in}}%
\pgfpathlineto{\pgfqpoint{7.472599in}{0.633600in}}%
\pgfpathlineto{\pgfqpoint{7.462242in}{0.633600in}}%
\pgfpathlineto{\pgfqpoint{7.462242in}{0.696643in}}%
\pgfpathlineto{\pgfqpoint{7.462242in}{0.696643in}}%
\pgfpathclose%
\pgfpathmoveto{\pgfqpoint{7.462242in}{0.721193in}}%
\pgfpathlineto{\pgfqpoint{7.472599in}{0.721193in}}%
\pgfpathlineto{\pgfqpoint{7.472599in}{0.708062in}}%
\pgfpathlineto{\pgfqpoint{7.462242in}{0.708062in}}%
\pgfpathlineto{\pgfqpoint{7.462242in}{0.721193in}}%
\pgfpathlineto{\pgfqpoint{7.462242in}{0.721193in}}%
\pgfpathclose%
\pgfpathmoveto{\pgfqpoint{7.539640in}{0.694229in}}%
\pgfpathlineto{\pgfqpoint{7.539640in}{0.684538in}}%
\pgfpathquadraticcurveto{\pgfqpoint{7.535245in}{0.686970in}}{\pgfqpoint{7.530832in}{0.688177in}}%
\pgfpathquadraticcurveto{\pgfqpoint{7.526419in}{0.689384in}}{\pgfqpoint{7.521916in}{0.689384in}}%
\pgfpathquadraticcurveto{\pgfqpoint{7.511830in}{0.689384in}}{\pgfqpoint{7.506246in}{0.682989in}}%
\pgfpathquadraticcurveto{\pgfqpoint{7.500680in}{0.676613in}}{\pgfqpoint{7.500680in}{0.665067in}}%
\pgfpathquadraticcurveto{\pgfqpoint{7.500680in}{0.653521in}}{\pgfqpoint{7.506246in}{0.647127in}}%
\pgfpathquadraticcurveto{\pgfqpoint{7.511830in}{0.640751in}}{\pgfqpoint{7.521916in}{0.640751in}}%
\pgfpathquadraticcurveto{\pgfqpoint{7.526419in}{0.640751in}}{\pgfqpoint{7.530832in}{0.641958in}}%
\pgfpathquadraticcurveto{\pgfqpoint{7.535245in}{0.643164in}}{\pgfqpoint{7.539640in}{0.645596in}}%
\pgfpathlineto{\pgfqpoint{7.539640in}{0.636014in}}%
\pgfpathquadraticcurveto{\pgfqpoint{7.535299in}{0.633996in}}{\pgfqpoint{7.530652in}{0.632988in}}%
\pgfpathquadraticcurveto{\pgfqpoint{7.526023in}{0.631961in}}{\pgfqpoint{7.520782in}{0.631961in}}%
\pgfpathquadraticcurveto{\pgfqpoint{7.506534in}{0.631961in}}{\pgfqpoint{7.498140in}{0.640913in}}%
\pgfpathquadraticcurveto{\pgfqpoint{7.489765in}{0.649865in}}{\pgfqpoint{7.489765in}{0.665067in}}%
\pgfpathquadraticcurveto{\pgfqpoint{7.489765in}{0.680486in}}{\pgfqpoint{7.498230in}{0.689312in}}%
\pgfpathquadraticcurveto{\pgfqpoint{7.506714in}{0.698156in}}{\pgfqpoint{7.521466in}{0.698156in}}%
\pgfpathquadraticcurveto{\pgfqpoint{7.526239in}{0.698156in}}{\pgfqpoint{7.530796in}{0.697165in}}%
\pgfpathquadraticcurveto{\pgfqpoint{7.535354in}{0.696192in}}{\pgfqpoint{7.539640in}{0.694229in}}%
\pgfpathlineto{\pgfqpoint{7.539640in}{0.694229in}}%
\pgfpathclose%
\pgfpathmoveto{\pgfqpoint{7.560300in}{0.647902in}}%
\pgfpathlineto{\pgfqpoint{7.572170in}{0.647902in}}%
\pgfpathlineto{\pgfqpoint{7.572170in}{0.638211in}}%
\pgfpathlineto{\pgfqpoint{7.562948in}{0.620199in}}%
\pgfpathlineto{\pgfqpoint{7.555689in}{0.620199in}}%
\pgfpathlineto{\pgfqpoint{7.560300in}{0.638211in}}%
\pgfpathlineto{\pgfqpoint{7.560300in}{0.647902in}}%
\pgfpathlineto{\pgfqpoint{7.560300in}{0.647902in}}%
\pgfpathclose%
\pgfusepath{fill}%
\end{pgfscope}%
\begin{pgfscope}%
\pgfsetbuttcap%
\pgfsetroundjoin%
\definecolor{currentfill}{rgb}{0.309804,0.372549,0.419608}%
\pgfsetfillcolor{currentfill}%
\pgfsetlinewidth{0.000000pt}%
\definecolor{currentstroke}{rgb}{0.309804,0.372549,0.419608}%
\pgfsetstrokecolor{currentstroke}%
\pgfsetdash{}{0pt}%
\pgfpathmoveto{\pgfqpoint{7.734360in}{0.671660in}}%
\pgfpathlineto{\pgfqpoint{7.734360in}{0.633600in}}%
\pgfpathlineto{\pgfqpoint{7.724003in}{0.633600in}}%
\pgfpathlineto{\pgfqpoint{7.724003in}{0.671317in}}%
\pgfpathquadraticcurveto{\pgfqpoint{7.724003in}{0.680269in}}{\pgfqpoint{7.720509in}{0.684700in}}%
\pgfpathquadraticcurveto{\pgfqpoint{7.717014in}{0.689149in}}{\pgfqpoint{7.710044in}{0.689149in}}%
\pgfpathquadraticcurveto{\pgfqpoint{7.701650in}{0.689149in}}{\pgfqpoint{7.696805in}{0.683800in}}%
\pgfpathquadraticcurveto{\pgfqpoint{7.691959in}{0.678468in}}{\pgfqpoint{7.691959in}{0.669228in}}%
\pgfpathlineto{\pgfqpoint{7.691959in}{0.633600in}}%
\pgfpathlineto{\pgfqpoint{7.681548in}{0.633600in}}%
\pgfpathlineto{\pgfqpoint{7.681548in}{0.696643in}}%
\pgfpathlineto{\pgfqpoint{7.691959in}{0.696643in}}%
\pgfpathlineto{\pgfqpoint{7.691959in}{0.686844in}}%
\pgfpathquadraticcurveto{\pgfqpoint{7.695688in}{0.692536in}}{\pgfqpoint{7.700713in}{0.695346in}}%
\pgfpathquadraticcurveto{\pgfqpoint{7.705757in}{0.698156in}}{\pgfqpoint{7.712349in}{0.698156in}}%
\pgfpathquadraticcurveto{\pgfqpoint{7.723211in}{0.698156in}}{\pgfqpoint{7.728776in}{0.691437in}}%
\pgfpathquadraticcurveto{\pgfqpoint{7.734360in}{0.684718in}}{\pgfqpoint{7.734360in}{0.671660in}}%
\pgfpathlineto{\pgfqpoint{7.734360in}{0.671660in}}%
\pgfpathclose%
\pgfpathmoveto{\pgfqpoint{7.779426in}{0.689384in}}%
\pgfpathquadraticcurveto{\pgfqpoint{7.771105in}{0.689384in}}{\pgfqpoint{7.766260in}{0.682881in}}%
\pgfpathquadraticcurveto{\pgfqpoint{7.761414in}{0.676379in}}{\pgfqpoint{7.761414in}{0.665067in}}%
\pgfpathquadraticcurveto{\pgfqpoint{7.761414in}{0.653756in}}{\pgfqpoint{7.766224in}{0.647253in}}%
\pgfpathquadraticcurveto{\pgfqpoint{7.771051in}{0.640751in}}{\pgfqpoint{7.779426in}{0.640751in}}%
\pgfpathquadraticcurveto{\pgfqpoint{7.787712in}{0.640751in}}{\pgfqpoint{7.792539in}{0.647271in}}%
\pgfpathquadraticcurveto{\pgfqpoint{7.797385in}{0.653810in}}{\pgfqpoint{7.797385in}{0.665067in}}%
\pgfpathquadraticcurveto{\pgfqpoint{7.797385in}{0.676271in}}{\pgfqpoint{7.792539in}{0.682827in}}%
\pgfpathquadraticcurveto{\pgfqpoint{7.787712in}{0.689384in}}{\pgfqpoint{7.779426in}{0.689384in}}%
\pgfpathlineto{\pgfqpoint{7.779426in}{0.689384in}}%
\pgfpathclose%
\pgfpathmoveto{\pgfqpoint{7.779426in}{0.698156in}}%
\pgfpathquadraticcurveto{\pgfqpoint{7.792936in}{0.698156in}}{\pgfqpoint{7.800645in}{0.689366in}}%
\pgfpathquadraticcurveto{\pgfqpoint{7.808372in}{0.680594in}}{\pgfqpoint{7.808372in}{0.665067in}}%
\pgfpathquadraticcurveto{\pgfqpoint{7.808372in}{0.649595in}}{\pgfqpoint{7.800645in}{0.640769in}}%
\pgfpathquadraticcurveto{\pgfqpoint{7.792936in}{0.631961in}}{\pgfqpoint{7.779426in}{0.631961in}}%
\pgfpathquadraticcurveto{\pgfqpoint{7.765863in}{0.631961in}}{\pgfqpoint{7.758172in}{0.640769in}}%
\pgfpathquadraticcurveto{\pgfqpoint{7.750499in}{0.649595in}}{\pgfqpoint{7.750499in}{0.665067in}}%
\pgfpathquadraticcurveto{\pgfqpoint{7.750499in}{0.680594in}}{\pgfqpoint{7.758172in}{0.689366in}}%
\pgfpathquadraticcurveto{\pgfqpoint{7.765863in}{0.698156in}}{\pgfqpoint{7.779426in}{0.698156in}}%
\pgfpathlineto{\pgfqpoint{7.779426in}{0.698156in}}%
\pgfpathclose%
\pgfpathmoveto{\pgfqpoint{7.835786in}{0.714547in}}%
\pgfpathlineto{\pgfqpoint{7.835786in}{0.696643in}}%
\pgfpathlineto{\pgfqpoint{7.857113in}{0.696643in}}%
\pgfpathlineto{\pgfqpoint{7.857113in}{0.688591in}}%
\pgfpathlineto{\pgfqpoint{7.835786in}{0.688591in}}%
\pgfpathlineto{\pgfqpoint{7.835786in}{0.654368in}}%
\pgfpathquadraticcurveto{\pgfqpoint{7.835786in}{0.646659in}}{\pgfqpoint{7.837894in}{0.644461in}}%
\pgfpathquadraticcurveto{\pgfqpoint{7.840001in}{0.642264in}}{\pgfqpoint{7.846486in}{0.642264in}}%
\pgfpathlineto{\pgfqpoint{7.857113in}{0.642264in}}%
\pgfpathlineto{\pgfqpoint{7.857113in}{0.633600in}}%
\pgfpathlineto{\pgfqpoint{7.846486in}{0.633600in}}%
\pgfpathquadraticcurveto{\pgfqpoint{7.834490in}{0.633600in}}{\pgfqpoint{7.829933in}{0.638067in}}%
\pgfpathquadraticcurveto{\pgfqpoint{7.825375in}{0.642552in}}{\pgfqpoint{7.825375in}{0.654368in}}%
\pgfpathlineto{\pgfqpoint{7.825375in}{0.688591in}}%
\pgfpathlineto{\pgfqpoint{7.817774in}{0.688591in}}%
\pgfpathlineto{\pgfqpoint{7.817774in}{0.696643in}}%
\pgfpathlineto{\pgfqpoint{7.825375in}{0.696643in}}%
\pgfpathlineto{\pgfqpoint{7.825375in}{0.714547in}}%
\pgfpathlineto{\pgfqpoint{7.835786in}{0.714547in}}%
\pgfpathlineto{\pgfqpoint{7.835786in}{0.714547in}}%
\pgfpathclose%
\pgfusepath{fill}%
\end{pgfscope}%
\begin{pgfscope}%
\pgfsetbuttcap%
\pgfsetroundjoin%
\definecolor{currentfill}{rgb}{0.309804,0.372549,0.419608}%
\pgfsetfillcolor{currentfill}%
\pgfsetlinewidth{0.000000pt}%
\definecolor{currentstroke}{rgb}{0.309804,0.372549,0.419608}%
\pgfsetstrokecolor{currentstroke}%
\pgfsetdash{}{0pt}%
\pgfpathmoveto{\pgfqpoint{7.968178in}{0.665283in}}%
\pgfpathquadraticcurveto{\pgfqpoint{7.955624in}{0.665283in}}{\pgfqpoint{7.950779in}{0.662419in}}%
\pgfpathquadraticcurveto{\pgfqpoint{7.945933in}{0.659556in}}{\pgfqpoint{7.945933in}{0.652621in}}%
\pgfpathquadraticcurveto{\pgfqpoint{7.945933in}{0.647109in}}{\pgfqpoint{7.949572in}{0.643867in}}%
\pgfpathquadraticcurveto{\pgfqpoint{7.953210in}{0.640643in}}{\pgfqpoint{7.959442in}{0.640643in}}%
\pgfpathquadraticcurveto{\pgfqpoint{7.968070in}{0.640643in}}{\pgfqpoint{7.973276in}{0.646749in}}%
\pgfpathquadraticcurveto{\pgfqpoint{7.978481in}{0.652855in}}{\pgfqpoint{7.978481in}{0.662978in}}%
\pgfpathlineto{\pgfqpoint{7.978481in}{0.665283in}}%
\pgfpathlineto{\pgfqpoint{7.968178in}{0.665283in}}%
\pgfpathlineto{\pgfqpoint{7.968178in}{0.665283in}}%
\pgfpathclose%
\pgfpathmoveto{\pgfqpoint{7.988838in}{0.669570in}}%
\pgfpathlineto{\pgfqpoint{7.988838in}{0.633600in}}%
\pgfpathlineto{\pgfqpoint{7.978481in}{0.633600in}}%
\pgfpathlineto{\pgfqpoint{7.978481in}{0.643164in}}%
\pgfpathquadraticcurveto{\pgfqpoint{7.974933in}{0.637437in}}{\pgfqpoint{7.969637in}{0.634699in}}%
\pgfpathquadraticcurveto{\pgfqpoint{7.964342in}{0.631961in}}{\pgfqpoint{7.956687in}{0.631961in}}%
\pgfpathquadraticcurveto{\pgfqpoint{7.947014in}{0.631961in}}{\pgfqpoint{7.941286in}{0.637401in}}%
\pgfpathquadraticcurveto{\pgfqpoint{7.935576in}{0.642840in}}{\pgfqpoint{7.935576in}{0.651954in}}%
\pgfpathquadraticcurveto{\pgfqpoint{7.935576in}{0.662582in}}{\pgfqpoint{7.942691in}{0.667985in}}%
\pgfpathquadraticcurveto{\pgfqpoint{7.949824in}{0.673389in}}{\pgfqpoint{7.963946in}{0.673389in}}%
\pgfpathlineto{\pgfqpoint{7.978481in}{0.673389in}}%
\pgfpathlineto{\pgfqpoint{7.978481in}{0.674416in}}%
\pgfpathquadraticcurveto{\pgfqpoint{7.978481in}{0.681566in}}{\pgfqpoint{7.973780in}{0.685475in}}%
\pgfpathquadraticcurveto{\pgfqpoint{7.969079in}{0.689384in}}{\pgfqpoint{7.960577in}{0.689384in}}%
\pgfpathquadraticcurveto{\pgfqpoint{7.955174in}{0.689384in}}{\pgfqpoint{7.950040in}{0.688087in}}%
\pgfpathquadraticcurveto{\pgfqpoint{7.944925in}{0.686790in}}{\pgfqpoint{7.940206in}{0.684196in}}%
\pgfpathlineto{\pgfqpoint{7.940206in}{0.693779in}}%
\pgfpathquadraticcurveto{\pgfqpoint{7.945879in}{0.695976in}}{\pgfqpoint{7.951229in}{0.697057in}}%
\pgfpathquadraticcurveto{\pgfqpoint{7.956579in}{0.698156in}}{\pgfqpoint{7.961640in}{0.698156in}}%
\pgfpathquadraticcurveto{\pgfqpoint{7.975329in}{0.698156in}}{\pgfqpoint{7.982084in}{0.691059in}}%
\pgfpathquadraticcurveto{\pgfqpoint{7.988838in}{0.683980in}}{\pgfqpoint{7.988838in}{0.669570in}}%
\pgfpathlineto{\pgfqpoint{7.988838in}{0.669570in}}%
\pgfpathclose%
\pgfusepath{fill}%
\end{pgfscope}%
\begin{pgfscope}%
\pgfsetbuttcap%
\pgfsetroundjoin%
\definecolor{currentfill}{rgb}{0.309804,0.372549,0.419608}%
\pgfsetfillcolor{currentfill}%
\pgfsetlinewidth{0.000000pt}%
\definecolor{currentstroke}{rgb}{0.309804,0.372549,0.419608}%
\pgfsetstrokecolor{currentstroke}%
\pgfsetdash{}{0pt}%
\pgfpathmoveto{\pgfqpoint{8.059676in}{0.696643in}}%
\pgfpathlineto{\pgfqpoint{8.070646in}{0.696643in}}%
\pgfpathlineto{\pgfqpoint{8.090351in}{0.643741in}}%
\pgfpathlineto{\pgfqpoint{8.110056in}{0.696643in}}%
\pgfpathlineto{\pgfqpoint{8.121026in}{0.696643in}}%
\pgfpathlineto{\pgfqpoint{8.097376in}{0.633600in}}%
\pgfpathlineto{\pgfqpoint{8.083308in}{0.633600in}}%
\pgfpathlineto{\pgfqpoint{8.059676in}{0.696643in}}%
\pgfpathlineto{\pgfqpoint{8.059676in}{0.696643in}}%
\pgfpathclose%
\pgfpathmoveto{\pgfqpoint{8.163985in}{0.665283in}}%
\pgfpathquadraticcurveto{\pgfqpoint{8.151430in}{0.665283in}}{\pgfqpoint{8.146585in}{0.662419in}}%
\pgfpathquadraticcurveto{\pgfqpoint{8.141740in}{0.659556in}}{\pgfqpoint{8.141740in}{0.652621in}}%
\pgfpathquadraticcurveto{\pgfqpoint{8.141740in}{0.647109in}}{\pgfqpoint{8.145378in}{0.643867in}}%
\pgfpathquadraticcurveto{\pgfqpoint{8.149017in}{0.640643in}}{\pgfqpoint{8.155249in}{0.640643in}}%
\pgfpathquadraticcurveto{\pgfqpoint{8.163877in}{0.640643in}}{\pgfqpoint{8.169082in}{0.646749in}}%
\pgfpathquadraticcurveto{\pgfqpoint{8.174288in}{0.652855in}}{\pgfqpoint{8.174288in}{0.662978in}}%
\pgfpathlineto{\pgfqpoint{8.174288in}{0.665283in}}%
\pgfpathlineto{\pgfqpoint{8.163985in}{0.665283in}}%
\pgfpathlineto{\pgfqpoint{8.163985in}{0.665283in}}%
\pgfpathclose%
\pgfpathmoveto{\pgfqpoint{8.184645in}{0.669570in}}%
\pgfpathlineto{\pgfqpoint{8.184645in}{0.633600in}}%
\pgfpathlineto{\pgfqpoint{8.174288in}{0.633600in}}%
\pgfpathlineto{\pgfqpoint{8.174288in}{0.643164in}}%
\pgfpathquadraticcurveto{\pgfqpoint{8.170739in}{0.637437in}}{\pgfqpoint{8.165444in}{0.634699in}}%
\pgfpathquadraticcurveto{\pgfqpoint{8.160148in}{0.631961in}}{\pgfqpoint{8.152493in}{0.631961in}}%
\pgfpathquadraticcurveto{\pgfqpoint{8.142821in}{0.631961in}}{\pgfqpoint{8.137093in}{0.637401in}}%
\pgfpathquadraticcurveto{\pgfqpoint{8.131383in}{0.642840in}}{\pgfqpoint{8.131383in}{0.651954in}}%
\pgfpathquadraticcurveto{\pgfqpoint{8.131383in}{0.662582in}}{\pgfqpoint{8.138498in}{0.667985in}}%
\pgfpathquadraticcurveto{\pgfqpoint{8.145630in}{0.673389in}}{\pgfqpoint{8.159752in}{0.673389in}}%
\pgfpathlineto{\pgfqpoint{8.174288in}{0.673389in}}%
\pgfpathlineto{\pgfqpoint{8.174288in}{0.674416in}}%
\pgfpathquadraticcurveto{\pgfqpoint{8.174288in}{0.681566in}}{\pgfqpoint{8.169587in}{0.685475in}}%
\pgfpathquadraticcurveto{\pgfqpoint{8.164885in}{0.689384in}}{\pgfqpoint{8.156384in}{0.689384in}}%
\pgfpathquadraticcurveto{\pgfqpoint{8.150980in}{0.689384in}}{\pgfqpoint{8.145847in}{0.688087in}}%
\pgfpathquadraticcurveto{\pgfqpoint{8.140731in}{0.686790in}}{\pgfqpoint{8.136012in}{0.684196in}}%
\pgfpathlineto{\pgfqpoint{8.136012in}{0.693779in}}%
\pgfpathquadraticcurveto{\pgfqpoint{8.141686in}{0.695976in}}{\pgfqpoint{8.147035in}{0.697057in}}%
\pgfpathquadraticcurveto{\pgfqpoint{8.152385in}{0.698156in}}{\pgfqpoint{8.157446in}{0.698156in}}%
\pgfpathquadraticcurveto{\pgfqpoint{8.171136in}{0.698156in}}{\pgfqpoint{8.177890in}{0.691059in}}%
\pgfpathquadraticcurveto{\pgfqpoint{8.184645in}{0.683980in}}{\pgfqpoint{8.184645in}{0.669570in}}%
\pgfpathlineto{\pgfqpoint{8.184645in}{0.669570in}}%
\pgfpathclose%
\pgfpathmoveto{\pgfqpoint{8.242500in}{0.686970in}}%
\pgfpathquadraticcurveto{\pgfqpoint{8.240753in}{0.687979in}}{\pgfqpoint{8.238699in}{0.688447in}}%
\pgfpathquadraticcurveto{\pgfqpoint{8.236646in}{0.688933in}}{\pgfqpoint{8.234178in}{0.688933in}}%
\pgfpathquadraticcurveto{\pgfqpoint{8.225388in}{0.688933in}}{\pgfqpoint{8.220687in}{0.683223in}}%
\pgfpathquadraticcurveto{\pgfqpoint{8.215986in}{0.677514in}}{\pgfqpoint{8.215986in}{0.666814in}}%
\pgfpathlineto{\pgfqpoint{8.215986in}{0.633600in}}%
\pgfpathlineto{\pgfqpoint{8.205575in}{0.633600in}}%
\pgfpathlineto{\pgfqpoint{8.205575in}{0.696643in}}%
\pgfpathlineto{\pgfqpoint{8.215986in}{0.696643in}}%
\pgfpathlineto{\pgfqpoint{8.215986in}{0.686844in}}%
\pgfpathquadraticcurveto{\pgfqpoint{8.219264in}{0.692590in}}{\pgfqpoint{8.224488in}{0.695364in}}%
\pgfpathquadraticcurveto{\pgfqpoint{8.229729in}{0.698156in}}{\pgfqpoint{8.237222in}{0.698156in}}%
\pgfpathquadraticcurveto{\pgfqpoint{8.238285in}{0.698156in}}{\pgfqpoint{8.239582in}{0.698011in}}%
\pgfpathquadraticcurveto{\pgfqpoint{8.240879in}{0.697885in}}{\pgfqpoint{8.242446in}{0.697597in}}%
\pgfpathlineto{\pgfqpoint{8.242500in}{0.686970in}}%
\pgfpathlineto{\pgfqpoint{8.242500in}{0.686970in}}%
\pgfpathclose%
\pgfpathmoveto{\pgfqpoint{8.253361in}{0.696643in}}%
\pgfpathlineto{\pgfqpoint{8.263718in}{0.696643in}}%
\pgfpathlineto{\pgfqpoint{8.263718in}{0.633600in}}%
\pgfpathlineto{\pgfqpoint{8.253361in}{0.633600in}}%
\pgfpathlineto{\pgfqpoint{8.253361in}{0.696643in}}%
\pgfpathlineto{\pgfqpoint{8.253361in}{0.696643in}}%
\pgfpathclose%
\pgfpathmoveto{\pgfqpoint{8.253361in}{0.721193in}}%
\pgfpathlineto{\pgfqpoint{8.263718in}{0.721193in}}%
\pgfpathlineto{\pgfqpoint{8.263718in}{0.708062in}}%
\pgfpathlineto{\pgfqpoint{8.253361in}{0.708062in}}%
\pgfpathlineto{\pgfqpoint{8.253361in}{0.721193in}}%
\pgfpathlineto{\pgfqpoint{8.253361in}{0.721193in}}%
\pgfpathclose%
\pgfpathmoveto{\pgfqpoint{8.314044in}{0.665283in}}%
\pgfpathquadraticcurveto{\pgfqpoint{8.301490in}{0.665283in}}{\pgfqpoint{8.296644in}{0.662419in}}%
\pgfpathquadraticcurveto{\pgfqpoint{8.291799in}{0.659556in}}{\pgfqpoint{8.291799in}{0.652621in}}%
\pgfpathquadraticcurveto{\pgfqpoint{8.291799in}{0.647109in}}{\pgfqpoint{8.295438in}{0.643867in}}%
\pgfpathquadraticcurveto{\pgfqpoint{8.299076in}{0.640643in}}{\pgfqpoint{8.305308in}{0.640643in}}%
\pgfpathquadraticcurveto{\pgfqpoint{8.313936in}{0.640643in}}{\pgfqpoint{8.319141in}{0.646749in}}%
\pgfpathquadraticcurveto{\pgfqpoint{8.324347in}{0.652855in}}{\pgfqpoint{8.324347in}{0.662978in}}%
\pgfpathlineto{\pgfqpoint{8.324347in}{0.665283in}}%
\pgfpathlineto{\pgfqpoint{8.314044in}{0.665283in}}%
\pgfpathlineto{\pgfqpoint{8.314044in}{0.665283in}}%
\pgfpathclose%
\pgfpathmoveto{\pgfqpoint{8.334704in}{0.669570in}}%
\pgfpathlineto{\pgfqpoint{8.334704in}{0.633600in}}%
\pgfpathlineto{\pgfqpoint{8.324347in}{0.633600in}}%
\pgfpathlineto{\pgfqpoint{8.324347in}{0.643164in}}%
\pgfpathquadraticcurveto{\pgfqpoint{8.320799in}{0.637437in}}{\pgfqpoint{8.315503in}{0.634699in}}%
\pgfpathquadraticcurveto{\pgfqpoint{8.310207in}{0.631961in}}{\pgfqpoint{8.302552in}{0.631961in}}%
\pgfpathquadraticcurveto{\pgfqpoint{8.292880in}{0.631961in}}{\pgfqpoint{8.287152in}{0.637401in}}%
\pgfpathquadraticcurveto{\pgfqpoint{8.281442in}{0.642840in}}{\pgfqpoint{8.281442in}{0.651954in}}%
\pgfpathquadraticcurveto{\pgfqpoint{8.281442in}{0.662582in}}{\pgfqpoint{8.288557in}{0.667985in}}%
\pgfpathquadraticcurveto{\pgfqpoint{8.295690in}{0.673389in}}{\pgfqpoint{8.309811in}{0.673389in}}%
\pgfpathlineto{\pgfqpoint{8.324347in}{0.673389in}}%
\pgfpathlineto{\pgfqpoint{8.324347in}{0.674416in}}%
\pgfpathquadraticcurveto{\pgfqpoint{8.324347in}{0.681566in}}{\pgfqpoint{8.319646in}{0.685475in}}%
\pgfpathquadraticcurveto{\pgfqpoint{8.314945in}{0.689384in}}{\pgfqpoint{8.306443in}{0.689384in}}%
\pgfpathquadraticcurveto{\pgfqpoint{8.301039in}{0.689384in}}{\pgfqpoint{8.295906in}{0.688087in}}%
\pgfpathquadraticcurveto{\pgfqpoint{8.290790in}{0.686790in}}{\pgfqpoint{8.286071in}{0.684196in}}%
\pgfpathlineto{\pgfqpoint{8.286071in}{0.693779in}}%
\pgfpathquadraticcurveto{\pgfqpoint{8.291745in}{0.695976in}}{\pgfqpoint{8.297095in}{0.697057in}}%
\pgfpathquadraticcurveto{\pgfqpoint{8.302444in}{0.698156in}}{\pgfqpoint{8.307506in}{0.698156in}}%
\pgfpathquadraticcurveto{\pgfqpoint{8.321195in}{0.698156in}}{\pgfqpoint{8.327949in}{0.691059in}}%
\pgfpathquadraticcurveto{\pgfqpoint{8.334704in}{0.683980in}}{\pgfqpoint{8.334704in}{0.669570in}}%
\pgfpathlineto{\pgfqpoint{8.334704in}{0.669570in}}%
\pgfpathclose%
\pgfpathmoveto{\pgfqpoint{8.408446in}{0.671660in}}%
\pgfpathlineto{\pgfqpoint{8.408446in}{0.633600in}}%
\pgfpathlineto{\pgfqpoint{8.398089in}{0.633600in}}%
\pgfpathlineto{\pgfqpoint{8.398089in}{0.671317in}}%
\pgfpathquadraticcurveto{\pgfqpoint{8.398089in}{0.680269in}}{\pgfqpoint{8.394594in}{0.684700in}}%
\pgfpathquadraticcurveto{\pgfqpoint{8.391100in}{0.689149in}}{\pgfqpoint{8.384129in}{0.689149in}}%
\pgfpathquadraticcurveto{\pgfqpoint{8.375736in}{0.689149in}}{\pgfqpoint{8.370890in}{0.683800in}}%
\pgfpathquadraticcurveto{\pgfqpoint{8.366045in}{0.678468in}}{\pgfqpoint{8.366045in}{0.669228in}}%
\pgfpathlineto{\pgfqpoint{8.366045in}{0.633600in}}%
\pgfpathlineto{\pgfqpoint{8.355634in}{0.633600in}}%
\pgfpathlineto{\pgfqpoint{8.355634in}{0.696643in}}%
\pgfpathlineto{\pgfqpoint{8.366045in}{0.696643in}}%
\pgfpathlineto{\pgfqpoint{8.366045in}{0.686844in}}%
\pgfpathquadraticcurveto{\pgfqpoint{8.369774in}{0.692536in}}{\pgfqpoint{8.374799in}{0.695346in}}%
\pgfpathquadraticcurveto{\pgfqpoint{8.379842in}{0.698156in}}{\pgfqpoint{8.386435in}{0.698156in}}%
\pgfpathquadraticcurveto{\pgfqpoint{8.397296in}{0.698156in}}{\pgfqpoint{8.402862in}{0.691437in}}%
\pgfpathquadraticcurveto{\pgfqpoint{8.408446in}{0.684718in}}{\pgfqpoint{8.408446in}{0.671660in}}%
\pgfpathlineto{\pgfqpoint{8.408446in}{0.671660in}}%
\pgfpathclose%
\pgfpathmoveto{\pgfqpoint{8.474460in}{0.694229in}}%
\pgfpathlineto{\pgfqpoint{8.474460in}{0.684538in}}%
\pgfpathquadraticcurveto{\pgfqpoint{8.470065in}{0.686970in}}{\pgfqpoint{8.465652in}{0.688177in}}%
\pgfpathquadraticcurveto{\pgfqpoint{8.461239in}{0.689384in}}{\pgfqpoint{8.456736in}{0.689384in}}%
\pgfpathquadraticcurveto{\pgfqpoint{8.446650in}{0.689384in}}{\pgfqpoint{8.441066in}{0.682989in}}%
\pgfpathquadraticcurveto{\pgfqpoint{8.435500in}{0.676613in}}{\pgfqpoint{8.435500in}{0.665067in}}%
\pgfpathquadraticcurveto{\pgfqpoint{8.435500in}{0.653521in}}{\pgfqpoint{8.441066in}{0.647127in}}%
\pgfpathquadraticcurveto{\pgfqpoint{8.446650in}{0.640751in}}{\pgfqpoint{8.456736in}{0.640751in}}%
\pgfpathquadraticcurveto{\pgfqpoint{8.461239in}{0.640751in}}{\pgfqpoint{8.465652in}{0.641958in}}%
\pgfpathquadraticcurveto{\pgfqpoint{8.470065in}{0.643164in}}{\pgfqpoint{8.474460in}{0.645596in}}%
\pgfpathlineto{\pgfqpoint{8.474460in}{0.636014in}}%
\pgfpathquadraticcurveto{\pgfqpoint{8.470119in}{0.633996in}}{\pgfqpoint{8.465472in}{0.632988in}}%
\pgfpathquadraticcurveto{\pgfqpoint{8.460843in}{0.631961in}}{\pgfqpoint{8.455602in}{0.631961in}}%
\pgfpathquadraticcurveto{\pgfqpoint{8.441354in}{0.631961in}}{\pgfqpoint{8.432960in}{0.640913in}}%
\pgfpathquadraticcurveto{\pgfqpoint{8.424585in}{0.649865in}}{\pgfqpoint{8.424585in}{0.665067in}}%
\pgfpathquadraticcurveto{\pgfqpoint{8.424585in}{0.680486in}}{\pgfqpoint{8.433050in}{0.689312in}}%
\pgfpathquadraticcurveto{\pgfqpoint{8.441534in}{0.698156in}}{\pgfqpoint{8.456286in}{0.698156in}}%
\pgfpathquadraticcurveto{\pgfqpoint{8.461059in}{0.698156in}}{\pgfqpoint{8.465616in}{0.697165in}}%
\pgfpathquadraticcurveto{\pgfqpoint{8.470173in}{0.696192in}}{\pgfqpoint{8.474460in}{0.694229in}}%
\pgfpathlineto{\pgfqpoint{8.474460in}{0.694229in}}%
\pgfpathclose%
\pgfpathmoveto{\pgfqpoint{8.546401in}{0.667715in}}%
\pgfpathlineto{\pgfqpoint{8.546401in}{0.662654in}}%
\pgfpathlineto{\pgfqpoint{8.498777in}{0.662654in}}%
\pgfpathquadraticcurveto{\pgfqpoint{8.499461in}{0.651954in}}{\pgfqpoint{8.505225in}{0.646353in}}%
\pgfpathquadraticcurveto{\pgfqpoint{8.510989in}{0.640751in}}{\pgfqpoint{8.521292in}{0.640751in}}%
\pgfpathquadraticcurveto{\pgfqpoint{8.527254in}{0.640751in}}{\pgfqpoint{8.532856in}{0.642210in}}%
\pgfpathquadraticcurveto{\pgfqpoint{8.538457in}{0.643669in}}{\pgfqpoint{8.543987in}{0.646605in}}%
\pgfpathlineto{\pgfqpoint{8.543987in}{0.636806in}}%
\pgfpathquadraticcurveto{\pgfqpoint{8.538403in}{0.634447in}}{\pgfqpoint{8.532549in}{0.633204in}}%
\pgfpathquadraticcurveto{\pgfqpoint{8.526696in}{0.631961in}}{\pgfqpoint{8.520679in}{0.631961in}}%
\pgfpathquadraticcurveto{\pgfqpoint{8.505585in}{0.631961in}}{\pgfqpoint{8.496777in}{0.640733in}}%
\pgfpathquadraticcurveto{\pgfqpoint{8.487969in}{0.649523in}}{\pgfqpoint{8.487969in}{0.664509in}}%
\pgfpathquadraticcurveto{\pgfqpoint{8.487969in}{0.679981in}}{\pgfqpoint{8.496327in}{0.689059in}}%
\pgfpathquadraticcurveto{\pgfqpoint{8.504685in}{0.698156in}}{\pgfqpoint{8.518878in}{0.698156in}}%
\pgfpathquadraticcurveto{\pgfqpoint{8.531595in}{0.698156in}}{\pgfqpoint{8.538998in}{0.689960in}}%
\pgfpathquadraticcurveto{\pgfqpoint{8.546401in}{0.681783in}}{\pgfqpoint{8.546401in}{0.667715in}}%
\pgfpathlineto{\pgfqpoint{8.546401in}{0.667715in}}%
\pgfpathclose%
\pgfpathmoveto{\pgfqpoint{8.536044in}{0.670759in}}%
\pgfpathquadraticcurveto{\pgfqpoint{8.535936in}{0.679243in}}{\pgfqpoint{8.531289in}{0.684304in}}%
\pgfpathquadraticcurveto{\pgfqpoint{8.526641in}{0.689384in}}{\pgfqpoint{8.518986in}{0.689384in}}%
\pgfpathquadraticcurveto{\pgfqpoint{8.510322in}{0.689384in}}{\pgfqpoint{8.505117in}{0.684484in}}%
\pgfpathquadraticcurveto{\pgfqpoint{8.499911in}{0.679585in}}{\pgfqpoint{8.499119in}{0.670687in}}%
\pgfpathlineto{\pgfqpoint{8.536044in}{0.670759in}}%
\pgfpathlineto{\pgfqpoint{8.536044in}{0.670759in}}%
\pgfpathclose%
\pgfusepath{fill}%
\end{pgfscope}%
\begin{pgfscope}%
\pgfsetbuttcap%
\pgfsetroundjoin%
\definecolor{currentfill}{rgb}{0.309804,0.372549,0.419608}%
\pgfsetfillcolor{currentfill}%
\pgfsetlinewidth{0.000000pt}%
\definecolor{currentstroke}{rgb}{0.309804,0.372549,0.419608}%
\pgfsetstrokecolor{currentstroke}%
\pgfsetdash{}{0pt}%
\pgfpathmoveto{\pgfqpoint{8.676156in}{0.665860in}}%
\pgfpathquadraticcurveto{\pgfqpoint{8.676156in}{0.677117in}}{\pgfqpoint{8.671509in}{0.683296in}}%
\pgfpathquadraticcurveto{\pgfqpoint{8.666880in}{0.689492in}}{\pgfqpoint{8.658486in}{0.689492in}}%
\pgfpathquadraticcurveto{\pgfqpoint{8.650164in}{0.689492in}}{\pgfqpoint{8.645517in}{0.683296in}}%
\pgfpathquadraticcurveto{\pgfqpoint{8.640870in}{0.677117in}}{\pgfqpoint{8.640870in}{0.665860in}}%
\pgfpathquadraticcurveto{\pgfqpoint{8.640870in}{0.654656in}}{\pgfqpoint{8.645517in}{0.648460in}}%
\pgfpathquadraticcurveto{\pgfqpoint{8.650164in}{0.642264in}}{\pgfqpoint{8.658486in}{0.642264in}}%
\pgfpathquadraticcurveto{\pgfqpoint{8.666880in}{0.642264in}}{\pgfqpoint{8.671509in}{0.648460in}}%
\pgfpathquadraticcurveto{\pgfqpoint{8.676156in}{0.654656in}}{\pgfqpoint{8.676156in}{0.665860in}}%
\pgfpathlineto{\pgfqpoint{8.676156in}{0.665860in}}%
\pgfpathclose%
\pgfpathmoveto{\pgfqpoint{8.686513in}{0.641417in}}%
\pgfpathquadraticcurveto{\pgfqpoint{8.686513in}{0.625332in}}{\pgfqpoint{8.679362in}{0.617479in}}%
\pgfpathquadraticcurveto{\pgfqpoint{8.672229in}{0.609626in}}{\pgfqpoint{8.657477in}{0.609626in}}%
\pgfpathquadraticcurveto{\pgfqpoint{8.652020in}{0.609626in}}{\pgfqpoint{8.647174in}{0.610436in}}%
\pgfpathquadraticcurveto{\pgfqpoint{8.642329in}{0.611247in}}{\pgfqpoint{8.637772in}{0.612940in}}%
\pgfpathlineto{\pgfqpoint{8.637772in}{0.623009in}}%
\pgfpathquadraticcurveto{\pgfqpoint{8.642329in}{0.620541in}}{\pgfqpoint{8.646778in}{0.619370in}}%
\pgfpathquadraticcurveto{\pgfqpoint{8.651227in}{0.618182in}}{\pgfqpoint{8.655838in}{0.618182in}}%
\pgfpathquadraticcurveto{\pgfqpoint{8.666033in}{0.618182in}}{\pgfqpoint{8.671094in}{0.623495in}}%
\pgfpathquadraticcurveto{\pgfqpoint{8.676156in}{0.628809in}}{\pgfqpoint{8.676156in}{0.639562in}}%
\pgfpathlineto{\pgfqpoint{8.676156in}{0.644695in}}%
\pgfpathquadraticcurveto{\pgfqpoint{8.672950in}{0.639112in}}{\pgfqpoint{8.667942in}{0.636356in}}%
\pgfpathquadraticcurveto{\pgfqpoint{8.662935in}{0.633600in}}{\pgfqpoint{8.655946in}{0.633600in}}%
\pgfpathquadraticcurveto{\pgfqpoint{8.644364in}{0.633600in}}{\pgfqpoint{8.637268in}{0.642426in}}%
\pgfpathquadraticcurveto{\pgfqpoint{8.630171in}{0.651270in}}{\pgfqpoint{8.630171in}{0.665860in}}%
\pgfpathquadraticcurveto{\pgfqpoint{8.630171in}{0.680486in}}{\pgfqpoint{8.637268in}{0.689312in}}%
\pgfpathquadraticcurveto{\pgfqpoint{8.644364in}{0.698156in}}{\pgfqpoint{8.655946in}{0.698156in}}%
\pgfpathquadraticcurveto{\pgfqpoint{8.662935in}{0.698156in}}{\pgfqpoint{8.667942in}{0.695400in}}%
\pgfpathquadraticcurveto{\pgfqpoint{8.672950in}{0.692644in}}{\pgfqpoint{8.676156in}{0.687078in}}%
\pgfpathlineto{\pgfqpoint{8.676156in}{0.696643in}}%
\pgfpathlineto{\pgfqpoint{8.686513in}{0.696643in}}%
\pgfpathlineto{\pgfqpoint{8.686513in}{0.641417in}}%
\pgfpathlineto{\pgfqpoint{8.686513in}{0.641417in}}%
\pgfpathclose%
\pgfpathmoveto{\pgfqpoint{8.706794in}{0.658475in}}%
\pgfpathlineto{\pgfqpoint{8.706794in}{0.696643in}}%
\pgfpathlineto{\pgfqpoint{8.717151in}{0.696643in}}%
\pgfpathlineto{\pgfqpoint{8.717151in}{0.658871in}}%
\pgfpathquadraticcurveto{\pgfqpoint{8.717151in}{0.649919in}}{\pgfqpoint{8.720628in}{0.645434in}}%
\pgfpathquadraticcurveto{\pgfqpoint{8.724122in}{0.640967in}}{\pgfqpoint{8.731111in}{0.640967in}}%
\pgfpathquadraticcurveto{\pgfqpoint{8.739487in}{0.640967in}}{\pgfqpoint{8.744350in}{0.646317in}}%
\pgfpathquadraticcurveto{\pgfqpoint{8.749231in}{0.651666in}}{\pgfqpoint{8.749231in}{0.660906in}}%
\pgfpathlineto{\pgfqpoint{8.749231in}{0.696643in}}%
\pgfpathlineto{\pgfqpoint{8.759588in}{0.696643in}}%
\pgfpathlineto{\pgfqpoint{8.759588in}{0.633600in}}%
\pgfpathlineto{\pgfqpoint{8.749231in}{0.633600in}}%
\pgfpathlineto{\pgfqpoint{8.749231in}{0.643291in}}%
\pgfpathquadraticcurveto{\pgfqpoint{8.745467in}{0.637545in}}{\pgfqpoint{8.740477in}{0.634753in}}%
\pgfpathquadraticcurveto{\pgfqpoint{8.735506in}{0.631961in}}{\pgfqpoint{8.728913in}{0.631961in}}%
\pgfpathquadraticcurveto{\pgfqpoint{8.718052in}{0.631961in}}{\pgfqpoint{8.712414in}{0.638715in}}%
\pgfpathquadraticcurveto{\pgfqpoint{8.706794in}{0.645470in}}{\pgfqpoint{8.706794in}{0.658475in}}%
\pgfpathlineto{\pgfqpoint{8.706794in}{0.658475in}}%
\pgfpathclose%
\pgfpathmoveto{\pgfqpoint{8.732858in}{0.698156in}}%
\pgfpathlineto{\pgfqpoint{8.732858in}{0.698156in}}%
\pgfpathlineto{\pgfqpoint{8.732858in}{0.698156in}}%
\pgfpathclose%
\pgfpathmoveto{\pgfqpoint{8.809572in}{0.665283in}}%
\pgfpathquadraticcurveto{\pgfqpoint{8.797017in}{0.665283in}}{\pgfqpoint{8.792172in}{0.662419in}}%
\pgfpathquadraticcurveto{\pgfqpoint{8.787327in}{0.659556in}}{\pgfqpoint{8.787327in}{0.652621in}}%
\pgfpathquadraticcurveto{\pgfqpoint{8.787327in}{0.647109in}}{\pgfqpoint{8.790965in}{0.643867in}}%
\pgfpathquadraticcurveto{\pgfqpoint{8.794604in}{0.640643in}}{\pgfqpoint{8.800836in}{0.640643in}}%
\pgfpathquadraticcurveto{\pgfqpoint{8.809464in}{0.640643in}}{\pgfqpoint{8.814669in}{0.646749in}}%
\pgfpathquadraticcurveto{\pgfqpoint{8.819875in}{0.652855in}}{\pgfqpoint{8.819875in}{0.662978in}}%
\pgfpathlineto{\pgfqpoint{8.819875in}{0.665283in}}%
\pgfpathlineto{\pgfqpoint{8.809572in}{0.665283in}}%
\pgfpathlineto{\pgfqpoint{8.809572in}{0.665283in}}%
\pgfpathclose%
\pgfpathmoveto{\pgfqpoint{8.830232in}{0.669570in}}%
\pgfpathlineto{\pgfqpoint{8.830232in}{0.633600in}}%
\pgfpathlineto{\pgfqpoint{8.819875in}{0.633600in}}%
\pgfpathlineto{\pgfqpoint{8.819875in}{0.643164in}}%
\pgfpathquadraticcurveto{\pgfqpoint{8.816326in}{0.637437in}}{\pgfqpoint{8.811031in}{0.634699in}}%
\pgfpathquadraticcurveto{\pgfqpoint{8.805735in}{0.631961in}}{\pgfqpoint{8.798080in}{0.631961in}}%
\pgfpathquadraticcurveto{\pgfqpoint{8.788408in}{0.631961in}}{\pgfqpoint{8.782680in}{0.637401in}}%
\pgfpathquadraticcurveto{\pgfqpoint{8.776970in}{0.642840in}}{\pgfqpoint{8.776970in}{0.651954in}}%
\pgfpathquadraticcurveto{\pgfqpoint{8.776970in}{0.662582in}}{\pgfqpoint{8.784085in}{0.667985in}}%
\pgfpathquadraticcurveto{\pgfqpoint{8.791217in}{0.673389in}}{\pgfqpoint{8.805339in}{0.673389in}}%
\pgfpathlineto{\pgfqpoint{8.819875in}{0.673389in}}%
\pgfpathlineto{\pgfqpoint{8.819875in}{0.674416in}}%
\pgfpathquadraticcurveto{\pgfqpoint{8.819875in}{0.681566in}}{\pgfqpoint{8.815174in}{0.685475in}}%
\pgfpathquadraticcurveto{\pgfqpoint{8.810472in}{0.689384in}}{\pgfqpoint{8.801971in}{0.689384in}}%
\pgfpathquadraticcurveto{\pgfqpoint{8.796567in}{0.689384in}}{\pgfqpoint{8.791434in}{0.688087in}}%
\pgfpathquadraticcurveto{\pgfqpoint{8.786318in}{0.686790in}}{\pgfqpoint{8.781599in}{0.684196in}}%
\pgfpathlineto{\pgfqpoint{8.781599in}{0.693779in}}%
\pgfpathquadraticcurveto{\pgfqpoint{8.787273in}{0.695976in}}{\pgfqpoint{8.792622in}{0.697057in}}%
\pgfpathquadraticcurveto{\pgfqpoint{8.797972in}{0.698156in}}{\pgfqpoint{8.803033in}{0.698156in}}%
\pgfpathquadraticcurveto{\pgfqpoint{8.816723in}{0.698156in}}{\pgfqpoint{8.823477in}{0.691059in}}%
\pgfpathquadraticcurveto{\pgfqpoint{8.830232in}{0.683980in}}{\pgfqpoint{8.830232in}{0.669570in}}%
\pgfpathlineto{\pgfqpoint{8.830232in}{0.669570in}}%
\pgfpathclose%
\pgfpathmoveto{\pgfqpoint{8.888087in}{0.686970in}}%
\pgfpathquadraticcurveto{\pgfqpoint{8.886340in}{0.687979in}}{\pgfqpoint{8.884286in}{0.688447in}}%
\pgfpathquadraticcurveto{\pgfqpoint{8.882233in}{0.688933in}}{\pgfqpoint{8.879765in}{0.688933in}}%
\pgfpathquadraticcurveto{\pgfqpoint{8.870975in}{0.688933in}}{\pgfqpoint{8.866274in}{0.683223in}}%
\pgfpathquadraticcurveto{\pgfqpoint{8.861573in}{0.677514in}}{\pgfqpoint{8.861573in}{0.666814in}}%
\pgfpathlineto{\pgfqpoint{8.861573in}{0.633600in}}%
\pgfpathlineto{\pgfqpoint{8.851162in}{0.633600in}}%
\pgfpathlineto{\pgfqpoint{8.851162in}{0.696643in}}%
\pgfpathlineto{\pgfqpoint{8.861573in}{0.696643in}}%
\pgfpathlineto{\pgfqpoint{8.861573in}{0.686844in}}%
\pgfpathquadraticcurveto{\pgfqpoint{8.864851in}{0.692590in}}{\pgfqpoint{8.870075in}{0.695364in}}%
\pgfpathquadraticcurveto{\pgfqpoint{8.875316in}{0.698156in}}{\pgfqpoint{8.882809in}{0.698156in}}%
\pgfpathquadraticcurveto{\pgfqpoint{8.883872in}{0.698156in}}{\pgfqpoint{8.885169in}{0.698011in}}%
\pgfpathquadraticcurveto{\pgfqpoint{8.886466in}{0.697885in}}{\pgfqpoint{8.888033in}{0.697597in}}%
\pgfpathlineto{\pgfqpoint{8.888087in}{0.686970in}}%
\pgfpathlineto{\pgfqpoint{8.888087in}{0.686970in}}%
\pgfpathclose%
\pgfpathmoveto{\pgfqpoint{8.927605in}{0.665283in}}%
\pgfpathquadraticcurveto{\pgfqpoint{8.915051in}{0.665283in}}{\pgfqpoint{8.910206in}{0.662419in}}%
\pgfpathquadraticcurveto{\pgfqpoint{8.905360in}{0.659556in}}{\pgfqpoint{8.905360in}{0.652621in}}%
\pgfpathquadraticcurveto{\pgfqpoint{8.905360in}{0.647109in}}{\pgfqpoint{8.908999in}{0.643867in}}%
\pgfpathquadraticcurveto{\pgfqpoint{8.912637in}{0.640643in}}{\pgfqpoint{8.918870in}{0.640643in}}%
\pgfpathquadraticcurveto{\pgfqpoint{8.927497in}{0.640643in}}{\pgfqpoint{8.932703in}{0.646749in}}%
\pgfpathquadraticcurveto{\pgfqpoint{8.937908in}{0.652855in}}{\pgfqpoint{8.937908in}{0.662978in}}%
\pgfpathlineto{\pgfqpoint{8.937908in}{0.665283in}}%
\pgfpathlineto{\pgfqpoint{8.927605in}{0.665283in}}%
\pgfpathlineto{\pgfqpoint{8.927605in}{0.665283in}}%
\pgfpathclose%
\pgfpathmoveto{\pgfqpoint{8.948265in}{0.669570in}}%
\pgfpathlineto{\pgfqpoint{8.948265in}{0.633600in}}%
\pgfpathlineto{\pgfqpoint{8.937908in}{0.633600in}}%
\pgfpathlineto{\pgfqpoint{8.937908in}{0.643164in}}%
\pgfpathquadraticcurveto{\pgfqpoint{8.934360in}{0.637437in}}{\pgfqpoint{8.929064in}{0.634699in}}%
\pgfpathquadraticcurveto{\pgfqpoint{8.923769in}{0.631961in}}{\pgfqpoint{8.916114in}{0.631961in}}%
\pgfpathquadraticcurveto{\pgfqpoint{8.906441in}{0.631961in}}{\pgfqpoint{8.900713in}{0.637401in}}%
\pgfpathquadraticcurveto{\pgfqpoint{8.895003in}{0.642840in}}{\pgfqpoint{8.895003in}{0.651954in}}%
\pgfpathquadraticcurveto{\pgfqpoint{8.895003in}{0.662582in}}{\pgfqpoint{8.902118in}{0.667985in}}%
\pgfpathquadraticcurveto{\pgfqpoint{8.909251in}{0.673389in}}{\pgfqpoint{8.923373in}{0.673389in}}%
\pgfpathlineto{\pgfqpoint{8.937908in}{0.673389in}}%
\pgfpathlineto{\pgfqpoint{8.937908in}{0.674416in}}%
\pgfpathquadraticcurveto{\pgfqpoint{8.937908in}{0.681566in}}{\pgfqpoint{8.933207in}{0.685475in}}%
\pgfpathquadraticcurveto{\pgfqpoint{8.928506in}{0.689384in}}{\pgfqpoint{8.920004in}{0.689384in}}%
\pgfpathquadraticcurveto{\pgfqpoint{8.914601in}{0.689384in}}{\pgfqpoint{8.909467in}{0.688087in}}%
\pgfpathquadraticcurveto{\pgfqpoint{8.904352in}{0.686790in}}{\pgfqpoint{8.899633in}{0.684196in}}%
\pgfpathlineto{\pgfqpoint{8.899633in}{0.693779in}}%
\pgfpathquadraticcurveto{\pgfqpoint{8.905306in}{0.695976in}}{\pgfqpoint{8.910656in}{0.697057in}}%
\pgfpathquadraticcurveto{\pgfqpoint{8.916006in}{0.698156in}}{\pgfqpoint{8.921067in}{0.698156in}}%
\pgfpathquadraticcurveto{\pgfqpoint{8.934756in}{0.698156in}}{\pgfqpoint{8.941511in}{0.691059in}}%
\pgfpathquadraticcurveto{\pgfqpoint{8.948265in}{0.683980in}}{\pgfqpoint{8.948265in}{0.669570in}}%
\pgfpathlineto{\pgfqpoint{8.948265in}{0.669570in}}%
\pgfpathclose%
\pgfpathmoveto{\pgfqpoint{9.022007in}{0.671660in}}%
\pgfpathlineto{\pgfqpoint{9.022007in}{0.633600in}}%
\pgfpathlineto{\pgfqpoint{9.011650in}{0.633600in}}%
\pgfpathlineto{\pgfqpoint{9.011650in}{0.671317in}}%
\pgfpathquadraticcurveto{\pgfqpoint{9.011650in}{0.680269in}}{\pgfqpoint{9.008156in}{0.684700in}}%
\pgfpathquadraticcurveto{\pgfqpoint{9.004661in}{0.689149in}}{\pgfqpoint{8.997691in}{0.689149in}}%
\pgfpathquadraticcurveto{\pgfqpoint{8.989297in}{0.689149in}}{\pgfqpoint{8.984452in}{0.683800in}}%
\pgfpathquadraticcurveto{\pgfqpoint{8.979607in}{0.678468in}}{\pgfqpoint{8.979607in}{0.669228in}}%
\pgfpathlineto{\pgfqpoint{8.979607in}{0.633600in}}%
\pgfpathlineto{\pgfqpoint{8.969196in}{0.633600in}}%
\pgfpathlineto{\pgfqpoint{8.969196in}{0.696643in}}%
\pgfpathlineto{\pgfqpoint{8.979607in}{0.696643in}}%
\pgfpathlineto{\pgfqpoint{8.979607in}{0.686844in}}%
\pgfpathquadraticcurveto{\pgfqpoint{8.983335in}{0.692536in}}{\pgfqpoint{8.988360in}{0.695346in}}%
\pgfpathquadraticcurveto{\pgfqpoint{8.993404in}{0.698156in}}{\pgfqpoint{8.999996in}{0.698156in}}%
\pgfpathquadraticcurveto{\pgfqpoint{9.010858in}{0.698156in}}{\pgfqpoint{9.016423in}{0.691437in}}%
\pgfpathquadraticcurveto{\pgfqpoint{9.022007in}{0.684718in}}{\pgfqpoint{9.022007in}{0.671660in}}%
\pgfpathlineto{\pgfqpoint{9.022007in}{0.671660in}}%
\pgfpathclose%
\pgfpathmoveto{\pgfqpoint{9.052898in}{0.714547in}}%
\pgfpathlineto{\pgfqpoint{9.052898in}{0.696643in}}%
\pgfpathlineto{\pgfqpoint{9.074224in}{0.696643in}}%
\pgfpathlineto{\pgfqpoint{9.074224in}{0.688591in}}%
\pgfpathlineto{\pgfqpoint{9.052898in}{0.688591in}}%
\pgfpathlineto{\pgfqpoint{9.052898in}{0.654368in}}%
\pgfpathquadraticcurveto{\pgfqpoint{9.052898in}{0.646659in}}{\pgfqpoint{9.055005in}{0.644461in}}%
\pgfpathquadraticcurveto{\pgfqpoint{9.057113in}{0.642264in}}{\pgfqpoint{9.063597in}{0.642264in}}%
\pgfpathlineto{\pgfqpoint{9.074224in}{0.642264in}}%
\pgfpathlineto{\pgfqpoint{9.074224in}{0.633600in}}%
\pgfpathlineto{\pgfqpoint{9.063597in}{0.633600in}}%
\pgfpathquadraticcurveto{\pgfqpoint{9.051601in}{0.633600in}}{\pgfqpoint{9.047044in}{0.638067in}}%
\pgfpathquadraticcurveto{\pgfqpoint{9.042487in}{0.642552in}}{\pgfqpoint{9.042487in}{0.654368in}}%
\pgfpathlineto{\pgfqpoint{9.042487in}{0.688591in}}%
\pgfpathlineto{\pgfqpoint{9.034886in}{0.688591in}}%
\pgfpathlineto{\pgfqpoint{9.034886in}{0.696643in}}%
\pgfpathlineto{\pgfqpoint{9.042487in}{0.696643in}}%
\pgfpathlineto{\pgfqpoint{9.042487in}{0.714547in}}%
\pgfpathlineto{\pgfqpoint{9.052898in}{0.714547in}}%
\pgfpathlineto{\pgfqpoint{9.052898in}{0.714547in}}%
\pgfpathclose%
\pgfpathmoveto{\pgfqpoint{9.141770in}{0.667715in}}%
\pgfpathlineto{\pgfqpoint{9.141770in}{0.662654in}}%
\pgfpathlineto{\pgfqpoint{9.094146in}{0.662654in}}%
\pgfpathquadraticcurveto{\pgfqpoint{9.094830in}{0.651954in}}{\pgfqpoint{9.100594in}{0.646353in}}%
\pgfpathquadraticcurveto{\pgfqpoint{9.106358in}{0.640751in}}{\pgfqpoint{9.116661in}{0.640751in}}%
\pgfpathquadraticcurveto{\pgfqpoint{9.122623in}{0.640751in}}{\pgfqpoint{9.128225in}{0.642210in}}%
\pgfpathquadraticcurveto{\pgfqpoint{9.133827in}{0.643669in}}{\pgfqpoint{9.139356in}{0.646605in}}%
\pgfpathlineto{\pgfqpoint{9.139356in}{0.636806in}}%
\pgfpathquadraticcurveto{\pgfqpoint{9.133773in}{0.634447in}}{\pgfqpoint{9.127919in}{0.633204in}}%
\pgfpathquadraticcurveto{\pgfqpoint{9.122065in}{0.631961in}}{\pgfqpoint{9.116049in}{0.631961in}}%
\pgfpathquadraticcurveto{\pgfqpoint{9.100954in}{0.631961in}}{\pgfqpoint{9.092146in}{0.640733in}}%
\pgfpathquadraticcurveto{\pgfqpoint{9.083339in}{0.649523in}}{\pgfqpoint{9.083339in}{0.664509in}}%
\pgfpathquadraticcurveto{\pgfqpoint{9.083339in}{0.679981in}}{\pgfqpoint{9.091696in}{0.689059in}}%
\pgfpathquadraticcurveto{\pgfqpoint{9.100054in}{0.698156in}}{\pgfqpoint{9.114247in}{0.698156in}}%
\pgfpathquadraticcurveto{\pgfqpoint{9.126964in}{0.698156in}}{\pgfqpoint{9.134367in}{0.689960in}}%
\pgfpathquadraticcurveto{\pgfqpoint{9.141770in}{0.681783in}}{\pgfqpoint{9.141770in}{0.667715in}}%
\pgfpathlineto{\pgfqpoint{9.141770in}{0.667715in}}%
\pgfpathclose%
\pgfpathmoveto{\pgfqpoint{9.131413in}{0.670759in}}%
\pgfpathquadraticcurveto{\pgfqpoint{9.131305in}{0.679243in}}{\pgfqpoint{9.126658in}{0.684304in}}%
\pgfpathquadraticcurveto{\pgfqpoint{9.122011in}{0.689384in}}{\pgfqpoint{9.114355in}{0.689384in}}%
\pgfpathquadraticcurveto{\pgfqpoint{9.105692in}{0.689384in}}{\pgfqpoint{9.100486in}{0.684484in}}%
\pgfpathquadraticcurveto{\pgfqpoint{9.095281in}{0.679585in}}{\pgfqpoint{9.094488in}{0.670687in}}%
\pgfpathlineto{\pgfqpoint{9.131413in}{0.670759in}}%
\pgfpathlineto{\pgfqpoint{9.131413in}{0.670759in}}%
\pgfpathclose%
\pgfpathmoveto{\pgfqpoint{9.212702in}{0.667715in}}%
\pgfpathlineto{\pgfqpoint{9.212702in}{0.662654in}}%
\pgfpathlineto{\pgfqpoint{9.165078in}{0.662654in}}%
\pgfpathquadraticcurveto{\pgfqpoint{9.165762in}{0.651954in}}{\pgfqpoint{9.171526in}{0.646353in}}%
\pgfpathquadraticcurveto{\pgfqpoint{9.177290in}{0.640751in}}{\pgfqpoint{9.187593in}{0.640751in}}%
\pgfpathquadraticcurveto{\pgfqpoint{9.193555in}{0.640751in}}{\pgfqpoint{9.199157in}{0.642210in}}%
\pgfpathquadraticcurveto{\pgfqpoint{9.204758in}{0.643669in}}{\pgfqpoint{9.210288in}{0.646605in}}%
\pgfpathlineto{\pgfqpoint{9.210288in}{0.636806in}}%
\pgfpathquadraticcurveto{\pgfqpoint{9.204704in}{0.634447in}}{\pgfqpoint{9.198850in}{0.633204in}}%
\pgfpathquadraticcurveto{\pgfqpoint{9.192997in}{0.631961in}}{\pgfqpoint{9.186980in}{0.631961in}}%
\pgfpathquadraticcurveto{\pgfqpoint{9.171886in}{0.631961in}}{\pgfqpoint{9.163078in}{0.640733in}}%
\pgfpathquadraticcurveto{\pgfqpoint{9.154270in}{0.649523in}}{\pgfqpoint{9.154270in}{0.664509in}}%
\pgfpathquadraticcurveto{\pgfqpoint{9.154270in}{0.679981in}}{\pgfqpoint{9.162628in}{0.689059in}}%
\pgfpathquadraticcurveto{\pgfqpoint{9.170986in}{0.698156in}}{\pgfqpoint{9.185179in}{0.698156in}}%
\pgfpathquadraticcurveto{\pgfqpoint{9.197896in}{0.698156in}}{\pgfqpoint{9.205299in}{0.689960in}}%
\pgfpathquadraticcurveto{\pgfqpoint{9.212702in}{0.681783in}}{\pgfqpoint{9.212702in}{0.667715in}}%
\pgfpathlineto{\pgfqpoint{9.212702in}{0.667715in}}%
\pgfpathclose%
\pgfpathmoveto{\pgfqpoint{9.202345in}{0.670759in}}%
\pgfpathquadraticcurveto{\pgfqpoint{9.202237in}{0.679243in}}{\pgfqpoint{9.197590in}{0.684304in}}%
\pgfpathquadraticcurveto{\pgfqpoint{9.192942in}{0.689384in}}{\pgfqpoint{9.185287in}{0.689384in}}%
\pgfpathquadraticcurveto{\pgfqpoint{9.176623in}{0.689384in}}{\pgfqpoint{9.171418in}{0.684484in}}%
\pgfpathquadraticcurveto{\pgfqpoint{9.166212in}{0.679585in}}{\pgfqpoint{9.165420in}{0.670687in}}%
\pgfpathlineto{\pgfqpoint{9.202345in}{0.670759in}}%
\pgfpathlineto{\pgfqpoint{9.202345in}{0.670759in}}%
\pgfpathclose%
\pgfpathmoveto{\pgfqpoint{9.231164in}{0.647902in}}%
\pgfpathlineto{\pgfqpoint{9.243052in}{0.647902in}}%
\pgfpathlineto{\pgfqpoint{9.243052in}{0.633600in}}%
\pgfpathlineto{\pgfqpoint{9.231164in}{0.633600in}}%
\pgfpathlineto{\pgfqpoint{9.231164in}{0.647902in}}%
\pgfpathlineto{\pgfqpoint{9.231164in}{0.647902in}}%
\pgfpathclose%
\pgfusepath{fill}%
\end{pgfscope}%
\begin{pgfscope}%
\pgfsetbuttcap%
\pgfsetroundjoin%
\definecolor{currentfill}{rgb}{0.309804,0.372549,0.419608}%
\pgfsetfillcolor{currentfill}%
\pgfsetlinewidth{0.000000pt}%
\definecolor{currentstroke}{rgb}{0.309804,0.372549,0.419608}%
\pgfsetstrokecolor{currentstroke}%
\pgfsetdash{}{0pt}%
\pgfpathmoveto{\pgfqpoint{9.352701in}{0.717645in}}%
\pgfpathlineto{\pgfqpoint{9.405836in}{0.717645in}}%
\pgfpathlineto{\pgfqpoint{9.405836in}{0.708062in}}%
\pgfpathlineto{\pgfqpoint{9.364066in}{0.708062in}}%
\pgfpathlineto{\pgfqpoint{9.364066in}{0.683187in}}%
\pgfpathlineto{\pgfqpoint{9.404089in}{0.683187in}}%
\pgfpathlineto{\pgfqpoint{9.404089in}{0.673623in}}%
\pgfpathlineto{\pgfqpoint{9.364066in}{0.673623in}}%
\pgfpathlineto{\pgfqpoint{9.364066in}{0.643164in}}%
\pgfpathlineto{\pgfqpoint{9.406845in}{0.643164in}}%
\pgfpathlineto{\pgfqpoint{9.406845in}{0.633600in}}%
\pgfpathlineto{\pgfqpoint{9.352701in}{0.633600in}}%
\pgfpathlineto{\pgfqpoint{9.352701in}{0.717645in}}%
\pgfpathlineto{\pgfqpoint{9.352701in}{0.717645in}}%
\pgfpathclose%
\pgfpathmoveto{\pgfqpoint{9.477507in}{0.696643in}}%
\pgfpathlineto{\pgfqpoint{9.454703in}{0.665968in}}%
\pgfpathlineto{\pgfqpoint{9.478678in}{0.633600in}}%
\pgfpathlineto{\pgfqpoint{9.466465in}{0.633600in}}%
\pgfpathlineto{\pgfqpoint{9.448111in}{0.658367in}}%
\pgfpathlineto{\pgfqpoint{9.429775in}{0.633600in}}%
\pgfpathlineto{\pgfqpoint{9.417544in}{0.633600in}}%
\pgfpathlineto{\pgfqpoint{9.442041in}{0.666580in}}%
\pgfpathlineto{\pgfqpoint{9.419634in}{0.696643in}}%
\pgfpathlineto{\pgfqpoint{9.431846in}{0.696643in}}%
\pgfpathlineto{\pgfqpoint{9.448561in}{0.674181in}}%
\pgfpathlineto{\pgfqpoint{9.465276in}{0.696643in}}%
\pgfpathlineto{\pgfqpoint{9.477507in}{0.696643in}}%
\pgfpathlineto{\pgfqpoint{9.477507in}{0.696643in}}%
\pgfpathclose%
\pgfpathmoveto{\pgfqpoint{9.521979in}{0.665283in}}%
\pgfpathquadraticcurveto{\pgfqpoint{9.509424in}{0.665283in}}{\pgfqpoint{9.504579in}{0.662419in}}%
\pgfpathquadraticcurveto{\pgfqpoint{9.499734in}{0.659556in}}{\pgfqpoint{9.499734in}{0.652621in}}%
\pgfpathquadraticcurveto{\pgfqpoint{9.499734in}{0.647109in}}{\pgfqpoint{9.503372in}{0.643867in}}%
\pgfpathquadraticcurveto{\pgfqpoint{9.507011in}{0.640643in}}{\pgfqpoint{9.513243in}{0.640643in}}%
\pgfpathquadraticcurveto{\pgfqpoint{9.521871in}{0.640643in}}{\pgfqpoint{9.527076in}{0.646749in}}%
\pgfpathquadraticcurveto{\pgfqpoint{9.532282in}{0.652855in}}{\pgfqpoint{9.532282in}{0.662978in}}%
\pgfpathlineto{\pgfqpoint{9.532282in}{0.665283in}}%
\pgfpathlineto{\pgfqpoint{9.521979in}{0.665283in}}%
\pgfpathlineto{\pgfqpoint{9.521979in}{0.665283in}}%
\pgfpathclose%
\pgfpathmoveto{\pgfqpoint{9.542639in}{0.669570in}}%
\pgfpathlineto{\pgfqpoint{9.542639in}{0.633600in}}%
\pgfpathlineto{\pgfqpoint{9.532282in}{0.633600in}}%
\pgfpathlineto{\pgfqpoint{9.532282in}{0.643164in}}%
\pgfpathquadraticcurveto{\pgfqpoint{9.528733in}{0.637437in}}{\pgfqpoint{9.523438in}{0.634699in}}%
\pgfpathquadraticcurveto{\pgfqpoint{9.518142in}{0.631961in}}{\pgfqpoint{9.510487in}{0.631961in}}%
\pgfpathquadraticcurveto{\pgfqpoint{9.500814in}{0.631961in}}{\pgfqpoint{9.495087in}{0.637401in}}%
\pgfpathquadraticcurveto{\pgfqpoint{9.489377in}{0.642840in}}{\pgfqpoint{9.489377in}{0.651954in}}%
\pgfpathquadraticcurveto{\pgfqpoint{9.489377in}{0.662582in}}{\pgfqpoint{9.496492in}{0.667985in}}%
\pgfpathquadraticcurveto{\pgfqpoint{9.503624in}{0.673389in}}{\pgfqpoint{9.517746in}{0.673389in}}%
\pgfpathlineto{\pgfqpoint{9.532282in}{0.673389in}}%
\pgfpathlineto{\pgfqpoint{9.532282in}{0.674416in}}%
\pgfpathquadraticcurveto{\pgfqpoint{9.532282in}{0.681566in}}{\pgfqpoint{9.527581in}{0.685475in}}%
\pgfpathquadraticcurveto{\pgfqpoint{9.522879in}{0.689384in}}{\pgfqpoint{9.514378in}{0.689384in}}%
\pgfpathquadraticcurveto{\pgfqpoint{9.508974in}{0.689384in}}{\pgfqpoint{9.503840in}{0.688087in}}%
\pgfpathquadraticcurveto{\pgfqpoint{9.498725in}{0.686790in}}{\pgfqpoint{9.494006in}{0.684196in}}%
\pgfpathlineto{\pgfqpoint{9.494006in}{0.693779in}}%
\pgfpathquadraticcurveto{\pgfqpoint{9.499680in}{0.695976in}}{\pgfqpoint{9.505029in}{0.697057in}}%
\pgfpathquadraticcurveto{\pgfqpoint{9.510379in}{0.698156in}}{\pgfqpoint{9.515440in}{0.698156in}}%
\pgfpathquadraticcurveto{\pgfqpoint{9.529130in}{0.698156in}}{\pgfqpoint{9.535884in}{0.691059in}}%
\pgfpathquadraticcurveto{\pgfqpoint{9.542639in}{0.683980in}}{\pgfqpoint{9.542639in}{0.669570in}}%
\pgfpathlineto{\pgfqpoint{9.542639in}{0.669570in}}%
\pgfpathclose%
\pgfpathmoveto{\pgfqpoint{9.609338in}{0.694229in}}%
\pgfpathlineto{\pgfqpoint{9.609338in}{0.684538in}}%
\pgfpathquadraticcurveto{\pgfqpoint{9.604943in}{0.686970in}}{\pgfqpoint{9.600530in}{0.688177in}}%
\pgfpathquadraticcurveto{\pgfqpoint{9.596117in}{0.689384in}}{\pgfqpoint{9.591614in}{0.689384in}}%
\pgfpathquadraticcurveto{\pgfqpoint{9.581527in}{0.689384in}}{\pgfqpoint{9.575943in}{0.682989in}}%
\pgfpathquadraticcurveto{\pgfqpoint{9.570377in}{0.676613in}}{\pgfqpoint{9.570377in}{0.665067in}}%
\pgfpathquadraticcurveto{\pgfqpoint{9.570377in}{0.653521in}}{\pgfqpoint{9.575943in}{0.647127in}}%
\pgfpathquadraticcurveto{\pgfqpoint{9.581527in}{0.640751in}}{\pgfqpoint{9.591614in}{0.640751in}}%
\pgfpathquadraticcurveto{\pgfqpoint{9.596117in}{0.640751in}}{\pgfqpoint{9.600530in}{0.641958in}}%
\pgfpathquadraticcurveto{\pgfqpoint{9.604943in}{0.643164in}}{\pgfqpoint{9.609338in}{0.645596in}}%
\pgfpathlineto{\pgfqpoint{9.609338in}{0.636014in}}%
\pgfpathquadraticcurveto{\pgfqpoint{9.604997in}{0.633996in}}{\pgfqpoint{9.600350in}{0.632988in}}%
\pgfpathquadraticcurveto{\pgfqpoint{9.595720in}{0.631961in}}{\pgfqpoint{9.590479in}{0.631961in}}%
\pgfpathquadraticcurveto{\pgfqpoint{9.576231in}{0.631961in}}{\pgfqpoint{9.567838in}{0.640913in}}%
\pgfpathquadraticcurveto{\pgfqpoint{9.559462in}{0.649865in}}{\pgfqpoint{9.559462in}{0.665067in}}%
\pgfpathquadraticcurveto{\pgfqpoint{9.559462in}{0.680486in}}{\pgfqpoint{9.567928in}{0.689312in}}%
\pgfpathquadraticcurveto{\pgfqpoint{9.576411in}{0.698156in}}{\pgfqpoint{9.591163in}{0.698156in}}%
\pgfpathquadraticcurveto{\pgfqpoint{9.595937in}{0.698156in}}{\pgfqpoint{9.600494in}{0.697165in}}%
\pgfpathquadraticcurveto{\pgfqpoint{9.605051in}{0.696192in}}{\pgfqpoint{9.609338in}{0.694229in}}%
\pgfpathlineto{\pgfqpoint{9.609338in}{0.694229in}}%
\pgfpathclose%
\pgfpathmoveto{\pgfqpoint{9.637599in}{0.714547in}}%
\pgfpathlineto{\pgfqpoint{9.637599in}{0.696643in}}%
\pgfpathlineto{\pgfqpoint{9.658925in}{0.696643in}}%
\pgfpathlineto{\pgfqpoint{9.658925in}{0.688591in}}%
\pgfpathlineto{\pgfqpoint{9.637599in}{0.688591in}}%
\pgfpathlineto{\pgfqpoint{9.637599in}{0.654368in}}%
\pgfpathquadraticcurveto{\pgfqpoint{9.637599in}{0.646659in}}{\pgfqpoint{9.639706in}{0.644461in}}%
\pgfpathquadraticcurveto{\pgfqpoint{9.641814in}{0.642264in}}{\pgfqpoint{9.648298in}{0.642264in}}%
\pgfpathlineto{\pgfqpoint{9.658925in}{0.642264in}}%
\pgfpathlineto{\pgfqpoint{9.658925in}{0.633600in}}%
\pgfpathlineto{\pgfqpoint{9.648298in}{0.633600in}}%
\pgfpathquadraticcurveto{\pgfqpoint{9.636302in}{0.633600in}}{\pgfqpoint{9.631745in}{0.638067in}}%
\pgfpathquadraticcurveto{\pgfqpoint{9.627188in}{0.642552in}}{\pgfqpoint{9.627188in}{0.654368in}}%
\pgfpathlineto{\pgfqpoint{9.627188in}{0.688591in}}%
\pgfpathlineto{\pgfqpoint{9.619587in}{0.688591in}}%
\pgfpathlineto{\pgfqpoint{9.619587in}{0.696643in}}%
\pgfpathlineto{\pgfqpoint{9.627188in}{0.696643in}}%
\pgfpathlineto{\pgfqpoint{9.627188in}{0.714547in}}%
\pgfpathlineto{\pgfqpoint{9.637599in}{0.714547in}}%
\pgfpathlineto{\pgfqpoint{9.637599in}{0.714547in}}%
\pgfpathclose%
\pgfpathmoveto{\pgfqpoint{9.667319in}{0.669786in}}%
\pgfpathlineto{\pgfqpoint{9.697651in}{0.669786in}}%
\pgfpathlineto{\pgfqpoint{9.697651in}{0.660564in}}%
\pgfpathlineto{\pgfqpoint{9.667319in}{0.660564in}}%
\pgfpathlineto{\pgfqpoint{9.667319in}{0.669786in}}%
\pgfpathlineto{\pgfqpoint{9.667319in}{0.669786in}}%
\pgfpathclose%
\pgfpathmoveto{\pgfqpoint{9.713070in}{0.658475in}}%
\pgfpathlineto{\pgfqpoint{9.713070in}{0.696643in}}%
\pgfpathlineto{\pgfqpoint{9.723427in}{0.696643in}}%
\pgfpathlineto{\pgfqpoint{9.723427in}{0.658871in}}%
\pgfpathquadraticcurveto{\pgfqpoint{9.723427in}{0.649919in}}{\pgfqpoint{9.726903in}{0.645434in}}%
\pgfpathquadraticcurveto{\pgfqpoint{9.730397in}{0.640967in}}{\pgfqpoint{9.737386in}{0.640967in}}%
\pgfpathquadraticcurveto{\pgfqpoint{9.745762in}{0.640967in}}{\pgfqpoint{9.750625in}{0.646317in}}%
\pgfpathquadraticcurveto{\pgfqpoint{9.755506in}{0.651666in}}{\pgfqpoint{9.755506in}{0.660906in}}%
\pgfpathlineto{\pgfqpoint{9.755506in}{0.696643in}}%
\pgfpathlineto{\pgfqpoint{9.765863in}{0.696643in}}%
\pgfpathlineto{\pgfqpoint{9.765863in}{0.633600in}}%
\pgfpathlineto{\pgfqpoint{9.755506in}{0.633600in}}%
\pgfpathlineto{\pgfqpoint{9.755506in}{0.643291in}}%
\pgfpathquadraticcurveto{\pgfqpoint{9.751742in}{0.637545in}}{\pgfqpoint{9.746752in}{0.634753in}}%
\pgfpathquadraticcurveto{\pgfqpoint{9.741781in}{0.631961in}}{\pgfqpoint{9.735189in}{0.631961in}}%
\pgfpathquadraticcurveto{\pgfqpoint{9.724327in}{0.631961in}}{\pgfqpoint{9.718689in}{0.638715in}}%
\pgfpathquadraticcurveto{\pgfqpoint{9.713070in}{0.645470in}}{\pgfqpoint{9.713070in}{0.658475in}}%
\pgfpathlineto{\pgfqpoint{9.713070in}{0.658475in}}%
\pgfpathclose%
\pgfpathmoveto{\pgfqpoint{9.739133in}{0.698156in}}%
\pgfpathlineto{\pgfqpoint{9.739133in}{0.698156in}}%
\pgfpathlineto{\pgfqpoint{9.739133in}{0.698156in}}%
\pgfpathclose%
\pgfpathmoveto{\pgfqpoint{9.839605in}{0.671660in}}%
\pgfpathlineto{\pgfqpoint{9.839605in}{0.633600in}}%
\pgfpathlineto{\pgfqpoint{9.829248in}{0.633600in}}%
\pgfpathlineto{\pgfqpoint{9.829248in}{0.671317in}}%
\pgfpathquadraticcurveto{\pgfqpoint{9.829248in}{0.680269in}}{\pgfqpoint{9.825754in}{0.684700in}}%
\pgfpathquadraticcurveto{\pgfqpoint{9.822259in}{0.689149in}}{\pgfqpoint{9.815289in}{0.689149in}}%
\pgfpathquadraticcurveto{\pgfqpoint{9.806895in}{0.689149in}}{\pgfqpoint{9.802050in}{0.683800in}}%
\pgfpathquadraticcurveto{\pgfqpoint{9.797204in}{0.678468in}}{\pgfqpoint{9.797204in}{0.669228in}}%
\pgfpathlineto{\pgfqpoint{9.797204in}{0.633600in}}%
\pgfpathlineto{\pgfqpoint{9.786793in}{0.633600in}}%
\pgfpathlineto{\pgfqpoint{9.786793in}{0.696643in}}%
\pgfpathlineto{\pgfqpoint{9.797204in}{0.696643in}}%
\pgfpathlineto{\pgfqpoint{9.797204in}{0.686844in}}%
\pgfpathquadraticcurveto{\pgfqpoint{9.800933in}{0.692536in}}{\pgfqpoint{9.805958in}{0.695346in}}%
\pgfpathquadraticcurveto{\pgfqpoint{9.811002in}{0.698156in}}{\pgfqpoint{9.817594in}{0.698156in}}%
\pgfpathquadraticcurveto{\pgfqpoint{9.828456in}{0.698156in}}{\pgfqpoint{9.834021in}{0.691437in}}%
\pgfpathquadraticcurveto{\pgfqpoint{9.839605in}{0.684718in}}{\pgfqpoint{9.839605in}{0.671660in}}%
\pgfpathlineto{\pgfqpoint{9.839605in}{0.671660in}}%
\pgfpathclose%
\pgfpathmoveto{\pgfqpoint{9.860247in}{0.696643in}}%
\pgfpathlineto{\pgfqpoint{9.870604in}{0.696643in}}%
\pgfpathlineto{\pgfqpoint{9.870604in}{0.633600in}}%
\pgfpathlineto{\pgfqpoint{9.860247in}{0.633600in}}%
\pgfpathlineto{\pgfqpoint{9.860247in}{0.696643in}}%
\pgfpathlineto{\pgfqpoint{9.860247in}{0.696643in}}%
\pgfpathclose%
\pgfpathmoveto{\pgfqpoint{9.860247in}{0.721193in}}%
\pgfpathlineto{\pgfqpoint{9.870604in}{0.721193in}}%
\pgfpathlineto{\pgfqpoint{9.870604in}{0.708062in}}%
\pgfpathlineto{\pgfqpoint{9.860247in}{0.708062in}}%
\pgfpathlineto{\pgfqpoint{9.860247in}{0.721193in}}%
\pgfpathlineto{\pgfqpoint{9.860247in}{0.721193in}}%
\pgfpathclose%
\pgfpathmoveto{\pgfqpoint{9.916697in}{0.689384in}}%
\pgfpathquadraticcurveto{\pgfqpoint{9.908375in}{0.689384in}}{\pgfqpoint{9.903530in}{0.682881in}}%
\pgfpathquadraticcurveto{\pgfqpoint{9.898685in}{0.676379in}}{\pgfqpoint{9.898685in}{0.665067in}}%
\pgfpathquadraticcurveto{\pgfqpoint{9.898685in}{0.653756in}}{\pgfqpoint{9.903494in}{0.647253in}}%
\pgfpathquadraticcurveto{\pgfqpoint{9.908321in}{0.640751in}}{\pgfqpoint{9.916697in}{0.640751in}}%
\pgfpathquadraticcurveto{\pgfqpoint{9.924983in}{0.640751in}}{\pgfqpoint{9.929810in}{0.647271in}}%
\pgfpathquadraticcurveto{\pgfqpoint{9.934655in}{0.653810in}}{\pgfqpoint{9.934655in}{0.665067in}}%
\pgfpathquadraticcurveto{\pgfqpoint{9.934655in}{0.676271in}}{\pgfqpoint{9.929810in}{0.682827in}}%
\pgfpathquadraticcurveto{\pgfqpoint{9.924983in}{0.689384in}}{\pgfqpoint{9.916697in}{0.689384in}}%
\pgfpathlineto{\pgfqpoint{9.916697in}{0.689384in}}%
\pgfpathclose%
\pgfpathmoveto{\pgfqpoint{9.916697in}{0.698156in}}%
\pgfpathquadraticcurveto{\pgfqpoint{9.930206in}{0.698156in}}{\pgfqpoint{9.937915in}{0.689366in}}%
\pgfpathquadraticcurveto{\pgfqpoint{9.945643in}{0.680594in}}{\pgfqpoint{9.945643in}{0.665067in}}%
\pgfpathquadraticcurveto{\pgfqpoint{9.945643in}{0.649595in}}{\pgfqpoint{9.937915in}{0.640769in}}%
\pgfpathquadraticcurveto{\pgfqpoint{9.930206in}{0.631961in}}{\pgfqpoint{9.916697in}{0.631961in}}%
\pgfpathquadraticcurveto{\pgfqpoint{9.903134in}{0.631961in}}{\pgfqpoint{9.895443in}{0.640769in}}%
\pgfpathquadraticcurveto{\pgfqpoint{9.887770in}{0.649595in}}{\pgfqpoint{9.887770in}{0.665067in}}%
\pgfpathquadraticcurveto{\pgfqpoint{9.887770in}{0.680594in}}{\pgfqpoint{9.895443in}{0.689366in}}%
\pgfpathquadraticcurveto{\pgfqpoint{9.903134in}{0.698156in}}{\pgfqpoint{9.916697in}{0.698156in}}%
\pgfpathlineto{\pgfqpoint{9.916697in}{0.698156in}}%
\pgfpathclose%
\pgfpathmoveto{\pgfqpoint{10.015224in}{0.671660in}}%
\pgfpathlineto{\pgfqpoint{10.015224in}{0.633600in}}%
\pgfpathlineto{\pgfqpoint{10.004867in}{0.633600in}}%
\pgfpathlineto{\pgfqpoint{10.004867in}{0.671317in}}%
\pgfpathquadraticcurveto{\pgfqpoint{10.004867in}{0.680269in}}{\pgfqpoint{10.001372in}{0.684700in}}%
\pgfpathquadraticcurveto{\pgfqpoint{9.997878in}{0.689149in}}{\pgfqpoint{9.990907in}{0.689149in}}%
\pgfpathquadraticcurveto{\pgfqpoint{9.982513in}{0.689149in}}{\pgfqpoint{9.977668in}{0.683800in}}%
\pgfpathquadraticcurveto{\pgfqpoint{9.972823in}{0.678468in}}{\pgfqpoint{9.972823in}{0.669228in}}%
\pgfpathlineto{\pgfqpoint{9.972823in}{0.633600in}}%
\pgfpathlineto{\pgfqpoint{9.962412in}{0.633600in}}%
\pgfpathlineto{\pgfqpoint{9.962412in}{0.696643in}}%
\pgfpathlineto{\pgfqpoint{9.972823in}{0.696643in}}%
\pgfpathlineto{\pgfqpoint{9.972823in}{0.686844in}}%
\pgfpathquadraticcurveto{\pgfqpoint{9.976551in}{0.692536in}}{\pgfqpoint{9.981577in}{0.695346in}}%
\pgfpathquadraticcurveto{\pgfqpoint{9.986620in}{0.698156in}}{\pgfqpoint{9.993213in}{0.698156in}}%
\pgfpathquadraticcurveto{\pgfqpoint{10.004074in}{0.698156in}}{\pgfqpoint{10.009640in}{0.691437in}}%
\pgfpathquadraticcurveto{\pgfqpoint{10.015224in}{0.684718in}}{\pgfqpoint{10.015224in}{0.671660in}}%
\pgfpathlineto{\pgfqpoint{10.015224in}{0.671660in}}%
\pgfpathclose%
\pgfusepath{fill}%
\end{pgfscope}%
\begin{pgfscope}%
\pgfsetbuttcap%
\pgfsetroundjoin%
\definecolor{currentfill}{rgb}{0.309804,0.372549,0.419608}%
\pgfsetfillcolor{currentfill}%
\pgfsetlinewidth{0.000000pt}%
\definecolor{currentstroke}{rgb}{0.309804,0.372549,0.419608}%
\pgfsetstrokecolor{currentstroke}%
\pgfsetdash{}{0pt}%
\pgfpathmoveto{\pgfqpoint{10.170012in}{0.694787in}}%
\pgfpathlineto{\pgfqpoint{10.170012in}{0.684989in}}%
\pgfpathquadraticcurveto{\pgfqpoint{10.165635in}{0.687240in}}{\pgfqpoint{10.160898in}{0.688357in}}%
\pgfpathquadraticcurveto{\pgfqpoint{10.156178in}{0.689492in}}{\pgfqpoint{10.151099in}{0.689492in}}%
\pgfpathquadraticcurveto{\pgfqpoint{10.143390in}{0.689492in}}{\pgfqpoint{10.139535in}{0.687132in}}%
\pgfpathquadraticcurveto{\pgfqpoint{10.135681in}{0.684773in}}{\pgfqpoint{10.135681in}{0.680035in}}%
\pgfpathquadraticcurveto{\pgfqpoint{10.135681in}{0.676433in}}{\pgfqpoint{10.138436in}{0.674380in}}%
\pgfpathquadraticcurveto{\pgfqpoint{10.141192in}{0.672326in}}{\pgfqpoint{10.149532in}{0.670471in}}%
\pgfpathlineto{\pgfqpoint{10.153080in}{0.669678in}}%
\pgfpathquadraticcurveto{\pgfqpoint{10.164104in}{0.667319in}}{\pgfqpoint{10.168751in}{0.663014in}}%
\pgfpathquadraticcurveto{\pgfqpoint{10.173398in}{0.658709in}}{\pgfqpoint{10.173398in}{0.651000in}}%
\pgfpathquadraticcurveto{\pgfqpoint{10.173398in}{0.642210in}}{\pgfqpoint{10.166445in}{0.637076in}}%
\pgfpathquadraticcurveto{\pgfqpoint{10.159493in}{0.631961in}}{\pgfqpoint{10.147334in}{0.631961in}}%
\pgfpathquadraticcurveto{\pgfqpoint{10.142273in}{0.631961in}}{\pgfqpoint{10.136779in}{0.632952in}}%
\pgfpathquadraticcurveto{\pgfqpoint{10.131286in}{0.633942in}}{\pgfqpoint{10.125215in}{0.635906in}}%
\pgfpathlineto{\pgfqpoint{10.125215in}{0.646605in}}%
\pgfpathquadraticcurveto{\pgfqpoint{10.130961in}{0.643615in}}{\pgfqpoint{10.136527in}{0.642120in}}%
\pgfpathquadraticcurveto{\pgfqpoint{10.142093in}{0.640643in}}{\pgfqpoint{10.147569in}{0.640643in}}%
\pgfpathquadraticcurveto{\pgfqpoint{10.154882in}{0.640643in}}{\pgfqpoint{10.158808in}{0.643146in}}%
\pgfpathquadraticcurveto{\pgfqpoint{10.162753in}{0.645650in}}{\pgfqpoint{10.162753in}{0.650207in}}%
\pgfpathquadraticcurveto{\pgfqpoint{10.162753in}{0.654422in}}{\pgfqpoint{10.159907in}{0.656674in}}%
\pgfpathquadraticcurveto{\pgfqpoint{10.157079in}{0.658925in}}{\pgfqpoint{10.147442in}{0.661014in}}%
\pgfpathlineto{\pgfqpoint{10.143840in}{0.661861in}}%
\pgfpathquadraticcurveto{\pgfqpoint{10.134222in}{0.663878in}}{\pgfqpoint{10.129935in}{0.668075in}}%
\pgfpathquadraticcurveto{\pgfqpoint{10.125666in}{0.672272in}}{\pgfqpoint{10.125666in}{0.679585in}}%
\pgfpathquadraticcurveto{\pgfqpoint{10.125666in}{0.688483in}}{\pgfqpoint{10.131970in}{0.693310in}}%
\pgfpathquadraticcurveto{\pgfqpoint{10.138274in}{0.698156in}}{\pgfqpoint{10.149874in}{0.698156in}}%
\pgfpathquadraticcurveto{\pgfqpoint{10.155602in}{0.698156in}}{\pgfqpoint{10.160663in}{0.697309in}}%
\pgfpathquadraticcurveto{\pgfqpoint{10.165743in}{0.696480in}}{\pgfqpoint{10.170012in}{0.694787in}}%
\pgfpathlineto{\pgfqpoint{10.170012in}{0.694787in}}%
\pgfpathclose%
\pgfpathmoveto{\pgfqpoint{10.218536in}{0.665283in}}%
\pgfpathquadraticcurveto{\pgfqpoint{10.205982in}{0.665283in}}{\pgfqpoint{10.201137in}{0.662419in}}%
\pgfpathquadraticcurveto{\pgfqpoint{10.196291in}{0.659556in}}{\pgfqpoint{10.196291in}{0.652621in}}%
\pgfpathquadraticcurveto{\pgfqpoint{10.196291in}{0.647109in}}{\pgfqpoint{10.199930in}{0.643867in}}%
\pgfpathquadraticcurveto{\pgfqpoint{10.203568in}{0.640643in}}{\pgfqpoint{10.209801in}{0.640643in}}%
\pgfpathquadraticcurveto{\pgfqpoint{10.218428in}{0.640643in}}{\pgfqpoint{10.223634in}{0.646749in}}%
\pgfpathquadraticcurveto{\pgfqpoint{10.228839in}{0.652855in}}{\pgfqpoint{10.228839in}{0.662978in}}%
\pgfpathlineto{\pgfqpoint{10.228839in}{0.665283in}}%
\pgfpathlineto{\pgfqpoint{10.218536in}{0.665283in}}%
\pgfpathlineto{\pgfqpoint{10.218536in}{0.665283in}}%
\pgfpathclose%
\pgfpathmoveto{\pgfqpoint{10.239196in}{0.669570in}}%
\pgfpathlineto{\pgfqpoint{10.239196in}{0.633600in}}%
\pgfpathlineto{\pgfqpoint{10.228839in}{0.633600in}}%
\pgfpathlineto{\pgfqpoint{10.228839in}{0.643164in}}%
\pgfpathquadraticcurveto{\pgfqpoint{10.225291in}{0.637437in}}{\pgfqpoint{10.219995in}{0.634699in}}%
\pgfpathquadraticcurveto{\pgfqpoint{10.214700in}{0.631961in}}{\pgfqpoint{10.207045in}{0.631961in}}%
\pgfpathquadraticcurveto{\pgfqpoint{10.197372in}{0.631961in}}{\pgfqpoint{10.191644in}{0.637401in}}%
\pgfpathquadraticcurveto{\pgfqpoint{10.185934in}{0.642840in}}{\pgfqpoint{10.185934in}{0.651954in}}%
\pgfpathquadraticcurveto{\pgfqpoint{10.185934in}{0.662582in}}{\pgfqpoint{10.193049in}{0.667985in}}%
\pgfpathquadraticcurveto{\pgfqpoint{10.200182in}{0.673389in}}{\pgfqpoint{10.214304in}{0.673389in}}%
\pgfpathlineto{\pgfqpoint{10.228839in}{0.673389in}}%
\pgfpathlineto{\pgfqpoint{10.228839in}{0.674416in}}%
\pgfpathquadraticcurveto{\pgfqpoint{10.228839in}{0.681566in}}{\pgfqpoint{10.224138in}{0.685475in}}%
\pgfpathquadraticcurveto{\pgfqpoint{10.219437in}{0.689384in}}{\pgfqpoint{10.210935in}{0.689384in}}%
\pgfpathquadraticcurveto{\pgfqpoint{10.205532in}{0.689384in}}{\pgfqpoint{10.200398in}{0.688087in}}%
\pgfpathquadraticcurveto{\pgfqpoint{10.195283in}{0.686790in}}{\pgfqpoint{10.190564in}{0.684196in}}%
\pgfpathlineto{\pgfqpoint{10.190564in}{0.693779in}}%
\pgfpathquadraticcurveto{\pgfqpoint{10.196237in}{0.695976in}}{\pgfqpoint{10.201587in}{0.697057in}}%
\pgfpathquadraticcurveto{\pgfqpoint{10.206937in}{0.698156in}}{\pgfqpoint{10.211998in}{0.698156in}}%
\pgfpathquadraticcurveto{\pgfqpoint{10.225687in}{0.698156in}}{\pgfqpoint{10.232442in}{0.691059in}}%
\pgfpathquadraticcurveto{\pgfqpoint{10.239196in}{0.683980in}}{\pgfqpoint{10.239196in}{0.669570in}}%
\pgfpathlineto{\pgfqpoint{10.239196in}{0.669570in}}%
\pgfpathclose%
\pgfpathmoveto{\pgfqpoint{10.309606in}{0.684538in}}%
\pgfpathquadraticcurveto{\pgfqpoint{10.313497in}{0.691527in}}{\pgfqpoint{10.318900in}{0.694841in}}%
\pgfpathquadraticcurveto{\pgfqpoint{10.324304in}{0.698156in}}{\pgfqpoint{10.331617in}{0.698156in}}%
\pgfpathquadraticcurveto{\pgfqpoint{10.341469in}{0.698156in}}{\pgfqpoint{10.346819in}{0.691257in}}%
\pgfpathquadraticcurveto{\pgfqpoint{10.352169in}{0.684376in}}{\pgfqpoint{10.352169in}{0.671660in}}%
\pgfpathlineto{\pgfqpoint{10.352169in}{0.633600in}}%
\pgfpathlineto{\pgfqpoint{10.341758in}{0.633600in}}%
\pgfpathlineto{\pgfqpoint{10.341758in}{0.671317in}}%
\pgfpathquadraticcurveto{\pgfqpoint{10.341758in}{0.680378in}}{\pgfqpoint{10.338533in}{0.684755in}}%
\pgfpathquadraticcurveto{\pgfqpoint{10.335327in}{0.689149in}}{\pgfqpoint{10.328753in}{0.689149in}}%
\pgfpathquadraticcurveto{\pgfqpoint{10.320701in}{0.689149in}}{\pgfqpoint{10.316018in}{0.683800in}}%
\pgfpathquadraticcurveto{\pgfqpoint{10.311353in}{0.678468in}}{\pgfqpoint{10.311353in}{0.669228in}}%
\pgfpathlineto{\pgfqpoint{10.311353in}{0.633600in}}%
\pgfpathlineto{\pgfqpoint{10.300942in}{0.633600in}}%
\pgfpathlineto{\pgfqpoint{10.300942in}{0.671317in}}%
\pgfpathquadraticcurveto{\pgfqpoint{10.300942in}{0.680432in}}{\pgfqpoint{10.297736in}{0.684791in}}%
\pgfpathquadraticcurveto{\pgfqpoint{10.294530in}{0.689149in}}{\pgfqpoint{10.287829in}{0.689149in}}%
\pgfpathquadraticcurveto{\pgfqpoint{10.279886in}{0.689149in}}{\pgfqpoint{10.275203in}{0.683782in}}%
\pgfpathquadraticcurveto{\pgfqpoint{10.270538in}{0.678414in}}{\pgfqpoint{10.270538in}{0.669228in}}%
\pgfpathlineto{\pgfqpoint{10.270538in}{0.633600in}}%
\pgfpathlineto{\pgfqpoint{10.260127in}{0.633600in}}%
\pgfpathlineto{\pgfqpoint{10.260127in}{0.696643in}}%
\pgfpathlineto{\pgfqpoint{10.270538in}{0.696643in}}%
\pgfpathlineto{\pgfqpoint{10.270538in}{0.686844in}}%
\pgfpathquadraticcurveto{\pgfqpoint{10.274086in}{0.692644in}}{\pgfqpoint{10.279039in}{0.695400in}}%
\pgfpathquadraticcurveto{\pgfqpoint{10.283993in}{0.698156in}}{\pgfqpoint{10.290801in}{0.698156in}}%
\pgfpathquadraticcurveto{\pgfqpoint{10.297682in}{0.698156in}}{\pgfqpoint{10.302491in}{0.694661in}}%
\pgfpathquadraticcurveto{\pgfqpoint{10.307300in}{0.691185in}}{\pgfqpoint{10.309606in}{0.684538in}}%
\pgfpathlineto{\pgfqpoint{10.309606in}{0.684538in}}%
\pgfpathclose%
\pgfpathmoveto{\pgfqpoint{10.382825in}{0.643056in}}%
\pgfpathlineto{\pgfqpoint{10.382825in}{0.609626in}}%
\pgfpathlineto{\pgfqpoint{10.372414in}{0.609626in}}%
\pgfpathlineto{\pgfqpoint{10.372414in}{0.696643in}}%
\pgfpathlineto{\pgfqpoint{10.382825in}{0.696643in}}%
\pgfpathlineto{\pgfqpoint{10.382825in}{0.687078in}}%
\pgfpathquadraticcurveto{\pgfqpoint{10.386104in}{0.692698in}}{\pgfqpoint{10.391075in}{0.695418in}}%
\pgfpathquadraticcurveto{\pgfqpoint{10.396064in}{0.698156in}}{\pgfqpoint{10.402981in}{0.698156in}}%
\pgfpathquadraticcurveto{\pgfqpoint{10.414473in}{0.698156in}}{\pgfqpoint{10.421641in}{0.689041in}}%
\pgfpathquadraticcurveto{\pgfqpoint{10.428828in}{0.679927in}}{\pgfqpoint{10.428828in}{0.665067in}}%
\pgfpathquadraticcurveto{\pgfqpoint{10.428828in}{0.650207in}}{\pgfqpoint{10.421641in}{0.641075in}}%
\pgfpathquadraticcurveto{\pgfqpoint{10.414473in}{0.631961in}}{\pgfqpoint{10.402981in}{0.631961in}}%
\pgfpathquadraticcurveto{\pgfqpoint{10.396064in}{0.631961in}}{\pgfqpoint{10.391075in}{0.634699in}}%
\pgfpathquadraticcurveto{\pgfqpoint{10.386104in}{0.637437in}}{\pgfqpoint{10.382825in}{0.643056in}}%
\pgfpathlineto{\pgfqpoint{10.382825in}{0.643056in}}%
\pgfpathclose%
\pgfpathmoveto{\pgfqpoint{10.418075in}{0.665067in}}%
\pgfpathquadraticcurveto{\pgfqpoint{10.418075in}{0.676487in}}{\pgfqpoint{10.413374in}{0.682989in}}%
\pgfpathquadraticcurveto{\pgfqpoint{10.408673in}{0.689492in}}{\pgfqpoint{10.400459in}{0.689492in}}%
\pgfpathquadraticcurveto{\pgfqpoint{10.392228in}{0.689492in}}{\pgfqpoint{10.387526in}{0.682989in}}%
\pgfpathquadraticcurveto{\pgfqpoint{10.382825in}{0.676487in}}{\pgfqpoint{10.382825in}{0.665067in}}%
\pgfpathquadraticcurveto{\pgfqpoint{10.382825in}{0.653648in}}{\pgfqpoint{10.387526in}{0.647145in}}%
\pgfpathquadraticcurveto{\pgfqpoint{10.392228in}{0.640643in}}{\pgfqpoint{10.400459in}{0.640643in}}%
\pgfpathquadraticcurveto{\pgfqpoint{10.408673in}{0.640643in}}{\pgfqpoint{10.413374in}{0.647145in}}%
\pgfpathquadraticcurveto{\pgfqpoint{10.418075in}{0.653648in}}{\pgfqpoint{10.418075in}{0.665067in}}%
\pgfpathlineto{\pgfqpoint{10.418075in}{0.665067in}}%
\pgfpathclose%
\pgfpathmoveto{\pgfqpoint{10.445994in}{0.721193in}}%
\pgfpathlineto{\pgfqpoint{10.456351in}{0.721193in}}%
\pgfpathlineto{\pgfqpoint{10.456351in}{0.633600in}}%
\pgfpathlineto{\pgfqpoint{10.445994in}{0.633600in}}%
\pgfpathlineto{\pgfqpoint{10.445994in}{0.721193in}}%
\pgfpathlineto{\pgfqpoint{10.445994in}{0.721193in}}%
\pgfpathclose%
\pgfpathmoveto{\pgfqpoint{10.531948in}{0.667715in}}%
\pgfpathlineto{\pgfqpoint{10.531948in}{0.662654in}}%
\pgfpathlineto{\pgfqpoint{10.484324in}{0.662654in}}%
\pgfpathquadraticcurveto{\pgfqpoint{10.485008in}{0.651954in}}{\pgfqpoint{10.490772in}{0.646353in}}%
\pgfpathquadraticcurveto{\pgfqpoint{10.496536in}{0.640751in}}{\pgfqpoint{10.506839in}{0.640751in}}%
\pgfpathquadraticcurveto{\pgfqpoint{10.512801in}{0.640751in}}{\pgfqpoint{10.518403in}{0.642210in}}%
\pgfpathquadraticcurveto{\pgfqpoint{10.524005in}{0.643669in}}{\pgfqpoint{10.529534in}{0.646605in}}%
\pgfpathlineto{\pgfqpoint{10.529534in}{0.636806in}}%
\pgfpathquadraticcurveto{\pgfqpoint{10.523951in}{0.634447in}}{\pgfqpoint{10.518097in}{0.633204in}}%
\pgfpathquadraticcurveto{\pgfqpoint{10.512243in}{0.631961in}}{\pgfqpoint{10.506227in}{0.631961in}}%
\pgfpathquadraticcurveto{\pgfqpoint{10.491132in}{0.631961in}}{\pgfqpoint{10.482324in}{0.640733in}}%
\pgfpathquadraticcurveto{\pgfqpoint{10.473516in}{0.649523in}}{\pgfqpoint{10.473516in}{0.664509in}}%
\pgfpathquadraticcurveto{\pgfqpoint{10.473516in}{0.679981in}}{\pgfqpoint{10.481874in}{0.689059in}}%
\pgfpathquadraticcurveto{\pgfqpoint{10.490232in}{0.698156in}}{\pgfqpoint{10.504425in}{0.698156in}}%
\pgfpathquadraticcurveto{\pgfqpoint{10.517142in}{0.698156in}}{\pgfqpoint{10.524545in}{0.689960in}}%
\pgfpathquadraticcurveto{\pgfqpoint{10.531948in}{0.681783in}}{\pgfqpoint{10.531948in}{0.667715in}}%
\pgfpathlineto{\pgfqpoint{10.531948in}{0.667715in}}%
\pgfpathclose%
\pgfpathmoveto{\pgfqpoint{10.521591in}{0.670759in}}%
\pgfpathquadraticcurveto{\pgfqpoint{10.521483in}{0.679243in}}{\pgfqpoint{10.516836in}{0.684304in}}%
\pgfpathquadraticcurveto{\pgfqpoint{10.512189in}{0.689384in}}{\pgfqpoint{10.504533in}{0.689384in}}%
\pgfpathquadraticcurveto{\pgfqpoint{10.495870in}{0.689384in}}{\pgfqpoint{10.490664in}{0.684484in}}%
\pgfpathquadraticcurveto{\pgfqpoint{10.485459in}{0.679585in}}{\pgfqpoint{10.484666in}{0.670687in}}%
\pgfpathlineto{\pgfqpoint{10.521591in}{0.670759in}}%
\pgfpathlineto{\pgfqpoint{10.521591in}{0.670759in}}%
\pgfpathclose%
\pgfusepath{fill}%
\end{pgfscope}%
\begin{pgfscope}%
\pgfsetbuttcap%
\pgfsetroundjoin%
\definecolor{currentfill}{rgb}{0.309804,0.372549,0.419608}%
\pgfsetfillcolor{currentfill}%
\pgfsetlinewidth{0.000000pt}%
\definecolor{currentstroke}{rgb}{0.309804,0.372549,0.419608}%
\pgfsetstrokecolor{currentstroke}%
\pgfsetdash{}{0pt}%
\pgfpathmoveto{\pgfqpoint{10.646060in}{0.665283in}}%
\pgfpathquadraticcurveto{\pgfqpoint{10.633506in}{0.665283in}}{\pgfqpoint{10.628661in}{0.662419in}}%
\pgfpathquadraticcurveto{\pgfqpoint{10.623815in}{0.659556in}}{\pgfqpoint{10.623815in}{0.652621in}}%
\pgfpathquadraticcurveto{\pgfqpoint{10.623815in}{0.647109in}}{\pgfqpoint{10.627454in}{0.643867in}}%
\pgfpathquadraticcurveto{\pgfqpoint{10.631092in}{0.640643in}}{\pgfqpoint{10.637324in}{0.640643in}}%
\pgfpathquadraticcurveto{\pgfqpoint{10.645952in}{0.640643in}}{\pgfqpoint{10.651158in}{0.646749in}}%
\pgfpathquadraticcurveto{\pgfqpoint{10.656363in}{0.652855in}}{\pgfqpoint{10.656363in}{0.662978in}}%
\pgfpathlineto{\pgfqpoint{10.656363in}{0.665283in}}%
\pgfpathlineto{\pgfqpoint{10.646060in}{0.665283in}}%
\pgfpathlineto{\pgfqpoint{10.646060in}{0.665283in}}%
\pgfpathclose%
\pgfpathmoveto{\pgfqpoint{10.666720in}{0.669570in}}%
\pgfpathlineto{\pgfqpoint{10.666720in}{0.633600in}}%
\pgfpathlineto{\pgfqpoint{10.656363in}{0.633600in}}%
\pgfpathlineto{\pgfqpoint{10.656363in}{0.643164in}}%
\pgfpathquadraticcurveto{\pgfqpoint{10.652815in}{0.637437in}}{\pgfqpoint{10.647519in}{0.634699in}}%
\pgfpathquadraticcurveto{\pgfqpoint{10.642224in}{0.631961in}}{\pgfqpoint{10.634569in}{0.631961in}}%
\pgfpathquadraticcurveto{\pgfqpoint{10.624896in}{0.631961in}}{\pgfqpoint{10.619168in}{0.637401in}}%
\pgfpathquadraticcurveto{\pgfqpoint{10.613458in}{0.642840in}}{\pgfqpoint{10.613458in}{0.651954in}}%
\pgfpathquadraticcurveto{\pgfqpoint{10.613458in}{0.662582in}}{\pgfqpoint{10.620573in}{0.667985in}}%
\pgfpathquadraticcurveto{\pgfqpoint{10.627706in}{0.673389in}}{\pgfqpoint{10.641827in}{0.673389in}}%
\pgfpathlineto{\pgfqpoint{10.656363in}{0.673389in}}%
\pgfpathlineto{\pgfqpoint{10.656363in}{0.674416in}}%
\pgfpathquadraticcurveto{\pgfqpoint{10.656363in}{0.681566in}}{\pgfqpoint{10.651662in}{0.685475in}}%
\pgfpathquadraticcurveto{\pgfqpoint{10.646961in}{0.689384in}}{\pgfqpoint{10.638459in}{0.689384in}}%
\pgfpathquadraticcurveto{\pgfqpoint{10.633056in}{0.689384in}}{\pgfqpoint{10.627922in}{0.688087in}}%
\pgfpathquadraticcurveto{\pgfqpoint{10.622807in}{0.686790in}}{\pgfqpoint{10.618087in}{0.684196in}}%
\pgfpathlineto{\pgfqpoint{10.618087in}{0.693779in}}%
\pgfpathquadraticcurveto{\pgfqpoint{10.623761in}{0.695976in}}{\pgfqpoint{10.629111in}{0.697057in}}%
\pgfpathquadraticcurveto{\pgfqpoint{10.634461in}{0.698156in}}{\pgfqpoint{10.639522in}{0.698156in}}%
\pgfpathquadraticcurveto{\pgfqpoint{10.653211in}{0.698156in}}{\pgfqpoint{10.659966in}{0.691059in}}%
\pgfpathquadraticcurveto{\pgfqpoint{10.666720in}{0.683980in}}{\pgfqpoint{10.666720in}{0.669570in}}%
\pgfpathlineto{\pgfqpoint{10.666720in}{0.669570in}}%
\pgfpathclose%
\pgfpathmoveto{\pgfqpoint{10.688047in}{0.721193in}}%
\pgfpathlineto{\pgfqpoint{10.698404in}{0.721193in}}%
\pgfpathlineto{\pgfqpoint{10.698404in}{0.633600in}}%
\pgfpathlineto{\pgfqpoint{10.688047in}{0.633600in}}%
\pgfpathlineto{\pgfqpoint{10.688047in}{0.721193in}}%
\pgfpathlineto{\pgfqpoint{10.688047in}{0.721193in}}%
\pgfpathclose%
\pgfpathmoveto{\pgfqpoint{10.720072in}{0.696643in}}%
\pgfpathlineto{\pgfqpoint{10.730429in}{0.696643in}}%
\pgfpathlineto{\pgfqpoint{10.730429in}{0.633600in}}%
\pgfpathlineto{\pgfqpoint{10.720072in}{0.633600in}}%
\pgfpathlineto{\pgfqpoint{10.720072in}{0.696643in}}%
\pgfpathlineto{\pgfqpoint{10.720072in}{0.696643in}}%
\pgfpathclose%
\pgfpathmoveto{\pgfqpoint{10.720072in}{0.721193in}}%
\pgfpathlineto{\pgfqpoint{10.730429in}{0.721193in}}%
\pgfpathlineto{\pgfqpoint{10.730429in}{0.708062in}}%
\pgfpathlineto{\pgfqpoint{10.720072in}{0.708062in}}%
\pgfpathlineto{\pgfqpoint{10.720072in}{0.721193in}}%
\pgfpathlineto{\pgfqpoint{10.720072in}{0.721193in}}%
\pgfpathclose%
\pgfpathmoveto{\pgfqpoint{10.793580in}{0.665860in}}%
\pgfpathquadraticcurveto{\pgfqpoint{10.793580in}{0.677117in}}{\pgfqpoint{10.788933in}{0.683296in}}%
\pgfpathquadraticcurveto{\pgfqpoint{10.784304in}{0.689492in}}{\pgfqpoint{10.775910in}{0.689492in}}%
\pgfpathquadraticcurveto{\pgfqpoint{10.767588in}{0.689492in}}{\pgfqpoint{10.762941in}{0.683296in}}%
\pgfpathquadraticcurveto{\pgfqpoint{10.758294in}{0.677117in}}{\pgfqpoint{10.758294in}{0.665860in}}%
\pgfpathquadraticcurveto{\pgfqpoint{10.758294in}{0.654656in}}{\pgfqpoint{10.762941in}{0.648460in}}%
\pgfpathquadraticcurveto{\pgfqpoint{10.767588in}{0.642264in}}{\pgfqpoint{10.775910in}{0.642264in}}%
\pgfpathquadraticcurveto{\pgfqpoint{10.784304in}{0.642264in}}{\pgfqpoint{10.788933in}{0.648460in}}%
\pgfpathquadraticcurveto{\pgfqpoint{10.793580in}{0.654656in}}{\pgfqpoint{10.793580in}{0.665860in}}%
\pgfpathlineto{\pgfqpoint{10.793580in}{0.665860in}}%
\pgfpathclose%
\pgfpathmoveto{\pgfqpoint{10.803937in}{0.641417in}}%
\pgfpathquadraticcurveto{\pgfqpoint{10.803937in}{0.625332in}}{\pgfqpoint{10.796786in}{0.617479in}}%
\pgfpathquadraticcurveto{\pgfqpoint{10.789653in}{0.609626in}}{\pgfqpoint{10.774901in}{0.609626in}}%
\pgfpathquadraticcurveto{\pgfqpoint{10.769444in}{0.609626in}}{\pgfqpoint{10.764598in}{0.610436in}}%
\pgfpathquadraticcurveto{\pgfqpoint{10.759753in}{0.611247in}}{\pgfqpoint{10.755196in}{0.612940in}}%
\pgfpathlineto{\pgfqpoint{10.755196in}{0.623009in}}%
\pgfpathquadraticcurveto{\pgfqpoint{10.759753in}{0.620541in}}{\pgfqpoint{10.764202in}{0.619370in}}%
\pgfpathquadraticcurveto{\pgfqpoint{10.768651in}{0.618182in}}{\pgfqpoint{10.773262in}{0.618182in}}%
\pgfpathquadraticcurveto{\pgfqpoint{10.783457in}{0.618182in}}{\pgfqpoint{10.788518in}{0.623495in}}%
\pgfpathquadraticcurveto{\pgfqpoint{10.793580in}{0.628809in}}{\pgfqpoint{10.793580in}{0.639562in}}%
\pgfpathlineto{\pgfqpoint{10.793580in}{0.644695in}}%
\pgfpathquadraticcurveto{\pgfqpoint{10.790374in}{0.639112in}}{\pgfqpoint{10.785366in}{0.636356in}}%
\pgfpathquadraticcurveto{\pgfqpoint{10.780359in}{0.633600in}}{\pgfqpoint{10.773370in}{0.633600in}}%
\pgfpathquadraticcurveto{\pgfqpoint{10.761788in}{0.633600in}}{\pgfqpoint{10.754692in}{0.642426in}}%
\pgfpathquadraticcurveto{\pgfqpoint{10.747595in}{0.651270in}}{\pgfqpoint{10.747595in}{0.665860in}}%
\pgfpathquadraticcurveto{\pgfqpoint{10.747595in}{0.680486in}}{\pgfqpoint{10.754692in}{0.689312in}}%
\pgfpathquadraticcurveto{\pgfqpoint{10.761788in}{0.698156in}}{\pgfqpoint{10.773370in}{0.698156in}}%
\pgfpathquadraticcurveto{\pgfqpoint{10.780359in}{0.698156in}}{\pgfqpoint{10.785366in}{0.695400in}}%
\pgfpathquadraticcurveto{\pgfqpoint{10.790374in}{0.692644in}}{\pgfqpoint{10.793580in}{0.687078in}}%
\pgfpathlineto{\pgfqpoint{10.793580in}{0.696643in}}%
\pgfpathlineto{\pgfqpoint{10.803937in}{0.696643in}}%
\pgfpathlineto{\pgfqpoint{10.803937in}{0.641417in}}%
\pgfpathlineto{\pgfqpoint{10.803937in}{0.641417in}}%
\pgfpathclose%
\pgfpathmoveto{\pgfqpoint{10.877697in}{0.671660in}}%
\pgfpathlineto{\pgfqpoint{10.877697in}{0.633600in}}%
\pgfpathlineto{\pgfqpoint{10.867340in}{0.633600in}}%
\pgfpathlineto{\pgfqpoint{10.867340in}{0.671317in}}%
\pgfpathquadraticcurveto{\pgfqpoint{10.867340in}{0.680269in}}{\pgfqpoint{10.863845in}{0.684700in}}%
\pgfpathquadraticcurveto{\pgfqpoint{10.860351in}{0.689149in}}{\pgfqpoint{10.853380in}{0.689149in}}%
\pgfpathquadraticcurveto{\pgfqpoint{10.844987in}{0.689149in}}{\pgfqpoint{10.840141in}{0.683800in}}%
\pgfpathquadraticcurveto{\pgfqpoint{10.835296in}{0.678468in}}{\pgfqpoint{10.835296in}{0.669228in}}%
\pgfpathlineto{\pgfqpoint{10.835296in}{0.633600in}}%
\pgfpathlineto{\pgfqpoint{10.824885in}{0.633600in}}%
\pgfpathlineto{\pgfqpoint{10.824885in}{0.696643in}}%
\pgfpathlineto{\pgfqpoint{10.835296in}{0.696643in}}%
\pgfpathlineto{\pgfqpoint{10.835296in}{0.686844in}}%
\pgfpathquadraticcurveto{\pgfqpoint{10.839025in}{0.692536in}}{\pgfqpoint{10.844050in}{0.695346in}}%
\pgfpathquadraticcurveto{\pgfqpoint{10.849093in}{0.698156in}}{\pgfqpoint{10.855686in}{0.698156in}}%
\pgfpathquadraticcurveto{\pgfqpoint{10.866547in}{0.698156in}}{\pgfqpoint{10.872113in}{0.691437in}}%
\pgfpathquadraticcurveto{\pgfqpoint{10.877697in}{0.684718in}}{\pgfqpoint{10.877697in}{0.671660in}}%
\pgfpathlineto{\pgfqpoint{10.877697in}{0.671660in}}%
\pgfpathclose%
\pgfpathmoveto{\pgfqpoint{10.947422in}{0.684538in}}%
\pgfpathquadraticcurveto{\pgfqpoint{10.951312in}{0.691527in}}{\pgfqpoint{10.956716in}{0.694841in}}%
\pgfpathquadraticcurveto{\pgfqpoint{10.962120in}{0.698156in}}{\pgfqpoint{10.969433in}{0.698156in}}%
\pgfpathquadraticcurveto{\pgfqpoint{10.979285in}{0.698156in}}{\pgfqpoint{10.984635in}{0.691257in}}%
\pgfpathquadraticcurveto{\pgfqpoint{10.989984in}{0.684376in}}{\pgfqpoint{10.989984in}{0.671660in}}%
\pgfpathlineto{\pgfqpoint{10.989984in}{0.633600in}}%
\pgfpathlineto{\pgfqpoint{10.979573in}{0.633600in}}%
\pgfpathlineto{\pgfqpoint{10.979573in}{0.671317in}}%
\pgfpathquadraticcurveto{\pgfqpoint{10.979573in}{0.680378in}}{\pgfqpoint{10.976349in}{0.684755in}}%
\pgfpathquadraticcurveto{\pgfqpoint{10.973143in}{0.689149in}}{\pgfqpoint{10.966569in}{0.689149in}}%
\pgfpathquadraticcurveto{\pgfqpoint{10.958517in}{0.689149in}}{\pgfqpoint{10.953834in}{0.683800in}}%
\pgfpathquadraticcurveto{\pgfqpoint{10.949169in}{0.678468in}}{\pgfqpoint{10.949169in}{0.669228in}}%
\pgfpathlineto{\pgfqpoint{10.949169in}{0.633600in}}%
\pgfpathlineto{\pgfqpoint{10.938758in}{0.633600in}}%
\pgfpathlineto{\pgfqpoint{10.938758in}{0.671317in}}%
\pgfpathquadraticcurveto{\pgfqpoint{10.938758in}{0.680432in}}{\pgfqpoint{10.935552in}{0.684791in}}%
\pgfpathquadraticcurveto{\pgfqpoint{10.932345in}{0.689149in}}{\pgfqpoint{10.925645in}{0.689149in}}%
\pgfpathquadraticcurveto{\pgfqpoint{10.917702in}{0.689149in}}{\pgfqpoint{10.913018in}{0.683782in}}%
\pgfpathquadraticcurveto{\pgfqpoint{10.908353in}{0.678414in}}{\pgfqpoint{10.908353in}{0.669228in}}%
\pgfpathlineto{\pgfqpoint{10.908353in}{0.633600in}}%
\pgfpathlineto{\pgfqpoint{10.897942in}{0.633600in}}%
\pgfpathlineto{\pgfqpoint{10.897942in}{0.696643in}}%
\pgfpathlineto{\pgfqpoint{10.908353in}{0.696643in}}%
\pgfpathlineto{\pgfqpoint{10.908353in}{0.686844in}}%
\pgfpathquadraticcurveto{\pgfqpoint{10.911902in}{0.692644in}}{\pgfqpoint{10.916855in}{0.695400in}}%
\pgfpathquadraticcurveto{\pgfqpoint{10.921808in}{0.698156in}}{\pgfqpoint{10.928617in}{0.698156in}}%
\pgfpathquadraticcurveto{\pgfqpoint{10.935498in}{0.698156in}}{\pgfqpoint{10.940307in}{0.694661in}}%
\pgfpathquadraticcurveto{\pgfqpoint{10.945116in}{0.691185in}}{\pgfqpoint{10.947422in}{0.684538in}}%
\pgfpathlineto{\pgfqpoint{10.947422in}{0.684538in}}%
\pgfpathclose%
\pgfpathmoveto{\pgfqpoint{11.064555in}{0.667715in}}%
\pgfpathlineto{\pgfqpoint{11.064555in}{0.662654in}}%
\pgfpathlineto{\pgfqpoint{11.016931in}{0.662654in}}%
\pgfpathquadraticcurveto{\pgfqpoint{11.017615in}{0.651954in}}{\pgfqpoint{11.023379in}{0.646353in}}%
\pgfpathquadraticcurveto{\pgfqpoint{11.029143in}{0.640751in}}{\pgfqpoint{11.039446in}{0.640751in}}%
\pgfpathquadraticcurveto{\pgfqpoint{11.045408in}{0.640751in}}{\pgfqpoint{11.051010in}{0.642210in}}%
\pgfpathquadraticcurveto{\pgfqpoint{11.056611in}{0.643669in}}{\pgfqpoint{11.062141in}{0.646605in}}%
\pgfpathlineto{\pgfqpoint{11.062141in}{0.636806in}}%
\pgfpathquadraticcurveto{\pgfqpoint{11.056557in}{0.634447in}}{\pgfqpoint{11.050703in}{0.633204in}}%
\pgfpathquadraticcurveto{\pgfqpoint{11.044849in}{0.631961in}}{\pgfqpoint{11.038833in}{0.631961in}}%
\pgfpathquadraticcurveto{\pgfqpoint{11.023739in}{0.631961in}}{\pgfqpoint{11.014931in}{0.640733in}}%
\pgfpathquadraticcurveto{\pgfqpoint{11.006123in}{0.649523in}}{\pgfqpoint{11.006123in}{0.664509in}}%
\pgfpathquadraticcurveto{\pgfqpoint{11.006123in}{0.679981in}}{\pgfqpoint{11.014481in}{0.689059in}}%
\pgfpathquadraticcurveto{\pgfqpoint{11.022839in}{0.698156in}}{\pgfqpoint{11.037032in}{0.698156in}}%
\pgfpathquadraticcurveto{\pgfqpoint{11.049749in}{0.698156in}}{\pgfqpoint{11.057152in}{0.689960in}}%
\pgfpathquadraticcurveto{\pgfqpoint{11.064555in}{0.681783in}}{\pgfqpoint{11.064555in}{0.667715in}}%
\pgfpathlineto{\pgfqpoint{11.064555in}{0.667715in}}%
\pgfpathclose%
\pgfpathmoveto{\pgfqpoint{11.054198in}{0.670759in}}%
\pgfpathquadraticcurveto{\pgfqpoint{11.054090in}{0.679243in}}{\pgfqpoint{11.049442in}{0.684304in}}%
\pgfpathquadraticcurveto{\pgfqpoint{11.044795in}{0.689384in}}{\pgfqpoint{11.037140in}{0.689384in}}%
\pgfpathquadraticcurveto{\pgfqpoint{11.028476in}{0.689384in}}{\pgfqpoint{11.023271in}{0.684484in}}%
\pgfpathquadraticcurveto{\pgfqpoint{11.018065in}{0.679585in}}{\pgfqpoint{11.017273in}{0.670687in}}%
\pgfpathlineto{\pgfqpoint{11.054198in}{0.670759in}}%
\pgfpathlineto{\pgfqpoint{11.054198in}{0.670759in}}%
\pgfpathclose%
\pgfpathmoveto{\pgfqpoint{11.133974in}{0.671660in}}%
\pgfpathlineto{\pgfqpoint{11.133974in}{0.633600in}}%
\pgfpathlineto{\pgfqpoint{11.123617in}{0.633600in}}%
\pgfpathlineto{\pgfqpoint{11.123617in}{0.671317in}}%
\pgfpathquadraticcurveto{\pgfqpoint{11.123617in}{0.680269in}}{\pgfqpoint{11.120122in}{0.684700in}}%
\pgfpathquadraticcurveto{\pgfqpoint{11.116628in}{0.689149in}}{\pgfqpoint{11.109657in}{0.689149in}}%
\pgfpathquadraticcurveto{\pgfqpoint{11.101263in}{0.689149in}}{\pgfqpoint{11.096418in}{0.683800in}}%
\pgfpathquadraticcurveto{\pgfqpoint{11.091573in}{0.678468in}}{\pgfqpoint{11.091573in}{0.669228in}}%
\pgfpathlineto{\pgfqpoint{11.091573in}{0.633600in}}%
\pgfpathlineto{\pgfqpoint{11.081162in}{0.633600in}}%
\pgfpathlineto{\pgfqpoint{11.081162in}{0.696643in}}%
\pgfpathlineto{\pgfqpoint{11.091573in}{0.696643in}}%
\pgfpathlineto{\pgfqpoint{11.091573in}{0.686844in}}%
\pgfpathquadraticcurveto{\pgfqpoint{11.095301in}{0.692536in}}{\pgfqpoint{11.100327in}{0.695346in}}%
\pgfpathquadraticcurveto{\pgfqpoint{11.105370in}{0.698156in}}{\pgfqpoint{11.111963in}{0.698156in}}%
\pgfpathquadraticcurveto{\pgfqpoint{11.122824in}{0.698156in}}{\pgfqpoint{11.128390in}{0.691437in}}%
\pgfpathquadraticcurveto{\pgfqpoint{11.133974in}{0.684718in}}{\pgfqpoint{11.133974in}{0.671660in}}%
\pgfpathlineto{\pgfqpoint{11.133974in}{0.671660in}}%
\pgfpathclose%
\pgfpathmoveto{\pgfqpoint{11.164864in}{0.714547in}}%
\pgfpathlineto{\pgfqpoint{11.164864in}{0.696643in}}%
\pgfpathlineto{\pgfqpoint{11.186191in}{0.696643in}}%
\pgfpathlineto{\pgfqpoint{11.186191in}{0.688591in}}%
\pgfpathlineto{\pgfqpoint{11.164864in}{0.688591in}}%
\pgfpathlineto{\pgfqpoint{11.164864in}{0.654368in}}%
\pgfpathquadraticcurveto{\pgfqpoint{11.164864in}{0.646659in}}{\pgfqpoint{11.166972in}{0.644461in}}%
\pgfpathquadraticcurveto{\pgfqpoint{11.169079in}{0.642264in}}{\pgfqpoint{11.175564in}{0.642264in}}%
\pgfpathlineto{\pgfqpoint{11.186191in}{0.642264in}}%
\pgfpathlineto{\pgfqpoint{11.186191in}{0.633600in}}%
\pgfpathlineto{\pgfqpoint{11.175564in}{0.633600in}}%
\pgfpathquadraticcurveto{\pgfqpoint{11.163567in}{0.633600in}}{\pgfqpoint{11.159010in}{0.638067in}}%
\pgfpathquadraticcurveto{\pgfqpoint{11.154453in}{0.642552in}}{\pgfqpoint{11.154453in}{0.654368in}}%
\pgfpathlineto{\pgfqpoint{11.154453in}{0.688591in}}%
\pgfpathlineto{\pgfqpoint{11.146852in}{0.688591in}}%
\pgfpathlineto{\pgfqpoint{11.146852in}{0.696643in}}%
\pgfpathlineto{\pgfqpoint{11.154453in}{0.696643in}}%
\pgfpathlineto{\pgfqpoint{11.154453in}{0.714547in}}%
\pgfpathlineto{\pgfqpoint{11.164864in}{0.714547in}}%
\pgfpathlineto{\pgfqpoint{11.164864in}{0.714547in}}%
\pgfpathclose%
\pgfusepath{fill}%
\end{pgfscope}%
\begin{pgfscope}%
\pgfsetbuttcap%
\pgfsetroundjoin%
\definecolor{currentfill}{rgb}{0.309804,0.372549,0.419608}%
\pgfsetfillcolor{currentfill}%
\pgfsetlinewidth{0.000000pt}%
\definecolor{currentstroke}{rgb}{0.309804,0.372549,0.419608}%
\pgfsetstrokecolor{currentstroke}%
\pgfsetdash{}{0pt}%
\pgfpathmoveto{\pgfqpoint{11.326461in}{0.687078in}}%
\pgfpathlineto{\pgfqpoint{11.326461in}{0.721193in}}%
\pgfpathlineto{\pgfqpoint{11.336818in}{0.721193in}}%
\pgfpathlineto{\pgfqpoint{11.336818in}{0.633600in}}%
\pgfpathlineto{\pgfqpoint{11.326461in}{0.633600in}}%
\pgfpathlineto{\pgfqpoint{11.326461in}{0.643056in}}%
\pgfpathquadraticcurveto{\pgfqpoint{11.323201in}{0.637437in}}{\pgfqpoint{11.318212in}{0.634699in}}%
\pgfpathquadraticcurveto{\pgfqpoint{11.313240in}{0.631961in}}{\pgfqpoint{11.306252in}{0.631961in}}%
\pgfpathquadraticcurveto{\pgfqpoint{11.294832in}{0.631961in}}{\pgfqpoint{11.287645in}{0.641075in}}%
\pgfpathquadraticcurveto{\pgfqpoint{11.280476in}{0.650207in}}{\pgfqpoint{11.280476in}{0.665067in}}%
\pgfpathquadraticcurveto{\pgfqpoint{11.280476in}{0.679927in}}{\pgfqpoint{11.287645in}{0.689041in}}%
\pgfpathquadraticcurveto{\pgfqpoint{11.294832in}{0.698156in}}{\pgfqpoint{11.306252in}{0.698156in}}%
\pgfpathquadraticcurveto{\pgfqpoint{11.313240in}{0.698156in}}{\pgfqpoint{11.318212in}{0.695418in}}%
\pgfpathquadraticcurveto{\pgfqpoint{11.323201in}{0.692698in}}{\pgfqpoint{11.326461in}{0.687078in}}%
\pgfpathlineto{\pgfqpoint{11.326461in}{0.687078in}}%
\pgfpathclose%
\pgfpathmoveto{\pgfqpoint{11.291176in}{0.665067in}}%
\pgfpathquadraticcurveto{\pgfqpoint{11.291176in}{0.653648in}}{\pgfqpoint{11.295877in}{0.647145in}}%
\pgfpathquadraticcurveto{\pgfqpoint{11.300578in}{0.640643in}}{\pgfqpoint{11.308791in}{0.640643in}}%
\pgfpathquadraticcurveto{\pgfqpoint{11.317005in}{0.640643in}}{\pgfqpoint{11.321724in}{0.647145in}}%
\pgfpathquadraticcurveto{\pgfqpoint{11.326461in}{0.653648in}}{\pgfqpoint{11.326461in}{0.665067in}}%
\pgfpathquadraticcurveto{\pgfqpoint{11.326461in}{0.676487in}}{\pgfqpoint{11.321724in}{0.682989in}}%
\pgfpathquadraticcurveto{\pgfqpoint{11.317005in}{0.689492in}}{\pgfqpoint{11.308791in}{0.689492in}}%
\pgfpathquadraticcurveto{\pgfqpoint{11.300578in}{0.689492in}}{\pgfqpoint{11.295877in}{0.682989in}}%
\pgfpathquadraticcurveto{\pgfqpoint{11.291176in}{0.676487in}}{\pgfqpoint{11.291176in}{0.665067in}}%
\pgfpathlineto{\pgfqpoint{11.291176in}{0.665067in}}%
\pgfpathclose%
\pgfpathmoveto{\pgfqpoint{11.382587in}{0.689384in}}%
\pgfpathquadraticcurveto{\pgfqpoint{11.374266in}{0.689384in}}{\pgfqpoint{11.369420in}{0.682881in}}%
\pgfpathquadraticcurveto{\pgfqpoint{11.364575in}{0.676379in}}{\pgfqpoint{11.364575in}{0.665067in}}%
\pgfpathquadraticcurveto{\pgfqpoint{11.364575in}{0.653756in}}{\pgfqpoint{11.369384in}{0.647253in}}%
\pgfpathquadraticcurveto{\pgfqpoint{11.374212in}{0.640751in}}{\pgfqpoint{11.382587in}{0.640751in}}%
\pgfpathquadraticcurveto{\pgfqpoint{11.390873in}{0.640751in}}{\pgfqpoint{11.395700in}{0.647271in}}%
\pgfpathquadraticcurveto{\pgfqpoint{11.400545in}{0.653810in}}{\pgfqpoint{11.400545in}{0.665067in}}%
\pgfpathquadraticcurveto{\pgfqpoint{11.400545in}{0.676271in}}{\pgfqpoint{11.395700in}{0.682827in}}%
\pgfpathquadraticcurveto{\pgfqpoint{11.390873in}{0.689384in}}{\pgfqpoint{11.382587in}{0.689384in}}%
\pgfpathlineto{\pgfqpoint{11.382587in}{0.689384in}}%
\pgfpathclose%
\pgfpathmoveto{\pgfqpoint{11.382587in}{0.698156in}}%
\pgfpathquadraticcurveto{\pgfqpoint{11.396096in}{0.698156in}}{\pgfqpoint{11.403806in}{0.689366in}}%
\pgfpathquadraticcurveto{\pgfqpoint{11.411533in}{0.680594in}}{\pgfqpoint{11.411533in}{0.665067in}}%
\pgfpathquadraticcurveto{\pgfqpoint{11.411533in}{0.649595in}}{\pgfqpoint{11.403806in}{0.640769in}}%
\pgfpathquadraticcurveto{\pgfqpoint{11.396096in}{0.631961in}}{\pgfqpoint{11.382587in}{0.631961in}}%
\pgfpathquadraticcurveto{\pgfqpoint{11.369024in}{0.631961in}}{\pgfqpoint{11.361333in}{0.640769in}}%
\pgfpathquadraticcurveto{\pgfqpoint{11.353660in}{0.649595in}}{\pgfqpoint{11.353660in}{0.665067in}}%
\pgfpathquadraticcurveto{\pgfqpoint{11.353660in}{0.680594in}}{\pgfqpoint{11.361333in}{0.689366in}}%
\pgfpathquadraticcurveto{\pgfqpoint{11.369024in}{0.698156in}}{\pgfqpoint{11.382587in}{0.698156in}}%
\pgfpathlineto{\pgfqpoint{11.382587in}{0.698156in}}%
\pgfpathclose%
\pgfpathmoveto{\pgfqpoint{11.482627in}{0.667715in}}%
\pgfpathlineto{\pgfqpoint{11.482627in}{0.662654in}}%
\pgfpathlineto{\pgfqpoint{11.435003in}{0.662654in}}%
\pgfpathquadraticcurveto{\pgfqpoint{11.435687in}{0.651954in}}{\pgfqpoint{11.441451in}{0.646353in}}%
\pgfpathquadraticcurveto{\pgfqpoint{11.447215in}{0.640751in}}{\pgfqpoint{11.457518in}{0.640751in}}%
\pgfpathquadraticcurveto{\pgfqpoint{11.463480in}{0.640751in}}{\pgfqpoint{11.469082in}{0.642210in}}%
\pgfpathquadraticcurveto{\pgfqpoint{11.474683in}{0.643669in}}{\pgfqpoint{11.480213in}{0.646605in}}%
\pgfpathlineto{\pgfqpoint{11.480213in}{0.636806in}}%
\pgfpathquadraticcurveto{\pgfqpoint{11.474629in}{0.634447in}}{\pgfqpoint{11.468775in}{0.633204in}}%
\pgfpathquadraticcurveto{\pgfqpoint{11.462921in}{0.631961in}}{\pgfqpoint{11.456905in}{0.631961in}}%
\pgfpathquadraticcurveto{\pgfqpoint{11.441811in}{0.631961in}}{\pgfqpoint{11.433003in}{0.640733in}}%
\pgfpathquadraticcurveto{\pgfqpoint{11.424195in}{0.649523in}}{\pgfqpoint{11.424195in}{0.664509in}}%
\pgfpathquadraticcurveto{\pgfqpoint{11.424195in}{0.679981in}}{\pgfqpoint{11.432553in}{0.689059in}}%
\pgfpathquadraticcurveto{\pgfqpoint{11.440911in}{0.698156in}}{\pgfqpoint{11.455104in}{0.698156in}}%
\pgfpathquadraticcurveto{\pgfqpoint{11.467821in}{0.698156in}}{\pgfqpoint{11.475224in}{0.689960in}}%
\pgfpathquadraticcurveto{\pgfqpoint{11.482627in}{0.681783in}}{\pgfqpoint{11.482627in}{0.667715in}}%
\pgfpathlineto{\pgfqpoint{11.482627in}{0.667715in}}%
\pgfpathclose%
\pgfpathmoveto{\pgfqpoint{11.472270in}{0.670759in}}%
\pgfpathquadraticcurveto{\pgfqpoint{11.472162in}{0.679243in}}{\pgfqpoint{11.467515in}{0.684304in}}%
\pgfpathquadraticcurveto{\pgfqpoint{11.462867in}{0.689384in}}{\pgfqpoint{11.455212in}{0.689384in}}%
\pgfpathquadraticcurveto{\pgfqpoint{11.446548in}{0.689384in}}{\pgfqpoint{11.441343in}{0.684484in}}%
\pgfpathquadraticcurveto{\pgfqpoint{11.436137in}{0.679585in}}{\pgfqpoint{11.435345in}{0.670687in}}%
\pgfpathlineto{\pgfqpoint{11.472270in}{0.670759in}}%
\pgfpathlineto{\pgfqpoint{11.472270in}{0.670759in}}%
\pgfpathclose%
\pgfpathmoveto{\pgfqpoint{11.539815in}{0.694787in}}%
\pgfpathlineto{\pgfqpoint{11.539815in}{0.684989in}}%
\pgfpathquadraticcurveto{\pgfqpoint{11.535438in}{0.687240in}}{\pgfqpoint{11.530701in}{0.688357in}}%
\pgfpathquadraticcurveto{\pgfqpoint{11.525982in}{0.689492in}}{\pgfqpoint{11.520903in}{0.689492in}}%
\pgfpathquadraticcurveto{\pgfqpoint{11.513193in}{0.689492in}}{\pgfqpoint{11.509339in}{0.687132in}}%
\pgfpathquadraticcurveto{\pgfqpoint{11.505484in}{0.684773in}}{\pgfqpoint{11.505484in}{0.680035in}}%
\pgfpathquadraticcurveto{\pgfqpoint{11.505484in}{0.676433in}}{\pgfqpoint{11.508240in}{0.674380in}}%
\pgfpathquadraticcurveto{\pgfqpoint{11.510996in}{0.672326in}}{\pgfqpoint{11.519336in}{0.670471in}}%
\pgfpathlineto{\pgfqpoint{11.522884in}{0.669678in}}%
\pgfpathquadraticcurveto{\pgfqpoint{11.533907in}{0.667319in}}{\pgfqpoint{11.538554in}{0.663014in}}%
\pgfpathquadraticcurveto{\pgfqpoint{11.543202in}{0.658709in}}{\pgfqpoint{11.543202in}{0.651000in}}%
\pgfpathquadraticcurveto{\pgfqpoint{11.543202in}{0.642210in}}{\pgfqpoint{11.536249in}{0.637076in}}%
\pgfpathquadraticcurveto{\pgfqpoint{11.529296in}{0.631961in}}{\pgfqpoint{11.517138in}{0.631961in}}%
\pgfpathquadraticcurveto{\pgfqpoint{11.512077in}{0.631961in}}{\pgfqpoint{11.506583in}{0.632952in}}%
\pgfpathquadraticcurveto{\pgfqpoint{11.501089in}{0.633942in}}{\pgfqpoint{11.495019in}{0.635906in}}%
\pgfpathlineto{\pgfqpoint{11.495019in}{0.646605in}}%
\pgfpathquadraticcurveto{\pgfqpoint{11.500765in}{0.643615in}}{\pgfqpoint{11.506331in}{0.642120in}}%
\pgfpathquadraticcurveto{\pgfqpoint{11.511896in}{0.640643in}}{\pgfqpoint{11.517372in}{0.640643in}}%
\pgfpathquadraticcurveto{\pgfqpoint{11.524685in}{0.640643in}}{\pgfqpoint{11.528612in}{0.643146in}}%
\pgfpathquadraticcurveto{\pgfqpoint{11.532556in}{0.645650in}}{\pgfqpoint{11.532556in}{0.650207in}}%
\pgfpathquadraticcurveto{\pgfqpoint{11.532556in}{0.654422in}}{\pgfqpoint{11.529711in}{0.656674in}}%
\pgfpathquadraticcurveto{\pgfqpoint{11.526883in}{0.658925in}}{\pgfqpoint{11.517246in}{0.661014in}}%
\pgfpathlineto{\pgfqpoint{11.513644in}{0.661861in}}%
\pgfpathquadraticcurveto{\pgfqpoint{11.504025in}{0.663878in}}{\pgfqpoint{11.499738in}{0.668075in}}%
\pgfpathquadraticcurveto{\pgfqpoint{11.495469in}{0.672272in}}{\pgfqpoint{11.495469in}{0.679585in}}%
\pgfpathquadraticcurveto{\pgfqpoint{11.495469in}{0.688483in}}{\pgfqpoint{11.501774in}{0.693310in}}%
\pgfpathquadraticcurveto{\pgfqpoint{11.508078in}{0.698156in}}{\pgfqpoint{11.519678in}{0.698156in}}%
\pgfpathquadraticcurveto{\pgfqpoint{11.525406in}{0.698156in}}{\pgfqpoint{11.530467in}{0.697309in}}%
\pgfpathquadraticcurveto{\pgfqpoint{11.535546in}{0.696480in}}{\pgfqpoint{11.539815in}{0.694787in}}%
\pgfpathlineto{\pgfqpoint{11.539815in}{0.694787in}}%
\pgfpathclose%
\pgfusepath{fill}%
\end{pgfscope}%
\begin{pgfscope}%
\pgfsetbuttcap%
\pgfsetroundjoin%
\definecolor{currentfill}{rgb}{0.309804,0.372549,0.419608}%
\pgfsetfillcolor{currentfill}%
\pgfsetlinewidth{0.000000pt}%
\definecolor{currentstroke}{rgb}{0.309804,0.372549,0.419608}%
\pgfsetstrokecolor{currentstroke}%
\pgfsetdash{}{0pt}%
\pgfpathmoveto{\pgfqpoint{11.674971in}{0.671660in}}%
\pgfpathlineto{\pgfqpoint{11.674971in}{0.633600in}}%
\pgfpathlineto{\pgfqpoint{11.664614in}{0.633600in}}%
\pgfpathlineto{\pgfqpoint{11.664614in}{0.671317in}}%
\pgfpathquadraticcurveto{\pgfqpoint{11.664614in}{0.680269in}}{\pgfqpoint{11.661120in}{0.684700in}}%
\pgfpathquadraticcurveto{\pgfqpoint{11.657625in}{0.689149in}}{\pgfqpoint{11.650655in}{0.689149in}}%
\pgfpathquadraticcurveto{\pgfqpoint{11.642261in}{0.689149in}}{\pgfqpoint{11.637416in}{0.683800in}}%
\pgfpathquadraticcurveto{\pgfqpoint{11.632571in}{0.678468in}}{\pgfqpoint{11.632571in}{0.669228in}}%
\pgfpathlineto{\pgfqpoint{11.632571in}{0.633600in}}%
\pgfpathlineto{\pgfqpoint{11.622160in}{0.633600in}}%
\pgfpathlineto{\pgfqpoint{11.622160in}{0.696643in}}%
\pgfpathlineto{\pgfqpoint{11.632571in}{0.696643in}}%
\pgfpathlineto{\pgfqpoint{11.632571in}{0.686844in}}%
\pgfpathquadraticcurveto{\pgfqpoint{11.636299in}{0.692536in}}{\pgfqpoint{11.641324in}{0.695346in}}%
\pgfpathquadraticcurveto{\pgfqpoint{11.646368in}{0.698156in}}{\pgfqpoint{11.652960in}{0.698156in}}%
\pgfpathquadraticcurveto{\pgfqpoint{11.663822in}{0.698156in}}{\pgfqpoint{11.669387in}{0.691437in}}%
\pgfpathquadraticcurveto{\pgfqpoint{11.674971in}{0.684718in}}{\pgfqpoint{11.674971in}{0.671660in}}%
\pgfpathlineto{\pgfqpoint{11.674971in}{0.671660in}}%
\pgfpathclose%
\pgfpathmoveto{\pgfqpoint{11.720038in}{0.689384in}}%
\pgfpathquadraticcurveto{\pgfqpoint{11.711716in}{0.689384in}}{\pgfqpoint{11.706871in}{0.682881in}}%
\pgfpathquadraticcurveto{\pgfqpoint{11.702025in}{0.676379in}}{\pgfqpoint{11.702025in}{0.665067in}}%
\pgfpathquadraticcurveto{\pgfqpoint{11.702025in}{0.653756in}}{\pgfqpoint{11.706835in}{0.647253in}}%
\pgfpathquadraticcurveto{\pgfqpoint{11.711662in}{0.640751in}}{\pgfqpoint{11.720038in}{0.640751in}}%
\pgfpathquadraticcurveto{\pgfqpoint{11.728323in}{0.640751in}}{\pgfqpoint{11.733150in}{0.647271in}}%
\pgfpathquadraticcurveto{\pgfqpoint{11.737996in}{0.653810in}}{\pgfqpoint{11.737996in}{0.665067in}}%
\pgfpathquadraticcurveto{\pgfqpoint{11.737996in}{0.676271in}}{\pgfqpoint{11.733150in}{0.682827in}}%
\pgfpathquadraticcurveto{\pgfqpoint{11.728323in}{0.689384in}}{\pgfqpoint{11.720038in}{0.689384in}}%
\pgfpathlineto{\pgfqpoint{11.720038in}{0.689384in}}%
\pgfpathclose%
\pgfpathmoveto{\pgfqpoint{11.720038in}{0.698156in}}%
\pgfpathquadraticcurveto{\pgfqpoint{11.733547in}{0.698156in}}{\pgfqpoint{11.741256in}{0.689366in}}%
\pgfpathquadraticcurveto{\pgfqpoint{11.748983in}{0.680594in}}{\pgfqpoint{11.748983in}{0.665067in}}%
\pgfpathquadraticcurveto{\pgfqpoint{11.748983in}{0.649595in}}{\pgfqpoint{11.741256in}{0.640769in}}%
\pgfpathquadraticcurveto{\pgfqpoint{11.733547in}{0.631961in}}{\pgfqpoint{11.720038in}{0.631961in}}%
\pgfpathquadraticcurveto{\pgfqpoint{11.706474in}{0.631961in}}{\pgfqpoint{11.698783in}{0.640769in}}%
\pgfpathquadraticcurveto{\pgfqpoint{11.691110in}{0.649595in}}{\pgfqpoint{11.691110in}{0.665067in}}%
\pgfpathquadraticcurveto{\pgfqpoint{11.691110in}{0.680594in}}{\pgfqpoint{11.698783in}{0.689366in}}%
\pgfpathquadraticcurveto{\pgfqpoint{11.706474in}{0.698156in}}{\pgfqpoint{11.720038in}{0.698156in}}%
\pgfpathlineto{\pgfqpoint{11.720038in}{0.698156in}}%
\pgfpathclose%
\pgfpathmoveto{\pgfqpoint{11.776398in}{0.714547in}}%
\pgfpathlineto{\pgfqpoint{11.776398in}{0.696643in}}%
\pgfpathlineto{\pgfqpoint{11.797724in}{0.696643in}}%
\pgfpathlineto{\pgfqpoint{11.797724in}{0.688591in}}%
\pgfpathlineto{\pgfqpoint{11.776398in}{0.688591in}}%
\pgfpathlineto{\pgfqpoint{11.776398in}{0.654368in}}%
\pgfpathquadraticcurveto{\pgfqpoint{11.776398in}{0.646659in}}{\pgfqpoint{11.778505in}{0.644461in}}%
\pgfpathquadraticcurveto{\pgfqpoint{11.780612in}{0.642264in}}{\pgfqpoint{11.787097in}{0.642264in}}%
\pgfpathlineto{\pgfqpoint{11.797724in}{0.642264in}}%
\pgfpathlineto{\pgfqpoint{11.797724in}{0.633600in}}%
\pgfpathlineto{\pgfqpoint{11.787097in}{0.633600in}}%
\pgfpathquadraticcurveto{\pgfqpoint{11.775101in}{0.633600in}}{\pgfqpoint{11.770544in}{0.638067in}}%
\pgfpathquadraticcurveto{\pgfqpoint{11.765987in}{0.642552in}}{\pgfqpoint{11.765987in}{0.654368in}}%
\pgfpathlineto{\pgfqpoint{11.765987in}{0.688591in}}%
\pgfpathlineto{\pgfqpoint{11.758385in}{0.688591in}}%
\pgfpathlineto{\pgfqpoint{11.758385in}{0.696643in}}%
\pgfpathlineto{\pgfqpoint{11.765987in}{0.696643in}}%
\pgfpathlineto{\pgfqpoint{11.765987in}{0.714547in}}%
\pgfpathlineto{\pgfqpoint{11.776398in}{0.714547in}}%
\pgfpathlineto{\pgfqpoint{11.776398in}{0.714547in}}%
\pgfpathclose%
\pgfusepath{fill}%
\end{pgfscope}%
\begin{pgfscope}%
\pgfsetbuttcap%
\pgfsetroundjoin%
\definecolor{currentfill}{rgb}{0.309804,0.372549,0.419608}%
\pgfsetfillcolor{currentfill}%
\pgfsetlinewidth{0.000000pt}%
\definecolor{currentstroke}{rgb}{0.309804,0.372549,0.419608}%
\pgfsetstrokecolor{currentstroke}%
\pgfsetdash{}{0pt}%
\pgfpathmoveto{\pgfqpoint{11.934061in}{0.667715in}}%
\pgfpathlineto{\pgfqpoint{11.934061in}{0.662654in}}%
\pgfpathlineto{\pgfqpoint{11.886436in}{0.662654in}}%
\pgfpathquadraticcurveto{\pgfqpoint{11.887121in}{0.651954in}}{\pgfqpoint{11.892885in}{0.646353in}}%
\pgfpathquadraticcurveto{\pgfqpoint{11.898649in}{0.640751in}}{\pgfqpoint{11.908952in}{0.640751in}}%
\pgfpathquadraticcurveto{\pgfqpoint{11.914914in}{0.640751in}}{\pgfqpoint{11.920515in}{0.642210in}}%
\pgfpathquadraticcurveto{\pgfqpoint{11.926117in}{0.643669in}}{\pgfqpoint{11.931647in}{0.646605in}}%
\pgfpathlineto{\pgfqpoint{11.931647in}{0.636806in}}%
\pgfpathquadraticcurveto{\pgfqpoint{11.926063in}{0.634447in}}{\pgfqpoint{11.920209in}{0.633204in}}%
\pgfpathquadraticcurveto{\pgfqpoint{11.914355in}{0.631961in}}{\pgfqpoint{11.908339in}{0.631961in}}%
\pgfpathquadraticcurveto{\pgfqpoint{11.893245in}{0.631961in}}{\pgfqpoint{11.884437in}{0.640733in}}%
\pgfpathquadraticcurveto{\pgfqpoint{11.875629in}{0.649523in}}{\pgfqpoint{11.875629in}{0.664509in}}%
\pgfpathquadraticcurveto{\pgfqpoint{11.875629in}{0.679981in}}{\pgfqpoint{11.883987in}{0.689059in}}%
\pgfpathquadraticcurveto{\pgfqpoint{11.892344in}{0.698156in}}{\pgfqpoint{11.906538in}{0.698156in}}%
\pgfpathquadraticcurveto{\pgfqpoint{11.919255in}{0.698156in}}{\pgfqpoint{11.926658in}{0.689960in}}%
\pgfpathquadraticcurveto{\pgfqpoint{11.934061in}{0.681783in}}{\pgfqpoint{11.934061in}{0.667715in}}%
\pgfpathlineto{\pgfqpoint{11.934061in}{0.667715in}}%
\pgfpathclose%
\pgfpathmoveto{\pgfqpoint{11.923704in}{0.670759in}}%
\pgfpathquadraticcurveto{\pgfqpoint{11.923595in}{0.679243in}}{\pgfqpoint{11.918948in}{0.684304in}}%
\pgfpathquadraticcurveto{\pgfqpoint{11.914301in}{0.689384in}}{\pgfqpoint{11.906646in}{0.689384in}}%
\pgfpathquadraticcurveto{\pgfqpoint{11.897982in}{0.689384in}}{\pgfqpoint{11.892777in}{0.684484in}}%
\pgfpathquadraticcurveto{\pgfqpoint{11.887571in}{0.679585in}}{\pgfqpoint{11.886779in}{0.670687in}}%
\pgfpathlineto{\pgfqpoint{11.923704in}{0.670759in}}%
\pgfpathlineto{\pgfqpoint{11.923704in}{0.670759in}}%
\pgfpathclose%
\pgfpathmoveto{\pgfqpoint{11.957260in}{0.665067in}}%
\pgfpathquadraticcurveto{\pgfqpoint{11.957260in}{0.653648in}}{\pgfqpoint{11.961961in}{0.647145in}}%
\pgfpathquadraticcurveto{\pgfqpoint{11.966663in}{0.640643in}}{\pgfqpoint{11.974876in}{0.640643in}}%
\pgfpathquadraticcurveto{\pgfqpoint{11.983090in}{0.640643in}}{\pgfqpoint{11.987809in}{0.647145in}}%
\pgfpathquadraticcurveto{\pgfqpoint{11.992546in}{0.653648in}}{\pgfqpoint{11.992546in}{0.665067in}}%
\pgfpathquadraticcurveto{\pgfqpoint{11.992546in}{0.676487in}}{\pgfqpoint{11.987809in}{0.682989in}}%
\pgfpathquadraticcurveto{\pgfqpoint{11.983090in}{0.689492in}}{\pgfqpoint{11.974876in}{0.689492in}}%
\pgfpathquadraticcurveto{\pgfqpoint{11.966663in}{0.689492in}}{\pgfqpoint{11.961961in}{0.682989in}}%
\pgfpathquadraticcurveto{\pgfqpoint{11.957260in}{0.676487in}}{\pgfqpoint{11.957260in}{0.665067in}}%
\pgfpathlineto{\pgfqpoint{11.957260in}{0.665067in}}%
\pgfpathclose%
\pgfpathmoveto{\pgfqpoint{11.992546in}{0.643056in}}%
\pgfpathquadraticcurveto{\pgfqpoint{11.989286in}{0.637437in}}{\pgfqpoint{11.984296in}{0.634699in}}%
\pgfpathquadraticcurveto{\pgfqpoint{11.979325in}{0.631961in}}{\pgfqpoint{11.972336in}{0.631961in}}%
\pgfpathquadraticcurveto{\pgfqpoint{11.960917in}{0.631961in}}{\pgfqpoint{11.953730in}{0.641075in}}%
\pgfpathquadraticcurveto{\pgfqpoint{11.946561in}{0.650207in}}{\pgfqpoint{11.946561in}{0.665067in}}%
\pgfpathquadraticcurveto{\pgfqpoint{11.946561in}{0.679927in}}{\pgfqpoint{11.953730in}{0.689041in}}%
\pgfpathquadraticcurveto{\pgfqpoint{11.960917in}{0.698156in}}{\pgfqpoint{11.972336in}{0.698156in}}%
\pgfpathquadraticcurveto{\pgfqpoint{11.979325in}{0.698156in}}{\pgfqpoint{11.984296in}{0.695418in}}%
\pgfpathquadraticcurveto{\pgfqpoint{11.989286in}{0.692698in}}{\pgfqpoint{11.992546in}{0.687078in}}%
\pgfpathlineto{\pgfqpoint{11.992546in}{0.696643in}}%
\pgfpathlineto{\pgfqpoint{12.002903in}{0.696643in}}%
\pgfpathlineto{\pgfqpoint{12.002903in}{0.609626in}}%
\pgfpathlineto{\pgfqpoint{11.992546in}{0.609626in}}%
\pgfpathlineto{\pgfqpoint{11.992546in}{0.643056in}}%
\pgfpathlineto{\pgfqpoint{11.992546in}{0.643056in}}%
\pgfpathclose%
\pgfpathmoveto{\pgfqpoint{12.023185in}{0.658475in}}%
\pgfpathlineto{\pgfqpoint{12.023185in}{0.696643in}}%
\pgfpathlineto{\pgfqpoint{12.033542in}{0.696643in}}%
\pgfpathlineto{\pgfqpoint{12.033542in}{0.658871in}}%
\pgfpathquadraticcurveto{\pgfqpoint{12.033542in}{0.649919in}}{\pgfqpoint{12.037018in}{0.645434in}}%
\pgfpathquadraticcurveto{\pgfqpoint{12.040512in}{0.640967in}}{\pgfqpoint{12.047501in}{0.640967in}}%
\pgfpathquadraticcurveto{\pgfqpoint{12.055877in}{0.640967in}}{\pgfqpoint{12.060740in}{0.646317in}}%
\pgfpathquadraticcurveto{\pgfqpoint{12.065621in}{0.651666in}}{\pgfqpoint{12.065621in}{0.660906in}}%
\pgfpathlineto{\pgfqpoint{12.065621in}{0.696643in}}%
\pgfpathlineto{\pgfqpoint{12.075978in}{0.696643in}}%
\pgfpathlineto{\pgfqpoint{12.075978in}{0.633600in}}%
\pgfpathlineto{\pgfqpoint{12.065621in}{0.633600in}}%
\pgfpathlineto{\pgfqpoint{12.065621in}{0.643291in}}%
\pgfpathquadraticcurveto{\pgfqpoint{12.061857in}{0.637545in}}{\pgfqpoint{12.056867in}{0.634753in}}%
\pgfpathquadraticcurveto{\pgfqpoint{12.051896in}{0.631961in}}{\pgfqpoint{12.045304in}{0.631961in}}%
\pgfpathquadraticcurveto{\pgfqpoint{12.034442in}{0.631961in}}{\pgfqpoint{12.028804in}{0.638715in}}%
\pgfpathquadraticcurveto{\pgfqpoint{12.023185in}{0.645470in}}{\pgfqpoint{12.023185in}{0.658475in}}%
\pgfpathlineto{\pgfqpoint{12.023185in}{0.658475in}}%
\pgfpathclose%
\pgfpathmoveto{\pgfqpoint{12.049248in}{0.698156in}}%
\pgfpathlineto{\pgfqpoint{12.049248in}{0.698156in}}%
\pgfpathlineto{\pgfqpoint{12.049248in}{0.698156in}}%
\pgfpathclose%
\pgfpathmoveto{\pgfqpoint{12.125962in}{0.665283in}}%
\pgfpathquadraticcurveto{\pgfqpoint{12.113408in}{0.665283in}}{\pgfqpoint{12.108562in}{0.662419in}}%
\pgfpathquadraticcurveto{\pgfqpoint{12.103717in}{0.659556in}}{\pgfqpoint{12.103717in}{0.652621in}}%
\pgfpathquadraticcurveto{\pgfqpoint{12.103717in}{0.647109in}}{\pgfqpoint{12.107355in}{0.643867in}}%
\pgfpathquadraticcurveto{\pgfqpoint{12.110994in}{0.640643in}}{\pgfqpoint{12.117226in}{0.640643in}}%
\pgfpathquadraticcurveto{\pgfqpoint{12.125854in}{0.640643in}}{\pgfqpoint{12.131059in}{0.646749in}}%
\pgfpathquadraticcurveto{\pgfqpoint{12.136265in}{0.652855in}}{\pgfqpoint{12.136265in}{0.662978in}}%
\pgfpathlineto{\pgfqpoint{12.136265in}{0.665283in}}%
\pgfpathlineto{\pgfqpoint{12.125962in}{0.665283in}}%
\pgfpathlineto{\pgfqpoint{12.125962in}{0.665283in}}%
\pgfpathclose%
\pgfpathmoveto{\pgfqpoint{12.146622in}{0.669570in}}%
\pgfpathlineto{\pgfqpoint{12.146622in}{0.633600in}}%
\pgfpathlineto{\pgfqpoint{12.136265in}{0.633600in}}%
\pgfpathlineto{\pgfqpoint{12.136265in}{0.643164in}}%
\pgfpathquadraticcurveto{\pgfqpoint{12.132717in}{0.637437in}}{\pgfqpoint{12.127421in}{0.634699in}}%
\pgfpathquadraticcurveto{\pgfqpoint{12.122125in}{0.631961in}}{\pgfqpoint{12.114470in}{0.631961in}}%
\pgfpathquadraticcurveto{\pgfqpoint{12.104798in}{0.631961in}}{\pgfqpoint{12.099070in}{0.637401in}}%
\pgfpathquadraticcurveto{\pgfqpoint{12.093360in}{0.642840in}}{\pgfqpoint{12.093360in}{0.651954in}}%
\pgfpathquadraticcurveto{\pgfqpoint{12.093360in}{0.662582in}}{\pgfqpoint{12.100475in}{0.667985in}}%
\pgfpathquadraticcurveto{\pgfqpoint{12.107608in}{0.673389in}}{\pgfqpoint{12.121729in}{0.673389in}}%
\pgfpathlineto{\pgfqpoint{12.136265in}{0.673389in}}%
\pgfpathlineto{\pgfqpoint{12.136265in}{0.674416in}}%
\pgfpathquadraticcurveto{\pgfqpoint{12.136265in}{0.681566in}}{\pgfqpoint{12.131564in}{0.685475in}}%
\pgfpathquadraticcurveto{\pgfqpoint{12.126863in}{0.689384in}}{\pgfqpoint{12.118361in}{0.689384in}}%
\pgfpathquadraticcurveto{\pgfqpoint{12.112957in}{0.689384in}}{\pgfqpoint{12.107824in}{0.688087in}}%
\pgfpathquadraticcurveto{\pgfqpoint{12.102708in}{0.686790in}}{\pgfqpoint{12.097989in}{0.684196in}}%
\pgfpathlineto{\pgfqpoint{12.097989in}{0.693779in}}%
\pgfpathquadraticcurveto{\pgfqpoint{12.103663in}{0.695976in}}{\pgfqpoint{12.109013in}{0.697057in}}%
\pgfpathquadraticcurveto{\pgfqpoint{12.114362in}{0.698156in}}{\pgfqpoint{12.119424in}{0.698156in}}%
\pgfpathquadraticcurveto{\pgfqpoint{12.133113in}{0.698156in}}{\pgfqpoint{12.139867in}{0.691059in}}%
\pgfpathquadraticcurveto{\pgfqpoint{12.146622in}{0.683980in}}{\pgfqpoint{12.146622in}{0.669570in}}%
\pgfpathlineto{\pgfqpoint{12.146622in}{0.669570in}}%
\pgfpathclose%
\pgfpathmoveto{\pgfqpoint{12.178197in}{0.714547in}}%
\pgfpathlineto{\pgfqpoint{12.178197in}{0.696643in}}%
\pgfpathlineto{\pgfqpoint{12.199524in}{0.696643in}}%
\pgfpathlineto{\pgfqpoint{12.199524in}{0.688591in}}%
\pgfpathlineto{\pgfqpoint{12.178197in}{0.688591in}}%
\pgfpathlineto{\pgfqpoint{12.178197in}{0.654368in}}%
\pgfpathquadraticcurveto{\pgfqpoint{12.178197in}{0.646659in}}{\pgfqpoint{12.180305in}{0.644461in}}%
\pgfpathquadraticcurveto{\pgfqpoint{12.182412in}{0.642264in}}{\pgfqpoint{12.188896in}{0.642264in}}%
\pgfpathlineto{\pgfqpoint{12.199524in}{0.642264in}}%
\pgfpathlineto{\pgfqpoint{12.199524in}{0.633600in}}%
\pgfpathlineto{\pgfqpoint{12.188896in}{0.633600in}}%
\pgfpathquadraticcurveto{\pgfqpoint{12.176900in}{0.633600in}}{\pgfqpoint{12.172343in}{0.638067in}}%
\pgfpathquadraticcurveto{\pgfqpoint{12.167786in}{0.642552in}}{\pgfqpoint{12.167786in}{0.654368in}}%
\pgfpathlineto{\pgfqpoint{12.167786in}{0.688591in}}%
\pgfpathlineto{\pgfqpoint{12.160185in}{0.688591in}}%
\pgfpathlineto{\pgfqpoint{12.160185in}{0.696643in}}%
\pgfpathlineto{\pgfqpoint{12.167786in}{0.696643in}}%
\pgfpathlineto{\pgfqpoint{12.167786in}{0.714547in}}%
\pgfpathlineto{\pgfqpoint{12.178197in}{0.714547in}}%
\pgfpathlineto{\pgfqpoint{12.178197in}{0.714547in}}%
\pgfpathclose%
\pgfpathmoveto{\pgfqpoint{12.267069in}{0.667715in}}%
\pgfpathlineto{\pgfqpoint{12.267069in}{0.662654in}}%
\pgfpathlineto{\pgfqpoint{12.219445in}{0.662654in}}%
\pgfpathquadraticcurveto{\pgfqpoint{12.220130in}{0.651954in}}{\pgfqpoint{12.225893in}{0.646353in}}%
\pgfpathquadraticcurveto{\pgfqpoint{12.231657in}{0.640751in}}{\pgfqpoint{12.241960in}{0.640751in}}%
\pgfpathquadraticcurveto{\pgfqpoint{12.247922in}{0.640751in}}{\pgfqpoint{12.253524in}{0.642210in}}%
\pgfpathquadraticcurveto{\pgfqpoint{12.259126in}{0.643669in}}{\pgfqpoint{12.264656in}{0.646605in}}%
\pgfpathlineto{\pgfqpoint{12.264656in}{0.636806in}}%
\pgfpathquadraticcurveto{\pgfqpoint{12.259072in}{0.634447in}}{\pgfqpoint{12.253218in}{0.633204in}}%
\pgfpathquadraticcurveto{\pgfqpoint{12.247364in}{0.631961in}}{\pgfqpoint{12.241348in}{0.631961in}}%
\pgfpathquadraticcurveto{\pgfqpoint{12.226254in}{0.631961in}}{\pgfqpoint{12.217446in}{0.640733in}}%
\pgfpathquadraticcurveto{\pgfqpoint{12.208638in}{0.649523in}}{\pgfqpoint{12.208638in}{0.664509in}}%
\pgfpathquadraticcurveto{\pgfqpoint{12.208638in}{0.679981in}}{\pgfqpoint{12.216995in}{0.689059in}}%
\pgfpathquadraticcurveto{\pgfqpoint{12.225353in}{0.698156in}}{\pgfqpoint{12.239547in}{0.698156in}}%
\pgfpathquadraticcurveto{\pgfqpoint{12.252263in}{0.698156in}}{\pgfqpoint{12.259666in}{0.689960in}}%
\pgfpathquadraticcurveto{\pgfqpoint{12.267069in}{0.681783in}}{\pgfqpoint{12.267069in}{0.667715in}}%
\pgfpathlineto{\pgfqpoint{12.267069in}{0.667715in}}%
\pgfpathclose%
\pgfpathmoveto{\pgfqpoint{12.256712in}{0.670759in}}%
\pgfpathquadraticcurveto{\pgfqpoint{12.256604in}{0.679243in}}{\pgfqpoint{12.251957in}{0.684304in}}%
\pgfpathquadraticcurveto{\pgfqpoint{12.247310in}{0.689384in}}{\pgfqpoint{12.239655in}{0.689384in}}%
\pgfpathquadraticcurveto{\pgfqpoint{12.230991in}{0.689384in}}{\pgfqpoint{12.225785in}{0.684484in}}%
\pgfpathquadraticcurveto{\pgfqpoint{12.220580in}{0.679585in}}{\pgfqpoint{12.219787in}{0.670687in}}%
\pgfpathlineto{\pgfqpoint{12.256712in}{0.670759in}}%
\pgfpathlineto{\pgfqpoint{12.256712in}{0.670759in}}%
\pgfpathclose%
\pgfusepath{fill}%
\end{pgfscope}%
\begin{pgfscope}%
\pgfsetbuttcap%
\pgfsetroundjoin%
\definecolor{currentfill}{rgb}{0.309804,0.372549,0.419608}%
\pgfsetfillcolor{currentfill}%
\pgfsetlinewidth{0.000000pt}%
\definecolor{currentstroke}{rgb}{0.309804,0.372549,0.419608}%
\pgfsetstrokecolor{currentstroke}%
\pgfsetdash{}{0pt}%
\pgfpathmoveto{\pgfqpoint{12.411637in}{0.667715in}}%
\pgfpathlineto{\pgfqpoint{12.411637in}{0.662654in}}%
\pgfpathlineto{\pgfqpoint{12.364013in}{0.662654in}}%
\pgfpathquadraticcurveto{\pgfqpoint{12.364697in}{0.651954in}}{\pgfqpoint{12.370461in}{0.646353in}}%
\pgfpathquadraticcurveto{\pgfqpoint{12.376225in}{0.640751in}}{\pgfqpoint{12.386528in}{0.640751in}}%
\pgfpathquadraticcurveto{\pgfqpoint{12.392490in}{0.640751in}}{\pgfqpoint{12.398092in}{0.642210in}}%
\pgfpathquadraticcurveto{\pgfqpoint{12.403694in}{0.643669in}}{\pgfqpoint{12.409223in}{0.646605in}}%
\pgfpathlineto{\pgfqpoint{12.409223in}{0.636806in}}%
\pgfpathquadraticcurveto{\pgfqpoint{12.403640in}{0.634447in}}{\pgfqpoint{12.397786in}{0.633204in}}%
\pgfpathquadraticcurveto{\pgfqpoint{12.391932in}{0.631961in}}{\pgfqpoint{12.385916in}{0.631961in}}%
\pgfpathquadraticcurveto{\pgfqpoint{12.370821in}{0.631961in}}{\pgfqpoint{12.362013in}{0.640733in}}%
\pgfpathquadraticcurveto{\pgfqpoint{12.353206in}{0.649523in}}{\pgfqpoint{12.353206in}{0.664509in}}%
\pgfpathquadraticcurveto{\pgfqpoint{12.353206in}{0.679981in}}{\pgfqpoint{12.361563in}{0.689059in}}%
\pgfpathquadraticcurveto{\pgfqpoint{12.369921in}{0.698156in}}{\pgfqpoint{12.384114in}{0.698156in}}%
\pgfpathquadraticcurveto{\pgfqpoint{12.396831in}{0.698156in}}{\pgfqpoint{12.404234in}{0.689960in}}%
\pgfpathquadraticcurveto{\pgfqpoint{12.411637in}{0.681783in}}{\pgfqpoint{12.411637in}{0.667715in}}%
\pgfpathlineto{\pgfqpoint{12.411637in}{0.667715in}}%
\pgfpathclose%
\pgfpathmoveto{\pgfqpoint{12.401280in}{0.670759in}}%
\pgfpathquadraticcurveto{\pgfqpoint{12.401172in}{0.679243in}}{\pgfqpoint{12.396525in}{0.684304in}}%
\pgfpathquadraticcurveto{\pgfqpoint{12.391878in}{0.689384in}}{\pgfqpoint{12.384222in}{0.689384in}}%
\pgfpathquadraticcurveto{\pgfqpoint{12.375559in}{0.689384in}}{\pgfqpoint{12.370353in}{0.684484in}}%
\pgfpathquadraticcurveto{\pgfqpoint{12.365148in}{0.679585in}}{\pgfqpoint{12.364355in}{0.670687in}}%
\pgfpathlineto{\pgfqpoint{12.401280in}{0.670759in}}%
\pgfpathlineto{\pgfqpoint{12.401280in}{0.670759in}}%
\pgfpathclose%
\pgfpathmoveto{\pgfqpoint{12.468826in}{0.694787in}}%
\pgfpathlineto{\pgfqpoint{12.468826in}{0.684989in}}%
\pgfpathquadraticcurveto{\pgfqpoint{12.464449in}{0.687240in}}{\pgfqpoint{12.459711in}{0.688357in}}%
\pgfpathquadraticcurveto{\pgfqpoint{12.454992in}{0.689492in}}{\pgfqpoint{12.449913in}{0.689492in}}%
\pgfpathquadraticcurveto{\pgfqpoint{12.442204in}{0.689492in}}{\pgfqpoint{12.438349in}{0.687132in}}%
\pgfpathquadraticcurveto{\pgfqpoint{12.434494in}{0.684773in}}{\pgfqpoint{12.434494in}{0.680035in}}%
\pgfpathquadraticcurveto{\pgfqpoint{12.434494in}{0.676433in}}{\pgfqpoint{12.437250in}{0.674380in}}%
\pgfpathquadraticcurveto{\pgfqpoint{12.440006in}{0.672326in}}{\pgfqpoint{12.448346in}{0.670471in}}%
\pgfpathlineto{\pgfqpoint{12.451894in}{0.669678in}}%
\pgfpathquadraticcurveto{\pgfqpoint{12.462918in}{0.667319in}}{\pgfqpoint{12.467565in}{0.663014in}}%
\pgfpathquadraticcurveto{\pgfqpoint{12.472212in}{0.658709in}}{\pgfqpoint{12.472212in}{0.651000in}}%
\pgfpathquadraticcurveto{\pgfqpoint{12.472212in}{0.642210in}}{\pgfqpoint{12.465259in}{0.637076in}}%
\pgfpathquadraticcurveto{\pgfqpoint{12.458306in}{0.631961in}}{\pgfqpoint{12.446148in}{0.631961in}}%
\pgfpathquadraticcurveto{\pgfqpoint{12.441087in}{0.631961in}}{\pgfqpoint{12.435593in}{0.632952in}}%
\pgfpathquadraticcurveto{\pgfqpoint{12.430099in}{0.633942in}}{\pgfqpoint{12.424029in}{0.635906in}}%
\pgfpathlineto{\pgfqpoint{12.424029in}{0.646605in}}%
\pgfpathquadraticcurveto{\pgfqpoint{12.429775in}{0.643615in}}{\pgfqpoint{12.435341in}{0.642120in}}%
\pgfpathquadraticcurveto{\pgfqpoint{12.440907in}{0.640643in}}{\pgfqpoint{12.446382in}{0.640643in}}%
\pgfpathquadraticcurveto{\pgfqpoint{12.453695in}{0.640643in}}{\pgfqpoint{12.457622in}{0.643146in}}%
\pgfpathquadraticcurveto{\pgfqpoint{12.461567in}{0.645650in}}{\pgfqpoint{12.461567in}{0.650207in}}%
\pgfpathquadraticcurveto{\pgfqpoint{12.461567in}{0.654422in}}{\pgfqpoint{12.458721in}{0.656674in}}%
\pgfpathquadraticcurveto{\pgfqpoint{12.455893in}{0.658925in}}{\pgfqpoint{12.446256in}{0.661014in}}%
\pgfpathlineto{\pgfqpoint{12.442654in}{0.661861in}}%
\pgfpathquadraticcurveto{\pgfqpoint{12.433035in}{0.663878in}}{\pgfqpoint{12.428748in}{0.668075in}}%
\pgfpathquadraticcurveto{\pgfqpoint{12.424480in}{0.672272in}}{\pgfqpoint{12.424480in}{0.679585in}}%
\pgfpathquadraticcurveto{\pgfqpoint{12.424480in}{0.688483in}}{\pgfqpoint{12.430784in}{0.693310in}}%
\pgfpathquadraticcurveto{\pgfqpoint{12.437088in}{0.698156in}}{\pgfqpoint{12.448688in}{0.698156in}}%
\pgfpathquadraticcurveto{\pgfqpoint{12.454416in}{0.698156in}}{\pgfqpoint{12.459477in}{0.697309in}}%
\pgfpathquadraticcurveto{\pgfqpoint{12.464557in}{0.696480in}}{\pgfqpoint{12.468826in}{0.694787in}}%
\pgfpathlineto{\pgfqpoint{12.468826in}{0.694787in}}%
\pgfpathclose%
\pgfpathmoveto{\pgfqpoint{12.498942in}{0.714547in}}%
\pgfpathlineto{\pgfqpoint{12.498942in}{0.696643in}}%
\pgfpathlineto{\pgfqpoint{12.520268in}{0.696643in}}%
\pgfpathlineto{\pgfqpoint{12.520268in}{0.688591in}}%
\pgfpathlineto{\pgfqpoint{12.498942in}{0.688591in}}%
\pgfpathlineto{\pgfqpoint{12.498942in}{0.654368in}}%
\pgfpathquadraticcurveto{\pgfqpoint{12.498942in}{0.646659in}}{\pgfqpoint{12.501049in}{0.644461in}}%
\pgfpathquadraticcurveto{\pgfqpoint{12.503157in}{0.642264in}}{\pgfqpoint{12.509641in}{0.642264in}}%
\pgfpathlineto{\pgfqpoint{12.520268in}{0.642264in}}%
\pgfpathlineto{\pgfqpoint{12.520268in}{0.633600in}}%
\pgfpathlineto{\pgfqpoint{12.509641in}{0.633600in}}%
\pgfpathquadraticcurveto{\pgfqpoint{12.497645in}{0.633600in}}{\pgfqpoint{12.493088in}{0.638067in}}%
\pgfpathquadraticcurveto{\pgfqpoint{12.488531in}{0.642552in}}{\pgfqpoint{12.488531in}{0.654368in}}%
\pgfpathlineto{\pgfqpoint{12.488531in}{0.688591in}}%
\pgfpathlineto{\pgfqpoint{12.480930in}{0.688591in}}%
\pgfpathlineto{\pgfqpoint{12.480930in}{0.696643in}}%
\pgfpathlineto{\pgfqpoint{12.488531in}{0.696643in}}%
\pgfpathlineto{\pgfqpoint{12.488531in}{0.714547in}}%
\pgfpathlineto{\pgfqpoint{12.498942in}{0.714547in}}%
\pgfpathlineto{\pgfqpoint{12.498942in}{0.714547in}}%
\pgfpathclose%
\pgfpathmoveto{\pgfqpoint{12.533885in}{0.696643in}}%
\pgfpathlineto{\pgfqpoint{12.544242in}{0.696643in}}%
\pgfpathlineto{\pgfqpoint{12.544242in}{0.633600in}}%
\pgfpathlineto{\pgfqpoint{12.533885in}{0.633600in}}%
\pgfpathlineto{\pgfqpoint{12.533885in}{0.696643in}}%
\pgfpathlineto{\pgfqpoint{12.533885in}{0.696643in}}%
\pgfpathclose%
\pgfpathmoveto{\pgfqpoint{12.533885in}{0.721193in}}%
\pgfpathlineto{\pgfqpoint{12.544242in}{0.721193in}}%
\pgfpathlineto{\pgfqpoint{12.544242in}{0.708062in}}%
\pgfpathlineto{\pgfqpoint{12.533885in}{0.708062in}}%
\pgfpathlineto{\pgfqpoint{12.533885in}{0.721193in}}%
\pgfpathlineto{\pgfqpoint{12.533885in}{0.721193in}}%
\pgfpathclose%
\pgfpathmoveto{\pgfqpoint{12.614994in}{0.684538in}}%
\pgfpathquadraticcurveto{\pgfqpoint{12.618885in}{0.691527in}}{\pgfqpoint{12.624288in}{0.694841in}}%
\pgfpathquadraticcurveto{\pgfqpoint{12.629692in}{0.698156in}}{\pgfqpoint{12.637005in}{0.698156in}}%
\pgfpathquadraticcurveto{\pgfqpoint{12.646858in}{0.698156in}}{\pgfqpoint{12.652207in}{0.691257in}}%
\pgfpathquadraticcurveto{\pgfqpoint{12.657557in}{0.684376in}}{\pgfqpoint{12.657557in}{0.671660in}}%
\pgfpathlineto{\pgfqpoint{12.657557in}{0.633600in}}%
\pgfpathlineto{\pgfqpoint{12.647146in}{0.633600in}}%
\pgfpathlineto{\pgfqpoint{12.647146in}{0.671317in}}%
\pgfpathquadraticcurveto{\pgfqpoint{12.647146in}{0.680378in}}{\pgfqpoint{12.643922in}{0.684755in}}%
\pgfpathquadraticcurveto{\pgfqpoint{12.640715in}{0.689149in}}{\pgfqpoint{12.634141in}{0.689149in}}%
\pgfpathquadraticcurveto{\pgfqpoint{12.626090in}{0.689149in}}{\pgfqpoint{12.621406in}{0.683800in}}%
\pgfpathquadraticcurveto{\pgfqpoint{12.616741in}{0.678468in}}{\pgfqpoint{12.616741in}{0.669228in}}%
\pgfpathlineto{\pgfqpoint{12.616741in}{0.633600in}}%
\pgfpathlineto{\pgfqpoint{12.606330in}{0.633600in}}%
\pgfpathlineto{\pgfqpoint{12.606330in}{0.671317in}}%
\pgfpathquadraticcurveto{\pgfqpoint{12.606330in}{0.680432in}}{\pgfqpoint{12.603124in}{0.684791in}}%
\pgfpathquadraticcurveto{\pgfqpoint{12.599918in}{0.689149in}}{\pgfqpoint{12.593217in}{0.689149in}}%
\pgfpathquadraticcurveto{\pgfqpoint{12.585274in}{0.689149in}}{\pgfqpoint{12.580591in}{0.683782in}}%
\pgfpathquadraticcurveto{\pgfqpoint{12.575926in}{0.678414in}}{\pgfqpoint{12.575926in}{0.669228in}}%
\pgfpathlineto{\pgfqpoint{12.575926in}{0.633600in}}%
\pgfpathlineto{\pgfqpoint{12.565515in}{0.633600in}}%
\pgfpathlineto{\pgfqpoint{12.565515in}{0.696643in}}%
\pgfpathlineto{\pgfqpoint{12.575926in}{0.696643in}}%
\pgfpathlineto{\pgfqpoint{12.575926in}{0.686844in}}%
\pgfpathquadraticcurveto{\pgfqpoint{12.579474in}{0.692644in}}{\pgfqpoint{12.584428in}{0.695400in}}%
\pgfpathquadraticcurveto{\pgfqpoint{12.589381in}{0.698156in}}{\pgfqpoint{12.596189in}{0.698156in}}%
\pgfpathquadraticcurveto{\pgfqpoint{12.603070in}{0.698156in}}{\pgfqpoint{12.607879in}{0.694661in}}%
\pgfpathquadraticcurveto{\pgfqpoint{12.612689in}{0.691185in}}{\pgfqpoint{12.614994in}{0.684538in}}%
\pgfpathlineto{\pgfqpoint{12.614994in}{0.684538in}}%
\pgfpathclose%
\pgfpathmoveto{\pgfqpoint{12.706856in}{0.665283in}}%
\pgfpathquadraticcurveto{\pgfqpoint{12.694302in}{0.665283in}}{\pgfqpoint{12.689456in}{0.662419in}}%
\pgfpathquadraticcurveto{\pgfqpoint{12.684611in}{0.659556in}}{\pgfqpoint{12.684611in}{0.652621in}}%
\pgfpathquadraticcurveto{\pgfqpoint{12.684611in}{0.647109in}}{\pgfqpoint{12.688250in}{0.643867in}}%
\pgfpathquadraticcurveto{\pgfqpoint{12.691888in}{0.640643in}}{\pgfqpoint{12.698120in}{0.640643in}}%
\pgfpathquadraticcurveto{\pgfqpoint{12.706748in}{0.640643in}}{\pgfqpoint{12.711954in}{0.646749in}}%
\pgfpathquadraticcurveto{\pgfqpoint{12.717159in}{0.652855in}}{\pgfqpoint{12.717159in}{0.662978in}}%
\pgfpathlineto{\pgfqpoint{12.717159in}{0.665283in}}%
\pgfpathlineto{\pgfqpoint{12.706856in}{0.665283in}}%
\pgfpathlineto{\pgfqpoint{12.706856in}{0.665283in}}%
\pgfpathclose%
\pgfpathmoveto{\pgfqpoint{12.727516in}{0.669570in}}%
\pgfpathlineto{\pgfqpoint{12.727516in}{0.633600in}}%
\pgfpathlineto{\pgfqpoint{12.717159in}{0.633600in}}%
\pgfpathlineto{\pgfqpoint{12.717159in}{0.643164in}}%
\pgfpathquadraticcurveto{\pgfqpoint{12.713611in}{0.637437in}}{\pgfqpoint{12.708315in}{0.634699in}}%
\pgfpathquadraticcurveto{\pgfqpoint{12.703020in}{0.631961in}}{\pgfqpoint{12.695364in}{0.631961in}}%
\pgfpathquadraticcurveto{\pgfqpoint{12.685692in}{0.631961in}}{\pgfqpoint{12.679964in}{0.637401in}}%
\pgfpathquadraticcurveto{\pgfqpoint{12.674254in}{0.642840in}}{\pgfqpoint{12.674254in}{0.651954in}}%
\pgfpathquadraticcurveto{\pgfqpoint{12.674254in}{0.662582in}}{\pgfqpoint{12.681369in}{0.667985in}}%
\pgfpathquadraticcurveto{\pgfqpoint{12.688502in}{0.673389in}}{\pgfqpoint{12.702623in}{0.673389in}}%
\pgfpathlineto{\pgfqpoint{12.717159in}{0.673389in}}%
\pgfpathlineto{\pgfqpoint{12.717159in}{0.674416in}}%
\pgfpathquadraticcurveto{\pgfqpoint{12.717159in}{0.681566in}}{\pgfqpoint{12.712458in}{0.685475in}}%
\pgfpathquadraticcurveto{\pgfqpoint{12.707757in}{0.689384in}}{\pgfqpoint{12.699255in}{0.689384in}}%
\pgfpathquadraticcurveto{\pgfqpoint{12.693851in}{0.689384in}}{\pgfqpoint{12.688718in}{0.688087in}}%
\pgfpathquadraticcurveto{\pgfqpoint{12.683602in}{0.686790in}}{\pgfqpoint{12.678883in}{0.684196in}}%
\pgfpathlineto{\pgfqpoint{12.678883in}{0.693779in}}%
\pgfpathquadraticcurveto{\pgfqpoint{12.684557in}{0.695976in}}{\pgfqpoint{12.689907in}{0.697057in}}%
\pgfpathquadraticcurveto{\pgfqpoint{12.695256in}{0.698156in}}{\pgfqpoint{12.700318in}{0.698156in}}%
\pgfpathquadraticcurveto{\pgfqpoint{12.714007in}{0.698156in}}{\pgfqpoint{12.720762in}{0.691059in}}%
\pgfpathquadraticcurveto{\pgfqpoint{12.727516in}{0.683980in}}{\pgfqpoint{12.727516in}{0.669570in}}%
\pgfpathlineto{\pgfqpoint{12.727516in}{0.669570in}}%
\pgfpathclose%
\pgfpathmoveto{\pgfqpoint{12.759091in}{0.714547in}}%
\pgfpathlineto{\pgfqpoint{12.759091in}{0.696643in}}%
\pgfpathlineto{\pgfqpoint{12.780418in}{0.696643in}}%
\pgfpathlineto{\pgfqpoint{12.780418in}{0.688591in}}%
\pgfpathlineto{\pgfqpoint{12.759091in}{0.688591in}}%
\pgfpathlineto{\pgfqpoint{12.759091in}{0.654368in}}%
\pgfpathquadraticcurveto{\pgfqpoint{12.759091in}{0.646659in}}{\pgfqpoint{12.761199in}{0.644461in}}%
\pgfpathquadraticcurveto{\pgfqpoint{12.763306in}{0.642264in}}{\pgfqpoint{12.769791in}{0.642264in}}%
\pgfpathlineto{\pgfqpoint{12.780418in}{0.642264in}}%
\pgfpathlineto{\pgfqpoint{12.780418in}{0.633600in}}%
\pgfpathlineto{\pgfqpoint{12.769791in}{0.633600in}}%
\pgfpathquadraticcurveto{\pgfqpoint{12.757794in}{0.633600in}}{\pgfqpoint{12.753237in}{0.638067in}}%
\pgfpathquadraticcurveto{\pgfqpoint{12.748680in}{0.642552in}}{\pgfqpoint{12.748680in}{0.654368in}}%
\pgfpathlineto{\pgfqpoint{12.748680in}{0.688591in}}%
\pgfpathlineto{\pgfqpoint{12.741079in}{0.688591in}}%
\pgfpathlineto{\pgfqpoint{12.741079in}{0.696643in}}%
\pgfpathlineto{\pgfqpoint{12.748680in}{0.696643in}}%
\pgfpathlineto{\pgfqpoint{12.748680in}{0.714547in}}%
\pgfpathlineto{\pgfqpoint{12.759091in}{0.714547in}}%
\pgfpathlineto{\pgfqpoint{12.759091in}{0.714547in}}%
\pgfpathclose%
\pgfpathmoveto{\pgfqpoint{12.818459in}{0.689384in}}%
\pgfpathquadraticcurveto{\pgfqpoint{12.810138in}{0.689384in}}{\pgfqpoint{12.805293in}{0.682881in}}%
\pgfpathquadraticcurveto{\pgfqpoint{12.800447in}{0.676379in}}{\pgfqpoint{12.800447in}{0.665067in}}%
\pgfpathquadraticcurveto{\pgfqpoint{12.800447in}{0.653756in}}{\pgfqpoint{12.805257in}{0.647253in}}%
\pgfpathquadraticcurveto{\pgfqpoint{12.810084in}{0.640751in}}{\pgfqpoint{12.818459in}{0.640751in}}%
\pgfpathquadraticcurveto{\pgfqpoint{12.826745in}{0.640751in}}{\pgfqpoint{12.831572in}{0.647271in}}%
\pgfpathquadraticcurveto{\pgfqpoint{12.836418in}{0.653810in}}{\pgfqpoint{12.836418in}{0.665067in}}%
\pgfpathquadraticcurveto{\pgfqpoint{12.836418in}{0.676271in}}{\pgfqpoint{12.831572in}{0.682827in}}%
\pgfpathquadraticcurveto{\pgfqpoint{12.826745in}{0.689384in}}{\pgfqpoint{12.818459in}{0.689384in}}%
\pgfpathlineto{\pgfqpoint{12.818459in}{0.689384in}}%
\pgfpathclose%
\pgfpathmoveto{\pgfqpoint{12.818459in}{0.698156in}}%
\pgfpathquadraticcurveto{\pgfqpoint{12.831969in}{0.698156in}}{\pgfqpoint{12.839678in}{0.689366in}}%
\pgfpathquadraticcurveto{\pgfqpoint{12.847405in}{0.680594in}}{\pgfqpoint{12.847405in}{0.665067in}}%
\pgfpathquadraticcurveto{\pgfqpoint{12.847405in}{0.649595in}}{\pgfqpoint{12.839678in}{0.640769in}}%
\pgfpathquadraticcurveto{\pgfqpoint{12.831969in}{0.631961in}}{\pgfqpoint{12.818459in}{0.631961in}}%
\pgfpathquadraticcurveto{\pgfqpoint{12.804896in}{0.631961in}}{\pgfqpoint{12.797205in}{0.640769in}}%
\pgfpathquadraticcurveto{\pgfqpoint{12.789532in}{0.649595in}}{\pgfqpoint{12.789532in}{0.665067in}}%
\pgfpathquadraticcurveto{\pgfqpoint{12.789532in}{0.680594in}}{\pgfqpoint{12.797205in}{0.689366in}}%
\pgfpathquadraticcurveto{\pgfqpoint{12.804896in}{0.698156in}}{\pgfqpoint{12.818459in}{0.698156in}}%
\pgfpathlineto{\pgfqpoint{12.818459in}{0.698156in}}%
\pgfpathclose%
\pgfpathmoveto{\pgfqpoint{12.901099in}{0.686970in}}%
\pgfpathquadraticcurveto{\pgfqpoint{12.899352in}{0.687979in}}{\pgfqpoint{12.897299in}{0.688447in}}%
\pgfpathquadraticcurveto{\pgfqpoint{12.895245in}{0.688933in}}{\pgfqpoint{12.892778in}{0.688933in}}%
\pgfpathquadraticcurveto{\pgfqpoint{12.883988in}{0.688933in}}{\pgfqpoint{12.879286in}{0.683223in}}%
\pgfpathquadraticcurveto{\pgfqpoint{12.874585in}{0.677514in}}{\pgfqpoint{12.874585in}{0.666814in}}%
\pgfpathlineto{\pgfqpoint{12.874585in}{0.633600in}}%
\pgfpathlineto{\pgfqpoint{12.864174in}{0.633600in}}%
\pgfpathlineto{\pgfqpoint{12.864174in}{0.696643in}}%
\pgfpathlineto{\pgfqpoint{12.874585in}{0.696643in}}%
\pgfpathlineto{\pgfqpoint{12.874585in}{0.686844in}}%
\pgfpathquadraticcurveto{\pgfqpoint{12.877863in}{0.692590in}}{\pgfqpoint{12.883087in}{0.695364in}}%
\pgfpathquadraticcurveto{\pgfqpoint{12.888329in}{0.698156in}}{\pgfqpoint{12.895822in}{0.698156in}}%
\pgfpathquadraticcurveto{\pgfqpoint{12.896884in}{0.698156in}}{\pgfqpoint{12.898181in}{0.698011in}}%
\pgfpathquadraticcurveto{\pgfqpoint{12.899478in}{0.697885in}}{\pgfqpoint{12.901045in}{0.697597in}}%
\pgfpathlineto{\pgfqpoint{12.901099in}{0.686970in}}%
\pgfpathlineto{\pgfqpoint{12.901099in}{0.686970in}}%
\pgfpathclose%
\pgfusepath{fill}%
\end{pgfscope}%
\begin{pgfscope}%
\pgfsetbuttcap%
\pgfsetroundjoin%
\definecolor{currentfill}{rgb}{0.309804,0.372549,0.419608}%
\pgfsetfillcolor{currentfill}%
\pgfsetlinewidth{0.000000pt}%
\definecolor{currentstroke}{rgb}{0.309804,0.372549,0.419608}%
\pgfsetstrokecolor{currentstroke}%
\pgfsetdash{}{0pt}%
\pgfpathmoveto{\pgfqpoint{12.979269in}{0.696643in}}%
\pgfpathlineto{\pgfqpoint{12.989626in}{0.696643in}}%
\pgfpathlineto{\pgfqpoint{13.002577in}{0.647451in}}%
\pgfpathlineto{\pgfqpoint{13.015455in}{0.696643in}}%
\pgfpathlineto{\pgfqpoint{13.027668in}{0.696643in}}%
\pgfpathlineto{\pgfqpoint{13.040618in}{0.647451in}}%
\pgfpathlineto{\pgfqpoint{13.053515in}{0.696643in}}%
\pgfpathlineto{\pgfqpoint{13.063872in}{0.696643in}}%
\pgfpathlineto{\pgfqpoint{13.047373in}{0.633600in}}%
\pgfpathlineto{\pgfqpoint{13.035161in}{0.633600in}}%
\pgfpathlineto{\pgfqpoint{13.021597in}{0.685277in}}%
\pgfpathlineto{\pgfqpoint{13.007980in}{0.633600in}}%
\pgfpathlineto{\pgfqpoint{12.995750in}{0.633600in}}%
\pgfpathlineto{\pgfqpoint{12.979269in}{0.696643in}}%
\pgfpathlineto{\pgfqpoint{12.979269in}{0.696643in}}%
\pgfpathclose%
\pgfpathmoveto{\pgfqpoint{13.133489in}{0.667715in}}%
\pgfpathlineto{\pgfqpoint{13.133489in}{0.662654in}}%
\pgfpathlineto{\pgfqpoint{13.085865in}{0.662654in}}%
\pgfpathquadraticcurveto{\pgfqpoint{13.086549in}{0.651954in}}{\pgfqpoint{13.092313in}{0.646353in}}%
\pgfpathquadraticcurveto{\pgfqpoint{13.098077in}{0.640751in}}{\pgfqpoint{13.108380in}{0.640751in}}%
\pgfpathquadraticcurveto{\pgfqpoint{13.114342in}{0.640751in}}{\pgfqpoint{13.119944in}{0.642210in}}%
\pgfpathquadraticcurveto{\pgfqpoint{13.125546in}{0.643669in}}{\pgfqpoint{13.131075in}{0.646605in}}%
\pgfpathlineto{\pgfqpoint{13.131075in}{0.636806in}}%
\pgfpathquadraticcurveto{\pgfqpoint{13.125492in}{0.634447in}}{\pgfqpoint{13.119638in}{0.633204in}}%
\pgfpathquadraticcurveto{\pgfqpoint{13.113784in}{0.631961in}}{\pgfqpoint{13.107768in}{0.631961in}}%
\pgfpathquadraticcurveto{\pgfqpoint{13.092673in}{0.631961in}}{\pgfqpoint{13.083865in}{0.640733in}}%
\pgfpathquadraticcurveto{\pgfqpoint{13.075058in}{0.649523in}}{\pgfqpoint{13.075058in}{0.664509in}}%
\pgfpathquadraticcurveto{\pgfqpoint{13.075058in}{0.679981in}}{\pgfqpoint{13.083415in}{0.689059in}}%
\pgfpathquadraticcurveto{\pgfqpoint{13.091773in}{0.698156in}}{\pgfqpoint{13.105966in}{0.698156in}}%
\pgfpathquadraticcurveto{\pgfqpoint{13.118683in}{0.698156in}}{\pgfqpoint{13.126086in}{0.689960in}}%
\pgfpathquadraticcurveto{\pgfqpoint{13.133489in}{0.681783in}}{\pgfqpoint{13.133489in}{0.667715in}}%
\pgfpathlineto{\pgfqpoint{13.133489in}{0.667715in}}%
\pgfpathclose%
\pgfpathmoveto{\pgfqpoint{13.123132in}{0.670759in}}%
\pgfpathquadraticcurveto{\pgfqpoint{13.123024in}{0.679243in}}{\pgfqpoint{13.118377in}{0.684304in}}%
\pgfpathquadraticcurveto{\pgfqpoint{13.113730in}{0.689384in}}{\pgfqpoint{13.106074in}{0.689384in}}%
\pgfpathquadraticcurveto{\pgfqpoint{13.097411in}{0.689384in}}{\pgfqpoint{13.092205in}{0.684484in}}%
\pgfpathquadraticcurveto{\pgfqpoint{13.087000in}{0.679585in}}{\pgfqpoint{13.086207in}{0.670687in}}%
\pgfpathlineto{\pgfqpoint{13.123132in}{0.670759in}}%
\pgfpathlineto{\pgfqpoint{13.123132in}{0.670759in}}%
\pgfpathclose%
\pgfpathmoveto{\pgfqpoint{13.150492in}{0.696643in}}%
\pgfpathlineto{\pgfqpoint{13.160849in}{0.696643in}}%
\pgfpathlineto{\pgfqpoint{13.160849in}{0.633600in}}%
\pgfpathlineto{\pgfqpoint{13.150492in}{0.633600in}}%
\pgfpathlineto{\pgfqpoint{13.150492in}{0.696643in}}%
\pgfpathlineto{\pgfqpoint{13.150492in}{0.696643in}}%
\pgfpathclose%
\pgfpathmoveto{\pgfqpoint{13.150492in}{0.721193in}}%
\pgfpathlineto{\pgfqpoint{13.160849in}{0.721193in}}%
\pgfpathlineto{\pgfqpoint{13.160849in}{0.708062in}}%
\pgfpathlineto{\pgfqpoint{13.150492in}{0.708062in}}%
\pgfpathlineto{\pgfqpoint{13.150492in}{0.721193in}}%
\pgfpathlineto{\pgfqpoint{13.150492in}{0.721193in}}%
\pgfpathclose%
\pgfpathmoveto{\pgfqpoint{13.224000in}{0.665860in}}%
\pgfpathquadraticcurveto{\pgfqpoint{13.224000in}{0.677117in}}{\pgfqpoint{13.219353in}{0.683296in}}%
\pgfpathquadraticcurveto{\pgfqpoint{13.214724in}{0.689492in}}{\pgfqpoint{13.206330in}{0.689492in}}%
\pgfpathquadraticcurveto{\pgfqpoint{13.198008in}{0.689492in}}{\pgfqpoint{13.193361in}{0.683296in}}%
\pgfpathquadraticcurveto{\pgfqpoint{13.188714in}{0.677117in}}{\pgfqpoint{13.188714in}{0.665860in}}%
\pgfpathquadraticcurveto{\pgfqpoint{13.188714in}{0.654656in}}{\pgfqpoint{13.193361in}{0.648460in}}%
\pgfpathquadraticcurveto{\pgfqpoint{13.198008in}{0.642264in}}{\pgfqpoint{13.206330in}{0.642264in}}%
\pgfpathquadraticcurveto{\pgfqpoint{13.214724in}{0.642264in}}{\pgfqpoint{13.219353in}{0.648460in}}%
\pgfpathquadraticcurveto{\pgfqpoint{13.224000in}{0.654656in}}{\pgfqpoint{13.224000in}{0.665860in}}%
\pgfpathlineto{\pgfqpoint{13.224000in}{0.665860in}}%
\pgfpathclose%
\pgfpathmoveto{\pgfqpoint{13.234357in}{0.641417in}}%
\pgfpathquadraticcurveto{\pgfqpoint{13.234357in}{0.625332in}}{\pgfqpoint{13.227206in}{0.617479in}}%
\pgfpathquadraticcurveto{\pgfqpoint{13.220073in}{0.609626in}}{\pgfqpoint{13.205321in}{0.609626in}}%
\pgfpathquadraticcurveto{\pgfqpoint{13.199864in}{0.609626in}}{\pgfqpoint{13.195018in}{0.610436in}}%
\pgfpathquadraticcurveto{\pgfqpoint{13.190173in}{0.611247in}}{\pgfqpoint{13.185616in}{0.612940in}}%
\pgfpathlineto{\pgfqpoint{13.185616in}{0.623009in}}%
\pgfpathquadraticcurveto{\pgfqpoint{13.190173in}{0.620541in}}{\pgfqpoint{13.194622in}{0.619370in}}%
\pgfpathquadraticcurveto{\pgfqpoint{13.199071in}{0.618182in}}{\pgfqpoint{13.203682in}{0.618182in}}%
\pgfpathquadraticcurveto{\pgfqpoint{13.213877in}{0.618182in}}{\pgfqpoint{13.218939in}{0.623495in}}%
\pgfpathquadraticcurveto{\pgfqpoint{13.224000in}{0.628809in}}{\pgfqpoint{13.224000in}{0.639562in}}%
\pgfpathlineto{\pgfqpoint{13.224000in}{0.644695in}}%
\pgfpathquadraticcurveto{\pgfqpoint{13.220794in}{0.639112in}}{\pgfqpoint{13.215786in}{0.636356in}}%
\pgfpathquadraticcurveto{\pgfqpoint{13.210779in}{0.633600in}}{\pgfqpoint{13.203790in}{0.633600in}}%
\pgfpathquadraticcurveto{\pgfqpoint{13.192209in}{0.633600in}}{\pgfqpoint{13.185112in}{0.642426in}}%
\pgfpathquadraticcurveto{\pgfqpoint{13.178015in}{0.651270in}}{\pgfqpoint{13.178015in}{0.665860in}}%
\pgfpathquadraticcurveto{\pgfqpoint{13.178015in}{0.680486in}}{\pgfqpoint{13.185112in}{0.689312in}}%
\pgfpathquadraticcurveto{\pgfqpoint{13.192209in}{0.698156in}}{\pgfqpoint{13.203790in}{0.698156in}}%
\pgfpathquadraticcurveto{\pgfqpoint{13.210779in}{0.698156in}}{\pgfqpoint{13.215786in}{0.695400in}}%
\pgfpathquadraticcurveto{\pgfqpoint{13.220794in}{0.692644in}}{\pgfqpoint{13.224000in}{0.687078in}}%
\pgfpathlineto{\pgfqpoint{13.224000in}{0.696643in}}%
\pgfpathlineto{\pgfqpoint{13.234357in}{0.696643in}}%
\pgfpathlineto{\pgfqpoint{13.234357in}{0.641417in}}%
\pgfpathlineto{\pgfqpoint{13.234357in}{0.641417in}}%
\pgfpathclose%
\pgfpathmoveto{\pgfqpoint{13.308117in}{0.671660in}}%
\pgfpathlineto{\pgfqpoint{13.308117in}{0.633600in}}%
\pgfpathlineto{\pgfqpoint{13.297760in}{0.633600in}}%
\pgfpathlineto{\pgfqpoint{13.297760in}{0.671317in}}%
\pgfpathquadraticcurveto{\pgfqpoint{13.297760in}{0.680269in}}{\pgfqpoint{13.294265in}{0.684700in}}%
\pgfpathquadraticcurveto{\pgfqpoint{13.290771in}{0.689149in}}{\pgfqpoint{13.283800in}{0.689149in}}%
\pgfpathquadraticcurveto{\pgfqpoint{13.275407in}{0.689149in}}{\pgfqpoint{13.270561in}{0.683800in}}%
\pgfpathquadraticcurveto{\pgfqpoint{13.265716in}{0.678468in}}{\pgfqpoint{13.265716in}{0.669228in}}%
\pgfpathlineto{\pgfqpoint{13.265716in}{0.633600in}}%
\pgfpathlineto{\pgfqpoint{13.255305in}{0.633600in}}%
\pgfpathlineto{\pgfqpoint{13.255305in}{0.721193in}}%
\pgfpathlineto{\pgfqpoint{13.265716in}{0.721193in}}%
\pgfpathlineto{\pgfqpoint{13.265716in}{0.686844in}}%
\pgfpathquadraticcurveto{\pgfqpoint{13.269445in}{0.692536in}}{\pgfqpoint{13.274470in}{0.695346in}}%
\pgfpathquadraticcurveto{\pgfqpoint{13.279513in}{0.698156in}}{\pgfqpoint{13.286106in}{0.698156in}}%
\pgfpathquadraticcurveto{\pgfqpoint{13.296967in}{0.698156in}}{\pgfqpoint{13.302533in}{0.691437in}}%
\pgfpathquadraticcurveto{\pgfqpoint{13.308117in}{0.684718in}}{\pgfqpoint{13.308117in}{0.671660in}}%
\pgfpathlineto{\pgfqpoint{13.308117in}{0.671660in}}%
\pgfpathclose%
\pgfpathmoveto{\pgfqpoint{13.339008in}{0.714547in}}%
\pgfpathlineto{\pgfqpoint{13.339008in}{0.696643in}}%
\pgfpathlineto{\pgfqpoint{13.360334in}{0.696643in}}%
\pgfpathlineto{\pgfqpoint{13.360334in}{0.688591in}}%
\pgfpathlineto{\pgfqpoint{13.339008in}{0.688591in}}%
\pgfpathlineto{\pgfqpoint{13.339008in}{0.654368in}}%
\pgfpathquadraticcurveto{\pgfqpoint{13.339008in}{0.646659in}}{\pgfqpoint{13.341115in}{0.644461in}}%
\pgfpathquadraticcurveto{\pgfqpoint{13.343222in}{0.642264in}}{\pgfqpoint{13.349707in}{0.642264in}}%
\pgfpathlineto{\pgfqpoint{13.360334in}{0.642264in}}%
\pgfpathlineto{\pgfqpoint{13.360334in}{0.633600in}}%
\pgfpathlineto{\pgfqpoint{13.349707in}{0.633600in}}%
\pgfpathquadraticcurveto{\pgfqpoint{13.337711in}{0.633600in}}{\pgfqpoint{13.333154in}{0.638067in}}%
\pgfpathquadraticcurveto{\pgfqpoint{13.328597in}{0.642552in}}{\pgfqpoint{13.328597in}{0.654368in}}%
\pgfpathlineto{\pgfqpoint{13.328597in}{0.688591in}}%
\pgfpathlineto{\pgfqpoint{13.320995in}{0.688591in}}%
\pgfpathlineto{\pgfqpoint{13.320995in}{0.696643in}}%
\pgfpathlineto{\pgfqpoint{13.328597in}{0.696643in}}%
\pgfpathlineto{\pgfqpoint{13.328597in}{0.714547in}}%
\pgfpathlineto{\pgfqpoint{13.339008in}{0.714547in}}%
\pgfpathlineto{\pgfqpoint{13.339008in}{0.714547in}}%
\pgfpathclose%
\pgfpathmoveto{\pgfqpoint{13.414136in}{0.694787in}}%
\pgfpathlineto{\pgfqpoint{13.414136in}{0.684989in}}%
\pgfpathquadraticcurveto{\pgfqpoint{13.409759in}{0.687240in}}{\pgfqpoint{13.405022in}{0.688357in}}%
\pgfpathquadraticcurveto{\pgfqpoint{13.400303in}{0.689492in}}{\pgfqpoint{13.395224in}{0.689492in}}%
\pgfpathquadraticcurveto{\pgfqpoint{13.387514in}{0.689492in}}{\pgfqpoint{13.383660in}{0.687132in}}%
\pgfpathquadraticcurveto{\pgfqpoint{13.379805in}{0.684773in}}{\pgfqpoint{13.379805in}{0.680035in}}%
\pgfpathquadraticcurveto{\pgfqpoint{13.379805in}{0.676433in}}{\pgfqpoint{13.382561in}{0.674380in}}%
\pgfpathquadraticcurveto{\pgfqpoint{13.385317in}{0.672326in}}{\pgfqpoint{13.393656in}{0.670471in}}%
\pgfpathlineto{\pgfqpoint{13.397205in}{0.669678in}}%
\pgfpathquadraticcurveto{\pgfqpoint{13.408228in}{0.667319in}}{\pgfqpoint{13.412875in}{0.663014in}}%
\pgfpathquadraticcurveto{\pgfqpoint{13.417523in}{0.658709in}}{\pgfqpoint{13.417523in}{0.651000in}}%
\pgfpathquadraticcurveto{\pgfqpoint{13.417523in}{0.642210in}}{\pgfqpoint{13.410570in}{0.637076in}}%
\pgfpathquadraticcurveto{\pgfqpoint{13.403617in}{0.631961in}}{\pgfqpoint{13.391459in}{0.631961in}}%
\pgfpathquadraticcurveto{\pgfqpoint{13.386398in}{0.631961in}}{\pgfqpoint{13.380904in}{0.632952in}}%
\pgfpathquadraticcurveto{\pgfqpoint{13.375410in}{0.633942in}}{\pgfqpoint{13.369340in}{0.635906in}}%
\pgfpathlineto{\pgfqpoint{13.369340in}{0.646605in}}%
\pgfpathquadraticcurveto{\pgfqpoint{13.375086in}{0.643615in}}{\pgfqpoint{13.380652in}{0.642120in}}%
\pgfpathquadraticcurveto{\pgfqpoint{13.386217in}{0.640643in}}{\pgfqpoint{13.391693in}{0.640643in}}%
\pgfpathquadraticcurveto{\pgfqpoint{13.399006in}{0.640643in}}{\pgfqpoint{13.402933in}{0.643146in}}%
\pgfpathquadraticcurveto{\pgfqpoint{13.406877in}{0.645650in}}{\pgfqpoint{13.406877in}{0.650207in}}%
\pgfpathquadraticcurveto{\pgfqpoint{13.406877in}{0.654422in}}{\pgfqpoint{13.404031in}{0.656674in}}%
\pgfpathquadraticcurveto{\pgfqpoint{13.401204in}{0.658925in}}{\pgfqpoint{13.391567in}{0.661014in}}%
\pgfpathlineto{\pgfqpoint{13.387965in}{0.661861in}}%
\pgfpathquadraticcurveto{\pgfqpoint{13.378346in}{0.663878in}}{\pgfqpoint{13.374059in}{0.668075in}}%
\pgfpathquadraticcurveto{\pgfqpoint{13.369790in}{0.672272in}}{\pgfqpoint{13.369790in}{0.679585in}}%
\pgfpathquadraticcurveto{\pgfqpoint{13.369790in}{0.688483in}}{\pgfqpoint{13.376095in}{0.693310in}}%
\pgfpathquadraticcurveto{\pgfqpoint{13.382399in}{0.698156in}}{\pgfqpoint{13.393999in}{0.698156in}}%
\pgfpathquadraticcurveto{\pgfqpoint{13.399727in}{0.698156in}}{\pgfqpoint{13.404788in}{0.697309in}}%
\pgfpathquadraticcurveto{\pgfqpoint{13.409867in}{0.696480in}}{\pgfqpoint{13.414136in}{0.694787in}}%
\pgfpathlineto{\pgfqpoint{13.414136in}{0.694787in}}%
\pgfpathclose%
\pgfusepath{fill}%
\end{pgfscope}%
\begin{pgfscope}%
\pgfsetbuttcap%
\pgfsetroundjoin%
\definecolor{currentfill}{rgb}{0.309804,0.372549,0.419608}%
\pgfsetfillcolor{currentfill}%
\pgfsetlinewidth{0.000000pt}%
\definecolor{currentstroke}{rgb}{0.309804,0.372549,0.419608}%
\pgfsetstrokecolor{currentstroke}%
\pgfsetdash{}{0pt}%
\pgfpathmoveto{\pgfqpoint{13.537286in}{0.689384in}}%
\pgfpathquadraticcurveto{\pgfqpoint{13.528964in}{0.689384in}}{\pgfqpoint{13.524119in}{0.682881in}}%
\pgfpathquadraticcurveto{\pgfqpoint{13.519274in}{0.676379in}}{\pgfqpoint{13.519274in}{0.665067in}}%
\pgfpathquadraticcurveto{\pgfqpoint{13.519274in}{0.653756in}}{\pgfqpoint{13.524083in}{0.647253in}}%
\pgfpathquadraticcurveto{\pgfqpoint{13.528910in}{0.640751in}}{\pgfqpoint{13.537286in}{0.640751in}}%
\pgfpathquadraticcurveto{\pgfqpoint{13.545571in}{0.640751in}}{\pgfqpoint{13.550399in}{0.647271in}}%
\pgfpathquadraticcurveto{\pgfqpoint{13.555244in}{0.653810in}}{\pgfqpoint{13.555244in}{0.665067in}}%
\pgfpathquadraticcurveto{\pgfqpoint{13.555244in}{0.676271in}}{\pgfqpoint{13.550399in}{0.682827in}}%
\pgfpathquadraticcurveto{\pgfqpoint{13.545571in}{0.689384in}}{\pgfqpoint{13.537286in}{0.689384in}}%
\pgfpathlineto{\pgfqpoint{13.537286in}{0.689384in}}%
\pgfpathclose%
\pgfpathmoveto{\pgfqpoint{13.537286in}{0.698156in}}%
\pgfpathquadraticcurveto{\pgfqpoint{13.550795in}{0.698156in}}{\pgfqpoint{13.558504in}{0.689366in}}%
\pgfpathquadraticcurveto{\pgfqpoint{13.566231in}{0.680594in}}{\pgfqpoint{13.566231in}{0.665067in}}%
\pgfpathquadraticcurveto{\pgfqpoint{13.566231in}{0.649595in}}{\pgfqpoint{13.558504in}{0.640769in}}%
\pgfpathquadraticcurveto{\pgfqpoint{13.550795in}{0.631961in}}{\pgfqpoint{13.537286in}{0.631961in}}%
\pgfpathquadraticcurveto{\pgfqpoint{13.523723in}{0.631961in}}{\pgfqpoint{13.516031in}{0.640769in}}%
\pgfpathquadraticcurveto{\pgfqpoint{13.508358in}{0.649595in}}{\pgfqpoint{13.508358in}{0.665067in}}%
\pgfpathquadraticcurveto{\pgfqpoint{13.508358in}{0.680594in}}{\pgfqpoint{13.516031in}{0.689366in}}%
\pgfpathquadraticcurveto{\pgfqpoint{13.523723in}{0.698156in}}{\pgfqpoint{13.537286in}{0.698156in}}%
\pgfpathlineto{\pgfqpoint{13.537286in}{0.698156in}}%
\pgfpathclose%
\pgfpathmoveto{\pgfqpoint{13.619926in}{0.686970in}}%
\pgfpathquadraticcurveto{\pgfqpoint{13.618178in}{0.687979in}}{\pgfqpoint{13.616125in}{0.688447in}}%
\pgfpathquadraticcurveto{\pgfqpoint{13.614072in}{0.688933in}}{\pgfqpoint{13.611604in}{0.688933in}}%
\pgfpathquadraticcurveto{\pgfqpoint{13.602814in}{0.688933in}}{\pgfqpoint{13.598113in}{0.683223in}}%
\pgfpathquadraticcurveto{\pgfqpoint{13.593412in}{0.677514in}}{\pgfqpoint{13.593412in}{0.666814in}}%
\pgfpathlineto{\pgfqpoint{13.593412in}{0.633600in}}%
\pgfpathlineto{\pgfqpoint{13.583001in}{0.633600in}}%
\pgfpathlineto{\pgfqpoint{13.583001in}{0.696643in}}%
\pgfpathlineto{\pgfqpoint{13.593412in}{0.696643in}}%
\pgfpathlineto{\pgfqpoint{13.593412in}{0.686844in}}%
\pgfpathquadraticcurveto{\pgfqpoint{13.596690in}{0.692590in}}{\pgfqpoint{13.601913in}{0.695364in}}%
\pgfpathquadraticcurveto{\pgfqpoint{13.607155in}{0.698156in}}{\pgfqpoint{13.614648in}{0.698156in}}%
\pgfpathquadraticcurveto{\pgfqpoint{13.615711in}{0.698156in}}{\pgfqpoint{13.617008in}{0.698011in}}%
\pgfpathquadraticcurveto{\pgfqpoint{13.618304in}{0.697885in}}{\pgfqpoint{13.619872in}{0.697597in}}%
\pgfpathlineto{\pgfqpoint{13.619926in}{0.686970in}}%
\pgfpathlineto{\pgfqpoint{13.619926in}{0.686970in}}%
\pgfpathclose%
\pgfusepath{fill}%
\end{pgfscope}%
\begin{pgfscope}%
\pgfsetbuttcap%
\pgfsetroundjoin%
\definecolor{currentfill}{rgb}{0.309804,0.372549,0.419608}%
\pgfsetfillcolor{currentfill}%
\pgfsetlinewidth{0.000000pt}%
\definecolor{currentstroke}{rgb}{0.309804,0.372549,0.419608}%
\pgfsetstrokecolor{currentstroke}%
\pgfsetdash{}{0pt}%
\pgfpathmoveto{\pgfqpoint{0.988790in}{0.504515in}}%
\pgfpathlineto{\pgfqpoint{0.988790in}{0.499454in}}%
\pgfpathlineto{\pgfqpoint{0.941166in}{0.499454in}}%
\pgfpathquadraticcurveto{\pgfqpoint{0.941850in}{0.488754in}}{\pgfqpoint{0.947614in}{0.483153in}}%
\pgfpathquadraticcurveto{\pgfqpoint{0.953378in}{0.477551in}}{\pgfqpoint{0.963681in}{0.477551in}}%
\pgfpathquadraticcurveto{\pgfqpoint{0.969643in}{0.477551in}}{\pgfqpoint{0.975245in}{0.479010in}}%
\pgfpathquadraticcurveto{\pgfqpoint{0.980846in}{0.480469in}}{\pgfqpoint{0.986376in}{0.483405in}}%
\pgfpathlineto{\pgfqpoint{0.986376in}{0.473606in}}%
\pgfpathquadraticcurveto{\pgfqpoint{0.980792in}{0.471247in}}{\pgfqpoint{0.974938in}{0.470004in}}%
\pgfpathquadraticcurveto{\pgfqpoint{0.969084in}{0.468761in}}{\pgfqpoint{0.963068in}{0.468761in}}%
\pgfpathquadraticcurveto{\pgfqpoint{0.947974in}{0.468761in}}{\pgfqpoint{0.939166in}{0.477533in}}%
\pgfpathquadraticcurveto{\pgfqpoint{0.930358in}{0.486323in}}{\pgfqpoint{0.930358in}{0.501309in}}%
\pgfpathquadraticcurveto{\pgfqpoint{0.930358in}{0.516781in}}{\pgfqpoint{0.938716in}{0.525859in}}%
\pgfpathquadraticcurveto{\pgfqpoint{0.947074in}{0.534956in}}{\pgfqpoint{0.961267in}{0.534956in}}%
\pgfpathquadraticcurveto{\pgfqpoint{0.973984in}{0.534956in}}{\pgfqpoint{0.981387in}{0.526760in}}%
\pgfpathquadraticcurveto{\pgfqpoint{0.988790in}{0.518583in}}{\pgfqpoint{0.988790in}{0.504515in}}%
\pgfpathlineto{\pgfqpoint{0.988790in}{0.504515in}}%
\pgfpathclose%
\pgfpathmoveto{\pgfqpoint{0.978433in}{0.507559in}}%
\pgfpathquadraticcurveto{\pgfqpoint{0.978325in}{0.516043in}}{\pgfqpoint{0.973678in}{0.521104in}}%
\pgfpathquadraticcurveto{\pgfqpoint{0.969030in}{0.526184in}}{\pgfqpoint{0.961375in}{0.526184in}}%
\pgfpathquadraticcurveto{\pgfqpoint{0.952711in}{0.526184in}}{\pgfqpoint{0.947506in}{0.521284in}}%
\pgfpathquadraticcurveto{\pgfqpoint{0.942300in}{0.516385in}}{\pgfqpoint{0.941508in}{0.507487in}}%
\pgfpathlineto{\pgfqpoint{0.978433in}{0.507559in}}%
\pgfpathlineto{\pgfqpoint{0.978433in}{0.507559in}}%
\pgfpathclose%
\pgfpathmoveto{\pgfqpoint{1.045978in}{0.531587in}}%
\pgfpathlineto{\pgfqpoint{1.045978in}{0.521789in}}%
\pgfpathquadraticcurveto{\pgfqpoint{1.041601in}{0.524040in}}{\pgfqpoint{1.036864in}{0.525157in}}%
\pgfpathquadraticcurveto{\pgfqpoint{1.032145in}{0.526292in}}{\pgfqpoint{1.027066in}{0.526292in}}%
\pgfpathquadraticcurveto{\pgfqpoint{1.019356in}{0.526292in}}{\pgfqpoint{1.015502in}{0.523932in}}%
\pgfpathquadraticcurveto{\pgfqpoint{1.011647in}{0.521573in}}{\pgfqpoint{1.011647in}{0.516835in}}%
\pgfpathquadraticcurveto{\pgfqpoint{1.011647in}{0.513233in}}{\pgfqpoint{1.014403in}{0.511180in}}%
\pgfpathquadraticcurveto{\pgfqpoint{1.017159in}{0.509126in}}{\pgfqpoint{1.025498in}{0.507271in}}%
\pgfpathlineto{\pgfqpoint{1.029047in}{0.506478in}}%
\pgfpathquadraticcurveto{\pgfqpoint{1.040070in}{0.504119in}}{\pgfqpoint{1.044717in}{0.499814in}}%
\pgfpathquadraticcurveto{\pgfqpoint{1.049365in}{0.495509in}}{\pgfqpoint{1.049365in}{0.487800in}}%
\pgfpathquadraticcurveto{\pgfqpoint{1.049365in}{0.479010in}}{\pgfqpoint{1.042412in}{0.473876in}}%
\pgfpathquadraticcurveto{\pgfqpoint{1.035459in}{0.468761in}}{\pgfqpoint{1.023301in}{0.468761in}}%
\pgfpathquadraticcurveto{\pgfqpoint{1.018240in}{0.468761in}}{\pgfqpoint{1.012746in}{0.469752in}}%
\pgfpathquadraticcurveto{\pgfqpoint{1.007252in}{0.470742in}}{\pgfqpoint{1.001182in}{0.472706in}}%
\pgfpathlineto{\pgfqpoint{1.001182in}{0.483405in}}%
\pgfpathquadraticcurveto{\pgfqpoint{1.006928in}{0.480415in}}{\pgfqpoint{1.012494in}{0.478920in}}%
\pgfpathquadraticcurveto{\pgfqpoint{1.018059in}{0.477443in}}{\pgfqpoint{1.023535in}{0.477443in}}%
\pgfpathquadraticcurveto{\pgfqpoint{1.030848in}{0.477443in}}{\pgfqpoint{1.034775in}{0.479946in}}%
\pgfpathquadraticcurveto{\pgfqpoint{1.038719in}{0.482450in}}{\pgfqpoint{1.038719in}{0.487007in}}%
\pgfpathquadraticcurveto{\pgfqpoint{1.038719in}{0.491222in}}{\pgfqpoint{1.035873in}{0.493474in}}%
\pgfpathquadraticcurveto{\pgfqpoint{1.033046in}{0.495725in}}{\pgfqpoint{1.023409in}{0.497814in}}%
\pgfpathlineto{\pgfqpoint{1.019807in}{0.498661in}}%
\pgfpathquadraticcurveto{\pgfqpoint{1.010188in}{0.500678in}}{\pgfqpoint{1.005901in}{0.504875in}}%
\pgfpathquadraticcurveto{\pgfqpoint{1.001632in}{0.509072in}}{\pgfqpoint{1.001632in}{0.516385in}}%
\pgfpathquadraticcurveto{\pgfqpoint{1.001632in}{0.525283in}}{\pgfqpoint{1.007937in}{0.530110in}}%
\pgfpathquadraticcurveto{\pgfqpoint{1.014241in}{0.534956in}}{\pgfqpoint{1.025841in}{0.534956in}}%
\pgfpathquadraticcurveto{\pgfqpoint{1.031569in}{0.534956in}}{\pgfqpoint{1.036630in}{0.534109in}}%
\pgfpathquadraticcurveto{\pgfqpoint{1.041709in}{0.533280in}}{\pgfqpoint{1.045978in}{0.531587in}}%
\pgfpathlineto{\pgfqpoint{1.045978in}{0.531587in}}%
\pgfpathclose%
\pgfpathmoveto{\pgfqpoint{1.076095in}{0.551347in}}%
\pgfpathlineto{\pgfqpoint{1.076095in}{0.533443in}}%
\pgfpathlineto{\pgfqpoint{1.097421in}{0.533443in}}%
\pgfpathlineto{\pgfqpoint{1.097421in}{0.525391in}}%
\pgfpathlineto{\pgfqpoint{1.076095in}{0.525391in}}%
\pgfpathlineto{\pgfqpoint{1.076095in}{0.491168in}}%
\pgfpathquadraticcurveto{\pgfqpoint{1.076095in}{0.483459in}}{\pgfqpoint{1.078202in}{0.481261in}}%
\pgfpathquadraticcurveto{\pgfqpoint{1.080309in}{0.479064in}}{\pgfqpoint{1.086794in}{0.479064in}}%
\pgfpathlineto{\pgfqpoint{1.097421in}{0.479064in}}%
\pgfpathlineto{\pgfqpoint{1.097421in}{0.470400in}}%
\pgfpathlineto{\pgfqpoint{1.086794in}{0.470400in}}%
\pgfpathquadraticcurveto{\pgfqpoint{1.074798in}{0.470400in}}{\pgfqpoint{1.070241in}{0.474867in}}%
\pgfpathquadraticcurveto{\pgfqpoint{1.065684in}{0.479352in}}{\pgfqpoint{1.065684in}{0.491168in}}%
\pgfpathlineto{\pgfqpoint{1.065684in}{0.525391in}}%
\pgfpathlineto{\pgfqpoint{1.058082in}{0.525391in}}%
\pgfpathlineto{\pgfqpoint{1.058082in}{0.533443in}}%
\pgfpathlineto{\pgfqpoint{1.065684in}{0.533443in}}%
\pgfpathlineto{\pgfqpoint{1.065684in}{0.551347in}}%
\pgfpathlineto{\pgfqpoint{1.076095in}{0.551347in}}%
\pgfpathlineto{\pgfqpoint{1.076095in}{0.551347in}}%
\pgfpathclose%
\pgfpathmoveto{\pgfqpoint{1.111038in}{0.533443in}}%
\pgfpathlineto{\pgfqpoint{1.121395in}{0.533443in}}%
\pgfpathlineto{\pgfqpoint{1.121395in}{0.470400in}}%
\pgfpathlineto{\pgfqpoint{1.111038in}{0.470400in}}%
\pgfpathlineto{\pgfqpoint{1.111038in}{0.533443in}}%
\pgfpathlineto{\pgfqpoint{1.111038in}{0.533443in}}%
\pgfpathclose%
\pgfpathmoveto{\pgfqpoint{1.111038in}{0.557993in}}%
\pgfpathlineto{\pgfqpoint{1.121395in}{0.557993in}}%
\pgfpathlineto{\pgfqpoint{1.121395in}{0.544862in}}%
\pgfpathlineto{\pgfqpoint{1.111038in}{0.544862in}}%
\pgfpathlineto{\pgfqpoint{1.111038in}{0.557993in}}%
\pgfpathlineto{\pgfqpoint{1.111038in}{0.557993in}}%
\pgfpathclose%
\pgfpathmoveto{\pgfqpoint{1.192147in}{0.521338in}}%
\pgfpathquadraticcurveto{\pgfqpoint{1.196038in}{0.528327in}}{\pgfqpoint{1.201441in}{0.531641in}}%
\pgfpathquadraticcurveto{\pgfqpoint{1.206845in}{0.534956in}}{\pgfqpoint{1.214158in}{0.534956in}}%
\pgfpathquadraticcurveto{\pgfqpoint{1.224010in}{0.534956in}}{\pgfqpoint{1.229360in}{0.528057in}}%
\pgfpathquadraticcurveto{\pgfqpoint{1.234710in}{0.521176in}}{\pgfqpoint{1.234710in}{0.508460in}}%
\pgfpathlineto{\pgfqpoint{1.234710in}{0.470400in}}%
\pgfpathlineto{\pgfqpoint{1.224299in}{0.470400in}}%
\pgfpathlineto{\pgfqpoint{1.224299in}{0.508117in}}%
\pgfpathquadraticcurveto{\pgfqpoint{1.224299in}{0.517178in}}{\pgfqpoint{1.221074in}{0.521555in}}%
\pgfpathquadraticcurveto{\pgfqpoint{1.217868in}{0.525949in}}{\pgfqpoint{1.211294in}{0.525949in}}%
\pgfpathquadraticcurveto{\pgfqpoint{1.203242in}{0.525949in}}{\pgfqpoint{1.198559in}{0.520600in}}%
\pgfpathquadraticcurveto{\pgfqpoint{1.193894in}{0.515268in}}{\pgfqpoint{1.193894in}{0.506028in}}%
\pgfpathlineto{\pgfqpoint{1.193894in}{0.470400in}}%
\pgfpathlineto{\pgfqpoint{1.183483in}{0.470400in}}%
\pgfpathlineto{\pgfqpoint{1.183483in}{0.508117in}}%
\pgfpathquadraticcurveto{\pgfqpoint{1.183483in}{0.517232in}}{\pgfqpoint{1.180277in}{0.521591in}}%
\pgfpathquadraticcurveto{\pgfqpoint{1.177071in}{0.525949in}}{\pgfqpoint{1.170370in}{0.525949in}}%
\pgfpathquadraticcurveto{\pgfqpoint{1.162427in}{0.525949in}}{\pgfqpoint{1.157744in}{0.520582in}}%
\pgfpathquadraticcurveto{\pgfqpoint{1.153079in}{0.515214in}}{\pgfqpoint{1.153079in}{0.506028in}}%
\pgfpathlineto{\pgfqpoint{1.153079in}{0.470400in}}%
\pgfpathlineto{\pgfqpoint{1.142668in}{0.470400in}}%
\pgfpathlineto{\pgfqpoint{1.142668in}{0.533443in}}%
\pgfpathlineto{\pgfqpoint{1.153079in}{0.533443in}}%
\pgfpathlineto{\pgfqpoint{1.153079in}{0.523644in}}%
\pgfpathquadraticcurveto{\pgfqpoint{1.156627in}{0.529444in}}{\pgfqpoint{1.161580in}{0.532200in}}%
\pgfpathquadraticcurveto{\pgfqpoint{1.166534in}{0.534956in}}{\pgfqpoint{1.173342in}{0.534956in}}%
\pgfpathquadraticcurveto{\pgfqpoint{1.180223in}{0.534956in}}{\pgfqpoint{1.185032in}{0.531461in}}%
\pgfpathquadraticcurveto{\pgfqpoint{1.189841in}{0.527985in}}{\pgfqpoint{1.192147in}{0.521338in}}%
\pgfpathlineto{\pgfqpoint{1.192147in}{0.521338in}}%
\pgfpathclose%
\pgfpathmoveto{\pgfqpoint{1.284009in}{0.502083in}}%
\pgfpathquadraticcurveto{\pgfqpoint{1.271454in}{0.502083in}}{\pgfqpoint{1.266609in}{0.499219in}}%
\pgfpathquadraticcurveto{\pgfqpoint{1.261764in}{0.496356in}}{\pgfqpoint{1.261764in}{0.489421in}}%
\pgfpathquadraticcurveto{\pgfqpoint{1.261764in}{0.483909in}}{\pgfqpoint{1.265402in}{0.480667in}}%
\pgfpathquadraticcurveto{\pgfqpoint{1.269041in}{0.477443in}}{\pgfqpoint{1.275273in}{0.477443in}}%
\pgfpathquadraticcurveto{\pgfqpoint{1.283901in}{0.477443in}}{\pgfqpoint{1.289106in}{0.483549in}}%
\pgfpathquadraticcurveto{\pgfqpoint{1.294312in}{0.489655in}}{\pgfqpoint{1.294312in}{0.499778in}}%
\pgfpathlineto{\pgfqpoint{1.294312in}{0.502083in}}%
\pgfpathlineto{\pgfqpoint{1.284009in}{0.502083in}}%
\pgfpathlineto{\pgfqpoint{1.284009in}{0.502083in}}%
\pgfpathclose%
\pgfpathmoveto{\pgfqpoint{1.304669in}{0.506370in}}%
\pgfpathlineto{\pgfqpoint{1.304669in}{0.470400in}}%
\pgfpathlineto{\pgfqpoint{1.294312in}{0.470400in}}%
\pgfpathlineto{\pgfqpoint{1.294312in}{0.479964in}}%
\pgfpathquadraticcurveto{\pgfqpoint{1.290763in}{0.474237in}}{\pgfqpoint{1.285468in}{0.471499in}}%
\pgfpathquadraticcurveto{\pgfqpoint{1.280172in}{0.468761in}}{\pgfqpoint{1.272517in}{0.468761in}}%
\pgfpathquadraticcurveto{\pgfqpoint{1.262845in}{0.468761in}}{\pgfqpoint{1.257117in}{0.474201in}}%
\pgfpathquadraticcurveto{\pgfqpoint{1.251407in}{0.479640in}}{\pgfqpoint{1.251407in}{0.488754in}}%
\pgfpathquadraticcurveto{\pgfqpoint{1.251407in}{0.499382in}}{\pgfqpoint{1.258522in}{0.504785in}}%
\pgfpathquadraticcurveto{\pgfqpoint{1.265655in}{0.510189in}}{\pgfqpoint{1.279776in}{0.510189in}}%
\pgfpathlineto{\pgfqpoint{1.294312in}{0.510189in}}%
\pgfpathlineto{\pgfqpoint{1.294312in}{0.511216in}}%
\pgfpathquadraticcurveto{\pgfqpoint{1.294312in}{0.518366in}}{\pgfqpoint{1.289611in}{0.522275in}}%
\pgfpathquadraticcurveto{\pgfqpoint{1.284910in}{0.526184in}}{\pgfqpoint{1.276408in}{0.526184in}}%
\pgfpathquadraticcurveto{\pgfqpoint{1.271004in}{0.526184in}}{\pgfqpoint{1.265871in}{0.524887in}}%
\pgfpathquadraticcurveto{\pgfqpoint{1.260755in}{0.523590in}}{\pgfqpoint{1.256036in}{0.520996in}}%
\pgfpathlineto{\pgfqpoint{1.256036in}{0.530579in}}%
\pgfpathquadraticcurveto{\pgfqpoint{1.261710in}{0.532776in}}{\pgfqpoint{1.267059in}{0.533857in}}%
\pgfpathquadraticcurveto{\pgfqpoint{1.272409in}{0.534956in}}{\pgfqpoint{1.277470in}{0.534956in}}%
\pgfpathquadraticcurveto{\pgfqpoint{1.291160in}{0.534956in}}{\pgfqpoint{1.297914in}{0.527859in}}%
\pgfpathquadraticcurveto{\pgfqpoint{1.304669in}{0.520780in}}{\pgfqpoint{1.304669in}{0.506370in}}%
\pgfpathlineto{\pgfqpoint{1.304669in}{0.506370in}}%
\pgfpathclose%
\pgfpathmoveto{\pgfqpoint{1.378411in}{0.508460in}}%
\pgfpathlineto{\pgfqpoint{1.378411in}{0.470400in}}%
\pgfpathlineto{\pgfqpoint{1.368054in}{0.470400in}}%
\pgfpathlineto{\pgfqpoint{1.368054in}{0.508117in}}%
\pgfpathquadraticcurveto{\pgfqpoint{1.368054in}{0.517069in}}{\pgfqpoint{1.364559in}{0.521500in}}%
\pgfpathquadraticcurveto{\pgfqpoint{1.361065in}{0.525949in}}{\pgfqpoint{1.354094in}{0.525949in}}%
\pgfpathquadraticcurveto{\pgfqpoint{1.345701in}{0.525949in}}{\pgfqpoint{1.340855in}{0.520600in}}%
\pgfpathquadraticcurveto{\pgfqpoint{1.336010in}{0.515268in}}{\pgfqpoint{1.336010in}{0.506028in}}%
\pgfpathlineto{\pgfqpoint{1.336010in}{0.470400in}}%
\pgfpathlineto{\pgfqpoint{1.325599in}{0.470400in}}%
\pgfpathlineto{\pgfqpoint{1.325599in}{0.533443in}}%
\pgfpathlineto{\pgfqpoint{1.336010in}{0.533443in}}%
\pgfpathlineto{\pgfqpoint{1.336010in}{0.523644in}}%
\pgfpathquadraticcurveto{\pgfqpoint{1.339738in}{0.529336in}}{\pgfqpoint{1.344764in}{0.532146in}}%
\pgfpathquadraticcurveto{\pgfqpoint{1.349807in}{0.534956in}}{\pgfqpoint{1.356400in}{0.534956in}}%
\pgfpathquadraticcurveto{\pgfqpoint{1.367261in}{0.534956in}}{\pgfqpoint{1.372827in}{0.528237in}}%
\pgfpathquadraticcurveto{\pgfqpoint{1.378411in}{0.521518in}}{\pgfqpoint{1.378411in}{0.508460in}}%
\pgfpathlineto{\pgfqpoint{1.378411in}{0.508460in}}%
\pgfpathclose%
\pgfpathmoveto{\pgfqpoint{1.440535in}{0.523878in}}%
\pgfpathlineto{\pgfqpoint{1.440535in}{0.557993in}}%
\pgfpathlineto{\pgfqpoint{1.450891in}{0.557993in}}%
\pgfpathlineto{\pgfqpoint{1.450891in}{0.470400in}}%
\pgfpathlineto{\pgfqpoint{1.440535in}{0.470400in}}%
\pgfpathlineto{\pgfqpoint{1.440535in}{0.479856in}}%
\pgfpathquadraticcurveto{\pgfqpoint{1.437274in}{0.474237in}}{\pgfqpoint{1.432285in}{0.471499in}}%
\pgfpathquadraticcurveto{\pgfqpoint{1.427314in}{0.468761in}}{\pgfqpoint{1.420325in}{0.468761in}}%
\pgfpathquadraticcurveto{\pgfqpoint{1.408905in}{0.468761in}}{\pgfqpoint{1.401718in}{0.477875in}}%
\pgfpathquadraticcurveto{\pgfqpoint{1.394549in}{0.487007in}}{\pgfqpoint{1.394549in}{0.501867in}}%
\pgfpathquadraticcurveto{\pgfqpoint{1.394549in}{0.516727in}}{\pgfqpoint{1.401718in}{0.525841in}}%
\pgfpathquadraticcurveto{\pgfqpoint{1.408905in}{0.534956in}}{\pgfqpoint{1.420325in}{0.534956in}}%
\pgfpathquadraticcurveto{\pgfqpoint{1.427314in}{0.534956in}}{\pgfqpoint{1.432285in}{0.532218in}}%
\pgfpathquadraticcurveto{\pgfqpoint{1.437274in}{0.529498in}}{\pgfqpoint{1.440535in}{0.523878in}}%
\pgfpathlineto{\pgfqpoint{1.440535in}{0.523878in}}%
\pgfpathclose%
\pgfpathmoveto{\pgfqpoint{1.405249in}{0.501867in}}%
\pgfpathquadraticcurveto{\pgfqpoint{1.405249in}{0.490448in}}{\pgfqpoint{1.409950in}{0.483945in}}%
\pgfpathquadraticcurveto{\pgfqpoint{1.414651in}{0.477443in}}{\pgfqpoint{1.422865in}{0.477443in}}%
\pgfpathquadraticcurveto{\pgfqpoint{1.431078in}{0.477443in}}{\pgfqpoint{1.435797in}{0.483945in}}%
\pgfpathquadraticcurveto{\pgfqpoint{1.440535in}{0.490448in}}{\pgfqpoint{1.440535in}{0.501867in}}%
\pgfpathquadraticcurveto{\pgfqpoint{1.440535in}{0.513287in}}{\pgfqpoint{1.435797in}{0.519789in}}%
\pgfpathquadraticcurveto{\pgfqpoint{1.431078in}{0.526292in}}{\pgfqpoint{1.422865in}{0.526292in}}%
\pgfpathquadraticcurveto{\pgfqpoint{1.414651in}{0.526292in}}{\pgfqpoint{1.409950in}{0.519789in}}%
\pgfpathquadraticcurveto{\pgfqpoint{1.405249in}{0.513287in}}{\pgfqpoint{1.405249in}{0.501867in}}%
\pgfpathlineto{\pgfqpoint{1.405249in}{0.501867in}}%
\pgfpathclose%
\pgfpathmoveto{\pgfqpoint{1.512421in}{0.531587in}}%
\pgfpathlineto{\pgfqpoint{1.512421in}{0.521789in}}%
\pgfpathquadraticcurveto{\pgfqpoint{1.508044in}{0.524040in}}{\pgfqpoint{1.503307in}{0.525157in}}%
\pgfpathquadraticcurveto{\pgfqpoint{1.498588in}{0.526292in}}{\pgfqpoint{1.493508in}{0.526292in}}%
\pgfpathquadraticcurveto{\pgfqpoint{1.485799in}{0.526292in}}{\pgfqpoint{1.481944in}{0.523932in}}%
\pgfpathquadraticcurveto{\pgfqpoint{1.478090in}{0.521573in}}{\pgfqpoint{1.478090in}{0.516835in}}%
\pgfpathquadraticcurveto{\pgfqpoint{1.478090in}{0.513233in}}{\pgfqpoint{1.480846in}{0.511180in}}%
\pgfpathquadraticcurveto{\pgfqpoint{1.483602in}{0.509126in}}{\pgfqpoint{1.491941in}{0.507271in}}%
\pgfpathlineto{\pgfqpoint{1.495490in}{0.506478in}}%
\pgfpathquadraticcurveto{\pgfqpoint{1.506513in}{0.504119in}}{\pgfqpoint{1.511160in}{0.499814in}}%
\pgfpathquadraticcurveto{\pgfqpoint{1.515807in}{0.495509in}}{\pgfqpoint{1.515807in}{0.487800in}}%
\pgfpathquadraticcurveto{\pgfqpoint{1.515807in}{0.479010in}}{\pgfqpoint{1.508855in}{0.473876in}}%
\pgfpathquadraticcurveto{\pgfqpoint{1.501902in}{0.468761in}}{\pgfqpoint{1.489744in}{0.468761in}}%
\pgfpathquadraticcurveto{\pgfqpoint{1.484682in}{0.468761in}}{\pgfqpoint{1.479189in}{0.469752in}}%
\pgfpathquadraticcurveto{\pgfqpoint{1.473695in}{0.470742in}}{\pgfqpoint{1.467625in}{0.472706in}}%
\pgfpathlineto{\pgfqpoint{1.467625in}{0.483405in}}%
\pgfpathquadraticcurveto{\pgfqpoint{1.473371in}{0.480415in}}{\pgfqpoint{1.478936in}{0.478920in}}%
\pgfpathquadraticcurveto{\pgfqpoint{1.484502in}{0.477443in}}{\pgfqpoint{1.489978in}{0.477443in}}%
\pgfpathquadraticcurveto{\pgfqpoint{1.497291in}{0.477443in}}{\pgfqpoint{1.501217in}{0.479946in}}%
\pgfpathquadraticcurveto{\pgfqpoint{1.505162in}{0.482450in}}{\pgfqpoint{1.505162in}{0.487007in}}%
\pgfpathquadraticcurveto{\pgfqpoint{1.505162in}{0.491222in}}{\pgfqpoint{1.502316in}{0.493474in}}%
\pgfpathquadraticcurveto{\pgfqpoint{1.499488in}{0.495725in}}{\pgfqpoint{1.489852in}{0.497814in}}%
\pgfpathlineto{\pgfqpoint{1.486249in}{0.498661in}}%
\pgfpathquadraticcurveto{\pgfqpoint{1.476631in}{0.500678in}}{\pgfqpoint{1.472344in}{0.504875in}}%
\pgfpathquadraticcurveto{\pgfqpoint{1.468075in}{0.509072in}}{\pgfqpoint{1.468075in}{0.516385in}}%
\pgfpathquadraticcurveto{\pgfqpoint{1.468075in}{0.525283in}}{\pgfqpoint{1.474379in}{0.530110in}}%
\pgfpathquadraticcurveto{\pgfqpoint{1.480684in}{0.534956in}}{\pgfqpoint{1.492283in}{0.534956in}}%
\pgfpathquadraticcurveto{\pgfqpoint{1.498011in}{0.534956in}}{\pgfqpoint{1.503073in}{0.534109in}}%
\pgfpathquadraticcurveto{\pgfqpoint{1.508152in}{0.533280in}}{\pgfqpoint{1.512421in}{0.531587in}}%
\pgfpathlineto{\pgfqpoint{1.512421in}{0.531587in}}%
\pgfpathclose%
\pgfpathmoveto{\pgfqpoint{1.533747in}{0.484702in}}%
\pgfpathlineto{\pgfqpoint{1.545635in}{0.484702in}}%
\pgfpathlineto{\pgfqpoint{1.545635in}{0.470400in}}%
\pgfpathlineto{\pgfqpoint{1.533747in}{0.470400in}}%
\pgfpathlineto{\pgfqpoint{1.533747in}{0.484702in}}%
\pgfpathlineto{\pgfqpoint{1.533747in}{0.484702in}}%
\pgfpathclose%
\pgfusepath{fill}%
\end{pgfscope}%
\begin{pgfscope}%
\pgfsetbuttcap%
\pgfsetroundjoin%
\definecolor{currentfill}{rgb}{0.309804,0.372549,0.419608}%
\pgfsetfillcolor{currentfill}%
\pgfsetlinewidth{0.000000pt}%
\definecolor{currentstroke}{rgb}{0.309804,0.372549,0.419608}%
\pgfsetstrokecolor{currentstroke}%
\pgfsetdash{}{0pt}%
\pgfpathmoveto{\pgfqpoint{1.653990in}{0.510531in}}%
\pgfpathlineto{\pgfqpoint{1.653990in}{0.479748in}}%
\pgfpathlineto{\pgfqpoint{1.672236in}{0.479748in}}%
\pgfpathquadraticcurveto{\pgfqpoint{1.681404in}{0.479748in}}{\pgfqpoint{1.685817in}{0.483549in}}%
\pgfpathquadraticcurveto{\pgfqpoint{1.690248in}{0.487349in}}{\pgfqpoint{1.690248in}{0.495167in}}%
\pgfpathquadraticcurveto{\pgfqpoint{1.690248in}{0.503056in}}{\pgfqpoint{1.685817in}{0.506785in}}%
\pgfpathquadraticcurveto{\pgfqpoint{1.681404in}{0.510531in}}{\pgfqpoint{1.672236in}{0.510531in}}%
\pgfpathlineto{\pgfqpoint{1.653990in}{0.510531in}}%
\pgfpathlineto{\pgfqpoint{1.653990in}{0.510531in}}%
\pgfpathclose%
\pgfpathmoveto{\pgfqpoint{1.653990in}{0.545096in}}%
\pgfpathlineto{\pgfqpoint{1.653990in}{0.519771in}}%
\pgfpathlineto{\pgfqpoint{1.670831in}{0.519771in}}%
\pgfpathquadraticcurveto{\pgfqpoint{1.679153in}{0.519771in}}{\pgfqpoint{1.683224in}{0.522887in}}%
\pgfpathquadraticcurveto{\pgfqpoint{1.687312in}{0.526022in}}{\pgfqpoint{1.687312in}{0.532434in}}%
\pgfpathquadraticcurveto{\pgfqpoint{1.687312in}{0.538792in}}{\pgfqpoint{1.683224in}{0.541944in}}%
\pgfpathquadraticcurveto{\pgfqpoint{1.679153in}{0.545096in}}{\pgfqpoint{1.670831in}{0.545096in}}%
\pgfpathlineto{\pgfqpoint{1.653990in}{0.545096in}}%
\pgfpathlineto{\pgfqpoint{1.653990in}{0.545096in}}%
\pgfpathclose%
\pgfpathmoveto{\pgfqpoint{1.642624in}{0.554445in}}%
\pgfpathlineto{\pgfqpoint{1.671678in}{0.554445in}}%
\pgfpathquadraticcurveto{\pgfqpoint{1.684683in}{0.554445in}}{\pgfqpoint{1.691707in}{0.549041in}}%
\pgfpathquadraticcurveto{\pgfqpoint{1.698750in}{0.543637in}}{\pgfqpoint{1.698750in}{0.533677in}}%
\pgfpathquadraticcurveto{\pgfqpoint{1.698750in}{0.525949in}}{\pgfqpoint{1.695148in}{0.521392in}}%
\pgfpathquadraticcurveto{\pgfqpoint{1.691545in}{0.516835in}}{\pgfqpoint{1.684556in}{0.515719in}}%
\pgfpathquadraticcurveto{\pgfqpoint{1.692950in}{0.513917in}}{\pgfqpoint{1.697597in}{0.508189in}}%
\pgfpathquadraticcurveto{\pgfqpoint{1.702244in}{0.502480in}}{\pgfqpoint{1.702244in}{0.493924in}}%
\pgfpathquadraticcurveto{\pgfqpoint{1.702244in}{0.482666in}}{\pgfqpoint{1.694589in}{0.476524in}}%
\pgfpathquadraticcurveto{\pgfqpoint{1.686934in}{0.470400in}}{\pgfqpoint{1.672794in}{0.470400in}}%
\pgfpathlineto{\pgfqpoint{1.642624in}{0.470400in}}%
\pgfpathlineto{\pgfqpoint{1.642624in}{0.554445in}}%
\pgfpathlineto{\pgfqpoint{1.642624in}{0.554445in}}%
\pgfpathclose%
\pgfpathmoveto{\pgfqpoint{1.745690in}{0.526184in}}%
\pgfpathquadraticcurveto{\pgfqpoint{1.737368in}{0.526184in}}{\pgfqpoint{1.732523in}{0.519681in}}%
\pgfpathquadraticcurveto{\pgfqpoint{1.727678in}{0.513179in}}{\pgfqpoint{1.727678in}{0.501867in}}%
\pgfpathquadraticcurveto{\pgfqpoint{1.727678in}{0.490556in}}{\pgfqpoint{1.732487in}{0.484053in}}%
\pgfpathquadraticcurveto{\pgfqpoint{1.737314in}{0.477551in}}{\pgfqpoint{1.745690in}{0.477551in}}%
\pgfpathquadraticcurveto{\pgfqpoint{1.753975in}{0.477551in}}{\pgfqpoint{1.758803in}{0.484071in}}%
\pgfpathquadraticcurveto{\pgfqpoint{1.763648in}{0.490610in}}{\pgfqpoint{1.763648in}{0.501867in}}%
\pgfpathquadraticcurveto{\pgfqpoint{1.763648in}{0.513071in}}{\pgfqpoint{1.758803in}{0.519627in}}%
\pgfpathquadraticcurveto{\pgfqpoint{1.753975in}{0.526184in}}{\pgfqpoint{1.745690in}{0.526184in}}%
\pgfpathlineto{\pgfqpoint{1.745690in}{0.526184in}}%
\pgfpathclose%
\pgfpathmoveto{\pgfqpoint{1.745690in}{0.534956in}}%
\pgfpathquadraticcurveto{\pgfqpoint{1.759199in}{0.534956in}}{\pgfqpoint{1.766908in}{0.526166in}}%
\pgfpathquadraticcurveto{\pgfqpoint{1.774635in}{0.517394in}}{\pgfqpoint{1.774635in}{0.501867in}}%
\pgfpathquadraticcurveto{\pgfqpoint{1.774635in}{0.486395in}}{\pgfqpoint{1.766908in}{0.477569in}}%
\pgfpathquadraticcurveto{\pgfqpoint{1.759199in}{0.468761in}}{\pgfqpoint{1.745690in}{0.468761in}}%
\pgfpathquadraticcurveto{\pgfqpoint{1.732127in}{0.468761in}}{\pgfqpoint{1.724435in}{0.477569in}}%
\pgfpathquadraticcurveto{\pgfqpoint{1.716762in}{0.486395in}}{\pgfqpoint{1.716762in}{0.501867in}}%
\pgfpathquadraticcurveto{\pgfqpoint{1.716762in}{0.517394in}}{\pgfqpoint{1.724435in}{0.526166in}}%
\pgfpathquadraticcurveto{\pgfqpoint{1.732127in}{0.534956in}}{\pgfqpoint{1.745690in}{0.534956in}}%
\pgfpathlineto{\pgfqpoint{1.745690in}{0.534956in}}%
\pgfpathclose%
\pgfpathmoveto{\pgfqpoint{1.802050in}{0.551347in}}%
\pgfpathlineto{\pgfqpoint{1.802050in}{0.533443in}}%
\pgfpathlineto{\pgfqpoint{1.823376in}{0.533443in}}%
\pgfpathlineto{\pgfqpoint{1.823376in}{0.525391in}}%
\pgfpathlineto{\pgfqpoint{1.802050in}{0.525391in}}%
\pgfpathlineto{\pgfqpoint{1.802050in}{0.491168in}}%
\pgfpathquadraticcurveto{\pgfqpoint{1.802050in}{0.483459in}}{\pgfqpoint{1.804157in}{0.481261in}}%
\pgfpathquadraticcurveto{\pgfqpoint{1.806265in}{0.479064in}}{\pgfqpoint{1.812749in}{0.479064in}}%
\pgfpathlineto{\pgfqpoint{1.823376in}{0.479064in}}%
\pgfpathlineto{\pgfqpoint{1.823376in}{0.470400in}}%
\pgfpathlineto{\pgfqpoint{1.812749in}{0.470400in}}%
\pgfpathquadraticcurveto{\pgfqpoint{1.800753in}{0.470400in}}{\pgfqpoint{1.796196in}{0.474867in}}%
\pgfpathquadraticcurveto{\pgfqpoint{1.791639in}{0.479352in}}{\pgfqpoint{1.791639in}{0.491168in}}%
\pgfpathlineto{\pgfqpoint{1.791639in}{0.525391in}}%
\pgfpathlineto{\pgfqpoint{1.784038in}{0.525391in}}%
\pgfpathlineto{\pgfqpoint{1.784038in}{0.533443in}}%
\pgfpathlineto{\pgfqpoint{1.791639in}{0.533443in}}%
\pgfpathlineto{\pgfqpoint{1.791639in}{0.551347in}}%
\pgfpathlineto{\pgfqpoint{1.802050in}{0.551347in}}%
\pgfpathlineto{\pgfqpoint{1.802050in}{0.551347in}}%
\pgfpathclose%
\pgfpathmoveto{\pgfqpoint{1.889409in}{0.508460in}}%
\pgfpathlineto{\pgfqpoint{1.889409in}{0.470400in}}%
\pgfpathlineto{\pgfqpoint{1.879052in}{0.470400in}}%
\pgfpathlineto{\pgfqpoint{1.879052in}{0.508117in}}%
\pgfpathquadraticcurveto{\pgfqpoint{1.879052in}{0.517069in}}{\pgfqpoint{1.875557in}{0.521500in}}%
\pgfpathquadraticcurveto{\pgfqpoint{1.872063in}{0.525949in}}{\pgfqpoint{1.865092in}{0.525949in}}%
\pgfpathquadraticcurveto{\pgfqpoint{1.856699in}{0.525949in}}{\pgfqpoint{1.851853in}{0.520600in}}%
\pgfpathquadraticcurveto{\pgfqpoint{1.847008in}{0.515268in}}{\pgfqpoint{1.847008in}{0.506028in}}%
\pgfpathlineto{\pgfqpoint{1.847008in}{0.470400in}}%
\pgfpathlineto{\pgfqpoint{1.836597in}{0.470400in}}%
\pgfpathlineto{\pgfqpoint{1.836597in}{0.557993in}}%
\pgfpathlineto{\pgfqpoint{1.847008in}{0.557993in}}%
\pgfpathlineto{\pgfqpoint{1.847008in}{0.523644in}}%
\pgfpathquadraticcurveto{\pgfqpoint{1.850737in}{0.529336in}}{\pgfqpoint{1.855762in}{0.532146in}}%
\pgfpathquadraticcurveto{\pgfqpoint{1.860805in}{0.534956in}}{\pgfqpoint{1.867398in}{0.534956in}}%
\pgfpathquadraticcurveto{\pgfqpoint{1.878259in}{0.534956in}}{\pgfqpoint{1.883825in}{0.528237in}}%
\pgfpathquadraticcurveto{\pgfqpoint{1.889409in}{0.521518in}}{\pgfqpoint{1.889409in}{0.508460in}}%
\pgfpathlineto{\pgfqpoint{1.889409in}{0.508460in}}%
\pgfpathclose%
\pgfusepath{fill}%
\end{pgfscope}%
\begin{pgfscope}%
\pgfsetbuttcap%
\pgfsetroundjoin%
\definecolor{currentfill}{rgb}{0.309804,0.372549,0.419608}%
\pgfsetfillcolor{currentfill}%
\pgfsetlinewidth{0.000000pt}%
\definecolor{currentstroke}{rgb}{0.309804,0.372549,0.419608}%
\pgfsetstrokecolor{currentstroke}%
\pgfsetdash{}{0pt}%
\pgfpathmoveto{\pgfqpoint{1.998144in}{0.502083in}}%
\pgfpathquadraticcurveto{\pgfqpoint{1.985589in}{0.502083in}}{\pgfqpoint{1.980744in}{0.499219in}}%
\pgfpathquadraticcurveto{\pgfqpoint{1.975899in}{0.496356in}}{\pgfqpoint{1.975899in}{0.489421in}}%
\pgfpathquadraticcurveto{\pgfqpoint{1.975899in}{0.483909in}}{\pgfqpoint{1.979537in}{0.480667in}}%
\pgfpathquadraticcurveto{\pgfqpoint{1.983176in}{0.477443in}}{\pgfqpoint{1.989408in}{0.477443in}}%
\pgfpathquadraticcurveto{\pgfqpoint{1.998036in}{0.477443in}}{\pgfqpoint{2.003241in}{0.483549in}}%
\pgfpathquadraticcurveto{\pgfqpoint{2.008447in}{0.489655in}}{\pgfqpoint{2.008447in}{0.499778in}}%
\pgfpathlineto{\pgfqpoint{2.008447in}{0.502083in}}%
\pgfpathlineto{\pgfqpoint{1.998144in}{0.502083in}}%
\pgfpathlineto{\pgfqpoint{1.998144in}{0.502083in}}%
\pgfpathclose%
\pgfpathmoveto{\pgfqpoint{2.018804in}{0.506370in}}%
\pgfpathlineto{\pgfqpoint{2.018804in}{0.470400in}}%
\pgfpathlineto{\pgfqpoint{2.008447in}{0.470400in}}%
\pgfpathlineto{\pgfqpoint{2.008447in}{0.479964in}}%
\pgfpathquadraticcurveto{\pgfqpoint{2.004898in}{0.474237in}}{\pgfqpoint{1.999603in}{0.471499in}}%
\pgfpathquadraticcurveto{\pgfqpoint{1.994307in}{0.468761in}}{\pgfqpoint{1.986652in}{0.468761in}}%
\pgfpathquadraticcurveto{\pgfqpoint{1.976979in}{0.468761in}}{\pgfqpoint{1.971252in}{0.474201in}}%
\pgfpathquadraticcurveto{\pgfqpoint{1.965542in}{0.479640in}}{\pgfqpoint{1.965542in}{0.488754in}}%
\pgfpathquadraticcurveto{\pgfqpoint{1.965542in}{0.499382in}}{\pgfqpoint{1.972656in}{0.504785in}}%
\pgfpathquadraticcurveto{\pgfqpoint{1.979789in}{0.510189in}}{\pgfqpoint{1.993911in}{0.510189in}}%
\pgfpathlineto{\pgfqpoint{2.008447in}{0.510189in}}%
\pgfpathlineto{\pgfqpoint{2.008447in}{0.511216in}}%
\pgfpathquadraticcurveto{\pgfqpoint{2.008447in}{0.518366in}}{\pgfqpoint{2.003745in}{0.522275in}}%
\pgfpathquadraticcurveto{\pgfqpoint{1.999044in}{0.526184in}}{\pgfqpoint{1.990543in}{0.526184in}}%
\pgfpathquadraticcurveto{\pgfqpoint{1.985139in}{0.526184in}}{\pgfqpoint{1.980005in}{0.524887in}}%
\pgfpathquadraticcurveto{\pgfqpoint{1.974890in}{0.523590in}}{\pgfqpoint{1.970171in}{0.520996in}}%
\pgfpathlineto{\pgfqpoint{1.970171in}{0.530579in}}%
\pgfpathquadraticcurveto{\pgfqpoint{1.975845in}{0.532776in}}{\pgfqpoint{1.981194in}{0.533857in}}%
\pgfpathquadraticcurveto{\pgfqpoint{1.986544in}{0.534956in}}{\pgfqpoint{1.991605in}{0.534956in}}%
\pgfpathquadraticcurveto{\pgfqpoint{2.005294in}{0.534956in}}{\pgfqpoint{2.012049in}{0.527859in}}%
\pgfpathquadraticcurveto{\pgfqpoint{2.018804in}{0.520780in}}{\pgfqpoint{2.018804in}{0.506370in}}%
\pgfpathlineto{\pgfqpoint{2.018804in}{0.506370in}}%
\pgfpathclose%
\pgfpathmoveto{\pgfqpoint{2.092545in}{0.533443in}}%
\pgfpathlineto{\pgfqpoint{2.069742in}{0.502768in}}%
\pgfpathlineto{\pgfqpoint{2.093716in}{0.470400in}}%
\pgfpathlineto{\pgfqpoint{2.081504in}{0.470400in}}%
\pgfpathlineto{\pgfqpoint{2.063150in}{0.495167in}}%
\pgfpathlineto{\pgfqpoint{2.044813in}{0.470400in}}%
\pgfpathlineto{\pgfqpoint{2.032583in}{0.470400in}}%
\pgfpathlineto{\pgfqpoint{2.057079in}{0.503380in}}%
\pgfpathlineto{\pgfqpoint{2.034672in}{0.533443in}}%
\pgfpathlineto{\pgfqpoint{2.046885in}{0.533443in}}%
\pgfpathlineto{\pgfqpoint{2.063600in}{0.510981in}}%
\pgfpathlineto{\pgfqpoint{2.080315in}{0.533443in}}%
\pgfpathlineto{\pgfqpoint{2.092545in}{0.533443in}}%
\pgfpathlineto{\pgfqpoint{2.092545in}{0.533443in}}%
\pgfpathclose%
\pgfpathmoveto{\pgfqpoint{2.158740in}{0.504515in}}%
\pgfpathlineto{\pgfqpoint{2.158740in}{0.499454in}}%
\pgfpathlineto{\pgfqpoint{2.111116in}{0.499454in}}%
\pgfpathquadraticcurveto{\pgfqpoint{2.111800in}{0.488754in}}{\pgfqpoint{2.117564in}{0.483153in}}%
\pgfpathquadraticcurveto{\pgfqpoint{2.123328in}{0.477551in}}{\pgfqpoint{2.133631in}{0.477551in}}%
\pgfpathquadraticcurveto{\pgfqpoint{2.139593in}{0.477551in}}{\pgfqpoint{2.145195in}{0.479010in}}%
\pgfpathquadraticcurveto{\pgfqpoint{2.150797in}{0.480469in}}{\pgfqpoint{2.156326in}{0.483405in}}%
\pgfpathlineto{\pgfqpoint{2.156326in}{0.473606in}}%
\pgfpathquadraticcurveto{\pgfqpoint{2.150743in}{0.471247in}}{\pgfqpoint{2.144889in}{0.470004in}}%
\pgfpathquadraticcurveto{\pgfqpoint{2.139035in}{0.468761in}}{\pgfqpoint{2.133019in}{0.468761in}}%
\pgfpathquadraticcurveto{\pgfqpoint{2.117924in}{0.468761in}}{\pgfqpoint{2.109117in}{0.477533in}}%
\pgfpathquadraticcurveto{\pgfqpoint{2.100309in}{0.486323in}}{\pgfqpoint{2.100309in}{0.501309in}}%
\pgfpathquadraticcurveto{\pgfqpoint{2.100309in}{0.516781in}}{\pgfqpoint{2.108666in}{0.525859in}}%
\pgfpathquadraticcurveto{\pgfqpoint{2.117024in}{0.534956in}}{\pgfqpoint{2.131217in}{0.534956in}}%
\pgfpathquadraticcurveto{\pgfqpoint{2.143934in}{0.534956in}}{\pgfqpoint{2.151337in}{0.526760in}}%
\pgfpathquadraticcurveto{\pgfqpoint{2.158740in}{0.518583in}}{\pgfqpoint{2.158740in}{0.504515in}}%
\pgfpathlineto{\pgfqpoint{2.158740in}{0.504515in}}%
\pgfpathclose%
\pgfpathmoveto{\pgfqpoint{2.148383in}{0.507559in}}%
\pgfpathquadraticcurveto{\pgfqpoint{2.148275in}{0.516043in}}{\pgfqpoint{2.143628in}{0.521104in}}%
\pgfpathquadraticcurveto{\pgfqpoint{2.138981in}{0.526184in}}{\pgfqpoint{2.131326in}{0.526184in}}%
\pgfpathquadraticcurveto{\pgfqpoint{2.122662in}{0.526184in}}{\pgfqpoint{2.117456in}{0.521284in}}%
\pgfpathquadraticcurveto{\pgfqpoint{2.112251in}{0.516385in}}{\pgfqpoint{2.111458in}{0.507487in}}%
\pgfpathlineto{\pgfqpoint{2.148383in}{0.507559in}}%
\pgfpathlineto{\pgfqpoint{2.148383in}{0.507559in}}%
\pgfpathclose%
\pgfpathmoveto{\pgfqpoint{2.215929in}{0.531587in}}%
\pgfpathlineto{\pgfqpoint{2.215929in}{0.521789in}}%
\pgfpathquadraticcurveto{\pgfqpoint{2.211552in}{0.524040in}}{\pgfqpoint{2.206814in}{0.525157in}}%
\pgfpathquadraticcurveto{\pgfqpoint{2.202095in}{0.526292in}}{\pgfqpoint{2.197016in}{0.526292in}}%
\pgfpathquadraticcurveto{\pgfqpoint{2.189307in}{0.526292in}}{\pgfqpoint{2.185452in}{0.523932in}}%
\pgfpathquadraticcurveto{\pgfqpoint{2.181597in}{0.521573in}}{\pgfqpoint{2.181597in}{0.516835in}}%
\pgfpathquadraticcurveto{\pgfqpoint{2.181597in}{0.513233in}}{\pgfqpoint{2.184353in}{0.511180in}}%
\pgfpathquadraticcurveto{\pgfqpoint{2.187109in}{0.509126in}}{\pgfqpoint{2.195449in}{0.507271in}}%
\pgfpathlineto{\pgfqpoint{2.198997in}{0.506478in}}%
\pgfpathquadraticcurveto{\pgfqpoint{2.210021in}{0.504119in}}{\pgfqpoint{2.214668in}{0.499814in}}%
\pgfpathquadraticcurveto{\pgfqpoint{2.219315in}{0.495509in}}{\pgfqpoint{2.219315in}{0.487800in}}%
\pgfpathquadraticcurveto{\pgfqpoint{2.219315in}{0.479010in}}{\pgfqpoint{2.212362in}{0.473876in}}%
\pgfpathquadraticcurveto{\pgfqpoint{2.205410in}{0.468761in}}{\pgfqpoint{2.193251in}{0.468761in}}%
\pgfpathquadraticcurveto{\pgfqpoint{2.188190in}{0.468761in}}{\pgfqpoint{2.182696in}{0.469752in}}%
\pgfpathquadraticcurveto{\pgfqpoint{2.177202in}{0.470742in}}{\pgfqpoint{2.171132in}{0.472706in}}%
\pgfpathlineto{\pgfqpoint{2.171132in}{0.483405in}}%
\pgfpathquadraticcurveto{\pgfqpoint{2.176878in}{0.480415in}}{\pgfqpoint{2.182444in}{0.478920in}}%
\pgfpathquadraticcurveto{\pgfqpoint{2.188010in}{0.477443in}}{\pgfqpoint{2.193485in}{0.477443in}}%
\pgfpathquadraticcurveto{\pgfqpoint{2.200798in}{0.477443in}}{\pgfqpoint{2.204725in}{0.479946in}}%
\pgfpathquadraticcurveto{\pgfqpoint{2.208670in}{0.482450in}}{\pgfqpoint{2.208670in}{0.487007in}}%
\pgfpathquadraticcurveto{\pgfqpoint{2.208670in}{0.491222in}}{\pgfqpoint{2.205824in}{0.493474in}}%
\pgfpathquadraticcurveto{\pgfqpoint{2.202996in}{0.495725in}}{\pgfqpoint{2.193359in}{0.497814in}}%
\pgfpathlineto{\pgfqpoint{2.189757in}{0.498661in}}%
\pgfpathquadraticcurveto{\pgfqpoint{2.180138in}{0.500678in}}{\pgfqpoint{2.175852in}{0.504875in}}%
\pgfpathquadraticcurveto{\pgfqpoint{2.171583in}{0.509072in}}{\pgfqpoint{2.171583in}{0.516385in}}%
\pgfpathquadraticcurveto{\pgfqpoint{2.171583in}{0.525283in}}{\pgfqpoint{2.177887in}{0.530110in}}%
\pgfpathquadraticcurveto{\pgfqpoint{2.184191in}{0.534956in}}{\pgfqpoint{2.195791in}{0.534956in}}%
\pgfpathquadraticcurveto{\pgfqpoint{2.201519in}{0.534956in}}{\pgfqpoint{2.206580in}{0.534109in}}%
\pgfpathquadraticcurveto{\pgfqpoint{2.211660in}{0.533280in}}{\pgfqpoint{2.215929in}{0.531587in}}%
\pgfpathlineto{\pgfqpoint{2.215929in}{0.531587in}}%
\pgfpathclose%
\pgfusepath{fill}%
\end{pgfscope}%
\begin{pgfscope}%
\pgfsetbuttcap%
\pgfsetroundjoin%
\definecolor{currentfill}{rgb}{0.309804,0.372549,0.419608}%
\pgfsetfillcolor{currentfill}%
\pgfsetlinewidth{0.000000pt}%
\definecolor{currentstroke}{rgb}{0.309804,0.372549,0.419608}%
\pgfsetstrokecolor{currentstroke}%
\pgfsetdash{}{0pt}%
\pgfpathmoveto{\pgfqpoint{2.324687in}{0.523878in}}%
\pgfpathlineto{\pgfqpoint{2.324687in}{0.557993in}}%
\pgfpathlineto{\pgfqpoint{2.335044in}{0.557993in}}%
\pgfpathlineto{\pgfqpoint{2.335044in}{0.470400in}}%
\pgfpathlineto{\pgfqpoint{2.324687in}{0.470400in}}%
\pgfpathlineto{\pgfqpoint{2.324687in}{0.479856in}}%
\pgfpathquadraticcurveto{\pgfqpoint{2.321427in}{0.474237in}}{\pgfqpoint{2.316438in}{0.471499in}}%
\pgfpathquadraticcurveto{\pgfqpoint{2.311466in}{0.468761in}}{\pgfqpoint{2.304477in}{0.468761in}}%
\pgfpathquadraticcurveto{\pgfqpoint{2.293058in}{0.468761in}}{\pgfqpoint{2.285871in}{0.477875in}}%
\pgfpathquadraticcurveto{\pgfqpoint{2.278702in}{0.487007in}}{\pgfqpoint{2.278702in}{0.501867in}}%
\pgfpathquadraticcurveto{\pgfqpoint{2.278702in}{0.516727in}}{\pgfqpoint{2.285871in}{0.525841in}}%
\pgfpathquadraticcurveto{\pgfqpoint{2.293058in}{0.534956in}}{\pgfqpoint{2.304477in}{0.534956in}}%
\pgfpathquadraticcurveto{\pgfqpoint{2.311466in}{0.534956in}}{\pgfqpoint{2.316438in}{0.532218in}}%
\pgfpathquadraticcurveto{\pgfqpoint{2.321427in}{0.529498in}}{\pgfqpoint{2.324687in}{0.523878in}}%
\pgfpathlineto{\pgfqpoint{2.324687in}{0.523878in}}%
\pgfpathclose%
\pgfpathmoveto{\pgfqpoint{2.289401in}{0.501867in}}%
\pgfpathquadraticcurveto{\pgfqpoint{2.289401in}{0.490448in}}{\pgfqpoint{2.294102in}{0.483945in}}%
\pgfpathquadraticcurveto{\pgfqpoint{2.298804in}{0.477443in}}{\pgfqpoint{2.307017in}{0.477443in}}%
\pgfpathquadraticcurveto{\pgfqpoint{2.315231in}{0.477443in}}{\pgfqpoint{2.319950in}{0.483945in}}%
\pgfpathquadraticcurveto{\pgfqpoint{2.324687in}{0.490448in}}{\pgfqpoint{2.324687in}{0.501867in}}%
\pgfpathquadraticcurveto{\pgfqpoint{2.324687in}{0.513287in}}{\pgfqpoint{2.319950in}{0.519789in}}%
\pgfpathquadraticcurveto{\pgfqpoint{2.315231in}{0.526292in}}{\pgfqpoint{2.307017in}{0.526292in}}%
\pgfpathquadraticcurveto{\pgfqpoint{2.298804in}{0.526292in}}{\pgfqpoint{2.294102in}{0.519789in}}%
\pgfpathquadraticcurveto{\pgfqpoint{2.289401in}{0.513287in}}{\pgfqpoint{2.289401in}{0.501867in}}%
\pgfpathlineto{\pgfqpoint{2.289401in}{0.501867in}}%
\pgfpathclose%
\pgfpathmoveto{\pgfqpoint{2.356388in}{0.533443in}}%
\pgfpathlineto{\pgfqpoint{2.366745in}{0.533443in}}%
\pgfpathlineto{\pgfqpoint{2.366745in}{0.470400in}}%
\pgfpathlineto{\pgfqpoint{2.356388in}{0.470400in}}%
\pgfpathlineto{\pgfqpoint{2.356388in}{0.533443in}}%
\pgfpathlineto{\pgfqpoint{2.356388in}{0.533443in}}%
\pgfpathclose%
\pgfpathmoveto{\pgfqpoint{2.356388in}{0.557993in}}%
\pgfpathlineto{\pgfqpoint{2.366745in}{0.557993in}}%
\pgfpathlineto{\pgfqpoint{2.366745in}{0.544862in}}%
\pgfpathlineto{\pgfqpoint{2.356388in}{0.544862in}}%
\pgfpathlineto{\pgfqpoint{2.356388in}{0.557993in}}%
\pgfpathlineto{\pgfqpoint{2.356388in}{0.557993in}}%
\pgfpathclose%
\pgfpathmoveto{\pgfqpoint{2.380993in}{0.533443in}}%
\pgfpathlineto{\pgfqpoint{2.391962in}{0.533443in}}%
\pgfpathlineto{\pgfqpoint{2.411668in}{0.480541in}}%
\pgfpathlineto{\pgfqpoint{2.431373in}{0.533443in}}%
\pgfpathlineto{\pgfqpoint{2.442342in}{0.533443in}}%
\pgfpathlineto{\pgfqpoint{2.418692in}{0.470400in}}%
\pgfpathlineto{\pgfqpoint{2.404625in}{0.470400in}}%
\pgfpathlineto{\pgfqpoint{2.380993in}{0.533443in}}%
\pgfpathlineto{\pgfqpoint{2.380993in}{0.533443in}}%
\pgfpathclose%
\pgfpathmoveto{\pgfqpoint{2.456644in}{0.533443in}}%
\pgfpathlineto{\pgfqpoint{2.467001in}{0.533443in}}%
\pgfpathlineto{\pgfqpoint{2.467001in}{0.470400in}}%
\pgfpathlineto{\pgfqpoint{2.456644in}{0.470400in}}%
\pgfpathlineto{\pgfqpoint{2.456644in}{0.533443in}}%
\pgfpathlineto{\pgfqpoint{2.456644in}{0.533443in}}%
\pgfpathclose%
\pgfpathmoveto{\pgfqpoint{2.456644in}{0.557993in}}%
\pgfpathlineto{\pgfqpoint{2.467001in}{0.557993in}}%
\pgfpathlineto{\pgfqpoint{2.467001in}{0.544862in}}%
\pgfpathlineto{\pgfqpoint{2.456644in}{0.544862in}}%
\pgfpathlineto{\pgfqpoint{2.456644in}{0.557993in}}%
\pgfpathlineto{\pgfqpoint{2.456644in}{0.557993in}}%
\pgfpathclose%
\pgfpathmoveto{\pgfqpoint{2.530152in}{0.523878in}}%
\pgfpathlineto{\pgfqpoint{2.530152in}{0.557993in}}%
\pgfpathlineto{\pgfqpoint{2.540509in}{0.557993in}}%
\pgfpathlineto{\pgfqpoint{2.540509in}{0.470400in}}%
\pgfpathlineto{\pgfqpoint{2.530152in}{0.470400in}}%
\pgfpathlineto{\pgfqpoint{2.530152in}{0.479856in}}%
\pgfpathquadraticcurveto{\pgfqpoint{2.526891in}{0.474237in}}{\pgfqpoint{2.521902in}{0.471499in}}%
\pgfpathquadraticcurveto{\pgfqpoint{2.516931in}{0.468761in}}{\pgfqpoint{2.509942in}{0.468761in}}%
\pgfpathquadraticcurveto{\pgfqpoint{2.498522in}{0.468761in}}{\pgfqpoint{2.491336in}{0.477875in}}%
\pgfpathquadraticcurveto{\pgfqpoint{2.484167in}{0.487007in}}{\pgfqpoint{2.484167in}{0.501867in}}%
\pgfpathquadraticcurveto{\pgfqpoint{2.484167in}{0.516727in}}{\pgfqpoint{2.491336in}{0.525841in}}%
\pgfpathquadraticcurveto{\pgfqpoint{2.498522in}{0.534956in}}{\pgfqpoint{2.509942in}{0.534956in}}%
\pgfpathquadraticcurveto{\pgfqpoint{2.516931in}{0.534956in}}{\pgfqpoint{2.521902in}{0.532218in}}%
\pgfpathquadraticcurveto{\pgfqpoint{2.526891in}{0.529498in}}{\pgfqpoint{2.530152in}{0.523878in}}%
\pgfpathlineto{\pgfqpoint{2.530152in}{0.523878in}}%
\pgfpathclose%
\pgfpathmoveto{\pgfqpoint{2.494866in}{0.501867in}}%
\pgfpathquadraticcurveto{\pgfqpoint{2.494866in}{0.490448in}}{\pgfqpoint{2.499567in}{0.483945in}}%
\pgfpathquadraticcurveto{\pgfqpoint{2.504268in}{0.477443in}}{\pgfqpoint{2.512482in}{0.477443in}}%
\pgfpathquadraticcurveto{\pgfqpoint{2.520695in}{0.477443in}}{\pgfqpoint{2.525414in}{0.483945in}}%
\pgfpathquadraticcurveto{\pgfqpoint{2.530152in}{0.490448in}}{\pgfqpoint{2.530152in}{0.501867in}}%
\pgfpathquadraticcurveto{\pgfqpoint{2.530152in}{0.513287in}}{\pgfqpoint{2.525414in}{0.519789in}}%
\pgfpathquadraticcurveto{\pgfqpoint{2.520695in}{0.526292in}}{\pgfqpoint{2.512482in}{0.526292in}}%
\pgfpathquadraticcurveto{\pgfqpoint{2.504268in}{0.526292in}}{\pgfqpoint{2.499567in}{0.519789in}}%
\pgfpathquadraticcurveto{\pgfqpoint{2.494866in}{0.513287in}}{\pgfqpoint{2.494866in}{0.501867in}}%
\pgfpathlineto{\pgfqpoint{2.494866in}{0.501867in}}%
\pgfpathclose%
\pgfpathmoveto{\pgfqpoint{2.615781in}{0.504515in}}%
\pgfpathlineto{\pgfqpoint{2.615781in}{0.499454in}}%
\pgfpathlineto{\pgfqpoint{2.568157in}{0.499454in}}%
\pgfpathquadraticcurveto{\pgfqpoint{2.568842in}{0.488754in}}{\pgfqpoint{2.574606in}{0.483153in}}%
\pgfpathquadraticcurveto{\pgfqpoint{2.580370in}{0.477551in}}{\pgfqpoint{2.590673in}{0.477551in}}%
\pgfpathquadraticcurveto{\pgfqpoint{2.596635in}{0.477551in}}{\pgfqpoint{2.602236in}{0.479010in}}%
\pgfpathquadraticcurveto{\pgfqpoint{2.607838in}{0.480469in}}{\pgfqpoint{2.613368in}{0.483405in}}%
\pgfpathlineto{\pgfqpoint{2.613368in}{0.473606in}}%
\pgfpathquadraticcurveto{\pgfqpoint{2.607784in}{0.471247in}}{\pgfqpoint{2.601930in}{0.470004in}}%
\pgfpathquadraticcurveto{\pgfqpoint{2.596076in}{0.468761in}}{\pgfqpoint{2.590060in}{0.468761in}}%
\pgfpathquadraticcurveto{\pgfqpoint{2.574966in}{0.468761in}}{\pgfqpoint{2.566158in}{0.477533in}}%
\pgfpathquadraticcurveto{\pgfqpoint{2.557350in}{0.486323in}}{\pgfqpoint{2.557350in}{0.501309in}}%
\pgfpathquadraticcurveto{\pgfqpoint{2.557350in}{0.516781in}}{\pgfqpoint{2.565708in}{0.525859in}}%
\pgfpathquadraticcurveto{\pgfqpoint{2.574065in}{0.534956in}}{\pgfqpoint{2.588259in}{0.534956in}}%
\pgfpathquadraticcurveto{\pgfqpoint{2.600975in}{0.534956in}}{\pgfqpoint{2.608378in}{0.526760in}}%
\pgfpathquadraticcurveto{\pgfqpoint{2.615781in}{0.518583in}}{\pgfqpoint{2.615781in}{0.504515in}}%
\pgfpathlineto{\pgfqpoint{2.615781in}{0.504515in}}%
\pgfpathclose%
\pgfpathmoveto{\pgfqpoint{2.605424in}{0.507559in}}%
\pgfpathquadraticcurveto{\pgfqpoint{2.605316in}{0.516043in}}{\pgfqpoint{2.600669in}{0.521104in}}%
\pgfpathquadraticcurveto{\pgfqpoint{2.596022in}{0.526184in}}{\pgfqpoint{2.588367in}{0.526184in}}%
\pgfpathquadraticcurveto{\pgfqpoint{2.579703in}{0.526184in}}{\pgfqpoint{2.574498in}{0.521284in}}%
\pgfpathquadraticcurveto{\pgfqpoint{2.569292in}{0.516385in}}{\pgfqpoint{2.568500in}{0.507487in}}%
\pgfpathlineto{\pgfqpoint{2.605424in}{0.507559in}}%
\pgfpathlineto{\pgfqpoint{2.605424in}{0.507559in}}%
\pgfpathclose%
\pgfusepath{fill}%
\end{pgfscope}%
\begin{pgfscope}%
\pgfsetbuttcap%
\pgfsetroundjoin%
\definecolor{currentfill}{rgb}{0.309804,0.372549,0.419608}%
\pgfsetfillcolor{currentfill}%
\pgfsetlinewidth{0.000000pt}%
\definecolor{currentstroke}{rgb}{0.309804,0.372549,0.419608}%
\pgfsetstrokecolor{currentstroke}%
\pgfsetdash{}{0pt}%
\pgfpathmoveto{\pgfqpoint{2.735782in}{0.501867in}}%
\pgfpathquadraticcurveto{\pgfqpoint{2.735782in}{0.513287in}}{\pgfqpoint{2.731081in}{0.519789in}}%
\pgfpathquadraticcurveto{\pgfqpoint{2.726380in}{0.526292in}}{\pgfqpoint{2.718166in}{0.526292in}}%
\pgfpathquadraticcurveto{\pgfqpoint{2.709935in}{0.526292in}}{\pgfqpoint{2.705234in}{0.519789in}}%
\pgfpathquadraticcurveto{\pgfqpoint{2.700532in}{0.513287in}}{\pgfqpoint{2.700532in}{0.501867in}}%
\pgfpathquadraticcurveto{\pgfqpoint{2.700532in}{0.490448in}}{\pgfqpoint{2.705234in}{0.483945in}}%
\pgfpathquadraticcurveto{\pgfqpoint{2.709935in}{0.477443in}}{\pgfqpoint{2.718166in}{0.477443in}}%
\pgfpathquadraticcurveto{\pgfqpoint{2.726380in}{0.477443in}}{\pgfqpoint{2.731081in}{0.483945in}}%
\pgfpathquadraticcurveto{\pgfqpoint{2.735782in}{0.490448in}}{\pgfqpoint{2.735782in}{0.501867in}}%
\pgfpathlineto{\pgfqpoint{2.735782in}{0.501867in}}%
\pgfpathclose%
\pgfpathmoveto{\pgfqpoint{2.700532in}{0.523878in}}%
\pgfpathquadraticcurveto{\pgfqpoint{2.703811in}{0.529498in}}{\pgfqpoint{2.708782in}{0.532218in}}%
\pgfpathquadraticcurveto{\pgfqpoint{2.713771in}{0.534956in}}{\pgfqpoint{2.720688in}{0.534956in}}%
\pgfpathquadraticcurveto{\pgfqpoint{2.732180in}{0.534956in}}{\pgfqpoint{2.739349in}{0.525841in}}%
\pgfpathquadraticcurveto{\pgfqpoint{2.746535in}{0.516727in}}{\pgfqpoint{2.746535in}{0.501867in}}%
\pgfpathquadraticcurveto{\pgfqpoint{2.746535in}{0.487007in}}{\pgfqpoint{2.739349in}{0.477875in}}%
\pgfpathquadraticcurveto{\pgfqpoint{2.732180in}{0.468761in}}{\pgfqpoint{2.720688in}{0.468761in}}%
\pgfpathquadraticcurveto{\pgfqpoint{2.713771in}{0.468761in}}{\pgfqpoint{2.708782in}{0.471499in}}%
\pgfpathquadraticcurveto{\pgfqpoint{2.703811in}{0.474237in}}{\pgfqpoint{2.700532in}{0.479856in}}%
\pgfpathlineto{\pgfqpoint{2.700532in}{0.470400in}}%
\pgfpathlineto{\pgfqpoint{2.690121in}{0.470400in}}%
\pgfpathlineto{\pgfqpoint{2.690121in}{0.557993in}}%
\pgfpathlineto{\pgfqpoint{2.700532in}{0.557993in}}%
\pgfpathlineto{\pgfqpoint{2.700532in}{0.523878in}}%
\pgfpathlineto{\pgfqpoint{2.700532in}{0.523878in}}%
\pgfpathclose%
\pgfpathmoveto{\pgfqpoint{2.789927in}{0.464546in}}%
\pgfpathquadraticcurveto{\pgfqpoint{2.785550in}{0.453288in}}{\pgfqpoint{2.781371in}{0.449866in}}%
\pgfpathquadraticcurveto{\pgfqpoint{2.777210in}{0.446426in}}{\pgfqpoint{2.770239in}{0.446426in}}%
\pgfpathlineto{\pgfqpoint{2.761954in}{0.446426in}}%
\pgfpathlineto{\pgfqpoint{2.761954in}{0.455090in}}%
\pgfpathlineto{\pgfqpoint{2.768042in}{0.455090in}}%
\pgfpathquadraticcurveto{\pgfqpoint{2.772311in}{0.455090in}}{\pgfqpoint{2.774670in}{0.457125in}}%
\pgfpathquadraticcurveto{\pgfqpoint{2.777048in}{0.459142in}}{\pgfqpoint{2.779912in}{0.466689in}}%
\pgfpathlineto{\pgfqpoint{2.781767in}{0.471409in}}%
\pgfpathlineto{\pgfqpoint{2.756280in}{0.533443in}}%
\pgfpathlineto{\pgfqpoint{2.767249in}{0.533443in}}%
\pgfpathlineto{\pgfqpoint{2.786955in}{0.484143in}}%
\pgfpathlineto{\pgfqpoint{2.806660in}{0.533443in}}%
\pgfpathlineto{\pgfqpoint{2.817629in}{0.533443in}}%
\pgfpathlineto{\pgfqpoint{2.789927in}{0.464546in}}%
\pgfpathlineto{\pgfqpoint{2.789927in}{0.464546in}}%
\pgfpathclose%
\pgfusepath{fill}%
\end{pgfscope}%
\begin{pgfscope}%
\pgfsetbuttcap%
\pgfsetroundjoin%
\definecolor{currentfill}{rgb}{0.309804,0.372549,0.419608}%
\pgfsetfillcolor{currentfill}%
\pgfsetlinewidth{0.000000pt}%
\definecolor{currentstroke}{rgb}{0.309804,0.372549,0.419608}%
\pgfsetstrokecolor{currentstroke}%
\pgfsetdash{}{0pt}%
\pgfpathmoveto{\pgfqpoint{2.898079in}{0.551347in}}%
\pgfpathlineto{\pgfqpoint{2.898079in}{0.533443in}}%
\pgfpathlineto{\pgfqpoint{2.919405in}{0.533443in}}%
\pgfpathlineto{\pgfqpoint{2.919405in}{0.525391in}}%
\pgfpathlineto{\pgfqpoint{2.898079in}{0.525391in}}%
\pgfpathlineto{\pgfqpoint{2.898079in}{0.491168in}}%
\pgfpathquadraticcurveto{\pgfqpoint{2.898079in}{0.483459in}}{\pgfqpoint{2.900186in}{0.481261in}}%
\pgfpathquadraticcurveto{\pgfqpoint{2.902294in}{0.479064in}}{\pgfqpoint{2.908778in}{0.479064in}}%
\pgfpathlineto{\pgfqpoint{2.919405in}{0.479064in}}%
\pgfpathlineto{\pgfqpoint{2.919405in}{0.470400in}}%
\pgfpathlineto{\pgfqpoint{2.908778in}{0.470400in}}%
\pgfpathquadraticcurveto{\pgfqpoint{2.896782in}{0.470400in}}{\pgfqpoint{2.892225in}{0.474867in}}%
\pgfpathquadraticcurveto{\pgfqpoint{2.887668in}{0.479352in}}{\pgfqpoint{2.887668in}{0.491168in}}%
\pgfpathlineto{\pgfqpoint{2.887668in}{0.525391in}}%
\pgfpathlineto{\pgfqpoint{2.880067in}{0.525391in}}%
\pgfpathlineto{\pgfqpoint{2.880067in}{0.533443in}}%
\pgfpathlineto{\pgfqpoint{2.887668in}{0.533443in}}%
\pgfpathlineto{\pgfqpoint{2.887668in}{0.551347in}}%
\pgfpathlineto{\pgfqpoint{2.898079in}{0.551347in}}%
\pgfpathlineto{\pgfqpoint{2.898079in}{0.551347in}}%
\pgfpathclose%
\pgfpathmoveto{\pgfqpoint{2.985438in}{0.508460in}}%
\pgfpathlineto{\pgfqpoint{2.985438in}{0.470400in}}%
\pgfpathlineto{\pgfqpoint{2.975081in}{0.470400in}}%
\pgfpathlineto{\pgfqpoint{2.975081in}{0.508117in}}%
\pgfpathquadraticcurveto{\pgfqpoint{2.975081in}{0.517069in}}{\pgfqpoint{2.971587in}{0.521500in}}%
\pgfpathquadraticcurveto{\pgfqpoint{2.968092in}{0.525949in}}{\pgfqpoint{2.961122in}{0.525949in}}%
\pgfpathquadraticcurveto{\pgfqpoint{2.952728in}{0.525949in}}{\pgfqpoint{2.947883in}{0.520600in}}%
\pgfpathquadraticcurveto{\pgfqpoint{2.943037in}{0.515268in}}{\pgfqpoint{2.943037in}{0.506028in}}%
\pgfpathlineto{\pgfqpoint{2.943037in}{0.470400in}}%
\pgfpathlineto{\pgfqpoint{2.932626in}{0.470400in}}%
\pgfpathlineto{\pgfqpoint{2.932626in}{0.557993in}}%
\pgfpathlineto{\pgfqpoint{2.943037in}{0.557993in}}%
\pgfpathlineto{\pgfqpoint{2.943037in}{0.523644in}}%
\pgfpathquadraticcurveto{\pgfqpoint{2.946766in}{0.529336in}}{\pgfqpoint{2.951791in}{0.532146in}}%
\pgfpathquadraticcurveto{\pgfqpoint{2.956835in}{0.534956in}}{\pgfqpoint{2.963427in}{0.534956in}}%
\pgfpathquadraticcurveto{\pgfqpoint{2.974288in}{0.534956in}}{\pgfqpoint{2.979854in}{0.528237in}}%
\pgfpathquadraticcurveto{\pgfqpoint{2.985438in}{0.521518in}}{\pgfqpoint{2.985438in}{0.508460in}}%
\pgfpathlineto{\pgfqpoint{2.985438in}{0.508460in}}%
\pgfpathclose%
\pgfpathmoveto{\pgfqpoint{3.006080in}{0.533443in}}%
\pgfpathlineto{\pgfqpoint{3.016437in}{0.533443in}}%
\pgfpathlineto{\pgfqpoint{3.016437in}{0.470400in}}%
\pgfpathlineto{\pgfqpoint{3.006080in}{0.470400in}}%
\pgfpathlineto{\pgfqpoint{3.006080in}{0.533443in}}%
\pgfpathlineto{\pgfqpoint{3.006080in}{0.533443in}}%
\pgfpathclose%
\pgfpathmoveto{\pgfqpoint{3.006080in}{0.557993in}}%
\pgfpathlineto{\pgfqpoint{3.016437in}{0.557993in}}%
\pgfpathlineto{\pgfqpoint{3.016437in}{0.544862in}}%
\pgfpathlineto{\pgfqpoint{3.006080in}{0.544862in}}%
\pgfpathlineto{\pgfqpoint{3.006080in}{0.557993in}}%
\pgfpathlineto{\pgfqpoint{3.006080in}{0.557993in}}%
\pgfpathclose%
\pgfpathmoveto{\pgfqpoint{3.078291in}{0.531587in}}%
\pgfpathlineto{\pgfqpoint{3.078291in}{0.521789in}}%
\pgfpathquadraticcurveto{\pgfqpoint{3.073914in}{0.524040in}}{\pgfqpoint{3.069176in}{0.525157in}}%
\pgfpathquadraticcurveto{\pgfqpoint{3.064457in}{0.526292in}}{\pgfqpoint{3.059378in}{0.526292in}}%
\pgfpathquadraticcurveto{\pgfqpoint{3.051669in}{0.526292in}}{\pgfqpoint{3.047814in}{0.523932in}}%
\pgfpathquadraticcurveto{\pgfqpoint{3.043959in}{0.521573in}}{\pgfqpoint{3.043959in}{0.516835in}}%
\pgfpathquadraticcurveto{\pgfqpoint{3.043959in}{0.513233in}}{\pgfqpoint{3.046715in}{0.511180in}}%
\pgfpathquadraticcurveto{\pgfqpoint{3.049471in}{0.509126in}}{\pgfqpoint{3.057811in}{0.507271in}}%
\pgfpathlineto{\pgfqpoint{3.061359in}{0.506478in}}%
\pgfpathquadraticcurveto{\pgfqpoint{3.072383in}{0.504119in}}{\pgfqpoint{3.077030in}{0.499814in}}%
\pgfpathquadraticcurveto{\pgfqpoint{3.081677in}{0.495509in}}{\pgfqpoint{3.081677in}{0.487800in}}%
\pgfpathquadraticcurveto{\pgfqpoint{3.081677in}{0.479010in}}{\pgfqpoint{3.074724in}{0.473876in}}%
\pgfpathquadraticcurveto{\pgfqpoint{3.067771in}{0.468761in}}{\pgfqpoint{3.055613in}{0.468761in}}%
\pgfpathquadraticcurveto{\pgfqpoint{3.050552in}{0.468761in}}{\pgfqpoint{3.045058in}{0.469752in}}%
\pgfpathquadraticcurveto{\pgfqpoint{3.039564in}{0.470742in}}{\pgfqpoint{3.033494in}{0.472706in}}%
\pgfpathlineto{\pgfqpoint{3.033494in}{0.483405in}}%
\pgfpathquadraticcurveto{\pgfqpoint{3.039240in}{0.480415in}}{\pgfqpoint{3.044806in}{0.478920in}}%
\pgfpathquadraticcurveto{\pgfqpoint{3.050372in}{0.477443in}}{\pgfqpoint{3.055847in}{0.477443in}}%
\pgfpathquadraticcurveto{\pgfqpoint{3.063160in}{0.477443in}}{\pgfqpoint{3.067087in}{0.479946in}}%
\pgfpathquadraticcurveto{\pgfqpoint{3.071032in}{0.482450in}}{\pgfqpoint{3.071032in}{0.487007in}}%
\pgfpathquadraticcurveto{\pgfqpoint{3.071032in}{0.491222in}}{\pgfqpoint{3.068186in}{0.493474in}}%
\pgfpathquadraticcurveto{\pgfqpoint{3.065358in}{0.495725in}}{\pgfqpoint{3.055721in}{0.497814in}}%
\pgfpathlineto{\pgfqpoint{3.052119in}{0.498661in}}%
\pgfpathquadraticcurveto{\pgfqpoint{3.042500in}{0.500678in}}{\pgfqpoint{3.038214in}{0.504875in}}%
\pgfpathquadraticcurveto{\pgfqpoint{3.033945in}{0.509072in}}{\pgfqpoint{3.033945in}{0.516385in}}%
\pgfpathquadraticcurveto{\pgfqpoint{3.033945in}{0.525283in}}{\pgfqpoint{3.040249in}{0.530110in}}%
\pgfpathquadraticcurveto{\pgfqpoint{3.046553in}{0.534956in}}{\pgfqpoint{3.058153in}{0.534956in}}%
\pgfpathquadraticcurveto{\pgfqpoint{3.063881in}{0.534956in}}{\pgfqpoint{3.068942in}{0.534109in}}%
\pgfpathquadraticcurveto{\pgfqpoint{3.074022in}{0.533280in}}{\pgfqpoint{3.078291in}{0.531587in}}%
\pgfpathlineto{\pgfqpoint{3.078291in}{0.531587in}}%
\pgfpathclose%
\pgfusepath{fill}%
\end{pgfscope}%
\begin{pgfscope}%
\pgfsetbuttcap%
\pgfsetroundjoin%
\definecolor{currentfill}{rgb}{0.309804,0.372549,0.419608}%
\pgfsetfillcolor{currentfill}%
\pgfsetlinewidth{0.000000pt}%
\definecolor{currentstroke}{rgb}{0.309804,0.372549,0.419608}%
\pgfsetstrokecolor{currentstroke}%
\pgfsetdash{}{0pt}%
\pgfpathmoveto{\pgfqpoint{3.165157in}{0.479856in}}%
\pgfpathlineto{\pgfqpoint{3.165157in}{0.446426in}}%
\pgfpathlineto{\pgfqpoint{3.154746in}{0.446426in}}%
\pgfpathlineto{\pgfqpoint{3.154746in}{0.533443in}}%
\pgfpathlineto{\pgfqpoint{3.165157in}{0.533443in}}%
\pgfpathlineto{\pgfqpoint{3.165157in}{0.523878in}}%
\pgfpathquadraticcurveto{\pgfqpoint{3.168436in}{0.529498in}}{\pgfqpoint{3.173407in}{0.532218in}}%
\pgfpathquadraticcurveto{\pgfqpoint{3.178396in}{0.534956in}}{\pgfqpoint{3.185313in}{0.534956in}}%
\pgfpathquadraticcurveto{\pgfqpoint{3.196805in}{0.534956in}}{\pgfqpoint{3.203974in}{0.525841in}}%
\pgfpathquadraticcurveto{\pgfqpoint{3.211160in}{0.516727in}}{\pgfqpoint{3.211160in}{0.501867in}}%
\pgfpathquadraticcurveto{\pgfqpoint{3.211160in}{0.487007in}}{\pgfqpoint{3.203974in}{0.477875in}}%
\pgfpathquadraticcurveto{\pgfqpoint{3.196805in}{0.468761in}}{\pgfqpoint{3.185313in}{0.468761in}}%
\pgfpathquadraticcurveto{\pgfqpoint{3.178396in}{0.468761in}}{\pgfqpoint{3.173407in}{0.471499in}}%
\pgfpathquadraticcurveto{\pgfqpoint{3.168436in}{0.474237in}}{\pgfqpoint{3.165157in}{0.479856in}}%
\pgfpathlineto{\pgfqpoint{3.165157in}{0.479856in}}%
\pgfpathclose%
\pgfpathmoveto{\pgfqpoint{3.200407in}{0.501867in}}%
\pgfpathquadraticcurveto{\pgfqpoint{3.200407in}{0.513287in}}{\pgfqpoint{3.195706in}{0.519789in}}%
\pgfpathquadraticcurveto{\pgfqpoint{3.191005in}{0.526292in}}{\pgfqpoint{3.182791in}{0.526292in}}%
\pgfpathquadraticcurveto{\pgfqpoint{3.174560in}{0.526292in}}{\pgfqpoint{3.169859in}{0.519789in}}%
\pgfpathquadraticcurveto{\pgfqpoint{3.165157in}{0.513287in}}{\pgfqpoint{3.165157in}{0.501867in}}%
\pgfpathquadraticcurveto{\pgfqpoint{3.165157in}{0.490448in}}{\pgfqpoint{3.169859in}{0.483945in}}%
\pgfpathquadraticcurveto{\pgfqpoint{3.174560in}{0.477443in}}{\pgfqpoint{3.182791in}{0.477443in}}%
\pgfpathquadraticcurveto{\pgfqpoint{3.191005in}{0.477443in}}{\pgfqpoint{3.195706in}{0.483945in}}%
\pgfpathquadraticcurveto{\pgfqpoint{3.200407in}{0.490448in}}{\pgfqpoint{3.200407in}{0.501867in}}%
\pgfpathlineto{\pgfqpoint{3.200407in}{0.501867in}}%
\pgfpathclose%
\pgfpathmoveto{\pgfqpoint{3.256983in}{0.502083in}}%
\pgfpathquadraticcurveto{\pgfqpoint{3.244429in}{0.502083in}}{\pgfqpoint{3.239584in}{0.499219in}}%
\pgfpathquadraticcurveto{\pgfqpoint{3.234738in}{0.496356in}}{\pgfqpoint{3.234738in}{0.489421in}}%
\pgfpathquadraticcurveto{\pgfqpoint{3.234738in}{0.483909in}}{\pgfqpoint{3.238377in}{0.480667in}}%
\pgfpathquadraticcurveto{\pgfqpoint{3.242015in}{0.477443in}}{\pgfqpoint{3.248247in}{0.477443in}}%
\pgfpathquadraticcurveto{\pgfqpoint{3.256875in}{0.477443in}}{\pgfqpoint{3.262081in}{0.483549in}}%
\pgfpathquadraticcurveto{\pgfqpoint{3.267286in}{0.489655in}}{\pgfqpoint{3.267286in}{0.499778in}}%
\pgfpathlineto{\pgfqpoint{3.267286in}{0.502083in}}%
\pgfpathlineto{\pgfqpoint{3.256983in}{0.502083in}}%
\pgfpathlineto{\pgfqpoint{3.256983in}{0.502083in}}%
\pgfpathclose%
\pgfpathmoveto{\pgfqpoint{3.277643in}{0.506370in}}%
\pgfpathlineto{\pgfqpoint{3.277643in}{0.470400in}}%
\pgfpathlineto{\pgfqpoint{3.267286in}{0.470400in}}%
\pgfpathlineto{\pgfqpoint{3.267286in}{0.479964in}}%
\pgfpathquadraticcurveto{\pgfqpoint{3.263738in}{0.474237in}}{\pgfqpoint{3.258442in}{0.471499in}}%
\pgfpathquadraticcurveto{\pgfqpoint{3.253147in}{0.468761in}}{\pgfqpoint{3.245492in}{0.468761in}}%
\pgfpathquadraticcurveto{\pgfqpoint{3.235819in}{0.468761in}}{\pgfqpoint{3.230091in}{0.474201in}}%
\pgfpathquadraticcurveto{\pgfqpoint{3.224381in}{0.479640in}}{\pgfqpoint{3.224381in}{0.488754in}}%
\pgfpathquadraticcurveto{\pgfqpoint{3.224381in}{0.499382in}}{\pgfqpoint{3.231496in}{0.504785in}}%
\pgfpathquadraticcurveto{\pgfqpoint{3.238629in}{0.510189in}}{\pgfqpoint{3.252750in}{0.510189in}}%
\pgfpathlineto{\pgfqpoint{3.267286in}{0.510189in}}%
\pgfpathlineto{\pgfqpoint{3.267286in}{0.511216in}}%
\pgfpathquadraticcurveto{\pgfqpoint{3.267286in}{0.518366in}}{\pgfqpoint{3.262585in}{0.522275in}}%
\pgfpathquadraticcurveto{\pgfqpoint{3.257884in}{0.526184in}}{\pgfqpoint{3.249382in}{0.526184in}}%
\pgfpathquadraticcurveto{\pgfqpoint{3.243979in}{0.526184in}}{\pgfqpoint{3.238845in}{0.524887in}}%
\pgfpathquadraticcurveto{\pgfqpoint{3.233730in}{0.523590in}}{\pgfqpoint{3.229010in}{0.520996in}}%
\pgfpathlineto{\pgfqpoint{3.229010in}{0.530579in}}%
\pgfpathquadraticcurveto{\pgfqpoint{3.234684in}{0.532776in}}{\pgfqpoint{3.240034in}{0.533857in}}%
\pgfpathquadraticcurveto{\pgfqpoint{3.245383in}{0.534956in}}{\pgfqpoint{3.250445in}{0.534956in}}%
\pgfpathquadraticcurveto{\pgfqpoint{3.264134in}{0.534956in}}{\pgfqpoint{3.270889in}{0.527859in}}%
\pgfpathquadraticcurveto{\pgfqpoint{3.277643in}{0.520780in}}{\pgfqpoint{3.277643in}{0.506370in}}%
\pgfpathlineto{\pgfqpoint{3.277643in}{0.506370in}}%
\pgfpathclose%
\pgfpathmoveto{\pgfqpoint{3.298970in}{0.533443in}}%
\pgfpathlineto{\pgfqpoint{3.309327in}{0.533443in}}%
\pgfpathlineto{\pgfqpoint{3.309327in}{0.470400in}}%
\pgfpathlineto{\pgfqpoint{3.298970in}{0.470400in}}%
\pgfpathlineto{\pgfqpoint{3.298970in}{0.533443in}}%
\pgfpathlineto{\pgfqpoint{3.298970in}{0.533443in}}%
\pgfpathclose%
\pgfpathmoveto{\pgfqpoint{3.298970in}{0.557993in}}%
\pgfpathlineto{\pgfqpoint{3.309327in}{0.557993in}}%
\pgfpathlineto{\pgfqpoint{3.309327in}{0.544862in}}%
\pgfpathlineto{\pgfqpoint{3.298970in}{0.544862in}}%
\pgfpathlineto{\pgfqpoint{3.298970in}{0.557993in}}%
\pgfpathlineto{\pgfqpoint{3.298970in}{0.557993in}}%
\pgfpathclose%
\pgfpathmoveto{\pgfqpoint{3.367524in}{0.523770in}}%
\pgfpathquadraticcurveto{\pgfqpoint{3.365777in}{0.524779in}}{\pgfqpoint{3.363723in}{0.525247in}}%
\pgfpathquadraticcurveto{\pgfqpoint{3.361670in}{0.525733in}}{\pgfqpoint{3.359202in}{0.525733in}}%
\pgfpathquadraticcurveto{\pgfqpoint{3.350412in}{0.525733in}}{\pgfqpoint{3.345711in}{0.520023in}}%
\pgfpathquadraticcurveto{\pgfqpoint{3.341010in}{0.514314in}}{\pgfqpoint{3.341010in}{0.503614in}}%
\pgfpathlineto{\pgfqpoint{3.341010in}{0.470400in}}%
\pgfpathlineto{\pgfqpoint{3.330599in}{0.470400in}}%
\pgfpathlineto{\pgfqpoint{3.330599in}{0.533443in}}%
\pgfpathlineto{\pgfqpoint{3.341010in}{0.533443in}}%
\pgfpathlineto{\pgfqpoint{3.341010in}{0.523644in}}%
\pgfpathquadraticcurveto{\pgfqpoint{3.344288in}{0.529390in}}{\pgfqpoint{3.349512in}{0.532164in}}%
\pgfpathquadraticcurveto{\pgfqpoint{3.354753in}{0.534956in}}{\pgfqpoint{3.362246in}{0.534956in}}%
\pgfpathquadraticcurveto{\pgfqpoint{3.363309in}{0.534956in}}{\pgfqpoint{3.364606in}{0.534811in}}%
\pgfpathquadraticcurveto{\pgfqpoint{3.365903in}{0.534685in}}{\pgfqpoint{3.367470in}{0.534397in}}%
\pgfpathlineto{\pgfqpoint{3.367524in}{0.523770in}}%
\pgfpathlineto{\pgfqpoint{3.367524in}{0.523770in}}%
\pgfpathclose%
\pgfpathmoveto{\pgfqpoint{3.388184in}{0.554445in}}%
\pgfpathlineto{\pgfqpoint{3.388184in}{0.523194in}}%
\pgfpathlineto{\pgfqpoint{3.378619in}{0.523194in}}%
\pgfpathlineto{\pgfqpoint{3.378619in}{0.554445in}}%
\pgfpathlineto{\pgfqpoint{3.388184in}{0.554445in}}%
\pgfpathlineto{\pgfqpoint{3.388184in}{0.554445in}}%
\pgfpathclose%
\pgfpathmoveto{\pgfqpoint{3.450254in}{0.531587in}}%
\pgfpathlineto{\pgfqpoint{3.450254in}{0.521789in}}%
\pgfpathquadraticcurveto{\pgfqpoint{3.445877in}{0.524040in}}{\pgfqpoint{3.441140in}{0.525157in}}%
\pgfpathquadraticcurveto{\pgfqpoint{3.436420in}{0.526292in}}{\pgfqpoint{3.431341in}{0.526292in}}%
\pgfpathquadraticcurveto{\pgfqpoint{3.423632in}{0.526292in}}{\pgfqpoint{3.419777in}{0.523932in}}%
\pgfpathquadraticcurveto{\pgfqpoint{3.415923in}{0.521573in}}{\pgfqpoint{3.415923in}{0.516835in}}%
\pgfpathquadraticcurveto{\pgfqpoint{3.415923in}{0.513233in}}{\pgfqpoint{3.418678in}{0.511180in}}%
\pgfpathquadraticcurveto{\pgfqpoint{3.421434in}{0.509126in}}{\pgfqpoint{3.429774in}{0.507271in}}%
\pgfpathlineto{\pgfqpoint{3.433322in}{0.506478in}}%
\pgfpathquadraticcurveto{\pgfqpoint{3.444346in}{0.504119in}}{\pgfqpoint{3.448993in}{0.499814in}}%
\pgfpathquadraticcurveto{\pgfqpoint{3.453640in}{0.495509in}}{\pgfqpoint{3.453640in}{0.487800in}}%
\pgfpathquadraticcurveto{\pgfqpoint{3.453640in}{0.479010in}}{\pgfqpoint{3.446687in}{0.473876in}}%
\pgfpathquadraticcurveto{\pgfqpoint{3.439735in}{0.468761in}}{\pgfqpoint{3.427576in}{0.468761in}}%
\pgfpathquadraticcurveto{\pgfqpoint{3.422515in}{0.468761in}}{\pgfqpoint{3.417021in}{0.469752in}}%
\pgfpathquadraticcurveto{\pgfqpoint{3.411528in}{0.470742in}}{\pgfqpoint{3.405457in}{0.472706in}}%
\pgfpathlineto{\pgfqpoint{3.405457in}{0.483405in}}%
\pgfpathquadraticcurveto{\pgfqpoint{3.411203in}{0.480415in}}{\pgfqpoint{3.416769in}{0.478920in}}%
\pgfpathquadraticcurveto{\pgfqpoint{3.422335in}{0.477443in}}{\pgfqpoint{3.427811in}{0.477443in}}%
\pgfpathquadraticcurveto{\pgfqpoint{3.435123in}{0.477443in}}{\pgfqpoint{3.439050in}{0.479946in}}%
\pgfpathquadraticcurveto{\pgfqpoint{3.442995in}{0.482450in}}{\pgfqpoint{3.442995in}{0.487007in}}%
\pgfpathquadraticcurveto{\pgfqpoint{3.442995in}{0.491222in}}{\pgfqpoint{3.440149in}{0.493474in}}%
\pgfpathquadraticcurveto{\pgfqpoint{3.437321in}{0.495725in}}{\pgfqpoint{3.427684in}{0.497814in}}%
\pgfpathlineto{\pgfqpoint{3.424082in}{0.498661in}}%
\pgfpathquadraticcurveto{\pgfqpoint{3.414464in}{0.500678in}}{\pgfqpoint{3.410177in}{0.504875in}}%
\pgfpathquadraticcurveto{\pgfqpoint{3.405908in}{0.509072in}}{\pgfqpoint{3.405908in}{0.516385in}}%
\pgfpathquadraticcurveto{\pgfqpoint{3.405908in}{0.525283in}}{\pgfqpoint{3.412212in}{0.530110in}}%
\pgfpathquadraticcurveto{\pgfqpoint{3.418516in}{0.534956in}}{\pgfqpoint{3.430116in}{0.534956in}}%
\pgfpathquadraticcurveto{\pgfqpoint{3.435844in}{0.534956in}}{\pgfqpoint{3.440905in}{0.534109in}}%
\pgfpathquadraticcurveto{\pgfqpoint{3.445985in}{0.533280in}}{\pgfqpoint{3.450254in}{0.531587in}}%
\pgfpathlineto{\pgfqpoint{3.450254in}{0.531587in}}%
\pgfpathclose%
\pgfusepath{fill}%
\end{pgfscope}%
\begin{pgfscope}%
\pgfsetbuttcap%
\pgfsetroundjoin%
\definecolor{currentfill}{rgb}{0.309804,0.372549,0.419608}%
\pgfsetfillcolor{currentfill}%
\pgfsetlinewidth{0.000000pt}%
\definecolor{currentstroke}{rgb}{0.309804,0.372549,0.419608}%
\pgfsetstrokecolor{currentstroke}%
\pgfsetdash{}{0pt}%
\pgfpathmoveto{\pgfqpoint{3.562640in}{0.531587in}}%
\pgfpathlineto{\pgfqpoint{3.562640in}{0.521789in}}%
\pgfpathquadraticcurveto{\pgfqpoint{3.558263in}{0.524040in}}{\pgfqpoint{3.553526in}{0.525157in}}%
\pgfpathquadraticcurveto{\pgfqpoint{3.548807in}{0.526292in}}{\pgfqpoint{3.543727in}{0.526292in}}%
\pgfpathquadraticcurveto{\pgfqpoint{3.536018in}{0.526292in}}{\pgfqpoint{3.532164in}{0.523932in}}%
\pgfpathquadraticcurveto{\pgfqpoint{3.528309in}{0.521573in}}{\pgfqpoint{3.528309in}{0.516835in}}%
\pgfpathquadraticcurveto{\pgfqpoint{3.528309in}{0.513233in}}{\pgfqpoint{3.531065in}{0.511180in}}%
\pgfpathquadraticcurveto{\pgfqpoint{3.533821in}{0.509126in}}{\pgfqpoint{3.542160in}{0.507271in}}%
\pgfpathlineto{\pgfqpoint{3.545709in}{0.506478in}}%
\pgfpathquadraticcurveto{\pgfqpoint{3.556732in}{0.504119in}}{\pgfqpoint{3.561379in}{0.499814in}}%
\pgfpathquadraticcurveto{\pgfqpoint{3.566026in}{0.495509in}}{\pgfqpoint{3.566026in}{0.487800in}}%
\pgfpathquadraticcurveto{\pgfqpoint{3.566026in}{0.479010in}}{\pgfqpoint{3.559074in}{0.473876in}}%
\pgfpathquadraticcurveto{\pgfqpoint{3.552121in}{0.468761in}}{\pgfqpoint{3.539963in}{0.468761in}}%
\pgfpathquadraticcurveto{\pgfqpoint{3.534901in}{0.468761in}}{\pgfqpoint{3.529408in}{0.469752in}}%
\pgfpathquadraticcurveto{\pgfqpoint{3.523914in}{0.470742in}}{\pgfqpoint{3.517844in}{0.472706in}}%
\pgfpathlineto{\pgfqpoint{3.517844in}{0.483405in}}%
\pgfpathquadraticcurveto{\pgfqpoint{3.523590in}{0.480415in}}{\pgfqpoint{3.529156in}{0.478920in}}%
\pgfpathquadraticcurveto{\pgfqpoint{3.534721in}{0.477443in}}{\pgfqpoint{3.540197in}{0.477443in}}%
\pgfpathquadraticcurveto{\pgfqpoint{3.547510in}{0.477443in}}{\pgfqpoint{3.551437in}{0.479946in}}%
\pgfpathquadraticcurveto{\pgfqpoint{3.555381in}{0.482450in}}{\pgfqpoint{3.555381in}{0.487007in}}%
\pgfpathquadraticcurveto{\pgfqpoint{3.555381in}{0.491222in}}{\pgfqpoint{3.552535in}{0.493474in}}%
\pgfpathquadraticcurveto{\pgfqpoint{3.549707in}{0.495725in}}{\pgfqpoint{3.540071in}{0.497814in}}%
\pgfpathlineto{\pgfqpoint{3.536469in}{0.498661in}}%
\pgfpathquadraticcurveto{\pgfqpoint{3.526850in}{0.500678in}}{\pgfqpoint{3.522563in}{0.504875in}}%
\pgfpathquadraticcurveto{\pgfqpoint{3.518294in}{0.509072in}}{\pgfqpoint{3.518294in}{0.516385in}}%
\pgfpathquadraticcurveto{\pgfqpoint{3.518294in}{0.525283in}}{\pgfqpoint{3.524599in}{0.530110in}}%
\pgfpathquadraticcurveto{\pgfqpoint{3.530903in}{0.534956in}}{\pgfqpoint{3.542503in}{0.534956in}}%
\pgfpathquadraticcurveto{\pgfqpoint{3.548230in}{0.534956in}}{\pgfqpoint{3.553292in}{0.534109in}}%
\pgfpathquadraticcurveto{\pgfqpoint{3.558371in}{0.533280in}}{\pgfqpoint{3.562640in}{0.531587in}}%
\pgfpathlineto{\pgfqpoint{3.562640in}{0.531587in}}%
\pgfpathclose%
\pgfpathmoveto{\pgfqpoint{3.611165in}{0.502083in}}%
\pgfpathquadraticcurveto{\pgfqpoint{3.598610in}{0.502083in}}{\pgfqpoint{3.593765in}{0.499219in}}%
\pgfpathquadraticcurveto{\pgfqpoint{3.588920in}{0.496356in}}{\pgfqpoint{3.588920in}{0.489421in}}%
\pgfpathquadraticcurveto{\pgfqpoint{3.588920in}{0.483909in}}{\pgfqpoint{3.592558in}{0.480667in}}%
\pgfpathquadraticcurveto{\pgfqpoint{3.596197in}{0.477443in}}{\pgfqpoint{3.602429in}{0.477443in}}%
\pgfpathquadraticcurveto{\pgfqpoint{3.611057in}{0.477443in}}{\pgfqpoint{3.616262in}{0.483549in}}%
\pgfpathquadraticcurveto{\pgfqpoint{3.621468in}{0.489655in}}{\pgfqpoint{3.621468in}{0.499778in}}%
\pgfpathlineto{\pgfqpoint{3.621468in}{0.502083in}}%
\pgfpathlineto{\pgfqpoint{3.611165in}{0.502083in}}%
\pgfpathlineto{\pgfqpoint{3.611165in}{0.502083in}}%
\pgfpathclose%
\pgfpathmoveto{\pgfqpoint{3.631825in}{0.506370in}}%
\pgfpathlineto{\pgfqpoint{3.631825in}{0.470400in}}%
\pgfpathlineto{\pgfqpoint{3.621468in}{0.470400in}}%
\pgfpathlineto{\pgfqpoint{3.621468in}{0.479964in}}%
\pgfpathquadraticcurveto{\pgfqpoint{3.617919in}{0.474237in}}{\pgfqpoint{3.612624in}{0.471499in}}%
\pgfpathquadraticcurveto{\pgfqpoint{3.607328in}{0.468761in}}{\pgfqpoint{3.599673in}{0.468761in}}%
\pgfpathquadraticcurveto{\pgfqpoint{3.590001in}{0.468761in}}{\pgfqpoint{3.584273in}{0.474201in}}%
\pgfpathquadraticcurveto{\pgfqpoint{3.578563in}{0.479640in}}{\pgfqpoint{3.578563in}{0.488754in}}%
\pgfpathquadraticcurveto{\pgfqpoint{3.578563in}{0.499382in}}{\pgfqpoint{3.585678in}{0.504785in}}%
\pgfpathquadraticcurveto{\pgfqpoint{3.592811in}{0.510189in}}{\pgfqpoint{3.606932in}{0.510189in}}%
\pgfpathlineto{\pgfqpoint{3.621468in}{0.510189in}}%
\pgfpathlineto{\pgfqpoint{3.621468in}{0.511216in}}%
\pgfpathquadraticcurveto{\pgfqpoint{3.621468in}{0.518366in}}{\pgfqpoint{3.616767in}{0.522275in}}%
\pgfpathquadraticcurveto{\pgfqpoint{3.612066in}{0.526184in}}{\pgfqpoint{3.603564in}{0.526184in}}%
\pgfpathquadraticcurveto{\pgfqpoint{3.598160in}{0.526184in}}{\pgfqpoint{3.593027in}{0.524887in}}%
\pgfpathquadraticcurveto{\pgfqpoint{3.587911in}{0.523590in}}{\pgfqpoint{3.583192in}{0.520996in}}%
\pgfpathlineto{\pgfqpoint{3.583192in}{0.530579in}}%
\pgfpathquadraticcurveto{\pgfqpoint{3.588866in}{0.532776in}}{\pgfqpoint{3.594215in}{0.533857in}}%
\pgfpathquadraticcurveto{\pgfqpoint{3.599565in}{0.534956in}}{\pgfqpoint{3.604627in}{0.534956in}}%
\pgfpathquadraticcurveto{\pgfqpoint{3.618316in}{0.534956in}}{\pgfqpoint{3.625070in}{0.527859in}}%
\pgfpathquadraticcurveto{\pgfqpoint{3.631825in}{0.520780in}}{\pgfqpoint{3.631825in}{0.506370in}}%
\pgfpathlineto{\pgfqpoint{3.631825in}{0.506370in}}%
\pgfpathclose%
\pgfpathmoveto{\pgfqpoint{3.645730in}{0.533443in}}%
\pgfpathlineto{\pgfqpoint{3.656700in}{0.533443in}}%
\pgfpathlineto{\pgfqpoint{3.676405in}{0.480541in}}%
\pgfpathlineto{\pgfqpoint{3.696110in}{0.533443in}}%
\pgfpathlineto{\pgfqpoint{3.707080in}{0.533443in}}%
\pgfpathlineto{\pgfqpoint{3.683430in}{0.470400in}}%
\pgfpathlineto{\pgfqpoint{3.669362in}{0.470400in}}%
\pgfpathlineto{\pgfqpoint{3.645730in}{0.533443in}}%
\pgfpathlineto{\pgfqpoint{3.645730in}{0.533443in}}%
\pgfpathclose%
\pgfpathmoveto{\pgfqpoint{3.775310in}{0.504515in}}%
\pgfpathlineto{\pgfqpoint{3.775310in}{0.499454in}}%
\pgfpathlineto{\pgfqpoint{3.727686in}{0.499454in}}%
\pgfpathquadraticcurveto{\pgfqpoint{3.728370in}{0.488754in}}{\pgfqpoint{3.734134in}{0.483153in}}%
\pgfpathquadraticcurveto{\pgfqpoint{3.739898in}{0.477551in}}{\pgfqpoint{3.750201in}{0.477551in}}%
\pgfpathquadraticcurveto{\pgfqpoint{3.756163in}{0.477551in}}{\pgfqpoint{3.761765in}{0.479010in}}%
\pgfpathquadraticcurveto{\pgfqpoint{3.767366in}{0.480469in}}{\pgfqpoint{3.772896in}{0.483405in}}%
\pgfpathlineto{\pgfqpoint{3.772896in}{0.473606in}}%
\pgfpathquadraticcurveto{\pgfqpoint{3.767312in}{0.471247in}}{\pgfqpoint{3.761458in}{0.470004in}}%
\pgfpathquadraticcurveto{\pgfqpoint{3.755604in}{0.468761in}}{\pgfqpoint{3.749588in}{0.468761in}}%
\pgfpathquadraticcurveto{\pgfqpoint{3.734494in}{0.468761in}}{\pgfqpoint{3.725686in}{0.477533in}}%
\pgfpathquadraticcurveto{\pgfqpoint{3.716878in}{0.486323in}}{\pgfqpoint{3.716878in}{0.501309in}}%
\pgfpathquadraticcurveto{\pgfqpoint{3.716878in}{0.516781in}}{\pgfqpoint{3.725236in}{0.525859in}}%
\pgfpathquadraticcurveto{\pgfqpoint{3.733594in}{0.534956in}}{\pgfqpoint{3.747787in}{0.534956in}}%
\pgfpathquadraticcurveto{\pgfqpoint{3.760504in}{0.534956in}}{\pgfqpoint{3.767907in}{0.526760in}}%
\pgfpathquadraticcurveto{\pgfqpoint{3.775310in}{0.518583in}}{\pgfqpoint{3.775310in}{0.504515in}}%
\pgfpathlineto{\pgfqpoint{3.775310in}{0.504515in}}%
\pgfpathclose%
\pgfpathmoveto{\pgfqpoint{3.764953in}{0.507559in}}%
\pgfpathquadraticcurveto{\pgfqpoint{3.764845in}{0.516043in}}{\pgfqpoint{3.760197in}{0.521104in}}%
\pgfpathquadraticcurveto{\pgfqpoint{3.755550in}{0.526184in}}{\pgfqpoint{3.747895in}{0.526184in}}%
\pgfpathquadraticcurveto{\pgfqpoint{3.739231in}{0.526184in}}{\pgfqpoint{3.734026in}{0.521284in}}%
\pgfpathquadraticcurveto{\pgfqpoint{3.728820in}{0.516385in}}{\pgfqpoint{3.728028in}{0.507487in}}%
\pgfpathlineto{\pgfqpoint{3.764953in}{0.507559in}}%
\pgfpathlineto{\pgfqpoint{3.764953in}{0.507559in}}%
\pgfpathclose%
\pgfpathmoveto{\pgfqpoint{3.833795in}{0.523878in}}%
\pgfpathlineto{\pgfqpoint{3.833795in}{0.557993in}}%
\pgfpathlineto{\pgfqpoint{3.844152in}{0.557993in}}%
\pgfpathlineto{\pgfqpoint{3.844152in}{0.470400in}}%
\pgfpathlineto{\pgfqpoint{3.833795in}{0.470400in}}%
\pgfpathlineto{\pgfqpoint{3.833795in}{0.479856in}}%
\pgfpathquadraticcurveto{\pgfqpoint{3.830535in}{0.474237in}}{\pgfqpoint{3.825546in}{0.471499in}}%
\pgfpathquadraticcurveto{\pgfqpoint{3.820574in}{0.468761in}}{\pgfqpoint{3.813586in}{0.468761in}}%
\pgfpathquadraticcurveto{\pgfqpoint{3.802166in}{0.468761in}}{\pgfqpoint{3.794979in}{0.477875in}}%
\pgfpathquadraticcurveto{\pgfqpoint{3.787810in}{0.487007in}}{\pgfqpoint{3.787810in}{0.501867in}}%
\pgfpathquadraticcurveto{\pgfqpoint{3.787810in}{0.516727in}}{\pgfqpoint{3.794979in}{0.525841in}}%
\pgfpathquadraticcurveto{\pgfqpoint{3.802166in}{0.534956in}}{\pgfqpoint{3.813586in}{0.534956in}}%
\pgfpathquadraticcurveto{\pgfqpoint{3.820574in}{0.534956in}}{\pgfqpoint{3.825546in}{0.532218in}}%
\pgfpathquadraticcurveto{\pgfqpoint{3.830535in}{0.529498in}}{\pgfqpoint{3.833795in}{0.523878in}}%
\pgfpathlineto{\pgfqpoint{3.833795in}{0.523878in}}%
\pgfpathclose%
\pgfpathmoveto{\pgfqpoint{3.798509in}{0.501867in}}%
\pgfpathquadraticcurveto{\pgfqpoint{3.798509in}{0.490448in}}{\pgfqpoint{3.803211in}{0.483945in}}%
\pgfpathquadraticcurveto{\pgfqpoint{3.807912in}{0.477443in}}{\pgfqpoint{3.816125in}{0.477443in}}%
\pgfpathquadraticcurveto{\pgfqpoint{3.824339in}{0.477443in}}{\pgfqpoint{3.829058in}{0.483945in}}%
\pgfpathquadraticcurveto{\pgfqpoint{3.833795in}{0.490448in}}{\pgfqpoint{3.833795in}{0.501867in}}%
\pgfpathquadraticcurveto{\pgfqpoint{3.833795in}{0.513287in}}{\pgfqpoint{3.829058in}{0.519789in}}%
\pgfpathquadraticcurveto{\pgfqpoint{3.824339in}{0.526292in}}{\pgfqpoint{3.816125in}{0.526292in}}%
\pgfpathquadraticcurveto{\pgfqpoint{3.807912in}{0.526292in}}{\pgfqpoint{3.803211in}{0.519789in}}%
\pgfpathquadraticcurveto{\pgfqpoint{3.798509in}{0.513287in}}{\pgfqpoint{3.798509in}{0.501867in}}%
\pgfpathlineto{\pgfqpoint{3.798509in}{0.501867in}}%
\pgfpathclose%
\pgfusepath{fill}%
\end{pgfscope}%
\begin{pgfscope}%
\pgfsetbuttcap%
\pgfsetroundjoin%
\definecolor{currentfill}{rgb}{0.309804,0.372549,0.419608}%
\pgfsetfillcolor{currentfill}%
\pgfsetlinewidth{0.000000pt}%
\definecolor{currentstroke}{rgb}{0.309804,0.372549,0.419608}%
\pgfsetstrokecolor{currentstroke}%
\pgfsetdash{}{0pt}%
\pgfpathmoveto{\pgfqpoint{3.908564in}{0.554445in}}%
\pgfpathlineto{\pgfqpoint{3.979658in}{0.554445in}}%
\pgfpathlineto{\pgfqpoint{3.979658in}{0.544862in}}%
\pgfpathlineto{\pgfqpoint{3.949830in}{0.544862in}}%
\pgfpathlineto{\pgfqpoint{3.949830in}{0.470400in}}%
\pgfpathlineto{\pgfqpoint{3.938410in}{0.470400in}}%
\pgfpathlineto{\pgfqpoint{3.938410in}{0.544862in}}%
\pgfpathlineto{\pgfqpoint{3.908564in}{0.544862in}}%
\pgfpathlineto{\pgfqpoint{3.908564in}{0.554445in}}%
\pgfpathlineto{\pgfqpoint{3.908564in}{0.554445in}}%
\pgfpathclose%
\pgfpathmoveto{\pgfqpoint{3.983152in}{0.554445in}}%
\pgfpathlineto{\pgfqpoint{3.994626in}{0.554445in}}%
\pgfpathlineto{\pgfqpoint{4.012296in}{0.483405in}}%
\pgfpathlineto{\pgfqpoint{4.029912in}{0.554445in}}%
\pgfpathlineto{\pgfqpoint{4.042701in}{0.554445in}}%
\pgfpathlineto{\pgfqpoint{4.060370in}{0.483405in}}%
\pgfpathlineto{\pgfqpoint{4.077986in}{0.554445in}}%
\pgfpathlineto{\pgfqpoint{4.089532in}{0.554445in}}%
\pgfpathlineto{\pgfqpoint{4.068422in}{0.470400in}}%
\pgfpathlineto{\pgfqpoint{4.054120in}{0.470400in}}%
\pgfpathlineto{\pgfqpoint{4.036396in}{0.543349in}}%
\pgfpathlineto{\pgfqpoint{4.018492in}{0.470400in}}%
\pgfpathlineto{\pgfqpoint{4.004191in}{0.470400in}}%
\pgfpathlineto{\pgfqpoint{3.983152in}{0.554445in}}%
\pgfpathlineto{\pgfqpoint{3.983152in}{0.554445in}}%
\pgfpathclose%
\pgfpathmoveto{\pgfqpoint{4.104608in}{0.554445in}}%
\pgfpathlineto{\pgfqpoint{4.152899in}{0.554445in}}%
\pgfpathlineto{\pgfqpoint{4.152899in}{0.544862in}}%
\pgfpathlineto{\pgfqpoint{4.115974in}{0.544862in}}%
\pgfpathlineto{\pgfqpoint{4.115974in}{0.520096in}}%
\pgfpathlineto{\pgfqpoint{4.149296in}{0.520096in}}%
\pgfpathlineto{\pgfqpoint{4.149296in}{0.510531in}}%
\pgfpathlineto{\pgfqpoint{4.115974in}{0.510531in}}%
\pgfpathlineto{\pgfqpoint{4.115974in}{0.470400in}}%
\pgfpathlineto{\pgfqpoint{4.104608in}{0.470400in}}%
\pgfpathlineto{\pgfqpoint{4.104608in}{0.554445in}}%
\pgfpathlineto{\pgfqpoint{4.104608in}{0.554445in}}%
\pgfpathclose%
\pgfpathmoveto{\pgfqpoint{4.170911in}{0.554445in}}%
\pgfpathlineto{\pgfqpoint{4.224047in}{0.554445in}}%
\pgfpathlineto{\pgfqpoint{4.224047in}{0.544862in}}%
\pgfpathlineto{\pgfqpoint{4.182277in}{0.544862in}}%
\pgfpathlineto{\pgfqpoint{4.182277in}{0.519987in}}%
\pgfpathlineto{\pgfqpoint{4.222300in}{0.519987in}}%
\pgfpathlineto{\pgfqpoint{4.222300in}{0.510423in}}%
\pgfpathlineto{\pgfqpoint{4.182277in}{0.510423in}}%
\pgfpathlineto{\pgfqpoint{4.182277in}{0.479964in}}%
\pgfpathlineto{\pgfqpoint{4.225056in}{0.479964in}}%
\pgfpathlineto{\pgfqpoint{4.225056in}{0.470400in}}%
\pgfpathlineto{\pgfqpoint{4.170911in}{0.470400in}}%
\pgfpathlineto{\pgfqpoint{4.170911in}{0.554445in}}%
\pgfpathlineto{\pgfqpoint{4.170911in}{0.554445in}}%
\pgfpathclose%
\pgfusepath{fill}%
\end{pgfscope}%
\begin{pgfscope}%
\pgfsetbuttcap%
\pgfsetroundjoin%
\definecolor{currentfill}{rgb}{0.309804,0.372549,0.419608}%
\pgfsetfillcolor{currentfill}%
\pgfsetlinewidth{0.000000pt}%
\definecolor{currentstroke}{rgb}{0.309804,0.372549,0.419608}%
\pgfsetstrokecolor{currentstroke}%
\pgfsetdash{}{0pt}%
\pgfpathmoveto{\pgfqpoint{4.357910in}{0.551689in}}%
\pgfpathlineto{\pgfqpoint{4.357910in}{0.540593in}}%
\pgfpathquadraticcurveto{\pgfqpoint{4.351444in}{0.543691in}}{\pgfqpoint{4.345698in}{0.545204in}}%
\pgfpathquadraticcurveto{\pgfqpoint{4.339952in}{0.546736in}}{\pgfqpoint{4.334603in}{0.546736in}}%
\pgfpathquadraticcurveto{\pgfqpoint{4.325326in}{0.546736in}}{\pgfqpoint{4.320283in}{0.543133in}}%
\pgfpathquadraticcurveto{\pgfqpoint{4.315240in}{0.539531in}}{\pgfqpoint{4.315240in}{0.532884in}}%
\pgfpathquadraticcurveto{\pgfqpoint{4.315240in}{0.527300in}}{\pgfqpoint{4.318590in}{0.524454in}}%
\pgfpathquadraticcurveto{\pgfqpoint{4.321940in}{0.521627in}}{\pgfqpoint{4.331288in}{0.519879in}}%
\pgfpathlineto{\pgfqpoint{4.338151in}{0.518474in}}%
\pgfpathquadraticcurveto{\pgfqpoint{4.350868in}{0.516043in}}{\pgfqpoint{4.356920in}{0.509937in}}%
\pgfpathquadraticcurveto{\pgfqpoint{4.362972in}{0.503831in}}{\pgfqpoint{4.362972in}{0.493600in}}%
\pgfpathquadraticcurveto{\pgfqpoint{4.362972in}{0.481369in}}{\pgfqpoint{4.354776in}{0.475065in}}%
\pgfpathquadraticcurveto{\pgfqpoint{4.346599in}{0.468761in}}{\pgfqpoint{4.330784in}{0.468761in}}%
\pgfpathquadraticcurveto{\pgfqpoint{4.324822in}{0.468761in}}{\pgfqpoint{4.318086in}{0.470112in}}%
\pgfpathquadraticcurveto{\pgfqpoint{4.311367in}{0.471463in}}{\pgfqpoint{4.304162in}{0.474111in}}%
\pgfpathlineto{\pgfqpoint{4.304162in}{0.485818in}}%
\pgfpathquadraticcurveto{\pgfqpoint{4.311079in}{0.481946in}}{\pgfqpoint{4.317725in}{0.479964in}}%
\pgfpathquadraticcurveto{\pgfqpoint{4.324372in}{0.478001in}}{\pgfqpoint{4.330784in}{0.478001in}}%
\pgfpathquadraticcurveto{\pgfqpoint{4.340511in}{0.478001in}}{\pgfqpoint{4.345806in}{0.481820in}}%
\pgfpathquadraticcurveto{\pgfqpoint{4.351102in}{0.485656in}}{\pgfqpoint{4.351102in}{0.492753in}}%
\pgfpathquadraticcurveto{\pgfqpoint{4.351102in}{0.498931in}}{\pgfqpoint{4.347301in}{0.502426in}}%
\pgfpathquadraticcurveto{\pgfqpoint{4.343501in}{0.505920in}}{\pgfqpoint{4.334837in}{0.507667in}}%
\pgfpathlineto{\pgfqpoint{4.327902in}{0.509018in}}%
\pgfpathquadraticcurveto{\pgfqpoint{4.315186in}{0.511540in}}{\pgfqpoint{4.309494in}{0.516943in}}%
\pgfpathquadraticcurveto{\pgfqpoint{4.303820in}{0.522347in}}{\pgfqpoint{4.303820in}{0.531984in}}%
\pgfpathquadraticcurveto{\pgfqpoint{4.303820in}{0.543133in}}{\pgfqpoint{4.311673in}{0.549545in}}%
\pgfpathquadraticcurveto{\pgfqpoint{4.319526in}{0.555958in}}{\pgfqpoint{4.333306in}{0.555958in}}%
\pgfpathquadraticcurveto{\pgfqpoint{4.339232in}{0.555958in}}{\pgfqpoint{4.345356in}{0.554877in}}%
\pgfpathquadraticcurveto{\pgfqpoint{4.351498in}{0.553814in}}{\pgfqpoint{4.357910in}{0.551689in}}%
\pgfpathlineto{\pgfqpoint{4.357910in}{0.551689in}}%
\pgfpathclose%
\pgfpathmoveto{\pgfqpoint{4.380714in}{0.554445in}}%
\pgfpathlineto{\pgfqpoint{4.433850in}{0.554445in}}%
\pgfpathlineto{\pgfqpoint{4.433850in}{0.544862in}}%
\pgfpathlineto{\pgfqpoint{4.392079in}{0.544862in}}%
\pgfpathlineto{\pgfqpoint{4.392079in}{0.519987in}}%
\pgfpathlineto{\pgfqpoint{4.432102in}{0.519987in}}%
\pgfpathlineto{\pgfqpoint{4.432102in}{0.510423in}}%
\pgfpathlineto{\pgfqpoint{4.392079in}{0.510423in}}%
\pgfpathlineto{\pgfqpoint{4.392079in}{0.479964in}}%
\pgfpathlineto{\pgfqpoint{4.434858in}{0.479964in}}%
\pgfpathlineto{\pgfqpoint{4.434858in}{0.470400in}}%
\pgfpathlineto{\pgfqpoint{4.380714in}{0.470400in}}%
\pgfpathlineto{\pgfqpoint{4.380714in}{0.554445in}}%
\pgfpathlineto{\pgfqpoint{4.380714in}{0.554445in}}%
\pgfpathclose%
\pgfpathmoveto{\pgfqpoint{4.455752in}{0.530002in}}%
\pgfpathlineto{\pgfqpoint{4.467622in}{0.530002in}}%
\pgfpathlineto{\pgfqpoint{4.467622in}{0.515719in}}%
\pgfpathlineto{\pgfqpoint{4.455752in}{0.515719in}}%
\pgfpathlineto{\pgfqpoint{4.455752in}{0.530002in}}%
\pgfpathlineto{\pgfqpoint{4.455752in}{0.530002in}}%
\pgfpathclose%
\pgfpathmoveto{\pgfqpoint{4.455752in}{0.484702in}}%
\pgfpathlineto{\pgfqpoint{4.467622in}{0.484702in}}%
\pgfpathlineto{\pgfqpoint{4.467622in}{0.475011in}}%
\pgfpathlineto{\pgfqpoint{4.458400in}{0.456999in}}%
\pgfpathlineto{\pgfqpoint{4.451141in}{0.456999in}}%
\pgfpathlineto{\pgfqpoint{4.455752in}{0.475011in}}%
\pgfpathlineto{\pgfqpoint{4.455752in}{0.484702in}}%
\pgfpathlineto{\pgfqpoint{4.455752in}{0.484702in}}%
\pgfpathclose%
\pgfusepath{fill}%
\end{pgfscope}%
\begin{pgfscope}%
\pgfsetbuttcap%
\pgfsetroundjoin%
\definecolor{currentfill}{rgb}{0.309804,0.372549,0.419608}%
\pgfsetfillcolor{currentfill}%
\pgfsetlinewidth{0.000000pt}%
\definecolor{currentstroke}{rgb}{0.309804,0.372549,0.419608}%
\pgfsetstrokecolor{currentstroke}%
\pgfsetdash{}{0pt}%
\pgfpathmoveto{\pgfqpoint{4.564641in}{0.551347in}}%
\pgfpathlineto{\pgfqpoint{4.564641in}{0.533443in}}%
\pgfpathlineto{\pgfqpoint{4.585968in}{0.533443in}}%
\pgfpathlineto{\pgfqpoint{4.585968in}{0.525391in}}%
\pgfpathlineto{\pgfqpoint{4.564641in}{0.525391in}}%
\pgfpathlineto{\pgfqpoint{4.564641in}{0.491168in}}%
\pgfpathquadraticcurveto{\pgfqpoint{4.564641in}{0.483459in}}{\pgfqpoint{4.566749in}{0.481261in}}%
\pgfpathquadraticcurveto{\pgfqpoint{4.568856in}{0.479064in}}{\pgfqpoint{4.575341in}{0.479064in}}%
\pgfpathlineto{\pgfqpoint{4.585968in}{0.479064in}}%
\pgfpathlineto{\pgfqpoint{4.585968in}{0.470400in}}%
\pgfpathlineto{\pgfqpoint{4.575341in}{0.470400in}}%
\pgfpathquadraticcurveto{\pgfqpoint{4.563345in}{0.470400in}}{\pgfqpoint{4.558788in}{0.474867in}}%
\pgfpathquadraticcurveto{\pgfqpoint{4.554230in}{0.479352in}}{\pgfqpoint{4.554230in}{0.491168in}}%
\pgfpathlineto{\pgfqpoint{4.554230in}{0.525391in}}%
\pgfpathlineto{\pgfqpoint{4.546629in}{0.525391in}}%
\pgfpathlineto{\pgfqpoint{4.546629in}{0.533443in}}%
\pgfpathlineto{\pgfqpoint{4.554230in}{0.533443in}}%
\pgfpathlineto{\pgfqpoint{4.554230in}{0.551347in}}%
\pgfpathlineto{\pgfqpoint{4.564641in}{0.551347in}}%
\pgfpathlineto{\pgfqpoint{4.564641in}{0.551347in}}%
\pgfpathclose%
\pgfpathmoveto{\pgfqpoint{4.652000in}{0.508460in}}%
\pgfpathlineto{\pgfqpoint{4.652000in}{0.470400in}}%
\pgfpathlineto{\pgfqpoint{4.641643in}{0.470400in}}%
\pgfpathlineto{\pgfqpoint{4.641643in}{0.508117in}}%
\pgfpathquadraticcurveto{\pgfqpoint{4.641643in}{0.517069in}}{\pgfqpoint{4.638149in}{0.521500in}}%
\pgfpathquadraticcurveto{\pgfqpoint{4.634655in}{0.525949in}}{\pgfqpoint{4.627684in}{0.525949in}}%
\pgfpathquadraticcurveto{\pgfqpoint{4.619290in}{0.525949in}}{\pgfqpoint{4.614445in}{0.520600in}}%
\pgfpathquadraticcurveto{\pgfqpoint{4.609600in}{0.515268in}}{\pgfqpoint{4.609600in}{0.506028in}}%
\pgfpathlineto{\pgfqpoint{4.609600in}{0.470400in}}%
\pgfpathlineto{\pgfqpoint{4.599189in}{0.470400in}}%
\pgfpathlineto{\pgfqpoint{4.599189in}{0.557993in}}%
\pgfpathlineto{\pgfqpoint{4.609600in}{0.557993in}}%
\pgfpathlineto{\pgfqpoint{4.609600in}{0.523644in}}%
\pgfpathquadraticcurveto{\pgfqpoint{4.613328in}{0.529336in}}{\pgfqpoint{4.618354in}{0.532146in}}%
\pgfpathquadraticcurveto{\pgfqpoint{4.623397in}{0.534956in}}{\pgfqpoint{4.629990in}{0.534956in}}%
\pgfpathquadraticcurveto{\pgfqpoint{4.640851in}{0.534956in}}{\pgfqpoint{4.646417in}{0.528237in}}%
\pgfpathquadraticcurveto{\pgfqpoint{4.652000in}{0.521518in}}{\pgfqpoint{4.652000in}{0.508460in}}%
\pgfpathlineto{\pgfqpoint{4.652000in}{0.508460in}}%
\pgfpathclose%
\pgfpathmoveto{\pgfqpoint{4.726571in}{0.504515in}}%
\pgfpathlineto{\pgfqpoint{4.726571in}{0.499454in}}%
\pgfpathlineto{\pgfqpoint{4.678947in}{0.499454in}}%
\pgfpathquadraticcurveto{\pgfqpoint{4.679631in}{0.488754in}}{\pgfqpoint{4.685395in}{0.483153in}}%
\pgfpathquadraticcurveto{\pgfqpoint{4.691159in}{0.477551in}}{\pgfqpoint{4.701462in}{0.477551in}}%
\pgfpathquadraticcurveto{\pgfqpoint{4.707424in}{0.477551in}}{\pgfqpoint{4.713026in}{0.479010in}}%
\pgfpathquadraticcurveto{\pgfqpoint{4.718627in}{0.480469in}}{\pgfqpoint{4.724157in}{0.483405in}}%
\pgfpathlineto{\pgfqpoint{4.724157in}{0.473606in}}%
\pgfpathquadraticcurveto{\pgfqpoint{4.718573in}{0.471247in}}{\pgfqpoint{4.712719in}{0.470004in}}%
\pgfpathquadraticcurveto{\pgfqpoint{4.706865in}{0.468761in}}{\pgfqpoint{4.700849in}{0.468761in}}%
\pgfpathquadraticcurveto{\pgfqpoint{4.685755in}{0.468761in}}{\pgfqpoint{4.676947in}{0.477533in}}%
\pgfpathquadraticcurveto{\pgfqpoint{4.668139in}{0.486323in}}{\pgfqpoint{4.668139in}{0.501309in}}%
\pgfpathquadraticcurveto{\pgfqpoint{4.668139in}{0.516781in}}{\pgfqpoint{4.676497in}{0.525859in}}%
\pgfpathquadraticcurveto{\pgfqpoint{4.684855in}{0.534956in}}{\pgfqpoint{4.699048in}{0.534956in}}%
\pgfpathquadraticcurveto{\pgfqpoint{4.711765in}{0.534956in}}{\pgfqpoint{4.719168in}{0.526760in}}%
\pgfpathquadraticcurveto{\pgfqpoint{4.726571in}{0.518583in}}{\pgfqpoint{4.726571in}{0.504515in}}%
\pgfpathlineto{\pgfqpoint{4.726571in}{0.504515in}}%
\pgfpathclose%
\pgfpathmoveto{\pgfqpoint{4.716214in}{0.507559in}}%
\pgfpathquadraticcurveto{\pgfqpoint{4.716106in}{0.516043in}}{\pgfqpoint{4.711459in}{0.521104in}}%
\pgfpathquadraticcurveto{\pgfqpoint{4.706811in}{0.526184in}}{\pgfqpoint{4.699156in}{0.526184in}}%
\pgfpathquadraticcurveto{\pgfqpoint{4.690492in}{0.526184in}}{\pgfqpoint{4.685287in}{0.521284in}}%
\pgfpathquadraticcurveto{\pgfqpoint{4.680081in}{0.516385in}}{\pgfqpoint{4.679289in}{0.507487in}}%
\pgfpathlineto{\pgfqpoint{4.716214in}{0.507559in}}%
\pgfpathlineto{\pgfqpoint{4.716214in}{0.507559in}}%
\pgfpathclose%
\pgfusepath{fill}%
\end{pgfscope}%
\begin{pgfscope}%
\pgfsetbuttcap%
\pgfsetroundjoin%
\definecolor{currentfill}{rgb}{0.309804,0.372549,0.419608}%
\pgfsetfillcolor{currentfill}%
\pgfsetlinewidth{0.000000pt}%
\definecolor{currentstroke}{rgb}{0.309804,0.372549,0.419608}%
\pgfsetstrokecolor{currentstroke}%
\pgfsetdash{}{0pt}%
\pgfpathmoveto{\pgfqpoint{4.811720in}{0.479856in}}%
\pgfpathlineto{\pgfqpoint{4.811720in}{0.446426in}}%
\pgfpathlineto{\pgfqpoint{4.801309in}{0.446426in}}%
\pgfpathlineto{\pgfqpoint{4.801309in}{0.533443in}}%
\pgfpathlineto{\pgfqpoint{4.811720in}{0.533443in}}%
\pgfpathlineto{\pgfqpoint{4.811720in}{0.523878in}}%
\pgfpathquadraticcurveto{\pgfqpoint{4.814998in}{0.529498in}}{\pgfqpoint{4.819969in}{0.532218in}}%
\pgfpathquadraticcurveto{\pgfqpoint{4.824959in}{0.534956in}}{\pgfqpoint{4.831875in}{0.534956in}}%
\pgfpathquadraticcurveto{\pgfqpoint{4.843367in}{0.534956in}}{\pgfqpoint{4.850536in}{0.525841in}}%
\pgfpathquadraticcurveto{\pgfqpoint{4.857723in}{0.516727in}}{\pgfqpoint{4.857723in}{0.501867in}}%
\pgfpathquadraticcurveto{\pgfqpoint{4.857723in}{0.487007in}}{\pgfqpoint{4.850536in}{0.477875in}}%
\pgfpathquadraticcurveto{\pgfqpoint{4.843367in}{0.468761in}}{\pgfqpoint{4.831875in}{0.468761in}}%
\pgfpathquadraticcurveto{\pgfqpoint{4.824959in}{0.468761in}}{\pgfqpoint{4.819969in}{0.471499in}}%
\pgfpathquadraticcurveto{\pgfqpoint{4.814998in}{0.474237in}}{\pgfqpoint{4.811720in}{0.479856in}}%
\pgfpathlineto{\pgfqpoint{4.811720in}{0.479856in}}%
\pgfpathclose%
\pgfpathmoveto{\pgfqpoint{4.846970in}{0.501867in}}%
\pgfpathquadraticcurveto{\pgfqpoint{4.846970in}{0.513287in}}{\pgfqpoint{4.842268in}{0.519789in}}%
\pgfpathquadraticcurveto{\pgfqpoint{4.837567in}{0.526292in}}{\pgfqpoint{4.829354in}{0.526292in}}%
\pgfpathquadraticcurveto{\pgfqpoint{4.821122in}{0.526292in}}{\pgfqpoint{4.816421in}{0.519789in}}%
\pgfpathquadraticcurveto{\pgfqpoint{4.811720in}{0.513287in}}{\pgfqpoint{4.811720in}{0.501867in}}%
\pgfpathquadraticcurveto{\pgfqpoint{4.811720in}{0.490448in}}{\pgfqpoint{4.816421in}{0.483945in}}%
\pgfpathquadraticcurveto{\pgfqpoint{4.821122in}{0.477443in}}{\pgfqpoint{4.829354in}{0.477443in}}%
\pgfpathquadraticcurveto{\pgfqpoint{4.837567in}{0.477443in}}{\pgfqpoint{4.842268in}{0.483945in}}%
\pgfpathquadraticcurveto{\pgfqpoint{4.846970in}{0.490448in}}{\pgfqpoint{4.846970in}{0.501867in}}%
\pgfpathlineto{\pgfqpoint{4.846970in}{0.501867in}}%
\pgfpathclose%
\pgfpathmoveto{\pgfqpoint{4.874888in}{0.557993in}}%
\pgfpathlineto{\pgfqpoint{4.885245in}{0.557993in}}%
\pgfpathlineto{\pgfqpoint{4.885245in}{0.470400in}}%
\pgfpathlineto{\pgfqpoint{4.874888in}{0.470400in}}%
\pgfpathlineto{\pgfqpoint{4.874888in}{0.557993in}}%
\pgfpathlineto{\pgfqpoint{4.874888in}{0.557993in}}%
\pgfpathclose%
\pgfpathmoveto{\pgfqpoint{4.931339in}{0.526184in}}%
\pgfpathquadraticcurveto{\pgfqpoint{4.923017in}{0.526184in}}{\pgfqpoint{4.918172in}{0.519681in}}%
\pgfpathquadraticcurveto{\pgfqpoint{4.913326in}{0.513179in}}{\pgfqpoint{4.913326in}{0.501867in}}%
\pgfpathquadraticcurveto{\pgfqpoint{4.913326in}{0.490556in}}{\pgfqpoint{4.918136in}{0.484053in}}%
\pgfpathquadraticcurveto{\pgfqpoint{4.922963in}{0.477551in}}{\pgfqpoint{4.931339in}{0.477551in}}%
\pgfpathquadraticcurveto{\pgfqpoint{4.939624in}{0.477551in}}{\pgfqpoint{4.944451in}{0.484071in}}%
\pgfpathquadraticcurveto{\pgfqpoint{4.949297in}{0.490610in}}{\pgfqpoint{4.949297in}{0.501867in}}%
\pgfpathquadraticcurveto{\pgfqpoint{4.949297in}{0.513071in}}{\pgfqpoint{4.944451in}{0.519627in}}%
\pgfpathquadraticcurveto{\pgfqpoint{4.939624in}{0.526184in}}{\pgfqpoint{4.931339in}{0.526184in}}%
\pgfpathlineto{\pgfqpoint{4.931339in}{0.526184in}}%
\pgfpathclose%
\pgfpathmoveto{\pgfqpoint{4.931339in}{0.534956in}}%
\pgfpathquadraticcurveto{\pgfqpoint{4.944848in}{0.534956in}}{\pgfqpoint{4.952557in}{0.526166in}}%
\pgfpathquadraticcurveto{\pgfqpoint{4.960284in}{0.517394in}}{\pgfqpoint{4.960284in}{0.501867in}}%
\pgfpathquadraticcurveto{\pgfqpoint{4.960284in}{0.486395in}}{\pgfqpoint{4.952557in}{0.477569in}}%
\pgfpathquadraticcurveto{\pgfqpoint{4.944848in}{0.468761in}}{\pgfqpoint{4.931339in}{0.468761in}}%
\pgfpathquadraticcurveto{\pgfqpoint{4.917775in}{0.468761in}}{\pgfqpoint{4.910084in}{0.477569in}}%
\pgfpathquadraticcurveto{\pgfqpoint{4.902411in}{0.486395in}}{\pgfqpoint{4.902411in}{0.501867in}}%
\pgfpathquadraticcurveto{\pgfqpoint{4.902411in}{0.517394in}}{\pgfqpoint{4.910084in}{0.526166in}}%
\pgfpathquadraticcurveto{\pgfqpoint{4.917775in}{0.534956in}}{\pgfqpoint{4.931339in}{0.534956in}}%
\pgfpathlineto{\pgfqpoint{4.931339in}{0.534956in}}%
\pgfpathclose%
\pgfpathmoveto{\pgfqpoint{4.987699in}{0.551347in}}%
\pgfpathlineto{\pgfqpoint{4.987699in}{0.533443in}}%
\pgfpathlineto{\pgfqpoint{5.009025in}{0.533443in}}%
\pgfpathlineto{\pgfqpoint{5.009025in}{0.525391in}}%
\pgfpathlineto{\pgfqpoint{4.987699in}{0.525391in}}%
\pgfpathlineto{\pgfqpoint{4.987699in}{0.491168in}}%
\pgfpathquadraticcurveto{\pgfqpoint{4.987699in}{0.483459in}}{\pgfqpoint{4.989806in}{0.481261in}}%
\pgfpathquadraticcurveto{\pgfqpoint{4.991913in}{0.479064in}}{\pgfqpoint{4.998398in}{0.479064in}}%
\pgfpathlineto{\pgfqpoint{5.009025in}{0.479064in}}%
\pgfpathlineto{\pgfqpoint{5.009025in}{0.470400in}}%
\pgfpathlineto{\pgfqpoint{4.998398in}{0.470400in}}%
\pgfpathquadraticcurveto{\pgfqpoint{4.986402in}{0.470400in}}{\pgfqpoint{4.981845in}{0.474867in}}%
\pgfpathquadraticcurveto{\pgfqpoint{4.977288in}{0.479352in}}{\pgfqpoint{4.977288in}{0.491168in}}%
\pgfpathlineto{\pgfqpoint{4.977288in}{0.525391in}}%
\pgfpathlineto{\pgfqpoint{4.969686in}{0.525391in}}%
\pgfpathlineto{\pgfqpoint{4.969686in}{0.533443in}}%
\pgfpathlineto{\pgfqpoint{4.977288in}{0.533443in}}%
\pgfpathlineto{\pgfqpoint{4.977288in}{0.551347in}}%
\pgfpathlineto{\pgfqpoint{4.987699in}{0.551347in}}%
\pgfpathlineto{\pgfqpoint{4.987699in}{0.551347in}}%
\pgfpathclose%
\pgfpathmoveto{\pgfqpoint{5.032891in}{0.551347in}}%
\pgfpathlineto{\pgfqpoint{5.032891in}{0.533443in}}%
\pgfpathlineto{\pgfqpoint{5.054217in}{0.533443in}}%
\pgfpathlineto{\pgfqpoint{5.054217in}{0.525391in}}%
\pgfpathlineto{\pgfqpoint{5.032891in}{0.525391in}}%
\pgfpathlineto{\pgfqpoint{5.032891in}{0.491168in}}%
\pgfpathquadraticcurveto{\pgfqpoint{5.032891in}{0.483459in}}{\pgfqpoint{5.034998in}{0.481261in}}%
\pgfpathquadraticcurveto{\pgfqpoint{5.037106in}{0.479064in}}{\pgfqpoint{5.043590in}{0.479064in}}%
\pgfpathlineto{\pgfqpoint{5.054217in}{0.479064in}}%
\pgfpathlineto{\pgfqpoint{5.054217in}{0.470400in}}%
\pgfpathlineto{\pgfqpoint{5.043590in}{0.470400in}}%
\pgfpathquadraticcurveto{\pgfqpoint{5.031594in}{0.470400in}}{\pgfqpoint{5.027037in}{0.474867in}}%
\pgfpathquadraticcurveto{\pgfqpoint{5.022480in}{0.479352in}}{\pgfqpoint{5.022480in}{0.491168in}}%
\pgfpathlineto{\pgfqpoint{5.022480in}{0.525391in}}%
\pgfpathlineto{\pgfqpoint{5.014879in}{0.525391in}}%
\pgfpathlineto{\pgfqpoint{5.014879in}{0.533443in}}%
\pgfpathlineto{\pgfqpoint{5.022480in}{0.533443in}}%
\pgfpathlineto{\pgfqpoint{5.022480in}{0.551347in}}%
\pgfpathlineto{\pgfqpoint{5.032891in}{0.551347in}}%
\pgfpathlineto{\pgfqpoint{5.032891in}{0.551347in}}%
\pgfpathclose%
\pgfpathmoveto{\pgfqpoint{5.121763in}{0.504515in}}%
\pgfpathlineto{\pgfqpoint{5.121763in}{0.499454in}}%
\pgfpathlineto{\pgfqpoint{5.074139in}{0.499454in}}%
\pgfpathquadraticcurveto{\pgfqpoint{5.074823in}{0.488754in}}{\pgfqpoint{5.080587in}{0.483153in}}%
\pgfpathquadraticcurveto{\pgfqpoint{5.086351in}{0.477551in}}{\pgfqpoint{5.096654in}{0.477551in}}%
\pgfpathquadraticcurveto{\pgfqpoint{5.102616in}{0.477551in}}{\pgfqpoint{5.108218in}{0.479010in}}%
\pgfpathquadraticcurveto{\pgfqpoint{5.113820in}{0.480469in}}{\pgfqpoint{5.119349in}{0.483405in}}%
\pgfpathlineto{\pgfqpoint{5.119349in}{0.473606in}}%
\pgfpathquadraticcurveto{\pgfqpoint{5.113766in}{0.471247in}}{\pgfqpoint{5.107912in}{0.470004in}}%
\pgfpathquadraticcurveto{\pgfqpoint{5.102058in}{0.468761in}}{\pgfqpoint{5.096042in}{0.468761in}}%
\pgfpathquadraticcurveto{\pgfqpoint{5.080947in}{0.468761in}}{\pgfqpoint{5.072140in}{0.477533in}}%
\pgfpathquadraticcurveto{\pgfqpoint{5.063332in}{0.486323in}}{\pgfqpoint{5.063332in}{0.501309in}}%
\pgfpathquadraticcurveto{\pgfqpoint{5.063332in}{0.516781in}}{\pgfqpoint{5.071689in}{0.525859in}}%
\pgfpathquadraticcurveto{\pgfqpoint{5.080047in}{0.534956in}}{\pgfqpoint{5.094240in}{0.534956in}}%
\pgfpathquadraticcurveto{\pgfqpoint{5.106957in}{0.534956in}}{\pgfqpoint{5.114360in}{0.526760in}}%
\pgfpathquadraticcurveto{\pgfqpoint{5.121763in}{0.518583in}}{\pgfqpoint{5.121763in}{0.504515in}}%
\pgfpathlineto{\pgfqpoint{5.121763in}{0.504515in}}%
\pgfpathclose%
\pgfpathmoveto{\pgfqpoint{5.111406in}{0.507559in}}%
\pgfpathquadraticcurveto{\pgfqpoint{5.111298in}{0.516043in}}{\pgfqpoint{5.106651in}{0.521104in}}%
\pgfpathquadraticcurveto{\pgfqpoint{5.102004in}{0.526184in}}{\pgfqpoint{5.094349in}{0.526184in}}%
\pgfpathquadraticcurveto{\pgfqpoint{5.085685in}{0.526184in}}{\pgfqpoint{5.080479in}{0.521284in}}%
\pgfpathquadraticcurveto{\pgfqpoint{5.075274in}{0.516385in}}{\pgfqpoint{5.074481in}{0.507487in}}%
\pgfpathlineto{\pgfqpoint{5.111406in}{0.507559in}}%
\pgfpathlineto{\pgfqpoint{5.111406in}{0.507559in}}%
\pgfpathclose%
\pgfpathmoveto{\pgfqpoint{5.180248in}{0.523878in}}%
\pgfpathlineto{\pgfqpoint{5.180248in}{0.557993in}}%
\pgfpathlineto{\pgfqpoint{5.190605in}{0.557993in}}%
\pgfpathlineto{\pgfqpoint{5.190605in}{0.470400in}}%
\pgfpathlineto{\pgfqpoint{5.180248in}{0.470400in}}%
\pgfpathlineto{\pgfqpoint{5.180248in}{0.479856in}}%
\pgfpathquadraticcurveto{\pgfqpoint{5.176988in}{0.474237in}}{\pgfqpoint{5.171999in}{0.471499in}}%
\pgfpathquadraticcurveto{\pgfqpoint{5.167028in}{0.468761in}}{\pgfqpoint{5.160039in}{0.468761in}}%
\pgfpathquadraticcurveto{\pgfqpoint{5.148619in}{0.468761in}}{\pgfqpoint{5.141432in}{0.477875in}}%
\pgfpathquadraticcurveto{\pgfqpoint{5.134263in}{0.487007in}}{\pgfqpoint{5.134263in}{0.501867in}}%
\pgfpathquadraticcurveto{\pgfqpoint{5.134263in}{0.516727in}}{\pgfqpoint{5.141432in}{0.525841in}}%
\pgfpathquadraticcurveto{\pgfqpoint{5.148619in}{0.534956in}}{\pgfqpoint{5.160039in}{0.534956in}}%
\pgfpathquadraticcurveto{\pgfqpoint{5.167028in}{0.534956in}}{\pgfqpoint{5.171999in}{0.532218in}}%
\pgfpathquadraticcurveto{\pgfqpoint{5.176988in}{0.529498in}}{\pgfqpoint{5.180248in}{0.523878in}}%
\pgfpathlineto{\pgfqpoint{5.180248in}{0.523878in}}%
\pgfpathclose%
\pgfpathmoveto{\pgfqpoint{5.144963in}{0.501867in}}%
\pgfpathquadraticcurveto{\pgfqpoint{5.144963in}{0.490448in}}{\pgfqpoint{5.149664in}{0.483945in}}%
\pgfpathquadraticcurveto{\pgfqpoint{5.154365in}{0.477443in}}{\pgfqpoint{5.162579in}{0.477443in}}%
\pgfpathquadraticcurveto{\pgfqpoint{5.170792in}{0.477443in}}{\pgfqpoint{5.175511in}{0.483945in}}%
\pgfpathquadraticcurveto{\pgfqpoint{5.180248in}{0.490448in}}{\pgfqpoint{5.180248in}{0.501867in}}%
\pgfpathquadraticcurveto{\pgfqpoint{5.180248in}{0.513287in}}{\pgfqpoint{5.175511in}{0.519789in}}%
\pgfpathquadraticcurveto{\pgfqpoint{5.170792in}{0.526292in}}{\pgfqpoint{5.162579in}{0.526292in}}%
\pgfpathquadraticcurveto{\pgfqpoint{5.154365in}{0.526292in}}{\pgfqpoint{5.149664in}{0.519789in}}%
\pgfpathquadraticcurveto{\pgfqpoint{5.144963in}{0.513287in}}{\pgfqpoint{5.144963in}{0.501867in}}%
\pgfpathlineto{\pgfqpoint{5.144963in}{0.501867in}}%
\pgfpathclose%
\pgfusepath{fill}%
\end{pgfscope}%
\begin{pgfscope}%
\pgfsetbuttcap%
\pgfsetroundjoin%
\definecolor{currentfill}{rgb}{0.309804,0.372549,0.419608}%
\pgfsetfillcolor{currentfill}%
\pgfsetlinewidth{0.000000pt}%
\definecolor{currentstroke}{rgb}{0.309804,0.372549,0.419608}%
\pgfsetstrokecolor{currentstroke}%
\pgfsetdash{}{0pt}%
\pgfpathmoveto{\pgfqpoint{5.320500in}{0.502660in}}%
\pgfpathquadraticcurveto{\pgfqpoint{5.320500in}{0.513917in}}{\pgfqpoint{5.315852in}{0.520096in}}%
\pgfpathquadraticcurveto{\pgfqpoint{5.311223in}{0.526292in}}{\pgfqpoint{5.302830in}{0.526292in}}%
\pgfpathquadraticcurveto{\pgfqpoint{5.294508in}{0.526292in}}{\pgfqpoint{5.289861in}{0.520096in}}%
\pgfpathquadraticcurveto{\pgfqpoint{5.285214in}{0.513917in}}{\pgfqpoint{5.285214in}{0.502660in}}%
\pgfpathquadraticcurveto{\pgfqpoint{5.285214in}{0.491456in}}{\pgfqpoint{5.289861in}{0.485260in}}%
\pgfpathquadraticcurveto{\pgfqpoint{5.294508in}{0.479064in}}{\pgfqpoint{5.302830in}{0.479064in}}%
\pgfpathquadraticcurveto{\pgfqpoint{5.311223in}{0.479064in}}{\pgfqpoint{5.315852in}{0.485260in}}%
\pgfpathquadraticcurveto{\pgfqpoint{5.320500in}{0.491456in}}{\pgfqpoint{5.320500in}{0.502660in}}%
\pgfpathlineto{\pgfqpoint{5.320500in}{0.502660in}}%
\pgfpathclose%
\pgfpathmoveto{\pgfqpoint{5.330857in}{0.478217in}}%
\pgfpathquadraticcurveto{\pgfqpoint{5.330857in}{0.462132in}}{\pgfqpoint{5.323706in}{0.454279in}}%
\pgfpathquadraticcurveto{\pgfqpoint{5.316573in}{0.446426in}}{\pgfqpoint{5.301821in}{0.446426in}}%
\pgfpathquadraticcurveto{\pgfqpoint{5.296363in}{0.446426in}}{\pgfqpoint{5.291518in}{0.447236in}}%
\pgfpathquadraticcurveto{\pgfqpoint{5.286673in}{0.448047in}}{\pgfqpoint{5.282116in}{0.449740in}}%
\pgfpathlineto{\pgfqpoint{5.282116in}{0.459809in}}%
\pgfpathquadraticcurveto{\pgfqpoint{5.286673in}{0.457341in}}{\pgfqpoint{5.291122in}{0.456170in}}%
\pgfpathquadraticcurveto{\pgfqpoint{5.295571in}{0.454982in}}{\pgfqpoint{5.300182in}{0.454982in}}%
\pgfpathquadraticcurveto{\pgfqpoint{5.310377in}{0.454982in}}{\pgfqpoint{5.315438in}{0.460295in}}%
\pgfpathquadraticcurveto{\pgfqpoint{5.320500in}{0.465609in}}{\pgfqpoint{5.320500in}{0.476362in}}%
\pgfpathlineto{\pgfqpoint{5.320500in}{0.481495in}}%
\pgfpathquadraticcurveto{\pgfqpoint{5.317293in}{0.475912in}}{\pgfqpoint{5.312286in}{0.473156in}}%
\pgfpathquadraticcurveto{\pgfqpoint{5.307279in}{0.470400in}}{\pgfqpoint{5.300290in}{0.470400in}}%
\pgfpathquadraticcurveto{\pgfqpoint{5.288708in}{0.470400in}}{\pgfqpoint{5.281611in}{0.479226in}}%
\pgfpathquadraticcurveto{\pgfqpoint{5.274515in}{0.488070in}}{\pgfqpoint{5.274515in}{0.502660in}}%
\pgfpathquadraticcurveto{\pgfqpoint{5.274515in}{0.517286in}}{\pgfqpoint{5.281611in}{0.526112in}}%
\pgfpathquadraticcurveto{\pgfqpoint{5.288708in}{0.534956in}}{\pgfqpoint{5.300290in}{0.534956in}}%
\pgfpathquadraticcurveto{\pgfqpoint{5.307279in}{0.534956in}}{\pgfqpoint{5.312286in}{0.532200in}}%
\pgfpathquadraticcurveto{\pgfqpoint{5.317293in}{0.529444in}}{\pgfqpoint{5.320500in}{0.523878in}}%
\pgfpathlineto{\pgfqpoint{5.320500in}{0.533443in}}%
\pgfpathlineto{\pgfqpoint{5.330857in}{0.533443in}}%
\pgfpathlineto{\pgfqpoint{5.330857in}{0.478217in}}%
\pgfpathlineto{\pgfqpoint{5.330857in}{0.478217in}}%
\pgfpathclose%
\pgfpathmoveto{\pgfqpoint{5.380858in}{0.502083in}}%
\pgfpathquadraticcurveto{\pgfqpoint{5.368304in}{0.502083in}}{\pgfqpoint{5.363459in}{0.499219in}}%
\pgfpathquadraticcurveto{\pgfqpoint{5.358613in}{0.496356in}}{\pgfqpoint{5.358613in}{0.489421in}}%
\pgfpathquadraticcurveto{\pgfqpoint{5.358613in}{0.483909in}}{\pgfqpoint{5.362252in}{0.480667in}}%
\pgfpathquadraticcurveto{\pgfqpoint{5.365890in}{0.477443in}}{\pgfqpoint{5.372122in}{0.477443in}}%
\pgfpathquadraticcurveto{\pgfqpoint{5.380750in}{0.477443in}}{\pgfqpoint{5.385956in}{0.483549in}}%
\pgfpathquadraticcurveto{\pgfqpoint{5.391161in}{0.489655in}}{\pgfqpoint{5.391161in}{0.499778in}}%
\pgfpathlineto{\pgfqpoint{5.391161in}{0.502083in}}%
\pgfpathlineto{\pgfqpoint{5.380858in}{0.502083in}}%
\pgfpathlineto{\pgfqpoint{5.380858in}{0.502083in}}%
\pgfpathclose%
\pgfpathmoveto{\pgfqpoint{5.401518in}{0.506370in}}%
\pgfpathlineto{\pgfqpoint{5.401518in}{0.470400in}}%
\pgfpathlineto{\pgfqpoint{5.391161in}{0.470400in}}%
\pgfpathlineto{\pgfqpoint{5.391161in}{0.479964in}}%
\pgfpathquadraticcurveto{\pgfqpoint{5.387613in}{0.474237in}}{\pgfqpoint{5.382317in}{0.471499in}}%
\pgfpathquadraticcurveto{\pgfqpoint{5.377022in}{0.468761in}}{\pgfqpoint{5.369367in}{0.468761in}}%
\pgfpathquadraticcurveto{\pgfqpoint{5.359694in}{0.468761in}}{\pgfqpoint{5.353966in}{0.474201in}}%
\pgfpathquadraticcurveto{\pgfqpoint{5.348256in}{0.479640in}}{\pgfqpoint{5.348256in}{0.488754in}}%
\pgfpathquadraticcurveto{\pgfqpoint{5.348256in}{0.499382in}}{\pgfqpoint{5.355371in}{0.504785in}}%
\pgfpathquadraticcurveto{\pgfqpoint{5.362504in}{0.510189in}}{\pgfqpoint{5.376625in}{0.510189in}}%
\pgfpathlineto{\pgfqpoint{5.391161in}{0.510189in}}%
\pgfpathlineto{\pgfqpoint{5.391161in}{0.511216in}}%
\pgfpathquadraticcurveto{\pgfqpoint{5.391161in}{0.518366in}}{\pgfqpoint{5.386460in}{0.522275in}}%
\pgfpathquadraticcurveto{\pgfqpoint{5.381759in}{0.526184in}}{\pgfqpoint{5.373257in}{0.526184in}}%
\pgfpathquadraticcurveto{\pgfqpoint{5.367854in}{0.526184in}}{\pgfqpoint{5.362720in}{0.524887in}}%
\pgfpathquadraticcurveto{\pgfqpoint{5.357605in}{0.523590in}}{\pgfqpoint{5.352885in}{0.520996in}}%
\pgfpathlineto{\pgfqpoint{5.352885in}{0.530579in}}%
\pgfpathquadraticcurveto{\pgfqpoint{5.358559in}{0.532776in}}{\pgfqpoint{5.363909in}{0.533857in}}%
\pgfpathquadraticcurveto{\pgfqpoint{5.369258in}{0.534956in}}{\pgfqpoint{5.374320in}{0.534956in}}%
\pgfpathquadraticcurveto{\pgfqpoint{5.388009in}{0.534956in}}{\pgfqpoint{5.394764in}{0.527859in}}%
\pgfpathquadraticcurveto{\pgfqpoint{5.401518in}{0.520780in}}{\pgfqpoint{5.401518in}{0.506370in}}%
\pgfpathlineto{\pgfqpoint{5.401518in}{0.506370in}}%
\pgfpathclose%
\pgfpathmoveto{\pgfqpoint{5.432859in}{0.479856in}}%
\pgfpathlineto{\pgfqpoint{5.432859in}{0.446426in}}%
\pgfpathlineto{\pgfqpoint{5.422448in}{0.446426in}}%
\pgfpathlineto{\pgfqpoint{5.422448in}{0.533443in}}%
\pgfpathlineto{\pgfqpoint{5.432859in}{0.533443in}}%
\pgfpathlineto{\pgfqpoint{5.432859in}{0.523878in}}%
\pgfpathquadraticcurveto{\pgfqpoint{5.436138in}{0.529498in}}{\pgfqpoint{5.441109in}{0.532218in}}%
\pgfpathquadraticcurveto{\pgfqpoint{5.446098in}{0.534956in}}{\pgfqpoint{5.453015in}{0.534956in}}%
\pgfpathquadraticcurveto{\pgfqpoint{5.464507in}{0.534956in}}{\pgfqpoint{5.471676in}{0.525841in}}%
\pgfpathquadraticcurveto{\pgfqpoint{5.478862in}{0.516727in}}{\pgfqpoint{5.478862in}{0.501867in}}%
\pgfpathquadraticcurveto{\pgfqpoint{5.478862in}{0.487007in}}{\pgfqpoint{5.471676in}{0.477875in}}%
\pgfpathquadraticcurveto{\pgfqpoint{5.464507in}{0.468761in}}{\pgfqpoint{5.453015in}{0.468761in}}%
\pgfpathquadraticcurveto{\pgfqpoint{5.446098in}{0.468761in}}{\pgfqpoint{5.441109in}{0.471499in}}%
\pgfpathquadraticcurveto{\pgfqpoint{5.436138in}{0.474237in}}{\pgfqpoint{5.432859in}{0.479856in}}%
\pgfpathlineto{\pgfqpoint{5.432859in}{0.479856in}}%
\pgfpathclose%
\pgfpathmoveto{\pgfqpoint{5.468109in}{0.501867in}}%
\pgfpathquadraticcurveto{\pgfqpoint{5.468109in}{0.513287in}}{\pgfqpoint{5.463408in}{0.519789in}}%
\pgfpathquadraticcurveto{\pgfqpoint{5.458707in}{0.526292in}}{\pgfqpoint{5.450493in}{0.526292in}}%
\pgfpathquadraticcurveto{\pgfqpoint{5.442262in}{0.526292in}}{\pgfqpoint{5.437561in}{0.519789in}}%
\pgfpathquadraticcurveto{\pgfqpoint{5.432859in}{0.513287in}}{\pgfqpoint{5.432859in}{0.501867in}}%
\pgfpathquadraticcurveto{\pgfqpoint{5.432859in}{0.490448in}}{\pgfqpoint{5.437561in}{0.483945in}}%
\pgfpathquadraticcurveto{\pgfqpoint{5.442262in}{0.477443in}}{\pgfqpoint{5.450493in}{0.477443in}}%
\pgfpathquadraticcurveto{\pgfqpoint{5.458707in}{0.477443in}}{\pgfqpoint{5.463408in}{0.483945in}}%
\pgfpathquadraticcurveto{\pgfqpoint{5.468109in}{0.490448in}}{\pgfqpoint{5.468109in}{0.501867in}}%
\pgfpathlineto{\pgfqpoint{5.468109in}{0.501867in}}%
\pgfpathclose%
\pgfusepath{fill}%
\end{pgfscope}%
\begin{pgfscope}%
\pgfsetbuttcap%
\pgfsetroundjoin%
\definecolor{currentfill}{rgb}{0.309804,0.372549,0.419608}%
\pgfsetfillcolor{currentfill}%
\pgfsetlinewidth{0.000000pt}%
\definecolor{currentstroke}{rgb}{0.309804,0.372549,0.419608}%
\pgfsetstrokecolor{currentstroke}%
\pgfsetdash{}{0pt}%
\pgfpathmoveto{\pgfqpoint{5.556330in}{0.533443in}}%
\pgfpathlineto{\pgfqpoint{5.566687in}{0.533443in}}%
\pgfpathlineto{\pgfqpoint{5.566687in}{0.470400in}}%
\pgfpathlineto{\pgfqpoint{5.556330in}{0.470400in}}%
\pgfpathlineto{\pgfqpoint{5.556330in}{0.533443in}}%
\pgfpathlineto{\pgfqpoint{5.556330in}{0.533443in}}%
\pgfpathclose%
\pgfpathmoveto{\pgfqpoint{5.556330in}{0.557993in}}%
\pgfpathlineto{\pgfqpoint{5.566687in}{0.557993in}}%
\pgfpathlineto{\pgfqpoint{5.566687in}{0.544862in}}%
\pgfpathlineto{\pgfqpoint{5.556330in}{0.544862in}}%
\pgfpathlineto{\pgfqpoint{5.556330in}{0.557993in}}%
\pgfpathlineto{\pgfqpoint{5.556330in}{0.557993in}}%
\pgfpathclose%
\pgfpathmoveto{\pgfqpoint{5.628541in}{0.531587in}}%
\pgfpathlineto{\pgfqpoint{5.628541in}{0.521789in}}%
\pgfpathquadraticcurveto{\pgfqpoint{5.624164in}{0.524040in}}{\pgfqpoint{5.619427in}{0.525157in}}%
\pgfpathquadraticcurveto{\pgfqpoint{5.614707in}{0.526292in}}{\pgfqpoint{5.609628in}{0.526292in}}%
\pgfpathquadraticcurveto{\pgfqpoint{5.601919in}{0.526292in}}{\pgfqpoint{5.598064in}{0.523932in}}%
\pgfpathquadraticcurveto{\pgfqpoint{5.594210in}{0.521573in}}{\pgfqpoint{5.594210in}{0.516835in}}%
\pgfpathquadraticcurveto{\pgfqpoint{5.594210in}{0.513233in}}{\pgfqpoint{5.596965in}{0.511180in}}%
\pgfpathquadraticcurveto{\pgfqpoint{5.599721in}{0.509126in}}{\pgfqpoint{5.608061in}{0.507271in}}%
\pgfpathlineto{\pgfqpoint{5.611609in}{0.506478in}}%
\pgfpathquadraticcurveto{\pgfqpoint{5.622633in}{0.504119in}}{\pgfqpoint{5.627280in}{0.499814in}}%
\pgfpathquadraticcurveto{\pgfqpoint{5.631927in}{0.495509in}}{\pgfqpoint{5.631927in}{0.487800in}}%
\pgfpathquadraticcurveto{\pgfqpoint{5.631927in}{0.479010in}}{\pgfqpoint{5.624974in}{0.473876in}}%
\pgfpathquadraticcurveto{\pgfqpoint{5.618022in}{0.468761in}}{\pgfqpoint{5.605863in}{0.468761in}}%
\pgfpathquadraticcurveto{\pgfqpoint{5.600802in}{0.468761in}}{\pgfqpoint{5.595308in}{0.469752in}}%
\pgfpathquadraticcurveto{\pgfqpoint{5.589815in}{0.470742in}}{\pgfqpoint{5.583745in}{0.472706in}}%
\pgfpathlineto{\pgfqpoint{5.583745in}{0.483405in}}%
\pgfpathquadraticcurveto{\pgfqpoint{5.589490in}{0.480415in}}{\pgfqpoint{5.595056in}{0.478920in}}%
\pgfpathquadraticcurveto{\pgfqpoint{5.600622in}{0.477443in}}{\pgfqpoint{5.606098in}{0.477443in}}%
\pgfpathquadraticcurveto{\pgfqpoint{5.613411in}{0.477443in}}{\pgfqpoint{5.617337in}{0.479946in}}%
\pgfpathquadraticcurveto{\pgfqpoint{5.621282in}{0.482450in}}{\pgfqpoint{5.621282in}{0.487007in}}%
\pgfpathquadraticcurveto{\pgfqpoint{5.621282in}{0.491222in}}{\pgfqpoint{5.618436in}{0.493474in}}%
\pgfpathquadraticcurveto{\pgfqpoint{5.615608in}{0.495725in}}{\pgfqpoint{5.605972in}{0.497814in}}%
\pgfpathlineto{\pgfqpoint{5.602369in}{0.498661in}}%
\pgfpathquadraticcurveto{\pgfqpoint{5.592751in}{0.500678in}}{\pgfqpoint{5.588464in}{0.504875in}}%
\pgfpathquadraticcurveto{\pgfqpoint{5.584195in}{0.509072in}}{\pgfqpoint{5.584195in}{0.516385in}}%
\pgfpathquadraticcurveto{\pgfqpoint{5.584195in}{0.525283in}}{\pgfqpoint{5.590499in}{0.530110in}}%
\pgfpathquadraticcurveto{\pgfqpoint{5.596803in}{0.534956in}}{\pgfqpoint{5.608403in}{0.534956in}}%
\pgfpathquadraticcurveto{\pgfqpoint{5.614131in}{0.534956in}}{\pgfqpoint{5.619192in}{0.534109in}}%
\pgfpathquadraticcurveto{\pgfqpoint{5.624272in}{0.533280in}}{\pgfqpoint{5.628541in}{0.531587in}}%
\pgfpathlineto{\pgfqpoint{5.628541in}{0.531587in}}%
\pgfpathclose%
\pgfusepath{fill}%
\end{pgfscope}%
\begin{pgfscope}%
\pgfsetbuttcap%
\pgfsetroundjoin%
\definecolor{currentfill}{rgb}{0.309804,0.372549,0.419608}%
\pgfsetfillcolor{currentfill}%
\pgfsetlinewidth{0.000000pt}%
\definecolor{currentstroke}{rgb}{0.309804,0.372549,0.419608}%
\pgfsetstrokecolor{currentstroke}%
\pgfsetdash{}{0pt}%
\pgfpathmoveto{\pgfqpoint{5.746058in}{0.508460in}}%
\pgfpathlineto{\pgfqpoint{5.746058in}{0.470400in}}%
\pgfpathlineto{\pgfqpoint{5.735701in}{0.470400in}}%
\pgfpathlineto{\pgfqpoint{5.735701in}{0.508117in}}%
\pgfpathquadraticcurveto{\pgfqpoint{5.735701in}{0.517069in}}{\pgfqpoint{5.732207in}{0.521500in}}%
\pgfpathquadraticcurveto{\pgfqpoint{5.728712in}{0.525949in}}{\pgfqpoint{5.721742in}{0.525949in}}%
\pgfpathquadraticcurveto{\pgfqpoint{5.713348in}{0.525949in}}{\pgfqpoint{5.708503in}{0.520600in}}%
\pgfpathquadraticcurveto{\pgfqpoint{5.703657in}{0.515268in}}{\pgfqpoint{5.703657in}{0.506028in}}%
\pgfpathlineto{\pgfqpoint{5.703657in}{0.470400in}}%
\pgfpathlineto{\pgfqpoint{5.693246in}{0.470400in}}%
\pgfpathlineto{\pgfqpoint{5.693246in}{0.533443in}}%
\pgfpathlineto{\pgfqpoint{5.703657in}{0.533443in}}%
\pgfpathlineto{\pgfqpoint{5.703657in}{0.523644in}}%
\pgfpathquadraticcurveto{\pgfqpoint{5.707386in}{0.529336in}}{\pgfqpoint{5.712411in}{0.532146in}}%
\pgfpathquadraticcurveto{\pgfqpoint{5.717455in}{0.534956in}}{\pgfqpoint{5.724047in}{0.534956in}}%
\pgfpathquadraticcurveto{\pgfqpoint{5.734908in}{0.534956in}}{\pgfqpoint{5.740474in}{0.528237in}}%
\pgfpathquadraticcurveto{\pgfqpoint{5.746058in}{0.521518in}}{\pgfqpoint{5.746058in}{0.508460in}}%
\pgfpathlineto{\pgfqpoint{5.746058in}{0.508460in}}%
\pgfpathclose%
\pgfpathmoveto{\pgfqpoint{5.791124in}{0.526184in}}%
\pgfpathquadraticcurveto{\pgfqpoint{5.782803in}{0.526184in}}{\pgfqpoint{5.777957in}{0.519681in}}%
\pgfpathquadraticcurveto{\pgfqpoint{5.773112in}{0.513179in}}{\pgfqpoint{5.773112in}{0.501867in}}%
\pgfpathquadraticcurveto{\pgfqpoint{5.773112in}{0.490556in}}{\pgfqpoint{5.777921in}{0.484053in}}%
\pgfpathquadraticcurveto{\pgfqpoint{5.782749in}{0.477551in}}{\pgfqpoint{5.791124in}{0.477551in}}%
\pgfpathquadraticcurveto{\pgfqpoint{5.799410in}{0.477551in}}{\pgfqpoint{5.804237in}{0.484071in}}%
\pgfpathquadraticcurveto{\pgfqpoint{5.809082in}{0.490610in}}{\pgfqpoint{5.809082in}{0.501867in}}%
\pgfpathquadraticcurveto{\pgfqpoint{5.809082in}{0.513071in}}{\pgfqpoint{5.804237in}{0.519627in}}%
\pgfpathquadraticcurveto{\pgfqpoint{5.799410in}{0.526184in}}{\pgfqpoint{5.791124in}{0.526184in}}%
\pgfpathlineto{\pgfqpoint{5.791124in}{0.526184in}}%
\pgfpathclose%
\pgfpathmoveto{\pgfqpoint{5.791124in}{0.534956in}}%
\pgfpathquadraticcurveto{\pgfqpoint{5.804633in}{0.534956in}}{\pgfqpoint{5.812343in}{0.526166in}}%
\pgfpathquadraticcurveto{\pgfqpoint{5.820070in}{0.517394in}}{\pgfqpoint{5.820070in}{0.501867in}}%
\pgfpathquadraticcurveto{\pgfqpoint{5.820070in}{0.486395in}}{\pgfqpoint{5.812343in}{0.477569in}}%
\pgfpathquadraticcurveto{\pgfqpoint{5.804633in}{0.468761in}}{\pgfqpoint{5.791124in}{0.468761in}}%
\pgfpathquadraticcurveto{\pgfqpoint{5.777561in}{0.468761in}}{\pgfqpoint{5.769870in}{0.477569in}}%
\pgfpathquadraticcurveto{\pgfqpoint{5.762197in}{0.486395in}}{\pgfqpoint{5.762197in}{0.501867in}}%
\pgfpathquadraticcurveto{\pgfqpoint{5.762197in}{0.517394in}}{\pgfqpoint{5.769870in}{0.526166in}}%
\pgfpathquadraticcurveto{\pgfqpoint{5.777561in}{0.534956in}}{\pgfqpoint{5.791124in}{0.534956in}}%
\pgfpathlineto{\pgfqpoint{5.791124in}{0.534956in}}%
\pgfpathclose%
\pgfpathmoveto{\pgfqpoint{5.847484in}{0.551347in}}%
\pgfpathlineto{\pgfqpoint{5.847484in}{0.533443in}}%
\pgfpathlineto{\pgfqpoint{5.868811in}{0.533443in}}%
\pgfpathlineto{\pgfqpoint{5.868811in}{0.525391in}}%
\pgfpathlineto{\pgfqpoint{5.847484in}{0.525391in}}%
\pgfpathlineto{\pgfqpoint{5.847484in}{0.491168in}}%
\pgfpathquadraticcurveto{\pgfqpoint{5.847484in}{0.483459in}}{\pgfqpoint{5.849592in}{0.481261in}}%
\pgfpathquadraticcurveto{\pgfqpoint{5.851699in}{0.479064in}}{\pgfqpoint{5.858184in}{0.479064in}}%
\pgfpathlineto{\pgfqpoint{5.868811in}{0.479064in}}%
\pgfpathlineto{\pgfqpoint{5.868811in}{0.470400in}}%
\pgfpathlineto{\pgfqpoint{5.858184in}{0.470400in}}%
\pgfpathquadraticcurveto{\pgfqpoint{5.846188in}{0.470400in}}{\pgfqpoint{5.841630in}{0.474867in}}%
\pgfpathquadraticcurveto{\pgfqpoint{5.837073in}{0.479352in}}{\pgfqpoint{5.837073in}{0.491168in}}%
\pgfpathlineto{\pgfqpoint{5.837073in}{0.525391in}}%
\pgfpathlineto{\pgfqpoint{5.829472in}{0.525391in}}%
\pgfpathlineto{\pgfqpoint{5.829472in}{0.533443in}}%
\pgfpathlineto{\pgfqpoint{5.837073in}{0.533443in}}%
\pgfpathlineto{\pgfqpoint{5.837073in}{0.551347in}}%
\pgfpathlineto{\pgfqpoint{5.847484in}{0.551347in}}%
\pgfpathlineto{\pgfqpoint{5.847484in}{0.551347in}}%
\pgfpathclose%
\pgfusepath{fill}%
\end{pgfscope}%
\begin{pgfscope}%
\pgfsetbuttcap%
\pgfsetroundjoin%
\definecolor{currentfill}{rgb}{0.309804,0.372549,0.419608}%
\pgfsetfillcolor{currentfill}%
\pgfsetlinewidth{0.000000pt}%
\definecolor{currentstroke}{rgb}{0.309804,0.372549,0.419608}%
\pgfsetstrokecolor{currentstroke}%
\pgfsetdash{}{0pt}%
\pgfpathmoveto{\pgfqpoint{5.951204in}{0.551347in}}%
\pgfpathlineto{\pgfqpoint{5.951204in}{0.533443in}}%
\pgfpathlineto{\pgfqpoint{5.972530in}{0.533443in}}%
\pgfpathlineto{\pgfqpoint{5.972530in}{0.525391in}}%
\pgfpathlineto{\pgfqpoint{5.951204in}{0.525391in}}%
\pgfpathlineto{\pgfqpoint{5.951204in}{0.491168in}}%
\pgfpathquadraticcurveto{\pgfqpoint{5.951204in}{0.483459in}}{\pgfqpoint{5.953311in}{0.481261in}}%
\pgfpathquadraticcurveto{\pgfqpoint{5.955419in}{0.479064in}}{\pgfqpoint{5.961903in}{0.479064in}}%
\pgfpathlineto{\pgfqpoint{5.972530in}{0.479064in}}%
\pgfpathlineto{\pgfqpoint{5.972530in}{0.470400in}}%
\pgfpathlineto{\pgfqpoint{5.961903in}{0.470400in}}%
\pgfpathquadraticcurveto{\pgfqpoint{5.949907in}{0.470400in}}{\pgfqpoint{5.945350in}{0.474867in}}%
\pgfpathquadraticcurveto{\pgfqpoint{5.940793in}{0.479352in}}{\pgfqpoint{5.940793in}{0.491168in}}%
\pgfpathlineto{\pgfqpoint{5.940793in}{0.525391in}}%
\pgfpathlineto{\pgfqpoint{5.933192in}{0.525391in}}%
\pgfpathlineto{\pgfqpoint{5.933192in}{0.533443in}}%
\pgfpathlineto{\pgfqpoint{5.940793in}{0.533443in}}%
\pgfpathlineto{\pgfqpoint{5.940793in}{0.551347in}}%
\pgfpathlineto{\pgfqpoint{5.951204in}{0.551347in}}%
\pgfpathlineto{\pgfqpoint{5.951204in}{0.551347in}}%
\pgfpathclose%
\pgfpathmoveto{\pgfqpoint{6.038563in}{0.508460in}}%
\pgfpathlineto{\pgfqpoint{6.038563in}{0.470400in}}%
\pgfpathlineto{\pgfqpoint{6.028206in}{0.470400in}}%
\pgfpathlineto{\pgfqpoint{6.028206in}{0.508117in}}%
\pgfpathquadraticcurveto{\pgfqpoint{6.028206in}{0.517069in}}{\pgfqpoint{6.024712in}{0.521500in}}%
\pgfpathquadraticcurveto{\pgfqpoint{6.021217in}{0.525949in}}{\pgfqpoint{6.014247in}{0.525949in}}%
\pgfpathquadraticcurveto{\pgfqpoint{6.005853in}{0.525949in}}{\pgfqpoint{6.001008in}{0.520600in}}%
\pgfpathquadraticcurveto{\pgfqpoint{5.996162in}{0.515268in}}{\pgfqpoint{5.996162in}{0.506028in}}%
\pgfpathlineto{\pgfqpoint{5.996162in}{0.470400in}}%
\pgfpathlineto{\pgfqpoint{5.985751in}{0.470400in}}%
\pgfpathlineto{\pgfqpoint{5.985751in}{0.557993in}}%
\pgfpathlineto{\pgfqpoint{5.996162in}{0.557993in}}%
\pgfpathlineto{\pgfqpoint{5.996162in}{0.523644in}}%
\pgfpathquadraticcurveto{\pgfqpoint{5.999891in}{0.529336in}}{\pgfqpoint{6.004916in}{0.532146in}}%
\pgfpathquadraticcurveto{\pgfqpoint{6.009960in}{0.534956in}}{\pgfqpoint{6.016552in}{0.534956in}}%
\pgfpathquadraticcurveto{\pgfqpoint{6.027413in}{0.534956in}}{\pgfqpoint{6.032979in}{0.528237in}}%
\pgfpathquadraticcurveto{\pgfqpoint{6.038563in}{0.521518in}}{\pgfqpoint{6.038563in}{0.508460in}}%
\pgfpathlineto{\pgfqpoint{6.038563in}{0.508460in}}%
\pgfpathclose%
\pgfpathmoveto{\pgfqpoint{6.113133in}{0.504515in}}%
\pgfpathlineto{\pgfqpoint{6.113133in}{0.499454in}}%
\pgfpathlineto{\pgfqpoint{6.065509in}{0.499454in}}%
\pgfpathquadraticcurveto{\pgfqpoint{6.066194in}{0.488754in}}{\pgfqpoint{6.071957in}{0.483153in}}%
\pgfpathquadraticcurveto{\pgfqpoint{6.077721in}{0.477551in}}{\pgfqpoint{6.088024in}{0.477551in}}%
\pgfpathquadraticcurveto{\pgfqpoint{6.093986in}{0.477551in}}{\pgfqpoint{6.099588in}{0.479010in}}%
\pgfpathquadraticcurveto{\pgfqpoint{6.105190in}{0.480469in}}{\pgfqpoint{6.110720in}{0.483405in}}%
\pgfpathlineto{\pgfqpoint{6.110720in}{0.473606in}}%
\pgfpathquadraticcurveto{\pgfqpoint{6.105136in}{0.471247in}}{\pgfqpoint{6.099282in}{0.470004in}}%
\pgfpathquadraticcurveto{\pgfqpoint{6.093428in}{0.468761in}}{\pgfqpoint{6.087412in}{0.468761in}}%
\pgfpathquadraticcurveto{\pgfqpoint{6.072318in}{0.468761in}}{\pgfqpoint{6.063510in}{0.477533in}}%
\pgfpathquadraticcurveto{\pgfqpoint{6.054702in}{0.486323in}}{\pgfqpoint{6.054702in}{0.501309in}}%
\pgfpathquadraticcurveto{\pgfqpoint{6.054702in}{0.516781in}}{\pgfqpoint{6.063059in}{0.525859in}}%
\pgfpathquadraticcurveto{\pgfqpoint{6.071417in}{0.534956in}}{\pgfqpoint{6.085611in}{0.534956in}}%
\pgfpathquadraticcurveto{\pgfqpoint{6.098327in}{0.534956in}}{\pgfqpoint{6.105730in}{0.526760in}}%
\pgfpathquadraticcurveto{\pgfqpoint{6.113133in}{0.518583in}}{\pgfqpoint{6.113133in}{0.504515in}}%
\pgfpathlineto{\pgfqpoint{6.113133in}{0.504515in}}%
\pgfpathclose%
\pgfpathmoveto{\pgfqpoint{6.102776in}{0.507559in}}%
\pgfpathquadraticcurveto{\pgfqpoint{6.102668in}{0.516043in}}{\pgfqpoint{6.098021in}{0.521104in}}%
\pgfpathquadraticcurveto{\pgfqpoint{6.093374in}{0.526184in}}{\pgfqpoint{6.085719in}{0.526184in}}%
\pgfpathquadraticcurveto{\pgfqpoint{6.077055in}{0.526184in}}{\pgfqpoint{6.071849in}{0.521284in}}%
\pgfpathquadraticcurveto{\pgfqpoint{6.066644in}{0.516385in}}{\pgfqpoint{6.065851in}{0.507487in}}%
\pgfpathlineto{\pgfqpoint{6.102776in}{0.507559in}}%
\pgfpathlineto{\pgfqpoint{6.102776in}{0.507559in}}%
\pgfpathclose%
\pgfusepath{fill}%
\end{pgfscope}%
\begin{pgfscope}%
\pgfsetbuttcap%
\pgfsetroundjoin%
\definecolor{currentfill}{rgb}{0.309804,0.372549,0.419608}%
\pgfsetfillcolor{currentfill}%
\pgfsetlinewidth{0.000000pt}%
\definecolor{currentstroke}{rgb}{0.309804,0.372549,0.419608}%
\pgfsetstrokecolor{currentstroke}%
\pgfsetdash{}{0pt}%
\pgfpathmoveto{\pgfqpoint{6.198282in}{0.479856in}}%
\pgfpathlineto{\pgfqpoint{6.198282in}{0.446426in}}%
\pgfpathlineto{\pgfqpoint{6.187871in}{0.446426in}}%
\pgfpathlineto{\pgfqpoint{6.187871in}{0.533443in}}%
\pgfpathlineto{\pgfqpoint{6.198282in}{0.533443in}}%
\pgfpathlineto{\pgfqpoint{6.198282in}{0.523878in}}%
\pgfpathquadraticcurveto{\pgfqpoint{6.201561in}{0.529498in}}{\pgfqpoint{6.206532in}{0.532218in}}%
\pgfpathquadraticcurveto{\pgfqpoint{6.211521in}{0.534956in}}{\pgfqpoint{6.218438in}{0.534956in}}%
\pgfpathquadraticcurveto{\pgfqpoint{6.229930in}{0.534956in}}{\pgfqpoint{6.237099in}{0.525841in}}%
\pgfpathquadraticcurveto{\pgfqpoint{6.244285in}{0.516727in}}{\pgfqpoint{6.244285in}{0.501867in}}%
\pgfpathquadraticcurveto{\pgfqpoint{6.244285in}{0.487007in}}{\pgfqpoint{6.237099in}{0.477875in}}%
\pgfpathquadraticcurveto{\pgfqpoint{6.229930in}{0.468761in}}{\pgfqpoint{6.218438in}{0.468761in}}%
\pgfpathquadraticcurveto{\pgfqpoint{6.211521in}{0.468761in}}{\pgfqpoint{6.206532in}{0.471499in}}%
\pgfpathquadraticcurveto{\pgfqpoint{6.201561in}{0.474237in}}{\pgfqpoint{6.198282in}{0.479856in}}%
\pgfpathlineto{\pgfqpoint{6.198282in}{0.479856in}}%
\pgfpathclose%
\pgfpathmoveto{\pgfqpoint{6.233532in}{0.501867in}}%
\pgfpathquadraticcurveto{\pgfqpoint{6.233532in}{0.513287in}}{\pgfqpoint{6.228831in}{0.519789in}}%
\pgfpathquadraticcurveto{\pgfqpoint{6.224130in}{0.526292in}}{\pgfqpoint{6.215916in}{0.526292in}}%
\pgfpathquadraticcurveto{\pgfqpoint{6.207685in}{0.526292in}}{\pgfqpoint{6.202984in}{0.519789in}}%
\pgfpathquadraticcurveto{\pgfqpoint{6.198282in}{0.513287in}}{\pgfqpoint{6.198282in}{0.501867in}}%
\pgfpathquadraticcurveto{\pgfqpoint{6.198282in}{0.490448in}}{\pgfqpoint{6.202984in}{0.483945in}}%
\pgfpathquadraticcurveto{\pgfqpoint{6.207685in}{0.477443in}}{\pgfqpoint{6.215916in}{0.477443in}}%
\pgfpathquadraticcurveto{\pgfqpoint{6.224130in}{0.477443in}}{\pgfqpoint{6.228831in}{0.483945in}}%
\pgfpathquadraticcurveto{\pgfqpoint{6.233532in}{0.490448in}}{\pgfqpoint{6.233532in}{0.501867in}}%
\pgfpathlineto{\pgfqpoint{6.233532in}{0.501867in}}%
\pgfpathclose%
\pgfpathmoveto{\pgfqpoint{6.290108in}{0.502083in}}%
\pgfpathquadraticcurveto{\pgfqpoint{6.277554in}{0.502083in}}{\pgfqpoint{6.272709in}{0.499219in}}%
\pgfpathquadraticcurveto{\pgfqpoint{6.267863in}{0.496356in}}{\pgfqpoint{6.267863in}{0.489421in}}%
\pgfpathquadraticcurveto{\pgfqpoint{6.267863in}{0.483909in}}{\pgfqpoint{6.271502in}{0.480667in}}%
\pgfpathquadraticcurveto{\pgfqpoint{6.275140in}{0.477443in}}{\pgfqpoint{6.281372in}{0.477443in}}%
\pgfpathquadraticcurveto{\pgfqpoint{6.290000in}{0.477443in}}{\pgfqpoint{6.295206in}{0.483549in}}%
\pgfpathquadraticcurveto{\pgfqpoint{6.300411in}{0.489655in}}{\pgfqpoint{6.300411in}{0.499778in}}%
\pgfpathlineto{\pgfqpoint{6.300411in}{0.502083in}}%
\pgfpathlineto{\pgfqpoint{6.290108in}{0.502083in}}%
\pgfpathlineto{\pgfqpoint{6.290108in}{0.502083in}}%
\pgfpathclose%
\pgfpathmoveto{\pgfqpoint{6.310768in}{0.506370in}}%
\pgfpathlineto{\pgfqpoint{6.310768in}{0.470400in}}%
\pgfpathlineto{\pgfqpoint{6.300411in}{0.470400in}}%
\pgfpathlineto{\pgfqpoint{6.300411in}{0.479964in}}%
\pgfpathquadraticcurveto{\pgfqpoint{6.296863in}{0.474237in}}{\pgfqpoint{6.291567in}{0.471499in}}%
\pgfpathquadraticcurveto{\pgfqpoint{6.286272in}{0.468761in}}{\pgfqpoint{6.278617in}{0.468761in}}%
\pgfpathquadraticcurveto{\pgfqpoint{6.268944in}{0.468761in}}{\pgfqpoint{6.263216in}{0.474201in}}%
\pgfpathquadraticcurveto{\pgfqpoint{6.257506in}{0.479640in}}{\pgfqpoint{6.257506in}{0.488754in}}%
\pgfpathquadraticcurveto{\pgfqpoint{6.257506in}{0.499382in}}{\pgfqpoint{6.264621in}{0.504785in}}%
\pgfpathquadraticcurveto{\pgfqpoint{6.271754in}{0.510189in}}{\pgfqpoint{6.285875in}{0.510189in}}%
\pgfpathlineto{\pgfqpoint{6.300411in}{0.510189in}}%
\pgfpathlineto{\pgfqpoint{6.300411in}{0.511216in}}%
\pgfpathquadraticcurveto{\pgfqpoint{6.300411in}{0.518366in}}{\pgfqpoint{6.295710in}{0.522275in}}%
\pgfpathquadraticcurveto{\pgfqpoint{6.291009in}{0.526184in}}{\pgfqpoint{6.282507in}{0.526184in}}%
\pgfpathquadraticcurveto{\pgfqpoint{6.277104in}{0.526184in}}{\pgfqpoint{6.271970in}{0.524887in}}%
\pgfpathquadraticcurveto{\pgfqpoint{6.266855in}{0.523590in}}{\pgfqpoint{6.262135in}{0.520996in}}%
\pgfpathlineto{\pgfqpoint{6.262135in}{0.530579in}}%
\pgfpathquadraticcurveto{\pgfqpoint{6.267809in}{0.532776in}}{\pgfqpoint{6.273159in}{0.533857in}}%
\pgfpathquadraticcurveto{\pgfqpoint{6.278508in}{0.534956in}}{\pgfqpoint{6.283570in}{0.534956in}}%
\pgfpathquadraticcurveto{\pgfqpoint{6.297259in}{0.534956in}}{\pgfqpoint{6.304014in}{0.527859in}}%
\pgfpathquadraticcurveto{\pgfqpoint{6.310768in}{0.520780in}}{\pgfqpoint{6.310768in}{0.506370in}}%
\pgfpathlineto{\pgfqpoint{6.310768in}{0.506370in}}%
\pgfpathclose%
\pgfpathmoveto{\pgfqpoint{6.332095in}{0.533443in}}%
\pgfpathlineto{\pgfqpoint{6.342452in}{0.533443in}}%
\pgfpathlineto{\pgfqpoint{6.342452in}{0.470400in}}%
\pgfpathlineto{\pgfqpoint{6.332095in}{0.470400in}}%
\pgfpathlineto{\pgfqpoint{6.332095in}{0.533443in}}%
\pgfpathlineto{\pgfqpoint{6.332095in}{0.533443in}}%
\pgfpathclose%
\pgfpathmoveto{\pgfqpoint{6.332095in}{0.557993in}}%
\pgfpathlineto{\pgfqpoint{6.342452in}{0.557993in}}%
\pgfpathlineto{\pgfqpoint{6.342452in}{0.544862in}}%
\pgfpathlineto{\pgfqpoint{6.332095in}{0.544862in}}%
\pgfpathlineto{\pgfqpoint{6.332095in}{0.557993in}}%
\pgfpathlineto{\pgfqpoint{6.332095in}{0.557993in}}%
\pgfpathclose%
\pgfpathmoveto{\pgfqpoint{6.400649in}{0.523770in}}%
\pgfpathquadraticcurveto{\pgfqpoint{6.398902in}{0.524779in}}{\pgfqpoint{6.396848in}{0.525247in}}%
\pgfpathquadraticcurveto{\pgfqpoint{6.394795in}{0.525733in}}{\pgfqpoint{6.392327in}{0.525733in}}%
\pgfpathquadraticcurveto{\pgfqpoint{6.383537in}{0.525733in}}{\pgfqpoint{6.378836in}{0.520023in}}%
\pgfpathquadraticcurveto{\pgfqpoint{6.374135in}{0.514314in}}{\pgfqpoint{6.374135in}{0.503614in}}%
\pgfpathlineto{\pgfqpoint{6.374135in}{0.470400in}}%
\pgfpathlineto{\pgfqpoint{6.363724in}{0.470400in}}%
\pgfpathlineto{\pgfqpoint{6.363724in}{0.533443in}}%
\pgfpathlineto{\pgfqpoint{6.374135in}{0.533443in}}%
\pgfpathlineto{\pgfqpoint{6.374135in}{0.523644in}}%
\pgfpathquadraticcurveto{\pgfqpoint{6.377413in}{0.529390in}}{\pgfqpoint{6.382637in}{0.532164in}}%
\pgfpathquadraticcurveto{\pgfqpoint{6.387878in}{0.534956in}}{\pgfqpoint{6.395371in}{0.534956in}}%
\pgfpathquadraticcurveto{\pgfqpoint{6.396434in}{0.534956in}}{\pgfqpoint{6.397731in}{0.534811in}}%
\pgfpathquadraticcurveto{\pgfqpoint{6.399028in}{0.534685in}}{\pgfqpoint{6.400595in}{0.534397in}}%
\pgfpathlineto{\pgfqpoint{6.400649in}{0.523770in}}%
\pgfpathlineto{\pgfqpoint{6.400649in}{0.523770in}}%
\pgfpathclose%
\pgfpathmoveto{\pgfqpoint{6.462899in}{0.504515in}}%
\pgfpathlineto{\pgfqpoint{6.462899in}{0.499454in}}%
\pgfpathlineto{\pgfqpoint{6.415275in}{0.499454in}}%
\pgfpathquadraticcurveto{\pgfqpoint{6.415959in}{0.488754in}}{\pgfqpoint{6.421723in}{0.483153in}}%
\pgfpathquadraticcurveto{\pgfqpoint{6.427487in}{0.477551in}}{\pgfqpoint{6.437790in}{0.477551in}}%
\pgfpathquadraticcurveto{\pgfqpoint{6.443752in}{0.477551in}}{\pgfqpoint{6.449354in}{0.479010in}}%
\pgfpathquadraticcurveto{\pgfqpoint{6.454956in}{0.480469in}}{\pgfqpoint{6.460485in}{0.483405in}}%
\pgfpathlineto{\pgfqpoint{6.460485in}{0.473606in}}%
\pgfpathquadraticcurveto{\pgfqpoint{6.454901in}{0.471247in}}{\pgfqpoint{6.449048in}{0.470004in}}%
\pgfpathquadraticcurveto{\pgfqpoint{6.443194in}{0.468761in}}{\pgfqpoint{6.437178in}{0.468761in}}%
\pgfpathquadraticcurveto{\pgfqpoint{6.422083in}{0.468761in}}{\pgfqpoint{6.413275in}{0.477533in}}%
\pgfpathquadraticcurveto{\pgfqpoint{6.404467in}{0.486323in}}{\pgfqpoint{6.404467in}{0.501309in}}%
\pgfpathquadraticcurveto{\pgfqpoint{6.404467in}{0.516781in}}{\pgfqpoint{6.412825in}{0.525859in}}%
\pgfpathquadraticcurveto{\pgfqpoint{6.421183in}{0.534956in}}{\pgfqpoint{6.435376in}{0.534956in}}%
\pgfpathquadraticcurveto{\pgfqpoint{6.448093in}{0.534956in}}{\pgfqpoint{6.455496in}{0.526760in}}%
\pgfpathquadraticcurveto{\pgfqpoint{6.462899in}{0.518583in}}{\pgfqpoint{6.462899in}{0.504515in}}%
\pgfpathlineto{\pgfqpoint{6.462899in}{0.504515in}}%
\pgfpathclose%
\pgfpathmoveto{\pgfqpoint{6.452542in}{0.507559in}}%
\pgfpathquadraticcurveto{\pgfqpoint{6.452434in}{0.516043in}}{\pgfqpoint{6.447787in}{0.521104in}}%
\pgfpathquadraticcurveto{\pgfqpoint{6.443140in}{0.526184in}}{\pgfqpoint{6.435484in}{0.526184in}}%
\pgfpathquadraticcurveto{\pgfqpoint{6.426821in}{0.526184in}}{\pgfqpoint{6.421615in}{0.521284in}}%
\pgfpathquadraticcurveto{\pgfqpoint{6.416410in}{0.516385in}}{\pgfqpoint{6.415617in}{0.507487in}}%
\pgfpathlineto{\pgfqpoint{6.452542in}{0.507559in}}%
\pgfpathlineto{\pgfqpoint{6.452542in}{0.507559in}}%
\pgfpathclose%
\pgfpathmoveto{\pgfqpoint{6.521384in}{0.523878in}}%
\pgfpathlineto{\pgfqpoint{6.521384in}{0.557993in}}%
\pgfpathlineto{\pgfqpoint{6.531741in}{0.557993in}}%
\pgfpathlineto{\pgfqpoint{6.531741in}{0.470400in}}%
\pgfpathlineto{\pgfqpoint{6.521384in}{0.470400in}}%
\pgfpathlineto{\pgfqpoint{6.521384in}{0.479856in}}%
\pgfpathquadraticcurveto{\pgfqpoint{6.518124in}{0.474237in}}{\pgfqpoint{6.513135in}{0.471499in}}%
\pgfpathquadraticcurveto{\pgfqpoint{6.508163in}{0.468761in}}{\pgfqpoint{6.501175in}{0.468761in}}%
\pgfpathquadraticcurveto{\pgfqpoint{6.489755in}{0.468761in}}{\pgfqpoint{6.482568in}{0.477875in}}%
\pgfpathquadraticcurveto{\pgfqpoint{6.475399in}{0.487007in}}{\pgfqpoint{6.475399in}{0.501867in}}%
\pgfpathquadraticcurveto{\pgfqpoint{6.475399in}{0.516727in}}{\pgfqpoint{6.482568in}{0.525841in}}%
\pgfpathquadraticcurveto{\pgfqpoint{6.489755in}{0.534956in}}{\pgfqpoint{6.501175in}{0.534956in}}%
\pgfpathquadraticcurveto{\pgfqpoint{6.508163in}{0.534956in}}{\pgfqpoint{6.513135in}{0.532218in}}%
\pgfpathquadraticcurveto{\pgfqpoint{6.518124in}{0.529498in}}{\pgfqpoint{6.521384in}{0.523878in}}%
\pgfpathlineto{\pgfqpoint{6.521384in}{0.523878in}}%
\pgfpathclose%
\pgfpathmoveto{\pgfqpoint{6.486099in}{0.501867in}}%
\pgfpathquadraticcurveto{\pgfqpoint{6.486099in}{0.490448in}}{\pgfqpoint{6.490800in}{0.483945in}}%
\pgfpathquadraticcurveto{\pgfqpoint{6.495501in}{0.477443in}}{\pgfqpoint{6.503714in}{0.477443in}}%
\pgfpathquadraticcurveto{\pgfqpoint{6.511928in}{0.477443in}}{\pgfqpoint{6.516647in}{0.483945in}}%
\pgfpathquadraticcurveto{\pgfqpoint{6.521384in}{0.490448in}}{\pgfqpoint{6.521384in}{0.501867in}}%
\pgfpathquadraticcurveto{\pgfqpoint{6.521384in}{0.513287in}}{\pgfqpoint{6.516647in}{0.519789in}}%
\pgfpathquadraticcurveto{\pgfqpoint{6.511928in}{0.526292in}}{\pgfqpoint{6.503714in}{0.526292in}}%
\pgfpathquadraticcurveto{\pgfqpoint{6.495501in}{0.526292in}}{\pgfqpoint{6.490800in}{0.519789in}}%
\pgfpathquadraticcurveto{\pgfqpoint{6.486099in}{0.513287in}}{\pgfqpoint{6.486099in}{0.501867in}}%
\pgfpathlineto{\pgfqpoint{6.486099in}{0.501867in}}%
\pgfpathclose%
\pgfpathmoveto{\pgfqpoint{6.547862in}{0.506586in}}%
\pgfpathlineto{\pgfqpoint{6.578195in}{0.506586in}}%
\pgfpathlineto{\pgfqpoint{6.578195in}{0.497364in}}%
\pgfpathlineto{\pgfqpoint{6.547862in}{0.497364in}}%
\pgfpathlineto{\pgfqpoint{6.547862in}{0.506586in}}%
\pgfpathlineto{\pgfqpoint{6.547862in}{0.506586in}}%
\pgfpathclose%
\pgfpathmoveto{\pgfqpoint{6.636158in}{0.523878in}}%
\pgfpathlineto{\pgfqpoint{6.636158in}{0.557993in}}%
\pgfpathlineto{\pgfqpoint{6.646515in}{0.557993in}}%
\pgfpathlineto{\pgfqpoint{6.646515in}{0.470400in}}%
\pgfpathlineto{\pgfqpoint{6.636158in}{0.470400in}}%
\pgfpathlineto{\pgfqpoint{6.636158in}{0.479856in}}%
\pgfpathquadraticcurveto{\pgfqpoint{6.632898in}{0.474237in}}{\pgfqpoint{6.627908in}{0.471499in}}%
\pgfpathquadraticcurveto{\pgfqpoint{6.622937in}{0.468761in}}{\pgfqpoint{6.615948in}{0.468761in}}%
\pgfpathquadraticcurveto{\pgfqpoint{6.604528in}{0.468761in}}{\pgfqpoint{6.597342in}{0.477875in}}%
\pgfpathquadraticcurveto{\pgfqpoint{6.590173in}{0.487007in}}{\pgfqpoint{6.590173in}{0.501867in}}%
\pgfpathquadraticcurveto{\pgfqpoint{6.590173in}{0.516727in}}{\pgfqpoint{6.597342in}{0.525841in}}%
\pgfpathquadraticcurveto{\pgfqpoint{6.604528in}{0.534956in}}{\pgfqpoint{6.615948in}{0.534956in}}%
\pgfpathquadraticcurveto{\pgfqpoint{6.622937in}{0.534956in}}{\pgfqpoint{6.627908in}{0.532218in}}%
\pgfpathquadraticcurveto{\pgfqpoint{6.632898in}{0.529498in}}{\pgfqpoint{6.636158in}{0.523878in}}%
\pgfpathlineto{\pgfqpoint{6.636158in}{0.523878in}}%
\pgfpathclose%
\pgfpathmoveto{\pgfqpoint{6.600872in}{0.501867in}}%
\pgfpathquadraticcurveto{\pgfqpoint{6.600872in}{0.490448in}}{\pgfqpoint{6.605573in}{0.483945in}}%
\pgfpathquadraticcurveto{\pgfqpoint{6.610274in}{0.477443in}}{\pgfqpoint{6.618488in}{0.477443in}}%
\pgfpathquadraticcurveto{\pgfqpoint{6.626701in}{0.477443in}}{\pgfqpoint{6.631421in}{0.483945in}}%
\pgfpathquadraticcurveto{\pgfqpoint{6.636158in}{0.490448in}}{\pgfqpoint{6.636158in}{0.501867in}}%
\pgfpathquadraticcurveto{\pgfqpoint{6.636158in}{0.513287in}}{\pgfqpoint{6.631421in}{0.519789in}}%
\pgfpathquadraticcurveto{\pgfqpoint{6.626701in}{0.526292in}}{\pgfqpoint{6.618488in}{0.526292in}}%
\pgfpathquadraticcurveto{\pgfqpoint{6.610274in}{0.526292in}}{\pgfqpoint{6.605573in}{0.519789in}}%
\pgfpathquadraticcurveto{\pgfqpoint{6.600872in}{0.513287in}}{\pgfqpoint{6.600872in}{0.501867in}}%
\pgfpathlineto{\pgfqpoint{6.600872in}{0.501867in}}%
\pgfpathclose%
\pgfpathmoveto{\pgfqpoint{6.667859in}{0.533443in}}%
\pgfpathlineto{\pgfqpoint{6.678216in}{0.533443in}}%
\pgfpathlineto{\pgfqpoint{6.678216in}{0.470400in}}%
\pgfpathlineto{\pgfqpoint{6.667859in}{0.470400in}}%
\pgfpathlineto{\pgfqpoint{6.667859in}{0.533443in}}%
\pgfpathlineto{\pgfqpoint{6.667859in}{0.533443in}}%
\pgfpathclose%
\pgfpathmoveto{\pgfqpoint{6.667859in}{0.557993in}}%
\pgfpathlineto{\pgfqpoint{6.678216in}{0.557993in}}%
\pgfpathlineto{\pgfqpoint{6.678216in}{0.544862in}}%
\pgfpathlineto{\pgfqpoint{6.667859in}{0.544862in}}%
\pgfpathlineto{\pgfqpoint{6.667859in}{0.557993in}}%
\pgfpathlineto{\pgfqpoint{6.667859in}{0.557993in}}%
\pgfpathclose%
\pgfpathmoveto{\pgfqpoint{6.770637in}{0.557993in}}%
\pgfpathlineto{\pgfqpoint{6.770637in}{0.549365in}}%
\pgfpathlineto{\pgfqpoint{6.760730in}{0.549365in}}%
\pgfpathquadraticcurveto{\pgfqpoint{6.755164in}{0.549365in}}{\pgfqpoint{6.752985in}{0.547114in}}%
\pgfpathquadraticcurveto{\pgfqpoint{6.750823in}{0.544862in}}{\pgfqpoint{6.750823in}{0.539008in}}%
\pgfpathlineto{\pgfqpoint{6.750823in}{0.533443in}}%
\pgfpathlineto{\pgfqpoint{6.767881in}{0.533443in}}%
\pgfpathlineto{\pgfqpoint{6.767881in}{0.525391in}}%
\pgfpathlineto{\pgfqpoint{6.750823in}{0.525391in}}%
\pgfpathlineto{\pgfqpoint{6.750823in}{0.470400in}}%
\pgfpathlineto{\pgfqpoint{6.740412in}{0.470400in}}%
\pgfpathlineto{\pgfqpoint{6.740412in}{0.525391in}}%
\pgfpathlineto{\pgfqpoint{6.711989in}{0.525391in}}%
\pgfpathlineto{\pgfqpoint{6.711989in}{0.470400in}}%
\pgfpathlineto{\pgfqpoint{6.701578in}{0.470400in}}%
\pgfpathlineto{\pgfqpoint{6.701578in}{0.525391in}}%
\pgfpathlineto{\pgfqpoint{6.691671in}{0.525391in}}%
\pgfpathlineto{\pgfqpoint{6.691671in}{0.533443in}}%
\pgfpathlineto{\pgfqpoint{6.701578in}{0.533443in}}%
\pgfpathlineto{\pgfqpoint{6.701578in}{0.537837in}}%
\pgfpathquadraticcurveto{\pgfqpoint{6.701578in}{0.548357in}}{\pgfqpoint{6.706477in}{0.553166in}}%
\pgfpathquadraticcurveto{\pgfqpoint{6.711377in}{0.557993in}}{\pgfqpoint{6.722004in}{0.557993in}}%
\pgfpathlineto{\pgfqpoint{6.731802in}{0.557993in}}%
\pgfpathlineto{\pgfqpoint{6.731802in}{0.549365in}}%
\pgfpathlineto{\pgfqpoint{6.721896in}{0.549365in}}%
\pgfpathquadraticcurveto{\pgfqpoint{6.716330in}{0.549365in}}{\pgfqpoint{6.714132in}{0.547114in}}%
\pgfpathquadraticcurveto{\pgfqpoint{6.711989in}{0.544862in}}{\pgfqpoint{6.711989in}{0.539008in}}%
\pgfpathlineto{\pgfqpoint{6.711989in}{0.533443in}}%
\pgfpathlineto{\pgfqpoint{6.740412in}{0.533443in}}%
\pgfpathlineto{\pgfqpoint{6.740412in}{0.537837in}}%
\pgfpathquadraticcurveto{\pgfqpoint{6.740412in}{0.548357in}}{\pgfqpoint{6.745311in}{0.553166in}}%
\pgfpathquadraticcurveto{\pgfqpoint{6.750211in}{0.557993in}}{\pgfqpoint{6.760856in}{0.557993in}}%
\pgfpathlineto{\pgfqpoint{6.770637in}{0.557993in}}%
\pgfpathlineto{\pgfqpoint{6.770637in}{0.557993in}}%
\pgfpathclose%
\pgfpathmoveto{\pgfqpoint{6.833229in}{0.504515in}}%
\pgfpathlineto{\pgfqpoint{6.833229in}{0.499454in}}%
\pgfpathlineto{\pgfqpoint{6.785605in}{0.499454in}}%
\pgfpathquadraticcurveto{\pgfqpoint{6.786289in}{0.488754in}}{\pgfqpoint{6.792053in}{0.483153in}}%
\pgfpathquadraticcurveto{\pgfqpoint{6.797817in}{0.477551in}}{\pgfqpoint{6.808120in}{0.477551in}}%
\pgfpathquadraticcurveto{\pgfqpoint{6.814082in}{0.477551in}}{\pgfqpoint{6.819684in}{0.479010in}}%
\pgfpathquadraticcurveto{\pgfqpoint{6.825285in}{0.480469in}}{\pgfqpoint{6.830815in}{0.483405in}}%
\pgfpathlineto{\pgfqpoint{6.830815in}{0.473606in}}%
\pgfpathquadraticcurveto{\pgfqpoint{6.825231in}{0.471247in}}{\pgfqpoint{6.819377in}{0.470004in}}%
\pgfpathquadraticcurveto{\pgfqpoint{6.813523in}{0.468761in}}{\pgfqpoint{6.807507in}{0.468761in}}%
\pgfpathquadraticcurveto{\pgfqpoint{6.792413in}{0.468761in}}{\pgfqpoint{6.783605in}{0.477533in}}%
\pgfpathquadraticcurveto{\pgfqpoint{6.774797in}{0.486323in}}{\pgfqpoint{6.774797in}{0.501309in}}%
\pgfpathquadraticcurveto{\pgfqpoint{6.774797in}{0.516781in}}{\pgfqpoint{6.783155in}{0.525859in}}%
\pgfpathquadraticcurveto{\pgfqpoint{6.791513in}{0.534956in}}{\pgfqpoint{6.805706in}{0.534956in}}%
\pgfpathquadraticcurveto{\pgfqpoint{6.818423in}{0.534956in}}{\pgfqpoint{6.825826in}{0.526760in}}%
\pgfpathquadraticcurveto{\pgfqpoint{6.833229in}{0.518583in}}{\pgfqpoint{6.833229in}{0.504515in}}%
\pgfpathlineto{\pgfqpoint{6.833229in}{0.504515in}}%
\pgfpathclose%
\pgfpathmoveto{\pgfqpoint{6.822872in}{0.507559in}}%
\pgfpathquadraticcurveto{\pgfqpoint{6.822764in}{0.516043in}}{\pgfqpoint{6.818117in}{0.521104in}}%
\pgfpathquadraticcurveto{\pgfqpoint{6.813469in}{0.526184in}}{\pgfqpoint{6.805814in}{0.526184in}}%
\pgfpathquadraticcurveto{\pgfqpoint{6.797150in}{0.526184in}}{\pgfqpoint{6.791945in}{0.521284in}}%
\pgfpathquadraticcurveto{\pgfqpoint{6.786739in}{0.516385in}}{\pgfqpoint{6.785947in}{0.507487in}}%
\pgfpathlineto{\pgfqpoint{6.822872in}{0.507559in}}%
\pgfpathlineto{\pgfqpoint{6.822872in}{0.507559in}}%
\pgfpathclose%
\pgfpathmoveto{\pgfqpoint{6.886761in}{0.523770in}}%
\pgfpathquadraticcurveto{\pgfqpoint{6.885014in}{0.524779in}}{\pgfqpoint{6.882960in}{0.525247in}}%
\pgfpathquadraticcurveto{\pgfqpoint{6.880907in}{0.525733in}}{\pgfqpoint{6.878439in}{0.525733in}}%
\pgfpathquadraticcurveto{\pgfqpoint{6.869649in}{0.525733in}}{\pgfqpoint{6.864948in}{0.520023in}}%
\pgfpathquadraticcurveto{\pgfqpoint{6.860247in}{0.514314in}}{\pgfqpoint{6.860247in}{0.503614in}}%
\pgfpathlineto{\pgfqpoint{6.860247in}{0.470400in}}%
\pgfpathlineto{\pgfqpoint{6.849836in}{0.470400in}}%
\pgfpathlineto{\pgfqpoint{6.849836in}{0.533443in}}%
\pgfpathlineto{\pgfqpoint{6.860247in}{0.533443in}}%
\pgfpathlineto{\pgfqpoint{6.860247in}{0.523644in}}%
\pgfpathquadraticcurveto{\pgfqpoint{6.863525in}{0.529390in}}{\pgfqpoint{6.868749in}{0.532164in}}%
\pgfpathquadraticcurveto{\pgfqpoint{6.873990in}{0.534956in}}{\pgfqpoint{6.881483in}{0.534956in}}%
\pgfpathquadraticcurveto{\pgfqpoint{6.882546in}{0.534956in}}{\pgfqpoint{6.883843in}{0.534811in}}%
\pgfpathquadraticcurveto{\pgfqpoint{6.885140in}{0.534685in}}{\pgfqpoint{6.886707in}{0.534397in}}%
\pgfpathlineto{\pgfqpoint{6.886761in}{0.523770in}}%
\pgfpathlineto{\pgfqpoint{6.886761in}{0.523770in}}%
\pgfpathclose%
\pgfpathmoveto{\pgfqpoint{6.949011in}{0.504515in}}%
\pgfpathlineto{\pgfqpoint{6.949011in}{0.499454in}}%
\pgfpathlineto{\pgfqpoint{6.901387in}{0.499454in}}%
\pgfpathquadraticcurveto{\pgfqpoint{6.902071in}{0.488754in}}{\pgfqpoint{6.907835in}{0.483153in}}%
\pgfpathquadraticcurveto{\pgfqpoint{6.913599in}{0.477551in}}{\pgfqpoint{6.923902in}{0.477551in}}%
\pgfpathquadraticcurveto{\pgfqpoint{6.929864in}{0.477551in}}{\pgfqpoint{6.935466in}{0.479010in}}%
\pgfpathquadraticcurveto{\pgfqpoint{6.941067in}{0.480469in}}{\pgfqpoint{6.946597in}{0.483405in}}%
\pgfpathlineto{\pgfqpoint{6.946597in}{0.473606in}}%
\pgfpathquadraticcurveto{\pgfqpoint{6.941013in}{0.471247in}}{\pgfqpoint{6.935160in}{0.470004in}}%
\pgfpathquadraticcurveto{\pgfqpoint{6.929306in}{0.468761in}}{\pgfqpoint{6.923289in}{0.468761in}}%
\pgfpathquadraticcurveto{\pgfqpoint{6.908195in}{0.468761in}}{\pgfqpoint{6.899387in}{0.477533in}}%
\pgfpathquadraticcurveto{\pgfqpoint{6.890579in}{0.486323in}}{\pgfqpoint{6.890579in}{0.501309in}}%
\pgfpathquadraticcurveto{\pgfqpoint{6.890579in}{0.516781in}}{\pgfqpoint{6.898937in}{0.525859in}}%
\pgfpathquadraticcurveto{\pgfqpoint{6.907295in}{0.534956in}}{\pgfqpoint{6.921488in}{0.534956in}}%
\pgfpathquadraticcurveto{\pgfqpoint{6.934205in}{0.534956in}}{\pgfqpoint{6.941608in}{0.526760in}}%
\pgfpathquadraticcurveto{\pgfqpoint{6.949011in}{0.518583in}}{\pgfqpoint{6.949011in}{0.504515in}}%
\pgfpathlineto{\pgfqpoint{6.949011in}{0.504515in}}%
\pgfpathclose%
\pgfpathmoveto{\pgfqpoint{6.938654in}{0.507559in}}%
\pgfpathquadraticcurveto{\pgfqpoint{6.938546in}{0.516043in}}{\pgfqpoint{6.933899in}{0.521104in}}%
\pgfpathquadraticcurveto{\pgfqpoint{6.929252in}{0.526184in}}{\pgfqpoint{6.921596in}{0.526184in}}%
\pgfpathquadraticcurveto{\pgfqpoint{6.912933in}{0.526184in}}{\pgfqpoint{6.907727in}{0.521284in}}%
\pgfpathquadraticcurveto{\pgfqpoint{6.902521in}{0.516385in}}{\pgfqpoint{6.901729in}{0.507487in}}%
\pgfpathlineto{\pgfqpoint{6.938654in}{0.507559in}}%
\pgfpathlineto{\pgfqpoint{6.938654in}{0.507559in}}%
\pgfpathclose%
\pgfpathmoveto{\pgfqpoint{7.018430in}{0.508460in}}%
\pgfpathlineto{\pgfqpoint{7.018430in}{0.470400in}}%
\pgfpathlineto{\pgfqpoint{7.008073in}{0.470400in}}%
\pgfpathlineto{\pgfqpoint{7.008073in}{0.508117in}}%
\pgfpathquadraticcurveto{\pgfqpoint{7.008073in}{0.517069in}}{\pgfqpoint{7.004578in}{0.521500in}}%
\pgfpathquadraticcurveto{\pgfqpoint{7.001084in}{0.525949in}}{\pgfqpoint{6.994113in}{0.525949in}}%
\pgfpathquadraticcurveto{\pgfqpoint{6.985720in}{0.525949in}}{\pgfqpoint{6.980874in}{0.520600in}}%
\pgfpathquadraticcurveto{\pgfqpoint{6.976029in}{0.515268in}}{\pgfqpoint{6.976029in}{0.506028in}}%
\pgfpathlineto{\pgfqpoint{6.976029in}{0.470400in}}%
\pgfpathlineto{\pgfqpoint{6.965618in}{0.470400in}}%
\pgfpathlineto{\pgfqpoint{6.965618in}{0.533443in}}%
\pgfpathlineto{\pgfqpoint{6.976029in}{0.533443in}}%
\pgfpathlineto{\pgfqpoint{6.976029in}{0.523644in}}%
\pgfpathquadraticcurveto{\pgfqpoint{6.979758in}{0.529336in}}{\pgfqpoint{6.984783in}{0.532146in}}%
\pgfpathquadraticcurveto{\pgfqpoint{6.989826in}{0.534956in}}{\pgfqpoint{6.996419in}{0.534956in}}%
\pgfpathquadraticcurveto{\pgfqpoint{7.007280in}{0.534956in}}{\pgfqpoint{7.012846in}{0.528237in}}%
\pgfpathquadraticcurveto{\pgfqpoint{7.018430in}{0.521518in}}{\pgfqpoint{7.018430in}{0.508460in}}%
\pgfpathlineto{\pgfqpoint{7.018430in}{0.508460in}}%
\pgfpathclose%
\pgfpathmoveto{\pgfqpoint{7.084444in}{0.531029in}}%
\pgfpathlineto{\pgfqpoint{7.084444in}{0.521338in}}%
\pgfpathquadraticcurveto{\pgfqpoint{7.080049in}{0.523770in}}{\pgfqpoint{7.075636in}{0.524977in}}%
\pgfpathquadraticcurveto{\pgfqpoint{7.071223in}{0.526184in}}{\pgfqpoint{7.066720in}{0.526184in}}%
\pgfpathquadraticcurveto{\pgfqpoint{7.056633in}{0.526184in}}{\pgfqpoint{7.051050in}{0.519789in}}%
\pgfpathquadraticcurveto{\pgfqpoint{7.045484in}{0.513413in}}{\pgfqpoint{7.045484in}{0.501867in}}%
\pgfpathquadraticcurveto{\pgfqpoint{7.045484in}{0.490321in}}{\pgfqpoint{7.051050in}{0.483927in}}%
\pgfpathquadraticcurveto{\pgfqpoint{7.056633in}{0.477551in}}{\pgfqpoint{7.066720in}{0.477551in}}%
\pgfpathquadraticcurveto{\pgfqpoint{7.071223in}{0.477551in}}{\pgfqpoint{7.075636in}{0.478758in}}%
\pgfpathquadraticcurveto{\pgfqpoint{7.080049in}{0.479964in}}{\pgfqpoint{7.084444in}{0.482396in}}%
\pgfpathlineto{\pgfqpoint{7.084444in}{0.472814in}}%
\pgfpathquadraticcurveto{\pgfqpoint{7.080103in}{0.470796in}}{\pgfqpoint{7.075456in}{0.469788in}}%
\pgfpathquadraticcurveto{\pgfqpoint{7.070827in}{0.468761in}}{\pgfqpoint{7.065586in}{0.468761in}}%
\pgfpathquadraticcurveto{\pgfqpoint{7.051338in}{0.468761in}}{\pgfqpoint{7.042944in}{0.477713in}}%
\pgfpathquadraticcurveto{\pgfqpoint{7.034569in}{0.486665in}}{\pgfqpoint{7.034569in}{0.501867in}}%
\pgfpathquadraticcurveto{\pgfqpoint{7.034569in}{0.517286in}}{\pgfqpoint{7.043034in}{0.526112in}}%
\pgfpathquadraticcurveto{\pgfqpoint{7.051518in}{0.534956in}}{\pgfqpoint{7.066270in}{0.534956in}}%
\pgfpathquadraticcurveto{\pgfqpoint{7.071043in}{0.534956in}}{\pgfqpoint{7.075600in}{0.533965in}}%
\pgfpathquadraticcurveto{\pgfqpoint{7.080157in}{0.532992in}}{\pgfqpoint{7.084444in}{0.531029in}}%
\pgfpathlineto{\pgfqpoint{7.084444in}{0.531029in}}%
\pgfpathclose%
\pgfpathmoveto{\pgfqpoint{7.156385in}{0.504515in}}%
\pgfpathlineto{\pgfqpoint{7.156385in}{0.499454in}}%
\pgfpathlineto{\pgfqpoint{7.108761in}{0.499454in}}%
\pgfpathquadraticcurveto{\pgfqpoint{7.109445in}{0.488754in}}{\pgfqpoint{7.115209in}{0.483153in}}%
\pgfpathquadraticcurveto{\pgfqpoint{7.120973in}{0.477551in}}{\pgfqpoint{7.131276in}{0.477551in}}%
\pgfpathquadraticcurveto{\pgfqpoint{7.137238in}{0.477551in}}{\pgfqpoint{7.142840in}{0.479010in}}%
\pgfpathquadraticcurveto{\pgfqpoint{7.148441in}{0.480469in}}{\pgfqpoint{7.153971in}{0.483405in}}%
\pgfpathlineto{\pgfqpoint{7.153971in}{0.473606in}}%
\pgfpathquadraticcurveto{\pgfqpoint{7.148387in}{0.471247in}}{\pgfqpoint{7.142533in}{0.470004in}}%
\pgfpathquadraticcurveto{\pgfqpoint{7.136679in}{0.468761in}}{\pgfqpoint{7.130663in}{0.468761in}}%
\pgfpathquadraticcurveto{\pgfqpoint{7.115569in}{0.468761in}}{\pgfqpoint{7.106761in}{0.477533in}}%
\pgfpathquadraticcurveto{\pgfqpoint{7.097953in}{0.486323in}}{\pgfqpoint{7.097953in}{0.501309in}}%
\pgfpathquadraticcurveto{\pgfqpoint{7.097953in}{0.516781in}}{\pgfqpoint{7.106311in}{0.525859in}}%
\pgfpathquadraticcurveto{\pgfqpoint{7.114669in}{0.534956in}}{\pgfqpoint{7.128862in}{0.534956in}}%
\pgfpathquadraticcurveto{\pgfqpoint{7.141579in}{0.534956in}}{\pgfqpoint{7.148982in}{0.526760in}}%
\pgfpathquadraticcurveto{\pgfqpoint{7.156385in}{0.518583in}}{\pgfqpoint{7.156385in}{0.504515in}}%
\pgfpathlineto{\pgfqpoint{7.156385in}{0.504515in}}%
\pgfpathclose%
\pgfpathmoveto{\pgfqpoint{7.146028in}{0.507559in}}%
\pgfpathquadraticcurveto{\pgfqpoint{7.145920in}{0.516043in}}{\pgfqpoint{7.141273in}{0.521104in}}%
\pgfpathquadraticcurveto{\pgfqpoint{7.136625in}{0.526184in}}{\pgfqpoint{7.128970in}{0.526184in}}%
\pgfpathquadraticcurveto{\pgfqpoint{7.120306in}{0.526184in}}{\pgfqpoint{7.115101in}{0.521284in}}%
\pgfpathquadraticcurveto{\pgfqpoint{7.109895in}{0.516385in}}{\pgfqpoint{7.109103in}{0.507487in}}%
\pgfpathlineto{\pgfqpoint{7.146028in}{0.507559in}}%
\pgfpathlineto{\pgfqpoint{7.146028in}{0.507559in}}%
\pgfpathclose%
\pgfusepath{fill}%
\end{pgfscope}%
\begin{pgfscope}%
\pgfsetbuttcap%
\pgfsetroundjoin%
\definecolor{currentfill}{rgb}{0.309804,0.372549,0.419608}%
\pgfsetfillcolor{currentfill}%
\pgfsetlinewidth{0.000000pt}%
\definecolor{currentstroke}{rgb}{0.309804,0.372549,0.419608}%
\pgfsetstrokecolor{currentstroke}%
\pgfsetdash{}{0pt}%
\pgfpathmoveto{\pgfqpoint{7.260829in}{0.551347in}}%
\pgfpathlineto{\pgfqpoint{7.260829in}{0.533443in}}%
\pgfpathlineto{\pgfqpoint{7.282155in}{0.533443in}}%
\pgfpathlineto{\pgfqpoint{7.282155in}{0.525391in}}%
\pgfpathlineto{\pgfqpoint{7.260829in}{0.525391in}}%
\pgfpathlineto{\pgfqpoint{7.260829in}{0.491168in}}%
\pgfpathquadraticcurveto{\pgfqpoint{7.260829in}{0.483459in}}{\pgfqpoint{7.262936in}{0.481261in}}%
\pgfpathquadraticcurveto{\pgfqpoint{7.265044in}{0.479064in}}{\pgfqpoint{7.271528in}{0.479064in}}%
\pgfpathlineto{\pgfqpoint{7.282155in}{0.479064in}}%
\pgfpathlineto{\pgfqpoint{7.282155in}{0.470400in}}%
\pgfpathlineto{\pgfqpoint{7.271528in}{0.470400in}}%
\pgfpathquadraticcurveto{\pgfqpoint{7.259532in}{0.470400in}}{\pgfqpoint{7.254975in}{0.474867in}}%
\pgfpathquadraticcurveto{\pgfqpoint{7.250418in}{0.479352in}}{\pgfqpoint{7.250418in}{0.491168in}}%
\pgfpathlineto{\pgfqpoint{7.250418in}{0.525391in}}%
\pgfpathlineto{\pgfqpoint{7.242817in}{0.525391in}}%
\pgfpathlineto{\pgfqpoint{7.242817in}{0.533443in}}%
\pgfpathlineto{\pgfqpoint{7.250418in}{0.533443in}}%
\pgfpathlineto{\pgfqpoint{7.250418in}{0.551347in}}%
\pgfpathlineto{\pgfqpoint{7.260829in}{0.551347in}}%
\pgfpathlineto{\pgfqpoint{7.260829in}{0.551347in}}%
\pgfpathclose%
\pgfusepath{fill}%
\end{pgfscope}%
\begin{pgfscope}%
\pgfsetbuttcap%
\pgfsetroundjoin%
\definecolor{currentfill}{rgb}{0.309804,0.372549,0.419608}%
\pgfsetfillcolor{currentfill}%
\pgfsetlinewidth{0.000000pt}%
\definecolor{currentstroke}{rgb}{0.309804,0.372549,0.419608}%
\pgfsetstrokecolor{currentstroke}%
\pgfsetdash{}{0pt}%
\pgfpathmoveto{\pgfqpoint{7.388078in}{0.531587in}}%
\pgfpathlineto{\pgfqpoint{7.388078in}{0.521789in}}%
\pgfpathquadraticcurveto{\pgfqpoint{7.383701in}{0.524040in}}{\pgfqpoint{7.378964in}{0.525157in}}%
\pgfpathquadraticcurveto{\pgfqpoint{7.374244in}{0.526292in}}{\pgfqpoint{7.369165in}{0.526292in}}%
\pgfpathquadraticcurveto{\pgfqpoint{7.361456in}{0.526292in}}{\pgfqpoint{7.357601in}{0.523932in}}%
\pgfpathquadraticcurveto{\pgfqpoint{7.353747in}{0.521573in}}{\pgfqpoint{7.353747in}{0.516835in}}%
\pgfpathquadraticcurveto{\pgfqpoint{7.353747in}{0.513233in}}{\pgfqpoint{7.356502in}{0.511180in}}%
\pgfpathquadraticcurveto{\pgfqpoint{7.359258in}{0.509126in}}{\pgfqpoint{7.367598in}{0.507271in}}%
\pgfpathlineto{\pgfqpoint{7.371146in}{0.506478in}}%
\pgfpathquadraticcurveto{\pgfqpoint{7.382170in}{0.504119in}}{\pgfqpoint{7.386817in}{0.499814in}}%
\pgfpathquadraticcurveto{\pgfqpoint{7.391464in}{0.495509in}}{\pgfqpoint{7.391464in}{0.487800in}}%
\pgfpathquadraticcurveto{\pgfqpoint{7.391464in}{0.479010in}}{\pgfqpoint{7.384511in}{0.473876in}}%
\pgfpathquadraticcurveto{\pgfqpoint{7.377559in}{0.468761in}}{\pgfqpoint{7.365400in}{0.468761in}}%
\pgfpathquadraticcurveto{\pgfqpoint{7.360339in}{0.468761in}}{\pgfqpoint{7.354845in}{0.469752in}}%
\pgfpathquadraticcurveto{\pgfqpoint{7.349352in}{0.470742in}}{\pgfqpoint{7.343281in}{0.472706in}}%
\pgfpathlineto{\pgfqpoint{7.343281in}{0.483405in}}%
\pgfpathquadraticcurveto{\pgfqpoint{7.349027in}{0.480415in}}{\pgfqpoint{7.354593in}{0.478920in}}%
\pgfpathquadraticcurveto{\pgfqpoint{7.360159in}{0.477443in}}{\pgfqpoint{7.365635in}{0.477443in}}%
\pgfpathquadraticcurveto{\pgfqpoint{7.372947in}{0.477443in}}{\pgfqpoint{7.376874in}{0.479946in}}%
\pgfpathquadraticcurveto{\pgfqpoint{7.380819in}{0.482450in}}{\pgfqpoint{7.380819in}{0.487007in}}%
\pgfpathquadraticcurveto{\pgfqpoint{7.380819in}{0.491222in}}{\pgfqpoint{7.377973in}{0.493474in}}%
\pgfpathquadraticcurveto{\pgfqpoint{7.375145in}{0.495725in}}{\pgfqpoint{7.365508in}{0.497814in}}%
\pgfpathlineto{\pgfqpoint{7.361906in}{0.498661in}}%
\pgfpathquadraticcurveto{\pgfqpoint{7.352288in}{0.500678in}}{\pgfqpoint{7.348001in}{0.504875in}}%
\pgfpathquadraticcurveto{\pgfqpoint{7.343732in}{0.509072in}}{\pgfqpoint{7.343732in}{0.516385in}}%
\pgfpathquadraticcurveto{\pgfqpoint{7.343732in}{0.525283in}}{\pgfqpoint{7.350036in}{0.530110in}}%
\pgfpathquadraticcurveto{\pgfqpoint{7.356340in}{0.534956in}}{\pgfqpoint{7.367940in}{0.534956in}}%
\pgfpathquadraticcurveto{\pgfqpoint{7.373668in}{0.534956in}}{\pgfqpoint{7.378729in}{0.534109in}}%
\pgfpathquadraticcurveto{\pgfqpoint{7.383809in}{0.533280in}}{\pgfqpoint{7.388078in}{0.531587in}}%
\pgfpathlineto{\pgfqpoint{7.388078in}{0.531587in}}%
\pgfpathclose%
\pgfpathmoveto{\pgfqpoint{7.418194in}{0.551347in}}%
\pgfpathlineto{\pgfqpoint{7.418194in}{0.533443in}}%
\pgfpathlineto{\pgfqpoint{7.439520in}{0.533443in}}%
\pgfpathlineto{\pgfqpoint{7.439520in}{0.525391in}}%
\pgfpathlineto{\pgfqpoint{7.418194in}{0.525391in}}%
\pgfpathlineto{\pgfqpoint{7.418194in}{0.491168in}}%
\pgfpathquadraticcurveto{\pgfqpoint{7.418194in}{0.483459in}}{\pgfqpoint{7.420301in}{0.481261in}}%
\pgfpathquadraticcurveto{\pgfqpoint{7.422409in}{0.479064in}}{\pgfqpoint{7.428893in}{0.479064in}}%
\pgfpathlineto{\pgfqpoint{7.439520in}{0.479064in}}%
\pgfpathlineto{\pgfqpoint{7.439520in}{0.470400in}}%
\pgfpathlineto{\pgfqpoint{7.428893in}{0.470400in}}%
\pgfpathquadraticcurveto{\pgfqpoint{7.416897in}{0.470400in}}{\pgfqpoint{7.412340in}{0.474867in}}%
\pgfpathquadraticcurveto{\pgfqpoint{7.407783in}{0.479352in}}{\pgfqpoint{7.407783in}{0.491168in}}%
\pgfpathlineto{\pgfqpoint{7.407783in}{0.525391in}}%
\pgfpathlineto{\pgfqpoint{7.400182in}{0.525391in}}%
\pgfpathlineto{\pgfqpoint{7.400182in}{0.533443in}}%
\pgfpathlineto{\pgfqpoint{7.407783in}{0.533443in}}%
\pgfpathlineto{\pgfqpoint{7.407783in}{0.551347in}}%
\pgfpathlineto{\pgfqpoint{7.418194in}{0.551347in}}%
\pgfpathlineto{\pgfqpoint{7.418194in}{0.551347in}}%
\pgfpathclose%
\pgfpathmoveto{\pgfqpoint{7.481795in}{0.502083in}}%
\pgfpathquadraticcurveto{\pgfqpoint{7.469240in}{0.502083in}}{\pgfqpoint{7.464395in}{0.499219in}}%
\pgfpathquadraticcurveto{\pgfqpoint{7.459550in}{0.496356in}}{\pgfqpoint{7.459550in}{0.489421in}}%
\pgfpathquadraticcurveto{\pgfqpoint{7.459550in}{0.483909in}}{\pgfqpoint{7.463188in}{0.480667in}}%
\pgfpathquadraticcurveto{\pgfqpoint{7.466827in}{0.477443in}}{\pgfqpoint{7.473059in}{0.477443in}}%
\pgfpathquadraticcurveto{\pgfqpoint{7.481687in}{0.477443in}}{\pgfqpoint{7.486892in}{0.483549in}}%
\pgfpathquadraticcurveto{\pgfqpoint{7.492098in}{0.489655in}}{\pgfqpoint{7.492098in}{0.499778in}}%
\pgfpathlineto{\pgfqpoint{7.492098in}{0.502083in}}%
\pgfpathlineto{\pgfqpoint{7.481795in}{0.502083in}}%
\pgfpathlineto{\pgfqpoint{7.481795in}{0.502083in}}%
\pgfpathclose%
\pgfpathmoveto{\pgfqpoint{7.502455in}{0.506370in}}%
\pgfpathlineto{\pgfqpoint{7.502455in}{0.470400in}}%
\pgfpathlineto{\pgfqpoint{7.492098in}{0.470400in}}%
\pgfpathlineto{\pgfqpoint{7.492098in}{0.479964in}}%
\pgfpathquadraticcurveto{\pgfqpoint{7.488549in}{0.474237in}}{\pgfqpoint{7.483254in}{0.471499in}}%
\pgfpathquadraticcurveto{\pgfqpoint{7.477958in}{0.468761in}}{\pgfqpoint{7.470303in}{0.468761in}}%
\pgfpathquadraticcurveto{\pgfqpoint{7.460631in}{0.468761in}}{\pgfqpoint{7.454903in}{0.474201in}}%
\pgfpathquadraticcurveto{\pgfqpoint{7.449193in}{0.479640in}}{\pgfqpoint{7.449193in}{0.488754in}}%
\pgfpathquadraticcurveto{\pgfqpoint{7.449193in}{0.499382in}}{\pgfqpoint{7.456308in}{0.504785in}}%
\pgfpathquadraticcurveto{\pgfqpoint{7.463441in}{0.510189in}}{\pgfqpoint{7.477562in}{0.510189in}}%
\pgfpathlineto{\pgfqpoint{7.492098in}{0.510189in}}%
\pgfpathlineto{\pgfqpoint{7.492098in}{0.511216in}}%
\pgfpathquadraticcurveto{\pgfqpoint{7.492098in}{0.518366in}}{\pgfqpoint{7.487397in}{0.522275in}}%
\pgfpathquadraticcurveto{\pgfqpoint{7.482696in}{0.526184in}}{\pgfqpoint{7.474194in}{0.526184in}}%
\pgfpathquadraticcurveto{\pgfqpoint{7.468790in}{0.526184in}}{\pgfqpoint{7.463657in}{0.524887in}}%
\pgfpathquadraticcurveto{\pgfqpoint{7.458541in}{0.523590in}}{\pgfqpoint{7.453822in}{0.520996in}}%
\pgfpathlineto{\pgfqpoint{7.453822in}{0.530579in}}%
\pgfpathquadraticcurveto{\pgfqpoint{7.459496in}{0.532776in}}{\pgfqpoint{7.464845in}{0.533857in}}%
\pgfpathquadraticcurveto{\pgfqpoint{7.470195in}{0.534956in}}{\pgfqpoint{7.475257in}{0.534956in}}%
\pgfpathquadraticcurveto{\pgfqpoint{7.488946in}{0.534956in}}{\pgfqpoint{7.495700in}{0.527859in}}%
\pgfpathquadraticcurveto{\pgfqpoint{7.502455in}{0.520780in}}{\pgfqpoint{7.502455in}{0.506370in}}%
\pgfpathlineto{\pgfqpoint{7.502455in}{0.506370in}}%
\pgfpathclose%
\pgfpathmoveto{\pgfqpoint{7.534030in}{0.551347in}}%
\pgfpathlineto{\pgfqpoint{7.534030in}{0.533443in}}%
\pgfpathlineto{\pgfqpoint{7.555357in}{0.533443in}}%
\pgfpathlineto{\pgfqpoint{7.555357in}{0.525391in}}%
\pgfpathlineto{\pgfqpoint{7.534030in}{0.525391in}}%
\pgfpathlineto{\pgfqpoint{7.534030in}{0.491168in}}%
\pgfpathquadraticcurveto{\pgfqpoint{7.534030in}{0.483459in}}{\pgfqpoint{7.536138in}{0.481261in}}%
\pgfpathquadraticcurveto{\pgfqpoint{7.538245in}{0.479064in}}{\pgfqpoint{7.544729in}{0.479064in}}%
\pgfpathlineto{\pgfqpoint{7.555357in}{0.479064in}}%
\pgfpathlineto{\pgfqpoint{7.555357in}{0.470400in}}%
\pgfpathlineto{\pgfqpoint{7.544729in}{0.470400in}}%
\pgfpathquadraticcurveto{\pgfqpoint{7.532733in}{0.470400in}}{\pgfqpoint{7.528176in}{0.474867in}}%
\pgfpathquadraticcurveto{\pgfqpoint{7.523619in}{0.479352in}}{\pgfqpoint{7.523619in}{0.491168in}}%
\pgfpathlineto{\pgfqpoint{7.523619in}{0.525391in}}%
\pgfpathlineto{\pgfqpoint{7.516018in}{0.525391in}}%
\pgfpathlineto{\pgfqpoint{7.516018in}{0.533443in}}%
\pgfpathlineto{\pgfqpoint{7.523619in}{0.533443in}}%
\pgfpathlineto{\pgfqpoint{7.523619in}{0.551347in}}%
\pgfpathlineto{\pgfqpoint{7.534030in}{0.551347in}}%
\pgfpathlineto{\pgfqpoint{7.534030in}{0.551347in}}%
\pgfpathclose%
\pgfpathmoveto{\pgfqpoint{7.568974in}{0.533443in}}%
\pgfpathlineto{\pgfqpoint{7.579331in}{0.533443in}}%
\pgfpathlineto{\pgfqpoint{7.579331in}{0.470400in}}%
\pgfpathlineto{\pgfqpoint{7.568974in}{0.470400in}}%
\pgfpathlineto{\pgfqpoint{7.568974in}{0.533443in}}%
\pgfpathlineto{\pgfqpoint{7.568974in}{0.533443in}}%
\pgfpathclose%
\pgfpathmoveto{\pgfqpoint{7.568974in}{0.557993in}}%
\pgfpathlineto{\pgfqpoint{7.579331in}{0.557993in}}%
\pgfpathlineto{\pgfqpoint{7.579331in}{0.544862in}}%
\pgfpathlineto{\pgfqpoint{7.568974in}{0.544862in}}%
\pgfpathlineto{\pgfqpoint{7.568974in}{0.557993in}}%
\pgfpathlineto{\pgfqpoint{7.568974in}{0.557993in}}%
\pgfpathclose%
\pgfpathmoveto{\pgfqpoint{7.641184in}{0.531587in}}%
\pgfpathlineto{\pgfqpoint{7.641184in}{0.521789in}}%
\pgfpathquadraticcurveto{\pgfqpoint{7.636808in}{0.524040in}}{\pgfqpoint{7.632070in}{0.525157in}}%
\pgfpathquadraticcurveto{\pgfqpoint{7.627351in}{0.526292in}}{\pgfqpoint{7.622272in}{0.526292in}}%
\pgfpathquadraticcurveto{\pgfqpoint{7.614563in}{0.526292in}}{\pgfqpoint{7.610708in}{0.523932in}}%
\pgfpathquadraticcurveto{\pgfqpoint{7.606853in}{0.521573in}}{\pgfqpoint{7.606853in}{0.516835in}}%
\pgfpathquadraticcurveto{\pgfqpoint{7.606853in}{0.513233in}}{\pgfqpoint{7.609609in}{0.511180in}}%
\pgfpathquadraticcurveto{\pgfqpoint{7.612365in}{0.509126in}}{\pgfqpoint{7.620705in}{0.507271in}}%
\pgfpathlineto{\pgfqpoint{7.624253in}{0.506478in}}%
\pgfpathquadraticcurveto{\pgfqpoint{7.635276in}{0.504119in}}{\pgfqpoint{7.639924in}{0.499814in}}%
\pgfpathquadraticcurveto{\pgfqpoint{7.644571in}{0.495509in}}{\pgfqpoint{7.644571in}{0.487800in}}%
\pgfpathquadraticcurveto{\pgfqpoint{7.644571in}{0.479010in}}{\pgfqpoint{7.637618in}{0.473876in}}%
\pgfpathquadraticcurveto{\pgfqpoint{7.630665in}{0.468761in}}{\pgfqpoint{7.618507in}{0.468761in}}%
\pgfpathquadraticcurveto{\pgfqpoint{7.613446in}{0.468761in}}{\pgfqpoint{7.607952in}{0.469752in}}%
\pgfpathquadraticcurveto{\pgfqpoint{7.602458in}{0.470742in}}{\pgfqpoint{7.596388in}{0.472706in}}%
\pgfpathlineto{\pgfqpoint{7.596388in}{0.483405in}}%
\pgfpathquadraticcurveto{\pgfqpoint{7.602134in}{0.480415in}}{\pgfqpoint{7.607700in}{0.478920in}}%
\pgfpathquadraticcurveto{\pgfqpoint{7.613266in}{0.477443in}}{\pgfqpoint{7.618741in}{0.477443in}}%
\pgfpathquadraticcurveto{\pgfqpoint{7.626054in}{0.477443in}}{\pgfqpoint{7.629981in}{0.479946in}}%
\pgfpathquadraticcurveto{\pgfqpoint{7.633926in}{0.482450in}}{\pgfqpoint{7.633926in}{0.487007in}}%
\pgfpathquadraticcurveto{\pgfqpoint{7.633926in}{0.491222in}}{\pgfqpoint{7.631080in}{0.493474in}}%
\pgfpathquadraticcurveto{\pgfqpoint{7.628252in}{0.495725in}}{\pgfqpoint{7.618615in}{0.497814in}}%
\pgfpathlineto{\pgfqpoint{7.615013in}{0.498661in}}%
\pgfpathquadraticcurveto{\pgfqpoint{7.605394in}{0.500678in}}{\pgfqpoint{7.601107in}{0.504875in}}%
\pgfpathquadraticcurveto{\pgfqpoint{7.596839in}{0.509072in}}{\pgfqpoint{7.596839in}{0.516385in}}%
\pgfpathquadraticcurveto{\pgfqpoint{7.596839in}{0.525283in}}{\pgfqpoint{7.603143in}{0.530110in}}%
\pgfpathquadraticcurveto{\pgfqpoint{7.609447in}{0.534956in}}{\pgfqpoint{7.621047in}{0.534956in}}%
\pgfpathquadraticcurveto{\pgfqpoint{7.626775in}{0.534956in}}{\pgfqpoint{7.631836in}{0.534109in}}%
\pgfpathquadraticcurveto{\pgfqpoint{7.636916in}{0.533280in}}{\pgfqpoint{7.641184in}{0.531587in}}%
\pgfpathlineto{\pgfqpoint{7.641184in}{0.531587in}}%
\pgfpathclose%
\pgfpathmoveto{\pgfqpoint{7.671301in}{0.551347in}}%
\pgfpathlineto{\pgfqpoint{7.671301in}{0.533443in}}%
\pgfpathlineto{\pgfqpoint{7.692627in}{0.533443in}}%
\pgfpathlineto{\pgfqpoint{7.692627in}{0.525391in}}%
\pgfpathlineto{\pgfqpoint{7.671301in}{0.525391in}}%
\pgfpathlineto{\pgfqpoint{7.671301in}{0.491168in}}%
\pgfpathquadraticcurveto{\pgfqpoint{7.671301in}{0.483459in}}{\pgfqpoint{7.673408in}{0.481261in}}%
\pgfpathquadraticcurveto{\pgfqpoint{7.675516in}{0.479064in}}{\pgfqpoint{7.682000in}{0.479064in}}%
\pgfpathlineto{\pgfqpoint{7.692627in}{0.479064in}}%
\pgfpathlineto{\pgfqpoint{7.692627in}{0.470400in}}%
\pgfpathlineto{\pgfqpoint{7.682000in}{0.470400in}}%
\pgfpathquadraticcurveto{\pgfqpoint{7.670004in}{0.470400in}}{\pgfqpoint{7.665447in}{0.474867in}}%
\pgfpathquadraticcurveto{\pgfqpoint{7.660890in}{0.479352in}}{\pgfqpoint{7.660890in}{0.491168in}}%
\pgfpathlineto{\pgfqpoint{7.660890in}{0.525391in}}%
\pgfpathlineto{\pgfqpoint{7.653289in}{0.525391in}}%
\pgfpathlineto{\pgfqpoint{7.653289in}{0.533443in}}%
\pgfpathlineto{\pgfqpoint{7.660890in}{0.533443in}}%
\pgfpathlineto{\pgfqpoint{7.660890in}{0.551347in}}%
\pgfpathlineto{\pgfqpoint{7.671301in}{0.551347in}}%
\pgfpathlineto{\pgfqpoint{7.671301in}{0.551347in}}%
\pgfpathclose%
\pgfpathmoveto{\pgfqpoint{7.706244in}{0.533443in}}%
\pgfpathlineto{\pgfqpoint{7.716601in}{0.533443in}}%
\pgfpathlineto{\pgfqpoint{7.716601in}{0.470400in}}%
\pgfpathlineto{\pgfqpoint{7.706244in}{0.470400in}}%
\pgfpathlineto{\pgfqpoint{7.706244in}{0.533443in}}%
\pgfpathlineto{\pgfqpoint{7.706244in}{0.533443in}}%
\pgfpathclose%
\pgfpathmoveto{\pgfqpoint{7.706244in}{0.557993in}}%
\pgfpathlineto{\pgfqpoint{7.716601in}{0.557993in}}%
\pgfpathlineto{\pgfqpoint{7.716601in}{0.544862in}}%
\pgfpathlineto{\pgfqpoint{7.706244in}{0.544862in}}%
\pgfpathlineto{\pgfqpoint{7.706244in}{0.557993in}}%
\pgfpathlineto{\pgfqpoint{7.706244in}{0.557993in}}%
\pgfpathclose%
\pgfpathmoveto{\pgfqpoint{7.783643in}{0.531029in}}%
\pgfpathlineto{\pgfqpoint{7.783643in}{0.521338in}}%
\pgfpathquadraticcurveto{\pgfqpoint{7.779248in}{0.523770in}}{\pgfqpoint{7.774835in}{0.524977in}}%
\pgfpathquadraticcurveto{\pgfqpoint{7.770422in}{0.526184in}}{\pgfqpoint{7.765919in}{0.526184in}}%
\pgfpathquadraticcurveto{\pgfqpoint{7.755832in}{0.526184in}}{\pgfqpoint{7.750248in}{0.519789in}}%
\pgfpathquadraticcurveto{\pgfqpoint{7.744682in}{0.513413in}}{\pgfqpoint{7.744682in}{0.501867in}}%
\pgfpathquadraticcurveto{\pgfqpoint{7.744682in}{0.490321in}}{\pgfqpoint{7.750248in}{0.483927in}}%
\pgfpathquadraticcurveto{\pgfqpoint{7.755832in}{0.477551in}}{\pgfqpoint{7.765919in}{0.477551in}}%
\pgfpathquadraticcurveto{\pgfqpoint{7.770422in}{0.477551in}}{\pgfqpoint{7.774835in}{0.478758in}}%
\pgfpathquadraticcurveto{\pgfqpoint{7.779248in}{0.479964in}}{\pgfqpoint{7.783643in}{0.482396in}}%
\pgfpathlineto{\pgfqpoint{7.783643in}{0.472814in}}%
\pgfpathquadraticcurveto{\pgfqpoint{7.779302in}{0.470796in}}{\pgfqpoint{7.774655in}{0.469788in}}%
\pgfpathquadraticcurveto{\pgfqpoint{7.770025in}{0.468761in}}{\pgfqpoint{7.764784in}{0.468761in}}%
\pgfpathquadraticcurveto{\pgfqpoint{7.750536in}{0.468761in}}{\pgfqpoint{7.742143in}{0.477713in}}%
\pgfpathquadraticcurveto{\pgfqpoint{7.733767in}{0.486665in}}{\pgfqpoint{7.733767in}{0.501867in}}%
\pgfpathquadraticcurveto{\pgfqpoint{7.733767in}{0.517286in}}{\pgfqpoint{7.742233in}{0.526112in}}%
\pgfpathquadraticcurveto{\pgfqpoint{7.750716in}{0.534956in}}{\pgfqpoint{7.765468in}{0.534956in}}%
\pgfpathquadraticcurveto{\pgfqpoint{7.770242in}{0.534956in}}{\pgfqpoint{7.774799in}{0.533965in}}%
\pgfpathquadraticcurveto{\pgfqpoint{7.779356in}{0.532992in}}{\pgfqpoint{7.783643in}{0.531029in}}%
\pgfpathlineto{\pgfqpoint{7.783643in}{0.531029in}}%
\pgfpathclose%
\pgfpathmoveto{\pgfqpoint{7.803114in}{0.484702in}}%
\pgfpathlineto{\pgfqpoint{7.815002in}{0.484702in}}%
\pgfpathlineto{\pgfqpoint{7.815002in}{0.470400in}}%
\pgfpathlineto{\pgfqpoint{7.803114in}{0.470400in}}%
\pgfpathlineto{\pgfqpoint{7.803114in}{0.484702in}}%
\pgfpathlineto{\pgfqpoint{7.803114in}{0.484702in}}%
\pgfpathclose%
\pgfusepath{fill}%
\end{pgfscope}%
\begin{pgfscope}%
\pgfsetbuttcap%
\pgfsetroundjoin%
\definecolor{currentfill}{rgb}{0.309804,0.372549,0.419608}%
\pgfsetfillcolor{currentfill}%
\pgfsetlinewidth{0.000000pt}%
\definecolor{currentstroke}{rgb}{0.309804,0.372549,0.419608}%
\pgfsetstrokecolor{currentstroke}%
\pgfsetdash{}{0pt}%
\pgfpathmoveto{\pgfqpoint{7.968590in}{0.547960in}}%
\pgfpathlineto{\pgfqpoint{7.968590in}{0.535982in}}%
\pgfpathquadraticcurveto{\pgfqpoint{7.962844in}{0.541332in}}{\pgfqpoint{7.956342in}{0.543962in}}%
\pgfpathquadraticcurveto{\pgfqpoint{7.949839in}{0.546609in}}{\pgfqpoint{7.942526in}{0.546609in}}%
\pgfpathquadraticcurveto{\pgfqpoint{7.928117in}{0.546609in}}{\pgfqpoint{7.920461in}{0.537801in}}%
\pgfpathquadraticcurveto{\pgfqpoint{7.912806in}{0.528994in}}{\pgfqpoint{7.912806in}{0.512332in}}%
\pgfpathquadraticcurveto{\pgfqpoint{7.912806in}{0.495725in}}{\pgfqpoint{7.920461in}{0.486917in}}%
\pgfpathquadraticcurveto{\pgfqpoint{7.928117in}{0.478109in}}{\pgfqpoint{7.942526in}{0.478109in}}%
\pgfpathquadraticcurveto{\pgfqpoint{7.949839in}{0.478109in}}{\pgfqpoint{7.956342in}{0.480757in}}%
\pgfpathquadraticcurveto{\pgfqpoint{7.962844in}{0.483405in}}{\pgfqpoint{7.968590in}{0.488754in}}%
\pgfpathlineto{\pgfqpoint{7.968590in}{0.476866in}}%
\pgfpathquadraticcurveto{\pgfqpoint{7.962628in}{0.472814in}}{\pgfqpoint{7.955945in}{0.470778in}}%
\pgfpathquadraticcurveto{\pgfqpoint{7.949281in}{0.468761in}}{\pgfqpoint{7.941860in}{0.468761in}}%
\pgfpathquadraticcurveto{\pgfqpoint{7.922767in}{0.468761in}}{\pgfqpoint{7.911780in}{0.480433in}}%
\pgfpathquadraticcurveto{\pgfqpoint{7.900810in}{0.492123in}}{\pgfqpoint{7.900810in}{0.512332in}}%
\pgfpathquadraticcurveto{\pgfqpoint{7.900810in}{0.532596in}}{\pgfqpoint{7.911780in}{0.544268in}}%
\pgfpathquadraticcurveto{\pgfqpoint{7.922767in}{0.555958in}}{\pgfqpoint{7.941860in}{0.555958in}}%
\pgfpathquadraticcurveto{\pgfqpoint{7.949389in}{0.555958in}}{\pgfqpoint{7.956053in}{0.553958in}}%
\pgfpathquadraticcurveto{\pgfqpoint{7.962736in}{0.551959in}}{\pgfqpoint{7.968590in}{0.547960in}}%
\pgfpathlineto{\pgfqpoint{7.968590in}{0.547960in}}%
\pgfpathclose%
\pgfpathmoveto{\pgfqpoint{8.010126in}{0.526184in}}%
\pgfpathquadraticcurveto{\pgfqpoint{8.001804in}{0.526184in}}{\pgfqpoint{7.996959in}{0.519681in}}%
\pgfpathquadraticcurveto{\pgfqpoint{7.992114in}{0.513179in}}{\pgfqpoint{7.992114in}{0.501867in}}%
\pgfpathquadraticcurveto{\pgfqpoint{7.992114in}{0.490556in}}{\pgfqpoint{7.996923in}{0.484053in}}%
\pgfpathquadraticcurveto{\pgfqpoint{8.001750in}{0.477551in}}{\pgfqpoint{8.010126in}{0.477551in}}%
\pgfpathquadraticcurveto{\pgfqpoint{8.018411in}{0.477551in}}{\pgfqpoint{8.023239in}{0.484071in}}%
\pgfpathquadraticcurveto{\pgfqpoint{8.028084in}{0.490610in}}{\pgfqpoint{8.028084in}{0.501867in}}%
\pgfpathquadraticcurveto{\pgfqpoint{8.028084in}{0.513071in}}{\pgfqpoint{8.023239in}{0.519627in}}%
\pgfpathquadraticcurveto{\pgfqpoint{8.018411in}{0.526184in}}{\pgfqpoint{8.010126in}{0.526184in}}%
\pgfpathlineto{\pgfqpoint{8.010126in}{0.526184in}}%
\pgfpathclose%
\pgfpathmoveto{\pgfqpoint{8.010126in}{0.534956in}}%
\pgfpathquadraticcurveto{\pgfqpoint{8.023635in}{0.534956in}}{\pgfqpoint{8.031344in}{0.526166in}}%
\pgfpathquadraticcurveto{\pgfqpoint{8.039071in}{0.517394in}}{\pgfqpoint{8.039071in}{0.501867in}}%
\pgfpathquadraticcurveto{\pgfqpoint{8.039071in}{0.486395in}}{\pgfqpoint{8.031344in}{0.477569in}}%
\pgfpathquadraticcurveto{\pgfqpoint{8.023635in}{0.468761in}}{\pgfqpoint{8.010126in}{0.468761in}}%
\pgfpathquadraticcurveto{\pgfqpoint{7.996563in}{0.468761in}}{\pgfqpoint{7.988872in}{0.477569in}}%
\pgfpathquadraticcurveto{\pgfqpoint{7.981198in}{0.486395in}}{\pgfqpoint{7.981198in}{0.501867in}}%
\pgfpathquadraticcurveto{\pgfqpoint{7.981198in}{0.517394in}}{\pgfqpoint{7.988872in}{0.526166in}}%
\pgfpathquadraticcurveto{\pgfqpoint{7.996563in}{0.534956in}}{\pgfqpoint{8.010126in}{0.534956in}}%
\pgfpathlineto{\pgfqpoint{8.010126in}{0.534956in}}%
\pgfpathclose%
\pgfpathmoveto{\pgfqpoint{8.105320in}{0.521338in}}%
\pgfpathquadraticcurveto{\pgfqpoint{8.109211in}{0.528327in}}{\pgfqpoint{8.114614in}{0.531641in}}%
\pgfpathquadraticcurveto{\pgfqpoint{8.120018in}{0.534956in}}{\pgfqpoint{8.127331in}{0.534956in}}%
\pgfpathquadraticcurveto{\pgfqpoint{8.137184in}{0.534956in}}{\pgfqpoint{8.142533in}{0.528057in}}%
\pgfpathquadraticcurveto{\pgfqpoint{8.147883in}{0.521176in}}{\pgfqpoint{8.147883in}{0.508460in}}%
\pgfpathlineto{\pgfqpoint{8.147883in}{0.470400in}}%
\pgfpathlineto{\pgfqpoint{8.137472in}{0.470400in}}%
\pgfpathlineto{\pgfqpoint{8.137472in}{0.508117in}}%
\pgfpathquadraticcurveto{\pgfqpoint{8.137472in}{0.517178in}}{\pgfqpoint{8.134248in}{0.521555in}}%
\pgfpathquadraticcurveto{\pgfqpoint{8.131041in}{0.525949in}}{\pgfqpoint{8.124467in}{0.525949in}}%
\pgfpathquadraticcurveto{\pgfqpoint{8.116416in}{0.525949in}}{\pgfqpoint{8.111732in}{0.520600in}}%
\pgfpathquadraticcurveto{\pgfqpoint{8.107067in}{0.515268in}}{\pgfqpoint{8.107067in}{0.506028in}}%
\pgfpathlineto{\pgfqpoint{8.107067in}{0.470400in}}%
\pgfpathlineto{\pgfqpoint{8.096656in}{0.470400in}}%
\pgfpathlineto{\pgfqpoint{8.096656in}{0.508117in}}%
\pgfpathquadraticcurveto{\pgfqpoint{8.096656in}{0.517232in}}{\pgfqpoint{8.093450in}{0.521591in}}%
\pgfpathquadraticcurveto{\pgfqpoint{8.090244in}{0.525949in}}{\pgfqpoint{8.083543in}{0.525949in}}%
\pgfpathquadraticcurveto{\pgfqpoint{8.075600in}{0.525949in}}{\pgfqpoint{8.070917in}{0.520582in}}%
\pgfpathquadraticcurveto{\pgfqpoint{8.066252in}{0.515214in}}{\pgfqpoint{8.066252in}{0.506028in}}%
\pgfpathlineto{\pgfqpoint{8.066252in}{0.470400in}}%
\pgfpathlineto{\pgfqpoint{8.055841in}{0.470400in}}%
\pgfpathlineto{\pgfqpoint{8.055841in}{0.533443in}}%
\pgfpathlineto{\pgfqpoint{8.066252in}{0.533443in}}%
\pgfpathlineto{\pgfqpoint{8.066252in}{0.523644in}}%
\pgfpathquadraticcurveto{\pgfqpoint{8.069800in}{0.529444in}}{\pgfqpoint{8.074753in}{0.532200in}}%
\pgfpathquadraticcurveto{\pgfqpoint{8.079707in}{0.534956in}}{\pgfqpoint{8.086515in}{0.534956in}}%
\pgfpathquadraticcurveto{\pgfqpoint{8.093396in}{0.534956in}}{\pgfqpoint{8.098205in}{0.531461in}}%
\pgfpathquadraticcurveto{\pgfqpoint{8.103015in}{0.527985in}}{\pgfqpoint{8.105320in}{0.521338in}}%
\pgfpathlineto{\pgfqpoint{8.105320in}{0.521338in}}%
\pgfpathclose%
\pgfpathmoveto{\pgfqpoint{8.217608in}{0.521338in}}%
\pgfpathquadraticcurveto{\pgfqpoint{8.221498in}{0.528327in}}{\pgfqpoint{8.226902in}{0.531641in}}%
\pgfpathquadraticcurveto{\pgfqpoint{8.232306in}{0.534956in}}{\pgfqpoint{8.239619in}{0.534956in}}%
\pgfpathquadraticcurveto{\pgfqpoint{8.249471in}{0.534956in}}{\pgfqpoint{8.254821in}{0.528057in}}%
\pgfpathquadraticcurveto{\pgfqpoint{8.260171in}{0.521176in}}{\pgfqpoint{8.260171in}{0.508460in}}%
\pgfpathlineto{\pgfqpoint{8.260171in}{0.470400in}}%
\pgfpathlineto{\pgfqpoint{8.249760in}{0.470400in}}%
\pgfpathlineto{\pgfqpoint{8.249760in}{0.508117in}}%
\pgfpathquadraticcurveto{\pgfqpoint{8.249760in}{0.517178in}}{\pgfqpoint{8.246535in}{0.521555in}}%
\pgfpathquadraticcurveto{\pgfqpoint{8.243329in}{0.525949in}}{\pgfqpoint{8.236755in}{0.525949in}}%
\pgfpathquadraticcurveto{\pgfqpoint{8.228703in}{0.525949in}}{\pgfqpoint{8.224020in}{0.520600in}}%
\pgfpathquadraticcurveto{\pgfqpoint{8.219355in}{0.515268in}}{\pgfqpoint{8.219355in}{0.506028in}}%
\pgfpathlineto{\pgfqpoint{8.219355in}{0.470400in}}%
\pgfpathlineto{\pgfqpoint{8.208944in}{0.470400in}}%
\pgfpathlineto{\pgfqpoint{8.208944in}{0.508117in}}%
\pgfpathquadraticcurveto{\pgfqpoint{8.208944in}{0.517232in}}{\pgfqpoint{8.205738in}{0.521591in}}%
\pgfpathquadraticcurveto{\pgfqpoint{8.202532in}{0.525949in}}{\pgfqpoint{8.195831in}{0.525949in}}%
\pgfpathquadraticcurveto{\pgfqpoint{8.187888in}{0.525949in}}{\pgfqpoint{8.183205in}{0.520582in}}%
\pgfpathquadraticcurveto{\pgfqpoint{8.178539in}{0.515214in}}{\pgfqpoint{8.178539in}{0.506028in}}%
\pgfpathlineto{\pgfqpoint{8.178539in}{0.470400in}}%
\pgfpathlineto{\pgfqpoint{8.168128in}{0.470400in}}%
\pgfpathlineto{\pgfqpoint{8.168128in}{0.533443in}}%
\pgfpathlineto{\pgfqpoint{8.178539in}{0.533443in}}%
\pgfpathlineto{\pgfqpoint{8.178539in}{0.523644in}}%
\pgfpathquadraticcurveto{\pgfqpoint{8.182088in}{0.529444in}}{\pgfqpoint{8.187041in}{0.532200in}}%
\pgfpathquadraticcurveto{\pgfqpoint{8.191995in}{0.534956in}}{\pgfqpoint{8.198803in}{0.534956in}}%
\pgfpathquadraticcurveto{\pgfqpoint{8.205684in}{0.534956in}}{\pgfqpoint{8.210493in}{0.531461in}}%
\pgfpathquadraticcurveto{\pgfqpoint{8.215302in}{0.527985in}}{\pgfqpoint{8.217608in}{0.521338in}}%
\pgfpathlineto{\pgfqpoint{8.217608in}{0.521338in}}%
\pgfpathclose%
\pgfpathmoveto{\pgfqpoint{8.305237in}{0.526184in}}%
\pgfpathquadraticcurveto{\pgfqpoint{8.296915in}{0.526184in}}{\pgfqpoint{8.292070in}{0.519681in}}%
\pgfpathquadraticcurveto{\pgfqpoint{8.287225in}{0.513179in}}{\pgfqpoint{8.287225in}{0.501867in}}%
\pgfpathquadraticcurveto{\pgfqpoint{8.287225in}{0.490556in}}{\pgfqpoint{8.292034in}{0.484053in}}%
\pgfpathquadraticcurveto{\pgfqpoint{8.296861in}{0.477551in}}{\pgfqpoint{8.305237in}{0.477551in}}%
\pgfpathquadraticcurveto{\pgfqpoint{8.313523in}{0.477551in}}{\pgfqpoint{8.318350in}{0.484071in}}%
\pgfpathquadraticcurveto{\pgfqpoint{8.323195in}{0.490610in}}{\pgfqpoint{8.323195in}{0.501867in}}%
\pgfpathquadraticcurveto{\pgfqpoint{8.323195in}{0.513071in}}{\pgfqpoint{8.318350in}{0.519627in}}%
\pgfpathquadraticcurveto{\pgfqpoint{8.313523in}{0.526184in}}{\pgfqpoint{8.305237in}{0.526184in}}%
\pgfpathlineto{\pgfqpoint{8.305237in}{0.526184in}}%
\pgfpathclose%
\pgfpathmoveto{\pgfqpoint{8.305237in}{0.534956in}}%
\pgfpathquadraticcurveto{\pgfqpoint{8.318746in}{0.534956in}}{\pgfqpoint{8.326455in}{0.526166in}}%
\pgfpathquadraticcurveto{\pgfqpoint{8.334183in}{0.517394in}}{\pgfqpoint{8.334183in}{0.501867in}}%
\pgfpathquadraticcurveto{\pgfqpoint{8.334183in}{0.486395in}}{\pgfqpoint{8.326455in}{0.477569in}}%
\pgfpathquadraticcurveto{\pgfqpoint{8.318746in}{0.468761in}}{\pgfqpoint{8.305237in}{0.468761in}}%
\pgfpathquadraticcurveto{\pgfqpoint{8.291674in}{0.468761in}}{\pgfqpoint{8.283983in}{0.477569in}}%
\pgfpathquadraticcurveto{\pgfqpoint{8.276309in}{0.486395in}}{\pgfqpoint{8.276309in}{0.501867in}}%
\pgfpathquadraticcurveto{\pgfqpoint{8.276309in}{0.517394in}}{\pgfqpoint{8.283983in}{0.526166in}}%
\pgfpathquadraticcurveto{\pgfqpoint{8.291674in}{0.534956in}}{\pgfqpoint{8.305237in}{0.534956in}}%
\pgfpathlineto{\pgfqpoint{8.305237in}{0.534956in}}%
\pgfpathclose%
\pgfpathmoveto{\pgfqpoint{8.403763in}{0.508460in}}%
\pgfpathlineto{\pgfqpoint{8.403763in}{0.470400in}}%
\pgfpathlineto{\pgfqpoint{8.393406in}{0.470400in}}%
\pgfpathlineto{\pgfqpoint{8.393406in}{0.508117in}}%
\pgfpathquadraticcurveto{\pgfqpoint{8.393406in}{0.517069in}}{\pgfqpoint{8.389912in}{0.521500in}}%
\pgfpathquadraticcurveto{\pgfqpoint{8.386418in}{0.525949in}}{\pgfqpoint{8.379447in}{0.525949in}}%
\pgfpathquadraticcurveto{\pgfqpoint{8.371053in}{0.525949in}}{\pgfqpoint{8.366208in}{0.520600in}}%
\pgfpathquadraticcurveto{\pgfqpoint{8.361363in}{0.515268in}}{\pgfqpoint{8.361363in}{0.506028in}}%
\pgfpathlineto{\pgfqpoint{8.361363in}{0.470400in}}%
\pgfpathlineto{\pgfqpoint{8.350952in}{0.470400in}}%
\pgfpathlineto{\pgfqpoint{8.350952in}{0.533443in}}%
\pgfpathlineto{\pgfqpoint{8.361363in}{0.533443in}}%
\pgfpathlineto{\pgfqpoint{8.361363in}{0.523644in}}%
\pgfpathquadraticcurveto{\pgfqpoint{8.365091in}{0.529336in}}{\pgfqpoint{8.370117in}{0.532146in}}%
\pgfpathquadraticcurveto{\pgfqpoint{8.375160in}{0.534956in}}{\pgfqpoint{8.381753in}{0.534956in}}%
\pgfpathquadraticcurveto{\pgfqpoint{8.392614in}{0.534956in}}{\pgfqpoint{8.398180in}{0.528237in}}%
\pgfpathquadraticcurveto{\pgfqpoint{8.403763in}{0.521518in}}{\pgfqpoint{8.403763in}{0.508460in}}%
\pgfpathlineto{\pgfqpoint{8.403763in}{0.508460in}}%
\pgfpathclose%
\pgfusepath{fill}%
\end{pgfscope}%
\begin{pgfscope}%
\pgfsetbuttcap%
\pgfsetroundjoin%
\definecolor{currentfill}{rgb}{0.309804,0.372549,0.419608}%
\pgfsetfillcolor{currentfill}%
\pgfsetlinewidth{0.000000pt}%
\definecolor{currentstroke}{rgb}{0.309804,0.372549,0.419608}%
\pgfsetstrokecolor{currentstroke}%
\pgfsetdash{}{0pt}%
\pgfpathmoveto{\pgfqpoint{8.544933in}{0.508460in}}%
\pgfpathlineto{\pgfqpoint{8.544933in}{0.470400in}}%
\pgfpathlineto{\pgfqpoint{8.534576in}{0.470400in}}%
\pgfpathlineto{\pgfqpoint{8.534576in}{0.508117in}}%
\pgfpathquadraticcurveto{\pgfqpoint{8.534576in}{0.517069in}}{\pgfqpoint{8.531082in}{0.521500in}}%
\pgfpathquadraticcurveto{\pgfqpoint{8.527587in}{0.525949in}}{\pgfqpoint{8.520617in}{0.525949in}}%
\pgfpathquadraticcurveto{\pgfqpoint{8.512223in}{0.525949in}}{\pgfqpoint{8.507378in}{0.520600in}}%
\pgfpathquadraticcurveto{\pgfqpoint{8.502532in}{0.515268in}}{\pgfqpoint{8.502532in}{0.506028in}}%
\pgfpathlineto{\pgfqpoint{8.502532in}{0.470400in}}%
\pgfpathlineto{\pgfqpoint{8.492121in}{0.470400in}}%
\pgfpathlineto{\pgfqpoint{8.492121in}{0.557993in}}%
\pgfpathlineto{\pgfqpoint{8.502532in}{0.557993in}}%
\pgfpathlineto{\pgfqpoint{8.502532in}{0.523644in}}%
\pgfpathquadraticcurveto{\pgfqpoint{8.506261in}{0.529336in}}{\pgfqpoint{8.511286in}{0.532146in}}%
\pgfpathquadraticcurveto{\pgfqpoint{8.516330in}{0.534956in}}{\pgfqpoint{8.522922in}{0.534956in}}%
\pgfpathquadraticcurveto{\pgfqpoint{8.533783in}{0.534956in}}{\pgfqpoint{8.539349in}{0.528237in}}%
\pgfpathquadraticcurveto{\pgfqpoint{8.544933in}{0.521518in}}{\pgfqpoint{8.544933in}{0.508460in}}%
\pgfpathlineto{\pgfqpoint{8.544933in}{0.508460in}}%
\pgfpathclose%
\pgfpathmoveto{\pgfqpoint{8.589999in}{0.526184in}}%
\pgfpathquadraticcurveto{\pgfqpoint{8.581678in}{0.526184in}}{\pgfqpoint{8.576832in}{0.519681in}}%
\pgfpathquadraticcurveto{\pgfqpoint{8.571987in}{0.513179in}}{\pgfqpoint{8.571987in}{0.501867in}}%
\pgfpathquadraticcurveto{\pgfqpoint{8.571987in}{0.490556in}}{\pgfqpoint{8.576796in}{0.484053in}}%
\pgfpathquadraticcurveto{\pgfqpoint{8.581624in}{0.477551in}}{\pgfqpoint{8.589999in}{0.477551in}}%
\pgfpathquadraticcurveto{\pgfqpoint{8.598285in}{0.477551in}}{\pgfqpoint{8.603112in}{0.484071in}}%
\pgfpathquadraticcurveto{\pgfqpoint{8.607957in}{0.490610in}}{\pgfqpoint{8.607957in}{0.501867in}}%
\pgfpathquadraticcurveto{\pgfqpoint{8.607957in}{0.513071in}}{\pgfqpoint{8.603112in}{0.519627in}}%
\pgfpathquadraticcurveto{\pgfqpoint{8.598285in}{0.526184in}}{\pgfqpoint{8.589999in}{0.526184in}}%
\pgfpathlineto{\pgfqpoint{8.589999in}{0.526184in}}%
\pgfpathclose%
\pgfpathmoveto{\pgfqpoint{8.589999in}{0.534956in}}%
\pgfpathquadraticcurveto{\pgfqpoint{8.603508in}{0.534956in}}{\pgfqpoint{8.611218in}{0.526166in}}%
\pgfpathquadraticcurveto{\pgfqpoint{8.618945in}{0.517394in}}{\pgfqpoint{8.618945in}{0.501867in}}%
\pgfpathquadraticcurveto{\pgfqpoint{8.618945in}{0.486395in}}{\pgfqpoint{8.611218in}{0.477569in}}%
\pgfpathquadraticcurveto{\pgfqpoint{8.603508in}{0.468761in}}{\pgfqpoint{8.589999in}{0.468761in}}%
\pgfpathquadraticcurveto{\pgfqpoint{8.576436in}{0.468761in}}{\pgfqpoint{8.568745in}{0.477569in}}%
\pgfpathquadraticcurveto{\pgfqpoint{8.561072in}{0.486395in}}{\pgfqpoint{8.561072in}{0.501867in}}%
\pgfpathquadraticcurveto{\pgfqpoint{8.561072in}{0.517394in}}{\pgfqpoint{8.568745in}{0.526166in}}%
\pgfpathquadraticcurveto{\pgfqpoint{8.576436in}{0.534956in}}{\pgfqpoint{8.589999in}{0.534956in}}%
\pgfpathlineto{\pgfqpoint{8.589999in}{0.534956in}}%
\pgfpathclose%
\pgfpathmoveto{\pgfqpoint{8.672639in}{0.523770in}}%
\pgfpathquadraticcurveto{\pgfqpoint{8.670892in}{0.524779in}}{\pgfqpoint{8.668839in}{0.525247in}}%
\pgfpathquadraticcurveto{\pgfqpoint{8.666785in}{0.525733in}}{\pgfqpoint{8.664317in}{0.525733in}}%
\pgfpathquadraticcurveto{\pgfqpoint{8.655528in}{0.525733in}}{\pgfqpoint{8.650826in}{0.520023in}}%
\pgfpathquadraticcurveto{\pgfqpoint{8.646125in}{0.514314in}}{\pgfqpoint{8.646125in}{0.503614in}}%
\pgfpathlineto{\pgfqpoint{8.646125in}{0.470400in}}%
\pgfpathlineto{\pgfqpoint{8.635714in}{0.470400in}}%
\pgfpathlineto{\pgfqpoint{8.635714in}{0.533443in}}%
\pgfpathlineto{\pgfqpoint{8.646125in}{0.533443in}}%
\pgfpathlineto{\pgfqpoint{8.646125in}{0.523644in}}%
\pgfpathquadraticcurveto{\pgfqpoint{8.649403in}{0.529390in}}{\pgfqpoint{8.654627in}{0.532164in}}%
\pgfpathquadraticcurveto{\pgfqpoint{8.659868in}{0.534956in}}{\pgfqpoint{8.667362in}{0.534956in}}%
\pgfpathquadraticcurveto{\pgfqpoint{8.668424in}{0.534956in}}{\pgfqpoint{8.669721in}{0.534811in}}%
\pgfpathquadraticcurveto{\pgfqpoint{8.671018in}{0.534685in}}{\pgfqpoint{8.672585in}{0.534397in}}%
\pgfpathlineto{\pgfqpoint{8.672639in}{0.523770in}}%
\pgfpathlineto{\pgfqpoint{8.672639in}{0.523770in}}%
\pgfpathclose%
\pgfpathmoveto{\pgfqpoint{8.683500in}{0.533443in}}%
\pgfpathlineto{\pgfqpoint{8.693857in}{0.533443in}}%
\pgfpathlineto{\pgfqpoint{8.693857in}{0.470400in}}%
\pgfpathlineto{\pgfqpoint{8.683500in}{0.470400in}}%
\pgfpathlineto{\pgfqpoint{8.683500in}{0.533443in}}%
\pgfpathlineto{\pgfqpoint{8.683500in}{0.533443in}}%
\pgfpathclose%
\pgfpathmoveto{\pgfqpoint{8.683500in}{0.557993in}}%
\pgfpathlineto{\pgfqpoint{8.693857in}{0.557993in}}%
\pgfpathlineto{\pgfqpoint{8.693857in}{0.544862in}}%
\pgfpathlineto{\pgfqpoint{8.683500in}{0.544862in}}%
\pgfpathlineto{\pgfqpoint{8.683500in}{0.557993in}}%
\pgfpathlineto{\pgfqpoint{8.683500in}{0.557993in}}%
\pgfpathclose%
\pgfpathmoveto{\pgfqpoint{8.711023in}{0.533443in}}%
\pgfpathlineto{\pgfqpoint{8.760214in}{0.533443in}}%
\pgfpathlineto{\pgfqpoint{8.760214in}{0.523986in}}%
\pgfpathlineto{\pgfqpoint{8.721272in}{0.478668in}}%
\pgfpathlineto{\pgfqpoint{8.760214in}{0.478668in}}%
\pgfpathlineto{\pgfqpoint{8.760214in}{0.470400in}}%
\pgfpathlineto{\pgfqpoint{8.709618in}{0.470400in}}%
\pgfpathlineto{\pgfqpoint{8.709618in}{0.479856in}}%
\pgfpathlineto{\pgfqpoint{8.748578in}{0.525175in}}%
\pgfpathlineto{\pgfqpoint{8.711023in}{0.525175in}}%
\pgfpathlineto{\pgfqpoint{8.711023in}{0.533443in}}%
\pgfpathlineto{\pgfqpoint{8.711023in}{0.533443in}}%
\pgfpathclose%
\pgfpathmoveto{\pgfqpoint{8.800453in}{0.526184in}}%
\pgfpathquadraticcurveto{\pgfqpoint{8.792132in}{0.526184in}}{\pgfqpoint{8.787286in}{0.519681in}}%
\pgfpathquadraticcurveto{\pgfqpoint{8.782441in}{0.513179in}}{\pgfqpoint{8.782441in}{0.501867in}}%
\pgfpathquadraticcurveto{\pgfqpoint{8.782441in}{0.490556in}}{\pgfqpoint{8.787250in}{0.484053in}}%
\pgfpathquadraticcurveto{\pgfqpoint{8.792078in}{0.477551in}}{\pgfqpoint{8.800453in}{0.477551in}}%
\pgfpathquadraticcurveto{\pgfqpoint{8.808739in}{0.477551in}}{\pgfqpoint{8.813566in}{0.484071in}}%
\pgfpathquadraticcurveto{\pgfqpoint{8.818411in}{0.490610in}}{\pgfqpoint{8.818411in}{0.501867in}}%
\pgfpathquadraticcurveto{\pgfqpoint{8.818411in}{0.513071in}}{\pgfqpoint{8.813566in}{0.519627in}}%
\pgfpathquadraticcurveto{\pgfqpoint{8.808739in}{0.526184in}}{\pgfqpoint{8.800453in}{0.526184in}}%
\pgfpathlineto{\pgfqpoint{8.800453in}{0.526184in}}%
\pgfpathclose%
\pgfpathmoveto{\pgfqpoint{8.800453in}{0.534956in}}%
\pgfpathquadraticcurveto{\pgfqpoint{8.813962in}{0.534956in}}{\pgfqpoint{8.821672in}{0.526166in}}%
\pgfpathquadraticcurveto{\pgfqpoint{8.829399in}{0.517394in}}{\pgfqpoint{8.829399in}{0.501867in}}%
\pgfpathquadraticcurveto{\pgfqpoint{8.829399in}{0.486395in}}{\pgfqpoint{8.821672in}{0.477569in}}%
\pgfpathquadraticcurveto{\pgfqpoint{8.813962in}{0.468761in}}{\pgfqpoint{8.800453in}{0.468761in}}%
\pgfpathquadraticcurveto{\pgfqpoint{8.786890in}{0.468761in}}{\pgfqpoint{8.779199in}{0.477569in}}%
\pgfpathquadraticcurveto{\pgfqpoint{8.771526in}{0.486395in}}{\pgfqpoint{8.771526in}{0.501867in}}%
\pgfpathquadraticcurveto{\pgfqpoint{8.771526in}{0.517394in}}{\pgfqpoint{8.779199in}{0.526166in}}%
\pgfpathquadraticcurveto{\pgfqpoint{8.786890in}{0.534956in}}{\pgfqpoint{8.800453in}{0.534956in}}%
\pgfpathlineto{\pgfqpoint{8.800453in}{0.534956in}}%
\pgfpathclose%
\pgfpathmoveto{\pgfqpoint{8.898980in}{0.508460in}}%
\pgfpathlineto{\pgfqpoint{8.898980in}{0.470400in}}%
\pgfpathlineto{\pgfqpoint{8.888623in}{0.470400in}}%
\pgfpathlineto{\pgfqpoint{8.888623in}{0.508117in}}%
\pgfpathquadraticcurveto{\pgfqpoint{8.888623in}{0.517069in}}{\pgfqpoint{8.885128in}{0.521500in}}%
\pgfpathquadraticcurveto{\pgfqpoint{8.881634in}{0.525949in}}{\pgfqpoint{8.874663in}{0.525949in}}%
\pgfpathquadraticcurveto{\pgfqpoint{8.866270in}{0.525949in}}{\pgfqpoint{8.861424in}{0.520600in}}%
\pgfpathquadraticcurveto{\pgfqpoint{8.856579in}{0.515268in}}{\pgfqpoint{8.856579in}{0.506028in}}%
\pgfpathlineto{\pgfqpoint{8.856579in}{0.470400in}}%
\pgfpathlineto{\pgfqpoint{8.846168in}{0.470400in}}%
\pgfpathlineto{\pgfqpoint{8.846168in}{0.533443in}}%
\pgfpathlineto{\pgfqpoint{8.856579in}{0.533443in}}%
\pgfpathlineto{\pgfqpoint{8.856579in}{0.523644in}}%
\pgfpathquadraticcurveto{\pgfqpoint{8.860308in}{0.529336in}}{\pgfqpoint{8.865333in}{0.532146in}}%
\pgfpathquadraticcurveto{\pgfqpoint{8.870377in}{0.534956in}}{\pgfqpoint{8.876969in}{0.534956in}}%
\pgfpathquadraticcurveto{\pgfqpoint{8.887830in}{0.534956in}}{\pgfqpoint{8.893396in}{0.528237in}}%
\pgfpathquadraticcurveto{\pgfqpoint{8.898980in}{0.521518in}}{\pgfqpoint{8.898980in}{0.508460in}}%
\pgfpathlineto{\pgfqpoint{8.898980in}{0.508460in}}%
\pgfpathclose%
\pgfpathmoveto{\pgfqpoint{8.959807in}{0.531587in}}%
\pgfpathlineto{\pgfqpoint{8.959807in}{0.521789in}}%
\pgfpathquadraticcurveto{\pgfqpoint{8.955430in}{0.524040in}}{\pgfqpoint{8.950693in}{0.525157in}}%
\pgfpathquadraticcurveto{\pgfqpoint{8.945974in}{0.526292in}}{\pgfqpoint{8.940894in}{0.526292in}}%
\pgfpathquadraticcurveto{\pgfqpoint{8.933185in}{0.526292in}}{\pgfqpoint{8.929330in}{0.523932in}}%
\pgfpathquadraticcurveto{\pgfqpoint{8.925476in}{0.521573in}}{\pgfqpoint{8.925476in}{0.516835in}}%
\pgfpathquadraticcurveto{\pgfqpoint{8.925476in}{0.513233in}}{\pgfqpoint{8.928232in}{0.511180in}}%
\pgfpathquadraticcurveto{\pgfqpoint{8.930987in}{0.509126in}}{\pgfqpoint{8.939327in}{0.507271in}}%
\pgfpathlineto{\pgfqpoint{8.942875in}{0.506478in}}%
\pgfpathquadraticcurveto{\pgfqpoint{8.953899in}{0.504119in}}{\pgfqpoint{8.958546in}{0.499814in}}%
\pgfpathquadraticcurveto{\pgfqpoint{8.963193in}{0.495509in}}{\pgfqpoint{8.963193in}{0.487800in}}%
\pgfpathquadraticcurveto{\pgfqpoint{8.963193in}{0.479010in}}{\pgfqpoint{8.956240in}{0.473876in}}%
\pgfpathquadraticcurveto{\pgfqpoint{8.949288in}{0.468761in}}{\pgfqpoint{8.937130in}{0.468761in}}%
\pgfpathquadraticcurveto{\pgfqpoint{8.932068in}{0.468761in}}{\pgfqpoint{8.926574in}{0.469752in}}%
\pgfpathquadraticcurveto{\pgfqpoint{8.921081in}{0.470742in}}{\pgfqpoint{8.915011in}{0.472706in}}%
\pgfpathlineto{\pgfqpoint{8.915011in}{0.483405in}}%
\pgfpathquadraticcurveto{\pgfqpoint{8.920757in}{0.480415in}}{\pgfqpoint{8.926322in}{0.478920in}}%
\pgfpathquadraticcurveto{\pgfqpoint{8.931888in}{0.477443in}}{\pgfqpoint{8.937364in}{0.477443in}}%
\pgfpathquadraticcurveto{\pgfqpoint{8.944677in}{0.477443in}}{\pgfqpoint{8.948603in}{0.479946in}}%
\pgfpathquadraticcurveto{\pgfqpoint{8.952548in}{0.482450in}}{\pgfqpoint{8.952548in}{0.487007in}}%
\pgfpathquadraticcurveto{\pgfqpoint{8.952548in}{0.491222in}}{\pgfqpoint{8.949702in}{0.493474in}}%
\pgfpathquadraticcurveto{\pgfqpoint{8.946874in}{0.495725in}}{\pgfqpoint{8.937238in}{0.497814in}}%
\pgfpathlineto{\pgfqpoint{8.933635in}{0.498661in}}%
\pgfpathquadraticcurveto{\pgfqpoint{8.924017in}{0.500678in}}{\pgfqpoint{8.919730in}{0.504875in}}%
\pgfpathquadraticcurveto{\pgfqpoint{8.915461in}{0.509072in}}{\pgfqpoint{8.915461in}{0.516385in}}%
\pgfpathquadraticcurveto{\pgfqpoint{8.915461in}{0.525283in}}{\pgfqpoint{8.921765in}{0.530110in}}%
\pgfpathquadraticcurveto{\pgfqpoint{8.928069in}{0.534956in}}{\pgfqpoint{8.939669in}{0.534956in}}%
\pgfpathquadraticcurveto{\pgfqpoint{8.945397in}{0.534956in}}{\pgfqpoint{8.950459in}{0.534109in}}%
\pgfpathquadraticcurveto{\pgfqpoint{8.955538in}{0.533280in}}{\pgfqpoint{8.959807in}{0.531587in}}%
\pgfpathlineto{\pgfqpoint{8.959807in}{0.531587in}}%
\pgfpathclose%
\pgfusepath{fill}%
\end{pgfscope}%
\begin{pgfscope}%
\pgfsetbuttcap%
\pgfsetroundjoin%
\definecolor{currentfill}{rgb}{0.309804,0.372549,0.419608}%
\pgfsetfillcolor{currentfill}%
\pgfsetlinewidth{0.000000pt}%
\definecolor{currentstroke}{rgb}{0.309804,0.372549,0.419608}%
\pgfsetstrokecolor{currentstroke}%
\pgfsetdash{}{0pt}%
\pgfpathmoveto{\pgfqpoint{9.068487in}{0.502083in}}%
\pgfpathquadraticcurveto{\pgfqpoint{9.055933in}{0.502083in}}{\pgfqpoint{9.051088in}{0.499219in}}%
\pgfpathquadraticcurveto{\pgfqpoint{9.046242in}{0.496356in}}{\pgfqpoint{9.046242in}{0.489421in}}%
\pgfpathquadraticcurveto{\pgfqpoint{9.046242in}{0.483909in}}{\pgfqpoint{9.049881in}{0.480667in}}%
\pgfpathquadraticcurveto{\pgfqpoint{9.053519in}{0.477443in}}{\pgfqpoint{9.059752in}{0.477443in}}%
\pgfpathquadraticcurveto{\pgfqpoint{9.068379in}{0.477443in}}{\pgfqpoint{9.073585in}{0.483549in}}%
\pgfpathquadraticcurveto{\pgfqpoint{9.078790in}{0.489655in}}{\pgfqpoint{9.078790in}{0.499778in}}%
\pgfpathlineto{\pgfqpoint{9.078790in}{0.502083in}}%
\pgfpathlineto{\pgfqpoint{9.068487in}{0.502083in}}%
\pgfpathlineto{\pgfqpoint{9.068487in}{0.502083in}}%
\pgfpathclose%
\pgfpathmoveto{\pgfqpoint{9.089147in}{0.506370in}}%
\pgfpathlineto{\pgfqpoint{9.089147in}{0.470400in}}%
\pgfpathlineto{\pgfqpoint{9.078790in}{0.470400in}}%
\pgfpathlineto{\pgfqpoint{9.078790in}{0.479964in}}%
\pgfpathquadraticcurveto{\pgfqpoint{9.075242in}{0.474237in}}{\pgfqpoint{9.069946in}{0.471499in}}%
\pgfpathquadraticcurveto{\pgfqpoint{9.064651in}{0.468761in}}{\pgfqpoint{9.056996in}{0.468761in}}%
\pgfpathquadraticcurveto{\pgfqpoint{9.047323in}{0.468761in}}{\pgfqpoint{9.041595in}{0.474201in}}%
\pgfpathquadraticcurveto{\pgfqpoint{9.035885in}{0.479640in}}{\pgfqpoint{9.035885in}{0.488754in}}%
\pgfpathquadraticcurveto{\pgfqpoint{9.035885in}{0.499382in}}{\pgfqpoint{9.043000in}{0.504785in}}%
\pgfpathquadraticcurveto{\pgfqpoint{9.050133in}{0.510189in}}{\pgfqpoint{9.064255in}{0.510189in}}%
\pgfpathlineto{\pgfqpoint{9.078790in}{0.510189in}}%
\pgfpathlineto{\pgfqpoint{9.078790in}{0.511216in}}%
\pgfpathquadraticcurveto{\pgfqpoint{9.078790in}{0.518366in}}{\pgfqpoint{9.074089in}{0.522275in}}%
\pgfpathquadraticcurveto{\pgfqpoint{9.069388in}{0.526184in}}{\pgfqpoint{9.060886in}{0.526184in}}%
\pgfpathquadraticcurveto{\pgfqpoint{9.055483in}{0.526184in}}{\pgfqpoint{9.050349in}{0.524887in}}%
\pgfpathquadraticcurveto{\pgfqpoint{9.045234in}{0.523590in}}{\pgfqpoint{9.040515in}{0.520996in}}%
\pgfpathlineto{\pgfqpoint{9.040515in}{0.530579in}}%
\pgfpathquadraticcurveto{\pgfqpoint{9.046188in}{0.532776in}}{\pgfqpoint{9.051538in}{0.533857in}}%
\pgfpathquadraticcurveto{\pgfqpoint{9.056888in}{0.534956in}}{\pgfqpoint{9.061949in}{0.534956in}}%
\pgfpathquadraticcurveto{\pgfqpoint{9.075638in}{0.534956in}}{\pgfqpoint{9.082393in}{0.527859in}}%
\pgfpathquadraticcurveto{\pgfqpoint{9.089147in}{0.520780in}}{\pgfqpoint{9.089147in}{0.506370in}}%
\pgfpathlineto{\pgfqpoint{9.089147in}{0.506370in}}%
\pgfpathclose%
\pgfpathmoveto{\pgfqpoint{9.147002in}{0.523770in}}%
\pgfpathquadraticcurveto{\pgfqpoint{9.145255in}{0.524779in}}{\pgfqpoint{9.143202in}{0.525247in}}%
\pgfpathquadraticcurveto{\pgfqpoint{9.141148in}{0.525733in}}{\pgfqpoint{9.138681in}{0.525733in}}%
\pgfpathquadraticcurveto{\pgfqpoint{9.129891in}{0.525733in}}{\pgfqpoint{9.125190in}{0.520023in}}%
\pgfpathquadraticcurveto{\pgfqpoint{9.120488in}{0.514314in}}{\pgfqpoint{9.120488in}{0.503614in}}%
\pgfpathlineto{\pgfqpoint{9.120488in}{0.470400in}}%
\pgfpathlineto{\pgfqpoint{9.110077in}{0.470400in}}%
\pgfpathlineto{\pgfqpoint{9.110077in}{0.533443in}}%
\pgfpathlineto{\pgfqpoint{9.120488in}{0.533443in}}%
\pgfpathlineto{\pgfqpoint{9.120488in}{0.523644in}}%
\pgfpathquadraticcurveto{\pgfqpoint{9.123767in}{0.529390in}}{\pgfqpoint{9.128990in}{0.532164in}}%
\pgfpathquadraticcurveto{\pgfqpoint{9.134232in}{0.534956in}}{\pgfqpoint{9.141725in}{0.534956in}}%
\pgfpathquadraticcurveto{\pgfqpoint{9.142788in}{0.534956in}}{\pgfqpoint{9.144084in}{0.534811in}}%
\pgfpathquadraticcurveto{\pgfqpoint{9.145381in}{0.534685in}}{\pgfqpoint{9.146948in}{0.534397in}}%
\pgfpathlineto{\pgfqpoint{9.147002in}{0.523770in}}%
\pgfpathlineto{\pgfqpoint{9.147002in}{0.523770in}}%
\pgfpathclose%
\pgfpathmoveto{\pgfqpoint{9.209252in}{0.504515in}}%
\pgfpathlineto{\pgfqpoint{9.209252in}{0.499454in}}%
\pgfpathlineto{\pgfqpoint{9.161628in}{0.499454in}}%
\pgfpathquadraticcurveto{\pgfqpoint{9.162313in}{0.488754in}}{\pgfqpoint{9.168077in}{0.483153in}}%
\pgfpathquadraticcurveto{\pgfqpoint{9.173840in}{0.477551in}}{\pgfqpoint{9.184143in}{0.477551in}}%
\pgfpathquadraticcurveto{\pgfqpoint{9.190105in}{0.477551in}}{\pgfqpoint{9.195707in}{0.479010in}}%
\pgfpathquadraticcurveto{\pgfqpoint{9.201309in}{0.480469in}}{\pgfqpoint{9.206839in}{0.483405in}}%
\pgfpathlineto{\pgfqpoint{9.206839in}{0.473606in}}%
\pgfpathquadraticcurveto{\pgfqpoint{9.201255in}{0.471247in}}{\pgfqpoint{9.195401in}{0.470004in}}%
\pgfpathquadraticcurveto{\pgfqpoint{9.189547in}{0.468761in}}{\pgfqpoint{9.183531in}{0.468761in}}%
\pgfpathquadraticcurveto{\pgfqpoint{9.168437in}{0.468761in}}{\pgfqpoint{9.159629in}{0.477533in}}%
\pgfpathquadraticcurveto{\pgfqpoint{9.150821in}{0.486323in}}{\pgfqpoint{9.150821in}{0.501309in}}%
\pgfpathquadraticcurveto{\pgfqpoint{9.150821in}{0.516781in}}{\pgfqpoint{9.159179in}{0.525859in}}%
\pgfpathquadraticcurveto{\pgfqpoint{9.167536in}{0.534956in}}{\pgfqpoint{9.181730in}{0.534956in}}%
\pgfpathquadraticcurveto{\pgfqpoint{9.194446in}{0.534956in}}{\pgfqpoint{9.201849in}{0.526760in}}%
\pgfpathquadraticcurveto{\pgfqpoint{9.209252in}{0.518583in}}{\pgfqpoint{9.209252in}{0.504515in}}%
\pgfpathlineto{\pgfqpoint{9.209252in}{0.504515in}}%
\pgfpathclose%
\pgfpathmoveto{\pgfqpoint{9.198895in}{0.507559in}}%
\pgfpathquadraticcurveto{\pgfqpoint{9.198787in}{0.516043in}}{\pgfqpoint{9.194140in}{0.521104in}}%
\pgfpathquadraticcurveto{\pgfqpoint{9.189493in}{0.526184in}}{\pgfqpoint{9.181838in}{0.526184in}}%
\pgfpathquadraticcurveto{\pgfqpoint{9.173174in}{0.526184in}}{\pgfqpoint{9.167969in}{0.521284in}}%
\pgfpathquadraticcurveto{\pgfqpoint{9.162763in}{0.516385in}}{\pgfqpoint{9.161970in}{0.507487in}}%
\pgfpathlineto{\pgfqpoint{9.198895in}{0.507559in}}%
\pgfpathlineto{\pgfqpoint{9.198895in}{0.507559in}}%
\pgfpathclose%
\pgfusepath{fill}%
\end{pgfscope}%
\begin{pgfscope}%
\pgfsetbuttcap%
\pgfsetroundjoin%
\definecolor{currentfill}{rgb}{0.309804,0.372549,0.419608}%
\pgfsetfillcolor{currentfill}%
\pgfsetlinewidth{0.000000pt}%
\definecolor{currentstroke}{rgb}{0.309804,0.372549,0.419608}%
\pgfsetstrokecolor{currentstroke}%
\pgfsetdash{}{0pt}%
\pgfpathmoveto{\pgfqpoint{9.273580in}{0.495275in}}%
\pgfpathlineto{\pgfqpoint{9.273580in}{0.533443in}}%
\pgfpathlineto{\pgfqpoint{9.283937in}{0.533443in}}%
\pgfpathlineto{\pgfqpoint{9.283937in}{0.495671in}}%
\pgfpathquadraticcurveto{\pgfqpoint{9.283937in}{0.486719in}}{\pgfqpoint{9.287413in}{0.482234in}}%
\pgfpathquadraticcurveto{\pgfqpoint{9.290908in}{0.477767in}}{\pgfqpoint{9.297896in}{0.477767in}}%
\pgfpathquadraticcurveto{\pgfqpoint{9.306272in}{0.477767in}}{\pgfqpoint{9.311135in}{0.483117in}}%
\pgfpathquadraticcurveto{\pgfqpoint{9.316016in}{0.488466in}}{\pgfqpoint{9.316016in}{0.497706in}}%
\pgfpathlineto{\pgfqpoint{9.316016in}{0.533443in}}%
\pgfpathlineto{\pgfqpoint{9.326373in}{0.533443in}}%
\pgfpathlineto{\pgfqpoint{9.326373in}{0.470400in}}%
\pgfpathlineto{\pgfqpoint{9.316016in}{0.470400in}}%
\pgfpathlineto{\pgfqpoint{9.316016in}{0.480091in}}%
\pgfpathquadraticcurveto{\pgfqpoint{9.312252in}{0.474345in}}{\pgfqpoint{9.307263in}{0.471553in}}%
\pgfpathquadraticcurveto{\pgfqpoint{9.302291in}{0.468761in}}{\pgfqpoint{9.295699in}{0.468761in}}%
\pgfpathquadraticcurveto{\pgfqpoint{9.284837in}{0.468761in}}{\pgfqpoint{9.279200in}{0.475515in}}%
\pgfpathquadraticcurveto{\pgfqpoint{9.273580in}{0.482270in}}{\pgfqpoint{9.273580in}{0.495275in}}%
\pgfpathlineto{\pgfqpoint{9.273580in}{0.495275in}}%
\pgfpathclose%
\pgfpathmoveto{\pgfqpoint{9.299643in}{0.534956in}}%
\pgfpathlineto{\pgfqpoint{9.299643in}{0.534956in}}%
\pgfpathlineto{\pgfqpoint{9.299643in}{0.534956in}}%
\pgfpathclose%
\pgfpathmoveto{\pgfqpoint{9.400115in}{0.508460in}}%
\pgfpathlineto{\pgfqpoint{9.400115in}{0.470400in}}%
\pgfpathlineto{\pgfqpoint{9.389758in}{0.470400in}}%
\pgfpathlineto{\pgfqpoint{9.389758in}{0.508117in}}%
\pgfpathquadraticcurveto{\pgfqpoint{9.389758in}{0.517069in}}{\pgfqpoint{9.386264in}{0.521500in}}%
\pgfpathquadraticcurveto{\pgfqpoint{9.382770in}{0.525949in}}{\pgfqpoint{9.375799in}{0.525949in}}%
\pgfpathquadraticcurveto{\pgfqpoint{9.367405in}{0.525949in}}{\pgfqpoint{9.362560in}{0.520600in}}%
\pgfpathquadraticcurveto{\pgfqpoint{9.357715in}{0.515268in}}{\pgfqpoint{9.357715in}{0.506028in}}%
\pgfpathlineto{\pgfqpoint{9.357715in}{0.470400in}}%
\pgfpathlineto{\pgfqpoint{9.347304in}{0.470400in}}%
\pgfpathlineto{\pgfqpoint{9.347304in}{0.533443in}}%
\pgfpathlineto{\pgfqpoint{9.357715in}{0.533443in}}%
\pgfpathlineto{\pgfqpoint{9.357715in}{0.523644in}}%
\pgfpathquadraticcurveto{\pgfqpoint{9.361443in}{0.529336in}}{\pgfqpoint{9.366469in}{0.532146in}}%
\pgfpathquadraticcurveto{\pgfqpoint{9.371512in}{0.534956in}}{\pgfqpoint{9.378104in}{0.534956in}}%
\pgfpathquadraticcurveto{\pgfqpoint{9.388966in}{0.534956in}}{\pgfqpoint{9.394531in}{0.528237in}}%
\pgfpathquadraticcurveto{\pgfqpoint{9.400115in}{0.521518in}}{\pgfqpoint{9.400115in}{0.508460in}}%
\pgfpathlineto{\pgfqpoint{9.400115in}{0.508460in}}%
\pgfpathclose%
\pgfpathmoveto{\pgfqpoint{9.466130in}{0.531029in}}%
\pgfpathlineto{\pgfqpoint{9.466130in}{0.521338in}}%
\pgfpathquadraticcurveto{\pgfqpoint{9.461735in}{0.523770in}}{\pgfqpoint{9.457322in}{0.524977in}}%
\pgfpathquadraticcurveto{\pgfqpoint{9.452909in}{0.526184in}}{\pgfqpoint{9.448406in}{0.526184in}}%
\pgfpathquadraticcurveto{\pgfqpoint{9.438319in}{0.526184in}}{\pgfqpoint{9.432735in}{0.519789in}}%
\pgfpathquadraticcurveto{\pgfqpoint{9.427169in}{0.513413in}}{\pgfqpoint{9.427169in}{0.501867in}}%
\pgfpathquadraticcurveto{\pgfqpoint{9.427169in}{0.490321in}}{\pgfqpoint{9.432735in}{0.483927in}}%
\pgfpathquadraticcurveto{\pgfqpoint{9.438319in}{0.477551in}}{\pgfqpoint{9.448406in}{0.477551in}}%
\pgfpathquadraticcurveto{\pgfqpoint{9.452909in}{0.477551in}}{\pgfqpoint{9.457322in}{0.478758in}}%
\pgfpathquadraticcurveto{\pgfqpoint{9.461735in}{0.479964in}}{\pgfqpoint{9.466130in}{0.482396in}}%
\pgfpathlineto{\pgfqpoint{9.466130in}{0.472814in}}%
\pgfpathquadraticcurveto{\pgfqpoint{9.461789in}{0.470796in}}{\pgfqpoint{9.457142in}{0.469788in}}%
\pgfpathquadraticcurveto{\pgfqpoint{9.452513in}{0.468761in}}{\pgfqpoint{9.447271in}{0.468761in}}%
\pgfpathquadraticcurveto{\pgfqpoint{9.433023in}{0.468761in}}{\pgfqpoint{9.424630in}{0.477713in}}%
\pgfpathquadraticcurveto{\pgfqpoint{9.416254in}{0.486665in}}{\pgfqpoint{9.416254in}{0.501867in}}%
\pgfpathquadraticcurveto{\pgfqpoint{9.416254in}{0.517286in}}{\pgfqpoint{9.424720in}{0.526112in}}%
\pgfpathquadraticcurveto{\pgfqpoint{9.433204in}{0.534956in}}{\pgfqpoint{9.447956in}{0.534956in}}%
\pgfpathquadraticcurveto{\pgfqpoint{9.452729in}{0.534956in}}{\pgfqpoint{9.457286in}{0.533965in}}%
\pgfpathquadraticcurveto{\pgfqpoint{9.461843in}{0.532992in}}{\pgfqpoint{9.466130in}{0.531029in}}%
\pgfpathlineto{\pgfqpoint{9.466130in}{0.531029in}}%
\pgfpathclose%
\pgfpathmoveto{\pgfqpoint{9.536557in}{0.508460in}}%
\pgfpathlineto{\pgfqpoint{9.536557in}{0.470400in}}%
\pgfpathlineto{\pgfqpoint{9.526200in}{0.470400in}}%
\pgfpathlineto{\pgfqpoint{9.526200in}{0.508117in}}%
\pgfpathquadraticcurveto{\pgfqpoint{9.526200in}{0.517069in}}{\pgfqpoint{9.522706in}{0.521500in}}%
\pgfpathquadraticcurveto{\pgfqpoint{9.519212in}{0.525949in}}{\pgfqpoint{9.512241in}{0.525949in}}%
\pgfpathquadraticcurveto{\pgfqpoint{9.503847in}{0.525949in}}{\pgfqpoint{9.499002in}{0.520600in}}%
\pgfpathquadraticcurveto{\pgfqpoint{9.494157in}{0.515268in}}{\pgfqpoint{9.494157in}{0.506028in}}%
\pgfpathlineto{\pgfqpoint{9.494157in}{0.470400in}}%
\pgfpathlineto{\pgfqpoint{9.483746in}{0.470400in}}%
\pgfpathlineto{\pgfqpoint{9.483746in}{0.557993in}}%
\pgfpathlineto{\pgfqpoint{9.494157in}{0.557993in}}%
\pgfpathlineto{\pgfqpoint{9.494157in}{0.523644in}}%
\pgfpathquadraticcurveto{\pgfqpoint{9.497885in}{0.529336in}}{\pgfqpoint{9.502911in}{0.532146in}}%
\pgfpathquadraticcurveto{\pgfqpoint{9.507954in}{0.534956in}}{\pgfqpoint{9.514546in}{0.534956in}}%
\pgfpathquadraticcurveto{\pgfqpoint{9.525408in}{0.534956in}}{\pgfqpoint{9.530974in}{0.528237in}}%
\pgfpathquadraticcurveto{\pgfqpoint{9.536557in}{0.521518in}}{\pgfqpoint{9.536557in}{0.508460in}}%
\pgfpathlineto{\pgfqpoint{9.536557in}{0.508460in}}%
\pgfpathclose%
\pgfpathmoveto{\pgfqpoint{9.585857in}{0.502083in}}%
\pgfpathquadraticcurveto{\pgfqpoint{9.573302in}{0.502083in}}{\pgfqpoint{9.568457in}{0.499219in}}%
\pgfpathquadraticcurveto{\pgfqpoint{9.563612in}{0.496356in}}{\pgfqpoint{9.563612in}{0.489421in}}%
\pgfpathquadraticcurveto{\pgfqpoint{9.563612in}{0.483909in}}{\pgfqpoint{9.567250in}{0.480667in}}%
\pgfpathquadraticcurveto{\pgfqpoint{9.570888in}{0.477443in}}{\pgfqpoint{9.577121in}{0.477443in}}%
\pgfpathquadraticcurveto{\pgfqpoint{9.585748in}{0.477443in}}{\pgfqpoint{9.590954in}{0.483549in}}%
\pgfpathquadraticcurveto{\pgfqpoint{9.596160in}{0.489655in}}{\pgfqpoint{9.596160in}{0.499778in}}%
\pgfpathlineto{\pgfqpoint{9.596160in}{0.502083in}}%
\pgfpathlineto{\pgfqpoint{9.585857in}{0.502083in}}%
\pgfpathlineto{\pgfqpoint{9.585857in}{0.502083in}}%
\pgfpathclose%
\pgfpathmoveto{\pgfqpoint{9.606516in}{0.506370in}}%
\pgfpathlineto{\pgfqpoint{9.606516in}{0.470400in}}%
\pgfpathlineto{\pgfqpoint{9.596160in}{0.470400in}}%
\pgfpathlineto{\pgfqpoint{9.596160in}{0.479964in}}%
\pgfpathquadraticcurveto{\pgfqpoint{9.592611in}{0.474237in}}{\pgfqpoint{9.587316in}{0.471499in}}%
\pgfpathquadraticcurveto{\pgfqpoint{9.582020in}{0.468761in}}{\pgfqpoint{9.574365in}{0.468761in}}%
\pgfpathquadraticcurveto{\pgfqpoint{9.564692in}{0.468761in}}{\pgfqpoint{9.558964in}{0.474201in}}%
\pgfpathquadraticcurveto{\pgfqpoint{9.553255in}{0.479640in}}{\pgfqpoint{9.553255in}{0.488754in}}%
\pgfpathquadraticcurveto{\pgfqpoint{9.553255in}{0.499382in}}{\pgfqpoint{9.560369in}{0.504785in}}%
\pgfpathquadraticcurveto{\pgfqpoint{9.567502in}{0.510189in}}{\pgfqpoint{9.581624in}{0.510189in}}%
\pgfpathlineto{\pgfqpoint{9.596160in}{0.510189in}}%
\pgfpathlineto{\pgfqpoint{9.596160in}{0.511216in}}%
\pgfpathquadraticcurveto{\pgfqpoint{9.596160in}{0.518366in}}{\pgfqpoint{9.591458in}{0.522275in}}%
\pgfpathquadraticcurveto{\pgfqpoint{9.586757in}{0.526184in}}{\pgfqpoint{9.578255in}{0.526184in}}%
\pgfpathquadraticcurveto{\pgfqpoint{9.572852in}{0.526184in}}{\pgfqpoint{9.567718in}{0.524887in}}%
\pgfpathquadraticcurveto{\pgfqpoint{9.562603in}{0.523590in}}{\pgfqpoint{9.557884in}{0.520996in}}%
\pgfpathlineto{\pgfqpoint{9.557884in}{0.530579in}}%
\pgfpathquadraticcurveto{\pgfqpoint{9.563558in}{0.532776in}}{\pgfqpoint{9.568907in}{0.533857in}}%
\pgfpathquadraticcurveto{\pgfqpoint{9.574257in}{0.534956in}}{\pgfqpoint{9.579318in}{0.534956in}}%
\pgfpathquadraticcurveto{\pgfqpoint{9.593007in}{0.534956in}}{\pgfqpoint{9.599762in}{0.527859in}}%
\pgfpathquadraticcurveto{\pgfqpoint{9.606516in}{0.520780in}}{\pgfqpoint{9.606516in}{0.506370in}}%
\pgfpathlineto{\pgfqpoint{9.606516in}{0.506370in}}%
\pgfpathclose%
\pgfpathmoveto{\pgfqpoint{9.680258in}{0.508460in}}%
\pgfpathlineto{\pgfqpoint{9.680258in}{0.470400in}}%
\pgfpathlineto{\pgfqpoint{9.669901in}{0.470400in}}%
\pgfpathlineto{\pgfqpoint{9.669901in}{0.508117in}}%
\pgfpathquadraticcurveto{\pgfqpoint{9.669901in}{0.517069in}}{\pgfqpoint{9.666407in}{0.521500in}}%
\pgfpathquadraticcurveto{\pgfqpoint{9.662913in}{0.525949in}}{\pgfqpoint{9.655942in}{0.525949in}}%
\pgfpathquadraticcurveto{\pgfqpoint{9.647548in}{0.525949in}}{\pgfqpoint{9.642703in}{0.520600in}}%
\pgfpathquadraticcurveto{\pgfqpoint{9.637858in}{0.515268in}}{\pgfqpoint{9.637858in}{0.506028in}}%
\pgfpathlineto{\pgfqpoint{9.637858in}{0.470400in}}%
\pgfpathlineto{\pgfqpoint{9.627447in}{0.470400in}}%
\pgfpathlineto{\pgfqpoint{9.627447in}{0.533443in}}%
\pgfpathlineto{\pgfqpoint{9.637858in}{0.533443in}}%
\pgfpathlineto{\pgfqpoint{9.637858in}{0.523644in}}%
\pgfpathquadraticcurveto{\pgfqpoint{9.641586in}{0.529336in}}{\pgfqpoint{9.646612in}{0.532146in}}%
\pgfpathquadraticcurveto{\pgfqpoint{9.651655in}{0.534956in}}{\pgfqpoint{9.658247in}{0.534956in}}%
\pgfpathquadraticcurveto{\pgfqpoint{9.669109in}{0.534956in}}{\pgfqpoint{9.674674in}{0.528237in}}%
\pgfpathquadraticcurveto{\pgfqpoint{9.680258in}{0.521518in}}{\pgfqpoint{9.680258in}{0.508460in}}%
\pgfpathlineto{\pgfqpoint{9.680258in}{0.508460in}}%
\pgfpathclose%
\pgfpathmoveto{\pgfqpoint{9.742382in}{0.502660in}}%
\pgfpathquadraticcurveto{\pgfqpoint{9.742382in}{0.513917in}}{\pgfqpoint{9.737735in}{0.520096in}}%
\pgfpathquadraticcurveto{\pgfqpoint{9.733106in}{0.526292in}}{\pgfqpoint{9.724712in}{0.526292in}}%
\pgfpathquadraticcurveto{\pgfqpoint{9.716391in}{0.526292in}}{\pgfqpoint{9.711743in}{0.520096in}}%
\pgfpathquadraticcurveto{\pgfqpoint{9.707096in}{0.513917in}}{\pgfqpoint{9.707096in}{0.502660in}}%
\pgfpathquadraticcurveto{\pgfqpoint{9.707096in}{0.491456in}}{\pgfqpoint{9.711743in}{0.485260in}}%
\pgfpathquadraticcurveto{\pgfqpoint{9.716391in}{0.479064in}}{\pgfqpoint{9.724712in}{0.479064in}}%
\pgfpathquadraticcurveto{\pgfqpoint{9.733106in}{0.479064in}}{\pgfqpoint{9.737735in}{0.485260in}}%
\pgfpathquadraticcurveto{\pgfqpoint{9.742382in}{0.491456in}}{\pgfqpoint{9.742382in}{0.502660in}}%
\pgfpathlineto{\pgfqpoint{9.742382in}{0.502660in}}%
\pgfpathclose%
\pgfpathmoveto{\pgfqpoint{9.752739in}{0.478217in}}%
\pgfpathquadraticcurveto{\pgfqpoint{9.752739in}{0.462132in}}{\pgfqpoint{9.745588in}{0.454279in}}%
\pgfpathquadraticcurveto{\pgfqpoint{9.738456in}{0.446426in}}{\pgfqpoint{9.723704in}{0.446426in}}%
\pgfpathquadraticcurveto{\pgfqpoint{9.718246in}{0.446426in}}{\pgfqpoint{9.713401in}{0.447236in}}%
\pgfpathquadraticcurveto{\pgfqpoint{9.708555in}{0.448047in}}{\pgfqpoint{9.703998in}{0.449740in}}%
\pgfpathlineto{\pgfqpoint{9.703998in}{0.459809in}}%
\pgfpathquadraticcurveto{\pgfqpoint{9.708555in}{0.457341in}}{\pgfqpoint{9.713004in}{0.456170in}}%
\pgfpathquadraticcurveto{\pgfqpoint{9.717453in}{0.454982in}}{\pgfqpoint{9.722064in}{0.454982in}}%
\pgfpathquadraticcurveto{\pgfqpoint{9.732259in}{0.454982in}}{\pgfqpoint{9.737321in}{0.460295in}}%
\pgfpathquadraticcurveto{\pgfqpoint{9.742382in}{0.465609in}}{\pgfqpoint{9.742382in}{0.476362in}}%
\pgfpathlineto{\pgfqpoint{9.742382in}{0.481495in}}%
\pgfpathquadraticcurveto{\pgfqpoint{9.739176in}{0.475912in}}{\pgfqpoint{9.734169in}{0.473156in}}%
\pgfpathquadraticcurveto{\pgfqpoint{9.729161in}{0.470400in}}{\pgfqpoint{9.722173in}{0.470400in}}%
\pgfpathquadraticcurveto{\pgfqpoint{9.710591in}{0.470400in}}{\pgfqpoint{9.703494in}{0.479226in}}%
\pgfpathquadraticcurveto{\pgfqpoint{9.696397in}{0.488070in}}{\pgfqpoint{9.696397in}{0.502660in}}%
\pgfpathquadraticcurveto{\pgfqpoint{9.696397in}{0.517286in}}{\pgfqpoint{9.703494in}{0.526112in}}%
\pgfpathquadraticcurveto{\pgfqpoint{9.710591in}{0.534956in}}{\pgfqpoint{9.722173in}{0.534956in}}%
\pgfpathquadraticcurveto{\pgfqpoint{9.729161in}{0.534956in}}{\pgfqpoint{9.734169in}{0.532200in}}%
\pgfpathquadraticcurveto{\pgfqpoint{9.739176in}{0.529444in}}{\pgfqpoint{9.742382in}{0.523878in}}%
\pgfpathlineto{\pgfqpoint{9.742382in}{0.533443in}}%
\pgfpathlineto{\pgfqpoint{9.752739in}{0.533443in}}%
\pgfpathlineto{\pgfqpoint{9.752739in}{0.478217in}}%
\pgfpathlineto{\pgfqpoint{9.752739in}{0.478217in}}%
\pgfpathclose%
\pgfpathmoveto{\pgfqpoint{9.828012in}{0.504515in}}%
\pgfpathlineto{\pgfqpoint{9.828012in}{0.499454in}}%
\pgfpathlineto{\pgfqpoint{9.780388in}{0.499454in}}%
\pgfpathquadraticcurveto{\pgfqpoint{9.781072in}{0.488754in}}{\pgfqpoint{9.786836in}{0.483153in}}%
\pgfpathquadraticcurveto{\pgfqpoint{9.792600in}{0.477551in}}{\pgfqpoint{9.802903in}{0.477551in}}%
\pgfpathquadraticcurveto{\pgfqpoint{9.808865in}{0.477551in}}{\pgfqpoint{9.814467in}{0.479010in}}%
\pgfpathquadraticcurveto{\pgfqpoint{9.820069in}{0.480469in}}{\pgfqpoint{9.825598in}{0.483405in}}%
\pgfpathlineto{\pgfqpoint{9.825598in}{0.473606in}}%
\pgfpathquadraticcurveto{\pgfqpoint{9.820015in}{0.471247in}}{\pgfqpoint{9.814161in}{0.470004in}}%
\pgfpathquadraticcurveto{\pgfqpoint{9.808307in}{0.468761in}}{\pgfqpoint{9.802291in}{0.468761in}}%
\pgfpathquadraticcurveto{\pgfqpoint{9.787196in}{0.468761in}}{\pgfqpoint{9.778388in}{0.477533in}}%
\pgfpathquadraticcurveto{\pgfqpoint{9.769581in}{0.486323in}}{\pgfqpoint{9.769581in}{0.501309in}}%
\pgfpathquadraticcurveto{\pgfqpoint{9.769581in}{0.516781in}}{\pgfqpoint{9.777938in}{0.525859in}}%
\pgfpathquadraticcurveto{\pgfqpoint{9.786296in}{0.534956in}}{\pgfqpoint{9.800489in}{0.534956in}}%
\pgfpathquadraticcurveto{\pgfqpoint{9.813206in}{0.534956in}}{\pgfqpoint{9.820609in}{0.526760in}}%
\pgfpathquadraticcurveto{\pgfqpoint{9.828012in}{0.518583in}}{\pgfqpoint{9.828012in}{0.504515in}}%
\pgfpathlineto{\pgfqpoint{9.828012in}{0.504515in}}%
\pgfpathclose%
\pgfpathmoveto{\pgfqpoint{9.817655in}{0.507559in}}%
\pgfpathquadraticcurveto{\pgfqpoint{9.817547in}{0.516043in}}{\pgfqpoint{9.812900in}{0.521104in}}%
\pgfpathquadraticcurveto{\pgfqpoint{9.808253in}{0.526184in}}{\pgfqpoint{9.800597in}{0.526184in}}%
\pgfpathquadraticcurveto{\pgfqpoint{9.791934in}{0.526184in}}{\pgfqpoint{9.786728in}{0.521284in}}%
\pgfpathquadraticcurveto{\pgfqpoint{9.781523in}{0.516385in}}{\pgfqpoint{9.780730in}{0.507487in}}%
\pgfpathlineto{\pgfqpoint{9.817655in}{0.507559in}}%
\pgfpathlineto{\pgfqpoint{9.817655in}{0.507559in}}%
\pgfpathclose%
\pgfpathmoveto{\pgfqpoint{9.886497in}{0.523878in}}%
\pgfpathlineto{\pgfqpoint{9.886497in}{0.557993in}}%
\pgfpathlineto{\pgfqpoint{9.896854in}{0.557993in}}%
\pgfpathlineto{\pgfqpoint{9.896854in}{0.470400in}}%
\pgfpathlineto{\pgfqpoint{9.886497in}{0.470400in}}%
\pgfpathlineto{\pgfqpoint{9.886497in}{0.479856in}}%
\pgfpathquadraticcurveto{\pgfqpoint{9.883237in}{0.474237in}}{\pgfqpoint{9.878248in}{0.471499in}}%
\pgfpathquadraticcurveto{\pgfqpoint{9.873276in}{0.468761in}}{\pgfqpoint{9.866288in}{0.468761in}}%
\pgfpathquadraticcurveto{\pgfqpoint{9.854868in}{0.468761in}}{\pgfqpoint{9.847681in}{0.477875in}}%
\pgfpathquadraticcurveto{\pgfqpoint{9.840512in}{0.487007in}}{\pgfqpoint{9.840512in}{0.501867in}}%
\pgfpathquadraticcurveto{\pgfqpoint{9.840512in}{0.516727in}}{\pgfqpoint{9.847681in}{0.525841in}}%
\pgfpathquadraticcurveto{\pgfqpoint{9.854868in}{0.534956in}}{\pgfqpoint{9.866288in}{0.534956in}}%
\pgfpathquadraticcurveto{\pgfqpoint{9.873276in}{0.534956in}}{\pgfqpoint{9.878248in}{0.532218in}}%
\pgfpathquadraticcurveto{\pgfqpoint{9.883237in}{0.529498in}}{\pgfqpoint{9.886497in}{0.523878in}}%
\pgfpathlineto{\pgfqpoint{9.886497in}{0.523878in}}%
\pgfpathclose%
\pgfpathmoveto{\pgfqpoint{9.851212in}{0.501867in}}%
\pgfpathquadraticcurveto{\pgfqpoint{9.851212in}{0.490448in}}{\pgfqpoint{9.855913in}{0.483945in}}%
\pgfpathquadraticcurveto{\pgfqpoint{9.860614in}{0.477443in}}{\pgfqpoint{9.868827in}{0.477443in}}%
\pgfpathquadraticcurveto{\pgfqpoint{9.877041in}{0.477443in}}{\pgfqpoint{9.881760in}{0.483945in}}%
\pgfpathquadraticcurveto{\pgfqpoint{9.886497in}{0.490448in}}{\pgfqpoint{9.886497in}{0.501867in}}%
\pgfpathquadraticcurveto{\pgfqpoint{9.886497in}{0.513287in}}{\pgfqpoint{9.881760in}{0.519789in}}%
\pgfpathquadraticcurveto{\pgfqpoint{9.877041in}{0.526292in}}{\pgfqpoint{9.868827in}{0.526292in}}%
\pgfpathquadraticcurveto{\pgfqpoint{9.860614in}{0.526292in}}{\pgfqpoint{9.855913in}{0.519789in}}%
\pgfpathquadraticcurveto{\pgfqpoint{9.851212in}{0.513287in}}{\pgfqpoint{9.851212in}{0.501867in}}%
\pgfpathlineto{\pgfqpoint{9.851212in}{0.501867in}}%
\pgfpathclose%
\pgfpathmoveto{\pgfqpoint{9.919658in}{0.484702in}}%
\pgfpathlineto{\pgfqpoint{9.931546in}{0.484702in}}%
\pgfpathlineto{\pgfqpoint{9.931546in}{0.470400in}}%
\pgfpathlineto{\pgfqpoint{9.919658in}{0.470400in}}%
\pgfpathlineto{\pgfqpoint{9.919658in}{0.484702in}}%
\pgfpathlineto{\pgfqpoint{9.919658in}{0.484702in}}%
\pgfpathclose%
\pgfusepath{fill}%
\end{pgfscope}%
\begin{pgfscope}%
\pgfsetbuttcap%
\pgfsetroundjoin%
\definecolor{currentfill}{rgb}{0.309804,0.372549,0.419608}%
\pgfsetfillcolor{currentfill}%
\pgfsetlinewidth{0.000000pt}%
\definecolor{currentstroke}{rgb}{0.309804,0.372549,0.419608}%
\pgfsetstrokecolor{currentstroke}%
\pgfsetdash{}{0pt}%
\pgfpathmoveto{\pgfqpoint{10.052405in}{0.554445in}}%
\pgfpathlineto{\pgfqpoint{10.100696in}{0.554445in}}%
\pgfpathlineto{\pgfqpoint{10.100696in}{0.544862in}}%
\pgfpathlineto{\pgfqpoint{10.063771in}{0.544862in}}%
\pgfpathlineto{\pgfqpoint{10.063771in}{0.520096in}}%
\pgfpathlineto{\pgfqpoint{10.097094in}{0.520096in}}%
\pgfpathlineto{\pgfqpoint{10.097094in}{0.510531in}}%
\pgfpathlineto{\pgfqpoint{10.063771in}{0.510531in}}%
\pgfpathlineto{\pgfqpoint{10.063771in}{0.470400in}}%
\pgfpathlineto{\pgfqpoint{10.052405in}{0.470400in}}%
\pgfpathlineto{\pgfqpoint{10.052405in}{0.554445in}}%
\pgfpathlineto{\pgfqpoint{10.052405in}{0.554445in}}%
\pgfpathclose%
\pgfpathmoveto{\pgfqpoint{10.109864in}{0.533443in}}%
\pgfpathlineto{\pgfqpoint{10.120221in}{0.533443in}}%
\pgfpathlineto{\pgfqpoint{10.120221in}{0.470400in}}%
\pgfpathlineto{\pgfqpoint{10.109864in}{0.470400in}}%
\pgfpathlineto{\pgfqpoint{10.109864in}{0.533443in}}%
\pgfpathlineto{\pgfqpoint{10.109864in}{0.533443in}}%
\pgfpathclose%
\pgfpathmoveto{\pgfqpoint{10.109864in}{0.557993in}}%
\pgfpathlineto{\pgfqpoint{10.120221in}{0.557993in}}%
\pgfpathlineto{\pgfqpoint{10.120221in}{0.544862in}}%
\pgfpathlineto{\pgfqpoint{10.109864in}{0.544862in}}%
\pgfpathlineto{\pgfqpoint{10.109864in}{0.557993in}}%
\pgfpathlineto{\pgfqpoint{10.109864in}{0.557993in}}%
\pgfpathclose%
\pgfpathmoveto{\pgfqpoint{10.141890in}{0.557993in}}%
\pgfpathlineto{\pgfqpoint{10.152247in}{0.557993in}}%
\pgfpathlineto{\pgfqpoint{10.152247in}{0.470400in}}%
\pgfpathlineto{\pgfqpoint{10.141890in}{0.470400in}}%
\pgfpathlineto{\pgfqpoint{10.141890in}{0.557993in}}%
\pgfpathlineto{\pgfqpoint{10.141890in}{0.557993in}}%
\pgfpathclose%
\pgfpathmoveto{\pgfqpoint{10.173915in}{0.557993in}}%
\pgfpathlineto{\pgfqpoint{10.184272in}{0.557993in}}%
\pgfpathlineto{\pgfqpoint{10.184272in}{0.470400in}}%
\pgfpathlineto{\pgfqpoint{10.173915in}{0.470400in}}%
\pgfpathlineto{\pgfqpoint{10.173915in}{0.557993in}}%
\pgfpathlineto{\pgfqpoint{10.173915in}{0.557993in}}%
\pgfpathclose%
\pgfusepath{fill}%
\end{pgfscope}%
\begin{pgfscope}%
\pgfsetbuttcap%
\pgfsetroundjoin%
\definecolor{currentfill}{rgb}{0.309804,0.372549,0.419608}%
\pgfsetfillcolor{currentfill}%
\pgfsetlinewidth{0.000000pt}%
\definecolor{currentstroke}{rgb}{0.309804,0.372549,0.419608}%
\pgfsetstrokecolor{currentstroke}%
\pgfsetdash{}{0pt}%
\pgfpathmoveto{\pgfqpoint{10.249767in}{0.495275in}}%
\pgfpathlineto{\pgfqpoint{10.249767in}{0.533443in}}%
\pgfpathlineto{\pgfqpoint{10.260124in}{0.533443in}}%
\pgfpathlineto{\pgfqpoint{10.260124in}{0.495671in}}%
\pgfpathquadraticcurveto{\pgfqpoint{10.260124in}{0.486719in}}{\pgfqpoint{10.263601in}{0.482234in}}%
\pgfpathquadraticcurveto{\pgfqpoint{10.267095in}{0.477767in}}{\pgfqpoint{10.274084in}{0.477767in}}%
\pgfpathquadraticcurveto{\pgfqpoint{10.282459in}{0.477767in}}{\pgfqpoint{10.287323in}{0.483117in}}%
\pgfpathquadraticcurveto{\pgfqpoint{10.292204in}{0.488466in}}{\pgfqpoint{10.292204in}{0.497706in}}%
\pgfpathlineto{\pgfqpoint{10.292204in}{0.533443in}}%
\pgfpathlineto{\pgfqpoint{10.302561in}{0.533443in}}%
\pgfpathlineto{\pgfqpoint{10.302561in}{0.470400in}}%
\pgfpathlineto{\pgfqpoint{10.292204in}{0.470400in}}%
\pgfpathlineto{\pgfqpoint{10.292204in}{0.480091in}}%
\pgfpathquadraticcurveto{\pgfqpoint{10.288439in}{0.474345in}}{\pgfqpoint{10.283450in}{0.471553in}}%
\pgfpathquadraticcurveto{\pgfqpoint{10.278479in}{0.468761in}}{\pgfqpoint{10.271886in}{0.468761in}}%
\pgfpathquadraticcurveto{\pgfqpoint{10.261025in}{0.468761in}}{\pgfqpoint{10.255387in}{0.475515in}}%
\pgfpathquadraticcurveto{\pgfqpoint{10.249767in}{0.482270in}}{\pgfqpoint{10.249767in}{0.495275in}}%
\pgfpathlineto{\pgfqpoint{10.249767in}{0.495275in}}%
\pgfpathclose%
\pgfpathmoveto{\pgfqpoint{10.275831in}{0.534956in}}%
\pgfpathlineto{\pgfqpoint{10.275831in}{0.534956in}}%
\pgfpathlineto{\pgfqpoint{10.275831in}{0.534956in}}%
\pgfpathclose%
\pgfpathmoveto{\pgfqpoint{10.364072in}{0.531587in}}%
\pgfpathlineto{\pgfqpoint{10.364072in}{0.521789in}}%
\pgfpathquadraticcurveto{\pgfqpoint{10.359696in}{0.524040in}}{\pgfqpoint{10.354958in}{0.525157in}}%
\pgfpathquadraticcurveto{\pgfqpoint{10.350239in}{0.526292in}}{\pgfqpoint{10.345160in}{0.526292in}}%
\pgfpathquadraticcurveto{\pgfqpoint{10.337451in}{0.526292in}}{\pgfqpoint{10.333596in}{0.523932in}}%
\pgfpathquadraticcurveto{\pgfqpoint{10.329741in}{0.521573in}}{\pgfqpoint{10.329741in}{0.516835in}}%
\pgfpathquadraticcurveto{\pgfqpoint{10.329741in}{0.513233in}}{\pgfqpoint{10.332497in}{0.511180in}}%
\pgfpathquadraticcurveto{\pgfqpoint{10.335253in}{0.509126in}}{\pgfqpoint{10.343593in}{0.507271in}}%
\pgfpathlineto{\pgfqpoint{10.347141in}{0.506478in}}%
\pgfpathquadraticcurveto{\pgfqpoint{10.358164in}{0.504119in}}{\pgfqpoint{10.362812in}{0.499814in}}%
\pgfpathquadraticcurveto{\pgfqpoint{10.367459in}{0.495509in}}{\pgfqpoint{10.367459in}{0.487800in}}%
\pgfpathquadraticcurveto{\pgfqpoint{10.367459in}{0.479010in}}{\pgfqpoint{10.360506in}{0.473876in}}%
\pgfpathquadraticcurveto{\pgfqpoint{10.353553in}{0.468761in}}{\pgfqpoint{10.341395in}{0.468761in}}%
\pgfpathquadraticcurveto{\pgfqpoint{10.336334in}{0.468761in}}{\pgfqpoint{10.330840in}{0.469752in}}%
\pgfpathquadraticcurveto{\pgfqpoint{10.325346in}{0.470742in}}{\pgfqpoint{10.319276in}{0.472706in}}%
\pgfpathlineto{\pgfqpoint{10.319276in}{0.483405in}}%
\pgfpathquadraticcurveto{\pgfqpoint{10.325022in}{0.480415in}}{\pgfqpoint{10.330588in}{0.478920in}}%
\pgfpathquadraticcurveto{\pgfqpoint{10.336154in}{0.477443in}}{\pgfqpoint{10.341629in}{0.477443in}}%
\pgfpathquadraticcurveto{\pgfqpoint{10.348942in}{0.477443in}}{\pgfqpoint{10.352869in}{0.479946in}}%
\pgfpathquadraticcurveto{\pgfqpoint{10.356814in}{0.482450in}}{\pgfqpoint{10.356814in}{0.487007in}}%
\pgfpathquadraticcurveto{\pgfqpoint{10.356814in}{0.491222in}}{\pgfqpoint{10.353968in}{0.493474in}}%
\pgfpathquadraticcurveto{\pgfqpoint{10.351140in}{0.495725in}}{\pgfqpoint{10.341503in}{0.497814in}}%
\pgfpathlineto{\pgfqpoint{10.337901in}{0.498661in}}%
\pgfpathquadraticcurveto{\pgfqpoint{10.328282in}{0.500678in}}{\pgfqpoint{10.323995in}{0.504875in}}%
\pgfpathquadraticcurveto{\pgfqpoint{10.319727in}{0.509072in}}{\pgfqpoint{10.319727in}{0.516385in}}%
\pgfpathquadraticcurveto{\pgfqpoint{10.319727in}{0.525283in}}{\pgfqpoint{10.326031in}{0.530110in}}%
\pgfpathquadraticcurveto{\pgfqpoint{10.332335in}{0.534956in}}{\pgfqpoint{10.343935in}{0.534956in}}%
\pgfpathquadraticcurveto{\pgfqpoint{10.349663in}{0.534956in}}{\pgfqpoint{10.354724in}{0.534109in}}%
\pgfpathquadraticcurveto{\pgfqpoint{10.359804in}{0.533280in}}{\pgfqpoint{10.364072in}{0.531587in}}%
\pgfpathlineto{\pgfqpoint{10.364072in}{0.531587in}}%
\pgfpathclose%
\pgfpathmoveto{\pgfqpoint{10.437868in}{0.504515in}}%
\pgfpathlineto{\pgfqpoint{10.437868in}{0.499454in}}%
\pgfpathlineto{\pgfqpoint{10.390244in}{0.499454in}}%
\pgfpathquadraticcurveto{\pgfqpoint{10.390929in}{0.488754in}}{\pgfqpoint{10.396692in}{0.483153in}}%
\pgfpathquadraticcurveto{\pgfqpoint{10.402456in}{0.477551in}}{\pgfqpoint{10.412759in}{0.477551in}}%
\pgfpathquadraticcurveto{\pgfqpoint{10.418721in}{0.477551in}}{\pgfqpoint{10.424323in}{0.479010in}}%
\pgfpathquadraticcurveto{\pgfqpoint{10.429925in}{0.480469in}}{\pgfqpoint{10.435455in}{0.483405in}}%
\pgfpathlineto{\pgfqpoint{10.435455in}{0.473606in}}%
\pgfpathquadraticcurveto{\pgfqpoint{10.429871in}{0.471247in}}{\pgfqpoint{10.424017in}{0.470004in}}%
\pgfpathquadraticcurveto{\pgfqpoint{10.418163in}{0.468761in}}{\pgfqpoint{10.412147in}{0.468761in}}%
\pgfpathquadraticcurveto{\pgfqpoint{10.397053in}{0.468761in}}{\pgfqpoint{10.388245in}{0.477533in}}%
\pgfpathquadraticcurveto{\pgfqpoint{10.379437in}{0.486323in}}{\pgfqpoint{10.379437in}{0.501309in}}%
\pgfpathquadraticcurveto{\pgfqpoint{10.379437in}{0.516781in}}{\pgfqpoint{10.387794in}{0.525859in}}%
\pgfpathquadraticcurveto{\pgfqpoint{10.396152in}{0.534956in}}{\pgfqpoint{10.410346in}{0.534956in}}%
\pgfpathquadraticcurveto{\pgfqpoint{10.423062in}{0.534956in}}{\pgfqpoint{10.430465in}{0.526760in}}%
\pgfpathquadraticcurveto{\pgfqpoint{10.437868in}{0.518583in}}{\pgfqpoint{10.437868in}{0.504515in}}%
\pgfpathlineto{\pgfqpoint{10.437868in}{0.504515in}}%
\pgfpathclose%
\pgfpathmoveto{\pgfqpoint{10.427511in}{0.507559in}}%
\pgfpathquadraticcurveto{\pgfqpoint{10.427403in}{0.516043in}}{\pgfqpoint{10.422756in}{0.521104in}}%
\pgfpathquadraticcurveto{\pgfqpoint{10.418109in}{0.526184in}}{\pgfqpoint{10.410454in}{0.526184in}}%
\pgfpathquadraticcurveto{\pgfqpoint{10.401790in}{0.526184in}}{\pgfqpoint{10.396584in}{0.521284in}}%
\pgfpathquadraticcurveto{\pgfqpoint{10.391379in}{0.516385in}}{\pgfqpoint{10.390586in}{0.507487in}}%
\pgfpathlineto{\pgfqpoint{10.427511in}{0.507559in}}%
\pgfpathlineto{\pgfqpoint{10.427511in}{0.507559in}}%
\pgfpathclose%
\pgfpathmoveto{\pgfqpoint{10.495057in}{0.531587in}}%
\pgfpathlineto{\pgfqpoint{10.495057in}{0.521789in}}%
\pgfpathquadraticcurveto{\pgfqpoint{10.490680in}{0.524040in}}{\pgfqpoint{10.485943in}{0.525157in}}%
\pgfpathquadraticcurveto{\pgfqpoint{10.481224in}{0.526292in}}{\pgfqpoint{10.476144in}{0.526292in}}%
\pgfpathquadraticcurveto{\pgfqpoint{10.468435in}{0.526292in}}{\pgfqpoint{10.464580in}{0.523932in}}%
\pgfpathquadraticcurveto{\pgfqpoint{10.460726in}{0.521573in}}{\pgfqpoint{10.460726in}{0.516835in}}%
\pgfpathquadraticcurveto{\pgfqpoint{10.460726in}{0.513233in}}{\pgfqpoint{10.463482in}{0.511180in}}%
\pgfpathquadraticcurveto{\pgfqpoint{10.466237in}{0.509126in}}{\pgfqpoint{10.474577in}{0.507271in}}%
\pgfpathlineto{\pgfqpoint{10.478125in}{0.506478in}}%
\pgfpathquadraticcurveto{\pgfqpoint{10.489149in}{0.504119in}}{\pgfqpoint{10.493796in}{0.499814in}}%
\pgfpathquadraticcurveto{\pgfqpoint{10.498443in}{0.495509in}}{\pgfqpoint{10.498443in}{0.487800in}}%
\pgfpathquadraticcurveto{\pgfqpoint{10.498443in}{0.479010in}}{\pgfqpoint{10.491490in}{0.473876in}}%
\pgfpathquadraticcurveto{\pgfqpoint{10.484538in}{0.468761in}}{\pgfqpoint{10.472380in}{0.468761in}}%
\pgfpathquadraticcurveto{\pgfqpoint{10.467318in}{0.468761in}}{\pgfqpoint{10.461824in}{0.469752in}}%
\pgfpathquadraticcurveto{\pgfqpoint{10.456331in}{0.470742in}}{\pgfqpoint{10.450261in}{0.472706in}}%
\pgfpathlineto{\pgfqpoint{10.450261in}{0.483405in}}%
\pgfpathquadraticcurveto{\pgfqpoint{10.456007in}{0.480415in}}{\pgfqpoint{10.461572in}{0.478920in}}%
\pgfpathquadraticcurveto{\pgfqpoint{10.467138in}{0.477443in}}{\pgfqpoint{10.472614in}{0.477443in}}%
\pgfpathquadraticcurveto{\pgfqpoint{10.479927in}{0.477443in}}{\pgfqpoint{10.483853in}{0.479946in}}%
\pgfpathquadraticcurveto{\pgfqpoint{10.487798in}{0.482450in}}{\pgfqpoint{10.487798in}{0.487007in}}%
\pgfpathquadraticcurveto{\pgfqpoint{10.487798in}{0.491222in}}{\pgfqpoint{10.484952in}{0.493474in}}%
\pgfpathquadraticcurveto{\pgfqpoint{10.482124in}{0.495725in}}{\pgfqpoint{10.472488in}{0.497814in}}%
\pgfpathlineto{\pgfqpoint{10.468885in}{0.498661in}}%
\pgfpathquadraticcurveto{\pgfqpoint{10.459267in}{0.500678in}}{\pgfqpoint{10.454980in}{0.504875in}}%
\pgfpathquadraticcurveto{\pgfqpoint{10.450711in}{0.509072in}}{\pgfqpoint{10.450711in}{0.516385in}}%
\pgfpathquadraticcurveto{\pgfqpoint{10.450711in}{0.525283in}}{\pgfqpoint{10.457015in}{0.530110in}}%
\pgfpathquadraticcurveto{\pgfqpoint{10.463319in}{0.534956in}}{\pgfqpoint{10.474919in}{0.534956in}}%
\pgfpathquadraticcurveto{\pgfqpoint{10.480647in}{0.534956in}}{\pgfqpoint{10.485709in}{0.534109in}}%
\pgfpathquadraticcurveto{\pgfqpoint{10.490788in}{0.533280in}}{\pgfqpoint{10.495057in}{0.531587in}}%
\pgfpathlineto{\pgfqpoint{10.495057in}{0.531587in}}%
\pgfpathclose%
\pgfusepath{fill}%
\end{pgfscope}%
\begin{pgfscope}%
\pgfsetbuttcap%
\pgfsetroundjoin%
\definecolor{currentfill}{rgb}{0.309804,0.372549,0.419608}%
\pgfsetfillcolor{currentfill}%
\pgfsetlinewidth{0.000000pt}%
\definecolor{currentstroke}{rgb}{0.309804,0.372549,0.419608}%
\pgfsetstrokecolor{currentstroke}%
\pgfsetdash{}{0pt}%
\pgfpathmoveto{\pgfqpoint{10.578157in}{0.479856in}}%
\pgfpathlineto{\pgfqpoint{10.578157in}{0.446426in}}%
\pgfpathlineto{\pgfqpoint{10.567746in}{0.446426in}}%
\pgfpathlineto{\pgfqpoint{10.567746in}{0.533443in}}%
\pgfpathlineto{\pgfqpoint{10.578157in}{0.533443in}}%
\pgfpathlineto{\pgfqpoint{10.578157in}{0.523878in}}%
\pgfpathquadraticcurveto{\pgfqpoint{10.581436in}{0.529498in}}{\pgfqpoint{10.586407in}{0.532218in}}%
\pgfpathquadraticcurveto{\pgfqpoint{10.591396in}{0.534956in}}{\pgfqpoint{10.598313in}{0.534956in}}%
\pgfpathquadraticcurveto{\pgfqpoint{10.609805in}{0.534956in}}{\pgfqpoint{10.616974in}{0.525841in}}%
\pgfpathquadraticcurveto{\pgfqpoint{10.624160in}{0.516727in}}{\pgfqpoint{10.624160in}{0.501867in}}%
\pgfpathquadraticcurveto{\pgfqpoint{10.624160in}{0.487007in}}{\pgfqpoint{10.616974in}{0.477875in}}%
\pgfpathquadraticcurveto{\pgfqpoint{10.609805in}{0.468761in}}{\pgfqpoint{10.598313in}{0.468761in}}%
\pgfpathquadraticcurveto{\pgfqpoint{10.591396in}{0.468761in}}{\pgfqpoint{10.586407in}{0.471499in}}%
\pgfpathquadraticcurveto{\pgfqpoint{10.581436in}{0.474237in}}{\pgfqpoint{10.578157in}{0.479856in}}%
\pgfpathlineto{\pgfqpoint{10.578157in}{0.479856in}}%
\pgfpathclose%
\pgfpathmoveto{\pgfqpoint{10.613407in}{0.501867in}}%
\pgfpathquadraticcurveto{\pgfqpoint{10.613407in}{0.513287in}}{\pgfqpoint{10.608706in}{0.519789in}}%
\pgfpathquadraticcurveto{\pgfqpoint{10.604005in}{0.526292in}}{\pgfqpoint{10.595791in}{0.526292in}}%
\pgfpathquadraticcurveto{\pgfqpoint{10.587560in}{0.526292in}}{\pgfqpoint{10.582859in}{0.519789in}}%
\pgfpathquadraticcurveto{\pgfqpoint{10.578157in}{0.513287in}}{\pgfqpoint{10.578157in}{0.501867in}}%
\pgfpathquadraticcurveto{\pgfqpoint{10.578157in}{0.490448in}}{\pgfqpoint{10.582859in}{0.483945in}}%
\pgfpathquadraticcurveto{\pgfqpoint{10.587560in}{0.477443in}}{\pgfqpoint{10.595791in}{0.477443in}}%
\pgfpathquadraticcurveto{\pgfqpoint{10.604005in}{0.477443in}}{\pgfqpoint{10.608706in}{0.483945in}}%
\pgfpathquadraticcurveto{\pgfqpoint{10.613407in}{0.490448in}}{\pgfqpoint{10.613407in}{0.501867in}}%
\pgfpathlineto{\pgfqpoint{10.613407in}{0.501867in}}%
\pgfpathclose%
\pgfpathmoveto{\pgfqpoint{10.669983in}{0.502083in}}%
\pgfpathquadraticcurveto{\pgfqpoint{10.657429in}{0.502083in}}{\pgfqpoint{10.652584in}{0.499219in}}%
\pgfpathquadraticcurveto{\pgfqpoint{10.647738in}{0.496356in}}{\pgfqpoint{10.647738in}{0.489421in}}%
\pgfpathquadraticcurveto{\pgfqpoint{10.647738in}{0.483909in}}{\pgfqpoint{10.651377in}{0.480667in}}%
\pgfpathquadraticcurveto{\pgfqpoint{10.655015in}{0.477443in}}{\pgfqpoint{10.661247in}{0.477443in}}%
\pgfpathquadraticcurveto{\pgfqpoint{10.669875in}{0.477443in}}{\pgfqpoint{10.675081in}{0.483549in}}%
\pgfpathquadraticcurveto{\pgfqpoint{10.680286in}{0.489655in}}{\pgfqpoint{10.680286in}{0.499778in}}%
\pgfpathlineto{\pgfqpoint{10.680286in}{0.502083in}}%
\pgfpathlineto{\pgfqpoint{10.669983in}{0.502083in}}%
\pgfpathlineto{\pgfqpoint{10.669983in}{0.502083in}}%
\pgfpathclose%
\pgfpathmoveto{\pgfqpoint{10.690643in}{0.506370in}}%
\pgfpathlineto{\pgfqpoint{10.690643in}{0.470400in}}%
\pgfpathlineto{\pgfqpoint{10.680286in}{0.470400in}}%
\pgfpathlineto{\pgfqpoint{10.680286in}{0.479964in}}%
\pgfpathquadraticcurveto{\pgfqpoint{10.676738in}{0.474237in}}{\pgfqpoint{10.671442in}{0.471499in}}%
\pgfpathquadraticcurveto{\pgfqpoint{10.666147in}{0.468761in}}{\pgfqpoint{10.658492in}{0.468761in}}%
\pgfpathquadraticcurveto{\pgfqpoint{10.648819in}{0.468761in}}{\pgfqpoint{10.643091in}{0.474201in}}%
\pgfpathquadraticcurveto{\pgfqpoint{10.637381in}{0.479640in}}{\pgfqpoint{10.637381in}{0.488754in}}%
\pgfpathquadraticcurveto{\pgfqpoint{10.637381in}{0.499382in}}{\pgfqpoint{10.644496in}{0.504785in}}%
\pgfpathquadraticcurveto{\pgfqpoint{10.651629in}{0.510189in}}{\pgfqpoint{10.665750in}{0.510189in}}%
\pgfpathlineto{\pgfqpoint{10.680286in}{0.510189in}}%
\pgfpathlineto{\pgfqpoint{10.680286in}{0.511216in}}%
\pgfpathquadraticcurveto{\pgfqpoint{10.680286in}{0.518366in}}{\pgfqpoint{10.675585in}{0.522275in}}%
\pgfpathquadraticcurveto{\pgfqpoint{10.670884in}{0.526184in}}{\pgfqpoint{10.662382in}{0.526184in}}%
\pgfpathquadraticcurveto{\pgfqpoint{10.656979in}{0.526184in}}{\pgfqpoint{10.651845in}{0.524887in}}%
\pgfpathquadraticcurveto{\pgfqpoint{10.646730in}{0.523590in}}{\pgfqpoint{10.642010in}{0.520996in}}%
\pgfpathlineto{\pgfqpoint{10.642010in}{0.530579in}}%
\pgfpathquadraticcurveto{\pgfqpoint{10.647684in}{0.532776in}}{\pgfqpoint{10.653034in}{0.533857in}}%
\pgfpathquadraticcurveto{\pgfqpoint{10.658383in}{0.534956in}}{\pgfqpoint{10.663445in}{0.534956in}}%
\pgfpathquadraticcurveto{\pgfqpoint{10.677134in}{0.534956in}}{\pgfqpoint{10.683889in}{0.527859in}}%
\pgfpathquadraticcurveto{\pgfqpoint{10.690643in}{0.520780in}}{\pgfqpoint{10.690643in}{0.506370in}}%
\pgfpathlineto{\pgfqpoint{10.690643in}{0.506370in}}%
\pgfpathclose%
\pgfpathmoveto{\pgfqpoint{10.711970in}{0.533443in}}%
\pgfpathlineto{\pgfqpoint{10.722327in}{0.533443in}}%
\pgfpathlineto{\pgfqpoint{10.722327in}{0.470400in}}%
\pgfpathlineto{\pgfqpoint{10.711970in}{0.470400in}}%
\pgfpathlineto{\pgfqpoint{10.711970in}{0.533443in}}%
\pgfpathlineto{\pgfqpoint{10.711970in}{0.533443in}}%
\pgfpathclose%
\pgfpathmoveto{\pgfqpoint{10.711970in}{0.557993in}}%
\pgfpathlineto{\pgfqpoint{10.722327in}{0.557993in}}%
\pgfpathlineto{\pgfqpoint{10.722327in}{0.544862in}}%
\pgfpathlineto{\pgfqpoint{10.711970in}{0.544862in}}%
\pgfpathlineto{\pgfqpoint{10.711970in}{0.557993in}}%
\pgfpathlineto{\pgfqpoint{10.711970in}{0.557993in}}%
\pgfpathclose%
\pgfpathmoveto{\pgfqpoint{10.780524in}{0.523770in}}%
\pgfpathquadraticcurveto{\pgfqpoint{10.778777in}{0.524779in}}{\pgfqpoint{10.776723in}{0.525247in}}%
\pgfpathquadraticcurveto{\pgfqpoint{10.774670in}{0.525733in}}{\pgfqpoint{10.772202in}{0.525733in}}%
\pgfpathquadraticcurveto{\pgfqpoint{10.763412in}{0.525733in}}{\pgfqpoint{10.758711in}{0.520023in}}%
\pgfpathquadraticcurveto{\pgfqpoint{10.754010in}{0.514314in}}{\pgfqpoint{10.754010in}{0.503614in}}%
\pgfpathlineto{\pgfqpoint{10.754010in}{0.470400in}}%
\pgfpathlineto{\pgfqpoint{10.743599in}{0.470400in}}%
\pgfpathlineto{\pgfqpoint{10.743599in}{0.533443in}}%
\pgfpathlineto{\pgfqpoint{10.754010in}{0.533443in}}%
\pgfpathlineto{\pgfqpoint{10.754010in}{0.523644in}}%
\pgfpathquadraticcurveto{\pgfqpoint{10.757288in}{0.529390in}}{\pgfqpoint{10.762512in}{0.532164in}}%
\pgfpathquadraticcurveto{\pgfqpoint{10.767753in}{0.534956in}}{\pgfqpoint{10.775246in}{0.534956in}}%
\pgfpathquadraticcurveto{\pgfqpoint{10.776309in}{0.534956in}}{\pgfqpoint{10.777606in}{0.534811in}}%
\pgfpathquadraticcurveto{\pgfqpoint{10.778903in}{0.534685in}}{\pgfqpoint{10.780470in}{0.534397in}}%
\pgfpathlineto{\pgfqpoint{10.780524in}{0.523770in}}%
\pgfpathlineto{\pgfqpoint{10.780524in}{0.523770in}}%
\pgfpathclose%
\pgfpathmoveto{\pgfqpoint{10.842774in}{0.504515in}}%
\pgfpathlineto{\pgfqpoint{10.842774in}{0.499454in}}%
\pgfpathlineto{\pgfqpoint{10.795150in}{0.499454in}}%
\pgfpathquadraticcurveto{\pgfqpoint{10.795834in}{0.488754in}}{\pgfqpoint{10.801598in}{0.483153in}}%
\pgfpathquadraticcurveto{\pgfqpoint{10.807362in}{0.477551in}}{\pgfqpoint{10.817665in}{0.477551in}}%
\pgfpathquadraticcurveto{\pgfqpoint{10.823627in}{0.477551in}}{\pgfqpoint{10.829229in}{0.479010in}}%
\pgfpathquadraticcurveto{\pgfqpoint{10.834831in}{0.480469in}}{\pgfqpoint{10.840360in}{0.483405in}}%
\pgfpathlineto{\pgfqpoint{10.840360in}{0.473606in}}%
\pgfpathquadraticcurveto{\pgfqpoint{10.834776in}{0.471247in}}{\pgfqpoint{10.828923in}{0.470004in}}%
\pgfpathquadraticcurveto{\pgfqpoint{10.823069in}{0.468761in}}{\pgfqpoint{10.817053in}{0.468761in}}%
\pgfpathquadraticcurveto{\pgfqpoint{10.801958in}{0.468761in}}{\pgfqpoint{10.793150in}{0.477533in}}%
\pgfpathquadraticcurveto{\pgfqpoint{10.784342in}{0.486323in}}{\pgfqpoint{10.784342in}{0.501309in}}%
\pgfpathquadraticcurveto{\pgfqpoint{10.784342in}{0.516781in}}{\pgfqpoint{10.792700in}{0.525859in}}%
\pgfpathquadraticcurveto{\pgfqpoint{10.801058in}{0.534956in}}{\pgfqpoint{10.815251in}{0.534956in}}%
\pgfpathquadraticcurveto{\pgfqpoint{10.827968in}{0.534956in}}{\pgfqpoint{10.835371in}{0.526760in}}%
\pgfpathquadraticcurveto{\pgfqpoint{10.842774in}{0.518583in}}{\pgfqpoint{10.842774in}{0.504515in}}%
\pgfpathlineto{\pgfqpoint{10.842774in}{0.504515in}}%
\pgfpathclose%
\pgfpathmoveto{\pgfqpoint{10.832417in}{0.507559in}}%
\pgfpathquadraticcurveto{\pgfqpoint{10.832309in}{0.516043in}}{\pgfqpoint{10.827662in}{0.521104in}}%
\pgfpathquadraticcurveto{\pgfqpoint{10.823015in}{0.526184in}}{\pgfqpoint{10.815359in}{0.526184in}}%
\pgfpathquadraticcurveto{\pgfqpoint{10.806696in}{0.526184in}}{\pgfqpoint{10.801490in}{0.521284in}}%
\pgfpathquadraticcurveto{\pgfqpoint{10.796285in}{0.516385in}}{\pgfqpoint{10.795492in}{0.507487in}}%
\pgfpathlineto{\pgfqpoint{10.832417in}{0.507559in}}%
\pgfpathlineto{\pgfqpoint{10.832417in}{0.507559in}}%
\pgfpathclose%
\pgfpathmoveto{\pgfqpoint{10.901259in}{0.523878in}}%
\pgfpathlineto{\pgfqpoint{10.901259in}{0.557993in}}%
\pgfpathlineto{\pgfqpoint{10.911616in}{0.557993in}}%
\pgfpathlineto{\pgfqpoint{10.911616in}{0.470400in}}%
\pgfpathlineto{\pgfqpoint{10.901259in}{0.470400in}}%
\pgfpathlineto{\pgfqpoint{10.901259in}{0.479856in}}%
\pgfpathquadraticcurveto{\pgfqpoint{10.897999in}{0.474237in}}{\pgfqpoint{10.893010in}{0.471499in}}%
\pgfpathquadraticcurveto{\pgfqpoint{10.888038in}{0.468761in}}{\pgfqpoint{10.881050in}{0.468761in}}%
\pgfpathquadraticcurveto{\pgfqpoint{10.869630in}{0.468761in}}{\pgfqpoint{10.862443in}{0.477875in}}%
\pgfpathquadraticcurveto{\pgfqpoint{10.855274in}{0.487007in}}{\pgfqpoint{10.855274in}{0.501867in}}%
\pgfpathquadraticcurveto{\pgfqpoint{10.855274in}{0.516727in}}{\pgfqpoint{10.862443in}{0.525841in}}%
\pgfpathquadraticcurveto{\pgfqpoint{10.869630in}{0.534956in}}{\pgfqpoint{10.881050in}{0.534956in}}%
\pgfpathquadraticcurveto{\pgfqpoint{10.888038in}{0.534956in}}{\pgfqpoint{10.893010in}{0.532218in}}%
\pgfpathquadraticcurveto{\pgfqpoint{10.897999in}{0.529498in}}{\pgfqpoint{10.901259in}{0.523878in}}%
\pgfpathlineto{\pgfqpoint{10.901259in}{0.523878in}}%
\pgfpathclose%
\pgfpathmoveto{\pgfqpoint{10.865974in}{0.501867in}}%
\pgfpathquadraticcurveto{\pgfqpoint{10.865974in}{0.490448in}}{\pgfqpoint{10.870675in}{0.483945in}}%
\pgfpathquadraticcurveto{\pgfqpoint{10.875376in}{0.477443in}}{\pgfqpoint{10.883589in}{0.477443in}}%
\pgfpathquadraticcurveto{\pgfqpoint{10.891803in}{0.477443in}}{\pgfqpoint{10.896522in}{0.483945in}}%
\pgfpathquadraticcurveto{\pgfqpoint{10.901259in}{0.490448in}}{\pgfqpoint{10.901259in}{0.501867in}}%
\pgfpathquadraticcurveto{\pgfqpoint{10.901259in}{0.513287in}}{\pgfqpoint{10.896522in}{0.519789in}}%
\pgfpathquadraticcurveto{\pgfqpoint{10.891803in}{0.526292in}}{\pgfqpoint{10.883589in}{0.526292in}}%
\pgfpathquadraticcurveto{\pgfqpoint{10.875376in}{0.526292in}}{\pgfqpoint{10.870675in}{0.519789in}}%
\pgfpathquadraticcurveto{\pgfqpoint{10.865974in}{0.513287in}}{\pgfqpoint{10.865974in}{0.501867in}}%
\pgfpathlineto{\pgfqpoint{10.865974in}{0.501867in}}%
\pgfpathclose%
\pgfpathmoveto{\pgfqpoint{10.927737in}{0.506586in}}%
\pgfpathlineto{\pgfqpoint{10.958070in}{0.506586in}}%
\pgfpathlineto{\pgfqpoint{10.958070in}{0.497364in}}%
\pgfpathlineto{\pgfqpoint{10.927737in}{0.497364in}}%
\pgfpathlineto{\pgfqpoint{10.927737in}{0.506586in}}%
\pgfpathlineto{\pgfqpoint{10.927737in}{0.506586in}}%
\pgfpathclose%
\pgfpathmoveto{\pgfqpoint{11.019815in}{0.501867in}}%
\pgfpathquadraticcurveto{\pgfqpoint{11.019815in}{0.513287in}}{\pgfqpoint{11.015114in}{0.519789in}}%
\pgfpathquadraticcurveto{\pgfqpoint{11.010413in}{0.526292in}}{\pgfqpoint{11.002199in}{0.526292in}}%
\pgfpathquadraticcurveto{\pgfqpoint{10.993968in}{0.526292in}}{\pgfqpoint{10.989267in}{0.519789in}}%
\pgfpathquadraticcurveto{\pgfqpoint{10.984566in}{0.513287in}}{\pgfqpoint{10.984566in}{0.501867in}}%
\pgfpathquadraticcurveto{\pgfqpoint{10.984566in}{0.490448in}}{\pgfqpoint{10.989267in}{0.483945in}}%
\pgfpathquadraticcurveto{\pgfqpoint{10.993968in}{0.477443in}}{\pgfqpoint{11.002199in}{0.477443in}}%
\pgfpathquadraticcurveto{\pgfqpoint{11.010413in}{0.477443in}}{\pgfqpoint{11.015114in}{0.483945in}}%
\pgfpathquadraticcurveto{\pgfqpoint{11.019815in}{0.490448in}}{\pgfqpoint{11.019815in}{0.501867in}}%
\pgfpathlineto{\pgfqpoint{11.019815in}{0.501867in}}%
\pgfpathclose%
\pgfpathmoveto{\pgfqpoint{10.984566in}{0.523878in}}%
\pgfpathquadraticcurveto{\pgfqpoint{10.987844in}{0.529498in}}{\pgfqpoint{10.992815in}{0.532218in}}%
\pgfpathquadraticcurveto{\pgfqpoint{10.997804in}{0.534956in}}{\pgfqpoint{11.004721in}{0.534956in}}%
\pgfpathquadraticcurveto{\pgfqpoint{11.016213in}{0.534956in}}{\pgfqpoint{11.023382in}{0.525841in}}%
\pgfpathquadraticcurveto{\pgfqpoint{11.030569in}{0.516727in}}{\pgfqpoint{11.030569in}{0.501867in}}%
\pgfpathquadraticcurveto{\pgfqpoint{11.030569in}{0.487007in}}{\pgfqpoint{11.023382in}{0.477875in}}%
\pgfpathquadraticcurveto{\pgfqpoint{11.016213in}{0.468761in}}{\pgfqpoint{11.004721in}{0.468761in}}%
\pgfpathquadraticcurveto{\pgfqpoint{10.997804in}{0.468761in}}{\pgfqpoint{10.992815in}{0.471499in}}%
\pgfpathquadraticcurveto{\pgfqpoint{10.987844in}{0.474237in}}{\pgfqpoint{10.984566in}{0.479856in}}%
\pgfpathlineto{\pgfqpoint{10.984566in}{0.470400in}}%
\pgfpathlineto{\pgfqpoint{10.974155in}{0.470400in}}%
\pgfpathlineto{\pgfqpoint{10.974155in}{0.557993in}}%
\pgfpathlineto{\pgfqpoint{10.984566in}{0.557993in}}%
\pgfpathlineto{\pgfqpoint{10.984566in}{0.523878in}}%
\pgfpathlineto{\pgfqpoint{10.984566in}{0.523878in}}%
\pgfpathclose%
\pgfpathmoveto{\pgfqpoint{11.072159in}{0.526184in}}%
\pgfpathquadraticcurveto{\pgfqpoint{11.063837in}{0.526184in}}{\pgfqpoint{11.058992in}{0.519681in}}%
\pgfpathquadraticcurveto{\pgfqpoint{11.054146in}{0.513179in}}{\pgfqpoint{11.054146in}{0.501867in}}%
\pgfpathquadraticcurveto{\pgfqpoint{11.054146in}{0.490556in}}{\pgfqpoint{11.058956in}{0.484053in}}%
\pgfpathquadraticcurveto{\pgfqpoint{11.063783in}{0.477551in}}{\pgfqpoint{11.072159in}{0.477551in}}%
\pgfpathquadraticcurveto{\pgfqpoint{11.080444in}{0.477551in}}{\pgfqpoint{11.085271in}{0.484071in}}%
\pgfpathquadraticcurveto{\pgfqpoint{11.090117in}{0.490610in}}{\pgfqpoint{11.090117in}{0.501867in}}%
\pgfpathquadraticcurveto{\pgfqpoint{11.090117in}{0.513071in}}{\pgfqpoint{11.085271in}{0.519627in}}%
\pgfpathquadraticcurveto{\pgfqpoint{11.080444in}{0.526184in}}{\pgfqpoint{11.072159in}{0.526184in}}%
\pgfpathlineto{\pgfqpoint{11.072159in}{0.526184in}}%
\pgfpathclose%
\pgfpathmoveto{\pgfqpoint{11.072159in}{0.534956in}}%
\pgfpathquadraticcurveto{\pgfqpoint{11.085668in}{0.534956in}}{\pgfqpoint{11.093377in}{0.526166in}}%
\pgfpathquadraticcurveto{\pgfqpoint{11.101104in}{0.517394in}}{\pgfqpoint{11.101104in}{0.501867in}}%
\pgfpathquadraticcurveto{\pgfqpoint{11.101104in}{0.486395in}}{\pgfqpoint{11.093377in}{0.477569in}}%
\pgfpathquadraticcurveto{\pgfqpoint{11.085668in}{0.468761in}}{\pgfqpoint{11.072159in}{0.468761in}}%
\pgfpathquadraticcurveto{\pgfqpoint{11.058595in}{0.468761in}}{\pgfqpoint{11.050904in}{0.477569in}}%
\pgfpathquadraticcurveto{\pgfqpoint{11.043231in}{0.486395in}}{\pgfqpoint{11.043231in}{0.501867in}}%
\pgfpathquadraticcurveto{\pgfqpoint{11.043231in}{0.517394in}}{\pgfqpoint{11.050904in}{0.526166in}}%
\pgfpathquadraticcurveto{\pgfqpoint{11.058595in}{0.534956in}}{\pgfqpoint{11.072159in}{0.534956in}}%
\pgfpathlineto{\pgfqpoint{11.072159in}{0.534956in}}%
\pgfpathclose%
\pgfpathmoveto{\pgfqpoint{11.142694in}{0.526184in}}%
\pgfpathquadraticcurveto{\pgfqpoint{11.134373in}{0.526184in}}{\pgfqpoint{11.129527in}{0.519681in}}%
\pgfpathquadraticcurveto{\pgfqpoint{11.124682in}{0.513179in}}{\pgfqpoint{11.124682in}{0.501867in}}%
\pgfpathquadraticcurveto{\pgfqpoint{11.124682in}{0.490556in}}{\pgfqpoint{11.129491in}{0.484053in}}%
\pgfpathquadraticcurveto{\pgfqpoint{11.134319in}{0.477551in}}{\pgfqpoint{11.142694in}{0.477551in}}%
\pgfpathquadraticcurveto{\pgfqpoint{11.150980in}{0.477551in}}{\pgfqpoint{11.155807in}{0.484071in}}%
\pgfpathquadraticcurveto{\pgfqpoint{11.160652in}{0.490610in}}{\pgfqpoint{11.160652in}{0.501867in}}%
\pgfpathquadraticcurveto{\pgfqpoint{11.160652in}{0.513071in}}{\pgfqpoint{11.155807in}{0.519627in}}%
\pgfpathquadraticcurveto{\pgfqpoint{11.150980in}{0.526184in}}{\pgfqpoint{11.142694in}{0.526184in}}%
\pgfpathlineto{\pgfqpoint{11.142694in}{0.526184in}}%
\pgfpathclose%
\pgfpathmoveto{\pgfqpoint{11.142694in}{0.534956in}}%
\pgfpathquadraticcurveto{\pgfqpoint{11.156203in}{0.534956in}}{\pgfqpoint{11.163913in}{0.526166in}}%
\pgfpathquadraticcurveto{\pgfqpoint{11.171640in}{0.517394in}}{\pgfqpoint{11.171640in}{0.501867in}}%
\pgfpathquadraticcurveto{\pgfqpoint{11.171640in}{0.486395in}}{\pgfqpoint{11.163913in}{0.477569in}}%
\pgfpathquadraticcurveto{\pgfqpoint{11.156203in}{0.468761in}}{\pgfqpoint{11.142694in}{0.468761in}}%
\pgfpathquadraticcurveto{\pgfqpoint{11.129131in}{0.468761in}}{\pgfqpoint{11.121440in}{0.477569in}}%
\pgfpathquadraticcurveto{\pgfqpoint{11.113767in}{0.486395in}}{\pgfqpoint{11.113767in}{0.501867in}}%
\pgfpathquadraticcurveto{\pgfqpoint{11.113767in}{0.517394in}}{\pgfqpoint{11.121440in}{0.526166in}}%
\pgfpathquadraticcurveto{\pgfqpoint{11.129131in}{0.534956in}}{\pgfqpoint{11.142694in}{0.534956in}}%
\pgfpathlineto{\pgfqpoint{11.142694in}{0.534956in}}%
\pgfpathclose%
\pgfpathmoveto{\pgfqpoint{11.199054in}{0.551347in}}%
\pgfpathlineto{\pgfqpoint{11.199054in}{0.533443in}}%
\pgfpathlineto{\pgfqpoint{11.220381in}{0.533443in}}%
\pgfpathlineto{\pgfqpoint{11.220381in}{0.525391in}}%
\pgfpathlineto{\pgfqpoint{11.199054in}{0.525391in}}%
\pgfpathlineto{\pgfqpoint{11.199054in}{0.491168in}}%
\pgfpathquadraticcurveto{\pgfqpoint{11.199054in}{0.483459in}}{\pgfqpoint{11.201162in}{0.481261in}}%
\pgfpathquadraticcurveto{\pgfqpoint{11.203269in}{0.479064in}}{\pgfqpoint{11.209753in}{0.479064in}}%
\pgfpathlineto{\pgfqpoint{11.220381in}{0.479064in}}%
\pgfpathlineto{\pgfqpoint{11.220381in}{0.470400in}}%
\pgfpathlineto{\pgfqpoint{11.209753in}{0.470400in}}%
\pgfpathquadraticcurveto{\pgfqpoint{11.197757in}{0.470400in}}{\pgfqpoint{11.193200in}{0.474867in}}%
\pgfpathquadraticcurveto{\pgfqpoint{11.188643in}{0.479352in}}{\pgfqpoint{11.188643in}{0.491168in}}%
\pgfpathlineto{\pgfqpoint{11.188643in}{0.525391in}}%
\pgfpathlineto{\pgfqpoint{11.181042in}{0.525391in}}%
\pgfpathlineto{\pgfqpoint{11.181042in}{0.533443in}}%
\pgfpathlineto{\pgfqpoint{11.188643in}{0.533443in}}%
\pgfpathlineto{\pgfqpoint{11.188643in}{0.551347in}}%
\pgfpathlineto{\pgfqpoint{11.199054in}{0.551347in}}%
\pgfpathlineto{\pgfqpoint{11.199054in}{0.551347in}}%
\pgfpathclose%
\pgfpathmoveto{\pgfqpoint{11.274183in}{0.531587in}}%
\pgfpathlineto{\pgfqpoint{11.274183in}{0.521789in}}%
\pgfpathquadraticcurveto{\pgfqpoint{11.269806in}{0.524040in}}{\pgfqpoint{11.265069in}{0.525157in}}%
\pgfpathquadraticcurveto{\pgfqpoint{11.260350in}{0.526292in}}{\pgfqpoint{11.255270in}{0.526292in}}%
\pgfpathquadraticcurveto{\pgfqpoint{11.247561in}{0.526292in}}{\pgfqpoint{11.243706in}{0.523932in}}%
\pgfpathquadraticcurveto{\pgfqpoint{11.239852in}{0.521573in}}{\pgfqpoint{11.239852in}{0.516835in}}%
\pgfpathquadraticcurveto{\pgfqpoint{11.239852in}{0.513233in}}{\pgfqpoint{11.242608in}{0.511180in}}%
\pgfpathquadraticcurveto{\pgfqpoint{11.245363in}{0.509126in}}{\pgfqpoint{11.253703in}{0.507271in}}%
\pgfpathlineto{\pgfqpoint{11.257252in}{0.506478in}}%
\pgfpathquadraticcurveto{\pgfqpoint{11.268275in}{0.504119in}}{\pgfqpoint{11.272922in}{0.499814in}}%
\pgfpathquadraticcurveto{\pgfqpoint{11.277569in}{0.495509in}}{\pgfqpoint{11.277569in}{0.487800in}}%
\pgfpathquadraticcurveto{\pgfqpoint{11.277569in}{0.479010in}}{\pgfqpoint{11.270617in}{0.473876in}}%
\pgfpathquadraticcurveto{\pgfqpoint{11.263664in}{0.468761in}}{\pgfqpoint{11.251506in}{0.468761in}}%
\pgfpathquadraticcurveto{\pgfqpoint{11.246444in}{0.468761in}}{\pgfqpoint{11.240951in}{0.469752in}}%
\pgfpathquadraticcurveto{\pgfqpoint{11.235457in}{0.470742in}}{\pgfqpoint{11.229387in}{0.472706in}}%
\pgfpathlineto{\pgfqpoint{11.229387in}{0.483405in}}%
\pgfpathquadraticcurveto{\pgfqpoint{11.235133in}{0.480415in}}{\pgfqpoint{11.240698in}{0.478920in}}%
\pgfpathquadraticcurveto{\pgfqpoint{11.246264in}{0.477443in}}{\pgfqpoint{11.251740in}{0.477443in}}%
\pgfpathquadraticcurveto{\pgfqpoint{11.259053in}{0.477443in}}{\pgfqpoint{11.262979in}{0.479946in}}%
\pgfpathquadraticcurveto{\pgfqpoint{11.266924in}{0.482450in}}{\pgfqpoint{11.266924in}{0.487007in}}%
\pgfpathquadraticcurveto{\pgfqpoint{11.266924in}{0.491222in}}{\pgfqpoint{11.264078in}{0.493474in}}%
\pgfpathquadraticcurveto{\pgfqpoint{11.261250in}{0.495725in}}{\pgfqpoint{11.251614in}{0.497814in}}%
\pgfpathlineto{\pgfqpoint{11.248011in}{0.498661in}}%
\pgfpathquadraticcurveto{\pgfqpoint{11.238393in}{0.500678in}}{\pgfqpoint{11.234106in}{0.504875in}}%
\pgfpathquadraticcurveto{\pgfqpoint{11.229837in}{0.509072in}}{\pgfqpoint{11.229837in}{0.516385in}}%
\pgfpathquadraticcurveto{\pgfqpoint{11.229837in}{0.525283in}}{\pgfqpoint{11.236141in}{0.530110in}}%
\pgfpathquadraticcurveto{\pgfqpoint{11.242446in}{0.534956in}}{\pgfqpoint{11.254045in}{0.534956in}}%
\pgfpathquadraticcurveto{\pgfqpoint{11.259773in}{0.534956in}}{\pgfqpoint{11.264835in}{0.534109in}}%
\pgfpathquadraticcurveto{\pgfqpoint{11.269914in}{0.533280in}}{\pgfqpoint{11.274183in}{0.531587in}}%
\pgfpathlineto{\pgfqpoint{11.274183in}{0.531587in}}%
\pgfpathclose%
\pgfpathmoveto{\pgfqpoint{11.304299in}{0.551347in}}%
\pgfpathlineto{\pgfqpoint{11.304299in}{0.533443in}}%
\pgfpathlineto{\pgfqpoint{11.325626in}{0.533443in}}%
\pgfpathlineto{\pgfqpoint{11.325626in}{0.525391in}}%
\pgfpathlineto{\pgfqpoint{11.304299in}{0.525391in}}%
\pgfpathlineto{\pgfqpoint{11.304299in}{0.491168in}}%
\pgfpathquadraticcurveto{\pgfqpoint{11.304299in}{0.483459in}}{\pgfqpoint{11.306407in}{0.481261in}}%
\pgfpathquadraticcurveto{\pgfqpoint{11.308514in}{0.479064in}}{\pgfqpoint{11.314998in}{0.479064in}}%
\pgfpathlineto{\pgfqpoint{11.325626in}{0.479064in}}%
\pgfpathlineto{\pgfqpoint{11.325626in}{0.470400in}}%
\pgfpathlineto{\pgfqpoint{11.314998in}{0.470400in}}%
\pgfpathquadraticcurveto{\pgfqpoint{11.303002in}{0.470400in}}{\pgfqpoint{11.298445in}{0.474867in}}%
\pgfpathquadraticcurveto{\pgfqpoint{11.293888in}{0.479352in}}{\pgfqpoint{11.293888in}{0.491168in}}%
\pgfpathlineto{\pgfqpoint{11.293888in}{0.525391in}}%
\pgfpathlineto{\pgfqpoint{11.286287in}{0.525391in}}%
\pgfpathlineto{\pgfqpoint{11.286287in}{0.533443in}}%
\pgfpathlineto{\pgfqpoint{11.293888in}{0.533443in}}%
\pgfpathlineto{\pgfqpoint{11.293888in}{0.551347in}}%
\pgfpathlineto{\pgfqpoint{11.304299in}{0.551347in}}%
\pgfpathlineto{\pgfqpoint{11.304299in}{0.551347in}}%
\pgfpathclose%
\pgfpathmoveto{\pgfqpoint{11.375771in}{0.523770in}}%
\pgfpathquadraticcurveto{\pgfqpoint{11.374024in}{0.524779in}}{\pgfqpoint{11.371971in}{0.525247in}}%
\pgfpathquadraticcurveto{\pgfqpoint{11.369918in}{0.525733in}}{\pgfqpoint{11.367450in}{0.525733in}}%
\pgfpathquadraticcurveto{\pgfqpoint{11.358660in}{0.525733in}}{\pgfqpoint{11.353959in}{0.520023in}}%
\pgfpathquadraticcurveto{\pgfqpoint{11.349258in}{0.514314in}}{\pgfqpoint{11.349258in}{0.503614in}}%
\pgfpathlineto{\pgfqpoint{11.349258in}{0.470400in}}%
\pgfpathlineto{\pgfqpoint{11.338847in}{0.470400in}}%
\pgfpathlineto{\pgfqpoint{11.338847in}{0.533443in}}%
\pgfpathlineto{\pgfqpoint{11.349258in}{0.533443in}}%
\pgfpathlineto{\pgfqpoint{11.349258in}{0.523644in}}%
\pgfpathquadraticcurveto{\pgfqpoint{11.352536in}{0.529390in}}{\pgfqpoint{11.357759in}{0.532164in}}%
\pgfpathquadraticcurveto{\pgfqpoint{11.363001in}{0.534956in}}{\pgfqpoint{11.370494in}{0.534956in}}%
\pgfpathquadraticcurveto{\pgfqpoint{11.371557in}{0.534956in}}{\pgfqpoint{11.372854in}{0.534811in}}%
\pgfpathquadraticcurveto{\pgfqpoint{11.374150in}{0.534685in}}{\pgfqpoint{11.375717in}{0.534397in}}%
\pgfpathlineto{\pgfqpoint{11.375771in}{0.523770in}}%
\pgfpathlineto{\pgfqpoint{11.375771in}{0.523770in}}%
\pgfpathclose%
\pgfpathmoveto{\pgfqpoint{11.415290in}{0.502083in}}%
\pgfpathquadraticcurveto{\pgfqpoint{11.402736in}{0.502083in}}{\pgfqpoint{11.397890in}{0.499219in}}%
\pgfpathquadraticcurveto{\pgfqpoint{11.393045in}{0.496356in}}{\pgfqpoint{11.393045in}{0.489421in}}%
\pgfpathquadraticcurveto{\pgfqpoint{11.393045in}{0.483909in}}{\pgfqpoint{11.396684in}{0.480667in}}%
\pgfpathquadraticcurveto{\pgfqpoint{11.400322in}{0.477443in}}{\pgfqpoint{11.406554in}{0.477443in}}%
\pgfpathquadraticcurveto{\pgfqpoint{11.415182in}{0.477443in}}{\pgfqpoint{11.420388in}{0.483549in}}%
\pgfpathquadraticcurveto{\pgfqpoint{11.425593in}{0.489655in}}{\pgfqpoint{11.425593in}{0.499778in}}%
\pgfpathlineto{\pgfqpoint{11.425593in}{0.502083in}}%
\pgfpathlineto{\pgfqpoint{11.415290in}{0.502083in}}%
\pgfpathlineto{\pgfqpoint{11.415290in}{0.502083in}}%
\pgfpathclose%
\pgfpathmoveto{\pgfqpoint{11.435950in}{0.506370in}}%
\pgfpathlineto{\pgfqpoint{11.435950in}{0.470400in}}%
\pgfpathlineto{\pgfqpoint{11.425593in}{0.470400in}}%
\pgfpathlineto{\pgfqpoint{11.425593in}{0.479964in}}%
\pgfpathquadraticcurveto{\pgfqpoint{11.422045in}{0.474237in}}{\pgfqpoint{11.416749in}{0.471499in}}%
\pgfpathquadraticcurveto{\pgfqpoint{11.411454in}{0.468761in}}{\pgfqpoint{11.403798in}{0.468761in}}%
\pgfpathquadraticcurveto{\pgfqpoint{11.394126in}{0.468761in}}{\pgfqpoint{11.388398in}{0.474201in}}%
\pgfpathquadraticcurveto{\pgfqpoint{11.382688in}{0.479640in}}{\pgfqpoint{11.382688in}{0.488754in}}%
\pgfpathquadraticcurveto{\pgfqpoint{11.382688in}{0.499382in}}{\pgfqpoint{11.389803in}{0.504785in}}%
\pgfpathquadraticcurveto{\pgfqpoint{11.396936in}{0.510189in}}{\pgfqpoint{11.411057in}{0.510189in}}%
\pgfpathlineto{\pgfqpoint{11.425593in}{0.510189in}}%
\pgfpathlineto{\pgfqpoint{11.425593in}{0.511216in}}%
\pgfpathquadraticcurveto{\pgfqpoint{11.425593in}{0.518366in}}{\pgfqpoint{11.420892in}{0.522275in}}%
\pgfpathquadraticcurveto{\pgfqpoint{11.416191in}{0.526184in}}{\pgfqpoint{11.407689in}{0.526184in}}%
\pgfpathquadraticcurveto{\pgfqpoint{11.402285in}{0.526184in}}{\pgfqpoint{11.397152in}{0.524887in}}%
\pgfpathquadraticcurveto{\pgfqpoint{11.392036in}{0.523590in}}{\pgfqpoint{11.387317in}{0.520996in}}%
\pgfpathlineto{\pgfqpoint{11.387317in}{0.530579in}}%
\pgfpathquadraticcurveto{\pgfqpoint{11.392991in}{0.532776in}}{\pgfqpoint{11.398341in}{0.533857in}}%
\pgfpathquadraticcurveto{\pgfqpoint{11.403690in}{0.534956in}}{\pgfqpoint{11.408752in}{0.534956in}}%
\pgfpathquadraticcurveto{\pgfqpoint{11.422441in}{0.534956in}}{\pgfqpoint{11.429196in}{0.527859in}}%
\pgfpathquadraticcurveto{\pgfqpoint{11.435950in}{0.520780in}}{\pgfqpoint{11.435950in}{0.506370in}}%
\pgfpathlineto{\pgfqpoint{11.435950in}{0.506370in}}%
\pgfpathclose%
\pgfpathmoveto{\pgfqpoint{11.467291in}{0.479856in}}%
\pgfpathlineto{\pgfqpoint{11.467291in}{0.446426in}}%
\pgfpathlineto{\pgfqpoint{11.456880in}{0.446426in}}%
\pgfpathlineto{\pgfqpoint{11.456880in}{0.533443in}}%
\pgfpathlineto{\pgfqpoint{11.467291in}{0.533443in}}%
\pgfpathlineto{\pgfqpoint{11.467291in}{0.523878in}}%
\pgfpathquadraticcurveto{\pgfqpoint{11.470569in}{0.529498in}}{\pgfqpoint{11.475541in}{0.532218in}}%
\pgfpathquadraticcurveto{\pgfqpoint{11.480530in}{0.534956in}}{\pgfqpoint{11.487447in}{0.534956in}}%
\pgfpathquadraticcurveto{\pgfqpoint{11.498939in}{0.534956in}}{\pgfqpoint{11.506107in}{0.525841in}}%
\pgfpathquadraticcurveto{\pgfqpoint{11.513294in}{0.516727in}}{\pgfqpoint{11.513294in}{0.501867in}}%
\pgfpathquadraticcurveto{\pgfqpoint{11.513294in}{0.487007in}}{\pgfqpoint{11.506107in}{0.477875in}}%
\pgfpathquadraticcurveto{\pgfqpoint{11.498939in}{0.468761in}}{\pgfqpoint{11.487447in}{0.468761in}}%
\pgfpathquadraticcurveto{\pgfqpoint{11.480530in}{0.468761in}}{\pgfqpoint{11.475541in}{0.471499in}}%
\pgfpathquadraticcurveto{\pgfqpoint{11.470569in}{0.474237in}}{\pgfqpoint{11.467291in}{0.479856in}}%
\pgfpathlineto{\pgfqpoint{11.467291in}{0.479856in}}%
\pgfpathclose%
\pgfpathmoveto{\pgfqpoint{11.502541in}{0.501867in}}%
\pgfpathquadraticcurveto{\pgfqpoint{11.502541in}{0.513287in}}{\pgfqpoint{11.497840in}{0.519789in}}%
\pgfpathquadraticcurveto{\pgfqpoint{11.493139in}{0.526292in}}{\pgfqpoint{11.484925in}{0.526292in}}%
\pgfpathquadraticcurveto{\pgfqpoint{11.476694in}{0.526292in}}{\pgfqpoint{11.471992in}{0.519789in}}%
\pgfpathquadraticcurveto{\pgfqpoint{11.467291in}{0.513287in}}{\pgfqpoint{11.467291in}{0.501867in}}%
\pgfpathquadraticcurveto{\pgfqpoint{11.467291in}{0.490448in}}{\pgfqpoint{11.471992in}{0.483945in}}%
\pgfpathquadraticcurveto{\pgfqpoint{11.476694in}{0.477443in}}{\pgfqpoint{11.484925in}{0.477443in}}%
\pgfpathquadraticcurveto{\pgfqpoint{11.493139in}{0.477443in}}{\pgfqpoint{11.497840in}{0.483945in}}%
\pgfpathquadraticcurveto{\pgfqpoint{11.502541in}{0.490448in}}{\pgfqpoint{11.502541in}{0.501867in}}%
\pgfpathlineto{\pgfqpoint{11.502541in}{0.501867in}}%
\pgfpathclose%
\pgfusepath{fill}%
\end{pgfscope}%
\begin{pgfscope}%
\pgfsetbuttcap%
\pgfsetroundjoin%
\definecolor{currentfill}{rgb}{0.309804,0.372549,0.419608}%
\pgfsetfillcolor{currentfill}%
\pgfsetlinewidth{0.000000pt}%
\definecolor{currentstroke}{rgb}{0.309804,0.372549,0.419608}%
\pgfsetstrokecolor{currentstroke}%
\pgfsetdash{}{0pt}%
\pgfpathmoveto{\pgfqpoint{11.665370in}{0.508460in}}%
\pgfpathlineto{\pgfqpoint{11.665370in}{0.470400in}}%
\pgfpathlineto{\pgfqpoint{11.655013in}{0.470400in}}%
\pgfpathlineto{\pgfqpoint{11.655013in}{0.508117in}}%
\pgfpathquadraticcurveto{\pgfqpoint{11.655013in}{0.517069in}}{\pgfqpoint{11.651519in}{0.521500in}}%
\pgfpathquadraticcurveto{\pgfqpoint{11.648025in}{0.525949in}}{\pgfqpoint{11.641054in}{0.525949in}}%
\pgfpathquadraticcurveto{\pgfqpoint{11.632660in}{0.525949in}}{\pgfqpoint{11.627815in}{0.520600in}}%
\pgfpathquadraticcurveto{\pgfqpoint{11.622970in}{0.515268in}}{\pgfqpoint{11.622970in}{0.506028in}}%
\pgfpathlineto{\pgfqpoint{11.622970in}{0.470400in}}%
\pgfpathlineto{\pgfqpoint{11.612559in}{0.470400in}}%
\pgfpathlineto{\pgfqpoint{11.612559in}{0.533443in}}%
\pgfpathlineto{\pgfqpoint{11.622970in}{0.533443in}}%
\pgfpathlineto{\pgfqpoint{11.622970in}{0.523644in}}%
\pgfpathquadraticcurveto{\pgfqpoint{11.626698in}{0.529336in}}{\pgfqpoint{11.631724in}{0.532146in}}%
\pgfpathquadraticcurveto{\pgfqpoint{11.636767in}{0.534956in}}{\pgfqpoint{11.643360in}{0.534956in}}%
\pgfpathquadraticcurveto{\pgfqpoint{11.654221in}{0.534956in}}{\pgfqpoint{11.659787in}{0.528237in}}%
\pgfpathquadraticcurveto{\pgfqpoint{11.665370in}{0.521518in}}{\pgfqpoint{11.665370in}{0.508460in}}%
\pgfpathlineto{\pgfqpoint{11.665370in}{0.508460in}}%
\pgfpathclose%
\pgfpathmoveto{\pgfqpoint{11.710437in}{0.526184in}}%
\pgfpathquadraticcurveto{\pgfqpoint{11.702115in}{0.526184in}}{\pgfqpoint{11.697270in}{0.519681in}}%
\pgfpathquadraticcurveto{\pgfqpoint{11.692425in}{0.513179in}}{\pgfqpoint{11.692425in}{0.501867in}}%
\pgfpathquadraticcurveto{\pgfqpoint{11.692425in}{0.490556in}}{\pgfqpoint{11.697234in}{0.484053in}}%
\pgfpathquadraticcurveto{\pgfqpoint{11.702061in}{0.477551in}}{\pgfqpoint{11.710437in}{0.477551in}}%
\pgfpathquadraticcurveto{\pgfqpoint{11.718722in}{0.477551in}}{\pgfqpoint{11.723550in}{0.484071in}}%
\pgfpathquadraticcurveto{\pgfqpoint{11.728395in}{0.490610in}}{\pgfqpoint{11.728395in}{0.501867in}}%
\pgfpathquadraticcurveto{\pgfqpoint{11.728395in}{0.513071in}}{\pgfqpoint{11.723550in}{0.519627in}}%
\pgfpathquadraticcurveto{\pgfqpoint{11.718722in}{0.526184in}}{\pgfqpoint{11.710437in}{0.526184in}}%
\pgfpathlineto{\pgfqpoint{11.710437in}{0.526184in}}%
\pgfpathclose%
\pgfpathmoveto{\pgfqpoint{11.710437in}{0.534956in}}%
\pgfpathquadraticcurveto{\pgfqpoint{11.723946in}{0.534956in}}{\pgfqpoint{11.731655in}{0.526166in}}%
\pgfpathquadraticcurveto{\pgfqpoint{11.739382in}{0.517394in}}{\pgfqpoint{11.739382in}{0.501867in}}%
\pgfpathquadraticcurveto{\pgfqpoint{11.739382in}{0.486395in}}{\pgfqpoint{11.731655in}{0.477569in}}%
\pgfpathquadraticcurveto{\pgfqpoint{11.723946in}{0.468761in}}{\pgfqpoint{11.710437in}{0.468761in}}%
\pgfpathquadraticcurveto{\pgfqpoint{11.696874in}{0.468761in}}{\pgfqpoint{11.689183in}{0.477569in}}%
\pgfpathquadraticcurveto{\pgfqpoint{11.681509in}{0.486395in}}{\pgfqpoint{11.681509in}{0.501867in}}%
\pgfpathquadraticcurveto{\pgfqpoint{11.681509in}{0.517394in}}{\pgfqpoint{11.689183in}{0.526166in}}%
\pgfpathquadraticcurveto{\pgfqpoint{11.696874in}{0.534956in}}{\pgfqpoint{11.710437in}{0.534956in}}%
\pgfpathlineto{\pgfqpoint{11.710437in}{0.534956in}}%
\pgfpathclose%
\pgfpathmoveto{\pgfqpoint{11.793077in}{0.523770in}}%
\pgfpathquadraticcurveto{\pgfqpoint{11.791329in}{0.524779in}}{\pgfqpoint{11.789276in}{0.525247in}}%
\pgfpathquadraticcurveto{\pgfqpoint{11.787223in}{0.525733in}}{\pgfqpoint{11.784755in}{0.525733in}}%
\pgfpathquadraticcurveto{\pgfqpoint{11.775965in}{0.525733in}}{\pgfqpoint{11.771264in}{0.520023in}}%
\pgfpathquadraticcurveto{\pgfqpoint{11.766563in}{0.514314in}}{\pgfqpoint{11.766563in}{0.503614in}}%
\pgfpathlineto{\pgfqpoint{11.766563in}{0.470400in}}%
\pgfpathlineto{\pgfqpoint{11.756152in}{0.470400in}}%
\pgfpathlineto{\pgfqpoint{11.756152in}{0.533443in}}%
\pgfpathlineto{\pgfqpoint{11.766563in}{0.533443in}}%
\pgfpathlineto{\pgfqpoint{11.766563in}{0.523644in}}%
\pgfpathquadraticcurveto{\pgfqpoint{11.769841in}{0.529390in}}{\pgfqpoint{11.775064in}{0.532164in}}%
\pgfpathquadraticcurveto{\pgfqpoint{11.780306in}{0.534956in}}{\pgfqpoint{11.787799in}{0.534956in}}%
\pgfpathquadraticcurveto{\pgfqpoint{11.788862in}{0.534956in}}{\pgfqpoint{11.790159in}{0.534811in}}%
\pgfpathquadraticcurveto{\pgfqpoint{11.791456in}{0.534685in}}{\pgfqpoint{11.793023in}{0.534397in}}%
\pgfpathlineto{\pgfqpoint{11.793077in}{0.523770in}}%
\pgfpathlineto{\pgfqpoint{11.793077in}{0.523770in}}%
\pgfpathclose%
\pgfpathmoveto{\pgfqpoint{11.851004in}{0.521338in}}%
\pgfpathquadraticcurveto{\pgfqpoint{11.854894in}{0.528327in}}{\pgfqpoint{11.860298in}{0.531641in}}%
\pgfpathquadraticcurveto{\pgfqpoint{11.865702in}{0.534956in}}{\pgfqpoint{11.873015in}{0.534956in}}%
\pgfpathquadraticcurveto{\pgfqpoint{11.882867in}{0.534956in}}{\pgfqpoint{11.888217in}{0.528057in}}%
\pgfpathquadraticcurveto{\pgfqpoint{11.893566in}{0.521176in}}{\pgfqpoint{11.893566in}{0.508460in}}%
\pgfpathlineto{\pgfqpoint{11.893566in}{0.470400in}}%
\pgfpathlineto{\pgfqpoint{11.883155in}{0.470400in}}%
\pgfpathlineto{\pgfqpoint{11.883155in}{0.508117in}}%
\pgfpathquadraticcurveto{\pgfqpoint{11.883155in}{0.517178in}}{\pgfqpoint{11.879931in}{0.521555in}}%
\pgfpathquadraticcurveto{\pgfqpoint{11.876725in}{0.525949in}}{\pgfqpoint{11.870151in}{0.525949in}}%
\pgfpathquadraticcurveto{\pgfqpoint{11.862099in}{0.525949in}}{\pgfqpoint{11.857416in}{0.520600in}}%
\pgfpathquadraticcurveto{\pgfqpoint{11.852751in}{0.515268in}}{\pgfqpoint{11.852751in}{0.506028in}}%
\pgfpathlineto{\pgfqpoint{11.852751in}{0.470400in}}%
\pgfpathlineto{\pgfqpoint{11.842340in}{0.470400in}}%
\pgfpathlineto{\pgfqpoint{11.842340in}{0.508117in}}%
\pgfpathquadraticcurveto{\pgfqpoint{11.842340in}{0.517232in}}{\pgfqpoint{11.839134in}{0.521591in}}%
\pgfpathquadraticcurveto{\pgfqpoint{11.835928in}{0.525949in}}{\pgfqpoint{11.829227in}{0.525949in}}%
\pgfpathquadraticcurveto{\pgfqpoint{11.821284in}{0.525949in}}{\pgfqpoint{11.816600in}{0.520582in}}%
\pgfpathquadraticcurveto{\pgfqpoint{11.811935in}{0.515214in}}{\pgfqpoint{11.811935in}{0.506028in}}%
\pgfpathlineto{\pgfqpoint{11.811935in}{0.470400in}}%
\pgfpathlineto{\pgfqpoint{11.801524in}{0.470400in}}%
\pgfpathlineto{\pgfqpoint{11.801524in}{0.533443in}}%
\pgfpathlineto{\pgfqpoint{11.811935in}{0.533443in}}%
\pgfpathlineto{\pgfqpoint{11.811935in}{0.523644in}}%
\pgfpathquadraticcurveto{\pgfqpoint{11.815484in}{0.529444in}}{\pgfqpoint{11.820437in}{0.532200in}}%
\pgfpathquadraticcurveto{\pgfqpoint{11.825390in}{0.534956in}}{\pgfqpoint{11.832199in}{0.534956in}}%
\pgfpathquadraticcurveto{\pgfqpoint{11.839080in}{0.534956in}}{\pgfqpoint{11.843889in}{0.531461in}}%
\pgfpathquadraticcurveto{\pgfqpoint{11.848698in}{0.527985in}}{\pgfqpoint{11.851004in}{0.521338in}}%
\pgfpathlineto{\pgfqpoint{11.851004in}{0.521338in}}%
\pgfpathclose%
\pgfpathmoveto{\pgfqpoint{11.942866in}{0.502083in}}%
\pgfpathquadraticcurveto{\pgfqpoint{11.930311in}{0.502083in}}{\pgfqpoint{11.925466in}{0.499219in}}%
\pgfpathquadraticcurveto{\pgfqpoint{11.920621in}{0.496356in}}{\pgfqpoint{11.920621in}{0.489421in}}%
\pgfpathquadraticcurveto{\pgfqpoint{11.920621in}{0.483909in}}{\pgfqpoint{11.924259in}{0.480667in}}%
\pgfpathquadraticcurveto{\pgfqpoint{11.927898in}{0.477443in}}{\pgfqpoint{11.934130in}{0.477443in}}%
\pgfpathquadraticcurveto{\pgfqpoint{11.942758in}{0.477443in}}{\pgfqpoint{11.947963in}{0.483549in}}%
\pgfpathquadraticcurveto{\pgfqpoint{11.953169in}{0.489655in}}{\pgfqpoint{11.953169in}{0.499778in}}%
\pgfpathlineto{\pgfqpoint{11.953169in}{0.502083in}}%
\pgfpathlineto{\pgfqpoint{11.942866in}{0.502083in}}%
\pgfpathlineto{\pgfqpoint{11.942866in}{0.502083in}}%
\pgfpathclose%
\pgfpathmoveto{\pgfqpoint{11.963526in}{0.506370in}}%
\pgfpathlineto{\pgfqpoint{11.963526in}{0.470400in}}%
\pgfpathlineto{\pgfqpoint{11.953169in}{0.470400in}}%
\pgfpathlineto{\pgfqpoint{11.953169in}{0.479964in}}%
\pgfpathquadraticcurveto{\pgfqpoint{11.949620in}{0.474237in}}{\pgfqpoint{11.944325in}{0.471499in}}%
\pgfpathquadraticcurveto{\pgfqpoint{11.939029in}{0.468761in}}{\pgfqpoint{11.931374in}{0.468761in}}%
\pgfpathquadraticcurveto{\pgfqpoint{11.921701in}{0.468761in}}{\pgfqpoint{11.915974in}{0.474201in}}%
\pgfpathquadraticcurveto{\pgfqpoint{11.910264in}{0.479640in}}{\pgfqpoint{11.910264in}{0.488754in}}%
\pgfpathquadraticcurveto{\pgfqpoint{11.910264in}{0.499382in}}{\pgfqpoint{11.917378in}{0.504785in}}%
\pgfpathquadraticcurveto{\pgfqpoint{11.924511in}{0.510189in}}{\pgfqpoint{11.938633in}{0.510189in}}%
\pgfpathlineto{\pgfqpoint{11.953169in}{0.510189in}}%
\pgfpathlineto{\pgfqpoint{11.953169in}{0.511216in}}%
\pgfpathquadraticcurveto{\pgfqpoint{11.953169in}{0.518366in}}{\pgfqpoint{11.948467in}{0.522275in}}%
\pgfpathquadraticcurveto{\pgfqpoint{11.943766in}{0.526184in}}{\pgfqpoint{11.935265in}{0.526184in}}%
\pgfpathquadraticcurveto{\pgfqpoint{11.929861in}{0.526184in}}{\pgfqpoint{11.924727in}{0.524887in}}%
\pgfpathquadraticcurveto{\pgfqpoint{11.919612in}{0.523590in}}{\pgfqpoint{11.914893in}{0.520996in}}%
\pgfpathlineto{\pgfqpoint{11.914893in}{0.530579in}}%
\pgfpathquadraticcurveto{\pgfqpoint{11.920567in}{0.532776in}}{\pgfqpoint{11.925916in}{0.533857in}}%
\pgfpathquadraticcurveto{\pgfqpoint{11.931266in}{0.534956in}}{\pgfqpoint{11.936327in}{0.534956in}}%
\pgfpathquadraticcurveto{\pgfqpoint{11.950016in}{0.534956in}}{\pgfqpoint{11.956771in}{0.527859in}}%
\pgfpathquadraticcurveto{\pgfqpoint{11.963526in}{0.520780in}}{\pgfqpoint{11.963526in}{0.506370in}}%
\pgfpathlineto{\pgfqpoint{11.963526in}{0.506370in}}%
\pgfpathclose%
\pgfpathmoveto{\pgfqpoint{11.984852in}{0.557993in}}%
\pgfpathlineto{\pgfqpoint{11.995209in}{0.557993in}}%
\pgfpathlineto{\pgfqpoint{11.995209in}{0.470400in}}%
\pgfpathlineto{\pgfqpoint{11.984852in}{0.470400in}}%
\pgfpathlineto{\pgfqpoint{11.984852in}{0.557993in}}%
\pgfpathlineto{\pgfqpoint{11.984852in}{0.557993in}}%
\pgfpathclose%
\pgfusepath{fill}%
\end{pgfscope}%
\begin{pgfscope}%
\pgfsetbuttcap%
\pgfsetroundjoin%
\definecolor{currentfill}{rgb}{0.309804,0.372549,0.419608}%
\pgfsetfillcolor{currentfill}%
\pgfsetlinewidth{0.000000pt}%
\definecolor{currentstroke}{rgb}{0.309804,0.372549,0.419608}%
\pgfsetstrokecolor{currentstroke}%
\pgfsetdash{}{0pt}%
\pgfpathmoveto{\pgfqpoint{12.088485in}{0.551347in}}%
\pgfpathlineto{\pgfqpoint{12.088485in}{0.533443in}}%
\pgfpathlineto{\pgfqpoint{12.109812in}{0.533443in}}%
\pgfpathlineto{\pgfqpoint{12.109812in}{0.525391in}}%
\pgfpathlineto{\pgfqpoint{12.088485in}{0.525391in}}%
\pgfpathlineto{\pgfqpoint{12.088485in}{0.491168in}}%
\pgfpathquadraticcurveto{\pgfqpoint{12.088485in}{0.483459in}}{\pgfqpoint{12.090593in}{0.481261in}}%
\pgfpathquadraticcurveto{\pgfqpoint{12.092700in}{0.479064in}}{\pgfqpoint{12.099184in}{0.479064in}}%
\pgfpathlineto{\pgfqpoint{12.109812in}{0.479064in}}%
\pgfpathlineto{\pgfqpoint{12.109812in}{0.470400in}}%
\pgfpathlineto{\pgfqpoint{12.099184in}{0.470400in}}%
\pgfpathquadraticcurveto{\pgfqpoint{12.087188in}{0.470400in}}{\pgfqpoint{12.082631in}{0.474867in}}%
\pgfpathquadraticcurveto{\pgfqpoint{12.078074in}{0.479352in}}{\pgfqpoint{12.078074in}{0.491168in}}%
\pgfpathlineto{\pgfqpoint{12.078074in}{0.525391in}}%
\pgfpathlineto{\pgfqpoint{12.070473in}{0.525391in}}%
\pgfpathlineto{\pgfqpoint{12.070473in}{0.533443in}}%
\pgfpathlineto{\pgfqpoint{12.078074in}{0.533443in}}%
\pgfpathlineto{\pgfqpoint{12.078074in}{0.551347in}}%
\pgfpathlineto{\pgfqpoint{12.088485in}{0.551347in}}%
\pgfpathlineto{\pgfqpoint{12.088485in}{0.551347in}}%
\pgfpathclose%
\pgfpathmoveto{\pgfqpoint{12.177357in}{0.504515in}}%
\pgfpathlineto{\pgfqpoint{12.177357in}{0.499454in}}%
\pgfpathlineto{\pgfqpoint{12.129733in}{0.499454in}}%
\pgfpathquadraticcurveto{\pgfqpoint{12.130418in}{0.488754in}}{\pgfqpoint{12.136181in}{0.483153in}}%
\pgfpathquadraticcurveto{\pgfqpoint{12.141945in}{0.477551in}}{\pgfqpoint{12.152248in}{0.477551in}}%
\pgfpathquadraticcurveto{\pgfqpoint{12.158210in}{0.477551in}}{\pgfqpoint{12.163812in}{0.479010in}}%
\pgfpathquadraticcurveto{\pgfqpoint{12.169414in}{0.480469in}}{\pgfqpoint{12.174944in}{0.483405in}}%
\pgfpathlineto{\pgfqpoint{12.174944in}{0.473606in}}%
\pgfpathquadraticcurveto{\pgfqpoint{12.169360in}{0.471247in}}{\pgfqpoint{12.163506in}{0.470004in}}%
\pgfpathquadraticcurveto{\pgfqpoint{12.157652in}{0.468761in}}{\pgfqpoint{12.151636in}{0.468761in}}%
\pgfpathquadraticcurveto{\pgfqpoint{12.136542in}{0.468761in}}{\pgfqpoint{12.127734in}{0.477533in}}%
\pgfpathquadraticcurveto{\pgfqpoint{12.118926in}{0.486323in}}{\pgfqpoint{12.118926in}{0.501309in}}%
\pgfpathquadraticcurveto{\pgfqpoint{12.118926in}{0.516781in}}{\pgfqpoint{12.127283in}{0.525859in}}%
\pgfpathquadraticcurveto{\pgfqpoint{12.135641in}{0.534956in}}{\pgfqpoint{12.149835in}{0.534956in}}%
\pgfpathquadraticcurveto{\pgfqpoint{12.162551in}{0.534956in}}{\pgfqpoint{12.169954in}{0.526760in}}%
\pgfpathquadraticcurveto{\pgfqpoint{12.177357in}{0.518583in}}{\pgfqpoint{12.177357in}{0.504515in}}%
\pgfpathlineto{\pgfqpoint{12.177357in}{0.504515in}}%
\pgfpathclose%
\pgfpathmoveto{\pgfqpoint{12.167000in}{0.507559in}}%
\pgfpathquadraticcurveto{\pgfqpoint{12.166892in}{0.516043in}}{\pgfqpoint{12.162245in}{0.521104in}}%
\pgfpathquadraticcurveto{\pgfqpoint{12.157598in}{0.526184in}}{\pgfqpoint{12.149943in}{0.526184in}}%
\pgfpathquadraticcurveto{\pgfqpoint{12.141279in}{0.526184in}}{\pgfqpoint{12.136073in}{0.521284in}}%
\pgfpathquadraticcurveto{\pgfqpoint{12.130868in}{0.516385in}}{\pgfqpoint{12.130075in}{0.507487in}}%
\pgfpathlineto{\pgfqpoint{12.167000in}{0.507559in}}%
\pgfpathlineto{\pgfqpoint{12.167000in}{0.507559in}}%
\pgfpathclose%
\pgfpathmoveto{\pgfqpoint{12.234546in}{0.531587in}}%
\pgfpathlineto{\pgfqpoint{12.234546in}{0.521789in}}%
\pgfpathquadraticcurveto{\pgfqpoint{12.230169in}{0.524040in}}{\pgfqpoint{12.225432in}{0.525157in}}%
\pgfpathquadraticcurveto{\pgfqpoint{12.220712in}{0.526292in}}{\pgfqpoint{12.215633in}{0.526292in}}%
\pgfpathquadraticcurveto{\pgfqpoint{12.207924in}{0.526292in}}{\pgfqpoint{12.204069in}{0.523932in}}%
\pgfpathquadraticcurveto{\pgfqpoint{12.200215in}{0.521573in}}{\pgfqpoint{12.200215in}{0.516835in}}%
\pgfpathquadraticcurveto{\pgfqpoint{12.200215in}{0.513233in}}{\pgfqpoint{12.202970in}{0.511180in}}%
\pgfpathquadraticcurveto{\pgfqpoint{12.205726in}{0.509126in}}{\pgfqpoint{12.214066in}{0.507271in}}%
\pgfpathlineto{\pgfqpoint{12.217614in}{0.506478in}}%
\pgfpathquadraticcurveto{\pgfqpoint{12.228638in}{0.504119in}}{\pgfqpoint{12.233285in}{0.499814in}}%
\pgfpathquadraticcurveto{\pgfqpoint{12.237932in}{0.495509in}}{\pgfqpoint{12.237932in}{0.487800in}}%
\pgfpathquadraticcurveto{\pgfqpoint{12.237932in}{0.479010in}}{\pgfqpoint{12.230979in}{0.473876in}}%
\pgfpathquadraticcurveto{\pgfqpoint{12.224027in}{0.468761in}}{\pgfqpoint{12.211868in}{0.468761in}}%
\pgfpathquadraticcurveto{\pgfqpoint{12.206807in}{0.468761in}}{\pgfqpoint{12.201313in}{0.469752in}}%
\pgfpathquadraticcurveto{\pgfqpoint{12.195820in}{0.470742in}}{\pgfqpoint{12.189750in}{0.472706in}}%
\pgfpathlineto{\pgfqpoint{12.189750in}{0.483405in}}%
\pgfpathquadraticcurveto{\pgfqpoint{12.195495in}{0.480415in}}{\pgfqpoint{12.201061in}{0.478920in}}%
\pgfpathquadraticcurveto{\pgfqpoint{12.206627in}{0.477443in}}{\pgfqpoint{12.212103in}{0.477443in}}%
\pgfpathquadraticcurveto{\pgfqpoint{12.219416in}{0.477443in}}{\pgfqpoint{12.223342in}{0.479946in}}%
\pgfpathquadraticcurveto{\pgfqpoint{12.227287in}{0.482450in}}{\pgfqpoint{12.227287in}{0.487007in}}%
\pgfpathquadraticcurveto{\pgfqpoint{12.227287in}{0.491222in}}{\pgfqpoint{12.224441in}{0.493474in}}%
\pgfpathquadraticcurveto{\pgfqpoint{12.221613in}{0.495725in}}{\pgfqpoint{12.211977in}{0.497814in}}%
\pgfpathlineto{\pgfqpoint{12.208374in}{0.498661in}}%
\pgfpathquadraticcurveto{\pgfqpoint{12.198756in}{0.500678in}}{\pgfqpoint{12.194469in}{0.504875in}}%
\pgfpathquadraticcurveto{\pgfqpoint{12.190200in}{0.509072in}}{\pgfqpoint{12.190200in}{0.516385in}}%
\pgfpathquadraticcurveto{\pgfqpoint{12.190200in}{0.525283in}}{\pgfqpoint{12.196504in}{0.530110in}}%
\pgfpathquadraticcurveto{\pgfqpoint{12.202808in}{0.534956in}}{\pgfqpoint{12.214408in}{0.534956in}}%
\pgfpathquadraticcurveto{\pgfqpoint{12.220136in}{0.534956in}}{\pgfqpoint{12.225197in}{0.534109in}}%
\pgfpathquadraticcurveto{\pgfqpoint{12.230277in}{0.533280in}}{\pgfqpoint{12.234546in}{0.531587in}}%
\pgfpathlineto{\pgfqpoint{12.234546in}{0.531587in}}%
\pgfpathclose%
\pgfpathmoveto{\pgfqpoint{12.264662in}{0.551347in}}%
\pgfpathlineto{\pgfqpoint{12.264662in}{0.533443in}}%
\pgfpathlineto{\pgfqpoint{12.285988in}{0.533443in}}%
\pgfpathlineto{\pgfqpoint{12.285988in}{0.525391in}}%
\pgfpathlineto{\pgfqpoint{12.264662in}{0.525391in}}%
\pgfpathlineto{\pgfqpoint{12.264662in}{0.491168in}}%
\pgfpathquadraticcurveto{\pgfqpoint{12.264662in}{0.483459in}}{\pgfqpoint{12.266770in}{0.481261in}}%
\pgfpathquadraticcurveto{\pgfqpoint{12.268877in}{0.479064in}}{\pgfqpoint{12.275361in}{0.479064in}}%
\pgfpathlineto{\pgfqpoint{12.285988in}{0.479064in}}%
\pgfpathlineto{\pgfqpoint{12.285988in}{0.470400in}}%
\pgfpathlineto{\pgfqpoint{12.275361in}{0.470400in}}%
\pgfpathquadraticcurveto{\pgfqpoint{12.263365in}{0.470400in}}{\pgfqpoint{12.258808in}{0.474867in}}%
\pgfpathquadraticcurveto{\pgfqpoint{12.254251in}{0.479352in}}{\pgfqpoint{12.254251in}{0.491168in}}%
\pgfpathlineto{\pgfqpoint{12.254251in}{0.525391in}}%
\pgfpathlineto{\pgfqpoint{12.246650in}{0.525391in}}%
\pgfpathlineto{\pgfqpoint{12.246650in}{0.533443in}}%
\pgfpathlineto{\pgfqpoint{12.254251in}{0.533443in}}%
\pgfpathlineto{\pgfqpoint{12.254251in}{0.551347in}}%
\pgfpathlineto{\pgfqpoint{12.264662in}{0.551347in}}%
\pgfpathlineto{\pgfqpoint{12.264662in}{0.551347in}}%
\pgfpathclose%
\pgfpathmoveto{\pgfqpoint{12.339791in}{0.531587in}}%
\pgfpathlineto{\pgfqpoint{12.339791in}{0.521789in}}%
\pgfpathquadraticcurveto{\pgfqpoint{12.335414in}{0.524040in}}{\pgfqpoint{12.330677in}{0.525157in}}%
\pgfpathquadraticcurveto{\pgfqpoint{12.325957in}{0.526292in}}{\pgfqpoint{12.320878in}{0.526292in}}%
\pgfpathquadraticcurveto{\pgfqpoint{12.313169in}{0.526292in}}{\pgfqpoint{12.309314in}{0.523932in}}%
\pgfpathquadraticcurveto{\pgfqpoint{12.305460in}{0.521573in}}{\pgfqpoint{12.305460in}{0.516835in}}%
\pgfpathquadraticcurveto{\pgfqpoint{12.305460in}{0.513233in}}{\pgfqpoint{12.308215in}{0.511180in}}%
\pgfpathquadraticcurveto{\pgfqpoint{12.310971in}{0.509126in}}{\pgfqpoint{12.319311in}{0.507271in}}%
\pgfpathlineto{\pgfqpoint{12.322859in}{0.506478in}}%
\pgfpathquadraticcurveto{\pgfqpoint{12.333883in}{0.504119in}}{\pgfqpoint{12.338530in}{0.499814in}}%
\pgfpathquadraticcurveto{\pgfqpoint{12.343177in}{0.495509in}}{\pgfqpoint{12.343177in}{0.487800in}}%
\pgfpathquadraticcurveto{\pgfqpoint{12.343177in}{0.479010in}}{\pgfqpoint{12.336224in}{0.473876in}}%
\pgfpathquadraticcurveto{\pgfqpoint{12.329272in}{0.468761in}}{\pgfqpoint{12.317113in}{0.468761in}}%
\pgfpathquadraticcurveto{\pgfqpoint{12.312052in}{0.468761in}}{\pgfqpoint{12.306558in}{0.469752in}}%
\pgfpathquadraticcurveto{\pgfqpoint{12.301065in}{0.470742in}}{\pgfqpoint{12.294995in}{0.472706in}}%
\pgfpathlineto{\pgfqpoint{12.294995in}{0.483405in}}%
\pgfpathquadraticcurveto{\pgfqpoint{12.300740in}{0.480415in}}{\pgfqpoint{12.306306in}{0.478920in}}%
\pgfpathquadraticcurveto{\pgfqpoint{12.311872in}{0.477443in}}{\pgfqpoint{12.317348in}{0.477443in}}%
\pgfpathquadraticcurveto{\pgfqpoint{12.324661in}{0.477443in}}{\pgfqpoint{12.328587in}{0.479946in}}%
\pgfpathquadraticcurveto{\pgfqpoint{12.332532in}{0.482450in}}{\pgfqpoint{12.332532in}{0.487007in}}%
\pgfpathquadraticcurveto{\pgfqpoint{12.332532in}{0.491222in}}{\pgfqpoint{12.329686in}{0.493474in}}%
\pgfpathquadraticcurveto{\pgfqpoint{12.326858in}{0.495725in}}{\pgfqpoint{12.317222in}{0.497814in}}%
\pgfpathlineto{\pgfqpoint{12.313619in}{0.498661in}}%
\pgfpathquadraticcurveto{\pgfqpoint{12.304001in}{0.500678in}}{\pgfqpoint{12.299714in}{0.504875in}}%
\pgfpathquadraticcurveto{\pgfqpoint{12.295445in}{0.509072in}}{\pgfqpoint{12.295445in}{0.516385in}}%
\pgfpathquadraticcurveto{\pgfqpoint{12.295445in}{0.525283in}}{\pgfqpoint{12.301749in}{0.530110in}}%
\pgfpathquadraticcurveto{\pgfqpoint{12.308053in}{0.534956in}}{\pgfqpoint{12.319653in}{0.534956in}}%
\pgfpathquadraticcurveto{\pgfqpoint{12.325381in}{0.534956in}}{\pgfqpoint{12.330442in}{0.534109in}}%
\pgfpathquadraticcurveto{\pgfqpoint{12.335522in}{0.533280in}}{\pgfqpoint{12.339791in}{0.531587in}}%
\pgfpathlineto{\pgfqpoint{12.339791in}{0.531587in}}%
\pgfpathclose%
\pgfpathmoveto{\pgfqpoint{12.362306in}{0.484702in}}%
\pgfpathlineto{\pgfqpoint{12.374176in}{0.484702in}}%
\pgfpathlineto{\pgfqpoint{12.374176in}{0.475011in}}%
\pgfpathlineto{\pgfqpoint{12.364954in}{0.456999in}}%
\pgfpathlineto{\pgfqpoint{12.357695in}{0.456999in}}%
\pgfpathlineto{\pgfqpoint{12.362306in}{0.475011in}}%
\pgfpathlineto{\pgfqpoint{12.362306in}{0.484702in}}%
\pgfpathlineto{\pgfqpoint{12.362306in}{0.484702in}}%
\pgfpathclose%
\pgfusepath{fill}%
\end{pgfscope}%
\begin{pgfscope}%
\pgfsetbuttcap%
\pgfsetroundjoin%
\definecolor{currentfill}{rgb}{0.309804,0.372549,0.419608}%
\pgfsetfillcolor{currentfill}%
\pgfsetlinewidth{0.000000pt}%
\definecolor{currentstroke}{rgb}{0.309804,0.372549,0.419608}%
\pgfsetstrokecolor{currentstroke}%
\pgfsetdash{}{0pt}%
\pgfpathmoveto{\pgfqpoint{12.449533in}{0.533443in}}%
\pgfpathlineto{\pgfqpoint{12.459890in}{0.533443in}}%
\pgfpathlineto{\pgfqpoint{12.472840in}{0.484251in}}%
\pgfpathlineto{\pgfqpoint{12.485719in}{0.533443in}}%
\pgfpathlineto{\pgfqpoint{12.497931in}{0.533443in}}%
\pgfpathlineto{\pgfqpoint{12.510882in}{0.484251in}}%
\pgfpathlineto{\pgfqpoint{12.523779in}{0.533443in}}%
\pgfpathlineto{\pgfqpoint{12.534136in}{0.533443in}}%
\pgfpathlineto{\pgfqpoint{12.517637in}{0.470400in}}%
\pgfpathlineto{\pgfqpoint{12.505424in}{0.470400in}}%
\pgfpathlineto{\pgfqpoint{12.491861in}{0.522077in}}%
\pgfpathlineto{\pgfqpoint{12.478244in}{0.470400in}}%
\pgfpathlineto{\pgfqpoint{12.466014in}{0.470400in}}%
\pgfpathlineto{\pgfqpoint{12.449533in}{0.533443in}}%
\pgfpathlineto{\pgfqpoint{12.449533in}{0.533443in}}%
\pgfpathclose%
\pgfpathmoveto{\pgfqpoint{12.549824in}{0.533443in}}%
\pgfpathlineto{\pgfqpoint{12.560181in}{0.533443in}}%
\pgfpathlineto{\pgfqpoint{12.560181in}{0.470400in}}%
\pgfpathlineto{\pgfqpoint{12.549824in}{0.470400in}}%
\pgfpathlineto{\pgfqpoint{12.549824in}{0.533443in}}%
\pgfpathlineto{\pgfqpoint{12.549824in}{0.533443in}}%
\pgfpathclose%
\pgfpathmoveto{\pgfqpoint{12.549824in}{0.557993in}}%
\pgfpathlineto{\pgfqpoint{12.560181in}{0.557993in}}%
\pgfpathlineto{\pgfqpoint{12.560181in}{0.544862in}}%
\pgfpathlineto{\pgfqpoint{12.549824in}{0.544862in}}%
\pgfpathlineto{\pgfqpoint{12.549824in}{0.557993in}}%
\pgfpathlineto{\pgfqpoint{12.549824in}{0.557993in}}%
\pgfpathclose%
\pgfpathmoveto{\pgfqpoint{12.592099in}{0.551347in}}%
\pgfpathlineto{\pgfqpoint{12.592099in}{0.533443in}}%
\pgfpathlineto{\pgfqpoint{12.613425in}{0.533443in}}%
\pgfpathlineto{\pgfqpoint{12.613425in}{0.525391in}}%
\pgfpathlineto{\pgfqpoint{12.592099in}{0.525391in}}%
\pgfpathlineto{\pgfqpoint{12.592099in}{0.491168in}}%
\pgfpathquadraticcurveto{\pgfqpoint{12.592099in}{0.483459in}}{\pgfqpoint{12.594206in}{0.481261in}}%
\pgfpathquadraticcurveto{\pgfqpoint{12.596314in}{0.479064in}}{\pgfqpoint{12.602798in}{0.479064in}}%
\pgfpathlineto{\pgfqpoint{12.613425in}{0.479064in}}%
\pgfpathlineto{\pgfqpoint{12.613425in}{0.470400in}}%
\pgfpathlineto{\pgfqpoint{12.602798in}{0.470400in}}%
\pgfpathquadraticcurveto{\pgfqpoint{12.590802in}{0.470400in}}{\pgfqpoint{12.586245in}{0.474867in}}%
\pgfpathquadraticcurveto{\pgfqpoint{12.581688in}{0.479352in}}{\pgfqpoint{12.581688in}{0.491168in}}%
\pgfpathlineto{\pgfqpoint{12.581688in}{0.525391in}}%
\pgfpathlineto{\pgfqpoint{12.574087in}{0.525391in}}%
\pgfpathlineto{\pgfqpoint{12.574087in}{0.533443in}}%
\pgfpathlineto{\pgfqpoint{12.581688in}{0.533443in}}%
\pgfpathlineto{\pgfqpoint{12.581688in}{0.551347in}}%
\pgfpathlineto{\pgfqpoint{12.592099in}{0.551347in}}%
\pgfpathlineto{\pgfqpoint{12.592099in}{0.551347in}}%
\pgfpathclose%
\pgfpathmoveto{\pgfqpoint{12.679458in}{0.508460in}}%
\pgfpathlineto{\pgfqpoint{12.679458in}{0.470400in}}%
\pgfpathlineto{\pgfqpoint{12.669101in}{0.470400in}}%
\pgfpathlineto{\pgfqpoint{12.669101in}{0.508117in}}%
\pgfpathquadraticcurveto{\pgfqpoint{12.669101in}{0.517069in}}{\pgfqpoint{12.665607in}{0.521500in}}%
\pgfpathquadraticcurveto{\pgfqpoint{12.662112in}{0.525949in}}{\pgfqpoint{12.655141in}{0.525949in}}%
\pgfpathquadraticcurveto{\pgfqpoint{12.646748in}{0.525949in}}{\pgfqpoint{12.641903in}{0.520600in}}%
\pgfpathquadraticcurveto{\pgfqpoint{12.637057in}{0.515268in}}{\pgfqpoint{12.637057in}{0.506028in}}%
\pgfpathlineto{\pgfqpoint{12.637057in}{0.470400in}}%
\pgfpathlineto{\pgfqpoint{12.626646in}{0.470400in}}%
\pgfpathlineto{\pgfqpoint{12.626646in}{0.557993in}}%
\pgfpathlineto{\pgfqpoint{12.637057in}{0.557993in}}%
\pgfpathlineto{\pgfqpoint{12.637057in}{0.523644in}}%
\pgfpathquadraticcurveto{\pgfqpoint{12.640786in}{0.529336in}}{\pgfqpoint{12.645811in}{0.532146in}}%
\pgfpathquadraticcurveto{\pgfqpoint{12.650855in}{0.534956in}}{\pgfqpoint{12.657447in}{0.534956in}}%
\pgfpathquadraticcurveto{\pgfqpoint{12.668308in}{0.534956in}}{\pgfqpoint{12.673874in}{0.528237in}}%
\pgfpathquadraticcurveto{\pgfqpoint{12.679458in}{0.521518in}}{\pgfqpoint{12.679458in}{0.508460in}}%
\pgfpathlineto{\pgfqpoint{12.679458in}{0.508460in}}%
\pgfpathclose%
\pgfpathmoveto{\pgfqpoint{12.724524in}{0.526184in}}%
\pgfpathquadraticcurveto{\pgfqpoint{12.716203in}{0.526184in}}{\pgfqpoint{12.711357in}{0.519681in}}%
\pgfpathquadraticcurveto{\pgfqpoint{12.706512in}{0.513179in}}{\pgfqpoint{12.706512in}{0.501867in}}%
\pgfpathquadraticcurveto{\pgfqpoint{12.706512in}{0.490556in}}{\pgfqpoint{12.711321in}{0.484053in}}%
\pgfpathquadraticcurveto{\pgfqpoint{12.716149in}{0.477551in}}{\pgfqpoint{12.724524in}{0.477551in}}%
\pgfpathquadraticcurveto{\pgfqpoint{12.732810in}{0.477551in}}{\pgfqpoint{12.737637in}{0.484071in}}%
\pgfpathquadraticcurveto{\pgfqpoint{12.742482in}{0.490610in}}{\pgfqpoint{12.742482in}{0.501867in}}%
\pgfpathquadraticcurveto{\pgfqpoint{12.742482in}{0.513071in}}{\pgfqpoint{12.737637in}{0.519627in}}%
\pgfpathquadraticcurveto{\pgfqpoint{12.732810in}{0.526184in}}{\pgfqpoint{12.724524in}{0.526184in}}%
\pgfpathlineto{\pgfqpoint{12.724524in}{0.526184in}}%
\pgfpathclose%
\pgfpathmoveto{\pgfqpoint{12.724524in}{0.534956in}}%
\pgfpathquadraticcurveto{\pgfqpoint{12.738033in}{0.534956in}}{\pgfqpoint{12.745743in}{0.526166in}}%
\pgfpathquadraticcurveto{\pgfqpoint{12.753470in}{0.517394in}}{\pgfqpoint{12.753470in}{0.501867in}}%
\pgfpathquadraticcurveto{\pgfqpoint{12.753470in}{0.486395in}}{\pgfqpoint{12.745743in}{0.477569in}}%
\pgfpathquadraticcurveto{\pgfqpoint{12.738033in}{0.468761in}}{\pgfqpoint{12.724524in}{0.468761in}}%
\pgfpathquadraticcurveto{\pgfqpoint{12.710961in}{0.468761in}}{\pgfqpoint{12.703270in}{0.477569in}}%
\pgfpathquadraticcurveto{\pgfqpoint{12.695597in}{0.486395in}}{\pgfqpoint{12.695597in}{0.501867in}}%
\pgfpathquadraticcurveto{\pgfqpoint{12.695597in}{0.517394in}}{\pgfqpoint{12.703270in}{0.526166in}}%
\pgfpathquadraticcurveto{\pgfqpoint{12.710961in}{0.534956in}}{\pgfqpoint{12.724524in}{0.534956in}}%
\pgfpathlineto{\pgfqpoint{12.724524in}{0.534956in}}%
\pgfpathclose%
\pgfpathmoveto{\pgfqpoint{12.769573in}{0.495275in}}%
\pgfpathlineto{\pgfqpoint{12.769573in}{0.533443in}}%
\pgfpathlineto{\pgfqpoint{12.779930in}{0.533443in}}%
\pgfpathlineto{\pgfqpoint{12.779930in}{0.495671in}}%
\pgfpathquadraticcurveto{\pgfqpoint{12.779930in}{0.486719in}}{\pgfqpoint{12.783406in}{0.482234in}}%
\pgfpathquadraticcurveto{\pgfqpoint{12.786900in}{0.477767in}}{\pgfqpoint{12.793889in}{0.477767in}}%
\pgfpathquadraticcurveto{\pgfqpoint{12.802265in}{0.477767in}}{\pgfqpoint{12.807128in}{0.483117in}}%
\pgfpathquadraticcurveto{\pgfqpoint{12.812009in}{0.488466in}}{\pgfqpoint{12.812009in}{0.497706in}}%
\pgfpathlineto{\pgfqpoint{12.812009in}{0.533443in}}%
\pgfpathlineto{\pgfqpoint{12.822366in}{0.533443in}}%
\pgfpathlineto{\pgfqpoint{12.822366in}{0.470400in}}%
\pgfpathlineto{\pgfqpoint{12.812009in}{0.470400in}}%
\pgfpathlineto{\pgfqpoint{12.812009in}{0.480091in}}%
\pgfpathquadraticcurveto{\pgfqpoint{12.808245in}{0.474345in}}{\pgfqpoint{12.803255in}{0.471553in}}%
\pgfpathquadraticcurveto{\pgfqpoint{12.798284in}{0.468761in}}{\pgfqpoint{12.791692in}{0.468761in}}%
\pgfpathquadraticcurveto{\pgfqpoint{12.780830in}{0.468761in}}{\pgfqpoint{12.775192in}{0.475515in}}%
\pgfpathquadraticcurveto{\pgfqpoint{12.769573in}{0.482270in}}{\pgfqpoint{12.769573in}{0.495275in}}%
\pgfpathlineto{\pgfqpoint{12.769573in}{0.495275in}}%
\pgfpathclose%
\pgfpathmoveto{\pgfqpoint{12.795636in}{0.534956in}}%
\pgfpathlineto{\pgfqpoint{12.795636in}{0.534956in}}%
\pgfpathlineto{\pgfqpoint{12.795636in}{0.534956in}}%
\pgfpathclose%
\pgfpathmoveto{\pgfqpoint{12.853942in}{0.551347in}}%
\pgfpathlineto{\pgfqpoint{12.853942in}{0.533443in}}%
\pgfpathlineto{\pgfqpoint{12.875268in}{0.533443in}}%
\pgfpathlineto{\pgfqpoint{12.875268in}{0.525391in}}%
\pgfpathlineto{\pgfqpoint{12.853942in}{0.525391in}}%
\pgfpathlineto{\pgfqpoint{12.853942in}{0.491168in}}%
\pgfpathquadraticcurveto{\pgfqpoint{12.853942in}{0.483459in}}{\pgfqpoint{12.856049in}{0.481261in}}%
\pgfpathquadraticcurveto{\pgfqpoint{12.858156in}{0.479064in}}{\pgfqpoint{12.864641in}{0.479064in}}%
\pgfpathlineto{\pgfqpoint{12.875268in}{0.479064in}}%
\pgfpathlineto{\pgfqpoint{12.875268in}{0.470400in}}%
\pgfpathlineto{\pgfqpoint{12.864641in}{0.470400in}}%
\pgfpathquadraticcurveto{\pgfqpoint{12.852645in}{0.470400in}}{\pgfqpoint{12.848088in}{0.474867in}}%
\pgfpathquadraticcurveto{\pgfqpoint{12.843531in}{0.479352in}}{\pgfqpoint{12.843531in}{0.491168in}}%
\pgfpathlineto{\pgfqpoint{12.843531in}{0.525391in}}%
\pgfpathlineto{\pgfqpoint{12.835929in}{0.525391in}}%
\pgfpathlineto{\pgfqpoint{12.835929in}{0.533443in}}%
\pgfpathlineto{\pgfqpoint{12.843531in}{0.533443in}}%
\pgfpathlineto{\pgfqpoint{12.843531in}{0.551347in}}%
\pgfpathlineto{\pgfqpoint{12.853942in}{0.551347in}}%
\pgfpathlineto{\pgfqpoint{12.853942in}{0.551347in}}%
\pgfpathclose%
\pgfusepath{fill}%
\end{pgfscope}%
\begin{pgfscope}%
\pgfsetbuttcap%
\pgfsetroundjoin%
\definecolor{currentfill}{rgb}{0.309804,0.372549,0.419608}%
\pgfsetfillcolor{currentfill}%
\pgfsetlinewidth{0.000000pt}%
\definecolor{currentstroke}{rgb}{0.309804,0.372549,0.419608}%
\pgfsetstrokecolor{currentstroke}%
\pgfsetdash{}{0pt}%
\pgfpathmoveto{\pgfqpoint{13.011944in}{0.521338in}}%
\pgfpathquadraticcurveto{\pgfqpoint{13.015835in}{0.528327in}}{\pgfqpoint{13.021239in}{0.531641in}}%
\pgfpathquadraticcurveto{\pgfqpoint{13.026642in}{0.534956in}}{\pgfqpoint{13.033955in}{0.534956in}}%
\pgfpathquadraticcurveto{\pgfqpoint{13.043808in}{0.534956in}}{\pgfqpoint{13.049158in}{0.528057in}}%
\pgfpathquadraticcurveto{\pgfqpoint{13.054507in}{0.521176in}}{\pgfqpoint{13.054507in}{0.508460in}}%
\pgfpathlineto{\pgfqpoint{13.054507in}{0.470400in}}%
\pgfpathlineto{\pgfqpoint{13.044096in}{0.470400in}}%
\pgfpathlineto{\pgfqpoint{13.044096in}{0.508117in}}%
\pgfpathquadraticcurveto{\pgfqpoint{13.044096in}{0.517178in}}{\pgfqpoint{13.040872in}{0.521555in}}%
\pgfpathquadraticcurveto{\pgfqpoint{13.037666in}{0.525949in}}{\pgfqpoint{13.031091in}{0.525949in}}%
\pgfpathquadraticcurveto{\pgfqpoint{13.023040in}{0.525949in}}{\pgfqpoint{13.018357in}{0.520600in}}%
\pgfpathquadraticcurveto{\pgfqpoint{13.013692in}{0.515268in}}{\pgfqpoint{13.013692in}{0.506028in}}%
\pgfpathlineto{\pgfqpoint{13.013692in}{0.470400in}}%
\pgfpathlineto{\pgfqpoint{13.003281in}{0.470400in}}%
\pgfpathlineto{\pgfqpoint{13.003281in}{0.508117in}}%
\pgfpathquadraticcurveto{\pgfqpoint{13.003281in}{0.517232in}}{\pgfqpoint{13.000074in}{0.521591in}}%
\pgfpathquadraticcurveto{\pgfqpoint{12.996868in}{0.525949in}}{\pgfqpoint{12.990168in}{0.525949in}}%
\pgfpathquadraticcurveto{\pgfqpoint{12.982224in}{0.525949in}}{\pgfqpoint{12.977541in}{0.520582in}}%
\pgfpathquadraticcurveto{\pgfqpoint{12.972876in}{0.515214in}}{\pgfqpoint{12.972876in}{0.506028in}}%
\pgfpathlineto{\pgfqpoint{12.972876in}{0.470400in}}%
\pgfpathlineto{\pgfqpoint{12.962465in}{0.470400in}}%
\pgfpathlineto{\pgfqpoint{12.962465in}{0.533443in}}%
\pgfpathlineto{\pgfqpoint{12.972876in}{0.533443in}}%
\pgfpathlineto{\pgfqpoint{12.972876in}{0.523644in}}%
\pgfpathquadraticcurveto{\pgfqpoint{12.976424in}{0.529444in}}{\pgfqpoint{12.981378in}{0.532200in}}%
\pgfpathquadraticcurveto{\pgfqpoint{12.986331in}{0.534956in}}{\pgfqpoint{12.993140in}{0.534956in}}%
\pgfpathquadraticcurveto{\pgfqpoint{13.000020in}{0.534956in}}{\pgfqpoint{13.004830in}{0.531461in}}%
\pgfpathquadraticcurveto{\pgfqpoint{13.009639in}{0.527985in}}{\pgfqpoint{13.011944in}{0.521338in}}%
\pgfpathlineto{\pgfqpoint{13.011944in}{0.521338in}}%
\pgfpathclose%
\pgfpathmoveto{\pgfqpoint{13.074086in}{0.495275in}}%
\pgfpathlineto{\pgfqpoint{13.074086in}{0.533443in}}%
\pgfpathlineto{\pgfqpoint{13.084443in}{0.533443in}}%
\pgfpathlineto{\pgfqpoint{13.084443in}{0.495671in}}%
\pgfpathquadraticcurveto{\pgfqpoint{13.084443in}{0.486719in}}{\pgfqpoint{13.087920in}{0.482234in}}%
\pgfpathquadraticcurveto{\pgfqpoint{13.091414in}{0.477767in}}{\pgfqpoint{13.098403in}{0.477767in}}%
\pgfpathquadraticcurveto{\pgfqpoint{13.106778in}{0.477767in}}{\pgfqpoint{13.111642in}{0.483117in}}%
\pgfpathquadraticcurveto{\pgfqpoint{13.116523in}{0.488466in}}{\pgfqpoint{13.116523in}{0.497706in}}%
\pgfpathlineto{\pgfqpoint{13.116523in}{0.533443in}}%
\pgfpathlineto{\pgfqpoint{13.126880in}{0.533443in}}%
\pgfpathlineto{\pgfqpoint{13.126880in}{0.470400in}}%
\pgfpathlineto{\pgfqpoint{13.116523in}{0.470400in}}%
\pgfpathlineto{\pgfqpoint{13.116523in}{0.480091in}}%
\pgfpathquadraticcurveto{\pgfqpoint{13.112758in}{0.474345in}}{\pgfqpoint{13.107769in}{0.471553in}}%
\pgfpathquadraticcurveto{\pgfqpoint{13.102798in}{0.468761in}}{\pgfqpoint{13.096205in}{0.468761in}}%
\pgfpathquadraticcurveto{\pgfqpoint{13.085344in}{0.468761in}}{\pgfqpoint{13.079706in}{0.475515in}}%
\pgfpathquadraticcurveto{\pgfqpoint{13.074086in}{0.482270in}}{\pgfqpoint{13.074086in}{0.495275in}}%
\pgfpathlineto{\pgfqpoint{13.074086in}{0.495275in}}%
\pgfpathclose%
\pgfpathmoveto{\pgfqpoint{13.100150in}{0.534956in}}%
\pgfpathlineto{\pgfqpoint{13.100150in}{0.534956in}}%
\pgfpathlineto{\pgfqpoint{13.100150in}{0.534956in}}%
\pgfpathclose%
\pgfpathmoveto{\pgfqpoint{13.148206in}{0.557993in}}%
\pgfpathlineto{\pgfqpoint{13.158563in}{0.557993in}}%
\pgfpathlineto{\pgfqpoint{13.158563in}{0.470400in}}%
\pgfpathlineto{\pgfqpoint{13.148206in}{0.470400in}}%
\pgfpathlineto{\pgfqpoint{13.148206in}{0.557993in}}%
\pgfpathlineto{\pgfqpoint{13.148206in}{0.557993in}}%
\pgfpathclose%
\pgfpathmoveto{\pgfqpoint{13.190481in}{0.551347in}}%
\pgfpathlineto{\pgfqpoint{13.190481in}{0.533443in}}%
\pgfpathlineto{\pgfqpoint{13.211807in}{0.533443in}}%
\pgfpathlineto{\pgfqpoint{13.211807in}{0.525391in}}%
\pgfpathlineto{\pgfqpoint{13.190481in}{0.525391in}}%
\pgfpathlineto{\pgfqpoint{13.190481in}{0.491168in}}%
\pgfpathquadraticcurveto{\pgfqpoint{13.190481in}{0.483459in}}{\pgfqpoint{13.192588in}{0.481261in}}%
\pgfpathquadraticcurveto{\pgfqpoint{13.194696in}{0.479064in}}{\pgfqpoint{13.201180in}{0.479064in}}%
\pgfpathlineto{\pgfqpoint{13.211807in}{0.479064in}}%
\pgfpathlineto{\pgfqpoint{13.211807in}{0.470400in}}%
\pgfpathlineto{\pgfqpoint{13.201180in}{0.470400in}}%
\pgfpathquadraticcurveto{\pgfqpoint{13.189184in}{0.470400in}}{\pgfqpoint{13.184627in}{0.474867in}}%
\pgfpathquadraticcurveto{\pgfqpoint{13.180070in}{0.479352in}}{\pgfqpoint{13.180070in}{0.491168in}}%
\pgfpathlineto{\pgfqpoint{13.180070in}{0.525391in}}%
\pgfpathlineto{\pgfqpoint{13.172469in}{0.525391in}}%
\pgfpathlineto{\pgfqpoint{13.172469in}{0.533443in}}%
\pgfpathlineto{\pgfqpoint{13.180070in}{0.533443in}}%
\pgfpathlineto{\pgfqpoint{13.180070in}{0.551347in}}%
\pgfpathlineto{\pgfqpoint{13.190481in}{0.551347in}}%
\pgfpathlineto{\pgfqpoint{13.190481in}{0.551347in}}%
\pgfpathclose%
\pgfpathmoveto{\pgfqpoint{13.225424in}{0.533443in}}%
\pgfpathlineto{\pgfqpoint{13.235781in}{0.533443in}}%
\pgfpathlineto{\pgfqpoint{13.235781in}{0.470400in}}%
\pgfpathlineto{\pgfqpoint{13.225424in}{0.470400in}}%
\pgfpathlineto{\pgfqpoint{13.225424in}{0.533443in}}%
\pgfpathlineto{\pgfqpoint{13.225424in}{0.533443in}}%
\pgfpathclose%
\pgfpathmoveto{\pgfqpoint{13.225424in}{0.557993in}}%
\pgfpathlineto{\pgfqpoint{13.235781in}{0.557993in}}%
\pgfpathlineto{\pgfqpoint{13.235781in}{0.544862in}}%
\pgfpathlineto{\pgfqpoint{13.225424in}{0.544862in}}%
\pgfpathlineto{\pgfqpoint{13.225424in}{0.557993in}}%
\pgfpathlineto{\pgfqpoint{13.225424in}{0.557993in}}%
\pgfpathclose%
\pgfpathmoveto{\pgfqpoint{13.267465in}{0.479856in}}%
\pgfpathlineto{\pgfqpoint{13.267465in}{0.446426in}}%
\pgfpathlineto{\pgfqpoint{13.257054in}{0.446426in}}%
\pgfpathlineto{\pgfqpoint{13.257054in}{0.533443in}}%
\pgfpathlineto{\pgfqpoint{13.267465in}{0.533443in}}%
\pgfpathlineto{\pgfqpoint{13.267465in}{0.523878in}}%
\pgfpathquadraticcurveto{\pgfqpoint{13.270743in}{0.529498in}}{\pgfqpoint{13.275714in}{0.532218in}}%
\pgfpathquadraticcurveto{\pgfqpoint{13.280704in}{0.534956in}}{\pgfqpoint{13.287620in}{0.534956in}}%
\pgfpathquadraticcurveto{\pgfqpoint{13.299112in}{0.534956in}}{\pgfqpoint{13.306281in}{0.525841in}}%
\pgfpathquadraticcurveto{\pgfqpoint{13.313468in}{0.516727in}}{\pgfqpoint{13.313468in}{0.501867in}}%
\pgfpathquadraticcurveto{\pgfqpoint{13.313468in}{0.487007in}}{\pgfqpoint{13.306281in}{0.477875in}}%
\pgfpathquadraticcurveto{\pgfqpoint{13.299112in}{0.468761in}}{\pgfqpoint{13.287620in}{0.468761in}}%
\pgfpathquadraticcurveto{\pgfqpoint{13.280704in}{0.468761in}}{\pgfqpoint{13.275714in}{0.471499in}}%
\pgfpathquadraticcurveto{\pgfqpoint{13.270743in}{0.474237in}}{\pgfqpoint{13.267465in}{0.479856in}}%
\pgfpathlineto{\pgfqpoint{13.267465in}{0.479856in}}%
\pgfpathclose%
\pgfpathmoveto{\pgfqpoint{13.302715in}{0.501867in}}%
\pgfpathquadraticcurveto{\pgfqpoint{13.302715in}{0.513287in}}{\pgfqpoint{13.298013in}{0.519789in}}%
\pgfpathquadraticcurveto{\pgfqpoint{13.293312in}{0.526292in}}{\pgfqpoint{13.285099in}{0.526292in}}%
\pgfpathquadraticcurveto{\pgfqpoint{13.276867in}{0.526292in}}{\pgfqpoint{13.272166in}{0.519789in}}%
\pgfpathquadraticcurveto{\pgfqpoint{13.267465in}{0.513287in}}{\pgfqpoint{13.267465in}{0.501867in}}%
\pgfpathquadraticcurveto{\pgfqpoint{13.267465in}{0.490448in}}{\pgfqpoint{13.272166in}{0.483945in}}%
\pgfpathquadraticcurveto{\pgfqpoint{13.276867in}{0.477443in}}{\pgfqpoint{13.285099in}{0.477443in}}%
\pgfpathquadraticcurveto{\pgfqpoint{13.293312in}{0.477443in}}{\pgfqpoint{13.298013in}{0.483945in}}%
\pgfpathquadraticcurveto{\pgfqpoint{13.302715in}{0.490448in}}{\pgfqpoint{13.302715in}{0.501867in}}%
\pgfpathlineto{\pgfqpoint{13.302715in}{0.501867in}}%
\pgfpathclose%
\pgfpathmoveto{\pgfqpoint{13.330633in}{0.557993in}}%
\pgfpathlineto{\pgfqpoint{13.340990in}{0.557993in}}%
\pgfpathlineto{\pgfqpoint{13.340990in}{0.470400in}}%
\pgfpathlineto{\pgfqpoint{13.330633in}{0.470400in}}%
\pgfpathlineto{\pgfqpoint{13.330633in}{0.557993in}}%
\pgfpathlineto{\pgfqpoint{13.330633in}{0.557993in}}%
\pgfpathclose%
\pgfpathmoveto{\pgfqpoint{13.362659in}{0.533443in}}%
\pgfpathlineto{\pgfqpoint{13.373016in}{0.533443in}}%
\pgfpathlineto{\pgfqpoint{13.373016in}{0.470400in}}%
\pgfpathlineto{\pgfqpoint{13.362659in}{0.470400in}}%
\pgfpathlineto{\pgfqpoint{13.362659in}{0.533443in}}%
\pgfpathlineto{\pgfqpoint{13.362659in}{0.533443in}}%
\pgfpathclose%
\pgfpathmoveto{\pgfqpoint{13.362659in}{0.557993in}}%
\pgfpathlineto{\pgfqpoint{13.373016in}{0.557993in}}%
\pgfpathlineto{\pgfqpoint{13.373016in}{0.544862in}}%
\pgfpathlineto{\pgfqpoint{13.362659in}{0.544862in}}%
\pgfpathlineto{\pgfqpoint{13.362659in}{0.557993in}}%
\pgfpathlineto{\pgfqpoint{13.362659in}{0.557993in}}%
\pgfpathclose%
\pgfpathmoveto{\pgfqpoint{13.440057in}{0.531029in}}%
\pgfpathlineto{\pgfqpoint{13.440057in}{0.521338in}}%
\pgfpathquadraticcurveto{\pgfqpoint{13.435662in}{0.523770in}}{\pgfqpoint{13.431249in}{0.524977in}}%
\pgfpathquadraticcurveto{\pgfqpoint{13.426836in}{0.526184in}}{\pgfqpoint{13.422333in}{0.526184in}}%
\pgfpathquadraticcurveto{\pgfqpoint{13.412247in}{0.526184in}}{\pgfqpoint{13.406663in}{0.519789in}}%
\pgfpathquadraticcurveto{\pgfqpoint{13.401097in}{0.513413in}}{\pgfqpoint{13.401097in}{0.501867in}}%
\pgfpathquadraticcurveto{\pgfqpoint{13.401097in}{0.490321in}}{\pgfqpoint{13.406663in}{0.483927in}}%
\pgfpathquadraticcurveto{\pgfqpoint{13.412247in}{0.477551in}}{\pgfqpoint{13.422333in}{0.477551in}}%
\pgfpathquadraticcurveto{\pgfqpoint{13.426836in}{0.477551in}}{\pgfqpoint{13.431249in}{0.478758in}}%
\pgfpathquadraticcurveto{\pgfqpoint{13.435662in}{0.479964in}}{\pgfqpoint{13.440057in}{0.482396in}}%
\pgfpathlineto{\pgfqpoint{13.440057in}{0.472814in}}%
\pgfpathquadraticcurveto{\pgfqpoint{13.435716in}{0.470796in}}{\pgfqpoint{13.431069in}{0.469788in}}%
\pgfpathquadraticcurveto{\pgfqpoint{13.426440in}{0.468761in}}{\pgfqpoint{13.421199in}{0.468761in}}%
\pgfpathquadraticcurveto{\pgfqpoint{13.406951in}{0.468761in}}{\pgfqpoint{13.398557in}{0.477713in}}%
\pgfpathquadraticcurveto{\pgfqpoint{13.390182in}{0.486665in}}{\pgfqpoint{13.390182in}{0.501867in}}%
\pgfpathquadraticcurveto{\pgfqpoint{13.390182in}{0.517286in}}{\pgfqpoint{13.398647in}{0.526112in}}%
\pgfpathquadraticcurveto{\pgfqpoint{13.407131in}{0.534956in}}{\pgfqpoint{13.421883in}{0.534956in}}%
\pgfpathquadraticcurveto{\pgfqpoint{13.426656in}{0.534956in}}{\pgfqpoint{13.431213in}{0.533965in}}%
\pgfpathquadraticcurveto{\pgfqpoint{13.435770in}{0.532992in}}{\pgfqpoint{13.440057in}{0.531029in}}%
\pgfpathlineto{\pgfqpoint{13.440057in}{0.531029in}}%
\pgfpathclose%
\pgfpathmoveto{\pgfqpoint{13.458069in}{0.533443in}}%
\pgfpathlineto{\pgfqpoint{13.468426in}{0.533443in}}%
\pgfpathlineto{\pgfqpoint{13.468426in}{0.470400in}}%
\pgfpathlineto{\pgfqpoint{13.458069in}{0.470400in}}%
\pgfpathlineto{\pgfqpoint{13.458069in}{0.533443in}}%
\pgfpathlineto{\pgfqpoint{13.458069in}{0.533443in}}%
\pgfpathclose%
\pgfpathmoveto{\pgfqpoint{13.458069in}{0.557993in}}%
\pgfpathlineto{\pgfqpoint{13.468426in}{0.557993in}}%
\pgfpathlineto{\pgfqpoint{13.468426in}{0.544862in}}%
\pgfpathlineto{\pgfqpoint{13.458069in}{0.544862in}}%
\pgfpathlineto{\pgfqpoint{13.458069in}{0.557993in}}%
\pgfpathlineto{\pgfqpoint{13.458069in}{0.557993in}}%
\pgfpathclose%
\pgfpathmoveto{\pgfqpoint{13.500344in}{0.551347in}}%
\pgfpathlineto{\pgfqpoint{13.500344in}{0.533443in}}%
\pgfpathlineto{\pgfqpoint{13.521670in}{0.533443in}}%
\pgfpathlineto{\pgfqpoint{13.521670in}{0.525391in}}%
\pgfpathlineto{\pgfqpoint{13.500344in}{0.525391in}}%
\pgfpathlineto{\pgfqpoint{13.500344in}{0.491168in}}%
\pgfpathquadraticcurveto{\pgfqpoint{13.500344in}{0.483459in}}{\pgfqpoint{13.502451in}{0.481261in}}%
\pgfpathquadraticcurveto{\pgfqpoint{13.504559in}{0.479064in}}{\pgfqpoint{13.511043in}{0.479064in}}%
\pgfpathlineto{\pgfqpoint{13.521670in}{0.479064in}}%
\pgfpathlineto{\pgfqpoint{13.521670in}{0.470400in}}%
\pgfpathlineto{\pgfqpoint{13.511043in}{0.470400in}}%
\pgfpathquadraticcurveto{\pgfqpoint{13.499047in}{0.470400in}}{\pgfqpoint{13.494490in}{0.474867in}}%
\pgfpathquadraticcurveto{\pgfqpoint{13.489933in}{0.479352in}}{\pgfqpoint{13.489933in}{0.491168in}}%
\pgfpathlineto{\pgfqpoint{13.489933in}{0.525391in}}%
\pgfpathlineto{\pgfqpoint{13.482332in}{0.525391in}}%
\pgfpathlineto{\pgfqpoint{13.482332in}{0.533443in}}%
\pgfpathlineto{\pgfqpoint{13.489933in}{0.533443in}}%
\pgfpathlineto{\pgfqpoint{13.489933in}{0.551347in}}%
\pgfpathlineto{\pgfqpoint{13.500344in}{0.551347in}}%
\pgfpathlineto{\pgfqpoint{13.500344in}{0.551347in}}%
\pgfpathclose%
\pgfpathmoveto{\pgfqpoint{13.561513in}{0.464546in}}%
\pgfpathquadraticcurveto{\pgfqpoint{13.557136in}{0.453288in}}{\pgfqpoint{13.552957in}{0.449866in}}%
\pgfpathquadraticcurveto{\pgfqpoint{13.548797in}{0.446426in}}{\pgfqpoint{13.541826in}{0.446426in}}%
\pgfpathlineto{\pgfqpoint{13.533540in}{0.446426in}}%
\pgfpathlineto{\pgfqpoint{13.533540in}{0.455090in}}%
\pgfpathlineto{\pgfqpoint{13.539628in}{0.455090in}}%
\pgfpathquadraticcurveto{\pgfqpoint{13.543897in}{0.455090in}}{\pgfqpoint{13.546257in}{0.457125in}}%
\pgfpathquadraticcurveto{\pgfqpoint{13.548635in}{0.459142in}}{\pgfqpoint{13.551498in}{0.466689in}}%
\pgfpathlineto{\pgfqpoint{13.553354in}{0.471409in}}%
\pgfpathlineto{\pgfqpoint{13.527867in}{0.533443in}}%
\pgfpathlineto{\pgfqpoint{13.538836in}{0.533443in}}%
\pgfpathlineto{\pgfqpoint{13.558541in}{0.484143in}}%
\pgfpathlineto{\pgfqpoint{13.578247in}{0.533443in}}%
\pgfpathlineto{\pgfqpoint{13.589216in}{0.533443in}}%
\pgfpathlineto{\pgfqpoint{13.561513in}{0.464546in}}%
\pgfpathlineto{\pgfqpoint{13.561513in}{0.464546in}}%
\pgfpathclose%
\pgfusepath{fill}%
\end{pgfscope}%
\begin{pgfscope}%
\pgfsetbuttcap%
\pgfsetroundjoin%
\definecolor{currentfill}{rgb}{0.309804,0.372549,0.419608}%
\pgfsetfillcolor{currentfill}%
\pgfsetlinewidth{0.000000pt}%
\definecolor{currentstroke}{rgb}{0.309804,0.372549,0.419608}%
\pgfsetstrokecolor{currentstroke}%
\pgfsetdash{}{0pt}%
\pgfpathmoveto{\pgfqpoint{0.963519in}{0.338883in}}%
\pgfpathquadraticcurveto{\pgfqpoint{0.950964in}{0.338883in}}{\pgfqpoint{0.946119in}{0.336019in}}%
\pgfpathquadraticcurveto{\pgfqpoint{0.941274in}{0.333156in}}{\pgfqpoint{0.941274in}{0.326221in}}%
\pgfpathquadraticcurveto{\pgfqpoint{0.941274in}{0.320709in}}{\pgfqpoint{0.944912in}{0.317467in}}%
\pgfpathquadraticcurveto{\pgfqpoint{0.948551in}{0.314243in}}{\pgfqpoint{0.954783in}{0.314243in}}%
\pgfpathquadraticcurveto{\pgfqpoint{0.963411in}{0.314243in}}{\pgfqpoint{0.968616in}{0.320349in}}%
\pgfpathquadraticcurveto{\pgfqpoint{0.973822in}{0.326455in}}{\pgfqpoint{0.973822in}{0.336578in}}%
\pgfpathlineto{\pgfqpoint{0.973822in}{0.338883in}}%
\pgfpathlineto{\pgfqpoint{0.963519in}{0.338883in}}%
\pgfpathlineto{\pgfqpoint{0.963519in}{0.338883in}}%
\pgfpathclose%
\pgfpathmoveto{\pgfqpoint{0.984179in}{0.343170in}}%
\pgfpathlineto{\pgfqpoint{0.984179in}{0.307200in}}%
\pgfpathlineto{\pgfqpoint{0.973822in}{0.307200in}}%
\pgfpathlineto{\pgfqpoint{0.973822in}{0.316764in}}%
\pgfpathquadraticcurveto{\pgfqpoint{0.970273in}{0.311037in}}{\pgfqpoint{0.964978in}{0.308299in}}%
\pgfpathquadraticcurveto{\pgfqpoint{0.959682in}{0.305561in}}{\pgfqpoint{0.952027in}{0.305561in}}%
\pgfpathquadraticcurveto{\pgfqpoint{0.942354in}{0.305561in}}{\pgfqpoint{0.936627in}{0.311001in}}%
\pgfpathquadraticcurveto{\pgfqpoint{0.930917in}{0.316440in}}{\pgfqpoint{0.930917in}{0.325554in}}%
\pgfpathquadraticcurveto{\pgfqpoint{0.930917in}{0.336182in}}{\pgfqpoint{0.938031in}{0.341585in}}%
\pgfpathquadraticcurveto{\pgfqpoint{0.945164in}{0.346989in}}{\pgfqpoint{0.959286in}{0.346989in}}%
\pgfpathlineto{\pgfqpoint{0.973822in}{0.346989in}}%
\pgfpathlineto{\pgfqpoint{0.973822in}{0.348016in}}%
\pgfpathquadraticcurveto{\pgfqpoint{0.973822in}{0.355166in}}{\pgfqpoint{0.969120in}{0.359075in}}%
\pgfpathquadraticcurveto{\pgfqpoint{0.964419in}{0.362984in}}{\pgfqpoint{0.955918in}{0.362984in}}%
\pgfpathquadraticcurveto{\pgfqpoint{0.950514in}{0.362984in}}{\pgfqpoint{0.945380in}{0.361687in}}%
\pgfpathquadraticcurveto{\pgfqpoint{0.940265in}{0.360390in}}{\pgfqpoint{0.935546in}{0.357796in}}%
\pgfpathlineto{\pgfqpoint{0.935546in}{0.367379in}}%
\pgfpathquadraticcurveto{\pgfqpoint{0.941220in}{0.369576in}}{\pgfqpoint{0.946569in}{0.370657in}}%
\pgfpathquadraticcurveto{\pgfqpoint{0.951919in}{0.371756in}}{\pgfqpoint{0.956980in}{0.371756in}}%
\pgfpathquadraticcurveto{\pgfqpoint{0.970669in}{0.371756in}}{\pgfqpoint{0.977424in}{0.364659in}}%
\pgfpathquadraticcurveto{\pgfqpoint{0.984179in}{0.357580in}}{\pgfqpoint{0.984179in}{0.343170in}}%
\pgfpathlineto{\pgfqpoint{0.984179in}{0.343170in}}%
\pgfpathclose%
\pgfpathmoveto{\pgfqpoint{1.046987in}{0.360678in}}%
\pgfpathlineto{\pgfqpoint{1.046987in}{0.394793in}}%
\pgfpathlineto{\pgfqpoint{1.057344in}{0.394793in}}%
\pgfpathlineto{\pgfqpoint{1.057344in}{0.307200in}}%
\pgfpathlineto{\pgfqpoint{1.046987in}{0.307200in}}%
\pgfpathlineto{\pgfqpoint{1.046987in}{0.316656in}}%
\pgfpathquadraticcurveto{\pgfqpoint{1.043727in}{0.311037in}}{\pgfqpoint{1.038737in}{0.308299in}}%
\pgfpathquadraticcurveto{\pgfqpoint{1.033766in}{0.305561in}}{\pgfqpoint{1.026777in}{0.305561in}}%
\pgfpathquadraticcurveto{\pgfqpoint{1.015358in}{0.305561in}}{\pgfqpoint{1.008171in}{0.314675in}}%
\pgfpathquadraticcurveto{\pgfqpoint{1.001002in}{0.323807in}}{\pgfqpoint{1.001002in}{0.338667in}}%
\pgfpathquadraticcurveto{\pgfqpoint{1.001002in}{0.353527in}}{\pgfqpoint{1.008171in}{0.362641in}}%
\pgfpathquadraticcurveto{\pgfqpoint{1.015358in}{0.371756in}}{\pgfqpoint{1.026777in}{0.371756in}}%
\pgfpathquadraticcurveto{\pgfqpoint{1.033766in}{0.371756in}}{\pgfqpoint{1.038737in}{0.369018in}}%
\pgfpathquadraticcurveto{\pgfqpoint{1.043727in}{0.366298in}}{\pgfqpoint{1.046987in}{0.360678in}}%
\pgfpathlineto{\pgfqpoint{1.046987in}{0.360678in}}%
\pgfpathclose%
\pgfpathmoveto{\pgfqpoint{1.011701in}{0.338667in}}%
\pgfpathquadraticcurveto{\pgfqpoint{1.011701in}{0.327248in}}{\pgfqpoint{1.016402in}{0.320745in}}%
\pgfpathquadraticcurveto{\pgfqpoint{1.021104in}{0.314243in}}{\pgfqpoint{1.029317in}{0.314243in}}%
\pgfpathquadraticcurveto{\pgfqpoint{1.037531in}{0.314243in}}{\pgfqpoint{1.042250in}{0.320745in}}%
\pgfpathquadraticcurveto{\pgfqpoint{1.046987in}{0.327248in}}{\pgfqpoint{1.046987in}{0.338667in}}%
\pgfpathquadraticcurveto{\pgfqpoint{1.046987in}{0.350087in}}{\pgfqpoint{1.042250in}{0.356589in}}%
\pgfpathquadraticcurveto{\pgfqpoint{1.037531in}{0.363092in}}{\pgfqpoint{1.029317in}{0.363092in}}%
\pgfpathquadraticcurveto{\pgfqpoint{1.021104in}{0.363092in}}{\pgfqpoint{1.016402in}{0.356589in}}%
\pgfpathquadraticcurveto{\pgfqpoint{1.011701in}{0.350087in}}{\pgfqpoint{1.011701in}{0.338667in}}%
\pgfpathlineto{\pgfqpoint{1.011701in}{0.338667in}}%
\pgfpathclose%
\pgfpathmoveto{\pgfqpoint{1.078688in}{0.370243in}}%
\pgfpathlineto{\pgfqpoint{1.089045in}{0.370243in}}%
\pgfpathlineto{\pgfqpoint{1.089045in}{0.306065in}}%
\pgfpathquadraticcurveto{\pgfqpoint{1.089045in}{0.294033in}}{\pgfqpoint{1.084452in}{0.288629in}}%
\pgfpathquadraticcurveto{\pgfqpoint{1.079877in}{0.283226in}}{\pgfqpoint{1.069682in}{0.283226in}}%
\pgfpathlineto{\pgfqpoint{1.065738in}{0.283226in}}%
\pgfpathlineto{\pgfqpoint{1.065738in}{0.291998in}}%
\pgfpathlineto{\pgfqpoint{1.068512in}{0.291998in}}%
\pgfpathquadraticcurveto{\pgfqpoint{1.074419in}{0.291998in}}{\pgfqpoint{1.076545in}{0.294736in}}%
\pgfpathquadraticcurveto{\pgfqpoint{1.078688in}{0.297455in}}{\pgfqpoint{1.078688in}{0.306065in}}%
\pgfpathlineto{\pgfqpoint{1.078688in}{0.370243in}}%
\pgfpathlineto{\pgfqpoint{1.078688in}{0.370243in}}%
\pgfpathclose%
\pgfpathmoveto{\pgfqpoint{1.078688in}{0.394793in}}%
\pgfpathlineto{\pgfqpoint{1.089045in}{0.394793in}}%
\pgfpathlineto{\pgfqpoint{1.089045in}{0.381662in}}%
\pgfpathlineto{\pgfqpoint{1.078688in}{0.381662in}}%
\pgfpathlineto{\pgfqpoint{1.078688in}{0.394793in}}%
\pgfpathlineto{\pgfqpoint{1.078688in}{0.394793in}}%
\pgfpathclose%
\pgfpathmoveto{\pgfqpoint{1.109651in}{0.332075in}}%
\pgfpathlineto{\pgfqpoint{1.109651in}{0.370243in}}%
\pgfpathlineto{\pgfqpoint{1.120008in}{0.370243in}}%
\pgfpathlineto{\pgfqpoint{1.120008in}{0.332471in}}%
\pgfpathquadraticcurveto{\pgfqpoint{1.120008in}{0.323519in}}{\pgfqpoint{1.123485in}{0.319034in}}%
\pgfpathquadraticcurveto{\pgfqpoint{1.126979in}{0.314567in}}{\pgfqpoint{1.133968in}{0.314567in}}%
\pgfpathquadraticcurveto{\pgfqpoint{1.142343in}{0.314567in}}{\pgfqpoint{1.147207in}{0.319917in}}%
\pgfpathquadraticcurveto{\pgfqpoint{1.152088in}{0.325266in}}{\pgfqpoint{1.152088in}{0.334506in}}%
\pgfpathlineto{\pgfqpoint{1.152088in}{0.370243in}}%
\pgfpathlineto{\pgfqpoint{1.162445in}{0.370243in}}%
\pgfpathlineto{\pgfqpoint{1.162445in}{0.307200in}}%
\pgfpathlineto{\pgfqpoint{1.152088in}{0.307200in}}%
\pgfpathlineto{\pgfqpoint{1.152088in}{0.316891in}}%
\pgfpathquadraticcurveto{\pgfqpoint{1.148323in}{0.311145in}}{\pgfqpoint{1.143334in}{0.308353in}}%
\pgfpathquadraticcurveto{\pgfqpoint{1.138363in}{0.305561in}}{\pgfqpoint{1.131770in}{0.305561in}}%
\pgfpathquadraticcurveto{\pgfqpoint{1.120909in}{0.305561in}}{\pgfqpoint{1.115271in}{0.312315in}}%
\pgfpathquadraticcurveto{\pgfqpoint{1.109651in}{0.319070in}}{\pgfqpoint{1.109651in}{0.332075in}}%
\pgfpathlineto{\pgfqpoint{1.109651in}{0.332075in}}%
\pgfpathclose%
\pgfpathmoveto{\pgfqpoint{1.135715in}{0.371756in}}%
\pgfpathlineto{\pgfqpoint{1.135715in}{0.371756in}}%
\pgfpathlineto{\pgfqpoint{1.135715in}{0.371756in}}%
\pgfpathclose%
\pgfpathmoveto{\pgfqpoint{1.223956in}{0.368387in}}%
\pgfpathlineto{\pgfqpoint{1.223956in}{0.358589in}}%
\pgfpathquadraticcurveto{\pgfqpoint{1.219579in}{0.360840in}}{\pgfqpoint{1.214842in}{0.361957in}}%
\pgfpathquadraticcurveto{\pgfqpoint{1.210123in}{0.363092in}}{\pgfqpoint{1.205044in}{0.363092in}}%
\pgfpathquadraticcurveto{\pgfqpoint{1.197334in}{0.363092in}}{\pgfqpoint{1.193480in}{0.360732in}}%
\pgfpathquadraticcurveto{\pgfqpoint{1.189625in}{0.358373in}}{\pgfqpoint{1.189625in}{0.353635in}}%
\pgfpathquadraticcurveto{\pgfqpoint{1.189625in}{0.350033in}}{\pgfqpoint{1.192381in}{0.347980in}}%
\pgfpathquadraticcurveto{\pgfqpoint{1.195137in}{0.345926in}}{\pgfqpoint{1.203477in}{0.344071in}}%
\pgfpathlineto{\pgfqpoint{1.207025in}{0.343278in}}%
\pgfpathquadraticcurveto{\pgfqpoint{1.218048in}{0.340919in}}{\pgfqpoint{1.222696in}{0.336614in}}%
\pgfpathquadraticcurveto{\pgfqpoint{1.227343in}{0.332309in}}{\pgfqpoint{1.227343in}{0.324600in}}%
\pgfpathquadraticcurveto{\pgfqpoint{1.227343in}{0.315810in}}{\pgfqpoint{1.220390in}{0.310676in}}%
\pgfpathquadraticcurveto{\pgfqpoint{1.213437in}{0.305561in}}{\pgfqpoint{1.201279in}{0.305561in}}%
\pgfpathquadraticcurveto{\pgfqpoint{1.196218in}{0.305561in}}{\pgfqpoint{1.190724in}{0.306552in}}%
\pgfpathquadraticcurveto{\pgfqpoint{1.185230in}{0.307542in}}{\pgfqpoint{1.179160in}{0.309506in}}%
\pgfpathlineto{\pgfqpoint{1.179160in}{0.320205in}}%
\pgfpathquadraticcurveto{\pgfqpoint{1.184906in}{0.317215in}}{\pgfqpoint{1.190472in}{0.315720in}}%
\pgfpathquadraticcurveto{\pgfqpoint{1.196038in}{0.314243in}}{\pgfqpoint{1.201513in}{0.314243in}}%
\pgfpathquadraticcurveto{\pgfqpoint{1.208826in}{0.314243in}}{\pgfqpoint{1.212753in}{0.316746in}}%
\pgfpathquadraticcurveto{\pgfqpoint{1.216697in}{0.319250in}}{\pgfqpoint{1.216697in}{0.323807in}}%
\pgfpathquadraticcurveto{\pgfqpoint{1.216697in}{0.328022in}}{\pgfqpoint{1.213852in}{0.330274in}}%
\pgfpathquadraticcurveto{\pgfqpoint{1.211024in}{0.332525in}}{\pgfqpoint{1.201387in}{0.334614in}}%
\pgfpathlineto{\pgfqpoint{1.197785in}{0.335461in}}%
\pgfpathquadraticcurveto{\pgfqpoint{1.188166in}{0.337478in}}{\pgfqpoint{1.183879in}{0.341675in}}%
\pgfpathquadraticcurveto{\pgfqpoint{1.179610in}{0.345872in}}{\pgfqpoint{1.179610in}{0.353185in}}%
\pgfpathquadraticcurveto{\pgfqpoint{1.179610in}{0.362083in}}{\pgfqpoint{1.185915in}{0.366910in}}%
\pgfpathquadraticcurveto{\pgfqpoint{1.192219in}{0.371756in}}{\pgfqpoint{1.203819in}{0.371756in}}%
\pgfpathquadraticcurveto{\pgfqpoint{1.209547in}{0.371756in}}{\pgfqpoint{1.214608in}{0.370909in}}%
\pgfpathquadraticcurveto{\pgfqpoint{1.219688in}{0.370080in}}{\pgfqpoint{1.223956in}{0.368387in}}%
\pgfpathlineto{\pgfqpoint{1.223956in}{0.368387in}}%
\pgfpathclose%
\pgfpathmoveto{\pgfqpoint{1.254073in}{0.388147in}}%
\pgfpathlineto{\pgfqpoint{1.254073in}{0.370243in}}%
\pgfpathlineto{\pgfqpoint{1.275399in}{0.370243in}}%
\pgfpathlineto{\pgfqpoint{1.275399in}{0.362191in}}%
\pgfpathlineto{\pgfqpoint{1.254073in}{0.362191in}}%
\pgfpathlineto{\pgfqpoint{1.254073in}{0.327968in}}%
\pgfpathquadraticcurveto{\pgfqpoint{1.254073in}{0.320259in}}{\pgfqpoint{1.256180in}{0.318061in}}%
\pgfpathquadraticcurveto{\pgfqpoint{1.258288in}{0.315864in}}{\pgfqpoint{1.264772in}{0.315864in}}%
\pgfpathlineto{\pgfqpoint{1.275399in}{0.315864in}}%
\pgfpathlineto{\pgfqpoint{1.275399in}{0.307200in}}%
\pgfpathlineto{\pgfqpoint{1.264772in}{0.307200in}}%
\pgfpathquadraticcurveto{\pgfqpoint{1.252776in}{0.307200in}}{\pgfqpoint{1.248219in}{0.311667in}}%
\pgfpathquadraticcurveto{\pgfqpoint{1.243662in}{0.316152in}}{\pgfqpoint{1.243662in}{0.327968in}}%
\pgfpathlineto{\pgfqpoint{1.243662in}{0.362191in}}%
\pgfpathlineto{\pgfqpoint{1.236061in}{0.362191in}}%
\pgfpathlineto{\pgfqpoint{1.236061in}{0.370243in}}%
\pgfpathlineto{\pgfqpoint{1.243662in}{0.370243in}}%
\pgfpathlineto{\pgfqpoint{1.243662in}{0.388147in}}%
\pgfpathlineto{\pgfqpoint{1.254073in}{0.388147in}}%
\pgfpathlineto{\pgfqpoint{1.254073in}{0.388147in}}%
\pgfpathclose%
\pgfpathmoveto{\pgfqpoint{1.338099in}{0.358138in}}%
\pgfpathquadraticcurveto{\pgfqpoint{1.341990in}{0.365127in}}{\pgfqpoint{1.347394in}{0.368441in}}%
\pgfpathquadraticcurveto{\pgfqpoint{1.352797in}{0.371756in}}{\pgfqpoint{1.360110in}{0.371756in}}%
\pgfpathquadraticcurveto{\pgfqpoint{1.369963in}{0.371756in}}{\pgfqpoint{1.375312in}{0.364857in}}%
\pgfpathquadraticcurveto{\pgfqpoint{1.380662in}{0.357976in}}{\pgfqpoint{1.380662in}{0.345260in}}%
\pgfpathlineto{\pgfqpoint{1.380662in}{0.307200in}}%
\pgfpathlineto{\pgfqpoint{1.370251in}{0.307200in}}%
\pgfpathlineto{\pgfqpoint{1.370251in}{0.344917in}}%
\pgfpathquadraticcurveto{\pgfqpoint{1.370251in}{0.353978in}}{\pgfqpoint{1.367027in}{0.358355in}}%
\pgfpathquadraticcurveto{\pgfqpoint{1.363821in}{0.362749in}}{\pgfqpoint{1.357246in}{0.362749in}}%
\pgfpathquadraticcurveto{\pgfqpoint{1.349195in}{0.362749in}}{\pgfqpoint{1.344512in}{0.357400in}}%
\pgfpathquadraticcurveto{\pgfqpoint{1.339847in}{0.352068in}}{\pgfqpoint{1.339847in}{0.342828in}}%
\pgfpathlineto{\pgfqpoint{1.339847in}{0.307200in}}%
\pgfpathlineto{\pgfqpoint{1.329436in}{0.307200in}}%
\pgfpathlineto{\pgfqpoint{1.329436in}{0.344917in}}%
\pgfpathquadraticcurveto{\pgfqpoint{1.329436in}{0.354032in}}{\pgfqpoint{1.326229in}{0.358391in}}%
\pgfpathquadraticcurveto{\pgfqpoint{1.323023in}{0.362749in}}{\pgfqpoint{1.316323in}{0.362749in}}%
\pgfpathquadraticcurveto{\pgfqpoint{1.308379in}{0.362749in}}{\pgfqpoint{1.303696in}{0.357382in}}%
\pgfpathquadraticcurveto{\pgfqpoint{1.299031in}{0.352014in}}{\pgfqpoint{1.299031in}{0.342828in}}%
\pgfpathlineto{\pgfqpoint{1.299031in}{0.307200in}}%
\pgfpathlineto{\pgfqpoint{1.288620in}{0.307200in}}%
\pgfpathlineto{\pgfqpoint{1.288620in}{0.370243in}}%
\pgfpathlineto{\pgfqpoint{1.299031in}{0.370243in}}%
\pgfpathlineto{\pgfqpoint{1.299031in}{0.360444in}}%
\pgfpathquadraticcurveto{\pgfqpoint{1.302579in}{0.366244in}}{\pgfqpoint{1.307533in}{0.369000in}}%
\pgfpathquadraticcurveto{\pgfqpoint{1.312486in}{0.371756in}}{\pgfqpoint{1.319295in}{0.371756in}}%
\pgfpathquadraticcurveto{\pgfqpoint{1.326175in}{0.371756in}}{\pgfqpoint{1.330985in}{0.368261in}}%
\pgfpathquadraticcurveto{\pgfqpoint{1.335794in}{0.364785in}}{\pgfqpoint{1.338099in}{0.358138in}}%
\pgfpathlineto{\pgfqpoint{1.338099in}{0.358138in}}%
\pgfpathclose%
\pgfpathmoveto{\pgfqpoint{1.455232in}{0.341315in}}%
\pgfpathlineto{\pgfqpoint{1.455232in}{0.336254in}}%
\pgfpathlineto{\pgfqpoint{1.407608in}{0.336254in}}%
\pgfpathquadraticcurveto{\pgfqpoint{1.408293in}{0.325554in}}{\pgfqpoint{1.414057in}{0.319953in}}%
\pgfpathquadraticcurveto{\pgfqpoint{1.419821in}{0.314351in}}{\pgfqpoint{1.430123in}{0.314351in}}%
\pgfpathquadraticcurveto{\pgfqpoint{1.436086in}{0.314351in}}{\pgfqpoint{1.441687in}{0.315810in}}%
\pgfpathquadraticcurveto{\pgfqpoint{1.447289in}{0.317269in}}{\pgfqpoint{1.452819in}{0.320205in}}%
\pgfpathlineto{\pgfqpoint{1.452819in}{0.310406in}}%
\pgfpathquadraticcurveto{\pgfqpoint{1.447235in}{0.308047in}}{\pgfqpoint{1.441381in}{0.306804in}}%
\pgfpathquadraticcurveto{\pgfqpoint{1.435527in}{0.305561in}}{\pgfqpoint{1.429511in}{0.305561in}}%
\pgfpathquadraticcurveto{\pgfqpoint{1.414417in}{0.305561in}}{\pgfqpoint{1.405609in}{0.314333in}}%
\pgfpathquadraticcurveto{\pgfqpoint{1.396801in}{0.323123in}}{\pgfqpoint{1.396801in}{0.338109in}}%
\pgfpathquadraticcurveto{\pgfqpoint{1.396801in}{0.353581in}}{\pgfqpoint{1.405159in}{0.362659in}}%
\pgfpathquadraticcurveto{\pgfqpoint{1.413516in}{0.371756in}}{\pgfqpoint{1.427710in}{0.371756in}}%
\pgfpathquadraticcurveto{\pgfqpoint{1.440426in}{0.371756in}}{\pgfqpoint{1.447829in}{0.363560in}}%
\pgfpathquadraticcurveto{\pgfqpoint{1.455232in}{0.355383in}}{\pgfqpoint{1.455232in}{0.341315in}}%
\pgfpathlineto{\pgfqpoint{1.455232in}{0.341315in}}%
\pgfpathclose%
\pgfpathmoveto{\pgfqpoint{1.444875in}{0.344359in}}%
\pgfpathquadraticcurveto{\pgfqpoint{1.444767in}{0.352843in}}{\pgfqpoint{1.440120in}{0.357904in}}%
\pgfpathquadraticcurveto{\pgfqpoint{1.435473in}{0.362984in}}{\pgfqpoint{1.427818in}{0.362984in}}%
\pgfpathquadraticcurveto{\pgfqpoint{1.419154in}{0.362984in}}{\pgfqpoint{1.413949in}{0.358084in}}%
\pgfpathquadraticcurveto{\pgfqpoint{1.408743in}{0.353185in}}{\pgfqpoint{1.407951in}{0.344287in}}%
\pgfpathlineto{\pgfqpoint{1.444875in}{0.344359in}}%
\pgfpathlineto{\pgfqpoint{1.444875in}{0.344359in}}%
\pgfpathclose%
\pgfpathmoveto{\pgfqpoint{1.524651in}{0.345260in}}%
\pgfpathlineto{\pgfqpoint{1.524651in}{0.307200in}}%
\pgfpathlineto{\pgfqpoint{1.514294in}{0.307200in}}%
\pgfpathlineto{\pgfqpoint{1.514294in}{0.344917in}}%
\pgfpathquadraticcurveto{\pgfqpoint{1.514294in}{0.353869in}}{\pgfqpoint{1.510800in}{0.358300in}}%
\pgfpathquadraticcurveto{\pgfqpoint{1.507306in}{0.362749in}}{\pgfqpoint{1.500335in}{0.362749in}}%
\pgfpathquadraticcurveto{\pgfqpoint{1.491941in}{0.362749in}}{\pgfqpoint{1.487096in}{0.357400in}}%
\pgfpathquadraticcurveto{\pgfqpoint{1.482251in}{0.352068in}}{\pgfqpoint{1.482251in}{0.342828in}}%
\pgfpathlineto{\pgfqpoint{1.482251in}{0.307200in}}%
\pgfpathlineto{\pgfqpoint{1.471840in}{0.307200in}}%
\pgfpathlineto{\pgfqpoint{1.471840in}{0.370243in}}%
\pgfpathlineto{\pgfqpoint{1.482251in}{0.370243in}}%
\pgfpathlineto{\pgfqpoint{1.482251in}{0.360444in}}%
\pgfpathquadraticcurveto{\pgfqpoint{1.485979in}{0.366136in}}{\pgfqpoint{1.491005in}{0.368946in}}%
\pgfpathquadraticcurveto{\pgfqpoint{1.496048in}{0.371756in}}{\pgfqpoint{1.502640in}{0.371756in}}%
\pgfpathquadraticcurveto{\pgfqpoint{1.513502in}{0.371756in}}{\pgfqpoint{1.519067in}{0.365037in}}%
\pgfpathquadraticcurveto{\pgfqpoint{1.524651in}{0.358318in}}{\pgfqpoint{1.524651in}{0.345260in}}%
\pgfpathlineto{\pgfqpoint{1.524651in}{0.345260in}}%
\pgfpathclose%
\pgfpathmoveto{\pgfqpoint{1.555542in}{0.388147in}}%
\pgfpathlineto{\pgfqpoint{1.555542in}{0.370243in}}%
\pgfpathlineto{\pgfqpoint{1.576868in}{0.370243in}}%
\pgfpathlineto{\pgfqpoint{1.576868in}{0.362191in}}%
\pgfpathlineto{\pgfqpoint{1.555542in}{0.362191in}}%
\pgfpathlineto{\pgfqpoint{1.555542in}{0.327968in}}%
\pgfpathquadraticcurveto{\pgfqpoint{1.555542in}{0.320259in}}{\pgfqpoint{1.557650in}{0.318061in}}%
\pgfpathquadraticcurveto{\pgfqpoint{1.559757in}{0.315864in}}{\pgfqpoint{1.566241in}{0.315864in}}%
\pgfpathlineto{\pgfqpoint{1.576868in}{0.315864in}}%
\pgfpathlineto{\pgfqpoint{1.576868in}{0.307200in}}%
\pgfpathlineto{\pgfqpoint{1.566241in}{0.307200in}}%
\pgfpathquadraticcurveto{\pgfqpoint{1.554245in}{0.307200in}}{\pgfqpoint{1.549688in}{0.311667in}}%
\pgfpathquadraticcurveto{\pgfqpoint{1.545131in}{0.316152in}}{\pgfqpoint{1.545131in}{0.327968in}}%
\pgfpathlineto{\pgfqpoint{1.545131in}{0.362191in}}%
\pgfpathlineto{\pgfqpoint{1.537530in}{0.362191in}}%
\pgfpathlineto{\pgfqpoint{1.537530in}{0.370243in}}%
\pgfpathlineto{\pgfqpoint{1.545131in}{0.370243in}}%
\pgfpathlineto{\pgfqpoint{1.545131in}{0.388147in}}%
\pgfpathlineto{\pgfqpoint{1.555542in}{0.388147in}}%
\pgfpathlineto{\pgfqpoint{1.555542in}{0.388147in}}%
\pgfpathclose%
\pgfpathmoveto{\pgfqpoint{1.591945in}{0.321502in}}%
\pgfpathlineto{\pgfqpoint{1.603833in}{0.321502in}}%
\pgfpathlineto{\pgfqpoint{1.603833in}{0.307200in}}%
\pgfpathlineto{\pgfqpoint{1.591945in}{0.307200in}}%
\pgfpathlineto{\pgfqpoint{1.591945in}{0.321502in}}%
\pgfpathlineto{\pgfqpoint{1.591945in}{0.321502in}}%
\pgfpathclose%
\pgfpathmoveto{\pgfqpoint{1.675575in}{0.381896in}}%
\pgfpathlineto{\pgfqpoint{1.675575in}{0.316548in}}%
\pgfpathlineto{\pgfqpoint{1.689318in}{0.316548in}}%
\pgfpathquadraticcurveto{\pgfqpoint{1.706718in}{0.316548in}}{\pgfqpoint{1.714788in}{0.324420in}}%
\pgfpathquadraticcurveto{\pgfqpoint{1.722857in}{0.332309in}}{\pgfqpoint{1.722857in}{0.349312in}}%
\pgfpathquadraticcurveto{\pgfqpoint{1.722857in}{0.366190in}}{\pgfqpoint{1.714788in}{0.374043in}}%
\pgfpathquadraticcurveto{\pgfqpoint{1.706718in}{0.381896in}}{\pgfqpoint{1.689318in}{0.381896in}}%
\pgfpathlineto{\pgfqpoint{1.675575in}{0.381896in}}%
\pgfpathlineto{\pgfqpoint{1.675575in}{0.381896in}}%
\pgfpathclose%
\pgfpathmoveto{\pgfqpoint{1.664209in}{0.391245in}}%
\pgfpathlineto{\pgfqpoint{1.687571in}{0.391245in}}%
\pgfpathquadraticcurveto{\pgfqpoint{1.711996in}{0.391245in}}{\pgfqpoint{1.723415in}{0.381086in}}%
\pgfpathquadraticcurveto{\pgfqpoint{1.734853in}{0.370927in}}{\pgfqpoint{1.734853in}{0.349312in}}%
\pgfpathquadraticcurveto{\pgfqpoint{1.734853in}{0.327572in}}{\pgfqpoint{1.723361in}{0.317377in}}%
\pgfpathquadraticcurveto{\pgfqpoint{1.711888in}{0.307200in}}{\pgfqpoint{1.687571in}{0.307200in}}%
\pgfpathlineto{\pgfqpoint{1.664209in}{0.307200in}}%
\pgfpathlineto{\pgfqpoint{1.664209in}{0.391245in}}%
\pgfpathlineto{\pgfqpoint{1.664209in}{0.391245in}}%
\pgfpathclose%
\pgfpathmoveto{\pgfqpoint{1.781180in}{0.338883in}}%
\pgfpathquadraticcurveto{\pgfqpoint{1.768626in}{0.338883in}}{\pgfqpoint{1.763781in}{0.336019in}}%
\pgfpathquadraticcurveto{\pgfqpoint{1.758935in}{0.333156in}}{\pgfqpoint{1.758935in}{0.326221in}}%
\pgfpathquadraticcurveto{\pgfqpoint{1.758935in}{0.320709in}}{\pgfqpoint{1.762574in}{0.317467in}}%
\pgfpathquadraticcurveto{\pgfqpoint{1.766212in}{0.314243in}}{\pgfqpoint{1.772444in}{0.314243in}}%
\pgfpathquadraticcurveto{\pgfqpoint{1.781072in}{0.314243in}}{\pgfqpoint{1.786278in}{0.320349in}}%
\pgfpathquadraticcurveto{\pgfqpoint{1.791483in}{0.326455in}}{\pgfqpoint{1.791483in}{0.336578in}}%
\pgfpathlineto{\pgfqpoint{1.791483in}{0.338883in}}%
\pgfpathlineto{\pgfqpoint{1.781180in}{0.338883in}}%
\pgfpathlineto{\pgfqpoint{1.781180in}{0.338883in}}%
\pgfpathclose%
\pgfpathmoveto{\pgfqpoint{1.801840in}{0.343170in}}%
\pgfpathlineto{\pgfqpoint{1.801840in}{0.307200in}}%
\pgfpathlineto{\pgfqpoint{1.791483in}{0.307200in}}%
\pgfpathlineto{\pgfqpoint{1.791483in}{0.316764in}}%
\pgfpathquadraticcurveto{\pgfqpoint{1.787935in}{0.311037in}}{\pgfqpoint{1.782639in}{0.308299in}}%
\pgfpathquadraticcurveto{\pgfqpoint{1.777344in}{0.305561in}}{\pgfqpoint{1.769689in}{0.305561in}}%
\pgfpathquadraticcurveto{\pgfqpoint{1.760016in}{0.305561in}}{\pgfqpoint{1.754288in}{0.311001in}}%
\pgfpathquadraticcurveto{\pgfqpoint{1.748578in}{0.316440in}}{\pgfqpoint{1.748578in}{0.325554in}}%
\pgfpathquadraticcurveto{\pgfqpoint{1.748578in}{0.336182in}}{\pgfqpoint{1.755693in}{0.341585in}}%
\pgfpathquadraticcurveto{\pgfqpoint{1.762826in}{0.346989in}}{\pgfqpoint{1.776947in}{0.346989in}}%
\pgfpathlineto{\pgfqpoint{1.791483in}{0.346989in}}%
\pgfpathlineto{\pgfqpoint{1.791483in}{0.348016in}}%
\pgfpathquadraticcurveto{\pgfqpoint{1.791483in}{0.355166in}}{\pgfqpoint{1.786782in}{0.359075in}}%
\pgfpathquadraticcurveto{\pgfqpoint{1.782081in}{0.362984in}}{\pgfqpoint{1.773579in}{0.362984in}}%
\pgfpathquadraticcurveto{\pgfqpoint{1.768176in}{0.362984in}}{\pgfqpoint{1.763042in}{0.361687in}}%
\pgfpathquadraticcurveto{\pgfqpoint{1.757927in}{0.360390in}}{\pgfqpoint{1.753207in}{0.357796in}}%
\pgfpathlineto{\pgfqpoint{1.753207in}{0.367379in}}%
\pgfpathquadraticcurveto{\pgfqpoint{1.758881in}{0.369576in}}{\pgfqpoint{1.764231in}{0.370657in}}%
\pgfpathquadraticcurveto{\pgfqpoint{1.769581in}{0.371756in}}{\pgfqpoint{1.774642in}{0.371756in}}%
\pgfpathquadraticcurveto{\pgfqpoint{1.788331in}{0.371756in}}{\pgfqpoint{1.795086in}{0.364659in}}%
\pgfpathquadraticcurveto{\pgfqpoint{1.801840in}{0.357580in}}{\pgfqpoint{1.801840in}{0.343170in}}%
\pgfpathlineto{\pgfqpoint{1.801840in}{0.343170in}}%
\pgfpathclose%
\pgfpathmoveto{\pgfqpoint{1.863352in}{0.368387in}}%
\pgfpathlineto{\pgfqpoint{1.863352in}{0.358589in}}%
\pgfpathquadraticcurveto{\pgfqpoint{1.858975in}{0.360840in}}{\pgfqpoint{1.854238in}{0.361957in}}%
\pgfpathquadraticcurveto{\pgfqpoint{1.849518in}{0.363092in}}{\pgfqpoint{1.844439in}{0.363092in}}%
\pgfpathquadraticcurveto{\pgfqpoint{1.836730in}{0.363092in}}{\pgfqpoint{1.832875in}{0.360732in}}%
\pgfpathquadraticcurveto{\pgfqpoint{1.829021in}{0.358373in}}{\pgfqpoint{1.829021in}{0.353635in}}%
\pgfpathquadraticcurveto{\pgfqpoint{1.829021in}{0.350033in}}{\pgfqpoint{1.831776in}{0.347980in}}%
\pgfpathquadraticcurveto{\pgfqpoint{1.834532in}{0.345926in}}{\pgfqpoint{1.842872in}{0.344071in}}%
\pgfpathlineto{\pgfqpoint{1.846420in}{0.343278in}}%
\pgfpathquadraticcurveto{\pgfqpoint{1.857444in}{0.340919in}}{\pgfqpoint{1.862091in}{0.336614in}}%
\pgfpathquadraticcurveto{\pgfqpoint{1.866738in}{0.332309in}}{\pgfqpoint{1.866738in}{0.324600in}}%
\pgfpathquadraticcurveto{\pgfqpoint{1.866738in}{0.315810in}}{\pgfqpoint{1.859785in}{0.310676in}}%
\pgfpathquadraticcurveto{\pgfqpoint{1.852833in}{0.305561in}}{\pgfqpoint{1.840674in}{0.305561in}}%
\pgfpathquadraticcurveto{\pgfqpoint{1.835613in}{0.305561in}}{\pgfqpoint{1.830119in}{0.306552in}}%
\pgfpathquadraticcurveto{\pgfqpoint{1.824626in}{0.307542in}}{\pgfqpoint{1.818556in}{0.309506in}}%
\pgfpathlineto{\pgfqpoint{1.818556in}{0.320205in}}%
\pgfpathquadraticcurveto{\pgfqpoint{1.824301in}{0.317215in}}{\pgfqpoint{1.829867in}{0.315720in}}%
\pgfpathquadraticcurveto{\pgfqpoint{1.835433in}{0.314243in}}{\pgfqpoint{1.840909in}{0.314243in}}%
\pgfpathquadraticcurveto{\pgfqpoint{1.848222in}{0.314243in}}{\pgfqpoint{1.852148in}{0.316746in}}%
\pgfpathquadraticcurveto{\pgfqpoint{1.856093in}{0.319250in}}{\pgfqpoint{1.856093in}{0.323807in}}%
\pgfpathquadraticcurveto{\pgfqpoint{1.856093in}{0.328022in}}{\pgfqpoint{1.853247in}{0.330274in}}%
\pgfpathquadraticcurveto{\pgfqpoint{1.850419in}{0.332525in}}{\pgfqpoint{1.840783in}{0.334614in}}%
\pgfpathlineto{\pgfqpoint{1.837180in}{0.335461in}}%
\pgfpathquadraticcurveto{\pgfqpoint{1.827562in}{0.337478in}}{\pgfqpoint{1.823275in}{0.341675in}}%
\pgfpathquadraticcurveto{\pgfqpoint{1.819006in}{0.345872in}}{\pgfqpoint{1.819006in}{0.353185in}}%
\pgfpathquadraticcurveto{\pgfqpoint{1.819006in}{0.362083in}}{\pgfqpoint{1.825310in}{0.366910in}}%
\pgfpathquadraticcurveto{\pgfqpoint{1.831614in}{0.371756in}}{\pgfqpoint{1.843214in}{0.371756in}}%
\pgfpathquadraticcurveto{\pgfqpoint{1.848942in}{0.371756in}}{\pgfqpoint{1.854003in}{0.370909in}}%
\pgfpathquadraticcurveto{\pgfqpoint{1.859083in}{0.370080in}}{\pgfqpoint{1.863352in}{0.368387in}}%
\pgfpathlineto{\pgfqpoint{1.863352in}{0.368387in}}%
\pgfpathclose%
\pgfpathmoveto{\pgfqpoint{1.935635in}{0.345260in}}%
\pgfpathlineto{\pgfqpoint{1.935635in}{0.307200in}}%
\pgfpathlineto{\pgfqpoint{1.925278in}{0.307200in}}%
\pgfpathlineto{\pgfqpoint{1.925278in}{0.344917in}}%
\pgfpathquadraticcurveto{\pgfqpoint{1.925278in}{0.353869in}}{\pgfqpoint{1.921783in}{0.358300in}}%
\pgfpathquadraticcurveto{\pgfqpoint{1.918289in}{0.362749in}}{\pgfqpoint{1.911318in}{0.362749in}}%
\pgfpathquadraticcurveto{\pgfqpoint{1.902924in}{0.362749in}}{\pgfqpoint{1.898079in}{0.357400in}}%
\pgfpathquadraticcurveto{\pgfqpoint{1.893234in}{0.352068in}}{\pgfqpoint{1.893234in}{0.342828in}}%
\pgfpathlineto{\pgfqpoint{1.893234in}{0.307200in}}%
\pgfpathlineto{\pgfqpoint{1.882823in}{0.307200in}}%
\pgfpathlineto{\pgfqpoint{1.882823in}{0.394793in}}%
\pgfpathlineto{\pgfqpoint{1.893234in}{0.394793in}}%
\pgfpathlineto{\pgfqpoint{1.893234in}{0.360444in}}%
\pgfpathquadraticcurveto{\pgfqpoint{1.896962in}{0.366136in}}{\pgfqpoint{1.901988in}{0.368946in}}%
\pgfpathquadraticcurveto{\pgfqpoint{1.907031in}{0.371756in}}{\pgfqpoint{1.913624in}{0.371756in}}%
\pgfpathquadraticcurveto{\pgfqpoint{1.924485in}{0.371756in}}{\pgfqpoint{1.930051in}{0.365037in}}%
\pgfpathquadraticcurveto{\pgfqpoint{1.935635in}{0.358318in}}{\pgfqpoint{1.935635in}{0.345260in}}%
\pgfpathlineto{\pgfqpoint{1.935635in}{0.345260in}}%
\pgfpathclose%
\pgfpathmoveto{\pgfqpoint{2.010205in}{0.341315in}}%
\pgfpathlineto{\pgfqpoint{2.010205in}{0.336254in}}%
\pgfpathlineto{\pgfqpoint{1.962581in}{0.336254in}}%
\pgfpathquadraticcurveto{\pgfqpoint{1.963265in}{0.325554in}}{\pgfqpoint{1.969029in}{0.319953in}}%
\pgfpathquadraticcurveto{\pgfqpoint{1.974793in}{0.314351in}}{\pgfqpoint{1.985096in}{0.314351in}}%
\pgfpathquadraticcurveto{\pgfqpoint{1.991058in}{0.314351in}}{\pgfqpoint{1.996660in}{0.315810in}}%
\pgfpathquadraticcurveto{\pgfqpoint{2.002262in}{0.317269in}}{\pgfqpoint{2.007791in}{0.320205in}}%
\pgfpathlineto{\pgfqpoint{2.007791in}{0.310406in}}%
\pgfpathquadraticcurveto{\pgfqpoint{2.002207in}{0.308047in}}{\pgfqpoint{1.996354in}{0.306804in}}%
\pgfpathquadraticcurveto{\pgfqpoint{1.990500in}{0.305561in}}{\pgfqpoint{1.984484in}{0.305561in}}%
\pgfpathquadraticcurveto{\pgfqpoint{1.969389in}{0.305561in}}{\pgfqpoint{1.960581in}{0.314333in}}%
\pgfpathquadraticcurveto{\pgfqpoint{1.951773in}{0.323123in}}{\pgfqpoint{1.951773in}{0.338109in}}%
\pgfpathquadraticcurveto{\pgfqpoint{1.951773in}{0.353581in}}{\pgfqpoint{1.960131in}{0.362659in}}%
\pgfpathquadraticcurveto{\pgfqpoint{1.968489in}{0.371756in}}{\pgfqpoint{1.982682in}{0.371756in}}%
\pgfpathquadraticcurveto{\pgfqpoint{1.995399in}{0.371756in}}{\pgfqpoint{2.002802in}{0.363560in}}%
\pgfpathquadraticcurveto{\pgfqpoint{2.010205in}{0.355383in}}{\pgfqpoint{2.010205in}{0.341315in}}%
\pgfpathlineto{\pgfqpoint{2.010205in}{0.341315in}}%
\pgfpathclose%
\pgfpathmoveto{\pgfqpoint{1.999848in}{0.344359in}}%
\pgfpathquadraticcurveto{\pgfqpoint{1.999740in}{0.352843in}}{\pgfqpoint{1.995093in}{0.357904in}}%
\pgfpathquadraticcurveto{\pgfqpoint{1.990446in}{0.362984in}}{\pgfqpoint{1.982790in}{0.362984in}}%
\pgfpathquadraticcurveto{\pgfqpoint{1.974127in}{0.362984in}}{\pgfqpoint{1.968921in}{0.358084in}}%
\pgfpathquadraticcurveto{\pgfqpoint{1.963715in}{0.353185in}}{\pgfqpoint{1.962923in}{0.344287in}}%
\pgfpathlineto{\pgfqpoint{1.999848in}{0.344359in}}%
\pgfpathlineto{\pgfqpoint{1.999848in}{0.344359in}}%
\pgfpathclose%
\pgfpathmoveto{\pgfqpoint{2.068690in}{0.360678in}}%
\pgfpathlineto{\pgfqpoint{2.068690in}{0.394793in}}%
\pgfpathlineto{\pgfqpoint{2.079047in}{0.394793in}}%
\pgfpathlineto{\pgfqpoint{2.079047in}{0.307200in}}%
\pgfpathlineto{\pgfqpoint{2.068690in}{0.307200in}}%
\pgfpathlineto{\pgfqpoint{2.068690in}{0.316656in}}%
\pgfpathquadraticcurveto{\pgfqpoint{2.065430in}{0.311037in}}{\pgfqpoint{2.060441in}{0.308299in}}%
\pgfpathquadraticcurveto{\pgfqpoint{2.055469in}{0.305561in}}{\pgfqpoint{2.048481in}{0.305561in}}%
\pgfpathquadraticcurveto{\pgfqpoint{2.037061in}{0.305561in}}{\pgfqpoint{2.029874in}{0.314675in}}%
\pgfpathquadraticcurveto{\pgfqpoint{2.022705in}{0.323807in}}{\pgfqpoint{2.022705in}{0.338667in}}%
\pgfpathquadraticcurveto{\pgfqpoint{2.022705in}{0.353527in}}{\pgfqpoint{2.029874in}{0.362641in}}%
\pgfpathquadraticcurveto{\pgfqpoint{2.037061in}{0.371756in}}{\pgfqpoint{2.048481in}{0.371756in}}%
\pgfpathquadraticcurveto{\pgfqpoint{2.055469in}{0.371756in}}{\pgfqpoint{2.060441in}{0.369018in}}%
\pgfpathquadraticcurveto{\pgfqpoint{2.065430in}{0.366298in}}{\pgfqpoint{2.068690in}{0.360678in}}%
\pgfpathlineto{\pgfqpoint{2.068690in}{0.360678in}}%
\pgfpathclose%
\pgfpathmoveto{\pgfqpoint{2.033405in}{0.338667in}}%
\pgfpathquadraticcurveto{\pgfqpoint{2.033405in}{0.327248in}}{\pgfqpoint{2.038106in}{0.320745in}}%
\pgfpathquadraticcurveto{\pgfqpoint{2.042807in}{0.314243in}}{\pgfqpoint{2.051020in}{0.314243in}}%
\pgfpathquadraticcurveto{\pgfqpoint{2.059234in}{0.314243in}}{\pgfqpoint{2.063953in}{0.320745in}}%
\pgfpathquadraticcurveto{\pgfqpoint{2.068690in}{0.327248in}}{\pgfqpoint{2.068690in}{0.338667in}}%
\pgfpathquadraticcurveto{\pgfqpoint{2.068690in}{0.350087in}}{\pgfqpoint{2.063953in}{0.356589in}}%
\pgfpathquadraticcurveto{\pgfqpoint{2.059234in}{0.363092in}}{\pgfqpoint{2.051020in}{0.363092in}}%
\pgfpathquadraticcurveto{\pgfqpoint{2.042807in}{0.363092in}}{\pgfqpoint{2.038106in}{0.356589in}}%
\pgfpathquadraticcurveto{\pgfqpoint{2.033405in}{0.350087in}}{\pgfqpoint{2.033405in}{0.338667in}}%
\pgfpathlineto{\pgfqpoint{2.033405in}{0.338667in}}%
\pgfpathclose%
\pgfpathmoveto{\pgfqpoint{2.137028in}{0.394793in}}%
\pgfpathlineto{\pgfqpoint{2.147385in}{0.394793in}}%
\pgfpathlineto{\pgfqpoint{2.147385in}{0.307200in}}%
\pgfpathlineto{\pgfqpoint{2.137028in}{0.307200in}}%
\pgfpathlineto{\pgfqpoint{2.137028in}{0.394793in}}%
\pgfpathlineto{\pgfqpoint{2.137028in}{0.394793in}}%
\pgfpathclose%
\pgfpathmoveto{\pgfqpoint{2.169054in}{0.370243in}}%
\pgfpathlineto{\pgfqpoint{2.179411in}{0.370243in}}%
\pgfpathlineto{\pgfqpoint{2.179411in}{0.307200in}}%
\pgfpathlineto{\pgfqpoint{2.169054in}{0.307200in}}%
\pgfpathlineto{\pgfqpoint{2.169054in}{0.370243in}}%
\pgfpathlineto{\pgfqpoint{2.169054in}{0.370243in}}%
\pgfpathclose%
\pgfpathmoveto{\pgfqpoint{2.169054in}{0.394793in}}%
\pgfpathlineto{\pgfqpoint{2.179411in}{0.394793in}}%
\pgfpathlineto{\pgfqpoint{2.179411in}{0.381662in}}%
\pgfpathlineto{\pgfqpoint{2.169054in}{0.381662in}}%
\pgfpathlineto{\pgfqpoint{2.169054in}{0.394793in}}%
\pgfpathlineto{\pgfqpoint{2.169054in}{0.394793in}}%
\pgfpathclose%
\pgfpathmoveto{\pgfqpoint{2.253495in}{0.345260in}}%
\pgfpathlineto{\pgfqpoint{2.253495in}{0.307200in}}%
\pgfpathlineto{\pgfqpoint{2.243138in}{0.307200in}}%
\pgfpathlineto{\pgfqpoint{2.243138in}{0.344917in}}%
\pgfpathquadraticcurveto{\pgfqpoint{2.243138in}{0.353869in}}{\pgfqpoint{2.239644in}{0.358300in}}%
\pgfpathquadraticcurveto{\pgfqpoint{2.236149in}{0.362749in}}{\pgfqpoint{2.229179in}{0.362749in}}%
\pgfpathquadraticcurveto{\pgfqpoint{2.220785in}{0.362749in}}{\pgfqpoint{2.215940in}{0.357400in}}%
\pgfpathquadraticcurveto{\pgfqpoint{2.211094in}{0.352068in}}{\pgfqpoint{2.211094in}{0.342828in}}%
\pgfpathlineto{\pgfqpoint{2.211094in}{0.307200in}}%
\pgfpathlineto{\pgfqpoint{2.200683in}{0.307200in}}%
\pgfpathlineto{\pgfqpoint{2.200683in}{0.370243in}}%
\pgfpathlineto{\pgfqpoint{2.211094in}{0.370243in}}%
\pgfpathlineto{\pgfqpoint{2.211094in}{0.360444in}}%
\pgfpathquadraticcurveto{\pgfqpoint{2.214823in}{0.366136in}}{\pgfqpoint{2.219848in}{0.368946in}}%
\pgfpathquadraticcurveto{\pgfqpoint{2.224892in}{0.371756in}}{\pgfqpoint{2.231484in}{0.371756in}}%
\pgfpathquadraticcurveto{\pgfqpoint{2.242345in}{0.371756in}}{\pgfqpoint{2.247911in}{0.365037in}}%
\pgfpathquadraticcurveto{\pgfqpoint{2.253495in}{0.358318in}}{\pgfqpoint{2.253495in}{0.345260in}}%
\pgfpathlineto{\pgfqpoint{2.253495in}{0.345260in}}%
\pgfpathclose%
\pgfpathmoveto{\pgfqpoint{2.328065in}{0.341315in}}%
\pgfpathlineto{\pgfqpoint{2.328065in}{0.336254in}}%
\pgfpathlineto{\pgfqpoint{2.280441in}{0.336254in}}%
\pgfpathquadraticcurveto{\pgfqpoint{2.281126in}{0.325554in}}{\pgfqpoint{2.286890in}{0.319953in}}%
\pgfpathquadraticcurveto{\pgfqpoint{2.292653in}{0.314351in}}{\pgfqpoint{2.302956in}{0.314351in}}%
\pgfpathquadraticcurveto{\pgfqpoint{2.308918in}{0.314351in}}{\pgfqpoint{2.314520in}{0.315810in}}%
\pgfpathquadraticcurveto{\pgfqpoint{2.320122in}{0.317269in}}{\pgfqpoint{2.325652in}{0.320205in}}%
\pgfpathlineto{\pgfqpoint{2.325652in}{0.310406in}}%
\pgfpathquadraticcurveto{\pgfqpoint{2.320068in}{0.308047in}}{\pgfqpoint{2.314214in}{0.306804in}}%
\pgfpathquadraticcurveto{\pgfqpoint{2.308360in}{0.305561in}}{\pgfqpoint{2.302344in}{0.305561in}}%
\pgfpathquadraticcurveto{\pgfqpoint{2.287250in}{0.305561in}}{\pgfqpoint{2.278442in}{0.314333in}}%
\pgfpathquadraticcurveto{\pgfqpoint{2.269634in}{0.323123in}}{\pgfqpoint{2.269634in}{0.338109in}}%
\pgfpathquadraticcurveto{\pgfqpoint{2.269634in}{0.353581in}}{\pgfqpoint{2.277992in}{0.362659in}}%
\pgfpathquadraticcurveto{\pgfqpoint{2.286349in}{0.371756in}}{\pgfqpoint{2.300543in}{0.371756in}}%
\pgfpathquadraticcurveto{\pgfqpoint{2.313259in}{0.371756in}}{\pgfqpoint{2.320662in}{0.363560in}}%
\pgfpathquadraticcurveto{\pgfqpoint{2.328065in}{0.355383in}}{\pgfqpoint{2.328065in}{0.341315in}}%
\pgfpathlineto{\pgfqpoint{2.328065in}{0.341315in}}%
\pgfpathclose%
\pgfpathmoveto{\pgfqpoint{2.317708in}{0.344359in}}%
\pgfpathquadraticcurveto{\pgfqpoint{2.317600in}{0.352843in}}{\pgfqpoint{2.312953in}{0.357904in}}%
\pgfpathquadraticcurveto{\pgfqpoint{2.308306in}{0.362984in}}{\pgfqpoint{2.300651in}{0.362984in}}%
\pgfpathquadraticcurveto{\pgfqpoint{2.291987in}{0.362984in}}{\pgfqpoint{2.286781in}{0.358084in}}%
\pgfpathquadraticcurveto{\pgfqpoint{2.281576in}{0.353185in}}{\pgfqpoint{2.280783in}{0.344287in}}%
\pgfpathlineto{\pgfqpoint{2.317708in}{0.344359in}}%
\pgfpathlineto{\pgfqpoint{2.317708in}{0.344359in}}%
\pgfpathclose%
\pgfpathmoveto{\pgfqpoint{2.347717in}{0.321502in}}%
\pgfpathlineto{\pgfqpoint{2.359587in}{0.321502in}}%
\pgfpathlineto{\pgfqpoint{2.359587in}{0.307200in}}%
\pgfpathlineto{\pgfqpoint{2.347717in}{0.307200in}}%
\pgfpathlineto{\pgfqpoint{2.347717in}{0.321502in}}%
\pgfpathlineto{\pgfqpoint{2.347717in}{0.321502in}}%
\pgfpathclose%
\pgfpathmoveto{\pgfqpoint{2.347717in}{0.366802in}}%
\pgfpathlineto{\pgfqpoint{2.359587in}{0.366802in}}%
\pgfpathlineto{\pgfqpoint{2.359587in}{0.352519in}}%
\pgfpathlineto{\pgfqpoint{2.347717in}{0.352519in}}%
\pgfpathlineto{\pgfqpoint{2.347717in}{0.366802in}}%
\pgfpathlineto{\pgfqpoint{2.347717in}{0.366802in}}%
\pgfpathclose%
\pgfpathmoveto{\pgfqpoint{2.474468in}{0.341315in}}%
\pgfpathlineto{\pgfqpoint{2.474468in}{0.336254in}}%
\pgfpathlineto{\pgfqpoint{2.426844in}{0.336254in}}%
\pgfpathquadraticcurveto{\pgfqpoint{2.427528in}{0.325554in}}{\pgfqpoint{2.433292in}{0.319953in}}%
\pgfpathquadraticcurveto{\pgfqpoint{2.439056in}{0.314351in}}{\pgfqpoint{2.449359in}{0.314351in}}%
\pgfpathquadraticcurveto{\pgfqpoint{2.455321in}{0.314351in}}{\pgfqpoint{2.460923in}{0.315810in}}%
\pgfpathquadraticcurveto{\pgfqpoint{2.466525in}{0.317269in}}{\pgfqpoint{2.472054in}{0.320205in}}%
\pgfpathlineto{\pgfqpoint{2.472054in}{0.310406in}}%
\pgfpathquadraticcurveto{\pgfqpoint{2.466471in}{0.308047in}}{\pgfqpoint{2.460617in}{0.306804in}}%
\pgfpathquadraticcurveto{\pgfqpoint{2.454763in}{0.305561in}}{\pgfqpoint{2.448747in}{0.305561in}}%
\pgfpathquadraticcurveto{\pgfqpoint{2.433653in}{0.305561in}}{\pgfqpoint{2.424845in}{0.314333in}}%
\pgfpathquadraticcurveto{\pgfqpoint{2.416037in}{0.323123in}}{\pgfqpoint{2.416037in}{0.338109in}}%
\pgfpathquadraticcurveto{\pgfqpoint{2.416037in}{0.353581in}}{\pgfqpoint{2.424394in}{0.362659in}}%
\pgfpathquadraticcurveto{\pgfqpoint{2.432752in}{0.371756in}}{\pgfqpoint{2.446946in}{0.371756in}}%
\pgfpathquadraticcurveto{\pgfqpoint{2.459662in}{0.371756in}}{\pgfqpoint{2.467065in}{0.363560in}}%
\pgfpathquadraticcurveto{\pgfqpoint{2.474468in}{0.355383in}}{\pgfqpoint{2.474468in}{0.341315in}}%
\pgfpathlineto{\pgfqpoint{2.474468in}{0.341315in}}%
\pgfpathclose%
\pgfpathmoveto{\pgfqpoint{2.464111in}{0.344359in}}%
\pgfpathquadraticcurveto{\pgfqpoint{2.464003in}{0.352843in}}{\pgfqpoint{2.459356in}{0.357904in}}%
\pgfpathquadraticcurveto{\pgfqpoint{2.454709in}{0.362984in}}{\pgfqpoint{2.447054in}{0.362984in}}%
\pgfpathquadraticcurveto{\pgfqpoint{2.438390in}{0.362984in}}{\pgfqpoint{2.433184in}{0.358084in}}%
\pgfpathquadraticcurveto{\pgfqpoint{2.427979in}{0.353185in}}{\pgfqpoint{2.427186in}{0.344287in}}%
\pgfpathlineto{\pgfqpoint{2.464111in}{0.344359in}}%
\pgfpathlineto{\pgfqpoint{2.464111in}{0.344359in}}%
\pgfpathclose%
\pgfpathmoveto{\pgfqpoint{2.497668in}{0.338667in}}%
\pgfpathquadraticcurveto{\pgfqpoint{2.497668in}{0.327248in}}{\pgfqpoint{2.502369in}{0.320745in}}%
\pgfpathquadraticcurveto{\pgfqpoint{2.507070in}{0.314243in}}{\pgfqpoint{2.515284in}{0.314243in}}%
\pgfpathquadraticcurveto{\pgfqpoint{2.523497in}{0.314243in}}{\pgfqpoint{2.528216in}{0.320745in}}%
\pgfpathquadraticcurveto{\pgfqpoint{2.532954in}{0.327248in}}{\pgfqpoint{2.532954in}{0.338667in}}%
\pgfpathquadraticcurveto{\pgfqpoint{2.532954in}{0.350087in}}{\pgfqpoint{2.528216in}{0.356589in}}%
\pgfpathquadraticcurveto{\pgfqpoint{2.523497in}{0.363092in}}{\pgfqpoint{2.515284in}{0.363092in}}%
\pgfpathquadraticcurveto{\pgfqpoint{2.507070in}{0.363092in}}{\pgfqpoint{2.502369in}{0.356589in}}%
\pgfpathquadraticcurveto{\pgfqpoint{2.497668in}{0.350087in}}{\pgfqpoint{2.497668in}{0.338667in}}%
\pgfpathlineto{\pgfqpoint{2.497668in}{0.338667in}}%
\pgfpathclose%
\pgfpathmoveto{\pgfqpoint{2.532954in}{0.316656in}}%
\pgfpathquadraticcurveto{\pgfqpoint{2.529693in}{0.311037in}}{\pgfqpoint{2.524704in}{0.308299in}}%
\pgfpathquadraticcurveto{\pgfqpoint{2.519733in}{0.305561in}}{\pgfqpoint{2.512744in}{0.305561in}}%
\pgfpathquadraticcurveto{\pgfqpoint{2.501324in}{0.305561in}}{\pgfqpoint{2.494137in}{0.314675in}}%
\pgfpathquadraticcurveto{\pgfqpoint{2.486969in}{0.323807in}}{\pgfqpoint{2.486969in}{0.338667in}}%
\pgfpathquadraticcurveto{\pgfqpoint{2.486969in}{0.353527in}}{\pgfqpoint{2.494137in}{0.362641in}}%
\pgfpathquadraticcurveto{\pgfqpoint{2.501324in}{0.371756in}}{\pgfqpoint{2.512744in}{0.371756in}}%
\pgfpathquadraticcurveto{\pgfqpoint{2.519733in}{0.371756in}}{\pgfqpoint{2.524704in}{0.369018in}}%
\pgfpathquadraticcurveto{\pgfqpoint{2.529693in}{0.366298in}}{\pgfqpoint{2.532954in}{0.360678in}}%
\pgfpathlineto{\pgfqpoint{2.532954in}{0.370243in}}%
\pgfpathlineto{\pgfqpoint{2.543311in}{0.370243in}}%
\pgfpathlineto{\pgfqpoint{2.543311in}{0.283226in}}%
\pgfpathlineto{\pgfqpoint{2.532954in}{0.283226in}}%
\pgfpathlineto{\pgfqpoint{2.532954in}{0.316656in}}%
\pgfpathlineto{\pgfqpoint{2.532954in}{0.316656in}}%
\pgfpathclose%
\pgfpathmoveto{\pgfqpoint{2.563592in}{0.332075in}}%
\pgfpathlineto{\pgfqpoint{2.563592in}{0.370243in}}%
\pgfpathlineto{\pgfqpoint{2.573949in}{0.370243in}}%
\pgfpathlineto{\pgfqpoint{2.573949in}{0.332471in}}%
\pgfpathquadraticcurveto{\pgfqpoint{2.573949in}{0.323519in}}{\pgfqpoint{2.577426in}{0.319034in}}%
\pgfpathquadraticcurveto{\pgfqpoint{2.580920in}{0.314567in}}{\pgfqpoint{2.587909in}{0.314567in}}%
\pgfpathquadraticcurveto{\pgfqpoint{2.596284in}{0.314567in}}{\pgfqpoint{2.601148in}{0.319917in}}%
\pgfpathquadraticcurveto{\pgfqpoint{2.606029in}{0.325266in}}{\pgfqpoint{2.606029in}{0.334506in}}%
\pgfpathlineto{\pgfqpoint{2.606029in}{0.370243in}}%
\pgfpathlineto{\pgfqpoint{2.616386in}{0.370243in}}%
\pgfpathlineto{\pgfqpoint{2.616386in}{0.307200in}}%
\pgfpathlineto{\pgfqpoint{2.606029in}{0.307200in}}%
\pgfpathlineto{\pgfqpoint{2.606029in}{0.316891in}}%
\pgfpathquadraticcurveto{\pgfqpoint{2.602264in}{0.311145in}}{\pgfqpoint{2.597275in}{0.308353in}}%
\pgfpathquadraticcurveto{\pgfqpoint{2.592304in}{0.305561in}}{\pgfqpoint{2.585711in}{0.305561in}}%
\pgfpathquadraticcurveto{\pgfqpoint{2.574850in}{0.305561in}}{\pgfqpoint{2.569212in}{0.312315in}}%
\pgfpathquadraticcurveto{\pgfqpoint{2.563592in}{0.319070in}}{\pgfqpoint{2.563592in}{0.332075in}}%
\pgfpathlineto{\pgfqpoint{2.563592in}{0.332075in}}%
\pgfpathclose%
\pgfpathmoveto{\pgfqpoint{2.589656in}{0.371756in}}%
\pgfpathlineto{\pgfqpoint{2.589656in}{0.371756in}}%
\pgfpathlineto{\pgfqpoint{2.589656in}{0.371756in}}%
\pgfpathclose%
\pgfpathmoveto{\pgfqpoint{2.666370in}{0.338883in}}%
\pgfpathquadraticcurveto{\pgfqpoint{2.653815in}{0.338883in}}{\pgfqpoint{2.648970in}{0.336019in}}%
\pgfpathquadraticcurveto{\pgfqpoint{2.644125in}{0.333156in}}{\pgfqpoint{2.644125in}{0.326221in}}%
\pgfpathquadraticcurveto{\pgfqpoint{2.644125in}{0.320709in}}{\pgfqpoint{2.647763in}{0.317467in}}%
\pgfpathquadraticcurveto{\pgfqpoint{2.651401in}{0.314243in}}{\pgfqpoint{2.657634in}{0.314243in}}%
\pgfpathquadraticcurveto{\pgfqpoint{2.666262in}{0.314243in}}{\pgfqpoint{2.671467in}{0.320349in}}%
\pgfpathquadraticcurveto{\pgfqpoint{2.676673in}{0.326455in}}{\pgfqpoint{2.676673in}{0.336578in}}%
\pgfpathlineto{\pgfqpoint{2.676673in}{0.338883in}}%
\pgfpathlineto{\pgfqpoint{2.666370in}{0.338883in}}%
\pgfpathlineto{\pgfqpoint{2.666370in}{0.338883in}}%
\pgfpathclose%
\pgfpathmoveto{\pgfqpoint{2.687030in}{0.343170in}}%
\pgfpathlineto{\pgfqpoint{2.687030in}{0.307200in}}%
\pgfpathlineto{\pgfqpoint{2.676673in}{0.307200in}}%
\pgfpathlineto{\pgfqpoint{2.676673in}{0.316764in}}%
\pgfpathquadraticcurveto{\pgfqpoint{2.673124in}{0.311037in}}{\pgfqpoint{2.667829in}{0.308299in}}%
\pgfpathquadraticcurveto{\pgfqpoint{2.662533in}{0.305561in}}{\pgfqpoint{2.654878in}{0.305561in}}%
\pgfpathquadraticcurveto{\pgfqpoint{2.645205in}{0.305561in}}{\pgfqpoint{2.639477in}{0.311001in}}%
\pgfpathquadraticcurveto{\pgfqpoint{2.633768in}{0.316440in}}{\pgfqpoint{2.633768in}{0.325554in}}%
\pgfpathquadraticcurveto{\pgfqpoint{2.633768in}{0.336182in}}{\pgfqpoint{2.640882in}{0.341585in}}%
\pgfpathquadraticcurveto{\pgfqpoint{2.648015in}{0.346989in}}{\pgfqpoint{2.662137in}{0.346989in}}%
\pgfpathlineto{\pgfqpoint{2.676673in}{0.346989in}}%
\pgfpathlineto{\pgfqpoint{2.676673in}{0.348016in}}%
\pgfpathquadraticcurveto{\pgfqpoint{2.676673in}{0.355166in}}{\pgfqpoint{2.671971in}{0.359075in}}%
\pgfpathquadraticcurveto{\pgfqpoint{2.667270in}{0.362984in}}{\pgfqpoint{2.658768in}{0.362984in}}%
\pgfpathquadraticcurveto{\pgfqpoint{2.653365in}{0.362984in}}{\pgfqpoint{2.648231in}{0.361687in}}%
\pgfpathquadraticcurveto{\pgfqpoint{2.643116in}{0.360390in}}{\pgfqpoint{2.638397in}{0.357796in}}%
\pgfpathlineto{\pgfqpoint{2.638397in}{0.367379in}}%
\pgfpathquadraticcurveto{\pgfqpoint{2.644071in}{0.369576in}}{\pgfqpoint{2.649420in}{0.370657in}}%
\pgfpathquadraticcurveto{\pgfqpoint{2.654770in}{0.371756in}}{\pgfqpoint{2.659831in}{0.371756in}}%
\pgfpathquadraticcurveto{\pgfqpoint{2.673520in}{0.371756in}}{\pgfqpoint{2.680275in}{0.364659in}}%
\pgfpathquadraticcurveto{\pgfqpoint{2.687030in}{0.357580in}}{\pgfqpoint{2.687030in}{0.343170in}}%
\pgfpathlineto{\pgfqpoint{2.687030in}{0.343170in}}%
\pgfpathclose%
\pgfpathmoveto{\pgfqpoint{2.708356in}{0.394793in}}%
\pgfpathlineto{\pgfqpoint{2.718713in}{0.394793in}}%
\pgfpathlineto{\pgfqpoint{2.718713in}{0.307200in}}%
\pgfpathlineto{\pgfqpoint{2.708356in}{0.307200in}}%
\pgfpathlineto{\pgfqpoint{2.708356in}{0.394793in}}%
\pgfpathlineto{\pgfqpoint{2.708356in}{0.394793in}}%
\pgfpathclose%
\pgfpathmoveto{\pgfqpoint{2.740382in}{0.370243in}}%
\pgfpathlineto{\pgfqpoint{2.750738in}{0.370243in}}%
\pgfpathlineto{\pgfqpoint{2.750738in}{0.307200in}}%
\pgfpathlineto{\pgfqpoint{2.740382in}{0.307200in}}%
\pgfpathlineto{\pgfqpoint{2.740382in}{0.370243in}}%
\pgfpathlineto{\pgfqpoint{2.740382in}{0.370243in}}%
\pgfpathclose%
\pgfpathmoveto{\pgfqpoint{2.740382in}{0.394793in}}%
\pgfpathlineto{\pgfqpoint{2.750738in}{0.394793in}}%
\pgfpathlineto{\pgfqpoint{2.750738in}{0.381662in}}%
\pgfpathlineto{\pgfqpoint{2.740382in}{0.381662in}}%
\pgfpathlineto{\pgfqpoint{2.740382in}{0.394793in}}%
\pgfpathlineto{\pgfqpoint{2.740382in}{0.394793in}}%
\pgfpathclose%
\pgfpathmoveto{\pgfqpoint{2.782656in}{0.388147in}}%
\pgfpathlineto{\pgfqpoint{2.782656in}{0.370243in}}%
\pgfpathlineto{\pgfqpoint{2.803982in}{0.370243in}}%
\pgfpathlineto{\pgfqpoint{2.803982in}{0.362191in}}%
\pgfpathlineto{\pgfqpoint{2.782656in}{0.362191in}}%
\pgfpathlineto{\pgfqpoint{2.782656in}{0.327968in}}%
\pgfpathquadraticcurveto{\pgfqpoint{2.782656in}{0.320259in}}{\pgfqpoint{2.784763in}{0.318061in}}%
\pgfpathquadraticcurveto{\pgfqpoint{2.786871in}{0.315864in}}{\pgfqpoint{2.793355in}{0.315864in}}%
\pgfpathlineto{\pgfqpoint{2.803982in}{0.315864in}}%
\pgfpathlineto{\pgfqpoint{2.803982in}{0.307200in}}%
\pgfpathlineto{\pgfqpoint{2.793355in}{0.307200in}}%
\pgfpathquadraticcurveto{\pgfqpoint{2.781359in}{0.307200in}}{\pgfqpoint{2.776802in}{0.311667in}}%
\pgfpathquadraticcurveto{\pgfqpoint{2.772245in}{0.316152in}}{\pgfqpoint{2.772245in}{0.327968in}}%
\pgfpathlineto{\pgfqpoint{2.772245in}{0.362191in}}%
\pgfpathlineto{\pgfqpoint{2.764644in}{0.362191in}}%
\pgfpathlineto{\pgfqpoint{2.764644in}{0.370243in}}%
\pgfpathlineto{\pgfqpoint{2.772245in}{0.370243in}}%
\pgfpathlineto{\pgfqpoint{2.772245in}{0.388147in}}%
\pgfpathlineto{\pgfqpoint{2.782656in}{0.388147in}}%
\pgfpathlineto{\pgfqpoint{2.782656in}{0.388147in}}%
\pgfpathclose%
\pgfpathmoveto{\pgfqpoint{2.843825in}{0.301346in}}%
\pgfpathquadraticcurveto{\pgfqpoint{2.839448in}{0.290088in}}{\pgfqpoint{2.835270in}{0.286666in}}%
\pgfpathquadraticcurveto{\pgfqpoint{2.831109in}{0.283226in}}{\pgfqpoint{2.824138in}{0.283226in}}%
\pgfpathlineto{\pgfqpoint{2.815852in}{0.283226in}}%
\pgfpathlineto{\pgfqpoint{2.815852in}{0.291890in}}%
\pgfpathlineto{\pgfqpoint{2.821941in}{0.291890in}}%
\pgfpathquadraticcurveto{\pgfqpoint{2.826209in}{0.291890in}}{\pgfqpoint{2.828569in}{0.293925in}}%
\pgfpathquadraticcurveto{\pgfqpoint{2.830947in}{0.295942in}}{\pgfqpoint{2.833811in}{0.303489in}}%
\pgfpathlineto{\pgfqpoint{2.835666in}{0.308209in}}%
\pgfpathlineto{\pgfqpoint{2.810179in}{0.370243in}}%
\pgfpathlineto{\pgfqpoint{2.821148in}{0.370243in}}%
\pgfpathlineto{\pgfqpoint{2.840853in}{0.320943in}}%
\pgfpathlineto{\pgfqpoint{2.860559in}{0.370243in}}%
\pgfpathlineto{\pgfqpoint{2.871528in}{0.370243in}}%
\pgfpathlineto{\pgfqpoint{2.843825in}{0.301346in}}%
\pgfpathlineto{\pgfqpoint{2.843825in}{0.301346in}}%
\pgfpathclose%
\pgfpathmoveto{\pgfqpoint{2.870862in}{0.321502in}}%
\pgfpathlineto{\pgfqpoint{2.882750in}{0.321502in}}%
\pgfpathlineto{\pgfqpoint{2.882750in}{0.307200in}}%
\pgfpathlineto{\pgfqpoint{2.870862in}{0.307200in}}%
\pgfpathlineto{\pgfqpoint{2.870862in}{0.321502in}}%
\pgfpathlineto{\pgfqpoint{2.870862in}{0.321502in}}%
\pgfpathclose%
\pgfpathmoveto{\pgfqpoint{2.993506in}{0.388489in}}%
\pgfpathlineto{\pgfqpoint{2.993506in}{0.377393in}}%
\pgfpathquadraticcurveto{\pgfqpoint{2.987040in}{0.380491in}}{\pgfqpoint{2.981294in}{0.382004in}}%
\pgfpathquadraticcurveto{\pgfqpoint{2.975548in}{0.383536in}}{\pgfqpoint{2.970199in}{0.383536in}}%
\pgfpathquadraticcurveto{\pgfqpoint{2.960922in}{0.383536in}}{\pgfqpoint{2.955879in}{0.379933in}}%
\pgfpathquadraticcurveto{\pgfqpoint{2.950836in}{0.376331in}}{\pgfqpoint{2.950836in}{0.369684in}}%
\pgfpathquadraticcurveto{\pgfqpoint{2.950836in}{0.364100in}}{\pgfqpoint{2.954186in}{0.361254in}}%
\pgfpathquadraticcurveto{\pgfqpoint{2.957536in}{0.358427in}}{\pgfqpoint{2.966884in}{0.356679in}}%
\pgfpathlineto{\pgfqpoint{2.973747in}{0.355274in}}%
\pgfpathquadraticcurveto{\pgfqpoint{2.986464in}{0.352843in}}{\pgfqpoint{2.992516in}{0.346737in}}%
\pgfpathquadraticcurveto{\pgfqpoint{2.998568in}{0.340631in}}{\pgfqpoint{2.998568in}{0.330400in}}%
\pgfpathquadraticcurveto{\pgfqpoint{2.998568in}{0.318169in}}{\pgfqpoint{2.990372in}{0.311865in}}%
\pgfpathquadraticcurveto{\pgfqpoint{2.982195in}{0.305561in}}{\pgfqpoint{2.966380in}{0.305561in}}%
\pgfpathquadraticcurveto{\pgfqpoint{2.960418in}{0.305561in}}{\pgfqpoint{2.953681in}{0.306912in}}%
\pgfpathquadraticcurveto{\pgfqpoint{2.946963in}{0.308263in}}{\pgfqpoint{2.939758in}{0.310911in}}%
\pgfpathlineto{\pgfqpoint{2.939758in}{0.322618in}}%
\pgfpathquadraticcurveto{\pgfqpoint{2.946675in}{0.318746in}}{\pgfqpoint{2.953321in}{0.316764in}}%
\pgfpathquadraticcurveto{\pgfqpoint{2.959968in}{0.314801in}}{\pgfqpoint{2.966380in}{0.314801in}}%
\pgfpathquadraticcurveto{\pgfqpoint{2.976107in}{0.314801in}}{\pgfqpoint{2.981402in}{0.318620in}}%
\pgfpathquadraticcurveto{\pgfqpoint{2.986698in}{0.322456in}}{\pgfqpoint{2.986698in}{0.329553in}}%
\pgfpathquadraticcurveto{\pgfqpoint{2.986698in}{0.335731in}}{\pgfqpoint{2.982897in}{0.339226in}}%
\pgfpathquadraticcurveto{\pgfqpoint{2.979097in}{0.342720in}}{\pgfqpoint{2.970433in}{0.344467in}}%
\pgfpathlineto{\pgfqpoint{2.963498in}{0.345818in}}%
\pgfpathquadraticcurveto{\pgfqpoint{2.950781in}{0.348340in}}{\pgfqpoint{2.945090in}{0.353743in}}%
\pgfpathquadraticcurveto{\pgfqpoint{2.939416in}{0.359147in}}{\pgfqpoint{2.939416in}{0.368784in}}%
\pgfpathquadraticcurveto{\pgfqpoint{2.939416in}{0.379933in}}{\pgfqpoint{2.947269in}{0.386345in}}%
\pgfpathquadraticcurveto{\pgfqpoint{2.955122in}{0.392758in}}{\pgfqpoint{2.968902in}{0.392758in}}%
\pgfpathquadraticcurveto{\pgfqpoint{2.974828in}{0.392758in}}{\pgfqpoint{2.980952in}{0.391677in}}%
\pgfpathquadraticcurveto{\pgfqpoint{2.987094in}{0.390614in}}{\pgfqpoint{2.993506in}{0.388489in}}%
\pgfpathlineto{\pgfqpoint{2.993506in}{0.388489in}}%
\pgfpathclose%
\pgfpathmoveto{\pgfqpoint{3.026108in}{0.388147in}}%
\pgfpathlineto{\pgfqpoint{3.026108in}{0.370243in}}%
\pgfpathlineto{\pgfqpoint{3.047435in}{0.370243in}}%
\pgfpathlineto{\pgfqpoint{3.047435in}{0.362191in}}%
\pgfpathlineto{\pgfqpoint{3.026108in}{0.362191in}}%
\pgfpathlineto{\pgfqpoint{3.026108in}{0.327968in}}%
\pgfpathquadraticcurveto{\pgfqpoint{3.026108in}{0.320259in}}{\pgfqpoint{3.028216in}{0.318061in}}%
\pgfpathquadraticcurveto{\pgfqpoint{3.030323in}{0.315864in}}{\pgfqpoint{3.036808in}{0.315864in}}%
\pgfpathlineto{\pgfqpoint{3.047435in}{0.315864in}}%
\pgfpathlineto{\pgfqpoint{3.047435in}{0.307200in}}%
\pgfpathlineto{\pgfqpoint{3.036808in}{0.307200in}}%
\pgfpathquadraticcurveto{\pgfqpoint{3.024811in}{0.307200in}}{\pgfqpoint{3.020254in}{0.311667in}}%
\pgfpathquadraticcurveto{\pgfqpoint{3.015697in}{0.316152in}}{\pgfqpoint{3.015697in}{0.327968in}}%
\pgfpathlineto{\pgfqpoint{3.015697in}{0.362191in}}%
\pgfpathlineto{\pgfqpoint{3.008096in}{0.362191in}}%
\pgfpathlineto{\pgfqpoint{3.008096in}{0.370243in}}%
\pgfpathlineto{\pgfqpoint{3.015697in}{0.370243in}}%
\pgfpathlineto{\pgfqpoint{3.015697in}{0.388147in}}%
\pgfpathlineto{\pgfqpoint{3.026108in}{0.388147in}}%
\pgfpathlineto{\pgfqpoint{3.026108in}{0.388147in}}%
\pgfpathclose%
\pgfpathmoveto{\pgfqpoint{3.059989in}{0.332075in}}%
\pgfpathlineto{\pgfqpoint{3.059989in}{0.370243in}}%
\pgfpathlineto{\pgfqpoint{3.070346in}{0.370243in}}%
\pgfpathlineto{\pgfqpoint{3.070346in}{0.332471in}}%
\pgfpathquadraticcurveto{\pgfqpoint{3.070346in}{0.323519in}}{\pgfqpoint{3.073822in}{0.319034in}}%
\pgfpathquadraticcurveto{\pgfqpoint{3.077317in}{0.314567in}}{\pgfqpoint{3.084306in}{0.314567in}}%
\pgfpathquadraticcurveto{\pgfqpoint{3.092681in}{0.314567in}}{\pgfqpoint{3.097544in}{0.319917in}}%
\pgfpathquadraticcurveto{\pgfqpoint{3.102426in}{0.325266in}}{\pgfqpoint{3.102426in}{0.334506in}}%
\pgfpathlineto{\pgfqpoint{3.102426in}{0.370243in}}%
\pgfpathlineto{\pgfqpoint{3.112783in}{0.370243in}}%
\pgfpathlineto{\pgfqpoint{3.112783in}{0.307200in}}%
\pgfpathlineto{\pgfqpoint{3.102426in}{0.307200in}}%
\pgfpathlineto{\pgfqpoint{3.102426in}{0.316891in}}%
\pgfpathquadraticcurveto{\pgfqpoint{3.098661in}{0.311145in}}{\pgfqpoint{3.093672in}{0.308353in}}%
\pgfpathquadraticcurveto{\pgfqpoint{3.088701in}{0.305561in}}{\pgfqpoint{3.082108in}{0.305561in}}%
\pgfpathquadraticcurveto{\pgfqpoint{3.071247in}{0.305561in}}{\pgfqpoint{3.065609in}{0.312315in}}%
\pgfpathquadraticcurveto{\pgfqpoint{3.059989in}{0.319070in}}{\pgfqpoint{3.059989in}{0.332075in}}%
\pgfpathlineto{\pgfqpoint{3.059989in}{0.332075in}}%
\pgfpathclose%
\pgfpathmoveto{\pgfqpoint{3.086053in}{0.371756in}}%
\pgfpathlineto{\pgfqpoint{3.086053in}{0.371756in}}%
\pgfpathlineto{\pgfqpoint{3.086053in}{0.371756in}}%
\pgfpathclose%
\pgfpathmoveto{\pgfqpoint{3.175591in}{0.360678in}}%
\pgfpathlineto{\pgfqpoint{3.175591in}{0.394793in}}%
\pgfpathlineto{\pgfqpoint{3.185948in}{0.394793in}}%
\pgfpathlineto{\pgfqpoint{3.185948in}{0.307200in}}%
\pgfpathlineto{\pgfqpoint{3.175591in}{0.307200in}}%
\pgfpathlineto{\pgfqpoint{3.175591in}{0.316656in}}%
\pgfpathquadraticcurveto{\pgfqpoint{3.172331in}{0.311037in}}{\pgfqpoint{3.167342in}{0.308299in}}%
\pgfpathquadraticcurveto{\pgfqpoint{3.162370in}{0.305561in}}{\pgfqpoint{3.155382in}{0.305561in}}%
\pgfpathquadraticcurveto{\pgfqpoint{3.143962in}{0.305561in}}{\pgfqpoint{3.136775in}{0.314675in}}%
\pgfpathquadraticcurveto{\pgfqpoint{3.129606in}{0.323807in}}{\pgfqpoint{3.129606in}{0.338667in}}%
\pgfpathquadraticcurveto{\pgfqpoint{3.129606in}{0.353527in}}{\pgfqpoint{3.136775in}{0.362641in}}%
\pgfpathquadraticcurveto{\pgfqpoint{3.143962in}{0.371756in}}{\pgfqpoint{3.155382in}{0.371756in}}%
\pgfpathquadraticcurveto{\pgfqpoint{3.162370in}{0.371756in}}{\pgfqpoint{3.167342in}{0.369018in}}%
\pgfpathquadraticcurveto{\pgfqpoint{3.172331in}{0.366298in}}{\pgfqpoint{3.175591in}{0.360678in}}%
\pgfpathlineto{\pgfqpoint{3.175591in}{0.360678in}}%
\pgfpathclose%
\pgfpathmoveto{\pgfqpoint{3.140305in}{0.338667in}}%
\pgfpathquadraticcurveto{\pgfqpoint{3.140305in}{0.327248in}}{\pgfqpoint{3.145007in}{0.320745in}}%
\pgfpathquadraticcurveto{\pgfqpoint{3.149708in}{0.314243in}}{\pgfqpoint{3.157921in}{0.314243in}}%
\pgfpathquadraticcurveto{\pgfqpoint{3.166135in}{0.314243in}}{\pgfqpoint{3.170854in}{0.320745in}}%
\pgfpathquadraticcurveto{\pgfqpoint{3.175591in}{0.327248in}}{\pgfqpoint{3.175591in}{0.338667in}}%
\pgfpathquadraticcurveto{\pgfqpoint{3.175591in}{0.350087in}}{\pgfqpoint{3.170854in}{0.356589in}}%
\pgfpathquadraticcurveto{\pgfqpoint{3.166135in}{0.363092in}}{\pgfqpoint{3.157921in}{0.363092in}}%
\pgfpathquadraticcurveto{\pgfqpoint{3.149708in}{0.363092in}}{\pgfqpoint{3.145007in}{0.356589in}}%
\pgfpathquadraticcurveto{\pgfqpoint{3.140305in}{0.350087in}}{\pgfqpoint{3.140305in}{0.338667in}}%
\pgfpathlineto{\pgfqpoint{3.140305in}{0.338667in}}%
\pgfpathclose%
\pgfpathmoveto{\pgfqpoint{3.233518in}{0.301346in}}%
\pgfpathquadraticcurveto{\pgfqpoint{3.229141in}{0.290088in}}{\pgfqpoint{3.224962in}{0.286666in}}%
\pgfpathquadraticcurveto{\pgfqpoint{3.220802in}{0.283226in}}{\pgfqpoint{3.213831in}{0.283226in}}%
\pgfpathlineto{\pgfqpoint{3.205545in}{0.283226in}}%
\pgfpathlineto{\pgfqpoint{3.205545in}{0.291890in}}%
\pgfpathlineto{\pgfqpoint{3.211633in}{0.291890in}}%
\pgfpathquadraticcurveto{\pgfqpoint{3.215902in}{0.291890in}}{\pgfqpoint{3.218262in}{0.293925in}}%
\pgfpathquadraticcurveto{\pgfqpoint{3.220640in}{0.295942in}}{\pgfqpoint{3.223503in}{0.303489in}}%
\pgfpathlineto{\pgfqpoint{3.225359in}{0.308209in}}%
\pgfpathlineto{\pgfqpoint{3.199872in}{0.370243in}}%
\pgfpathlineto{\pgfqpoint{3.210841in}{0.370243in}}%
\pgfpathlineto{\pgfqpoint{3.230546in}{0.320943in}}%
\pgfpathlineto{\pgfqpoint{3.250252in}{0.370243in}}%
\pgfpathlineto{\pgfqpoint{3.261221in}{0.370243in}}%
\pgfpathlineto{\pgfqpoint{3.233518in}{0.301346in}}%
\pgfpathlineto{\pgfqpoint{3.233518in}{0.301346in}}%
\pgfpathclose%
\pgfpathmoveto{\pgfqpoint{3.337935in}{0.347115in}}%
\pgfpathquadraticcurveto{\pgfqpoint{3.329829in}{0.347115in}}{\pgfqpoint{3.325182in}{0.342774in}}%
\pgfpathquadraticcurveto{\pgfqpoint{3.320553in}{0.338433in}}{\pgfqpoint{3.320553in}{0.330850in}}%
\pgfpathquadraticcurveto{\pgfqpoint{3.320553in}{0.323249in}}{\pgfqpoint{3.325182in}{0.318908in}}%
\pgfpathquadraticcurveto{\pgfqpoint{3.329829in}{0.314567in}}{\pgfqpoint{3.337935in}{0.314567in}}%
\pgfpathquadraticcurveto{\pgfqpoint{3.346040in}{0.314567in}}{\pgfqpoint{3.350705in}{0.318926in}}%
\pgfpathquadraticcurveto{\pgfqpoint{3.355388in}{0.323303in}}{\pgfqpoint{3.355388in}{0.330850in}}%
\pgfpathquadraticcurveto{\pgfqpoint{3.355388in}{0.338433in}}{\pgfqpoint{3.350741in}{0.342774in}}%
\pgfpathquadraticcurveto{\pgfqpoint{3.346112in}{0.347115in}}{\pgfqpoint{3.337935in}{0.347115in}}%
\pgfpathlineto{\pgfqpoint{3.337935in}{0.347115in}}%
\pgfpathclose%
\pgfpathmoveto{\pgfqpoint{3.326569in}{0.351942in}}%
\pgfpathquadraticcurveto{\pgfqpoint{3.319256in}{0.353743in}}{\pgfqpoint{3.315167in}{0.358751in}}%
\pgfpathquadraticcurveto{\pgfqpoint{3.311097in}{0.363776in}}{\pgfqpoint{3.311097in}{0.370981in}}%
\pgfpathquadraticcurveto{\pgfqpoint{3.311097in}{0.381050in}}{\pgfqpoint{3.318265in}{0.386904in}}%
\pgfpathquadraticcurveto{\pgfqpoint{3.325452in}{0.392758in}}{\pgfqpoint{3.337935in}{0.392758in}}%
\pgfpathquadraticcurveto{\pgfqpoint{3.350489in}{0.392758in}}{\pgfqpoint{3.357640in}{0.386904in}}%
\pgfpathquadraticcurveto{\pgfqpoint{3.364791in}{0.381050in}}{\pgfqpoint{3.364791in}{0.370981in}}%
\pgfpathquadraticcurveto{\pgfqpoint{3.364791in}{0.363776in}}{\pgfqpoint{3.360702in}{0.358751in}}%
\pgfpathquadraticcurveto{\pgfqpoint{3.356631in}{0.353743in}}{\pgfqpoint{3.349372in}{0.351942in}}%
\pgfpathquadraticcurveto{\pgfqpoint{3.357586in}{0.350033in}}{\pgfqpoint{3.362161in}{0.344449in}}%
\pgfpathquadraticcurveto{\pgfqpoint{3.366754in}{0.338883in}}{\pgfqpoint{3.366754in}{0.330850in}}%
\pgfpathquadraticcurveto{\pgfqpoint{3.366754in}{0.318620in}}{\pgfqpoint{3.359297in}{0.312081in}}%
\pgfpathquadraticcurveto{\pgfqpoint{3.351840in}{0.305561in}}{\pgfqpoint{3.337935in}{0.305561in}}%
\pgfpathquadraticcurveto{\pgfqpoint{3.324047in}{0.305561in}}{\pgfqpoint{3.316572in}{0.312081in}}%
\pgfpathquadraticcurveto{\pgfqpoint{3.309115in}{0.318620in}}{\pgfqpoint{3.309115in}{0.330850in}}%
\pgfpathquadraticcurveto{\pgfqpoint{3.309115in}{0.338883in}}{\pgfqpoint{3.313726in}{0.344449in}}%
\pgfpathquadraticcurveto{\pgfqpoint{3.318355in}{0.350033in}}{\pgfqpoint{3.326569in}{0.351942in}}%
\pgfpathlineto{\pgfqpoint{3.326569in}{0.351942in}}%
\pgfpathclose%
\pgfpathmoveto{\pgfqpoint{3.322408in}{0.369900in}}%
\pgfpathquadraticcurveto{\pgfqpoint{3.322408in}{0.363380in}}{\pgfqpoint{3.326479in}{0.359723in}}%
\pgfpathquadraticcurveto{\pgfqpoint{3.330568in}{0.356067in}}{\pgfqpoint{3.337935in}{0.356067in}}%
\pgfpathquadraticcurveto{\pgfqpoint{3.345266in}{0.356067in}}{\pgfqpoint{3.349390in}{0.359723in}}%
\pgfpathquadraticcurveto{\pgfqpoint{3.353533in}{0.363380in}}{\pgfqpoint{3.353533in}{0.369900in}}%
\pgfpathquadraticcurveto{\pgfqpoint{3.353533in}{0.376439in}}{\pgfqpoint{3.349390in}{0.380095in}}%
\pgfpathquadraticcurveto{\pgfqpoint{3.345266in}{0.383752in}}{\pgfqpoint{3.337935in}{0.383752in}}%
\pgfpathquadraticcurveto{\pgfqpoint{3.330568in}{0.383752in}}{\pgfqpoint{3.326479in}{0.380095in}}%
\pgfpathquadraticcurveto{\pgfqpoint{3.322408in}{0.376439in}}{\pgfqpoint{3.322408in}{0.369900in}}%
\pgfpathlineto{\pgfqpoint{3.322408in}{0.369900in}}%
\pgfpathclose%
\pgfpathmoveto{\pgfqpoint{3.421079in}{0.332075in}}%
\pgfpathlineto{\pgfqpoint{3.421079in}{0.370243in}}%
\pgfpathlineto{\pgfqpoint{3.431436in}{0.370243in}}%
\pgfpathlineto{\pgfqpoint{3.431436in}{0.332471in}}%
\pgfpathquadraticcurveto{\pgfqpoint{3.431436in}{0.323519in}}{\pgfqpoint{3.434912in}{0.319034in}}%
\pgfpathquadraticcurveto{\pgfqpoint{3.438406in}{0.314567in}}{\pgfqpoint{3.445395in}{0.314567in}}%
\pgfpathquadraticcurveto{\pgfqpoint{3.453771in}{0.314567in}}{\pgfqpoint{3.458634in}{0.319917in}}%
\pgfpathquadraticcurveto{\pgfqpoint{3.463515in}{0.325266in}}{\pgfqpoint{3.463515in}{0.334506in}}%
\pgfpathlineto{\pgfqpoint{3.463515in}{0.370243in}}%
\pgfpathlineto{\pgfqpoint{3.473872in}{0.370243in}}%
\pgfpathlineto{\pgfqpoint{3.473872in}{0.307200in}}%
\pgfpathlineto{\pgfqpoint{3.463515in}{0.307200in}}%
\pgfpathlineto{\pgfqpoint{3.463515in}{0.316891in}}%
\pgfpathquadraticcurveto{\pgfqpoint{3.459751in}{0.311145in}}{\pgfqpoint{3.454762in}{0.308353in}}%
\pgfpathquadraticcurveto{\pgfqpoint{3.449790in}{0.305561in}}{\pgfqpoint{3.443198in}{0.305561in}}%
\pgfpathquadraticcurveto{\pgfqpoint{3.432336in}{0.305561in}}{\pgfqpoint{3.426699in}{0.312315in}}%
\pgfpathquadraticcurveto{\pgfqpoint{3.421079in}{0.319070in}}{\pgfqpoint{3.421079in}{0.332075in}}%
\pgfpathlineto{\pgfqpoint{3.421079in}{0.332075in}}%
\pgfpathclose%
\pgfpathmoveto{\pgfqpoint{3.447142in}{0.371756in}}%
\pgfpathlineto{\pgfqpoint{3.447142in}{0.371756in}}%
\pgfpathlineto{\pgfqpoint{3.447142in}{0.371756in}}%
\pgfpathclose%
\pgfpathmoveto{\pgfqpoint{3.535384in}{0.368387in}}%
\pgfpathlineto{\pgfqpoint{3.535384in}{0.358589in}}%
\pgfpathquadraticcurveto{\pgfqpoint{3.531007in}{0.360840in}}{\pgfqpoint{3.526270in}{0.361957in}}%
\pgfpathquadraticcurveto{\pgfqpoint{3.521551in}{0.363092in}}{\pgfqpoint{3.516471in}{0.363092in}}%
\pgfpathquadraticcurveto{\pgfqpoint{3.508762in}{0.363092in}}{\pgfqpoint{3.504907in}{0.360732in}}%
\pgfpathquadraticcurveto{\pgfqpoint{3.501053in}{0.358373in}}{\pgfqpoint{3.501053in}{0.353635in}}%
\pgfpathquadraticcurveto{\pgfqpoint{3.501053in}{0.350033in}}{\pgfqpoint{3.503809in}{0.347980in}}%
\pgfpathquadraticcurveto{\pgfqpoint{3.506564in}{0.345926in}}{\pgfqpoint{3.514904in}{0.344071in}}%
\pgfpathlineto{\pgfqpoint{3.518452in}{0.343278in}}%
\pgfpathquadraticcurveto{\pgfqpoint{3.529476in}{0.340919in}}{\pgfqpoint{3.534123in}{0.336614in}}%
\pgfpathquadraticcurveto{\pgfqpoint{3.538770in}{0.332309in}}{\pgfqpoint{3.538770in}{0.324600in}}%
\pgfpathquadraticcurveto{\pgfqpoint{3.538770in}{0.315810in}}{\pgfqpoint{3.531817in}{0.310676in}}%
\pgfpathquadraticcurveto{\pgfqpoint{3.524865in}{0.305561in}}{\pgfqpoint{3.512707in}{0.305561in}}%
\pgfpathquadraticcurveto{\pgfqpoint{3.507645in}{0.305561in}}{\pgfqpoint{3.502151in}{0.306552in}}%
\pgfpathquadraticcurveto{\pgfqpoint{3.496658in}{0.307542in}}{\pgfqpoint{3.490588in}{0.309506in}}%
\pgfpathlineto{\pgfqpoint{3.490588in}{0.320205in}}%
\pgfpathquadraticcurveto{\pgfqpoint{3.496334in}{0.317215in}}{\pgfqpoint{3.501899in}{0.315720in}}%
\pgfpathquadraticcurveto{\pgfqpoint{3.507465in}{0.314243in}}{\pgfqpoint{3.512941in}{0.314243in}}%
\pgfpathquadraticcurveto{\pgfqpoint{3.520254in}{0.314243in}}{\pgfqpoint{3.524180in}{0.316746in}}%
\pgfpathquadraticcurveto{\pgfqpoint{3.528125in}{0.319250in}}{\pgfqpoint{3.528125in}{0.323807in}}%
\pgfpathquadraticcurveto{\pgfqpoint{3.528125in}{0.328022in}}{\pgfqpoint{3.525279in}{0.330274in}}%
\pgfpathquadraticcurveto{\pgfqpoint{3.522451in}{0.332525in}}{\pgfqpoint{3.512815in}{0.334614in}}%
\pgfpathlineto{\pgfqpoint{3.509212in}{0.335461in}}%
\pgfpathquadraticcurveto{\pgfqpoint{3.499594in}{0.337478in}}{\pgfqpoint{3.495307in}{0.341675in}}%
\pgfpathquadraticcurveto{\pgfqpoint{3.491038in}{0.345872in}}{\pgfqpoint{3.491038in}{0.353185in}}%
\pgfpathquadraticcurveto{\pgfqpoint{3.491038in}{0.362083in}}{\pgfqpoint{3.497342in}{0.366910in}}%
\pgfpathquadraticcurveto{\pgfqpoint{3.503646in}{0.371756in}}{\pgfqpoint{3.515246in}{0.371756in}}%
\pgfpathquadraticcurveto{\pgfqpoint{3.520974in}{0.371756in}}{\pgfqpoint{3.526036in}{0.370909in}}%
\pgfpathquadraticcurveto{\pgfqpoint{3.531115in}{0.370080in}}{\pgfqpoint{3.535384in}{0.368387in}}%
\pgfpathlineto{\pgfqpoint{3.535384in}{0.368387in}}%
\pgfpathclose%
\pgfpathmoveto{\pgfqpoint{3.609180in}{0.341315in}}%
\pgfpathlineto{\pgfqpoint{3.609180in}{0.336254in}}%
\pgfpathlineto{\pgfqpoint{3.561556in}{0.336254in}}%
\pgfpathquadraticcurveto{\pgfqpoint{3.562240in}{0.325554in}}{\pgfqpoint{3.568004in}{0.319953in}}%
\pgfpathquadraticcurveto{\pgfqpoint{3.573768in}{0.314351in}}{\pgfqpoint{3.584071in}{0.314351in}}%
\pgfpathquadraticcurveto{\pgfqpoint{3.590033in}{0.314351in}}{\pgfqpoint{3.595635in}{0.315810in}}%
\pgfpathquadraticcurveto{\pgfqpoint{3.601236in}{0.317269in}}{\pgfqpoint{3.606766in}{0.320205in}}%
\pgfpathlineto{\pgfqpoint{3.606766in}{0.310406in}}%
\pgfpathquadraticcurveto{\pgfqpoint{3.601182in}{0.308047in}}{\pgfqpoint{3.595328in}{0.306804in}}%
\pgfpathquadraticcurveto{\pgfqpoint{3.589474in}{0.305561in}}{\pgfqpoint{3.583458in}{0.305561in}}%
\pgfpathquadraticcurveto{\pgfqpoint{3.568364in}{0.305561in}}{\pgfqpoint{3.559556in}{0.314333in}}%
\pgfpathquadraticcurveto{\pgfqpoint{3.550748in}{0.323123in}}{\pgfqpoint{3.550748in}{0.338109in}}%
\pgfpathquadraticcurveto{\pgfqpoint{3.550748in}{0.353581in}}{\pgfqpoint{3.559106in}{0.362659in}}%
\pgfpathquadraticcurveto{\pgfqpoint{3.567464in}{0.371756in}}{\pgfqpoint{3.581657in}{0.371756in}}%
\pgfpathquadraticcurveto{\pgfqpoint{3.594374in}{0.371756in}}{\pgfqpoint{3.601777in}{0.363560in}}%
\pgfpathquadraticcurveto{\pgfqpoint{3.609180in}{0.355383in}}{\pgfqpoint{3.609180in}{0.341315in}}%
\pgfpathlineto{\pgfqpoint{3.609180in}{0.341315in}}%
\pgfpathclose%
\pgfpathmoveto{\pgfqpoint{3.598823in}{0.344359in}}%
\pgfpathquadraticcurveto{\pgfqpoint{3.598715in}{0.352843in}}{\pgfqpoint{3.594067in}{0.357904in}}%
\pgfpathquadraticcurveto{\pgfqpoint{3.589420in}{0.362984in}}{\pgfqpoint{3.581765in}{0.362984in}}%
\pgfpathquadraticcurveto{\pgfqpoint{3.573101in}{0.362984in}}{\pgfqpoint{3.567896in}{0.358084in}}%
\pgfpathquadraticcurveto{\pgfqpoint{3.562690in}{0.353185in}}{\pgfqpoint{3.561898in}{0.344287in}}%
\pgfpathlineto{\pgfqpoint{3.598823in}{0.344359in}}%
\pgfpathlineto{\pgfqpoint{3.598823in}{0.344359in}}%
\pgfpathclose%
\pgfpathmoveto{\pgfqpoint{3.666368in}{0.368387in}}%
\pgfpathlineto{\pgfqpoint{3.666368in}{0.358589in}}%
\pgfpathquadraticcurveto{\pgfqpoint{3.661991in}{0.360840in}}{\pgfqpoint{3.657254in}{0.361957in}}%
\pgfpathquadraticcurveto{\pgfqpoint{3.652535in}{0.363092in}}{\pgfqpoint{3.647456in}{0.363092in}}%
\pgfpathquadraticcurveto{\pgfqpoint{3.639746in}{0.363092in}}{\pgfqpoint{3.635892in}{0.360732in}}%
\pgfpathquadraticcurveto{\pgfqpoint{3.632037in}{0.358373in}}{\pgfqpoint{3.632037in}{0.353635in}}%
\pgfpathquadraticcurveto{\pgfqpoint{3.632037in}{0.350033in}}{\pgfqpoint{3.634793in}{0.347980in}}%
\pgfpathquadraticcurveto{\pgfqpoint{3.637549in}{0.345926in}}{\pgfqpoint{3.645888in}{0.344071in}}%
\pgfpathlineto{\pgfqpoint{3.649437in}{0.343278in}}%
\pgfpathquadraticcurveto{\pgfqpoint{3.660460in}{0.340919in}}{\pgfqpoint{3.665107in}{0.336614in}}%
\pgfpathquadraticcurveto{\pgfqpoint{3.669755in}{0.332309in}}{\pgfqpoint{3.669755in}{0.324600in}}%
\pgfpathquadraticcurveto{\pgfqpoint{3.669755in}{0.315810in}}{\pgfqpoint{3.662802in}{0.310676in}}%
\pgfpathquadraticcurveto{\pgfqpoint{3.655849in}{0.305561in}}{\pgfqpoint{3.643691in}{0.305561in}}%
\pgfpathquadraticcurveto{\pgfqpoint{3.638630in}{0.305561in}}{\pgfqpoint{3.633136in}{0.306552in}}%
\pgfpathquadraticcurveto{\pgfqpoint{3.627642in}{0.307542in}}{\pgfqpoint{3.621572in}{0.309506in}}%
\pgfpathlineto{\pgfqpoint{3.621572in}{0.320205in}}%
\pgfpathquadraticcurveto{\pgfqpoint{3.627318in}{0.317215in}}{\pgfqpoint{3.632884in}{0.315720in}}%
\pgfpathquadraticcurveto{\pgfqpoint{3.638449in}{0.314243in}}{\pgfqpoint{3.643925in}{0.314243in}}%
\pgfpathquadraticcurveto{\pgfqpoint{3.651238in}{0.314243in}}{\pgfqpoint{3.655165in}{0.316746in}}%
\pgfpathquadraticcurveto{\pgfqpoint{3.659109in}{0.319250in}}{\pgfqpoint{3.659109in}{0.323807in}}%
\pgfpathquadraticcurveto{\pgfqpoint{3.659109in}{0.328022in}}{\pgfqpoint{3.656263in}{0.330274in}}%
\pgfpathquadraticcurveto{\pgfqpoint{3.653436in}{0.332525in}}{\pgfqpoint{3.643799in}{0.334614in}}%
\pgfpathlineto{\pgfqpoint{3.640197in}{0.335461in}}%
\pgfpathquadraticcurveto{\pgfqpoint{3.630578in}{0.337478in}}{\pgfqpoint{3.626291in}{0.341675in}}%
\pgfpathquadraticcurveto{\pgfqpoint{3.622022in}{0.345872in}}{\pgfqpoint{3.622022in}{0.353185in}}%
\pgfpathquadraticcurveto{\pgfqpoint{3.622022in}{0.362083in}}{\pgfqpoint{3.628327in}{0.366910in}}%
\pgfpathquadraticcurveto{\pgfqpoint{3.634631in}{0.371756in}}{\pgfqpoint{3.646231in}{0.371756in}}%
\pgfpathquadraticcurveto{\pgfqpoint{3.651959in}{0.371756in}}{\pgfqpoint{3.657020in}{0.370909in}}%
\pgfpathquadraticcurveto{\pgfqpoint{3.662099in}{0.370080in}}{\pgfqpoint{3.666368in}{0.368387in}}%
\pgfpathlineto{\pgfqpoint{3.666368in}{0.368387in}}%
\pgfpathclose%
\pgfpathmoveto{\pgfqpoint{3.721810in}{0.332075in}}%
\pgfpathlineto{\pgfqpoint{3.721810in}{0.370243in}}%
\pgfpathlineto{\pgfqpoint{3.732167in}{0.370243in}}%
\pgfpathlineto{\pgfqpoint{3.732167in}{0.332471in}}%
\pgfpathquadraticcurveto{\pgfqpoint{3.732167in}{0.323519in}}{\pgfqpoint{3.735643in}{0.319034in}}%
\pgfpathquadraticcurveto{\pgfqpoint{3.739137in}{0.314567in}}{\pgfqpoint{3.746126in}{0.314567in}}%
\pgfpathquadraticcurveto{\pgfqpoint{3.754502in}{0.314567in}}{\pgfqpoint{3.759365in}{0.319917in}}%
\pgfpathquadraticcurveto{\pgfqpoint{3.764246in}{0.325266in}}{\pgfqpoint{3.764246in}{0.334506in}}%
\pgfpathlineto{\pgfqpoint{3.764246in}{0.370243in}}%
\pgfpathlineto{\pgfqpoint{3.774603in}{0.370243in}}%
\pgfpathlineto{\pgfqpoint{3.774603in}{0.307200in}}%
\pgfpathlineto{\pgfqpoint{3.764246in}{0.307200in}}%
\pgfpathlineto{\pgfqpoint{3.764246in}{0.316891in}}%
\pgfpathquadraticcurveto{\pgfqpoint{3.760482in}{0.311145in}}{\pgfqpoint{3.755492in}{0.308353in}}%
\pgfpathquadraticcurveto{\pgfqpoint{3.750521in}{0.305561in}}{\pgfqpoint{3.743929in}{0.305561in}}%
\pgfpathquadraticcurveto{\pgfqpoint{3.733067in}{0.305561in}}{\pgfqpoint{3.727429in}{0.312315in}}%
\pgfpathquadraticcurveto{\pgfqpoint{3.721810in}{0.319070in}}{\pgfqpoint{3.721810in}{0.332075in}}%
\pgfpathlineto{\pgfqpoint{3.721810in}{0.332075in}}%
\pgfpathclose%
\pgfpathmoveto{\pgfqpoint{3.747873in}{0.371756in}}%
\pgfpathlineto{\pgfqpoint{3.747873in}{0.371756in}}%
\pgfpathlineto{\pgfqpoint{3.747873in}{0.371756in}}%
\pgfpathclose%
\pgfpathmoveto{\pgfqpoint{3.848345in}{0.345260in}}%
\pgfpathlineto{\pgfqpoint{3.848345in}{0.307200in}}%
\pgfpathlineto{\pgfqpoint{3.837988in}{0.307200in}}%
\pgfpathlineto{\pgfqpoint{3.837988in}{0.344917in}}%
\pgfpathquadraticcurveto{\pgfqpoint{3.837988in}{0.353869in}}{\pgfqpoint{3.834494in}{0.358300in}}%
\pgfpathquadraticcurveto{\pgfqpoint{3.830999in}{0.362749in}}{\pgfqpoint{3.824029in}{0.362749in}}%
\pgfpathquadraticcurveto{\pgfqpoint{3.815635in}{0.362749in}}{\pgfqpoint{3.810790in}{0.357400in}}%
\pgfpathquadraticcurveto{\pgfqpoint{3.805944in}{0.352068in}}{\pgfqpoint{3.805944in}{0.342828in}}%
\pgfpathlineto{\pgfqpoint{3.805944in}{0.307200in}}%
\pgfpathlineto{\pgfqpoint{3.795533in}{0.307200in}}%
\pgfpathlineto{\pgfqpoint{3.795533in}{0.370243in}}%
\pgfpathlineto{\pgfqpoint{3.805944in}{0.370243in}}%
\pgfpathlineto{\pgfqpoint{3.805944in}{0.360444in}}%
\pgfpathquadraticcurveto{\pgfqpoint{3.809673in}{0.366136in}}{\pgfqpoint{3.814698in}{0.368946in}}%
\pgfpathquadraticcurveto{\pgfqpoint{3.819742in}{0.371756in}}{\pgfqpoint{3.826334in}{0.371756in}}%
\pgfpathquadraticcurveto{\pgfqpoint{3.837196in}{0.371756in}}{\pgfqpoint{3.842761in}{0.365037in}}%
\pgfpathquadraticcurveto{\pgfqpoint{3.848345in}{0.358318in}}{\pgfqpoint{3.848345in}{0.345260in}}%
\pgfpathlineto{\pgfqpoint{3.848345in}{0.345260in}}%
\pgfpathclose%
\pgfpathmoveto{\pgfqpoint{3.905516in}{0.360570in}}%
\pgfpathquadraticcurveto{\pgfqpoint{3.903768in}{0.361579in}}{\pgfqpoint{3.901715in}{0.362047in}}%
\pgfpathquadraticcurveto{\pgfqpoint{3.899662in}{0.362533in}}{\pgfqpoint{3.897194in}{0.362533in}}%
\pgfpathquadraticcurveto{\pgfqpoint{3.888404in}{0.362533in}}{\pgfqpoint{3.883703in}{0.356823in}}%
\pgfpathquadraticcurveto{\pgfqpoint{3.879002in}{0.351114in}}{\pgfqpoint{3.879002in}{0.340414in}}%
\pgfpathlineto{\pgfqpoint{3.879002in}{0.307200in}}%
\pgfpathlineto{\pgfqpoint{3.868591in}{0.307200in}}%
\pgfpathlineto{\pgfqpoint{3.868591in}{0.370243in}}%
\pgfpathlineto{\pgfqpoint{3.879002in}{0.370243in}}%
\pgfpathlineto{\pgfqpoint{3.879002in}{0.360444in}}%
\pgfpathquadraticcurveto{\pgfqpoint{3.882280in}{0.366190in}}{\pgfqpoint{3.887503in}{0.368964in}}%
\pgfpathquadraticcurveto{\pgfqpoint{3.892745in}{0.371756in}}{\pgfqpoint{3.900238in}{0.371756in}}%
\pgfpathquadraticcurveto{\pgfqpoint{3.901301in}{0.371756in}}{\pgfqpoint{3.902598in}{0.371611in}}%
\pgfpathquadraticcurveto{\pgfqpoint{3.903895in}{0.371485in}}{\pgfqpoint{3.905462in}{0.371197in}}%
\pgfpathlineto{\pgfqpoint{3.905516in}{0.360570in}}%
\pgfpathlineto{\pgfqpoint{3.905516in}{0.360570in}}%
\pgfpathclose%
\pgfpathmoveto{\pgfqpoint{3.967766in}{0.341315in}}%
\pgfpathlineto{\pgfqpoint{3.967766in}{0.336254in}}%
\pgfpathlineto{\pgfqpoint{3.920141in}{0.336254in}}%
\pgfpathquadraticcurveto{\pgfqpoint{3.920826in}{0.325554in}}{\pgfqpoint{3.926590in}{0.319953in}}%
\pgfpathquadraticcurveto{\pgfqpoint{3.932354in}{0.314351in}}{\pgfqpoint{3.942657in}{0.314351in}}%
\pgfpathquadraticcurveto{\pgfqpoint{3.948619in}{0.314351in}}{\pgfqpoint{3.954220in}{0.315810in}}%
\pgfpathquadraticcurveto{\pgfqpoint{3.959822in}{0.317269in}}{\pgfqpoint{3.965352in}{0.320205in}}%
\pgfpathlineto{\pgfqpoint{3.965352in}{0.310406in}}%
\pgfpathquadraticcurveto{\pgfqpoint{3.959768in}{0.308047in}}{\pgfqpoint{3.953914in}{0.306804in}}%
\pgfpathquadraticcurveto{\pgfqpoint{3.948060in}{0.305561in}}{\pgfqpoint{3.942044in}{0.305561in}}%
\pgfpathquadraticcurveto{\pgfqpoint{3.926950in}{0.305561in}}{\pgfqpoint{3.918142in}{0.314333in}}%
\pgfpathquadraticcurveto{\pgfqpoint{3.909334in}{0.323123in}}{\pgfqpoint{3.909334in}{0.338109in}}%
\pgfpathquadraticcurveto{\pgfqpoint{3.909334in}{0.353581in}}{\pgfqpoint{3.917692in}{0.362659in}}%
\pgfpathquadraticcurveto{\pgfqpoint{3.926049in}{0.371756in}}{\pgfqpoint{3.940243in}{0.371756in}}%
\pgfpathquadraticcurveto{\pgfqpoint{3.952960in}{0.371756in}}{\pgfqpoint{3.960363in}{0.363560in}}%
\pgfpathquadraticcurveto{\pgfqpoint{3.967766in}{0.355383in}}{\pgfqpoint{3.967766in}{0.341315in}}%
\pgfpathlineto{\pgfqpoint{3.967766in}{0.341315in}}%
\pgfpathclose%
\pgfpathmoveto{\pgfqpoint{3.957409in}{0.344359in}}%
\pgfpathquadraticcurveto{\pgfqpoint{3.957301in}{0.352843in}}{\pgfqpoint{3.952653in}{0.357904in}}%
\pgfpathquadraticcurveto{\pgfqpoint{3.948006in}{0.362984in}}{\pgfqpoint{3.940351in}{0.362984in}}%
\pgfpathquadraticcurveto{\pgfqpoint{3.931687in}{0.362984in}}{\pgfqpoint{3.926482in}{0.358084in}}%
\pgfpathquadraticcurveto{\pgfqpoint{3.921276in}{0.353185in}}{\pgfqpoint{3.920484in}{0.344287in}}%
\pgfpathlineto{\pgfqpoint{3.957409in}{0.344359in}}%
\pgfpathlineto{\pgfqpoint{3.957409in}{0.344359in}}%
\pgfpathclose%
\pgfpathmoveto{\pgfqpoint{4.024954in}{0.368387in}}%
\pgfpathlineto{\pgfqpoint{4.024954in}{0.358589in}}%
\pgfpathquadraticcurveto{\pgfqpoint{4.020577in}{0.360840in}}{\pgfqpoint{4.015840in}{0.361957in}}%
\pgfpathquadraticcurveto{\pgfqpoint{4.011121in}{0.363092in}}{\pgfqpoint{4.006041in}{0.363092in}}%
\pgfpathquadraticcurveto{\pgfqpoint{3.998332in}{0.363092in}}{\pgfqpoint{3.994478in}{0.360732in}}%
\pgfpathquadraticcurveto{\pgfqpoint{3.990623in}{0.358373in}}{\pgfqpoint{3.990623in}{0.353635in}}%
\pgfpathquadraticcurveto{\pgfqpoint{3.990623in}{0.350033in}}{\pgfqpoint{3.993379in}{0.347980in}}%
\pgfpathquadraticcurveto{\pgfqpoint{3.996135in}{0.345926in}}{\pgfqpoint{4.004474in}{0.344071in}}%
\pgfpathlineto{\pgfqpoint{4.008023in}{0.343278in}}%
\pgfpathquadraticcurveto{\pgfqpoint{4.019046in}{0.340919in}}{\pgfqpoint{4.023693in}{0.336614in}}%
\pgfpathquadraticcurveto{\pgfqpoint{4.028340in}{0.332309in}}{\pgfqpoint{4.028340in}{0.324600in}}%
\pgfpathquadraticcurveto{\pgfqpoint{4.028340in}{0.315810in}}{\pgfqpoint{4.021388in}{0.310676in}}%
\pgfpathquadraticcurveto{\pgfqpoint{4.014435in}{0.305561in}}{\pgfqpoint{4.002277in}{0.305561in}}%
\pgfpathquadraticcurveto{\pgfqpoint{3.997215in}{0.305561in}}{\pgfqpoint{3.991722in}{0.306552in}}%
\pgfpathquadraticcurveto{\pgfqpoint{3.986228in}{0.307542in}}{\pgfqpoint{3.980158in}{0.309506in}}%
\pgfpathlineto{\pgfqpoint{3.980158in}{0.320205in}}%
\pgfpathquadraticcurveto{\pgfqpoint{3.985904in}{0.317215in}}{\pgfqpoint{3.991470in}{0.315720in}}%
\pgfpathquadraticcurveto{\pgfqpoint{3.997035in}{0.314243in}}{\pgfqpoint{4.002511in}{0.314243in}}%
\pgfpathquadraticcurveto{\pgfqpoint{4.009824in}{0.314243in}}{\pgfqpoint{4.013751in}{0.316746in}}%
\pgfpathquadraticcurveto{\pgfqpoint{4.017695in}{0.319250in}}{\pgfqpoint{4.017695in}{0.323807in}}%
\pgfpathquadraticcurveto{\pgfqpoint{4.017695in}{0.328022in}}{\pgfqpoint{4.014849in}{0.330274in}}%
\pgfpathquadraticcurveto{\pgfqpoint{4.012021in}{0.332525in}}{\pgfqpoint{4.002385in}{0.334614in}}%
\pgfpathlineto{\pgfqpoint{3.998783in}{0.335461in}}%
\pgfpathquadraticcurveto{\pgfqpoint{3.989164in}{0.337478in}}{\pgfqpoint{3.984877in}{0.341675in}}%
\pgfpathquadraticcurveto{\pgfqpoint{3.980608in}{0.345872in}}{\pgfqpoint{3.980608in}{0.353185in}}%
\pgfpathquadraticcurveto{\pgfqpoint{3.980608in}{0.362083in}}{\pgfqpoint{3.986913in}{0.366910in}}%
\pgfpathquadraticcurveto{\pgfqpoint{3.993217in}{0.371756in}}{\pgfqpoint{4.004817in}{0.371756in}}%
\pgfpathquadraticcurveto{\pgfqpoint{4.010544in}{0.371756in}}{\pgfqpoint{4.015606in}{0.370909in}}%
\pgfpathquadraticcurveto{\pgfqpoint{4.020685in}{0.370080in}}{\pgfqpoint{4.024954in}{0.368387in}}%
\pgfpathlineto{\pgfqpoint{4.024954in}{0.368387in}}%
\pgfpathclose%
\pgfpathmoveto{\pgfqpoint{4.055071in}{0.388147in}}%
\pgfpathlineto{\pgfqpoint{4.055071in}{0.370243in}}%
\pgfpathlineto{\pgfqpoint{4.076397in}{0.370243in}}%
\pgfpathlineto{\pgfqpoint{4.076397in}{0.362191in}}%
\pgfpathlineto{\pgfqpoint{4.055071in}{0.362191in}}%
\pgfpathlineto{\pgfqpoint{4.055071in}{0.327968in}}%
\pgfpathquadraticcurveto{\pgfqpoint{4.055071in}{0.320259in}}{\pgfqpoint{4.057178in}{0.318061in}}%
\pgfpathquadraticcurveto{\pgfqpoint{4.059285in}{0.315864in}}{\pgfqpoint{4.065770in}{0.315864in}}%
\pgfpathlineto{\pgfqpoint{4.076397in}{0.315864in}}%
\pgfpathlineto{\pgfqpoint{4.076397in}{0.307200in}}%
\pgfpathlineto{\pgfqpoint{4.065770in}{0.307200in}}%
\pgfpathquadraticcurveto{\pgfqpoint{4.053774in}{0.307200in}}{\pgfqpoint{4.049217in}{0.311667in}}%
\pgfpathquadraticcurveto{\pgfqpoint{4.044660in}{0.316152in}}{\pgfqpoint{4.044660in}{0.327968in}}%
\pgfpathlineto{\pgfqpoint{4.044660in}{0.362191in}}%
\pgfpathlineto{\pgfqpoint{4.037058in}{0.362191in}}%
\pgfpathlineto{\pgfqpoint{4.037058in}{0.370243in}}%
\pgfpathlineto{\pgfqpoint{4.044660in}{0.370243in}}%
\pgfpathlineto{\pgfqpoint{4.044660in}{0.388147in}}%
\pgfpathlineto{\pgfqpoint{4.055071in}{0.388147in}}%
\pgfpathlineto{\pgfqpoint{4.055071in}{0.388147in}}%
\pgfpathclose%
\pgfpathmoveto{\pgfqpoint{4.126543in}{0.360570in}}%
\pgfpathquadraticcurveto{\pgfqpoint{4.124796in}{0.361579in}}{\pgfqpoint{4.122742in}{0.362047in}}%
\pgfpathquadraticcurveto{\pgfqpoint{4.120689in}{0.362533in}}{\pgfqpoint{4.118221in}{0.362533in}}%
\pgfpathquadraticcurveto{\pgfqpoint{4.109431in}{0.362533in}}{\pgfqpoint{4.104730in}{0.356823in}}%
\pgfpathquadraticcurveto{\pgfqpoint{4.100029in}{0.351114in}}{\pgfqpoint{4.100029in}{0.340414in}}%
\pgfpathlineto{\pgfqpoint{4.100029in}{0.307200in}}%
\pgfpathlineto{\pgfqpoint{4.089618in}{0.307200in}}%
\pgfpathlineto{\pgfqpoint{4.089618in}{0.370243in}}%
\pgfpathlineto{\pgfqpoint{4.100029in}{0.370243in}}%
\pgfpathlineto{\pgfqpoint{4.100029in}{0.360444in}}%
\pgfpathquadraticcurveto{\pgfqpoint{4.103307in}{0.366190in}}{\pgfqpoint{4.108531in}{0.368964in}}%
\pgfpathquadraticcurveto{\pgfqpoint{4.113772in}{0.371756in}}{\pgfqpoint{4.121265in}{0.371756in}}%
\pgfpathquadraticcurveto{\pgfqpoint{4.122328in}{0.371756in}}{\pgfqpoint{4.123625in}{0.371611in}}%
\pgfpathquadraticcurveto{\pgfqpoint{4.124922in}{0.371485in}}{\pgfqpoint{4.126489in}{0.371197in}}%
\pgfpathlineto{\pgfqpoint{4.126543in}{0.360570in}}%
\pgfpathlineto{\pgfqpoint{4.126543in}{0.360570in}}%
\pgfpathclose%
\pgfpathmoveto{\pgfqpoint{4.137404in}{0.370243in}}%
\pgfpathlineto{\pgfqpoint{4.147761in}{0.370243in}}%
\pgfpathlineto{\pgfqpoint{4.147761in}{0.307200in}}%
\pgfpathlineto{\pgfqpoint{4.137404in}{0.307200in}}%
\pgfpathlineto{\pgfqpoint{4.137404in}{0.370243in}}%
\pgfpathlineto{\pgfqpoint{4.137404in}{0.370243in}}%
\pgfpathclose%
\pgfpathmoveto{\pgfqpoint{4.137404in}{0.394793in}}%
\pgfpathlineto{\pgfqpoint{4.147761in}{0.394793in}}%
\pgfpathlineto{\pgfqpoint{4.147761in}{0.381662in}}%
\pgfpathlineto{\pgfqpoint{4.137404in}{0.381662in}}%
\pgfpathlineto{\pgfqpoint{4.137404in}{0.394793in}}%
\pgfpathlineto{\pgfqpoint{4.137404in}{0.394793in}}%
\pgfpathclose%
\pgfpathmoveto{\pgfqpoint{4.214802in}{0.367829in}}%
\pgfpathlineto{\pgfqpoint{4.214802in}{0.358138in}}%
\pgfpathquadraticcurveto{\pgfqpoint{4.210407in}{0.360570in}}{\pgfqpoint{4.205994in}{0.361777in}}%
\pgfpathquadraticcurveto{\pgfqpoint{4.201581in}{0.362984in}}{\pgfqpoint{4.197078in}{0.362984in}}%
\pgfpathquadraticcurveto{\pgfqpoint{4.186992in}{0.362984in}}{\pgfqpoint{4.181408in}{0.356589in}}%
\pgfpathquadraticcurveto{\pgfqpoint{4.175842in}{0.350213in}}{\pgfqpoint{4.175842in}{0.338667in}}%
\pgfpathquadraticcurveto{\pgfqpoint{4.175842in}{0.327121in}}{\pgfqpoint{4.181408in}{0.320727in}}%
\pgfpathquadraticcurveto{\pgfqpoint{4.186992in}{0.314351in}}{\pgfqpoint{4.197078in}{0.314351in}}%
\pgfpathquadraticcurveto{\pgfqpoint{4.201581in}{0.314351in}}{\pgfqpoint{4.205994in}{0.315558in}}%
\pgfpathquadraticcurveto{\pgfqpoint{4.210407in}{0.316764in}}{\pgfqpoint{4.214802in}{0.319196in}}%
\pgfpathlineto{\pgfqpoint{4.214802in}{0.309614in}}%
\pgfpathquadraticcurveto{\pgfqpoint{4.210461in}{0.307596in}}{\pgfqpoint{4.205814in}{0.306588in}}%
\pgfpathquadraticcurveto{\pgfqpoint{4.201185in}{0.305561in}}{\pgfqpoint{4.195944in}{0.305561in}}%
\pgfpathquadraticcurveto{\pgfqpoint{4.181696in}{0.305561in}}{\pgfqpoint{4.173302in}{0.314513in}}%
\pgfpathquadraticcurveto{\pgfqpoint{4.164927in}{0.323465in}}{\pgfqpoint{4.164927in}{0.338667in}}%
\pgfpathquadraticcurveto{\pgfqpoint{4.164927in}{0.354086in}}{\pgfqpoint{4.173392in}{0.362912in}}%
\pgfpathquadraticcurveto{\pgfqpoint{4.181876in}{0.371756in}}{\pgfqpoint{4.196628in}{0.371756in}}%
\pgfpathquadraticcurveto{\pgfqpoint{4.201401in}{0.371756in}}{\pgfqpoint{4.205958in}{0.370765in}}%
\pgfpathquadraticcurveto{\pgfqpoint{4.210515in}{0.369792in}}{\pgfqpoint{4.214802in}{0.367829in}}%
\pgfpathlineto{\pgfqpoint{4.214802in}{0.367829in}}%
\pgfpathclose%
\pgfpathmoveto{\pgfqpoint{4.243063in}{0.388147in}}%
\pgfpathlineto{\pgfqpoint{4.243063in}{0.370243in}}%
\pgfpathlineto{\pgfqpoint{4.264390in}{0.370243in}}%
\pgfpathlineto{\pgfqpoint{4.264390in}{0.362191in}}%
\pgfpathlineto{\pgfqpoint{4.243063in}{0.362191in}}%
\pgfpathlineto{\pgfqpoint{4.243063in}{0.327968in}}%
\pgfpathquadraticcurveto{\pgfqpoint{4.243063in}{0.320259in}}{\pgfqpoint{4.245171in}{0.318061in}}%
\pgfpathquadraticcurveto{\pgfqpoint{4.247278in}{0.315864in}}{\pgfqpoint{4.253763in}{0.315864in}}%
\pgfpathlineto{\pgfqpoint{4.264390in}{0.315864in}}%
\pgfpathlineto{\pgfqpoint{4.264390in}{0.307200in}}%
\pgfpathlineto{\pgfqpoint{4.253763in}{0.307200in}}%
\pgfpathquadraticcurveto{\pgfqpoint{4.241766in}{0.307200in}}{\pgfqpoint{4.237209in}{0.311667in}}%
\pgfpathquadraticcurveto{\pgfqpoint{4.232652in}{0.316152in}}{\pgfqpoint{4.232652in}{0.327968in}}%
\pgfpathlineto{\pgfqpoint{4.232652in}{0.362191in}}%
\pgfpathlineto{\pgfqpoint{4.225051in}{0.362191in}}%
\pgfpathlineto{\pgfqpoint{4.225051in}{0.370243in}}%
\pgfpathlineto{\pgfqpoint{4.232652in}{0.370243in}}%
\pgfpathlineto{\pgfqpoint{4.232652in}{0.388147in}}%
\pgfpathlineto{\pgfqpoint{4.243063in}{0.388147in}}%
\pgfpathlineto{\pgfqpoint{4.243063in}{0.388147in}}%
\pgfpathclose%
\pgfpathmoveto{\pgfqpoint{4.331935in}{0.341315in}}%
\pgfpathlineto{\pgfqpoint{4.331935in}{0.336254in}}%
\pgfpathlineto{\pgfqpoint{4.284311in}{0.336254in}}%
\pgfpathquadraticcurveto{\pgfqpoint{4.284996in}{0.325554in}}{\pgfqpoint{4.290760in}{0.319953in}}%
\pgfpathquadraticcurveto{\pgfqpoint{4.296523in}{0.314351in}}{\pgfqpoint{4.306826in}{0.314351in}}%
\pgfpathquadraticcurveto{\pgfqpoint{4.312788in}{0.314351in}}{\pgfqpoint{4.318390in}{0.315810in}}%
\pgfpathquadraticcurveto{\pgfqpoint{4.323992in}{0.317269in}}{\pgfqpoint{4.329522in}{0.320205in}}%
\pgfpathlineto{\pgfqpoint{4.329522in}{0.310406in}}%
\pgfpathquadraticcurveto{\pgfqpoint{4.323938in}{0.308047in}}{\pgfqpoint{4.318084in}{0.306804in}}%
\pgfpathquadraticcurveto{\pgfqpoint{4.312230in}{0.305561in}}{\pgfqpoint{4.306214in}{0.305561in}}%
\pgfpathquadraticcurveto{\pgfqpoint{4.291120in}{0.305561in}}{\pgfqpoint{4.282312in}{0.314333in}}%
\pgfpathquadraticcurveto{\pgfqpoint{4.273504in}{0.323123in}}{\pgfqpoint{4.273504in}{0.338109in}}%
\pgfpathquadraticcurveto{\pgfqpoint{4.273504in}{0.353581in}}{\pgfqpoint{4.281862in}{0.362659in}}%
\pgfpathquadraticcurveto{\pgfqpoint{4.290219in}{0.371756in}}{\pgfqpoint{4.304413in}{0.371756in}}%
\pgfpathquadraticcurveto{\pgfqpoint{4.317129in}{0.371756in}}{\pgfqpoint{4.324532in}{0.363560in}}%
\pgfpathquadraticcurveto{\pgfqpoint{4.331935in}{0.355383in}}{\pgfqpoint{4.331935in}{0.341315in}}%
\pgfpathlineto{\pgfqpoint{4.331935in}{0.341315in}}%
\pgfpathclose%
\pgfpathmoveto{\pgfqpoint{4.321578in}{0.344359in}}%
\pgfpathquadraticcurveto{\pgfqpoint{4.321470in}{0.352843in}}{\pgfqpoint{4.316823in}{0.357904in}}%
\pgfpathquadraticcurveto{\pgfqpoint{4.312176in}{0.362984in}}{\pgfqpoint{4.304521in}{0.362984in}}%
\pgfpathquadraticcurveto{\pgfqpoint{4.295857in}{0.362984in}}{\pgfqpoint{4.290651in}{0.358084in}}%
\pgfpathquadraticcurveto{\pgfqpoint{4.285446in}{0.353185in}}{\pgfqpoint{4.284653in}{0.344287in}}%
\pgfpathlineto{\pgfqpoint{4.321578in}{0.344359in}}%
\pgfpathlineto{\pgfqpoint{4.321578in}{0.344359in}}%
\pgfpathclose%
\pgfpathmoveto{\pgfqpoint{4.390421in}{0.360678in}}%
\pgfpathlineto{\pgfqpoint{4.390421in}{0.394793in}}%
\pgfpathlineto{\pgfqpoint{4.400778in}{0.394793in}}%
\pgfpathlineto{\pgfqpoint{4.400778in}{0.307200in}}%
\pgfpathlineto{\pgfqpoint{4.390421in}{0.307200in}}%
\pgfpathlineto{\pgfqpoint{4.390421in}{0.316656in}}%
\pgfpathquadraticcurveto{\pgfqpoint{4.387161in}{0.311037in}}{\pgfqpoint{4.382171in}{0.308299in}}%
\pgfpathquadraticcurveto{\pgfqpoint{4.377200in}{0.305561in}}{\pgfqpoint{4.370211in}{0.305561in}}%
\pgfpathquadraticcurveto{\pgfqpoint{4.358791in}{0.305561in}}{\pgfqpoint{4.351605in}{0.314675in}}%
\pgfpathquadraticcurveto{\pgfqpoint{4.344436in}{0.323807in}}{\pgfqpoint{4.344436in}{0.338667in}}%
\pgfpathquadraticcurveto{\pgfqpoint{4.344436in}{0.353527in}}{\pgfqpoint{4.351605in}{0.362641in}}%
\pgfpathquadraticcurveto{\pgfqpoint{4.358791in}{0.371756in}}{\pgfqpoint{4.370211in}{0.371756in}}%
\pgfpathquadraticcurveto{\pgfqpoint{4.377200in}{0.371756in}}{\pgfqpoint{4.382171in}{0.369018in}}%
\pgfpathquadraticcurveto{\pgfqpoint{4.387161in}{0.366298in}}{\pgfqpoint{4.390421in}{0.360678in}}%
\pgfpathlineto{\pgfqpoint{4.390421in}{0.360678in}}%
\pgfpathclose%
\pgfpathmoveto{\pgfqpoint{4.355135in}{0.338667in}}%
\pgfpathquadraticcurveto{\pgfqpoint{4.355135in}{0.327248in}}{\pgfqpoint{4.359836in}{0.320745in}}%
\pgfpathquadraticcurveto{\pgfqpoint{4.364537in}{0.314243in}}{\pgfqpoint{4.372751in}{0.314243in}}%
\pgfpathquadraticcurveto{\pgfqpoint{4.380964in}{0.314243in}}{\pgfqpoint{4.385684in}{0.320745in}}%
\pgfpathquadraticcurveto{\pgfqpoint{4.390421in}{0.327248in}}{\pgfqpoint{4.390421in}{0.338667in}}%
\pgfpathquadraticcurveto{\pgfqpoint{4.390421in}{0.350087in}}{\pgfqpoint{4.385684in}{0.356589in}}%
\pgfpathquadraticcurveto{\pgfqpoint{4.380964in}{0.363092in}}{\pgfqpoint{4.372751in}{0.363092in}}%
\pgfpathquadraticcurveto{\pgfqpoint{4.364537in}{0.363092in}}{\pgfqpoint{4.359836in}{0.356589in}}%
\pgfpathquadraticcurveto{\pgfqpoint{4.355135in}{0.350087in}}{\pgfqpoint{4.355135in}{0.338667in}}%
\pgfpathlineto{\pgfqpoint{4.355135in}{0.338667in}}%
\pgfpathclose%
\pgfpathmoveto{\pgfqpoint{4.447555in}{0.391245in}}%
\pgfpathlineto{\pgfqpoint{4.518649in}{0.391245in}}%
\pgfpathlineto{\pgfqpoint{4.518649in}{0.381662in}}%
\pgfpathlineto{\pgfqpoint{4.488821in}{0.381662in}}%
\pgfpathlineto{\pgfqpoint{4.488821in}{0.307200in}}%
\pgfpathlineto{\pgfqpoint{4.477401in}{0.307200in}}%
\pgfpathlineto{\pgfqpoint{4.477401in}{0.381662in}}%
\pgfpathlineto{\pgfqpoint{4.447555in}{0.381662in}}%
\pgfpathlineto{\pgfqpoint{4.447555in}{0.391245in}}%
\pgfpathlineto{\pgfqpoint{4.447555in}{0.391245in}}%
\pgfpathclose%
\pgfpathmoveto{\pgfqpoint{4.522144in}{0.391245in}}%
\pgfpathlineto{\pgfqpoint{4.533617in}{0.391245in}}%
\pgfpathlineto{\pgfqpoint{4.551287in}{0.320205in}}%
\pgfpathlineto{\pgfqpoint{4.568903in}{0.391245in}}%
\pgfpathlineto{\pgfqpoint{4.581692in}{0.391245in}}%
\pgfpathlineto{\pgfqpoint{4.599362in}{0.320205in}}%
\pgfpathlineto{\pgfqpoint{4.616978in}{0.391245in}}%
\pgfpathlineto{\pgfqpoint{4.628523in}{0.391245in}}%
\pgfpathlineto{\pgfqpoint{4.607413in}{0.307200in}}%
\pgfpathlineto{\pgfqpoint{4.593112in}{0.307200in}}%
\pgfpathlineto{\pgfqpoint{4.575388in}{0.380149in}}%
\pgfpathlineto{\pgfqpoint{4.557484in}{0.307200in}}%
\pgfpathlineto{\pgfqpoint{4.543182in}{0.307200in}}%
\pgfpathlineto{\pgfqpoint{4.522144in}{0.391245in}}%
\pgfpathlineto{\pgfqpoint{4.522144in}{0.391245in}}%
\pgfpathclose%
\pgfpathmoveto{\pgfqpoint{4.643600in}{0.391245in}}%
\pgfpathlineto{\pgfqpoint{4.691890in}{0.391245in}}%
\pgfpathlineto{\pgfqpoint{4.691890in}{0.381662in}}%
\pgfpathlineto{\pgfqpoint{4.654965in}{0.381662in}}%
\pgfpathlineto{\pgfqpoint{4.654965in}{0.356896in}}%
\pgfpathlineto{\pgfqpoint{4.688288in}{0.356896in}}%
\pgfpathlineto{\pgfqpoint{4.688288in}{0.347331in}}%
\pgfpathlineto{\pgfqpoint{4.654965in}{0.347331in}}%
\pgfpathlineto{\pgfqpoint{4.654965in}{0.307200in}}%
\pgfpathlineto{\pgfqpoint{4.643600in}{0.307200in}}%
\pgfpathlineto{\pgfqpoint{4.643600in}{0.391245in}}%
\pgfpathlineto{\pgfqpoint{4.643600in}{0.391245in}}%
\pgfpathclose%
\pgfpathmoveto{\pgfqpoint{4.709902in}{0.391245in}}%
\pgfpathlineto{\pgfqpoint{4.763038in}{0.391245in}}%
\pgfpathlineto{\pgfqpoint{4.763038in}{0.381662in}}%
\pgfpathlineto{\pgfqpoint{4.721268in}{0.381662in}}%
\pgfpathlineto{\pgfqpoint{4.721268in}{0.356787in}}%
\pgfpathlineto{\pgfqpoint{4.761291in}{0.356787in}}%
\pgfpathlineto{\pgfqpoint{4.761291in}{0.347223in}}%
\pgfpathlineto{\pgfqpoint{4.721268in}{0.347223in}}%
\pgfpathlineto{\pgfqpoint{4.721268in}{0.316764in}}%
\pgfpathlineto{\pgfqpoint{4.764047in}{0.316764in}}%
\pgfpathlineto{\pgfqpoint{4.764047in}{0.307200in}}%
\pgfpathlineto{\pgfqpoint{4.709902in}{0.307200in}}%
\pgfpathlineto{\pgfqpoint{4.709902in}{0.391245in}}%
\pgfpathlineto{\pgfqpoint{4.709902in}{0.391245in}}%
\pgfpathclose%
\pgfpathmoveto{\pgfqpoint{4.784941in}{0.366802in}}%
\pgfpathlineto{\pgfqpoint{4.796811in}{0.366802in}}%
\pgfpathlineto{\pgfqpoint{4.796811in}{0.352519in}}%
\pgfpathlineto{\pgfqpoint{4.784941in}{0.352519in}}%
\pgfpathlineto{\pgfqpoint{4.784941in}{0.366802in}}%
\pgfpathlineto{\pgfqpoint{4.784941in}{0.366802in}}%
\pgfpathclose%
\pgfpathmoveto{\pgfqpoint{4.784941in}{0.321502in}}%
\pgfpathlineto{\pgfqpoint{4.796811in}{0.321502in}}%
\pgfpathlineto{\pgfqpoint{4.796811in}{0.311811in}}%
\pgfpathlineto{\pgfqpoint{4.787589in}{0.293799in}}%
\pgfpathlineto{\pgfqpoint{4.780330in}{0.293799in}}%
\pgfpathlineto{\pgfqpoint{4.784941in}{0.311811in}}%
\pgfpathlineto{\pgfqpoint{4.784941in}{0.321502in}}%
\pgfpathlineto{\pgfqpoint{4.784941in}{0.321502in}}%
\pgfpathclose%
\pgfpathmoveto{\pgfqpoint{4.857764in}{0.370243in}}%
\pgfpathlineto{\pgfqpoint{4.868121in}{0.370243in}}%
\pgfpathlineto{\pgfqpoint{4.868121in}{0.307200in}}%
\pgfpathlineto{\pgfqpoint{4.857764in}{0.307200in}}%
\pgfpathlineto{\pgfqpoint{4.857764in}{0.370243in}}%
\pgfpathlineto{\pgfqpoint{4.857764in}{0.370243in}}%
\pgfpathclose%
\pgfpathmoveto{\pgfqpoint{4.857764in}{0.394793in}}%
\pgfpathlineto{\pgfqpoint{4.868121in}{0.394793in}}%
\pgfpathlineto{\pgfqpoint{4.868121in}{0.381662in}}%
\pgfpathlineto{\pgfqpoint{4.857764in}{0.381662in}}%
\pgfpathlineto{\pgfqpoint{4.857764in}{0.394793in}}%
\pgfpathlineto{\pgfqpoint{4.857764in}{0.394793in}}%
\pgfpathclose%
\pgfpathmoveto{\pgfqpoint{4.900039in}{0.388147in}}%
\pgfpathlineto{\pgfqpoint{4.900039in}{0.370243in}}%
\pgfpathlineto{\pgfqpoint{4.921365in}{0.370243in}}%
\pgfpathlineto{\pgfqpoint{4.921365in}{0.362191in}}%
\pgfpathlineto{\pgfqpoint{4.900039in}{0.362191in}}%
\pgfpathlineto{\pgfqpoint{4.900039in}{0.327968in}}%
\pgfpathquadraticcurveto{\pgfqpoint{4.900039in}{0.320259in}}{\pgfqpoint{4.902146in}{0.318061in}}%
\pgfpathquadraticcurveto{\pgfqpoint{4.904253in}{0.315864in}}{\pgfqpoint{4.910738in}{0.315864in}}%
\pgfpathlineto{\pgfqpoint{4.921365in}{0.315864in}}%
\pgfpathlineto{\pgfqpoint{4.921365in}{0.307200in}}%
\pgfpathlineto{\pgfqpoint{4.910738in}{0.307200in}}%
\pgfpathquadraticcurveto{\pgfqpoint{4.898742in}{0.307200in}}{\pgfqpoint{4.894185in}{0.311667in}}%
\pgfpathquadraticcurveto{\pgfqpoint{4.889628in}{0.316152in}}{\pgfqpoint{4.889628in}{0.327968in}}%
\pgfpathlineto{\pgfqpoint{4.889628in}{0.362191in}}%
\pgfpathlineto{\pgfqpoint{4.882026in}{0.362191in}}%
\pgfpathlineto{\pgfqpoint{4.882026in}{0.370243in}}%
\pgfpathlineto{\pgfqpoint{4.889628in}{0.370243in}}%
\pgfpathlineto{\pgfqpoint{4.889628in}{0.388147in}}%
\pgfpathlineto{\pgfqpoint{4.900039in}{0.388147in}}%
\pgfpathlineto{\pgfqpoint{4.900039in}{0.388147in}}%
\pgfpathclose%
\pgfpathmoveto{\pgfqpoint{4.975167in}{0.368387in}}%
\pgfpathlineto{\pgfqpoint{4.975167in}{0.358589in}}%
\pgfpathquadraticcurveto{\pgfqpoint{4.970790in}{0.360840in}}{\pgfqpoint{4.966053in}{0.361957in}}%
\pgfpathquadraticcurveto{\pgfqpoint{4.961334in}{0.363092in}}{\pgfqpoint{4.956255in}{0.363092in}}%
\pgfpathquadraticcurveto{\pgfqpoint{4.948545in}{0.363092in}}{\pgfqpoint{4.944691in}{0.360732in}}%
\pgfpathquadraticcurveto{\pgfqpoint{4.940836in}{0.358373in}}{\pgfqpoint{4.940836in}{0.353635in}}%
\pgfpathquadraticcurveto{\pgfqpoint{4.940836in}{0.350033in}}{\pgfqpoint{4.943592in}{0.347980in}}%
\pgfpathquadraticcurveto{\pgfqpoint{4.946348in}{0.345926in}}{\pgfqpoint{4.954688in}{0.344071in}}%
\pgfpathlineto{\pgfqpoint{4.958236in}{0.343278in}}%
\pgfpathquadraticcurveto{\pgfqpoint{4.969259in}{0.340919in}}{\pgfqpoint{4.973906in}{0.336614in}}%
\pgfpathquadraticcurveto{\pgfqpoint{4.978554in}{0.332309in}}{\pgfqpoint{4.978554in}{0.324600in}}%
\pgfpathquadraticcurveto{\pgfqpoint{4.978554in}{0.315810in}}{\pgfqpoint{4.971601in}{0.310676in}}%
\pgfpathquadraticcurveto{\pgfqpoint{4.964648in}{0.305561in}}{\pgfqpoint{4.952490in}{0.305561in}}%
\pgfpathquadraticcurveto{\pgfqpoint{4.947429in}{0.305561in}}{\pgfqpoint{4.941935in}{0.306552in}}%
\pgfpathquadraticcurveto{\pgfqpoint{4.936441in}{0.307542in}}{\pgfqpoint{4.930371in}{0.309506in}}%
\pgfpathlineto{\pgfqpoint{4.930371in}{0.320205in}}%
\pgfpathquadraticcurveto{\pgfqpoint{4.936117in}{0.317215in}}{\pgfqpoint{4.941683in}{0.315720in}}%
\pgfpathquadraticcurveto{\pgfqpoint{4.947248in}{0.314243in}}{\pgfqpoint{4.952724in}{0.314243in}}%
\pgfpathquadraticcurveto{\pgfqpoint{4.960037in}{0.314243in}}{\pgfqpoint{4.963964in}{0.316746in}}%
\pgfpathquadraticcurveto{\pgfqpoint{4.967908in}{0.319250in}}{\pgfqpoint{4.967908in}{0.323807in}}%
\pgfpathquadraticcurveto{\pgfqpoint{4.967908in}{0.328022in}}{\pgfqpoint{4.965063in}{0.330274in}}%
\pgfpathquadraticcurveto{\pgfqpoint{4.962235in}{0.332525in}}{\pgfqpoint{4.952598in}{0.334614in}}%
\pgfpathlineto{\pgfqpoint{4.948996in}{0.335461in}}%
\pgfpathquadraticcurveto{\pgfqpoint{4.939377in}{0.337478in}}{\pgfqpoint{4.935090in}{0.341675in}}%
\pgfpathquadraticcurveto{\pgfqpoint{4.930821in}{0.345872in}}{\pgfqpoint{4.930821in}{0.353185in}}%
\pgfpathquadraticcurveto{\pgfqpoint{4.930821in}{0.362083in}}{\pgfqpoint{4.937126in}{0.366910in}}%
\pgfpathquadraticcurveto{\pgfqpoint{4.943430in}{0.371756in}}{\pgfqpoint{4.955030in}{0.371756in}}%
\pgfpathquadraticcurveto{\pgfqpoint{4.960758in}{0.371756in}}{\pgfqpoint{4.965819in}{0.370909in}}%
\pgfpathquadraticcurveto{\pgfqpoint{4.970898in}{0.370080in}}{\pgfqpoint{4.975167in}{0.368387in}}%
\pgfpathlineto{\pgfqpoint{4.975167in}{0.368387in}}%
\pgfpathclose%
\pgfpathmoveto{\pgfqpoint{5.068200in}{0.360570in}}%
\pgfpathquadraticcurveto{\pgfqpoint{5.066453in}{0.361579in}}{\pgfqpoint{5.064400in}{0.362047in}}%
\pgfpathquadraticcurveto{\pgfqpoint{5.062346in}{0.362533in}}{\pgfqpoint{5.059878in}{0.362533in}}%
\pgfpathquadraticcurveto{\pgfqpoint{5.051089in}{0.362533in}}{\pgfqpoint{5.046387in}{0.356823in}}%
\pgfpathquadraticcurveto{\pgfqpoint{5.041686in}{0.351114in}}{\pgfqpoint{5.041686in}{0.340414in}}%
\pgfpathlineto{\pgfqpoint{5.041686in}{0.307200in}}%
\pgfpathlineto{\pgfqpoint{5.031275in}{0.307200in}}%
\pgfpathlineto{\pgfqpoint{5.031275in}{0.370243in}}%
\pgfpathlineto{\pgfqpoint{5.041686in}{0.370243in}}%
\pgfpathlineto{\pgfqpoint{5.041686in}{0.360444in}}%
\pgfpathquadraticcurveto{\pgfqpoint{5.044964in}{0.366190in}}{\pgfqpoint{5.050188in}{0.368964in}}%
\pgfpathquadraticcurveto{\pgfqpoint{5.055429in}{0.371756in}}{\pgfqpoint{5.062923in}{0.371756in}}%
\pgfpathquadraticcurveto{\pgfqpoint{5.063985in}{0.371756in}}{\pgfqpoint{5.065282in}{0.371611in}}%
\pgfpathquadraticcurveto{\pgfqpoint{5.066579in}{0.371485in}}{\pgfqpoint{5.068146in}{0.371197in}}%
\pgfpathlineto{\pgfqpoint{5.068200in}{0.360570in}}%
\pgfpathlineto{\pgfqpoint{5.068200in}{0.360570in}}%
\pgfpathclose%
\pgfpathmoveto{\pgfqpoint{5.130450in}{0.341315in}}%
\pgfpathlineto{\pgfqpoint{5.130450in}{0.336254in}}%
\pgfpathlineto{\pgfqpoint{5.082826in}{0.336254in}}%
\pgfpathquadraticcurveto{\pgfqpoint{5.083510in}{0.325554in}}{\pgfqpoint{5.089274in}{0.319953in}}%
\pgfpathquadraticcurveto{\pgfqpoint{5.095038in}{0.314351in}}{\pgfqpoint{5.105341in}{0.314351in}}%
\pgfpathquadraticcurveto{\pgfqpoint{5.111303in}{0.314351in}}{\pgfqpoint{5.116905in}{0.315810in}}%
\pgfpathquadraticcurveto{\pgfqpoint{5.122507in}{0.317269in}}{\pgfqpoint{5.128036in}{0.320205in}}%
\pgfpathlineto{\pgfqpoint{5.128036in}{0.310406in}}%
\pgfpathquadraticcurveto{\pgfqpoint{5.122453in}{0.308047in}}{\pgfqpoint{5.116599in}{0.306804in}}%
\pgfpathquadraticcurveto{\pgfqpoint{5.110745in}{0.305561in}}{\pgfqpoint{5.104729in}{0.305561in}}%
\pgfpathquadraticcurveto{\pgfqpoint{5.089635in}{0.305561in}}{\pgfqpoint{5.080827in}{0.314333in}}%
\pgfpathquadraticcurveto{\pgfqpoint{5.072019in}{0.323123in}}{\pgfqpoint{5.072019in}{0.338109in}}%
\pgfpathquadraticcurveto{\pgfqpoint{5.072019in}{0.353581in}}{\pgfqpoint{5.080376in}{0.362659in}}%
\pgfpathquadraticcurveto{\pgfqpoint{5.088734in}{0.371756in}}{\pgfqpoint{5.102928in}{0.371756in}}%
\pgfpathquadraticcurveto{\pgfqpoint{5.115644in}{0.371756in}}{\pgfqpoint{5.123047in}{0.363560in}}%
\pgfpathquadraticcurveto{\pgfqpoint{5.130450in}{0.355383in}}{\pgfqpoint{5.130450in}{0.341315in}}%
\pgfpathlineto{\pgfqpoint{5.130450in}{0.341315in}}%
\pgfpathclose%
\pgfpathmoveto{\pgfqpoint{5.120093in}{0.344359in}}%
\pgfpathquadraticcurveto{\pgfqpoint{5.119985in}{0.352843in}}{\pgfqpoint{5.115338in}{0.357904in}}%
\pgfpathquadraticcurveto{\pgfqpoint{5.110691in}{0.362984in}}{\pgfqpoint{5.103036in}{0.362984in}}%
\pgfpathquadraticcurveto{\pgfqpoint{5.094372in}{0.362984in}}{\pgfqpoint{5.089166in}{0.358084in}}%
\pgfpathquadraticcurveto{\pgfqpoint{5.083961in}{0.353185in}}{\pgfqpoint{5.083168in}{0.344287in}}%
\pgfpathlineto{\pgfqpoint{5.120093in}{0.344359in}}%
\pgfpathlineto{\pgfqpoint{5.120093in}{0.344359in}}%
\pgfpathclose%
\pgfpathmoveto{\pgfqpoint{5.187639in}{0.368387in}}%
\pgfpathlineto{\pgfqpoint{5.187639in}{0.358589in}}%
\pgfpathquadraticcurveto{\pgfqpoint{5.183262in}{0.360840in}}{\pgfqpoint{5.178525in}{0.361957in}}%
\pgfpathquadraticcurveto{\pgfqpoint{5.173805in}{0.363092in}}{\pgfqpoint{5.168726in}{0.363092in}}%
\pgfpathquadraticcurveto{\pgfqpoint{5.161017in}{0.363092in}}{\pgfqpoint{5.157162in}{0.360732in}}%
\pgfpathquadraticcurveto{\pgfqpoint{5.153308in}{0.358373in}}{\pgfqpoint{5.153308in}{0.353635in}}%
\pgfpathquadraticcurveto{\pgfqpoint{5.153308in}{0.350033in}}{\pgfqpoint{5.156063in}{0.347980in}}%
\pgfpathquadraticcurveto{\pgfqpoint{5.158819in}{0.345926in}}{\pgfqpoint{5.167159in}{0.344071in}}%
\pgfpathlineto{\pgfqpoint{5.170707in}{0.343278in}}%
\pgfpathquadraticcurveto{\pgfqpoint{5.181731in}{0.340919in}}{\pgfqpoint{5.186378in}{0.336614in}}%
\pgfpathquadraticcurveto{\pgfqpoint{5.191025in}{0.332309in}}{\pgfqpoint{5.191025in}{0.324600in}}%
\pgfpathquadraticcurveto{\pgfqpoint{5.191025in}{0.315810in}}{\pgfqpoint{5.184072in}{0.310676in}}%
\pgfpathquadraticcurveto{\pgfqpoint{5.177120in}{0.305561in}}{\pgfqpoint{5.164961in}{0.305561in}}%
\pgfpathquadraticcurveto{\pgfqpoint{5.159900in}{0.305561in}}{\pgfqpoint{5.154406in}{0.306552in}}%
\pgfpathquadraticcurveto{\pgfqpoint{5.148913in}{0.307542in}}{\pgfqpoint{5.142842in}{0.309506in}}%
\pgfpathlineto{\pgfqpoint{5.142842in}{0.320205in}}%
\pgfpathquadraticcurveto{\pgfqpoint{5.148588in}{0.317215in}}{\pgfqpoint{5.154154in}{0.315720in}}%
\pgfpathquadraticcurveto{\pgfqpoint{5.159720in}{0.314243in}}{\pgfqpoint{5.165196in}{0.314243in}}%
\pgfpathquadraticcurveto{\pgfqpoint{5.172508in}{0.314243in}}{\pgfqpoint{5.176435in}{0.316746in}}%
\pgfpathquadraticcurveto{\pgfqpoint{5.180380in}{0.319250in}}{\pgfqpoint{5.180380in}{0.323807in}}%
\pgfpathquadraticcurveto{\pgfqpoint{5.180380in}{0.328022in}}{\pgfqpoint{5.177534in}{0.330274in}}%
\pgfpathquadraticcurveto{\pgfqpoint{5.174706in}{0.332525in}}{\pgfqpoint{5.165069in}{0.334614in}}%
\pgfpathlineto{\pgfqpoint{5.161467in}{0.335461in}}%
\pgfpathquadraticcurveto{\pgfqpoint{5.151849in}{0.337478in}}{\pgfqpoint{5.147562in}{0.341675in}}%
\pgfpathquadraticcurveto{\pgfqpoint{5.143293in}{0.345872in}}{\pgfqpoint{5.143293in}{0.353185in}}%
\pgfpathquadraticcurveto{\pgfqpoint{5.143293in}{0.362083in}}{\pgfqpoint{5.149597in}{0.366910in}}%
\pgfpathquadraticcurveto{\pgfqpoint{5.155901in}{0.371756in}}{\pgfqpoint{5.167501in}{0.371756in}}%
\pgfpathquadraticcurveto{\pgfqpoint{5.173229in}{0.371756in}}{\pgfqpoint{5.178290in}{0.370909in}}%
\pgfpathquadraticcurveto{\pgfqpoint{5.183370in}{0.370080in}}{\pgfqpoint{5.187639in}{0.368387in}}%
\pgfpathlineto{\pgfqpoint{5.187639in}{0.368387in}}%
\pgfpathclose%
\pgfpathmoveto{\pgfqpoint{5.217755in}{0.388147in}}%
\pgfpathlineto{\pgfqpoint{5.217755in}{0.370243in}}%
\pgfpathlineto{\pgfqpoint{5.239081in}{0.370243in}}%
\pgfpathlineto{\pgfqpoint{5.239081in}{0.362191in}}%
\pgfpathlineto{\pgfqpoint{5.217755in}{0.362191in}}%
\pgfpathlineto{\pgfqpoint{5.217755in}{0.327968in}}%
\pgfpathquadraticcurveto{\pgfqpoint{5.217755in}{0.320259in}}{\pgfqpoint{5.219862in}{0.318061in}}%
\pgfpathquadraticcurveto{\pgfqpoint{5.221970in}{0.315864in}}{\pgfqpoint{5.228454in}{0.315864in}}%
\pgfpathlineto{\pgfqpoint{5.239081in}{0.315864in}}%
\pgfpathlineto{\pgfqpoint{5.239081in}{0.307200in}}%
\pgfpathlineto{\pgfqpoint{5.228454in}{0.307200in}}%
\pgfpathquadraticcurveto{\pgfqpoint{5.216458in}{0.307200in}}{\pgfqpoint{5.211901in}{0.311667in}}%
\pgfpathquadraticcurveto{\pgfqpoint{5.207344in}{0.316152in}}{\pgfqpoint{5.207344in}{0.327968in}}%
\pgfpathlineto{\pgfqpoint{5.207344in}{0.362191in}}%
\pgfpathlineto{\pgfqpoint{5.199743in}{0.362191in}}%
\pgfpathlineto{\pgfqpoint{5.199743in}{0.370243in}}%
\pgfpathlineto{\pgfqpoint{5.207344in}{0.370243in}}%
\pgfpathlineto{\pgfqpoint{5.207344in}{0.388147in}}%
\pgfpathlineto{\pgfqpoint{5.217755in}{0.388147in}}%
\pgfpathlineto{\pgfqpoint{5.217755in}{0.388147in}}%
\pgfpathclose%
\pgfpathmoveto{\pgfqpoint{5.289227in}{0.360570in}}%
\pgfpathquadraticcurveto{\pgfqpoint{5.287480in}{0.361579in}}{\pgfqpoint{5.285427in}{0.362047in}}%
\pgfpathquadraticcurveto{\pgfqpoint{5.283373in}{0.362533in}}{\pgfqpoint{5.280906in}{0.362533in}}%
\pgfpathquadraticcurveto{\pgfqpoint{5.272116in}{0.362533in}}{\pgfqpoint{5.267414in}{0.356823in}}%
\pgfpathquadraticcurveto{\pgfqpoint{5.262713in}{0.351114in}}{\pgfqpoint{5.262713in}{0.340414in}}%
\pgfpathlineto{\pgfqpoint{5.262713in}{0.307200in}}%
\pgfpathlineto{\pgfqpoint{5.252302in}{0.307200in}}%
\pgfpathlineto{\pgfqpoint{5.252302in}{0.370243in}}%
\pgfpathlineto{\pgfqpoint{5.262713in}{0.370243in}}%
\pgfpathlineto{\pgfqpoint{5.262713in}{0.360444in}}%
\pgfpathquadraticcurveto{\pgfqpoint{5.265992in}{0.366190in}}{\pgfqpoint{5.271215in}{0.368964in}}%
\pgfpathquadraticcurveto{\pgfqpoint{5.276457in}{0.371756in}}{\pgfqpoint{5.283950in}{0.371756in}}%
\pgfpathquadraticcurveto{\pgfqpoint{5.285012in}{0.371756in}}{\pgfqpoint{5.286309in}{0.371611in}}%
\pgfpathquadraticcurveto{\pgfqpoint{5.287606in}{0.371485in}}{\pgfqpoint{5.289173in}{0.371197in}}%
\pgfpathlineto{\pgfqpoint{5.289227in}{0.360570in}}%
\pgfpathlineto{\pgfqpoint{5.289227in}{0.360570in}}%
\pgfpathclose%
\pgfpathmoveto{\pgfqpoint{5.300089in}{0.370243in}}%
\pgfpathlineto{\pgfqpoint{5.310446in}{0.370243in}}%
\pgfpathlineto{\pgfqpoint{5.310446in}{0.307200in}}%
\pgfpathlineto{\pgfqpoint{5.300089in}{0.307200in}}%
\pgfpathlineto{\pgfqpoint{5.300089in}{0.370243in}}%
\pgfpathlineto{\pgfqpoint{5.300089in}{0.370243in}}%
\pgfpathclose%
\pgfpathmoveto{\pgfqpoint{5.300089in}{0.394793in}}%
\pgfpathlineto{\pgfqpoint{5.310446in}{0.394793in}}%
\pgfpathlineto{\pgfqpoint{5.310446in}{0.381662in}}%
\pgfpathlineto{\pgfqpoint{5.300089in}{0.381662in}}%
\pgfpathlineto{\pgfqpoint{5.300089in}{0.394793in}}%
\pgfpathlineto{\pgfqpoint{5.300089in}{0.394793in}}%
\pgfpathclose%
\pgfpathmoveto{\pgfqpoint{5.377487in}{0.367829in}}%
\pgfpathlineto{\pgfqpoint{5.377487in}{0.358138in}}%
\pgfpathquadraticcurveto{\pgfqpoint{5.373092in}{0.360570in}}{\pgfqpoint{5.368679in}{0.361777in}}%
\pgfpathquadraticcurveto{\pgfqpoint{5.364266in}{0.362984in}}{\pgfqpoint{5.359763in}{0.362984in}}%
\pgfpathquadraticcurveto{\pgfqpoint{5.349676in}{0.362984in}}{\pgfqpoint{5.344092in}{0.356589in}}%
\pgfpathquadraticcurveto{\pgfqpoint{5.338526in}{0.350213in}}{\pgfqpoint{5.338526in}{0.338667in}}%
\pgfpathquadraticcurveto{\pgfqpoint{5.338526in}{0.327121in}}{\pgfqpoint{5.344092in}{0.320727in}}%
\pgfpathquadraticcurveto{\pgfqpoint{5.349676in}{0.314351in}}{\pgfqpoint{5.359763in}{0.314351in}}%
\pgfpathquadraticcurveto{\pgfqpoint{5.364266in}{0.314351in}}{\pgfqpoint{5.368679in}{0.315558in}}%
\pgfpathquadraticcurveto{\pgfqpoint{5.373092in}{0.316764in}}{\pgfqpoint{5.377487in}{0.319196in}}%
\pgfpathlineto{\pgfqpoint{5.377487in}{0.309614in}}%
\pgfpathquadraticcurveto{\pgfqpoint{5.373146in}{0.307596in}}{\pgfqpoint{5.368499in}{0.306588in}}%
\pgfpathquadraticcurveto{\pgfqpoint{5.363870in}{0.305561in}}{\pgfqpoint{5.358628in}{0.305561in}}%
\pgfpathquadraticcurveto{\pgfqpoint{5.344380in}{0.305561in}}{\pgfqpoint{5.335987in}{0.314513in}}%
\pgfpathquadraticcurveto{\pgfqpoint{5.327611in}{0.323465in}}{\pgfqpoint{5.327611in}{0.338667in}}%
\pgfpathquadraticcurveto{\pgfqpoint{5.327611in}{0.354086in}}{\pgfqpoint{5.336077in}{0.362912in}}%
\pgfpathquadraticcurveto{\pgfqpoint{5.344561in}{0.371756in}}{\pgfqpoint{5.359312in}{0.371756in}}%
\pgfpathquadraticcurveto{\pgfqpoint{5.364086in}{0.371756in}}{\pgfqpoint{5.368643in}{0.370765in}}%
\pgfpathquadraticcurveto{\pgfqpoint{5.373200in}{0.369792in}}{\pgfqpoint{5.377487in}{0.367829in}}%
\pgfpathlineto{\pgfqpoint{5.377487in}{0.367829in}}%
\pgfpathclose%
\pgfpathmoveto{\pgfqpoint{5.405748in}{0.388147in}}%
\pgfpathlineto{\pgfqpoint{5.405748in}{0.370243in}}%
\pgfpathlineto{\pgfqpoint{5.427074in}{0.370243in}}%
\pgfpathlineto{\pgfqpoint{5.427074in}{0.362191in}}%
\pgfpathlineto{\pgfqpoint{5.405748in}{0.362191in}}%
\pgfpathlineto{\pgfqpoint{5.405748in}{0.327968in}}%
\pgfpathquadraticcurveto{\pgfqpoint{5.405748in}{0.320259in}}{\pgfqpoint{5.407855in}{0.318061in}}%
\pgfpathquadraticcurveto{\pgfqpoint{5.409963in}{0.315864in}}{\pgfqpoint{5.416447in}{0.315864in}}%
\pgfpathlineto{\pgfqpoint{5.427074in}{0.315864in}}%
\pgfpathlineto{\pgfqpoint{5.427074in}{0.307200in}}%
\pgfpathlineto{\pgfqpoint{5.416447in}{0.307200in}}%
\pgfpathquadraticcurveto{\pgfqpoint{5.404451in}{0.307200in}}{\pgfqpoint{5.399894in}{0.311667in}}%
\pgfpathquadraticcurveto{\pgfqpoint{5.395337in}{0.316152in}}{\pgfqpoint{5.395337in}{0.327968in}}%
\pgfpathlineto{\pgfqpoint{5.395337in}{0.362191in}}%
\pgfpathlineto{\pgfqpoint{5.387736in}{0.362191in}}%
\pgfpathlineto{\pgfqpoint{5.387736in}{0.370243in}}%
\pgfpathlineto{\pgfqpoint{5.395337in}{0.370243in}}%
\pgfpathlineto{\pgfqpoint{5.395337in}{0.388147in}}%
\pgfpathlineto{\pgfqpoint{5.405748in}{0.388147in}}%
\pgfpathlineto{\pgfqpoint{5.405748in}{0.388147in}}%
\pgfpathclose%
\pgfpathmoveto{\pgfqpoint{5.494620in}{0.341315in}}%
\pgfpathlineto{\pgfqpoint{5.494620in}{0.336254in}}%
\pgfpathlineto{\pgfqpoint{5.446996in}{0.336254in}}%
\pgfpathquadraticcurveto{\pgfqpoint{5.447680in}{0.325554in}}{\pgfqpoint{5.453444in}{0.319953in}}%
\pgfpathquadraticcurveto{\pgfqpoint{5.459208in}{0.314351in}}{\pgfqpoint{5.469511in}{0.314351in}}%
\pgfpathquadraticcurveto{\pgfqpoint{5.475473in}{0.314351in}}{\pgfqpoint{5.481075in}{0.315810in}}%
\pgfpathquadraticcurveto{\pgfqpoint{5.486676in}{0.317269in}}{\pgfqpoint{5.492206in}{0.320205in}}%
\pgfpathlineto{\pgfqpoint{5.492206in}{0.310406in}}%
\pgfpathquadraticcurveto{\pgfqpoint{5.486622in}{0.308047in}}{\pgfqpoint{5.480768in}{0.306804in}}%
\pgfpathquadraticcurveto{\pgfqpoint{5.474914in}{0.305561in}}{\pgfqpoint{5.468898in}{0.305561in}}%
\pgfpathquadraticcurveto{\pgfqpoint{5.453804in}{0.305561in}}{\pgfqpoint{5.444996in}{0.314333in}}%
\pgfpathquadraticcurveto{\pgfqpoint{5.436188in}{0.323123in}}{\pgfqpoint{5.436188in}{0.338109in}}%
\pgfpathquadraticcurveto{\pgfqpoint{5.436188in}{0.353581in}}{\pgfqpoint{5.444546in}{0.362659in}}%
\pgfpathquadraticcurveto{\pgfqpoint{5.452904in}{0.371756in}}{\pgfqpoint{5.467097in}{0.371756in}}%
\pgfpathquadraticcurveto{\pgfqpoint{5.479814in}{0.371756in}}{\pgfqpoint{5.487217in}{0.363560in}}%
\pgfpathquadraticcurveto{\pgfqpoint{5.494620in}{0.355383in}}{\pgfqpoint{5.494620in}{0.341315in}}%
\pgfpathlineto{\pgfqpoint{5.494620in}{0.341315in}}%
\pgfpathclose%
\pgfpathmoveto{\pgfqpoint{5.484263in}{0.344359in}}%
\pgfpathquadraticcurveto{\pgfqpoint{5.484155in}{0.352843in}}{\pgfqpoint{5.479508in}{0.357904in}}%
\pgfpathquadraticcurveto{\pgfqpoint{5.474860in}{0.362984in}}{\pgfqpoint{5.467205in}{0.362984in}}%
\pgfpathquadraticcurveto{\pgfqpoint{5.458541in}{0.362984in}}{\pgfqpoint{5.453336in}{0.358084in}}%
\pgfpathquadraticcurveto{\pgfqpoint{5.448130in}{0.353185in}}{\pgfqpoint{5.447338in}{0.344287in}}%
\pgfpathlineto{\pgfqpoint{5.484263in}{0.344359in}}%
\pgfpathlineto{\pgfqpoint{5.484263in}{0.344359in}}%
\pgfpathclose%
\pgfpathmoveto{\pgfqpoint{5.553105in}{0.360678in}}%
\pgfpathlineto{\pgfqpoint{5.553105in}{0.394793in}}%
\pgfpathlineto{\pgfqpoint{5.563462in}{0.394793in}}%
\pgfpathlineto{\pgfqpoint{5.563462in}{0.307200in}}%
\pgfpathlineto{\pgfqpoint{5.553105in}{0.307200in}}%
\pgfpathlineto{\pgfqpoint{5.553105in}{0.316656in}}%
\pgfpathquadraticcurveto{\pgfqpoint{5.549845in}{0.311037in}}{\pgfqpoint{5.544856in}{0.308299in}}%
\pgfpathquadraticcurveto{\pgfqpoint{5.539884in}{0.305561in}}{\pgfqpoint{5.532896in}{0.305561in}}%
\pgfpathquadraticcurveto{\pgfqpoint{5.521476in}{0.305561in}}{\pgfqpoint{5.514289in}{0.314675in}}%
\pgfpathquadraticcurveto{\pgfqpoint{5.507120in}{0.323807in}}{\pgfqpoint{5.507120in}{0.338667in}}%
\pgfpathquadraticcurveto{\pgfqpoint{5.507120in}{0.353527in}}{\pgfqpoint{5.514289in}{0.362641in}}%
\pgfpathquadraticcurveto{\pgfqpoint{5.521476in}{0.371756in}}{\pgfqpoint{5.532896in}{0.371756in}}%
\pgfpathquadraticcurveto{\pgfqpoint{5.539884in}{0.371756in}}{\pgfqpoint{5.544856in}{0.369018in}}%
\pgfpathquadraticcurveto{\pgfqpoint{5.549845in}{0.366298in}}{\pgfqpoint{5.553105in}{0.360678in}}%
\pgfpathlineto{\pgfqpoint{5.553105in}{0.360678in}}%
\pgfpathclose%
\pgfpathmoveto{\pgfqpoint{5.517819in}{0.338667in}}%
\pgfpathquadraticcurveto{\pgfqpoint{5.517819in}{0.327248in}}{\pgfqpoint{5.522521in}{0.320745in}}%
\pgfpathquadraticcurveto{\pgfqpoint{5.527222in}{0.314243in}}{\pgfqpoint{5.535435in}{0.314243in}}%
\pgfpathquadraticcurveto{\pgfqpoint{5.543649in}{0.314243in}}{\pgfqpoint{5.548368in}{0.320745in}}%
\pgfpathquadraticcurveto{\pgfqpoint{5.553105in}{0.327248in}}{\pgfqpoint{5.553105in}{0.338667in}}%
\pgfpathquadraticcurveto{\pgfqpoint{5.553105in}{0.350087in}}{\pgfqpoint{5.548368in}{0.356589in}}%
\pgfpathquadraticcurveto{\pgfqpoint{5.543649in}{0.363092in}}{\pgfqpoint{5.535435in}{0.363092in}}%
\pgfpathquadraticcurveto{\pgfqpoint{5.527222in}{0.363092in}}{\pgfqpoint{5.522521in}{0.356589in}}%
\pgfpathquadraticcurveto{\pgfqpoint{5.517819in}{0.350087in}}{\pgfqpoint{5.517819in}{0.338667in}}%
\pgfpathlineto{\pgfqpoint{5.517819in}{0.338667in}}%
\pgfpathclose%
\pgfpathmoveto{\pgfqpoint{5.631458in}{0.316656in}}%
\pgfpathlineto{\pgfqpoint{5.631458in}{0.283226in}}%
\pgfpathlineto{\pgfqpoint{5.621047in}{0.283226in}}%
\pgfpathlineto{\pgfqpoint{5.621047in}{0.370243in}}%
\pgfpathlineto{\pgfqpoint{5.631458in}{0.370243in}}%
\pgfpathlineto{\pgfqpoint{5.631458in}{0.360678in}}%
\pgfpathquadraticcurveto{\pgfqpoint{5.634736in}{0.366298in}}{\pgfqpoint{5.639708in}{0.369018in}}%
\pgfpathquadraticcurveto{\pgfqpoint{5.644697in}{0.371756in}}{\pgfqpoint{5.651614in}{0.371756in}}%
\pgfpathquadraticcurveto{\pgfqpoint{5.663105in}{0.371756in}}{\pgfqpoint{5.670274in}{0.362641in}}%
\pgfpathquadraticcurveto{\pgfqpoint{5.677461in}{0.353527in}}{\pgfqpoint{5.677461in}{0.338667in}}%
\pgfpathquadraticcurveto{\pgfqpoint{5.677461in}{0.323807in}}{\pgfqpoint{5.670274in}{0.314675in}}%
\pgfpathquadraticcurveto{\pgfqpoint{5.663105in}{0.305561in}}{\pgfqpoint{5.651614in}{0.305561in}}%
\pgfpathquadraticcurveto{\pgfqpoint{5.644697in}{0.305561in}}{\pgfqpoint{5.639708in}{0.308299in}}%
\pgfpathquadraticcurveto{\pgfqpoint{5.634736in}{0.311037in}}{\pgfqpoint{5.631458in}{0.316656in}}%
\pgfpathlineto{\pgfqpoint{5.631458in}{0.316656in}}%
\pgfpathclose%
\pgfpathmoveto{\pgfqpoint{5.666708in}{0.338667in}}%
\pgfpathquadraticcurveto{\pgfqpoint{5.666708in}{0.350087in}}{\pgfqpoint{5.662007in}{0.356589in}}%
\pgfpathquadraticcurveto{\pgfqpoint{5.657306in}{0.363092in}}{\pgfqpoint{5.649092in}{0.363092in}}%
\pgfpathquadraticcurveto{\pgfqpoint{5.640860in}{0.363092in}}{\pgfqpoint{5.636159in}{0.356589in}}%
\pgfpathquadraticcurveto{\pgfqpoint{5.631458in}{0.350087in}}{\pgfqpoint{5.631458in}{0.338667in}}%
\pgfpathquadraticcurveto{\pgfqpoint{5.631458in}{0.327248in}}{\pgfqpoint{5.636159in}{0.320745in}}%
\pgfpathquadraticcurveto{\pgfqpoint{5.640860in}{0.314243in}}{\pgfqpoint{5.649092in}{0.314243in}}%
\pgfpathquadraticcurveto{\pgfqpoint{5.657306in}{0.314243in}}{\pgfqpoint{5.662007in}{0.320745in}}%
\pgfpathquadraticcurveto{\pgfqpoint{5.666708in}{0.327248in}}{\pgfqpoint{5.666708in}{0.338667in}}%
\pgfpathlineto{\pgfqpoint{5.666708in}{0.338667in}}%
\pgfpathclose%
\pgfpathmoveto{\pgfqpoint{5.693564in}{0.332075in}}%
\pgfpathlineto{\pgfqpoint{5.693564in}{0.370243in}}%
\pgfpathlineto{\pgfqpoint{5.703921in}{0.370243in}}%
\pgfpathlineto{\pgfqpoint{5.703921in}{0.332471in}}%
\pgfpathquadraticcurveto{\pgfqpoint{5.703921in}{0.323519in}}{\pgfqpoint{5.707397in}{0.319034in}}%
\pgfpathquadraticcurveto{\pgfqpoint{5.710892in}{0.314567in}}{\pgfqpoint{5.717880in}{0.314567in}}%
\pgfpathquadraticcurveto{\pgfqpoint{5.726256in}{0.314567in}}{\pgfqpoint{5.731119in}{0.319917in}}%
\pgfpathquadraticcurveto{\pgfqpoint{5.736001in}{0.325266in}}{\pgfqpoint{5.736001in}{0.334506in}}%
\pgfpathlineto{\pgfqpoint{5.736001in}{0.370243in}}%
\pgfpathlineto{\pgfqpoint{5.746358in}{0.370243in}}%
\pgfpathlineto{\pgfqpoint{5.746358in}{0.307200in}}%
\pgfpathlineto{\pgfqpoint{5.736001in}{0.307200in}}%
\pgfpathlineto{\pgfqpoint{5.736001in}{0.316891in}}%
\pgfpathquadraticcurveto{\pgfqpoint{5.732236in}{0.311145in}}{\pgfqpoint{5.727247in}{0.308353in}}%
\pgfpathquadraticcurveto{\pgfqpoint{5.722275in}{0.305561in}}{\pgfqpoint{5.715683in}{0.305561in}}%
\pgfpathquadraticcurveto{\pgfqpoint{5.704822in}{0.305561in}}{\pgfqpoint{5.699184in}{0.312315in}}%
\pgfpathquadraticcurveto{\pgfqpoint{5.693564in}{0.319070in}}{\pgfqpoint{5.693564in}{0.332075in}}%
\pgfpathlineto{\pgfqpoint{5.693564in}{0.332075in}}%
\pgfpathclose%
\pgfpathmoveto{\pgfqpoint{5.719628in}{0.371756in}}%
\pgfpathlineto{\pgfqpoint{5.719628in}{0.371756in}}%
\pgfpathlineto{\pgfqpoint{5.719628in}{0.371756in}}%
\pgfpathclose%
\pgfpathmoveto{\pgfqpoint{5.812949in}{0.338667in}}%
\pgfpathquadraticcurveto{\pgfqpoint{5.812949in}{0.350087in}}{\pgfqpoint{5.808247in}{0.356589in}}%
\pgfpathquadraticcurveto{\pgfqpoint{5.803546in}{0.363092in}}{\pgfqpoint{5.795333in}{0.363092in}}%
\pgfpathquadraticcurveto{\pgfqpoint{5.787101in}{0.363092in}}{\pgfqpoint{5.782400in}{0.356589in}}%
\pgfpathquadraticcurveto{\pgfqpoint{5.777699in}{0.350087in}}{\pgfqpoint{5.777699in}{0.338667in}}%
\pgfpathquadraticcurveto{\pgfqpoint{5.777699in}{0.327248in}}{\pgfqpoint{5.782400in}{0.320745in}}%
\pgfpathquadraticcurveto{\pgfqpoint{5.787101in}{0.314243in}}{\pgfqpoint{5.795333in}{0.314243in}}%
\pgfpathquadraticcurveto{\pgfqpoint{5.803546in}{0.314243in}}{\pgfqpoint{5.808247in}{0.320745in}}%
\pgfpathquadraticcurveto{\pgfqpoint{5.812949in}{0.327248in}}{\pgfqpoint{5.812949in}{0.338667in}}%
\pgfpathlineto{\pgfqpoint{5.812949in}{0.338667in}}%
\pgfpathclose%
\pgfpathmoveto{\pgfqpoint{5.777699in}{0.360678in}}%
\pgfpathquadraticcurveto{\pgfqpoint{5.780977in}{0.366298in}}{\pgfqpoint{5.785948in}{0.369018in}}%
\pgfpathquadraticcurveto{\pgfqpoint{5.790938in}{0.371756in}}{\pgfqpoint{5.797854in}{0.371756in}}%
\pgfpathquadraticcurveto{\pgfqpoint{5.809346in}{0.371756in}}{\pgfqpoint{5.816515in}{0.362641in}}%
\pgfpathquadraticcurveto{\pgfqpoint{5.823702in}{0.353527in}}{\pgfqpoint{5.823702in}{0.338667in}}%
\pgfpathquadraticcurveto{\pgfqpoint{5.823702in}{0.323807in}}{\pgfqpoint{5.816515in}{0.314675in}}%
\pgfpathquadraticcurveto{\pgfqpoint{5.809346in}{0.305561in}}{\pgfqpoint{5.797854in}{0.305561in}}%
\pgfpathquadraticcurveto{\pgfqpoint{5.790938in}{0.305561in}}{\pgfqpoint{5.785948in}{0.308299in}}%
\pgfpathquadraticcurveto{\pgfqpoint{5.780977in}{0.311037in}}{\pgfqpoint{5.777699in}{0.316656in}}%
\pgfpathlineto{\pgfqpoint{5.777699in}{0.307200in}}%
\pgfpathlineto{\pgfqpoint{5.767288in}{0.307200in}}%
\pgfpathlineto{\pgfqpoint{5.767288in}{0.394793in}}%
\pgfpathlineto{\pgfqpoint{5.777699in}{0.394793in}}%
\pgfpathlineto{\pgfqpoint{5.777699in}{0.360678in}}%
\pgfpathlineto{\pgfqpoint{5.777699in}{0.360678in}}%
\pgfpathclose%
\pgfpathmoveto{\pgfqpoint{5.840867in}{0.394793in}}%
\pgfpathlineto{\pgfqpoint{5.851224in}{0.394793in}}%
\pgfpathlineto{\pgfqpoint{5.851224in}{0.307200in}}%
\pgfpathlineto{\pgfqpoint{5.840867in}{0.307200in}}%
\pgfpathlineto{\pgfqpoint{5.840867in}{0.394793in}}%
\pgfpathlineto{\pgfqpoint{5.840867in}{0.394793in}}%
\pgfpathclose%
\pgfpathmoveto{\pgfqpoint{5.872893in}{0.370243in}}%
\pgfpathlineto{\pgfqpoint{5.883250in}{0.370243in}}%
\pgfpathlineto{\pgfqpoint{5.883250in}{0.307200in}}%
\pgfpathlineto{\pgfqpoint{5.872893in}{0.307200in}}%
\pgfpathlineto{\pgfqpoint{5.872893in}{0.370243in}}%
\pgfpathlineto{\pgfqpoint{5.872893in}{0.370243in}}%
\pgfpathclose%
\pgfpathmoveto{\pgfqpoint{5.872893in}{0.394793in}}%
\pgfpathlineto{\pgfqpoint{5.883250in}{0.394793in}}%
\pgfpathlineto{\pgfqpoint{5.883250in}{0.381662in}}%
\pgfpathlineto{\pgfqpoint{5.872893in}{0.381662in}}%
\pgfpathlineto{\pgfqpoint{5.872893in}{0.394793in}}%
\pgfpathlineto{\pgfqpoint{5.872893in}{0.394793in}}%
\pgfpathclose%
\pgfpathmoveto{\pgfqpoint{5.945104in}{0.368387in}}%
\pgfpathlineto{\pgfqpoint{5.945104in}{0.358589in}}%
\pgfpathquadraticcurveto{\pgfqpoint{5.940727in}{0.360840in}}{\pgfqpoint{5.935990in}{0.361957in}}%
\pgfpathquadraticcurveto{\pgfqpoint{5.931270in}{0.363092in}}{\pgfqpoint{5.926191in}{0.363092in}}%
\pgfpathquadraticcurveto{\pgfqpoint{5.918482in}{0.363092in}}{\pgfqpoint{5.914627in}{0.360732in}}%
\pgfpathquadraticcurveto{\pgfqpoint{5.910773in}{0.358373in}}{\pgfqpoint{5.910773in}{0.353635in}}%
\pgfpathquadraticcurveto{\pgfqpoint{5.910773in}{0.350033in}}{\pgfqpoint{5.913528in}{0.347980in}}%
\pgfpathquadraticcurveto{\pgfqpoint{5.916284in}{0.345926in}}{\pgfqpoint{5.924624in}{0.344071in}}%
\pgfpathlineto{\pgfqpoint{5.928172in}{0.343278in}}%
\pgfpathquadraticcurveto{\pgfqpoint{5.939196in}{0.340919in}}{\pgfqpoint{5.943843in}{0.336614in}}%
\pgfpathquadraticcurveto{\pgfqpoint{5.948490in}{0.332309in}}{\pgfqpoint{5.948490in}{0.324600in}}%
\pgfpathquadraticcurveto{\pgfqpoint{5.948490in}{0.315810in}}{\pgfqpoint{5.941537in}{0.310676in}}%
\pgfpathquadraticcurveto{\pgfqpoint{5.934585in}{0.305561in}}{\pgfqpoint{5.922426in}{0.305561in}}%
\pgfpathquadraticcurveto{\pgfqpoint{5.917365in}{0.305561in}}{\pgfqpoint{5.911871in}{0.306552in}}%
\pgfpathquadraticcurveto{\pgfqpoint{5.906378in}{0.307542in}}{\pgfqpoint{5.900308in}{0.309506in}}%
\pgfpathlineto{\pgfqpoint{5.900308in}{0.320205in}}%
\pgfpathquadraticcurveto{\pgfqpoint{5.906053in}{0.317215in}}{\pgfqpoint{5.911619in}{0.315720in}}%
\pgfpathquadraticcurveto{\pgfqpoint{5.917185in}{0.314243in}}{\pgfqpoint{5.922661in}{0.314243in}}%
\pgfpathquadraticcurveto{\pgfqpoint{5.929974in}{0.314243in}}{\pgfqpoint{5.933900in}{0.316746in}}%
\pgfpathquadraticcurveto{\pgfqpoint{5.937845in}{0.319250in}}{\pgfqpoint{5.937845in}{0.323807in}}%
\pgfpathquadraticcurveto{\pgfqpoint{5.937845in}{0.328022in}}{\pgfqpoint{5.934999in}{0.330274in}}%
\pgfpathquadraticcurveto{\pgfqpoint{5.932171in}{0.332525in}}{\pgfqpoint{5.922535in}{0.334614in}}%
\pgfpathlineto{\pgfqpoint{5.918932in}{0.335461in}}%
\pgfpathquadraticcurveto{\pgfqpoint{5.909314in}{0.337478in}}{\pgfqpoint{5.905027in}{0.341675in}}%
\pgfpathquadraticcurveto{\pgfqpoint{5.900758in}{0.345872in}}{\pgfqpoint{5.900758in}{0.353185in}}%
\pgfpathquadraticcurveto{\pgfqpoint{5.900758in}{0.362083in}}{\pgfqpoint{5.907062in}{0.366910in}}%
\pgfpathquadraticcurveto{\pgfqpoint{5.913366in}{0.371756in}}{\pgfqpoint{5.924966in}{0.371756in}}%
\pgfpathquadraticcurveto{\pgfqpoint{5.930694in}{0.371756in}}{\pgfqpoint{5.935755in}{0.370909in}}%
\pgfpathquadraticcurveto{\pgfqpoint{5.940835in}{0.370080in}}{\pgfqpoint{5.945104in}{0.368387in}}%
\pgfpathlineto{\pgfqpoint{5.945104in}{0.368387in}}%
\pgfpathclose%
\pgfpathmoveto{\pgfqpoint{6.017387in}{0.345260in}}%
\pgfpathlineto{\pgfqpoint{6.017387in}{0.307200in}}%
\pgfpathlineto{\pgfqpoint{6.007030in}{0.307200in}}%
\pgfpathlineto{\pgfqpoint{6.007030in}{0.344917in}}%
\pgfpathquadraticcurveto{\pgfqpoint{6.007030in}{0.353869in}}{\pgfqpoint{6.003535in}{0.358300in}}%
\pgfpathquadraticcurveto{\pgfqpoint{6.000041in}{0.362749in}}{\pgfqpoint{5.993070in}{0.362749in}}%
\pgfpathquadraticcurveto{\pgfqpoint{5.984676in}{0.362749in}}{\pgfqpoint{5.979831in}{0.357400in}}%
\pgfpathquadraticcurveto{\pgfqpoint{5.974986in}{0.352068in}}{\pgfqpoint{5.974986in}{0.342828in}}%
\pgfpathlineto{\pgfqpoint{5.974986in}{0.307200in}}%
\pgfpathlineto{\pgfqpoint{5.964575in}{0.307200in}}%
\pgfpathlineto{\pgfqpoint{5.964575in}{0.394793in}}%
\pgfpathlineto{\pgfqpoint{5.974986in}{0.394793in}}%
\pgfpathlineto{\pgfqpoint{5.974986in}{0.360444in}}%
\pgfpathquadraticcurveto{\pgfqpoint{5.978714in}{0.366136in}}{\pgfqpoint{5.983740in}{0.368946in}}%
\pgfpathquadraticcurveto{\pgfqpoint{5.988783in}{0.371756in}}{\pgfqpoint{5.995376in}{0.371756in}}%
\pgfpathquadraticcurveto{\pgfqpoint{6.006237in}{0.371756in}}{\pgfqpoint{6.011803in}{0.365037in}}%
\pgfpathquadraticcurveto{\pgfqpoint{6.017387in}{0.358318in}}{\pgfqpoint{6.017387in}{0.345260in}}%
\pgfpathlineto{\pgfqpoint{6.017387in}{0.345260in}}%
\pgfpathclose%
\pgfpathmoveto{\pgfqpoint{6.091957in}{0.341315in}}%
\pgfpathlineto{\pgfqpoint{6.091957in}{0.336254in}}%
\pgfpathlineto{\pgfqpoint{6.044333in}{0.336254in}}%
\pgfpathquadraticcurveto{\pgfqpoint{6.045017in}{0.325554in}}{\pgfqpoint{6.050781in}{0.319953in}}%
\pgfpathquadraticcurveto{\pgfqpoint{6.056545in}{0.314351in}}{\pgfqpoint{6.066848in}{0.314351in}}%
\pgfpathquadraticcurveto{\pgfqpoint{6.072810in}{0.314351in}}{\pgfqpoint{6.078412in}{0.315810in}}%
\pgfpathquadraticcurveto{\pgfqpoint{6.084013in}{0.317269in}}{\pgfqpoint{6.089543in}{0.320205in}}%
\pgfpathlineto{\pgfqpoint{6.089543in}{0.310406in}}%
\pgfpathquadraticcurveto{\pgfqpoint{6.083959in}{0.308047in}}{\pgfqpoint{6.078105in}{0.306804in}}%
\pgfpathquadraticcurveto{\pgfqpoint{6.072252in}{0.305561in}}{\pgfqpoint{6.066235in}{0.305561in}}%
\pgfpathquadraticcurveto{\pgfqpoint{6.051141in}{0.305561in}}{\pgfqpoint{6.042333in}{0.314333in}}%
\pgfpathquadraticcurveto{\pgfqpoint{6.033525in}{0.323123in}}{\pgfqpoint{6.033525in}{0.338109in}}%
\pgfpathquadraticcurveto{\pgfqpoint{6.033525in}{0.353581in}}{\pgfqpoint{6.041883in}{0.362659in}}%
\pgfpathquadraticcurveto{\pgfqpoint{6.050241in}{0.371756in}}{\pgfqpoint{6.064434in}{0.371756in}}%
\pgfpathquadraticcurveto{\pgfqpoint{6.077151in}{0.371756in}}{\pgfqpoint{6.084554in}{0.363560in}}%
\pgfpathquadraticcurveto{\pgfqpoint{6.091957in}{0.355383in}}{\pgfqpoint{6.091957in}{0.341315in}}%
\pgfpathlineto{\pgfqpoint{6.091957in}{0.341315in}}%
\pgfpathclose%
\pgfpathmoveto{\pgfqpoint{6.081600in}{0.344359in}}%
\pgfpathquadraticcurveto{\pgfqpoint{6.081492in}{0.352843in}}{\pgfqpoint{6.076845in}{0.357904in}}%
\pgfpathquadraticcurveto{\pgfqpoint{6.072197in}{0.362984in}}{\pgfqpoint{6.064542in}{0.362984in}}%
\pgfpathquadraticcurveto{\pgfqpoint{6.055878in}{0.362984in}}{\pgfqpoint{6.050673in}{0.358084in}}%
\pgfpathquadraticcurveto{\pgfqpoint{6.045467in}{0.353185in}}{\pgfqpoint{6.044675in}{0.344287in}}%
\pgfpathlineto{\pgfqpoint{6.081600in}{0.344359in}}%
\pgfpathlineto{\pgfqpoint{6.081600in}{0.344359in}}%
\pgfpathclose%
\pgfpathmoveto{\pgfqpoint{6.150442in}{0.360678in}}%
\pgfpathlineto{\pgfqpoint{6.150442in}{0.394793in}}%
\pgfpathlineto{\pgfqpoint{6.160799in}{0.394793in}}%
\pgfpathlineto{\pgfqpoint{6.160799in}{0.307200in}}%
\pgfpathlineto{\pgfqpoint{6.150442in}{0.307200in}}%
\pgfpathlineto{\pgfqpoint{6.150442in}{0.316656in}}%
\pgfpathquadraticcurveto{\pgfqpoint{6.147182in}{0.311037in}}{\pgfqpoint{6.142193in}{0.308299in}}%
\pgfpathquadraticcurveto{\pgfqpoint{6.137221in}{0.305561in}}{\pgfqpoint{6.130233in}{0.305561in}}%
\pgfpathquadraticcurveto{\pgfqpoint{6.118813in}{0.305561in}}{\pgfqpoint{6.111626in}{0.314675in}}%
\pgfpathquadraticcurveto{\pgfqpoint{6.104457in}{0.323807in}}{\pgfqpoint{6.104457in}{0.338667in}}%
\pgfpathquadraticcurveto{\pgfqpoint{6.104457in}{0.353527in}}{\pgfqpoint{6.111626in}{0.362641in}}%
\pgfpathquadraticcurveto{\pgfqpoint{6.118813in}{0.371756in}}{\pgfqpoint{6.130233in}{0.371756in}}%
\pgfpathquadraticcurveto{\pgfqpoint{6.137221in}{0.371756in}}{\pgfqpoint{6.142193in}{0.369018in}}%
\pgfpathquadraticcurveto{\pgfqpoint{6.147182in}{0.366298in}}{\pgfqpoint{6.150442in}{0.360678in}}%
\pgfpathlineto{\pgfqpoint{6.150442in}{0.360678in}}%
\pgfpathclose%
\pgfpathmoveto{\pgfqpoint{6.115156in}{0.338667in}}%
\pgfpathquadraticcurveto{\pgfqpoint{6.115156in}{0.327248in}}{\pgfqpoint{6.119858in}{0.320745in}}%
\pgfpathquadraticcurveto{\pgfqpoint{6.124559in}{0.314243in}}{\pgfqpoint{6.132772in}{0.314243in}}%
\pgfpathquadraticcurveto{\pgfqpoint{6.140986in}{0.314243in}}{\pgfqpoint{6.145705in}{0.320745in}}%
\pgfpathquadraticcurveto{\pgfqpoint{6.150442in}{0.327248in}}{\pgfqpoint{6.150442in}{0.338667in}}%
\pgfpathquadraticcurveto{\pgfqpoint{6.150442in}{0.350087in}}{\pgfqpoint{6.145705in}{0.356589in}}%
\pgfpathquadraticcurveto{\pgfqpoint{6.140986in}{0.363092in}}{\pgfqpoint{6.132772in}{0.363092in}}%
\pgfpathquadraticcurveto{\pgfqpoint{6.124559in}{0.363092in}}{\pgfqpoint{6.119858in}{0.356589in}}%
\pgfpathquadraticcurveto{\pgfqpoint{6.115156in}{0.350087in}}{\pgfqpoint{6.115156in}{0.338667in}}%
\pgfpathlineto{\pgfqpoint{6.115156in}{0.338667in}}%
\pgfpathclose%
\pgfpathmoveto{\pgfqpoint{6.269719in}{0.370243in}}%
\pgfpathlineto{\pgfqpoint{6.269719in}{0.307200in}}%
\pgfpathlineto{\pgfqpoint{6.259308in}{0.307200in}}%
\pgfpathlineto{\pgfqpoint{6.259308in}{0.362191in}}%
\pgfpathlineto{\pgfqpoint{6.230885in}{0.362191in}}%
\pgfpathlineto{\pgfqpoint{6.230885in}{0.307200in}}%
\pgfpathlineto{\pgfqpoint{6.220474in}{0.307200in}}%
\pgfpathlineto{\pgfqpoint{6.220474in}{0.362191in}}%
\pgfpathlineto{\pgfqpoint{6.210567in}{0.362191in}}%
\pgfpathlineto{\pgfqpoint{6.210567in}{0.370243in}}%
\pgfpathlineto{\pgfqpoint{6.220474in}{0.370243in}}%
\pgfpathlineto{\pgfqpoint{6.220474in}{0.374637in}}%
\pgfpathquadraticcurveto{\pgfqpoint{6.220474in}{0.384940in}}{\pgfqpoint{6.225337in}{0.389858in}}%
\pgfpathquadraticcurveto{\pgfqpoint{6.230218in}{0.394793in}}{\pgfqpoint{6.240287in}{0.394793in}}%
\pgfpathlineto{\pgfqpoint{6.250698in}{0.394793in}}%
\pgfpathlineto{\pgfqpoint{6.250698in}{0.386165in}}%
\pgfpathlineto{\pgfqpoint{6.240791in}{0.386165in}}%
\pgfpathquadraticcurveto{\pgfqpoint{6.235225in}{0.386165in}}{\pgfqpoint{6.233046in}{0.383914in}}%
\pgfpathquadraticcurveto{\pgfqpoint{6.230885in}{0.381662in}}{\pgfqpoint{6.230885in}{0.375808in}}%
\pgfpathlineto{\pgfqpoint{6.230885in}{0.370243in}}%
\pgfpathlineto{\pgfqpoint{6.269719in}{0.370243in}}%
\pgfpathlineto{\pgfqpoint{6.269719in}{0.370243in}}%
\pgfpathclose%
\pgfpathmoveto{\pgfqpoint{6.259308in}{0.394667in}}%
\pgfpathlineto{\pgfqpoint{6.269719in}{0.394667in}}%
\pgfpathlineto{\pgfqpoint{6.269719in}{0.381554in}}%
\pgfpathlineto{\pgfqpoint{6.259308in}{0.381554in}}%
\pgfpathlineto{\pgfqpoint{6.259308in}{0.394667in}}%
\pgfpathlineto{\pgfqpoint{6.259308in}{0.394667in}}%
\pgfpathclose%
\pgfpathmoveto{\pgfqpoint{6.301636in}{0.388147in}}%
\pgfpathlineto{\pgfqpoint{6.301636in}{0.370243in}}%
\pgfpathlineto{\pgfqpoint{6.322963in}{0.370243in}}%
\pgfpathlineto{\pgfqpoint{6.322963in}{0.362191in}}%
\pgfpathlineto{\pgfqpoint{6.301636in}{0.362191in}}%
\pgfpathlineto{\pgfqpoint{6.301636in}{0.327968in}}%
\pgfpathquadraticcurveto{\pgfqpoint{6.301636in}{0.320259in}}{\pgfqpoint{6.303744in}{0.318061in}}%
\pgfpathquadraticcurveto{\pgfqpoint{6.305851in}{0.315864in}}{\pgfqpoint{6.312336in}{0.315864in}}%
\pgfpathlineto{\pgfqpoint{6.322963in}{0.315864in}}%
\pgfpathlineto{\pgfqpoint{6.322963in}{0.307200in}}%
\pgfpathlineto{\pgfqpoint{6.312336in}{0.307200in}}%
\pgfpathquadraticcurveto{\pgfqpoint{6.300339in}{0.307200in}}{\pgfqpoint{6.295782in}{0.311667in}}%
\pgfpathquadraticcurveto{\pgfqpoint{6.291225in}{0.316152in}}{\pgfqpoint{6.291225in}{0.327968in}}%
\pgfpathlineto{\pgfqpoint{6.291225in}{0.362191in}}%
\pgfpathlineto{\pgfqpoint{6.283624in}{0.362191in}}%
\pgfpathlineto{\pgfqpoint{6.283624in}{0.370243in}}%
\pgfpathlineto{\pgfqpoint{6.291225in}{0.370243in}}%
\pgfpathlineto{\pgfqpoint{6.291225in}{0.388147in}}%
\pgfpathlineto{\pgfqpoint{6.301636in}{0.388147in}}%
\pgfpathlineto{\pgfqpoint{6.301636in}{0.388147in}}%
\pgfpathclose%
\pgfpathmoveto{\pgfqpoint{6.373217in}{0.370243in}}%
\pgfpathlineto{\pgfqpoint{6.383574in}{0.370243in}}%
\pgfpathlineto{\pgfqpoint{6.383574in}{0.307200in}}%
\pgfpathlineto{\pgfqpoint{6.373217in}{0.307200in}}%
\pgfpathlineto{\pgfqpoint{6.373217in}{0.370243in}}%
\pgfpathlineto{\pgfqpoint{6.373217in}{0.370243in}}%
\pgfpathclose%
\pgfpathmoveto{\pgfqpoint{6.373217in}{0.394793in}}%
\pgfpathlineto{\pgfqpoint{6.383574in}{0.394793in}}%
\pgfpathlineto{\pgfqpoint{6.383574in}{0.381662in}}%
\pgfpathlineto{\pgfqpoint{6.373217in}{0.381662in}}%
\pgfpathlineto{\pgfqpoint{6.373217in}{0.394793in}}%
\pgfpathlineto{\pgfqpoint{6.373217in}{0.394793in}}%
\pgfpathclose%
\pgfpathmoveto{\pgfqpoint{6.445427in}{0.368387in}}%
\pgfpathlineto{\pgfqpoint{6.445427in}{0.358589in}}%
\pgfpathquadraticcurveto{\pgfqpoint{6.441050in}{0.360840in}}{\pgfqpoint{6.436313in}{0.361957in}}%
\pgfpathquadraticcurveto{\pgfqpoint{6.431594in}{0.363092in}}{\pgfqpoint{6.426515in}{0.363092in}}%
\pgfpathquadraticcurveto{\pgfqpoint{6.418805in}{0.363092in}}{\pgfqpoint{6.414951in}{0.360732in}}%
\pgfpathquadraticcurveto{\pgfqpoint{6.411096in}{0.358373in}}{\pgfqpoint{6.411096in}{0.353635in}}%
\pgfpathquadraticcurveto{\pgfqpoint{6.411096in}{0.350033in}}{\pgfqpoint{6.413852in}{0.347980in}}%
\pgfpathquadraticcurveto{\pgfqpoint{6.416608in}{0.345926in}}{\pgfqpoint{6.424947in}{0.344071in}}%
\pgfpathlineto{\pgfqpoint{6.428496in}{0.343278in}}%
\pgfpathquadraticcurveto{\pgfqpoint{6.439519in}{0.340919in}}{\pgfqpoint{6.444166in}{0.336614in}}%
\pgfpathquadraticcurveto{\pgfqpoint{6.448814in}{0.332309in}}{\pgfqpoint{6.448814in}{0.324600in}}%
\pgfpathquadraticcurveto{\pgfqpoint{6.448814in}{0.315810in}}{\pgfqpoint{6.441861in}{0.310676in}}%
\pgfpathquadraticcurveto{\pgfqpoint{6.434908in}{0.305561in}}{\pgfqpoint{6.422750in}{0.305561in}}%
\pgfpathquadraticcurveto{\pgfqpoint{6.417689in}{0.305561in}}{\pgfqpoint{6.412195in}{0.306552in}}%
\pgfpathquadraticcurveto{\pgfqpoint{6.406701in}{0.307542in}}{\pgfqpoint{6.400631in}{0.309506in}}%
\pgfpathlineto{\pgfqpoint{6.400631in}{0.320205in}}%
\pgfpathquadraticcurveto{\pgfqpoint{6.406377in}{0.317215in}}{\pgfqpoint{6.411943in}{0.315720in}}%
\pgfpathquadraticcurveto{\pgfqpoint{6.417508in}{0.314243in}}{\pgfqpoint{6.422984in}{0.314243in}}%
\pgfpathquadraticcurveto{\pgfqpoint{6.430297in}{0.314243in}}{\pgfqpoint{6.434224in}{0.316746in}}%
\pgfpathquadraticcurveto{\pgfqpoint{6.438168in}{0.319250in}}{\pgfqpoint{6.438168in}{0.323807in}}%
\pgfpathquadraticcurveto{\pgfqpoint{6.438168in}{0.328022in}}{\pgfqpoint{6.435322in}{0.330274in}}%
\pgfpathquadraticcurveto{\pgfqpoint{6.432495in}{0.332525in}}{\pgfqpoint{6.422858in}{0.334614in}}%
\pgfpathlineto{\pgfqpoint{6.419256in}{0.335461in}}%
\pgfpathquadraticcurveto{\pgfqpoint{6.409637in}{0.337478in}}{\pgfqpoint{6.405350in}{0.341675in}}%
\pgfpathquadraticcurveto{\pgfqpoint{6.401081in}{0.345872in}}{\pgfqpoint{6.401081in}{0.353185in}}%
\pgfpathquadraticcurveto{\pgfqpoint{6.401081in}{0.362083in}}{\pgfqpoint{6.407386in}{0.366910in}}%
\pgfpathquadraticcurveto{\pgfqpoint{6.413690in}{0.371756in}}{\pgfqpoint{6.425290in}{0.371756in}}%
\pgfpathquadraticcurveto{\pgfqpoint{6.431018in}{0.371756in}}{\pgfqpoint{6.436079in}{0.370909in}}%
\pgfpathquadraticcurveto{\pgfqpoint{6.441158in}{0.370080in}}{\pgfqpoint{6.445427in}{0.368387in}}%
\pgfpathlineto{\pgfqpoint{6.445427in}{0.368387in}}%
\pgfpathclose%
\pgfpathmoveto{\pgfqpoint{6.538460in}{0.360570in}}%
\pgfpathquadraticcurveto{\pgfqpoint{6.536713in}{0.361579in}}{\pgfqpoint{6.534660in}{0.362047in}}%
\pgfpathquadraticcurveto{\pgfqpoint{6.532606in}{0.362533in}}{\pgfqpoint{6.530138in}{0.362533in}}%
\pgfpathquadraticcurveto{\pgfqpoint{6.521349in}{0.362533in}}{\pgfqpoint{6.516647in}{0.356823in}}%
\pgfpathquadraticcurveto{\pgfqpoint{6.511946in}{0.351114in}}{\pgfqpoint{6.511946in}{0.340414in}}%
\pgfpathlineto{\pgfqpoint{6.511946in}{0.307200in}}%
\pgfpathlineto{\pgfqpoint{6.501535in}{0.307200in}}%
\pgfpathlineto{\pgfqpoint{6.501535in}{0.370243in}}%
\pgfpathlineto{\pgfqpoint{6.511946in}{0.370243in}}%
\pgfpathlineto{\pgfqpoint{6.511946in}{0.360444in}}%
\pgfpathquadraticcurveto{\pgfqpoint{6.515224in}{0.366190in}}{\pgfqpoint{6.520448in}{0.368964in}}%
\pgfpathquadraticcurveto{\pgfqpoint{6.525689in}{0.371756in}}{\pgfqpoint{6.533183in}{0.371756in}}%
\pgfpathquadraticcurveto{\pgfqpoint{6.534245in}{0.371756in}}{\pgfqpoint{6.535542in}{0.371611in}}%
\pgfpathquadraticcurveto{\pgfqpoint{6.536839in}{0.371485in}}{\pgfqpoint{6.538406in}{0.371197in}}%
\pgfpathlineto{\pgfqpoint{6.538460in}{0.360570in}}%
\pgfpathlineto{\pgfqpoint{6.538460in}{0.360570in}}%
\pgfpathclose%
\pgfpathmoveto{\pgfqpoint{6.600710in}{0.341315in}}%
\pgfpathlineto{\pgfqpoint{6.600710in}{0.336254in}}%
\pgfpathlineto{\pgfqpoint{6.553086in}{0.336254in}}%
\pgfpathquadraticcurveto{\pgfqpoint{6.553770in}{0.325554in}}{\pgfqpoint{6.559534in}{0.319953in}}%
\pgfpathquadraticcurveto{\pgfqpoint{6.565298in}{0.314351in}}{\pgfqpoint{6.575601in}{0.314351in}}%
\pgfpathquadraticcurveto{\pgfqpoint{6.581563in}{0.314351in}}{\pgfqpoint{6.587165in}{0.315810in}}%
\pgfpathquadraticcurveto{\pgfqpoint{6.592767in}{0.317269in}}{\pgfqpoint{6.598296in}{0.320205in}}%
\pgfpathlineto{\pgfqpoint{6.598296in}{0.310406in}}%
\pgfpathquadraticcurveto{\pgfqpoint{6.592713in}{0.308047in}}{\pgfqpoint{6.586859in}{0.306804in}}%
\pgfpathquadraticcurveto{\pgfqpoint{6.581005in}{0.305561in}}{\pgfqpoint{6.574989in}{0.305561in}}%
\pgfpathquadraticcurveto{\pgfqpoint{6.559895in}{0.305561in}}{\pgfqpoint{6.551087in}{0.314333in}}%
\pgfpathquadraticcurveto{\pgfqpoint{6.542279in}{0.323123in}}{\pgfqpoint{6.542279in}{0.338109in}}%
\pgfpathquadraticcurveto{\pgfqpoint{6.542279in}{0.353581in}}{\pgfqpoint{6.550636in}{0.362659in}}%
\pgfpathquadraticcurveto{\pgfqpoint{6.558994in}{0.371756in}}{\pgfqpoint{6.573188in}{0.371756in}}%
\pgfpathquadraticcurveto{\pgfqpoint{6.585904in}{0.371756in}}{\pgfqpoint{6.593307in}{0.363560in}}%
\pgfpathquadraticcurveto{\pgfqpoint{6.600710in}{0.355383in}}{\pgfqpoint{6.600710in}{0.341315in}}%
\pgfpathlineto{\pgfqpoint{6.600710in}{0.341315in}}%
\pgfpathclose%
\pgfpathmoveto{\pgfqpoint{6.590353in}{0.344359in}}%
\pgfpathquadraticcurveto{\pgfqpoint{6.590245in}{0.352843in}}{\pgfqpoint{6.585598in}{0.357904in}}%
\pgfpathquadraticcurveto{\pgfqpoint{6.580951in}{0.362984in}}{\pgfqpoint{6.573296in}{0.362984in}}%
\pgfpathquadraticcurveto{\pgfqpoint{6.564632in}{0.362984in}}{\pgfqpoint{6.559426in}{0.358084in}}%
\pgfpathquadraticcurveto{\pgfqpoint{6.554221in}{0.353185in}}{\pgfqpoint{6.553428in}{0.344287in}}%
\pgfpathlineto{\pgfqpoint{6.590353in}{0.344359in}}%
\pgfpathlineto{\pgfqpoint{6.590353in}{0.344359in}}%
\pgfpathclose%
\pgfpathmoveto{\pgfqpoint{6.627962in}{0.388147in}}%
\pgfpathlineto{\pgfqpoint{6.627962in}{0.370243in}}%
\pgfpathlineto{\pgfqpoint{6.649289in}{0.370243in}}%
\pgfpathlineto{\pgfqpoint{6.649289in}{0.362191in}}%
\pgfpathlineto{\pgfqpoint{6.627962in}{0.362191in}}%
\pgfpathlineto{\pgfqpoint{6.627962in}{0.327968in}}%
\pgfpathquadraticcurveto{\pgfqpoint{6.627962in}{0.320259in}}{\pgfqpoint{6.630070in}{0.318061in}}%
\pgfpathquadraticcurveto{\pgfqpoint{6.632177in}{0.315864in}}{\pgfqpoint{6.638662in}{0.315864in}}%
\pgfpathlineto{\pgfqpoint{6.649289in}{0.315864in}}%
\pgfpathlineto{\pgfqpoint{6.649289in}{0.307200in}}%
\pgfpathlineto{\pgfqpoint{6.638662in}{0.307200in}}%
\pgfpathquadraticcurveto{\pgfqpoint{6.626666in}{0.307200in}}{\pgfqpoint{6.622109in}{0.311667in}}%
\pgfpathquadraticcurveto{\pgfqpoint{6.617551in}{0.316152in}}{\pgfqpoint{6.617551in}{0.327968in}}%
\pgfpathlineto{\pgfqpoint{6.617551in}{0.362191in}}%
\pgfpathlineto{\pgfqpoint{6.609950in}{0.362191in}}%
\pgfpathlineto{\pgfqpoint{6.609950in}{0.370243in}}%
\pgfpathlineto{\pgfqpoint{6.617551in}{0.370243in}}%
\pgfpathlineto{\pgfqpoint{6.617551in}{0.388147in}}%
\pgfpathlineto{\pgfqpoint{6.627962in}{0.388147in}}%
\pgfpathlineto{\pgfqpoint{6.627962in}{0.388147in}}%
\pgfpathclose%
\pgfpathmoveto{\pgfqpoint{6.691563in}{0.338883in}}%
\pgfpathquadraticcurveto{\pgfqpoint{6.679009in}{0.338883in}}{\pgfqpoint{6.674164in}{0.336019in}}%
\pgfpathquadraticcurveto{\pgfqpoint{6.669318in}{0.333156in}}{\pgfqpoint{6.669318in}{0.326221in}}%
\pgfpathquadraticcurveto{\pgfqpoint{6.669318in}{0.320709in}}{\pgfqpoint{6.672957in}{0.317467in}}%
\pgfpathquadraticcurveto{\pgfqpoint{6.676595in}{0.314243in}}{\pgfqpoint{6.682827in}{0.314243in}}%
\pgfpathquadraticcurveto{\pgfqpoint{6.691455in}{0.314243in}}{\pgfqpoint{6.696661in}{0.320349in}}%
\pgfpathquadraticcurveto{\pgfqpoint{6.701866in}{0.326455in}}{\pgfqpoint{6.701866in}{0.336578in}}%
\pgfpathlineto{\pgfqpoint{6.701866in}{0.338883in}}%
\pgfpathlineto{\pgfqpoint{6.691563in}{0.338883in}}%
\pgfpathlineto{\pgfqpoint{6.691563in}{0.338883in}}%
\pgfpathclose%
\pgfpathmoveto{\pgfqpoint{6.712223in}{0.343170in}}%
\pgfpathlineto{\pgfqpoint{6.712223in}{0.307200in}}%
\pgfpathlineto{\pgfqpoint{6.701866in}{0.307200in}}%
\pgfpathlineto{\pgfqpoint{6.701866in}{0.316764in}}%
\pgfpathquadraticcurveto{\pgfqpoint{6.698318in}{0.311037in}}{\pgfqpoint{6.693022in}{0.308299in}}%
\pgfpathquadraticcurveto{\pgfqpoint{6.687727in}{0.305561in}}{\pgfqpoint{6.680072in}{0.305561in}}%
\pgfpathquadraticcurveto{\pgfqpoint{6.670399in}{0.305561in}}{\pgfqpoint{6.664671in}{0.311001in}}%
\pgfpathquadraticcurveto{\pgfqpoint{6.658961in}{0.316440in}}{\pgfqpoint{6.658961in}{0.325554in}}%
\pgfpathquadraticcurveto{\pgfqpoint{6.658961in}{0.336182in}}{\pgfqpoint{6.666076in}{0.341585in}}%
\pgfpathquadraticcurveto{\pgfqpoint{6.673209in}{0.346989in}}{\pgfqpoint{6.687331in}{0.346989in}}%
\pgfpathlineto{\pgfqpoint{6.701866in}{0.346989in}}%
\pgfpathlineto{\pgfqpoint{6.701866in}{0.348016in}}%
\pgfpathquadraticcurveto{\pgfqpoint{6.701866in}{0.355166in}}{\pgfqpoint{6.697165in}{0.359075in}}%
\pgfpathquadraticcurveto{\pgfqpoint{6.692464in}{0.362984in}}{\pgfqpoint{6.683962in}{0.362984in}}%
\pgfpathquadraticcurveto{\pgfqpoint{6.678559in}{0.362984in}}{\pgfqpoint{6.673425in}{0.361687in}}%
\pgfpathquadraticcurveto{\pgfqpoint{6.668310in}{0.360390in}}{\pgfqpoint{6.663590in}{0.357796in}}%
\pgfpathlineto{\pgfqpoint{6.663590in}{0.367379in}}%
\pgfpathquadraticcurveto{\pgfqpoint{6.669264in}{0.369576in}}{\pgfqpoint{6.674614in}{0.370657in}}%
\pgfpathquadraticcurveto{\pgfqpoint{6.679964in}{0.371756in}}{\pgfqpoint{6.685025in}{0.371756in}}%
\pgfpathquadraticcurveto{\pgfqpoint{6.698714in}{0.371756in}}{\pgfqpoint{6.705469in}{0.364659in}}%
\pgfpathquadraticcurveto{\pgfqpoint{6.712223in}{0.357580in}}{\pgfqpoint{6.712223in}{0.343170in}}%
\pgfpathlineto{\pgfqpoint{6.712223in}{0.343170in}}%
\pgfpathclose%
\pgfpathmoveto{\pgfqpoint{6.733550in}{0.370243in}}%
\pgfpathlineto{\pgfqpoint{6.743907in}{0.370243in}}%
\pgfpathlineto{\pgfqpoint{6.743907in}{0.307200in}}%
\pgfpathlineto{\pgfqpoint{6.733550in}{0.307200in}}%
\pgfpathlineto{\pgfqpoint{6.733550in}{0.370243in}}%
\pgfpathlineto{\pgfqpoint{6.733550in}{0.370243in}}%
\pgfpathclose%
\pgfpathmoveto{\pgfqpoint{6.733550in}{0.394793in}}%
\pgfpathlineto{\pgfqpoint{6.743907in}{0.394793in}}%
\pgfpathlineto{\pgfqpoint{6.743907in}{0.381662in}}%
\pgfpathlineto{\pgfqpoint{6.733550in}{0.381662in}}%
\pgfpathlineto{\pgfqpoint{6.733550in}{0.394793in}}%
\pgfpathlineto{\pgfqpoint{6.733550in}{0.394793in}}%
\pgfpathclose%
\pgfpathmoveto{\pgfqpoint{6.817991in}{0.345260in}}%
\pgfpathlineto{\pgfqpoint{6.817991in}{0.307200in}}%
\pgfpathlineto{\pgfqpoint{6.807634in}{0.307200in}}%
\pgfpathlineto{\pgfqpoint{6.807634in}{0.344917in}}%
\pgfpathquadraticcurveto{\pgfqpoint{6.807634in}{0.353869in}}{\pgfqpoint{6.804139in}{0.358300in}}%
\pgfpathquadraticcurveto{\pgfqpoint{6.800645in}{0.362749in}}{\pgfqpoint{6.793674in}{0.362749in}}%
\pgfpathquadraticcurveto{\pgfqpoint{6.785281in}{0.362749in}}{\pgfqpoint{6.780435in}{0.357400in}}%
\pgfpathquadraticcurveto{\pgfqpoint{6.775590in}{0.352068in}}{\pgfqpoint{6.775590in}{0.342828in}}%
\pgfpathlineto{\pgfqpoint{6.775590in}{0.307200in}}%
\pgfpathlineto{\pgfqpoint{6.765179in}{0.307200in}}%
\pgfpathlineto{\pgfqpoint{6.765179in}{0.370243in}}%
\pgfpathlineto{\pgfqpoint{6.775590in}{0.370243in}}%
\pgfpathlineto{\pgfqpoint{6.775590in}{0.360444in}}%
\pgfpathquadraticcurveto{\pgfqpoint{6.779319in}{0.366136in}}{\pgfqpoint{6.784344in}{0.368946in}}%
\pgfpathquadraticcurveto{\pgfqpoint{6.789387in}{0.371756in}}{\pgfqpoint{6.795980in}{0.371756in}}%
\pgfpathquadraticcurveto{\pgfqpoint{6.806841in}{0.371756in}}{\pgfqpoint{6.812407in}{0.365037in}}%
\pgfpathquadraticcurveto{\pgfqpoint{6.817991in}{0.358318in}}{\pgfqpoint{6.817991in}{0.345260in}}%
\pgfpathlineto{\pgfqpoint{6.817991in}{0.345260in}}%
\pgfpathclose%
\pgfpathmoveto{\pgfqpoint{6.892561in}{0.341315in}}%
\pgfpathlineto{\pgfqpoint{6.892561in}{0.336254in}}%
\pgfpathlineto{\pgfqpoint{6.844937in}{0.336254in}}%
\pgfpathquadraticcurveto{\pgfqpoint{6.845621in}{0.325554in}}{\pgfqpoint{6.851385in}{0.319953in}}%
\pgfpathquadraticcurveto{\pgfqpoint{6.857149in}{0.314351in}}{\pgfqpoint{6.867452in}{0.314351in}}%
\pgfpathquadraticcurveto{\pgfqpoint{6.873414in}{0.314351in}}{\pgfqpoint{6.879016in}{0.315810in}}%
\pgfpathquadraticcurveto{\pgfqpoint{6.884618in}{0.317269in}}{\pgfqpoint{6.890147in}{0.320205in}}%
\pgfpathlineto{\pgfqpoint{6.890147in}{0.310406in}}%
\pgfpathquadraticcurveto{\pgfqpoint{6.884564in}{0.308047in}}{\pgfqpoint{6.878710in}{0.306804in}}%
\pgfpathquadraticcurveto{\pgfqpoint{6.872856in}{0.305561in}}{\pgfqpoint{6.866840in}{0.305561in}}%
\pgfpathquadraticcurveto{\pgfqpoint{6.851745in}{0.305561in}}{\pgfqpoint{6.842938in}{0.314333in}}%
\pgfpathquadraticcurveto{\pgfqpoint{6.834130in}{0.323123in}}{\pgfqpoint{6.834130in}{0.338109in}}%
\pgfpathquadraticcurveto{\pgfqpoint{6.834130in}{0.353581in}}{\pgfqpoint{6.842487in}{0.362659in}}%
\pgfpathquadraticcurveto{\pgfqpoint{6.850845in}{0.371756in}}{\pgfqpoint{6.865038in}{0.371756in}}%
\pgfpathquadraticcurveto{\pgfqpoint{6.877755in}{0.371756in}}{\pgfqpoint{6.885158in}{0.363560in}}%
\pgfpathquadraticcurveto{\pgfqpoint{6.892561in}{0.355383in}}{\pgfqpoint{6.892561in}{0.341315in}}%
\pgfpathlineto{\pgfqpoint{6.892561in}{0.341315in}}%
\pgfpathclose%
\pgfpathmoveto{\pgfqpoint{6.882204in}{0.344359in}}%
\pgfpathquadraticcurveto{\pgfqpoint{6.882096in}{0.352843in}}{\pgfqpoint{6.877449in}{0.357904in}}%
\pgfpathquadraticcurveto{\pgfqpoint{6.872802in}{0.362984in}}{\pgfqpoint{6.865146in}{0.362984in}}%
\pgfpathquadraticcurveto{\pgfqpoint{6.856483in}{0.362984in}}{\pgfqpoint{6.851277in}{0.358084in}}%
\pgfpathquadraticcurveto{\pgfqpoint{6.846072in}{0.353185in}}{\pgfqpoint{6.845279in}{0.344287in}}%
\pgfpathlineto{\pgfqpoint{6.882204in}{0.344359in}}%
\pgfpathlineto{\pgfqpoint{6.882204in}{0.344359in}}%
\pgfpathclose%
\pgfpathmoveto{\pgfqpoint{6.951046in}{0.360678in}}%
\pgfpathlineto{\pgfqpoint{6.951046in}{0.394793in}}%
\pgfpathlineto{\pgfqpoint{6.961403in}{0.394793in}}%
\pgfpathlineto{\pgfqpoint{6.961403in}{0.307200in}}%
\pgfpathlineto{\pgfqpoint{6.951046in}{0.307200in}}%
\pgfpathlineto{\pgfqpoint{6.951046in}{0.316656in}}%
\pgfpathquadraticcurveto{\pgfqpoint{6.947786in}{0.311037in}}{\pgfqpoint{6.942797in}{0.308299in}}%
\pgfpathquadraticcurveto{\pgfqpoint{6.937826in}{0.305561in}}{\pgfqpoint{6.930837in}{0.305561in}}%
\pgfpathquadraticcurveto{\pgfqpoint{6.919417in}{0.305561in}}{\pgfqpoint{6.912230in}{0.314675in}}%
\pgfpathquadraticcurveto{\pgfqpoint{6.905061in}{0.323807in}}{\pgfqpoint{6.905061in}{0.338667in}}%
\pgfpathquadraticcurveto{\pgfqpoint{6.905061in}{0.353527in}}{\pgfqpoint{6.912230in}{0.362641in}}%
\pgfpathquadraticcurveto{\pgfqpoint{6.919417in}{0.371756in}}{\pgfqpoint{6.930837in}{0.371756in}}%
\pgfpathquadraticcurveto{\pgfqpoint{6.937826in}{0.371756in}}{\pgfqpoint{6.942797in}{0.369018in}}%
\pgfpathquadraticcurveto{\pgfqpoint{6.947786in}{0.366298in}}{\pgfqpoint{6.951046in}{0.360678in}}%
\pgfpathlineto{\pgfqpoint{6.951046in}{0.360678in}}%
\pgfpathclose%
\pgfpathmoveto{\pgfqpoint{6.915761in}{0.338667in}}%
\pgfpathquadraticcurveto{\pgfqpoint{6.915761in}{0.327248in}}{\pgfqpoint{6.920462in}{0.320745in}}%
\pgfpathquadraticcurveto{\pgfqpoint{6.925163in}{0.314243in}}{\pgfqpoint{6.933377in}{0.314243in}}%
\pgfpathquadraticcurveto{\pgfqpoint{6.941590in}{0.314243in}}{\pgfqpoint{6.946309in}{0.320745in}}%
\pgfpathquadraticcurveto{\pgfqpoint{6.951046in}{0.327248in}}{\pgfqpoint{6.951046in}{0.338667in}}%
\pgfpathquadraticcurveto{\pgfqpoint{6.951046in}{0.350087in}}{\pgfqpoint{6.946309in}{0.356589in}}%
\pgfpathquadraticcurveto{\pgfqpoint{6.941590in}{0.363092in}}{\pgfqpoint{6.933377in}{0.363092in}}%
\pgfpathquadraticcurveto{\pgfqpoint{6.925163in}{0.363092in}}{\pgfqpoint{6.920462in}{0.356589in}}%
\pgfpathquadraticcurveto{\pgfqpoint{6.915761in}{0.350087in}}{\pgfqpoint{6.915761in}{0.338667in}}%
\pgfpathlineto{\pgfqpoint{6.915761in}{0.338667in}}%
\pgfpathclose%
\pgfpathmoveto{\pgfqpoint{7.059570in}{0.368387in}}%
\pgfpathlineto{\pgfqpoint{7.059570in}{0.358589in}}%
\pgfpathquadraticcurveto{\pgfqpoint{7.055193in}{0.360840in}}{\pgfqpoint{7.050456in}{0.361957in}}%
\pgfpathquadraticcurveto{\pgfqpoint{7.045736in}{0.363092in}}{\pgfqpoint{7.040657in}{0.363092in}}%
\pgfpathquadraticcurveto{\pgfqpoint{7.032948in}{0.363092in}}{\pgfqpoint{7.029093in}{0.360732in}}%
\pgfpathquadraticcurveto{\pgfqpoint{7.025238in}{0.358373in}}{\pgfqpoint{7.025238in}{0.353635in}}%
\pgfpathquadraticcurveto{\pgfqpoint{7.025238in}{0.350033in}}{\pgfqpoint{7.027994in}{0.347980in}}%
\pgfpathquadraticcurveto{\pgfqpoint{7.030750in}{0.345926in}}{\pgfqpoint{7.039090in}{0.344071in}}%
\pgfpathlineto{\pgfqpoint{7.042638in}{0.343278in}}%
\pgfpathquadraticcurveto{\pgfqpoint{7.053662in}{0.340919in}}{\pgfqpoint{7.058309in}{0.336614in}}%
\pgfpathquadraticcurveto{\pgfqpoint{7.062956in}{0.332309in}}{\pgfqpoint{7.062956in}{0.324600in}}%
\pgfpathquadraticcurveto{\pgfqpoint{7.062956in}{0.315810in}}{\pgfqpoint{7.056003in}{0.310676in}}%
\pgfpathquadraticcurveto{\pgfqpoint{7.049051in}{0.305561in}}{\pgfqpoint{7.036892in}{0.305561in}}%
\pgfpathquadraticcurveto{\pgfqpoint{7.031831in}{0.305561in}}{\pgfqpoint{7.026337in}{0.306552in}}%
\pgfpathquadraticcurveto{\pgfqpoint{7.020844in}{0.307542in}}{\pgfqpoint{7.014773in}{0.309506in}}%
\pgfpathlineto{\pgfqpoint{7.014773in}{0.320205in}}%
\pgfpathquadraticcurveto{\pgfqpoint{7.020519in}{0.317215in}}{\pgfqpoint{7.026085in}{0.315720in}}%
\pgfpathquadraticcurveto{\pgfqpoint{7.031651in}{0.314243in}}{\pgfqpoint{7.037127in}{0.314243in}}%
\pgfpathquadraticcurveto{\pgfqpoint{7.044439in}{0.314243in}}{\pgfqpoint{7.048366in}{0.316746in}}%
\pgfpathquadraticcurveto{\pgfqpoint{7.052311in}{0.319250in}}{\pgfqpoint{7.052311in}{0.323807in}}%
\pgfpathquadraticcurveto{\pgfqpoint{7.052311in}{0.328022in}}{\pgfqpoint{7.049465in}{0.330274in}}%
\pgfpathquadraticcurveto{\pgfqpoint{7.046637in}{0.332525in}}{\pgfqpoint{7.037000in}{0.334614in}}%
\pgfpathlineto{\pgfqpoint{7.033398in}{0.335461in}}%
\pgfpathquadraticcurveto{\pgfqpoint{7.023780in}{0.337478in}}{\pgfqpoint{7.019493in}{0.341675in}}%
\pgfpathquadraticcurveto{\pgfqpoint{7.015224in}{0.345872in}}{\pgfqpoint{7.015224in}{0.353185in}}%
\pgfpathquadraticcurveto{\pgfqpoint{7.015224in}{0.362083in}}{\pgfqpoint{7.021528in}{0.366910in}}%
\pgfpathquadraticcurveto{\pgfqpoint{7.027832in}{0.371756in}}{\pgfqpoint{7.039432in}{0.371756in}}%
\pgfpathquadraticcurveto{\pgfqpoint{7.045160in}{0.371756in}}{\pgfqpoint{7.050221in}{0.370909in}}%
\pgfpathquadraticcurveto{\pgfqpoint{7.055301in}{0.370080in}}{\pgfqpoint{7.059570in}{0.368387in}}%
\pgfpathlineto{\pgfqpoint{7.059570in}{0.368387in}}%
\pgfpathclose%
\pgfpathmoveto{\pgfqpoint{7.133365in}{0.341315in}}%
\pgfpathlineto{\pgfqpoint{7.133365in}{0.336254in}}%
\pgfpathlineto{\pgfqpoint{7.085741in}{0.336254in}}%
\pgfpathquadraticcurveto{\pgfqpoint{7.086426in}{0.325554in}}{\pgfqpoint{7.092190in}{0.319953in}}%
\pgfpathquadraticcurveto{\pgfqpoint{7.097954in}{0.314351in}}{\pgfqpoint{7.108257in}{0.314351in}}%
\pgfpathquadraticcurveto{\pgfqpoint{7.114219in}{0.314351in}}{\pgfqpoint{7.119820in}{0.315810in}}%
\pgfpathquadraticcurveto{\pgfqpoint{7.125422in}{0.317269in}}{\pgfqpoint{7.130952in}{0.320205in}}%
\pgfpathlineto{\pgfqpoint{7.130952in}{0.310406in}}%
\pgfpathquadraticcurveto{\pgfqpoint{7.125368in}{0.308047in}}{\pgfqpoint{7.119514in}{0.306804in}}%
\pgfpathquadraticcurveto{\pgfqpoint{7.113660in}{0.305561in}}{\pgfqpoint{7.107644in}{0.305561in}}%
\pgfpathquadraticcurveto{\pgfqpoint{7.092550in}{0.305561in}}{\pgfqpoint{7.083742in}{0.314333in}}%
\pgfpathquadraticcurveto{\pgfqpoint{7.074934in}{0.323123in}}{\pgfqpoint{7.074934in}{0.338109in}}%
\pgfpathquadraticcurveto{\pgfqpoint{7.074934in}{0.353581in}}{\pgfqpoint{7.083292in}{0.362659in}}%
\pgfpathquadraticcurveto{\pgfqpoint{7.091649in}{0.371756in}}{\pgfqpoint{7.105843in}{0.371756in}}%
\pgfpathquadraticcurveto{\pgfqpoint{7.118559in}{0.371756in}}{\pgfqpoint{7.125962in}{0.363560in}}%
\pgfpathquadraticcurveto{\pgfqpoint{7.133365in}{0.355383in}}{\pgfqpoint{7.133365in}{0.341315in}}%
\pgfpathlineto{\pgfqpoint{7.133365in}{0.341315in}}%
\pgfpathclose%
\pgfpathmoveto{\pgfqpoint{7.123008in}{0.344359in}}%
\pgfpathquadraticcurveto{\pgfqpoint{7.122900in}{0.352843in}}{\pgfqpoint{7.118253in}{0.357904in}}%
\pgfpathquadraticcurveto{\pgfqpoint{7.113606in}{0.362984in}}{\pgfqpoint{7.105951in}{0.362984in}}%
\pgfpathquadraticcurveto{\pgfqpoint{7.097287in}{0.362984in}}{\pgfqpoint{7.092082in}{0.358084in}}%
\pgfpathquadraticcurveto{\pgfqpoint{7.086876in}{0.353185in}}{\pgfqpoint{7.086084in}{0.344287in}}%
\pgfpathlineto{\pgfqpoint{7.123008in}{0.344359in}}%
\pgfpathlineto{\pgfqpoint{7.123008in}{0.344359in}}%
\pgfpathclose%
\pgfpathmoveto{\pgfqpoint{7.160384in}{0.316656in}}%
\pgfpathlineto{\pgfqpoint{7.160384in}{0.283226in}}%
\pgfpathlineto{\pgfqpoint{7.149973in}{0.283226in}}%
\pgfpathlineto{\pgfqpoint{7.149973in}{0.370243in}}%
\pgfpathlineto{\pgfqpoint{7.160384in}{0.370243in}}%
\pgfpathlineto{\pgfqpoint{7.160384in}{0.360678in}}%
\pgfpathquadraticcurveto{\pgfqpoint{7.163662in}{0.366298in}}{\pgfqpoint{7.168633in}{0.369018in}}%
\pgfpathquadraticcurveto{\pgfqpoint{7.173623in}{0.371756in}}{\pgfqpoint{7.180539in}{0.371756in}}%
\pgfpathquadraticcurveto{\pgfqpoint{7.192031in}{0.371756in}}{\pgfqpoint{7.199200in}{0.362641in}}%
\pgfpathquadraticcurveto{\pgfqpoint{7.206387in}{0.353527in}}{\pgfqpoint{7.206387in}{0.338667in}}%
\pgfpathquadraticcurveto{\pgfqpoint{7.206387in}{0.323807in}}{\pgfqpoint{7.199200in}{0.314675in}}%
\pgfpathquadraticcurveto{\pgfqpoint{7.192031in}{0.305561in}}{\pgfqpoint{7.180539in}{0.305561in}}%
\pgfpathquadraticcurveto{\pgfqpoint{7.173623in}{0.305561in}}{\pgfqpoint{7.168633in}{0.308299in}}%
\pgfpathquadraticcurveto{\pgfqpoint{7.163662in}{0.311037in}}{\pgfqpoint{7.160384in}{0.316656in}}%
\pgfpathlineto{\pgfqpoint{7.160384in}{0.316656in}}%
\pgfpathclose%
\pgfpathmoveto{\pgfqpoint{7.195633in}{0.338667in}}%
\pgfpathquadraticcurveto{\pgfqpoint{7.195633in}{0.350087in}}{\pgfqpoint{7.190932in}{0.356589in}}%
\pgfpathquadraticcurveto{\pgfqpoint{7.186231in}{0.363092in}}{\pgfqpoint{7.178018in}{0.363092in}}%
\pgfpathquadraticcurveto{\pgfqpoint{7.169786in}{0.363092in}}{\pgfqpoint{7.165085in}{0.356589in}}%
\pgfpathquadraticcurveto{\pgfqpoint{7.160384in}{0.350087in}}{\pgfqpoint{7.160384in}{0.338667in}}%
\pgfpathquadraticcurveto{\pgfqpoint{7.160384in}{0.327248in}}{\pgfqpoint{7.165085in}{0.320745in}}%
\pgfpathquadraticcurveto{\pgfqpoint{7.169786in}{0.314243in}}{\pgfqpoint{7.178018in}{0.314243in}}%
\pgfpathquadraticcurveto{\pgfqpoint{7.186231in}{0.314243in}}{\pgfqpoint{7.190932in}{0.320745in}}%
\pgfpathquadraticcurveto{\pgfqpoint{7.195633in}{0.327248in}}{\pgfqpoint{7.195633in}{0.338667in}}%
\pgfpathlineto{\pgfqpoint{7.195633in}{0.338667in}}%
\pgfpathclose%
\pgfpathmoveto{\pgfqpoint{7.252210in}{0.338883in}}%
\pgfpathquadraticcurveto{\pgfqpoint{7.239655in}{0.338883in}}{\pgfqpoint{7.234810in}{0.336019in}}%
\pgfpathquadraticcurveto{\pgfqpoint{7.229965in}{0.333156in}}{\pgfqpoint{7.229965in}{0.326221in}}%
\pgfpathquadraticcurveto{\pgfqpoint{7.229965in}{0.320709in}}{\pgfqpoint{7.233603in}{0.317467in}}%
\pgfpathquadraticcurveto{\pgfqpoint{7.237242in}{0.314243in}}{\pgfqpoint{7.243474in}{0.314243in}}%
\pgfpathquadraticcurveto{\pgfqpoint{7.252102in}{0.314243in}}{\pgfqpoint{7.257307in}{0.320349in}}%
\pgfpathquadraticcurveto{\pgfqpoint{7.262513in}{0.326455in}}{\pgfqpoint{7.262513in}{0.336578in}}%
\pgfpathlineto{\pgfqpoint{7.262513in}{0.338883in}}%
\pgfpathlineto{\pgfqpoint{7.252210in}{0.338883in}}%
\pgfpathlineto{\pgfqpoint{7.252210in}{0.338883in}}%
\pgfpathclose%
\pgfpathmoveto{\pgfqpoint{7.272870in}{0.343170in}}%
\pgfpathlineto{\pgfqpoint{7.272870in}{0.307200in}}%
\pgfpathlineto{\pgfqpoint{7.262513in}{0.307200in}}%
\pgfpathlineto{\pgfqpoint{7.262513in}{0.316764in}}%
\pgfpathquadraticcurveto{\pgfqpoint{7.258964in}{0.311037in}}{\pgfqpoint{7.253669in}{0.308299in}}%
\pgfpathquadraticcurveto{\pgfqpoint{7.248373in}{0.305561in}}{\pgfqpoint{7.240718in}{0.305561in}}%
\pgfpathquadraticcurveto{\pgfqpoint{7.231045in}{0.305561in}}{\pgfqpoint{7.225317in}{0.311001in}}%
\pgfpathquadraticcurveto{\pgfqpoint{7.219608in}{0.316440in}}{\pgfqpoint{7.219608in}{0.325554in}}%
\pgfpathquadraticcurveto{\pgfqpoint{7.219608in}{0.336182in}}{\pgfqpoint{7.226722in}{0.341585in}}%
\pgfpathquadraticcurveto{\pgfqpoint{7.233855in}{0.346989in}}{\pgfqpoint{7.247977in}{0.346989in}}%
\pgfpathlineto{\pgfqpoint{7.262513in}{0.346989in}}%
\pgfpathlineto{\pgfqpoint{7.262513in}{0.348016in}}%
\pgfpathquadraticcurveto{\pgfqpoint{7.262513in}{0.355166in}}{\pgfqpoint{7.257811in}{0.359075in}}%
\pgfpathquadraticcurveto{\pgfqpoint{7.253110in}{0.362984in}}{\pgfqpoint{7.244609in}{0.362984in}}%
\pgfpathquadraticcurveto{\pgfqpoint{7.239205in}{0.362984in}}{\pgfqpoint{7.234071in}{0.361687in}}%
\pgfpathquadraticcurveto{\pgfqpoint{7.228956in}{0.360390in}}{\pgfqpoint{7.224237in}{0.357796in}}%
\pgfpathlineto{\pgfqpoint{7.224237in}{0.367379in}}%
\pgfpathquadraticcurveto{\pgfqpoint{7.229911in}{0.369576in}}{\pgfqpoint{7.235260in}{0.370657in}}%
\pgfpathquadraticcurveto{\pgfqpoint{7.240610in}{0.371756in}}{\pgfqpoint{7.245671in}{0.371756in}}%
\pgfpathquadraticcurveto{\pgfqpoint{7.259360in}{0.371756in}}{\pgfqpoint{7.266115in}{0.364659in}}%
\pgfpathquadraticcurveto{\pgfqpoint{7.272870in}{0.357580in}}{\pgfqpoint{7.272870in}{0.343170in}}%
\pgfpathlineto{\pgfqpoint{7.272870in}{0.343170in}}%
\pgfpathclose%
\pgfpathmoveto{\pgfqpoint{7.330725in}{0.360570in}}%
\pgfpathquadraticcurveto{\pgfqpoint{7.328977in}{0.361579in}}{\pgfqpoint{7.326924in}{0.362047in}}%
\pgfpathquadraticcurveto{\pgfqpoint{7.324871in}{0.362533in}}{\pgfqpoint{7.322403in}{0.362533in}}%
\pgfpathquadraticcurveto{\pgfqpoint{7.313613in}{0.362533in}}{\pgfqpoint{7.308912in}{0.356823in}}%
\pgfpathquadraticcurveto{\pgfqpoint{7.304211in}{0.351114in}}{\pgfqpoint{7.304211in}{0.340414in}}%
\pgfpathlineto{\pgfqpoint{7.304211in}{0.307200in}}%
\pgfpathlineto{\pgfqpoint{7.293800in}{0.307200in}}%
\pgfpathlineto{\pgfqpoint{7.293800in}{0.370243in}}%
\pgfpathlineto{\pgfqpoint{7.304211in}{0.370243in}}%
\pgfpathlineto{\pgfqpoint{7.304211in}{0.360444in}}%
\pgfpathquadraticcurveto{\pgfqpoint{7.307489in}{0.366190in}}{\pgfqpoint{7.312712in}{0.368964in}}%
\pgfpathquadraticcurveto{\pgfqpoint{7.317954in}{0.371756in}}{\pgfqpoint{7.325447in}{0.371756in}}%
\pgfpathquadraticcurveto{\pgfqpoint{7.326510in}{0.371756in}}{\pgfqpoint{7.327807in}{0.371611in}}%
\pgfpathquadraticcurveto{\pgfqpoint{7.329104in}{0.371485in}}{\pgfqpoint{7.330671in}{0.371197in}}%
\pgfpathlineto{\pgfqpoint{7.330725in}{0.360570in}}%
\pgfpathlineto{\pgfqpoint{7.330725in}{0.360570in}}%
\pgfpathclose%
\pgfpathmoveto{\pgfqpoint{7.370243in}{0.338883in}}%
\pgfpathquadraticcurveto{\pgfqpoint{7.357689in}{0.338883in}}{\pgfqpoint{7.352844in}{0.336019in}}%
\pgfpathquadraticcurveto{\pgfqpoint{7.347998in}{0.333156in}}{\pgfqpoint{7.347998in}{0.326221in}}%
\pgfpathquadraticcurveto{\pgfqpoint{7.347998in}{0.320709in}}{\pgfqpoint{7.351637in}{0.317467in}}%
\pgfpathquadraticcurveto{\pgfqpoint{7.355275in}{0.314243in}}{\pgfqpoint{7.361507in}{0.314243in}}%
\pgfpathquadraticcurveto{\pgfqpoint{7.370135in}{0.314243in}}{\pgfqpoint{7.375341in}{0.320349in}}%
\pgfpathquadraticcurveto{\pgfqpoint{7.380546in}{0.326455in}}{\pgfqpoint{7.380546in}{0.336578in}}%
\pgfpathlineto{\pgfqpoint{7.380546in}{0.338883in}}%
\pgfpathlineto{\pgfqpoint{7.370243in}{0.338883in}}%
\pgfpathlineto{\pgfqpoint{7.370243in}{0.338883in}}%
\pgfpathclose%
\pgfpathmoveto{\pgfqpoint{7.390903in}{0.343170in}}%
\pgfpathlineto{\pgfqpoint{7.390903in}{0.307200in}}%
\pgfpathlineto{\pgfqpoint{7.380546in}{0.307200in}}%
\pgfpathlineto{\pgfqpoint{7.380546in}{0.316764in}}%
\pgfpathquadraticcurveto{\pgfqpoint{7.376998in}{0.311037in}}{\pgfqpoint{7.371702in}{0.308299in}}%
\pgfpathquadraticcurveto{\pgfqpoint{7.366407in}{0.305561in}}{\pgfqpoint{7.358752in}{0.305561in}}%
\pgfpathquadraticcurveto{\pgfqpoint{7.349079in}{0.305561in}}{\pgfqpoint{7.343351in}{0.311001in}}%
\pgfpathquadraticcurveto{\pgfqpoint{7.337641in}{0.316440in}}{\pgfqpoint{7.337641in}{0.325554in}}%
\pgfpathquadraticcurveto{\pgfqpoint{7.337641in}{0.336182in}}{\pgfqpoint{7.344756in}{0.341585in}}%
\pgfpathquadraticcurveto{\pgfqpoint{7.351889in}{0.346989in}}{\pgfqpoint{7.366010in}{0.346989in}}%
\pgfpathlineto{\pgfqpoint{7.380546in}{0.346989in}}%
\pgfpathlineto{\pgfqpoint{7.380546in}{0.348016in}}%
\pgfpathquadraticcurveto{\pgfqpoint{7.380546in}{0.355166in}}{\pgfqpoint{7.375845in}{0.359075in}}%
\pgfpathquadraticcurveto{\pgfqpoint{7.371144in}{0.362984in}}{\pgfqpoint{7.362642in}{0.362984in}}%
\pgfpathquadraticcurveto{\pgfqpoint{7.357238in}{0.362984in}}{\pgfqpoint{7.352105in}{0.361687in}}%
\pgfpathquadraticcurveto{\pgfqpoint{7.346990in}{0.360390in}}{\pgfqpoint{7.342270in}{0.357796in}}%
\pgfpathlineto{\pgfqpoint{7.342270in}{0.367379in}}%
\pgfpathquadraticcurveto{\pgfqpoint{7.347944in}{0.369576in}}{\pgfqpoint{7.353294in}{0.370657in}}%
\pgfpathquadraticcurveto{\pgfqpoint{7.358643in}{0.371756in}}{\pgfqpoint{7.363705in}{0.371756in}}%
\pgfpathquadraticcurveto{\pgfqpoint{7.377394in}{0.371756in}}{\pgfqpoint{7.384149in}{0.364659in}}%
\pgfpathquadraticcurveto{\pgfqpoint{7.390903in}{0.357580in}}{\pgfqpoint{7.390903in}{0.343170in}}%
\pgfpathlineto{\pgfqpoint{7.390903in}{0.343170in}}%
\pgfpathclose%
\pgfpathmoveto{\pgfqpoint{7.422479in}{0.388147in}}%
\pgfpathlineto{\pgfqpoint{7.422479in}{0.370243in}}%
\pgfpathlineto{\pgfqpoint{7.443805in}{0.370243in}}%
\pgfpathlineto{\pgfqpoint{7.443805in}{0.362191in}}%
\pgfpathlineto{\pgfqpoint{7.422479in}{0.362191in}}%
\pgfpathlineto{\pgfqpoint{7.422479in}{0.327968in}}%
\pgfpathquadraticcurveto{\pgfqpoint{7.422479in}{0.320259in}}{\pgfqpoint{7.424586in}{0.318061in}}%
\pgfpathquadraticcurveto{\pgfqpoint{7.426693in}{0.315864in}}{\pgfqpoint{7.433178in}{0.315864in}}%
\pgfpathlineto{\pgfqpoint{7.443805in}{0.315864in}}%
\pgfpathlineto{\pgfqpoint{7.443805in}{0.307200in}}%
\pgfpathlineto{\pgfqpoint{7.433178in}{0.307200in}}%
\pgfpathquadraticcurveto{\pgfqpoint{7.421182in}{0.307200in}}{\pgfqpoint{7.416625in}{0.311667in}}%
\pgfpathquadraticcurveto{\pgfqpoint{7.412067in}{0.316152in}}{\pgfqpoint{7.412067in}{0.327968in}}%
\pgfpathlineto{\pgfqpoint{7.412067in}{0.362191in}}%
\pgfpathlineto{\pgfqpoint{7.404466in}{0.362191in}}%
\pgfpathlineto{\pgfqpoint{7.404466in}{0.370243in}}%
\pgfpathlineto{\pgfqpoint{7.412067in}{0.370243in}}%
\pgfpathlineto{\pgfqpoint{7.412067in}{0.388147in}}%
\pgfpathlineto{\pgfqpoint{7.422479in}{0.388147in}}%
\pgfpathlineto{\pgfqpoint{7.422479in}{0.388147in}}%
\pgfpathclose%
\pgfpathmoveto{\pgfqpoint{7.511350in}{0.341315in}}%
\pgfpathlineto{\pgfqpoint{7.511350in}{0.336254in}}%
\pgfpathlineto{\pgfqpoint{7.463726in}{0.336254in}}%
\pgfpathquadraticcurveto{\pgfqpoint{7.464411in}{0.325554in}}{\pgfqpoint{7.470175in}{0.319953in}}%
\pgfpathquadraticcurveto{\pgfqpoint{7.475939in}{0.314351in}}{\pgfqpoint{7.486242in}{0.314351in}}%
\pgfpathquadraticcurveto{\pgfqpoint{7.492204in}{0.314351in}}{\pgfqpoint{7.497805in}{0.315810in}}%
\pgfpathquadraticcurveto{\pgfqpoint{7.503407in}{0.317269in}}{\pgfqpoint{7.508937in}{0.320205in}}%
\pgfpathlineto{\pgfqpoint{7.508937in}{0.310406in}}%
\pgfpathquadraticcurveto{\pgfqpoint{7.503353in}{0.308047in}}{\pgfqpoint{7.497499in}{0.306804in}}%
\pgfpathquadraticcurveto{\pgfqpoint{7.491645in}{0.305561in}}{\pgfqpoint{7.485629in}{0.305561in}}%
\pgfpathquadraticcurveto{\pgfqpoint{7.470535in}{0.305561in}}{\pgfqpoint{7.461727in}{0.314333in}}%
\pgfpathquadraticcurveto{\pgfqpoint{7.452919in}{0.323123in}}{\pgfqpoint{7.452919in}{0.338109in}}%
\pgfpathquadraticcurveto{\pgfqpoint{7.452919in}{0.353581in}}{\pgfqpoint{7.461277in}{0.362659in}}%
\pgfpathquadraticcurveto{\pgfqpoint{7.469634in}{0.371756in}}{\pgfqpoint{7.483828in}{0.371756in}}%
\pgfpathquadraticcurveto{\pgfqpoint{7.496544in}{0.371756in}}{\pgfqpoint{7.503947in}{0.363560in}}%
\pgfpathquadraticcurveto{\pgfqpoint{7.511350in}{0.355383in}}{\pgfqpoint{7.511350in}{0.341315in}}%
\pgfpathlineto{\pgfqpoint{7.511350in}{0.341315in}}%
\pgfpathclose%
\pgfpathmoveto{\pgfqpoint{7.500993in}{0.344359in}}%
\pgfpathquadraticcurveto{\pgfqpoint{7.500885in}{0.352843in}}{\pgfqpoint{7.496238in}{0.357904in}}%
\pgfpathquadraticcurveto{\pgfqpoint{7.491591in}{0.362984in}}{\pgfqpoint{7.483936in}{0.362984in}}%
\pgfpathquadraticcurveto{\pgfqpoint{7.475272in}{0.362984in}}{\pgfqpoint{7.470067in}{0.358084in}}%
\pgfpathquadraticcurveto{\pgfqpoint{7.464861in}{0.353185in}}{\pgfqpoint{7.464069in}{0.344287in}}%
\pgfpathlineto{\pgfqpoint{7.500993in}{0.344359in}}%
\pgfpathlineto{\pgfqpoint{7.500993in}{0.344359in}}%
\pgfpathclose%
\pgfpathmoveto{\pgfqpoint{7.528354in}{0.394793in}}%
\pgfpathlineto{\pgfqpoint{7.538711in}{0.394793in}}%
\pgfpathlineto{\pgfqpoint{7.538711in}{0.307200in}}%
\pgfpathlineto{\pgfqpoint{7.528354in}{0.307200in}}%
\pgfpathlineto{\pgfqpoint{7.528354in}{0.394793in}}%
\pgfpathlineto{\pgfqpoint{7.528354in}{0.394793in}}%
\pgfpathclose%
\pgfpathmoveto{\pgfqpoint{7.586605in}{0.301346in}}%
\pgfpathquadraticcurveto{\pgfqpoint{7.582228in}{0.290088in}}{\pgfqpoint{7.578049in}{0.286666in}}%
\pgfpathquadraticcurveto{\pgfqpoint{7.573889in}{0.283226in}}{\pgfqpoint{7.566918in}{0.283226in}}%
\pgfpathlineto{\pgfqpoint{7.558632in}{0.283226in}}%
\pgfpathlineto{\pgfqpoint{7.558632in}{0.291890in}}%
\pgfpathlineto{\pgfqpoint{7.564720in}{0.291890in}}%
\pgfpathquadraticcurveto{\pgfqpoint{7.568989in}{0.291890in}}{\pgfqpoint{7.571349in}{0.293925in}}%
\pgfpathquadraticcurveto{\pgfqpoint{7.573727in}{0.295942in}}{\pgfqpoint{7.576590in}{0.303489in}}%
\pgfpathlineto{\pgfqpoint{7.578446in}{0.308209in}}%
\pgfpathlineto{\pgfqpoint{7.552959in}{0.370243in}}%
\pgfpathlineto{\pgfqpoint{7.563928in}{0.370243in}}%
\pgfpathlineto{\pgfqpoint{7.583633in}{0.320943in}}%
\pgfpathlineto{\pgfqpoint{7.603339in}{0.370243in}}%
\pgfpathlineto{\pgfqpoint{7.614308in}{0.370243in}}%
\pgfpathlineto{\pgfqpoint{7.586605in}{0.301346in}}%
\pgfpathlineto{\pgfqpoint{7.586605in}{0.301346in}}%
\pgfpathclose%
\pgfpathmoveto{\pgfqpoint{7.613641in}{0.321502in}}%
\pgfpathlineto{\pgfqpoint{7.625530in}{0.321502in}}%
\pgfpathlineto{\pgfqpoint{7.625530in}{0.307200in}}%
\pgfpathlineto{\pgfqpoint{7.613641in}{0.307200in}}%
\pgfpathlineto{\pgfqpoint{7.613641in}{0.321502in}}%
\pgfpathlineto{\pgfqpoint{7.613641in}{0.321502in}}%
\pgfpathclose%
\pgfpathmoveto{\pgfqpoint{7.736286in}{0.388489in}}%
\pgfpathlineto{\pgfqpoint{7.736286in}{0.377393in}}%
\pgfpathquadraticcurveto{\pgfqpoint{7.729820in}{0.380491in}}{\pgfqpoint{7.724074in}{0.382004in}}%
\pgfpathquadraticcurveto{\pgfqpoint{7.718328in}{0.383536in}}{\pgfqpoint{7.712979in}{0.383536in}}%
\pgfpathquadraticcurveto{\pgfqpoint{7.703702in}{0.383536in}}{\pgfqpoint{7.698659in}{0.379933in}}%
\pgfpathquadraticcurveto{\pgfqpoint{7.693615in}{0.376331in}}{\pgfqpoint{7.693615in}{0.369684in}}%
\pgfpathquadraticcurveto{\pgfqpoint{7.693615in}{0.364100in}}{\pgfqpoint{7.696966in}{0.361254in}}%
\pgfpathquadraticcurveto{\pgfqpoint{7.700316in}{0.358427in}}{\pgfqpoint{7.709664in}{0.356679in}}%
\pgfpathlineto{\pgfqpoint{7.716527in}{0.355274in}}%
\pgfpathquadraticcurveto{\pgfqpoint{7.729243in}{0.352843in}}{\pgfqpoint{7.735296in}{0.346737in}}%
\pgfpathquadraticcurveto{\pgfqpoint{7.741348in}{0.340631in}}{\pgfqpoint{7.741348in}{0.330400in}}%
\pgfpathquadraticcurveto{\pgfqpoint{7.741348in}{0.318169in}}{\pgfqpoint{7.733152in}{0.311865in}}%
\pgfpathquadraticcurveto{\pgfqpoint{7.724975in}{0.305561in}}{\pgfqpoint{7.709160in}{0.305561in}}%
\pgfpathquadraticcurveto{\pgfqpoint{7.703198in}{0.305561in}}{\pgfqpoint{7.696461in}{0.306912in}}%
\pgfpathquadraticcurveto{\pgfqpoint{7.689743in}{0.308263in}}{\pgfqpoint{7.682538in}{0.310911in}}%
\pgfpathlineto{\pgfqpoint{7.682538in}{0.322618in}}%
\pgfpathquadraticcurveto{\pgfqpoint{7.689455in}{0.318746in}}{\pgfqpoint{7.696101in}{0.316764in}}%
\pgfpathquadraticcurveto{\pgfqpoint{7.702748in}{0.314801in}}{\pgfqpoint{7.709160in}{0.314801in}}%
\pgfpathquadraticcurveto{\pgfqpoint{7.718887in}{0.314801in}}{\pgfqpoint{7.724182in}{0.318620in}}%
\pgfpathquadraticcurveto{\pgfqpoint{7.729478in}{0.322456in}}{\pgfqpoint{7.729478in}{0.329553in}}%
\pgfpathquadraticcurveto{\pgfqpoint{7.729478in}{0.335731in}}{\pgfqpoint{7.725677in}{0.339226in}}%
\pgfpathquadraticcurveto{\pgfqpoint{7.721877in}{0.342720in}}{\pgfqpoint{7.713213in}{0.344467in}}%
\pgfpathlineto{\pgfqpoint{7.706278in}{0.345818in}}%
\pgfpathquadraticcurveto{\pgfqpoint{7.693561in}{0.348340in}}{\pgfqpoint{7.687870in}{0.353743in}}%
\pgfpathquadraticcurveto{\pgfqpoint{7.682196in}{0.359147in}}{\pgfqpoint{7.682196in}{0.368784in}}%
\pgfpathquadraticcurveto{\pgfqpoint{7.682196in}{0.379933in}}{\pgfqpoint{7.690049in}{0.386345in}}%
\pgfpathquadraticcurveto{\pgfqpoint{7.697902in}{0.392758in}}{\pgfqpoint{7.711682in}{0.392758in}}%
\pgfpathquadraticcurveto{\pgfqpoint{7.717608in}{0.392758in}}{\pgfqpoint{7.723732in}{0.391677in}}%
\pgfpathquadraticcurveto{\pgfqpoint{7.729874in}{0.390614in}}{\pgfqpoint{7.736286in}{0.388489in}}%
\pgfpathlineto{\pgfqpoint{7.736286in}{0.388489in}}%
\pgfpathclose%
\pgfpathmoveto{\pgfqpoint{7.768888in}{0.388147in}}%
\pgfpathlineto{\pgfqpoint{7.768888in}{0.370243in}}%
\pgfpathlineto{\pgfqpoint{7.790215in}{0.370243in}}%
\pgfpathlineto{\pgfqpoint{7.790215in}{0.362191in}}%
\pgfpathlineto{\pgfqpoint{7.768888in}{0.362191in}}%
\pgfpathlineto{\pgfqpoint{7.768888in}{0.327968in}}%
\pgfpathquadraticcurveto{\pgfqpoint{7.768888in}{0.320259in}}{\pgfqpoint{7.770996in}{0.318061in}}%
\pgfpathquadraticcurveto{\pgfqpoint{7.773103in}{0.315864in}}{\pgfqpoint{7.779587in}{0.315864in}}%
\pgfpathlineto{\pgfqpoint{7.790215in}{0.315864in}}%
\pgfpathlineto{\pgfqpoint{7.790215in}{0.307200in}}%
\pgfpathlineto{\pgfqpoint{7.779587in}{0.307200in}}%
\pgfpathquadraticcurveto{\pgfqpoint{7.767591in}{0.307200in}}{\pgfqpoint{7.763034in}{0.311667in}}%
\pgfpathquadraticcurveto{\pgfqpoint{7.758477in}{0.316152in}}{\pgfqpoint{7.758477in}{0.327968in}}%
\pgfpathlineto{\pgfqpoint{7.758477in}{0.362191in}}%
\pgfpathlineto{\pgfqpoint{7.750876in}{0.362191in}}%
\pgfpathlineto{\pgfqpoint{7.750876in}{0.370243in}}%
\pgfpathlineto{\pgfqpoint{7.758477in}{0.370243in}}%
\pgfpathlineto{\pgfqpoint{7.758477in}{0.388147in}}%
\pgfpathlineto{\pgfqpoint{7.768888in}{0.388147in}}%
\pgfpathlineto{\pgfqpoint{7.768888in}{0.388147in}}%
\pgfpathclose%
\pgfpathmoveto{\pgfqpoint{7.802769in}{0.332075in}}%
\pgfpathlineto{\pgfqpoint{7.802769in}{0.370243in}}%
\pgfpathlineto{\pgfqpoint{7.813126in}{0.370243in}}%
\pgfpathlineto{\pgfqpoint{7.813126in}{0.332471in}}%
\pgfpathquadraticcurveto{\pgfqpoint{7.813126in}{0.323519in}}{\pgfqpoint{7.816602in}{0.319034in}}%
\pgfpathquadraticcurveto{\pgfqpoint{7.820097in}{0.314567in}}{\pgfqpoint{7.827086in}{0.314567in}}%
\pgfpathquadraticcurveto{\pgfqpoint{7.835461in}{0.314567in}}{\pgfqpoint{7.840324in}{0.319917in}}%
\pgfpathquadraticcurveto{\pgfqpoint{7.845206in}{0.325266in}}{\pgfqpoint{7.845206in}{0.334506in}}%
\pgfpathlineto{\pgfqpoint{7.845206in}{0.370243in}}%
\pgfpathlineto{\pgfqpoint{7.855563in}{0.370243in}}%
\pgfpathlineto{\pgfqpoint{7.855563in}{0.307200in}}%
\pgfpathlineto{\pgfqpoint{7.845206in}{0.307200in}}%
\pgfpathlineto{\pgfqpoint{7.845206in}{0.316891in}}%
\pgfpathquadraticcurveto{\pgfqpoint{7.841441in}{0.311145in}}{\pgfqpoint{7.836452in}{0.308353in}}%
\pgfpathquadraticcurveto{\pgfqpoint{7.831480in}{0.305561in}}{\pgfqpoint{7.824888in}{0.305561in}}%
\pgfpathquadraticcurveto{\pgfqpoint{7.814027in}{0.305561in}}{\pgfqpoint{7.808389in}{0.312315in}}%
\pgfpathquadraticcurveto{\pgfqpoint{7.802769in}{0.319070in}}{\pgfqpoint{7.802769in}{0.332075in}}%
\pgfpathlineto{\pgfqpoint{7.802769in}{0.332075in}}%
\pgfpathclose%
\pgfpathmoveto{\pgfqpoint{7.828833in}{0.371756in}}%
\pgfpathlineto{\pgfqpoint{7.828833in}{0.371756in}}%
\pgfpathlineto{\pgfqpoint{7.828833in}{0.371756in}}%
\pgfpathclose%
\pgfpathmoveto{\pgfqpoint{7.918371in}{0.360678in}}%
\pgfpathlineto{\pgfqpoint{7.918371in}{0.394793in}}%
\pgfpathlineto{\pgfqpoint{7.928728in}{0.394793in}}%
\pgfpathlineto{\pgfqpoint{7.928728in}{0.307200in}}%
\pgfpathlineto{\pgfqpoint{7.918371in}{0.307200in}}%
\pgfpathlineto{\pgfqpoint{7.918371in}{0.316656in}}%
\pgfpathquadraticcurveto{\pgfqpoint{7.915111in}{0.311037in}}{\pgfqpoint{7.910122in}{0.308299in}}%
\pgfpathquadraticcurveto{\pgfqpoint{7.905150in}{0.305561in}}{\pgfqpoint{7.898161in}{0.305561in}}%
\pgfpathquadraticcurveto{\pgfqpoint{7.886742in}{0.305561in}}{\pgfqpoint{7.879555in}{0.314675in}}%
\pgfpathquadraticcurveto{\pgfqpoint{7.872386in}{0.323807in}}{\pgfqpoint{7.872386in}{0.338667in}}%
\pgfpathquadraticcurveto{\pgfqpoint{7.872386in}{0.353527in}}{\pgfqpoint{7.879555in}{0.362641in}}%
\pgfpathquadraticcurveto{\pgfqpoint{7.886742in}{0.371756in}}{\pgfqpoint{7.898161in}{0.371756in}}%
\pgfpathquadraticcurveto{\pgfqpoint{7.905150in}{0.371756in}}{\pgfqpoint{7.910122in}{0.369018in}}%
\pgfpathquadraticcurveto{\pgfqpoint{7.915111in}{0.366298in}}{\pgfqpoint{7.918371in}{0.360678in}}%
\pgfpathlineto{\pgfqpoint{7.918371in}{0.360678in}}%
\pgfpathclose%
\pgfpathmoveto{\pgfqpoint{7.883085in}{0.338667in}}%
\pgfpathquadraticcurveto{\pgfqpoint{7.883085in}{0.327248in}}{\pgfqpoint{7.887786in}{0.320745in}}%
\pgfpathquadraticcurveto{\pgfqpoint{7.892488in}{0.314243in}}{\pgfqpoint{7.900701in}{0.314243in}}%
\pgfpathquadraticcurveto{\pgfqpoint{7.908915in}{0.314243in}}{\pgfqpoint{7.913634in}{0.320745in}}%
\pgfpathquadraticcurveto{\pgfqpoint{7.918371in}{0.327248in}}{\pgfqpoint{7.918371in}{0.338667in}}%
\pgfpathquadraticcurveto{\pgfqpoint{7.918371in}{0.350087in}}{\pgfqpoint{7.913634in}{0.356589in}}%
\pgfpathquadraticcurveto{\pgfqpoint{7.908915in}{0.363092in}}{\pgfqpoint{7.900701in}{0.363092in}}%
\pgfpathquadraticcurveto{\pgfqpoint{7.892488in}{0.363092in}}{\pgfqpoint{7.887786in}{0.356589in}}%
\pgfpathquadraticcurveto{\pgfqpoint{7.883085in}{0.350087in}}{\pgfqpoint{7.883085in}{0.338667in}}%
\pgfpathlineto{\pgfqpoint{7.883085in}{0.338667in}}%
\pgfpathclose%
\pgfpathmoveto{\pgfqpoint{7.976298in}{0.301346in}}%
\pgfpathquadraticcurveto{\pgfqpoint{7.971921in}{0.290088in}}{\pgfqpoint{7.967742in}{0.286666in}}%
\pgfpathquadraticcurveto{\pgfqpoint{7.963582in}{0.283226in}}{\pgfqpoint{7.956611in}{0.283226in}}%
\pgfpathlineto{\pgfqpoint{7.948325in}{0.283226in}}%
\pgfpathlineto{\pgfqpoint{7.948325in}{0.291890in}}%
\pgfpathlineto{\pgfqpoint{7.954413in}{0.291890in}}%
\pgfpathquadraticcurveto{\pgfqpoint{7.958682in}{0.291890in}}{\pgfqpoint{7.961042in}{0.293925in}}%
\pgfpathquadraticcurveto{\pgfqpoint{7.963419in}{0.295942in}}{\pgfqpoint{7.966283in}{0.303489in}}%
\pgfpathlineto{\pgfqpoint{7.968139in}{0.308209in}}%
\pgfpathlineto{\pgfqpoint{7.942651in}{0.370243in}}%
\pgfpathlineto{\pgfqpoint{7.953621in}{0.370243in}}%
\pgfpathlineto{\pgfqpoint{7.973326in}{0.320943in}}%
\pgfpathlineto{\pgfqpoint{7.993031in}{0.370243in}}%
\pgfpathlineto{\pgfqpoint{8.004001in}{0.370243in}}%
\pgfpathlineto{\pgfqpoint{7.976298in}{0.301346in}}%
\pgfpathlineto{\pgfqpoint{7.976298in}{0.301346in}}%
\pgfpathclose%
\pgfpathmoveto{\pgfqpoint{8.087649in}{0.381338in}}%
\pgfpathlineto{\pgfqpoint{8.058938in}{0.336470in}}%
\pgfpathlineto{\pgfqpoint{8.087649in}{0.336470in}}%
\pgfpathlineto{\pgfqpoint{8.087649in}{0.381338in}}%
\pgfpathlineto{\pgfqpoint{8.087649in}{0.381338in}}%
\pgfpathclose%
\pgfpathmoveto{\pgfqpoint{8.084659in}{0.391245in}}%
\pgfpathlineto{\pgfqpoint{8.098961in}{0.391245in}}%
\pgfpathlineto{\pgfqpoint{8.098961in}{0.336470in}}%
\pgfpathlineto{\pgfqpoint{8.110957in}{0.336470in}}%
\pgfpathlineto{\pgfqpoint{8.110957in}{0.327013in}}%
\pgfpathlineto{\pgfqpoint{8.098961in}{0.327013in}}%
\pgfpathlineto{\pgfqpoint{8.098961in}{0.307200in}}%
\pgfpathlineto{\pgfqpoint{8.087649in}{0.307200in}}%
\pgfpathlineto{\pgfqpoint{8.087649in}{0.327013in}}%
\pgfpathlineto{\pgfqpoint{8.049716in}{0.327013in}}%
\pgfpathlineto{\pgfqpoint{8.049716in}{0.337983in}}%
\pgfpathlineto{\pgfqpoint{8.084659in}{0.391245in}}%
\pgfpathlineto{\pgfqpoint{8.084659in}{0.391245in}}%
\pgfpathclose%
\pgfpathmoveto{\pgfqpoint{8.210294in}{0.367829in}}%
\pgfpathlineto{\pgfqpoint{8.210294in}{0.358138in}}%
\pgfpathquadraticcurveto{\pgfqpoint{8.205899in}{0.360570in}}{\pgfqpoint{8.201486in}{0.361777in}}%
\pgfpathquadraticcurveto{\pgfqpoint{8.197073in}{0.362984in}}{\pgfqpoint{8.192570in}{0.362984in}}%
\pgfpathquadraticcurveto{\pgfqpoint{8.182483in}{0.362984in}}{\pgfqpoint{8.176900in}{0.356589in}}%
\pgfpathquadraticcurveto{\pgfqpoint{8.171334in}{0.350213in}}{\pgfqpoint{8.171334in}{0.338667in}}%
\pgfpathquadraticcurveto{\pgfqpoint{8.171334in}{0.327121in}}{\pgfqpoint{8.176900in}{0.320727in}}%
\pgfpathquadraticcurveto{\pgfqpoint{8.182483in}{0.314351in}}{\pgfqpoint{8.192570in}{0.314351in}}%
\pgfpathquadraticcurveto{\pgfqpoint{8.197073in}{0.314351in}}{\pgfqpoint{8.201486in}{0.315558in}}%
\pgfpathquadraticcurveto{\pgfqpoint{8.205899in}{0.316764in}}{\pgfqpoint{8.210294in}{0.319196in}}%
\pgfpathlineto{\pgfqpoint{8.210294in}{0.309614in}}%
\pgfpathquadraticcurveto{\pgfqpoint{8.205953in}{0.307596in}}{\pgfqpoint{8.201306in}{0.306588in}}%
\pgfpathquadraticcurveto{\pgfqpoint{8.196677in}{0.305561in}}{\pgfqpoint{8.191435in}{0.305561in}}%
\pgfpathquadraticcurveto{\pgfqpoint{8.177188in}{0.305561in}}{\pgfqpoint{8.168794in}{0.314513in}}%
\pgfpathquadraticcurveto{\pgfqpoint{8.160418in}{0.323465in}}{\pgfqpoint{8.160418in}{0.338667in}}%
\pgfpathquadraticcurveto{\pgfqpoint{8.160418in}{0.354086in}}{\pgfqpoint{8.168884in}{0.362912in}}%
\pgfpathquadraticcurveto{\pgfqpoint{8.177368in}{0.371756in}}{\pgfqpoint{8.192120in}{0.371756in}}%
\pgfpathquadraticcurveto{\pgfqpoint{8.196893in}{0.371756in}}{\pgfqpoint{8.201450in}{0.370765in}}%
\pgfpathquadraticcurveto{\pgfqpoint{8.206007in}{0.369792in}}{\pgfqpoint{8.210294in}{0.367829in}}%
\pgfpathlineto{\pgfqpoint{8.210294in}{0.367829in}}%
\pgfpathclose%
\pgfpathmoveto{\pgfqpoint{8.228306in}{0.394793in}}%
\pgfpathlineto{\pgfqpoint{8.238663in}{0.394793in}}%
\pgfpathlineto{\pgfqpoint{8.238663in}{0.307200in}}%
\pgfpathlineto{\pgfqpoint{8.228306in}{0.307200in}}%
\pgfpathlineto{\pgfqpoint{8.228306in}{0.394793in}}%
\pgfpathlineto{\pgfqpoint{8.228306in}{0.394793in}}%
\pgfpathclose%
\pgfpathmoveto{\pgfqpoint{8.259269in}{0.332075in}}%
\pgfpathlineto{\pgfqpoint{8.259269in}{0.370243in}}%
\pgfpathlineto{\pgfqpoint{8.269626in}{0.370243in}}%
\pgfpathlineto{\pgfqpoint{8.269626in}{0.332471in}}%
\pgfpathquadraticcurveto{\pgfqpoint{8.269626in}{0.323519in}}{\pgfqpoint{8.273102in}{0.319034in}}%
\pgfpathquadraticcurveto{\pgfqpoint{8.276597in}{0.314567in}}{\pgfqpoint{8.283586in}{0.314567in}}%
\pgfpathquadraticcurveto{\pgfqpoint{8.291961in}{0.314567in}}{\pgfqpoint{8.296824in}{0.319917in}}%
\pgfpathquadraticcurveto{\pgfqpoint{8.301706in}{0.325266in}}{\pgfqpoint{8.301706in}{0.334506in}}%
\pgfpathlineto{\pgfqpoint{8.301706in}{0.370243in}}%
\pgfpathlineto{\pgfqpoint{8.312063in}{0.370243in}}%
\pgfpathlineto{\pgfqpoint{8.312063in}{0.307200in}}%
\pgfpathlineto{\pgfqpoint{8.301706in}{0.307200in}}%
\pgfpathlineto{\pgfqpoint{8.301706in}{0.316891in}}%
\pgfpathquadraticcurveto{\pgfqpoint{8.297941in}{0.311145in}}{\pgfqpoint{8.292952in}{0.308353in}}%
\pgfpathquadraticcurveto{\pgfqpoint{8.287980in}{0.305561in}}{\pgfqpoint{8.281388in}{0.305561in}}%
\pgfpathquadraticcurveto{\pgfqpoint{8.270527in}{0.305561in}}{\pgfqpoint{8.264889in}{0.312315in}}%
\pgfpathquadraticcurveto{\pgfqpoint{8.259269in}{0.319070in}}{\pgfqpoint{8.259269in}{0.332075in}}%
\pgfpathlineto{\pgfqpoint{8.259269in}{0.332075in}}%
\pgfpathclose%
\pgfpathmoveto{\pgfqpoint{8.285333in}{0.371756in}}%
\pgfpathlineto{\pgfqpoint{8.285333in}{0.371756in}}%
\pgfpathlineto{\pgfqpoint{8.285333in}{0.371756in}}%
\pgfpathclose%
\pgfpathmoveto{\pgfqpoint{8.373574in}{0.368387in}}%
\pgfpathlineto{\pgfqpoint{8.373574in}{0.358589in}}%
\pgfpathquadraticcurveto{\pgfqpoint{8.369197in}{0.360840in}}{\pgfqpoint{8.364460in}{0.361957in}}%
\pgfpathquadraticcurveto{\pgfqpoint{8.359741in}{0.363092in}}{\pgfqpoint{8.354661in}{0.363092in}}%
\pgfpathquadraticcurveto{\pgfqpoint{8.346952in}{0.363092in}}{\pgfqpoint{8.343098in}{0.360732in}}%
\pgfpathquadraticcurveto{\pgfqpoint{8.339243in}{0.358373in}}{\pgfqpoint{8.339243in}{0.353635in}}%
\pgfpathquadraticcurveto{\pgfqpoint{8.339243in}{0.350033in}}{\pgfqpoint{8.341999in}{0.347980in}}%
\pgfpathquadraticcurveto{\pgfqpoint{8.344755in}{0.345926in}}{\pgfqpoint{8.353094in}{0.344071in}}%
\pgfpathlineto{\pgfqpoint{8.356643in}{0.343278in}}%
\pgfpathquadraticcurveto{\pgfqpoint{8.367666in}{0.340919in}}{\pgfqpoint{8.372313in}{0.336614in}}%
\pgfpathquadraticcurveto{\pgfqpoint{8.376961in}{0.332309in}}{\pgfqpoint{8.376961in}{0.324600in}}%
\pgfpathquadraticcurveto{\pgfqpoint{8.376961in}{0.315810in}}{\pgfqpoint{8.370008in}{0.310676in}}%
\pgfpathquadraticcurveto{\pgfqpoint{8.363055in}{0.305561in}}{\pgfqpoint{8.350897in}{0.305561in}}%
\pgfpathquadraticcurveto{\pgfqpoint{8.345836in}{0.305561in}}{\pgfqpoint{8.340342in}{0.306552in}}%
\pgfpathquadraticcurveto{\pgfqpoint{8.334848in}{0.307542in}}{\pgfqpoint{8.328778in}{0.309506in}}%
\pgfpathlineto{\pgfqpoint{8.328778in}{0.320205in}}%
\pgfpathquadraticcurveto{\pgfqpoint{8.334524in}{0.317215in}}{\pgfqpoint{8.340090in}{0.315720in}}%
\pgfpathquadraticcurveto{\pgfqpoint{8.345655in}{0.314243in}}{\pgfqpoint{8.351131in}{0.314243in}}%
\pgfpathquadraticcurveto{\pgfqpoint{8.358444in}{0.314243in}}{\pgfqpoint{8.362371in}{0.316746in}}%
\pgfpathquadraticcurveto{\pgfqpoint{8.366315in}{0.319250in}}{\pgfqpoint{8.366315in}{0.323807in}}%
\pgfpathquadraticcurveto{\pgfqpoint{8.366315in}{0.328022in}}{\pgfqpoint{8.363469in}{0.330274in}}%
\pgfpathquadraticcurveto{\pgfqpoint{8.360641in}{0.332525in}}{\pgfqpoint{8.351005in}{0.334614in}}%
\pgfpathlineto{\pgfqpoint{8.347403in}{0.335461in}}%
\pgfpathquadraticcurveto{\pgfqpoint{8.337784in}{0.337478in}}{\pgfqpoint{8.333497in}{0.341675in}}%
\pgfpathquadraticcurveto{\pgfqpoint{8.329228in}{0.345872in}}{\pgfqpoint{8.329228in}{0.353185in}}%
\pgfpathquadraticcurveto{\pgfqpoint{8.329228in}{0.362083in}}{\pgfqpoint{8.335533in}{0.366910in}}%
\pgfpathquadraticcurveto{\pgfqpoint{8.341837in}{0.371756in}}{\pgfqpoint{8.353437in}{0.371756in}}%
\pgfpathquadraticcurveto{\pgfqpoint{8.359164in}{0.371756in}}{\pgfqpoint{8.364226in}{0.370909in}}%
\pgfpathquadraticcurveto{\pgfqpoint{8.369305in}{0.370080in}}{\pgfqpoint{8.373574in}{0.368387in}}%
\pgfpathlineto{\pgfqpoint{8.373574in}{0.368387in}}%
\pgfpathclose%
\pgfpathmoveto{\pgfqpoint{8.403691in}{0.388147in}}%
\pgfpathlineto{\pgfqpoint{8.403691in}{0.370243in}}%
\pgfpathlineto{\pgfqpoint{8.425017in}{0.370243in}}%
\pgfpathlineto{\pgfqpoint{8.425017in}{0.362191in}}%
\pgfpathlineto{\pgfqpoint{8.403691in}{0.362191in}}%
\pgfpathlineto{\pgfqpoint{8.403691in}{0.327968in}}%
\pgfpathquadraticcurveto{\pgfqpoint{8.403691in}{0.320259in}}{\pgfqpoint{8.405798in}{0.318061in}}%
\pgfpathquadraticcurveto{\pgfqpoint{8.407905in}{0.315864in}}{\pgfqpoint{8.414390in}{0.315864in}}%
\pgfpathlineto{\pgfqpoint{8.425017in}{0.315864in}}%
\pgfpathlineto{\pgfqpoint{8.425017in}{0.307200in}}%
\pgfpathlineto{\pgfqpoint{8.414390in}{0.307200in}}%
\pgfpathquadraticcurveto{\pgfqpoint{8.402394in}{0.307200in}}{\pgfqpoint{8.397837in}{0.311667in}}%
\pgfpathquadraticcurveto{\pgfqpoint{8.393280in}{0.316152in}}{\pgfqpoint{8.393280in}{0.327968in}}%
\pgfpathlineto{\pgfqpoint{8.393280in}{0.362191in}}%
\pgfpathlineto{\pgfqpoint{8.385678in}{0.362191in}}%
\pgfpathlineto{\pgfqpoint{8.385678in}{0.370243in}}%
\pgfpathlineto{\pgfqpoint{8.393280in}{0.370243in}}%
\pgfpathlineto{\pgfqpoint{8.393280in}{0.388147in}}%
\pgfpathlineto{\pgfqpoint{8.403691in}{0.388147in}}%
\pgfpathlineto{\pgfqpoint{8.403691in}{0.388147in}}%
\pgfpathclose%
\pgfpathmoveto{\pgfqpoint{8.492563in}{0.341315in}}%
\pgfpathlineto{\pgfqpoint{8.492563in}{0.336254in}}%
\pgfpathlineto{\pgfqpoint{8.444938in}{0.336254in}}%
\pgfpathquadraticcurveto{\pgfqpoint{8.445623in}{0.325554in}}{\pgfqpoint{8.451387in}{0.319953in}}%
\pgfpathquadraticcurveto{\pgfqpoint{8.457151in}{0.314351in}}{\pgfqpoint{8.467454in}{0.314351in}}%
\pgfpathquadraticcurveto{\pgfqpoint{8.473416in}{0.314351in}}{\pgfqpoint{8.479017in}{0.315810in}}%
\pgfpathquadraticcurveto{\pgfqpoint{8.484619in}{0.317269in}}{\pgfqpoint{8.490149in}{0.320205in}}%
\pgfpathlineto{\pgfqpoint{8.490149in}{0.310406in}}%
\pgfpathquadraticcurveto{\pgfqpoint{8.484565in}{0.308047in}}{\pgfqpoint{8.478711in}{0.306804in}}%
\pgfpathquadraticcurveto{\pgfqpoint{8.472857in}{0.305561in}}{\pgfqpoint{8.466841in}{0.305561in}}%
\pgfpathquadraticcurveto{\pgfqpoint{8.451747in}{0.305561in}}{\pgfqpoint{8.442939in}{0.314333in}}%
\pgfpathquadraticcurveto{\pgfqpoint{8.434131in}{0.323123in}}{\pgfqpoint{8.434131in}{0.338109in}}%
\pgfpathquadraticcurveto{\pgfqpoint{8.434131in}{0.353581in}}{\pgfqpoint{8.442489in}{0.362659in}}%
\pgfpathquadraticcurveto{\pgfqpoint{8.450846in}{0.371756in}}{\pgfqpoint{8.465040in}{0.371756in}}%
\pgfpathquadraticcurveto{\pgfqpoint{8.477757in}{0.371756in}}{\pgfqpoint{8.485160in}{0.363560in}}%
\pgfpathquadraticcurveto{\pgfqpoint{8.492563in}{0.355383in}}{\pgfqpoint{8.492563in}{0.341315in}}%
\pgfpathlineto{\pgfqpoint{8.492563in}{0.341315in}}%
\pgfpathclose%
\pgfpathmoveto{\pgfqpoint{8.482206in}{0.344359in}}%
\pgfpathquadraticcurveto{\pgfqpoint{8.482097in}{0.352843in}}{\pgfqpoint{8.477450in}{0.357904in}}%
\pgfpathquadraticcurveto{\pgfqpoint{8.472803in}{0.362984in}}{\pgfqpoint{8.465148in}{0.362984in}}%
\pgfpathquadraticcurveto{\pgfqpoint{8.456484in}{0.362984in}}{\pgfqpoint{8.451279in}{0.358084in}}%
\pgfpathquadraticcurveto{\pgfqpoint{8.446073in}{0.353185in}}{\pgfqpoint{8.445281in}{0.344287in}}%
\pgfpathlineto{\pgfqpoint{8.482206in}{0.344359in}}%
\pgfpathlineto{\pgfqpoint{8.482206in}{0.344359in}}%
\pgfpathclose%
\pgfpathmoveto{\pgfqpoint{8.546095in}{0.360570in}}%
\pgfpathquadraticcurveto{\pgfqpoint{8.544347in}{0.361579in}}{\pgfqpoint{8.542294in}{0.362047in}}%
\pgfpathquadraticcurveto{\pgfqpoint{8.540241in}{0.362533in}}{\pgfqpoint{8.537773in}{0.362533in}}%
\pgfpathquadraticcurveto{\pgfqpoint{8.528983in}{0.362533in}}{\pgfqpoint{8.524282in}{0.356823in}}%
\pgfpathquadraticcurveto{\pgfqpoint{8.519581in}{0.351114in}}{\pgfqpoint{8.519581in}{0.340414in}}%
\pgfpathlineto{\pgfqpoint{8.519581in}{0.307200in}}%
\pgfpathlineto{\pgfqpoint{8.509170in}{0.307200in}}%
\pgfpathlineto{\pgfqpoint{8.509170in}{0.370243in}}%
\pgfpathlineto{\pgfqpoint{8.519581in}{0.370243in}}%
\pgfpathlineto{\pgfqpoint{8.519581in}{0.360444in}}%
\pgfpathquadraticcurveto{\pgfqpoint{8.522859in}{0.366190in}}{\pgfqpoint{8.528082in}{0.368964in}}%
\pgfpathquadraticcurveto{\pgfqpoint{8.533324in}{0.371756in}}{\pgfqpoint{8.540817in}{0.371756in}}%
\pgfpathquadraticcurveto{\pgfqpoint{8.541880in}{0.371756in}}{\pgfqpoint{8.543177in}{0.371611in}}%
\pgfpathquadraticcurveto{\pgfqpoint{8.544474in}{0.371485in}}{\pgfqpoint{8.546041in}{0.371197in}}%
\pgfpathlineto{\pgfqpoint{8.546095in}{0.360570in}}%
\pgfpathlineto{\pgfqpoint{8.546095in}{0.360570in}}%
\pgfpathclose%
\pgfpathmoveto{\pgfqpoint{8.597141in}{0.368387in}}%
\pgfpathlineto{\pgfqpoint{8.597141in}{0.358589in}}%
\pgfpathquadraticcurveto{\pgfqpoint{8.592764in}{0.360840in}}{\pgfqpoint{8.588027in}{0.361957in}}%
\pgfpathquadraticcurveto{\pgfqpoint{8.583308in}{0.363092in}}{\pgfqpoint{8.578228in}{0.363092in}}%
\pgfpathquadraticcurveto{\pgfqpoint{8.570519in}{0.363092in}}{\pgfqpoint{8.566664in}{0.360732in}}%
\pgfpathquadraticcurveto{\pgfqpoint{8.562810in}{0.358373in}}{\pgfqpoint{8.562810in}{0.353635in}}%
\pgfpathquadraticcurveto{\pgfqpoint{8.562810in}{0.350033in}}{\pgfqpoint{8.565566in}{0.347980in}}%
\pgfpathquadraticcurveto{\pgfqpoint{8.568322in}{0.345926in}}{\pgfqpoint{8.576661in}{0.344071in}}%
\pgfpathlineto{\pgfqpoint{8.580210in}{0.343278in}}%
\pgfpathquadraticcurveto{\pgfqpoint{8.591233in}{0.340919in}}{\pgfqpoint{8.595880in}{0.336614in}}%
\pgfpathquadraticcurveto{\pgfqpoint{8.600527in}{0.332309in}}{\pgfqpoint{8.600527in}{0.324600in}}%
\pgfpathquadraticcurveto{\pgfqpoint{8.600527in}{0.315810in}}{\pgfqpoint{8.593575in}{0.310676in}}%
\pgfpathquadraticcurveto{\pgfqpoint{8.586622in}{0.305561in}}{\pgfqpoint{8.574464in}{0.305561in}}%
\pgfpathquadraticcurveto{\pgfqpoint{8.569402in}{0.305561in}}{\pgfqpoint{8.563909in}{0.306552in}}%
\pgfpathquadraticcurveto{\pgfqpoint{8.558415in}{0.307542in}}{\pgfqpoint{8.552345in}{0.309506in}}%
\pgfpathlineto{\pgfqpoint{8.552345in}{0.320205in}}%
\pgfpathquadraticcurveto{\pgfqpoint{8.558091in}{0.317215in}}{\pgfqpoint{8.563656in}{0.315720in}}%
\pgfpathquadraticcurveto{\pgfqpoint{8.569222in}{0.314243in}}{\pgfqpoint{8.574698in}{0.314243in}}%
\pgfpathquadraticcurveto{\pgfqpoint{8.582011in}{0.314243in}}{\pgfqpoint{8.585938in}{0.316746in}}%
\pgfpathquadraticcurveto{\pgfqpoint{8.589882in}{0.319250in}}{\pgfqpoint{8.589882in}{0.323807in}}%
\pgfpathquadraticcurveto{\pgfqpoint{8.589882in}{0.328022in}}{\pgfqpoint{8.587036in}{0.330274in}}%
\pgfpathquadraticcurveto{\pgfqpoint{8.584208in}{0.332525in}}{\pgfqpoint{8.574572in}{0.334614in}}%
\pgfpathlineto{\pgfqpoint{8.570969in}{0.335461in}}%
\pgfpathquadraticcurveto{\pgfqpoint{8.561351in}{0.337478in}}{\pgfqpoint{8.557064in}{0.341675in}}%
\pgfpathquadraticcurveto{\pgfqpoint{8.552795in}{0.345872in}}{\pgfqpoint{8.552795in}{0.353185in}}%
\pgfpathquadraticcurveto{\pgfqpoint{8.552795in}{0.362083in}}{\pgfqpoint{8.559099in}{0.366910in}}%
\pgfpathquadraticcurveto{\pgfqpoint{8.565404in}{0.371756in}}{\pgfqpoint{8.577003in}{0.371756in}}%
\pgfpathquadraticcurveto{\pgfqpoint{8.582731in}{0.371756in}}{\pgfqpoint{8.587793in}{0.370909in}}%
\pgfpathquadraticcurveto{\pgfqpoint{8.592872in}{0.370080in}}{\pgfqpoint{8.597141in}{0.368387in}}%
\pgfpathlineto{\pgfqpoint{8.597141in}{0.368387in}}%
\pgfpathclose%
\pgfpathmoveto{\pgfqpoint{8.657086in}{0.316764in}}%
\pgfpathlineto{\pgfqpoint{8.675656in}{0.316764in}}%
\pgfpathlineto{\pgfqpoint{8.675656in}{0.380888in}}%
\pgfpathlineto{\pgfqpoint{8.655446in}{0.376835in}}%
\pgfpathlineto{\pgfqpoint{8.655446in}{0.387192in}}%
\pgfpathlineto{\pgfqpoint{8.675548in}{0.391245in}}%
\pgfpathlineto{\pgfqpoint{8.686914in}{0.391245in}}%
\pgfpathlineto{\pgfqpoint{8.686914in}{0.316764in}}%
\pgfpathlineto{\pgfqpoint{8.705484in}{0.316764in}}%
\pgfpathlineto{\pgfqpoint{8.705484in}{0.307200in}}%
\pgfpathlineto{\pgfqpoint{8.657086in}{0.307200in}}%
\pgfpathlineto{\pgfqpoint{8.657086in}{0.316764in}}%
\pgfpathlineto{\pgfqpoint{8.657086in}{0.316764in}}%
\pgfpathclose%
\pgfpathmoveto{\pgfqpoint{8.759701in}{0.381338in}}%
\pgfpathlineto{\pgfqpoint{8.730989in}{0.336470in}}%
\pgfpathlineto{\pgfqpoint{8.759701in}{0.336470in}}%
\pgfpathlineto{\pgfqpoint{8.759701in}{0.381338in}}%
\pgfpathlineto{\pgfqpoint{8.759701in}{0.381338in}}%
\pgfpathclose%
\pgfpathmoveto{\pgfqpoint{8.756711in}{0.391245in}}%
\pgfpathlineto{\pgfqpoint{8.771012in}{0.391245in}}%
\pgfpathlineto{\pgfqpoint{8.771012in}{0.336470in}}%
\pgfpathlineto{\pgfqpoint{8.783008in}{0.336470in}}%
\pgfpathlineto{\pgfqpoint{8.783008in}{0.327013in}}%
\pgfpathlineto{\pgfqpoint{8.771012in}{0.327013in}}%
\pgfpathlineto{\pgfqpoint{8.771012in}{0.307200in}}%
\pgfpathlineto{\pgfqpoint{8.759701in}{0.307200in}}%
\pgfpathlineto{\pgfqpoint{8.759701in}{0.327013in}}%
\pgfpathlineto{\pgfqpoint{8.721767in}{0.327013in}}%
\pgfpathlineto{\pgfqpoint{8.721767in}{0.337983in}}%
\pgfpathlineto{\pgfqpoint{8.756711in}{0.391245in}}%
\pgfpathlineto{\pgfqpoint{8.756711in}{0.391245in}}%
\pgfpathclose%
\pgfpathmoveto{\pgfqpoint{8.836252in}{0.352519in}}%
\pgfpathquadraticcurveto{\pgfqpoint{8.844412in}{0.350771in}}{\pgfqpoint{8.848987in}{0.345242in}}%
\pgfpathquadraticcurveto{\pgfqpoint{8.853580in}{0.339730in}}{\pgfqpoint{8.853580in}{0.331624in}}%
\pgfpathquadraticcurveto{\pgfqpoint{8.853580in}{0.319196in}}{\pgfqpoint{8.845024in}{0.312369in}}%
\pgfpathquadraticcurveto{\pgfqpoint{8.836469in}{0.305561in}}{\pgfqpoint{8.820708in}{0.305561in}}%
\pgfpathquadraticcurveto{\pgfqpoint{8.815430in}{0.305561in}}{\pgfqpoint{8.809829in}{0.306606in}}%
\pgfpathquadraticcurveto{\pgfqpoint{8.804227in}{0.307650in}}{\pgfqpoint{8.798265in}{0.309740in}}%
\pgfpathlineto{\pgfqpoint{8.798265in}{0.320709in}}%
\pgfpathquadraticcurveto{\pgfqpoint{8.802984in}{0.317953in}}{\pgfqpoint{8.808604in}{0.316548in}}%
\pgfpathquadraticcurveto{\pgfqpoint{8.814242in}{0.315143in}}{\pgfqpoint{8.820384in}{0.315143in}}%
\pgfpathquadraticcurveto{\pgfqpoint{8.831065in}{0.315143in}}{\pgfqpoint{8.836667in}{0.319358in}}%
\pgfpathquadraticcurveto{\pgfqpoint{8.842268in}{0.323573in}}{\pgfqpoint{8.842268in}{0.331624in}}%
\pgfpathquadraticcurveto{\pgfqpoint{8.842268in}{0.339063in}}{\pgfqpoint{8.837063in}{0.343242in}}%
\pgfpathquadraticcurveto{\pgfqpoint{8.831857in}{0.347439in}}{\pgfqpoint{8.822581in}{0.347439in}}%
\pgfpathlineto{\pgfqpoint{8.812783in}{0.347439in}}%
\pgfpathlineto{\pgfqpoint{8.812783in}{0.356787in}}%
\pgfpathlineto{\pgfqpoint{8.823031in}{0.356787in}}%
\pgfpathquadraticcurveto{\pgfqpoint{8.831407in}{0.356787in}}{\pgfqpoint{8.835856in}{0.360138in}}%
\pgfpathquadraticcurveto{\pgfqpoint{8.840305in}{0.363488in}}{\pgfqpoint{8.840305in}{0.369792in}}%
\pgfpathquadraticcurveto{\pgfqpoint{8.840305in}{0.376259in}}{\pgfqpoint{8.835712in}{0.379717in}}%
\pgfpathquadraticcurveto{\pgfqpoint{8.831137in}{0.383193in}}{\pgfqpoint{8.822581in}{0.383193in}}%
\pgfpathquadraticcurveto{\pgfqpoint{8.817898in}{0.383193in}}{\pgfqpoint{8.812548in}{0.382167in}}%
\pgfpathquadraticcurveto{\pgfqpoint{8.807199in}{0.381158in}}{\pgfqpoint{8.800786in}{0.379032in}}%
\pgfpathlineto{\pgfqpoint{8.800786in}{0.389155in}}%
\pgfpathquadraticcurveto{\pgfqpoint{8.807271in}{0.390957in}}{\pgfqpoint{8.812927in}{0.391857in}}%
\pgfpathquadraticcurveto{\pgfqpoint{8.818582in}{0.392758in}}{\pgfqpoint{8.823590in}{0.392758in}}%
\pgfpathquadraticcurveto{\pgfqpoint{8.836541in}{0.392758in}}{\pgfqpoint{8.844070in}{0.386868in}}%
\pgfpathquadraticcurveto{\pgfqpoint{8.851617in}{0.380996in}}{\pgfqpoint{8.851617in}{0.370981in}}%
\pgfpathquadraticcurveto{\pgfqpoint{8.851617in}{0.363992in}}{\pgfqpoint{8.847618in}{0.359183in}}%
\pgfpathquadraticcurveto{\pgfqpoint{8.843619in}{0.354374in}}{\pgfqpoint{8.836252in}{0.352519in}}%
\pgfpathlineto{\pgfqpoint{8.836252in}{0.352519in}}%
\pgfpathclose%
\pgfpathmoveto{\pgfqpoint{8.955583in}{0.338667in}}%
\pgfpathquadraticcurveto{\pgfqpoint{8.955583in}{0.350087in}}{\pgfqpoint{8.950882in}{0.356589in}}%
\pgfpathquadraticcurveto{\pgfqpoint{8.946181in}{0.363092in}}{\pgfqpoint{8.937967in}{0.363092in}}%
\pgfpathquadraticcurveto{\pgfqpoint{8.929735in}{0.363092in}}{\pgfqpoint{8.925034in}{0.356589in}}%
\pgfpathquadraticcurveto{\pgfqpoint{8.920333in}{0.350087in}}{\pgfqpoint{8.920333in}{0.338667in}}%
\pgfpathquadraticcurveto{\pgfqpoint{8.920333in}{0.327248in}}{\pgfqpoint{8.925034in}{0.320745in}}%
\pgfpathquadraticcurveto{\pgfqpoint{8.929735in}{0.314243in}}{\pgfqpoint{8.937967in}{0.314243in}}%
\pgfpathquadraticcurveto{\pgfqpoint{8.946181in}{0.314243in}}{\pgfqpoint{8.950882in}{0.320745in}}%
\pgfpathquadraticcurveto{\pgfqpoint{8.955583in}{0.327248in}}{\pgfqpoint{8.955583in}{0.338667in}}%
\pgfpathlineto{\pgfqpoint{8.955583in}{0.338667in}}%
\pgfpathclose%
\pgfpathmoveto{\pgfqpoint{8.920333in}{0.360678in}}%
\pgfpathquadraticcurveto{\pgfqpoint{8.923611in}{0.366298in}}{\pgfqpoint{8.928583in}{0.369018in}}%
\pgfpathquadraticcurveto{\pgfqpoint{8.933572in}{0.371756in}}{\pgfqpoint{8.940489in}{0.371756in}}%
\pgfpathquadraticcurveto{\pgfqpoint{8.951980in}{0.371756in}}{\pgfqpoint{8.959149in}{0.362641in}}%
\pgfpathquadraticcurveto{\pgfqpoint{8.966336in}{0.353527in}}{\pgfqpoint{8.966336in}{0.338667in}}%
\pgfpathquadraticcurveto{\pgfqpoint{8.966336in}{0.323807in}}{\pgfqpoint{8.959149in}{0.314675in}}%
\pgfpathquadraticcurveto{\pgfqpoint{8.951980in}{0.305561in}}{\pgfqpoint{8.940489in}{0.305561in}}%
\pgfpathquadraticcurveto{\pgfqpoint{8.933572in}{0.305561in}}{\pgfqpoint{8.928583in}{0.308299in}}%
\pgfpathquadraticcurveto{\pgfqpoint{8.923611in}{0.311037in}}{\pgfqpoint{8.920333in}{0.316656in}}%
\pgfpathlineto{\pgfqpoint{8.920333in}{0.307200in}}%
\pgfpathlineto{\pgfqpoint{8.909922in}{0.307200in}}%
\pgfpathlineto{\pgfqpoint{8.909922in}{0.394793in}}%
\pgfpathlineto{\pgfqpoint{8.920333in}{0.394793in}}%
\pgfpathlineto{\pgfqpoint{8.920333in}{0.360678in}}%
\pgfpathlineto{\pgfqpoint{8.920333in}{0.360678in}}%
\pgfpathclose%
\pgfpathmoveto{\pgfqpoint{8.982439in}{0.332075in}}%
\pgfpathlineto{\pgfqpoint{8.982439in}{0.370243in}}%
\pgfpathlineto{\pgfqpoint{8.992796in}{0.370243in}}%
\pgfpathlineto{\pgfqpoint{8.992796in}{0.332471in}}%
\pgfpathquadraticcurveto{\pgfqpoint{8.992796in}{0.323519in}}{\pgfqpoint{8.996272in}{0.319034in}}%
\pgfpathquadraticcurveto{\pgfqpoint{8.999767in}{0.314567in}}{\pgfqpoint{9.006755in}{0.314567in}}%
\pgfpathquadraticcurveto{\pgfqpoint{9.015131in}{0.314567in}}{\pgfqpoint{9.019994in}{0.319917in}}%
\pgfpathquadraticcurveto{\pgfqpoint{9.024876in}{0.325266in}}{\pgfqpoint{9.024876in}{0.334506in}}%
\pgfpathlineto{\pgfqpoint{9.024876in}{0.370243in}}%
\pgfpathlineto{\pgfqpoint{9.035233in}{0.370243in}}%
\pgfpathlineto{\pgfqpoint{9.035233in}{0.307200in}}%
\pgfpathlineto{\pgfqpoint{9.024876in}{0.307200in}}%
\pgfpathlineto{\pgfqpoint{9.024876in}{0.316891in}}%
\pgfpathquadraticcurveto{\pgfqpoint{9.021111in}{0.311145in}}{\pgfqpoint{9.016122in}{0.308353in}}%
\pgfpathquadraticcurveto{\pgfqpoint{9.011150in}{0.305561in}}{\pgfqpoint{9.004558in}{0.305561in}}%
\pgfpathquadraticcurveto{\pgfqpoint{8.993697in}{0.305561in}}{\pgfqpoint{8.988059in}{0.312315in}}%
\pgfpathquadraticcurveto{\pgfqpoint{8.982439in}{0.319070in}}{\pgfqpoint{8.982439in}{0.332075in}}%
\pgfpathlineto{\pgfqpoint{8.982439in}{0.332075in}}%
\pgfpathclose%
\pgfpathmoveto{\pgfqpoint{9.008503in}{0.371756in}}%
\pgfpathlineto{\pgfqpoint{9.008503in}{0.371756in}}%
\pgfpathlineto{\pgfqpoint{9.008503in}{0.371756in}}%
\pgfpathclose%
\pgfpathmoveto{\pgfqpoint{9.056559in}{0.370243in}}%
\pgfpathlineto{\pgfqpoint{9.066916in}{0.370243in}}%
\pgfpathlineto{\pgfqpoint{9.066916in}{0.307200in}}%
\pgfpathlineto{\pgfqpoint{9.056559in}{0.307200in}}%
\pgfpathlineto{\pgfqpoint{9.056559in}{0.370243in}}%
\pgfpathlineto{\pgfqpoint{9.056559in}{0.370243in}}%
\pgfpathclose%
\pgfpathmoveto{\pgfqpoint{9.056559in}{0.394793in}}%
\pgfpathlineto{\pgfqpoint{9.066916in}{0.394793in}}%
\pgfpathlineto{\pgfqpoint{9.066916in}{0.381662in}}%
\pgfpathlineto{\pgfqpoint{9.056559in}{0.381662in}}%
\pgfpathlineto{\pgfqpoint{9.056559in}{0.394793in}}%
\pgfpathlineto{\pgfqpoint{9.056559in}{0.394793in}}%
\pgfpathclose%
\pgfpathmoveto{\pgfqpoint{9.088585in}{0.394793in}}%
\pgfpathlineto{\pgfqpoint{9.098942in}{0.394793in}}%
\pgfpathlineto{\pgfqpoint{9.098942in}{0.307200in}}%
\pgfpathlineto{\pgfqpoint{9.088585in}{0.307200in}}%
\pgfpathlineto{\pgfqpoint{9.088585in}{0.394793in}}%
\pgfpathlineto{\pgfqpoint{9.088585in}{0.394793in}}%
\pgfpathclose%
\pgfpathmoveto{\pgfqpoint{9.162092in}{0.360678in}}%
\pgfpathlineto{\pgfqpoint{9.162092in}{0.394793in}}%
\pgfpathlineto{\pgfqpoint{9.172449in}{0.394793in}}%
\pgfpathlineto{\pgfqpoint{9.172449in}{0.307200in}}%
\pgfpathlineto{\pgfqpoint{9.162092in}{0.307200in}}%
\pgfpathlineto{\pgfqpoint{9.162092in}{0.316656in}}%
\pgfpathquadraticcurveto{\pgfqpoint{9.158832in}{0.311037in}}{\pgfqpoint{9.153843in}{0.308299in}}%
\pgfpathquadraticcurveto{\pgfqpoint{9.148871in}{0.305561in}}{\pgfqpoint{9.141883in}{0.305561in}}%
\pgfpathquadraticcurveto{\pgfqpoint{9.130463in}{0.305561in}}{\pgfqpoint{9.123276in}{0.314675in}}%
\pgfpathquadraticcurveto{\pgfqpoint{9.116107in}{0.323807in}}{\pgfqpoint{9.116107in}{0.338667in}}%
\pgfpathquadraticcurveto{\pgfqpoint{9.116107in}{0.353527in}}{\pgfqpoint{9.123276in}{0.362641in}}%
\pgfpathquadraticcurveto{\pgfqpoint{9.130463in}{0.371756in}}{\pgfqpoint{9.141883in}{0.371756in}}%
\pgfpathquadraticcurveto{\pgfqpoint{9.148871in}{0.371756in}}{\pgfqpoint{9.153843in}{0.369018in}}%
\pgfpathquadraticcurveto{\pgfqpoint{9.158832in}{0.366298in}}{\pgfqpoint{9.162092in}{0.360678in}}%
\pgfpathlineto{\pgfqpoint{9.162092in}{0.360678in}}%
\pgfpathclose%
\pgfpathmoveto{\pgfqpoint{9.126806in}{0.338667in}}%
\pgfpathquadraticcurveto{\pgfqpoint{9.126806in}{0.327248in}}{\pgfqpoint{9.131508in}{0.320745in}}%
\pgfpathquadraticcurveto{\pgfqpoint{9.136209in}{0.314243in}}{\pgfqpoint{9.144422in}{0.314243in}}%
\pgfpathquadraticcurveto{\pgfqpoint{9.152636in}{0.314243in}}{\pgfqpoint{9.157355in}{0.320745in}}%
\pgfpathquadraticcurveto{\pgfqpoint{9.162092in}{0.327248in}}{\pgfqpoint{9.162092in}{0.338667in}}%
\pgfpathquadraticcurveto{\pgfqpoint{9.162092in}{0.350087in}}{\pgfqpoint{9.157355in}{0.356589in}}%
\pgfpathquadraticcurveto{\pgfqpoint{9.152636in}{0.363092in}}{\pgfqpoint{9.144422in}{0.363092in}}%
\pgfpathquadraticcurveto{\pgfqpoint{9.136209in}{0.363092in}}{\pgfqpoint{9.131508in}{0.356589in}}%
\pgfpathquadraticcurveto{\pgfqpoint{9.126806in}{0.350087in}}{\pgfqpoint{9.126806in}{0.338667in}}%
\pgfpathlineto{\pgfqpoint{9.126806in}{0.338667in}}%
\pgfpathclose%
\pgfpathmoveto{\pgfqpoint{9.247722in}{0.341315in}}%
\pgfpathlineto{\pgfqpoint{9.247722in}{0.336254in}}%
\pgfpathlineto{\pgfqpoint{9.200098in}{0.336254in}}%
\pgfpathquadraticcurveto{\pgfqpoint{9.200782in}{0.325554in}}{\pgfqpoint{9.206546in}{0.319953in}}%
\pgfpathquadraticcurveto{\pgfqpoint{9.212310in}{0.314351in}}{\pgfqpoint{9.222613in}{0.314351in}}%
\pgfpathquadraticcurveto{\pgfqpoint{9.228575in}{0.314351in}}{\pgfqpoint{9.234177in}{0.315810in}}%
\pgfpathquadraticcurveto{\pgfqpoint{9.239779in}{0.317269in}}{\pgfqpoint{9.245308in}{0.320205in}}%
\pgfpathlineto{\pgfqpoint{9.245308in}{0.310406in}}%
\pgfpathquadraticcurveto{\pgfqpoint{9.239725in}{0.308047in}}{\pgfqpoint{9.233871in}{0.306804in}}%
\pgfpathquadraticcurveto{\pgfqpoint{9.228017in}{0.305561in}}{\pgfqpoint{9.222001in}{0.305561in}}%
\pgfpathquadraticcurveto{\pgfqpoint{9.206906in}{0.305561in}}{\pgfqpoint{9.198099in}{0.314333in}}%
\pgfpathquadraticcurveto{\pgfqpoint{9.189291in}{0.323123in}}{\pgfqpoint{9.189291in}{0.338109in}}%
\pgfpathquadraticcurveto{\pgfqpoint{9.189291in}{0.353581in}}{\pgfqpoint{9.197648in}{0.362659in}}%
\pgfpathquadraticcurveto{\pgfqpoint{9.206006in}{0.371756in}}{\pgfqpoint{9.220199in}{0.371756in}}%
\pgfpathquadraticcurveto{\pgfqpoint{9.232916in}{0.371756in}}{\pgfqpoint{9.240319in}{0.363560in}}%
\pgfpathquadraticcurveto{\pgfqpoint{9.247722in}{0.355383in}}{\pgfqpoint{9.247722in}{0.341315in}}%
\pgfpathlineto{\pgfqpoint{9.247722in}{0.341315in}}%
\pgfpathclose%
\pgfpathmoveto{\pgfqpoint{9.237365in}{0.344359in}}%
\pgfpathquadraticcurveto{\pgfqpoint{9.237257in}{0.352843in}}{\pgfqpoint{9.232610in}{0.357904in}}%
\pgfpathquadraticcurveto{\pgfqpoint{9.227963in}{0.362984in}}{\pgfqpoint{9.220308in}{0.362984in}}%
\pgfpathquadraticcurveto{\pgfqpoint{9.211644in}{0.362984in}}{\pgfqpoint{9.206438in}{0.358084in}}%
\pgfpathquadraticcurveto{\pgfqpoint{9.201233in}{0.353185in}}{\pgfqpoint{9.200440in}{0.344287in}}%
\pgfpathlineto{\pgfqpoint{9.237365in}{0.344359in}}%
\pgfpathlineto{\pgfqpoint{9.237365in}{0.344359in}}%
\pgfpathclose%
\pgfpathmoveto{\pgfqpoint{9.301254in}{0.360570in}}%
\pgfpathquadraticcurveto{\pgfqpoint{9.299507in}{0.361579in}}{\pgfqpoint{9.297454in}{0.362047in}}%
\pgfpathquadraticcurveto{\pgfqpoint{9.295400in}{0.362533in}}{\pgfqpoint{9.292933in}{0.362533in}}%
\pgfpathquadraticcurveto{\pgfqpoint{9.284143in}{0.362533in}}{\pgfqpoint{9.279441in}{0.356823in}}%
\pgfpathquadraticcurveto{\pgfqpoint{9.274740in}{0.351114in}}{\pgfqpoint{9.274740in}{0.340414in}}%
\pgfpathlineto{\pgfqpoint{9.274740in}{0.307200in}}%
\pgfpathlineto{\pgfqpoint{9.264329in}{0.307200in}}%
\pgfpathlineto{\pgfqpoint{9.264329in}{0.370243in}}%
\pgfpathlineto{\pgfqpoint{9.274740in}{0.370243in}}%
\pgfpathlineto{\pgfqpoint{9.274740in}{0.360444in}}%
\pgfpathquadraticcurveto{\pgfqpoint{9.278018in}{0.366190in}}{\pgfqpoint{9.283242in}{0.368964in}}%
\pgfpathquadraticcurveto{\pgfqpoint{9.288484in}{0.371756in}}{\pgfqpoint{9.295977in}{0.371756in}}%
\pgfpathquadraticcurveto{\pgfqpoint{9.297039in}{0.371756in}}{\pgfqpoint{9.298336in}{0.371611in}}%
\pgfpathquadraticcurveto{\pgfqpoint{9.299633in}{0.371485in}}{\pgfqpoint{9.301200in}{0.371197in}}%
\pgfpathlineto{\pgfqpoint{9.301254in}{0.360570in}}%
\pgfpathlineto{\pgfqpoint{9.301254in}{0.360570in}}%
\pgfpathclose%
\pgfpathmoveto{\pgfqpoint{9.352301in}{0.368387in}}%
\pgfpathlineto{\pgfqpoint{9.352301in}{0.358589in}}%
\pgfpathquadraticcurveto{\pgfqpoint{9.347924in}{0.360840in}}{\pgfqpoint{9.343186in}{0.361957in}}%
\pgfpathquadraticcurveto{\pgfqpoint{9.338467in}{0.363092in}}{\pgfqpoint{9.333388in}{0.363092in}}%
\pgfpathquadraticcurveto{\pgfqpoint{9.325679in}{0.363092in}}{\pgfqpoint{9.321824in}{0.360732in}}%
\pgfpathquadraticcurveto{\pgfqpoint{9.317969in}{0.358373in}}{\pgfqpoint{9.317969in}{0.353635in}}%
\pgfpathquadraticcurveto{\pgfqpoint{9.317969in}{0.350033in}}{\pgfqpoint{9.320725in}{0.347980in}}%
\pgfpathquadraticcurveto{\pgfqpoint{9.323481in}{0.345926in}}{\pgfqpoint{9.331821in}{0.344071in}}%
\pgfpathlineto{\pgfqpoint{9.335369in}{0.343278in}}%
\pgfpathquadraticcurveto{\pgfqpoint{9.346393in}{0.340919in}}{\pgfqpoint{9.351040in}{0.336614in}}%
\pgfpathquadraticcurveto{\pgfqpoint{9.355687in}{0.332309in}}{\pgfqpoint{9.355687in}{0.324600in}}%
\pgfpathquadraticcurveto{\pgfqpoint{9.355687in}{0.315810in}}{\pgfqpoint{9.348734in}{0.310676in}}%
\pgfpathquadraticcurveto{\pgfqpoint{9.341781in}{0.305561in}}{\pgfqpoint{9.329623in}{0.305561in}}%
\pgfpathquadraticcurveto{\pgfqpoint{9.324562in}{0.305561in}}{\pgfqpoint{9.319068in}{0.306552in}}%
\pgfpathquadraticcurveto{\pgfqpoint{9.313574in}{0.307542in}}{\pgfqpoint{9.307504in}{0.309506in}}%
\pgfpathlineto{\pgfqpoint{9.307504in}{0.320205in}}%
\pgfpathquadraticcurveto{\pgfqpoint{9.313250in}{0.317215in}}{\pgfqpoint{9.318816in}{0.315720in}}%
\pgfpathquadraticcurveto{\pgfqpoint{9.324382in}{0.314243in}}{\pgfqpoint{9.329857in}{0.314243in}}%
\pgfpathquadraticcurveto{\pgfqpoint{9.337170in}{0.314243in}}{\pgfqpoint{9.341097in}{0.316746in}}%
\pgfpathquadraticcurveto{\pgfqpoint{9.345042in}{0.319250in}}{\pgfqpoint{9.345042in}{0.323807in}}%
\pgfpathquadraticcurveto{\pgfqpoint{9.345042in}{0.328022in}}{\pgfqpoint{9.342196in}{0.330274in}}%
\pgfpathquadraticcurveto{\pgfqpoint{9.339368in}{0.332525in}}{\pgfqpoint{9.329731in}{0.334614in}}%
\pgfpathlineto{\pgfqpoint{9.326129in}{0.335461in}}%
\pgfpathquadraticcurveto{\pgfqpoint{9.316510in}{0.337478in}}{\pgfqpoint{9.312224in}{0.341675in}}%
\pgfpathquadraticcurveto{\pgfqpoint{9.307955in}{0.345872in}}{\pgfqpoint{9.307955in}{0.353185in}}%
\pgfpathquadraticcurveto{\pgfqpoint{9.307955in}{0.362083in}}{\pgfqpoint{9.314259in}{0.366910in}}%
\pgfpathquadraticcurveto{\pgfqpoint{9.320563in}{0.371756in}}{\pgfqpoint{9.332163in}{0.371756in}}%
\pgfpathquadraticcurveto{\pgfqpoint{9.337891in}{0.371756in}}{\pgfqpoint{9.342952in}{0.370909in}}%
\pgfpathquadraticcurveto{\pgfqpoint{9.348032in}{0.370080in}}{\pgfqpoint{9.352301in}{0.368387in}}%
\pgfpathlineto{\pgfqpoint{9.352301in}{0.368387in}}%
\pgfpathclose%
\pgfpathmoveto{\pgfqpoint{9.450287in}{0.360678in}}%
\pgfpathlineto{\pgfqpoint{9.450287in}{0.394793in}}%
\pgfpathlineto{\pgfqpoint{9.460644in}{0.394793in}}%
\pgfpathlineto{\pgfqpoint{9.460644in}{0.307200in}}%
\pgfpathlineto{\pgfqpoint{9.450287in}{0.307200in}}%
\pgfpathlineto{\pgfqpoint{9.450287in}{0.316656in}}%
\pgfpathquadraticcurveto{\pgfqpoint{9.447026in}{0.311037in}}{\pgfqpoint{9.442037in}{0.308299in}}%
\pgfpathquadraticcurveto{\pgfqpoint{9.437066in}{0.305561in}}{\pgfqpoint{9.430077in}{0.305561in}}%
\pgfpathquadraticcurveto{\pgfqpoint{9.418657in}{0.305561in}}{\pgfqpoint{9.411470in}{0.314675in}}%
\pgfpathquadraticcurveto{\pgfqpoint{9.404302in}{0.323807in}}{\pgfqpoint{9.404302in}{0.338667in}}%
\pgfpathquadraticcurveto{\pgfqpoint{9.404302in}{0.353527in}}{\pgfqpoint{9.411470in}{0.362641in}}%
\pgfpathquadraticcurveto{\pgfqpoint{9.418657in}{0.371756in}}{\pgfqpoint{9.430077in}{0.371756in}}%
\pgfpathquadraticcurveto{\pgfqpoint{9.437066in}{0.371756in}}{\pgfqpoint{9.442037in}{0.369018in}}%
\pgfpathquadraticcurveto{\pgfqpoint{9.447026in}{0.366298in}}{\pgfqpoint{9.450287in}{0.360678in}}%
\pgfpathlineto{\pgfqpoint{9.450287in}{0.360678in}}%
\pgfpathclose%
\pgfpathmoveto{\pgfqpoint{9.415001in}{0.338667in}}%
\pgfpathquadraticcurveto{\pgfqpoint{9.415001in}{0.327248in}}{\pgfqpoint{9.419702in}{0.320745in}}%
\pgfpathquadraticcurveto{\pgfqpoint{9.424403in}{0.314243in}}{\pgfqpoint{9.432617in}{0.314243in}}%
\pgfpathquadraticcurveto{\pgfqpoint{9.440830in}{0.314243in}}{\pgfqpoint{9.445549in}{0.320745in}}%
\pgfpathquadraticcurveto{\pgfqpoint{9.450287in}{0.327248in}}{\pgfqpoint{9.450287in}{0.338667in}}%
\pgfpathquadraticcurveto{\pgfqpoint{9.450287in}{0.350087in}}{\pgfqpoint{9.445549in}{0.356589in}}%
\pgfpathquadraticcurveto{\pgfqpoint{9.440830in}{0.363092in}}{\pgfqpoint{9.432617in}{0.363092in}}%
\pgfpathquadraticcurveto{\pgfqpoint{9.424403in}{0.363092in}}{\pgfqpoint{9.419702in}{0.356589in}}%
\pgfpathquadraticcurveto{\pgfqpoint{9.415001in}{0.350087in}}{\pgfqpoint{9.415001in}{0.338667in}}%
\pgfpathlineto{\pgfqpoint{9.415001in}{0.338667in}}%
\pgfpathclose%
\pgfpathmoveto{\pgfqpoint{9.535916in}{0.341315in}}%
\pgfpathlineto{\pgfqpoint{9.535916in}{0.336254in}}%
\pgfpathlineto{\pgfqpoint{9.488292in}{0.336254in}}%
\pgfpathquadraticcurveto{\pgfqpoint{9.488977in}{0.325554in}}{\pgfqpoint{9.494741in}{0.319953in}}%
\pgfpathquadraticcurveto{\pgfqpoint{9.500505in}{0.314351in}}{\pgfqpoint{9.510808in}{0.314351in}}%
\pgfpathquadraticcurveto{\pgfqpoint{9.516770in}{0.314351in}}{\pgfqpoint{9.522371in}{0.315810in}}%
\pgfpathquadraticcurveto{\pgfqpoint{9.527973in}{0.317269in}}{\pgfqpoint{9.533503in}{0.320205in}}%
\pgfpathlineto{\pgfqpoint{9.533503in}{0.310406in}}%
\pgfpathquadraticcurveto{\pgfqpoint{9.527919in}{0.308047in}}{\pgfqpoint{9.522065in}{0.306804in}}%
\pgfpathquadraticcurveto{\pgfqpoint{9.516211in}{0.305561in}}{\pgfqpoint{9.510195in}{0.305561in}}%
\pgfpathquadraticcurveto{\pgfqpoint{9.495101in}{0.305561in}}{\pgfqpoint{9.486293in}{0.314333in}}%
\pgfpathquadraticcurveto{\pgfqpoint{9.477485in}{0.323123in}}{\pgfqpoint{9.477485in}{0.338109in}}%
\pgfpathquadraticcurveto{\pgfqpoint{9.477485in}{0.353581in}}{\pgfqpoint{9.485843in}{0.362659in}}%
\pgfpathquadraticcurveto{\pgfqpoint{9.494200in}{0.371756in}}{\pgfqpoint{9.508394in}{0.371756in}}%
\pgfpathquadraticcurveto{\pgfqpoint{9.521110in}{0.371756in}}{\pgfqpoint{9.528513in}{0.363560in}}%
\pgfpathquadraticcurveto{\pgfqpoint{9.535916in}{0.355383in}}{\pgfqpoint{9.535916in}{0.341315in}}%
\pgfpathlineto{\pgfqpoint{9.535916in}{0.341315in}}%
\pgfpathclose%
\pgfpathmoveto{\pgfqpoint{9.525559in}{0.344359in}}%
\pgfpathquadraticcurveto{\pgfqpoint{9.525451in}{0.352843in}}{\pgfqpoint{9.520804in}{0.357904in}}%
\pgfpathquadraticcurveto{\pgfqpoint{9.516157in}{0.362984in}}{\pgfqpoint{9.508502in}{0.362984in}}%
\pgfpathquadraticcurveto{\pgfqpoint{9.499838in}{0.362984in}}{\pgfqpoint{9.494633in}{0.358084in}}%
\pgfpathquadraticcurveto{\pgfqpoint{9.489427in}{0.353185in}}{\pgfqpoint{9.488635in}{0.344287in}}%
\pgfpathlineto{\pgfqpoint{9.525559in}{0.344359in}}%
\pgfpathlineto{\pgfqpoint{9.525559in}{0.344359in}}%
\pgfpathclose%
\pgfpathmoveto{\pgfqpoint{9.593105in}{0.368387in}}%
\pgfpathlineto{\pgfqpoint{9.593105in}{0.358589in}}%
\pgfpathquadraticcurveto{\pgfqpoint{9.588728in}{0.360840in}}{\pgfqpoint{9.583991in}{0.361957in}}%
\pgfpathquadraticcurveto{\pgfqpoint{9.579272in}{0.363092in}}{\pgfqpoint{9.574192in}{0.363092in}}%
\pgfpathquadraticcurveto{\pgfqpoint{9.566483in}{0.363092in}}{\pgfqpoint{9.562628in}{0.360732in}}%
\pgfpathquadraticcurveto{\pgfqpoint{9.558774in}{0.358373in}}{\pgfqpoint{9.558774in}{0.353635in}}%
\pgfpathquadraticcurveto{\pgfqpoint{9.558774in}{0.350033in}}{\pgfqpoint{9.561530in}{0.347980in}}%
\pgfpathquadraticcurveto{\pgfqpoint{9.564286in}{0.345926in}}{\pgfqpoint{9.572625in}{0.344071in}}%
\pgfpathlineto{\pgfqpoint{9.576174in}{0.343278in}}%
\pgfpathquadraticcurveto{\pgfqpoint{9.587197in}{0.340919in}}{\pgfqpoint{9.591844in}{0.336614in}}%
\pgfpathquadraticcurveto{\pgfqpoint{9.596491in}{0.332309in}}{\pgfqpoint{9.596491in}{0.324600in}}%
\pgfpathquadraticcurveto{\pgfqpoint{9.596491in}{0.315810in}}{\pgfqpoint{9.589539in}{0.310676in}}%
\pgfpathquadraticcurveto{\pgfqpoint{9.582586in}{0.305561in}}{\pgfqpoint{9.570428in}{0.305561in}}%
\pgfpathquadraticcurveto{\pgfqpoint{9.565366in}{0.305561in}}{\pgfqpoint{9.559873in}{0.306552in}}%
\pgfpathquadraticcurveto{\pgfqpoint{9.554379in}{0.307542in}}{\pgfqpoint{9.548309in}{0.309506in}}%
\pgfpathlineto{\pgfqpoint{9.548309in}{0.320205in}}%
\pgfpathquadraticcurveto{\pgfqpoint{9.554055in}{0.317215in}}{\pgfqpoint{9.559620in}{0.315720in}}%
\pgfpathquadraticcurveto{\pgfqpoint{9.565186in}{0.314243in}}{\pgfqpoint{9.570662in}{0.314243in}}%
\pgfpathquadraticcurveto{\pgfqpoint{9.577975in}{0.314243in}}{\pgfqpoint{9.581901in}{0.316746in}}%
\pgfpathquadraticcurveto{\pgfqpoint{9.585846in}{0.319250in}}{\pgfqpoint{9.585846in}{0.323807in}}%
\pgfpathquadraticcurveto{\pgfqpoint{9.585846in}{0.328022in}}{\pgfqpoint{9.583000in}{0.330274in}}%
\pgfpathquadraticcurveto{\pgfqpoint{9.580172in}{0.332525in}}{\pgfqpoint{9.570536in}{0.334614in}}%
\pgfpathlineto{\pgfqpoint{9.566933in}{0.335461in}}%
\pgfpathquadraticcurveto{\pgfqpoint{9.557315in}{0.337478in}}{\pgfqpoint{9.553028in}{0.341675in}}%
\pgfpathquadraticcurveto{\pgfqpoint{9.548759in}{0.345872in}}{\pgfqpoint{9.548759in}{0.353185in}}%
\pgfpathquadraticcurveto{\pgfqpoint{9.548759in}{0.362083in}}{\pgfqpoint{9.555063in}{0.366910in}}%
\pgfpathquadraticcurveto{\pgfqpoint{9.561368in}{0.371756in}}{\pgfqpoint{9.572967in}{0.371756in}}%
\pgfpathquadraticcurveto{\pgfqpoint{9.578695in}{0.371756in}}{\pgfqpoint{9.583757in}{0.370909in}}%
\pgfpathquadraticcurveto{\pgfqpoint{9.588836in}{0.370080in}}{\pgfqpoint{9.593105in}{0.368387in}}%
\pgfpathlineto{\pgfqpoint{9.593105in}{0.368387in}}%
\pgfpathclose%
\pgfpathmoveto{\pgfqpoint{9.622987in}{0.316656in}}%
\pgfpathlineto{\pgfqpoint{9.622987in}{0.283226in}}%
\pgfpathlineto{\pgfqpoint{9.612576in}{0.283226in}}%
\pgfpathlineto{\pgfqpoint{9.612576in}{0.370243in}}%
\pgfpathlineto{\pgfqpoint{9.622987in}{0.370243in}}%
\pgfpathlineto{\pgfqpoint{9.622987in}{0.360678in}}%
\pgfpathquadraticcurveto{\pgfqpoint{9.626265in}{0.366298in}}{\pgfqpoint{9.631237in}{0.369018in}}%
\pgfpathquadraticcurveto{\pgfqpoint{9.636226in}{0.371756in}}{\pgfqpoint{9.643143in}{0.371756in}}%
\pgfpathquadraticcurveto{\pgfqpoint{9.654635in}{0.371756in}}{\pgfqpoint{9.661803in}{0.362641in}}%
\pgfpathquadraticcurveto{\pgfqpoint{9.668990in}{0.353527in}}{\pgfqpoint{9.668990in}{0.338667in}}%
\pgfpathquadraticcurveto{\pgfqpoint{9.668990in}{0.323807in}}{\pgfqpoint{9.661803in}{0.314675in}}%
\pgfpathquadraticcurveto{\pgfqpoint{9.654635in}{0.305561in}}{\pgfqpoint{9.643143in}{0.305561in}}%
\pgfpathquadraticcurveto{\pgfqpoint{9.636226in}{0.305561in}}{\pgfqpoint{9.631237in}{0.308299in}}%
\pgfpathquadraticcurveto{\pgfqpoint{9.626265in}{0.311037in}}{\pgfqpoint{9.622987in}{0.316656in}}%
\pgfpathlineto{\pgfqpoint{9.622987in}{0.316656in}}%
\pgfpathclose%
\pgfpathmoveto{\pgfqpoint{9.658237in}{0.338667in}}%
\pgfpathquadraticcurveto{\pgfqpoint{9.658237in}{0.350087in}}{\pgfqpoint{9.653536in}{0.356589in}}%
\pgfpathquadraticcurveto{\pgfqpoint{9.648835in}{0.363092in}}{\pgfqpoint{9.640621in}{0.363092in}}%
\pgfpathquadraticcurveto{\pgfqpoint{9.632390in}{0.363092in}}{\pgfqpoint{9.627688in}{0.356589in}}%
\pgfpathquadraticcurveto{\pgfqpoint{9.622987in}{0.350087in}}{\pgfqpoint{9.622987in}{0.338667in}}%
\pgfpathquadraticcurveto{\pgfqpoint{9.622987in}{0.327248in}}{\pgfqpoint{9.627688in}{0.320745in}}%
\pgfpathquadraticcurveto{\pgfqpoint{9.632390in}{0.314243in}}{\pgfqpoint{9.640621in}{0.314243in}}%
\pgfpathquadraticcurveto{\pgfqpoint{9.648835in}{0.314243in}}{\pgfqpoint{9.653536in}{0.320745in}}%
\pgfpathquadraticcurveto{\pgfqpoint{9.658237in}{0.327248in}}{\pgfqpoint{9.658237in}{0.338667in}}%
\pgfpathlineto{\pgfqpoint{9.658237in}{0.338667in}}%
\pgfpathclose%
\pgfpathmoveto{\pgfqpoint{9.686156in}{0.370243in}}%
\pgfpathlineto{\pgfqpoint{9.696513in}{0.370243in}}%
\pgfpathlineto{\pgfqpoint{9.696513in}{0.307200in}}%
\pgfpathlineto{\pgfqpoint{9.686156in}{0.307200in}}%
\pgfpathlineto{\pgfqpoint{9.686156in}{0.370243in}}%
\pgfpathlineto{\pgfqpoint{9.686156in}{0.370243in}}%
\pgfpathclose%
\pgfpathmoveto{\pgfqpoint{9.686156in}{0.394793in}}%
\pgfpathlineto{\pgfqpoint{9.696513in}{0.394793in}}%
\pgfpathlineto{\pgfqpoint{9.696513in}{0.381662in}}%
\pgfpathlineto{\pgfqpoint{9.686156in}{0.381662in}}%
\pgfpathlineto{\pgfqpoint{9.686156in}{0.394793in}}%
\pgfpathlineto{\pgfqpoint{9.686156in}{0.394793in}}%
\pgfpathclose%
\pgfpathmoveto{\pgfqpoint{9.728430in}{0.388147in}}%
\pgfpathlineto{\pgfqpoint{9.728430in}{0.370243in}}%
\pgfpathlineto{\pgfqpoint{9.749757in}{0.370243in}}%
\pgfpathlineto{\pgfqpoint{9.749757in}{0.362191in}}%
\pgfpathlineto{\pgfqpoint{9.728430in}{0.362191in}}%
\pgfpathlineto{\pgfqpoint{9.728430in}{0.327968in}}%
\pgfpathquadraticcurveto{\pgfqpoint{9.728430in}{0.320259in}}{\pgfqpoint{9.730538in}{0.318061in}}%
\pgfpathquadraticcurveto{\pgfqpoint{9.732645in}{0.315864in}}{\pgfqpoint{9.739130in}{0.315864in}}%
\pgfpathlineto{\pgfqpoint{9.749757in}{0.315864in}}%
\pgfpathlineto{\pgfqpoint{9.749757in}{0.307200in}}%
\pgfpathlineto{\pgfqpoint{9.739130in}{0.307200in}}%
\pgfpathquadraticcurveto{\pgfqpoint{9.727133in}{0.307200in}}{\pgfqpoint{9.722576in}{0.311667in}}%
\pgfpathquadraticcurveto{\pgfqpoint{9.718019in}{0.316152in}}{\pgfqpoint{9.718019in}{0.327968in}}%
\pgfpathlineto{\pgfqpoint{9.718019in}{0.362191in}}%
\pgfpathlineto{\pgfqpoint{9.710418in}{0.362191in}}%
\pgfpathlineto{\pgfqpoint{9.710418in}{0.370243in}}%
\pgfpathlineto{\pgfqpoint{9.718019in}{0.370243in}}%
\pgfpathlineto{\pgfqpoint{9.718019in}{0.388147in}}%
\pgfpathlineto{\pgfqpoint{9.728430in}{0.388147in}}%
\pgfpathlineto{\pgfqpoint{9.728430in}{0.388147in}}%
\pgfpathclose%
\pgfpathmoveto{\pgfqpoint{9.817302in}{0.341315in}}%
\pgfpathlineto{\pgfqpoint{9.817302in}{0.336254in}}%
\pgfpathlineto{\pgfqpoint{9.769678in}{0.336254in}}%
\pgfpathquadraticcurveto{\pgfqpoint{9.770363in}{0.325554in}}{\pgfqpoint{9.776127in}{0.319953in}}%
\pgfpathquadraticcurveto{\pgfqpoint{9.781890in}{0.314351in}}{\pgfqpoint{9.792193in}{0.314351in}}%
\pgfpathquadraticcurveto{\pgfqpoint{9.798155in}{0.314351in}}{\pgfqpoint{9.803757in}{0.315810in}}%
\pgfpathquadraticcurveto{\pgfqpoint{9.809359in}{0.317269in}}{\pgfqpoint{9.814889in}{0.320205in}}%
\pgfpathlineto{\pgfqpoint{9.814889in}{0.310406in}}%
\pgfpathquadraticcurveto{\pgfqpoint{9.809305in}{0.308047in}}{\pgfqpoint{9.803451in}{0.306804in}}%
\pgfpathquadraticcurveto{\pgfqpoint{9.797597in}{0.305561in}}{\pgfqpoint{9.791581in}{0.305561in}}%
\pgfpathquadraticcurveto{\pgfqpoint{9.776487in}{0.305561in}}{\pgfqpoint{9.767679in}{0.314333in}}%
\pgfpathquadraticcurveto{\pgfqpoint{9.758871in}{0.323123in}}{\pgfqpoint{9.758871in}{0.338109in}}%
\pgfpathquadraticcurveto{\pgfqpoint{9.758871in}{0.353581in}}{\pgfqpoint{9.767229in}{0.362659in}}%
\pgfpathquadraticcurveto{\pgfqpoint{9.775586in}{0.371756in}}{\pgfqpoint{9.789780in}{0.371756in}}%
\pgfpathquadraticcurveto{\pgfqpoint{9.802496in}{0.371756in}}{\pgfqpoint{9.809899in}{0.363560in}}%
\pgfpathquadraticcurveto{\pgfqpoint{9.817302in}{0.355383in}}{\pgfqpoint{9.817302in}{0.341315in}}%
\pgfpathlineto{\pgfqpoint{9.817302in}{0.341315in}}%
\pgfpathclose%
\pgfpathmoveto{\pgfqpoint{9.806945in}{0.344359in}}%
\pgfpathquadraticcurveto{\pgfqpoint{9.806837in}{0.352843in}}{\pgfqpoint{9.802190in}{0.357904in}}%
\pgfpathquadraticcurveto{\pgfqpoint{9.797543in}{0.362984in}}{\pgfqpoint{9.789888in}{0.362984in}}%
\pgfpathquadraticcurveto{\pgfqpoint{9.781224in}{0.362984in}}{\pgfqpoint{9.776018in}{0.358084in}}%
\pgfpathquadraticcurveto{\pgfqpoint{9.770813in}{0.353185in}}{\pgfqpoint{9.770020in}{0.344287in}}%
\pgfpathlineto{\pgfqpoint{9.806945in}{0.344359in}}%
\pgfpathlineto{\pgfqpoint{9.806945in}{0.344359in}}%
\pgfpathclose%
\pgfpathmoveto{\pgfqpoint{9.895367in}{0.362984in}}%
\pgfpathquadraticcurveto{\pgfqpoint{9.887045in}{0.362984in}}{\pgfqpoint{9.882200in}{0.356481in}}%
\pgfpathquadraticcurveto{\pgfqpoint{9.877355in}{0.349979in}}{\pgfqpoint{9.877355in}{0.338667in}}%
\pgfpathquadraticcurveto{\pgfqpoint{9.877355in}{0.327356in}}{\pgfqpoint{9.882164in}{0.320853in}}%
\pgfpathquadraticcurveto{\pgfqpoint{9.886991in}{0.314351in}}{\pgfqpoint{9.895367in}{0.314351in}}%
\pgfpathquadraticcurveto{\pgfqpoint{9.903653in}{0.314351in}}{\pgfqpoint{9.908480in}{0.320871in}}%
\pgfpathquadraticcurveto{\pgfqpoint{9.913325in}{0.327410in}}{\pgfqpoint{9.913325in}{0.338667in}}%
\pgfpathquadraticcurveto{\pgfqpoint{9.913325in}{0.349871in}}{\pgfqpoint{9.908480in}{0.356427in}}%
\pgfpathquadraticcurveto{\pgfqpoint{9.903653in}{0.362984in}}{\pgfqpoint{9.895367in}{0.362984in}}%
\pgfpathlineto{\pgfqpoint{9.895367in}{0.362984in}}%
\pgfpathclose%
\pgfpathmoveto{\pgfqpoint{9.895367in}{0.371756in}}%
\pgfpathquadraticcurveto{\pgfqpoint{9.908876in}{0.371756in}}{\pgfqpoint{9.916585in}{0.362966in}}%
\pgfpathquadraticcurveto{\pgfqpoint{9.924312in}{0.354194in}}{\pgfqpoint{9.924312in}{0.338667in}}%
\pgfpathquadraticcurveto{\pgfqpoint{9.924312in}{0.323195in}}{\pgfqpoint{9.916585in}{0.314369in}}%
\pgfpathquadraticcurveto{\pgfqpoint{9.908876in}{0.305561in}}{\pgfqpoint{9.895367in}{0.305561in}}%
\pgfpathquadraticcurveto{\pgfqpoint{9.881804in}{0.305561in}}{\pgfqpoint{9.874113in}{0.314369in}}%
\pgfpathquadraticcurveto{\pgfqpoint{9.866439in}{0.323195in}}{\pgfqpoint{9.866439in}{0.338667in}}%
\pgfpathquadraticcurveto{\pgfqpoint{9.866439in}{0.354194in}}{\pgfqpoint{9.874113in}{0.362966in}}%
\pgfpathquadraticcurveto{\pgfqpoint{9.881804in}{0.371756in}}{\pgfqpoint{9.895367in}{0.371756in}}%
\pgfpathlineto{\pgfqpoint{9.895367in}{0.371756in}}%
\pgfpathclose%
\pgfpathmoveto{\pgfqpoint{9.993893in}{0.345260in}}%
\pgfpathlineto{\pgfqpoint{9.993893in}{0.307200in}}%
\pgfpathlineto{\pgfqpoint{9.983536in}{0.307200in}}%
\pgfpathlineto{\pgfqpoint{9.983536in}{0.344917in}}%
\pgfpathquadraticcurveto{\pgfqpoint{9.983536in}{0.353869in}}{\pgfqpoint{9.980042in}{0.358300in}}%
\pgfpathquadraticcurveto{\pgfqpoint{9.976548in}{0.362749in}}{\pgfqpoint{9.969577in}{0.362749in}}%
\pgfpathquadraticcurveto{\pgfqpoint{9.961183in}{0.362749in}}{\pgfqpoint{9.956338in}{0.357400in}}%
\pgfpathquadraticcurveto{\pgfqpoint{9.951493in}{0.352068in}}{\pgfqpoint{9.951493in}{0.342828in}}%
\pgfpathlineto{\pgfqpoint{9.951493in}{0.307200in}}%
\pgfpathlineto{\pgfqpoint{9.941082in}{0.307200in}}%
\pgfpathlineto{\pgfqpoint{9.941082in}{0.370243in}}%
\pgfpathlineto{\pgfqpoint{9.951493in}{0.370243in}}%
\pgfpathlineto{\pgfqpoint{9.951493in}{0.360444in}}%
\pgfpathquadraticcurveto{\pgfqpoint{9.955221in}{0.366136in}}{\pgfqpoint{9.960247in}{0.368946in}}%
\pgfpathquadraticcurveto{\pgfqpoint{9.965290in}{0.371756in}}{\pgfqpoint{9.971883in}{0.371756in}}%
\pgfpathquadraticcurveto{\pgfqpoint{9.982744in}{0.371756in}}{\pgfqpoint{9.988310in}{0.365037in}}%
\pgfpathquadraticcurveto{\pgfqpoint{9.993893in}{0.358318in}}{\pgfqpoint{9.993893in}{0.345260in}}%
\pgfpathlineto{\pgfqpoint{9.993893in}{0.345260in}}%
\pgfpathclose%
\pgfpathmoveto{\pgfqpoint{10.014535in}{0.394793in}}%
\pgfpathlineto{\pgfqpoint{10.024892in}{0.394793in}}%
\pgfpathlineto{\pgfqpoint{10.024892in}{0.307200in}}%
\pgfpathlineto{\pgfqpoint{10.014535in}{0.307200in}}%
\pgfpathlineto{\pgfqpoint{10.014535in}{0.394793in}}%
\pgfpathlineto{\pgfqpoint{10.014535in}{0.394793in}}%
\pgfpathclose%
\pgfpathmoveto{\pgfqpoint{10.072787in}{0.301346in}}%
\pgfpathquadraticcurveto{\pgfqpoint{10.068410in}{0.290088in}}{\pgfqpoint{10.064231in}{0.286666in}}%
\pgfpathquadraticcurveto{\pgfqpoint{10.060070in}{0.283226in}}{\pgfqpoint{10.053099in}{0.283226in}}%
\pgfpathlineto{\pgfqpoint{10.044814in}{0.283226in}}%
\pgfpathlineto{\pgfqpoint{10.044814in}{0.291890in}}%
\pgfpathlineto{\pgfqpoint{10.050902in}{0.291890in}}%
\pgfpathquadraticcurveto{\pgfqpoint{10.055171in}{0.291890in}}{\pgfqpoint{10.057530in}{0.293925in}}%
\pgfpathquadraticcurveto{\pgfqpoint{10.059908in}{0.295942in}}{\pgfqpoint{10.062772in}{0.303489in}}%
\pgfpathlineto{\pgfqpoint{10.064627in}{0.308209in}}%
\pgfpathlineto{\pgfqpoint{10.039140in}{0.370243in}}%
\pgfpathlineto{\pgfqpoint{10.050109in}{0.370243in}}%
\pgfpathlineto{\pgfqpoint{10.069815in}{0.320943in}}%
\pgfpathlineto{\pgfqpoint{10.089520in}{0.370243in}}%
\pgfpathlineto{\pgfqpoint{10.100489in}{0.370243in}}%
\pgfpathlineto{\pgfqpoint{10.072787in}{0.301346in}}%
\pgfpathlineto{\pgfqpoint{10.072787in}{0.301346in}}%
\pgfpathclose%
\pgfpathmoveto{\pgfqpoint{10.161677in}{0.388147in}}%
\pgfpathlineto{\pgfqpoint{10.161677in}{0.370243in}}%
\pgfpathlineto{\pgfqpoint{10.183003in}{0.370243in}}%
\pgfpathlineto{\pgfqpoint{10.183003in}{0.362191in}}%
\pgfpathlineto{\pgfqpoint{10.161677in}{0.362191in}}%
\pgfpathlineto{\pgfqpoint{10.161677in}{0.327968in}}%
\pgfpathquadraticcurveto{\pgfqpoint{10.161677in}{0.320259in}}{\pgfqpoint{10.163784in}{0.318061in}}%
\pgfpathquadraticcurveto{\pgfqpoint{10.165891in}{0.315864in}}{\pgfqpoint{10.172376in}{0.315864in}}%
\pgfpathlineto{\pgfqpoint{10.183003in}{0.315864in}}%
\pgfpathlineto{\pgfqpoint{10.183003in}{0.307200in}}%
\pgfpathlineto{\pgfqpoint{10.172376in}{0.307200in}}%
\pgfpathquadraticcurveto{\pgfqpoint{10.160380in}{0.307200in}}{\pgfqpoint{10.155823in}{0.311667in}}%
\pgfpathquadraticcurveto{\pgfqpoint{10.151266in}{0.316152in}}{\pgfqpoint{10.151266in}{0.327968in}}%
\pgfpathlineto{\pgfqpoint{10.151266in}{0.362191in}}%
\pgfpathlineto{\pgfqpoint{10.143664in}{0.362191in}}%
\pgfpathlineto{\pgfqpoint{10.143664in}{0.370243in}}%
\pgfpathlineto{\pgfqpoint{10.151266in}{0.370243in}}%
\pgfpathlineto{\pgfqpoint{10.151266in}{0.388147in}}%
\pgfpathlineto{\pgfqpoint{10.161677in}{0.388147in}}%
\pgfpathlineto{\pgfqpoint{10.161677in}{0.388147in}}%
\pgfpathclose%
\pgfpathmoveto{\pgfqpoint{10.249036in}{0.345260in}}%
\pgfpathlineto{\pgfqpoint{10.249036in}{0.307200in}}%
\pgfpathlineto{\pgfqpoint{10.238679in}{0.307200in}}%
\pgfpathlineto{\pgfqpoint{10.238679in}{0.344917in}}%
\pgfpathquadraticcurveto{\pgfqpoint{10.238679in}{0.353869in}}{\pgfqpoint{10.235184in}{0.358300in}}%
\pgfpathquadraticcurveto{\pgfqpoint{10.231690in}{0.362749in}}{\pgfqpoint{10.224719in}{0.362749in}}%
\pgfpathquadraticcurveto{\pgfqpoint{10.216326in}{0.362749in}}{\pgfqpoint{10.211480in}{0.357400in}}%
\pgfpathquadraticcurveto{\pgfqpoint{10.206635in}{0.352068in}}{\pgfqpoint{10.206635in}{0.342828in}}%
\pgfpathlineto{\pgfqpoint{10.206635in}{0.307200in}}%
\pgfpathlineto{\pgfqpoint{10.196224in}{0.307200in}}%
\pgfpathlineto{\pgfqpoint{10.196224in}{0.394793in}}%
\pgfpathlineto{\pgfqpoint{10.206635in}{0.394793in}}%
\pgfpathlineto{\pgfqpoint{10.206635in}{0.360444in}}%
\pgfpathquadraticcurveto{\pgfqpoint{10.210363in}{0.366136in}}{\pgfqpoint{10.215389in}{0.368946in}}%
\pgfpathquadraticcurveto{\pgfqpoint{10.220432in}{0.371756in}}{\pgfqpoint{10.227025in}{0.371756in}}%
\pgfpathquadraticcurveto{\pgfqpoint{10.237886in}{0.371756in}}{\pgfqpoint{10.243452in}{0.365037in}}%
\pgfpathquadraticcurveto{\pgfqpoint{10.249036in}{0.358318in}}{\pgfqpoint{10.249036in}{0.345260in}}%
\pgfpathlineto{\pgfqpoint{10.249036in}{0.345260in}}%
\pgfpathclose%
\pgfpathmoveto{\pgfqpoint{10.306206in}{0.360570in}}%
\pgfpathquadraticcurveto{\pgfqpoint{10.304459in}{0.361579in}}{\pgfqpoint{10.302406in}{0.362047in}}%
\pgfpathquadraticcurveto{\pgfqpoint{10.300352in}{0.362533in}}{\pgfqpoint{10.297885in}{0.362533in}}%
\pgfpathquadraticcurveto{\pgfqpoint{10.289095in}{0.362533in}}{\pgfqpoint{10.284393in}{0.356823in}}%
\pgfpathquadraticcurveto{\pgfqpoint{10.279692in}{0.351114in}}{\pgfqpoint{10.279692in}{0.340414in}}%
\pgfpathlineto{\pgfqpoint{10.279692in}{0.307200in}}%
\pgfpathlineto{\pgfqpoint{10.269281in}{0.307200in}}%
\pgfpathlineto{\pgfqpoint{10.269281in}{0.370243in}}%
\pgfpathlineto{\pgfqpoint{10.279692in}{0.370243in}}%
\pgfpathlineto{\pgfqpoint{10.279692in}{0.360444in}}%
\pgfpathquadraticcurveto{\pgfqpoint{10.282970in}{0.366190in}}{\pgfqpoint{10.288194in}{0.368964in}}%
\pgfpathquadraticcurveto{\pgfqpoint{10.293436in}{0.371756in}}{\pgfqpoint{10.300929in}{0.371756in}}%
\pgfpathquadraticcurveto{\pgfqpoint{10.301991in}{0.371756in}}{\pgfqpoint{10.303288in}{0.371611in}}%
\pgfpathquadraticcurveto{\pgfqpoint{10.304585in}{0.371485in}}{\pgfqpoint{10.306152in}{0.371197in}}%
\pgfpathlineto{\pgfqpoint{10.306206in}{0.360570in}}%
\pgfpathlineto{\pgfqpoint{10.306206in}{0.360570in}}%
\pgfpathclose%
\pgfpathmoveto{\pgfqpoint{10.368456in}{0.341315in}}%
\pgfpathlineto{\pgfqpoint{10.368456in}{0.336254in}}%
\pgfpathlineto{\pgfqpoint{10.320832in}{0.336254in}}%
\pgfpathquadraticcurveto{\pgfqpoint{10.321516in}{0.325554in}}{\pgfqpoint{10.327280in}{0.319953in}}%
\pgfpathquadraticcurveto{\pgfqpoint{10.333044in}{0.314351in}}{\pgfqpoint{10.343347in}{0.314351in}}%
\pgfpathquadraticcurveto{\pgfqpoint{10.349309in}{0.314351in}}{\pgfqpoint{10.354911in}{0.315810in}}%
\pgfpathquadraticcurveto{\pgfqpoint{10.360513in}{0.317269in}}{\pgfqpoint{10.366043in}{0.320205in}}%
\pgfpathlineto{\pgfqpoint{10.366043in}{0.310406in}}%
\pgfpathquadraticcurveto{\pgfqpoint{10.360459in}{0.308047in}}{\pgfqpoint{10.354605in}{0.306804in}}%
\pgfpathquadraticcurveto{\pgfqpoint{10.348751in}{0.305561in}}{\pgfqpoint{10.342735in}{0.305561in}}%
\pgfpathquadraticcurveto{\pgfqpoint{10.327641in}{0.305561in}}{\pgfqpoint{10.318833in}{0.314333in}}%
\pgfpathquadraticcurveto{\pgfqpoint{10.310025in}{0.323123in}}{\pgfqpoint{10.310025in}{0.338109in}}%
\pgfpathquadraticcurveto{\pgfqpoint{10.310025in}{0.353581in}}{\pgfqpoint{10.318382in}{0.362659in}}%
\pgfpathquadraticcurveto{\pgfqpoint{10.326740in}{0.371756in}}{\pgfqpoint{10.340934in}{0.371756in}}%
\pgfpathquadraticcurveto{\pgfqpoint{10.353650in}{0.371756in}}{\pgfqpoint{10.361053in}{0.363560in}}%
\pgfpathquadraticcurveto{\pgfqpoint{10.368456in}{0.355383in}}{\pgfqpoint{10.368456in}{0.341315in}}%
\pgfpathlineto{\pgfqpoint{10.368456in}{0.341315in}}%
\pgfpathclose%
\pgfpathmoveto{\pgfqpoint{10.358099in}{0.344359in}}%
\pgfpathquadraticcurveto{\pgfqpoint{10.357991in}{0.352843in}}{\pgfqpoint{10.353344in}{0.357904in}}%
\pgfpathquadraticcurveto{\pgfqpoint{10.348697in}{0.362984in}}{\pgfqpoint{10.341042in}{0.362984in}}%
\pgfpathquadraticcurveto{\pgfqpoint{10.332378in}{0.362984in}}{\pgfqpoint{10.327172in}{0.358084in}}%
\pgfpathquadraticcurveto{\pgfqpoint{10.321967in}{0.353185in}}{\pgfqpoint{10.321174in}{0.344287in}}%
\pgfpathlineto{\pgfqpoint{10.358099in}{0.344359in}}%
\pgfpathlineto{\pgfqpoint{10.358099in}{0.344359in}}%
\pgfpathclose%
\pgfpathmoveto{\pgfqpoint{10.439388in}{0.341315in}}%
\pgfpathlineto{\pgfqpoint{10.439388in}{0.336254in}}%
\pgfpathlineto{\pgfqpoint{10.391764in}{0.336254in}}%
\pgfpathquadraticcurveto{\pgfqpoint{10.392448in}{0.325554in}}{\pgfqpoint{10.398212in}{0.319953in}}%
\pgfpathquadraticcurveto{\pgfqpoint{10.403976in}{0.314351in}}{\pgfqpoint{10.414279in}{0.314351in}}%
\pgfpathquadraticcurveto{\pgfqpoint{10.420241in}{0.314351in}}{\pgfqpoint{10.425843in}{0.315810in}}%
\pgfpathquadraticcurveto{\pgfqpoint{10.431445in}{0.317269in}}{\pgfqpoint{10.436974in}{0.320205in}}%
\pgfpathlineto{\pgfqpoint{10.436974in}{0.310406in}}%
\pgfpathquadraticcurveto{\pgfqpoint{10.431391in}{0.308047in}}{\pgfqpoint{10.425537in}{0.306804in}}%
\pgfpathquadraticcurveto{\pgfqpoint{10.419683in}{0.305561in}}{\pgfqpoint{10.413667in}{0.305561in}}%
\pgfpathquadraticcurveto{\pgfqpoint{10.398572in}{0.305561in}}{\pgfqpoint{10.389765in}{0.314333in}}%
\pgfpathquadraticcurveto{\pgfqpoint{10.380957in}{0.323123in}}{\pgfqpoint{10.380957in}{0.338109in}}%
\pgfpathquadraticcurveto{\pgfqpoint{10.380957in}{0.353581in}}{\pgfqpoint{10.389314in}{0.362659in}}%
\pgfpathquadraticcurveto{\pgfqpoint{10.397672in}{0.371756in}}{\pgfqpoint{10.411865in}{0.371756in}}%
\pgfpathquadraticcurveto{\pgfqpoint{10.424582in}{0.371756in}}{\pgfqpoint{10.431985in}{0.363560in}}%
\pgfpathquadraticcurveto{\pgfqpoint{10.439388in}{0.355383in}}{\pgfqpoint{10.439388in}{0.341315in}}%
\pgfpathlineto{\pgfqpoint{10.439388in}{0.341315in}}%
\pgfpathclose%
\pgfpathmoveto{\pgfqpoint{10.429031in}{0.344359in}}%
\pgfpathquadraticcurveto{\pgfqpoint{10.428923in}{0.352843in}}{\pgfqpoint{10.424276in}{0.357904in}}%
\pgfpathquadraticcurveto{\pgfqpoint{10.419629in}{0.362984in}}{\pgfqpoint{10.411974in}{0.362984in}}%
\pgfpathquadraticcurveto{\pgfqpoint{10.403310in}{0.362984in}}{\pgfqpoint{10.398104in}{0.358084in}}%
\pgfpathquadraticcurveto{\pgfqpoint{10.392899in}{0.353185in}}{\pgfqpoint{10.392106in}{0.344287in}}%
\pgfpathlineto{\pgfqpoint{10.429031in}{0.344359in}}%
\pgfpathlineto{\pgfqpoint{10.429031in}{0.344359in}}%
\pgfpathclose%
\pgfpathmoveto{\pgfqpoint{10.503277in}{0.388147in}}%
\pgfpathlineto{\pgfqpoint{10.503277in}{0.370243in}}%
\pgfpathlineto{\pgfqpoint{10.524604in}{0.370243in}}%
\pgfpathlineto{\pgfqpoint{10.524604in}{0.362191in}}%
\pgfpathlineto{\pgfqpoint{10.503277in}{0.362191in}}%
\pgfpathlineto{\pgfqpoint{10.503277in}{0.327968in}}%
\pgfpathquadraticcurveto{\pgfqpoint{10.503277in}{0.320259in}}{\pgfqpoint{10.505385in}{0.318061in}}%
\pgfpathquadraticcurveto{\pgfqpoint{10.507492in}{0.315864in}}{\pgfqpoint{10.513976in}{0.315864in}}%
\pgfpathlineto{\pgfqpoint{10.524604in}{0.315864in}}%
\pgfpathlineto{\pgfqpoint{10.524604in}{0.307200in}}%
\pgfpathlineto{\pgfqpoint{10.513976in}{0.307200in}}%
\pgfpathquadraticcurveto{\pgfqpoint{10.501980in}{0.307200in}}{\pgfqpoint{10.497423in}{0.311667in}}%
\pgfpathquadraticcurveto{\pgfqpoint{10.492866in}{0.316152in}}{\pgfqpoint{10.492866in}{0.327968in}}%
\pgfpathlineto{\pgfqpoint{10.492866in}{0.362191in}}%
\pgfpathlineto{\pgfqpoint{10.485265in}{0.362191in}}%
\pgfpathlineto{\pgfqpoint{10.485265in}{0.370243in}}%
\pgfpathlineto{\pgfqpoint{10.492866in}{0.370243in}}%
\pgfpathlineto{\pgfqpoint{10.492866in}{0.388147in}}%
\pgfpathlineto{\pgfqpoint{10.503277in}{0.388147in}}%
\pgfpathlineto{\pgfqpoint{10.503277in}{0.388147in}}%
\pgfpathclose%
\pgfpathmoveto{\pgfqpoint{10.574749in}{0.360570in}}%
\pgfpathquadraticcurveto{\pgfqpoint{10.573002in}{0.361579in}}{\pgfqpoint{10.570949in}{0.362047in}}%
\pgfpathquadraticcurveto{\pgfqpoint{10.568895in}{0.362533in}}{\pgfqpoint{10.566428in}{0.362533in}}%
\pgfpathquadraticcurveto{\pgfqpoint{10.557638in}{0.362533in}}{\pgfqpoint{10.552937in}{0.356823in}}%
\pgfpathquadraticcurveto{\pgfqpoint{10.548235in}{0.351114in}}{\pgfqpoint{10.548235in}{0.340414in}}%
\pgfpathlineto{\pgfqpoint{10.548235in}{0.307200in}}%
\pgfpathlineto{\pgfqpoint{10.537824in}{0.307200in}}%
\pgfpathlineto{\pgfqpoint{10.537824in}{0.370243in}}%
\pgfpathlineto{\pgfqpoint{10.548235in}{0.370243in}}%
\pgfpathlineto{\pgfqpoint{10.548235in}{0.360444in}}%
\pgfpathquadraticcurveto{\pgfqpoint{10.551514in}{0.366190in}}{\pgfqpoint{10.556737in}{0.368964in}}%
\pgfpathquadraticcurveto{\pgfqpoint{10.561979in}{0.371756in}}{\pgfqpoint{10.569472in}{0.371756in}}%
\pgfpathquadraticcurveto{\pgfqpoint{10.570535in}{0.371756in}}{\pgfqpoint{10.571831in}{0.371611in}}%
\pgfpathquadraticcurveto{\pgfqpoint{10.573128in}{0.371485in}}{\pgfqpoint{10.574695in}{0.371197in}}%
\pgfpathlineto{\pgfqpoint{10.574749in}{0.360570in}}%
\pgfpathlineto{\pgfqpoint{10.574749in}{0.360570in}}%
\pgfpathclose%
\pgfpathmoveto{\pgfqpoint{10.636999in}{0.341315in}}%
\pgfpathlineto{\pgfqpoint{10.636999in}{0.336254in}}%
\pgfpathlineto{\pgfqpoint{10.589375in}{0.336254in}}%
\pgfpathquadraticcurveto{\pgfqpoint{10.590060in}{0.325554in}}{\pgfqpoint{10.595824in}{0.319953in}}%
\pgfpathquadraticcurveto{\pgfqpoint{10.601587in}{0.314351in}}{\pgfqpoint{10.611890in}{0.314351in}}%
\pgfpathquadraticcurveto{\pgfqpoint{10.617852in}{0.314351in}}{\pgfqpoint{10.623454in}{0.315810in}}%
\pgfpathquadraticcurveto{\pgfqpoint{10.629056in}{0.317269in}}{\pgfqpoint{10.634586in}{0.320205in}}%
\pgfpathlineto{\pgfqpoint{10.634586in}{0.310406in}}%
\pgfpathquadraticcurveto{\pgfqpoint{10.629002in}{0.308047in}}{\pgfqpoint{10.623148in}{0.306804in}}%
\pgfpathquadraticcurveto{\pgfqpoint{10.617294in}{0.305561in}}{\pgfqpoint{10.611278in}{0.305561in}}%
\pgfpathquadraticcurveto{\pgfqpoint{10.596184in}{0.305561in}}{\pgfqpoint{10.587376in}{0.314333in}}%
\pgfpathquadraticcurveto{\pgfqpoint{10.578568in}{0.323123in}}{\pgfqpoint{10.578568in}{0.338109in}}%
\pgfpathquadraticcurveto{\pgfqpoint{10.578568in}{0.353581in}}{\pgfqpoint{10.586926in}{0.362659in}}%
\pgfpathquadraticcurveto{\pgfqpoint{10.595283in}{0.371756in}}{\pgfqpoint{10.609477in}{0.371756in}}%
\pgfpathquadraticcurveto{\pgfqpoint{10.622193in}{0.371756in}}{\pgfqpoint{10.629596in}{0.363560in}}%
\pgfpathquadraticcurveto{\pgfqpoint{10.636999in}{0.355383in}}{\pgfqpoint{10.636999in}{0.341315in}}%
\pgfpathlineto{\pgfqpoint{10.636999in}{0.341315in}}%
\pgfpathclose%
\pgfpathmoveto{\pgfqpoint{10.626642in}{0.344359in}}%
\pgfpathquadraticcurveto{\pgfqpoint{10.626534in}{0.352843in}}{\pgfqpoint{10.621887in}{0.357904in}}%
\pgfpathquadraticcurveto{\pgfqpoint{10.617240in}{0.362984in}}{\pgfqpoint{10.609585in}{0.362984in}}%
\pgfpathquadraticcurveto{\pgfqpoint{10.600921in}{0.362984in}}{\pgfqpoint{10.595715in}{0.358084in}}%
\pgfpathquadraticcurveto{\pgfqpoint{10.590510in}{0.353185in}}{\pgfqpoint{10.589717in}{0.344287in}}%
\pgfpathlineto{\pgfqpoint{10.626642in}{0.344359in}}%
\pgfpathlineto{\pgfqpoint{10.626642in}{0.344359in}}%
\pgfpathclose%
\pgfpathmoveto{\pgfqpoint{10.682660in}{0.338883in}}%
\pgfpathquadraticcurveto{\pgfqpoint{10.670106in}{0.338883in}}{\pgfqpoint{10.665260in}{0.336019in}}%
\pgfpathquadraticcurveto{\pgfqpoint{10.660415in}{0.333156in}}{\pgfqpoint{10.660415in}{0.326221in}}%
\pgfpathquadraticcurveto{\pgfqpoint{10.660415in}{0.320709in}}{\pgfqpoint{10.664054in}{0.317467in}}%
\pgfpathquadraticcurveto{\pgfqpoint{10.667692in}{0.314243in}}{\pgfqpoint{10.673924in}{0.314243in}}%
\pgfpathquadraticcurveto{\pgfqpoint{10.682552in}{0.314243in}}{\pgfqpoint{10.687758in}{0.320349in}}%
\pgfpathquadraticcurveto{\pgfqpoint{10.692963in}{0.326455in}}{\pgfqpoint{10.692963in}{0.336578in}}%
\pgfpathlineto{\pgfqpoint{10.692963in}{0.338883in}}%
\pgfpathlineto{\pgfqpoint{10.682660in}{0.338883in}}%
\pgfpathlineto{\pgfqpoint{10.682660in}{0.338883in}}%
\pgfpathclose%
\pgfpathmoveto{\pgfqpoint{10.703320in}{0.343170in}}%
\pgfpathlineto{\pgfqpoint{10.703320in}{0.307200in}}%
\pgfpathlineto{\pgfqpoint{10.692963in}{0.307200in}}%
\pgfpathlineto{\pgfqpoint{10.692963in}{0.316764in}}%
\pgfpathquadraticcurveto{\pgfqpoint{10.689415in}{0.311037in}}{\pgfqpoint{10.684119in}{0.308299in}}%
\pgfpathquadraticcurveto{\pgfqpoint{10.678824in}{0.305561in}}{\pgfqpoint{10.671168in}{0.305561in}}%
\pgfpathquadraticcurveto{\pgfqpoint{10.661496in}{0.305561in}}{\pgfqpoint{10.655768in}{0.311001in}}%
\pgfpathquadraticcurveto{\pgfqpoint{10.650058in}{0.316440in}}{\pgfqpoint{10.650058in}{0.325554in}}%
\pgfpathquadraticcurveto{\pgfqpoint{10.650058in}{0.336182in}}{\pgfqpoint{10.657173in}{0.341585in}}%
\pgfpathquadraticcurveto{\pgfqpoint{10.664306in}{0.346989in}}{\pgfqpoint{10.678427in}{0.346989in}}%
\pgfpathlineto{\pgfqpoint{10.692963in}{0.346989in}}%
\pgfpathlineto{\pgfqpoint{10.692963in}{0.348016in}}%
\pgfpathquadraticcurveto{\pgfqpoint{10.692963in}{0.355166in}}{\pgfqpoint{10.688262in}{0.359075in}}%
\pgfpathquadraticcurveto{\pgfqpoint{10.683561in}{0.362984in}}{\pgfqpoint{10.675059in}{0.362984in}}%
\pgfpathquadraticcurveto{\pgfqpoint{10.669655in}{0.362984in}}{\pgfqpoint{10.664522in}{0.361687in}}%
\pgfpathquadraticcurveto{\pgfqpoint{10.659406in}{0.360390in}}{\pgfqpoint{10.654687in}{0.357796in}}%
\pgfpathlineto{\pgfqpoint{10.654687in}{0.367379in}}%
\pgfpathquadraticcurveto{\pgfqpoint{10.660361in}{0.369576in}}{\pgfqpoint{10.665711in}{0.370657in}}%
\pgfpathquadraticcurveto{\pgfqpoint{10.671060in}{0.371756in}}{\pgfqpoint{10.676122in}{0.371756in}}%
\pgfpathquadraticcurveto{\pgfqpoint{10.689811in}{0.371756in}}{\pgfqpoint{10.696566in}{0.364659in}}%
\pgfpathquadraticcurveto{\pgfqpoint{10.703320in}{0.357580in}}{\pgfqpoint{10.703320in}{0.343170in}}%
\pgfpathlineto{\pgfqpoint{10.703320in}{0.343170in}}%
\pgfpathclose%
\pgfpathmoveto{\pgfqpoint{10.734895in}{0.388147in}}%
\pgfpathlineto{\pgfqpoint{10.734895in}{0.370243in}}%
\pgfpathlineto{\pgfqpoint{10.756222in}{0.370243in}}%
\pgfpathlineto{\pgfqpoint{10.756222in}{0.362191in}}%
\pgfpathlineto{\pgfqpoint{10.734895in}{0.362191in}}%
\pgfpathlineto{\pgfqpoint{10.734895in}{0.327968in}}%
\pgfpathquadraticcurveto{\pgfqpoint{10.734895in}{0.320259in}}{\pgfqpoint{10.737003in}{0.318061in}}%
\pgfpathquadraticcurveto{\pgfqpoint{10.739110in}{0.315864in}}{\pgfqpoint{10.745595in}{0.315864in}}%
\pgfpathlineto{\pgfqpoint{10.756222in}{0.315864in}}%
\pgfpathlineto{\pgfqpoint{10.756222in}{0.307200in}}%
\pgfpathlineto{\pgfqpoint{10.745595in}{0.307200in}}%
\pgfpathquadraticcurveto{\pgfqpoint{10.733599in}{0.307200in}}{\pgfqpoint{10.729041in}{0.311667in}}%
\pgfpathquadraticcurveto{\pgfqpoint{10.724484in}{0.316152in}}{\pgfqpoint{10.724484in}{0.327968in}}%
\pgfpathlineto{\pgfqpoint{10.724484in}{0.362191in}}%
\pgfpathlineto{\pgfqpoint{10.716883in}{0.362191in}}%
\pgfpathlineto{\pgfqpoint{10.716883in}{0.370243in}}%
\pgfpathlineto{\pgfqpoint{10.724484in}{0.370243in}}%
\pgfpathlineto{\pgfqpoint{10.724484in}{0.388147in}}%
\pgfpathlineto{\pgfqpoint{10.734895in}{0.388147in}}%
\pgfpathlineto{\pgfqpoint{10.734895in}{0.388147in}}%
\pgfpathclose%
\pgfpathmoveto{\pgfqpoint{10.818922in}{0.358138in}}%
\pgfpathquadraticcurveto{\pgfqpoint{10.822813in}{0.365127in}}{\pgfqpoint{10.828216in}{0.368441in}}%
\pgfpathquadraticcurveto{\pgfqpoint{10.833620in}{0.371756in}}{\pgfqpoint{10.840933in}{0.371756in}}%
\pgfpathquadraticcurveto{\pgfqpoint{10.850786in}{0.371756in}}{\pgfqpoint{10.856135in}{0.364857in}}%
\pgfpathquadraticcurveto{\pgfqpoint{10.861485in}{0.357976in}}{\pgfqpoint{10.861485in}{0.345260in}}%
\pgfpathlineto{\pgfqpoint{10.861485in}{0.307200in}}%
\pgfpathlineto{\pgfqpoint{10.851074in}{0.307200in}}%
\pgfpathlineto{\pgfqpoint{10.851074in}{0.344917in}}%
\pgfpathquadraticcurveto{\pgfqpoint{10.851074in}{0.353978in}}{\pgfqpoint{10.847850in}{0.358355in}}%
\pgfpathquadraticcurveto{\pgfqpoint{10.844643in}{0.362749in}}{\pgfqpoint{10.838069in}{0.362749in}}%
\pgfpathquadraticcurveto{\pgfqpoint{10.830018in}{0.362749in}}{\pgfqpoint{10.825334in}{0.357400in}}%
\pgfpathquadraticcurveto{\pgfqpoint{10.820669in}{0.352068in}}{\pgfqpoint{10.820669in}{0.342828in}}%
\pgfpathlineto{\pgfqpoint{10.820669in}{0.307200in}}%
\pgfpathlineto{\pgfqpoint{10.810258in}{0.307200in}}%
\pgfpathlineto{\pgfqpoint{10.810258in}{0.344917in}}%
\pgfpathquadraticcurveto{\pgfqpoint{10.810258in}{0.354032in}}{\pgfqpoint{10.807052in}{0.358391in}}%
\pgfpathquadraticcurveto{\pgfqpoint{10.803846in}{0.362749in}}{\pgfqpoint{10.797145in}{0.362749in}}%
\pgfpathquadraticcurveto{\pgfqpoint{10.789202in}{0.362749in}}{\pgfqpoint{10.784519in}{0.357382in}}%
\pgfpathquadraticcurveto{\pgfqpoint{10.779854in}{0.352014in}}{\pgfqpoint{10.779854in}{0.342828in}}%
\pgfpathlineto{\pgfqpoint{10.779854in}{0.307200in}}%
\pgfpathlineto{\pgfqpoint{10.769443in}{0.307200in}}%
\pgfpathlineto{\pgfqpoint{10.769443in}{0.370243in}}%
\pgfpathlineto{\pgfqpoint{10.779854in}{0.370243in}}%
\pgfpathlineto{\pgfqpoint{10.779854in}{0.360444in}}%
\pgfpathquadraticcurveto{\pgfqpoint{10.783402in}{0.366244in}}{\pgfqpoint{10.788355in}{0.369000in}}%
\pgfpathquadraticcurveto{\pgfqpoint{10.793309in}{0.371756in}}{\pgfqpoint{10.800117in}{0.371756in}}%
\pgfpathquadraticcurveto{\pgfqpoint{10.806998in}{0.371756in}}{\pgfqpoint{10.811807in}{0.368261in}}%
\pgfpathquadraticcurveto{\pgfqpoint{10.816617in}{0.364785in}}{\pgfqpoint{10.818922in}{0.358138in}}%
\pgfpathlineto{\pgfqpoint{10.818922in}{0.358138in}}%
\pgfpathclose%
\pgfpathmoveto{\pgfqpoint{10.936055in}{0.341315in}}%
\pgfpathlineto{\pgfqpoint{10.936055in}{0.336254in}}%
\pgfpathlineto{\pgfqpoint{10.888431in}{0.336254in}}%
\pgfpathquadraticcurveto{\pgfqpoint{10.889115in}{0.325554in}}{\pgfqpoint{10.894879in}{0.319953in}}%
\pgfpathquadraticcurveto{\pgfqpoint{10.900643in}{0.314351in}}{\pgfqpoint{10.910946in}{0.314351in}}%
\pgfpathquadraticcurveto{\pgfqpoint{10.916908in}{0.314351in}}{\pgfqpoint{10.922510in}{0.315810in}}%
\pgfpathquadraticcurveto{\pgfqpoint{10.928112in}{0.317269in}}{\pgfqpoint{10.933641in}{0.320205in}}%
\pgfpathlineto{\pgfqpoint{10.933641in}{0.310406in}}%
\pgfpathquadraticcurveto{\pgfqpoint{10.928058in}{0.308047in}}{\pgfqpoint{10.922204in}{0.306804in}}%
\pgfpathquadraticcurveto{\pgfqpoint{10.916350in}{0.305561in}}{\pgfqpoint{10.910334in}{0.305561in}}%
\pgfpathquadraticcurveto{\pgfqpoint{10.895240in}{0.305561in}}{\pgfqpoint{10.886432in}{0.314333in}}%
\pgfpathquadraticcurveto{\pgfqpoint{10.877624in}{0.323123in}}{\pgfqpoint{10.877624in}{0.338109in}}%
\pgfpathquadraticcurveto{\pgfqpoint{10.877624in}{0.353581in}}{\pgfqpoint{10.885981in}{0.362659in}}%
\pgfpathquadraticcurveto{\pgfqpoint{10.894339in}{0.371756in}}{\pgfqpoint{10.908533in}{0.371756in}}%
\pgfpathquadraticcurveto{\pgfqpoint{10.921249in}{0.371756in}}{\pgfqpoint{10.928652in}{0.363560in}}%
\pgfpathquadraticcurveto{\pgfqpoint{10.936055in}{0.355383in}}{\pgfqpoint{10.936055in}{0.341315in}}%
\pgfpathlineto{\pgfqpoint{10.936055in}{0.341315in}}%
\pgfpathclose%
\pgfpathmoveto{\pgfqpoint{10.925698in}{0.344359in}}%
\pgfpathquadraticcurveto{\pgfqpoint{10.925590in}{0.352843in}}{\pgfqpoint{10.920943in}{0.357904in}}%
\pgfpathquadraticcurveto{\pgfqpoint{10.916296in}{0.362984in}}{\pgfqpoint{10.908641in}{0.362984in}}%
\pgfpathquadraticcurveto{\pgfqpoint{10.899977in}{0.362984in}}{\pgfqpoint{10.894771in}{0.358084in}}%
\pgfpathquadraticcurveto{\pgfqpoint{10.889566in}{0.353185in}}{\pgfqpoint{10.888773in}{0.344287in}}%
\pgfpathlineto{\pgfqpoint{10.925698in}{0.344359in}}%
\pgfpathlineto{\pgfqpoint{10.925698in}{0.344359in}}%
\pgfpathclose%
\pgfpathmoveto{\pgfqpoint{11.005474in}{0.345260in}}%
\pgfpathlineto{\pgfqpoint{11.005474in}{0.307200in}}%
\pgfpathlineto{\pgfqpoint{10.995117in}{0.307200in}}%
\pgfpathlineto{\pgfqpoint{10.995117in}{0.344917in}}%
\pgfpathquadraticcurveto{\pgfqpoint{10.995117in}{0.353869in}}{\pgfqpoint{10.991623in}{0.358300in}}%
\pgfpathquadraticcurveto{\pgfqpoint{10.988128in}{0.362749in}}{\pgfqpoint{10.981158in}{0.362749in}}%
\pgfpathquadraticcurveto{\pgfqpoint{10.972764in}{0.362749in}}{\pgfqpoint{10.967919in}{0.357400in}}%
\pgfpathquadraticcurveto{\pgfqpoint{10.963073in}{0.352068in}}{\pgfqpoint{10.963073in}{0.342828in}}%
\pgfpathlineto{\pgfqpoint{10.963073in}{0.307200in}}%
\pgfpathlineto{\pgfqpoint{10.952662in}{0.307200in}}%
\pgfpathlineto{\pgfqpoint{10.952662in}{0.370243in}}%
\pgfpathlineto{\pgfqpoint{10.963073in}{0.370243in}}%
\pgfpathlineto{\pgfqpoint{10.963073in}{0.360444in}}%
\pgfpathquadraticcurveto{\pgfqpoint{10.966802in}{0.366136in}}{\pgfqpoint{10.971827in}{0.368946in}}%
\pgfpathquadraticcurveto{\pgfqpoint{10.976871in}{0.371756in}}{\pgfqpoint{10.983463in}{0.371756in}}%
\pgfpathquadraticcurveto{\pgfqpoint{10.994324in}{0.371756in}}{\pgfqpoint{10.999890in}{0.365037in}}%
\pgfpathquadraticcurveto{\pgfqpoint{11.005474in}{0.358318in}}{\pgfqpoint{11.005474in}{0.345260in}}%
\pgfpathlineto{\pgfqpoint{11.005474in}{0.345260in}}%
\pgfpathclose%
\pgfpathmoveto{\pgfqpoint{11.036365in}{0.388147in}}%
\pgfpathlineto{\pgfqpoint{11.036365in}{0.370243in}}%
\pgfpathlineto{\pgfqpoint{11.057691in}{0.370243in}}%
\pgfpathlineto{\pgfqpoint{11.057691in}{0.362191in}}%
\pgfpathlineto{\pgfqpoint{11.036365in}{0.362191in}}%
\pgfpathlineto{\pgfqpoint{11.036365in}{0.327968in}}%
\pgfpathquadraticcurveto{\pgfqpoint{11.036365in}{0.320259in}}{\pgfqpoint{11.038472in}{0.318061in}}%
\pgfpathquadraticcurveto{\pgfqpoint{11.040580in}{0.315864in}}{\pgfqpoint{11.047064in}{0.315864in}}%
\pgfpathlineto{\pgfqpoint{11.057691in}{0.315864in}}%
\pgfpathlineto{\pgfqpoint{11.057691in}{0.307200in}}%
\pgfpathlineto{\pgfqpoint{11.047064in}{0.307200in}}%
\pgfpathquadraticcurveto{\pgfqpoint{11.035068in}{0.307200in}}{\pgfqpoint{11.030511in}{0.311667in}}%
\pgfpathquadraticcurveto{\pgfqpoint{11.025954in}{0.316152in}}{\pgfqpoint{11.025954in}{0.327968in}}%
\pgfpathlineto{\pgfqpoint{11.025954in}{0.362191in}}%
\pgfpathlineto{\pgfqpoint{11.018353in}{0.362191in}}%
\pgfpathlineto{\pgfqpoint{11.018353in}{0.370243in}}%
\pgfpathlineto{\pgfqpoint{11.025954in}{0.370243in}}%
\pgfpathlineto{\pgfqpoint{11.025954in}{0.388147in}}%
\pgfpathlineto{\pgfqpoint{11.036365in}{0.388147in}}%
\pgfpathlineto{\pgfqpoint{11.036365in}{0.388147in}}%
\pgfpathclose%
\pgfpathmoveto{\pgfqpoint{11.066085in}{0.343386in}}%
\pgfpathlineto{\pgfqpoint{11.096417in}{0.343386in}}%
\pgfpathlineto{\pgfqpoint{11.096417in}{0.334164in}}%
\pgfpathlineto{\pgfqpoint{11.066085in}{0.334164in}}%
\pgfpathlineto{\pgfqpoint{11.066085in}{0.343386in}}%
\pgfpathlineto{\pgfqpoint{11.066085in}{0.343386in}}%
\pgfpathclose%
\pgfpathmoveto{\pgfqpoint{11.149427in}{0.360570in}}%
\pgfpathquadraticcurveto{\pgfqpoint{11.147680in}{0.361579in}}{\pgfqpoint{11.145627in}{0.362047in}}%
\pgfpathquadraticcurveto{\pgfqpoint{11.143573in}{0.362533in}}{\pgfqpoint{11.141105in}{0.362533in}}%
\pgfpathquadraticcurveto{\pgfqpoint{11.132316in}{0.362533in}}{\pgfqpoint{11.127614in}{0.356823in}}%
\pgfpathquadraticcurveto{\pgfqpoint{11.122913in}{0.351114in}}{\pgfqpoint{11.122913in}{0.340414in}}%
\pgfpathlineto{\pgfqpoint{11.122913in}{0.307200in}}%
\pgfpathlineto{\pgfqpoint{11.112502in}{0.307200in}}%
\pgfpathlineto{\pgfqpoint{11.112502in}{0.370243in}}%
\pgfpathlineto{\pgfqpoint{11.122913in}{0.370243in}}%
\pgfpathlineto{\pgfqpoint{11.122913in}{0.360444in}}%
\pgfpathquadraticcurveto{\pgfqpoint{11.126191in}{0.366190in}}{\pgfqpoint{11.131415in}{0.368964in}}%
\pgfpathquadraticcurveto{\pgfqpoint{11.136656in}{0.371756in}}{\pgfqpoint{11.144150in}{0.371756in}}%
\pgfpathquadraticcurveto{\pgfqpoint{11.145212in}{0.371756in}}{\pgfqpoint{11.146509in}{0.371611in}}%
\pgfpathquadraticcurveto{\pgfqpoint{11.147806in}{0.371485in}}{\pgfqpoint{11.149373in}{0.371197in}}%
\pgfpathlineto{\pgfqpoint{11.149427in}{0.360570in}}%
\pgfpathlineto{\pgfqpoint{11.149427in}{0.360570in}}%
\pgfpathclose%
\pgfpathmoveto{\pgfqpoint{11.182173in}{0.362984in}}%
\pgfpathquadraticcurveto{\pgfqpoint{11.173852in}{0.362984in}}{\pgfqpoint{11.169006in}{0.356481in}}%
\pgfpathquadraticcurveto{\pgfqpoint{11.164161in}{0.349979in}}{\pgfqpoint{11.164161in}{0.338667in}}%
\pgfpathquadraticcurveto{\pgfqpoint{11.164161in}{0.327356in}}{\pgfqpoint{11.168970in}{0.320853in}}%
\pgfpathquadraticcurveto{\pgfqpoint{11.173798in}{0.314351in}}{\pgfqpoint{11.182173in}{0.314351in}}%
\pgfpathquadraticcurveto{\pgfqpoint{11.190459in}{0.314351in}}{\pgfqpoint{11.195286in}{0.320871in}}%
\pgfpathquadraticcurveto{\pgfqpoint{11.200131in}{0.327410in}}{\pgfqpoint{11.200131in}{0.338667in}}%
\pgfpathquadraticcurveto{\pgfqpoint{11.200131in}{0.349871in}}{\pgfqpoint{11.195286in}{0.356427in}}%
\pgfpathquadraticcurveto{\pgfqpoint{11.190459in}{0.362984in}}{\pgfqpoint{11.182173in}{0.362984in}}%
\pgfpathlineto{\pgfqpoint{11.182173in}{0.362984in}}%
\pgfpathclose%
\pgfpathmoveto{\pgfqpoint{11.182173in}{0.371756in}}%
\pgfpathquadraticcurveto{\pgfqpoint{11.195682in}{0.371756in}}{\pgfqpoint{11.203391in}{0.362966in}}%
\pgfpathquadraticcurveto{\pgfqpoint{11.211119in}{0.354194in}}{\pgfqpoint{11.211119in}{0.338667in}}%
\pgfpathquadraticcurveto{\pgfqpoint{11.211119in}{0.323195in}}{\pgfqpoint{11.203391in}{0.314369in}}%
\pgfpathquadraticcurveto{\pgfqpoint{11.195682in}{0.305561in}}{\pgfqpoint{11.182173in}{0.305561in}}%
\pgfpathquadraticcurveto{\pgfqpoint{11.168610in}{0.305561in}}{\pgfqpoint{11.160919in}{0.314369in}}%
\pgfpathquadraticcurveto{\pgfqpoint{11.153246in}{0.323195in}}{\pgfqpoint{11.153246in}{0.338667in}}%
\pgfpathquadraticcurveto{\pgfqpoint{11.153246in}{0.354194in}}{\pgfqpoint{11.160919in}{0.362966in}}%
\pgfpathquadraticcurveto{\pgfqpoint{11.168610in}{0.371756in}}{\pgfqpoint{11.182173in}{0.371756in}}%
\pgfpathlineto{\pgfqpoint{11.182173in}{0.371756in}}%
\pgfpathclose%
\pgfpathmoveto{\pgfqpoint{11.228284in}{0.394793in}}%
\pgfpathlineto{\pgfqpoint{11.238641in}{0.394793in}}%
\pgfpathlineto{\pgfqpoint{11.238641in}{0.307200in}}%
\pgfpathlineto{\pgfqpoint{11.228284in}{0.307200in}}%
\pgfpathlineto{\pgfqpoint{11.228284in}{0.394793in}}%
\pgfpathlineto{\pgfqpoint{11.228284in}{0.394793in}}%
\pgfpathclose%
\pgfpathmoveto{\pgfqpoint{11.260310in}{0.394793in}}%
\pgfpathlineto{\pgfqpoint{11.270667in}{0.394793in}}%
\pgfpathlineto{\pgfqpoint{11.270667in}{0.307200in}}%
\pgfpathlineto{\pgfqpoint{11.260310in}{0.307200in}}%
\pgfpathlineto{\pgfqpoint{11.260310in}{0.394793in}}%
\pgfpathlineto{\pgfqpoint{11.260310in}{0.394793in}}%
\pgfpathclose%
\pgfpathmoveto{\pgfqpoint{11.316760in}{0.362984in}}%
\pgfpathquadraticcurveto{\pgfqpoint{11.308438in}{0.362984in}}{\pgfqpoint{11.303593in}{0.356481in}}%
\pgfpathquadraticcurveto{\pgfqpoint{11.298748in}{0.349979in}}{\pgfqpoint{11.298748in}{0.338667in}}%
\pgfpathquadraticcurveto{\pgfqpoint{11.298748in}{0.327356in}}{\pgfqpoint{11.303557in}{0.320853in}}%
\pgfpathquadraticcurveto{\pgfqpoint{11.308384in}{0.314351in}}{\pgfqpoint{11.316760in}{0.314351in}}%
\pgfpathquadraticcurveto{\pgfqpoint{11.325046in}{0.314351in}}{\pgfqpoint{11.329873in}{0.320871in}}%
\pgfpathquadraticcurveto{\pgfqpoint{11.334718in}{0.327410in}}{\pgfqpoint{11.334718in}{0.338667in}}%
\pgfpathquadraticcurveto{\pgfqpoint{11.334718in}{0.349871in}}{\pgfqpoint{11.329873in}{0.356427in}}%
\pgfpathquadraticcurveto{\pgfqpoint{11.325046in}{0.362984in}}{\pgfqpoint{11.316760in}{0.362984in}}%
\pgfpathlineto{\pgfqpoint{11.316760in}{0.362984in}}%
\pgfpathclose%
\pgfpathmoveto{\pgfqpoint{11.316760in}{0.371756in}}%
\pgfpathquadraticcurveto{\pgfqpoint{11.330269in}{0.371756in}}{\pgfqpoint{11.337978in}{0.362966in}}%
\pgfpathquadraticcurveto{\pgfqpoint{11.345706in}{0.354194in}}{\pgfqpoint{11.345706in}{0.338667in}}%
\pgfpathquadraticcurveto{\pgfqpoint{11.345706in}{0.323195in}}{\pgfqpoint{11.337978in}{0.314369in}}%
\pgfpathquadraticcurveto{\pgfqpoint{11.330269in}{0.305561in}}{\pgfqpoint{11.316760in}{0.305561in}}%
\pgfpathquadraticcurveto{\pgfqpoint{11.303197in}{0.305561in}}{\pgfqpoint{11.295506in}{0.314369in}}%
\pgfpathquadraticcurveto{\pgfqpoint{11.287832in}{0.323195in}}{\pgfqpoint{11.287832in}{0.338667in}}%
\pgfpathquadraticcurveto{\pgfqpoint{11.287832in}{0.354194in}}{\pgfqpoint{11.295506in}{0.362966in}}%
\pgfpathquadraticcurveto{\pgfqpoint{11.303197in}{0.371756in}}{\pgfqpoint{11.316760in}{0.371756in}}%
\pgfpathlineto{\pgfqpoint{11.316760in}{0.371756in}}%
\pgfpathclose%
\pgfpathmoveto{\pgfqpoint{11.361808in}{0.332075in}}%
\pgfpathlineto{\pgfqpoint{11.361808in}{0.370243in}}%
\pgfpathlineto{\pgfqpoint{11.372165in}{0.370243in}}%
\pgfpathlineto{\pgfqpoint{11.372165in}{0.332471in}}%
\pgfpathquadraticcurveto{\pgfqpoint{11.372165in}{0.323519in}}{\pgfqpoint{11.375642in}{0.319034in}}%
\pgfpathquadraticcurveto{\pgfqpoint{11.379136in}{0.314567in}}{\pgfqpoint{11.386125in}{0.314567in}}%
\pgfpathquadraticcurveto{\pgfqpoint{11.394500in}{0.314567in}}{\pgfqpoint{11.399364in}{0.319917in}}%
\pgfpathquadraticcurveto{\pgfqpoint{11.404245in}{0.325266in}}{\pgfqpoint{11.404245in}{0.334506in}}%
\pgfpathlineto{\pgfqpoint{11.404245in}{0.370243in}}%
\pgfpathlineto{\pgfqpoint{11.414602in}{0.370243in}}%
\pgfpathlineto{\pgfqpoint{11.414602in}{0.307200in}}%
\pgfpathlineto{\pgfqpoint{11.404245in}{0.307200in}}%
\pgfpathlineto{\pgfqpoint{11.404245in}{0.316891in}}%
\pgfpathquadraticcurveto{\pgfqpoint{11.400480in}{0.311145in}}{\pgfqpoint{11.395491in}{0.308353in}}%
\pgfpathquadraticcurveto{\pgfqpoint{11.390520in}{0.305561in}}{\pgfqpoint{11.383927in}{0.305561in}}%
\pgfpathquadraticcurveto{\pgfqpoint{11.373066in}{0.305561in}}{\pgfqpoint{11.367428in}{0.312315in}}%
\pgfpathquadraticcurveto{\pgfqpoint{11.361808in}{0.319070in}}{\pgfqpoint{11.361808in}{0.332075in}}%
\pgfpathlineto{\pgfqpoint{11.361808in}{0.332075in}}%
\pgfpathclose%
\pgfpathmoveto{\pgfqpoint{11.387872in}{0.371756in}}%
\pgfpathlineto{\pgfqpoint{11.387872in}{0.371756in}}%
\pgfpathlineto{\pgfqpoint{11.387872in}{0.371756in}}%
\pgfpathclose%
\pgfpathmoveto{\pgfqpoint{11.446177in}{0.388147in}}%
\pgfpathlineto{\pgfqpoint{11.446177in}{0.370243in}}%
\pgfpathlineto{\pgfqpoint{11.467504in}{0.370243in}}%
\pgfpathlineto{\pgfqpoint{11.467504in}{0.362191in}}%
\pgfpathlineto{\pgfqpoint{11.446177in}{0.362191in}}%
\pgfpathlineto{\pgfqpoint{11.446177in}{0.327968in}}%
\pgfpathquadraticcurveto{\pgfqpoint{11.446177in}{0.320259in}}{\pgfqpoint{11.448285in}{0.318061in}}%
\pgfpathquadraticcurveto{\pgfqpoint{11.450392in}{0.315864in}}{\pgfqpoint{11.456877in}{0.315864in}}%
\pgfpathlineto{\pgfqpoint{11.467504in}{0.315864in}}%
\pgfpathlineto{\pgfqpoint{11.467504in}{0.307200in}}%
\pgfpathlineto{\pgfqpoint{11.456877in}{0.307200in}}%
\pgfpathquadraticcurveto{\pgfqpoint{11.444880in}{0.307200in}}{\pgfqpoint{11.440323in}{0.311667in}}%
\pgfpathquadraticcurveto{\pgfqpoint{11.435766in}{0.316152in}}{\pgfqpoint{11.435766in}{0.327968in}}%
\pgfpathlineto{\pgfqpoint{11.435766in}{0.362191in}}%
\pgfpathlineto{\pgfqpoint{11.428165in}{0.362191in}}%
\pgfpathlineto{\pgfqpoint{11.428165in}{0.370243in}}%
\pgfpathlineto{\pgfqpoint{11.435766in}{0.370243in}}%
\pgfpathlineto{\pgfqpoint{11.435766in}{0.388147in}}%
\pgfpathlineto{\pgfqpoint{11.446177in}{0.388147in}}%
\pgfpathlineto{\pgfqpoint{11.446177in}{0.388147in}}%
\pgfpathclose%
\pgfpathmoveto{\pgfqpoint{11.554286in}{0.360570in}}%
\pgfpathquadraticcurveto{\pgfqpoint{11.552539in}{0.361579in}}{\pgfqpoint{11.550486in}{0.362047in}}%
\pgfpathquadraticcurveto{\pgfqpoint{11.548432in}{0.362533in}}{\pgfqpoint{11.545965in}{0.362533in}}%
\pgfpathquadraticcurveto{\pgfqpoint{11.537175in}{0.362533in}}{\pgfqpoint{11.532474in}{0.356823in}}%
\pgfpathquadraticcurveto{\pgfqpoint{11.527772in}{0.351114in}}{\pgfqpoint{11.527772in}{0.340414in}}%
\pgfpathlineto{\pgfqpoint{11.527772in}{0.307200in}}%
\pgfpathlineto{\pgfqpoint{11.517361in}{0.307200in}}%
\pgfpathlineto{\pgfqpoint{11.517361in}{0.370243in}}%
\pgfpathlineto{\pgfqpoint{11.527772in}{0.370243in}}%
\pgfpathlineto{\pgfqpoint{11.527772in}{0.360444in}}%
\pgfpathquadraticcurveto{\pgfqpoint{11.531051in}{0.366190in}}{\pgfqpoint{11.536274in}{0.368964in}}%
\pgfpathquadraticcurveto{\pgfqpoint{11.541516in}{0.371756in}}{\pgfqpoint{11.549009in}{0.371756in}}%
\pgfpathquadraticcurveto{\pgfqpoint{11.550071in}{0.371756in}}{\pgfqpoint{11.551368in}{0.371611in}}%
\pgfpathquadraticcurveto{\pgfqpoint{11.552665in}{0.371485in}}{\pgfqpoint{11.554232in}{0.371197in}}%
\pgfpathlineto{\pgfqpoint{11.554286in}{0.360570in}}%
\pgfpathlineto{\pgfqpoint{11.554286in}{0.360570in}}%
\pgfpathclose%
\pgfpathmoveto{\pgfqpoint{11.616536in}{0.341315in}}%
\pgfpathlineto{\pgfqpoint{11.616536in}{0.336254in}}%
\pgfpathlineto{\pgfqpoint{11.568912in}{0.336254in}}%
\pgfpathquadraticcurveto{\pgfqpoint{11.569597in}{0.325554in}}{\pgfqpoint{11.575360in}{0.319953in}}%
\pgfpathquadraticcurveto{\pgfqpoint{11.581124in}{0.314351in}}{\pgfqpoint{11.591427in}{0.314351in}}%
\pgfpathquadraticcurveto{\pgfqpoint{11.597389in}{0.314351in}}{\pgfqpoint{11.602991in}{0.315810in}}%
\pgfpathquadraticcurveto{\pgfqpoint{11.608593in}{0.317269in}}{\pgfqpoint{11.614123in}{0.320205in}}%
\pgfpathlineto{\pgfqpoint{11.614123in}{0.310406in}}%
\pgfpathquadraticcurveto{\pgfqpoint{11.608539in}{0.308047in}}{\pgfqpoint{11.602685in}{0.306804in}}%
\pgfpathquadraticcurveto{\pgfqpoint{11.596831in}{0.305561in}}{\pgfqpoint{11.590815in}{0.305561in}}%
\pgfpathquadraticcurveto{\pgfqpoint{11.575721in}{0.305561in}}{\pgfqpoint{11.566913in}{0.314333in}}%
\pgfpathquadraticcurveto{\pgfqpoint{11.558105in}{0.323123in}}{\pgfqpoint{11.558105in}{0.338109in}}%
\pgfpathquadraticcurveto{\pgfqpoint{11.558105in}{0.353581in}}{\pgfqpoint{11.566462in}{0.362659in}}%
\pgfpathquadraticcurveto{\pgfqpoint{11.574820in}{0.371756in}}{\pgfqpoint{11.589014in}{0.371756in}}%
\pgfpathquadraticcurveto{\pgfqpoint{11.601730in}{0.371756in}}{\pgfqpoint{11.609133in}{0.363560in}}%
\pgfpathquadraticcurveto{\pgfqpoint{11.616536in}{0.355383in}}{\pgfqpoint{11.616536in}{0.341315in}}%
\pgfpathlineto{\pgfqpoint{11.616536in}{0.341315in}}%
\pgfpathclose%
\pgfpathmoveto{\pgfqpoint{11.606179in}{0.344359in}}%
\pgfpathquadraticcurveto{\pgfqpoint{11.606071in}{0.352843in}}{\pgfqpoint{11.601424in}{0.357904in}}%
\pgfpathquadraticcurveto{\pgfqpoint{11.596777in}{0.362984in}}{\pgfqpoint{11.589122in}{0.362984in}}%
\pgfpathquadraticcurveto{\pgfqpoint{11.580458in}{0.362984in}}{\pgfqpoint{11.575252in}{0.358084in}}%
\pgfpathquadraticcurveto{\pgfqpoint{11.570047in}{0.353185in}}{\pgfqpoint{11.569254in}{0.344287in}}%
\pgfpathlineto{\pgfqpoint{11.606179in}{0.344359in}}%
\pgfpathlineto{\pgfqpoint{11.606179in}{0.344359in}}%
\pgfpathclose%
\pgfpathmoveto{\pgfqpoint{11.675022in}{0.339460in}}%
\pgfpathquadraticcurveto{\pgfqpoint{11.675022in}{0.350717in}}{\pgfqpoint{11.670375in}{0.356896in}}%
\pgfpathquadraticcurveto{\pgfqpoint{11.665745in}{0.363092in}}{\pgfqpoint{11.657352in}{0.363092in}}%
\pgfpathquadraticcurveto{\pgfqpoint{11.649030in}{0.363092in}}{\pgfqpoint{11.644383in}{0.356896in}}%
\pgfpathquadraticcurveto{\pgfqpoint{11.639736in}{0.350717in}}{\pgfqpoint{11.639736in}{0.339460in}}%
\pgfpathquadraticcurveto{\pgfqpoint{11.639736in}{0.328256in}}{\pgfqpoint{11.644383in}{0.322060in}}%
\pgfpathquadraticcurveto{\pgfqpoint{11.649030in}{0.315864in}}{\pgfqpoint{11.657352in}{0.315864in}}%
\pgfpathquadraticcurveto{\pgfqpoint{11.665745in}{0.315864in}}{\pgfqpoint{11.670375in}{0.322060in}}%
\pgfpathquadraticcurveto{\pgfqpoint{11.675022in}{0.328256in}}{\pgfqpoint{11.675022in}{0.339460in}}%
\pgfpathlineto{\pgfqpoint{11.675022in}{0.339460in}}%
\pgfpathclose%
\pgfpathmoveto{\pgfqpoint{11.685379in}{0.315017in}}%
\pgfpathquadraticcurveto{\pgfqpoint{11.685379in}{0.298932in}}{\pgfqpoint{11.678228in}{0.291079in}}%
\pgfpathquadraticcurveto{\pgfqpoint{11.671095in}{0.283226in}}{\pgfqpoint{11.656343in}{0.283226in}}%
\pgfpathquadraticcurveto{\pgfqpoint{11.650885in}{0.283226in}}{\pgfqpoint{11.646040in}{0.284036in}}%
\pgfpathquadraticcurveto{\pgfqpoint{11.641195in}{0.284847in}}{\pgfqpoint{11.636638in}{0.286540in}}%
\pgfpathlineto{\pgfqpoint{11.636638in}{0.296609in}}%
\pgfpathquadraticcurveto{\pgfqpoint{11.641195in}{0.294141in}}{\pgfqpoint{11.645644in}{0.292970in}}%
\pgfpathquadraticcurveto{\pgfqpoint{11.650093in}{0.291782in}}{\pgfqpoint{11.654704in}{0.291782in}}%
\pgfpathquadraticcurveto{\pgfqpoint{11.664899in}{0.291782in}}{\pgfqpoint{11.669960in}{0.297095in}}%
\pgfpathquadraticcurveto{\pgfqpoint{11.675022in}{0.302409in}}{\pgfqpoint{11.675022in}{0.313162in}}%
\pgfpathlineto{\pgfqpoint{11.675022in}{0.318295in}}%
\pgfpathquadraticcurveto{\pgfqpoint{11.671816in}{0.312712in}}{\pgfqpoint{11.666808in}{0.309956in}}%
\pgfpathquadraticcurveto{\pgfqpoint{11.661801in}{0.307200in}}{\pgfqpoint{11.654812in}{0.307200in}}%
\pgfpathquadraticcurveto{\pgfqpoint{11.643230in}{0.307200in}}{\pgfqpoint{11.636133in}{0.316026in}}%
\pgfpathquadraticcurveto{\pgfqpoint{11.629037in}{0.324870in}}{\pgfqpoint{11.629037in}{0.339460in}}%
\pgfpathquadraticcurveto{\pgfqpoint{11.629037in}{0.354086in}}{\pgfqpoint{11.636133in}{0.362912in}}%
\pgfpathquadraticcurveto{\pgfqpoint{11.643230in}{0.371756in}}{\pgfqpoint{11.654812in}{0.371756in}}%
\pgfpathquadraticcurveto{\pgfqpoint{11.661801in}{0.371756in}}{\pgfqpoint{11.666808in}{0.369000in}}%
\pgfpathquadraticcurveto{\pgfqpoint{11.671816in}{0.366244in}}{\pgfqpoint{11.675022in}{0.360678in}}%
\pgfpathlineto{\pgfqpoint{11.675022in}{0.370243in}}%
\pgfpathlineto{\pgfqpoint{11.685379in}{0.370243in}}%
\pgfpathlineto{\pgfqpoint{11.685379in}{0.315017in}}%
\pgfpathlineto{\pgfqpoint{11.685379in}{0.315017in}}%
\pgfpathclose%
\pgfpathmoveto{\pgfqpoint{11.706723in}{0.370243in}}%
\pgfpathlineto{\pgfqpoint{11.717080in}{0.370243in}}%
\pgfpathlineto{\pgfqpoint{11.717080in}{0.307200in}}%
\pgfpathlineto{\pgfqpoint{11.706723in}{0.307200in}}%
\pgfpathlineto{\pgfqpoint{11.706723in}{0.370243in}}%
\pgfpathlineto{\pgfqpoint{11.706723in}{0.370243in}}%
\pgfpathclose%
\pgfpathmoveto{\pgfqpoint{11.706723in}{0.394793in}}%
\pgfpathlineto{\pgfqpoint{11.717080in}{0.394793in}}%
\pgfpathlineto{\pgfqpoint{11.717080in}{0.381662in}}%
\pgfpathlineto{\pgfqpoint{11.706723in}{0.381662in}}%
\pgfpathlineto{\pgfqpoint{11.706723in}{0.394793in}}%
\pgfpathlineto{\pgfqpoint{11.706723in}{0.394793in}}%
\pgfpathclose%
\pgfpathmoveto{\pgfqpoint{11.763173in}{0.362984in}}%
\pgfpathquadraticcurveto{\pgfqpoint{11.754852in}{0.362984in}}{\pgfqpoint{11.750006in}{0.356481in}}%
\pgfpathquadraticcurveto{\pgfqpoint{11.745161in}{0.349979in}}{\pgfqpoint{11.745161in}{0.338667in}}%
\pgfpathquadraticcurveto{\pgfqpoint{11.745161in}{0.327356in}}{\pgfqpoint{11.749970in}{0.320853in}}%
\pgfpathquadraticcurveto{\pgfqpoint{11.754798in}{0.314351in}}{\pgfqpoint{11.763173in}{0.314351in}}%
\pgfpathquadraticcurveto{\pgfqpoint{11.771459in}{0.314351in}}{\pgfqpoint{11.776286in}{0.320871in}}%
\pgfpathquadraticcurveto{\pgfqpoint{11.781131in}{0.327410in}}{\pgfqpoint{11.781131in}{0.338667in}}%
\pgfpathquadraticcurveto{\pgfqpoint{11.781131in}{0.349871in}}{\pgfqpoint{11.776286in}{0.356427in}}%
\pgfpathquadraticcurveto{\pgfqpoint{11.771459in}{0.362984in}}{\pgfqpoint{11.763173in}{0.362984in}}%
\pgfpathlineto{\pgfqpoint{11.763173in}{0.362984in}}%
\pgfpathclose%
\pgfpathmoveto{\pgfqpoint{11.763173in}{0.371756in}}%
\pgfpathquadraticcurveto{\pgfqpoint{11.776682in}{0.371756in}}{\pgfqpoint{11.784391in}{0.362966in}}%
\pgfpathquadraticcurveto{\pgfqpoint{11.792119in}{0.354194in}}{\pgfqpoint{11.792119in}{0.338667in}}%
\pgfpathquadraticcurveto{\pgfqpoint{11.792119in}{0.323195in}}{\pgfqpoint{11.784391in}{0.314369in}}%
\pgfpathquadraticcurveto{\pgfqpoint{11.776682in}{0.305561in}}{\pgfqpoint{11.763173in}{0.305561in}}%
\pgfpathquadraticcurveto{\pgfqpoint{11.749610in}{0.305561in}}{\pgfqpoint{11.741919in}{0.314369in}}%
\pgfpathquadraticcurveto{\pgfqpoint{11.734246in}{0.323195in}}{\pgfqpoint{11.734246in}{0.338667in}}%
\pgfpathquadraticcurveto{\pgfqpoint{11.734246in}{0.354194in}}{\pgfqpoint{11.741919in}{0.362966in}}%
\pgfpathquadraticcurveto{\pgfqpoint{11.749610in}{0.371756in}}{\pgfqpoint{11.763173in}{0.371756in}}%
\pgfpathlineto{\pgfqpoint{11.763173in}{0.371756in}}%
\pgfpathclose%
\pgfpathmoveto{\pgfqpoint{11.861700in}{0.345260in}}%
\pgfpathlineto{\pgfqpoint{11.861700in}{0.307200in}}%
\pgfpathlineto{\pgfqpoint{11.851343in}{0.307200in}}%
\pgfpathlineto{\pgfqpoint{11.851343in}{0.344917in}}%
\pgfpathquadraticcurveto{\pgfqpoint{11.851343in}{0.353869in}}{\pgfqpoint{11.847848in}{0.358300in}}%
\pgfpathquadraticcurveto{\pgfqpoint{11.844354in}{0.362749in}}{\pgfqpoint{11.837383in}{0.362749in}}%
\pgfpathquadraticcurveto{\pgfqpoint{11.828990in}{0.362749in}}{\pgfqpoint{11.824144in}{0.357400in}}%
\pgfpathquadraticcurveto{\pgfqpoint{11.819299in}{0.352068in}}{\pgfqpoint{11.819299in}{0.342828in}}%
\pgfpathlineto{\pgfqpoint{11.819299in}{0.307200in}}%
\pgfpathlineto{\pgfqpoint{11.808888in}{0.307200in}}%
\pgfpathlineto{\pgfqpoint{11.808888in}{0.370243in}}%
\pgfpathlineto{\pgfqpoint{11.819299in}{0.370243in}}%
\pgfpathlineto{\pgfqpoint{11.819299in}{0.360444in}}%
\pgfpathquadraticcurveto{\pgfqpoint{11.823028in}{0.366136in}}{\pgfqpoint{11.828053in}{0.368946in}}%
\pgfpathquadraticcurveto{\pgfqpoint{11.833096in}{0.371756in}}{\pgfqpoint{11.839689in}{0.371756in}}%
\pgfpathquadraticcurveto{\pgfqpoint{11.850550in}{0.371756in}}{\pgfqpoint{11.856116in}{0.365037in}}%
\pgfpathquadraticcurveto{\pgfqpoint{11.861700in}{0.358318in}}{\pgfqpoint{11.861700in}{0.345260in}}%
\pgfpathlineto{\pgfqpoint{11.861700in}{0.345260in}}%
\pgfpathclose%
\pgfpathmoveto{\pgfqpoint{11.922527in}{0.368387in}}%
\pgfpathlineto{\pgfqpoint{11.922527in}{0.358589in}}%
\pgfpathquadraticcurveto{\pgfqpoint{11.918150in}{0.360840in}}{\pgfqpoint{11.913413in}{0.361957in}}%
\pgfpathquadraticcurveto{\pgfqpoint{11.908693in}{0.363092in}}{\pgfqpoint{11.903614in}{0.363092in}}%
\pgfpathquadraticcurveto{\pgfqpoint{11.895905in}{0.363092in}}{\pgfqpoint{11.892050in}{0.360732in}}%
\pgfpathquadraticcurveto{\pgfqpoint{11.888196in}{0.358373in}}{\pgfqpoint{11.888196in}{0.353635in}}%
\pgfpathquadraticcurveto{\pgfqpoint{11.888196in}{0.350033in}}{\pgfqpoint{11.890951in}{0.347980in}}%
\pgfpathquadraticcurveto{\pgfqpoint{11.893707in}{0.345926in}}{\pgfqpoint{11.902047in}{0.344071in}}%
\pgfpathlineto{\pgfqpoint{11.905595in}{0.343278in}}%
\pgfpathquadraticcurveto{\pgfqpoint{11.916619in}{0.340919in}}{\pgfqpoint{11.921266in}{0.336614in}}%
\pgfpathquadraticcurveto{\pgfqpoint{11.925913in}{0.332309in}}{\pgfqpoint{11.925913in}{0.324600in}}%
\pgfpathquadraticcurveto{\pgfqpoint{11.925913in}{0.315810in}}{\pgfqpoint{11.918960in}{0.310676in}}%
\pgfpathquadraticcurveto{\pgfqpoint{11.912008in}{0.305561in}}{\pgfqpoint{11.899849in}{0.305561in}}%
\pgfpathquadraticcurveto{\pgfqpoint{11.894788in}{0.305561in}}{\pgfqpoint{11.889294in}{0.306552in}}%
\pgfpathquadraticcurveto{\pgfqpoint{11.883801in}{0.307542in}}{\pgfqpoint{11.877730in}{0.309506in}}%
\pgfpathlineto{\pgfqpoint{11.877730in}{0.320205in}}%
\pgfpathquadraticcurveto{\pgfqpoint{11.883476in}{0.317215in}}{\pgfqpoint{11.889042in}{0.315720in}}%
\pgfpathquadraticcurveto{\pgfqpoint{11.894608in}{0.314243in}}{\pgfqpoint{11.900084in}{0.314243in}}%
\pgfpathquadraticcurveto{\pgfqpoint{11.907396in}{0.314243in}}{\pgfqpoint{11.911323in}{0.316746in}}%
\pgfpathquadraticcurveto{\pgfqpoint{11.915268in}{0.319250in}}{\pgfqpoint{11.915268in}{0.323807in}}%
\pgfpathquadraticcurveto{\pgfqpoint{11.915268in}{0.328022in}}{\pgfqpoint{11.912422in}{0.330274in}}%
\pgfpathquadraticcurveto{\pgfqpoint{11.909594in}{0.332525in}}{\pgfqpoint{11.899957in}{0.334614in}}%
\pgfpathlineto{\pgfqpoint{11.896355in}{0.335461in}}%
\pgfpathquadraticcurveto{\pgfqpoint{11.886737in}{0.337478in}}{\pgfqpoint{11.882450in}{0.341675in}}%
\pgfpathquadraticcurveto{\pgfqpoint{11.878181in}{0.345872in}}{\pgfqpoint{11.878181in}{0.353185in}}%
\pgfpathquadraticcurveto{\pgfqpoint{11.878181in}{0.362083in}}{\pgfqpoint{11.884485in}{0.366910in}}%
\pgfpathquadraticcurveto{\pgfqpoint{11.890789in}{0.371756in}}{\pgfqpoint{11.902389in}{0.371756in}}%
\pgfpathquadraticcurveto{\pgfqpoint{11.908117in}{0.371756in}}{\pgfqpoint{11.913178in}{0.370909in}}%
\pgfpathquadraticcurveto{\pgfqpoint{11.918258in}{0.370080in}}{\pgfqpoint{11.922527in}{0.368387in}}%
\pgfpathlineto{\pgfqpoint{11.922527in}{0.368387in}}%
\pgfpathclose%
\pgfpathmoveto{\pgfqpoint{11.943853in}{0.321502in}}%
\pgfpathlineto{\pgfqpoint{11.955741in}{0.321502in}}%
\pgfpathlineto{\pgfqpoint{11.955741in}{0.307200in}}%
\pgfpathlineto{\pgfqpoint{11.943853in}{0.307200in}}%
\pgfpathlineto{\pgfqpoint{11.943853in}{0.321502in}}%
\pgfpathlineto{\pgfqpoint{11.943853in}{0.321502in}}%
\pgfpathclose%
\pgfusepath{fill}%
\end{pgfscope}%
\end{pgfpicture}%
\makeatother%
\endgroup%

\end{lrbox}
\noindent\resizebox{!}{6.85in}{\usebox{\writeupfigure}}
\end{figure}
\end{landscape}
\restoregeometry
\clearpage
\begin{thebibliography}{9}
\bibitem{csa}
Callaway, Brantly, and Pedro H. C. Sant'Anna. 2021.
``Difference-in-Differences with Multiple Time Periods.''
\emph{Journal of Econometrics} 225(2): 200--230.
\href{https://doi.org/10.1016/j.jeconom.2020.12.001}{doi:10.1016/j.jeconom.2020.12.001}.

\bibitem{bjs}
Borusyak, Kirill, Xavier Jaravel, and Jann Spiess. 2024.
``Revisiting Event-Study Designs: Robust and Efficient Estimation.''
\emph{Review of Economic Studies} 91(6): 3253--3285.
\href{https://doi.org/10.1093/restud/rdae007}{doi:10.1093/restud/rdae007}.
\end{thebibliography}

\appendix
\section{Search frame, eligibility, and the implemented sample}
\label{app:sample}

\subsection{Bibliographic frame and screening}
The frozen search frame was assembled on 1 October 2026 from Crossref journal-article records for the \emph{American Economic Review} (AER), \emph{Econometrica}, \emph{Journal of Political Economy} (JPE), \emph{Quarterly Journal of Economics} (QJE), and \emph{Review of Economic Studies} (ReStud). Of 5,121 retrieved records, 5,014 fall within the configured 2016--2025 publication window; 107 date spillovers or unresolved-date records were retained outside that denominator. The dating rule prioritizes the print publication date, followed by published, issued, and online dates. Missing month or day components are assigned January or the first day for filtering. The retrieval contained no duplicate losses under the DOI/paper-identity rule.

Records were reviewed in a frozen hash-based order. All 5,014 positions have a recorded disposition, but screening depth varies: a first-pass exclusion is not a full-text replication audit. Text searches and external catalogs supplied leads, followed by inspection of equations, main-text tables and figures, and replication materials when necessary. The frame has not been independently reconciled against every journal's table of contents; complete coverage of this frozen frame therefore does not establish exhaustive coverage of all eligible published work.

Of the 53 advancing papers, 50 entered through a qualifying TWFE result and three through a modern-estimator-only result. The screening inventory at that checkpoint contains 1,117 candidate scalar specifications or event-study vectors across the 57 advancing or held papers: 842 TWFE scalars, 216 TWFE vectors, 11 modern scalars, and 48 modern vectors. Of these candidates, 1,067 belong to advancing papers and 50 to hold-only papers; 248 candidate rows were held, including 198 within advancing papers. 

\begin{table}[ht]
\centering
\caption{Screening flow and current analysis coverage}
\label{tab:flow}
\begin{tabular}{p{0.76\linewidth}r}
\toprule
Stage & Count\\
\midrule
Retrieved bibliographic records & 5,121\\
Out-of-window or unresolved-date records & 107\\
In-window records with screening dispositions & 5,014\\
\quad First-pass exclusions & 3,419\\
\quad Equation-level exclusions & 1,531\\
\quad Exclusions following full-text audit & 7\\
\quad Papers held for verification & 4\\
\quad Papers advanced to package/data review & 53\\
\midrule
Implemented studies / distinct specifications & 22 / 335\\
Studies / specifications retained in at least one panel & 20 / 275\\
Additional retained pre-treatment averages & 22\\
Retained CSA / BJS estimator pairs & 263 / 288\\
\bottomrule
\end{tabular}
\par\vspace{0.35em}
\begin{minipage}{\linewidth}\footnotesize
Screening dispositions are the validated 3 October checkpoint. Implementation and figure counts refer to the 8 October release. Later stages are not mutually exclusive exclusion categories: a paper can advance while individual specifications remain unresolved.
\end{minipage}
\end{table}

\subsection{Specification eligibility and amendments}
Eligible specifications have one binary treatment at a stable assignment-unit level, unit or group effects, and calendar-time effects. At least two distinct first-adoption dates must be observed among treated units in the usable sample. Panels and repeated cross sections are both eligible.

The adjustment rule permits fixed pretreatment characteristics, including flexible functions of those characteristics, interacted with calendar-time effects, as well as fixed subgroup-by-time effects. Ordinary time-varying controls, changing subgroup membership, unit-specific trends, additional treatments, treatment-heterogeneity interactions, triple differences, instrumental variables, nonlinear outcome models, and stacked event-study specifications are outside the screen. 

The target inventory includes all qualifying main-text outcomes, transformations, samples, weighting choices, and specifications in both tables and figures. Appendix-only results do not provide an entry. A duplicate regression appearing in a table and a figure is one specification with multiple source locations. An inference-only variant is not a new point-estimator specification. A dynamic regression is inventoried as one coefficient vector. A static effect displayed beside that same vector is a linked estimand unless it uses a distinct sample or weighting definition. A small number of published TWFE specifications include treatment reversals; we drop these observations. 




\section{Estimation samples, estimators, and event-study summaries}
\label{app:methods}

\subsection{Baseline TWFE and outcome normalization}
For a static specification, the reference model has the form
\begin{equation}
Y_{it}=\alpha_i+\lambda_t+\beta D_{it}+q_{it}'\gamma+u_{it},
\end{equation}
where $q_{it}$ contains only the permitted baseline-by-time or fixed-subgroup-by-time adjustments. The original outcome transformation, analysis weights, fixed-effect identifiers, and eligible sample are preserved. In repeated cross sections, $i$ denotes the stable group used by the source specification where appropriate. Conflicting source records are not automatically collapsed or recoded; treatment-history identifiers and fixed-effect identifiers can differ when the source design requires it.

The baseline comparison TWFE imposes no CSA/BJS support restriction. It retains always-treated observations and observations lacking a measured untreated outcome. Approved post-adoption zero-treatment cells are omitted as above. Consequently, some baseline TWFE estimates differ from the published estimate. Event averages, modern-only entry, and the fixed-effect exception in Appendix~\ref{app:exceptions} further cause these to differ.

Let $\mathcal O$ be the complete original author-estimation sample for a specification. The outcome scale is
\begin{equation}
s_Y^2=\frac{\sum_{i\in\mathcal O}(Y_i-\bar Y_{\mathcal O})^2}{|\mathcal O|-1}.
\end{equation}
All comparison estimates use $Y_i/s_Y$, without subtracting the mean. Regression analysis weights do not enter this standard deviation. The original sample is used before reversal or estimator-support omissions, and aligned estimators retain this same scale. Previously logged or otherwise transformed outcomes remain in their published definition. 

\subsection{CSA comparisons}
The core implementation constructs cohort-time effects using weighted cross-sectional means. For treatment cohort $g$, current period $t\geq g$, and an allowed baseline stratum $h$, let $b(g)$ be the latest observed calendar period before adoption and let $C_t$ denote never-treated or not-yet-treated observations, with first adoption strictly after $t$. Without continuous adjustment, the comparison is
\begin{equation}
\widehat{ATT}_{hgt}=
  (\bar Y_{hgt}-\bar Y_{hg b(g)})
 -(\bar Y_{hC_t t}-\bar Y_{hC_t b(g)}).
\end{equation}
Each of the four means must have positive analysis-weight support. A supported stratum is retained even if another stratum for that cohort and time is unsupported. Effects are aggregated using the current treated-observation analysis-weight mass of supported cells.

This is a repeated-cross-section implementation even when the source data are a panel: it does not require every treated observation to have its own observed baseline outcome. Cohort-stratum baseline support can suffice. Its interpretation therefore requires the relevant composition and parallel-trends conditions; it is not universally the same statistic as a native balanced-panel CSA estimator. For allowed continuous baseline covariates interacted with time, the conditional branch uses outcome-regression DiD of the repeated-cross-section form, evaluating untreated change at pooled pre/post treated covariate means. 

\subsection{BJS comparisons}
BJS fits the allowed untreated-outcome model using untreated observations only, imputes the untreated outcome at identifiable treated observations, and averages observed-minus-imputed outcomes using treated analysis weights. In the basic fixed-effect case, identification requires the relevant unit and time components to be connected through the untreated-observation graph. With additional permitted adjustments, prediction estimability is checked in the fitted design span.

BJS does not inherit CSA's four-cell requirement, and its targets are not intersected with CSA's. Conversely, CSA's cohort baseline support does not guarantee an identifiable BJS unit-specific untreated prediction. The supported treated populations can therefore differ for reasons beyond the removal of always-treated observations. The row masks, support counts, and retained treated-weight shares are recorded at the specification-estimator level.

\subsection{Exact-union sample alignment}
For each estimator and contrast, the aligned comparison refits the full registered TWFE specification on the exact union of rows actually used by that modern estimator. For CSA, this includes treated-current, treated-baseline, donor-current, and donor-baseline rows for the selected effects and reference components. For BJS, it includes supported treated target rows and the eligible untreated rows used to fit its nuisance model. 

The modern estimator is retained on its own target and construction; TWFE keeps the original specification, permitted adjustment span, analysis weights, outcome scale, and dependence rule when refitted on the union. In particular, alignment does not force a common distribution of implicit treatment-effect weights. 

\subsection{Event studies and reference bins}
The complete joint event-study regression is retained. Its post-treatment summary is a linear contrast $\hat\theta=\sum_{k\in K}v_k\hat\beta_k$, with equal weights across registered published coefficients or bins unless different weights are explicitly recorded. An open-ended bin receives one bin's weight; it is not split into individual years. Reported horizon coefficients do not count as separate specifications.

The comparison window is the structurally estimable intersection of registered horizons across the estimators included in that comparison. Default comparisons use TWFE, CSA, and BJS; a genuinely single-modern comparison uses that estimator and TWFE and is labeled accordingly. Registered weights are renormalized over retained horizons. Availability is determined before bootstrapping from support and contrast estimability, not effect size, significance, estimated standard errors, or draw success. Individual omitted horizons and their reasons are preserved.

For aligned event studies, the baseline common window is the starting point. Exact-union masks are rebuilt and horizons unidentifiable in either corresponding TWFE fit are removed until the window stabilizes. Unpooled nuisance event indicators remain in the joint regression. One adjusted specification in Study~17 averages years 0--6 in the baseline family and 0--2 in the aligned family; its original registered range was 0--8. All other currently retained comparisons keep the same horizon set across families. Thus that point's movement reflects both row alignment and a changed horizon.

Published reference bins are preserved. When the reference includes treated periods, modern effects are contrasted against the corresponding estimated reference quantity rather than relabeled as an untreated-baseline ATT. Its uncertainty and covariance enter every draw. This applies to Study~18's omitted interval $(-10,0]$, which includes adoption year zero.

Pre-treatment coefficients are pooled separately over common, strictly pre-treatment terms, excluding the omitted term and any bin containing treatment or post-treatment observations. In the core event implementation, CSA uses its pre-adoption baseline and donors untreated at both comparison dates; BJS pre coefficients come from a separate untreated-only joint event regression. The mixed-reference construction described above is a distinct contrast. These averages are retained as diagnostics and enter Figures~\ref{fig:baseline}--\ref{fig:aligned} as separate points from the post-treatment averages. Both pre and post averages use square markers, while static specifications use circles.

Comparable pre-treatment averages are available for 22 of the 28 event-study specifications, across five studies. The remaining six specifications have no comparable modern pre-treatment estimate or no separately registered pre-treatment coefficient. Pre-treatment comparisons use their own pooled-contrast leverage measures and paired bootstrap covariance, with the same screening and inference rules as the other points. Missing pre-treatment estimates are not assigned zero. For ordinary untreated-only lead regressions, aligned TWFE and BJS may be algebraically identical. A pooled pre average is not an omnibus pre-trend test, and a mixed-reference pre average is not a pure untreated-baseline placebo.

\section{Joint inference and figure construction}
\label{app:inference}

\subsection{Positive-weight cluster bootstrap}
For each declared dependence block containing $G$ clusters, draw independent $E_c^{(b)}\sim\operatorname{Exp}(1)$ and set
\begin{equation}
W_c^{(b)}=\frac{G E_c^{(b)}}{\sum_{d=1}^G E_d^{(b)}},\qquad b=1,\ldots,999.
\end{equation}
The same cluster multiplier weights every observation and compatible estimator/specification coordinate in the block. Original analysis weights are multiplied by $W_c^{(b)}$; cluster totals are not equalized. Estimator fits, nuisance fits, aggregation masses, and the original-sample outcome scale are recalculated in each draw. Shared underlying innovations and recorded normalizations couple subsamples when their native cluster frames differ. Published dependence definitions, including observation-level robust inference where applicable, are preserved.

For the outcome scale alone, the draw-specific frequency-weight continuation is
\begin{equation}
 s_Y^{2(b)}=
 \frac{\sum_{i\in\mathcal O}W_{c(i)}^{(b)}(Y_i-\bar Y_W^{(b)})^2}
 {\sum_{i\in\mathcal O}W_{c(i)}^{(b)}-1},\qquad
 \bar Y_W^{(b)}=\frac{\sum_{i\in\mathcal O}W_{c(i)}^{(b)}Y_i}
 {\sum_{i\in\mathcal O}W_{c(i)}^{(b)}}.
\end{equation}
Regression analysis weights are absent from this scale calculation. This convention returns the ordinary sample variance at unit multipliers. For event summaries, the fixed contrast, including any reference-bin subtraction, is applied in each draw. Marginal coefficient standard errors are never averaged to obtain an event-average standard error.

Strictly positive multipliers preserve the observed support instead of deleting clusters. No trimming, weight floors, ridge penalties, redraws, or failed-draw deletion are used. A coordinate's standard error requires all 999 finite draws; a covariance requires completeness for both coordinates. Point estimates are the unit-multiplier fits, not bootstrap means. Covariances are sample covariances centered on draw means with divisor 998 and no additional finite-cluster multiplier.

The source release has 23 declared covariance blocks across 22 studies. Cross-specification covariance is available within a shared block. Study~2 uses two distinct published clustering partitions and retains two separate blocks; their cross-block covariance is unavailable, not zero. Between-study covariance is not estimated. Equal numerical draw indices from different blocks must not be treated as paired draws. These distinctions do not prevent the within-pair inference used by the figures.

\subsection{Paired tests and plotting units}
For modern estimator $M$, define $\hat\Delta_M=\hat\theta_M-\hat\theta_T$. Its standard error uses the shared-draw covariance:
\begin{equation}
\widehat{\SE}(\hat\Delta_M)=
\left[\widehat V_M+\widehat V_T-2\widehat C_{MT}\right]^{1/2}.
\end{equation}
A filled marker requires $|\hat\Delta_M/\widehat{\SE}(\hat\Delta_M)|>1.96$ and no inference-warning flag. Plotted points are
\begin{equation}
x=\frac{\hat\theta_T}{\widehat{\SE}(\hat\theta_T)},\qquad
y=\frac{\hat\theta_M}{\widehat{\SE}(\hat\theta_T)}.
\end{equation}
The denominator is the saved bootstrap standard error for that pair's baseline or aligned TWFE coordinate in outcome-SD units; its bootstrap calculation includes variation in the outcome scale. The vertical coordinate is not the modern estimator's own $t$ statistic, and $y-x$ is not the paired-difference $t$ statistic. Dividing both point estimates by a common positive number preserves equality and signs. This final plotting transformation does not change the underlying inference.

Each coordinate is clipped separately at $-5$ or $5$ for display. Capped points use diamonds, with an outer diamond surrounding capped event-study squares; a diamond's apparent distance from the diagonal need not equal its actual gap, and points capped at the same corner can differ substantially. The uncapped estimates, rather than plotted coordinates, determine signs and tests. Colors identify anonymous studies, with publication years retained and four chronological legend columns. Anonymous study IDs are assigned from the 22-study source release and are not renumbered after screening.

\section{Concentration screen and descriptive-count definitions}
\label{app:screen}

The screen measures how an estimator's design weights are concentrated across its published dependence clusters. For each actual scalar or pooled contrast, express the raw-outcome statistic as $\hat\theta=\sum_i a_iY_i$. Contributions from repeated appearances of the same original observation are first netted at that observation before squaring. Define
\begin{equation}
L_c=\frac{\sum_{i\in c}a_i^2}{\sum_i a_i^2},\qquad
G_{\mathrm{eff}}=\frac{1}{\sum_c L_c^2},\qquad
L_{\max}=\max_c L_c.
\end{equation}
The common outcome scale cancels from these shares. 

An estimator pair is excluded if either plotted estimator has $G_{\mathrm{eff}}<5$ or $L_{\max}>0.5$. A retained pair receives an inference warning if either estimator, or their nonidentity difference, has $G_{\mathrm{eff}}<10$ or $L_{\max}>0.25$. Earlier inference warnings are also retained. Warning-marked points are crosses.


\begin{table}[ht]
\centering
\caption{Counts underlying the figure annotations}
\small
\begin{tabular}{lrrrr}
\toprule
& \multicolumn{2}{c}{Baseline} & \multicolumn{2}{c}{Aligned}\\
\cmidrule(lr){2-3}\cmidrule(lr){4-5}
& CSA & BJS & CSA & BJS\\
\midrule
Available before screen & 353 & 357 & 353 & 357\\
Excluded by screen & 90 & 69 & 90 & 69\\
Retained pairs & 263 & 288 & 263 & 288\\
\quad Included pre-treatment averages & 22 & 22 & 22 & 22\\
Warning-marked pairs & 60 & 75 & 58 & 68\\
Opposite nonzero signs & 32 & 20 & 40 & 16\\
Significant paired gap (filled) & 22 & 42 & 21 & 30\\
Significant gap and opposite signs & 3 & 3 & 0 & 3\\
At least one capped coordinate & 32 & 43 & 31 & 41\\
\bottomrule
\end{tabular}
\end{table}

For all three annotated shares, the denominator is the number of retained estimator pairs, including warning-marked pairs. ``Sign-reversals'' means $\hat\theta_T\hat\theta_M<0$. ``Significant differences'' counts filled points, and ``Significant + Reversal'' is the intersection of those two categories. Warning-marked points can contribute to sign reversals but cannot contribute to either filled-point numerator. Their absence from the significant count is not evidence of a precise zero gap. No share is weighted to give each study equal influence, and no pooled independent-specification test is implied.

\section{Implementation exceptions and reproduction records}
\label{app:exceptions}

The following records explain departures from a literal ``published TWFE versus modern estimator on the same sample'' description. They use only the anonymous study IDs shown in the exhibits.

\paragraph{Fixed-effect specification (Study 8).}
The literal source code omits an intercept and one indicator from each fixed-effect set, imposing an additional restriction at the omitted-category intersection. Its coefficients and robust standard errors reproduce the published rounding and are preserved separately. The plotted TWFE comparator uses the full unrestricted additive fixed-effect span under the recorded implementation exception. The bootstrap follows the published observation-level robust dependence rule using 1,048 state-year rows, rather than eight state clusters.

\paragraph{Conflicting records and dependence partitions (Study 2).}
The analysis preserves the original rows and fixed-effect identifiers even where a fixed-effect-unit/time cell contains conflicting treatment or outcome records. Separately audited finer treatment histories determine adoption cohorts, including a joint firm-identifier history where needed. CSA uses these histories, while BJS retains the published fixed-effect structure. The published state and state-by-year clustering partitions remain separate covariance blocks.


\paragraph{Native modern estimates and constructed TWFE (Studies 16, 20, and 21).}
Study~16 contributes both published TWFE level specifications and published CSA log-outcome specifications. The latter receive constructed TWFE comparisons. The comparison CSA uses four-cell repeated-cross-section means with period-specific published analysis weights; its separately preserved native CSA uses a panel formulation. Dependence is clustered by county. Study~20's native CSA variants include never-treated and not-yet-treated donors and their native baseline conventions; the comparison implementation uses its registered not-yet-treated donor convention. Source variants yielding the same normalized comparison are counted once. Study~21's native BJS event studies receive constructed joint TWFE comparisons with the authored indicators preserved.

\paragraph{Event windows and mixed references (Studies 17 and 18).}
Study~17 contains the sole retained event comparison whose common post window changes between figure families, as detailed in Appendix~\ref{app:methods}. Its strictly pre-treatment summary uses years $-3$ and $-2$. Study~18 preserves the published mixed pre/post reference bin in the modern contrasts and their bootstrap draws. Neither implementation replaces missing horizons with zero. 



\end{document}
